\documentclass[10pt,oneside]{book}
\usepackage[margin=1in]{geometry}
\usepackage{amsmath,amssymb,amsthm,mathtools,mathrsfs}
\usepackage{fontspec}
\setmonofont[Path=fonts/,Scale=0.85]{JetBrainsMono.ttf}
\usepackage{newunicodechar}
\usepackage{booktabs,longtable,array,enumitem}
\usepackage{xcolor}
\usepackage[colorlinks=true,linkcolor=blue!50!black,urlcolor=blue!50!black,citecolor=blue!50!black]{hyperref}
\setlist{itemsep=2pt,topsep=4pt}
\sloppy\emergencystretch=5em

\newtheorem{theorem}{Theorem}[chapter]
\newtheorem{proposition}[theorem]{Proposition}
\newtheorem{lemma}[theorem]{Lemma}
\newtheorem{corollary}[theorem]{Corollary}
\theoremstyle{definition}
\newtheorem{definition}[theorem]{Definition}
\newtheorem{example}[theorem]{Example}
\theoremstyle{remark}
\newtheorem{remark}[theorem]{Remark}

\newcommand{\Lean}[1]{\texttt{#1}}
\newcommand{\leanref}[3]{\par{\footnotesize\raggedright\noindent\textsf{Lean:}\ \texttt{#1}\ \ (\texttt{#2}, lines #3)\par}}
\newcommand{\Q}{\mathbf{Q}}\newcommand{\Z}{\mathbf{Z}}\newcommand{\C}{\mathbf{C}}
\newcommand{\R}{\mathbf{R}}\newcommand{\N}{\mathbf{N}}\newcommand{\F}{\mathbf{F}}
\newcommand{\Sha}{\mbox{\fontencoding{U}\fontfamily{wncy}\selectfont Sh}}
\renewcommand{\Sha}{\mathrm{Sha}}
\renewcommand{\Re}{\operatorname{Re}}\renewcommand{\Im}{\operatorname{Im}}
\DeclareMathOperator{\Gal}{Gal}\DeclareMathOperator{\GL}{GL}\DeclareMathOperator{\U}{U}
\DeclareMathOperator{\Frob}{Frob}
\newcommand{\ASK}{\mathrm{ASK}}\newcommand{\Main}{\mathrm{Main}}
\newcommand{\Ga}{\mathrm{Ga}}\newcommand{\Fo}{\mathrm{Fo}}\newcommand{\Ee}{\mathrm{Ee}}
\newcommand{\sg}{\mathrm{sg}}\newcommand{\gainD}{\mathrm{gainD}}

% common Unicode characters that may appear in running text
\newunicodechar{ψ}{\ensuremath{\psi}}\newunicodechar{χ}{\ensuremath{\chi}}\newunicodechar{ε}{\ensuremath{\varepsilon}}
\newunicodechar{ℓ}{\ensuremath{\ell}}\newunicodechar{≤}{\ensuremath{\le}}\newunicodechar{≥}{\ensuremath{\ge}}
\newunicodechar{→}{\ensuremath{\to}}\newunicodechar{∈}{\ensuremath{\in}}\newunicodechar{×}{\ensuremath{\times}}
\newunicodechar{θ}{\ensuremath{\theta}}\newunicodechar{ρ}{\ensuremath{\rho}}\newunicodechar{σ}{\ensuremath{\sigma}}
\newunicodechar{μ}{\ensuremath{\mu}}\newunicodechar{λ}{\ensuremath{\lambda}}\newunicodechar{π}{\ensuremath{\pi}}
\newunicodechar{Φ}{\ensuremath{\Phi}}\newunicodechar{Λ}{\ensuremath{\Lambda}}\newunicodechar{ζ}{\ensuremath{\zeta}}
\newunicodechar{γ}{\ensuremath{\gamma}}\newunicodechar{δ}{\ensuremath{\delta}}\newunicodechar{η}{\ensuremath{\eta}}
\newunicodechar{α}{\ensuremath{\alpha}}\newunicodechar{β}{\ensuremath{\beta}}\newunicodechar{τ}{\ensuremath{\tau}}
\newunicodechar{ω}{\ensuremath{\omega}}\newunicodechar{ξ}{\ensuremath{\xi}}\newunicodechar{κ}{\ensuremath{\kappa}}
\newunicodechar{≠}{\ensuremath{\ne}}\newunicodechar{∞}{\ensuremath{\infty}}\newunicodechar{⊂}{\ensuremath{\subset}}
\newunicodechar{⊆}{\ensuremath{\subseteq}}\newunicodechar{∑}{\ensuremath{\sum}}\newunicodechar{√}{\ensuremath{\surd}}\newunicodechar{ℝ}{\ensuremath{\mathbb R}}

% Printed contents list parts and chapters; PDF bookmarks keep sections.
\setcounter{tocdepth}{0}
\AtBeginDocument{\hypersetup{bookmarksdepth=2}}


\newcommand{\leanrefs}[1]{\par{\footnotesize\raggedright\noindent\textsf{Lean:}\ #1\par}}

\begin{document}

\begin{titlepage}\centering\vspace*{3cm}{\LARGE\bfseries The Riemann hypothesis and GRH for Dirichlet $L$-functions\par}\vspace{1em}{\Large A readable account of the Lean proof in Levent Alpöge's archive\par}\vspace{2em}{\large rh-lean-roadmap\par}\vspace{1em}{November 2026\par}\end{titlepage}

\tableofcontents

\chapter*{Introduction}\addcontentsline{toc}{chapter}{Introduction}\markboth{INTRODUCTION}{INTRODUCTION}\label{ch:intro}

On 20 November 2026 Levent Alpöge announced a proof, found with Claude, of the Riemann hypothesis and of the generalized Riemann hypothesis for Dirichlet $L$-functions. The proof is a Lean~4 archive of 175{,}187 files and 17.4 million lines, without comments, checked by several independent kernels \cite{sutherland-2026}; on 21 November 2026 the Lean FRO and the Mathlib Initiative certified in a joint statement that the statement and the proof of the Riemann hypothesis are correct, the statement being the one in Mathlib \cite{leanfro-2026}. This book gives the argument of the archive as mathematics: every statement the proof needs, in the order in which the proof uses it, with a proof written for a reader and the location of the Lean declaration that proves it.

\section*{The theorem}

The archive proves, for Mathlib's definitions of $\zeta$ and of $L(s,\chi)$:
\begin{itemize}
\item if $\zeta(s)=0$, $s\neq1$ and $s$ is not a negative even integer, then $\Re s=\tfrac12$;
\item for every Dirichlet character $\chi$ of any modulus, $L(s,\chi)\ne0$ when $\Re s>\tfrac12$ (for the principal character, $s=1$ is excluded, where Mathlib assigns a value to the pole);
\item every zero of $L(s,\chi)$ with $0<\Re s<1$ lies on $\Re s=\tfrac12$, for all $\chi$ and for primitive $\chi$.
\end{itemize}
All five statements follow from one theorem: no Dirichlet $L$-function vanishes to the right of the critical line.

\section*{The argument in six steps}

\begin{enumerate}
\item \emph{The critical line} (Part~I). For primitive $\chi$ the functional equation reflects zeros in $0<\Re s<\tfrac12$ to $\tfrac12<\Re s<1$, and imprimitive $L$-functions differ from primitive ones by Euler factors without zeros in $\Re s>0$. It suffices to show $L(s,\chi)\ne0$ for $\Re s>\tfrac12$ and primitive $\chi$.
\item \emph{A window criterion} (Part~I). Fix a primitive character $\psi$ of order $o$ and let $\mathcal X$ be the $o$ primitive characters induced by its powers. A Landau-type argument shows: if a weighted, centred prime sum $G_m(y)$ satisfies $G_m(y)\le y^{\theta}$ on whole intervals $[Y,Y^{3/2}]$ for arbitrarily large $Y$, then no $L(s,\chi)$, $\chi\in\mathcal X$, vanishes for $\Re s>\theta$. Only a one-sided bound is needed, and only on a sequence of windows.
\item \emph{The prime sum over Frobenius classes} (Part~II). For each point $x$ of a window the prime sum is distributed over the Frobenius classes of a finite abelian Galois group. Orthogonality of characters turns the count in each class into a sum over characters; averaging over a cyclic subgroup leaves the target family $\mathcal X$ together with three families of companion characters twisted by the quadratic characters of two auxiliary primes $q,d$. The auxiliary primes are chosen for each window by a mean-square estimate over candidate primes and a zero-density estimate, so that every twisted companion obeys a square-root bound on the whole window.
\item \emph{Multiplicity from theta lifts} (Parts~III and~IV). The proof follows Stepanov's method for curves over finite fields, in which an auxiliary function vanishes to high order at many points. Over the integers the archive builds the multiplicity from a chain of about $\sqrt x$ theta lifts between definite unitary groups. Each lift applies a differential operator, given by the weights, to a local theta kernel; the factorials it produces, minus what the denominators of a formal-group expansion can cancel, give a computable lower bound for the divisibility of the periods by each prime, valid on a whole lattice. The input lattice is then divided by these prime powers, and the pairing stays integral against a global lattice glued across Hecke characters.
\item \emph{Two bounds on one slope} (Part~V). Division raises the average slope of the lattices by a sum of logarithms of primes, which is the prime sum of step~3 up to a small loss. The size of the period maps at the complex places, through a slope inequality for the whole family of maps, bounds the same slope from above. This is Stepanov's inequality ``multiplicity times number of points is at most the degree''. The main terms on the arithmetic and geometric sides are the same number, and the chain length $\sqrt x$ puts the error at square-root size.
\item \emph{Extra divisibility} (Part~VI). Divisibility of special values that does not come from the chain is bounded by an estimate for the total length of Tate--Shafarevich-type groups over a family of CM subfields of cyclotomic fields.
\end{enumerate}

\section*{The three inputs}

In Lean, steps~3 to~6 start from three statements, proved elsewhere in the archive:
\begin{itemize}
\item \textsf{Green}: the classification of the cuspidal characters of $\GL_d(\F_p)$ by irreducible polynomials (Part~VII);
\item \textsf{VWin}: the window bound for Tate--Shafarevich-type groups (Part~VI);
\item \textsf{HypsAlt}: a structure of 46 fields asserting that one family of data satisfies all the conditions of the theta chain at once (Part~VIII).
\end{itemize}
The proofs of \textsf{VWin} and \textsf{HypsAlt} use Galois representations attached to polarized cuspidal automorphic representations and a local Langlands correspondence; Part~IX describes the forms in which the archive proves them.

\section*{How the book is written}

\begin{itemize}
\item The arguments made for this proof are given with proofs. Results from the literature (prime number theorems in progressions, zero-density estimates, Siegel's mass formula, Galois representations, local Langlands, Kato's Euler system, foundations of étale cohomology) are stated in the form the archive uses, with references, and a sentence records how the archive's version compares with the published one.
\item Statements are numbered by node identifiers such as \ref{B3.12}, which group them by part. The Index of Lean declarations at the end of the appendix lists, for each statement by its number, the archive declarations that prove it, with their modules and line numbers. The archive often contains several versions of a result; the one cited is the version on the path of the final proof. The archive version is the one with SHA-256 \Lean{9631b4f0\ldots c7fc}.
\item Lean engineering (coercions, bookkeeping, routine inequalities) is absorbed into the proofs.
\item The main text keeps every statement and the proofs of moderate length. Longer proofs are given in the appendix; in the main text an outline of each says how the proof goes and which results it uses.
\item Chapter~\ref{ch:p1c4} proves the window criterion, Theorem~\ref{B1.31}, with a Gaussian kernel; this proof is shorter than the archive's and is not formalised. The archive's proof is given in Chapter~\ref{ch:p1c5} of the appendix. This is the only place where the main text departs from the route of the formal proof.
\item Chapter~\ref{ch:p3c0} collects the standard facts on the Weil representation and theta lifts that Parts~III, IV and~VIII use.
\end{itemize}

\section*{Reading order}

Part~I and the first chapters of Parts~II and~V give the frame of the argument. Parts~III, IV and~V contain the new mechanism: the local divisibility of one theta lift, its accumulation along the chain and the slope inequality that turns it into a bound for prime sums. Parts~VI to~IX treat the inputs.

\section*{About this text}

The text was written by several AI models, including Claude, working from the Lean source. An independent author, who is not a specialist in this area, set the goals and the editorial rules (plain mathematical English, no hedging, every statement tied to the Lean declaration that proves it) and organised the review: models working in isolation drafted and revised the chapters, and other models read anonymised copies of the whole book as referees, in ten rounds; every point they raised was settled against the Lean source. No human expert has reviewed the text. Where the text and the Lean archive differ, the archive is right; the proofs of a few deep lemmas in Parts V, VI and VIII are still given in outline, and the Lean declarations cited after these statements prove them. The aim is a readable account for mathematicians; shorter versions of the same argument, down to one for readers without mathematics, are listed in the repository. Alpöge works at Anthropic; this text is independent of him and of Anthropic. The mathematics is Alpöge's archive; the exposition is ours.


\part{The five statements and the window criterion}

% ---- p1c1 ----
\chapter{Five statements and the critical line}\label{ch:p1c1}

Let $\mathbf1_N$ be the principal character modulo $N\ge1$, and let
$\chi^\ast$ be the primitive character inducing $\chi$. Write $L(s,\chi)$
for the continued Dirichlet $L$-function. For modulus one it is $\zeta(s)$.
With $G(s)=\pi^{-s/2}\Gamma(s/2)$, define the value at $1$ by
\[
 L(1,\mathbf1_1)=\zeta(1)=
 \lim_{s\to1}\left(\zeta(s)-\frac1{G(s)(s-1)}\right)
 =\frac{\gamma-\log(4\pi)}2<0.
\]
Here $\gamma$ is Euler's constant. The constant term of $\zeta$ at $1$
is $\gamma$, and $G(1)=1$, $G'(1)=-(\gamma+\log(4\pi))/2$. Thus the
subtraction in the display adds $G'(1)$ to the constant term. The sign
follows from $\gamma<1<\log(4\pi)$ \cite[\S\S5.4(ii), 25.2(ii)]{olver-2010}.

\section{The five statements}

\begin{theorem}[Riemann hypothesis]\label{B1.01}
Every zero $s$ of $\zeta$ with $s\ne1$ and
$s\notin\{-2,-4,-6,\ldots\}$ satisfies $\Re s=1/2$.
\end{theorem}
\begin{proof}
Apply Lemma~\ref{B1.06} to Theorem~\ref{B1.09}.
\end{proof}

\begin{theorem}[Explicit zeta statement]\label{B1.02}
For every $s\in\C$,
\[
 \zeta(s)=0,\qquad
 \neg\exists n\in\N\;\;s=-2(n+1),\qquad s\ne1
 \quad\Longrightarrow\quad \Re s=\tfrac12 .
\]
Here $0\in\N$.
\end{theorem}
\begin{proof}
The integers $-2(n+1)$ for $n\in\N$ are the negative even integers
excluded in Theorem~\ref{B1.01}.
\end{proof}

\begin{theorem}[Nonvanishing for every character]\label{B1.03}
Let $\chi$ be a complex Dirichlet character modulo $N\ge1$ and let
$s\in\C$. Then
\[
 \Re s>\tfrac12,\qquad
 (\chi\ne\mathbf1_N\ \text{or}\ s\ne1)
 \quad\Longrightarrow\quad L(s,\chi)\ne0 .
\]
\end{theorem}
\begin{proof}
Apply Theorem~\ref{B1.09}.
\end{proof}

\begin{theorem}[Zeros in the strip for every character]\label{B1.04}
Let $\chi$ be a complex Dirichlet character modulo $N\ge1$.
Every zero $s$ of $L(\cdot,\chi)$ with $0<\Re s<1$ satisfies $\Re s=1/2$.
\end{theorem}
\begin{proof}
Apply Lemma~\ref{B1.07} to Theorem~\ref{B1.09}.
\end{proof}

\begin{theorem}[Zeros in the strip for primitive characters]\label{B1.05}
Let $\chi$ be primitive modulo $N\ge1$. Every zero $s$ of
$L(\cdot,\chi)$ with $0<\Re s<1$ satisfies $\Re s=1/2$.
\end{theorem}
\begin{proof}
Apply Lemma~\ref{B1.08} to Theorem~\ref{B1.09}.
\end{proof}

\section{Reduction to a right half-plane}

Write $H$ for the assertion of Theorem~\ref{B1.03}.

\begin{lemma}[The character of modulus one]\label{B1.06}
The assertion $H$ implies the Riemann hypothesis.
\end{lemma}
\begin{proof}
Apply $H$ to $L(s,\mathbf1_1)=\zeta(s)$ for $\Re s>1/2$ and
$s\ne1$, then use Lemma~\ref{B1.11}.
\end{proof}

\begin{lemma}[Reflection for an arbitrary character]\label{B1.07}
Assume $H$. For every character $\chi$ modulo $N\ge1$, a zero
$s$ of $L(\cdot,\chi)$ with $0<\Re s<1$ satisfies $\Re s=1/2$.
\end{lemma}
\begin{proof}
Lemma~\ref{B1.12} gives $L(s,\chi^\ast)=0$. Apply
Lemma~\ref{B1.08} to $\chi^\ast$.
\end{proof}

\begin{lemma}[Reflection for a primitive character]\label{B1.08}
Assume $H$. For every primitive character $\chi$ modulo $N\ge1$,
a zero $s$ of $L(\cdot,\chi)$ with $0<\Re s<1$ satisfies $\Re s=1/2$.
\end{lemma}
\begin{proof}
In the strip, $H$ excludes $\Re s>1/2$. If $\Re s<1/2$,
Lemma~\ref{B1.13} gives a zero at $1-s$ for $\chi^{-1}$, also
excluded by $H$.
\end{proof}

\begin{theorem}[Common input to the five statements]\label{B1.09}
The assertion $H$ holds.
\end{theorem}
\begin{proof}
Apply Theorem~\ref{B1.15}.
\end{proof}

\begin{definition}[Two nonvanishing assertions]\label{B1.10}
Let $\mathcal P$ consist of pairs $(N,\psi)$ with $N\ge1$ and $\psi$
primitive modulo $N$. Say that $\psi$ is zero-free above $\theta\in\R$
if $L(s,\psi)\ne0$ whenever $\Re s>\theta$. Write $H_{\rm prim}$ for the
assertion that every $(N,\psi)\in\mathcal P$ is zero-free above $1/2$,
including the value at $1$.
\end{definition}

\begin{lemma}[Zeta reflection from the right]\label{B1.11}
Suppose that $\zeta(s)\ne0$ for $\Re s>1/2$ and $s\ne1$.
Then the Riemann hypothesis holds.
\end{lemma}
\begin{proof}
Suppose $\zeta(s)=0$ and $\Re s<1/2$. Put $t=1-s$.
Since $\zeta(0)=-1/2$, one has $t\ne1$ \cite[\S25.6]{olver-2010}.
The reflection formula at $\Re t>1/2$ gives
\[
 \zeta(1-t)=2(2\pi)^{-t}\Gamma(t)\cos(\pi t/2)\zeta(t)
\]
\cite[\S12.8, (12)]{apostol-1976}\cite[(25.4.1), \S5.2(i)]{olver-2010}.
By the hypothesis and gamma nonvanishing, $\cos(\pi t/2)=0$.
Thus $t$ is an odd integer greater than $1$, so $s$ is negative and even.
\end{proof}

\begin{lemma}[Euler factors and the conductor]\label{B1.12}
Let $\chi$ be a character modulo $N\ge1$ with conductor $f$. If
$\chi^\ast\ne\mathbf1_f$ or $s\ne1$, then
\[
 L(s,\chi)=L(s,\chi^\ast)\prod_{p\mid N}(1-\chi^\ast(p)p^{-s}).
\]
For any character $\varphi$ and prime $p$, one has
$|\varphi(p)p^{-s}|\le p^{-\Re s}<1$ if $\Re s>0$.
In particular $L(s,\chi)=0$ in $0<\Re s<1$ implies $L(s,\chi^\ast)=0$.
\end{lemma}
\begin{proof}
Compare Euler products on $\Re s>1$ and continue on the stated domain
\cite[\S12.10]{apostol-1976}\cite[(25.15.4)]{olver-2010}.
The factors at $p\mid f$ equal $1$. Since $\varphi(p)$ is zero
or a root of unity, the bound follows from $|p^{-s}|=p^{-\Re s}$.
\end{proof}

\begin{lemma}[Reflection of primitive zeros]\label{B1.13}
Let $\psi$ be primitive modulo $N\ge1$. Put $a_\psi=0$ for even
$\psi$ and $a_\psi=1$ for odd $\psi$. Write
$\Gamma_{\R}(z)=\pi^{-z/2}\Gamma(z/2)$, with its meromorphic continuation.
In the strip $0<\Re s<1$ put
$\mathcal L_\psi(s)=\Gamma_{\R}(s+a_\psi)L(s,\psi)$.
Then $L(s,\psi)=0$ implies $L(1-s,\psi^{-1})=0$.
\end{lemma}
\begin{proof}
In this normalization the functional equation is
\[
 \mathcal L_{\psi^{-1}}(1-s)=N^{s-1/2}W(\psi^{-1})\mathcal L_\psi(s),
 \qquad
 W(\psi)=i^{-a_\psi}N^{-1/2}\sum_{u\bmod N}\psi(u)e^{2\pi iu/N}.
\]
For $N>1$ use \cite[\S12.10, (16)]{apostol-1976}. For $N=1$ use
the completed zeta equation, with $W(\mathbf1_1)=1$
\cite[(25.4.3)--(25.4.4)]{olver-2010}. The two gamma arguments have
positive real part, so their values are finite and nonzero
\cite[\S5.2(i)]{olver-2010}.
\end{proof}

\begin{lemma}[Passage to all characters]\label{B1.14}
The assertion $H_{\rm prim}$ implies $H$.
\end{lemma}
\begin{proof}
Fix $N,\chi,s$ as in $H$. If $\chi^\ast$ is principal, then $\chi$
is principal and $s\ne1$. Lemma~\ref{B1.12} expresses $L(s,\chi)$
as the product of the nonzero primitive value and nonzero Euler factors.
\end{proof}

\begin{definition}[The five toric assertions]\label{B1.14k}
Let $E$ be CM, $F=E^+$ and $(\phi_i)_{i\in I}$ a CM type. Write
$D_F,D_E$ for the field discriminants, $f=|I|=[F:\Q]$, $k=2s$
for an integer $s\ge1$, and $n=kf$. For integers $\kappa_i\ge0$
put $K_\kappa=\sum_i\kappa_i$ and $M=n+2K_\kappa$. Write $\operatorname{Nm}$ for the absolute ideal norm. For a number field $L$, write $\operatorname{Art}_L$ for idèlic reciprocity normalized to send a uniformizer to geometric Frobenius modulo inertia. Arithmetic Frobenius acts on the residue field by $z\mapsto z^{Nv}$ at a prime $v$, and geometric Frobenius is its inverse. For a finite idèle $t$ abbreviate $\operatorname{Art}_L((1,t))$ to $\operatorname{Art}_L(t)$.

A base conductor $\mathfrak f$ is a proper nonzero integral ideal of
$F$. Each prime divisor splits in $E/F$, has exponent at most two
and lies over an odd rational prime unramified in $E$ with residue
degree one at every place of $E$ above it. For a finite set $Q$ of
split primes of $F$ coprime to $\mathfrak f$, put
$\mathfrak f_Q=\mathfrak f\prod_{\mathfrak p\in Q}\mathfrak p$ and
$O_Q=\mathcal O_F+\mathfrak f_Q\mathcal O_E$. The characters $\chi$
are unitary, trivial on the id\`eles of $F$, with infinity component
$\prod_i(\phi_i/\bar\phi_i)^{s+\kappa_i}$. At base level their
conductor exponents above $\mathfrak p$ are those of $\mathfrak f$.
At auxiliary level they may also have exponent at most one above $Q$.

Use the toric class sums of \ref{B8.21a}. Let $\mathbb A_{E,f}$ be the
finite adeles and $\widehat O_Q$ the finite completion of $O_Q$. Put
\[
 \operatorname{Pic}(O_Q)=E^\times\backslash
 \mathbb A_{E,f}^\times/\widehat O_Q^\times,\qquad
 Y_y=E\cap y\widehat O_Q .
\]
For a finite id\`ele $y$, let $|y|_{E,f}$ be the product of its
normalized local absolute values. Define $S_y$ by continuation to
$\sigma=0$ of
\[
 \sum_{t\in(Y_y\setminus0)/\mathcal O_F^\times}
 \prod_i\bar\phi_i(t)^{\kappa_i}\phi_i(t)^{-k-\kappa_i}
 |\phi_i(t)|^{-2\sigma}.
\]
The ordinary class term is $S_y/(\chi_f(y)|y|_{E,f}^s)$, where
$\chi_f$ is the finite part of $\chi$. The ordinary sum adds these
terms, and the twisted sum multiplies the term at $[y]$ by
$\psi([y])^{-1}$ for a character $\psi$ of $\operatorname{Pic}(O_Q)$.
For the positive integer $d$ in the quadratic datum put
\[
 C_0=|D_F|^{(k-1)/2}\prod_i\Gamma(k+\kappa_i)
 (2\pi/\sqrt d)^{K_\kappa},\qquad
 x_\Omega(T)=\frac{C_0T}{\Omega^M},\qquad
 \beta_\Lambda(T)=x_\Omega(T)^{12}\quad
 (\Omega^{12}=\operatorname{disc}\Lambda).
\]
Here $\Omega\ne0$ and $\operatorname{disc}\Lambda$ is the Weierstrass
discriminant of the lattice $\Lambda$. For an ordinary ideal class
$C$, let $\zeta_F(z,C)$ be the continuation of the sum of
$(\operatorname{Nm}\mathfrak a)^{-z}$ over integral ideals in $C$.
When all $z(C)=\zeta_F(1-k,C)$ are rational, put
$e_\ell=\max_C\max(0,-v_\ell(z(C)/2^f))$, with zero contribution
from a zero value and $v_\ell(\ell)=1$.

The assertion $\mathscr J$ is the conjunction of the following statements.
\begin{enumerate}
\item Suppose $E/\Q$ is abelian, $E/F$ is unramified at finite places,
$2\mid D_E$, $f\ge2$, and $d>4$ is prime. Suppose $w^2=-d$ in $E$,
$\phi_i(w)=i_{\C}\sqrt d$ for all $i$, and $\prod_i\phi_i(\xi)=1$ for
each root of unity $\xi$ of $E$. Here $i_{\C}$ has positive imaginary
part. For each $s\ge1$ choose rational $z(C)$ for all classes at once.
For each $\kappa$, base conductor and base character use
$\Lambda=\mathcal O_{\Q(\sqrt{-d})}\subset\C$. At a rational prime
$\ell$, let $a_\ell$ be the maximum base conductor exponent above $\ell$.
Put $c_\ell=a_\ell$ if $a_\ell>0$ and
$c_\ell=v_\ell(|D_E|)/(2f)$ otherwise. For each auxiliary set away from
$\ell$, every class term and every Picard-twisted sum $T$ has algebraic
$\beta_\Lambda(T)$, and at every embedding above $\ell$,
\begin{equation}\label{eq:p1c1-toric-denominator}
 \begin{split}
 v_\ell(\beta_\Lambda(T))&\ge-12b_\ell,\\
 b_\ell&=e_\ell+Ma_\ell+
 \begin{cases}
 0,&c_\ell=0,\\
 nc_\ell+K_\kappa(2c_\ell+1+\mathbf1_{\ell=2}
 +\tfrac12\mathbf1_{\ell=d}),&c_\ell>0 .
 \end{cases}
 \end{split}
\end{equation}
Equivalently $\ell^r\beta_\Lambda(T)^j$ is integral above $\ell$
whenever $j\ge1$, $r\in\Z$ and $r/j\ge12b_\ell$.

\item Let $K$ be imaginary quadratic of discriminant $-d<-4$, with a
fixed embedding and integer lattice $\Lambda_K$. For each automorphism
$g$ of $\C$ fixing $K$, choose a finite idèle $t_g$ with $\operatorname{Art}_K(t_g)^{-1}=g|_{K^{\rm ab}}$, its inverse ideal
$A_g$, and the lattice $\Lambda_g$ for that ideal, so that
$j(\Lambda_g)=g j(\Lambda_K)$. The classes of $A_g$ index the conjugate
$j$-values bijectively. In every finite Galois $L/\Q$ containing $K$
and $j(\Lambda_K)$ each class occurs $[L:\Q]/(2h_K)$ times among the
embeddings over $K$, where $h_K$ is the class number.
The choices work for every CM $E\supset K$ Galois over $\Q$ with
trivial relative different over $F$, each type of size at least two
extending the fixed embedding, all $s,\kappa_i\ge1$, and all admissible
conductors and ordinary sums at auxiliary level. Write $T'$ for the
sum whose weights are permuted by $g$ on the type and whose character
$\chi'$ has these infinity weights and the $g$-conjugate finite part.
With $A=\operatorname{Nm}A_g$ and
$l_\chi=\chi'_f(t_g)|t_g|_{E,f}^s$, using the image of $t_g$ in the
id\`eles of $E$, one has algebraicity of $\beta_{\Lambda_K}(T)$ and
\begin{equation}\label{eq:p1c1-ordinary-transport}
 \begin{split}
 g\beta_{\Lambda_K}(T)
 &=\frac{(C_0 A^{-K_\kappa}l_\chi^{-1}T')^{12}}
 {(\operatorname{disc}\Lambda_g)^M},\\
 |g\beta_{\Lambda_K}(T)|
 &=\frac{(C_0|T'|)^{12}}
 {(A^6|\operatorname{disc}\Lambda_g|)^M}.
 \end{split}
\end{equation}
For two such sums of the same weight with $T_2\ne0$, the ratio is
algebraic, $T'_2\ne0$, and
$g(T_1/T_2)=(l_{\chi_1}^{-1}T'_1)/(l_{\chi_2}^{-1}T'_2)$.

\item Let $K$ be quadratic with $\delta^2=-d$, $d>0$, and a fixed
complex embedding. There is $\Omega_0\ne0$ chosen before the CM
extension, type or twist. For every such $E/K$ unramified over $F$
at finite places, type extending the fixed embedding, $s,\kappa_i\ge1$,
and Picard twist $\psi$ of a base character at $\mathfrak f_Q$,
the twisted sum $T$ satisfies $x_{\Omega_0}(T)\in\overline{\Q}$.
Fix a compatible embedding of an algebraic closure of $K$ and an automorphism $g$ of $\C$ fixing the embedded $K$. Take any finite idèle $t$ satisfying $\operatorname{Art}_K(t)^{-1}=g|_{K^{\rm ab}}$. There is $\Omega_1\ne0$ common to all these sums. With
$\psi'=g\psi$, $T'$ the conjugate sum, $l_\chi=\chi'_f(t)|t|_{E,f}^s$,
and $[t]$ the class of the image of $t$ in $\operatorname{Pic}(O_Q)$,
\begin{equation}\label{eq:p1c1-twisted-transport}
 g x_{\Omega_0}(T)=
 x_{\Omega_1}\bigl(|t|_{K,f}^{-K_\kappa}
 [l_\chi\psi'([t])]^{-1}T'\bigr).
\end{equation}

\item For every CM $E$ and $s\ge1$ the rational $z(C)$ exist, and
\begin{equation}\label{eq:p1c1-class-denominator}
 \begin{split}
 e_\ell&\le f v_\ell(2)+
 \min_{A\ge2,\,\ell\nmid A}v_\ell(A^n-1),\\
 \min_{A\ge2,\,\ell\nmid A}v_\ell(A^n-1)
 &=\begin{cases}
 2+v_2(n),&\ell=2,\\
 1+v_\ell(n),&\ell\text{ odd},\ \ell-1\mid n,\\
 0,&\ell\text{ odd},\ \ell-1\nmid n .
 \end{cases}
 \end{split}
\end{equation}
For $n\ge16$ the support is finite and
$\sum_\ell e_\ell\log\ell\le n\log(n+1)$.

\item Use the recursive size data of \ref{B8.19}, the admissible level-character systems of \ref{B8.20}, and the character grids of \ref{B8.21}.
There is $C_{F,0}$ such that for each $\varepsilon$ and $C_F\ge C_{F,0}$
there are $C'_0,x_0$ valid for all $C\ge C'_0$ and every recursive size datum of \ref{B8.19}
at $x\ge x_0$. Take an admissible level-character system of \ref{B8.20} with reference index $\mu_0$, a grid of \ref{B8.21} of parameters $(1,C_F,b)$, its Weil system and a full choice of tests. In an admissible level-character system and grid of \ref{B8.20}--\ref{B8.21}, at each level index $t$ and grid position $(i,j)$, the grid character $\chi_{t,i,j}=\chi_{0,i,j}\xi_t$ differs from the reference character $\chi_{0,i,j}$ at $\mu_0$ by a Hecke character twist $\xi_t$ of $E$ (with $\xi_{\mu_0}=1$) satisfying four conditions: (a) $\xi_t$ is trivial on $\mathbb A_F^\times$, has finite order (hence trivial infinity component), and is tame at every finite place of $E$; (b) at each finite place $v$ of $E$ where $\xi_t$ ramifies and $\chi_{0,i,j}$ is unramified, the residue field has prime order $q_v\mid Q_n(t)$ with $v\in W_t$ of depth $d_v\ge1$, and the set of such places is stable under conjugation of $E/F$; (c) at each finite place $v$ of $E$ where both $\xi_t$ and $\chi_{0,i,j}$ ramify, the residue field has prime order $q_v$, $v$ is unramified over $q_v$, and $\chi_{0,i,j}$ has conductor exponent $2$; and (d) for $h$ equal to half the sum of the grid ranks at $(i,j)$ and $(i+1,j+1)$, the fitted gap $\kappa_\tau(h)>0$ of \ref{B8.21} gives common infinity exponent $m_\tau=\kappa_\tau(h)+h>h\ge e_0$ at every CM-type embedding $\tau$, where $e_0$ is the lower range parameter below. Let $K=\iota_D(K_D)\subset\C$ be the image of the coefficient field of the point datum $D$. Impose the following coefficient-field condition. The field $K$ contains all finite local component values of the reference row characters, of every level twist and of every layer character. At places where the respective character is unramified, it contains the uniformizer value of the auxiliary character $\eta$ divided by $\sqrt{Nw}$ and the uniformizer values of the three auxiliary characters $\chi_\Psi,\delta,\delta'$ of the level system. Here $Nw$ is the residue cardinality. The field also contains all eighth roots of unity. For each level index $t$, let $W_t=W_0\cup Q(t)$, where $W_0$ is the base set of finite places and $Q(t)$ is the set specified by the level at $t$. For $w\in W_t$, put $d_w=1$ off $W_0$ and let $d_w$ on $W_0$ be the larger of one and the grid depth at the place below $w$. Require $K$ to contain all $(Nw)^{d_w}$th roots of unity for $w\in W_t$, and $\sqrt{Nw}$ for every $w$ in any $W_t$ or in the set of places of the test choice. On each supported term, the additive phases with nonzero test and translated factors belong to $K$. In every coset expansion of least depth, the top and row jet values of every supported summand belong to $K$ at every prescribed class and for all jet indices.
At each pair of grid positions, require the ratio of the Hecke value at $t$ to its reference value at $\mu_0$ to lie in $K$ whenever the latter is nonzero, at every integer from $1$ to $1+k_0+m(u+r+7)$. Here $k_0$ is the starting rank, $m$ the chain length, $u$ the rank increment parameter and $r=k_0+m-1$. At every pair admitting a bottom triple, require a choice of such a triple and two nonzero periods, one for each side. For every $t$, finite adelic target point and monomial test for which the corresponding integral converges absolutely, its value divided by the period of its side lies in $K$.
At each pair of grid positions let $e_0$ be the lower range parameter
and let $L_t,L_0$ be the Hecke functions of the character at $t$
and the reference character. Put $Q_n(t)=\prod_{w\in W_t\setminus W_0}Nw$.
For every integer $1-e_0\le s\le0$ there is $u\in K_D$, a unit
outside the primes dividing $Q_n(t)$, such that
\begin{equation}\label{eq:p1c1-grid-reflection}
 L_t(s)L_0(1-s)=\iota_D(u)L_t(1-s)L_0(s).
\end{equation}
\end{enumerate}
\end{definition}

\section{The common arithmetic input}

Let $\mathcal G$ be Green's assertion of Definition~\ref{B7.06}.
For each prime $p$, rank $r\ge1$ and compatible label datum, it
parametrizes irreducible cuspidal characters of $\GL_r(\F_p)$ by
monic irreducible polynomials of degree $r$ with nonzero constant
term and gives the degrees and elliptic values.

Let $\mathcal V$ be the window assertion of Definition~\ref{B6.03}
with exponent $401$ for every primitive $\psi$. Given
$0<\varepsilon<1/52$ and a page function $a$ as in \ref{B6.02},
there are $J\ge1$ and $X_0$ such that for each $X\ge X_0$ and each
admissible pair $q,d$ at scale $X^{1/J}$ one can choose
$Y\ge X^{1/J}$ with $Y^2\le X$. At each $x\in[Y,Y^2]$, for each
field $E$ specified by $(\psi,a(x),q,d)$ and all choices of
\ref{B6.01}, the lengths $s_{\ell,e}(E)$ are finite for every
prime $\ell$ and $1\le e<t=\lfloor x^{401}\rfloor$, and
\[
 \sum_{\ell\in S}\log\ell\sum_{1\le e<t}s_{\ell,e}(E)
 \le \exp(x^\varepsilon/2)
 \qquad(S\text{ a finite set of primes}).
\]

Let $\mathcal D(D,O,c,\beta,C,C_H)$ be the definite-side condition
of \ref{B8.35}. Here $D,O$ are the point datum and reference
realization, $c,\beta$ are the selected cut and bundle, and $C,C_H$
are the bounds. For each linked pair of sides it quantifies over
compatible numerical parameters and coordinate submodules with
division identifications agreeing with $\beta$. For every word
it requires a factored column, metric frames and a height table
satisfying the terminal and layer inequalities on the inverse
image of the input lattice under the division identification.

A family satisfies the local-error estimate if there are $C,x_0\in\R$ such that, at every point with $x\ge x_0$ where the prescribed cohomology lengths are finite for every prime and every $1\le e<t$, inequality~\eqref{eq:B8-linkN} holds in every layer whenever its weighted length series is summable. Write $\mathcal N(\mathcal C)$ if, for every sufficiently small $\varepsilon>0$, one can choose the size thresholds of Definition~\ref{B8.39} so that every family in $\mathcal C$ of parameter $\varepsilon$ meeting those thresholds at $C_H$ satisfies the local-error estimate. The thresholds precede the choice of the family and $C_H$. Let $\mathcal H$ consist of the common data of Definition~\ref{B8.02}, with $\mathcal D$ as the definite-side condition and $\mathcal N$ as the condition on family classes. An element specifies a class $\mathcal A$ of families and selects a compatible bundle at each eligible point of each member. All the conditions of that definition use this selection. Let $W_3(\mathcal A)$ be the subclass with the twisting, size and design conditions at parameter $3$ specified in Definition~\ref{B8.02}. For each eligible family $F_m$ in this subclass and each eligible $C_H$, one constant $C\ge0$ and one point threshold give both the indefinite bounds and $\mathcal D$ for the selected bundles.
Write $x_1(F_m)$ for the lower endpoint of the family data and $I_{F_m}(x)$ for the set indexing its points at $x$. Beyond the threshold for the period estimates, they hold at every $x>x_1(F_m)$ for all $i\in I_{F_m}(x)$. The local-error condition is $\mathcal N(W_3(\mathcal A))$.

\subsection{Evaluation at CM lattices}

Fix a rational prime $\ell$ and an embedding
$\lambda:\overline{\Q}\cap\C\longrightarrow\C_\ell$. Write
$v=v_\ell\circ\lambda$, with $v(0)=+\infty$, and put
$R=\{z\in\overline{\Q}\cap\C:v(z)\ge0\}$. A valuation bound
includes algebraicity. For $\tau$ in the upper half-plane set
$p=2\pi i_{\C}$, $q=\exp(p\tau)$, and
$\Delta=q\prod_{j\ge1}(1-q^j)^{24}$. Use $j=q^{-1}+744+\cdots$.

\begin{lemma}[Evaluation of a finite family of modular forms]\label{B1.14a}
Let $N\ge1$ and $w\ge0$ be integers. Let $(H_i)_{i\in I}$ be a finite
family of holomorphic functions on the upper half-plane permuted by
the weight-$w$ slash action of $\mathrm{SL}_2(\Z)$. Suppose they have
convergent expansions in nonnegative powers of $q^{1/N}$ and all
coefficients have valuation at least $-r$ for a common $r\in\R$.
If $v(j(\tau))\ge0$ and $\omega^{12}=\Delta(\tau)$, then
$v(H_i(\tau)/\omega^w)\ge-r$ for every $i$.
\end{lemma}
\begin{proof}[Outline]
Scale at valuation \(t>r\) and take elementary symmetric functions of twelfth powers. Divide by powers of \(\Delta\) and cancel cusp principal parts with \(j\) to get a monic equation over the valuation ring, then let \(t\) decrease to \(r\). The complete proof is in \S\ref{prf:B1.14a}.
\end{proof}

For a lattice $\Lambda$ write $D(\Lambda)=g_2(\Lambda)^3-27g_3(\Lambda)^2$
and let $\wp_\Lambda$ be its Weierstrass function. The quasi-period
homomorphism of the Weierstrass zeta function extends to
$u\mapsto s_2(\Lambda)u+b\bar u$ on $\C$. For $\Lambda_\tau=\Z\tau+\Z$,
\[
 D(\Lambda_\tau)=p^{12}\Delta(\tau),\qquad
 s_2(\Lambda_\tau)=G_2(\tau)-\pi/\Im\tau ,
\]
where $G_2$ has constant term $\pi^2/3$ \cite[Ch.~2]{shimura-1971}.

\begin{lemma}[The elliptic correction at a CM point]\label{B1.14b}
Let $d>4$ be prime and suppose $w_0=i_{\C}\sqrt d$ carries $\Lambda$
into itself. Let $\Omega^{12}=D(\Lambda)$ and take a nonzero
$P\in\C/\Lambda$ killed by a positive integer prime to $\ell$.
Then $U=(s_2(\Lambda)+\wp_\Lambda(P))/\Omega^2$ satisfies
\[
 v(U)\ge-\begin{cases}
 1,&\ell=2,\ d\equiv3\pmod4,\\
 2,&\ell=2,\ d\equiv1\pmod4,\\
 1/2,&\ell=d,\\
 0,&\text{otherwise}.
 \end{cases}
\]
\end{lemma}
\begin{proof}[Outline]
Use torsion expansions and \ref{B1.14a} for integrality of \(12X(P)\) and \(X(P)-X(P')\). Cancel poles in the isogeny sum and integrate over a period to relate \(s_2\) to kernel values, then use \(w_0\) and, where needed, \(1\pm w_0\) to bound \(v(U)\). The complete proof is in \S\ref{prf:B1.14b}.
\end{proof}

\begin{lemma}[CM coordinates for all orders of differentiation]\label{B1.14c}
Use the field, type, square root $w$, integer $s$, partial zeta values
and base conductor in the first assertion of Definition~\ref{B1.14k}.
Take $Q$ away from $\ell$ and a finite id\`ele $y$ equal to one above
$\ell$. Put $Y=Y_y$, $a=a_\ell$, $c=c_\ell$, $e=e_\ell$ and
$\Lambda=\mathcal O_{\Q(\sqrt{-d})}\subset\C$. For $|z_i|<1$ put
$u_i(t,z)=\phi_i(t)+z_i\overline{\phi_i(t)}$ and let $g(z)$ be
$c_k=((k-1)!)^f|D_F|^{(k-1)/2}$ times the sum of
$\prod_i u_i(t,z)^{-k}$ over $(Y\setminus0)/\mathcal O_F^\times$.
At $k=2$ insert $\prod_i|u_i(t,z)|^{-2\sigma}$, sum for
$\Re\sigma>0$ and continue to $\sigma=0$.

For each $\Omega^{12}=D(\Lambda)$ there are nonzero $a_i,b_i\in\C$,
chosen before the derivative orders, such that for every
$\kappa\in\N^I$ and any ordering of the derivatives,
\begin{equation}\label{eq:p1c1-coordinate-estimates}
 \begin{split}
 v(a_i/\Omega)&\ge-a,\qquad
 v\left(\frac{2\pi}{a_i b_i\sqrt d}\right)
 \ge-c-\tfrac12\mathbf1_{\ell=d},\\
 v\left(\partial^\kappa g(0)\prod_i
 \frac{(b_i/a_i)^{\kappa_i}}{a_i^k}\right)
 &\ge-e-(n+K_\kappa)c .
 \end{split}
\end{equation}
\end{lemma}
\begin{proof}[Outline]
Clear lattice denominators and correct by torsion to form modular families, then use Poisson summation, the partial zeta functional equation and normal convergence. Apply \ref{B1.14a} at the CM point and \ref{B1.14b} to the change of direction, then rescale for the bounds. The complete proof is in \S\ref{prf:B1.14c}.
\end{proof}

\begin{lemma}[Toric bounds, transport and reflection]\label{B1.15a}
The elliptic main theorem of complex multiplication, the Hilbert
class field theorem for $j$, Siegel--Klingen rationality with
Deligne--Ribet integrality, the partial zeta functional equation
and the Hecke--Tate functional equation give $\mathscr J$ of
Definition~\ref{B1.14k}.
\end{lemma}
\begin{proof}[Outline]
Uses Siegel--Klingen, Deligne--Ribet, and Lemma~\ref{B1.14c} for toric denominator bounds, and elliptic complex multiplication, Lemma~\ref{B1.14b}, and the Hilbert class field for transport. Local epsilon-factor formulas and Hecke--Tate functional equations give reflection. The complete proof is in \S\ref{prf:B1.15a}.
\end{proof}

\begin{definition}[The unitary and local theta inputs]\label{B1.15ac}
Let $P$ and $F$ be the following products of propositions. Each
entry quantifies over the indicated fields, ranks, levels, measures
and representations. A witness chosen in an entry has the scope
of that entry.

The following entries occur in both products.
\begin{enumerate}
\item Over a nonarchimedean local field, smooth irreducibles of
$\GL$ are admissible, Jacquet modules and parabolic induction
preserve finite length, and normalized induction commutes with
contragredients. The normalized product of two smooth irreducible
unitarizable representations is irreducible. With normalized
induction as multiplication, the direct sum of Grothendieck groups
is polynomial over $\Z$ on the classes of segments.

\item For a unitary group the automorphic forms at fixed open
level, finite-dimensional continuous compact type and
infinitesimal character form a finite-dimensional space.
Smooth uniformly moderate cuspidal forms have rapid decay.
Convolution by a smooth compactly supported test is trace
class on the cuspidal space for bi-invariant Haar measure.
On a closed real matrix group such a test is a finite sum
of convolutions $g_i*h_i$, with $g_i$ supported in a prescribed
identity neighborhood and $h_i$ in the support of the test.
An irreducible unitary adelic representation of a nondegenerate
CM Hermitian group has a restricted tensor factorization with
unramified unit vectors. For the special unitary group of a
nondegenerate CM Hermitian space of dimension at least two,
the rational points times the subgroup at the specified
places are dense when that subgroup is noncompact.

\item For $\mathrm U(p,q)$ the Harish--Chandra parameters cover the
discrete series with their compact types. Each permissible parameter
gives a discrete series in the abstract group and in the matrix
realization of its signature. For a parameter $\lambda$ of rank $a$,
the compact-dual normalization gives formal degree
$\prod_{1\le i<j\le a}|\lambda_i-\lambda_j|/(j-i)$.
The Harish--Chandra module assertions include analyticity of
compact-finite vectors. A $(\mathfrak g,K)$ map from the compact-finite
vectors of an irreducible unitary representation to a smooth
moderate globalization extends uniquely to an equivariant map on
smooth vectors, with closed image there. A nonzero map controls
each source Sobolev seminorm by a continuous target seminorm as
in \ref{B8.38a}. An irreducible strongly continuous unitary
representation of a locally compact group with regular two-sided
Haar measure and one nonzero square-integrable coefficient has
square-integrable coefficients and satisfies Schur orthogonality
with one positive formal degree.

\item The prescribed restrictions of the Hecke splitting
characters, with a standard additive character, give the
Weil action for a CM unitary pair. For ranks $a,b$ with $a\ge1$
and with $B$ anisotropic or $a+\operatorname{Witt}(B)<b$,
the absolute theta integral is bounded by a continuous Schwartz
seminorm. The theta kernel is integrable, and for $2a<b$ its
mean equals the absolutely convergent Siegel sum with rational
lattice of covolume one. The Witt index is the dimension of
a maximal totally isotropic subspace. The doubling input gives
meromorphic continuation and finite pole order for flat sections,
the basic integral identity for $\Re s>a/2$, and continuation
of the cuspidal zeta integral.
For a characteristic-zero nonarchimedean field with ring of
integers $\mathcal O$ and residue cardinality $q$, normalize
$K=\GL_a(\mathcal O)$ to volume one. For an unramified
representation with Satake parameters $\alpha$ and normalized
zonal coefficient, there is $s_0$ such that for every
$\Re s>s_0$ the spherical doubling integral converges absolutely
and equals
\[
 L(\alpha,s+1/2)L(\alpha^{-1},s+1/2)
 \prod_{j=1}^{a}(1-q^{-2s-j}).
\]
Here $L(\alpha,z)=\prod_i(1-\alpha_i q^{-z})^{-1}$.
\end{enumerate}

The additional entries of $P$ are as follows.
\begin{enumerate}
\item Regular algebraic cuspidal representations of positive rank
over a totally complex field are cohomological. Over a CM field
the two conjugate infinitesimal weights are reflections about one
half-integral centre. The polarized Galois assertion is that of
\ref{B9.02}, with the weights of \ref{B9.04}. For above-$\ell$
comparison in conjugate self-dual rank $a\ge2$ the polarization is
$\chi_{E/F}^a$, where $\chi_{E/F}$ is the quadratic character of
the CM extension. Choose the representation, weights and jump
multisets before the prolongation and correspondence of
\ref{B9.47}. For each correspondence $R$ of Definition~\ref{B9.47}, apply the following inertia test to the parameter $y=(\rho,N)=R(\pi_w\otimes|\det|_w^{(1-a)/2})$ at $w\mid\ell$. Require $N=0$ and $\det(X-\rho(sigma))=(X-1)^a$ for every $\sigma$ in wild inertia $P_w$. Write $I_w$ for inertia and $q_w$ for the residue cardinality. Put $d_w=2$ if $\ell$ splits completely in $E$ and either $a=1$ or $\ell>16a^2$, and put $d_w=1$ otherwise. For every lift $\gamma\in I_w$ of a topological generator of $I_w/P_w$, require $z^{q_w-1}=1$ for each root $z$ of $\det(X-\rho(\gamma))$ of multiplicity at least $d_w$. Under this test, the period parameter in Definition~\ref{B9.06} has rank $a$, zero monodromy and the same inertia characteristic polynomials as $y$.
For $E/E^+$ unramified at every finite place and the standard
definite unitary group of rank $n\ge1$, definite isobaric base
change gives one family of components for all places as in
\ref{B6.04a}. The $\GL$ entries give Jacquet--Shalika rigidity
and fixed-level, compact-type and infinitesimal-character
finiteness. For a unitary cuspidal representation, at each
unramified place $w$ with Satake parameters, every such parameter
$\alpha$ satisfies $q_w^{-1/2}<|\alpha|<q_w^{1/2}$, where $q_w$
is the residue cardinality \cite{jacquet-shalika-1981}.

\item The finite general linear group inputs are \ref{B7.06}
and the regular formula of \ref{B7.05}. At a nonarchimedean
place use the double-coset filtration of the Jacquet module
of an induced representation and irreducibility of products on
disjoint unramified cuspidal lines. First-congruence-fixed
vectors fixed by the integral unipotent radical inject into the
Jacquet module. A first-congruence-fixed vector in an irreducible
implies first-congruence-fixed vectors in every member of its
cuspidal support. In characteristic zero generic unitarizable
irreducibles are complementary products, and a standard ordering
of multisegment factors gives an irreducible cosocle.

\item There is a global Artin map for each number field and
algebraic closure. For this map each algebraic Hecke character
has a unique continuous $\ell$-adic realization with its
prescribed infinity type. Above $\ell$ the character comparison
gives the rank-one Weil--Deligne parameter for each compatible
local component and every admissible correspondence of
\ref{B9.47} with the specified period data. Continuous de Rham
representations are potentially semistable. Integral crystalline
lattices have prismatic realizations and equivariant maps lift
uniquely to maps of crystals. Ray class fields realize ray groups.
A character of ideals prime to the conductor, trivial on the
ray-principal ideals, gives the unique open-kernel Galois
character with the prescribed Frobenius values. At
$v^n\Vert\mathfrak m$ its inertia image consists of the values
on principal ideals generated by elements prime to $v$ and
congruent to one modulo $\mathfrak m/v^n$. The prime-existence
entries are the finite Galois specialization and the
open-normal-subgroup assertion of \ref{B8.42p}.

\item For a regular integral parameter translated by $(a-1)/2$,
the discrete series has unique lowest compact type with weight
$\lambda+\rho_n-\rho_c$, where $\rho_n,\rho_c$ are the half-sums
of the positive noncompact and compact roots of its chamber.
Other weights arise by adding sums of positive roots. The
classification exhausts square-integrable irreducibles up to
order-preserving relabeling within each sign. The analytic
entry includes admissible coefficients and density for the
Levi submodule. The two convergent Siegel--Weil fields are
those of \ref{B8.13}. The first uses $a<b\le2a$ and
$a+\operatorname{Witt}(B)<b$. The second includes the volume
and nonsingular Gram assertions. In the Schr\"odinger model
a Fock vector in the convergent range has theta integral
equal to the holomorphic Siegel value at the special point.
The flat Siegel series continues with finite pole order.


\item For every prime $d>3$ with $d\equiv3\pmod4$, let
$K=\Q(\sqrt{-d})$ and $F_0=\Q(j(\mathcal O_K))$. There is a CM
elliptic curve $A/F_0$ with $j(A)=j(\mathcal O_K)$ and good
reduction away from $d$. At each finite place of $F_0$ there is
a local minimal equation of discriminant a local unit times
$-d^3$. For an imaginary quadratic $K$ of discriminant
less than $-4$, let $\Lambda$ be its integer lattice. For
every positive integer $n$, the ray class field of $(n)$ is
\[
 K\bigl(j(\Lambda),\, W_\Lambda(P):
 0\ne P\in(\C/\Lambda)[n]\bigr),\qquad
 W_\Lambda(P)=
 \frac{g_2(\Lambda)g_3(\Lambda)}{D(\Lambda)}\,\wp_\Lambda(P).
\]
The generator $j(\Lambda)$ can equivalently be divided by $1728$.
In particular at $n=1$ the field is $K(j(\Lambda))$.

\item The nonendoscopic hypotheses in $P$ have five factors. Its local
part includes normalized induction and Jacquet finiteness,
unitary and disjoint-product irreducibility, the geometric
lemma, segment and square-integrable classification, the
tempered product and generic unlinked criteria, multisegment
classification, exhaustion and standard cosocle, the
Moy--Prasad Jacquet isomorphism, local correspondence and
converse, depth-zero induction, degree-zero Brauer induction,
Chevalley units and the regular Green formula. Its global part
includes decay and finiteness, cuspidal spectrum, Jacquet--Shalika,
isobaric transfer, Gelfand--Kazhdan duality, polarized and
dual local--global compatibility, purity and cohomologicality,
reciprocity, ray fields, above-$\ell$ character and Weil--Deligne
comparison, and potential semistability. Its cohomology part
gives normalized Poitou--Tate, local and global finiteness
and Euler characteristics, local duality and unramified
orthogonality, Greenberg--Wiles, local finiteness, duality and
Euler characteristic over the completed field, and local
Bloch--Kato. The remaining factors are $\mathscr J$ of
\ref{B1.14k} and the Hecke--Tate functional equation.

\item The analytic and de Rham hypotheses in $P$ include the Weil action, existence, analyticity and classification for discrete series, trace class, convolution and restricted tensor factorizations, the Weil--Siegel, doubling and unramified integral formulas, and partial Hecke continuation. Its de Rham
count says that for a continuous rank-$a$ representation de
Rham at $w$, the filtration jumps at each compatible embedding
inducing $w$ have total multiplicity $a$.
For the small-prime estimate, $P$ includes the seven hypotheses $\mathscr H_2$ of Definition~\ref{B8.41h}. The three further local cohomology entries of
$P$ are those of \ref{B8.41p} for a finite discrete
$\Z/\ell^r$-module with decomposition-group action and a
valuation subring of the chosen algebraic closure.
\end{enumerate}

The additional entries of $F$ are as follows.
\begin{enumerate}
\item For every prime $g$ and rank $N$ there are polynomial-indexed
finite-$\GL$ series characters, with the series law for $N>0$.
Over a finite field $k$ there is a cuspidal irreducible in every
rank $r\ge1$. Let $P_r$ be the mirabolic subgroup, consisting of
matrices with last row $(0,\dots,0,1)$, and let $U_r$ be the upper
unitriangular subgroup. For every nontrivial additive character
$\psi_0$ of $k$, the restriction of each irreducible cuspidal
representation of $\GL_r(k)$ to $P_r$ is
$\operatorname{Ind}_{U_r}^{P_r}\psi_U$, where
$\psi_U(u)=\psi_0(\sum_{i<r}u_{i,i+1})$.
Weyl interlacing holds for the complex general linear and compact
inclusions of successive ranks. For every $q\ge1$ and every
reduced class $a\bmod q$ the least prime in the class is at most
$Cq^L$ for absolute $C,L>0$.

\item The $j$-value of any ideal lattice of an imaginary
quadratic field generates its Hilbert class field. For a finite idèle $t$ with ideal $\mathfrak b$, the automorphism $\operatorname{Art}_K(t)^{-1}$ of the Hilbert class field carries the value at $\mathfrak a$ to that at $\mathfrak b^{-1}\mathfrak a$. The values form one orbit over
the quadratic field. Nondegenerate Hermitian matrices of equal
rank over a quadratic extension of a nonarchimedean field of
characteristic different from two are congruent exactly when
their determinant quotient is a norm. The entries include the
local-to-global classification over a CM field and the Hasse
principle over a quadratic extension of number fields.
Quadratic class field theory gives the local characters, norm
theorem, exact sequence, product formula and scalar extension
compatibility.

\item For admissible nonarchimedean $\GL$ representations the
Iwahori projection onto the maximal compact torus invariants
of the upper-unipotent Jacquet module is bijective. Over
characteristic zero the shuffle algebra on an ordered
alphabet is polynomial on its Lyndon words. For the convergent
Rallis formula use $a<b\le2a$ with $B$ anisotropic or
$a+\operatorname{Witt}(B)<b$, compatible doubled Weil data, a
cuspidal source, vanishing lower lifts and a compact-finite
doubled test. The lifts are square integrable and their
Petersson pairing divided by target volume is the doubling
zeta value of \ref{B8.13}.

\item With finitely many Euler factors removed, the Hecke
$L$-function of a unitary character outside the norm-power family
is entire and nonzero on $\Re s\ge1$. For the trivial character the
corresponding assertion applies to $(s-1)\zeta^T(s)$, where $T$
is the set of omitted places. At full level every nontrivial
unitary character, including the nontrivial norm powers, has
continuation continuous and nonzero at one. There is an
absolute $C>0$ such that for every primitive Dirichlet character
$\chi$ of conductor $q>1$,
\begin{equation}\label{eq:p1c1-log-derivative}
 \left|\frac{L'}{L}(1,\chi)
 -\sum_{\substack{0<\Re\rho<1\\|\Im\rho|\le1}}
 \frac{\operatorname{ord}_\rho L}{1-\rho}\right|
 \le C\log(3q).
\end{equation}
The sum runs over distinct zeros with the indicated multiplicities.

\item The compact-finite subspace of an irreducible unitary
realization is admissible, analytic, dense, Lie-stable and
irreducible as a $(\mathfrak g,K)$-module. For $\mathrm U(p,q)$
the discrete-series entries give the lowest compact type with
weight $\lambda+\rho_n-\rho_c$, equivalence for the same
parameter, the holomorphic lowest-weight property and
equivalence after changing either frame of the same Hermitian
matrix. Globalization gives admissible closure and infinitesimal
equivalence. Smooth uniformly moderate cusp forms and their
invariant derivatives are bounded. At fixed open level,
growth exponent $A$ and decay order $M$, a finite set $L$ of
derivative words and a uniform constant bound the decay
seminorm by the sum of their moderate-growth seminorms.
Let $d$ be the greatest word length in $L$, with $d=0$ if $L$ is empty. Write $H$ for the group height and $H_{\rm q}$ for the quotient height. With $D_l$ denoting differentiation along the word $l$, the estimate gives
\[
 \sup_g |\phi(g)|H_{\rm q}(g)^M
 \le C\sum_{|l|\le d}\sup_g\frac{|D_l\phi(g)|}{H(g)^A},
\]
where the sum runs over all words of length at most $d$ in the chosen Lie algebra generators, including the empty word.

\item Over a nonarchimedean field of characteristic different
from two, smooth irreducible Heisenberg representations with a
prescribed nontrivial continuous central character exist and
are equivalent. The two real metaplectic realizations give the
twofold cover of the real symplectic group and its strongly
continuous unitary and Schwartz actions intertwining Heisenberg
translations. They supply the derived quadratic operators, the
Gaussian coefficient identity and the Gaussian eigenline over
the maximal compact. The general linear Howe decomposition
applies to complex general linear and unitary inclusions.
Kudla's splitting over a characteristic-zero nonarchimedean
field gives a smooth Schr\"odinger action for the split
unitary member. For a Levi element $a$, the action is
\[
 (\omega(a)\varphi)(x)=
 \chi(\det a)|N\det a|^{\dim V/2}\varphi(xa),
 \qquad H(w)\omega(g)=\omega(g)H(wg).
\]
Here $V$ is the partner space, $N$ denotes the quadratic norm
in the first formula, and $H(w)$ is the Heisenberg translation
for the row $w$ in the doubled matrix space. The action is
isometric for every additive Haar measure on the Schr\"odinger
matrix space. For the standard spaces $E^N$ and $E^{N'}$ with identity Gram
matrices and standard lattices, which are definite at every
archimedean place, and even product of ranks $NN'$, the rational
theta transformation is pullback by the rational pair
$(\gamma,\gamma')$ times
\[
 \prod_\tau\tau(\det\gamma)^{-a_\tau}
 \tau(\det\gamma')^{-a'_\tau},
 \qquad 2a_\tau=\epsilon'n'_\tau-N',\qquad
 2a'_\tau=\epsilon n_\tau+N .
\]
Here $\epsilon,\epsilon'$ are the two Hermitian signs and
$n_\tau,n'_\tau$ the splitting-character infinity exponents.
For a CM unitary pair with compatible Hecke characters there
is also an adelic Weil action, with a realization for each
good model. In the real polynomial oscillator model with
compatible Cartan frames, stable subspaces containing the
first member's coinvariant relations have bounded strict-chain
length. The sum of two proper such subspaces is proper.

\item For a nondegenerate CM Hermitian space with Witt frame
there are a compact subgroup $K$, a compact subset $\Omega$
of the modulus-one minimal parabolic, and a positive cone
parameter giving the Iwasawa and rational Siegel-set coverings.
The Siegel set $\mathfrak S$ is closed and has finite Haar volume.
The conjugates $a^{-1}pa$ for $p\in\Omega$ and cone elements $a$
stay in one compact set. For a fixed Haar measure there is $C$
such that for every nonnegative $f$ and measurable nonnegative
$h$ with $f(pak)\le h(t)$ on the Siegel set coordinates,
\[
 \int_{\mathfrak S} f(g)\,dg
 \le C\int_{\operatorname{cone}}h(t)
 \bigl(\delta_P(t)\textstyle\prod_i t_i\bigr)^{-1}\,dt .
\]
Here $t$ is the positive torus coordinate of $a$ and $\delta_P$
is the parabolic modulus in these coordinates. There are
$C',A$ with $H(g)\le C'H(\gamma g)^A$ for every $g\in\mathfrak S$
and every rational $\gamma$, where $H$ is the group height.
Smooth moderate cusp forms decay on $\mathfrak S$ to every
polynomial order in $H$.
\end{enumerate}
\end{definition}

\begin{lemma}[Simplicity of external tensor products]\label{B1.15ba}
Let $(A_i)_{i=1}^r$ be unital associative $\C$-algebras, and for each $i$ let $M_i$ be a simple left $A_i$-module of at most countable dimension over $\C$. Then $\operatorname{End}_{A_i}(M_i)=\C$, and the external tensor product $\bigotimes_{i=1}^r M_i$ is a simple left $\bigotimes_{i=1}^r A_i$-module. More generally, let $(A_v)_{v\in V}$ be a family of unital associative $\C$-algebras, let each $M_v$ be a simple left $A_v$-module of at most countable dimension over $\C$, and for each $v$ outside a finite subset $S_0\subset V$ fix a nonzero vector $s_v\in M_v$. Then the restricted tensor product $\bigotimes'_{v\in V}M_v$ formed with respect to $(s_v)_{v\notin S_0}$ is simple under the restricted tensor algebra.
\end{lemma}
\begin{proof}[Outline]
Countable dimension and Schur's lemma force scalar commutants, so each algebra prescribes images of finite independent tuples. Isolating pure tensors gives simplicity for finite products, and passing to a common finite set \(S\supseteq S_0\) yields the restricted case. The complete proof is in \S\ref{prf:B1.15ba}.
\end{proof}

\begin{proposition}[Inputs for the unitary and local theta arguments]\label{B1.15b}
For every triple $O$ of global and local inputs as in \ref{B9.01}, both products
$P,F$ of Definition~\ref{B1.15ac}, as used in Theorem~\ref{B8.01},
hold.
\end{proposition}
\begin{proof}[Outline]
Verifies each component of \(P\) and \(F\) in \ref{B1.15ac} from published automorphic, Galois, and theta theorems alongside direct proofs of doubling, trace class, Jensen--Blaschke bounds, Poisson summation, \ref{B1.15a}, and \ref{B1.15ba}. The complete proof is in \S\ref{prf:B1.15b}.
\end{proof}

\begin{theorem}[Construction of the common input]\label{B1.15}
Let $\mathcal R$ be the Galois representation assertion of
Definitions~\ref{B9.04}--\ref{B9.06}, with the parameter laws
of Definition~\ref{B9.47}. Then $\mathcal R$ implies $H$.
In particular, $H$ holds by Theorem~\ref{B9.02}.
\end{theorem}
\begin{proof}[Outline]
Obtain \(P,F\) by \ref{B1.15b}, the theta identity and tests by \ref{B8.11} and \ref{B8.18}, and the bounds from \ref{B8.24} and \ref{B8.41}. Use \ref{B8.01} to obtain an element of \(\mathcal H\), then apply \ref{B1.16}. The complete proof is in \S\ref{prf:B1.15}.
\end{proof}

\begin{theorem}[Three inputs give all characters]\label{B1.16}
Assume $\mathcal G$ and $\mathcal V$ and suppose $\mathcal H$
has an element. Then $H$ holds.
\end{theorem}
\begin{proof}
Apply Theorem~\ref{B1.17} and Lemma~\ref{B1.14}.
\end{proof}

\begin{theorem}[Three inputs give primitive nonvanishing]\label{B1.17}
Assume $\mathcal G$ and $\mathcal V$ and suppose $\mathcal H$
has an element. Then $H_{\rm prim}$ holds.
\end{theorem}
\begin{proof}[Outline]
Use \ref{B7.54} with $\mathcal G$ and $\mathcal H$ for connectedness with exponents $(2,0)$, then \ref{B5.33} for chain bounds under \ref{B6.01}. For primitive $\psi$, use \ref{B2.61} with $\mathcal V$ to get hosts with $\lambda=3/2$, then \ref{B1.18} for nonvanishing. The complete proof is in \S\ref{prf:B1.17}.
\end{proof}


% ---- p1c2 ----
\chapter{How arithmetic windows enter the zero-free criterion}\label{ch:arithmetic-windows}

Use the $L$-value convention of \ref{B1.10}. Write
$\mathcal H_\lambda(S)$ if, for every $Y_0\in\R$, some $Y\geq Y_0$
satisfies $[Y,Y^\lambda]\subset S$. In each window condition we may take
$Y\geq\max(Y_0,1)$, since $Y_0$ is arbitrary.

\section{The arithmetic interface}

A family $\mathcal F$ at $\varepsilon$ has scales $W(\mathcal F)$ and
objects $D_{x,i}$, $i\in I_{\mathcal F}(x)$.
Write $x_1(\mathcal F)>1$ for the threshold in the family data.
Write $E_{\mathcal L}(\mathcal F)$ if there are $C,x_0\in\R$ such that, for every $x\geq x_0$ and $i\in I_{\mathcal F}(x)$, the host and top set of $D_{x,i}$ satisfy the arithmetic requirement of Definition~\ref{B2.58} at $(\mathcal L,x,\varepsilon,C)$. This fixes the subgroup index at $o_{\mathcal L}$, represents the retained characters in the positive half of the host characters, and matches the designated character, when present, with the $q$-sign character. When no character is designated, the primitive representative of the $q$-sign lies outside the retained set. The top set must satisfy the coset counts \eqref{eq:top-shares} with tolerance $x^{-1+\varepsilon}$. For each primitive representative $\chi$ of an unretained positive host character or of a negative host character in Definition~\ref{B2.58}, the bounds are $|g_\chi(x)|\leq x^{1/2+\varepsilon}$ and $|(L'/L)(1,\chi)|\leq C\log x$.
For a relation $T$ on families and the real parameters below,
write $\mathcal T_{e_D,e_L}(\mathcal F)$ for
\[
\exists K,c_{HH},e_K,C_L,x_0\in\R\
\forall x\geq x_0,\quad x\in W(\mathcal F)
\Longrightarrow T(\mathcal F;e_D,e_L,K,c_{HH},e_K,C_L,x).
\]
The comparison $\mathcal C_{\mathcal L}$ asserts, for every
$e_D,e_L\in\R$, $\varepsilon>0$, and $\mathcal F$,
\[
\begin{gathered}
E_{\mathcal L}(\mathcal F)\land\mathcal T_{e_D,e_L}(\mathcal F)
\Longrightarrow\exists x_1\ \forall x\geq x_1,\\
x\in W(\mathcal F)\Longrightarrow
G_m^{\mathcal L}(x)\leq x^{1/2+(e_D+e_L+7)\varepsilon},
\end{gathered} \tag{1}
\]
with $G_m^{\mathcal L}$ from \ref{B1.21}.
Write $\mathcal W_{\mathcal L}(\lambda)$ for
\[
\begin{gathered}
\exists e_D,e_L\in\R\ \exists\varepsilon_1>0\
\forall\,0<\varepsilon<\varepsilon_1\ \exists\mathcal F,\\
E_{\mathcal L}(\mathcal F),\qquad
\mathcal H_\lambda(W(\mathcal F)),\qquad
\mathcal T_{e_D,e_L}(\mathcal F).
\end{gathered} \tag{2}
\]

\begin{definition}[Host conditions and the cyclic requirement]\label{def:p1c2-host}
At a scale $s=(x,\varepsilon,J)$ put $b=\exp(x^\varepsilon)$ and $\delta_1=x^{-\varepsilon/8}$. Let $E$ be the host field, $K=\Q(\sqrt{-d})$, and $N=Q_Sqd$, where $Q_S$ is the conductor parameter of the host. In its frame write $n=|A|$, $n_0=|A_0|$ and $o=[A:A_0]$. The host conditions with constants $C_{T_2},C_8,C_5,c_6,C_b$ are as follows.

The host has thickness constant $C_{T_2}$. Its thick set $T_2$ satisfies \eqref{eq:thick-size-bounds} with $C_T=C_{T_2}$ and $4L_w|T_2|\le n_0$, where
\[
 L_w=10\lceil405\log x\rceil+2\lceil30/\varepsilon\rceil+22.
\]
Its top set $TT\subset A$ satisfies $\delta_1 n/2\le |TT|\le 2\delta_1 n$ and the weighted bound \eqref{eq:bad-top-budget} with constant $C_b$. The conductor and frame satisfy $Q_S\le C_8 n$ and $J,o\le C_8$. The auxiliary ideal in $K$ is $\mathcal O_K$, and the host's splitting bound is at least $\lfloor x^{\varepsilon/(5J)}\rfloor$.

The level primes form the set $\{g_0,g_1,p,p',q_0\}$ and split in $K$. For $i=0,1$ one has $(kb)^8\le g_i\le 2(kb)^8$, $g_i\equiv1\pmod N$, and there is $\pi_i\in\mathcal O_K$ with $\pi_i\bar\pi_i=g_i$ and $\pi_i\equiv1\pmod{N\mathcal O_E}$. The exceptional prime sets $\mathcal B_r(g_0)$ and $\mathcal B_r(g_1)$ of Theorem~\ref{B2.51} are disjoint. The primes $p,p'$ split completely in $E$ and are $1$ modulo $\operatorname{lcm}(15,N)$. The prime $q_0$ is $1$ modulo $N$ and differs from $p,p'$. All three are at most $N^{C_8}$.

There is one $Y\in[x^{1/J},x]$ at which $q,d$ satisfy the prime ranges, splitting conditions and bound on $\log N$ of Lemma~\ref{B2.55}, the six quadratic twists satisfy the conductor and zero-free box conditions there, and the real host characters whose conductor is divisible by $q$ or $d$ satisfy the conductor and box bounds of Theorem~\ref{B2.52} at the heights specified before Theorem~\ref{B2.56}. For the primitive representatives of the two quadratic signs and their product, require
\[
 |g_\chi(x)|\le x^{1/2+\varepsilon},
 \qquad |(L'/L)(1,\chi)|\le C_5\log x .
\]
For the quadratic character modulo $d$, the absolute value of the real part of its logarithmic derivative at $1$ is at most $x^\varepsilon$. Every primitive host character satisfies \eqref{eq:thin-character-bound} with constant $c_6$.

A \emph{host requirement} is a condition on the scale, its frame, the arithmetic datum on that frame, a finite top set and an auxiliary rank $t\in\N$. At a host we use its own frame, arithmetic datum, top set and rank. The \emph{length bound} at this host means the finiteness and the weighted inequality of Definition~\ref{B6.01} for $E,t$ and $B=\sqrt b$. For a primitive $\psi$, the \emph{cyclic requirement} $A_\psi$ is the conjunction of this length bound and the arithmetic requirement of Definition~\ref{B2.58} for $TT$ at $(\mathcal L_\psi,x,\varepsilon,4)$, where $\mathcal L_\psi$ is the level of the powers of $\psi$ from Definition~\ref{B1.21}.
\end{definition}

\begin{theorem}[Arithmetic input to the window criterion]\label{B1.18}
Fix $\lambda>1$. Assume the family selection principle stated below. Suppose that for every primitive character $\psi$ of positive modulus there is a supply of hosts on $\lambda$-windows satisfying the cyclic requirement $A_\psi$ of Definition~\ref{def:p1c2-host}. Then
$L(s,\psi)\neq0$ for every such $\psi$ and $\Re s>1/2$.

A supply of hosts on $\lambda$-windows for a requirement $A$ means that for each $0<\varepsilon<1/52$ there are real constants $J,C_8\geq1$, $C_5,c_6,C_b>0$, $C_{T_2},x_0\in\R$, and a set $W\subset\R$ with $\mathcal H_\lambda(W)$ for which the following holds.
For every $x\in W$ with $x\geq x_0$, one can choose rank data at $s=(x,\varepsilon,J)$ satisfying the three ranges for $n,m,k$ in \eqref{eq:host-rank-ranges} and the identity $r+t=8fr+4+T_On_W$ of Theorem~\ref{B8.59}, where $f=2n$, $r=k+m-1$, and $t$ is the auxiliary rank. On these ranks there must be a host and a prime $q_0$ satisfying Definition~\ref{def:p1c2-host} with the chosen constants, including the same thickness constant $C_{T_2}$ at every point. The scale, frame, arithmetic datum, top set and auxiliary rank of this host must satisfy $A$.

For family selection use the arithmetic parameter families of the construction in Theorem~\ref{B8.05}. Their indexed models carry the design and lattice data with common parameter constants. In (1) and (2), take $T$ to mean that at $x$ one base model and one pair of integer weight functions admit a model at every translation $u_0\in A$, with the prescribed offset and those same top weights. The selected models have the same frame, $m,f,q,d$, and realise the global, local and arithmetic chain data with all the estimates in the conclusion of Theorem~\ref{B5.33}, with division, loss and height exponents at most $e_D,e_L,c_{HH}$ and window constants $K,e_K,C_L$.
For each $\lambda'>1$ and request $A$, let
$\mathcal C_{\lambda',A}$ be the families with $\lambda'$-windows
satisfying $A$ for every $x>x_1(\mathcal F)$ and every
$i\in I_{\mathcal F}(x)$. The family selection principle is the following implication for every such $\lambda'$ and $A$. Suppose there is a supply of hosts for $A$ on $\lambda'$-windows. Suppose also that each $\mathcal F\in\mathcal C_{\lambda',A}$ has a threshold $x_V$ such that the length bound of Definition~\ref{def:p1c2-host} holds at every $x\geq x_V$ and every $i\in I_{\mathcal F}(x)$, for the field, auxiliary rank and weight of that model. Then there are $e_D,e_L\in\R$ and $\varepsilon_1>0$ such that each $0<\varepsilon<\varepsilon_1$ admits a family at $\varepsilon$ in $\mathcal C_{\lambda',A}$ satisfying $\mathcal T_{e_D,e_L}$.
\end{theorem}

\begin{proof}[Outline]
At each sufficiently large allowed scale, cyclicity gives the arithmetic condition of \ref{B2.58} with \(C=4\) and the length bound of \ref{B6.01}. Family selection gives (2), \ref{B2.20} gives (1) at the same index, and \ref{B1.19} and \ref{B1.22} give the conclusion. The complete proof is in \S\ref{prf:B1.18}.
\end{proof}

\begin{lemma}[Covering primitive characters by cyclic levels]\label{B1.19}
Fix $\lambda>1$ and a property $Q$ of primitive characters.
For every $\psi$ satisfying $Q$, assume (1), the supply (2) at
$\lambda$, and \ref{B1.22} for every window exponent $>1$ at
$\mathcal L_\psi$. Then $L(s,\psi)\neq0$ whenever
$Q(\psi)$ and $\Re s>1/2$.
Each $\psi$ may use a containing level with these properties and its own
$\lambda>1$.
\end{lemma}

\begin{proof}
Apply \ref{B1.20} to the assigned level.
The cyclic level contains $\psi$, its own primitive reduction.
\end{proof}

\begin{lemma}[Removing the exponent loss]\label{B1.20}
Fix $\lambda\in\R$, a family type with its map to the underlying
family data, and a relation $T$.
Assume (1), (2), and the displayed implication of Theorem~\ref{B1.22} for this $\lambda$ and every $0<\theta<1$. Then every kept character has nonvanishing $L$-values
throughout $\Re s>1/2$.
\end{lemma}

\begin{proof}[Outline]
Choose \(1/2<\theta<\min(\Re s,1)\) and \(\varepsilon\) with \(1/2+(e_D+e_L+7)\varepsilon\leq\theta\). On windows from (2) past the threshold in (1), comparison gives \(G_m^{\mathcal L}(y)\leq y^\theta\), so the window implication gives nonvanishing at \(s\). The complete proof is in \S\ref{prf:B1.20}.
\end{proof}

\section{Levels and positive weighted families}

\begin{definition}[Levels and cyclic levels]\label{B1.21}
A level has a finite kept set $\mathcal X_{\mathcal L}$ of primitive
characters with positive moduli. The set contains the modulus-one
character and is closed under inversion. The level also specifies
$o_{\mathcal L}\in\N$ and may designate a real odd character
$\xi\in\mathcal X_{\mathcal L}$. Put
\[
\begin{gathered}
w_{\mathcal L}(\chi)=
\begin{cases}
1,&\chi=\xi\text{ when }\xi\text{ is specified},\\
(3o_{\mathcal L}+1)^{-1},&\text{otherwise},
\end{cases}\\
G_m^{\mathcal L}(y)=G_z(y)+
\sum_{\chi\in\mathcal X_{\mathcal L}\setminus\{1\}}
w_{\mathcal L}(\chi)\Re g_\chi(y).
\end{gathered}
\]
Use $G_z,g_\chi$ from \ref{B1.26}. Write
$\mathrm{ZF}_{\mathcal L}(\theta)$ if $L(s,\chi)\neq0$ for every
$\chi\in\mathcal X_{\mathcal L}$ and every $s\in\C$ with $\Re s>\theta$.

A subgroup $\Xi$ of characters modulo $M\geq1$ gives a level
by primitive reduction, with $o_{\mathcal L}=|\Xi|$ and no distinguished
character. Reduction preserves inversion and sends the principal
character to modulus one, so the axioms hold.
For $\Xi=\langle\psi\rangle$, write $\mathcal L_\psi$. If $\psi$
has order $o$, each nonprincipal weight is $1/(3o+1)$.
\end{definition}

\begin{theorem}[The level window criterion]\label{B1.22}
For a level $\mathcal L$, $\lambda>1$ and $0<\theta<1$,
\[
\mathcal H_\lambda(\{y\in[1,\infty):G_m^{\mathcal L}(y)\leq y^\theta\})
\Longrightarrow\mathrm{ZF}_{\mathcal L}(\theta).
\]


\end{theorem}

\begin{proof}
Apply \ref{B1.24} with $\delta=\lambda-1$ and $C=1$.
The identities in \ref{B1.23} transfer the window bound and
the nonvanishing conclusion.
\end{proof}

\begin{lemma}[Realising a level as a positive weighted family]\label{B1.23}
Every level $\mathcal L$ admits a family $F$ of \ref{B1.25} with
$G_m(F,y)=G_m^{\mathcal L}(y)$ for all $y\in\R$ and
$\mathrm{ZF}(F;\theta)\Longleftrightarrow
\mathrm{ZF}_{\mathcal L}(\theta)$ for all $\theta\in\R$.
\end{lemma}

\begin{proof}[Outline]
Index \(F\) by \(\mathcal{X}_{\mathcal{L}}\) with weight \(1\) at modulus \(1\) and \(w_{\mathcal{L}}(\chi)>0\) elsewhere for admissibility. Dropping \(n=0\) and matching principal and nonprincipal terms via \ref{B1.26} identifies \(G_m\), with identical nonvanishing indices. The complete proof is in \S\ref{prf:B1.23}.
\end{proof}

\begin{theorem}[Pointwise power-window criterion]\label{B1.24}
Let $F$ be a positive weighted character family,
$0<\theta<1$, $\delta>0$, and $C\in\R$. If
\[
\forall Y_0\in\R\ \exists Y\geq Y_0\
\forall y\in[Y,Y^{1+\delta}],\quad G_m(F,y)\leq Cy^\theta,
\tag{3}
\]
then $\mathrm{ZF}(F;\theta)$.
\end{theorem}

\begin{proof}
By \ref{B1.27}, (3) gives the integral window condition.
Apply \ref{B1.28}.
\end{proof}

\begin{definition}[Positive weighted character families]\label{B1.25}
A family $F$ has a finite index set $I$, positive integer moduli $N_i$,
complex Dirichlet characters $\chi_i$ modulo $N_i$, and weights $a_i>0$,
with $N_i=1$ or $\chi_i\neq1$ for each $i$.
Characters may be imprimitive.
For $\theta\in\R$, write $\mathrm{ZF}(F;\theta)$ if
$L(s,\chi_i)\neq0$ for every $i\in I$ and every
$s\in\C$ with $\Re s>\theta$.
\end{definition}

\begin{definition}[Centred smoothed prime sums]\label{B1.26}
Let $\Lambda$ be the von Mangoldt function and $\gamma$ Euler's
constant. For $y\in\R$, set
\[
\begin{gathered}
S_\chi(y)=\sum_{1\leq n\leq\lfloor y\rfloor_+}
\Lambda(n)\chi(n)\max(y/n-1,0),\\
G_\chi(y)=
\begin{cases}
S_\chi(y)-y\log y+(1+\gamma)y,&\chi=1,\\
S_\chi(y)+yL'(1,\chi)/L(1,\chi),&\chi\neq1.
\end{cases}
\end{gathered}
\]
Put $G_m(F,y)=\sum_{i\in I}a_i\Re G_{\chi_i}(y)$.
Here $\lfloor y\rfloor_+=\max(\lfloor y\rfloor,0)$.
Use $\log y=\log|y|$ for $y<0$ and $\log0=0$.
Write $G_z=G_1$ at modulus one and $g_\chi=G_\chi$ for nontrivial
primitive $\chi$. Put $g_F(u)=G_m(F,e^u)$, $u\in\R$.
\end{definition}

\section{From pointwise bounds to integral windows}

\begin{lemma}[An integrable upper envelope]\label{B1.27}
For $\theta,\delta\in\R$, let $\mathrm{WH}(F;\theta,\delta)$
mean that every $\eta>0$ and $Y_0\in\R$ admit $Y\geq Y_0$ and
$R:\R\to\R$ with $y\mapsto R(y)y^{-\theta-1}\in L^1(J_Y)$,
where $J_Y=[Y,Y^{1+\delta}]$, and
\[
\forall y\in J_Y,\quad 0\leq R(y),\quad G_m(F,y)\leq R(y),
\qquad
\int_{J_Y}R(y)y^{-\theta-1}\,dy\leq Y^\eta. \tag{4}
\]
For every $\theta,C\in\R$ and $\delta>0$, (3) implies
$\mathrm{WH}(F;\theta,\delta)$.
\end{lemma}

\begin{proof}
Fix $\eta>0,Y_0$. Choose a window in (3) with $Y\geq\max(Y_0,1)$
and $|C|\delta\log Y\leq Y^\eta$, possible since
$(\log Y)/Y^\eta\to0$.
Set $R(y)=|C|y^\theta$ for $y>0$ and $R(y)=0$ for $y\leq0$.
On $J_Y$, $R$ is nonnegative and dominates $G_m(F,y)$.
The integrand $R(y)y^{-\theta-1}=|C|/y$ is continuous there, and
\[
\int_{J_Y}R(y)y^{-\theta-1}\,dy
=|C|\log\frac{Y^{1+\delta}}Y
=|C|\delta\log Y\leq Y^\eta.
\]
\end{proof}

\begin{theorem}[Integral window criterion for a character family]\label{B1.28}
For a positive weighted character family $F$, $0<\theta<1$, and
$\delta>0$, one has
$\mathrm{WH}(F;\theta,\delta)\Longrightarrow\mathrm{ZF}(F;\theta)$.
\end{theorem}

\begin{proof}
By \ref{B1.30}, the Laplace transform of $g_F$ on $u>0$ continues
meromorphically to $\Re s>0$ as $\Phi$, with the pole data of \ref{B1.33}.
Its pole set $Z_F$ contains all zeros of the family $L$-functions
with positive real part.
Apply \ref{B1.29} and then \ref{B1.31} to obtain
$\Re\rho\leq\theta$ for $\rho\in Z_F$.
A zero $s$ with $\Re s>\theta$ contradicts \ref{B1.32} if
$\Re s\geq1$, and the bound on $Z_F$ if $\Re s<1$.
\end{proof}

\begin{lemma}[Logarithmic windows]\label{B1.29}
For $g:\R\to\R$ and $\theta,\delta\in\R$, let
$\mathrm{WH}_U(g;\theta,\delta)$ mean that every $\eta>0$ and $U_0\in\R$
admit $U\geq U_0$ and $R_U:\R\to\R$ with
$u\mapsto R_U(u)e^{-\theta u}\in L^1(I_U)$,
where $I_U=[U,(1+\delta)U]$, and
\[
\forall u\in I_U,\quad 0\leq R_U(u),\quad g(u)\leq R_U(u),
\qquad
\int_{I_U}R_U(u)e^{-\theta u}\,du\leq e^{\eta U}. \tag{5}
\]
For every character family $F$, $\theta\in\R$, and $\delta>0$,
$\mathrm{WH}(F;\theta,\delta)\Longrightarrow
\mathrm{WH}_U(g_F;\theta,\delta)$.
\end{lemma}

\begin{proof}[Outline]
Choose \(Y\ge\max(e^{U_0},1)\) with envelope \(R\) from the hypothesis, and set \(U=\log Y\) and \(R_U(u)=R(e^u)\) so \(\exp\) maps \(I_U\) bijectively onto \(J_Y\). Folland's change-of-variables theorem for \(u\mapsto e^u\) yields integrability and the bound \(e^{\eta U}\). The complete proof is in \S\ref{prf:B1.29}.
\end{proof}


% ---- p1c3 ----
\chapter{Zeros become poles}\label{ch:zeros-become-poles}

Throughout, characters have positive integer modulus. Fix a family
\(F=(\chi_i,a_i)_{i\in I}\) as in Definition~\ref{B1.25} and use the sums of
Definition~\ref{B1.26}. Put
\[
 g(u)=G_m(F,e^u)=\sum_{i\in I}a_i\Re G_{\chi_i}(e^u),
 \qquad
 \mathcal M H(s)=\int_1^\infty H(y)y^{-s-1}\,dy.
\]
Put \(D_\chi(s)=L'(s,\chi)/L(s,\chi)\) where defined.
Use real logarithms in powers.

\section{The analytic interface}

\begin{theorem}[Pole data for a character family]\label{B1.30}
Choose the functions of \ref{B1.34} for \(\chi_i,\chi_i^{-1}\).
The sum
\[
 \Phi(s)=\sum_{i\in I}\frac{a_i}{2}
 \bigl(\Phi_{\chi_i}(s)+\Phi_{\chi_i^{-1}}(s)\bigr)
\]
gives pole data for \(g,Z_F\) as in \ref{B1.33}.
\end{theorem}
\begin{proof}[Outline]
Measurability and growth follow from \ref{B1.38}, while \ref{B1.45} and \ref{B1.37} allow termwise integration on \(\Re s>2\). Applying \ref{B1.34} and \ref{B1.45} over finite \(I\) isolates \(Z_F\) and gives analyticity off \(Z_F\), with residues from \ref{B1.46}. The complete proof is in \S\ref{prf:B1.30}.
\end{proof}

\begin{theorem}[Integral windows exclude poles]\label{B1.31}
Let \(g:\R\to\R\), \(\Phi:\C\to\C\), and \(Z\subset\C\)
satisfy \ref{B1.33}. Fix \(0<\theta<1\) and \(\delta>0\).
Assume the logarithmic window condition
\(\mathrm{WH}_U(g;\theta,\delta)\) of Lemma~\ref{B1.29}.
Then \(\Re\rho\le\theta\) for every \(\rho\in Z\).
\end{theorem}
\begin{proof}
This is Proposition~\ref{B1.54e}, proved in Chapter~\ref{ch:p1c4}. Chapter~\ref{ch:p1c5} of the appendix gives the proof that the archive formalises.
\end{proof}

\begin{lemma}[Nonvanishing on the closed half-plane]\label{B1.32}
For \(i\in I\) and \(\Re s\ge1\), \(L(s,\chi_i)\ne0\),
including the assigned modulus-one value at \(1\).
\end{lemma}
Davenport gives nonvanishing at regular points
\cite[Sections~4, 13--14]{davenport-2000}.
At modulus one and \(s=1\), \ref{B1.10} gives
\(\zeta(1)=(\gamma-\log(4\pi))/2\ne0\).
% [Springer edition record](https://link.springer.com/book/9780387950976)

\begin{definition}[Zero set and pole data]\label{B1.33}
Define
\[
 Z_F=\{\rho\in\C:\Re\rho>0,\ 
 \exists i\in I,\ L(\rho,\chi_i)=0
 \ \hbox{or}\ L(\bar\rho,\chi_i)=0\}.
\]
For \(g:\R\to\R\), \(\Phi:\C\to\C\), and \(Z\subset\C\),
\emph{pole data} require Borel measurability of \(g\) and some
\(C_0\in\R\) with
\[
 |g(u)|\le C_0e^{2u}\quad(u\ge0),\qquad
 \Phi(s)=\int_0^\infty g(u)e^{-su}\,du\quad(\Re s>2).
\]
Every point of \(Z\) has positive real part.
\(\Phi\) is holomorphic on \(\{s:\Re s>0\}\setminus Z\).
The set \(Z\) has no accumulation in \(\Re s>0\).
Each \(\rho\in Z\) admits \(c_\rho\ne0\) and \(h_\rho\)
analytic near \(\rho\) with
\[
 \Phi(s)=\frac{c_\rho}{s-\rho}+h_\rho(s)
 \quad(0<|s-\rho|\ \hbox{sufficiently small}).
\]
If \(\rho\) is real, \(c_\rho\) is a positive real number.
\end{definition}

\section{One character and the point \(1\)}

\begin{lemma}[Continuation for one character]\label{B1.34}
Let \(\chi\) have modulus \(N\ge1\), with \(N=1\) or
\(\chi\ne1\). There exists \(\Phi_\chi:\C\to\C\) such that
\[
 \Phi_\chi(s)=\int_0^\infty G_\chi(e^u)e^{-su}\,du
 \quad(\Re s>2).
\]
For each \(\rho\ne0\), \(\Phi_\chi\) is analytic at \(\rho\) if
\(L(\rho,\chi)\ne0\), and there are \(m\in\N\) and \(h\) analytic near
\(\rho\) with
\[
 m=0\ \Longleftrightarrow\ L(\rho,\chi)\ne0,\qquad
 \Phi_\chi(s)=\frac{c_\chi(\rho)}{s-\rho}+h(s)
\]
on a punctured neighbourhood of \(\rho\).
Here \(c_\chi(\rho)=m/[\rho(1-\rho)]\) if \(\rho\ne1\),
and \(c_\chi(1)=0\). At a zero, \(m\) is its multiplicity.
The zeros of \(L(\,\cdot\,,\chi)\) have no accumulation point in \(\C\).
\end{lemma}
\begin{proof}[Outline]
For \(\chi\ne 1\), apply \ref{B1.36} with \(b=0\) to \(L(\cdot,\chi)\) and obtain the integral from \ref{B1.41} and \ref{B1.37}. For \(\chi=1\), apply \ref{B1.36} with \(b=1\) to \((s-1)\zeta(s)\) to match \ref{B1.43} via the Laurent expansion at \(1\), then apply \ref{B1.37}. The complete proof is in \S\ref{prf:B1.34}.
\end{proof}

\begin{lemma}[Local logarithmic derivative]\label{B1.35}
Let \(f\) be analytic near \(\rho\) with a nonzero germ.
Then \(m=\operatorname{ord}_\rho f\in\N\), with
\(m=0\iff f(\rho)\ne0\), and some \(h\) analytic at \(\rho\) satisfies
\(f'(s)/f(s)=m/(s-\rho)+h(s)\) on a punctured neighbourhood of \(\rho\).
\end{lemma}
This is the logarithmic derivative form of the standard local factorization
theorem \cite[Chapter~4]{ahlfors-1979}.

\begin{lemma}[Removing the apparent singularity]\label{B1.36}
Let \(f:\C\to\C\) be entire, \(f(1)\ne0\), and \(b\in\C\).
Write \(D=f'/f\), with division by zero assigned the value zero, and put
\[
 N(s)=D(s)-sD(1)+b(1-s),\qquad
 \Phi_f(s)=-\frac{N(s)}{s(s-1)}\quad(s\ne1).
\]
It extends analytically at \(1\), with
\(\Phi_f(1)=D(1)+b-D'(1)\).
It is analytic where \(s\ne0\) and \(f(s)\ne0\).
At a nonzero zero \(\rho\) of order \(m\),
\[
 \Phi_f(s)=\frac{m}{\rho(1-\rho)(s-\rho)}+h_\rho(s)
\]
on a punctured neighbourhood of \(\rho\), with \(h_\rho\) analytic at
\(\rho\).
The zeros of \(f\) are isolated.
\end{lemma}
\begin{proof}[Outline]
Since \(N(1)=0\), \(N(s)/(s-1)\) extends analytically to \(1\), and the identity theorem with \(f(1)\ne 0\) isolates zeros. At a zero \(\rho\ne 0,1\), \ref{B1.35} expands \(D\) and splitting \(-1/[s(s-1)]\) at \(\rho\) yields the simple pole with residue \(m/[\rho(1-\rho)]\). The complete proof is in \S\ref{prf:B1.36}.
\end{proof}

\section{The Mellin calculations}

\begin{lemma}[Exponential substitution]\label{B1.37}
For \(G:\R\to\C\) and \(s\in\C\),
\[
 \int_0^\infty G(e^u)e^{-su}\,du
 =\int_1^\infty G(y)y^{-s-1}\,dy,
\]
where the Bochner integral has value zero for a nonintegrable integrand.
Both integrands are absolutely integrable for \(G=G_\chi\) and \(\Re s>2\).
\end{lemma}
The smooth inverse maps \(y=e^u\) and \(u=\log y\) preserve Lebesgue
measurability. Change of variables for the absolute values
\cite[Theorem~2.47(a)]{folland-1999} gives equivalence of integrability.
In the integrable case the same substitution gives the identity.
For \(G_\chi\), \ref{B1.38} gives the majorant
\(C_\chi e^{-(\Re s-2)u}\) when \(\Re s>2\).

\begin{lemma}[Measurability and quadratic growth]\label{B1.38}
For every \(\chi\), \(G_\chi:\R\to\C\) is Borel measurable,
and some \(C_\chi\in\R\) satisfies
\[
 |G_\chi(y)|\le C_\chi y^2\qquad(y\ge1).
\]
For \(f:\N\to\C\) with \(|f(n)|\le1\),
\[
 \left|\sum_{1\le n\le y}\Lambda(n)f(n)(y/n-1)\right|
 \le2y^2\qquad(y\ge1).
\]
\end{lemma}
\begin{proof}
For \(t\in\R\), write \(t_+=\max(t,0)\).
Since \(0\le\Lambda(n)\le n\) and \((y/n-1)_+\le y/n\), each
summand has norm at most \(y\), and there are at most \(2y\) terms.
Use \(0\le\log y\le y\) to take \(C_\chi=3+|1+\gamma|\)
for principal \(\chi\) and \(C_\chi=2+|D_\chi(1)|\) otherwise.
On each Borel level set of \(\lfloor y\rfloor_+\), the ramp is a
finite sum of Borel functions, and the centring terms are Borel.
\end{proof}

\begin{lemma}[Transform of a single ramp]\label{B1.39}
For \(n\in\N\), \(n\ge1\), and \(s\in\C\) with
\(\sigma=\Re s>1\),
\[
 \int_1^\infty (y/n-1)_+y^{-s-1}\,dy
 =\frac{n^{-s}}{s(s-1)},\qquad
 \int_1^\infty |(y/n-1)_+y^{-s-1}|\,dy
 =\frac{n^{-\sigma}}{\sigma(\sigma-1)}.
\]
\end{lemma}
By \cite[Equation~(5.9.1), \(\mu=\nu=1\)]{olver-2010},
the substitution \(y=ne^u\) gives
\(n^{-s}((s-1)^{-1}-s^{-1})\) for the first integral.
Taking the modulus replaces \(s\) by \(\sigma\) in the integrand.
% [DLMF (5.9.1)](https://dlmf.nist.gov/5.9.E1)

\begin{lemma}[Termwise transformation]\label{B1.40}
Let \(f:\N\to\C\), \(\Re s=\sigma>1\), and suppose
\(\sum_{n\ge1}|f(n)|\Lambda(n)n^{-\sigma}<\infty\).
If \(E:(1,\infty)\to\C\) is integrable, then
\[
 \begin{split}
 &\int_1^\infty
 \left(\sum_{1\le n\le y}\Lambda(n)f(n)(y/n-1)y^{-s-1}
 +E(y)\right)\,dy\\
 &\hspace{1cm}=
 \frac{\sum_{n\ge1}f(n)\Lambda(n)n^{-s}}{s(s-1)}
 +\int_1^\infty E(y)\,dy .
 \end{split}
\]
\end{lemma}
By \ref{B1.39}, the sum of the \(L^1\)-norms of the ramp terms is
\(\sum_{n\ge1}|f(n)|\Lambda(n)n^{-\sigma}/[\sigma(\sigma-1)]\).
Folland's Theorem~2.25 \cite{folland-1999} gives the stated interchange.

\begin{lemma}[Transform for a nonprincipal character]\label{B1.41}
For nonprincipal \(\chi\) and \(\Re s>1\),
\[
 \mathcal M G_\chi(s)
 =-\frac{D_\chi(s)-sD_\chi(1)}{s(s-1)}.
\]
\end{lemma}
\begin{proof}
Logarithmic differentiation of the Euler product
\cite[Section~11.5, p.~231]{apostol-1976} gives
\[
 \sum_{n\ge1}\chi(n)\Lambda(n)n^{-s}=-D_\chi(s).
\]
In \ref{B1.40} take \(E(y)=D_\chi(1)y^{-s}\) and use
\(\int_1^\infty y^{-s}\,dy=(s-1)^{-1}\).
\end{proof}
% [Euler product](https://dlmf.nist.gov/25.15.E2)

\begin{lemma}[The logarithmic centring term]\label{B1.42}
For \(\Re s>1\), the function \((\log y)y^{-s}\) is
absolutely integrable on \((1,\infty)\), and
\[
 \lim_{y\to\infty}(\log y)y^{1-s}=0,\qquad
 \int_1^\infty(\log y)y^{-s}\,dy=\frac1{(s-1)^2}.
\]
\end{lemma}
With \(y=e^u\), equation~(5.9.1) of \cite{olver-2010} with
\(\mu=1,\nu=2\) gives \(\int_0^\infty ue^{-(s-1)u}\,du=(s-1)^{-2}\).
Since \(\Re(s-1)>0\), the integral converges absolutely and
\(ue^{-(s-1)u}\to0\).

\begin{lemma}[Transform for the modulus-one character]\label{B1.43}
For \(1\) modulo \(1\) and \(\Re s>1\),
\[
 \mathcal M G_1(s)=
 -\frac{\zeta'(s)/\zeta(s)+1/(s-1)+1-(1+\gamma)s}{s(s-1)}.
\]
\end{lemma}
\begin{proof}
Here
\(G_1(y)=\sum_{1\le n\le y}\Lambda(n)(y/n-1)
-y\log y+(1+\gamma)y\).
Use \(\sum_{n\ge1}\Lambda(n)n^{-s}=-\zeta'(s)/\zeta(s)\).
Applying \ref{B1.40} to the ramp and \ref{B1.42} to the
logarithmic term gives
\[
 \mathcal M G_1(s)=
 -\frac{\zeta'(s)/\zeta(s)}{s(s-1)}
 -\frac1{(s-1)^2}+\frac{1+\gamma}{s-1}.
\]
The identity \(s/(s-1)=1+1/(s-1)\) gives the stated numerator.
\end{proof}

\section{Conjugation and residues}

\begin{lemma}[Conjugation of the Hurwitz function]\label{B1.44}
Let \(a\in\R/\Z\). Write \(H_a\) for the archive's Hurwitz
zeta function with parameter modulo \(1\).
For every \(s\in\C\),
\[
 H_a(\bar s)=\overline{H_a(s)},\qquad
 \Gamma_{\R}(\bar s)=\overline{\Gamma_{\R}(s)},\qquad
 \Gamma_{\R}(s)=\pi^{-s/2}\Gamma(s/2).
\]
The difference \(H_a(s)-\Gamma_{\R}(s)^{-1}/(s-1)\) is entire
with value \(H_a(1)\) at \(1\).
\end{lemma}
\begin{proof}[Outline]
Use gamma conjugation for \(J\) and the Hurwitz residue to subtract \(J(s)/(s-1)\), obtaining an entire extension with value \(H_a(1)\) at \(1\). Extend conjugation from the Hurwitz series by the identity theorem, then restore the polar term off \(1\). The complete proof is in \S\ref{prf:B1.44}.
\end{proof}
% [DLMF 25.11](https://dlmf.nist.gov/25.11)

\begin{lemma}[Inverse characters and conjugation]\label{B1.45}
For every \(\chi\), \(s\in\C\), and \(y\in\R\),
\[
 L(\bar s,\chi^{-1})=\overline{L(s,\chi)},\qquad
 G_{\chi^{-1}}(y)=\overline{G_\chi(y)}.
\]
Consequently \(L(\rho,\chi^{-1})=0\) if and only if
\(L(\bar\rho,\chi)=0\).
For nonprincipal \(\chi\), also
\(D_{\chi^{-1}}(1)=\overline{D_\chi(1)}\).
\end{lemma}
\begin{proof}[Outline]
Conjugate the Hurwitz expansion of \(L(s,\chi)\) via \ref{B1.44} and \(\overline{\chi(j)}=\chi^{-1}(j)\), then differentiate at \(1\) for nonprincipal \(\chi\) to get \(D_{\chi^{-1}}(1)=\overline{D_\chi(1)}\). Real ramp weights and centring terms give the \(G_\chi\) identity. The complete proof is in \S\ref{prf:B1.45}.
\end{proof}
% [Hurwitz expansion](https://dlmf.nist.gov/25.15.E3)

\begin{lemma}[Positive weighted multiplicities]\label{B1.46}
For \(\rho\in Z_F\), let \(m_i(\rho),m_i^-(\rho)\) be the
multiplicities at \(\rho\) for \(\chi_i,\chi_i^{-1}\), respectively,
with multiplicity zero at nonzeros. The function in \ref{B1.30}
has residue
\[
 c_\rho=\frac{S_\rho}{\rho(1-\rho)},\qquad
 S_\rho=\sum_{i\in I}\frac{a_i}{2}
 \bigl(m_i(\rho)+m_i^-(\rho)\bigr)>0.
\]
Thus \(c_\rho\ne0\). If \(\rho=\beta\in\R\),
then \(0<\beta<1\) and \(c_\beta>0\).
\end{lemma}
\begin{proof}
By \ref{B1.32} and \ref{B1.45}, \(0<\Re\rho<1\).
Sum the expansions of \ref{B1.34} on a common punctured
neighbourhood. By \ref{B1.34} and \ref{B1.45}, an index \(i_0\)
witnessing \(\rho\in Z_F\) contributes at least \(a_{i_0}/2>0\)
to \(S_\rho\). The other terms are nonnegative. The denominator is
nonzero, and for real \(\rho=\beta\) it is positive.
\end{proof}


% ---- p1c4 ----
\chapter{A Gaussian window excludes poles}\label{ch:p1c4}

This chapter proves Theorem~\ref{B1.31} with a Gaussian kernel. The archive proves the same theorem by estimates for high derivatives of the Laplace transform; that proof is given in Chapter~\ref{ch:p1c5} of the appendix. The two proofs share Lemmas~\ref{B1.54} and~\ref{B1.55}. The Gaussian proof below is not formalised in the archive; Lemmas~\ref{B1.54a}--\ref{B1.54d} and Proposition~\ref{B1.54e} therefore carry no Lean reference.

We use the pole data $(g,\Phi,Z)$ of Definition~\ref{B1.33}. Write
\[
 H=\{s\in\C:\Re s>0\},\qquad
 B(z,r)=\{w\in\C:|w-z|<r\},\qquad
 \overline B(z,r)=\{w\in\C:|w-z|\le r\}.
\]
At $\rho\in Z$, denote the coefficient of $(s-\rho)^{-1}$ by
$c_\rho$. It is nonzero, and it is positive real when $\rho$ is real.

\begin{lemma}[The absolute-convergence region]\label{B1.54}
For every $\sigma_0>2$ there is $M\in\R$ such that
$|\Phi(s)|\le M$ whenever $\Re s\ge\sigma_0$.
More precisely, a growth bound $|g(u)|\le C_0e^{2u}$ for $u\ge0$
gives $M=C_0/(\sigma_0-2)$. Every $\rho\in Z$ satisfies $\Re\rho\le2$.
\end{lemma}

The Laplace majorant is $C_0e^{-(\sigma_0-2)u}$
\cite[Chapter~II, \S3]{widder-1941}.
If $\Re\rho>2$, take $\sigma_0=(2+\Re\rho)/2$.
Local boundedness gives $(s-\rho)\Phi(s)\to0$. The pole expansion
gives the limit $c_\rho\ne0$.
This is the bounded case of the removable-singularity criterion
\cite[Chapter~IV, \S3.1, Theorem~7]{ahlfors-1979}.

\begin{lemma}[Local finiteness]\label{B1.55}
Let $Z\subset\C$. Suppose every $z\in H$ has a punctured
neighbourhood disjoint from $Z$. For every $z\in\C$ and $r\in\R$
with $r<\Re z$, the set $Z\cap\overline B(z,r)$ is finite.
\end{lemma}

This is the closed-disc case of the Bolzano--Weierstrass theorem
\cite[Chapter~III, \S1.4, Theorem~7]{ahlfors-1979}.

Throughout, \((g,\Phi,Z)\) satisfies the pole data of
Definition~\ref{B1.33} with a constant \(C_0\ge0\); thus \(g\) is real,
Borel measurable, \(|g(u)|\le C_0e^{2u}\) for \(u\ge0\), and
\(\Phi(s)=\int_0^\infty g(u)e^{-su}\,du\) for \(\Re s>2\).
For \(\rho\in Z\) let \(c_\rho\neq0\) be the coefficient of
\((s-\rho)^{-1}\) in~\ref{B1.33}. Write \(H=\{s\in\C:\Re s>0\}\) and fix
\(0<\theta<1\). For \(A>0\), \(U>0\), \(v\in\R\) and \(s\in\C\) put
\[
 I_A(U,v)=\sqrt{\frac{A}{\pi U}}\int_0^\infty g(u)\,e^{-\theta u}
          e^{-A(u-v)^2/U}\,du,
 \qquad
 E_{U,v}(s)=\exp\Bigl((s-\theta)v+\frac{(s-\theta)^2U}{4A}\Bigr).
\]
The integral converges absolutely, its integrand being bounded by
\(C_0e^{2u-A(u-v)^2/U}\). For \(s=\sigma+it\) one has
\[
 |E_{U,v}(s)|=\exp\Bigl((\sigma-\theta)v+
   \frac{\bigl((\sigma-\theta)^2-t^2\bigr)U}{4A}\Bigr),
 \qquad E_{U,v}(\bar s)=\overline{E_{U,v}(s)} .\tag{G}
\]

\begin{lemma}[Basic properties of the pole data]\label{B1.54a}
\begin{enumerate}
\item \(|\Phi(s)|\le C_0\) for \(\Re s\ge3\), and \(\Re\rho\le2\) for
every \(\rho\in Z\).
\item Every compact \(K\subset H\) meets \(Z\) in a finite set.
\item If \(\rho\in Z\), then \(\bar\rho\in Z\) and
\(c_{\bar\rho}=\overline{c_\rho}\).
\end{enumerate}
\end{lemma}

\begin{proof}
(a) is Lemma~\ref{B1.54} with \(\sigma_0=3\).
(b) Cover \(K\) by finitely many closed discs
\(\overline B(z,r_z)\), \(z\in K\), with \(r_z=\tfrac12\min_K\Re<\Re z\),
and apply Lemma~\ref{B1.55} to each.

(c) Put \(Z^*=Z\cup\bar Z\). Since conjugation preserves \(H\), no point
of \(H\) is an accumulation point of \(Z^*\). By (b), \(Z^*\) is countable,
and \(\Omega=H\setminus Z^*\) is open and invariant under conjugation.
It is connected. Given \(p,q\in\Omega\), the paths
\([p,w]\cup[w,q]\), with \(w\) on the perpendicular bisector of
\([p,q]\) inside \(H\), lie in \(H\) by convexity. Two of them meet only
at \(p\) and \(q\), and there are uncountably many, so one of them misses
\(Z^*\).

The function \(\Psi(s)=\overline{\Phi(\bar s)}\) is holomorphic on
\(\Omega\). For \(\Re s>2\), since \(g\) is real,
\(\Psi(s)=\int_0^\infty g(u)e^{-su}\,du=\Phi(s)\). By the identity
theorem, \(\Phi=\Psi\) on \(\Omega\).

Let \(\rho\in Z\) and choose \(r>0\) such that
\(B(\rho,r)\setminus\{\rho\}\subset\Omega\) and
\(\Phi=c_\rho/(s-\rho)+h_\rho\) there. For
\(0<|s-\bar\rho|<r\) we get
\(\Phi(s)=\overline{c_\rho}/(s-\bar\rho)+\overline{h_\rho(\bar s)}\),
and the last term is bounded near \(\bar\rho\). Hence \(\Phi\) is
unbounded near \(\bar\rho\). Since \(\Phi\) is holomorphic on
\(H\setminus Z\), this forces \(\bar\rho\in Z\). Subtracting the
expansion at \(\bar\rho\) from this one, the function
\((c_{\bar\rho}-\overline{c_\rho})/(s-\bar\rho)\) is bounded near
\(\bar\rho\), so \(c_{\bar\rho}=\overline{c_\rho}\).
\end{proof}

\begin{lemma}[Gaussian representation]\label{B1.54b}
For \(A,U>0\) and \(v\in\R\), the integral below converges absolutely and
\[
 I_A(U,v)=\frac1{2\pi}\int_{\R}\Phi(3+it)\,E_{U,v}(3+it)\,dt .
\]
\end{lemma}

\begin{proof}
Put \(T=A/U\). For \(b\in\R\),
\(\int_\R e^{-t^2/(4T)}e^{ibt}\,dt=2\sqrt{\pi T}\,e^{-Tb^2}\).
For real \(c,y\) the exponent \((c+it)y+(c+it)^2/(4T)\) equals
\(cy+c^2/(4T)-t^2/(4T)+it\bigl(y+c/(2T)\bigr)\). Hence
\[
 \frac1{2\pi}\int_\R\exp\Bigl((c+it)y+\frac{(c+it)^2}{4T}\Bigr)dt
 =\sqrt{T/\pi}\,e^{-Ty^2}.
\]
Take \(c=3-\theta\), \(y=v-u\), \(s=3+it\). The exponent becomes
\((s-\theta)(v-u)+(s-\theta)^2U/(4A)\); multiplying by \(e^{-\theta u}\)
gives, for every real \(u\),
\[
 \sqrt{T/\pi}\,e^{-\theta u}e^{-T(u-v)^2}
 =\frac1{2\pi}\int_\R E_{U,v}(3+it)\,e^{-(3+it)u}\,dt .
\]
By (G), \(t\mapsto|E_{U,v}(3+it)|\) is integrable. Moreover
\(|g(u)e^{-(3+it)u}|\le C_0e^{-u}\) for \(u\ge0\), so
\((u,t)\mapsto g(u)E_{U,v}(3+it)e^{-(3+it)u}\) is integrable on
\([0,\infty)\times\R\). Multiply by \(g(u)\), integrate over \(u\ge0\),
apply Fubini's theorem \cite[Theorem~2.37]{folland-1999}, and use
\(\Phi(3+it)=\int_0^\infty g(u)e^{-(3+it)u}\,du\).
\end{proof}

\begin{lemma}[One-sided bound on a window]\label{B1.54c}
Let \(\delta>0\), \(d=\min(\delta,1)\), \(A\ge512/d^2\), \(U\ge1\), and
\(\eta>0\). Let \(R_U\) satisfy~(5) of Lemma~\ref{B1.29} on
\(I_U=[U,(1+\delta)U]\). If \(v\le2U\), \(v-U\ge dU/4\) and
\((1+d)U-v\ge dU/4\), then
\[
 I_A(U,v)\le\sqrt{\frac{A}{\pi U}}\,e^{\eta U}+\sqrt2\,C_0e^{-4U}.
\]
\end{lemma}

\begin{proof}
Put \(T=A/U\), \(W=[U,(1+d)U]\subset I_U\), and
\(k(u)=\sqrt{T/\pi}\,e^{-T(u-v)^2}\in(0,\sqrt{T/\pi}\,]\).

On \(W\) we have \(g(u)e^{-\theta u}k(u)\le\sqrt{T/\pi}\,R_U(u)e^{-\theta u}\).
Indeed, if \(g(u)\le0\) the left side is at most \(0\); otherwise
\(0<g(u)\le R_U(u)\). Since \(R_U\ge0\) on \(I_U\),
\[
 \int_W g(u)e^{-\theta u}k(u)\,du\le\sqrt{T/\pi}
 \int_{I_U}R_U(u)e^{-\theta u}\,du\le\sqrt{T/\pi}\,e^{\eta U}.
\]

Now let \(u\ge0\) with \(u\notin W\), and put \(r=|u-v|\); then
\(r\ge dU/4\). From \(u\le v+r\le 2U+r\) and \(\theta\ge0\) we get
\((2-\theta)u\le4U+2r\). Also
\(Tr^2/4\ge (A/U)(dU/4)^2/4=Ad^2U/64\ge8U\) and
\(Tr/4\ge Ad/16\ge32/d\ge2\), so \(Tr^2/4\ge2r\). Therefore
\[
 |g(u)|e^{-\theta u}e^{-Tr^2}\le C_0e^{4U+2r-Tr^2}
 \le C_0e^{-4U}e^{-Tr^2/2}.
\]
Since \(\int_\R\sqrt{T/\pi}\,e^{-T(u-v)^2/2}\,du=\sqrt2\), the part of
\(I_A(U,v)\) outside \(W\) is at most \(\sqrt2\,C_0e^{-4U}\).
Adding the two parts gives the claim.
\end{proof}

\begin{lemma}[Residue expansion]\label{B1.54d}
Let \(A\ge1\), \(R>0\) with \(R^2\ge9+32A\), and \(a\in(\theta/2,\theta)\).
Put \(Q=\{\sigma+it:\ a\le\sigma\le3,\ |t|\le R\}\) and assume
\(Z\cap\partial Q=\varnothing\). Then \(P=Z\cap Q\) is finite and lies in
the interior of \(Q\). Moreover there is \(K\ge0\) such that, with
\(\kappa=\min\{2,(\theta-a)/2\}\), for all \(U\ge1\) and
\(v\in[U,2U]\),
\[
 \Bigl|I_A(U,v)-\sum_{\rho\in P}c_\rho E_{U,v}(\rho)\Bigr|
 \le Ke^{-\kappa U}.
\]
\end{lemma}

\begin{proof}
\(P\) is finite by Lemma~\ref{B1.54a}(b). It lies in
\(\Re s\le2\) by (a), hence in the interior of \(Q\). Since \(Z\) is
closed in \(H\) and misses \(\partial Q\), \(\Phi\) is holomorphic near
\(\partial Q\) and on \(Q^\circ\setminus P\). At \(\rho\in P\) the
function \(\Phi E_{U,v}\) has a simple pole with residue
\(c_\rho E_{U,v}(\rho)\). The residue theorem
\cite[Chapter~4, \S5.1]{ahlfors-1979} gives
\[
 \int_{3-iR}^{3+iR}\Phi E_{U,v}\,ds
 =2\pi i\sum_{\rho\in P}c_\rho E_{U,v}(\rho)
 +\int_{a-iR}^{a+iR}+\int_{a+iR}^{3+iR}-\int_{a-iR}^{3-iR}\Phi E_{U,v}\,ds,
\]
all along straight segments. Let \(K_1=\max_{\partial Q}|\Phi|\).
Use (G) with \(U\le v\le2U\) and \(|\sigma-\theta|\le3\) on \(Q\).

\emph{Horizontal edges} (\(|t|=R\)). Here
\((\sigma-\theta)v\le6U\) and \((\sigma-\theta)^2\le9\), so
\(|E_{U,v}|\le\exp(6U+(9-R^2)U/(4A))\le e^{-2U}\). Each edge
contributes at most \(3K_1e^{-2U}\).

\emph{Left edge} (\(\sigma=a\)). Here \(0<\theta-a<1/2\) and \(A\ge1\),
so \((a-\theta)v+(a-\theta)^2U/(4A)\le-(\theta-a)U+(\theta-a)U/8\).
Hence \(|E_{U,v}|\le e^{-(\theta-a)U/2}\), and the edge contributes at
most \(2RK_1e^{-(\theta-a)U/2}\).

\emph{Right tail} (\(\sigma=3\), \(|t|\ge R\)). Here
\(|\Phi(3+it)|\le C_0\) and
\(|E_{U,v}|\le e^{-2U}\exp(-(t^2-R^2)U/(4A))\), because
\(9-t^2\le-32A-(t^2-R^2)\). From \(t^2-R^2\ge2R(|t|-R)\),
\(\int_{|t|\ge R}e^{-(t^2-R^2)U/(4A)}\,dt\le4A/(RU)\le4A/R\).

Split the integral of Lemma~\ref{B1.54b} at \(|t|=R\), and note
that \(\frac1{2\pi}\int_{-R}^R\Phi E_{U,v}(3+it)\,dt
=\frac1{2\pi i}\int_{3-iR}^{3+iR}\Phi E_{U,v}\,ds\). This gives the
claim with \(K=(6K_1+2RK_1+4AC_0/R)/(2\pi)\).
\end{proof}

\begin{proposition}[The window criterion]\label{B1.54e}
Under the hypotheses of Theorem~\ref{B1.31}, \(\Re\rho\le\theta\) for every \(\rho\in Z\).
\end{proposition}

\begin{proof}
Suppose \(\rho_0=\beta_0+i\gamma_0\in Z\) with \(\beta_0>\theta\).

\emph{Constants.} Put \(d=\min(\delta,1)\) and
\(A=\max\{512/d^2,\ \gamma_0^2/(\beta_0-\theta)\}\). Let
\(R_1=1+|\gamma_0|+\sqrt{9+32A}\). By
Lemma~\ref{B1.54a}(b), the set
\(Z\cap\{\theta/2\le\Re s\le3,\ |\Im s|\le R_1\}\) is finite. Choose
\(R\in[R_1-1,R_1]\) different from every \(|\Im\rho|\) there, and
\(a\in(\theta/2,\theta)\) different from every \(\Re\rho\) there. Then
\(R>|\gamma_0|\), \(R^2\ge9+32A\), and \(Z\cap\partial Q=\varnothing\),
because \(\Re\rho\le2<3\). Lemma~\ref{B1.54d} applies, and
\(\rho_0\in P\) because \(a<\theta<\beta_0\le2\).

\emph{Rates.} For \(\rho=\beta+i\omega\in P\) and \(\lambda\in\R\) put
\(L_\rho(\lambda)=(\beta-\theta)\lambda+((\beta-\theta)^2-\omega^2)/(4A)\).
By (G), for \(v=\lambda U+h\),
\[
 E_{U,v}(\rho)=e^{UL_\rho(\lambda)+(\beta-\theta)h}e^{i\Theta},\qquad
 \Theta=\omega(\lambda U+h)+(\beta-\theta)\omega U/(2A). \tag{R}
\]
Each \(L_\rho\) is affine in \(\lambda\), and \(L_\rho\equiv L_{\rho'}\)
if and only if \(\beta=\beta'\) and \(|\omega|=|\omega'|\). As \(P\) is
finite, we may fix \(\lambda\in(1+d/3,\,1+2d/3)\) at which no two
non-identical \(L_\rho\) coincide. Let \(M=\max_{\rho\in P}L_\rho(\lambda)\)
and let \(P^*\) be the set of maximisers. All \(\rho\in P^*\) share one
real part \(\beta^*\) and one value \(|\Im\rho|=\omega^*\ge0\).
By Lemma~\ref{B1.54a}(c) and the symmetry of \(Q\), either
\(P^*=\{\beta^*\}\), or \(P^*=\{\rho^*,\overline{\rho^*}\}\) with
\(\rho^*=\beta^*+i\omega^*\), \(\omega^*>0\) and
\(c_{\overline{\rho^*}}=\overline{c_{\rho^*}}\).

Since \(\lambda>1\) and \(A\ge\gamma_0^2/(\beta_0-\theta)\),
\[
 M\ge L_{\rho_0}(\lambda)\ge(\beta_0-\theta)-\gamma_0^2/(4A)
 \ge\tfrac34(\beta_0-\theta)>0.
\]
For \(\beta\ge a\) we have
\(\lambda+(\beta-\theta)/(4A)>1-1/(8A)>0\). Hence \(\beta\le\theta\)
implies
\(L_\rho(\lambda)\le(\beta-\theta)\bigl(\lambda+(\beta-\theta)/(4A)\bigr)\le0<M\),
and so \(\beta^*>\theta\). Fix \(\varepsilon>0\) with
\(L_\rho(\lambda)\le M-\varepsilon\) for every \(\rho\in P\setminus P^*\).

\emph{Moving the centre.} Let \(H_*=0\) if \(\omega^*=0\) and
\(H_*=2\pi/\omega^*\) otherwise. For \(U>0\), let \(h_U=0\) if
\(\omega^*=0\). Otherwise let \(h_U\in[0,H_*)\) be the unique number with
\(\arg c_{\rho^*}+\omega^*(\lambda U+h_U)+(\beta^*-\theta)\omega^*U/(2A)
\in2\pi\mathbf Z\). Put \(v_U=\lambda U+h_U\).

\emph{Lower bound.} By (R), the choice of \(h_U\), and (G),
\[
 \sum_{\rho\in P^*}c_\rho E_{U,v_U}(\rho)=
 \begin{cases}
 c_{\beta^*}e^{MU},&\omega^*=0,\\
 2\Re\bigl(c_{\rho^*}E_{U,v_U}(\rho^*)\bigr)
 =2|c_{\rho^*}|e^{MU+(\beta^*-\theta)h_U},&\omega^*>0.
 \end{cases}
\]
This is a real number \(\ge c_*e^{MU}\), where \(c_*=c_{\beta^*}>0\)
(a real pole has positive residue, by~\ref{B1.33}) or
\(c_*=2|c_{\rho^*}|>0\). For \(\rho\in P\setminus P^*\), (R) gives
\(|c_\rho E_{U,v_U}(\rho)|\le|c_\rho|e^{2H_*}e^{(M-\varepsilon)U}\).

Let \(U\ge\max\{1,3H_*\}\). Then \(U\le v_U\le\tfrac53U+H_*\le2U\), so
Lemma~\ref{B1.54d} applies. Since \(I_A(U,v_U)\) and the sum over
\(P^*\) are real,
\[
 I_A(U,v_U)\ge c_*e^{MU}-e^{2H_*}\Bigl(\sum_{\rho\in P}|c_\rho|\Bigr)
 e^{(M-\varepsilon)U}-Ke^{-\kappa U}.
\]
Hence there is \(U_1\) with \(I_A(U,v_U)\ge\tfrac12c_*e^{MU}\) for every
\(U\ge U_1\).

\emph{Upper bound.} Let \(U\ge\max\{1,3H_*,12H_*/d\}\). Then
\(v_U\le2U\), \(v_U-U\ge(\lambda-1)U>dU/4\), and
\((1+d)U-v_U\ge(1+d-\lambda)U-H_*>dU/3-H_*\ge dU/4\).

\emph{Contradiction.} Put \(\eta=M/2\). Choose
\(U_0\ge\max\{U_1,1,3H_*,12H_*/d\}\) such that
\(\tfrac12c_*e^{MU}>\sqrt{A/\pi}\,e^{MU/2}+\sqrt2\,C_0\) for all
\(U\ge U_0\). The hypothesis \(\mathrm{WH}_U(g;\theta,\delta)\) gives
\(U\ge U_0\) and \(R_U\) satisfying~(5). By Lemma~\ref{B1.54c}
and the lower bound,
\[
 \tfrac12c_*e^{MU}\le I_A(U,v_U)\le
 \sqrt{\tfrac{A}{\pi U}}\,e^{MU/2}+\sqrt2\,C_0e^{-4U}
 <\tfrac12c_*e^{MU},
\]
a contradiction. Hence \(\Re\rho\le\theta\) for every \(\rho\in Z\).
\end{proof}

\begin{remark}
The proof uses no growth bound for \(\Phi\) in vertical strips: only \(|\Phi|\le C_0\) on \(\Re s=3\) and boundedness on the fixed compact set \(\partial Q\). Positivity of residues is used only at the single real pole \(\beta^*\).
\end{remark}

\part{The arithmetic side: prime sums and auxiliary primes}

% ---- p2c1 ----
\chapter{Counting primes class by class}\label{ch:p2c1}

Fix a frame with finite abelian group $A$ and distinguished element
$a_0$, and put $A_0=\langle a_0\rangle$. The arithmetic data in
Alp\"oge's archive \cite{alpoge-2026} attach a Frobenius class and an
inertia subgroup to each rational prime. Character orthogonality
expresses the weighted prime-power count as a finite sum. After replacing
the weights on thick positions by a common weight, we average over
$A_0$. This bounds the real part of a weighted sum over the surviving
characters.

\section{The host and its local classes}

\begin{definition}[Host data used in the prime count]\label{B2.01}
For a frame $\phi$, the host data $H$ used in the prime count consist of the auxiliary primes and modulus below, together with a cyclotomic subfield and its Artin map satisfying the following conditions.
The frame has $e\in\N$ and a surjective homomorphism
$h:A\to(\Z/2^{e+2}\Z)^\times/\{\pm1\}$ whose restriction to
$A_0$ is injective. Write
$\mathcal G_\phi=(\{\pm1\}\times\{\pm1\})\times A$,
with the last factor written multiplicatively in host coordinates.
The host supplies primes $q,d$ and a positive integer $Q_S$ such that
\[
 \begin{gathered}
 e\ge2,\qquad q\equiv3\pmod8,\qquad d\equiv7\pmod{16},
 \qquad \left(\frac{-d}{q}\right)=1,\\
 2^{e+2}\mid Q_S,\qquad (Q_S,qd)=1,\qquad
 \left(\frac{-d}{p}\right)=1
 \quad\text{for every odd prime }p\mid Q_S.
 \end{gathered}
\]
Set $\mathfrak m=Q_Sqd$. Choose a primitive $\mathfrak m$th root of unity $\zeta_{\mathfrak m}$. The field data include a CM field $E$ Galois over $\Q$, an embedding $E\hookrightarrow\Q(\zeta_{\mathfrak m})$,
an isomorphism
$\operatorname{Gal}(E/\Q)\simeq\mathcal G_\phi$, and an element
$s_d\in E$ with
\[
 s_d^2=-d,\qquad
 \sigma(s_d)=\sigma_d s_d
 \quad\text{if }\sigma\longleftrightarrow((\sigma_q,\sigma_d),a).
\]
Write $\operatorname{Art}:(\Z/\mathfrak m\Z)^\times\to\mathcal G_\phi$ for the arithmetic ideal Artin homomorphism. At $p\nmid\mathfrak m$ this is the restriction of $\operatorname{Art}_{\Q}(t_p)^{-1}$, where the idèle $t_p$ has component $p$ at $p$ and $1$ elsewhere.
If $\operatorname{Art}(u)=((\sigma_q,\sigma_d),a)$, then $h(a)=[u\bmod 2^{e+2}]$.
It is compatible with restriction of the cyclotomic automorphism
$\zeta_{\mathfrak m}\mapsto\zeta_{\mathfrak m}^{\,u}$ and sends $-1$
to $((-1,-1),1)$. It also satisfies
\[
 \begin{split}
 u\equiv1\pmod{Q_Sq}&\ \Longrightarrow\
 \operatorname{Art}(u)\in\{1,\iota_d\},
 \qquad \iota_d=((1,-1),1),\\
 u\equiv1\pmod{Q_Sd}&\ \Longrightarrow\
 \operatorname{Art}(u)^2=1.
 \end{split}
\]
The extension $E/\Q$ has degree $4|A|$. The Artin map is surjective,
with $d$-coordinate $(u/d)$ at $u$. For each prime $p\nmid\mathfrak m$, arithmetic Frobenius corresponds to $\operatorname{Art}(p)$ under the Galois isomorphism. More generally, let
$\mathfrak m^{(p)}$ be the largest divisor of $\mathfrak m$ prime
to $p$. The inertia subgroup at $p$ is the Artin image of the
units congruent to $1$ modulo $\mathfrak m^{(p)}$. The Frobenius
coset is the Artin image of the units congruent to $p$ modulo
$\mathfrak m^{(p)}$.
\end{definition}

For the sums we use additive coordinates. Put
\[
 E_0=(\Z/2\Z)^2,\qquad G=E_0\times A,\qquad
 \Phi=\{((v_d,v_q),h):v_d=0\},\qquad f=|\Phi|=2|A|.
\]
Write $\epsilon_d=((1,0),0)$, $c=((1,1),0)$ and
$\tau(t)=((0,0),t)$.
The character group $\widehat G$ consists of homomorphisms
$G\to\C^\times$. Write $\mathbf1$ for the trivial character.
Let
$\chi_d((v_d,v_q),h)=(-1)^{v_d}$ and
$\chi_q((v_d,v_q),h)=(-1)^{v_q}$.

\begin{lemma}[Changing coordinates and characters]\label{B2.02}
For every frame $\phi$ there is a group isomorphism
\[
 T:\mathcal G_\phi\longrightarrow G,\qquad
 T((\sigma_q,\sigma_d),a)
 =((\overline{\sigma_d},\overline{\sigma_q}),a),
 \qquad \overline{1}=0,\quad\overline{-1}=1.
\]
Here the last coordinate is converted to additive notation.
Put $c_\phi=((-1,-1),1)$. For a host $H$, its Galois isomorphism sends complex conjugation to $c_\phi$.
We have $T(\iota_d)=\epsilon_d$ and $T(c_\phi)=c$.
For $t\in A$, the host element $\tau_1(t)=((1,1),t)$ satisfies
$T(\tau_1(t))=\tau(t)$.
The kernel of the $d$-sign character is $\Psi_d=\{((\sigma_q,1),a):\sigma_q\in\{\pm1\},\ a\in A\}$. The map $T$ carries $\Psi_d$ onto $\Phi$.
Composition with $T$ gives a bijection
$\widehat G\to\widehat{\mathcal G_\phi}$,
$\alpha\mapsto\theta_\alpha=\alpha\circ T$.
Pullback sends the trivial character to the trivial character.
At $((\sigma_q,\sigma_d),a)$, the pullbacks of $\chi_d$ and
$\chi_q$ take the values $\sigma_d$ and $\sigma_q$, respectively.
Then $\theta_\alpha$ is trivial on $\tau_1(A_0)$ exactly when
$\alpha$ is trivial on $\tau(A_0)$. The set $X_+$ of such host characters with value $1$ at $\iota_d$ therefore corresponds to the $\alpha$ trivial on $\tau(A_0)$ with $\alpha(\epsilon_d)=1$.
\end{lemma}

\begin{proof}[Outline]
Converting \(\{\pm 1\}\) to \(\mathbb{Z}/2\mathbb{Z}\), swapping signs, and writing \(A\) additively defines \(T\), with element images and \(\Psi_d\mapsto\Phi\) checked directly. Pullback gives the character bijection and matches the \(\tau_1(A_0)\) and \(\iota_d\) conditions. The complete proof is in \S\ref{prf:B2.02}.
\end{proof}

Let $I_p\subseteq G$ be the inertia subgroup, and choose a Frobenius
representative $F_p\in G$. At a ramified prime its class is
$F_p+I_p$.

\begin{lemma}[The three local cases]\label{B2.03}
For any host, let
$S=\{q\}\cup\{p:p\text{ is prime and }p\mid Q_S\}$.
Then $2,q\in S$, $d\notin S$, and
\[
 \begin{array}{ll}
 p\text{ prime},\ p\notin S,\ p\ne d:& I_p=\{0\},\\
 p\in S:& I_p\subseteq\Phi,\quad F_p\in\Phi,\\
 p=d:& I_d=\{0,\epsilon_d\}.
 \end{array}
\]
Consequently every coset $iF_d+I_d$ contains exactly one element
outside $\Phi$.
\end{lemma}

\begin{proof}[Outline]
Host conditions give trivial inertia outside \(S\cup\{d\}\) and, with reciprocity, sign \(1\) for inertia and Frobenius in \(S\). At \(d\), the Chinese remainder theorem and Lemma~\ref{B2.02} give \(I_d=\{0,\epsilon_d\}\), with \(\epsilon_d\) exchanging the halves. The complete proof is in \S\ref{prf:B2.03}.
\end{proof}

\section{The exact prime-power identity}

Let $A$ be any finite abelian group. An arithmetic datum $D$ consists of $q,d\in\N$ and, for each
$n\in\N$, a subgroup $I_n\le G$ and an element $F_n\in G$.
The local hypotheses mean that $d$ is prime and there is a finite
$S\subseteq\N$ with $2,q\in S$, $d\notin S$, and the three
local conditions displayed in Lemma~\ref{B2.03}.
Write $[z]_+=\max(z,0)$ and use $a/0=0$ throughout.

\begin{definition}[Local values, sums, and centring]\label{B2.04}
For $\alpha\in\widehat G$, a prime $p$, and $i\ge1$, set
\[
 V_\alpha(p^i)=\frac1{|I_p|}
 \sum_{h\in I_p}\alpha(iF_p+h),\qquad
 v_\alpha(n)=\prod_{p\mid n}V_\alpha(p^{v_p(n)})\quad(n\ge1).
\]
At $n=1$ the product is empty. Take the empty product at $n=0$ as well, so $v_\alpha(0)=v_\alpha(1)=1$.
Let $\ell:\widehat G\to\C$ be arbitrary. For every $y\in\R$ define
\[
 G_\alpha(y)=
 \begin{cases}
 G_z(y),&\alpha=\mathbf1,\\[2pt]
 \displaystyle\sum_{1\le n\le y}\Lambda(n)v_\alpha(n)(y/n-1)+y\ell_\alpha,
 &\alpha\ne\mathbf1.
 \end{cases}
\]
Here $G_z$ and the total real logarithm are as in
Definition~\ref{B1.26}. A sum over $1\le n\le y$ is empty for $y<1$.
For $t\in A$, define the partner map $r_t(g)=c+g+\tau(t)$.
It maps $G\setminus\Phi$ bijectively onto $\Phi$.
Given $a\in\R$, $w:G\to\R$, and $N\in\N$, put
\[
\begin{split}
 \ASK_N(a,t,w)={}&
 \sum_{\substack{p\le N,\ p\ {\rm prime}\\
 p\ne2,q,d,\ F_p\notin\Phi}}
 2\log p
 \sum_{\substack{1\le i\le N\\iF_p\notin\Phi}}
 [a/p^i-w(r_t(iF_p))]_+\\
 &+\log d\sum_{i=1}^N
 \sum_{\substack{g\in iF_d+I_d\\g\notin\Phi}}
 [a/d^i-w(r_t(g))]_+ .
\end{split}
\]
For the same $a,t,w$, define the main term
\[
 \Main(a,w)
 =a\left(\log a-1-\gamma+\ell_{\chi_d}
 -\frac1f\sum_{\rho\in\Phi}\log w(\rho)\right)
\]
and the centred character sum
\[
 \Delta(a,t,w)
 =\frac1f\sum_{g\notin\Phi}\sum_{\alpha\in\widehat G}
 \overline{\alpha(g)}\,w(r_t(g))\,
 G_\alpha\left(\frac{a}{w(r_t(g))}\right).
\]
\end{definition}

The value $V_\alpha(p^i)$ is unchanged when $F_p$ is replaced
by another representative of its inertia coset. Character
orthogonality on $I_p$ also gives
\[
 V_\alpha(p^i)=
 \begin{cases}
 \alpha(iF_p),&\alpha|_{I_p}=\mathbf1,\\
 0,&\alpha|_{I_p}\ne\mathbf1.
 \end{cases}
\]
In particular $|v_\alpha(n)|\le1$.

\begin{lemma}[Orthogonality on an inertia coset]\label{B2.05}
For every arithmetic datum $D$, every prime $p$, every integer
$i\ge1$, and every $g\in G$,
\[
 \frac1f\sum_{\alpha\in\widehat G}
 \overline{\alpha(g)}\,v_\alpha(p^i)
 =\frac{2}{|I_p|}\,\mathbf1_{\,g\in iF_p+I_p}.
\]
Under the local conditions of Lemma~\ref{B2.03}, the coefficient
on the indicated coset is $2$ at an unramified prime and $1$ at $d$.

\end{lemma}

\begin{proof}
The product defining $v_\alpha(p^i)$ has the single factor $V_\alpha(p^i)$. Since $|\alpha(g)|=1$, the left side is
\[
 \frac1{f|I_p|}\sum_{h\in I_p}
 \sum_{\alpha\in\widehat G}\alpha(iF_p+h-g).
\]
The inner sum is $|G|$ when $iF_p+h=g$ and is $0$ otherwise.
There is one such $h$ precisely when $g\in iF_p+I_p$.
Finally $|G|=2f$.
\end{proof}

\begin{lemma}[Expansion of one layer]\label{B2.06}
Assume the local conditions of Lemma~\ref{B2.03}.
Let $a\in\R$, $t\in A$, and $w:G\to\R$ satisfy
$w(\rho)>0$ for $\rho\in\Phi$. Let $N\in\N$ satisfy
\[
 a/w(\rho)\le N\qquad(\rho\in\Phi).
\]
Then
\[
 \ASK_N(a,t,w)
 =\frac1f\sum_{g\notin\Phi}\sum_{\alpha\in\widehat G}
 \overline{\alpha(g)}
 \sum_{n=1}^N\Lambda(n)v_\alpha(n)
 [a/n-w(r_t(g))]_+ .
\]
The equality is in $\C$, with the real left side embedded in $\C$.
\end{lemma}

\begin{proof}[Outline]
Expand into prime powers and use Lemma~\ref{B2.05} to sum over the parts of inertia cosets outside \(\Phi\) for \(p^i\le N\). Primes in \(S\) vanish, \(d\) has weight \(\log d\), and other primes contribute only at odd powers with \(F_p\notin\Phi\). The complete proof is in \S\ref{prf:B2.06}.
\end{proof}

\begin{theorem}[Exact balance for one layer]\label{B2.07}
Under the hypotheses of Lemma~\ref{B2.06}, assume further that
$a>0$. For every choice of $\ell:\widehat G\to\C$,
\[
 \ASK_N(a,t,w)-\Main(a,w)=\Delta(a,t,w).
\]
\end{theorem}

\begin{proof}[Outline]
Write each weighted sum as \(WG_\alpha(a/W)\) plus a remainder in Lemma~\ref{B2.06}. On \(\epsilon_d+\Phi\), orthogonality leaves the trivial and \(\chi_d\) terms, and reindexing the trivial term by \(r_t\) gives \(\operatorname{MAIN}(a,w)\). The complete proof is in \S\ref{prf:B2.07}.
\end{proof}

\section{Averaging the thick contribution}

Let $A_0$ be any subgroup of $A$. Define
\[
 P=\{\alpha\in\widehat G:\alpha(\tau(s))=1
 \text{ for every }s\in A_0\},
 \qquad e_1(t)=((1,0),t).
\]
For a finite subset $T_2\subseteq A$, put
$\delta=|T_2|/|A|$.
For $b,y\in\R$ and $u,t\in A$, define the uniform thick sum
\[
 \Delta_{\rm thick}(u,t;y)
 =\frac1f\sum_{s\in T_2}\sum_{\alpha\in\widehat G}
 \alpha(-e_1(s+u-t))\,bG_\alpha(y).
\]

\begin{lemma}[Averaging over the subgroup]\label{B2.08}
Let $s\in A_0$, $t\in A$, and $\alpha\in\widehat G$. Then
\[
 \sum_{u\in A_0}\alpha(-e_1(s+u-t))
 =\begin{cases}
 |A_0|\alpha(e_1(t)),&\alpha\in P,\\
 0,&\alpha\notin P.
 \end{cases}
\]
Consequently, if $T_2\subseteq A_0$, then for all $b,y\in\R$
and $t\in A$,
\[
 \sum_{u\in A_0}\Delta_{\rm thick}(u,t;y)
 =|A_0|\,b\,\frac{\delta}{2}
 \sum_{\alpha\in P}\alpha(e_1(t))G_\alpha(y).
\]
\end{lemma}

\begin{proof}[Outline]
Using \(-e_1(s+u-t)=e_1(t)+\tau(-s)+\tau(-u)\) factors out \(\alpha(e_1(t))\alpha(\tau(-s))\) from \(\sum_{u\in A_0}\alpha(\tau(-u))\), equal to \(|A_0|\) on \(P\) and \(0\) off \(P\). Summing \(bG_\alpha(y)/f\) times this over \(s\in T_2,\alpha\in\widehat{G}\) gives the sum. The complete proof is in \S\ref{prf:B2.08}.
\end{proof}

For $u\in A$, the thick positions form
$\mathcal T_u=\{((0,1),s+u):s\in T_2\}\subseteq\Phi$.
Fix a scale $X'>0$ and an error bound $\theta\ge0$. The weight comparison uses the estimate
\begin{equation}\label{eq:b2-thin}
 |G_\alpha(y)|\le y\theta
 \qquad(\alpha\in\widehat G,\ y\ge X'/3),
\end{equation}
and a bound for changing a positive weight,
\begin{equation}\label{eq:b2-lip}
 \left|vG_\alpha(a/v)-v'G_\alpha(a/v')\right|
 \le \frac{3a}{\min(v,v')}|v-v'|
 \qquad(a,v,v'>0,\ \alpha\in\widehat G).
\end{equation}


\begin{lemma}[Replacing the weights on thick positions]\label{B2.09}
Let $T_2\subseteq A$ and $u,t\in A$. Suppose
\[
 \begin{gathered}
 x\ge3,\quad b>0,\quad b'>0,\quad \eta\ge0,\quad
 X'>0,\quad\theta\ge0,\qquad b'X'\le xb,\\
 |a/b-x|\le1.
 \end{gathered}
\]
Assume \eqref{eq:b2-thin} and \eqref{eq:b2-lip}, and suppose
$w:G\to\R$ satisfies
\[
 \begin{array}{ll}
 \rho\in\Phi:&w(\rho)>0,\\
 \rho\in\mathcal T_u:&b\le w(\rho)\le b+\eta+1,\\
 \rho\in\Phi\setminus\mathcal T_u:&w(\rho)\le2b'.
 \end{array}
\]
Then
\[
 \Re\Delta_{\rm thick}(u,t;a/b)
 \le\Re\Delta(a,t,w)
 +2f\left(3(x+1)(\eta+1)+(x+1)b\theta\right).
\]
\end{lemma}

\begin{proof}[Outline]
Use \(r_t^{-1}\) to locate thick positions and replace their weights by \(b\) using \eqref{eq:b2-lip}. Use \(a/w(\rho)\ge X'/3\) and \eqref{eq:b2-thin} to bound discarded terms, then sum absolute errors over positions and characters, normalize, and take real parts. The complete proof is in \S\ref{prf:B2.09}.
\end{proof}

\begin{theorem}[From bounds on translates to an average bound]\label{B2.10}
Assume the local conditions of Lemma~\ref{B2.03}.
Let $A_0\subseteq A$ be a subgroup, let $TT,T_2\subseteq A$,
and let $m,N\in\N$. Choose $a:A\to\R$ and write $a_t=a(t)$.
Choose weights $w_{u,j}:G\to\R$ for $u\in A$ and $0\le j<m$.
Suppose $T_2\subseteq A_0$, $x\ge3$, $b,b',X'>0$,
$\eta,\theta\ge0$, and $b'X'\le xb$.
Assume $|a_t/b-x|\le1$ for every $t\in TT$.
For each $u\in A_0$ and $0\le j<m$, assume the position bounds
of Lemma~\ref{B2.09} for $w_{u,j}$, and require
$a_t/w_{u,j}(\rho)\le N$ for every $t\in TT$ and $\rho\in\Phi$.
Assume \eqref{eq:b2-thin} and the two cost inequalities
$2f(x+1)\theta\le1/x$ and $6f(x+1)(\eta+1)\le b/x$.
Define
\[
 \operatorname{gap}(u)=
 \sum_{t\in TT}\sum_{j=0}^{m-1}
 \left(\ASK_N(a_t,t,w_{u,j})-\Re\Main(a_t,w_{u,j})\right),
 \qquad Q=|TT|mb .
\]
For any $\mathscr E,Y_*\in\R$, if
$\operatorname{gap}(u)\le\mathscr E$ for every $u\in A_0$ and
\[
 |TT|Y_*
 \le \sum_{t\in TT}\Re
 \sum_{\alpha\in P}
 \alpha(e_1(t))G_\alpha(a_t/b),
\]
then
\[
 Q\,\frac{\delta}{2}\,Y_*
 \le\mathscr E+\frac{2Q}{x}.
\]
\end{theorem}

\begin{proof}[Outline]
Use \(\Psi(z)\le3z\) \cite[Theorem~12]{rosser-schoenfeld-1962} to prove \eqref{eq:b2-lip} and apply \ref{B2.07} and \ref{B2.09} layerwise with the cost bounds. Sum using the gap bound, average over \(A_0\) using \ref{B2.08}, and apply the \(Y_*\) lower bound. The complete proof is in \S\ref{prf:B2.10}.
\end{proof}


% ---- p2c2 ----
\chapter{Reading off the target sum}\label{ch:p2c2}

Use the arithmetic data of \ref{B2.04} on \(G=E_0\times A\). For \(\alpha\in\widehat G\) and \(n\ge1\), write \(V_\alpha(n)=\prod_{p\mid n}V_\alpha(p^{v_p(n)})\) for the product of the local inertia averages defined there. Write \(1\) for the trivial character, so \(G_1=G_z\). For the subgroup \(A_0\le A\), let \(P\) be the set of characters trivial on \(\{0\}\times A_0\), as in \ref{B2.08}, and put \(e_1(t)=\epsilon_d+(0,t)\). For \(K_p\subseteq\widehat G\) and \(c\in\R\), set \(w_c(\alpha)=1\) when \(\alpha|_{\{0\}\times A}=1\), and \(w_c(\alpha)=c\) otherwise. Define
\[
 G_m^A(y)=\sum_{\alpha\in K_p}w_c(\alpha)\Re G_\alpha(y),
 \qquad
 \Sigma_A(y)=\sum_{\alpha\in P\setminus K_p}|G_\alpha(y)|.
\]

\section{A common point of evaluation}

We say that the prime sums satisfy the increment bounds with constant \(C_d\ge0\) if, for all \(y,z\ge1\),
\begin{align}
 |G_\alpha(y)-G_\alpha(z)|
 &\le |y-z|\bigl(\log\max(y,z)+C_d+|\ell_\alpha|\bigr)
 &&(\alpha\ne1),\label{eq:p2c2-deriv}\\
 |G_z(y)-G_z(z)|&\le C_d|y-z|.\label{eq:p2c2-zeta-lip}
\end{align}

\begin{lemma}[Moving one evaluation point]\label{B2.11}
Assume the increment bounds \eqref{eq:p2c2-deriv} and \eqref{eq:p2c2-zeta-lip} with \(C_d\ge0\). For \(x\ge3\), \(|y-x|\le1\), and \(\alpha\ne1\),
\[
 |G_\alpha(y)-G_\alpha(x)|\le\log(x+1)+C_d+|\ell_\alpha|,
 \qquad
 |G_\alpha(x)|\le x(\log x+C_d+|\ell_\alpha|).
\]
For the trivial character, \(|G_z(y)-G_z(x)|\le C_d\).
\end{lemma}
\begin{proof}
Since \(y\ge2\) and \(\max(y,x)\le x+1\), \eqref{eq:p2c2-deriv} gives the first bound. By \ref{B2.04}, \(G_\alpha(1)=\ell_\alpha\) for \(\alpha\ne1\). Thus \eqref{eq:p2c2-deriv} at \(x,1\) gives
\[
 |G_\alpha(x)|\le(x-1)(\log x+C_d+|\ell_\alpha|)+|\ell_\alpha|
 \le x(\log x+C_d+|\ell_\alpha|).
\]
The trivial-character bound follows from \eqref{eq:p2c2-zeta-lip}.
\end{proof}

\begin{lemma}[Comparison for each character]\label{B2.12}
Let \(TT\subseteq A\), \(n_T=|TT|\), \(x\ge3\), and \(|y_t-x|\le1\) for \(t\in TT\). Assume the increment bounds \eqref{eq:p2c2-deriv} and \eqref{eq:p2c2-zeta-lip} with \(C_d\ge0\). Let \(K_p\subseteq\widehat G\), \(c,L,R\in\R\), and \(\rho,K_o\ge0\). Suppose \(1\in K_p\), \(\alpha(\epsilon_d)=1\) for \(\alpha\in K_p\), and
\begin{align*}
 \left|\sum_{t\in TT}\alpha(0,t)-cn_T\right|&\le\rho n_T
 &&(\alpha\in K_p,\ \alpha|_{\{0\}\times A}\ne1),\\
 |\ell_\alpha|&\le K_o&&(\alpha\in K_p,\ \alpha\ne1),\\
 |\ell_\alpha|&\le L&&(\alpha\in P\setminus K_p).
\end{align*}
Require \(C_d\le R\) and
\[
 \begin{aligned}
 \log(x+1)+C_d+K_o+\rho x(\log x+C_d+K_o)&\le R,\\
 \log(x+1)+C_d+L&\le R.
 \end{aligned}
\]
Set \(v_\alpha=w_c(\alpha)\Re G_\alpha(x)\) for \(\alpha\in K_p\), and \(v_\alpha=-|G_\alpha(x)|\) otherwise. Then
\[
 n_T(v_\alpha-R)\le
 \sum_{t\in TT}\Re\bigl(\alpha(e_1(t))G_\alpha(y_t)\bigr)
 \qquad(\alpha\in P).
\]
If \(K_p\subseteq P\), then \(\sum_{\alpha\in P}v_\alpha=G_m^A(x)-\Sigma_A(x)\).
\end{lemma}
\begin{proof}[Outline]
Use \ref{B2.11} to sum the trivial bounds, control point shifts and character-sum errors for retained terms, and bound the remaining real parts using modulus-one values. If \(K_p\subseteq P\), the terms sum to \(G_m^A(x)-\Sigma_A(x)\). The complete proof is in \S\ref{prf:B2.12}.
\end{proof}

\section{The geometric bound and the arithmetic comparison}

Write \(n=|A|\), \(n_2=|T_2|\), \(Q=n_Tmb\), and \(\delta_2=n_2/n\), with \(\operatorname{gap}\) as in \ref{B2.10}. For rank \(k>0\) and top-weight scale \(a_E\ge0\), put
\begin{equation}\label{eq:p2c2-bracket}
 \mathcal B=n_Ta_E\left(1+\frac{m^2}{k}\right)
 +m(2nmbx^\varepsilon)
 +m\frac{n_2kbx^\varepsilon}{n(nx^{-\varepsilon})}.
\end{equation}
Write \(\ASK_D(N,a,t,w)\) for the sum \(\ASK_N(a,t,w)\) of \ref{B2.04} with datum \(D\). Write \(\ASK_{\mathscr D}(u)\) for the prime sum attached to translation \(u\) in \(\mathscr D\), and \(\Main_{\mathscr D}(u)\) for the geometric main term of \ref{B5.34}.

\begin{theorem}[The bound at each translation]\label{B2.13}
Let \(\mathscr D\) be a datum of \ref{B5.43}, with finite translation and top sets \(U,\mathcal T\), finite nonempty word sets \(W(u)\) for \(u\in U\), and \(m,f\in\N\). Fix arithmetic data \(D,\ell\), a cutoff \(N\in\N\), sets \(TT,T_2\subseteq A\), and weights \(a_t=a(t)\) and \(w_{s,j}:G\to\R\) for \(t,s\in A\) and \(0\le j<m\). Choose maps
\[
 \operatorname{tr}:U\to A,\qquad e_T:\mathcal T\simeq TT,\qquad
 e_P:\{r\in\N:r<f\}\simeq\Phi .
\]
For \(\tau\in\mathcal T\), require the geometric top weight at \(\tau\) to equal \(a_{e_T(\tau)}\). Write \(\omega_{u,j,r}\) for the geometric weight at translation \(u\), layer \(j\), and index \(r\), and require
\[
 \omega_{u,j,r}=w_{\operatorname{tr}(u),j}(e_P(r))
 \qquad(u\in U,\ 0\le j<m,\ 0\le r<f).
\]
In the summand of the geometric main term of \ref{B5.34}, require the two constants following \(\log a_t-1\) to be \(-\gamma\) and \(+\Re\ell_{\chi_d}\), respectively, where \(\gamma\) is Euler's constant and \(\chi_d\) is defined in \ref{B2.02}. Assume
\[
 \ASK_{\mathscr D}(u)=
 \sum_{t\in TT}\sum_{j=0}^{m-1}\ASK_D(N,a_t,t,w_{\operatorname{tr}(u),j})
 \qquad(u\in U).
\]
Let \(\mathcal C\) be the chain constants and \(X'',e_K,K\in\R\). Assume \(\mathrm{H}1\) through \(\mathrm{H}11\) of \ref{B5.43} and the window hypotheses of \ref{B5.45} through its trace bound. Set \(C_{\rm err}=e_{\rm Div}+e_{\rm Loss}+3\). Let \(x,\varepsilon\) be the window parameters of \(\mathscr D\), \(b\) its thick weight scale, \(k\) its starting rank, and \(a_E\) its top-weight scale, denoted by \(a\) in \ref{B5.45}. Suppose
\[
 \begin{gathered}
 B_0\le2nmbx^\varepsilon,\qquad
 N_c\le\frac{n_2kbx^\varepsilon}{n(nx^{-\varepsilon})},\\
 \log x\le x^\varepsilon,\qquad
 200(1+K)^2(1+c_{HH})(1+e_K)\le x^\varepsilon,
 \end{gathered}
\]
where \(c_{HH}\) is the height exponent in the chain constants \(\mathcal C\). Then for every \(u\in U\),
\[
 \operatorname{gap}(\operatorname{tr}(u))
 =\ASK_{\mathscr D}(u)-\Main_{\mathscr D}(u)
 \le x^{C_{\rm err}\varepsilon}\mathcal B.
\]
\end{theorem}
\begin{proof}[Outline]
By \ref{B5.43}, the difference is at most \(E_8(u)\) of \ref{B5.42}, bounded by \(x^{C_{\rm err}\varepsilon}\mathcal{B}\) via \ref{B5.45} and the \(B_0,N_c\) bounds. Reindexing via \ref{B5.46} matches the arithmetic main sum, giving \(\operatorname{gap}\) by \ref{B2.10}. The complete proof is in \S\ref{prf:B2.13}.
\end{proof}

For the family of \ref{B5.32a}, take \(\varepsilon>0\) sufficiently small and \(x\) above the common threshold. Its realisation in \ref{B5.33} uses \(c_{HH}=C_H\) and supplies \(\mathrm H9\) of \ref{B5.43}. The block Gram estimates of \ref{B5.31a} pass to the word family by the isometries \(J_v\) of \ref{B5.31}. For each \(u\in U\), let \(\eta(u)\) and \(e(u)\) be the frame parameters of \(W(u)\). Apply \ref{B5.32} at the height \(\mathcal H=x^{c_{HH}}b\) of \ref{B5.33} to obtain
\[
 (\eta(u)+4\sqrt{e(u)})\mathcal H
 \le(\bar\eta+4\sqrt{\bar e})\mathcal H\le6b.
\]

For \(M',K_o,C_E,c\in\R\), the character conditions on \(K_p\) and \(TT\) are
\begin{equation}\label{eq:p2c2-chars}
 \begin{gathered}
 1\in K_p,\quad K_p\subseteq P,\quad
 \alpha(\epsilon_d)=1\ (\alpha\in K_p),\quad |P|\le M',\\
 \left|\sum_{t\in TT}\alpha(0,t)-cn_T\right|
 \le M'x^{-1+\varepsilon}n_T
 \quad(\alpha\in K_p,\ \alpha|_{\{0\}\times A}\ne1),\\
 |\ell_\alpha|\le K_o\quad(\alpha\in K_p,\ \alpha\ne1),\qquad
 |\ell_\alpha|\le C_E\log x\quad(\alpha\in P\setminus K_p).
 \end{gathered}
\end{equation}

Montgomery and Vaughan give \(\sum_{n\le u}\Lambda(n)/n=\log u+O(1)\) for \(u\ge2\) \cite[\S2.2, Theorem~2.7(a)]{montgomery-vaughan-2007}. 

\begin{lemma}[Explicit bounds for the Mangoldt sum]\label{lem:p2c2-mertens}
For every integer \(r\ge1\),
\[
 \log r-1\le\sum_{1\le n\le r}\frac{\Lambda(n)}n
 \le\log r+\log4+4.
\]
\end{lemma}
\begin{proof}
Put \(\Psi(r)=\sum_{1\le n\le r}\Lambda(n)\). Sum \(\log j=\sum_{n\mid j}\Lambda(n)\) for \(1\le j\le r\) to get \(\log(r!)=\sum_{n\le r}\Lambda(n)\lfloor r/n\rfloor\). Thus
\[
 \log(r!)\le r\sum_{n\le r}\frac{\Lambda(n)}n
 \le\log(r!)+\Psi(r).
\]
The exponential series gives \(r^r/r!\le e^r\), and \(r!\le r^r\).
Rosser and Schoenfeld give \(\Psi(r)<1.03883r\) for \(r>0\) \cite[Theorem~12]{rosser-schoenfeld-1962}. Hence \(\Psi(r)\le(\log4+4)r\). Substitute and divide by \(r\).
\end{proof}

\begin{theorem}[Arithmetic comparison]\label{B2.14}
Fix \(C,C_E,K_o,c,M'\in\R\) and \(\varepsilon>0\), with \(C\ge0\) and \(|c|\le1\). For some \(x_1\) depending on these constants, the following holds for \(x\ge x_1\).

Let \(A\) be a finite abelian group, \(A_0\le A\), \(K_p\subseteq\widehat G\), and \(S\subset\N\) finite. Let \(m,N\in\N\) and take arithmetic data \(D,\ell,TT,T_2,a_t,w_{s,j}\) as above. Assume \(d\) is prime, \(2,q\in S\), and \(d\notin S\). Require \(I_d=\{0,\epsilon_d\}\) and \(I_p=\{0\}\) for primes \(p\notin S\cup\{d\}\). For every \(p\in S\), require \(I_p\subseteq\Phi\) and \(F_p\in\Phi\).

Let \(k>0\) and \(a_E\ge0\). For real \(b,b',\eta,X',\theta\), require \(b,b',X'>0\), \(\eta,\theta\ge0\), and \(b'X'\le xb\). Impose the weight, cost, and thin-value hypotheses of \ref{B2.10}, including \(a_t/w_{s,j}(\rho)\le N\) for every \(s\in A_0\), \(0\le j<m\), \(t\in TT\), and \(\rho\in\Phi\). Require \(T_2\subseteq A_0\), \eqref{eq:p2c2-chars}, and
\begin{equation}\label{eq:p2c2-sizes}
 \begin{gathered}
 a_E\le2xb,\qquad \tfrac12nx^{-\varepsilon/8}\le n_T,\qquad
 x^{-\varepsilon}\le\delta_2,\\
 x^{1/2-\varepsilon}\le m\le x^{1/2+\varepsilon/2},\qquad
 xm\le k\le n^2x^{1/2-\varepsilon}.
 \end{gathered}
\end{equation}
If \(\operatorname{gap}(s)\le x^{C\varepsilon}\mathcal B\) for every \(s\in A_0\), then
\(G_m^A(x)\le x^{1/2+(C+3)\varepsilon}+\Sigma_A(x)\).
\end{theorem}
\begin{proof}[Outline]
Bound increments using \ref{B2.04}, \ref{lem:p2c2-mertens}, and the derivative of \ref{B1.26}. Sum the bounds of \ref{B2.12} at scaled top points, use the threshold to control the error, and combine the gap bound with \ref{B2.10} to apply \ref{B5.55}. The complete proof is in \S\ref{prf:B2.14}.
\end{proof}

\section{Reading the characters on the host}

Let \(\phi\) be a frame and \(H\) a host satisfying \ref{B2.01}. Use the arithmetic datum \(D\) obtained from its Artin map in the coordinates of \ref{B2.02}. Let \(\theta_\alpha=\alpha\circ T\), where \(T\) is the coordinate isomorphism of \ref{B2.02}. For a host character \(\theta\), let \(\pi_H(\theta)\) be the primitive Dirichlet character inducing \(\theta\circ\operatorname{art}\) on the units modulo \(Q_Sqd\), and put \(\chi_\alpha=\pi_H(\theta_\alpha)\). A function \(\ell\) is compatible if, for every \(\alpha\ne1\), each primitive Dirichlet character of positive modulus with values \(V_\alpha(n)\) for all \(n\ge1\) has logarithmic derivative \(\ell_\alpha\) at \(1\).

\begin{lemma}[Primitive representatives]\label{B2.15}
The map \(\pi_H\) is injective and sends \(1\) to the primitive trivial character. If \(\ell\) is compatible, then
\[
 \ell_\alpha=\frac{L'(1,\chi_\alpha)}{L(1,\chi_\alpha)}\quad(\alpha\ne1),
 \qquad
 G_\alpha(y)=g_{\chi_\alpha}(y)\quad(\alpha\in\widehat G,\ y\in\R),
\]
where \(g_1=G_z\) and \(g_\chi\) for \(\chi\ne1\) is defined in \ref{B1.26}.
\end{lemma}
\begin{proof}[Outline]
Pull back via \ref{B2.02} and use conductor and inertia to identify primitive values and logarithmic derivatives. Artin surjectivity gives injectivity, the trivial term gives \(G_z\), and the other formulas agree by \ref{B2.04} and \ref{B1.26}. The complete proof is in \S\ref{prf:B2.15}.
\end{proof}

\begin{lemma}[The number of surviving characters]\label{B2.16}
Let \(A\) be the finite abelian group of the frame \(\phi\), let \(a_0\) be its distinguished element, and put \(A_0=\langle a_0\rangle\) and \(o_\phi=|A|/\operatorname{ord}(a_0)\). Then \([A:A_0]=o_\phi\) and \(|P|=4o_\phi\). For any subgroup \(A_0\) of a finite abelian group \(A\),
\[
 P\simeq\widehat{E_0\times(A/A_0)},\qquad |P|=4[A:A_0].
\]
\end{lemma}
\begin{proof}
The characters in \(P\) are constant on the fibres of \(E_0\times A\to E_0\times(A/A_0)\), so pullback gives the isomorphism. For a finite abelian group \(B\), \(|\widehat B|=|B|\) \cite[Theorems~7 and~9]{serre-1977}. Apply this with \(B=E_0\times(A/A_0)\). The frame gives \(o_\phi=|A|/\operatorname{ord}(a_0)=[A:A_0]\).
\end{proof}

Write \(\mathcal K\) for the finite family of primitive Dirichlet characters retained at level \(\Lambda\), and \(o\in\N\) for the integer specified by the level. The family contains \(1\) and is closed under inversion. Its designated character \(\chi_{\rm sp}\), when present, is a real odd member of \(\mathcal K\). For \(\chi\ne1\) put \(\kappa_\chi=L'(1,\chi)/L(1,\chi)\). Put \(\kappa_1=0\), the archive's value from the total derivative at the pole of \(\zeta\). Set
\[
 c=(3o+1)^{-1},\qquad K_H=\{\alpha:\chi_\alpha\in\mathcal K\},
 \qquad K_o=\sum_{\chi\in\mathcal K}|\kappa_\chi|.
\]
The retained host characters form the set \(\pi_H^{-1}(\mathcal K)\). Put \(\iota_d=((1,-1),1)\), the host element with additive coordinate \(\epsilon_d\). As in \ref{B2.58}, let \(X_+\) and \(X_-\) consist of the host characters trivial on \(A_0\) that take the value \(+1\) and \(-1\), respectively, at \(\iota_d\). The pointwise bounds in \ref{B2.58} concern the set \(\mathcal B_H(\mathcal K)=(X_+\setminus\pi_H^{-1}(\mathcal K))\cup X_-\).

For \(\alpha\in P\setminus K_H\), the host character \(\theta_\alpha\) is trivial on \(A_0\) by \ref{B2.02}. Since \(\iota_d^2=1\), we have \(\theta_\alpha(\iota_d)\in\{1,-1\}\). In the first case \(\theta_\alpha\in X_+\) and \(\chi_\alpha\) lies outside \(\mathcal K\) by the definition of \(K_H\). In the second case \(\theta_\alpha\in X_-\). Thus
\begin{equation}\label{eq:p2c2-bounded}
 \theta_\alpha\in\mathcal B_H(\mathcal K)\qquad(\alpha\in P\setminus K_H).
\end{equation}

\begin{lemma}[Coset shares give the required weights]\label{B2.17}
Assume the arithmetic host conditions of \ref{B2.58} for level \(\Lambda\), host \(H\), and top set \(TT\), with parameters \(x>0\) and \(\varepsilon,C\in\R\). Let \(\ell\) be compatible. Then the character conditions \eqref{eq:p2c2-chars} hold with \(K_p=K_H\), \(c=(3o+1)^{-1}\), \(M'=4o\), \(C_E=C\), and \(K_o\) as above.
\end{lemma}
\begin{proof}[Outline]
Identify retained characters by \ref{B2.58} and \ref{B2.15}, and get \(|P|=4o\) by \ref{B2.16} and \ref{B2.58}. Use quotient-character cancellation in the coset sums, the share errors in \ref{B2.58}, and the bounds transferred via \ref{B2.15}. The complete proof is in \S\ref{prf:B2.17}.
\end{proof}

\begin{theorem}[The retained sum is the target sum]\label{B2.18}
Let \(\Lambda\) be a level and suppose that \(H\) and \(TT\) satisfy the arithmetic host conditions of \ref{B2.58} for \(\Lambda\) at \(x,\varepsilon,C\in\R\). Let \(\ell\) be compatible. With \(K_p=K_H\) and \(c=(3o+1)^{-1}\), for every \(y\in\R\),
\[
 G_m^A(y)=G_{m,\Lambda}(y)
 :=G_z(y)+\sum_{\chi\in\mathcal K\setminus\{1\}}
 w_\Lambda(\chi)\Re g_\chi(y).
\]
Here \(w_\Lambda(\chi)=1\) for the designated \(\chi_{\rm sp}\) when present, and \(w_\Lambda(\chi)=c\) otherwise.
For a primitive Dirichlet character \(\psi\) of order \(o\), let \(\mathcal X\) consist of the primitive representatives of \(\psi^j\) for \(0\le j<o\). At the cyclic level of \(\psi\), the formula gives
\[
 G_m(y)=G_z(y)+\frac1{3o+1}\sum_{\chi\in\mathcal X\setminus\{1\}}\Re g_\chi(y).
\]
\end{theorem}
\begin{proof}[Outline]
Use \ref{B2.02}, \ref{B2.15}, and \ref{B2.58} to identify \(K_H\) with \(\mathcal K\), preserving weight \(1\) on the trivial term and any designated term and \(c\) elsewhere. Replace \(G_\alpha\) by \(g_{\chi_\alpha}\) via \ref{B2.15} and reindex. The complete proof is in \S\ref{prf:B2.18}.
\end{proof}

\section{The remaining characters and the final bound}

\begin{theorem}[The sum over the remaining characters]\label{B2.19}
Under the hypotheses of \ref{B2.17},
\[
 \Sigma_A(x)\le4o\,x^{1/2+\varepsilon},\qquad K_p=K_H.
\]
\end{theorem}
\begin{proof}
For \(\alpha\in P\setminus K_H\), \eqref{eq:p2c2-bounded} puts \(\theta_\alpha\) in the range of the pointwise bound of \ref{B2.58}. Compatibility gives \(G_\alpha(x)=g_{\chi_\alpha}(x)\) by \ref{B2.15}. By \ref{B2.16} and the index condition in \ref{B2.58}, \(|P|=4o\). Since \(x>0\),
\[
 \Sigma_A(x)=\sum_{\alpha\in P\setminus K_H}|g_{\chi_\alpha}(x)|
 \le |P\setminus K_H|\,x^{1/2+\varepsilon}
 \le 4o\,x^{1/2+\varepsilon}.
\]
\end{proof}

For the cyclic level of a primitive \(\psi\) of conductor \(M\), take \(0<\varepsilon<1/52\) and a sufficiently large window point \(x\ge1\) as in \ref{B2.59}. Choose \(Y,q,d,e\) by \ref{B2.48}, so that \(Y\le x\) and \((q,d)\) is a good pair in the sense of \ref{B2.42} for modulus \(M\), exponent \(\varepsilon\), and interval \([Y,x]\). Use these primes in the host of \ref{B2.46} and its completion in \ref{B2.59}. For \(\alpha\in P\setminus K_H\), \ref{B2.45} expresses \(\chi_\alpha\) on the units modulo \(Q_Sqd\) as a twist of a power \(\xi\) of \(\psi\) by the quadratic characters at \(q,d\), with exponents \(i,j\in\{0,1\}\). The pair \((i,j)\) is nonzero. Indeed, \(M\mid Q_S\), so the zero pair would give \(\chi_\alpha=\xi^{\rm prim}\in\mathcal X\) by \ref{B2.33}, and hence \(\alpha\in K_H\). Since \(Q_S\mid\operatorname{lcm}(2^{e+2},M)\), the twist identity holds on the units modulo \(2Mqd\). The bound in \ref{B2.42} at \(y=x\) gives
\[
 |g_{\chi_\alpha}(x)|\le x^{1/2+\varepsilon/2}\le x^{1/2+\varepsilon},
\]
as in \ref{B2.47}.

\begin{definition}[Chain data at a host]\label{def:p2c2-core-at-host}
Let \(H\) be a raw host on a frame \(\phi\) with group \(A\), and let \(TT\subseteq A\). We say that \(H\) and \(TT\) admit chain data at \(x,\varepsilon\) with bounds \(e_D,e_L,K,c_{HH},e_K\) if the following requirements hold. There is a chain datum \(\mathscr D\) with constants \(\mathcal C\) and \(X''\in\R\), a translation map \(\operatorname{tr}\) and bijections \(e_T,e_P\) as in \ref{B2.13}, and a compatible \(\ell\). The data also include a cutoff \(N\in\N\), a set \(T_2\subseteq A_0\), weights \(a,w\) as in \ref{B2.13}, and real parameters \(\eta,X',\theta\). Require the window point and exponent of \(\mathscr D\) to be \(x\) and \(\varepsilon\), respectively. Write \(c_{HH}(\mathcal C)\) for the height exponent of this chain and require
\[
 e_{\rm Div}\le e_D,\qquad e_{\rm Loss}\le e_L,\qquad c_{HH}(\mathcal C)\le c_{HH}.
\]
For the arithmetic datum of \(H\), require the identifications, \(\mathrm{H}1\) through \(\mathrm{H}11\), and the bounds on \(B_0,N_c\) of \ref{B2.13}. Impose the window hypotheses of \ref{B5.45} through its trace bound.
In \ref{B2.14} use the same \(m,T_2,a,w,N\). Take \(b,b'\) to be the thick and thin weight scales of \(\mathscr D\), \(k\) its starting rank, and \(a_E\) its top-weight scale. Impose the weight, cost, thin-value, and cutoff conditions of \ref{B2.14} for the arithmetic datum of \(H\), the same \(\ell\), and \(\eta,X',\theta\). The translation map covers \(A_0\). The window hypotheses of \ref{B5.45} already give \(m\le x^{1/2+\varepsilon/2}\) and \(xm\le k\) in \eqref{eq:p2c2-sizes}. Require the remaining five inequalities in \eqref{eq:p2c2-sizes}, namely \(a_E\le2xb\), \(\tfrac12 n x^{-\varepsilon/8}\le n_T\), \(x^{-\varepsilon}\le\delta_2\), \(x^{1/2-\varepsilon}\le m\), and \(k\le n^2x^{1/2-\varepsilon}\).
\end{definition}

\begin{theorem}[The target prime-sum bound]\label{B2.20}
For every level \(\Lambda\), fixed \(e_D,e_L,K,c_{HH},e_K,C\in\R\), and \(\varepsilon>0\), there exists \(x_1\in\R\) such that for all \(x\ge x_1\), every frame \(\phi\), every host \(H\) satisfying \ref{B2.01}, and every top set \(TT\) in the group of \(\phi\), the bound
\[
 G_{m,\Lambda}(x)\le x^{1/2+(e_D+e_L+7)\varepsilon}
\]
holds whenever \(H\) and \(TT\) satisfy the arithmetic host conditions of \ref{B2.58} for \(\Lambda,x,\varepsilon,C\) and admit chain data in the sense of \ref{def:p2c2-core-at-host} at \(x,\varepsilon\) with bounds \(e_D,e_L,K,c_{HH},e_K\).
\end{theorem}
\begin{proof}[Outline]
Take the chain from \ref{def:p2c2-core-at-host} and bound gaps via \ref{B5.45} and \ref{B2.13}. With \ref{B2.03} and \ref{B2.17}, apply \ref{B2.14}, identify its left side by \ref{B2.18}, and use \ref{B2.19} to bound the remaining sum, absorbing \(1+4o\). The complete proof is in \S\ref{prf:B2.20}.
\end{proof}


% ---- p2c3 ----
\chapter{Choosing the auxiliary primes \(q\) and \(d\)}
\label{ch:p2c3}

Fix a positive integer \(M\), \(0<\varepsilon<1/52\), and \(A\geq1\).
For every sufficiently large \(Y\), we choose primes \(q,d\) such that the bounds of Definition~\ref{B2.42} hold throughout \([Y,Y^A]\) for every character modulo \(M\) and every \((i,j)\in\{0,1\}^2\setminus\{(0,0)\}\).
We first choose \(q\equiv3\pmod8\).
For this \(q\), we choose \(d\equiv7\pmod{16}\) subject to congruences involving \(q\) and all small odd primes.
At each stage we discard candidates with a large finite character sum, and candidates whose twisted \(L\)-function has a zero in a specified rectangle.
A mean-square estimate controls the first set, and a zero-density estimate controls the second.

\section{Candidate primes and the analytic inputs}

\begin{definition}[Candidate sets and finite twists]
\label{B2.21}
For \(U>1\), \(m_0\in\N_{>0}\), and \(a_0\in\N\), put
\[
 P(m_0,a_0,U)
 =
 \{p\text{ prime}:\sqrt U<p\leq U,\ p\equiv a_0\pmod{m_0}\}.
\]

The archive also defines \(P\) for real \(U\), with \(P=\varnothing\) for \(U\leq0\).

For \(y,Z\in\R\), \(q\in\N\), and \(\kappa:\N\to\C\), define
\(c_y(n)=\Lambda(n)\min\{1,y/n\}\) for \(n\geq1\), and
\[
 S(Z,y,q,\kappa)
 =
 \sum_{1\leq n\leq\lfloor Z\rfloor_+}
 c_y(n)\kappa(n)\left(\frac nq\right).
\]
Here \(\lfloor Z\rfloor_+=\max\{0,\lfloor Z\rfloor\}\). All real upper limits in integer sums have this meaning.
For general \(q\), the symbol is the product of the Legendre symbols over its prime factors, counted with multiplicity, with value \(1\) at \(q=0\).
The factor at \(2\) is \(0\) for even \(n\) and \(1\) for odd \(n\).
The selections use odd prime denominators.
% Source: [Mathlib's Jacobi symbol](https://leanprover-community.github.io/mathlib4_docs/Mathlib/NumberTheory/LegendreSymbol/JacobiSymbol.html).

For \(\alpha>0\), \(Y>1\), and \(U>1\), call
\((m_0,a_0,U)\) \emph{admitted at \((\alpha,Y)\)} if
\[
 \begin{gathered}
 m_0>0,\qquad 8\mid m_0,\qquad (a_0,m_0)=1,\\
 \frac14(\log Y)^2\log m_0\leq\log U,\qquad
 Y^{\alpha/6}\leq\log U
 \leq Y^{\alpha/5}(\log Y)^2.
 \end{gathered}
\]
For a finite set \(P\subset\N\) and a set \(I\subseteq\R\), the \emph{mean-square bound} on \(P\times I\) with cutoff \(Z\) and constant \(C\) means that
\[
 \sum_{q\in P}|S(Z,y,q,\kappa)|^2
 \leq Cy\log y\,|P|
\]
for every \(y\in I\) and every \(\kappa:\N\to\C\) with
\(|\kappa(n)|\leq1\) for all \(n\).
For arbitrary \(I\), take \(\log 0=0\) and \(\log y=\log|y|\) for \(y<0\).

For \(Q>0\), let
\(\mathcal R(Q)=\{s\in\C:\Re s>399/400,\ |\Im s|\leq Q\}\).
Let \(\mathcal Z(Q)\) consist of the Dirichlet characters \(\chi\) of positive modulus \(N\) for which
\[
 \chi\ne1,\qquad N\leq Q,\qquad
 L(s,\chi)\ne0\quad(s\in\mathcal R(Q)).
\]
For such a nonprincipal character, use the centred sum
\[
 g_\chi(y)
 =
 \sum_{1\leq n\leq y}
 \Lambda(n)\chi(n)\left(\frac yn-1\right)
 +y\frac{L'}{L}(1,\chi).
\]
The \emph{analytic box} \(\mathcal B(Q,C,\tau,Y,X)\) consists of the characters \(\chi\in\mathcal Z(Q)\) for which \(\left|(L'/L)(1,\chi)\right|\leq C\log Y\) and \(|g_\chi(y)|\leq y^{1/2+\tau}\) for every \(y\in[Y,X]\).


A primitive representative of a function \(f:\N\to\C\) at a positive level \(L\) means a primitive character \(\chi\), of some positive conductor, such that
\(\chi(n)=f(n)\) whenever \((n,L)=1\).

\end{definition}

Set \(z=Y^{\varepsilon/5}\) and \(z'_0=z(\log Y)^2\).
When a window exponent \(K\geq1\) has been fixed, the finite sums will have upper limit \(Z=(Y^K)^{400}\).

\begin{theorem}[Deuring--Heilbronn repulsion]
\label{B2.22}
There exist absolute constants \(c_0,c_1,c_2>0\) and an integer
\(B\geq0\) with the following property.
Let \(M_0\geq3\) be an integer, let \(T\geq2\), and put
\(\mathcal L=\log(M_0T)\).
Let \(\chi_1,\chi\) be characters modulo \(M_0\).
Suppose
\[
 \chi_1^2=1,\quad \chi_1\ne1,\quad
 L(\beta_1,\chi_1)=0,\quad
 1-\frac{c_1}{\mathcal L}<\beta_1<1.
\]
If \(L(\rho,\chi)=0\), with the pole pair \((1,1)\) excluded, and
\[
 \Re\rho\geq1-c_0,\qquad
 |\Im\rho|\leq T,\qquad
 (\chi,\rho)\ne(\chi_1,\beta_1),
\]
then
\[
 1-\Re\rho
 \geq
 \frac{c_2}{\mathcal L}
 \log\!\left(
 \frac{c_1}{(1-\beta_1)\mathcal L^B}
 \right).
\]
\end{theorem}

Corollary~1.1 of the revised paper of Benli, Goel, Twiss, and Zaman \cite{benli-2026} gives explicit constants for \(M_0>400{,}000\), \(T\geq4\), and \(\Re\rho>1/2\).
Take \(T'=\max\{T,4\}\) and \(c_0<1/2\).
Choose \(c_1\leq1/10\). Then
\(1-\beta_1<c_1/\mathcal L\leq1/(10\log M_0)\), as required by the corollary.
In this interval \(\beta_1\) is a simple zero of \(\prod_{\chi\bmod M_0}L(s,\chi)\) by the zero-free region in \cite[(1.2)]{benli-2026}.
All factors are holomorphic at \(\beta_1\), so \(\chi_1\) is the unique character whose \(L\)-function vanishes there.
Nonvanishing at \(1\) for nonprincipal characters and the exclusion of the pole pair give \(\rho\ne1\).
Thus every pair \((\chi,\rho)\) in the statement gives a zero \(\rho\ne\beta_1\) of the product.
Its denominator \(10\log M_0+\log T'+107\) is comparable to \(\mathcal L\), so the corollary gives the bound with \(B=1\) after decreasing \(c_1\) and \(c_2\).
For the finitely many smaller moduli, nonvanishing at \(1\) for nonprincipal real characters lets us decrease \(c_1\).
The archive allows a general nonnegative integer exponent \(B\).
% Source: [Benli--Goel--Twiss--Zaman, explicit Deuring--Heilbronn repulsion](https://arxiv.org/html/2410.06082v3).

\begin{theorem}[Relative prime number theorem in progressions]
\label{B2.23}
There are absolute real constants \(c,A_0,C_0\), with \(c>0\), such that, for every positive integer \(m_0\) divisible by \(8\), one can choose either no exceptional pair or a pair \((\chi_1,\beta_1)\), where \(\chi_1\) is a character modulo \(m_0\), satisfying
\[
 \chi_1^2=1,\qquad \chi_1\ne1,\qquad
 L(\beta_1,\chi_1)=0,\qquad
 1-\frac1{8\log m_0}<\beta_1<1.
\]
The choice of no exceptional pair means that no pair with these properties exists.

Put
\[
 \operatorname{li}_2(t)=\int_2^t\frac{du}{\log u}.
\]
Let \(\pi(t;m_0,a_0)\) count primes \(p\leq t\) with
\(p\equiv a_0\pmod{m_0}\).
For \((a_0,m_0)=1\), define
\[
 H_{a_0}(t)=
 \begin{cases}
 \operatorname{li}_2(t)/\varphi(m_0),
       &\text{if there is no exceptional pair},\\[2mm]
 \bigl(\operatorname{li}_2(t)
       -\chi_1(a_0)\operatorname{li}_2(t^{\beta_1})\bigr)/
       \varphi(m_0),
       &\text{otherwise}.
 \end{cases}
\]
Then, whenever \(t\geq2\) and
\(\log t\geq C_0\log m_0\log(\log m_0+2)\), one has
\[
 \left|\pi(t;m_0,a_0)-H_{a_0}(t)\right|
 \leq
 A_0\left(
 e^{-c\log t/\log m_0}
 +e^{-\sqrt{c\log t}}
 \right)H_{a_0}(t).
\]
\end{theorem}

\begin{proof}[Outline]
Bounds \(\sum_\chi|\psi(y,\chi)-\mathrm{main}(y,\chi)|\) in a near range via the zero-free region and in a far range via a smoothed explicit formula, zero density, and Theorem~\ref{B2.22}. Orthogonality and Abel summation then give the relative bound on \(\pi(t;m_0,a_0)\). The complete proof is in \S\ref{prf:B2.23}.
\end{proof}

\begin{lemma}[The exceptional zero at a fixed modulus]\label{B2.23a}
There is an absolute \(c>0\) such that for every positive integer \(m\) there is at most one pair \((\chi,s)\), with \(\chi\) a Dirichlet character modulo \(m\), for which
\[
 L(s,\chi)=0,\qquad
 \Re s\geq 1-\frac{c}{\log(m(|\Im s|+2))}.
\]
If the pair exists, \(\chi\) is real and nonprincipal, and \(s\) is a simple real zero.
\end{lemma}

\begin{proof}[Outline]
McCurley's zero-free region at \(k=3m\) with \(c=1/50\) contains the stated region and allows at most one simple real zero of a real nonprincipal character. Injective induction from \(m\) to \(k\) and nonvanishing Euler factors on \(\Re s>0\) transfer this to modulus \(m\). The complete proof is in \S\ref{prf:B2.23a}.
\end{proof}

\begin{theorem}[Exceptional characters across moduli]
\label{B2.24}
There is an absolute \(c_{\mathrm P}>0\) such that the following holds.
Let \(D\geq3\), and let \(\chi_i\) be nonprincipal real characters modulo positive integers \(m_i\leq D\), for \(i=1,2\).
If
\(L(\beta_i,\chi_i)=0\) and
\(1-c_{\mathrm P}/\log D<\beta_i<1\), then
\[
 \chi_1(n)=\chi_2(n)\qquad((n,m_1m_2)=1).
\]
Moreover, for every \(\nu>0\), there is \(c_\nu>0\) such that every real zero \(\beta\) of a nonprincipal real character modulo a positive integer \(m_0\) satisfies
\(1-\beta\geq c_\nu m_0^{-\nu}\).
\end{theorem}

The first assertion is Page's exceptional-character theorem
\cite{page-1935}, expressed as equality of induced character values.
Equality on integers coprime to \(m_1m_2\) gives equality at that common level, hence equality of the primitive representatives.
Conversely, a common primitive representative gives the displayed equality.
To prove the first assertion of Theorem~\ref{B2.24}, take \(c_{\mathrm P}=c/3\) in terms of the constant of Lemma~\ref{B2.23a}. Induce both characters to modulus \(m=m_1m_2\). Induction multiplies each \(L\)-function by finitely many Euler factors, so each induced character retains its zero \(\beta_i\). Since \(2m\leq D^3\), both zeros satisfy \(\beta_i>1-c/\log(2m)\). Uniqueness in Lemma~\ref{B2.23a} identifies the induced characters, so \(\chi_1(n)=\chi_2(n)\) whenever \((n,m)=1\).

Davenport \cite[Chapter~21]{davenport-2000} gives Siegel's theorem.
Thorner--Zaman \cite[Remark~1.2, (1.2)]{thorner-2024} state the bound for the exceptional zero with an effectively computable \(c_\nu\) for \(\nu\geq1/2\).
Zeros outside the exceptional interval satisfy the same bound after decreasing \(c_\nu\), since
\(1/\log(3m_0)\gg_\nu m_0^{-\nu}\).
% Sources: [Page, publication record](https://academic.oup.com/plms/article-abstract/s2-39/1/116/1475562); [Thorner--Zaman, equation (1.2)](https://arxiv.org/pdf/2108.10878).

\begin{lemma}[Size of every candidate set]
\label{B2.25}
For every \(\alpha>0\), there is \(Y_0=Y_0(\alpha)\) such that, for \(Y\geq Y_0\), every progression admitted at \((\alpha,Y)\) satisfies
\(|P(m_0,a_0,U)|\geq U^{9/10}\).
\end{lemma}

\begin{proof}[Outline]
Use Theorem~\ref{B2.24} to bound the exceptional subtraction in Theorem~\ref{B2.23}. Admissibility controls the error and leaves at least \(U^{9/10}\) primes above \(\sqrt U\). The complete proof is in \S\ref{prf:B2.25}.
\end{proof}

\section{Quadratic cancellation and the mean square}

\begin{lemma}[One exceptional squarefree kernel]
\label{B2.26}
Fix \(\alpha,A,B>0\).
For all sufficiently large \(Y\), uniformly over progressions
\((m_0,a_0,U)\) admitted at \((\alpha,Y)\), there is \(n_*\in\N\) such that
\[
 \left|\sum_{q\in P(m_0,a_0,U)}
             \left(\frac qs\right)\right|
 \leq Y^{-B}|P(m_0,a_0,U)|
 \tag{3}\label{eq:aux-kernel-cancel}
\]
for every odd squarefree integer \(s\) satisfying
\(1<s\leq Y^A\), \((s,m_0)=1\), and \(s\ne n_*\).
\end{lemma}

\begin{proof}[Outline]
Theorem~\ref{B2.24} and CRT allow at most one exceptional kernel. In Theorem~\ref{B2.23}, orthogonality cancels ordinary and mismatched terms, distance from \(1\) bounds any surviving term, and Lemma~\ref{B2.25} controls the small-prime loss. The complete proof is in \S\ref{prf:B2.26}.
\end{proof}

For \(N\in\N\), define
\[
 s_p(y,N)=
 \sum_{\substack{a\geq1\\p^a\leq N}}
 \Lambda(p^a)\min\{1,y/p^a\}
 \quad(p\text{ prime}),
\]
and put \(s_p(y,N)=0\) for nonprime \(p\).
Let
\[
 R_0(y,m_0,N)=
 \sum_{\substack{p^a\leq N\\
                  a\text{ even or }p\mid m_0}}
 \Lambda(p^a)\min\{1,y/p^a\}.
\]
Each prime power in this last sum is counted once.

\begin{lemma}[Expansion of the quadratic moment]
\label{B2.27}
Let \(P\) be a finite set of primes, let \(8\mid m_0\) and
\((a_0,m_0)=1\), and set \(N=\lfloor Z\rfloor_+\).
Suppose \(y,\eta\geq0\), \(|\kappa(n)|\leq1\) for all \(n\), and
\(q\equiv a_0\pmod{m_0}\), \(q>N^2\) for every \(q\in P\).
Suppose \(s_p(y,N)\leq B\) for every \(p\), and suppose that for some
\(n_*\in\N\),
\[
 \left|\sum_{q\in P}\left(\frac qs\right)\right|
 \leq\eta|P|
\]
whenever
\[
 1<s\leq N^2,\qquad
 s\text{ odd and squarefree},\qquad
 (s,m_0)=1,\qquad s\ne n_*.
\]
Then
\[
 \begin{split}
 \sum_{q\in P}|S(Z,y,q,\kappa)|^2
 \leq |P|\bigg\{&
 \sum_{\substack{p\leq N\\p\text{ prime}}}s_p(y,N)^2
 +R_0(y,m_0,N)^2\\
 &+2B^2+2BR_0(y,m_0,N)
 +\eta\left(\sum_{n=1}^Nc_y(n)\right)^2
 \bigg\}.
 \end{split}
\]
\end{lemma}

\begin{proof}[Outline]
Expand the square and use quadratic reciprocity to group pairs by the odd squarefree kernel of \(nn'\) with primes dividing \(m_0\) removed. Count pairs with kernel \(1\) or the exceptional kernel separately, and apply cancellation to the rest. The complete proof is in \S\ref{prf:B2.27}.
\end{proof}

\begin{lemma}[Prime-power coefficient estimates]
\label{B2.28}
There are absolute positive constants \(C_1,C_2,C_3,C_4\) such that, for every \(y\geq3\) and \(N\in\N\),
\[
 s_p(y,N)\leq C_2\log y\quad(p\in\N),
 \qquad
 \sum_{\substack{p\leq N\\p\text{ prime}}}s_p(y,N)^2
 \leq C_3y\log y.
 \tag{10}\label{eq:aux-prime-power-bounds}
\]
If \(w\geq1\) and every prime \(p\mid m_0\) with \(p\leq N\) satisfies \(p\leq w\), then
\[
 R_0(y,m_0,N)\leq C_1(\sqrt y+w\log y).
 \tag{11}\label{eq:aux-R0-bound}
\]
For every \(y\geq0\) and \(Z\geq2\),
\[
 \sum_{1\leq n\leq Z}c_y(n)\leq C_4y\log Z.
 \tag{12}\label{eq:aux-total-coefficient}
\]
One may take \(C_2=3\).
\end{lemma}

\begin{proof}[Outline]
Split prime-power sums at \(p^a=y\) and bound the tails geometrically. Use Chebyshev's bound and partial summation for the squares, group even powers and primes dividing \(m_0\) for \(R_0\), and use Mertens's bound for the total mass. The complete proof is in \S\ref{prf:B2.28}.
\end{proof}

\begin{theorem}[Uniform mean square over candidate primes]
\label{B2.29}
There is an absolute \(C_{\mathrm m}>0\) with the following property.
Fix \(0<\alpha\leq1\) and \(K\geq1\).
For all sufficiently large \(Y\), let
\((m_0,a_0,U)\) be admitted at \((\alpha,Y)\), and assume
\[
 p\mid m_0,\ p\text{ prime}
 \quad\Longrightarrow\quad
 p\leq Y^{\alpha/5}\ \text{or}\ p>(Y^K)^{400}.
 \tag{13}\label{eq:aux-modulus-gap}
\]
Then, with \(Z=(Y^K)^{400}\),
\[
 \sum_{q\in P(m_0,a_0,U)}
 |S(Z,y,q,\kappa)|^2
 \leq C_{\mathrm m}y\log y\,|P(m_0,a_0,U)|
\]
for every \(Y\leq y\leq Y^K\) and every
\(\kappa:\N\to\C\) satisfying \(|\kappa(n)|\leq1\) for all \(n\).
\end{theorem}

\begin{proof}[Outline]
Insert the cancellation from Lemma~\ref{B2.26} and the coefficient bounds from Lemma~\ref{B2.28} into Lemma~\ref{B2.27}. Admissibility supplies the prime cutoff and the gap in the prime divisors of \(m_0\). The complete proof is in \S\ref{prf:B2.29}.
\end{proof}

\section{Bounds throughout a window and zero-density exclusions}

\begin{theorem}[Candidates exceeding the bound somewhere in a window]
\label{B2.30}
Fix \(C\geq0\), \(\tau>0\), and \(\eta>0\).
There is \(Y_0=Y_0(C,\tau,\eta)\) such that, for \(Y\geq Y_0\), the following holds for every real \(X,Z\), finite sets \(B\subseteq P\subset\N\), and function
\(\kappa:\N\to\C\) with \(|\kappa(n)|\leq1\) for all \(n\).
Assume the mean-square bound of Definition~\ref{B2.21} on \(P\times[Y,X]\) with cutoff \(Z\) and constant \(C\).
If every \(q\in B\) has a point \(y\in[Y,X]\) such that
\[
 |S(Z,y,q,\kappa)|>\tfrac12y^{1/2+\tau},
\]
then \(|B|\leq\eta|P|\).
\end{theorem}

\begin{proof}[Outline]
The mean-square bound applies to \(vS_q'(v)\) by restricting to \(n>v\). On each dyadic block, use the fundamental theorem of calculus and Cauchy–Schwarz to bound the proportion exceeding the threshold by \(O(T^{-\tau})\), then sum over the blocks. The complete proof is in \S\ref{prf:B2.30}.
\end{proof}

\begin{theorem}[Primitive characters with a zero near \(1\)]
\label{B2.31}
There is an absolute \(C_{\mathrm z}>0\) such that, for every \(Q\geq3\), every finite set \(\mathcal E\) of primitive Dirichlet characters satisfying
\[
 \operatorname{cond}(\chi)\leq Q,\qquad
 L(s,\chi)=0\text{ for some }s\in\mathcal R(Q)
\]
has \(|\mathcal E|\leq C_{\mathrm z}Q^{1/4}\).
A character is counted once, regardless of how many zeros it has in the rectangle.
\end{theorem}

\begin{proof}[Outline]
Multiply the partial sum of \(L(\rho_\chi,\chi)\) by a M\"obius mollifier and use partial summation to detect each zero on a dyadic block. Hal\'asz--Montgomery duality with the divisor bound and character-pair summation, summed over blocks, yields the bound. The complete proof is in \S\ref{prf:B2.31}.
\end{proof}

In the literature, the same conclusion also follows from the \(4/5\leq\sigma\leq 1\) clause of Montgomery's Theorem~12.2 \cite[Theorem~12.2]{montgomery-1971}, which states that for \(4/5\leq\sigma\leq1\) and \(T\geq2\),
\[
 \sum_{\substack{\chi\ \mathrm{primitive}\\
                  \operatorname{cond}(\chi)\leq Q}}
 N(\sigma,T,\chi)
 \ll (Q^2T)^{2(1-\sigma)/\sigma}(\log QT)^{14},
\]
where \(N(\sigma,T,\chi)\) counts nontrivial zeros \(\rho\) with \(\Re\rho\geq\sigma\) and \(|\Im\rho|\leq T\). At \(\sigma=399/400\) and \(T=Q\), this gives \(O(Q^{6/399}(\log Q)^{14})=O(Q^{1/4})\).

\begin{lemma}[From bad characters to bad candidate primes]
\label{B2.32}
There is an absolute \(C_{\mathrm b}\geq0\) with the following property.
Fix real numbers \(Q\geq3\) and \(F\), and let \(B\) be a finite set of primes.
Suppose that every \(q\in B\) satisfies \(q>F\) and admits a primitive character \(\chi_q\) of conductor
\(N_q=f_qq\), where \(f_q\in\N_{>0}\), \(f_q\leq F\), and \(N_q\leq Q\), whose \(L\)-function has a zero in \(\mathcal R(Q)\).
Then \(|B|\leq C_{\mathrm b}Q^{1/4}\).
\end{lemma}

\begin{proof}
Choose one witness \((N_q,\chi_q)\) for each \(q\).
The resulting map to primitive characters with their conductors is injective.
Indeed, equality of two witnesses gives \(f_qq=f_{q'}q'\).
The prime \(q\) divides \(f_{q'}q'\).
If \(q\mid f_{q'}\), then \(q\leq f_{q'}\leq F\), contrary to \(q>F\).
Thus \(q\mid q'\), so \(q=q'\).
Apply Theorem~\ref{B2.31} to the image.
\end{proof}

\begin{lemma}[Primitive twists and uniqueness of representatives]
\label{B2.33}
Let \(\xi\) be a primitive character of positive conductor \(f\), and let \(q\) be an odd prime with \(q\nmid f\).
Then there is a nonprincipal primitive character \(\chi\) of conductor \(fq\) such that
\[
 \chi(n)=\xi(n)\left(\frac nq\right)
 \qquad(n\in\N).
 \tag{18}\label{eq:aux-primitive-twist}
\]
Also, if primitive characters \(\chi,\chi'\) of positive conductors
\(N,N'\) satisfy \(\chi(n)=\chi'(n)\) whenever \((n,L)=1\), for some positive integer \(L\), then their conductors and characters agree.
\end{lemma}

\begin{proof}[Outline]
The Chinese remainder theorem combines the primitive factors of coprime conductors into a nonprincipal character of conductor \(fq\). For uniqueness, induce both representatives to \(NN'L\) and descend their common character to \(\gcd(N,N')\), then use primitivity. The complete proof is in \S\ref{prf:B2.33}.
\end{proof}

\section{Recovering the centred prime sum}

For a nonprincipal character \(\chi\) of positive modulus, define
\[
 D_\chi=-\frac{L'}{L}(1,\chi),\qquad
 T_\chi(V)=D_\chi-\sum_{1\leq n\leq V}
                         \frac{\Lambda(n)\chi(n)}n.
\]


\begin{lemma}[Tail estimate from a zero-free rectangle]
\label{B2.34}
Fix \(0<\alpha\leq1\) and \(K\geq1\).
There is \(Y_0=Y_0(\alpha,K)\) such that, for \(Y\geq Y_0\), every character \(\chi\) of positive modulus \(N\) satisfying
\[
 Y^{\alpha/6}\leq\log Q\leq Y^{\alpha/4},
 \qquad \chi\in\mathcal Z(Q)
\]
obeys
\(\left|Y^KT_\chi((Y^K)^{400})\right|\leq Y^\alpha\).
\end{lemma}

\begin{proof}[Outline]
Shifts the integrated Perron contour for \(\mathcal S_\chi(t)\) to \(1-1/400+\delta\) and bounds \(L'/L\) via convexity and Landau's Borel--Carath\'eodory lemma on the zero-free rectangle. Comparing the contour and finite-sum differences bounds \(|T_\chi(V)|\). The complete proof is in \S\ref{prf:B2.34}.
\end{proof}

\begin{lemma}[Logarithmic derivatives, including the constant \(4\)]
\label{B2.35}
There is an absolute \(C_{\mathrm L}>0\) such that, for every
\(0<\alpha\leq1\), there is \(Y_0=Y_0(\alpha)\) such that, for every \(Y\geq Y_0\) and every character \(\chi\) of positive modulus,
\[
 \left|\frac{L'}L(1,\chi)\right|\leq C_{\mathrm L}\log Y
 \tag{22}\label{eq:aux-kappa-C}
\]
whenever \(Y^{\alpha/6}\leq\log Q\leq Y^{\alpha/4}\) and \(\chi\in\mathcal Z(Q)\).

More precisely, fix \(0<\varepsilon<1/52\) and \(A_1>0\).
There is \(Y_1=Y_1(\varepsilon,A_1)\) such that, for \(Y\geq Y_1\),
\[
 Y^{\varepsilon/6}\leq\log Q
 \leq A_1Y^{\varepsilon/5}(\log Y)^2,
 \qquad \chi\in\mathcal Z(Q)
\]
imply
\[
 \left|\frac{L'}L(1,\chi)\right|\leq4\log Y.
 \tag{23}\label{eq:aux-kappa-four}
\]
\end{lemma}

\begin{proof}[Outline]
Lemma~\ref{B2.34} bounds the tail at \(Y^{400}\) and Mertens bounds the finite sum. For the explicit constant, (\ref{eq:aux-tail-general}) bounds the tail at \((\log Q)^{1000}\), giving \(1000\log\log Q+O(1)\) and hence the \(4\log Y\) bound. The complete proof is in \S\ref{prf:B2.35}.
\end{proof}

\begin{lemma}[From a finite sum to the analytic box]
\label{B2.36}
There is an absolute \(C>0\) such that the following holds.
Fix \(0<\alpha\leq1/4\), \(\tau>0\), and \(K\geq1\).
For all sufficiently large \(Y\), suppose \(\chi\) is a character of positive modulus satisfying
\[
 Y^{\alpha/6}\leq\log Q\leq Y^{\alpha/4},
 \qquad \chi\in\mathcal Z(Q),
\]
and suppose that, for some \(q\in\N\) and \(\kappa:\N\to\C\),
\(\chi(n)=\kappa(n)(n/q)\) for every \(n\in\N\).
If
\[
 |S((Y^K)^{400},y,q,\kappa)|
 \leq\tfrac12y^{1/2+\tau}
 \qquad(Y\leq y\leq Y^K),
\]
then \(\chi\in\mathcal B(Q,C,\tau,Y,Y^K)\).
\end{lemma}

\begin{proof}[Outline]
With \(Z=(Y^K)^{400}\), splitting at \(n=y\) gives \(g_\chi(y)=S(Z,y,q,\kappa)-yT_\chi(Z)\). Bounding \(|yT_\chi(Z)|\le Y^\alpha\le\frac{1}{2}y^{1/2+\tau}\) via \ref{B2.34} and the triangle inequality gives the window bound on \(g_\chi\), plus the logarithmic derivative bound. The complete proof is in \S\ref{prf:B2.36}.
\end{proof}

\section{The two selections}

\begin{theorem}[One selection stage]
\label{B2.37}
There is an absolute \(C>0\) with the following property.
Fix \(0<\alpha\leq1/4\), \(\tau,\eta>0\), \(K\geq1\), and \(R\in\N\).
There is \(Y_0=Y_0(\alpha,\tau,K,\eta,R)\) such that the following holds for \(Y\geq Y_0\).
Let \((m_0,a_0,U)\) be admitted at \((\alpha,Y)\), and assume
\eqref{eq:aux-modulus-gap}.
Let \(1\leq F\leq\sqrt U\), and let \(\Xi\) be a finite set of primitive characters such that
\(|\Xi|\leq R\) and \(\operatorname{cond}(\xi)\leq F\) for \(\xi\in\Xi\).
Then \(P=P(m_0,a_0,U)\) is nonempty.

Let \(B\subseteq P\). Suppose that for every \(q\in B\) there are \(\xi\in\Xi\) and a positive integer \(L\) for which some primitive representative of \(n\mapsto\xi(n)(n/q)\) at level \(L\) lies outside \(\mathcal B(FU,C,\tau,Y,Y^K)\). Then \(|B|\leq\eta|P|\).
\end{theorem}

\begin{proof}[Outline]
Use Lemma~\ref{B2.25} and Lemma~\ref{B2.32} to bound zero failures, and Theorems~\ref{B2.29} and \ref{B2.30} to bound finite-sum failures over \(\Xi\). For each remaining prime, \ref{B2.33} identifies the primitive twist and \ref{B2.36} gives the box bounds. The complete proof is in \S\ref{prf:B2.37}.
\end{proof}

Put
\[
 P_1(z)=P(8,3,e^z),\qquad
 \mathfrak m(z,q)=q\prod_{\substack{3\leq p\leq z\\p\text{ prime}}}p.
\]
Define
\[
 \begin{split}
 P_2(z,z',q)=\{d\text{ prime}:\
 &e^{z'/2}<d\leq e^{z'},\\
 &d\equiv7\pmod{16},\
 \mathfrak m(z,q)\mid d+1\}.
 \end{split}
 \tag{29}\label{eq:aux-second-set}
\]
\begin{definition}[Selection conditions for the two primes]\label{def:aux-twist-selection}
Fix a positive integer \(N_0\) and real parameters \(C,\tau,Y,X,z\). The \emph{first selection condition} on \(q\) requires that, for every character \(\xi\) modulo \(N_0\), every primitive representative of \(n\mapsto\xi(n)(n/q)\) at level \(N_0q\) belong to \(\mathcal B(N_0e^z,C,\tau,Y,X)\).
For a given \(q\) and real \(z'\), the \emph{second selection condition} on \(d\) requires that, for every character \(\xi\) modulo \(N_0\) and each \(a\in\{0,1\}\), every primitive representative of \(n\mapsto\xi(n)(n/q)^a(n/d)\) at level \(N_0qd\) belong to \(\mathcal B(N_0e^{z+z'},C,\tau,Y,X)\).
\end{definition}

\begin{lemma}[First selection]
\label{B2.38}
There is an absolute \(C_1>0\) with the following property.
Fix \(0<\varepsilon<1/52\), \(\tau,\eta>0\), \(K\geq1\), and \(N_0\in\N_{>0}\).
There is \(Y_0=Y_0(\varepsilon,\tau,K,\eta,N_0)\) such that, for \(Y\geq Y_0\), with \(z=Y^{\varepsilon/5}\), the set \(P_1(z)\) is nonempty.
Every subset \(B\subseteq P_1(z)\) whose members fail the first selection condition of Definition~\ref{def:aux-twist-selection} on \([Y,Y^K]\) with constant \(C_1\) satisfies \(|B|\leq\eta|P_1(z)|\).
\end{lemma}

\begin{proof}[Outline]
Apply Theorem~\ref{B2.37} to primitive representatives modulo \(N_0\) with \(m_0=8\), \(a_0=3\), \(U=e^z\), and \(F=N_0\). The parameters meet the hypotheses for large \(Y\), and agreement on units modulo \(N_0q\) transfers the first selection condition. The complete proof is in \S\ref{prf:B2.38}.
\end{proof}

\begin{lemma}[Second selection, uniformly in \(q\)]
\label{B2.39}
There is an absolute \(C_2>0\) with the following property.
Fix \(0<\varepsilon<1/52\), \(\tau,\eta>0\), \(K\geq1\), and \(N_0\in\N_{>0}\).
There is \(Y_0=Y_0(\varepsilon,\tau,K,\eta,N_0)\) such that the following holds for \(Y\geq Y_0\).
Put \(z=Y^{\varepsilon/5}\).
For every \(z(\log Y)^2\leq z'\leq K^2z(\log Y)^2\) and every
\(q\in P_1(z)\), the set \(P_2(z,z',q)\) is nonempty.
Every subset \(B\subseteq P_2(z,z',q)\) whose members fail the second selection condition of Definition~\ref{def:aux-twist-selection} on \([Y,Y^K]\) with constant \(C_2\) satisfies \(|B|\leq\eta|P_2(z,z',q)|\).
\end{lemma}

\begin{proof}[Outline]
Use CRT and the primorial bound to verify the progression and size conditions. Apply Theorem~\ref{B2.37} to primitive representatives and their \(q\)-twists from Lemma~\ref{B2.33}, and transfer the second selection condition by agreement on units modulo \(N_0qd\). The complete proof is in \S\ref{prf:B2.39}.
\end{proof}

\begin{theorem}[Proportions in the two selections]
\label{B2.40}
There is an absolute \(C>0\) such that, for fixed
\(0<\varepsilon<1/52\), \(\tau,\eta>0\), \(K\geq1\), and \(N_0\in\N_{>0}\), the following holds for all sufficiently large \(Y\).
Put \(z=Y^{\varepsilon/5}\).
For every real \(X\leq Y^K\) and
\(z(\log Y)^2\leq z'\leq K^2z(\log Y)^2\), there is
\(G_1\subseteq P_1(z)\) such that
\[
 |P_1(z)|>0,\qquad |G_1|\geq(1-\eta)|P_1(z)|.
\]
Every \(q\in G_1\) satisfies the first selection condition on \([Y,X]\) with constant \(C\).
For each such \(q\), there is \(G_2(q)\subseteq P_2(z,z',q)\) such that
\[
 |P_2(z,z',q)|>0,\qquad
 |G_2(q)|\geq(1-\eta)|P_2(z,z',q)|,
\]
and every \(d\in G_2(q)\) satisfies the second selection condition for this \(q\) on the same window with constant \(C\).
\end{theorem}

\begin{proof}[Outline]
Apply \ref{B2.38} and \ref{B2.39} using the larger of their constants and restrict from \([Y,Y^K]\) to \([Y,X]\). Defining \(G_1\) and \(G_2(q)\) as the complements of the failing candidates in \(P_1(z)\) and \(P_2(z,z',q)\) yields nonemptiness and both proportion bounds. The complete proof is in \S\ref{prf:B2.40}.
\end{proof}

\begin{theorem}[A single pair satisfying both stages]
\label{B2.41}
There is an absolute \(C>0\) such that the following holds.
Fix \(0<\varepsilon<1/52\), \(\tau>0\), \(K\geq1\), and \(N_0\in\N_{>0}\).
For every sufficiently large \(Y\), every real \(X\leq Y^K\), and every
\(z'_0\leq z'\leq K^2z'_0\), where
\(z=Y^{\varepsilon/5}\) and \(z'_0=z(\log Y)^2\), there are primes \(q,d\) with \(N_0<q<d\) satisfying
\[
 \begin{gathered}
 q\equiv3\pmod8,\qquad
 z/2<\log q\leq z,\\
 d\equiv7\pmod{16},\qquad
 z'/2<\log d\leq z',\\
 p\mid d+1
 \quad\text{for every odd prime }p
       \text{ with }p\leq z\text{ or }p=q,
 \end{gathered}
 \tag{34}\label{eq:aux-pair-rows}
\]
and both selection conditions of Definition~\ref{def:aux-twist-selection} on \([Y,X]\) with constant \(C\).
\end{theorem}

\begin{proof}[Outline]
Apply \ref{B2.40} with \(\eta=1/2\) so \(G_1\) and \(G_2(q)\) are nonempty, and pick \(q\in G_1\) and \(d\in G_2(q)\) to obtain the congruence, size, and box conditions. For large \(Y\), \(e^{z/2}>N_0\) and \(z'\ge 2z\) ensure \(N_0<q<d\). The complete proof is in \S\ref{prf:B2.41}.
\end{proof}

\section{The pair required by the arithmetic construction}

\begin{definition}[Window pairs and good pairs]
\label{B2.42}
For real \(\varepsilon,Y,X\), put \(z=Y^{\varepsilon/5}\) and \(z'_0=z(\log Y)^2\).
A \emph{window pair} \((q,d)\) consists of primes satisfying
\eqref{eq:aux-pair-rows} with \(z'=z'_0\).

For \(i,j\in\{0,1\}\), put
\[
 f_{\xi,i,j}(n)=
 \xi(n)\left(\frac nq\right)^i
        \left(\frac nd\right)^j,\qquad
 H_j=8Me^ze^{jz'_0}.
\]

A window pair \((q,d)\) is \emph{good for modulus \(M\) on \([Y,X]\) with exponent \(\varepsilon\)} if \((M,qd)=1\) and the following two conditions hold.

For every character \(\xi\) modulo \(M\), every
\((i,j)\in\{0,1\}^2\setminus\{(0,0)\}\), and every primitive representative
\(\chi\) of \(f_{\xi,i,j}\) at level \(2Mqd\),
\[
 |g_\chi(y)|\leq y^{1/2+\varepsilon/2},
 \qquad
 \left|\frac{L'}L(1,\chi)\right|\leq4\log y
 \quad(Y\leq y\leq X).
 \tag{35}\label{eq:aux-pair-good-bounds}
\]
For every real primitive character \(\rho\) whose conductor divides \(8M\), every
\((i,j)\in\{0,1\}^2\setminus\{(0,0)\}\), and every primitive representative
\(\chi\) of \(f_{\rho,i,j}\) at level \(2Mqd\),
\[
 L(s,\chi)\ne0
 \quad
 \bigl(\Re s>399/400,\ |\Im s|\leq H_j\bigr).
 \tag{36}\label{eq:aux-pair-good-boxes}
\]
\end{definition}

\begin{theorem}[Good auxiliary primes for every sufficiently large window]
\label{B2.43}
For every positive integer \(M\), every \(0<\varepsilon<1/52\), and every
\(A\geq1\), there is \(Y_1\) such that, for every \(Y\geq Y_1\), there exists a good pair of primes \((q,d)\) for modulus \(M\) on \([Y,Y^A]\) with exponent \(\varepsilon\), in the sense of Definition~\ref{B2.42}.
\end{theorem}

\begin{proof}[Outline]
Use \ref{B2.41} for the window pair and second selection, and \ref{B2.33} for the twist in first selection. Induce to \(8M\) for \ref{B2.42} and the centred sum in \ref{B2.21}, using the boxes for zero-free rectangles and \ref{eq:aux-kappa-four} for \(4\log y\). The complete proof is in \S\ref{prf:B2.43}.
\end{proof}


% ---- p2c4 ----
\chapter{Building the arithmetic host of a window}\label{ch:arithmetic-window-host}

Fix a primitive Dirichlet character $\psi$ of conductor $M\geq1$.
Write $\Xi=\langle\psi\rangle$, $o_\psi=|\Xi|$, and
$\mathcal X=\{\xi^{\mathrm{prim}}:\xi\in\Xi\}$ for the retained family
of the cyclic level $\mathcal L_\psi$ of Definition~\ref{B1.21}.

In the frame of \ref{B2.01}, put $n=|A|$, $n_0=|A_0|$, and
$o_\phi=[A:A_0]$, where $A_0=\langle a_0\rangle$. For a scale
$s=(x,\varepsilon,J)$ with $x>1$, $0<\varepsilon<1/52$ and real $J\geq1$,
write $b=\exp(x^\varepsilon)$, $\delta_1=x^{-\varepsilon/8}$ and
$\varepsilon_z=\varepsilon/J$.
The rank conditions on $R$ at scale $s$ require
\begin{equation}\label{eq:host-rank-ranges}
 \begin{gathered}
 x^{2+8\varepsilon}\leq n\leq2x^{2+8\varepsilon},\qquad
 \sqrt{x\delta_1}\leq m\leq2\sqrt{x\delta_1},\\
 x^{9/2+10\varepsilon}\leq k\leq2x^{9/2+10\varepsilon}.
 \end{gathered}
\end{equation}
Here $f=2n$, $r=k+m-1$, and the auxiliary rank $t$ satisfies
$r+t=8fr+4+T_On_W$, with $T_O,n_W$ as in \ref{B8.59}.

\section{The field and its characters}

For a subgroup $\Xi$ of the characters modulo $M$, let
\[
 B=(\Z/M\Z)^\times/\operatorname{ann}(\Xi),\qquad
 \operatorname{ann}(\Xi)=\{u:\xi(u)=1\text{ for every }\xi\in\Xi\}.
\]
Finite character duality gives $\widehat B=\Xi$. For $e\geq0$, put
$H_e=(\Z/2^{e+2}\Z)^\times/\{\pm1\}$ and $D_e=\operatorname{lcm}(2^{e+2},M)$.
Let $A^{(e)}$ be the image of $(\Z/D_e\Z)^\times$ in $H_e\times B$.
Choose $Q_S\mid D_e$ minimal among multiples of $2^{e+2}$ through
which this map factors. In the frame at $e$, put $A=A^{(e)}$ and
$A_0=\ker(A^{(e)}\to B)$. Let $a_0$ generate $A_0$.

\begin{lemma}[The order of the frame and the Artin kernel]\label{B2.44}
For every $M\geq1$ and subgroup $\Xi$ of the Dirichlet characters modulo
$M$, there are integers $1\leq c\leq|\Xi|$ and $e_0$ such that
$|A^{(e)}|=c2^e$ for $e\geq e_0$. At every $e$, $[A:A_0]=|\Xi|$.

Take $\Xi=\langle\psi\rangle$. Suppose $e\geq2$, the primes $q,d$
satisfy
\begin{equation}\label{eq:host-prime-hyp}
 q\equiv3\pmod8,\qquad d\equiv7\pmod{16},\qquad
 \left(\frac{-d}{q}\right)=1,
\end{equation}
and $(M,qd)=1$ and $(-d/\ell)=1$ for every odd prime $\ell\mid M$.
In the resulting cyclic host, $M\mid Q_S$. For every integer $a\geq0$
coprime to $Q_Sqd$,
\begin{equation}\label{eq:cyclic-artin-kernel}
 \begin{split}
 \operatorname{Art}(a)=1\quad\Longleftrightarrow\quad&
 a\equiv1\text{ or }-1\pmod{2^{e+2}},\\
 &\psi(a)=1,\qquad
 \left(\frac aq\right)=\left(\frac ad\right)=1 .
 \end{split}
\end{equation}
\end{lemma}

\begin{proof}[Outline]
Project onto \(H_e\) and \(B\) to get the order and index, and use finiteness of \(B\) to stabilize \(C_e\). The Artin kernel imposes both signs and the joint image, and primitivity gives \(M\mid Q_S\). The complete proof is in \S\ref{prf:B2.44}.
\end{proof}

\begin{lemma}[Characters trivial on the averaging subgroup]\label{B2.45}
Under the hypotheses of the second part of \ref{B2.44}, every host
character $\theta$ trivial on $A_0$ has a representation
\begin{equation}\label{eq:host-twist-form}
 \theta(\operatorname{Art}(a))
   =\xi(a)\left(\frac aq\right)^i\left(\frac ad\right)^j,
 \qquad (a,Q_Sqd)=1,
\end{equation}
where $\xi\in\Xi$ and $i,j\in\{0,1\}$. On the inertia generator $\epsilon_d=((1,-1),1)$ at $d$, one has $\theta(\epsilon_d)=(-1)^j$.

\end{lemma}

\begin{proof}
The restriction of $\theta$ to $A$ factors through $A/A_0\simeq B$
and defines $\xi\in\widehat B=\Xi$. Factor
\[
 ((s_q,s_d),z)=
 ((s_q,1),w(s_q))\,((1,s_d),1)\,((1,1),z\,w(s_q)^{-1}).
\]
The first two factors give the characters $s_q\mapsto s_q^i$ and
$s_d\mapsto s_d^j$. Substitution into \eqref{eq:artin-with-sign} gives
\eqref{eq:host-twist-form}. Evaluating the second sign factor at $-1$ gives $\theta(\epsilon_d)=(-1)^j$.
\end{proof}

For $B_0\in\N$, impose the following field conditions at $B_0$. The first coordinate of $\operatorname{Art}(a)$ is $(a/q)$, and $E$ contains a square root of $-q$ on which each Galois automorphism acts by its first sign. For each prime $\ell\mid Q_S$ there must be a unit $u$
modulo $Q_Sqd$ with $u\equiv1\bmod Q_Sqd/\ell$ and
$\operatorname{Art}(u)\ne1$. Require also $Q_S\leq n^2$,
$o_\phi^{64}\leq2^e$, $q\ne3$, $(-d/p)=1$ for odd primes $p\leq B_0$,
and $4\ell\leq B_0$ for odd primes $\ell\mid Q_S$.

\begin{theorem}[Constructing the field]\label{B2.46}
For every primitive $\psi$, there are integers $c\geq1$ and $e_0$
with the following property. Let $e\geq e_0$ and let $q,d$ be primes
satisfying \eqref{eq:host-prime-hyp}, with $q\ne3$ and $(M,qd)=1$.
Suppose $B_0\geq4M$ and $(-d/p)=1$ for every odd prime $p\leq B_0$.
There are a frame $\phi$ at the prescribed exponent $e$ and a host arithmetic datum $H$ on $\phi$ with the prescribed primes $q,d$, such that $n=c2^e$, $o_\phi=o_\psi$, and
\begin{equation}\label{eq:field-conductor-bounds}
 Q_S\mid\operatorname{lcm}(2^{e+2},M),\qquad Q_S\leq4Mn.
\end{equation}
The field conditions at $B_0$ hold, $M\mid Q_S$, and the Artin kernel is given by \eqref{eq:cyclic-artin-kernel}. Every host character trivial on $A_0$ has the twist representation and inertia value of Lemma~\ref{B2.45}.
Every member of $\mathcal X$ is represented by a host character
trivial on $A_0$ with value $1$ on $\epsilon_d$. The primitive representative of the $q$-sign character $\theta_q((s_q,s_d),z)=s_q$ lies outside $\mathcal X$.
\end{theorem}

\begin{proof}[Outline]
Use \ref{B2.44} for the frame order and take the Artin kernel's fixed field, using CRT and Gauss sums for \ref{B2.01}. Apply \ref{B2.03} to factor the different, use minimality and the witness at \(2\) for local conditions, and use \ref{B2.44} and \ref{B2.45} for twists. The complete proof is in \S\ref{prf:B2.46}.
\end{proof}

Let $X_+$ and $X_-$ denote the host characters trivial on $A_0$ with
values $1$ and $-1$ on $\epsilon_d$, respectively. A character is
\emph{retained} if its primitive Dirichlet representative belongs to
$\mathcal X$. Let $\mathcal B$ consist of the nonretained characters in $X_+$ and all characters in $X_-$.

\begin{theorem}[A nonzero twist for every remaining character]\label{B2.47}
Let $\varepsilon>0$, $x\geq1$, and $Y\leq x$. Suppose the twist representation and inertia value of Lemma~\ref{B2.45} hold for $H$, $M\mid Q_S$, and every
prime divisor of $Q_S$ divides $2M$. Assume $(q,d)$ is a good pair for modulus $M$ on $[Y,x]$ with parameter $\varepsilon$, as in Definition~\ref{B2.42}.
For every $\theta\in\mathcal B$ and every primitive
Dirichlet representative $\chi$ of $\theta$,
\[
 |g_\chi(x)|\leq x^{1/2+\varepsilon},
 \qquad \left|\frac{L'}{L}(1,\chi)\right|\leq4\log x.
\]
\end{theorem}

\begin{proof}[Outline]
Lemma~\ref{B2.45} writes \(\theta\circ\operatorname{Art}\) as a twist by \((a/q)^i(a/d)^j\), and membership in \(\mathcal B\) rules out \((i,j)=(0,0)\). Thus \(\chi\) represents \(f_{\xi,i,j}\) at level \(2Mqd\), and Definition~\ref{B2.42} at \(y=x\) yields both bounds. The complete proof is in \S\ref{prf:B2.47}.
\end{proof}

\section{Matching the two windows}

For the local data of Definition~\ref{B6.01}, write $s_{p,j}(E)$ for the $\Z_p$-length of $\Sha_{p,j}(E)$. Write $V_E(t)\leq B$ if, for every algebraic closure of $E$ and every choice of those data, the lengths are finite for all primes $p$ and $1\leq j<t$, and
\begin{equation}\label{eq:vv-budget}
 \sum_{p\in S}\log p\sum_{j=1}^{t-1}
          s_{p,j}(E)\leq B
 \quad\text{for every finite set of primes }S.
\end{equation}
The field condition of Definition~\ref{B6.02} requires an embedding $j:E\hookrightarrow\Q(\zeta_N)$ for some positive common
multiple $N$ of $2^{e+2},M,q,d$. The automorphism
$\zeta_N\mapsto\zeta_N^a$ fixes $j(E)$ pointwise exactly for the
unit residues in \eqref{eq:cyclic-artin-kernel}. Call this the
prescribed $\psi$-kernel.

For each $0<\varepsilon<1/52$ and page function $e$ of \ref{B6.02},
the cohomological window hypothesis of Definition~\ref{B6.03} for $\psi$ with exponent $401$ gives $J\geq1$ and $X_0$.
If $X\geq X_0$ and $(q,d)$ is a window pair at $X^{1/J}$ for $\varepsilon$ in the sense of Definition~\ref{B6.02}, some $W\geq X^{1/J}$ satisfies $W^2\leq X$. For every
$u\in[W,W^2]$ and every number field $E$ with the prescribed
$\psi$-kernel at $(e(u),q,d)$, we have
$V_E(\lfloor u^{401}\rfloor)\leq\sqrt{\exp(u^\varepsilon)}$.

\begin{lemma}[Rescaling a window]\label{B2.48a}
Let \(\varepsilon>0\), \(c\ge1\), \(\lambda=c^{1/(2+8\varepsilon)}\), and \(W\ge\max(1,\lambda^3,\lambda^{400})\). Put \(Z=\lambda W\). For \(Z\le x\le Z^{3/2}\), put \(v=x/\lambda\). Then
\[
 W\le v\le x\le W^2,\qquad
 \lfloor x^{400}\rfloor\le\lfloor v^{401}\rfloor,\qquad
 e^{v^\varepsilon/2}\le e^{x^\varepsilon/2}.
\]
If \(v^{2+8\varepsilon}\le 2^a\le 2v^{2+8\varepsilon}\), then
\(x^{2+8\varepsilon}\le c2^a\le 2x^{2+8\varepsilon}\).
\end{lemma}
\begin{proof}
We have \(\lambda\ge1\) and \(x\ge\lambda W\), so \(W\le v\le x\). Since \(\lambda^3\le W\), squaring gives \((\lambda W)^3\le W^4\), and hence \(x\le W^2\). Now \(\lambda^{400}\le W\le v\) gives \(x^{400}=\lambda^{400}v^{400}\le v^{401}\). Take floors. The exponential bound follows from \(v\le x\). Multiply the bounds for \(2^a\) by \(\lambda^{2+8\varepsilon}=c\).
\end{proof}

\begin{theorem}[A common analytic and cohomological window]\label{B2.48}
Let $0<\varepsilon<1/52$ and $c\geq1$. Assume the cohomological window hypothesis for $\psi$ with exponent $401$ and the pair-existence theorem
\ref{B2.43}. There is $J\geq1$ such that, for every $y\in\R$, there
is $Z\geq y$ for which every $x\in[Z,Z^{3/2}]$ admits
$Y\in\R$ and $q,d,e\in\N$ satisfying
\begin{equation}\label{eq:window-datum}
 \begin{gathered}
 x^{1/J}\leq Y\leq x,\qquad (q,d)\text{ is a good pair for }(M,\varepsilon)\text{ on }[Y,x],\\
 x^{2+8\varepsilon}\leq c2^e\leq2x^{2+8\varepsilon}.
 \end{gathered}
\end{equation}
Every field with the prescribed $\psi$-kernel at $(e,q,d)$ then satisfies
\begin{equation}\label{eq:window-vv}
 V_E(\lfloor x^{400}\rfloor)\leq\sqrt{\exp(x^\varepsilon)}.
\end{equation}

\end{theorem}

\begin{proof}[Outline]
Choose the page function by rounding up and apply \ref{B6.03} to a good pair supplied by \ref{B2.43}. Use \ref{B2.48a} to rescale the window, restrict the pair's interval, and transfer the bound by shortening the nonnegative sum for the data of \ref{B6.01}. The complete proof is in \S\ref{prf:B2.48}.
\end{proof}

\section{Primes for the level structure}

\begin{theorem}[A polynomial bound for a progression prime]\label{B2.49}
There is an absolute constant $A\geq1$ such that, for every integer
$N\geq2$, there is a prime $p$ with $p\equiv1\pmod N$ and $p\leq N^A$.
\end{theorem}

Linnik's theorem follows from the relative prime-counting estimate of
Thorner and Zaman \cite[\S1.2]{thorner-2024}. Apply it to the class
$1\bmod N$.

A coefficient family has a finite additive abelian group $G$ and parameters
$D\geq\max(4,|G|)$ and $D^{-1}\leq\kappa\leq4$. Its coefficients
$a_g(n),\lambda_g(n)$ are nonnegative for $g\in G$, $n\in\N$.
Require $a_g(0)=0$, $a_g(1)=\mathbf1_{g=0}$ and
\begin{equation}\label{eq:host-family-counts}
 \begin{gathered}
 \left|\sum_{1\leq n\leq y}a_g(n)-\kappa y\right|\leq10\sqrt y\quad(y\geq1),\\
 \left|\sum_n a_g(n)e^{-nt}-\kappa/t\right|\leq D^{10}t\quad(0<t\leq1).
 \end{gathered}
\end{equation}
With $\Lambda(n)$ the von Mangoldt function, require
$\sum_g\lambda_g(n)\leq2\Lambda(n)$ and
\[
 a_g(n)\log n=\sum_{uv=n}\sum_{j\in G}\lambda_j(u)a_{g-j}(v)\qquad(n\geq1).
\]
For $\chi\in\widehat G$, put $a_\chi(n)=\sum_g\chi(g)a_g(n)$ and
denote the trivial character by $1$. For $\chi^2=1$, require the
Dirichlet convolution $a_1*a_\chi$ to be nonnegative and
$(a_1*a_\chi)(n^2)\geq a_1(n)$.  Write $\widehat\lambda_\chi(n)=\sum_g\chi(g)\lambda_g(n)$,
$\Psi_\chi(y)=\sum_{1\leq n\leq y}\widehat\lambda_\chi(n)$ and
$\ell(t)=\log(D(|t|+2))$.

\begin{lemma}[The possible exceptional zero]\label{B2.49a}
For a family as above let $F_\chi(s)$ denote the continuation to
$\Re s>1/2$ of $(s-1)^{\mathbf1_{\chi=1}}\sum_{n\geq1}a_\chi(n)n^{-s}$.
There are absolute $0<c_0\leq1/20$, $a>0$ and $0<c_9\leq1$ such that,
in every family, the zeros of all the $F_\chi$ in
$\Re s\geq1-c_0/\ell(\Im s)$ consist of at most one pair $(\chi,\beta)$.
For this pair $\chi^2=1$ and $\beta$ is a simple real zero with
$3/4\leq\beta\leq1-c_9D^{-a}$.
\end{lemma}

\begin{proof}[Outline]
Derive the local logarithmic derivative by convolution and partial summation, and use trigonometric positivity and orthogonality to isolate a possible simple real zero of a real character. Theta estimates, positivity, and Cauchy's bound give the zero gap. The complete proof is in \S\ref{prf:B2.49a}.
\end{proof}

\begin{lemma}[A weighted estimate at one endpoint]\label{B2.49b}
Fix $0<c_0\leq1/20$. For any family and character $\chi$, suppose
the set $Z_\chi$ of zeros of $F_\chi$ in $\Re s\geq1-c_0/\ell(\Im s)$
is empty or consists of a simple real zero $\beta<1$. For $u\geq2$
and $\log D\leq\sqrt{\log u}$,
\begin{equation}\label{eq:host-endpoint-weight}
 \left|\Psi_\chi(u)-\mathbf1_{\chi=1}u+\sum_{\beta\in Z_\chi}\frac{u^\beta}{\beta}\right|
 \leq C u e^{-c_7\sqrt{\log u}},
\end{equation}
where $C,c_7>0$ depend only on $c_0$.
\end{lemma}

\begin{proof}[Outline]
Use Chebyshev for bounded \(\sqrt{\log u}\) and truncated Perron otherwise. Shift the contour using the local logarithmic derivative with the exceptional term removed, taking residues at the principal pole and any crossed exceptional zero and absorbing an uncrossed term. The complete proof is in \S\ref{prf:B2.49b}.
\end{proof}

\begin{lemma}[Principal-class weight in a short interval]\label{B2.49c}
There are absolute $A,c,\eta>0$ such that for every family, $X\geq1$,
$\log D\leq c\sqrt{\log X}$ and $e^{-c\sqrt{\log X}}\leq h\leq1$ imply
\[
 \sum_{X<n\leq X(1+h)}\lambda_0(n)\geq\eta hX/D^A .
\]
\end{lemma}

\begin{proof}[Outline]
Use the zero region and gap from \ref{B2.49a} in \ref{B2.49b}, apply orthogonality, and subtract at the two endpoints. Bound the exceptional integral by the gap and the endpoint errors by the hypotheses, then divide by \(|G|\). The complete proof is in \S\ref{prf:B2.49c}.
\end{proof}

For $d\equiv3\pmod4$, call $p$ a norm prime for $(d,N)$ if $p$ is prime and there exist integers $u,v$ such that
\begin{equation}\label{eq:norm-one-coordinates}
 4p=u^2+dv^2,\qquad N\mid v,\qquad 2N\mid u-2-v.
\end{equation}
Write $\mathcal P(d,N;X,X')$ for the norm primes for $(d,N)$ in $(X,X']$.

\begin{theorem}[Norm primes in a short interval]\label{B2.50}
There are absolute $A,c,c'>0$ as follows.
Let $d\geq100$ be prime with $d\equiv3\pmod4$, $N\in\N$ with $N\geq3$,
and $X,h\in\R$ with
\[
 X\geq1,\qquad
 \log(4dN^2)\leq c\sqrt{\log X},\qquad
 e^{-c\sqrt{\log X}}\leq h\leq1.
\]
Then
\begin{equation}\label{eq:short-norm-primes}
 |\mathcal P(d,N;X,X(1+h))|
 \geq\frac{c'hX}{(4dN^2)^A\log X}.
\end{equation}
\end{theorem}

\begin{proof}[Outline]
Use Minkowski's bound, lattice counts, Poisson summation, and ideal factorization to apply \ref{B2.49c} to the principal ray class. Remove higher prime powers by Chebyshev's bound and read off the coordinates from generators congruent to \(1\) modulo \(N\). The complete proof is in \S\ref{prf:B2.50}.
\end{proof}

For the separation of $g_0,g_1$, define
\[
 \mathcal B_r(g)=\{g\}\cup\bigcup_{j=1}^r
 \{\ell\text{ prime}:g^j<\ell^{\,4v_\ell(g^j-1)}\}.
\]

\begin{theorem}[Choosing the five level primes]\label{B2.51}
There is an absolute $C_8\geq1$ such that the following holds for every
$0<\varepsilon<1/52$ and all sufficiently large scales $s$.
Suppose $R$ satisfies the rank conditions of Theorem~\ref{B8.59} at $s$, $N,d\in\N$, and
$d\geq100$ is prime with $d\equiv3\pmod4$. Assume $N>0$, $d\mid N$
and $\log N\leq3x^{\varepsilon/5}(\log x)^2$.
There are primes $g_0,g_1,p,p',q_0$ with
\[
 (kb)^8\leq g_i\leq2(kb)^8,\quad
 g_i\text{ is a norm prime for }(d,N)\quad(i=0,1),\qquad
 g_0\ne g_1,\quad \mathcal B_r(g_0)\cap\mathcal B_r(g_1)=\varnothing.
\]
The primes $p,p',q_0$ are pairwise distinct and at most $N^{C_8}$.
They satisfy $p\equiv p'\equiv1\pmod{15N}$, $q_0\equiv1\pmod N$,
and $p,p'\ne g_i$ for $i=0,1$.
\end{theorem}

\begin{proof}[Outline]
Use \ref{B2.50} for a pool of norm primes and count roots modulo prime powers to exclude candidates whose sets intersect that of \(g_0\). Apply \ref{B2.49} with moduli \(15N\), \(15Np\), and \(Npp'\) to get distinct smaller progression primes. The complete proof is in \S\ref{prf:B2.51}.
\end{proof}

\section{Real characters and the thin positions}

Extend the window-pair parameters of \ref{B2.42} to real $Y$.
Put $a=\varepsilon/5$. For $Y>0$, let $z=Y^a$. For $Y<0$, let
$z=|Y|^a\cos(\pi a)$. For $Y\ne0$, put $z'=z(\log|Y|)^2$.
At $Y=0$, take $z=\mathbf1_{a=0}$ and $z'=0$.
Set $X_1(Y)=\exp z$, $X_2(Y)=\exp z'$ and
$H_M(j)=8MX_1(Y)X_2(Y)^j$ for $j\in\{0,1\}$.
Write $\mathcal R(T)=\{s\in\C:\Re s>399/400,\ |\Im s|\leq T\}$. Let $Q_{\mathrm{odd}}$
be the odd part of $Q_S$.

\begin{theorem}[Real-character boxes and the thin-position bound]\label{B2.52}
Fix a raw host $H$ on a frame $\phi$, and denote its auxiliary primes by $q,d$. Let $M\geq1$ be an integer and
$\varepsilon,Y,X\in\R$. Suppose $Q_{\mathrm{odd}}\mid M$ and every
prime divisor of $Q_S$ divides $2M$. If $(q,d)$ satisfies the good-pair conditions of Definition~\ref{B2.42} for $(M,\varepsilon,Y,X)$, then every
primitive real character $\chi$ of the host field whose conductor
$N_\chi$ is divisible by $q$ or $d$ satisfies
\[
 L(s,\chi)\ne0\quad(s\in\mathcal R(H_M(j))),\qquad
 N_\chi\leq8Q_{\mathrm{odd}}X_1(Y)X_2(Y)^j,\qquad
 j=\mathbf1_{d\mid N_\chi}.
\]
In particular, these characters have no real zero $\beta>399/400$.

Take $c_6=1$. For every $B\in\R$ and $0<\varepsilon<1/52$, there is
$x_0(B,\varepsilon)$ such that the following holds for every $x\geq x_0$,
every frame $\phi$, and every raw host $H$ on $\phi$.
Suppose $Q_{\mathrm{odd}}\leq B$ and
$\log(Q_Sqd)\leq2x^{\varepsilon/5}(\log x)^2$, and suppose every
primitive real host character with conductor divisible by $q$ or $d$
has no real zero greater than $399/400$. Then every primitive host
character $\chi$ satisfies
\begin{equation}\label{eq:thin-character-bound}
 |G_\chi(y)|\leq y\exp(-c_6x^{\varepsilon/8})
 \qquad\left(y\geq\frac{\exp(x^{\varepsilon/3})}{3}\right),
\end{equation}
where $G_\chi$ is as in \ref{B1.26}, with $G_1=G_z$.
\end{theorem}

\begin{proof}[Outline]
Use \ref{B2.33} and \ref{B2.42} for the twist conductor and initial bounds. Combine real-zero exclusion with the classical region, bound the explicit formula's zero sum by the functional equation and zero counts, and use partial summation and \ref{B1.26} for tail decay. The complete proof is in \S\ref{prf:B2.52}.
\end{proof}

\section{The thick set and the top set}

In the multiplicative host coordinates,
$\tau_2(a)=((-1,1),\iota(a))$, where $\iota$ identifies $A$ with its
multiplicative copy. The thick places are $\tau_2(u_0+sa_0)$ for
$s\in T_2$, with $u_0$ the host's offset in $A$.

\begin{theorem}[A thick set without short relations]\label{B2.53}
There is an absolute $C_T>0$. Fix $0<\varepsilon<1/52$ and real
$C_o,C_L$. For all sufficiently large $x$, let a frame $\phi$ and
an integer $L\geq2$ satisfy $L\leq C_L\log x$, $o_\phi\leq C_o$,
and $x^{2+8\varepsilon}/2\leq n\leq C_o x^{2+8\varepsilon}$.
There is a finite set $T_2\subset\N$ with $4Ls<n_0$ for $s\in T_2$,
$4L|T_2|\leq n_0$, and
\begin{equation}\label{eq:thick-size-bounds}
 |T_2|\geq
 \frac{n_0}{L}\exp\left(-C_T\sqrt{\log n_0\log L}\right),
 \qquad |T_2|\geq x^{-\varepsilon}n.
\end{equation}
For distinct $s_1,s_2,s_3\in T_2$ and integers $i,j$ with
$0<|i|,|j|\leq L$ and $i\ne j$,
\begin{equation}\label{eq:no-short-relation}
 j(s_2-s_1)\ne i(s_3-s_1).
\end{equation}
\end{theorem}

\begin{proof}[Outline]
Encode integer points on a sphere in base \(4Ld\) so that a short relation forces three distinct points to be collinear. Choose the dimension and digit range for the size cutoff, then delete \(0\) and use the frame bounds for both counts. The complete proof is in \S\ref{prf:B2.53}.
\end{proof}

For $\gamma\in A$, take $\lambda_\gamma=4/(3o_\phi+1)$ on $A_0$ and
$3/(3o_\phi+1)$ elsewhere. The coset shares of $TT$ have error at most $\rho$ if
\begin{equation}\label{eq:top-shares}
 \left||TT\cap(\gamma+A_0)|-\lambda_\gamma|TT|\right|
       \leq\rho|TT|\qquad(\gamma\in A).
\end{equation}

Put $X'=\exp(x^{\varepsilon/3})$, and let $h_p\in A$ be the host
Frobenius coordinate for an unramified prime $p$. The inert-class
primes of the host are those satisfying
\[
 p\nmid Q_Sqd,\qquad
 \left(\frac{-d}{p}\right)=-1,\qquad
 \left(\frac pq\right)=-1.
\]
Let $\mathcal P_s$ consist of these primes with $p<\lceil X'\rceil$
and $\operatorname{ord}(h_p)\geq x^{\varepsilon/4}$.
Let $b_p(TT)$ count the members of $TT$ whose $h_p$-orbit fails the
good-orbit condition below. The required bound on these counts is
\begin{equation}\label{eq:bad-top-budget}
 \sum_{p\in\mathcal P_s}\frac{\log p}{p}\,b_p(TT)
 \leq4C_bx^{-2}\delta_1 n
          \sum_{p\in\mathcal P_s}\frac{\log p}{p}
\end{equation}
where $TT$ is the top set of $H$.

For a selected point $t$, let $t_0$ be the first element of
$t+\langle h_p\rangle$ in the frame's enumeration of $A$. Mark
coordinate $i$ when $t_0+ih_p\in TT$. Define
\[
 \begin{gathered}
 H_\varepsilon=2+8\varepsilon,\qquad
 K_\varepsilon=\left\lceil8(H_\varepsilon+2)/\varepsilon\right\rceil,
 \qquad R_\varepsilon=K_\varepsilon+1,\qquad u=\lceil1/\delta_1\rceil,\\
 \ell=(6R_\varepsilon+1)u,\qquad
 M_0=\left\lceil(H_\varepsilon+3)\log x\,10^{6R_\varepsilon+1}\right\rceil,
 \qquad G'=(M_0+5)\ell,\\
 N_1=1+\left\lceil\max(6\mathrm e G'\delta_1,(H_\varepsilon+3)\log_2x)\right\rceil,
 \qquad
 L_2=2\sup_{\substack{p\ {\rm prime}\\p\ {\rm inert\ class}}}
              \left\lceil\log_p((kx^{26})^4)\right\rceil .
 \end{gathered}
\]
On an orbit of length $N_{\rm orb}$, the marked set must have at most
$K_\varepsilon$ points in every cyclic window of length $L_2$ and at
most $N_1$ in every cyclic window of length $G'$. Divide the orbit into
consecutive full tiles of length $\ell$. A tile is marked when its
odd-numbered cells of length $u$ are empty and each even-numbered cell
contains exactly one marked point, of even orbit coordinate. Every
run of $M_0+1$ tiles, cyclically among the full tiles, must contain
a marked tile.

\begin{theorem}[Top positions with the required shares]\label{B2.54}
Fix $0<\varepsilon<1/52$ and $C_8\in\R$. There are $C_b>0$ and $x_0$
such that the following holds. Suppose $R$ satisfies the rank conditions at $s=(x,\varepsilon,J)$ and the primes of $H$ satisfy the size, splitting and modulus conditions at $Y$ stated before Lemma~\ref{B2.55}. If
\[
 \begin{gathered}
 x\geq x_0,\qquad J,o_\phi\leq C_8,\qquad Q_S\leq C_8n,\qquad x^{1/J}\leq Y\leq x,
 \end{gathered}
\]
then some $TT\subset A$ satisfies $\delta_1n/2\leq|TT|\leq2\delta_1n$ and the coset share bound \eqref{eq:top-shares} with $\rho=x^{-1+\varepsilon}$. For this choice of top set, the weighted counts of points on bad orbits satisfy \eqref{eq:bad-top-budget}.
\end{theorem}

\begin{proof}[Outline]
Sample independently at the two prescribed rates and use union bounds for crowded windows and independence across tiles to bound bad orbits conditional on selection. Combine the expected weighted bad count with the second moment of coset counts to choose the set. The complete proof is in \S\ref{prf:B2.54}.
\end{proof}

\section{Completing a host over the chosen field}


The size, splitting and modulus conditions at $Y$ require that $q,d$ be primes satisfying \eqref{eq:host-prime-hyp} and
\[
 \begin{gathered}
 \sqrt{X_1(Y)}<q\leq X_1(Y),\qquad
 \sqrt{X_2(Y)}<d\leq X_2(Y),\\
 \left(\frac{-d}{p}\right)_{K}=1
       \quad(p\text{ prime},\ p\leq Y^{\varepsilon/5}),\qquad
 \log(Q_Sqd)\leq2Y^{\varepsilon/5}(\log Y)^2.
 \end{gathered}
\]
The subscript $K$ denotes the Kronecker symbol, so the splitting
condition includes $p=2$. The pointwise bounds for the three sign characters require $|g_\chi(x)|$ to be at most
$x^{1/2+\varepsilon}$ and $|(L'/L)(1,\chi)|$ to be at most $C\log x$ for
every primitive character $\chi$ that agrees on units modulo $8qd$
with either sign character $(\,\cdot\,/q)$, $(\,\cdot\,/d)$ or with
their product. For the six quadratic twists, take those three characters and their products with
$\chi_8(n)=(2/n)_K$ of conductor $8$. For each of the six, every
primitive character $\chi$ that agrees with it on units modulo
$8qd$ must be real, have conductor $N_\chi\leq T$, and satisfy $L(s,\chi)\ne0$ throughout $\mathcal R(T)$. Take $T=8X_1(Y)$ for the two twists involving
$q$ alone and $T=8X_1(Y)X_2(Y)$ for the four involving $d$.

\begin{lemma}[The numerical conditions supplied by a good pair]\label{B2.55}
Fix $0<\varepsilon<1/52$ and $C_0\in\R$.
For all sufficiently large $x$, suppose $R$ satisfies the rank conditions at $s$ and $H$ is a host with $J,o_\phi\leq C_0$ and $Q_S\leq C_0n$.
If $M\geq1$, $x^{1/J}\leq Y\leq x$, and $(q,d)$ is a good pair for $(M,\varepsilon)$ on $[Y,x]$, then the size, splitting and modulus conditions at $Y$ hold. The three sign characters satisfy the pointwise bounds with $C=4$, and the six quadratic twists satisfy the stated reality, conductor and nonvanishing conditions. The host also satisfies
$q,d>4tr+1$.
\end{lemma}

\begin{proof}[Outline]
Use \ref{B2.42} and \(d\equiv7\pmod{16}\) for splitting, and rank and modulus bounds for the conductor budget. Identify the six twists by reciprocity and \ref{B2.33}, apply \ref{B2.42} at \(M\) and \(8M\) for bounds and boxes, and use the host for \(q,d>4tr+1\). The complete proof is in \S\ref{prf:B2.55}.
\end{proof}

Put $N=Q_Sqd$. The following conditions on a completed host at scale $s$ use constants $C_8,C_5,c_6,C_b$ and an auxiliary prime $q_0$.
Require the size bounds and $\mathcal B_r$-disjointness of \ref{B2.51}.
Require the thick-set conditions of \ref{B2.53} at
$L_w=10\lceil405\log x\rceil+2\lceil30/\varepsilon\rceil+22$, with
$C_T$ in \eqref{eq:thick-size-bounds} equal to the thickness constant chosen for the host. Require the top-set size bounds of Theorem~\ref{B2.54} and the weighted orbit bound \eqref{eq:bad-top-budget} with constant $C_b$.
Require \eqref{eq:thin-character-bound} with constant $c_6$ and
the pointwise bounds for the three sign characters with $C=C_5$.

The primes $p,p'$ must split completely in $E$, and all level primes
must split in $K=\Q(\sqrt{-d})$. There must be $\pi_i\in\mathcal O_K$ with
\begin{equation}\label{eq:level-norm-and-congruence}
 \pi_i\overline{\pi_i}=g_i,\qquad
 \pi_i\equiv1\pmod{Q_Sqd\,\mathcal O_E}\quad(i=0,1).
\end{equation}
Require $p,p'\equiv1\pmod{\operatorname{lcm}(15,N)}$,
$q_0\equiv1\pmod N$, $q_0\notin\{p,p'\}$, and
require the set of level primes to be exactly $\{g_0,g_1,p,p',q_0\}$.
For $\chi_d=(\,\cdot\,/d)$ require $|\Re(L'/L)(1,\chi_d)|\leq x^\varepsilon$.
Take $\mathfrak c_K=\mathcal O_K$, splitting bound at least
$\lfloor x^{\varepsilon_z/5}\rfloor$, $J,o_\phi\leq C_8$, and
$Q_S\leq C_8n$. Require a single $Y\in[x^{1/J},x]$ for which the size, splitting and modulus conditions hold, the six quadratic twists satisfy the conditions preceding Lemma~\ref{B2.55}, and the host characters satisfy the boxes of
\ref{B2.52} at height $H_{Q_{\mathrm{odd}}}(j)$ for the primitive real
host characters with conductor divisible by $q$ or $d$. Here
$j=\mathbf1_{d\mid N_\chi}$ and the conductor is at most this height.

A choice record specifies a finite level set containing four distinct
primes $g_0,g_1,p,p'$, with $p,p'\equiv1\bmod15$. Every member of
the level set is coprime to $Q_S$. The record includes a nonzero
integral ideal $\mathfrak c_K$ contained in no $m\mathcal O_K$ for
an integer $m>1$. Each prime ideal factor has norm a rational
prime that splits in $K$. The norm of $\mathfrak c_K$ is coprime
to $2,3,5,q,d,Q_S$ and to each member of the level set.
The remaining entries are finite sets $TT\subseteq A$ and
$T_2\subseteq\N$, an offset $u_0\in A$, an integer $L_w\geq2$ and
a real constant $C_{T_2}$. They satisfy $4L_ws<n_0$ for every
$s\in T_2$ and \eqref{eq:no-short-relation} with $L=L_w$.

Call the ranks $R$ admissible for $H_0$ if the following holds for every $B_0\in\N$ and every choice of the data just described. Whenever $H_0$ satisfies the field conditions at $B_0$ preceding Theorem~\ref{B2.46} and $q,d>4tr+1$, there is a host on $R$ extending $H_0$ with those choices and splitting bound $B_0$.

\begin{theorem}[Completing a fixed arithmetic datum]\label{B2.56}
Fix $0<\varepsilon<1/52$ and $C_0\geq1$. There are real constants
$C_{T_2},C_8,c_6,C_b,x_0$, with $C_8\geq C_0$, $c_6>0$ and $C_b>0$,
with the following property.
Let $R$ satisfy the rank conditions at $s=(x,\varepsilon,J)$ and be admissible for an arithmetic datum $H_0$ on its frame. Let $M\geq1$ be an integer and $Y\in\R$. Suppose $(q,d)$ is a good pair for $(M,\varepsilon)$ on $[Y,x]$, the field conditions on $H_0$ hold at $B_0=\lfloor Y^{\varepsilon/5}\rfloor$, and
\[
 \begin{gathered}
 x\geq x_0,\qquad J,o_\phi\leq C_0,\qquad Q_S\leq C_0n,\qquad M\leq C_0,\\
 Q_S\mid\operatorname{lcm}(2^{e+2},M),\qquad x^{1/J}\leq Y\leq x.
 \end{gathered}
\]
There are a host $H$ extending $H_0$ with thickness constant $C_{T_2}$ and a prime $q_0$ satisfying the completed-host conditions with constants $C_8,4,c_6,C_b$. Its top set satisfies \eqref{eq:top-shares} with $\rho=x^{-1+\varepsilon}$.
\end{theorem}

\begin{proof}[Outline]
Bound \(\log N\) via Lemma~\ref{B2.55}, then use Theorems~\ref{B2.51} and~\ref{B2.53} with admissibility to form a host. Theorem~\ref{B2.54} supplies the top set, and Lemma~\ref{B2.55} with Theorem~\ref{B2.52} verifies the character bounds. The complete proof is in \S\ref{prf:B2.56}.
\end{proof}

\begin{theorem}[Admissible ranks on a prescribed frame]\label{B2.57}
For every $0<\varepsilon<1/52$ there is $x_0$ such that for
$s=(x,\varepsilon,J)$ with $x\geq x_0$, a frame $\phi$ satisfying
$x^{2+8\varepsilon}\leq |A|\leq2x^{2+8\varepsilon}$, and any
arithmetic datum $H_0$ on $\phi$, there is a rank datum $R$ on
$\phi$ satisfying the rank conditions at $s$ and admissible for $H_0$.
\end{theorem}

\begin{proof}[Outline]
Choose ranks by \ref{B8.59}, choose a CM embedding using the field data, and identify \(\operatorname{Gal}(E/K)\) by the trivial \(d\)-sign. Choose one prime of \(K\) per conjugate pair to factor the different for \ref{B2.46} and set the splitting bound to \(B_0\). The complete proof is in \S\ref{prf:B2.57}.
\end{proof}

\section{The arithmetic conditions and the window supply}

\begin{definition}[The arithmetic requirement at a host]\label{B2.58}
Let $\mathcal L$ be a level, $\phi$ a frame, and $H$ a raw host on
$\phi$. Take $TT\subset A$ and $x,\varepsilon,C\in\R$.
Let $S=\mathcal X_{\mathcal L}$ be the retained family of Definition~\ref{B1.21}, and let $K_H(S)$ be the set of host characters whose primitive representative belongs to $S$.
Set $\mathcal B_H(S)=(X_+\setminus K_H(S))\cup X_-$.
The arithmetic requirement for $\mathcal L$ at $H$ consists of the following conditions. The index $[A:A_0]$ equals the level parameter $o_{\mathcal L}$, every $\chi\in S$ has a representative in $X_+$, and $TT$ satisfies \eqref{eq:top-shares} with $\rho=x^{-1+\varepsilon}$.
If the level designates a character, it must be the primitive representative of the $q$-sign character $\theta_q$. If the level has no designated character, the primitive representative of $\theta_q$ must lie outside $S$. For every $\theta\in\mathcal B_H(S)$ and every primitive
representative $\chi$ of $\theta$, require
$|g_\chi(x)|\leq x^{1/2+\varepsilon}$ and $|(L'/L)(1,\chi)|\leq C\log x$.

An arithmetic datum $H_0$ with field $E$ satisfies the cyclic host requirement for $\psi$ at scale $s$, top set $TT$ and auxiliary rank $t$ if its underlying field and Artin data satisfy the arithmetic requirement above for $\mathcal L_\psi$ with $C=4$, and $V_E(t)\leq\sqrt b$.
\end{definition}

Say that hosts satisfying a condition $P$ exist on windows of exponent $\lambda$ if for
each $0<\varepsilon<1/52$ there are constants $C_{T_2},C_8,C_5,c_6,C_b,x_0$,
with $C_8\geq1$ and $C_5,c_6,C_b>0$, a number $J\geq1$, and a set
$\mathcal W\subset\R$ such that
\begin{equation}\label{eq:window-set-property}
 \forall y\in\R\ \exists Z\geq y,\qquad
 [Z,Z^\lambda]\subset\mathcal W .
\end{equation}
For every scale $s=(x,\varepsilon,J)$ with $x\geq x_0$ in $\mathcal W$,
there must be ranks $R$ satisfying the rank conditions at $s$, a host $H$ with thickness constant $C_{T_2}$, and a prime $q_0$ satisfying the completed-host conditions of Theorem~\ref{B2.56} with constants $C_8,C_5,c_6,C_b$ and the additional condition $P(C_5,s,R,H,q_0)$.



\begin{theorem}[Assembling the data on one window]\label{B2.59}\label{B2.60}
Fix a primitive $\psi$ and assume the cohomological window hypothesis of Definition~\ref{B6.03} for $\psi$ with exponent $401$.
For every $0<\varepsilon<1/52$ assume the completion assertion of
\ref{B2.56} for every $C_0\geq1$. Suppose also that for each
$0<\varepsilon<1/52$ there is $x_0(\varepsilon)$ such that, at
every scale $s=(x,\varepsilon,J)$ with $x\geq x_0(\varepsilon)$,
every datum $H_0$ with $M\mid Q_S$ and \eqref{eq:cyclic-artin-kernel},
on a frame with
\[
 x^{2+8\varepsilon}\leq n\leq2x^{2+8\varepsilon},\qquad o_\phi=o_\psi,
\]
admits ranks on that frame satisfying the rank conditions at $s$ and admissible for $H_0$. Then hosts satisfying the cyclic host requirement of Definition~\ref{B2.58} for $\psi$ exist on windows of exponent $3/2$, with $C_5=4$. By \ref{B2.56}, the conclusion holds under the
window and rank-supply hypotheses.
\end{theorem}

\begin{proof}[Outline]
Theorem~\ref{B2.48} with \ref{B2.43} supplies windows with good pairs, and Theorem~\ref{B2.46} builds the frame and datum \(H_0\). Rank supply, completion (via Theorem~\ref{B2.56}), and Theorem~\ref{B2.47} then verify Definition~\ref{B2.58}. The complete proof is in \S\ref{prf:B2.60}.
\end{proof}

\begin{theorem}[Arithmetic hosts on arbitrarily distant windows]\label{B2.61}
Let $\psi$ be a primitive Dirichlet character satisfying the cohomological window hypothesis of Definition~\ref{B6.03} with exponent $401$.
For each $0<\varepsilon<1/52$, one choice of $J\geq1$ and of the host constants works on arbitrarily distant intervals $[Z,Z^{3/2}]$.
Every $x$ in those intervals admits ranks satisfying the rank conditions at $(x,\varepsilon,J)$ and a completed host with the chosen constants, $C_5=4$, and the cyclic host requirement of Definition~\ref{B2.58} for its arithmetic datum, top set and auxiliary rank.
\end{theorem}

\begin{proof}[Outline]
Apply Theorem~\ref{B2.60} with rank supply from Theorem~\ref{B2.57}, good pairs from Theorems~\ref{B2.43} and~\ref{B2.48}, and completion from Theorem~\ref{B2.56}. The hosts inherit the character bounds of Lemma~\ref{B2.55} and Theorems~\ref{B2.52} and~\ref{B2.47}. The complete proof is in \S\ref{prf:B2.61}.
\end{proof}


\part{Divisibility from one theta lift}

% ---- p3c0 ----
\chapter{Background: the Weil representation and theta lifts}\label{ch:p3c0}

\section{Lattice theta series}

\begin{theorem}[Weighted theta series]\label{B3.00a}
Let \(L\) be a positive definite even lattice of rank \(d\), with bilinear form \((\, ,\, )\) and dual \(L^\vee\). Let \(N_L\) be its level, the least positive integer such that \(N_L(y,y)/2\in\Z\) for every \(y\in L^\vee\). If \(P\) is a homogeneous harmonic polynomial of degree \(j\) on \(L\otimes\R\), then
\[
 \theta_{L,P}(z)=\sum_{x\in L}P(x)\exp\bigl(\pi i z(x,x)\bigr),
 \qquad \Im z>0,
\]
is a holomorphic modular form of weight \(d/2+j\) on \(\Gamma_0(N_L)\), with the lattice character and, in odd rank, the half-integral weight multiplier. It is cuspidal if \(j>0\). For even \(d\) the character is the Kronecker symbol \(\bigl(\frac{(-1)^{d/2}\det L}{\,\cdot\,}\bigr)\), where \(\det L\) is the Gram determinant. \cite{elkies-2019}
\end{theorem}

The transformation \(z\mapsto z+1\) uses the evenness of \(L\). For \(z\mapsto-1/z\), Poisson summation on the cosets of \(L\) in \(L^\vee\) gives the finite Fourier transform on \(L^\vee/L\). Harmonicity makes the Fourier transform of the polynomial times the Gaussian another Gaussian times the same polynomial, with the extra factor of degree \(j\) in the weight. The coset series have nonnegative exponents, and their constant terms vanish for \(j>0\) because \(P(0)=0\). Summing an acted-on Schwartz function over rational points is the corresponding construction in the Weil representation.

\section{The oscillator representation}

Let \(k\) be a local field of characteristic zero and \(S\) a symplectic \(k\)-space with form \(\langle\, ,\,\rangle\). Its Heisenberg group is \(H(S)=S\times k\), with multiplication
\[
 (s,t)(s',t')=(s+s',t+t'+\tfrac12\langle s,s'\rangle).
\]
Fix a nontrivial unitary additive character \(\psi_0\) of \(k\).

\begin{theorem}[Stone--von Neumann and Weil]\label{B3.00b}
There is, up to unitary equivalence, a unique irreducible unitary representation \(\rho_{\psi_0}\) of \(H(S)\) with central character \(\psi_0\). For \(S=X\oplus X^\vee\) with \(\langle(x,\xi),(y,\eta)\rangle=\eta(x)-\xi(y)\), its Schr\"odinger realization on \(L^2(X)\) is
\[
 (\rho_{\psi_0}(x,\xi;t)\phi)(u)
   =\psi_0\bigl(t+\xi(u)+\tfrac12\xi(x)\bigr)\phi(u+x).
\]
For each \(g\in\operatorname{Sp}(S)\) the intertwiner carrying \(\rho_{\psi_0}(s,t)\) to \(\rho_{\psi_0}(gs,t)\) is unique up to a unit scalar. These intertwiners give a projective representation of \(\operatorname{Sp}(S)\) and a genuine representation \(\omega_{\psi_0}\) of its metaplectic double cover \(\operatorname{Mp}(S)\). The nontrivial element of the covering kernel acts by \(-1\). The operators preserve the Schwartz space \(\mathcal S(X)\), with Schwartz--Bruhat functions at finite places. \cite{weil-1964,kudla-1996}
\end{theorem}

\begin{theorem}[Schr\"odinger operators]\label{B3.00c}
Use the preceding polarization. For a symmetric map \(b:X\to X^\vee\), let \(n_b(x,\xi)=(x,\xi-bx)\). For \(a\in\GL(X)\) let \(m_a(x,\xi)=(ax,\xi\circ a^{-1})\). These give the unipotent radical and Levi of the Siegel parabolic stabilizing \(X^\vee\). In dual bases let \(w(x,\xi)=(\xi,-x)\). The operators for these elements, with \(dv\) self-dual for \(\psi_0\) and the coordinate pairing, are
\[
 \begin{aligned}
 \omega_{\psi_0}(n_b)\phi(u)&=\psi_0\bigl(\tfrac12(bu)(u)\bigr)\phi(u),\\
 \omega_{\psi_0}(\widetilde m_a)\phi(u)
   &=c_{\psi_0}(\widetilde m_a)|\det a|_k^{-1/2}\phi(a^{-1}u),\\
 \omega_{\psi_0}(\widetilde w)\phi(u)
   &=c_{\psi_0}(\widetilde w)\int_X\phi(v)\psi_0(u\cdot v)\,dv .
 \end{aligned}
\]
The first line uses the canonical unipotent splitting. The scalars in the other two lines are the Weil-index factors for the chosen lifts. Changing a lift changes its scalar by a sign. \cite{weil-1964,kudla-1996}
\end{theorem}

\section{Unitary pairs and local theta lifting}

A reductive dual pair consists of reductive subgroups that are mutual centralizers in a symplectic group. Let \(E/F\) be a CM extension with involution \(z\mapsto\bar z\). For a Hermitian space \(V\) and a skew-Hermitian space \(W\) over \(E\), the trace to \(F\) of the product of the two forms defines a symplectic form on \(S=V\otimes_E W\). The tensor actions give the dual pair \((\U(V),\U(W))\); here \(\U\) denotes the isometry group.

\begin{theorem}[Splitting and local Howe duality]\label{B3.00d}
At a place \(v\) of \(F\), let \(\eta_v\) be the quadratic character of \(F_v^\times\) for \(E_v/F_v\), trivial at a split place. Choose unitary characters \(\chi_V,\chi_W\) of \(E_v^\times\) with restrictions \(\eta_v^{\dim_E V}\) and \(\eta_v^{\dim_E W}\) to \(F_v^\times\), respectively. They give splittings over the two unitary groups in the circle extension obtained from the metaplectic cover by extending its central kernel \(\{\pm1\}\) to the unit circle. The splitting over \(\U(V)\) uses \(\chi_W\), and that over \(\U(W)\) uses \(\chi_V\). The resulting Weil representation is a representation of the product of the two unitary groups. \cite{kudla-1994,gan-takeda-2016}

For an irreducible admissible representation \(\pi\) of \(\U(V)(F_v)\), the maximal quotient of this Weil representation on which \(\U(V)(F_v)\) acts by copies of \(\pi\) has the form \(\pi\otimes\Theta(\pi)\). At archimedean places this is understood on Harish-Chandra modules. If \(\Theta(\pi)\ne0\), its maximal semisimple quotient \(\theta(\pi)\) is irreducible. Distinct \(\pi\) with nonzero small lifts have distinct small lifts. For the fixed pair and splitting data the correspondence is symmetric in the two groups. \cite{howe-1989,waldspurger-1990,gan-takeda-2016,prasad-1996}
\end{theorem}

Howe established the archimedean duality. Over nonarchimedean fields Waldspurger proved it in odd residue characteristic, and Gan and Takeda proved it for symplectic--orthogonal and unitary pairs over quadratic field extensions in arbitrary residue characteristic. At a split place the unitary pair becomes a general linear pair. Definite real unitary groups are compact, and their types can be read in the Fock realization of the same Weil representation.

\section{Adelic theta lifts and their periods}

Write \(\mathbb A=\mathbb A_F\) and \([G]=G(F)\backslash G(\mathbb A)\). Choose a nontrivial character of \(\mathbb A/F\) and global splitting characters on \(E^\times\backslash\mathbb A_E^\times\) with the prescribed quadratic restrictions. With these data the restricted tensor product of the local Weil representations and the rational splitting give the kernel
\[
 K_\phi(g,h)=\sum_{x\in X(F)}(\omega(g,h)\phi)(x),
 \qquad \phi\in\mathcal S(X(\mathbb A)),
\]
where \(X\) is a rational Lagrangian in the ambient symplectic space. It is automorphic in \(g\) and \(h\). For a cusp form \(f\) on \([G]\), the lift to the other group is the automorphic function
\[
 T_\phi f(h)=\int_{[G]}K_\phi(g,h)\overline{f(g)}\,dg.
\]
The integral is absolutely convergent. For an irreducible cuspidal representation \(\pi\), its lift is the span as \(f\in\pi\) and \(\phi\) vary. \cite{weil-1964,prasad-1996}

\begin{theorem}[Tower property]\label{B3.00e}
Fix the source and a target Witt tower \(V_r=V_0\oplus\mathbb H_E^r\), with \(V_0\) anisotropic and \(\mathbb H_E\) a hyperbolic plane. Use compatible splitting data. For \(1\le j\le r\), the constant terms along the parabolic stabilizing a \(j\)-dimensional isotropic subspace, restricted to the \(\U(V_{r-j})\) factor of its Levi, span the lift to \(\U(V_{r-j})\). The test function for this constant term is the restriction at zero in the variables of the removed isotropic part in a mixed Schr\"odinger model. Thus the first nonzero lift in the tower is cuspidal, and the lifts above it are nonzero. On an anisotropic target every automorphic lift is cuspidal. \cite{rallis-1984,prasad-1996,gan-qiu-takeda-2014}
\end{theorem}

For the norm of a lift from \(\U(A)\) to \(\U(B)\), the doubled group is \(\U(A\oplus A^-)\), where \(A^-\) has the negative form. Its diagonal isotropic subspace has dimension \(a=\dim_E A\). Write \(b=\dim_E B\), and let \(r_B\) be the rational Witt index of \(B\). The doubled Weil representation pairs this split group with \(\U(B)\). Evaluation at zero, \(g\mapsto(\omega_d(g,1)\phi)(0)\), determines the Siegel--Weil section at \(s_0=(b-a)/2\). Denote its flat extension by \(F_\phi(s)\). It is a section of normalized Siegel induction from \(\chi_B(\det t)|\det t|_{\mathbb A_E}^{s}\) on the Levi, with the unnormalized exponent \(s+a/2\). Let \(\mathcal E(g,s,F_\phi)\) be the Eisenstein series of this section.

\begin{theorem}[Siegel--Weil in the convergent and definite ranges]\label{B3.00f}
If \(b-r_B>a\), the doubled theta integral converges absolutely and \(\mathcal E(g,s,F_\phi)\) is holomorphic at \(s=s_0\). With the rational additive covolume equal to one,
\[
 \mathcal E(g,s_0,F_\phi)
 =\frac{1}{\operatorname{vol}[\U(B)]}\int_{[\U(B)]}K_\phi(g,h)\,dh ,
\]
where the kernel uses the doubled Weil representation \cite{weil-1965,yamana-2013}. For definite \(B\) with \(1\le b\le a\) (and \(a\ge1\)), the Eisenstein series is holomorphic at \(s_0=(b-a)/2\) and its value is twice this average \cite[Theorem~7.1 and \S3.1]{gan-qiu-takeda-2014}.
\end{theorem}

Weil proved the first identity in the range \(b>2a\), and Yamana established it throughout the convergent range \(b-r_B>a\) \cite[Theorem~2.2]{yamana-2013}; this identity is used in \ref{B8.13} with the additional bound \(b\le2a\). For definite \(B\) and \(1\le b\le a\), the point \(s_0=(b-a)/2\le0\) lies in the second-term range; this case is the anisotropic unitary second-term identity of Gan, Qiu and Takeda \cite[Theorem~7.1 and \S3.1]{gan-qiu-takeda-2014} and is restated in Theorem~\ref{B8.12}. In general divergent ranges the regularized Siegel--Weil formula relates Laurent coefficients of the Siegel Eisenstein series to those of the regularized theta integral; Kudla and Rallis established the regularized first-term identity for orthogonal--symplectic pairs, and Gan, Qiu and Takeda proved the second-term identity and inner product formula including unitary pairs \cite{kudla-rallis-1994,gan-qiu-takeda-2014}.

\begin{theorem}[Rallis inner product in the convergent unitary case]\label{B3.00g}
Suppose \(a<b\le2a\), \(b-r_B>a\), and all lower lifts of the irreducible cuspidal representation \(\pi\) of \(\U(A)(\mathbb A)\) in the target Witt tower vanish. Take \(f_1,f_2\in\pi\) and Schwartz vectors \(\phi_1,\phi_2\) whose doubled vector \(\delta\) is \(K\)-finite. In compatible Schr\"odinger coordinates for the doubled space, \(\delta\) is \(\phi_1\boxtimes\overline{\phi_2}\), with the second coordinate identified using the negative form. For the doubling embedding \(\iota\) use the convention in which the doubled kernel at \(\iota(g_1,g_2)\) has the factor \(\chi_B(\det g_2)^{-1}\) times the product of the first kernel and the conjugate of the second. Set
\[
 Z(s)=\int_{[\U(A)]^2}
 \mathcal E(\iota(g_1,g_2),s,F_\delta)
 \,\overline{f_1(g_1)}f_2(g_2)\chi_B(\det g_2)\,dg_1\,dg_2 ,
\]
initially in a right half-plane. Both lifts belong to \(L^2([\U(B)])\), and the constant Laurent coefficient of the continued integral satisfies
\[
 \operatorname{CT}_{s=s_0} Z(s)
   =\frac{\langle T_{\phi_1} f_1,T_{\phi_2} f_2\rangle_{[\U(B)]}}
           {\operatorname{vol}[\U(B)]}.
\]
For factorizable data the unfolded integrals are local matrix-coefficient integrals. In the standard doubling normalization the unramified integrals give the twisted standard \(L\)-factor at \(s+1/2\) divided by explicit Hecke \(L\)-factors. \cite{rallis-1987,gan-qiu-takeda-2014}
\end{theorem}

\begin{theorem}[Seesaw identity]\label{B3.00h}
Let \((G_1,G_2)\) and \((H_1,H_2)\) be dual pairs with \(H_1\subseteq G_1\) and \(G_2\subseteq H_2\). Use compatible Weil data so that the two kernels restrict to a common kernel on \([H_1]\times[G_2]\). For forms \(f\) on \([G_2]\) and \(u\) on \([H_1]\), suppose the common kernel times \(\overline f\,\overline u\) is absolutely integrable on that product. With the lift convention above,
\[
 \int_{[H_1]} (T_{G_2\to G_1}f)(h)\,\overline{u(h)}\,dh
 =
 \int_{[G_2]} (T_{H_1\to H_2}u)(g)\,\overline{f(g)}\,dg .
\]
The two sides use the same test vector under the common kernel identification. \cite{kudla-1996,prasad-1996}
\end{theorem}

\section{Use in the proof}

Part III uses the Weil model for the dependence on a lattice indicator in \ref{B3.35a} and the Hecke relation in \ref{B3.36}. Taylor expansion of the corresponding theta sections in formal-group coordinates writes the local kernels as \(\exp(-\Phi)U\) with the CM pairing phase and a factor integral after logarithmic substitution, as in \ref{B3.05} and \ref{B8.56c}. Part IV turns these kernel coefficients into lift bounds in \ref{B4.10a} and whole-column bounds in \ref{B4.09}. In Part VIII, Poisson summation reappears in \ref{B8.10e}, lower lifts and local matching enter \ref{B8.11} and \ref{B8.16}, and the doubled identities give \ref{B8.12} and \ref{B8.13}. The seesaw in \ref{B8.15} exchanges a lift against a period and supplies test functionals that separate the rows.


% ---- p3c1 ----
\chapter{The local kernel and the differential operator}\label{ch:p3c1}

Fix a prime $p$. For a multi-index $n$, write $|n|=\sum_x n_x$ and $n!=\prod_x n_x$. Write $[Z^n]f$ for the coefficient of $Z^n$ in $f$. Each series substituted for an indeterminate has zero constant term.

\section{Local data and weight diagrams}

\begin{definition}[The Weyl weight lattice \(H_j(L)\)]\label{B3.01a}
Let \(V\) be a finite set of places, and let \(\iota:O\to K\) be an injective homomorphism from a discrete valuation ring \(O\) into a field \(K\). Fix integers \(k_1,t\in\N\) and put \(N_j=k_1+1+j(t+1)\) for \(j\ge0\). At each \(v\in V\), let \(n'_v\ge2\) be even, let \(W_v\) be a discrete valuation ring with fraction field \(Q_v\) and an automorphism \(\sigma_v\) equipped with an embedding \(\iota_v:Q_v\to K\) carrying \(W_v\) into \(\iota(O)\), and index the \(n'_v\) embeddings \(\iota_v\circ\sigma_v^s\) by cyclic positions \(s\in\{0,\dots,n'_v-1\}\). Fix nonincreasing initial weights \(\lambda_{v,s}=(\lambda_{v,s,1}\ge\cdots\ge\lambda_{v,s,k_1+1})\in\Z^{k_1+1}\) and shifts \(w_{v,l}(s)\in\Z\) (\(l\ge1\)) with \(w_{v,1}(s)\ge-\lambda_{v,s,k_1+1}\) and \(w_{v,l}(s)\le w_{v,l+1}(s)\).

At level \(j\ge0\) and block \((v,s)\), append \(t+1\) copies of \(-w_{v,b}(s)\) to \(\lambda_{v,s}\) for each \(b=1,\dots,j\) in order; if \(s\) is odd, reverse the resulting list of length \(N_j\) and negate its entries. Denote the last entry of the resulting nonincreasing list by \(d_{j,v,s}\), and subtract \(d_{j,v,s}\) from every entry to form a partition \(\kappa_{j,v,s}\). In an \(N_j\times N_j\) matrix \(Y_{v,s}\) of indeterminates, let \(H_{\kappa_{j,v,s}}(Y_{v,s})\) be the product of the initial principal minors of the column heights of \(\kappa_{j,v,s}\). For rows \(a,b\in\{1,\dots,N_j\}\) and \(r\ge0\), define the divided-power operator \(E_{a,b}^{(r)}\) on \(O[Y_{v,s}]\) by extracting the \(T^r\)-coefficient of \(P(Y_{v,s}+TE_{b,a}Y_{v,s})\), where \(E_{b,a}\) is the matrix unit, so that the substitution adds \(T\) times row \(a\) to row \(b\). Define the unframed Weyl lattice \(\Delta_O(\kappa_{j,v,s})\subset O[Y_{v,s}]\) as the \(O\)-span of finite words in the operators \(E_{a,b}^{(r)}\) applied to \(H_{\kappa_{j,v,s}}\).

For a tuple \(L=(L_v)_{v\in V}\) of free \(W_v\)-lattices \(L_v=G_v W_v^{N_j}\subset Q_v^{N_j}\) with \(G_v\in\GL_{N_j}(Q_v)\), let \(A_{v,s}=(\iota_v\circ\sigma_v^s)(G_v)\in\GL_{N_j}(K)\), and let \(\Delta_O(\kappa_{j,v,s};A_{v,s})\subset K[Y_{v,s}]\) consist of the polynomials \(F(Y_{v,s})=F_0(A_{v,s}^{-1}Y_{v,s})\) with \(F_0\in\iota(\Delta_O(\kappa_{j,v,s}))\). The Weyl weight lattice \(H_j(L)\subset K[(Y_{v,s})_{v,s}]\) is the \(O\)-span (via \(\iota\)) of the products
\[
 \left(\prod_{v,s}(\det A_{v,s})^{-d_{j,v,s}}\right)\prod_{v,s}F_{v,s}(Y_{v,s}),
 \qquad F_{v,s}\in\Delta_O(\kappa_{j,v,s};A_{v,s}).
\]
Changing \(G_v\) on the right by \(\GL_{N_j}(W_v)\) multiplies \(A_{v,s}\) on the right by an element of \(\GL_{N_j}(\iota(O))\), which has determinant in \(\iota(O^\times)\) and preserves \(\iota(\Delta_O(\kappa_{j,v,s}))\), so \(H_j(L)\) depends only on \(L\).
\end{definition}

\begin{definition}[Local systems and lattice bounds]\label{B3.01}
Let $V$ be a finite set of places. Let $O$ be a discrete valuation ring with uniformizer $\varpi$ and coefficient map $\iota:O\to K$ to a field. At level $j$ let $\mathcal M_j$ be the $K$-space of polynomial-valued functions on components and tuples of local submodules. Its $O$-submodule $\Lambda_j$ consists of functions supported on lattice tuples with $F(L)\in H_j(L)$, the Weyl weight lattice of Definition~\ref{B3.01a}, and $F(\gamma L)=\chi(\gamma)\gamma F(L)$ for the component's level group and character. Take $K$-linear lifts $K_j:\mathcal M_j\to\mathcal M_{j+1}$.

Fix $k_1,t,c,m,r,e_O\in\N$ and $d_G:\N\to\N$. Require $p$ to be associated to $\varpi^{e_O}$. At each lattice tuple in levels $0\le j\le m$, require a finite stabilizer $\Gamma_{j,L}$ with $v_p(|\Gamma_{j,L}|)\le d_G(j)$. For $f\in\mathcal M_j$ and $g\in\R$, write $v_j(f)\ge g$ if $f\in\iota(\varpi)^{\lceil e_Og\rceil}\Lambda_j$.

At each $v\in V$ require $p\ge5$ along with the bounds $m,r\ge1$ and $2r\le k_1+1$. Use cyclic positions $s\bmod n'_v$ with $n'_v\ge2$ even and marked set $T_v$. The integers $\lambda_{v,s,1}\ge\cdots\ge\lambda_{v,s,k_1+1}$ have absolute value at most $c$, except for the prescribed first and last entries $a_v(s),-b_v(s)$ at marked positions. The integer shifts satisfy
\[
 w_{v,1}(s)\ge\max\{c,-\lambda_{v,s,k_1+1}\},\qquad
 w_{v,l}(s)\le w_{v,l+1}(s)\quad(l\ge1).
\]
Use the trace-dual bases of \ref{B3.08}.
\end{definition}

\begin{definition}[Layer weights and remaining columns]\label{B3.02}
Suppress $v$. Put $N_j=k_1+1+j(t+1)$. The list $\nu_{j,s}$ starts with $\lambda_{s,i}+w_j(s)$ for $1\le i\le k_1+1$, followed by $t+1$ copies of $w_j(s)-w_b(s)$ for each $b=1,\ldots,j$ in order. Set $\nu_{j,s,i}=0$ off this list. Put
\[
 \alpha_s=\mathbf1_{s\in T}(\lambda_{s,1}-\lambda_{s,2}),\qquad
 C_{j,s}=\nu_{j,s,1}-\alpha_s,\qquad L_{j,s}=\sum_i\nu_{j,s,i}-\alpha_s.
\]
For $j\ge1$, the rows form a partition by \ref{B3.01}. Deleting $\alpha_s$ height-one columns leaves $C_{j,s}$ columns and $L_{j,s}$ boxes.
\end{definition}

\begin{definition}[Digit loss]\label{B3.03}
Set $\ell_p(z)=0$ for $z\le0$ and $\ell_p(z)=1+\lfloor\log_p z\rfloor$ for $z>0$. For $1\le j\le m$, put $r_0=N_{j-1}$ and $r_1=N_j$. The digit loss is
\[
 \delta_j=\sum_s\bigl(r_1\ell_p(\alpha_s)+r_0r_1\ell_p(L_{j,s})\bigr)
 +\sum_{s\in T}\ell_p(\alpha_s)^2.
\]

\end{definition}

\section{Formal parameters and admissible kernels}

\begin{definition}[Parameters and phase equivalence]\label{B3.04}
Let $W=W(\overline{\F}_p)$ and $K_0=\operatorname{Frac}(W)$. Write $\sigma$ for Witt Frobenius. Set
\[
 \ell(T)=\sum_{b\ge0}p^{-b}T^{p^{2b}},\quad
 f^\ell(X)=f((\ell(X_i))_i),\quad
 [X^n]\exp(f)=\sum_{k=0}^{|n|}\frac{[X^n]f^k}{k!}\quad(f(0)=0).
\]
Write $f\simeq g$ for $f,g\in K_0[[X]]$ if $(f-g)(0)=0$ and $\exp((f-g)^\ell)\in W[[X]]$. Use $f^\ell$ also over extensions of $K_0$.
\end{definition}

\begin{definition}[The local theta condition]\label{B3.05}
The pairing laws for $b\in XYK_0[[X,Y]]$ are additivity in each variable modulo $\simeq$ and $b(\eta X,Y)\simeq b(X,\sigma(\eta)Y)$ for $\eta\in W$ with $\sigma^2\eta=\eta$. Put
\[
 E_v=\{0,\ldots,n'_v/2-1\}\times\{0,\ldots,N_j-1\}
 \times\{0,\ldots,N_{j+1}-1\},\qquad E_A=\coprod_v E_v.
\]
The field $K_c$ contains $K$ and $K_0$. Its valuation ring $O_c$ contains $W$ and satisfies $K\cap O_c=\iota(O)$. For $z\in K_c^\times$ require $p^n/z\in O_c$ for some $n\in\N$, and if $z^{-1}\notin O_c$, require $z^a/p\in O_c$ for some $a\in\N$.

For each source $L$ at level $j$ and target $L'$ at level $j+1$, including their components, the kernel data specify a set of frames from $L$ to $L'$. In the arithmetic model a frame consists of a basis of each lattice at every place. For a frame $\phi$, let $\Theta_{A,\phi}\in K_c[[X,Y]]$ be the generator on $E_A$ and $\Theta_v$ its restriction by setting the other variables to zero. The theta condition on either series, with index set $E$, requires the pairing laws and a factorization with $\Phi\in K_0[[X,Y]]$,
\[
 \Theta=\exp(-\Phi)U,\qquad U^\ell\in O_c[[X,Y]],\qquad
 \Phi(-X,-Y)\simeq\Phi,\qquad \Phi\simeq\sum_{e\in E}b(X_e,Y_e).
\]
Let $\mathcal T_j$ substitute the unscaled forms of \ref{B3.08}. The data include a matrix $B_\phi$ and a fixed-coordinate series satisfying
$\Theta^F_{j,L,L'}=(\mathcal T_j\Theta_{A,\phi})(B_\phi Z)$
for every $\phi$.
\end{definition}

\begin{lemma}[Integral derivatives of the CM pairing]\label{B3.05b}
For every prime $p$, the pairing laws of \ref{B3.05} imply $m!n![X^mY^n]b\in W$ for all $m,n\ge0$.
\end{lemma}

\begin{proof}[Outline]
The product and chain rules give integral derivatives at zero for a phase equivalent to zero. Differentiate the additivity defects for exponents at least two and use the conjugation defect and the unit \(\eta-\sigma(\eta)\) for \([XY]b\). The complete proof is in \S\ref{prf:B3.05b}.
\end{proof}

\begin{definition}[The rational weight space at level \(j\)]\label{def:p3c1-rational-weight}
At each \((v,s)\), append \(t+1\) copies of \(-w_{v,l}(s)\) to the initial weight \(\lambda_{v,s}\) for each \(l=1,\ldots,j\), in order. At odd \(s\), reverse the resulting list and negate its entries. Subtract the last entry from every entry of the resulting weight to obtain a partition \(\kappa_{v,s}\).

In an \(N_j\times N_j\) matrix of variables for this block, take the polynomials over \(K\) that are homogeneous of degree \(\kappa_{v,s,a}\) in column \(a\) for each \(a\), and are invariant under adding a multiple of any earlier column to a later one. The level-\(j\) rational weight space is the \(K\)-span of products of one such polynomial from each block, in their separate matrix variables.
\end{definition}

\begin{definition}[The arithmetic differential operator]\label{B3.05aa}
Keep the data of Definitions~\ref{B3.01}--\ref{B3.05}. For \(j\ge0\), put \(n=N_j\) and \(n'=N_{j+1}\), write \(Y=(Y_{v,s})_{v,s}\) and \(Y'=(Y'_{v,s})_{v,s}\) for the \(n\times n\) source and \(n'\times n'\) target variable matrices, and introduce an \(n\times n'\) matrix \(d_{v,s}\) of symbol variables in each block \((v,s)\), with \(d=(d_{v,s})_{v,s}\). In block \((v,s)\), put \(c=\nu_{j+1,s,1}\), \(e=\nu_{j+1,s,n}\), and \(b=\sum_{a=1}^{n'}\nu_{j+1,s,a}\).
\begin{enumerate}
\item Define a \(K\)-algebra homomorphism from \(K[Y]\) to \(K_c[Y'][d]\) by substituting, in each block \((v,s)\), the \(n\times n\) initial-column submatrix \(D_{v,s}=(d_{v,s}Y'_{v,s})_{:,1:n}\) for \(Y_{v,s}\) when \(s\) is even, and \(\operatorname{adj}T_{v,s}\) for \(Y_{v,s}\) when \(s\) is odd, where \(T_{v,s}=(\operatorname{adj}Y'_{v,s})_{\mathrm{last}\ n\ \mathrm{rows}}d_{v,s}^{\mathsf t}\).
\item Form the rational monomial \(\prod_{v,s}M_{v,s}=N_+/N_-\) by setting \(M_{v,s}=(\det D_{v,s})^e\) for even \(s\) and \(M_{v,s}=(\det Y'_{v,s})^{c-b}(\det T_{v,s})^{c-nc+b}\) for odd \(s\), placing factors with nonnegative exponent in the polynomial numerator \(N_+\in K_c[Y'][d]\) and factors with negative exponent (with negated exponent) in the polynomial denominator \(N_-\in K_c[Y'][d]\). Evaluating at \(Y'_{v,s}=I_{n'}\) with \(d_{v,s}^{\mathsf t}\) the inclusion of the first \(n\) rows (for \(s\) even) or last \(n\) rows (for \(s\) odd) gives \(\det D_{v,s}=1\) and \(\det T_{v,s}=1\), so \(N_-\ne0\) in the integral domain \(K_c[Y'][d]\) and multiplication by \(N_-\) is injective. Choose a \(K\)-linear left inverse \(\Pi_j:K_c[Y'][d]\to K_c[Y'][d]\) to multiplication by \(N_-\), and define the symbol \(R_G(d)\in K_c[Y'][d]\) of \(G\in K[Y]\) as \(\Pi_j\) applied to \(N_+\) times the image of \(G\) under step~(1).
\item Let \(Z=(Z_{v,s,a,a'})_{v,s,\,0\le a<n,\,0\le a'<n'}\) be the kernel variables matching the entries of \(d\). For \(G\in K[Y]\) and \(f\in K_c[[Z]]\), the arithmetic differential operator \(\mathcal D_j:K[Y]\times K_c[[Z]]\to K_c[Y']\) is the \(K\)-bilinear map
\[
 \mathcal D_j(G,f)=
 \sum_{\substack{\alpha\in\operatorname{supp} R_G\\
 \forall v,s:\ \sum_{a,a'}\alpha_{v,s,a,a'}=\sum_{a=1}^{n'}\nu_{j+1,s,a}}}
 [d^\alpha]R_G\,\alpha![Z^\alpha]f.
\]
\end{enumerate}
\end{definition}

\begin{definition}[Admissibility of the kernel and lifts]\label{B3.05a}
At each level $0\le j\le m$, fix a finite set of component--lattice pairs containing one representative of each level-group orbit in each component. The kernel data give integer lengths $\ell_{v,j}(L)$ on tuples. A lower edge of index $i$ at $v$ decreases $\ell_{v,j}$ by $i$ and preserves the other lengths. The data also give a $K$-bilinear polynomial-valued operator $\mathcal D_j(G,f)$ on the ambient source polynomials and $K_c[[Z]]$. It depends on $f$ through the coefficients of degree $\sum_a\nu_{j+1,s,a}$ in each block $(v,s)$. Set
\[
 \Delta_v(l)=\sum_{s\ {\rm even}}w_{v,l}(s)-\sum_{s\ {\rm odd}}w_{v,l}(s),\qquad
 D_j(L,L')=\prod_v p^{\Delta_v(j+1)(\ell_{v,j}(L)-\ell_{v,j+1}(L'))},\qquad
 \mathcal K_{j,L,L'}(G)=D_j(L,L')\,\mathcal D_j(G,\Theta^F_{j,L,L'}).
\]
For successive levels with $0\le j<m$, the following conditions use lattice pairs as sources and targets. A framed condition applies to every frame unless it specifies a choice. Write $V_j=K\Lambda_j$.
\begin{enumerate}
\item The prime-only clause is the Dieudonn\'e--Dwork criterion for every odd prime, proved in Lemma~\ref{B3.06}. In every frame, the generator $\Theta_{A,\phi}$ and each restriction $\Theta_v$ satisfy the theta condition of Definition~\ref{B3.05} on $E_A$ and $E_v$, respectively. The fixed series $b$ satisfies the pairing laws of Definition~\ref{B3.05} for every choice of the place set, including the empty set.
\item For every column choice $I$, all coefficients of the polynomial in \eqref{eq:p3c3-product-evaluation} at $\mathbf i=\mathbf0$ belong to $O_c$. The lift satisfies $K_j(V_j)\subseteq V_{j+1}$.
\item At each level $j<m$ there is a scalar $c_K\in\{0\}\cup O_c^\times$ such that, for every $F\in V_j$ and every target lattice pair $L'$, $(K_jF)(L')$ is the sum of differential kernels over source orbit representatives in \eqref{eq:p3c3-class-expansion} if $j+1<m$, and the double sum in \eqref{eq:p3c3-terminal-expansion} if $j+1=m$. At the terminal step the finite index sets and source pairs $L_{xy}$ serve for all $F\in V_j$ and targets, as do the positive integer denominators $w_x$ with $v_p(w_x)\le d_G(j)$. With these same terminal data, summing the kernel over source neighbours of index $i$ gives the action of the source Hecke operator on $F$, as in \eqref{eq:p3c3-hecke-adjoint}, for every $v$ and $i\ge0$.
\item For each $v$ and $j<m$ there are units $a,a'\in O^\times$ such that \eqref{eq:p3c3-indicator} expresses the fixed-coordinate theta series through the indicator of $xL_v\subseteq L'_v$ for every lattice pair $L,L'$, using the given lengths. The index $\omega=\omega_{v,j}(L,L')$ of the additive map in this formula depends only on the two components and the lattice entries away from $v$. For each such index, $\Xi_\omega$ maps the additive group of rational-valued functions on the $N_{j+1}\times N_j$ matrices over the local field to $K_c[[Z]]$. The differential kernel is invariant under a change of source representative. For every element $\gamma$ of the source level group and every $G$ in the rational weight space of Definition~\ref{def:p3c1-rational-weight},
\[
 \mathcal K_{j,\gamma L,L'}(\chi(\gamma)\gamma G)
 =\mathcal K_{j,L,L'}(G).
\]

\item At each level $0\le j\le m$, fix a place $v$ and an edge index $i\ge0$. For every $K_c$-valued function on these edges invariant under the simultaneous level-group action on both endpoints, summation by the source orbit and summation by the target orbit give the same result when each representative is weighted by the reciprocal of its stabilizer order, as in \eqref{eq:p3c3-orbit-adjoint}.
\item For an abstract operator $\mathcal D_j$ require the following frame condition, with length factor $D_j(L,L'')$. Lemma~\ref{B3.10a} proves it for the arithmetic operator of Definition~\ref{B3.05aa}. For both the test coefficients and the scalars in the target lattice expansion, interpret a lower bound $g$ as in Definition~\ref{B3.12}. For each $L,L'$ and each simultaneous lower neighbour $L''$ of $L'$ with indices $0\le i_v\le r$, one frame from $L$ to $L''$ must give this implication for all $f\in K_c[[Z]]$, $g\in\R$ and $G\in H_j(L)$.
\item The operators $y_v$ of Definition~\ref{B3.35} commute pairwise on $\mathcal M_0$, and each maps $\Lambda_0$ into itself. For every $0\le j<m$ and every $v$, all coefficients in the single-place test of Definition~\ref{B3.12} belong to $O_c$.
\end{enumerate}
\end{definition}

\begin{lemma}[Dieudonn\'e--Dwork criterion]\label{B3.06}
Let $p$ be odd and let $d\ge0$. For $\Psi\in K_0[[X_1,\ldots,X_d]]$ with $\Psi(0)=0$,
\[
 \exp(\Psi)\in W[[X]]
 \quad\Longleftrightarrow\quad
 (\sigma\Psi)(X_1^p,\ldots,X_d^p)-p\Psi(X)\in pW[[X]].
\]
\end{lemma}

For odd $p$ this extends the one-variable criterion over $\Q_p$ to finitely many variables with Witt Frobenius. The one-variable criterion also holds at $p=2$ \cite[Lemma 1]{krattenthaler-2015}.

\begin{proof}[Outline]
Legendre's formula identifies \(p(X)W[[X]]\) with \(1+p(X)W[[X]]\), rewriting the condition as \(\varphi F/F^p\in 1+p(X)W[[X]]\) for \(F=\exp\Psi\). Necessity holds modulo \(p\); sufficiency inducts on degree \(h\ge 1\) by dividing \(F\) by its integral truncation below \(h\). The complete proof is in \S\ref{prf:B3.06}.
\end{proof}

\begin{definition}[An expansion with controlled denominators]\label{B3.07}
Let $R'\subset K'$ with $K'/K_0$ a field extension. Use variables $Z_{s,a,b}$ indexed by $s\bmod n'$ with $0\le a<r_0$ and $0\le b<r_1$. Put
\[
 S_{s,a,b}=\sum_{l\ge0}(-1)^l Z_{s-2l,a,b}^{p^{2l}},\qquad
 H_u=\sum_{s,a,b}u_s S_{s,a,b}^p S_{s+1,a,b}\quad(u_s\in K').
\]
For $f,H\in K'[[Z]]$, write $f\in\mathcal C_{R'}(H)$ if there are $U_i\in R'[[Z]]$ with $[Z^n]f=\sum_{i=0}^{|n|}p^{-i}[Z^n](H^iU_i)$ for every $n$. The normed-field version requires $\|[Z^n]U_i\|\le1$ for all $i,n$.
\end{definition}

\section{Reduction of the phase}

\begin{theorem}[The explicit odd-power phase]\label{B3.08}
Let $p\ge5$ and let $b$ satisfy the pairing laws of \ref{B3.05}. There is $c\in W$ with $\sigma^2c=c$ such that $\exp(b^\ell-G_c)\in1+(X,Y)W[[X,Y]]$, where
\[
 A(X,Y)=\sum_{0\le i\le l}\frac{X^{p^{2i}}Y^{p^{2l+1}}}{p^{i+l+1}},
 \qquad G_c=cA(X,Y)-\sigma(c)A(Y,X).
\]
For even $n'\ge2$ and $r_0,r_1\ge0$, take bases $\omega_r,\omega_r^*$ of $W(\F_{p^{n'}})$ over $W(\F_{p^2})$, indexed by $0\le r<n'/2$, with $\sum_{h=0}^{n'/2-1}\sigma^{2h}(\omega_r\omega_{r'}^*)=\delta_{r,r'}$. In the variables of \ref{B3.07} set
\[
 L^X_{r,a,b}=\sum_{s\ {\rm even}}\sigma^s(\omega_r^*)Z_{s,a,b},
 \qquad L^Y_{r,a,b}=\sum_{s\ {\rm odd}}\sigma^s(\omega_r)Z_{s,a,b}.
\]
Let $K'/K_0$ contain $\pi$ with $\pi^{p^2-1}=p$. Put $R_\pi=W[\pi]$ and $\mu(T)=\ell^{-1}(\pi T)/\pi$. Then $\mu\in T\Z[[T]]$ and $\pi^{(p+1)i}/i!\in R_\pi$ for $i\ge0$. For every $\Phi\simeq\sum_e b(X_e,Y_e)$ with $e=(r,a,b)$ as above,
\[
 \exp(-\Phi(\pi L^X,\pi L^Y))=\exp(-\pi^{p+1}H_u/p)V,\qquad
 V\in1+(Z)R_\pi[[Z]],\qquad
 u_s=\begin{cases}-\sigma(c)&s\ {\rm even},\\ c&s\ {\rm odd}.\end{cases}
\]
\end{theorem}

\begin{proof}[Outline]
Use Hazewinkel's lemma, root-of-unity averaging, and coefficient induction to identify the integral Frobenius difference with \(D_\sigma(B-G_c)\) and apply \ref{B3.06}. Trace duality gives \(H_u\), and Legendre's formula makes the remaining exponential an integral unit. The complete proof is in \S\ref{prf:B3.08}.
\end{proof}

\section{Kernel expansions}

\begin{theorem}[Single-place and simultaneous expansions]\label{B3.09}
Assume \ref{B3.05a} and fix $0\le j<m$. Fix a lattice pair and a frame. Let $K'/K_c$ contain $\pi$ with $\pi^{p^2-1}=p$, and let $R'\subset K'$ contain $W[\pi]$ and $O_c$. Write $Z_v$ for the variables at $v$ and $Z=(Z_v)_v$.

At each place one can choose $u_{v,s}\in R'$ so that $\Theta_v(\pi L_v^X,\pi L_v^Y)\in\mathcal C_{R'}(H'_v)$ whenever $H'_v=H_{u_v}+pG_v$ with $G_v\in(Z_v)R'[[Z_v]]$. For the generator one can choose all $u_{v,s}\in R'$ so that for every family $G_v\in(Z)R'[[Z]]$, with $H'_v=H_{u_v}(Z_v)+pG_v(Z)$, there are $U_{\mathbf i}\in R'[[Z]]$ for $\mathbf i:V\to\N$ such that
\[
 [Z^n]\Theta_A(\pi L^X,\pi L^Y)
 =\sum_{\substack{\mathbf i:V\to\N\\i_v\le |n|\ \forall v}}
 [Z^n]\left(U_{\mathbf i}\prod_v p^{-i_v}(H'_v)^{i_v}\right)
 \qquad\text{for every }n .
\]
\end{theorem}

\begin{proof}[Outline]
Use \ref{B3.05} and \ref{B3.04} for the zero phase at the empty set, then apply \ref{B3.08} place by place and absorb the phase correction into the integral cofactor. Expand using the factorial bound of \ref{B3.08} and the zero constant terms of the \(H'_v\). The complete proof is in \S\ref{prf:B3.09}.
\end{proof}

\section{The polynomial and its coefficients}

\begin{definition}[Products of column minors]\label{B3.10}
For $j\ge0$ put $r_0=N_j$, $r_1=N_{j+1}$. Use $Z_{v,s,a,b}$ with $0\le a<r_0$, $0\le b<r_1$, and auxiliary variables $\tau_{v,s,h,a}$ with $h\in\N$. A column choice assigns $I_{v,s,c}:\N\to\{0,\ldots,r_1-1\}$ to each column $0\le c<\nu_{j+1,s,1}$. In its block put $h_c=\#\{a:c<\nu_{j+1,s,a}\}$, with $1\le a\le r_1$. For $\mathbf i=(i_v)_v\in\N^V$ set $q_c=\#\{h'<h_c:I_{v,s,c}(h')<i_v\}$ and
\[
 \begin{aligned}
 D_{v,s,c}&=\det\left(\sum_{a<r_0}\tau_{v,s,h,a}Z_{v,s,a,I_{v,s,c}(h')}\right)_{h,h'<h_c},
 &P_{I,v}&=\prod_{s,c}D_{v,s,c},\\
 e_{v,s,c}&=\begin{cases}q_c&s\ {\rm odd},\\
 \max\{\min(h_c,i_v)-q_c,0\}&s\ {\rm even},\end{cases}
 &P_{I,\mathbf i}&=\prod_{v,s,c}p^{e_{v,s,c}}D_{v,s,c}.
 \end{aligned}
\]
Thus $P_{I,\mathbf0}=\prod_vP_{I,v}$. At $v$ the degree in each of $Z$ and $\tau$ is the sum of column heights $\sum_s(\alpha_s+L_{j+1,s})$.
\end{definition}

\begin{lemma}[Contraction of a complementary column]\label{B3.10b}
Work over a field of characteristic zero. For a square matrix \(Y\) of size \(s\), write \(C_h(Y,z)\) for the matrix with first \(s-h\) columns taken from \(Y\) and last \(h\) columns taken from the beginning of \(z\). Fix heights \(0\le h_k\le s\), \(1\le k\le c\). Let \(H\) be the product of the initial principal minors of heights \(s-h_k\). For polynomials in \(Y\) use the pairing
\[
 [P,Q]=\sum_\alpha\alpha![Y^\alpha]P[Y^\alpha]Q.
\]
Let \(I_k\) be a list of \(s-h_k\) distinct row indices, let \(E_{I_k}\) have the corresponding unit vectors as columns in the indicated order, and put
\(Q_I(Y)=\prod_k\det Y_{I_k,\,1:s-h_k}\).
Then
\[
 \frac{\bigl[\prod_k\det C_{h_k}(Y,z),Q_I(Y)\bigr]}{[H,H]}
 =\prod_k\det[\,E_{I_k}\mid z_{:,1:h_k}\,].
\]
The identity holds over the polynomial ring in the entries of \(z\), including after a substitution that forces \(z\) to have deficient rank.
\end{lemma}

\begin{proof}[Outline]
Row degrees reduce the normalized pairing with \(H\) to evaluation at the identity. Reverse and factor out the generic column matrix, use adjointness and expand along the unit columns, then specialize the resulting polynomial identity. The complete proof is in \S\ref{prf:B3.10b}.
\end{proof}

\begin{lemma}[A frame for simultaneous lower neighbours]\label{B3.10a}
Let $K=\operatorname{Frac}(O)$ and $K_c=\mathbb C_p$ with the given embeddings of $K$ and $K_0$, where $K\hookrightarrow\mathbb C_p$ maps $O$ into $\mathcal O_{\mathbb C_p}$. Use the arithmetic operator $\mathcal D_j$ of Definition~\ref{B3.05aa}. At $v$ let $W_v$ be the local Witt ring with residue field $\F_{p^{n'_v}}$ and uniformizer $p$. Index its embeddings into $O$ by $s$. For $A=GW_v^n$ the arithmetic length is $s(A)=-\operatorname{ord}_p\det G$, independent of the basis since changes of basis have unit determinant.

Let $0\le j<m$ and let $L,L'$ be component--lattice pairs at $j,j+1$. Fix $0\le i_v\le r$ and take $L''$ in the target component with $pL'_v\subseteq L''_v\subseteq L'_v$ and $\dim_{W_v/p}(L'_v/L''_v)=i_v$. Let $D(L,L'')$ denote $D_j(L,L'')$ with the arithmetic lengths, and put
\[
 N'_{v,j+1}(i)=\sum_{s\ {\rm even}}\sum_{a=1}^i\lambda_{v,s,a}
 +i\sum_{s\ {\rm odd}}w_{v,j+1}(s).
\]
There is a frame $\phi$ from $L$ to $L''$ that serves for all $f\in K_c[[Z]]$ and $g\in\R$. If every coefficient of $\mathcal D_{P_{I,\mathbf i}}(f)$ from \ref{B3.11} has valuation at least $g$ for every $I$, then for every $G\in H_j(L)$,
\[
 \left(\prod_v p^{N'_{v,j+1}(i_v)}\right)D(L,L'')\,\mathcal D_j(G,f(B_\phi Z))
 \in \left\{\sum_{a=1}^{k}t_a z_a:k\ge0,\ t_a\in K_c,\ v_p(t_a)\ge g,\
 z_a\in H_{j+1}(L')\right\}.
\]
Here $v_p(p)=1$ and the target polynomials are extended to $K_c$.
\end{lemma}

\begin{proof}[Outline]
Builds \(\phi\) from compatible bases of \(L_v,L'_v,L''_v\) and evaluates the symbol of \ref{B3.05aa} on Weyl generators from \ref{B3.01a}. Framed pairings matching \ref{B3.10} give a unit scale, and integral Schur--Weyl duality yields the \(H_{j+1}(L')\) sum. The complete proof is in \S\ref{prf:B3.10a}.
\end{proof}

\begin{definition}[Differential evaluation]\label{B3.11}
For $P\in\Z[\tau][Z]$ and $f\in K_c[[Z]]$ use
\[
 \mathcal D_P(f)=\sum_\alpha[Z^\alpha]P\,\alpha![Z^\alpha]f=(P(\partial_Z)f)(0).
\]
The equality follows from $(\partial_Z^\alpha f)(0)=\alpha![Z^\alpha]f$. The same finite contraction applies to symbol polynomials with target polynomial coefficients.
\end{definition}

\begin{definition}[The single-place coefficient bound]\label{B3.12}
For $z\in K_c$ and $g\in\R$ define
\[
 \nu_c(z)\ge g\quad\Longleftrightarrow\quad
 \forall a\in\Z\ \forall b\in\N_{>0},\quad a\le bg\ \Longrightarrow\ z^bp^{-a}\in O_c .
\]
For the ring of a real valuation $v$ with $v(p)=1$, this means $v(z)\ge g$ by density of the rationals. For $j\ge0$ and $v\in V$, put $\Phi_{j,v}=\sum_{e\in E_v}b(X_{v,e},Y_{v,e})$ in the full variable set. The single-place coefficient bound $\mathcal L_{v,j}(g)$ means that
\[
 \nu_c\!\left(
 [\tau^n]\mathcal D_{P_{I,v}}\bigl(\mathcal T_j(\exp(-\Phi_{j,v})U)\bigr)\right)\ge g
\]
for every $U\in K_c[[X,Y]]$ with $U^\ell\in O_c[[X,Y]]$ and for all choices of $I$ and auxiliary multi-index $n$. The condition at $g$ implies the condition at every $g'\le g$.
\end{definition}


% ---- p3c2 ----
\chapter{Entry estimates: factorials against denominators}\label{ch:b3-entry-estimates}

The preceding chapter associates a differential polynomial to the
target weight diagram and controls the denominators of the scaled
theta kernel. The entry estimates bound their differential pairing at one step of
the theta chain.

Fix a prime \(p>2\) and a positive integer \(\nu\). Put
\[
 S=\Z/(2\nu)\Z,\qquad
 \rho_p=\frac1{p^2-1},\qquad
 \gamma_p=\frac{p}{p^2-1}.
\]
The parity of a position is the parity of its representative in
\(\{0,\ldots,2\nu-1\}\). Addition to a position takes place in \(S\).
The constants satisfy
\begin{equation}\label{eq:b3-entry-geometric}
 \sum_{l\ge0}p^{-(2l+1)}=\gamma_p,
 \qquad
 \frac1{p-1}=\gamma_p+\rho_p.
\end{equation}

For a nonarchimedean normed field \(K\), put
\(\mathcal O_K=\{a\in K:|a|\le1\}\). A polynomial or power series is
integral when its coefficients lie in \(\mathcal O_K\). When
\(|p|=p^{-1}\), write \(v_p(a)=-\log_p|a|\) for \(a\ne0\) and
\(v_p(0)=+\infty\). For \(a\in K\) (with \(a=0\) allowed) and \(x\in\R\), the bound
\(|a|\le p^{-x}\) is equivalent to \(v_p(a)\ge x\). All power series operations are formal.

\section{The two local exponents}

For the weight diagram at a place \(v\) in \ref{B3.02}, take
\(2\nu=n'_v\), the number of positions. Identify each position with
its residue class in \(S\) and let \(T\) be the image of the marked set
\(T_v\). Fix a layer \(1\le j\le m\), where \(m\) is the chain length.
At \(s\), let \(\alpha_s\) count the distinguished height-one columns.
The remaining columns number \(C_s\) and contain \(L_s\) boxes.
Thus \(\alpha_s=0\) for \(s\notin T\). Take \(r_0=N_{j-1}\) and
\(r_1=N_j\), with \(N_j\) as in \ref{B3.02}. In the entry estimates below, the
marked set \(T\subseteq S\) and the ranks \(r_0,r_1\in\N\) are arbitrary.

For \(a\in\N\), set
\[
 \ell_p(a)=
 \begin{cases}
  0,&a=0,\\
  1+\lfloor\log_p a\rfloor,&a>0.
 \end{cases}
\]
We use the error from \ref{B3.03} in the form
\begin{equation}\label{eq:b3-entry-error}
 \delta=
 \sum_{s\in S}
 \bigl(r_1\ell_p(\alpha_s)+r_0r_1\ell_p(L_s)\bigr)
 +\sum_{s\in T}\ell_p(\alpha_s)^2.
\end{equation}

\begin{definition}[The ordinary local exponent]\label{B3.13}
For \(s\in S\), let
\[
 D_s=
 \bigl\{(s'-s)_{\mathrm{rep}}:
       s'\in T,\ s'\text{ has parity opposite to }s\bigr\},
\]
where \((s'-s)_{\mathrm{rep}}\) denotes the representative in
\(\{0,\ldots,2\nu-1\}\). Define
\begin{equation}\label{eq:b3-entry-ga}
 g_A=
 \sum_{s\in T}
 \sum_{\substack{d\ge1,\ d\ {\rm odd}\\d\notin D_s}}
       [\alpha_sp^{-d}-C_{s+d}]_+
 -\delta,
 \qquad [x]_+=\max(x,0).
\end{equation}
For these weight data, \(g_A=g_{0,v}(j)\).
\end{definition}

Every \(d\in D_s\) is the first positive distance from \(s\) to its
endpoint. The sum omits this odd integer \(d\) and includes
\(d+2\nu\). The bound \([\alpha_sp^{-d}-C_{s+d}]_+\le\alpha_sp^{-d}\)
gives convergence.

The identification with the local system follows by substituting
\[
 \alpha_s=\lambda_{s,1}-\lambda_{s,2},\quad C_s=\nu_{j,s,2}
 \quad(s\in T),
 \qquad
 \alpha_s=0,\quad C_s=\nu_{j,s,1}\quad(s\notin T).
\]
The error \eqref{eq:b3-entry-error} is then \(\delta_j\).

\begin{definition}[The enhanced local exponent]\label{B3.14}
Write
\[
 T_{\mathrm{odd}}=\{s\in T:s\text{ is odd}\},\qquad
 T_{\mathrm{ev}}=T\setminus T_{\mathrm{odd}},\qquad
 A_{\mathrm{ev}}=\sum_{s\in T_{\mathrm{ev}}}\alpha_s.
\]
Define
\begin{equation}\label{eq:b3-entry-gb}
 g_B=
 \gamma_p\sum_{s\in T}\alpha_s
 -\sum_{s\in T_{\mathrm{odd}}}
      \sum_{\substack{d\ge1\\d\ {\rm odd}}}
         \min(\alpha_sp^{-d},C_{s+d})
 -\delta.
\end{equation}
For these weight data, \(g_B=g_v(j)\) and
\(A_{\mathrm{ev}}=\mathrm{alEven}(v)\).
\end{definition}

The series converges by
\(0\le\min(\alpha_sp^{-d},C_{s+d})\le\alpha_sp^{-d}\).

\section{Differential pairings and column bounds}

For a finite variable set \(\mathcal V\), a multi-index
\(n\in\N^{\mathcal V}\), a polynomial \(P\), and a power series \(F\), put
\[
 |n|=\sum_{x\in\mathcal V}n_x,\qquad
 n!=\prod_{x\in\mathcal V}n_x!,\qquad
 \langle P,F\rangle
   =\sum_n [z^n]P\,n!\,[z^n]F.
\]
The sum in this differential evaluation (\ref{B3.11}) runs over
the monomials of \(P\). Write \(\mathbf1_x\) for the unit multi-index
at \(x\).

\begin{lemma}[Linear factors and Hasse derivatives]\label{B3.15}
Let \(K\) be a field, and let \(\mathcal V,\mathcal K\) be finite sets.
For directions \(b_k\in K^{\mathcal V}\), define
\[
 L_k(z)=\sum_{x\in\mathcal V}b_{k,x}z_x,\qquad
 D_b^{[\beta]}P
  =[t^\beta]P\left(z+\sum_{k\in\mathcal K}t_kb_k\right)
 \quad(\beta\in\N^{\mathcal K}).
\]
For every polynomial \(P\) and power series \(F\),
\begin{equation}\label{eq:b3-hasse-pair}
 \left\langle P,\left(\prod_kL_k^{\beta_k}\right)F\right\rangle
   =\left(\prod_k\beta_k!\right)
       \left\langle D_b^{[\beta]}P,F\right\rangle.
\end{equation}
For the integrality and degree assertions, \(\mathcal V\) may be
arbitrary. Let \(K\) be an ultrametric normed field. Integral \(P\) and
integral directions give integral \(D_b^{[\beta]}P\). If \(P\) is
homogeneous of degree \(d\), then the derivative is homogeneous of
degree \(d-|\beta|\) for \(|\beta|\le d\), and is zero for \(|\beta|>d\).
\end{lemma}

\begin{proof}[Outline]
Check adjunction of linear multiplication and directional differentiation on monomials, then iterate and use the translation expansion to identify the iterated derivative with the Hasse derivative times the factorials. The expansion preserves integrality and gives the degree drop. The complete proof is in \S\ref{prf:B3.15}.
\end{proof}

Let \(\Sigma\) be a set of positions and use variables
\[
 z_{s,a,b},\qquad
 s\in\Sigma,\quad 0\le a<r_0,\quad 0\le b<r_1.
\]
For \(c=1,\ldots,C\), choose a position \(s_c\), a height \(h_c\in\N\),
a matrix \(M_c=(M_{c,i,a})_{0\le i<h_c,\ 0\le a<r_0}\), and a map
\[
 I_c:\{0,\ldots,h_c-1\}\longrightarrow\{0,\ldots,r_1-1\}.
\]
Put
\begin{equation}\label{eq:b3-low-minors}
 \Delta_c(z)=
 \det\left(\sum_{a<r_0}M_{c,i,a}z_{s_c,a,I_c(j)}
                  \right)_{0\le i,j<h_c},
 \qquad Q_0(z)=\prod_{c=1}^C\Delta_c(z).
\end{equation}
A height-zero determinant is \(1\). The maps \(I_c\) may repeat column
labels. A direction at \(s\) with row vector \(\xi\in K^{r_0}\)
and column label \(b\) is
\[
 B(s,\xi,b)_{t,a,b'}=
 \begin{cases}
  \xi_a,&t=s,\ b'=b,\\
  0,&\text{otherwise}.
 \end{cases}
\]

\begin{lemma}[Vanishing when columns are exhausted]\label{B3.16}
Let \(K\) be any field. For a finite family of directions
\(b_k=B(\operatorname{pos}(k),\xi_k,\operatorname{col}(k))\)
indexed by \(\mathcal K\), let \(\beta\in\N^{\mathcal K}\).
A direction is used when its multiplicity is positive.
Each of the following conditions forces \(D_b^{[\beta]}\Delta_c=0\).
\begin{enumerate}
 \item Some \(\beta_k\ge2\).
 \item Distinct used directions \(k,k'\) have the same position and
       either \(\xi_k=\xi_{k'}\) or
       \(\operatorname{col}(k)=\operatorname{col}(k')\).
 \item A used direction has position different from \(s_c\), or has
       column outside \(\operatorname{im}I_c\).
\end{enumerate}
Let \(\mathcal K_0\subseteq\mathcal K\) consist of directions at a
single position \(s_0\). If \(D_b^{[\beta]}Q_0\ne0\),
the following bounds hold.
\begin{equation}\label{eq:b3-hits-cap}
 \sum_{k\in\mathcal K_0}\beta_k
 \le \#\{c:s_c=s_0\}
 \quad\text{if all }\xi_k\ (k\in\mathcal K_0)\text{ are equal},
\end{equation}
and, for every \(J\subseteq\{0,\ldots,r_1-1\}\),
\begin{equation}\label{eq:b3-hits-columns}
 \sum_{k\in\mathcal K_0}\beta_k
 \le \sum_{c:s_c=s_0}\#(\operatorname{im}I_c\cap J)
 \quad\text{if }\operatorname{col}(k)\in J
       \text{ for }k\in\mathcal K_0.
\end{equation}
\end{lemma}

\begin{proof}[Outline]
Repeated labels give zero. Column multilinearity allows each available label in \(J\) at most once, and alternation allows a shared row vector at most once per determinant when the Hasse product formula distributes the directions. The complete proof is in \S\ref{prf:B3.16}.
\end{proof}

\section{Four parts of the substituted kernel}

On the cyclic position set \(S\), the series from \ref{B3.07} is
\begin{equation}\label{eq:b3-hh}
 S_{s,a,b}(z)=\sum_{l\ge0}(-1)^lz_{s-2l,a,b}^{p^{2l}},
 \qquad
 H(z)=\sum_{s\in S}\sum_{a<r_0,\ b<r_1}
               u_sS_{s,a,b}(z)^pS_{s+1,a,b}(z),
\end{equation}
where \(u_s\in\mathcal O_K\). Introduce top variables \(t_{s,b}\) and remaining variables \(y_{s,a,b}\).
Let \(\xi_s=(\xi_{s,a})_{a<r_0}\) be integral vectors. Make the substitution
\begin{equation}\label{eq:b3-top-low-sub}
 z_{s,a,b}\longmapsto
 y_{s,a,b}+\mathbf1_{s\in T}\xi_{s,a}t_{s,b}.
\end{equation}
Write \(h(t,y)\) for the resulting series. The remaining degree of
a monomial is its total degree in \(y\).

For \(E\in\N\), index the pure-top terms by
\[
 \phi=(v,l,l',a,b),\qquad
 v\in S,\quad 0\le l,l'<E,\quad a<r_0,\quad b<r_1.
\]
Set \(s_\phi=v-2l\), \(t_\phi=v+1-2l'\), and define
\begin{align}
 d^A_\phi
   &=p^{2l+1}\mathbf1_{(s_\phi,b)}
       +p^{2l'}\mathbf1_{(t_\phi,b)},\label{eq:b3-da}\\
 a^A_\phi
   &=\mathbf1_{s_\phi,t_\phi\in T}\,
      u_v(-1)^{l+l'}
      \xi_{s_\phi,a}^{p^{2l+1}}\xi_{t_\phi,a}^{p^{2l'}}.
       \label{eq:b3-aa}
\end{align}
For the second family, let
\(\kappa=(s,l,b)\) with \(s\in S\), \(0\le l<E\), \(b<r_1\), and put
\begin{align}
 d^B_\kappa&=p^{2l+1}\mathbf1_{(s,b)},\label{eq:b3-db}\\
 a^B_\kappa&=\mathbf1_{s\in T}(-1)^lu_{s+2l},\label{eq:b3-ab}\\
 L_\kappa(y)&=
   \sum_{a<r_0}\xi_{s,a}^{p^{2l+1}}y_{s+2l+1,a,b}.
       \label{eq:b3-b-direction}
\end{align}
The exponents \(d^A_\phi,d^B_\kappa\) index top monomials. The direction
\(b_\kappa\) of \(L_\kappa\) is \(B(s+2l+1,(\xi_{s,a}^{p^{2l+1}})_a,b)\).

\begin{lemma}[The four-part decomposition]\label{B3.17}
Assume \(K\) is nonarchimedean, \(p>2\) is prime, \(|p|=p^{-1}\),
and all \(u_s,\xi_{s,a}\) are integral. Let \(D,E\in\N\) satisfy
\(D<p^{2E}\). There are integral power series \(Z_{\ge p}\) and \(Y\)
such that every monomial of \(Z_{\ge p}\) with nonzero coefficient has
remaining degree at least \(p\), and
\begin{equation}\label{eq:b3-four-parts}
 h(t,y)\equiv
  \sum_\phi a^A_\phi t^{d^A_\phi}
  +\sum_\kappa a^B_\kappa t^{d^B_\kappa}L_\kappa(y)
  +Z_{\ge p}(t,y)+pY(t,y)
       \pmod{\text{total degree }>D}.
\end{equation}
The families and coefficients are given by
\eqref{eq:b3-da}--\eqref{eq:b3-b-direction}.
\end{lemma}

\begin{proof}[Outline]
Keep \(l<E\) through degree \(D\) and use Frobenius modulo \(p\) to separate pure top terms, terms linear in \(y\), and terms of \(y\)-degree at least \(p\). Reindex and sum the linear terms to get \(L_\kappa\), with remainder \(pY\) for integral \(Y\). The complete proof is in \S\ref{prf:B3.17}.
\end{proof}

For a series \(h\) with zero constant term, write \(g\in\mathcal C_p(h)\)
if there are integral series \(U_i\), \(i\in\N\), with
\begin{equation}\label{eq:b3-class-recall}
 [z^n]g=\sum_{i=0}^{|n|}p^{-i}[z^n](h^iU_i)
 \quad\text{for every multi-index }n.
\end{equation}
This is the expansion class of \ref{B3.07}. For finitely many source
variables, let \(\phi\) be a substitution by integral series with zero
constant term. If \(g\in\mathcal C_p(h)\), then
\(g\circ\phi\in\mathcal C_p(h\circ\phi)\). The witnesses are
\(U_i\circ\phi\), since substitution preserves integrality and lower
bounds on total degree.

\begin{lemma}[The factorial bound for a nonzero pairing]\label{B3.18}
Let \(\mathcal T,\mathcal L,\mathcal A,\mathcal K\) be finite sets,
and let \(K\) be a nonarchimedean normed field with \(|p|=p^{-1}\),
where \(p\) is prime. Use variables \(t_\tau\), \(\tau\in\mathcal T\),
and \(y_x\), \(x\in\mathcal L\).

Fix a total degree \(D\in\N\).
Suppose \(h\) has zero constant term, \(g\in\mathcal C_p(h)\), and,
in total degrees at most \(D\),
\[
 h=
 \sum_{\phi\in\mathcal A}a^A_\phi t^{d^A_\phi}
 +\sum_{\kappa\in\mathcal K}
          a^B_\kappa t^{d^B_\kappa}L_\kappa(y)
 +Z_{\ge p}+pY.
\]
Here \(d^A_\phi,d^B_\kappa\in\N^{\mathcal T}\) and
\(L_\kappa(y)=\sum_x b_{\kappa,x}y_x\). The scalars \(a^A_\phi,a^B_\kappa\),
the coordinates \(b_{\kappa,x}\), and the series \(Z_{\ge p},Y\) are
integral. Every nonzero monomial of \(Z_{\ge p}\) has \(y\)-degree
at least \(p\).

Let \(Q_0(y)\) be an integral homogeneous polynomial of degree \(d_0\),
and let \(e\in\N^{\mathcal T}\) satisfy \(|e|+d_0=D\).
If \(J=\langle t^eQ_0(y),g\rangle\ne0\), there exist multiplicities
\(m_\phi,\beta_\kappa\in\N\), an integer \(i_2\ge0\), and a multi-index
\(n\in\operatorname{supp}D_b^{[\beta]}Q_0\) such that
\begin{gather}
 m_\phi>0\Longrightarrow a^A_\phi\ne0,\qquad
 \beta_\kappa>0\Longrightarrow a^B_\kappa\ne0,
       \label{eq:b3-term-active}\\
 \sum_\phi m_\phi d^A_\phi
       +\sum_\kappa\beta_\kappa d^B_\kappa\le e,
 \qquad pi_2\le|n|,\label{eq:b3-term-consumption}\\
 |J|\le
 p^{\,\sum_\phi m_\phi+\sum_\kappa\beta_\kappa+i_2}
       |e!|\,
       \left|\left(\sum_\kappa\beta_\kappa\right)!\right|\,
       |n!|.\label{eq:b3-term-factorials}
\end{gather}
The inequality between top multi-indices is coordinatewise.
\end{lemma}

\begin{proof}[Outline]
Expand \(h^i\) by the multinomial formula and use \ref{B3.15} to move the linear factors to Hasse derivatives and collect \(B!\). The selected monomials give the exponent bounds, \(pY\) cancels the matching denominators, and the ultrametric inequality gives the bound. The complete proof is in \S\ref{prf:B3.18}.
\end{proof}

\section{The arithmetic contribution of the factorials}

\begin{lemma}[A factorial valuation bound]\label{B3.19}
For every prime \(p\) and every \(a\in\N\), the arithmetic valuation
satisfies
\begin{equation}\label{eq:b3-factorial-digit}
 v_p(a!)\ge\frac{a}{p-1}-\ell_p(a).
\end{equation}
\end{lemma}

Legendre's formula gives \(v_p(a!)=(a-s_p(a))/(p-1)\), where \(s_p(a)\)
is the base-\(p\) digit sum \cite[Ch.~I, \S2, Ex.~14, p.~7]{koblitz-1984}.
Substituting \(s_p(a)\le(p-1)\ell_p(a)\) in Legendre's formula gives the bound.


\begin{lemma}[The net contribution of three factorial families]\label{B3.20}
Let \(S\) be a finite set, \(T\subseteq S\), \(p\) prime, and
\(r_0,r_1\in\N\). Let \(\alpha,L:S\to\N\), with \(\alpha_s=0\)
outside \(T\). Suppose nonnegative integer arrays
\[
 e_{s,b}\ (b<r_1),\qquad
 a_{s,a,b}\ (a<r_0,\ b<r_1),\qquad
 i_{s,l}\ (l\in\N)
\]
and integers \(i_0,i_2\ge0\) satisfy
\begin{align}
 &\sum_{b<r_1}e_{s,b}=\alpha_s,\qquad
   \sum_{a,b}a_{s,a,b}\le L_s,\label{eq:b3-star-local}\\
 &i_{s,l}=0\quad\text{if }s\notin T
                   \text{ or }l\notin E,\label{eq:b3-star-support}\\
 &p^{2l+1}i_{s,l}\le\alpha_s,\label{eq:b3-star-source}\\
 &\sum_{s,a,b}a_{s,a,b}+\sum_{s\in S}\sum_{l\in E}i_{s,l}
        =\sum_sL_s,\qquad
   pi_2\le\sum_sL_s,\label{eq:b3-star-total}
\end{align}
where \(E\subseteq\N\) is finite. Define \(\delta\) by
\eqref{eq:b3-entry-error}, and set
\[
 \mathcal F=
 \sum_{s,b}v_p(e_{s,b}!)
 +\sum_{s\in S}\sum_{l\in E}v_p(i_{s,l}!)
 +\sum_{s,a,b}v_p(a_{s,a,b}!).
\]
Then
\begin{equation}\label{eq:b3-star-arith}
\begin{aligned}
 &\mathcal F-\rho_p\sum_{s\in S}(\alpha_s+L_s)
       -\left(i_0+\sum_{s\in S}\sum_{l\in E}i_{s,l}+i_2\right)\\
 &\qquad\ge \gamma_p\sum_{s\in T}\alpha_s
       -i_0-\sum_{s\in T}\sum_{l\in E}i_{s,l}-\delta.
\end{aligned}
\end{equation}
\end{lemma}

\begin{proof}[Outline]
Apply \ref{B3.19} to the three factorial families, controlling digit errors and the number of positive indices by the local degree bounds. Use the degree identity, \(c_p-\rho_p=\gamma_p\), and \(i_2\le\gamma_pL\) to obtain the estimate. The complete proof is in \S\ref{prf:B3.20}.
\end{proof}

\section{A common estimate for every nonzero entry}

Let \(K\) be an ultrametric normed field with \(|p|=p^{-1}\), where
\(p>2\) is prime, and let \(\nu\ge1\). Use variables \(z_{s,a,b}\)
with \(s\in S\), \(a<r_0\), \(b<r_1\). Take the determinant product
\(Q_0\) of \eqref{eq:b3-low-minors} with positions \(s_c\in S\)
and integral frame entries \(M_{c,i,a}\).

Choose integers \(e_{s,b}\ge0\), supported on \(T\), and integral
vectors \(\xi_s\in\mathcal O_K^{r_0}\). Define
\begin{align}
 P_{\mathrm{ent}}(z)
  &=\prod_{s\in S,\ b<r_1}
       \left(\sum_{a<r_0}\xi_{s,a}z_{s,a,b}\right)^{e_{s,b}}
       Q_0(z),\label{eq:b3-entry-polynomial}\\
 \alpha_s&=\sum_{b<r_1}e_{s,b},\qquad
 C_s=\#\{c:s_c=s\},\qquad
 L_s=\sum_{c:s_c=s}h_c,\label{eq:b3-entry-counts}\\
 D&=\sum_s(\alpha_s+L_s).\label{eq:b3-entry-degree}
\end{align}
Let \(H\) be \eqref{eq:b3-hh}, with integral \(u_s\), and take
\(F\in\mathcal C_p(H)\). Under these entry hypotheses, the entry is
\(\mathcal J=\langle P_{\mathrm{ent}},F\rangle\).

\begin{theorem}[The common entry estimate]\label{B3.21}
Under the entry hypotheses, either \(\mathcal J=0\), or there are
an integer \(E\ge0\) and nonnegative multiplicities \(m_\phi,\beta_\kappa\)
for the families of \ref{B3.17} such that
\begin{gather}
 m_\phi>0\Longrightarrow s_\phi,t_\phi\in T,\qquad
 \beta_{(s,l,b)}>0\Longrightarrow s\in T,
       \label{eq:b3-entry-support}\\
 \sum_\phi m_\phi d^A_\phi
       +\sum_\kappa\beta_\kappa d^B_\kappa\le e,
       \label{eq:b3-entry-budget}\\
 D_b^{[\beta]}Q_0\ne0,\label{eq:b3-entry-hasse-nonzero}
\end{gather}
and
\begin{equation}\label{eq:b3-entry-star}
 v_p(\mathcal J)\ge
 \gamma_p\sum_{s\in T}\alpha_s
     -\sum_\phi m_\phi-\sum_\kappa\beta_\kappa
     -\delta+\rho_pD.
\end{equation}
Here \(b_\kappa\) is the direction in \eqref{eq:b3-b-direction}.
\end{theorem}

\begin{proof}[Outline]
Use the chain rule to substitute on the series side, then apply \ref{B3.17} and \ref{B3.18} to obtain the multiplicities and a Hasse derivative monomial. Group linear factors by source and level, then apply \ref{B3.20} using factorial divisibility and the degree bounds. The complete proof is in \S\ref{prf:B3.21}.
\end{proof}

\section{First distances and the ordinary gain}

\begin{lemma}[Charging losses to first distances]\label{B3.22}
Fix \(n,a,a',d\in\N\) and \(s,s'\in\Z/n\Z\).
For \(n=0\), interpret \(\Z/0\Z\) as \(\Z\).
Suppose \(a'<a\), \(s+a=s'+a'\), and \(d\) is the least positive
integer with \(s+d=s'\). Then \(d\le a\).

For the scalar estimate, take a finite subset \(T\) of \(S=\Z/n\Z\),
where \(n\in\N\). Let \(p>1\) and \(\delta\) be real and let
\(\alpha,C:S\to\R_{\ge0}\). Give each \(s\) a finite
set \(D_s\subseteq\N\). Suppose nonnegative real arrays
\(i_{s,l},\operatorname{ch}_{s,l}\) are supported in a common finite
set of indices \(l\in\N\) and satisfy
\begin{align}
 &\operatorname{ch}_{s,l}=0
       &&\text{if }2l+1\notin D_s,\label{eq:b3-charge-support}\\
 &\operatorname{ch}_{s,l}+i_{s,l}
        \le\alpha_sp^{-(2l+1)}
       &&\text{if }s\in T,\ 2l+1\in D_s,\label{eq:b3-charge-excluded}\\
 &i_{s,l}\le
       \min(\alpha_sp^{-(2l+1)},C_{s+2l+1})
       &&\text{if }s\in T,\ 2l+1\notin D_s.
       \label{eq:b3-charge-free}
\end{align}
If a real number \(V\) satisfies
\[
 V\ge
 \gamma_p\sum_{s\in T}\alpha_s
 -\sum_{s\in T,l}\operatorname{ch}_{s,l}
 -\sum_{s\in T,l}i_{s,l}-\delta,
\]
then
\begin{equation}\label{eq:b3-charge-conclusion}
 V\ge
 \sum_{s\in T}
   \sum_{\substack{d\ge1,\ d\ {\rm odd}\\d\notin D_s}}
        [\alpha_sp^{-d}-C_{s+d}]_+-\delta.
\end{equation}
\end{lemma}

\begin{proof}[Outline]
Minimality gives \(d\le a-a'\le a\). The unused allowance is nonnegative at excluded odd distances and at least \([x_l-C_{s+2l+1}]_+\) at the others, so sum using \(\sum_l x_l=\gamma_p\alpha_s\) and subtract \(\delta\). The complete proof is in \S\ref{prf:B3.22}.
\end{proof}

\begin{theorem}[The ordinary entry bound]\label{B3.23}
Under the entry hypotheses of \ref{B3.21}, with \(D_s\) as in
\ref{B3.13},
\begin{equation}\label{eq:b3-entry-a1}
 |\mathcal J|\le p^{-(g_A+\rho_pD)}.
\end{equation}
\end{theorem}

\begin{proof}[Outline]
For the multiplicities from \ref{B3.21}, bound linear terms by source degree and target determinant count. Charge pure top terms to the first positive distance from the endpoint with larger shift, then use \ref{B3.22} to obtain the gain \(g_A\) in the valuation bound. The complete proof is in \S\ref{prf:B3.23}.
\end{proof}

\section{Parity and the enhanced entry bounds}

\begin{theorem}[The second unmodified entry bound]\label{B3.24}
Under the entry hypotheses,
\begin{equation}\label{eq:b3-entry-a2}
 |\mathcal J|\le p^{-(g_B-A_{\mathrm{ev}}+\rho_pD)}.
\end{equation}
\end{theorem}

\begin{proof}[Outline]
For the multiplicities from \ref{B3.21}, sum the even top degrees to get \(A+B_{\mathrm{ev}}\le A_{\mathrm{ev}}\). Bound odd-source terms by source degree and target determinant count, then substitute into the common valuation estimate. The complete proof is in \S\ref{prf:B3.24}.
\end{proof}

\begin{theorem}[The prescribed-column entry bound]\label{B3.25}
Under the entry hypotheses, let
\(J\subseteq\{0,\ldots,r_1-1\}\) be arbitrary. Define the nonnegative
integer
\begin{align}
 N_J={}&
 \sum_{\substack{s\ {\rm even}\\b\notin J}}e_{s,b}
 +\sum_{\substack{s\ {\rm odd}\\b\in J}}e_{s,b}
 +\sum_{c:\ s_c\ {\rm odd}}\#(\operatorname{im}I_c\cap J).
       \label{eq:b3-entry-column-payment}
\end{align}
Then
\begin{equation}\label{eq:b3-entry-b}
 |p^{N_J}\mathcal J|\le p^{-(g_B+\rho_pD)}.
\end{equation}
\end{theorem}

\begin{proof}[Outline]
For \ref{B3.21}, sum top degrees at even labels outside \(J\) and odd labels in \(J\), then use column bounds at odd targets to show \(N_J\) covers pure top and even-source terms. Bound odd-source terms at source and target and add \(N_J\) to the valuation bound. The complete proof is in \S\ref{prf:B3.25}.
\end{proof}

\section{From entries to coefficients of the frame jet}

At each \(s\) take \(c_s\in\N\) columns of heights
\(h_{s,c}\in\N\), \(h_{s,c}\le r_0\) for \(0\le c<c_s\).
Choose maps
\[
 I_{s,c}:\{0,\ldots,h_{s,c}-1\}
                  \longrightarrow\{0,\ldots,r_1-1\}.
\]
Use independent frame variables
\(\tau_{s,i,a}\), \(s\in S\), \(i,a<r_0\), and define
\begin{equation}\label{eq:b3-frame-polynomial}
 \mathcal P(\tau,z)=
 \prod_{s\in S}\prod_{c<c_s}
  \det\left(
       \sum_{a<r_0}\tau_{s,i,a}z_{s,a,I_{s,c}(j)}
      \right)_{0\le i,j<h_{s,c}}.
\end{equation}
For a power series \(\Theta(z)\), its frame jet is the polynomial
\begin{equation}\label{eq:b3-frame-jet}
 \mathcal J_\Theta(\tau)
   =\sum_n n!\,[z^n]\Theta\,[z^n]\mathcal P(\tau,z).
\end{equation}
This is the coefficientwise version of the operator in
\ref{B3.10} and \ref{B3.11}. For these diagrams put
\begin{align}
 \alpha_s&=
   \begin{cases}
    \#\{c<c_s:h_{s,c}=1\},&s\in T,\\
    0,&s\notin T,
   \end{cases}\label{eq:b3-gauss-alpha}\\
 C_s&=c_s-\alpha_s,\qquad
 L_s=\sum_{c<c_s}h_{s,c}-\alpha_s,\qquad
 D=\sum_s(\alpha_s+L_s)
   =\sum_s\sum_{c<c_s}h_{s,c}.
       \label{eq:b3-gauss-counts}
\end{align}
Use these counts in \(\delta,g_A,g_B,A_{\mathrm{ev}}\).

\begin{theorem}[Frame-jet coefficients of the unscaled kernel]\label{B3.26}
Let \(K\) be a characteristic-zero nonarchimedean normed field.
Fix a prime \(p>2\) with \(\lvert p\rvert=p^{-1}\).
Take \(\pi\in K\) satisfying \(\pi^{p^2-1}=p\). Take \(\nu\ge1\), \(r_0,r_1\in\N\),
\(u_s\in\mathcal O_K\), and \(H\) as in \eqref{eq:b3-hh}. Suppose
\(\Theta(\pi z)\in\mathcal C_p(H)\). For the diagrams and maps in
\eqref{eq:b3-frame-polynomial}, each coefficient
\(a_\eta=[\tau^\eta]\mathcal J_\Theta\) satisfies
\begin{align}
 |a_\eta|&\le p^{-g_A},\label{eq:b3-kernel-a1}\\
 |a_\eta|&\le p^{-(g_B-A_{\mathrm{ev}})}.\label{eq:b3-kernel-a2}
\end{align}
For every integer \(i\ge1\), set
\[
 J_i=\{b:0\le b<r_1,\ b<i\},\qquad
 q_{s,c}(i)=\#(\operatorname{im}I_{s,c}\cap J_i),
\]
and
\begin{equation}\label{eq:b3-kernel-edge-exponent}
 M_i=
 \sum_{\substack{s\in S\\s\ {\rm even}}}
        \sum_{c<c_s}\bigl(\min(h_{s,c},i)-q_{s,c}(i)\bigr)
 +\sum_{\substack{s\in S\\s\ {\rm odd}}}
        \sum_{c<c_s}q_{s,c}(i).
\end{equation}
Then
\begin{equation}\label{eq:b3-kernel-b}
 |p^{M_i}a_\eta|\le p^{-g_B}.
\end{equation}
The differences in \eqref{eq:b3-gauss-counts} and
\eqref{eq:b3-kernel-edge-exponent} are nonnegative integers.
\end{theorem}

\begin{proof}[Outline]
Use the Gauss norm on \(K(\tau)\) to bound frame coefficients and separate marked factors to apply \ref{B3.23} and \ref{B3.24}. For prescribed columns, use \(N_{J_i}\le M_i\) and \ref{B3.25}, then divide all bounds by \(|\pi^D|\) to undo the degree-\(D\) scaling. The complete proof is in \S\ref{prf:B3.26}.
\end{proof}

For the target diagrams, \ref{B3.27} uses the ordinary coefficient
bound. The bounds with exponents \(g_B-A_{\mathrm{ev}}\) and \(g_B\)
give the two estimates of \ref{B3.28}.


% ---- p3c3 ----
\chapter{One lift on a lattice and the Hecke commutation}\label{ch:b3-one-lift}

Fix a local system as in Definition~\ref{B3.01} satisfying Definition~\ref{B3.05a} \cite{alpoge-2026}. Write \(\iota:O\to K\) for the coefficient map. For a uniformizer \(\varpi\) of \(O\), put \(e_O=\operatorname{ord}_\varpi(p)\). Put \(k=k_1+1\) and \(N_j=k+j(t+1)\). At level \(0\le j\le m\), the space \(\mathcal H_j\) consists of polynomial-valued functions on components and tuples of submodules, with variables \(X_{v,s,a,b}\) indexed by the places \(v\), their embeddings \(s\), and \(1\le a,b\le N_j\). A form in \(\Lambda_j\) is level-equivariant. At a component--lattice pair \(L\) it takes values in the weight lattice \(H_j(L)\) of Lemma~\ref{B3.10a}, and it vanishes elsewhere. Put \(V_j=K\Lambda_j\) and use
\[
 v_j(f)\ge g\quad\Longleftrightarrow\quad
 f\in\iota(\varpi)^{\lceil e_Og\rceil}\Lambda_j
 \qquad(g\in\R).
 \tag{3.27.1}\label{eq:p3c3-form-valuation}
\]
For the lifts \(K_j\) of Definition~\ref{B3.01}, put \(K^{[0]}=1\) and \(K^{[b]}=K_{b-1}\cdots K_0\) for \(1\le b\le m\).

Let \(\mathcal P\) be the finite set of local positions. At \(v\in\mathcal P\), the discrete valuation ring \(W_v\) has uniformizer \(p\) and residue field of order \(p^{n'_v}\). The embeddings are indexed in cyclic order by \(s=0,\ldots,n'_v-1\), where \(n'_v\ge2\) is even. Put \(\eta_v=n'_v/2\). Write \(\lambda_{v,s,1}\ge\cdots\ge\lambda_{v,s,k}\) for the initial weight and \(w_{v,b}(s)\) for the level weights of Definition~\ref{B3.01}. If \(\mathcal P\ne\varnothing\), then \(p\ge5\) and \(1\le r\le k/2\). At step \(0\le j<m\) use Definitions~\ref{B3.01}, \ref{B3.02}, \ref{B3.13} and~\ref{B3.14} to set
\[
 A_v=g_{A,v}(j+1),\qquad B_v=g_{B,v}(j+1),\qquad
 E_v=\sum_{\substack{s\in T_v\\s\ {\rm even}}}\alpha_{v,s},\qquad
 R_j(v)=\sum_{v'\ne v}A_{v'}-d_G(j).
 \tag{3.27.2}\label{eq:p3c3-gains}
\]
Write \(T_v\) for the set of marked positions and \(\alpha_{v,s}\) for the digit error. The bound on the \(p\)-valuation of the stabilizer orders is \(d_G(j)\). We normalize \(v_p(p)=1\).

\section{From the abstract kernel to the actual weights}

Let \(K_c\supseteq K\) and \(O_c\) be the kernel coefficient field and ring of Definition~\ref{B3.05}. The bound of Definition~\ref{B3.12} is
\[
 \nu_c(z)\ge g\quad\Longleftrightarrow\quad
 z^bp^{-a}\in O_c\quad\text{for all }a\in\Z,\ b\in\Z_{\ge1}
 \text{ with }a\le bg .
 \tag{3.27.3}\label{eq:p3c3-coeff-valuation}
\]
When \(e_O>0\), \(O_c\cap K=\iota(O)\) gives \(p^{-1}\notin O_c\). For \(z\ne0\) set \(v(z)=\sup\{a/b:z^bp^{-a}\in O_c,\ a\in\Z,\ b\ge1\}\). The rank-one condition applied to \(z,z^{-1}\) bounds this supremum above and below. The set of eligible rational exponents is downward closed, since \(p^n\in O_c\) exactly for integers \(n\ge0\). Clearing denominators gives \(v(xy)\ge v(x)+v(y)\). The valuation-ring dichotomy for exponents above \(v(x)\) gives \(v(x^{-1})\ge-v(x)\). The product inequality applied to \(xx^{-1}\) and to \((xy)y^{-1}\) gives equality in both laws. The second rank-one condition gives \(v(z)>0\) for nonzero \(z\) in the maximal ideal. Thus \(O_c=\{v\ge0\}\) and \(v(p)=1\). With \(v(0)=+\infty\), \eqref{eq:p3c3-coeff-valuation} is equivalent to \(v(z)\ge g\). Factoring out the summand of smaller value proves the ultrametric inequality. When \(e_O=0\), \(p\) is a unit of \(O_c\) and \eqref{eq:p3c3-coeff-valuation} is equivalent to \(z\in O_c\).

Write \(U^\ell(X)=U(\ell(X))\) for coordinatewise logarithmic substitution and call \(U\) integral after this substitution if \(U^\ell\in O_c[[X]]\). For the pairing \(b\) of Definition~\ref{B3.05}, put \(\Phi_v=\sum_{e\in E_v^{\rm pair}}b(X_{v,e},Y_{v,e})\), with \(E_v^{\rm pair}\) the pairing index set at \(v\). Write \(\mathcal C_j=\mathcal T_j\) for the substitution of Definition~\ref{B3.05}. With the determinant polynomial \(P_{I,v}\) of Definition~\ref{B3.10}, put
\[
 \mathcal J_{I,v}(U)=
 \left.P_{I,v}(\partial_Z)\mathcal C_j(\exp(-\Phi_v)U)\right|_{Z=0}.
 \tag{3.27.4}
\]
The contraction leaves the auxiliary polynomial variables \(\tau\) of Definition~\ref{B3.11}.

\begin{theorem}[The ordinary coefficient bound for the actual weight diagram]\label{B3.27}
For \(v\in\mathcal P\), \(0\le j<m\), every \(U\) integral after logarithmic substitution, every column choice \(I\), and every auxiliary multi-index \(n\),
\[
 \nu_c([\tau^n]\mathcal J_{I,v}(U))\ge A_v.
 \tag{3.27.5}\label{eq:p3c3-coeff-ordinary}
\]
Call \(g\in\R\) an ordinary coefficient bound at \((j,v)\) if \(\nu_c([\tau^n]\mathcal J_{I,v}(U))\ge g\) for every such \(U,I,n\).
\end{theorem}

\begin{proof}[Outline]
Delete the marked height-one columns from the diagram of Definition \ref{B3.02}. Scale by \(\pi\) and use \ref{B3.08} and the phase of \ref{B3.07} to expand in \(p^{-d}H_u^d\) with integral multipliers, then apply \ref{B3.26} for the norm bound \(p^{-A_v}\). The complete proof is in \S\ref{prf:B3.27}.
\end{proof}

For \(1\le i\le r\), let \(d_{s,c}(I,i)=\#\{a:0\le a<h_{s,c},\ I(s,c,a)<i\}\), with target labels starting at zero. The shifted factor of Definition~\ref{B3.10} is
\[
 P_{I,v}^{(i)}=p^{\sum_{s,c}\epsilon_{s,c}(I,i)}P_{I,v},\qquad
 \epsilon_{s,c}(I,i)=
 \begin{cases}
 d_{s,c}(I,i),&s\ {\rm odd},\\
 [\min(h_{s,c},i)-d_{s,c}(I,i)]_+,&s\ {\rm even}.
 \end{cases}
 \tag{3.28.1}\label{eq:p3c3-column-shift}
\]

\begin{theorem}[The shifted ordinary bound and the selected-column bound]\label{B3.28}
For \(v,j,I,n,U\) as in Theorem~\ref{B3.27},
\[
 \nu_c([\tau^n]\mathcal J_{I,v}(U))\ge B_v-E_v.
 \tag{3.28.2}
\]
For \(1\le i\le r\),
\[
 \nu_c\!\left([\tau^n]\left.
 P_{I,v}^{(i)}(\partial_Z)\mathcal C_j(\exp(-\Phi_v)U)\right|_{Z=0}\right)
 \ge B_v.
 \tag{3.28.3}\label{eq:p3c3-selected}
\]
\end{theorem}

\begin{proof}
Apply the second ordinary estimate of Theorem~\ref{B3.26} to the diagram and scaled expansion of Theorem~\ref{B3.27}. Its exponent is \(B_v-E_v\). A determinant with repeated target labels vanishes. Otherwise \(d_{s,c}(I,i)\le\min(h_{s,c},i)\) counts distinct labels below \(i\). Thus \eqref{eq:p3c3-column-shift} supplies the selected-column exponent in Theorem~\ref{B3.26}, which gives \eqref{eq:p3c3-selected}.
\end{proof}

\section{Products of positions and the class expansion}

Put \(P_{I,\mathbf i}=\prod_vP_{I,v}^{(i_v)}\) for \(\mathbf i=(i_v)_v\), \(0\le i_v\le r\), with \(P_{I,v}^{(0)}=P_{I,v}\). For the framed series \(\Theta=\Theta_{j,L,L',\phi}\) of Definition~\ref{B3.05a}, put
\[
 \mathcal J_{I,\mathbf i}(\Theta)
 =\left.P_{I,\mathbf i}(\partial_Z)\mathcal C_j(\Theta)\right|_{Z=0}.
 \tag{3.29.1}\label{eq:p3c3-product-evaluation}
\]

\begin{theorem}[Addition of the coefficient exponents]\label{B3.29}
For \(0\le j<m\), every source and target lattice pair and every frame \(\phi\), there is a series \(U_1\) integral after logarithmic substitution such that
\[
 \Theta=\left(\prod_v\exp(-\Phi_v)\right)U_1.
 \tag{3.29.2}
\]
If \(q_v\in\Q\) satisfy
\[
 \begin{aligned}
 i_v=0&\ \Longrightarrow\ q_v\text{ is an ordinary coefficient bound at }(j,v)\ \text{or}\ q_v\le B_v-E_v,\\
 i_v>0&\ \Longrightarrow\ q_v\le B_v ,
 \end{aligned}
 \tag{3.29.3}\label{eq:p3c3-choice-exponents}
\]
then \(\nu_c([\tau^n]\mathcal J_{I,\mathbf i}(\Theta))\ge\sum_v q_v\) for every \(I,n\).
\end{theorem}

\begin{proof}[Outline]
Use phase equivalence in \ref{B3.05a} to absorb the phase difference into an integral unit. Pull back logarithmic monomials by the inverse logarithm and contract by position, adding local bounds from \ref{B3.28} in products and preserving them in the finite sum. The complete proof is in \S\ref{prf:B3.29}.
\end{proof}

An index-\(i\) edge of full \(W_v\)-lattices \(A,B\) satisfies
\[
 pB\subseteq A\subseteq B,\qquad |B/A|=(p^{n'_v})^i.
 \tag{3.35.1}\label{eq:p3c3-edge}
\]
At level \(0\le b\le m\) put
\[
 (T_{v,i}^{(b)}f)(c,\mathbf L)=
 \sum_{A\to_i L_v} f(c,\mathbf L[v\leftarrow A]).
 \tag{3.35.2}\label{eq:p3c3-T}
\]
Here \(A\to_i L_v\) denotes an edge as above, and the operator on forms vanishes off lattice tuples. For \(0\le i\le r\) set
\[
 \begin{split}
 N'_{v,b}(i)&=\sum_{s\ {\rm even}}\sum_{a=1}^{i}\lambda_{v,s,a}
 +i\sum_{s\ {\rm odd}}w_{v,b}(s),\\
 \widetilde T_{v,0}^{(b)}&=1,\qquad
 \widetilde T_{v,i}^{(b)}=p^{N'_{v,b}(i)}T_{v,i}^{(b)}\quad(i>0).
 \end{split}
 \tag{3.35.4}\label{eq:p3c3-normalized-T}
\]
For a list \(v_1,\ldots,v_a\) of distinct places containing the support of \(\mathbf i\), set
\[
 \widetilde T_{\mathbf i}^{(j+1)}
 =\widetilde T_{v_1,i_{v_1}}^{(j+1)}\cdots
 \widetilde T_{v_a,i_{v_a}}^{(j+1)}.
 \tag{3.30.1}\label{eq:p3c3-product-T}
\]
 Write \(\operatorname{ker}_{j,L,L'}(G)\) for the differential kernel of Definition~\ref{B3.05a} applied to the source polynomial \(G\).

\begin{lemma}[The loss in summing lattice classes]\label{B3.30}
Let \(0\le j<m\), \(0\le i_v\le r\), and use a list as in \eqref{eq:p3c3-product-T}. Suppose that for every source and target lattice pair, every frame and every column choice \(I\), all coefficients in \eqref{eq:p3c3-product-evaluation} have lower bound \(g\in\R\) in the sense of \eqref{eq:p3c3-coeff-valuation}. Then for every \(f\in\Lambda_j\),
\[
 v_{j+1}(\widetilde T_{\mathbf i}^{(j+1)}K_jf)\ge g-d_G(j).
 \tag{3.30.2}\label{eq:p3c3-class-bound}
\]
The common class factor \(c_K\) belongs to \(O_c\) and is a unit if nonzero. Each nonterminal denominator is a stabilizer order \(|\Gamma_{j,L}|\), and each terminal denominator \(w_x\) is a positive integer. Both have \(p\)-valuation at most \(d_G(j)\).
\end{lemma}

\begin{proof}[Outline]
Use the class frames of \ref{B3.05a} and \ref{B3.10a} to put neighbours in the target fibre, then account for class denominators to get \(g-d_G(j)\). Eliminate dependent vectors using coefficients of least uniformizer order and descend to \(K\) by a coordinate minor. The complete proof is in \S\ref{prf:B3.30}.
\end{proof}

\begin{lemma}[Real exponents and the lattice endpoint]\label{B3.31}
Let \(0\le j<m\), \(0\le i_v\le r\), and let \(q_v\in\R\) satisfy \eqref{eq:p3c3-choice-exponents}. For a list as in \eqref{eq:p3c3-product-T} and every \(f\in\Lambda_j\),
\[
 v_{j+1}(\widetilde T_{\mathbf i}^{(j+1)}K_jf)\ge\sum_v q_v-d_G(j).
 \tag{3.31.1}
\]
\end{lemma}

\begin{proof}
Put \(s=\sum_vq_v\). For each \(s'<s\), choose rationals \(q'_v\le q_v\) with \(\sum_vq'_v\ge s'\). Theorem~\ref{B3.29} and Lemma~\ref{B3.30} give the bound \(s'-d_G(j)\). With \(x=e_O(s-d_G(j))\), take \(\delta=(x-\lceil x\rceil+1)/(e_O+1)>0\). Then \(\lceil x\rceil-1<x-e_O\delta\le\lceil x\rceil\). The bound at \(s'=s-\delta\) therefore gives the stated lattice by \eqref{eq:p3c3-form-valuation}.
\end{proof}

\section{The three single-step estimates}

\begin{theorem}[The ordinary lift of an integral input]\label{B3.32}
For \(0\le j<m\) and \(f\in\Lambda_j\),
\[
 v_{j+1}(K_jf)\ge\sum_vA_v-d_G(j).
 \tag{3.32.1}
\]
For every lattice pair, frame, column choice \(I\), and multi-index \(n\), the coefficient \([\tau^n]\mathcal J_{I,\mathbf0}(\Theta)\) has bound \(\sum_v A_v\). The integral-kernel hypothesis also gives
\[
 v_{j+1}(K_jf)\ge-d_G(j).
 \tag{3.32.2}\label{eq:p3c3-integral-kernel}
\]
\end{theorem}

\begin{proof}[Outline]
Take \(i_v=0,q_v=A_v\) by \ref{B3.27} and apply \ref{B3.31} to the empty list; tails in \ref{B3.13} give \(A_v\in\mathbb{Q}\), so \ref{B3.29} gives the coefficient bound. At \(\mathbf{i}=\mathbf{0}\), \ref{B3.05a} gives bound \(0\) and \ref{B3.30} gives the integral-kernel bound. The complete proof is in \S\ref{prf:B3.32}.
\end{proof}

\begin{theorem}[The distinguished-position bound for an ordinary lift]\label{B3.33}
For \(v\in\mathcal P\), \(0\le j<m\), and \(f\in\Lambda_j\),
\[
 v_{j+1}(K_jf)\ge B_v-E_v+R_j(v).
 \tag{3.33.1}
\]
\end{theorem}

\begin{proof}
Set \(q_v=B_v-E_v\) by Theorem~\ref{B3.28} and \(q_{v'}=A_{v'}\) for \(v'\ne v\). Apply Lemma~\ref{B3.31} with all indices zero.
\end{proof}

At level \(b\) and place \(v\), let \(\kappa_s\) be the row obtained by appending to \(\lambda_{v,s}\), for \(l=1,\ldots,b\) in order, \(t+1\) copies of \(-w_{v,l}(s)\). The actual weight is \(\kappa_s\) at even \(s\) and its negative reverse at odd \(s\). Put
\[
 \mathcal A_i(\kappa)=
 \sum_{s\ {\rm even}}\sum_{a=1}^{i}\kappa_{s,a}
 -\sum_{s\ {\rm odd}}\sum_{a=N_b-i+1}^{N_b}\kappa_{s,a}.
 \tag{3.34.4}\label{eq:p3c3-edge-exponent}
\]

\begin{theorem}[An output Hecke operator]\label{B3.34}
Fix \(v\in\mathcal P\). For \(0\le j<m\) and \(1\le i\le r\),
\[
 \widetilde T_{v,i}^{(j+1)}\Lambda_{j+1}\subseteq\Lambda_{j+1},\qquad
 v_{j+1}(\widetilde T_{v,i}^{(j+1)}K_jf)\ge B_v+R_j(v)
 \quad(f\in\Lambda_j).
 \tag{3.34.1--2}
\]
At any level \(0\le b\le m\), for \(0\le i\le N_b\) one has \(p^{\mathcal A_i(\kappa)}T_{v,i}^{(b)}\Lambda_b\subseteq\Lambda_b\).
\end{theorem}

\begin{proof}[Outline]
Apply \ref{B3.31} with \(B_v\) at \(v\) and \(A_{v'}\) elsewhere. Bound the rescaling exponent using row degrees and the determinant twist, tensor over embeddings, and sum over edges as in \ref{B3.30}, using the even and odd row bounds for the level estimate. The complete proof is in \S\ref{prf:B3.34}.
\end{proof}

\section{The operators and the neighbour count}

\begin{definition}[The input operator and the output polynomial]\label{B3.35}
Use \eqref{eq:p3c3-edge}, \eqref{eq:p3c3-T} and \eqref{eq:p3c3-normalized-T}. At the initial weight put
\[
 \mathcal A_v=\sum_{s\ {\rm even}}\sum_{a=1}^{r}\lambda_{v,s,a}
 -\sum_{s\ {\rm odd}}\sum_{a=k-r+1}^{k}\lambda_{v,s,a},
 \qquad y_v=p^{\mathcal A_v}T_{v,r}^{(0)}.
 \tag{3.35.3}
\]
For \(u_l,\mathfrak m_l\in K\), \(u_l\ne0\), \(1\le l\le b\), set
\[
 U_b=\prod_{l=1}^bu_l,\qquad
 \beta_{v,a}=p^{\eta_v(t-2a)}\ (0\le a\le t),\qquad
 \mathcal F_{v,b}=(\mathfrak m_l\beta_{v,a})_{1\le l\le b,\ 0\le a\le t}.
 \tag{3.35.5}
\]
With \(h_\ell\) the complete homogeneous symmetric polynomial and \(h_0=1\), put
\[
 \begin{split}
 c_{v,b,\ell}&=
 p^{N'_{v,b}(r)-N'_{v,b}(r-\ell)+\eta_v\ell(N_b-2r+\ell)}
 h_\ell(\mathcal F_{v,b})\quad(0\le\ell\le r),\\
 \widetilde T_{v,i}^{\,u,(b)}&=U_b^{-i}\widetilde T_{v,i}^{(b)},\qquad
 Z_{v,b}=\sum_{\ell=0}^r(-1)^\ell c_{v,b,\ell}\widetilde T_{v,r-\ell}^{\,u,(b)} .
 \end{split}
 \tag{3.35.6}\label{eq:p3c3-output-polynomial}
\]
\end{definition}

\begin{lemma}[The theta series as a function of a lattice indicator]\label{B3.35a}
Fix \(v\in\mathcal P\) and \(0\le j<m\). Put \(N=N_j,\ N'=N_{j+1},\ \eta=\eta_v\) and \(M=\operatorname{Mat}_{N'\times N}(\operatorname{Frac}W_v)\). For source and target lattice pairs \(L,L'\), set
\[
 I_{L,L'}(x)=\mathbf1_{xL_v\subseteq L'_v},\qquad
 s=\ell_{v,j}(L),\quad s'=\ell_{v,j+1}(L'),\quad d=Ns'-N's.
\]
Here \(\ell_{v,j}\) and \(\ell_{v,j+1}\) are the integer lengths of Definition~\ref{B3.05a} at the source and target levels. A lower edge of index \(h\) decreases the length at its place by \(h\) and preserves the lengths at the other places. Replacing either \(v\)-entry preserves \(\omega(L,L')\). There are \(a,a'\in K^\times\) of valuation zero, common to all lattice pairs at \(v,j\), and for each \(\omega\) an additive map \(\Xi_\omega:(M\to\Q)\to K_c[[Z_j]]\) with
\[
 \Theta^F_{L,L'}=a^s(a')^{s'}p^{-\eta d}\,\Xi_{\omega(L,L')}(I_{L,L'}).
 \tag{3.35.7}\label{eq:p3c3-indicator}
\]
Here \(\Theta^F_{L,L'}\) is the theta series in the fixed coordinates \(Z_j\). Lemma~\ref{B3.35a-arithmetic} gives the arithmetic charts and shows how the units and the representative law yield this identity on the selected chain.
\end{lemma}

\begin{proof}[Outline]
In an abstract system, item 4 of \ref{B3.05a} supplies the units and additive maps directly. For an arithmetic system, taking signed determinant lengths and the data outside \(v\) lets \ref{B3.35a-arithmetic} choose the units and identify chart and representative values. The complete proof is in \S\ref{prf:B3.35a}.
\end{proof}

\begin{lemma}[Changing either lattice of the arithmetic section]\label{B3.35d}
Let \(E/K_0\) be the host field over its imaginary quadratic subfield, and let \(\Psi\) consist of the embeddings of \(E\) inducing the chosen embedding of \(K_0\). Write \(L_n^0=\mathcal O_E^n\). For full lattices \(A,B\) of the same rank, let \(\mathfrak d(A,B)\) be the fractional ideal whose localization at a finite prime is generated by the determinant of a change from a basis of \(A\) to one of \(B\). Put \(c(A)=\chi(\mathfrak d(L_n^0,A))\). Here \(\chi\) is the character of norm-one ideals in Lemma~\ref{B8.56a}, and \(A\) is self-dual for the standard Hermitian form.

For self-dual \(A,B\) of even ranks \(n,n'\), respectively, put
\[
 P(A,B)=\{(\delta\tau(X))_{\tau\in\Psi}:
 X\in\mathfrak c\mathcal O_E\,\mathfrak b^{-1}
                 \operatorname{Hom}_{\mathcal O_E}(A,B)\},
 \qquad C(A,B)=c(A)^{n'/2}c(B)^{-n/2}.
\]
Here \(\mathfrak c,\delta\) are the curve ideal and period multiplier, and
\(\mathfrak b\bar{\mathfrak b}=\mathfrak D_{E/K_0}\) as in Lemma~\ref{B4.22c}.
For period lattices \(P,Q\) write
\[
 I_{P,Q}\theta=[Q:P\cap Q]^{-1/2}
                   \sum_{x\in Q/(P\cap Q)}W_x\theta ,
\]
where \(W_x\) is the theta translation including the semicharacter of \(Q\).
For the arithmetic sections \(e_{A,B}\), the section
\(\theta_{A,B}=C(A,B)^{-1}e_{A,B}\) changes by \(I_{P,Q}\) when either endpoint changes. Moreover, rational unitary changes \(g,g'\) of the two endpoints give
\(e_{gA,g'B}(Z)=e_{A,B}((g')^{-1}Zg)\), with the matrices read at the embeddings in \(\Psi\).
\end{lemma}

\begin{proof}[Outline]
Adapt the Weil moves \ref{eq:B4.22c-source} and \ref{eq:B4.22c-target} to the chosen character and divide by the new determinant factor. The ratio of the two determinant factors is invariant under rational unitary changes, so \ref{eq:B4.22c-equiv} gives the pullback. The complete proof is in \S\ref{prf:B3.35d}.
\end{proof}

\begin{definition}[Finite test of a decorated pair]\label{B3.35f}
Fix a word, a twist, chart components at stages \(i,i+1\), and their selected local Schwartz vectors. At a prime \(P\) of the word level, let \(l_P,l'_P\) be the source and target residue functionals on the corresponding lattices \(N_P,N'_P\). Give \(Z\) the factor
\[
 \mathbf1_{\{ZN_P\subseteq N'_P,\ l'_P Z=l_P\}} .
\]
At each type or twist prime \(P\), read \(Z\) in the adapted source and target frames as \(X=f'^{-1}Zf\). Put \(q=|k(P)|\), and write \(\psi\) for the selected additive character. The factor at \(P\) is as follows.

At a matrix step put \(a=N_i,\ b=N_{i+1},\ d=m-1-i,\ d'=d-1\). Use indices starting at zero. Complete \(X\) to a \((b+d')\times(a+d)\) matrix \(D\). Its first \(a\) columns are \(\binom X0\). In an extra column \(a+j\), the free rows are all rows if \(j=d'\), and the rows \(l<b+j\) if \(1\le j<d'\). In all other extra-column entries put \(\mathbf1_{l=b+j}\). Denote the set of free entries by \(F\). Let \(A,B\) be the support and period depths, and set
\[
 D_{A,B}=\{q^{-A}u:0\le u<q^{A+B},\ u\in\Z\}.
\]
For \(t\in D_{A,B}^{F}\) let \(D(t)\) be the completion with free entries \(t\). Let \(\Phi_{h,h'}\) be the selected vector transported by the source and target component matrices \(h,h'\). In enlarged dimensions \(a_1,b_1\), with local step character \(\mu\), valuation \(v(q)=1\) and chosen square root \(s_q\) of \(q\), this transport is
\[
 \Phi_{h,h'}(D)=
 \frac{\mu(\det h)}{\mu(\det h')}
 s_q^{-b_1v(\det h)+a_1v(\det h')}\Phi(h'^{-1}Dh).
\]
The local factor is
\[
 \frac{q^{-B|F|}}{\operatorname{vol}(U)}
 \sum_{t\in D_{A,B}^{F}}\Phi_{h,h'}(D(t))\,
 \psi\left(-\sum_{j=0}^{d'-1}D(t)_{b+j,a+j+1}\right).
\]
Here \(U\) is the intersection of the transported target level with the unipotent group fixing the first \(b\) coordinates pointwise, fixing the first extra basis vector and acting trivially on the graded lines of the extra-coordinate flag. Its Haar measure gives the stabilizer of the reference lattice in this unipotent group volume one.

At \(i=m-2\), take a vector \(y'\) in the marked target coset and express it in the target frame. The local factor is \(\Phi_{h,1}([X\mid y'])\). At \(i=m-1\), write \(X=\binom{X_0}{\eta}\), take \(y\) in the marked source coset in the source frame, and use
\[
 \Phi_{\rm top}(X_0)\psi(-\eta y)
 \begin{cases}
 \mathbf1_{\eta\in\mathcal O_P^a},&P\text{ a type prime},\\
 q^{-Ma}\mathbf1_{q^M\eta\in\mathcal O_P^a},&P\text{ a twist prime}.
 \end{cases}
\]
Here \(M\) is the twist depth, and \(\Phi_{\rm top}\) is the selected vector for the target with its last row removed. The product of all the factors, read in the fixed complex embedding, is \(\xi_{y,y'}(Z)\) for the adelic pair \(y,y'\). The vectors used to represent the two marked cosets are taken at their respective stages.
\end{definition}

\begin{lemma}[Class characters and the finite test in arithmetic charts]\label{B3.35e}
Fix the chart data of a constructed column chain, a word and a twist. Let \(H\subset\C\) be its global coefficient field. Denote the compact level groups at consecutive components by \(U_j,U_{j+1}\) and their characters by \(\lambda_j,\lambda_{j+1}\). Write \(\mu_{j+1}\) for the step character on norm-one finite idèles and \(\chi_I\) for the determinant-ideal character on these idèles. Put \(b_l=\lfloor l(t+1)/2\rfloor\). If level \(l\) has a component, then
\[
                 \chi_I(x)^{b_l}\in H^\times
                 \qquad(x\bar x=1).
\]
Let \(S\) be the primes of \(E\) used for the finite test, and \(c>0\) its level. The test \(\xi_{y,y'}(Z)\) at finite adelic points \(y,y'\) in fixed components depends on their components at the places below \(S\). For rational unitary \(g,g'\), expressed in standard orthogonal coordinates, and \(k\in U_j,\ k'\in U_{j+1}\), it satisfies
\[
\begin{aligned}
 \xi_{gy,g'y'}(Z)&=\xi_{y,y'}((g')^{-1}Zg),\\
 \mu_{j+1}(\det k')\lambda_j(k)\,\xi_{yk,y'k'}(Z)
 &=\mu_{j+1}(\det k)\lambda_{j+1}(k')\,\xi_{y,y'}(Z).
\end{aligned}
\]
If \(M=\mathfrak c\mathcal O_E\,\mathfrak b^{-1}\operatorname{Hom}(A,B)\) for the two kernel lattices, the test has the same value at \(Z_1,Z_2\) whenever \(Z_1-Z_2\in cM_Q\) at every \(Q\in S\).
\end{lemma}

\begin{proof}[Outline]
Use uniformizer values and the principal-ideal law to place the character powers in the coefficient field. Evaluate the local tests in transported frames using \ref{B3.35f} and reindex the Fourier sums to verify congruence invariance and the two transformation laws. The complete proof is in \S\ref{prf:B3.35e}.
\end{proof}

\begin{lemma}[One finite operator for a family of lattice pairs]\label{B3.35g}
In the arithmetic charts of Lemma~\ref{B3.35e}, fix the components at both ends and a finite family of self-dual pairs differing only at one split pair \(P,\bar P\). Let \(S\) be the support of the finite test, disjoint from \(\bar S\), and let \(c>0\) be its level. Assume \(P,\bar P\notin S\cup\bar S\) and both are prime to \(c\). The operation of applying the finite test to the arithmetic section is the restriction, at every pair in the family, of one complex linear operator on functions of the period variable. It takes entire arithmetic sections to functions with algebraic Taylor coefficients in the normalized differential coordinates of the curve.
\end{lemma}

\begin{proof}[Outline]
Use the Chinese remainder theorem and \ref{B3.35e} to get one Fourier sum for every pair. Algebraicity follows from torsion translates, the fourth-power normalization \ref{B4.22e}, endpoint averages \ref{B3.35d}, and the test values in \ref{B3.35f}. The complete proof is in \S\ref{prf:B3.35g}.
\end{proof}

\begin{lemma}[The indicator identity on the arithmetic charts]\label{B3.35a-arithmetic}
At a permitted site of Theorem~\ref{B3.45}, fix the selected chain on its prescribed filling. Use its compatible coefficient embeddings, the curve and parameter of Lemma~\ref{B8.56f}, and the finite test of Lemma~\ref{B3.35e}. Take a positive test level \(c\) and a test support disjoint from its conjugate. The split pair at each moving position lies off this support and away from the primes dividing \(c\). For the selected even ranks and self-dual kernel lattices, the chart series has \eqref{eq:p3c3-indicator}, with the length of a lattice \(GW_v^n\) at \(v\) given by \(-\operatorname{ord}_p\det G\) and with \(\omega\) recording both components and all entries away from \(v\). It satisfies the source transformation \eqref{eq:p3c3-chart-action}, and hence agrees with the series transported from source representatives.
\end{lemma}

\begin{proof}[Outline]
Compute the units from character valuations, then use finite Fourier sums and \ref{B3.35d} to map indicators to volume-scaled sections. The common test and algebraic coefficients of \ref{B3.35g} transfer relations to chart series, whose source law follows from \ref{B3.35e}. The complete proof is in \S\ref{prf:B3.35a-arithmetic}.
\end{proof}

\begin{lemma}[Evaluation at a fixed lattice pair]\label{B3.35b}
Use the notation of Lemma~\ref{B3.35a} and write \(s_w,s'_w\) for the lengths at \(w\). Set \(D(L,L')=\prod_w p^{\Delta_w(j+1)(s_w-s'_w)}\). For a source polynomial \(G\) and \(b:M\to\Q\), define
\[
 J_{L,L'}(b;G)=D(L,L')a^s(a')^{s'}p^{-\eta d}
 \mathcal D_j(G,\Xi_{\omega(L,L')}(b)),
\]
where the differential evaluation \(\mathcal D_j\) of Lemma~\ref{B3.10a} is \(K\)-linear in each variable. Then \(\operatorname{ker}_{j,L,L'}(G)=J_{L,L'}(I_{L,L'};G)\). For finitely many \(b_i:M\to\Q\) and \(n_i\in\N\),
\[
 J_{L,L'}\Bigl(\sum_i n_i b_i;G\Bigr)=\sum_i n_i J_{L,L'}(b_i;G).
\]
\end{lemma}

\begin{proof}
Substitute \eqref{eq:p3c3-indicator} in \(D(L,L')\mathcal D_j(G,\Theta^F_{L,L'})\), the fixed-coordinate expression for the kernel. Additivity gives the multiplicity formula.
\end{proof}

\begin{lemma}[Changing the endpoint in a weighted edge sum]\label{B3.35c}
Fix \(v\in\mathcal P\) and \(h\ge0\). At \(0\le j\le m\), write \(A\to_h B\) for an index-\(h\) edge at \(v\) with the same component and entries away from \(v\). Put \(w(A)=|\Gamma_{j,A}|\). For any \(K_c\)-valued function \(\phi\) invariant under simultaneous level translation,
\[
 \sum_{A\in\mathcal R_j}\frac1{w(A)}\sum_{B:\,A\to_h B}\phi(A,B)
 =\sum_{B\in\mathcal R_j}\frac1{w(B)}\sum_{A:\,A\to_h B}\phi(A,B).
 \tag{3.35.20}\label{eq:p3c3-orbit-adjoint}
\]
For \(j<m,\ f\in V_j\), and a target lattice pair \(L'\), the class data of Lemma~\ref{B3.30} satisfy
\[
 \sum_x\frac1{w_x}\sum_y\sum_{L_{x,y}\to_h B}
 \operatorname{ker}_{j,B,L'}(f(L_{x,y}))
 =\sum_x\frac1{w_x}\sum_y
 \operatorname{ker}_{j,L_{x,y},L'}((T_{v,h}^{(j)}f)(L_{x,y})).
 \tag{3.35.21}\label{eq:p3c3-hecke-adjoint}
\]
In \eqref{eq:p3c3-class-expansion}, \(x\) runs over representatives with \(w_x=w(L_x)\), and \(y\) has one value.
\end{lemma}

\begin{proof}[Outline]
Divide each orbit-pair sum by the stabilizer orders and invert the level elements to swap endpoints. Apply coefficientwise using \ref{B3.34} for the target fibre and \ref{B3.05a} for equivariance and terminal neighbours, then use linearity to obtain the input Hecke sum. The complete proof is in \S\ref{prf:B3.35c}.
\end{proof}

\begin{lemma}[Neighbour counting and the one-step Hecke relation]\label{B3.36}
Let \(W\) be a discrete valuation ring with uniformizer \(p\). Its fraction field is \(Q_W\), and its residue field has order \(q=p^{n'}\). Let \(N,g,i\) be nonnegative integers. Take full lattices \(A\subset Q_W^N,\ B\subset Q_W^{N+g}\) and a linear map \(x:Q_W^N\to Q_W^{N+g}\). Then
\[
 \begin{split}
 &\#\{B_1:pB\subseteq B_1\subseteq B,\ |B/B_1|=q^i,\ xA\subseteq B_1\}\\
 &\quad=\sum_{\ell=0}^{\min(i,g)} q^{(i-\ell)(g-\ell)}{g\brack\ell}_q
 \#\{A_1:A\subseteq A_1\subseteq p^{-1}A,\ |A_1/A|=q^{i-\ell},\
 xA_1\subseteq B\}.
 \end{split}
 \tag{3.36.1}\label{eq:p3c3-neighbour-count}
\]
Here \({g\brack\ell}_q\) is the Gaussian binomial coefficient. For \(v\in\mathcal P\) put
\(\mathcal E_{v,i}^{(b)}=p^{-\eta_vi(N_b-i)}T_{v,i}^{(b)}\).
There are \(u_l,\mathfrak a_l\in K^\times\), \(1\le l\le m\), of valuations \(0,\Delta_v(l)\), respectively, such that for \(0\le j<m,\ i\ge0,\ f\in V_j\),
\[
 \mathcal E_{v,i}^{(j+1)}K_jf
 =u_{j+1}^{\,i}K_j\left(
 \sum_{\ell=0}^i e_\ell((\mathfrak a_{j+1}\beta_{v,a})_{a=0}^{t})
 \mathcal E_{v,i-\ell}^{(j)}f\right).
 \tag{3.36.4}\label{eq:p3c3-one-step}
\]
Here \(e_\ell\) is elementary symmetric, with \(e_0=1\) and \(e_\ell=0\) for \(\ell>t+1\).
\end{lemma}

\begin{proof}[Outline]
Count graphs with prescribed intersection and image for the Gaussian convolution. Evaluate the indicator identity via \ref{B3.35a} and \ref{B3.35b}, use \ref{B3.35c} for Hecke sums preserving \(V_j\) by \ref{B3.34}, then rewrite Gaussian coefficients and collect powers. The complete proof is in \S\ref{prf:B3.36}.
\end{proof}

\begin{definition}[The exact commutation exponent]\label{B3.37}
For \(v\in\mathcal P\) and \(b\ge0\), put
\[
 P_v^{\Z}(b)=
 \sum_{s\ {\rm odd}}\left(rw_{v,b}(s)+\sum_{a=k-r+1}^{k}\lambda_{v,s,a}\right)
 +\eta_vrb(t+1).
 \tag{3.37.1}
\]
Define \(P_v(b)\) by the same sum in \(\R\), with last term \(n'_v rb(t+1)/2\). Since \(n'_v=2\eta_v\), the image of \(P_v^{\Z}(b)\) in \(\R\) is \(P_v(b)\).
\end{definition}

\section{Commutation along a prefix}

\begin{theorem}[The prefix Hecke commutation and its integral coefficients]\label{B3.38}
For each \(v\in\mathcal P\) there are \(u_l,\mathfrak m_l\in K^\times\), \(1\le l\le m\), with
\[
 v_p(u_l)=0,\qquad v_p(\mathfrak m_l)=\Delta_v(l),
 \tag{3.38.1}\label{eq:p3c3-prefix-unit}
\]
such that for \(1\le b\le m\) and \(f\in V_0\),
\[
 Z_{v,b}K^{[b]}f=p^{P_v^{\Z}(b)}K^{[b]}(y_vf).
 \tag{3.38.2}\label{eq:p3c3-prefix-commutation}
\]
The coefficients satisfy
\[
 c_{v,b,\ell}\in O\quad(0\le\ell\le r),\qquad v_p(c_{v,b,r})\ge E_v .
 \tag{3.38.3}\label{eq:p3c3-prefix-coefficients}
\]
At \(b=j+1\), every \(f\in\Lambda_j\) satisfies
\[
 v_{j+1}(Z_{v,j+1}K_jf)\ge B_v+R_j(v).
 \tag{3.38.4}
\]
\end{theorem}

\begin{proof}[Outline]
Normalize and iterate \ref{B3.36}, then extract the coefficient of \(z^r\) using the reciprocal product. Use \ref{B3.01} for integrality and extra identity-term valuation, then apply \ref{B3.34} to positive-index terms and \ref{B3.33} to the identity term. The complete proof is in \S\ref{prf:B3.38}.
\end{proof}


% ---- p3c4 ----
\chapter{Induction along the chain}\label{ch:chain-induction}

Let \(X\) be a local system at a prime \(p\) as in \ref{B3.01}, with finite place set \(\mathcal P\) and chain length \(m\). The weight conditions give \(m\ge1\) if \(\mathcal P\ne\varnothing\). Let \(K\) be the coefficient field and \(O\) the coefficient DVR with uniformizer \(\varpi\) and coefficient map \(\iota:O\to K\). Write \((p)=(\varpi^e)\) in \(O\), where \(e=e_O\ge0\).

Let \(\Lambda_j\subset M_j\) be the integral lattice in the space of forms at level \(j\), and write \(K\Lambda_j\) for its \(K\)-span. For the lifts \(K_j\), put
\[
 K^{[0]}=1,\qquad K^{[b]}=K_{b-1}\cdots K_0:M_0\longrightarrow M_b.
\]
Write \(\nu_j(F)\ge h\) for \(F\in\varpi^{\lceil eh\rceil}\Lambda_j\), and \(\nu_K(a)\ge h\) for \(a\in\varpi^{\lceil eh\rceil}O\), with scalars mapped into \(K\) by \(\iota\). Write \(\nu_K(a)=h\) when \(h\) is the greatest bound for \(a\). If \(e>0\) and \(K=\operatorname{Frac}(O)\), the scalar relation uses the valuation normalized by \(\nu_K(p)=1\). Write \(d_j=d_G(j)\) for the stabilizer loss and \([h]_+=\max(h,0)\).

For \(v\in\mathcal P\) and \(1\le b\le m\), write \(A_v(b)=g_{A,v}(b)\) and \(B_v(b)=g_{B,v}(b)\) as in \ref{B3.13}--\ref{B3.14}. Let \(P_v(b)\) be the real image of the integer exponent \(P_v^{\Z}(b)\) of \ref{B3.37}. For \(s\in T_v\), put \(\alpha_{v,s}=\lambda_{v,s,1}-\lambda_{v,s,2}\) as in \ref{B3.02}, and set \(E_v=\sum_{s\in T_v,\ s\ {\rm even}}\alpha_{v,s}\).

Fix kernel data \(Y\) for \(X\). The lifts \(K_j\) and input operators \(y_v\) are \(K\)-linear. So are the normalized output operators \(\widetilde T_{v,i,b}\) of \ref{B3.35}. The \(y_v\) commute. Write \(\mathcal L_{v,j}(h)\) for the single-place bound of Definition~\ref{B3.12}: for every series allowed there and every column choice, each frame coefficient \(z\) of the differentiated local kernel at \((j,v)\) satisfies \(z^bp^{-a}\in O_c\) whenever \(a\in\Z\), \(b\in\N_{>0}\), and \(a\le bh\).

For each place, choose sequences \((u_{v,a})_{a=1}^m\) and \((\mathfrak m_{v,a})_{a=1}^m\) in \(K\). For \(1\le b\le m\), use the coefficients \(c_{v,b,\ell}\) of (3.35.6) in \ref{B3.35} for these sequences. Put \(U_v(b)=\prod_{a=1}^b u_{v,a}\), and let \(r\) be the system parameter. The operator of \ref{B3.35} is
\[
 Z_{v,b}=\sum_{\ell=0}^r(-1)^\ell c_{v,b,\ell}
 U_v(b)^{-(r-\ell)}\widetilde T_{v,r-\ell,b},
 \qquad \widetilde T_{v,0,b}=1.
\]
At a fixed layer \(1\le b\le m\), the single-place relation is
\begin{equation}\label{eq:chain-prefix-relation}
 Z_{v,b}K^{[b]}F
 =p^{P_v^{\Z}(b)}K^{[b]}y_vF
 \qquad(v\in\mathcal P,\ F\in K\Lambda_0).
\end{equation}

\section{Several places and the input lattice}

For an ordered list \(L=(v_1,\ldots,v_a)\), put
\[
 Y_L=y_{v_1}\cdots y_{v_a},\qquad
 Z_{L,b}=Z_{v_1,b}\cdots Z_{v_a,b}.
\]
The rightmost operator acts first, and empty products are identities. For a subset \(C\) of the places in \(L\), let \(Z_{C,b}=Z_{L|_C,b}\), where \(L|_C\) retains the order in \(L\). For a set \(C\) of places use \(Y_C=\prod_{v\in C}y_v\), which is well defined by commutativity.

\begin{lemma}[Simultaneous operators at distinct places]\label{B3.39}
Let \(L\) be a list of distinct places. Fix \(1\le b\le m\). Suppose \(p\ne0\) in \(K\) and every \(y_v\) preserves \(K\Lambda_0\). Assume \eqref{eq:chain-prefix-relation} at this layer. Then for \(F\in K\Lambda_0\),
\begin{equation}\label{eq:chain-multiple-commutation}
 Z_{L,b}K^{[b]}F=p^{\sum_{v\in L}P_v^{\Z}(b)}K^{[b]}Y_LF.
\end{equation}
For the estimate, fix \(0\le j<m\). Assume \(K=\operatorname{Frac}(O)\). At this step, assume the joint coefficient estimate of \ref{B3.29} under (3.29.3) and the coefficient-to-lattice implication of \ref{B3.30}. Suppose that for every \(v\in\mathcal P\),
\[
 u_{v,a}\in O^\times\quad(1\le a\le m),\qquad
 \nu_K(c_{v,j+1,\ell})\ge0\quad(0\le\ell\le r),\qquad
 \nu_K(c_{v,j+1,r})\ge E_v.
\]
Let \(q_v\in\R\) for every \(v\in\mathcal P\). Suppose \(q_v\le B_v(j+1)\) for \(v\in L\) and \(\mathcal L_{v,j}(q_v)\) holds for \(v\notin L\). Then every \(F\in\Lambda_j\) satisfies
\begin{equation}\label{eq:chain-multiple-bound}
 \nu_{j+1}\bigl(Z_{L,j+1}K_jF\bigr)\ge\sum_{v\in\mathcal P}q_v-d_j.
\end{equation}
\end{lemma}

\begin{proof}[Outline]
Induct on the places using the single-place identity. Expand the output, lower the bound by \(E_v\) at zero indices, and apply the real-exponent argument of \ref{B3.31} with \ref{B3.29} and \ref{B3.30}, restoring the subtracted amounts with the scalar coefficient. The complete proof is in \S\ref{prf:B3.39}.
\end{proof}

\begin{definition}[The lattice allowing enhanced estimates]\label{B3.40}
For \(V\subseteq\mathcal P\) and \(S:\mathcal P\to\R\), let \(I(S,V)\subseteq M_0\) be the \(O\)-span of the vectors
\begin{equation}\label{eq:chain-generator}
 y_{v_1}^m\cdots y_{v_a}^m
 \left(\iota(\varpi)^{\,\sum_{v\in V\setminus C}\lceil eS_v\rceil}F_0\right),
 \qquad C\subseteq V,\quad F_0\in\Lambda_0,
\end{equation}
where \((v_1,\ldots,v_a)\) ranges over enumerations of \(C\) without repetition.
\end{definition}

Commutativity gives \(y_{v_1}^m\cdots y_{v_a}^m=Y_C^m\).

\section{Prefix induction and the maximum at each layer}

For Lemma~\ref{B3.41} and Theorem~\ref{B3.42}, assume \(K=\operatorname{Frac}(O)\) and \(p\ne0\) in \(K\). Require \(y_v(\Lambda_0)\subseteq\Lambda_0\) for every \(v\). Assume the joint coefficient estimate and coefficient-to-lattice implication used in Lemma~\ref{B3.39} at every step. Use sequences satisfying \eqref{eq:chain-prefix-relation} at every layer and the unit and coefficient bounds of Lemma~\ref{B3.39} at every step.

\begin{lemma}[Prefix estimate for an arbitrary selection rule]\label{B3.41}
Let \(L\) be a list of distinct places. Let \(G^+,G^0:\mathcal P\times\N\to\R\) satisfy, for \(0\le j<m\),
\[
 G^+_v(j+1)\le B_v(j+1)\quad(v\in L),\qquad
 \mathcal L_{v,j}\bigl(G^0_v(j+1)\bigr)\quad(v\in\mathcal P).
\]
For each \(0\le j<m\), choose a subset \(C_j\) of the places in \(L\) and put
\begin{equation}\label{eq:chain-step-budget}
 \beta_j=
 \sum_{v\in C_j}G^+_v(j+1)+\sum_{v\notin C_j}G^0_v(j+1)
 -d_j-\sum_{v\in C_j}P_v(j+1).
\end{equation}
Then, for every \(0\le b\le m\) and \(F\in\Lambda_0\),
\begin{equation}\label{eq:chain-prefix-bound}
 \nu_b\bigl(K^{[b]}Y_L^bF\bigr)\ge\sum_{j=0}^{b-1}\beta_j.
\end{equation}
\end{lemma}

\begin{proof}[Outline]
Induct on the prefix length, commuting the complementary operators into the integral input. Use \ref{B3.39} to replace the selected operators and divide out the resulting power of \(p\), then apply the ceiling inequality to gain \(\beta_b\). The complete proof is in \S\ref{prf:B3.41}.
\end{proof}

\begin{theorem}[Chain bound with external numerical estimates]\label{B3.42}
Fix \(V\subseteq\mathcal P\) and \(S:\mathcal P\to\R\). Let \(G^+,G^0,P^*:\mathcal P\times\N\to\R\). Suppose
\begin{align*}
 G^+_v(b)&\le B_v(b),&
 P_v(b)&\le P^*_v(b)
 &&(v\in V,\ 1\le b\le m),\\
 \mathcal L_{v,j}&\bigl(G^0_v(j+1)\bigr)
 &&&&(v\in\mathcal P,\ 0\le j<m),\\
 \sum_{b=1}^m[G^+_v(b)-P^*_v(b)-G^0_v(b)]_+&\le S_v
 &&&&(v\in V).
\end{align*}
Then every \(F\in I(S,V)\) satisfies
\begin{equation}\label{eq:chain-final-bound}
 \nu_m(K^{[m]}F)\ge T,
\end{equation}
where
\begin{equation}\label{eq:chain-final-budget}
 T=\sum_{b=1}^m\left(
 \sum_{v\in V}\max\{G^+_v(b)-P^*_v(b),G^0_v(b)\}
 +\sum_{v\notin V}G^0_v(b)\right)-\sum_{j=0}^{m-1}d_j.
\end{equation}
\end{theorem}

\begin{proof}[Outline]
Apply \ref{B3.41} to each generator of \ref{B3.40}, selecting places in \(C\) with \(G_v^+-P_v^*\ge G_v^0\). Use \(P_v\le P_v^*\) for the bound on \(C\), cover rounded shortfalls off \(C\) with the uniformizer powers, and extend by linearity. The complete proof is in \S\ref{prf:B3.42}.
\end{proof}

\section{The scalar budget and the local theorem}

\begin{lemma}[Budget for enhancement]\label{B3.43}
Fix \(p,n',k_1,t,m,c,r\in\N\) and \(T\subseteq\{0,\ldots,n'-1\}\). Let the functions
\[
 a,b:\{0,\ldots,n'-1\}\to\Z,\quad
 \lambda:\{0,\ldots,n'-1\}\times\{1,\ldots,k_1+1\}\to\Z,\quad
 w:\N\times\{0,\ldots,n'-1\}\to\Z
\]
satisfy the standard weight conditions of \ref{B3.01}. In particular \(p\ge5\) is prime, \(n'\ge2\) is even, \(r\ge1\), \(2r\le k_1+1\), and \(m\ge1\). With the gains \(g_A(j),g_B(j)\) of \ref{B3.13}--\ref{B3.14} and the price \(P(j)\) of \ref{B3.37}, set
\[
 S=\sum_{j=1}^m[g_B(j)-P(j)-g_A(j)]_+,\qquad
 \gamma_p=\frac{p}{p^2-1}.
\]
Then \(0\le S\le m\gamma_p\sum_{s\in T}\alpha_s\).
\end{lemma}

\begin{proof}[Outline]
Dominance and monotonicity make capacities in \ref{B3.02} nonnegative, so \ref{B3.13}, \ref{B3.14}, and \ref{B3.03} give \(g_A(j)\ge-\delta_j\) and \(g_B(j)\le\gamma_p\sum_{s\in T}\alpha_s-\delta_j\). Since \(P(j)\ge 0\) by \ref{B3.37}, summing over \(j\) gives the bound. The complete proof is in \S\ref{prf:B3.43}.
\end{proof}

\begin{theorem}[Chain bound for the original local parameters]\label{B3.44}
Assume \(\mathcal L_{v,j}(A_v(j+1))\) for every \(v\in\mathcal P\) and \(0\le j<m\). Assume also that for every choice of \(G^+,G^0,P^*,S,V\) satisfying the four conditions of Theorem~\ref{B3.42}, inequality~\eqref{eq:chain-final-bound} holds for every \(F\in I(S,V)\). For each \(v\), set
\[
 S_v=\sum_{b=1}^m[B_v(b)-P_v(b)-A_v(b)]_+.
\]
Then \(S_v\le m\gamma_p\sum_{s\in T_v}\alpha_{v,s}\), and for every \(V\subseteq\mathcal P\) and \(F\in I(S,V)\),
\begin{equation}\label{eq:chain-original-bound}
 \nu_m(K^{[m]}F)\ge
 \sum_{b=1}^m\left(
 \sum_{v\in V}\max\{B_v(b)-P_v(b),A_v(b)\}
 +\sum_{v\notin V}A_v(b)\right)-\sum_{j=0}^{m-1}d_j.
\end{equation}
\end{theorem}

\begin{proof}
Apply \ref{B3.43} at each place. In the assumed external bound, take \(G^+_v=B_v\) and \(G^0_v=A_v\). Set \(P^*_v=P_v\). The local coefficient condition holds by hypothesis, and the three inequalities of \ref{B3.42} hold with equality. Substitution into \eqref{eq:chain-final-budget} gives \eqref{eq:chain-original-bound}.
\end{proof}

\begin{definition}[Permitted pairs and the prescribed initial model]\label{B3.45a}
Fix the host of Definition~\ref{B2.01}, a word \(u\) with level prime \(T_u\), and a finite coefficient place \(w\) of residue characteristic \(p\). Call \((u,w)\) permitted if \(p\ge5\) is of inert class as in Definition~\ref{B3.46}, lies outside the host level, does not divide \(T_u\), and the second component of every arithmetic Frobenius at every host prime above \(p\) has even order. At a permitted pair write \(G_{u,w}\) for the initial model in the arithmetic datum. It gives the coefficient map \(O\to K\), the local places and weights, and the finite component set. Each component has a level group acting on its local lattices and a character with values in \(O^\times\). The model also gives the nonnegative integer used to compare the two integral weight lattices.
\end{definition}

\begin{theorem}[The local-system theorem]\label{B3.45}
 For every prime \(p\), each local system \(X\) with kernel data \(Y\) satisfying Definition~\ref{B3.05a} has the following eleven properties.

 The first is the basic lift bound \(-d_j\) of \ref{B3.32}. The next two are the single-place and all-place product expansions of \ref{B3.09}. The fourth and fifth are the ordinary lift bound \(\sum_v A_v(j+1)-d_j\) and the full differentiated-kernel coefficient bound \(\sum_v A_v(j+1)\) of \ref{B3.32}. The sixth is the family of ordinary local coefficient bounds of \ref{B3.27}.

The seventh is the distinguished-place bound of \ref{B3.33} for the unmodified lift. The eighth asserts that, for \(1\le i\le r\), \(\widetilde T_{v,i,j+1}\) preserves \(\Lambda_{j+1}\) and satisfies the bound of \ref{B3.34} on \(K_j\Lambda_j\). For each \(v\), the ninth assertion supplies sequences \(u_{v,a},\mathfrak m_{v,a}\) such that, for \(1\le a\le m\),
\[
 \nu_K(u_{v,a})=0,\qquad
 \nu_K(\mathfrak m_{v,a})
 =\sum_{s\ {\rm even}}w_{v,a}(s)-\sum_{s\ {\rm odd}}w_{v,a}(s),
\]
where the sums run over the local positions \(0\le s<n'_v\). For these sequences it gives the prefix identity and coefficient estimates of \ref{B3.38} and the bound
\[
 \nu_{j+1}(Z_{v,j+1}K_jF)\ge
 B_v(j+1)+\sum_{v'\ne v}A_{v'}(j+1)-d_j
 \qquad(v\in\mathcal P,\ 0\le j<m,\ F\in\Lambda_j).
\]
The tenth assertion gives the two bounds in \ref{B3.44}. The eleventh gives \eqref{eq:chain-final-bound} for every \(F\in I(S,V)\) whenever the four conditions of Theorem~\ref{B3.42} hold. For every choice of local weights satisfying the hypotheses of Lemma~\ref{B3.43}, the enhancement budget also satisfies \(S\le m\gamma_p\sum_{s\in T}\alpha_s\).

At each permitted pair of Definition~\ref{B3.45a}, the eleven assertions hold for every admissible system at \(p=\operatorname{char}k(w)\) with initial model \(G_{u,w}\).
\end{theorem}

\begin{proof}[Outline]
Use \ref{B3.12} for the fraction field and \ref{B3.05a} for lifts in \ref{B3.38}. Use \ref{B3.42} with \ref{B3.29}–\ref{B3.30} for \(e>0\), iterate \ref{B3.32} for \(e=0\), combine \ref{B3.27} with the chain bound for \ref{B3.44}, and use \ref{B3.43} for the budget. The complete proof is in \S\ref{prf:B3.45}.
\end{proof}

% ---- p3c5 ----
\chapter{Global Frobenius-orbit parameters}\label{ch:global-frobenius-orbit-parameters}

Fix arithmetic data \(D\) and size data \(L\) in Alp\"oge's construction \cite{alpoge-2026}. Let \(X\) consist of the coset representatives and offsets of \ref{B3.46}--\ref{B3.47}, a real share \(\kappa_I(p,n',S,\mathbf a,\mathbf b)\) and an exponent \(e_p\in\N\). Here \(p,n'\in\N\), \(S\subseteq I_{n'}=\{i\in\N:i<n'\}\), and \(\mathbf a,\mathbf b:I_{n'}\to\Z\). Write \(H\) for the additive group \(A\) of the frame in \ref{B2.01}, \(n=|H|\) and \(f=2n\). Let \(k\) be the starting rank, \(m\) the number of lifts and \(t+1\) the rank increase per lift. The rank after \(j\) lifts is \(N_j=k+j(t+1)\), and \(f\) is the number of embedding indices. Let \(TT\subseteq H\) be the host's set of top positions. Write \(a_s,b_s\) for the two integer weights prescribed by \(L\) at \(s\in TT\), extended by zero on \(H\setminus TT\).

Let \(x>1\) and \(0<\varepsilon<1/52\) be the window parameters. Set
\[
 b=\exp(x^\varepsilon),\quad a_\ast=xb,\quad
 X'=\exp(x^{\varepsilon/3}),\quad b'=\frac{a_\ast}{X'},\quad
 w''=\frac{b'}{xn},\quad X''=\frac{a_\ast}{w''}.
\]
Let \(\Delta>0\) be one quarter of the spacing between consecutive entries of the lower interlacing row of the arithmetic Harish--Chandra parameter. Put \(S=4k\Delta\) and \(c=4S\). Let \(\eta>0\) be the offset width in \ref{B3.47}, and put \(p_0=k^2x^{26}\), \(\delta_1=x^{-\varepsilon/8}\) and \([y]_+=\max(y,0)\). The top-weight window and \(k\ge5\) give
\[
 a_\ast\le a_s\le a_\ast+\frac{a_\ast}{kx^{13}}\le\frac65a_\ast,
 \qquad b_s\ge b'\qquad(s\in TT).
 \tag{0}\label{eq:p3c5-top-window}
\]

\section{Orbits and global weight diagrams}

\begin{definition}[Frobenius orbits and marked positions]\label{B3.46}
Use the host isomorphism \(\Gal(E/\Q)\simeq G=\{\pm1\}^2\times H\) of \ref{B2.01}, writing the last coordinate additively. For each rational prime \(p\), let \(\mathcal F_p\subseteq G\) be the set of images of arithmetic Frobenius elements at primes of \(E\) above \(p\). Fix \(g_p\in\mathcal F_p\) and let \(h=h_p\) be its \(H\)-coordinate.

Write \(q,d,Q_S\) for the host parameters. For \(p\nmid Q_Sqd\), take the first sign coordinate of the Artin image of \(p\bmod Q_Sqd\). At the other primes set this sign to zero. For \(p\nmid Q_S\) with \(p\ne q,d\) and \(\bigl(\frac{-d}{p}\bigr)\ne1\), the signal and inert classes have signs \(1\) and \(-1\), respectively. With \(\operatorname{ord}\) denoting additive order, set
\[
 (\Delta_p,n_p)=
 \begin{cases}
 (2h,2\operatorname{ord}(2h)),&p\text{ of signal class},\\
 (h,\operatorname{ord}(h)),&\text{otherwise}.
 \end{cases}
\]
For \(a,s\in H\), \(X\) supplies a representative \(\beta_a(s)\in s+\langle a\rangle\) with \(\beta_a(s+a)=\beta_a(s)\). Put
\[
 \mathcal R_p=\beta_{\Delta_p}(H),\qquad
 T(t_0)=\{s\in TT:\beta_{\Delta_p}(s)=t_0\},\qquad N_v(t_0)=|T(t_0)|.
\]
For \(s\in H\) and \(e\in\N\), the labelled partner is
\[
 z_e(s)=\bigl(\mathbf 1_{\{p\text{ of signal class},\ e\text{ odd}\}},s+eh\bigr)
 \in\{0,1\}\times H.
\]
Call \(s\) even on its orbit if \(\beta_{\Delta_p}(s)+ih=s\) for some even integer \(0\le i<n_p\).
\end{definition}

\begin{definition}[Global capacities and logarithmic error]\label{B3.47}
A position \((0,s)\) is marked if \(s\in TT\). On the second sheet the thick positions are \((1,u_0+ia_0)\) for \(i\in T_2\), where \(a_0\) is the distinguished element of \ref{B2.01}, \(T_2\subset\N\) is the host's thick index set and \(u_0\in H\) is its chosen translation. The remaining positions are thin. For \(j\in\N\), \(X\) supplies offsets \(o_j(z)\) in
\[
 [b'/x,b'/x+\eta],\qquad [b,b+\eta+1],\qquad [w'',w''+\eta]
\]
at marked, thick, and thin positions, respectively. Define
\[
 W_j(e,s)=\mathbf 1_{\{e=0,\ s\in TT\}}b_s+o_j(e,s),
 \qquad
 \kappa_j=\frac{k+(2j-1)(t+1)}2,\qquad B_j(z)=W_j(z)-\kappa_j .
\]
At layer \(j\) put \(\bar\alpha_s=a_s-2c\), \(\bar C_{s,e}=B_j(z_e(s))+c\), and
\[
 z_p=\log_p(2a_\ast),\qquad
 \bar\delta(t_0)=K_{\delta,1}n_pN_m^2z_p+K_{\delta,2}N_v(t_0)z_p^2.
\]
The real constants \(K_{\delta,1}\) and \(K_{\delta,2}\) weight the period and marked-position terms in this error. Write \(K\) for the family of global real constants and \(K_{\mathrm{Pr}}\) for its price coefficient in \eqref{eq:p3c5-price}.
\end{definition}

\begin{definition}[Ordinary global gain]\label{B3.48}
For \(s,s'\in H\), let \(d_p(s,s')\) be the least \(e\ge1\) with \(z_e(s)=(0,s')\), taking \(d_p(s,s')=\infty\) if there is no such \(e\). Let \(\mathcal E_p(s)\) consist of the finite odd distances \(d_p(s,s')\) for \(s'\in TT\setminus\{s\}\). Define
\[
 g_0(t_0,j)=
 \sum_{s\in T(t_0)}
 \sum_{\substack{e\ge1,\ e\ {\rm odd}\\e\notin\mathcal E_p(s)}}
 \left[\frac{\bar\alpha_s}{p^e}-\bar C_{s,e}\right]_+-\bar\delta(t_0).
\]
\end{definition}

\begin{definition}[Enhanced global gain]\label{B3.49}
Let \(T^V(t_0)\) consist of the marks in \(T(t_0)\) that are even on their orbit in the sense of \ref{B3.46}. Put \(\gamma_p=p/(p^2-1)=\sum_{e\ge1,\ e\ {\rm odd}}p^{-e}\). Define
\[
 g_U(t_0,j)=\gamma_p\sum_{s\in T(t_0)}\bar\alpha_s
 -\sum_{s\in T(t_0)\setminus T^V(t_0)}
 \sum_{\substack{e\ge1\\e\ {\rm odd}}}
 \min\left(\frac{\bar\alpha_s}{p^e},\bar C_{s,e}\right)-\bar\delta(t_0).
\]
\end{definition}

Both gains are defined for \(t_0\in H\) and \(j\in\N\), with nonsummable series assigned sum zero. For \(C\ge0\),
\[
 \max\{[u/p^e-C]_+,|\min(u/p^e,C)|\}\le |u|p^{-e}.
 \tag{1a}\label{eq:p3c5-summable}
\]
If \(\bar C_{s,e}\ge0\) for every \(s\in T(t_0)\setminus T^V(t_0)\) and every positive odd \(e\), \eqref{eq:p3c5-summable} gives absolute convergence. Summing \(v-\min(v,C)=[v-C]_+\) over these indices gives
\[
 g_U(t_0,j)+\bar\delta(t_0)
 =\gamma_p\sum_{s\in T^V(t_0)}\bar\alpha_s
 +\sum_{s\in T(t_0)\setminus T^V(t_0)}
 \sum_{\substack{e\ge1\\e\ {\rm odd}}}
 \left[\frac{\bar\alpha_s}{p^e}-\bar C_{s,e}\right]_+ .
 \tag{1}\label{eq:p3c5-enhanced}
\]

\section{Prices and selected orbits}

Put \(h_G=2+8\varepsilon\), \(K_G=\lceil8(h_G+2)/\varepsilon\rceil\), \(r_G=K_G+1\), and \(M_{\mathrm{Pr}}=w''+\eta+c+N_m\).

\begin{definition}[Embedding and filling identifications]\label{def:p3c5-orbit-model}
The host labels the embedding indices \(\tau\) bijectively by the elements of \(G\) with second sign \(+1\). Write \(\ell(\tau)\) for this label and \(\ell(e,s)=(((-1)^e,1),s)\) for the label of a position on sheet \(e\in\{0,1\}\).

An \emph{embedding identification} gives, for each word \(u\) and coefficient place \(w\) permitted by \ref{B3.45}, a bijection from the local positions \((v,i)\) to the embedding indices. The base weight at \((v,i)\) has \(k\) entries and equals the arithmetic level-zero weight for \(u\) at the corresponding embedding. Let \(p=\operatorname{char}k(w)\). For every prime of \(E\) above \(p\) and every arithmetic Frobenius there, let \(g\in G\) be its image. The period at each \(v\) is the order of the \(H\)-coordinate of \(g\). Advancing \(i\) by one multiplies the embedding label by \(g\) and by the image of complex conjugation.

A \emph{filling identification} for these bijections requires the local filling to have length \(m\), parameter \(t\) and row bound at most \(c\). The edge rank is \(r_G\) for \(p<p_0\) and \(1\) for \(p\ge p_0\). The position map \(t_0\mapsto v(t_0)\) is injective on \(\mathcal R_p\). A local position with label \(\ell(1,s)\) is unmarked and has shift \(B_j(1,s)\) at every layer \(1\le j\le m\). For \(t_0\in\mathcal R_p\), the period at \(v(t_0)\) is \(\operatorname{ord}(h_p)\). The orbit coordinate at this representative is \(\iota_{t_0}:\Z/n_p\Z\to H\), \(\iota_{t_0}(i)=t_0+ih_p\). Write \(\iota(i)\) when \(t_0\) is fixed. Its image is the coset of \(t_0\), and equality of two values means that the indices differ by a multiple of \(\operatorname{ord}(h_p)=n_p\). At index \(i\) put \(s=\iota(i)\). The label is \(\ell(0,s)\), the shift is \(B_j(0,s)\), and the position is marked exactly when \(s\in TT\). At such a mark the two local weight parameters are \(a_s,b_s\), so the extreme base-weight entries are \(a_s,-b_s\). A \emph{choice of ties} for the orbit data and a local model \(Y\) is a choice of the embedding bijections at all permitted pairs for which the filling identification above holds. Call \(Y\) with these identifications a \emph{tied model}. The quantifier over ties ranges over all such choices of bijections.
\end{definition}

\begin{definition}[Size bounds at a window]\label{def:p3c5-size}
Use the following notation for the two sets of numerical hypotheses in this chapter:
\[
 B=x^\varepsilon,\quad \theta=\varepsilon/100,\quad U=x^\theta,\quad
 \nu=2n\delta_1,\quad R_0=n/B,\quad r=k+m-1,\quad N=N_m,\quad
 M=Q_Sqd .
\]
Let \(\Lambda\) be the real size parameter for the bottom-weight window, \(\delta_2=|T_2|/n\) the thick density and \(L_w\) the length used to exclude short relations in \(T_2\). Put
\[
 a_0'=\frac{kB}{m\delta_2},\qquad
 l_*=\frac{mb}{nt^2(N+t)B^2+kn^2(S+t+B)},\qquad H_0=\frac{\eta}{m}.
\]
Write \(d_E\) for the root discriminant of \(E\), \(\varepsilon_0=\varepsilon/J\) for the host exponent with prescribed \(J\ge1\), and \(C_t\) for the prescribed exponent in \(t\le x^{C_t}\). Let \(\chi_{-d}\) be the quadratic character modulo \(d\). Write \(b_i(\tau),e_i(\tau)\), \(0\le i<m\), for the two rows of layer parameters at \(\tau\); these \(b_i(\tau)\) are distinct from the bottom weights \(b_s\) on \(TT\). Let \(p_1^*,p_2^*\) denote the host's two auxiliary level primes congruent to \(1\bmod 15\), and let \(\mathcal W\) be the set of words.

The following bounds are common to the size and parameter hypotheses:
\[
\begin{gathered}
 k\le\sqrt{x}\,nR_0,\quad \sqrt{k/m}\le R_0,\quad
 \max\{m^2B,m\sqrt{x}\log x\}\le k,\\
 \log f\le\sqrt{m/2-1},\quad (\log f)^2\le m^2B,\quad
 M^{36}t^{150}\le\sqrt b,\quad x^{\varepsilon_0/6}\le\log q,\quad
 |\Re(L'/L)(1,\chi_{-d})|\le B,\\
 |TT|\le\nu,\quad b_s\le (1+\nu B\Lambda/k)b'\quad(s\in TT),\quad
 4S\le\Lambda b'/k,\quad N m^{10}\le\Delta,\quad
 \log\Delta\le x^{\varepsilon/4}/2,\quad kS\le w'',\\
 t\le\eta\le b/(t^3fk^3),\quad \log\eta\le x^{\varepsilon/4},\quad
 t\eta\le w''/x,\quad m(t+1)\eta\le w'',\quad
 m(S+\eta)\le w''B,\\
 N\eta\le\delta_2mw'',\qquad
 \eta^2rB\le\delta_2m(w'')^2,\qquad
 r(b'+S)(\log x)a_0'\le\delta_2mb,\\
 b'+S+\eta+b/l_*<l_*,\qquad 2\eta+t+8b/l_*+7<l_*,\\
 \max\{12b/l_*+2t+8,\,ktB,\,4N^3+2t+6,\,
             2(r+mt)^3+t+3,\,Nm^{10}\}\le H_0 .
\end{gathered}
\]
At every embedding require \(b_{m-1}(\tau)\le2b_0(\tau)\). For all embeddings \(\tau,\tau'\) and \(0\le i,i'<m\), require \(e_i(\tau)\le X''e_{i'}(\tau')\). Also require \(\log p_a^*\le B\) for \(a=1,2\) and \(|\mathcal W|\le2^{k-3}\).

The \emph{size bounds} impose the common bounds and the following additional conditions. Each of \(n,m,k\) lies between \(U^{-1}\) and \(U\) times the respective quantity
\[
 x^{2+8\varepsilon},\qquad \sqrt{x\delta_1},\qquad x^{9/2+10\varepsilon}.
\]
For the positive constant \(C_\Gamma\) in the size data, normalized by
\(\log\Gamma_u=C_\Gamma m^5 fB\) for the prescribed column parameter \(\Gamma_u\), require
\[
\begin{gathered}
 \max\{Um^5 f B^2,\,8C_\Gamma m^5 f B\}\le k,\quad
 Um^6B\le k\log k,\quad Um^3B\log d_E\le k(\log k)^2,\\
 U\log\log b\le rf\log r,\quad UB\le f\log k,\quad
 f\le\min\{m^2n,r^2\},\quad 1\le t\le Ux^{C_t}.
\end{gathered}
\]
For \(a=q,d\) require \(x^{\varepsilon_0/5-\varepsilon_0/100}\le\log a\le x^{\varepsilon/5}U\). The remaining bounds are
\[
\begin{gathered}
 1\le d_E\le Q_S\sqrt{qd}/2,\quad U^{-1}n\delta_1\le|TT|,\quad
 B\le\Lambda\le UB,\quad
 4\nu B\Lambda/k+U(\log x)\Lambda/k\le x^{-2},\\
 B^{-1}\le\delta_2,\quad L_w\le(\log x)^2,\quad
 \max\{Ur^{19}f^5B,\,Ukm^{9/2}\}\le H_0,\quad UmN\le\Delta,\\
 Umk\Delta\le w'',\quad 4Uk^2\Delta\le\Lambda b',\quad
 Um\log(d_E/2)\le k,\quad \sqrt b\le w'' .
\end{gathered}
\]
For the two fixed primes require \((kb)^8\le g_i\le2g_j\) and \(\log g_i\le UB\) for \(i,j\in\{0,1\}\). The host's prescribed splitting bound for odd primes is at least \(11\).
\end{definition}

\begin{proposition}[The global price and its two bounds]\label{B3.50}
Define
\[
 \Pi_p(t_0,j)=
 \begin{cases}
 \displaystyle\sum_{s\in T(t_0)\setminus T^V(t_0)}(B_j(0,s)-b_s)
 +K_{\mathrm{Pr}}n_pM_{\mathrm{Pr}},&p\ge p_0,\\
 r_G\bigl(N_v(t_0)(2b'+c)+n_pM_{\mathrm{Pr}}\bigr),&p<p_0 .
 \end{cases} \tag{2}\label{eq:p3c5-price}
\]
Assume the size bounds of Definition~\ref{def:p3c5-size}. Put \(R=2n\delta_1x^\varepsilon\Lambda/k\). The size clauses used here and in \ref{B3.53} give
\[
 \begin{gathered}
 \Lambda\ge x^\varepsilon,\qquad b_s\le(1+R)b'\quad(s\in TT),\qquad
 4R+x^\theta(\log x)\Lambda/k\le x^{-2},\\
 1\le t\le\eta\le b/(t^3fk^3),\qquad
 m(t+1)\eta\le w'',\qquad kS\le w''.
 \end{gathered}
 \tag{2a}\label{eq:p3c5-size-clauses}
\]
Let \(Y\) satisfy the embedding and filling identifications of Definition~\ref{def:p3c5-orbit-model}, and take a word \(u\) and place \(w\) permitted by \ref{B3.45}, with \(p=\operatorname{char}k(w)\). For \(t_0\in\mathcal R_p\), let \(P_v(j)\) be the exchange exponent of \ref{B3.37} at \(v=v(t_0)\). If \(p\ge p_0\) and \(K_{\mathrm{Pr}}\ge1\), the Hecke rank is \(1\) and \(P_v(j)\le\Pi_p(t_0,j)\) for \(1\le j\le m\). If \(p<p_0\), \(r_G<k\), and \(x\ge2\), the rank is \(r_G\) and the same bound holds.
\end{proposition}

\begin{proof}[Outline]
Identify the orbit via \ref{def:p3c5-orbit-model} and use the exchange sum of \ref{B3.37}. For large primes, evaluate marked terms and bound unmarked terms by \ref{B3.47} and \ref{B3.01}, while for small primes use the size bounds at marks and bound the rows after the first. The complete proof is in \S\ref{prf:B3.50}.
\end{proof}

Put \(n'=\operatorname{ord}(h)\), \(S_{t_0}=\{0\le i<n':t_0+ih\in TT\}\), and
\[
 \begin{gathered}
 u_1=\lceil1/\delta_1\rceil,\qquad
 \ell=(6r_G+1)u_1,\qquad M_0=\lceil(h_G+3)(\log x)10^{6r_G+1}\rceil,\\
 G'=(M_0+5)\ell,\qquad
 N_{\mathrm{win}}=\left\lceil\max\{6\mathrm eG'\delta_1,(h_G+3)\log_2x\}\right\rceil+1 .
 \end{gathered} \tag{4}\label{eq:p3c5-window}
\]
Let \(L_1(p)=\lceil\log_p((kx^{26})^4)\rceil\) and \(L_2=2\sup_{p\ \text{of inert class}}L_1(p)\), with the supremum over primes. Call \(t_0\) good if \(S_{t_0}\subseteq\Z/n'\Z\) satisfies the window and tile conditions preceding Theorem~\ref{B2.54} with \((K_\varepsilon,R_\varepsilon,u,N_1)=(K_G,r_G,u_1,N_{\mathrm{win}})\) and the parameters of \eqref{eq:p3c5-window}.

\begin{definition}[Selection and counted gain]\label{B3.51}
Let \(\mathcal G_p\) be the good representatives in \(\mathcal R_p\). Write \(K_{\kappa,II}\) for the small-prime share coefficient and \(K_\Lambda,K_\rho,K_N\) for the global coefficients in the bandwidth formula below. Set \(\Sigma_p=m\sum_{t_0\in\mathcal G_p}\kappa_v(t_0)\), where
\[
 \kappa_v(t_0)=
 \begin{cases}
 \min\{1,K_{\kappa,II}\,|S_{t_0}|r_G^2 kx^{-26}\},&p<p_0,\\
 \kappa_I\bigl(p,n',S_{t_0},(a_{t_0+ih})_i,(b_{t_0+ih})_i\bigr),&p\ge p_0 .
 \end{cases}
\]
At an inert-class prime, \(N_v(t_0)=|S_{t_0}|\) for \(t_0\in\mathcal R_p\). A candidate is an inert-class prime \(p<X'\) with \(n'\ge x^{\varepsilon/4}\). For \(p\ge p_0\) impose one further condition. Every positive integral power of \(p\) must lie outside the closed interval \([X'(1-\rho_B),X'(1+\rho_B)]\), where
\[
 N_\ast=2n\delta_1,\qquad
 \mathcal L=K_\Lambda N_{\mathrm{win}}^2\log a_\ast,\qquad
 \rho_B=\frac{4N_\ast x^\varepsilon\mathcal L}{k}
 +\frac{K_\rho K_N(\log x)\mathcal L}{k}.
\]
Set
\[
 V_p=
 \begin{cases}
 \mathcal G_p,&p\text{ a candidate and }\Sigma_p\le x^{-\varepsilon},\\
 \varnothing,&\text{otherwise}.
 \end{cases}\tag{5}\label{eq:p3c5-selection}
\]
For \(1\le j\le m\), the counted gain is
\[
 g^{\rm count}_{p,j}=
 \begin{cases}
 0,&(p\text{ of inert class and }h=0)\text{ or }p>X'',\\
 \displaystyle\sum_{t_0\in V_p}
 \max\{g_U(t_0,j)-\Pi_p(t_0,j),g_0(t_0,j),0\}
 +\sum_{t_0\in\mathcal R_p\setminus V_p}[g_0(t_0,j)]_+,
 &\text{otherwise}.
 \end{cases}\tag{6}\label{eq:p3c5-gain}
\]
\end{definition}

\section{Identifying the global data}

\begin{definition}[Filling compatibility]\label{B3.52}
At a finite coefficient-field place \(w\) with \(p=\operatorname{char}k(w)\), call the fillings compatible if the following identities hold. Use the labels \(\ell(\tau)\) and \(\ell(z)\) of Definition~\ref{def:p3c5-orbit-model}. For \(1\le j\le m\), let \(w_j(\tau)\) be entry \(j-1\), numbered from zero, of the exponent row of the arithmetic Harish--Chandra parameter at \(\tau\). Then
\[
 W_j(z)=w_j(\tau)
 \quad(1\le j\le m,\ \ell(\tau)=\ell(z)), \tag{7}
\]
and
\[
 \omega_{j,p}(s,e)=W_j(z_e(s))
 \qquad(j\in\N,\ s\in TT,\ e\in\N). \tag{8}
\]
Here \(\omega_{j,p}(s,e)\) is the partner weight prescribed by \(L\). Regard the prescribed set \(T_w\subseteq TT\) of top positions at \(w\) as a subset of \(H\), and require
\[
 T_w=
 \begin{cases}
 TT,&p\text{ of signal class or }p\ge X',\\
 \{s\in TT:\beta_{\Delta_p}(s)\in V_p\},&\text{otherwise}.
 \end{cases}\tag{9}
\]
Let \(\mathcal S_L\) be the prescribed finite support of the size data and \(g_j(w)\) their gain before primes are discarded. For \(w\in\mathcal S_L\) and \(p\) of inert or signal class require
\[
 g_j(w)=g^{\rm count}_{p,j}\qquad(1\le j\le m). \tag{10}
\]
Write \(\tau_w,I_w\) for the exponent and index budgets prescribed at \(w\). Require
\[
 \tau_w=\tau_P(p),\qquad I_w=I_P(p),\qquad
 \Lambda=\mathcal L. \tag{11}
\]
Both budgets vanish if \(V_p=\varnothing\). Otherwise put
\[
 \begin{aligned}
 S_v(t_0)&=\left\lceil\sum_{j=1}^m
 [g_U(t_0,j)-\Pi_p(t_0,j)-[g_0(t_0,j)]_+]_+\right\rceil,\\
 \tau_P(p)&=e_p+
 m\sum_{t_0\in V_p}\bigl([\lceil t_v(t_0)\rceil]_++D_p+1\bigr),\\
 I_P(p)&=2\Sigma_p\max_{t_0\in V_p}S_v(t_0).
 \end{aligned}
\]
Use \((t_v(t_0),D_p)=(t_{II},D_{II})\) for \(p<p_0\) and \((t_I,D_I)\) for \(p\ge p_0\), from \ref{B3.52a} with length \(n'=n_p\), marks \(S_{t_0}\), and weights \(\mathbf a=(a_{t_0+ih})_i,\ \mathbf b=(b_{t_0+ih})_i\).
\end{definition}

\begin{lemma}[The two local presentations]\label{B3.52a}
Use \(A=a_\ast\), \(r=r_G\), \(z=\log_p A\), \(q_p=p^{G'}\), \(c_0=2b'+2c\), and \(U=40n'\delta_1\). Let \(K_1,K_2\) be the coefficients of the two error terms in \(t_I\) below, with the same \(K_2\) in \(c_6\). Write \(K_{\rm DI},K_{\rm DII},K_{\rm tII}\) for the coefficients in \(D_I,D_{II},t_{II}\), respectively. For \(N\) marks on a cycle of length \(n'\), the bounds and costs are
\[
\begin{aligned}
 D_I&=\lceil K_{\rm DI}Ux^{26}(N_{\rm win}^5+k^3)z\rceil_+,\\
 D_{II}&=\lceil K_{\rm DII}U(N_{\rm win}^3+r^5+k^2)z+2rn'c+rUc_0\rceil_+,\\
 t_{II}&=K_{\rm tII}rN(\mathcal LA/(pk)+Ap^{-L_1(p)}+c_0+kq_p)+2rn'c,\\
 c_6&=8(D_I+1)+K_2Ukz+4Uc+2n'c,\\
 t_I&=J+K_1NN_{\rm win}(A/p^2+\mathcal LA/(pk))
 +K_2N(q_pc_6+kq_p+kz+c)+2n'c .
\end{aligned}
\]
Here \(\lceil y\rceil_+=[\lceil y\rceil]_+\). Each odd mark \(y\) contributes \(\min(b_y,a_{y-d_y}/p^{d_y})\) to \(J\), where \(d_y\) is the least positive distance back to an even mark within one period. Its contribution is zero if there is no even mark. Parity uses representatives \(0,\ldots,n'-1\).

Prescribe the index exponents by \(e_p=\lceil C_e f k^2\log_p a_\ast\rceil_+\) for all \(p\in\N\), with a fixed \(C_e\in\R\). At \(p=0,1\) take the logarithm to base \(p\) to be zero. For every \(p,n'\in\N\), \(S\subseteq I_{n'}\) and \(\mathbf a,\mathbf b:I_{n'}\to\Z\), prescribe \(\kappa_I(p,n',S,\mathbf a,\mathbf b)\) as the infimum of the nonnegative \(\kappa\) passing the following test for every eligible model and dominant base weight. This infimum passes the test. The test takes \(r',t_\ast,D_0,c_*\in\N\) with \(4(D_0+1)\le c_*\). For this share use
\[
 (r',t_\ast,D_0,c_*)=(1,\lceil t_I\rceil_+,D_I,8(D_I+1)).
\]
For \(p<p_0\), \ref{B4.07a} uses \((r_G,\lceil t_{II}\rceil_+,D_{II},8(D_{II}+1))\) and the fraction of \ref{B3.51}.

The models use complete DVRs \(W\) with uniformizer \(p\ge5\) prime, residue order \(p^{n'}\), \(n'\ge2\) even, and Frobenius lift \(\sigma\) with \(\sigma^{n'}=1\). Take all finite extensions \(W\to O\) of complete DVRs with compatible fraction-field inclusion, uniformizer \(\omega\) satisfying \(\omega^2=p\), residue order \(p^{2n'}\), and square roots in \(O\) of all rational integers. A dominant base weight \(\lambda_{s,a}\) has \(k\) entries at each \(s\), numbered from zero. At \(s\in S\) its extremes are \(a_s,-b_s\). The other entries at a mark and all entries off \(S\) have absolute value at most \(\max(0,\lfloor c\rfloor)\). If there is no eligible model, the share is zero.

Use \(\lambda^{\rm act}_{s,a}=\lambda_{s,a}\) at even \(s\) and \(\lambda^{\rm act}_{s,a}=-\lambda_{s,k-1-a}\) at odd \(s\). Let \(V\) be the tensor product over \(\operatorname{Frac}(O)\) of the dual Weyl spaces for these weights. For a full \(W\)-lattice \(\Lambda\subset\operatorname{Frac}(W)^k\), let \(H_\Lambda\) be the integral dual Weyl tensor lattice on its Frobenius translates over \(O\). Products of the tableau bases in transported frames, with their determinant twists, form a field basis of \(V\) and span \(H_\Lambda\) over \(O\). Indeed two frames of the same \(O\)-lattice differ by \(U\in\mathrm{GL}_k(O)\). Substitution by \(U\) and \(U^{-1}\) preserves the integral column-invariant polynomials, and the determinant twist for \(U\) is a unit. The intersection over frames in each dual Weyl factor therefore equals its value in any one frame.

A rank-\(r'\) edge \(e:\Lambda'\to\Lambda\) is a pair of full \(W\)-lattices with \(p\Lambda\subseteq\Lambda'\subseteq\Lambda\) and \(\dim_{W/pW}(\Lambda/\Lambda')=r'\). Form \(H_e\) on \(\Lambda'_O+\omega\Lambda_O\) at each translate, where the subscript denotes extension to \(O\). Put
\[
 A_{r'}=\sum_{s=0}^{n'-1}\sum_{0\le a<\min(r',k)}\lambda^{\rm act}_{s,a},\qquad
 R_e=\omega^{-A_{r'}}H_e,\qquad
 R'_e=\omega^{-A_{r'}}H_e^\vee ,
\]
where \({}^\vee\) denotes the \(O\)-linear dual lattice. For every target \(\Lambda\) there must be a finite list \((e_i,F_i)_{i=1}^M\), with repetitions allowed, such that \(e_i:\Lambda_i\to\Lambda\) and
\[
 F_i\in\omega^{-t_\ast}H_{\Lambda_i}\cap\bigcap_{e:\Lambda_i\to *}R_e,\qquad
 F_i\in\omega^{c_*}R_e\quad(e:\Lambda_i\to *,\ e\ne e_i).
\]
All edges have rank \(r'\). For some \(h_0\in\N\) with \(h_0\ge(1-\kappa)\dim V\), there must be \(B\in O^{h_0\times M}\) and an \(O\)-linear map
\[
 U_0:O^{h_0}\longrightarrow \omega^{-t_\ast}H_\Lambda^\vee\cap\bigcap_{e:*\to\Lambda}R'_e
\]
such that
\[
 \sum_{i=1}^{M} B_{ji}\,(U_0\xi)(p^{A_{r'}}F_i)
 =p^{D_0}\xi_j\qquad(\xi\in O^{h_0},\ 1\le j\le h_0).
\]
\end{lemma}

\begin{proof}
Substitute \eqref{eq:p3c5-window} into the local budgets. The share is the infimum \(s\) of the passing nonnegative fractions. The empty list with \(h_0=0\) passes at \(1\), so \(s\in[0,1]\). For each model, weight and target, put \(q=(1-s)\dim V\). If \(q>0\), choose a passing \(\kappa\) with \((1-\kappa)\dim V>\lceil q\rceil-1\). Its witness has integer rank at least \(\lceil q\rceil\), so it passes at \(s\). For \(q=0\) use the empty list.
\end{proof}

\begin{definition}[Arithmetic parameter conditions]\label{def:p3c5-parameters}
Use the notation and the common bounds of Definition~\ref{def:p3c5-size}. The \emph{parameter conditions} specify positive real constants in the following inequalities. A repeated constant has the same value throughout. First require
\[
\begin{gathered}
 x^{2C_t+5}<\min(q,d),\qquad C_k m^5fB^2\le k,\qquad
 C_6m^6B\le k\log k,\qquad C_fB\le f\log k,\\
 \log\log b\le c_Lrf\log r,\qquad
 4\nu B\Lambda/k+C_\rho(\log x)\Lambda/k\le x^{-2},\qquad
 4C_{\Delta G}k^2\Delta\le\Lambda b'.
\end{gathered}
\]
For each integer \(l\ge11\) some \(j\in\N\) satisfies \(a_0'\le(l-1)^j\) and \(2j+2\le L_w\). Require \(L_w\le C_L\log x\) and \(j_0\le C_j\log x\), where \(j_0\) is the prescribed jet depth.

Write \(u_i(\tau),v_i(\tau)\) for the upper and lower interlacing rows at \(\tau\). If \(\ell(\tau)=\ell(0,s)\) for \(s\in TT\), require
\(b_0(\tau)=b_s+\lfloor b'/x\rfloor\). At every other embedding require
\[
 u_0(\tau)-v_0(\tau)\le2S,\qquad
 u_{k-1}(\tau)=v_{k-2}(\tau)-\tfrac12 .
\]
For every embedding require, over the indicated row indices,
\[
\begin{gathered}
 |b_i(\tau)-e_j(\tau)|\ge 2b/l_*+t/2+3
       \quad(0\le i,j<m),\\
 |e_i(\tau)-b_i(\tau)-H_0/2|\le2
       \quad(0\le i<m),\\
 H_0\le b_{i+1}(\tau)-b_i(\tau)\le H_0+1
       \quad(0\le i<m-1).
\end{gathered}
\]
Each consecutive difference \(v_i(\tau)-v_{i+1}(\tau)\), the sum \(v_{k-2}(\tau)+b_0(\tau)\), and each consecutive difference \(b_{i+1}(\tau)-b_i(\tau)\) must exceed \(C_{\rm gap}km^4+Nm^{10}\).

Let \(\eta_F\) be the parameter in the division data governing the size of the retained fraction. Require \(\eta_F=C_{\eta F}x^{-C_F}\) and
\[
\begin{gathered}
 \max\{k^{C_w},b^{c_w}\}\le w'',\\
 \max\{k^{C_\Delta}(C_\Delta m)^3/\eta_F^{C_\Delta},\
       C_\Delta mN,\
       k^{C_\Delta}(C_\Delta m^4 B)^4,\
       (NfB)^{C_{\rm count}}\}\le\Delta,\\
 \max\{C_r r^{19}f^5B,\ (NfB)^{C_{\rm count}},\
       C_hkm^{9/2},\ 120rm^{3/2},\ 16mN\}\le H_0,\\
 t+2\bigl(k+m(t+r+7)\bigr)^3+1\le H_0,\qquad
 fN^{C_N}\le H_0,\\
 (4f+10)r\le C_Wr\bigl(f+f^2(k-1)^3B^2/m\bigr).
\end{gathered}
\]
There is a real \(Y\ge(kb)^8\) with \(N^4M^{C_Y}\le Y^{1/4}\), \(\log Y\le C_{YL}B\), and \(Y\le g_i\le2Y\) for \(i=0,1\). Both \(g_i\) are \(1\bmod M\). For each \(g_i\) there is a \(\pi_i\in\mathcal O_E\cap\Q(\sqrt{-d})\) with \(\pi_i\bar\pi_i=g_i\). Both \(p_a^*\) split completely in \(E\), are \(1\bmod15\), and are distinct from each \(g_i\). Every host level prime is of split quadratic class.

For the prescribed conductor ideal \(\mathfrak c\) of the imaginary quadratic field require \(\log N\mathfrak c\le B\) and \(N\mathfrak c\le C_{\mathfrak c}d\). For every root of unity \(\zeta\in E\) and each \(0\le i<m\), require
\[
 \prod_\tau \tau_E(\zeta)^{e_i(\tau)}=1 .
\]
Here \(\tau_E\) is the host embedding labelled by \(\ell(\tau)\) and \(e_i(\tau)\) is integral.
\end{definition}

\begin{proposition}[The index between the two integral weight lattices]\label{B3.52b}
Let \(O\) be the coefficient DVR of a permitted model, with fraction field \(F\) and residue characteristic \(p\). For a dominant integral weight \(\lambda=(\lambda_0,\ldots,\lambda_{k-1})\), put \(\nu_i=\lambda_i-\lambda_{k-1}\) and \(l=\nu_0\). In the polynomial ring over \(O\) on a \(k\) by \(k\) matrix \(Y\), let \(A_\lambda\) consist of polynomials homogeneous of degree \(\nu_j\) in column \(j\) and invariant under adding an indeterminate multiple of column \(i\) to column \(j\) for \(i<j\). Let \(h_\lambda\) be the product of initial principal minors, one of size equal to the height of each column of \(\nu\). Let \(D_\lambda\) be the \(O\)-span of successive coefficients of elementary row substitutions on \(h_\lambda\), with the empty word included. These are the dual Weyl and Weyl lattices before the determinant twist. For \(k\ge5\), put \(d_p(\lambda)=\binom{k}{2}\lfloor\log_p(l+k-1)\rfloor\). Then
\[
 p^{d_p(\lambda)}A_\lambda\subseteq D_\lambda\subseteq A_\lambda.
 \tag{11a}\label{eq:p3c5-factor-index}
\]
The inclusions hold after common transport to a lattice frame and the determinant twist for \(\lambda\).

For the arithmetic specialization assume the parameter conditions of Definition~\ref{def:p3c5-parameters} and the embedding identification of Definition~\ref{def:p3c5-orbit-model}. At a permitted model these identify the pairs \((v,s)\) of an orbit and a position on it with the \(f\) embedding indices and match the base weights with the arithmetic level-zero weights. Let \(M_{\rm na}\) and \(M\) be the level-zero function lattices of \ref{B3.01} with fibres given by the dual Weyl and Weyl tensor lattices, respectively. If \(C_e\ge1\) and \(e_p=\lceil C_e f k^2\log_p a_\ast\rceil_+\), then
\[
 p^{e_p}M_{\rm na}\subseteq M\subseteq M_{\rm na}.
 \tag{11b}\label{eq:p3c5-total-index}
\]
\end{proposition}

\begin{proof}[Outline]
Compare ranks by tableau reduction and \ref{B3.10a} and bound denominators of the projectors that produce a scalar. Transport and twist by \ref{B3.10a}, then sum weight-spread bounds with \(R\) from \ref{B3.50} to get the lattice inclusions of \ref{B3.01}. The complete proof is in \S\ref{prf:B3.52b}.
\end{proof}

\section{The numerical comparison}

At \(v=v(t_0)\) in a permitted tied model use \(\iota(i)\) of \ref{B3.50}. Write \(\alpha_i,C_i,L_i,\delta_v\) for the column count, capacity, remaining box count and digit error of \ref{B3.02}--\ref{B3.03}, and \(g_{A,v}(j),g_{B,v}(j)\) for the gains of \ref{B3.13}--\ref{B3.14}.

\begin{theorem}[Comparison when \(z_p\ge1\)]\label{B3.53}
Assume the size bounds of Definition~\ref{def:p3c5-size}, \(K_{\delta,1}\ge5\), \(K_{\delta,2}\ge4\), \(K_{\mathrm{Pr}}\ge1\), \(r_G<k\) and \(x\ge2\). For every choice of ties, permitted pair \((u,w)\), \(t_0\in\mathcal R_p\), and \(1\le j\le m\), the condition \(z_p\ge1\) implies
\[
 g_U(t_0,j)\le g_{B,v}(j),\qquad
 g_0(t_0,j)\le g_{A,v}(j),\qquad P_v(j)\le\Pi_p(t_0,j).
 \tag{12}\label{eq:p3c5-small-residue}
\]
\end{theorem}

\begin{proof}[Outline]
Compare counts and capacities via \ref{B3.01}, \ref{B3.02} and match exclusions via \ref{B3.13}. Compare positive parts termwise using \ref{B3.14} at odd marks, subtract digit errors bounded via \ref{B3.47} and \ref{B3.03}, and use \ref{B3.50} for the price. The complete proof is in \S\ref{prf:B3.53}.
\end{proof}

\begin{theorem}[Zero counted gain when \(z_p<1\)]\label{B3.54}
Assume the parameter conditions of Definition~\ref{def:p3c5-parameters}, nonnegativity of the global constants \(K\), and \(K_{\mathrm{Pr}}\ge1\). For every coefficient-field place \(w\) with \(p=\operatorname{char}k(w)\), \(t_0\in\mathcal R_p\), and \(1\le j\le m\),
\[
 z_p<1\quad\Longrightarrow\quad
 g_U(t_0,j)-\Pi_p(t_0,j)\le0,\qquad g_0(t_0,j)\le0.
 \tag{17}\label{eq:p3c5-large-residue}
\]
\end{theorem}

\begin{proof}[Outline]
For \(z_p<1\), the ordinary summands vanish by \ref{B3.47} and the enhanced gain is at most \(N_v(t_0)\) by \ref{B3.49}. The price is at least \(N_v(t_0)\) by \ref{B3.46} and marked prices for large primes and by size bounds for small primes, so both counted gains vanish. The complete proof is in \S\ref{prf:B3.54}.
\end{proof}

\begin{definition}[Twisted chains and their comparison]\label{B3.54k}
For each word \(u\) and coefficient place \(w\) satisfying the permission conditions of \ref{B3.45}, the local datum \(Y\) specifies a filling with levels and a chain, a position map \(v\), and an \(\mathcal O\)-submodule \(V^{\rm ram}\) of the local function space. Write \(h_g\) for these permission hypotheses. A family of twisted chains for \(Y\) assigns levels on this filling and a chain on those levels to each tuple \(\mu\) at each permitted pair \((u,w)\). At the tuple \(\mu_w\) attached to \(w\), this pair equals the pair in \(Y\). Write \(Y_\mu\) for \(Y\) with the levels and chain replaced by this pair. Conversely, given levels and chains for all tuples on the fixed fillings, use the pairs at \(\mu_w\) as the local data at \(w\). The given family is then a family of twisted chains for these data.

Let \(\iota\) be the comparison datum with maps \(\iota_{u,w,h_g}\) from the arithmetic initial spaces to the local models. Call \(\iota\) a \emph{tied arithmetic comparison} for these data if the following three conditions hold. In the first condition it suffices that there exist one choice of ties.
\begin{enumerate}
\item The embedding and filling identifications of Definition~\ref{def:p3c5-orbit-model} hold with the same bijections at every permitted pair.   

      
\item At each permitted \((u,w,h_g)\), the underlying set of \(V^{\rm ram}\) is the \(K_f\)-linear span of the image under \(\iota_{u,w,h_g}\) of the arithmetic ramified submodule. Here \(K_f\) is the local coefficient field.
\item For \(1\le j\le m\) and every \(w\), the stored \(F_j(w)\) equals the value in \ref{B4.02}.
\end{enumerate}
For each tuple \(\mu\), let \(Q'(\mu)\) be the finite set of residue characteristics excluded by the division data, and fix the twisted forms \(T\). A family of twisted chains satisfies the \emph{terminal comparison} when it has the following valuation property, proved under the stage and lattice hypotheses of Lemma~\ref{B4.13a}, at each permitted pair \((u,w)\) and each used tuple \(\mu\) with \(p_w=\operatorname{char}k(w)\notin Q'(\mu)\). The used tuples are the reference tuple, the two distinguished twists, and the tuples \(\mu_{w'}\) attached to coefficient places \(w'\). For each arithmetic initial vector \(x\) and \(s\in\R\), if the lift of \(\iota_{u,w,h_g}(x)\) along \(Y_\mu\) has terminal valuation at least \(s\), the corresponding column has a rational representative with terminal lattice valuation at least \(s\). Use the reference column and lattice at the reference tuple and those of \(T\) at the other tuples, as in \ref{B4.06}.
\end{definition}

\section{Substitution into the integral theta chain}

\begin{theorem}[Global exponents in the local chain]\label{B3.55}
Fix \(D,L,X,Y,K\) with the embedding and filling identifications of Definition~\ref{def:p3c5-orbit-model}, and fix a pair \((u,w)\) permitted by \ref{B3.45}, with \(p=\operatorname{char}k(w)\). Assume \eqref{eq:p3c5-small-residue} for every choice of ties and permitted pair with \(z_p\ge1\), and \eqref{eq:p3c5-large-residue} for every coefficient-field place with \(z_p<1\). Both comparisons range over all \(t_0\in\mathcal R_p\) and \(1\le j\le m\).

Fix a complementary budget \(\mathrm{low}\in\R\). Assume there is a standard system as in \ref{B3.01} on the local model at \((u,w)\), with the prescribed filling, levels and chain and the trace-dual Witt bases of \ref{B3.08}. Use its uniformizer \(\varpi\) of \(\mathcal O\), exponent \(e_{\mathcal O}\) with \((p)=(\varpi^{e_{\mathcal O}})\), and stabilizer bounds \(d_G(j)\). Assume the system has kernel data as in \ref{B3.12} and that \(\mathcal O\to K_f\) is injective. Require the ordinary coefficient bound of \ref{B3.27} and the conditional chain inequality of \ref{B3.42}. Also require the coefficient bound of \ref{B3.12} with exponent zero at every position and every step \(0\le j<m\), as in \ref{B3.05a}. Assume there are real numbers \(F_n(v,j)\) giving the coefficient bound of \ref{B3.12} at step \(j-1\) and position \(v\) for \(v\notin v(\mathcal R_p)\) and \(1\le j\le m\), with
\[
 \mathrm{low}\le\sum_{j=1}^m\sum_{v\notin v(\mathcal R_p)}F_n(v,j)
 -\sum_{j=0}^{m-1}d_G(j).
 \tag{20}\label{eq:p3c5-complement}
\]

Let \(K^{[m]}=K_{m-1}\cdots K_0\) and let \(\Lambda_m\) be the terminal integral lattice. For \(V'\subseteq V_s\subseteq\mathcal R_p\) and \(F\) in the initial integral lattice \(M\), put \(y_{V'}F=(\prod_{t_0\in V'}y_{v(t_0)}^{\,m})F\), in the prescribed finite-set order of \ref{B3.40}. Write \(v_{\Lambda_m}(z)\ge h\) for \(z\in\varpi^{\lceil e_{\mathcal O}h\rceil}\Lambda_m\). Then
\[
 v_{\Lambda_m}\bigl(K^{[m]}y_{V'}F\bigr)
 \ge \mathrm{low}+\sum_{j=1}^m
 \left(\sum_{t_0\in V'}\max\{g_U(t_0,j)-\Pi_p(t_0,j),g_0(t_0,j),0\}
 +\sum_{t_0\in\mathcal R_p\setminus V'}[g_0(t_0,j)]_+\right).
 \tag{21}\label{eq:p3c5-local-final}
\]
\end{theorem}

\begin{proof}[Outline]
Combine bounds in \ref{B3.12} and use global gains and price or zero inputs in \ref{B3.42} according to \(z_p\). Use injectivity to preserve generation in \ref{B3.40}, sum layers with stabilizer losses, scale and round up, then split the ordinary sum as in \ref{B4.09}. The complete proof is in \S\ref{prf:B3.55}.
\end{proof}

\section{Passage to the global service condition}

\begin{lemma}[Hecke images for the chosen twisted chains]\label{B3.56}
Take a member at or above the threshold for the first closed condition of Theorem~\ref{B4.16}, at \(x>x_1\). Assume the error hypotheses of \ref{B4.03} for this member. Assume the comparisons of \ref{B3.53} at every permitted model with $\log_p(2a_\ast)\ge1$ and those of \ref{B3.54} at every coefficient place with $\log_p(2a_\ast)<1$. Choose the levels and twisted chains supplied by that condition, with its matching admissible systems and its complementary bounds at \(\sigma=1\). For every used twist \(\mu\), every embedding and filling identification for its local data as in \ref{B3.54k}, and every permitted pair \((u,w)\) with \(p_w\notin Q'(\mu)\), inequality \eqref{eq:p3c5-local-final} holds for every indicated subset of representatives and every initial integral vector. Its complementary budget is the first quantity in \eqref{eq:p4c2-budgets} at \(\sigma=1\).
\end{lemma}
\begin{proof}[Outline]
Use the first clause for the complementary budget, \ref{B3.45} for the ordinary and chain bounds, and \ref{B3.05a} for the zero bound. Apply \ref{B3.55} using the fraction field and the comparisons in both prime ranges, then split the ordinary sum as in \ref{B4.09}. The complete proof is in \S\ref{prf:B3.56}.
\end{proof}

Suppose the fillings are compatible at every finite place in the sense of \ref{B3.52} and the discarded set of \ref{B4.04} contains every inert-class prime for which \(h_p\) is nonzero of odd order. For \(p=\operatorname{char}k(w)\) and \(1\le j\le m\), Lemma~\ref{B4.25} bounds \(\gainD_j(w)\) above by the gain in \eqref{eq:p4c2-served-gain}. Theorem~\ref{B4.29} uses this inequality together with the size conditions of \ref{B4.15} and both column bounds of \ref{B4.14} at normalization \(1\), with their stated hypotheses.


\part{Divisibility along the chain and division of lattices}

% ---- p4c1 ----
\chapter{What each layer records}\label{ch:layer-records}

Fix chain data $D$ of initial rank $k$ and length $m$, with size data $L$ \cite{alpoge-2026}. Each lift increases the rank by $t_{\rm ch}+1$, where $t_{\rm ch}$ is odd. Let $x>1$ be the window point and $0<\varepsilon<1/52$ its exponent in the size datum. Put $b=\exp(x^\varepsilon)$. At a finite place $w$ of the coefficient field $K$, let $p$ be the characteristic of the residue field and $e_w=v_w(p)$ the ramification index, with $v_w$ normalized to take value one on a uniformizer. In a formula by cases, use the first
applicable case. Let $\mathcal C,\mathcal S$ be the inert and signal
classes of \ref{B3.46}. The chain identifies its embedding indices with the elements of $G=\{\pm1\}^2\times A$ whose second sign is $+1$. For $z=(\epsilon,t)\in Z=\{0,1\}\times A$, let $\tau(z)$ be the index corresponding to $(((-1)^\epsilon,1),t)$. Write
$w_j(z)=w_j(\tau(z))$ for extension weight $j-1$, where $1\le j\le m$.
Set $X'=e^{x^{\varepsilon/3}}$.

\begin{definition}[The orbit base value]\label{B4.01}
For integers $p\ge2$, $r\ge1$ and capacities $d_s\in\R_{\ge0}$ on $\Z/r\Z$,
let $\mathcal P$ consist of arrays $x_{s,i}\in\R_{\ge0}$ for positive odd
$i<r$ with $\sum_i(p^ix_{s,i}+x_{s-i,i})\le d_s$. Put
\[
 F_{p,r}(\mathbf d)=\frac p{p^2-1}\sum_s d_s
 -\sup_{x\in\mathcal P}\sum_{s,i}x_{s,i}.
\]
Use the filling compatibility of \ref{B3.52} for the chain weights.
For a prime $p$, take $h_p,n_p$ from \ref{B3.46}. Set
$P_p(\epsilon,t)=(\epsilon+\mathbf1_{p\in\mathcal S},t+h_p)$ with the
first coordinate modulo $2$, and put
\[
 F_{\rm orb}(j,p,z)=F_{p,n_p}\!\left(\bigl((k-1)w_j(P_p^sz)\bigr)_{s\in\Z/n_p\Z}\right).
\]
\end{definition}
\begin{lemma}[Attainment and change of orbit origin]\label{B4.01d}
For $p\ge2$, $r\ge1$ and nonnegative capacities, the supremum in \ref{B4.01} is attained. For the filled capacities, $F_{\rm orb}(j,p,P_p^a z)=F_{\rm orb}(j,p,z)$ for every $a\ge0$.
\end{lemma}
\begin{proof}[Outline]
Use the periods from \ref{B3.46} and the weight bounds from \ref{B3.47} to obtain nonnegative capacities and maximize flow on the compact feasible set. Changing the orbit origin cyclically permutes the data, preserving the capacity sum and maximum flow. The complete proof is in \S\ref{prf:B4.01d}.
\end{proof}

\begin{lemma}[Duality and changing the demands on an orbit]\label{B4.01a}
For fixed $p,r$ and a capacity vector $D$, write $F(D)=F_{p,r}(D)$ as in Definition~\ref{B4.01}, and put $\gamma=p/(p^2-1)$.
For integers $p\ge3$, $r\ge1$ and capacities $d_s\ge0$, the maximum flow
in \ref{B4.01} equals the minimum of $\sum_s d_sy_s$ over
$0\le y_s\le1$ with $p^iy_s+y_{s+i}\ge1$ for every positive odd $i<r$.
Every minimizer satisfies
\[
 \sum_{s:\,y_s\ge2/p}d_s\le p^{-1}\sum_s d_s.
\]
Suppose $p\ge5$ is prime, $r\ge2$ is even, $T\subseteq\Z/r\Z$,
$k\in\N$ with $k\ge3$, $w_s>0$ and $0\le e_s\le4w_s$. Set
$W_s=(k-1)w_s$ and $D_s=W_s+e_s-2w_s\mathbf1_T(s)$. Then
\[
 F(D)\ge F(W)-\frac{11}{p}\sum_s w_s.
\]
\end{lemma}
\begin{proof}[Outline]
Dualize the loads of \ref{B4.01} using attainment from \ref{B4.01d}, and use the gap identity to cap the dual at \(1\) and force incoming saturation at \(y_s\ge 2/p\). Test the dual at \(W+e\) and use flow monotonicity to bound the two losses. The complete proof is in \S\ref{prf:B4.01a}.
\end{proof}

\begin{lemma}[Weights of the letter logarithm and of a phase]\label{B4.01c}
Let $K'$ be a characteristic-zero nonarchimedean normed field, and let $p$
be an odd prime with $0<|p|<1$. Normalize the additive valuation by
$v(p)=1$. Put $\rho=1/(p^2-1)$ and $q_*=1/(p-1)$, and give $cZ^n$
weight $v(c)+\rho|n|$. For a self-map $P$ of a finite index set, the
substitution
\[
 L_i(Z)=\sum_{h\ge0}p^{-h}Z_{P^hi}^{p^{2h}}
\]
has a formal inverse $g$ with zero constant terms. Every term of $g_i$
has weight at least $\rho$.

Let $r\ge2$ be even and take $P(s)=s-2$ on $\Z/r\Z$. Let
$\beta\in K_0[[Z_s:s\in\Z/r\Z]]$, where $K_0\subseteq K'$ is a subfield
with $v(K_0^\times)\subseteq\Z$. Let $X=(Z_s)_{s\ {\rm even}}$ and
$Y=(Z_s)_{s\ {\rm odd}}$. Fix a primitive $(p^r-1)$-st root
$\zeta\in K'$. Assume that $\beta$ is invariant under
$Z_s\mapsto\zeta^{(-1)^sp^s}Z_s$, with $s$ represented in $[0,r-1]$,
and that every monomial in its support has positive degree in each
parity block and bidegree different from $(1,1)$. For a copy $X_1$ of
$X$, put
\[
 \Delta_X\beta=\beta(X+X_1,Y)-\beta(X,Y)-\beta(X_1,Y),
\]
and define $\Delta_Y\beta$ by splitting $Y$ in the same way. Suppose
$\exp(\Delta_X\beta)$ and $\exp(\Delta_Y\beta)$ have integral coefficients
and constant one after substitution by $L$ in each copy.
Every nonzero term of $\beta$ has weight at least $q_*-1$. Each term
of weight below $q_*$ has the form
\[
 cZ_s^{p^u}Z_t^{p^v},\qquad
 s,t\text{ of opposite parity},\quad u,v\in\Z_{\ge0},\quad
 (u,v)\ne(0,0),\quad s+u=t+v\pmod r .
\]
\end{lemma}
\begin{proof}[Outline]
Invert the letter logarithm by degree and use integrality and torus invariance to bound both coboundary weights below by \(q_*\). Binomial splitting gives the bound \(q_*-1\) for exceptional pairs of powers, with \(s+u\equiv t+v\pmod r\) by torus invariance. The complete proof is in \S\ref{prf:B4.01c}.
\end{proof}

\begin{proposition}[Cyclic and one-letter bounds]\label{B4.01b}
Let $p\ge5$ be prime and $N_0,N_1\in\N$ with $N_0<N_1$. On the even
cycle $\Z/r\Z$, $r\ge2$, take Young diagrams $\nu_s$ with at most $N_0$
rows. Put
\[
 \ell_{s,ab}=\sum_{h\ge0}\frac{Z_{s-2h,ab}^{p^{2h}}}{p^h}.
\]
Suppose a formal kernel over $\C_p$ has an expansion
\[
 \Theta(\xi\ell)=U\exp\!\left(\epsilon\sum_{a,b}
 \beta(\ell_{\bullet,ab})\right),
 \qquad |\xi|_p=1,\quad\epsilon=\pm1,\quad |[Z^n]U|_p\le1
\]
for every exponent $n$, with $(a,b)$ ranging over the $N_0N_1$ matrix
entries. Use the phase hypotheses of \ref{B4.01c} over $\C_p$ with
$K_0=\widehat{\Q_p^{\rm nr}}$, the completed maximal unramified extension
of $\Q_p$. Its valuation is integer-valued on nonzero elements, since
each such element has an approximation of the same valuation in a
finite unramified extension.

Let $A_s,B_s$ be the first two row lengths, with missing rows of length
zero. For $T\subseteq\Z/r\Z$ put $D_s=|\nu_s|-A_s\mathbf1_T(s)$, where
$|\nu_s|$ counts boxes. Put $c_s=B_s$ on $T$ and $c_s=A_s$ elsewhere.
All distances $e$ in the following formulas are positive odd integers.
For $s\in T$ set
\[
 h_s=\inf\{e:s+e\in T\},\qquad
 M_s=\max_{p^e\le|\nu_{s-e}|}c_{s-e},
\]
with empty values $\infty$ and zero, respectively. Put
\[
 \begin{aligned}
 H_s&=\sum_{e<h_s}(A_sp^{-e}-c_{s+e})_+
 +\sum_{e>h_s}A_sp^{-e},\\
 J_s&=\sum_{e\ne h_s}((A_s-B_s)p^{-e}-c_{s+e})_+ .
 \end{aligned}
\]
For each diagram column of height $h$, choose a strictly increasing map
$\{1,\ldots,h\}\to\{1,\ldots,N_1\}$ and form $P_I$ by \ref{B3.10}.
For a real size bound $a\ge p$ with $|\nu_s|\le a$ for all $s$ and
$2\log_p a+2\le N_0-1$, each coefficient
$[\tau^\mu]\mathcal D_{P_I}(\Theta)$ of \ref{B3.11} has valuation at least
\[
 F(D)+\sum_{s\in T}(H_s-\gamma M_s)-rN_0N_1(\log_p a+1),
\]
where $F,\gamma$ are as in \ref{B4.01a}. If the capacities $D_s$ have the
form in the comparison of \ref{B4.01a}, one may replace
$F(D)+\sum_{s\in T}H_s$ by
$F(W)-11\sum_s w_s/p+\sum_{s\in T}J_s$.

For one letter, let $\nu$ have at most $N_0$ rows. Use
$\ell_{ab}=\sum_{h\ge0} Z_{ab}^{p^{2h}}/p^h$ and assume the same form of
expansion with phase
$(\epsilon/2)\sum_{a,b}\beta(\ell_{ab},\xi'\ell_{ab})$,
where $|\xi'|_p=1$ and $\beta$ satisfies the phase hypotheses of
\ref{B4.01c} at $r=2$.
For $a\ge p$ and $|\nu|\le a$, each frame coefficient of the evaluation
\ref{B3.11} for such column choices has valuation at least
\[
 \rho|\nu|-N_0N_1(\log_p a+1),\qquad \rho=\frac1{p^2-1}.
\]
\end{proposition}
\begin{proof}[Outline]
Use \ref{B3.26} for the frame and \ref{B3.11} for row degrees, then \ref{B4.01c} and the dual of \ref{B4.01a} to raise phase weights to \(q_*\), using period two for one letter. Apply Legendre and \ref{B3.19} for digit loss and the dual to compare with \(F(W)\). The complete proof is in \S\ref{prf:B4.01b}.
\end{proof}

\begin{definition}[The two integer weight rows]\label{B4.02z}
Fix an embedding type \(\tau\). Number coordinates from zero. The interlacing coordinates are integers \(v_0,\ldots,v_{k-2}\). Write \(b_0,\ldots,b_{m-1}\) and \(e_0,\ldots,e_{m-1}\) for its two sequences of auxiliary integer weights, with \(e_{j-1}=w_j(\tau)\). These \(b_i\) are indexed weights at \(\tau\). The scale \(b\) is common to all types. Put \(t=t_{\rm ch}\) and
\[
 (s_0,\ldots,s_{k+m-2})=(v_0,\ldots,v_{k-2},-b_0,\ldots,-b_{m-1}).
\]
For \(1\le j\le m\), put \(P=k+j-1\), \(Q=m-j\). The prescribed rows at layer \(j\) are
\[
 \begin{aligned}
 \kappa^{\rm big}_{j,\tau,i}
    &=s_i+e_{j-1}-(P-1-i)-(t+1)/2 &&(0\le i<P),\\
 \kappa^{\rm small}_{j,\tau,c}
    &=e_{j-1}-b_{j+c}+(t+1)/2+c &&(0\le c<Q).
 \end{aligned}
\]
The weight compatibility condition requires equality with these rows coordinate by coordinate. Both rows are integral since \(t\) is odd. When \(j=m\) the second row is empty.
\end{definition}

\begin{lemma}[Spacing of the auxiliary weights]\label{B4.03a}
At every type the weight parameters of Definition~\ref{B4.02z} satisfy
\[
 e_i-b_i\ge(t+3)/2\quad(0\le i<m),\qquad
 b_{i+1}-e_i\ge(t+3)/2\quad(0\le i<m-1).
\]
In particular \(e_{i+1}-e_i\ge t+3\). The initial column coordinates \(u_0,\ldots,u_{k-1}\) are half-integers in \(\Z+1/2\) with \(u_i>v_i>u_{i+1}\) and \(u_{k-1}>-b_0\).
\end{lemma}
\begin{proof}
The two displayed inequalities and the strict interlacing are the defining inequalities on the parameters at an embedding. Apply the first inequality at \(i+1\) and the second at \(i\). Adding them cancels \(b_{i+1}\).
\end{proof}

\begin{lemma}[The two row weights and the column degrees]\label{B4.02a}
Write \(t=t_{\rm ch}\), \(N_l=k+l(t+1)\), and number coordinates from zero. Fix an embedding index \(\rho\). The Harish--Chandra parameters at \(\rho\) consist of half-integers \(U_0,\ldots,U_{k-1}\) for the word, integers \(v_0,\ldots,v_{k-2}\) for the common row, and integers \(b_q,e_q\) for \(0\le q<m\). Their separation conditions are
\[
 U_i>v_i>U_{i+1},\qquad U_{k-1}>-b_0,\qquad
 e_q-b_q\ge (t+3)/2,\qquad b_{q+1}-e_q\ge(t+3)/2,
\]
where a condition involving \(q+1\) is imposed for \(q+1<m\).
The entries \(v_i,b_q,e_q\) and the two endpoints of \(U\) are independent of the word.

For \(1\le j\le m\), put \(l=j-1\), \(p=k+j-1\), \(w=w_{j\rho}=e_l\), and let \(z\) be the concatenation
\[
 (v_0,\ldots,v_{k-2},-b_0,\ldots,-b_{m-1}).
\]
The two signed row tuples are
\[
 \kappa^{\rm big}_{j\rho i}=z_i+w-(p-1-i)-(t+1)/2
       \quad(0\le i<p),\qquad
 \kappa^{\rm small}_{j\rho c}=-b_{j+c}+w+(t+1)/2+c
       \quad(0\le c<m-j).
\]
The first tuple is weakly decreasing with entries at least one. The second is weakly decreasing with entries at most minus one.

The column degrees at this layer form a weakly decreasing tuple of positive integers of length \(N_l\). In order, they are
\[
 \nu_{uj\rho i}=U_i+e_l-N_l-t/2+i \quad(0\le i<k),
 \qquad
 \nu_{uj\rho,k+q(t+1)+d}=e_l-e_q-(l-q)(t+1)
 \quad(0\le q<l,\ 0\le d\le t).
\]
\end{lemma}

\begin{proof}[Outline]
Simplify the row shifts using the oddness of \(t\), then use interlacing and half-gap bounds for row order and endpoint signs. Substitute the source parameters to obtain the column head and tail, and use the same bounds for their order and positivity. The complete proof is in \S\ref{prf:B4.02a}.
\end{proof}

\begin{lemma}[Weights of a selected word]\label{B4.02b}
Take a point of the construction of \ref{B8.05} beyond its numerical cutoff. Let \(n=|A|\), \(L=x^\varepsilon\), \(X'=\exp(x^{\varepsilon/3})\), \(b=e^L\), \(a=xb\), \(b'=a/X'\), \(w''=b/(nX')\). With the exponents of \ref{B8.02ac} put
\[
 \delta=\lceil x^{c_D}\rceil,\qquad h_0=2\lceil x^{c_h}\rceil,\qquad
 S_c=4k\delta,\qquad \eta=mh_0 .
\]
At each \(\rho\) there is an integer \(B_\rho\) such that
\[
 b_q=B_\rho+qh_0,\quad e_q=B_\rho+(q+1/2)h_0,\quad
 v_i=4\delta(k-2-i).
\]
At a marked position \(\rho=(0,s)\), \(s\in TT\), one has
\(B_\rho=b_s+\lfloor b'/x\rfloor\). At a thin position \(B_\rho=\lceil w''\rceil\). At a thick position \(B_\rho=\lceil b\rceil+d_\rho\), where the integer \(d_\rho\) is zero except possibly at one designated thick position and \(0\le d_\rho<|\mu(E)|\). Here \(\mu(E)\) is the group of roots of unity of the host field.

If \(\xi_i\in\{0,1\}\) is bit \(i\) of the word, then
\(U_{i+1}=v_i-1/2-\xi_i\delta\) for \(0\le i<k-2\). Its endpoints are
\[
 (U_0,U_{k-1})=
 \begin{cases}
 (a_s+(k-1)/2,-b_s-(k-1)/2),&\rho=(0,s),\ s\in TT,\\
 (v_0+S_c+1/2,-1/2),&\text{otherwise}.
 \end{cases}
\]
In particular the two half-gaps of \ref{B4.02a} are both \(h_0/2\), and each actual tail degree in block \(q<j-1\) is \((j-1-q)(h_0-t-1)\).
\end{lemma}

\begin{proof}[Outline]
Take the common base and endpoints from \ref{B8.05}, use the bit formula for positive gaps, and use \ref{B8.02ac} and \ref{B8.59} for marked endpoint inequalities. Combine the rank and exponent bounds for \(h_0\ge t+3\) and use \ref{B4.02a} for tail degrees. The complete proof is in \S\ref{prf:B4.02b}.
\end{proof}

\begin{definition}[The base term of a layer]\label{B4.02}
Put $n=|A|$ and $X''=nxX'$. For the integer row tuples
$\kappa^{\rm big}_{j,\tau},\kappa^{\rm small}_{j,\tau}$ of lengths
$k+j-1,m-j$, respectively, use the row compatibility of Definition~\ref{B4.02z} whenever the prescribed weights are required. Set
\[
 \begin{aligned}
 F_d(j)&=-\tfrac{k-1}{2}\sum_{t\in A}\min\{w_j(0,t),w_j(1,t)\},\\
 D_j^{\rm Fock}&=\sum_\tau\left(\sum_a|\kappa^{\rm big}_{j,\tau,a}|
 +\sum_a|\kappa^{\rm small}_{j,\tau,a}|\right),
 & F_g(j)&=-D_j^{\rm Fock}/2 .
 \end{aligned}
\]
Write $F^L_j(w)$ for the base term that the size datum assigns to layer $j$ at $w$. The base-term condition requires $F^L_j(w)=F_j(w)$ for every $1\le j\le m$ and every finite place $w$, where
\[
 F_j(w)=\begin{cases}
 F_g(j),&p\in\{g_0,g_1\},\\
 F_d(j),&p=d,\\
 \displaystyle\sum_\tau(k-1)w_j(\tau)/(p^2-1),
 &p\in\mathcal C,\ h_p=0,\ p\le X'',\\
 \displaystyle n_p^{-1}\sum_{z\in Z}F_{\rm orb}(j,p,z),
 &p\in\mathcal C\cup\mathcal S,\ p\le X'',\\
 0,&\text{otherwise}.
 \end{cases}
\]
\end{definition}
\begin{lemma}[Counting the orbit contribution]\label{B4.02y}
If $T$ contains exactly one representative of each $P_p$-orbit in $Z$, then $n_p^{-1}\sum_{z\in Z}F_{\rm orb}(j,p,z)=\sum_{z\in T}F_{\rm orb}(j,p,z)$. The host prime $d$ has the ramified class corresponding to $d$.
\end{lemma}
\begin{proof}
By \ref{B3.46} each orbit has length $n_p$. Away from the signal class a return takes exactly a multiple of the order of $h_p$. In the signal class it takes an even number $2a$ of steps, with $a(2h_p)=0$. Lemma~\ref{B4.01d} makes the summand constant on each of these orbits. At $p=d$ the tests for a prime dividing $Q_S$ and for $p=q$ both fail, since $d\nmid Q_S$ and $q\ne d$. The next test, $p=d$, gives the stated class.
\end{proof}

\begin{definition}[The loss charged to a column]\label{B4.03}
Let $M_u>0$ be the stabilizer denominator of column $u$ and $u_0$ the reference word. Put $\delta_{u,w}=v_p(M_u)$, where $v_p(p)=1$.
Write $T_u$ for the level prime of $u$, $E$ for the host field, $F=E^+$
for its real subfield, and $\tau(y)$ for the image of $y\in E$ in $K$ at
embedding index $\tau$. Use $Q_S,q,d$ of \ref{B3.46} and the fixed
auxiliary primes $g_0,g_1$. The datum $X$ specifies $\pi_i$ for $i=0,1$
and $\pi_{T,u}$ for each $u$, all in $\mathcal O_E\cap\Q(\sqrt{-d})$, with
\[
 \pi_i\bar\pi_i=g_i,\qquad
 \pi_{T,u}\bar\pi_{T,u}=T_u,\qquad
 \pi_{T,u}\equiv1\pmod{Q_Sqd(g_0g_1)^2\mathcal O_E}.
\]
It injects the local level indices $a$ of $u$ into the primes $v_u(a)$
of $F$ above $T_u$. At a prime $v$ of $F$ let $F_v$ be the completion
and $\mathfrak p_v$ its prime ideal in $\mathcal O_F$. The local level
data specify a splitting $E\otimes_F F_v\simeq F_v\times F_v$ for each
$u,i$ and $v\mid g_i$. For each $u$ and $v\mid T_u$ they specify a finite
module $\mathcal M_{u,v}$ over the subring of $E\otimes_F F_v$ generated
by $\mathcal O_E\otimes1$ and $1\otimes\mathcal O_{F_v}$. The column data
supply the seen flag $\epsilon_u(w)\in\{0,1\}$ and an embedding index
$\tau_u(w,a)$ for each local level index $a$.

For real constants $c_g,c_{\rm dig},c_6$, the local loss conditions are the following seven requirements.
\begin{enumerate}
\item $c_g\ge10$, $c_{\rm dig}\ge3$ and $c_6\ge12$.
\item For every $u,i$ and $v\mid g_i$, the first component of
$\bar\pi_i\otimes1$ under the specified splitting belongs to the
maximal ideal of $\mathcal O_{F_v}$.
\item For each $u,a$ and $v=v_u(a)$, every integral lift of
$\pi_{T,u}\otimes1$ in the indicated subring annihilates
$\mathcal M_{u,v}$.
\item For each $u$ and $w\mid T_u$, $\epsilon_u(w)=1$ exactly when
$v_w(\tau(\pi_{T,u}))>0$ for every embedding index $\tau$.
\item For each $u$, $w\mid T_u$ and $a$, put $\tau=\tau_u(w,a)$.
For every $z\in\mathcal O_F$,
$z\in\mathfrak p_{v_u(a)}\Longleftrightarrow v_w(\tau(z))>0$.
\item For every word $u$ and tuple $\mu$, the twisted chain has denominator
$M_{u_0}$. This integer is divisible by the stabilizer order of every class
at every stage $0\le i\le m$ of each twisted chain.
\item Write $E^L_{u,j}(w)$ for the loss that the size datum assigns to word $u$ at layer $j$ and place $w$. Require $E^L_{u,j}(w)=E_{u,j}(w)$ for every word $u$, every $1\le j\le m$ and every finite place $w$, with $E_{u,j}(w)$ given below.
\end{enumerate}

Put $N_i=k+i(t_{\rm ch}+1)$, $f=2|A|$ and $a_*=xb$. Use the parameter $\eta>0$ from the weight intervals of Definition~\ref{B3.47}. The row scale in the orbit error term is $S_c=k(v_0-v_1)$, where $v_0-v_1$ is the common spacing of consecutive interlacing coordinates.
For $0\le i<N_{j-1}$ let $s_{u,\tau,i}$ be weight $i$ of column $u$
at $\tau$. For $i<k$ use its initial weight. For
$i=k+\ell(t_{\rm ch}+1)+r$, $0\le r\le t_{\rm ch}$, use
$-w^u_{\ell+1}(\tau)+(t_{\rm ch}-2r)/2$, where $w^u_{\ell+1}(\tau)$ is
extension weight $\ell$ of column $u$ at $\tau$. Put
\[
 \nu_{u,j,\tau,i}=s_{u,\tau,i}+w_j(\tau)
 -\frac{N_{j-1}+N_j-1-2i}{2},
 \qquad A_{u,j,\tau}=\sum_{i<N_{j-1}}\nu_{u,j,\tau,i}.
\]
For fixed $u,j$, write $A_{\epsilon,t},\nu_{\epsilon,t,i}$ for these
quantities at $\tau(\epsilon,t)$. Put
\[
\begin{aligned}
 G_{i,u,j,w}&=F_g(j)+e_w^{-1}\sum_\tau A_{u,j,\tau}v_w(\tau(\bar\pi_i))
 +c_gf(N_m^2+m^2\eta),\\
 c_{u,j,t}&=A_{0,t}+\mathbf1_{t\in TT}(\nu_{0,t,1}-\nu_{0,t,0}+\nu_{1,t,0}),\\
 D_j^*&=F_d(j)+\tfrac12\max_u\sum_{t\in A}c_{u,j,t},\\
 C_{j,p}^*&=c_{\rm dig}fN_m^2\log_p a_*+\frac{\sum_\tau w_j(\tau)}{p^2-1}.
\end{aligned}
\]
Set $Z_A=\{(\epsilon,t)\in Z:p\notin\mathcal S,\ \epsilon=0\}$ and
$Z_{BC}=Z\setminus Z_A$. Put $Z_{\rm top}=Z_{BC}\cap(\{0\}\times TT)$.
Let $R_{BC}(j,p)$ be (4.26) of \ref{B4.11} with $BC=Z_{BC}$,
$TT_{BC}=Z_{\rm top}$, $\gamma(p)=p/(p^2-1)$, $c_{e,6}=c_6$ and
$c_{\mathrm{dl},i}=K_{\delta,i}$ from \ref{B3.47}. Define
\[
 I_{j,p}^*=\max\left\{0,n_p^{-1}\sum_{z\in Z_A}F_{\rm orb}(j,p,z)\right\}
 +R_{BC}(j,p)+X_{j,p}^* .
\]
Here $X_{j,p}^*=(k-1)p^{-3}\sum_{z\in Z_{BC}}w_j(z)
+c_6fN_m^2\log_p a_*$ if $p\in\mathcal C$ and $h_p$ has odd order
greater than $1$, and is zero otherwise.

Put $\mathrm{Top}=\{\tau(0,t):t\in TT\}$.
For integral highest weights $\lambda^u_{i,\tau,r}$ put
\[
 W_{u,\tau}=\max_{i,r}
 \left|\lambda^u_{i,\tau,r}-\left\lfloor i(t_{\rm ch}+1)/2\right\rfloor\right|
\]
over $0\le i\le m$ and $0\le r<N_i$, except for $r=0$ on $\mathrm{Top}$.
Use zero for an empty maximum. Set
\[
 T_{u,w}^*=\epsilon_u(w)\bigl(2\sum_\tau W_{u,\tau}+fN_m+
 |\mathrm{Top}|(t_{\rm ch}+1)(m+1)/4\bigr).
\]
Then the loss is
\[
 E_{u,j}(w)=\begin{cases}
 T_{u,w}^*,&p=T_u,\\
 \delta_{u,w}+G_{i,u,j,w},&p=g_i,\ i=0,1,\\
 \delta_{u,w}+D_j^*,&p=d,\\
 \delta_{u,w}+C_{j,p}^*,&p\in\mathcal C,\ h_p=0,\ p\le X'',\\
 \delta_{u,w}+I_{j,p}^*,&p\in\mathcal C\cup\mathcal S,\ p\le X'',\\
 \delta_{u,w},&\text{otherwise}.
 \end{cases}
\]
\end{definition}
\begin{lemma}[The common stabilizer denominator]\label{B4.03b}
Under condition~(6) of \ref{B4.03}, every word has $M_u=M_{u_0}$, and every twisted chain has this denominator.
\end{lemma}
\begin{proof}
Condition~(6) gives the assertion for each twisted chain. The twist by the reference tuple $\mu_0$ gives the original chain of
$u$. Its denominator $M_u$ therefore equals $M_{u_0}$ by condition (6)
of \ref{B4.03}.
\end{proof}
The ordinary and enhanced column bounds of \ref{B4.14} apply at
$p\notin\{T_u,g_0,g_1\}$ for used twists $\mu$ with $p\notin Q'(\mu)$.
The enhanced bound also requires matching local systems and the terminal estimate.
For $p=g_i\ne T_u$, see \ref{B4.23} and \ref{B4.24}.

\begin{definition}[The gain retained after selection]\label{B4.04}
Use $\rho_B,\Sigma_p$ and candidacy from \ref{B3.51}. Put
$o_p=\operatorname{ord}(h_p)$ and define $B_p$ by
\[
 \begin{aligned}
 B_p\iff{}&
 o_p<x^{\varepsilon/4}\\
 &{}\lor\bigl(p\ge k^2x^{26}\ \land\
 \exists i\in\N_{>0},\ |p^i-X'|\le\rho_BX'\bigr)\\
 &{}\lor\bigl(p\text{ is a candidate}\ \land\ \Sigma_p>x^{-\varepsilon}\bigr).
 \end{aligned}
\]
For $C_r\in\R$, define a set $\mathcal N_{C_r}$ of rational primes by
\[
 \begin{aligned}
 p\in\mathcal N_{C_r}\iff{}&
 p>X''\ \lor\ p\notin\mathcal C\cup\mathcal S\\
 &{}\lor\bigl(p\in\mathcal C\ \land\
 ((o_p>1\land o_p\text{ is odd})\lor(p<X'(1+C_r\rho_B)\land B_p))\bigr).
 \end{aligned}
\]
The prime selection condition requires the set $\mathcal D$ of discarded primes to equal $\mathcal N_{C_r}\setminus\{T_u:u\text{ is a word}\}$.
At a place $w$ above $p$, the retained gain $\gainD_j(w)$ is zero if $p\in\mathcal D$. Otherwise it equals the original gain assigned by the size datum to layer $j$ at $w$.
\end{definition}
\begin{lemma}[Gain at the exceptional primes]\label{B4.04a}
If $p=T_u$ for any word, the retained gain equals the original gain. At $p=g_0,g_1$ both gains vanish. At a supported place of inert or signal class satisfying the filling identities of \ref{B3.52}, the original gain is the prime gain in that definition for $1\le j\le m$.
\end{lemma}
\begin{proof}
Membership in the discarded set requires $p\ne T_v$ for every word $v$. Thus a level prime uses the second branch of \ref{B4.04}. The size datum prescribes zero original gain at each fixed auxiliary prime, and either branch then gives zero. The last assertion is the gain identity of \ref{B3.52} on its specified support.
\end{proof}

\begin{definition}[The net budget of a column]\label{B4.05}
Set
\[
 \sg(u,w)=\sum_{j=1}^m\bigl(\gainD_j(w)+F^L_j(w)-E^L_{u,j}(w)\bigr).
\]
Let $S_L$ denote the finite support specified by the size datum. Its original gain, stored base term and stored loss vanish at every $w\notin S_L$, for every layer and word.
\end{definition}
\begin{lemma}[Value and support of the budget]\label{B4.05a}
For every $u$ and $w\notin S_L$ one has $\sg(u,w)=0$. When the base and loss identifications of \ref{B4.02} and \ref{B4.03} hold, the budget at any $w$ is $\sum_{j=1}^m(\gainD_j(w)+F_j(w)-E_{u,j}(w))$.
\end{lemma}
\begin{proof}
Outside $S_L$, either branch defining the retained gain gives zero because the original gain vanishes there. The other two summands vanish by the support conditions. For the second assertion substitute the two identifications into the finite sum of \ref{B4.05}.
\end{proof}



% ---- p4c2 ----
\chapter{From local lifts to bounds on whole columns}\label{ch:local-to-column}

Let \(p\) be the characteristic of the residue field at a finite place \(w\) of \(K\). Normalise \(v_w(K^\times)=\Z\) and put \(e_w=v_w(p)\). For a word \(u\), let \(T_u\) be its level prime and \(\delta_{u,w}=v_p(D\Gamma_u)\) its stabiliser loss. The chain of \(m\ge2\) lifts has spaces \(V_{u,i}\) with lattices \(\Lambda_{u,i}\) and complex embeddings \(\iota_{u,i}\). Superscripts \(\mu\) denote twists, and \(e^\mu_0\) identifies the initial lattices. The reference index is \(\mu_0\). The indices in use are \(\mu_0,\mu_\ast,\mu'_\ast,\mu_w\), with excluded sets \(Q'(\mu)\).

Put \(a_\ast=xb\) and \(f=2|A|\). Write \(t_{\rm ch}+1\) for the increase in rank at each lift, so that the rank at stage \(i\) is \(N_i=k+i(t_{\rm ch}+1)\). The even starting rank and odd \(t_{\rm ch}\) give even \(N_i\). Use the weights \(w_j(z)\), base terms \(F_j(w)\) and cutoff \(X''\) of \ref{B4.02}. Write \(E^{\rm ev}_{u,j}(w)\) for the explicit loss given by the case formula of \ref{B4.03}, with the constants specified there and in \ref{B3.47} \cite{alpoge-2026}. The error hypotheses of \ref{B4.03} impose its bounds on the error constants, its level and place conditions, the common stabiliser denominator, and the identity \(E_{u,j}(w)=E^{\rm ev}_{u,j}(w)\) for every \(u,w\) and \(1\le j\le m\).

\section{Valuation bounds and the input lattice}

\begin{definition}[Bounds for vectors and columns]\label{B4.06}
Write \(\Lambda_{(w)}=\mathcal O_{K,(w)}\Lambda\). For \(s\in\R\) put \(\Lambda_w[s]=\{az:z\in\Lambda_{(w)},\ a\in K,\ a=0\text{ or }v_w(a)\ge e_ws\}\). Let \(\mathcal C_{u,\mu,w}(s)\) consist of those \(x\in V_{u,0}\) whose complex column is the image of a rational vector in \(\Lambda_w[s]\), with \(\Lambda\) the terminal lattice of the column. At \(\mu_0\) the column is \(C_u(\iota_{u,0}x)\) and the lattice is \(\Lambda_{u,m}\). At \(\mu\ne\mu_0\) they are \(C^\mu_u(\iota^\mu_{u,0}e^\mu_0x)\) and \(\Lambda^\mu_{u,m}\). The embedding of \(e^\mu_0x\) is the image of \(\iota_{u,0}x\) under the initial complex twist isomorphism.
\end{definition}

At \((u,w)\), write \(h_g\) for the prime conditions in \ref{B3.45} under which the local model is prescribed. Thus \(p\ge5\) is an inert-class prime outside the host level, with \(p\nmid T_u\). The second Frobenius component has even order at every host prime above \(p\). Write \(\mathcal O,\mathcal K,M\) for its DVR, fraction field and input lattice. Use \(R_p\) from \ref{B3.46} and put \(W_p=V_p\) for the selected set of \ref{B3.51}. At the top weights \(a_L,b_L\) of \ref{B3.48}--\ref{B3.50}, set
\[
 A_{t,j}=[g_{0,t}(j)]_+,\qquad H_{t,j}=\max\{g_{U,t}(j)-\operatorname{Pr}_t(j),A_{t,j}\},
 \qquad [a]_+=\max(a,0).
\]
Let \(Y_U=\prod_{t\in U}y_{\iota(t)}^m\) for \(U\subseteq R_p\), in the model's order, with the neighbour \(y_v\) of \ref{B3.35} and the position map \(\iota\). Put \(S_t=\lceil\sum_{j=1}^m(H_{t,j}-A_{t,j})\rceil\) and \(n(B)=\sum_{t\in W_p\setminus B}S_t\).

\begin{definition}[The lattice supporting the enhanced estimate]\label{B4.07}
In a permitted model, let \(V^{\rm ram}\) be its ramified submodule and \(\tau_p\) the budget of \ref{B3.52}. Set
\[
 N_P=\operatorname{span}_{\mathcal O}
 \{Y_B(p^{n(B)}F): B\subseteq W_p,\ F\in M\}
 \cap p^{[\lceil\tau_p\rceil]_+}M\cap V^{\rm ram}.
 \tag{4.6}\label{eq:p4c2-np}
\]
For the rational ramified subspace \(V^{\rm ram}_{u,0}\), put \(S^{\rm ram}_{u,w}=(\Lambda_{u,0})_{(w)}\cap V^{\rm ram}_{u,0}\). Let \(\mathcal N_{u,w}\) consist of the rational initial vectors \(x\) for which \(c_{u,w,h_g}(x)\in N_P\) for some permitted model with initial comparison \(c_{u,w,h_g}\), together with all of \(S^{\rm ram}_{u,w}\) when \(W_p=\varnothing\).
\end{definition}

\begin{lemma}[The rational lattice below the enhanced lattice]\label{B4.07c}
At a permitted \((u,w)\) use the tied arithmetic comparison \(c=c_{u,w,h_g}\) into the space \(\mathscr P\) of component-indexed polynomial functions of lattice tuples. Put \(A=\mathcal O_{K,(w)}\), \(V=V_{u,0}\), \(L_w=(\Lambda_{u,0})_{(w)}\) and \(V^r=V^{\rm ram}_{u,0}\). The coefficient embedding \(i:K\hookrightarrow\mathcal K\) has \(i^{-1}(\mathcal O)=A\). The comparison is \(i\)-semilinear with \(\mathcal K\)-independent images of a \(K\)-basis. It gives \(M=\mathcal O c(L_w)\) and \(V^{\rm ram}=\mathcal K c(V^r)\), with the latter viewed as an \(\mathcal O\)-submodule of \(\mathscr P\).

Write \(R_w=L_w\cap V^r\) and, for \(W_p\ne\varnothing\), \(J_w=\{z\in V^r:c(z)\in N_P\}\). Then \(\mathcal O c(R_w)=M\cap V^{\rm ram}\), \(\mathcal O c(J_w)=N_P\), and for \(a=\sum_{t\in W_p}S_t\), \(b_0=[\lceil\tau_p\rceil]_+\),
\[
 p^{a+b_0}R_w\subseteq J_w\subseteq p^{b_0}R_w.
 \tag{4.6a}\label{eq:p4c2-rational-lattice}
\]
The member, \(u\) and \(w\) determine these modules.
\end{lemma}
\begin{proof}[Outline]
Transport the class polynomials to lattice tuples and use the tableau degree bound for integral scaling in \ref{B3.52b}. Retract to \(K\), commute intersections with flat base change, and compare basis coordinates to get \(J_w=C_w\). The complete proof is in \S\ref{prf:B4.07c}.
\end{proof}

\begin{lemma}[An individual edge in the neighbour sum]\label{B4.07b}
Fix integers \(q_v\ge0\) in a permitted model. For the base and intermediate tensor fibres \(H_L,H_a\) of \ref{B3.52a}, put \(A_v=\mathcal A_v\), \(R_a=\omega^{-A_v}H_a\), \(R'_a=\omega^{-A_v}H_a^\vee\), with \({}^\vee\) denoting integral duality. Define
\[
 U_L=\omega^{-q_v}H_L\cap\bigcap_{a:L\to *}R_a,\qquad
 P_L=\left(\omega^{-q_v}H_L^\vee\cap\bigcap_{a:*\to L}R'_a\right)^\vee .
\]
Then
\[
 \omega^{q_v}U_L\subseteq H_L\subseteq U_L,\qquad
 \omega^{q_v}H_L\subseteq P_L\subseteq H_L,\qquad
 p^{A_v}R_a\subseteq P_L\quad(a:*\to L).
 \tag{4.6d}\label{eq:p4c2-fibres}
\]
Let \(L_B\) be the equivariant function lattice with product fibre \(\mathscr L_B\) given by \(P\) on \(B\) and \(H\) elsewhere. For \(v\in B\), replace \(P\) at \(v\) by \(U\) to obtain \(L^+_{B,v}\) and \(\mathscr L^+_{B,v}\). Assume trivial level stabilisers and fix a full lattice tuple \(\mathbf L\), so every \(\mathbf L_v\) is a lattice. Denote the equivariant orbit function with value \(\xi\) at \(\mathbf L\) in component \(i\) by \(\delta_{i,\mathbf L}(\xi)\), and write \(\mathbf L[S]=\mathbf L[v\leftarrow S]\). Then \(y_vL^+_{B,v}\subseteq L_B\). If \(S\to L_v\), \(\xi\in\mathscr L^+_{B,v}(\mathbf L[S])\), and \(p^{A_v}\xi\in p^E\mathscr L_B(\mathbf L[L'])\) for every other edge \(S\to L'\), where \(E\in\N\), then
\[
 y_v\delta_{i,\mathbf L[S]}(\xi)\equiv
 \delta_{i,\mathbf L}(p^{A_v}\xi)\pmod{p^E L_B}.
 \tag{4.6e}\label{eq:p4c2-edge}
\]
\end{lemma}
\begin{proof}[Outline]
Use the half-step inclusions in the fibres of \ref{B3.52a} and the defining dual pairing to prove integrality. In the orbit function, all other edges vanish modulo \(p^E\), and the trivial stabilizer gives unique transport of the chosen edge. The complete proof is in \S\ref{prf:B4.07b}.
\end{proof}

\begin{proposition}[Colength of the enhanced lattice]\label{B4.07a}
In a permitted model put \(R=M\cap V^{\rm ram}\), \(W=W_p\), \(e=v_{\mathcal O}(p)\) and \(d_M=\operatorname{rk}_{\mathcal O}M\). Assume finitely many level classes with trivial stabilisers, commuting neighbours preserving \(M,M_{\rm na}\), saturation of \(R\) in \(M\), and the following inclusions for the prescribed lattice exponent \(b_p\in\N\),
\[
 p^{b_p}M_{\rm na}\subseteq M\subseteq M_{\rm na}.
 \tag{4.6f}\label{eq:p4c2-sandwich}
\]
In the arithmetic specialisation of \ref{B3.52b}, the orbit positions index the \(f\) embeddings and their \(k\) weight entries agree with the level-zero weights. Under the parameter bounds of that proposition, its comparison of Weyl and dual Weyl lattices gives these inclusions with \(b_p=[\lceil C_e f k^2\log_p a_\ast\rceil]_+\) for \(C_e\ge1\). Saturation follows from the ramified subspace comparison of \ref{B4.07c}.

Assume that the position map is injective on \(W\) and identify \(W\) with its image. At each \(v\in W\) require the presentation of \ref{B3.52a} with rank \(r\), \(t_\ast=q_v=[\lceil t_v\rceil]_+\), \(D_0=D_p\) and share \(\kappa_v\). Its off-edge test is \(F_i\in\omega^{c_\ast}R_{a'}\) at every other edge from the source of \(a_i\), with \(c_\ast\in\N\) and \(c_\ast\ge4(D_p+1)\) when \(W\ne\varnothing\). The test at \(c_\ast=8(D_p+1)\) suffices. The tied models use \(r=r_G\) for \(p<p_0\) and \(r=1\) otherwise. If \(m\ge1\) and \(m\sum_{v\in W}\kappa_v\le1/2\), then
\[
 N_P\subseteq p^{\tau_p}R,\qquad
 \operatorname{length}_{\mathcal O}(R/N_P)\le e d_M(\tau_p+\mathrm{io}_p)
\]
for the budgets of \ref{B3.52}.
\end{proposition}
\begin{proof}[Outline]
Use Smith bases, \ref{B4.07b}, \ref{B3.52a}, and Nakayama's lemma to bound the number of exceptional tensors. Descending induction bounds the tails of the Smith exponents, and summing these bounds with saturation and the budgets of \ref{B3.52} gives the result. The complete proof is in \S\ref{prf:B4.07a}.
\end{proof}

\begin{definition}[The two column budgets]\label{B4.08}
For \(\sigma\in\R\), denote the base and enhanced column budgets by \(B^0_{u,w}(\sigma)\) and \(B^+_{u,w}(\sigma)\). With \(g^{\rm s}_{p,j}\) the gain used in the column bound, put
\[
 B^0_{u,w}(\sigma)=\sum_{j=1}^m(F_j(w)-E^{\rm ev}_{u,j}(w))+\sigma\delta_{u,w},
 \qquad
 B^+_{u,w}(\sigma)=B^0_{u,w}(\sigma)+\sum_{j=1}^m g^{\rm s}_{p,j}.
 \tag{4.7}\label{eq:p4c2-budgets}
\]
Let \(\mathcal E\) be the set of primes dividing \(N_{\rm cond}g_0g_1\), where \(N_{\rm cond}\) is the host conductor. Let \(\mathcal O_{\rm odd}\) be the inert-class primes at which \(h_P\ne0\) has odd additive order. Denote the counted gain of \ref{B3.51} by \(g^{\rm c}_{p,j}\). Then
\[
 g^{\rm s}_{p,j}=
 \begin{cases}
 g^{\rm c}_{p,j},&p\text{ inert-class or signal},\ p\notin\mathcal E\cup\mathcal O_{\rm odd},\\
 0,&\text{otherwise}.
 \end{cases}
 \tag{4.8}\label{eq:p4c2-served-gain}
\]
The counted gain of \ref{B3.51} vanishes for \(p>X''\) or inert class with \(h_P=0\). Elsewhere it is
\[
 g^{\rm c}_{p,j}=\sum_{t\in W_p}H_{t,j}+\sum_{t\in R_p\setminus W_p}A_{t,j}.
 \tag{4.9}\label{eq:p4c2-counted-gain}
\]
\end{definition}

\section{Hecke images and their integral span}

Use the local data \(Y\) and the common identification of orbit positions with embeddings in \ref{B3.54k}(1), with the filling, edge ranks, marked positions and weights specified there. Let \(W\) assign levels and a chain on this filling to every tuple \(\mu\) at every permitted \((u,w)\). Its levels and chain at \(\mu_w\) must agree with those of \(Y\). Write \(Y_\mu\) for \(Y\) with the levels and chain assigned by \(W\) at \(\mu\).
With \(K_\mu^{[m]}\) the local composite and \(\varpi\) a uniformiser of \(\mathcal O\), let \(\mathcal L_\mu(s)\) consist of the local inputs \(F\) satisfying
\[
 K_\mu^{[m]}F\in\varpi^{\lceil v_{\mathcal O}(p)s\rceil}\Lambda^{\rm loc}_{\mu,m}.
 \tag{4.10}\label{eq:p4c2-vkge}
\]

\begin{theorem}[The whole-chain estimate on Hecke images]\label{B4.09}
Fix \(Y,W,\sigma\) and assume the error hypotheses of \ref{B4.03}. For each used \(\mu\), identification for \(Y_\mu\) as in \ref{B3.54k}(1), and permitted \((u,w,h_g)\) with \(p\notin Q'(\mu)\), assume a matching admissible system and kernel of \ref{B3.45} on the prescribed model, filling, levels and chain. Write \(d_G(i)\) for its stabiliser losses and \(\mathcal V\) for the local positions, with representative injection \(\iota:R_p\to\mathcal V\). Require \(f_{v,j}\in\R\) such that the differentiated coefficients at step \(j-1\) and position \(v\) satisfy the lower bound \(f_{v,j}\) of \ref{B3.12} for \(v\in\mathcal V\setminus\iota(R_p)\) and \(1\le j\le m\), and
\[
 B^0_{u,w}(\sigma)\le
 \sum_{j=1}^m\sum_{v\notin\iota(R_p)}f_{v,j}-\sum_{i=0}^{m-1}d_G(i).
 \tag{4.11}\label{eq:p4c2-complement}
\]
For every tie for \(Y\) and permitted place with \(\log_p(2a_\ast)\ge1\), assume the comparisons of \ref{B3.53} for \(t\in R_p\), \(1\le j\le m\),
\[
 g_{U,t}(j)\le g^{\rm loc}_{B,\iota(t)}(j),\qquad
 g_{0,t}(j)\le g^{\rm loc}_{A,\iota(t)}(j),\qquad
 P^{\rm loc}_{\iota(t)}(j)\le\operatorname{Pr}_t(j).
 \tag{4.12}\label{eq:p4c2-small-prime}
\]
At every coefficient-field place with \(\log_p(2a_\ast)<1\), assume the comparisons of \ref{B3.54} for the same \(t,j\),
\[
 g_{U,t}(j)-\operatorname{Pr}_t(j)\le0,\qquad g_{0,t}(j)\le0.
 \tag{4.13}\label{eq:p4c2-large-prime}
\]
Then for each of the indicated twists, ties and permitted places outside \(Q'(\mu)\), every \(V'\subseteq V\subseteq R_p\) and \(F\in M\) satisfy
\[
 Y_{V'}F\in\mathcal L_\mu\!\left(B^0_{u,w}(\sigma)+
 \sum_{j=1}^m\left(\sum_{t\in V'}H_{t,j}+\sum_{t\in V\setminus V'}A_{t,j}
 +\sum_{t\in R_p\setminus V}A_{t,j}\right)\right).
 \tag{4.14}\label{eq:p4c2-hecke}
\]
\end{theorem}
\begin{proof}[Outline]
The matching system supplies \ref{B3.45} and \ref{B3.42}, while item 7 of \ref{B3.05a} gives the exponent-zero bound of \ref{B3.12} everywhere. Carrying the tie and comparisons to the twist, apply \ref{B3.55} with lower exponent \(B^0_{u,w}(\sigma)\) and split at \(V\). The complete proof is in \S\ref{prf:B4.09}.
\end{proof}

\begin{lemma}[Extension from generators to \(N_P\)]\label{B4.10}
Under \ref{B4.09}, for each used \(\mu\), identification for \(Y_\mu\) as in \ref{B3.54k}(1), and permitted \((u,w,h_g)\) with \(p\notin Q'(\mu)\), every \(F\in N_P\) satisfies
\[
 F\in\mathcal L_\mu\!\left(B^0_{u,w}(\sigma)+\sum_{j=1}^m g^{\rm c}_{p,j}\right).
 \tag{4.17}\label{eq:p4c2-local-np}
\]
The lattice \(N_P\) is common to these twists.
\end{lemma}
\begin{proof}[Outline]
Apply the Hecke-image bound with \(V=W_p,V'=B\), multiply by \(p^{n(B)}\), and use \(\sum_j H_{t,j}\le\sum_j A_{t,j}+S_t\) on omitted representatives to dominate \(\sum_j g^{\rm c}_{p,j}\). Extend to \(N_P\) by \(\mathcal{O}\)-linearity, with common defining data across twists. The complete proof is in \S\ref{prf:B4.10}.
\end{proof}

\begin{lemma}[Testing a local form lattice at its classes]\label{B4.10b}
Fix a word \(u\), a used tuple \(\mu\), and \(0\leq i<m\). Let \(w\) have residue characteristic \(p\) outside \(\{T_u,g_0,g_1\}\cup Q'(\mu)\). Put \(A=\mathcal O_{K,(w)}\), and fix an embedding \(\iota_w:K\hookrightarrow\C_p\) inducing \(w\).

At a class representative, express the value of a form in the local lattice frame, with its determinant twist. Remove the determinant factor to obtain its polynomial part. In the source frame include the full-height columns prescribed by the source shift. A reading of this value is its normalized Weyl pairing with a product of flag minors of the resulting diagram, one minor for each column, as in \ref{B3.10a}. Use the target diagram and target determinant twist for the readings at stage \(i+1\). A rational class value means the evaluation of a weight polynomial over \(K\), with the prescribed determinant twist.

At either stage, a complex class function belongs to the image of the localized form lattice if and only if it belongs to the complex space of forms at that stage and each of its class values has a rational representative all of whose readings have absolute value at most one. For a rational target representative every reading has the absolute value of an element of \(\iota_w(K)\).
\end{lemma}
\begin{proof}[Outline]
Choose rational frames by approximation and use \ref{B3.10a} to test integrality on twisted flag-minor products. Project and polarize to recover readings in the coefficient field, then clear stabilizer denominators by a scalar that is a unit at \(w\). The complete proof is in \S\ref{prf:B4.10b}.
\end{proof}

\begin{lemma}[From kernel entries to a lift]\label{B4.10a}
Fix \(u\) and a used twist \(\mu\). A place \(w\) is level-free when \(p\notin\{T_u,g_0,g_1\}\) and \(p\notin Q'(\mu)\). Fix such a place. At \(0\le i<m\) let \(\nu_z\) be the entry diagram at \(z\), with at most \(N_i\) rows, and \(\Theta\) the kernel series for a pair of classes at \(i,i+1\). A column choice \(I\) selects \(h\) increasing target indices from \(\{1,\ldots,N_{i+1}\}\) for each diagram column of height \(h\). Let \(G_{\nu,I}\) be the product of minors of \(X_zZ_z\) on the first \(h\) rows and the chosen columns, as in \ref{B3.10}. An entry bound \(B\) means that the absolute values of all coefficients of \(P_I=\left.G_{\nu,I}(X,\partial_Z)\Theta\right|_{Z=0}\) are at most \(p^{-B}\), for every pair and \(I\).

Suppose the diagrams and series agree with the kernel datum, the class-pair scalars have norm one, the pair maps preserve evaluated rational spaces, and the lift carries the complex space of forms at stage \(i\) into that at stage \(i+1\). Use the local class readings of \ref{B4.10b}. Assume the denominator condition of \ref{B4.03}, and use \(\sigma=0\), or \(\sigma=1\) with trivial initial stabilisers. An entry bound \(B\in\R\) gives a rational lift of every integral input with target lattice bound
\[
 B-\bigl(1-\sigma\mathbf1_{\{i=0\}}\bigr)\delta_{u,w}.
 \tag{4.20}\label{eq:p4c2-entry-lift}
\]
At \(\mu_0\) the same bound holds for the reference lift.
\end{lemma}
\begin{proof}[Outline]
Cauchy--Binet and straightening bound coefficients, and \ref{B4.10b} bounds pair terms and rounds exponents. Clear denominators by \ref{B4.03}, apply the target test of \ref{B4.10b}, and transport by stage isomorphisms, using trivial stabilizers at \(\sigma=1\). The complete proof is in \S\ref{prf:B4.10a}.
\end{proof}

\section{Orbit bounds and ordinary places}

Let \(P\) be the shift of period \(n_P\) from \ref{B4.01}. The entry datum \(\mathscr E\) gives diagrams \(\nu_{u,i,z}\) with at most \(N_i\) rows, prescribed by \ref{B4.03} for \(i<m\), and a series \(\Theta\) over \(\C_p\) for each \(u,\mu,w,i\) and consecutive class pair. The letter data at \(w\) consist of a scalar \(\xi_w\) and, for each position \(z\), a scalar \(\epsilon_{w,z}\) and a series \(\beta_{w,z}\) in letters indexed by \(\Z/n_P\Z\). A \emph{legal representative set} at \(p\) is a subset \(R\subseteq\{0,1\}\times A\) such that for each position \(z\) there is exactly one \(r\in R\) with \(P^e r=z\) for some integer \(e\ge0\). Thus it contains one point of each orbit of the shift in Definition~\ref{B4.01}. For every such \(R\), admissibility requires \(|\xi_w|_p=1\), signs \(\epsilon_{w,r}=\pm1\) for \(r\in R\), and the four letter conditions of \ref{B4.01c} for each \(\beta_{w,r}\). Its coefficients lie in \(\widehat{\Q_p^{\rm nr}}\), and it is invariant under \(Z_s\mapsto a^{(-1)^sp^s}Z_s\) for \(a^{p^{n_P}-1}=1\). Every monomial with nonzero coefficient has positive degree in both parity blocks and bidegree different from \((1,1)\). The exponentials of the two additive coboundaries have integral coefficients and constant term one after letter-logarithm substitution, as in \ref{B4.01c}.

Use \(\ell_{z,ab}(Z)=\sum_{h\ge0}p^{-h}Z_{P^{-2h}z,ab}^{p^{2h}}\), with \(a,b\) indexing the \(N_iN_{i+1}\) matrix entries. An even cyclic expansion has the form of \ref{B4.01b}, with \(\Theta(\xi_w\ell)\) equal to an integral multiplier times the exponential of the sum of the phases from \(R\). At inert class with \(h_P=0\), each kernel instead has a common unit \(\xi_0\) scaling \(\ell\) on the left, and its own units \(\xi_z\), signs \(\epsilon_z\) and two-letter phases \(\beta_z\) satisfying the coefficient-field, root-of-unity invariance, support and coboundary conditions of Lemma~\ref{B4.01c} at period two. The exponent on the right is \(\sum_z(\epsilon_z/2)\sum_{a,b}\beta_z(\ell_{z,ab},\xi_z\ell_{z,ab})\). For inert class of nonzero odd order, the required expansion is
\[
 \Theta(\xi\ell(Z))=U(Z)\exp\left(
 \sum_{r\in R}\frac{\epsilon_r}{2}\sum_{a,b}
 \beta_r((\ell_{P^sr,ab}(Z))_{s\in\Z/(2n_P)})\right),
 \tag{4.20a}\label{eq:p4c2-odd-product}
\]
with \(U\) integral, \(|\xi|_p=1\), signs \(\epsilon_r\), and a finite set \(R\) of positions. For each kernel the phases satisfy the conditions of \ref{B4.01c} at period \(2n_P\), over one field \(K_0\subseteq\C_p\) with integer-valued \(v_p\) and for one primitive \((p^{2n_P}-1)\)-st root.

\begin{definition}[Admissible orbit data]\label{def:p4c2-orbit}
Call the orbit data admissible if the following five requirements hold and, under the common denominator condition, \eqref{eq:p4c2-entry-lift} holds for every used \(\mu\), level-free \(w\), \(0\le i<m\) and \(B\in\R\). 
\begin{enumerate}
\item At every signal prime and every even-order inert prime require \(p\ge5\) and even \(n_P\ge2\). Every inert-class prime is at least five and \(g_0,g_1\ne d\).
\item At every level-free place on the signal and even-order inert branches, each legal \(R\) has admissible letter data. For every class pair at \(i<m\) of even rank product the cyclic expansion uses the same \(\xi_w,\epsilon_{w,r},\beta_{w,r}\), for each \(R\). At level-free places on the zero-shift and nonzero odd-order inert branches, require the respective expansions above for every such class pair.
\item Require the three estimates of \ref{B4.01b} for all admissible data with its size and rank hypotheses, in the following form with constants \(C_6,C\). Multiply each dimension and logarithm error by \(C\). In the refined cyclic bound use \(C_6p^{-1}\sum_s w_s\) for the weight loss, and keep \(\gamma(p)=p/(p^2-1)\) for the reverse-partner loss. Also require, at every inert-class \(w\) with \(5\le p\le X''\), \(h_P\ne0\) and even \(n_P\ge2\), the bound \(A_{t,i+1}\) for every \(u,i<m,t\in R_p\) and every local kernel with the expansion at \((0,t)\) for the specified letter data there. Proposition~\ref{B8.56} supplies this last hypothesis for the selected arithmetic family.
\item At every level-free place with \(p\ne d\), require entry exponent zero for even rank product. At level-free places of discriminant class \(p=d\), require \eqref{eq:p8c5-ramified} for every class pair at \(i<m\).
\item For every \(u,i<m,z\), put \(j=i+1\), \(w_z=w_j(z)\), \(\nu_z=\nu_{u,i,z}\), and number the rows from one. Set \(\operatorname{cap}_z=\nu_{z,2}\) on the top and \(\nu_{z,1}\) elsewhere. Require
\[
 \begin{gathered}
 w_z>0,\quad |\nu_z|\le a_\ast^2,\quad
 2\log_5(a_\ast^2)+2\le k-1,\quad X''\le a_\ast,\quad a_\ast>1,\\
 0\le|\nu_z|-(k-1)w_z\le4w_z\quad\text{off the top},\\
 0\le|\nu_z|-\nu_{z,1}-(k-3)w_z\le4w_z,\quad
 \nu_{z,1}\ge2w_z\quad\text{on the top},\\
 \operatorname{cap}_z\le b+\eta+1+3S_0,\qquad
 \operatorname{cap}_z\le B_j(z)+c_G,\qquad
 \bar\alpha_t\le\nu_{(0,t),1}-\nu_{(0,t),2}\quad(t\in TT).
 \end{gathered}
 \tag{4.20b}\label{eq:p4c2-shape}
\]
Here \(B_j(z)=\operatorname{fill}_j(z)-\kappa_j\), \(\kappa_j=(k+(2j-1)(t_{\rm ch}+1))/2\), and the filling, \(c_G,\bar\alpha_t\) are as in \ref{B3.47}--\ref{B3.52}. Here \(\eta\) is the weight-window width and \(S_0\) the weight-spread parameter of \ref{B4.03}. The constants satisfy
\[
 C\ge0,\quad C_6\le c_{e,6},\quad c_{\mathrm{dl},1},c_{\mathrm{dl},2}\ge0,\quad
 3C\le\min(1+c_{\mathrm{dl},1},c_{\rm dig}).
 \tag{4.20c}\label{eq:p4c2-constants}
\]
\end{enumerate}
\end{definition}
The two restriction laws, the two addition laws and the four remaining orbit estimates are stated in \ref{B4.11a}(1)--(5). The one-letter kernel bound in (4) also uses item 3 above.

\begin{lemma}[Four consequences of the orbit estimates]\label{B4.11}
Assume the admissibility conditions of Definition~\ref{def:p4c2-orbit} and the conclusions of \ref{B4.11a}. Let \(g_{0,S}(p,j)=\sum_{t\in R_p}A_{t,j}\) except at \(p\in\mathcal O_{\rm odd}\), where it is zero. On the signal and even-order inert branches the representatives are respectively \(\{0\}\times R_p\) and \(\{0,1\}\times R_p\). Put \(g_{0,r,j}=A_{t,j}\) for \(r=(0,t)\) and \(g_{0,r,j}=0\) for \(r=(1,t)\). In (i)--(iii) assume the base-term identity of \ref{B4.02} and the error hypotheses of \ref{B4.03}, and take a used \(\mu\) with \(p\notin Q'(\mu)\) and \(x\in S^{\rm ram}_{u,w}\).
\begin{enumerate}
\item[(i)] For \(p\ne T_u\), \(p\notin\mathcal E\), \(p\le X''\), and \(p\) signal or inert-class with \(h_P\ne0\), one has
\[
 x\in\mathcal C_{u,\mu,w}\!\left(B^0_{u,w}(\sigma)+\sum_jg_{0,S}(p,j)\right).
 \tag{4.21}\label{eq:p4c2-ordinary-gain}
\]
\item[(ii)] For \(p\ne T_u\), \(p\notin\mathcal E\), \(p\le X''\), and \(p\) inert-class with \(h_P=0\), the column exponent is at least
\[
 \sum_j(F_j(w)-C^*_{j,p}-\delta_{u,w})+\sigma\delta_{u,w}.
 \tag{4.22}\label{eq:p4c2-zero-shift}
\]
\item[(iii)] For \(p\ne T_u\) in the discriminant class, the column exponent is at least
\[
 \sum_j(F_j(w)-D^*_j-\delta_{u,w})+\sigma\delta_{u,w}.
 \tag{4.23}\label{eq:p4c2-disc-column}
\]
\item[(iv)] Under the denominator condition, at \(p\notin\{T_u,d,g_0,g_1\}\) each integral input at \(0\le i<m\) has a lift bound of
\[
 -\bigl(1-\sigma\mathbf1_{\{i=0\}}\bigr)\delta_{u,w}.
 \tag{4.24}\label{eq:p4c2-ordinary-step}
\]
This holds for the reference chain and each used nonreference twist outside \(Q'(\mu)\).
\end{enumerate}
Use the errors of \ref{B4.03}. Write \(BC\) for kind-\(BC\) positions, \(TT_{BC}\) for their top positions, and \(F_{\mathrm{Orb},j}(p,z)\) for the orbit base of \ref{B4.01}--\ref{B4.02}. Let \(\mathrm{err}_{r,j}\) be the following error restricted to the orbit of \(r\):
\[
 \begin{aligned}
 &\gamma(p)\#TT_{BC}(b+\eta+1+3S_0)+\frac{c_{e,6}}p\sum_{z\in BC}w_j(z)\\
 &\quad+\#BC\,N_{j-1}N_j(\log_p(4n_PN_ma_\ast)+1)
 +c_{\mathrm{dl},1}\#BC\,N_m^2\log_p(2a_\ast)
 +c_{\mathrm{dl},2}\#TT_{BC}(\log_p(2a_\ast))^2.
 \end{aligned}
 \tag{4.26}\label{eq:p4c2-orbit-error}
\]
\end{lemma}
\begin{proof}[Outline]
Use the kind-\(A\) bound and \ref{B4.11a}(1)--(2) at signal and even-order inert primes, and (5), (4), (3) at odd-order inert, zero-shift inert, and discriminant primes with errors from \ref{B4.03}. Lift by \eqref{eq:p4c2-entry-lift} and sum by \ref{B4.12}. The complete proof is in \S\ref{prf:B4.11}.
\end{proof}

\begin{proposition}[The orbit entry estimates]\label{B4.11a}\label{B4.16a}
Assume the cyclic and one-letter estimates with constants \(C_6,C\) in Definition~\ref{def:p4c2-orbit}(3), the conditions \eqref{eq:p4c2-shape}--\eqref{eq:p4c2-constants}, and the discriminant expansion \eqref{eq:p8c5-ramified} at level-free places in the discriminant class. In (2)--(5), \(u\) is any word, \(0\le i<m\), \(j=i+1\), and \(\rho(p)=1/(p^2-1)\). At inert primes let \(F_{\rm inert}\) be the base term of \ref{B4.02} at \(j,p\).
\begin{enumerate}
\item For every \(i\in\N\) and prime \(p\ge5\) with even \(n_P\ge2\), the cyclic expansion restricts to each orbit for every representative system and admissible letter datum. For \(p\ge5\) and \(n_P=1\) the one-position expansion restricts to each position. For \(i<m\) uniform local exponent bounds add to a full kernel bound in each case.
\item For \(5\le p\le X''\) with even \(n_P\ge2\), of signal class or inert class with \(h_P\ne0\), each kind-\(BC\) representative \(r\) has exponent at least \(g_{0,r,j}+F_{\mathrm{Orb},j}(p,r)-\mathrm{err}_{r,j}\). This holds for every unit \(\xi\), sign, phase satisfying the four letter conditions and kernel with the corresponding local expansion.
\item At a level-free place above \(d\), every class pair has entry exponent at least \(-\frac12\sum_t c_{u,j,t}\), using the contraction counts of \ref{B4.03}.
\item At inert-class primes \(5\le p\le X''\) with \(h_P=0\), one has \(F_{\rm inert}-C^*_{j,p}\le B_C\), where
\[
 B_C=\rho(p)\sum_z|\nu_{u,i,z}|-CfN_iN_{i+1}(\log_p(a_\ast^2)+1).
\]
Every kernel with the one-position expansion has exponent at least \(B_C\).
\item At nonzero odd-order inert primes \(5\le p\le X''\), each kernel with \eqref{eq:p4c2-odd-product} has exponent at least
\[
 L_{\rm odd}=\frac1{n_P}\sum_{z\in BC}F_{\mathrm{Orb},j}(p,z)
 -\left(\frac{k-1}{p^3}+\frac8p\right)\sum_{z\in BC}w_j(z)
 -3fN_m^2\log_pa_\ast.
\]
The large-error condition of \ref{B4.03} gives \(F_{\rm inert}-I^*_{j,p}\le L_{\rm odd}\).
\end{enumerate}
\end{proposition}
\begin{proof}[Outline]
Factor evaluations by \ref{B3.11} and apply cyclic bounds with \ref{B3.48} and \ref{B4.01b}, or the one-letter bound at zero shift. Use \ref{B4.03} and the integral residual series at the discriminant, and at odd order raise weights from \ref{B4.01a} and apply \ref{B4.01c}. The complete proof is in \S\ref{prf:B4.16a}.
\end{proof}

\begin{lemma}[Adding the bounds at the remaining primes]\label{B4.12}
In either chain let \(K^{[r]}\) be the first \(r\) lifts and \(V_i,\Lambda_i,\iota_i\) its rational data. If \(n\le m\) and step \(i<n\) has bound \(b_i\in\R\) on every vector of \(\Lambda_i\), then for \(x\in(\Lambda_0)_{(w)}\) and \(0\le r\le n\) there is \(y_r\in V_r\) with
\[
 \iota_r(y_r)=K^{[r]}(\iota_0x),\qquad
 y_r\in(\Lambda_r)_w[\textstyle\sum_{i<r}b_i].
 \tag{4.28}\label{eq:p4c2-chain}
\]
In particular for the primes and twists of \ref{B4.11}(iv), and \(x\in S^{\rm ram}_{u,w}\), one has
\[
 x\in\mathcal C_{u,\mu,w}(-(m-\sigma)\delta_{u,w}).
 \tag{4.29}\label{eq:p4c2-other-column}
\]
\end{lemma}
\begin{proof}
The vectors of a given bound form an \(\mathcal O_{K,(w)}\)-module by factoring the scalar of least valuation in a sum. Thus step bounds extend to localised inputs. Start at \(y_0=x\). Given \(y_r=az\), apply the next bound to \(z\) and multiply its representative by \(a\). Linearity gives the correct embedded output and the scalar valuations add. Start the twisted column at \(e^\mu_0x\) and sum \eqref{eq:p4c2-ordinary-step} to obtain \eqref{eq:p4c2-other-column}.
\end{proof}

\begin{lemma}[The base column bound at an ordinary place]\label{B4.13}
Assume the base and error identities and \ref{B4.11}--\ref{B4.12}. At \(p\notin\{T_u,g_0,g_1\}\), for any used \(\mu\) with \(p\notin Q'(\mu)\) and \(x\in S^{\rm ram}_{u,w}\),
\[
 x\in\mathcal C_{u,\mu,w}(B^0_{u,w}(\sigma)).
 \tag{4.30}\label{eq:p4c2-base-column}
\]
\end{lemma}
\begin{proof}
At the discriminant prime the error is \(\delta_{u,w}+D^*_j\), so use \eqref{eq:p4c2-disc-column}. Inert-class and signal primes are coprime to the conductor. For these primes up to \(X''\), apply \eqref{eq:p4c2-zero-shift} when the inert shift vanishes and \eqref{eq:p4c2-ordinary-gain} otherwise. At remaining places, \(p\ne d\), and \ref{B4.02}--\ref{B4.03} give \(F_j(w)=0\), \(E^{\rm ev}_{u,j}(w)=\delta_{u,w}\). Apply \eqref{eq:p4c2-other-column}.
\end{proof}

\begin{definition}[Arithmetic chart conditions at positive levels]\label{def:p4c2-charts}
Fix a word, a twist and a permitted place \(w\) over \(p\). Put \(A_w=\mathcal O_{K,(w)}\) and write \(\jmath:K\hookrightarrow\mathcal K\) for the coefficient embedding. At each stage \(1\le j\le m\), let \(V_j^{\rm rat}\) be the twisted rational form space, \(\Lambda^{\rm rat}_{j,(w)}\) its localised lattice and \(L_j^{\rm loc}\) the lattice of integral local forms.

In each chart require the following comparison data. The Hermitian form is nondegenerate and has identity Gram matrix in the chosen frame. The local splittings above \(p\) exchange their two factors under conjugation. The component representatives are unitary and have identity component at the specified places above \(p\). Every finite adelic unitary point has a presentation \(\gamma y_a k\) with \(\gamma\) rational unitary, \(y_a\) a component representative and \(k\) unitary and in the prescribed level away from \(p\). The exceptional places are away from \(p\). A chart stabiliser acts through a rational unitary matrix whose conjugate by \(y_a\) belongs to the level away from \(p\). In the chosen frame its local matrix gives the prescribed action on lattice tuples.

The position labels identify the weights and the embeddings on each block. At position \(s\) the embedding is the \(s\)-th Frobenius translate of the local embedding, with conjugation on the host-field argument for odd \(s\). The local embeddings restrict on the real host field to the completion maps, and the local field identifications carry the specified valuation rings onto the completion integer rings. A local unitary matrix whose matrix and inverse are integral in the chosen frame belongs to the local level. The level character and the twisting character at class representatives have values in the complex image of \(K\). For corresponding stabiliser and level elements, the local stabiliser character times the image under \(\jmath\) of the level character is one.

Require \(w\) to belong to the domain of the chart's integral lattice test. Thus \(x\in\Lambda^{\rm rat}_{j,(w)}\) exactly when, at every class and in every host-field basis of its dual lattice at \(p\), the class value is represented by a weight polynomial in the prescribed \(A_w\)-lattice, with the frame and determinant twist of \ref{B4.10b}. The scalar normalising this twist is the complex image of an \(A_w\)-unit. For each full tuple representing the class, its compared value belongs to the local Weyl fibre exactly when the normalised class value passes the integral test.

Finally, in each component a finite set of full lattice tuples meets every orbit under the level group, and
\[
 \dim_K V_j^{\rm rat}=\operatorname{rk}_{\mathcal O}L_j^{\rm loc}
 \qquad(1\le j\le m).
\]
\end{definition}

\begin{lemma}[Descent of a bound on the terminal vector]\label{B4.13a}
Let \(c\) be the input comparison, \(T\) the twisted forms and \(L_C(\mu,u,w,h_g)\) levels and chains on the filling. For each used \(\mu\) and permitted \((u,w,h_g)\) with \(p\notin Q'(\mu)\), suppose that at each stage \(j\) there is a map \(c_j\) from the twisted rational space to the local form space, semilinear for the coefficient embedding. Require \(c_0(e^\mu_0x)=c(x)\) and \(c_{j+1}(y)=K_j^{\rm loc}c_j(z)\) whenever \(y\) represents the complex lift of \(z\) at \(j<m\). For every rational terminal vector \(y\) and real \(s\), assume that a bound \(s\) on \(c_m(y)\) in the local terminal lattice gives the same bound on \(y\) in the twisted rational lattice. Each twisted lift on a rational input has a rational representative. Assume the reference identification of the initial isomorphism at \(\mu_0\).  For those indices and all \(x\in V_{u,0}\) and \(s\in\R\),
\[
 c_{u,w,h_g}(x)\in\mathcal L_\mu(s)\ \Longrightarrow
 x\in\mathcal C_{u,\mu,w}(s).
 \tag{4.31}\label{eq:p4c2-descend-terminal}
\]
For the arithmetic chart comparisons, assume the positive-level chart conditions of Definition~\ref{def:p4c2-charts}, rationality of the polynomial representations in the identified embedding blocks, and bases of the localised rational lattices. Then the terminal descent holds.
\end{lemma}
\begin{proof}[Outline]
In the charts of Definition~\ref{def:p4c2-charts}, use rational frames, flat base change, and coefficient-field retraction to identify the local lattice with the extended rational lattice. Descend the bound by coordinate valuations and transport it by the terminal isomorphism. The complete proof is in \S\ref{prf:B4.13a}.
\end{proof}

\begin{theorem}[The complete column estimates at ordinary places]\label{B4.14}
Fix \(\sigma\in\R\) and assume the terminal comparison \eqref{eq:p4c2-descend-terminal} for \(W\). The column theorem uses the two specialisations \(\sigma=0\) and \(\sigma=1\). Under the error hypotheses of \ref{B4.03}, assume the systems and comparisons of \ref{B4.09}. Assume \ref{B4.11}--\ref{B4.12} with their respective hypotheses. For every used twist and permitted model, assume the terminal estimate of \ref{B3.42} for every matching admissible system and kernel. This assumption also covers \(p\in Q'(\mu)\).

At \(p\notin\{T_u,g_0,g_1\}\) and for a used \(\mu\) with \(p\notin Q'(\mu)\), the filling and embedding identifications, the ramified-span equality and the base-term identity of \ref{B3.54k}(1)--(3), together with the error hypotheses of \ref{B4.03}, give, for all \(x\in V_{u,0}\),
\[
 x\in\mathcal N_{u,w}\ \Longrightarrow
 x\in\mathcal C_{u,\mu,w}(B^+_{u,w}(\sigma)).
 \tag{4.32}\label{eq:p4c2-enhanced-column}
\]
At the same places and twists, the base-term identity of \ref{B4.02} and the error hypotheses of \ref{B4.03} give, for \(x\in S^{\rm ram}_{u,w}\),
\[
 x\in\mathcal C_{u,\mu,w}(B^0_{u,w}(\sigma)).
 \tag{4.33}\label{eq:p4c2-ordinary-column}
\]
The served gain vanishes under any of
\[
 p\in\mathcal E,\quad p\text{ outside the inert and signal classes},
 \quad(p\text{ inert-class and }h_P=0),\quad p>X'',\quad p\in\mathcal O_{\rm odd}.
 \tag{4.34}\label{eq:p4c2-zero-gain}
\]
For every permitted model and used \(\mu\) with \(p\notin Q'(\mu)\), and each identification for \(Y_\mu\) as in \ref{B3.54k}(1), the error hypotheses of \ref{B4.03} give the local bound of \ref{B4.10}.
\end{theorem}
\begin{proof}[Outline]
Use \ref{B4.13}, \ref{B4.08}, and \ref{B4.09}--\ref{B4.10} for ordinary, zero-gain, and local bounds. Transport ties and descend by \eqref{eq:p4c2-descend-terminal}, and handle the empty set by \ref{B3.54k}(3) or \eqref{eq:p4c2-ordinary-gain} according to served gain. The complete proof is in \S\ref{prf:B4.14}.
\end{proof}

\section{Uniform constants and service for a family}

Let \(\mathfrak F_\varepsilon\) denote the parameter families of \ref{B1.18} at exponent \(\varepsilon\). Fix \(\mathcal F\in\mathfrak F_\varepsilon\), with members \(D_{x,i},L_{x,i}\) for \(i\in I_x\) and design and size hypotheses for \(x>x_1\). For each member put \(n=|A|\), \(Y=X''(n)\), and write \(\overline F_{A,j}(p)=n_P^{-1}\sum_{z\in Z_A}F_{\mathrm{Orb},j}(p,z)\) for the contribution from the kind-\(A\) positions \(Z_A\) of \ref{B4.03}. The uniform stabiliser bound means that there are \(C_D,x_D\in\R\) such that \(\log D\Gamma_u\le C_D fN_m\log(fN_m)\) for every \(x\ge x_D\), \(i\in I_x\) and word \(u\).

\begin{lemma}[Uniform size constants]\label{B4.15}
Put \(\theta=\varepsilon/100\), \(z=w''(n)\), \(h=t_{\rm ch}+1\) and let \(\Delta\) be the spacing parameter for the shifted weights. The size conditions give
\[
 \begin{gathered}
 x^{2+8\varepsilon-\theta}\le n\le x^{2+8\varepsilon+\theta},\qquad
 x^{9/2+10\varepsilon-\theta}\le k\le x^{9/2+10\varepsilon+\theta},\\
 t_{\rm ch}\eta\le z/x,\qquad mh\eta\le z,\qquad kS_0\le z,\qquad
 N_mm^{10}\le\Delta,\qquad 2\log\Delta\le x^{\varepsilon/4}.
 \end{gathered}
\]
 Use also \(z\le b'/x\) from the family scales.

Suppose \(\mathcal F\) satisfies the uniform stabiliser bound. There are \(C\ge0\) and \(x_0\in\R\) such that \eqref{eq:p4c2-size-contractions}--\eqref{eq:p4c2-size-orbits} hold simultaneously for every \(x\ge x_0\) with \(x>x_1\) and every \(i\in I_x\). For all \(u,1\le j\le m\) the first two bounds are
\[
 \sum_{t\in A}c_{u,j,t}\le2\#TT\,b+2fkb',
 \tag{4.38}\label{eq:p4c2-size-contractions}
\]
\[
 \log D\Gamma_u\le CfN_m\log(fN_m).
 \tag{4.39}\label{eq:p4c2-size-denominator}
\]
For \(1\le j\le m\) and every finite set \(S\) of inert-class or signal primes \(p\le Y\), excluding zero-shift inert primes,
\[
 \sum_{p\in S}\log p\,(I^*_{j,p}-[\overline F_{A,j}(p)]_+)
 \le Cfb\log Y.
 \tag{4.40}\label{eq:p4c2-size-orbits}
\]
\end{lemma}
\begin{proof}[Outline]
Interlacing and weight windows bound the contractions of \ref{B4.03}, and the stabilizer hypothesis bounds denominators. Apply Rosser--Schoenfeld to the harmonic error sum and use \(p\equiv\pm1\pmod{4q_e}\) at odd Frobenius order to bound the cubic sum. The complete proof is in \S\ref{prf:B4.15}.
\end{proof}

A property \(P\) is closed on the family class \(\mathscr C\) if
\[
 \begin{gathered}
 \exists\varepsilon_1>0\quad\forall\,0<\varepsilon\le\varepsilon_1\quad
 \exists\mathcal K_{\rm size}\quad\forall\mathcal F\in\mathfrak F_\varepsilon,\ C_H\in\R,\\
 \mathcal F\in\mathscr C\ \land\ (\mathcal F,C_H)\text{ meets the thresholds }\mathcal K_{\rm size}
 \ \Longrightarrow\ P(\mathcal F,C_H).
 \end{gathered}
 \tag{4.46}\label{eq:p4c2-closed}
\]
The preliminary constants have a fixed lower threshold. The middle threshold may depend on the outer constants. The final threshold may depend on both the outer and the middle constants, with \(C_H\) included among the latter. Finite conjunctions use the smallest \(\varepsilon_1\) and the componentwise largest thresholds. Let \(P_{x,i}(C)\) mean that \eqref{eq:p4c2-size-contractions}--\eqref{eq:p4c2-size-orbits} hold with constant \(C\), and that the enhanced and ordinary column estimates of \ref{B4.14} hold at \(\sigma=1\) under their respective hypotheses. Say that the family satisfies the uniform column and size bounds if
\[
 \exists C\ge0\ \exists x_0\in\R\ \forall x\in\R,\qquad
 (x\ge x_0\land x>x_1)\ \Longrightarrow\ \forall i\in I_x,\ P_{x,i}(C).
 \tag{4.47}\label{eq:p4c2-servf}
\]

\begin{theorem}[Uniform service for the family]\label{B4.16}
Assume that the uniform stabiliser bound and each of the following two conditions are closed on \(\mathscr C\) in the sense of \eqref{eq:p4c2-closed}.

First, beyond a family threshold each member admits levels and chains \(LC(\mu)\) on its filling. Form \(Y'\) by taking the levels and chain at \(\mu_w\) for each \(w\), and let \(W\) use the pair at \(\mu\) for each twist as in \ref{B3.54k}. Require the terminal comparison \eqref{eq:p4c2-descend-terminal}. Under the error hypotheses of \ref{B4.03}, for every used \(\mu\), identification for \(Y_\mu\) as in \ref{B3.54k}(1) and permitted \((u,w,h_g)\) with \(p\notin Q'(\mu)\), require the matching system and kernel of \ref{B4.09} and the complementary bounds of \eqref{eq:p4c2-complement} at \(\sigma=1\).

Second, beyond a family threshold each member admits kernel data \(E\). For every \(\mu\) and every level-free \((u,\mu,w)\), the class-pair scalars have norm one at each \(i<m\). For each used \(\mu\) at such a place, the pair maps carry evaluated rational vectors to evaluated rational vectors. Each reference or used twisted lift sends the complex form space at \(i\) into that at \(i+1\). In each of these chains every class at stage zero has trivial stabiliser. Choose a common system of letter data \(\xi_w,\epsilon_{w,z},\beta_{w,z}\) as above. Use the weights and \(E\) for the diagrams and series of \(\mathscr E\). Require the expansion, jet integrality, discriminant and shape conditions of Definition~\ref{def:p4c2-orbit}(2), (4) and \eqref{eq:p4c2-shape}. Require the bound without a neighbour operator of item 3 for the specified data, for every inert-class \(w\) with \(5\le p\le X''\), \(h_P\ne0\), even \(n_P\ge2\), every \(u,0\le i<m,t\in R_p\), and every kernel with the prescribed local expansion at \((0,t)\).

Then the uniform column and size bounds \eqref{eq:p4c2-servf} are closed on \(\mathscr C\). The conclusion also holds for a prescribed selection under the same closed stabiliser and twisted-chain conditions. In the second closed condition, require that for every family in \(\mathscr C\), every \(x\) at or above its selection threshold with \(x>x_1\), and every \(i\in I_x\), the selected kernel satisfies the stated requirements.
\end{theorem}
\begin{proof}[Outline]
Use \ref{B4.15} and splitting for \ref{B3.53}--\ref{B3.54}, and \ref{B3.45} for \ref{B3.42} within \ref{B4.09}--\ref{B4.10}. Use \ref{B4.10a}, \ref{B4.11a}, \ref{B4.11}--\ref{B4.12}, and \ref{B4.14} for column bounds while preserving \ref{B3.54k} and \ref{B4.07}. The complete proof is in \S\ref{prf:B4.16}.
\end{proof}

The period estimate of \ref{B4.19} gives exponent \(e_w(s-\delta_{u,w})\) from a column bound \(s\). Theorems \ref{B4.26} and \ref{B4.27} apply it to \eqref{eq:p4c2-enhanced-column} and \eqref{eq:p4c2-ordinary-column}, respectively.


% ---- p4c3 ----
\chapter{From columns to periods}\label{ch:columns-to-periods}

Fix the arithmetic member \(D\), its size data \(L\) and its twist exclusion data \(D_x\) as in \ref{B3.54k}~\cite{alpoge-2026}. Write \(K\) for its coefficient field, \(\iota:K\hookrightarrow\C\) for the chosen embedding and \(\mathcal O\) for its ring of integers.
For a finite place \(w\), write \(p_w,\mathfrak p_w,e_w\) for its
residue characteristic, prime ideal and ramification index.
Normalize \(v_w\) by \(v_w(p_w)=e_w\), and put
\(\mathcal O_{(w)}=\mathcal O_{K,\mathfrak p_w}\).
For a word \(u\), let \(T_u\) be its level prime, \(\Delta_u\) its
stabilizer denominator and \(\delta_{u,w}=v_{p_w}(\Delta_u)\).
Here \(v_p\) on positive integers is the exponent of \(p\). Let \(E\) be the host CM field with quadratic subfield \(K_d=\Q(\sqrt{-d})\). The \(f=[E:K_d]\) types \(\tau\) are those of \ref{B4.03}. The chain has
\(m\) lifts and ranks \(N_i=k+i(t+1)\), with \(t\) the odd chain parameter.

Fix the twisted forms \(T\). The used tuples are the reference
tuple \(\mu_0\) and the selected tuples \(\mu_w\), with
\(p_w\notin Q'(\mu_w)\), where \(Q'(\mu)\) is the finite exclusion set attached to \(\mu\) in \ref{B3.54k}. At the reference tuple this set is empty.
Write \(\Lambda^\mu_{u,j}\) for the lattice at stage \(j\) of the
\(\mu\)-chain and \(\jmath^\mu_{u,j}\) for its complex embedding.
Superscript \(0\) denotes the reference chain. Let \(V_u\) be the rational ramified input space of the word \(u\), let \(j_u:V_u\hookrightarrow V^0_{u,0}\) be its initial inclusion, and let \(N_u\) be its input lattice. Write \(M_{H,u}=M_u^{\mathrm{ram}}=j_u^{-1}(\Lambda^0_{u,0})\) for the original ramified lattice.
For \(\phi\in V_u\), write \(P_u^\mu(\phi,H)\) for the column period at \(\mu\) tested against \(H\), with complex input the image of \(\phi\) under the \(\iota\)-semilinear realization of \(V_u\). Lemma~\ref{B4.19a} expresses this period in terms of \(j_u\). An admissible
test lies in \(\widetilde M_\mu\cap V_0(\mu)\), where
\(\widetilde M_\mu\) is the test lattice and \(V_0(\mu)\) is the image
of the self-adjoint projector on the test space at \(\mu\) given by the central idempotent of the type. As in \ref{B4.05} put
\[
 S_{u,w}=\sg(u,w),\qquad B_{u,w}=\sum_{j=1}^m(F_j(w)-E_{u,j}(w)).
\]
For an \(\mathcal O\)-submodule \(\Lambda\), write
\(\Lambda_{(w)}=\mathcal O_{(w)}\Lambda\).

\section{Localizing the terminal pairing}

\begin{lemma}[A pairing bound survives localization]\label{B4.17}
Let \(\Lambda\) be an \(\mathcal O\)-submodule of a \(K\)-space \(V\),
let \(g:V\to\C\) be \(\iota\)-semilinear, and fix \(b\in\R\) and a
finite place \(w\). If every \(x\in\Lambda\) has \(g(x)=\iota(a_x)\)
with \(a_x\in K\) and \(a_x=0\) or \(v_w(a_x)\ge b\), then the same
holds on \(\Lambda_{(w)}\).
\end{lemma}
\begin{proof}
For \(z=\sum_\nu c_\nu x_\nu\) with \(c_\nu\in\mathcal O_{(w)}\),
semilinearity gives \(g(z)=\iota(\sum_\nu c_\nu a_{x_\nu})\).
Each nonzero summand has valuation at least \(b\).
\end{proof}

Let \(\mathcal T_u^\mu(z,H)\) be the pairing obtained by evaluating
and lowering \(\jmath^\mu_{u,m}z\) and testing against \(H\) pulled
back to the chain test space. The reference point hypothesis is item~(1) of \ref{B8.09} for every word. Evaluation is injective on class labels and induces bijections on their stabilizers. Evaluation followed by lowering takes the terminal forms into the point space and has a rational realization carrying the terminal lattice into the point lattice. This realization satisfies \eqref{eq:p8c1-point-norm} at every infinite place, and the complex map satisfies the squared-norm bound there stated. The twisted point hypothesis requires an integral rational realization of evaluation followed by lowering for each word and each used \(\mu\ne\mu_0\), as in item~(2) of \ref{B8.09}. Use the common stabilizer denominator of \ref{B4.03}: the denominator of every twisted chain equals \(\Delta_{u_0}\) for the distinguished word \(u_0\), and every class stabilizer order at every stage of every such chain divides it. Under this condition the reference period hypothesis requires, for each word, point-lattice vector and admissible test, one representative \(y\in K\) of the point period with \(y=0\) or \(v_w(y)\ge-e_w\delta_{u,w}\) at every finite \(w\).
Under the same denominator condition the twisted period hypothesis requires one representative for each such vector and test at each used \(\mu\ne\mu_0\), with the bound simultaneous over all \(w\) for which \(p_w\notin Q'(\mu)\). These are items~(3) and~(4) of \ref{B8.09}.

\begin{lemma}[The terminal denominator]\label{B4.18}
Assume the common stabilizer denominator and the reference and twisted point and period hypotheses above. Fix \(u,w\) and take \(\mu=\mu_0\) or a used \(\mu\ne\mu_0\)
with \(p_w\notin Q'(\mu)\). For \(z\in(\Lambda^\mu_{u,m})_{(w)}\)
and admissible \(H\) at \(\mu\), there is \(a\in K\) with
\[
 \iota(a)=\mathcal T_u^\mu(z,H),\qquad
 a=0\ \text{or}\ v_w(a)\ge-e_w\delta_{u,w}.
\]
\end{lemma}
\begin{proof}
For an integral terminal vector, the rational point map represents
its evaluated and lowered vector. Apply the period hypothesis to
its point-lattice image, then \ref{B4.17} in the terminal variable.
\end{proof}

\begin{lemma}[Class denominators on the constructed columns]\label{B4.18a}
Assume the common stabilizer denominator of \ref{B4.03}, with each point stabilizer injecting into the corresponding terminal stabilizer. Suppose that for every word \(u\) and used tuple \(\mu\), evaluation followed by lowering on the \(\mu\)-column is given in local class coordinates at each finite place \(w\) by an integral interlacing map, multiplied on \(\mu_0\) under reference frame changes by inverse values of a terminal level character of finite order taking values in \(\iota(K)\). Suppose that at each finite place \(w\) the local terminal lattice spans the terminal \(K\)-space, the point embedding into the complex point space is injective and characterizes vectors of the localized point lattice by \(w\)-integrality of their local class coordinates, and the point lattice is the intersection of its localizations across finite places. Suppose further that every admissible test \(H\) is an \(\mathcal O\)-linear combination of row generators whose rational row lifts have class polynomials in the \(w\)-local Weyl dual of the point lattice at every finite \(w\) for \(\mu_0\) and at every \(w\) with \(p_w\notin Q'(\mu)\) for other used \(\mu\), and that the point period is the Petersson class sum weighted by the reciprocal point stabilizer orders. Then for every word and used tuple, the rational point
map takes the terminal lattice into the point lattice. For a
point-lattice vector \(x\) and admissible \(H\), the point period
has a representative \(y=\sum_c b_c/h_c\in K\), where \(h_c\)
is the point stabilizer order at class \(c\). For \(\mu_0\) at every
\(w\), and for other used \(\mu\) at every \(w\) with
\(p_w\notin Q'(\mu)\), all \(b_c\) and \(\Delta_u y\) are integral
at \(w\).
\end{lemma}
\begin{proof}[Outline]
Integral interlacing and root-of-unity frame factors extend linearly from the terminal lattice and glue to the point lattice. Pairing with \(w\)-local Weyl-dual row lifts gives \(w\)-integral \(b_c\), and \(h_c\mid\Delta_u\) makes \(\Delta_u y\) integral at \(w\). The complete proof is in \S\ref{prf:B4.18a}.
\end{proof}

A column has bound \(s\) at \(w\) in the sense of \ref{B4.06}
if its terminal image is represented by \(az\), with
\(z\in(\Lambda^\mu_{u,m})_{(w)}\) and \(a\in K\) either zero or
of valuation at least \(e_ws\). For \(\mu\ne\mu_0\), first transport the initial reference vector by the comparison of \ref{B4.06}, then apply the twisted lifts.

\begin{lemma}[The two evaluations of a period]\label{B4.19a}
Write \(j:V_u\to V^0_{u,0}\) for the inclusion of the rational ramified subspace. Let \(E_0\) be the identity on reference initial class functions and \(E_\mu\) the initial class-function comparison for \(\mu\ne\mu_0\). Write \(C^\mu\) for the composite of the lifts on the corresponding chain. For a complex terminal class function \(Y\), let \(t^\mu(Y,H)\) be its evaluation followed by lowering and pairing with the pullback of \(H\) to that chain. Then
\[
 P_u^\mu(\phi,H)=t^\mu\!\left(C^\mu E_\mu\jmath^0_{u,0}(j\phi),H\right).
\]
If \(Y=\jmath^\mu_{u,m}(z)\), the right terminal pairing is \(\mathcal T_u^\mu(z,H)\).
\end{lemma}
\begin{proof}[Outline]
Apply the complex input identification to recover \(\jmath^0_{u,0}(j\phi)\), then substitute into the reference and twisted definitions. On an embedded rational terminal vector this is the pairing of \ref{B4.18}. The complete proof is in \S\ref{prf:B4.19a}.
\end{proof}

\begin{theorem}[From a column bound to a period bound]\label{B4.19}
Assume the hypotheses of \ref{B4.18}. Fix a word \(u\), a finite
place \(w\), and \(\phi\in V_u\). Take \(s\in\R\) and a used
\(\mu\) with \(p_w\notin Q'(\mu)\). If the \(\mu\)-column from
the initial vector \(j\phi\) of \ref{B4.19a} has bound \(s\), then for every admissible
\(H\) at \(\mu\) there is \(y\in K\) with
\[
 \iota(y)=P_u^\mu(\phi,H),\qquad
 y=0\ \text{or}\ v_w(y)\ge e_w(s-\delta_{u,w}).
\]
\end{theorem}
\begin{proof}[Outline]
Write the column image as \(\jmath^\mu_{u,m}(az)\) with \(v_w(a)\ge e_ws\) by \ref{B4.06} and \ref{B4.19a}, and apply \ref{B4.18} to \(z,H\) for \(y_0\in K\) with \(v_w(y_0)\ge -e_w\delta_{u,w}\). Linearity in \ref{B4.19a} gives \(y=ay_0\), and adding valuations yields the bound. The complete proof is in \S\ref{prf:B4.19}.
\end{proof}

\section{The level prime}

\begin{lemma}[The budget above \(T_u\)]\label{B4.20}
Suppose the level primes belong to the split class and the base terms and losses have the values in \ref{B4.02} and \ref{B4.03}. If \(p_w=T_u\), then the local size parameter \(\tau(w)\) of \ref{B3.52} vanishes and, for \(1\le j\le m\),
\[
 g^{\rm raw}_j(w)=F_j(w)=0,\qquad E_{u,j}(w)=T^*_{u,w},
 \qquad S_{u,w}=-mT^*_{u,w},
\]
where \(g^{\rm raw}_j(w)\) is the gain before screening in \ref{B4.04} and \(T^*_{u,w}\) is the level-prime loss defined in \ref{B4.03}.
\end{lemma}
\begin{proof}
The level prime \(T_u\) belongs to the split class and is distinct from both fixed auxiliary primes. Outside the inert and signal classes the size data give \(\tau(w)=g^{\rm raw}_j(w)=0\). The split class and the fixed-prime exclusion give \(F_j(w)=0\) by \ref{B4.02}. Apply \ref{B4.03} to the loss and \ref{B4.04} to the screened gain.
\end{proof}

\begin{lemma}[The level prime exceeds the stabilizer cutoff]\label{B4.21b}
Write \(t\) for the odd chain parameter and \(N_j=k+j(t+1)\) for the rank after \(j\) lifts. The rank conditions on a member are
\[
 2\le m,\qquad m+3\le k,\qquad f\le k,\qquad
 k+m-1+5\le t,\qquad 3m<t .
\]
Its size condition at \(x>x_1\) includes
\[
 (Q_Sqd)^{36}t^{150}\le\sqrt b .
\]
If the level prime of a word \(u\) satisfies \(b^{m/2}\le T_u\), then \(2fN_m+1<T_u\).
\end{lemma}
\begin{proof}
The rank conditions give \(t\ge7\), \(m\le t\) and \(f\le k\le t\). Thus
\[
 N_m\le t^2+2t,\qquad
 2fN_m+1\le2t^3+4t^2+1\le7t^3<t^5 .
\]
The host conductor factor \(Q_S\) is positive. Indeed, the level prime is coprime to \(Q_S\), whereas its greatest common divisor with zero would be \(T_u>1\). Since \(q,d\) are primes, \(Q_Sqd\ge1\). The displayed size condition now gives \(t^{300}\le b\). In particular \(b\ge1\), so \(t^5\le b\le b^{m/2}\le T_u\).
\end{proof}

\begin{definition}[The design class]\label{def:p4c3-design}
Use the arithmetic families of \ref{B1.18}, with size and column data as in \ref{B4.16}. For a family class \(\mathcal C\) and \(C_0\in\R\), write \(\mathcal D(\mathcal C,C_0)\) for the subclass with the following design conditions. All point cutoffs are uniform in the member index.

The family has the good twists used in the design construction of Chapter~\ref{ch:p8c1} and satisfies its choice conditions with \(C_0\), some \(C_H\in\R\) and positive constants \(C_5,c_6,C_b\). These are the lower bound on the thick set, the bounds on the number of top positions and on bad orbits, the estimates for the three quadratic characters and the thin character sums, and the two bounds on word sizes and word count. In the latter two bounds the required inequalities are
\[
 C_0^{-1}\le \frac{d_u\Gamma^2}{D_*}\le C_0
 \quad\hbox{for every word }u,\qquad
 C_0\Gamma^2x^{C_H}\le \#\{\text{words}\},
\]
where \(d_u\) is the dimension attached to the word, \(D_*\) the reference dimension scale and \(\Gamma=\exp(C_\Gamma m^5fL_0)\), with \(L_0\) the logarithmic scale of the member.

Beyond a family cutoff, for every member and every word \(u\), each prime divisor of \(\Delta_u\) is at most \(2fN_m+1\), and
\[
 b^{m/2}\le T_u,\qquad
 \log T_u\le C_TmL_0,\qquad
 \log\Delta_u\le C fN_m\log(fN_m).
\]
Here \(C_T\) is the family's level-size constant, and one \(C\in\R\) works for the family.

For a class \(\mathcal A\), write \(\mathcal A_{\rm ch}\) for the subclass used in the construction before \ref{B8.01}. It requires \(C_{YL}\le9\), the zero-free boxes and quadratic-prime conditions of \ref{B8.47}, tame distinctness at both fixed primes, and eventual divisibility of every \(\Delta_u\) by the following convention:
\[
 \Delta_u\mid M(2fN_m),\qquad
 M(a)=\prod_{\substack{p\le a+1\\p\ {\rm prime}}}
 p^{\sum_{j\ge0}\lfloor a/(p^j(p-1))\rfloor}.
\]
At a fixed prime \(g\), tame distinctness requires the tame parameters to be squares and the list consisting of \(1\), the parameters and their inverses to have distinct entries. Their prime-to-\(\ell\) parts must also remain distinct for each prime \(\ell\ne g\) outside the exceptional set
\(\{\ell:\exists\,1\le d\le r,\ g^d<\ell^{4v_\ell(g^d-1)}\}\), where \(r\) is the rank parameter. The tame, box and quadratic-prime conditions hold eventually at every member.
\end{definition}

\begin{lemma}[Eventual coprimality of the level prime]\label{B4.21}
Fix a family class \(\mathcal C\) and \(C_0,\varepsilon\in\R\).
For every arithmetic family \(\mathfrak F\) of parameter \(\varepsilon\) in the design class \(\mathcal D(\mathcal C,C_0)\) of Definition~\ref{def:p4c3-design}, there is \(x_0\in\R\) such that
\(T_u\nmid\Delta_u\) for every word of every member \(D_{x,i}\)
with \(x\ge x_0\), \(x>x_1\) and \(i\in I_x\), where \(x_1\) is the family's initial threshold and \(I_x\) its index set at \(x\).
This eventual assertion is closed on \(\mathcal D(\mathcal C,C_0)\) in the sense of \eqref{eq:p4c2-closed}.
\end{lemma}
\begin{proof}[Outline]
Taking \(x_0=\max(x_s,x_T)\) from the uniform stabilizer and level cutoffs lets \ref{B4.21b} give \(T_u>2fN_m+1\), excluding \(T_u\mid\Delta_u\) by primality of \(T_u\). Setting the \(\varepsilon\)-cutoff to \(1\) and size thresholds to \(0\) in \ref{B4.16} gives closedness. The complete proof is in \S\ref{prf:B4.21}.
\end{proof}

\begin{proposition}[One input lattice for a word]\label{B4.21a}
Fix a word \(u\) of an arithmetic member. Write \(j=j_u\) and let \(V^r=j(V_u)\subseteq V^0=V^0_{u,0}\). Put
\(L^0=\Lambda^0_{u,0}\) and \(S^{\mathrm{ram}}_{u,w}=(L^0)_{(w)}\cap V^r\).
Define \(\mathcal O_{(w)}\)-submodules \(B_w\subseteq V_u\) as follows.

At a permitted place with \(W_{p_w}\ne\varnothing\), take
\(B_w=j^{-1}J_w\) using \ref{B4.07c}. Equivalently \(x\in B_w\)
if \(c(jx)\in N_P\) for the comparison map and enhanced lattice of
that model. The permitted prime conditions in \ref{B4.07} certify
the model at the given \(u,w\).

At \(p_w=T_u\), let \(I^{\rm top}_{u,w}\) consist of the level indices
\(v\) with top type \(\tau_w(v)\) when \(w\) is seen in the sense
of \ref{B4.03}. Take \(I^{\rm top}_{u,w}=\varnothing\) otherwise. The
upward operator \(U_v\) on zeroth columns sums the values at the
upper neighbors indexed by residue vectors of marking one, with
the values pulled back by the frame maps. Put
\[
 U=\prod_{v\in I^{\rm top}_{u,w}}U_v,\qquad
 b_v=-\lambda^u_{0,\tau_w(v),k-1},\qquad
 a_0=m\sum_{v\in I^{\rm top}_{u,w}} b_v ,
\]
where the product is in the level order and the weight coordinates
start at zero. Let \(N_{\lambda,u,w}\) consist of rational
\(y\in V^0\) for which
\[
 \jmath^0_{u,0}(y)=T_u^{a_0}U^m\jmath^0_{u,0}(z)
 \quad\text{for some }z\in S^{\mathrm{ram}}_{u,w}.
\]
At these places put \(B_w=j^{-1}N_{\lambda,u,w}\). At the remaining
places put \(B_w=(M_{H,u})_{(w)}\). A permitted prime does not
divide \(T_u\), so the first two prescriptions have disjoint domains.

The intersection \(N_u=\bigcap_{w<\infty}B_w\) is a full
\(\mathcal O\)-lattice in \(V_u\) and \((N_u)_{(w)}=B_w\) at
every finite place.  Away from the level prime one has
\[
 (N_u)_{(w)}\subseteq(M_{H,u})_{(w)}\qquad(p_w\ne T_u).
\] There is equality when \(p_w\ne T_u\) and
\(W_{p_w}=\varnothing\).
\end{proposition}
\begin{proof}[Outline]
Use \ref{B4.07c} at selected places and invert the upward operators by \ref{B5.05a} at level places to get commensurable lattices. The exceptions are finite by \ref{B3.51}, so glue by \ref{B4.35}. The complete proof is in \S\ref{prf:B4.21a}.
\end{proof}

\begin{theorem}[Periods above the level prime]\label{B4.22}
Assume the hypotheses of \ref{B4.18}, the level-prime identities of \ref{B4.20} and the description \((N_u)_{(w)}=j_u^{-1}N_{\lambda,u,w}\) of \ref{B4.21a} at every \(w\mid T_u\). Suppose \(T_u\nmid\Delta_u\) for every word.
Assume both sharp column estimates of \ref{B4.22a}: scaling the reference column by \(\iota(c)\) makes it integral for \(v_w(c)\ge e_wmT^*_{u,w}\), and each used nonreference column outside its exclusion set has bound \(-mT^*_{u,w}\), in both cases on \(N_{\lambda,u,w}\). For \(u,w\) with \(p_w=T_u\),
\(\phi\in N_u\) and admissible \(H\) at \(\mu_w\), there is \(y\in K\)
with
\[
 \iota(y)=P_u^{\mu_w}(\phi,H),\qquad
 y=0\ \text{or}\ v_w(y)\ge e_w S_{u,w}.
\]
\end{theorem}
\begin{proof}[Outline]
Using \(n_0=mT^*_{u,w}\in\mathbb{N}\) from \ref{B4.03}, apply \ref{B4.22a} (with \(c=T_u^{n_0}\) at \(\mu_0\)) to get column bound \(-n_0\) on \(j_u\phi\). Apply \ref{B4.19} with \(\delta_{u,w}=0\), and use \ref{B4.20} over the \(m\) lifts to identify \(S_{u,w}=-n_0\). The complete proof is in \S\ref{prf:B4.22}.
\end{proof}

\begin{lemma}[Binomial moments and simultaneous approximation]\label{B4.22d}
Let \(p\) be prime, \(r\geq0\), and \(b_\nu(c)=\prod_{j=1}^r\binom{c_j}{\nu_j}\) for \(\nu\in\N^r\) and \(c\in\Z_p^r\). Given \(A_\nu\in\mathcal O_{\C_p}\), there is a unique continuous \(\C_p\)-linear functional \(\mu\) on \(C(\Z_p^r,\C_p)\) with \(\mu(b_\nu)=A_\nu\). Its norm is at most one. For any finite family \(g_i\) of continuous functions and any \(\eta>0\), there are finitely many \(c\in\N^r\) and coefficients \(d_c\in\mathcal O_{\C_p}\), common to the family, with
\[
 \left|\mu(g_i)-\sum_c d_c g_i(c)\right|_p<\eta\qquad\text{for every }i .
\]
\end{lemma}
\begin{proof}[Outline]
Finite differences give the binomial basis norm identity, and Mahler interpolation gives density, so the moment functional extends with norm at most one. Approximate all tests by step functions on one residue partition. The complete proof is in \S\ref{prf:B4.22d}.
\end{proof}

\begin{lemma}[Standard isometry and the two determinant factors]\label{B4.22e}
Use the host \(E/K_d\) of \ref{B2.01}, the CM type \(\Sigma\), and the normalized curve lattice \(\Lambda=\delta\mathfrak c\) of \ref{B5.07x}. Put \(f=[E:K_d]\). In a coordinate product put
\[
 H(z,w)=\frac{\sum_i z_i\bar w_i}{\operatorname{area}(\Lambda)},\quad
 B(z,w)=\Im H(z,w),\quad
 \alpha(z)=\exp\!\left(\frac{\pi i\sqrt d}{2}H(z,z)\right).
\]
Define \(U_yF(z)=F(z+y)\exp(-\pi H(z,y)-\pi H(y,y)/2)\) and \(W_y=\alpha(y)U_y\). For commensurable self-dual period lattices \(P,Q\) with their indicated semicharacters, write
\[
 I_{P,Q}F=[Q:P\cap Q]^{-1/2}\sum_{y\in Q/(P\cap Q)}W_yF
\]
on \(P\)-fixed sections. The square root is positive real. Let
\(P_1=\{(\delta\tau(x))_{\tau\in\Sigma}:x\in
\mathfrak c\mathcal O_E\mathfrak b^{-1}\}\), where
\(\mathfrak b\bar{\mathfrak b}=\mathfrak D_{E/K_d}\).
Let \(\Theta\) be the one-entry theta section with the fourth-power normalization described below. For odd \(p\) split in \(K_d\) and prime to \(N\mathfrak c\), there is an isometry \(g\) such that \(Q=g(P_1^4)\) agrees with \(P=\Lambda^{4f}\) locally at \(p\) and
\[
 \Theta^{\otimes4}(z)=e\,(I_{P,Q}\theta_\Lambda^{\otimes4f})(gz),
 \qquad |\jmath(e)|_p=1 .
\]
For even ranks \(a,b\), put \(F_0=\Theta^{\otimes ab}\) on the standard matrix lattice. Under a unitary source change \(\gamma\), its pullback equals its normalized lattice average times
\((\prod_\tau\tau(\overline{\det\gamma}))^{b/2}\).
Under a unitary target change \(\gamma'\), the factor is
\((\prod_\tau\tau(\det\gamma'))^{a/2}\). The pullback in both variables is \(F(z)\mapsto F(\gamma'^{*}z\gamma)\), with matrices read at the types.
\end{lemma}
\begin{proof}[Outline]
Use four squares and Cayley approximation to match the lattices at \(p\). Gaussian integration and finite Gauss sums compare pullbacks with theta averages, giving a \(p\)-adic unit normalization and, by repetition across rows or columns, the endpoint determinant factors. The complete proof is in \S\ref{prf:B4.22e}.
\end{proof}

\begin{lemma}[Integral moments of the matrix theta vector at a split prime]\label{B4.22c}
Let \(K_d\) and the selected CM type \(\Sigma\) of \(E\) be those
of the host of \ref{B2.01}, and fix a field isomorphism
\(\jmath:\C\simeq\C_p\). Suppose \(p\ne2\) splits in \(K_d\).
Let \(\mathfrak c\) be the nonzero curve ideal, of norm prime to
\(2dp\), and \(\Lambda=\delta\mathfrak c\subset\C\) its period
lattice with the discriminant normalization of \ref{B5.07x}.
Write \(\mathfrak b\bar{\mathfrak b}=\mathfrak D_{E/K_d}\) for the
host factorization of the relative different. For self-dual lattices
\(L,L'\) in the standard Hermitian spaces \(E^a,E^b\) with \(a,b\)
even, put
\[
 M(L,L')=(\mathfrak c\mathcal O_E)\mathfrak b^{-1}
 \operatorname{Hom}_{\mathcal O_E}(L,L'),\qquad
 \Lambda(L,L')=\{(\delta\tau(x))_{\tau\in\Sigma}:x\in M(L,L')\}.
\]
Let \(\omega=\Omega_0/\jmath(\delta_0)\), where \(\delta_0\) is the
discriminant root of the principal curve lattice and \(\Omega_0\)
its normalized split period. For a holomorphic function \(F\)
and a list \(v=(v_1,\ldots,v_r)\) of period directions, put
\[
 A_\nu(F;v)=\jmath\left(
 \left[\prod_{j=1}^r\binom{D_{\jmath^{-1}(\omega)v_j}}{\nu_j}\right]F(0)\right),
 \qquad
 \binom Dn=\frac{D(D-1)\cdots(D-n+1)}{n!}.
\]
Here \(D_v\) denotes directional differentiation and \(\binom D0=1\).
Suppose \(|A_\nu|_p\le1\) for the normalized even theta of \(\Lambda\)
and its prime-to-\(p\) torsion translates, for all lists in \(\Lambda\)
and all \(\nu\in\N^r\). Then the same bound holds for lists in
\(\Lambda(L,L')\) and prime-to-\(p\) torsion translates of the
character-normalized matrix theta \(e_{L,L'}\) constructed below.
For each list the derivatives in the scaled directions are power
moments of a \(\C_p\)-valued measure of norm at most one on \(\Z_p^r\).
\end{lemma}
\begin{proof}[Outline]
Apply \ref{B4.22e} to the fourth tensor power to get integrality for the standard matrix theta vector from the elliptic case. Prove lattice-move equivariance via \ref{B4.22e} and \ref{B2.01}, align \(L,L'\) above \(p\) by unitary maps, and get the measure from \ref{B4.22d}. The complete proof is in \S\ref{prf:B4.22c}.
\end{proof}

\begin{lemma}[Translation of the normalized elliptic theta section]\label{B4.22f}
Let \(d\) be the host prime, \(d\equiv7\bmod16\), and let the odd prime \(p\) split in \(K_d\). On the principal curve with lattice \(\mathcal O_{K_d}\), let \(\theta\) be the even theta section with \(\theta(0)=1\) and zero at the nonzero half-period \(1/2\). Use \(U,W,H\) of \ref{B4.22e} for this lattice. For a complex analytic germ \(F\), write \(\widehat F\) for its Taylor expansion transported by \(\jmath:\C\simeq\C_p\) and composed with \(\omega\log(1+Z)\), in the normalized ordinary multiplicative chart. Write \(S_\zeta f(Z)=f(\zeta(1+Z)-1)\) on bounded series when \(|\zeta-1|_p<1\).

Suppose \(s\ge1\), \(p\nmid s\), \(st\in\mathcal O_{K_d}\), \(n\ge0\), and
\[
 y=sv/p^n,\qquad v\in\mathcal O_{K_d},\qquad
 |\jmath(\bar v)|_p\le p^{-2n}.
\]
Put \(\zeta_y=\jmath(\exp(-2\pi i\,\Im H(1,y)))\). Then
\[
 \zeta_y^{p^n}=1,\qquad
 \widehat{W_yU_t\theta}=S_{\zeta_y}\widehat{U_t\theta}.
\]
\end{lemma}
\begin{proof}[Outline]
Normalize the multiplicative chart by the period pairing. The elliptic theta quotient gives integral coefficients and a functional equation forcing the translation ratio to be one, and elliptic translation extends the identity to \(U_t\theta\). The complete proof is in \S\ref{prf:B4.22f}.
\end{proof}

\begin{lemma}[One measure for all allowed matrix translations]\label{B4.22g}
Let \(M=M(N,N')\) and \(e=e_{N,N'}\) be as in \ref{B4.22c} for the constructed family. Fix a prime-to-\(p\) translate \(t\), and choose \(N_0>0\) prime to \(p\) with \(N_0t\) a period. Use the embedding of rational matrices into the analytic period space in \ref{B4.22c}. Let \((\beta_j)_{j=1}^r\) in \(M\) be a \(K_d\)-basis that generates the period module locally at both primes over \(p\). Put \(x(c)=\sum_jc_j\beta_j\) for \(c\in\N^r\). Write
\[
 [y,x]=\exp(-2\pi i B(x,y)),\qquad
 M_x(\nu)=\prod_{\tau,k,l}
 \bigl(u_\tau\iota_w(\tau(x^\dagger)_{kl})\bigr)^{\nu_{\tau kl}},
\]
where \(x^\dagger\) is the Hermitian adjoint, \(\iota_w:K\hookrightarrow\C_p\) induces \(w\), and \(u_\tau\) is the differential generator at type \(\tau\) divided, after embedding, by the ideal-curve period. Write \(\partial_c^\nu\) for the coordinate derivative with the differential generators inserted. If \(v\) is the quotient of the conjugate level factors with valuation \(-1\) in the chosen split component and \(+1\) in the other, then there is a single measure \(\mu_t\) of norm at most one such that for every \(y\in N_0vM\) and every jet index \(\nu\),
\begin{equation}\label{eq:B4.22b-shared}
 \jmath\bigl(\partial_c^\nu(W_y W_t e)(0)\bigr)
 =\jmath(\alpha(t))\int_{\Z_p^r}\jmath([y,x(c)])M_{x(c)}(\nu)\,d\mu_t(c).
\end{equation}
The integrand denotes the continuous extension from \(\N^r\). The same statement holds for a self-dual companion pair with its own fixed \(t\).
\end{lemma}
\begin{proof}[Outline]
Show prime-to-\(p\) shifts preserve the Gauss norm via the special-fibre valuation and \ref{B4.22f}, so Lemma~\ref{B4.22c} applies to \(U_t e\). Expand \(e\) in an orthogonal frame, push forward the measure of Lemma~\ref{B4.22d}, and transport \(W_y\) using \ref{B4.22f}. The complete proof is in \S\ref{prf:B4.22g}.
\end{proof}

\begin{lemma}[The two finite quotients in the split-level moment calculation]\label{B4.22h}
Let \(p=\varpi\bar\varpi\) be a split level prime as in Lemma~\ref{B4.22g}. Write \(M=M(N,N')\) for its period module. At the selected primes over \(\varpi\) let the target have surjective residue markings. Replace the target lattice there by the kernels of the markings, at their conjugates by the Hermitian duals of these kernels, and keep all other completions. Denote the resulting period module by \(M'\) and put \(I=M\cap M'\). Let \(B\) be the half-period module obtained by multiplying \(M\) by \(pg\) at the selected primes and row primes, by \((pg)^{-1}\) at their conjugates, and by \(1\) elsewhere. Here \(g\) is supported on the rational row primes, which differ from \(p\). Let \(c_A\) be supported on the same row primes, and let \(F=fa\), where \(a\) is the source rank and \(f\) is the number of selected level primes. The residue fields there have order \(p\).

Put \(H_S=\operatorname{Hom}_S(N,N')\), \(Z_0=pgH_S\), \(r=c_Ag\bar\varpi\). The curve and different factors in \(M\) are units at the test primes. Then
\[
 [M:I]=[M':I]=p^\ell,\quad 0\le\ell\le F,\qquad
 |rI/(rI\cap Z_0)|=[I:\varpi M].
\]
If \(n\mid[M:M\cap B]\) and \(J=rI/(rI\cap Z_0)\), then
\(v_p(n)\le\ell+v_p(|J|)\). In \(B/(B\cap M)\) the annihilator of \(rI\) under \(e^{2\pi i B_H(-,-)}\), where \(B_H=\Im H\), has prime-to-\(p\) part the image of \(pB\) and \(p\)-part the image of \(M'\cap B\). The latter has order \(p^\ell\).
\end{lemma}

\begin{proof}[Outline]
Use the local marking equations, the preimage of \(Z_0\), and the chain of lattice indices to obtain the index and valuation bounds. Duality identifies the \(p\)-part of the annihilator, and integrality identifies its prime-to-\(p\) part. The complete proof is in \S\ref{prf:B4.22h}.
\end{proof}

\begin{lemma}[The point kernel at the two halves of the split level]\label{B4.22b}
Use the constructed theta family with rational jets and unramified
level characters. Fix a used tuple \(\mu\), a word \(u\) and a stage
\(0\le i<m\). Write \(p=T_u=\varpi\bar\varpi\),
\(a=N_i,\ b=N_{i+1}\) and \(s=(b-a)/2=(t+1)/2\) in this lemma.
Fix decorated points \(d,d'\) with self-dual lattices \(N,N'\)
of ranks \(a,b\), markings \(l_P,l'_P\) at \(P\mid\varpi\) and
adapted frames at the selected row primes \(S_{\rm row}\).
Write \(\operatorname{Hom}_S(N,N')\) for rational \(E\)-linear
maps integral on the local lattices at the primes in \(S\).
Put
\[
 S=S_{\rm row}\cup\{P:P\mid\varpi\},\quad
 q_0=\prod_{P\in S_{\rm row}}|k(P)|,\quad
 c_A=q_0^{A_{\rm row}},\quad g=q_0^{B_{\rm row}},\quad c_B=pg ,
\]
where \(k(P)\) is the residue field and \(A_{\rm row},B_{\rm row}\)
are the support and period exponents of the row selection.
The row test \(\Xi_0:\operatorname{Hom}_E(E^a,E^b)\to K\) has
support in \(c_A^{-1}\operatorname{Hom}_{S_{\rm row}}(N,N')\),
is integral at \(w\mid p\), and is invariant under
\(g\operatorname{Hom}_{S_{\rm row}}(N,N')\).
The set \(S\) is disjoint from its conjugate. By \ref{B8.54} there
are \(f\) primes over \(\varpi\), each of residue field \(\F_p\),
and \(p\) is a unit at the row primes.

For the decorated source \(d\), let \(L_d\subseteq N\) be the sublattice cut out by \(l_P=0\) at every \(P\mid\varpi\), and write \(L_d^\vee\) for its Hermitian dual. Define \(L_{d'}\subseteq N'\) using the target markings \(l'_P\), and write \(L_{d'}^\vee\) for its dual.
Let \(S_w(M)\) denote the local Weyl value lattice on \(M\)
with the weight at its stage. The ordinary source condition uses \(S_w(L_d^\vee)\), which is the condition on \(N\) at a seen place. Inserting all target level indices at such a place gives \(S_w(L_{d'})\). For the point kernel \(k(d,d')\) with row test
\(\Xi_0\), one has
\[
 \begin{aligned}
 p^{fa-\sum_\tau h'_\tau}k(d,d') S_w(L_d^\vee)
 &\subseteq S_w(L_{d'}) &&(w\text{ seen}),\\
 k(d,d') S_w(L_d^\vee)
 &\subseteq p^{fs}S_w(L_{d'}^\vee) &&(w\text{ unseen}),
 \end{aligned}
\]
where \(h'_\tau\) is the last coordinate of the target point weight \(\lambda'_\tau\).
The size hypothesis is \(2\,\operatorname{spread}(\lambda'_\tau)<p\)
at every type, where spread is the largest minus the smallest
coordinate plus \(b\).
With data away from \(P\) fixed, the upper-neighbor relation
\(e\uparrow_P e'\) at \(P\mid\varpi\) gives
\[
 \sum_{d'\uparrow_P e'} k(d,e')=p^s\sum_{e\uparrow_P d} k(e,d').
\]
\end{lemma}
\begin{proof}[Outline]
Bounds coset theta jet tables via \ref{B4.22c}, \ref{B4.22d}, \ref{B4.22g}, and \ref{B4.22h}, and passes to Weyl lattices via Jantzen's sum formula. Unfolding neighbor sums onto a common lattice using \ref{B4.22c} and \ref{B4.22f} proves the neighbor identity. The complete proof is in \S\ref{prf:B4.22b}.
\end{proof}

\begin{proposition}[The sharp column exponent at the level]\label{B4.22a}
Take the constructed columns of \ref{B8.09} with their rational
jets, weight signs and compatible level-zero lattices. Use the
size conditions of \ref{B4.22b} and \(T_u>2fN_m+1\).
Fix \(u,w\) with \(p_w=T_u\). For \(x\in N_{\lambda,u,w}\)
and \(c\in K^\times\) with \(v_w(c)\ge e_wmT^*_{u,w}\),
the reference column of \(x\) scaled by \(\iota(c)\) belongs
to the image of its terminal local lattice. For a used
\(\mu\ne\mu_0\) with \(p_w\notin Q'(\mu)\), the twisted column
of \(x\) has bound \(-mT^*_{u,w}\).
\end{proposition}
\begin{proof}[Outline]
Use \ref{B4.22b} with unit stabilizer orders to commute upward sums through lifts and iterate to \(mT^*_{u,w}\) by \ref{B4.03}. Transport the input of \ref{B4.21a} through the level-zero identification and read the twisted bound by \ref{B4.06}. The complete proof is in \S\ref{prf:B4.22a}.
\end{proof}

\section{The two fixed primes}

Fix local error data \(X_e\). With the notation of \ref{B4.03} put
\[
 \nu_{i,u,j,w}=e_w^{-1}\sum_\tau A_{u,j,\tau}v_w(\tau(\bar\pi_i)),
 \qquad \mathcal E_G=f\bigl(mN_m(N_m+t)+m^3\eta\bigr).
\]
The two fixed-prime hypotheses use the same constant \(C_1\in\R\).
For \(p_w=g_i\ne T_u\), the tuple bound requires, for every used
\(\mu\), \(\phi\in M_u^{\mathrm{ram}}\) and
\(H\in\widetilde M_\mu\), a representative \(y\in K\) of
\(P_u^\mu(\phi,H)\) with \(y=0\) or
\(v_w(y)/e_w\ge-\sum_j\nu_{i,u,j,w}-C_1\mathcal E_G\).
At the same places the budget inequality is \(B_{u,w}\le-\sum_j\nu_{i,u,j,w}-C_1\mathcal E_G\).

\begin{lemma}[The bound on the original lattice at \(g_i\)]\label{B4.23}
For \(p_w=g_i\), \(i\in\{0,1\}\), one has \(S_{u,w}=B_{u,w}\).
If \(p_w\ne T_u\) and the fixed-prime hypotheses hold, then for
\(\phi\in M_u^{\mathrm{ram}}\) and \(H\in\widetilde M_{\mu_w}\)
there is \(y\in K\) representing \(P_u^{\mu_w}(\phi,H)\) with
\(y=0\) or \(v_w(y)\ge e_w B_{u,w}\).
\end{lemma}
\begin{proof}
By \ref{B4.04} the screened gain vanishes above \(g_i\).
Apply the tuple bound at \(\mu_w\) and the budget inequality at \(w\) to obtain \(v_w(y)/e_w\ge B_{u,w}\) for \(y\ne0\).
\end{proof}

\begin{proposition}[The common constant at the two fixed primes]\label{B4.23a}
Use the constructed rows of \ref{B8.06}. For each used \(\mu\), assume that one choice of host, character and CM theta data realizes the period comparison for every word, as follows. Use the induced CM type and the abelian unramified extension of the host construction, with trivial auxiliary ideal. Choose the primes over the real subfield and the characters with the prescribed infinity weights and Weil character. Use the discriminant twelfth root, theta character and differential generators of that CM realization for all the word comparisons. It identifies the rational column and row lattices with the inputs and their rational pairing with the complex period. It gives rational normalized jets
on the summation support at every positive divisor of the
product of the support denominator and test modulus.
Put \(D_{u,j}=\sum_\tau A_{u,j,\tau}\). Suppose for \(p=g_0,g_1\)
and all \(u,\ 1\le j\le m\) that
\[
 2fN_m+1<p,\qquad 2D_{u,j}+2<p,\qquad D_{u,j}+N_m<p .
\]
Then the tuple bound holds with \(C_1=1\).
If the base terms have the values in \ref{B4.02} and the local error data satisfy \ref{B4.03}, the fixed-prime budget inequality holds for
\(C_1\ge0\) with \(2C_1\le c_g\), the fixed-prime error
coefficient of \ref{B4.03}. For \(0<\varepsilon\le1/100\),
a selected family with middle size threshold \(2\) satisfies
both bounds with \(C_1=1\) beyond a cutoff uniform in the
member index.
\end{proposition}
\begin{proof}[Outline]
Reduces the Hurwitz jet bound to elliptic theta translates via \ref{B4.22c} and \ref{B4.22e}--\ref{B4.22g} by formal-group root extraction. Combines \ref{B4.22a}--\ref{B4.22b} and Lemma~\ref{B3.10a} for \(C_1=1\), then \ref{B4.02}--\ref{B4.03} for the budget and family bounds. The complete proof is in \S\ref{prf:B4.23a}.
\end{proof}

\begin{theorem}[Periods above the fixed primes]\label{B4.24}
Assume the two fixed-prime inequalities above \ref{B4.23} and the inclusion \((N_u)_{(w)}\subseteq(M_{H,u})_{(w)}\) away from \(T_u\). Suppose that for every word, tuple and test, multiplying the complex column input by any \(c\in\C\) multiplies its period by \(c\). If \(p_w=g_i\ne T_u\), then for every
\(\phi\in N_u\) and \(H\in\widetilde M_{\mu_w}\) there is \(y\in K\)
representing \(P_u^{\mu_w}(\phi,H)\) with \(y=0\) or
\(v_w(y)\ge e_wS_{u,w}\).
\end{theorem}
\begin{proof}
By the localized inclusion choose \(\rho\in\mathcal O\setminus\mathfrak p_w\)
with \(\rho\phi\in M_u^{\mathrm{ram}}\). Apply \ref{B4.23}. The complex input for \(\rho\phi\) is \(\iota(\rho)\) times that for \(\phi\), so division by \(\rho\) gives a representative for \(\phi\). Its valuation is unchanged since \(v_w(\rho)=0\).
\end{proof}

\section{All remaining finite places}

Use the filling and lattice comparison data of \ref{B4.07}--\ref{B4.14}. Write \(G_{p,j}\) for the gain in \eqref{eq:p4c2-served-gain}.

\begin{lemma}[Comparison of the two gains]\label{B4.25}
If the filling identities of \ref{B3.52} hold at every finite place and the discarded set \(\mathcal D\) of \ref{B4.04} contains every inert-class prime \(p\) with \(h_p\) of odd order greater than one,
then \(\gainD_j(w)\le G_{p_w,j}\) for every \(w\) and \(1\le j\le m\).
\end{lemma}
\begin{proof}[Outline]
Nonnegativity \(G_{p,j}\ge 0\) from \ref{B3.51} covers all places where the screened or raw gain vanishes. At remaining inert and signal places outside \(\mathcal{D}\), the filling identity of \ref{B4.04} and excluding odd-order inert primes match both gains to the counted gain. The complete proof is in \S\ref{prf:B4.25}.
\end{proof}

Substituting the loss formula of \ref{B4.03} into \eqref{eq:p4c2-budgets} at normalization \(1\) gives the base and enhanced column exponents
\[
 s^{\rm base}_{u,w}=B_{u,w}+\delta_{u,w},\qquad
 s^{\rm enh}_{u,w}=\delta_{u,w}+
 \sum_{j=1}^m(G_{p_w,j}+F_j(w)-E_{u,j}(w)).
\]

\begin{theorem}[The enhanced period bound away from \(T_u,g_0,g_1\)]\label{B4.26}
Assume the point and period hypotheses of \ref{B4.19}, the local error conditions of \ref{B4.03} and the three compatibility conditions of \ref{B3.54k} on the local fillings, the ramified space and the base terms. Assume the enhanced column estimate \eqref{eq:p4c2-enhanced-column} at normalization \(1\), the localized inclusion of \ref{B4.21a} away from \(T_u\), and the gain comparison of \ref{B4.25}. Assume that every \(w\) with
\(W_{p_w}\ne\varnothing\) and \(p_w\ne T_u\) satisfies the
model conditions of \ref{B4.07}. 
For every such place assume the selected-place description of \ref{B4.21a},
\[
 \phi\in(N_u)_{(w)}
 \quad\Longleftrightarrow\quad
 c_{u,w}(j_u\phi)\in N_{P,u,w}
 \qquad(\phi\in V_u).
\]
Here \(c_{u,w}\) is the comparison map and \(N_{P,u,w}\) the enhanced lattice of the permitted model in \ref{B4.07}.
If \(p_w\notin\{T_u,g_0,g_1\}\), then for \(\phi\in N_u\)
and admissible \(H\) at \(\mu_w\) there is \(y\in K\) representing
\(P_u^{\mu_w}(\phi,H)\) with \(y=0\) or \(v_w(y)\ge e_wS_{u,w}\).
\end{theorem}
\begin{proof}[Outline]
Use the local lattice conditions of \ref{B4.21a} to apply the enhanced column estimate, then bound its exponent below by \(S_{u,w}+\delta_{u,w}\) using \ref{B4.25}. With the denominator of \ref{B4.03}, \ref{B4.19} subtracts \(\delta_{u,w}\) in passing to the period. The complete proof is in \S\ref{prf:B4.26}.
\end{proof}

\begin{theorem}[The complementary bound on the original lattice]\label{B4.27}
Assume the point and period hypotheses of \ref{B4.19}, the base-term formula of \ref{B4.02}, the local error conditions of \ref{B4.03} and the ordinary column estimate \eqref{eq:p4c2-ordinary-column} at normalization \(1\).
If \(p_w\notin\{T_u,g_0,g_1\}\), then for
\(\phi\in M_u^{\mathrm{ram}}\) and admissible \(H\) at \(\mu_w\)
there is \(y\in K\) representing \(P_u^{\mu_w}(\phi,H)\) with
\(y=0\) or \(v_w(y)\ge e_wB_{u,w}\).
\end{theorem}
\begin{proof}
The input image lies in \(S^{\mathrm{ram}}_{u,w}\), so \eqref{eq:p4c2-ordinary-column} gives exponent \(s^{\rm base}_{u,w}=B_{u,w}+\delta_{u,w}\). With the common stabilizer denominator of \ref{B4.03}, Theorem~\ref{B4.19} gives exponent \(B_{u,w}\) for the period.
\end{proof}

\section{The all-place conclusion}

\begin{lemma}[Rounding the valuation budget]\label{B4.28}
Suppose that for every word \(u\), every finite place \(w\),
every \(\phi\in N_u\), and every admissible \(H\) at \(\mu_w\),
there is a representative \(y\in K\) of \(P_u^{\mu_w}(\phi,H)\)
with \(y=0\) or \(v_w(y)\ge e_wS_{u,w}\).
The same representative satisfies \(y=0\) or
\(v_w(y)\ge\lfloor e_wS_{u,w}\rfloor\). 
\end{lemma}
\begin{proof}
The floor is at most \(e_wS_{u,w}\).
\end{proof}

The complementary period bound requires that for every \(u,w\) with \(p_w\ne T_u\), every \(\phi\in M_u^{\mathrm{ram}}\) and every admissible \(H\) at \(\mu_w\), there be a representative \(y\in K\) of \(P_u^{\mu_w}(\phi,H)\) with \(y=0\) or \(v_w(y)\ge e_wB_{u,w}\).

\begin{theorem}[Divisibility for fixed arithmetic data]\label{B4.29}
Fix \(D,L\) and the local data for the size and column estimates in \ref{B4.16}.
Assume the point and period hypotheses of \ref{B4.18}, the local error conditions of \ref{B4.03} and the compatibility conditions of \ref{B3.54k}. For some \(C\in\R\) assume the size estimates of \ref{B4.15} with constant \(C\) and both column estimates of \ref{B4.14} at normalization \(1\).
Require the level-prime hypotheses of \ref{B4.22}, the two fixed-prime inequalities with one \(C_1\in\R\), and the scalar homogeneity of the column period in \ref{B4.24}.
Assume the localized inclusion away from \(T_u\) in \ref{B4.21a}, the permitted-model and selected-lattice hypotheses of \ref{B4.26}, and the filling and screening hypotheses of \ref{B4.25}.
Then the rounded bound of \ref{B4.28} holds on \(N_u\) at every finite place, and the complementary period bound holds on \(M_u^{\mathrm{ram}}\) away from \(T_u\).
\end{theorem}
\begin{proof}
For \(N_u\), apply \ref{B4.22} above \(T_u\) and \ref{B4.24} above the fixed primes distinct from \(T_u\). Apply \ref{B4.26} elsewhere. Lemma~\ref{B4.28} gives the rounded bound at all finite places.
For \(M_u^{\mathrm{ram}}\), apply \ref{B4.23} at fixed
primes away from \(T_u\). At the remaining places away from \(T_u\), the base-term identity in \ref{B3.54k} and the ordinary column estimate in \ref{B4.14} allow us to apply \ref{B4.27}.
\end{proof}

\section{Uniformity in the selected family}

For an arithmetic family \(\mathfrak F\) of parameter \(\varepsilon\), an assertion holds eventually if there is \(x_0\in\R\) such that it holds for every \(x\ge x_0\) with \(x>x_1\) and every \(i\in I_x\).
A family meets the thresholds \((\mathcal K,C_H)\) when its constants satisfy the staged lower bounds of Definition~\ref{B8.39}. Use the closedness quantifiers of that definition.

\begin{theorem}[A common threshold for the family]\label{B4.30}
Let \(\mathfrak F\) be an arithmetic family of parameter \(\varepsilon>0\).
Assume that each of the reference and twisted point hypotheses and each of the reference and twisted period hypotheses of \ref{B4.18} holds eventually. Suppose that for one \(C\ge0\) the size estimates of \ref{B4.15} with constant \(C\) and both column estimates of \ref{B4.14} at normalization \(1\) hold eventually. Assume also that both sharp level-prime column bounds of \ref{B4.22a} hold eventually for all words, and that eventually \(T_u\nmid\Delta_u\) for all words. Assume that for this
family there is one \(C_1\in\R\) for which both fixed-prime
hypotheses hold eventually. Then the rounded and complementary period bounds of \ref{B4.29} hold eventually with a common cutoff.

For closedness on a class \(\mathcal C\), count the two sharp
column bounds as one input and the existence of \(C_1\) for
both eventual fixed-prime bounds as one input. If these eight
inputs are closed, the conjunction of the two eventual conclusions
is closed. 
\end{theorem}
\begin{proof}[Outline]
Take the largest cutoff and verify \ref{B4.07} using \ref{B3.51}. Use \ref{B4.21a}, \ref{B4.02}–\ref{B4.04}, \ref{B4.14}, and \ref{B4.20} to apply \ref{B4.29}, with closedness by the finite conjunction rule of \ref{B8.39}. The complete proof is in \S\ref{prf:B4.30}.
\end{proof}

\begin{theorem}[The selected family has divisible periods]\label{B4.31}
For each family in \(\mathcal A\), fix a cutoff and, at every member beyond it with \(x>x_1\), a compatible twist realization as in \ref{B8.02}, including its rows, comparison maps, self-dual host and Weil equivariance. Call this choice the twist selection. Let \(\mathcal D_3=\mathcal D(\mathcal A_{\rm ch},3)\), where \(\mathcal A_{\rm ch}\) is the character subclass in Definition~\ref{def:p4c3-design}.
Assume the eventual point and period hypotheses and both
sharp-column hypotheses of \ref{B4.30} are closed on
\(\mathcal D_3\). Assume the existence of a common \(C_1\in\R\)
for the two eventual fixed-prime hypotheses is closed there.
On \(\mathcal D_3\) assume both closed hypotheses of the selection form of \ref{B4.16}: eventual existence of compatible twisted chains with the local system and complementary bounds, and the orbit and jet conditions there stated for the kernel of each selected member beyond its selection cutoff.

Then there is \(\varepsilon_1>0\) such that for every
\(0<\varepsilon\le\varepsilon_1\) there are staged lower thresholds \(\mathcal K\) with the following property. For every arithmetic family \(\mathfrak F\in\mathcal D_3\) of parameter \(\varepsilon\) and every \(C_H\in\R\) for which the family meets \((\mathcal K,C_H)\),
eventually for each word \(u\), each finite place \(w\),
each \(\phi\in N_u\) and each admissible \(H\) at \(\mu_w\),
there is \(y\in K\) with
\[
 \iota(y)=P_u^{\mu_w}(\phi,H),\qquad
 y=0\ \text{or}\
 v_w(y)\ge\left\lfloor e_w\sum_{j=1}^{m}
 \bigl(\gainD_j(w)+F_j(w)-E_{u,j}(w)\bigr)\right\rfloor .
\]
All period data are those of the member \(D_{x,i},L_{x,i}\).
\end{theorem}
\begin{proof}[Outline]
Use \ref{def:p4c3-design} and \ref{B4.16} for closedness of \ref{B4.15} and \ref{B4.14}, and \ref{B4.21} for coprimality. Apply \ref{B4.30} and take the closedness quantifiers in order, using the constructed rows of \ref{B8.09} and \ref{B8.06}. The complete proof is in \S\ref{prf:B4.31}.
\end{proof}


% ---- p4c4 ----
\chapter{Gluing row lattices across Hecke characters}\label{ch:hecke-row-gluing}

\section{Characters and the participating family}

\begin{definition}[Allowed finite-order characters]\label{B4.32}
Fix a number-field extension $E/F$ and a coefficient field $K$. Let $S\subseteq\N$. With $\mathbb A_{F,f}$ the finite adele ring of $F$, put
\[
 \mathcal A=E\otimes_F\mathbb A_{F,f},\qquad
 \mathcal C=\operatorname{Hom}(\mathcal A^\times,K^\times).
\]
Let $R\subseteq\mathcal A$ be the subring generated by the images of $\mathcal O_E$ and $\prod_{v<\infty}\mathcal O_{F_v}$. For an ideal $J\subseteq R$, let $U(J)$ be the image in $\mathcal A^\times$ of $\ker(R^\times\to(R/J)^\times)$.

Let $v$ range over the nonzero prime ideals of $\mathcal O_F$. Write $Nv=\#(\mathcal O_F/v)$ and let $p(v)$ be the residue characteristic. Set
\[
 U_v=\{a\in\mathcal A^\times:a_{v'}=1\text{ for every }v'\ne v\},
 \qquad I_v=R^\times\cap U_v,\qquad
 W_v=U(\sqrt{p(v)R})\cap U_v.
\]
Here $a_{v'}$ denotes the projection to $(E\otimes_F F_{v'})^\times$. We identify $R^\times$ with its image in $\mathcal A^\times$.

Call $v$ good for $S$ if there are distinct primes $\mathfrak P,\mathfrak P'$ of $\mathcal O_E$ such that
\[
 v\mathcal O_E=\mathfrak P\mathfrak P',\qquad
 Nv\text{ is prime},\qquad Nv\notin v^2,\qquad p(v)\notin S.
\]
In the third condition, $Nv$ is viewed as an element of $\mathcal O_F$.

Let $\mathcal C_{\mathrm{fin}}$ consist of the finite-order characters $\tau\in\mathcal C$ such that $\tau(e\otimes1)=1$ for every $e\in E^\times$ and $U(NR)\subseteq\ker\tau$ for some positive integer $N$. Define the subgroup $\mathcal T(E/F,K,S)$ of $\mathcal C_{\mathrm{fin}}$ by requiring triviality on the base ideles $1\otimes a$, $a\in\mathbb A_{F,f}^\times$, and on every $W_v$. Require also trivial restriction to $I_v$ whenever $v$ fails the good-place condition. The reference character is $1$.
\end{definition}

\begin{definition}[The family used by the division construction]\label{B4.33}
Let $K$ be a number field and $TU\subseteq\mathcal T(E/F,K,S)$ a finite subset. The row and column types form $\mathcal T_U=\{1\}\cup TU$. The division kernel datum gives a bijection $e:T\to\mathcal T_U$ from its index set $T$. Put $t_0=e^{-1}(1)$ and write $\tau_t$ for the underlying character of $e(t)$.

For the arithmetic family, let $U$ be the chosen set of labels and take $TU=\{\tau_a:a\in U\}$, where $\tau_a$ is the finite character of the label $a$. The reference label has finite character $1$, and the map $a\mapsto\tau_a$ is injective. The reference label therefore supplies the unique label for $t_0$. For each $t\ne t_0$, the character $\tau_t$ belongs to $TU$ and has a unique preimage in $U$. Write $\theta(t)$ for the resulting label and $\sigma:K\hookrightarrow\C$ for the arithmetic embedding. Character compatibility gives
\[
 \chi_{\theta(t)}(a)=\sigma(\tau_t(a_f))
\]
for every idele $a$ of $E$, with finite component $a_f\in\mathcal A^\times$.
\end{definition}

\section{The common row lattice}

Put $\mathcal O=\mathcal O_K$. For a finite place $w$ of $K$, write $\mathfrak p_w$ for its prime ideal. Localize each $\mathcal O$-submodule $L$ inside its ambient $K$-space by
\[
 L_{(w)}=\mathcal O_{\mathfrak p_w}L
 =\{x:ax\in L\text{ for some }a\in\mathcal O\setminus\mathfrak p_w\}.
\]
A full lattice is a finitely generated $\mathcal O$-submodule whose $K$-span is the ambient space.

\begin{definition}[Local choice and global intersection]\label{B4.34}
Fix $r\in\N$. For each $t\in T$, take $n_t,d_t\in\N$ and an $\mathcal O$-submodule $M_t\subseteq K^{n_t}$. Choose $K$-linear maps
\[
 \pi_t:K^{n_t}\longrightarrow K^{d_t},\qquad
 \iota_t:K^{d_t}\longrightarrow K^r.
\]
Put $\Lambda_t=\pi_t(M_t)$ and $L_t=\iota_t(\Lambda_t)$. Choose $t_w\in T$ at each finite place and put
\[
 RR=\bigcap_{w<\infty}(L_{t_w})_{(w)}\subseteq K^r.
\]
\end{definition}

\begin{theorem}[Gluing finitely many local changes]\label{B4.35}
Fix $r\in\N$. Let $L_0$ and $L(w)$ be full lattices in $K^r$, with $L(w)$ given at each finite place $w$. Suppose $(L(w))_{(w)}=(L_0)_{(w)}$ outside a finite set. Then $G=\bigcap_w(L(w))_{(w)}$ is a full lattice and $G_{(w)}=(L(w))_{(w)}$ at every finite place.
\end{theorem}

This is the existence and local recovery part of Voight's lattice correspondence specialized to $\mathcal O$ and $K^r$ \cite[Theorem 9.4.9]{voight-2021}. The inverse in that correspondence takes the intersection of the local lattices \cite[Lemma 9.4.6]{voight-2021}.

\begin{proposition}[Recovery of the chosen row lattice]\label{B4.36}
In Definition \ref{B4.34}, suppose each $\Lambda_t$ is full and each $\iota_t$ is surjective. Suppose $\{w:t_w\ne t_0\}$ is finite. Then $RR$ is a full lattice and
\[
 RR_{(w)}=(\iota_{t_w}(\Lambda_{t_w}))_{(w)}
 \qquad(w<\infty).
\]
\end{proposition}
\begin{proof}
The images of generators of $\Lambda_t$ generate $L_t$, and $KL_t=\iota_t(K\Lambda_t)=K^r$. Apply \ref{B4.35} with $L_0=L_{t_0}$ and $L(w)=L_{t_w}$.
\end{proof}

\begin{lemma}[The type at conjugate places]\label{B4.36a}
Use the family of \ref{B4.33}. Let $\mathrm{Tw}$ be the image in the finite places of $K$ of a finite conjugation-stable set $W$ of arithmetic places. Write $\bar w$ for the conjugate of $w$ under the arithmetic involution of $K$ and $p_w$ for the residue characteristic of $w$. Let $\delta,\delta'$ be the auxiliary Hecke characters, each of order $15$ on the local units at every place of $E$ above its respective rational prime $p_\delta,p_{\delta'}$. For $w\in\mathrm{Tw}$ take Hecke characters $\psi_w$ unramified above both primes with $\psi_{\bar w}=\psi_w$. Put $\xi_w=\delta'\psi_w$ if $p_w=p_\delta$ and $\xi_w=\delta\psi_w$ otherwise. Assume the full arithmetic column family contains $\xi_w$ for every $w\in\mathrm{Tw}$.

Use the selector $w\mapsto t_w$ from the proof in both coordinate systems. In each system suppose every $\Lambda_t$ is full and every $\iota_t$ is an isomorphism. Then
\[
 t_{\bar w}=t_w,\qquad w\in\mathrm{Tw}\ \Longleftrightarrow\ t_w\ne t_0.
\]
In either system, fix $w$ and put $t=t_w$. Then $RR_{(z)}=(L_t)_{(z)}$ for $z=w,\bar w$.
\end{lemma}
\begin{proof}[Outline]
Choose \(t_w\) by \(\xi_w\) via \ref{B4.33} on \(\mathrm{Tw}\) and take \(t_0\) elsewhere. Conjugation invariance, injectivity and order \(15\) give \(t_{\bar w}=t_w\ne t_0\) on \(\mathrm{Tw}\), and \ref{B4.36} recovers both localizations of \(L_{t_w}\). The complete proof is in \S\ref{prf:B4.36a}.
\end{proof}

\section{Passing the period bound to coordinate rows}

Let $V_\mu$ be the complex row space for the arithmetic label $\mu$. Put $V_0(\mu)=\operatorname{im}(e_{\mathrm{op},\mu})$ for the idempotent row operator $e_{\mathrm{op},\mu}$. Let $A\subseteq\mathcal O$ be the image of the arithmetic coefficient ring. The test set $\widetilde M_\mu$ is the $\sigma(A)$-span of a finite set $\mathcal G_\mu$.

\begin{lemma}[Integral decomposition of test vectors]\label{B4.37}
Fix $\mu$ and suppose
\[
 \sigma(a)y\in\widetilde M_\mu
 \qquad(a\in\mathcal O,\ y\in\widetilde M_\mu).
\]
For every $x\in V_0(\mu)$ admitting an expression
\[
 x=\sum_{v\in\mathcal G_\mu}\sigma(c_v)v,\qquad c_v\in\mathcal O,
\]
there are a finite set $S_x\subseteq\widetilde M_\mu\cap V_0(\mu)$ and coefficients $b_y\in\mathcal O$ such that
\[
 x=\sum_{y\in S_x}\sigma(b_y)y.
\]

\end{lemma}
\begin{proof}
Scalar stability and closure under sums give $x\in\widetilde M_\mu$. Take $S_x=\{x\}$ and $b_x=1$.
\end{proof}

For a finite place $w$ and $s\in\R$, put
\[
 \mathcal V_{w,s}=\{0\}\cup\{a\in K^\times:v_w(a)\ge s\}.
\]
Here $v_w(a)$ is the exponent of $\mathfrak p_w$ in $a\mathcal O$. For $a=b/c\in K^\times$ with $b,c\in\mathcal O$, both $v_w$ and the additive valuation used in \ref{B4.28} give $\operatorname{ord}_{\mathfrak p_w}(b)-\operatorname{ord}_{\mathfrak p_w}(c)$. The set $\mathcal V_{w,s}$ is an $\mathcal O$-submodule. It is the fractional ideal $(\mathfrak p_w\mathcal O_{\mathfrak p_w})^{\lceil s\rceil}$ of the discrete valuation ring $\mathcal O_{\mathfrak p_w}$.

\begin{proposition}[Period bounds on the coordinate row lattice]\label{B4.38}
Fix a column $u$. Let $N_u\subseteq V_u$ be the input lattice of Proposition~\ref{B4.21a} in the rational ramified input space of $u$, and let $E_u:K^{a_u}\to V_u$ be the column bridge. For each $t\in T$, let $J_t:K^{n_t}\to V_{\theta(t)}$ be the row bridge.
Write $P_{u,t}(\phi,H)\in K$ for the coordinate bilinear pairing. As in Chapter~\ref{ch:columns-to-periods}, $P_u^\mu(v,y)$ denotes the complex period obtained by pairing the image of $v\in V_u$ under the complex embedding $j_N$ with the row $y\in V_\mu$ at label $\mu$. Assume that for every $t\in T$ the identity
\[
 \sigma(P_{u,t}(\phi,H))=P_u^{\theta(t)}(E_u\phi,J_tH)
\]
holds for all $\phi\in K^{a_u}$ and $H\in K^{n_t}$. Assume also that for every $t\in T$ and $H\in M_t$,
\[
 J_tH\in V_0(\theta(t)),\qquad
 J_tH=\sum_{v\in\mathcal G_{\theta(t)}}\sigma(c_v)v
 \quad\text{for some }c_v\in\mathcal O.
\]
Fix a type $t$, a finite place $w$, a threshold $s\in\R$ and a subset $N\subseteq K^{a_u}$. Suppose the integral decomposition in the conclusion of Lemma~\ref{B4.37} holds for $\theta(t)$ and, for every $x\in N$ and $y\in\widetilde M_{\theta(t)}\cap V_0(\theta(t))$, there is $a\in\mathcal V_{w,s}$ with
\[
 \sigma(a)=P_u^{\theta(t)}(E_ux,y).
\]
Then
\[
 P_{u,t}(\phi,H)\in\mathcal V_{w,s}
 \qquad
 (\phi\in\operatorname{span}_{\mathcal O}N,\ H\in M_t).
\]

Let $\mu_0$ be the reference label and $\mu_\star,\mu'_\star$ the two auxiliary labels. Let $\mu_w$ be the arithmetic label used in the period bound of \ref{B4.29}. Assume the column bridge satisfies $\theta(t_w)=\mu_w$ at every finite place $w$. Suppose the integral decomposition in the conclusion of Lemma~\ref{B4.37} holds for every $\mu\in\{\mu_0,\mu_\star,\mu'_\star\}\cup\{\mu_w:w<\infty\}$.
If $\mathcal N_u\subseteq K^{a_u}$ is an $\mathcal O$-submodule with $E_u(\mathcal N_u)\subseteq N_u$ and the periods on $N_u$ satisfy the rounded valuation bound of Theorem~\ref{B4.29} at every finite place, then
\[
 P_{u,t_w}(\phi,H)=0
 \quad\text{or}\quad
 v_w(P_{u,t_w}(\phi,H))\ge s_{u,w},
 \qquad
 s_{u,w}=\lfloor e_w\sg(u,w)\rfloor
\]
for every finite $w$ and for all $\phi\in\mathcal N_u$ and $H\in M_{t_w}$. Here $e_w$ is the ramification index of $w$ over its rational residue characteristic and $\sg(u,w)$ is the budget $S_{u,w}$ of \ref{B4.28}.
\end{proposition}
\begin{proof}[Outline]
The integral decomposition and bridge identity express the period as an \(\mathcal O\)-linear combination of representatives in \(\mathcal V_{w,s}\). Extend by column linearity to \(\operatorname{span}_{\mathcal O}N\) and use \ref{B4.29} for the local representatives. The complete proof is in \S\ref{prf:B4.38}.
\end{proof}


% ---- p4c5 ----
\chapter{The pairing after division}\label{ch:pairing-after-division}

Let $K$ be the coefficient number field and put $\mathcal O=\mathcal O_K$.
For a finite place $w$, write $\mathfrak p_w$ for its prime ideal and
$v_w:K^\times\to\Z$ for the normalized additive valuation. For an
$\mathcal O$-submodule $M$ of a $K$-vector space, use the localization
\[
 M_{(w)}=\{x:\text{there is }t\in\mathcal O\setminus\mathfrak p_w
                    \text{ with }tx\in M\}.
\]
A lattice means a finitely generated $\mathcal O$-submodule spanning its
ambient $K$-vector space. For $R\subset W$ put
$R^\vee=\{\lambda\in W^*:\lambda(R)\subset\mathcal O\}$.

\section{The divided column module}

Fix the arithmetic data $D$ and the size data $L$.
For each word $u$ and finite place $w$, let
\[
 S_{u,w}=\sg(u,w),\qquad
 s_{u,w}=\lfloor e_w S_{u,w}\rfloor ,
 \qquad e_w=\text{the ramification index of }w\text{ over }\Q.
\]
Definition~\ref{B4.05} gives the budget $S_{u,w}$, and
Lemma~\ref{B4.28} passes from the period bound with threshold $e_wS_{u,w}$
to the bound with integer threshold $s_{u,w}$.
By Lemma~\ref{B4.05a}, both quantities vanish outside the finite set $S_L$ of Definition~\ref{B4.05}.
For each word $u$, let $V_u$ be its column space over $K$ and $N_u\subset V_u$ the input $\mathcal O$-submodule specified by $L$.

\begin{definition}[Division ideal and divided submodule]\label{B4.39}
For each word $u$, define
\[
 I_u=\prod_{w\in S_L}\mathfrak p_w^{-s_{u,w}},
 \qquad
 I_uN_u=\operatorname{span}_{\mathcal O}
       \{a\phi:a\in I_u,\ \phi\in N_u\}.
\]
The division condition for $(D,L)$ is the equality
$N'_u=I_uN_u$ for every $u$, where $N'_u\subset V_u$ is the $\mathcal O$-submodule specified by the division data of $D$ for the word $u$.
\end{definition}

At a place with $s_{u,w}>0$, this operation permits additional
denominators of order $s_{u,w}$. A negative $s_{u,w}$ imposes additional
divisibility.

\begin{proposition}[Realization of local shifts]\label{B4.40}
Let $N\subset K^r$ be a lattice, with $r\in\N$.
Fix a finitely supported integer function $w\mapsto\tau_w$. Put
$I=\prod_w\mathfrak p_w^{-\tau_w}$. Then $IN$ is a lattice. For every
$w$ there exists $a_w\in K^\times$ with $v_w(a_w)=-\tau_w$, and for
every such choice
\[
                  (IN)_{(w)}=(a_wN)_{(w)}.
\]
Moreover, any $\mathcal O$-submodule $N'\subset K^r$ with these
localizations equals $IN$.

More generally, if $J$ is any nonzero fractional ideal and
$a\in K^\times$ satisfies $v_w(a)=v_w(J)$, where $v_w(J)$ is the
exponent of $\mathfrak p_w$ in $J$, then $(JN)_{(w)}=(aN)_{(w)}$
for every $\mathcal O$-submodule $N\subset K^r$.
\end{proposition}

Local principality follows from the discrete valuation ring criterion
\cite[Lemma~10.120.17, Tag~034X]{stacks-2026}.
Voight gives the local criterion and reconstruction for lattices over a
Dedekind domain \cite[Corollary~9.4.7 and Theorem~9.4.9]{voight-2021}.
The archive specializes reconstruction to prescribed shifts over
$\mathcal O_K$ and proves uniqueness among all submodules. For uniqueness,
apply \cite[Lemma~10.23.1, Tag~00HN]{stacks-2026} to $(N'+IN)/N'$ and
$(N'+IN)/IN$.

For arbitrary $\mathcal O$-submodules $N,N'$ and any integer function
$\tau$, call $N'$ a \emph{shift of $N$ by $\tau$} if for every $w$
there is $a_w\in K^\times$ of valuation $-\tau_w$ such that
$N'_{(w)}=(a_wN)_{(w)}$.

\begin{lemma}[The arithmetic and coordinate shifts agree]\label{B4.41}
Fix $u$ and a $K$-linear isomorphism $e:K^r\to V_u$.
Let $N=e^{-1}(N_u)$ and let $N^c\subset K^r$ be an $\mathcal O$-submodule
with $e(N^c)=N'_u$.
Assume the division condition of Definition~\ref{B4.39}.
Let $g_u$ be the real-valued function on finite places given by $g_u(w)=s_{u,w}$.
For any real-valued function $g$ on the finite places, put $d_w=\lfloor g(w)\rfloor$. When $d$ has finite support, write $N(g)=(\prod_{d_w\ne0}\mathfrak p_w^{-d_w})N$.
For $g=g_u$ one has $d_w=s_{u,w}$ at every $w$, and
\[
 N^c=N(g_u)=I_uN.
\]
In particular $N^c$ is the shift of $N$ by $w\mapsto s_{u,w}$.
\end{lemma}

\begin{proof}
Since $s_{u,w}$ is an integer, the $w$-factor of the shift ideal is
$\mathfrak p_w^{-s_{u,w}}$. Their product is $I_u$ by Definition~\ref{B4.39}.
By $K$-linearity, $e(I_uN)=I_ue(N)=I_uN_u=N'_u=e(N^c)$.
Since $e$ is injective, $I_uN=N^c$. At $w$, a power of a local uniformizer gives
$a_w\in K^\times$ with $v_w(a_w)=-s_{u,w}=v_w(I_u)$.
Proposition~\ref{B4.40} then gives $N^c_{(w)}=(a_wN)_{(w)}$.
\end{proof}

\section{Cancellation and local integrality}

\begin{lemma}[Cancellation on a shifted column]\label{B4.42}
Fix a finite place $w$ and $s\in\Z$. Let $N,N'\subset K^r$, $r\in\N$,
be $\mathcal O$-submodules with
$N'_{(w)}=(aN)_{(w)}$ for some $a\in K^\times$ of valuation $-s$.
Let $f:K^r\to K$ be $K$-linear. If
\[
 \forall\psi\in N,\qquad f(\psi)=0
                     \quad\text{or}\quad v_w(f(\psi))\ge s,
\]
then for every $\phi\in N'$ one has
$f(\phi)=0$ or $v_w(f(\phi))\ge0$.
\end{lemma}

\begin{proof}
For $\phi\in N'$, the equality of localizations gives
$t\in\mathcal O\setminus\mathfrak p_w$ and $\psi\in N$ such that
$t\phi=a\psi$. Hence
\[
                    t f(\phi)=a f(\psi),\qquad v_w(t)=0.
\]
If $f(\psi)=0$, then $f(\phi)=0$. If $f(\phi)\ne0$, both sides are
nonzero and
\[
 v_w(f(\phi))=v_w(a)+v_w(f(\psi))-v_w(t)
             =-s+v_w(f(\psi))\ge0.
\]
\end{proof}

\begin{lemma}[From a local row pairing to global integers]\label{B4.43}
Fix $w$, integers $n,h\ge0$, an $\mathcal O$-submodule $M\subset K^n$,
and $K$-linear maps $\pi:K^n\to K^h$, $\iota:K^h\to K^h$.
Let $b:K^h\to K$ and $P:K^n\to K$ be $K$-linear.
Fix $\zeta\in K^\times$ with $v_w(\zeta)=0$. Suppose
$P(F)=\zeta\,b(\iota\pi F)$ for every $F\in K^n$.
For any $c\in\Z$, if $P(F)=0$ or $v_w(P(F))\ge c$ for every $F\in M$,
then
\[
 x\in(\iota\pi M)_{(w)}
 \quad\Longrightarrow\quad
 b(x)=0\ \text{or}\ v_w(b(x))\ge c.
\]

For the global assertion, fix $h$, a $K$-linear functional $b:K^h\to K$,
and an $\mathcal O$-submodule $R\subset K^h$. Suppose the local hypotheses
hold at every $w$ with this $b$, source dimension $n_w$, and $c=0$.
If $R_{(w)}=(\iota_w\pi_wM_w)_{(w)}$ for every $w$, then
$b(R)\subset\mathcal O$.
\end{lemma}

\begin{proof}[Outline]
Writing \(tx=\iota\pi F\) with \(t\in\mathcal{O}\setminus\mathfrak{p}_w,F\in M\) gives \(P(F)=\zeta tb(x)\), so \(v_w(\zeta)=v_w(t)=0\) gives \(v_w(b(x))\ge c\) if \(b(x)\ne 0\). Globally, \(v_w(b(x))\ge 0\) everywhere makes \((b(x))\) integral, so \(b(x)\in\mathcal{O}\). The complete proof is in \S\ref{prf:B4.43}.
\end{proof}

\section{The pairing into the common dual}

\begin{theorem}[Integral pairing after division]\label{B4.44}
Fix $r,h\in\N$ and a $K$-linear map $B:K^r\to(K^h)^*$.
Take $\mathcal O$-submodules $N,N'\subset K^r$ and $R\subset K^h$.
For each finite place $w$, fix $n_w\in\N$ and an $\mathcal O$-submodule
$M_w\subset K^{n_w}$. Let
\[
 \pi_w:K^{n_w}\to K^h,\quad \iota_w:K^h\to K^h,\quad
 P_w:K^r\times K^{n_w}\to K
\]
be $K$-linear maps and a $K$-bilinear pairing, respectively.
Suppose $\zeta_w\in K^\times$ has $v_w(\zeta_w)=0$ and
\[
 P_w(\phi,F)=\zeta_w B(\phi)(\iota_w\pi_wF)
                   \quad(\phi\in K^r,\ F\in K^{n_w}),
 \qquad
 R_{(w)}=(\iota_w\pi_wM_w)_{(w)}.
\]
Let $\sigma_w\in\R$ and $\tau_w\in\Z$ be arbitrary functions of $w$.
Suppose $N'$ is a shift of $N$ by $\tau$, with
$\tau_w\le\lceil\sigma_w\rceil$ at every $w$. Assume at each $w$ that
\[
 \phi\in N,\ F\in M_w
 \quad\Longrightarrow\quad
 P_w(\phi,F)=0\ \text{or}\ v_w(P_w(\phi,F))\ge\sigma_w.
\]
Then $B(N')\subset R^\vee$.
\end{theorem}

\begin{proof}
Fix $w$ and $F\in M_w$. Each nonzero value of
$f_F(\phi)=P_w(\phi,F)$ on $N$ has integer valuation, hence
\[
 v_w(f_F(\phi))\ge\lceil\sigma_w\rceil\ge\tau_w.
\]
Apply Lemma~\ref{B4.42} to $f_F$. Fix any $\phi\in N'$.
The lemma gives $P_w(\phi,F)=0$ or $v_w(P_w(\phi,F))\ge0$
for every $w$ and $F\in M_w$.
Lemma~\ref{B4.43} with $b=B(\phi)$ gives $B(\phi)(R)\subset\mathcal O$.
\end{proof}

\begin{theorem}[The selected characters give one dual lattice]\label{B4.45}
Use the column record and the row data of Proposition~\ref{B4.36},
with selected index $t_w$ and common coordinate row lattice $R^c\subset K^h$.
Put $A_w=\iota_{t_w}\pi_{t_w}$, where $\pi_t$ is the row map of
Definition~\ref{B4.34}. Proposition~\ref{B4.36} gives
$R^c_{(w)}=(A_wM_{t_w})_{(w)}$.
For a column index $i$, let $B^c_i:K^{a_i}\to(K^h)^*$ be its coordinate
column map. The column record gives
\[
 P_{i,t_w}(\phi,F)=\zeta_{i,t_w}B^c_i(\phi)(A_wF),
 \qquad \zeta_{i,t_w}\in K^\times,\quad v_w(\zeta_{i,t_w})=0.
\]
Let $N,N'\subset K^{a_i}$ be $\mathcal O$-submodules and let
$\sigma_w\in\R$, $\tau_w\in\Z$. Suppose $N'$ is a shift of $N$ by
$\tau$, with $\tau_w\le\lceil\sigma_w\rceil$ at every $w$. Assume at each $w$ that
\[
 \phi\in N,\ F\in M_{t_w}
 \quad\Longrightarrow\quad
 P_{i,t_w}(\phi,F)=0\ \text{or}\ v_w(P_{i,t_w}(\phi,F))\ge\sigma_w.
\]
Then $B^c_i(N')\subset (R^c)^\vee$.

For the arithmetic application, assume $\theta(t_w)=\mu_w$ at every
finite place $w$, with the arithmetic labels of Definition~\ref{B4.33}.
For each arithmetic word $u$, let $i(u)$ be the coordinate column index assigned to $u$ by the division data, and put $B^c_u=B^c_{i(u)}$ and $a_u=a_{i(u)}$. Let $B_u:V_u\to W^*$
be the period map and $RR\subset W$ the common row lattice in the
arithmetic row space $W$. Fix $K$-linear isomorphisms $E_u:K^{a_u}\to V_u$
and $E_R:K^h\to W$, and an $\mathcal O$-submodule $N^c_u\subset K^{a_u}$.
Assume
\[
 E_u(N^c_u)=N'_u,\qquad E_R(R^c)=RR,\qquad
 B_u(E_u x)(E_R y)=B^c_u(x)(y)\quad(x\in K^{a_u},\ y\in K^h).
\]
Under the division condition of Definition~\ref{B4.39} and the hypotheses
of Proposition~\ref{B4.38} for $E_u$ and column $i(u)$,
one has $B_u(N'_u)\subset RR^\vee$.
\end{theorem}

\begin{proof}[Outline]
Apply \ref{B4.44} with \(R=R^c\), \(B=B^c_i\), \(M_w=M_{t_w}\), and \(\pi_w=A_w\) for the coordinate claim. Setting \(N=E_u^{-1}(N_u)\), obtain the pairing bound from \ref{B4.38} and the shift \(N^c_u\) from \ref{B4.41}, then transport the inclusion via \(E_u\) and \(E_R\). The complete proof is in \S\ref{prf:B4.45}.
\end{proof}


% ---- p4c6 ----
\chapter{Slope change and the data for step 5}\label{ch:p4c6}

Write \(n_K=[K:\Q]\) and \(\mathcal O=\mathcal O_K\). At a finite place \(w\), write \(\mathfrak p_w\) for its prime ideal and \(p_w\)
for its residue characteristic. Normalize the valuation \(v_w\) by
\(v_w(\mathfrak p_w)=1\). Put \(e_w=v_w(p_w)\) and
\(\ell_w=\log N\mathfrak p_w/n_K\). Use \(S_{u,w}=\sg(u,w)\) from
\ref{B4.05} and \(s_{u,w}=\lfloor e_wS_{u,w}\rfloor,N_u,I_u,N'_u\)
from \ref{B4.38}--\ref{B4.39}. The column modules \(N_u,N'_u\) use the same basis and frames. Write \(S\) for the prescribed finite set of places outside which the local gains, base terms and losses vanish, as in \ref{B4.05}.

\section{The exact change of slope}

Let \(V\) be a finite-dimensional space over a totally complex number
field \(K\), with basis \(b\) and invertible frames \(X_v\) at infinity.
Use one embedding per infinite place. The real realization \(\rho_{b,X}(y)\)
consists of the real and imaginary parts of \(X_v^{-1}v([y]_b)\).
The norm at \(v\) is the Euclidean norm in this frame. For a nonzero
finitely generated \(\mathcal O\)-submodule \(M\subset V\), put
\[
 \mu_{b,X}(M)=2\,\operatorname{slope}(\rho_{b,X}M)-\log 2+
 \frac{\log|\operatorname{disc}K|}{n_K},
\]
where \(\operatorname{slope}\) is minus the log covolume in the real span
divided by the \(\Z\)-rank. Let \(\mu_{\max}(M)\) be the supremum of
\(\mu(L)\) over all nonzero \(\mathcal O\)-submodules \(L\subseteq M\),
with induced metrics. A full lattice spans its ambient \(K\)-space.
Its dual \(M^\vee=\{f\in V^\vee:f(M)\subset\mathcal O\}\) has the dual
metrics.

\begin{lemma}[Multiplication by an integral ideal]\label{B4.46}
For \(M\) as above and a nonzero integral ideal \(\mathfrak b\),
\[
 \mu_{b,X}(\mathfrak b M)=\mu_{b,X}(M)-\frac2{n_K}\log N\mathfrak b.
\]
\end{lemma}

\begin{lemma}[Multiplication by a fractional ideal]\label{B4.47}
Let \(I\) be a nonzero fractional ideal and \(M\subset V\) an
\(\mathcal O\)-submodule. If \(N=IM\) is nonzero and finitely generated,
then so is \(M\), and
\[
 \mu_{b,X}(N)=\mu_{b,X}(M)-\frac2{n_K}\log N(I),
\]
where \(N(I)\in\Q_{>0}\) is the absolute fractional-ideal norm.
\end{lemma}

These are the ideal-line cases of Bost's Hermitian tensor formula
\cite[Proposition 5.5(3)]{fonseca-2019}, with twice the Arakelov slope.
The ideal line has absolute-value metrics and normalized degree
\(-\log N(I)/n_K\). Multiplication identifies the tensor product with
\(IM\), and \(M=I^{-1}N\).

\begin{lemma}[Norm of the division ideal]\label{B4.48}
For every word \(u\), \(I_u\ne0\) and
\[
 \log N(I_u)=-\sum_{w\in S}s_{u,w}\log N\mathfrak p_w .
\]
\end{lemma}

\begin{proof}
The nonzero factors \(\mathfrak p_w^{-s_{u,w}}\) of \ref{B4.39} have
norms \((N\mathfrak p_w)^{-s_{u,w}}\). Multiply and take logarithms.
\end{proof}

\begin{theorem}[Slope difference in the arithmetic data]\label{B4.49}
Assume \(N'_u=I_uN_u\) for every word \(u\), as in \ref{B4.39}. For any word \(u\) with
\(N'_u\) nonzero and finitely generated and with invertible frames,
\[
 \Delta\mu_u:=\mu(N'_u)-\mu(N_u)
 =2\sum_{w\in S}\ell_ws_{u,w}.
\]
\end{theorem}

\begin{proof}
Apply \ref{B4.47} to \(N'_u=I_uN_u\) and substitute \ref{B4.48}.
\end{proof}

For a full lattice \(N\subset K^d\) over a totally complex field put
\[
 \operatorname{lcv}_X(N)=\frac2{n_K}\log\!\left(
 2^{n_Kd/2}|\operatorname{disc}K|^{-d/2}\operatorname{covol}_X(N)\right),
\]
with Euclidean covolume in the frames. Whenever the localizations
\(A_{(w)},A'_{(w)}\) are full lattices in the same local space, let
\(\operatorname{idx}_w(A,A')\) be the valuation of the determinant
of a basis of \(A'_{(w)}\) relative to one of \(A_{(w)}\).
A shift of \(N\) by \(\tau_w\in\Z\) at \(w\) means
\(N'_{(w)}=a_wN_{(w)}\) with \(a_w\in K^\times\) and
\(v_w(a_w)=-\tau_w\).

\begin{theorem}[Slope difference in coordinates]\label{B4.50}
Let \(K\) be a number field with involution and \(N\subset K^d\)
a full lattice, \(d\ge0\). If a submodule \(N'\subset K^d\) is a shift of
\(N\) by \(\tau_w\) at \(w\), then \(\operatorname{idx}_w(N,N')=-d\tau_w\).
For totally complex \(K\) and invertible frames \(X_v\),
\(\mu_X(N)=-\operatorname{lcv}_X(N)/d\) if \(d>0\). If \(N'\) is full
and the shift holds at every finite place, then for \(d>0\) the shifts
have finite support and
\[
 \operatorname{lcv}_X(N')-\operatorname{lcv}_X(N)=-2d\sum_w\ell_w\tau_w,
 \qquad
 \mu_X(N')-\mu_X(N)=2\sum_w\ell_w\tau_w .
\]
For \(d=0\) both normalized log covolumes are zero.
\end{theorem}

\begin{proof}
The local determinant is \(a_w^d\). For full lattices put
\(N_0=N\cap N'\). The covolume ratio is \([N:N_0]/[N':N_0]\).
Its local length difference is \(\operatorname{idx}_w(N,N')\), so the
normalized log covolume difference is
\(2\sum_w\ell_w\operatorname{idx}_w(N,N')\), with finite support.
The real rank \(n_Kd\) gives the slope formula. In dimension zero
the covolume and its normalizing factor are both \(1\).
\end{proof}

\section{Rounding and the rational primes}

Put \(A_u=\sum_{w\in S}2\ell_we_wS_{u,w}\).

\begin{lemma}[Rounding error]\label{B4.51}
Under the hypotheses of \ref{B4.49},
\[
 |\Delta\mu_u-A_u|\le
 2\sum_{\substack{w\in S\\S_{u,w}\ne0}}\ell_w .
\]
\end{lemma}

\begin{proof}
By \ref{B4.49}, the difference sums
\(2\ell_w(\lfloor e_wS_{u,w}\rfloor-e_wS_{u,w})\).
Each summand vanishes for \(S_{u,w}=0\) and has absolute value
at most \(2\ell_w\).
\end{proof}

\begin{lemma}[Total weight above a rational prime]\label{B4.52}
For a number field \(K\) and a rational prime \(p\), let \(W_p\) be the
places above \(p\). With \(f_w=[\mathcal O/\mathfrak p_w:\F_p]\),
\[
 \sum_{w\in W_p}e_w f_w=n_K,\qquad
 \sum_{w\in W_p}e_w\ell_w=\log p .
\]
\end{lemma}

This is the ramification degree formula
\cite[Remark 15.113.6, Tag 09E8]{stacks-2026b} for \(\Z_{(p)}\) in
\(K/\Q\), with \(N\mathfrak p_w=p^{f_w}\) in the weighted form.

\begin{lemma}[Support of the gain and loss in a parameter family]\label{B4.52a}
Fix a parameter family at \(\varepsilon\) as in \ref{B4.15}. Write \(x_1\) for its threshold and \(I_x\) for the set indexing its members at \(x\). Take any \(x>x_1\) and \(i\in I_x\). Use \(X''\)
from \ref{B3.51} and the stabilizer denominator \(M_u\) from \ref{B4.03}.
At each place \(w\) with \(p_w>X''\), the raw gain vanishes for every
\(1\le j\le m\). If also \(p_w\notin\{T_u,g_0,g_1,d\}\) and
\(p_w\nmid M_u\), then \(E_{u,j}(w)=0\).
We call the simultaneous vanishing of the raw gain and the loss for every \(u,j,w\) satisfying these three prime conditions the support condition above \(X''\).
\end{lemma}

\begin{proof}
The raw gain vanishes outside \(S\) and outside the inert and signal classes. There the filling equality of \ref{B3.52} gives the
prime gain of \ref{B3.51}, which vanishes above \(X''\).
In the formula for the loss \(E_{u,j}(w)\) in \ref{B4.03}, the exceptional terms occur only at \(p_w\in\{T_u,g_0,g_1,d\}\). The fixed inert branch and the remaining
inert-or-signal branch require \(p_w\le X''\).
The remaining term is \(v_{p_w}(M_u)=0\).
\end{proof}

\begin{theorem}[Rational-prime bound for the error]\label{B4.53}
Retain the hypotheses of \ref{B4.49}. Assume \(X''\ge0\) and the support condition above \(X''\) of \ref{B4.52a}. If \(p_w>X''\),
\(p_w\notin\{T_u,g_0,g_1,d\}\) and \(p_w\nmid M_u\), then \(S_{u,w}=0\).
Moreover,
\[
 |\Delta\mu_u-A_u|\le
 2\bigl(X''\log4+\log T_u+\log g_0+\log g_1+\log d+\log M_u\bigr).
\]
\end{theorem}

\begin{proof}[Outline]
Outside \(P_u\), the budget in \ref{B4.05} gives \(S_{u,w}=0\). Apply \ref{B4.51} and \ref{B4.52}, then bound the product of small primes by binomial induction and include the exceptional primes and prime divisors of \(M_u\). The complete proof is in \S\ref{prf:B4.53}.
\end{proof}

\begin{definition}[Size bounds for the auxiliary primes]\label{def:p4c6-prime-sizes}
At scale \(x>1\) with \(0<\varepsilon<1/52\) and window parameter \(J\ge1\), put \(\varepsilon_z=\varepsilon/J\), \(\theta=\varepsilon/100\) and \(b=\exp(x^\varepsilon)\). Let \(\chi_d\) be the quadratic character of \(\Q(\sqrt{-d})\), and put \(D_1=\Re(L'(1,\chi_d)/L(1,\chi_d))\). The bounds for the host primes \(q,d\) are
\[
 \begin{gathered}
 x^{\varepsilon_z/6}\le\log q,\qquad
 x^{19\varepsilon_z/100}\le\log p\le x^{\varepsilon/5+\theta}
 \quad(p=q,d),\\
 |D_1|\le x^\varepsilon,\qquad
 1\le\operatorname{rd}(E)\le \tfrac12 Q_S\sqrt{qd},
 \end{gathered}
\]
where \(E\) is the host field and \(Q_S\) is the host modulus of \ref{B3.46}. For the two fixed primes require, for all \(i,j\in\{0,1\}\),
\[
 (kb)^8\le g_i\le 2g_j,\qquad \log g_i\le x^{\theta+\varepsilon}.
\]
\end{definition}

\begin{theorem}[Uniform error in the family parameters]\label{B4.54}
Assume the auxiliary-prime bounds of Definition~\ref{def:p4c6-prime-sizes} and the support condition above \(X''\) of \ref{B4.52a}. For every \(u\), assume \(N'_u=I_uN_u\) as in \ref{B4.39}, with the invertible column frames of the arithmetic datum. Suppose every \(N'_u\) is finitely generated and
has positive \(K\)-rank. Let \(C_T,C_{\rm stab}\in\R\) satisfy
\(\log T_u\le C_TX''\) and \(\log M_u\le C_{\rm stab} X''\) for every \(u\).
Then
\[
 |\Delta\mu_u-A_u|\le 2(\log4+3+C_T+C_{\rm stab})X''
 \qquad\text{for every }u.
\]
\end{theorem}

\begin{proof}
The positive rank of each \(N'_u\) makes it nonzero, as required in \ref{B4.49}. With \(\theta=\varepsilon/100\)
the size assumptions give \(\log g_i\le x^{\theta+\varepsilon}\) for
\(i=0,1\) and \(\log d\le x^{\varepsilon/5+\theta}\).
Since \(0<\varepsilon<1/52\) and \(x>1\), both powers are at most
\(x\le X''=X'xn\), where \(X',n\ge1\). Apply \ref{B4.53}.
\end{proof}

\begin{lemma}[The index of a projected theta lattice at one place]\label{B4.54a}
Work over the local ring at a finite place \(v\), with valuation
\(v(\cdot)\) and index \([B:B']_v=\operatorname{idx}_v(B,B')\).
For a subspace \(V\) put \(B_V=B\cap V\).
Let a finite group \(P\) act by \(\rho\) on an \(n\)-dimensional
\(K\)-space. Let \(U\le P\) have order \(h\) and index \(s\), with
character \(\psi:U\to\mathcal O^\times\). For each \(g\notin U\)
assume some \(u\in U\cap gUg^{-1}\) satisfies
\(\psi(u)\ne\psi(g^{-1}ug)\). Let \(W\) be the \(\psi\)-eigenspace.
For an idempotent \(e\) with image \(V_0\), put \(Z=W\cap V_0\)
and \(d_Z=\dim Z\).

Take full lattices \(\Theta_1\subseteq\Theta\subseteq B\), with \(B,\Theta\)
stable under \(P\). For an integer \(E\ge0\) and \(r\in\R\) assume
\begin{equation}\label{eq:p4c6-theta-cut}
 [B_V:(\Theta_1)_V]_v\le E\dim V+r\qquad\text{for every subspace }V .
\end{equation}
Suppose \(\Theta\subseteq\sigma\Theta_K\) for a full lattice \(\Theta_K\)
and \(\sigma\in K^\times\), and take a lattice \(\Theta_0\) with
\(h(\Theta_K)_W\subseteq\Theta_0\subseteq W\).
Take a module \(\Theta_J\) with \(e\Theta_0\subseteq e\Theta_J\)
and full lattices \(e\Theta_J,X_3,A\) in \(Z\) with \(X_3\subseteq A\)
and
\begin{equation}\label{eq:p4c6-theta-comparison}
 [B_Z:X_3]_v\le [B_Z:e\Theta_J]_v+e_v\xi,\qquad \xi\in\R .
\end{equation}
Then
\[
 [B_Z:A]_v\le d_Z\{E-v(\sigma)+v(h)\}+\frac{r+n v(h)}s+e_v\xi .
\]
\end{lemma}

\begin{proof}[Outline]
Use the subspace index bound and character averaging to bound \(s[B_Z:B'_Z]_v\) by \(r+nv(h)\) via maps with composite \(h\). Scale \(B'_Z\) into \(e\Theta_J\) and compare indices using \(X_3\subseteq A\). The complete proof is in \S\ref{prf:B4.54a}.
\end{proof}

\section{The row and column data}

Fix a totally complex field \(K\) with involution \(a\mapsto\bar a\).
Choose a subfield \(K_0\subseteq K\) such that for every embedding
\(\sigma:K\hookrightarrow\C\) and \(a\in K_0\),
\(\sigma(\bar a)=\overline{\sigma(a)}\). At the chosen infinite place
\(v_0\) require the identity on all of \(K\).
Use common types \(t\) with base type \(t_0\). Let \(\mathcal I\) be the
finite spectral index set and index the layers by \(1\le j\le m\).
On each definite side, reference and selected, let \(M_t\) be the
test lattice, \(P_t\) the surjective quotient projection and
\(\Lambda_t=P_tM_t\). Write \(T_{jt}\) for the layer lattice with
Hermitian matrix \(H_{jt}\) and frames \(X_{jt,v}\).
All these lattices are full. The test matrices
\(H_{V,t}\) have entries in \(K_0\) and are positive definite at every
infinite place. They induce the quotient metrics.
Every \(X_{jt,v}\) is invertible, and
\[
 (X_{jt,v_0}^{-1})^*X_{jt,v_0}^{-1}=v_0(H_{jt}).
\]
At \(t_0\) the quotient type subspaces form a direct sum spanning the
quotient, and the layer type subspaces are in direct sum.

\begin{definition}[Reference rows and selected columns]\label{def:p4c6-core}
Fix the selected pairings \(P_{u,t}\) and column frames \(X_{\mathrm c,v}\)
for each word \(u\) of dimension \(d_u\). Fix a row subspace \(\mathcal R\subseteq\mathcal H\) of a complex
inner-product space \(\mathcal H\). Take complex-linear maps
\(P_u^{\rm per}:\C^{d_u}\to\mathcal H\).
The common real data are \(F_{0,j,w},l_{3,j},l_{6,j},E_{1,j},
h_{\rm r},\lambda_{\rm r},\delta\), with \(F_{0,j,w}\) the reference base
term. Choose the normalized cut-exponent bound \(J_0(t,w,j)\) of
\ref{B5.08x} for every type, finite place and layer.

On each side choose a selector \(w\mapsto t_w\) with \(t_{\bar w}=t_w\),
where \(\bar w\) is the conjugate place. The set
\(\mathrm{Tw}=\{w:t_w\ne t_0\}\) is finite. Write
\(J_{0,j,w}=J_0(t_w,w,j)\) in the index bounds.
The base quotient has positive dimension. Choose a row space \(K^d\)
with frames \(X_{E,v}\), spectral subspaces \(E_\alpha\) and maps
\(\iota_t\) from the definite quotient spaces. For each layer choose
a space \(K^{d_j}\) with frames \(X_{W,j,v}\), subspaces \(W_{j,\alpha}\),
an \(\mathcal O\)-submodule \(M_j\) and maps \(\theta_{jt}\) from the
definite layer spaces. Take real factors \(\kappa_v,\lambda_{jt\alpha v}\)
and real local errors \(e^0_{wj}\). Denote the reference and selected
row dimensions by \(D_R,D_{\rm c}\).

On the reference side \(H_{j,t_0}\) has entries in \(K_0\) and induces
the definite layer metrics at every infinite place. The type subspaces
at \(t_0\) span each layer. All reference row and layer frames are invertible. Each \(M_j\) is full and
the maps \(\iota_t,\theta_{jt}\) are isomorphisms. At \(t_0\) these maps send the definite type
subspaces onto \(E_\alpha,W_{j,\alpha}\), respectively, with
\(\dim W_{j,\alpha}=\dim E_\alpha\).
The factors \(\kappa_v,\lambda_{jt\alpha v}\) are positive. For \(t=t_0\)
or \(v=v_0\), the pullbacks of the row and layer Hermitian forms under
\(\iota_t,\theta_{jt}\) equal the definite forms times
\(\kappa_v\prod_j\lambda_{jt\alpha v}\) and \(\lambda_{jt\alpha v}\),
respectively, on pairs whose first argument lies in the preimage
of the indicated spectral subspace.

Require \(\operatorname{lcv}_{X_{W,j}}(M_j)/D_R\le l_{3,j}\)
in each reference layer. With \(A_{jt}=\theta_{jt}(T_{jt})\) and
\(e'_{jw}=e_we^0_{wj}\), impose the three inequalities of \ref{B8.38b}
on the reference data. Require nonnegativity and \eqref{eq:B8-paired-index} at each place in
\(\mathrm{Tw}\). Impose \eqref{eq:B8-paired-error} in each layer.
Under the hypotheses of \ref{B8.40a}, the window bound gives constants \(C_N',x_0\) such that, for every family member at \(x\ge x_0\) and every layer \(j\), the weighted sum \(\eta_j\) of the local cut errors over the twisted places in \(S\) satisfies \(\eta_j\le D_0 C_N' x^\varepsilon(B_0+N_c)\), where \(D_0\) is the reference row dimension. Apply row \(3\) of
\ref{B8.38} beyond the selection and row cutoffs with the common row
constant of \ref{B8.38b}. Then \(e'_{jw}=R_w/m\).

On the selected side the row frames are invertible. The maps \(\iota_t\)
are bijective and their type images at \(t_0\) are the \(E_\alpha\).
Require positive factors and the same rule for the row metric.
Put \(\lambda_{j,\alpha,v}=\lambda_{j,t_0,\alpha,v}\) and require
\begin{equation}\label{eq:p4c6-selected-spectrum}
 l_{6,j}\le\log\lambda_{j,\alpha,v}\qquad(j,\alpha,v).
\end{equation}
For each word take a \(K\)-linear map
\(B_u^{\rm c}:K^{d_u}\to(K^{D_{\rm c}})^\vee\) and factors
\(\zeta_{ut}\in K^\times\) with \(\zeta_{u,t_0}=1\) and
\(v_w(\zeta_{u,t_w})=0\). The definite evaluation is
\[
 P_{u,t}(\phi,F)=\zeta_{ut}B_u^{\rm c}(\phi)(\iota_tP_tF).
\]
At \(v_0\) take a conjugate-linear isometry \(s_E\) from the complex row
coordinates onto \(\mathcal R\), of dimension \(D_{\rm c}\).
Require a row-frame isometry \(V\) and, for each \(u\), a column-frame
isometry \(U_u\) and \(a'_u\in\C^\times\) such that
\[
 (B_u^{\rm c})_{v_0}(x)(y)
 =a'_u\langle s_E(Vy),P_u^{\rm per}(U_ux)\rangle .
\]
The inner product is conjugate-linear in its first argument.

The selectors and factors \(\kappa_v,\lambda_{jt\alpha v}\) agree
across the two sides. Off \(\mathrm{Tw}\), both \(F_{0,j,w}\) and
\[
 \operatorname{idx}_w(M_j,A_{j,t_w})
 +\operatorname{idx}_{\bar w}(M_j,A_{j,t_w})
 +\operatorname{idx}_w(T_{j,t_w},T_{j,t_w}^{\vee_{H_{j,t_w}}})
\]
vanish in every reference layer. Here \(\vee_H\) is the lattice dual
for the \(K\)-valued Hermitian form \(H\), with integral pairings.
Let \(R_0,R_{\mathrm c}\) be the glued row lattices of \ref{B4.36},
formed from \(\iota_{t_w}\Lambda_{t_w}\) on their respective sides.
Write \(\langle\cdot\rangle_v\) for the mean over infinite places. Put
\[
 \bar\kappa=\langle\log\kappa_v\rangle_v,\qquad
 \mu_0=\mu(R_0^\vee)-\bar\kappa,\qquad \mu_C=\mu(R_{\mathrm c}^\vee).
\]
Require \(\mu_{\max}(R_0)\le h_{\mathrm r}-\bar\kappa\).
For \(0\le\delta\le1/3\), require the passage inequalities of
\ref{B4.55a} for \(R_1=R_{\mathrm c}\), \(\lambda=\lambda_{\mathrm r}\)
and every \(U\ge\mu_{\max}(R_0)\).
\end{definition}

In the arithmetic construction \ref{B4.36a} gives the selector condition.
The first row of \ref{B8.38} gives passage beyond the selection and
row cutoffs.

\begin{definition}[Definite row and layer conditions at denominator \(s\)]\label{def:p4c6-definite-blocks}
Let \(s\) be a positive integer. The distinguished embedding takes the involution on \(K\) to complex conjugation. The set of tuple labels is finite. Each retained \(K\)-valued character of the row group \(G\) is the character of an irreducible representation over \(\overline K\).

Impose the following conditions on the row space for every tuple and on the space for every layer and tuple. The coefficient pairing is nondegenerate and Hermitian, and each class stabilizer preserves it. At a class \(c\), write \(A_c\) for its full coefficient lattice, \(A_c^\vee\) for its dual under this pairing and \(\Delta_c\subseteq A_c\) for its specified stabilizer-stable submodule. Require
\[
 sA_c^\vee\subseteq\Delta_c,\qquad s\Delta_c\subseteq A_c^\vee .
\]
Forms take stabilizer-invariant values at the classes. Their pairing is the sum of the coefficient pairings weighted by the reciprocal stabilizer orders. Its Gram matrix in the chosen basis is positive definite at every infinite place. The Gram entries and the matrices of \(G\) in this basis belong to the subfield of \(a\in K\) satisfying \(\sigma(\bar a)=\overline{\sigma(a)}\) for every embedding \(\sigma\). The group preserves the class-sum pairing and both modules of forms defined by taking values in \(A_c\) and in \(\Delta_c\), respectively.

The ordinary and inert operators together commute pairwise and commute with \(G\). Each operator has split characteristic polynomial over \(K\). The place maps for the ordinary and inert operators are each injective. At each ordinary place \(P\), its norm parameter \(q_P\) is positive. The block length is at least one, the list of prescribed blocks is nonempty, and each block parameter \(\nu(P)\) is nonzero at every ordinary place.
\end{definition}

\begin{definition}[Division inputs for a word]\label{def:p4c6-division}
Use the reference rows and selected columns of Definition~\ref{def:p4c6-core}. Fix a word \(u\) and the numerical data \(R\) of \ref{B5.08}. Write \(\beta_j\) for the row coefficient at layer \(j\) and \(w_{j,\rho}\) for the weight at position \(\rho\). Put
\[
 H=H^*=H_P+\sum_j(H_{B,j}-H_{X,j}),\qquad
 R_*=\text{the right-hand side of \ref{B5.08}} .
\]
Require \(d_u>0\), invertible column frames at infinity and a full lattice
\(N\subset K^{d_u}\). Take \(\sigma_w\in\R\) with finitely supported
floors \(\tau_w=\lfloor\sigma_w\rfloor\).

Start the selected column chain of \ref{B5.07} in frames \(X_{\mathrm c,v}\). Its maps \(K_i\) have invertible source and target
frames, and \(i(j)\) assigns a step to each layer. The terminal map
\(P_u^{\rm term}\) goes from the last space, with frames
\(X_{\mathrm{last},v}\), to the base test dual. Its composite with
the chain is \(P_{u,t_0}\). Choose invertible base test frames \(X_{0,v}\)
with \((X_{0,v}^{-1})^*X_{0,v}^{-1}=v(H_{V,t_0})\).
Use the height table of \ref{B8.35} for \(R\) under the rank, scale
and constant conditions of \ref{B5.07a}. Write \(h_P=A\) and
\(k_{j,v}=s_{j,v}\) for its terminal and layer entries. Use the coefficient maps and weights of \ref{B8.35}. Let \(a_i\) count the coefficient groups at step \(i\), let \(h_i^{\rm cl}\) bound their cardinalities, and let \(t_{\max}(i,v)\) bound the squared norms of the transitions from step \(i\) to step \(i+1\) at \(v\). Put \(C_{i,v}=a_i a_{i+1}h_i^{\rm cl}h_{i+1}^{\rm cl}t_{\max}(i,v)\). Require
\begin{equation}\label{eq:p4c6-height-inputs}
 \begin{gathered}
 \|P_u^{\rm term}\|_{X_{\mathrm{last},v}\to X_{0,v}^\vee}^2\le e^{h_P},
 \qquad \log C_{i,v}\le\sum_{j:i(j)=i} k_{j,v}
 \quad(i,v),\\
 \sum_j H_{X,j}\le\sum_j l_{6,j}.
 \end{gathered}
\end{equation}
For the coefficient bound use the convention \(\log0=0\).

For every reference type \(t\) and finite place \(w\), choose a basis
of \(\Lambda_t\) over the semilocal ring at \(w,\bar w\). Write
\(G_{t,w}\) for the quotient Hermitian Gram determinant in this basis
as in \ref{B5.08a}, and \(G_w=G_{t_w,w}\). Take a finite
\(S_y\supseteq\mathrm{Tw}\) with \(v_w(G_w)=0\) off \(S_y\).
Put \(J_{j,w}=\operatorname{idx}_w(T,T^{\vee_{H_{j,t_w}}})\) for
the reference layer \(T=T_{j,t_w}\). Choose \(L_G,E_3,E_{2,j},g_w,\eta_{w,j}\)
satisfying
\begin{equation}\label{eq:p4c6-gram-inputs}
 \begin{gathered}
 v_w(G_w)/D_R\le g_w\quad(w\in S_y),\qquad
 \sum_{w\in S_y}\ell_wg_w\le L_G,\\
 J_{j,w}/D_R\le e_wJ_{0,j,w}+\eta_{w,j}/D_R
 \quad(w\in\mathrm{Tw},\ 1\le j\le m),\\
 \sum_{w\in\mathrm{Tw}}\ell_w\eta_{w,j}/D_R\le E_{2,j}
 \quad(1\le j\le m),\qquad
 \sum_{w\in S_y}2\ell_we_w/p_w^2\le E_3 .
 \end{gathered}
\end{equation}
Set
\[
 F_*=\sum_j\sum_{w\in\mathrm{Tw}}\ell_we_w F_{0,j,w},\qquad
 B_*=\sum_j l_{3,j}+2F_*+L_G+\sum_j(E_{1,j}+E_{2,j})+E_3 .
\]
Choose \(E_{l3},Q_{\rm pass}\) with
\begin{equation}\label{eq:p4c6-row-budgets}
 \begin{gathered}
 \sum_j l_{3,j}\le\sum_{j,\rho}\beta_j w_{j,\rho}\log\Omega+E_{l3},
 \qquad
 E_b:=E_{l3}+L_G+\sum_j(E_{1,j}+E_{2,j})+E_3
 \le c_emx^\varepsilon(B_0+N_c),\\
 0\le\delta,\qquad 1\le Q_{\rm pass},\qquad 3\delta Q_{\rm pass}\le1,\\
 B_*+h_{\mathrm r}+28\lambda_{\mathrm r}\le Q_{\rm pass},\qquad
 \lambda_{\mathrm r}\le\lambda_0,\qquad 28\le C_0 .
 \end{gathered}
\end{equation}

\end{definition}

In the three-sided construction impose the conditions of Definition~\ref{def:p4c6-definite-blocks} with denominator \(s\). Assume that in every layer the transferred labels of any two distinct base spectral types differ in an ordinary Hecke polynomial or an inert eigenvalue at a place of that layer. Require a stable selection contained in the prescribed blocks. The selected side supplies the chain bounds by
\ref{B8.36} and the height comparison by \ref{B5.07b}.
On the reference side take \(S_y=\mathrm{Tw}\cup
\{w:v_w(G_{t_0,w})\ne0\}\). In a fixed \(K\)-basis the nonzero
Gram determinant has finite valuation support. The lattice generated
by that basis and \(\Lambda_{t_0}\) agree locally off a finite set
by commensurability. A semilocal change of basis with determinant
\(c\) changes the Gram valuation by \(v_w(c)+v_{\bar w}(c)\).
Outside the finite exceptional set and its conjugate this change is
zero. Thus \(S_y\) is finite and \(v_w(G_w)=0\) off \(S_y\).

Use the positive integer denominators \(s,d_g\) and the reference
layer denominators \(d_{j,t}\mid d_g\) of \ref{B5.08x}.
For the row group \(G\) assume \(\log|G|\le\lambda_0\).
At twisted places require the cut exponent
\(J^{\rm cut}_{jw}\le e_wJ_0(t_w,w,j)\) and the cut lengths \(a_{jw}\)
with \(\sum_{w\in\mathrm{Tw}}\ell_wa_{jw}\le D_Rb_j\) in each layer.
By the local bounds of \ref{B5.08x}, take
\(g_w=v_w(|G|s)\) and \(\eta_{w,j}=a_{jw}
+D_R(v_w(d_{j,t_w})+v_w(s))\).
For a positive integer \(q\), \ref{B4.52} gives
\(\sum_w\ell_wv_w(q)=\log q\) by factoring \(q\) into rational primes.
Consequently,
\[
 \sum_{w\in S_y}\ell_wg_w\le\log(|G|s)\le\lambda_0+\log s,\qquad
 \sum_{w\in\mathrm{Tw}}\ell_w\eta_{w,j}/D_R
 \le b_j+\log d_g+\log s .
\]
These right-hand sides give \(L_G\) and \(E_{2,j}\) in
\eqref{eq:p4c6-gram-inputs}. The denominator \(p_w/\sqrt{e_w}\) gives
the summand \(2\ell_we_w/p_w^2\), so \ref{B5.08x} gives \(E_3=1\).
Under \ref{B5.09a}, take \(Q_{\rm pass}=Q_0\) for sufficiently large
selected members.

\begin{theorem}[Four conclusions in coordinates]\label{B4.55}
Use \ref{def:p4c6-core} and \ref{def:p4c6-division}. Write \(B=B_u^{\rm c}\)
with operator norms in the column and dual selected row frames.
At every finite \(w\) suppose \(P_{u,t_w}(\phi,F)\in\mathcal V_{w,\sigma_w}\)
for all \(\phi\in N\) and all selected \(F\in M_{t_w}\), where
\(\mathcal V_{w,\sigma_w}\) consists of zero and elements of valuation
at least \(\sigma_w\). Let \(N^\sigma\) be the shift by \(\tau_w\)
from \ref{B4.40}. Then \(B(N^\sigma)\subseteq R_{\mathrm c}^\vee\),
and there are \(c_v>0\) with \(\|B\|_v^2\le c_v\) and
\(\langle\log c_v\rangle_v=H-\bar\kappa\). Moreover,
\[
 \mu_C-\bar\kappa+H-2F_*\le R_*,
 \qquad
 \mu_{X_{\mathrm c}}(N^\sigma)-\mu_{X_{\mathrm c}}(N)
 =2\sum_w\ell_w\tau_w .
\]
\end{theorem}

\begin{proof}[Outline]
Use \ref{B4.40}, \ref{B4.45}, and \ref{B4.50} for fullness, inclusion, and slope. Adjust the log norm bounds from \ref{B5.07a} and \ref{B5.07} to mean \(H-\bar\kappa\), then use \ref{B5.08x} and the row budgets and apply \ref{B5.08} with the adjusted slope. The complete proof is in \S\ref{prf:B4.55}.
\end{proof}

\begin{lemma}[Identification of the row base term]\label{B4.55b}
Use the finite support \(S\) in the arithmetic sum of \ref{B5.02}. Let \(\iota\) send each arithmetic place \(w\) to the place of the row coefficient field represented by \(\mathfrak p_w\). Use this bijection and the layer bijection \(j\mapsto a(j)\) to \(\{1,\ldots,m\}\).
Assume the reference row is untwisted off \(\mathrm{Tw}\). If
\begin{equation}\label{eq:p4c6-base-tie}
 F_{0,j,y}=F_{a(j)}(\iota^{-1}y)\qquad\text{for all }y,j,
\end{equation}
then \(2F_*=\Fo\). On the fixed arithmetic scale
\eqref{eq:p4c6-base-tie} holds for the weight tuples of \ref{B4.02a} and the finite base terms of \ref{B4.02}.
\end{lemma}

\begin{proof}[Outline]
Use \ref{B4.02a}, \ref{B3.46}, and \ref{B3.52} to match weights and Frobenius orbits, hence the capacities of \ref{B4.01} and base terms of \ref{B4.02}. Sum over the common support via the place and layer bijections, using the doubled weights of \ref{B5.02}. The complete proof is in \S\ref{prf:B4.55b}.
\end{proof}

\begin{lemma}[Passage between character levels]\label{B4.55a}
Let \(K\) be totally complex. Let a finite group \(G\) of order \(M\)
act unitarily at infinity on a metrized \(K\)-space and preserve a full
lattice. Assume the retained irreducible types have scalar endomorphisms and
distinct characters. Their central projectors must sum to the identity. For \(i=0,1\), take a linear \(K\)-valued character
\(\chi_i\) on a subgroup of order \(h_i\). In type \(s\) let their
multiplicities be \(1,n_s\), respectively, with
\begin{equation}\label{eq:p4c6-multiplicities}
 (1-\delta)\bar n\le n_s\le(1+\delta)\bar n,\qquad
 \bar n>0,\quad 0\le\delta\le1/3 .
\end{equation}
Let \(P_i\) be the ambient lattice cut by the \(\chi_i\)-space.
Assume the row embedding identifies \(R_i\) with a nonzero lattice
between \(h_iP_i\) and \(P_i\), in the ambient frames multiplied
by \(\sqrt{h_i}\). For \(U\ge\mu_{\max}(R_0)\) and \(\lambda\ge\log M\),
\[
 \begin{aligned}
 \mu_{\max}(R_1)&\le U+14\lambda,\\
 \mu(R_1^\vee)&\le\mu(R_0^\vee)
 +3\delta\bigl(\mu(R_0^\vee)+U+28\lambda\bigr)+28\lambda .
 \end{aligned}
\]
\end{lemma}

\begin{proof}[Outline]
Compare row slopes with character cuts using \ref{B4.46} and clear projector denominators. Use integral isotypic maps and their adjoints to bound maximum slopes and compare degrees, then use multiplicity bounds for the degree deficit and duality for the mean slope. The complete proof is in \S\ref{prf:B4.55a}.
\end{proof}

\begin{theorem}[Return to the original period data]\label{B4.56}
Fix the arithmetic data and local budgets of \ref{B4.05}, and constants \(C,C_{\mathrm{abs}}\in\R\). For each arithmetic word
\(u\) choose its selected column and an input of \ref{def:p4c6-division}
for a coordinate lattice \(\mathcal N_u\subset K^{d_u}\), numerical data \(R\)
and \(\sigma_w=s_{u,w}\) under the place identification.
Use the bridges \(E_u,J_t\) and period identity of \ref{B4.38}.
For each selected type \(t\) and \(F\in M_t\) require
\(J_tF\in V_0(\theta(t))\) and an integral expression for \(J_tF\)
in the generators \(\mathcal G_{\theta(t)}\) via the coefficient
embedding. Require \(\theta(t_w)=\mu_w\) at every finite place, where
\(\mu_w\) is the arithmetic label of \ref{B4.29}.

The column isometry \(E_u\) identifies the shift of \(\mathcal N_u\)
by the floors of \(\sigma_w\) with \(N'_u\). A row isometry \(E_R\)
identifies \(R_{\mathrm c}\) with \(RR\). The pairing identity of
\ref{B4.45} reads \(B_u(E_u\phi)(E_Ry)=B_u^{\rm c}(\phi)(y)\).
Assume the period bound of \ref{B4.29}: for every \(u,w\), every \(\phi\in N_u\) and every \(F\in\widetilde M_{\mu_w}\cap V_0(\mu_w)\), the period of \(\phi\) against \(F\) is the image of some \(y\in K\) under the coefficient embedding, with \(y=0\) or \(v_w(y)\ge s_{u,w}\). Assume also \(N'_u=I_uN_u\) for every \(u\), as in \ref{B4.39}. For each of the reference and two auxiliary labels, and each label \(\mu_w\), assume the integral decomposition of \ref{B4.37}: every vector in \(V_0(\mu)\) in the \(\mathcal O\)-span of \(\mathcal G_\mu\) is a finite \(\mathcal O\)-linear combination of vectors in \(\widetilde M_\mu\cap V_0(\mu)\), with scalars acting through the coefficient embedding.

Put \(a=xe^{x^\varepsilon}\), \(n_T=\#TT\) and
\(L_a=\log(a\,\mathrm{rd}(E))\), with \(\mathrm{rd}(E)\) the root
discriminant of the host field. For top positions \(t\in TT\)
with weights \(a_t\), put
\[
 \mathrm{MAIN}(u)=\sum_{t\in TT}
 \bigl[m(\varphi(a_t)-a_t\log\Omega)-a_t\log T_u\bigr],
 \qquad
 \varphi(z)=z\log z-z\ (z>0),\quad \varphi(0)=0 .
\]
For the coordinate bounds \(H_u,R_{*,u}\) of \ref{def:p4c6-division}
and the recorded arithmetic height bound \(h_u\), assume
\begin{equation}\label{eq:p4c6-package-comparisons}
 H_u-\bar\kappa=h_u,\qquad
 2F_*=\Fo,\qquad
 R_{*,u}\le\mathrm{MAIN}(u)+Cn_TaL_a+
 Cm(B_0+N_c)x^{C_{\mathrm{abs}}\varepsilon}.
\end{equation}
On the fixed scale the middle equality follows from \ref{B4.55b}
with the tuples of \ref{B4.02a} and the base terms of \ref{B4.02}.
Then for every word \(u\),
\[
 \begin{gathered}
 B_u(N'_u)\subseteq RR^\vee,\qquad
 B_u\ne0\ \Longrightarrow\
 h(B_u):=\langle\log\|B_u\|_v^2\rangle_v\le h_u,\\
 \mu(RR^\vee)+h_u-\Fo
 \le\mathrm{MAIN}(u)+Cn_TaL_a+
 Cm(B_0+N_c)x^{C_{\mathrm{abs}}\varepsilon}.
 \end{gathered}
\]
\end{theorem}

\begin{proof}[Outline]
Identify the shift ideal by \ref{B4.05}, \ref{B4.40}, \ref{B4.39}, and \ref{B4.48}. Decompose by \ref{B4.37}, bound valuations by \ref{B4.29}, and transfer them by \ref{B4.38}, then transport \ref{B4.55} by pairing and isometry and substitute the package comparisons. The complete proof is in \S\ref{prf:B4.56}.
\end{proof}


\part{Two bounds on one slope}

% ---- p5c1 ----
\chapter{The lattice after division}\label{ch:lattice-after-division}

Fix \(x>1\), \(0<\varepsilon<1/52\), and a translation \(u_0\) with finite word set \(W(u_0)\). The host CM field is \(E\), with \(f=[E:\Q]/2\) column positions \(\rho\). The coefficient field \(K\) is a totally complex number field with a distinguished infinite place \(v_0\), given by an embedding \(\sigma_0\), and a ring involution \(a\mapsto\bar a\) satisfying \(\sigma_0(\bar a)=\overline{\sigma_0(a)}\); the field \(K\) itself need not be a CM field. For a finite place \(w\) of \(K\), write \(\mathfrak p_w,p_w,e_w\) for its prime ideal, residue characteristic and ramification index. Let \(S\) be the finite support of the local data, and put
\[
 \ell_w=\frac{\log N\mathfrak p_w}{[K:\Q]},\qquad c_w=2e_w\ell_w .
\]
Write \(\langle\cdot\rangle_\infty\) for the mean over infinite places. In the metrics of Part~IV, a nonzero map \(B\) has height \(h(B)=\langle\log\|B\|_v^2\rangle_\infty\). Slope differences use twice the normalized Arakelov convention. The height \(h(B)\) is twice the archimedean contribution to the height of a linear map in \cite[\S2.6]{bost-chen-2013}.

For a word \(u\), let \(N_u^0\) denote the input lattice \(N_u\) of \ref{B4.39} and let \(I_u\) be the fractional ideal of \ref{B5.03}. Set \(N'_u=I_uN_u^0\). Let \(M_u\) be the Hermitian lattice of column forms for \(u\) before the first lift. The \(K\)-linear column inclusion identifies the word space with the subspace of the zeroth column killed by all marked diamond sums. Let \(M_{H,u}\) be the inverse image of \(M_u\) under this inclusion, with the metric of \(N_u^0\). If \(V_u,W\) are the spans of \(N'_u,RR\), the period map is \(B_u:V_u\to W^\vee\). Put \(C^\vee=RR^\vee=\{\ell\in W^\vee:\ell(RR)\subseteq\mathcal O_K\}\) and \(\mu_C=\mu(C^\vee)\).

Use
\[
 L=x^\varepsilon,\quad b=e^L,\quad a=xb,\quad B_0=fmbL,\quad
 L_a=\log(a\,\mathrm{rd}E),\quad n_T=\#TT,\quad S_a=\sum_{t\in TT}a_t.
\]
Here \(\mathrm{rd}E\) is the root discriminant. The set \(TT\) consists of the top positions, with integer top weights \(a\le a_t\le2a\). The chain has \(m\) lifts and starts at rank \(k\). Its ranks are \(N_j=k+j(t_{\rm ch}+1)\) for \(0\le j\le m\), with \(t_{\rm ch}\) as in \ref{B4.03}. Write \(T_u\) for the word's level prime, and use the auxiliary size \(b'\) and prime cutoff \(X''\) of the window.

\section{Prime powers and finite-place budgets}

\begin{definition}[The required prime-power sum]\label{B5.01}
For the host group \(A\), set
\[
 G=(\Z/2\Z)^2\times A,\qquad
 \Phi=\{((e_1,e_2),t)\in G:e_1=0\},\qquad
 r_t(g)=((1,1),0)+g+((0,0),t).
\]
Let \(F_p\) be the chosen Frobenius element and \(I_p\le G\) the inertia subgroup. The exceptional primes are \(q_1,d\). For \(N\in\N\), \(a_t\in\R\) and \(w_{u_0j}:G\to\R\), put \([z]_+=\max(z,0)\) and
\[
\begin{aligned}
 A_{tj}^{(N)}
 ={}&\sum_{\substack{p\le N,\ p\ {\rm prime}\\p\notin\{2,q_1,d\}\\ F_p\notin\Phi}}
 2\log p\sum_{\substack{1\le i\le N\\iF_p\notin\Phi}}
 \left[\frac{a_t}{p^i}-w_{u_0j}(r_t(iF_p))\right]_+\\
 &+\log d\sum_{i=1}^{N}\sum_{\substack{g\in iF_d+I_d\\g\notin\Phi}}
 \left[\frac{a_t}{d^i}-w_{u_0j}(r_t(g))\right]_+ .
\end{aligned}
\]
Then \(\ASK(u_0)=\sum_{t\in TT}\sum_{j=0}^{m-1}A_{tj}^{(N)}\), with \(j\) denoting layer \(j+1\).
\end{definition}
\begin{lemma}[Cutoff for a positive partner weight]\label{B5.01a}
In Definition~\ref{B5.01}, if the partner weight is at least one and $a_t\le2a$, every strictly positive summand has both its prime and its exponent at most $\lceil2a\rceil$.
\end{lemma}
\begin{proof}
The first binary coordinate forces \(i\) odd in the ordinary-prime sum. A positive term with partner weight \(w\ge1\) requires \(p^i<a_t/w\), with \(p=d\) in the exceptional sum. Thus \(N=\lceil2a\rceil\) includes all positive terms for such weights.
\end{proof}

\begin{definition}[Local multiplicities and their global totals]\label{B5.02}
Use the gain \(g_{jw}\), base term \(F_{jw}\) and loss \(E_{ujw}\) of layer \(1\le j\le m\) from Parts~III and~IV. Set \(\gainD_j(w)=0\) when \(p_w\) belongs to the discarded set \(\mathcal D\) of \ref{B4.04}, and \(\gainD_j(w)=g_{jw}\) elsewhere. Put
\[
 \sg(u,w)=\sum_j(\gainD_j(w)+F_{jw}-E_{ujw}).
\]
For the containment budget \(\tau_w\) and excess colength \(\mathrm{io}_w\) of \ref{B3.52}, define
\[
\begin{aligned}
 \Ga&=\sum_w c_w\left(\sum_j\gainD_j(w)-\tau_w\right),
 &\mathcal T&=\sum_w c_w\tau_w,\\
 \Fo&=\sum_{j,w}c_wF_{jw},
 &\Ee_u&=\sum_{j,w}c_wE_{ujw},
 &\mathrm{IO}&=\sum_w e_w\ell_w\,\mathrm{io}_w .
\end{aligned}
\]
The place sums are over \(S\).
\end{definition}
\begin{lemma}[Sum of the local budgets]\label{B5.02a}
The totals \(\Ga,\Fo\) are common to the words, and \(\sum_w c_w\sg(u,w)=\Ga+\Fo-\Ee_u+\mathcal T\).
\end{lemma}
\begin{proof}
The gain, containment and base terms in Definition~\ref{B5.02} are independent of the word. Interchange the finite layer and place sums and add back \(\mathcal T\).
\end{proof}

\begin{definition}[The ideal prescribed by period valuations]\label{B5.03}
At \(w\) use the selected character \(\mu_w\), row lattice \(\widetilde M_{\mu_w}\), test subspace \(V_0(\mu_w)\) and column inclusion \(j_N\) of Part~IV. Normalize \(v_w\) to be one on uniformizers. The unrounded bound requires that for every \(\phi\in N_u^0\) and \(F\in\widetilde M_{\mu_w}\cap V_0(\mu_w)\), the column period at character \(\mu_w\), evaluated on \(j_N\phi\) and \(F\), be the image under the chosen embedding of some \(y\in K\) with \(y=0\) or \(v_w(y)\ge e_w\sg(u,w)\). Put
\[
 s_{uw}=\lfloor e_w\sg(u,w)\rfloor,\qquad I_u=\prod_{w\in S}\mathfrak p_w^{-s_{uw}}.
\]
\end{definition}

By \ref{B4.39}, the ideal of Definition~\ref{B5.03} gives \(N'_u=I_uN_u^0=\operatorname{span}_{\mathcal O_K}\{a\phi:a\in I_u,\ \phi\in N_u^0\}\) for every word \(u\). The unrounded bound gives the threshold \(s_{uw}\). \ref{B4.29} gives the rounded bound and \ref{B4.38} extends it to coordinate rows. The shift of \ref{B4.40} is \(N'_u\) by \ref{B4.41}. For the pairing in \ref{B4.44}, \ref{B4.45} gives \(B_u(N'_u)\subseteq RR^\vee\).

\section{The two slopes of the divided lattice}

\begin{lemma}[The exact change of mean slope]\label{B5.04}
Let \(N\subseteq M\) be lattices of rank \(d>0\) in the same Hermitian space over \(K\), with \(M\) finitely generated over \(\mathcal O_K\). Write \(\lambda_w=\operatorname{length}_w(M_{\mathfrak p_w}/N_{\mathfrak p_w})\). Then
\[
 \mu(N)=\mu(M)-\frac2d\sum_w\ell_w\lambda_w.
\]
If \(N_u^0\subseteq M_{H,u}\) have quotient supported on \(S\), write \(\ell^{\rm col}_{uw}=e_w^{-1}\operatorname{length}_{\mathcal O_{K,\mathfrak p_w}}((M_{H,u}/N_u^0)_{\mathfrak p_w})\) for its normalized local length. For \(d_u=\operatorname{rk}_K N'_u>0\),
\[
 \mu(N'_u)=\mu(M_{H,u})+2\sum_{w\in S}\ell_ws_{uw}
 -\frac2{d_u}\sum_{w\in S}\ell_we_w\,\ell^{\rm col}_{uw}.
\]
\end{lemma}
\begin{proof}
The covolume ratio is \(\prod_w(N\mathfrak p_w)^{\lambda_w}\), by the cardinalities of the local composition factors. Divide its negative doubled logarithm by \(d[K:\Q]\). Apply the formula to \(N_u^0\subseteq M_{H,u}\) and use \ref{B4.47} and \ref{B4.48} for the twist by \(I_u\).
\end{proof}

\begin{lemma}[A local comparison for the maximum slope]\label{B5.05}
Suppose \(N_u^0,M_{H,u}\) span the same space of dimension \(d_u>0\), and \((N_u^0)_{\mathfrak p_w}\subseteq\mathfrak p_w^{k_w}(M_{H,u})_{\mathfrak p_w}\) at every finite place, with \(k_w\in\N\) supported on \(S\). Then
\[
 \mu_{\max}(N'_u)\le\mu_{\max}(M_{H,u})+2\sum_{w\in S}\ell_w(s_{uw}-k_w).
\]
\end{lemma}
\begin{proof}
For \(J=\prod_{w\in S}\mathfrak p_w^{k_w}\), the hypotheses give \(N'_u\subseteq I_uJM_{H,u}\) with its induced metric. Apply \ref{B4.47} after dividing each sublattice by \(I_uJ\). Take the supremum.
\end{proof}

\begin{lemma}[Numerical conditions at a split level]\label{B5.05b}
Fix a word \(u\). Put \(T=T_u\), \(N=N_m\), \(a=xb\), and let \(a_{\max}\) be the maximum of zero and the leading marked weights of this word. For \(0\le j\le m\), write \(\lambda_{j,\tau,i}\), \(0\le i<N_j\), for the integer highest weight at type \(\tau\). The level conditions on these weights are
\[
 \begin{gathered}
 2fN+1<T,\qquad 4a_{\max}<T,\\
 2(\lambda_{j,\tau,i}-\lambda_{j,\tau,h}+N_j)<T
 \quad(0\le j\le m,\ 0\le i,h<N_j),\\
 -\lambda_{0,\tau,k-1}\ge0\qquad(\tau\ \text{marked}).
 \end{gathered}
\]
They hold whenever \(x>1,\ m\ge4,\ 1\le f\le N,\ 3N\le b,\ 24x\le b,\ 0<b'\le b,\ T\ge b^{m/2}\), and the weights satisfy
\[
 a_{\max}\le2xb,\qquad
 -(2b+b')\le\lambda_{j,\tau,i}\le 2xb+(b')^2/64,
 \qquad \lambda_{0,\tau,k-1}\le0\quad(\tau\ \text{marked})
\]
with the last inequality required only at the marked types and the two-sided bound required at every type.
\end{lemma}
\begin{proof}
We have \(T\ge b^2\), \(b>24\) and \(2fN+1\le2N^2+1<9N^2\le b^2\). Also \(4a_{\max}\le8xb\le b^2/3\). At any two indices of the same weight the two-sided estimate gives
\[
 2(\lambda_{j,\tau,i}-\lambda_{j,\tau,h}+N_j)
 \le 4xb+b^2/32+8b
 \le (1/6+1/32+1/3)b^2<b^2 .
\]
Here \(N_j\le N\le b\). The final sign condition is the bottom-entry hypothesis.
\end{proof}

\begin{lemma}[Cost at the level prime]\label{B5.05a}
Use the split level models with the trivial class stabilizers of \ref{B8.54} and the weight inequalities of \ref{B5.05b}. For a selected word at type \(\rho\), the first \(k\) integral weight coordinates are \(U_i-(k-1)/2+i\), numbered from zero. Lift \(q+1\) appends \(t_{\rm ch}+1\) copies of \(s_{q+1}-w_{q+1,\rho}\), with \(s_j\) as in \ref{B5.08b}. At a mark \(\rho=(0,s)\) the last initial coordinate is \(-b_s\). The endpoint bounds for the selected words in \ref{B5.11c} and the bounds \(s_j\le w_{j\rho}\le2b\) in \ref{B5.08b} give \(|\lambda_{j,\rho,i}|\le2a\) eventually. Thus \(4a_{\max}\le8a\) and \(2(\lambda_{j,\rho,i}-\lambda_{j,\rho,h}+N_j)\le8a+2N_m\). By \ref{B8.51} and the rank bounds of \ref{B8.59}, eventually \(m\ge4\) and \(T_u\ge b^{m/2}\ge b^2>\max\{8a+2N_m,2fN_m+1\}\). Assume $a_{\max}\le2a$, $k\le N_m\le b'$ and, at every active marked type, $b_v\le2b'$ and $W_{u,\tau}\le3b'$ in the notation of \ref{B4.03}, where $b_v=-\lambda_{0,\tau,k-1}$. Put \(T=T_u,\ H=T-1\), let $d$ be the dimension of the common kernel of the marked diamond sums in the zeroth column, and let \(d_M=\dim_\C M^0_{u,\C}\) be the zeroth column dimension. Then
\[
 d=d_M(1-1/H)^f,\qquad
 \mu(M_{H,u})=\mu(M_u)-\frac{f\log H}{H-1},\qquad
 \mu_{\max}(M_{H,u})\le\mu_{\max}(M_u).
\]
The containment of \ref{B5.05} holds at \(w\mid T\) with exponent zero. For every finite \(F\subseteq\{w:w\mid T\}\) and \(d>0\),
\[
 \sum_{w\in F}\frac{2\ell_we_w\,\ell^{\rm col}_{uw}}d
 \le36n_Tm(a/k+b')\log T.
\]
\end{lemma}
\begin{proof}[Outline]
Trivial class stabilizers and pro-\(g_i\) kernels at two primes make the diamond action free, giving the rank and slope formulas from orbit splittings. Up-down operator identities on the diamond kernel, \ref{B4.03}, and \ref{B4.52} bound the colength sum. The complete proof is in \S\ref{prf:B5.05a}.
\end{proof}

\begin{theorem}[Mean and maximum slopes after division]\label{B5.06}
Fix \(u\) and use the comparisons of \ref{B5.04}--\ref{B5.05a}, with \(d=d_u>0,\ d_M>0,\ T_u>2,\ k_w=e_w\tau_w\). Put \(\delta=1-d/d_M,\ S_T=\{w\in S:p_w=T_u\}\), and
\[
 \sigma'_w=s_{uw}/e_w,\quad
 R_u=\sum_{w\in S}c_w(\sg(u,w)-\sigma'_w),\quad
 P_u=\sum_{w\in S_T}c_w\,\ell^{\rm col}_{uw}/d .
\]
Assume \(\tau_w\in\N,\ \mathrm{io}_w\ge0,\ \tau_w=0\) on \(S_T\), and \(\ell^{\rm col}_{uw}\le d_M(\tau_w+\mathrm{io}_w)\) off \(S_T\). Suppose
\[
\begin{gathered}
 \mathcal T\le b^2,\quad 2f/T_u\le1/2,\quad
 4fb^2/T_u+f\log(T_u-1)/(T_u-2)\le1,\\
 P_u\le c_1n_Tm(a/k+b')\log T_u,\quad R_u\le c_1X'',\quad
 c_1\ge4,\quad B_0,\mathrm{IO}\ge0.
\end{gathered}
\]
For \(J_u=\Ga+\Fo-\Ee_u\) and \(\mathrm{Lo}_u=c_1\mathrm{IO}+c_1n_Tm(a/k+b')\log T_u+c_1X''+1\), one has
\[
 \mu(M_u)+J_u-\mathrm{Lo}_u\le\mu(N'_u),\qquad
 \mu_{\max}(N'_u)\le\mu_{\max}(M_u)+J_u+c_1mB_0 .
\]
\end{theorem}
\begin{proof}[Outline]
Use the index formula \ref{B5.04} and the level correction \ref{B5.05a} to bound the mean slope after rounding and colength losses. Apply \ref{B5.05} to the maximum slope and sum the selected-prime budgets using \ref{B5.20a}. The complete proof is in \S\ref{prf:B5.06}.
\end{proof}

\section{The height of a chain}

\begin{theorem}[The archimedean chain estimate]\label{B5.07}
Let \(q\) map the test space to its Hermitian quotient \(Q\), and let \(\iota:Q\to W\) be the row isomorphism at \(\mu_0\). Suppose the spectral spaces \(Q_p\), \(p\in I\), span \(Q\), with \(I\) finite and nonempty. For \(z\in Q_p,\ y\in Q\) suppose
\begin{equation}\label{eq:B5-chain-metric}
 \langle\iota z,\iota y\rangle_v=c_{pv}\langle z,y\rangle_v,\qquad
 c_{pv}=\kappa_v\prod_{j=1}^m\lambda_{jpv},\qquad
 \kappa_v,\lambda_{jpv}>0.
\end{equation}
Suppose \(B_u\ne0\) and \((\iota q)^\vee B_u\) factors as kernel maps \(K_i\) followed by \(P_u^{\rm term}\), with layer \(j\) assigned to \(i(j)\). At each infinite place \(v\), assume
\[
 \|P_u^{\rm term}\|_v^2\le e^{h_P},\qquad
 \|K_i\|_v^2\le\exp\left(\sum_{j:i(j)=i}h_{K,jv}\right).
\]
Then
\[
 \log\|B_u\|_v^2\le z_v=h_P+\sum_jh_{K,jv}-\log\kappa_v
 -\min_{p\in I}\sum_j\log\lambda_{jpv}.
\]
If \(\sum_jL_{6,j}\le\langle\min_{p\in I}\sum_j\log\lambda_{jpv}\rangle_\infty\), then
\[
 h(B_u)\le\langle z_v\rangle_\infty
 \le h_P+\sum_j\langle h_{K,jv}\rangle_\infty-\sum_jL_{6,j}
 -\langle\log\kappa_v\rangle_\infty .
\]

In the tables of \ref{B8.35}--\ref{B8.36} and the realization of \ref{B8.38}, take the common size constant \(C_{\rm sz}\). At an embedding \(\sigma:K\hookrightarrow\C\), the row metric scalar is \(\kappa_{[\sigma]}\), where \([\sigma]\) is the infinite place of \(\sigma\). Let \(x_1\) be the family cutoff. At each \(x>x_1\) and each index \(i\in I_x\), the construction assigns every word \(u\) its height allowance \(h_u\) by the equality below. With these metrics the period satisfies
\[
 h(B_u)\le h_u
 =H_{P,u}+\sum_jH_{K,uj}-\sum_jL_{6,j}-\langle\log\kappa_v\rangle_\infty.
\]
Use the width \(S_c\), the block allowance \(\eta\), and the degrees \(\nu_{uj\rho i}\) of \ref{B4.02a}--\ref{B4.03}. Put \(n_V=k+m-1+t_{\rm ch},\ v_E=\mathrm{rd}E/2\), and \(\varphi(t)=t\log|t|-t\) for \(t\ne0\), with \(\varphi(0)=0\). For the quadratic character \(\chi_d\), put \(D_1=\Re(L'(1,\chi_d)/L(1,\chi_d))\) and \(\log\Omega=\gamma+\log(\pi/v_E)-D_1\). The budgets are
\[
\begin{aligned}
 H_{P,u}&=-S_a\log T_u+C_{\rm sz}
 \bigl(n_Tb'N_m+f(m\eta+S_c)N_m+fN_m^2\bigr)L_a,\\
 H_{K,uj}&=\sum_\rho K_{uj\rho}+C_{\rm sz}fN_m^2(L_a+m^3L),\\
 K_{uj\rho}&=\sum_{i=0}^{N_{j-1}-1}
 \bigl(\varphi(\nu_{uj\rho i})-\nu_{uj\rho i}\log\Omega\bigr),\\
 L_{6,j}&=\sum_\rho\log\operatorname{Hook}(a^+_{j\rho})-C_{\rm sz}fn_V^2L',
 \qquad L'=\log(b\,\mathrm{rd}E\,g_0g_1n_V).
\end{aligned}
\]
Here \(g_0,g_1\) are the fixed theta primes, and \(a^+_{j\rho,i}=\kappa^{\rm big}_{j,\rho,i}\) of \ref{B4.02} for \(0\le i<k+j-1\). For a partition \(\alpha\) with \(p\) parts including trailing zeros,
\[
 \log\operatorname{Hook}(\alpha)
 =\sum_{i=0}^{p-1}\log\Gamma(\alpha_i+p-i)
 -\sum_{0\le i<l<p}\log(\alpha_i-\alpha_l+l-i).
\]
\end{theorem}
\begin{proof}[Outline]
The spectral row weights bound the smallest squared singular value of the dual row map. Multiply the terminal and kernel norm bounds from \ref{B5.07a}, then average the logarithms and use the spectral lower bound of \ref{B5.07b}. The complete proof is in \S\ref{prf:B5.07}.
\end{proof}

\begin{lemma}[The last branching map and the class sum]\label{B5.07c}
Let \(N=n+1\). At each type \(\tau\) let \(\lambda,\sigma\) be dominant integral weights of ranks \(N,n\) with common determinant shift and \(\lambda_i\ge\sigma_i\ge\lambda_{i+1}\). Put
\[
 A_i=\sigma_i-\lambda_{i+1},\quad B_i=\lambda_i-\sigma_i,\quad
 D_{0,\tau}=\sum_iA_i,\quad n_\tau=|\lambda|-|\sigma|=\lambda_1-D_{0,\tau},
\]
where bars denote sums of signed entries. In the Fock variables let \(P_i\) be the leading minor of size \(i\) and \(Z_i\) the same minor with row \(i\) replaced by row \(N\). After removing the determinant shift put
\[
 h_\lambda=\prod_iP_i^{A_i+B_i},\quad
 \xi=\prod_iP_i^{A_i}Z_i^{B_i},\quad
 \kappa_\tau=\|\xi\|_F^2/\|h_\lambda\|_F^2 .
\]
The Fock metric has orthogonal monomials of squared norm \(\|X^\alpha\|_F^2=\prod_s\alpha_s!\). Set
\[
 c_\tau=\prod_{i=1}^n\prod_{v\in I_i\setminus\{\sigma_i\}}
 \prod_{j=1}^n(v-i-\sigma_j+j),\qquad
 I_i=[\lambda_{i+1},\min(\lambda_i,\lambda_{i+1}+D_{0,\tau})]\cap\Z.
\]
With \(|t_\tau|^2=T>0\), the squared branching norm in unit coefficient frames is \(C(T)=T^{-\sum_\tau n_\tau}\prod_\tau c_\tau^2\kappa_\tau\). In inverse-stabilizer class metrics the associated period on \(\psi,F\) has absolute square at most \(C(T)\|\psi\|^2\|F\|^2\).
\end{lemma}
\begin{proof}[Outline]
Extract the last-row coefficient and use multiplicity-one branching to identify the highest vector. Tensor the adjoint norm identities, scale by the CM generator of \ref{B4.03}, and apply Cauchy--Schwarz to the coefficient pairings and weighted class sum. The complete proof is in \S\ref{prf:B5.07c}.
\end{proof}

\begin{lemma}[Vanishing of the constant terms of the level kernel]\label{B5.07d}
Let \(E/F\) be the CM extension of the selected realization. Let \(V\) be its definite Hermitian space of dimension \(n\) and \(W\) its Hermitian space of dimension \(r\). Suppose that at one marked split place the finite Schwartz function, in the prescribed lattice frames, is supported on matrices \(X\in\operatorname{Mat}_{n,r}(\mathcal O)\) whose first \(r\) rows form an invertible matrix modulo the maximal ideal. For the reference theta kernel, every lift from the definite class space has zero constant term along every proper rational parabolic of \(\U(W)\).
\end{lemma}
\begin{proof}[Outline]
Average along a rational isotropic line and use Heisenberg covariance under transvections to eliminate nonconstant Fourier modes. Residue-field modulations at the marked split place eliminate the constant mode, so the parabolic constant term vanishes. The complete proof is in \S\ref{prf:B5.07d}.
\end{proof}

\begin{lemma}[The spherical integral for the self-dual lattice vector]\label{B5.07e}
Let \(w\) be a finite unmarked place of \(F\) in the reference tuple, let \(Q_w=Nw\), and give the stabilizer \(K_w\) of the self-dual lattice in \(\U(W)(F_w)\) volume one. Put \(n=r+t\), \(t\ge r+5\), and \(u=(1+t)/2\). Let \(\phi_w^\circ(g)\) be the Weil coefficient of the lattice vacuum divided by its value at the identity, in the convention linear in the translated vector. Let \(F_1\) be a \(K_w\)-fixed simultaneous Hecke eigenvector in a unitary representation with untwisted base-change Satake multisets \(B_P\), \(P\mid w\). Let \(\ell\) be a linear functional bounded on its translates. Write \(R(g)\) for translation, \(\chi_P\) for the splitting character of \(V\) evaluated at a uniformizer at \(P\), and \(\epsilon_w=1\) in the split case, \(-1\) in the inert case. Then, with
\[
 d_w=\prod_{h=t+1}^{n}(1-\epsilon_w^h Q_w^{-h})^{-1},
\]
one has the absolutely convergent identity
\begin{equation}\label{eq:B5-spherical-cosets}
 \sum_{gK_w}\phi_w^\circ(g)\,\overline{\ell(R(g)F_1)}
 =
 \frac{1}{d_w}
 \prod_{P\mid w}\prod_{\beta\in B_P}
       (1-\chi_P^{-1}\beta (NP)^{-u})^{-1}\,
 \overline{\ell(F_1)} .
\end{equation}
\end{lemma}
\begin{proof}[Outline]
Match the split and inert Cartan vacuum coefficients to the spherical doubling section. Apply the unramified doubling formula of Eischen, Harris, Li and Skinner to the Hecke eigenvector, using the trivial spherical representation to majorize the coset sum. The complete proof is in \S\ref{prf:B5.07e}.
\end{proof}

\begin{lemma}[The full theta lift and its projection at a coefficient embedding]\label{B5.07y}
Fix \(1\le j\le m\) at \(\mu_0\) in a normalized primitive realization with the standard definite lattice, an admissible choice of member parameters, and the extremal vectors and level action of \ref{B8.55}. Assume the rapid-decay, Siegel--Weil, doubling, finite-length and formal-degree inputs of \ref{B1.15ac}. The twist-prime set at \(\mu_0\) is empty. Each marked prime \(P\) is prime to \(Q_Sqd\,N(\mathfrak c_K)\), where \(Q_S\) is the excluded-prime product and \(q,d\) are the auxiliary primes. The host datum specifies \(\mathfrak c_K\) as a nonzero integral ideal of \(\mathcal O_{K_d}\). Put \(r=k+m-1,\ t=t_{\rm ch},\ n=n_V=r+t\). Assume \(t\ge r+5,\ k\ge5,\ r\le2k,\ n\ge8,\ g_0,g_1\ge(kb)^8\). At each real label the shift satisfies \(0<e_j\le2b\) and the cloud weights are in \([0,b]\). The auxiliary entries satisfy \(|b_l-e_j|\le b\).

Let \(\sigma:K\hookrightarrow\C\) agree with the distinguished embedding on \(K_d\). At \(\rho\), let \(\tau\) be the label induced by \(\sigma\), and put \(p_\rho=k+j-1,\ q_\rho=m-j\). With indices starting at zero, the two partitions are
\[
 a^+_{\sigma,\rho,i}=\kappa^{\rm big}_{j,\tau,i}\quad(0\le i<p_\rho),
 \qquad
 a^-_{\sigma,\rho,i}=-\kappa^{\rm small}_{j,\tau,q_\rho-1-i}\quad(0\le i<q_\rho).
\]
Fix the conjugate splitting characters, the prescribed lowest-type kernel and principal level at the marked places on the rank-\(r\) side. Let \(Y_\sigma\) be the multiplicity space of the prescribed discrete series in cusp forms at this level. For a system \(c\) of spherical eigenvalues off the marked places, put \(Y_\sigma(c)=\bigcap_A\ker(A-c(A)I)\), over spherical double-coset operators \(A\). There is a full lift \(T\) to \(Y_\sigma\), and its projection \(S=eT\) by the retained-type projector is the prescribed projected lift. For \(Y_\sigma(c)\ne0\) one has
\[
 TT^*z=l_c z\quad(z\in Y_\sigma(c)),\qquad
 \log l_c\ge\sum_\rho\log\bigl(\operatorname{Hook}(a^+_{\sigma,\rho})
 \operatorname{Hook}(a^-_{\sigma,\rho})\bigr)-8fn^2L',
 \qquad l_c>0.
\]
In particular, for \(H_j^+=\sum_\rho\log\operatorname{Hook}(a^+_{j\rho})\) and \(D_j=H_j^+-8fn^2L'\), one has \(l_c\ge e^{D_j}\) and \(\|Sx\|^2\ge e^{D_j}\|x\|^2\) for \(x\perp\ker S\).
\end{lemma}
\begin{proof}[Outline]
Construct the cuspidal lift using \ref{B5.07d}. Unfold its adjoint pairing by \ref{B8.12}, use \ref{B5.07e} and the hook formula \ref{B5.07} to compute the scalar, then bound the mass, Euler and formal-degree factors and compress to the retained types. The complete proof is in \S\ref{prf:B5.07y}.
\end{proof}

\begin{lemma}[Terminal and kernel budgets]\label{B5.07a}
Let \(\mathcal P\) be the set of types in the height table of \ref{B8.35}. Assume \(N=N_m\ge k+m,\ k,m\ge1,\ |\mathcal P|\le f,\ 1\le L\le L_a\) and the table's lower bounds on \(C_h=C_{\rm sz}\). Its terminal and layer entries satisfy \(A\le H_{P,u},\ \langle s_{jv}\rangle_\infty\le H_{K,uj}\). In the frames of \ref{B8.36},
\[
 \|P_u^{\rm term}\|_v^2\le e^A,\qquad
 \|K_i\|_v^2\le\exp\left(\sum_{j:i(j)=i}s_{jv}\right).
\]
\end{lemma}
\begin{proof}[Outline]
Sum the terminal table bounds with the level term kept exact and use \ref{B5.07c} in the selected evaluation. Average the kernel table using \ref{B5.07x} and \ref{B5.07}, then apply weighted Cauchy–Schwarz to the frame coefficients. The complete proof is in \S\ref{prf:B5.07a}.
\end{proof}

\begin{lemma}[The coefficient of a definite class]\label{B5.07fa}
Fix a layer of the reference tuple of \ref{B8.55}. Let \(E\) be the CM field of the unitary pair, \(L_0=\mathcal O_E^n\) the standard definite lattice, and \(r=k+m-1\) the rank on the other side. For a definite class \(a\), let \(L_a^h\) be the lattice of its adelic representative. The lattice used in the theta section is \(L_a=\overline{L_a^h}\), with conjugation in the standard coordinates. Choose its clearing integer \(N_a>0\), so that \(N_aL_0\subset L_a\cap L_0\), and put
\[
 m_a=[L_a\cap L_0:N_aL_0]^r,\qquad
 Z_a=[L_a:N_aL_0]^r,\qquad I_a=[L_a:L_a\cap L_0]^r.
\]
At each moving split pair choose the half specified by the representative. Let \(d_a\) be its determinant on those halves, extended by one elsewhere. Write \(\epsilon_a\) for the coefficient of this class in the algebraic theta section, before the lattice average and the class measure. If \(\widetilde\sigma\) is the embedding of the comparison field over \(\sigma\), and \(\chi_{W,\sigma}\) is the splitting character on the target side, then
\[
 \widetilde\sigma(\epsilon_a)\,\frac{Z_a}{m_a}
       =\chi_{W,\sigma}(d_a)^{-1}\sqrt{I_a}.
\]
\end{lemma}
\begin{proof}[Outline]
Use split elementary divisors and the odd rank from \ref{B8.58}--\ref{B8.59} to compute the lattice index and the rational factor in the normalization of \ref{B8.55}. Evaluate the determinant value under the embedding and substitute into the class coefficient. The complete proof is in \S\ref{prf:B5.07fa}.
\end{proof}

\begin{lemma}[The rational projected lift and its analytic realizations]\label{B5.07f}
At the reference tuple of \ref{B8.55}, fix layer \(j\), with zero-based index \(j-1\). Let \(X\) be its \(K\)-space of definite class forms and \(Y=K^{n_W}\) the coordinate space of its rational cusp frame. Write \(j_\sigma,J_\sigma\) for their semilinear realizations, and use the subscript \(0\) for the distinguished embedding. With base phase one there is a unique \(K\)-linear map \(\theta:X\to Y\) satisfying \(J_0\theta=S_0j_0\). At each embedding \(\sigma\) agreeing with the distinguished embedding on \(K_d\), there are a unitary component multiplication \(U_\sigma\), a central level translation \(g_\sigma\), and a scalar \(|\xi_\sigma|=1\) such that
\begin{equation}\label{eq:B5-conjugate-lift}
             U_\sigma J_\sigma\theta x
                  =\xi_\sigma g_\sigma S_\sigma j_\sigma x
                   \qquad(x\in X).
\end{equation}
Here \(S_\sigma\) is the projected lift for the conjugate splitting characters and the prescribed extremal tensor.
\end{lemma}
\begin{proof}[Outline]
Project the rational class expansion of \ref{B8.55} using \ref{B5.07y}. At each embedding, match Fourier weights by central translation, vacuum coefficients by \ref{B5.07fa}, and extremal polynomials up to phase by Schur's lemma, then apply the projector. The complete proof is in \S\ref{prf:B5.07f}.
\end{proof}

\begin{lemma}[The two Gram forms on a rational spectral part]\label{B5.07g}
Fix a layer \(1\le j\le m\) at the reference tuple \(\mu_0\) of the primitive realization in \ref{B8.38}. Use the phase-one lift \(\theta:X\to Y\) of \ref{B5.07f} and the sizes of \ref{B5.07y}. Write \(Z_j=K^{d_j}\) for the reference layer coordinates and \(B:Z_j\to X\) for their inclusion as definite class forms by the orthogonal section of the layer quotient. Let \(N\) be the order of the reference level subgroup.

A spherical Hecke system is carried by \(p\in I\) if its dual Hecke space on the initial row is nonzero and the reference row adjoint takes that space into the complex span of the initial inert type indexed by \(p\). Let \(W_{jp}\) be the sum, at the distinguished embedding, of the retained layer spaces of the systems carried by \(p\). Put
\[
 Z_{jp}=\{x\in Z_j:S_0j_0Bx\in W_{jp}\}.
\]
Here \(j_0\) realizes a definite class form and \(S_0\) is the projected lift. In particular \(Z_{jp}\) is a \(K\)-space.

For \(v\ne v_0\), let \(\sigma\) represent \(v\). Let \(H\) be the matrix of the definite class pairing in a \(K\)-basis of \(X\), with each class weighted by the reciprocal of its stabilizer order. Define
\[
 M_v(x,x')=\frac1N\,\overline{\sigma(Bx)}^{\,t}\sigma(H)\sigma(Bx'),
 \qquad
 y_v(x)=\frac1{\sqrt N}\,J_\sigma\theta Bx.
\]
The paired row at \(j,p,v\) consists of the two forms
\(\langle y_v(x),y_v(x')\rangle\) and \(M_v(x,x')\), with
\(x\in Z_{jp}\) and \(x'\in Z_j\). Call the part occurring at \(v\) if
\(M_v(x,x)\ne0\) for some \(x\in Z_{jp}\).

Suppose that \(j_\tau BZ_j\subseteq(\ker S_\tau)^\perp\) for each embedding \(\tau\) agreeing with the distinguished embedding on \(K_d\). On an occurring part, take an ordered \(K\)-basis of \(Z_{jp}\) and let \(b\) be its first member with \(M_v(b,b)\ne0\). Then
\[
 \lambda_{jpv}=\frac{\|y_v(b)\|^2}{M_v(b,b)}>0,\qquad
 \langle y_v(x),y_v(x')\rangle=\lambda_{jpv}M_v(x,x')
 \quad(x\in Z_{jp},\ x'\in Z_j).
\]
This determines the same scalar for every choice of basis. On an empty part both forms vanish on \(Z_{jp}\times Z_j\), and the selected scalar is \(\lambda_{jpv}=\exp L_{6,j}\) at the fixed scale of \ref{B5.07}.
\end{lemma}
\begin{proof}[Outline]
Use \ref{B8.38} to isolate the rational Hecke part. The inert vacuum transfer \ref{B8.14} and the split triangular relations give a common target Hecke system, so the scalar of \ref{B5.07y} relates the two Gram forms by adjunction. The complete proof is in \S\ref{prf:B5.07g}.
\end{proof}

\begin{lemma}[Subtraction on the base spectrum]\label{B5.07b}
Use a normalized primitive realization with the separated types and level-one Hecke decomposition of \ref{B8.08}. Assume the hypotheses of \ref{B5.07y} and the class expansion and level action of \ref{B8.55}. Use the definite block conditions of Definition~\ref{def:p4c6-definite-blocks}. Equip the totally complex coefficient field \(K\) with its ring involution \(a\mapsto\bar a\) compatible with complex conjugation at the distinguished infinite place \(v_0\) (so \(\sigma_0(\bar a)=\overline{\sigma_0(a)}\) at the base embedding \(\sigma_0\), and \(M^*=\bar M^{\,t}\)), and let \(K_0=\{a\in K:\tau(\bar a)=\overline{\tau(a)}\text{ for all }\tau:K\hookrightarrow\C\}\) be its maximal CM subfield. For the surjective reference block quotient \(P\), let \(H_1\) be the block Gram matrix, \(E_1\) its orthogonal section and \(G_1=E_1^*H_1E_1\). Require \(H_1,G_1\) to have entries in \(K_0\), with \(\sigma(H_1)>0\) at every embedding. At the representative \(\sigma\) of each infinite place \(v\), let \(X_v\) be the orthonormal frame of the quotient metric induced by \(\sigma(P)\) from \(\sigma(H_1)\). Require \((X_v^{-1})^*X_v^{-1}=\sigma(G_1)\). Use the conjugate pairing on these rational quotient coordinates. At an embedding representing a nonbase place, the realized squared norm on a rational vector in an occurring block is its primitive scalar times the corresponding diagonal value of that pairing. Take the layer sizes from \ref{B5.07y}. Take phase one for the rational reference lift and \(C_{\rm sz}\ge8\). Use the spaces \(W_{jp}\) of \ref{B5.07g}. At \(v_0\) take the prescribed scalar of \(S_0S_0^*\) on \(W_{jp}\) when this space is nonzero and put \(\lambda_{jpv_0}=e^{L_{6,j}}\) when it is zero. At \(v\ne v_0\) take \(\lambda_{jpv}\) from the rational part and its Gram forms as in \ref{B5.07g}. Then
\[
 L_{6,j}\le\log\lambda_{jpv},\qquad
 \sum_jL_{6,j}\le\min_{p\in I}\sum_j\log\lambda_{jpv}
\]
at every infinite \(v\). The left sum equals \(\sum_jH_{X,j}\) in \ref{B5.08}.
\end{lemma}
\begin{proof}[Outline]
Compares \(L_{6,j}\) of \ref{B5.07} with \(D_j\) at \(v_0\). Uses Definition~\ref{def:p4c6-definite-blocks} and the quotient frame to show \(\tau(B^*)=\tau(B)^*\), transports kernel orthogonality to all \(v\ne v_0\) for \(\lambda_{jpv}\ge e^{D_j}\), and sums over \(j\). The complete proof is in \S\ref{prf:B5.07b}.
\end{proof}

\begin{lemma}[Metric of a CM generator]\label{B5.07x}
Let \(K_d=\Q(\sqrt{-d})\) have discriminant \(-d<-4\). Take a nonzero integral \(\mathcal O_{K_d}\)-ideal \(\mathfrak c\) of norm prime to \(2d\), and put \(\Lambda=\delta\mathfrak c\), with \(\delta^{12}=\operatorname{disc}(\mathfrak c)\) its Weierstrass discriminant. Let \(M\subseteq\C\) be a number field containing \(g_2(\Lambda),g_3(\Lambda)\) and the prescribed images of \(E\). If \(s_\tau\in M^\times\) generates the host differential ideal at \(\tau\) over \(\mathcal O_M\), then
\[
 \frac1{[M:\Q]}\sum_{\sigma:M\hookrightarrow\C}
 \log\frac{\operatorname{area}(\Lambda_\sigma)}{\pi|\sigma s_\tau|^2}
 =\gamma+\log(\pi/v_E)-D_1,
\]
where \(\Lambda_\sigma\) has invariants \(\sigma g_2(\Lambda),\sigma g_3(\Lambda)\).
\end{lemma}
\begin{proof}[Outline]
Use Chowla--Selberg and the equal class multiplicities of \ref{B1.15a} to average log areas. Remove the unit term by the product formula, subtract twice the mean log norm of the differential generator, and use the relative discriminant identity for \ref{B5.07}. The complete proof is in \S\ref{prf:B5.07x}.
\end{proof}

\section{Combining the divisor and height budgets}

Let \(J,\mathcal P\) be finite with \(\#J=m,\ \#\mathcal P\le f\), and \(\mathcal U\subseteq\mathcal P\) of size \(n_T\). Use \(x>1,\ \varepsilon>0,\ L=x^\varepsilon,\ b=e^L,\ a=xb,\ B_0=fmbL\), and take \(r_d\ge1,\ \log r_d,\log x\le L,\ L_a=\log(ar_d)\). Assume
\[
\begin{gathered}
 N_c\ge0,\quad m,f,k\ge1,\quad k+m\le n_*,\\
 S_c,t_c,\eta,b',w'',\lambda_0,n_V,L_p\ge0,\quad n_V\le n_*,\\
 m\lambda_0,\ n_*^2,\ m^3n_*^2L,\ m\eta n_*L,\ kn_*\le b,\\
 kS_c,\ m(t_c+1)\eta\le w'',\quad n_*w''\le b,\quad b'n_*,\ mkb'\le a,\\
 1\le C_0\le c_e,\quad 0\le C_h\le C_0,\quad C_hL_p\le C_0L_a.
\end{gathered}
\]
Fix \(r\in\N\), real \(\log T,\log\Omega\) with \(|\log\Omega|\le L_a\), and weights \(a_\rho\). Use \(\varphi\) from \ref{B5.07}. For \(j,\rho\) take row entries \(\nu_{j\rho i}\) for \(0\le i<r\) and log hooks \(\log H_{j\rho}\). Put \(I_\rho=\mathbf1_{\mathcal U}(\rho)\). Assume \(0\le w_{j\rho}\le2b,\ |\beta_j-k|\le m\) with \(\beta_j=(k+j-1)-(m-j)\) in arithmetic layer numbering, and
\[
\begin{aligned}
 \sum_i\varphi(\nu_{j\rho i})-\log H_{j\rho}&\le I_\rho\varphi(a_\rho)+e_{j\rho},\\
 \sum_i\nu_{j\rho i}&=I_\rho a_\rho+kw_{j\rho}+\Delta_{j\rho},\\
 e_{j\rho}&\le C_0(w_{j\rho}+S_c+n_*^2+m(t_c+1)\eta)L_a,\\
 |\Delta_{j\rho}|&\le C_0(w_{j\rho}+k(n_*+S_c)+m(t_c+1)\eta+I_\rho kb').
\end{aligned}
\]
Set
\[
\begin{aligned}
 H_P&=-\sum_{\rho\in\mathcal U}a_\rho\log T
 +C_h(n_Tb'n_*+f(m\eta+S_c)n_*+fn_*^2)L_a,\\
 H_{B,j}&=\sum_{\rho,i}(\varphi(\nu_{j\rho i})-\nu_{j\rho i}\log\Omega)
 +C_hfn_*^2(L_a+m^3L),\\
 H_{X,j}&=\sum_\rho\log H_{j\rho}-C_hfn_V^2L_p,\qquad
 H^*=H_P+\sum_j(H_{B,j}-H_{X,j}),\\
 \mathcal D(T)&=\sum_{\rho\in\mathcal U}
 \bigl(m(\varphi(a_\rho)-a_\rho\log\Omega)-a_\rho\log T\bigr).
\end{aligned}
\]

\begin{lemma}[An upper bound for the extension weights]\label{B5.08d}
At a selected point with the weights of Lemma~\ref{B4.02b}, write \(e_q(\rho)\) for the extension weight at position \(\rho\), \(0\le q<m\). For all sufficiently large \(x\), uniformly in the position,
\[
 e_q(\rho)\le e_{m-1}(\rho)\le 2b .
\]
\end{lemma}
\begin{proof}[Outline]
Use the increasing weights in \ref{B4.02b} and bound their last term at each kind of position. The totient bound on roots of unity and the size estimates \ref{B8.59} and \ref{B8.35a} absorb the thick-position adjustment into half the step. The complete proof is in \S\ref{prf:B5.08d}.
\end{proof}

\begin{lemma}[Wordwise numerical conditions]\label{B5.08b}
Write \(t=t_{\rm ch}\), \(N=N_m\), \(n_V=k+m-1+t\), \(r_d=\operatorname{rd}E\), \(A_*=ar_d\), \(L_a=\log A_*\), \(L_p=\log(br_dg_0g_1n_V)\). For \(1\le j\le m\) set
\[
 s_j=\{k+(2j-1)(t+1)\}/2,\qquad B=m(t+1)\eta .
\]
For a fixed word use the rows of \ref{B4.02a}. The wordwise numerical conditions on \(N_c,\lambda_0,C_0,c_e,C_h,S_w,C_1\) are the following scalar and coordinate bounds:
\[
\begin{gathered}
 \log r_d,\log x\le L,\quad |\log\Omega|\le L_a,\quad
 N_c,\lambda_0\ge0,\quad m\lambda_0,N^2,m^3N^2L,m\eta NL,kN\le b,\\
 kS_c,B\le w'',\quad Nw''\le b,\quad b'N,mkb'\le a,\\
 1\le C_0\le c_e,\quad 0\le C_h\le C_0,\quad C_hL_p\le C_0L_a,\quad
 2\le C_1,\quad 2S_w\le C_1S_c,\quad C_1+1\le C_0,\\
 s_j\le w_{j\rho}\le2b,\quad w_{j\rho}+S_w\le A_*,\quad
 a_s+w_{j,(0,s)}\le A_* \quad(s\in TT),\\
 \nu_{uj\rho i}\ge0\quad(i<N_{j-1}),\qquad
 \sum_{k\le i<N_{j-1}}\nu_{uj\rho i}\le B,\qquad
 U_i>v_i>U_{i+1}\quad(i<k-1).
\end{gathered}
\]
Here and below the \(U_i,v_i\) are taken at the position in question. Off the marked set the endpoint conditions are
\[
 U_{k-1}+(k-1)/2\ge0,\qquad U_0-(k-1)/2\le S_w .
\]
At a mark they are \(a_s,b_s\ge0\),
\(v_{k-2}+(k-2)/2\ge0\), and \(v_0-(k-2)/2\le S_w\).

These conditions hold for every word of the selected family for all sufficiently large \(x\). Take \(C_h=C_{\rm sz}\ge0\), \(S_w=(1+|C_c|)S_c\), \(C_1=2(1+|C_c|)\), where \(C_c=2\) in \ref{B8.01k}. Take \(C_0\) sufficiently large in terms of the fixed constants and \(c_e\ge C_0\). Here
\[
 N_c=|T_2|kbL^2/n^2,\qquad \lambda_0=\log M,
\]
with \(T_2\) the thick host labels, \(n=|A|\), and \(M\) the order of the finite row group.
\end{lemma}

\begin{proof}[Outline]
Use the rank and size bounds \ref{B8.59} and \ref{B8.35a} and count matrices to control the scalar conditions. Use \ref{B3.47} and \ref{B5.08d} for the weight range and \ref{B4.02a} and \ref{B4.02b} for the degree, tail and interlacing bounds. The complete proof is in \S\ref{prf:B5.08b}.
\end{proof}

\begin{lemma}[Cancellation of the kernel factorials against the row hook]\label{B5.08c}
Fix a word and impose the numerical conditions of \ref{B5.08b}. At layer \(j\) and position \(\rho\), put \(I_\rho=1\) if \(\rho=(0,s)\) with \(s\in TT\), and \(I_\rho=0\) otherwise. Write \(a_\rho=a_s\) at a mark and \(a_\rho=0\) elsewhere. Set \(\alpha_i=\kappa^{\rm big}_{j\rho i}\), \(H_{j\rho}=\operatorname{Hook}(\alpha)\), and abbreviate \(\nu_{uj\rho i}\) to \(\nu_i\). Define the two real remainders by
\[
 e_{j\rho}=\sum_{i<N_{j-1}}\varphi(\nu_i)-\log H_{j\rho}
                 -I_\rho\varphi(a_\rho),\qquad
 \Delta_{j\rho}=\sum_{i<N_{j-1}}\nu_i-I_\rho a_\rho-kw_{j\rho}.
\]
They satisfy
\[
 \begin{aligned}
 e_{j\rho}&\le C_0(w_{j\rho}+S_c+N^2+B)L_a,\\
 |\Delta_{j\rho}|&\le C_0\bigl(w_{j\rho}+k(N+S_c)+B+I_\rho kb'\bigr).
 \end{aligned}
\]
Extending the degree tuple by zeros up to length \(N\) preserves both sums. Thus all four row hypotheses preceding \ref{B5.08} hold for the selected words with \(r=n_*=N\).
\end{lemma}

\begin{proof}[Outline]
Remove the shifts of \ref{B4.02a} and bound the hook product of \ref{B5.07} by row factorials. Telescope the paired differences, retain the marked leading term, and sum the shifted degrees for the second remainder using \ref{B5.08b}. The complete proof is in \S\ref{prf:B5.08c}.
\end{proof}

\begin{lemma}[The divisor slope and the height budget]\label{B5.08}
For the general arrays specified above, suppose that the four row bounds hold. For the selected arrays use \ref{B5.08b} and \ref{B5.08c}. Suppose also that
\[
\begin{gathered}
 \widetilde\mu_C\le\mu_0+L_{\rm pass},\qquad
 \mu_0\le\sum_{j,\rho}\beta_jw_{j\rho}\log\Omega+2F_*+E_b,\\
 L_{\rm pass}\le1+C_0m\lambda_0,\qquad E_b\le c_emL(B_0+N_c).
\end{gathered}
\]
Then
\[
 \widetilde\mu_C+H^*-2F_*\le\mathrm{rhs}
 :=\mathcal D(T)+(50C_0+10)n_TaL_a+60c_emL(B_0+N_c).
\]
\end{lemma}
\begin{proof}[Outline]
Insert the degree identity and the remainders of \ref{B5.08c}, collect the marked and terminal terms into \(\mathcal D(T)\), and sum the residual allowances. Use \ref{B5.08x} and \eqref{eq:p4c6-row-budgets} for the reference divisor bound. The complete proof is in \S\ref{prf:B5.08}.
\end{proof}

\begin{lemma}[The two local indices of a change of row type]\label{B5.08a}
Use the reference row system of \ref{B4.55}, with base type \(t_0\) and glued lattice \(R_0\). Write \(L_t=\iota_t\Lambda_t\) for the image of the quotient lattice, \(A_{jt}=\theta_{jt}T_{jt}\) for the layer image in the space of \(M_j\), and \(\sharp\) for Hermitian duality. For commensurable lattices put
\[
 [A:B]_w=\operatorname{length}_w(A/(A\cap B))
 -\operatorname{length}_w(B/(A\cap B)).
\]
Here the totally complex field \(K\) carries the ring involution \(a\mapsto\bar a\) satisfying \(\sigma_0(\bar a)=\overline{\sigma_0(a)}\) at the distinguished infinite place \(v_0\). Let \(\bar w\) be the finite place of \(K\) pulled back from \(w\) under \(a\mapsto\bar a\) (so \(v_{\bar w}(a)=v_w(\bar a)\)), and let \(G_{t,w}\in K\) represent the quotient Gram determinant at \(v_0\) in a semilocal basis at \(w,\bar w\). Define
\[
 C_w(t)=v_w(G_{t,w})+\sum_j\bigl([M_j:A_{jt}]_w+[M_j:A_{jt}]_{\bar w}
 +[T_{jt}:T_{jt}^\sharp]_w\bigr).
\]
For the type \(t_w\) of \ref{B4.36a},
\[
 [L_{t_0}:R_0]_w+[L_{t_0}:R_0]_{\bar w}=C_w(t_w)-C_w(t_0).
\]
\end{lemma}
\begin{proof}[Outline]
Express the local indices by determinant valuations in semilocal bases. Factor the row determinant ratio into layer ratios and a common power by \ref{def:p4c6-core}, then divide the two type identities and take valuations at conjugate places. The complete proof is in \S\ref{prf:B5.08a}.
\end{proof}

\begin{lemma}[The reference row divisor]\label{B5.08x}
Take the well-formed reference side of \ref{B4.55} and the selector of \ref{B4.36a}. Use the untwisted identities with the reference-row clauses of \ref{B8.38b}. Take the bounds from \eqref{eq:p4c6-gram-inputs}. Write \(D_C=\operatorname{rk}R_0>0,\ \bar\kappa_0=\langle\log\kappa_{\rm base}\rangle_\infty\). Let \(\operatorname{lcv}\) be log covolume in the slope normalization and \(l_{3,j}\) bound \(\operatorname{lcv}(M_j)/D_C\). With \(F_*,L_G,E_{1,j},E_{2,j},E_3\) of \ref{B4.55},
\[
 \mu(R_0^\vee)-\bar\kappa_0
 \le\sum_jl_{3,j}+2F_*+L_G+\sum_j(E_{1,j}+E_{2,j})+E_3 .
\]
\end{lemma}
\begin{proof}[Outline]
Use \ref{B5.08a}, exterior Cauchy--Schwarz, and the product formula to cancel the base terms. Bound the rest by \ref{B8.38b} and \eqref{eq:p4c6-gram-inputs} using the denominator, cut, and weights of \ref{B4.55}, \ref{B8.40a}, and \ref{B4.52}, then sum by parts. The complete proof is in \S\ref{prf:B5.08x}.
\end{proof}

\section{The four conclusions of division}

\begin{theorem}[The division package]\label{B5.09}
Use the coordinate lattices and valuation bounds of \ref{B4.55} for \(N_u^0,N'_u\), with \(s_{uw}=\lfloor\sigma_w\rfloor\) of finite support. Take \(H^*,\mathrm{rhs}\) from \ref{B5.08} and the factorization and norm bounds of \ref{B5.07}, with mean log budget after spectral subtraction and before metric correction at most \(H^*\). Put \(\bar\kappa=\langle\log\kappa_v\rangle_\infty\) and \(\mu_0=\mu(R_0^\vee)-\langle\log\kappa_{\rm base}\rangle_\infty\), bounded above by \(B_*\) from \ref{B5.08x}. Use the size and \(\mu_0\) bounds of \ref{B5.08} and assume
\[
\begin{gathered}
 \delta_0\ge0,\quad Q_0\ge1,\quad 3\delta_0Q_0\le1,\quad
 B_*+H_0+28\lambda\le Q_0,\quad 0\le\lambda\le\lambda_0,\quad C_0\ge28,\\
 \mu(RR^\vee)-\bar\kappa\le\mu_0+3\delta_0(\mu_0+H_0+28\lambda)+28\lambda.
\end{gathered}
\]
The four conclusions of \ref{B4.55} hold for \(B_u,N'_u,RR\) with \(H=H^*\) and \(R_*=\mathrm{rhs}\).
\end{theorem}
\begin{proof}[Outline]
Apply \ref{B5.08} with \(\widetilde{\mu}_C=\mu_C-\bar{\kappa}\) after bounding passage cost by \(1+C_0m\lambda_0\), and shift and exponentiate \(z_v\) of \ref{B5.07} to mean \(H^*-\bar{\kappa}\). Then apply \ref{B4.55}, using \ref{B4.38} and \ref{B5.03} for \(\sigma_w=s_{uw}\). The complete proof is in \S\ref{prf:B5.09}.
\end{proof}

\begin{lemma}[The unit budget for row passage]\label{B5.09a}
Use the selected family's error bounds \eqref{eq:p4c6-row-budgets}. Put \(g=\min(g_0,g_1),\ r=k+m-1\), with \(M\) the finite row group order, and assume \(4fr\le g\). Set
\[
 A_3=\sum_{j,\rho}\beta_jw_{j\rho}\log\Omega,\quad
 B=A_3+\Fo+c_emL(B_0+N_c),\quad Q_0=\max\{1,B+x^{C_H}b+28\log M\}.
\]
For each retained character \(\chi\), let \(n(\chi)\) be the multiplicity of the selected subgroup character and put \(\bar n=\prod_wq_w^k\), with \(q_w\) the residue size at a selected place. For fixed family constants and sufficiently large \(x\), \(|n(\chi)/\bar n-1|\le\delta_0=8fr/g\). With \(H_0=x^{C_H}b,\ \lambda=\log M\), one has
\[
 B_*+H_0+28\lambda\le Q_0,\qquad 1\le Q_0\le3b^2,\qquad 3\delta_0Q_0\le1.
\]
For \(P=1+f+k+m+t_{\rm ch}+N_m+L\), one has \(\Fo\le256P^9b\).
\end{lemma}
\begin{proof}[Outline]
Use \ref{B7.47b} and \ref{B7.48} to bound multiplicities by subset counts and \ref{B4.55b} for the row budget. Bound passage by the sizes, inert terms, and grouped prime sum of \ref{B5.08}, \ref{B4.02}, \ref{B4.01}, \ref{B5.06}, and \ref{B4.52}. The complete proof is in \S\ref{prf:B5.09a}.
\end{proof}

\begin{definition}[Pointwise division conditions]\label{B5.10b}
Use the selected row lattice \(RR\) in the row space \(W\), the divided word lattices \(N'_u\) and the periods \(B_u\) of \ref{B4.56}. Put \(d_u=\operatorname{rk}N'_u>0\), \(\mu_C=\mu(RR^\vee)\), \(S_a=\sum_{t\in TT}a_t\), and
\[
\begin{aligned}
 s_{\rm fr}&=\frac{\sum_u d_u(\mu(N'_u)-h_u)}{\sum_u d_u},&
 H_u&=h_u+S_a\log T_u,\\
 \mathcal D(T)&=m\sum_{t\in TT}\bigl(\varphi(a_t)-a_t\log\Omega\bigr)
                 -S_a\log T,&
 E_{\rm div}&=c_{\rm Div}
 \bigl(n_TaL_a+m(B_0+N_c)x^{e_{\rm Div}\varepsilon}\bigr).
\end{aligned}
\]
Here the heights and \(\varphi\) are as in \ref{B5.07}. The pointwise division conditions with constants \(c_{\rm Div},e_{\rm Div},C_H\) are the following.

The lattice \(RR\) is finitely generated over \(\mathcal O_K\) and spans \(W\). The frame satisfies \(s_{\rm fr}+\mu_{\max}(RR)\le x^{C_H}b\). For each word \(u\) and layer \(j\), the raw gain \(g_{jw}\) and loss \(E_{ujw}\) vanish at every \(w\) with \(p_w>X''\), \(p_w\notin\{T_u,g_0,g_1,d\}\) and \(p_w\nmid D_{\Gamma,u}\).

The complex period on rational vectors is the distinguished embedding of the \(K\)-valued pairing. Write \(\mathcal B_u(\phi,F)\) for the inner product of the realized row vector \(F\) with the period image of the realized column vector \(\phi\), in that order. There is a common row isometry \(V\) and, for each \(u\), a column isometry \(U_u\) and a scalar \(a_u\in\C^\times\) such that \(B_u^\C(\phi,F)=a_u\mathcal B_u(U_u\phi,VF)\).

For every word \(u\),
\[
 B_u(N'_u)\subseteq RR^\vee,\qquad
 B_u\ne0\ \Longrightarrow\ h(B_u)\le h_u,\qquad
 \mu_C+h_u-\Fo\le\mathcal D(T_u)+E_{\rm div}.
\]
Finally, for every word \(u\), layer \(j\), and second word \(u'\),
\[
 \left|\sum_{w\in S}c_wE_{ujw}-S_a\log(g_0g_1)\right|\le c_{\rm Div}B_0,
 \qquad
 |H_u-H_{u'}|\le c_{\rm Div}(n_TaL_a+mB_0).
\]
\end{definition}

\begin{theorem}[Division data for the actual period family]\label{B5.10}
Fix a selected family in the range of \(\varepsilon\) and with the ordered thresholds of \ref{B5.10a}. Use full divided column lattices of positive rank and a finitely generated row lattice spanning \(W\). Assume the rational row and column isometries of \ref{B4.56} carry the coordinate lattices onto these lattices and preserve their pairings. Assume the rationality and complex isometries for the period in Definition~\ref{B5.10b}, with a common row isometry for all words.

Use the finite support and fixed weight tuples of \ref{B4.52a} and \ref{B4.55b}, and the coordinate valuation and lattice hypotheses of \ref{B4.55}. For every word use its selected level \(T_u\), the definite tables and frame factorizations of \ref{B8.35}--\ref{B8.36}, and the primitive reference realization with the Hecke, metric and size hypotheses of \ref{B5.07a} and \ref{B5.07b}. Give it the height allowance of \ref{B5.07} on the same positions and layers. Assume the wordwise size and interlacing hypotheses of \ref{B5.08b}--\ref{B5.08c}, the reference covolume and Gram bounds of \ref{B5.08x}, the paired-index and cut hypotheses of \ref{B8.38b} and \ref{B8.40a}, and the numerical passage hypotheses of \ref{B5.09a}. Assume also the row-passage hypotheses of \ref{B4.55a} (the finite group action, scalar endomorphisms with distinct characters, and intermediate lattices \(h_iP_i\subseteq R_i\subseteq P_i\)) and the base maximum-slope bound \(U\ge\mu_{\max}(R_0)\) of \ref{B8.23}, which together with the multiplicity bound of \ref{B5.09a} and the choice of \(C_H\) in \ref{B5.10a} supply both the maximum-slope bound for \(RR\) in \ref{B5.10a} and the passage inequality \(\mu(RR^\vee)-\bar\kappa\le\mu_0+3\delta_0(\mu_0+H_0+28\lambda)+28\lambda\) of \ref{B5.09}.

For the frame clause via \ref{B5.10a}, assume the functional, row, level and stabilizer hypotheses of \ref{B8.22}, \ref{B8.23} and \ref{B8.38}, together with the hypotheses of \ref{B5.06} and \ref{B5.12}: namely, the local comparisons and containments of \ref{B5.04}--\ref{B5.05a} and \ref{B5.20a}, the integrality and colength conditions \(\tau_w\in\N\), \(\mathrm{io}_w\ge0\), \(\tau_w=0\) on \(S_T\), and \(\ell^{\rm col}_{uw}\le d_M(\tau_w+\mathrm{io}_w)\) off \(S_T\), the bounds on \(\mathcal T\), \(T_u\), \(P_u\) and \(R_u\) in \ref{B5.06}, the counting window of \ref{B5.11}, the column models beyond a common threshold, and \(|D_1|\le x^\varepsilon\) eventually. Take the hypotheses of \ref{B5.13} (the seven error conditions of \ref{B4.03}, the eight clauses of \ref{B5.12k}, and equality of stored and evaluated losses) and the size bounds of \ref{B5.14}. All these data belong to the same selected member.

Let \(c_E\) be the constant of \ref{B5.13}, choose \(C_H\) as in \ref{B5.10a}, and take \(e_{\rm Div}\ge1\) and \(c_{\rm Div}\ge\max(50C_0+10,60c_e,c_E,20+6/\varepsilon)\). Beyond a common cutoff, every selected member satisfies the pointwise division conditions of Definition~\ref{B5.10b}.
\end{theorem}
\begin{proof}[Outline]
Applies \ref{B5.09} to the coordinate pairing via \ref{B8.35}, \ref{B5.07a}, \ref{B5.07b}, and \ref{B5.08x}, transferring bounds through \ref{B4.56}. Supplies the remaining clauses of \ref{B5.10b} from \ref{B4.52a}, \ref{B5.10a}, \ref{B5.13}, and \ref{B5.14}. The complete proof is in \S\ref{prf:B5.10}.
\end{proof}

\begin{lemma}[Choosing the height exponent]\label{B5.10a}
Take \(0<\varepsilon\le\varepsilon_1\) in the chapter's range for sufficiently small \(\varepsilon_1>0\). Use the selected class with the functional and row hypotheses of \ref{B8.22}, \ref{B8.23} and \ref{B8.38}, and the level and stabilizer bounds. Apply \ref{B5.06}, \ref{B5.07}, \ref{B5.12} and \ref{B5.13} with full divided lattices of positive rank. Fix the outer constants \(c_{\rm out}\), including the chain bound \(C_t\). Set \(M=\max(C_t,5)+3\). For the exponent \(F_0=F_0(\varepsilon,c_{\rm out})\) below, take \(C_H\ge1+\max\{4M+1,9M+2,F_0,9M+1\}\). Take the thresholds of \ref{B8.39} for the row and size estimates. Assume every preliminary constant meets the preliminary lower bound, every middle constant (including \(C_H\)) meets the middle bound evaluated at \(c_{\rm out}\), and every final constant meets the final bound evaluated at the concatenated outer and middle lists. Then there is an \(x_0>x_1\) such that \(s_{\rm fr}+\mu_{\max}(RR)\le x^{C_H}b\) for every \(x\ge x_0\) and every \(i\in I_x\).
\end{lemma}
\begin{proof}[Outline]
Bound row slopes by \ref{B8.23} and \ref{B4.55a}, and divided slopes by \ref{B5.06} and \ref{B5.12}, summing \ref{B3.51} and \ref{B3.52} via \ref{B4.52a} and \ref{B4.52}. Average with heights from \eqref{eq:B8-five-size}, \ref{B8.35}--\ref{B8.37}, and \ref{B4.03a}. The complete proof is in \S\ref{prf:B5.10a}.
\end{proof}

\section{The auxiliary prime and the original lattice}

Use the starting volume of Lemma~\ref{B8.51b} and the genus masses of \ref{B8.48}. Set
\[
 d_u^\star=\operatorname{mass}_k\,\operatorname{idx}_{\Pi}(g_0,g_1)
 \prod_\rho\dim V(u,\rho),\qquad
 D^\star=\operatorname{vol}(Y_0)\operatorname{idx}_{\Sigma}(g_0,g_1)
 \prod_\rho d_{\Sigma,\rho}.
\]
The indices are the indicated finite-level indices. Write \(V(u,\rho)\) for the column representation and \(d_{\Sigma,\rho}\) for the row spectral degree. Let \(C_\Gamma\) be the family's size constant for the count and put \(\Gamma_u=\exp(C_\Gamma m^5fL)\), the same for every word. Set \(X_u=(D^\star/d_u^\star)^{1/(fk)}\), \(T_u^\star=X_u\Gamma_u^{-2/(fk)}\) and \(\ell_g=\log(g_0g_1)\).

\begin{lemma}[Separated parameters for the count]\label{B5.11c}
At all sufficiently large selected members put \(C_Y=C_0=\max(5,C_{YL}+1)\), \(r=k+m-1\), \(v_E=\operatorname{rd}E/2\). Use the zero-based indices of \ref{B4.02b}. The scalar conditions for the count are
\[
 \pi\le v_E,\quad \log(2v_E),\log(2k)\le L,\quad
 f,k\le x^{C_0},\quad m^2 L\le k,\quad44\le m\le\min(k,x),
 \quad(kb)^8\le g_i,\quad \log g_i\le C_YL .
\]
The interlacing row has spacing \(4\delta\), \(\delta\ge1\), and belongs to \([0,S_c]\). The interior entries of \(U\) are the bit choices in \ref{B4.02b}. For all words \(U_0-U_{k-1}\le2a\), \(U_0-v_0\ge S_c\). Off the marks \(U_0-v_0\le2S_c\) and \(U_{k-1}=v_{k-2}-1/2\). At the marks \(U_0\le2a\) and \(-a\le U_{k-1}\le-S_c\). Finally,
\[
 \begin{gathered}
 0<b_q<e_q\le (1+1/x)b_q,\qquad e_{m-1}\le2b,\qquad
 S_c\le b_0/x,\\
 e_q<b_{q+1},\qquad b_{q+1}-b_q\ge r^{19}f^5L/C_0
 \quad(0\le q<m-1).
 \end{gathered}
\]
\end{lemma}

\begin{proof}[Outline]
Use \ref{B8.35a} and \ref{B2.51} for size and prime bounds and \ref{B4.02b} for the word endpoints. Compare the step with the base for the relative weight bound and with the rank scales for the gap, using \ref{B5.08d} for the final weight. The complete proof is in \S\ref{prf:B5.11c}.
\end{proof}

\begin{theorem}[Three counting bounds for the auxiliary primes]\label{B5.11}
Suppose \(1\le L\le L_a,\ k\ge3,\ f\ge1,\ m\ge1\), \(\Gamma_u>1\) is common to the words with \(\log\Gamma_u/(fk)\le1\), and \(T_u^\star\le T_u\le(1+1/(fk))T_u^\star\). Write \(\mathcal W=\sum_{j=1}^m f^{-1}\sum_\rho\log(w_{j\rho}v_E/\pi)\). For a common \(C\) satisfying
\begin{equation}\label{eq:B5-count-torus}
 \mathcal W+m\ell_g-C(L+m^2L/k)\le\log X_u\le CmL,\qquad
 |\log X_u-\log X_{u'}|\le C\log k,
\end{equation}
put \(c_{\rm Count}=\max(C,0)+2\) and \(\lambda_5=\mathcal W+m\ell_g-c_{\rm Count}(L_a+m^2L_a/k)\). Then
\[
 \lambda_5\le\log T_u\le c_{\rm Count}mL,\qquad
 |\log T_u-\log T_{u'}|\le c_{\rm Count}\log k .
\]
For an admissible family with \ref{B8.49} and \ref{B8.50}, the mass and spectral estimates give common \(c_{\rm Count},x_0\) on all members at \(x\ge x_0\).
\end{theorem}
\begin{proof}[Outline]
Use \ref{B5.11c} and \ref{B8.50a} to compare row and column dimensions. Combine the volume and mass formulas \ref{B8.51b} and \ref{B8.48} with the wordwise Vandermonde comparison, then transfer the bounds from \(X_u\) to \(T_u\) using \ref{B8.51a}. The complete proof is in \S\ref{prf:B5.11}.
\end{proof}

\begin{lemma}[Uniform level ratios in the spectral row]\label{B5.11b}
Use the rank, numerical and weight hypotheses of the selected spectral row of \ref{B8.10}, with column-count factor \(C_0\ge1\) and \(C_\ell\ge40\), with \(C_\ell\) evaluated for this family as in Theorem~\ref{B8.10}. Beyond the rank and mass thresholds, every ordered pair of words and every type \(\rho\) satisfy
\[
 T_u/T_{u'}\le C_0e^{7/4},\qquad
 \Delta_\rho(u')\le(C_0e^{7/4})^k\Delta_\rho(u).
\]
The distinct-word row uses \(C_0=3\).
\end{lemma}
\begin{proof}[Outline]
Expand the log Vandermonde ratio, telescope the linear terms, and bound the remainder and edges using \ref{B5.05a} at the edge. Compare count products from \ref{B8.54} and \eqref{eq:B5-count-product} using \ref{B8.48} and shared row bounds, then take the \(kf\)-th root. The complete proof is in \S\ref{prf:B5.11b}.
\end{proof}

\begin{lemma}[Slopes before summing the column weights]\label{B5.11a}
For a word with a column model, write the integral highest weight at \(\rho\) as \(\lambda_i=\nu_i+t_\rho,\ 1\le i\le k\), with \(\nu_1\ge\cdots\ge\nu_k=0\). Put \(d_\rho=\sum_i\nu_i,\ s_\rho=\nu_1,\ l=\log T_u\). Then
\[
\begin{aligned}
 \mu(M_u)&\ge-\sum_\rho\bigl[(d_\rho/k+t_\rho)l+d_\rho\log k+
 k^2\log(s_\rho+k)\bigr],\\
 \mu_{\max}(M_u)&\le\sum_\rho\bigl[4d_\rho\log(kv_E)+k^2\log(s_\rho+k)+|t_\rho|l\bigr].
\end{aligned}
\]
\end{lemma}
\begin{proof}[Outline]
Compute the mean slope using the dual index, determinant character, Cauchy--Binet, and the hook formula of \ref{B5.07}. Bound maximum slope using transference for short columns, the product formula for the twist, and coordinate minors before taking the orthogonal sum. The complete proof is in \S\ref{prf:B5.11a}.
\end{proof}

\begin{theorem}[Mean and maximum slopes of the original lattice]\label{B5.12}
For every word and column model, the zeroth-layer lattice \(Q_{u,c,\rho}\) at a class \(c\) and column embedding \(\rho\) has slope \(-\log T_u/k\). Suppose an admissible family satisfies \ref{B5.11} and has column models beyond a common threshold. Assume also that \(|D_1|\le x^\varepsilon\) eventually. Then for common \(c_{\rm Slope}\ge0,x_0\) and every member at \(x\ge x_0\), every word satisfies
\[
 -\mu(M_u)\le c_{\rm Slope}n_TaL_a(1+m/k),\qquad
 \mu_{\max}(M_u)\le c_{\rm Slope}n_TaL_a .
\]
\end{theorem}
\begin{proof}[Outline]
Apply \ref{B5.11a} to the zeroth-layer lattice and sum the degree, determinant, and weight terms, retaining the leading top entries. Absorb the remaining terms by growth bounds and use the level estimate for \(m/k\) in the mean slope. The complete proof is in \S\ref{prf:B5.12}.
\end{proof}

\begin{definition}[The second complement at a point]\label{B5.12k}
Use the weights and losses of \ref{B4.02} and \ref{B4.03}. For a word \(u\), layer \(1\le j\le m\) and label \(\tau\), let \(A_{u,j,\tau}\) sum the entries of the layer diagram, and put \(A_{uj}=\sum_\tau A_{u,j,\tau}\). Use the Fock degree \(D_j^{\rm Fock}\) of \ref{B4.02} and write \(w_{j\tau}=w_j(\tau)\). Let \(c'_g\) be the fixed-prime error coefficient. Define
\[
 R_j=2\sum_\tau w_{j\tau}
 +C_{\deg}\bigl(f(kS_c+mt_{\rm ch}\eta+kN_m)+|TT|b'/x\bigr),
 \qquad H_g=c'_g f(N_m^2+m^2\eta).
\]
Let \(n_{iujw}=e_w^{-1}\sum_\tau A_{u,j,\tau}v_w(\tau(\bar\pi_i))\) be the normalized norm contribution, and let \(E^T_{u,w}=T_{u,w}^*\) be the level loss of \ref{B4.03}. Its full allowance \(E^\pi_u\) is the parenthesis of \(T_{u,w}^*\), used in \ref{B4.22a}.  For each signal or residue-class-zero prime \(p\le X''\), let \(J^*_{j,p}=C^*_{j,p}\) if \(p\) is in residue class zero and \(h_p=0\), and \(J^*_{j,p}=I^*_{j,p}\) otherwise, with the two losses as in \ref{B4.03}.

For real constants \(C_{\deg},C_g,C_T,C_{\rm stab},C_d,C_I\), the point predicate requires each word and layer to satisfy the following clauses at \(i=0,1\).
\begin{enumerate}
 \item \(T_u\ne g_i\).
 \item On the complete fibre, \(\sum_{w\mid g_i}c_wn_{iujw}=A_{uj}\log g_i\).
 \item \(\left|A_{uj}-D_j^{\rm Fock}-S_a\right|\le R_j\).
 \item \((R_j+2v_{g_i}(D_{\Gamma,u})+2H_g)\log g_i\le C_gB_0\).
 \item For every finite set \(F\) of places over \(T_u\),
 \[
 0\le\sum_{w\in F}c_wE^T_{u,w}\le E^\pi_u\log T_u\le C_TB_0.
 \]
 \item \(2\log D_{\Gamma,u}\le C_{\rm stab} B_0\).
 \item \(2\log d\,|D_j^*|\le C_dB_0\).
 \item For every finite set \(P\) of signal or residue-class-zero primes \(p\le X''\),
 \[
 \sum_{p\in P}2\log p\,|J^*_{j,p}|\le C_I B_0 .
 \]
\end{enumerate}
For a family with cutoff \(x_1\), require that there exist six real constants \(C_{\deg},C_g,C_T,C_{\rm stab},C_d,C_I\), with \(C_g\ge0\), and a real \(x_0\) such that all eight clauses hold for every \(x\ge x_0\) with \(x>x_1\) and every \(i\in I_x\). The hypothesis in \ref{B8.02} requires this property on \(W_3\) with the quantifiers of \ref{B8.39}. For each sufficiently small \(\varepsilon>0\), one choice of thresholds works for every family in \(W_3\) and every \(C_H\) satisfying the threshold inequalities.
\end{definition}

\section{Comparing words}

\begin{lemma}[The layerwise \(E\)-comparison]\label{B5.13}
Use the seven error conditions of \ref{B4.03} and the eight clauses of \ref{B5.12k} on the selected family above their common cutoff. The stored loss is the evaluated loss and the fixed primes are outside the level fibre. With \(c_E=2C_g+C_T+C_{\rm stab}+C_d+C_I\), each word and layer satisfies
\[
 \left|\sum_{w\in S}c_wE_{ujw}-S_a\ell_g\right|\le c_EB_0,\qquad
 \sum_{\substack{w\in S\\p_w\notin\{g_0,g_1\}}}|c_wE_{ujw}|\le c_EB_0.
\]
Consequently \(|\Ee_u-mS_a\ell_g|\le mc_EB_0\).
\end{lemma}
\begin{proof}[Outline]
Combine \ref{B4.03}, \ref{B4.02}, and \ref{B4.52} to evaluate the fixed-prime sums. Apply \ref{B5.12k} to bound their deviations, the level contribution, and the remaining terms grouped by prime, then sum over layers. The complete proof is in \S\ref{prf:B5.13}.
\end{proof}

\begin{lemma}[The difference of adjusted heights]\label{B5.14}
For the heights of \ref{B5.07} put \(H_u=h_u+S_a\log T_u\). For every pair of words of the fixed member,
\[
 H_u-H_{u'}=\sum_{j=1}^m\sum_\rho
 (K_{uj\rho}-K_{u'j\rho}).
\]
Use the column degrees of \ref{B4.02a} and suppose that the following size bounds hold. For selected members they follow from \ref{B5.08b} and \ref{B8.35a}, since \(S_c+w_{j\rho}\le3b\).
\[
 S_c+w_{j\rho}\le6x^6b,\qquad kS_c\le b,\qquad
 |D_1|\le L,\qquad 0\le\log\mathrm{rd}E\le L .
\]
Then \(|H_u-H_{u'}|\le(20+6/\varepsilon)(n_TaL_a+mB_0)\).
\end{lemma}
\begin{proof}[Outline]
Add \(S_a\log T_u\) to the heights of \ref{B5.07} and cancel common terms. Use \ref{B4.02a} to bound the varying head degrees, then integrate \(\log s-\log\Omega\) between them and sum using \(kS_c\le b\). The complete proof is in \S\ref{prf:B5.14}.
\end{proof}


% ---- p5c2 ----
\chapter{Prime sums and net gain}\label{ch:prime-sums-net-gain}

At a window point \(x\) of Alp\"oge's construction \cite{alpoge-2026},
let \(n_T=|TT|\) count the top positions and \(m\) the layers. Put
\(Q=n_Tmb\), with \(b=e^{x^\varepsilon}\) and \(0<\varepsilon<1/52\).
The common weights of \ref{B5.22} give bounds of order \(Q\) for the
deficit \(\ASK-\Ga\) and the index \(\mathrm{IO}\) of
\ref{B5.01}--\ref{B5.02}.

\section{The deficit at a rational prime}

Let \(J\) be the layer set and \(P_s\) the finite set of ordinary
primes. For \(j\in J\) and \(t\in TT\) let \(A_{jtp}\) be the
contribution before \(2\log p\) and \(D_{jt}\) the contribution at
\(d\), including \(\log d\). Write
\[
 \ASK=2\sum_{j,t,p\in P_s}\log p\,A_{jtp}+\ASK_d,\qquad
 \ASK_d=\sum_{j,t}D_{jt},
\]
where the first sum is over \(J\times TT\times P_s\).

\begin{lemma}[Decomposition over rational primes]\label{B5.15}
Give a finite set of places \(W\) weights \(\lambda_w\) with
\(\sum_{w\mid p}\lambda_w=\log p\) for \(p\in P_s\). Suppose the gains
and deductions have constant values \(g_{jp}\) and \(\tau_p\) on each
fiber and vanish off \(P_s\). Set
\[
 \Ga=2\sum_{w\in W}\lambda_w\left(\sum_j g_{j,p(w)}-\tau_{p(w)}\right),
 \qquad
 \Delta_p=\sum_{j,t}A_{jtp}-\sum_j g_{jp}+\tau_p .
\]
Then
\[
 \ASK-\Ga=2\sum_{p\in P_s}\log p\,\Delta_p+\ASK_d .
 \tag{5.15.1}\label{5.15.1}
\]
Let \(x_1\) be the starting scale of the admissible family \(F_m\). For \(x\ge\max(x_1+1,3)\), the model sums use the
signal and inert-class primes of \ref{B3.46} up to
\(B=\max(\lfloor a+a/(kx^{13})\rfloor,\lfloor X''\rfloor)\), with the
scales of \eqref{5.20.1}. In a consistent
presentation, the model sum equals \(\ASK_N\) of \ref{B5.01} for any
integer cutoff \(N\ge a_t/w_j(\rho)\) for every \(t\in TT\),
\(j=0,\ldots,m-1\), and \(\rho\in\Phi\).
\end{lemma}

\begin{proof}[Outline]
Group places by prime using residue degrees, ramification indices and fiber constancy from \ref{B3.52}. The bounds in \ref{B5.17a} and \ref{B5.21a} remove ordinary terms above \(B\), and the matching places in \ref{B5.21a} identify the weights at each eligible cutoff. The complete proof is in \S\ref{prf:B5.15}.
\end{proof}

\begin{lemma}[Covered positions and discarded primes]\label{B5.16}
Partition \(P_s\) into served primes \(S\) and discarded classes
\(P_1,P_2,P_h\). At \(p\in S\) suppose \(TT_p\subseteq TT\) and
\[
 \sum_{j\in J,\ t\in TT_p}(A_{jtp}-\ell_p)-mr_p
 \le\sum_jg_{jp}-\tau_p . \tag{5.16.1}\label{5.16.1}
\]
If gains and deductions vanish at discarded primes, then
\[
 \ASK-\Ga\le
 2\sum_{p\in S}\log p\left(\sum_{j,\ t\notin TT_p}A_{jtp}
 +m|TT_p|\ell_p+mr_p\right)
 +2\sum_{p\in P_1\cup P_2\cup P_h}\log p\sum_{j,t}A_{jtp}
 +\ASK_d . \tag{5.16.3}\label{5.16.3}
\]
\end{lemma}

\begin{proof}
Use \eqref{5.16.1} in \(\Delta_p\) at served primes and the vanishing
hypothesis at discarded primes, then apply \eqref{5.15.1}.
\end{proof}

In the arithmetic model \(P_h\) consists of discarded primes above
\(X''\). Among the remaining discarded primes, \(P_2\) consists of the
candidates of \ref{B3.51} with \(\Sigma_p>\theta=x^{-\varepsilon}\).
The others form \(P_1\).

\section{The finite-sum loss bound}

Write \([u]_+=\max(u,0)\) and \(\log_p a=\log a/\log p\). With
nonnegative constants \(C,f,S_k,\ell_x,\mathcal L,Q_1,Q_2\),
positive \(a,k,K_0\), and a real function \(L_1\), put
\[
\begin{aligned}
 F_p&=Cf\{N_0^2\log_p a+(\log_p a)^2\},&
 C_p&=CS_k\log_p(2a),\\
 R_p&=\begin{cases}Q_1&p<p_0,\\Q_2&p\ge p_0,\end{cases}&
 O_p&=C\ell_x\left(\frac{\mathcal L a}{pk}+
 \begin{cases}ap^{-L_1(p)}&p<p_0,\\ a/p^2&p\ge p_0.\end{cases}\right)
\end{aligned}\tag{5.17.1}\label{5.17.1}
\]

\begin{theorem}[The losses from all prime classes]\label{B5.17}
Assume \ref{B5.15} and \ref{B5.16}, \(a,X'\ge1\), \(X''\ge2\),
\(X''w''=a\), \(p_0\ge3\), \(d\ge4\), and
\(0\le a_t\le a(1+\vartheta)\) with \(0\le\vartheta\le1/2\). Let
\[
 A_{jtp}=\sum_{i\in I_{jtp}}[a_t/p^i-w_{jtpi}]_+,\qquad
 D_{jt}=\log d\sum_{i\in I^d_{jt}}[a_t/d^i-w^d_{jti}]_+ .
 \tag{5.17.2}\label{5.17.2}
\]
The exponent sets are finite, with \(i\) positive odd and
\(w_{jtpi}\ge w''\) in the ordinary sum and \(i\ge1,\ w^d_{jti}\ge0\)
in the ramified sum. Suppose
\[
 \begin{gathered}
 p\le X''\ (p\in S),\quad p>X''\ (p\in P_h),\\
 r_p\le F_p,\qquad \ell_p\le C_p+\mathbf1_{\{p<X'\}}(R_p+O_p)
 \quad(p\in S),\\
 K_0\le p^{L_1(p)-1}\quad(p\in P_s,\ p<p_0).
 \end{gathered}\tag{5.17.3}\label{5.17.3}
\]
Then \(\ASK-\Ga\le E_{\rm fl}+E_{\rm cp}+E_{\rm pr}+E_{\rm or}
+E_{\rm out}+E_1+E_2+E_h+E_d\), where
\[
\begin{aligned}
 E_{\rm fl}&=2mCfX''\{N_0^2\log a+2(\log a)^2\},&
 E_{\rm cp}&=2CS_k\log(2a)\,mn_TX'',\\
 E_{\rm pr}&=2\log4\,mn_T(p_0Q_1+X'Q_2),\\
 E_{\rm or}&=2C\ell_x a\,mn_T
 \left\{\left(\frac{\mathcal L}{k}+\frac1{K_0}\right)
 (\log X'+\log4)+\frac{2(\log(p_0-1)+1)}{p_0-1}\right\},\\
 E_{\rm out}&=4ma\sum_{p\in S}\frac{\log p}{p}|TT\setminus TT_p|,&
 E_i&=4mn_Ta\sum_{p\in P_i}\frac{\log p}{p}\quad(i=1,2),\\
 E_h&=2mn_T\log(2X'')\{a\vartheta^2+w''\vartheta\},&
 E_d&=2mn_Ta\,\frac{\log d}{d}.
\end{aligned}\tag{5.17.5}\label{5.17.5}
\]
\end{theorem}

\begin{proof}[Outline]
Bound losses by geometric series and counts of served primes and covered positions, keeping only the first power per short interval above \(X''\). Use binomial and factorial counts for the prime sums and integrate the inverse-square tail to bound price and orbit errors. The complete proof is in \S\ref{prf:B5.17}.
\end{proof}

\begin{lemma}[Lower bound for an ordinary partner]\label{B5.17a}
In the model of \ref{B3.46}--\ref{B3.47}, the weight at position \(z\) in lift \(\ell\) is \(B_\ell(z)+\kappa_\ell\). It is at least \(w''=b'/(xn)\) for every lift and every \(z\in\{0,1\}\times H\). Each ordinary exponent \(i\) in \eqref{5.17.2} is odd and has weight \(B_\ell(z_i(t))+\kappa_\ell\).
\end{lemma}
\begin{proof}
The marked, thick, and thin lower bounds in \ref{B3.47} are
\(b'/x,b,w''\), respectively, and \(w''\le b'/x\le b\). The
partner formula and exponent restriction are those of \ref{B3.46}
and \ref{B3.52}.
\end{proof}

\section{Exceptional primes and exceptional positions}

\begin{lemma}[Small order and narrow prime-power bands]\label{B5.18}
Let \(e,J_0\ge1\) be integers, \(Y_0\ge1\), \(X'\ge2\), \(p_0\ge8\),
and \(0\le\rho\le1/2\). Write \(h_e(s)\) for the class of an odd
integer in \((\Z/2^{e+2}\Z)^\times/\langle-1\rangle\). Suppose every
prime in a finite set \(T\) is odd and satisfies at least one of
\[
 p\le Y_0,\quad\operatorname{ord}(h_e(p))\le J_0,
 \qquad
 p\ge p_0,\quad X'(1-\rho)\le p^i\le X'(1+\rho)
 \quad\text{for some }i\ge1.
\]
Then, with \(C_{\rm ord}=12\),
\[
 \sum_{p\in T}\frac{\log p}{p}\le
 C_{\rm ord}J_0\frac{\log^2(2^eY_0)}{2^e}
 +\lfloor\log X'\rfloor(4\rho+p_0^{-1})\log(2X').
 \tag{5.18.1}\label{5.18.1}
\]
For \(P_1\) under these hypotheses with \(Y_0=\max(Y,X'')\) and
\(J_0=\lfloor x^{\varepsilon/4}\rfloor\), write
\(E_1\le E_{\rm ord}+E_{\rm band}\) for \eqref{5.18.1} multiplied
by \(4mn_Ta\), in the order of its two terms.
\label{5.18.2}
\end{lemma}

\begin{proof}[Outline]
Use successive differences of squares to put small-order classes in two residue classes and bound their contribution by a harmonic sum. For each band exponent, take roots of the interval endpoints, count the integers, and sum the resulting contribution bounds. The complete proof is in \S\ref{prf:B5.18}.
\end{proof}

For candidate primes \(\mathcal P_c\) write
\[
 R_W=\sum_{p\in\mathcal P_c}\frac{\log p}{p}
 (2\mathbf1_{\{\Sigma_p>\theta\}}+4\Sigma_p).
 \tag{5.19.1}\label{5.19.1}
\]

\begin{lemma}[Bad positions, second-class primes, and local indices]\label{B5.19}
For each prime \(p\) in a finite set \(\mathcal P_R\), choose a set \(U_p\subseteq TT\) and put \(U_p=\varnothing\) for \(p\notin\mathcal P_R\). Suppose \(TT\setminus TT_p\subseteq U_p\) at every served prime and \(\sum_{p\in\mathcal P_R}(\log p/p)|U_p|\le R_{\rm out}\).
Then \(E_{\rm out}\le4amR_{\rm out}\). If the candidate shares are
nonnegative and \(P_2\subseteq\{p\in\mathcal P_c:\Sigma_p>\theta\}\),
then \(E_2\le2amn_TR_W\). If
\[
 \mathrm{io}_p\le2\Sigma_p\,\frac{2mn_Ta}{p}\qquad(p\in\mathcal P_c),
 \tag{5.19.2}\label{5.19.2}
\]
then, for arbitrary real shares,
\[
 \sum_{p\in\mathcal P_c}\log p\,\mathrm{io}_p\le amn_TR_W .
 \tag{5.19.3}\label{5.19.3}
\]
The last two implications hold with any \(T\ge0\) in place of \(n_T\)
throughout and \(R\ge R_W\) on the right.
\end{lemma}
\begin{proof}
Apply the bad-position count to \(E_{\rm out}\). For \(E_2\) use
the indicator in \eqref{5.19.1}. Summing \eqref{5.19.2} gives
\eqref{5.19.3} since the indicator terms are nonnegative.
\end{proof}

\section{Losses on the scale \(Q\)}

For \(x\ge1,\ 0<\varepsilon\le1\) put
\[
\begin{gathered}
 z=x^\varepsilon,\quad u=x^{\varepsilon/3},\quad b=e^z,\quad
 a=xb,\quad X'=e^u,\quad b'=a/X',\quad f=2n,\quad
 w''=b'/(xn),\quad X''=xnX',\\
 \delta_1=x^{-\varepsilon/8},\quad\theta=x^{-\varepsilon},\quad
 \ell_x=\log x,\quad S_k=S_c+k+\eta,\quad Q_1=Cr_Gb'\ell_x,\\
 Q_2=C\{b'/x+\eta+S_c+(w''+t_0)\log x/\delta_1\}+b_0 .
\end{gathered}\tag{5.20.1}\label{5.20.1}
\]
For \(\nu\ge0,\ K_s\ge1\) the size conditions are
\[
\begin{gathered}
 x^2\le2^{e_2}\le n\le x^\nu,\quad x^2\le k\le x^\nu,\quad
 0\le N_0,t_0,r_G\le x^\nu,\quad x^4+1\le p_0\le x^\nu,\\
 0\le\mathcal L\le x^{2\varepsilon},\quad
 0\le S_c,\eta\le e^{x^{\varepsilon/4}},\quad
 0\le\rho\le x^{-2},\quad0\le\vartheta\le x^{-1},\\
 K_0,d\ge x^2,\quad 0\le C,C_{\rm ord}\le K_s,\quad
 1\le Y\le2X',\quad 0\le b_0\le e^{z/2},
\end{gathered}\tag{5.20.2}\label{5.20.2}
\]
with \(e_2,p_0,d\) integers. Assume
\[
 100K_sx^{4\varepsilon}\le x,\qquad
 100K_sx^{4\nu+4}\le e^u,\qquad
 100K_sx^{4\nu+4}e^{u+x^{\varepsilon/4}}\le e^{z/2}.
 \tag{5.20.3}\label{5.20.3}
\]
Set \(\sigma_x=100K_sx^{4\varepsilon-1}\).
\label{5.20.4}

\begin{theorem}[Joint loss bounds in a size window]\label{B5.20}
Take \(C_{\rm ord}=12\), \(TT\ne\varnothing\), and \(P_s\) consisting
of odd primes. Assume \eqref{5.20.1}--\eqref{5.20.3}. For each family
member \(\upsilon\) assume the hypotheses of \ref{B5.17} and
\ref{B5.19}, with common nonnegative shares and bad-position predicate,
and \ref{B5.18} for \(P_1\) with \(e=e_2,\ Y_0=\max(Y,X'')\) and
order \(<x^{\varepsilon/4}\) in the small-order case. If
\[
 R_{\rm out}\le n_Tx^{-2+\varepsilon},\qquad R_W\le x^{-9},
 \tag{5.20.5}\label{5.20.5}
\]
then, for \(Q=n_T|J|b\),
\[
 \begin{aligned}
 \ASK_\upsilon-\Ga_\upsilon&\le(2C\log4+13\sigma_x)Q\le(3C+13)Q,\\
 \sum_{p\in\mathcal P_c}\log p\,\mathrm{io}_{\upsilon,p}
 &\le\sigma_xQ\le Q .
 \end{aligned}\tag{5.20.7}\label{5.20.7}
\]
For fixed \(K_s,\nu\) and \(0<\varepsilon<1/4\), \eqref{5.20.3}
holds eventually.
\end{theorem}

\begin{proof}[Outline]
Use the served-prime and index bounds if \(m=0\). Otherwise divide losses by \(Q\), insert window sizes and use \ref{B5.18} and \ref{B5.19} to bound the excess over \(2C\log 4\) by multiples of \(\sigma_x\), with size inequalities following from growth comparisons. The complete proof is in \S\ref{prf:B5.20}.
\end{proof}

\begin{lemma}[Splitting an integral pairing among its summands]\label{B5.20k}
Let \(O\) be a discrete valuation ring and \(0\ne z\in O\).
If a finite family \(A_j\in M_n(O)\) satisfies \(\sum_j A_j=zI_n\),
there are integral rectangular matrices \(P_j,Q_j\) with
\(P_jA_jQ_j=zI_{n_j}\) and \(\sum_jn_j\ge n\).
\end{lemma}
\begin{proof}
In Smith bases \cite[\S4.5]{ash-2007} retain the \(n_j\) entries
\(a\mid z\), and scale their source vectors by \(z/a\). Let \(W\)
span the retained target vectors and \(Z=zO^n\). The other entries
belong to \(\mathfrak m zO\), where \(\mathfrak m\) is the maximal
ideal. Hence \(Z\subset W+\mathfrak m Z\). Nakayama's lemma applied
to \((W+Z)/W\) gives \(Z\subset W\) \cite[Prop.~2.6]{atiyah-1969}.
Taking ranks gives the claim.
\end{proof}

\section{Local projection estimates}

Use the lattices \(H_\Lambda,R_e,R'_e\) and pairing
\([F,G]=G(p^{A_{r'}}F)\) of \ref{B3.52a}. The test at rank \(r'\),
depth \(t\), denominator \(D\), and order \(c_*\) asks for a pairing
matrix \(p^DI\) of rank at least \((1-\kappa)\dim V\) under the
stated lattice conditions.

\begin{lemma}[Separating the residue planes]\label{B5.20d}
Let \(1\le r\le k\), \(F=\mathbf F_{p^{n'}}\), \(n'\ge1\), and
\(0\le z_0\le\cdots\le z_b=z_0+n'\), \(b\ge1\). Put
\(Q_j=p^{z_{j+1}-z_j}\), \(I_0=\{1,\ldots,r\}\).
Let \(H\) be a rank-\(r\) row matrix of a plane. In frame \(j\)
take row \(l\) to be \(\sum_i T_{jli}H_{i,*}^{[p^{u_{j,l}}]}\),
raised entrywise to \(p^{z_j-u_{j,l}}\), with \(0\le u_{j,l}\le z_j\)
and independent \(T_{jli}\). Write \(M_{j,I}\) for its minor on \(I\)
and \(\Delta_j=M_{j,I_0}\). The polynomial
\[
 L=\Delta_0\prod_I\left(\prod_{j<b}\Delta_j^{Q_j-1}
 -\prod_{j<b}M_{j,I}^{Q_j-1}\right)
\]
over the one-column replacements \(I\) of \(I_0\) vanishes on
every competing plane and restricts on the coordinate plane to
\(\Delta_0\prod_{j<b}\Delta_j^{r(k-r)(Q_j-1)}\ne0\).
\end{lemma}
\begin{proof}[Outline]
Factor minors via \(d_I^{p^{z_j}}\); for \(d_{I_0}\ne 0\), row-reducing to \(I_r\) kills a competing plane's replacement factor via \(d^{p^{z_b}-p^{z_0}}=1\) on \(F^\times\). On the coordinate plane replacements vanish with unit diagonals, and lifting \(L^h\) gives order \(h\). The complete proof is in \S\ref{prf:B5.20d}.
\end{proof}

\begin{lemma}[Assembling tests on incoming edges]\label{B5.20e}
Let \(t,D\in\N\), \(c_*\ge4(D+1)\), and take a finite set \(E\)
of incoming edges at a target \(\Lambda\). For each \(e\in E\)
let \(F_{ei}\), \(1\le i\le n_e\), satisfy the source conditions of
\ref{B3.52a} at \(e\). Suppose for every \(u\in O^{n_e}\) with a
unit coordinate there is a target vector of depth \(t\), integral
in every \(R'_{e'}\), whose value on \(\sum_i u_iF_{ei}\) divides
\(p^D\) and which belongs to \(\omega^{c_*}R'_{e'}\) for
\(e'\in E\setminus\{e\}\). The test has retained rank \(\sum_e n_e\).
\end{lemma}
\begin{proof}
On a primitive row of length \(h=\sum_e n_e\), use the test for a
block with a unit coordinate. The other blocks contribute multiples
of \(p^{D+1}\) since \(\omega^2=p\). Thus every primitive row takes
a value dividing \(p^D\) on the image of the target lattice in
\(O^h\). Apply this to the Smith coordinate rows. The image contains
\(p^DO^h\), whose coordinate vectors give the required lifts.
\end{proof}

\begin{lemma}[A window of integral Gram matrices]\label{B5.20f}
Let \(A(T)=\sum_{i=0}^JC_i\binom Ti\), \(C_i\in M_d(\Z_p)\). Suppose
\(\det A\ne0\) has degree \(s\) and leading coefficient \(b\) with
\(v(b)+s/(p-1)\le dL\), where \(v(p)=1\). Let \(h\in\N\) and \(H,J+1\le p^h\),
\(H>0\), and let \(\beta_u\) count Smith factors of valuation at
least \(E>0\), including zero factors, on \(H\) consecutive
nonnegative integers. Then
\[
 H^{-1}\sum_u\beta_u\le d\left(J/H+(L+Jh/H)/E\right).
\]
If \(A\) is the matrix of a pairing multiplied by a nonzero
integral scalar, the retained Smith directions give integral
pairs of valuation below \(E\) for the original pairing.
\end{lemma}
\begin{proof}[Outline]
Prove integrality after substitution and use the Smith rows to bound the determinant's Gauss norm. Count integers near each root by residue classes up to height \(h\) and by single points above it, then sum and subtract the scalar valuation. The complete proof is in \S\ref{prf:B5.20f}.
\end{proof}

\begin{lemma}[Returning from the enlarged weight box]\label{B5.20g}
Let \(|S|=N\ge1\), and let \(I_a,I_b\subseteq\N\) be integer intervals
of sizes \(w_a,w_b>0\). Enlarge their upper endpoints by \(J+c\)
for \(J,c\in\N\). Write \(\mathbb E_+\) for the enlarged product
average and \(\mathbb E_\tau\) for the average on
\((I_a+\tau)^S\times(I_b+\tau)^S\). Suppose \(f,g,q\ge0\) satisfy
\[
 \sum_{j\in\{0,\ldots,J\}^S} f(A-j,B-j)\prod_{s\in S}q(A_s,B_s,j_s)
 \le g(A,B)
\]
on the enlarged box, with natural subtraction. For \(\tau=0,c\)
suppose \(\sum_{i=0}^J q(\alpha+i,\beta+i,i)\ge\delta\ge0\) on the
translated coordinate intervals and
\((w_a+J+c)^N(w_b+J+c)^N\le2\delta^Nw_a^Nw_b^N\).
Then \(\mathbb E_\tau f\le2\mathbb E_+g\). On the original box let
\(f_-=f(A-c,B-c)\) if the shifted pair belongs to it and \(f_-=1\)
otherwise. If the latter event has probability at most \(1/Y\) and
\(0\le e\le1/Y\), then
\[
 \mathbb E_0(f_-+2f(A+c,B+c)+e)\le6\mathbb E_+g+2/Y.
\]
\end{lemma}
\begin{proof}
For each \(j\) translate the pairs from the \(\tau\)-box by \(j\).
At each pair their total weight is at least \(\delta^N\) by
factorization over the coordinates. Sum and compare cardinalities.
The downward shift is injective wherever it stays in the original
box. Thus \(\mathbb E_0 f_-\le\mathbb E_0f+1/Y\), and the second claim follows.
\end{proof}

\begin{definition}[Paths and their labels]\label{B5.20l}
A signed letter of degree \(b\) denotes \(\operatorname{Sym}^b V\) for a vector letter and \(\operatorname{Sym}^b V^\vee\) for a covector letter. A path in rank \(k\) starts at the zero row. At a vector letter it passes from a decreasing integral row \(\mu\) to a row \(\nu\) with
\[
 \nu_1\ge\mu_1\ge\nu_2\ge\mu_2\ge\cdots\ge\nu_k\ge\mu_k,
 \qquad \sum_j(\nu_j-\mu_j)=b.
\]
At a covector letter interchange \(\mu,\nu\) in these conditions. An ordinary restriction step from rank \(k\) to rank \(k-1\) chooses a row \(\eta\) with \(\mu_j\ge\eta_j\ge\mu_{j+1}\).

Fix \(k=r+K\). After the signed letters take \(r\) restriction steps and denote the resulting \(K\)-row by \(\lambda\). A label is a semistandard tableau of shape \(\lambda-\lambda_K\) on \(1,\ldots,K\), with the determinant character \(\det^{\lambda_K}\) understood. Equivalently, put \(h(c)=\#\{a:\lambda_a\ge c\}\) for \(c\in\Z\). Choose sets \(I_c\subseteq\{1,\ldots,K\}\) of size \(h(c)\) for \(c>0\) and \(K-h(c)\) for \(c\le0\). Put \(J_c=I_c\) for \(c>0\) and \(J_c=I_c^{\,c}\) for \(c\le0\). The tableau condition is
\[
 |J_{c+1}\cap[1,t]|\le |J_c\cap[1,t]|\qquad(1\le t\le K).
\]
The sets \(I_c\) are empty outside a finite interval. In particular the multiplicity of the path is \(\dim V_\lambda\).
\end{definition}

\begin{lemma}[Loss at the prefix cuts]\label{B5.20m}
Take a word of \(n_0\) signed letters followed by \(r\) restriction steps, and put \(t=n_0+r\), \(k=r+K\), \(q=K-2t\). Suppose \(n_0\ge2r\), the first and last \(r\) signed letters are vectors, \(t+r\le K\), and \(q\ge k/2\). Let \(k\le b_i\le2a\) be the signed-letter degrees, where \(a\ge3\), \(p\ge13\). Degrees of successive letters of the same sign increase by a factor at least \(p^2/5\).

At a restriction step use the shorter row for computing sums of positive or negative parts. To refer to a coordinate in rank \(k\), restore the deleted leading positions with the common value \(b_*=\sum_{i\le n_0}b_i\). Count a restriction step as a vector letter when locating the last \(r\) vector letters before a cut. Before step \(i>r\), denote the oldest of these \(r\) letters by \(o(i)\). It is a signed letter. Give the cut a scale \(v_i>0\) such that every earlier signed-letter degree is at most \(4v_i/p\), and put \(\xi_i=L v_i/(pk)\), where \(L>0\).

Discard a path if any of the following occurs. A covector step removes more than \(\xi_i\) from the sum of positive parts. A vector step after the first \(r\), including a restriction step, removes more than \(\xi_i\) from the sum of negative parts in absolute value. At a cut before step \(i>r\), coordinate \(r\) of the restored row is less than \(2b_{o(i)}/3\). The fraction discarded, counting labels, is at most
\[
 \exp\left(11t^2\log(1+2at)-\frac{\min(L,k)}{32}\right).
\]
\end{lemma}

\begin{proof}[Outline]
Pieri rules and interlacing give the path count in \ref{B5.20l} and propagate tail deficits. Concavity and summation by parts bound the central Weyl factors for each failed test, and bounding the other factors and counting rows gives the discarded fraction. The complete proof is in \S\ref{prf:B5.20m}.
\end{proof}

\begin{lemma}[Labels on a shared frame row]\label{B5.20n}
Let \(Q,b,d\) be nonnegative integers with \(Q>0,\ d\le2b<Q\), and let \(r\ge2\). In every prescribed residue class for \(d+\sum_{i=1}^r v_i\) modulo \(Q\), there are at least
\[
 Q^{r-1}(1-2rb/Q)
\]
choices with \(0\le v_i<Q-2b\). For each choice, \(e=d+\sum_i v_i\) satisfies \(0\le e<rQ\).
\end{lemma}
\begin{proof}
With coordinates in \([0,Q)\) the count is \(Q^{r-1}\), since the congruence determines the last coordinate. For a specified coordinate and a specified value in \([Q-2b,Q)\), a different coordinate can be solved for. Thus each specified value removes \(Q^{r-2}\) choices. Sum over the \(2b\) values and the \(r\) coordinates. Finally \(d+\sum_i v_i\le2b+r(Q-2b-1)<rQ\).
\end{proof}

\begin{lemma}[Binomial approximation to a product of digit units]\label{B5.20ha}
Let \(p\ge5\) be prime. For each element \(i\) of a finite set \(I\), fix a power \(P_i\) of \(p\) and an integer \(0\le d_i<P_i\). Put
\[
 U(b)=\prod_{i\in I}
 \frac{(d_i+P_i b_i)!}{d_i!\,b_i!\,(P_i!)^{b_i}}
 \qquad(b\in\N^I).
\]
The values of \(U\) are \(p\)-adic units. For either sign and every integer \(c\ge1\) there are coefficients \(a_\nu\in\Z_p\), indexed by \(\nu\in\N^I\) with \(\sum_i\nu_i\le4c\), such that
\[
 U(b)^{\pm1}\equiv
 \sum_\nu a_\nu\prod_i\binom{b_i}{\nu_i}\pmod {p^c}
 \qquad(b\in\N^I).
\]
\end{lemma}
\begin{proof}[Outline]
Split the factorials into blocks and expand the unit factors or their inverses in an auxiliary variable. Bound coefficient degrees by finite differences, then specialize at \(p\), reduce modulo \(p^c\), and use the product binomial basis. The complete proof is in \S\ref{prf:B5.20ha}.
\end{proof}

\begin{lemma}[Uniform competing-edge orders of columns, bordered factors, and ties]\label{B5.20he}
Let \(p\ge5\) be prime, let \(n'\ge1\), let \(1\le r<k\) with \(K=k-r\), and let \(O\) be a discrete valuation ring over \(W(\mathbf F_{p^{n'}})\) with uniformizer \(\omega\) satisfying \(\omega^2=p\). Let \(\sigma\) be the Frobenius automorphism of \(W(\mathbf F_{p^{n'}})\). For \(g\in\GL_k(W(\mathbf F_{p^{n'}}))\) and \(a\in\Z\), write \(g_a=\sigma^a(g)\) for its \(a\)-th Frobenius conjugate, and define the \(g\)-twisted vector row \(v^{(g)}_a=X^{\rm vec}_a g_a\) and covector row \(w^{(g)}_a=X^{\rm cov}_a g_a^{-\mathsf t}\). For any \(g'\in\GL_k(W(\mathbf F_{p^{n'}}))\), let the edge substitution \(s_{g'}\) send \(X^{\rm vec}_a\) to \(X^{\rm vec}_a\operatorname{diag}(I_r,\omega I_K)(g'_a)^{-1}\) and \(X^{\rm cov}_a\) to \(X^{\rm cov}_a\operatorname{diag}(\omega I_r,I_K)(g'_a)^{\mathsf t}\). Define \(\operatorname{ord}_{g'}(F)=\sup\{m\in\N:s_{g'}(F)\in\omega^m O[X]\}\in\N\cup\{\infty\}\) (so \(\operatorname{ord}_{g'}(F)\ge m\) means \(s_{g'}(F)\in\omega^m O[X]\)), and write \(X^{\rm vec}_{a,b}\) (\(r<b\le k\)) and \(X^{\rm cov}_{a,b}\) (\(1\le b\le r\)) for the normal variables.
\begin{enumerate}
\item Let \(L\) be a list of \(N\ge1\) signed letters with nondecreasing integer positions \(x_0\le\cdots\le x_{N-1}\le z\) such that distinct letters of the same sign have strictly increasing positions. Let \(P\) be a path of rank \(k\) on \(L\) with column heights \(h_t=\#\{a:\mu^{(t)}_a\ge c\}\) for a fixed column \(c\in\Z\), and call step \(i\) (\(0\le i<N\)) used in \(c\) if \(h_{i+1}\ne h_i\). Call a used step \(i\) costly if it is a vector step from height \(h_i\ge r\) or a covector step from height \(h_i\le r\), and suppose \(x_i<x_j\) whenever \(i<j\), \(i\) is used in \(c\), and \(j\) is costly in \(c\). Indexing \(K^{(t)}\) by subsets \(I\subseteq\{1,\ldots,k\}\) of size \(h_t\) (for \(c>0\)) or \(k-h_t\) (for \(c\le0\)), define \(K^{(t)}\) recursively over \(O[X]\) from \(K^{(0)}_\varnothing=1\): at the first used letter use power \(1\), at each subsequent used letter \(i\) raise each coordinate polynomial \(K_I\) to the polynomial power \(p^{x_i-y}\) (where \(y\) is the preceding used position) and, writing \(u=v^{(g)}_{x_i}\) at a vector letter and \(u=w^{(g)}_{x_i}\) at a covector letter, update by right exterior multiplication \(\sum_{b\in I}(-1)^{\#\{j\in I:j>b\}}u_bK_{I\setminus\{b\}}\) (for \(c>0\), vector, or \(c\le0\), covector) or right contraction \(\sum_{b\notin I}(-1)^{\#\{j\in I:j>b\}}u_bK_{I\cup\{b\}}\) (for \(c>0\), covector, or \(c\le0\), vector), and after the last used position \(y_{\rm last}\) raise each coordinate to \(p^{z-y_{\rm last}}\) to obtain \(K_c(g;J)\) for each terminal index \(J\) (with \(K_c(g;\varnothing)=1\) if no step is used in \(c\)). Let \(\mathrm{oc}(c)=\sum_{i\text{ costly in }c}p^{z-x_i}\). Then:
\begin{enumerate}
\item \(K_c(g;J)\in O[X]\) is homogeneous of multidegree \(p^{z-x_i}\) at every used letter \(i\) (and \(0\) at unused letters);
\item \(\operatorname{ord}_{g'}(K_c(g;J))\ge\mathrm{oc}(c)\) for every \(g'\in\GL_k(W(\mathbf F_{p^{n'}}))\);
\item under the chosen substitution \(s_g\),
\[
 s_g\bigl(K_c(g;J)\bigr)-\omega^{\mathrm{oc}(c)}\operatorname{pow}_z\bigl(K^{\rm sep}_c(J)\bigr)
 \in\omega^{\mathrm{oc}(c)+1}O[X],
\]
where \(K^{\rm sep}_c(J)\) is the unpowered separated column polynomial built by the same recursion with power \(1\) at every step using only nonnormal variables at non-costly steps and normal variables at costly steps, and \(\operatorname{pow}_z\) is the entrywise segment substitution sending each variable at letter \(i\) to its \(p^{z-x_i}\)-th power;
\item \(\mathrm{oc}(c)\le 2p^{z-x_0}\).
\end{enumerate}
\item Let \(y_1,\ldots,y_r<x\) be \(r\) distinct vector positions strictly preceding \(x\). Let \(A_x(g)\) be the \(r\times k\) matrix with entries \((v^{(g)}_{y_a,c})^{p^{x-y_a}}\), let \(M_x(g)\) be its first \(r\) columns, and put \(D_x(g)=\det M_x(g)\). At a vector letter at \(x\) and \(r<b\le k\), let \(B_{x,b}(g)\) be the \((r+1)\times(r+1)\) minor on columns \(\{1,\ldots,r,b\}\) obtained by bordering \(A_x(g)\) below with \(v^{(g)}_x\). At a covector letter at \(x\) and \(1\le b\le r\), let \(P_{x,b}(g)\) be the \(b\)-th entry of \(\operatorname{adj}(M_x(g))\,A_x(g)\,(w^{(g)}_x)^{\mathsf t}\). Then \(B_{x,b}(g),P_{x,b}(g)\in O[X]\), and for every \(g'\in\GL_k(W(\mathbf F_{p^{n'}}))\),
\[
 \operatorname{ord}_{g'}\bigl(B_{x,b}(g)\bigr)\ge1,\qquad
 \operatorname{ord}_{g'}\bigl(P_{x,b}(g)\bigr)\ge1.
\]
Under \(s_g\),
\[
 s_g\bigl(B_{x,b}(g)\bigr)-\omega s_g(D_x(g))X^{\rm vec}_{x,b}\in\omega^2 O[X],\qquad
 s_g\bigl(P_{x,b}(g)\bigr)-\omega s_g(D_x(g))X^{\rm cov}_{x,b}\in\omega^2 O[X],
\]
and every term of either difference has normal degree equal to its \(\omega\)-order, at least \(p\).
\item At any position \(x\) carrying both a vector row and a covector row, for \(E_x=\sum_{c=1}^k v^{(g)}_{x,c}w^{(g)}_{x,c}\) and any \(P=p^e\) (\(e\ge0\)), every \(g'\in\GL_k(W(\mathbf F_{p^{n'}}))\) satisfies
\[
 s_{g'}(E_x^P)=\omega^P\left(\sum_{c=1}^k X^{\rm vec}_{x,c}X^{\rm cov}_{x,c}\right)^P\in\omega^P O[X],
\]
and \(s_g(E_x^P)-\omega^P\sum_{c=1}^k(X^{\rm vec}_{x,c})^P(X^{\rm cov}_{x,c})^P\in\omega^{P+2}O[X]\).
\item For any \(g'\in\GL_k(W(\mathbf F_{p^{n'}}))\), the order \(\operatorname{ord}_{g'}\) is additive on products in \(O[X]\). In particular, any product of \(\omega^m\), column factors \(K_c(g;J)\) of orders \(a_c\), \(q\) bordered or pairing factors, tie factors \(E_x^{P_h}\), integral nonnormal factors, and a localizer power of order at least \(c_0\) on competing planes has order at least \(m+\sum_c a_c+q+\sum_h P_h+c_0\) under every competing \(s_{g'}\).
\end{enumerate}
\end{lemma}
\begin{proof}
Fix \(g'\in\GL_k(W(\mathbf F_{p^{n'}}))\) and put \(H=(g')^{-1}g\in\GL_k(W(\mathbf F_{p^{n'}}))\) and \(H_a=\sigma^a(H)=(g'_a)^{-1}g_a\). Write \(u_1(a),\ldots,u_k(a)\in O^k\) for the rows of \(H_a\) mapped into \(O\), and \(u_1^\vee(a),\ldots,u_k^\vee(a)\in(O^k)^\vee\) for the rows of \((H_a^{-1})^{\mathsf t}\). Let \(\widetilde W_a\subseteq O^k\) be the rank-\(r\) submodule spanned by \(u_1(a),\ldots,u_r(a)\) and \(\widetilde W_a^\perp\subseteq(O^k)^\vee\) the rank-\(K\) submodule spanned by \(u_{r+1}^\vee(a),\ldots,u_k^\vee(a)\). Since \(H_aH_a^{-1}=I_k\), \(\widetilde W_a^\perp\) annihilates \(\widetilde W_a\) identically over \(O\). Expanding \(s_{g'}\) on the twisted rows gives
\[
 s_{g'}(v^{(g)}_a)=v_0(a)+\omega v_1(a),\qquad
 s_{g'}(w^{(g)}_a)=w_0(a)+\omega w_1(a),
\]
where \(v_0(a)=\sum_{b=1}^r X^{\rm vec}_{a,b}u_b(a)\in\widetilde W_a\otimes_O O[X]\), \(v_1(a)=\sum_{b=r+1}^k X^{\rm vec}_{a,b}u_b(a)\), \(w_0(a)=\sum_{b=r+1}^k X^{\rm cov}_{a,b}u_b^\vee(a)\in\widetilde W_a^\perp\otimes_O O[X]\), and \(w_1(a)=\sum_{b=1}^r X^{\rm cov}_{a,b}u_b^\vee(a)\). All entries lie in \(O[X]\).

We record two congruences in \(O[X]\) for any integer \(e\ge1\). First, because \(p=\omega^2\) and \(p\ge5\), every term of the binomial expansion of \((a+\omega b)^p-a^p\) is divisible by \(p\omega=\omega^3\) or \(\omega^p\), hence by \(\omega^2\). Induction on \(e\) gives \((a+\omega b)^{p^e}\equiv a^{p^e}\pmod{\omega^2 O[X]}\). Second, for any finite sum in \(O[X]\), \((\sum_\alpha c_\alpha)^{p^e}\equiv\sum_\alpha c_\alpha^{p^e}\pmod{pO[X]}\), and in \(W(\mathbf F_{p^{n'}})\) one has \(\sigma^e(c)\equiv c^{p^e}\pmod p\). Since \(p\) is odd, raising the Leibniz expansion of a minor to \(p^e\) preserves every permutation sign and yields
\[
 \bigl(\det((H_y)_{S,I})\bigr)^{p^e}
 \equiv\det\bigl(((H_y)_{S,I})^{[p^e]}\bigr)
 \equiv\det\bigl((H_{y+e})_{S,I}\bigr)\pmod{pO}
\]
for every \(h\)-subsets \(S,I\subseteq\{1,\ldots,k\}\), and the same congruence holds with \((H_y^{-1})^{\mathsf t}\) in place of \(H_y\).

For part~(1), integrality \(K_c(g;J)\in O[X]\) is immediate from the definition, and the telescoping product of powers after step \(i\) gives multidegree \(p^{z-x_i}\) at each used letter \(i\). Suppose first that \(c>0\). For \(0\le h\le k\) and \(a\in\Z\), let \(M_{a,h}\subseteq\bigwedge^h(O[X]^k)\) be the \(O[X]\)-span of \(u_S(a)=u_{s_1}(a)\wedge\cdots\wedge u_{s_h}(a)\) over subsets \(S\) of size \(h\) with \(S\subseteq\{1,\ldots,r\}\) if \(h\le r\) and \(\{1,\ldots,r\}\subseteq S\) if \(h>r\). Writing any \(A\in M_{y,h}\) as \(\sum_S f_S u_S(y)\), the two congruences above show that for every \(e\ge1\) and \(B\in\bigwedge^h(O[X]^k)\), the coordinatewise \(p^e\)-th power satisfies
\[
 (A+\omega B)^{[p^e]}\equiv A^{[p^e]}\equiv\widetilde A\pmod{\omega^2 O[X]},
 \qquad \widetilde A=\sum_S f_S^{p^e}u_S(y+e)\in M_{y+e,h}.
\]
If no letter is used in \(c\), then \(K_c(g;\varnothing)=1\) and \(\mathrm{oc}(c)=0\), so (1a)--(1d) are immediate. Otherwise we prove by induction over the used letters \(i_1<\cdots<i_m\) (\(m\ge1\)) in column \(c\) that after step \(i_t\) (at position \(y_t=x_{i_t}\) and height \(h'_t\)), letting \(d_t=\sum_{s\le t,\ i_s\text{ costly}}p^{y_t-y_s}\), one has \(s_{g'}(K^{(i_t+1)})=\omega^{d_t}Q_t\) with \(Q_t\in M_{y_t,h'_t}+\omega\bigwedge^{h'_t}(O[X]^k)\). Before the first used letter, \(h_0=0<r\) and \(K^{(0)}=1\in M_{y,0}\); the first used step must rise from \(0\) to \(1\), is not costly, and gives \(v_0(y_1)+\omega v_1(y_1)\in M_{y_1,1}+\omega O[X]^k\). At step \(i_t\) (\(t\ge2\)) with \(e=y_t-y_{t-1}\ge0\), raising each coordinate of \(s_{g'}(K^{(i_{t-1}+1)})=\omega^{d_{t-1}}Q_{t-1}\) to \(p^e\) produces \(\omega^{d_{t-1}p^e}Q_{t-1}^{[p^e]}\).
Write \(\iota_w\) for right contraction (matching the coordinate formula in the statement).
\begin{itemize}
\item If step \(i_t\) is a non-costly vector rise (\(h'_{t-1}<r\)), then \(Q_{t-1}^{[p^e]}=\widetilde A+\omega C\) with \(\widetilde A\in M_{y_t,h'_{t-1}}\), and right wedging with \(v_0(y_t)+\omega v_1(y_t)\) yields \(\widetilde A\wedge v_0(y_t)+\omega(\cdots)\) with \(\widetilde A\wedge v_0(y_t)\in M_{y_t,h'_t}\).
\item If step \(i_t\) is a costly vector rise (\(h'_{t-1}\ge r\)), the hypothesis gives \(e\ge1\), so \(Q_{t-1}^{[p^e]}=\widetilde A+\omega^2 C\) with \(\widetilde A\in M_{y_t,h'_{t-1}}\). Because every basis wedge in \(M_{y_t,h'_{t-1}}\) already contains \(u_1(y_t)\wedge\cdots\wedge u_r(y_t)\) and \(v_0(y_t)\in\widetilde W_{y_t}\otimes_O O[X]\), one has \(\widetilde A\wedge v_0(y_t)=0\) identically in \(O[X]\). Thus \((\widetilde A+\omega^2 C)\wedge(v_0(y_t)+\omega v_1(y_t))=\omega(\widetilde A\wedge v_1(y_t)+\omega(\cdots))\), increasing the \(\omega\)-order by \(1\) to \(d_t=d_{t-1}p^e+1\) with \(\widetilde A\wedge v_1(y_t)\in M_{y_t,h'_t}\).
\item If step \(i_t\) is a non-costly covector fall (\(h'_{t-1}>r\)), then \(Q_{t-1}^{[p^e]}=\widetilde A+\omega C\) with \(\widetilde A\in M_{y_t,h'_{t-1}}\). Since \(w_0(y_t)\in\widetilde W_{y_t}^\perp\otimes_O O[X]\) annihilates \(\widetilde W_{y_t}\), right contracting with \(w_0(y_t)+\omega w_1(y_t)\) preserves the factor \(u_1(y_t)\wedge\cdots\wedge u_r(y_t)\) in the leading term, giving \(\iota_{w_0(y_t)}(\widetilde A)\in M_{y_t,h'_t}\) modulo \(\omega\).
\item If step \(i_t\) is a costly covector fall (\(h'_{t-1}\le r\)), the hypothesis gives \(e\ge1\), so \(Q_{t-1}^{[p^e]}=\widetilde A+\omega^2 C\) with \(\widetilde A\in M_{y_t,h'_{t-1}}=\bigwedge^{h'_{t-1}}(\widetilde W_{y_t}\otimes_O O[X])\). Because \(w_0(y_t)\) annihilates \(\widetilde W_{y_t}\), one has \(\iota_{w_0(y_t)}(\widetilde A)=0\) identically in \(O[X]\). Right contracting with \(w_0(y_t)+\omega w_1(y_t)\) gives \(\omega(\iota_{w_1(y_t)}(\widetilde A)+\omega(\cdots))\), increasing the \(\omega\)-order by \(1\) to \(d_t=d_{t-1}p^e+1\) with \(\iota_{w_1(y_t)}(\widetilde A)\in M_{y_t,h'_t}\).
\end{itemize}
After the last used step \(y_{\rm last}=y_m\), raising each coordinate of \(\omega^{d_m}Q_m\) to \(p^{z-y_{\rm last}}\) gives \(\omega^{d_m p^{z-y_{\rm last}}}Q_m^{[p^{z-y_{\rm last}}]}\), where \(d_m p^{z-y_{\rm last}}=\mathrm{oc}(c)\). This proves (1b) for \(c>0\). For \(c\le0\), replacing \((O^k,\widetilde W_a,r)\) by \(((O^k)^\vee,\widetilde W_a^\perp,K)\) and height \(h_t\) by dual height \(k-h_t\) interchanges vector and covector steps and gives the identical induction.

To prove (1c), specialize the preceding induction to \(g'=g\), where \(H=I_k\), so \(v_0(a),w_0(a)\) are the nonnormal coordinate tuples and \(v_1(a),w_1(a)\) are the normal coordinate tuples. The four bullets show that at each used step \(i_t\), the scaled reduction \(\overline Q_t\in(O/\omega O)[X]\) is obtained from \(\overline Q_{t-1}^{[p^e]}\) by applying the nonnormal update at non-costly steps and the normal update at costly steps. Because \(O/\omega O\) has characteristic \(p\) and the coefficients of \(K^{\rm sep}_c\) lie in \(\mathbf F_p\), raising any polynomial in \(\mathbf F_p[X]\subseteq(O/\omega O)[X]\) to the \(p^e\)-th power coincides with raising every variable entrywise to the \(p^e\)-th power. Induction over the used steps therefore gives \(\overline Q_t=\operatorname{pow}_{y_t}(\overline{K^{\rm sep,(i_t+1)}_c})\), and the final raising by \(p^{z-y_{\rm last}}\) gives \(\overline{Q_m^{[p^{z-y_{\rm last}}]}}=\operatorname{pow}_z(\overline{K^{\rm sep}_c(J)})\). Multiplying by \(\omega^{\mathrm{oc}(c)}\) proves (1c). For (1d), the strict-increase hypothesis forces all costly positions in \(c\) to be distinct integers in \([x_0,z]\), so \(\mathrm{oc}(c)\le\sum_{m=0}^{z-x_0}p^m\le\frac{p}{p-1}p^{z-x_0}\le 2p^{z-x_0}\).

For part~(2), each \(x-y_a\ge1\), so the congruences above applied to each row \(s_{g'}(v^{(g)}_{y_a})=v_0(y_a)+\omega v_1(y_a)\) give
\[
 s_{g'}(A_x(g))\equiv T_x H_{x,[1..r]}\pmod{\omega^2 O[X]},
\]
where \((T_x)_{a,b'}=(X^{\rm vec}_{y_a,b'})^{p^{x-y_a}}\) (\(1\le a,b'\le r\)) and \(H_{x,[1..r]}\) consists of the first \(r\) rows of \(H_x\). Modulo \(\omega\), the first \(r\) rows of \(s_{g'}(B_{x,b}(g))\) and its \((r+1)\)-st row \(v_0(x)\) all lie in the \(r\)-dimensional subspace \(\widetilde W_x\otimes(O/\omega O)[X]\), so they are linearly dependent and \(s_{g'}(B_{x,b}(g))\equiv0\pmod{\omega O[X]}\). Likewise, modulo \(\omega\) every row of \(s_{g'}(A_x(g))\) lies in \(\widetilde W_x\otimes(O/\omega O)[X]\) while \(s_{g'}(w^{(g)}_x)\equiv w_0(x)\in\widetilde W_x^\perp\otimes(O/\omega O)[X]\), so \(s_{g'}(A_x(g)(w^{(g)}_x)^{\mathsf t})\equiv0\pmod{\omega O[X]}\), hence \(s_{g'}(P_{x,b}(g))\equiv0\pmod{\omega O[X]}\). Under \(s_g\) (where \(H_x=I_k\)), \(s_g(M_x(g))=T_x\) exactly, \(s_g(D_x(g))=\det T_x\), and for \(c>r\) entry \((a,c)\) of \(s_g(A_x(g))\) is \(\omega^{p^{x-y_a}}(X^{\rm vec}_{y_a,c})^{p^{x-y_a}}\). Expanding \(s_g(B_{x,b}(g))\) along its last column gives \(\omega s_g(D_x(g))X^{\rm vec}_{x,b}\) plus terms of the form \(\pm\omega^{p^{x-y_a}}(X^{\rm vec}_{y_a,b})^{p^{x-y_a}}\) times nonnormal cofactors; expanding \(s_g(P_{x,b}(g))\) gives \(\omega s_g(D_x(g))X^{\rm cov}_{x,b}\) plus terms \(\sum_{a,c>r}(\operatorname{adj}T_x)_{b,a}\omega^{p^{x-y_a}}(X^{\rm vec}_{y_a,c})^{p^{x-y_a}}X^{\rm cov}_{x,c}\). Since \(p^{x-y_a}\ge p\ge5\), both congruences hold modulo \(\omega^2 O[X]\) and every error term has normal degree equal to its \(\omega\)-order.

For part~(3), the contraction \(E_x=\sum_{c=1}^k X^{\rm vec}_{x,c}X^{\rm cov}_{x,c}\) is invariant under \(g_x\) and \(g'_x\), while \(\operatorname{diag}(I_r,\omega I_K)\operatorname{diag}(\omega I_r,I_K)=\omega I_k\); hence \(s_{g'}(E_x)=\omega\sum_{c=1}^k X^{\rm vec}_{x,c}X^{\rm cov}_{x,c}\) identically for every \(g'\). Raising to \(P=p^e\) gives the stated identity, and the multinomial congruence modulo \(p=\omega^2\) gives the congruence modulo \(\omega^{P+2}O[X]\).

For part~(4), since \(O/\omega O\) is a field, the ideal \(\omega O[X]\) is prime in \(O[X]\), so the Gauss valuation \(v_\omega\) on \(O[X]\) is additive on products. Because \(s_{g'}\) is a ring homomorphism, \(\operatorname{ord}_{g'}(FG)=\operatorname{ord}_{g'}(F)+\operatorname{ord}_{g'}(G)\) for all \(F,G\in O[X]\). Multiplying the lower bounds from parts~(1)--(3) with the order \(m\) of \(\omega^m\), order \(\ge0\) of integral nonnormal factors, and order \(\ge c_0\) of the localizer yields the stated bound.
\end{proof}

\begin{lemma}[The plane count and the retained rank]\label{B5.20h}
Use \ref{B3.52a} with \(p\ge13\), \(2\le r<k\), and a marked set
\(S\) of size \(N\) on an even cycle of length \(n'\). Every cyclic
window of length \(L_2\) has at most \(r-1\) marks, and every window
of length \(G'\) has at most \(N_1\). Let
\(u_1=\lceil1/\delta\rceil\) for the positive marking parameter
\(\delta\). Divide the cycle into consecutive full tiles of length
\((6r+1)u_1\), and partition each full tile into \(6r+1\) cells of
length \(u_1\). In a marked tile the even cells each have one even
mark and the odd cells have none. In the cyclic order of full
tiles every run of \(M_0+1\) full tiles meets a marked tile, where
\(G'=(M_0+5)(6r+1)u_1\). Suppose \(1\le L_1\), \(2L_1\le L_2<u_1\),
and \(3G'\le n'\).

Let \(a_0\ge1,\ Y\ge r,\ c_b,c_t,D,c_*\in\N\), with \(a_0,Y\)
allowed to be real. The dominant base weight has unspecified
entries of absolute value at most \(c_b\) and marked extremes
\((a_s,-b_s)\) with
\(a_0+c_b\le a_s\le2a_0-c_b\), \(0\le b_s\), and \(b_s+2c_b\le c_t\).
Put \(o=rNc_t+2rn'c_b\). For \(C\) above an absolute threshold
assume
\[
\begin{gathered}
 N_1^2\le rN,\quad N_1+3r\le k,\quad CN_1^2\log a_0\le k,\quad
 256k^2NY^2\le p^{L_1-1},\\
 CNk^2\log_p a_0\le D,\quad 8(D+1)\le c_*,\quad
 Ck^3r^2Np^{5G'}(c_*+o)Y\le a_0 .
\end{gathered}
\]
At every target the test passes with denominator \(D\), order \(c_*\),
fraction \(\min(1,CNr^2k/Y)\), and integer depth at least
\[
 CrN\left(CN_1^2\log(a_0)\frac{a_0}{pk}
 +a_0p^{-L_1}+c_t+kp^{G'}\right)+2rn'c_b .
 \tag{5.20.15}\label{eq:p5c2-plane-depth}
\]
\end{lemma}

\begin{proof}[Outline]
Counts segment symbols via \ref{B5.20l}, \ref{B5.20m}, and \ref{B5.20n}, lifting them with \ref{B5.20d} and \ref{B5.20he}. Builds dual tests via isotypic projections and \ref{B5.20ha}, \ref{B5.20hb}, \ref{B5.20hc}, \ref{B5.20hd}, then sums over edges using \ref{B5.20e}. The complete proof is in \S\ref{prf:B5.20h}.
\end{proof}

\begin{lemma}[The plane coefficient after normal projection]\label{B5.20hb}
Use the segment construction and size hypotheses of Lemma~\ref{B5.20h}. In segment \(j\), list the letters after the left frame in chronological order as \(x_0,\ldots,x_{l+r-1}\), where \(l\) is the number of interior letters. A normal prefix of length \(u\) consists of the first \(u\) letters of this list. Fix a simultaneous normal isotypic projection, and denote the first entry of its highest weight at this prefix by \(k_{ju}\). Put
\[
 s_u=\left\lfloor\frac{b_{x_u}}{Yk^2}\right\rfloor .
\]
Let \(H_{ju}\) be the maximum of the following numbers, including the last two entries under the indicated conditions:
\[
 \frac{b_{x_u}}2,\qquad
 k_{j,u+1}+m_*+s_u\quad(x_u\text{ even}),\qquad
 k_{j,u-1}+m_*+s_u\quad(u>0,\ x_{u-1}\text{ odd}).
\]

At each physical row first select the coordinate residues modulo its non-left-frame power \(M_x\), divide out the corresponding monomial, and replace \(X^{M_x}\) by \(X\) coordinatewise. Select the plane and normal row degrees, apply the joint projection in the normal variables, and take any normal monomial coefficient. Write \(F\) for the resulting plane polynomial. This construction applies to the flat component of every primitive combination in Lemma~\ref{B5.20h}, including after division by its scalar of exact order.

Number the shared frame slots in reverse chronological order, starting at \(1\). Write \(P_t=p^{z_{j+1}-z_j}\) for the power at a boundary slot \(t\) of frame \(j\). A boundary digit component of \(F\) with label \(\rho\) keeps exponents \(\rho_t+P_t\beta_t\), coordinate by coordinate, and writes \(\beta_t\) on a new vector letter. Interior exponents are unchanged. Then the following bounds hold.
\begin{enumerate}
\item For \(0\le u\le l\), on the preceding frame's quotient letters together with the first \(u\) interior letters, every boundary digit component has least signed weight at least \(H_{ju}-s_u/8\).
\item For \(1\le m<r\), put \(u=l+r-m\). On the original plane alphabet, the least signed weight of \(F\) on the interior and slots \(h>m\) of frame \(j\) is at least \(H_{ju}-3s_u/8\).
\end{enumerate}
Here the least signed weight is taken after a generic change of basis on the indicated letters, as in Lemma~\ref{B5.20h}. In particular, if \(F=P_c\) is normalized by \(v_{\rm F}(P_c)=0\), \(m_0\) attains this minimum, and \(P_{\rm link}\) is the product of the projections for the second set of bounds, then
\[
 \langle P_c,P_{\rm link}X^{m_0}\rangle
       =[X^{m_0}]P_c\,m_0!,
 \qquad
 v\bigl(\langle P_c,P_{\rm link}X^{m_0}\rangle\bigr)
       =\tfrac12v(m_0!).
\]
\end{lemma}
\begin{proof}[Outline]
Establishes the chain and link bounds on undifferentiated symbol products by separating shared boundary letters and bounding left-frame degree loss. Propagates both bounds by induction through residual operators from \ref{B5.20h}, then uses self-adjoint link projectors. The complete proof is in \S\ref{prf:B5.20hb}.
\end{proof}

\begin{lemma}[An integral preimage preserving the links]\label{B5.20hc}
In the alphabet of Lemma~\ref{B5.20hb}, let \(I_j\) denote the interior of segment \(j\), and let \(d_{I_j}=\sum_{i\in I_j}d_i\) for the prescribed plane row degrees \(d_i\). At each shared slot \(t\) introduce a label letter \(R_t\), a quotient letter \(W_t\), and a surplus letter \(S_t\), all vectors. The recombination \(\Phi\) sends their coordinates to \(X_t,X_t^{P_t},X_t\), respectively, coordinatewise, and fixes the interiors. Call a monomial allowable if its recombination has the prescribed row degrees, each coordinate exponent on \(R_t\) is less than \(P_t\), and
\[
 \deg S_{j,1}=0,\qquad \deg S_{j,h}\le d_{I_j}\quad(1\le h\le r).
\]
An allowable polynomial is supported on allowable monomials.

Let \(T\) be an integral polynomial on the original plane alphabet of the prescribed multidegree, with least weight at least \(\theta_{jm}\) on \(I_j\) and the shared slots \(h>m\), for every \(j\) and \(1\le m<r\). There is an integral allowable \(Z\) with \(\Phi Z=T\). Its least weight on
\[
 \{W_{j-1,h}:1\le h\le r\}\ \cup\ I_j\ \cup\
 \{S_{j,h}:h>m\}
\]
is at least \(\theta_{jm}\) for every \(j,m\).
\end{lemma}
\begin{proof}[Outline]
Builds an integral basis of path polynomials via \ref{B5.20l}, \ref{B5.20h}, and \ref{B5.20m} whose retained paths span each link cut. Reassigns interior-using boundary columns to surplus letters and splits label exponents into quotients and remainders. The complete proof is in \S\ref{prf:B5.20hc}.
\end{proof}

\begin{lemma}[The detecting pairing on a boundary digit block]\label{B5.20hd}
Use the alphabets of Lemmas~\ref{B5.20hb} and~\ref{B5.20hc}. Normalize \(v(p)=1\). For an exponent array \(b\) write \(b!=\prod_i b_i!\), and use the bilinear form for which the monomials are orthogonal and \(X^b\) has square \(b!\). Write \(D_fX^b=f(b)X^b\) for a diagonal operator. Let \(F\) have the prescribed original row degrees and \(v_{\rm F}(F)=0\). Suppose all its boundary digit components satisfy the chain bounds \(\theta_{ji}\) of Lemma~\ref{B5.20hb}(1). Denote the row factorial unit \(u\) of Lemma~\ref{B5.20h} by \(u(b)\), and put
\[
 \Lambda(Z)=\langle F,D_u\Phi Z\rangle,\qquad
 n_0=\tfrac12v\Bigl(\prod_i d_i!\Bigr),\qquad
 k'=\sum_i r(1+\log_p(d_i+r)).
\]

Fix an integer \(c'\ge1\). For a chain cut \(i\) in segment \(j\), let \(\sigma_{ji}\) be a slack satisfying
\[
 r\,d_{I_{j-1}}/P_t+13rc'\le\sigma_{ji}
 \qquad(t\text{ in frame }j-1).
\]
Let \(P_{\rm pre}\) be the product of the factorial-orthogonal prefix projections on the chain cuts \(\{W_{j-1,h}:1\le h\le r\}\cup\{x_0,\ldots,x_{i-1}\}\) (with \(x_0,\ldots,x_{i-1}\) the first \(i\) interior letters of segment \(j\)) at thresholds \(\theta_{ji}-\sigma_{ji}\), with total loss at most \(D_m\) in \(v_{\rm F}\). On an enlarged monomial with label exponents \(\rho\) and quotient exponents \(\beta\), let \(\alpha\) keep its interior exponents and have boundary exponents \(\rho_t+P_t\beta_t\). Define
\[
 \Gamma=u(\alpha)\prod_{t,a}
 \frac{(\rho_{ta}+P_t\beta_{ta})!}
      {\rho_{ta}!\,\beta_{ta}!\,(P_t!)^{\beta_{ta}}}.
\]
Let \(g_+,g_-\) be integral diagonal multipliers congruent on allowable monomials to \(\Gamma,\Gamma^{-1}\), respectively, modulo \(p^{c'}\). For every integral allowable \(Z\),
\[
 v\!\left(\Lambda(D_{g_-}P_{\rm pre}D_{g_+}Z)-\Lambda(Z)\right)
 \ge c'+n_0-2k'-D_m .
\]
\end{lemma}
\begin{proof}[Outline]
Works blockwise, rewriting \(\Lambda\) via divided derivatives, approximating \(D_u,D_{u^{-1}}\) by \ref{B5.20ha}, and expanding boundary derivatives in quotient variables. Self-adjointness of \(P_{\rm pre}\) and carry bounds from \ref{B3.19} control the error. The complete proof is in \S\ref{prf:B5.20hd}.
\end{proof}


\begin{definition}[The lattice at an internal tie]\label{B5.20o}
Work over \(\Q_p\). On signed letters \(X_i=(X_{i1},\ldots,X_{ir})\) of degrees \(d_i\) use the form
\[
 [X^\alpha,X^\beta]_0=\mathbf1_{\alpha=\beta}
 \frac{\prod_{i,c}\alpha_{ic}!}{\prod_i d_i!}.
\]
The action on a prefix of length \(j\) is generated by
\[
 E_{ab}^{j}=\sum_{\substack{i\le j\\i\ {\rm vector}}}
 X_{ia}\partial_{ib}
 -\sum_{\substack{i\le j\\i\ {\rm covector}}}X_{ib}\partial_{ia}.
\]
A prefix with threshold \(M_j\) contributes the largest \(E_{ab}^{j}\)-stable subspace in the span of monomials of signed first-coordinate weight at most \(M_j\). An absent threshold means \(+\infty\).

Take a prefix of \(m+1\) letters ending in a vector \(W\) of degree \(o\), and put \(S=\sum_i d_i\), \(\rho=r\). For an integral weight \(w\), intersect the imposed prefix spaces with the full diagonal weight space \(w\) and with the kernels of \(E^{m+1}_{ab}\) for \(a<b<r\). Intersect with the integral coefficient lattice to obtain \(M_w\). Put \(x=W_\rho\),
\[
 M_{\ge e}=\{f\in M_w:x^e\mid f\},\qquad
 G_e=\{x^e[x^e]f:f\in M_{\ge e}\}.
\]
Here \([x^e]\) keeps all other variables. Thus \(G_e\) is the image lattice of the restriction map, with the power \(x^e\) restored.
\end{definition}

\begin{lemma}[Denominators for types and highest vectors]\label{B5.20t}
Let \(A\) be a finite-dimensional split integral-weight \(\mathfrak{gl}_h\)-module over \(\Q_p\), with semisimple diagonal action. Let \(L\) be a full \(\Z_p\)-lattice preserved by weight extraction and by every divided root operator \(E_{ab}^{(j)}=E_{ab}^j/j!\), \(a\ne b\). Suppose that, for any two weights occurring in \(A\), their entries at each fixed coordinate differ by at most \(l\). Put
\[
 d=\binom h2\lfloor\log_p(l+h-1)\rfloor .
\]
For \(h=1\) put \(d=0\). Projection onto any sum of isotypic components takes \(L\) into \(p^{-d}L\). In an isotypic component of highest weight \(\lambda\), let \(\Delta_\lambda\) be the integral span of divided-root words applied to its integral highest vectors. Then
\[
 p^d(L\cap A_\lambda)\subseteq\Delta_\lambda .
\]
\end{lemma}
\begin{proof}[Outline]
Extract weight zero and use the row recursion to bound the denominator of the invariant projection on maps. Averaging rank-one maps gives the type projectors, and averaging highest-weight extraction gives the divided-root generation bound. The complete proof is in \S\ref{prf:B5.20t}.
\end{proof}

\begin{lemma}[Restriction at the last vector]\label{B5.20u}
Fix the first \(m\) signed letters and their degrees as in \ref{B5.20o}, and let the last letter be a vector \(W\) in \(r\ge2\) coordinates. For each \(d\ge0\), let \(Y_d\) be the rational polynomial space of last degree \(d\) satisfying all the prefix caps, before imposing a weight or a highest-vector condition. Write \(R\) for evaluation at \(W_r=0\), and \(I=\{1,\ldots,r-1\}\).

Define \(X_d\) on the hyperplane \(W_r=0\) as follows. Keep every cap on a proper prefix. For each cap \(C\) containing \(W\), impose the cap \(C\) on the first \(m\) letters for \(\mathfrak{gl}_r\), and on all \(m+1\) letters for \(\mathfrak{gl}(I)\). The latter condition retains the types with first highest-weight entry at most \(C\). Then
\[
 RY_d=X_d .
\]
If \(S\) is the sum of the first \(m\) degrees and \(d\), and \(L\) denotes the integral coefficient lattice in each polynomial space, then
\[
 p^{2b}(X_d\cap L)\subseteq R(Y_d\cap L),
 \qquad b=\binom r2\lfloor\log_p(4S+r-1)\rfloor .
\]
\end{lemma}
\begin{proof}[Outline]
Differentiate to check the caps under restriction, then prove surjectivity by induction on degree using a highest-vector kernel calculation. Apply \ref{B5.20t} twice for a larger Lie algebra action whose highest vectors lift in degree zero. The complete proof is in \S\ref{prf:B5.20u}.
\end{proof}

\begin{lemma}[The graded image and its pairing partners]\label{B5.20v}
Use \ref{B5.20o} with \(r\ge2\), and put
\[
 e_1=\binom r2\lfloor\log_p(S+r-1)\rfloor,\qquad
 e_2=\binom r2\lfloor\log_p(4S+r-1)\rfloor .
\]
For \(0\le e\le o\), use the space \(X_{o-e}\) of \ref{B5.20u} after subtracting \(e\) from each cap containing \(W\). In this space take the full weight \(w-e\epsilon_r\) and the highest vectors for \(\mathfrak{gl}_{r-1}\). Denote the result by \(H\). Then \(\Q_pG_e=W_r^eH\). If \(m>0\), then
\[
 p^{e_1+2e_2}\bigl(\Q_pG_e\cap P_{\Z_p}\bigr)\subseteq G_e .
\]
The orthogonal projection onto \(\Q_pG_e\) for \([\ ,\ ]_0\) takes \(P_{\Z_p}\) into \(p^{-(m+1)e_1}P_{\Z_p}\). For \(m=0\) both inclusions hold with exponent zero. For \(e>o\) the graded image is zero.
\end{lemma}
\begin{proof}[Outline]
Divide by \(W_r^e\) and use \ref{B5.20u} with the shifted caps to obtain the restricted image and lift. Use \ref{B5.20t} to bound denominators for the highest-vector projection of the lift and for the commuting projections that give the orthogonal projector. The complete proof is in \S\ref{prf:B5.20v}.
\end{proof}

\begin{lemma}[Detecting a primitive vector of the graded lattice]\label{B5.20p}
In \ref{B5.20o}, write \(V_e=\Q_pG_e\), and let \(P_{\Z_p}\) be the integral coefficient lattice in the specified multidegree. Suppose
\[
 p^c(V_e\cap P_{\Z_p})\subseteq G_e .
\]
Suppose also that, for every monomial \(X^\alpha\) of that multidegree, there is \(y_\alpha\in V_e\) with \(p^D y_\alpha\in P_{\Z_p}\) and
\([v,y_\alpha]_0=[v,X^\alpha]_0\) for all \(v\in V_e\). In any basis of \(G_e\) its Gram determinant is nonzero and has valuation at most
\[
 (\operatorname{rank}G_e)(D+2c).
\]
\end{lemma}
\begin{proof}[Outline]
Use the lattice inclusion to bound the least coefficient valuation by \(c\), then pair with its scaled monomial partner to get a factorial pairing of valuation at most \(D+2c\). Apply this test in Smith bases to bound the diagonal valuations and the determinant. The complete proof is in \S\ref{prf:B5.20p}.
\end{proof}

\begin{lemma}[Leading determinant at an internal tie]\label{B5.20q}
Use \ref{B5.20o}, with \(r\ge2\). Adjoin a covector \(Y\) of degree \(n\ge S\). Write \(\eta=Y_\rho\), \(E=\sum_c W_cY_c\), and \(F=\sum_c\partial_{W_c}\partial_{Y_c}\). For a nonnegative shift \(z\) put
\[
 T_z=\sum_{l=0}^o
 \frac{(-1)^l E^lF^l}{(l!)^2\binom{o+n+r+z-1}{l}} .
\]
Define \(f[n]\) by the following substitution on the prefix, followed by multiplication by \(\eta^{n-S}\). On a vector \(V\), replace \(V_\rho\) by \(\sum_cV_cY_c\) and \(V_c\) by \(\eta V_c\) for \(c\ne\rho\). On a covector \(U\), replace \(U_\rho\) by \(\eta U_\rho\) and \(U_c\) by \(\eta U_c-U_\rho Y_c\) for \(c\ne\rho\). Write \(\langle\ ,\ \rangle\) for the unnormalized factorial form, with monomial diagonal \(\alpha!\).

Choose a basis \(f_b\) of \(M_w\) adapted to \(M_{\ge e}\). Let \(e_b\) be the order of \(f_b\) in \(x\) and \(g_b=x^{e_b}[x^{e_b}]f_b\). For each \(e\), write \(d_e=\#\{b:e_b=e\}\) and \(H_e=([g_b,g_{b'}]_0)_{e_b=e_{b'}=e}\). For an integer \(u\ge0\), the matrix with values
\[
 A_{bb'}(n)=
 \frac{p^u\binom{o+n+r+z-1}{o}}
 {(n-S)!\,S!\prod_i d_i!}\,
 \langle T_z f_b[n], f_{b'}[n]\rangle
\]
is polynomial in \(n\). If the \(H_e\) are nonsingular, its determinant has degree and leading coefficient
\[
 s_A=\sum_b(S+o-e_b),\qquad
 b_A=\prod_{e:d_e>0}
 \left(\frac{p^u e!\binom{z+e}{e}}{o!\,S!}\right)^{d_e}\det H_e .
\]
\end{lemma}
\begin{proof}[Outline]
Expand in filtration bases and get leading coefficients from the real harmonic eigenvalues of \(T_z\). Orthogonality and Cauchy–Schwarz give strict mixed-order degree bounds, so over \(\mathbb Q_p\) the leading determinant coefficient is the product of block determinants. The complete proof is in \S\ref{prf:B5.20q}.
\end{proof}

\begin{lemma}[Increasing the covector degree]\label{B5.20r}
Use the prefix and substitution of \ref{B5.20q}. Fix \(0\ne f\in M_w\), the degrees on the remaining letters, and all prefix thresholds. Let \(Z_n\) be the intersection of the threshold spaces with the span of products \(v b\), where \(v\) is in the module generated by \(f[n]\) on the prefix including \(Y\), and \(b\) is a polynomial in the remaining letters of the specified multidegree. If \(S\le n\le n'\), then \(\dim Z_n\le\dim Z_{n'}\).
\end{lemma}
\begin{proof}[Outline]
Multiply by \(Y_\rho^{n'-n}\) and project with the prefix Casimir to map the highest line onto the new highest line. The map shifts every weight by the same character and commutes with the raising operators. A maximal weight in a nonzero kernel would therefore give a highest vector. The unique highest line of the source maps nontrivially, so the map is injective. Extending it coefficientwise over the suffix monomials preserves the prefix thresholds. The complete proof is in \S\ref{prf:B5.20r}.
\end{proof}

\begin{lemma}[Integral contraction of the middle rows]\label{B5.20s}
Let a marked factor have \(k\) rows and put \(h=k-2\). Write \(c\) for the bound on its middle entries. Write \([P,Q]=(P,0,\ldots,0,-Q)\) for a flat weight. Suppose \(\lambda\) has extremes \(P+c,-Q-c\), where \(P,Q\ge0\), and put \(\delta=\lambda-[P,Q]\). In a selected edge frame, there are maps
\[
 i:V_{[P,Q]}\longrightarrow V_{[P+2c,Q+2c]},\quad
 U_b:V_{[P,Q]}\longrightarrow V_\lambda,\quad
 L_a:V_\lambda\longrightarrow V_{[P+2c,Q+2c]}
\]
indexed by \(\dim V_{(\lambda_2,\ldots,\lambda_{k-1})}\) middle directions, with \(L_aU_b=\delta_{ab}p^{k_b}i\). In the range \(2c\le 8a_1,\ k\ge7\),
\[
 0\le k_b\le\binom h2\lfloor\log_p(2c+h-1)\rfloor .
\]
Multiplication by the vector defining \(i\) gives a saturated inclusion of the selected normalized edge lattices.
\end{lemma}
\begin{proof}[Outline]
Use \ref{B3.52b}, complementary minors and Smith reduction to pair the middle lattices, then lift the tableaux to define the contraction. Highest-tableau evaluation gives a unit leading coefficient for the multiplier of \(i\), so adapted-frame reduction proves saturation. The complete proof is in \S\ref{prf:B5.20s}.
\end{proof}

\begin{lemma}[From prefix Smith factors to the level lattice]\label{B5.20w}
Consider a signed word of \(L\) letters in \(r\ge2\) coordinates. Mark \(t\) disjoint adjacent vector--covector pairs, in order. Impose prefix caps ending at the vectors of marked pairs or at the end of the word. Let \(Y\) be their intersection in fixed multidegree and let \(M=Y\cap\Z_p[X]\). Use the normalized factorial form \(B_0\) of \ref{B5.20o} on the word. For pair \(i\), use the operator \(T_{z_i}\) of \ref{B5.20q} on its two letters and call it \(\Phi_i\).

Let \(X\ge1\) bound every letter degree and let \(S_{\rm tot}\) be their sum. Put
\[
 a_0=\binom r2\lfloor\log_p(S_{\rm tot}+r-1)\rfloor .
\]
Suppose \((L+2)a_0\le\kappa\), \(Lr\lfloor\log_p X\rfloor\le\kappa\), and the degree \(n_i\) of each marked covector is at least the total degree of its preceding prefix. At each such prefix and full weight \(w\), form the matrix of \ref{B5.20q} with \(u=\lceil\kappa\rceil\), using a basis of \(M_w\). Impose the caps through the marked vector in this prefix. Write \(b_{iw}\) for the number of its Smith entries of valuation greater than \(E\ge0\), counting zero as valuation \(+\infty\). For each nonzero \(M_w\) choose one nonzero \(f\in M_w\), and write \(z_{iw}\) for the dimension of the space \(Z_{n_i}\) of \ref{B5.20r} for this choice.

For the form \(B(v,v')=B_0(\Phi_1\cdots\Phi_t v,v')\) on \(M\), the number of Smith entries of valuation greater than \(\lfloor t(E+\kappa)+\kappa\rfloor\) is at most \(\sum_{i,w} b_{iw}z_{iw}\).
\end{lemma}
\begin{proof}[Outline]
Use \ref{B5.20q} to identify prefix highest vectors and retain Smith entries of valuation at most \(E\) at each tie. Extend the pairing identities using \ref{B5.20t} to clear denominators, and bound the number of large Smith entries by the codimension of the common fixed space. The complete proof is in \S\ref{prf:B5.20w}.
\end{proof}

\begin{lemma}[The rank-one localizing polynomial]\label{B5.20x}
Let \(s_j\) be the successive even starts in cyclic order modulo \(n'\), with gaps \(G_j\) and at least two starts. Write \(X_{ja}\), \(0\le a<k\), for the vector coordinates there. In the frame in which the distinguished residue line is coordinate zero, put \(h_j=p^{G_j}-1\) and fix one start \(j_0\). The polynomial
\[
 R=X_{j_0,0}\prod_{a=1}^{k-1}
       \left(\prod_jX_{j0}^{h_j}-\prod_jX_{ja}^{h_j}\right)
\]
has order at least one at every other residue line. At the distinguished line its reduction after the edge substitution is
\[
 X_{j_0,0}\prod_j X_{j0}^{(k-1)h_j}.
\]
Consequently \(R^c\) supplies order \(c\) at competing lines and uses precisely \(c((k-1)h_j+\mathbf1_{j=j_0})\) degrees at start \(j\).
\end{lemma}
\begin{proof}[Outline]
Modulo the uniformizer, edge substitution scales \(X_{ja}\) by \(u_a^{p^{s_j}}\), and \(u^{p^{n'}}=u\) gives \(\prod_j(u^{p^{s_j}})^{h_j}=1\) for \(u\ne 0\), killing \(R\) on competing lines. At \(u_0=1\) this gives the monomial, and \(R^c\) gives order \(c\) and multidegrees. The complete proof is in \S\ref{prf:B5.20x}.
\end{proof}

\begin{lemma}[The depth constant for one segment]\label{B5.20ia}
Consider a gap \(G\ge2\) with \(s\) interior positions at increasing positive odd distances \(d_i<G\). Write \(P_i=p^{G-d_i}\), \(Q\ge p^G\), \(p\ge5\), and \(k\ge2\). Let \(a_1,b',L'\) be positive, \(a_S\ge a_1\) the starting weight, \(b_i\le2b'\) the interior weights, and \(b_+\le2b'\) the terminal weight. Let \(P_iH_i\le b_i\), and let the capacities satisfy
\[
 C_0\le a_1/p^G,\qquad C_{i+1}=C_i+H_i-e_i,\qquad H_i,e_i\ge0 .
\]
At a position with \(p^{d_i}<a_1/b'\) suppose \(b_i\le a_S/p^{d_i}\) and \(e_i\le (1+4/p^2)H_i+L'C_i/k\). At the other positions suppose \(b_i\ge a_1/p^{d_i}\) and \(e_i\le C_i+L'H_i/k\). Suppose also that
\[
 o_+\le C_s+(L'/k+4/p^2)\sum_i2a_1/P_i+L'b_+/k.
\]
Let \(D_i\) be the sum of the nonzero-coordinate vector digits and the coordinate-zero covector digit, each digit in \([0,P_i)\). Put
\[
 R_*=a_1/p^2+b'/p+L'a_1/(pk)+L'b'/k+kQ .
\]
The order of the block satisfies
\[
 o_++\sum_i(P_ie_i+D_i)
 \le\sum_i\min(b_i,a_S/p^{d_i})+30(s+1)R_* .
\]
\end{lemma}
\begin{proof}[Outline]
Unroll the capacity recursion and bound the earlier contributions by a geometric sum using the spacing of odd positions. Bound the excess over the minimum in each range, then add the digit and terminal orders. The complete proof is in \S\ref{prf:B5.20ia}.
\end{proof}

\begin{lemma}[An integral binomial expansion at an internal tie]\label{B5.20ib}
Use the matrix of \ref{B5.20q}, with \(m+1\) prefix letters, and let \(X\ge1\) satisfy \(2S+1\le X\). Put \(\ell_X=1+\log_p X\). If
\[
 u\ge ((m+2)r+r+1)\ell_X ,
\]
each entry of \(A(n)\) is a linear combination over \(\Z_p\) of \(\binom nj\), \(0\le j\le S+o\).
\end{lemma}
\begin{proof}[Outline]
Use \ref{B5.20q} to bound the entry degrees and \ref{B3.19} to show the scaled values at \(S,\ldots,3S\) are integral. Take forward differences there and translate the shifted binomial basis. The complete proof is in \S\ref{prf:B5.20ib}.
\end{proof}

\begin{lemma}[Diagonal derivatives on a block with ties]\label{B5.20ic}
Take a word of vector and covector letters in \(r\) coordinates over the fraction field of a discrete valuation ring \(O\) of characteristic zero. Impose upper bounds \(b_J\) on signed first-coordinate weights at specified prefixes \(J\), taking the largest stable subspace at each prefix as in \ref{B5.20o}. Pair disjoint consecutive vector and covector letters at internal ties. Include a final vector tied to a separate covector \(\xi\). At tie \(h\) write
\[
 E_h=\sum_a W_{ha}Y_{ha},\qquad H_h=E_h+Z_h S'_h,
\]
where \(Z_h,S'_h\) are separate scalar variables, and fix a shift \(z_h\ge0\).
For fixed full row degrees let \(L(b)\) be the \(O\)-span of
\[
 \prod_h H_h^{j_h} Z_h^{z_h}\, y\,\xi^\alpha .
\]
Here \(j_h\) ranges from zero to the lesser of the paired degrees, \(y\) has the degrees obtained by subtracting \(j_h\) from both letters of each pair, with the terminal covector recorded separately by \(\xi^\alpha\), and \(y\) is integral and satisfies all the prefix bounds.

For a multi-index \(\nu\) supported on interior letter coordinates and internal scalar variables, let \(T_\nu\) act on a monomial \(X^\mu\) by multiplication by \(\prod_v\binom{\mu_v}{\nu_v}\). Then
\[
 T_\nu L(b)\subseteq L\bigl((b_J+|\nu|)_J\bigr).
\]
Furthermore, an integral polynomial in the fraction-field span of \(L(b)\) belongs to \(L(b)\).
\end{lemma}
\begin{proof}[Outline]
Extract coefficients after scaling one variable at a time to bound the increase in prefix weights, expanding the affected tie factors. For saturation, set \(Z_h=1\) and replace \(S'_h\) by \(S'_h-E_h\) to recover integral reduced coefficients. The complete proof is in \S\ref{prf:B5.20ic}.
\end{proof}

\begin{lemma}[The integral rank-one column estimates]\label{B5.20i}
Let \(S\) have \(N\) marks on an even cycle of length \(n'\), with at
least two even marks and successive even gaps at most \(G'\).
Each gap has at most \(N_1\) interior marks. Suppose \(N_1\ge1\),
\(2N_1+2\le k-1\), \(k^2Y^2\le p<a/b'\), \(Y\ge1\), \(b'>0\),
and \(p\) prime. Let \(\mathcal L=C_0N_1^2\log a\), \(\Delta\ge0\),
\(c_b,D\in\N\), and put \(c_*=8(D+1)\),
\(c_4=c_*+2Nc_b+2n'c_b\), \(\rho_{\rm loc}=2HN_1\mathcal L/k+4\Delta/b'\).
Assume
\[
\begin{gathered}
 H\le Y_b,\quad 2Y_bN\mathcal L<k,\quad 8\log(NY)\le\mathcal L,\quad
 Y_bN\mathcal L b'/k\le\Delta\le b'/2,\\
 kY^2Y_bN\mathcal L\log a\le a/b',\quad HN\mathcal L^2\le k,\quad
 \rho_{\rm loc}<1/3,\quad c_b\le\mathcal L b'/k,\\
 HNY^2(N_1^5+k^3)\log_p a\le D,\qquad
 Hk^3Np^{4G'}(c_*+(N+n')c_b)Y\le b'.
\end{gathered}\tag{5.20.17}\label{eq:p5c2-rank-one-size}
\]
Suppose \((a/b')[1-\rho_{\rm loc},1+\rho_{\rm loc}]\) contains no
odd positive power of \(p\). Take \(C_0\ge T_0\), \(H\ge T_2(C_0)\),
and \(C=C_*(C_0)\) for suitable absolute \(T_0>0\) and positive
functions \(T_2,C_*\). Use the nonempty integer intervals
\[
 I_a=[\lceil a\rceil,\lfloor a+a/(kY)\rfloor],\qquad
 I_b=[\lceil b'\rceil,\lfloor b'+\Delta\rfloor].
\]
Put \(a_1=\lceil a\rceil\), \(L'=C_0N_1^2\log a_1\),
\(J_0=L'b'/k\), and \(J_b=\lfloor J_0\rfloor\).
For \(u\ge0\) let \(P(u)=\binom{u+k-1}{k-1}\), with \(P(u)=0\)
for negative \(u\), and set
\[
 d(u,v)=P(u)P(v)-P(u-1)P(v-1),\qquad
 q(u,v,i)=d(u-i,v-i)/(P(u)P(v))
\]
for \(0\le i\le\min(u,v)\), and zero otherwise. For each local model
there are nonnegative \(f,g\) satisfying the pointwise inequality
of \ref{B5.20g} on the enlarged box with \(J=J_b,\ c=c_b\), and
\[
 \mathbb E_+g\le(C_A+C_s)NN_1/Y .
 \tag{5.20.18}\label{eq:p5c2-rank-one-mean}
\]
For each original \(A,B\) and every dominant base weight with
those extremes and other entries of absolute value at most \(c_b\),
the rank-one test passes at fraction
\(f_-+2f(A+c_b,B+c_b)+12Nkc_b/b'\), order \(c_*\), denominator \(D\),
and the ceiling of
\[
\begin{split}
 J_{\rm odd}&+CNN_1(a/p^2+b'/p+\mathcal L a/(pk)+\mathcal L b'/k)\\
 &+CN\{p^{G'}(c_*+CNk\log_p a+4Nc_b+2n'c_b)
 +kp^{G'}+k\log_p a+c_b\}+2n'c_b ,
\end{split}\tag{5.20.19}\label{eq:p5c2-rank-one-depth}
\]
where \(J_{\rm odd}\) is the odd-mark sum in \ref{B3.52a}.
\end{lemma}

\begin{proof}[Outline]
Bounds segment Gram losses via \ref{B5.20f} and \ref{B5.20w}, lifting paired block families with \ref{B5.20he}, \ref{B5.20ia}, \ref{B5.20ic}, and \ref{B5.20x}. Splits into harmonic summands via \ref{B5.20k}, then inserts middle rows by \ref{B5.20s} and \ref{B5.20e}. The complete proof is in \S\ref{prf:B5.20i}.
\end{proof}

\begin{lemma}[The comparison on one good orbit]\label{B5.20aa}
Fix a family in the design class of \ref{B8.02}. Write \(a=xb\), \(c=4S_c\), and \(L=\log x\). At a candidate prime \(p\), put \(z=\log_p(2a)\) and
\[
 H_-=\frac{\mathcal La}{pk}+ap^{-L_1^{\rm loc}},
 \qquad H_+=\frac{\mathcal La}{pk}+\frac{a}{p^2},
 \qquad L_1^{\rm loc}=\left\lceil\log_p((kx^{26})^4)\right\rceil .
\]
There is a family constant \(C_*\ge3\) for which the following holds at all sufficiently large scales. Define
\[
 \ell_p^{\rm int}=C_*(c+\eta)z+
 \begin{cases}
 C_*r_Gb'L+C_*LH_-,&p<p_0,\\
 b'/x+\eta+C_*w''L/\delta_1+b_0+C_*LH_+,&p_0\le p<X',
 \end{cases}
 \qquad b_0=\exp(x^{\varepsilon/2}).
\]
Let \(v\) be a good orbit with representative \(t_0\), let \(T=T(t_0)\) be its set of marks as in \ref{B3.46}, and fix \(1\le j\le m\). Write \(P_{v,j}=\Pi_p(t_0,j)\) for the global price in \eqref{eq:p3c5-price}. For \(s\in T\) put
\[
 A_s=\sum_{\substack{h\ge1\\h\ {\rm odd}}}
       [a_s/p^h-W_j(z_h(s))]_+ .
\]
With the depth \(t_v\) and denominator \(D_p\) of \ref{B3.52}, one has
\[
 \sum_{s\in T}(A_s-\ell_p^{\rm int})
 \le g_U(t_0,j)+\bar\delta(t_0)-P_{v,j}
       -\lceil t_v\rceil_+-D_p-1 .
 \tag{5.20.21}\label{eq:p5c2-selected-comparison}
\]
\end{lemma}

\begin{proof}[Outline]
Split the odd-power sum into positive parts and minima and bound the change to the data of \ref{B3.49}. Use \ref{B3.50} and \ref{B3.52a} to bound the costs, charging odd marks above the cutoff to distinct minima at preceding even marks by \ref{B3.47}. The complete proof is in \S\ref{prf:B5.20aa}.
\end{proof}

\begin{theorem}[Local comparison on the design class]\label{B5.20a}
Given the twist selection and family duties of \ref{B8.02}, the
following assertion is closed on \(\mathcal D_3(\mathcal A)\) in the
sense of \ref{B8.39}. For each such family \(F_m\), let \(I(x)\) index its models at scale \(x\). There are \(C\ge0\) and \(x_0\) for which the following bounds hold for every \(x\ge x_0\) and every \(i\in I(x)\).
For every word \(u\) and place \(w\) with \(p(w)\ne T_u\), the
containment of \ref{B5.05} holds at exponent
\(\lfloor e_w\tau_w\rfloor_+\), and
\(\ell_{u,w}\le d_M(\tau_w+\mathrm{io}_w)\), where \(e_w\ell_{u,w}\) is the length of \((M_{H,u})_{(w)}/(N_u)_{(w)}\) as in \ref{B5.04}, and \(d_M\) is as in \ref{B4.07a}. If \(V_p=\varnothing\), these three quantities
\(\tau_w,\mathrm{io}_w,\ell_{u,w}\) vanish.

For \(v\in V_p\), write \(N_v=|TT\cap v|\) and \(S_v\) for the
ceiling in \ref{B3.52}. Then \(S_v\le2mN_va/p\).
At a signal or inert-class prime \(p\le X''\) the gain is between
zero and \(4n_Ta/p\) per lift, and \eqref{5.16.1} holds with
\(r_p=F_p\) and
\(\ell_p=C_p+\mathbf1_{\{p\ {\rm inert\ class},\ p<X'\}}(R_p+O_p)\).
Use the row bound \(S_c\) of \ref{B4.03} and set \(N_0=N_m,\ t_0=t\),
\(L_1(p)=1+\lceil\log_p(k^2x^{78})\rceil\), and
\(b_0=b^{x^{-\varepsilon/2}}\) in \eqref{5.17.1} and \eqref{5.20.1}.
For a candidate with \(\Sigma_p\le x^{-\varepsilon}\) the selection
\(V_p\) is the full set of good representatives. On each good
candidate orbit the mean of \(\kappa_v\) over its local box is
at most \(N_vx^{-12}\), pointwise below \(p_0\).
\end{theorem}

\begin{proof}[Outline]
Extend the local tests by scalars, with loss bounds from \ref{B5.20b} and \ref{B5.20c}. Apply \ref{B4.07a} and \ref{B4.07c} to bound the local length and prove containment, then sum the orbit comparisons of \ref{B5.20aa} to bound the gain. The complete proof is in \S\ref{prf:B5.20a}.
\end{proof}

\begin{lemma}[Numerical thresholds for the weight box]\label{B5.20y}
For each \(0<\varepsilon<1/52\) there is a threshold datum \(\mathcal K_W(\varepsilon)\) in the sense of \ref{B8.39} with the following property for every family of parameter \(\varepsilon\) and every \(CH\in\R\) passing its size tests. Write \(C_0\) for the middle orbit constant multiplying \(N_1^2\log a_1\) in \ref{B5.20i}, and \(C_N\) for the last-list coefficient allowed in \(N_{\rm win}\le C_N\log x\). Then \(C_0>0\), the coefficient inequalities \eqref{eq:p5c2-family-thresholds} hold, and \(N_{\rm win}\le C_N\log x\) at all sufficiently large indexed points. The tests use the numerical lists of the family and \(CH\).
\end{lemma}
\begin{proof}[Outline]
Bound \(N_{\rm win}\) via \ref{B3.51}. Use \(T_0\) from \ref{B5.20i} for the middle threshold and maximize over \(A_\varepsilon\) and the doubled \(T_2\) and \(C_*\) values on the eligible list for the last threshold. The complete proof is in \S\ref{prf:B5.20y}.
\end{proof}

\begin{lemma}[Size bounds on good candidate orbits]\label{B5.20j}
Fix a family and a real \(CH\) passing the numerical tests of \ref{B5.20y}. Set
\(L_1^{\rm loc}=\lceil\log_p((kx^{26})^4)\rceil\),
\(N_1=N_{\rm win}\). Eventually each good candidate orbit of
length \(n'\) with \(N\) marks has even \(n'\) and
\[
 \begin{gathered}
 2L_1^{\rm loc}\le L_2\le2\lceil256\log x\rceil<u_1,\qquad
 3G'\le x^{\varepsilon/4}\le n',\\
 N\le40n'\delta_1,\qquad n'\le NC_{31}\log x/\delta_1,\qquad
 N_1\le C_N\log x .
 \end{gathered}
\]
Successive even marks give at least two cyclic segments of
length at most \(G'\) with at most \(N_1\) interior marks.
Here \(C_N\) is the window coefficient of \ref{B5.20y}, and \(C_{31}\) is a family constant.
\end{lemma}

\begin{proof}[Outline]
Bound \(L_2,G',N_1\) from parameter bounds, and exclude odd \(n'\) since an odd order divides \(o\le x^{\varepsilon/100}\), contradicting \(n'\ge x^{\varepsilon/4}\). The run condition bounds even gaps by \(G'\), giving \(n'\le NG'\), two gaps, and the \(N\) bound. The complete proof is in \S\ref{prf:B5.20j}.
\end{proof}

\begin{lemma}[Projection at rank \(r_G\)]\label{B5.20b}
Under the family hypotheses of \ref{B5.20a}, eventually every good
candidate orbit below \(p_0\) passes \ref{B3.52a} at \(r=r_G\),
depth \(\lceil t_{II}\rceil_+\), denominator \(D_{II}\), and fraction
\(\min(1,C_{\rm frac}Nr^2k/x^{26})\).
\end{lemma}
\begin{proof}[Outline]
Use the orbit bounds of \ref{B5.20j} and the family thresholds to check the hypotheses of \ref{B5.20h} with \(Y=x^{26}\). Compare its required depth and retained rank with those of the rank test. The complete proof is in \S\ref{prf:B5.20b}.
\end{proof}

\begin{lemma}[Mean projection loss at rank one]\label{B5.20c}
For every family and real \(CH\) passing the numerical tests of \ref{B5.20y}, eventually each good
candidate orbit at \(p\ge p_0\) has mean least passing fraction
in \ref{B3.52a} at most \(C_mNN_1/Y\) in its uniform integer
box, with \(N=N_v,\ N_1=N_{\rm win},\ Y=x^{13}\).
\end{lemma}

\begin{proof}[Outline]
Check the size bounds for \ref{B5.20i} using \ref{B5.20j} and the band exclusion of \ref{B3.51}, then return to the original box by \ref{B5.20g}. Transport the pairings to the local models and compare depths before averaging the bound on the least passing fraction. The complete proof is in \S\ref{prf:B5.20c}.
\end{proof}

\section{Arithmetic models and common weights}

For a model \(i\in I(x)\) define
\[
 \mathrm{IO}_i=\sum_{w\in\operatorname{supp}_i}
 \frac{e_w\log N\mathfrak p_w}{[K:\Q]}\,\mathrm{io}_{i,w}.
 \tag{5.21.1}\label{eq:p5c2-local-index}
\]
A consistent presentation uses the model's frame, layer count,
host arithmetic datum, and weighted top labels. Its \(w_j\)
reads lift \(j+1\) at the code of an embedding.

\begin{theorem}[Loss bounds for a small-share model]\label{B5.21}
For each inert-class prime \(p\), let \(b_p(TT)\) count the top positions whose \(h_p\)-orbit fails the window or tile condition of \ref{B2.54}. Suppose that for all sufficiently large \(x\) and every model \(i\in I(x)\) of \(F_m\),
\[
 \frac{\delta_1n}{2}\le n_T\le2\delta_1n,\qquad
 \sum_{\substack{p<X'\ {\rm of\ inert\ class}\\ \operatorname{ord}(h_p)\ge x^{\varepsilon/4}}}
 \frac{\log p}{p}\,b_p(TT)\le n_Tx^{-2+\varepsilon}.
\]  Assume \ref{B5.20a} with
fixed \(C_G\) and \(b_0=b^{\delta(x)}\) with \(\delta(x)\to0\).
There are \(c_{\rm Loss}\ge0,x_0\) such that for
\(x\ge x_0,\ i\in I(x),\ R_W(i)\le x^{-9}\),
\[
 \mathrm{IO}_i\le c_{\rm Loss}Q_i,\qquad Q_i=|TT_i|m_i b_i .
 \tag{5.21.2}\label{eq:p5c2-index-bound}
\]
For every consistent presentation with an integer cutoff satisfying
\[
 a_t/w_j(\rho)\le N_{\rm cut}
 \quad(j=0,\ldots,m_i-1,\ t\in TT_i,\ \rho\in\Phi),
 \tag{5.21.3}\label{eq:p5c2-cutoff}
\]
one has
\[
 \ASK_{N_{\rm cut}}-\Ga_i\le c_{\rm Loss}Q_i .
 \tag{5.21.4}\label{eq:p5c2-gain-bound}
\]
\end{theorem}

\begin{proof}[Outline]
Check windows by \ref{B5.17a}, discards by \ref{B4.04} and served primes by \ref{B5.20a}. Use \ref{B5.18} and bad orbits for losses, \ref{B3.52} for the index, \ref{B5.15} for prime weights, \ref{B5.21a} for cutoff existence and \ref{B5.20} for both loss bounds. The complete proof is in \S\ref{prf:B5.21}.
\end{proof}

\begin{lemma}[Places and node sizes]\label{B5.21a}
For a model \(D\) in its own frame, the archimedean indices map
bijectively to \(\Phi\). At \(x\ge\max(x_1+1,3)\) in an admissible
family, assign the model weights at lift \(j+1\) for \(0\le j<m\)
to matched indices and assign weight one to unmatched positions.
Then \(w_j(\rho)\ge1\) on \(\Phi\). At thick positions
\(b\le w_j(\rho)\le b+\eta_{\rm sz}+1\), and elsewhere on \(\Phi\),
\(w_j(\rho)\le2b'\).
\end{lemma}

\begin{proof}[Outline]
Composing \(e_\tau\) with \(T\) of \ref{B2.02} maps the archimedean indices bijectively onto \(\Phi\), since the image has size \(f=2n=|\Phi|\). Parameter inequalities of \ref{B3.47} place mark and thin weight intervals in \([1,2b']\), while thick labels give the thick interval. The complete proof is in \S\ref{prf:B5.21a}.
\end{proof}

For the top labels \(\mathcal T(x)\subset\N\), an admissible
integer pair for a base model satisfies
\[
 a\le a_v(t)\le a+a/(kx^{13}),\qquad
 b'\le b_v(t)\le b_{\rm top}(x)\qquad(t\in\mathcal T(x)).
 \tag{5.22.1}\label{eq:p5c2-common-box}
\]
For \(x>x_1\), \(b_{\rm top}(x)=b'+2n\delta_1x^\varepsilon\mathcal L b'/k\).
A model realizes the pair when its top weights agree with it.

\begin{lemma}[A single weight vector for the translations]\label{B5.22}
Assume \ref{B5.21} and the following common-weight hypothesis. At each sufficiently large window point \(x\) there are a base model \(i_0\in I(x)\) and functions \(a_v,b_v:\N\to\Z\) satisfying \eqref{eq:p5c2-common-box} for \(i_0\). Some model in \(I(x)\) realizes this pair, and every model \(i\) realizing it satisfies \(R_W(i)\le x^{-9}\). Then for some \(c_{\rm Loss}\ge0,x_1\) the same pair
works at every translation and every realization satisfies
\eqref{eq:p5c2-index-bound} and \eqref{eq:p5c2-gain-bound} for window
points \(x\ge x_1\). The common-weight exponent is \(e_{\rm Loss}=0\).
\end{lemma}
\begin{proof}
Choose a pair supplied by the common-weight hypothesis. Every realization satisfies \(R_W(i)\le x^{-9}\), so \ref{B5.21} gives both loss bounds. The covering property carries the pair to every
integer translation \(0\le u_0<n\).
\end{proof}

\begin{proposition}[Small share in the integer weight box]\label{B5.22a}
There is \(\varepsilon_0>0\) such that for
\(0<\varepsilon\le\varepsilon_0\) there is a threshold \(\mathcal K\)
with the following property. For every family \(F_m\) of parameter \(\varepsilon\) and every \(CH\in\R\) for which the family passes the size tests of \ref{B8.39} at threshold \(\mathcal K\),
there is \(x_0(F_m,CH)\) such that for every \(x\ge x_0\) and every \(i\in I(x)\), one can choose integer functions \(A,B\) on
the model group, zero off \(TT\), satisfying \eqref{eq:p5c2-common-box}
on \(TT\) and
\[
 R_W(A,B)=\sum_{p\in\mathcal P_c}\frac{\log p}{p}
 \left(2\mathbf1_{\{\Sigma_p(A,B)>x^{-\varepsilon}\}}
 +4\Sigma_p(A,B)\right)\le x^{-9}.
\]
\end{proposition}

\begin{proof}[Outline]
Set the threshold by \ref{B5.20y} and average over the integer box using \ref{B3.51}, \ref{B5.20c} and \ref{B5.20j}. Injectivity and disjointness bound \(\mathbb E\Sigma_p\), and the excess-share and prime-sum bounds give a pair with \(R_W\le x^{-9}\). The complete proof is in \S\ref{prf:B5.22a}.
\end{proof}

\begin{lemma}[The weights of a model in the constructed family]\label{B5.22b}
Suppose that for all sufficiently large \(x\), every model \(i\in I(x)\) of \(F_m\) satisfies \(R_W(i)\le x^{-9}\) at its own top weights. Then \(F_m\) satisfies the common-weight hypothesis of \ref{B5.22}.
\end{lemma}
\begin{proof}[Outline]
Pick \(i_0\in I(x)\) beyond the thresholds, transport its top weights to \(\mathcal{T}(x)\), and extend by zero so \(i_0\) realizes the pair and every realization satisfies the share bound by hypothesis. For the constructed family, \ref{B8.02a} supplies \(R_W(i)\le x^{-9}\). The complete proof is in \S\ref{prf:B5.22b}.
\end{proof}


% ---- p5c3 ----
\chapter{A slope bound for families of maps}\label{ch:p5c3}

\section{Slopes and height}

Let \(K\) be a totally complex number field with ring of integers \(\mathcal O_K\). Put
\[
 n_K=[K:\Q],\qquad
 c_K=-\log 2+\frac{\log|\operatorname{disc}K|}{n_K}.
\]
A lattice in a finite dimensional \(K\)-space \(V\) is a finitely generated \(\mathcal O_K\)-submodule. Write \(r(M)=\dim_K\operatorname{span}_K M\). Choose one embedding \(K\hookrightarrow\C\) at each complex place \(v\) and give the scalar extension \(V_v\) a Hermitian norm. For a lattice \(L\), its diagonal image in \(\bigoplus_{v\mid\infty} V_v\), with squared norm \(\sum_v\|x_v\|_v^2\), has real rank \(\rho(L)=n_Kr(L)\).

For \(L\ne0\), let
\[
 \mu_{\mathrm{real}}(L)=-\frac{\log\operatorname{covol}(L)}{\rho(L)},
 \qquad \mu(L)=2\mu_{\mathrm{real}}(L)+c_K,
\]
where the covolume is taken in the real span of \(L\). The Harder--Narasimhan (HN) filtration uses induced and orthogonal quotient metrics. The maximum real slope is taken over nonzero sublattices and the minimum over nonzero quotients by saturated sublattices. Use the same affine change for the extreme slopes \(\mu_{\min}\) and \(\mu_{\max}\). For the zero lattice the real slopes are zero by convention and the normalized slopes are \(c_K\).

For a nonzero \(K\)-linear map \(\phi:V\to W\), its height is
\[
 h(\phi)=\frac{1}{\#\{v\mid\infty\}}
 \sum_{v\mid\infty}\log\|\phi_v\|_{\mathrm{op}}^2.
\]

\begin{proposition}[The slope inequality for one map]\label{B5.23}
Let \(V,W\) be finite dimensional \(K\)-spaces with the preceding metrics and \(A\subseteq V\) a lattice. Let \(\phi:V\to W\) be a \(K\)-linear map with \(\phi(A)\ne0\). Then
\[
 \mu_{\min}(A)-h(\phi)\le\mu_{\min}(\phi(A)).
\]
For every lattice \(B\subseteq W\) containing \(\phi(A)\),
\[
 \mu_{\min}(A)\le\mu_{\max}(B)+h(\phi).
\]
If \(\phi\) is injective on \(A\), then also
\[
 \mu_{\max}(A)\le\mu_{\max}(B)+h(\phi).
\]
\end{proposition}

This is the totally complex integral case of the slope inequality of Bost and Chen, in doubled slopes \cite[Propositions~1.2 and 2.10]{bost-chen-2013}. By Proposition~1.2 the maximal slope is invariant under finite extension. Duality gives the same assertion for the minimal slope, so Proposition~2.10 applies over \(K\). Restrict the source to the span of \(A\). The target is the span of \(\phi(A)\) for the surjective inequality and the span of \(B\) for the other two. The lattice inclusions bound the finite-place norms by one. Restriction weakly decreases the archimedean norms, so twice the full height is at most \(h(\phi)\).

\begin{proposition}[Counting slopes with the HN filtration]\label{B5.24}
Let \(L\) be a lattice, and let \(\lambda_1(L),\ldots,\lambda_{\rho(L)}(L)\) be its real HN slopes, with a factor of rank \(j\) over \(K\) contributing \(n_Kj\) entries. For \(\theta\in\R\), let \(L_{\ge\theta}\) be the saturated HN step containing the slopes at least \(\theta\). Then
\[
 \#\{j:\lambda_j(L)<\theta\}=\rho(L)-\rho(L_{\ge\theta}).
\]
For \(\rho(L)>0\) and \(a\in\R\),
\[
 \sum_{j=1}^{\rho(L)}(2\lambda_j(L)+a)
 =\rho(L)(2\mu_{\mathrm{real}}(L)+a).
\]
For \(L\ne0\), the least and greatest entries of \(2\lambda_j(L)+c_K\) are \(\mu_{\min}(L)\) and \(\mu_{\max}(L)\). If \(L_{\ge\theta}\ne0\), its minimum real slope is at least \(\theta\). Every nonzero sublattice \(M\subseteq L\) with minimum real slope at least \(\theta\) lies in \(L_{\ge\theta}\).
\end{proposition}

Gaudron's Proposition~5.13 gives the canonical filtration. Lemmas~5.10 and~5.11 identify its extreme slopes in the Hermitian case \cite{gaudron-2008}. These are the Hermitian formulas in the normalization by real rank and covolume.

\begin{lemma}[The share below a threshold]\label{B5.25}
For each \(i\) in a finite set \(I\), let \(x_{ij}\in\R\) for \(1\le j\le m_i\), where \(m_i\) is a nonnegative integer. Put \(R=\sum_i m_i\). Suppose \(U,\tau\ge0\), \(s\in\R\), and
\[
 \sum_{i,j}x_{ij}\ge Rs,\qquad x_{ij}\le s+U.
\]
Define
\[
 z_U(\tau)=
 \begin{cases}
 U/(U+\tau),&U>0,\\
 0,&U=0.
 \end{cases}
\]
Then
\[
 \sum_i\#\{j:x_{ij}<s-\tau\}\le Rz_U(\tau).
\]
\end{lemma}

\begin{proof}
Let \(q\) be the number of entries below \(s-\tau\). For \(U>0\),
\[
 Rs\le\sum_{i,j}x_{ij}
 \le q(s-\tau)+(R-q)(s+U)=R(s+U)-q(U+\tau).
\]
If \(U=0\), the hypotheses force every entry to equal \(s\).
\end{proof}

\section{The image of a family at one complex place}

Let \(V_i\) for \(i\) in a finite set \(I\), and \(Z\), be finite dimensional complex Hilbert spaces. Let \(\widetilde B_i:V_i\to Z\) be complex-linear and \(c_i>0\). Put
\[
 d_i=\dim_\C V_i,\qquad R=\sum_i d_i,\qquad D_Z=\dim_\C Z,\qquad
 P_i=c_i\widetilde B_i^*\widetilde B_i,\qquad
 Q_Z=\sum_i c_i\widetilde B_i\widetilde B_i^*.
\]
Write \(\|T\|_{\mathrm{HS}}^2=\operatorname{tr}(T^*T)\). The two frame estimates are
\begin{align}
 \sum_i\|P_i-1\|_{\mathrm{HS}}^2&\le eR,
 &0\le e\le\tfrac14, \label{eq:p5c3-frame-source}\\
 D_Z\operatorname{tr}(Q_Z^2)&\le
 (1+\eta)(\operatorname{tr}Q_Z)^2,
 &\eta\ge0. \label{eq:p5c3-frame-target}
\end{align}

\begin{theorem}[Codimension of the images]\label{B5.26}
Suppose \(R>0\) and \eqref{eq:p5c3-frame-source}--\eqref{eq:p5c3-frame-target} hold. Let \(V'_i\subseteq V_i\) be complex subspaces with \(\sum_i(d_i-\dim_\C V'_i)\le zR\). Then
\[
 \operatorname{codim}_Z\sum_i\widetilde B_iV'_i
 \le(\eta+4z+4\sqrt e)D_Z.
\]
\end{theorem}

\begin{proof}[Outline]
Use Cauchy–Schwarz and the source frame bound to control the trace on the orthogonal complement of the restricted images. The target frame bound gives \(q-\sqrt{\eta q}\le2(z+\sqrt{ze})\), which is estimated separately for \(z\ge1/4\) and \(z<1/4\). The complete proof is in \S\ref{prf:B5.26}.
\end{proof}

\begin{lemma}[From complex dimension to \(K\)-rank]\label{B5.27}
Let \(B_i:V_i\to W\) be \(K\)-linear maps between finite dimensional \(K\)-spaces. Put \(d_i=\dim_KV_i\) and \(D_C=\dim_KW\). Assume \(R=\sum_i d_i>0\). Suppose that at a complex place \(v\), after a common invertible change of target coordinates and an invertible change in each source, there are lifts
\[
 \widetilde B_i:\C^{d_i}\longrightarrow Z=\C^{D_C}\oplus\C^{D_0}
\]
whose projections onto the first summand represent \(B_{i,v}\). Assume \eqref{eq:p5c3-frame-source}--\eqref{eq:p5c3-frame-target} with \(D_Z=D_C+D_0\). If \(\mathcal O_K\)-submodules \(N_i^\tau\subseteq V_i\) satisfy \(\sum_i(d_i-r(N_i^\tau))\le zR\), then
\[
 D_C-r\left(\sum_i B_iN_i^\tau\right)
 \le(\eta+4z+4\sqrt e)D_Z .
\]
\end{lemma}

\begin{proof}
Let \(V'_i\) be the complex span of \(N_i^\tau\) in the source coordinates, so \(\dim_\C V'_i=r(N_i^\tau)\). The projection sends \(S=\sum_i\widetilde B_iV'_i\) onto the complex span of \(\sum_i B_iN_i^\tau\) in the target coordinates. The kernel of the projection has dimension \(D_0\), so \(D_Z-\dim_\C S\ge D_C-r(\sum_i B_iN_i^\tau)\). Apply \ref{B5.26}.
\end{proof}

\section{Integration of the slope count}

\begin{proposition}[A finite layer-cake inequality]\label{B5.28}
Let \(D\ge1\) be an integer, \(s\in\R\), \(H\ge0\), and \(x_1,\ldots,x_D\ge s-H\). Let \(w:[0,H]\to\R\). Suppose
\[
 \#\{j:x_j<s-\tau\}\le Dw(\tau)\qquad(0\le\tau\le H)
\]
and \(\min(1,w)\) is integrable on \([0,H]\). Then
\[
 s-\int_0^H\min(1,w(\tau))\,d\tau
 \le \frac1D\sum_{j=1}^D x_j.
\]
\end{proposition}

The counting-measure form of Lieb and Loss \cite[Theorem~1.13]{lieb-loss-2001} gives
\[
 \sum_j(s-x_j)_+=\int_0^H\#\{j:x_j<s-\tau\}\,d\tau.
\]
Use \(s-x_j\le(s-x_j)_+\) and bound the integrand by \(D\min(1,w(\tau))\).

For \(U,H\ge0\), define
\[
 \varphi(U,H)=
 \begin{cases}
 U\log(1+H/U),&U>0,\\
 0,&U=0,
 \end{cases}
\]
and for \(0\le\delta<1\) and \(\eta,e\ge0\) put
\begin{equation}\label{eq:p5c3-frame-error}
 E_{\mathrm{fr}}(U,H;\delta,\eta,e)
 =\frac{4\varphi(U,H)+(\eta+4\sqrt e)H}{1-\delta}.
\end{equation}

\begin{theorem}[The strong family slope inequality]\label{B5.29}\label{B5.30}
Let \(I\) be finite and nonempty. Let \(N_i\) be full lattices in finite dimensional \(K\)-spaces \(V_i\), and \(C^\vee\) a full lattice in a finite dimensional \(K\)-space \(W\). Give the spaces Hermitian norms at the complex places. Let \(B_i:V_i\to W\) be \(K\)-linear. Choose \(h_i\in\R\) for each \(i\in I\) and suppose
\[
 B_iN_i\subseteq C^\vee,\qquad B_i\ne0\ \Longrightarrow\ h(B_i)\le h_i.
\]
Set \(d_i=\dim_KV_i\) and \(R=\sum_i d_i>0\). At one complex place suppose the maps admit lifts to a common \(Z\) as in \ref{B5.27} with the frame estimates \eqref{eq:p5c3-frame-source}--\eqref{eq:p5c3-frame-target}. Write \(D_C=\dim_K W\) and \(D_Z=\dim_\C Z\), and suppose \(0\le\delta<1\) and \((1-\delta)D_Z\le D_C\). Let \(s\in\R\), \(U,H\ge0\), and assume
\[
 Rs\le\sum_i d_i(\mu(N_i)-h_i),\qquad
 \mu_{\max}(N_i)-h_i\le s+U\quad(i\in I),\qquad
 s-\mu_{\min}(C^\vee)\le H.
\]
Then
\[
 s\le\mu(C^\vee)+E_{\mathrm{fr}}(U,H;\delta,\eta,e).
\]
In particular, under the same hypotheses,
\[
 s\le\mu(C^\vee)+\log\# I+E_{\mathrm{fr}}(U,H;\delta,\eta,e).
\]
\end{theorem}

\begin{proof}[Outline]
Bound omitted source ranks by \ref{B5.24} and \ref{B5.25} and place retained images by \ref{B5.23}. Convert the target deficit from \ref{B5.27} to a slope count by \ref{B5.24}, integrate by \ref{B5.28}, and use \(\log\#I\ge0\) for the second bound. The complete proof is in \S\ref{prf:B5.30}.
\end{proof}

\section{The frame supplied by the period maps}

Fix \(C_0\ge1\) and \(C_H>0\). Take \(x>1\) with \(\sigma=x^{-C_H}\le1/20\). At translation \(u_0\), let \(\mathcal W=W(u_0)\) be the set of words. The period maps \(B_u\) of \ref{B5.10} have divided lattices \(N'_u\) and target \(RR^\vee\). Put \(d_u=r(N'_u)\) and \(R=\sum_ud_u\). Write \(M^0_{u,\C}\) for the unrestricted level-zero column space and \(f_u=\dim_\C M^0_{u,\C}\). At one complex place suppose
\[
 \widetilde B_u=A_uJ_u,\qquad A_u:\C^{f_u}\to Z,\qquad
 J_u:\C^{d_u}\to\C^{f_u},\quad J_u^*J_u=1.
\]
For \(c_u>0\), put \(K_{uv}=\sqrt{c_uc_v}\,A_u^*A_v\) and \(\mathcal R=\sum_u f_u\).

The frame input consists of \(D_*>0\), \(\eta_0,\eta_1\in\R\), \(e_0\ge0\) and \(0\le d_e\le1/4\) with
\begin{align}
 \|K_{uv}\|_{\mathrm{HS}}^2
 &\le(1+\eta_0)\frac{f_uf_v}{D_*}
 &&(u\ne v), \label{eq:p5c3-cross}\\
 \|K_{uu}-1\|_{\mathrm{HS}}^2
 &\le e_0f_u
 &&(u\in\mathcal W), \label{eq:p5c3-diagonal}\\
 D_Z&\le(1+\eta_1)D_*,
 &\sum_u(f_u-d_u)&\le d_e\mathcal R, \label{eq:p5c3-rank-loss}\\
 -1\le\eta_0&\le\sigma/2,
 &\eta_1&\le\sigma/2,\qquad \sqrt{e_0}\le\sigma/2.
 \label{eq:p5c3-small-input}
\end{align}
Fix \(e_B\ge2e_0\). In the period model put \(L=x^\varepsilon\) and \(b=e^L\), with \(0<\varepsilon<1/52\). Let \(m\) be the chain length and \(k\) the starting rank. Write \(f\) for the number of complex places of the host CM field. With \(\Gamma=\exp(C_\Gamma m^5fL)\) for \(C_\Gamma>0\), the counting inputs are
\begin{equation}\label{eq:p5c3-code-input}
 C_0^{-1}\le \frac{f_u\Gamma^2}{D_*}\le C_0
 \quad(u\in\mathcal W),\qquad
 \#\mathcal W\ge C_0\Gamma^2x^{C_H}.
\end{equation}

\begin{theorem}[Application to the divided period family]\label{B5.31}
Assume \eqref{eq:p5c3-cross}--\eqref{eq:p5c3-code-input}. Let \(N'_u\) and \(RR^\vee\) be full lattices with \(d_u>0\) for each word. Suppose
\[
 B_uN'_u\subseteq RR^\vee,\qquad
 B_u\ne0\ \Longrightarrow\ h(B_u)\le h_u,
\]
and the projections of \(\widetilde B_u\) represent \(B_u\) in invertible coordinates at the chosen place as in \ref{B5.27}. Suppose \(0\le\delta<1\) and \((1-\delta)D_Z\le r(RR^\vee)\).

In the period model take \(D_*=D^\star\) from \ref{B5.11} and \(\delta=1-D_{\rm rows}/D_{\rm all}\), where \(D_{\rm all}=D_Z\) and \(D_{\rm rows}=r(RR^\vee)\) are the full and retained row counts.

Define
\[
 s=\frac{\sum_ud_u(\mu(N'_u)-h_u)}R,\qquad
 U=\max_u\{\mu_{\max}(N'_u)-h_u\}-s .
\]
Then the frame estimates hold with
\[
 \eta=5\sigma+3\sqrt{d_e},\qquad e=2e_0\le e_B,\qquad 0\le e\le\tfrac14 .
\]
For every \(H>0\) with \(s-\mu_{\min}(RR^\vee)\le H\),
\[
 s\le\mu(RR^\vee)+\log\#\mathcal W+E_{\mathrm{fr}}(U,H;\delta,\eta,e).
\]
\end{theorem}

\begin{proof}[Outline]
Compress the source defects and combine the Gram second moment, compressed trace lower bound and dimension bound to get \(\eta=5\sigma+3\sqrt{d_e}\). Use the mean condition of \ref{B5.29} and the slope bound from \ref{B5.24} to apply \ref{B5.30}. The complete proof is in \S\ref{prf:B5.31}.
\end{proof}

\begin{lemma}[Gram estimates in the arithmetic coordinates]\label{B5.31a}
Use the selected realization of \ref{B8.02} under the source and size hypotheses of \eqref{eq:p8c1-distinct-row} in \ref{B8.10}. Use the level models of \ref{B8.37} and the row system of \ref{B8.52}. Use the Bessel identification of \ref{B5.10} for the period maps. Fix the outer parameters of \ref{B8.10}, so the source tuple in \ref{B8.02c} is
\[
 (C',C_\ell,C_{YL}+1,\tfrac14,3,\tfrac12,C_{Tu},C_0),
 \qquad C_0=3.
\]
At \(A=1\) take the signed-estimate constants \(C_B,C_4\) of \ref{B8.10d}. For the joint choice in \ref{B5.32a}, take the common middle value \(C_\Gamma=C_H=C_F\) above both \(1+\max(C_B,C_4)\) and the middle threshold of \ref{B5.10a}. The exponent \(F_0\) there depends on \(\varepsilon\) and the already fixed outer constants. Use this middle value in \(\Gamma\) and in the column-size conditions. Then choose \(C',C_\ell\) above the final thresholds of the spectral rows in \ref{B8.10}, including the same-word thresholds of \ref{B8.10y} and \ref{B8.10d}. Require \(C_\ell\ge40\) for \ref{B5.11b}. Use the absolute convergence and ties of \ref{B8.10} for this realization.

Uniformly for large \(x\), one may take
\[
 \eta_0=b^{-6},\qquad
 e_0=2C_0/\Gamma+2e^{-kf\log k},\qquad
 \eta_1=e^{-c_7rf\log r},\qquad d_e=2f/b
\]
in \eqref{eq:p5c3-cross}--\eqref{eq:p5c3-rank-loss}, with \(r=k+m-1\) and \(c_7=1/16\). Under Green's parametrization of cuspidal characters by irreducible polynomials (Definition~\ref{B7.06}) and the size and local-law conditions of \ref{B7.53}, the target defect of \ref{B5.31} satisfies
\begin{equation}\label{eq:p5c3-row-defect}
 0\le\delta\le e^{-x^{\alpha/7}},\qquad \alpha=\varepsilon/J
\end{equation}
with \(J\ge1\) the parameter of \ref{B7.52}.
\end{lemma}

\begin{proof}[Outline]
Normalize the identity cell and bound the remaining Gram terms using \ref{B8.10b}, \ref{B8.10y}, and \ref{B8.10d}. Transfer via \ref{B5.10}, bound row defects by \ref{B8.52} and \ref{B7.52}, and estimate rank loss using \ref{B5.05a}. The complete proof is in \S\ref{prf:B5.31a}.
\end{proof}

\section{Parameters common to the family}

\begin{theorem}[Uniform frame error and the height interval]\label{B5.32}
Fix an admissible family and constants \(C_H>0\), \(C_0\ge1\). Suppose that for all sufficiently large \(x\) the full lattices and period-map data of \ref{B5.31} and the counting inputs hold uniformly, with \(\sigma=x^{-C_H}\). Suppose the members admit the frame input of \ref{B5.31} with a common \(d_e=d_e(x)\) and \(2e_0\le e_B=e_B(x)\), where
\[
 0\le d_e,e_B\le\tfrac14,\qquad
 (3\sqrt{d_e}+4\sqrt{e_B})x^{C_H}\le1 .
\]
Then beyond a common threshold \ref{B5.31} applies to every member with \(0\le\eta\le\bar\eta\) and \(0\le e\le\bar e\), where
\[
 \bar\eta=5x^{-C_H}+3\sqrt{d_e},\qquad
 \bar e=e_B,\qquad (\bar\eta+4\sqrt{\bar e})x^{C_H}\le6 .
\]
In the period model with the division data of \ref{B5.10}, write \(RR_{u_0}\) for the full positive-rank lattice of translation \(u_0\). All translations use \(H=x^{C_H}b\), where \(b=e^{x^\varepsilon}\). For each \(u_0\),
\[
 s(u_0)\le\mu(RR_{u_0}^\vee)+\log\#W(u_0)
 +\frac{4\varphi(U(u_0),x^{C_H}b)+6b}{1-\delta(u_0)} .
\]
\end{theorem}

\begin{proof}[Outline]
Above the cutoffs and \(20^{1/C_H}\), \ref{B5.31} gives \(\eta=\bar{\eta}\), \(e\le\bar{e}\), and \((\bar{\eta}+4\sqrt{\bar{e}})x^{C_H}\le 6\). Arakelov slope duality and \ref{B5.10} give \(s-\mu_{\min}(RR^\vee)\le x^{C_H}b\), so \ref{B5.31} with \(H=x^{C_H}b\) yields the bound. The complete proof is in \S\ref{prf:B5.32}.
\end{proof}

\begin{lemma}[The rank in the selected level window]\label{B5.32b}
Consider a family of the working class \(W_1(\mathcal A)\) of \ref{B8.02} satisfying the closed class-number condition of \ref{B8.54} and carrying the symmetric columns of item~5 of \ref{B8.09}. Use the word levels in the window of \ref{B5.11}. Let \(f_u\) be the dimension of the lowest, level-zero column form space for the word \(u\), \(f=[F:\Q]\), and \(k\) the starting rank. Let \(D^\star>0\) be the row main term and \(d_u^\star>0\) the product of the definite mass, the pre-level index and the dimensions of the archimedean coefficient factors for \(u\). Write \(\Gamma\) for the common positive factor used to define the word windows. Beyond a single family cutoff, for every point and word,
\[
 f_u=d_u^\star(T_u^k-1)^f,\qquad
 \frac12\le\frac{f_u\Gamma^2}{D^\star}\le\exp(1).
\]
\end{lemma}
\begin{proof}[Outline]
Apply \ref{B8.54} at level zero using \ref{B2.56}, \ref{B8.51a} and \ref{B2.51}. Evaluate the symmetric column of \ref{B8.09} at class representatives and use the class count of \ref{B8.54} to get \(f_u=d_u^\star(T_u^k-1)^f\), then apply the selected prime window. The complete proof is in \S\ref{prf:B5.32b}.
\end{proof}

\begin{corollary}[The code and the shared precision]\label{B5.32a}
Under the design, selected-column and class-number hypotheses of \ref{B8.02}, for every sufficiently small \(\varepsilon>0\) and every prescribed threshold record, the construction supplies a family and an exponent \(C_H>0\) such that at each constructed point and translation \(\#W(u_0)=\lceil3\Gamma^2x^{C_H}\rceil\). Beyond a common threshold, \(\#W(u_0)\le 2^{k-3}\) for every translation. This word count and the normalized rank bound of \ref{B5.32b} give \eqref{eq:p5c3-code-input} with \(C_0=3\). For the family of \ref{B5.31a} the common upper bounds in \ref{B5.32} may be taken, for large \(x\), as
\[
 d_e(x)=E^2,\qquad e_B(x)=2E^2,\qquad E=\exp(-x^{\varepsilon/8}).
\]
\end{corollary}

\begin{proof}[Outline]
Choose the family and code by \ref{B8.05}, bound the word count by \(2^{k-3}\) using the growth of \(k\), and apply the rank estimate \ref{B5.32b}. Use \ref{B5.31a} with \(E=\exp(-x^{\varepsilon/8})\) to obtain the common bounds in \ref{B5.32}. The complete proof is in \S\ref{prf:B5.32a}.
\end{proof}


% ---- p5c4 ----
\chapter{The pointwise bound and its eight error terms}\label{ch:p5c4}

\section{Common data and their realisation}

In the numerical statements \(x>0\) and \(\varepsilon\in\R\). In a window family \(0<\varepsilon<1/52\). Let \(U_0\) be a finite set of translations with a finite nonempty word set \(W(u_0)\) at each \(u_0\). Let \(TT\) be the finite set of top positions. Write \(n_T=|TT|\). Let \(m\ge0\) and \(f>0\) be integers. Index the layers by \(0\le j<m\) and the \(f\) archimedean positions by \(\rho\). In the arithmetic model Lemma~\ref{B5.21a} identifies these positions with \(\Phi\) from Definition~\ref{B5.01}. Write \(w_{u_0j\rho}\) for their weights and \(a_t\ge0\) for the top weights, which are common to all translations. Put
\[
 S_a=\sum_{t\in TT}a_t,\qquad Q=n_Tmb,\qquad
 \langle z_\rho\rangle_\rho=\frac1f\sum_{\rho=1}^f z_\rho .
\]
At a fixed translation let \(M_u\) be the zeroth-layer column lattice and \(N'_u\) the divided lattice. Write
\[
 d_u=\operatorname{rk}N'_u,\qquad
 \mu_u=\mu(N'_u),\quad \mu_u^+=\mu_{\max}(N'_u),\qquad
 \mu_{M,u}=\mu(M_u),\quad \mu_{M,u}^+=\mu_{\max}(M_u).
\]
The target \(C^\vee=RR^\vee\) has mean slope \(\mu_C\) and minimum slope \(\mu_{C,\min}\). Let \(T_u>0\) be the level, given in the arithmetic model by the word's auxiliary prime, and \(h_u\) the height allowance. The numerical statements allow arbitrary real slopes and positive real \(d_u\). Define
\[
 \langle z_u\rangle_d=\frac{\sum_{u\in W(u_0)}d_uz_u}
 {\sum_{u\in W(u_0)}d_u},\qquad
 s=\langle\mu_u-h_u\rangle_d,\qquad
 U=\max_{u\in W(u_0)}(\mu_u^+-h_u)-s .
\]

Use the scales \(a,b,b',k,L,L_a,B_0,N_c\) of the preceding chapters, with \(k>0\). Take \(X''\in\R\) common to the translations. In the window model \(G=E_0\times A\) as in \ref{B2.02}, with \(n=|A|\). Here \(X''=a/w''\) with \(w''=b'/(xn)\). Put \(\ell_g=\log(g_0g_1)\) for the two fixed auxiliary primes. Use the net gain \(\Ga\), base term \(\Fo\), total layer loss \(\Ee_u\) and local-index total \(\mathrm{IO}\) of Definition~\ref{B5.02}, and the aggregated sum \(\ASK(u_0)\) of Definition~\ref{B5.01}. At each translation \(\Ga,\Fo,\mathrm{IO}\) are common to its words.

For real constants \(c_{\rm Div},e_{\rm Div},c_E,c_H,c_{\rm Count},c_{\rm Slope},c_1,c_{\rm Loss},e_{\rm Loss},c_{HH}\), put
\begin{align}
 \log\Omega&=\gamma+\log(\pi/v)-L_pL,\nonumber\\
 S_{\rm div}&=m\sum_{t\in TT}
 (a_t\log a_t-a_t-a_t\log\Omega),\label{eq:p5c4-Sdiv}\\
 E_{\rm div}&=c_{\rm Div}n_TaL_a+
 c_{\rm Div}m(B_0+N_c)x^{e_{\rm Div}\varepsilon},\nonumber\\
 P&=c_{\rm Loss}Qx^{e_{\rm Loss}\varepsilon},\label{eq:p5c4-P}\\
 \mathrm{Lo}_u&=c_1\mathrm{IO}
 +c_1n_Tm(a/k+b')\log T_u+c_1X''+1,\qquad
 J=\langle\mathrm{Lo}_u\rangle_d .\label{eq:p5c4-Lo}
\end{align}
Here \(v>0\) and \(\gamma,L_pL\in\R\). In the arithmetic model \(v=\mathrm{rd}(E)/2\) and \(\gamma\) is Euler's constant. The parameter \(L_pL\) then equals \(\Re\bigl(L'(1,\chi_d)/L(1,\chi_d)\bigr)\), where \(\chi_d\) is the quadratic character of discriminant \(-d\), as in Lemma~\ref{B5.07x}. Set \(a_t\log a_t=0\) at \(a_t=0\).

Uniform family estimates have common constants and thresholds for all \(i\in I_x\). The frame bounds may depend on \(x\). For \(e_L\ge0\), the common-weight loss hypothesis supplies \(c_{\rm Loss}\ge0\) and, at every sufficiently large window point, a base model \(i_0\) and functions \(a_v,b_v:\N\to\Z\) admissible for \(i_0\) as in \ref{B5.22}. Some model \(i_1\) realises both functions on the top labels. Every model realising them satisfies the losses of \ref{B5.21} with \(P\) from \eqref{eq:p5c4-P} and \(e_{\rm Loss}=e_L\).

Use the following hypotheses at each translation. The local frame parameters are \(\delta,\eta,e\), with common bounds \(\bar\delta,\bar\eta,\bar e\), and the local height parameter is \(\mathcal H\). The counting threshold \(\lambda_5\) of \ref{B5.11} is expressed in the common data in \ref{B5.34}.
\begin{enumerate}[label=\(\mathrm{H}\arabic*\)]
\item\label{hyp:p5c4-setting}
\[
 x,b,k,v>0,\quad a,b',c_1\ge0,\quad
 a_t\ge0,\quad w_{u_0j\rho}>0,\quad d_u,T_u>0,\quad
 \mu_u\le\mu_u^+,\qquad
 \bar\delta<1,\quad\bar\eta\ge0,\quad 0\le\bar e\le\tfrac14.
\]
\item\label{hyp:p5c4-div}
The division and height estimate is
\[
 \mu_C+h_u-\Fo\le S_{\rm div}-S_a\log T_u+E_{\rm div}.
\]
\item\label{hyp:p5c4-E}
The total layer loss satisfies
\[
 |\Ee_u-mS_a\ell_g|\le mc_EB_0 .
\]
\item\label{hyp:p5c4-H}
For any two words,
\[
 |(h_u+S_a\log T_u)-(h_{u'}+S_a\log T_{u'})|
 \le c_H(n_TaL_a+mB_0).
\]
\item\label{hyp:p5c4-count}
The levels satisfy
\[
 \lambda_5(u_0)\le\log T_u\le c_{\rm Count}mL,\qquad
 |\log T_u-\log T_{u'}|\le c_{\rm Count}\log k .
\]
\item\label{hyp:p5c4-M}
The zeroth-layer column slopes satisfy
\[
 -c_{\rm Slope}n_TaL_a(1+m/k)\le\mu_{M,u},\qquad
 \mu_{M,u}^+\le c_{\rm Slope}n_TaL_a .
\]
\item\label{hyp:p5c4-N}
The divided lattice satisfies
\begin{align*}
 \mu_u^+&\le\mu_{M,u}^++\Ga+\Fo-\Ee_u+c_1mB_0,\\
 \mu_{M,u}+\Ga+\Fo-\Ee_u-\mathrm{Lo}_u&\le\mu_u .
\end{align*}
\item\label{hyp:p5c4-loss}
\[
 \ASK(u_0)-\Ga(u_0)\le P,\qquad \mathrm{IO}(u_0)\le P .
\]
\item\label{hyp:p5c4-frame}
For \(u>0\) and \(H>0\) put \(\vartheta(u,H)=u\log(1+H/u)\), and set \(\vartheta(0,H)=0\). If \(0\le\delta<1\), \(\eta\ge0\), \(0\le e\le1/4\), and \(\mathcal H>0\) satisfies \(s-\mu_{C,\min}\le\mathcal H\), then
\[
 s\le\mu_C+\log|W(u_0)|
 +\frac{4\vartheta(U,\mathcal H)+(\eta+4\sqrt e)\mathcal H}
 {1-\delta}.
\]
\item\label{hyp:p5c4-HH}
\[
 s-\mu_{C,\min}\le\mathcal H,\qquad
 0<\mathcal H\le x^{c_{HH}}b .
\]
\item\label{hyp:p5c4-parameters}
\[
 0\le\delta\le\bar\delta,\qquad
 0\le\eta\le\bar\eta,\qquad 0\le e\le\bar e .
\]
\end{enumerate}

\begin{theorem}[Common realisation at a window point]\label{B5.33}
Let \(\mathcal F\) be a window family as in \ref{B5.22}. Fix \(C_{\rm abs},C_H,e_L\ge0\), \(C_5\in\R\) and \(c_6>0\). Assume the uniform estimates of \ref{B5.06} and \ref{B5.10}--\ref{B5.14}, with division exponent \(C_{\rm abs}\) and height exponent \(C_H\) as in \ref{B5.10a} for the selected design. At all sufficiently large \(x\), assume that each \(i\in I_x\) has \(d_u>0\) and \(\mu(N'_u)\le\mu_{\max}(N'_u)\) for every word, and \(\mu_{\min}(RR^\vee)=-\mu_{\max}(RR)\).

For all sufficiently large \(x\) assume common bounds \(\bar\eta=\bar\eta(x)\) and \(\bar e=\bar e(x)\) with
\begin{equation}\label{eq:p5c4-frame-budget}
 \bar\eta\ge0,\qquad 0\le\bar e\le\tfrac14,\qquad
 (\bar\eta+4\sqrt{\bar e})x^{C_H}\le6 .
\end{equation}
For each \(i\in I_x\) assume that some \(0\le\eta_i\le\bar\eta\) and \(0\le e_i\le\bar e\) give \ref{hyp:p5c4-frame} for every \(\mathcal H=H>0\) with \(s_i-\mu_{C,\min,i}\le H\). Use the model's defect \(\delta=\delta_i=1-D_{{\rm rows},i}/D_{{\rm all},i}\), as in Theorem~\ref{B5.31}. Assume the defect bound of \ref{B7.52} for these defects and the common-weight loss hypothesis with exponent \(e_L\).

Under the small-\(\varepsilon\) design hypotheses of Corollary~\ref{B5.32a}, use its \(d_e,e_B\) in Theorem~\ref{B5.32} to take \(\bar\eta=5x^{-C_H}+3\sqrt{d_e}\) and \(\bar e=e_B\). For the same target, \ref{B8.52}--\ref{B8.53} supply the good-label injection and relative cylinder counts on the compatible row system of \ref{B7.52}. The local idempotents use one common nontrivial additive character at the places above each distinguished prime. With the weight laws of \ref{B7.50} and \(C_s=6\), \ref{B7.53} gives the defect bound.

Assume that for all sufficiently large \(x\) and every \(i\in I_x\), its host satisfies the thin-character and three sign-character hypotheses of Lemma~\ref{B5.33a} with the fixed constants \(c_6,C_5\), and its top set satisfies
\begin{equation}\label{eq:p5c4-top-density}
 \tfrac12 nx^{-\varepsilon/8}\le n_T\le2nx^{-\varepsilon/8}.
\end{equation}
For the fixed \(c_6\), assume an exponent \(e_K\ge0\) and a threshold common to the models above which the size bounds of \ref{B5.45} on \(m,k,L_a,b',X'',|W(u_0)|\) and the chain-length lower bound of \ref{B5.50} hold. With \(X'=\exp(x^{\varepsilon/3})\), \(\theta=\exp(-c_6x^{\varepsilon/8})\), \(n_2\) the number of thick labels and \(\eta_{\rm sz}\ge0\) the common width, assume also
\begin{equation}\label{eq:p5c4-size-input}
 \begin{gathered}
 2f(x+1)\theta\le1/x,\qquad 6f(x+1)(\eta_{\rm sz}+1)\le b/x,\\
 n_2/n\ge x^{-\varepsilon},\qquad k\le n^2x^{1/2-\varepsilon}.
 \end{gathered}
\end{equation}

Then there are real constants \(K,c_{HH},e_K,C_L,x_0\) such that, at each window point \(x\ge x_0\), there are common integer weight functions \(a_v,b_v\) and a model realising them at each translation. One can take
\[
 c_{HH}=C_H,\qquad C_L=C_5,\qquad
 e_{\rm Div}=C_{\rm abs},\qquad e_{\rm Loss}=e_L .
\]
The selected models share the frame and global scales, including \(m,f,q,d\), and satisfy \ref{hyp:p5c4-setting}--\ref{hyp:p5c4-parameters}. With \(b=e^{x^\varepsilon}\), \(a=xb\) and \(N=\lceil2a\rceil\), their weights satisfy
\[
 a\le a_t\le2a,\qquad |a_t/b-x|\le1,\qquad
 a_t/w_{u_0j\rho}\le N\quad(t\in TT,\ 0\le j<m,\ \rho\in\Phi).
\]
They give the arithmetic data of \ref{B5.46}--\ref{B5.49}, with \(\ell\) from \ref{B5.33a} and \(b'X'=xb\). They satisfy the windows of \ref{B5.45} and the chain-length bounds of \ref{B5.50}. They also satisfy \eqref{eq:p5c4-top-density} and \eqref{eq:p5c4-size-input}. In particular \(1\le k\le x^{e_K}\).

If \(S_2\) is the common set of integer thick labels and \(a_0\) generates \(A_0\), put \(T_2=\{sa_0:s\in S_2\}\). The host's label-range parameter \(L_w\ge2\) gives \(0\le s\le4L_ws<\operatorname{ord}(a_0)\) for \(s\in S_2\), so \(n_2=|T_2|=|S_2|\). At the thick positions \(((0,1),s_2+u_0)\) with \(s_2\in T_2\), the weights lie in \([b,b+\eta_{\rm sz}+1]\). They lie in \((0,2b']\) at the other positions of \(\Phi\). For the family of \ref{B5.22} one may take \(e_L=0\).
\end{theorem}

\begin{proof}[Outline]
\ref{B5.22} aligns weights and gives \(e_L=0\). Use \ref{B5.10}, \ref{B5.13}, \ref{B5.14}, \ref{B5.11}, \ref{B5.12}, \ref{B5.06} for words, \ref{B7.52} and duality for frames, budgets for windows, \ref{B5.21a} for weights, \ref{B2.03} locally and \ref{B5.33a} for constants. The complete proof is in \S\ref{prf:B5.33}.
\end{proof}

\begin{lemma}[Character constants for a host]\label{B5.33a}
Let \(F_m\) be an admissible family of exponent \(\varepsilon\), and fix \(c_6,C_5\in\R\). For every \(x\ge0\) and model at \(x\), take the arithmetic datum of its host in the same frame, with the \(q\)-coordinate condition of \ref{B2.46}. There is a function \(\ell:\widehat G\to\C\) with \(L_pL=\Re\ell(\chi_C)\), where \(\chi_C\) is the \(d\)-sign. If \(\xi\ne\mathbf1\) and a primitive Dirichlet character \(\chi\) of positive modulus satisfies \(\chi(n)=\prod_{p\mid n}V_\xi(p^{v_p(n)})\) for every \(n\ge1\), with the local averages \(V_\xi\) of Definition~\ref{B2.04}, then \(\ell(\xi)=(L'/L)(1,\chi)\).

Use this \(\ell\) in \(G_\xi\) from \ref{B2.04}. Set \(X'=\exp(x^{\varepsilon/3})\) and \(\theta=\exp(-c_6x^{\varepsilon/8})\). If the primitive characters of the host satisfy \(|g_\chi(y)|\le y\theta\) for \(y\ge X'/3\) as in \ref{B8.47}, then \(|G_\xi(y)|\le y\theta\) on that range for every \(\xi\). Suppose the primitive representatives of \(({-q}/n)\), \(({-d}/n)\) and their product on units modulo \(8qd\) satisfy
\begin{equation}\label{eq:p5c4-qd-input}
 |g_\chi(x)|\le x^{1/2+\varepsilon},\qquad
 \left|\frac{L'(1,\chi)}{L(1,\chi)}\right|\le C_5\log x .
\end{equation}
Then \(|G_\xi(x)|\le x^{1/2+\varepsilon}\) and \(|\ell(\xi)|\le C_5\log x\) for the two sign characters and their product. The hosts constructed in Theorem~\ref{B2.59} satisfy \eqref{eq:p5c4-qd-input} with \(C_5=4\) at sufficiently large window points.
\end{lemma}

\begin{proof}[Outline]
Identify primitive representatives using Frobenius and quadratic reciprocity, then transfer bounds via \ref{B2.15}. Select primes in stages \ref{B2.38} and \ref{B2.39}, using \ref{B2.36} for window bounds and \ref{B2.35} to bound logarithmic derivatives. The complete proof is in \S\ref{prf:B5.33a}.
\end{proof}

\section{The main term}

\begin{definition}[Counting threshold and geometric main term]\label{B5.34}
With the convention of \eqref{eq:p5c4-Sdiv}, put
\begin{align}
 r&=c_{\rm Count}(L_a+m^2L_a/k),\label{eq:p5c4-r}\\
 \lambda_5(u_0)
 &=\sum_{j=0}^{m-1}
 \left\langle\log(w_{u_0j\rho}v/\pi)\right\rangle_\rho
 +m\ell_g-r,\label{eq:p5c4-lambda}\\
 \Main(u_0)
 &=\sum_{j=0}^{m-1}\sum_{t\in TT}
 a_t\left(\log a_t-1-\gamma+L_pL
 -\left\langle\log w_{u_0j\rho}\right\rangle_\rho\right).
 \label{eq:p5c4-main}
\end{align}
\end{definition}

\begin{lemma}[Exact cancellation of the geometric main terms]\label{B5.35}
For \(v>0\), \(f>0\) and \(w_{u_0j\rho}>0\),
\[
 S_{\rm div}-S_a\lambda_5(u_0)+mS_a\ell_g
 =\Main(u_0)+S_ar .
\]
\end{lemma}

\begin{proof}
Put \(Z=\sum_j\langle\log w_{u_0j\rho}\rangle_\rho\) and \(B=\sum_t a_t(\log a_t-1-\gamma+L_pL)\). The definitions give
\[
 S_{\rm div}=mB-mS_a\log(\pi/v),\quad
 \lambda_5(u_0)=Z-m\log(\pi/v)+m\ell_g-r,\quad
 \Main(u_0)=mB-S_aZ .
\]
Substitution gives \(S_{\rm div}-S_a\lambda_5(u_0)+mS_a\ell_g=mB-S_aZ+S_ar\).
\end{proof}

With the window-model parameter \(\log\Gamma\) equal to \(C_\Gamma m^5fL\) as in \eqref{eq:p5c3-code-input}, the earlier remainder and threshold are
\[
 r_{\rm old}=c_{\rm Count}(L+m^2L/k)+\frac{2\log\Gamma}{fk},
 \qquad \lambda_{\rm old}(u_0)=\lambda_5(u_0)+r-r_{\rm old}.
\]
Substitution in \ref{B5.35} gives
\[
 S_{\rm div}-S_a\lambda_{\rm old}(u_0)+mS_a\ell_g
 =\Main(u_0)+S_ar_{\rm old}.
\]

\section{The lower bound and the slope spread}

\begin{lemma}[Lower bound for the height-adjusted mean slope]\label{B5.36}
Assume \ref{hyp:p5c4-setting}--\ref{hyp:p5c4-E} and \ref{hyp:p5c4-count}--\ref{hyp:p5c4-N}. Put
\[
 K_5(u_0)=c_{\rm Slope}n_TaL_a(1+m/k)
 +m(S_a\ell_g+c_EB_0)
 +S_{\rm div}-S_a\lambda_5(u_0)+E_{\rm div}.
\]
Then \(\Ga+\mu_C-K_5(u_0)-J\le s\) at every translation.
\end{lemma}

\begin{proof}[Outline]
Start from \ref{hyp:p5c4-N}, bound the height-adjusted forcing term below using \ref{hyp:p5c4-div} and \ref{hyp:p5c4-count}, and apply the error and slope bounds of \ref{hyp:p5c4-E} and \ref{hyp:p5c4-M}. Taking the \(d\)-weighted mean gives the lower bound for \(s\). The complete proof is in \S\ref{prf:B5.36}.
\end{proof}

\begin{lemma}[Variation between two words]\label{B5.37}
Assume \ref{hyp:p5c4-setting} and \ref{hyp:p5c4-E}--\ref{hyp:p5c4-count}. Put
\[
 V=2mc_EB_0+c_H(n_TaL_a+mB_0)
 +S_ac_{\rm Count}\log k .
\]
Then \((-\Ee_u-h_u)-(-\Ee_{u'}-h_{u'})\le V\) for all \(u,u'\in W(u_0)\).
\end{lemma}

\begin{proof}
By \ref{hyp:p5c4-E}, \(\Ee_{u'}-\Ee_u\le2mc_EB_0\). With \(H_u=h_u+S_a\log T_u\),
\[
 h_{u'}-h_u
 =H_{u'}-H_u+S_a(\log T_u-\log T_{u'}).
\]
The first term is at most \(c_H(n_TaL_a+mB_0)\) by \ref{hyp:p5c4-H}. Since \(S_a\ge0\), \ref{hyp:p5c4-count} bounds the second term by \(S_ac_{\rm Count}\log k\).
\end{proof}

\begin{lemma}[Maximum slope above the weighted mean]\label{B5.38}
Assume \ref{hyp:p5c4-setting} and \ref{hyp:p5c4-E}--\ref{hyp:p5c4-N}. Put
\[
 K_3=c_{\rm Slope}n_TaL_a(2+m/k)+V+c_1mB_0 .
\]
Then \(0\le U\le K_3+J\).
\end{lemma}

\begin{proof}[Outline]
Subtracting \ref{hyp:p5c4-N} at \(u\) and \(u'\) cancels common terms, and \ref{hyp:p5c4-M} with \ref{B5.37} bounds the difference by \(K_3+\mathrm{Lo}_{u'}\). Averaging over \(u'\) and maximizing over \(u\) gives \(U\le K_3+J\), while \ref{hyp:p5c4-setting} gives \(0\le U\). The complete proof is in \S\ref{prf:B5.38}.
\end{proof}

\section{The logarithmic correction and the upper bound}

\begin{lemma}[Uniform logarithmic correction]\label{B5.39}
Let \(u,H,X,b,U_*\in\R\) satisfy \(0\le u\le U_*\), \(0<H\le Xb\) and \(X,b>0\). With \(\vartheta\) as in \ref{hyp:p5c4-frame},
\[
 \vartheta(u,H)\le (U_*+b)\log(1+X).
\]
\end{lemma}

\begin{proof}
Put \(L_X=\log(1+X)>0\). The bound holds at \(u=0\). For \(u>0\), use \(\log z\le z-1\) at \(z=(1+H/u)/(1+X)\) and multiply by \(u\). Then
\[
 \vartheta(u,H)
 \le uL_X+\frac{H}{1+X}-u\left(1-\frac1{1+X}\right)
 \le U_*L_X+\frac{Xb}{1+X}.
\]
The same inequality at \(z=(1+X)^{-1}\) gives \(X/(1+X)\le L_X\).
\end{proof}

\begin{theorem}[Upper bound for the net gain]\label{B5.40}
Assume \ref{hyp:p5c4-setting}--\ref{hyp:p5c4-N} and \ref{hyp:p5c4-frame}--\ref{hyp:p5c4-parameters}. Write
\[
 X=x^{c_{HH}}>0,\qquad
 K_0=S_ar+E_{\rm div}+mc_EB_0+c_{\rm Slope}n_TaL_a(1+m/k).
\]
At every translation,
\begin{align*}
 \Ga(u_0)\le{}&
 \Main(u_0)+\log|W(u_0)|+K_0+J\\
 &+\frac{4\log(1+X)(K_3+J+b)+(\bar\eta+4\sqrt{\bar e})Xb}
 {1-\bar\delta}.
\end{align*}
\end{theorem}

\begin{proof}[Outline]
Combine \ref{B5.36} with \ref{hyp:p5c4-frame}, cancel \(\mu_C\), and rewrite \(K_5(u_0)\) by \ref{B5.35}. Bound the numerator using \ref{B5.38}, \ref{hyp:p5c4-HH}, \ref{B5.39}, and \ref{hyp:p5c4-parameters}, then replace \(1/(1-\delta)\) by \(1/(1-\bar{\delta})\). The complete proof is in \S\ref{prf:B5.40}.
\end{proof}

\begin{lemma}[Uniform bound for the mean quotient loss]\label{B5.41}
Assume \ref{hyp:p5c4-setting} and \ref{hyp:p5c4-count}. Then
\[
 J\le c_1\mathrm{IO}
 +c_1n_Tm(a/k+b')c_{\rm Count}mL+c_1X''+1 .
\]
If also \ref{hyp:p5c4-loss} holds, then \(J\le M_{\rm Lo}\), where
\[
 M_{\rm Lo}=c_1P+c_1n_Tm(a/k+b')c_{\rm Count}mL+c_1X''+1 .
\]
\end{lemma}

\begin{proof}
By \ref{hyp:p5c4-setting}, the coefficient of \(\log T_u\) in \eqref{eq:p5c4-Lo} is nonnegative. Use \(\log T_u\le c_{\rm Count}mL\) and take the weighted mean. Since \(c_1\ge0\), \(\mathrm{IO}\le P\) gives the second bound.
\end{proof}

\section{The eight error terms}

\begin{definition}[Error at one translation]\label{B5.42}
Put \(\bar U=K_3+M_{\rm Lo}\) and \(X=x^{c_{HH}}\). The error at \(u_0\) is
\begin{align}
 E_8(u_0)&=
 \underbrace{S_ac_{\rm Count}(L_a+m^2L_a/k)}_{\text{\rm 1}}
 \nonumber\\
 &+\underbrace{c_{\rm Div}n_TaL_a+
 c_{\rm Div}m(B_0+N_c)x^{e_{\rm Div}\varepsilon}}
 _{\text{\rm 2}}
 \nonumber\\
 &+\underbrace{mc_EB_0}_{\text{\rm 3}}
 +\underbrace{c_{\rm Slope}n_TaL_a(1+m/k)}_{\text{\rm 4}}
 \nonumber\\
 &+\underbrace{\log|W(u_0)|}_{\text{\rm 5}}
 +\underbrace{c_{\rm Loss}n_Tmbx^{e_{\rm Loss}\varepsilon}}
 _{\text{\rm 6}}
 \nonumber\\
 &+\underbrace{M_{\rm Lo}}_{\text{\rm 7}}
 +\underbrace{\frac{4\log(1+X)(\bar U+b)
 +(\bar\eta+4\sqrt{\bar e})Xb}
 {1-\bar\delta}}_{\text{\rm 8}} .
 \label{eq:p5c4-eight}
\end{align}
In terms \(7\) and \(8\) the expanded quantities are
\begin{align*}
 M_{\rm Lo}={}&c_1c_{\rm Loss}n_Tmbx^{e_{\rm Loss}\varepsilon}
 +c_1n_Tm(a/k+b')c_{\rm Count}mL+c_1X''+1,\\
 \bar U={}&c_{\rm Slope}n_TaL_a(2+m/k)
 +2mc_EB_0+c_H(n_TaL_a+mB_0)\\
 &+S_ac_{\rm Count}\log k+c_1mB_0+M_{\rm Lo}.
\end{align*}
The family identifies each word set with \(W(x)\), so term \(5\) is constant across the selected translations.
\end{definition}

The kernel in Theorem~\ref{B5.29} has integral \(\int_0^{\mathcal H}U/(U+\tau)\,d\tau=\vartheta(U,\mathcal H)\), with value zero at \(U=0\). Since \(W(u_0)\) is nonempty, \(\log|W(u_0)|\ge0\). Adding this quantity to the first bound of Theorem~\ref{B5.29} gives its second bound and term \(5\) above. Use \(M_{\rm Lo}\) to bound \(J\) in both \ref{B5.36} and \ref{B5.38}.

\begin{theorem}[The pointwise bound for \(\ASK\)]\label{B5.43}
Assume \ref{hyp:p5c4-setting}--\ref{hyp:p5c4-parameters}. Then for every \(u_0\in U_0\),
\[
 \boxed{\ \ASK(u_0)\le\Main(u_0)+E_8(u_0)\ } .
\]
At every sufficiently large window point this holds for the models selected in Theorem~\ref{B5.33}.
\end{theorem}

\begin{proof}
The coefficient of \(J\) in \ref{B5.40} is \(1+4\log(1+X)/(1-\bar\delta)>0\). Substitute \(J\le M_{\rm Lo}\) from \ref{B5.41} and apply \ref{hyp:p5c4-loss}. Then
\begin{align*}
 \ASK-\Main
 &\le P+\Ga-\Main\\
 &\le P+\log|W(u_0)|+K_0+M_{\rm Lo}\\
 &\quad+\frac{4\log(1+X)(K_3+M_{\rm Lo}+b)
 +(\bar\eta+4\sqrt{\bar e})Xb}{1-\bar\delta}.
\end{align*}
Expand \(K_0\) and substitute \(\bar U=K_3+M_{\rm Lo}\) to obtain \eqref{eq:p5c4-eight}.
\end{proof}

Under the weight and constant identifications of Lemma~\ref{B5.46}, \(\Main(u_0)\) is the sum of the real parts of the arithmetic main terms over top positions and layers. Under the window hypotheses of Theorem~\ref{B5.45}, the error satisfies
\[
 E_8(u_0)\le x^{(e_{\rm Div}+e_{\rm Loss}+3)\varepsilon}
 \bigl[n_Ta(1+m^2/k)+m(B_0+N_c)\bigr].
\]


% ---- p5c5 ----
\chapter{Four terms and the square-root bound}\label{ch:p5c5}

For the chains of Alp\"oge's construction \cite{alpoge-2026},
Theorem~\ref{B5.43} bounds \(\ASK(u_0)-\Main(u_0)\) by the error
\(E_8(u_0)\) of Definition~\ref{B5.42}.

Write \(n_T\) for the number of top positions, \(m\) for the chain
length, \(k\) for the initial rank, and \(a\) for the common scale of the
top weights. Let \(B_0,N_c\) be the parameters in the second term of
Definition~\ref{B5.42}, and let \(e_{\rm Div},e_{\rm Loss}\) be the
division and loss exponents. Set
\begin{equation}\label{eq:four-bracket}
 \mathcal E=n_Ta\left(1+\frac{m^2}{k}\right)+m(B_0+N_c),
 \qquad C_{\rm err}=e_{\rm Div}+e_{\rm Loss}+3.
\end{equation}

\section{Combining the eight terms}

\begin{theorem}[A real-variable error bound]\label{B5.44}
Let all the parameters below be real. Suppose
\[
\begin{gathered}
 y,y_D,y_L\ge1,\qquad n_T,a\ge1,\qquad b,m\ge0,\qquad k>0,\\
 b\le a,\quad m\le k,\quad mb\le a,\quad
 0\le L,L_a\le2y,\\
 b'\ge0,\quad kb'\le a,\quad 0\le X''\le a,\quad
 0\le S_a\le2n_Ta,\quad B_0,N_c\ge0,\\
 \ell_W\le a,\quad 0\le\ell_k\le gy,\quad
 0\le\ell_H\le hy,\quad 0\le d_I\le2,\quad
 \tau\le Ky\,n_Ta.
\end{gathered}
\]
Assume that each of
\[
 c_{\rm Div},c_E,c_H,c_{\rm Count},c_{\rm Slope},c_1,c_{\rm Loss}
\]
lies in \([0,K]\), that \(h,g\ge1\), and that
\[
 200(1+K)^2hg\le y.
\]
Define
\[
\begin{aligned}
 P_0&=c_{\rm Loss}n_Tmb\,y_L,\\
 J_0&=c_1P_0+
 c_1n_Tm(a/k+b')\,c_{\rm Count}mL+c_1X''+1,\\
 V_0&=2mc_EB_0+c_H(n_TaL_a+mB_0)+S_ac_{\rm Count}\ell_k,\\
 \bar U&=c_{\rm Slope}n_TaL_a(2+m/k)+V_0+c_1mB_0+J_0.
\end{aligned}
\]
Let
\[
\begin{aligned}
 R_1&=S_ac_{\rm Count}(L_a+m^2L_a/k),\\
 R_2&=c_{\rm Div}n_TaL_a+c_{\rm Div}m(B_0+N_c)y_D,\\
 R_3&=mc_EB_0,\qquad
 R_4=c_{\rm Slope}n_TaL_a(1+m/k),\\
 R_5&=\ell_W,\qquad R_6=P_0,\qquad R_7=J_0,\\
 R_8&=d_I\{4\ell_H(\bar U+b)+\tau\}.
\end{aligned}
\]
Then
\[
 \sum_{i=1}^8R_i\le y_Dy_Ly^3
 \left\{n_Ta(1+m^2/k)+m(B_0+N_c)\right\}.
\]
\end{theorem}

\begin{proof}[Outline]
Set \(P=y_Dy_LyE\), bound the first six terms by multiples of \(P\), and use \(b'\le a/k\) to bound \(J_0\) in the estimates for \(\bar U+b\) and \(R_8\). The sum is at most \(200(1+K)^2hgyP\le y^2P\) by the hypothesis on \(y\). The complete proof is in \S\ref{prf:B5.44}.
\end{proof}

\begin{theorem}[The error bound in a chain window]\label{B5.45}
For the geometric data of Definition~\ref{B5.42}, take finite sets
\(U_0,\mathrm{Top}\). For each \(u_0\in U_0\), let \(W(u_0)\) be
finite and nonempty. Let \(m,f\in\N\), \(n_T=\#\mathrm{Top}\), and
\(S_a=\sum_{t\in\mathrm{Top}}a_t\).
Assume that the window data satisfy
\[
\begin{gathered}
 x\ge1,\quad 0<\varepsilon\le1,\quad b\ge1,\quad xb\le a,\quad n_T\ge1,\\
 m\le x^{1/2+\varepsilon/2},\quad xm\le k,\quad
 e_K\ge0,\quad 1\le k\le x^{e_K},\quad k\le a,\\
 0\le L,L_a\le2x^\varepsilon,\quad b'\ge0,\quad kb'\le a,\quad
 0\le X''\le a,\\
 0\le a_t\le2a\quad(t\in\mathrm{Top}),\qquad B_0,N_c\ge0,\\
 \log\#W(u_0)\le(k-3)\log2\quad(u_0\in U_0),\\
 \bar\delta\le\tfrac12,\qquad \bar\eta,\bar e\ge0.
\end{gathered}
\]
Suppose the seven coefficients of Theorem~\ref{B5.44} lie in \([0,K]\),
\(e_{\rm Div},e_{\rm Loss},c_{HH}\ge0\), and
\[
 (\bar\eta+4\sqrt{\bar e})x^{c_{HH}}b
 \le Kx^\varepsilon n_Ta.
\]
Assume also
\[
 \log x\le x^\varepsilon,\qquad
 200(1+K)^2(1+c_{HH})(1+e_K)\le x^\varepsilon.
\]
Then, for every \(u_0\in U_0\),
\[
 E_8(u_0)\le x^{C_{\rm err}\varepsilon}\mathcal E.
\]
\end{theorem}

\begin{proof}[Outline]
The window hypotheses allow Theorem~\ref{B5.44} with \(y=x^\varepsilon\), \(y_D=x^{e_{\rm Div}\varepsilon}\), and \(y_L=x^{e_{\rm Loss}\varepsilon}\). By Definition~\ref{B5.42}, the sum is \(E_8(u_0)\), and \(y_Dy_Ly^3=x^{C_{\rm err}\varepsilon}\). The complete proof is in \S\ref{prf:B5.45}.
\end{proof}

\section{The arithmetic difference and its average}

Let \(H\) be a finite abelian group, written additively, and put
\[
 n=\#H,\qquad \Gamma=(\Z/2\Z)^2\times H,\qquad
 \Phi=\{((0,e),h):e\in\Z/2\Z,\ h\in H\}.
\]
Thus \(f=\#\Phi=2n\). Identify \(H\) with \(\{(0,0)\}\times H\)
inside \(\Gamma\). Write
\[
 r_t(g)=((1,1),0)+g+((0,0),t),\quad
 e_1(t)=((1,0),t),\quad
 \chi_C((e_0,e_1),h)=(-1)^{e_0}.
\]
Characters of \(\Gamma\) take values in \(\C^\times\), and
\(\mathbf1\) denotes the trivial character.
Use the arithmetic datum \(D\) and centred functions \(G_\psi\)
of Definition~\ref{B2.04}, with \(\ell:\widehat\Gamma\to\C\) and
\(G_{\mathbf1}=G_z\). Fix \(m,N\in\N\), top weights \(a_t\in\R\)
for \(t\in H\), and local weights \(w_{s,j}:\Gamma\to\R\) for
\(s\in H\) and \(0\le j<m\).
Write \(\ASK(D,N,A,t,w)\) for \(\ASK_N(A,t,w)\) of
Definition~\ref{B2.04}. Let \(\gamma\) be Euler's constant and use
the real logarithm \(\log r=\log|r|\) for \(r\ne0\), with \(\log0=0\).
For \(A\in\R\) and \(w:\Gamma\to\R\), the one-layer main term is
\begin{equation}\label{eq:arith-main}
 \mathrm{MAIN}(\ell,A,w)
 =A\left(\log A-1-\gamma+\ell(\chi_C)
 -\frac1f\sum_{\rho\in\Phi}\log w(\rho)\right).
\end{equation}

Take the geometric data of Definition~\ref{B5.34} on \(U_0\),
with top set \(\mathrm{Top}\) and \(f\) places per layer.
Fix \(\mathrm{tr}:U_0\to H\) and \(TT\subseteq H\). Choose bijections
\[
 e_T:\mathrm{Top}\longrightarrow TT,\qquad
 e_P:\{0,\ldots,f-1\}\longrightarrow\Phi.
\]
In the models of Theorem~\ref{B5.33}, the place bijection of
Lemma~\ref{B5.21a} identifies the sum of log weights over \(\Phi\)
with that over the archimedean places at lift \(j+1\) for each
translation. In the construction of Theorem~\ref{B5.33}, choose
\(\ell\) at the base model \(i_0\) by Lemma~\ref{B5.33a} and use this
choice for every translation.

Write \(a^{\rm geom}_\tau,w^{\rm geom}_{u_0j}(\rho),
\gamma_{\rm geom},L_pL\) for the geometric top and local weights and the constants in Definition~\ref{B5.34}. Assume
\[
\begin{aligned}
 a^{\rm geom}_\tau&=a_{e_T(\tau)},&
 w^{\rm geom}_{u_0j}(\rho)&=w_{\mathrm{tr}(u_0),j}(e_P\rho),\\
 \gamma_{\rm geom}&=\gamma,& L_pL&=\Re\ell(\chi_C)
\end{aligned}
\]
for \(\tau\in\mathrm{Top}\), \(u_0\in U_0\), \(0\le j<m\), and
\(0\le\rho<f\). Write \(\ASK_{\rm geom}\) for the geometric prime sum
of Theorem~\ref{B5.43}, and assume
\[
 \ASK_{\rm geom}(u_0)=
 \sum_{t\in TT}\sum_{j=0}^{m-1}
 \ASK(D,N,a_t,t,w_{\mathrm{tr}(u_0),j})
 \qquad(u_0\in U_0).
\]

\begin{lemma}[Identification of the main term]\label{B5.46}
Under these identifications, every \(u_0\in U_0\) satisfies
\[
 \Main(u_0)=
 \sum_{t\in TT}\sum_{j=0}^{m-1}
 \Re\mathrm{MAIN}(\ell,a_t,w_{\mathrm{tr}(u_0),j}).
\]
\end{lemma}

\begin{proof}
By Definition~\ref{B5.34}, the geometric main term is
\[
 \sum_{j=0}^{m-1}\sum_{\tau\in\mathrm{Top}}a^{\rm geom}_\tau
 \left(\log a^{\rm geom}_\tau-1-\gamma_{\rm geom}+L_pL
 -\frac1f\sum_{\rho=0}^{f-1}
 \log w^{\rm geom}_{u_0j}(\rho)\right).
\]
Reindex the logarithm sum by \(e_P\) and the top sum by \(e_T\).
The constant identities make each summand the real part of
\eqref{eq:arith-main}. Exchange the two outer sums.
\end{proof}

Let \(T_2\subseteq H\), and put \(n_T=|TT|\) and \(n_2=|T_2|\).
For \(x>0\), \(\varepsilon>0\), common weight \(b>0\), initial
rank \(k>0\), and top-weight scale \(a_E\in\R\), define
\[
\begin{aligned}
 B_*(x)&=2nmbx^\varepsilon,&
 N_*(x)&=\frac{n_2kbx^\varepsilon}{n(nx^{-\varepsilon})},\\
 \mathcal B_x(a_E)&=
 n_Ta_E(1+m^2/k)+m\{B_*(x)+N_*(x)\}.
\end{aligned}
\]
For \(s\in H\), set
\[
 \mathrm{gap}(s)=
 \sum_{t\in TT}\sum_{j=0}^{m-1}
 \{\ASK(D,N,a_t,t,w_{s,j})
 -\Re\mathrm{MAIN}(\ell,a_t,w_{s,j})\}.
\]

\begin{theorem}[The geometric difference is the arithmetic difference]
\label{B5.47}
Assume the identifications stated before Lemma~\ref{B5.46}.
Suppose the eleven hypotheses of Theorem~\ref{B5.43}, the window
hypotheses of Theorem~\ref{B5.45}, and
\[
 B_0\le B_*(x),\qquad N_c\le N_*(x)
\]
hold for the same geometric data. Then every \(u_0\in U_0\) satisfies
\[
 \mathrm{gap}(\mathrm{tr}(u_0))
 =\ASK_{\rm geom}(u_0)-\Main(u_0)
 \le x^{C_{\rm err}\varepsilon}\mathcal B_x(a).
\]
\end{theorem}

\begin{proof}
The assumed identity for \(\ASK_{\rm geom}\) and
Lemma~\ref{B5.46} give the equality. By
Theorems~\ref{B5.43} and~\ref{B5.45}, the geometric difference is at most
\(E_8(u_0)\le x^{C_{\rm err}\varepsilon}\mathcal E\).
By \(e_T\), the top cardinalities agree. Substituting the bounds on
\(B_0,N_c\) in \eqref{eq:four-bracket} gives
\(\mathcal E\le\mathcal B_x(a)\).
\end{proof}

\begin{theorem}[The one-layer Fourier identity]\label{B5.48}
Let \(D\) have Frobenius elements \(F_p\in\Gamma\), inertia subgroups
\(I_p\le\Gamma\), and distinguished integers \(q_1,d\).
Take a finite set \(S\subseteq\N\) with \(2,q_1\in S\), and assume
\(d\) is prime and \(d\notin S\). Impose the local conditions of
Lemma~\ref{B2.03}
\[
\begin{gathered}
 I_p=\{0\}\quad\text{for primes }p\notin S\cup\{d\},\\
 I_p\subseteq\Phi,\quad F_p\in\Phi\quad(p\in S),\qquad
 I_d=\{0,((1,0),0)\}.
\end{gathered}
\]
Let \(\ell:\widehat\Gamma\to\C\), \(A>0\), \(t\in H\),
\(N\in\N\), and \(w:\Gamma\to\R\).
If
\[
 w(\rho)>0,\qquad A/w(\rho)\le N\qquad(\rho\in\Phi),
\]
then
\[
 \ASK(D,N,A,t,w)-\mathrm{MAIN}(\ell,A,w)
 =\frac1f\sum_{g\notin\Phi}\sum_{\psi\in\widehat\Gamma}
 \overline{\psi(g)}\,w(r_t(g))
 G_\psi\!\left(\frac A{w(r_t(g))}\right)
\]
as an equality in \(\C\).
\end{theorem}

\begin{proof}
Theorem~\ref{B2.07} with \(a=A\) identifies the difference with
\(\mathrm{De}(A,t,w)\). Expand this term by Definition~\ref{B2.04}.
\end{proof}

For \(A_0\le H\), write
\[
 \Psi_A=\{\psi\in\widehat\Gamma:
 \psi((0,0),s)=1\ \text{for every }s\in A_0\}.
\]
A position \(\rho\) is thick at \(u_0\) if
\(\rho=((0,1),s+u_0)\) for some \(s\in T_2\).
Put \(Q=n_Tmb\) and \(\delta_2=n_2/n\).

\begin{lemma}[Averaging over a subgroup]\label{B5.49}
Let the host conditions of Theorem~\ref{B5.48} hold, and let
\(TT,T_2\subseteq H\), \(T_2\subseteq A_0\le H\), \(m,N\in\N\).
Suppose \(x\ge3\), \(b,b',X'>0\), \(\eta,\theta\ge0\), and
\(b'X'\le xb\). Assume that for every \(t\in TT\),
\[
 |a_t/b-x|\le1.
\]
For every \(u_0\in A_0\) and \(0\le j<m\), assume
\[
\begin{gathered}
 w_{u_0j}(\rho)>0\quad(\rho\in\Phi),\\
 a_t/w_{u_0j}(\rho)\le N\quad((t,\rho)\in TT\times\Phi),\\
 b\le w_{u_0j}(\rho)\le b+\eta+1\quad\text{at thick positions},\\
 w_{u_0j}(\rho)\le2b'\quad\text{at the other positions of }\Phi.
\end{gathered}
\]
Suppose the thin-scale estimate and the replacement costs satisfy
\[
\begin{gathered}
 |G_\psi(y)|\le y\theta
 \quad(\psi\in\widehat\Gamma,\ y\ge X'/3),\\
 2f(x+1)\theta\le1/x,\qquad
 6f(x+1)(\eta+1)\le b/x.
\end{gathered}
\]
Let \(E,Z\in\R\). If
\[
 \mathrm{gap}(u_0)\le E\quad(u_0\in A_0)
\]
and
\[
 n_TZ\le
 \sum_{t\in TT}\Re\sum_{\psi\in\Psi_A}
 \psi(e_1(t))G_\psi(a_t/b),
\]
then
\[
 Q\frac{\delta_2}{2}Z\le E+\frac{2Q}{x}.
\]
\end{lemma}

\begin{proof}
With \(A=H\), the surviving set of Lemma~\ref{B2.08} is
\(P=\Psi_A\), and \(\delta=\delta_2\). Apply Theorem~\ref{B2.10}
with \(\mathscr E=E\) and \(Y_*=Z\).
\end{proof}

\section{The four normalized terms}

\begin{lemma}[Both ends of the chain-length window]\label{B5.50}
Let \(m\in\N\) be the length of a chain at a scale \(x>1\),
\(0<\varepsilon<1/52\). Suppose its size condition is
\[
 x^{-\varepsilon/100}\sqrt{x\,x^{-\varepsilon/8}}
 \le m\le
 x^{\varepsilon/100}\sqrt{x\,x^{-\varepsilon/8}}.
\]
Then
\[
 x^{1/2-\varepsilon}\le m\le x^{1/2+\varepsilon/2}.
\]
\end{lemma}

\begin{proof}
The two exponents \(1/2-\varepsilon/16\pm\varepsilon/100\)
lie in \([1/2-\varepsilon,1/2+\varepsilon/2]\). Since \(x>1\),
the bounds follow by monotonicity.
\end{proof}

\begin{lemma}[The top-weight term]\label{B5.51}
Let \(x,X,Y,a_E,n_T,m,b\in\R\), with \(n_T,m,b\ge0\).
If \(x\le Xm\), \(X\le Y\), and \(a_E\le2xb\), then
\[
 n_Ta_E\le2Y(n_Tmb).
\]
In particular, for \(x\ge1\), \(\varepsilon>0\), and
\(m\ge x^{1/2-\varepsilon}\), one may take
\(X=x^{1/2+\varepsilon}\) and
\(Y=x^{1/2+13\varepsilon/8}\).
\end{lemma}

\begin{proof}
The assumptions give \(x\le Xm\le Ym\), so
\[
 n_Ta_E\le2n_Txb\le2n_TYmb.
\]
For the specified powers,
\(x=x^{1/2+\varepsilon}x^{1/2-\varepsilon}\le Xm\)
and \(X\le Y\).
When \(Q>0\), the intermediate bound reads
\(n_Ta_E/Q\le2x/m\).
\end{proof}

\begin{lemma}[The inverse-rank term]\label{B5.52}
Let \(x>0\), \(Y\ge1\), \(n_T,m,b\ge0\), \(k>0\), and
\(a_E\in\R\). If \(xm\le k\) and \(a_E\le2xb\), then
\[
 n_Ta_E\frac{m^2}{k}\le2Y(n_Tmb).
\]
\end{lemma}

\begin{proof}
As \(m\ge0\) and \(xm\le k\), one has \(m^2/k\le m/x\).
Therefore
\[
 n_Ta_E\frac{m^2}{k}
 \le n_T(2xb)\frac m x=2Q\le2YQ.
\]
When \(Q>0\), the sharper normalized estimate is
\(n_Ta_Em^2/(kQ)\le2xm/k\le2\).
\end{proof}

\begin{lemma}[The layer-position term]\label{B5.53}
Let \(x>0\), \(\varepsilon\in\R\), \(m,b,n\ge0\), and
\(n_T\in\R\). Set \(Y=x^{1/2+13\varepsilon/8}\).
If
\[
 m\le x^{1/2+\varepsilon/2},\qquad
 \tfrac12 nx^{-\varepsilon/8}\le n_T,
\]
then
\[
 m(2nmbx^\varepsilon)\le4Y(n_Tmb).
\]
\end{lemma}

\begin{proof}
The density assumption gives \(n_T\ge0\) and
\(n\le2n_Tx^{\varepsilon/8}\). Thus
\[
\begin{aligned}
 2nm x^\varepsilon
 &\le4n_Tm x^{\varepsilon+\varepsilon/8}\\
 &\le4n_Tx^{1/2+\varepsilon/2+\varepsilon+\varepsilon/8}
 =4n_TY.
\end{aligned}
\]
Multiply by \(mb\ge0\).
\end{proof}

\begin{lemma}[The thick-set rank term]\label{B5.54}
Let \(x\ge1\), \(\varepsilon>0\), \(m,b\ge0\), \(n>0\), \(k\ge0\),
and \(n_T,n_2\in\R\).
Set \(Y=x^{1/2+13\varepsilon/8}\). If
\[
 \tfrac12 nx^{-\varepsilon/8}\le n_T,\qquad
 n_2\le n,\qquad k\le n^2x^{1/2-\varepsilon},
\]
then
\[
 m\,\frac{n_2kbx^\varepsilon}{n(nx^{-\varepsilon})}
 \le2Y(n_Tmb).
\]
\end{lemma}

\begin{proof}[Outline]
Rewrite the left side as \((mb)(n_2k/n^2)x^{2\varepsilon}\) and use \(n_2k\le nk\), the upper bound on \(k\), and the lower density bound on \(n_T\) to bound \((n_2k/n^2)x^{2\varepsilon}\le 2n_TY\). Multiplying by \(mb\ge 0\) yields the stated inequality. The complete proof is in \S\ref{prf:B5.54}.
\end{proof}

\section{The square-root estimate}

\begin{theorem}[The size-window bound]\label{B5.55}
Let \(C\ge0\), \(\varepsilon>0\), \(x\ge3\), and
\(n,n_T,n_2,m\in\N\), with \(n>0\). Let \(b>0\),
\(a_E\ge0\), and \(k,E,G_x\in\R\).
Suppose
\[
\begin{gathered}
 a_E\le2xb,\qquad
 \tfrac12 nx^{-\varepsilon/8}\le n_T,\qquad
 x^{-\varepsilon}\le\delta_2=\frac{n_2}{n},\qquad n_2\le n,\\
 x^{1/2-\varepsilon}\le m\le x^{1/2+\varepsilon/2},\qquad
 xm\le k\le n^2x^{1/2-\varepsilon},\\
 E\le x^{C\varepsilon}\mathcal B_x(a_E),\qquad
 26\le x^{3\varepsilon/8}.
\end{gathered}
\]
If
\[
 Q\frac{\delta_2}{2}
 \left(G_x-4x^{1/2+\varepsilon}\right)
 \le E+\frac{2Q}{x},
 \qquad Q=n_Tmb,
\]
then
\[
 G_x\le x^{1/2+(C+3)\varepsilon}.
\]
\end{theorem}

\begin{proof}[Outline]
Summing \ref{B5.51} through \ref{B5.54} gives \(\mathcal{B}_x(a_E)\le 10YQ\), hence \(E/Q\le 10PY\) for \(P=x^{C\varepsilon},Y=x^{1/2+13\varepsilon/8}\). Multiplying by \(2/Q\) and using \(\delta_2^{-1}\le x^\varepsilon\) and \(26\le x^{3\varepsilon/8}\) yields the \(G_x\) bound. The complete proof is in \S\ref{prf:B5.55}.
\end{proof}

For a finite set \(K_p\subseteq\Psi_A\) of retained characters and
\(c\in\R\), define
\[
 \omega_c(\psi)=
 \begin{cases}
 1,&\psi\text{ is trivial on }H,\\
 c,&\text{otherwise},
 \end{cases}
 \qquad
 G_m^A(x)=\sum_{\psi\in K_p}\omega_c(\psi)\Re G_\psi(x),
\]
and
\[
 \mathrm{Sig}_A(x)=
 \sum_{\psi\in\Psi_A\setminus K_p}|G_\psi(x)|.
\]
In the notation of Theorem~\ref{B2.14}, \(A=H\), \(P=\Psi_A\),
and the remainder is \(\Sigma_A=\mathrm{Sig}_A\).

\begin{theorem}[Bound for the retained character sum]\label{B5.56}
Fix real numbers \(C,\varepsilon,C_E,K_o,c,M\), with
\(C\ge0\), \(\varepsilon>0\), and \(|c|\le1\).
There is \(x_1\ge3\), depending on these constants, such that
the following holds for every \(x\ge x_1\).

Let \(H\) be a finite abelian group, \(A_0\le H\), \(TT\subseteq H\),
and \(T_2\subseteq A_0\). Impose all the assumptions of
Lemma~\ref{B5.49} preceding the introduction of \(E\) and \(Z\).
Let \(a_E\ge0\), \(k>0\), and suppose
\[
\begin{gathered}
 a_E\le2xb,\qquad
 \tfrac12 nx^{-\varepsilon/8}\le n_T,\qquad
 x^{-\varepsilon}\le n_2/n,\\
 x^{1/2-\varepsilon}\le m\le x^{1/2+\varepsilon/2},\qquad
 xm\le k\le n^2x^{1/2-\varepsilon}.
\end{gathered}
\]
Let \(K_p\subseteq\Psi_A\) contain \(\mathbf1\). Suppose
\[
 \#\Psi_A\le M,\qquad
 \psi(((1,0),0))=1\quad(\psi\in K_p),
\]
and assume the character bounds
\[
\begin{aligned}
 \left|\sum_{t\in TT}\psi((0,0),t)-cn_T\right|
 &\le Mx^{-1+\varepsilon}n_T
 &&(\psi\in K_p,\ \psi|_H\ne\mathbf1),\\
 |\ell(\psi)|&\le K_o
 &&(\psi\in K_p\setminus\{\mathbf1\}),\\
 |\ell(\psi)|&\le C_E\log x
 &&(\psi\in\Psi_A\setminus K_p).
\end{aligned}
\]
If every \(u_0\in A_0\) satisfies
\[
 \mathrm{gap}(u_0)\le x^{C\varepsilon}\mathcal B_x(a_E),
\]
then
\[
 G_m^A(x)\le
 x^{1/2+(C+3)\varepsilon}+\mathrm{Sig}_A(x).
\]
\end{theorem}

\begin{proof}
Take \(x_1\) as in \eqref{eq:p2c2-arithmetic-threshold} with \(M'=M\).
Equation~\eqref{eq:p2c2-top-lower} applies with
\(G_*=G_m^A(x)-\mathrm{Sig}_A(x)\).
Use Lemma~\ref{B5.49} with \(Z=G_*-4x^{1/2+\varepsilon}\) and
\(E=x^{C\varepsilon}\mathcal B_x(a_E)\). Since \(n_2\le n\),
Theorem~\ref{B5.55} then gives \(G_*\le x^{1/2+(C+3)\varepsilon}\).
\end{proof}

Fix a level \(\Lambda\) with integer parameter \(o=o_\Lambda\) as in Definition~\ref{B1.21}, and fix real
\(e_D,e_L,K,c_{HH},e_K,C_E\) and \(\varepsilon>0\).
Take the threshold of Theorem~\ref{B5.56} with
\(C=\max(e_D+e_L,0)+3\), \(M=4o\), \(c=(3o+1)^{-1}\), and
\(K_o\) as in Lemma~\ref{B2.17}. Enlarge the threshold to \(x_0\)
so that for all \(x\ge x_0\),
\begin{equation}\label{eq:p5c5-large-x}
 \log x\le x^\varepsilon,\qquad
 200(1+K)^2(1+c_{HH})(1+e_K)\le x^\varepsilon,\qquad
 1+4o\le x^\varepsilon.
\end{equation}
Such a threshold exists because \(x^\varepsilon\to\infty\) and
\(\log x/x^\varepsilon\to0\).

Let \(x\ge x_0\), take a frame \(\phi\) and a host \(H_{\rm host}\)
as in Definition~\ref{B2.01}. Write \(H\) for the finite abelian group of the frame and \(A_0\) for its distinguished subgroup. Fix \(TT\subseteq H\). Suppose that \(H_{\rm host}\) and \(TT\) admit chain data at \(x,\varepsilon\) with bounds \(e_D,e_L,K,c_{HH},e_K\) in the sense of Definition~\ref{def:p2c2-core-at-host}. Suppose also that they satisfy the arithmetic requirement of Definition~\ref{B2.58} for the level \(\Lambda\) at \(x,\varepsilon\) with constant \(C_E\).
Use the arithmetic datum of the host and choose chain data as above, with cutoff \(N\), local weights \(w\) and compatible \(\ell\). The models of
Theorem~\ref{B5.33} take \(N=\lceil2a\rceil\).

The window hypotheses give \(e_K\ge0\). The chain's height exponent
lies in \([0,c_{HH}]\), so the second inequality of
\eqref{eq:p5c5-large-x} holds with the chain's exponent.
The chosen chain data satisfy the remaining hypotheses of Theorem~\ref{B5.47}.
Since \(0\le e_{\rm Div}\le e_D\) and \(0\le e_{\rm Loss}\le e_L\),
we have \(C=e_D+e_L+3\) and \(C_{\rm err}\le C\).
The translation map covers \(A_0\), so for each \(s\in A_0\) we may choose \(u\in U_0\) with \(\mathrm{tr}(u)=s\). By Theorem~\ref{B5.47} and
\(\mathcal B_x(a)\ge0\),
\[
 \mathrm{gap}(s)\le x^{C_{\rm err}\varepsilon}\mathcal B_x(a)
 \le x^{C\varepsilon}\mathcal B_x(a).
\]

Take \(K_p=K_H\), the characters whose primitive representatives on
\(H_{\rm host}\) belong to the retained family of \(\Lambda\).
Lemma~\ref{B2.17} supplies the character conditions of
Theorem~\ref{B5.56} with this same \(\ell\).
Lemma~\ref{B2.03} gives the local conditions, and the chosen chain data satisfy the size, weight, cutoff, thin-scale and cost conditions. Apply Theorem~\ref{B5.56}
with \(a_E=a\) to get
\[
 G_m^A(x)\le x^{1/2+(e_D+e_L+6)\varepsilon}+\mathrm{Sig}_A(x).
\]
For this host, Theorem~\ref{B2.18} at \(y=x\) gives
\(G_m^A(x)=G_{m,\Lambda}(x)\). With \(K_p=K_H\),
\(\mathrm{Sig}_A(x)=\Sigma_A(x)\) term by term. Since \(x\ge3\), the arithmetic requirement of Definition~\ref{B2.58} and the compatibility of the chosen \(\ell\) give the hypotheses of Theorem~\ref{B2.19}. Thus
\(\mathrm{Sig}_A(x)\le4o\,x^{1/2+\varepsilon}\).

Since \(e_D+e_L\ge0\), the last inequality of
\eqref{eq:p5c5-large-x} gives
\[
 G_{m,\Lambda}(x)\le(1+4o)x^{1/2+(e_D+e_L+6)\varepsilon}
 \le x^{1/2+(e_D+e_L+7)\varepsilon}.
\]


\part{The window bound VWin}

% ---- p6c1 ----
\chapter{The window bound and its inputs}\label{ch:window-inputs}

\begin{definition}[Local kernels and the total bound]\label{B6.01}
For a number field $E$ and a prime $\ell$, let $S_\ell=\{v:v\mid\ell\}$ and let $G_{E,S_\ell}$ be the Galois group of the maximal extension of $E$ unramified at finite places outside $S_\ell$. For $k\in\Z$ put
\[
 \Sha_{\ell,k}(E)=\ker\left(
 H^1(G_{E,S_\ell},\Q_\ell/\Z_\ell(-k))
 \longrightarrow\prod_{v\in S_\ell}H^1(G_{E_v},\Q_\ell/\Z_\ell(-k))
 \right).
\]
Use continuous cohomology with discrete coefficients and write $s_{\ell,k}(E)$ for the $\Z_\ell$-length of the kernel. For $t\in\N$ and $B\in\R$, say that $E$ satisfies the total bound up to $t$ with bound $B$ if the following conditions hold. For every algebraic closure $\Omega/E$ and every family of local data in this common closure, with places above $S_\ell$ and descended cyclotomic characters chosen for each $\ell$, the lengths are finite for all $\ell$ and $1\le k<t$, and
\begin{equation}\label{eq:p6c1-total-bound}
 \sum_{\ell\in S}\log\ell\sum_{1\le k<t}s_{\ell,k}(E)\le B
\end{equation}
for every finite set $S$ of rational primes.
\end{definition}

\begin{definition}[Fields, prime pairs and page functions]\label{B6.02}
Fix a primitive Dirichlet character $\psi$ modulo $M$. Let $a\ge0$ be an integer and let $q,d$ be primes. Call $E$ a field associated with $(\psi,a,q,d)$ if $E\simeq_\Q\Q(\zeta_N)^H$ for some positive common multiple $N$ of $2^{a+2},M,q,d$, where $\zeta_N$ is primitive and
\[
 H=\{u\in(\Z/N\Z)^\times:u\equiv\pm1\pmod{2^{a+2}},\
 \psi(u)=1,\ (u/q)=(u/d)=1\}.
\]
The last two conditions use Jacobi symbols. For $0<\varepsilon<1/52$ and $U>0$, put $A=U^{\varepsilon/5}$. Call $(q,d)$ a window pair for $\varepsilon$ at $U$ if
\[
 \begin{gathered}
 q\equiv3\pmod8,\qquad d\equiv7\pmod{16},\\
 A/2<\log q\le A,\qquad A(\log U)^2/2<\log d\le A(\log U)^2,\\
 q\mid d+1,\qquad p\mid d+1\quad\text{for every odd prime }p\le A .
 \end{gathered}
\]
A page function $a:\R\to\N$ satisfies $x^{2+8\varepsilon}\le2^{a(x)}\le2x^{2+8\varepsilon}$ for every $x\ge1$.
\end{definition}

\begin{definition}[The window assertion]\label{B6.03}
For $C_t\in\R$ put $t(x)=\lfloor x^{C_t}\rfloor_\N$. The window assertion for $\psi$ with exponent $C_t$ requires that each $0<\varepsilon<1/52$ and each page function $a$ for this $\varepsilon$ admit reals $J\ge1$ and $X_0$ with the following property. For every $X\ge X_0$ and every window pair $(q,d)$ for $\varepsilon$ at $X^{1/J}$, there is a $Y\ge X^{1/J}$ with $Y^2\le X$ such that, for all $x\in[Y,Y^2]$, every field associated with $(\psi,a(x),q,d)$ satisfies the total bound of Definition~\ref{B6.01} up to $t(x)$ with bound $e^{x^\varepsilon/2}$. The window bound is the assertion with $C_t=401$ for every primitive character $\psi$.
\end{definition}

\begin{theorem}[The representation input gives the window]\label{B6.04}
The window bound of Definition~\ref{B6.03} follows from the four groups in Definition~\ref{def:p6c1-classical-inputs}, the published results in Remark~\ref{B6.05ca}, and weak unitary descent in ranks two through four as in Lemma~\ref{B9.36a}. In particular it follows from the representation assertion $\mathcal R$ of Theorem~\ref{B9.02}, with \ref{B9.04}--\ref{B9.06}, using Proposition~\ref{B6.05} to supply the classical inputs. The base change in that assembly is Lemma~\ref{B6.04a}, with the character of Lemma~\ref{B9.36c}. The trace comparisons used in attachment and descent are \ref{B6.04b}--\ref{B6.04c}. The local statements used for attachment include the unlinked exponent bound of Lemma~\ref{B9.38a} and the weak correspondence of \ref{B9.46}--\ref{B9.47}.
\end{theorem}
\begin{proof}[Outline]
Assemble the automorphic data via \ref{B6.05} from \(\mathcal{R}\), the transfer \ref{B6.04a}, and the published inputs of \ref{B6.05ca}. Supplying the descent step in ranks two through four by \ref{B9.36a} lets \ref{B6.06} yield \ref{B6.03} with exponent \(401\). The complete proof is in \S\ref{prf:B6.04}.
\end{proof}

\section{The trace comparison}

Let $E/E^+$ be CM and unramified at every finite place, and let $U$ be the unitary group of the standard positive definite Hermitian form of rank $n\ge1$. At each infinite place $v$ of $E^+$, choose an embedding of $E$ above $v$ and a dominant integral weight $\lambda_v$ of rank $n$. For a dominant integral weight $l$ of rank $r$, write
\[
 H(l)=\{l_j+(r-1)/2-j:0\le j<r\}.
\]
Use the rational twist $\theta_J(g)=J_r(\bar g^{-1})^tJ_r^{-1}$ in rank $r$, where $(J_r)_{ij}=(-1)^i1_{i+j=r-1}$ with indices starting at zero. The rational twisted norm is $g\theta_J(g)$. At a split place the local twist is the identity and the local norm polynomial of $y$ is its ordinary characteristic polynomial. Fix a section $\Sigma$ of the finite places of $E$ over $E^+$. Give the standard integral groups and the maximal compact subgroups of the local centralizers volume one.

Let $\mu$ be a unitary id\`ele class character of $E$ restricting to the quadratic character of $E/E^+$. Assume $\mu_w=(-1)^{\operatorname{ord}_w}$ at inert places and $\mu_v(z)=(z/|z|)^{e_v}$ at infinity, with $e_v$ odd. For monic local polynomials $Q_1,Q_2$ of degrees $a,b$ with separable product, write $O_w^{Q_1}(Q_1Q_2;f_w)$ for the sum of unitary orbital integrals with polynomial $Q_1Q_2$. At an inert place require both root multisets to be stable under conjugate inversion. Weight the orbit of $\gamma$ there by $(-1)^{\operatorname{ord}_w\det(h|_{\ker Q_1(\gamma)})}$, where $h$ is the Hermitian form. A change of basis multiplies the determinant by a norm of even valuation. At a split place the sign is one. Write $SO_w(Q_1,Q_2;\phi_w)$ for the sum of twisted orbital integrals with the indicated norm polynomials. Put
\[
 R=\operatorname{Res}(Q_1,Q_2),\qquad
 D=((-1)^aQ_1(0))^b((-1)^bQ_2(0))^a.
\]
At a split place take $\Delta_w=\mu_w((-1)^{ab}D)|R^2/D|_w^{1/2}$ when $D\ne0$, and zero when $D=0$. At an inert place take $\Delta_w=(-1)^{\operatorname{ord}_w R}|R|_w^{1/2}$. The absolute value sends a uniformizer to the inverse residue cardinality. Matching means
\begin{equation}\label{eq:p6c1-orbital-match}
 SO_w(Q_1,Q_2;\phi_w)
 =\Delta_w(Q_1,Q_2)O_w^{Q_1}(Q_1Q_2;f_w).
\end{equation}
The empty block has polynomial one, and a zeroth power is one.

A monic separable polynomial $P$ over $E$ is elliptic if its roots have modulus one at every complex embedding. Write $h_l(P)$ for the product of the conjugated weight characters evaluated on the roots at the chosen embeddings, and $r(P)$ for the number of irreducible factors. For a finite adelic centralizer $C$ with rational subgroup $\Gamma$ and maximal compact subgroup $C_c$, put
\[
 \operatorname{mass}(C)=
 \frac{\#(\Gamma\backslash C/C_c)}{\#(\Gamma\cap C_c)} .
\]
Let $\kappa(P)$ be the rational twisted centralizer mass divided by $2^{r(P)}$, with value zero for an empty twisted class and one in rank zero. Define
\[
 \mathcal S_{a,b}(l_1,l_2;\phi)=
 \sum_{P_1,P_2}\kappa(P_1)\kappa(P_2)h_{l_1}(P_1)h_{l_2}(P_2)
 \prod_{w\in\Sigma}SO_w(P_1,P_2;\phi_w),
\]
over ordered coprime elliptic polynomials of degrees $a,b$. Let $G_U(\lambda,f)$ sum over rational unitary classes of separable polynomial $P$, with summand the centralizer mass times $h_\lambda(P)$ times the product of finite orbital integrals.

\begin{lemma}[The trace identity for two blocks]\label{B6.04b}\label{B6.04y}
For each $a+b=n$ there is a finite collection of pairs of weights $(l_{1,I},l_{2,I})$ and signs $s_I\in\{1,-1\}$. At each $v$, $I_v$ selects $a$ entries of $H(\lambda_v)$. The infinitesimal characters of the two weights are the selected entries shifted by $-be_v/2$ and the complementary entries shifted by $-ae_v/2$. Each of $(n,0)$ and $(0,n)$ has one term, with sign $+1$ and weight $\lambda$ on the positive-rank block.

Let $f=\bigotimes_{w\in\Sigma} f_w$ be a compactly supported Hecke test at open level $K'=\prod_w K'_w\subseteq\prod_w K_w$, where $K_w$ is the standard integral subgroup of $U$. Assume $K'_w=K_w$ and $f_w=1_{K_w}$ almost everywhere. For each $a+b=n$ take compact locally constant tests $\phi_{a,b,w}$ on the two general linear blocks, equal to the product integral unit almost everywhere. If the tests satisfy \eqref{eq:p6c1-orbital-match} at each $w\in\Sigma$, then
\begin{equation}\label{eq:p6c1-two-block-trace}
 G_U(\lambda,f)=
 \sum_{a+b=n}\sum_I s_I\,\mathcal S_{a,b}(l_{1,I},l_{2,I};\phi_{a,b}).
\end{equation}
If the nonzero set of $f_{w_0}$ consists of elements with separable characteristic polynomial at some $w_0\in\Sigma$, then $G_U(\lambda,f)=d_\lambda^{-1}\operatorname{tr}(f)$ on the forms of weight $\lambda$ and level $K'$, where $d_\lambda$ is the weight representation's dimension.
\end{lemma}
\begin{proof}[Outline]
Regroup the Hecke trace by rational conjugacy and use \ref{B6.04e} for unit orbital factors. Norm reciprocity and Hermitian classification select the rational terms, common centralizers give the masses, and the resultant formula and Weyl minors give the signed block weights. The complete proof is in \S\ref{prf:B6.04y}.
\end{proof}

At a split place identify $U$ with $\GL_n(E_w)$. Let $P$ be the upper block parabolic of type $(a,b)$ with modulus $\delta_P$. For $m=\operatorname{diag}(m_1,m_2)$ the transfer is
\begin{equation}\label{eq:p6c1-split-match}
 \mu_w((-1)^{ab}\det(m_1)^b\det(m_2)^a)\,
 \delta_P(m)^{1/2}\int_{M_{a,b}(E_w)}\int_{\GL_n(\mathcal O_w)}
 f_w(k^{-1}mn(X)k)\,dk\,dX,
 \qquad n(X)=\begin{pmatrix}1&X\\0&1\end{pmatrix}.
\end{equation}
The compact group and the integral matrix lattice have volume one. In the orbital integral the change from the unipotent conjugating coordinate to the off-diagonal coordinate is $Y\mapsto Y-m_1^{-1}Ym_2$. On triangularizing the two blocks its determinant is $\pm R/\det(m_1)^b$. Multiplying the absolute value of this determinant by $\delta_P(m)^{1/2}$ gives the factor $|R^2/D|_w^{1/2}$ in \eqref{eq:p6c1-orbital-match}.

\section{The spherical unit at an inert place}

Let $L/F$ be an unramified quadratic extension of nonarchimedean local fields of characteristic zero. Write $O$ for the integers of $L$ and $q$ for the residue cardinality of $F$. Write $\bar z$ for conjugation and let $v$ be the valuation with value one on a uniformizer. Put $\eta(z)=(-1)^{v(z)}$ for $z\ne0$.

Let $B$ be a finite commutative $L$-algebra with involution $\sigma$ extending conjugation and an element $g$. Write $P_g$ for the characteristic polynomial of multiplication by $g$. If $P_g$ is separable, every root occurs in the minimal polynomial, so $B=L[g]\simeq L[X]/P_g$ is a product of fields. Let $D_B$ be the product of their integer rings. Put
\[
 C_B=(B^\times)^\sigma/\{u\sigma(u):u\in B^\times\},\qquad
 T_B=\{u\in B^\times:u\sigma(u)=1\},\qquad T_B^0=T_B\cap D_B^\times .
\]
For $c\in(B^\times)^\sigma$, let $h_c(y,z)=\operatorname{Tr}_{B/L}(cy\sigma z)$ and let $X_B(c)$ consist of full $O$-lattices stable under $g$ and self-dual for $h_c$. A full lattice is finitely generated over $O$ and spans $B$ over $L$. Define
\[
 N_B(c)=\sum_{[M]\in T_B\backslash X_B(c)}[T_B^0:(T_B)_M],
 \qquad S(B)=\sum_{[c]\in C_B}N_B(c),\qquad
 (T_B)_M=\{u\in T_B:uM=M\}.
\]
A lattice automorphism and its inverse are integral, so $(T_B)_M\subset T_B^0$. Multiplication by $u$ takes lattices for $cu\sigma(u)$ to lattices for $c$. On $B=B_1\times B_2$ with the componentwise involution and generator, put
\[
 K(B_1,B_2)=\sum_{[(c_1,c_2)]\in C_B}\eta(\det h_{c_1})N_B(c_1,c_2).
\]
The determinant uses an $L$-basis of $B_1$. A basis change multiplies it by a norm and preserves valuation parity. In dimension zero the determinant and count are one.

\begin{proposition}[Spherical matching at an inert place]\label{B6.04e}
Let $E/E^+$ be CM and unramified at finite places. At an inert place $w$ put $L=E_w$ and $F=L^{\bar{\ }}$. Let $q_0$ be the residue cardinality of $F$. Let $a+b=n$ with $a,b\ge0$. Take monic $Q_1,Q_2\in L[X]$ of degrees $a,b$ with separable product and
\[
 X^{\deg Q_i}\overline{Q_i}(X^{-1})=\overline{Q_i(0)}Q_i(X).
\]
In $B_i=L[X]/Q_i$ let $g_i$ be the image of $X$ and extend conjugation by $\sigma_i(g_i)=g_i^{-1}$. Then
\begin{equation}\label{eq:p6c1-unit-match}
 (-q_0)^{-v(\operatorname{Res}(Q_1,Q_2))}K(B_1,B_2)=S(B_1)S(B_2).
\end{equation}
In \eqref{eq:p6c1-orbital-match} this matches $1_{K_w}$ with $1_{\GL_a(O)\times\GL_b(O)}$ on the twisted general linear groups.
\end{proposition}

\begin{lemma}[Two lattice counts for a marked vector]\label{B6.04ea}
Suppose $\sigma(g)=g$ and $P_g$ is separable. Let $r=\dim_L B$. For a full $O$-lattice $M$ put $M^*=\{z:\operatorname{Tr}_{B/L}(zM)\subset O\}$. If $m_1,\ldots,m_r$ is an $O$-basis of $M$, put
\[
 s_g(M)=\eta\bigl(\det(\operatorname{Tr}_{B/L}(g^i m_j))_{0\le i<r,\,1\le j\le r}\bigr).
\]
The power basis and the trace pairing give a nonzero determinant, and its sign is independent of the integral basis. Let $G_B(u)$ sum $s_g(M)$ over full lattices satisfying $gM\subset M$, $\sigma M=M$, $u\in M$ and $1\in M^*$. For $c\in(B^\times)^\sigma$ and $w'\in B^\times$,
\begin{equation}\label{eq:p6c1-marked}
 G_B(cw'\sigma w')=\#\{M\in X_B(c):w'\in M\}.
\end{equation}
\end{lemma}
\begin{proof}[Outline]
Apply Beuzart-Plessis's linearized Jacquet–Rallis theorem, checking regularity by the moment matrix. Lattice descent gives the linear trace determinant, and orthonormal bases identify the Hermitian cosets with marked self-dual stable lattices. The complete proof is in \S\ref{prf:B6.04ea}.
\end{proof}

\begin{lemma}[Lattices with prescribed intersection and projection]\label{B6.04eb}
Let $M_i\subset B_i$ be full $g_i$-stable lattices and suppose $R=\operatorname{Res}(P_{g_1},P_{g_2})\ne0$. The number of $(g_1,g_2)$-stable lattices $M\subset B_1\times B_2$ with $M\cap B_1=M_1$ and $\operatorname{pr}_2M=M_2$ is $q^{2v(R)}$. If $\sigma_i(g_i)=g_i$ and $\sigma_i M_i=M_i$, the number also satisfying $\sigma M=M$ is $q^{v(R)}$.
\end{lemma}
\begin{proof}[Outline]
Parametrizing lattices by \(f\in V/H\) turns stability into \(D(f)=g_1f-fg_2\in H\) with \(\det D=R\), giving \(q^{2v(R)}\) by elementary divisors. For self-adjoint lattices, \(\tau(f)=\sigma_1f\sigma_2\) and \(z\) of \ref{B6.04ea} split classes as \(a+zb\), giving \(q^{v(R)}\). The complete proof is in \S\ref{prf:B6.04eb}.
\end{proof}

\begin{lemma}[The self-adjoint identity]\label{B6.04ec}
Let $\sigma_i(g_i)=g_i$ and suppose $P_{g_1}P_{g_2}$ is separable. For $R=\operatorname{Res}(P_{g_1},P_{g_2})$,
\begin{equation}\label{eq:p6c1-lie-count}
 K(B_1,B_2)=(-q)^{v(R)}S(B_1)S(B_2).
\end{equation}
\end{lemma}
\begin{proof}[Outline]
Use conductor bounds to recover weighted lattice counts by normalizing marked counts, and alternate the tests over norm classes. Away from the endpoints, the basis determinant and lifts in \ref{B6.04eb} give \((-q)^{v(R)}\), and normalization removes the endpoint terms. The complete proof is in \S\ref{prf:B6.04ec}.
\end{proof}

\begin{lemma}[Passage from a norm-one generator to a self-adjoint generator]\label{B6.04ed}
Suppose $g_i\sigma_i(g_i)=1$ and $P_{g_1}P_{g_2}$ is separable. Then
\begin{equation}\label{eq:p6c1-group-count}
 K(B_1,B_2)=(-q)^{v(\operatorname{Res}(P_{g_1},P_{g_2}))}S(B_1)S(B_2).
\end{equation}
\end{lemma}
\begin{proof}[Outline]
Split by fixed idempotents and induct, using dual pairs and \ref{B6.04eb} for exchanged factors. Apply \ref{B6.04ec} to a self-adjoint generator of the local order over its unramified coefficient ring, then compare by trace duality and odd residue degree. The complete proof is in \S\ref{prf:B6.04ed}.
\end{proof}

\begin{proof}[Proof of Proposition~\ref{B6.04e}]
The reciprocal identity and monicity give $Q_i(0)\ne0$ and define the involution. Lemma~\ref{B6.04ed} proves \eqref{eq:p6c1-unit-match}.

For the unitary integral a regular element of polynomial $Q_1Q_2$ makes its space a rank-one module over $B=B_1\times B_2$, with adjoint $\sigma$ on $B$. Apply the trace pairing to $h(y,1)$. The adjoint relation gives $h=h_c$ with $c$ fixed and invertible. Changing the module generator multiplies $c$ by a norm, and the centralizer is $T_B$. An integral coset corresponds to a self-dual stable lattice. Every contributing form is standard by the orthonormal basis argument in \ref{B6.04ea}. With integral group and $T_B^0$ of volume one, the coset contributes $[T_B^0:(T_B)_M]$. Since $\ker Q_1(g)=B_1$, its sign is $\eta(\det h_{c_1})$. The signed unitary integral is $K(B_1,B_2)$.

For a twisted block of rank $r>0$ put $\epsilon=(-1)^{r-1}$. A twisted element $\delta$ corresponds to the form $b(y,z)=y^tH\bar z$ with $H=((\delta J_r)^{-1})^t$. Twisted conjugation by $k$ replaces $H$ by $k^tH\bar k$. For $g=\delta\theta_J(\delta)$, $\bar H^t=\epsilon g^tH$, or $\overline{b(y,z)}=\epsilon b(gz,y)$. This also gives preservation by $g$. In the rank-one $B$-module the trace pairing writes $b(y,z)=\operatorname{Tr}_{B/L}(\beta y\sigma z)$, and the relation becomes $\sigma\beta=\epsilon g\beta$. Conversely this relation gives norm $g$ on recovering $\delta$ from $H$. The coefficients for a twisted class differ by norms.

For integral $g$ there is a solution $\beta_0$ invertible in $R_1=O[g]$. Set $a=\epsilon g^{-1}$ and take $\beta_0=z+a\sigma z$. On two exchanged maximal ideals of $R_1$ prescribe residues $1,0$ for $z$. At a fixed maximal ideal use $1$ when $1+a$ has nonzero residue, and otherwise use a scalar in $O$ whose residue differs from its conjugate. In the latter case $\beta_0$ reduces to $z-\sigma z$. The Chinese remainder theorem gives one $z$ with these residues. Thus the twisted classes are represented by $\beta_0c$ for $[c]\in C_B$.

Since $J_r\in\GL_r(O)$, an integral twisted representative is equivalent to a Gram matrix invertible over $O$. Then $\epsilon g$ preserves the lattice by the relation between $H$ and $\bar H^t$. On a stable lattice both $\beta_0$ and $\beta_0^{-1}$ act as automorphisms. Thus the integral test for $\beta_0c$ amounts to $g$-stability and $h_c$-self-duality. The centralizer and stabilizer weights are those of $T_B$, so the stable twisted integral is $S(B)$. A nonintegral generator gives zero since neither an integral twisted representative nor a full stable lattice can have that norm. The rank-zero integral is one.

The two twisted integrals multiply. Since $E_w$ has residue cardinality $q_0^2$, the factor in \eqref{eq:p6c1-orbital-match} is $(-1)^{v(R)}(q_0^2)^{-v(R)/2}=(-q_0)^{-v(R)}$.
\end{proof}

\section{Trace powers and base change}

\begin{lemma}[A bound for convolution powers on a block]\label{B6.04c}
Let $E/E^+$ be as in \ref{B6.04b} and $a,b\ge0$. On the two general linear factors use $\theta=\theta_0$, where $\theta_0(g)=(\bar g^{-1})^t$. Fix Haar measures, a finite matrix multiplicity, compactly supported continuous end kernels with convolution operators $T_1,T_2$ on the rational quotient, a matrix $M$, and $r_0,r_1\ge0$. There exist $C\ge0$ and $N\in\N$ with the following property.

For $k\in\N$ let $V$ be bounded and suppose $T_1VT_2$ has a compactly supported continuous convolution kernel $\Phi$. Put $h=\operatorname{tr}(\Phi M)$ and assume it has a common open finite level on both sides and logarithmic support size at most $r_0+kr_1$ in each factor. Logarithmic size sums the positive logarithms of the local maximum entry norms of $g$ and $g^{-1}$, with squared norms at complex places. In each factor every rational twisted class meeting the nonzero set has separable norm polynomial.

For a matrix $g$ put $W^{\perp_g}=g(W^{\perp_{\rm std}})$, with $\perp_{\rm std}$ the standard Hermitian orthogonal complement. At every infinite place require each slice to be very cuspidal. This means invariance under twisted conjugation by the compact unitary group and
\[
 \int_N h(ug)\,du=0
 \quad\text{for }0\ne W_1\subseteq W_2,\quad
 W_1^{\perp_g}=W_2,\quad W_2^{\perp_g}=W_1 .
\]
Fix the other variables of the slice. Let $N$ be the unipotent radical for $W_1\subseteq W_2$. In rank $d$ it consists of $1+X$ with $X(\C^d)\subseteq W_2$ and $XW_2\subseteq W_1$, subject to $XW_1=0$. Then the regular elliptic twisted geometric sum $\mathcal E(h)$ satisfies
\begin{equation}\label{eq:p6c1-growth}
 |\mathcal E(h)|\le C(1+k)^N\|V\|.
\end{equation}
\end{lemma}
\begin{proof}[Outline]
Adelic reduction gives a compact cut containing the elliptic terms, and alternating flag weights and very cuspidal integration cancel the other terms. Cauchy–Schwarz bounds the diagonal by \(\|V\|\) times a translate count whose integral grows polynomially by unfolding. The complete proof is in \S\ref{prf:B6.04c}.
\end{proof}

\begin{lemma}[The complex Lefschetz integral]\label{B6.04l}
In rank $r$ over $\C$, use $\theta_0$ and a dominant integral weight $l$. Put $\inf(l)=(H(l),-H(l))$ for the two central parameter multisets. The wave multipliers on the exterior algebra with weight coefficients give compactly supported kernels $\phi$ with the very cuspidal property of \ref{B6.04c}. If $\operatorname{val}$ is the multiplier value as a function of the two central parameters, then for $d=\operatorname{diag}(z_i)$ with $|z_i|=1$ and distinct $z_i^2$,
\[
 \int_{\GL_r(\C)}\phi(x^{-1}d\theta_0(x))\,dx
 =\operatorname{val}(\inf(l))\,\overline{\operatorname{ch}_l(z_1^2,\ldots,z_r^2)}.
\]
Here $\operatorname{ch}_l$ is the Weyl character, and the kernel uses the same Haar measure as the integral. For two blocks the formula is the product of the two integrals.
\end{lemma}
\begin{proof}[Outline]
Finite propagation gives compact support, and exterior contraction gives vanishing on flag radicals. On the Hermitian cone the cochain identity reduces the supertrace to degree zero, where the multiplier at \(\inf(l)\) and the flip give the conjugated Weyl character. The complete proof is in \S\ref{prf:B6.04l}.
\end{proof}

\begin{lemma}[Rational and stable geometric coefficients]\label{B6.04m}
For a product finite test $\phi$ of ranks $a,b$ as in \ref{B6.04b}, assume the product of the two norm polynomials is separable wherever $\phi_{w_0}\ne0$ for some $w_0\in\Sigma$. Let $\mathcal G_{a,b}(l_1,l_2;\phi)$ be the sum over rational twisted classes with elliptic norm polynomials in each block. Each summand is weighted by the centralizer mass. The other factors are the two weight characters and the class's finite orbital integrals. With $\Sigma_\infty$ the set of complex places of $E$,
\[
 \mathcal G_{a,b}(l_1,l_2;\phi)
 =2^{(a+b)|\Sigma_\infty|}\mathcal S_{a,b}(l_1,l_2;\phi).
\]
For an adelic product kernel $h$ with finite tests giving $\phi$ under the split and inert conventions and archimedean integrals of \ref{B6.04l}, write $\nu_v$ for the product of the multiplier values at the two weights at $v$. The same rational sum with the adelic orbital measures is
\[
 \mathcal E(h)=c_{a,b}\Bigl(\prod_{v\in\Sigma_\infty}\nu_v\Bigr)
 \mathcal G_{a,b}(l_1,l_2;\phi)
\]
for a nonzero $c_{a,b}$ depending on the fixed Haar measures.
\end{lemma}
\begin{proof}[Outline]
Use \ref{B6.04e} and \ref{B6.04b} to describe fibres by norm symbols, then count real signs by reciprocity and identify classes by the Hasse norm theorem. Combine both block counts with their masses and the archimedean multipliers of \ref{B6.04l}. The complete proof is in \S\ref{prf:B6.04m}.
\end{proof}

\begin{lemma}[The algebra of variable comparison tests]\label{B6.04da}
Use the field \(E\), the character \(\mu\) and the two-block comparison of \ref{B6.04b}. Write \(V\) for the space of definite unitary forms of the chosen weight and finite level. For each term \(s\) of the comparison, of ranks \(a,b\), write \(H_s\) for the space of square-integrable functions on
\[
 (\GL_a(E)\times\GL_b(E))\backslash
 (\GL_a(\mathbf A_E)\times\GL_b(\mathbf A_E))
\]
with the finite matrix multiplicity of the test. Convolution acts on the right. For a compact open subgroup \(L\) write \(e_L=\mathbf1_L/\operatorname{vol}(L)\).

Choose a nonempty finite set \(T'\) of inert spherical testing places containing the prescribed inert places, a finite set \(S_0\) of split testing places containing the prescribed split places, and the remaining finite exceptional set \(J\). Take these sets disjoint. A variable finite test may vary at a chosen split place outside \(T'\cup S_0\cup J\) where \(\mu\) is unramified. The following operators are taken as tuples on \(V\) and on all the \(H_s\).

The reference operator \(a_0\) is the identity on \(V\). On \(H_s\) its finite kernel is the diagonal matrix of level idempotents. Its infinite kernel is the product of the undifferentiated wave kernels of \ref{B6.04l}, normalized to have product of values one at the weights of \(s\).

At a variable split place \(w\), put \(K_n=\GL_n(\mathcal O_w)\), choose a uniformizer \(\varpi_w\), and take the characteristic function of
\[
 K_n\operatorname{diag}(\underbrace{\varpi_w,\ldots,\varpi_w}_{j},
                  1,\ldots,1)K_n,\qquad 0\le j\le n.
\]
On \(V\) use its spherical action. On \(H_s\) replace the reference idempotent at either member of the split pair, in its chosen coordinates, by the transfer \eqref{eq:p6c1-split-match} divided by \(\mu_w((-1)^{ab})\). Keep the other factors of the reference test. For an infinite generator choose a block, an infinite place, one of its two rank factors and a central differential operator from either complex central action of \ref{B6.04l}. Apply that operator to the reference kernel in the indicated factor. On \(V\) and on the other blocks take its scalar value at the indicated weight.

Let \(\mathcal A\) be the unital complex algebra in
\[
 \operatorname{End}(V)\times\prod_s\mathcal B(H_s)
\]
generated by these tuples and their componentwise adjoints. Write \(\rho,R_s\) for the coordinate maps. Fix one separating test at each place of \(T'\cup S_0\). Denote by \(\beta_s\) the convolution on \(H_s\) with their transferred matrix kernels, the reference level idempotents in the other finite factors, and the reference kernel at infinity. On \(V\), let \(u\) act in each automorphic summand by the product of the prescribed \(m\)-fold local convolutions at the exceptional places, where \(m\ge1\). Then \(\mathcal A\) is commutative, \(u\) commutes with \(\rho(\mathcal A)\), and \(\beta_s\) commutes with \(R_s(\mathcal A)\).
\end{lemma}
\begin{proof}[Outline]
The level idempotent acts as identity, transpose gives spherical commutativity, and even time convolution commutes with central operators, with the same identities for adjoints. The variable operators act by scalars on the summands preserved by \(u\). The complete proof is in \S\ref{prf:B6.04da}.
\end{proof}

\begin{lemma}[A spherical vector from an ordered translation character]\label{B6.04db}
Let \(L\) be a finite extension of \(\Q_p\), with ring of integers \(O_L\) and uniformizer \(\varpi\). Write \(K_j=\GL_j(O_L)\), and let \(I_j\subset K_j\) be the inverse image of the upper triangular subgroup over the residue field. Let \(P\subset\GL_j(L)\) be the upper block parabolic with block sizes \(j_1,\ldots,j_s\). For smooth irreducible admissible complex representations \(\sigma_i\) of \(\GL_{j_i}(L)\), form the normalized induction
\[
 M=\operatorname{Ind}_{P}^{\GL_j(L)}(\sigma_1\boxtimes\cdots\boxtimes\sigma_s).
\]
Give \(I_j\) volume one. For \(x_1>\cdots>x_j\) integral, put \(t_x=\operatorname{diag}(\varpi^{x_1},\ldots,\varpi^{x_j})\) and
\[
 A_x=\delta_B(t_x)^{1/2}{\bf1}_{I_jt_xI_j},
 \qquad \delta_B(t)=\prod_{i<i'}|t_i/t_{i'}|_L,
\]
where \(B\) is the upper triangular Borel. Suppose there is a nonzero vector in \(M^{I_j}\) on which every \(A_x\) acts by \(\prod_i\xi_i^{x_i}\), with \(\xi_i\in\C^\times\) and
\[
 |\xi_1|\ge\cdots\ge|\xi_j|.
\]
Then each \(\sigma_i\) has a \(K_{j_i}\)-fixed vector. The union, with multiplicities, of their normalized Satake multisets is \(\{\xi_1,\ldots,\xi_j\}\).
\end{lemma}
\begin{proof}[Outline]
Read the ordered character in a subquotient's Jacquet module by Iwahori projection. Frobenius reciprocity and spherical Langlands classification give a fixed vector, which lifts by averaging, and the normalized constant term recovers the block Satake multiset. The complete proof is in \S\ref{prf:B6.04db}.
\end{proof}

\begin{lemma}[Finiteness of the discrete inducing data]\label{B6.04u}
Let \(E\) be a CM field. In each rank \(1\le d\le n\), fix a compact open finite level, finitely many maximal compact types at infinity, and a bounded set of infinitesimal parameters. Normalize unitary representations of \(\GL_d(\mathbf A_E)\) to have trivial character on the positive real split center by twisting by \(|\det|^{it}\). There are finitely many isomorphism classes in the discrete spectrum with these normalized parameters, a fixed vector at the given finite level, and one of the specified types. The same conclusion holds for tuples of such representations of total rank at most \(n\).
\end{lemma}
\begin{proof}[Outline]
Apply Müller's adelic counting theorem modulo the split centre. Bounded infinitesimal parameters bound the Casimir eigenvalue, and finiteness of the cuspidal and residual counts gives finite lists for each type and hence for tuples. The complete proof is in \S\ref{prf:B6.04u}.
\end{proof}

\begin{lemma}[Translation tests at an inert place]\label{B6.04dc}
Let \(L/L^+\) be the unramified quadratic extension at an inert place in \ref{B6.04b}. Write \(q=\#k_L\), choose a uniformiser \(\varpi\) of \(L^+\), and use an integral Witt basis for the unitary group. Write \(I_U,I_r\) for the upper Iwahori subgroups of the unitary group and of \(\GL_r(L)\), respectively. Give their maximal compact subgroups volume one. For an integral vector \(z\) put \(t_z=\operatorname{diag}(\varpi^{z_i})\).

Take \(x_0>\cdots>x_{n-1}\), put \(i^*=n-1-i\), and set \(h_i=x_i-x_{i^*}\). Assume the positive entries of \(h\) are at least two and their consecutive gaps are at least two. The unitary test is
\[
 u_x=\operatorname{vol}(I_U)^{-1}{\bf1}_{I_Ut_hI_U}.
\]
For \(a+b=n\), partition the indices into sets \(A,B\) of sizes \(a,b\), each invariant under \(i\mapsto i^*\). Write \(x_A,x_B\) in their inherited order. Put
\[
 v_A(x)=\sum_{i\in A,\ j\in B}\min(h_i,h_j),\qquad
 d_{r,z}=\operatorname{vol}(I_r)^{-1}{\bf1}_{I_rt_zI_r}.
\]
Then \(u_x\) matches
\[
 \sum_{A}(-\sqrt q)^{-v_A(x)}\,d_{a,x_A}\otimes d_{b,x_B}
\]
in \eqref{eq:p6c1-orbital-match}. For a finite linear combination of the \(u_x\), its \(k\)-th convolution power, \(k\ge1\), matches the sum over \(A\) of the \(k\)-th convolution powers of the corresponding linear combinations for that same \(A\).
\end{lemma}
\begin{proof}[Outline]
Diagonalize the translation cosets to identify the contributing orbits and compute the resultant factor of \ref{B6.04b}. Additivity of the exponents and resultant valuations in the chamber gives the convolution powers for each fixed assignment. The complete proof is in \S\ref{prf:B6.04dc}.
\end{proof}

\begin{lemma}[The trace-power estimate]\label{B6.04d}
In the definite unitary comparison for an automorphic constituent $\pi$, prescribe a set $S$ of at most one split place and a finite set $T$ of spherical inert places. There are test levels and separating tests as in Lemma~\ref{B6.04da}, with $S\subseteq S_0$ and $T\subseteq T'$. Use its algebra $\mathcal A$, with coordinate actions $\rho$ and $R_s$. Put $D_s=\beta_s^m$, where $m\ge1$ is the repetition exponent of the fixed tests. For each $g\in\mathcal A$ there are $C_s\ge0$, $N_s\in\N$ and $t_s(k)\in\C$ such that for all $k\ge1$,
\begin{equation}\label{eq:p6c1-trace-power}
 \operatorname{tr}((u\rho(g))^k)=\sum_s t_s(k),\qquad
 |t_s(k)|\le C_s(1+k)^{N_s}\|(D_sR_s(g))^k\|.
\end{equation}
The finite level $K$ and finite-dimensional archimedean type spaces $W$ for the final candidates can be chosen from $\pi,S$ before $T$ is chosen.
\end{lemma}
\begin{proof}[Outline]
Fix levels and types from \ref{B6.04da}, then separate candidates by inert interpolation and split Hecke projections. Use \ref{B6.04dc}, \ref{B6.04e}, \ref{B6.04l}, and \ref{B6.04m} in the two-block trace identity for ordered words, then apply the growth bound. The complete proof is in \S\ref{prf:B6.04d}.
\end{proof}

\begin{lemma}[One base change for all places]\label{B6.04a}
Let $F/F^+$ be CM and unramified at finite places, and let $\pi$ be automorphic on the standard definite unitary group of rank $n\ge1$. There is a family of discrete automorphic components on $\GL_{n_i}(\mathbf A_F)$ with $\sum_i n_i=n$ and $n_i=a_id_i$, $a_i,d_i\ge1$. In unnormalized parabolic induction, component $i$ embeds in the induction of
\[
 \rho_i|\det|^{(a_i-1)((d_i-1)/2-j)},\qquad 0\le j<d_i,
\]
where $\rho_i$ is normalized, unitary, conjugate self-dual and cuspidal on $\GL_{a_i}(\mathbf A_F)$.

At each finite place $w$ the induced product of the local components has an irreducible subquotient $P_w$, spherical if $\pi$ is unramified there. At an inert spherical place the Satake datum $(\alpha_i)_{1\le i\le\lfloor n/2\rfloor}$ of $\pi$ gives $P_w$ the multiset $\{\alpha_i,\alpha_i^{-1}\}_i$, with an extra entry $1$ in odd rank. At a split place $P_w$ occurs in the local restriction of $\pi$. At each embedding $\tau:F\hookrightarrow\C$, the union of component infinitesimal characters is $H(\lambda_\tau)$ for the corresponding weight of $\pi$.
\end{lemma}

\begin{lemma}[Galois representations for definite unitary base change]\label{B6.04aa}
Let $F/F^+$ be CM and unramified at every finite place, and let $\pi$ be automorphic on the standard definite unitary group of rank $n\ge1$ with dominant integral weight $\lambda=(\lambda_\tau)_\tau$. Let the discrete components on $\GL_{n_i}(\mathbf A_F)$ ($n_i=a_id_i$, $\sum_i n_i=n$) with normalized unitary conjugate self-dual cuspidal bases $\rho_i$ on $\GL_{a_i}(\mathbf A_F)$ be given by Lemma~\ref{B6.04a}. Fix a prime $\ell$, an isomorphism $\iota:\bar{\Q}_\ell\simeq\C$, and a prolongation of each finite place of $F$.

For each component $i$ and each $0\le j<d_i$, the infinitesimal weights of $\rho_i$ at every embedding $\tau:F\hookrightarrow\C$ are distinct and lie in $(n-d_i)/2+\Z$. Moreover, there exists a continuous semisimple representation $r_{i,j}:G_F\to\GL_{a_i}(\bar{\Q}_\ell)$ with the following properties.
\begin{enumerate}
\item At every finite place $w\nmid\ell$ and every map $R$ of \ref{B9.47}, $\mathrm{WD}_\iota(r_{i,j}|_{G_{F_w}})$ satisfies the local comparison of Definition~\ref{B9.05} with
\[
 R\!\left(\bigl(\rho_{i,w}|\det|_w^{(d_i-1)/2-j}\bigr)\otimes|\det|_w^{(1-n)/2}\right).
\]
In particular, if $\rho_{i,w}$ is spherical with Satake multiset $B_w(\rho_i)$, then $r_{i,j}$ is unramified at $w$ and inverse arithmetic Frobenius has eigenvalues $q_w^{(n-d_i)/2+j}B_w(\rho_i)$.
\item At every finite place $w\mid\ell$, if $\rho_{i,w}$ is spherical, then $r_{i,j}|_{G_{F_w}}$ is crystalline and de Rham. If the local parameter above has zero monodromy and satisfies the toral or full tame test of Definition~\ref{B9.05}, then $r_{i,j}|_{G_{F_w}}$ is potentially semistable and de Rham and satisfies the local comparison of Definition~\ref{B9.06}. In each case its filtration jumps at an embedding $\tau:F\hookrightarrow\bar{\Q}_\ell$ are $(n-d_i)/2+j-z=(n-1)/2-(z+(d_i-1)/2-j)$ as $z$ runs over the infinitesimal weights of $\rho_i$ at $\iota\circ\tau$.
\end{enumerate}
Consequently $r_\pi=\bigoplus_i\bigoplus_{j=0}^{d_i-1}r_{i,j}:G_F\to\GL_n(\bar{\Q}_\ell)$ is a continuous semisimple representation: at every finite place $w$ where $\pi$ is unramified (with spherical transfer $P_w$ of Lemma~\ref{B6.04a} and Satake multiset $B_w(P_w)$), $r_\pi$ is unramified with inverse arithmetic Frobenius characteristic polynomial $P_{w,B_w(P_w)}(X)=\prod_{\alpha\in B_w(P_w)}(X-q_w^{(n-1)/2}\alpha)$ when $w\nmid\ell$, and crystalline when $w\mid\ell$, with filtration jumps $\{k-\lambda_{\iota\circ\tau,k}:0\le k<n\}$ at each $\tau$.
\end{lemma}
\begin{proof}[Outline]
Extracts the weights of \(\rho_i\) from \ref{B6.04a} and applies \ref{B9.02} directly in the even case or after a Hecke twist from \ref{B9.36c} in the odd case. Untwisting via \ref{B6.05c}, Chebotarev density, and summing give \(r_\pi\) and \ref{B9.05}--\ref{B9.06}. The complete proof is in \S\ref{prf:B6.04aa}.
\end{proof}

\begin{lemma}[A nonzero block character]\label{B6.04g}
Use the data of \ref{B6.04d}, with $uv=v$ for a nonzero vector $v$ and $\rho(g)v=p(g)v$ for $g\in\mathcal A$. Put $B_s=D_s^*D_s$ and let $\mathcal C_s$ be the commutative unital $C^*$-algebra generated by $R_s(\mathcal A)$ and $B_s$. For some $s$ there is a character $\chi$ of $\mathcal C_s$ with $\chi\circ R_s=p$ and $\chi(B_s)>0$.
\end{lemma}
\begin{proof}[Outline]
Under the vanishing assumption, use \ref{B6.04da}, a finite sum of squares and the Gelfand norm to get a test equal to one at \(p\) that contracts every block. The trace-power bound and Vandermonde inversion contradict the eigenvalue \(1\). The complete proof is in \S\ref{prf:B6.04g}.
\end{proof}

\begin{lemma}[Spectral realization with the reference convolution]\label{B6.04fa}
In a block of \ref{B6.04da} put \(A_0=R_s(a_0)\), \(X=(\beta_s^m)^*\beta_s^m\), and let \(\chi\) be a character of \(\mathcal C_s\) with \(\chi(A_0)\ne0\). There exist characters \(\eta_j\) converging to \(\chi\) on every variable generator and on \(X\), with \(\eta_j(A_0)\ne0\). For each \(j\) there is a nonzero continuous vector-valued function \(f_j\) on the rational quotient of the block on which \(A_0\), the variable generators and \(X\) act by their respective values at \(\eta_j\). For the reduction height \(H\) of \eqref{eq:p6c1-reduction}, the functions \(H^{-M}f_j\) are square-integrable for a common integer \(M\).
\end{lemma}
\begin{proof}[Outline]
Bound weighted norms of periodized kernel rows by adelic reduction. Cyclic spectral representation and Fubini give weighted square-integrable eigenvectors near \(\chi\), and convolution by the reference kernel gives continuous representatives along a convergent sequence. The complete proof is in \S\ref{prf:B6.04fa}.
\end{proof}

\begin{proposition}[Reading the block character]\label{B6.04f}
For a block $s$ of ranks $a,b$ and a character $\chi$ from \ref{B6.04g}, there are discrete families $P_a,P_b$ of those total ranks, with the shifted infinitesimal characters of \ref{B6.04b}. In their product realization there is a nonzero vector $F$ of the prescribed level and compact types. For each variable finite generator $t$ of Lemma~\ref{B6.04da}, its finite kernel acts on $F$ by
\[
 \frac{\chi(R_s(t))}{\chi(R_s(a_0))}.
\]
Let $b_s$ be the finite matrix kernel of $\beta_s$, write $*$ for convolution, and put $b_s^\dagger(x)_{ij}=\overline{b_s(x^{-1})_{ji}}$. The finite kernel $(b_s^\dagger)^{*m}*b_s^{*m}$ of $X=B_s$ acts by $\chi(X)/\chi(R_s(a_0))^{2m}$. In particular $b_s^{*m}F\ne0$, where a kernel acts by right convolution.
\end{proposition}
\begin{proof}[Outline]
Approximate the character of \ref{B6.04g} by joint eigenvectors in Langlands induction fibres. The weights from \ref{B6.04b} and \ref{B6.04l} bound the inducing data and force the limiting twists to vanish, and the finite kernel eigenvalue equations pass to the limit. The complete proof is in \S\ref{prf:B6.04f}.
\end{proof}

\begin{proof}[Proof of Lemma~\ref{B6.04a}]
Apply Lemma~\ref{B9.36c} with $S=\varnothing$ and $e_v=1$ at every chosen embedding to obtain $\mu$.

Fix the prescribed split set $S$ and a finite inert spherical set $T$. Apply \ref{B6.04d}, \ref{B6.04g} and \ref{B6.04f}. For the resulting block put
\[
 D_T=(P_a\otimes(\mu^b\circ\det))\boxplus
 (P_b\otimes(\mu^a\circ\det)),
\]
where $\boxplus$ concatenates discrete components. The twists add $be_v/2$ and $ae_v/2$ to the respective infinitesimal multisets, so their union is $H(\lambda_v)$. Fix an inert testing place. The finite kernel factors over the places, and its matrix at this place is diagonal in the assignments. If every diagonal entry had zero $m$th power on the corresponding vectors, the global $m$th power would vanish. Choose a surviving entry. It has an allowed assignment, since all other entries are zero. On the local product representation the corresponding Iwahori fixed space has dimension at most $m$ by the choice in \ref{B6.04d}. To pass to that space, hold the other local variables fixed and apply the separating equivariant maps from the induced realization to the local product of the inducing representations. These maps intertwine each coset sum, so at least one detects the nonzero power. Write $Q$ for the resulting linear combination of positive-chamber translations. Since $Q^m\ne0$, Cayley--Hamilton gives an eigenvalue $z\ne0$ of $Q$. Its eigenspace is invariant under every positive-chamber translation by their addition law. In it choose a nonzero common invariant subspace of least dimension. The eigenspace there of any one translation is nonzero and invariant under the others, hence is the whole subspace. A common eigenvector therefore gives a chamber character at which the separating test has value $z\ne0$.
Interpolation identifies the assigned character with the ordered target. To read the individual factors, use normalized chamber translations on the two Iwahori spaces. On their tensor product the eigencharacter for a pair of translations is the product of the prescribed two character values. Simultaneous triangularization on each finite-dimensional factor shows that a pair of its diagonal characters has those same products for every pair of strictly dominant translations. Each individual character is thereby determined up to a constant, and applying its multiplication law to a translation and its square makes that constant one. Each character occurs on an eigenline in its joint generalized eigenspace. After the modulus normalizations of the translations, an allowed assignment takes the target coordinates in their original order within each block. The remaining transfer signs have absolute value one, as do the unramified twists. Thus the two untwisted lists separately have nonincreasing absolute values. Apply Lemma~\ref{B6.04db} in each positive-rank factor. Every inducing constituent at $w$ has a hyperspecial fixed vector; a rank-zero block contributes the empty list. The unramified twists preserve those vectors. Their tensor gives a hyperspecial fixed section of the joined induction by Iwasawa decomposition. Undoing the chamber multipliers gives the two block Satake multisets multiplied by $\mu_w(\varpi_w)^b$ and $\mu_w(\varpi_w)^a$, respectively.

For fixed $S$ there are finitely many rank partitions and partitions of the prescribed infinitesimal multisets. A $K$-fixed vector in an induction gives a uniform open level on the inducing factors by restriction to the Levi in the finitely many compact double cosets. Normalize the split center and fix the levels and compact types. Cuspidal finiteness gives finitely many constituents in a bounded infinitesimal range \cite{harish-chandra-1968}. For discrete constituents the residual classification reduces finiteness to the cuspidal case \cite{moeglin-waldspurger-1989}. The ranks bound the string lengths, which determine the shifts of the cuspidal parameters. Restriction to the Levi bounds the levels. Hence $D_T$ ranges over finitely many classes up to permutation and local equivalence. If every class failed at some inert spherical place, picking one such place for each would give a finite $T$ with no match. Thus one family for each $S$ matches at all inert spherical places and at $S$.

Take $S=\varnothing$ to obtain a family $D$. For each split $w$ repeat with $S=\{w\}$. The two families have the same Satake union almost everywhere. M{\oe}glin and Waldspurger write the discrete constituent on $\GL_{ad}$ as a string of length $d$ on a unitary cuspidal representation of rank $a$ \cite{moeglin-waldspurger-1989}. For $t_j=(d-1)/2-j$, the normalized embedding uses $-t_j$ and the square root of the parabolic modulus contributes $at_j$ on block $j$. Thus the unnormalized exponents are $(a-1)t_j$.

Jacquet--Shalika rigidity recovers the cuspidal support from the cofinite Satake union \cite{jacquet-shalika-1981}. For each unitary cusp form the shifts are centered intervals of step one. Their multiset union determines the string lengths by removing a longest interval and repeating. This identifies the local inducing factors of the two families. Permuting them preserves the class of the finite-length induction in the local Grothendieck group, so the split subquotient occurs for $D$. At inert spherical places take the spherical subquotient, and at the remaining inert places take an irreducible subquotient of the nonzero induction.

The cofinite Satake union is invariant under conjugate duality by inert reciprocity and inverse transpose at split pairs. Rigidity pairs the strings with their conjugate-dual strings of the same length. At each embedding the paired unitary cusp parameters are complex conjugates and are real because their real string shifts lie in $H(\lambda_\tau)$. Distinct paired strings would produce a repeated entry there, whereas successive entries differ by $\lambda_j-\lambda_{j+1}+1\ge1$. Thus each string is conjugate self-dual.
\end{proof}



\section{From representations to the window}

\begin{proposition}[Induction and cuspidal support]\label{B6.05a}
Let \(K\) be a nonarchimedean local field. All representations in this proposition are smooth complex representations. Write \(G_n=\GL_n(K)\), \(\nu(g)=|\det g|_K\), and \(\pi_1\times\pi_2\) for normalized parabolic induction. The cuspidal support of an irreducible representation is the multiset of supercuspidal inducing factors, up to isomorphism and permutation.

\begin{enumerate}
\item Irreducible representations of \(G_n\) are admissible. Normalized Jacquet restriction and normalized parabolic induction preserve finite length. Smooth duality commutes with normalized induction. For inducing data \(V^\vee,V\), the pairing of the induced spaces is obtained by integrating the pointwise pairing over \(\GL_n(\mathcal O_K)\).
\item Let \(P,Q\) be standard parabolics with Levi factors \(M,L\). The normalized \(Q\)-Jacquet module of a representation induced from \(M\) has a filtration indexed by \(W_L\backslash W_{G_n}/W_M\), where \(W\) denotes the Weyl group. For a minimal representative \(w\), the corresponding quotient is obtained by Jacquet restriction in \(M\) along \(M\cap w^{-1}Qw\), transport by \(w\), and induction in \(L\) along \(L\cap wPw^{-1}\). For an exterior tensor product on \(M\), each cell multiplicity is the sum, over the constituent tensors of its Jacquet restriction, of their multiplicities times the corresponding induction multiplicities in \(L\).
\item A segment is a list \(\Delta=(\rho,\nu\rho,\ldots,\nu^{d-1}\rho)\), \(d\ge1\), with \(\rho\) supercuspidal of positive rank. Let \(Z(\Delta)\) and \(L(\Delta)\) be the irreducible subrepresentation and irreducible quotient, respectively, of the product in the displayed order. The direct sum of the Grothendieck groups in all ranks, with multiplication given by induction, is the polynomial ring over \(\Z\) on the classes \(Z(\Delta)\). On the Levi with \(d\) equal blocks the Jacquet module of \(Z(\Delta)\) is the tensor of the displayed factors in order.
\item If \(\pi_1,\pi_2\) are irreducible and unitarizable, then \(\pi_1\times\pi_2\) is irreducible.
\item Let \(S_i\) be the cuspidal support of an irreducible \(\pi_i\). Suppose that for every \(\rho\in S_1\) and \(\rho'\in S_2\) of the same rank,
\[
             \rho\not\simeq \rho'\nu^j\qquad(j=-1,0,1).
\]
Then \(\pi_1\times\pi_2\) is irreducible.
\item Suppose \({\rm char}(K)=0\). An irreducible of positive rank is essentially square-integrable if and only if it is \(L(\Delta)\) for a segment. In particular every square-integrable irreducible has this form. An irreducible of positive rank is tempered if and only if it is a normalized product of square-integrable irreducibles.
\item Two segments are linked if their union is a segment and neither contains the other. In characteristic zero every generic irreducible is a product of the \(L(\Delta)\) for pairwise unlinked segments.
\item In characteristic zero every irreducible has a unique Langlands multisegment \(\mathfrak m\). Its representation \(L(\mathfrak m)\) is the irreducible cosocle of the corresponding standard module, the product of the segment factors in Langlands order. For any two multisegments, \(L(\mathfrak m+\mathfrak n)\) occurs as a subquotient of \(L(\mathfrak m)\times L(\mathfrak n)\).
\end{enumerate}
\end{proposition}
\begin{proof}[Outline]
Use the geometric lemma for the Jacquet filtration, Bernstein's unitary product theorem and the segment and tempered classifications. The support condition excludes linkage across the factors, and multisegment multiplication supplies the constituent indexed by the sum. The complete proof is in \S\ref{prf:B6.05a}.
\end{proof}

\begin{proposition}[Jacquet fixed vectors, finite groups and units]\label{B6.05b}
Keep the notation of Proposition~\ref{B6.05a}. Put \(K_n^1=\ker(\GL_n(\mathcal O_K)\to\GL_n(k_K))\), where \(k_K\) is the residue field.
\begin{enumerate}
\item Let \(P=MU\) be a standard parabolic of \(G_n\). Put \(J=K_n^1(U\cap\GL_n(\mathcal O_K))\) and \(J_M=M\cap K_n^1\). For an admissible \(V\), the quotient map to the \(U\)-coinvariants gives a bijection \(V^J\to(V_U)^{J_M}\).
\item Suppose \({\rm char}(K)=0\). If every member of the cuspidal support of an irreducible \(V\) has a nonzero first-congruence-fixed vector in its own rank, then \(V^{K_n^1}\ne0\).
\item Suppose \({\rm char}(K)=0\). Fix a nontrivial continuous unitary additive character \(\psi\) of \(K\). Let \(n\ge2\) and let \(\pi_1,\pi_2\) be irreducible generic representations of \(G_n\) with the same central character. If \(\gamma(s,\pi_1\times\tau,\psi)=\gamma(s,\pi_2\times\tau,\psi)\) as rational functions of \(q^{-s}\), for every irreducible generic \(\tau\) of \(G_r\) and every \(1\le r\le\lfloor n/2\rfloor\), then \(\pi_1\simeq\pi_2\). Here \(q\) is the residue cardinality. In the application, \(n=m+1\), the two representations are supercuspidal, and for every irreducible \(\tau\) of \(G_m\) their Whittaker functional equations have a common scalar for the same \(\psi\) and Haar measure. These data imply the hypotheses just stated.
\item For a complex representation \(V\) of a finite group \(G\), there are subgroups \(H_i\), linear characters \(\lambda_i\) of \(H_i\), and integers \(a_i\) such that
\[
 \operatorname{tr}_V-(\dim V)1_G
   =\sum_i a_i\operatorname{Ind}_{H_i}^{G}(\lambda_i-1_{H_i}).
\]
\item For a prime \(p\) and \(d\ge1\), regular characters of \(\F_{p^d}^{\times}\), modulo \(p\)-power conjugation, parametrize the irreducible cuspidal characters of \(\GL_d(\F_p)\), with the degree and elliptic values of \ref{B7.05}.
\item For a number field \(E\) and \(m,b\ge1\), there is an integer \(a\ge1\) prime to \(b\) such that \(u\in\mathcal O_E^\times\) and \(u-1\in a\mathcal O_E\) imply \(u=v^m\) for some \(v\in\mathcal O_E^\times\).
\end{enumerate}
\end{proposition}
\begin{proof}[Outline]
Use Moy–Prasad and \ref{B6.05a} for Jacquet fixed vectors, then derive the local converse hypotheses from rank \(m\) by multiplicativity and ramified twists. Apply Brauer induction, \ref{B7.31}, and Chevalley's congruence theorem for the remaining assertions. The complete proof is in \S\ref{prf:B6.05b}.
\end{proof}

\begin{proposition}[Automorphic and character inputs]\label{B6.05c}
Use normalized Satake parameters, and let \(q_v\) denote the residue cardinality at a finite place \(v\).
\begin{enumerate}
\item On the unitary groups of the Hermitian spaces in the host, a smooth cuspidal automorphic function of uniform moderate growth is rapidly decreasing. If \(h\) is the adelic height and \(h_0(g)=\inf_{\gamma\in U(F)}h(\gamma g)\), then for each integer \(N\ge0\), \(\sup_g h_0(g)^N|\phi(g)|<\infty\).
\item On these groups fix an open finite level, a continuous finite-dimensional compact type at infinity and the eigenvalues of the archimedean central differential operators. The corresponding automorphic forms span a finite-dimensional space. The same assertion holds for \(\GL_a\) over a totally complex number field with the infinitesimal character fixed at every embedding.
\item On \(\GL_a\) over a totally complex field the unitary cuspidal space, with the split central direction removed, is a discrete Hilbert sum of irreducibles. Its smooth compact-finite constituents are cuspidal automorphic representations. A smooth cuspidal automorphic form whose absolute value is invariant under the split centre is bounded.
\item Over a totally complex field, let two finite lists of cuspidal representations of positive ranks have the same union, with multiplicity, of their Satake multisets at almost every finite place. The lists agree up to isomorphism and permutation. For a unitary cuspidal representation and each spherical \(v\), every Satake root satisfies \(q_v^{-1/2}<|\alpha|<q_v^{1/2}\).
\item For a cuspidal representation of \(\GL_a\) over a totally complex field, inverse transpose realizes the contragredient. The induced bilinear pairing of the representation with its inverse-transpose transform is nondegenerate in both variables.
\item For a polarized regular algebraic cuspidal input as in \ref{B9.02}, the attached representation has the weights and local comparisons of \ref{B9.04}--\ref{B9.06}.
\item Let \(E\) be totally complex, \(a\ge1\), and \(\pi\) cuspidal on \(\GL_a(\mathbf A_E)\). Let \(m_\tau\) be its multiset of infinitesimal weights at \(\tau:E\hookrightarrow\C\). If \(m_\tau\subset (a-1)/2+\Z\) at every embedding, there is one \(w\in\Z\) such that \(m_{\bar\tau}=w-m_\tau\) for all \(\tau\). If in addition each multiset has distinct entries, \(\pi\) has nonzero relative Lie algebra cohomology for an algebraic coefficient system.
\item For each number field \(E\), global reciprocity gives an open quotient map from the idèles onto \(G_E^{\rm ab}\). Its kernel is the closure of the product of \(E^\times\) and the connected archimedean subgroup. Local units map onto the inertia image, a uniformizer maps to inverse arithmetic Frobenius modulo inertia, and the closure of the local image is the decomposition image. For a nonzero integral modulus \(\mathfrak f\), the narrow ray class field is unramified away from \(\mathfrak f\) at finite places. Its Frobenius elements away from \(\mathfrak f\) generate the Galois group. An ideal prime to \(\mathfrak f\) pairs trivially with all its characters precisely when it has a totally positive generator congruent to one modulo \(\mathfrak f\).
\item Fix \(\iota:\overline{\Q}_\ell\simeq\C\). Let \(\chi\) be an algebraic Hecke character of \(E\) with angular and norm exponents given on the connected archimedean idèles by its integral infinity type \(n_\tau\), and let \(\mu\) be its continuous \(\ell\)-adic realization with this reciprocity convention. At \(v\mid\ell\) it is de Rham and potentially crystalline, crystalline when \(\chi_v\) is unramified, with filtration jump \(-n_{\iota\tau}\) at each induced embedding \(\tau:E\hookrightarrow\overline{\Q}_\ell\). Its Weil--Deligne character is \(\chi_v\) composed with local reciprocity in the Weil-to-multiplicative-group direction, with zero monodromy. These assertions hold for the compatible period identifications of \ref{B9.06}.
\item For \(K/\Q_\ell\) finite and a continuous finite-dimensional \(\Q_\ell\)-representation \(V\) of \(G_K\), if \(V\) is de Rham, it becomes semistable over a finite extension of \(K\).
\end{enumerate}
\end{proposition}
\begin{proof}[Outline]
Use rapid decrease and finiteness, compare reciprocal Satake roots by isobaric uniqueness, and apply \ref{B9.02} and Clozel's purity and cohomology results. Reciprocity, rank-one period comparison, and Berger's monodromy theorem give the character and local Galois assertions. The complete proof is in \S\ref{prf:B6.05c}.
\end{proof}

\begin{proposition}[Finite cohomology and its local conditions]\label{B6.05d}
Let \(r\ge1\), let \(\ell\) be prime, put \(R=\Z/\ell^r\), and set \(A^\dagger=\operatorname{Hom}_R(A,\mu_{\ell^r})\) for a finite discrete \(R\)-module \(A\) with continuous Galois action.
\begin{enumerate}
\item At a finite place \(v\) of a number field \(F\), write \(H_v^i(A)=H^i(F_v,A)\). These groups are finite in degrees \(0,1,2\) and vanish in higher degrees. Cup product and the local invariant give perfect pairings of \(H_v^i(A)\) with \(H_v^{2-i}(A^\dagger)\). One has
\[
 \frac{|H_v^1(A)|}{|H_v^0(A)|\,|H_v^2(A)|}
 =
 \begin{cases}
 |A|^{[F_v:\Q_\ell]},&v\mid\ell,\\
 1,&v\nmid\ell .
 \end{cases}
\]
At \(\ell\) all three assertions also hold for the absolute Galois group of each finite subextension \(K/\Q_\ell\) of the chosen algebraic closure, with \([K:\Q_\ell]\) in the formula.
\item Suppose \(v\nmid\ell\) and inertia acts trivially on \(A\). Put
\(H^1_{\rm ur}(F_v,A)=\ker(H_v^1(A)\to H^1(I_v,A))\).
Then \(H^1_{\rm ur}(F_v,A)\) and \(H^1_{\rm ur}(F_v,A^\dagger)\) are each the exact annihilator of the other.
\item Let \(S\) be a finite set of finite places containing every place above \(\ell\), and let \(A\) be a module for \(G_{F,S}\), the Galois group allowing ramification at \(S\) and infinity. Write \(h_i=|H^i(G_{F,S},A)|\) for \(0\le i\le2\), and \(a_\infty=\prod_{v\mid\infty}|H^0(F_v,A)|\). These three global groups are finite, and
\[
             h_0h_2|A|^{[F:\Q]}=h_1a_\infty .
\]
For odd \(\ell\), cohomology vanishes in degrees at least three. For odd \(\ell\) or totally complex \(F\), the normalized nine-term sequence, the duality of the localization kernels, and the higher vanishing are those of \ref{B6.09}.
\item Suppose \(\ell\) is odd. For each \(v\in S\) choose an \(R\)-submodule \(L_v\subset H^1_v(A)\), and let \(L_v^\perp\) be its orthogonal complement for \(A^\dagger\). Let \(S_L(A)\) consist of the global classes with restrictions in \(L_v\) at \(S\), using \(G_{F,S}\). Then
\[
 \frac{|S_L(A)|}{|S_{L^\perp}(A^\dagger)|}
 =\frac{|H^0(F,A)|}{|H^0(F,A^\dagger)|\,a_\infty}
   \prod_{v\in S}\frac{|L_v|}{|H^0(F_v,A)|}.
\]
\item Let \(K/\Q_\ell\) be finite and let \(V\) be a finite-dimensional continuous de Rham representation over \(\Q_\ell\). Define \(H_f^1(K,V)\) as in \ref{B8.41h}. Then
\[
 \dim H_f^1(K,V)=\dim H^0(K,V)
        +\dim_{\Q_\ell}\bigl(D_{\rm dR}(V)/{\rm Fil}^0\bigr).
\]
Under the Tate pairing its annihilator is \(H_f^1(K,V^\vee(1))\). If \(T\subset V\) is a stable full lattice, take the inverse image of the finite condition for \(T\) and the image of the rational finite condition for \(V^\vee(1)/T^\vee(1)\). These conditions too are exact annihilators.
\end{enumerate}
\end{proposition}
\begin{proof}[Outline]
Use inflation, local duality and the local Euler characteristic of \ref{B8.41p} to prove unramified orthogonality. Apply \ref{B6.09} and the Euler characteristics to count the Selmer groups, and use Bloch–Kato for the de Rham local condition. The complete proof is in \S\ref{prf:B6.05d}.
\end{proof}

\begin{proposition}[Abelian schemes and periods]\label{B6.05e}
\begin{enumerate}
\item Let \(R\to R_0=R/I\) be a surjection of commutative rings with divided powers on \(I\). Suppose there is an integer \(N\) for which every product of divided powers of elements of \(I\) of total divided-power degree at least \(N\) vanishes. For an abelian scheme \(A_0/R_0\), evaluate its first crystalline homology on this thickening. Any lift identifies this module with its first de Rham homology. Given one lift, every direct summand reducing to the Hodge submodule is the Hodge submodule of a lift with the specified reduction. A homomorphism between special fibres lifts uniquely between two specified lifts if its crystalline map preserves the chosen filtrations.
\item For an abelian scheme \(A/R\) of relative dimension \(g\), \(H^1_{\rm dR}(A/R)\) is projective of rank \(2g\). All its classes are primitive for the group law. Cup product identifies its exterior square with \(H^2_{\rm dR}(A/R)\). The Hodge sequence is exact, with projective end terms of rank \(g\), and commutes with arbitrary base change. In homology the Hodge submodule is isotropic for the pairing of a line bundle, and a principal polarization gives a perfect pairing.
\item The dual abelian scheme and its Poincaré bundle give biduality and commute with base change. On a thickening as in (1), a smooth separated quasi-compact lift of a scheme computes its crystalline cohomology by its relative de Rham complex, in particular in degrees one and two.
\item For a complex abelian variety with uniformization \(\C^g/\Lambda\), integration defines \(\Lambda\to H^{\rm dR}_1(A)\). The image spans, and the projection to \({\rm Lie}(A)\) is the lattice inclusion in the uniformizing coordinates. The map commutes with endomorphisms. If \(H\) is the Appell--Humbert form of a line bundle, the pullback of its de Rham pairing is \(\epsilon\,2\pi i\,\Im H\) on \(\Lambda\), with the same \(\epsilon\in\{1,-1\}\) for all such data.
\item For an abelian scheme over a smooth affine complex base, pull the family back to a holomorphic parameter neighborhood of \(0\). A given period lattice at \(0\) extends locally to holomorphic lattice vectors in a local trivialization of the Lie bundle. This continuation commutes with endomorphisms. A continued cycle paired with a relative de Rham class has a holomorphic period \(P(z)\). If \(t_1,\ldots,t_d\) lie in the augmentation ideal of a divided-power thickening over \(\C\), and every product of \(s+1\) of them vanishes, evaluation on the base map given by this \(s\)-jet gives the crystalline period
\[
                  \sum_{|\alpha|\le s}
                       \frac{\partial^\alpha P(0)}{\alpha!}t^\alpha .
\]
\item Let \(A\) be an abelian variety over the fraction field of a discrete valuation ring with perfect residue field. Good reduction is equivalent to trivial inertia on \(T_p A\) for some prime \(p\) invertible in the ring. Equivalently all torsion of order invertible in the ring is unramified, or this holds for infinitely many such orders. For each invertible prime \(p\), potential good reduction is equivalent to finite inertia image on \(T_p A\). If the residue field is finite and \(A\) has complex multiplication, it has potential good reduction.
\end{enumerate}
\end{proposition}
\begin{proof}[Outline]
Use Grothendieck's deformation theorem, de Rham cohomology and the Poincaré bundle, with divided-power contraction for crystalline comparison. Obtain periods by integration and Gauss–Manin Taylor transport, and apply the Serre–Tate reduction criteria. The complete proof is in \S\ref{prf:B6.05e}.
\end{proof}

\begin{proposition}[Complex multiplication and special values]\label{B6.05f}
\begin{enumerate}
\item Let \(K\) be imaginary quadratic, \(\mathfrak a\) a nonzero fractional ideal, and fix an analytic uniformization by \(\C/\mathfrak a\). If \(\sigma\in\operatorname{Aut}(\C/K)\) restricts to \(\operatorname{Art}_K(s)^{-1}\) on \(K^{\rm ab}\), with \(\operatorname{Art}_K\) normalized as in Definition~\ref{B1.14k}, its conjugate has uniformization by \(\C/s_f^{-1}\mathfrak a\). On torsion labels the transport is multiplication by the inverse finite idèle. Here the lattice notation means the fractional ideal with the indicated finite completions.
\item The field \(K(j(\C/\mathcal O_K))\) is the Hilbert class field. Its degree is the class number, and its conjugates of \(j(\C/\mathcal O_K)\) are \(j(\C/\mathfrak a)\), equal for two fractional ideals exactly when their classes agree.
\item Let \(F\) be totally real, \(\mathfrak f\ne0\) an integral modulus, and \(b\) a rational function on its narrow ray classes. Use the Dirichlet series on integral ideals prime to \(\mathfrak f\), weighted by \(b\), and its continuation \(L(s,b)\). Then \(L(1-k,b)\in\Q\) for \(k\ge1\). If \(\mathfrak f=(\ell^n)\), \(n\ge1\), and \(b\) takes values in \(\Z_\ell\), then, for each integral ideal \(\mathfrak c\) prime to \(\ell\), with \(b_{\mathfrak c}(\mathfrak a)=b(\mathfrak c\mathfrak a)\),
\[
           L(1-k,b)-(\mathrm N\mathfrak c)^k L(1-k,b_{\mathfrak c})
              \in\Z_\ell .
\]
\item For an ordinary ideal class \(C\) of a totally real field \(F\) of degree \(n\), put \(Z_C(s)=|D_F|^{s/2}\Gamma_{\R}(s)^n\zeta_C(s)\), with \(\Gamma_{\R}(s)=\pi^{-s/2}\Gamma(s/2)\). If \(CC'=[\mathfrak d_F]\), the class of the different, then \(Z_C(s)=Z_{C'}(1-s)\) wherever the two gamma factors are finite.
\item Let \(\chi\) be a unitary Hecke character of a totally complex field with a choice of one embedding per complex place. Suppose its angular exponents \(n_\tau\in\Z\) include a nonzero exponent. In the completion by the local gamma factors, with the local Tate factors taken at the conductor of \(\chi\), one has
\[
 \Lambda(s,\chi)=(-i)^{\sum_\tau |n_\tau|}
       \prod_{v\nmid\infty}\epsilon_v(s,\chi_v)\,\Lambda(1-s,\chi^{-1})
\]
in the additive-character convention of \ref{B1.15a}.
\end{enumerate}
\end{proposition}
\begin{proof}[Outline]
Apply idèlic complex multiplication to lattices, torsion and \(j\)-invariants, and extend ray-class rationality and integral differences by linearity. Use the partial-zeta equation for dual lattice classes and multiply Tate's local functional equations. The complete proof is in \S\ref{prf:B6.05f}.
\end{proof}

\begin{definition}[Classical inputs for the window bound]\label{def:p6c1-classical-inputs}
Write \(h_N\) for the conjunction of the following four groups, in the precise forms of \ref{B6.05a}--\ref{B6.05f}.
\begin{enumerate}
\item For general linear groups over nonarchimedean local fields, use Bernstein--Zelevinsky finiteness and the geometric lemma, the Moy--Prasad isomorphism for Jacquet modules, and the classification by segments and multisegments. This includes the polynomial description of the Grothendieck ring, the classification of essentially square-integrable representations by segments, irreducibility for products of unitarizable irreducibles and for the separated supports of \ref{B6.05a}(5), tempered induction from square-integrable factors, and the genericity criterion for unlinked segments. Use also Langlands exhaustion and the unique irreducible quotient of a standard module, the local converse theorem and depth-zero induction. The finite-group inputs are Brauer induction for characters of degree zero and Green's parametrization of cuspidal characters by Frobenius orbits of regular characters. The unit congruence input says that for each number field and positive integers \(m,b\), there is a positive integer \(a\) prime to \(b\) such that every unit congruent to \(1\) modulo \(a\) is an \(m\)th power of a unit.
\item For the proof of \ref{B6.04f} also use Langlands' spectral decomposition of the full adelic $L^2$ space into normalized inductions from discrete Levi data, integrated over the unitary characters of the real split centre \cite[Thm.~7.2]{arthur-2005}. For automorphic representations use rapid decay of cusp forms, Harish--Chandra finiteness, finiteness of spaces of general linear automorphic forms at fixed data, the cuspidal spectral decomposition, Jacquet--Shalika multiplicity one and Gelfand--Kazhdan duality. Include the representation assertion of Theorem~\ref{B9.02} and Clozel's purity and cohomology theorems. For characters use global Artin reciprocity, ray class fields and the local comparison for the \(\ell\)-adic realizations of algebraic Hecke characters. A de Rham representation over a finite extension of \(\Q_\ell\) is potentially semistable.
\item For Galois cohomology with finite coefficients use local and global finiteness and Euler characteristic formulas, local Tate duality and the Poitou--Tate sequence of Theorem~\ref{B6.09}. At a finite place away from \(\ell\) where the coefficient module is unramified, its unramified condition and the unramified condition for its Tate dual are exact annihilators by \ref{B6.05d}(2). At the places above \(\ell\) use the local conditions of the Selmer problem and their annihilators as in \ref{B6.05d}(4),(5). Include the Greenberg--Wiles formula and the Bloch--Kato local conditions. Local finiteness, duality and the local Euler characteristic formula are required also over each finite extension of \(\Q_\ell\) in the chosen algebraic closure.
\item For abelian schemes use deformation across nilpotent divided-power thickenings, the Hodge sequence and base change for de Rham cohomology, duality and crystalline comparison. For complex abelian varieties use comparison with the period lattice and its Gauss--Manin variation in families. At a discrete valuation use the good and potential good reduction criteria. The remaining assertions are elliptic CM reciprocity, generation of the Hilbert class field by \(j\), rationality and integral translated differences for totally real ray-class values, and the two functional equations used in Lemma~\ref{B1.15a}. Their precise forms are \ref{B6.05e} and \ref{B6.05f}.
\end{enumerate}
\end{definition}

\begin{remark}[Published inputs used in the proof of the window bound]\label{B6.05ca}
In addition to the four groups in Definition~\ref{def:p6c1-classical-inputs}, the proofs of the following numbered results use the indicated published theorems.

\begin{enumerate}
\item For a normalized cuspidal newform \(f\) of weight \(h\ge2\), let \(V_\lambda(f)\) be its representation in Kato's normalization at a finite place \(\lambda\) of its coefficient field. Let \(K/\Q\) be finite abelian, \(1\le r\le h-1\), \(2r\ne h\), and \(T\subset V_\lambda(f)(r)=V\) a \(G_\Q\)-stable lattice. The subgroup of \(H^1(K,V/T)\) whose localization at each place \(w\) belongs to the image of \(H_f^1(K_w,V)\) is finite. For fixed \(K,f,r\) it is zero at almost all \(\lambda\), for every such lattice. Lemma~\ref{B6.43b} uses finiteness with \(K\) quadratic, \(h=2e+1\) and \(r=e\ge1\). Thus \(2r=h-1\), including at \(e=1\). \cite[\S14.1 and Thm.~14.2(1)]{kato-2004}

\item Let \(L\) be CM with maximal totally real subfield \(L^+\), let \(n\ge2\), and let \((\pi,\chi)\) be a regular algebraic polarized cuspidal automorphic representation of \(\GL_n(\mathbf A_L)\). Fix \(\iota:\overline{\Q}_\ell\simeq\C\). If \(r_{\pi,\iota}\) is irreducible on \(G_{L(\mu_{\ell^\infty})}\), then
\[
 H_f^1(L^+,\operatorname{ad}r_{\pi,\iota})=0 .
\]
The coefficient space is the full matrix algebra with conjugation on \(G_L\) and the negative transpose action determined by the polarization at complex conjugation. In Theorem~\ref{B6.53}, the cuspidal witness supplies \(\pi\), and the image of \(G_{L(\mu_{\ell^\infty})}\) is exactly the subgroup in its third hypothesis. \cite[Thm.~6.1]{thorne-2022}

\item For every finite abelian number field \(L\) and prime \(\ell\), the class-group Iwasawa module of the cyclotomic \(\Z_\ell\)-extension has \(\mu_\ell=0\). The finite-layer form used in Theorem~\ref{B6.17} is a bound, independent of \(j\), on the \(\F_\ell\)-dimension of
\(\mathrm{Cl}(L(\mu_{\ell^j}))/\mathrm{Cl}(L(\mu_{\ell^j}))^\ell\).
The bases in that theorem contain \(\mu_\ell\) for odd \(\ell\), and contain \(i\) for \(\ell=2\). Lemma~\ref{B6.18} applies the bound to these abelian bases in the proof of \ref{B6.07}. \cite{ferrero-washington-1979}

\item Let \(L/F\) be an unramified quadratic extension of nonarchimedean local fields of characteristic zero, and let \(\eta\) be its quadratic character. Embed \(\GL_r\) in the upper left corner of \(\GL_{r+1}\). For a relatively regular semisimple \(Y\in M_{r+1}(F)\), put
\(\Delta(Y)=\eta(\det(e^*,e^*Y,\ldots,e^*Y^r))\), using the indicated rows and the last coordinate row \(e^*\). Then
\[
 \Delta(Y)\int_{\GL_r(F)}
 {\bf1}_{M_{r+1}(\mathcal O_F)}(h^{-1}Yh)\eta(\det h)\,dh
 =
 \int_{U_r(F)}
 {\bf1}_{M_{r+1}(\mathcal O_L)}(u^{-1}Xu)\,du
\]
if \(Y\) is \(\GL_r(L)\)-conjugate to a self-adjoint \(X\) for the standard Hermitian form. The right side is zero if there is no such match. The integral subgroups on both sides have volume one. In Lemma~\ref{B6.04ea} the nonzero determinant of the moment matrix makes the marked column and row cyclic and cocyclic, respectively. These are the relative regularity tests for the matrices used there. \cite[Thm.~1]{beuzart-plessis-2021}

\item Fix an algebraic closure of a number field \(F\) and a finite set \(S\) of places. Let \(\mathcal L\) be a family of cyclic extensions of prime degree such that for every finite extension \(D/F\) and every finite \(v\notin S\), some \(L\in\mathcal L\) is linearly disjoint from \(D\) and splits \(v\) completely. Suppose given continuous semisimple representations \(r_L:G_L\to\GL_n(\overline{\Q}_\ell)\), invariant up to isomorphism under \(\Gal(L/F)\), whose restrictions from any two members to their compositum are isomorphic. There is a continuous semisimple representation over \(F\) restricting to \(r_L\) at every member. For Lemma~\ref{B6.06c} the members are quadratic. Infinitely many split at the prescribed place, and the Galois closure of \(D\) contains only finitely many quadratic fields. This verifies generality with \(S=\varnothing\). \cite[\S1, Lem.~1]{sorensen-2020}

\item Let \(O\) be a DVR, let \(M\) be a finite free \(O\)-module, and let a commuting family of endomorphisms preserve \(M\). A simultaneous system of eigenvalues on a nonzero vector of \(M/\mathfrak m_OM\) lifts, after finite extension of the fraction field and localization at a prime above \(\mathfrak m_O\), to the system of a characteristic-zero simultaneous eigenvector. In Lemma~\ref{B6.48ba} the lattice is the integral Fourier lattice of cusp forms at the indicated weight, level and character. Its Hecke operators commute, the Fourier coefficient at one makes the residual eigenvector nonzero, and the congruences include the operators at the level primes. \cite[Lem.~6.11]{deligne-serre-1974}

\item For a quaternion algebra \(D/F\), the global Jacquet--Langlands correspondence identifies the non-character discrete automorphic representations of \(D^\times\) with the cuspidal representations of \(\GL_2(\mathbf A_F)\) which are essentially square-integrable at the ramified places of \(D\). It preserves central characters and identifies the local actions at split places. In Lemma~\ref{B6.49b}, \(F\) is real quadratic and \(D\) ramifies at exactly the two real places. A holomorphic discrete series of weight \(k\) corresponds there, with the positive center removed, to
\(\operatorname{Sym}^{k-2}\otimes\operatorname{nrd}^{\,1-k/2}\), where \(\operatorname{nrd}\) is reduced norm. Both weights exceed two, so quaternionic characters contribute zero. Hence the traces agree for every finite Hecke word at the stated level, including words at the two places over \(\ell\). The congruence between different weights is the further calculation of that lemma. \cite[\S16, Thm.~16.1 and (16.1.4)--(16.1.8)]{jacquet-langlands-1970}

\item For \(G=\operatorname{Res}_{E/\Q}\GL_d\), remove the positive real split center, fix an open compact subgroup at the finite places and an irreducible maximal compact type at infinity. For each real bound on the absolute value of the Casimir eigenvalue, the sum of the multiplicities of that level and type in the discrete spectrum under the bound is finite, including the residual spectrum. This is the finiteness consequence of the cuspidal and residual counting estimates of \cite[Cor.~1.3 and Thm.~1.4, (1.17)--(1.18)]{mueller-2024}. In Lemma~\ref{B6.04u}, \(E\) is CM, the list of compact types is finite, and the Casimir is a fixed polynomial on the bounded set of infinitesimal parameters.

\item In Langlands's spectral decomposition for \(\GL_r\) over a number field, with the positive split central character fixed and unitary, the automorphic \(L^2\) space is the Hilbert sum, over associate classes of Levi subgroups, of the Eisenstein direct integrals formed from their discrete automorphic spectra. The inductions are normalized. Their parameters run over the imaginary dual of the real split center of the Levi modulo that of the group, with the Weyl identifications given by normalized intertwining operators. For a Levi \(\prod_i\GL_{r_i}\), these parameters twist the factors by \(|\det|^{it_i}\), \(t_i\in\R\), subject to the fixed central character. In Proposition~\ref{B6.04f} the block convolutions act in the regular automorphic \(L^2\) representation used in \ref{B6.04fa}. The decomposition applies to that representation before extraction of generalized eigenfunctions. A nonzero spectral projection for the commuting block operators is nonzero on a set of fibers of positive measure. Thus the spectral characters used for approximation may be taken from the unitary inducing parameters. \cite[\S7, Thm.~7.7]{langlands-1976}\cite[Thms.~6.1 and~7.8]{delorme-2023}
\end{enumerate}

The Galois attachment used here is Theorem~\ref{B9.02}. For comparison, Chenevier and Harris take a cuspidal \(\Pi\) over a CM field with \(\Pi^c\simeq\Pi^\vee\) and nonzero relative Lie algebra cohomology with irreducible algebraic coefficients. Their Theorem~3.2.3 gives a continuous semisimple \(\ell\)-adic representation. Away from \(\ell\) its Frobenius-semisimplified Weil--Deligne representation has the semisimple Weil representation of the local parameter of \(\Pi_v|\det|_v^{(1-n)/2}\), with monodromy bounded by that parameter in the order of part~(a\(^{\prime}\)). Above \(\ell\) it is de Rham, with filtration jumps \(\mu_i+n-i\), \(1\le i\le n\), for the highest weight at the corresponding embedding in the authors' coefficient convention. It is crystalline at spherical places, with the characteristic polynomial of crystalline Frobenius prescribed by the same normalized local parameter, and semistable at Iwahori-spherical places. The witness in Chapter~\ref{ch:vwin-automorphic} uses the cuspidal hypothesis. Its local comparisons are those of \ref{B9.04}--\ref{B9.06}. \cite[Hyp.~1.1, (1.6), Thm.~3.2.3]{chenevier-harris-2013}
\end{remark}

\begin{proposition}[Assembling the four groups of inputs]\label{B6.05}
With the published inputs of Remark~\ref{B6.05ca}, the representation assertion $\mathcal R$ of Theorem~\ref{B9.02} implies the window bound of Definition~\ref{B6.03}.  
\end{proposition}
\begin{proof}[Outline]
Combine \(\mathcal R\) with \ref{B6.04a}, \ref{B6.05a}, \ref{B6.05b}, \ref{B6.05c}, \ref{B6.05d}, \ref{B6.05e}, and \ref{B6.05f}. Use \ref{B9.36a} for weak descent in ranks \(2,3,4\), then apply \ref{B6.06} to obtain \ref{B6.03}. The complete proof is in \S\ref{prf:B6.05}.
\end{proof}

\begin{lemma}[Finite multiplicity in a fixed cuspidal constituent]\label{B6.06a}
Let \(F\) be totally real, let \(\pi\) be an irreducible cuspidal constituent for \(\GL_a/F\), and let \(U\) be an open subgroup of the finite adelic group. Put \(K_\infty=\prod_{v\mid\infty}\mathrm O(a)\). For a finite-dimensional complex representation \(\xi\) of \(K_\infty\), the images of all \(K_\infty\)-equivariant maps \(\xi\to\pi^U\) lie in a common finite-dimensional subspace of \(\pi\).
\end{lemma}
\begin{proof}[Outline]
Use the Capelli operators to fix the central ideal and quotient by the common evaluation kernel to get finitely many compact types. Apply Harish-Chandra finiteness at a compact open level with those types and central ideal, then intersect with \(\pi\). The complete proof is in \S\ref{prf:B6.06a}.
\end{proof}

\begin{proposition}[Coefficients, conjugation and the constituent count]\label{B6.06b}
Let \(E\) be totally complex, \(a\ge1\), \(G=\GL_a\), and \(\Xi\) an irreducible algebraic coefficient representation for \(\operatorname{Res}_{E/\Q}G\) over \(\C\).
\begin{enumerate}
\item A cuspidal \(\pi\) whose infinitesimal multiset at every embedding is the negative of that of \(\Xi\) has nonzero relative Lie algebra cohomology with coefficients in \(\Xi\).
\item For a regular algebraic cuspidal \(\pi\) and \(\sigma\in\operatorname{Aut}(\C)\), there is a cuspidal \(\pi'\) with finite part the scalar conjugate \({}^{\sigma}\pi_f\). Its infinitesimal multiset at \(\sigma\tau\) equals that of \(\pi\) at \(\tau\).
\item For a compact open \(U\subset G(\mathbf A_{E,f})\), put \(Y=G(\mathbf A_{E,f})/U\) and give \(\operatorname{Map}(Y,\Xi)\) the action \((\gamma f)(y)=\Xi(\gamma)f(\gamma^{-1}y)\) of \(G(E)\). If \(S\) is a finite set of pairwise nonisomorphic cuspidal representations, all cohomological for this same \(\Xi\), and each has a nonzero \(U\)-fixed finite vector, then
\[
          |S|\le\sum_{j\ge0}
          \dim_{\C}H^j\bigl(G(E),\operatorname{Map}(Y,\Xi)\bigr).
\]
\end{enumerate}
\end{proposition}
\begin{proof}[Outline]
Use Clozel's cohomology calculation and conjugation theorem. Inject independent tensors of archimedean classes and finite fixed vectors from the cuspidal summands, and identify the target by Shapiro's lemma on the finite-level orbits. The complete proof is in \S\ref{prf:B6.06b}.
\end{proof}

\begin{lemma}[Patching at split completions]\label{B6.06c}
Let \(F\) be totally real, \(G_F=\Gal(\overline F/F)\), and let \(\mathcal L\) be a family of quadratic CM extensions in \(\overline F\). Suppose that for every finite place \(v\) there are infinitely many members split at \(v\). For \(L\in\mathcal L\) let \(r_L:G_L\to\GL_a(\overline{\Q}_\ell)\) be continuous and semisimple. Suppose \(r_L\) is invariant up to isomorphism under \(\Gal(L/F)\), and \(r_L,r_{L'}\) have isomorphic restrictions to \(G_{LL'}\) for each pair. There is a continuous semisimple \(r\) of \(G_F\) with \(r|_{G_L}\simeq r_L\) for all but finitely many \(L\).
\end{lemma}
\begin{proof}[Outline]
Given finite \(D/F\) and \(v\), pick a member of \(\mathcal{L}\) split at \(v\) outside the finitely many quadratic subfields of the Galois closure of \(D\) to get linear disjointness. Sorensen's prime-degree patching lemma then yields \(r\) from invariance and compatibility. The complete proof is in \S\ref{prf:B6.06c}.
\end{proof}

\begin{proposition}[Totally real and cyclic inputs]\label{B6.06d}
\begin{enumerate}
\item For cuspidal representations of positive rank over a totally real field \(F\), equality of Satake multisets almost everywhere implies isomorphism. Inverse transpose gives contragredient duality, and finite-level compact-type admissibility holds in the form of \ref{B6.06a}. Let \(f\) be a smooth cuspidal automorphic form of moderate growth, fixed by an open compact subgroup at the finite places. Assume that its translates by the maximal compact group at infinity span a finite-dimensional space, that \(Z(U(\mathfrak g_{\infty,\C}))f\) is finite dimensional, and that the Lie algebra of the real split centre acts on \(f\) by a scalar. Then \(f\) is a finite sum of vectors in irreducible cuspidal constituents.
\item For a quadratic totally complex extension \(L/F\) of a totally real field and a regular algebraic cuspidal \(\pi\) on \(\GL_2(\mathbf A_F)\), either \(\pi\) has the self-twist by the quadratic character of \(L/F\), or its base change is cuspidal. The lift has the same infinitesimal multiset over each real place and copies the local representation at a split place. At common unramified places its Satake roots are the residue-degree powers of the roots downstairs.
\item For a normalized cuspidal newform \(f\) of weight \(h\ge2\) and level \(M\), and \(j\in\Z\), there is a cuspidal representation of \(\GL_2(\mathbf A_\Q)\) of infinitesimal character \(\{j+\tfrac12,j+\tfrac32-h\}\). At every \(p\nmid M\) it is spherical, and its Satake multiset \(B_p\) satisfies
\[
      \prod_{\alpha\in B_p}(X-p^{j+1/2}\alpha)
          =X^2-a_p(f)X+\omega_f(p)p^{h-1},
\]
where \(\omega_f\) is the nebentype.
\item Let \(L/K\) be cyclic of prime degree with \(K\) totally complex and \(a\ge1\). A regular algebraic cuspidal representation of \(\GL_a(\mathbf A_K)\) has cuspidal base change in the absence of a nontrivial self-twist by a character of \(\Gal(L/K)\), viewed by reciprocity as a Hecke character. A regular algebraic cuspidal representation over \(L\) whose Satake multisets agree at conjugate places almost everywhere descends to a cuspidal representation over \(K\). Infinitesimal characters pull back on embeddings and the Satake rule is the residue-degree power rule. Two cuspidal descents of the same cuspidal system differ by a character of \(\Gal(L/K)\), including the trivial character.
\item If \(K/\Q\) is imaginary quadratic, a cuspidal \(\GL_2(\mathbf A_K)\) representation with distinct half-integral infinitesimal weights of integer difference at least two at each embedding, and with the cofinite Satake invariance in (4), descends with the given infinitesimal weights and with Satake matching at almost every split rational prime.
\item Over every totally real field of even degree there is a totally definite quaternion algebra split at every finite place. Its finite-place matrix identifications may be chosen integral almost everywhere.
\end{enumerate}
\end{proposition}
\begin{proof}[Outline]
Use strong multiplicity one for duality, \ref{B6.06a} for admissibility, and Harish-Chandra finiteness for finite cuspidal decomposition. Apply cyclic base change for lifting and descent, adelize and norm-twist the newform, and prescribe local quaternion invariants. The complete proof is in \S\ref{prf:B6.06d}.
\end{proof}

\begin{proposition}[The totally real rank-two attachment]\label{B6.06e}
Let \(F\) be totally real and \(\pi\) regular algebraic cuspidal on \(\GL_2(\mathbf A_F)\). Suppose \(\pi^\vee\simeq\pi\otimes\chi^{-1}\), where the Hecke character \(\chi\) has the same sign at \(-1\) at all real places. Fix \(\ell\) and \(\iota:\overline{\Q}_\ell\simeq\C\). There is a continuous semisimple \(r:G_F\to\GL_2(\overline{\Q}_\ell)\) with the following properties.
\begin{enumerate}
\item At every spherical \(v\nmid\ell\), \(r\) is unramified and its inverse arithmetic Frobenius roots are \(q_v^{1/2}B_v\) under \(\iota\), where \(B_v\) is the Satake multiset.
\item If the cuspidal support at \(v\) consists of characters with no pair \(\xi,\xi|\cdot|_v\), the local Weil--Deligne comparison is with the parameter of \(\pi_v|\det|_v^{-1/2}\) for every correspondence of \ref{B9.47}. Away from \(\ell\) the absolute values of the eigenvalues of each degree-one Weil element satisfy \(|\alpha|<q_v|\beta|\) for every pair. At \(\ell\) the representation is potentially semistable and de Rham there, and it is crystalline at a spherical place. The spherical components have the indicated unlinked character support.
\item At a place above \(\ell\) satisfying the support condition in (2), the jumps at \(\tau\) are the multiset \(\{\tfrac12-z:z\in m_{\iota\tau}\}\). There is one integer \(w\) for which each such multiset, and the multiset defined by the same formula at every embedding, is invariant under \(j\mapsto w-j\).
\end{enumerate}
\end{proposition}
\begin{proof}[Outline]
Take cuspidal CM lifts by \ref{B6.06d}(2), attach by \ref{B9.02} and patch by Chebotarev and \ref{B6.06c}. Use split completions for \ref{B9.05} and \ref{B9.06}, the rank-two unitary classification for the root bound and \ref{B9.04} for weight symmetry. The complete proof is in \S\ref{prf:B6.06e}.
\end{proof}

Write \(\mathscr H\) for \(h_N\), the weak unitary descent of \ref{B9.36a} in the ranks and at the levels specified in Proposition~\ref{B6.05}, and the assertions of \ref{B6.06a}, \ref{B6.06b}, \ref{B6.06d} and \ref{B6.06e}. In the rank-two descent entry one may instead supply directly the conclusion of the next lemma under its Galois hypotheses.

\begin{lemma}[Descent of a specified rank-two representation]\label{B6.06f}
Use the cyclic and split-prime inputs of \ref{B6.06d} and the attachments of \ref{B9.02} and \ref{B6.06e}. Let \(L/\Q\) be abelian CM and contain the fixed field \(K\) of an odd quadratic character. Let \(W:G_\Q\to\GL_2(\overline{\Q}_\ell)\) be continuous with determinant factoring through a finite cyclotomic extension. Suppose \(W|_{G_L}\) is irreducible and attached, via \(\iota\), to a regular algebraic cuspidal representation. Then \(W\) has a regular algebraic cuspidal attachment over \(\Q\).
\end{lemma}
\begin{proof}[Outline]
Descend by \ref{B6.06d}(4), polarize using its fibre characters and attach by \ref{B9.02}, then use Chebotarev and the intertwiner line to choose each twist. Over \(\mathbb Q\), use split-prime descent by \ref{B6.06d}(5) and attach by \ref{B6.06e} before the final twist. The complete proof is in \S\ref{prf:B6.06f}.
\end{proof}

\begin{proposition}[From automorphic inputs to counting]\label{B6.06}
Assume $h_N$ of Definition~\ref{def:p6c1-classical-inputs} and weak unitary descent in ranks two through four as in Lemma~\ref{B9.36a}. Then the additional hypotheses used in counting hold. On the cuspidal spectrum they are cohomological realization, conjugation of finite Hecke data and the bound by cohomology dimensions. Over totally real fields they include rank-two Galois attachment, strong multiplicity one, contragredient duality, admissibility, cuspidal semisimplicity and quadratic base change. The remaining hypotheses concern the representation attached to a newform, cyclic base change of prime degree over a totally complex field, rank-two descent and the existence over every totally real field of even degree of a totally definite quaternion algebra split at all finite places. These assertions have the forms of \ref{B6.06a}, \ref{B6.06b}, \ref{B6.06d} and \ref{B6.06e}. The descent for a specified Galois representation is \ref{B6.06f}. Together with $h_N$ and weak unitary descent, these assertions imply the window bound of Definition~\ref{B6.03}.
\end{proposition}
\begin{proof}[Outline]
Assemble the inputs from \ref{B6.06a}, \ref{B6.06b}, \ref{B6.06d}, \ref{B6.06e} and \ref{B6.06f}. Combine the small-twist bound \ref{B6.15} with the counts of \ref{B6.30}, \ref{B6.31} and \ref{B6.34} in \ref{B6.07} to obtain the window bound. The complete proof is in \S\ref{prf:B6.06}.
\end{proof}

\begin{proposition}[Four counts suffice]\label{B6.07}
Assume $h_N$ of Definition~\ref{def:p6c1-classical-inputs}, the finiteness and $H^2$ estimate of Theorem~\ref{B6.15} for some real constant $C_K$, and the four counts of \ref{B6.30}, \ref{B6.31} and \ref{B6.34}. Then the window bound of Definition~\ref{B6.03} holds.
\end{proposition}
\begin{proof}[Outline]
Apply \ref{B6.27} at \(C_t=401\), getting \ref{B6.18} from \ref{B6.17} and Ferrero-Washington vanishing. Obtain the odd-prime plus bound via \ref{B6.39} using \ref{B6.32}, \ref{B6.35}, \ref{B6.36}, the minus bound \ref{B6.13} by Poitou-Tate duality, and combine via \ref{B6.14}. The complete proof is in \S\ref{prf:B6.07}.
\end{proof}






% ---- p6c2 ----
\chapter{Three length bounds}\label{ch:vwin-three-length-bounds}
Write $s_\ell(E,k)=s_{\ell,k}(E)$ as in \ref{B6.01} and $n=[E:\Q]$.

\section{Complex conjugation and arithmetic duality}
\begin{proposition}[Splitting and addition of the bounds]\label{B6.08}
For CM $E$ and $k\in\N$, let $s_\ell^\pm(E,k)$ be the lengths of the
$(-1)^k,(-1)^{k+1}$ eigenspaces $\Sha_{\ell,k}^\pm(E)$ of complex
conjugation. For odd $\ell$, $s_\ell=s_\ell^++s_\ell^-$. Fix $t\in\N$.
For all local data in \ref{B6.01}, suppose $s_2$ and both odd lengths
are finite for $1\le k<t$. If for every finite set $T$ of odd primes
\[
\sum_{1\le k<t}s_2(E,k)\le B_2,\qquad
\sum_{\ell\in T}\log\ell\sum_{1\le k<t}s_\ell^\pm(E,k)\le B_\pm,
\]
then $E$ satisfies the total bound of Definition~\ref{B6.01} up to $t$ with bound $B_2\log 2+B_-+B_+$.
\end{proposition}
\begin{proof}
Conjugation $c$ permutes the places above $\ell$. At odd $\ell$,
$(1\pm(-1)^kc)/2$ split the kernel. Add lengths with weights $\log\ell$.
\end{proof}

\begin{theorem}[Finite-coefficient Poitou--Tate duality]\label{B6.09}
Let $K$ be a number field, $\ell$ prime and $r\ge1$, with $\ell\ne2$
or $K$ totally complex. Put $R=\Z/\ell^r\Z$. For each primitive
$\zeta\in\mu_{\ell^r}$ there exists a conjugation-compatible family
of isomorphisms $\mathrm{inv}_v:H^2(K_v,\mu_{\ell^r})\simeq R$
at finite places and their lifts to $\overline K$. The cup of a
uniformiser's Kummer class with the unramified character sending
arithmetic Frobenius to $1$ has invariant $1$.

Fix this family. Let $S$ be a finite set of finite places containing
all places above $\ell$, and choose a lift of each $v\in S$ to $\overline K$.
For any finite continuous $R[G_{K,S}]$-module $A$, write
$A^\dagger=\mathrm{Hom}_R(A,\mu_{\ell^r})$, $D_R=\mathrm{Hom}_R(-,R)$,
$P^i(A)=\bigoplus_{v\in S}H^i(K_v,A)$ and $H_G^i=H^i(G_{K,S},-)$.
Then
\[
\begin{aligned}
0&\to H_G^0(A)\to P^0(A)\to D_RH_G^2(A^\dagger)\\
 &\to H_G^1(A)\to P^1(A)\to D_RH_G^1(A^\dagger)\\
 &\to H_G^2(A)\to P^2(A)\to D_RH_G^0(A^\dagger)\to0
\end{aligned}
\]
is exact. The duality maps sum the local cup invariants, with the
restriction of the dual class first for $i=1$ and second for $i=0,2$.
The localization kernels $\Sha_S^i(A)$ are finite for $i=1,2$,
$\Sha_S^2(A)\simeq D_R\Sha_S^1(A^\dagger)$, and $H_G^i(A)=0$ for $i\ge3$.
\end{theorem}
This is the finite-coefficient form of Milne's nine-term sequence
in the range where the archimedean Tate terms vanish
\cite[Theorem~I.4.10]{milne-2006}.

\section{The negative component}
\begin{lemma}[Bernoulli numbers at a cyclotomic level]\label{p6c2:bernoulli}\label{B6.09a}
Let $\theta$ be a Dirichlet character modulo $N\ge1$, induced by
the primitive character $\chi$ of conductor $f$, and let $j\ge2$.
For $\eta$ modulo $m$, set $B_{j,\eta}=m^{j-1}\sum_{a=1}^m\eta(a)B_j(a/m)$,
where $B_j(X)$ is the Bernoulli polynomial. Put
$\mathrm{primorial}(h)=\prod_{p\le h,\ p\ {\rm prime}}p$. Then
$A_{N,j}(\theta)=N\mathrm{primorial}(j+1)B_{j,\chi}$ is an algebraic integer.
For $f\mid H$, the character $\chi_H$ induced by $\chi$ modulo $H$ satisfies
\[
B_{j,\chi_H}=B_{j,\chi}\prod_{\substack{p\mid H\\p\nmid f}}
(1-\chi(p)p^{j-1}).
\]
If $\theta(-1)=(-1)^j$, both $B_{j,\chi}$ and $B_{j,\theta}$ are nonzero.
For each complex embedding $\sigma$,
\[
 \max\bigl\{|\sigma A_{N,j}(\theta)|,\,
 N\mathrm{primorial}(j+1)|\sigma B_{j,\theta}|\bigr\}
 \le b(N,j):=16N^{j+1}j!.
 \tag{6.1}\label{p6c2:bern}
\]
\end{lemma}
\begin{proof}[Outline]
Von Staudt–Clausen clears denominators, and Bernoulli distribution gives the level factors. The functional equation, Euler product and level factors give nonvanishing, while Fourier and binomial estimates give the bound. The complete proof is in \S\ref{prf:B6.09a}.
\end{proof}

\begin{theorem}[A Bernoulli bound for one character]\label{B6.10}
Let $E\subseteq\Q(\mu_N)$ be CM, $\ell$ odd prime, $\ell\nmid n$,
and $O\subseteq\overline{\Q}_\ell$ a DVR finite over $\Z_\ell$ with
fraction field $K_O$. For $k\ge1$, let $\theta:\Gal(E/\Q)\to O^\times$
have conjugation sign $(-1)^{k+1}$ and level-$N$ pullback $\widetilde\theta$.
With $\varepsilon_\ell$ cyclotomic, let $C_\theta$ be the kernel of
restriction to finite inertia and decomposition at $\ell$ in
$H^1(\Q,(K_O/O)(\theta^{-1}\varepsilon_\ell^{-k}))$. Restriction identifies
$C_\theta$ with the $\theta$-part of $\Sha_{\ell,k}(E)\otimes_{\Z_\ell}O$.
Normalize $v_O$ by $v_O(\pi)=1$ for a uniformiser $\pi$. Then
\[
\mathrm{length}_O C_\theta\le v_O(A_{N,k+1}(\widetilde\theta)).
\tag{6.2}\label{p6c2:one-char}
\]
\end{theorem}
\begin{proof}[Outline]
Inflation–restriction identifies the character part. Jacobi-sum descent uses the level formula in \ref{p6c2:bernoulli} and cup products via \ref{B6.05d}(2) and \ref{B6.09}, while Kummer theory gives vanishing in the residual cyclotomic case. The complete proof is in \S\ref{prf:B6.10}.
\end{proof}

\begin{lemma}[Cyclotomic annihilation on unramified homomorphisms]\label{B6.10a}
Let $N\ge1$, $K=\Q(\mu_N)$ and $G=(\Z/N\Z)^\times$, with $\tau_b$
acting by $b$-th powers on $\mu_N$. For an odd positive integer $a$
prime to $N$, set
\[
j_b=\frac{a[a^{-1}b]_N-[b]_N}{N},\qquad E_b=\frac{a-1}{2}-j_b .
\]
For each ideal class $c$, written additively,
$\sum_{b\in G}E_b\tau_b^{-1}c=0$. If $N>2$ and $h:G_K\to X$ is a
locally constant homomorphism to an abelian group, zero on finite
inertia, then for every $u\in G_K$ and lifts $\sigma_b$ of $\tau_b$,
\[
\sum_{b\in G}E_bh(\sigma_b^{-1}u\sigma_b)=0.
\]
\end{lemma}
\begin{proof}[Outline]
Use the Jacobi-sum valuations in \ref{p6c2:jacobi} to show the carry products are principal, then telescope and translate to obtain the required principal ideal. Inverse Artin reciprocity transfers the identity to the unramified homomorphism. The complete proof is in \S\ref{prf:B6.10a}.
\end{proof}

\begin{lemma}[Summation away from the degree]\label{B6.11}
For CM $E\subseteq\Q(\mu_N)$, $N,k\ge1$, $s_\ell^-(E,k)$ is finite
at odd primes $\ell\nmid n$. For finite sets $T$ of these primes,
\[
\sum_{\ell\in T}s_\ell^-(E,k)\log\ell\le n\log b(N,k+1).
\tag{6.8}\label{p6c2:away-degree}
\]
\end{lemma}
\begin{proof}
Take the product $D$ of $A_{N,k+1}(\widetilde\theta)$ over characters
of $\Gal(E/\Q)$ of sign $(-1)^{k+1}$, where $\widetilde\theta$ is the
level-$N$ pullback. By \ref{p6c2:bernoulli}, $D$ is a Galois-invariant
nonzero algebraic integer and $\log|D|\le n\log b(N,k+1)$.
Sum \eqref{p6c2:one-char} in a splitting $O$ and divide by $v_O(\ell)$
to get $s_\ell^-\le v_\ell(D)$. Summing over $T$ proves the bound.
\end{proof}

\begin{proposition}[Primes dividing the degree]\label{B6.12}
Let $E\subseteq\Q(\mu_N)$ be CM, $N,k\ge1$, and $\ell$ odd prime.
For some $r\in\N$,
\[
\ell^r\Sha^-_{\ell,k}(E)=0,\qquad r\log\ell\le30n(k+1)\log(k\ell N).
\tag{6.9}\label{p6c2:exponent}
\]
The length is finite and
\[
s_\ell^-(E,k)\log\ell\le240\ell n^2(k+1)\log^2(k\ell N).
\tag{6.10}\label{p6c2:at-degree}
\]
For $M=(m,H_M)$, $m>0$, $H_M\le(\Z/m\Z)^\times$, put
$K_M=\Q(\mu_m)^{H_M}$ and $o=[K_M:\Q]$. Take primes
$q\equiv3\bmod8$, $d\equiv7\bmod16$ with $\gcd(qd,m)=1$.
For $a\in\N$, write $E_{M,a,q,d}$ for the compositum $K_M\,\Q(\mu_{2^{a+2}})^+\Q(\sqrt{-q},\sqrt{-d})$.
If $E\simeq_\Q E_{M,a,q,d}$, every odd prime dividing $n$ divides $o$.
\end{proposition}
\begin{proof}[Outline]
Bernoulli distribution and \ref{B6.10a} give the annihilator, and the adjugate with \ref{p6c2:bernoulli} and \ref{p6c2:det} gives \ref{p6c2:exponent}. Kummer theory, \ref{p6c2:class}, and torsion filtration give \ref{p6c2:at-degree}, and a Galois embedding controls odd primes. The complete proof is in \S\ref{prf:B6.12}.
\end{proof}

\section{The negative budget and the common window}
\begin{theorem}[Uniform negative budget]\label{B6.13}
Take $M,o,q,d$ as in Proposition~\ref{B6.12}, integers $a,t\ge2$, and a CM field $E\simeq_\Q E_{M,a,q,d}$. Let $F$ be its conductor, the
least level with $E\hookrightarrow\Q(\mu_F)$. The lengths
$s_\ell^-(E,k)$ are finite for odd $\ell$ and $1\le k<t$.
For each finite set $T$ of odd primes,
\[
\begin{split}
\sum_{\ell\in T}\log\ell\sum_{1\le k<t}s_\ell^-(E,k)
 &\le B_-(240,o,n,F,t),\\
B_-(C,o,n,F,t)
 &=3nt^2\log(tF)+C\sum_{\substack{\ell\mid o\\\ell\ {\rm odd\ prime}}}
 \ell n^2t^2\log^2(t\ell F).
\end{split}\tag{6.11}\label{p6c2:minus}
\]
For real $C$ and $o\in\N$, put
$K=3+|C|\sum_{\ell\mid o,\ \ell\ {\rm odd\ prime}}\ell^3$
with empty sum at $o=0$. If $F,t\ge1$, $n\le F$, and $\log(Ft)\le2W$,
then
\[
B_-(C,o,n,F,t)\le Ke^{8W}.
\tag{6.12}\label{p6c2:minus-exp}
\]
\end{theorem}
\begin{proof}
Sum \eqref{p6c2:away-degree} at $N=F$ using
$\log b(F,k+1)\le(2t+5)\log(tF)$ and $(t-1)(2t+5)\le3t^2$.
Sum \eqref{p6c2:at-degree} at odd $\ell\mid n$, all of which divide
$o$ by \ref{B6.12}. For $u=Ft\ge1$, $nt\le u$, so
\[
3nt^2\log u\le3u^4,\qquad
\ell n^2t^2\log^2(\ell u)\le\ell^3u^4,\qquad u^4\le e^{8W}.
\]
\end{proof}

\begin{definition}[Joint window condition]\label{p6c2:joint-window}
Fix $C_t\in\R$. The joint window condition at exponent $C_t$ requires the following for every primitive Dirichlet character $\psi$ of conductor $m$, every $0<\varepsilon<1/52$, every $A_1\in\R$, and all functions $a,t:\R\to\N$ satisfying
\[
 a(x)\ge2,\qquad 2^{a(x)}\le x^{A_1},\qquad 2\le t(x)\le x^{C_t}
 \qquad(x\ge2).
\]
There exist real constants $J_0>0$, $C$, $P$ and $X_1$ with the following property. Let $X\ge X_1$ and let $q\equiv3\pmod8$ and $d\equiv7\pmod{16}$ be primes with $\gcd(qd,m)=1$. If
\[
 \min(q,d)>\max(2X^{C_t}+1,P),\qquad
 \log(qd)\le X^{\varepsilon/(4J_0)},
\]
there is a $Y\ge X^{1/J_0}$ with $Y^2\le X$ such that both bounds below hold throughout $[Y,Y^2]$.

For each $x$ in this interval and every field $E$ associated with $(\psi,a(x),q,d)$ as in Definition~\ref{B6.02}, put $D_x=\exp(x^{4\varepsilon/5})$ and let $F$ be the conductor of $E$. For every choice of local data as in Definition~\ref{B6.01}, the lengths $s_2(E,k)$ and $s_\ell^+(E,k)$ for odd primes $\ell$ are finite for $1\le k<t(x)$. For every finite set $T$ of odd primes,
\[
 \sum_{1\le k<t(x)}s_2(E,k)\le D_x,\qquad
 \sum_{\ell\in T}\log\ell\sum_{1\le k<t(x)}s_\ell^+(E,k)
 \le (Ft(x))^C+CD_x .
\]
\end{definition}

\begin{theorem}[Absorption into the window bound]\label{B6.14}
Let $C_t\ge1$. The joint window condition of Definition~\ref{p6c2:joint-window}, supplied by Theorems~\ref{B6.27} and~\ref{B6.39}, and the negative bound of Theorem~\ref{B6.13} imply the window assertion of Definition~\ref{B6.03} for every primitive Dirichlet character $\psi$ with exponent $C_t$.
\end{theorem}
\begin{proof}[Outline]
Under \ref{B6.03} and \ref{p6c2:page}, use \ref{B6.02} for \ref{p6c2:pair} and the conductor bound \ref{p6c2:log}. Combine \ref{p6c2:minus} and \ref{p6c2:minus-exp} with \ref{eq:B6-two-sum} and \ref{eq:b6-pluswindow} via \ref{B6.08} and absorb the smaller powers of \(x\). The complete proof is in \S\ref{prf:B6.14}.
\end{proof}


% ---- p6c3 ----
\chapter{The cyclotomic tower and a common window}\label{ch:B6-cyclotomic-window}

For a number field $L$, a prime $\ell$ and $k\in\Z$, let $S_\ell$ be the places of $L$ above $\ell$. Put
\[
 s_\ell(L,k)=\operatorname{length}_{\Z_\ell}\Sha_{\ell,k}(L),
 \qquad
 V_\ell(L,k)=\operatorname{length}_{\Z_\ell}H^2(G_{L,S_\ell},\Z_\ell(k+1)).
\]
Both lengths may be infinite. The group $\Sha_{\ell,k}(L)$ is defined in \ref{B6.01}. Write $D_x=\exp(x^{4\varepsilon/5})$ for $x,\varepsilon>0$ and $j_0(\ell,R)=\lfloor\log_\ell R\rfloor+3$, with $\lfloor\log_\ell0\rfloor=0$. A cyclotomic layer $L_m/L_0$ of index $m$ has degree $\ell^m$. Let $h_N$ be the data of Proposition~\ref{B6.05}. We use its local finiteness and its Poitou--Tate and local Tate dualities.

\section{A bound for small twists}

\begin{theorem}[The small-twist $H^2$ bound]\label{B6.15}
There is $C_K\ge1$ such that whenever $L$ is totally complex and contains $\mu_\ell$ for odd $\ell$ or $i$ for $\ell=2$, for every $k\ge1$ and choice of local data the group $H^2(G_{L,S_\ell},\Z_\ell(k+1))$ is finite. With $n=[L:\Q]$ and $D_L=\operatorname{disc}(L)$,
\begin{equation}\label{eq:B6-small-H2}
 \log\#H^2(G_{L,S_\ell},\Z_\ell(k+1))
 \le\exp\!\left((2kn)^{C_K}\log(2|D_L|)\right).
\end{equation}
\end{theorem}

\begin{proof}[Outline]
Use finite coefficients, local Tate duality and \ref{B6.09} to compare localization kernels via \eqref{eq:B6-H2-local-correction}. Bound local invariants by cyclotomic indices and \eqref{eq:B6-lte}, then absorb their contribution using \eqref{eq:B6-small-Sha}. The complete proof is in \S\ref{prf:B6.15}.
\end{proof}

\begin{theorem}[Class-group counting for a small twist]\label{B6.16}
There is an absolute $C$ such that for every prime $\ell$, every totally complex $L$ with the roots of unity of \ref{B6.15}, and every $k\ge1$,
\begin{equation}\label{eq:B6-small-Sha}
 s_\ell(L,k)<\infty,\qquad
 s_\ell(L,k)\log\ell\le\exp\!\left((2k[L:\Q])^C\log(2|D_L|)\right).
\end{equation}
Let $F/L$ be cyclic over a totally complex $L$, with generator $\sigma$ of order $\ell^{\nu-c}$. Here $\ell$ is prime, $1\le c\le\nu$, and $c\ge2$ if $\ell=2$. Suppose $F$ contains a primitive $\ell^\nu$th root $\zeta$ satisfying $\sigma(\zeta)=\zeta^{1+\ell^c}$. Fix $k\ge1$ and $b,B\in\N$. Put
\[
 h=\operatorname{ord}(\sigma),\qquad a=(1+\ell^c)^k,\qquad
 \mathcal N_a(\alpha)=\prod_{j=0}^{h-1}\sigma^j(\alpha)^{a^j}.
\]
Let $\mathcal E_{\nu,k}$ be the functions $\varphi$ from the finite primes of $F$ to $\Z/\ell^\nu\Z$ satisfying
\[
 \sum_P\operatorname{ord}_P(\alpha)\varphi(P)=0
 \quad(\alpha\in F^\times),\qquad \varphi(\sigma^{-1}P)=a\varphi(P).
\]
Suppose every finite set of $\alpha\in F^\times$ whose $\mathcal N_a(\alpha)^{\ell^b}$ are distinct modulo $F^{\times\ell^\nu}$ has size at most $B$ if each $\alpha$ admits a divisor $\mathfrak d_\alpha$ with
\[
 \operatorname{ord}_P(\alpha)=\mathfrak d_\alpha(\sigma P)-a\mathfrak d_\alpha(P)
 \qquad(\text{all }P).
\]
Then $\mathcal E_{\nu,k}$ is finite and
\begin{equation}\label{eq:B6-finite-class-count}
 \#\mathcal E_{\nu,k}\,(\ell^\nu)^{r_2(L)}
 \le B\,\ell^{(b+v_\ell(k))g(L)+c+v_\ell(k)},\qquad
 g(L)=2[L:\Q]+\lfloor\log_2 h_L\rfloor+2.
\end{equation}
Here $h_L$ is the class number and $r_2(L)$ counts complex places.
\end{theorem}

\begin{proof}[Outline]
Reduce to weighted norms by the snake lemma, unit determinant and Hilbert 90. Separate classes by bar cycles and Chebotarev, bound Bockstein torsion by \ref{B6.16a} and \ref{B6.16c} and ranks by Borel stability, then bound Sha via \ref{B6.20} and finite-layer descent. The complete proof is in \S\ref{prf:B6.16}.
\end{proof}

\begin{lemma}[Integral stability over the ring of integers]\label{B6.16a}
Let $L$ be totally complex and $G_r=\GL_r(\mathcal O_L)$. For $i,r,r'\in\N$ with $2i+4\le r\le r'$, adjoining an identity block induces an isomorphism $H_i(G_r,\Z)\simeq H_i(G_{r'},\Z)$. The torsion at rank $r'$ is finite and has order $\#H_i(G_r,\Z)_{\mathrm{tors}}$.
\end{lemma}

\begin{proof}[Outline]
Use the Chinese remainder theorem to bound the Bass stable rank by two, then compose van der Kallen's integral stability isomorphisms in positive degrees. In degree zero stabilization is the identity, and \ref{B6.16c} gives torsion finiteness. The complete proof is in \S\ref{prf:B6.16a}.
\end{proof}

\begin{lemma}[An exact complex with bounded orbits]\label{B6.16b}
Let $L$ be totally complex, $n=[L:\Q]$, $D=|D_L|$, $N\in\N$ and $G=\GL_N(\mathcal O_L)$. Put $d=nN$, $I=D^Nd^d$, $R=N^NI^2$ and
\[
 U_p=I^{1+dN}\bigl((2R+1)^d+1\bigr)^{dN(p+1)}.
\]
There are $G$-sets $X_p$ of ordered vertex tuples with deletion faces whose augmented complex $\Z[X_\bullet]\to\Z$ is exact. In degree $p\ge0$ there are at most $U_p$ orbits, and each stabilizer has order at most $3^{nN^2}$. If $N\ge1$,
\begin{equation}\label{eq:B6-orbit-bound}
 U_p\le (2Dd)^{12(p+1)d^3N}.
\end{equation}
\end{lemma}

\begin{proof}[Outline]
Use vertex tuples minimizing a common functional and prove exactness on finite convex nerves by hyperplane cuts. Count orbits through common minimal vectors and active equations using Minkowski and Hadamard, and bound stabilizers by reduction modulo \(3\). The complete proof is in \S\ref{prf:B6.16b}.
\end{proof}

\begin{lemma}[Torsion at a fixed rank]\label{B6.16c}
With the notation of \ref{B6.16b}, let $m\ge0$, $A=U_m$ and $s=3^{nN^2}$. The torsion subgroup of $H_m(G,\Z)$ is finite of order $T\le(m+2)^A s^{mAs^m}$. If $k\ge1$, $1\le N\le4k+6$ and $m\le2k+1$, then
\begin{equation}\label{eq:B6-homology-torsion}
 \log T\le\exp\!\bigl((2kn)^{24}\log(2D)\bigr).
\end{equation}
\end{lemma}

\begin{proof}[Outline]
Shift homology along the exact complex of \ref{B6.16b} and bound stabilizer terms by Shapiro's lemma, transfer and bar ranks. Use Smith normal form and Hadamard for the final coinvariant torsion, then combine the bounds with \eqref{eq:B6-orbit-bound}. The complete proof is in \S\ref{prf:B6.16c}.
\end{proof}

\section{Generator bounds and quadratic extensions}

An abelian datum $M$ specifies $N_M\ge1$ and a subgroup of $(\Z/N_M\Z)^\times$, with fixed field $K_M$ of degree $n_M$. For primes $q,d$ with
\begin{equation}\label{eq:B6-prime-pair}
 q\equiv3\pmod8,\qquad d\equiv7\pmod{16},\qquad \gcd(qd,N_M)=1,
\end{equation}
put
\[
 K_{M,q,d}=K_M(i,\sqrt{-q},\sqrt{-d}),\qquad
 E_{M,a,q,d}=K_M\,\Q(\zeta_{2^{a+2}})^+\Q(\sqrt{-q})\,\Q(\sqrt{-d}).
\]
For a primitive Dirichlet character $\psi$ modulo $f$, write $M_\psi=(f,\ker\psi)$ for the abelian datum whose subgroup consists of the units $u$ modulo $f$ with $\psi(u)=1$. Then $E_{M_\psi,a,q,d}$ is the field associated with $(\psi,a,q,d)$ in Definition~\ref{B6.02}.

For a base $L_0$ with the roots of unity of \ref{B6.15}, put $L_\infty=L_0(\mu_{\ell^\infty})$. Let $\mathcal D_\ell(L_0)$ be the continuous $\Q_\ell/\Z_\ell$-valued characters of $G_{L_\infty}$ unramified everywhere and split above $\ell$. Its compact dual is $\mathrm{XP}=X_\ell(L_0)$. Omitting splitting gives a larger character group with dual $\mathrm{XX}$, so $\mathrm{XX}\twoheadrightarrow\mathrm{XP}$.

\begin{theorem}[Ferrero--Washington in finite layers]\label{B6.17}
Let $\ell$ be prime and $L/\Q$ finite abelian, with $\mu_\ell\subset L$ for odd $\ell$ and $i\in L$ for $\ell=2$. There is $r\in\N$ such that for every $L'=L(\mu_{\ell^n})$, $n\ge0$,
\[
 \#\bigl(\mathrm{Cl}(L')/\mathrm{Cl}(L')^\ell\bigr)\le\ell^r.
\]
\end{theorem}

Ferrero and Washington prove the vanishing of $\mu_\ell$ for cyclotomic $\Z_\ell$-extensions of abelian fields \cite{ferrero-washington-1979}. For these bases, this vanishing is equivalent to a uniform bound on finite-layer $\ell$-ranks.

\begin{theorem}[Uniform numbers of tower generators]\label{B6.18}
For odd prime $\ell$ and an abelian datum $A$ with $\mu_\ell\subset K_A$, some $R_\ell(A)\ge1$ bounds the $\Z_\ell$-generator numbers of both tower modules for every realization of $A$.

For each $M$, some $R_0(M)\ge1$ bounds both generator numbers at $2$ for $K_{M,q,d}$ uniformly in \eqref{eq:B6-prime-pair}. If $g$ bounds the unramified mod-$2$ character dimension in the infinite tower of $K_M(i)$, one can take $R_0(M)=4(g+6n_M)$.
\end{theorem}

\begin{proof}[Outline]
Bound finite-layer characters via \ref{B6.17} and class field theory, use \ref{B6.16} for descent, and apply compact Nakayama. At \(2\), use Frobenius indices for the tame contribution and bound each quadratic step by twice the invariant dimension. The complete proof is in \S\ref{prf:B6.18}.
\end{proof}

\begin{lemma}[The cost of a quadratic extension]\label{B6.19}
If $E\subset E'$ are number fields with $[E':E]\mid2$, then for every $k\in\Z$,
\begin{equation}\label{eq:B6-quadratic-loss}
 s_2(E,k)\le1+s_2(E',k).
\end{equation}
For $a\ge2$ and a prime pair \eqref{eq:B6-prime-pair}, every realization $E$ of $E_{M,a,q,d}$ admits a layer $L_m/K_{M,q,d}$ with $m\le a+2$ and $s_2(E,k)\le1+s_2(L_m,k)$ for every $k\ge1$. The base satisfies
\[
 [K_{M,q,d}:\Q]\le8n_M,\qquad K_{M,q,d}\hookrightarrow\Q(\zeta_{4N_Mqd}).
\]
\end{lemma}

\begin{proof}[Outline]
Restriction preserves Sha; for a quadratic extension, inflation-restriction yields a \(2\)-torsion subquotient of a locally cyclic group, bounding length loss by \(1\). Adjoining \(\mu_{2^{a+2}}\) to \(K_{M,q,d}\) gives \(L_m\) with \([L_m:E]\mid 2\) and the degree and embedding. The complete proof is in \S\ref{prf:B6.19}.
\end{proof}

\section{Twisted cokernels and their centres}

For a topological generator $\sigma$ of $\Gal(L_\infty/L_0)$, put $u=\chi_\ell(\sigma)$ and let $\tau$ be its action on $X_\ell(L_0)$. Then
\begin{equation}\label{eq:B6-unit-level}
 v_\ell(u-1)=c,\qquad c\ge1,\qquad c\ge2\ \text{if }\ell=2,
\end{equation}
with $\mu_{\ell^c}\subset L_0$. For an endomorphism $\tau$ of a $\Z_\ell$-module $X$, write
\[
 w_{\tau,u}(k,m)=\operatorname{length}_{\Z_\ell}X/((u^k\tau)^{\ell^m}-1)X.
\]

\begin{theorem}[Descent to a twisted cokernel]\label{B6.20}
Let $L_0$ contain the roots of unity in \ref{B6.15}. For the tower module and generator above, every layer $L_m/L_0$ and every $k\ge1$ satisfy
\begin{equation}\label{eq:B6-cokernel-descent}
 s_\ell(L_m,k)=w_{\tau,u}(k,m).
\end{equation}
\end{theorem}

\begin{proof}[Outline]
Surjectivity of \(u^{-k\ell^m}-1\) on \((\mathbb{Q}_\ell/\mathbb{Z}_\ell)(-k)\) and cohomological dimension \(1\) kill \(H^1,H^2\) on \(\Gamma_m\) and open subgroups, so Sha is \((\mathcal{D}_\ell(L_0)(-k))^{\Gamma_m}\). Duality gives \((u^k\tau)^{\ell^m}\)-coinvariants. The complete proof is in \S\ref{prf:B6.20}.
\end{proof}

\begin{theorem}[Centres of twisted cokernels]\label{B6.21}
Let $X$ have at most $R$ generators over $\Z_\ell$ and let $\tau$ be a $\Z_\ell$-linear endomorphism. Assume $u,c$ satisfy \eqref{eq:B6-unit-level}. For each $j\ge0$ there are $s_1,\ldots,s_R\in\Z_\ell$ with
\begin{equation}\label{eq:B6-centres-fixed}
 w_{\tau,u}(k,j)\le\sum_{i=1}^R(j+c+v_\ell(s_i-k))
 \qquad(k\in\N),
\end{equation}
where $v_\ell(0)=\infty$. If $\tau$ is unipotent modulo $\ell$, there are $s_1,\ldots,s_R\in\Z_\ell$ with
\begin{equation}\label{eq:B6-centres-all}
 w_{\tau,u}(k,m)\le\sum_{i=1}^R(\max(m,j_0)+c+v_\ell(s_i-k))
 \qquad(k,m\in\N),\qquad j_0=j_0(\ell,R).
\end{equation}
For these centres and $b\in\N$, if $|s_i-k|_\ell\ge\ell^{-b}$ for every $i$, then
\begin{equation}\label{eq:B6-centres-far}
 w_{\tau,u}(k,m)\le R(\max(m,j_0)+c+b).
\end{equation}
\end{theorem}

\begin{proof}[Outline]
Use Smith normal form on a free lift and \eqref{eq:B6-lte} to bound lengths by valuations and find a common centre of nested balls. At \(j_0\) use \eqref{eq:B6-centres-fixed} and the distance bound, then use binomial expansion and Nakayama above and divisibility below. The complete proof is in \S\ref{prf:B6.21}.
\end{proof}

\begin{theorem}[The two tower bounds]\label{B6.22}
Let $L_0$ be totally complex with the roots of unity of \ref{B6.15}, and let $X_\ell(L_0)$ have at most $R$ generators, with $R\ge1$. Use $c$ from \eqref{eq:B6-unit-level} and $j_0=j_0(\ell,R)$. There are $s_1,\ldots,s_R\in\Z_\ell$ such that for every $k\ge1$, $m,b\in\N$ and every layer $L_m/L_0$, if $|s_i-k|_\ell\ge\ell^{-b}$ for all $i$, then
\[
 s_\ell(L_m,k)\le R(\max(m,j_0)+c+b).
\]
For every $k\ge1$ and layer $L_m/L_0$, one also has
\begin{equation}\label{eq:B6-second-bound}
 s_\ell(L_m,k)\le V_\ell(L_{j_0},k)+R(m-j_0)_+.
\end{equation}
\end{theorem}

\begin{proof}
On $X_\ell(L_0)/\ell X_\ell(L_0)$ the continuous tower action factors through a cyclic group of order $\ell^a$. Thus $(\tau-1)^{\ell^a}=\tau^{\ell^a}-1=0$ there. Apply \ref{B6.20} and \ref{B6.21} for the first bound. For the second, apply \eqref{eq:B6-growth} and \eqref{eq:B6-H2-local-correction} at the layer of index $j_0$. Layers of equal index are isomorphic over $L_0$.
\end{proof}

\section{Exceptional twists}

\begin{lemma}[At most $R$ exceptions below each cutoff]\label{B6.23}
Fix the centres of \ref{B6.22} and the corresponding $c$. For each $Z\in\N$ some $\mathcal E_Z\subseteq\{1,\ldots,Z\}$ has size at most $R$. For $1\le k\le Z$ outside $\mathcal E_Z$ and any layer $L_m/L_0$,
\begin{equation}\label{eq:B6-first-bound}
 s_\ell(L_m,k)\le R\bigl(\max(m,j_0(\ell,R))+c+\lfloor\log_\ell Z\rfloor\bigr).
\end{equation}
\end{lemma}

\begin{proof}
For $Z>0$, put $b=\lfloor\log_\ell Z\rfloor$ and exclude the $1\le k\le Z$ with $|s_i-k|_\ell<\ell^{-b}$ for any $i$. Each centre excludes at most one integer, since the integers it excludes are congruent modulo $\ell^{b+1}>Z$. Apply \ref{B6.22}. For $Z=0$ take the empty set.
\end{proof}

\begin{lemma}[A common interval for finitely many exceptions]\label{B6.24}
Let $R\in\N$, $0<\eta\le C_t$, $X>1$, and $\mathcal E\subset\N$ with $\#\mathcal E\le R$. Put $J=2^{R+2}(C_t/\eta)^R$. There is $Y$ with $X^{1/J}\le Y$ and $Y^2\le X$ such that for every $Y\le x\le Y^2$,
\begin{equation}\label{eq:B6-good}
 k\in\mathcal E,\quad k\le x^{C_t}\quad\Longrightarrow\quad k\le x^\eta.
\end{equation}
\end{lemma}

\begin{proof}[Outline]
Set \(U=\log X\), \(B=C_t/\eta\), and \(T=\{\log k/C_t:k\in\mathcal{E},k\ge 1\}\), and induct on \(R\). Take \(y=U/2\) if \(T\) misses \((U/(2B),U]\); otherwise remove a point from that interval, apply induction with endpoint \(U/(2B)\), and set \(Y=e^y\). The complete proof is in \S\ref{prf:B6.24}.
\end{proof}

\section{Small exceptions and the pointwise budget}

Put $P_{C,\varepsilon}(x)=C(1+\log x)+2\exp(Cx^{3\varepsilon/4})$ and $\eta=\varepsilon/(2C_K)$, with $C_K$ from \ref{B6.15}.

\begin{lemma}[A pointwise bound and its sum]\label{B6.25}
Fix a prime $\ell$, $R,c_1,n_0\in\N$, $c_0,C_t,A_0\ge0$, and $J,\varepsilon>0$. There is $C\ge0$ depending on these and $C_K$ for which the bound below holds. Assume the roots of unity of \ref{B6.15} and \eqref{eq:B6-second-bound} for $L_0$. Suppose $\mathcal E$ satisfies \ref{B6.23} at $Z=\lfloor X^{C_t}\rfloor$ with $c\le c_1$, and $[L_{j_0}:\Q]\le n_0$. For $x\ge1$ in $[X^{1/J},X]$ satisfying \eqref{eq:B6-good} and $\log(2|D_{L_{j_0}}|)\le c_0x^{\varepsilon/4}$, each $1\le k\le x^{C_t}$ and $m\le A_0(1+\log x)$ satisfy
\begin{equation}\label{eq:B6-point-bound}
 s_\ell(L_m,k)<\infty,\qquad s_\ell(L_m,k)+1\le P_{C,\varepsilon}(x).
\end{equation}
For every $C_t,C\in\R$ and $\varepsilon>0$, some $x_0\ge2$ satisfies
\begin{equation}\label{eq:B6-final-absorption}
 x^{C_t}P_{C,\varepsilon}(x)\le D_x\qquad(x\ge x_0).
\end{equation}
\end{lemma}

\begin{proof}[Outline]
Use \eqref{eq:B6-first-bound} outside the exceptional set. Inside, combine \eqref{eq:B6-good} and \eqref{eq:B6-small-H2} at the fixed layer, transfer via \eqref{eq:B6-second-bound}, and use the bounds on \(\log X\) and \(m\) to get the logarithmic estimate. The complete proof is in \S\ref{prf:B6.25}.
\end{proof}

\begin{theorem}[Bounds at $2$ and at a fixed odd prime]\label{B6.26}
Fix $J,\varepsilon>0$, $C_t\ge0$ and $A_1\in\R$. Call $a,t:\R\to\N$ admissible if for every $x\ge2$,
\begin{equation}\label{eq:B6-admitted}
 a(x)\ge2,\qquad 2^{a(x)}\le x^{A_1},\qquad 2\le t(x)\le x^{C_t}.
\end{equation}
Each threshold $x_0$ below is uniform in admissible $a,t$ and $X\ge2$. The estimates hold for $x\in[X^{1/J},X]$, $x\ge x_0$, satisfying \eqref{eq:B6-good} for the corresponding exceptional set.

At $2$, fix $M$ as in \ref{B6.18}. There is $x_0\ge2$ such that for every pair \eqref{eq:B6-prime-pair} with $\log(qd)\le X^{\varepsilon/(4J)}$, a set $\mathcal E_2$ of size at most $R_0(M)$ gives, for every realization $E$ of $E_{M,a(x),q,d}$,
\begin{equation}\label{eq:B6-two-sum}
 s_2(E,k)<\infty\quad(1\le k<t(x)),\qquad
 \sum_{1\le k<t(x)}s_2(E,k)\le D_x.
\end{equation}
At an odd prime $\ell$, fix $A$ with $\mu_\ell\subset K_A$, put $R_\ell=R_\ell(A)$ and $j_\ell=j_0(\ell,R_\ell)$. There is $x_0\ge2$ such that for each $a,t,X$ some $\mathcal E_\ell$ of size at most $R_\ell$ gives, for every realization $L_0$ of $A$ and layer $L_{j_\ell}/L_0$,
\begin{equation}\label{eq:B6-odd-layer-sum}
 s_\ell(L_{j_\ell},k)<\infty\quad(1\le k<t(x)),\qquad
 \sum_{1\le k<t(x)}s_\ell(L_{j_\ell},k)\le D_x.
\end{equation}
\end{theorem}

\begin{proof}[Outline]
At \(2\), use \ref{B6.19} and the cyclotomic derivative for field bounds and choose exceptions via \ref{B6.18}, \ref{B6.22}, \ref{B6.23}. Transfer \ref{B6.25} by quadratic comparison at \(2\) and layer isomorphism at odd primes, then sum via \eqref{eq:B6-final-absorption}. The complete proof is in \S\ref{prf:B6.26}.
\end{proof}

\section{Selection of the common window}

\begin{theorem}[The common Iwasawa window]\label{B6.27}
Fix $C_t\ge1$, an abelian datum $M$, a finite set $B'$ of odd primes and, for each $\ell\in B'$, a datum $A_\ell$ with $\mu_\ell\subset K_{A_\ell}$. Fix $0<\varepsilon<1/52$ and $A_1\in\R$. With $C_K$ from \ref{B6.15} and the bounds of \ref{B6.18}, put
\[
 R=R_0(M)+\sum_{\ell\in B'}R_\ell(A_\ell),\qquad
 \eta=\varepsilon/(2C_K),\qquad J=2^{R+2}(C_t/\eta)^R.
\]
There is $X_0$ such that for every admissible $a,t$ of \ref{B6.26}, $X\ge X_0$ and pair \eqref{eq:B6-prime-pair} with
\begin{equation}\label{eq:B6-qd-window-size}
 \log(qd)\le X^{\varepsilon/(4J)},
\end{equation}
there is $Y$ with $X^{1/J}\le Y$ and $Y^2\le X$. For every $x\in[Y,Y^2]$, \eqref{eq:B6-two-sum} holds for every realization $E$ of $E_{M,a(x),q,d}$. For each $\ell\in B'$, \eqref{eq:B6-odd-layer-sum} holds for every realization of $A_\ell$ and its layers of index $j_0(\ell,R_\ell(A_\ell))$.
\end{theorem}

\begin{proof}
Let $x_1$ be the maximum of $2$ and the thresholds in \ref{B6.26}, and set $X_0=x_1^J$. The union $\mathcal E_2\cup\bigcup_{\ell\in B'}\mathcal E_\ell$ has size at most $R$. Since $0<\eta\le C_t$, \ref{B6.24} gives an interval $[Y,Y^2]$ satisfying \eqref{eq:B6-good} for the union. Throughout it $x\ge X^{1/J}\ge x_1$, so both estimates of \ref{B6.26} apply.
\end{proof}

At $C_t=401$, Proposition~\ref{B6.07} combines the window of Theorem~\ref{B6.27} with the automorphic counts to give the window bound of Definition~\ref{B6.03}.

For a primitive Dirichlet character $\psi$ of conductor $f$, take $M=M_\psi$ and $B'=\mathcal B_\psi$ from Definition~\ref{B6.28}. For each $\ell\in B'$, take $A_\ell=(\ell f,\{1\})$, so $K_{A_\ell}=\Q(\zeta_{\ell f})$. The dyadic estimate supplies the term at $2$ in \ref{B6.14}. The unweighted odd-prime sum supplies the hypothesis of \ref{B6.32}. The factor $\log\ell$ enters in \ref{B6.37}.



% ---- p6c4 ----
\chapter{Summing the plus component by classes}\label{ch:b6-classwise}

Write $\chi=\psi$ for the fixed primitive character of conductor $f$ and order $o$. Let $G_F$ denote the absolute Galois group of a field $F$. For a CM field $E$ and a finite set $S$ of odd primes, put
\[
 s^+_{\ell,k}(E)=\operatorname{length}_{\Z_\ell}\Sha^+_{\ell,k}(E),
 \qquad
 P^+(E,t;S)=\sum_{\ell\in S}\log\ell
 \sum_{1\le k<t}s^+_{\ell,k}(E).
\]
The plus component has complex conjugation eigenvalue $(-1)^k$. For the full group over a number field $L$, write $s_{\ell,k}(L)=\operatorname{length}_{\Z_\ell}\Sha_{\ell,k}(L)$. The local data are those of \ref{B6.01} and \ref{B6.08}.

\section{Character components of a host}

\begin{definition}[Hosts and character classes]\label{B6.28}
Let $a\ge2$ be an integer. Let $q,d$ be primes with $q\equiv3\bmod8$ and $d\equiv7\bmod16$. Assume $(qd,f)=1$. Put $N=\operatorname{lcm}(2^{a+2},f)qd$. Let $H\le(\Z/N\Z)^\times$ consist of the units $u$ satisfying
\[
 u\bmod 2^{a+2}\in\{1,-1\},\qquad \chi(u\bmod f)=1,\qquad
 u\bmod q,\ u\bmod d\ \text{are squares}.
\]
A field of this host is a realization of $\Q(\mu_N)^H$. The host is large for $t\in\N$ if $q,d>2t+1$ and $q,d>\max\{f,o,p_K,h_*,p_1\}$. Here $p_K$ is the least prime at least $7$ congruent to $3\bmod4$ and prime to $f$, and $h_*=\#\operatorname{Cl}(\Q(\sqrt{-p_K}))$. Write $o_{\rm odd}$ for the odd part of $o$. Let $p_1$ be the least prime congruent to $1\bmod o_{\rm odd}$ and prime to $2fp_K$. Set
\[
 \mathcal B_\chi=\{3\}\cup\{r:r\text{ is an odd prime dividing }ofp_Kh_*p_1\}.
\]

For an odd prime $\ell$ and a DVR $O\subset\overline{\Q}_\ell$ finite over $\Z_\ell$, define
\[
 X_{\ell,O}(k)=
 \{\theta:(\Z/N\Z)^\times\to O^\times:
       H\subseteq\ker\theta,\ \theta(-1)=(-1)^k\}.
\]
Write $\theta=\theta_0\kappa$, where $\theta_0$ has order prime to $\ell$ and $\kappa$ has $\ell$-power order. For a character $\eta$ modulo $N$, let $\eta_r$ be its Chinese remainder factor on $(\Z/r^{v_r(N)}\Z)^\times$, inflated to $(\Z/N\Z)^\times$ and trivial for $r\nmid N$. Put
\[
 \theta_P=\theta_0(\theta_0)_\ell^{-1},\qquad
 \theta_1=\theta_0(\theta_0)_2^{-1},\qquad
 c_2(\theta)=v_2(\operatorname{cond}\theta_2).
\]
Let $f_0$ be the multiplicative order of $\ell$ modulo $\operatorname{lcm}(4,\operatorname{ord}\theta_1)$, and put $n_0=v_2(\ell^{f_0}-1)$ and $R_1(\ell,\theta)=16o(\ell f)^2\,4^{n_0}$. For a host large for $t$, the triples $(\ell,k,\theta)$ with $1\le k<t$ and $\theta\in X_{\ell,O}(k)$ have the following classes.
\[
\begin{array}{c|l}
 \mathrm{I}&\ell\notin\mathcal B_\chi,\quad \ell-1\nmid2k\\
 \mathrm{II}&k\ge2,\quad\operatorname{ord}\theta_P\text{ even},\quad
       (\ell\in\mathcal B_\chi\text{ or }\ell-1\mid2k)\\
 \mathrm{III}&\ell\notin\mathcal B_\chi,\quad \ell-1\mid2k,\quad
       \theta\ne1,\quad\operatorname{ord}\theta\text{ odd}\\
 \mathrm{IV}&\ell\in\mathcal B_\chi,\quad k\ge2,\quad
       \operatorname{ord}\theta_P\text{ odd}\\
 \mathrm{V}&\ell\in\mathcal B_\chi,\quad k=1,\quad R_1(\ell,\theta)\le2^{c_2(\theta)}\\
 \mathrm{Vs}&\ell\in\mathcal B_\chi,\quad k=1,\quad 2^{c_2(\theta)}<R_1(\ell,\theta)\\
 \mathrm{Z}&\ell\notin\mathcal B_\chi,\quad\ell-1\mid2k,\quad\theta=1 .
\end{array}
\]
\end{definition}

\begin{lemma}[Orders and conductors of host characters]\label{B6.28a}
The host quotient has order dividing $2^{a+2}o$. At odd $r$, the order of $\theta_r$ divides $o$ if $r\mid f$ and divides $2$ otherwise. In particular,
\begin{equation}\label{eq:b6-host-orders}
 \operatorname{ord}\theta\mid 2o\,2^{c_2(\theta)},\qquad
 \operatorname{ord}\theta_1\mid2o,\qquad
 f_0\le F_0:=\varphi(4o),\qquad 2^{n_0}<\ell^{F_0}.
\end{equation}
If $\operatorname{ord}\theta$ is odd, then $\theta_2=1$. In this case $\operatorname{ord}\theta\mid o_{\rm odd}$ and $\operatorname{cond}\theta\mid f$.
For a host large for $t$, if $\ell\in\mathcal B_\chi$ then $\ell<q,d$ and $\operatorname{cond}\theta\mid f$ whenever $\operatorname{ord}\theta_P$ is odd. If $\ell\notin\mathcal B_\chi$, then $\theta_0=\theta$ and $\theta_\ell^2=1$. If also $\ell-1\mid2k$ and $1\le k<t$, then $5\le\ell\le2k+1<q,d$, $k\ge2$, $\ell\nmid N$ and $\theta_P=\theta$.
\end{lemma}
\begin{proof}[Outline]
Embed the host quotient in the dyadic, character-image and square-class factors to bound its order and conductor. Largeness makes the factors at \(2,q,d\) survive in \(\theta_P\) and vanish at odd order, and \(\ell-1\mid 2k\) excludes \(\ell\mid N\) outside the fixed set. The complete proof is in \S\ref{prf:B6.28a}.
\end{proof}

\begin{definition}[Character lengths]\label{B6.29}
Let $\varepsilon_\ell$ be the cyclotomic character of $G_\Q$, and write $D_p,I_p$ for a decomposition group and its inertia subgroup at $p$. Inflate $\theta$ via the action on $\mu_N$ and put
\[
 A_O=O\otimes_{\Z_\ell}(\Q_\ell/\Z_\ell),\qquad
 M_{\ell,k,\theta}=A_O(\theta^{-1}\varepsilon_\ell^{-k}).
\]
Let $\mathcal S_{\ell,k,\theta}\subseteq H^1(G_\Q,M_{\ell,k,\theta})$ consist of the classes vanishing on finite inertia and on decomposition groups at $p\mid\ell N$. Set $L_{\ell,k,\theta}=\operatorname{length}_O\mathcal S_{\ell,k,\theta}$ and $e_O=\operatorname{length}_O(O/\ell O)$. At finite length put $a_{\ell,k,\theta}=L_{\ell,k,\theta}/e_O$. A character bound $B$ means finite length and $a_{\ell,k,\theta}\log\ell\le B$.
\end{definition}

\begin{corollary}[Independence of the coefficient ring]\label{B6.29a}
If two coefficient DVRs give the same $\overline{\Q}_\ell$-valued character, the lengths in \ref{B6.29} are finite simultaneously and their normalized values agree when finite.
\end{corollary}
\begin{proof}
Compare both rings with the ring of integers of a common finite extension in $\overline{\Q}_\ell$. Finite free extension commutes with continuous cohomology and the restriction kernels because it commutes with continuous cochains and is exact. It multiplies both lengths by the ramification index. Faithful flatness preserves the distinction between finite and infinite length.
\end{proof}

\section{The automorphic counting bounds}

Assume the inputs $\mathscr H$ defined before Lemma~\ref{B6.06f}, and the bound of \ref{B6.15}. A bar denotes reduction modulo the coefficient maximal ideal, and $\omega$ denotes the Teichmüller lift of $\bar\varepsilon_\ell$. Write $C_\ell=\overline{\Q}_\ell$ and normalize $v_\ell$ by $v_\ell(\ell)=1$.

\begin{definition}[Cuspidal attachment with witness]\label{def:b6-cuspidal-witness}
Let $K$ be a CM field, let $r:G_K\to\GL_n(C_\ell)$ be continuous and irreducible, and fix $\iota:C_\ell\simeq\C$. A \emph{cuspidal attachment} of $r$ is a regular algebraic cuspidal automorphic representation $\pi$ of $\GL_n(\mathbb A_K)$ with the following property. At all but finitely many finite places $v$, both $r$ and $\pi$ are unramified. At those places $\iota$ carries the eigenvalue multiset of $r$ at inverse arithmetic Frobenius to $q_v^{(n-1)/2}$ times the Satake multiset of $\pi_v$, where $q_v$ is the residue cardinality.

Call a place $v\mid\ell$ \emph{unlinked} for $r$ if, at every prolongation of $v$, the decomposition group preserves a full flag whose quotient characters satisfy $\kappa_i\ne\kappa_j\varepsilon_\ell$ for all $i,j$. Call $v$ \emph{potentially unlinked} if at every prolongation one can choose a full stable flag for which $\kappa_i^m\ne(\kappa_j\varepsilon_\ell)^m$ for all $i,j$ and all integers $m\ge1$.

We say that $r$ has a \emph{cuspidal attachment with witness} if it has a cuspidal attachment and the following conditions hold. Every unlinked place above $\ell$ is potentially unlinked. For each cuspidal attachment $\pi$ of $r$ there is a real number $t_\pi$ such that, at every finite $v\nmid\ell$ and every potentially unlinked $v\mid\ell$, the cuspidal support of $\pi_v$ consists of characters of $K_v^\times$. Each such character $\xi$ satisfies
\[
 q_v^{\,t_\pi-1/2}<|\xi(\varpi_v)|<q_v^{\,t_\pi+1/2},
\]
where $\varpi_v$ is a uniformiser.
\end{definition}

\begin{lemma}[Restriction in prime rank]\label{B6.30aa}
Let \(G\) be a profinite group, \(H\) a closed normal subgroup, and \(r\) a continuous irreducible representation of \(G\) of dimension \(n=2\) or \(3\) over \(\overline{\Q}_\ell\), defined over a finite extension of \(\Q_\ell\). Suppose its residual semisimplification is a sum of characters of \(G\). Let \(D\le G\) preserve a complete flag, with quotient characters \(\lambda_1,\ldots,\lambda_n\). Suppose some \(u\in D\cap H\) has pairwise distinct \(\lambda_i(u)\), and some \(s\in D\) has \(\bar\lambda_a(s)\ne\bar\lambda_b(s)\). Then \(r|_H\) is irreducible. If every positive power of \(u\) has distinct eigenvalues, the same holds on \(H\cap U\) for every open subgroup \(U\le G\).
\end{lemma}
\begin{proof}[Outline]
Translates of a simple submodule give two reducible cases, scalar action or permuted lines. Distinct eigenvalues exclude scalars, and the flag and residual separation exclude permuted lines because their reduced characters coincide on the permutation kernel. The complete proof is in \S\ref{prf:B6.30aa}.
\end{proof}

\begin{lemma}[Irreducibility on the subgroups used at a marked point]\label{B6.30ac}
Under the character and field hypotheses of \ref{B6.40}, let \(r\) be a characteristic-zero ordinary marked point. In Type I its restriction to every open subgroup of \(G_K\) is irreducible. In Type II let \(W\) extend the irreducible two-dimensional corner as in \ref{B6.52}(3). Then \(W|_{G_K\cap\ker\varepsilon_\ell}\) is irreducible.
\end{lemma}
\begin{proof}[Outline]
In Type I, distinct inertia eigenvalues exclude scalars on the normal core, and residual separation from \ref{eq:p6c5-link-input-tests} gives a trace contradiction for permuted lines. In Type II, a projective centralizer and the inertia trace formula exclude reducibility. The complete proof is in \S\ref{prf:B6.30ac}.
\end{proof}

\begin{lemma}[The finite adjoint classes at a characteristic-zero point]\label{B6.30ab}
At a Type I point as in \ref{B6.30ac} the marked polarised adjoint has zero Bloch--Kato Selmer group over \(\Q\). At a Type II point the same holds for the polarised adjoint of its corner with determinant \(\eta_K\) and newform attachment as in \ref{B6.52}(3).
\end{lemma}
\begin{proof}[Outline]
Ordinary eigenvalues and trace identities prove irreducibility on the required subgroups. Use \ref{B6.30ab}, the attachments of \ref{B6.52} and base change in \ref{B6.06} to identify the polarised adjoints and apply Thorne's vanishing as in \ref{B6.53}. The complete proof is in \S\ref{prf:B6.30ab}.
\end{proof}

\begin{lemma}[Multiplicity in the projected ordinary space]\label{B6.30ad}
Assume \(\mathscr H\). Use the rank-three definite unitary comparison of \ref{B6.51a}, at its base level. Write \(L/F\) for the CM extension, \(O\) for the coefficient DVR, \(\Lambda\) for the weight ring, \(B_0\) for the universal ordinary determinant ring and \(M\) for the ordinary module of \ref{B6.48c}. Put \(J=\operatorname{Ann}_{B_0}M\). There is a \(\Lambda\)-linear idempotent \(e\) on \(M\) commuting with \(B_0\) with the following property.

For a characteristic-zero arithmetic character \(x\) of \(B_0/J\), put
\[
 M_x=\overline{\Q}_\ell\otimes_{\Lambda,x|_\Lambda}M,\qquad
 \mu_e(x)=\dim_{\overline{\Q}_\ell}(eM_x)[x]_{\rm gen}.
\]
Here the bracket denotes the simultaneous generalized eigenspace for the action of \(B_0\) with character \(x\). Call \(x\) irreducible if its determinant is the determinant of an absolutely irreducible representation. If \(x,y\) are irreducible arithmetic characters of the same strictly increasing labelled algebraic weight, and the finite parts of the ordered characters of \(y\) are pairwise distinct on the units at each place above \(\ell\), then
\[
                         1\leq\mu_e(x)=\mu_e(y).
\]
\end{lemma}
\begin{proof}[Outline]
Builds \(e\) from nondegenerate characters of \(U(k_0)\), finding rank \((q_0-1)^2\) on generic \(\tau^{K^1}\) via Whittaker models. Counts the fibre factors using \ref{B6.48d}, \ref{B6.04a}, \ref{B6.48a}, multiplicity one, Jacquet modules at \(v\mid\ell\), and tame types. The complete proof is in \S\ref{prf:B6.30ad}.
\end{proof}

\begin{lemma}[Fixed-weight derivations on the Hecke quotient]\label{B6.30ae}
In Lemma~\ref{B6.30ad}, let \(x\) be an irreducible arithmetic point of \(B_0/J\) of strictly increasing labelled algebraic weight. Every \(O\)-linear derivation at \(x\) on \(B_0/J\) which kills the values of the ordered characters on ordinary inertia is zero.
\end{lemma}
\begin{proof}[Outline]
Uses separated characters via \ref{B6.48c} and \ref{B6.48d} to detect \(B_0/J\) and pass to its image on \(eM\) by \ref{B6.30ad}. Comparing multiplicities of a separating element \(t\) via \ref{B6.30ad} gives \(x(f'(t))\ne0\), which kills fixed-weight derivations. The complete proof is in \S\ref{prf:B6.30ae}.
\end{proof}

\begin{lemma}[Local vanishing of the Hecke ideal on the matrix space]\label{B6.30af}
Assume \(\mathscr H\), and use the rank-three comparison over \(L\) in \ref{B6.51a} with coefficient DVR \(O\). Write \(A\) for the graded matrix-coordinate algebra, \(B=A_0\) for its weight-zero coefficient ring, and \(\phi:B_0\to B\) for the determinant map after twisting the matrices by \(\varepsilon_\ell^{-1}\). The grading here is diagonal change of basis with weights \(1,0,-1\). Let \(\xi:A\to\overline{\Q}_\ell\) be an integral \(O\)-algebra homomorphism, taking the maximal ideal of \(B\) to elements of positive valuation. Suppose its kernel is off the two triangular loci and the locus on which the middle coordinate splits from the outer two. For each \(a\in J\) there exists \(s\in A\) with
\[
                        \xi(s)\ne0,\qquad s\phi(a)=0.
\]
\end{lemma}
\begin{proof}[Outline]
Shows by \ref{B6.51a} that \(A_{\mathfrak q}\) has no embedded primes and \(\phi(J)\subset P\) at each minimal prime \(P\). Finds a point on \(P\) via \ref{B6.30ae} and proves regularity there using \ref{B6.48a}, \ref{B6.30aa}, and \ref{B6.54}, making \(A_P\) a field. The complete proof is in \S\ref{prf:B6.30af}.
\end{proof}

\begin{lemma}[Two fixed-weight directions at a crystalline link]\label{B6.30ag}
In the rank-three frame of \ref{B6.51a}, let \(\xi:R\to O\) be an integral evaluation of the complete framed ordinary matrix ring, after finite coefficient extension. Write \(\lambda_{v,i}\), \(0\le i\le2\), for the ordered local characters. If the representation is potentially crystalline at \(v\mid\ell\) and
\(\lambda_{v,i}=\varepsilon_\ell\lambda_{v,i+1}\) for some \(i=0,1\), then there are two \(O\)-linearly independent \(O\)-valued derivations at \(\xi\) which kill every ordinary character value on inertia.
\end{lemma}
\begin{proof}[Outline]
Uses crystallinity of the adjacent subquotient to build an extra local cocycle via \ref{B6.05d}. Global duality in \ref{B6.05d} and \ref{B6.51frame} then bounds normalized cocycles mod \(\ell^m\) by \(C|Z_m|\ge q^2\), forcing two independent \(O\)-valued derivations. The complete proof is in \S\ref{prf:B6.30ag}.
\end{proof}

\begin{lemma}[No cyclotomic link at the Type I endpoint]\label{B6.30ah}
Assume \(\mathscr H\) and the character and field conditions of \ref{B6.40}. Let \(r\) be an integral ordinary marked point as in \ref{B6.52}, irreducible on every open subgroup of \(G_K\). Suppose that at each \(w\mid\ell\) its rank-three restriction to \(G_K\) has a flag of level \((N;k,0,-k)\), with the local alternative that the flag is unlinked or the restriction is potentially crystalline. Then at each \(w\mid\ell\) the level flag is unlinked. In particular this holds for \(k=1\).
\end{lemma}
\begin{proof}[Outline]
Places a linked point in the frame of \ref{B6.52} and \ref{B6.51a}, where \ref{B6.30ag} gives two independent fixed-inertia derivations. Both kill \(B\) by \ref{B6.30af} and \ref{B6.30ae}, so a combination kills a weight-\(1\) monomial and the whole ring, a contradiction. The complete proof is in \S\ref{prf:B6.30ah}.
\end{proof}

\begin{corollary}[Representation inputs for the first count]\label{B6.30a}
For every datum satisfying the character and field conditions of \ref{B6.40}, the required Type I admissibility and adjoint rigidity follow from $\mathscr H$ and \ref{B6.15}. The same inputs give the required Type II block rigidity; in both types the vanishing of the adjoint Bloch--Kato Selmer group at the point is Lemma~\ref{B6.30ab}.
\end{corollary}
\begin{proof}[Outline]
In Type I, gets admissibility via \ref{B6.30aa}, \ref{B6.40n}, \ref{B6.30ah}, and \ref{B6.52}, and rigidity via \ref{B6.48k} and \ref{B6.53}. In Type II, applies \ref{B6.52}, \ref{B6.30aa}, and \ref{B6.53}, killing other adjoint summands by \ref{B6.43} and \ref{B6.40p}. The complete proof is in \S\ref{prf:B6.30a}.
\end{proof}

\begin{theorem}[Class I, summed over primes]\label{B6.30}
For every $\chi$ there are $C_{\mathrm I}\in\N$ and distinct primes $r_1,r_2$. Let a host be large for $t$ with $q,d>r_1,r_2$. Fix $1\le k<t$ and a finite set $S$ of odd primes with $\ell\notin\mathcal B_\chi$ and $\ell-1\nmid2k$ for each $\ell\in S$. For every choice of coefficient rings $(O_\ell)_{\ell\in S}$ the length is finite for every $\ell\in S$ and $\theta\in X_{\ell,O_\ell}(k)$, and
\[
 \sum_{\ell\in S}\log\ell\sum_{\theta\in X_{\ell,O_\ell}(k)}
 a_{\ell,k,\theta}\le(Nt)^{C_{\mathrm I}} .
\]
\end{theorem}
\begin{proof}[Outline]
Chooses primes \(r_1,r_2>30fp_K+1\) and verifies \ref{B6.40m} and \ref{B6.40} via \ref{B6.28a}, \ref{B6.05}, \ref{B6.09}, \ref{B6.15}, \ref{B6.30a}, and \ref{B6.41}. Sums over character kernels using the two-prime estimate in \ref{B6.42} to bound the total. The complete proof is in \S\ref{prf:B6.30}.
\end{proof}

\begin{remark}[The trivial character and \(K_{4j}(\Z)\)]\label{B6.30b}
Fix \(\chi\), and put \(B=\{2\}\cup\mathcal B_\chi\). For a finite abelian group \(D\), write \(D[\ell^\infty]\) for its \(\ell\)-primary subgroup. There are constants \(A_\chi,C_\chi>0\) such that, for every integer \(j\ge1\),
\[
 \log\!\left(\prod_{\ell\notin B}
             |K_{4j}(\Z)[\ell^\infty]|\right)
 \le A_\chi(1+j)^{C_\chi}.
\]
The primes with \(\ell-1\nmid4j\) occur in \ref{B6.30}. The primes outside \(B\) with \(\ell-1\mid4j\) have zero contribution by \ref{B6.36}.
\end{remark}
\begin{proof}[Outline]
Choose host primes by \ref{B2.23} and take the trivial character of \ref{B6.28}. Duality in \ref{B6.09} and Weibel's comparison identify the dual strict groups of \ref{B6.29} with the primary parts of \(K_{4j}(\mathbb Z)\), then \ref{B6.36} and \ref{B6.30} bound the sum. The complete proof is in \S\ref{prf:B6.30b}.
\end{proof}


\begin{definition}[The marked symplectic datum]\label{B6.31b}
Use the host, \(N,f,o,p_1,n_0,c_2(\theta)\) and character classes of Definition~\ref{B6.28}. Fix a host large for \(t\), an odd prime \(\ell\), a coefficient DVR \(O\) as in that definition, \(1\le k<t\), and \(\theta\in X_{\ell,O}(k)\) of class II or III. Regard Dirichlet characters as characters of \(G_{\Q}\) by their action on roots of unity. Write \(\varepsilon_\ell\) for the cyclotomic character and \(\omega\) for the Teichm\"uller lift of its reduction. Put
\begin{equation}\label{eq:b6-sp4-levels}
\begin{aligned}
 n_{\mathrm{II}}&=\left\lceil\log_2\max\{\ell N4^{n_0},\,2^{c_2(\theta)+6}o\ell^3f^2\}\right\rceil,\\
 n_{\mathrm{III}}&=\left\lceil\log_2\max\{\ell fp_1 4^{n_0},\,8\ell^2\}\right\rceil .
\end{aligned}
\end{equation}
In II take \((n',p',m')=(n_{\mathrm{II}},1,1)\). In III take
\((n',p',m')=(n_{\mathrm{III}},p_1,\operatorname{ord}\theta)\).
Choose \(O\)-valued characters \(\xi_2,\xi_0\), where \(\xi_2\) is primitive odd modulo \(2^{n'}\) and \(\xi_0\) has order \(m'\) modulo \(p'\). Set
\[
 \xi=\xi_0\xi_2,\quad \Psi=\varepsilon_\ell^k\theta,\quad
 \alpha=\varepsilon_\ell\omega^{-1}\xi,\quad \nu=\alpha\Psi,\quad
 (\delta_0,\delta_1,\delta_2,\delta_3)=(\nu,\Psi,\alpha,1).
\]
Write \(r_0=\operatorname{diag}(\delta_0,\delta_1,\delta_2,\delta_3)\) and \(\bar r\) for its reduction. Let \(S_D\) be the rational primes dividing \(2\ell Np'\). Choose a prolongation at each rational prime \(v\), with decomposition group \(D_v\) and inertia \(I_v\), and choose a complex conjugation. Set
\[
 T=\{\ell\}\cup\{v\in S_D\setminus\{\ell\}:\ell\mid\#\theta(I_v)\}.
\]

A good datum includes these choices and satisfies the following conditions. We have \(k\ge2\), \(\Psi\) is even, and \(\xi\) is odd of order prime to \(\ell\) and unramified at \(T\). The characters \(\theta,\xi\) are unramified outside \(S_D\). For \(v\in T\) there are arithmetic Frobenius lifts \(\sigma_v\), with \(\varepsilon_\ell(\sigma_\ell)=1\), such that
\(\overline{\xi\Psi^i}(\sigma_v)\ne1\) for every \(i\in\Z\).
There are markings
\[
 \phi=\lim_{a\to\infty}\sigma_\ell^{\,\ell^{a!}},\qquad
 g_0\in G_{\Q},\qquad s_0\in I_\ell\cap\ker\theta\cap\ker\omega
\]
such that \(\bar\alpha(g_0),\bar\Psi(g_0)\ne1\) and \(\varepsilon_\ell(s_0)\) has infinite order.

Finally, let \(\gamma\) be any of \(\alpha,\Psi/\alpha,\nu\) or their inverses. A continuous cocycle with values in the residue field of \(O\) acted on by \(\bar\gamma\) is a global coboundary if it is a coboundary on each \(D_v\) for \(v\in T\) and vanishes on \(I\cap\ker\bar r\) for every finite inertia group \(I\) away from \(T\).
\end{definition}

\begin{definition}[Allowed lifts and their points]\label{B6.31c}
Fix a good datum as in Definition~\ref{B6.31b}. Use the alternating form with antidiagonal \((1,1,-1,-1)\) in the indicated coordinate order. A matrix is block diagonal here if it preserves both coordinate planes with indices \(\{0,2\}\) and \(\{1,3\}\). Let \(R\) be a complete Noetherian local \(O\)-algebra whose structure map is local and induces an isomorphism on residue fields. An allowed lift over \(R\) is a continuous lift \(r:G_{\Q}\to\GL_4(R)\) of \(\bar r\) with similitude \(\nu\) and the following properties.

Away from \(T\), \(r\) is trivial on \(I\cap\ker\bar r\) for every finite inertia group \(I\). The matrix \(r(\phi)\) is block diagonal. The entries \(02,20\) of \(r(g_0)+\nu(g_0)r(g_0)^{-1}\) vanish.

At each \(v\in T\) require a complete Noetherian local \(O\)-domain \(B_v\) with the same residue field, a lift \(b_v\) of \(\bar r|_{D_v}\) of similitude \(\nu\), and \(O\)-algebra maps \(B_v\to O,\ B_v\to R\) taking \(b_v\) to \(r_0|_{D_v},r|_{D_v}\), respectively. For \(v\ne\ell\) require \(b_v=1\) on \(I_v\cap\ker\theta\). At \(\ell\) put \(s=s_0\), \(d_i=\delta_i(s)\) and \(P_i=\prod_{j\ne i}(b_\ell(s)-d_jI)\). The local equations are
\[
 \det(XI-b_\ell(s))=\prod_i(X-d_i),\qquad
 P_j b_\ell(g)P_i=0\quad(i<j,\ g\in D_\ell),\qquad
 P_i(b_\ell(h)-\delta_i(h)I)P_i=0\quad(h\in I_\ell).
\]
For any commutative \(O\)-algebra \(A\), an allowed \(A\)-point is the entrywise image of an allowed lift under an \(O\)-algebra map \(R\to A\). A block point has block diagonal image. A characteristic-zero point in the following statements has \(A=\overline{\Q}_\ell\), denoted by \(C_\ell\). An allowed dual-number lift of such a point is an allowed point over \(C_\ell[\epsilon]/(\epsilon^2)\) reducing to it.

Give entry \(ij\) weight \(v_i-v_j\), where
\((v_0,v_1,v_2,v_3)=((1,1),(0,1),(1,0),(0,0))\).
The weight-zero algebra of \(r\) is the \(O\)-subalgebra generated by products of entries \(r(g)_{ij}\) of total weight \((0,0)\), with the group elements allowed to vary among the factors.
\end{definition}

\begin{lemma}[Cocycles of negative weight]\label{B6.31a}
Let $\eta:G_\Q\to C_\ell^\times$ be a continuous character of finite order. For $e\ge1$, put $\lambda=\eta^{-1}\varepsilon_\ell^{-e}$. Every continuous $C_\ell(\lambda)$-valued cocycle that is a coboundary on all finite inertia groups and on $D_\ell$ is a global coboundary.
\end{lemma}
\begin{proof}[Outline]
Baire's theorem reduces the cocycle to finite coefficients, and \ref{B6.15} and \ref{B6.32a} bound the restriction image and kernel for torsion coefficients. Scale and reduce the cocycle, then evaluate at \(\lambda(g)\ne1\) to obtain a limiting cobounding scalar. The complete proof is in \S\ref{prf:B6.31a}.
\end{proof}

\begin{lemma}[Existence of the marked datum]\label{B6.31d}
For every host and character of class II or III as in Definition~\ref{B6.31b}, a finite extension of the coefficient DVR admits a good datum for this host and character, with the levels of \eqref{eq:b6-sp4-levels}.
\end{lemma}
\begin{proof}[Outline]
Use \ref{B6.28a} to choose the odd-modulus factor and pair it with a primitive odd dyadic character to build the diagonal lift. Verify Frobenius separation and the markings, and use \ref{B6.59} to obtain residual Selmer vanishing. The complete proof is in \S\ref{prf:B6.31d}.
\end{proof}

\begin{lemma}[Construction and formal smoothness of the local domains $B_\ell$ and $B_v$]\label{B6.31e}
Fix a good datum $\mathcal D$ as in Definition~\ref{B6.31b}, with $E=\operatorname{Frac}O$, and give entry $ij$ of $M_4$ bidegree $(l_i-l_j,m_i-m_j)$, where $l=(1,0,1,0)$ and $m=(1,1,0,0)$.
\begin{enumerate}
\item For each $v\in T\setminus\{\ell\}$, the quotient $D_v/(I_v\cap\ker\theta)$ is abelian, topologically generated by $\sigma_v$ and an inertia generator $\tau_v$ of finite order $m_v=\#\theta(I_v)$. Let $U_v$ be the universal lifting ring for continuous lifts $b_v:D_v/(I_v\cap\ker\theta)\to\mathrm{GSp}_4(R)$ of $\bar r|_{D_v}$ with similitude $\nu|_{D_v}$, with base point $x_{0,v}:U_v\to O$ specialising to $r_0|_{D_v}$. Put $\zeta=\theta(\tau_v)\in O^\times$, let $J_v\subset U_v$ be the closed ideal generated by the coefficients of $\det(XI_4-b_v(\tau_v))-(X-\zeta)^2(X-1)^2$, and let $I_v$ be the intersection of the minimal primes of $\operatorname{sat}_\ell(\sqrt{J_v})$ contained in $\ker x_{0,v}$, where $\operatorname{sat}_\ell(\mathfrak a)=\{a:\exists n\ge0,\ \ell^n a\in\mathfrak a\}$. Then $I_v$ is a single prime ideal invariant under $\mathrm{Sp}_4$-conjugation congruent to $I_4$, so $B_v=U_v/I_v$ is a complete Noetherian local $O$-domain with a bigraded homogeneous presentation. Over $E$, the completion $\widehat{B_v[1/\ell]}_{\ker x_{0,v}}$ is formally smooth of dimension $10$, with coordinate counts $(1,1,2)$ in first-degree-zero bidegrees $((0,-1),(0,1),(0,0))$ and $1$ in each of the six cross-block bidegrees $(1,0),(-1,0),(1,1),(-1,-1),(-1,1),(1,-1)$; each cross-block tangent direction is a coboundary in $B^1(D_v,\mathfrak{sp}_4)$.
\item At $\ell$, let $U_\ell$ be the universal lifting ring for continuous lifts $b_\ell:D_\ell\to\mathrm{GSp}_4(R)$ of $\bar r|_{D_\ell}$ with similitude $\nu|_{D_\ell}$, with base point $x_{0,\ell}:U_\ell\to O$. Let $J_\ell\subset U_\ell$ be the closed ideal generated by the local ordinary equations of Definition~\ref{B6.31c} (for $s=s_0$), and let $I_\ell$ be the intersection of the minimal primes of $\operatorname{sat}_\ell(\sqrt{J_\ell})$ contained in $\ker x_{0,\ell}$. Then $I_\ell$ is a single prime ideal invariant under $\mathrm{Sp}_4$-conjugation congruent to $I_4$, so $B_\ell=U_\ell/I_\ell$ is a complete Noetherian local $O$-domain with a bigraded homogeneous presentation. Over $E$, the completion $\widehat{B_\ell[1/\ell]}_{\ker x_{0,\ell}}$ is formally smooth of dimension $14$, with coordinate counts $(1,2,2)$ in first-degree-zero bidegrees $((0,-1),(0,1),(0,0))$, count $2$ in each of the three upper cross-block bidegrees $(1,0),(-1,1),(1,1)$, and count $1$ in each of the three lower cross-block bidegrees $(-1,0),(1,-1),(-1,-1)$.
\item In the framed presentation of \ref{B6.58} over $B=B_\ell\widehat\otimes_O\widehat\bigotimes_{v\in T\setminus\{\ell\}}B_v$ (with frame at $\ell$ fixed to $I_4$), imposing the six cross-block equations $(0,1),(0,3),(2,1),(1,0),(3,0),(1,2)$ on $r(\phi)=b_\ell(\phi)$ and on each frame matrix $h_v$ ($v\in T\setminus\{\ell\}$) eliminates $6+6(|T|-1)$ cross-block coordinates over $E$. The resulting formal chart over $E$ has $(1,2,2)$ local coordinates from $\ell$ and $(1,1,2)$ from each $v\in T\setminus\{\ell\}$ in first-degree-zero bidegrees $((0,-1),(0,1),(0,0))$, together with three coordinates of nonzero first degree, namely $u:(1,0)$, $v:(-1,1)$, and $w:(1,1)$. Its first-degree-zero formal subring over $E$ adds the two algebraically independent coordinates $uv$ (weight $1$) and $wv$ (weight $2$).
\end{enumerate}
\end{lemma}
\begin{proof}[Outline]
Uses Lagrangian eigenspaces at \(v\ne\ell\) and flag lifting at \(\ell\) to show \(B_v,B_\ell\) are formally smooth domains of dimensions \(10,14\). Eliminates cross-block coordinates in \ref{B6.58} via \ref{B6.09} and the formal implicit function theorem. The complete proof is in \S\ref{prf:B6.31e}.
\end{proof}

\begin{theorem}[Classes II and III]\label{B6.31}
For each $\chi$ there are $C_{\mathrm{II}},C_{\mathrm{III}}\in\N$. For every host large for $t$, every odd prime $\ell$, every coefficient ring $O$ and every $1\le k<t$, a character $\theta\in X_{\ell,O}(k)$ of class II or III has finite length and satisfies the respective bound
\[
 a_{\ell,k,\theta}\log\ell\le(\ell Nk)^{C_{\mathrm{II}}},\qquad a_{\ell,k,\theta}\log\ell\le(\ell Nk)^{C_{\mathrm{III}}}.
\]
\end{theorem}
\begin{proof}[Outline]
Uses \ref{B6.31d}, \ref{B6.31e}, and \ref{B6.58} to build a bigraded presentation, killing nonpositive tangents via \ref{B6.31a}. Proves two-sided points are irreducible (\ref{B6.49}) and rigid (\ref{B6.48}), then bounds lengths via \ref{B6.58}, \ref{B6.50}, and \ref{B6.29a}. The complete proof is in \S\ref{prf:B6.31}.
\end{proof}

\section{Class IV and its cyclotomic layer}

\begin{lemma}[Restriction to an abelian field]\label{B6.32a}
Let $k\ge1$ and $L_1/\Q$ be abelian of degree $n$ with $\theta|_{G_{L_1}}=1$. Then
\[
 L_{\ell,k,\theta}\le e_O\bigl(s_{\ell,k}(L_1)+n\nu_\ell(N,k)\bigr),
 \tag{6.IV}\label{eq:b6-transfer}
\]
where
\[
 \nu_\ell(N,k)=
 \begin{cases}
 \max\{v_\ell(N),1\}+v_\ell(k),&\ell\mid N\text{ or }\ell-1\mid k,\\
 0,&\text{otherwise}.
 \end{cases}
\]
The inertia invariants $M_{\ell,k,\theta}^{I_\ell}$ have length at most $e_O\nu_\ell(N,k)$, and $\nu_\ell(N,k)\le2Nk$.
\end{lemma}
\begin{proof}[Outline]
Inflation–restriction and scalar coboundaries bound the kernel length by \(n\) times the global invariant length. A cyclotomic inertia lift bounds inertia invariants, with a unit difference if \(\ell\nmid N\) and \(\ell-1\nmid k\), and restriction lands in the strict group. The complete proof is in \S\ref{prf:B6.32a}.
\end{proof}

\begin{theorem}[Transfer of the class IV sum]\label{B6.32}
Fix $\chi$, $\ell\in\mathcal B_\chi$ and $j\in\N$. Let $t\ge2$. Assume that for every realization $L_0$ of $\Q(\mu_{\ell f})$, cyclotomic layer $L_1/L_0$ of degree $\ell^j$ and choice of local data, the lengths for $1\le k<t$ are finite and $\sum_{1\le k<t}s_{\ell,k}(L_1)\le B$. For every host large for $t$ and every $O$-valued host character $\theta$ with $\operatorname{ord}\theta_P$ odd, the lengths $L_{\ell,k,\theta}$ are finite for $1\le k<t$ and
\[
 \sum_{1\le k<t}a_{\ell,k,\theta}\le B+(Nt)^{C_4},
 \qquad C_4=\ell f\ell^j+3 .
\]
\end{theorem}
\begin{proof}
By \ref{B6.28a}, $\operatorname{cond}\theta\mid f$. Take the field of \ref{B6.33}, of degree $n\le\ell f\ell^j$. After division by $e_O$ the sum of the errors in \eqref{eq:b6-transfer} is at most $2nNt^2\le(Nt)^{C_4}$.
\end{proof}

\begin{lemma}[A field that trivializes the transferred characters]\label{B6.33}
For every odd prime $\ell$ and $j\in\N$, each algebraic closure of $\Q$ contains abelian fields $L_0\subseteq L_1$ with $L_0\simeq\Q(\mu_{\ell f})$. The extension $L_1/L_0$ is a cyclotomic layer of degree $\ell^j$ and satisfies $[L_1:\Q]\le\ell f\ell^j$. For every integral domain $A$ and $N\ge1$, the Galois character associated to an $A$-valued Dirichlet character of level $N$ and conductor dividing $f$ is trivial on $G_{L_1}$.
\end{lemma}
\begin{proof}
Take $L_0=\Q(\mu_{\ell f})$ and $L_1=\Q(\mu_{\ell^{j+1}f})$. The totient ratio is $\ell^j$, and $[L_1:\Q]=\ell^j\varphi(\ell f)\le\ell^j\ell f$. The character factors through the action on $\mu_f\subseteq L_0$.
\end{proof}

\section{The two classes at twist one}

\begin{theorem}[Class V]\label{B6.34}
For each $\chi$ there is $C_{\mathrm V}\in\N$. For every host large for $t\ge2$, every $\ell\in\mathcal B_\chi$ and $\theta\in X_{\ell,O}(1)$ with $R_1(\ell,\theta)\le2^{c_2(\theta)}$,
\[
 L_{\ell,1,\theta}<\infty,\qquad a_{\ell,1,\theta}\log\ell\le(\ell N)^{C_{\mathrm V}}.
\]
\end{theorem}
\begin{proof}[Outline]
Builds the Gaussian frame ring (\ref{B6.51}, \ref{B6.52}), counts relations via \ref{B6.09}, and forces Type I points via \ref{B6.59}. Kills nonpositive tangents by \ref{B6.31a}, proves rigidity via \ref{B6.51a} and \ref{B6.53}, and bounds lengths by \ref{B6.58} and \ref{B6.41}. The complete proof is in \S\ref{prf:B6.34}.
\end{proof}

\begin{theorem}[Class Vs]\label{B6.35}
For each $\chi$ there is $C_{\mathrm{Vs}}\in\N$. For every host large for $t\ge2$, every $\ell\in\mathcal B_\chi$ and $\theta\in X_{\ell,O}(1)$ with $2^{c_2(\theta)}<R_1(\ell,\theta)$,
\[
 L_{\ell,1,\theta}<\infty,\qquad a_{\ell,1,\theta}\log\ell\le(\ell N)^{C_{\mathrm{Vs}}}.
\]
\end{theorem}
\begin{proof}[Outline]
Class Vs and the host order bounds give bounded degree after adjoining \(\mu_\ell\). Bound the discriminant by cyclotomic inclusion and transitivity, then apply \ref{B6.15} and absorb the transfer error by restriction to obtain a power bound in \(\ell N\). The complete proof is in \S\ref{prf:B6.35}.
\end{proof}

\section{The trivial character}

\begin{theorem}[Vanishing in class Z]\label{B6.36}
Let $\ell\ge5$ be prime and $k\ge2$ even with $\ell-1\mid2k$. For every $O$, the subgroup of $H^1(G_\Q,A_O(-k))$ vanishing on finite inertia and on a decomposition group at $\ell$ has length zero. In particular, $L_{\ell,k,1}=a_{\ell,k,1}=0$ in class Z.
\end{theorem}
\begin{proof}[Outline]
For \(\ell\)-torsion use Kronecker–Weber and decomposition at \(\ell\) if \(\ell-1\mid k\). Otherwise normalize on inertia and kill the restriction to \(\mathbb Q(\sqrt\ell)\) by Minkowski and divisor counting, using the range from \ref{B6.28a} and parity. The complete proof is in \S\ref{prf:B6.36}.
\end{proof}

\section{Summing back to the field}

\begin{lemma}[Descent to character lengths]\label{B6.36a}
Let $E$ be a CM field of the host with $n=[E:\Q]$. Assume $O$ contains the $\varphi(N)$-th roots of unity. For odd $\ell$ and $k\ge1$,
\begin{equation}\tag{6.R}\label{eq:b6-redall}
 e_Os^+_{\ell,k}(E)\le
 \sum_{\theta\in X_{\ell,O}(k)}
 \bigl(L_{\ell,k,\theta}+e_O\delta_\ell(N,n,k)\bigr),
\end{equation}
where $\omega(N)$ counts the distinct prime divisors of $N$ and
\begin{equation}\tag{6.L}\label{eq:b6-loss}
 \delta_\ell(N,n,k)=(n+n^2)\nu_\ell(N,k)+
                 \omega(N)(n+1)v_\ell(n)\le(Nk)^6 .
\end{equation}
If $\ell\nmid n$, one may take $\delta_\ell=0$.
\end{lemma}
\begin{proof}[Outline]
Factor the generator operators to bound the plus length by eigenspace lengths and twist to invariants. Bound restriction errors by transgression, \ref{B6.32a}, \(n\)-torsion and Frobenius cohomology, with zero cost for added conditions if \(\ell\nmid n\). The complete proof is in \S\ref{prf:B6.36a}.
\end{proof}

\begin{theorem}[The primes in the fixed set]\label{B6.37}
Fix $\chi$, $\ell\in\mathcal B_\chi$ and $j\in\N$. There is $C_{\ell,j}\in\N$ depending on these data. Let the host be large for $t\ge2$ and let $B\ge0$ satisfy the layer hypothesis of \ref{B6.32}. For every corresponding CM field $E$, the lengths for $1\le k<t$ are finite and
\[
 \log\ell\sum_{1\le k<t}s^+_{\ell,k}(E)\le(Nt)^{C_{\ell,j}}+C_{\ell,j}B .
\]
\end{theorem}
\begin{proof}[Outline]
Choose \(O\) by \ref{B6.36a}; characters of odd \(\operatorname{ord}\theta_P\) have conductor dividing \(f\) by \ref{B6.28a} and are bounded by \ref{B6.32}. The rest are bounded by \ref{B6.31}, \ref{B6.34}, and \ref{B6.35}, and summing over \(1\le k<t\) determines \(C_{\ell,j}\). The complete proof is in \S\ref{prf:B6.37}.
\end{proof}

\begin{theorem}[The primes outside the fixed set]\label{B6.38}
For each $\chi$ there are $C_{\mathrm{out}}\in\N$ and primes $r_1,r_2$. Let the host be large for $t\ge2$ with $q,d>r_1,r_2$. Let $E$ be a corresponding CM field and $S$ a finite set of odd primes disjoint from $\mathcal B_\chi$. Then the lengths for $\ell\in S$ and $1\le k<t$ are finite and
\[
 P^+(E,t;S)\le(Nt)^{C_{\mathrm{out}}}.
\]
\end{theorem}
\begin{proof}[Outline]
By \ref{B6.28a}, \ref{B6.36a} applies with zero loss and every non-class-I prime satisfies \(\ell\le 2t-1\). Summing class I via \ref{B6.30}, eliminating class Z by \ref{B6.36}, and bounding classes II and III by \ref{B6.31} yields a fixed power of \(Nt\). The complete proof is in \S\ref{prf:B6.38}.
\end{proof}

\begin{theorem}[The pointwise bound and its common window]\label{B6.39}
For fixed $\chi$ and $j:\N\to\N$ there are $C\in\N$ and primes $r_1,r_2$. Let the host be large for $t\ge2$ with $q,d>r_1,r_2$. Suppose $B\ge0$ satisfies the layer hypothesis of \ref{B6.32} with index $j(\ell)$ at every $\ell\in\mathcal B_\chi$. Then for every corresponding CM field $E$, odd prime $\ell$ and $1\le k<t$, the plus length is finite. For each finite set $S$ of odd primes,
\[
 P^+(E,t;S)\le(Nt)^C+CB .
 \tag{6.P}\label{eq:b6-point}
\]

Fix tower data $K=(R_0,R_\ell,C_K)$ with $C_K>0$. Use the field datum $M$ cut out by $\chi$ and $A_\ell=\Q(\mu_{\ell f})$ on $\mathcal B_\chi$. Set $j(\ell)=\lfloor\log_\ell R_\ell(A_\ell)\rfloor+3$ there, with the convention of \ref{B6.27}, and extend $j$ to $\N$. Take the corresponding $C,r_1,r_2$. For $C_t>0$, $0<\varepsilon<1/52$ and $A_1\in\R$, put
\[
 R=R_0(M)+\sum_{\ell\in\mathcal B_\chi}R_\ell(A_\ell),
 \qquad J=2^{R+2}(2C_KC_t/\varepsilon)^R,
 \qquad D_x=\exp(x^{4\varepsilon/5}).
\]
For the given $a,t,q,d,X$, a $K$-window is an interval $[Y,Y^2]$ with $X^{1/J}\le Y$ and $Y^2\le X$ on which \eqref{eq:B6-two-sum} holds for the corresponding fields and \eqref{eq:B6-odd-layer-sum} holds for all realizations of $A_\ell$ and their degree-$\ell^{j(\ell)}$ layers at each $\ell\in\mathcal B_\chi$.

There is $X_0$ such that for every $a,t$ admissible in \eqref{eq:B6-admitted}, every $X\ge X_0$, and every prime pair of \ref{B6.28} satisfying
\[
 q,d>2X^{C_t}+1,\qquad q,d>\max\{f,o,p_K,h_*,p_1,r_1,r_2\},
 \qquad\log(qd)\le X^{\varepsilon/(4J)},
\]
the following holds on every $K$-window. For all $x$ in the window, all corresponding CM fields $E$ and all finite sets $S$ of odd primes, the plus lengths for $1\le k<t(x)$ are finite and
\[
 P^+(E,t(x);S)\le\bigl(\operatorname{cond}(E)t(x)\bigr)^C+CD_x .
 \tag{6.W}\label{eq:b6-pluswindow}
\]
For the tower data of \ref{B6.27} and $C_t\ge1$, such a window exists after increasing $X_0$. For $C_t=401$, the choice $t(x)=\lfloor x^{401}\rfloor$ satisfies the two bounds on $t$ in \eqref{eq:B6-admitted}.
\end{theorem}
\begin{proof}[Outline]
Combine \ref{B6.37} and \ref{B6.38} for odd primes and use singletons for finiteness. On a common window use \(B=D_x\) and \(\operatorname{cond}(E)=N\) in the pointwise bound, with \ref{B6.27} supplying the tower window and bound at \(2\), and \(t=\lfloor x^{401}\rfloor\). The complete proof is in \S\ref{prf:B6.39}.
\end{proof}




% ---- p6c5 ----
\chapter{Automorphic counts and deformation rigidity}\label{ch:vwin-automorphic}

Let \(\mathscr H\) denote the inputs of \ref{B6.05} and \ref{B6.06}, with duality as in \ref{B6.09}. Put \(G_F=\Gal(\overline F/F)\), \(C_\ell=\overline{\Q}_\ell\), \(v_\ell(\ell)=1\), and let \(\varepsilon_\ell\) be the cyclotomic character. Coefficient DVRs \(O\subset C_\ell\) are finite over \(\Z_\ell\). Use the Selmer group \(\mathcal S_{\ell,k,\vartheta}\), length \(L_{\ell,k,\vartheta}\) and normalised length \(a_{\ell,k,\vartheta}\) of \ref{B6.29}.

A representation of level \((M;\mathbf m)\) is continuous and diagonal on inertia away from \(\ell\), with characters factoring through the action on \(\mu_M\). At each place above \(\ell\) it has a flag with decomposition-group characters \(\kappa_i\) satisfying \(\kappa_i|_I=\varepsilon_\ell^{m_i}\chi_i|_I\), where \(I\) is inertia and \(\chi_i\) factors through \(\mu_M\). The flag is unlinked if \(\kappa_i\ne\kappa_j\varepsilon_\ell\) for all \(i,j\). At each such place the flag must be unlinked or the representation potentially crystalline. Separated ordinary requires an unlinked flag with the prescribed inertia characters at every place above \(\ell\).

Fix \(\iota:C_\ell\simeq\C\). A place above \(\ell\) is unlinked if each prolongation admits an unlinked flag, and potentially unlinked if each admits a flag with \(\kappa_i^a\ne(\kappa_j\varepsilon_\ell)^a\) for all \(i,j\) and integers \(a\ge1\). \begin{definition}[Witness and automorphy conditions]\label{def:p6c5-witness}
Let \(r:G_F\to\GL_n(C_\ell)\) be a representation over a CM field \(F\). The witness condition at \(\iota\) requires every unlinked place above \(\ell\) to be potentially unlinked. It also requires that for each regular algebraic cuspidal representation \(\pi\) attached to \(r\) via \(\iota\), there be a single \(t\in\R\) such that the cuspidal support of \(\pi_v\) consists of characters at every finite \(v\nmid\ell\) and every potentially unlinked \(v\mid\ell\), with
\[
 q_v^{t-1/2}<|\xi(\varpi_v)|<q_v^{t+1/2}
\]
for each support character \(\xi\), where \(q_v\) is the residue cardinality and \(\varpi_v\) a uniformiser. A representation \(r\) of \(G_\Q\) satisfies the quadratic witness condition if \(r|_{G_K}\) satisfies the witness condition at \(\iota\) for every imaginary quadratic field \(K\) for which \(r|_{G_K}\) is irreducible. The quadratic automorphy condition at \(\iota\) requires that if \(r\) is continuous, then it satisfies the quadratic witness condition and each such restriction admits a regular algebraic cuspidal attachment via \(\iota\).
\end{definition}
 The attachment in \ref{B9.02} takes cuspidal input, as in \cite[Theorem~3.2.3]{chenevier-harris-2013}.

\section{Trace differences and the first count}

Fix an imaginary quadratic \(K\) of discriminant \(D_K\) and quadratic character \(\eta\). The family \(\mathcal A_3(K,\iota;k,N)\) consists of irreducible separated ordinary \(\rho:G_K\to\GL_3(C_\ell)\) of level \((N;(k,0,-k))\) with conjugate self-dual cuspidal attachment and witness for \(\rho\varepsilon_\ell^{-1}\). The family \(\mathcal A_{2+1}(K,\iota;k,N|D_K|)\) consists of representations isomorphic to \(W|_{G_K}\oplus1\), with the block in positions \(0,2\). Here \(W\) is irreducible over \(G_\Q\), satisfies the quadratic witness condition, has level \((N|D_K|;(k,-k))\) and has determinant of cyclotomic weight zero. The twist \(W\varepsilon_\ell^{-k}\) has a newform attachment of weight \(2k+1\).

\begin{definition}[The marked deformation problem and its local ring]\label{B6.40m}
Fix an odd prime \(\ell\), an integer \(k\ge1\), an even integer \(N\ge1\), an imaginary quadratic field \(K\) of discriminant \(D_K\) and class number \(h_K\), and a Dirichlet character \(\vartheta:(\Z/N)^\times\to O^\times\) with values in a finite DVR extension \(O\subset C_\ell\) of \(\Z_\ell\) having residue field \(\mathbf k=O/\mathfrak m_O\). Assume
\[
 \ell\nmid\operatorname{ord}\vartheta,\quad
 \operatorname{comp}_\ell(\vartheta)^2=1,\quad
 \vartheta(-1)=(-1)^k,\quad
 \ell\nmid |D_K|h_K,\quad
 (N,D_K)=1,\quad
 \ell-1\nmid2k,
\]
and, after inflating \(\vartheta\) and its components to \(G_\Q\) via the action on \(\mu_N\), \(\ker\overline\varepsilon_\ell\cap\ker\vartheta_P\not\subseteq G_K\), where \(\vartheta_P\) is the prime-to-\(\ell\) order part of \(\vartheta\) with its \(\ell\)-local factor removed. Let \(\eta:G_\Q\to\{\pm1\}\) have kernel \(G_K\), and put \(\Psi=\varepsilon_\ell^k\vartheta\).

Fix a complex conjugation \(c\in G_\Q\) (so \(c^2=1\), \(\eta(c)=-1\), and \(\Psi(c)=1\)) and an embedding \(j:\overline\Q\hookrightarrow\overline\Q_\ell\). This embedding selects a prolongation \(A_0\mid\ell\) and an injection \(\iota:G_{\Q_\ell}\hookrightarrow G_\Q\) with image \(D_\ell=D_{A_0}\) and \(\iota(I_{\Q_\ell})=I_\ell=I_{A_0}\). Let \(K_v\subset\overline\Q_\ell\) be the fixed field of \(\iota^{-1}(G_K)\) (so \(K_v=\Q_\ell\) if \(\ell\) splits in \(K\), and \(K_v/\Q_\ell\) is the unramified quadratic extension if \(\ell\) is inert in \(K\), since \(\ell\nmid D_K\)). Identify \(G_{\Q_\ell}\) with \(D_\ell\) via \(\iota\), and let \(K_1/K_v\) be the extension fixed by \(G_{K_v}\cap\ker\operatorname{comp}_\ell(\vartheta)\); since \(\operatorname{comp}_\ell(\vartheta)^2=1\), \(\ell\) is odd, and \(\vartheta\operatorname{comp}_\ell(\vartheta)^{-1}\) is unramified at \(\ell\), \(K_1/K_v\) is totally tamely ramified of degree \(1\) or \(2\) with \(I_{K_1}=\ker(\vartheta|_{I_\ell})\), and \(K_1/\Q_\ell\) is abelian with Galois group of exponent dividing \(2\).

Use zero-based matrix indices \(0,1,2\) and let \(J\in\GL_3(\Z)\) be the antidiagonal matrix \(J_{ab}=\delta_{a,2-b}\). For any commutative ring \(R\), define the group \(\mathcal G_3(R)=\GL_3(R)\rtimes\{\pm1\}\), where \(-1\) acts on \(\GL_3(R)\) by \(\theta(A)=J(A^{-1})^tJ^{-1}\). For \(x=(A,\epsilon)\in\mathcal G_3(R)\), write \(x_{\mathrm{mat}}=A\). Set \(\mathrm{wt}_1(i)=1-i\) and \(\mathrm{wt}(a,b)=\mathrm{wt}_1(a)-\mathrm{wt}_1(b)=b-a\). The base lift is the homomorphism
\[
 r_0:G_\Q\longrightarrow\mathcal G_3(O),\qquad
 r_0(g)=\bigl(\operatorname{diag}(\Psi(g),1,\Psi(g)^{-1}),\,\eta(g)\bigr),
\]
with residual representation \(\bar r_0:G_\Q\to\mathcal G_3(\mathbf k)\). A \emph{marking} is a pair \((c,g_1)\) with \(g_1\in G_K\) and \(\Psi(g_1)^2-1\in O^\times\).

For a complete Noetherian local \(O\)-algebra \(R\) whose structure map is local and induces an isomorphism \(O/\mathfrak m_O\simeq R/\mathfrak m_R\), a \emph{continuous lift} of \(\bar r_0\) over \(R\) is a homomorphism \(r:G_\Q\to\mathcal G_3(R)\) lifting \(\bar r_0\) whose matrix entries are \(\mathfrak m_R\)-adically continuous (so its sign coordinate is \(\eta\) since \(\ell\) is odd). The lift \(r\) is \emph{normalised} (or \emph{marked}) at \((c,g_1)\) if
\[
 r(c)_{\mathrm{mat}}=I_3,\qquad
 \bigl(r(g_1)_{\mathrm{mat}}-r(cg_1^{-1}c)_{\mathrm{mat}}\bigr)_{01}=0,\qquad
 \bigl(r(g_1)_{\mathrm{mat}}-r(cg_1^{-1}c)_{\mathrm{mat}}\bigr)_{10}=0.
\]
Let \(\mathcal K_{\mathrm{away}}\subset G_\Q\) consist of all elements \(g\in I_A\cap\ker\eta\cap\ker\vartheta\) as \(A\) ranges over all finite places of \(\overline\Q\) not above \(\ell\).

Let \(U\) be the universal unrestricted lifting ring for continuous lifts of \(\bar r_0|_{D_\ell}\) to \(\mathcal G_3(R)\), with universal lift \(r^{\mathrm{univ}}:D_\ell\to\mathcal G_3(U)\) and base point \(x_{0,\ell}:U\to O\) specialising \(r^{\mathrm{univ}}\) to \(r_0|_{D_\ell}\). Choose \(s_0\in I_\ell\cap\ker(\vartheta|_{D_\ell})\) such that \(\varepsilon_\ell(s_0)\) has infinite order, and put \(d_i=\Psi(s_0)^{1-i}=\varepsilon_\ell(s_0)^{k(1-i)}\in O^\times\) for \(i=0,1,2\). These three values are pairwise distinct. Set
\[
 P_i=\prod_{j\ne i}\bigl(r^{\mathrm{univ}}(s_0)_{\mathrm{mat}}-d_jI_3\bigr)\in M_3(U),\qquad
 c_i=\prod_{j\ne i}(d_i-d_j)\in O\setminus\{0\}.
\]
Let \(J_{\mathrm{ord}}\subset U\) be the closed ideal generated by the coefficients of \(\det(XI_3-r^{\mathrm{univ}}(s_0)_{\mathrm{mat}})-\prod_{i=0}^2(X-d_i)\), the entries of \(P_j r^{\mathrm{univ}}(g)_{\mathrm{mat}}P_i\) for \(0\le i<j\le2\) and \(g\in G_{K_v}\), and the entries of \(P_i(r^{\mathrm{univ}}(h)_{\mathrm{mat}}-\Psi(h)^{1-i}I_3)P_i\) for \(i=0,1,2\) and \(h\in I_\ell\). Fix an arithmetic Frobenius lift \(\varphi_v\in G_{K_v}\cap G_{K_1}\) (which exists because \(K_1/K_v\) is totally ramified), put \(t_i=\operatorname{tr}(P_i r^{\mathrm{univ}}(\varphi_v)_{\mathrm{mat}})\in U\), and define the cleared adjacent Frobenius differences
\[
 \delta_1=c_1t_0-\varepsilon_\ell(\varphi_v)c_0t_1\in U,\qquad
 \delta_2=c_2t_1-\varepsilon_\ell(\varphi_v)c_1t_2\in U.
\]
Define the saturated ideal
\[
 I_{\mathrm{sat}}=\{a\in U:\exists n\ge0,\ (\ell\delta_1\delta_2)^na\in\sqrt{J_{\mathrm{ord}}}\},
\]
and let \(I\subset U\) be the intersection of the minimal primes of \(I_{\mathrm{sat}}\) contained in \(\ker x_{0,\ell}\). The \emph{distinguished local ring} is \(B=U/I\). The \emph{marked global deformation ring} \(S\) is the universal complete Noetherian local \(O\)-algebra with residue field \(\mathbf k\) representing normalised continuous lifts \(r:G_\Q\to\mathcal G_3(R)\) of \(\bar r_0\) that vanish on \(\mathcal K_{\mathrm{away}}\) (meaning \(r(g)=1\) for \(g\in\mathcal K_{\mathrm{away}}\)) and whose restriction \(r|_{D_\ell}\) factors through \(B=U/I\).

The \emph{cohomological inputs for the marked problem} are local finiteness, local Tate duality, the local Euler characteristic formula, the local Bloch--Kato comparison, global finiteness, the global Euler characteristic formula, Poitou--Tate duality (\ref{B6.05} and \ref{B6.09}), and the bound of \ref{B6.15}.
\end{definition}

\begin{lemma}[Conjugation invariance and formal smoothness of the distinguished local ring]\label{B6.40n}
Assume the local cohomological inputs of Definition~\ref{B6.40m}, and keep the notation \(U,I,B=U/I,x_{0,\ell},K_v,K_1\) of that definition (for any \(k\ge1\), including \(k=1\), and whether \(\ell\) splits or is inert in \(K\)). Put \(E=\operatorname{Frac}O\).
\begin{enumerate}
\item The ideal \(I\subset U\) is prime and, for every complete Noetherian local \(O\)-algebra \(R\) with residue field \(\mathbf k\) as in Definition~\ref{B6.40m}, is invariant under conjugation in \(\mathcal G_3(R)\) by every \((h,1)\) with \(h\in\GL_3(R)\) congruent to \(I_3\bmod\mathfrak m_R\). Consequently \(B=U/I\) is a complete Noetherian local \(O\)-domain with a homogeneous power series presentation \(B=O[[X]]/(H)\) for the diagonal conjugation grading \(\mathrm{wt}(a,b)=b-a\), and for every \(Y\in M_3(O)\) the infinitesimal conjugation \((I_3+\epsilon Y,1)r_0|_{D_\ell}(I_3+\epsilon Y,1)^{-1}\) over \(O[\epsilon]/(\epsilon^2)\) defines an \(O\)-valued derivation on \(B\) at \(x_{0,\ell}\) killing \(H\).
\item The completion \(\widehat{B[1/\ell]}_{\ker x_{0,\ell}}\) is formally smooth over \(E\) of dimension \(12\) and coincides with the formal framed crystalline deformation ring \(R^{\mathrm{cris}}_{r_0}\) of \(r_0|_{D_\ell}:D_\ell\to\mathcal G_3(E)\) parameterising deformations whose restriction to \(G_{K_1}\) is crystalline. There is a graded \(O\)-algebra homomorphism \(\varphi:U\to E[[t_1,\ldots,t_{12}]]\) with \(\ker\varphi=I\), constant term \(x_{0,\ell}\), linear terms identifying the rational cotangent space of \(B\) at \(x_{0,\ell}\) with the span of the coordinates, and \(1,2,3,4,2\) coordinates in weights \(-2,-1,0,1,2\), respectively.
\item Every characteristic-zero point \(f:B\to C_\ell\) yields a representation \(r_f:D_\ell\to\mathcal G_3(C_\ell)\) that is crystalline on \(G_{K_1}\) and conjugate to an upper-triangular representation on \(G_{K_v}\) with inertia characters \((\Psi,1,\Psi^{-1})\). Every characteristic-zero tangent vector of \(B\) at \(f\) gives a cocycle in \(Z^1_f(D_\ell,\operatorname{ad}r_f)\).
\end{enumerate}
\end{lemma}
\begin{proof}[Outline]
Shows \(R^{\mathrm{cris}}_{r_0}\) is formally smooth of dimension \(12\) via Bloch--Kato and filtered \(\varphi\)-modules. Identifies \(I\cdot U^\wedge=I_{\mathrm{cris}}\) via \ref{B6.48h} and slope splitting, giving primeness, conjugation invariance, and crystallinity. The complete proof is in \S\ref{prf:B6.40n}.
\end{proof}

\begin{lemma}[Marked normal form and base vanishing]\label{B6.40p}
In the marked problem of Definition~\ref{B6.40m}, assume the cohomological inputs of that definition.
\begin{enumerate}
\item Every two-sided characteristic-zero point \(r:G_\Q\to\mathcal G_3(C_\ell)\) of \(S\) (meaning one for which neither all positive-weight nor all negative-weight matrix entries vanish on \(G_K\)) is either of Type~I (\(r|_{G_K}\) irreducible) or of Type~II (\(r|_{G_K}\) block diagonal on \(\{0,2\}\oplus\{1\}\) with irreducible corner on \(\{0,2\}\)). In Type~II, one has \(r(g)_{\mathrm{mat},11}=1\) and \(\det(\operatorname{corner}(r(g)_{\mathrm{mat}}))=1\) for all \(g\in G_K\), and defining \(W(g)=\operatorname{corner}(r(g)_{\mathrm{mat}})\) for \(g\in G_K\) and \(W(g)=\operatorname{corner}(r(g)_{\mathrm{mat}})\operatorname{diag}(-1,1)\) for \(g\notin G_K\) yields a continuous irreducible representation \(W:G_\Q\to\GL_2(C_\ell)\) extending the corner with \(\det W=\eta\) and \(\operatorname{tr}W(c)=0\).
\item Any two two-sided characteristic-zero points of \(S\) whose matrix restrictions to \(G_K\) are isomorphic lie in the same \(\mathbb G_m\)-orbit under diagonal conjugation by \(\operatorname{diag}(t,1,t^{-1})\).
\item At the base point \(r_0\), let \(V_j\) be the weight-\(j\) summand of \(\operatorname{ad}r_0\) over \(C_\ell\). Then \(H^1_f(\Q,V_j)=0\) for \(j=-2,-1,0\), and in particular \(H^1_f(\Q,C_\ell(\eta))=0\).
\end{enumerate}
\end{lemma}
\begin{proof}[Outline]
Uses the diagonal marking matrix at \(g_1\) to classify two-sided points, applying \ref{B6.40n} and class field theory in Type II and diagonal intertwiners for (2). Proves \(H^1_f(\Q,V_j)=0\) for \(j\le0\) via \ref{B6.31a} and class-group finiteness. The complete proof is in \S\ref{prf:B6.40p}.
\end{proof}

\begin{theorem}[A single-character trace bound]\label{B6.40}
Assume the cohomological inputs for the marked problem of Definition~\ref{B6.40m} and \ref{B6.15}, \ref{B6.43}, \ref{B6.52}, \ref{B6.58}. Assume adjoint rigidity and admissibility for Type I, irreducible on \(G_K\), and block adjoint rigidity for Type II, with an irreducible corner and a middle line on \(G_K\). Let \(\ell\) be an odd prime, \(k\ge1\), \(N\ge1\) even, and \(\vartheta:(\Z/N)^\times\to O^\times\) satisfy
\[
 \ell\nmid\operatorname{ord}\vartheta,\quad
 \operatorname{comp}_\ell(\vartheta)^2=1,\quad
 \vartheta(-1)=(-1)^k,\quad
 \ell\nmid |D_K|h_K,\quad (N,D_K)=1,\quad \ell-1\nmid2k .
\]
Here \(h_K\) is the class number and \(\operatorname{comp}_\ell\) is the modulus component at \(\ell\). Let \(\vartheta_P\) be the prime-to-\(\ell\) order part with its \(\ell\)-local component removed. Inflate the characters to \(G_\Q\) via the action on \(\mu_N\) and assume \(\ker\overline\varepsilon_\ell\cap\ker\vartheta_P\not\subseteq G_K\). Put \(\Psi=\varepsilon_\ell^k\vartheta\), fix \(\sigma\in G_K\), and set \(T=\Psi(\sigma)+1+\Psi(\sigma)^{-1}\).

If \(\operatorname{tr}\rho(\sigma)\ne T\) throughout \(\mathcal A_3\cup\mathcal A_{2+1}\), then \(L_{\ell,k,\vartheta}<\infty\), and there is a finite set \(X\subset\mathcal A_3\cup\mathcal A_{2+1}\) of pairwise nonisomorphic representations with
\[
 a_{\ell,k,\vartheta}\le 2\sum_{\rho\in X}v_\ell(\operatorname{tr}\rho(\sigma)-T).
 \tag{6.40.1}
\]
\end{theorem}

\begin{proof}[Outline]
Combines \ref{B6.40n} and \ref{B6.58} to present \(S\), counts obstructions by \ref{B6.09}, and bounds \(L_{\ell,k,\vartheta}\) via \ref{B6.40a}. Verifies finiteness (\ref{B6.52}) and tangent bounds (\ref{B6.40p}, \ref{B6.43}), then applies \ref{B6.58} and \ref{B6.29}. The complete proof is in \S\ref{prf:B6.40}.
\end{proof}

\begin{lemma}[Integral lift on the ambient algebra]\label{B6.40a}
In the marked problem of \ref{B6.40}, let \(S\) be the normalised global ring, \(x_0:S\to O\) the base point, \(q:A\to S\) the graded model map, and \(x=x_0q\). For \(0\ne b\in O\), \(\mathcal S_{\ell,k,\vartheta}[b]\) injects additively into the \(O/(b)\)-valued derivations at \(x\) on \(A\) that kill \(F_i\). Each image derivation lifts to an \(O\)-valued derivation at \(x\) supported in weight \(-1\).
\end{lemma}
\begin{proof}[Outline]
Normalizes each \(b\)-torsion Selmer cocycle at \(g_1\) to build a square-zero lift \((1+\epsilon tV)r_0\). Conjugation invariance in \ref{B6.40n} lifts it to an \(O\)-valued weight-\(-1\) derivation on \(A\) killing \(F_i\bmod b\), with \((1,0)\) entries recovering \(t\). The complete proof is in \S\ref{prf:B6.40a}.
\end{proof}

\begin{theorem}[An integer for three-dimensional trace differences]\label{B6.41}
Assume the automorphic inputs of \(\mathscr H\). For each imaginary quadratic \(K\) there are \(c,A\in\N\) as follows. Let \(k,N\ge1\), let \(p\ge7\) be prime with \(p\nmid N|D_K|\), and take \(w\mid p\) in \(K\). Put \(q=\mathrm Nw\). Let \((z_u)_{u\in\mathcal O}\) be a finite indexed family of roots of unity permuted by every automorphism of \(\C\). There is \(D\in\Z\setminus\{0\}\), chosen before \(\ell\), with
\[
 \log|D|\le c\,\#\mathcal O\,[N(k+1)]^A\log p .
 \tag{6.41.1}
\]
For every prime \(\ell\ne2,p\), every \(\iota\), every arithmetic Frobenius \(\sigma\) at \(w\), and every family of finite sets \(X_u\subset\mathcal A_3\cup\mathcal A_{2+1}\) of pairwise nonisomorphic representations, the differences
\[
 \Delta_{u,\rho}=\operatorname{tr}\rho(\sigma^{-1})
 -\iota^{-1}(q^{-k}z_u+1+q^kz_u^{-1})
\]
are nonzero and satisfy
\[
 \sum_u\sum_{\rho\in X_u}v_\ell(\Delta_{u,\rho})\le v_\ell(D).
 \tag{6.41.2}\label{B6.41-valuation}
\]
\end{theorem}
\begin{proof}[Outline]
Clear denominators in the differences of \ref{B6.41a}, use \ref{B6.41-traces} for nonvanishing and bound their integral Galois-invariant product by \ref{B6.41-labels}. Distinct factors and the unit denominator bound the selected valuations. The complete proof is in \S\ref{prf:B6.41-valuation}.
\end{proof}

\begin{lemma}[From Galois level to automorphic level and infinitesimal character]\label{B6.41b}
Assume the automorphic inputs of \(\mathscr H\). Let \(K\) be a CM field, \(n\ge2\), \(m=(m_1,\ldots,m_n)\in\Z^n\) strictly decreasing, \(M\ge1\), \(\ell\) an odd prime, and \(\iota:C_\ell\simeq\C\). Let \(r:G_K\to\GL_n(C_\ell)\) be a continuous irreducible representation of Galois level \((M;m)\) (away from \(\ell\), diagonal on inertia with characters factoring through the action on \(\mu_M\); at each \(w\mid\ell\), preserving a flag with inertia characters \(\varepsilon_\ell^{m_i}\chi_i\) where \(\chi_i\) factors through the action on \(\mu_M\)). Assume that \(r\) satisfies the witness condition of Definition~\ref{def:p6c5-witness} at \(\iota\) and that every place \(w\mid\ell\) of \(K\) is unlinked or potentially unlinked for \(r\). Let \(\pi\) be a regular algebraic cuspidal automorphic representation of \(\GL_n(\mathbb A_K)\) attached to \(r\) via \(\iota\) that is polarised, i.e.\ \(\pi^\vee\simeq\pi^c\otimes(\chi\circ N_{K/K^+}\circ\det)\) for an algebraic Hecke character \(\chi\) of \(K^+\) satisfying \(\chi_v(-1)=(-1)^n\) at every real place \(v\) of \(K^+\). Then:
\begin{enumerate}
\item At every complex embedding \(\varphi:K\hookrightarrow\C\), \(\pi\) has the infinitesimal character determined by the weights \(m\).
\item The finite part \(\pi_{\mathrm{fin}}\) has a nonzero vector fixed by the principal congruence subgroup \(K(M)=\ker(\GL_n(\widehat{\mathcal O}_K)\to\GL_n(\widehat{\mathcal O}_K/M\widehat{\mathcal O}_K))\).
\item At every finite place \(w\nmid\ell\) of \(K\) where \(\pi_w\) is spherical with Satake multiset \(B\), \(r\) is unramified at \(w\) and \(\iota\) carries the eigenvalue multiset of \(r(\sigma_w^{-1})\) to \(q_w^{(n-1)/2}B\).
\end{enumerate}
In particular:
\begin{enumerate}
\item[(a)] If \(\pi\) is conjugate self-dual (\(\pi^\vee\simeq\pi^c\)), then \(\pi\) is polarised by \(\eta_{K/K^+}^n\) and conclusions (1)--(3) hold.
\item[(b)] Let \(r_0:G_\Q\to\GL_n(C_\ell)\) be irreducible of Galois level \((M;m)\), satisfying the quadratic automorphy condition of Definition~\ref{def:p6c5-witness} at \(\iota\) and \(r_0(g)^tPr_0(g)=\theta(g)P\) for an invertible matrix \(P\) and a character \(\theta\) of cyclotomic weight \(\mathrm{wt}\in\Z\). Assume either that \(m_i\ne m_j+1\) for all \(i,j\), or that \(\ell\mid M\) and \(r_0\) is unlinked at \(\ell\). If \(K'\) is an imaginary quadratic field on which \(r=r_0|_{G_{K'}}\) remains irreducible, then \(r\) has Galois level \((M;m)\), every place of \(K'\) above \(\ell\) is unlinked or potentially unlinked for \(r\), and its cuspidal attachment \(\pi\) over \(K'\) is polarised and satisfies (1)--(3).
\end{enumerate}
\end{lemma}
\begin{proof}[Outline]
Definition~\ref{def:p6c5-witness} makes each \(\pi_w\) a zero-monodromy principal series, and local-global compatibility gives (1)--(3). Parts (a)--(b) verify potential unlinkedness and polarisation using \(\eta_{K/K^+}^n\) and strong multiplicity one. The complete proof is in \S\ref{prf:B6.41b}.
\end{proof}

\begin{lemma}[Hecke labels and their traces]\label{B6.41a}
In \ref{B6.41} the two families have finite, conjugation-stable label sets \(S_1,S_2\) with equivariant algebraic traces \(t(s)\). For each \(\ell,\iota,\sigma\), representations admit labels with \(t(s)=\iota(\operatorname{tr}\rho(\sigma^{-1}))\), distinct for distinct isomorphism classes in each family. With constants \(c_i,A_i\in\N\) depending on \(K\),
\[
 |t(s)|\le\sqrt q+2,\qquad q^kt(s)\text{ is integral},
 \tag{6.41.3}\label{B6.41-traces}
\]
\[
 \#S_1\le c_1[N(k+1)]^{A_1},\qquad
 \#S_2\le c_2[N|D_K|(k+1)]^{A_2}.
 \tag{6.41.4}\label{B6.41-labels}
\]
In rank \(n\ge2\), suppose the irreducible restrictions to the \(n^2+2\) auxiliary fields below have at most \(P(M,H)\) such CM labels per field at level \(M\) and weight bound \(H\). The rational label count is at most \(2^{n+1}(n^2+1)P(M,H)\).
\end{lemma}
\begin{proof}[Outline]
Lemma~\ref{B6.41b}(a) and Proposition~\ref{B6.06b} bound CM Hecke labels via level-\(3N\) cohomology and the Weyl dimension formula. Restriction to \(n^2+1\) quadratic fields via Lemma~\ref{B6.41b}(b) gives the rational count, and purity yields the trace bounds. The complete proof is in \S\ref{prf:B6.41-labels}.
\end{proof}

\begin{theorem}[One integer for a character orbit]\label{B6.42}
Assume \ref{B6.40} and \ref{B6.41}. For every number field \(K\) there are \(c_K,A_K\in\N\) as follows. Let \(k,N\ge1\), \(G=(\Z/N)^\times\), \(H\le G\), and \(p\nmid30N|D_K|\) prime. There is \(D\in\Z\setminus\{0\}\) with
\[
 \log|D|\le c_K\varphi([G:H])[N(k+1)]^{A_K}\log p.
 \tag{6.42.1}
\]
For every prime \(\ell\ne p\) and coefficient ring \(O\), suppose some \(O\)-valued character of kernel \(H\) satisfies the datum conditions of \ref{B6.40}. Then all characters in that \(O\)-valued orbit have finite length and
\[
 \sum_{\ker\vartheta=H}a_{\ell,k,\vartheta}\le2v_\ell(D).
 \tag{6.42.2}
\]
\end{theorem}
\begin{proof}[Outline]
Apply \ref{B6.41} to inverse character values, using a good character for all orbit values and primitive powers to preserve the datum. Match inverse Frobenius traces with \ref{B6.40}, sum by \ref{B6.41-valuation}, and use two auxiliary primes. The complete proof is in \S\ref{prf:B6.42}.
\end{proof}

\section{The two-dimensional vanishing input}

Write \(\overline\varepsilon_{\ell^r}\) for the reduction of \(\varepsilon_\ell\) modulo \(\ell^r\). For a prolongation \(A\) of the \(\ell\)-adic place, write \(D_A\) and \(I_A\) for its decomposition and inertia groups. Number the coordinates of two-dimensional column representations \(0,1\).

The normal form for \((k,S,A,W)\) requires an odd prime \(\ell\), \(k\ge1\), a continuous \(W:G_\Q\to\GL_2(C_\ell)\) with integral matrices over a finite coefficient ring, and a finite set \(S\) containing \(\ell\) and all ramified primes. Its determinant \(\delta\) is odd quadratic, and \(W\) is irreducible on \(\ker\varepsilon_\ell\cap\ker\delta\). On \(D_A\) it is upper triangular with diagonal characters \(\alpha_W,\beta\). For some \(h\ge1\),
\[
 \alpha_W(s)^h=\varepsilon_\ell(s)^{kh}\quad(s\in I_A),\qquad
 \alpha_W\ne\varepsilon_\ell\quad\text{on }D_A.
 \tag{6.44.1}\label{eq:p6c5-normal}
\]
A source for this normal form is a newform \(g\) of weight \(2k+1\) and level \(M_g\) prime to \(\ell\), with \(|\iota^{-1}(a_\ell(g))|_\ell=1\). It is attached to \(W\varepsilon_\ell^{-k}\delta^{-1}\xi\), where \(\xi\) factors through \(\mu_\ell\) and \(\beta\varepsilon_\ell^k\delta^{-1}\xi=1\) on \(I_A\). A \(W\)-cocycle \(y\) is unramified away from \(\ell\) if at each prime \(p\ne\ell\) its restriction to inertia at a prolongation of \(p\) is a coboundary. Its second coordinate is a coboundary on \(D_A\) precisely when
\[
 \exists b\in C_\ell,\quad
 y(s)_1=(\beta(s)-1)b\quad(s\in D_A).
 \tag{6.44.2}\label{eq:p6c5-ordinary}
\]

The spherical test at \((\iota,k)\) is imposed on every twist \(W\phi\) by a character factoring through \(\Gal(\Q(\mu_M)/\Q)\) for some \(M\ge1\). If \(W\phi\) has a \(D_A\)-stable flag at a prolongation \(A\mid\ell\) whose successive inertia characters are \(\varepsilon_\ell^k,\varepsilon_\ell^{-k}\), in that order, then the component at \(\ell\) of every regular algebraic cuspidal representation attached to \(W\phi\) via \(\iota\) has a nonzero \(\GL_2(\Z_\ell)\)-fixed vector.

Let \(L\) be a number field and let \(U\) be an \(\ell\)-adic representation of \(G_L\). For a finite place \(v\), choose a decomposition group \(G_{L_v}\) with inertia subgroup \(I_v\). We use the local conditions
\[
 H_f^1(L_v,U)=
 \begin{cases}
 \ker\bigl(H^1(G_{L_v},U)\longrightarrow H^1(I_v,U)\bigr),
       &v\nmid\ell,\\[2pt]
 \ker\bigl(H^1(G_{L_v},U)\longrightarrow
           H^1(G_{L_v},B_{\mathrm{cris}}\otimes_{\Q_\ell}U)\bigr),
       &v\mid\ell .
 \end{cases}
\]
Here \(B_{\mathrm{cris}}\) is Fontaine's crystalline period ring, with the diagonal Galois action on the tensor product. At an archimedean place the local condition is zero. Write \(H_f^1(L,U)\) for the classes in continuous \(H^1(G_L,U)\) whose localizations belong to these subspaces. Conjugating a decomposition group transports its local condition.

\begin{lemma}[Kato vanishing in odd weight and quadratic descent]\label{B6.43b}
Let \(\ell\) be odd, let \(e\ge1\) be an integer, and put \(h=2e+1\). Fix \(\iota:\overline{\Q}_\ell\simeq\C\). Let \(W\) be a continuous irreducible two-dimensional representation of \(G_\Q\) over \(\overline{\Q}_\ell\). Suppose \(W(-e)=W\otimes\varepsilon_\ell^{-e}\) has a weight-\(h\) newform attachment via \(\iota\), in the following sense.

There are an integer \(M\ge1\), a Dirichlet character \(\omega\bmod M\), and a cusp form \(f=\sum_{n\ge1}a_nq^n\) on \(\Gamma_1(M)\) of weight \(h\) and character \(\omega\). It satisfies \(a_1=1\) and \(T_pf=a_pf\) at every prime \(p\nmid M\). At a proper divisor of \(M\) there is no nonzero cusp form of weight \(h\) and any nebentype with these eigenvalues at all but finitely many primes. At all but finitely many primes \(p\nmid M\ell\), \(W(-e)\) is unramified and, for arithmetic Frobenius \(\sigma_p\),
\[
 \iota\!\left(\det\bigl(X-W(-e)(\sigma_p^{-1})\bigr)\right)
       =X^2-a_pX+\omega(p)p^{2e},
\]
where \(\iota\) acts on coefficients.

Let \(K/\Q\) be imaginary quadratic and let \(\theta\) be its quadratic character. Then
\[
 H_f^1(\Q,W)=H_f^1(\Q,W\otimes\theta)=0 .
\]
\end{lemma}
\begin{proof}[Outline]
Identify the primitive newform and match \(W\) with Kato's twist by Frobenius comparison. Kato's finiteness over \(K\) gives rational Selmer vanishing by divisibility and closedness of coboundaries, then restriction and corestriction give both descents. The complete proof is in \S\ref{prf:B6.43b}.
\end{proof}

\begin{theorem}[Kato vanishing for two quadratic companions]\label{B6.43}
Assume \(\mathscr H\). Let \(\ell\) be odd, \(k\ge1\), \(\ell-1\nmid2k\), and \(K\) imaginary quadratic with character \(\theta\). Suppose \(\ell\nmid |D_K|h_K\), \(N\ge1\) is even, and \((N,D_K)=1\). Let \(\vartheta:G_\Q\to C_\ell^\times\) factor through \(\mu_N\), have order prime to \(\ell\), and have square trivial on inertia at \(\ell\). Put \(\Psi=\varepsilon_\ell^k\vartheta\) and assume \(\Psi(c)=1\) at complex conjugation. Let \(W:G_\Q\to\GL_2(C_\ell)\) have level \((N|D_K|;(k,-k))\), determinant \(\theta\), and satisfy
\[
 \operatorname{tr}W(s)=\Psi(s)+\Psi(s)^{-1}\quad(s\in I_A,\ A\mid\ell),
 \tag{6.43.1}
\]
\[
 |\operatorname{tr}W(g)-\theta(g)\Psi(g)-\Psi(g)^{-1}|_\ell<1\quad(g\in G_\Q).
 \tag{6.43.2}
\]
There are a prolongation \(A\mid\ell\) and an integer \(M\ge1\) such that all \(s\in I_A\) fixing \(\mu_M\) have characteristic polynomials \((X-\varepsilon_\ell(s)^k)(X-\varepsilon_\ell(s)^{-k})\). At some \(A\mid\ell\), \(\theta|_{I_A}=1\) and an ordered flag has inertia characters \(\varepsilon_\ell^k\chi,\varepsilon_\ell^{-k}\chi\), with \(\chi\) factoring through \(\mu_\ell\). Assume \(W|_{G_K\cap\ker\varepsilon_\ell}\) is irreducible. If \(W\varepsilon_\ell^{-k}\) has a weight-\((2k+1)\) newform attachment via \(\iota\) as in Lemma~\ref{B6.43b}, and \(W\) satisfies the quadratic witness condition at this \(\iota\), then
\[
 H_f^1(\Q,W)=H_f^1(\Q,W\otimes\theta)=0.
 \tag{6.43.3}
\]
\end{theorem}
\begin{proof}
Take \(e=k\) in Lemma~\ref{B6.43b}. The level hypothesis gives continuity of \(W\). A \(G_\Q\)-stable line would be stable under \(G_K\cap\ker\varepsilon_\ell\), contrary to the assumed irreducibility there. The given newform is attached to \(W(-k)\) via the fixed \(\iota\), and the character of \(K/\Q\) is \(\theta\). The two conclusions of the lemma give (6.43.3).
\end{proof}

\begin{lemma}[A source of level prime to \(\ell\)]\label{B6.43a}
Assume \(\mathscr H\), \(\ell\) odd, and \(k\ge1\). Let \(V\) be continuous of dimension two with odd quadratic determinant \(\delta\), irreducible on \(\ker\varepsilon_\ell\cap\ker\delta\). Away from \(\ell\) let its inertia characters be diagonal of bounded cyclotomic modulus. Suppose a finite cyclotomic twist of \(V\varepsilon_\ell^{-k}\) has a weight-\((2k+1)\) newform attachment. At \(A\mid\ell\) suppose \(\delta|_{I_A}=1\) and an ordered flag has inertia characters \(\varepsilon_\ell^k\chi,\varepsilon_\ell^{-k}\chi\), with \(\chi\) factoring through \(\mu_\ell\). If \(V\) passes the spherical test, it is isomorphic to a representation \(V'\) in normal form at \(A\) for some \(S\), with a source newform \(g\).
\end{lemma}
\begin{proof}[Outline]
Twist and use the spherical test and strong multiplicity one to get level prime to \(\ell\). The flag gives root valuations \(0,2k\), and the recurrence with \ref{eq:p6c5-coeff} excludes a torsion unit root, giving \ref{eq:p6c5-normal} in an adapted stable lattice. The complete proof is in \S\ref{prf:B6.43a}.
\end{proof}

\begin{lemma}[Vanishing under the spherical test]\label{B6.44}
Assume the cohomological inputs of \(\mathscr H\) and take the arithmetic hypotheses and the newform attachment of \ref{B6.43}. If \(W\) passes the spherical test, then \(H_f^1(\Q,W)=H_f^1(\Q,W\otimes\det W)=0\).
\end{lemma}
\begin{proof}
The arithmetic hypotheses give \(k\ge1\), continuity of \(W\), and irreducibility on \(G_K\cap\ker\varepsilon_\ell\), hence on \(G_\Q\). Apply Lemma~\ref{B6.43b} to the given attachment of \(W(-k)\). Its quadratic character is \(\det W\).
\end{proof}


\begin{lemma}[Coefficient bounds and nonvanishing at the edge]\label{B6.46d}
Let \(g'\) be a cuspidal Hecke eigenform of weight \(2k+1\), \(k\ge1\), and nebentype \(\omega\) at level \(N\), normalised by \(a_1=1\) with \(a_n=0\) when \((n,N)>1\). Then
\[
 \sum_{n\le X}|a_n|^2\ll X^{2k+1},\qquad |a_n|\ll n^{k+1/2}.
 \tag{6.46.16}\label{eq:p6c5-coeff}
\]
The same bounds hold before depletion. For every Dirichlet character \(\xi\), the twisted \(L\)-value with factors at the level and modulus omitted has a nonzero right limit at each real \(w\ge k+1\).
\end{lemma}
\begin{proof}[Outline]
Parseval and partial summation give \ref{eq:p6c5-square}, and the good Euler factors give \ref{eq:p6c5-positivity}. An edge zero bounds the right side, while odd nebentype and divergence of reciprocal primes in the class \(-1\bmod M\) make the left side unbounded. The complete proof is in \S\ref{prf:B6.46d}.
\end{proof}


\section{Four-dimensional symplectic points}

Throughout this section use the good data of Definition~\ref{B6.31b} and the lifts and points of Definition~\ref{B6.31c}.
In characteristic \(\ell\), write \(\varepsilon_\ell\) for the reduction of the cyclotomic character when it takes values in the coefficient field.

\begin{lemma}[Separation in the marked ordinary flag]\label{B6.48j}
At a characteristic-zero point of Definition~\ref{B6.31c}, the local action at \(\ell\) has an ordered flag with characters \(\lambda_i\). On inertia \(\lambda_i=\delta_i\), and on \(D_\ell\) their reductions agree. If \(U\le D_\ell\) is open and contains \(\sigma_\ell\), then
\(\lambda_i|_U\ne(\varepsilon_\ell\lambda_j)|_U\) for \(i<j\).
\end{lemma}
\begin{proof}[Outline]
The marked projectors of \ref{B6.31c} give the flag, and integrality over the local Henselian ring identifies the reductions. Test the cyclotomic twist by inertia exponents within each plane and by the marked Frobenius of \ref{B6.31b} across planes. The complete proof is in \S\ref{prf:B6.48j}.
\end{proof}

\begin{lemma}[Crystallinity of an ordered flag]\label{B6.48h}
Let \(K,E\) be finite extensions of \(\Q_\ell\), let \(A\) be a finite-dimensional commutative local \(E\)-algebra with residue field \(E\), and let \(r:G_K\to\GL_n(A)\) be continuous and upper triangular. Write \(\kappa_i\) for its diagonal characters and \(\bar\kappa_i\) for their images in \(E^\times\). Suppose that \(\kappa_i|_{I_K}=\varepsilon_\ell^{w_i}|_{I_K}\) in \(A^\times\), where \(w_1>\cdots>w_n\) are integers. Suppose also that \(\bar\kappa_i\ne\varepsilon_\ell\bar\kappa_j\) for \(i<j\). Then \(A^n\), regarded as an \(E\)-representation, is crystalline.
\end{lemma}
\begin{proof}[Outline]
Induct from crystalline twists of unramified lines. At each step, the weight order and character separation let the Bloch–Kato dimension formula, local duality and the Euler characteristic show that every extension class is crystalline. The complete proof is in \S\ref{prf:B6.48h}.
\end{proof}

\begin{lemma}[Cyclotomic restriction at a symplectic point]\label{B6.48k}
Let \(r\) be an irreducible characteristic-zero point of a good datum of case II or III. Take the fitted tower used for the ordinary comparison, with imaginary quadratic field \(K\), totally real abelian field \(F\), and \(L=FK\). Suppose \(r|_{G_K}\) has a cuspidal attachment with witness. Then \(r|_{G_{L(\mu_{\ell^\infty})}}\) is absolutely irreducible. The same conclusion holds after the common scalar twist used to put the representation on the level group.
\end{lemma}
\begin{proof}[Outline]
Distinct residual characters exclude nontrivial self-twists trivial on \(G_L\). The local witness and \ref{B6.48j} give simple spectrum on the cyclotomic kernel, so reducibility gives, as in \ref{B6.30aa}, a quadratic self-twist excluded by the residual marking. The complete proof is in \S\ref{prf:B6.48k}.
\end{proof}

\begin{theorem}[Infinitesimal rigidity of an automorphic symplectic point]\label{B6.48}
Assume the unramified inputs from \(\mathscr H\) and \ref{B6.53}. Let \(D\) be of case II or III, with finite characters \(\vartheta,\xi,\omega\) factoring through \(\mu_N\) for some \(N\ge1\). Fix \(\iota:C_\ell\simeq\C\) and an irreducible characteristic-zero point \(r\) of \(D\). Suppose that \(r|_{G_K}\) has a regular algebraic cuspidal attachment with witness for every quadratic CM field \(K\) for which this restriction is irreducible. Every allowed dual-number lift \(r'\) of \(r\) has
\[
 r'(g)=r(g)+\epsilon(Xr(g)-r(g)X),\qquad X\in\mathfrak{sp}_4(C_\ell).
 \tag{6.48.1}
\]
\end{theorem}
\begin{proof}[Outline]
Write the lift as a symplectic adjoint cocycle and verify the finite local condition by inertia averaging and \ref{B6.48j}, \ref{B6.48h}. Use \ref{B6.48k} and \ref{B6.53} to obtain a fixed point over the tower, then average and project to get a symplectic coboundary. The complete proof is in \S\ref{prf:B6.48}.
\end{proof}

\begin{lemma}[The spherical and Iwahori comparison maps]\label{B6.48i}
At finitely many split auxiliary places \(j\), let the local group be \(\GL_n\) and the residue cardinality \(Q_j\). Take spherical and Iwahori fixed vectors in the same representation, with common level elsewhere. Write \(M,V\) for the resulting modules and \(i:M\hookrightarrow V\) for inclusion. Use the unnormalised coset sum \(T:V\to M\). Let \(u_{ji}\) denote the commuting Bernstein root operators. Suppose their elementary symmetric functions on \(M\) act by the signed coefficients of monic degree-\(n\) polynomials \(P_j\) over the spherical coefficient ring \(B\). Let \(D\) be the ordered splitting algebra over \(B\) of those polynomials, with roots \(x_{ji}\). Then there are maps
\[
 D\otimes_B M\ \underset{\Psi}{\overset{\Phi}{\rightleftarrows}}\ V
\]
intertwining root polynomials and the common coefficient operators. Each composite is the action on its source of
\[
 f(x)=\prod_j\prod_{a\ne b}(x_{ja}-x_{jb})(x_{ja}-Q_jx_{jb}).
\]
The assertion also holds on matching direct summands cut out by the ordinary and ordered residual idempotents.
\end{lemma}
\begin{proof}[Outline]
Define the forward map by root operators and the reverse map by coset sums. Use the Bernstein relation and an alternating polynomial identity to compute both composites, then scale the reverse map to obtain the stated multiplier. The complete proof is in \S\ref{prf:B6.48i}.
\end{proof}

\begin{lemma}[Recovering the ordered flag at a separated weight]\label{B6.48ca}
Let \(K/\Q_\ell\) be finite, let \(W_K\) be its Weil group, and use the reciprocity map \(a:W_K\to K^\times\) taking a lift of a uniformiser to an element of valuation one. Fix \(\iota:\overline{\Q}_\ell\simeq\C\). Let \(r\) be an \(n\)-dimensional potentially semistable representation of \(G_K\) over a finite extension of \(\Q_\ell\), allowing extension of the coefficient field. Suppose that its filtration jumps at each embedding \(\tau:K\hookrightarrow\overline{\Q}_\ell\) are
\[
 h_{0,\tau}<\cdots<h_{n-1,\tau}.
\]
Let \(\omega_i:K^\times\to\C^\times\) be smooth characters whose restrictions to \(\mathcal O_K^\times\) are pairwise distinct. Suppose the characteristic polynomial of the Weil action of \(r\) at each \(s\in W_K\), after applying \(\iota\), is
\[
 \prod_{i=0}^{n-1}\bigl(X-\omega_i(a(s))\bigr).
\]
Suppose there are continuous characters \(\kappa_i:G_K\to\overline{\Q}_\ell^\times\) satisfying
\[
 \kappa_i(s)=\iota^{-1}\!\bigl(\omega_i(a(s))\bigr)
                  \prod_\tau\tau(a(s))^{-h_{i,\tau}}
 \qquad(s\in W_K).
\]
Then \(r\) has a \(G_K\)-stable complete flag whose successive characters, starting with the subline, are \(\kappa_0,\ldots,\kappa_{n-1}\).
\end{lemma}
\begin{proof}[Outline]
Reciprocity as in \ref{B9.06} computes the rank-one parameters, and distinct inertia characters split the Weil module into lines. After twisting by the inverse of the first character, weak admissibility gives a trivial subline and induction gives the ordered flag. The complete proof is in \S\ref{prf:B6.48ca}.
\end{proof}

\begin{lemma}[The ordinary projector and finite-free module]\label{B6.48cb}
Let \(L/F\) be CM, with \(F=L^+\), unramified at finite places. Let \(\ell\) be odd and \(n\ge1\), and choose a prolongation at each finite place. Fix a distinguished place \(v_0\) of norm an odd prime \(p\ne\ell\) with \(p\nmid D_L\). Let the finite set of tame type places be disjoint from the places over \(\ell\) and omit \(v_0\). Assume that the places over \(\ell\), the tame type places and \(v_0\) split in \(L\), and let \(S\) be their union. Fix a coefficient DVR \(O\), residual eigenvalues at the good spherical and ordinary torus operators, and ordered residual local characters \(\psi_{vi}\), \(0\le i<n\), at each \(v\mid\ell\).

In the definite unitary group over \(F\), let the away level \(K\) have principal congruence level one at \(v_0\), spherical level outside \(S\), and a tame type \(\tau\) on a free coefficient lattice \(O^b\) whose inertial multiset at each tame type place consists of the Teichmüller lifts of the ordered residual characters. At each \(v\mid\ell\), let the \(i\)-th compact diagonal coordinate, \(0\le i<n\), have character \(\theta=(\psi_{vi}\varepsilon_\ell^i)^{-1}\) on the reciprocity lift over the weight ring \(A=\Lambda=O[[X_1,\ldots,X_{n[F:\Q]}]]\), where each variable is a universal pro-\(\ell\) inertia value times a fixed unit, minus one.

More generally, let \(A'\) be a complete Noetherian local \(O\)-algebra and \(\theta'\) a continuous character of the compact diagonal torus at \(\ell\) with values in \((A')^\times\). Consider \((A')^b\)-valued functions invariant on the left by rational points, transforming on the right by \(\tau^{-1}\) on \(K\) and \((\theta')^{-1}\) on the compact diagonal torus, invariant under the upper integral unipotents \(N_+\) at \(\ell\), and fixed modulo each precision by some lower unipotent congruence subgroup. For \(z=(\operatorname{diag}(\varpi_v^{n-1},\ldots,1))_{v\mid\ell}\) put
\[
 U_z f(g)=\sum_{u\in N_+/(zN_+z^{-1})} f(guz),
\]
and let \(Y(A',\theta')\) consist of the functions fixed modulo each precision by a positive power of \(U_z\), writing \(Y=Y(A,\theta)\).
\begin{enumerate}
\item Stabilizers at level \(K\) act trivially on the coefficient fibres. Factorial powers of \(U_z\) converge modulo every precision to a continuous \(A'\)-linear ordinary projector \(e\) onto \(Y(A',\theta')\).
\item Let \(i:Y(A',\theta')\to(A')^t\) evaluate at double-coset representatives at precision one, with coefficient coordinates included, and let \(\sigma:(A')^t\to Y(A',\theta')\) extend values by right translation and apply \(e\). Write \(V_j\) for the functions with all values in \(\mathfrak m_{A'}^j(A')^b\). Then \(Y(A',\theta')\cap V_j=\mathfrak m_{A'}^j Y(A',\theta')\) for all \(j\ge0\), \(\sigma i\) is congruent to the identity modulo \(\mathfrak m_{A'}Y(A',\theta')\) and is an automorphism, and \(Y(A',\theta')\) is finite free over \(A'\).
\item Let \(T_r\) run through the good spherical operators and the ordinary torus operators given at \(\ell\) by putting \(\varpi_v\) on the first \(j\) diagonal entries, \(1\le j\le n\). For lifts \(a_r\) of the specified residual eigenvalues, the submodule
\[
 M(K)=\{f\in Y:\ \forall r,a\ \exists h\ge1,\ (T_r-a_r)^h f\in\mathfrak m_A^a Y\}
\]
is a finite free \(A\)-module and an \(A\)-direct summand of \(Y\).
\end{enumerate}
\end{lemma}
\begin{proof}[Outline]
Principal level one at \(v_0\) trivializes fibre stabilizers, and factorial powers of \(U_z\) converge to the projector \(e\). Precision-one evaluation and Nakayama's lemma make \(Y(A',\theta')\) finite free, and lifted Hecke idempotents split off \(M(K)\). The complete proof is in \S\ref{prf:B6.48cb}.
\end{proof}

\begin{proposition}[The ordinary module at the three auxiliary levels]\label{B6.48c}
Assume \(\mathscr H\). Let \(L/F\) be CM, with \(F=L^+\), unramified at finite places. Let \(\ell\) be odd and \(n\ge1\), and choose a prolongation at each finite place. Fix a distinguished place \(v_0\) of norm an odd prime \(p\ne\ell\) with \(p\nmid D_L\). Let the finite set of tame type places be disjoint from the places over \(\ell\) and omit \(v_0\). Assume that the places over \(\ell\), the tame type places and \(v_0\) split in \(L\). Let \(S\) be their union and use the level group with \(S'=S\). Fix a coefficient DVR \(O\), pairwise distinct residual global characters with polarised product determinant, ordered residual local characters, and prescribed lifts on the \(\ell\)-primary torsion of the local abelianisations. Suppose an ordinary determinant pair with these data exists over a complete Noetherian local coefficient algebra with residue field \(k=O/\mathfrak m_O\).

The ordinary determinant pairs have a universal ring \(B\). In the definite unitary level model an ordinary Hecke module is finite free over the weight ring and has \(B\)-annihilator \(J\). The modules at the spherical and two Iwahori levels of a split auxiliary system satisfy \ref{B6.57}(1). For the systems of \ref{B6.56}, one may take \(E=2n!\) in the trace identity there.
\end{proposition}
\begin{proof}[Outline]
Builds \(B\) from determinant and character deformation rings, acting on \(M(K)\) of \ref{B6.48cb} via \ref{B6.04a} and \ref{B6.48ca} at separating weights. Sheet splitting, root-algebra nilpotence, and \ref{B6.48i} yield auxiliary-level freeness and trace identities. The complete proof is in \S\ref{prf:B6.48c}.
\end{proof}

\begin{lemma}[Control at an arithmetic weight]\label{B6.48d}
Use the definite group, tame coefficient lattice and level \(K\) of Lemma~\ref{B6.48cb}. Let \(A\) be its weight ring, \(\theta\) its compact diagonal character and \(Y\) its ordinary module before localization at the residual Hecke character. Let \(x:A\to\overline{\Q}_\ell\) be an arithmetic specialization. Write
\[
                 x\theta=\chi_\lambda\alpha
\]
on the compact diagonal torus, where \(\chi_\lambda\) is the lowest-weight character of the algebraic coefficient representation \(W_\lambda\), and \(\alpha\) has finite order. At each embedding, if \(h_0<\cdots<h_{n-1}\) are the ordinary weights, the exponents of \(\chi_\lambda\), in the order of the flag, are \(h_i-i\).

At sufficiently deep Iwahori level at \(\ell\), the classical ordinary space of weight \(W_\lambda\), tame type and level \(K\), and right diagonal character \(\alpha^{-1}\), is Hecke-equivariantly isomorphic to \(Y\otimes_{A,x}\overline{\Q}_\ell\). On the classical side the operator of a positive diagonal element \(a\) is its double-coset operator divided by \(\chi_\lambda(a)\). The isomorphism restricts to the residual summand \(M(K)\) of Lemma~\ref{B6.48cb}.
\end{lemma}
\begin{proof}[Outline]
Identifies \(O'\otimes_A Y\simeq Y'\) via \ref{B6.48cb} and Nakayama's lemma. Pairing with the lowest-weight functional \(p_0\) intertwines diagonal operators, giving injectivity by \(U_z\)-divisibility and surjectivity by lifting modulo \(\varpi\). The complete proof is in \S\ref{prf:B6.48d}.
\end{proof}

\begin{lemma}[Irreducible arithmetic points of the Hecke quotient]\label{B6.48a}
Assume \(\mathscr H\). In a definite unitary ordinary model of rank \(n\ge1\) over a CM field \(L\) with \(S'=S\), let \(J\) be the Hecke ideal in the universal ordinary determinant ring \(B\). Fix \(\iota\). Let \(x:B\to C_\ell\) kill \(J\) and have locally algebraic ordinary characters of strictly increasing weights at each embedding. Write \(D_x\) for the universal determinant law with its coefficients specialised by \(x\). Let \(p:G_L\to\Gamma_0\) be the quotient map to the CM subgroup of the level group. An irreducible representation \(r\) of \(G_L\) with determinant law \(p^*D_x\) has a regular algebraic conjugate self-dual cuspidal attachment.

Suppose \(r\) is continuous. Every unlinked place over \(\ell\) is potentially unlinked. For every regular algebraic cuspidal attachment \(\pi\) there is \(t\in\R\) such that at every finite place \(v\) of \(L\) away from the distinguished place of \(L^+\), with \(v\nmid\ell\) or \(r\) potentially unlinked at \(v\), the cuspidal support of \(\pi_v\) is toral and its characters \(\xi\) satisfy
\begin{equation}\tag{6.48.5}\label{eq:p6c5-toral-support}
 q_v^{t-1/2}<|\xi(\varpi_v)|<q_v^{t+1/2}.
\end{equation}
Here \(q_v\) is the residue cardinality and \(\varpi_v\) a uniformiser.
\end{lemma}
\begin{proof}[Outline]
Classical control \ref{B6.48d}, base change \ref{B6.04a} and Galois attachment \ref{B9.01} identify the cuspidal factor by Frobenius density. Read toral characters from Jacquet eigenvalues and use \ref{B9.38a} and \ref{B9.55a} to test their support and ratios. The complete proof is in \S\ref{prf:B6.48a}.
\end{proof}

\begin{lemma}[Coordinates for the lifting conditions]\label{B6.48ab}
Let \(F\) be a field, \(T\) a finite-dimensional \(F\)-space, \(g,D\ge0\), and \(U\subset T^g\) a subspace of codimension at most \(g\). Write \(\iota_j:T\to T^g\) for inclusion in slot \(j\) and \(W=\bigcap_j\iota_j^{-1}U\). If \(\dim W+g^2\le D\), there are \(D\) forms \(X_b\) spanning \(T^\vee\) and \(g^2\) independent vectors \(c_{ij}\in F^D\) with
\[
 \sum_{j,b}c_{ij,b}X_b(a_j)=0\qquad(a\in U,\ 1\le i\le g).
\]
\end{lemma}
\begin{proof}
Pad a basis of \(U^\perp\) with zeros to obtain \(g\) spanning forms \(L_i\). The forms \(L_i\iota_j\) have common kernel \(W\). Put these forms, with their repetitions and zeros, in the first \(g^2\) slots of \(X_b\). Fill the other slots with a basis of a complement of their span and zeros. Take \(c_{ij}\) to be the unit vector of its assigned slot. The row equation is \(L_i(a)=0\).
\end{proof}

\begin{lemma}[Dimension at an auxiliary-level augmentation]\label{B6.48ac}
Let \(L\to B\) be a local map of complete Noetherian local rings, with \(L\) regular of dimension \(d_0\). Let \(M\ne0\) be a finite \(B\)-module, let \(y_1,\ldots,y_h\in\mathfrak m_B\), and write \(Y=(y_1,\ldots,y_h)\). Suppose \(M/YM\) is finite free over \(L\) and multiplication gives an isomorphism
\[
 (M/YM)[T_1,\ldots,T_h]\ \longrightarrow\
 \bigoplus_{a\ge0}Y^aM/Y^{a+1}M,\qquad T_i\longmapsto [y_i].
\]
For each minimal prime \(Z\) of \(\operatorname{Supp}_B M\) one has
\(\dim B/Z=d_0+h\).
\end{lemma}
\begin{proof}[Outline]
Lift an augmentation basis using the graded isomorphism to make \(M\) free over a power series ring over \(L\). Each minimal support prime contracts to zero there, so integrality gives the dimension. The complete proof is in \S\ref{prf:B6.48ac}.
\end{proof}

\begin{lemma}[The component with a nonzero kernel vector]\label{B6.48ad}
Let \((P,\mathfrak m)\) be a regular local ring. Let \(g\ge1\), let \(Z=(z_{ij})\) be a \(g\)-by-\(g\) matrix and \(u=(u_i)\) a column over \(\mathfrak m\). Suppose their \(g^2+g\) entries have linearly independent images in \(\mathfrak m/\mathfrak m^2\). The radical of the ideal generated by the entries of \(Zu\) and by \(\det Z\) is prime. If a prime \(P_0\) contains \((Zu)\), omits some \(u_i\), and satisfies \(\dim P/P_0\ge\dim P-g\), then \(P_0\) is contained in every prime over \((Zu,\det Z)\).
\end{lemma}
\begin{proof}[Outline]
The projective incidence scheme of kernel lines containing \(u\) is regular with connected closed fibre, so its image is irreducible. The adjugate puts \(\det Z\) in \(P_0\), and the cut \(Z=\operatorname{diag}(u)\) gives the dimension bound forcing equality. The complete proof is in \S\ref{prf:B6.48ad}.
\end{proof}

\begin{lemma}[Comparison of two ordinary flags across a square-zero ideal]\label{B6.48aad}
Let \(R\) be a finite local algebra over an infinite field \(F\), with residue field \(F\), and let \(I\subset R\) satisfy \(I^2=0\). Let \(D\) be a group. For \(b=0,1\), let \(r_b:D\to\GL_n(R)\) and let \(\chi_{b,i}:D\to R^\times\), \(0\le i<n\), be characters. Suppose the representations and their respective characters agree modulo \(I\), and the residual characters are pairwise distinct. Suppose \(r_0\) is upper triangular in a basis \(w\), with ordered diagonal \(\chi_{0,i}\), and
\[
 \prod_{i=0}^{n-1}\bigl(r_1(g_i)-\chi_{1,i}(g_i)\bigr)=0
 \qquad(g_0,\ldots,g_{n-1}\in D),
\]
with the factors written in increasing order. Put
\[
 Z(g)=r_1(g)r_0(g)^{-1}-1,\qquad
 dY(g)=r_0(g)Yr_0(g)^{-1}-Y .
\]
There is \(Y\in M_n(I)\) such that \(Z-dY\) is upper triangular in the basis \(w\) and has \(i\)-th diagonal entry \(\chi_{1,i}/\chi_{0,i}-1\). If an \(F\)-subalgebra of \(M_n(R)\) contains the representation matrices and the scalar matrices of the character values for both systems, \(Y\) can be taken in that subalgebra.
\end{lemma}
\begin{proof}[Outline]
Use Lagrange projectors for a group-algebra element separating the residual characters, and use the ordered products to recover both flags. Conjugate by the sum of matching projector products and expand using \(I^2=0\). The complete proof is in \S\ref{prf:B6.48aad}.
\end{proof}

\begin{lemma}[The patched ordinary module over the auxiliary tower]\label{B6.48aae}
Assume \(\mathscr H\). Fix the definite ordinary unitary setting of Proposition~\ref{B6.48c} over a CM field \(L/L^+\) with \(d_+=[L^+:\Q]\), odd \(\ell\), \(n\ge1\), \(S'=S\), and a coefficient DVR \(O\) containing the images of \(\mathcal O_L\) under all \(\ell\)-adic embeddings. Fix a residual datum, its universal ordinary determinant ring \(B_0\), and the base ordinary module \(M_0\) of Proposition~\ref{B6.48c} localised at the residual system. Let \(s_1,\ldots,s_q\), \(q\ge0\), lie in the CM subgroup of the base level group, with characteristic polynomials split over the residue field of \(O\). Take a sequence of split auxiliary systems indexed by \(N\) as in Proposition~\ref{B6.56}, with auxiliary places \(w_{N,1},\ldots,w_{N,q}\) of norms \(Q_j=\mathrm Nw_{N,j}\equiv1\pmod{\ell^N}\), tame generators \(\tau_{N,j}\), and inverse Frobenii \(\phi_{N,j}\) converging to \(s_j\) in the finite quotients of the base group.

Divide each auxiliary level group by the kernel of the maximal pro-\(\ell\) quotient of the inverse image of the open base subgroup, take compatible surjections \(\pi_N:\Phi\to\Gamma_N\) from one finitely generated profinite group \(\Phi\), fix a nonprincipal ultrafilter on the levels \(N\), and let \(\Gamma_\infty=\Phi/K\) be the limit quotient defined by the ultrafilter, with limiting commuting elements \(\tau_j,\phi_j\) mapping to \(1,s_j\) at base level. Adjoin a power series variable \(s\), let \(B^{\rm sp}\) be the finite flat ordered-root splitting algebra over \(B_0[[s]]\) for the characteristic polynomials at \(s_1,\ldots,s_q\) with residual root lifts \(\widetilde u_{ji}\), and put \(d=nd_+\), \(\Lambda=O[[s,w_1,\ldots,w_d]]\), and \(S=\Lambda[[y_{ji}]]\) in the weight parameters \(w_b\) and \(nq\) variables \(y_{ji}\). Let \(R\) be the complete Noetherian local \(O\)-algebra for the continuous determinant on the CM subgroup of \(\Phi\), the ordered local characters with their prescribed torsion values, and the independent variables \(s,U_{ji}-\widetilde u_{ji},y_{ji}\), with augmentation \(\pi:R\to B^{\rm sp}\) sending \(y_{ji}\) to zero. Let \(\mathfrak a\subset R\) be generated by the polarisation relations and the traces, after left multiplication by group elements, of \([k]-1\) for \(k\in K\), the ordered ordinary products, the prescribed inertia differences, and
\[
 \prod_i([\phi_j]-U_{ji}),\qquad
 \prod_{b\ne i}([\phi_j]-U_{jb})([\tau_j]-1-y_{ji})
\]
for all lifts to \(\Phi\). Put \(Y=(y_{ji})\), \(\mathfrak l=\mathfrak m_\Lambda\), \(X=B^{\rm sp}\otimes_{B_0[[s]]}(M_0\otimes_O O[[s]])\) with \(Y\) acting as zero, and
\[
 f=\prod_{j=1}^q\prod_{i\ne b}(U_{ji}-U_{jb})^2\in R.
\]

At each level \(N\), with \(a_j=v_\ell(Q_j-1)\) and \(\Delta_N=\prod_{j,i}\Z/\ell^{a_j}\), let \(M_N\) be the ordinary residual Hecke and ordered-root summand of the deepest auxiliary module of Proposition~\ref{B6.48c} on which every \(U_{ji}-\widetilde u_{ji}\) is topologically nilpotent, after scalar extension by \(s\).
\begin{enumerate}
\item Each \(M_N\) is free over \(\Lambda[\Delta_N]=S/((1+y_{ji})^{\ell^{a_j}}-1:j,i)\). For fixed \(b,j\ge1\), the quotient \(M_N/(Y^b+\mathfrak l^j)M_N\) is free over \(S/(Y^b+\mathfrak l^j)\) at all sufficiently high \(N\).
\item For each \(j\ge1\), on a set of levels \(N\) in the ultrafilter there are \(R\)-linear maps between \(M_N/(Y+\mathfrak l^j)M_N\) and \(X/\mathfrak l^j X\) whose composites in both orders are multiplication by \(f\).
\item For \(b\ge1\) set
\[
 L_b=\prod_N M_N\big/\{(v_N):\
 \forall j,\ v_N\in(Y^b+\mathfrak l^j)M_N
 \text{ on a set in the ultrafilter}\},
\]
\(\mathcal B=\widehat{R[1/f]}^{\,Y}\), and \(\mathcal M=\varprojlim_b L_b[1/f]\). Then \(\mathcal M\) is a finite \(\mathcal B\)-module satisfying
\[
 \mathcal M/Y\mathcal M\simeq\mathcal B\otimes_R X,\qquad
 \operatorname{gr}_Y\mathcal M\simeq(\mathcal M/Y\mathcal M)[y_{ji}],
\]
and \(\mathfrak a\mathcal B\subseteq\sqrt{\operatorname{Ann}_{\mathcal B}\mathcal M}\).
\end{enumerate}
\end{lemma}
\begin{proof}[Outline]
Derives freeness of \(M_N\) from \ref{B6.48c} and compares \(M_N/(y)M_N\) with \(X\) via \ref{B6.48i}. Ultrafilter quotients \(L_b\), inversion of \(f\), and inverse limits produce \(\mathcal M\), its \(Y\)-graded isomorphism, and \(\mathfrak a\)-annihilation. The complete proof is in \S\ref{prf:B6.48aae}.
\end{proof}

\begin{lemma}[Generalisations of ordinary Hecke points]\label{B6.48aa}
Assume \(\mathscr H\). Fix the definite ordinary unitary setting over \(L/L^+\) with \(\ell\) odd, \(n\ge1\), \(S'=S\), \(\mu_\ell\not\subset L\), and a coefficient DVR \(O\) containing the images of \(\mathcal O_L\) under all \(\ell\)-adic embeddings. Fix a residual datum, its universal ordinary determinant ring \(B_0\), the weight algebra and \(\iota\). The Hecke ideal \(J_0\) of \ref{B6.48a} can be chosen to have the following properties.

Let \(H\) be the CM subgroup of the level group, \(c\) its chosen complex conjugation, and \(H'=H\cap\ker\varepsilon_\ell\), using the full \(\ell\)-adic character in this kernel. A curve consists of a finite extension \(k\) of the residue field of \(O\), a \(t\)-adically continuous \(\rho:H\to\GL_n(k[[t]])\), and an invertible symmetric polarising matrix \(P\) over \(k((t))\) with
\[
 \rho(s)^tP\rho(csc^{-1})=\varepsilon_\ell(s)^{1-n}P.
\]
The residual determinant is prescribed, and \(\rho\) kills the prescribed inertia kernels. At each selected \(v\mid\ell\) there is a flag over \(k((t))\) with ordered characters \(\psi_{v,i}\) in \(k[[t]]^\times\) of the prescribed reductions. The \(O\)-map \(f:B_0\to k[[t]]\) specialises the universal determinant and characters to these data. Require \(\psi_{v,i}\ne\psi_{v,j}\) for \(i\ne j\) and \(\psi_{v,i}\ne\varepsilon_\ell\psi_{v,j}\) for \(i<j\). Every open subset of \(H\) meeting \(H'\) must contain an \(s\in H'\) with separable characteristic polynomial. On the distinguished decomposition group \(D_0\), require
\[
 \{X\in M_n(k((t))):\ \varepsilon_\ell(s)\rho(s)X=X\rho(s)
 \text{ for all }s\in D_0\}=0.
\]

\textup{(a)} If \(\rho|_{H'}\) is generically absolutely irreducible and \(J_0\subseteq\ker f\), then every prime \(\mathfrak q\subseteq\ker f\) contains \(J_0\).

\textup{(b)} Let \(f_0,\ldots,f_h\), \(h\ge0\), be curves as above. Over the fraction field of each curve require a decomposition into two nonzero representations of \(H\), each polarised with multiplier \(\varepsilon_\ell^{1-n}\), absolutely irreducible and mutually nonisomorphic on \(H'\). A prime \(\mathfrak p\) is stable if the specialised determinant over \(\overline{\kappa(\mathfrak p)}\) is that of an irreducible representation.

Use the universal determinant \(D\) and ordered characters \(\psi^{\rm univ}_{v,i}\) in \(B_0[H]\). Take the Cayley--Hamilton relations given by the coefficients of \(\chi_D(z)(z)\) in the independent indeterminates \(T_j\) for every \(z=\sum_j T_jz_j\), \(z_j\in B_0[H]\). Adjoin the ordered products \(\prod_{i=1}^n([g_i]-\psi^{\rm univ}_{v,i}(g_i))\) at the ordinary groups, the differences \([s]-1\) on the prescribed inertia kernels, and their adjoints under \([s]^*=\varepsilon_\ell(s)^{1-n}[cs^{-1}c^{-1}]\). Let \(I\) be the two-sided ideal so generated and put
\[
 I^{\rm cl}=\bigcap_{b,V}
 \bigl(I+\mathfrak m_{B_0}^b B_0[H]+([v]-1:v\in V)\bigr),
\]
where \(b\ge1\) and \(V\) is open normal in \(H\). A nonstable prime \(\mathfrak p\) is thick if \(\overline{\kappa(\mathfrak p)}[H]/(I^{\rm cl})\) is nonsemisimple. Extend the ideal after taking the intersection. If consecutive curves share a thick prime below their kernels and a stable prime below \(\ker f_0\) contains \(J_0\), then every stable or thick prime below \(\ker f_h\) contains \(J_0\).
\end{lemma}
\begin{proof}[Outline]
Bounds auxiliary limit cohomology via \ref{B6.09} and presents each completed curve-prime ring using \ref{B6.48ab} and \ref{B6.48aad}. Patching \ref{B6.48c} via \ref{B6.48aae} and \ref{B6.48ac}, then comparing minimal primes by \ref{B6.48ad}, propagates \(J_0\). The complete proof is in \S\ref{prf:B6.48aa}.
\end{proof}

\begin{lemma}[Uniform finite-level bounds]\label{B6.48aaa}
Use the curve representation, its block decomposition and the ordinary conditions in the proof of \ref{B6.48aa}. Write \(A=k[[t]]\), \(F=k((t))\), \(A_m=A/t^m\), \(Q=\#k\), and \(d_+=[L^+:\Q]\). Write \(g\) for the dimension over \(F\) of ordinary odd adjoint classes. A single list of \(q\) auxiliary tests can be chosen for the given finite collection of curves. For each curve there are constants \(C_1,C_2\) such that, for every \(m\ge1\), at every sufficiently high level \(N\),
\[
 \#W_m\le Q^{mn(d_++q)+C_1},\qquad
 \#(E_m/t^e E_m)\le Q^{ge+C_2}\quad(1\le e\le m).
\]
Here the level threshold may depend on \(m\). Both constants are independent of that threshold.
\end{lemma}
\begin{proof}[Outline]
Choose trace tests detecting twisted dual classes modulo boundaries and realize them at auxiliary places via \ref{B6.56}. Smith normal form, Euler characteristics and finite duality in \ref{B6.05d} and \ref{B6.09} bound the ordinary classes and obstruction quotients. The complete proof is in \S\ref{prf:B6.48aaa}.
\end{proof}

\begin{lemma}[Ordinary classes on the limit group]\label{B6.48aab}
With the levels and common cover of \ref{B6.48aa}, the ordinary \(V_1(F)\)-classes on \(\Gamma_\infty\) have dimension at most \(D_*=n(d_++q)\). Given \(g+1\) tuples of \(g\) bounded ordinary even cocycles, some nonzero \(F\)-linear combination of the tuples is the tuple of primed components of a bounded ordinary \(V_2(F)\)-cocycle. Inflation identifies the ordinary odd cohomology over \(F\) on the base and the limit.
\end{lemma}
\begin{proof}[Outline]
Transfer cocycles and boundaries through common finite quotients, using \ref{B6.48aaa} and Smith normal form to bound classes and lift a nonzero combination through all depths. Odd cocycles vanish on the tame generators of \ref{B6.56} by disjoint block spectra and descend. The complete proof is in \S\ref{prf:B6.48aab}.
\end{proof}

\begin{lemma}[The completed module at a curve prime]\label{B6.48aac}
Use \(R,\Lambda,\mathfrak a,Y,f,X,\mathcal B,\mathcal M\) of Lemma~\ref{B6.48aae}. Let \(p\) be a prime of \(R\) with \(Y\subset p\), \(f\notin p\), and \(\Lambda/(p\cap\Lambda)=k[[t]]\) via \(s\mapsto t\) for a finite extension \(k\) of the residue field of \(O\). Put \(B'=\widehat{R_p}\), and write \(M'=B'\otimes_{\mathcal B}\mathcal M\). Multiplication induces
\[
 (B'\otimes_R X)[T_1,\ldots,T_{nq}]
       \ \simeq\ \operatorname{gr}_{Y}M'.
\]
Every minimal support prime after passage to \(C=B'/\mathfrak aB'\) has quotient dimension \(1+n(d_++q)\).
\end{lemma}
\begin{proof}[Outline]
Uses \ref{B6.48aae} along the ultrafilter to prove the graded isomorphism for \(\mathcal M\) and transfers it to \(M'\) by flat completion at \(p\). Applying \ref{B6.48ac} over \(\widehat{\Lambda_{p_\Lambda}}\) and passing to \(C\) gives dimension \(1+n(d_++q)\). The complete proof is in \S\ref{prf:B6.48aac}.
\end{proof}


\begin{lemma}[A curve avoiding the excluded loci]\label{B6.48bb}
Let \(R\) be a complete Noetherian local ring with finite residue field. Let \(I\) be an ideal, \(b\ge1\), and \(s\in\mathfrak m_R\), with \(\dim R/I\ge b+1\). Let \(A_i,B_j\) be two finite lists of ideals. Suppose \(\dim R/A_i<b\) for every \(i\), and \(\dim R/\mathfrak q<b\) for every prime \(\mathfrak q\supseteq B_j\) containing none of the \(A_i\). There is a prime \(\mathfrak p\supseteq I+(s)\) with \(\dim R/\mathfrak p=1\) containing none of the listed ideals. If \(R/(s)\) has characteristic \(\ell>0\), there is a local map \(R\to k'[[t]]\) with kernel \(\mathfrak p\), where \(k'\) is a finite field.
\end{lemma}
\begin{proof}[Outline]
Enlarge the exclusion list by minimal primes and take a component after cutting by \(s\). Repeated height-one cuts keep a product of elements from the excluded ideals nonzero, and in equal characteristic the final curve normalizes to \(k'[[t]]\). The complete proof is in \S\ref{prf:B6.48bb}.
\end{proof}

\begin{definition}[Four-dimensional frame cohomology]\label{B6.48be}
The \emph{four-dimensional frame cohomology} hypothesis is the conjunction of the following Galois cohomology inputs for finite discrete \(\ell\)-primary modules (with \(\ell\) odd): local finiteness, local Tate duality, the local Euler characteristic formula, unramified orthogonality, normalized Poitou--Tate duality (\ref{B6.09}), global finiteness, the global Euler characteristic formula, and the Greenberg--Wiles formula (\cite[Ch.~I, Theorems~2.1, 2.6, 2.8, 4.10, 5.1 and Corollaries~2.3, 4.15]{milne-2006}, \cite[Proposition~1.6]{wiles-1995}). These are contained in the standing input package \(\mathscr H\) of \ref{B6.05} and \ref{B6.09}.
\end{definition}

\begin{lemma}[Presentation bound for the four-dimensional framed ring]\label{B6.48bf}
Assume Definition~\ref{B6.48be}. Let \(T\) be a fitted deep tower of case II or III with totally real abelian field \(F\) of degree \(d=[F:\Q]\) and \(s\) places above \(\ell\), let \(\Lambda\subset\overline{\Q}_\ell\) be a coefficient DVR finite over \(\Z_\ell\) with residue field \(k_\Lambda=\Lambda/\mathfrak m_\Lambda\), let \(A\) be the matrix ring of \(T\) over \(\Lambda\), and let \(x:A\to k_\Lambda\) be a residual point satisfying the prescribed pairings, reduction, continuity, and either of the residual support shapes of Lemma~\ref{B6.48b}. Then the completed framed matrix ring \(R\) at \(x\) admits a presentation over \(\Lambda\) with \(t\) variables and \(r\) relations satisfying
\[
 t-r\ge 2d+2-2s .
\]
\end{lemma}
\begin{proof}[Outline]
Verify \(\operatorname{Hom}_\Gamma(\mathfrak{sp}_4(k'),k'(\bar\varepsilon_\ell))=0\) and apply the cone Euler characteristic of Definition~\ref{B6.48be} to get \(t_0-r_0\ge2d+2\) for \(R_0\). Fixing the local \(\ell\)-primary flag characters adds at most \(2s\) relations. The complete proof is in \S\ref{prf:B6.48bf}.
\end{proof}

\begin{lemma}[Minimal components through a residual point]\label{B6.48b}
Assume the four-dimensional frame cohomology of Definition~\ref{B6.48be}. Let \(T\) be a fitted deep tower of case II or III and \(F\) its totally real abelian field. Assume \(\vartheta|_{G_F}\) has order prime to \(\ell\). Choose a coefficient DVR \(O\subseteq\Lambda\subset\overline{\Q}_\ell\), finite over \(\Z_\ell\), and write \(k_\Lambda=\Lambda/\mathfrak m_\Lambda\). Let \(j:B\to A\) be the coefficient and matrix rings of the tower over \(\Lambda\), with their structure map. Let \(J\subset B\) be a potential modularity ideal for the tower with the following conditions. For some \(e\ge1\),
\begin{equation}\tag{6.48b.0}\label{eq:p6c5-potential-finite}
 \mathfrak m_B^e\subseteq J+\mathfrak m_\Lambda B+
 (\psi_{v,i}(s)-1:v\mid\ell,\ s\in I_v,\ 1\le i\le4).
\end{equation}
Here \(\psi_{v,i}\) are the universal flag characters. Put \(U_v=I_v\cap\ker\varepsilon_\ell\cap\ker\bar r_D\), with \(\bar r_D\) the residual diagonal representation. For every \(\Lambda\)-map \(y:A\to k'[[t]]\), \(k'\) finite, with \(yj(J)=0\), constant term killing \(\mathfrak m_B\), full generic matrix span and pairwise separation of its flag characters on each \(U_v\), require every prime of \(B\) below \(\ker(yj)\) to contain \(J\).

Put \(d=[F:\Q]\) and let \(s\) be the number of local flags. Suppose \(s\), the strict root cocycle dimensions and the global cocycle dimensions of weights \(0,\pm\psi\) are at most \(C\). For each odd root suppose the global and product-local cocycle dimensions are at most \(d+C\). Assume \(d\ge30(C+1)^2\).

Take a residual point \(x:A\to k_\Lambda\) satisfying the prescribed pairings, reduction and continuity. Use basis weights \((1,1),(0,1),(1,0),(0,0)\), and put \(\alpha=(1,0)\), \(\beta=(-1,1)\), \(\psi=(0,1)\), \(\theta=(1,1)\). Denote by \(E(w)\) the presence of an off-diagonal representation entry of weight \(w\). The local flag matrices \(n_v\) have ordered characters \(\bar\chi_i|_{G_v}\). The representation support is contained in \(\{\beta,-\alpha,-\theta,-\psi\}\) with flag support in \(\{-\alpha,-\psi,-\theta\}\), or the representation support is contained in \(\{-\alpha,-\beta,-\psi,-\theta\}\).

Let \(V\) mean support in \(\{\beta,-\alpha,-\theta\}\), flag support in \(\{-\alpha,-\psi,-\theta\}\) and \((n_v)_{20}\ne0\) for every flag. Let \(H\) mean support in \(\{-\beta,-\psi,-\theta\}\) and \((n_v)_{10}\ne0\) for every flag. Assume one of
\[
 \begin{gathered}
 E(-\psi)\land(E(\beta)\lor E(-\alpha)),\qquad
 E(-\alpha)\land E(-\beta),\qquad V\land E(-\alpha),\\
 V\land E(\beta)\land E(-\theta)
 \ \text{with }\bar\chi_0/\bar\chi_1\ne1\text{ on every }G_v,\qquad
 H\land(E(-\psi)\lor(E(-\beta)\land E(-\theta))).
 \end{gathered}
\]
If one minimal prime below \(\ker x\) contains \(JA\), every such prime does.
\end{lemma}
\begin{proof}[Outline]
In the completed local ring \(R\) at \(x\), Lemma~\ref{B6.48bf} and connectedness give adjacent minimal primes \(P,Q\) with \(\dim R/(P+Q)\ge b+1\). Bounding reducible and coincidence loci below \(b\), apply \ref{B6.48bb} to build a curve propagating \(JR\) from \(P\) to \(Q\). The complete proof is in \S\ref{prf:B6.48b}.
\end{proof}

\begin{lemma}[The reducible supports at the five residual points]\label{B6.48bc}
Use the hypotheses, basis and root notation of \ref{B6.48b}. In the completed matrix ring at \(x\), every reducible support locus has dimension at most \(1+d+20C\), with the possible exception, in the fifth alternative, of the support contained in \(\{\pm\beta,\pm\theta\}\). That exception forces equality of two flag characters on a whole local decomposition group.
\end{lemma}
\begin{proof}[Outline]
Use marked projectors to list proper supports and diagonal degeneration of the cases in \ref{B6.48b} to bound their tangent spaces by cocycle counts. In the exceptional support both blocks share the first flag eigenvalue, forcing a character equality. The complete proof is in \S\ref{prf:B6.48bc}.
\end{proof}

\begin{lemma}[Dimension at a coincidence of flag characters]\label{B6.48bd}
In \ref{B6.48b}, let \(q\) be a prime of \(B\) containing \(J\) and the uniformiser. If \(\psi_{v,i}=\psi_{v,j}\bmod q\) on \(U_v\) for some \(i\ne j\), then
\[
 \dim B/q\le 2d+1-d_v,\qquad d_v=d/s.
\]
Outside the small reducible loci of \ref{B6.48bc}, a prime in the completed matrix ring above \(q\) has quotient dimension at most \(\dim B/q+2+4s\).
\end{lemma}
\begin{proof}[Outline]
Use inertia generators and one cyclotomic value to count parameters after two flag characters coincide. Outside the small supports of \ref{B6.48bc}, traces bound the matrix parameters and lower unipotent coordinates bound the flag parameters. The complete proof is in \S\ref{prf:B6.48bd}.
\end{proof}


\begin{lemma}[A cuspidal Eisenstein congruence at an auxiliary prime]\label{B6.48ba}
Fix an odd prime \(\ell\) and \(\iota:C_\ell\simeq\C\). Let \(\phi\) be primitive of conductor \(D>1\) prime to \(\ell\). Let \(w\nmid D\) be prime with \(\ell\mid w-1\) and \(\phi(w)=1\), and let \(\theta\) have conductor \(w\) and order \(\ell\). Put \(H=Dw\) and \(\chi=\phi\theta\). For every sufficiently large \(h\) with \((-1)^h=\phi(-1)\) there is a normalised cuspidal eigenform of weight \(h\), level \(\Gamma_1(H)\) and character \(\chi\) satisfying
\[
 a_p\equiv1+\chi(p)p^{h-1}\pmod{\mathfrak m_{C_\ell}}
 \qquad(p\text{ prime}).
\]
Congruences are transported by \(\iota^{-1}\). The eigenform can be chosen so that conjugating its Fourier coefficients gives cusp forms of the conjugate characters.
\end{lemma}
\begin{proof}[Outline]
Average cusp Eisenstein series by the primitive character and use the geometric series for \ref{eq:p6c5-eis-constant}. Subtract an integral form to remove the constants, then lift the residual divisor-sum eigenvector by Deligne–Serre and normalise. The complete proof is in \S\ref{prf:B6.48ba}.
\end{proof}

\begin{lemma}[Finiteness from valuation points]\label{B6.49a}
Let \(O\) be a complete DVR with finite residue field and uniformiser \(\varpi\). Let \(R\) be a complete local \(O\)-algebra presented by finitely many coordinates \(z_i\in\mathfrak m_R\) and homogeneous relations, with the coordinates assigned weights in a free abelian group and \(O\) of weight zero. Suppose that every local \(O\)-map \(R\to A\) to a complete DVR with finite residue field that kills \(\varpi\) also kills every positive-degree monomial of weight zero. Then the weight-zero part of \(R\) is a finite \(O\)-module. For integer weights, the monomial condition holds if at every such point all zero-weight coordinates vanish and all positive-weight coordinates or all negative-weight coordinates vanish.
\end{lemma}
\begin{proof}[Outline]
A weight-zero monomial generator outside the uniformiser radical survives in a normalised curve quotient, contrary to the DVR hypothesis. Weight-zero projection gives nilpotence modulo the uniformiser, and successive approximation gives \(O\)-finiteness. The complete proof is in \S\ref{prf:B6.49a}.
\end{proof}

\begin{lemma}[Traces at two ordinary places]\label{B6.49b}
Let \(F\) be real quadratic, and let the odd prime \(\ell\) split as \(vv'\) in \(F\). Fix an identification \(\overline{\Q}_\ell\simeq\C\) and the resulting assignment of the two real embeddings to \(v,v'\). Fix a finite-order central character \(\omega\). At a finite place \(w\), write \(O_w\) for the local integers and take the subgroup of \(\GL_2(O_w)\) with bottom row congruent to \((0,1)\) modulo \(\mathfrak p_w^{c_w}\). Assume the nonnegative integers \(c_w\) have finite support, contain the conductor of \(\omega\) with its multiplicities, satisfy \(c_v=c_{v'}=1\), and satisfy \(c_u>0\) at some place \(u\) of norm two. Let \(h\ge4\), with \(\omega_\tau(-1)=(-1)^h\) at both real embeddings.

For \(t\ge0\), use the holomorphic cusp space of weight \((h+2t,h)\), central character \(\omega\), and the indicated finite level. At \(w\), let \(D_w\) be the sum of right translates for the double coset of \(\operatorname{diag}(\varpi_w,1)\), counting right cosets once. Define
\[
 T_w(t)=(\mathrm Nw)^{h/2-1}\ell^{\,t{\bf1}_{w=v}}D_w .
\]
For a word \(W\) in finite places, let \(T_W(t)\) be the ordered product and let \(b_w\) be the number of occurrences of \(w\). If \(r\ge1\), \(b_v,b_{v'}\ge r\), \(N\ge0\), and \((\ell-1)\ell^N\mid t\), both traces are \(\ell\)-integral and
\[
 \operatorname{tr}T_W(t)-\operatorname{tr}T_W(0)
 \in \ell^{\min(N,r)}\mathcal O_{\overline{\Q}_\ell}.
\]

Suppose the parallel space has dimension \(N>0\) and contains a joint eigenform with integral normalized eigenvalues \(s_w\), where \(s_v,s_{v'}\) are units. Then the mixed space contains a joint eigenform whose normalized eigenvalues reduce to the same \(s_w\) at every finite place.
\end{lemma}
\begin{proof}[Outline]
Use Jacquet–Langlands to express both traces by the same quaternionic coset sum and compare the weight factors using the unit roots. A polynomial in the Hecke operators then counts systems in the given residual class. The complete proof is in \S\ref{prf:B6.49b}.
\end{proof}

\begin{lemma}[The marked integral basis]\label{B6.49c}
Let \(D\) be a good datum of case II or III. Start with an integral representation \(r\) reducing to the diagonal of Definition~\ref{B6.31b}, with the prescribed similitude and inertia conditions, and an integral symplectic ordinary flag at \(\ell\) whose quotient characters have the prescribed reductions and inertia restrictions. Use a complete discrete valuation ring with finite residue field, allowing finite extensions. These data exist for every characteristic-zero point of Definition~\ref{B6.31c} and for the valuation specialisations of an allowed lift used in the finiteness argument.

Write \(A=\langle e_0,e_2\rangle\), \(B=\langle e_1,e_3\rangle\). Suppose the generic representation is absolutely irreducible, or preserves \(B\), has full matrix span on \(B\), and has no invariant complement to \(B\) over the algebraic closure. One can choose the order of the residual labels before choosing the initial tower. For its separating marking \(t\) there is an integral symplectic basis in which \(r(t)\) is diagonal, the reduction of \(r\) is diagonal, and the flag has a lower unitriangular symplectic representative \(n\). Every off-diagonal entry of \(n\) of weight different from \(-\psi\) belongs to the maximal ideal. In this basis the generic representation has entries of both weights \(\psi,-\psi\) and of at least one weight in \(\{\beta,-\alpha,-\theta\}\).
\end{lemma}
\begin{proof}[Outline]
Specialize the flag of \ref{B6.31c} by properness and order the labels using \ref{B6.31b}. Lift the marking's spectral projectors and factor the flag matrix using its unit minors, then use full span or nonsplitting to find the required entries. The complete proof is in \S\ref{prf:B6.49c}.
\end{proof}

\begin{lemma}[Choice of lattice over the rational base]\label{B6.49d}
Use the basis and marking of Lemma~\ref{B6.49c}. After a finite extension and diagonal conjugation by a symplectic similitude, write \(R(g)\) and \(N\) for the reductions of the representation and its lower unitriangular flag. Write \(d_i\) for the residual characters in the chosen order. Then \(R(t)\) is diagonal. The representation support lies in
\[
 S_P=\{\beta,-\alpha,-\theta,-\psi\},
\]
the flag support lies in \(S_f=\{-\alpha,-\psi,-\theta\}\), and \(E(\beta)\) holds. Either \(E(-\psi)\) holds, or the following conditions hold together:
\[
 \begin{gathered}
 \operatorname{supp}R\subseteq\{\beta,-\alpha,-\theta\},\qquad
 N_{20}\ne0,\qquad d_0|_{D_\ell}=d_2|_{D_\ell},\\
 \nexists\,m\quad R_{12}(s)=m\bigl(d_2(s)-d_1(s)\bigr)
                     \quad\text{for every }s\in D_\ell .
 \end{gathered}
\]
Here \(m\) ranges over the residue field, also after extending it.
\end{lemma}
\begin{proof}[Outline]
Maximize the two diagonal scaling exponents in order under the valuation constraints from \ref{B6.49c}. The tight inequalities and the coboundary condition in \ref{B6.31b}, applied to the flag equations, force the stated support and exclude the local scalar. The complete proof is in \S\ref{prf:B6.49d}.
\end{proof}

\begin{lemma}[The first point on the tower]\label{B6.49e}
Restrict the lattice of Lemma~\ref{B6.49d} to the totally real field \(F\) of the initial tower, or to a fitted enlargement. Transport its ordinary flag to the chosen prolongations over \(\ell\). Its reduction is a residual point of the tower and satisfies one of the following alternatives of Lemma~\ref{B6.48b}:
\[
 E(-\psi)\land E(\beta),\qquad V\land E(-\alpha).
\]
\end{lemma}
\begin{proof}[Outline]
Transport the flag of \ref{B6.49d} by triangular factorization and compute its \((2,0)\) entry. The separating marking keeps the relevant cocycles nonzero on restriction, and the excluded local coboundary forces an entry of weight \(-\alpha\) in the second case. The complete proof is in \S\ref{prf:B6.49e}.
\end{proof}

\begin{lemma}[The nonsplit companion of a block point]\label{B6.49f}
Let an integral pair as in Lemma~\ref{B6.49c} be block diagonal on \(A\oplus B\), with full generic matrix span on \(A\). There is an integral pair with the same characteristic polynomials at every group element and the same ordinary inertia characters, preserving \(B\), having full generic matrix span on \(B\), and having no invariant complement to \(B\).
\end{lemma}
\begin{proof}[Outline]
Finite Selmer counts and compactness give a nonzero integral cocycle in symplectic maps from \(A\) to \(B\) with the local flag condition. Twist in the square-zero subgroup and use the normalization at complex conjugation to rule out a splitting. The complete proof is in \S\ref{prf:B6.49f}.
\end{proof}

\begin{lemma}[Matrix entries in an open coset]\label{B6.49g}
Let \(G\) be profinite, and let \(r\) be a continuous representation over a finite extension of \(k((t))\), where \(k\) is finite. Compute matrix spans and Zariski closures over an algebraic closure of the coefficient field. Suppose \(r\) is absolutely irreducible of dimension \(n\in\{2,3,4\}\). Fix an eigenbasis of \(r(u)\) for an element \(u\in G\) whose eigenvalue ratios at distinct indices have infinite order. Suppose every open normal subgroup of \(G\) has noncommutative image.

In each open coset there is a matrix \(M\) for which all diagonal entries of \(M\) and \(M^{-1}\) are nonzero, with one possible alternative in dimension four. In that alternative the four lines divide into two pairs, the diagonal two-by-two blocks of \(M\) vanish, and every entry of both off-diagonal blocks is nonzero.
\end{lemma}
\begin{proof}[Outline]
Spectral projectors and irreducibility identify the span of the identity component as a product of full matrix algebras in equal-sized blocks. In each coset, Zariski density in a translate of the component gives simultaneous nonvanishing of the required entries. The complete proof is in \S\ref{prf:B6.49g}.
\end{proof}

\begin{lemma}[Rank-two Eisenstein lifts and imaginary quadratic twists]\label{B6.49h}
Assume \(\mathscr H\). Fix an odd prime \(\ell\) and \(\iota:C_\ell\simeq\C\).
\begin{enumerate}
\item Let \(\phi\) be a primitive Dirichlet character of conductor \(D>1\) prime to \(\ell\), let \(w\nmid\ell D\) be a prime with \(\ell\mid w-1\) and \(\phi(w)=1\), and let \(\theta\) be a Dirichlet character of conductor \(w\) and order \(\ell\). View Dirichlet characters as \(C_\ell^\times\)-valued characters of \(G_\Q\) via \(\iota^{-1}\). For every sufficiently large integer \(h\ge2\) with \((-1)^h=\phi(-1)\), put \(D_h=\varepsilon_\ell^{h-1}\phi\theta\) and \(\Phi_h=\varepsilon_\ell^{h-1}\phi\). Then Lemma~\ref{B6.48ba} and Propositions~\ref{B6.06d}(3) and~\ref{B6.06e} yield a normalized cuspidal newform of weight \(h\), level \(Dw\) and nebentype \(\phi\theta\), and a continuous irreducible representation \(\rho:G_\Q\to\GL_2(C_\ell)\) of determinant \(D_h\) with residual characteristic polynomial \((X-1)(X-\bar\Phi_h)\) on \(G_\Q\), inertia type \((1,\phi\theta)\) and unlinked principal-series or spherical cuspidal support away from \(\ell\), and an ordinary flag at \(\ell\) of inertia weights \((h-1,0)\) and decomposition-group residual characters \((\bar\Phi_h,1)\). If \(H\le G_\Q\) is open normal and \(\bar\Phi_h|_H\ne1\), then \(\rho|_H\) is irreducible; if \(\bar\Phi_h(s)^2\ne1\) for some \(s\in G_\Q\), then \(\rho\) is irreducible on every open subgroup of \(G_\Q\).
\item Let \(K/\Q\) be an imaginary quadratic field of discriminant \(D_K\) prime to \(\ell\), with odd quadratic character \(\eta_K\). There is an algebraic Hecke character of \(K\) of infinity type \((-1,0)\), unramified outside \(D_K\) and of prime-to-\(\ell\) inertia exponent at primes dividing \(D_K\), whose restriction to rational ideles is \(\eta_K|\cdot|^{-1}\). Its \(\ell\)-adic Galois character \(\kappa_0:G_K\to C_\ell^\times\) via \(\iota\) is unramified outside \(\ell D_K\), has prime-to-\(\ell\) inertia order away from \(\ell\), satisfies \(\kappa_0(g)\kappa_0(cgc^{-1})=\varepsilon_\ell(g)^{-1}\) for complex conjugation \(c\) and \(g\in G_K\), and has rational-idelic transfer \(\eta_K\varepsilon_\ell^{-1}\) from \(G_\Q\).
\end{enumerate}
\end{lemma}
\begin{proof}[Outline]
For (1), \ref{B6.48ba}, \ref{B6.06d}(3), and \ref{B6.06e} produce a newform and twisted Galois representation whose local data follow from congruences and inertia eigenlines. For (2), extending an idelic character of type \((-1,0)\) gives \(\kappa_0\). The complete proof is in \S\ref{prf:B6.49h}.
\end{proof}

\begin{lemma}[Twisted intertwiners and line permutations on a curve]\label{B6.49i}
Let \(A\) be a valuation domain with residue field \(k\), fraction field \(F\), and a fixed extension of its valuation to \(\overline F\). Let \(m\ge1\).
\begin{enumerate}
\item Let \(D_0\) be a group, \(\rho:D_0\to\GL_m(A)\) a representation, and \(\chi:D_0\to A^\times\) a character. Suppose there is \(s_0\in D_0\) whose residual characteristic polynomial \(\det(X-\overline{\rho(s_0)})\) has a single root in \(\overline k\) while \(\bar\chi(s_0)\ne1\). Then
\[
 \{X\in M_m(F):\ \chi(s)\rho(s)X=X\rho(s)\text{ for all }s\in D_0\}=0.
\]
\item Let \(H\) be a group and \(\rho:H\to\GL_m(A)\) a representation whose residual characteristic polynomial at every \(h\in H\) is \(\prod_{i=1}^m(X-\bar\chi_i(h))\) for characters \(\bar\chi_1,\ldots,\bar\chi_m:H\to\overline k^\times\). Suppose \(\rho\) permutes transitively a decomposition \(\overline F^m=L_1\oplus\cdots\oplus L_m\) into one-dimensional subspaces, and let \(P\le H\) be the kernel of this permutation action. Then the eigenvalues of \(\rho(p)\) on \(L_1,\ldots,L_m\) for \(p\in P\) are valuation units whose reductions all coincide, and \(\bar\chi_1|_P=\cdots=\bar\chi_m|_P\).
\end{enumerate}
\end{lemma}
\begin{proof}[Outline]
For (1), distinct residual roots make the characteristic polynomials of \(\chi(s_0)\rho(s_0)\) and \(\rho(s_0)\) coprime, so Cayley-Hamilton forces \(X=0\). For (2), trivial residual commutators equate reduced line eigenvalues across the transitive orbit. The complete proof is in \S\ref{prf:B6.49i}.
\end{proof}

\begin{theorem}[Finiteness and automorphy of symplectic points]\label{B6.49}
Assume \(\mathscr H\), take a host \(h\) large for \(t\) as in \ref{B6.28}, and let \(\ell\) be odd, \(0<k<t\), \(\vartheta\in X_\ell(k)\). For every good four-dimensional datum \(D\) of case II or III, the weight-zero algebra of each allowed lift is finite over \(O\). At every characteristic-zero block point \(r\),
\[
 \operatorname{span}_{C_\ell}\{r(g)_{\{0,2\},\{0,2\}}:g\in G_\Q\}
 \ne M_2(C_\ell).
 \tag{6.49.1}
\]
There is one \(\iota\) for which every irreducible characteristic-zero point has an attachment with witness on each irreducible quadratic CM restriction.
\end{theorem}
\begin{proof}[Outline]
Build towers and seeds in cases II (\ref{B6.49b}) and III (\ref{B6.49h}), linking residual points (\ref{B6.49c}--\ref{B6.49f}) by Selmer cocycles. Propagating \ref{B6.48aa} and \ref{B6.48b} yields finiteness (\ref{B6.49a}), automorphy (\ref{B6.48a}), and (6.49.1). The complete proof is in \S\ref{prf:B6.49}.
\end{proof}

\begin{theorem}[An integer for four-dimensional trace differences]\label{B6.50}
Assume \(\mathscr H\). There are \(c,A\in\N\) as follows. Let \(k\ge2\), \(M\ge1\), and \(p\ge5\) be prime with \(p\nmid M\). Let \(z_i^h=1\) for \(z_1,z_2,z_3\in\C\) and \(h\ge1\). There is \(D\in\Z\setminus\{0\}\) with
\[
 \log|D|\le c\varphi(h)[M(k+1)]^A\log p .
 \tag{6.50.1}
\]
For each prime \(\ell\ne2,p\) dividing \(M\), fix \(\iota\) and a Frobenius \(\sigma_p\). Let \(X\) be a finite pairwise nonisomorphic set of irreducible \(G_\Q\)-representations satisfying the quadratic automorphy condition at \(\iota\), of level \((M;(k+1,k,1,0))\) with unlinked flags at \(\ell\). Require for each \(r\in X\) that \(r(g)^tPr(g)=\varepsilon_\ell(g)^{k+1}\chi(g)P\) for some invertible \(P\) and finite cyclotomic \(\chi\). Then
\(\Delta_r=\operatorname{tr}r(\sigma_p)-\iota^{-1}(p^{k+1}z_1+p^kz_2+pz_3+1)\ne0\) and
\[
 \sum_{r\in X}v_\ell(\Delta_r)\le v_\ell(D).
 \tag{6.50.2}
\]
\end{theorem}
\begin{proof}[Outline]
Apply \ref{B6.41a} in rank \(4\) via Lemma~\ref{B6.41b}(b) for a conjugation-stable label set \(S\) with bounded integral traces \(t_s\). The Galois-invariant product \(D=\prod_{s,a}F_{s,a}\) is a nonzero integer whose size and distinct factors yield both bounds. The complete proof is in \S\ref{prf:B6.50}.
\end{proof}

\section{Three-dimensional endpoints over the Gaussian field}

The three-dimensional deformation group over a coefficient ring \(R\) is \(\GL_3(R)\rtimes\{1,-1\}\) for the polarising involution. On the sign kernel its matrix part is a representation. For a Gaussian host the sign is \(\eta_4\), the quadratic character of \(\Q(i)\). Let \(\vartheta_{\ell'}\) be the prime-to-\(\ell\) order part of \(\vartheta\), let \(c_2(\vartheta)\) be its conductor exponent at \(2\) before taking that part, and put
\[
 S_b=\{v\mid N:v\ne\ell,\ \operatorname{comp}_v(\kappa_\ell(\vartheta))\ne1\},
 \qquad T_b=\{\ell\}\cup S_b,
\]
where \(\kappa_\ell\) is as in \ref{B6.28}. A Gaussian host has \(\ell\ne2\), \(k\ge1\), \(2\mid N\), \(\vartheta(-1)=(-1)^k\), and an integer \(o_1\ge1\) with
\begin{equation}\tag{6.51.0}\label{eq:p6c5-gaussian-host}
 \operatorname{comp}_r(\vartheta_{\ell'})^{2o_1}=1\quad(r\text{ odd prime}),
 \qquad 8o_1v^2<2^{c_2(\vartheta)}\quad(v\in T_b),
 \qquad c_2(\vartheta)\ge4 .
\end{equation}
Write \(S\) for the places dividing \(\ell N\), \(\Gamma_{\Q,S}\) for the Galois quotient unramified outside \(S\), and \(\Psi=\varepsilon_\ell^k\vartheta\). The frame uses the ordinary flag at \(\ell\). At \(v\in S_b\) a basis congruent to the identity diagonalises \(D_v\cap\ker\eta_4\), with prescribed diagonal on inertia. At the remaining places the frame kills inertia on \(\ker\eta_4\cap\ker\bar\Psi\), where a bar denotes reduction in the coefficient residue field.

\begin{lemma}[Graded finiteness from DVR points]\label{B6.51fin}
Let \(O\) be a complete DVR with uniformiser \(\varpi\) and finite residue field. Let \(R\) have a presentation as a quotient of \(O[[X_1,\ldots,X_s]]\) whose kernel is stable under projection to each weight for an assignment of integer weights to the \(X_i\). Write \(z_i\) for their images in the maximal ideal of \(R\) and \(R_0\) for the image of the weight-zero series. Suppose that every local \(O\)-algebra map \(R\to A\) to a complete DVR with finite residue field and \(\varpi A=0\) kills the zero-weight coordinates and kills all coordinates of at least one strict weight sign. Then \(R_0\) is finite over \(O\).
\end{lemma}
\begin{proof}[Outline]
Use Dickson's lemma for finite monomial generation. Normalise a one-dimensional quotient and use the DVR sign condition to put each generator in the uniformiser radical, then project to weight zero and lift a residue basis to get \(O\)-finiteness. The complete proof is in \S\ref{prf:B6.51fin}.
\end{proof}

\begin{theorem}[The Gaussian three-dimensional lifting assertion]\label{B6.51}
Assume the marking, residue-characteristic polynomial identity, and characteristic-zero attachment of \ref{B6.52}(4). For a Gaussian host and its frame \(F\) on \(\Gamma_{\Q,S}\), and any \(\iota\), there is a marking \(g_1\) with
\begin{equation}\tag{6.51.1}\label{eq:p6c5-gaussian-marking}
 \eta_4(g_1)=1,\qquad \Psi(g_1)^2-1\in O^\times .
\end{equation}
Let \(u\) be a universal lift of the marked frame. Its weight-zero algebra \(R_0\) consists of the elements of its coefficient ring that take equal values at any two specialisations related by conjugation by \(\operatorname{diag}(t,1,t^{-1})\). Here the specialisations take values in the same complete Noetherian local \(O\)-algebra with residue field \(O/\mathfrak m_O\), and \(t\) is any unit of that algebra. Then \(R_0\) is finite over \(O\). At a \(C_\ell\)-point of the algebra generated by the weight parts, if the matrix part of the associated \(r\) is irreducible on \(G_K\), \(K=\Q(i)\), then \((r|_{G_K})_{\rm mat}\varepsilon_\ell^{-1}\) is conjugate self-dual automorphic with witness.
\end{theorem}
\begin{proof}[Outline]
Pull back the marking and polynomial from \ref{B6.52} and use residue triangularity and \ref{B6.51fin} for a characteristic-zero DVR model. A three-cycle trace Vandermonde test gives open-subgroup irreducibility, so apply \ref{B6.52} again. The complete proof is in \S\ref{prf:B6.51}.
\end{proof}

\begin{definition}[Residual linked-frame and seed conditions]\label{B6.51frame}
Let \(\Phi\) be a linked frame with group \(\Gamma\) over a finite field \(k\) of odd characteristic \(\ell\). Write \(\eta\) for the sign, \(\Gamma_0\) for its kernel, \(\bar\Psi\) for the residual character on \(\Gamma\), and \(\bar\varepsilon\) for the residual cyclotomic character. At an ordinary place \(v\) write \(I_v\) for inertia, \(N_v\) for its specified subgroup in the cyclotomic kernel, and \(d_v\) for the local degree. The local maps and cyclotomic character are continuous and the finite characters have open kernels. The marking \(g_0\) has \(\bar\Psi(g_0)^2\ne1\). The involutions \(c,c_w\) have sign \(-1\), with \(\bar\Psi(c)=1\). The infinite-place indices \(w\) number \(d=\sum_vd_v\). Each open subgroup of \(\Gamma\) has a finite space of continuous \(\F_\ell\)-characters.

Require \(\bar\varepsilon,\bar\varepsilon\bar\Psi^{\pm1}\) nontrivial on \(\Gamma_0\), \(\bar\varepsilon\bar\Psi^{\pm2}\eta\) nontrivial on \(\Gamma\), and \(\bar\varepsilon|_{\Gamma_0}\) distinct from every integral power of \(\bar\Psi|_{\Gamma_0}\). Each local group contains an element with \(\bar\Psi^2\ne1\), and \(\bar\Psi|_{N_v}=1\). For each \(v\), require \(d_v\) elements of \(I_v\cap\ker\bar\Psi\), all but one in \(N_v\), whose values determine the restriction to \(I_v\) of every continuous character of the local group into a finite abelian \(\ell\)-group. Require one positive exponent killing all such characters of \(\Gamma_0\) that are trivial on inertia at both lifts of every framed place.

For a finite field \(k'\) of characteristic \(\ell\), let \(V\) be a finite-dimensional \(k'\)-representation of \(\Gamma\) with open kernel. Let \(b_v\subseteq V\) be stable under the local group at \(v\). Write \(H_b^i\) for the cohomology of the cone of the restriction of continuous cochains to the local cochains with coefficients \(V/b_v\), shifted by \(-1\). If \(\operatorname{Hom}_{\Gamma}(V,k'(\bar\varepsilon))=0\), require \(H_b^2\) finite-dimensional and
\begin{equation}\tag{6.51.2}\label{eq:p6c5-frame-euler}
 \dim H_b^2+\dim H_b^0+\sum_v d_v\dim b_v
 \le \dim H_b^1+\sum_w\dim V^{c_w=1}.
\end{equation}
Take \(b_v=0\) for the strict condition. For a character \(\chi\), a strict class is a pair consisting of a locally constant cocycle \(y\) and scalars \(b'_v\) with \(y(g)=(\chi(g)-1)b'_v\) on the local groups, modulo the pairs \(((\chi-1)b,(b)_v)\). Require the strict groups for \(\bar\Psi^{-1}\eta\) and \(\bar\Psi^{-2}\eta\) to be nonzero. The local square test makes each \(b'_v\) unique for these characters. The bounds with constant \(C_0\) are
\begin{equation}\tag{6.51.3}\label{eq:p6c5-frame-cohomology-bounds}
 \dim H^1_{\rm str}(\Gamma,k(\bar\Psi^{-2}\eta))\le C_0,\quad
 \dim H^1_{\rm str}(\Gamma_0,k(\bar\Psi^{-1}))\le C_0,\quad
 \dim H^1(\Gamma,k)\le1+C_0 .
\end{equation}
The strict condition on \(\Gamma_0\) uses both local lifts.

The seed data are an imaginary quadratic \(K_0\subset L\) for the abelian CM field \(L/\Q\), an integer \(D'>0\) prime to \(\ell\), a prime \(w\nmid D'\), and an integer \(j\ge0\). The \(C_\ell^\times\)-valued characters \(\chi',\eta'\) of \(G_\Q\) factor through \(\Q(\mu_{D'})\) and take value one on the residue class of \(w\). Here \(\eta'\) is the lift of the sign with kernel \(G_{K_0}\). The character \(\theta\) factors through \(\Q(\mu_w)\), has order \(\ell\), and is trivial on \(G_L\). Require
\[
 \chi'(c)=(-1)^j,\qquad \chi'\ne\eta',\qquad
 \ell\nmid |D_{K_0}|\operatorname{ord}(\chi'),\qquad
 \overline{\varepsilon_\ell^j\chi'}=\bar\Psi
\]
on the absolute group for the last equality. Require \(s\in G_L\) with \(\overline{(\varepsilon_\ell^j\chi')(s)}^{\,2}\ne1\).
\end{definition}


\begin{lemma}[The obstruction bound after imposing the inertia relations]\label{B6.51framea}
Let \(F\) be totally real, let \(\ell\) be odd, and let \(S\) consist of the places of \(F\) above \(\ell\). Let \(K\) be an open normal subgroup of \(G_F\) of index prime to \(\ell\). Form \(\Gamma\) by dividing \(G_F\) by the closed normal subgroup generated by \(I_q\cap K\) at all \(q\notin S\), using all prolongations. For \(v\in S\) use the map \(G_{F_v}\to G_F\to\Gamma\) attached to a prolongation. Let \(V\) be a representation of \(\Gamma\) with open kernel over a finite field \(k'\) of characteristic \(\ell\), of finite dimension over \(k'\). For any locally stable \(b_v\subseteq V\), form \(H_b^i\) as in Definition~\ref{B6.51frame}. If \(\operatorname{Hom}_{\Gamma}(V,k'(\bar\varepsilon))=0\), these groups are finite-dimensional for \(0\le i\le2\) and satisfy \eqref{eq:p6c5-frame-euler}, with \(d_v=[F_v:\Q_\ell]\).
\end{lemma}
\begin{proof}[Outline]
Use inflation to bound the degree-two kernel by the added local terms. Combine \ref{B6.09}, global and local Euler characteristics and the exact sequence for the local quotients to bound the framed obstruction group. The complete proof is in \S\ref{prf:B6.51framea}.
\end{proof}

\begin{lemma}[A determinantal section of a power series ring]\label{B6.51aa}
Let \(O\) be a complete discrete valuation ring. Let \(a,b,D\) be nonnegative integers with \(a=b+D\), let \(g\ge1\), and put
\[
 S=O[[x_1,\ldots,x_a,u_1,\ldots,u_g]],\qquad
 u=(u_1,\ldots,u_g)^t .
\]
Let \(f_1,\ldots,f_b\in\mathfrak m_S\) and \(Z\in M_g(\mathfrak m_S)\). Put
\[
 A=S/(f_1,\ldots,f_b,(Zu)_1,\ldots,(Zu)_g,\det Z).
\]
Every irreducible component of \(\operatorname{Spec}A\) has dimension at least \(D+1\). For any two minimal primes of \(A\) there is a chain of minimal primes from one to the other such that
\(\dim A/(P+Q)\ge D\) for consecutive members \(P,Q\).
\end{lemma}
\begin{proof}[Outline]
Parametrize the universal singular matrix with a kernel vector to obtain a prime radical, then bound its height using the adjugate identity. Specialize to \(Y=Z\) and \(v=u\) and apply the height and connectedness theorems to the section. The complete proof is in \S\ref{prf:B6.51aa}.
\end{proof}

\begin{lemma}[Numbers of even and odd relations at a block point]\label{B6.51ab}
Consider the completed ordinary matrix deformation ring \(R\) at a residual block point of the rank-three frame of Definition~\ref{B6.51frame}. Use the coefficient DVR \(O'\) with the residue field of the point. Conjugation by the involution which acts as \(1\) on the outer two coordinates and as \(-1\) on the middle coordinate gives the even and odd parts. There is a presentation over \(O'\) with \(a\) even and \(g\) odd variables, \(a-(3d+1)\) even relations, and \(g\) odd relations. The odd variables generate the ideal defining the block locus. If the point is in the closure of the complement of that locus, then \(g>0\).
\end{lemma}
\begin{proof}[Outline]
Split the minimal presentation by parity and use the square-zero obstruction and marked tangent count to bound the relations via the Euler characteristic for the frame of \ref{B6.51frame}. The odd variables generate the block ideal. The complete proof is in \S\ref{prf:B6.51ab}.
\end{proof}

\begin{lemma}[Points on the linked CM frame]\label{B6.51a}
Assume \(\mathscr H\). Put \(\Psi=\varepsilon_\ell^e\vartheta\), \(e\ge1\), with \(N\) even, \(\ell N\mid M\), and every prime ramified in \(\eta\) dividing \(M\). Fix \(B\mid\ell,\iota\) and a Galois link for \((\eta,\Psi,c,g_1,B,M)\) to a frame over an abelian CM field \(L\) satisfying \ref{B6.51frame}. Assume \(d>3C_0+7\) and \(d_v>5\). For indices \(0\le i<n\) put \(w_{3,i}=1-i\) and \(w_{2,i}=1-2i\).

For \(n=2,3\), take an adically continuous framed lift over a coefficient DVR finite over \(O\), with the prescribed diagonal reduction and inertia characters pulled back from \(\Psi^{w_{n,i}}\). If its representation \(r\) on \(G_L\) is irreducible, \(r\kappa_n\) is conjugate self-dual automorphic with witness. Here \(\kappa_3=\varepsilon_\ell^{-1}\) and one algebraic \(\kappa_2\) works for all rank-two points. In rank three over a complete DVR \(A\) with finite residue field and \(\mathfrak m_OA=0\), impose the same residual condition and require the flag characters to kill the kernel of \(\bar\Psi\) in inertia. The characteristic polynomial at \(g\) on the sign kernel is then \(\prod_{i=0}^2(X-\bar\Psi(g)^{1-i})\).
\end{lemma}
\begin{proof}[Outline]
Builds rank-two and adjoint rank-three seeds via \ref{B6.49h} and \ref{B6.48c}, extracting separating curves with \ref{B6.48bb}, \ref{B6.49g}, and \ref{B6.49i}. Propagates Hecke containment via \ref{B6.51aa}, \ref{B6.51ab}, and \ref{B6.48aa}, then applies \ref{B6.48a}. The complete proof is in \S\ref{prf:B6.51a}.
\end{proof}

\begin{lemma}[The newform and local witness after descent]\label{B6.52a}
Fix \(\iota:\overline{\Q}_\ell\simeq\C\). Suppose \(k\ge1\), \(L/\Q\) and \(W\) satisfy Lemma~\ref{B6.06f}, and \(W|_{G_L}\otimes\kappa\) has an attachment with witness in the sense of Definition~\ref{def:p6c5-witness} for an algebraic Hecke character with realization \(\kappa\). Make the following local assumptions.
\begin{enumerate}
\item At one prolongation above \(\ell\), on the intersection of inertia with an open normal subgroup of \(G_\Q\), \(W\) has an ordered flag with characters \(\varepsilon_\ell^k,\varepsilon_\ell^{-k}\).
\item At \(p\ne\ell\), the action of inertia factors through a finite abelian quotient on which the decomposition group acts trivially by conjugation.
\end{enumerate}
Then \(W\varepsilon_\ell^{-k}\) is attached to a normalized weight-\((2k+1)\) newform, and \(W\) satisfies the quadratic witness condition.
\end{lemma}
\begin{proof}[Outline]
Descend by \ref{B6.06f}, use \ref{B9.47}, \ref{B9.06} and \ref{B9.04} to identify the weight, and obtain the newform via \ref{B6.06d}(3). Transfer the support bounds by local base change and use the inertia flag ratios for potential unlinking. The complete proof is in \S\ref{prf:B6.52a}.
\end{proof}

\begin{theorem}[The endpoints used in cases I and V]\label{B6.52}
Assume \(\mathscr H\). For each good imaginary quadratic datum of the first count, each \(\iota\) and each \(B\mid\ell\), there is one marking \(g_1\) for assertions 1--3, with \(\eta(g_1)=1\) and \(\bar\Psi(g_1)^2\ne1\). An ordinary point has sign \(\eta\), is adically continuous, and reduces to the diagonal base point \(r_0\), whose matrix part on \(\ker\eta\) is \(\operatorname{diag}(\Psi,1,\Psi^{-1})\). Its matrix at complex conjugation \(c\) is \(1\). The \(01,10\) entries of the matrix difference \(r(g_1)-r(cg_1^{-1}c)\) vanish. Away from \(\ell\) it kills inertia with \(\eta=\Psi=1\). At \(B\) a basis congruent to the identity makes it upper triangular on \(D_B\cap\ker\eta\) with diagonal \((\Psi,1,\Psi^{-1})\) on inertia.
\begin{enumerate}
\item On one finite-index \(U\le G_\Q\), every such point over a complete DVR \(A\) with finite residue field and \(\mathfrak m_OA=0\) has characteristic polynomial \(\prod_{i=0}^2(X-\bar\Psi(g)^{1-i})\) for \(g\in U\cap\ker\eta\).
\item For a point over \(O'\subset C_\ell\) finite over \(O\), irreducibility on every open subgroup of \(\ker\eta\) gives a conjugate self-dual cuspidal attachment with witness for its restriction to \(G_K\) twisted by \(\varepsilon_\ell^{-1}\), where \(K\) is the sign field.
\item For a Type II point over \(O'\) with trivial middle character on \(\ker\eta\), every representation \(W:G_\Q\to\GL_2(C_\ell)\) agreeing there with its irreducible corner and satisfying \(\operatorname{tr}W(c)=0\) has \(W\varepsilon_\ell^{-k}\) attached to a weight-\((2k+1)\) newform, and satisfies the quadratic witness condition.
\item For each Gaussian host satisfying \eqref{eq:p6c5-gaussian-host}, and each \(\iota,c,B\), one marking satisfying \eqref{eq:p6c5-gaussian-marking} makes assertions 1 and 2 hold for its frame. The lifting conclusion is \ref{B6.51}.
\end{enumerate}
\end{theorem}
\begin{proof}[Outline]
Build a marked CM tower using Selmer classes, \ref{B6.17} and the obstruction bound of \ref{B6.51framea}. Apply \ref{B6.51a} through the link for the residual polynomials and the two automorphic ranks, then use \ref{B6.52a} for the corner. The complete proof is in \S\ref{prf:B6.52}.
\end{proof}


\section{Adjoint vanishing and ordinary derivations}

A polarised ordinary frame for a CM field \(L\) gives a quotient \(\Gamma\) of \(G_{L^+}\), its CM subgroup \(\Gamma_0\), and a complex conjugation \(c\). Write \(D_v\) and \(I_v\) for the local decomposition and inertia groups, and pull back through the frame maps when restricting to them. At the distinguished place the residual diagonal characters coincide on \(D_0\), and \(\bar\varepsilon_\ell|_{D_0}\ne1\). Write \(I_0\) for its inertia. For a representation \(\rho\) with polarisation \(P\), let \(R(g)X=\rho(g)X\rho(g)^{-1}\) on \(\Gamma_0\), and set
\begin{equation}\tag{6.53.0}\label{eq:p6c5-adjoint}
 R(c)X=-P^{-1}X^tP .
\end{equation}
For any action \(R\), write \(d_RX(g)=R(g)X-X\).

An ordinary fixed-weight cocycle on \(\Gamma_0\) is a continuous matrix-valued function \(z\) satisfying
\begin{equation}\tag{6.54.0}\label{eq:p6c5-ordinary-cocycle}
 z(gh)=z(g)+R(g)z(h),\qquad
 Pz(cgc^{-1})=-z(g)^tP.
\end{equation}
It vanishes on the prescribed tame inertia kernels. At each \(v\mid\ell\) it is upper triangular in the ordered flag after subtraction of a local coboundary, with zero diagonal on \(I_v\). Put \(\Gamma'=\Gamma_0\cap\ker\varepsilon_\ell\). Determinant regularity means that there are finitely many \(s_j\in\Gamma'\) with separable characteristic polynomial such that a continuous cocycle \(u\) for \(\varepsilon_\ell R\) with \(Pu(cgc^{-1})=u(g)^tP\) is a twisted coboundary whenever all the tests \(\operatorname{tr}(u(s_j)\rho(s_j)^a)\), \(a\ge0\), vanish.

\begin{theorem}[Bloch--Kato vanishing for the extended adjoint]\label{B6.53}
Assume \(\mathscr H\). Take a CM Hida setting at the base level \(S\) of \ref{B6.48c}, the union of the places of \(L^+\) above \(\ell\), the tame type places and the distinguished place. Use the level frame for this set. Suppose \(L\) contains no primitive \(\ell\)-th root of unity and \(O\) contains the images of \(\mathcal O_L\) under all \(\ell\)-adic embeddings. Let \(2\le n\le4\), fix an admissible residual ordinary datum with determinant ring \(B\), and take \(\iota:C_\ell\simeq\C\), an arithmetic point \(x\) of \(B\), and \(\rho:\Gamma_0\to\GL_n(C_\ell)\) whose determinant law is the specialisation \(D_x\) of the universal determinant at \(x\). Require the following.
\begin{enumerate}
\item The representation \(\rho\) is continuous, is trivial on \(I_0\), and its pullback to \(G_L\) is conjugate self-dual cuspidal automorphic with witness.
\item At every \(v\mid\ell\) it has a flag whose successive characters \(\kappa_{v,i}\) are the specialisations at \(x\) of the universal ordinary characters, in the prescribed order.
\item The restriction \(\rho|_{\Gamma'}\) is irreducible.
\item The elements of \(\Gamma'\) at which \(\rho\) has separable characteristic polynomial are dense in \(\Gamma'\).
\item At each \(v\mid\ell\) there is an element of \(D_v\) mapping to \(\Gamma'\) on which every ratio of distinctly indexed flag characters has infinite order.
\item At each \(v\mid\ell\) there are strictly decreasing integers \(m_i\) and an integer \(h\ge1\) with
\begin{equation}\tag{6.53.1}\label{eq:p6c5-inertia-powers}
 (\kappa_{v,i}(s)\kappa_{v,j}(s)^{-1}\varepsilon_\ell(s)^{m_j-m_i})^h=1
 \quad(s\in I_v,\ i,j).
\end{equation}
\end{enumerate}
Let \(P\in\GL_n(C_\ell)\) satisfy \(\rho(g)^tP\rho(cgc^{-1})=\varepsilon_\ell(g)^{1-n}P\) for \(g\in\Gamma_0\), and let \(R\) be the extended adjoint action \eqref{eq:p6c5-adjoint}. If a continuous \(R\)-cocycle \(z\) on \(\Gamma\) has Bloch--Kato finite pullback to \(G_{L^+}\), then for some \(X\in M_n(C_\ell)\),
\begin{equation}\tag{6.53.2}\label{eq:p6c5-adjoint-coboundary}
 z(g)=d_RX(g)\qquad(g\in\Gamma).
\end{equation}
\end{theorem}
\begin{proof}[Outline]
Identify the pullback with its cuspidal attachment by Chebotarev and Brauer–Nesbitt and match the extended adjoints via the polarisation. Apply Thorne's vanishing using cyclotomic irreducibility, then descend the coboundary by surjectivity. The complete proof is in \S\ref{prf:B6.53}.
\end{proof}

\begin{lemma}[Cocycles with zero trace derivative]\label{B6.53a}
Let \(\rho:G\to\GL_n(K)\) over a field \(K\) have full matrix span. An adjoint cocycle \(z\) with \(\operatorname{tr}(z(g)\rho(g))=0\) for every \(g\) is a coboundary.
\end{lemma}
\begin{proof}[Outline]
Using the zero-trace condition at \(gh\), the cocycle law, full span, and the trace pairing shows \(\partial(\rho(g))=z(g)\rho(g)\) defines a derivation on matrices. Putting \(A=\sum_i\partial(E_{i1})E_{1i}\) gives \(\partial(E_{jk})=AE_{jk}-E_{jk}A\), hence \(z=d_R(-A)\). The complete proof is in \S\ref{prf:B6.53a}.
\end{proof}

\begin{proposition}[An ordinary point on the prescribed unitary Hecke module]\label{B6.54a}
In the Hida setting of \ref{B6.54}, let the level module and its annihilator \(J\) be as in \ref{B6.48c}. An arithmetic point \(x\) with the stated cuspidal attachment, distinguished inertia kernel and ordered unlinked ordinary characters occurs on that module at its prescribed weight and level. Its Hecke character has the determinant coefficients and ordered diamond characters of \(x\). In particular \(J\subseteq\ker x\).
\end{proposition}
\begin{proof}[Outline]
Use full span and traces to obtain unramifiedness outside the level, and local compatibility for sphericity. Descent \ref{B9.36a} and split comparison \ref{B6.04a} give the required eigenspace, whose determinant agrees on good Frobenii and hence everywhere. The complete proof is in \S\ref{prf:B6.54a}.
\end{proof}

\begin{theorem}[Ordinary cocycles are coboundaries]\label{B6.54}
Assume \(\mathscr H\), \(2\le n\le4\), an admissible datum in the Hida frame at the base level \(S\) of \ref{B6.48c}, an arithmetic point \(x\) of weight \(h\) and an \(\iota:C_\ell\simeq\C\). Let \(\rho:\Gamma_0\to\GL_n(C_\ell)\) be continuous with determinant law \(D_x\) and polarisation \(P\) as in \ref{B6.53}. Write \(\kappa_{v,i}\) for the ordered ordinary characters specialised at \(x\). Assume \(\rho(I_0)=1\), conjugate self-dual cuspidal automorphy with witness on pullback to \(G_L\), determinant regularity, and
\begin{equation}\tag{6.54.1}\label{eq:p6c5-unlinked}
 \operatorname{span}_{C_\ell}\rho(\Gamma_0)=M_n(C_\ell),
 \qquad \kappa_{v,i}\ne\kappa_{v,j}\varepsilon_\ell
 \quad(v\mid\ell,\ i,j).
\end{equation}
Every ordinary fixed-weight cocycle \(z\) has \(z(g)=d_RY(g)\) on \(\Gamma_0\) for some \(Y\in M_n(C_\ell)\).
\end{theorem}
\begin{proof}[Outline]
Realize the cocycle as an integral dual-number lift with zero inertia derivatives. Use \ref{B6.48c} and \ref{B6.54a} for Hecke control, \ref{B6.55} to kill trace derivatives, and \ref{B6.53a} to recover a coboundary. The complete proof is in \S\ref{prf:B6.54}.
\end{proof}

\begin{theorem}[Inertia derivatives determine trace derivatives]\label{B6.55}
Assume the rank-\(n\) inputs of \(\mathscr H\), \(n\ge1\), in the Hida frame at the base level \(S\) of \ref{B6.48c}. Let \(B\) be the universal ordinary determinant ring with pair \((D^{\rm univ},\psi^{\rm univ})\), let \(x:B\to O\) be an integral arithmetic point, and let \(J\subseteq\ker x\). Require the level-ideal condition of \ref{B6.57}\textup{(1)} for the base module \(M\), with \(D=n[L^+:\Q]\) weight parameters in \(\Lambda\) given by inertia elements. Require compatible level modules at every split auxiliary level with distinct new places outside \(S\) and the compatible map to the base. For every integral arithmetic point \(x':B\to O\), put \(I_{x'}=\ker(x'|_\Lambda)\). For every \(t\) in the \(O\)-subalgebra of \(B\) generated by the coefficients of the characteristic polynomials of \(D^{\rm univ}\) at \(g\in\Gamma_0\), and every integer \(k\ge1\), require
\[
 t^kM\subseteq I_{x'}M\quad\Longrightarrow\quad tM\subseteq I_{x'}M.
\]

Suppose the determinant law specialised at \(x\) is that of an adically continuous representation \(\rho_x:\Gamma_0\to\GL_n(O)\) with an integral invertible symmetric or alternating polarisation as in \ref{B6.53}. Its ordered flag characters are the specialisations \(\kappa_{v,i}=x\circ\psi^{\rm univ}_{v,i}\). The representation kills the designated inertia kernels, and its matrix span over \(O\) contains \(aM_n(O)\) for some \(a\ne0\). At \(v\mid\ell\) assume an integral flag, a local element with distinct ordered values, and \(\kappa_{v,i}\ne\kappa_{v,j}\varepsilon_\ell\) for \(i<j\). At the distinguished place require
\begin{equation}\tag{6.55.1}\label{eq:p6c5-distinguished}
 \bigl(\varepsilon_\ell(g)\rho_x(g)Z=Z\rho_x(g)\ \text{for all }g\in D_0\bigr)
 \ \Longrightarrow\ Z=0\qquad (Z\in M_n(O)).
\end{equation}
Assume determinant regularity at elements \(s_1,\ldots,s_{q_0}\in\Gamma'\) with distinct integral roots \(\alpha_{ji}\) at each \(s_j\). Let \(\tau^{\rm univ}(g)\) be the negative of the coefficient of \(X^{n-1}\) in the universal characteristic polynomial of \(g\). Every \(O\)-linear \(x\)-derivation \(\delta:B\to O\) killing the universal ordinary characters on inertia satisfies
\begin{equation}\tag{6.55.2}\label{eq:p6c5-inertia-to-trace}
 \delta(\tau^{\rm univ}(g))=0\qquad(g\in\Gamma_0).
\end{equation}
\end{theorem}
\begin{proof}[Outline]
Use the places of \ref{B6.56} and the bound of \ref{B6.55a}. Arithmetic control and \ref{B6.57a} make the local weight-fibre factor a field, giving an idempotent, and the vanishing weight derivatives let \ref{B6.57} kill the trace derivatives. The complete proof is in \S\ref{prf:B6.55}.
\end{proof}

\begin{theorem}[Taylor--Wiles auxiliary places]\label{B6.56}
Assume the local and global cohomological inputs. Take the CM Hida frame at the base level \(S\) of \ref{B6.48c} with residual characters \(\bar\chi_i\), and suppose each open subgroup of \(\Gamma\) has finitely many continuous \(\F_\ell\)-characters. For an open normal subgroup \(H\) of \(\Gamma\), write \(N_\ell(H)\) for the kernel of its maximal pro-\(\ell\) quotient. Put
\[
 H=\Gamma_0\cap\ker\bar\varepsilon_\ell\cap\bigcap_i\ker\bar\chi_i .
\]
For any \(s_1,\ldots,s_{q_0}\in\Gamma'\) there are auxiliary systems of distinct \(v_{r,j}\notin S\), split in \(L\), and places \(w_{r,j}\mid v_{r,j}\) with
\begin{equation}\tag{6.56.1}\label{eq:p6c5-aux-norm}
 \mathrm Nw_{r,j}\equiv1\bmod\ell^r .
\end{equation}
The enlarged level has a quotient map \(h_r:\Gamma_r\twoheadrightarrow\Gamma\), inverse arithmetic Frobenii \(\phi_{r,j}\), and tame elements \(\tau_{r,j}\). The differences \(d_{r,j}=h_r(\phi_{r,j})s_j^{-1}\) belong to \(H\), and for every open normal \(U\supseteq N_\ell(H)\) they belong to \(U\) for all sufficiently large \(r\). Moreover \(h_r(\tau_{r,j})=1\) and
\begin{equation}\tag{6.56.2}\label{eq:p6c5-tame-relation}
 \phi_{r,j}^{-1}\tau_{r,j}\phi_{r,j}
 \equiv\tau_{r,j}^{\mathrm Nw_{r,j}}\pmod{N_\ell(h_r^{-1}H)} .
\end{equation}
The kernel of \(h_r\) lies in the closed normal subgroup generated by the \(\tau_{r,j}\) and \(N_\ell(h_r^{-1}H)\). The groups \(h_r^{-1}H\) have uniformly bounded numbers of continuous \(\F_\ell\)-characters.
\end{theorem}
\begin{proof}[Outline]
Apply Chebotarev in finite level quotients using cofinal pro-\(\ell\) approximations and cyclotomic congruences. The new tame generators give the conjugation and normal generation relations, and their finitely many coset conjugates bound the character count. The complete proof is in \S\ref{prf:B6.56}.
\end{proof}

\begin{lemma}[The ordinary Selmer trace bound]\label{B6.55a}
Under \ref{B6.55} use the system of \ref{B6.56}. Put \(A_m=O/\mathfrak m_O^m\) and let \(\bar\rho_{r,m}\) be the pullback of \(\rho_x\) to the CM subgroup of \(\Gamma_r\), reduced to \(A_m\). Define \(\mathcal S_{r,m}^{\rm tr}\) to be the set of functions \(g\mapsto\operatorname{tr}(z(g)\bar\rho_{r,m}(g))\) for locally constant matrix cocycles over \(A_m\) satisfying \eqref{eq:p6c5-ordinary-cocycle} and vanishing on the prescribed tame inertia kernels. At \(v\mid\ell\) require that a local coboundary makes the cocycle upper triangular in the prescribed flag, with zero diagonal on
\[
 \mu_v=\{g\in D_v:\text{the image of }g\text{ in the topological
 abelianisation has finite }\ell\text{-power order}\}.
\]
At the auxiliary places use the full local condition. Put \(Q=\#(O/\mathfrak m_O)\), \(b=\binom n2\), \(d=[L^+:\Q]\) and \(s_\ell=\#\{v\mid\ell\text{ in }L^+\}\). There is a constant \(C\) such that for every \(m\ge1\) and all sufficiently large \(r\),
\begin{equation}\tag{6.55.3}\label{eq:p6c5-selmer-count}
 \#\mathcal S_{r,m}^{\rm tr}\,Q^{m((n+b)s_\ell+bd)}
 \le Q^{C+m((n+b)(s_\ell+d)+nq_0)}.
\end{equation}
In particular \(\#\mathcal S_{r,m}^{\rm tr}\le Q^{m(D+nq_0)+C}\).
\end{lemma}
\begin{proof}[Outline]
Use \ref{B6.09} and Euler characteristics to bound the local factors, with \ref{B6.56} identifying the auxiliary centralizers. Descend dual cocycles from auxiliary level and use determinant regularity and Smith reduction to bound their classes uniformly. The complete proof is in \S\ref{prf:B6.55a}.
\end{proof}

\begin{lemma}[Uniform trace and character comparisons]\label{B6.57a}
Use the integral representation, polarisation and separated flags of \ref{B6.55}, with \(\ell\ne2\). On an adapted level take a compatible ordinary determinant pair over a complete local \(O\)-algebra \(R\) whose specialisation at \(x_R:R\to O\) is the pullback of the determinant law of \(\rho_x\) together with its ordered characters \((\kappa_{v,i})\). At each designated tame place \(w\), assume that the pair's inertia kernel \(I_p(w)\) equals the frame's prescribed kernel \(I_u(w)\). An \(x_R\)-derivation into an \(O\)-module is \(O\)-linear and satisfies \(\delta(ab)=x_R(a)\delta(b)+x_R(b)\delta(a)\). There is a fixed \(c_1\in O\setminus\{0\}\), independent of level and \(m\), such that the trace derivatives of any \(A_m\)-valued \(x_R\)-derivation multiplied by \(c_1\) belong to \(\mathcal S_{r,m}^{\rm tr}\) of \ref{B6.55a}. There is a fixed \(c_2\in O\setminus\{0\}\), independent of level and target, that kills the character derivatives of every \(x_R\)-derivation that kills traces.
\end{lemma}
\begin{proof}[Outline]
Use rank-one Cayley–Hamilton corners to recover a scaled matrix derivative and average for polarisation. Differentiate tame and flag relations at a separating element and use eigenvalue interpolation to annihilate character derivatives when trace derivatives vanish. The complete proof is in \S\ref{prf:B6.57a}.
\end{proof}

\begin{lemma}[Counting derivations in finite quotients]\label{B6.57}
Let \(\ell\ne2\), \(n\ge1\), and let \(O\) be a complete coefficient DVR with finite residue field \(\mathbf k\). Fix an adapted system as in \ref{B6.56}, an ordinary determinant pair \(u\) over a complete Noetherian local \(O\)-algebra \(B\) with residue field \(\mathbf k\), an ideal \(J\), and an integral point \(x\) as in \ref{B6.55}. Suppose the determinant coefficients and characters generate \(B\) adically, \(J\subseteq\ker x\), and the following conditions hold.
\begin{enumerate}
\item The base module \(M\) is finite free over \(\Lambda=O[[X_1,\ldots,X_D]]\), has support \(V(J)\), and each \(X_k\) acts by its weight character times a fixed unit, minus one. At auxiliary level \(r\) the ordinary pair \(u_r\) over \(R\) maps to the pullback of \(u\) via \(\beta:R\to B\), with the same weight identities and the prescribed enlarged inertia kernels. The module \(M_1\) is finite over \(\Lambda\) and projective over \(\Lambda[\Delta]\), where \(\Delta=\prod_{j,i}\Z/\ell^{a_j}\), \(a_j\ge r\). Its standard generators act by \(1+y_{ji}\). The surjection \(M_1\to M_w\) intertwines the root operators and has kernel \(\sum_{j,i}y_{ji}M_1\).

Write \(Q_j=\mathrm Nw_{r,j}\), let \(u_{ji}\) be the root operators, and put
\[
 f=\prod_j\prod_{i\ne a}(u_{ji}-u_{ja})(u_{ji}-Q_ju_{ja}).
\]
The comparison maps between \(M_w\) and \(M\) tensored with the ordered splitting algebra of the base Frobenius polynomials intertwine the coefficient action via \(\beta\) and the root-polynomial actions. Both composites are multiplication by \(f\). For some fixed integer \(E\ge0\) and every \(j,i\), the operator
\begin{equation}\tag{6.57.0}\label{eq:p6c5-operator}
 f^E\operatorname{tr}_{u_r}\!\left[
 \left(\prod_{a\ne i}(\phi_{r,j}-u_{ja})\right)^{E+1}
 (\tau_{r,j}-1-y_{ji})\right]
\end{equation}
acts by zero on \(M_1\), with the trace extended linearly in the operator coefficients.
\item Some \(t_0\in B\) with \(x(t_0)\ne0\) satisfies \(t_0aM\subseteq\ker(x|_\Lambda)M\) for every \(a\in\ker x\).
\item The representation \(\rho_x\) kills \(N_\ell(H)\) and has distinct integral roots \(\alpha_{ji}\) at \(s_j\).
\item For some fixed nonnegative integer \(C\), at every \(m\ge1\) and all sufficiently large \(r\) the finite-layer Selmer trace set of \ref{B6.55a} satisfies
\begin{equation}\tag{6.57.1}\label{eq:p6c5-count-bound}
 \#\mathcal S_{r,m}^{\rm tr}\le |\mathbf k|^{m(D+nq_0)+C}.
\end{equation}
\end{enumerate}
Then every \(O\)-linear \(x\)-derivation \(\delta:B\to O\) killing the selected weight characters kills all universal traces.
\end{lemma}
\begin{proof}[Outline]
Use \ref{B6.56} and the comparison operators to lift weight and diamond derivations, then adjoin a putative nonzero trace derivative. By \ref{B6.57a} their combinations give Selmer traces with bounded loss of precision, exceeding the count bound. The complete proof is in \S\ref{prf:B6.57}.
\end{proof}
\section{The algebraic length comparison}

Let \(A=\bigoplus_j A_j\) be a graded \(O\)-algebra over a domain and \(x:A\to O\) an \(O\)-point. Put \(\mathfrak a=\ker x\) and \(S_x=\operatorname{sat}_O(\mathfrak a^2)\), the saturation under nonzero elements of \(O\). Let \(C_j\) be the image of \(A_j\cap\mathfrak a\) in \(A/S_x\). For a family \(F=(F_i)\) in \(\mathfrak a\), put
\[
 C_j(F)=C_j/(C_j\cap\langle[F_i]\rangle_O).
 \tag{6.58.0}
\]
A point is fixed if it kills every nonzero weight. It is two-sided if some element of positive weight and some element of negative weight have nonzero values. A point of one strict sign has a nonzero value of that sign and kills all weights of the opposite sign. The torus acts by multiplying a weight-\(j\) value by \(t^j\).

\begin{lemma}[Two-sided orbit slices, chart reduction and the Rees module]\label{B6.58c}
Let \(O\) be a complete characteristic-zero DVR with uniformiser \(\pi\). Let \(K\) be an algebraically closed extension of \(\operatorname{Frac}O\) with an extending valuation \(v\) normalised by \(v(\pi)=1\). Let \(A_0\) be complete Noetherian local with residue field that of \(O\). Assume \(A=\bigoplus_j A_j\) is generated over \(A_0\) by finitely many elements of weights \(-2,-1,1,2\), and \(x:A\to O\) is a fixed \(O\)-point. Let \(I\) be generated by \(F_1,\ldots,F_{r+1}\) of weight \(-1\), \(G_1,\ldots,G_s\) of weight \(-2\), \(E_1,\ldots,E_{r'}\) of weight \(1\), \(E'_1,\ldots,E'_{s'}\) of weight \(2\), and \(Z_1,\ldots,Z_m\in A_0\cap\ker x\). Suppose the following five orbit conditions hold.
\begin{enumerate}
\item \(r\ge1\). There is a graded map to a formal power series ring over \(\operatorname{Frac}O\) with constant term \(x\) and linear terms that identify the rational cotangent space with the span of the variables. In weights \(-2,-1,0,1,2\) the variable counts are \(s,r,m,r'+1,s'+1\), with no other weights. The \(F_i,G_i,Z_i\) span the rational cotangent spaces of weights \(-1,-2,0\), respectively.
\item The image of \(A_0\) in \(A/I\) is finite over \(O\).
\item The fibre over the \(A_0\)-restriction of a fixed \(K\)-point of \(A/I\) distinct from \(x_K\) contains points of at most one strict sign.
\item Every two-sided \(K\)-point of \(A/I\) has tangent space of dimension at most one.
\item An element \(a\in A_0\cap\ker x\) is nonzero at every two-sided \(K\)-point of \(A/I\).
\end{enumerate}
Then the following assertions hold.
\begin{itemize}
\item For every \(f\in A_{-2}\) and \(g\in A_2\), the slice \(f=1,\ g\ne0\) of \((A/I)\otimes_O K\) is finite and reduced over \(K\), and there is a finite set \(T\) of two-sided \(K\)-points of \(A/I\) meeting every two-sided orbit.
\item The quotient \(C_{-1}(F)\) has finite length. After an invertible \(O\)-linear change among \(F_1,\ldots,F_{r+1}\) preserving \(I\), there are \(y_1,\ldots,y_r\in A_{-1}\) and integers \(b_1,\ldots,b_r\ge0\) such that
\begin{equation}\tag{6.58.4}\label{eq:p6c5-cot-smith}
 F_i-\pi^{b_i}y_i\in S_x\quad(1\le i\le r),\qquad F_{r+1}\in S_x,
\end{equation}
and
\begin{equation}\tag{6.58.5}\label{eq:p6c5-cot-length}
 \operatorname{length}_O C_{-1}(F)=\sum_{i=1}^r b_i .
\end{equation}
\item Setting the weight-zero coordinates given by \(Z_1,\ldots,Z_m\) equal to zero by formal inversion defines a graded \(O\)-algebra map \(\phi\) from \(A\) to a formal power series domain over \(\operatorname{Frac}O\) in the remaining variables, and \(a\) may be replaced modulo \(I\) by an element of \(A_0\cap\ker x\) satisfying \(\phi(a)\ne0\).
\item Inside \(A[T_0,T_1,T_2]\), let \(R\) be the \(\N^3\)-graded \(A_0\)-subalgebra generated in the three unit multidegrees by \(A_{-2}T_0\), \(A_2T_1\), and \((A_0F_{r+1}^2+A_{-2}^2A_2)T_2\). Let \(J\subset R\) be the kernel of \(\phi\) on coefficients, put \(M=R/J\), and let \(\mathbf t\) be the homogeneous list consisting of \(F_i^2T_0\) (\(1\le i\le r\)) and \(G_iT_0\) in the first direction, \(E_i^2T_1\) and \(E'_iT_1\) in the second, and \(F_{r+1}^2T_2\) in the third. Let \(R_{++}\) be the product of the three ideals of \(R\) generated in unit multidegrees, call a homogeneous list admissible on a graded \(R\)-module \(N\) if a power of \(\mathfrak m_{A_0}R\,R_{++}\) lies in its ideal plus \(\operatorname{Ann}_R N\), and for an admissible list \(L\) write \(\chi_L(N)_d\) for the alternating sum of \(O\)-lengths of its Koszul homology in multidegree \(d\). For any homogeneous section \(s\) with \((s,L)\) admissible on \(N\) and coordinatewise sufficiently large \(d\),
\begin{equation}\tag{6.58.5a}\label{eq:p6c5-koszul-recursion}
 \chi_{(s,L)}(N)_d
 =\chi_L(N/sN)_d-\chi_L(0:_N s)_{d-\deg s}.
\end{equation}
\end{itemize}
\end{lemma}
\begin{proof}[Outline]
Finite generation of \(A_{-2},A_2\) and tangent dimension at most one make each slice \(f=1,g\ne0\) finite and reduced, giving a cover \(T\). Smith reduction computes \(C_{-1}(F)\), formal inversion builds \(\phi\), and Koszul exact sequences give the length recursion. The complete proof is in \S\ref{prf:B6.58c}.
\end{proof}

\begin{lemma}[From Selmer length to trace differences]\label{B6.58}
\emph{Framed presentation.} Let \(O\) be complete Noetherian local with finite residue field \(\mathbf k\) of odd characteristic \(\ell\). Let \(\Gamma\) be profinite, with finitely many continuous \(\F_\ell\)-characters on each open subgroup. Fix a block decomposition of a basis with weights \(w_i\) in a finite free lattice. An invertible symmetric or alternating form pairs the blocks by an involution and has constant total weight on its nonzero entries. Fix a residual representation \(\bar r\) and a lift \(r_0\). Both preserve the form up to similitude, permute the blocks and preserve the weight spaces. Allowed lifts differ from \(r_0\) by block-diagonal isometries congruent to the identity.

Fix finite families of continuous local maps \(\alpha_v:G_v\to\Gamma\) and \(\beta_q:G_q\to\Gamma\), with a distinguished frame \(v_0\) among the \(v\). Fix \(N_q\le G_q\) on which \(\bar r\beta_q\) is trivial. The local conditions at \(q\) hold at \(r_0\), are given by closed ideals, and are invariant under allowed conjugation and torus conjugation congruent to the identity. Local solutions lift across small extensions. If two allowed lifts have the same image across a small extension and the first is a solution, the second is a solution exactly when the two agree on \(N_q\).

Let \(V_j\) be the weight-\(j\) part of the block-diagonal Lie algebra of isometries over \(\mathbf k\), with action \(\operatorname{Ad}\bar r\). Write \(\delta b(g)=gb-b\) and \(du(g,h)=u(g)+gu(h)-u(gh)\). Let \(H^1_{\mathrm{fr},j}\) consist of pairs \((z,(b_v))\) with \(z\) a continuous cocycle in \(V_j\), \(z\alpha_v=\delta b_v\), \(b_{v_0}=0\), and \(z\beta_q|_{N_q}=0\). Let \(\operatorname{Ob}_j\) be the space of continuous triples \((c,(u_v),(u_q))\) with \(c\) a \(2\)-cocycle in \(V_j\) and \(du_v=\alpha_v^*c,\ du_q=\beta_q^*c\), modulo the triples
\[
 \bigl(du,(u\alpha_v-\delta b_v)_v,(u\beta_q-e_q)_q\bigr).
\]
Here \(u\) is a continuous global \(1\)-cochain, the \(b_v\) are local vectors, and the \(e_q\) are continuous local cocycles vanishing on \(N_q\). Assume each \(\operatorname{Ob}_j\) is finite-dimensional.

Let \(B\) be a complete Noetherian local \(O\)-algebra with residue field \(\mathbf k\) and a homogeneous power series presentation on finitely many variables. Give each \(G_v\) an allowed lift over \(B\) whose \(ab\)-entry has weight \(w_a-w_b\). The universal algebra of lifts framed against these local lifts, with the \(v_0\)-frame equal to the identity, has a homogeneous presentation relative to \(B\). The new variables are entries of \(r-r_0\) or of frame matrices minus the identity. The relations can be padded with zeros to give
\[
 \#\text{variables}_j=\dim_{\mathbf k}H^1_{\mathrm{fr},j},\qquad
 \#\text{relations}_j=\dim_{\mathbf k}\operatorname{Ob}_j .
 \tag{6.58.1}
\]

\emph{Cotangent comparison.} Let \(O\) be a DVR with finite residue field, \(x(A_j)=0\), and \(F_i\in A_j\). Suppose an \(O\)-module \(M\) injects additively into the \(O/(c)\)-valued \(x\)-derivations for \(0\ne c\in O\). The image derivations kill the \(F_i\) and lift to \(O\)-valued \(x\)-derivations supported in weight \(j\). With infinite lengths allowed,
\[
 \operatorname{length}_O M\le\operatorname{length}_O C_j(F).
 \tag{6.58.2}
\]

\emph{Orbit estimate.} Let \(O\) be a complete characteristic-zero DVR with uniformiser \(\pi\). Let \(K\) be an algebraically closed extension of \(\operatorname{Frac}O\) with an extending valuation \(v\) normalised by \(v(\pi)=1\). Let \(A_0\) be complete Noetherian local with residue field that of \(O\). Assume \(A\) is generated over \(A_0\) by finitely many elements of weights \(-2,-1,1,2\), and \(x:A\to O\) is fixed. Let \(I\) be generated by \(F_1,\ldots,F_{r+1}\) of weight \(-1\), \(G_1,\ldots,G_s\) of weight \(-2\), \(E_1,\ldots,E_{r'}\) of weight \(1\), \(E'_1,\ldots,E'_{s'}\) of weight \(2\), and \(Z_1,\ldots,Z_m\in A_0\cap\ker x\). Suppose the following conditions hold.
\begin{enumerate}
\item \(r\ge1\). There is a graded map to a formal power series ring over \(\operatorname{Frac}O\) with constant term \(x\) and linear terms that identify the rational cotangent space with the span of the variables. In weights \(-2,-1,0,1,2\) the variable counts are \(s,r,m,r'+1,s'+1\), with no other weights. The \(F_i,G_i,Z_i\) span the rational cotangent spaces of weights \(-1,-2,0\), respectively.
\item The image of \(A_0\) in \(A/I\) is finite over \(O\).
\item The fibre over the \(A_0\)-restriction of a fixed \(K\)-point of \(A/I\) distinct from \(x_K\) contains points of at most one strict sign.
\item Every two-sided \(K\)-point of \(A/I\) has tangent space of dimension at most one.
\item An element \(a\in A_0\cap\ker x\) is nonzero at every two-sided \(K\)-point of \(A/I\).
\end{enumerate}
Then \(C_{-1}(F)\) has finite length. There is a finite set \(T\) of two-sided \(K\)-points of \(A/I\) meeting every two-sided orbit. Every such \(T\) satisfies
\[
 \operatorname{length}_O C_{-1}(F)\le2\sum_{\xi\in T}v(\xi(a)).
 \tag{6.58.3}
\]
\end{lemma}
\begin{proof}[Outline]
Presents the framed universal ring on pro-\(\ell\) generators using \(H^1_{\mathrm{fr},j}\) and \(\operatorname{Ob}_j\) via small extensions. Derives the cotangent bound from \(x\)-derivations and the orbit bound from \ref{B6.58c} via Koszul Euler characteristics. The complete proof is in \S\ref{prf:B6.58}.
\end{proof}

\begin{lemma}[The lower intersection number]\label{B6.58a}
Use the chart, exponents, Rees algebra and admissibility notion of Lemma~\ref{B6.58c}. If \(\mathbf t\) is admissible on \(M/aM\), then in coordinatewise sufficiently large multidegree \(d\),
\begin{equation}\tag{6.58.6}\label{eq:p6c5-intersection-lower}
 2^{r+r'}\sum_i b_i\le\chi_{\mathbf t}(M/aM)_d .
\end{equation}
\end{lemma}
\begin{proof}[Outline]
Pass to leading terms in the negative and positive weights and compute their Hilbert series. Removing the Smith scalars gives the lower bound through the Koszul coefficients, and \ref{B6.58b} gives nonnegativity for the remaining terms. The complete proof is in \S\ref{prf:B6.58a}.
\end{proof}

\begin{lemma}[The upper intersection number]\label{B6.58b}
With the orbit hypotheses and construction of Lemma~\ref{B6.58c}, the list \((\mathbf t,a)\) is admissible on \(M\), and in sufficiently large multidegrees,
\begin{equation}\tag{6.58.7}\label{eq:p6c5-intersection-upper}
 \chi_{\mathbf t}(M/aM)_d\le\chi_{(a)}(M/\mathbf tM)_d
 \le2^{r+1+r'}\sum_{\xi\in T}v(\xi(a)).
\end{equation}
On any finite graded \(R\)-module an admissible list of nonnegative multidegrees has eventually nonnegative Koszul Euler characteristic.
\end{lemma}
\begin{proof}[Outline]
Classifies primes over \(J+(\mathbf t)\) avoiding \(R_{++}\) via \ref{B6.58c} to show admissibility of \((\mathbf t,a)\) and bound horizontal multiplicities by \(2^{r+1+r'}\). Prime filtrations bound \(\chi_{(a)}(M/\mathbf t M)_d\) and give Koszul nonnegativity. The complete proof is in \S\ref{prf:B6.58b}.
\end{proof}

\section{Residual vanishing from a Bernoulli unit}

For a character \(\gamma:G_\Q\to\mathbf k^\times\), let \(\mathbf k(\gamma)\) have action by \(\gamma\). Consider the following inertia condition on a cocycle \(z\) in this module. For every place \(A\), with \(J_A=I_A\cap\ker\gamma\), there is a subgroup \(H\le G_\Q\) such that
\begin{equation}\tag{6.59.0}\label{eq:p6c5-unrev}
 [J_A:J_A\cap H]\text{ is finite and prime to }\ell,\qquad z|_{J_A\cap H}=0 .
\end{equation}

\begin{lemma}[Residual Bernoulli-unit criterion]\label{B6.59}
Let \(\ell\) be odd, \(O\subset C_\ell\) a coefficient DVR and \(\operatorname{red}:O\to\mathbf k\) a ring map to a discrete field of characteristic \(\ell\). Let \(\vartheta:(\Z/N)^\times\to O^\times\), \(N\ge1\), be nontrivial of order prime to \(\ell\). For a nontrivial character \(\lambda\) at level \(L\), put \(B_{1,\lambda}=L^{-1}\sum_{t\in(\Z/L)^\times}[t]_L\lambda(t)\), where \([t]_L\) is the least nonnegative residue. Regard \(\bar\vartheta\) as the reduced Galois character. Assume
\[
 \bar\vartheta\ne\bar\varepsilon_\ell,\qquad |B_{1,\vartheta^{-1}}|_\ell=1.
 \tag{6.59.1}
\]
Every continuous cocycle \(z:G_\Q\to\mathbf k(\bar\vartheta)\) satisfying the inertia condition \eqref{eq:p6c5-unrev} has the form
\begin{equation}\tag{6.59.2}\label{eq:p6c5-bernoulli-coboundary}
 z(g)=(\bar\vartheta(g)-1)b\quad(g\in G_\Q),\qquad b\in\mathbf k.
\end{equation}

Let \(\phi\bmod f,\xi\bmod 2^n\) be primitive \(O\)-valued characters, with \(f\ge1\) odd, \(n\ge0\) and \(d\ge1\). Assume
\[
 4\mid\ell^d-1,\qquad \operatorname{ord}\phi\mid\ell^d-1,\qquad
 \phi(-1)\xi(-1)=-1,\qquad f4^{v_2(\ell^d-1)}\le2^n .
 \tag{6.59.3}
\]
Then \(|B_{1,\phi\xi}|_\ell=1\) at level \(f2^n\), and the cocycle conclusion holds for \(\gamma=\overline{\phi\xi}\) and its inverse.
\end{lemma}

\begin{proof}[Outline]
Use \ref{eq:p6c5-unrev} to make the cyclotomic restriction unramified and kill it by smoothing and Stickelberger \ref{B6.10a}. Commutativity gives the coboundary, and \ref{B6.59a} and last dyadic inertia give units and distinguish the residual characters. The complete proof is in \S\ref{prf:B6.59}.
\end{proof}

\begin{lemma}[An effective dyadic Bernoulli unit]\label{B6.59a}
Let \(\ell\) be prime, \(d\ge1\), \(f\ge1\) odd and \(n\ge0\). Let \(\phi,\xi\) be primitive characters over \(C_\ell\) of levels \(f,2^n\). Put \(Q=\ell^d\) and \(e=v_2(Q-1)\). If \(4\mid Q-1\), \(\operatorname{ord}\phi\mid Q-1\), \(\phi(-1)\xi(-1)=-1\), and \(f4^e\le2^n\), then \(|B_{1,\phi\xi}|_\ell=1\) at level \(f2^n\).
\end{lemma}
\begin{proof}[Outline]
Group the Bernoulli sum modulo \(2^n\) and use \ref{eq:p6c5-dyadic-shift} for integral Frobenius-fixed coefficients. If reduction vanishes, average to restrict the residues, subtract the \(1-s\) shift, and pair by the last dyadic involution to get \(-2\) by odd parity. The complete proof is in \S\ref{prf:B6.59a}.
\end{proof}

Count I uses the trace bound of \ref{B6.42} and the two vanishing statements of \ref{B6.43}. It also uses the imaginary quadratic endpoints of \ref{B6.52}.
Counts II and III use \ref{B6.49} and the Bernoulli-unit criterion of \ref{B6.59} for the six residual characters in \ref{B6.31}.
Count V uses \ref{B6.51} and applies \ref{B6.59} to the product character of \ref{B6.34}, whose reduction and inverse reduction are \(\eta_4\bar\Psi^2\) and \(\eta_4\bar\Psi^{-2}\).








% ---- p6c6 ----
\chapter{The exponent 401 and the downstream window}\label{ch:B6-6}

\begin{lemma}[The rank cutoff]\label{B6.60}\label{B6.61}
For a family satisfying the page normalization, $e_P=12$, so
$c_{TO}=28e_P+41=377$ and $C_t=c_{TO}+24=401$.

Let $0<\varepsilon<1/52$ and $x\ge2^{100}$. At scale $x$, put
$T_O=\lceil x^{377}\rceil$. Let $R$ be a rank datum with rank parameters
$m,k,t\in\N$. Write $n=|A|$, where $A$ is the finite abelian group of
its frame.
Put $f=2n$, $r=k+m-1$, and $L=x^\varepsilon$. Thus $m\ge2$ and
$m+3\le k$.
Define
\[
 n_W=8\left\lceil\frac{10r}{8}
 \left(f+\frac{f^2(k-1)^3L^2}{m}\right)\right\rceil .
\]
Assume
\[
\begin{aligned}
 x^{2+8\varepsilon}&\le n\le2x^{2+8\varepsilon},\\
 x^{1/2-\varepsilon/16}&\le m\le2x^{1/2-\varepsilon/16},\\
 x^{9/2+10\varepsilon}&\le k\le2x^{9/2+10\varepsilon},
 \qquad r+t=8fr+4+T_On_W.
\end{aligned}
\]
Then $t\le x^{400}$.
\end{lemma}
\begin{proof}[Outline]
Put \(u=x^{1/100}\ge 2\) and bound \(f,k,r,m,L^2\) by powers of \(u\) using \(\varepsilon<1/52\). Substituting these estimates into \(n_W\), \(T_O\), and the rank identity yields \(t\le 10708u^{39967}\le u^{40000}=x^{400}\) via \(10708\le 2^{33}\le u^{33}\). The complete proof is in \S\ref{prf:B6.61}.
\end{proof}

\begin{theorem}[Windows at arbitrarily large scales]\label{B6.62}
Fix a primitive character $\psi$, $0<\varepsilon<1/52$, and $c\ge1$.
Assume the window assertion for $\psi$ with exponent $401$ in Definition~\ref{B6.03} and the
prime-pair existence assertion of \ref{B2.43}.
Then there is $J\ge1$ such that, for every $h\in\R$, there is $Z\ge h$
with the following property. For every $x\in[Z,Z^{3/2}]$ there exist
$U\in\R$, primes $q,d$, and $a\in\N$ such that
\[
 x^{1/J}\le U\le x,\qquad
 (q,d)\text{ is a good pair for }(\operatorname{cond}\psi,\varepsilon)\text{ on }[U,x],
 \qquad
 x^{2+8\varepsilon}\le c2^a\le2x^{2+8\varepsilon}.
\]
Use the good-pair conditions of \ref{B2.42}. Every field associated with $(\psi,a,q,d)$ in Definition~\ref{B6.02} satisfies the total bound of Definition~\ref{B6.01} up to $\lfloor x^{400}\rfloor$ with budget $\exp(x^\varepsilon/2)$.
\end{theorem}
\begin{proof}
Apply Theorem~\ref{B2.48} to \(M=\operatorname{cond}\psi\), with the given \(\varepsilon,c\), and take \(y=h\). Its \(J,Z\) have the required quantifiers. At each \(x\) take \(U=Y\) and \(a=e\) in \eqref{eq:window-datum}, keeping its primes \(q,d\). Equation~\eqref{eq:window-vv} gives the assertion for every field of Definition~\ref{B6.02}, since \(\sqrt{\exp(x^\varepsilon)}=\exp(x^\varepsilon/2)\).
\end{proof}

In Theorem~\ref{B2.59}, the rank $t$ is a natural number, so
Lemma~\ref{B6.61} gives $t\le\lfloor x^{400}\rfloor$ at the scales used
there. The window estimate therefore covers every required index
$1\le j<t$.

\part{Input one: Green}

% ---- p7c1 ----
\chapter{Green's classification and its use}\label{ch:green}

\section{Green's theorem}
For a prime $p$ put $G_d=\GL_d(\F_p)$, $N_d=p^d-1$, and $\Delta_d=\prod_{1\le j<d}(p^j-1)$. Write $\mu_N$ for the group of $N$th roots of unity in the indicated field.

\begin{definition}[Cuspidality]\label{B7.01}
A cuspidal of $G_d$ for $d>0$ is an irreducible character with mean zero on the standard unipotent radicals of types $(a,d-a)$, $0<a<d$.
\end{definition}

\begin{definition}[Parameters]\label{B7.02}
Let $\mathcal P_d$ consist of monic irreducibles $P$ of degree $d$ over $\F_p$ with $P(0)\ne0$. For $[K:\F_p]=d$, a character $\theta\in\widehat{K^\times}$ is regular if the $\theta^{p^i}$, $0\le i<d$, are distinct.
\end{definition}

\begin{definition}[Label datum]\label{B7.03}
A label datum $L$ consists of a unital multiplicative map $t:\overline{\F}_p\to\C$ injective off zero and generators $u_d$ of $\mu_{N_d}(\overline{\F}_p)$ with $u_{de}^{N_{de}/N_d}=u_d$ for $d,e>0$. For $x\in\mu_{N_d}(\overline{\F}_p)$ set $\Theta_d(x,u_d^a)=t(x)^a$, with value zero off $\langle u_d\rangle$ in the second argument.
\end{definition}

\begin{definition}[Prime label]\label{B7.04}
An elliptic $y\in G_d$ has irreducible characteristic polynomial. A cuspidal $\chi$ has label $P\in\mathcal P_d$ if for every root $x$ of $P$, every elliptic $y$, and every characteristic root $z$ of $y$,
\begin{equation}\label{eq:green-label}
 \chi(y)=(-1)^{d-1}\sum_{i<d}\Theta_d(x,z^{p^i}).
\end{equation}
\end{definition}

\begin{definition}[Regular form]\label{B7.05}
For every prime $p$ and $[K:\F_p]=d>0$, there is a family $\tau_\theta$ indexed by $\widehat{K^\times}$ that induces a bijection from regular Frobenius orbits to cuspidals. For regular $\theta$, one has $\tau_\theta(1)=\Delta_d$ and
\[
 \tau_\theta(y)=(-1)^{d-1}\sum_{i<d}\theta(z^{p^i})
\]
for elliptic $y$ and each characteristic root $z\in K^\times$.
\end{definition}

\begin{definition}[Polynomial form]\label{B7.06}
Green asserts that for every prime $p$, datum $L$ and $d>0$, there is a family $\tau_P$ on $\F_p[X]$ that restricts to a bijection from $\mathcal P_d$ to cuspidals. For $P\in\mathcal P_d$, $\tau_P$ has label $P$ and degree $\Delta_d$.
\end{definition}

\subsection{Construction of the regular family}
Let $V$ have dimension $n$ over a finite field $F$ of size $v$. Put $H=\GL(V)$, let $c_g=\det(XI-g)$ for $g\in\operatorname{End}_F(V)$, and set $I(\phi,\psi)=\sum_{g\in H}\phi(g)\overline{\psi(g)}$ and $M=v^n-1$.

\begin{lemma}[Frobenius]\label{B7.07}
For irreducible $P\in\F_p[X]$ with root $x$, the roots are $(x^{p^i})_{0\le i<\deg P}$, and $x^{p^j}=x$ holds exactly when $\deg P\mid j$ for $j\ge0$. This is the prime-field case of the finite-field orbit theorem \cite[Ch.~2]{lidl-niederreiter-1997}.
\end{lemma}

\begin{lemma}[Lift]\label{B7.08}
There is a unital multiplicative map $t:\overline F\to\C$ injective off zero.
\end{lemma}
\begin{proof}
Take Teichmüller lifts in the Witt ring of $\overline F$. Reduction distinguishes the lifts, and the nonzero lifts are roots of unity. The part of the fraction field algebraic over $\Q$ therefore contains the lifts and embeds in $\C$.
\end{proof}

\begin{definition}[Root sums]\label{B7.09}
For $h\ge0$, set $s_h(P)=\sum_{P(\alpha)=0}t(\alpha)^h$, with multiplicities and $s_h(0)=0$. The multiset of roots of a monic product gives $s_h(PQ)=s_h(P)+s_h(Q)$ for monic $P,Q$. Also $s_h((X-1)^j)=j$.
\end{definition}

\begin{definition}[Subspace inversion]\label{B7.10}
For $s:F[X]\to\C$, set
\[
 \Psi_s(V,g)=s(c_g)-\sum_{\substack{W\subsetneq V\\gW\subseteq W}}\Psi_s(W,g|_W).
\]
Write $\Psi_h=\Psi_{s_h}$ and $R_h(g)=s_h(c_g)$.
\end{definition}

\begin{lemma}[Primitivity]\label{B7.11}
The recursion is conjugacy invariant. Call $s$ additive if $s(PQ)=s(P)+s(Q)$ for all monic $P,Q\in F[X]$. If $s$ is additive, then
\[
 \sum_{h:V/A\to A}\Psi_s(V,g+h)=0
 \qquad(0<A<V,\ gA\subseteq A),
\]
where $h$ is viewed as an endomorphism via quotient and inclusion. A class function on endomorphisms satisfying this identity is called primitive.
\end{lemma}
\begin{proof}[Outline]
Conjugation transports invariant subspaces. Average \ref{B7.10} over the off-diagonal maps, using induction for subspaces with partial intersection with \(A\) and counting lifts for disjoint ones, so the two surviving block sums add to \(s(c_g)\). The complete proof is in \S\ref{prf:B7.11}.
\end{proof}

\begin{lemma}[Values]\label{B7.12}
If $n>0$ and $s((X-1)^j)=j$ for all $j\ge0$, then $\Psi_s(1)=\prod_{1\le j<n}(1-v^j)$. For additive $s$, irreducible $c_g$ gives $\Psi_s(g)=s(c_g)$.
\end{lemma}
\begin{proof}
With $b_{n,a}$ the number of $a$-dimensional subspaces, counting maps $F^n\to F^j$ by their kernels gives, at $Y=v^j$,
\[
 \sum_{a=0}^n b_{n,a}\prod_{i=0}^{a-1}(Y-v^i)=Y^n.
\]
Differentiate the polynomial identity at $Y=1$ to obtain the recursion for the proposed values, with zero in dimension zero. For irreducible $c_g$, the sole proper invariant subspace is zero.
\end{proof}

\begin{lemma}[Virtual characters]\label{B7.13}
For $h\ge0$, both $R_h$ and $\Psi_h$ restrict to virtual characters of $H$.
\end{lemma}
\begin{proof}[Outline]
Over a finite extension containing the eigenvalues of \(H\), Green's Theorem 1 applied to \(\sum_i X_i^h\) shows \(R_h\) is a virtual character. Reading \ref{B7.10} as parabolic induction of \(\Psi_h\) tensored with the trivial character proves \(\Psi_h\) by induction. The complete proof is in \S\ref{prf:B7.13}.
\end{proof}

\begin{lemma}[Mass]\label{B7.14}
Write $M_\phi(P)=\sum_{c_g=P}\phi(g)$ over endomorphisms. For additive $s$ and monic $P$ of degree $n>0$, the mass for $\Psi_s$ is $|H|s(Q)/(v^n-1)$ if $P=Q^a$ for monic irreducible $Q$ and $a>0$, and zero otherwise.
\end{lemma}
\begin{proof}[Outline]
Count invariant subspaces in \ref{B7.10} and count vectors by their cyclic spans to get two convolution identities for the normalized masses. Use additivity and unique factorization to solve for the mass of \(\Psi_s\). The complete proof is in \S\ref{prf:B7.14}.
\end{proof}

\begin{lemma}[Torus pairing]\label{B7.15}
For primitive class $\phi$ and $n>0$, put $P_\alpha=\min_F(\alpha)$ and $a_\alpha=n/\deg P_\alpha$ for $\alpha\in\mu_M(\overline F)$. Then
\[
 I(\phi,R_l)=\sum_{\alpha\in\mu_M}a_\alpha M_\phi(P_\alpha^{a_\alpha})\overline{t(\alpha)^l},
 \qquad
 \sum_{\alpha\in\mu_M}t(\alpha)^h\overline{t(\alpha)^l}=M\mathbf1_{h\equiv l\bmod M}.
\]
\end{lemma}
\begin{proof}
Split a nonprimary polynomial into coprime blocks. The Sylvester map on the off-diagonal block is invertible, so all variations are conjugate and primitivity forces the value to vanish. For $c_g=P^a$ in $H$, the roots lie in $\mu_M$ since $\deg P\mid n$, and each contributes with multiplicity $a$. This gives the first identity. The second identity follows because the restriction of $t$ bijects the two groups of $M$th roots.
\end{proof}

\begin{lemma}[Mixed pairing]\label{B7.16}
For $n>0$, $I(\Psi_h,R_l)=|H|\#\{i<n:hv^i\equiv l\bmod M\}$.
\end{lemma}
\begin{proof}
Insert \ref{B7.14} and $a_\alpha s_h(P_\alpha)=\sum_{i<n}t(\alpha)^{hv^i}$ into \ref{B7.15}.
\end{proof}

\begin{proposition}[Expansion]\label{B7.17}
For $n>0$ and primitive class $\phi$, $n|H|\phi=\sum_{0\le l<M}I(\phi,R_l)\Psi_l$ on $H$.
\end{proposition}
\begin{proof}[Outline]
Use \ref{B7.15} and \ref{B7.16} to show the difference has zero pairings and hence zero primary masses by Fourier inversion. On each primary fibre, descend on the kernel dimension using \ref{B7.11} and use \ref{B7.14} to show that mass determines the values. The complete proof is in \S\ref{prf:B7.17}.
\end{proof}

\begin{lemma}[Inverted pairing]\label{B7.18}
For primitive $\phi$ and additive $s$, $I(\phi,\Psi_s)=I(\phi,s\circ c)$.
\end{lemma}
\begin{proof}
The term $W=0$ in \ref{B7.10} vanishes. For each $0<W<V$, fix the diagonal blocks of $g$ stabilizing $W$. The restriction to $W$ is constant while the off-diagonal block varies, so primitivity cancels the pairing.
\end{proof}

\begin{lemma}[Regular exponent]\label{B7.19}
If $n>0$ and $hv^i\not\equiv h\bmod M$ for $0<i<n$, then $I(\Psi_h,\Psi_h)=|H|$. Its restriction to $H$ is $(-1)^{n-1}$ times an irreducible character.
\end{lemma}
\begin{proof}
The norm follows from \ref{B7.18} and \ref{B7.16}. By \ref{B7.13} the irreducible coefficients are integers, so normalized norm one leaves a single coefficient $\pm1$. Its sign follows from \ref{B7.12}.
\end{proof}

\begin{lemma}[Cuspidality]\label{B7.20}
Over $\F_p$, the irreducible character of \ref{B7.19} is cuspidal.
\end{lemma}
\begin{proof}
Apply \ref{B7.11} at $g=1$, where the $1+h$ form the unipotent radical for $A$.
\end{proof}

\begin{theorem}[Regular family]\label{B7.21}
For $[K:\F_p]=d>0$, there exists a family $(\tau_\theta)_{\theta\in\widehat{K^\times}}$ with the following properties on regular parameters. Its members are cuspidal and have the values of \ref{B7.05}. They agree exactly within Frobenius orbits.
\end{theorem}
\begin{proof}
Embed $K$ in $\overline{\F}_p$ and write $\theta=t^h$ on $K^\times$ using \ref{B7.08}. Take $(-1)^{d-1}\Psi_h|_{G_d}$ for regular parameters by \ref{B7.19} and \ref{B7.20}, and zero elsewhere. The values follow from \ref{B7.07} and \ref{B7.12}. The pairings detect the orbit by \ref{B7.16} and determine the character by \ref{B7.17}.
\end{proof}

\begin{lemma}[Orbit injection]\label{B7.22}
For $[K:\F_p]=d>0$, an isomorphism $K^\times\simeq\widehat{K^\times}$ sends chosen roots of $P\in\mathcal P_d$ to characters $c(P)$ in distinct regular orbits.
\end{lemma}
\begin{proof}
Both groups are cyclic of order $N_d$. The isomorphism commutes with $p$th powers, and each chosen root has an orbit of size $d$ that determines its minimal polynomial by \ref{B7.07}.
\end{proof}

\begin{lemma}[Character injection]\label{B7.23}
The map $L_d(P)=\tau_{c(P)}$ injects $\mathcal P_d$ into the cuspidals of $G_d$.
\end{lemma}
\begin{proof}
Apply \ref{B7.21} to the orbits of \ref{B7.22}.
\end{proof}

Let $I(C)$ be the Harish--Chandra product of a cuspidal list $C$ of total degree $n$. For nontrivial additive $\eta$ of $\F_p$, set $\xi_\eta(u)=\eta(\sum_i u_{i,i+1})$ on the upper unitriangular subgroup $U_n$. Write $\operatorname{wh}_\eta$ for the linearly extended $\xi_\eta$-multiplicity. On a finite group $G$ use
\[
 \langle\phi,\psi\rangle=\frac1{|G|}\sum_{g\in G}\phi(g)\psi(g^{-1}).
\]
For characters of $H$ this agrees with $I(\phi,\psi)/|H|$.

\begin{theorem}[Generic constituent]\label{B7.24}
Every $I(C)$ has an irreducible constituent of positive Whittaker multiplicity.
\end{theorem}
\begin{proof}[Outline]
The finite-field Gelfand–Graev theorem gives each cuspidal factor Whittaker multiplicity one. In Mackey restriction to \(U_n\), Frobenius reciprocity gives one from the block-reversing double coset, so some irreducible constituent has positive Whittaker multiplicity. The complete proof is in \S\ref{prf:B7.24}.
\end{proof}

\begin{theorem}[Harish--Chandra pairing]\label{B7.25}
For equal total degrees, $\langle I(C),I(C')\rangle=0$ unless $C'$ permutes $C$. Distinct entries in $C$ give $\langle I(C),I(C)\rangle=1$. See \cite[pp.~107--148]{zelevinsky-1981} for the general linear groups over finite fields.
\end{theorem}
\begin{proof}[Outline]
Writing \(C'=(\sigma,C'_0)\), reciprocity and Mackey restriction expand the pairing over block distributions into the Levi factors. Cuspidality kills split blocks and multi-block products, leaving \(\sum_{i:C_i=\sigma}\langle I(C\setminus i),I(C'_0)\rangle\) for induction. The complete proof is in \S\ref{prf:B7.25}.
\end{proof}

\begin{theorem}[Gelfand--Graev norm]\label{B7.26}
For $n>0$, $\Gamma_\eta=\operatorname{Ind}_{U_n}^{G_n}\xi_\eta$ satisfies $\operatorname{wh}_\eta(\Gamma_\eta)=p^n-p^{n-1}$ and $\langle\chi,\Gamma_\eta\rangle=\operatorname{wh}_\eta(\chi)$ for class functions $\chi$. Over a finite field the Gelfand--Graev character is multiplicity free and has \((p-1)p^{n-1}\) constituents in the present prime-field case \cite[Prop.~9.4]{zelevinsky-1981}\cite[\S1, p.~76]{deriziotis-gotsis-1998}. Its self-pairing therefore has this value. Frobenius reciprocity identifies pairing against the induced character with the \(\xi_\eta\)-multiplicity on restriction to \(U_n\), and linearity gives the formula for class functions.
\end{theorem}

\begin{corollary}[Generic count]\label{B7.27}
For $n>0$ the number of generic irreducibles is at most $p^n-p^{n-1}$.
\end{corollary}
\begin{proof}
Each contributes at least one to the norm in \ref{B7.26}.
\end{proof}

\begin{lemma}[Support count]\label{B7.28}
For $n>0$ there are at most $p^n-p^{n-1}$ cuspidal multisets of total degree $n$.
\end{lemma}
\begin{proof}
The generic constituents supplied by \ref{B7.24} for distinct multisets are distinct by \ref{B7.25} and are counted by \ref{B7.27}.
\end{proof}

\begin{lemma}[Polynomial count]\label{B7.29}
There are $(p-1)p^{n-1}$ monic degree-$n$ polynomials with nonzero constant term for $n>0$.
\end{lemma}
\begin{proof}
There are $p-1$ choices for the constant coefficient and $p$ for each of the other $n-1$ lower coefficients.
\end{proof}

\begin{lemma}[Exhaustion]\label{B7.30}
Injective cuspidal families $L_e$ on $\mathcal P_e$ for all $e>0$ exhaust the cuspidals.
\end{lemma}
\begin{proof}
Replace each irreducible factor $P$ of a polynomial counted in \ref{B7.29} by $(\deg P,L_{\deg P}(P))$. This embeds the polynomials into the multisets of \ref{B7.28} and attains the bound. A missing degree-$n$ cuspidal would add a singleton.
\end{proof}

\begin{theorem}[Green's classification]\label{B7.31}
The regular form \ref{B7.05} holds. It takes the prime-field cuspidals and their values at the identity and at elliptic elements from Green's full classification \cite[\S\S7--8, Thm.~14]{green-1955}.
\end{theorem}

\section{Polynomial labels}
Fix $L$ and $\iota:K=\F_{p^d}\hookrightarrow\overline{\F}_p$, $d>0$.

\begin{lemma}[Torus]\label{B7.32}
The image of $\iota$ is the root set of $X^{p^d}-X$. The preimage $v$ of $u_d$ generates $K^\times$.
\end{lemma}
\begin{proof}
The $p^d$ elements of $\iota(K)$ solve the equation, and $u_d$ has order $p^d-1$.
\end{proof}

\begin{lemma}[Duality]\label{B7.33}
The map $x\mapsto\theta_x$, defined by $\theta_x(v^a)=t(x)^a$, is a bijection from $\mu_{N_d}(\overline{\F}_p)$ onto $\widehat{K^\times}$. It satisfies $\theta_x(k)=\Theta_d(x,\iota(k))$.
\end{lemma}
\begin{proof}
Evaluation at $v$ identifies the character group with $\mu_{N_d}(\C)$, which the $N_d$ distinct $t(x)$ exhaust. The identity follows from \ref{B7.03}.
\end{proof}

\begin{lemma}[Regularity]\label{B7.34}
Evaluation at $v$ gives $\theta_{x^a}=\theta_x^a$ for $x\in\mu_{N_d}$, $a\ge0$. Thus $\theta_x$ is regular exactly when the $x^{p^i}$, $i<d$, are distinct.
\end{lemma}

\begin{lemma}[Label roots]\label{B7.35}
Each $P\in\mathcal P_d$ has $d$ distinct roots in $\mu_{N_d}$ by \ref{B7.07} and $P(0)\ne0$.
\end{lemma}

\begin{lemma}[Root sums]\label{B7.36}
For irreducible $f$ with root $z$, summing a function along $(z^{p^i})_{i<\deg f}$ sums it over the roots. For $x,w\in\mu_{N_d}$ and $a\ge0$, $\Theta_d(x^a,w)=\Theta_d(x,w^a)$ by \ref{B7.03}.
\end{lemma}

\begin{theorem}[Relabelling]\label{B7.37}
The regular form \ref{B7.05} implies \ref{B7.06}.
\end{theorem}
\begin{proof}[Outline]
Choose a root giving a regular character by \ref{B7.34} and \ref{B7.35}. Use \ref{B7.32} and \ref{B7.33} to evaluate at elliptic elements and \ref{B7.36} for independence of the root, then identify the regular orbits with degree-\(d\) minimal polynomials. The complete proof is in \S\ref{prf:B7.37}.
\end{proof}
The unitary-group hypotheses use \ref{B7.31}. The polynomial form is used in Theorem~\ref{B1.03}.

\section{Series weights}

\begin{lemma}[Proper subfields]\label{B7.38}
For $p\ge5$, $d>0$, the set $B$ of elements of degree $<d$ in $\F_{p^d}^\times$ satisfies $2d|B|<N_d$.
\end{lemma}
\begin{proof}
For $d>1$, put $h=\lfloor d/2\rfloor$. Every proper subfield has degree at most $h$, so $|B|\le S_h=\sum_{e=1}^h(p^e-1)$. Induction gives $4S_h+4h+5\le5p^h$ and $5h\le p^h$, hence $4hS_h<p^{2h}-1$. This suffices for even $d$. For odd $d>1$, $2dS_h<8hS_h<2p^{2h}<N_d$. In degree one $B$ is empty.
\end{proof}

\begin{lemma}[Power sums]\label{B7.39}
Let two lists of length $d>0$ in $\mu_M(\C)$, $M>1$, have equal power sums outside $B\subset\Z/M$, with $2d|B|<M$. Then the first entry of the second list occurs in the first.
\end{lemma}
\begin{proof}
The multiplicity differences $c_\zeta$ have support of size at most $2d$ and Fourier support in $B$. Parseval and Cauchy--Schwarz give
\[
 M\sum_\zeta|c_\zeta|^2\le2d|B|\sum_\zeta|c_\zeta|^2.
\]
Thus all differences vanish.
\end{proof}

\begin{theorem}[Unique labels]\label{B7.40}
For $p\ge5$, two prime labels of one cuspidal agree. Under Green, two cuspidals with the same prime label agree.
\end{theorem}
\begin{proof}
Evaluate \eqref{eq:green-label} on multiplication by $u_d^a$ of degree $d$. The lifted root lists have equal power sums outside the exponents corresponding to $B$ in \ref{B7.38}. By \ref{B7.39} and injectivity of $t$, the two minimal polynomials share a root. Under Green, each cuspidal belongs to the family and its unique label fixes its parameter.
\end{proof}

\begin{definition}[Series]\label{B7.41}
An irreducible $\rho$ has label $f$ if it occurs in $I(C)$ for a cuspidal list whose prime labels multiply to $f$. Sum the corresponding central idempotents $z_\rho$ to get $E_f$. For $a\in\C[G_n]$ put $w_a(f)=\Re[aE_f]_1/\Re[a]_1$, where $[\,]_1$ is the identity coefficient and a zero denominator gives weight zero.
\end{definition}

\begin{theorem}[Squarefree series]\label{B7.42}
For $p\ge5$ under Green, each irreducible of $G_n$ has a unique monic degree-$n$ label with nonzero constant term. For squarefree such $f$, its series consists of $\rho_f=I((\deg P,\tau_P)_{P\mid f})$, over monic irreducible factors.
\end{theorem}
\begin{proof}
A minimal Levi with nonzero Harish--Chandra restriction supplies cuspidal factors by transitivity. Reciprocity gives occurrence in their product. The support is unique by \ref{B7.25}, and the prime labels are unique by \ref{B7.40}. For squarefree $f$, the factors are distinct and determined by $f$. Their product is invariant under permutation and has norm one by \ref{B7.25}.
\end{proof}

\begin{lemma}[Arithmetic relabelling]\label{B7.43}
Assume $p\ge5$ and Green. Let $O$ label simple cuspidal representations by $\mathcal P_d$ in each degree $d>0$, with equality equivalent to isomorphism. Let $E^O_{r,f}$ be the rank-$r$ series projector for the product $f$ of $O$-labels. There are inverse degreewise permutations $\sigma,\pi$ of the prime labels. On monic $f$ with $f(0)\ne0$ their multiplicative extensions preserve degree and squarefreeness, with $\pi_*\sigma_*f=f$ and $E^O_{r,f}=E_{\sigma_*f}$.
\end{lemma}
\begin{proof}
For a simple cuspidal $V_P$ of character $\tau_P$, set $\pi(P)=O(V_P)$. The isomorphism criterion and Green make $\pi$ well-defined and injective. Finiteness gives bijectivity. Its inverse $\sigma$ gives $\sigma_*f=\prod_P\sigma(P)^{v_P(f)}$. Both projectors select the same supports, and the substitution preserves the degrees and multiplicities.
\end{proof}

For a subgroup $H$ and a degree-one character $\xi$, put $e(H,\xi)=|H|^{-1}\sum_{u\in H}\xi(u)^{-1}u$. Set $e_U=e(U_n,\xi_\eta)$. For $n=k+2s+1$, $k\ge1$, $s\ge0$, use zero-based indices. Let $H_J\subset U_n$ impose $u_{a,s+k}=0$ for $s\le a<s+k$. Set $e_J=e(H_J,\eta\circ\phi_J)$, where $\phi_J$ sums consecutive entries along $(0,\dots,s-1,s+k,\dots,n-1)$ and $(s,\dots,s+k-1)$. Write $w_U=w_{e_U}$ and $w_J=w_{e_J}$.

\begin{lemma}[Idempotents]\label{B7.44}
The averages $e_U,e_J$ are idempotent with identity coefficients $p^{-b},p^{k-b}$, where $b=n(n-1)/2$. For idempotent $a$ with $\Re[a]_1\ge0$, $w_a(f)\ge0$. The coefficient $[az_\rho]_1$ is $\rho(1)/|G_n|$ times the rank of $a$ in the representation of $\rho$.
\end{lemma}
\begin{proof}
Each average projects onto its character space. The $k$ zero entries give $|H_J|=p^{b-k}$. Left multiplication by $az_\rho$ has trace $|G_n|[az_\rho]_1$ on the regular representation. Only the $\rho(1)$ copies of $\rho$ contribute, with trace equal to rank on each. This also gives nonnegativity.
\end{proof}

\begin{lemma}[Multiplicity one]\label{B7.45}
For a cuspidal character $\chi$ of $G_d$ with $\chi(1)=\Delta_d$, we have $\operatorname{wh}_\eta(\chi)=1$. Every complex irreducible cuspidal character occurs in the Gelfand--Graev character over a finite field, and each constituent occurs once \cite[Prop.~9.4]{zelevinsky-1981}\cite[\S1, p.~76]{deriziotis-gotsis-1998}. Apply reciprocity to the nondegenerate character $\xi_\eta$ of $U_d$. The published result allows every finite field and requires no degree hypothesis.
\end{lemma}

\begin{theorem}[Squarefree weight]\label{B7.46}
Under \ref{B7.42}, for monic squarefree $f$ of degree $n$ with $f(0)\ne0$, $w_U(f)=\prod_{P\mid f}(p^{\deg P}-1)^{-1}$ over monic irreducibles.
\end{theorem}
\begin{proof}
The parabolic index times the cuspidal degrees gives $\rho_f(1)=|G_n|p^{-b}\prod_{P\mid f}(p^{\deg P}-1)^{-1}$. The Whittaker multiplicity is one by \ref{B7.45} and \ref{B7.48}. Apply \ref{B7.44} to $\rho_f$.
\end{proof}

\begin{lemma}[Column exchange]\label{B7.47a}
Let $e'_J$ be the $\xi_\eta$-average on the subgroup of $U_n$ given by $u_{a,k}=0$ for $a<k$. For $n=k+2s+1$ and every class function $\chi$ of $G_n$, extended linearly to the group algebra, $\chi(e_J)=\chi(e'_J)$.
\end{lemma}
\begin{proof}[Outline]
Move the inserted chain to the front one step at a time. The nondegenerate commutator pairing gives Mackey self-pairings of one, and the common intersection gives a positive cross-pairing, so induction and reciprocity preserve the trace. The complete proof is in \S\ref{prf:B7.47a}.
\end{proof}

\begin{lemma}[Mirabolic trace]\label{B7.47b}
For $0\le i\le n$, let $D_i^\eta$ take Harish--Chandra restriction to $(i,n-i)$, then the Whittaker multiplicity space of the second block, with the same notation for its character. Under \ref{B7.47a}, for every class function $\chi$,
\begin{equation}\label{eq:mirabolic-trace}
 \chi(e_J)=\sum_{i<k}(p^{k-i}-p^{k-i-1})\operatorname{wh}_\eta(D_i^\eta\chi)
 +\operatorname{wh}_\eta(D_k^\eta\chi).
\end{equation}
\end{lemma}
\begin{proof}[Outline]
Decompose the mirabolic tail average by column-character orbits until the first zero step. Mackey restriction gives Harish–Chandra products whose Whittaker multiplicities use \ref{B7.48} and \ref{B7.26}, and \ref{B7.47a} transfers the trace. The complete proof is in \S\ref{prf:B7.47b}.
\end{proof}

\begin{theorem}[Lower comparison]\label{B7.47}
For any prime $p$, datum $L$, $n=k+2s+1$ with $k\ge1$, and every $f$, $(1-p^{-1})w_U(f)\le w_J(f)$.
\end{theorem}
\begin{proof}
For irreducible $\chi$ all terms of \ref{B7.47b} are nonnegative. The $i=0$ term is $(p^k-p^{k-1})\chi(e_U)$. Sum the coefficient formula of \ref{B7.44} over the series and divide by the identity-coefficient ratio $p^k$.
\end{proof}

\begin{lemma}[Derivative counts]\label{B7.48}
Let $C$ have total degree $n$ and cuspidal factors of Whittaker multiplicity one. Then $\operatorname{wh}_\eta(I(C))=1$, and $\operatorname{wh}_\eta(D_i^\eta I(C))$ counts subsets of positions of total degree $i$, for $0\le i\le n$.
\end{lemma}
\begin{proof}[Outline]
Mackey restriction distributes inducing blocks between the Levi factors, with cuspidality killing split blocks. For each surviving distribution, \ref{B7.24} gives \(1\) on the opposite cell while other cells vanish on a simple root subgroup where \(\xi_\eta\) is nontrivial. The complete proof is in \S\ref{prf:B7.48}.
\end{proof}

\begin{theorem}[Upper comparison]\label{B7.49}
Under \ref{B7.46}, if $n=k+2s+1$ and $k\ge1$, then $w_J(f)\le(1+p^{-1})^nw_U(f)$.
\end{theorem}
\begin{proof}
For the subset counts $c_i$ of \ref{B7.48}, the trace lemma \ref{B7.47b} and \ref{B7.44} give
\[
 \frac{w_J(f)}{w_U(f)}
 =p^{-k}\rho_f(e_J)\le\sum_{i\le k}c_ip^{-i}
 \le\prod_{P\mid f}(1+p^{-\deg P})\le(1+p^{-1})^n.
\]
\end{proof}

\begin{theorem}[Arithmetic weights]\label{B7.50}
Assume Green. Let $D$ be an arithmetic point datum with host CM field $E$, spectral rank $r$, and distinguished primes $p_0,p_1$. Fix an embedding $j_E:E\hookrightarrow\overline{\Q}$. At each $p_j$ use the normalized local types for this embedding: the Teichmüller lift and compatible generators give the label datum of \ref{B7.03}, and the arithmetic projector for a polynomial $f$ is the series projector for its cuspidal labelling as in \ref{B7.43}, with value zero off the monic degree-$r$ labels of nonzero constant term. At each type prime $p_j$ and for nontrivial $\eta$, the arithmetic $U,J$ weights are probabilities supported on monic degree-$r$ labels with nonzero constant term. They satisfy \ref{B7.47} on these labels and \ref{B7.46}, \ref{B7.49} on squarefree ones.
\end{theorem}
\begin{proof}[Outline]
Use \ref{B7.43} to relabel the projectors degreewise, preserving squarefreeness. The projectors sum to the identity by \ref{B7.42}, and \ref{B7.44} gives nonnegative weights whose numerators sum to the denominator. The complete proof is in \S\ref{prf:B7.50}.
\end{proof}

\section{Bad rows}
\subsection{Factor degrees at one place}
Let $K$ be a finite field of cardinality $u$. Write $\mathcal S_N(K)$ for the monic squarefree polynomials of degree $N$ with nonzero constant term. For monic $f$ put
\[
 w(f)=\prod_{P^a\parallel f}(u^{\deg P}-1)^{-a},
 \qquad M_N(E)=\sum_{f\in\mathcal S_N(K)\cap E}w(f),
\]
where $P$ runs over monic irreducibles. Write $\lambda(f)$ for the multiset of factor degrees, with multiplicities, and $\lambda(\sigma)$ for the multiset of cycle lengths, including fixed points.

\begin{lemma}[Comparison by factor degrees]\label{B7.50a}
For $N\ge0$ and any collection $\mathcal A$ of partitions of $N$,
\[
 M_N\{f:\lambda(f)\in\mathcal A\}
 \le \frac{|\{\sigma\in S_N:\lambda(\sigma)\in\mathcal A\}|}{N!}.
\]
If $h$ is monic with $\deg h\le N$, then $M_N\{f:h\mid f\}\le w(h)$. For $N\ge1$ the total mass lies between $1-2/u$ and $1$.
\end{lemma}
\begin{proof}[Outline]
Bound the number of degree-\(j\) irreducibles with nonzero constant term by \((u^j-1)/j\) to compare each factorization type with its cycle proportion. Divide by a prescribed divisor to bound its mass, and count excluded polynomials for the total mass lower bound. The complete proof is in \S\ref{prf:B7.50a}.
\end{proof}

\begin{lemma}[A divisor of prescribed degree]\label{B7.50b}
For every finite field $K$ and all integers $0<b<N$,
\[
 M_N\{f:\exists h\in K[X],\ h\mid f,\ \deg h=b\}\le 7/8.
\]
\end{lemma}
\begin{proof}[Outline]
Reduce by \ref{B7.50a} to uniform permutations, with \(X\) counting invariant \(b\)-sets (\(\mathbb{E}X=1\)) and \(Y\) those with fixed points inside and out. Inclusion-exclusion gives \(\mathbb{E}Y\ge 1/4\), and swapping gives \(2\mathbf{1}_{X>0}+Y\le 2X\), hence \(7/8\). The complete proof is in \S\ref{prf:B7.50b}.
\end{proof}

\begin{lemma}[Prime degrees in the upper half]\label{B7.50c}
For every integer $N\ge100$ let $\mathcal J_N=\{b\text{ prime}:N<2b,\ b+3\le N\}$. Then
\[
 H_N=\sum_{b\in\mathcal J_N}\frac1b\ge \frac1{5\log N}.
\]
\end{lemma}
\begin{proof}[Outline]
For \(N<20000\), tabulate reciprocal lower bounds using primes shared by each interval. For larger \(N\), bound the prime factors of the central binomial coefficient to count primes in the upper half, discarding the last three possible values. The complete proof is in \S\ref{prf:B7.50c}.
\end{proof}

\begin{lemma}[Missing a long prime-degree factor]\label{B7.50d}
For every finite field $K$ and integer $N\ge100$,
\[
 M_N\{f:\text{no irreducible factor of }f\text{ has degree in }\mathcal J_N\}
 \le 1-\frac1{5\log N}.
\]
\end{lemma}
\begin{proof}
For $b\in\mathcal J_N$ a $b$-cycle is unique if it occurs. Choosing and ordering its support gives $\binom Nb(b-1)!(N-b)!=N!/b$ permutations. The events for distinct $b$ are disjoint since the lengths exceed $N/2$. The complementary probability is $1-H_N$. Apply \ref{B7.50a} and \ref{B7.50c}.
\end{proof}

For prime $\ell$ and nonzero $x$ of finite order, write $x_{\ell'}$ for the prime-to-$\ell$ component of $x$ in $\langle x\rangle$. The projection is multiplicative and commutes with Frobenius. Put
\[
 B(u,N)=\{u\}\cup\{\ell\text{ prime}:\exists\,1\le j\le N,\ u^j<\ell^{4v_\ell(u^j-1)}\}.
\]

\begin{lemma}[Collisions of prime-to-$\ell$ parts]\label{B7.50e}
Take roots in an algebraic closure of $K$. For $u\ge4096$ and $N\ge0$, the $M_N$-mass of polynomials with distinct roots $x,y$ satisfying $x_{\ell'}=y_{\ell'}$ for some prime $\ell\notin B(u,N)$ is at most $23\log u/\sqrt u$.

Let $N\ge1$, $c\in K^\times$, and let $q,d$ be distinct odd primes with $\max(q,d)\ge16$ and $16,q,d\mid u-1$. The $M_N$-mass of polynomials with a root $x$ satisfying $x_{\ell'}=c_{\ell'}$ for some prime $\ell\ne u$ is at most
\[
 \frac1{4\min(q,d)}+\frac{1+\log N}{u-1}.
\]
\end{lemma}
\begin{proof}[Outline]
Count collisions within and across factors using Frobenius orbits and projection fibres, and sum the divisor bounds of \ref{B7.50a}. For a fixed target, count roots of prime-power binomials and use prime valuations to sum the bounds by degree. The complete proof is in \S\ref{prf:B7.50e}.
\end{proof}

\begin{corollary}[Exceptional factorisations in one row]\label{B7.50y}
Under the hypotheses of \ref{B7.50e} with $u=p\ge4096$ prime and $N=r\ge100$, let $P$ be a probability on the monic degree-$r$ polynomials over $\F_p$ with masses $w(f)$ on $\mathcal S_r(\F_p)$. The complement has probability at most $2/p$. On $\mathcal S_r(\F_p)$ the bounds of \ref{B7.50b}, \ref{B7.50d} and \ref{B7.50e} hold for $P$.
\end{corollary}
\begin{proof}
The restriction is $M_r$, whose total mass is at least $1-2/p$ by \ref{B7.50a}.
\end{proof}

\subsection{Arithmetic rows}
Write $\nu=f$ for the number of places in each type-prime tuple and $m$ for the chain length. For arithmetic row weights $w_{A,j,v}$ put
\[
 P_j(S)=\sum_{t\in S}\prod_{v\mid p_j}w_{A,j,v}(t_v),\qquad
 W(\nu,m,r,p,q,d)=\frac{3\nu}{p}+\frac{23\nu\log p}{\sqrt p}
 +2\nu m\left(\frac1{4\min(q,d)}+\frac{1+\log r}{p-1}\right)
 +\frac r2(7/8)^\nu+e^{-\nu/(5\log r)}.
\]
At $p=p_j$ a good tuple consists of squarefree admissible polynomials $t_v$. For prime $\ell\notin B(p,r)$ of \ref{B7.50e}, the map $x\mapsto x_{\ell'}$ is injective on the roots of each $t_v$. For every prime $\ell\ne p$, these parts avoid the prime-to-$\ell$ parts of the prescribed $\chi_{v,a}^{\pm1}\in\F_p^\times$, $1\le a\le m$. For each $1\le b\le r/2$ some $t_v$ has no degree-$b$ divisor. Some $t_v$ has an irreducible factor of degree in $\mathcal J_r$.

\begin{lemma}[Bad-label weight]\label{B7.51}
For an arithmetic point datum with host field $E$, assume the weight laws of \ref{B7.50} and $r\ge100$. Require both distinguished primes $p_j$ to be at least $4096$, to split completely in $E$ and in the ray class field $M_{\rm ray}$ of the host subfield $K$ for modulus $Q_Sqd$, and to satisfy $p_j\equiv1\pmod{Q_Sqd}$. Suppose that for each $j\in\{0,1\}$ there is a nontrivial additive character $\eta_j$ such that the row idempotent at every place above $p_j$ is either $e_U(\eta_j)$ or $e_J(\eta_j)$. Then every row weight is nonnegative. For each $j$ and every finite set $S$ of tuples bad at $p_j$,
\[
 0\le P_j(S)\le W(\nu,m,r,p_j,q,d).
\]
\end{lemma}
\begin{proof}[Outline]
Match the collision tests to \ref{B7.50e} and use \ref{B7.50y} with independence and a union bound for the product \(U\)-law. Apply the lower comparison of \ref{B7.50} on the complementary event to transfer the bound to the row laws. The complete proof is in \S\ref{prf:B7.51}.
\end{proof}

\begin{definition}[Compatibility of an arithmetic row system]\label{def:p7c1-compatible-rows}
Fix an arithmetic point with host CM field \(E\), spectral rank \(r\), and distinguished primes \(p_0,p_1\). Call a row system compatible with this point if the following conditions hold.

The real places of the row system's totally real field are in bijection with the archimedean indices of the point. Its rank, totally real degree and root-discriminant parameter agree with \(r,f,\operatorname{rd}(E)\). At corresponding infinite places its signature and Harish--Chandra parameter agree with those prescribed by the point. The coordinates of the parameter are strictly decreasing and belong to \((r-1)/2+\Z\). The local representation is discrete series with this parameter, and the chosen vector is a unit vector in its lowest \(K\)-type.

The level integer is \(p_0p_1\) and its distinguished prime is \(p_0\). If \(p_0\) splits completely in \(E\), it splits completely in the row system's CM field. The full row dimension is \(D_{\rm all}\) of the point. The product of the archimedean formal degrees, the volume of the fundamental domain and the real part of the full finite-adelic kernel at the identity equals \(D^\star\) of the point.

For each \(j=0,1\) there is a nontrivial additive character \(\eta_j\) of \(\F_{p_j}\) such that at every row place above \(p_j\) the local idempotent is \(e_J(\eta_j)\).
\end{definition}

\begin{theorem}[Defect bound]\label{B7.52}
Fix $0<\varepsilon\le1$, $c>0$, $C_s\in\R$ and a family $F_m$ with $J\ge1$ and $\alpha=\varepsilon/J$. Write $I_x$ for the indices of the members of $F_m$ at $x$. For all sufficiently large $x$ and each $i\in I_x$, suppose there is a compatible arithmetic row system $A$ in the sense of \ref{def:p7c1-compatible-rows} for which the weight laws of \ref{B7.50} hold and
\[
 p_j\ge e^{8x^\alpha},\qquad x^2\le\nu\le x^{C_s},\qquad
 m\le x^{C_s},\qquad100\le r\le x^{C_s},\qquad
 \min(q,d)\ge e^{x^{\alpha/6}} .
\]
Assume the map from retained to all-row indices preserves labels and is injective with image exactly the indices good at both primes. For every set $U$ of label pairs, assume that convolution with the kernel for $U$ has image of complex dimension $\#\{a:\text{the label pair of }a\text{ lies in }U\}$ on the multiplicity space of the prescribed archimedean representation. Here $a$ ranges over all row indices.

Let $\Omega$ be the set of all label pairs, $e$ the finite-adelic kernel of the full row system, and $e_T(U)$ the kernel for $U\subseteq\Omega$. Write $C_j(S)=\{(t_0,t_1)\in\Omega:t_j\in S\}$. The kernel identities are
\begin{equation}\label{eq:green-cylinder}
 e_T(\Omega)=e,\qquad
 \Re\bigl(e_T(C_j(S))(1)\bigr)
 =\Re e(1)\sum_{t\in S}\prod_{v\mid p_j}w_{A,j,v}(t_v)
 \quad(j=0,1,\ S\text{ finite}).
\end{equation}
For each cylinder let $N_j(S)$ count its indices. Put
\[
 D_j(S)=d_\infty\operatorname{vol}(\mathscr F)\Re\bigl(e_T(C_j(S))(1)\bigr),\qquad
 D_*=d_\infty\operatorname{vol}(\mathscr F)\Re e(1),
\]
where $d_\infty$ is the product of the archimedean formal degrees and $\mathscr F$ is the automorphic fundamental domain with the row system's Haar measure. Assume $D_*>0$ and, for every cylinder,
\[
 |N_j(S)-D_j(S)|\le\epsilon_C D_j(S),\qquad
 \epsilon_C=e^{-cr\nu\log r}.
\]
Eventually, uniformly in $i$, the relation $D_{\rm rows}=(1-\delta_{x,i})D_{\rm all}$ between the retained and total counts gives $0\le\delta_{x,i}\le e^{-x^{\alpha/7}}$.
\end{theorem}
\begin{proof}[Outline]
Bound discarded rows by the two bad cylinders and use relative counting and the cylinder identities to bound the defect by their weights. Use \ref{def:p7c1-compatible-rows} for the idempotents in \ref{B7.51} and the size bounds for uniform decay. The complete proof is in \S\ref{prf:B7.52}.
\end{proof}

\begin{theorem}[Closed defect bound]\label{B7.53}
Assume Green. Fix $c>0$, a family class $\mathcal A$ and a twist selection. Suppose the following assertion is closed on $\mathcal D_3(\mathcal A)$. There is an $x_c$ such that for every $x\ge x_c$ and every $i\in I_x$, one compatible row system satisfies the index, dimension and relative-count conditions of \ref{B7.52}. Then the defect bound $0\le\delta_{x,i}\le e^{-x^{\alpha/7}}$ of \ref{B7.52} is closed on this class in the sense of \ref{B8.39}.
\end{theorem}
\begin{proof}[Outline]
Keep the size thresholds of \ref{B8.39} and take a common point threshold. The selected types give \ref{B7.50} on the same rows, and the family sizes verify the numerical hypotheses of \ref{B7.52} with \(C_s=6\). The complete proof is in \S\ref{prf:B7.53}.
\end{proof}

\section{Connected families}

\begin{theorem}[Connected families]\label{B7.54}
Assume Green's polynomial classification of \ref{B7.06}. Under the common hypotheses of \ref{B8.02}, assume that the family bound for the weighted local lengths in \ref{B8.40} is closed on $\mathcal D_3(\mathcal A)$ in the sense of \ref{B8.39}. Fix $\lambda>1$ and a request $G$ for which hosts exist on $\lambda$-windows in the sense of \ref{B1.18}. Let $\mathcal E$ consist of the families with $\lambda$-windows for which every member beyond the family's initial cutoff satisfies $G$. Suppose that in each family in $\mathcal E$, all members at sufficiently large $x$ satisfy the finiteness and weighted length bound of \ref{B6.01} for their host field $E$, with $1\le k<t$ and $B=\sqrt b$. Then there is an $\varepsilon_1>0$ such that for every $0<\varepsilon<\varepsilon_1$ there is a family of exponent $\varepsilon$ in $\mathcal E$ with the following property. There are constants $K,c_{HH},e_K,C_L,x_0$ such that every window point $x\ge x_0$ admits the common chain realisation of \ref{B5.33} with division and loss exponents $(e_D,e_L)=(2,0)$.
\end{theorem}
\begin{proof}[Outline]
Combine the closed estimates by \ref{B8.39} and choose a family meeting their thresholds. Use the local lengths via \ref{B8.40a} and \ref{B5.10} for division exponent \(2\) and common weights via \ref{B5.22} for loss exponent \(0\), then apply \ref{B5.33}. The complete proof is in \S\ref{prf:B7.54}.
\end{proof}


\part{Input three: the witness for HypsAlt}

% ---- p8c1 ----
\chapter{A common witness for the 46 conditions}\label{ch:p8c1}

Let \(F_m\) be a family of parameter \(\varepsilon\), with point indices \(i\in I_x\) at \(x>F_m.x_1\). Write \(D,L,O\) for the arithmetic data, size data and objects at \((x,i)\), respectively. Put \(b=\exp(x^\varepsilon)\), and let \(\iota_D:K_D\hookrightarrow\C\) be the coefficient embedding. Use the two products \(P,F\) of \ref{B1.15ac}. The join used for the test vectors is the join in the fourth part of the toric input of \ref{B1.14k}.

\section{Existence and the common data}

Use the family quantifiers \(\operatorname{Closed}_{\mathcal C}\) of Definition~\ref{B8.39} and the subclass \(\mathcal A_{\rm ch}\) of Definition~\ref{def:p4c3-design}. Put
\[
 W_1(\mathcal A)=\{F_m\in\mathcal A_{\rm ch}:F_m\text{ satisfies the smoothness and twisting conditions below}\},
 \qquad W_3(\mathcal A)=\mathcal D_3(\mathcal A)=\mathcal D(\mathcal A_{\rm ch},3).
\]
For a frame with group \(A\), let \(n_H=|A|\) and \(f=2n_H\). The rank at layer \(j\) is \(N(j)=k+j(t+1)\), where \(t\) is the rank increment parameter. Membership in \(\mathcal A_{\rm ch}\) requires tame distinctness as in clause~4 of \ref{B8.01k}, the zero-free boxes and auxiliary-prime conditions of \ref{B8.47}, and \(C_{YL}\le9\). In each of the latter two pointwise tests the scale \(Y\) may be chosen in \([x^{1/J},x]\), separately for the two tests. Membership also requires that eventually, at every point and for every word \(u\), its stabilizer denominator \(D_{\Gamma,u}\) divide
\[
 M(2fN(m)),\qquad
 M(a)=\prod_{\substack{p\le a+1\\p\ {\rm prime}}}
 p^{\sum_{j\ge0}\lfloor a/(p^j(p-1))\rfloor}.
\]
Smoothness means that beyond one family cutoff, at every point and for every word \(u\), each prime divisor of \(D_{\Gamma,u}\) be at most \(2fN(m)+1\). For the twisting condition use the family's own theta data and clause~7 of \ref{B8.01k} at power one. Clause~7 in the duty of a member of \(\mathcal A\) imposes the tests for every power prime to the residue characteristic. Thus
\[
 W_3(\mathcal A)(F_m)\ \Longleftrightarrow\
 W_1(\mathcal A)(F_m)\land
 (\exists CH\in\R:\ F_m\text{ satisfies the choice bounds of Definition~\ref{def:p4c3-design} at }3,CH)
 \land\ (F_m\text{ satisfies the level and stabilizer size bounds of Definition~\ref{def:p4c3-design}}).
\]

Let \(h_L\) assert that for every number field \(E\), algebraic closure \(\Omega/E\), finite place \(w\) and valuation subring of \(\Omega\) above \(w\), there exists a weak local Langlands correspondence over \(E_w\) in the sense of Definition~\ref{B9.47}. Theorem~\ref{B9.46} gives this existence. Fix \(c_{\rm Dep}\in\R\) and a condition \(A\) on the base, frame and twist at a point. Let \(\mathcal N_{\rm list}\) be the norm-choice extension with the real expansion data and the law that rational map values give rational column maps. It also requires a character grid at \((1,c_{\rm Dep},b)\) (Definition~\ref{B8.21c}). The coefficient field contains the layer character values and a primitive root of order \((N v)^{d_v}\) at each grid place of depth \(d_v\). The signed local character lists and their theta characters agree with the grid in both directions. An origin consists of a preliminary construction, a state satisfying the same grid condition, and a lift preserving its Galois and form data and its signed norm choices. On an untwisted base its chosen places form a union of fibres over the fixed coefficient subfield. With the extensions taken at the same point indices, put
\[
\begin{aligned}
 \mathcal N_{\rm orig}&=\{n\in\mathcal N_{\rm list}:\text{an origin induces the given twisting set and norm choices}\},\\
 \mathcal N_A&=\{n\in\mathcal N_{\rm orig}:A\text{ holds on its base, frame and twist}\},\\
 p_A&=(R_{\rm row},\mathcal N_A,\mathbf1,\top).
\end{aligned}\tag{8.6}
\]
Here \(R_{\rm row}\) is the prescribed comparison of the constructed and arithmetic rows. In particular it identifies their complex row forms at the reference and twisted tuples, as in Lemma~\ref{lem:p8c1-row-forms}. The last two entries prescribe a singleton of extra data and allow every geometry. For parameters \(p\), call a member (Definition~\ref{B8.21b}) constructed if its columns and later stages come from one weighted construction with the prescribed row, norm, division and theta data, and its kernel entries are those of the local place matrices of that construction. Call it good if it also admits twisting data for its own theta realization satisfying the local twisting and character estimates, with disjoint auxiliary prime sets at nonreference places of different twist indices. A point is linked to a bundle at \(c\) if a weighted point with the same arithmetic datum and objects realizes that bundle: its recursion, spectral realization, rational rows, expansions, weights and selected side agree under transport, the normalization scalar is one, and the column characters agree. For \(p_A\) denote these two member conditions and this relation by \(Q_A,Q'_A,L_A\), respectively. Put \(p_0(h_L,c_{\rm Dep})=(R_{\rm row},\mathcal N_{\rm list},\mathbf1,\top)\). Say that \(A\) holds eventually on bases if the following holds for all sufficiently small \(0<\varepsilon<1/52\). Given the host constants with \(C_8\ge1\), take a threshold record for the standard family parameters, followed in order by thresholds for \(c_D,c_h,x\). In every good context beyond them one can prescribe finite algebraic elements in the two coefficient fields, required primes and finite sets of primes to avoid at each index. Every base over that context and every frame and twist meeting these requirements satisfy \(A\). The first coefficient requirement may depend on the preliminary base, frame and first field, the prime requirements on both fields, and the second coefficient requirement also on the twist characters. Put
\[
 \mathrm{NL}(\mathcal C)=\operatorname{Closed}_{\mathcal C}[(F_m,CH)\mapsto F_m\text{ satisfies the family local-error estimate of Definition~\ref{B8.40}}].
\]
Use the definite-side condition of Definition~\ref{B8.35}. In its terminal clause, quantify over all families of integral submodules of the face coordinate spaces and all compatible three-sided division data \(I\) agreeing under transport with the division datum of the selected bundle. For each word \(u\), let \(e_u\) identify its face coordinates with the column space. The lattice in the column-chain assertion of \ref{B8.35} at that face is \(e_u^{-1}(N_u)\), where \(N_u\) is the input lattice of the word. The finiteness premise of \ref{B8.40} ranges over all primes.

\begin{theorem}[Existence of the third input]\label{B8.01}
Assume the two collections of hypotheses \(P,F\) of Definition~\ref{B1.15ac}, the theta-integral comparison and local matching of Theorem~\ref{B8.11}, and the covolume bound formulated in Definition~\ref{B8.24}. For the application with \(O\) as in \ref{B9.01}, Theorem~\ref{B8.25} supplies this covolume bound from the premises established in Lemma~\ref{B8.25a}; the family assembly is Proposition~\ref{B8.33}, with \(C_0=C_1+C_2+1\). There is a common choice satisfying Definition~\ref{B8.02}. Its predicate on a class of families is the closure, in the sense of Definition~\ref{B8.39}, of the local-error estimate of Definition~\ref{B8.40}. Its predicate on the definite side is Definition~\ref{B8.35}.
\end{theorem}

\begin{proof}[Outline]
Combines \(P,F\) with \ref{B1.14k} and \ref{B1.15ac} to obtain separated rational tests, sets \(C_F=\max(C_{F0},C_{F1})\) to align the functional and grid extensions, and projects the augmented extension with the covolume and \(2,3\) bounds. The complete proof is in \S\ref{prf:B8.01}.
\end{proof}

\begin{definition}[Local data for the twisting tests]\label{B8.01ka}
Fix a point of the family. Let \(E\) be its CM field, let \(G_E=\operatorname{Gal}(\overline E/E)\), and let \(r=k+m-1\). A twist index \(a\) specifies a residue prime \(\ell=\ell_a\) and coefficient discrete valuation rings \(R_0\subset R\subset R_h\) for \(h\ge1\). Here \(R_h\) is finite over \(R\) and contains a designated primitive \(\ell^h\)-th root of unity. Write \(K=\operatorname{Frac}R\). On nonzero elements of the fraction field of any of these rings \(A\), use the valuation
\[
 v_A(z)=\frac{\operatorname{ord}_A(z)}{\operatorname{ord}_A(\ell)}.
\]

The twisting datum specifies a finite set \(\Sigma_a\) of finite places of \(E\). It contains every place above \(\ell,g_0,g_1\) and every place above
\[
 p(a)=
 \begin{cases}
 p_{\rm de},&\ell\ne p_{\rm de},\\
 p'_{\rm de},&\ell=p_{\rm de}.
 \end{cases}
\]
The Hecke tests use all places \(P\notin\Sigma_a\) with \(P\ne\bar P\). For choosing the conductor there is a further exclusion. Put \(d_E^*=Q_Sqd\) and \(n_{\mathfrak c}=N\mathfrak c_K\). Let \(p_*\) be the least prime not dividing \(d_E^*(p_{\rm de}-1)(p'_{\rm de}-1)n_{\mathfrak c}\), and let \(q_0\) be the least prime congruent to \(1\bmod p_*d_E^*\) that does not divide \(n_{\mathfrak c}\). Let \(\mathcal V\) be the set of places of \(E\) above
\[
 3,\ 5,\ g_0,\ g_1,\ p_{\rm de},\ p'_{\rm de},\ q_0,
 \quad\text{the word primes }T_u,\quad
 \text{and the prime divisors of }n_{\mathfrak c}.
\]

For \(1\le j\le m\) and \(\sigma\in\{+1,-1\}\), let \(V\) be the space of forms over \(K\) at \((a,j,\sigma)\). The finite group acting on it is
\[
 G_{\rm typ}=\prod_{i=0}^1\prod_{v\ {\rm designated\ above}\ g_i}
                    \GL_r(\F_{g_i}).
\]
Write \(d=|G_{\rm typ}|\). The point gives a subgroup \(U\le G_{\rm typ}\), a character \(\xi_{a,\sigma}\) of \(U\), and the action \(P\) of the type projector on \(V\). Set
\[
 V^\circ=P\{v\in V:uv=\xi_{a,\sigma}(u)v\text{ for every }u\in U\}.
\]
The point's theta data specify a set \(\mathcal D_{aj\sigma}\) of allowed inputs. Each \(z\) in it has a vector \(\theta(z)\in V\) and a finite set of level places \(\mathcal L(z)\). Define the projected integral span
\[
 H=\sum_{z\in\mathcal D_{aj\sigma}}R\,P\theta(z).
\]
The module for the basis test is \(dH\). For a conductor set \(Q\), level \(h\), and Galois character \(\eta:G_E\to R_h^\times\), each \(z\) has a set \(\Theta_\eta(z)\) of twisted theta vectors in \(R_h\otimes_R V\). Write \(P_h=1\otimes P\). The twisted module and its ambient module are
\[
 H_\eta=\sum_{z\in\mathcal D_{aj\sigma}}\ \sum_{y\in\Theta_\eta(z)}R_h P_h y,
 \qquad H_\eta^{\rm amb}=P_h\bigl(\theta_{\rm amb}(\Lambda_{\rm amb})\bigr),
\]
where \(\Lambda_{\rm amb}\) is the lattice in the auxiliary theta space at \((a,j,\sigma,h,Q,\eta)\) and \(\theta_{\rm amb}\) is its given linear map to \(R_h\otimes_R V\).

For two full lattices \(M,N\) over a coefficient discrete valuation ring \(A\) in the same space, define the signed index, with the test lattice in the first argument, by
\[
 i_A(M,N)=
 \frac{\operatorname{length}_A(M/(M\cap N))-
       \operatorname{length}_A(N/(M\cap N))}
      {\operatorname{ord}_A(\ell)}.
\]
Here a full lattice is a finitely generated submodule whose fraction-field span is the given space. For a finite injective extension \(R_h\to A\), write \(H_{\eta,A}\) for the \(A\)-span of the image of \(H_\eta\) in \(A\otimes_R V\), and \(H_A\) for the \(A\)-span of \(1\otimes H\).

For a finite conductor set \(Q\), let \(D_a(Q)\) be the finite abelian \(\ell\)-group of the twisting datum, with its surjective Artin map \(\alpha_Q:G_E\to D_a(Q)\). For \(Q\) disjoint from the places above \(\ell\), every finite anticyclotomic abelian \(\ell\)-quotient unramified outside \(Q\) factors through this map. For each abelian theta label \(\lambda\) the datum specifies a sign \(\sigma_\lambda\) and coefficients \(b_{Q,\lambda,\delta}\in K\) for \(\delta\in D_a(Q)\). For a character \(\gamma:D_a(Q)\to R_h^\times\), put
\[
 A_{a,Q,\lambda}(\gamma)
 =\sum_{\delta\in D_a(Q)} b_{Q,\lambda,\delta}\,
                              \gamma(\delta)^{\sigma_\lambda},
 \qquad
 c_0(a,\lambda)
 =v_R\!\left(\sum_{\delta\in D_a(\varnothing)}
                         b_{\varnothing,\lambda,\delta}\right).
\]
The first sum uses the coefficient embedding \(K\hookrightarrow\operatorname{Frac}R_h\). On the constructed members the empty-conductor sum is nonzero by Lemma~\ref{B8.07k}, applied at the selected index. For \(P\in Q\), use the chosen place of \(\overline E\) over \(P\) and its inertia group \(I_P\). For \(\eta=\gamma\circ\alpha_Q\) ramified at \(P\), put
\[
 t_P(\eta)=\inf_{\substack{g\in I_P\\\eta(g)\ne1}}
                         v_{R_h}(\eta(g)-1).
\]
The set in the infimum is nonempty and finite because the image of \(\eta\) is finite.
\end{definition}

\begin{definition}[The seven family conditions]\label{B8.01k}
Fix \(F_m\) of parameter \(\varepsilon\). At a point write \(N_m=N(m)\), and use the scale parameters \(X',\delta_1\) of \ref{B2.54}. The notation \(\lceil y\rceil_+\) denotes the nonnegative ceiling. The duty of \(F_m\) consists of the following assertions.
\begin{enumerate}
\item The parameters \((C_p,e_P,c_{TO},C_t,C_{\eta F},C_c,c_w,C_{YL})\) equal \((1,12,377,401,6,2,9/10,9)\). At every \(x>F_m.x_1\) and every point, the host conductor ideal \(\mathfrak c_K\subseteq\mathcal O_K\) equals \(\mathcal O_K\). The vectors \((D.\mathrm{hc}(\tau)).v\) agree at all embeddings \(\tau\), and \(g_i\le b^9\) for \(i=0,1\). The designated primes \(p_{\rm de},p'_{\rm de}\) are both below \(X'\), and
\[
 j_0=\lceil\log_5(n_HN_m/m)\rceil_+,\quad
 L_w=10\lceil(C_t+4)\log x\rceil_++2\lceil30/\varepsilon\rceil_++22,\quad
 T_O=\lceil x^{c_{TO}}\rceil_+.
\]
\item At the same points, \(n_H,m,k\) belong respectively to \([a,2a]\) for \(a=x^{2+8\varepsilon},\sqrt{x\delta_1},x^{9/2+10\varepsilon}\).
\item At the same points the first modulus satisfies \(q_0<X'\).
\item Beyond one cutoff, for each point, \(g=g_i\), and each designated place above \(g\), the indexed list \(1,\chi_j,\chi_j^{-1}\) of tame labels has distinct entries, and each \(\chi_j\) is a square in \(\F_g^\times\). For each prime \(\ell\) outside
\[
 \{g\}\cup\{\ell:\exists\,1\le d\le r,\ g^d<\ell^{4v_\ell(g^d-1)}\},
 \qquad r=k+m-1,
\]
the prime-to-\(\ell\) parts of the list are distinct.
\item Beyond one cutoff every point satisfies the small-share bound of \ref{B5.22b} for its own weights, as in \ref{B8.02a}.
\item Eventually at all \(x>F_m.x_1\) and every point, each tuple is \(\mu_0,\mu_*,\mu_*'\), or \(\mu_w\) for a finite coefficient place \(w\).
\item Eventually at all \(x>F_m.x_1\) and every point, at each \(w\) with \(\mu_w\ne\mu_0\) the twist index, set \(Q\), coefficient level \(k'\) and signed character agree at \(w,\bar w\). The sign selects the character or its inverse and reverses when \(w\ne\bar w\). Use Definition~\ref{B8.01ka} for the point's own theta data. Choose nonnegative integers \(N_0(a)\) and finite sets \(S(a,j,\sigma)\) of theta inputs for all indices, layers and signs before choosing \(w\) or a character power. At each index \(a\) selected by a nonreference \(w\), require the following for every \(1\le j\le m\) and each sign. The vectors \(dP\theta(z)\) for all allowed \(z\) belong to \(V^\circ\) and span it over \(K\). The set \(S(a,j,\sigma)\) is contained in \(\mathcal D_{aj\sigma}\), and the vectors \(dP\theta(z)\) indexed by this finite set form an \(R\)-basis of \(dH\).

Fix such a \(w\). Write \(a,Q,k',\gamma_w\) for its index, conductor set, coefficient level and signed character of \(D_a(Q)\). For every \(n\in\Z\) with \((|n|,\ell_a)=1\), take \(\gamma=\gamma_w^n\), \(\eta=\gamma\circ\alpha_Q\) and \(e=\ell_a^{-2}\). Require \(k'\ge1\). Each \(P\in Q\) satisfies \(P\ne\bar P\) and has absolute residue degree and ramification index one. Require \(Q\) to be disjoint from \(\Sigma_a\cup\mathcal V\) and from \(\mathcal L(z)\) for every \(z\in S(a,j,\sigma)\), all \(1\le j\le m\), and both signs. The character \(\eta\) has open kernel, goes to its inverse under complex conjugation, takes values in the group generated by the designated primitive \(\ell_a^{k'}\)-th root, is ramified exactly at \(Q\), and is trivial on every decomposition group above \(\ell_a\).

For this power there are \(e_{j\sigma}\ge0\) with \(\sum_{j=1}^m\sum_{\sigma=\pm1}e_{j\sigma}\le e\). At \((j,\sigma)\) use the character \(\eta^\sigma\). The module \(H_{\eta^\sigma}\) lies in \(H^{\rm amb}_{\eta^\sigma}\) and is a full lattice in the scalar extension of \(V^\circ\). For every finite injective extension of discrete valuation rings \(R_{k'}\to A\) and every full \(A\)-lattice \(M\) in that scalar extension over \(\operatorname{Frac}A\), require
\[
 i_A(M,H_{\eta^\sigma,A})\le i_A(M,H_A)+e_{j\sigma}.
\] A signature \(s\) is good if its prescribed Hecke eigenvalues occur on a nonzero vector in the image of its type idempotent at its layer and sign, for the Hecke operators of orders \(1,\ldots,r\) at every finite place \(P\ne\bar P\) outside \(\Sigma_a\). Let \(\rho_s\) be its integral rank-\(r\) representation and \(\chi_s\) its prescribed character before the auxiliary twist. Write \(O\) for the coefficient discrete valuation ring at level \(k'\), \(K\) for its fraction field, and \(\eta_s=\eta^{\sigma_s}\), where \(\sigma_s\) is the sign attached to the signature \(s\). For any such character \(\eta\), define \(L_s(\eta,Q)\) to be the length over \(O\), divided by \(\operatorname{ord}_O(\ell_a)\), of the following subgroup of \(H^1(E,(K/O)^r)\). The action on the coefficients is \(\rho_s\chi_s^{-1}\eta^{-1}\), with the original coefficients mapped to \(O\). Above \(\ell_a\) take the \(O\)-span of the reductions of the finite local cocycles for \(\rho_s\chi_s^{-1}\) with coordinates in the original fraction field, using the prescribed local coefficient embedding. Away from \(\ell_a\), allow arbitrary restrictions at \(Q\), at the places above \(g_0,g_1\), and at the places above \(p_{\rm de}\) if \(\ell_a\ne p_{\rm de}\), or above \(p'_{\rm de}\) otherwise. Require unramified restrictions at all remaining finite places. Put \(\nu_a=\min_{g\in G_E}v_{\ell_a}(\varepsilon_{\ell_a}(g)^{tr}-1)\), where \(\varepsilon_{\ell_a}\) is the cyclotomic character, \(t\) is the auxiliary rank and \(v_{\ell_a}(\ell_a)=1\). The requirement is \(L_s(\eta_s,Q)\le\nu_a+e_s\) for every good \(s\), with \(e_s\ge0\) and \(\sum_{s\ {\rm good}}e_s\le e\). For every abelian theta label \(\lambda\), require \(A_{a,Q,\lambda}(\gamma)\ne0\) and
\[
 -N_0(a)\le v_{R_{k'}}(A_{a,Q,\lambda}(\gamma))\le c_0(a,\lambda)+e.
\]
For the same Artin character require \(t\sum_{P\in Q}t_P(\eta)\le e\). The witnesses \(N_0,S\) are common to all places and powers at this point.
\end{enumerate}
\end{definition}

\begin{definition}[Pointwise design bounds]\label{def:p8c1-design-bounds}
Use the ranks and host sets of \ref{B2.56}. Write \(n_0\) for the size parameter of the second set \(T_2\), \(L_w\) for the window length and \(n\) for the size parameter in the top-count estimate of \ref{B5.21}. The thickness and top-count bounds are
\[
 |T_2|\ge \frac{n_0}{L_w}
 \exp\!\left(-C_{T_2}\sqrt{\log n_0\log L_w}\right),
 \qquad \tfrac12\delta_1 n\le |TT|\le2\delta_1n .
\]
Let \(\mathcal P\) consist of the residue-class-zero primes \(p<X'\) for which the Frobenius coordinate \(h_p\) has order at least \(x^{\varepsilon/4}\). Let \(b_p\) count the positions \(t\in TT\) whose orbit fails the good-orbit test of \ref{B5.19}. Require
\[
 \sum_{p\in\mathcal P}\frac{\log p}{p}b_p
 \le 4C_bx^{-2}\delta_1n\sum_{p\in\mathcal P}\frac{\log p}{p}.
\]
For every primitive \(\chi\) inducing, at modulus \(8qd\), the prescribed \(q\)-character, the \(d\)-character or their product, require
\[
 |g_\chi(x)|\le x^{1/2+\varepsilon},\qquad
 |(L'/L)(1,\chi)|\le C_5\log x .
\]
For the number \(M\) of words require
\[
 3\Gamma^2x^{CH}\le M\le3\Gamma^2x^{CH}+1,\qquad
 \Gamma=\exp(C_\Gamma m^5 f\log b).
\]
\end{definition}

\begin{definition}[Good finite-type labels]\label{def:p8c1-good-labels}
At either type prime \(g\), let \(f_v\in\F_g[X]\) be the polynomial assigned at each designated place \(v\) above \(g\), and let \(r=k+m-1\). Each \(f_v\) must be monic and squarefree of degree \(r\) with nonzero constant term. For every prime \(\ell\) outside the exceptional set in clause~4 of \ref{B8.01k}, distinct roots of \(f_v\) must have distinct prime-to-\(\ell\) parts. For every prime \(\ell\ne g\), the prime-to-\(\ell\) part of each root must differ from those of every prescribed tame character value at \(v\) and its inverse. Here roots are taken in \(\overline{\F}_g\), and the prime-to-\(\ell\) part is the component of order prime to \(\ell\) in its cyclic group.

For each integer \(1\le a\le r/2\), there must be a place \(v\) at which \(f_v\) has no divisor of degree \(a\). At some place there must be an irreducible factor of prime degree \(d\) with \(r<2d\) and \(d+3\le r\). A tuple is good if these conditions hold at both type primes.
\end{definition}

\begin{definition}[The forty-six conditions on the common data]\label{B8.02}\label{B8.02g}
Let \(N\) be a predicate on family classes. For an arithmetic datum \(D\), objects \(O\) at that datum, scale constant \(c\), bundle \(\beta\) at \((D,O,c)\) and real numbers \(C,CH\), let \(DR(D,O,c,\beta,C,CH)\) be a predicate. A common choice for \(N,DR\) is defined by items~1--46 below. To specify the domain of the selection in item~2, write \(x_1\) and \(C_{\rm sz}\) for a family's cutoff and size constant. For a class \(\mathcal A\) the selection has values
\[
 \sigma(F_m,h,x,h_x,i)\in\mathrm{Bd}(D_{x,i},O_{x,i},C_{\rm sz}),
 \qquad h:F_m\in\mathcal A,\quad h_x:x_1<x,\quad i\in I_x. \tag{8.2}
\]
Here \(\mathrm{Bd}(D,O,c)\) denotes the set of bundles at \((D,O,c)\). A bundle \(\beta\) consists of the recursion realization, its spectral realization, the arithmetic family of row sides with its identification with the recursion, two nonnegative integers \(d_0,h_0\) for which the word spacing is \(\Delta=d_0\) and the size parameter is \(\eta=mh_0\), the square-integrable columns, and their indefinite realization at the scale determined by \(d_0,h_0,c\). The selection of recursions takes the recursion component of (8.2) and has cutoff \(x_1\). At every point of (8.2) its cutoff condition \(x_1\le x\) follows from \(x_1<x\).

For closure on a class \(W\), choose \(\varepsilon_1>0\) and then, for each \(0<\varepsilon\le\varepsilon_1\), choose a threshold datum \(\mathcal K\). Quantify over every family of parameter \(\varepsilon\) in \(W\) and every \(CH\in\R\) at which it meets these thresholds. In the notation of Definition~\ref{B8.39}, this means \(C_P\ge\mathcal K_{\rm pre}\), every entry of the middle list \(\mathbf m(CH)\) is at least \(\mathcal K_{\rm mid}(\mathbf o)\), and every entry of the last list is at least \(\mathcal K_{\rm last}(\mathbf o\mathbin{\Vert}\mathbf m(CH))\). Here \(\mathbf o\) is the outer list and \(\Vert\) denotes concatenation. Point cutoffs inside the conclusion may depend on \(F_m,CH\) and are uniform in \(i\). Reference surjectivity, the cylinder count, the real-place comparison and coefficient-lattice stability use all \(x\) beyond their point cutoffs. Rows that use the selected bundle impose \(x>F_m.x_1\).

For item~30, let \(h\) be a membership of the family in \(W_3(\mathcal A)\). Put \(c=C_{\rm sz}\) and take \(\beta\) from (8.2) using the induced membership in \(\mathcal A\). Write \(\mathcal B\) for the nine indefinite-side conditions of \ref{B8.38} on the realization of \(\beta\) at \(C,CH\). The paired-index and functional bounds there assume that one constant works in \eqref{eq:B8-divinN} eventually at every point and every lift layer. The lower weight bound assumes that at each layer a single leaf map has the prescribed values and phase of \ref{B8.55}. The assertion used in item~30 is
\begin{equation}\tag{8.4}\label{eq:p8c1-simultaneous}
 \operatorname{Closed}_{W_3(\mathcal A)}
 \left[(F_m,CH)\mapsto\forall h:W_3(\mathcal A)(F_m)\ \exists C\ge0\
 \text{ eventually along the selection, }\ \mathcal B\land DR(D,O,c,\beta,C,CH)\right].
\end{equation}


In the following list a selected point ranges over all arguments of (8.2). Each condition at that point uses the same selected bundle. Within a family, ``eventually'' means that one cutoff works for all later scales and every point index at those scales.

\begin{enumerate}
\item Choose a class \(\mathcal A\) of families, allowing every real family parameter \(\varepsilon\).
\item At each selected point choose one bundle as in (8.2), at \(c=C_{\rm sz}\). The same choice supplies both row sides and the columns for all the requirements below.
\item Choose a function \(c_T:\R^2\to\R\).
\item Let \(M(n)\) be the integer defined before Theorem~\ref{B8.01}. Suppose that for every \(n\ge0\) and every finite subgroup \(H\le\GL_n(\Q)\) one has \(|H|\mid M(n)\). Fix \(\lambda>1\). Let \(G\) be any condition on a scale, an arithmetic frame, the host data on that frame, a subset of its finite group, and an auxiliary rank. In applying \(G\) to a host, use its top set and its own auxiliary rank.

Suppose hosts satisfying \(G\) exist on windows of exponent \(\lambda\) in the sense of \eqref{eq:window-set-property} and the paragraph following it. In particular, for each \(0<\varepsilon<1/52\) the host constants and the set of scales are chosen before the individual hosts. Then there exists \(\varepsilon_0>0\) such that for every \(0<\varepsilon\le\varepsilon_0\) and every threshold datum \(\mathcal K\) of Definition~\ref{B8.39}, there are \(F_m\in\mathcal A\) of parameter \(\varepsilon\) and numbers \(CH\ge0\), \(C_5,C_b>0\) with the following properties. The family meets \(\mathcal K\) at \(CH\), has \(C_F=CH\), and satisfies \(C_{Tu}\ge c_T(\varepsilon,C_{YL})\). Its window contains \([Y,Y^\lambda]\) for arbitrarily large \(Y\). At every \(x>x_1\) and \(i\in I_x\) its host satisfies \(G\).

The family belongs to \(\mathcal A_{\rm ch}\). Its chosen signed twist satisfies the tests of clause~7 of Definition~\ref{B8.01k} at power one with its own theta data. Eventually at every point the second-set and top-set bounds, bad-orbit estimate with \(C_b\), three character estimates with \(C_5\), and both word-count bounds of Definition~\ref{def:p8c1-design-bounds} hold. In particular the word-count interval is \([3\Gamma^2x^{CH},3\Gamma^2x^{CH}+1]\).
\item Every family in \(\mathcal A\) satisfies all seven clauses of Definition~\ref{B8.01k}, including the tests for every power prime to the residue characteristic in clause~7.
\item At each selected point the extremal vectors exist and the single-map identities of Proposition~\ref{B8.55} hold in every layer.
\item At each selected point, every place of the twisting set on the reference row side corresponds to an arithmetic place of the size support, and its tuple differs from the reference tuple.
\item At each selected point require the primitive comparison on the selected face and separation of the transferred reference types. More explicitly, the split Hecke systems on the rational rows and lifts obey their class-function laws. The inert operators have matrices over the maximal CM coefficient subfield, and their adjoints for the class sum belong to the indexed inert operators. On the complex row and layer spaces the respective Hecke actions intertwine the maps from class functions. The finite group acts unitarily and the row and layer maps are equivariant. The selected subspaces lie in their specified row and layer blocks.

There are finite systems of complex linear tests with common kernel zero on the row space and on each layer space. On the images of the rational theta maps their values lie in the distinguished image of the coefficient field. On the reference layer the tests can also be chosen rational on the images of the coordinate basis. For each intertwining system of layer Hecke operators, each occurring good joint eigenspace on the class-function side is contained in the range of the adjoint lift, as in the decomposition of Lemma~\ref{B5.07g}. The two faces have positive reference dimension. The layer scalars \(\lambda_{j,t,p,v}\) and the common scalar \(\kappa\) are positive. On a type summand the layer lift times its adjoint acts by \(\lambda_{j,t,p,v_0}\), and the composite row map times its adjoint acts by \(\kappa\prod_{j=1}^m\lambda_{j,t,p,v_0}\), with the corresponding row-type indexing. Here \(v_0\) is the distinguished embedding.

On each row side the level space is the intersection of the principal space with the prescribed character space. At principal level the tuple maps are injective with the same range as the reference map. The reference layer space is the intersection of the principal layer space with the prescribed inverse-character space. The row type summands on both sides give orthogonal direct sums. The reference layer type summands give orthogonal direct sums there, with the same dimensions as the corresponding reference row summands. The reference coefficient modules are lattices. At each coefficient embedding the rational layer coordinates preserve linear independence and span the complex space, and the transported group action is the specified action multiplied by its coefficient character. At the other infinite places the inner product of transported rational lift vectors satisfies the conjugate Gram identity of Lemma~\ref{B5.07g}. Finally, for two distinct occurring reference types and every layer, their transfers differ at a split Hecke place or at an indexed inert operator.
\item At each selected point choose one assignment \(a\) from finite types to row indices. For every layer \(j\) and index \(c\), the spaces of types \(\sigma\) with \(a(\sigma)=c\) form an internal direct sum equal to the complex span of the transported rational space of index \(c\). For a nonzero such component the arithmetic and row-side \(\lambda\)'s agree at the reference embedding for the reference tuple and every tuple attached to a finite place. The sum of the initial row counts over \(a^{-1}(c)\) is the dimension over the coefficient field of the initial row space of index \(c\).
\item At each selected point the comparison identifies the two integral \(RR\) lattices with their respective indefinite lattices, including their metrics. At each used tuple the embedded rational lift columns span the indicated complex image. In the same coordinates the projected theta maps agree. The local ideals and their numerical corrections agree at every layer and place, with twist correction zero outside the twisting set. The column pairing after division agrees on every pair of rational inputs.
\item At each selected point, for every word, used tuple and pair of rational inputs, applying the coefficient embedding to the definite evaluation gives the arithmetic column period on the transported inputs.
\item Closed on \(W_1(\mathcal A)\), if \(C_{Tu}\ge c_T(\varepsilon,C_{YL})\), then eventually, for every \(i\) and every word \(u\),
\[
 b^{m/2}\le T_u,\qquad \log T_u\le C_{Tu}m\log b .
\]
\item Closed on \(W_1(\mathcal A)\), take any \(x>x_1\) and any \(i\). Suppose \(g_0,g_1\) split completely in the real base field and \(2fk<g_a-1\) for some \(a\in\{0,1\}\). For each word whose \(T_u\) splits completely in that field, the initial class set has cardinality \(M_k I_{01}(T_u^k-1)^f\), in the notation of Theorem~\ref{B8.54}.
\item There are positive constants \(C_1,c,C_7,C_{11}\) common to all arithmetic data for the zero-density bound and zero-free regions of Theorem~\ref{B8.45}, the explicit formula \eqref{eq:p8c5-explicit} and the unit-window zero count \eqref{eq:p8c5-window}, in the respective ranges of Theorems~\ref{B8.45} and~\ref{B8.46}.
\item Under the rank and gap conditions of Lemma~\ref{B8.02y}, each real row function admits the compactly supported reproducing integral and the weighted absolute integral of that lemma. The product of the normalized real matrix coefficients is integrable.
\item On the automorphic \(L^2\) quotient, use a finite-adelic Haar measure, a fundamental domain and the right regular action. Let \(e\) be locally constant and supported in the specified compact subgroup. At each infinite place take an irreducible unitary representation of formal degree \(d_v\) and a unit vector whose conjugate diagonal coefficient times \(d_v\) is integrable. If \(R_e\) is finite convolution by \(e\) and \(T_\infty\) is convolution by the product of these real coefficients, then \(\operatorname{im}(R_eT_\infty)\) is finite-dimensional.
\item Closed on \(W_3(\mathcal A)\), the size hypotheses of \eqref{eq:p8c1-level-row} imply that inequality for every word, with the six parameters of Theorem~\ref{B8.10}.
\item Choose a function of finite real lists with values \((C_{\rm res},n_0,A,C_b,C,c,c_b)\), where \(n_0\in\N\) and the other coordinates are real.
\item At the outer parameter list of every family, \(C\ge0\) and \(A\ge2(C_{YL}+1)\).
\item Under the hypotheses and in the order of thresholds of Proposition~\ref{B8.10y}, the sum of the products of the norms for any finite set of same-row pairs of the indicated two minor types satisfies the bound of that proposition.
\item For the nongeneric corank-one pairs the signed sum satisfies the first estimate of Lemma~\ref{B8.10d}, with that lemma's point hypotheses and thresholds.
\item For the generic scalar corank-one pairs of an equal-word operator the signed sum satisfies the second estimate of Lemma~\ref{B8.10d}, with its additional counting parameter chosen before its thresholds.
\item For every real-valued function of the outer parameter list, the two bounds \eqref{eq:p8c1-loop-rows} are closed on \(W_3(\mathcal A)\), using the selection in item~2 and the point hypotheses of Theorem~\ref{B8.10}.
\item Closed on \(W_3(\mathcal A)\), for every membership in that class, distinct words in the resulting selected spectral system satisfy \eqref{eq:p8c1-distinct-row} under its source hypotheses.
\item Closed on \(W_3(\mathcal A)\), the finite-level idempotent estimate of Theorem~\ref{B8.52} holds with \(n_0,A,C,c,c_b\) from item~18 evaluated at the family's outer parameter list. The levels, measures and signature restriction are those in that estimate.
\item At every selected point, spaces of equal rank and admissible constant signatures have the same fundamental-domain volume for their prescribed measures. At rank \(\dim Y_0\), each positive definite space has the standard genus mass. Its finite-adelic class set modulo rational isometries and the stabilizer of its base lattice is finite. For any complete set of representatives, the sum of the reciprocal stabilizer orders equals its genus mass.
\item At every selected point choose a word-to-column assignment agreeing with the division assignment, complex column coordinate isomorphisms, a linear map from complex reference coordinates into the arithmetic ambient space, and a semilinear map of row spaces. The reference coordinate map covers the good reference subspace. On rational inputs, reference evaluation in these coordinates gives the embedded definite evaluation. The row map has image the good selected row space and multiplies norms by a fixed positive scalar. It intertwines reference theta evaluation up to one nonzero scalar. For each word the column coordinate map, with the prescribed conjugation, multiplies the norm by a positive scalar that may depend on the word. In each finite spectral realization the transported row-column pairing equals the model pairing up to the prescribed conjugation and one nonzero scalar per word.
\item Closed on \(W_3(\mathcal A)\), eventually at every point the image of the reference space \(V_0(\mu_0)\) under \(\iota_{\mu_0}\) is the entire target.
\item For every \(F_m\in\mathcal A\), eventually at every point and every tuple \(\mu\), multiplication by \(\iota_D(a)\) preserves \(\widetilde M_\mu\) for all \(a\in\mathcal O_{K_D}\).
\item Require \eqref{eq:p8c1-simultaneous}. In particular the same \(C\ge0\) serves for the nine indefinite assertions and the definite-side predicate at the same selected bundle.
\item Closed on \(W_3(\mathcal A)\), eventually at \(x>x_1\), reference point evaluation for every word satisfies item~1 of Theorem~\ref{B8.09}, including \eqref{eq:p8c1-point-norm}.
\item With the same family quantifiers, each used nonreference tuple has integral rational point evaluation for every word, as in item~2 of Theorem~\ref{B8.09}.
\item With the same family quantifiers, the common denominator hypothesis of Theorem~\ref{B8.09} gives one representative of each reference period satisfying \eqref{eq:p8c1-period-valuation} at every finite place.
\item With the same family quantifiers and denominator hypothesis, each used nonreference period has one representative satisfying \eqref{eq:p8c1-period-valuation} at every place outside its excluded residue primes.
\item With the same family quantifiers, the reference scaling and the twisted column bound at \(T_u\) in item~6 of Theorem~\ref{B8.09} hold for every word.
\item Closed on \(W_3(\mathcal A)\), there is one real \(C_1\) for the two eventual estimates of Lemmas~\ref{B4.23} and~\ref{B4.23a} at a place above \(g_a\ne T_u\). The lower bound for the normalized valuation of a nonzero period representative and the upper bound for \(\sum_j(F_j-E_{u,j})\) use the same quantity \(-\sum_j vNN_{aujw}-C_1\operatorname{errC}(D,L)\) of those lemmas. The period assertion includes the zero representative.
\item Closed on \(W_3(\mathcal A)\), there are \(C_{\deg},C_g,C_T,C_\Gamma,C_d,C_I,x_0\in\R\), with \(C_g\ge0\), such that all eight conditions of Definition~\ref{B5.12k} hold at every point with \(x\ge x_0\) and \(x>x_1\).
\item The following two assertions are closed on \(W_3(\mathcal A)\). Eventually at every \(x>x_1\) and every point index there are local lifts with the comparisons and the stabilizer and complementary-orbit bounds of \eqref{eq:p8c5-twist}, with initial saving one. For every membership in this class and every point above \(x_1\) at or beyond the selection cutoff, the rational chain maps, initial stabilizers and selected kernel entries satisfy the second assertion of Proposition~\ref{B8.56}. Its orbit phases and signs must work for the entries of that selection.
\item Choose \(c_{\rm Ic}\in\R\).
\item Require \(c_{\rm Ic}>0\).
\item Closed on \(W_3(\mathcal A)\), eventually at every point there is an arithmetic row system whose selected indices map injectively onto the good indices. Its real parameters, split level and finite projectors agree with those for the point. Its cylinder counts satisfy Theorem~\ref{B8.52} with error parameter \(\exp(-c_{\rm Ic}rf\log r)\), \(r=k+m-1\).
\item Closed on \(W_3(\mathcal A)\), for each membership in that class there is a cutoff such that at every selected point beyond it the following holds for all \(0\le j<m\). At the reference tuple take the level subspace for the selected subgroup and its inverse character. Include its coordinate vectors into the full space of layer forms, project to the good part, and apply orthogonal projection there to the part of the selected layer subspace lying in the good part. In coordinates for that last subspace, every vector over the coefficient field is the image of a level vector under this composite.
\item Closed on \(W_3(\mathcal A)\), the real measures, formal degrees, quotient volumes, symmetric columns and level-zero characters satisfy item~5 of Theorem~\ref{B8.09}, with the real-place family quantifiers there.
\item Take any arithmetic datum with the following volume normalization. For every Hermitian space \(W\) of admissible constant signature there exists a Haar measure with the compact-dual normalization at infinity and finite lattice stabilizer of measure one. For each such measure the rational unitary group admits a fundamental domain, and every such domain has measure \(\operatorname{vol}_D(W)\). For this datum require the unit-hyperplane assertion of Theorem~\ref{B8.49}, including \eqref{eq:p8c5-volume-ratio}, in each of its stated ranks and signatures.
\item For every arithmetic datum and every everywhere positive definite Hermitian space of rank two with its self-dual base lattice, the genus mass is \(M_2\) of Proposition~\ref{B8.50}.
\item Require \(N(W_3(\mathcal A))\).
\end{enumerate}
\end{definition}

\subsection*{Index of archive field names for the numbered items of Definition~\ref{B8.02}}
\begin{longtable}{rlrlrl}
\toprule
Item & Archive field & Item & Archive field & Item & Archive field\\
\midrule
\endhead
1 & \(\mathcal A\) & 17 & \Lean{hTF} & 33 & \Lean{aPer0}\\
2 & \(\sigma\) & 18 & \Lean{c30} & 34 & \Lean{aPerTw}\\
3 & \Lean{cT} & 19 & \Lean{hc30} & 35 & \Lean{aP10}\\
4 & \Lean{hEx} & 20 & \Lean{sTy} & 36 & \Lean{aAtG}\\
5 & \Lean{hduty} & 21 & \Lean{sNG} & 37 & \Lean{aC2}\\
6 & \Lean{hRB} & 22 & \Lean{sM} & 38 & \Lean{hServ}\\
7 & \Lean{hTw} & 23 & \Lean{sLp} & 39 & \Lean{cIc}\\
8 & \Lean{hPA} & 24 & \Lean{sO} & 40 & \Lean{hcI}\\
9 & \Lean{hLamP} & 25 & \Lean{h30} & 41 & \Lean{irrC}\\
10 & \Lean{hDivG} & 26 & \Lean{hInstR} & 42 & \Lean{zL}\\
11 & \Lean{hDcPtG} & 27 & \Lean{selJ} & 43 & \Lean{tie}\\
12 & \Lean{hTu} & 28 & \Lean{jOn} & 44 & \Lean{hU}\\
13 & \Lean{hn4} & 29 & \Lean{hMt} & 45 & \Lean{hG}\\
14 & \Lean{hLit} & 30 & \Lean{rRows} & 46 & \Lean{nRest}\\
15 & \Lean{hWu} & 31 & \Lean{aPt0} & & \\
16 & \Lean{hRF} & 32 & \Lean{aPtTw} & & \\
\bottomrule
\end{longtable}

\begin{definition}[Size and counting conditions for the spectral estimate]\label{B8.02c}
Use the point model of \ref{B8.10}. Write \(r_0=(m-2)/2\), \(n=k+m-1\), \(L_b=\log b>1\), \(g=g_0g_1\), and \(v_E=\operatorname{rd}(E)/2\), where \(E/F\) is the CM extension of the model and \(\operatorname{rd}(E)\) is its root discriminant. Let \(\mathcal D\) be the spacing of the two choices of an interior word coordinate and \(C_\Gamma\) the size constant from the point. At a place \(\tau\) of the CM type, use the \(k-1\) interlacing parameters \(v_i\), the \(m\) positive weight parameters \(B_j\), the decreasing length-\(n\) Harish--Chandra parameter \(\lambda\), and the length-\(k\) weight \(u_\tau=(u_i)\) of a word \(u\). Indices here start at one. A top place is a place over a chosen top position. There put \(B_*=-u_{0,k}\) for the base word \(u_0\).

The source condition \(\mathrm{Cnt.Hyp0}\) at \((C_0,C',C_\ell,C_g,c_b,C_w,c_T,C_T)\) has \(C_0\ge1\) and the scalar inequalities
\[
\begin{gathered}
 8C_\Gamma\le C',\quad C'm^5fL_b\le k,\quad m^4\le f\le k,\quad \log f\le\sqrt{r_0},\\
 0<C_\ell\le m,\quad C_\ell L_b\le f\log k,\quad C_\ell m^2L_b\le k,\qquad
 v_E\ge1,\quad C_\ell m\log v_E\le k,\\
 C_\ell m^3L_b\log\operatorname{rd}(E)\le k(\log k)^2,\quad
 \log\operatorname{rd}(E)\le L_b/C_\ell,\\
 (kb)^8\le g_i\le b^{C_g}\ (i=0,1),\qquad (kfL_b)^{C_\ell}\le\mathcal D,\qquad
 b^{c_Tm}\le T_u\le b^{C_Tm}\quad\text{for every }u .
\end{gathered}
\]
At every place and all applicable indices and words require
\[
\begin{gathered}
 C_\ell mk\mathcal D\le B_j,\quad
 \max(k^{C_\ell},b^{c_b})\le B_j\le b^2,\qquad B_m\le2B_1,\quad u_1\le b^{C_w},\\
 B_{j+1}-B_j\ge C_\ell n^{19}f^5L_b,\qquad
 \lambda_i-\lambda_{i+1}>C_\ell km^4+nm^{10},\\
 0\le v_i\le4k\mathcal D,\qquad u_1-v_1\ge4k\mathcal D.
\end{gathered}
\]
For each word and \(1\le i\le k-2\), the choice \(u_{i+1}=v_i-1/2\) or \(u_{i+1}=v_i-1/2-\mathcal D\) is constant across places. Off the top, \(u_1-v_1\le8k\mathcal D\) and \(u_k=v_{k-1}-1/2\). On the top, \(B_1-B_*\ge3\) and \(B_1/2\le B_*\). The interlacing is \(u_i>\lambda_i>u_{i+1}\) for \(i<k\) and \(u_k>\lambda_k\). Distinct words have distinct weight vectors at every place.

The distinct host primes \(g_0,g_1\) and each level prime \(T_u\) split completely in \(E\), with \(T_u\ne g_i\) and \(T_u\equiv1\pmod{g^2}\). Require \(\operatorname{disc}E=(\operatorname{disc}F)^2\) in the model's discriminant convention. With \(d_u^S,D^\star\) the column dimension and main term from \ref{B5.11}, require \(C_0^{-1}\le d_u^S(\Gamma_u^{\mathrm{sp}})^2/D^\star\le C_0\), where \(\Gamma_u^{\mathrm{sp}}\) denotes the real factor \(\Gamma_u\) in the spectral count of \ref{B5.11}. The source hypothesis of (8.22) is this condition with \(C_0=3\) and the constants in \ref{B8.10}.

For a pair \((u,u')\), \(\mathrm{Cnt.Stand}\) means \(r_0\ge1\), that \(T_u,T_{u'}\) are completely split primes at least three, and that \(g_0,g_1\) are distinct completely split primes with \((T_uT_{u'},g)=1\) and \(T_u\equiv T_{u'}\pmod g\). The condition \(\mathrm{Cnt.Hyp}(u,u';A)\) comprises \(\mathrm{Cnt.Hyp0}\) and \(\mathrm{Cnt.Stand}\) together with
\[
 \max(T_u/T_{u'},T_{u'}/T_u)\le A,\quad
 C_\ell m^6L_b\le k\log k,\quad
 (\log f)^2\le A m^2L_b,\quad \log f\le A r_0^2,
 \qquad A\in\R.
\]
It requires also \(V(u'_\tau)\le A^kV(u_\tau)\) at each \(\tau\) in both orders, where \(V(w)=\prod_{i<j}(w_i-w_j)\).
\end{definition}

\begin{lemma}[Counting conditions from the source sizes]\label{B8.02d}
Under \ref{B8.02c}, \(C'\ge1\) and \(C_\ell>2\), \(\mathrm{Cnt.Hyp}(u,u;1)\) holds for each word \(u\). For a pair of words and \(A\ge1\) it suffices to establish the two level-ratio bounds and the two Vandermonde bounds of \ref{B8.02c}. In the distinct-word row of \ref{B8.10}, \ref{B5.11b} gives these with \(A=3e^{7/4}\) under its size and mass thresholds.
\end{lemma}

\begin{proof}[Outline]
Derive \(m\ge 4,r_0\ge 1\), the modulo-\(g\) congruence, \(T_u\ne 2\), and coprimality from the hypotheses, and chain the scalar bounds of \ref{B8.02c}. The two comparisons are equalities for a repeated word and follow in both orders from \ref{B5.11b} for distinct words. The complete proof is in \S\ref{prf:B8.02d}.
\end{proof}

Let \(E/F\) be a CM extension unramified at every finite place, let \(O=\mathcal O_E\), and put
\[
 f=[F:\Q],\qquad \rho=|\operatorname{disc}E|^{1/(2f)}=2v_E,\qquad
 \ell=\log_2(2f),\qquad c=1+\ell.
\]
For an integral Hermitian lattice \(P\), write \(P^\#=\{z:h(z,P)\subset O\}\), with the dual taken in \(EP\), and put \(d(P)=|P^\#/P|\). All lattice indices below are cardinalities of additive quotients. Let \(M_t\) be the mass of the genus of \(O^t\) with its standard Hermitian form. Mass weights an isometry class by the reciprocal of its automorphism order.

\begin{lemma}[Integral lattices of a given discriminant order]\label{B8.02ba}
A projective \(O\)-lattice of rank \(t\ge1\) has at most \(i^{t+\ell}\) sublattices of index \(i\ge1\). In a fixed totally positive definite Hermitian \(E\)-space \(W\) of rank \(t\), one has, for every integer \(n\ge1\),
\[
 \sum_{\substack{[P]\text{ in }W\\P\text{ integral},\ d(P)=n}}
       \frac1{|\operatorname{Aut}P|}
 \le M_t n^{(t+\ell)/2}.
\]
\end{lemma}

\begin{proof}[Outline]
Count index-\(i\) inclusions via quotient ideals and residual hyperplanes. Enlarge to a maximal lattice, bound its mass by the unramified models and the Gan–Hanke–Yu formula, and count inclusions with inverse stabilizer weights using \(n=bi^2\). The complete proof is in \S\ref{prf:B8.02ba}.
\end{proof}

\begin{lemma}[Primitive intersections and their gluings]\label{B8.02bb}
Let \(V\) be totally positive definite of rank \(k=j+s\), where \(j,s\ge1\). Let \(\mathcal G\) be an isometry-invariant collection of full integral lattices in \(V\) such that \(L^\#/L\) is cyclic of order \(D\) for every \(L\in\mathcal G\). For \(a\ge1\) an integer put
\[
 N_{j,a}(L)=\#\{U\subset V:\ U\text{ an }E\text{-subspace},\
                  \dim_E U=j,\ d(L\cap U)=a\}.
\]
These intersections are precisely the primitive rank-\(j\) sublattices with discriminant order \(a\). Then
\[
 \sum_{[L]\in\mathcal G/{\rm U}(V)}\frac{N_{j,a}(L)}{|\operatorname{Aut}L|}
 \le M_jM_s\,a^{j/2}(aD)^{(9s+8+4\ell)/2}.
\]
\end{lemma}

\begin{proof}[Outline]
Bound rational orbits by Landherr's classification and \ref{B8.02ba}. Glue the two primitive intersections, bound the gluing generators by cyclicity of the ambient discriminant, and use inverse stabilizer weights and the masses from \ref{B8.02ba}. The complete proof is in \S\ref{prf:B8.02bb}.
\end{proof}

\begin{lemma}[Comparison of the two standard masses]\label{B8.02bc}
Suppose \(v_E\ge1\), \(k\ge3\) and \(1\le s<k\). Put
\[
 M_*=(20f\rho^2)^{2f(6+c)}.
\]
Then
\[
 M_sM_{k-s}\le M_k M_*\,\rho^f\,16^{skf}.
\]
\end{lemma}

\begin{proof}[Outline]
Use the trace-dual minimum and Banaszczyk's inequality to count short Gram matrices in ranks one and two, and \ref{B8.02bb} to count intermediate lattices. Multiply adjacent mass-ratio bounds, using the Euler factors of \ref{B8.48} in higher ranks. The complete proof is in \S\ref{prf:B8.02bc}.
\end{proof}

\begin{lemma}[Primitive sublattices in the split-level genus]\label{B8.02bd}
Suppose \(v_E\ge1\), and let \(L_0\) be a full integral lattice in a totally positive definite Hermitian \(E\)-space of rank \(k\ge3\). Suppose \(T\) is a rational prime splitting completely in \(E\) and \(L_0^\#/L_0\simeq O/(T)\). Put \(D=T^{2f}\), and let \(m_0\) be its genus mass. Then
\[
 m_0=M_k\left(\frac{T^k-1}{T-1}\right)^f.
\]
For all integers \(1\le s<k\), \(a\ge1\), one has
\[
 \frac1{m_0}\sum_{[L]\in\operatorname{gen}(L_0)}
       \frac{N_{k-s,a}(L)}{|\operatorname{Aut}L|}
 \le
 M_*\rho^f16^{skf}(a/D)^{k/2}(aD)^{5s+5+2c}.
\]
Here and below a sum over a genus takes one representative of each isometry class.
\end{lemma}

\begin{proof}[Outline]
Choose one prime per conjugate pair to form self-dual enlargements. Count their residual hyperplane fibres with inverse stabilizer weights and get the mass from \ref{B8.02ba}, then normalize the estimate of \ref{B8.02bb} by this mass and apply \ref{B8.02bc}. The complete proof is in \S\ref{prf:B8.02bd}.
\end{proof}

\begin{definition}[Embedded data at level \(g\)]\label{B8.02be}
Let \(L\) be a full integral lattice in a definite space \(V\), and let \(g\ge2\), \(r\ge0\) be integers. An embedded datum of index \(x\) consists of a full integral lattice \(K\subset V\) with \(d(K)=d(L)\) and an isometry \(K/gK\to L/gL\) for the forms modulo \(g\). Write \(\Gamma\subset K\oplus L\) for the inverse image of its graph, and put \(C=\{z:(z,z)\in\Gamma\}\). Require \([L:C]=x\) and require \(L/C\) to be generated by at most \(r\) elements. A triple consists of data \((K,\Gamma),(K',\Gamma')\) of indices \(x,x'\) and an \(E\)-linear isometry \(\eta\) of \(V\) taking \(K'\) to \(K\) such that \((y,z)\in\Gamma'\) if and only if \((\eta y,z)\in\Gamma\).
\end{definition}

\begin{lemma}[Fibre with prescribed fixed space]\label{B8.02bf}
Use \ref{B8.02be} with \(L^\#/L\) cyclic of order \(D\). Write \(k=\dim_E V\), \(X=xx'\), \(Y=XD^2\), and \(e=3(2r+1)\). For a subspace \(U\) of codimension \(s\), put \(a=d(L\cap U)\) and \(A_s=(2/g)^{2sf}DX\), with \(d(0)=1\). The number of triples for which \(\ker(1-\eta)=U\) is zero if \(a>A_s\). Otherwise it is at most
\[
 (A_s/a)^{k/2}X^{2s+2c+1}Y^e(4sf)^{4sf},
\]
where the last factor is one for \(s=0\). When \(s=k\) one also has the bound
\[
 X^{k/2}Y^{c+e}
 \left((2/g)^{kf}(4kf)^{4f}DX\right)^k .
\]
\end{lemma}

\begin{proof}[Outline]
Use projection as in \ref{B8.02bb} to bound the fixed lattice discriminant. Count lattices via \ref{B8.02ba}, lifts and graphs via \ref{B8.02bb}, and isometries via \ref{B8.02}, using \ref{B8.02ba} to count both lattices directly when the fixed space is zero. The complete proof is in \S\ref{prf:B8.02bf}.
\end{proof}

\begin{lemma}[Summation over the two length shells and the two Gram shells]\label{B8.10k}
Let \(r,j,h\) be positive integers, let \(H,E\ge1\), and let \(D>0\). On a measure space let \(S,S'\ge1\) be the two length maxima and \(Y,Y'>0\) the two Gram volumes. Suppose that, almost everywhere,
\[
 YY'D\le 2^j,\qquad D\le(SS')^r,\qquad
 \max(Y,Y')\le\max(S,S')^rE,\qquad
 1\le s_*\le\max\{1,(SS')^r/(2H)^{2r}\}.
\]
For \(i\ge1\) the \(i\)-th length shell is \((2^iH,2^{i+1}H]\), and the zeroth shell is \([1,2H]\). On the pair of shells \(i,i'\) put \(R=2^{i+i'+1}H\). Let \(A,B\) be nonnegative measurable functions. Assume that for every \(y>0\) the integral of \(A^2\) on these shells with \(y<Y\le2y\) is at most \(K R^e y^{2h}2^{-2qi}\). Assume the corresponding bound \(K'R^e y^{2h}2^{-2qi'}\) for \(B^2\) with \(y<Y'\le2y\). If \(e\ge0\) and \(q\ge e+4r+1\), then
\[
 \int AB\le
 16\sqrt{KK'}(2H)^{e+4r}E^2
 (2^j/D)^h s_*^{-(q-e-4r-1)/r}.
\]
\end{lemma}

\begin{proof}[Outline]
Fix the length shells and apply Cauchy–Schwarz to the squared-integral bounds on dyadic Gram-volume shells. Use the Gram-product constraint for all but one power, bound that power by the separate Gram shells, and sum four geometric series using the length decay. The complete proof is in \S\ref{prf:B8.10k}.
\end{proof}

\begin{lemma}[Normalization and propagation of length moments]\label{B8.02bg}
Let \(G\) be a local Cayley unitary group, let \(S\) be its Bessel subgroup, and let \(\nu\) be the prescribed measure on \(S\backslash G\). Let \((\sigma,\mathcal H)\) be the discrete-series representation and let \(\mathcal H^\infty\) denote its smooth vectors. Let \(\eta\) be the unitary Bessel representation of \(S\) on a finite-dimensional Hilbert space \(V\), and let \(B,L:\mathcal H^\infty\to V\) be continuous \(\eta\)-equivariant complex-linear maps. Choose a nonzero smooth vector \(v\) in the lowest compact type and put \(w=v/\|v\|\) and \(\Phi(g)=B(\sigma(g)v)\). Suppose there is \(c\in\C\) such that
\[
 B=cL\quad\hbox{on }\mathcal H^\infty,\qquad d>0,\qquad
 d\int_{S\backslash G}\|L(\sigma(g)w)\|^2\,d\nu(g)=1,
 \qquad 0<I_B:=\int_{S\backslash G}\|\Phi(g)\|^2\,d\nu(g)<\infty .
\]
Then
\[
 c\ne0,\qquad I_B=\frac{\|v\|^2|c|^2}{d},\qquad
 d\rho(g):=\frac{\|\Phi(g)\|^2}{I_B}\,d\nu(g)
 =d\|L(\sigma(g)w)\|^2\,d\nu(g),
\]
and \(\rho\) is a probability measure.

Let \(r,q\) be nonnegative integers. Let \(l_0,\ldots,l_r\) be positive measurable functions on this quotient, with \(l_r\ge1\) almost everywhere for \(\rho\), and put \(R_i=l_{i+1}/l_i\) for \(0\le i<r\). Suppose \(A,R\ge1\) and
\[
 \int l_0^{2q'}\,d\rho\le A^{q'},\qquad
 \int R_i^{2q'}\,d\rho\le R^{2q'}
 \quad
 (0\le q'\le(r+1)q,\ 0\le i<r).
\]
For every \(0\le a\le r\),
\[
 \int l_a^{2q}\,d\rho\le(A R^{2a})^q,\qquad
 \int l_a^{-2q}\,d\rho\le R^{2rq}.
\]
For \(b\ge1\), the second bound is \((b^{2r})^{2q}\) when \(R=b^2\).
\end{lemma}
\begin{proof}[Outline]
Linearity gives \(\Phi(g)=\|v\|cL(\sigma(g)w)\); squaring and integrating yield \(I_B\), \(c\ne0\), and the probability measure \(\rho\). Factoring \(l_a\) and \(l_a^{-1}\) through the ratios \(R_i\) and applying Hölder's inequality gives both moment bounds. The complete proof is in \S\ref{prf:B8.02bg}.
\end{proof}

\begin{lemma}[One ideal clearing the leading minors]\label{B8.02bh}
Let \(E\) be a CM field of degree \(2f\), let \(\mathfrak d\) be a nonzero conjugation-stable integral ideal, and let \(T\ge1\) be an integer. Fix \(0\le j\le r\), elements \(p_0=1,p_1,\ldots,p_j\in E\), and \(q_l=\sum_{a=0}^{l-1}(r-a)\). Suppose that
\[
 \mathfrak a_l=(p_l)\mathfrak d^{-2q_l}
\]
is an integral ideal whenever \(p_l\ne0\). Put \(S=\{l:1\le l\le j,\ p_l\ne0\}\), \(D_b=\prod_{l\in S}N(\mathfrak a_l)\), and \(N_T=T^f\). Then
\[
 \mathfrak c=(T)\prod_{l\in S}\mathfrak a_l\overline{\mathfrak a_l}
\]
is a nonzero conjugation-stable integral ideal contained in \((T)\), and
\[
 D_b\ge1,\qquad N(\mathfrak c)=(D_bN_T)^2,\qquad
 \frac{\mathfrak c\mathfrak d^{2q_l}}{p_l}\subset\mathcal O_E,
 \quad
 \frac{\mathfrak c\mathfrak d^{2q_l}}{\bar p_l}\subset\mathcal O_E
 \quad(0\le l\le j,\ p_l\ne0).
\]
If \(p_j\ne0\), then \(N(\mathfrak a_j)\ge1\).
\end{lemma}
\begin{proof}[Outline]
Containment and conjugation symmetry of \(\mathfrak a_l,\overline{\mathfrak a_l}\) give \(\bar{\mathfrak c}=\mathfrak c\subset(T)\). The identity \(\mathfrak a_l\mathfrak d^{2q_l}=(p_l)\) and norm multiplicativity yield the divisibilities and norm bounds. The complete proof is in \S\ref{prf:B8.02bh}.
\end{proof}

\begin{proposition}[The same-word estimate by leading minors]\label{B8.02b}\label{B8.10y}
Fix \(C_g,c_b,C_w,c_T,C_T,C_0\in\R\), with \(c_b,c_T>0\). There are \(C_H,C_h,C_p\ge0\) such that, for each \(C_s\ge C_p+4\), the following holds for sufficiently large \(C\), then sufficiently large \(C',C_\ell\). Use the tied models of \ref{B8.10}, with \(\mathrm{Cnt.Hyp0}\) and the unramified CM condition of \ref{B8.48}. Put \(r=r_S=(m-2)/2\), and assume
\[
\begin{gathered}
 C_\ell m^6\log b\le k\log k,\quad (kb)^8\le g_i\le b^{C_g},
 \quad C_\ell\log\operatorname{rd}(E)\le\log b,\quad \log f\le r^2,\\
 |\lambda_{\tau,a}-\lambda_{\tau,c}|>C_\ell km^3\ (a\ne c),\qquad
 T_w\le b^{C_Tm},\qquad (T_w,g_0g_1)=1 .
\end{gathered}
\]
For a relevant nonidentity term of word \(u\) with \(G(t)=0\), its type \((j,\delta)\) records the rank \(j\) of the reduced pairing and its leading pairing minors, with \(\delta_0=1\) and \(\delta_l=0\) for \(l>j\). Put \(q_l=\sum_{a=1}^l(r+1-a)\), \(\mathfrak d=\mathfrak D_{F/\Q}\mathcal O_E\), and \(H_1=b^{C_Hm}\). For \(\delta_l\ne0\), let \(d_{\tau l}=|\tau(\delta_l)|\operatorname{rd}(E)^{-2q_l}\) and \(d_l=N((\delta_l)\mathfrak d^{-2q_l})\), with \(N\) the absolute ideal norm. Set
\[
 D_\delta=d_j,\quad
 S_\delta=\prod_\tau\max\bigl(1,\max_{\substack{\delta_l\ne0\\1\le l\le j}}d_{\tau l}/(2H_1)^{2r}\bigr),
 \qquad H(c)=b^{cm^2f}.
\]
The product is over the CM type. Define
\[
 A(j,\delta)=(C/k)^{j(k-1-j)f}D_\delta^{-(k-1-j)/2}H(C)^{Cr^4}S_\delta^{-C_sk},
 \qquad M(j,\delta)=D_\delta^{k/2}H(C_h)^{C_hr^4}S_\delta^{C_pk}.
\]
For every finite set \(U\) of ordered same-row pairs of terms of types \((j,\delta),(j',\delta')\) in the diagonal kernel, with row vectors \(v_t,v_{t'}\),
\[
 \sum_{(t,t')\in U}\|v_t\|\|v_{t'}\|
 \le\dim(\pi_u)\operatorname{mass}_u(AM)(j,\delta)(AM)(j',\delta').
\]
Here \(\pi_u\) is the archimedean coefficient representation and \(\operatorname{mass}_u\) is the sum of inverse row-lattice stabilizer orders.
\end{proposition}

\begin{proof}[Outline]
Reduces frame pairs to normal form (\(j\ge1\)) and splits into amplitude and count bounds. Length moments (\ref{B8.02bg}) and Gram shells (\ref{B8.10k}) give the amplitude; \ref{B8.02bh}, \ref{B8.02bd}, and \ref{B8.02bf} bound fibres and genus triples. The complete proof is in \S\ref{prf:B8.10y}.
\end{proof}

\begin{lemma}[Normalization of a nonzero Bessel average]\label{B8.02yaa}
Let \(G\subset\GL_n(\C)\) be a local Cayley unitary group and let \(S=U\ltimes N\) be its Bessel subgroup. Fix Haar measures \(dg,da,dn\) such that multiplication carries \(da\,dn\) to a two-sided Haar measure on \(S\). Fix a measure \(\nu\) on \(S\backslash G\) and \(c_S>0\) such that
\[
 \int_{S\backslash G}\int_{U\times N}f(anx)\,da\,dn\,d\nu(x)
 =c_S\int_G f(g)\,dg
\]
for every nonnegative measurable \(f\).
Let \((\sigma,\mathcal H)\) be a continuous unitary representation of \(G\), let \((\pi,V)\) be a continuous finite-dimensional unitary representation of \(U\), and let \(\psi\) be the unitary Bessel character of \(N\). Fix a unit vector \(w\in\mathcal H\) and \(x_0\in G\), and define
\[
 m(y)=\langle\sigma(x_0)w,\sigma(y)w\rangle,\qquad
 a(x)=\int_{U\times N}|m(anx)|\,da\,dn,
\]
\[
 B(x)=\int_{U\times N}
          m(anx)\overline{\psi(n)}\pi(a)^{-1}\,da\,dn.
\]
Suppose \(a\) is \(S\)-invariant and
\[
 a(x)\le C\|x\|_{\rm op}^{M},\qquad
 \int_G |m(y)|\|y\|_{\rm op}^{M}\,dy<\infty
\]
for some \(C>0\), \(M\ge0\). Suppose \(B\) is continuous,
\[
 B(anx)=\psi(n)\pi(a)B(x),
\]
and \(B(x_1)\ne0\) for some \(x_1\in G\). There is \(\xi\in V\) such that \(A(x)=B(x)\xi\) satisfies
\[
 0<\mathcal N:=\int_{S\backslash G}\|A(x)\|^2\,d\nu(x)
 \le c_SC\|\xi\|^2\int_G|m(y)|\|y\|_{\rm op}^{M}\,dy<\infty .
\]
The function \(W^0=\mathcal N^{-1/2}A\) has quotient square norm one and the same equivariance. Every absolutely convergent identity
\[
 \int_G e(g)A(xg)\,dg=A(x)
\]
for a scalar function \(e\) also holds with \(A\) replaced by \(W^0\).
\end{lemma}
\begin{proof}[Outline]
Pick \(\xi\) with \(A(x_1)\ne0\); unitarity, Tonelli's theorem, and quotient integration give \(\mathcal N<\infty\). Disintegrating a neighborhood of \(x_1\) where \(\|A\|>\epsilon\) proves \(\mathcal N>0\), and scaling by \(\mathcal N^{-1/2}\) gives \(W^0\). The complete proof is in \S\ref{prf:B8.02yaa}.
\end{proof}

\begin{lemma}[A nonzero reproducing coefficient and the period normalization]\label{B8.02yab}
Let \(G\) be a local Cayley unitary group and let \(S=U\ltimes N\) be its Bessel subgroup. Fix Haar measures \(dg,da,dn\) such that multiplication carries \(da\,dn\) to a two-sided Haar measure on \(S\). Let \(\nu\) be the quotient measure determined by
\[
 \int_G f(g)\,dg
 =\int_{S\backslash G}\int_{U\times N}f(anx)\,da\,dn\,d\nu(x)
 \qquad(f\ge0\text{ measurable}).
\]
Let \((\sigma,\mathcal H)\) be a unitary discrete-series representation with formal degree \(d>0\), let \((\pi,V)\) be a finite-dimensional unitary representation of \(U\), and let \(\psi\) be the unitary Bessel character of \(N\). Fix a smooth unit vector \(w\in\mathcal H\), and put
\[
 \mathcal D=\operatorname{span}_{\C}\{\sigma(x)w:x\in G\}.
\]
For \(u,v\in\mathcal D\), define the operator-valued period
\[
 \mathcal P(u,v)
 =\int_{U\times N}
     \langle v,\sigma(an)u\rangle\overline{\psi(n)}\pi(a)^{-1}\,da\,dn.
\]
Suppose these integrals converge absolutely. Let \(\ell:\mathcal H^\infty\to V\) be a continuous complex-linear Bessel map, meaning that
\(\ell(\sigma(an)u)=\psi(n)\pi(a)\ell(u)\).
Suppose that
\[
 \mathcal P(u,v)\eta=\langle\ell(v),\eta\rangle\ell(u)
 \quad(u,v\in\mathcal D,\ \eta\in V)
\]
and that, for every \(v\in\mathcal D\),
\[
 \int_{S\backslash G}
       \|\ell(v)\|^2\|\ell(\sigma(x)w)\|^2\,d\nu(x)
 =\frac{\|\ell(v)\|^2}{d}.
\]
Let \(W:G\to V\) be continuous and satisfy
\[
 W(anx)=\psi(n)\pi(a)W(x),\qquad
 W(x)=d\int_G\overline{\langle w,\sigma(g)w\rangle}W(xg)\,dg,
\]
with an absolutely convergent integral for every \(x\).
Then \(W(x)=\kappa\ell(\sigma(x)w)\) for some \(\kappa\in\C\).
If \(W\) is nonzero, there is \(x_0\in G\) such that
\(\xi=\ell(\sigma(x_0)w)\ne0\). For this choice,
\[
 A(x):=\mathcal P(\sigma(x)w,\sigma(x_0)w)\xi
       =\|\xi\|^2\ell(\sigma(x)w),
 \qquad A(x_0)\ne0,
\]
\[
 \int_{S\backslash G}\|\ell(\sigma(x)w)\|^2\,d\nu(x)=d^{-1},
 \qquad
 \|A\|_{L^2(\nu)}^2=\frac{\|\xi\|^4}{d}.
\]
Consequently
\[
 W^0(x):=\frac{A(x)}{\|A\|_{L^2(\nu)}}
 =\sqrt d\,\ell(\sigma(x)w)
\]
has quotient square norm one, and \(W=(\kappa/\sqrt d)W^0\).
\end{lemma}
\begin{proof}[Outline]
Left translation, Fubini's theorem, quotient integration, and the rank-one period identity yield \(W=\kappa\Phi\). When \(W\ne0\), taking \(\xi=\Phi(x_0)\ne0\) in the period and norm identities gives the formulas for \(A\), \(\int\|\Phi\|^2\,d\nu\), and \(W^0\). The complete proof is in \S\ref{prf:B8.02yab}.
\end{proof}

\begin{proposition}[Real Bessel integrals and local proportionality]\label{B8.02ya}
Put \(n=k+m-1=2r_S+k+1\), and suppose \(n\ge1801\) and \(b_{\tau,j+1}-b_{\tau,j}\ge16n^{19}\) at every CM embedding \(\tau\) and every \(0\le j<m-1\). In a spectral realization \(X\) as in Theorem~\ref{B8.10}, fix a word \(u\) and an infinite place \(v\) of the totally real field \(F\) with corresponding CM embedding \(\tau_v\). Let \(G_v\simeq\U(p_v,q_v)\) (\(p_v+q_v=n\), with negative signature index \(q_v\)) be the local unitary group with Haar measure \(dg_v\), and let \((\sigma_v,\mathcal H_v)\) be the local discrete series representation of \(G_v\) in \(X\) with Harish--Chandra parameter \(\lambda_v\), formal degree \(d_v>0\), and unit lowest-\(K_v\)-type vector \(w_v\in\mathcal H_v\). Pass to the local Cayley signature coordinates on \(G_v\), and denote the Cayley group, Haar measure and Bessel subgroup of word \(u\) by \(G_v^0\), \(dg_v^0\) and \(S_{u,v}^0=U_{u,v}^0\ltimes N_{u,v}^0\le G_v^0\). Here \(N_{u,v}^0\) is the unipotent radical stabilizing the local isotropic flag of length \(r_S\) and orthogonal negative-norm vector \(z_{u,v}\), equipped with its generic unitary Bessel character \(\psi_{u,v}^0:N_{u,v}^0\to\C^\times\), and \(U_{u,v}^0\simeq\U(k)\) is the compact definite stabilizer acting unitarily on the finite-dimensional Hilbert space \(V_{u,v}\) by the irreducible representation \(\pi_{u,v}\) with compact Harish--Chandra parameter \(u_{\tau_v}=(a_1,\ldots,a_k)\). The local numerical data of \(X\) have even \(k\ge2\). Both \(u_{\tau_v}\) and \(\lambda_v=(\lambda_{v,1},\ldots,\lambda_{v,n})\) are strictly decreasing, with
\[
 a_i-\frac{k-1}{2}\in\Z,\qquad
 \lambda_{v,s}-\frac{n-1}{2}\in\Z,\qquad
 a_i-\lambda_{v,s}\in\Z+\frac12
 \quad(1\le i\le k,\ 1\le s\le n).
\]
Their interlacing is
\[
 a_i>\lambda_{v,i}>a_{i+1}\quad(1\le i<k),\qquad
 a_k>\lambda_{v,k},
 \qquad \lambda_{v,k+j}=-b_{\tau_v,j}\quad(0\le j<m).
\]
The simple-root character parameters of \(\psi_{u,v}^0\) are nonzero. Fix Haar measures \(da\) on \(U_{u,v}^0\) and \(dn\) on \(N_{u,v}^0\), let \(\nu_{u,v}^0\) be a quotient measure on \(S_{u,v}^0\backslash G_v^0\) disintegrating \(dg_v^0\) with respect to \(da\,dn\), and set \(N_v=4n^2\sqrt{2q_v}\).
\begin{enumerate}
\item For every \(x_0\in G_v^0\), the nonnegative integral
\[
 a_v(y;x_0)=\int_{U_{u,v}^0\times N_{u,v}^0}
 |\langle\sigma_v(x_0)w_v,\sigma_v(any)w_v\rangle|\,da\,dn
\]
is finite for all \(y\in G_v^0\) and satisfies \(a_v(y;x_0)\le C_v(x_0)\|y\|_{\rm op}^{N_v}\) for a constant \(C_v(x_0)>0\) independent of \(y\). For all \(x,x_0\in G_v^0\), the operator-valued Bessel integral
\begin{equation}\label{eq:p8c1-bessel-average}
 B_v(x,x_0)=\int_{U_{u,v}^0\times N_{u,v}^0}
 \langle\sigma_v(x_0)w_v,\sigma_v(anx)w_v\rangle\,\overline{\psi_{u,v}^0(n)}\,\pi_{u,v}(a)^{-1}\,da\,dn
\end{equation}
converges absolutely in \(\operatorname{End}_\C(V_{u,v})\), depends continuously on \(x\in G_v^0\), and satisfies \(B_v(a_0n_0x,x_0)=\psi_{u,v}^0(n_0)\pi_{u,v}(a_0)B_v(x,x_0)\) for every \((a_0,n_0)\in U_{u,v}^0\times N_{u,v}^0\).
\item There exist \(x_1,x_0\in G_v^0\) and \(\xi\in V_{u,v}\) such that \(B_v(x_1,x_0)\xi\ne0\).
\item With \(x_0,\xi\) as in~\textup{(2)}, put \(A_v(x)=B_v(x,x_0)\xi\). Then \(0<\int_{S_{u,v}^0\backslash G_v^0}\|A_v(x)\|^2\,d\nu_{u,v}^0(x)<\infty\), and if the global row function \(W_u\) of word \(u\) in \(X\) does not vanish identically on \(G_\infty\), the local factor \(W_{u,v}:G_v\to V_{u,v}\) of \(W_u\) transported to \(G_v^0\) is a scalar multiple of \(W_v^0=A_v/\|A_v\|_{L^2(\nu_{u,v}^0)}\).
\end{enumerate}
\end{proposition}

\begin{proof}[Outline]
Uses Cartan integration and dominated convergence for convergence and equivariance, then Bessel nonvanishing and \ref{B8.02yaa} for nonzero \(L^2\) norm. Schur orthogonality and Bessel multiplicity one identify \(W_v\) as a multiple of \(W_v^0\). The complete proof is in \S\ref{prf:B8.02ya}.
\end{proof}

\begin{lemma}[Reproduction and integrability of the real row function]\label{B8.02y}
Put \(n=k+m-1=2r_S+k+1\). Suppose \(n\ge1801\) and \(b_{\tau,j+1}-b_{\tau,j}\ge16n^{19}\) at every embedding and \(0\le j<m-1\). In a spectral realization \(X\) as in \ref{B8.10}, write \(G_\infty=\prod_vG_v\), let \(S_u\) be the real Bessel stabilizer for word \(u\), and let \(\nu_u\) be its quotient measure. With negative signature indices \(q_v\), put
\[
 N=4n^2\sqrt{2\max_v q_v},\qquad
 L_u(x)=\inf_{s\in S_u}\prod_v\|(sx)_v\|_{\rm op}.
\]
For each real row function \(W_u\) there is a bounded measurable \(a_u\) supported on a compact set containing the identity such that
\[
 W_u(x)=\int a_u(z)W_u(xz)\,dz,\qquad
 \int\|W_u(x)\|L_u(x)^N\,d\nu_u(x)<\infty.
\]
The convolution is absolutely convergent at each \(x\). The product \(\phi\) of the conjugate lowest-type coefficients times their formal degrees belongs to \(L^1(G_\infty)\).
\end{lemma}

\begin{proof}[Outline]
Cartan integration and matrix-coefficient decay give weighted \(L^1\) integrability of the lowest-type coefficient and \(\phi\in L^1(G_\infty)\). Unfold Proposition~\ref{B8.02ya} for the weighted \(W_u\) bound, and average Schur orthogonality near \(1\) to build \(a_u\). The complete proof is in \S\ref{prf:B8.02y}.
\end{proof}

\begin{proposition}[The fifth construction duty]\label{B8.02a}
For the class \(\mathcal A\) of \ref{B8.05}, every \(F_m\in\mathcal A\), \(x>F_m.x_1\), and \(i\in I_x\) has \(R_W(i)\le x^{-9}\) for the model's own weights, where \(R_W\) is the prime sum of \ref{B5.22b}. Thus \(x_W=F_m.x_1+1\) works in that assertion.
\end{proposition}

\begin{proof}[Outline]
Choose integer weights with prime sum at most \(x^{-9}\) by \ref{B5.22a}. Use \ref{B8.05} to prescribe them at all translations and impose a common cutoff, and \ref{B8.02ac} to identify the boxes, extended weights, and prime sum. The complete proof is in \S\ref{prf:B8.02a}.
\end{proof}

\begin{lemma}[A trial pair on a host and at a translated point]\label{B8.02ac}
Fix a scale \(s=(x,\varepsilon,J)\), ranks \(R\), a host \(H\), and the orbit constants \(K\) of \ref{B3.51}. Write \(A_H\) for the enumerated frame group. For exponents \(c_D,c_h\) put \(\Delta=\lceil x^{c_D}\rceil\), \(h_0=2\lceil x^{c_h}\rceil\), \(\mathrm{Sc}=4k\Delta\), and \(\eta=mh_0\). For integer functions \(\alpha,\beta\) on \(A_H\), define \(R_H(\alpha,\beta)\) by \eqref{5.19.1} at \(\theta=x^{-\varepsilon}\), with the host's prime classes, Frobenius coordinates and \(TT=H.TT\). Use the least frame label in each candidate orbit and the two shares of \ref{B3.51}, with \(c=4\mathrm{Sc}\) in \ref{B3.52a}.

Let \(z\) be a point constructed at a translation of \(H\) with prescribed top weights. Its scale, ranks, auxiliary primes \(q,d\), top labels, prime classes and Frobenius labels agree with those of \(H\). Suppose its size data give \(\mathrm{Sc},\eta\) as above and have orbit factor \(\mathcal L=K_\Lambda N_{\rm win}^2\log a_*\) as in \ref{B3.51}. Use the same representative convention at \(z\) and the local projection fraction of \ref{B3.52a}. Under the frame identification the integer boxes agree, and for every trial pair
\[
 R_z(\alpha,\beta)=R_H(\alpha,\beta).
\]
If the prescribed top weights are given by \(\alpha,\beta\) and the two functions vanish off \(TT\), the point's extended weights are exactly \(\alpha,\beta\).
\end{lemma}

\begin{proof}[Outline]
Common arithmetic data and matching frames give the same candidate-prime tests by \ref{B3.51}, small-prime shares directly, and large-prime shares by \ref{B3.52a}. The box parameters and top labels identify the endpoints and zero-extended weights. The complete proof is in \S\ref{prf:B8.02ac}.
\end{proof}

\section{From member parameters to the common data}

\begin{lemma}[Finite dimension of the projected row space]\label{B8.03a}
Let \(G=U(V)\) be the unitary group of a nondegenerate Hermitian space over the CM extension \(E/F\) of \ref{B8.02}. Give \(G(\mathbb A_F)\) product Haar measure. Write \(G_\infty\) and \(G_f\) for its infinite and finite factors, and put \(Y=G(F)\backslash G(\mathbb A_F)\). For an integrable kernel \(q\), write \(R_qh(x)=\int q(z)h(xz)\,dz\), using the factor on which \(q\) is defined.

At each real place \(v\), let \(\sigma_v\) be an irreducible unitary representation with formal degree \(d_v\) for the chosen measure. Take a unit vector \(w_v\) for which
\(\phi_v(g)=d_v\overline{\langle w_v,\sigma_v(g)w_v\rangle}\) is integrable. Put \(\phi=\prod_v\phi_v\). For every locally constant kernel \(e\) on \(G_f\) supported in the compact lattice stabilizer, the image of \(R_eR_\phi\) on \(L^2(Y)\) is finite dimensional.
\end{lemma}
\begin{proof}[Outline]
Use formal-degree orthogonality to project by convolution, then use height twists and Fourier inversion to show that unipotent averages vanish. Control the cusps by this vanishing and use equicontinuity on the compact part to bound dimension by finitely many evaluations. The complete proof is in \S\ref{prf:B8.03a}.
\end{proof}

\begin{theorem}[Assembly on an augmented extension]\label{B8.03}
Assume \(P\) and the separating rational row tests of \ref{B8.18}. Fix weak local correspondences \(h_L\) as in Theorem~\ref{B9.46}, \(c_{\rm Dep}\in\R\), and a condition \(A\) holding eventually on bases in the sense of (8.6). Assume the linked \(2,3\)-bound of \ref{B8.01} at \(c_{\rm Dep}\). Along \(p_A\) of (8.6), assume the covolume and functional rows closed on \(W_3(\mathcal A)\). For each eligible \(F_m,CH\), the covolume row gives \(C_0\ge0\) such that for each \(C\ge C_0\), the third-layer covolume bound of \ref{B8.34} at \(C,CH\) holds eventually at selected points. Under the family division estimate \eqref{eq:B8-divinN}, the functional row gives the bound of \ref{B8.23} with the same quantifiers. The point cutoffs may depend on \(C\). Then the common data of Definition~\ref{B8.02} exist with norm condition \(\mathrm{NL}\) and definite-side condition \ref{B8.35}.
\end{theorem}

\begin{proof}[Outline]
Transfers member data via \ref{B8.07}, \ref{B8.43h}, and \ref{B8.08}, and derives all member rows from \ref{B8.06} using \ref{B8.13}, \ref{B8.52}, \ref{B8.03a}, and \ref{B8.08b}. Invokes \ref{B8.04} with \ref{B8.44}, \ref{B8.48}, \ref{B8.49}, and \ref{B8.50}. The complete proof is in \S\ref{prf:B8.03}.
\end{proof}

\begin{theorem}[Filling the common record from all rows]\label{B8.04}
Let \(KP\) be a condition on a base and frame, \(Q,Q'\) member conditions, and \(L_k\) a bundle linkage. Fix \(N,DR\) as in \ref{B8.02}. Assume the following.
\begin{enumerate}
 \item The adjacent-rank mass formula of \ref{B8.48}.
 \item The uniform analytic estimates of \ref{B8.44}, the reproduction and weighted integral of \ref{B8.02y}, finite-dimensional row spaces as in \ref{B8.02}, the minor-label estimate of \ref{B8.10y}, and both signed upper bounds of \ref{B8.10d}.
 \item \ref{B8.49} for each datum satisfying the quotient-volume hypothesis and \ref{B8.50} for every datum.
 \item Assume Minkowski division implies the following four requirements. First, beyond the construction thresholds one can fix a preliminary base and frame such that any two prescribed finite sets of algebraic elements can be accommodated in the respective coefficient fields of a \(KP\)-base over that frame. Frame-preserving enlargement preserves \(KP\). Second, an eligible untwisted \(KP\)-base with the required primes admits a \(Q\)-member. Third, every such member supplies the norm-choice and theta data used in the next enlargement. Fourth, such a member admits an enlargement with the same frame and prime set to meet new avoidance and coefficient conditions, with a member satisfying \(Q,Q'\) and twisting compatibility.
 \item For \(KP,Q,Q',L_k\), assume the bundle supply defined before \ref{B8.05} and all 26 member conditions of \ref{B8.06}, using \(N\) for the norm condition and \(DR\) for the definite side.
\end{enumerate}
Then the common data of Definition~\ref{B8.02} exist for \(N,DR\), with
\[
 c_T(\varepsilon,y)=\max(11,16+3\max(5,y+1)). \tag{8.8}
\]
\end{theorem}

\begin{proof}[Outline]
Fix the counting constants in \ref{B8.06} and intersect the member rows with \ref{B8.51}. Use \ref{B8.05} to select a common family, choose joins pointwise with the closed thresholds, and fill the record of \ref{B8.02} using \ref{B8.05} and \ref{B8.02a}. The complete proof is in \S\ref{prf:B8.04}.
\end{proof}

\section{Choosing the family after imposing the rows}

Let \(L_k\) link points to bundles \(\mathrm{Bd}(D,O,c)\). Write \(Q\Rightarrow_{KP}R\) if, beyond the construction thresholds, every member satisfying \(Q\) on an eligible base and frame satisfying \(KP\) also satisfies \(R\), for every choice of the subsequent column, row, division and core data. Define the construction class \(\mathcal C\) as follows. Require the fixed family constants and all the pointwise equalities and rank bounds in clauses~1 and~2 of Definition~\ref{B8.01k}. At every \(x>x_1\) and every \(i\in I_x\), impose the following conditions on the point.
\begin{itemize}
\item The prime \(q_0\) of Definition~\ref{B8.01ka} belongs to the host level set, satisfies \(q_0\equiv1\pmod{p_*d_E^*}\), \(\log q_0\le x^\varepsilon\), \(q_0<X'\) and \(q_0\le l_J\), and differs from \(p_{\rm de},p'_{\rm de},g_0,g_1\) and every \(T_u\). Here \(l_J=mbx^\varepsilon/[k(m+t)\log(3(b+\mathrm{Sc}+\eta+t+1))]\), with the spacing and height parameters of Lemma~\ref{B8.02ac}. In particular this same prime supplies a host level prime congruent to one modulo \(d_E^*\) below \(X'\).
\item The tame labels pass clause~4 of Definition~\ref{B8.01k}. The zero-free boxes and auxiliary-prime tests of Theorem~\ref{B8.47} hold, with a scale in \([x^{1/J},x]\) for each test. For every positive integer \(M\) divisible by the orders of all finite subgroups of \(\GL_{2fN(m)}(\Q)\), and every word \(u\), require \(D_{\Gamma,u}\mid M\). The same divisibility holds for \(M=0\).
\item The point's own two weight functions belong to the boxes of Lemma~\ref{B5.22a} and satisfy its small-share bound \(R_W(i)\le x^{-9}\) as in Proposition~\ref{B8.02a}.
\item Each tuple is used as in clause~6 of Definition~\ref{B8.01k}. The point's own theta data pass the tests in clause~7 at every nonreference place for all the stated powers. These include the power-one test. For every tuple \(\mu\), the coefficient lattice \(\widetilde M_\mu\) is stable under multiplication by \(\iota_D(\mathcal O_{K_D})\).
\end{itemize}
For a condition \(U(D,O)\), let \(\mathcal C(U)\) consist of those families in \(\mathcal C\) for which \(U(D_{x,i},O_{x,i})\) holds at every \(x>x_1\) and \(i\in I_x\).
 A row \(\operatorname{Row}(Q,L_k,H)\) gives a bundle condition \(E(\kappa,D,O,c,\beta)\) on linked bundles of \(Q\)-points beyond the parameter and scale cutoffs. If each family is in \(\mathcal C\) with a linked parameter \(\kappa\), and each selected bundle is linked to a \(Q\)-point with that parameter and satisfies \(E\), the row gives \(H(\mathcal A,\sigma)\). Put
\[
 H\text{ holds after strengthening}\ \Longleftrightarrow\
 \exists R\,[(Q\Rightarrow_{KP}R)
 \land\operatorname{Row}(Q'\land R,L_k,H)].
 \tag{8.10}\label{eq:p8c1-row-large}
\]
For a member condition \(U\), write \(\mathsf B(U,L_k)\) if, beyond the point-construction thresholds, every \(U\)-point admits an \(L_k\)-linked bundle at every \(c\in\R\). Bundle supply for \(KP,Q,Q',L_k\) means that some \(R\) satisfies \(Q\Rightarrow_{KP}R\) and \(\mathsf B(Q'\land R,L_k)\).

\begin{lemma}[The scale before the word levels are chosen]\label{B8.05a}
Fix the construction constants and the exponents \(c_D,c_h\) of \ref{B8.02ac}, with \(c_D\ge0\) and \(c_h\ge402\). Use a good host of \ref{B2.56}, weights anywhere in the two prescribed boxes, and any prescribed code words. Write \(L=x^\varepsilon\), \(r=k+m-1\), and \(M_j\) for the standard definite genus mass in rank \(j\). For \(g=g_0,g_1\), let \(I_c(g)\) be the index of the upper unitriangular group in \(\GL_k(\F_g)\). Let \(I_r(g)\) be the index in \(\GL_r(\F_g)\) of the upper unitriangular group with the \(k\) entries above the diagonal in the last column of the middle block of size \(k+1\) set to zero. The middle block begins after \((m-2)/2\) coordinates.

For a decreasing parameter \(z\) of length \(j\), put
\[
 W_j(z)=\frac{\prod_{a<b}(z_a-z_b)}{\prod_{a=0}^{j-1}a!}.
\]
At the \(f\) embeddings of the CM type use the preliminary column parameter \(U_{v,\rho}\) of word \(v\), and the common row parameter \(\lambda_\rho=(v_{0,\rho},\ldots,v_{k-2,\rho},-b_{0,\rho},\ldots,-b_{m-1,\rho})\) in the coordinates of \ref{B4.02b}. The two pre-level factors are
\[
 \begin{aligned}
 d_v^\star&=M_k\prod_{g=g_0,g_1} I_c(g)^f
                         \prod_\rho W_k(U_{v,\rho}),\\
 D^\star&=M_r\prod_{g=g_0,g_1} I_r(g)^f
                         \prod_\rho W_r(\lambda_\rho).
 \end{aligned}
\]
With \(\Gamma=\exp(C_\Gamma m^5 f L)\), one has, for all sufficiently large \(x\),
\[
          \frac{1}{fk}\log\frac{D^\star}{d_v^\star\Gamma^2}
                    \ \ge\ mL
\]
uniformly in the prescribed weights and words.
\end{lemma}
\begin{proof}[Outline]
Combine adjacent-rank mass bounds from \ref{B8.48} with the two finite index ratios and the Vandermonde gap bounds from \ref{B4.02b}. The mass factorials cancel the Weyl denominators, and the prime lower bounds in \ref{B2.51} absorb the losses. The complete proof is in \S\ref{prf:B8.05a}.
\end{proof}

\begin{theorem}[A common family and selection]\label{B8.05}
Assume the mass formula of \ref{B8.48} and the volume ratio of \ref{B8.49} for every \(D\) satisfying its quotient-volume hypothesis. Fix \(KP,Q,Q',\mathrm{Bd},L_k,H\), and \(c_T:\R^2\to\R\). Assume item~4 of \ref{B8.04}, bundle supply for \(KP,Q,Q',L_k\), and \eqref{eq:p8c1-row-large} for \(H\). There exist \(\mathcal A,\sigma\) with the following properties.
\begin{enumerate}
 \item Minkowski division implies the design supply of Definition~\ref{B8.02} for \(\mathcal A,c_T\).
 \item Every \(F_m\in\mathcal A\) satisfies the seven duties of \ref{B8.01k}, with the following stronger clauses, and eventually all its coefficient lattices \(\widetilde M_\mu\) are stable under \(\iota_D(\mathcal O_{K_D})\). In clause~3, at each point above \(x_1=F_m.x_1\), the prime \(q_0\) is in the host level set and is \(1\) modulo the specified auxiliary-prime factor times the host conductor. It satisfies \(\log q_0\le L=x^\varepsilon\), \(q_0<X'\), \(q_0\le l_J\), and differs from \(p_{\rm de},p'_{\rm de},g_0,g_1\) and every \(T_u\). Here \(X'\) is the scale bound of \ref{B3.51} and \(l_J\) is the tame-generator limit. In clause~5 the model's own weights satisfy the box and share conditions of \ref{B8.02a} at every \(x>x_1\) and every point.
 \item Every \(F_m\in\mathcal A\) belongs to the construction class \(\mathcal C\) with a linked parameter \(\kappa\). At every \(x>F_m.x_1,\ i\in I_x\), the selected bundle comes from a linked \(Q'\)-point with that \(\kappa\).
 \item \(H(\mathcal A,\sigma)\) holds.
\end{enumerate}
\end{theorem}

\begin{proof}[Outline]
Intersect bundle and row conditions via \ref{B8.04} to define \(\mathcal A,\sigma\) and verify items 2--4 and \ref{B8.01k}. Under Minkowski division, Theorem~\ref{B2.56} hosts, \ref{B2.49}, and \ref{B5.22a} weights assemble the family of Definition~\ref{B8.02}. The complete proof is in \S\ref{prf:B8.05}.
\end{proof}

\section{The member rows}

For a property \(r\) of a member condition put
\[
 \operatorname{Lg1}(KP,Q,Q';r)\ \Longleftrightarrow\
 \exists R\,[(Q\Rightarrow_{KP}R)\land r(Q'\land R)].
 \tag{8.13}\label{eq:p8c1-lg1}
\]
Along \(p\), take \(Q,Q'\) to be constructed and good membership and \(L_k\) the point-to-bundle relation defined after (8.6). The base condition \(KP_x(B,f)\) has two clauses. For some field data \(\delta\), the designated subfield of \(\C\) attached to \(B\) is exactly the compositum of the host character's value field, the fields \(M'\) and the ray field of \(\delta\), the prescribed roots-of-unity field, the value field of the frame characters and the fields \(K_{\Psi,j}\) of \(\delta\) for \(1\le j\le m\). The standard embedding to level \(m\) must also send each starting frame vector to its counterpart.

\begin{lemma}[Row forms in the arithmetic coordinates]\label{lem:p8c1-row-forms}
Fix a constructed member with its row comparison. For a tuple \(\mu\), let \(S_\mu\) be the specified finite generating set of \(\widetilde M_\mu\) in its complex ambient space. Write \(\Lambda_\mu\) for the integral lattice of rational row forms before transport. For every word \(u\) there is a complex-linear identification \(e_{u,\mu}\) of the period row space with the ambient space. It is the prescribed reference identification at \(\mu_0\) and the twisted identification at every other tuple. If \(r_{u,\mu}(z)\) is the row form of a rational input \(z\), every element of \(S_\mu\) is \(e_{u,\mu}r_{u,\mu}(z)\) for some \(z\), and
\[
 z\in\Lambda_\mu
 \quad\Longleftrightarrow\quad
 e_{u,\mu}r_{u,\mu}(z)\in
 \sum_{v\in S_\mu}\iota_D(\mathcal O_{K_D})v .
\]
\end{lemma}

\begin{proof}
Fix \(\mu\). In the rational row coordinates of the construction each designated generator has a coordinate preimage. Choose one for each generator. The coordinate comparison realizes each such preimage by a rational row and preserves its ambient value. Conversely, an integral combination of the generators is realized by the same combination of their coordinate preimages, since the coordinate embedding is semilinear for \(\iota_D\). The coordinate lattice condition puts this vector in the integral lattice. Uniqueness of the rational form with its given complex realization identifies it with the original row. In the other direction the coordinate lattice condition sends every integral row into this span.

Now fix \(u\). Take the specified reference identification when \(\mu=\mu_0\) and the specified twisted identification otherwise. The row comparison says that pulling the ambient realization back by this identification gives the same form on every input as \(r_{u,\mu}\). Applying the identification again gives the ambient value just computed. The generator preimages and the lattice equivalence therefore hold for this word.
\end{proof}

\begin{theorem}[The twenty-six rows]\label{B8.06}
Assume global Artin reciprocity, the second convergent unitary Siegel--Weil formula of \ref{B8.13}, discrete-series cover and classification, Harish--Chandra modules, cusp decay, the Casselman--Wallach input, surjectivity onto the point space, and the lower-weight bound. Fix weak local correspondences \(h_L\) as in Theorem~\ref{B9.46}, a norm-choice extension \(\mathcal N\), an extra extension \(S\), and a geometry condition \(\Gamma\). For \(p=(R_{\rm row},\mathcal N,S,\Gamma)\), with the row comparison of (8.6), assume the open conditions of \ref{B8.08}, the covolume and functional implications on this extension, with the family and point quantifiers of \ref{B8.34} and \ref{B8.23}, respectively, and the linked row \(\mathcal A\mapsto NN(W_3(\mathcal A))\) for a predicate \(NN\) on family classes. Then all 26 member conditions hold for \(KP_x,Q,Q',L_k\), with \(NN\) as the norm condition and Definition~\ref{B8.35} on the definite side. Each has the form \eqref{eq:p8c1-lg1}, with closure on \(W_1\) or \(W_3\) where specified. The join type is nonempty, and the counting row gives both estimates with their constants.
\end{theorem}

\begin{proof}[Outline]
Take rows from \ref{B8.09}, \ref{B8.10}, \ref{B8.54}, \ref{B8.49} and link periods for \ref{B4.23}, \ref{B4.23a}. Check \ref{B5.12k} via \ref{B4.03}, \ref{B4.52}, intersect with \ref{B8.08}, \ref{B8.55}, \ref{B8.56}, \ref{B8.52}, \ref{B8.13}, and apply \ref{B8.38}. The complete proof is in \S\ref{prf:B8.06}.
\end{proof}

\begin{lemma}[An initial state with zero local contributions]\label{B8.07a}
Assume \(P\) and the rational test assertion of Theorem~\ref{B8.18}. Work beyond the starting-column thresholds of item~4 of Theorem~\ref{B8.04}. Fix an eligible starting column \(w\), including its base, frame and twist. There is a choice of row data over \(w\) and a compatible later state with the following properties. Its set of local Galois indices is empty, and its coordinate \(\eta_{j,v}\) is zero for every \(j\in\N\) and every finite coefficient place \(v\). In the untwisted case the selected set of places is a union of fibres of restriction to the fixed coefficient subfield.
\end{lemma}

\begin{proof}[Outline]
Pull back the untwisted reference support and use the rational tests of \ref{B8.18} with the multiplicity identity to lift the rows. Assemble the integral and spectral data with zero length coordinates and an empty local Galois index set. The complete proof is in \S\ref{prf:B8.07a}.
\end{proof}

\begin{lemma}[One cyclotomic specialization for the Selmer and theta polynomials]\label{B8.07b}
Let \(O_0\hookrightarrow O\) be an injective local homomorphism of discrete valuation rings of mixed characteristic \((0,\ell)\), with \(O_0\) of finite residue field. For either ring \(A\), write \(d_A=\operatorname{ord}_A(\ell)\), where a uniformizer has order one. On any discrete valuation ring extending \(O\) by an injective local map, use the valuation \(v_\ell\) with \(v_\ell(\ell)=1\).

Fix nonnegative integers \(R,D,b\). Write \(Z=(X_1,\ldots,X_R,Y_1,\ldots,Y_R)\). Suppose that for every sufficiently large integer \(M\) we have the following two finite families, with index sets and bounds independent of \(M\).
\begin{enumerate}
\item For each \(s\), the polynomial \(F_{M,s}\in O_0[Z]\) has a coefficient of total degree at most \(D\) and order at most \(\mu_s\).
\item For each \(y\), there are \(f_{M,y},A_{M,y},g_{M,y}\in O[Z]\). Suppressing the subscripts, suppose that
\[
 f-Ag\in\bigl(Y_i,(1+X_i)^{\ell^M}-1:1\le i\le R\bigr),
 \qquad \operatorname{ord}_O g(0)\le v_y .
\]
In the residue polynomial ring, suppose that
\[
 \overline A\equiv u\prod_{i=1}^R X_i^{n_i}
       \pmod{(Z_j^{\ell^M}:1\le j\le 2R)},
 \qquad u(0)\ne0,\quad 0\le n_i\le b .
\]
The polynomials and the exponents \(n_i\) may depend on \(M,y\).
\end{enumerate}
Here \(\mu_s,v_y\) are nonnegative integers. Given \(\delta>0\), there exist \(M\), \(1\le k\le M\), and integers \(0\le e_j<k\), with \(\ell^{e_j}/\varphi(\ell^k)\le\delta\), having the following property. For every primitive \(\ell^k\)-th root \(\zeta\) in a coefficient ring as above, set \(z_j=\zeta^{\ell^{e_j}}-1\). All the following values are nonzero and satisfy
\[
 v_\ell(F_{M,s}(z))\le \mu_s/d_{O_0}+\delta,\qquad
 v_\ell(f_{M,y}(z))\le v_y/d_O+\delta .
\]
\end{lemma}

\begin{proof}[Outline]
Use a diagonal subsequence and a weighted substitution to find coefficients of bounded degree and valuation, using a weighted valuation for the theta products. Choose a cyclotomic level avoiding ties for the least term valuation and evaluate at \(\zeta-1\). The complete proof is in \S\ref{prf:B8.07b}.
\end{proof}

\begin{lemma}[The Selmer length at the common character]\label{B8.07c}
Fix a twist index \(a\), write \(\ell=\ell_a\), and let \(\mathcal S\) be its finite set of good signatures. Write \(L_s(\eta,Q)\) for the normalized global Selmer length of Definition~\ref{B8.01k}, with the character acting on the coefficient module by \(\eta^{-1}\). Fix a finite set of forbidden places of the host field \(E\). Suppose the following data are given for every sufficiently large \(M\), with \(R\) independent of \(M\).

There are \(2R\) split primes \(P_j\) of \(E\) whose conjugate pairs are disjoint. Their union \(Q\) avoids the forbidden places and the places above \(\ell\); its places have residue degree and ramification index one over \(\Q\). Write \(G_E=\Gal(\overline E/E)\). For each \(j\) there is a continuous anticyclotomic character \(d_j:G_E\to\Z/\ell^M\), trivial at the decomposition groups above \(\ell\), unramified outside the pair of \(P_j\), and surjective on inertia at both places of that pair.

Put \(\Delta=(\Z/\ell^M)^{2R}\) and let \(b_j\) be its coordinate generators. For each \(s\in\mathcal S\), suppose there is a square matrix \(\Phi_s\) over \(O_0[\Delta]\) presenting the Selmer condition in the following sense. For each designated coefficient ring \(A\) at level \(k\) and roots \(u_j\in A\) of order dividing \(\ell^M\), put \(\eta(g)=\prod_j u_j^{d_j(g)}\). The Selmer group for \(s,\eta,Q\) is identified over \(A\) with the kernel of the matrix specialized by \([b_j]\mapsto u_j\) on copies of \(\operatorname{Frac}(A)/A\). Suppose \(\det\Phi_s\) has a polynomial representative \(F_{M,s}\) under \(Z_j\mapsto[b_j]-1\) satisfying the first hypothesis of Lemma~\ref{B8.07b} with \(\mu_s=d_{O_0}\nu_a\). Include in that lemma any finite family of theta polynomials satisfying its second hypothesis on these same primes.

Then for any \(e>0\) there is one ray character \(\xi\) of \(\ell\)-power order such that, for every \(h\in\Z\) prime to \(\ell\) and every choice of signs \(\sigma_s\in\{1,-1\}\),
\[
 L_s(\xi^{h\sigma_s},Q)\le \nu_a+\delta,\qquad
 \delta=\frac{e}{1+|\mathcal S|+2m+4Rt}.
\]
Here \(m\) is the chain length and \(t\) the auxiliary rank. At the same character the theta polynomials have the bounds of Lemma~\ref{B8.07b}, with error \(\delta\), for either sign of \(h\). The character is ramified exactly at \(Q\) and is trivial above \(\ell\). The sum of the signature errors is at most \(e\), as is \(t\) times the sum of the tame valuations.
\end{lemma}

\begin{proof}[Outline]
Use \ref{B8.07b} to build one ray character and bound its tame valuations. For signed powers prime to \(\ell\), elementary divisors give Selmer lengths as normalized determinant valuations, and the same specialization bounds the theta values. The complete proof is in \S\ref{prf:B8.07c}.
\end{proof}

\begin{definition}[Twisting data of an untwisted member]\label{B8.07d}
Let \(E\) be the CM host field. Call a twist index \(a\) eligible if restriction to the fixed coefficient subfield gives the place of \(a\) at some coefficient place \(w\) with \(\mu_w\ne\mu_0\). Write \(\ell=\ell_a\), let \(c\) act on \(G_E=\Gal(\overline E/E)\) by complex conjugation, and write \(\varepsilon_\ell\) for the cyclotomic character. For a good signature \(s\) put \(T_s=\rho_s\chi_s^{-1}\), with the conventions of Definition~\ref{B8.01k}. A parenthesized integer on a representation denotes a Tate twist.

The conditions on an eligible index are the following.
\begin{enumerate}
\item For each good \(s\) there is an integral isomorphism
\[
                 (T_s^\vee(1))^c\simeq T_s(t).
\]
At every place above \(\ell\), for each embedding of the local field, every jump of the de Rham filtration of \(T_s\) lies outside \([-1,t]\).
\item For each abelian theta label \(\lambda\), the sum of its untwisted ray coefficients at empty auxiliary conductor is nonzero. Its prescribed doubled integer parameter is \(2u\) with
\[
 1\le u\le k+m(t+b_r+1)+1,
 \qquad b_r=\text{the odd member of }\{r+6,r+7\}.
\]
\item The ranks and type primes satisfy \(100\le r\) and \(2r\le g_i\) for \(i=0,1\).
\end{enumerate}

The twisting data exported by the member consist of nonnegative integers \(N_0(a)\) and finite sets \(S(a,j,\sigma)\) of its allowed theta data, with the following properties.
At every \(w\) with \(\mu_w\ne\mu_0\), the chosen index satisfies the preceding conditions. For each \(1\le j\le m\) and sign \(\sigma\), the projected allowed theta vectors span the prescribed space over the fraction field. The projections indexed by \(S(a,j,\sigma)\) form a basis of their integral span taken over \emph{all} allowed data.

There is also a uniform assertion at every eligible index satisfying the three conditions. Suppose the projected vectors have the indicated full rank. Fix \(e>0\) and any finite sets of allowed data giving integral bases as above. For every finite set \(A\) of places of \(E\), there is a triple \((Q,k',\xi)\) for which all powers \(\xi^h\), \(h\in\Z\) prime to \(\ell\), pass the character, projected-lattice, Selmer, abelian theta and tame tests of Definition~\ref{B8.01k}(7), with error \(e\). The conductor set avoids \(A\), the host exclusions and the levels of the chosen basis vectors. The sign at a signature or a layer is applied to the tested power. The integers \(N_0(a)\) serve for every choice of \(e\), bases and \(A\).
\end{definition}

\begin{lemma}[Polarization at the Selmer normalization]\label{B8.07e}
Let \(T=T_s\) be the representation of an attached good signature, over its base coefficient discrete valuation ring \(O_0\). Suppose its residual image spans the full matrix algebra. In the unramified correspondence for \(T\), the attached conjugate self-dual cuspidal representation has shift \((1-t)/2\). For any finite-order anticyclotomic character \(\eta\) over an extension of \(O_0\), put \(V=T\eta^{-1}\). Then
\[
       (V^\vee(1))^c\simeq V(t).
\]
Over the fraction field \(V\) and \((V^\vee(1))^c\) are inequivalent. If their reductions are equivalent over an extension of the residue field, then \(\nu_a>0\).
\end{lemma}

\begin{proof}[Outline]
Compare split Frobenius eigenvalues and use \ref{B8.42p} to obtain equal traces. Full matrix span and the trace pairing give integral conjugacy, which survives anticyclotomic twisting, and determinants give the cyclotomic obstruction. The complete proof is in \S\ref{prf:B8.07e}.
\end{proof}

\begin{lemma}[Ray coordinates with prescribed Frobenius values]\label{B8.07f}
Fix the host, \(\ell\), and \(M\ge1\), and put \(n=\ell^M\). Let \(T_0\) be the places above \(p_{\rm de}\) if \(\ell\ne p_{\rm de}\), and above \(p'_{\rm de}\) otherwise. There are an open normal subgroup \(H\subset G_E\), containing the inertia groups at \(T_0\), and a finite exceptional set \(S\) with the following property. At every split prime \(P\notin S\) whose decomposition group lies in \(H\), there is an additive character \(d_P:G_E\to\Z/n\) with the local properties of Lemma~\ref{B8.07c}.

Let \(U\triangleleft G_E\) be open, let \(z\) belong to inertia at a place of \(T_0\), and let \(h_1,\ldots,h_R\in H\) fix \(\mu_{2n}\). Outside any additional finite set one can choose \(R\) primes \(P_i\) followed by \(R\) primes \(P'_i\), all with disjoint conjugate pairs, with local degree one over \(\Q\) and norms \(1\bmod 2n\). At the first primes the Frobenius conjugacy class modulo \(U\) is that of \(z\). At \(P'_i\) it is that of \(h_i\), and
\[
                    d_{P_j}(\Frob_{P'_i})=\delta_{ij}.
\]
All these primes have the indicated ray coordinates.
\end{lemma}

\begin{proof}[Outline]
Principalize ideals over the \(S\)-integers and use Artin reciprocity in \ref{B1.15ac} to form anticyclotomic ray coordinates. Apply \ref{B8.42p} to choose the first primes, then use their inertia to correct the second Frobenius targets before applying it again. The complete proof is in \S\ref{prf:B8.07f}.
\end{proof}

\begin{lemma}[The square matrix for the Selmer condition]\label{B8.07g}
Use a finite family of ray coordinates from Lemma~\ref{B8.07f}. Let \(Q\) be the union of their conductor pairs, \(\Delta\) the product of their cyclic targets, \(\Lambda=O_0[\Delta]\), and \(\kappa:G_E\to\Lambda^\times\) the inverse of their product in \(\Delta\subset\Lambda^\times\). Fix a good signature and write \(T=T_s\). For a discrete \(\ell\)-primary \(\Lambda\)-module \(B\), use \(T\otimes_{O_0}B\) with the additional action of \(\kappa\). Denote its Selmer group by \(S(B)\). Use the local conditions of Definition~\ref{B8.01k}, taking above \(\ell\) the \(B\)-linear combinations of integral finite cocycles, that is, finite sums of cocycles \(g\mapsto z(g)b\) with \(b\in B\).

Suppose the residual image of \(T\) spans the matrix algebra, \(r\ge2\), and \(T\) and its Tate dual are unramified outside the indicated exceptional places. At \(v\mid\ell\) suppose \(H^0(E_v,T\otimes\operatorname{Frac}(O_0))=0\) and the local Bloch--Kato dimension formula applies. Suppose also that the sums of the dimensions of \(D_{\rm dR}/\mathrm{Fil}^0\) at a pair of conjugate embeddings equal \(r\). Use the finite-coefficient cohomology hypotheses in Definition~\ref{B1.15ac}. There is a square matrix \(\Phi\) over \(\Lambda\), with entries in its maximal ideal, for which
\[
                         S(B)\simeq\ker(\Phi:B^a\longrightarrow B^a)
\]
naturally in \(B\).
\end{lemma}

\begin{proof}[Outline]
The local cocycle counts, Euler characteristics from \ref{B1.15ac}, and duality in \ref{B6.09} bound the obstruction to lifting Selmer classes. Additive duality gives a representing module, and that bound gives at most as many relations as generators. The complete proof is in \S\ref{prf:B8.07g}.
\end{proof}

\begin{lemma}[The coefficient detected by first-order variation]\label{B8.07h}
Let \(P_i,P'_i\), \(1\le i\le R\), be primes as in Lemma~\ref{B8.07f}. Use the universal matrix of Lemma~\ref{B8.07g} for a fixed signature. Suppose the residual local cohomology in degree one vanishes at the first pairs, and the residual action of \(T\) and the cyclotomic character are trivial on the decomposition groups at the second pairs. Put \(k_0=O_0/\mathfrak m_0\). Let \(S_0,S_0'\) be the residual Selmer spaces for \(T\) and its Tate dual at empty auxiliary conductor, with orthogonal local conditions. Use the dot product for the Tate pairing after choosing a generator of the residual Tate line.

Choose an arithmetic Frobenius \(\phi_i\) at each \(P'_i\). On cocycle representatives \(x\in S_0\), \(y\in S'_0\), define
\[
 B_i(x,y)=x(\phi_i)\cdot y(\phi_i)-x(c\phi_i c^{-1})\cdot y(c\phi_i c^{-1}).
\]
Suppose a linear combination \(\sum_i a_iB_i\) has left kernel of dimension at most \(b\), where \(b=0\) or \(1\). If \(4rR+2<\ell^M\), the following hold.
For \(b=0\), every polynomial representative of \(\det\Phi\) has a unit coefficient of total degree at most \(4rR\).
For \(b=1\) and a nonempty matrix, there is an entry of \(\operatorname{adj}\Phi\) for which the same conclusion holds for every representative.
\end{lemma}

\begin{proof}[Outline]
Use reciprocity in \ref{B6.09} and the ray coordinates of \ref{B8.07f} to bound the dimension of the kernel that lifts to first order. Apply elementary divisors to the matrix of \ref{B8.07g} to detect a low-degree coefficient of the determinant or a cofactor. The complete proof is in \S\ref{prf:B8.07h}.
\end{proof}

\begin{lemma}[The cyclotomic factor in the integral cocycle calculation]\label{B8.07i}
Let \(A\) be a commutative local \(O_0\)-algebra. On \(A^r\) let \(U\) be the representation obtained from \(T_s\) by an anticyclotomic character, and put \(U'=U^\vee(1)\). Suppose the matrices of \(U\) span \(M_r(A)\), and write
\[
 U'(cgc^{-1})=J\,\varepsilon_\ell(g)^t U(g)J^{-1}
\]
with \(J\in\GL_r(A)\). Let \(L\triangleleft G_E\) be stable under \(c\), with \(U\) and \(U'\) trivial on \(L\). Suppose there is \(z\) with \(\varepsilon_\ell(z)=1\), \(U(z)=d\,1_r\), and \(d^2-1\in A^\times\). For cocycles \(x,y\) of \(U,U'\), respectively, write \(\bar h=chc^{-1}\). If
\[
                  x(h)\cdot y(h)=x(\bar h)\cdot y(\bar h)
                         \qquad(h\in L),
\]
then, for all \(u,v\in L\) and all coordinates \(i,j\),
\[
                       \ell^{\nu_a}x(u)_i y(v)_j=0 \quad\text{in } A.
\]
\end{lemma}

\begin{proof}[Outline]
Polarize the cocycle identity and use the scalar at \(z\) to isolate the cross terms. Modulo their annihilator, the matrix-unit argument of \ref{B8.07e} gives determinant equality, so the cyclotomic factor from \ref{B8.01k} kills every coordinate product. The complete proof is in \S\ref{prf:B8.07i}.
\end{proof}

\begin{lemma}[A finite coefficient ring detecting a cofactor]\label{B8.07j}
Write \(d=\operatorname{ord}_{O_0}(\ell^{\nu_a})\), and let \(\varpi\) be a uniformizer. Suppose a polynomial \(G\in O_0[X_1,\ldots,X_N]\) has a unit coefficient of total degree at most \(w\) on a monomial involving none of the variables in a specified subset \(J\). Put \(L=(d+1)w\). There are a finite local \(O_0\)-algebra \(A\), coefficients \(a_i\in A\) zero for \(i\in J\) and units otherwise, and \(u\in A\) such that
\[
 \varpi^{d+1}=0,\qquad u^{2L+2}=0,\qquad
 \ell^{\nu_a}u^{2L+1}\ne0,\qquad G((a_i u)_i)\mid u^L .
\]
\end{lemma}

\begin{proof}[Outline]
Set the specified variables to zero and choose nonzero residue-field coordinates detecting a homogeneous part. Lift them to units in a finite free local algebra and adjoin nilpotent \(u\), using the residue order to get divisibility and the free basis for nonvanishing. The complete proof is in \S\ref{prf:B8.07j}.
\end{proof}

\begin{lemma}[Local conditions for the representations of a member]\label{B8.07k}
Assume the hypotheses of Theorem~\ref{B8.07}. Take its untwisted constructed point, with its chosen column, and an eligible index \(a\) of Definition~\ref{B8.07d}. Write \(O_0\) for its base coefficient ring, \(k_0\) for the residue field, \(\ell=\ell_a\), and \(T=T_s=\rho_s\chi_s^{-1}\) for a good signature. Put \(T'=T^\vee(1)\). Both representations are unramified outside the places over \(\ell\) and the exceptional places in Definition~\ref{B8.01k}(7). The local hypotheses of Lemma~\ref{B8.07g} and the three index conditions of Definition~\ref{B8.07d} hold.

Let \(D_0\) consist of the places over \(p_{\rm de}\) if \(\ell\ne p_{\rm de}\), and over \(p'_{\rm de}\) otherwise. There is a single inertia element \(z\) at \(D_0\), independent of the signature, such that
\[
 \varepsilon_\ell(z)=1,\qquad T_s(z)=d_s1_r,\qquad
 d_s^2-1\in O_0^\times .
\]
At the exceptional places away from \(\ell\), both residual representations have zero local cohomology in degrees zero and one.
\end{lemma}

\begin{proof}[Outline]
Use \ref{B7.51} and \ref{B8.42pa} for residual irreducibility and tame inertia for local vanishing. The attachments in \ref{B9.06} and \ref{B9.07} give the filtration, \ref{B8.07e} gives polarization, and \ref{B8.11h} gives the toric nonvanishing. The complete proof is in \S\ref{prf:B8.07k}.
\end{proof}

\begin{lemma}[Choosing the residual Frobenius tests]\label{B8.07l}
Fix \(a\) and a finite set of good signatures as in Lemma~\ref{B8.07k}. There is a number \(R_0\), independent of the ray depth \(M\), with the following property. In Lemma~\ref{B8.07f} the first \(R_0\) second primes can be prescribed so that the left-kernel hypothesis of Lemma~\ref{B8.07h} holds for every signature, with \(b=0\) if \(\nu_a=0\) and \(b=1\) otherwise. One may impose an additional open normal congruence on the first primes, requiring their Frobenius to be conjugate to the specified \(z\) modulo that subgroup, and enlarge the finite exclusion set.
\end{lemma}

\begin{proof}[Outline]
Duality in \ref{B8.41h} and \ref{B6.09} gives equal Selmer dimensions. Using \ref{B8.07k} and \ref{B8.07e}, show the restricted cocycle pairings span all forms or all alternating forms, then realize a finite set of tests at primes by \ref{B8.07f}. The complete proof is in \S\ref{prf:B8.07l}.
\end{proof}

\begin{lemma}[Preserving the tests when a pair is added]\label{B8.07m}
Suppose \(\nu=\nu_a>0\). Put \(I=(\ell^\nu)\mathfrak m_0\). Fix \(R\) first primes as in Lemma~\ref{B8.07l}. At a stage with \(n\le R\) second primes use variables \(X_0,\ldots,X_{R-1}\) for the first primes and \(Y_0,\ldots,Y_{n-1}\) for the second primes. The first ray coordinates at a second Frobenius are the corresponding coordinate vector. The second decomposition groups lie in the group \(L\) of Lemma~\ref{B8.07l}. Use the union of all conductor pairs for the relaxed Selmer condition.

For a signature and a positive integer \(c_0\), consider the following tests on \emph{every} square presentation of this condition. The determinant test requires every polynomial representative of the determinant to have a coefficient outside \(I\) in degree \(<c_0\). The cofactor test permits a matrix of size zero. For any positive size it requires an entry of the adjugate such that every representative has a unit coefficient in degree \(<c_0\) on a monomial avoiding \(X_n,\ldots,X_{R-1}\).

Suppose \(n<R\), \(M\ge\nu+2\), and \(c_0+2r\le\ell^{M-\nu-1}\). Adjoin any second prime satisfying the same conditions for index \(n\), with a new conductor pair. Each test already satisfied at \(c_0\) holds at \(c_0+2r\) for the enlarged family.
\end{lemma}

\begin{proof}[Outline]
Recover the old Selmer module by killing \(2r\) inertia generators annihilated by \(p\). The resulting Fitting ideal inclusions multiply old minors by \(p^{2r}\), whose leading term preserves detection within the enlarged degree cutoff. The complete proof is in \S\ref{prf:B8.07m}.
\end{proof}

\begin{lemma}[A dual class detected by an old cofactor]\label{B8.07n}
Use the stage \(n<R\) of Lemma~\ref{B8.07m}. Let \(d=\operatorname{ord}_{O_0}(\ell^\nu)\), \(w\ge0\), \(l=(d+1)w\), and suppose
\[
               M\ge\nu+2,\qquad 2l+2+2r\le\ell^{M-\nu-1}.
\]
Suppose a signature passes the cofactor test at \(w+1\) and fails the determinant test at \(2l+2\). There are a finite local specialization \(A\) of the coefficients, an old presentation \(\Psi\) specialized to \(A\), indices \(i,j\), and \(u\in A\), with
\[
 I A=0,\quad u^{2l+2}=0,\quad
 \ell^\nu u^{2l+1}\ne0,\quad \det\Psi=0,\quad
                         (\operatorname{adj}\Psi)_{ij}\mid u^l .
\]
Let \(x\) be the old Selmer class corresponding to column \(j\) of \(\operatorname{adj}\Psi\). There is a cocycle \(y\) for the specialized Tate dual with coefficients in \(A\), orthogonal to the old Selmer condition with additively dual coefficients, for which \(\ell^\nu u^l[y]\ne0\). On the subgroup \(L_1\) obtained from \(L\) by killing all old ray coordinates except the unused first coordinates \(n,\ldots,R-1\), some \(h\) satisfies \(B_h(x,y)\ne0\), with the formula of Lemma~\ref{B8.07h} read in \(A\).
\end{lemma}

\begin{proof}[Outline]
Use \ref{B8.07j} and \ref{B8.07m} to kill the determinant while detecting a cofactor. The order formula \ref{B6.09} forces a surviving dual class, and \ref{B8.07i} with the commutator argument of \ref{B8.07l} forces a nonzero pairing. The complete proof is in \S\ref{prf:B8.07n}.
\end{proof}

\begin{lemma}[The relation determinant at a new prime]\label{B8.07o}
Under the numerical hypotheses of Lemma~\ref{B8.07n}, a signature passing the cofactor test at \(w+1\) admits an extension by one second prime for which it passes the determinant test at \(2(d+1)w+2+2r\). The first primes retain their prescribed Frobenius classes. All conductor pairs remain disjoint and avoid the prescribed finite set.
\end{lemma}

\begin{proof}[Outline]
Carry forward a passing determinant by \ref{B8.07m}. Otherwise use \ref{B8.07n} and \ref{B8.42p} to choose a prime with nonzero pairing, then use tame relations and reciprocity \ref{B6.09} to make that pairing a coefficient of a relation determinant. The complete proof is in \S\ref{prf:B8.07o}.
\end{proof}

\begin{proposition}[Uniform degree on a common set of primes]\label{B8.07p}
For the actual good signatures, the Selmer polynomial hypothesis of Lemma~\ref{B8.07c} holds with \(\mu_s=\operatorname{ord}_{O_0}(\ell^{\nu_a})\). The numbers of primes and the degree bound are independent of \(M\), of the prescribed first-prime congruence, and of the finite exclusions.
\end{proposition}

\begin{proof}[Outline]
Start with the primes of \ref{B8.07l} and tests of \ref{B8.07h}. For positive \(\nu_a\), add one prime per signature by \ref{B8.07o} and preserve earlier tests by \ref{B8.07m}, then take polynomial representatives of the matrices from \ref{B8.07g}. The complete proof is in \S\ref{prf:B8.07p}.
\end{proof}

\begin{lemma}[The integral class form giving a specialized theta vector]\label{B8.07r}
Fix an eligible index of Definition~\ref{B8.07d}, a layer and a sign. Let \(O\) be the coefficient discrete valuation ring at this index. For an admissible twist \(t\), let \(O_t\) be its designated coefficient ring. Let \(M_t\) be the lattice of definite class forms whose value at each finite adelic point belongs to the coefficient lattice at that point. Write \(A_t\) for the theta map on class forms, with the reciprocal stabilizer orders in its class sum. Every allowed specialized theta generator belongs to \(A_t(M_t)\). This includes the trivial character at empty auxiliary conductor.
\end{lemma}

\begin{proof}[Outline]
Form an integral class form by summing the transported polynomials over the contributing rational elements. The dual kernel law makes their theta values equal, canceling the reciprocal stabilizer order, and a determinant twist unit gives the required generator. The complete proof is in \S\ref{prf:B8.07r}.
\end{proof}

\begin{lemma}[Removing the frame scalar from a lattice bound]\label{B8.07s}
Let \(O\) be a discrete valuation ring with fraction field \(K\subset\C\), and suppose the rational prime \(\ell\) lies in its maximal ideal. Let \(V\) be a finite-dimensional \(K\)-space, \(B\subset V\) an \(O\)-lattice, and \(r:V\to W\) a \(K\)-linear reading into a complex vector space. Suppose the readings of a \(K\)-basis of \(V\) are linearly independent over \(\C\). Let \(x_\alpha\in V\) be any family. If a fixed \(\lambda\in\C^\times\) satisfies
\[
                       \lambda r(x_\alpha)\in r(B)
                       \qquad\text{for every }\alpha ,
\]
there is \(N\ge0\) such that \(\ell^N x_\alpha\in B\) for every \(\alpha\). The same \(N\) works after adding coefficients with the same image under a map of finite ray groups.
\end{lemma}
\begin{proof}[Outline]
For \(x_\alpha\ne 0\), \(\lambda r(x_\alpha)=r(y)\) (\(y\in B\)) and \(\mathbb{C}\)-linear independence of basis readings give \(\lambda\in K^\times\), so \(\lambda x_\beta\in B\). Taking \(N\) with \(\ell^N/\lambda\in O\) gives \(\ell^N x_\beta\in B\), preserved by fibre sums. The complete proof is in \S\ref{prf:B8.07s}.
\end{proof}

\begin{proposition}[Twisting data of the constructed point]\label{B8.07q}
Under the hypotheses of Theorem~\ref{B8.07}, each of its untwisted constructed points has the twisting data of Definition~\ref{B8.07d} for its own norm choices and theta series.
\end{proposition}

\begin{proof}[Outline]
Bounds toric denominators by \ref{B1.15a} and builds full-rank projected spans via \ref{B8.07r} and \ref{B8.18}. Computes auxiliary-pair factors on ray polynomials via \ref{B8.43h}, \ref{B8.07s}, and \ref{B8.07k}, then applies \ref{B8.07b}, \ref{B8.07p}, and \ref{B8.07c}. The complete proof is in \S\ref{prf:B8.07q}.
\end{proof}


\begin{lemma}[The local indices and lattices of a twisted layer]\label{B8.07t}
Fix the rational row realization at the finite level of the member construction. Let \(K_0\) be its coefficient field, \(m\) its number of layers, and \(w\) a finite place of \(K_0\) of residue characteristic \(\ell\). Write \(\mu(w)\) for the tuple selected at \(w\) and \(\mu_0\) for the reference tuple. Suppose \(1\le j\le m\) and \(\mu(w)\ne\mu_0\). Let \(W\) be the rational space of the twisted layer at \((j,w)\).

Use the split Hecke operators on \(W\) of orders \(1,\ldots,n_V\), where \(n_V\) is the rank of the group acting on the layer. The places used are outside the conductor of the twist and away from \(\ell,g_0,g_1\) and \(p_{\rm de}\), with \(p'_{\rm de}\) in place of \(p_{\rm de}\) when \(\ell=p_{\rm de}\). Let \(K'\) be the local coefficient field chosen for these operators at \(w\), let \(O\) be its valuation ring, and put \(V=K'\otimes_{K_0}W\). Let \(\mathbb T\subseteq\operatorname{End}_{K'}(V)\) be the \(O\)-algebra generated by the indicated operators.

The indices over \((j,w)\) in the localized realization are precisely
\[
                  (j,w,\mathfrak m),\qquad
                  \mathfrak m\in\operatorname{Max}\mathbb T.
\]
For such an index \(i\), the coefficient ring \(O_i\) is \(O\). Let \(V_{\mathfrak m}\) be the sum of the joint generalized eigenspaces whose characters reduce with kernel \(\mathfrak m\). Write \(p_{\mathfrak m}\) for projection onto that sum along the other summands. If \(S\subset W\) is the inverse image of the specified finite set of integral layer generators under the rational row embedding, then the local module and its comparison with the row are
\[
 M_i=\sum_{s\in S}O\,p_{\mathfrak m}(1\otimes s),
 \qquad
 K'\otimes_O M_i \xrightarrow{\ \sim\ } V_{\mathfrak m},
 \qquad a\otimes z\longmapsto az.
\]
In these coordinates the map from \(W\) to the local module after inverting \(\ell\) is \(s\mapsto p_{\mathfrak m}(1\otimes s)\). In particular,
\[
          \sum_{i\in I_{j,w}}\operatorname{rank}_{O_i}M_i
                         =\dim_{K_0}W .
\]
The length coordinate of the localized member is the sum in Definition~\ref{B8.40} formed with these \(M_i\) and their Hecke actions. At a reference place the localized index set is empty.
\end{lemma}

\begin{proof}[Outline]
Split the commuting integral Hecke action and group the joint eigenspaces by residual maximal ideal. Chinese remainder projectors decompose the layer lattice, and the quotient in \ref{B8.40} gives the sum of the normalized local lengths. The complete proof is in \S\ref{prf:B8.07t}.
\end{proof}

\begin{theorem}[Member and bundle supply on the same extension]\label{B8.07}
Assume \(P\) and the separating rational row tests of \ref{B8.18}. For every choice \(h_L\) of weak local correspondences as in Theorem~\ref{B9.46}, \(c_{\mathrm{Dep}}\in\R\), and every \(A\) holding eventually on bases, the following hold on the parameters of (8.6):
\[
 \text{the mass ratio \eqref{eq:p8c5-mass-ratio}},\qquad \text{the member conditions for }KP_x,Q_A,Q'_A,
 \qquad \text{bundle supply for }KP_x,Q_A,Q'_A,L_A. \tag{8.14}
\]
The member assertion consists of the four conditions in item~4 of \ref{B8.04}.
\end{theorem}

\begin{proof}[Outline]
Apply \ref{B8.48} for the mass ratio and construct members using \ref{B8.07a} and \ref{B8.07t}. Use \ref{B8.07q} to choose twists before coefficient enlargement, then match the kernel and spectral data in the common frames to form linked bundles. The complete proof is in \S\ref{prf:B8.07}.
\end{proof}

\begin{lemma}[The relative scalar at a finite coefficient place]\label{B8.08a}
In a linked realization on the list extension, assume \(P\) and work beyond the cutoffs for the row comparison. Fix a word \(u\). For a tuple \(t\), let \(Q_t\) be its chosen set of twisting places of the host CM field, and let \(S_t\) be the set of rational primes below \(Q_t\). Write \(t_0=\mu_0\) for the reference tuple. The unnormalized comparison scalars \(c_t\ne0\) can be chosen so that
\[
       (c_t/c_{t_0})^N=r,\qquad
       N\in\N_{>0},\quad r\in\Q^\times,\quad
       v_p(r)=0\quad(p\notin S_t).
\]
For every finite coefficient place \(w\), its residue characteristic is outside \(S_{\mu_w}\).
\end{lemma}

\begin{proof}[Outline]
Cancel level indices against the averaging factors and rescale the kernel charts to remove the intermediate phases. The remaining scalar ratio is a root of unity, and the selected twist is unramified above the residue characteristic of the coefficient place. The complete proof is in \S\ref{prf:B8.08a}.
\end{proof}

\begin{lemma}[The range of the arithmetic row embedding]\label{B8.08d}
At a linked point of \textup{(8.6)}, take the spectral realization at the reference row level and its prescribed normalization. Let \(H\) be the image of the reference finite-level and lowest-type projectors on the automorphic \(L^2\) space. Let \(G\) be the sum of the good finite-type idempotents on \(H\), and \(E_{\rm ar}\) the complex arithmetic reference row space with its selected orthonormal basis. The conjugate of the model row embedding, after the prescribed identification of ambient spaces, maps \(E_{\rm ar}\) onto \(GH\). For each good label \(t\), its selected basis vectors span the whole single-label space.
\end{lemma}
\begin{proof}[Outline]
Use formal-degree orthogonality to identify level intertwiners by evaluation at a lowest-type vector, with cuspidality from \ref{B8.02y} and \ref{B8.03a}. Transport the good projector through the arithmetic coordinates and apply \ref{B7.44} label by label. The complete proof is in \S\ref{prf:B8.08d}.
\end{proof}

\begin{lemma}[First occurrence on the row chain]\label{B8.08g}
Put \(r=k+m-1\). The row spaces at consecutive stages \(j,j+1\) have ranks \(r+jt,r+(j+1)t\). Use the real theta correspondence in the convention of the polynomial model, taking the contragredient parameter on passing from the row representation to its dual. At a fixed real place subtract the half-exponent of the target splitting character from the Harish--Chandra parameter. Suppose the resulting two lists for the smaller member are
\[
                 (A,B\,;\,C,D),
\]
where \(A,C\) have positive entries, \(B,D\) have negative entries, all entries are distinct integers, and all have absolute value at least \((t+3)/2\). Here \(t\) is positive and odd. On the larger member the centered lists are
\[
        (A,I_t,D\,;\,C,B),\qquad
        I_t=((t-1)/2,(t-3)/2,\ldots,-(t-1)/2).
\]
For the smaller member as source, its local theta lift vanishes on removing any positive number of hyperbolic planes from the indicated target. The same holds with the larger member as source and the smaller member as target. Consequently a global lift with these data at a real place and with cuspidal source is cuspidal whenever the tower constant-term formula of Theorem~\ref{B3.00e} applies.
\end{lemma}

\begin{proof}[Outline]
Apply Atobe's discrete series criterion to identify first occurrence in both directions. Vanishing of the lower lifts gives cuspidality by the constant-term formula of \ref{B3.00e}. The complete proof is in \S\ref{prf:B8.08g}.
\end{proof}

\begin{lemma}[The value of the chain numerator]\label{B8.08e}
Let \(S\) contain the exceptional finite places of the CM field \(E\). At \(w\notin S\) let \(B_w\) be the multiset of bottom parameters, of cardinality at most \(r\), with \(|\alpha|\le (Nw)^{r/2}\) for \(\alpha\in B_w\). Let \(\mu_1,\ldots,\mu_m\) be the unitary splitting characters of the row chain, unramified off \(S\), with the prescribed distinct archimedean exponents at one CM embedding. For \(0\le j<m\), twist the bottom parameters by \(\mu_{j+1}^{-1}\) and adjoin, for \(1\le i\le j\), the local roots
\[
 (\mu_i/\mu_{j+1})(\varpi_w)(Nw)^{\,l-(t-1)/2},
       \qquad 0\le l<t.
\]
Here \(\varpi_w\) is a uniformizer and \(t\ge r+5\). The partial Euler product of this multiset, evaluated in the variable \(s+1/2\), has a holomorphic continuation to a neighborhood of \(s=t/2\) with nonzero value there.
\end{lemma}

\begin{proof}[Outline]
Apply Atobe's discrete-series criterion in both directions and use \ref{B3.00e} to deduce cuspidality from lower-lift vanishing. Factor the chain Euler product, using absolute convergence for the bottom factor and \ref{B8.11h} for the character ratios. The complete proof is in \S\ref{prf:B8.08e}.
\end{proof}

\begin{lemma}[The column dimension at a linked member]\label{B8.08j}
At every good linked member in the range of the size conditions of Theorem~\ref{B8.08}, let \(u\) be any word, \(T=T_u\) its level prime, and \(d\) the dimension of its unrestricted column space. Use the pre-level factor \(d_u^\star\) and the row factor \(D^\star\) of Lemma~\ref{B8.05a}. Put
\[
 T_*=\left(\frac{D^\star}{d_u^\star\Gamma^2}\right)^{1/(fk)}.
\]
Then
\[
 d=d_u^\star(T^k-1)^f,\qquad
 \frac12\le\frac{d\Gamma^2}{D^\star}\le\exp(1)<3.
\]
The unrestricted column space is equivariantly isomorphic to the space of complex functions on a finite free \((\F_T^\times)^f\)-set. In these coordinates the ramified space is the common kernel of the sums over the individual factors.
\end{lemma}

\begin{proof}[Outline]
Uses \ref{B8.54} away from \(T\) to get \(d_u^\star\) from \ref{B8.05a}, then adjoins \((T^k-1)^f\) free \((\F_T^\times)^f\)-orbit choices at \(T\) to model the space and ramified kernel. The level window for \(T_*\) yields the ratio bounds. The complete proof is in \S\ref{prf:B8.08j}.
\end{proof}

\begin{lemma}[The collision sum from finite fibres]\label{B8.08k}
Fix one unrestricted word space at a member of Lemma~\ref{B8.08j}. Write
\[
 r=(m-2)/2,\quad g=g_0g_1,\quad v=\operatorname{rd}(E)/2,\quad
 L=\log b,\quad N=T^f,\quad R=d/D^\star,\quad B=kf\log k.
\]
Use the size conditions of Definition~\ref{B8.02c} with width \(3\) and \(C_\ell\ge40\). For terms \(t,t'\) in the same row use the orbital vectors \(v_t,v_{t'}\) and the finite-type factor \(\Gamma_1\) of (8.19). Let \(\mathcal B\) consist of the ordered pairs for which both Gram invariants are nonzero and the corank is at least two.

The following finite-fibre test suffices for
\[
 \Gamma_1\sum_{(t,t')\in\mathcal B}|\langle v_t,v_{t'}\rangle|
 \le g^{-kf/2}\frac{d^2}{D^\star}+e^{-B/2}d .
\]

For the test, index the rows and columns of a reduced Gram matrix \(Z\) from \(0\) to \(r\). Write \(p_l(Z)\) for its leading minor of order \(l\), \(\rho=2v\), and
\[
 z_{l,\tau}=\rho^{-2\sum_{j=1}^l(r+1-j)}|\tau(p_l(Z))|,
 \qquad
 S_a(Z)=\prod_{\tau}\prod_{l=1}^{r+1}
                 \max\{1,z_{l,\tau}b^{-am^2}\}.
\]
The product over \(\tau\) uses one embedding of each conjugate pair. Put
\[
 G_Z=g^{-mf}\prod_\tau(z_{r+1,\tau}^2-T^2z_{r,\tau}^2),
 \qquad n_Z=\left(\frac{\exp(1)\pi}{kv}\right)^{kf}(G_Z/N)^k .
\]
Use labels on represented support with \(G_Z>0\). For constants \(c\ge\max(1,C_T)\), \(C_h,C_2,C_a\ge0\), write
\[
 h_Z=b^{c^2m^3f}S_c(Z)^4,\quad S_Z=S_{C_h}(Z),\quad
 e_s=c(s+m)m+c(1+\log_2(2f)).
\]
For each entry position \(a,b'\), suppose the vectors
\((\rho^{a+b'-2r}\tau(Z_{ab'}))_\tau\) belong to an additive lattice on which each nonzero vector has product of absolute values at least one. Suppose the product of the maxima of one and the absolute values of these coordinates is at most \(h_Z\).

For \(2\le s\le k\), let \(\Xi_s(h)\) be \(M_sM_{k-s}/M_k\) for \(s\le k-2\), and \((ck)^{kf}(hN)^{ck}\) for \(s=k-1,k\), with \(M_j\) the standard genus masses of Lemma~\ref{B8.02bc}. Divide the pairs of corank \(s\) into the nondegenerate and degenerate classes of \eqref{eq:p8c1-loop-nd-deg}, indexed by \(i=1,2\), respectively. Set \(b_1(Z)=1\). Set \(b_2(Z)=T^{-f/2}K(Z)\), where \(K(Z)=1\) if \(p_r(Z)=0\), and otherwise
\[
 K(Z)=\prod_\tau\frac{1+c_\tau}{\sqrt{c_\tau^2-1}},
 \qquad c_\tau=\frac{|\tau(p_{r+1}(Z)/p_r(Z))|}{T}.
\]
Use this expression on its represented support \(c_\tau>1\).

Suppose every finite fibre in class \(i\) with labels \(Z,Z'\) has normalized absolute contribution at most \(d\,W_i A_i(Z)A_i(Z')\), where the contribution includes \(\Gamma_1\), and \(W_i,A_i\) are nonnegative. Require
\[
 \begin{split}
 W_i&\le \Xi_s(h_Zh_{Z'})(h_Zh_{Z'})^{e_s}(csf)^{6sf}
 \left((2/g)^{(s-i)f}
       \frac{\sqrt{G_ZG_{Z'}}}{N}b_i(Z)b_i(Z')\right)^k,\\
 n_Z A_i(Z)^2 b_i(Z)^{2k}
 &\le h_Z^{2cr^2}\Gamma S_Z^{\,a_i-2q/(r+1)}F_i ,
 \end{split}
\]
where
\[
 a_1=C_2m^2,\quad a_2=C_2(m^2+2m^4),\quad
 F_1=2R,\quad F_2=\exp(2C_2m^6 fL)(C_a/k)^{kf}.
\]
It suffices to have the amplitude inequality for the single integer
\[
 q=\left\lceil(r+1)
 \left(4ck+4e_k+4cr^2+\tfrac12a_2+12(r+1)^2\right)\right\rceil .
\]
For these hypotheses a sufficient size threshold is
\[
 \begin{split}
 A&=c^2+2\max(C_h-c,0),\\
 C&=\max\{c,\,4+A(6c+2),\,2C_2,\,C_a,\,408\exp(1)\pi\},\\
 C'&\ge\max\{8C_\Gamma,\,3C_2,\,16(\exp(1)C)^{32}\}.
 \end{split}
\]
In particular \(q\le(9c+3)km^2\).
\end{lemma}

\begin{proof}[Outline]
Bounds the label sum \(\sum_Z S_Z^{-12(r+1)^2}\) by disjoint-disk volume comparison on the entry lattices, then combines the fibre and amplitude inequalities via \ref{B8.10l}. Estimates \(\Xi_s\) using \ref{B8.48} and \ref{B8.02bc} and sums over coranks \(2\le s\le k\). The complete proof is in \S\ref{prf:B8.08k}.
\end{proof}

\begin{lemma}[The nonzero-Gram part from three pair sums]\label{B8.08l}
Fix a word \(u\) and write its Gram operator as a sum \(K=\sum_t\operatorname{op}(t)\) that converges absolutely in Hilbert--Schmidt norm. Let the orbital vectors satisfy \(\sum_t\|v_t\|<\infty\). Each term \(t\) has a row \(\operatorname{row}(t)\), a finite-type index \(\iota(\operatorname{row}(t))\) and a Gram invariant \(G(t)\). Each ordered pair \((t,t')\) of terms in the same row has a corank \(s(t,t')\). Let \(E_{\rm gen}\) and \(E_{\rm ng}\) be the generic and the nongeneric terms with \(G(t)\ne0\). For a set \(E_1\) of terms put \(\operatorname{part}(E_1)=\sum_{t\in E_1}\operatorname{op}(t)\) and \(w_i(E_1)=\sum_{t\in E_1:\,\iota(\operatorname{row}(t))=i}v_t\). For a set \(U\) of ordered pairs of terms put
\[
 P(U)=\Gamma_1\sum_{(t,t')\in U}\Re\langle v_t,v_{t'}\rangle,\qquad
 |P|(U)=\Gamma_1\sum_{(t,t')\in U}|\langle v_t,v_{t'}\rangle| ,
\]
where \(\Gamma_1>0\) is the finite-type factor. Assume the following.
\begin{enumerate}
\item[(i)] The terms of the identity cell have \(G(t)=0\) and sum to the identity. Thus \(K-I=X+Z\), with \(X=\operatorname{part}(E_{\rm gen})+\operatorname{part}(E_{\rm ng})\) and \(Z\) the sum over the remaining terms.
\item[(ii)] For \(E_1=E_{\rm gen}\) and \(E_1=E_{\rm ng}\), the inequality of Lemma~\ref{B8.10f} holds with some \(P_g\ge0\):
\(\|\operatorname{part}(E_1)\|_{\rm HS}^2\le P_g\Gamma_1\sum_i\|w_i(E_1)\|^2\).
\item[(iii)] \(\langle v_t,v_{t'}\rangle=0\) whenever \(\operatorname{row}(t)\ne\operatorname{row}(t')\) and \(\iota(\operatorname{row}(t))=\iota(\operatorname{row}(t'))\).
\item[(iv)] \(s(t,t')\ge1\) for terms \(t,t'\) in the same row with \(G(t)\ne0\) and \(G(t')\ne0\).
\end{enumerate}
Let \(\mathcal B\) be the set of same-row ordered pairs with both Gram invariants nonzero and corank at least two, as in Lemma~\ref{B8.08k}. Let \(\mathcal U_{\rm gen}\) and \(\mathcal U_{\rm ng}\) be the same-row ordered pairs of corank one in \(E_{\rm gen}\) and in \(E_{\rm ng}\), respectively. If
\[
 |P|(\mathcal B)\le c_{\rm col},\qquad P(\mathcal U_{\rm gen})\le c_{\rm gen},\qquad P(\mathcal U_{\rm ng})\le c_{\rm ng},
\]
then
\[
 \|X\|_{\rm HS}^2\le2P_g(c_{\rm col}+c_{\rm gen}+c_{\rm ng}).
\]
Suppose further that \(P_g=(1+3/g_0)^f(1+3/g_1)^f\), with \((kb)^8\le g_0,g_1\), \(b\ge2\) and \(1\le f\le k\). Let \(d,D^\star,Y,B\) be as in Lemma~\ref{B8.08h}, with \(Y\ge1\), and let \(C_{\rm gen},C_{\rm ng}\ge0\). Suppose the three bounds above hold with
\[
 c_{\rm col}=(g_0g_1)^{-kf/2}\frac{d^2}{D^\star}+e^{-B/2}d,\qquad
 c_{\rm gen}=e^{C_{\rm gen}Y}\frac{d^2}{D^\star},\qquad
 c_{\rm ng}=e^{C_{\rm ng}Y}\frac{d^2}{D^\star}.
\]
Put \(C_1=\max(C_{\rm gen},C_{\rm ng})+1\) and \(a=C_1Y\). Then \(C_1\ge1\), and the first equal-word estimate of Lemma~\ref{B8.08h} holds:
\[
 \|X\|_{\rm HS}^2\le\frac{12}{5}\left((1+e^{a})\frac{d^2}{D^\star}+e^{-B/2}d\right).
\]
\end{lemma}

\begin{proof}[Outline]
Expands \(\sum_i\|w_i(E_1)\|^2\) into same-row pairs by (iii), splits by corank via (iv), and applies (ii) with \(\|x+y\|^2\le2\|x\|^2+2\|y\|^2\) and disjointness of \(E_{\rm gen},E_{\rm ng}\). Uses \(1+x\le e^x\) and \(2\le e^Y\) for the numerical bound. The complete proof is in \S\ref{prf:B8.08l}.
\end{proof}

\begin{lemma}[Orthogonality of orbital rows]\label{B8.08m}
Fix a tied equal-word model and a word \(u\). Let \(\mathcal R_u\) be the set of ordered pairs of row lattices. The orbital Hilbert space is the Hilbert direct sum of the spaces of normalized orbital operators indexed by \(\mathcal R_u\). For a term \(t\), write \(\rho_t=(\Lambda_t,\Lambda'_t)\) for its row, \(A_t\) for the Hilbert--Schmidt image of its normalized orbital amplitude, and \(w_t\in\R\) for its row weight. Its orbital vector is
\[
 v_t(\rho)=
 \begin{cases}
 w_tA_t,&\rho=\rho_t,\\
 0,&\rho\ne\rho_t.
 \end{cases}
\]
For terms in distinct rows,
\[
 \rho_t\ne\rho_{t'}\quad\Longrightarrow\quad\langle v_t,v_{t'}\rangle=0 .
\]
Consequently hypothesis~\textup{(iii)} of Lemma~\ref{B8.08l} holds for this model.
\end{lemma}
\begin{proof}
The inner product in the Hilbert direct sum is the sum of the inner products in its row coordinates. The supports of \(v_t\) and \(v_{t'}\) are contained in \(\{\rho_t\}\) and \(\{\rho_{t'}\}\), respectively. If these rows differ, at least one of the two vectors vanishes in each coordinate. Every coordinate inner product is therefore zero.
\end{proof}

\begin{lemma}[Full corank-one classes and their amplitudes]\label{B8.08n}
Use the term vectors, row map, corank function and signed pair sum of Lemma~\ref{B8.08l}, and assume \(\sum_t\|v_t\|<\infty\) and \(\Gamma_1>0\). Let \(E_{\rm gen}\) be the set of generic terms with nonzero Gram invariant. Suppose that
\[
 t\sim t'\quad\Longleftrightarrow\quad\operatorname{row}(t)=\operatorname{row}(t')\ \text{ and }\ s(t,t')=1
\]
is an equivalence relation on the terms with nonzero Gram invariant. Let \(\mathcal C\) be the set of equivalence classes of its restriction to \(E_{\rm gen}\), and let \(\mathcal U_{\rm gen}\) consist of all same-row ordered pairs in \(E_{\rm gen}\) of corank one.

Let \(\mathcal P\) and \(\mathcal N\) be sets of point labels and normal labels, with maps \(\gamma:E_{\rm gen}\to\mathcal P\) and \(\nu:E_{\rm gen}\to\mathcal N\). For \(c\in\mathcal C\), \(\gamma_0\in\mathcal P\) and \(\nu_0\in\mathcal N\) put
\[
 V(c,\gamma_0,\nu_0)=\sum_{\substack{t\in c\\ \gamma(t)=\gamma_0,\ \nu(t)=\nu_0}}v_t .
\]
Let \(J_2:\mathcal P\to[0,\infty)\) satisfy \(J_2(\gamma(t))=0\Rightarrow v_t=0\) for \(t\in E_{\rm gen}\). Define the normalized squared amplitude by
\[
 a_2(c,\gamma_0,\nu_0)=
 \begin{cases}
 \|V(c,\gamma_0,\nu_0)\|^2/J_2(\gamma_0),&J_2(\gamma_0)>0,\\
 0,&J_2(\gamma_0)=0.
 \end{cases}
\]
Then \(a_2\ge0\), and
\[
 P(\mathcal U_{\rm gen})=\Gamma_1\sum_{c\in\mathcal C}\Bigl\|\sum_{\gamma_0\in\mathcal P,\ \nu_0\in\mathcal N}V(c,\gamma_0,\nu_0)\Bigr\|^2,
 \qquad
 \|V(c,\gamma_0,\nu_0)\|^2=a_2(c,\gamma_0,\nu_0)\,J_2(\gamma_0).
\]
The represented point set \(\mathcal P_*=\{\gamma(t):t\in E_{\rm gen},\ v_t\ne0\}\) consists precisely of the point labels represented by a nonzero term in one of these class and normal fibres.

For the finite amplitude estimate, fix \(\gamma_0\in\mathcal P\) and a finite set \(S\subset\mathcal C\times\mathcal N\). Let \(r,c_1,c_0\) be nonnegative integers and put \(m=2r+2\). Let \(\mathcal V\) be a finite index set and \(\Omega\subset\mathcal V\). For each \(p\in\mathcal V\) let \(D_p\) be a finite set, let \(w_p:D_p\to[0,\infty)\) and \(a_p:D_p\to\C\), and let \(H_p,n_p,n_p^{\rm gen}\) be real numbers with \(H_p\ge1\), \(n_p,n_p^{\rm gen}\ge0\), and \(H_p\ge2\) for \(p\in\Omega\). Assume
\[
 \sum_{d\in D_p}w_p(d)|a_p(d)|^2\le H_p^{c_1(r+1)^2}n_p,
 \qquad
 \begin{cases}
 n_p\le8H_p^{c_0(r+1)^2+2m}n_p^{\rm gen},&p\in\Omega,\\
 n_p=n_p^{\rm gen},&p\notin\Omega.
 \end{cases}
\]
Suppose there are nonnegative weights \(\rho_c\), a number \(V_0\ge0\), a subset \(S_0\subset S\) and a map \(\mathbf d:S_0\to\prod_{p\in\mathcal V}D_p\) such that, for \((c,\nu_0)\in S\),
\[
 a_2(c,\gamma_0,\nu_0)=
 \begin{cases}
 \rho_c\Bigl|\prod_{p\in\mathcal V}a_p\bigl(\mathbf d(c,\nu_0)_p\bigr)\Bigr|^2,&(c,\nu_0)\in S_0,\\
 0,&(c,\nu_0)\notin S_0,
 \end{cases}
\]
and, for every tuple \(\mathbf e\in\prod_{p\in\mathcal V}D_p\),
\[
 \sum_{\substack{(c,\nu_0)\in S_0\\ \mathbf d(c,\nu_0)=\mathbf e}}\rho_c\le V_0\prod_{p\in\mathcal V}w_p(e_p).
\]
Put \(H_{\rm tot}=\prod_{p\in\mathcal V}H_p\), \(M_{\rm gen}=V_0\prod_{p\in\mathcal V}n_p^{\rm gen}\) and \(e_7=(c_1+c_0)(r+1)^2+2m+5\). Then
\[
 \sum_{(c,\nu_0)\in S}a_2(c,\gamma_0,\nu_0)\le H_{\rm tot}^{e_7}M_{\rm gen}.
\]
\end{lemma}
\begin{proof}[Outline]
Partitions \(\mathcal U_{\rm gen}\) into class squares, regrouping the pair sum by Cauchy--Schwarz to express \(P(\mathcal U_{\rm gen})\) via fibre norms. Groups \(S_0\) by tuples \(\mathbf e\) and combines the fibre and local weighted bounds to bound \(\sum a_2\). The complete proof is in \S\ref{prf:B8.08n}.
\end{proof}

\begin{lemma}[The unitary dimension and mass bound]\label{B8.08oa}
Fix a tied equal-word model of Theorem~\ref{B8.10} and a word \(u\). Let \(E/F\) be its CM extension, let \(g_0,g_1\) be its two auxiliary primes, assumed to split completely in \(E\), put \(g=g_0g_1\), assume \((T_u,g)=1\), and let \(\mathcal V_{\mathrm{typ}}\) be the finite set of places of \(F\) dividing \(g\), with residue fields \(\F_{q_v}\) for \(v\in\mathcal V_{\mathrm{typ}}\). Let \(k\) be the initial rank, put
\[
 \mathcal I=\prod_{v\in\mathcal V_{\mathrm{typ}}}\GL_k(\F_{q_v}),
\]
and let \(\mathcal U\le\mathcal I\) be the subgroup of tuples that are upper unitriangular at every \(v\in\mathcal V_{\mathrm{typ}}\). Let \(\Gamma_1\) be the finite-type factor of Lemma~\ref{B8.10f}, namely the product over \(v\in\mathcal V_{\mathrm{typ}}\) of the values at \(1\) of the normalized induced Whittaker characters on \(\GL_k(\F_{q_v})\). Let \(C_u\) be the complex column space of \(u\), let \(\pi_u\) be its archimedean coefficient representation, and put
\[
 d=\dim_\C C_u,\qquad n_\pi=\dim_\C\pi_u.
\]
Let \((V,h_V)\) be the Hermitian \(E\)-space of the model (with \(h_V\), and its adelic extension \(h_A\) below, conjugate-linear in the first variable and linear in the second), let \(F_u\) be the fixed flag of \(u\) with top isotropic subspace \(X_u\subset V\) of dimension \(r\) and rational flag stabilizer \(\operatorname{Stab}(F_u)\) (fixing the flag basis of \(X_u\) and the final vector modulo \(X_u\)), and let \(\mathcal R_0(u)\) be the model's finite set of representatives for the \(\operatorname{Stab}(F_u)\)-orbits of row lattices \(\Lambda\subset V\). For \(\Lambda\in\mathcal R_0(u)\), let \(\operatorname{Aut}_{F_u}(\Lambda)\) be the group of \(E\)-linear automorphisms of \(X_u^\perp/X_u\) induced by elements of \(\operatorname{Stab}(F_u)\) preserving \(\Lambda\), set \(w(\Lambda)=|\operatorname{Aut}_{F_u}(\Lambda)|^{-1}\) when \(\operatorname{Aut}_{F_u}(\Lambda)\) is finite and \(w(\Lambda)=0\) otherwise, and put
\[
 \mu_u=\operatorname{mass}_u=\sum_{\Lambda\in\mathcal R_0(u)}w(\Lambda).
\]
Then \(\Gamma_1=|\mathcal I/\mathcal U|\ge1\); in particular \(\Gamma_1>0\) and \(\Gamma_1\ge0\), and
\[
 \Gamma_1\mu_u n_\pi\le d.
\]
\end{lemma}

\begin{proof}[Outline]
Evaluates \ref{B8.10f} at \(1\) for \(\Gamma_1=|\mathcal I/\mathcal U|\). Lifts \(\mathcal I/\mathcal U\) into \(H_{\mathrm{int}}(u)\) via split-place frames, bounds each double-coset fibre by \(|\operatorname{Aut}_{F_u}(\Lambda)|\), and sums over \(\Lambda\in\mathcal R_0(u)\). The complete proof is in \S\ref{prf:B8.08oa}.
\end{proof}

\begin{lemma}[Vanishing of orbital vectors on nonrelevant terms]\label{B8.08ob}
Fix a tied equal-word model of Theorem~\ref{B8.10} and a word \(u\). Let \(E/F\) be its CM extension. For a term \(t\) of word \(u\), let \(F_u\) be the fixed flag of \(u\) and \(F'_t\) the transported flag of \(t\), let \(\mathfrak n(F_u)\) and \(\mathfrak n(F'_t)\) be their rational unipotent radical Lie algebras with root-sum functionals \(\ell_{F_u}\) and \(\ell_{F'_t}\) into \(E\), and call \(t\) relevant when \(\operatorname{Tr}_{E/F}\circ\ell_{F_u}\) and \(\operatorname{Tr}_{E/F}\circ\ell_{F'_t}\) agree on \(\mathfrak n(F_u)\cap\mathfrak n(F'_t)\). Let \(v_t\) be the orbital vector of \(t\) as in Lemma~\ref{B8.08m}. If \(t\) is not relevant, then \(v_t=0\). In particular, if \(\mathcal T_0\) is the set of nonidentity terms of word \(u\) with zero Gram invariant and \(\mathcal T_{\mathrm{rel}}\subseteq\mathcal T_0\) is its subset of relevant terms, then
\[
 v_t=0\qquad(t\in\mathcal T_0\setminus\mathcal T_{\mathrm{rel}}).
\]
\end{lemma}

\begin{proof}[Outline]
Writes \(v_t\) via \ref{B8.08m} as an orbital integral and shows \(\Gamma_{\mathrm{orb}}(t)\backslash R_t\) is compact with \(\chi_t\) trivial on \(\Gamma_{\mathrm{orb}}(t)\). For nonrelevant \(t\), exponentiation gives \(\chi_t(r_0)=-1\), so averaging kills \(v_t\). The complete proof is in \S\ref{prf:B8.08ob}.
\end{proof}

\begin{proposition}[The norm expansion on the three nonidentity cells]\label{B8.08oc}
Fix a tied equal-word model of Theorem~\ref{B8.10} (with \(g_0,g_1\) split completely in \(E\) and \((T_u,g)=1\)) and a word \(u\). Let \(C_u\) be the complex column space of \(u\), and let the orbital Hilbert space and inner product \(\langle\cdot,\cdot\rangle\) (conjugate-linear in the first variable and linear in the second) be those of Lemma~\ref{B8.08m}. Let \(\mathcal I=\prod_{v\in\mathcal V_{\mathrm{typ}}}\GL_k(\F_{q_v})\) and \(\mathcal U\le\mathcal I\) be as in Lemma~\ref{B8.08oa}. At each \(v\in\mathcal V_{\mathrm{typ}}\) with selected additive character \(\psi_v\) of \(\F_{q_v}\), let \(\psi_{\mathrm{nat},v}:\GL_k(\F_{q_v})\to\C\) equal \(\psi_v\bigl(\sum_{a=1}^{k-1}x_{a,a+1}\bigr)\) for upper unitriangular \(x\) and \(0\) otherwise, put \(\psi_{\mathcal I}(x)=\prod_{v\in\mathcal V_{\mathrm{typ}}}\psi_{\mathrm{nat},v}(x_v)\) for \(x\in\mathcal I\), and define the finite-type Gram matrix \(\mathcal G=(\mathcal G_{ij})_{i,j\in\mathcal I}\) by
\[
 \mathcal G_{ij}=\overline{\gamma_{\mathrm{fin}}(j^{-1}i)},\qquad
 \gamma_{\mathrm{fin}}(x)=|\mathcal U|^{-1}\sum_{y\in\mathcal I}\psi_{\mathcal I}(yxy^{-1}).
\]
For each term \(t\) of word \(u\), let \(O_t\) be its Hilbert--Schmidt operator on \(C_u\), let \(v_t\) be its orbital vector, and let \(\iota_t\in\mathcal I\) be its finite-type index. Assume
\[
 \sum_t\|O_t\|_{\mathrm{HS}}<\infty,\qquad
 \sum_t\|v_t\|<\infty.
\]
Let \(E_{\mathrm{gen}}\) and \(E_{\mathrm{ng}}\) be the sets of generic and nongeneric terms with nonzero Gram invariant, and let \(\mathcal T_0\) be the set of nonidentity terms with zero Gram invariant. For \(S\in\{E_{\mathrm{gen}},E_{\mathrm{ng}},\mathcal T_0\}\), put
\[
 \operatorname{part}(S)=\sum_{t\in S}O_t,\qquad
 w_i(S)=\sum_{\substack{t\in S\\\iota_t=i}}v_t\quad(i\in\mathcal I).
\]
Then
\[
 \|\operatorname{part}(S)\|_{\mathrm{HS}}^2
 =\Re\sum_{i,j\in\mathcal I}\mathcal G_{ij}\bigl\langle w_i(S),w_j(S)\bigr\rangle.
\]
In particular, for \(Z=\operatorname{part}(\mathcal T_0)\) and \(w_i=w_i(\mathcal T_0)\),
\[
 \|Z\|_{\mathrm{HS}}^2
 =\Re\sum_{i,j\in\mathcal I}\mathcal G_{ij}\langle w_i,w_j\rangle.
\]
\end{proposition}

\begin{proof}[Outline]
Uses block Parseval and bijections \(\Phi_\Lambda\) built via \ref{B8.08oa} to lift \(\|\operatorname{part}(S)\|_{\mathrm{HS}}^2\) to \(\mathcal I^2\) by stabilizer invariance. Summing the Whittaker phases over \(\mathcal I^2\) yields the Gram matrix \(\mathcal G_{ij}\). The complete proof is in \S\ref{prf:B8.08oc}.
\end{proof}

\begin{proposition}[Summation of the zero-Gram pair amplitudes]\label{B8.08o}
Fix \(C_H,C_h,C_p\ge0\) and \(C\ge2(C_H+1)\), and put \(C_s=C_p+4\). There is a constant \(C_{\ell,\mathrm{pair}}\ge16\), depending only on these constants, with the following property.

Let a tied equal-word model satisfy the source conditions of Definition~\ref{B8.02c}, the unramified CM hypothesis, and the parameter conditions of Proposition~\ref{B8.02b}, with \(C_\ell\ge C_{\ell,\mathrm{pair}}\). Assume the fixed-type estimate of Proposition~\ref{B8.02b} for the chosen constants. Write
\[
 r=(m-2)/2,\qquad L=\log b,\qquad \rho=\operatorname{rd}(E),\qquad H(c)=b^{cm^2f}.
\]
In particular, \(m=2r+2\), \(k\ge m+3\), \(L>1\), \(1\le f\le k\), and
\[
 C_\ell m^6L\le k\log k,\qquad
 C_\ell\log\rho\le L,\qquad (kb)^8\le g_0,g_1.
\]
Fix a word \(u\). Let \(C_u\) be its complex column space, let \(\pi_u\) be its archimedean coefficient representation, and put
\[
 d=\dim_\C C_u,\qquad n_\pi=\dim_\C\pi_u,\qquad
 \mu_u=\operatorname{mass}_u.
\]
Here \(\mu_u\) is the sum of the inverse row-lattice stabilizer orders, with the normalization of Proposition~\ref{B8.02b}. Let \(\Gamma_1\) be the finite-type factor of Lemma~\ref{B8.08l}, and put
\[
 P_g=(1+3/g_0)^f(1+3/g_1)^f.
\]
Let \(\mathcal T_0\) be the set of terms with zero Gram invariant outside the identity cell, and let \(\mathcal T_{\mathrm{rel}}\) be its relevant terms. For a term \(t\), write
\[
 \operatorname{row}(t)=(\Lambda_t,\Lambda'_t,\iota_t)
\]
for its complete row label, including its finite-type index \(\iota_t\). Let \(v_t\) be its orbital vector, and assume
\begin{equation}\label{eq:p8c1-zero-normalization}
 \Gamma_1\ge0,\qquad
 \Gamma_1\mu_u n_\pi\le d,\qquad
 v_t=0\quad(t\in\mathcal T_0\setminus\mathcal T_{\mathrm{rel}}).
\end{equation}
Define the set of ordered pairs
\[
 \mathcal Q_0=\{(t,t')\in\mathcal T_0^2:
                         \operatorname{row}(t)=\operatorname{row}(t')\}.
\]
Then the family \((\|v_t\|\|v_{t'}\|)_{(t,t')\in\mathcal Q_0}\) is summable, and
\begin{equation}\label{eq:p8c1-zero-pair-amplitudes}
 P_g\Gamma_1\sum_{(t,t')\in\mathcal Q_0}\|v_t\|\|v_{t'}\|
 \le d\exp(-kf\log k).
\end{equation}
\end{proposition}

\begin{proof}[Outline]
Sums \(S_\delta^{-3r}\) over rank-\(j\) labels via minor integrality (\ref{B8.02b}) and dyadic disk packing in \(\mathcal O_E\). Combines \ref{B8.02b}, \ref{B8.02c}, \(C_{\ell,\mathrm{pair}}\), and \(P_g\le2\) to bound the pair sum by \(d\exp(-kf\log k)\). The complete proof is in \S\ref{prf:B8.08o}.
\end{proof}

\begin{lemma}[From pair amplitudes to the zero-Gram operator]\label{B8.08p}
Fix a tied equal-word model satisfying the selected local hypotheses of Theorem~\ref{B8.10} (as used in Lemma~\ref{B8.10f}), and fix a word \(u\). Use the notation \(k,f,\mathcal T_0,\mathcal Q_0,v_t,\operatorname{row}(t),\iota_t,d,P_g,\Gamma_1\) of Proposition~\ref{B8.08o}. Let \(O_t\) be the Hilbert--Schmidt operator of a term \(t\), and assume
\[
 \sum_t\|O_t\|_{\mathrm{HS}}<\infty,\qquad
 \sum_t\|v_t\|<\infty,\qquad \Gamma_1\ge0.
\]
Let \(\mathcal I\) be the finite set of finite-type indices, and put
\[
 Z=\sum_{t\in\mathcal T_0}O_t,\qquad
 w_i=\sum_{\substack{t\in\mathcal T_0\\\iota_t=i}}v_t
 \quad(i\in\mathcal I).
\]
Let \(\mathcal G=(\mathcal G_{ij})_{i,j\in\mathcal I}\) be the finite-type Gram matrix. Assume the norm expansion
\begin{equation}\label{eq:p8c1-zero-expansion}
 \|Z\|_{\mathrm{HS}}^2
 =\Re\sum_{i,j\in\mathcal I}\mathcal G_{ij}\langle w_i,w_j\rangle.
\end{equation}
If \eqref{eq:p8c1-zero-pair-amplitudes} holds, then
\[
 \sum_{t\in\mathcal T_0}\|O_t\|_{\mathrm{HS}}<\infty,
 \qquad
 \|Z\|_{\mathrm{HS}}^2\le d e^{-kf\log k}.
\]
\end{lemma}

\begin{proof}[Outline]
Applies \(\mathcal G\preceq P_g\Gamma_1 I\) from \ref{B8.10f} to bound \(\|Z\|_{\mathrm{HS}}^2\le P_g\Gamma_1\sum_i\|w_i\|^2\). Expands \(\|w_i\|^2\) via row orthogonality from \ref{B8.08m} and Cauchy--Schwarz to reduce to the pair-amplitude bound. The complete proof is in \S\ref{prf:B8.08p}.
\end{proof}

\begin{theorem}[Uniform thresholds for the zero-Gram estimate]\label{B8.08q}
Fix \(C_g,c_b,C_w,c_T,C_T,C_0\in\R\), with \(c_b,c_T>0\). There exist thresholds \(C'_{\mathrm{req}},C_{\ell,\mathrm{req}}\) with the following property. For every \(C'\ge C'_{\mathrm{req}}\) and \(C_\ell\ge C_{\ell,\mathrm{req}}\), use a tied model of Proposition~\ref{B8.02b} satisfying the source conditions of Definition~\ref{B8.02c} and the unramified CM hypothesis. Fix any word \(u\), and use the term operators, vectors, and column space of Lemma~\ref{B8.08p}. Suppose that the absolute-convergence hypotheses of that lemma hold. Then, with
\[
 Z=\sum_{\substack{t\text{ outside the identity cell}\\G(t)=0}}O_t,
 \qquad d=\dim_\C C_u,
\]
one has
\begin{equation}\label{eq:p8c1-zero-gram-bound}
 \sum_{\substack{t\text{ outside the identity cell}\\G(t)=0}}
                  \|O_t\|_{\mathrm{HS}}<\infty,
 \qquad
 \|Z\|_{\mathrm{HS}}^2\le d e^{-kf\log k}.
\end{equation}
The thresholds depend on the fixed constants and are uniform in the datum, the tied model, and the word.
\end{theorem}

\begin{proof}[Outline]
Takes the maximum of thresholds from \ref{B8.02b} and \ref{B8.08o} and checks \ref{B8.02b} from \ref{B8.02c}. Combines \ref{B8.08oa}, \ref{B8.08ob}, \ref{B8.08o}, \ref{B8.08oc}, and \ref{B8.10f}, then concludes via \ref{B8.08p}. The complete proof is in \S\ref{prf:B8.08q}.
\end{proof}

\begin{lemma}[The numerical test on an unrestricted word space]\label{B8.08h}
Let \(C\) be the unrestricted complex column space for one word, \(d=\dim C\), and \(K=\alpha^{-1}P^*P\) its Gram operator, where \(P\) is the period map and \(\alpha>0\) the unfolding constant. Let \(k\ge5\) be the initial rank, \(m\ge2\) the chain length, \(2\le f\le k\) the degree of the real field, and \(L\ge1\) the logarithmic scale. Put
\[
       Y=m^5fL,\qquad a=C_1Y,\qquad B=kf\log k ,
\]
where \(C_1\ge1\). Suppose \(K-I=X+Z\), with \(X\) the sum over the nonzero Gram invariant after removing the identity cell, and \(Z\) the remaining nonidentity sum. Suppose the two equal-word estimates, in the normalization with main term \(D^\star>0\), are
\[
 \|X\|_{\rm HS}^2\le \frac{12}{5}
       \left((1+e^a)\frac{d^2}{D^\star}+e^{-B/2}d\right),
          \qquad \|Z\|_{\rm HS}^2\le e^{-B}d .
\]
Let \(d^\star>0\) be the pre-level dimension factor, let
\(\Gamma=\exp(C_\Gamma Y)\), \(C_\Gamma\ge C_1+1\), and set
\(T_*=(D^\star/(d^\star\Gamma^2))^{1/(fk)}\).
Assume the level prime \(T\) satisfies
\[
 d=d^\star(T^k-1)^f,\qquad
 T_*\le T\le(1+1/(fk))T_*,\qquad
 T\equiv1\pmod{(g_0g_1)^2},\qquad g_0\ge k^8 ,
\]
where \(g_0,g_1\) are the two type primes. Suppose the level model identifies \(C\) with functions on a finite free \((\F_T^\times)^f\)-set, with each diamond operator acting by translation in its coordinate. Let \(U\) be the common kernel of the \(f\) coordinate sums. Then
\[
 \|K-I\|_{\rm HS}^2\le e d,\qquad
 \dim U\ge(1-\delta)d,\qquad 2e<1-\delta ,
 \qquad
 e=8e^{-a}+6e^{-B/2},\quad \delta=\frac{f}{T-1}.
\]
\end{lemma}

\begin{proof}[Outline]
Use the level window to bound \(d/D^\star\) and combine the two Hilbert-Schmidt estimates. Count characters nontrivial in every coordinate on each free orbit, then use the level congruence and size bounds to show \(2e<1-\delta\). The complete proof is in \S\ref{prf:B8.08h}.
\end{proof}

\begin{lemma}[A nonzero bottom pairing]\label{B8.08f}
Let \(C_u\) be the complex column space of a word \(u\), \(d_0=\dim C_u>0\), \(P_u:C_u\to H\) its period map into the spectral Hilbert space, and \(\alpha_u>0\) its unfolding constant. Put \(K_u=\alpha_u^{-1}P_u^*P_u\). Let \(U_u^{\rm ram}\subset C_u\) be the ramified subspace in the column coordinates. Suppose, for real \(e,\delta\),
\[
       \|K_u-1\|_{\rm HS}^2\le e d_0,\qquad
       \dim U_u^{\rm ram}\ge(1-\delta)d_0,\qquad e<1-\delta .
\]
Then some \(\phi\in U_u^{\rm ram}\) satisfies \(P_u\phi\ne0\). At a linked realization with the full reference level of Lemma~\ref{B8.08d}, suppose the column identification intertwines the unrestricted period maps. Then a full reference multiplicity map has nonzero translated pairing with a model column representing \(\phi\).
\end{lemma}

\begin{proof}[Outline]
Bound the period kernel dimension by \(\|K_u-1\|_{\rm HS}^2\) to find a ramified column with nonzero period. Surjective full-level evaluation from \ref{B8.08d} and invertible central translation realize that period by a reference row. The complete proof is in \S\ref{prf:B8.08f}.
\end{proof}

\begin{lemma}[Comparison with a full reference row]\label{B8.08i}
Use the row and column comparisons at a linked realization, with the common finite translate and level measures of Lemma~\ref{B8.08a}. Fix a word \(u\). Let \(T_0\) be a full level-zero reference multiplicity map, with a preimage under descent on the reference chain. Suppose that its translated model pairing with a permitted column is \(b_0\ne0\). For a tuple \(t\), a column \(\phi\) in the complex realization of the permitted space and a rational row \(F\), let \(A_t(\phi,F)\) be the arithmetic top period in the complex model. Use the column change from the chain and the bottom column isometry to form the translated model pairing \(B_t(\phi,F)\) with the descended row. Before dividing the column change by its norm multiplier, there is a \(p\ne0\), independent of \(t,\phi,F\), such that
\[
                         c_t A_t(\phi,F)b_0=p B_t(\phi,F).
\]
Here \(c_t\) is the unnormalized scalar of Lemma~\ref{B8.08a}.
\end{lemma}

\begin{proof}[Outline]
Unfold each double theta step to descend relations among weighted periods, using cusp decay as in \ref{B8.15f}. Transfer the bottom relation to translated pairings and use the nonzero reference pairing to obtain the comparison scalar. The complete proof is in \S\ref{prf:B8.08i}.
\end{proof}

\begin{theorem}[The twelve open conditions]\label{B8.08}
Assume \(P\). For every choice \(h_L\) of weak local correspondences as in Theorem~\ref{B9.46} and \(c_{\mathrm{Dep}}\in\R\), the parameters \(p_0(h_L,c_{\mathrm{Dep}})\) satisfy the following linked rows.
\[
\begin{array}{c|p{0.76\textwidth}}
 \text{Primitive rows} & Primitive rows and separated types.\\
 \text{Finite types} & Finite types assigned to each row index and layer. The isotypic spaces form a direct sum equal to the transported row space and have prescribed \(\lambda\) on nonzero components at used tuples. The assigned counts sum to the base row dimension.\\
 \text{Divisor comparison} & Both \(RR\) lattices agree with their indefinite realizations, and rational lifts span at used tuples. The theta maps, numerical data and column pairings agree.\\
 \text{Definite evaluation} & At used tuples the embedded definite evaluation equals the arithmetic column period on transported inputs.\\
 \text{Join maps} & The join maps and identities of \ref{B8.02} at every selected point.\\
 \text{Recursion} & Compatible indefinite recursion at every point.\\
 \text{Label realization} & Choose a Hida realization \(X\), a good operator and an isometry from spectral to model space. The automorphic datum has the prescribed rank, degree, root discriminant, signatures, Harish--Chandra parameters, auxiliary levels and dimension. Its archimedean representations have the specified discrete series and lowest unit vectors, and its formal-degree and volume expression equals \(D^\star\). There is a complex-linear isometry from the arithmetic reference row space onto the complex conjugate of the full good subspace cut out by the operator. After the ambient identification, conjugating this isometry gives the spectral row embedding. At each type prime the assigned local polynomials are valid labels, each row label tuple is good in the sense of Definition~\ref{def:p8c1-good-labels}, and the realized row belongs to its labelled type space.\\
 \text{Good labels} & Eventually at every selected point, every finite-type row satisfies Definition~\ref{def:p8c1-good-labels} at both type primes, for the prescribed tame characters and their inverses.\\
 \text{Geometry} & Point summands with class-sum pinning, recursion ties and an evaluation grid.\\
 \text{Multiplicity} & Eventually, the multiplicity identity \eqref{eq:p8c1-multiplicity} with the unit conditions below.\\
 \text{Infinite-place factors} & Eventually, for each \(e:K_D\hookrightarrow\C\) at its infinite place \(v_e\), the factor of the division datum at \(e\) equals the selected row-side factor \(\kappa(v_e)\).\\
 \text{Definite side} & Eventual definite-side condition of \ref{B8.35} for each \(C\ge C_0\).\\
\end{array} \tag{8.16}
\]
For the multiplicity identity, use the definite face and complex realization of the selected bundle \(\beta\). Write \(K\) for the coefficient field of this face and \(\iota:K\hookrightarrow\C\) for its embedding. Let \(u\) range over its column indices and \(t\) over its tuples. Let \(\phi\in K^{a_u}\) be a column coordinate vector and \(F\in K^{n_t}\) a row coordinate vector, with \(a_u,n_t\) their respective dimensions. Write \(e_t\) for the rational row realization, \(j_t\) for its map to the complex realization, \(\Theta_t\) for the theta map, \(s_{EE}\) for the isometry on its output space, and \(P_u\) for the spectral column map. There are a complex-linear map \(V\), column matrices \(U_u\), and scalars \(a_u'\ne0\), \(\zeta_{u,t}\ne0\) such that for every \(u,t,\phi,F\),
\begin{equation}\tag{8.16a}\label{eq:p8c1-multiplicity}
 \zeta_{u,t}\,\iota\bigl(\mathrm{Pt}_{u,t}(\phi,F)\bigr)
 =a_u'\left\langle
 s_{EE}\bigl(V\Theta_t j_t e_t F\bigr),\,
 P_u\bigl(U_u\iota(\phi)\bigr)
 \right\rangle .
\end{equation}
Here \(\mathrm{Pt}_{u,t}\) is the definite evaluation pairing and \(\iota(\phi)\) is taken coordinatewise. The map \(V\) preserves the norm on the good theta subspace and sends it into itself. Each \(U_u\) preserves the prescribed column norm, and \(\zeta_{u,\mu_0}=1\). At each finite coefficient place \(w\), any \(z\in K\) with \(\iota(z)=\zeta_{u,\mu_w}\) has \(v_w(z)=0\).

The good-label, multiplicity, infinite-place and definite-side conditions are closed on \(W_3\). For the definite-side condition, choose \(C_0\ge0\) per eligible \(F_m,CH\), then a cutoff per \(C\ge C_0\). For extensions with the same \(h_L,S,\Gamma\), a proof, for every choice in \(N'\), that a choice in \(N\) exists at the same point indices transfers each linked row with conclusion \(H(\mathcal A,\sigma)\) from \(N\) to \(N'\).
\end{theorem}

\begin{proof}[Outline]
Combine the reference pairing of \ref{B8.08f} with descent in \ref{B8.08i} and scalar units from \ref{B8.08a}. Use \ref{B8.08d} for the full good subspace and \ref{B7.37} and \ref{B8.37} for the finite-type and definite-side rows. The complete proof is in \S\ref{prf:B8.08}.
\end{proof}

\begin{lemma}[The lower weight before the choice of a family]\label{B8.08b}
Assume \(P\) of Definition~\ref{B1.15ac}. On the class \(\mathcal F_{\rm cov}\) of Definition~\ref{B8.24}, the following assertion is closed in the sense of Definition~\ref{B8.39}. For each eligible \(F_m,C_H\) there is \(x_0\) such that, whenever \(x>x_1\) and \(x\ge x_0\), the following holds for every \(i\in I_x\) and every point in \(\mathcal P\) of Definition~\ref{def:B8-reference-points} over \((x,i)\). For every \(C\in\R\), every layer \(1\le j\le m\), every primitive block \(a\) and every infinite coefficient place \(v\),
\[
 L_{6,j}\le\log\lambda^{(1)}_{j,\mu_0,a,v}.
\]
Here the superscript denotes level one, and \(L_{6,j}\) is the sixth numerical entry at \(C,C_H\) in the notation of \ref{B5.07}.
\end{lemma}

\begin{proof}[Outline]
Check the size and prime conditions of \ref{B5.07y} at one family cutoff using the level-one realization of \ref{B8.08}. The lift comparison \ref{B5.07f} and scalar identities \ref{B5.07b} give the bound for every \(C\), whose coefficient in \(L_{6,j}\) is zero. The complete proof is in \S\ref{prf:B8.08b}.
\end{proof}

\begin{lemma}[A common cutoff for the reference index bounds]\label{B8.08c}
Use the level-zero realization at the points of \(\mathcal P\cap\mathcal R\) in Definition~\ref{def:B8-reference-points}. Use \(D_R,i_{jw},e_w,\ell_w,c_w,F_{0,j,w},J_{0,j,w}\) as in \ref{B8.38b}, and set \(U=L(B_0+N_c)\). Here \(D_R>0\), \(m\ge1\), and \(U\ge0\).

Suppose that on a family satisfying the division estimate \eqref{eq:B8-divinN} there are \(a\ge0\) and \(x_0\) with the following property. For \(x>x_1,\ x\ge x_0\), every index \(i\in I_x\) and every such point over \((x,i)\) admit numbers \(u_{jw}\ge0\) satisfying
\[
 \frac{1}{D_R}\sum_j i_{jw}
 \le \sum_j(2e_wF_{0,j,w}+u_{jw})+c_w
 \quad(w\in\mathrm{Tw}),\qquad
 \sum_{w\in\mathrm{Tw}}\ell_w(u_{jw}+e_wJ_{0,j,w})\le aU
 \quad(1\le j\le m).
\]
Suppose \(J_{0,j,w}\ge0\) on these points. Then the three reference inequalities
\eqref{eq:B8-paired-index}--\eqref{eq:B8-paired-error} hold at every one of the same points for every \(C\ge2a\). In particular, if the hypotheses hold with the family quantifiers of \ref{B8.39} on \(\mathcal F_{\rm cov}\), so does the reference index implication of \ref{B8.06}.
\end{lemma}

\begin{proof}[Outline]
Take the positive part of the excess index at each place and divide it equally among the layers. Averaging the weighted errors and adding the bound for the nonnegative reference terms gives all three inequalities for every \(C\ge2a\) at the original cutoff. The complete proof is in \S\ref{prf:B8.08c}.
\end{proof}


\section{Evaluation, periods, and the selected spectrum}

\begin{theorem}[Six arithmetic rows]\label{B8.09}
Fix construction conditions \(\mathcal S\) and a point-to-bundle relation \(L_k\). Take \(KP=KP_x\) and let \(Q,Q'\) mean constructed and good membership for \(\mathcal S\), respectively, in the sense defined after (8.6). The following \(\operatorname{Lg1}\)-rows are closed on \(W_3\) and eventual at selected points.
\begin{enumerate}
 \item For each word \(u\), reference evaluation is injective on the specified labels and bijective on their stabilizers. Its complex map \(\mathrm{Pt}_u\) takes \(F_i(m)\) into \(\mathrm{ptC}_u\). An integral rational realization \(P_u\) satisfies, for each infinite \(v\) and rational level-\(m\) vector \(z\),
 \begin{equation}\tag{8.17}\label{eq:p8c1-point-norm}
 q_{u,\mathrm{pt},v}(P_uz)\le\mathrm{ptB}(D,u)q_{u,m,v}(z),
 \qquad \mathrm{ptB}(D,u)=\prod_\tau T_u^{-\deg_\tau}\mathrm{beOn}_\tau^2.
 \end{equation}
 Here \(q_v\) is the real majorant diagonal, and \(\deg_\tau,\mathrm{beOn}_\tau\) are the chain degree and norm factor. The same constant bounds \(\|\mathrm{Pt}_u\phi\|^2/\|\phi\|^2\) for nonzero \(\phi\in F_i(m)\).
 \item Each used nonreference tuple has a rational integral twisted evaluation for every word.
 \item Assume the common denominator condition of item~6 of Definition~\ref{B4.03}: every twisted chain has denominator \(D_{\Gamma,u_0}\), and every class stabilizer order at every layer \(0\le j\le m\) of every twisted chain divides this integer. For each reference point-lattice vector and \(H\in\widetilde M_{\mu_0}\cap V_0(\mu_0)\), one \(y\in K_D\) represents the period through \(\iota_D\). At every finite \(w\), with residue characteristic \(p\) and ramification index \(e_w\),
 \begin{equation}\tag{8.18}\label{eq:p8c1-period-valuation}
 y=0\quad\text{or}\quad v_w(y)\ge-e_w v_p(D_{\Gamma,u}).
 \end{equation}
 \item Under the same common denominator condition, one representative of the twisted period for each used \(\mu\ne\mu_0\), point-lattice vector, and \(H\in\widetilde M_\mu\cap V_0(\mu)\) satisfies \eqref{eq:p8c1-period-valuation} at every \(w\) with \(p\notin Q_p(\mu)\).
 \item At each point, reindexing a signature basis identifies the scaled Haar measure with the prescribed compact-dual normalization. At every infinite place the prescribed discrete-series formal degree is attained and every discrete series with the specified parameters has that degree for this measure. The adelic quotient measures give the prescribed fundamental-domain volumes for all admissible signatures. The point also has a compatible symmetric column system. For each word the level-zero character has finite exponent and values in \(\iota_D(K_D)\).
 \item Let \(p=T_u\) and \(E^T_{u,w}=E^{\rm atT}_{u,w}\) as in \ref{B4.22}. For every input in \(N_{\lambda,u,w}\) and \(c\in K_D^\times\) with \(v_w(c)\ge e_wmE^T_{u,w}\), scaling the reference column by \(\iota_D(c)\) puts it in the image of \(S_m(w)\). For each used \(\mu\ne\mu_0\) with \(p\notin Q_p(\mu)\), the twisted column is represented by \(az\) under the terminal embedding, with \(z\) in the local twisted level-\(m\) lattice and \(a=0\) or \(v_w(a)\ge-e_wmE^T_{u,w}\).
\end{enumerate}
\end{theorem}

\begin{proof}[Outline]
Extracts the column system, lowering maps, and Petersson decomposition from good members, with local Weyl-dual row lifts from \ref{lem:p8c1-row-forms}. Applies \ref{B4.18a} for items 1 to 4, the real comparison for item 5, and \ref{B4.22a} for item 6. The complete proof is in \S\ref{prf:B8.09}.
\end{proof}

For the term vectors \(v_t\) of the spectral rows and a set \(S\) of term pairs, put
\[
 \operatorname{PairAbs}_{u,u'}(S)=\Gamma_1\sum_{(t,t')\in S}|\langle v_t,v_{t'}\rangle|, \tag{8.19}
\]
where \(\Gamma_1>0\) is the finite-type normalization. Among same-row pairs with both Gram values \(G\ne0\), let \(S_{\rm big}\) have corank at least two and \(S_{\rm nsc}\) consist of nonscalar generic pairs of corank one. Write \(d_w^C,D_C^\star\) for the model column dimensions and main term, and \(d_w^S,D^\star\) for their arithmetic counterparts. Set
\[
 \epsilon_L(D)=e^{-kf\log k/3},\qquad
 \Delta_{u,u'}=\frac{d_u^C d_{u'}^C}{D_C^\star}. \tag{8.20}
\]

\begin{theorem}[Three estimates for the selected spectral rows]\label{B8.10}
Assume the finite-volume and nonsingular Gram-density assertions of the second formula of \ref{B8.13}. For every \(KP,Q,Q',L_k\), the following estimates give \(\operatorname{Lg1}\)-rows closed on \(W_3\). Put
\[
 C'=C_k,\qquad C_\ell=\min\{C_f,C_w,C_\Delta,C_{\rm count},C_r,C_{\rm gap}\},
 \qquad p=(C_{YL}+1,\tfrac14,3,\tfrac12,C_{Tu},3).
\]
Use the family constants of Definition~\ref{def:p3c5-parameters}. With \(B=\log b\), the rank bounds there are \(C_km^5fB^2\le k\) and \(C_fB\le f\log k\). The exponent \(C_w\) occurs in \(k^{C_w}\le w''\). The spacing bound uses \(C_\Delta\) and the term \((NfB)^{C_{\rm count}}\). The height bound uses \(C_r r^{19}f^5B\le H_0\), and \(C_{\rm gap}\) occurs in the consecutive-gap bound. The third coordinate of \(p\) is the separate weight upper-bound parameter in Definition~\ref{B8.02c}; its value is \(3\).
Fix any real-valued function \(a_F\) of the outer list. For each eligible \(F_m,CH\) above the closed thresholds there is a common cutoff \(x_0\). Each assertion holds under its point hypotheses for all \(x\ge x_0\) with \(x>F_m.x_1\) at or above the selection cutoff, and every \(i\in I_x\).

For \eqref{eq:p8c1-loop-rows} assume \(\mathrm{Cnt.Hyp}(D,L,u,u';C',C_\ell,p,a_F(F_m.c_{\rm out}))\) of \ref{B8.02c}, absolute summability of \(\|\operatorname{op}(t)\|\) and \(\|v_t\|\), and the following compatibility of the Hida model \(X\) with its full data \(Y\). It identifies \(d_w^C=d_w^S,\ D_C^\star=D^\star\). It includes level-subgroup normalization and neatness for every word, the row-function conditions and integrable coefficient projector of \ref{B8.02y}, and the standard adelic additive character with compatible measures, local models and local arithmetic data. Require also \(T_w\bigl(\pi/(v_Eg_0g_1)\bigr)^{2r_S}\ge1\) for every word \(w\), with \(v_E=\operatorname{rd}(E)/2\) and \(r_S=m/2-1\). Then
\begin{equation}\tag{8.21}\label{eq:p8c1-loop-rows}
 \operatorname{PairAbs}_{u,u'}(S_{\rm big})\le\epsilon_L(D)\Delta_{u,u'},
 \qquad
 \operatorname{PairAbs}_{u,u'}(S_{\rm nsc})\le\epsilon_L(D)\Delta_{u,u'} .
\end{equation}
For a set of terms \(E_1\) write \(\operatorname{part}(E_1)=\sum_{t\in E_1}\operatorname{op}(t)\). Let \(K_0\) be the sum of these operators for the identity cell and the generic and nongeneric cells with \(G\ne0\). Under the source condition of \ref{B8.02c} at \(3,C',C_\ell,C_{YL}+1,\tfrac14,3,\tfrac12,C_{Tu}\), distinct words satisfy
\begin{equation}\tag{8.22}\label{eq:p8c1-distinct-row}
 \|K_0(u,u')\|_{\rm HS}^2\le (1+b^{-6})\frac{d_u^S d_{u'}^S}{D^\star}.
\end{equation}
Under \(\mathrm{Cnt.Hyp0}(D,L;C',C_\ell,p)\), every level prime satisfies
\begin{equation}\tag{8.23}\label{eq:p8c1-level-row}
 1\le T_w\left(\frac{\pi}{v_Eg_0g_1}\right)^{2r_S},
 \qquad r_S=m/2-1 .
\end{equation}
\end{theorem}

\begin{proof}[Outline]
Use the volume ratio of \ref{B8.49} and the critical-level bounds to force the level inequality. Apply \ref{B8.10a} for the loop estimates and use the ratios from \ref{B5.11b} with \ref{B8.10b} and \ref{B8.10d} for distinct words. The complete proof is in \S\ref{prf:B8.10}.
\end{proof}

\begin{lemma}[The terminal-flag Gram inequality]\label{B8.10c}
Use the selected local models and row representatives under the source and size hypotheses of \eqref{eq:p8c1-distinct-row}. The levels \(T,T'\) are congruent modulo \((g_0g_1)^2\). The scalar-partner coefficients \(c_{\rho p}\) and the terminal kernel \(K_{\rm top}\) use the same local characters and measures. Let \(\alpha>0\) be the product of local normalizers and \(s_d=\sup\mathrm{Dens}/\alpha\) over represented matrices. For finite sets \(Y\) of terminal flags and \(S\) of scalar-partner rows and for \(z\in\C^Y\),
\[
 \sum_{\rho\in S}\left|\sum_{p\in Y}\overline{z_p}c_{\rho p}\right|^2
 \le s_d\sum_{p,q\in Y}\overline{z_p}K_{\rm top}(p,q)z_q .
\]
\end{lemma}

\begin{proof}[Outline]
Expand the kernel into positive rank-one summands and use rational approximation and gluing to identify partner fibres. Shifted determinants fix the frame normalizers on each support, so the row multipliers bounded by \(s_d\) give the inequality termwise. The complete proof is in \S\ref{prf:B8.10c}.
\end{proof}

\begin{lemma}[The interpolation grid at a real place]\label{B8.10ca}
Put \(n=k+2r+1,\ K=3,\ N=3(r+1)\), with \(k\ge2\). Let \(T_p,T_q>0\) be the normalized final norms. Use two decreasing weights \(u,u'\) with coordinates in \(u_1+\Z\) and a decreasing spectral parameter \(\lambda\), with all coordinates bounded by \(b\ge\max(3,k)\) in absolute value. Require the common interlacing inequalities
\[
 \max(u_{i+1},u'_{i+1})+\tfrac12
 \le\min(u_i,u'_i)-\tfrac12\qquad(1\le i<k).
\]
For a cut at coordinate \(1\le l\le k\) use
\[
 Y=\{2-\max(u_l,u'_l)+5j:0\le j<N(2r+1)\}.
\]
Assume \(|u_l-u'_l|\ge30(r+1)^2+6\) and \(c>0,\ c+2\le b\). The adjacent gaps of the first \(k-1\) coordinates of \(\lambda\) are at least \(1+c\). For each of the two weights \(a=u,u'\), assume half-unit interlacing with these coordinates, or the following one-index variant. At \(x=-y\), for each \(1\le i<k\) outside one fixed exceptional index, require \(a_i\ge\lambda_i+1/2\) if \(\lambda_i>x\) and \(a_{i+1}\le\lambda_i-1/2\) otherwise. These conditions hold at each \(y\in Y\). Put
\[
 P_*(y)=\prod_{a=1}^n(|y+\lambda_a|+c),\qquad
 q(y)=-T_pT_q\prod_{i=1}^k(y+1/2+u_i)(y+1/2+u'_i),\qquad
 \Lambda^2=\max\{1,(4/3)^N b^{8r+20}/(T_pT_q)\}.
\]
A positive letter based at \(y\in Y\) has evaluation point \(y\), and a negative letter based there has evaluation point \(y+1\). Take lists \(w\) with distinct evaluation points separated by more than two. Define \(q_w(y)\) by multiplying \(q(y)\) by \((y-z+1)/(y-z)\) for each positive letter with point \(z\) and by \((z-y)/(z-y-1)\) for each negative one.

If \(|w|+3\le N\), there are three disjoint rows \(y_{aj}\) of \(d+1\) nodes in \(Y\) avoiding the bases of \(w\). Here \(d=2i\) for a wave at level \(0\le i\le r\), or \(d=2(i-s)+s\) for a lift of length \(1\le s\le i\le r\). For \(L_{aj}=\prod_{h\ne j}(y_{aj}-y_{ah})^{-1}\) they satisfy
\begin{equation}\label{eq:p8c1-wave-grid}
 \sum_{j_1,j_2,j_3}\prod_{a=1}^3 |L_{a j_a}|
 \left(\frac{P_*(y_{a j_a})P_*(y_{a j_a}+1)}{q_w(y_{a j_a})}\right)^{1/2}
 \le\Lambda^3 .
\end{equation}
\end{lemma}

\begin{proof}[Outline]
Use displacement and interlacing to bound the denominators and show \(q(y)>0\). Pair spectral factors with neighboring weight factors and bound the ratios by the gaps, then use three disjoint rows of unused nodes and the Lagrange coefficient sums. The complete proof is in \S\ref{prf:B8.10ca}.
\end{proof}

\begin{lemma}[Normalization in a determinant expansion]\label{B8.10o}
Let \(s=r+1,\ m=2s,\ n\ge s,\ b\ge\max(2,n)\), and \(C_h\ge0\). Write \(\mathcal J_\mu\) for finitely many vector-valued functions on a measured coordinate chart and \(I_\mu\) for their finite squared \(L^2\) norms, with the common point weight included in the measure. Let \(D\) be a scalar function on the chart. Start at \(\mu=0\). At each index reached in fewer than \(a\) steps suppose
\[
 D\mathcal J_\mu=\sum_{\substack{I\subset\{1,\ldots,n\}\\ |I|=s}}
              \tau_{\mu I}\mathcal J_{\mu+\mathbf1_I},
 \qquad
 |\tau_{\mu I}|\le H^s,\qquad H\le b^{(C_h+10)m}.
\]
Assume the identities hold at the represented points under consideration, that zero-energy functions vanish there, and that \(I_{\mu+\mathbf1_I}\le (\exp(1)n)^s I_\mu\) for \(I_\mu>0\). There is an expansion of \(D^a\mathcal J_0\) at those points in functions of energy at most \(I_0\), with sum of absolute coefficients at most \(b^{(C_h/2+6)m^2a}\).
\end{lemma}

\begin{proof}[Outline]
Rescale each positive-energy function to the initial energy and sum the absolute coefficients of its children using the energy ratio. Iterate the one-step bound over paths and combine terms at the same index. The complete proof is in \S\ref{prf:B8.10o}.
\end{proof}

\begin{definition}[Partial orbital pairings and Capelli-weighted energies]\label{B8.10ha}
Fix integers \(k\ge2\), \(r\ge0\), \(m=2r+2\), and \(n=2r+k+1\). A local Bessel datum of rank \((k,r)\) consists of the Cayley Hermitian matrix \(H_{k,r}\) on \(\C^n\) of signature \((k+r,r+1)\), its unitary group \(G=\U(H_{k,r})\), a discrete series representation \((\sigma,\mathcal H_\sigma)\) of \(G\) with decreasing Harish--Chandra parameter \(\lambda=(\lambda_1,\ldots,\lambda_n)\), two Bessel frame vectors \(z,z'\in\C^n\) in the perpendicular of the length-\(r\) isotropic flag with positive final norms \(T_p=-\Re h_{H_{k,r}}(z,z)>0\) and \(T_q=-\Re h_{H_{k,r}}(z',z')>0\), compact definite stabilizers \(U_z,U_{z'}\simeq\U(k)\) equipped with probability Haar measures \(da,da'\) and irreducible finite-dimensional unitary representations \((\pi,W_p),(\pi',W_q)\) of decreasing highest weights \(u,u'\in\R^k\), unipotent radicals \(N_z,N_{z'}\) carrying unipotent Haar measures \(dn,dn'\) of finite positive scale \(\kappa_N\), and Bessel characters \(\psi_z,\psi_{z'}\) of nonzero parameter \(c_\psi\in\C^\times\). Standardized coordinates mean \(c_\psi=-2i\) and standard unipotent Haar measure (\(\kappa_N=1\)).

For \(0\le j\le r+1\), a level-\(j\) block point \(\gamma=((t_1,s_1),\ldots,(t_d,s_d))\) is a list of pairs \((t_i,s_i)\in\C^\times\times\N_{\ge1}\) with \(\sum_{i=1}^d s_i=j\). Its block matrix \(M(\gamma)\) is the \(j\times j\) block-diagonal matrix with \(i\)-th block \(t_iJ_{s_i}\), where \(J_s\) has ones on its antidiagonal, and its order-\(l\) leading minor is \(D_l(\gamma)=\det M(\gamma)_{1\le a,b\le l}\) (\(1\le l\le j\), \(D_0(\gamma)=1\)). At top level \(j=r+1\), define
\[
 G_{\rm loc}(\gamma)=|D_{r+1}(\gamma)|^2-T_pT_q|D_r(\gamma)|^2.
\]
The level-\(j\) point weight \(p_{k,r}(T_p,T_q;\gamma)\) equals \(1\) at \(j=0\), and for \(j\ge1\) equals the product over the blocks \((t_i,s_i)\), starting after \(a_i=\sum_{h<i}s_h\) positions, of \(|t_i|^{2s_i(2(r+1-a_i-s_i)+k+s_i-2)}\), with the last factor replaced by \(\max(0,|t_d|^2-T_pT_q)^{k-1}\) when \(j=r+1\) and \(s_d=1\). On any top-level block point represented by an element \(g\in G\) whose frame-pairing matrix is \(M(\gamma)\), one has \(G_{\rm loc}(\gamma)\ge0\) and \(p_{k,r}(T_p,T_q;\gamma)\ge0\).

For a smooth operator-valued coefficient \(H:G\to\operatorname{Hom}_\C(W_q,W_p)\) and a top-level block point \(\gamma\), whenever \(g\in G\) represents \(M(\gamma)\), \(\psi_z\) and \(g\)-conjugated \(\psi_{z'}\) agree on the common unipotent stabilizer \(R_g=N_z\cap gN_{z'}g^{-1}\), \(\mu_R\) is the root-coordinate measure on \(R_g\), and \(\nu\) disintegrates \((da\,dn)\times(da'\,dn')\) along \(\mu_R\), define the orbital integral \(J(H;g,\nu)\in\operatorname{Hom}_\C(W_q,W_p)\) by integrating \(\psi_z(n)\overline{\psi_{z'}(n')}\,\pi(a)H(a^{-1}n^{-1}gn'a')\pi'(a')^{-1}\) against \(d\nu\). Let \(\|J(H;\gamma)\|_{\rm HS}^2\in[0,\infty]\) be the supremum of \(\|J(H;g,\nu)\|_{\rm HS}^2\) over all admissible \((g,\mu_R,\nu)\), with value \(0\) when none exists.

In standardized coordinates, for \(0\le j\le r+1\) and level-\(j\) block points \(\gamma,\gamma'\), the partial orbital pairing of two smooth rapidly decaying coefficients \(H_1,H_2\) integrates the first \(j\) flag positions against \(\psi_z,\overline{\psi_{z'}}\) and \(\pi,\pi'\), multiplies the two arguments by \(p_{k,r}(T_p,T_q;\gamma)^{1/2}\) and \(p_{k,r}(T_p,T_q;\gamma')^{1/2}\), and takes the inner product of the resulting functions on the residual corank-\((r+1-j)\) Bessel slice. At \(j=0\) (\(\gamma=\gamma'=[]\)) this is the sesquilinear Bessel pairing \(\beta(H_1,H_2)\) with empty-list weight \(1\); at \(j=r+1\) and \(\gamma=\gamma'\) with \(G_{\rm loc}(\gamma)>0\), its diagonal value \(E(H;\gamma)\) equals \(p_{k,r}(T_p,T_q;\gamma)\|J(H;\gamma)\|_{\rm HS}^2\).

Let \(\mathcal Z(\mathfrak g_\C)\) be the center of the universal enveloping algebra of the complexified Lie algebra \(\mathfrak g_\C\simeq\mathfrak{gl}_n(\C)\) of \(G\). For \(y\in\R\), the Capelli central element \(C(y)\in\mathcal Z(\mathfrak g_\C)\) acts on any representation of infinitesimal character \(\mu=(\mu_1,\ldots,\mu_n)\) by the scalar \(\prod_{a=1}^n(y+\mu_a)\). Using the grid \(Y\) and signed letters \((\epsilon,y_0)\) (\(\epsilon\in\{+,-\}\), \(y_0\in Y\), with evaluation point \(z(\epsilon,y_0)=y_0\) for \(\epsilon=+\) and \(y_0+1\) for \(\epsilon=-\)) of Lemma~\ref{B8.10ca}, for an admissible letter list \(w\) of length at most \(N=3(r+1)\) and a shift parameter \(c>0\), put
\[
 n_c(w)=\prod_{(\epsilon,y_0)\in w}P_*\bigl(z(\epsilon,y_0)\bigr),\qquad
 P_*(z)=\prod_{a=1}^n(|z+\lambda_a|+c),
\]
and let \(C_w=\prod_{(\epsilon,y_0)\in w}C(z(\epsilon,y_0))\in\mathcal Z(\mathfrak g_\C)\) denote the ordered product of the Capelli operators of \(w\). Inserting the letters of \(w\) into the level-\(j\) partial orbital pairing defines the Capelli-weighted partial orbital pairing, and multiplying \(H\) by a smooth Gram cutoff \(\chi(G_{\rm loc}/\epsilon)\) supported on \(\{G_{\rm loc}>0\}\) gives the cut Gram forms of the partial orbital pairings. Its diagonal energy \(E_j(w;H;\gamma)\) satisfies \(E_{r+1}([];H;\gamma)=E(H;\gamma)\) and \(E_0(w;H;[])=\Re\beta(C_wH,H)\). For \(K\) additional grid nodes \(\mathbf y=(y_{a j_a})_{a=1}^K\) avoiding the bases of \(w\), a superscript \(+,-\) on a partial pairing denotes inserting the positive letters at \(\mathbf y\) in the first argument and their negative companions in the second, contracted by the word form of \(w\).
\end{definition}

\begin{lemma}[Polynomial contractions between compact types]\label{B8.10hba}
Let \(k,r\in\N\), \(k\ge2\), and let \(W_b\subset\C^{2r+k+1}\) be the central coordinate subspace on positions \(r+1,\ldots,r+k+1\) for the Cayley Hermitian form \(h\). Let \(z\in W_b\) satisfy \(-\Re h(z,z)>0\), and put \(E_z=W_b\cap z^\perp\). Let \((\pi,V)\) be a nonzero finite-dimensional continuous irreducible unitary representation of \(U_z\), with compact parameter \(u=(u_1,\ldots,u_k)\) in the Capelli normalization, whose shifted coordinates are \(u_j-k/2+1/2\). Suppose that every \(u_j-u_{j+1}\) is a positive integer. Fix \(1\le l<k\), \(d\in\N\), and \(b\ge1\) such that
\[
 |u_j|\le b\quad(1\le j\le k),\qquad d\le b,\qquad
 2d+1\le u_l-u_{l+1}.
\]
There are compact types \((\pi_+,V_+)\), \((\pi_-,V_-)\) with parameters \(u+d\mathbf1_{l+1}\), \(u-d\mathbf1_l\), and maps
\[
 A_\pm:\C^{2r+k+1}\longrightarrow\operatorname{Hom}_{\C}(V_\pm,V)
\]
with the following properties. Each map is a finite sum of fixed operators multiplied by polynomials of total degree at most \(d\) in \(\omega\) and \(\overline\omega\). For \(a\in U_z\), \(t\in\R\), and \(\omega\in\C^{2r+k+1}\),
\[
 A_\pm(a\omega)\pi_\pm(a)=\pi(a)A_\pm(\omega),\qquad
 A_\pm(t\omega)=t^dA_\pm(\omega).
\]
For \(\omega\in E_z\),
\[
 \|A_\pm(\omega)\|^2\le(\Re h(\omega,\omega))^d,
\]
and for every \(v\in V\),
\[
 (\Re h(\omega,\omega))^d\|v\|^2
 \le(3b)^{2d}
 \bigl(\|A_+(\omega)^*v\|^2+\|A_-(\omega)^*v\|^2\bigr).
\]
\end{lemma}
\begin{proof}[Outline]
Compose one-step contractions into degree-\(d\) chains \(A_\pm\) diagonal on Capelli eigenspaces \(V_M\). The gap \(2d+1\) bounds one eigenvalue below by \((3b)^{-2d}\) on each \(V_M\), and equivariance with homogeneity extends the bound to \(E_z\). The complete proof is in \S\ref{prf:B8.10hba}.
\end{proof}

\begin{proposition}[Standardized minor-moment bound on the wave grid]\label{B8.10hb}
Fix \(C_h,C_a\in\R\) with \(C_a\ge0\). There exist \(C\ge0\) and \(C_{g,1}\ge1000\) such that for every \(C_g\ge C_{g,1}\) the following holds at any standardized local Bessel datum of Definition~\ref{B8.10ha} whose discrete series has the prescribed signature \(\varepsilon_{k,r}\), positive on the indices \(1,\dots,k-1\) and alternating \(+,-,+,\dots\) on the indices \(k,\dots,n\), with \(k\ge2\), \(m=2r+2\), \(n=2r+k+1\), scale \(b\ge2\), final norms \(1\le T_p,T_q\le b^{2|C_h|m}\), \(K\)-finite discrete-series coefficient \(F\), and big-cell section of squared section integral \(I_{\rm std}\). Assume that \(\lambda,u,u'\) have coordinates bounded by \(b\) in absolute value and satisfy the strict half-integral interlacing, adjacent gaps \(\lambda_s-\lambda_{s'}\ge C_gm^4\) (\(s<s'\)), opposite-sign gap at least \(C_gm^{9/2}(3(r+1)n^2+n^3)\), a splitting scale \(b_v\ge C_gn^7\) satisfying
\[
 0\le\lambda_a\le b_v/8\quad(1\le a<k),\qquad
 -2b_v\le\lambda_a\le-b_v\quad(k\le a\le n),
\]
and the equal-weight or distinct-weight gap hypothesis of Lemma~\ref{B8.10h}. Then for every \(1\le l\le r+1\), every integer \(0\le e\le C_am^3\), and every top-level block point \(\gamma\) with \(G_{\rm loc}(\gamma)\ne0\),
\begin{equation}\label{eq:p8c1-minor-moment}
 |D_l(\gamma)|^{2e}\,p_{k,r}(T_p,T_q;\gamma)\,\|J(F;\gamma)\|_{\rm HS}^2
 \le b^{Cm^5}\max\bigl(1,|G_{\rm loc}(\gamma)|^{-1}\bigr)^{40(r+1)^2}I_{\rm std}.
\end{equation}
\end{proposition}

\begin{proof}[Outline]
Cuts off \(G_{\rm loc}^{p_0}F\) by infinitesimal character and applies \ref{B8.10ca} (with \ref{B8.10hba} when \(u=u'\)) for wave-grid decay. A period-normalized shift family and exterior-power Pieri expansion then bound \(|D_l(\gamma)|^{2e}\). The complete proof is in \S\ref{prf:B8.10hb}.
\end{proof}

\begin{lemma}[Torus rescaling of local Bessel integrals]\label{B8.10hc}
Under the hypotheses of Proposition~\ref{B8.10hb} except that \(1\le T_p,T_q\le b^{|C_h|m}\), the Bessel character has parameter \(-2i\kappa_0\) with \(\kappa_0>0\), \(\kappa_0^{\pm1}\le b^{|C_h|}\), \(\kappa_0^{2r}T_p\ge1\), \(\kappa_0^{2r}T_q\ge1\), and the unipotent Haar measure has arbitrary finite positive scale \(\kappa_N\), diagonal torus conjugation with parameters \(\kappa=\kappa_0^{-1}\) and \(\tau=\kappa_0^r\) carries the datum to character parameter \(c_\psi=-2i\), with a finite positive unipotent Haar scale, and final norms \(\kappa_0^{2r}T_p,\kappa_0^{2r}T_q\). Put \(g_*(\kappa_0)=\kappa^{2r(r+1)}\tau^{4(r+1)}=\kappa_0^{2r(r+1)}\) and \(F_*(\kappa_0)=\max(\kappa_0,\kappa_0^{-1})^{2m^2}\max(\kappa_0^r,\kappa_0^{-r})^{2m}\). Then \(F_*(\kappa_0)\le b^{4|C_h|m^2}\), and every \(1\le l\le r+1\), integer \(0\le e\le C_am^3\), and top-level block point \(\gamma\) with \(G_{\rm loc}(\gamma)\ne0\) satisfy
\[
 |D_l(\gamma)|^{2e}\,p_{k,r}(T_p,T_q;\gamma)\,\|J(F;\gamma)\|_{\rm HS}^2
 \le b^{2|C_h|m^2e}\,F_*(\kappa_0)\,b^{Cm^5}
 \max\bigl(1,|g_*(\kappa_0)G_{\rm loc}(\gamma)|^{-1}\bigr)^{40(r+1)^2}I,
\]
where \(I\) is the squared big-cell section integral in the original coordinates and \(C\) is the constant of Proposition~\ref{B8.10hb}.
\end{lemma}

\begin{proof}[Outline]
Shows scaling \(dn=s\,dn_{\rm std}\) multiplies \(I\) and the orbital energy of \ref{B8.10ha} by \(s^4\). Diagonal torus conjugation (\(\kappa=\kappa_0^{-1},\tau=\kappa_0^r\)) reduces to parameter \(-2i\), where \ref{B8.10hb} and minor bounds apply. The complete proof is in \S\ref{prf:B8.10hc}.
\end{proof}

\begin{lemma}[The weighted sum of real orbital integrals]\label{B8.10h}
Fix \(A_w\ge0\) and \(C_h\in\R\). There are \(C_b,C_{g,0}\) such that the following holds for \(C_g\ge C_{g,0}\). Put \(r=r_S,\ m=2r+2,\ n=k+2r+1,\ \xi=\pi/(v_Eg_0g_1)\). Use the same \(T,T'\in\N\) and \(b\) at every real place, with
\[
 \xi>0,\quad \xi^{\pm1}\le b^{C_h},\quad
 1\le T,T'\le b^{C_hm},\qquad T\xi^{2r},\ T'\xi^{2r}\ge1 .
\]
At each place take the prescribed signature, standard character, positive finite section normalization, two Bessel frames of final norms \(T,T'\), a \(K\)-finite discrete-series coefficient, and a big-cell section. Use the Harish--Chandra parameter \(\lambda\) prescribed by the point, with coordinate gaps at least \(C_gm^4\). For each length-\(k\) weight \(u,u'\) require \(u_i>\lambda_i>u_{i+1}\) for \(1\le i<k\) and \(u_k>\lambda_k\), with the corresponding inequalities for \(u'\). The parameter and type coordinates have the prescribed integral or half-integral normalization, so the strict interlacing separates them by at least \(1/2\). All their coordinates are bounded by \(b\) in absolute value. Opposite signature signs have gap at least
\[
 C_gm^{9/2}\bigl(3(r+1)n^2+n^3\bigr).
\]
There is \(b_v\ge C_gn^7\) with \(\lambda_a\in[0,b_v/8]\) when \(\min_i u_i<\lambda_a<\max_i u_i\), and \(\lambda_a\in[-2b_v,-b_v]\) otherwise. If \(u=u'\), some consecutive gap \(u_l-u_{l+1}\ge2C_g(r+1)^2+1\) has a coordinate of the first group at or above \(u_l\) and one at or below \(u_{l+1}\). If \(u\ne u'\), some interval of length at least \(C_g(r+1)^2\) lies between \(-u_l\) and \(-u'_l\), at distance at least two from both endpoints.

Sum over rational block points \(\gamma\) with an integral representative and nonzero integral \(G=D_{r+1}\overline{D_{r+1}}-TT'D_r\overline{D_r}\). Write \(D_j\) for its leading minors and use the block ends in products over \(j\). At an infinite place \(v\), use the complex absolute value of the chosen embedding. Put
\[
 W_v=|G|_v\prod_j|D_j|_v^2,\qquad
 w_v=\prod_j(1+|D_j|_v^2)^2,\qquad
 W=\prod_v W_v,\quad w=\prod_v w_v .
\]
A block of length \(s\) with entry \(t\) has matrix \(tJ_s\), where \(J_s\) has ones on its antidiagonal. For a block starting after \(a\) positions its local entry is \(z=\operatorname{rd}(E)^{-(2r+3-2(a+1)-s)}v(t)\). The rotated entry is \((-1)^r(-i)^{2a+s-1}z\). The local point weight \(p_v\) is the product of \(|z|^{2s(2(r+1-a-s)+k+s-2)}\) over the blocks. For a represented final singleton, use the factor \((|z|^2-TT')^{k-1}\). Put \(p(\gamma)=\prod_v p_v(\gamma)\).

Write the coefficient as \(F(h)=P\sigma(h)Q\), with the images of \(Q\) and \(P^*\) compact-finite, and let \(\pi,\pi'\) be the two definite types. For each unitary representative \(h\) of the rotated matrix of frame pairings, require the two unipotent characters to agree on the common stabilizer. In compact and unipotent coordinates \((a,n),(a',n')\) the integrand of the orbital operator is
\[
 \psi_z(n)\overline{\psi_{z'}(n')}\,\pi(a)
 F(a^{-1}n^{-1}hn'a')\pi'(a')^{-1}.
\]
Here \(\psi_z,\psi_{z'}\) are the Bessel characters of the two frames with parameter \(-2i\xi\) in the arithmetic chart. Use each quotient measure that disintegrates the product Haar measure along the common stabilizer with the prescribed root-coordinate measure. Write \(\|J_v(\gamma)\|_{\rm HS}^2\) for the supremum of the squared norms over these choices, with value zero on empty support. Let \(I_v\) be the squared section integral in diagonal coordinates \((c_0,\ldots,c_{r-1},d)\), with Lebesgue density
\[
 {\bf1}_{|d|^2\ge TT'}(|d|^2-TT')^{k-1}
 \prod_{a=0}^{r-1}|c_a|^{2(2(r-a)+k-1)} .
\]
Then
\[
 \sum_\gamma w(\gamma)p(\gamma)W(\gamma)^{A_w(r+1)^2}
               \prod_v\|J_v(\gamma)\|_{\rm HS}^2
 \le b^{C_bm^5f}\prod_v I_v .
\]
\end{lemma}

\begin{proof}[Outline]
Dominate each local weight by one minor power, apply Lemma~\ref{B8.10hc} with thresholds from Proposition~\ref{B8.10hb}, and multiply over infinite places using \(\prod_v|G(\gamma)|_v\ge1\). Summing against \(\sum_\gamma w(\gamma)^{-1}\) via a polydisk count yields the bound. The complete proof is in \S\ref{prf:B8.10h}.
\end{proof}

\begin{lemma}[A character sum on the norm congruence]\label{B8.10j}
Let \(F\) be a nonarchimedean local field of characteristic zero, with valuation ring \(O\), uniformizer \(\varpi\), and residue cardinality \(q\). Let \(O_E\) be the integers of the unramified quadratic extension or \(O\times O\) with exchange involution. Write \(h\) for an integral Hermitian form, linear in the first variable, on an \(O_E\)-lattice \(U\). Suppose \(W\subseteq U\) contains \(\rho\ge4\) orthonormal vectors. Let \(a\ge0\), \(N\ge a+1\), and use probability measure on \(U/\varpi^N U\). Write \(Z\) for the congruence \(h(u,u)\in\varpi^aO\). For every additive character \(\chi\) of this quotient,
\[
 \left|q^a\mathbb E({\bf1}_Z\chi)-{\bf1}_{\chi=1}\right|
 \le \frac87q^{1-\rho}.
\]
The probability that \(Z\) holds and \(h(u,W)\) is a proper ideal of \(O_E\) is at most \(q^{-a}(16/7)q^{1-\rho}\).
\end{lemma}

\begin{proof}[Outline]
Expand the norm congruence in additive characters on the orthonormal coordinates. Square each coordinate sum and use trace orthogonality, multiply the bounds and sum by conductor, then restrict the coordinates to each maximal ideal to bound the proper-pairing event. The complete proof is in \S\ref{prf:B8.10j}.
\end{proof}

\begin{lemma}[Counting the two normal coefficients]\label{B8.10n}
Let \(E\) be a CM field of degree \(2f\), write \(\bar x\) for complex conjugation, and write \(N\) for the absolute ideal norm. For a nonzero fractional ideal \(J\) and \(a\in E^\times\),
\[
 \#\{x\in J:x\bar x=a\}\,(NJ)^2
          \le 8f^2\,|N_{E/\Q}a|.
\]
In particular let \(T,T'\ge3,\ P\ge1\), and suppose a set of labels injects into pairs \((x,x')\) with \(x\bar x=x'\bar x'=1\), \(x\in J,\ x'\in J'\). If
\[
 (NJ)^{-2}\le T^{2f}P^2,\qquad
 (NJ')^{-2}\le {T'}^{2f}P^2,
\]
the number of labels is at most \((TT')^{6f}P^6\).
\end{lemma}

\begin{proof}[Outline]
Count the ideals \((x)J^{-1}\) among the divisors of \((a)(J\bar J)^{-1}\) and bound each fibre by the number of roots of unity using their cyclotomic degrees. Apply the estimate to both coordinates and use injectivity of the labels. The complete proof is in \S\ref{prf:B8.10n}.
\end{proof}

\begin{definition}[Corank-one partner classes and normal data]\label{B8.10da}
Fix words \(u,v\) in a tied Hida model \((X,Y)\) on the Hermitian \(E\)-space \((V,H_E)\) of dimension \(n=k+2r+1\), where \(r=r_S\ge1\), \(m=2r+2\), \(f=[F:\Q]\), and \(g=g_0g_1\). For a double-coset term \(t\) of \((u,v)\) with nonzero Gram invariant \(G(t)\ne0\), let \(\operatorname{row}(t)=(\Lambda_t,\Lambda'_t,\iota_t)\) be its complete row label, \(F_u\) the fixed flag of \(u\), and \(F'_t\) the moving flag of \(t\). Their \(2r+2\) vectors span a nondegenerate subspace \(W_t=\operatorname{span}_E(F_u,F'_t)\subset V\) whose orthogonal complement \(W_t^\perp\) is totally positive definite of rank \(k-1\).

At each finite place \(p\) of \(F\) of residue cardinality \(q_p\), let \(\theta_p\in E_p^\times\) be the local frame-difference scale and \(\Lambda_{t,p}\) the \(p\)-adic row lattice of \(t\). The local saturation index (defect) \(\beta_p(t)\ge1\) is the index inside \(\Lambda_{t,p}\cap W_{t,p}\) of the submodule generated by the first frame and the \(\theta_p^{-1}\)-scaled differences of the two frames. The term \(t\) is generic (\(t\in E_{\rm gen}\)) if \(\beta_p(t)=1\) at every finite place \(p\), and nongeneric (\(t\in E_{\rm ng}\)) if \(\beta_p(t)>1\) at some finite place \(p\).

Two terms \(t,t'\) with \(G(t)\ne0\) and \(G(t')\ne0\) in the same complete row have corank \(s(t,t')=1\) when the common fixed subspace of their relative rational isometry in \(W_t^\perp\cap W_{t'}^\perp\) has dimension \(k-1\), so \(W_t=W_{t'}\) and the relative isometry acts as the identity on \(W_t^\perp\). As in Lemma~\ref{B8.08n}, this equivalence relation partitions \(E_{\rm ng}\) into nongeneric corank-one partner classes \(c\); for \(u=v\), the generic scalar-partner classes \(c\) inside \(E_{\rm gen}\) are the classes of Lemma~\ref{B8.10c}. Each \(t\in c\) has a rational block point \(\gamma=\gamma(t)\) of level \(r+1\) over \(E\) with leading minors \(D_l(\gamma)\), nonzero integral Gram invariant \(G(\gamma)=D_{r+1}(\gamma)\overline{D_{r+1}(\gamma)}-T_uT_vD_r(\gamma)\overline{D_r(\gamma)}\), Gram norm factor \(N_G(\gamma)=\prod_\tau|\tau(G(\gamma))|\ge1\) over the CM type, block-end minor norm factor \(N_D(\gamma)=\prod_\tau\prod_j|\tau(D_j(\gamma))|^2\ge1\) so \(W(\gamma)=N_G(\gamma)N_D(\gamma)\) in Lemma~\ref{B8.10h}, and finite-place heights \(H_p(\gamma)\ge1\) whose product \(H(\gamma)=\prod_p H_p(\gamma)\) satisfies \(H(\gamma)\le g^fN_G(\gamma)N_D(\gamma)\). Within \((c,\gamma)\), a term is indexed by its normal datum \(\nu=\nu(t)\), whose two norm-one coordinates \(x\in J,x'\in J'\) (\(x\bar x=x'\bar x'=1\)) lie in fractional ideals of \(E\) attached to \((c,\gamma)\), with vector subtotal
\[
 V(c,\gamma,\nu)=\sum_{\substack{t\in c\\ \gamma(t)=\gamma,\ \nu(t)=\nu}}v_t.
\]
At each finite place \(p\), write \(d\in D_p\) for the local normal data at \(\gamma\) with orbit weights \(w_d\ge0\), unsigned local reference mass \(n_p\ge0\), generic local mass \(n_p^{\rm gen}\ge0\), global generic mass \(n_{\rm gen}(\gamma)=V_0\prod_p n_p^{\rm gen}\) as in Lemma~\ref{B8.08n}, and local orbital character sums \(a_N(d)\) (over presentations with \(\beta_p>1\)) and \(a_Z(d)\) (over presentations with \(\beta_p=1\)). Let \(\Omega\) be the finite set of exceptional places where \(\beta_p>1\) or \(n_p\ne n_p^{\rm gen}\) can occur, so \(H_p(\gamma)\ge q_p\ge2\) for \(p\in\Omega\).
\end{definition}

\begin{lemma}[Changing the normalization of Gram densities]\label{B8.10dba}
Let \(w\) range over the archimedean product place \(\infty\) and the finite places of a number field. At each \(w\), let \(V_w\) be the local tuple space, \(\mathcal H_w\) its space of Hermitian Gram matrices, \(q_w:V_w\to\mathcal H_w\) the Gram map, and \(\mathcal H_w^{\rm nd}\) the open subset of nonsingular matrices. Fix measures \(\mu_w^0,\nu_w^0\) on \(V_w,\mathcal H_w\), a test function \(\Phi_w:V_w\to\C\), and a continuous function \(f_w:\mathcal H_w^{\rm nd}\to\C\) such that, for every compact \(K\subset\mathcal H_w^{\rm nd}\),
\[
 \int_{q_w^{-1}(K)}\Phi_w\,d\mu_w^0
 =\int_K f_w(S)\,d\nu_w^0(S).
\]
Let \(c_w,e_w>0\), and put
\[
 \mu_w=c_w\mu_w^0,\qquad \nu_w=e_w\nu_w^0,
 \qquad a_w=c_w/e_w,\qquad g_w=a_w f_w.
\]
Then \(g_w\) is continuous on \(\mathcal H_w^{\rm nd}\), and
\[
 \int_{q_w^{-1}(K)}\Phi_w\,d\mu_w
 =\int_K g_w(S)\,d\nu_w(S).
\]
Fix a nonsingular rational Gram matrix \(T\), and write \(T_w\) for its image at \(w\). Suppose the finite-place products \(\prod_v a_v\) and \(\prod_v f_v(T_v)\) converge as nets over finite sets of places, and
\[
 a_\infty\prod_v a_v=1.
\]
Then the finite-place product of the \(g_v(T_v)\) converges and
\[
 g_\infty(T_\infty)\prod_v g_v(T_v)
 =f_\infty(T_\infty)\prod_v f_v(T_v).
\]
\end{lemma}
\begin{proof}[Outline]
Rescales each local integral by \(a_w=c_w/e_w\), preserving continuity. Factoring finite partial products of \(g_v(T_v)\) into \(a_v\) and \(f_v(T_v)\) products, taking the net limit, and applying \(a_\infty\prod_v a_v=1\) yields the global identity. The complete proof is in \S\ref{prf:B8.10dba}.
\end{proof}

\begin{lemma}[Siegel--Weil generic mass comparison]\label{B8.10db}
Under the hypotheses of Lemma~\ref{B8.10d}, there is a constant \(C'_0\ge0\), depending only on the six standing constants (before the thresholds for \(C',C_\ell\)), such that at each represented block point \(\gamma\) one has \(n_p\le8H_p(\gamma)^{C'_0(r+1)^2+2m}n_p^{\rm gen}\) for \(p\in\Omega\) and \(n_p=n_p^{\rm gen}\) for \(p\notin\Omega\). Write \(p(\gamma)=G_k(\gamma)t(\gamma)\) for the real point weight of Lemma~\ref{B8.10h}, factored into its Gram-power factor \(G_k(\gamma)\) and level multiplier \(t(\gamma)\), and write \(\mathrm{dens}(\gamma)\) for the product of normalized local Siegel--Weil representation densities at \(\gamma\). Then \(t(\gamma)^{-1}\le H(\gamma)^{6m}\), \(\mathrm{dens}(\gamma)\le e_1H(\gamma)\) with \(0<e_1\le2\), and the global generic mass \(n_{\rm gen}(\gamma)\) satisfies
\begin{equation}\label{eq:p8c1-generic-mass-comparison}
 n_{\rm gen}(\gamma)\le A_{\rm re}G_k(\gamma)\,\mathrm{dens}(\gamma)\,H(\gamma)^{10(r+1)^2}
 \le A_{\rm re}e_1H(\gamma)^{e_N}p(\gamma),\qquad e_N=10(r+1)^2+6m+1,
\end{equation}
where \(A_{\rm re}>0\) is the archimedean mass-to-section normalization.
\end{lemma}

\begin{proof}[Outline]
Uses \ref{B8.02ya} for archimedean uniqueness and energy convergence, and \ref{B8.10dba} to normalize Gram-fibre densities into the theta-mean formula for \(2n_0<M_0\). Applying it to \(W_t^\perp\) with local comparisons gives the mass bound. The complete proof is in \S\ref{prf:B8.10db}.
\end{proof}

\begin{lemma}[Normalization of the rotated orbital factors]\label{B8.10dca}
Use the rank parameters and scale \(b>1\) of Lemma~\ref{B8.10h}, and fix \(C_1\in\R\). Let \(\tau\) run over a CM type of cardinality \(f>0\). At each \(\tau\), fix a representative of a rotated block point, a Haar measure \(\mu_{G,\tau}\) on the full local unitary group containing the two frames, a root-coordinate measure \(\mu_{R,\tau}\) on its common stabilizer, and a quotient measure \(\nu_{\tau,1}\) which disintegrates \(\mu_{G,\tau}\) along \(\mu_{R,\tau}\). Thus, for every nonnegative measurable function \(F\),
\[
 \int F(x)\,d\mu_{G,\tau}(x)
 =\int\!\int F(sx)\,d\mu_{R,\tau}(s)\,d\nu_{\tau,1}(x).
\]
Write
\[
 \mathcal A_\tau(\nu)=\int \Phi_\tau(x)\,d\nu(x)
\]
for the local row-function orbital operator, with its fixed normalization included in the kernel \(\Phi_\tau\), which takes values in linear maps between finite-dimensional unitary spaces. Let \(Q(\gamma)>0\) be the squared global orbital factor at the represented row point \(\gamma\), for a fixed pinned quotient measure. Suppose
\[
 P(\gamma)=\prod_\tau\|\mathcal A_\tau(\nu_{\tau,1})\|_{\rm HS}^2>0,
 \qquad
 \vartheta=\left(\frac{P(\gamma)}{Q(\gamma)}\right)^{1/(2f)},
 \qquad b^{-C_1m^2}\le\vartheta\le b^{C_1m^2}.
\]
Then \(\nu_{\tau,\vartheta}=\vartheta^{-1}\nu_{\tau,1}\) is a quotient measure for the stabilizer measure \(\vartheta\mu_{R,\tau}\), with the prescribed normalization window, and
\begin{equation}\label{eq:p8c1-selected-orbital-factorization}
 Q(\gamma)=\prod_\tau\|\mathcal A_\tau(\nu_{\tau,\vartheta})\|_{\rm HS}^2.
\end{equation}
For a subtotal \(V(c,\gamma,\nu)\), define
\(a_2(c,\gamma,\nu)=\|V(c,\gamma,\nu)\|^2/Q(\gamma)\)
using the same pinned quotient measure. Then
\[
 \|V(c,\gamma,\nu)\|^2
 =a_2(c,\gamma,\nu)
   \prod_\tau\|\mathcal A_\tau(\nu_{\tau,\vartheta})\|_{\rm HS}^2.
\]
\end{lemma}
\begin{proof}[Outline]
Canceling \(\vartheta,\vartheta^{-1}>0\) in disintegration verifies \(\nu_{\tau,\vartheta}\). Linearity scales each squared norm by \(\vartheta^{-2}\), so multiplying over \(\tau\) via \(\vartheta^{2f}=P(\gamma)/Q(\gamma)\) and by \(a_2\) yields both identities. The complete proof is in \S\ref{prf:B8.10dca}.
\end{proof}

\begin{proposition}[Rotated archimedean reduction for corank-one subtotals]\label{B8.10dc}
Fix the six standing constants of Proposition~\ref{B8.10y} and any real weight parameter \(A_w\ge0\). For both pinned nongeneric corank-one classes of \((u,v)\) under the base and pair standing conditions of Definition~\ref{B8.02c}, and equal-word generic scalar corank-one classes of \((u,u)\) under the same base and pair standing conditions, there are constants \(C_a,C_n\in\R\), \(C_b\ge0\), and thresholds for \(C',C_\ell\) with the following normalized estimates. On the support where the global squared orbital factor \(Q(\gamma)\) is positive, use the same pinned quotient measure for \(Q\) and \(a_2(c,\gamma,\nu)=\|V(c,\gamma,\nu)\|^2/Q(\gamma)\). For the selected local row-function orbital operators of Lemma~\ref{B8.10dca},
\[
 Q(\gamma)=\prod_\tau\|\mathcal A_\tau(\nu_{\tau,\vartheta})\|_{\rm HS}^2,
 \qquad
 \|V(c,\gamma,\nu)\|^2
 \le a_2(c,\gamma,\nu)C_J\prod_\tau\|J_\tau(\gamma)\|_{\rm HS}^2,
 \qquad C_J=\prod_\tau\bigl(K_\tau b^{C_am^2}\bigr),
\]
where \(K_\tau>0\) are the local Bessel-comparison constants and \(\|J_\tau(\gamma)\|_{\rm HS}^2\) denotes the supremum of Lemma~\ref{B8.10h}.
At every represented point, including \(Q(\gamma)=0\), the analytic comparison is
\[
 Q(\gamma)\le C_J\prod_\tau\|J_\tau(\gamma)\|_{\rm HS}^2.
\]
If \(M_\tau\) is the local row-function section integral, obtained by using the squared row-function orbital operator and ambient Haar measure with the diagonal-coordinate density of Lemma~\ref{B8.10h}, the section comparison is
\[
 \left(\prod_\tau K_\tau\right)\prod_\tau I_\tau
 \le\prod_\tau M_\tau
 \le b^{C_nm^2f}I_{\rm re},
\]
where \(I_{\rm re}\ge0\) is the normalized row integral of \(X\). Lemma~\ref{B8.10h} applies at exponent \(A_w\), with these normalization losses absorbed into \(C_b\).
\end{proposition}

\begin{proof}[Outline]
For \(Q(\gamma)>0\), Proposition~\ref{B8.02ya} and Lemma~\ref{B8.10h} give unit-root measures, and Lemma~\ref{B8.10dca} with local Bessel bounds yields the orbital bound. Section bounds cancel \(K_\tau\) and absorb all remaining factors into \(C_b\). The complete proof is in \S\ref{prf:B8.10dc}.
\end{proof}

\begin{lemma}[The two signed corank-one sums]\label{B8.10d}
For a set \(U\) of term pairs put \(P(U)=\Gamma_1\sum_{(t,t')\in U}\Re\langle v_t,v_{t'}\rangle\). Assume the second formula of \ref{B8.13} and fix the six standing constants of \ref{B8.10y}. There is \(C_4\ge0\), followed by thresholds for \(C',C_\ell\), such that, under the size and split-prime hypotheses of \ref{B8.02c} for \(u,v\), the Hida-model compatibility specified before \eqref{eq:p8c1-loop-rows}, and absolute convergence,
\[
 P(U_{\rm ng})\le 2^{-(k-2r_S-5)}e^{C_4m^5f\log b}\Delta_{u,v}.
\]
Here \(U_{\rm ng}\) consists of same-row nongeneric corank-one pairs. The pair standing conditions of \ref{B8.02c} include \(r_S\ge1\), complete splitting of the primes \(T_u,T_v\ge3\) and of the distinct primes \(g_0,g_1\), \((T_uT_v,g_0g_1)=1\), and \(T_u\equiv T_v\pmod{g_0g_1}\).

For every fixed additional real counting parameter \(A\), there is \(C_B\ge0\), followed by thresholds for \(C',C_\ell\), such that \(\mathrm{Cnt.Hyp}(u,u;A)\) and the same ties and convergence give
\[
 P(U_{\rm sc})\le e^{C_Bm^5f\log b}\Delta_{u,u}
\]
for same-row generic scalar corank-one pairs of the equal-word operator.
\end{lemma}

\begin{proof}[Outline]
Rewrite the sums as class squares via Lemma~\ref{B8.08n} and bound local defect sums using a symmetric isometry kernel and Lemma~\ref{B8.10j}. Count normal labels by Lemma~\ref{B8.10n}, apply Cauchy--Schwarz, and combine Lemma~\ref{B8.10db} with Proposition~\ref{B8.10dc}. The complete proof is in \S\ref{prf:B8.10d}.
\end{proof}

\begin{lemma}[Cancellation of the two represented-frame determinants]\label{B8.10l}
Let \(k,f\ge1\), \(s\ge2\), and \(g,v,N>0\). Let \(G_1,G_2>0\) be the normalized Gram products of two represented frames. Put
\[
 n_i=\left(\frac{\exp(1)\pi}{kv}\right)^{kf}
          (G_i/N)^k .
\]
Let \(M\ge0\) be the weight of a fibre of pairs, \(A_i\ge0\) bounds for the two amplitudes, and \(W\ge0\) the remaining counting factor. For \(d\in\{1,2\}\) and \(B_i>0\), suppose
\[
 M\le W\left((2/g)^{(s-d)f}
        \frac{\sqrt{G_1G_2}}{N}B_1B_2\right)^k,\qquad
 n_i A_i^2 B_i^{2k}\le L_i\quad(i=1,2).
\]
Then
\[
 MA_1A_2\le
 W(2/g)^{(s-d)kf}
 \left(\frac{kv}{\exp(1)\pi}\right)^{kf}\sqrt{L_1L_2}.
\]
\end{lemma}

\begin{proof}
Multiply the amplitude inequalities and take the nonnegative square root. The factor on \(A_1A_2\) is
\[
 (B_1B_2)^k\sqrt{n_1n_2}
 =(B_1B_2)^k
 \left(\frac{\exp(1)\pi}{kv}\right)^{kf}
 \left(\frac{\sqrt{G_1G_2}}{N}\right)^k .
\]
Division by its positive value and substitution into the fibre bound give the assertion.
\end{proof}

\begin{lemma}[The saving on the small side of an ordered pair]\label{B8.10m}
Let \(T,T',g>0\), and let \(v\) range over \(f\) places. At each place let \(s_{iv}>0\) and \(c_{iv}=\sqrt{1+s_{iv}^2}\) for \(i=1,2\). Suppose
\[
       \left(\frac{\sqrt{TT'}}g\right)^f
       \le \prod_v\frac{c_{1v}+c_{2v}}{s_{1v}s_{2v}} .
\]
Then for at least one \(i\), putting \(\nu_v=T's_{iv}^2\), one has
\[
       \prod_v\min(\nu_v,T')\le (4g\sqrt{T'/T})^f .
\]
The same conclusion for the single available side holds if its \(s_v>0\) satisfy \((\sqrt{TT'}/g)^f\le\prod_v s_v^{-1}\).

For any side satisfying this conclusion let \(k\ge1,\ r\ge0,\ m=2r+2\) be integers. Let \(1\le C\le k^{1/3}\), \(0\le\eta\le1/4\), and let \(T_{{\rm crit},v}>0\) satisfy
\[
  g^{-m}T_{{\rm crit},v}\ge k^{3r},\qquad
  \prod_v(T_{{\rm crit},v}/T')^k\le 2C_0\Gamma^2
\]
for \(C_0,\Gamma>0\). Put \(a_v=T_{{\rm crit},v}/T'\). The local minimum factors satisfy
\[
 \left(\prod_v\min\{1,\nu_v a_v g^{-m}
                    (C/k)^{2r(1-\eta)}\}\right)^k
 \le 2C_0\Gamma^2
       \left(4\sqrt{T'/T}\,g^{1-m}k^{-r}\right)^{kf}.
\]
\end{lemma}

\begin{proof}[Outline]
Bound the product of the two clipped local sizes and take the smaller side, using \(\min(1,s)^2\le s\) for a single side. Insert this bound into the product of local minima using the two bounds on the critical scales. The complete proof is in \S\ref{prf:B8.10m}.
\end{proof}

\begin{definition}[Reduced Gram labels and switch coordinates]\label{B8.10ab}
In a tied Hida model \((X,Y)\) under the hypotheses of \eqref{eq:p8c1-loop-rows}, fix words \(u,u'\) with level primes \(T_u,T_{u'}\) in the Hermitian \(E\)-space \((V,H_E)\) of dimension \(n=k+2r+1\), where \(r=r_S=(m-2)/2\ge1\), \(m=2r+2\), \(f=[F:\Q]\), \(v=\operatorname{rd}(E)/2\), \(\rho=2v=\operatorname{rd}(E)\), \(g=g_0g_1\), \(L_b=\log b\), and \(N=T_u^f\). Each double-coset term \(t\) for \((u,u')\) with nonzero Gram invariant \(G(t)\ne0\) has complete row \(\operatorname{row}(t)=(\Lambda_t,\Lambda'_t,\iota_t)\), row weight \(\operatorname{wt}(t)=\bigl(w(\Lambda_t)w(\Lambda'_t)\bigr)^{1/2}\), orbital vector \(v_t\), fixed flag \(F_u\) of rank \(r+1\), moving flag \(F'_t\) of rank \(r+1\), and \((r+1)\times(r+1)\) reduced Gram matrix \(Z=Z(t)=(Z_{ab'})_{0\le a,b'\le r}\) over \(E\) with leading minors \(p_l(Z)=\det(Z_{ab'})_{0\le a,b'<l}\) (\(1\le l\le r+1\), \(p_0(Z)=1\)). At each place \(\tau\) of the CM type, set \(d_{l,\tau}(Z)=\rho^{-2q_l}|\tau(p_l(Z))|\) with \(q_l=\sum_{i=1}^l(r+1-i)\). For \(a\ge0\), define
\[
 S_a(Z)=\prod_\tau\prod_{l=1}^{r+1}\max\{1,d_{l,\tau}(Z)b^{-am^2}\},\qquad
 G_Z=g^{-mf}\prod_\tau\bigl(d_{r+1,\tau}(Z)^2-T_uT_{u'}d_{r,\tau}(Z)^2\bigr),
\]
so \(G_Z>0\) on represented support. For corank \(s\ge2\), define the represented-frame normalizer \(n_Z^{(2)}=\bigl(\exp(1)\pi/(kv)\bigr)^{kf}(G_Z/N)^k\), and put \(K(Z)=1\) if \(p_r(Z)=0\) and \(K(Z)=\prod_\tau(1+c_\tau)/\sqrt{c_\tau^2-1}\) with \(c_\tau=|\tau(p_{r+1}(Z)/p_r(Z))|/\sqrt{T_uT_{u'}}>1\) if \(p_r(Z)\ne0\). For any real \(c\ge\max(1,C_{Tu})\), define the switch height \(h_Z=b^{c^2m^3f}S_c(Z)^4\).

For a same-row pair \((t,t')\) with \(G(t)\ne0\) and \(G(t')\ne0\), let \(s(t,t')\in\{1,\ldots,k\}\) be the codimension in rank \(k\) of the common fixed subspace of their relative rational isometry in \(W_t^\perp\cap W_{t'}^\perp\), where \(W_t=\operatorname{span}_E(F_u,F'_t)\) has dimension \(m=2r+2\) and totally positive definite orthogonal complement \(W_t^\perp\) of dimension \(k-1\). For \(2\le s\le k\), partition the corank-\(s\) pairs of \(S_{\rm big}\) into the nondegenerate class \(\mathcal P_{{\rm nd},s}\) and the degenerate class \(\mathcal P_{{\rm deg},s}\) as in Lemma~\ref{B8.08k}. When \(s(t,t')=1\), the common fixed subspace equals \(W_t^\perp=W_{t'}^\perp\), so \(W_t=W_{t'}\) and the relative isometry acts as the identity on \(W_t^\perp\); since \((F_u,F'_t)\) is an \(E\)-basis of \(W_t\), the transition entries expressing the partner moving flag \(F'_{t'}\subset W_t\) in the basis \((F_u,F'_t)\) are called the switch coordinates of \(t'\) relative to \(t\).
\end{definition}

\begin{proposition}[Higher-corank loop pair sums]\label{B8.10ac}
Under the hypotheses and thresholds of \eqref{eq:p8c1-loop-rows}, use the notation of Definition~\ref{B8.10ab}, put \(f_w=d_w^S\), \(R'=f_{u'}/D^\star\), and \(\Gamma=e^{C_\Gamma m^5fL_b}\), and choose \(C\ge\max(1,4\sqrt{a_F(F_m.c_{\rm out})},C_{Tu})\) above the count and amplitude thresholds. Set \(H=b^{Cm^2f}\), let \(\Xi_s=M_sM_{k-s}/M_k\) for \(2\le s\le k-2\) with the standard definite genus masses \(M_j\) of Lemma~\ref{B8.02bc}, and at the endpoint ranks put
\begin{equation}\tag{8.25}\label{eq:p8c1-endpoint-index}
 \Xi_s=(Ck)^{kf}(H^{Cm}N)^{Ck},\qquad s=k-1,k.
\end{equation}
Define \(F_s=(kv)^{kf}\Xi_s C^{kf}H^{C(s+m)m^2}(Csf)^{6sf}\Gamma\), \(U_s=(2/g)^{(s-1)kf}F_s\), and \(V_s=(2/g)^{(s-2)kf}F_se^{Cm^6fL_b}(C/k)^{kf}\). Then for every \(2\le s\le k\) and every finite subset \(U\),
\begin{equation}\label{eq:p8c1-loop-nd-deg}
 \begin{aligned}
 \Gamma_1\sum_{(t,t')\in U}|\langle v_t,v_{t'}\rangle|
 &\le f_u R'U_s &&(U\subset\mathcal P_{{\rm nd},s}),\\
 \Gamma_1\sum_{(t,t')\in U}|\langle v_t,v_{t'}\rangle|
 &\le f_u V_s &&(U\subset\mathcal P_{{\rm deg},s}).
 \end{aligned}
\end{equation}
\end{proposition}

\begin{proof}[Outline]
Pairs are grouped by Gram labels \((Z,Z')\) and rank-\(s\) lattice, combining fibre counts with amplitude bounds via Lemma~\ref{B8.10l}. A high-order moment \(q\) gives summable label decay, and Lemma~\ref{B8.08oa} completes both bounds. The complete proof is in \S\ref{prf:B8.10ac}.
\end{proof}

\begin{lemma}[The single-side product estimate]\label{B8.10ada}
Let \(T,T',g>0\), and let \(s_\tau>0\) be indexed by a set of \(f\) places. If
\[
 \left(\frac{\sqrt{TT'}}g\right)^f\prod_\tau s_\tau\le1,
\]
then
\[
 \prod_\tau\min(T's_\tau^2,T')
 \le\bigl(g\sqrt{T'/T}\bigr)^f
 \le\bigl(4g\sqrt{T'/T}\bigr)^f.
\]
\end{lemma}
\begin{proof}
For \(0<s\le1\), one has \(T's^2\le T's\); for \(s\ge1\), one has \(T'\le T's\). Thus \(\min(T's^2,T')\le T's\). Put \(q=\sqrt{TT'}/g\), so \(T'=g\sqrt{T'/T}\,q\). Multiplication gives
\[
 \prod_\tau\min(T's_\tau^2,T')
 \le(T')^f\prod_\tau s_\tau
 =\bigl(g\sqrt{T'/T}\bigr)^f q^f\prod_\tau s_\tau
 \le\bigl(g\sqrt{T'/T}\bigr)^f.
\]
The last inequality in the statement follows from positivity.
\end{proof}

\begin{proposition}[Represented-frame genus mean and switch partner bounds]\label{B8.10ad}
Under the hypotheses and thresholds of Proposition~\ref{B8.10ac}, there exist a constant \(c\ge\max(1,C_{Tu})\), determined before the family and point thresholds from \(C_{Tu}\) and the rank-independent constants of the switch and partner estimates, and positive local represented-frame factors \(\nu_\tau(T_u,T_{u'},Z)>0\) at the places \(\tau\) of the CM type, with product \(n_Z=\prod_\tau\nu_\tau(T_u,T_{u'},Z)>0\) and switch height \(h_Z=b^{c^2m^3f}S_c(Z)^4\), with the following properties.
\begin{enumerate}
\item On a row family of total mass \(M\), for every represented reduced Gram label \(Z\), the inverse-stabilizer-weighted sum of inequivalent first frames of label \(Z\) (the represented-frame genus mean) is at most \(M h_Z^{cr^2}n_Z\).
\item For a fixed first term \(t\) of label \(Z\), the switch coordinates of Definition~\ref{B8.10ab} map the corank-one partners \(t'\) of label \(Z'\) in the same row injectively into integral partner data of cardinality at most \(c^f(h_Zh_{Z'})^{cr^2}\). Consequently, every finite fibre of same-row corank-one pairs of labels \((Z,Z')\) has weighted cardinality bounded in the two orders by
\begin{equation}\tag{8.26}\label{eq:p8c1-switch-count}
 M h_Z^{c r^2}n_Z\,c^f(h_Zh_{Z'})^{c r^2},
 \qquad M h_{Z'}^{c r^2}n_{Z'}\,c^f(h_{Z'}h_Z)^{c r^2},
\end{equation}
and hence by their geometric mean \(M c^f(h_Zh_{Z'})^{3cr^2/2}\sqrt{n_Zn_{Z'}}\).
\item Each corank-one term of label \(Z\) has squared orbital amplitude per unit row weight and coefficient dimension carrying the inverse factor \(n_Z^{-1}\), and on every nonscalar pair \((t,t')\in S_{\rm nsc}\) at least one side satisfies the small-side saving of Lemma~\ref{B8.10m}. Enlarging \(C\) above the threshold determined by \(c\) and setting \(J=C^fH^{Cm^3}\Gamma e^{Cm^6fL_b}(Cg^{-(2r+1)}k^{-r})^{kf/2}\), every finite subset \(U\subset S_{\rm nsc}\) satisfies
\begin{equation}\label{eq:p8c1-loop-one}
 \Gamma_1\sum_{(t,t')\in U}|\langle v_t,v_{t'}\rangle|
 \le\max(1,C_0)f_u R' J.
\end{equation}
\end{enumerate}
\end{proposition}

\begin{proof}[Outline]
Gluing rank-\((k-1)\) complements via Lemma~\ref{B8.10c} and bounding switch coordinates gives the geometric-mean pair count. Splitting \(U\) by Lemma~\ref{B8.10m} cancels \(n_Z,n_{Z'}\), and high moments with Lemma~\ref{B8.08k} bound the sum. The complete proof is in \S\ref{prf:B8.10ad}.
\end{proof}

\begin{lemma}[Summing the loop pairs]\label{B8.10a}
Under the hypotheses and thresholds of \eqref{eq:p8c1-loop-rows},
\(\operatorname{PairAbs}_{u,u'}(S_{\rm big}\cup S_{\rm nsc})\le e^{-P/3}\Delta_{u,u'}\), where \(P=kf\log k\).
\end{lemma}

\begin{proof}[Outline]
Partition \(S_{\rm big}\cup S_{\rm nsc}\) and apply Propositions~\ref{B8.10ac} and~\ref{B8.10ad} with Definition~\ref{B8.02c}. Rescale \(\Gamma\) to \(\gamma\), bound mass factors via Lemma~\ref{B8.02bc} and rank thresholds, and sum the decaying bounds over \(2\le s\le k\). The complete proof is in \S\ref{prf:B8.10a}.
\end{proof}

\begin{lemma}[The two flag lattices and their nonzero frequencies]\label{B8.10aa}
Let \(E\) be the CM field, \(f=[E:\Q]/2,\ g=g_0g_1\), and \(\mathfrak d\) the absolute different of \(E\). For a block of length \(s\) starting at \(b_0\) put \(\nu=2r+3-2b_0-s\). For fractional ideals \(A,B\) write \([A,B]_s=\{t\in A:t-(-1)^s\bar t\in B\}\). The period lattice, inverse-difference lattice and character dual of the period lattice are respectively
\[
 L_{b_0s}=[g\mathfrak d^\nu,g^2\mathfrak d^\nu]_s,\quad
 L'_{b_0s}=[g^{-1}\mathfrak d^{-\nu},\mathfrak d^{-\nu}]_s,\quad
 L_{b_0s}^*=\{l\in E:\operatorname{Tr}_{E/\Q}(lt/g)\in\Z\ (t\in L_{b_0s})\}.
\]
Suppose \(J\) is a nonzero fractional ideal, \(L'\) an additive subgroup of \(J\), and \(z_v\ge0\) with \(\prod_vz_v\le \sqrt{N(J)}(a/\Theta)^f,\ a>0,\ \Theta>2a\). Here \(N(J)\) is the absolute ideal norm and \(v\) runs through one embedding of each conjugate pair. With an absolute constant \(C\), every finite \(S\subset L'\setminus\{0\}\) satisfies
\begin{equation}\label{eq:p8c1-lattice-sum}
 \sum_{l\in S}\prod_v\min\{1,(z_v/|v(l)|)^5\}\le [C(a/\Theta)^{5/2}]^f .
\end{equation}
\end{lemma}

\begin{proof}[Outline]
Use ideal norms to separate points in dyadic boxes, subdivide into cells with at most one point, and sum the fifth-power bounds. Identify the dual by the trace character and bound its ideal using the containment of \(g^2\mathfrak d^\nu\) in the period lattice. The complete proof is in \S\ref{prf:B8.10aa}.
\end{proof}

\begin{lemma}[The oscillatory integral at a merging pair of blocks]\label{B8.10g}
Let \(s=n+1\ge2\), and let \(J_j\) be the matrix with ones on the reverse diagonal in size \(j\). For \(\beta\in\C^\times\) put
\[
 \psi_\beta(z)=\exp(i\Re(\beta z)),\qquad
 d\mu_\beta(z)=|\beta|(2\pi)^{-1}\,d\Re z\,d\Im z .
\]
Use this measure on each independent complex entry. For a strictly upper triangular matrix \(Y\) put \(\ell(Y)=\sum_i Y_{i,i+1}\). Let \(\Phi\) be a smooth compactly supported function on the complex \(s\)-by-\(s\) matrices with values in a complex Banach space. Define
\[
 \begin{split}
 L_\Phi(z)&=\int\psi_\beta(\ell(Y))\Phi\bigl(zJ_s(1+Y)\bigr)\,dY,\\
 B_\Phi(t,c)&=\int\psi_\beta(a_n+\ell(V)+b_n)
 \Phi\left(
 \begin{pmatrix}1&0\\a&1\end{pmatrix}
 \begin{pmatrix}tJ_n&0\\0&c\end{pmatrix}
 \begin{pmatrix}1+V&b\\0&1\end{pmatrix}
 \right)\,da\,dV\,db .
 \end{split}
\]
Here \(a\) is a row of length \(n\), \(b\) a column of length \(n\), \(V\) is strictly upper triangular of size \(n\), and \(Y\) is strictly upper triangular of size \(s\). Entries are indexed from one, so \((aJ_n)_1=a_n\).

For \(D=ct^n\) in a fixed compact subset of \(\C^\times\), uniformly as \(t\to0\),
\[
 |t|^{n^2}B_\Phi(t,c)
 =\frac1s\sum_{z^s=(-1)^nD}
 |z|^{n^2}\psi_\beta(sz/t)L_\Phi(z)+O_{\Phi,D}(|t|).
\]
The error constant can be chosen for the entire compact set.
\end{lemma}

\begin{proof}[Outline]
Use the reversed matrix determinant to eliminate the first diagonal entry and locate critical points at \(z^s=(-1)^nD\). Integrate by parts away from the roots and use stationary phase with the determinant derivative and triangular change of variables to get \(L_\Phi(z)\). The complete proof is in \S\ref{prf:B8.10g}.
\end{proof}

\begin{lemma}[Boundary pairings in the inverse entry]\label{B8.10i}
Let \(s\ge2\), \(|\omega|=1\), and let \(\chi\) be a nontrivial continuous additive character of \(\C\). Use planar Lebesgue measure. For each of finitely many prefixes \(p\), let \(B_p(z)\) be a measurable Hilbert-space vector such that
\(b_p(z)=s|z|^{s-1}B_p(z)\) is square integrable on compact subsets of \(\C^\times\). Set
\[
 A_{p,x}(D)=\sum_{z^s=\omega D}\chi(szx)B_p(z).
\]
Suppose \(\sup_{x\ne0}\|U_{p,x}\|_2^2\le M_p<\infty\) and
\(\|U_{p,x}-A_{p,x}\|_{L^2(K)}\to0\) on each compact annulus \(K\) as \(|x|\to\infty\). To compare fibres at different entries use the unitary change of variable
\[
 (T_x V)(c)=|x|^{-(s-1)} V(c/x^{s-1}).
\]
Then \(\|b_p\|_2^2\le M_p\). For \(x'=x+\sigma\), each of the three pairings
\[
 \langle T_x A_{p,x},T_{x'} A_{p',x'}\rangle,\quad
 \langle T_x U_{p,x},T_{x'} A_{p',x'}\rangle,\quad
 \langle T_x A_{p,x},T_{x'} U_{p',x'}\rangle
\]
tends to
\[
 \int_\C\chi(\sigma z)\langle b_p(z),b_{p'}(z)\rangle\,dz.
\]
The inner products are linear in the second argument. The original and boundary fibres together have a positive semidefinite block Gram matrix.
\end{lemma}

\begin{proof}[Outline]
Apply Fourier decay to distinct-root terms in squared annular norms to get global square integrability. Use domination and continuity of dilations to pass to boundary fibres, keep matching-root frequencies, and realize both families in one Hilbert space for positivity. The complete proof is in \S\ref{prf:B8.10i}.
\end{proof}

\begin{lemma}[Regularity of the truncated group coefficient]\label{B8.10p}
Let \(G=\U(p,q)\) in coordinates with Hermitian matrix \(\operatorname{diag}(1_p,-1_q)\). Put \(\|g\|_F^2=\sum_{ij}|g_{ij}|^2\). Let \(\theta\) be a smooth real function equal to zero on \((-\infty,0]\) and to one on \([1,\infty)\). For a smooth coefficient \(F\) with values in a finite-dimensional normed space put
\[
 F_n(g)=\theta\left(\frac{4-\|g\|_F^2/(n+1)^2}{3}\right)F(g)
 \qquad(n\in\N).
\]
Then \(F_n=F\) for \(\|g\|_F\le n+1\), and \(F_n\) has compact support in \(\|g\|_F\le2(n+1)\). Fix \(n\), an integer \(d\ge0\) and a real \(A\). The quantities \(\|g\|_F^A\|L_{\mathbf u}R_{\mathbf v}F_n(g)\|\) are uniformly bounded over \(g\in G\) and all words of left and right invariant derivatives of total length at most \(d\) whose directions have norm at most one.
\end{lemma}

\begin{proof}[Outline]
Checking the cutoff argument gives \(F_n=F\) on the inner ball and support in the intersection of \(G\) with the outer closed ball, which is compact and avoids \(0\). Continuity in \(g\), multilinearity in the directions, and compactness yield the uniform bound. The complete proof is in \S\ref{prf:B8.10p}.
\end{proof}

\begin{lemma}[The distinct-word scalar contribution]\label{B8.10e}
Under the source and size hypotheses of \eqref{eq:p8c1-distinct-row}, the signed generic scalar corank-one sum for distinct words \(u,v\) is at most \(A\Delta_{u,v}\), where
\[
 A=(1+\delta_s)e_1,\qquad
 0\le\delta_s\le b^{-7},\qquad 1\le e_1\le\exp(12f\,2^{-(k-2r_S-2)}).
\]
\end{lemma}

\begin{proof}[Outline]
Cut at a differing bit and use \ref{B8.10c} to bound row squares by flag energies. Iterate the energy bound using interpolation in \ref{B8.10h}, lattice sums in \ref{B8.10aa} and limits from \ref{B8.10g} and \ref{B8.10i}, then convert by formal degree. The complete proof is in \S\ref{prf:B8.10e}.
\end{proof}

\begin{lemma}[The finite-type matrix]\label{B8.10f}
Under the selected local hypotheses of \ref{B8.10}, the finite-type Gram matrix is at most \(P_g\Gamma_1 I\) in positive semidefinite order, where \(P_g=\prod_{j=0}^1(1+3/g_j)^f\). For a cell \(E_1\) with absolute summability and the norm expansion, put \(w_i(E_1)=\sum_{t\in E_1:\,\iota(\mathrm{row}(t))=i}v_t\), where \(\iota\) gives the finite-type index of a row. Then
\[
 \|\operatorname{part}(E_1)\|_{\rm HS}^2\le P_g\Gamma_1\sum_i\|w_i(E_1)\|^2 .
\]
\end{lemma}

\begin{proof}[Outline]
Diagonalize the local Whittaker convolution using \ref{B7.23}, \ref{B7.24}, \ref{B7.25}, \ref{B7.29} and multiplicity one from \ref{B7.26}. Bound eigenvalues by last-column restrictions, tensor the matrix bounds, and pass to the limit by absolute summability. The complete proof is in \S\ref{prf:B8.10f}.
\end{proof}

\begin{lemma}[The scalar row and the full kernel]\label{B8.10b}
Under the source and size hypotheses of \eqref{eq:p8c1-distinct-row}, distinct words \(u,v\) satisfy the bound \((1+b^{-6})\Delta_{u,v}\) for the squared Hilbert--Schmidt norms of \(K_0\) and the full normalized Gram operator.
\end{lemma}

\begin{proof}[Outline]
Unfold absolutely by \ref{B8.02y} and discard the identity and zero-Gram terms by local pairings, irrelevance, and Schur's lemma. Use \ref{B8.10f} and row orthogonality, then combine bounds from \ref{B8.10e}, \ref{B8.10a}, and \ref{B8.10d} at the size thresholds. The complete proof is in \S\ref{prf:B8.10b}.
\end{proof}





% ---- p8c2 ----
\chapter{Theta correspondence, test vectors and evaluation}\label{ch:p8c2}

Let \(E/F\) be a CM extension with conjugation \(\sigma\) and adele rings \(\mathbb A_E,\mathbb A_F\). For a group \(H\) over \(F\) write \([H]=H(F)\backslash H(\mathbb A_F)\) and let \(\operatorname{avg}_{[H]}\) be integration divided by the quotient volume. Write \(\mathscr P\) for the product \(F\) of local theta hypotheses in Definition~\ref{B1.15ac}, and \(\mathscr J\) for the conjunction of the five toric assertions in Definition~\ref{B1.14k}. At scale \(x>1\) put \(L=x^\varepsilon=\log b\), with \(0<\varepsilon<1/52\).

\section{Matching theta integrals}

Let \(V,V'\) be Hermitian and \(W\) skew Hermitian, of dimensions \(a,a',a+a'\) and respective signatures \((a-p,p),(q,a'-q),(a+a'-p-q,p+q)\) at every infinite place. Require
\begin{equation}\label{eq:p8c2-triple-range}
 a'<a,\qquad a+a'\equiv0\pmod2,\qquad
 2(p+q)+1<\frac{a-a'}2,\qquad 2(p+q)+1<a'.
\end{equation}
Take the prescribed Cartan involutions and self-dual lattices with split frames. Assume \(E/\Q\) is Galois and \(E/F\) is unramified at finite places. The additive character \(\psi\) of \(\mathbb A_F/F\) has infinite part \(y\mapsto\exp(-2\pi i\operatorname{Tr}_{F_\infty/\R}y)\). Put \(T=S_0\sqcup Q\subseteq P\), where \(P\) consists of the prescribed split places with chosen embeddings. At inert places the residue characteristic is odd and \(Nv>3\), with \(Nv\) the residue cardinality of \(F_v\). At inert places and on \(P\), \(\psi_v\) has conductor \(\mathcal O_v\). Fix \(v_1\notin T\) above a rational prime split completely in \(E\).

The unitary Hecke characters \(\alpha,\beta\) restrict to \(\eta^{a'},\eta^a\), where \(\eta\) is quadratic for \(E/F\). Their angular exponents satisfy \(m_{\alpha,w}-m_{\beta,w}>a+a'\) at every infinite place \(w\) of \(F\). Both characters are unramified off \(T\), and \(\nu=\alpha/\beta\) is ramified at both split places of \(E\) above every \(v\in T\) (so its local conductor exponent satisfies \(c_v\ge1\) for every \(v\in T\), and in particular for every \(v\in S_0\subseteq T\)). Thus \(\nu|_{\mathbb A_F^\times}=1\) and \(\nu(\bar t)=\nu(t)^{-1}\). The Weil models \(\omega,\omega'\) on \(V\otimes_E W,V'\otimes_E W\) use the splitting pairs \((\beta,\beta/\alpha),(\alpha,\alpha/\beta)\).

For a split embedding \(t:E\hookrightarrow F_v\), let \(U_t\) be the set of primed Schwartz vectors whose image in some allowed barred model for \(t\) vanishes on matrices of less than full row rank. For the chosen embedding \(t_1\) at \(v_1\), let \(\mathcal S'_{t_1}\) be the span of restricted pure tensors whose \(v_1\)-factor belongs to \(U_{t_1}\). Put
\[
 I(\Phi,h)=\int_{[\U(V)]}\theta_\omega(g,h;\Phi)\,dg
\]
and define \(I'\) with \(V',\omega'\). The theta kernels sum the transformed tests over rational points of the Schr\"odinger spaces. Real polynomial vectors are the span of Heisenberg derivatives of Gaussians.

Matching maps \(p_\infty,p'_\infty\) on real polynomial vectors and \(p_v,p'_v\) on finite Schwartz vectors are surjections onto common nonzero irreducible quotients \(\Pi_\infty,\Pi_v\). They intertwine the target action, at infinity through its Lie algebra and maximal compact. They kill the finite source differences and the real source Lie action, and they are invariant under the real source maximal compact. Off \(T\) the two vacua have the same nonzero image. At \(v\in T\) use split and barred model maps \(\iota,\iota'\) sending vacua to integral-matrix indicators. Put \(D=a+a'\), \(X_0=(I_a\ 0)\), \(Y_0=(0\ I_{a'})\), and \(w(g)=\nu_v(\det g)|\det g|_v^{D/2}\) at the chosen embedding. Normalize \(K_j=\GL_j(\mathcal O_v)\) to volume one. In these coordinates the matching images agree when \(P(\iota\phi)=P'(\iota'\phi')\), where the two operators are given by
\begin{equation}\label{eq:p8c2-split-matching}
\begin{aligned}
 P(\phi)(h)
 &=\beta_v(\det h)|\det h|_v^{-a/2}
 \sum_{\mathbf n}^{\mathrm{fin}}
 \int_{K_a\varpi^{\mathbf n}K_a}\phi(gX_0h^{-t})w(g)\,dg,\\
 P'(\phi')(h)
 &=\alpha_v(\det h)|\det h|_v^{a'/2}
 \int_{\GL_{a'}(F_v)}\phi'(gY_0h)w(g)\,dg .
\end{aligned}
\end{equation}
Here \(h\in\GL_D(F_v)\), \(\varpi\) is a uniformizer, and the first operator uses the finite-support sum over Cartan double cosets indexed by weakly decreasing integer tuples.

A polynomial match has equal images under the tensors of the finite and real maps. The finite tensor includes \(T\) and every non-vacuum place. General matches are limits in the real Schwartz topology of polynomial matches with fixed finite factors and primed vectors in \(\mathcal S'_{t_1}\).

\begin{lemma}[Hecke nonvanishing at the required integers]\label{lem:p8c2-hecke-nonvanishing}\label{B8.11h}
Let \(\chi\) be a unitary Hecke character of a CM field \(E\). At each infinite place \(v\) let \(m_v\) be its angular exponent, with local factor \(\Gamma_\C(s+|m_v|/2)\), and suppose that some \(m_v\ne0\). Then \(L(k,\chi)\) is finite and nonzero at every integer \(k\ge1\), and at every integer \(k\le0\) for which \(k+|m_v|/2>0\) at all infinite places. At \(k\ge1\) one may omit any finite set of finite Euler factors.
\end{lemma}

\begin{proof}[Outline]
Use the Euler product for \(k>1\) and a nonnegative Taylor series to contradict a zero at \(1\). Tate's functional equation transfers nonvanishing to the allowed \(k\le0\), and omitting Euler factors multiplies the values at \(k\ge1\) by nonzero factors. The complete proof is in \S\ref{prf:B8.11h}.
\end{proof}

Fix a real place \(w\). Put
\[
 P=a-p,\quad Q=p,\quad P'=q,\quad Q'=a'-q,\qquad
 R=P+Q',\quad S=Q+P',\quad D=R+S.
\]
Thus the target has signature \((R,S)\). Write
\[
 c=\frac{m_{\alpha,w}-m_{\beta,w}}2,\qquad
 n=c-\frac{R-S}{2},\qquad m=c+\frac{R-S}{2}.
\]
Evaluating the restrictions of \(\alpha,\beta\) at the idele which is \(-1\) at \(w\) and \(1\) elsewhere gives
\(m_{\alpha,w}\equiv a'\) and \(m_{\beta,w}\equiv a\pmod 2\).
Since \(D\) is even, \(c\) is integral. The inequality \(2c>D\) gives \(n>S\) and \(m>R\).

For a source of signature \((b,e)\), use polynomial coordinates \(z_{ij}\) with \(b+e\) rows and \(D\) columns, in signature order. Denote this polynomial space by \(\mathcal F_{b,e}\). For a coordinate \(l=(i,j)\), define two operators
\[
 (A_l,B_l)=
 \begin{cases}
 (z_l,\partial_l),&\text{if the row and column have equal signs},\\
 (-\partial_l,z_l),&\text{if their signs differ}.
 \end{cases}
 \qquad H_{lt}=A_lB_t+\tfrac12\delta_{lt}.
\]
Here a variable in an operator formula means multiplication by that variable. The source and target matrix units act, before the character twists, by
\[
 D_1(E_{ih})=\sum_j H_{(i,j),(h,j)},\qquad
 D_2(E_{jk})=\sum_i H_{(i,j),(i,k)}.
\]
These actions commute, as follows from \([B_l,A_t]=\delta_{lt}\). Write \(\Theta_{b,e}(h)\) for the quotient of \(\mathcal F_{b,e}\) by the span of
\((D_1(Z)-h\operatorname{tr}Z)f\), for all source matrices \(Z\) and polynomials \(f\). Write \(t^{[j]}\) for a string of \(j\) coordinates equal to \(t\).

\begin{lemma}[Generators of the two real theta modules]\label{B8.11i}
Give \(T=\Theta_{P,Q}(c)\) the target action \(D_2\), and give \(T'=\Theta_{P',Q'}(-c)\) the target action \(D_2+c\operatorname{tr}\). Both modules are nonzero. Each is generated by an irreducible compact type of highest weight
\[
 \lambda=
 \big((d+n)^{[P]},d^{[Q']}\,;\,
             (m-d)^{[Q]},(-d)^{[P']}\big),
 \qquad d=(P-Q)/2.
\]
The compact Lie algebra here is \(\mathfrak k=\mathfrak{gl}_R(\C)\oplus\mathfrak{gl}_S(\C)\), with upper triangular positive roots in both blocks.
\end{lemma}
\begin{proof}[Outline]
Use a Fischer pairing to show that the determinant products survive the source relations. Joint harmonics give the compact types, \ref{B1.15b} gives generation, and the vacuum shifts and trace twist identify the highest weights. The complete proof is in \S\ref{prf:B8.11i}.
\end{proof}

\begin{lemma}[The first noncompact operators]\label{B8.11j}
In either module of Lemma~\ref{B8.11i}, let \(L\) denote the generating compact type and let \(\mathfrak p\) be the two off-diagonal blocks of \(\mathfrak{gl}_D(\C)\) relative to \(R,S\). Let \(C_K=\sum E_{ij}E_{ji}\), summed over ordered pairs in the same block. If \(C_K\) acts on \(L\) by \(b_0\), every eigenvalue on \(\mathfrak p L=\operatorname{span}\{Yv:Y\in\mathfrak p,\ v\in L\}\) is strictly greater than \(b_0+D\).
\end{lemma}
\begin{proof}[Outline]
Use the minimum degree from \ref{B8.11i} to represent \(\mathfrak pL\) by joint harmonics of degree \(d_0+2\). The possible weight changes increase the compact Casimir eigenvalue by \(D+2(\lambda_i-\lambda_j)>D\). The complete proof is in \S\ref{prf:B8.11j}.
\end{proof}

\begin{lemma}[A common irreducible real quotient]\label{B8.11k}
For the actions of Lemma~\ref{B8.11i}, the modules \(T,T'\) admit surjections onto a common nonzero irreducible \(\mathfrak{gl}_D(\C)\)-module, intertwining every target matrix unit.
\end{lemma}
\begin{proof}[Outline]
By \ref{B8.11i}, the module generated by the graph of the compact-type identification projects onto both modules. Filter by off-block degree and use \ref{B8.11j} to show that the graph is a full Casimir eigenspace, then quotient by the sum of proper submodules. The complete proof is in \S\ref{prf:B8.11k}.
\end{proof}

\begin{lemma}[Nonzero definite theta integrals]\label{B8.11l}
Assume Theorem~\ref{B8.12}, and keep the triple and Weil models of \eqref{eq:p8c2-triple-range} with \(p=q=0\). There exist a real-polynomial Schwartz vector \(\Phi\) on the unprimed side with \(I(\Phi,\cdot)\ne0\) and a real-polynomial Schwartz vector \(\Phi'\in\mathcal S'_{t_1}\) on the primed side with \(I'(\Phi',\cdot)\ne0\).
\end{lemma}
\begin{proof}[Outline]
Averages a real polynomial from \ref{B8.11i} with congruence indicators at \(T\) and doubles against the constant function to relate the squared theta norm to primitive Hecke \(L\)-values. Nonvanishing via \ref{lem:p8c2-hecke-nonvanishing} yields a witness on each side. The complete proof is in \S\ref{prf:B8.11l}.
\end{proof}

\begin{lemma}[Constancy of the reverse lift]\label{B8.11m}
Keep the triple and Weil models of \eqref{eq:p8c2-triple-range}, let \((S,\omega_S)\) be \((V,\omega)\) or \((V',\omega')\), and let \(P\) be a Schwartz vector on \((S\otimes_E W)(\mathbb A_F)\). Let \(u\) be a continuous automorphic function on \([\U(W)]\) such that \(h\mapsto u(h)\overline{\theta_{\omega_S}(x,h;P)}\) is integrable on \([\U(W)]\) for every \(x\in\U(S)(\mathbb A_F)\) and
\[
 T_Pu(x)=\int_{[\U(W)]}u(h)\overline{\theta_{\omega_S}(x,h;P)}\,dh
\]
is continuous on \([\U(S)]\). Suppose there is a finite set \(E_{\rm exc}\) of finite places of \(F\) containing \(T\cup\{v_1\}\) such that at every \(v\notin E_{\rm exc}\) above a rational prime split completely in \(E\), \(u=\sum_j U_j(\xi_j)\) for finitely many \(\xi_j\in\Pi_v\) and \(\U(W)(F_v)\)-equivariant linear maps \(U_j:\Pi_v\to R_v\) from the local quotient \(\Pi_v\) at \(v\) of \eqref{eq:p8c2-split-matching} and Lemma~\ref{B8.11f} into a subspace \(R_v\) of continuous functions on \([\U(W)]\) integrable against \(\overline{\theta_{\omega_S}(x,\cdot;P)}\) for all \(x\in\U(S)(\mathbb A_F)\). Then \(T_Pu\) is constant on \(\U(S)(\mathbb A_F)\).
\end{lemma}
\begin{proof}[Outline]
At good split places \(v\notin E_{\rm exc}\), covariant matrix distributions vanish on singular strata by stabilizer-twist mismatch, giving \(\U(S)(F_v)\)-invariance of \(T_Pu\). Strong approximation and Hecke \(L\)-function holomorphy at \(1\) force constancy. The complete proof is in \S\ref{prf:B8.11m}.
\end{proof}

\begin{theorem}[Matching and proportional theta integrals]\label{B8.11}
Assume Theorem~\ref{B8.12}. Matching maps exist for every triple above. Fix such maps. The integral \(I\) converges on Schwartz vectors. For \(g\in\U(V')(\mathbb A_F)\), let \(H_{\rm quot}(g)=\inf_{\gamma\in\U(V')(F)}\operatorname{ht}(\gamma g)\), using adelic height. For every \(\Phi'\in\mathcal S'_{t_1}\) and \(N\in\N\) there are \(C_N\in\R,\ k_N\in\N\) such that
\begin{equation}\label{eq:p8c2-full-rank-decay}
 |\theta_{\omega'}(g,h;\Phi')|H_{\rm quot}(g)^N
 \le C_N\operatorname{ht}(h)^{k_N}\qquad(g,h\text{ arbitrary}).
\end{equation}
The integral \(I'\) converges with moderate growth on this space and is nonzero for some test there. There is a unique \(c\in\C^\times\) such that
\begin{equation}\label{eq:p8c2-matched-integrals}
 I(\Phi,h)=cI'(\Phi',h)
\end{equation}
for every match and every \(h\).

At \(v\notin T\) above a prime split completely in \(E\), the kernels are the local source differences and \(p'_v(U_t)=\Pi_v\) for each \(t:E\hookrightarrow F_v\). A Schwartz functional killing the primed relations factors through \(p'_v\). The primed Weil action preserves \(\mathcal S'_{t_1}\). The sign triples in the step formulas below belong to \(\{\pm1\}^3\) and are uniform in the base fields.
\end{theorem}

For fractional ideals \(\mathfrak b,\mathfrak c\), signs \(e_1,e_2,s\) and \(z\in\Z\), put
\begin{equation}\label{eq:p8c2-step-factor}
 \mathfrak f(e_1,e_2,s;\mathfrak b,\mathfrak c;z)
 =\nu(\mathfrak b)^{e_1}\nu(\mathfrak c)^{e_2}
 \prod_{v\in T}\nu_v(s)\epsilon_v(\psi,\nu_v,z).
\end{equation}
Use the prescribed ideal character and \(\nu_v\) at the chosen embedding. For uniformizer \(\varpi_v\) and conductor exponent \(c_v\), set
\[
 \epsilon_v(\psi,\nu_v,z)
 =(Nv)^{-c_vz}\nu_v(\varpi_v)^{c_v}
 \sum_{u\in(\mathcal O_v/\varpi_v^{c_v})^\times}
 \nu_v(u)^{-1}\psi_v(u/\varpi_v^{c_v}).
\]
In the \(\beta\)-step write \(a'=b_0+2,\ \ell=(a-b_0-2)/2\). Keep \(V\) and remove adapted planes from \(V',W\) at frame positions \((0,b_0+1),(a,a+b_0+1)\). The lower indices are \((p,q-1)\). For lower comparison constant \(c_1\),
\begin{equation}\label{eq:p8c2-beta-step}
 cL^{S_0}(\nu,\ell+1)
 =\frac{c_1}{C\,\mathfrak f(e_{\beta,1},e_{\beta,2},s_\beta;\mathfrak b,\mathfrak c;-\ell)}.
\end{equation}
In the \(\alpha\)-step write \(q=0,\ a=b_0+2,\ \ell=(b_0+2-a')/2\). Keep \(V'\) and remove planes from \(V,W\) at \((0,b_0+1)\) on each side. The lower indices are \((p-1,0)\), and
\begin{equation}\label{eq:p8c2-alpha-step}
 c=C\,\mathfrak f(e_{\alpha,1},e_{\alpha,2},s_\alpha;\mathfrak b,\mathfrak c;1-\ell)
 L^{S_0}(\nu,\ell)c_1 .
\end{equation}
For the plane vectors \(x,x'\) in Lemma~\ref{B8.11b} use the fractional ideals \(\mathfrak b=\{s\in E:sx\in\Lambda_{\rm source}\}\) and \(\mathfrak c=\{s\in E:sx'\in\Lambda_W\}\). The \(L\)-function \(L^{S_0}(\nu,\cdot)\) omits places above \(S_0\); because \(\nu\) is ramified above every \(v\in T\supseteq S_0\), every primitive Euler factor above \(S_0\) equals \(1\), so \(L^{S_0}(\nu,s)=L_E(\nu,s)\) identically. The fixed models, measures, domains, planes and real maps determine \(C\ne0\) before the characters, Weil actions and finite matching maps are chosen.



\begin{proof}[Outline]
Builds local matching maps via \ref{B8.11i}, \ref{B8.11k}, and \ref{B1.15ba}. Inducts on \(p+q\), proving \(p=q=0\) via \ref{B8.11l}, \ref{B8.11m}, and Petersson proportionality, and positive \(p+q\) via Levi reductions and two-step contraction. The complete proof is in \S\ref{prf:B8.11}.
\end{proof}

\begin{lemma}[The Petersson norm of a single theta integral]\label{B8.11e}
Keep the triple preceding Theorem~\ref{B8.11}. For any Haar measure and rational fundamental domain on the target, \(I(\Phi,\cdot)\) is square integrable for polynomial tests, and \(I'(\Phi',\cdot)\) is square integrable for polynomial tests in \(\mathcal S'_{t_1}\). Both integrals converge absolutely and are continuous. At fixed allowed finite factors they depend continuously on the real Schwartz factor, locally uniformly in the target. The primed kernel satisfies \eqref{eq:p8c2-full-rank-decay} and the primed integral has moderate growth.
\end{lemma}

\begin{proof}[Outline]
Apply Poisson summation and count on the Siegel sets of \ref{B1.15b}. Pair tests with their reflections so full-rank minors give chamber decay, then integrate against chamber densities for convergence and square integrability. The complete proof is in \S\ref{prf:B8.11e}.
\end{proof}

\begin{lemma}[Source differences on split matrix spaces]\label{B8.11f}
In either barred split realization put \(k=F_v,\ d=a+a'\), and let \(n\) be the source dimension. Let \(\mathcal S=C_c^\infty(\operatorname{Mat}_{n,d}(k))\) and let \(\mathcal S^\circ\) be supported on rank \(n\). Write \(A(U)\) for the span of \(\omega(g,1)\phi-\phi\) with \(\phi\in U,\ g\in\GL_n(k)\). If the character \(\zeta\) in the action
\[
 \phi(Y)\longmapsto\zeta(\det g)|\det g|^{d/2}\phi({}^tgY)
\]
is nontrivial, then the natural maps are isomorphisms
\[
 \mathcal S^\circ/A(\mathcal S^\circ)\longrightarrow
 \mathcal S/A(\mathcal S)\longrightarrow\Pi_v .
\]
The second map is induced by the matching map. At a completely split good place the character is nontrivial for either embedding. At the selected primed place \(\mathcal S^\circ\) corresponds to \(\operatorname{span} U_{t_1}\).
\end{lemma}

\begin{proof}[Outline]
Row averaging and scaling give full-rank representatives. Weighted fibre integration realizes the coinvariants as an induction irreducible by \ref{B8.43h}, and the Hecke law at split primes gives nontriviality of the good-place character. The complete proof is in \S\ref{prf:B8.11f}.
\end{proof}

\begin{lemma}[Descent of the two integrals]\label{B8.11g}
Fix any matching maps for the triple preceding Theorem~\ref{B8.11}. An unprimed polynomial test matching zero has zero integral. A primed polynomial test matching zero has zero integral if it belongs to \(\mathcal S'_{t_1}\). On these domains each integral is determined by its image in the common restricted tensor product.
\end{lemma}

\begin{proof}[Outline]
Use the Petersson radical from \ref{B8.11e}, split coinvariants from \ref{B8.11f}, and compact averaging with \ref{B8.14a}. At infinity extend orthogonality from the generating type of \ref{B1.15b} by Lie adjointness, then descend factor by factor. The complete proof is in \S\ref{prf:B8.11g}.
\end{proof}

\begin{lemma}[Real vanishing after zero evaluation]\label{lem:p8c2-real-vanishing}
Use compatible models over the same good base, with source and target signatures \((r,s),(r',s')\) and target dimension \(d_1=r'+s'\). Let \(m_w\) be the source exponent at \(w\). If \(|m_w|>d_1\) everywhere at infinity and \(\min(r',s')<\min(r,s)\), the source integral vanishes on full-rank tests. If also \(d_1+2\min(r,s)<2(r+s)\), it vanishes on all Schwartz tests.
\end{lemma}

\begin{proof}
Fix the target point and finite factors. Equations \eqref{eq:p8c2-theta-moderate-density} and \eqref{eq:p8c2-theta-negative-sums} give the continuous invariant Schwartz functional \(\mathcal L\) in the two cases. In one Fock factor differentiation gives
\[
 \mathcal L(d\omega_{\rm pol}(Y)P)
 =-\tfrac12m_w\operatorname{tr}(Y)\mathcal L(P),
\]
where the polynomial action excludes the splitting scalar. For a rank-one \(Y=vw^t\) of trace one in one complexified compact source block, its derivative part on a matching target sign column \(c\) is \(\sum_{a,b}v_aw_bX_{ac}\partial_{bc}\). On the other columns it is \(-\sum_{a,b}v_aw_bX_{bc}\partial_{ac}\). The scalar part is half the difference of the column counts. The required derivative weight has sign opposite to \(m_w\). Negate the operator and exchange the rank-one vectors if necessary to make that weight positive. Choose the source block with \(\min(r',s')\) positive-action columns.

For a linear polynomial \(\ell\), choose a nonzero covector annihilating those columns of \(\ell\), and a vector pairing to one with it. The source block dimension exceeds the number of columns. The derivative now kills the positive-column part of \(\ell\); on the others it sends its rank-one projection \(\ell_-\) to \(-\ell_-\) and kills \(\ell-\ell_-\). The binomial expansion of \(\ell^j\) has nonpositive derivative weights, so each term is killed by \(\mathcal L\). Polarization and Schwartz continuity give the result.
\end{proof}

\begin{lemma}[Doubling a matrix coefficient in the definite case]\label{B8.11d}
Keep the matched models of Theorem~\ref{B8.11} with \(p=q=0\). Write \(G=\U(W), H'=\U(V')\), and use the Petersson product linear in its first argument. Give \(G\) any Haar measure and rational fundamental domain. Let \(Y\) be the polynomial matches \(y=(\Phi,\Phi')\) with \(\Phi'\in\mathcal S'_{t_1}\), and put \(J_y=I'(\Phi',\cdot)\), \(F_y=sI(\Phi,\cdot)+tJ_y\) for fixed \(s,t\in\C\). If
\[
 \langle F_y,F_z\rangle_{[G]}=\lambda\langle J_y,J_z\rangle_{[G]}
 \qquad(y,z\in Y),
\]
then for each polynomial primed test \(P\) and \(y\in Y\),
\[
 \|T_PF_y\|_{[H']}^2=\lambda\|T_PJ_y\|_{[H']}^2,\qquad
 T_Pf(x)=\int_{[G]}f(h)\overline{\theta_{\omega'}(x,h;P)}\,dh .
\]
\end{lemma}

\begin{proof}[Outline]
Unfold the doubled pairing into matrix coefficients and use Petersson proportionality to compare them. Continue to the definite Siegel–Weil point by \ref{B8.12}, where Fubini gives the squared reverse-lift norms with a common measure factor. The complete proof is in \S\ref{prf:B8.11d}.
\end{proof}

\begin{lemma}[Removing a hyperbolic plane]\label{B8.11a}
For \(q>0\) the \(\beta\)-planes can be made adapted by choosing Cartan involutions and split frames. Their complements have indices \((p,q-1)\). For \(q=0<p\) the \(\alpha\)-step gives \((p-1,0)\). Removal on the \(V\)-side is possible for every \(p>0\) with \(q\) fixed.
\end{lemma}
\begin{proof}[Outline]
Use Hasse–Minkowski to find rational isotropic lines and lattice duality with a trace-one correction to split off isotropic planes. Choose compatible frames and Cartan involutions and check the resulting indices, sizes and character restrictions. The complete proof is in \S\ref{prf:B8.11a}.
\end{proof}

\begin{lemma}[The two Levi identities]\label{B8.11b}
Fix an adapted step and the upper and lower real maps. Assume polynomial comparison and nonzero polynomial witnesses on both sides at smaller index sums, with the primed witness in \(\mathcal S'_{t_1}\). Let \(c_1\) be the lower comparison scalar. Write \(A+B,A'+B'\) for the constant terms along the removed target line, in order of the characters \(\alpha|\cdot|_E^{a'/2},\beta|\cdot|_E^{a/2}\) on \(\mathbb A_E^\times\). Put \(z=\ell+1\) for \(\beta\) and \(z=\ell\) for \(\alpha\), and set \(S=\mathfrak f(e_1,e_2,s;\mathfrak b,\mathfrak c;1-z)L^{S_0}(\nu,z)\) with uniform signs. There is \(C\ne0\) fixed before the characters, Weil actions and finite maps such that at each Levi lift \(m\) of \((1,h_1)\) from the lower group,
\[
 \begin{array}{ll}
 \beta:& CSB(m)=c_1 B'(m)\quad\text{for polynomial matches with }\Phi'\in\mathcal S'_{t_1},\\
 \alpha:& B'=0,\quad A(m)=CS c_1 A'(m)\quad\text{for every polynomial match}.
 \end{array}
\]
Some unprimed polynomial test has \(B(m)\ne0\) in the \(\beta\)-step. Some polynomial match has \(A'(m)\ne0\) in the \(\alpha\)-step.
\end{lemma}

\begin{proof}[Outline]
Unfolds the line constant term into zero evaluation and a mixed-polarization cone integral, keeping full-rank support at \(v_1\) via \ref{B8.11f}. Local factorization computes Euler and Gauss factors on standard tests; restriction uniqueness and lower comparison extend them. The complete proof is in \S\ref{prf:B8.11b}.
\end{proof}

\begin{lemma}[A line test for theta-integral differences]\label{B8.11c}
Set
\[
 \Sigma=\{I(\Phi)-I'(\Phi'):\Phi,\Phi'\text{ polynomial},\ \Phi'\in\mathcal S'_{t_1}\}.
\]
For nonzero totally isotropic rational \(U\subset W\), let \(C_U\) be the average over the radical of its stabilizer. If \(f\in\Sigma\) has zero constant term on the line of one rational hyperbolic plane, then \(C_Uf=0\) for every such \(U\). If also \(p+q>0\), then \(f=0\).
\end{lemma}

\begin{proof}[Outline]
Inducts on \(\dim_E U\), killing Whittaker coefficients via shear stabilizers (\(j\ge3\)) and smoothed \ref{B8.11e} bounds (\(j=2\)). For \(p+q>0\), shows \(f\) is orthogonal to both integrals via \ref{B8.11f}, \ref{B8.11g}, \ref{B8.11m}, and \ref{B8.11l}. The complete proof is in \S\ref{prf:B8.11c}.
\end{proof}

\section{The integral formulas used in the construction}

\begin{theorem}[Doubling continuation and the definite Siegel--Weil value]\label{B8.12}
Let \(Y\) be Hermitian or skew Hermitian of dimension \(n\ge1\) over the good base, with common signature and the prescribed Cartan data, self-dual lattice and split frames. Fix a doubling frame, with Iwasawa height \(a\) and standard maximal compact \(K\). Let \(\chi\) be a continuous Hecke character of \(\mathbb A_E^\times\), trivial on \(E^\times\), with \(\chi|_{\mathbb A_F^\times}=\eta^c,\ c\in\N\). For continuous right \(K\)-finite \(f\) satisfying \(f(px)=\chi(\det t)|\det t|_E^{s_1+n/2}f(x)\) on the Siegel parabolic with Levi coordinate \(t\), the series
\[
 E(s,x)=\sum_\gamma a(\gamma x)^{s-s_1}f(\gamma x)
\]
over rational Siegel-parabolic cosets is holomorphic in a right half-plane. Near each \(s_0\) some \(k\in\N\) gives an extension of \((s-s_0)^k E(s,\cdot)\) to continuous moderate-growth functions. Evaluation by a linear functional bounded on each class \(|h(x)|\le\operatorname{ht}(x)^N\) is holomorphic there.

Let \(A,Y\) have opposite Hermitian signs and dimensions \(m,n\ge1\) with the stated model conditions. Suppose \(A\) has common definite signature. Take \(\chi|_{\mathbb A_F^\times}=\eta^m\) and compatible doubling data with functional \(l\). For \(K\)-finite Schwartz \(\Phi\), the section \(x\mapsto l(\omega(x)\Phi)\) at \(s_0=(m-n)/2\) has holomorphic Eisenstein value
\begin{equation}\label{eq:p8c2-definite-sw}
 E(s_0,x;\Phi)=
 \begin{cases}
 \operatorname{avg}_{[\U(A)]}\theta(x,g;\Phi),&n<m,\\
 2\operatorname{avg}_{[\U(A)]}\theta(x,g;\Phi),&n\ge m
 \end{cases}
\end{equation}
for any Haar measure and fundamental domain.
\end{theorem}

The first assertion restricts Langlands's continuation theorem to unitary doubling sections \cite{langlands-1976}. The second is the anisotropic Siegel--Weil identity restricted to definite \(A\), with quotient-average normalization \cite[Theorem~7.1 and \S3.1]{gan-qiu-takeda-2014}.

\begin{theorem}[Convergent Siegel--Weil and Rallis formulas]\label{B8.13}
For the first and third assertions let \(A,B\) be Hermitian of dimensions \(a,b\) over the CM extension \(E/F\) of the same arithmetic point, and let \(r_B\) be the rational Witt index of \(B\). For the first assume
\begin{equation}\label{eq:p8c2-convergent-range}
 a<b\le2a,\qquad a+r_B<b.
\end{equation}
Use compatible Iwasawa and Cartan data with doubled Weil characters \(1,\chi_B\) and rational additive covolume one. For any Haar measure on \(\U(B)(\mathbb A_F)\) and allowed Fock vector \(\Phi\), the descended kernel is integrable at every doubled group point \(g\). Its quotient average is the holomorphic value \(E(g,f_\Phi,s_0)\) of the flat section at \(s_0=(b-a)/2\).

Over a CM extension take nonsingular Hermitian matrices \(J,T\) of sizes \(M,n\) with \(2n<M\). The quotient \([\U(J)]\) has finite volume. Take \(\Phi=\Phi_\infty\prod_{v<\infty}\Phi_v\) with \(\Phi_\infty\) smooth compactly supported and finite factors Schwartz--Bruhat, equal almost everywhere to the integral-tuple indicator in a fixed \(F\)-basis of \(E\). There are admissible Haar measures and fundamental domains and nontrivial characters of \(\mathbb A_F/F\) with self-dual measures on tuples and Hermitian matrices. For every such choice the Gram pushforwards have continuous densities \(g_\infty,g_v\) on the nonsingular locus, with
\begin{equation}\label{eq:p8c2-density}
 \operatorname{avg}_{[\U(J)]}
 \sum_{\operatorname{Gram}(X)=T}\Phi(g^{-1}X)
 =g_\infty(T)\prod_{v<\infty}g_v(T)
\end{equation}
and convergent product.

For the inner product require \(a<b\le2a\) and either \(r_B=0\) or \(a+r_B<b\). Use compatible Weil models \(\omega,\omega_d\) with Iwasawa and covolume-one data, right invariant Haar measures on the source, and a prescribed real realization and Cartan frame. Let \(\pi\) be irreducible cuspidal on \(\U(A)\). Its convergent lifts to the prescribed lower signatures obtained by removing any \(i\ge1\) hyperbolic planes from \(B\) must vanish for the compatible models and specified smooth forms. For allowed tests \(\Phi_1,\Phi_2\) require the doubled test
\[
 \delta(\Phi_1,\Phi_2)(x)=
 \Phi_1(\operatorname{res}_L x)\overline{\Phi_2(-\operatorname{res}_R x)}
\]
to be \(K\)-finite, with the two restrictions indicated. For \(f_1,f_2\) in the specified smooth cuspidal realization, the integrals
\(\theta_{\Phi_i}f_i(h)=\int_{[\U(A)]}\theta_\omega(g,h;\Phi_i)\overline{f_i(g)}\,dg\)
converge in \(L^2([\U(B)])\), and
\begin{equation}\label{eq:p8c2-rallis}
 \operatorname{CT}_{s=s_0} Z_{\chi_B^{-1}}(s,\delta(\Phi_1,\Phi_2),f_1,f_2)
 =\operatorname{vol}[\U(B)]^{-1}
 \langle\theta_{\Phi_1}f_1,\theta_{\Phi_2}f_2\rangle_{[\U(B)]}.
\end{equation}
Here \(\operatorname{CT}\) is the constant Laurent coefficient and \(Z_\mu\) uses \(\mu(\det g_2)^{-1}\) at the second group argument.
\end{theorem}

The first formula specializes Yamana's Theorem~2 to \(\delta=1,\ m=b,\ n=a,\ r=r_B\) and the stated rank range and tests \cite[p.~237]{yamana-2013}. The second restricts Weil's orbital formula to the indicated nonsingular unitary fibres with the measure of \S44, Theorem~2 \cite[\S52, Theorem~5, equation~(41)]{weil-1965}. The third is the convergent-range Rallis identity with lower-lift vanishing in quotient-average normalization \cite[\S3.1 and Propositions~11.1(ii) and~11.2]{gan-qiu-takeda-2014}.

For compact-finite doubled tests, approximate at fixed finite level in Schwartz seminorms and project onto the given compact types. The restrictions to \(K\) lie in a fixed finite-dimensional matrix-coefficient space and determine the flat sections by Iwasawa decomposition. Integration over \(\U(B)\) is continuous by the calculation of \eqref{eq:p8c2-cone-exponents} with tuple dimension \(2a\). Its partial exponents are
\begin{equation}\label{eq:p8c2-doubled-cone}
 2fj(a-b+j)<0\qquad(1\le j\le r_B),\qquad f=[F:\Q].
\end{equation}
At \(r_B=0\) the cone has no variables. The doubled self-pairing and this cone bound give
\begin{equation}\label{eq:p8c2-rallis-kernel-norm}
 \|\theta_\omega(g,\cdot;\Phi_i)\|_{L^2([\U(B)])}\le C_i\operatorname{ht}(g)^{N_i}.
\end{equation}
Use Cauchy--Schwarz on the target and rapid decay on the source for Fubini \cite{harish-chandra-1968}. The factor outside the doubled kernel is \(\chi_B(\det g_2)\), supplied by the indicated twist. The quotient average supplies the volume.

\section{Vacuum vectors and separation}

\begin{lemma}[Unramified vacuum generation]\label{B8.14}
Let \(L_{\rm loc}\) be a complete discretely valued field of characteristic zero with integers \(\mathcal O\), uniformizer \(\varpi\) and residue cardinality \(q^2\) for odd \(q\ge2\). Its involution \(\sigma\) preserves valuation and fixes \(\varpi\), and some \(u\in\mathcal O\) has \(\sigma u-u\in\mathcal O^\times\). Take Hermitian \(J_S\in\GL_n(\mathcal O), J_T\in\GL_{n'}(\mathcal O)\), and \(\delta\in\mathcal O^\times\) with \(\sigma\delta=-\delta\). Let the additive character \(\psi:L_{\rm loc}\to\C^\times\) be trivial on \(\mathcal O^\sigma\) and nontrivial on \(\varpi^{-1}\mathcal O^\sigma\).

On a complex space \(M\) take operators \(\rho(x)\) for \(x\in\operatorname{Mat}_{n'\times n}(L_{\rm loc})\) and a representation \(\omega\) of \(\U(J_S)\times\U(J_T)\). Write \(x^*={}^t\!\sigma(x)\) and \(H(x,y)=\delta\operatorname{tr}(J_S^{-1}y^*J_Tx)\). Assume
\[
 \rho(0)=1,\quad
 \rho(x)\rho(y)=\psi((H(x,y)+\sigma H(x,y))/2)\rho(x+y),\quad
 \omega(g,h)\rho(x)=\rho(hxg^{-1})\omega(g,h).
\]
The integral groups \(K_S,K_T\) fix the vectors fixed by all integral translations. Suppose one such vector \(s\ne0\) has translates spanning \(M\). Then
\[
 M^{K_S}\subseteq\operatorname{span}_\C\{\omega(1,h)s:h\in\U(J_T)\},\qquad
 M^{K_T}\subseteq\operatorname{span}_\C\{\omega(g,1)s:g\in\U(J_S)\}.
\]
\end{lemma}

\begin{proof}[Outline]
Decompose into lattice character lines and use stabilizer phases to find integral Gram representatives. Isolate the required congruences by averaging vacuum translates, use Witt extension and induction on lattice index for spanning, and apply the adjoint model. The complete proof is in \S\ref{prf:B8.14}.
\end{proof}

\begin{definition}[Pointwise size inequalities]\label{def:p8c2-point-sizes}
Use the scale and host of Proposition~\ref{B3.50}. Write \(n=f/2\), \(r=k+m-1\), \(M=N(m)=k+m(t+1)\), and put
\[
 \theta=\varepsilon/100,\quad u=x^\theta,\quad
 \varepsilon_z=\varepsilon/J,\quad \delta_1=x^{-\varepsilon/8},
 \quad \delta_2=|T_2|/n,\quad A=2n\delta_1 .
\]
Here \(J\ge1\) is the scale parameter, \(T_2\) is the thick set, and \(TT\) is the set of top positions. Put \(X'=\exp(x^{\varepsilon/3})\), \(b'=xb/X'\), \(w''=b'/(xn)\), \(X''=xb/w''\). Write \(\Lambda_G\) for the price parameter, \(\eta\) for the weight-window width, \(\Delta\) for the node separation parameter, and \(S_c=4k\Delta\). The constant \(C_\Gamma\) normalizes the member modulus by \(\log\Gamma_P=C_\Gamma m^5fL\). Set
\[
 A_1=\frac{kL}{m\delta_2},\qquad
 A_*=\frac{mb}{nt^2(M+t)L^2+kn^2(S_c+t+L)}.
\]
The pointwise size hypothesis in Proposition~\ref{B3.50} consists of the following inequalities. In the first display \(y\in[z/u,uz]\) means both endpoint inequalities.
\[
\begin{gathered}
 n\in[x^{2+8\varepsilon}/u,ux^{2+8\varepsilon}],\qquad
 m\in[\sqrt{x\delta_1}/u,u\sqrt{x\delta_1}],\qquad
 k\in[x^{9/2+10\varepsilon}/u,ux^{9/2+10\varepsilon}],\\
 \max(um^5fL^2,8C_\Gamma m^5 fL,m^2L,m\sqrt x\log x)\le k,\qquad
 k\le\sqrt x\,n^2x^{-\varepsilon},\qquad
 \sqrt{k/m}\le nx^{-\varepsilon},\\
 um^6L\le k\log k,\qquad um^3L\log\operatorname{rd}(E)\le k(\log k)^2,\qquad
 \log f\le\sqrt{m/2-1},\qquad(\log f)^2\le m^2L,\\
 u\log\log b\le rf\log r,\qquad uL\le f\log k,\qquad
 f\le\min(m^2n,r^2),\qquad 1\le t\le x^{C_t+\theta},\\
 (Q_Sqd)^{36}t^{150}\le\sqrt b .
\end{gathered}
\]
For the host primes, the top weights \(b_s\), and the quadratic character \(\chi_d\) of \(\Q(\sqrt{-d})\), require
\[
\begin{gathered}
 \log q\ge x^{\varepsilon_z/6},\qquad
 x^{\varepsilon_z/5-\varepsilon_z/100}\le\min(\log q,\log d),\qquad
 \max(\log q,\log d)\le x^{\varepsilon/5+\theta},\\
 \left|\Re\frac{L'(1,\chi_d)}{L(1,\chi_d)}\right|\le L,\qquad
 1\le\operatorname{rd}(E)\le Q_S\sqrt{qd}/2,\qquad
 n\delta_1/u\le|TT|\le A,\qquad L\le\Lambda_G\le uL,\\
 (4AL+u\log x)\Lambda_G/k\le x^{-2},\qquad
 b_s\le b'+AL\Lambda_G b'/k\quad(s\in TT),\qquad
 x^{-\varepsilon}\le\delta_2,\qquad L_w\le(\log x)^2 .
\end{gathered}
\]
Here \(L_w\) is the host word-length parameter. The width and spacing inequalities are
\[
\begin{gathered}
 t\le\eta\le b/(t^3fk^3),\qquad \log\eta\le x^{\varepsilon/4},\qquad
 t\eta\le w''/x,\qquad m(t+1)\eta\le w'',\qquad m(S_c+\eta)\le w''L,\\
 M\eta\le\delta_2mw'',\qquad \eta^2rL\le\delta_2m(w'')^2,\qquad
 r(b'+S_c)(\log x)A_1\le\delta_2mb,\\
 \max\{12b/A_*+2t+8,\ ktL,\ ur^{19}f^5L,\ 4M^3+2t+6,\
 ukm^{9/2},\ 2(r+mt)^3+t+3,\ Mm^{10}\}\le\eta/m,\\
 \max\{b'+S_c+\eta+b/A_*,\ 2\eta+t+8b/A_*+7\}<A_*,\\
 \max(Mm^{10},umM)\le\Delta,\qquad \log\Delta\le x^{\varepsilon/4}/2,\qquad
 \max(kS_c,umk\Delta,\sqrt b)\le w'',\\
 4kS_c\le\Lambda_G b',\qquad 4uk^2\Delta\le\Lambda_G b',\qquad
 um\log(\operatorname{rd}(E)/2)\le k .
\end{gathered}
\]
At every infinite position \(\tau\) require \(b_{\tau,m-1}\le2b_{\tau,0}\); for every two layer weights \(w_{j,\tau},w_{j',\tau'}\), \(1\le j,j'\le m\), require \(w_{j,\tau}\le X''w_{j',\tau'}\). For \(i,i'\in\{0,1\}\) require \((kb)^8\le g_i\le2g_{i'}\) and \(\log g_i\le uL\). Both \(\log p_{\rm de}\) and \(\log p'_{\rm de}\) are at most \(L\). There are at most \(2^{k-3}\) words, and the host cutoff up to which every odd prime splits in \(\Q(\sqrt{-d})\) is at least \(11\).
\end{definition}

\begin{corollary}[Two averages of an inert test]\label{B8.14a}
At an inert place of a triple preceding Theorem~\ref{B8.11}, let \(H,G\) be the source and target, \(K_H,K_G\) their integral groups, \(\mathcal S\) the local model and \(s\) its vacuum. Put \(A=\operatorname{span}_\C\{hx-x:h\in H,\ x\in\mathcal S\}\). Then \(\mathcal S=A+\operatorname{span}_\C Gs\). For each \(x\in\mathcal S\) there are \(k_1,\ldots,k_b\in K_G,\ b\ge1,\ c\in\C\) with \(\sum_i k_i x-cs\in A\), and \(k_i s=s\).
\end{corollary}

\begin{proof}[Outline]
Verify the unimodular frame, unramified character, and vacuum properties at the inert place, then apply \ref{B8.14} to both finite averages. The source average keeps \(x\bmod A\) in \(\operatorname{span}_{\mathbb{C}}Gs\); the target average is a scalar multiple of \(s\bmod A\). The complete proof is in \S\ref{prf:B8.14a}.
\end{proof}

For a real realization \(R\), let \(\mathcal V\) be the tensor product of the smooth archimedean representation spaces of the bottom row. Let \(\mathcal E_0\) be the subspace of its automorphic multiplicity maps at the trivial level index satisfying the prescribed level character and fixed by the type idempotent. Write \(t(w)\) for evaluation under tensor linearization. Let \(H\) be the top definite group, embedded by \(\iota\) in the last row group. Let \(N,\psi_N\) be the Bessel unipotent and character. For real test \(\rho\) and grid test \(\Phi\), let \(F_{\rho,\Phi,u}\) be the successive lift of \(u\). For coefficient action \(\rho_{\rm coef}\), dual \(\xi\) and column \(\Psi\), put \(\phi_{\xi,\Psi}(h)=\xi(\rho_{\rm coef}(h^{-1})\Psi(h_f))\). The period is
\[
 \mathcal B_{\rho,\xi,\Psi,\Phi}(u)
 =\int_{[H]}\int_{[N]}\phi_{\xi,\Psi}(h)
 F_{\rho,\Phi,u}(n\iota(h))\overline{\psi_N(n)}\,dn\,dh.
\]
The row size conditions comprise the pointwise size inequalities of Definition~\ref{def:p8c2-point-sizes}, the existence of compatible parameters with weight-gap constant \(10\), and the following additional bounds. Write \(r=k+m-1\) for the bottom row rank, \(f=[F:\Q]\), \(N(m)=k+m(t+1)\) for the last chain rank, \(S_c\) for the center bound and \(g_{\rm sep}\) for the prescribed grid spacing at gap parameter one. The upper bound for ranks across the whole grid is \(N_*=k+m(t+r+7)=k+m(t+k+m+6)\), where \(t\) is the auxiliary rank. The numerical bounds are
\begin{gather}
 m^4\le f,\quad 44\le m,\quad
 x^2m^3L\log\operatorname{rd}(E)\le k,\quad fmN(m)^2d\le b,
 \label{eq:p8c2-size-rows-a}\\
 fm\le(r+1)^3,\quad \log g_i\le9L+1,\quad \log q_0\le L,
 \label{eq:p8c2-size-rows-b}\\
 \log\operatorname{rd}(E)\le\log\operatorname{cond}(E)\le L/100,\quad
 S_c+(m+1)g_{\rm sep}+\tfrac12\le b .
 \label{eq:p8c2-size-rows-c}
\end{gather}
The conductor ideal condition is \(\mathfrak c_K=\mathcal O_K\). At each infinite place and weight position \(l\), \(b_{\tau,l}>0,\ e_{\tau,l}\le4b\).

\begin{lemma}[Reconstruction at a fixed Heisenberg level]\label{B8.15d}
Let \(S,T\) be two spaces in a grid Weil pair over \(E/F\). Put
\(X=\operatorname{Hom}_E(S,T)\), and write \(X_{\mathbb A}\) for its adelic points. Let \([x,u]\) be the alternating form of the pair and \(\psi_0\) its standard additive character of sign \(-1\). Write
\[
 H(u)f(x)=\psi_0([x,u]/2)f(x+u).
\]
A theta test in this model is a continuous function on \(X_{\mathbb A}\), smooth at infinity, fixed by \(H(U)\) for some open finite adelic lattice \(U\), and satisfying
\[
 f(x+\gamma)=\psi_0(-[\gamma,x]/2)f(x)
 \qquad(\gamma\in X(F)).
\]
Here and below \(X(F)\) denotes the rational points of the underlying \(F\)-space. Use Haar measure for which \(X(F)\) has covolume one. The pairing \((f,g)\) is the integral of \(f\overline g\) on \(X(F)\backslash X_{\mathbb A}\).

For a finite adelic lattice \(\Gamma\subset X_{\mathbb A}\), let \(D(\Gamma)\) consist of the theta tests fixed by \(H(2v)\) for every \(v\in\Gamma\), and by \(H(v)\) whenever \(v\in\Gamma\) is supported on one factor above a split finite place of \(F\). Let \(\Gamma^\circ\) consist of the \(u\in X_{\mathbb A}\) such that \(\psi_0([u,2v])=1\) for the first set of \(v\) and \(\psi_0([u,v])=1\) for the second.

Suppose that \(L\) is linear on theta tests, that \(u\mapsto L(H(u)f)\) is continuous and bounded for each test \(f\), and that
\[
 L\left(\int a(u)H(u)f\,d\nu(u)\right)
   =\int a(u)L(H(u)f)\,d\nu(u)
\]
whenever \(\nu\) is a countable sum of finite measures, \(a\in L^1(\nu)\), and the pointwise integral is a theta test. For every nonzero \(v_0\in D(\Gamma)\) and every \(f\in D(\Gamma)\),
\[
 (v_0,v_0)L(f)
   =\int_{\Gamma^\circ}(f,H(u)v_0)L(H(u)v_0)\,du .
\]
\end{lemma}

\begin{proof}[Outline]
Use the Heisenberg commutator for support and differentiation for decay. Fourier expansion and unfolding recover \(f\) up to the vacuum norm squared from its vacuum matrix coefficients, and superposition gives the formula for \(L\). The complete proof is in \S\ref{prf:B8.15d}.
\end{proof}

\begin{lemma}[The prescribed tests detect a depth-fixed period]\label{B8.15e}
Fix a row arrow \(i\), \(0\le i<m\). Write \(S=Y_{i+1}\), \(T=Y_i\) for the successive row spaces and \(\Lambda_S,\Lambda_T\) for their lattices. Put \(M=\operatorname{Hom}_{\mathcal O_E}(\Lambda_S,\Lambda_T)\). Let \(e\) select the marked factor at each of the split depth places, and let \(\mathfrak m\) be the product of the selected prime powers with their grid depths. The integral oscillator lattice \(\Lambda_X\) is the finite completion of the lattice generated by the maps
\[
 s'\longmapsto a\,t\,h_S(s,s'),
 \qquad s\in\Lambda_S,\quad t\in\Lambda_T,\quad
 a\in(\mathfrak c_K\mathcal O_E)\mathfrak b^{-1},
\]
where \(h_S(s,s')=\bar s^{\,t}J_Ss'\) in the grid frame, \(\mathfrak c_K\) is the host conductor ideal and \(\mathfrak b\) is the chosen factor of the relative different. The row depth lattice is
\[
 \Gamma_i=\{u\in\Lambda_X:e u\in\widehat{\mathfrak m M}\}.
\]
For \(\lambda\in M/\mathfrak m M\), let \(v_\lambda\) be the coset vacuum in the oscillator model for \((-S,T)\). Its selected split factor is the indicator of \(e\widehat{\mathfrak m M}+e\lambda\) on \(eX_{\mathbb A}\); elsewhere one uses the grid vacuum. For a complex polynomial \(P\) in the real-place creation operators \(a^+\), put
\[
 E_i(\lambda,P)=\overline{P(a^+)v_\lambda}.
\]
Conjugation here identifies tests for \((-S,T)\) with tests for \((S,T)\). Under the grid coprimalities
\(\gcd(p_{\rm de},Q_Sq)=\gcd(p'_{\rm de},Q_Sq)=1\), every \(E_i(\lambda,P)\) belongs to \(D(\Gamma_i)\). For every \(L\) as in Lemma~\ref{B8.15d},
\[
 \bigl(\forall\lambda,P\ \ L(E_i(\lambda,P))=0\bigr)
 \quad\Longrightarrow\quad
 \bigl(\forall f\in D(\Gamma_i)\ \ L(f)=0\bigr).
\]
\end{lemma}

\begin{proof}[Outline]
Finite translations permute the coset vacua up to scalars. At infinity the holomorphic vacuum translates and reproducing integral extend vanishing to all real translates, so \ref{B8.15d} gives vanishing on the depth space. The complete proof is in \S\ref{prf:B8.15e}.
\end{proof}

\begin{lemma}[Exchange of periods on a column chain]\label{B8.15f}
Fix a grid and the coprimalities of Lemma~\ref{B8.15e}. Use the row measures and any Haar measures and fundamental domains on the cores. Write \(Y_i,V_i\) for the row and core at height \(i\). Let \(F_0,\ldots,F_m\) be a chain as in Lemma~\ref{B8.15c}, including its stated depth condition on every downward test. Let \(u\) be a bottom row form for which every partial row lift with the prescribed tests is rapidly decreasing, including \(u\) itself. Write \(j:\U(Y_0)\to\U(V_0)\) for the stabilizer embedding. If
\[
 \int_{[\U(Y_0)]}F_0(j(h))u(h)\,dh\ne0 ,
\]
then there are polynomials \(P_i\) and cosets \(\lambda_i\) at all \(0\le i<m\) such that the top Bessel period of \(F_m\) against the row lift with these tests is nonzero.
\end{lemma}

\begin{proof}[Outline]
Use \ref{B8.15d} to recover a nonzero bottom Fourier–Jacobi period and propagate nonvanishing by unfolding and Fourier uniqueness. Use \ref{B8.15e} to select the prescribed tests at each step, then conjugate at the top to obtain the Bessel period. The complete proof is in \S\ref{prf:B8.15f}.
\end{proof}

\begin{definition}[Finite-group series at the two theta primes]\label{def:p8c2-series}
For \(g=g_0,g_1\), a labeling consists of a multiplicative map from \(\overline{\F}_g\) to \(\C\), injective on nonzero elements, and generators of \(\F_{g^a}^{\times}\) for \(a\ge1\) compatible under norm. Use the finite general linear group series of Definition~\ref{B1.15ac}. The series condition requires, for each \(g\), a labeling with the prescribed Teichmüller character such that, for every polynomial \(P\in\F_g[X]\), the given series element in \(\C[\GL_r(\F_g)]\) equals the sum of the central character idempotents of the irreducibles labeled by \(P\).
\end{definition}

\begin{theorem}[Real functionals separate rows]\label{B8.15}
Assume \(\mathscr P\) (which includes the absolutely convergent unitary Siegel--Weil formula for \(b>2a\) \cite{weil-1965}, the unramified local doubling zeta-integral formula for both split and inert unitary groups \cite[\S4.2.1]{eischen-harris-li-skinner-2020}\cite[p.~334, Remark~3]{lapid-rallis-2005}, and Harish--Chandra finiteness \cite{harish-chandra-1968}), Lemma~\ref{B8.14}, Lemma~\ref{B8.15j} (via \cite{liu-2022}), and the first formula of Theorem~\ref{B8.13}. For every \(\varepsilon\in\R,\ C_F\ge6510\) there is \(x_0\) such that the following holds for data of exponent \(\varepsilon\) at \(x\ge x_0\) satisfying the row size conditions above and the series condition of Definition~\ref{def:p8c2-series}. For a real realization \(R\), require the prescribed level characters at the base index. Identify its finite group with the product of the residue-field matrix groups at the chosen places above \(g_0,g_1\), so that reduction agrees componentwise, each place lies over the prime selected by the frame, and the identification carries the embedded coefficients of the level idempotent to those of the type idempotent. Every such realization with a grid at \((1,C_F,b)\) (Definition~\ref{B8.21c}), its Weil system and compatible row measures satisfies
\[
 \forall w\in\mathcal V\setminus\{0\}\quad \exists\rho,\xi\quad
 \forall t\in\mathcal E_0\setminus\{0\}\quad\exists\Psi,\Phi\quad
 \mathcal B_{\rho,\xi,\Psi,\Phi}(t(w))\ne0
\]
with convergence of the lifts and integrations.
\end{theorem}

\begin{proof}[Outline]
Use \ref{B8.15b} and \ref{B8.15c} to choose a nonzero column chain, with cuspidality and decay from \ref{B3.00e} and \ref{B1.15b}. Apply \ref{B8.15f} and use real Bessel multiplicity one to choose a real pair that works for every nonzero row map. The complete proof is in \S\ref{prf:B8.15}.
\end{proof}

\begin{lemma}[The two factors in the lowering projection]\label{B8.15g}
Let \(G=\U(p,q)\), let \(z\) be a positive coordinate unit vector in a diagonal signature frame, and let \(H=\U(p-1,q)\) be its stabilizer. Put \(n=p+q\), and suppose \(p-1\ge q\ge1\) and \(n\ge4\). Let \(\pi,\sigma\) be discrete series of \(G,H\) with decreasing Harish--Chandra parameters \(\lambda,\mu\). Suppose the entries of both parameters have absolute value at most \(B\), consecutive entries in each list differ by at least two, and entries of opposite signs in either signed list differ in absolute value by at least \(1000n^3\). Put \(X=8nB\). Write \(e_l u\) for the tuple with \(u\) in coordinate \(l\) and zero elsewhere.

On the smooth vectors of \((\C^n)^\vee\otimes\pi\), let \(p_s\) denote the projection to the discrete series \(\pi_s\) of parameter \(\lambda-e_s\), where \(e_s\) in a parameter denotes the \(s\)-th standard basis vector. These projections sum to the identity. There are component maps \(R_s\) to the unitary models of \(\pi_s\) and nonzero real constants \(c_s\) with the following properties. For every lowest compact type unit vector \(v\) of \(\pi\), every \(l\), and every bounded \(H\)-intertwiner \(P:\pi\to\sigma\),
\[
 |c_s|\,\|R_s(e_l v)\|^2\le X^{4n},\qquad
 P\bigl((p_s(e_z u))_z\bigr)=\gamma_s Pu\quad(u\in\pi^\infty),
 \qquad
 \gamma_s=
 \frac{\displaystyle\prod_{a=1}^{n-1}(\mu_a-\lambda_s+\tfrac12)}
      {\displaystyle\prod_{j\ne s}(\lambda_j-\lambda_s)}.
\]
Here \(e_z u\) means insertion in the coordinate of \(z\). In particular \(|\gamma_s|\le X^n\). For the same Haar measure on \(G\), the formal degrees satisfy \(d_\pi/d_{\pi_s}\le X^n\).

For later use, if \(D(A)\) is the complexified derived action of a matrix \(A\) and \(v\) is as above, then
\[
 \|D(A_1)\cdots D(A_t)v\|
 \le \bigl(8\sqrt n(B+S)\bigr)^t
 \quad\text{if }t\le S,\quad \|A_j\|_{\rm HS}\le1.
\]
\end{lemma}

\begin{proof}[Outline]
Decompose the tensor product by the coherent character identity and globalize by \ref{B8.38a}. The Capelli comatrix gives the projections and their stabilizer scalars, while compact Casimir bounds and \ref{B1.15b} control the norms and formal degrees. The complete proof is in \S\ref{prf:B8.15g}.
\end{proof}

\begin{lemma}[Haar volumes for a bounded column]\label{B8.15h}
Use the real groups and parameters of Lemma~\ref{B8.15g}. Give \(G,H\) Haar measures and give \(H\backslash G\) the corresponding quotient measure. Write \(B_H(1)=\{h:\|h-1\|_{\rm HS}<1\}\) and
\[
 D=\{Hg:\|g^{-1}z\|^2\le2X^{8n}\}.
\]
For a lowest \(H\)-type unit vector \(w\) of \(\sigma\), its diagonal coefficient is integrable, and
\[
 \frac1{d_\sigma}\le X^{67n^2}\operatorname{vol}(B_H(1)),
 \qquad
 d_\pi\operatorname{vol}(D)\operatorname{vol}(B_H(1))
 \le4X^{24n^3}.
\]
Every coset in \(D\) has a representative of Hilbert--Schmidt norm at most \(X^{10n}\).
\end{lemma}

\begin{proof}[Outline]
Combine Cartan decay from \ref{B8.53a} with ball coverings and Schur orthogonality from \ref{B1.15b}. Hyperbolic rotations give bounded coset representatives, and \ref{B8.15g} bounds the coefficient variation on a small ball for the quotient volume estimate. The complete proof is in \S\ref{prf:B8.15h}.
\end{proof}

\begin{lemma}[A column moment for the branching coefficient]\label{B8.15i}
In Lemma~\ref{B8.15g}, suppose the two signed parameters satisfy the discrete branching rule: in their merged decreasing list, adjacent signs agree exactly when the entries belong to different lists. Let \(i:\sigma\hookrightarrow\pi|_H\) be an isometric intertwiner. For lowest-type unit vectors \(w\) of \(\sigma\) and \(v_0\) of \(\pi\),
\[
 \int_G\|g^{-1}z\|^2\,|\langle iw,\pi(g)v_0\rangle|^2\,dg
 \le \frac{X^{6n+3}}{d_\pi}\le\frac{X^{8n}}{d_\pi}.
\]
\end{lemma}

\begin{proof}[Outline]
Use \ref{B8.15h} and Schur orthogonality to project onto the branching line, using multiplicity one. Apply the tensor projections and bounds of \ref{B8.15g} to the column-weighted coefficient and combine the squared-integral estimates by Cauchy–Schwarz. The complete proof is in \S\ref{prf:B8.15i}.
\end{proof}

\begin{lemma}[A coefficient for the bottom period]\label{B8.15a}
Let \(Y\) be the bottom row of rank \(r\), \(V=Y\perp Ez_0\) the bottom core of rank \(n=r+1\), and \(j:\U(Y)\to\U(V)\) the stabilizer embedding of the prescribed norm-one \(z_0\). Let \(\mathfrak m\) be the grid modulus. Under the size and level assumptions of Theorem~\ref{B8.15}, there is an absolute \(C\ge0\) such that eventually for each \(\varepsilon,C_F\), uniformly in the grid and discrete-series models, the following holds. Put \(L_{\rm loss}=Cfn^4\log b\). There is a nonzero pure smooth row tensor \(w_*\) and a core coefficient \(k\), the product of the real coefficients and the finite core congruence indicator. For every nonzero row map \(t\), put \(u=t(w_*)\). For the row Haar measure there is \(c\ne0\) with
\begin{equation}\label{eq:p8c2-row-reproduce}
 \int_{\U(Y)(\mathbb A_F)}k(j(y)^{-1}j(h))u(h)\,dh=c\,u(y).
\end{equation}
For \(y\in\U(Y)(\mathbb A_F)\), put \(N_\tau(z)=\|j(y)_\tau^{-1}z_\tau\|\) in the adapted Euclidean frames. For rational \(\gamma\in\U(V)(F)\) and \(z=\gamma z_0\),
\begin{equation}\label{eq:p8c2-row-coefficient}
 \int_{\U(Y)(\mathbb A_F)}|k(j(y)^{-1}\gamma j(h))|\,dh
 \prod_\tau N_\tau(z)^{2n+5}\le e^{L_{\rm loss}}|c|.
\end{equation}
The integrand vanishes unless \(z-z_0\in\mathfrak m\Lambda_{j(y)}\), where the lattice has finite completion \(j(y)_f\widehat\Lambda_V\). For \(X\) in the row Lie algebra in one adapted real frame,
\begin{equation}\label{eq:p8c2-row-variation}
 |u(y\exp(sX))-u(y)|\le e^{L_{\rm loss}}\sup|u|\,|s|\,\|X\|_{\rm HS}.
\end{equation}
\end{lemma}

\begin{proof}[Outline]
Use discrete branching, the column moment \ref{B8.15i} and the volumes \ref{B8.15h} to find a bounded translate with large row projection. A compact operator gives a reproducing eigenvector, with decay from \ref{B8.53a} and support from the finite congruence indicator. The complete proof is in \S\ref{prf:B8.15a}.
\end{proof}

\begin{lemma}[A bottom period at the prescribed vector]\label{B8.15b}
For each \(\varepsilon\in\R\) and \(C_F\ge6510\) there is a threshold for grids satisfying the row size conditions of Theorem~\ref{B8.15}. For any bottom row and core models of the prescribed parameters, any nonzero row multiplicity map \(t\) at row level and any nonzero \(w\) in its smooth tensor space, there are a core multiplicity map \(y\) at core level and an \(a\) in the span of pure local \(K\)-finite tensors with
\[
 \operatorname{Per}_0(y(a),t(w))=\int_{[\U(Y)]}y(a)(j(h))t(w)(h)\,dh\ne0.
\]
\end{lemma}

\begin{proof}[Outline]
Use \ref{B8.15a} and cuspidality to rule out short lattice vectors where the row form has large modulus. Unfold the translated Poincaré series so congruence depth makes the identity term dominate, then use cyclicity and compact-type approximation. The complete proof is in \S\ref{prf:B8.15b}.
\end{proof}

\begin{lemma}[The inert spherical Hecke identity in even ranks]\label{B8.15j}
Let \(L/K\) be an unramified quadratic extension of nonarchimedean local fields of characteristic zero, with \(Q\) elements in the residue field of \(K\), \(Q\) odd. Take the split skew-Hermitian space of dimension \(2a\) and the split Hermitian space of dimension \(2b\) over \(L\), with \(1\le a\le b\). Use self-dual lattices, an additive character of \(K\) of conductor \(\mathcal O_K\), and trivial splitting characters. Normalize Haar measure on each group so that the lattice stabilizer has volume one. Let \(\omega\) be the Weil representation and \(\phi_0\) its lattice vacuum.

Write \(\mathcal H_a,\mathcal H_b\) for the spherical Hecke algebras and \(\mathcal S_a,\mathcal S_b\) for their normalized Satake transforms. Use coordinates in which the base-change roots in dimension \(2a\) are \(t_1,t_1^{-1},\ldots,t_a,t_a^{-1}\). For \(h'\in\mathcal H_b\) there is an \(h\in\mathcal H_a\) with
\[
 \omega(1,h')\phi_0=\omega(h,1)\phi_0.
\]
For every such \(h\), with \(c=b-a\),
\[
 (\mathcal S_a h)(t_1,\ldots,t_a)
 =
 (\mathcal S_b h')(Q,Q^3,\ldots,Q^{2c-1},t_1,\ldots,t_a).
\]
The list preceding the \(t_j\)'s is empty for \(c=0\).
\end{lemma}

\begin{proof}[Outline]
Use \ref{B8.14} and compact averaging to turn the target action on the vacuum into source convolution. Liu's spherical-module theorem gives the Satake substitution adjoining \(Q,Q^3,\ldots,Q^{2(b-a)-1}\). The complete proof is in \S\ref{prf:B8.15j}.
\end{proof}

\begin{lemma}[Spherical parameters in the even core chain]\label{lem:p8c2-spherical-transfer}
Consider a gap-one grid and admissible level chain with Weil system and row measures, core ranks \(n=r+1,\ d=r+5,\ n_i=n+id\) (\(0\le i\le m\), \(r\) odd), and exceptional set \(S\) consisting of the base place above \(q_0\) and the places of \(F\) above \(g_0g_1p_{\rm de}p'_{\rm de}\). The two splitting characters in step \(i\) are \(\widetilde\chi_i=\chi_{m-i}^{-1}\), with \(\chi_{m-i}\) the grid character. At \(w\mid v\notin S\), the base-change Satake multiset \(B(w)\) of a nonzero conjugate bottom spherical cuspidal eigenfunction satisfies
\begin{equation}\label{eq:p8c2-bottom-satake}
 |B(w)|\le n,\qquad |\beta|\le(Nw)^{(n-1)/2}\quad(\beta\in B(w)).
\end{equation}
If the conjugate of a cuspidal value at stage \(i\) has multiset \(B_i(w)\), the conjugate of its nonzero upward vacuum lift has multiset
\[
 B_{i+1}(w)=B_i(w)\sqcup
 \{\widetilde\chi_i(w)(Nw)^{(d-1)/2-l}:0\le l<d\}.
\]
\end{lemma}

\begin{proof}
Cuspidal decay gives a bounded spherical eigenfunction. Spherical convolution on bounded functions has norm bounded by its \(L^1\)-mass. In the split chart the trivial Satake bound for this isometric action bounds each normalized root by \((Nv)^{(n-1)/2}\). At an inert place use the trivial bound in the unitary spherical algebra. For its inert coordinates \(\alpha_j\) the base-change multiset is
\[
 \{\alpha_j,\alpha_j^{-1}:1\le j\le\lfloor n/2\rfloor\}\sqcup\{1\}^{\,n\bmod2}.
\]
With \(Nw=(Nv)^2\) the inert bound gives \eqref{eq:p8c2-bottom-satake}.

Move Hecke convolutions through the Weil pairing by equivariance. Write \(\chi_S,\chi_T\) for the two step characters, \(r_0=n_i\) for the source rank, and \(q=Nw\). In the split matrix model the target generator of index \(k\), taken on inverse double cosets, acts on the vacuum as a sum of source generators of indices \(j\le k\). Count the lattice modifications by intersection dimension \(j\). The complementary choices contribute the Gaussian binomial \({d\brack k-j}_q\), and the lifts into the source block contribute \(q^{(r_0-j)(k-j)}\). The Weil multiplier changes their product by
\[
 (\chi_S(w)q^{r_0/2})^{-k}(\chi_T(w)q^{(r_0+d)/2})^j.
\]
Substitute the source eigenvalues \(q^{j(r_0-j)/2}e_j(B_{\rm old}^{-1})\), where \(e_j\) is elementary symmetric and the multiset is inverted entrywise. The identity
\[
 e_u(q^{(d-1)/2},q^{(d-3)/2},\ldots,q^{-(d-1)/2})
 =q^{-u(d-u)/2}{d\brack u}_q
\]
gives the sum \(q^{k(r_0+d-k)/2}e_k(B_{\rm new}^{-1})\), with
\[
 B_{\rm new}=\chi_S(w)\chi_T(w)^{-1}B_{\rm old}
 \sqcup\{\chi_S(w)q^{(d-1)/2-l}:0\le l<d\}.
\]
At an inert \(v\notin S\), put \(Q=Nv\), so \(Nw=Q^2\). Both core ranks are even: \(r\) is odd, \(n_i=r+1+i(r+5)\), and \(d=r+5\). The self-dual frames give the split spaces in these even dimensions over the unramified quadratic extension. Multiplying the source form by an anti-invariant unit gives the skew-Hermitian convention with the same integral stabilizer. The residue characteristic is odd and the local additive character has conductor \(\mathcal O_v\). Each splitting character is unramified and restricts to the even power of the quadratic character on \(F_v^\times\). It is therefore one on the units of \(E_w\) and on a common uniformizer, hence trivial on \(E_w^\times\). Apply Lemma~\ref{B8.15j} with \(2a=n_i\) and \(2b=n_i+d\). Move its matching source convolution across the invariant bilinear pairing. This replaces \(h(g)\) by \(h(g^{-1})\), which is the same function because the unitary Cartan double cosets are invariant under inversion. The eigenvalue of every target operator on the nonzero lift is consequently its Satake polynomial evaluated at the old coordinates with \(Q,Q^3,\ldots,Q^{d-1}\) adjoined. Base change adds those entries and their inverses. In terms of \(Nw\) these are exactly \((Nw)^{(d-1)/2-l}\), \(0\le l<d\), and \(\widetilde\chi_i(w)=1\). In the chain \(\chi_S=\chi_T=\widetilde\chi_i\). Pairing the ordinary kernel with the conjugate cusp form computes the conjugate of the upward lift with conjugated kernel, which gives the stated parameter convention.
\end{proof}

\begin{lemma}[Core values spanned by column chains]\label{B8.15c}
Assume the local inputs \(\mathscr P\) of Theorem~\ref{B8.15}, in particular the absolutely convergent unitary Siegel--Weil formula for \(b>2a\) \cite{weil-1965}, unramified local doubling for both split and inert unitary groups \cite[\S4.2.1]{eischen-harris-li-skinner-2020}\cite[p.~334, Remark~3]{lapid-rallis-2005}, and Harish--Chandra finiteness \cite{harish-chandra-1968}, together with Lemma~\ref{B8.14}, Lemma~\ref{B8.15j}, and the first formula of Theorem~\ref{B8.13}. Require \(p_{\rm de},p'_{\rm de}\) coprime to \(Q_Sq\), \(r>1800\), and \(b_{\tau,j+1}-b_{\tau,j}>2N_*^3\) at every infinite place \(\tau\) for \(j+1<m\). Use a gap-one grid, an admissible level chain, its Weil system and row measures. Choose Haar measures and fundamental domains for the core groups. For every bottom core model of the prescribed parameters and Cartan data, core multiplicity map \(y\) at core level and \(K\)-finite tensor \(a\), the value \(y(a)\) belongs to the span of the following \(F_0\). There are rapidly decreasing \(F_i\) on the core quotients with
\[
 F_m=\phi_{\xi,\Psi},\quad\Psi\in\mathcal C_{\rm top},\qquad
 F_i(g)=\int_{[\U(V_{i+1})]}\theta_i(h,g;P_i)F_{i+1}(h)\,dh
 \quad(0\le i<m).
\]
Here \(\mathcal C_{\rm top}\) is the space of automorphic forms on the definite top group with the prescribed coefficient representation, top level and trivial level character. Each test \(P_i\) is fixed by the depth-lattice Heisenberg translations, with selected components in \(\mathfrak m\widehat{\operatorname{Hom}}(\Lambda_S,\Lambda_T)\) for the pair's source and target lattices.
\end{lemma}

\begin{proof}[Outline]
Decomposes the bottom multiplicity space into Hecke eigenvectors and shows each upward lift nonzero via doubling, \ref{B8.13}, Siegel--Weil, and \ref{lem:p8c2-hecke-nonvanishing}. Distinct Satake multisets give downward surjectivity, and Lie stability spans \(K\)-finite vectors. The complete proof is in \S\ref{prf:B8.15c}.
\end{proof}

\section{Twisted square models}

Fix an arithmetic point with CM extension \(E/F\), and choose its realization \(\theta\) at the infinite places. Let \(W,W_1,X\) have ranks \(u,v,d=u+v\), with \(d\) even and signatures \((u-p,p),(v-q,q),(d-p-q,p+q)\). Require \(v+2(2(p+q)+1)<u\) and \(2(p+q)+1<v\). The characters \(\alpha,\beta\) restrict to \(\eta^v,\eta^u\) and have rational exponents \(a_\tau,b_\tau\) with \(2(a_\tau-b_\tau)<-d\). At each finite component both characters are unramified or their quotient is ramified at a split component prime to the different of \(F/\Q\). Write \(\mathcal P(W,X)\) for polynomial Fock vectors. Algebraic Fock functions are finite sums of \(\operatorname{Fock}(f,P)\) with permitted finite factors \(f\) and real polynomials \(P\). Let \(\mathcal R_{\rm fin}\) be the finite source differences and \(\theta\operatorname{Int}_\omega\) the source theta integral.

\begin{definition}[Twisted-square triples, seesaw transfer and adjoint model identifications]\label{B8.16c}
Write \(x\mapsto\bar x\) for conjugation on \(E/F\) and \(\xi_0\in E^\times\) (\(\bar\xi_0=-\xi_0\)) for the host imaginary scalar, a unit at every inert place of \(E/F\) satisfying \(\operatorname{Im}\tau(\xi_0)>0\) for every \(\tau\in\theta\). Assume \(W,W_1,X\) carry Cartan involutions and self-dual \(\mathcal O_E\)-lattices. A tuple \((W,W_1,X,\alpha,\beta)\) satisfying the rank, signature, exponent and ramification conditions above (where the unitary Hecke characters \(\alpha,\beta\) act on the unit circle at each \(\tau\in\theta\) with integer angular exponents \(2a_\tau,2b_\tau\)) is called a \emph{twisted-square triple}.

Fix frames \((A,A_1,B)\) for \((W,W_1,X)\) with Hermitian Gram matrices \(J_W,J_{W_1},J_X\), and write \(H_X(h)\) for the adelic matrix height of \(h\in\U(X)(\mathbb A_F)\) in frame \(B\). For Hermitian spaces \(M,X\) over \(E/F\) with frames \((A_M,B)\) (in particular for \(M\in\{W,W_1\}\)), an oscillator Weil representation \(\omega\) of sign \(\epsilon\in\{+,-\}\) and splitting pair \((\chi_M,\chi_X)\) in the model of Lemma~\ref{B8.15d} carries a continuous additive character \(\psi\) of \(\mathbb A_F/F\) of sign \(-1\) (infinite part \(y\mapsto\exp(-2\pi i\operatorname{Tr}_{F_\infty/\R}y)\)), acts on functions on \(\operatorname{Hom}_E(M,X)(\mathbb A_F)\) through \(J_M\) for \(\epsilon=+\) and \(-J_M\) for \(\epsilon=-\), and has theta kernel \(\theta_\omega(g,h;\Phi)=(\omega(g,h)\Phi)(0)\) and source theta integral \(\theta\operatorname{Int}_\omega(\Phi)(h)=\int_{\mathcal F}\theta_\omega(g,h;\Phi)\,d\nu(g)\) for a Haar measure \(\nu\) and \(\U(M)(F)\)-fundamental domain \(\mathcal F\) on \(\U(M)(\mathbb A_F)\). Permitted tests \(\Phi\) are theta tests in the model of Lemma~\ref{B8.15d}, and permitted finite functions \(f\) are smooth oscillator theta functions annihilated by the archimedean Fock annihilation operators for \((\epsilon M,X)\) in frames \((A_M,B)\). For \(P\in\mathcal P(W,X)\) and \(P'\in\mathcal P(W_1,X)\), write \(\operatorname{Fock}(f,P)\) and \(\operatorname{Fock}'(f',P')\) for the Fock functions of signs \(+\) and \(-\); their complex spans are the unprimed and primed algebraic Fock spaces, on which \(\mathcal R_{\rm fin}\) denotes the complex span of finite-adelic source differences \(\omega(u_f,1)\Psi-\Psi\). \emph{Matching data} \((\omega,\omega',\Pi,p,p')\) for \(\psi\) consist of oscillator Weil representations \(\omega,\omega'\) for \((W,X),(W_1,X)\) of signs \(+,-\) and splitting pairs \((\beta,\beta/\alpha),(\alpha,\alpha/\beta)\), a nonzero complex vector space \(\Pi\) jointly irreducible under the commuting actions of \(\U(X)(\mathbb A_{F,f})\) and the complexified target Lie algebras at all infinite places, and target-equivariant linear maps \(p,p'\) mapping the unprimed and primed algebraic Fock spaces onto \(\Pi\) and vanishing on \(\mathcal R_{\rm fin}\) and on the archimedean source Lie actions.

Given Haar measures and fundamental domains \((\nu,\mathcal F),(\nu_1,\mathcal F_1)\) on \(\U(W)(\mathbb A_F),\U(W_1)(\mathbb A_F)\), a \emph{seesaw transfer} consists of a bijection \(e:\{1,\ldots,d\}\simeq\{1,\ldots,u+v\}\), Hermitian spaces \(V,V'\) of ranks \(u,v\) and a skew-Hermitian space \(W^{\rm skew}\) of rank \(u+v\) with global frames (together with local split frames of their self-dual lattices) and Gram matrices
\[
 J_V=J_W,\qquad J_{V'}=-J_{W_1},\qquad J_{W^{\rm skew}}=\xi_0\,(J_X^t)^e
 \quad\bigl(((J_X^t)^e)_{ij}=(J_X^t)_{e^{-1}(i),e^{-1}(j)}\bigr),
\]
infinite embeddings of \(E\) conjugate to those of \(\theta\), good base data preceding Theorem~\ref{B8.11} with the same additive character \(\psi\), Schr\"odinger Weil models \((\tilde\omega,\tilde\omega')\) on \((V\otimes_E W^{\rm skew})(\mathbb A_F)\) and \((V'\otimes_E W^{\rm skew})(\mathbb A_F)\) forming a triple preceding Theorem~\ref{B8.11} with characters \((\alpha,\beta)\) and source measures and domains pushed forward from \((\nu,\mathcal F),(\nu_1,\mathcal F_1)\) along the frame matrix identifications \(\kappa_W:\U(W)\xrightarrow{\sim}\U(V)\) and \(\kappa_{W_1}:\U(W_1)\xrightarrow{\sim}\U(V')\), and the target isomorphism \(\kappa_X:\U(X)\xrightarrow{\sim}\U(W^{\rm skew})\) sending a matrix \(h\) in frame \(B\) to \((\bar h)^e\) (entrywise conjugation followed by reindexing by \(e\)). Given \(\mathbb A_F\)-linear topological symplectic isomorphisms
\[
 c_A:(V\otimes_E W^{\rm skew})(\mathbb A_F)\xrightarrow{\;\sim\;}\operatorname{Hom}_E(W,X)(\mathbb A_F),\qquad
 c_B:(V'\otimes_E W^{\rm skew})(\mathbb A_F)\xrightarrow{\;\sim\;}\operatorname{Hom}_E(W_1,X)(\mathbb A_F)
\]
preserving the rational subspaces, the \emph{adjoint model identifications} \((\mathcal W_A,\mathcal W_B)\) send Schr\"odinger tests \(\Phi_D\) on \(\mathbb A_F^{u\times(u+v)}\) and \(\Phi'_D\) on \(\mathbb A_F^{v\times(u+v)}\) to \(\Phi=\mathcal W_A(\Phi_D)\) and \(\Phi'=\mathcal W_B(\Phi'_D)\) defined by
\[
 \Phi(c_A(w))=\sum_{\xi\in F^{u\times(u+v)}}(\rho_\psi(w)\Phi_D)(\xi),\qquad
 \Phi'(c_B(w'))=\sum_{\xi'\in F^{v\times(u+v)}}(\rho'_\psi(w')\Phi'_D)(\xi'),
\]
where \(\rho_\psi,\rho'_\psi\) denote the Darboux Heisenberg actions for \(\psi\) on the two Schr\"odinger spaces.
\end{definition}

\begin{lemma}[Seesaw transfer of a twisted-square triple]\label{B8.16d}
Let \((W,W_1,X,\alpha,\beta)\) be a twisted-square triple with frames \((A,A_1,B)\), Haar measures and fundamental domains \((\nu,\mathcal F)\) on \(\U(W)(\mathbb A_F)\) and \((\nu_1,\mathcal F_1)\) on \(\U(W_1)(\mathbb A_F)\), and oscillator Weil representations \(\omega,\omega'\) on \(\operatorname{Hom}_E(W,X)(\mathbb A_F)\) and \(\operatorname{Hom}_E(W_1,X)(\mathbb A_F)\) of signs \(+,-\) and splitting pairs \((\beta,\beta/\alpha),(\alpha,\alpha/\beta)\) for an additive character \(\psi\) of \(\mathbb A_F/F\). Then there exist a seesaw transfer \((e,V,V',W^{\rm skew},\tilde\omega,\tilde\omega',\kappa_W,\kappa_{W_1},\kappa_X)\) and symplectic isomorphisms \((c_A,c_B)\) as in Definition~\ref{B8.16c} with the following properties. The spaces \(V,V',W^{\rm skew}\) have signatures \((u-p,p),(q,v-q),(d-p-q,p+q)\) satisfying \eqref{eq:p8c2-triple-range}, at every infinite place \(w\) of \(F\) with chosen embedding \(\tau\in\theta\) the angular exponents under \(\bar\tau\) satisfy \(m_{\alpha,w}-m_{\beta,w}=2(b_\tau-a_\tau)>d\), and there exist \(C'>0\) and \(N'\in\N\) with \(\operatorname{ht}(\kappa_X(h))\le C'H_X(h)^{N'}\) for all \(h\in\U(X)(\mathbb A_F)\). The adjoint model identifications \((\mathcal W_A,\mathcal W_B)\) map the full polynomial Schwartz spaces of \((\tilde\omega,\tilde\omega')\) \(\C\)-linearly and bijectively onto the unprimed and primed algebraic Fock spaces of \((\omega,\omega')\), intertwine the finite-adelic and archimedean Lie actions of source and target via \((\kappa_W,\kappa_X)\) and \((\kappa_{W_1},\kappa_X)\), and satisfy
\[
 \theta_\omega(g,h;\Phi)=\theta_{\tilde\omega}(\kappa_W(g),\kappa_X(h);\Phi_D),\qquad
 \theta_{\omega'}(g_1,h;\Phi_0)=\theta_{\tilde\omega'}(\kappa_{W_1}(g_1),\kappa_X(h);\Phi_{0D})
\]
for all \(g\in\U(W)(\mathbb A_F)\), \(g_1\in\U(W_1)(\mathbb A_F)\), \(h\in\U(X)(\mathbb A_F)\), and corresponding algebraic pairs \((\Phi_D,\Phi)\) and \((\Phi_{0D},\Phi_0)\). Under the unitary Siegel-set estimate of Proposition~\ref{B1.15b}, for every \(h\in\U(X)(\mathbb A_F)\), every corresponding unprimed algebraic pair \((\Phi_D,\Phi)\), and every corresponding primed algebraic pair \((\Phi_{0D},\Phi_0)\) with \(\Phi_{0D}\in\mathcal S'_{t_1}\), both integrals below converge absolutely and satisfy
\[
 \theta\operatorname{Int}_{\omega}(\Phi)(h)=I(\Phi_D,\kappa_X(h)),\qquad
 \theta\operatorname{Int}_{\omega'}(\Phi_0)(h)=I'(\Phi_{0D},\kappa_X(h)).
\]
\end{lemma}
\begin{proof}[Outline]
Transfers \((W,W_1,X)\) to \((V,V',W^{\rm skew})\) matching \ref{B8.16c} and \ref{B8.11}. Adjoint maps \(c_A,c_B\) induce bijections \((\mathcal W_A,\mathcal W_B)\), giving the kernel identities and, via \ref{B1.15b} and \ref{B8.11e}, the integral identities. The complete proof is in \S\ref{prf:B8.16d}.
\end{proof}

\begin{lemma}[The seesaw comparison scalar]\label{B8.16e}
Assume \(\mathscr P\). Keep the twisted-square triple \((W,W_1,X,\alpha,\beta)\) with frames \((A,A_1,B)\), matching data \((\omega,\omega',\Pi,p,p')\) for \(\psi\), and seesaw transfer \((e,V,V',W^{\rm skew},\tilde\omega,\tilde\omega',\kappa_W,\kappa_{W_1},\kappa_X)\) with adjoint model identifications \((\mathcal W_A,\mathcal W_B)\) from Lemma~\ref{B8.16d}. Fix matching maps \((q_{\rm ss},q'_{\rm ss})\) onto the jointly irreducible target quotient \(\Pi_{\rm ss}\) preceding Theorem~\ref{B8.11} for \((V,V',W^{\rm skew},\tilde\omega,\tilde\omega')\). Then there exists a unique nonzero scalar \(s\in\C^\times\), called the \emph{seesaw comparison scalar}, such that for every corresponding unprimed algebraic pair \((\Phi_D,\Phi)\) and every corresponding primed algebraic pair \((\Phi_{0D},\Phi_0)\) with \(p(\Phi)=p'(\Phi_0)\) in \(\Pi\) (and in particular whenever \(\Phi_{0D}\in\mathcal S'_{t_1}\), on which subclass \(s\) is already uniquely determined), the pair \((\Phi_D,s\Phi_{0D})\) is a polynomial match for \((q_{\rm ss},q'_{\rm ss})\).
\end{lemma}
\begin{proof}[Outline]
Transports \(p,p'\) via \ref{B8.16d} and factors \(\Pi,\Pi_{\rm ss}\) through one simple restricted tensor product using \ref{B1.15b}, \ref{B8.11f}, \ref{B8.14a}, \ref{B8.11g}, and \ref{B8.11}. Countable-dimensional Schur's lemma yields the unique scalar \(s\in\C^\times\). The complete proof is in \S\ref{prf:B8.16e}.
\end{proof}

\begin{theorem}[Convergence and matching for twisted squares]\label{B8.16}
Assume \(\mathscr P\), Theorem~\ref{B8.11} and the unitary Siegel-set estimate of Proposition~\ref{B1.15b}. For either Hermitian sign and compatible data for \(M,X\), suppose at some infinite place \(\operatorname{sig}(M)=(\operatorname{rank}M-p,p)\) and \(\operatorname{rank}X<2(\operatorname{rank}M-p)\). For each permitted \(\Phi\), Haar measure and fundamental domain there are \(C\in\R,\ N\in\N\) with
\begin{equation}\label{eq:p8c2-square-moderate}
 \int_{[\U(M)]}|\theta_\omega(g,h;\Phi)|\,dg\le C H_X(h)^N
 \qquad(h\text{ arbitrary})
\end{equation}
for the prescribed target height.

Fix the standard additive character of sign \(-1\), the frames and the exponents \(a_\tau,b_\tau\). There are surjections \(q,q'\) from \(\mathcal P(W,X),\mathcal P(W_1,X)\) onto a common \(P_{\rm real}\). For every twisted-square triple with these data, matching data exist with signs \(+,-\) and splitting pairs \((\beta,\beta/\alpha),(\alpha,\alpha/\beta)\). The common quotient \(\Pi\ne0\) is jointly irreducible for the finite adelic and real Lie actions. The maps \(p,p'\) on the stable algebraic Fock spaces intertwine these actions and kill the finite source differences and source Lie action. For each match there are an isomorphism \(e:P_{\rm real}\otimes_\C P_{\rm fin}\to\Pi\) and maps \(r,r'\) onto \(P_{\rm fin}\) on permitted finite functions with
\[
 p(\operatorname{Fock}(f,P))=e(q(P)\otimes r(f)),\qquad
 p'(\operatorname{Fock}'(f',P'))=e(q'(P')\otimes r'(f')) .
\]

For any additive character with a twisted-square triple, matching data and compatible measures and domains, there is \(c\ne0\) such that every primed algebraic \(\Phi'\) has an algebraic representative \(\Phi_0\bmod\mathcal R_{\rm fin}\) satisfying \eqref{eq:p8c2-square-moderate}. For every unprimed algebraic \(\Phi\),
\[
 p(\Phi)=p'(\Phi_0)\quad\Longrightarrow\quad
 \theta\operatorname{Int}_\omega(\Phi)=c\,\theta\operatorname{Int}_{\omega'}(\Phi_0).
\]
\end{theorem}

\begin{proof}[Outline]
Proves moderate growth on a Siegel set of \ref{B1.15b} via \ref{B8.11e}. Builds matching data and \(P_{\rm real}\otimes_\C P_{\rm fin}\simeq\Pi\) via \ref{B1.15ac}, \ref{B8.16d}, and \ref{B8.11}, then transfers representatives and proportionality via \ref{B8.16e}. The complete proof is in \S\ref{prf:B8.16}.
\end{proof}

\section{Fields of definition and finite tests}

A complete choice fixes a grid with Weil system, row measures, a nonzero smooth tensor with real tests, square constants and a test lattice. For preliminary datum \(P\) with frame \(f\), write \(\iota_X:K_X\hookrightarrow\C\) for the coefficient embedding of a completion \(X\). The coefficient-field condition requires in \(\iota_X(K_X)\) the finite local component values of the row characters \(\mu_{0,j}\), the level twists \(\xi_t\) and the layer characters \(\chi_i\) for \(1\le i\le m\). It requires the finite local values of \(\eta\) with the half-weight normalization and of \(\chi_\Psi,\delta,\delta'\) with weight zero. It contains the eighth roots of unity. At index \(t\) the finite test places are the chosen lifts of the exceptional places of \(F\), together with the twist places \(Q_t\). At each such \(w\) it contains the roots of unity of order \((Nw)^{d_w}\), where \(d_w\) is the prescribed base depth at the chosen base place, two at the other chosen lifts, and one elsewhere. It contains \(\sqrt{Nw}\) at every finite test place and at each additional place chosen by the grid to represent the support of the auxiliary ideal and its conjugate. Each phase and jet value on a support term belongs to the field.

For each \(t,i,j\) and integer \(1\le s\le N_*+1\), the ratio \(L_t(s)/L_{\mu_0}(s)\) for the corresponding layer characters belongs to the field whenever \(L_{\mu_0}(s)\ne0\). For each \((i,j)\) for which a bottom triple exists, require a choice of such a triple and two nonzero periods \(\Omega,\Omega'\). For every \(t\), finite target point and monomial test for which the corresponding integral exists absolutely, its value divided by its own side's period belongs to the field.

\begin{lemma}[A common period field for the two bottom integrals]\label{B8.17a}
Assume the local theta inputs \(\mathscr P\). Fix a grid with its Weil system, a finite set \(T\) of level indices satisfying the prescribed character and ramification conditions, and \(0\le i,j<m\). For a standard bottom triple at \((i,j)\), write \(A,A'\) for its two source spaces and \(W\) for its target. For \(t\in T\), \(h\in\U(W)(\mathbb A_{F,f})\) and a Fock monomial \(X^k\) on the respective source, write \(I_t(h,k)\) and \(I'_t(h,k)\) for the two theta integrals. There are \(\Omega,\Omega'\in\C^\times\) and a number field \(M\subset\C\) such that
\[
                   I_t(h,k)/\Omega\in M,\qquad
                   I'_t(h,k)/\Omega'\in M
\]
for every such test on its respective side.
\end{lemma}

\begin{proof}[Outline]
Reduce translates and monomials to finite sets by the class-number theorem and central rotation. Definite class sums give algebraic projections, and the CM Taylor formula gives normalized values whose compositum works for both families. The complete proof is in \S\ref{prf:B8.17a}.
\end{proof}

\begin{theorem}[A finite field requirement for a finite character set]\label{B8.17}
Assume \(\mathscr P,\mathscr J\), the Hecke--Tate functional equation and Theorems~\ref{B8.11}, \ref{B8.15}, \ref{B8.16}. There is \(C_*\in\R\) such that for each \(\varepsilon\in\R,\ C_F\ge C_*\) there is \(x_0\) with the following property. Given a scale \(s\) of exponent \(\varepsilon\) with \(x(s)\ge x_0\), \(P,f\) and a finite set \(\Xi\) of Hecke characters of \(E\), there is a finite \(S\subset\overline{\Q}\cap\C\). For each completion \(X\) with \(S\subseteq\iota_X(K_X)\), size datum over \(X\) at \(s\), and compatible Hida realization \(R\) of frame \(f\) with twisting places away from \(q_0\) and characters in \(\Xi\), there is a complete choice at \((1,C_F,b(s))\) satisfying the coefficient-field condition above.
\end{theorem}

\begin{proof}[Outline]
For each level key, use \ref{B8.11}, \ref{B8.16}, \ref{B8.15} and \ref{B8.17a} to obtain data and fields. Adjoin character, phase and jet values and toric ratios from \ref{B1.14k} proved in \ref{B1.15a}, then take the union of field bases and transport by level agreement. The complete proof is in \S\ref{prf:B8.17}.
\end{proof}

\begin{definition}[Data for the finite rational tests]\label{def:p8c2-test-stage}
The coefficient-field condition on the base and frame requires a choice of field data for which the embedded base coefficient field is exactly the compositum of the embedded host \(E\), the auxiliary field \(M'\) of the choice, its ray field, the frame character field and the fields of the layer characters for \(1\le j\le m\), together with the prescribed cyclotomic fields. Their orders are the order of the finite group for the remaining rows, \(g_i(g_i-1)\) for \(i=0,1\), and \(3^4,5^4\). The terminal inclusion \(X'\hookrightarrow X_m\) sends each frame vector of \(X'\) to the frame vector of \(X_m\) with the same index.

At \(g_i\), fix the labeling used for the frame and let \(\omega_i\in\F_{g_i}^{\times}\) be the label of its character. Each type place must lie above the frame's specified prime in the distinguished imaginary quadratic field. The Teichmüller character and every polynomial series idempotent are those of this fixed labeling, and at every indexed type place the layer label numbered \(j=0,\ldots,m-1\) is \(\omega_i^{2(j+1)}\). The indexed type places, taken over both primes, must be in bijection with the selected group places, preserving their underlying places of \(E\) and their specified primes in the quadratic field. The element defining the trivial type is a central idempotent of the complex group algebra of the product of the residue-field matrix groups.

The theta-map condition requires a nonzero archimedean theta intertwiner on each successive row arrow and on the terminal arrow, common to all level indices on that arrow. Its matrix coefficients satisfy the smooth tensor norm bound, and it intertwines the target action and kills the diagonal source relations for the compact groups and Lie algebras. At each level index there must be multiplicity maps on these arrows given by theta integration with the prescribed measures and domains. Their composite, followed by the bottom projection, must be realized by a map into the row Hilbert space. On each auxiliary layer there must also be a compact-equivariant polynomial test map taking some unit lowest-type tensor and unit coefficient vector to the distinguished polynomial. At each index its theta integral must take values in the prescribed automorphic multiplicity space and its projection must admit a lift to the corresponding Hilbert model.

For geometric evaluation, at every layer and compatible embedding require a semilinear map from the rational cuspidal forms to the Hilbert space whose evaluation at the chosen vector gives the geometric form. This condition uses the prescribed vector at the good embeddings and a unit vector giving the conjugate form at the other embeddings. At good embeddings evaluation must be injective. At every embedding the resulting realization must carry each finite linearly independent family to a complex linearly independent family and its image must span. Require a rational action of the finite group on the forms which preserves integrality at every valuation and under realization gives the geometric action with the prescribed character factor. The forms integral at all the specified valuations must be a full lattice, and the dimension of the geometric cuspidal space must equal the multiplicity dimension in the real model at each layer.
\end{definition}

\begin{theorem}[Finite rational test functionals]\label{B8.18}
Assume the inputs of Theorem~\ref{B8.17}, the unitary Siegel-set estimate of Proposition~\ref{B1.15b} and the finite-group cuspidal character input. In the eventual large-field construction take a base and frame with the coefficient field and frame normalization of Definition~\ref{def:p8c2-test-stage}. Require the set of twist primes to contain \(g_0,g_1,d\) and every rational prime at most \(X''(n)\). Fix a pinned choice of columns and any completion by the remaining row, division, local, type and geometric data. For any admissible fields and set of places, take a stage \(z\) satisfying the finite-group, theta-map and geometric evaluation conditions of Definition~\ref{def:p8c2-test-stage}. Let \(EE\) be the complex row space attached to the realization and face of \(z\). At each level index \(t\), let \(FR_t\) be the source over the coefficient field \(K\), let \(j_t\) realize it in the complex source, and let \(\Theta_t\) be the theta map from that source to \(EE\). For the stage embedding \(v_0:K\hookrightarrow\C\) there are \(N\in\N,\ \psi_i\in EE^\vee\) with
\begin{equation}\label{eq:p8c2-row-tests}
 \bigcap_{0\le i<N}\ker\psi_i=0,\qquad
 \psi_i(\Theta_t(j_t x))\in v_0(K)
 \quad(t\text{ arbitrary},\ 0\le i<N,\ x\in FR_t).
\end{equation}
\end{theorem}

\begin{proof}[Outline]
Use separation from \ref{B8.15} and spanning columns to get finitely many functionals with zero common kernel. Project to good summands and compare models by \ref{B8.11} and \ref{B8.16}, then normalize and transfer signs to obtain \(v_0(K)\)-valued evaluations. The complete proof is in \S\ref{prf:B8.18}.
\end{proof}

\section{Preserving a prescribed grid}

\begin{lemma}[Finite field requirements for full grids]\label{B8.20}\label{B8.19}
For a grid \(G\), let \(v_0\) be its base place, \(w(v_0)\) the chosen place above it, and \(\chi_{a+1}\) its layer characters. Put \(\gamma(G)=(v_0,w(v_0),(\chi_{a+1})_{0\le a<m})\). A full grid over \(R\) (Definition~\ref{B8.21d}) admits a compatible complete choice with the evaluation and square-field conditions. The bottom clause concerns the base index \(\mu_0\). Suppose there is a common pair \((\delta,t_0)\) admitting bottom pairs for every \(i,j<m\). Then one can choose a common pair and bottom witnesses for all \(i,j\) so that both norms of each witness lie in \(\iota(K)\). Put \(h_{ij}=\lfloor(r_{i+1,j+1}+r_{i,j})/2\rfloor\). For \(i,j<m\), \(1\le s'\le h_{ij}\) and either twist, the corresponding untwisted and twisted toric values \(x,y\) have \(y/x\in\iota(K)\).

Fix the scale \(s\), the constant \(C_F\), a completion \(X\), a frame \(f\) and its allowed set \(\Xi\) of twist characters. Among the realizations of Theorem~\ref{B8.17} take those with square-recursive size data (Definition~\ref{B8.21d}), the row size conditions of Theorem~\ref{B8.15} and the prescribed host with \(F=E^+\). For the quadratic character \(\chi_d\) of \(\Q(\sqrt{-d})\), require \(\bigl|\Re(L'(1,\chi_d)/L(1,\chi_d))\bigr|\le x^\varepsilon\). Say that prescribed grids extend at this completion if, for each such realization, every grid \(G_0\) at \((1,C_F,b(s))\) admits a full grid \(G\) at the same parameters with \(\gamma(G)=\gamma(G_0)\). The grid extension condition at \(C_F\) requires this property at every completion of the construction. The finite field requirement for prescribed grids is the assertion of Theorem~\ref{B8.17} with the conclusion that prescribed grids extend at every completion whose coefficient field contains the requested finite set.

Suppose at the requisite thresholds that each permitted grid admits a Weil system and initial choice extending to a complete choice over a number field \(N_1\subset\C\) with the same grid choice and Weil system. Suppose also that the square-field condition holds over a number field \(N_2\subset\C\) for each square-recursive datum, admissible level character system with base index and permitted grid at \((1,C_F,b)\). Assume level agreement transports the grid, Weil system, evaluations and field conditions. If level keys for finite character sets and grid keys at fixed gap are finite, then the finite field requirement for prescribed grids holds.
\end{lemma}

\begin{proof}
Take the union of bases of \(N_1N_2\) for representatives of the realized pairs of keys at gap one. In a coefficient field containing the union, transport the full representative to the realization with its pair of keys.
\end{proof}

\begin{lemma}[Algebraic ratios on a prescribed grid]\label{lem:p8c2-square-field}
Under the inputs of Theorem~\ref{B8.21}, for each permitted grid and square-recursive datum there is a number field \(N_2\subset\C\) satisfying the square-field condition of Lemma~\ref{B8.20}.
\end{lemma}

\begin{proof}
For a pair \(i,j\) and integer \(s\) in the required range, take \(\chi=\mu_{0,j}\xi_{\mu_0}/\chi_{m-i}\) and either grid twist \(\delta_c\). Here \(\mu_{0,j}\) is the row character, \(\xi_{\mu_0}\) the base twist and \(\chi_{m-i}\) a grid character. Let \(c_{\tau,a}\) be the floor of the midpoint of the adjacent grid weights at positions \(a-1,a\). Put \(\kappa_\tau=\max(0,e_{\tau,j+1}+c_{\tau,m-i}-s)\). The grid gives
\[
 s\le h_{ij}\le N_*,\qquad
 e_{\tau,j+1}+c_{\tau,m-i}\ge g_{\rm sep}>s .
\]
Thus \(\kappa_\tau>0\) and \(\kappa_\tau+s=e_{\tau,j+1}+c_{\tau,m-i}\). Unfolding the toric sums gives the primitive Euler products with the prescribed auxiliary factors in the convergence half-plane. Take the auxiliary set empty and continue to \(s\). The values are \(u=L(s,\chi), u_c=L(s,\chi\delta_c)\). Their weights give nontrivial angular type, so both values are nonzero by Lemma~\ref{lem:p8c2-hecke-nonvanishing}.

Let \(d\) be the host auxiliary prime and \(K_d=\Q(\sqrt{-d})\) embedded by the grid CM type. Use \(\Lambda=\mathcal O_{K_d}\) as the integral period lattice. With \(C_0,M\) as in Definition~\ref{B1.14k}, apply its ordinary transport assertion, proved in Lemma~\ref{B1.15a}. It gives
\[
 \beta(v)=\frac{(C_0v)^{12}}{\operatorname{disc}(\Lambda)^M}\in\overline{\Q}
 \qquad(v=u,u_c).
\]
Hence \((u_c/u)^{12}=\beta(u_c)/\beta(u)\) and the ratio is algebraic. Uniqueness of the toric sums gives finitely many ratios. If bottom witnesses exist with common \((\delta,t_0)\), choose one such family. For each class representative choose a prime at which it is integral. The coordinate integrality bound with smallness puts every coordinate \(a\) in \(\overline{\Q}\). The squared norm at a representative is a sum of \(a\bar a\), and the class sum weights it by the reciprocal of the finite stabilizer order. Each selected norm is algebraic. Adjoin the norms for the finitely many pairs and the toric ratios.
\end{proof}

\begin{theorem}[Eventual existence of evaluation grids]\label{B8.21}
Under the inputs of Theorem~\ref{B8.17} there is \(C_{F1}\in\R\) such that for every \(C_F\ge C_{F1}\) the grid extension condition of Lemma~\ref{B8.20} eventually holds in the large-field sense of Theorem~\ref{B8.07}.
\end{theorem}

\begin{proof}[Outline]
Use \ref{B8.17} and \ref{lem:p8c2-square-field} to obtain \(N_1,N_2\) and transport the conditions by level agreement. Ray class finiteness and the finite base-place choices give finitely many grid keys, so \ref{B8.20} gives the fourth-stage field request. The complete proof is in \S\ref{prf:B8.21}.
\end{proof}

\begin{lemma}[The finite height of the higher toric ratios]\label{B8.21a}
Use a fitted grid (Definition~\ref{B8.21d}) under the size hypotheses of Theorem~\ref{B8.22}, with coefficient field \(K\) and \(f=[E^+:\Q]\). At finite \(w\) of \(K\) let \(v_w\) be the uniformizer-one valuation and \(\ell_w=\log N\mathfrak p_w/[K:\Q]\). For \(i,j<m\) put \(h=h_{ij}\) as in Lemma~\ref{B8.20}. For \(1\le s\le h\) use \(u,u_c\) of Lemma~\ref{lem:p8c2-square-field}, with \(z_c=u_c/u\in K\) through the specified embedding. Put \(z_T=1,\ z_{T,c}=z_c\). Use the weights and normalization of Definition~\ref{B1.14k}, namely
\[
 \kappa_\tau=\max(0,e_{\tau,j+1}+c_{\tau,m-i}-s),\quad
 M=2sf+2\sum_\tau\kappa_\tau,\quad
 C_0=|D_F|^{(2s-1)/2}\prod_\tau\Gamma(2s+\kappa_\tau)
 (2\pi/\sqrt d)^{\sum_\tau\kappa_\tau}.
\]
For \(v=u,u_c\), let \(P_v\) be the monic minimal polynomial over \(K\) of \(\beta(v)=(C_0v)^{12}/\operatorname{disc}(\Lambda)^M\), with \(\Lambda=\mathcal O_{K_d}\). Put \(y_w(v)=v_w(P_v(0))/(12\deg P_v)\) and \(g_T(w)=\operatorname{trunc}_0y_w(u)\), truncating toward zero. Outside \(1\le s\le h\) use both ratios one and \(g_T=0\).

Beyond the joint grid and point thresholds, for every sufficiently large \(C_1\) the in-range ratios are nonzero, and for every integer \(s\) and \(z=z_T,z_{T,c}\) the sum has finite support and satisfies
\begin{equation}\label{eq:p8c2-toric-height}
 \sum_w\ell_w|v_w(z)+g_T(w)|\le B_0,\qquad B_0=C_1fb(\log b+\log x).
\end{equation}
\end{lemma}

\begin{proof}[Outline]
Use \ref{lem:p8c2-square-field} and \ref{B1.14c} and take denominator bounds from \ref{B1.14k} proved in \ref{B1.15a}. Use Hecke estimates, CM transport, Chowla–Selberg and the product formula to bound valuations, then scale by \(z_c^{12}\) with a common shift. The complete proof is in \S\ref{prf:B8.21a}.
\end{proof}

\section{The functional bound}

Let \(RR\) be the glued row lattice of positive rank with infinite frames \(X_{E,v}\) and normalizations \(\kappa_v>0\). The squared-norm condition is
\begin{equation}\label{eq:p8c2-fun}
 \mu_{\max,RR}\le x^{C_H}b-\langle\log\kappa\rangle_\infty .
\end{equation}
Here \(C_H\) is the family height exponent and \(\langle\log\kappa\rangle_\infty\) is the arithmetic mean of \(\log\kappa_v\) over the infinite places of the coefficient field.

\begin{definition}[Members and member points]\label{B8.21b}
Fix rank parameters consisting of the target character order \(o\ge1\), half-degree \(n=o\), real degree \(f=2n\), even starting rank \(k\ge2\), even chain length \(m\ge2\), and odd auxiliary rank \(t\ge1\), with bottom row rank \(r=k+m-1\), auxiliary rank \(n_V=r+t\), and column-layer ranks \(N(j)=k+j(t+1)\) (\(j\in\N\)), satisfying
\[
 2\le m,\quad m+3\le k,\quad r+5\le t,\quad 3m<t,\quad m\le f\le k,\quad m^2\le k+m(t+1).
\]
A Harish--Chandra parameter tuple at these ranks is a quadruple \(H=(\lambda^{\rm HC},v,b,e)\) with \(\lambda^{\rm HC}=(u_0,\dots,u_{k-1})\in(\Z+\tfrac12)^k\), \(v=(v_0,\dots,v_{k-2})\in\Z^{k-1}\subset\Q^{k-1}\), \(b=(b_0,\dots,b_{m-1})\in\Z^m\), and \(e=(e_0,\dots,e_{m-1})\in\Z^m\) satisfying \(u_{a+1}<v_a<u_a\) (\(0\le a<k-1\)), \(-b_0<u_{k-1}\), \(e_j-b_j\ge(t+3)/2\) (\(0\le j<m\)), and \(b_{j+1}-e_j\ge(t+3)/2\) (\(0\le j<m-1\)).

An \emph{arithmetic datum} \(D\) at these ranks consists of the following host, Hermitian-space, row-and-division, member, local-Galois, and theta-jet data:
\begin{enumerate}
\item A \emph{host datum} over an abelian CM extension \(E/\Q\) (unramified over its maximal totally real subfield \(F=E^+\) at finite places, with \([F:\Q]=f\), CM conjugation \(x\mapsto\bar x\), CM type \(\Psi\) of cardinality \(f\), and distinguished embedding \(\varphi_0\in\Psi\), \(\varphi_0:E\hookrightarrow\C\)) and imaginary quadratic subfield \(K_d=\Q(\sqrt{-d})\subset E\) (for a prime \(d>4\) with \(d\equiv3\pmod4\)): an auxiliary prime \(q\ne3\) with \(\sqrt{-q}\in E\) and \(4tr+1<q,d\); a positive integer factor \(Q_S\le n^2\) and split bound \(B_{\rm sp}\in\N\) such that every odd prime \(p\le B_{\rm sp}\) satisfies \(\bigl(\frac{-d}{p}\bigr)=1\) and every odd prime \(l\mid Q_S\) satisfies \(4l\le B_{\rm sp}\); an ideal \(\mathfrak b\subseteq\mathcal O_E\) preserved by Galois automorphisms projecting to \(\Psi\) and factoring the relative different as \(\mathfrak D_{E/K_d}=\mathfrak b\bar{\mathfrak b}\) with \(\mathfrak b+\bar{\mathfrak b}=\mathcal O_E\); a finite set \(\mathcal P_{\rm lev}\) of primes coprime to \(Q_S\) containing distinct theta primes \(g_0\ne g_1\) and distinct prime-twist primes \(p_{\rm de}\ne p'_{\rm de}\) with \(p_{\rm de}\equiv p'_{\rm de}\equiv1\pmod{15}\) and \(\{p_{\rm de},p'_{\rm de}\}\cap\{g_0,g_1\}=\varnothing\) (together with a split auxiliary prime \(q_0\notin\{g_0,g_1,p_{\rm de},p'_{\rm de}\}\)); a nonzero primitive ideal \(\mathfrak c_K\subseteq\mathcal O_{K_d}\) coprime in norm to \(30\,q\,d\,Q_S\) and \(\mathcal P_{\rm lev}\) whose prime factors in \(\mathcal O_{K_d}\) have prime norms splitting in \(K_d\); finite top-position and thick sets \(TT\) and \(T_2\subset\N\), word-length integer \(L_w\ge2\), and constant \(C_{T_2}\in\R\).
\item \emph{Hermitian spaces} over \(E/F\) (each a finite-dimensional \(E\)-vector space \(V\) with nondegenerate conjugate-symmetric \(E\)-sesquilinear form \(h_V\), \(E\)-basis, signature map, unimodular \(\mathcal O_E\)-lattice \(L_0(V)\), and volume \(\operatorname{vol}(V)>0\)): definite column spaces \(X^{\rm col}_j\) (\(0\le j\le m\)) of rank \(N(j)\), a subspace \(X'\hookrightarrow X^{\rm col}_m\) of rank \(r+mt\) isometric onto the orthogonal complement of an anisotropic vector \(z\in X^{\rm col}_m\), reference row spaces \(Y_j\) (\(0\le j\le m\)) of rank \(r+jt\) and signature \((r+jt-\lfloor(m-j)/2\rfloor,\lfloor(m-j)/2\rfloor)\) at every embedding \(E\hookrightarrow\C\) with isometric embeddings \(X^{\rm col}_j\hookrightarrow Y_j\) (\(0\le j<m\)), an isometry \(X'\xrightarrow{\sim}Y_m\), and a vector \(z_m\in Y_0\) of norm \(-T_{u_0}\) (with \(T_{u_0}\) the base level prime of item~(4)(b)) and \(r_S=(m-2)/2\) pairwise orthogonal hyperbolic planes in \(Y_0\) orthogonal to \(X^{\rm col}_0\) and \(z_m\), row-layer spaces \(W_j\) (\(1\le j\le m\)) of rank \(r\) finite-adelically isometric to \(Y_0\), and auxiliary spaces \(V^{\rm sp}_j\) (\(1\le j\le m\), written \(V^{\rm sp}_j\) to distinguish them from the core spaces \(V_i\) of Section~3) of rank \(n_V=r+t\) and sign \(\epsilon_V(j)\in\{\pm1\}\) making \(L_0(V^{\rm sp}_j)\) orthonormal for \(\epsilon_V(j)h_{V^{\rm sp}_j}\).
\item \emph{Coefficient field, row data, and division lattices}: a totally complex Galois number field \(K/\Q\) with involution \(x\mapsto x^*\), distinguished embedding \(\iota:K\hookrightarrow\C\) (\(\iota(x^*)=\overline{\iota(x)}\)), integral subring \(O_{\rm row}\to K\) containing \(\mathcal O_E\), base subfield \(K_0\subseteq K\), embedding \(E\hookrightarrow K_0\subseteq K\) pulling \(\iota\) back to \(\varphi_0\) and carrying \(\bar x\) to \(x^*\), and finite-place set \(\Sigma_f(K)\simeq\operatorname{Spec}_{\ne0}(\mathcal O_K)\) with valuations \(v_w\), ramification indices \(e_w\), and involution \(\bar w\); a finite CM-embedding index set \(\mathcal T\) (\(\mathcal T\xrightarrow{\sim}\Psi\), \(|\mathcal T|=f\)), finite word set \(\mathcal W\) with base word \(u_0\in\mathcal W\), and Harish--Chandra tuples \(H_{u,\tau}=(\lambda^{\rm HC}_{u,\tau},v_\tau,b_\tau,e_\tau)\) (\(u\in\mathcal W,\ \tau\in\mathcal T\)) agreeing with \(H_\tau=H_{u_0,\tau}\) on \(v_\tau,b_\tau,e_\tau\) and on the first and last entries \(u_{u,\tau,0},u_{u,\tau,k-1}\) of \(\lambda^{\rm HC}_{u,\tau}\); a tuple set \(\mathcal T_{\rm tup}\) with base index \(\mu_0\), prime-twist indices \(\mu_*,\mu'_*\), and conjugation-invariant place maps \(\mu_{\rm at},\mu_{\rm tw}:\Sigma_f(K)\to\mathcal T_{\rm tup}\) with \(\mu_{\rm at}(w)\in\{\mu_0,\mu_{\rm tw}(w)\}\); finite-dimensional complex inner-product spaces \(\mathcal E\) (dimension \(D_{\rm rows}=(1-\delta_p)D_{\rm all}\), \(0\le\delta_p\le1\)), \(\mathcal E_0\), ambient spaces \(\mathcal A_\mu\) with symmetric idempotents and \(O_{\rm row}\)-lattices, multiplicity spaces \(\mathcal M_{\mu,j}\) (\(\mathcal M_{\mu,0}\simeq\mathcal E\), \(\mathcal M_{\mu,m}\simeq\mathcal A_\mu\), \(\mathcal M_{\mu,j+1}\to\mathcal M_{\mu,j}\)), word column spaces \(\mathcal M_u\) with pairings \(\mathcal M_u\times\mathcal A_\mu\to\C\), \(\C\)-bilinear pairings \(B_u:\mathcal M_u\times\mathcal E\to\C\), units \(\zeta_{u,\mu}\in K^\times\) (\(\zeta_{u,\mu_0}=1\)), and layer spaces \(M_j\) and \(M_{t,\mu,j}\) (with \(K\)-forms, \(O_{\rm row}\)-lattices, symmetric idempotents, eigenvalues, integer local transition indices, and rational invariants \(\eta_j(w),L_{\sharp}(l,e)\in\Q_{\ge0}\)); and Hermitian \(\mathcal O_K\)-lattices \(RR\) and \(\mathcal N_u\) (\(u\in\mathcal W\), each a \(K\)-vector space \(V\) with \(K\)-basis \(\mathbf b\), invertible infinite-place frame matrices \(X_v\in\GL_d(\C)\), and \(\mathcal O_K\)-submodule \(\mathcal L\subset V\)) with \(K\)-linear period maps \(V(\mathcal N_u)\to V(RR)^\vee\), heights \(h_u\in\R\), parameter \(\eta_F\in\R\), and \(\iota\)-semilinear basis-spanning realizations \(j_{\mathcal E}:V(RR)\to\mathcal E\) (isometric at \(\iota\)) and \(j_{N,u}:V(\mathcal N_u)\to\mathcal M_u\).
\item At each word \(u\in\mathcal W\), a word-level \emph{member} \((C_u,A_u,P_u)\) together with \emph{member forms}:
\begin{enumerate}
\item a \emph{column chain} \(C_u\) assigning to each \(i\ge0\) a finite class set \(\mathcal C_{u,i}\) (with \(\mathcal O_E\)-lattices \(\Lambda_{u,i}(c)\subset X^{\rm col}_i\) in \(D\) for \(0\le i\le m\)), finite stabilizer groups \(\Gamma_{u,i}(c)\) (\(c\in\mathcal C_{u,i}\)), an integer \(D_{\Gamma,u}>0\) divisible by \(|\Gamma_{u,i}(c)|\) for all \(0\le i\le m\) and \(c\in\mathcal C_{u,i}\), fibre sets \(Y_{u,i}\) with factors \(d_{u,i}(y,\tau)\in\C^\times\), unit-norm determinant factors \(\rho_{u,i}(c),\rho'_{u,i}(c')\) for \(i<m\) given in \(D\) through an abelian determinant group \(D_t\) with characters \(\tilde\mu_{i+1},\chi_D:D_t\to\C^\times\) and class map \(\delta_{u,i}:\mathcal C_{u,i}\to D_t\) by
\[
 \rho_{u,i}(c)=\tilde\mu_{i+1}(\delta_{u,i}(c))^{-1}\chi_D(\delta_{u,i}(c))^{-N(i+1)/2},\qquad
 \rho'_{u,i}(c')=\tilde\mu_{i+1}(\delta_{u,i+1}(c'))\chi_D(\delta_{u,i+1}(c'))^{N(i)/2},
\]
unfolding maps \(\mathrm{unf}_{u,i}(c,c'):\C^{Y_{u,i}}\to\C^{Y_{u,i+1}}\), and shift exponents \(s^\sharp_{u,i}(\tau)=e_{\tau,i}-N(i+1)/2\in\Z\) and \(s^{\rm div}_{u,i}(\tau)=e_{\tau,i}-N(i)/2\in\Z\) (\(0\le i<m\));
\item a \emph{top-level column datum} \(A_u\) specifying a prime \(T_u\in\mathcal P_{\rm lev}\setminus\{p_{\rm de},p'_{\rm de}\}\) (with \(T_u=\pi_u\bar\pi_u\) for \(\pi_u\in\mathcal O_E\cap K_d\) satisfying \(\pi_u\equiv1\pmod{Q_S\,q\,d\,(g_0 g_1)^2\mathcal O_E}\), \(c_u\bar c_u T_u=T_{u_0}\) for \(c_u\in E^\times\), and \(T_u^\star\le T_u\le(1+(fk)^{-1})T_u^\star\)), a finite place set \(\mathcal V_{T,u}\) of cardinality \(f\) with finite fields \(\F_v\) of cardinality \(T_u\), nontrivial additive characters \(\psi_{T,v}:\F_v\to\C^\times\), actions \(\rho_{T,v}:\F_v^\times\to\GL_\C(\C^{\mathcal C_{u,0}\times Y_{u,0}})\), a finite-dimensional \(\rho_{T,v}(\F_v^\times)\)-stable subspace \(M_{u,0}\subset\C^{\mathcal C_{u,0}\times Y_{u,0}}\) (paired with its ramified subspace \(M_{{\rm ram},u}\simeq\mathcal M_u\)), an integer \(d_u^\star>0\), a common real factor \(D^\star>0\), and a common price parameter \(\Gamma_P>1\);
\item a \emph{period map} \(P_u\) on \(C_u\), consisting of complex vector spaces \(M_{{\rm t},u}\simeq\mathcal A_{\mu_0}\), \(W_{z,u}\), and \(W_{{\rm row},u}\), \(\C\)-linear maps \(\mathrm{ev}_{z,u}:\C^{\mathcal C_{u,m}\times Y_{u,m}}\to W_{z,u}\) and \(\mathrm{low}_u:W_{z,u}\to W_{{\rm row},u}\), and a \(\C\)-bilinear pairing \(\operatorname{per}_u:W_{{\rm row},u}\times M_{{\rm t},u}\to\C\), accompanied in \(D\) at every \(\mu\in\mathcal T_{\rm tup}\) by a twisted chain \(C_{u,\mu}\) (\(C_{u,\mu_0}=C_u\)), a base-layer isomorphism \(e_{0,u,\mu}:\C^{\mathcal C_{u,0}\times Y_{u,0}}\xrightarrow{\sim}\C^{\mathcal C_{u,\mu,0}\times Y_{u,\mu,0}}\), and a period map \(P_{u,\mu}\) with \(M_{{\rm t},u,\mu}\simeq\mathcal A_\mu\) realizing the \(\mu\)-th column-period pairing on \(\mathcal M_u\times\mathcal A_\mu\); and
\item \emph{member forms} on \((C_u,A_u,P_u)\): for each \(i\ge0\), \(v\in\mathcal V_{T,u}\), and \(c\in\mathcal C_{u,i}\), rank-\(N(i)\) \(\F_v\)-spaces \(U_{u,i,v,c},K_{u,i,v,c}\) with a nonzero covector on \(U_{u,i,v,c}\), a nonzero vector in \(K_{u,i,v,c}\), and raising/lowering maps from \(U_{u,i,v,c}\) and \(K_{u,i,v,c}^\vee\) to \(\mathcal C_{u,i}\) and \(\GL_\C(\C^{Y_{u,i}})\); \(\F_v^\times\)-dilations on \(\mathcal C_{u,i}\) and \(\C^{Y_{u,i}}\) preserving \(|\Gamma_{u,i}(c)|\) and inducing \(\rho_{T,v}\) at \(i=0\); a family of bijections \(\mathcal V_{T,u}\xrightarrow{\sim}\mathcal T\) indexed by \(w\in\Sigma_f(K)\); antitone integer dominant weights \(\lambda_{u,j}(\tau)\in\Z^{N(j)}\) given for \(0\le j\le m\) and \(0\le a<N(j)\) by \(\lambda_{u,j}(\tau)_a=u_{u,\tau,a}-\rho_{k,a}\) (\(a<k\), with rank-\(k\) half-sum vector \(\rho_{k,a}=(k-1)/2-a\)) and \(\lambda_{u,j}(\tau)_a=-e_{\tau,q}+\tfrac{k}{2}+\tfrac{(2q+1)(t+1)}{2}\) (\(k\le a<N(j),\ q=\lfloor(a-k)/(t+1)\rfloor\)), plus a weight in \(\Z^{N(m)-1}\); fibre subspaces \(F_{u,i}\subseteq\C^{\mathcal C_{u,i}\times Y_{u,i}}\) (\(F_{u,0}=M_{u,0}\)) with sesquilinear fibre forms \(\langle\cdot,\cdot\rangle_{V,u,i,c}\) and stabilizer-weighted inner product \(\langle\varphi,\psi\rangle_{C,u,i}=\sum_{c\in\mathcal C_{u,i}}|\Gamma_{u,i}(c)|^{-1}\langle\varphi(c),\psi(c)\rangle_{V,u,i,c}\); a finite point-class set mapping to \(\mathcal C_{u,m}\) with injective stabilizer maps, a point subspace \(W_{{\rm pt},u}\subseteq W_{{\rm row},u}\) with sesquilinear form \(\langle\cdot,\cdot\rangle_{P,u}\), and good Hermitian \(\mathcal O_K\)-lattices \(L_{u,i}\) and \(L_{{\rm pt},u}\) whose \(\iota\)-semilinear realizations carry their \(K\)-bases to \(\C\)-bases of \(F_{u,i}\) and \(W_{{\rm pt},u}\) and match at \(\iota\) the real parts of \(\langle\cdot,\cdot\rangle_{C,u,i}\) and \(\langle\cdot,\cdot\rangle_{P,u}\), with \(V(\mathcal N_u)\) embedding \(K\)-linearly and norm-isometrically onto the preimage of \(M_{{\rm ram},u}\) in \(V(L_{u,0})\).
\end{enumerate}
\item \emph{Split-group, \(p\)-adic Galois, Haar, and theta-jet packages}: a split-group package with integer spacing \(\Delta\in\Z_{>0}\) satisfying \(v_{\tau,l}-v_{\tau,l+1}=4\Delta\) (\(0\le l<k-2\)), an injective word coding \(\ell:\mathcal W\hookrightarrow\{0,1\}^{k-2}\) setting \(u_{u,\tau,l+1}=v_{\tau,l}-\tfrac12-\Delta\mathbf 1_{\ell_u(l)=1}\), and a Hilbert space \(H_S\) (\(\dim_\C H_S=D_{\rm all}\), containing an anti-unitarily isometric copy of \(\mathcal E\)) receiving \(M_{u,0}\); \(p\)-adic period-ring (\(A_{\inf},B_{\rm cris},B_{\rm dR}\)), Galois-bridge, local Hecke, base-change, twisted-module, and finite-group type packages with \(G_{0,g}\simeq\prod_{k=0,1}\GL_r(\Z/g_k\Z)^f\) and central series idempotents at \(g_0,g_1\); compatible Haar measures on \(\U(P,Q)\) (probability on \(\U(N,0)\)); global theta-jet situations at signal and inert primes between ranks \(N(j-1)\) and \(N(j)\) (\(1\le j\le m\)); and local level modules of rank \(\dim_\C M_{u,0}\) at inert primes \(p\ge5\) outside \(\mathcal P_{\rm lev}\) not dividing \(T_u\).
\end{enumerate}
A \emph{size datum} \(L_{\rm sz}\) over \(D\) fixes a scale \(s=(x,\varepsilon,J)\) (\(x>1,\ 0<\varepsilon<1/52,\ J\ge1,\ L=x^\varepsilon,\ b=e^L,\ a=xb,\ X'=\exp(x^{\varepsilon/3})\)), integer top weights \(a_t,b_t\in\Z\) (\(t\in TT\)) in \([a,a+a/(kx^{13})]\) and \([b',\infty)\) (with \(b'=xb/X'\) and \(X''(n)=xnX'\) as in Definition~\ref{def:p8c2-point-sizes}) fixing \(u_{u_0,\tau,0}=a_t+(k-1)/2\) and \(u_{u_0,\tau,k-1}=-b_t-(k-1)/2\) at top embeddings \(\tau=\tau_1(t)\), a finite support \(S_{\rm supp}\subset\Sigma_f(K)\) for the local division-gain, finite-factor, error, and twist-indicator functions \(G_{D,j},F_j,E_{u,j},\tau_{\rm loc}\) on \(\Sigma_f(K)\) and local transition indices (with \(G_{D,j}=0\) above \(g_0,g_1\), \(F_j(w)=0\) for \(1\le j\le m\) when \(\mu_{\rm at}(w)=\mu_0\) or when \(\operatorname{char}(w)>X''(n)\) outside \(\{g_0,g_1,d\}\), and \(F_j(\bar w)=F_j(w)\)), \(b_{\tau,j}\le b^2\), a drop set avoiding all \(T_u\), depth \(j_0\ge\log_5(n\,N(m)/m)\) with \(10j_0+20\le L_w\), \(\Re(L'(1,\chi_d)/L(1,\chi_d))\), input \(\mathcal O_K\)-submodules \(\mathrm{NuLat}_u\subseteq V(\mathcal N_u)\), thick density \(\delta_2=|T_2|/n\in(0,1]\), width \(\eta>0\), center bound \(S_c=4k\Delta>0\) bounding \(0\le v_{\tau,a},u_{u,\tau,a'}\le S_c\) (\(0<a'<k-1\)) and layer-weight windows of width \(\eta\) on \(w_j(\tau)=e_{\tau,j-1}\) (\(1\le j\le m\)), subsets \(TT_w\subseteq TT\), \(1\le t\le x^{C_t}\), \(\log\Gamma_P=C_\Gamma m^5fL\) (\(C_\Gamma>0\)), partition weights \(\ge1\), and local length ratios in \(\Q\).

A \emph{construction context} \(\mathcal K\) fixes a scale \(s\), family constant record, gap/spacing exponents \(c_D,c_h\in\R\), preliminary host-word-parameter datum \(P_{\rm pre}\) with \(\mathfrak c_K=\mathcal O_{K_d}\), auxiliary prime \(q_1\in\mathcal P_{\rm lev}\setminus\{g_0,g_1,p_{\rm de},p'_{\rm de}\}\), integer box weights \(a,b\), injective word coding \(c:\mathcal W\hookrightarrow\{0,1\}^{k-2}\) meeting the host-and-window checks, and level-prime choices \(T_u=\pi_u\bar\pi_u\) with \(d_u^\star,D^\star\) given by the standard unitary mass-index products and \(\Gamma_P=\exp(C_\Gamma m^5fL)\). Over \(\mathcal K\), a \emph{built point} \(\beta\) consists of: a seven-stage tuple (base field \(K\) with \(E\hookrightarrow K_s\subseteq K_0\subseteq K\) containing the group/prime cyclotomic roots and capitulation generators, word columns and member forms, remaining row spaces, division stage, split-group stage, local Galois-bridge stage, and Haar/theta-jet stage) assembling an arithmetic datum \(D\); a character frame over \(P_{\rm pre}\) and finite-place twisting data (a finite twist-prime set \(\mathcal P_{\rm tw}\), split prime-ideal sets \(Q(w)\subset\operatorname{Spec}_{\ne0}(\mathcal O_E)\) coprime to \(\operatorname{char}(w)\), exponents \(k(w)\in\N\), and conjugation-invariant tame Hecke characters \(\psi_w\) of \(E\) of order dividing \(\operatorname{char}(w)^{k(w)}\), trivial off \(\mathcal P_{\rm tw}\)); a core specialization fixing tuple twists \(\xi_\mu\) (\(\xi_{\mu_0}=1,\ \xi_{\mu_*}=\delta,\ \xi_{\mu'_*}=\delta'\)), active twist places \(W_{\rm tw}\) (forcing \(\{g_0,g_1,d\}\) and all primes \(\le X''(n)\) into \(\mathcal P_{\rm tw}\) and vanishing of lift base terms off \(W_{\rm tw}\)), \(\Delta\), \(\eta_F=6x^{-C_F}\), and Teichm\"{u}ller labels \(\omega^{2(j+1)}\) at \(g_0,g_1\); an object layer equipping \(D\) with glued-lattice characterization on \(RR\), Shimura/divisor models, automorphic/theta actions, local-to-global chain kernels, base-level row side, ambient lattices, and unitary genus-mass Haar gauges; and a ledger specialization equipping \(D\) with a size datum \(L_{\rm sz}\) at \((s,a,b)\), \(\mathrm{NuLat}_u\) comparison with \(M_{{\rm ram},u}\) away from \(T_u\), and local error/valuation/recursion ties. For a parameter predicate \(\kappa\), point predicate \(Q\), and bundle linkage \(\mathrm{Lk}\) at \(c\in\R\), a \emph{built member} \((D,O,\mathfrak b)\) is realized by \(\mathcal K\) satisfying \(\kappa\), a built point \(\beta\) over \(\mathcal K\) satisfying \(Q\) with datum \(D\) and object package \(O\), and a bundle \(\mathfrak b\) linked to \(\beta\) at \(c\) via \(\mathrm{Lk}\), supplying the notions of member and built member in Section~1 of Chapter~\ref{ch:p8c1}.

Along a family \(F_m\) at \(x>x_1\) and \(i\in I_x\), a \emph{member point} \(p\) over \((D,L_{\rm sz})\) packages: an indefinite row-side full Hecke datum over a finite group \(G_p\simeq G_{0,g}\), finite tuple set \(T\) with base index \(\mu_0\in T\), subgroups \(P_0,P_1\le G_p\) with characters \(\psi_0:P_0\to K^\times\) and \(\psi_1:P_1\to K^\times\), definite subspaces, two-level indefinite data, definite columns, Hilbert-space periods, and three-sided division ties; integers \(\Delta,h_0\in\N\) matching \(\Delta_D=\Delta\) and \(\eta=mh_0\), stabilizer divisibility into \(D_{\Gamma,u_0}\), \(F=E^+\), and host parameter-interlacing and \(100r^4\)/\(256r^4\) gap bounds; a projective-lattice core over \(K_s\subseteq K\) and level-\(3\) real realization \(R\) weakly characterized on the induced face, with standard sign-\(-1\) additive character \(\psi\), sign-\(-1\) residue-trace compatibility of \(\psi_0\) at \(P_g\), coprimality of \(Q_t\) to \(Q_S\,q\,N\mathfrak c_K\), unit phases, isometric top projection, orthonormal frames and unimodular \(\mathcal O_E\)-lattices on \(X'\) and on the definite spaces \(V^{\rm sp}_{i+1}\) (\(0\le i<m\), positive definite for \(\epsilon_V(i+1)h_{V^{\rm sp}_{i+1}}\)), and \(F\)-diagonal Gram matrices on \(W_{i+1}\) and \(V^{\rm sp}_{i+1}\); weight tuples \(\kappa_i\) and extremal vectors \(v_0(i,\sigma)\) (\(\sigma:K\hookrightarrow\C\)) with geometric cusp-form evaluation, conjugate, finite-group action, equivariance, and phase pins; residue-field matrix-group and series-idempotent pins at \(g_0,g_1\); and a tuple bijection \(\mathcal T_{\rm tup}\xrightarrow{\sim}T\) identifying \(\mu_0\) with the base tuple and sending \(\mu_{\rm at}(w)\) to a gluing index \(\mu(w)\in T\) at each finite place \(w\) of \(K\). Fixing a representative complex embedding \(v:K\hookrightarrow\C\) at each infinite place \(v\mid\infty\) of \(K\), an invertible frame matrix \(X_v\in\GL_d(\C)\) on \(K^d\) defines \(\|y\|_v=\|X_v^{-1}v(y)\|_2\) for the Euclidean norm \(\|\cdot\|_2\) on \(\C^d\), with positive definite metric \(H_v=(X_v^{-1})^\dagger X_v^{-1}\) (\(\|y\|_v^2=v(y)^\dagger H_v v(y)\)) and canonical frame \(X_v=H_v^{-1/2}\). For \(p\in\mathcal P\) of Definition~\ref{def:p8c3-admissible}, the level-zero data furnish at each \(t\in T\) a source evaluation \(\mathcal O_K\)-lattice \(M_t\subset K^{n_V(t)}\) with Hermitian matrix \(H_V(t)\in M_{n_V(t)}(K)\) having positive definite embeddings \(H_V(t)_v=v(H_V(t))\) and frames \(X_V(t)_v=H_V(t)_v^{-1/2}\) (\(v\mid\infty\)), a surjective quotient matrix \(P_t\in M_{d_Q(t)\times n_V(t)}(K)\) (\(P_{t,v}=v(P_t)\)) equipping \(K^{d_Q(t)}\) with the positive definite quotient Hermitian metric
\[
 H_Q(t)_v=\bigl(P_{t,v}H_V(t)_v^{-1}P_{t,v}^\dagger\bigr)^{-1}\in M_{d_Q(t)}(\C)
\]
(inducing \(\|q\|_v=\inf_{P_{t,v}x=q}\|x\|_v\)) and the \emph{direct quotient Hermitian frame}
\[
 X_Q(t)_v=H_Q(t)_v^{-1/2}=\bigl(P_{t,v}H_V(t)_v^{-1}P_{t,v}^\dagger\bigr)^{1/2}\in\GL_{d_Q(t)}(\C)
\]
satisfying \((X_Q(t)_v^{-1})^\dagger X_Q(t)_v^{-1}=H_Q(t)_v\), and a row-realization matrix \(\iota_t:K^{d_Q(t)}\to RR\otimes_{\mathcal O_K}K\) such that \(\iota_{\mu_0}\) is injective, \(\iota_{\mu_0}P_{\mu_0}\) is onto \(RR\otimes_{\mathcal O_K}K\), \(\mu(w)=\mu_0\) outside finitely many twist places, and \(RR_w=\iota_{\mu(w)}P_{\mu(w)}M_{\mu(w),w}\) at every finite place \(w\) of \(K\).
\end{definition}

\begin{definition}[Character grids and grids]\label{B8.21c}
Fix an arithmetic datum \(D\) over \(E/F\) (Definition~\ref{B8.21b}). A \emph{character frame} over \(D\) specifies prime ideals \(\bar\pi_k\subset\mathcal O_{K_d}\) above \(g_k\) (\(k=0,1\)), injective characters \(\omega_k:(\Z/g_k\Z)^\times\hookrightarrow\C^\times\), and unitary Hecke characters of \(E\): base row characters \(\mu_{0,j}\) (\(0\le j<m\)) trivial on \(\mathbb A_F^\times\), of angular infinity exponent \(-2e_{\tau,j}\) across \(\Psi\), tamely ramified everywhere, unramified outside primes of \(E\) above \(g_0g_1\), and sending any integer \(a\) coprime to \(g_k\) at any prime \(w\mid\bar\pi_k\) of \(E\) to \(\omega_k(a\bmod g_k)^{-2(j+1)}\); the everywhere unramified splitting character \(\eta\) restricting to the quadratic character of \(E/F\) on \(\mathbb A_F^\times\) with angular infinity exponent \(1\) across \(\Psi\); the everywhere unramified central character \(\chi_\Psi\) trivial on \(\mathbb A_F^\times\) with angular infinity exponent \(2\) across \(\Psi\); and finite-order prime-twist characters \(\delta,\delta'\) at \(p=p_{\rm de},p'_{\rm de}\) trivial on \(\mathbb A_F^\times\), tamely ramified everywhere, unramified outside primes of \(E\) above \(p\), and restricting on \(\mathcal O_{E,v}^\times\) at every prime \(v\mid p\) of \(E\) to a character of exact order \(15\). A \emph{level-character datum} \(\mathcal L\) over \(D\) indexed by \(T\) adjoins to a character frame level-twist Hecke characters \(\xi_t\) of \(E\) (\(t\in T\)), compatible unitary characters on \(\mathbb A_E^\times=(\mathbb A_F\otimes_F E)^\times\) for \(\mu_{0,j},\eta,\xi_t\), a finite set \(P_g\) of finite places of \(E\) above \(g_0,g_1\), finite ramification-place sets \(Q_t\) (\(t\in T\)) of finite places of \(E\), residue-field characters at finite places of \(E\), and reference \(\mathcal O_E\)-lattices \(\Lambda_{Y,j},\Lambda_{X'},\Lambda_{V,i},\Lambda_{W,i}\) with finite-adelic frames in \(\GL(\cdot,E\otimes_F\mathbb A_{F,f})\) on the reference row spaces \(Y_{\min(j,m)}\) (\(j\in\N\), equal for \(0\le j\le m\) to \(Y_j\) of rank \(r+jt\) and signature \((r+jt-\lfloor(m-j)/2\rfloor,\lfloor(m-j)/2\rfloor)\)), the definite space \(X'\) (rank \(r+mt\), signature \((r+mt,0)\)), the auxiliary spaces \(V^{\rm sp}_{i+1}\) (\(0\le i<m\), rank \(n_V=r+t\)), and the row-layer spaces \(W_{i+1}\) (\(0\le i<m\), rank \(r\)). Fix a real parameter triple \((c_{\rm gap},c_{\rm Dep},b)\).

A \emph{character grid} over \(\mathcal L\) at \((c_{\rm gap},c_{\rm Dep},b)\) consists of a grid place datum and a grid character datum at gap \(c_{\rm gap}\):
\begin{enumerate}
\item The \emph{grid place datum} specifies a finite base place \(v_0\) of \(F\) dividing the split auxiliary prime \(q_0\) of \(D\) (with \(q_0\notin\{g_0,g_1,p_{\rm de},p'_{\rm de}\}\), so \(v_0\notin S_1\) below) and a lift \(v\mapsto w(v)\) from finite places of \(F\) to finite places of \(E\). Let \(S_1\) be the finite set of prime ideals of \(\mathcal O_F\) dividing \(g_0g_1p_{\rm de}p'_{\rm de}\), put \(S_0=\{v_0\}\cup S_1\), and put \(W_0=\{w(v):v\in S_0\}\). For every \(v\in S_0\), \(w(v)\) lies over \(v\) and admits a \emph{split local idempotent} \(e\in\mathbb A_E=\mathbb A_F\otimes_F E\), defined by \(e^2=e\), \(e\bar e=0\), \(e+\bar e=\varepsilon_v\) (for the place idempotent \(\varepsilon_v\in\mathbb A_F\) equal to \(1\) at \(v\) and \(0\) at \(\infty\) and at all finite places \(v'\ne v\)), and the condition that a nonzero \(x\in\mathcal O_E\) lies in \(\mathfrak P_{w(v)}\) if and only if the finite part of \(ex^{-1}\) lies outside \(\widehat{\mathcal O}_E\); for \(v\mid g_k\) (\(k=0,1\)), \(w(v)\in P_g\); for each \(p\in\{p_{\rm de},p'_{\rm de}\}\), a single prime ideal \(\mathfrak p\subset\mathcal O_{K_d}\) lies below \(w(v)\) for all \(v\mid p\); and for every \(t\in T\) and \(\mathfrak P\in Q_t\) over \(v\) of \(F\), one has \(v\ne v_0\), while \(v\in S_1\) implies \(\mathfrak P=w(v)\). Each \(q\in S_0\) has prime ideal \(\mathfrak q\subset\mathcal O_F\) of prime-power absolute norm \(N\mathfrak q\) (also written \(q\) inside logarithms and \(w\)-adic valuations), selected prime ideal \(\mathfrak P_{w(q)}\subset\mathcal O_E\), grid conductor exponent \(c_{v_0}=1\) and \(c_q=2\) for \(q\in S_0\setminus\{v_0\}\), and prescribed positive integer grid depth \(d_{\rm pl}(q)\) at \((c_{\rm gap},c_{\rm Dep},b)\) (equal to the base depth at \(v_0\) and \(2\) on \(S_1\), as in Lemma~\ref{B8.15b} and Section~4). Put
\[
 Q_{\rm lev}=\prod_{q\in S_0}(N\mathfrak q)^{c_q},\quad
 c_S=\log Q_{\rm lev}=\sum_{q\in S_0}c_q\log N\mathfrak q,\quad
 \mathfrak f_F=\prod_{q\in S_0}\mathfrak q^{c_q}\subset\mathcal O_F,\quad
 \mathfrak m_E=\prod_{q\in S_0}\mathfrak P_{w(q)}^{c_q}\subset\mathcal O_E .
\]
\item The \emph{grid character datum} specifies Hecke characters \(\chi_1,\dots,\chi_m\) of \(E\) paired with continuous unitary characters on \(\mathbb A_E^\times\) agreeing at split finite places such that for each \(1\le i\le m\): \(\chi_i\) is trivial on \(F^\times\) at every finite and infinite place of \(F\); its archimedean component satisfies \(\chi_{i,\infty}(k)=\prod_{\tau\in\mathcal T}(\sigma_\tau(k)/|\sigma_\tau(k)|)^{2\,\mathrm{cen}(\tau,i)}\) (\(k\in E^\times\)), where \(\mathrm{cen}(\tau,i)=c_{\tau,i}\in\Z\) is the floor of the \(i\)-th adjacent-weight midpoint of the signed Harish--Chandra grid parameter attached to \(H_\tau\), \(d_s=r+6\), and gap \(c_{\rm gap}\) (as in Lemma~\ref{lem:p8c2-square-field}); for every finite place \(\mathfrak P\) of \(E\) not above \(S_0\), \(\chi_{i,\mathfrak P}(\mathcal O_{E,\mathfrak P}^\times)=\{1\}\); for \(\mathfrak P\) above \(S_1\), \(\chi_{i,\mathfrak P}(1+\mathfrak P^2\mathcal O_{E,\mathfrak P})=\{1\}\) and \(\chi_{i,\mathfrak P}(1+\mathfrak P\mathcal O_{E,\mathfrak P})\ne\{1\}\); for \(\mathfrak P\mid p_{\rm de}p'_{\rm de}\), \(\chi_{i,\mathfrak P}(x)=1\) whenever \(x\in E_{\mathfrak P}^\times\) satisfies \(x^{N\mathfrak P-1}=1\); and for \(\mathfrak P\) above \(v_0\), \(\chi_{i,\mathfrak P}(1+\mathfrak P\mathcal O_{E,\mathfrak P})=\{1\}\) and \(\chi_{i,\mathfrak P}(\mathcal O_{E,\mathfrak P}^\times)\ne\{1\}\).
\end{enumerate}

For Hermitian spaces \(A,B\) over \(E/F\) with Gram matrices \(J_A,J_B\) in their frames, an integer \(c\ge0\), and \(\epsilon\in\{\pm1\}\subset E\), a \emph{Hermitian orthogonal-extension decomposition} from \(A\) to \(B\) of codimension \(c\) and sign \(\epsilon\) consists of an isometric \(E\)-linear embedding \(\iota_{A,B}:A\hookrightarrow B\) (with frame matrix \(M_{\iota_{A,B}}\in M_{\operatorname{rank}B\times\operatorname{rank}A}(E)\) and \(\operatorname{rank}A+c=\operatorname{rank}B\)) and a matrix \(a\in M_{\operatorname{rank}B\times c}(E)\) satisfying \(M_{\iota_{A,B}}^\dagger J_B a=0\) and \(a^\dagger J_B a=G_{c,\epsilon}\), where \((\cdot)^\dagger\) is CM-conjugate transpose and the \(c\times c\) signed anti-diagonal Gram matrix \(G_{c,\epsilon}\) has \((p,q)\)-entry (\(0\le p,q<c\)) equal to \(\epsilon\) if \(p+q+1=c\) and \(p=q\), \(1\) if \(p+q+1=c\) and \(p\ne q\), and \(0\) if \(p+q+1\ne c\); it assigns to any \(\mathcal O_E\)-lattice \(\Lambda_A\) on \(A\) and \(c\)-tuple \((\mathfrak I_0,\dots,\mathfrak I_{c-1})\) of nonzero fractional ideals of \(\mathcal O_E\) the extension \(\mathcal O_E\)-lattice \(\iota_{A,B}(\Lambda_A)\oplus\bigoplus_{p=0}^{c-1}\mathfrak I_p a_p\) on \(B\) (with \(a_p\) the \(p\)-th column of \(a\)). For nonzero fractional ideals \(\mathfrak b_1,\mathfrak b_2\) of \(\mathcal O_E\) in \(E\) and \(c\ge0\), define the \emph{staircase fractional-ideal tuple} of length \(c\) by
\[
 \operatorname{stair}(\mathfrak b_1,\mathfrak b_2;c)_p
 =\mathfrak b_1^{c-p}\mathfrak b_2^{-(p+1)}\qquad(0\le p<c).
\]
Recall from Lemma~\ref{B8.15e} the host oscillator fractional ideal \(\mathfrak a_D=(\mathfrak c_K\mathcal O_E)\mathfrak b^{-1}\) of \(E\) and its CM conjugate \(\bar{\mathfrak a}_D\), and let \(\mathfrak a_{\rm st},\bar{\mathfrak a}_{\rm st}\) be the conjugate pair of first-column fractional ideals of \(\mathcal O_E\) attached to \((v_0,w)\) at \((c_{\rm gap},c_{\rm Dep},b)\) (as in Lemma~\ref{B8.15f}). Put \(n_{\rm rk}=r+1=k+m\), \(d_s=r+6=k+m+5\), \(r_{i,j}=r+jt+id_s\) (\(0\le i,j\le m\)),
\[
 n_i=n_{\rm rk}+i(d_s-1)=(r+1)+i(r+5)\quad(0\le i\le m),\qquad
 n'_j=n_m+j(t+1)\quad(0\le j\le m).
\]
A \emph{grid} \(G\) over \(\mathcal L\) at \((c_{\rm gap},c_{\rm Dep},b)\) adjoins to a character grid at \((c_{\rm gap},c_{\rm Dep},b)\) a vertex system of Hermitian spaces with \(\mathcal O_E\)-lattices indexed by: grid nodes \((i,j)\) (\(1\le i\le m,\ 0\le j\le m\)) of rank \(r_{i,j}\) and negative signature \(\lfloor(m-j)/2\rfloor\) (with row \(i=0\) given by the reference spaces \(Y_j\) (\(0\le j\le m\)) of \(\mathcal L\) and lattice \(\Lambda_{Y,j}\)); core nodes \(P_i\) (\(0\le i\le m\)) of rank \(n_i\) and negative signature \(\lfloor(m-i+1)/2\rfloor\); and positive definite top nodes \(X_j\) (\(1\le j\le m\)) of rank \(n'_j\) and negative signature \(0\) (with initial top node \(X_0=P_m\) of rank \(n'_0=n_m\)), linked by four families of Hermitian orthogonal-extension decompositions identifying each target \(\mathcal O_E\)-lattice with the corresponding extension lattice of its source: from \(Y_0\) to \(P_0\) of codimension \(1\) and sign \(+1\) with ideal \(\mathcal O_E\) on \(\Lambda_{Y,0}\); from \(P_i\) to \((i,0)\) (\(1\le i\le m\)) of codimension \(i-1\) and sign \(-1\) with tuple \(\operatorname{stair}(\mathfrak a_{\rm st},\bar{\mathfrak a}_{\rm st};i-1)\); from \(X_j\) to \((m,j)\) (\(1\le j<m\)) of codimension \(m-1-j\) and sign \(-1\) with tuple \(\operatorname{stair}(\mathfrak a_D,\bar{\mathfrak a}_D;m-1-j)\) (and at \(j=0\) the \(i=m\) decomposition from \(X_0=P_m\) to \((m,0)\)); and from \((m,m)\) to \(X_m\) of codimension \(1\) and sign \(+1\) with ideal \(\mathcal O_E\). At a member point \(p\) with scale weight \(b\) and \(C_F\in\R\), a grid on \(p\) at \(C_F\) is a grid over the level-character datum of \(p\)'s real realization at \((1,C_F,b)\).
\end{definition}

\begin{definition}[Definite CM theta pairs and coordinate integrality allowances]\label{B8.21ca}
Let \(D\) be an arithmetic datum with imaginary quadratic subfield \(K_d=\Q(\sqrt{-d})\subset E\), real degree \(f=[F:\Q]\), CM-embedding index set \(\mathcal T\), and distinguished embedding \(\varphi_0:E\hookrightarrow\C\).
\begin{enumerate}
\item An ideal \(\mathfrak c\subseteq\mathcal O_{K_d}\) is a \emph{curve ideal} for \(D\) if \(\mathfrak c\ne0\) and \(\gcd(N\mathfrak c,2d)=1\) (in particular \(\mathfrak c=\mathcal O_{K_d}\) is a curve ideal). For \(\mathfrak c\subseteq\mathcal O_{K_d}\) and \(\delta\in\C\), set
\[
 \Lambda(\mathfrak c,\delta)=\operatorname{span}_{\Z}\{\delta\,\varphi_0(c):c\in\mathfrak c\}\subset\C,\qquad
 \operatorname{area}(\Lambda)=\inf\bigl\{|\Im(\bar w_1 w_2)|>0:w_1,w_2\in\Lambda\bigr\},
\]
with Weierstrass invariants \(g_2(\Lambda),g_3(\Lambda)\in\C\), modular discriminant \(\operatorname{disc}(\Lambda)\), and Weierstrass function \(\wp_\Lambda\). A \emph{curve datum} \((\mathfrak c,\delta)\) over \(D\) consists of a curve ideal \(\mathfrak c\subseteq\mathcal O_{K_d}\) and a \emph{discriminant twelfth-root period} \(\delta\in\C\) defined by \(\delta^{12}=\operatorname{disc}(\Lambda(\mathfrak c,1))\). A point \(t_0\in\C\) is an \emph{unkilled two-torsion point} for \(\Lambda=\Lambda(\mathfrak c,\delta)\) if \(2t_0\in\Lambda\) and \(\frac{1\pm i\sqrt d}{2}t_0\notin\Lambda\); write \(\theta_0(\Lambda,t_0)\) for the associated normalized genus-one CM theta series and \(\nu(\Lambda,t_0)\) for its Gaussian weighted mean square, the average over a fundamental domain \(\mathcal F_\Lambda\) of \(\Lambda\) in \(\C\) of \(|\theta_0(\Lambda,t_0)(z)|^2\exp(-\pi|z|^2/\operatorname{area}(\Lambda))\). An integer \(Q\ge1\) is an \emph{admissible level modulus} for \(D\) if every rational prime \(p\mid Q\) satisfies \(p=2\) or \(\bigl(\frac{-d}{p}\bigr)=1\); for a rational prime \(p\), a level \(Q'\in\N\) is \emph{\(p\)-compatible with \(Q\)} if its \(p\)-adic valuation in \(\N\) (with \(v_p(0)=0\)) satisfies \(v_p(Q')\le v_p(Q)\).
\item For \(d_{\rm rk}\in\N\), let \(G_{\mathbb A}(d_{\rm rk})\le\GL_{d_{\rm rk}}(\mathbb A_{E,f})\) (\(\mathbb A_{E,f}=E\otimes_\Q\mathbb A_{\Q,f}\)) be the finite-adelic unitary group preserving the identity Hermitian form \(I_{d_{\rm rk}}\) over \(\mathbb A_{\Q,f}\) under \(E\)'s CM conjugation, with rational subgroup \(G_{\rm rat}(d_{\rm rk})\le G_{\mathbb A}(d_{\rm rk})\) generated by \(F\)-rational unitary matrices and integral subgroup \(G_{\rm int}(d_{\rm rk})\le G_{\mathbb A}(d_{\rm rk})\) generated by \(\widehat{\mathcal O}_E\)-integral unitary matrices; for a rational prime \(p\), an element \(h\in G_{\mathbb A}(d_{\rm rk})\) is \emph{integral over \(p\)} if \(hk\) is trivial above \(p\) for some \(k\in G_{\rm int}(d_{\rm rk})\). Given a curve datum \((\mathfrak c,\delta)\) over \(D\), an unkilled two-torsion point \(t_0\in\C\) for \(\Lambda=\Lambda(\mathfrak c,\delta)\), and \(Q\in\N\), a \emph{definite CM theta pair} \(P\) at \((\mathfrak c,\delta,t_0,Q)\) consists of:
\begin{enumerate}
\item source and target ranks \(a_P,d_P\in\N\) with \(a_P\le d_P\) and \(d_P\) even, and complex configuration space \(V_{\C}(a_P,d_P)=\prod_{\tau\in\mathcal T}M_{d_P\times a_P}(\C)\);
\item an \emph{odd-rank CM vacuum} \(e_0:V_{\C}(a_P,d_P)\to\C\) at \((\mathfrak c,\delta,t_0)\), defined as an element of the CM theta space of the standard lattices \(\mathcal O_E^{a_P},\mathcal O_E^{d_P}\) on the CM torus of \((\mathfrak c,\delta)\) having Taylor coefficients in \(\overline{\Q}\subset\C\), local factorization and unit-phase normalization at \(t_0\), and squared norm on the base period lattice \(\Lambda_0(a_P,d_P)\) equal to \(\nu(\Lambda,t_0)^{f\,a_Pd_P}\);
\item a \emph{compatible Weil representation} \(W=(\chi_V,\chi'_V,\epsilon_W,\epsilon'_W,\Omega_W)\) for \(e_0\ne0\), consisting of unitary Hecke data \(\chi_V,\chi'_V\) (with archimedean exponent tuple \((m_{W,\tau})_{\tau\in\mathcal T}\in\Z^{\mathcal T}\) for \(\chi'_V\)), signs \(\epsilon_W,\epsilon'_W\in\{\pm1\}\), and a \(\C\)-linear group action \(\Omega_W\) of \(G_{\mathbb A}(a_P)\times G_{\mathbb A}(d_P)\) preserving the CM theta space on the torus of \((\mathfrak c,\delta)\), locally intertwining the Heisenberg--Weyl translation operators \(\mathcal W_{\mathfrak c,\delta}(\pi_\delta(z))\) at torus points \(\pi_\delta(z)\) of \(\mathfrak c\)-congruent shifts \(z,z'\in M_{d_P\times a_P}(E)\), fixing \(e_0\) under integral elements supported at inert primes, and acting on \(e_0\) at elements \((g,g')\) supported at split primes (moving \(\mathcal O_E^{a_P},\mathcal O_E^{d_P}\) to \(L_1,L'_1\) with split-adjusted determinants \(x,x'\in\mathbb A_E^\times\) trivial mod inert primes) by \(\chi_V(x')^{\epsilon_W}\chi'_V(x)^{\epsilon'_W}\) times the lattice-change operator from \(\Lambda_0(a_P,d_P)\) to \(\Lambda(L_1,L'_1)\);
\item a degree tuple \(n_P=(n_{P,\tau})_{\tau\in\mathcal T}\in\N^{\mathcal T}\) satisfying \(2n_{P,\tau}=\epsilon'_W m_{W,\tau}-d_P\) for all \(\tau\in\mathcal T\);
\item a test function \(F_P:V_{\C}(a_P,d_P)\to\C\) invariant under \(\Omega_W(k,1)\) for all \(k\in G_{\rm int}(a_P)\) and \emph{preserved at level \(Q\) relative to \(e_0\)}, meaning that there exist \(k_{\rm pr}\in\N\), \(r_1,\dots,r_{k_{\rm pr}}\in\overline{\Q}\), and \(z_1,\dots,z_{k_{\rm pr}}\in M_{d_P\times a_P}(E)\) with \(Q\cdot\pi_\delta(z_s)\in\Lambda_0(a_P,d_P)\), \(Q^{a_P d_P}r_s\) integral over \(\Z\) (\(1\le s\le k_{\rm pr}\)), \(\bigl(\sum_{s=1}^{k_{\rm pr}}|\rho(r_s)|\bigr)^2\le Q^{a_P d_P}\) for every ring homomorphism \(\rho:\overline{\Q}\to\C\), and \(F_P(y)=\sum_{s=1}^{k_{\rm pr}}r_s\bigl(\mathcal W_{\mathfrak c,\delta}(\pi_\delta(z_s))e_0\bigr)(y)\); and
\item a subgroup \(J_P\le G_{\rm int}(d_P)\) of finite index \(1\le[G_{\rm int}(d_P):J_P]\le Q^{d_P^2+4}\) acting on \(F_P\) under \(\Omega_W(1,j)\) (\(j\in J_P\)) by scalars of modulus \(1\).
\end{enumerate}
For \(N\in\N\) and \(B\in\R\), \(P\) is \emph{\((N,B)\)-small} if \(d_P\le2N\) and \(n_{P,\tau}\le2B\) for all \(\tau\in\mathcal T\).
\item For a definite CM theta pair \(P\), write \(\mathrm{Cls}_D(d_P,J_P)=G_{\rm rat}(d_P)\backslash G_{\mathbb A}(d_P)/J_P\) for its finite double-coset class set, \(|\Gamma_D(J_P,h)|\in\N_{>0}\) for the finite stabilizer order of \(h\in G_{\mathbb A}(d_P)\), and \(\mathrm{Tup}_D(a_P,d_P,n_P)\) for the finite index set of \(a_P\times a_P\) minor tuples \(T\) of multidegree \(n_P\) (as in Lemma~\ref{B8.17a}), with normalized minor-jet \emph{class-tuple coordinate} \(c_P(h,T)\in\C\) of \(\Omega_W(1,h)F_P\) at \(0\), pointwise squared sum \(\mathrm{sq}_P(h)=\sum_{T\in\mathrm{Tup}_D(a_P,d_P,n_P)}|c_P(h,T)|^2\), and representative-independent \emph{norm}
\[
 \mathrm{norm}(P)=\sum_{q\in\mathrm{Cls}_D(d_P,J_P)}|\Gamma_D(J_P,h_q)|^{-1}\,\mathrm{sq}_P(h_q)\ge0\qquad(h_q\in q).
\]
On a grid \(G\) over \(\mathcal L\) at \((c_{\rm gap},c_{\rm Dep},b)\), for \(t\in T\), \(0\le i,j<m\), an ideal \(\mathfrak m\subseteq\mathcal O_E\), and \((\delta,t_0,Q)\), define the larger and smaller bottom source ranks and splitting characters by
\[
 \begin{gathered}
 a_+(i,j)=r_{i+1,j+1}-2\bigl\lfloor\tfrac{m-j-1}{2}\bigr\rfloor,\qquad
 a_-(i,j)=r_{i,j}-2\bigl\lfloor\tfrac{m-j}{2}\bigr\rfloor,\\
 \alpha_{t,i,j}=\mu_{0,j}\xi_t\eta^{r_{i,j}\bmod2},\qquad
 \beta_{i,j}=\chi_{m-i}\eta^{r_{i,j}\bmod2},\qquad
 \nu_{t,i,j}=\mu_{0,j}\xi_t\chi_{m-i}^{-1}.
 \end{gathered}
\]
A pair \((P,P')\) of definite CM theta pairs at \((\mathcal O_{K_d},\delta,t_0,Q)\) is a \emph{standard bottom pair} on \(G\) at \((t,i,j,\mathfrak m)\) if: \(a_P=a_+(i,j)\), \(a_{P'}=a_-(i,j)\), and \(d_P=d_{P'}=a_+(i,j)+a_-(i,j)\); on finite split ideles, \((\chi_{V,P}^{\epsilon_{W,P}},\chi'_{V,P}{}^{\epsilon'_{W,P}})=(\beta_{i,j}^{-1},\nu_{t,i,j}^{-1})\) and \((\chi_{V,P'}^{\epsilon_{W,P'}},\chi'_{V,P'}{}^{\epsilon'_{W,P'}})=(\alpha_{t,i,j},\nu_{t,i,j}^{-1})\); \(F_P=\sum_{r\in R}\Omega_{W,P}(r,1)\phi_{\rm cos}(P;\mathfrak m,0)\) and \(F_{P'}=\sum_{r'\in R'}\Omega_{W,P'}(r',1)\phi_{\rm cos}(P';\bar{\mathfrak m},a_+(i,j))\) for finite coset-representative sets \(R\subseteq G_{\rm int}(a_P)\) modulo \(\mathfrak m\) and \(R'\subseteq G_{\rm int}(a_{P'})\) modulo \(\bar{\mathfrak m}\) of the coset vacua at target-row offsets \(0\) and \(a_+(i,j)\); and both \(J_P\) and \(J_{P'}\) equal \(\{h\in G_{\rm int}(d_P):h_{u,v}\in\widehat{\mathfrak m}\text{ for all }0\le v<a_+(i,j)\le u<d_P\}\).
\item For \(C\in\R\), integers \(f',N,Q\ge1\), and \(B,L\in\R\), set
\[
 I(C;f',N,B,L,Q)=2Cf'N\bigl(\log(2f'N)+BL+2NB\log Q\bigr).
\]
Let \(S'\) be a finite set of rational primes and \(\delta_{\rm vp}:\N\to\R_{\ge0}\) a function vanishing outside \(S'\) with \(\sum_{p\in S'}\delta_{\rm vp}(p)\log p\le b_{\rm int}\). Say that \(\delta_{\rm vp}\) satisfies the \emph{\(p\)-adic coordinate integrality bound} for \((D,N,B,Q)\) if for every discriminant twelfth-root period \(\delta\in\C\) of \(\mathcal O_{K_d}\), unkilled two-torsion point \(t_0\in\C\) for \(\Lambda(\mathcal O_{K_d},\delta)\), rational prime \(p\) (\(|p|_p=p^{-1}\) on \(\C_p\)), \(p\)-compatible level \(Q'\in\N\), and \((N,B)\)-small definite CM theta pair \(P\) at \((\mathcal O_{K_d},\delta,t_0,Q')\): every class \(q\in\mathrm{Cls}_D(d_P,J_P)\) admits a representative \(h\in q\) integral over \(p\); every \(h\in G_{\mathbb A}(d_P)\) satisfies \(|\Gamma_D(J_P,h)|\ne0\) and \(v_p(|\Gamma_D(J_P,h)|)\le\delta_{\rm vp}(p)\); and for every \(h\in G_{\mathbb A}(d_P)\) integral over \(p\) and \(T\in\mathrm{Tup}_D(a_P,d_P,n_P)\), one has \(c_P(h,T)\in\overline{\Q}\) and
\[
 |\lambda(c_P(h,T))|_p\le p^{\delta_{\rm vp}(p)}
\]
for every ring homomorphism \(\lambda:\overline{\Q}\to\C_p\). Say that \((S',\delta_{\rm vp})\) satisfies the \emph{norm valuation bound} for \((D,N,B,Q)\) with budget \(b_{\rm int}\) if for every discriminant twelfth-root period \(\delta\in\C\) of \(\mathcal O_{K_d}\), unkilled two-torsion point \(t_0\in\C\) for \(\Lambda(\mathcal O_{K_d},\delta)\), totally complex number field \(K\) with \(\sigma:K\hookrightarrow\C\), index set \(\mathcal I\), families \((P_\alpha)_{\alpha\in\mathcal I},(P'_\alpha)_{\alpha\in\mathcal I}\) of \((N,B)\)-small definite CM theta pairs at \((\mathcal O_{K_d},\delta,t_0,Q)\) each having at least one nonzero coordinate, and \(z_{A,\alpha},z'_{A,\alpha}\in K\) with \(\sigma(z_{A,\alpha})=\mathrm{norm}(P_\alpha)\) and \(\sigma(z'_{A,\alpha})=\mathrm{norm}(P'_\alpha)\), the local allowance \(d_\ell(w)=\sum_{p\in S'}\delta_{\rm vp}(p)v_w(p)\ge0\) satisfies \(\sum_{w\in U}\ell_w d_\ell(w)\le b_{\rm int}\) for every finite set \(U\) of finite places of \(K\), together with \(z_{A,\alpha}\ne0\), \(z'_{A,\alpha}\ne0\), \(v_w(z_{A,\alpha})\ge-3d_\ell(w)\), and \(v_w(z'_{A,\alpha})\ge-3d_\ell(w)\) for all \(\alpha\in\mathcal I\) and finite \(w\). Finally, for a grid \(G\) at \((c_{\rm gap},c_{\rm Dep},b)\), define the two integer bottom grid weights at \((i,j,\tau)\) (\(0\le i,j<m,\ \tau\in\mathcal T\)) by
\[
 \begin{aligned}
 w_A(c_{\rm gap};i,j,\tau)&=\Bigl\lfloor\frac{2\,e^{\rm gr}_{\tau,j+1}-(r_{i,j}\bmod2)-a_-(i,j)}{2}\Bigr\rfloor\in\Z,\\
 w_{A'}(c_{\rm gap};i,j,\tau)&=\Bigl\lfloor\frac{-2\,\mathrm{cen}(c_{\rm gap},\tau,m-i)-(r_{i,j}\bmod2)+a_+(i,j)}{2}\Bigr\rfloor\in\Z,
 \end{aligned}
\]
where \(e^{\rm gr}_{\tau,j+1}\in\Z\) is the \((j+1)\)-st upper grid weight of \(\tau\) at gap \(c_{\rm gap}\) (as in Lemma~\ref{lem:p8c2-square-field}), and put
\[
 B_w(c_{\rm gap})=\max\Bigl(1,\;\max_{\substack{0\le i,j<m\\ \tau\in\mathcal T}}\max\bigl(|w_A(c_{\rm gap};i,j,\tau)|,|w_{A'}(c_{\rm gap};i,j,\tau)|\bigr)\Bigr)\in\N_{>0}\subset\R_{\ge1},
\]
abbreviated \(B_w=B_w(1)\) at gap \(1\); a function \(d_\ell\) from finite places of \(K\) to \(\R\) is a \emph{grid integrality allowance} on \(G\) at constant \(C\in\R\) if some \((S',\delta_{\rm vp})\) satisfies both the \(p\)-adic coordinate integrality bound and the norm valuation bound for \((D,N_*,B_w(c_{\rm gap}),Q_{\rm lev})\) with budget \(I(C;f,N_*,B_w(c_{\rm gap}),\log b,Q_{\rm lev})\) and \(d_\ell(w)=\sum_{p\in S'}\delta_{\rm vp}(p)v_w(p)\) for every finite place \(w\) of \(K\).
\end{enumerate}
\end{definition}

\begin{definition}[Square-recursive size data, full grids, and fitted grids]\label{B8.21d}
Let \(D\) be an arithmetic datum and \(L_{\rm sz}\) a size datum over \(D\) at scale \(s\) of exponent \(\varepsilon\), with \(L=x^\varepsilon\) and \(b=\exp(L)\).
\begin{enumerate}
\item The size datum \(L_{\rm sz}\) is \emph{square-recursive} if: \(L_{\rm sz}\) satisfies the row size conditions, comprising \eqref{eq:p8c2-size-rows-a}--\eqref{eq:p8c2-size-rows-c}, the pointwise size inequalities of Definition~\ref{def:p8c2-point-sizes}, the existence of compatible window/tower parameters with weight-gap constant \(10\), \(\mathfrak c_K=\mathcal O_{K_d}\), and \(b_{\tau,l}>0,\ e_{\tau,l}\le4b\) for all \(\tau\in\mathcal T\) and \(0\le l<m\); \(D\) satisfies \(F=E^+\) together with the host parameter-interlacing and gap conditions (\((t+3)/2\) bounds on \(e_{\tau,j}-b_{\tau,j}\) and \(b_{\tau,j+1}-e_{\tau,j}\), consecutive row-parameter and cloud-boundary gaps at least \(100r^4\), and node gaps at least \(\max(4N(m)^3+2t+6,256r^4)\)); and the quadratic character \(\chi_d\) of \(K_d\) satisfies \(\bigl|\Re(L'(1,\chi_d)/L(1,\chi_d))\bigr|\le x^\varepsilon\). On an indefinite row-side datum over a finite group \(G_p\) with subgroups \(P_0,P_1\le G_p\) and characters \(\psi_0:P_0\to K^\times,\ \psi_1:P_1\to K^\times\), the \emph{level-\(2\) face} \(F_{\rm fc}\) retains \((P_0,\psi_0,P_1,\psi_1)\) while setting the definite-column index set to \(\varnothing\), the period subspace to \(0\subseteq\C\), and all eleven auxiliary face scalars to \(0\). Given a real realization \(R\) over \(G_p\), tuple index set \(T\) with base index \(\mu_0\), and subfield \(K_s\subseteq K\), the \emph{standard sign-\(-1\) normalization at \((P_0,\psi_0)\)} requires \(R\)'s additive character \(\psi:\mathbb A_F/F\to\C^\times\) to have archimedean part \(x_\infty\mapsto\exp\bigl(-2\pi i\sum_\rho\theta_\rho(x_\infty\otimes1)\bigr)\) and, for every family of split local idempotents \((e_{\mathfrak P})_{\mathfrak P\in P_g}\) and every \(u\in P_0\),
\[
 \iota(\psi_0(u))=\prod_{\mathfrak P\in P_g}\psi_{e_{\mathfrak P}}\bigl((N\mathfrak P)^{-1}\,S_{\rm tr}(e_G(u)_{\mathfrak P})\bigr),
\]
where \(S_{\rm tr}(g_{\mathfrak P})\in\mathcal O_E\) lifts the superdiagonal residue sum of \(g_{\mathfrak P}\in\GL_r(\mathcal O_E/\mathfrak P)\). The triple \((L_{\rm sz},R,F_{\rm fc})\) satisfies the \emph{Hida row-setting conditions} if \(L_{\rm sz}\) satisfies the row size conditions, \(D\) satisfies the host conditions with \(F=E^+\), \(R\) satisfies the Hida level and residue-group type conditions of Theorem~\ref{B8.17} and Definition~\ref{def:p8c2-test-stage} (including \(Q_t\subseteq\operatorname{Ram}(\xi_t)\subseteq Q_t\cup\overline{Q}_t\) disjoint from \(q_0\), and realized along \(\mathcal P\) from the series condition of Definition~\ref{def:p8c2-series}), and \(F_{\rm fc}\) carries the standard sign-\(-1\) normalization at \((P_0,\psi_0)\). At the base level \(\ell_0=\ell(P_0,\psi_0)\) of \(F_{\rm fc}\), write \(M_{t,{\rm lv}}(F_{\rm fc})\subset K^{n_V(t)}\) for the source evaluation \(\mathcal O_K\)-lattice at \(t\in T\) (equal at \(p\in\mathcal P\) to \(M_t\) of Definition~\ref{B8.21b}), and for \(x\in K^{n_V(t)}\) define the row-source form \(\mathrm{form}_{F_{\rm fc}}(t,x)=\sum_{\alpha=1}^{d_{\rm amb}(t)}(v_{E,t}(\ell_0)x)_\alpha b_{t,\alpha}\in FR_t\) via the level-inclusion matrix \(v_{E,t}(\ell_0)\in M_{d_{\rm amb}(t)\times n_V(t)}(K)\) into the \(K\)-basis \((b_{t,\alpha})\) of \(FR_t\) in Theorem~\ref{B8.18}.
\item Let \(G\) be a grid at \((1,C_F,b)\) with \(W_0=\{w(v):v\in S_0\}\). A \emph{Weil system} \(W\) on \(G\) extending \(R\) assigns to each grid arrow \((S,T_{\rm sp})\) of \(G\) with sign \(\epsilon\in\{\pm1\}\) and unitary Hecke splitting characters \((\chi_S,\chi_{T_{\rm sp}})\) restricting on \(\mathbb A_F^\times\) to \((\eta^{\operatorname{rank}S},\eta^{\operatorname{rank}T_{\rm sp}})\) a continuous unitary oscillator Weil representation \(\omega_{S,T_{\rm sp}}\) of \(\U(S)(\mathbb A_F)\times\U(T_{\rm sp})(\mathbb A_F)\) on theta tests on \(\operatorname{Hom}_E(S,T_{\rm sp})(\mathbb A_F)\) for \(R\)'s sign-\(-1\) character \(\psi\) (in the model of Lemma~\ref{B8.15d} and Definition~\ref{B8.16c}, agreeing with \(R\) on row arrows) that intertwines Heisenberg shifts by \(hw_0g^{-1}\), fixes evaluation at \(0\) under \(\U(S)(F)\times\U(T_{\rm sp})(F)\), preserves \(L^2\) norm on the rational additive quotient, and acts at every split local idempotent \(e\in\mathbb A_E\) on \(e\)-invariant functionals by a positive real multiple of the local determinant factor of \((\chi_S,\chi_{T_{\rm sp}})\). For \(t\in T\) and any finite place \(w'\) of \(E\) over \(v\) of \(F\), set \(\mathcal P(G,t)=W_0\cup Q_t\), \(e_{\rm dep}(G,w')=d_{\rm pl}(v)\) if \(w'\in W_0\) and \(0\) otherwise, and \(d_G(w')=\max(e_{\rm dep}(G,w'),1)\). A \emph{finite-adelic lattice frame} of an \(\mathcal O_E\)-lattice \(\Lambda\subset E^d\) is an element \(\beta\in\GL_d(E\otimes_F\mathbb A_{F,f})\) satisfying \(\widehat{\Lambda}=\beta(\widehat{\mathcal O}_E^d)\). A \emph{grid choice} \(S_0^{\rm gr}\) on \(U\subseteq T\) consists of: finite-adelic lattice frames \(\beta_{i,j}\) on the grid nodes \((i,j)\) (\(0\le i,j\le m\), equal to \(\mathcal L\)'s frames at \(i=0\)) and \(\beta_{P,j}\) on the top nodes \(X_j\) (\(0\le j\le m\)) such that at every \(t\in U\) and \(w'\in\mathcal P(G,t)\) with uniformizer \(\varpi_{w'}\in E_{w'}\), for each \(0\le j<m\) (with \(d_j\) the codimension-\(c_j\) decomposition from \(X_j\) to \((m,j)\), \(c_j=m-1-j\)), \(\beta_{m,j,w'}=\bigl(M_{\iota_{d_j}}\beta_{P,j,w'}\;\big|\;a_{d_j}\operatorname{diag}(u_{j,w',0},\dots,u_{j,w',c_j-1})\bigr)\) where \(u_{j,w',p}=\epsilon_{\rm fr}(c_j,p)\varpi_{w'}^{\mathbf 1_{j=0}\,e_{\rm dep}(G,w')(p+1)}\) (\(\epsilon_{\rm fr}(c,p)=-1\) for \(2p+1>c\) and \(+1\) for \(2p+1\le c\)), while at \(j=m\) (for \(d_{\rm top}\) from \((m,m)\) to \(X_m\)) \(\beta_{P,m,w'}=\bigl(M_{\iota_{d_{\rm top}}}\beta_{m,m,w'}\;\big|\;a_{d_{\rm top}}\bigr)\); \emph{standard normalizations} \((N_{\rm std},Z_{\rm norm})\) on \(U\), comprising a finite rational prime set \(S_{\rm bad}\) containing all primes below \(S_0,\bigcup_{t\in U}Q_t\) and the prime divisors of \(Q_S\,q\,d\,g_0 g_1\,N\mathfrak c_K\), global \(E\)-frames on all nodes of \(G\) orthonormal on definite nodes and compatible with \(d_{\rm top}\), finite-adelic unitary elements supported at split places off \(S_{\rm bad}\) carrying those frames' \(\mathcal O_E\)-lattices to \(\widehat{\Lambda}_X\) on definite nodes, and a curve datum \((\mathfrak c_K,\delta_Z)\) with unkilled two-torsion point \(t_{0,Z}\) (Definition~\ref{B8.21ca}(1)), odd-rank CM vacua \(e_{0,\wp}\), oscillator vacua \(\phi_{0,\wp}\), and scalars \(c_{0,\wp}\in\C^\times\) across definite grid pairs \(\wp\) (independent of \(t\) on top pairs) equating evaluation at \(0\) of Fock monomials on \(\phi_{0,\wp}\) with \(c_{0,\wp}\) times the \(\delta_Z\)-period product and the \(t_{0,Z}\)-jet of \(e_{0,\wp}\) at \(0\); and a finite set \(P_A\) of places of \(E\) in the support of \(\mathfrak a_D\) or \(\bar{\mathfrak a}_D\) containing a unique place above each place of \(F\) below the support of \(\mathfrak a_D\). Writing \(\mathcal W_L(\lambda)=\bigotimes_{\tau\in\mathcal T}\nabla_L(\lambda(\tau))\) for the tensor product of Schur--Weyl highest-weight polynomial modules over a field \(L\), a \emph{sum choice} \(S\) on \(U\subseteq T\) extends \(S_0^{\rm gr}\) by a \(\C\)-linear map \(\Phi_{\rm co}:W_{c_X}\to\mathcal W_\C(\lambda_{\rm list}(0,m))\) from \(R\)'s top-column coefficient space to the listed-weight polynomial module at \((0,m)\) that acts, under a \(\iota\)-semilinear bridge family tie \(c_{0,R}:\mathcal W_K(\lambda_R)\to W_{c_X}\) (and \(c_{0,L,i}\) on \(V^{\rm sp}_{i+1}\)) intertwining the unitary actions, pairings, and integral lattices off \(S_b\cup Q_t\), by applying \(\iota\) to polynomial coefficients and reindexing \(\{1,\dots,\operatorname{rank}X'\}\simeq\{1,\dots,r_{0,m}\}\). For \(U=T\), \(S\) satisfies the \emph{evaluation-field condition} at \(\mu_0\) under \(\iota:K\hookrightarrow\C\) (that is, the coefficient-field condition preceding Theorem~\ref{B8.17}) if \(\iota(K)\) contains the normalized local values of \(\mu_{0,j}\) (\(0\le j<m\)), \(\xi_t\) (\(t\in T\)), and \(\chi_i\) (\(1\le i\le m\)) at weight \(0\), \(\eta\) at weight \(1/2\), and \(\chi_\Psi,\delta,\delta'\) at weight \(0\), all \(8\)-th roots of unity, all \((Nw')^{d_G(w')}\)-th roots of unity for \(w'\in W_0\cup\bigcup_{t\in T}Q_t\), \(\sqrt{Nw'}\) for \(w'\in P_A\cup\bigcup_{t\in T}(W_0\cup Q_t)\), the local Heisenberg phases on nonvanishing moved test terms at \(w'\in\mathcal P(G,t)\) and CM jet coefficients on supported coset-expansion classes, the Hecke \(L\)-value ratios \(L(s,\nu_{t,i,j})/L(s,\nu_{\mu_0,i,j})\) (\(0\le i,j<m,\ t\in T,\ 1\le s\le N_*+1\)) whenever \(L(s,\nu_{\mu_0,i,j})\ne0\), and the period-normalized bottom class sums of Lemma~\ref{B8.17a}. On \(X_0=P_m\) of rank \(n_p=n_m=(r+1)+m(r+5)\), let \(\mathrm{Cl}_+(G)\) be its finite congruence double-coset class set in \(\U(X_0)(\mathbb A_{F,f})\), let \(\mathcal D_{\rm test}(G)\) be the finite index set of prescribed row-coset grid test data, put \(\mathcal I_{\rm test}(G)=\mathcal D_{\rm test}(G)\times\mathrm{Cl}_+(G)\), let \(\lambda_{P,G}=(\lambda_{P,G}(\tau))_{\tau\in\mathcal T}\) be the antitone integer dominant weight tuple of \(X_0\) at gap \(1\) with shifted nonnegative tuple \(\nu_{P,G}(\tau)_a=\lambda_{P,G}(\tau)_a-\lambda_{P,G}(\tau)_{n_p-1}\in\N\), and put \(\mathcal I_{\rm coef}(G)=\prod_{\tau\in\mathcal T}\mathcal T_{\rm ss}(n_p,n_p;\nu_{P,G}(\tau))\), where \(\mathcal T_{\rm ss}(n_p,n_p;\nu)\) is the finite set of interlaced \(n_p\times n_p\) nonnegative integer matrices \((m_{a,q})\) with column sums \(\sum_{a=0}^{n_p-1}m_{a,q}=\nu_q\) indexing the Weyl basis of \(\nabla(\lambda_{P,G}(\tau))\). A \emph{complete choice} \(a\) of ambient dimension \(n\in\N\) adjoins to a sum choice \(S\) on \(T\): Haar measures and fundamental domains on \(\U(V_X)(\mathbb A_F)\) for all nodes \(X\) of \(G\) and on the top Bessel unipotent flag of \((m,0)\); a real test vector \(\rho\) and linear functional \(\xi\) on the dual highest-weight module of \(\lambda_{P,G}\) on \(X_0\); downward chain vectors (Fock polynomials on \((Y_{l+1},Y_l)\) for \(0\le l<m-1\) and on \((X',Y_{m-1})\), together with a smooth tensor in the bottom-row archimedean space \(M_X\)); real matching data on \(W\) (arrow vacua matching \(W\)'s row vacua at \(j=0\) and proper submodules of polynomial-route pairs at each \((i,j)\) surjecting onto both routes); class representatives \(x_c\in\U(X_0)(\mathbb A_{F,f})\) (\(c\in\mathrm{Cl}_+(G)\)); a bijection \(e_a:\{1,\dots,n\}\xrightarrow{\sim}\mathcal I_{\rm test}(G)\times\mathcal I_{\rm coef}(G)\); and seesaw constants \(c_{ij,t}\in\C\) (\(0\le i,j<m,\ t\in T\)), meeting the separation, rationality, and series conditions of Theorem~\ref{B8.18} and Definition~\ref{def:p8c2-series}; for \(t\in T\), \(x'\in FR_t\), and \(1\le i\le n\), write \(v_{\rm row}(a,t,x')_i\in\C\) for the \(i\)-th evaluated row-pairing sum (\(v_i(t,x')\) on the good summand in \eqref{eq:p8c2-test-formula}, indexed via \(e_a\)). A \emph{full grid} over \(R\) at \((1,C_F,b)\) is a grid \(G\) with compatible Weil system \(W\) over \(R\) admitting a complete choice satisfying both the evaluation-field condition and the square-field condition of Lemma~\ref{B8.20}.
\item At a member point \(p\) and constant \(C_F\), a grid \(G\) on \(p\) at \(C_F\) \emph{fits \(p\)} if: \(\iota(K_0)\) contains the unramified normalized values of \(p\)'s character frame (\(\mu_{0,j},\chi_\Psi,\delta,\delta'\) at weight \(0\) and \(\eta\) at weight \(1/2\)), the unramified weight-\(0\) values of \(\chi_1,\dots,\chi_m\), and a primitive \((N\mathfrak q)^{d_{\rm pl}(q)}\)-th root of unity for every \(q\in S_0\) (as in Section~1 of Chapter~\ref{ch:p8c1}); at every twisted finite place \(w\) of \(K\) (\(\mu_{\rm at}(w)\ne\mu_0\)), its twist index \(a_{\rm tw}=\mathrm{tw}(w)\) has residue prime \(\ell_{a_{\rm tw}}=\operatorname{char}(w)\), its underlying place lies under \(w\) or \(\bar w\), and \(\mathrm{tw}(w)\) agrees across paired twisted places whenever \(Q(w)\ne\varnothing\); for every such \(w\), every \(\mathfrak P\in Q(w)\) splits in \(E/F\), \(Q(w)\) is closed under CM conjugation, and \(Q(w)\) is disjoint from the excluded finite set \(S_0(a_{\rm tw})\) of places of \(E\) at \(a_{\rm tw}\) (which contains the places above \(\ell_{a_{\rm tw}},g_0,g_1\) and the selected prime-twist prime in \(\{p_{\rm de},p'_{\rm de}\}\setminus\{\ell_{a_{\rm tw}}\}\)); for every active twist index \(a_{\rm tw}\) there exists \(N_{\rm tw}(a_{\rm tw})\in\N\) such that for every twisted place \(w\) with \(\mathrm{tw}(w)=a_{\rm tw}\), integer \(m'\) coprime to \(\ell_{a_{\rm tw}}\), and abelian layer index \(\lambda\) at \(a_{\rm tw}\), the evaluated twisted toric element for \(\gamma_w^{m'}\) is nonzero with normalized \(\ell_{a_{\rm tw}}\)-adic valuation in \([-N_{\rm tw}(a_{\rm tw}),\,c_0(a_{\rm tw},\lambda)+\ell_{a_{\rm tw}}^{-2}]\) (clause~7 of Definition~\ref{B8.01k}, in the notation of Definition~\ref{B8.01ka}); at every twisted place \(w\) (with \(a_{\rm tw}=\mathrm{tw}(w)\)) and Frobenius element at a finite place \(v\notin Q(w)\) of \(E\) off \(\ell_{a_{\rm tw}}\), preimages in \(K\) of \(\psi_w(v)^{-1}\) match the local Artin character values; and at every active \(a_{\rm tw}\), \(S_0(a_{\rm tw})\) contains all places above the base conductor at \(w(v_0)\) and the abelian layer characters and toric elements at \(a_{\rm tw}\) match those of \(G\) across \(0\le j,a<m\) and \(1\le s\le N_*+1\). A \emph{fitted grid} at \((p,C_F)\) consists of a grid \(G\) fitting \(p\), a Weil system \(W\) on \(G\) over \(p\)'s real realization, and a complete choice on \((p,G,W)\) satisfying the evaluation-field condition and the square-field condition of Lemma~\ref{B8.20}. Put \(T_{\rm use}(p)=\{\mu_0\}\cup\{\mu(w):w\text{ a finite place of }K\}\). A \emph{fitted evaluation datum} \(d\) at \((p,C_F)\) packages \((G,W)\) with \(n\in\N\), matrices \(X'_v\in M_n(\C)\) (\(v\mid\infty\)), an \(\mathcal O_K\)-submodule \(M\subset K^n\), matrices \(V_t^0\in M_{n\times n_V(t)}(K)\) (\(t\in T\)), elements \(r_{ij,t}\in K\) (\(t\in T,\ 0\le i,j<m\)), a function \(d_\ell\) on finite places of \(K\), finite sets \(S_d,S_e\subset\N\), functions \(N_0(i,j,s,l)\in\R\) (\(0\le i,j<m,\ s\in\Z,\ l\in\N\)), bottom elements \(z_A(i,j),z_A'(i,j)\in K\), higher toric elements \(z_T(i,j,s),z_{T,0}(i,j,s),z_{T,1}(i,j,s)\in K\) with shifts \(g_T(i,j,s,w)\in\Z\), selector \(c(w)\in\{0,1\}\), and \(B_w'\in\R\), for which some complete choice \(a\) of ambient dimension \(n\) on \((p,G,W)\) satisfies:
\begin{enumerate}
\item \(a\) satisfies the separation, rationality, series, evaluation-field, and square-field conditions;
\item \((M,(X'_v)_{v\mid\infty})\) is the full-rank direct-sum Hermitian test \(\mathcal O_K\)-lattice in \(K^n\) (\(M\otimes_{\mathcal O_K}K=K^n\)) defined from the form lattice \(L_{\rm form}(G,x_c,e_a)\subset K^n\) and Hermitian form Gram matrix \(G_{\rm form}(G,e_a)\in\GL_n(K)\) assembled along \(e_a:\{1,\dots,n\}\xrightarrow{\sim}\mathcal I_{\rm test}(G)\times\mathcal I_{\rm coef}(G)\) from the class-representative lattices at \(x_c\) on \(X_0\) and the Weyl module of \(\lambda_{P,G}\) by taking the dual lattice \(M=L_{\rm form}(G,x_c,e_a)^\vee\) and frames satisfying \((X'_v{}^{-1})^\dagger X'_v{}^{-1}=\bigl(v(G_{\rm form}(G,e_a))^{-1}\bigr)^t\) at every \(v\mid\infty\);
\item for every \(t\in T_{\rm use}(p)\), \(x\in K^{n_V(t)}\), and \(1\le i\le n\), \(\iota((V_t^0 x)_i)=v_{\rm row}(a,t,e_{V,0}(p,t,x))_i\), where \(e_{V,0}(p,t,x)=\sum_{\alpha=1}^{d_{\rm amb}(t)}(v_{E,t}(\ell_0)x)_\alpha b_{t,\alpha}\in FR_t\) is the level-zero row-source image of \(x\) at \(\ell_0=\ell(P_0,\psi_0)\); and
\item for every \(0\le i,j<m\), there exist seesaw proportionality constants \(c_{ij}(t)\in\C\) (\(t\in T\), obtained via the seesaw comparison scalar of Definition~\ref{B8.16c}, Lemma~\ref{B8.16e}, and Theorem~\ref{B8.16}) with \(c_{ij}(\mu_0)\ne0\) such that at every \(t\in T_{\rm use}(p)\), on the \((i,j)\) grid square between \((i,j+1)\) and \((i+1,j)\) through \((i+1,j+1)\) and \((i,j)\), every \(t\)-admissible grid volume package, matched kernel pair \((\Theta,\Theta')\) with big and small theta kernels \(K_{\rm big}(W,\Theta,g,g_1;\cdot)\) and \(K_{\rm small}(W,\Theta',g,g_1)\), and convergent Hecke translation datum \((z,g,g_1)\) of formal degree \(\deg(z)\) and translation operator \(T_z\) satisfy
\[
 \deg(z)\int_{\mathcal F_{\rm big}}K_{\rm big}(W,\Theta,g,g_1;u)\,d\nu_{\rm big}(u)
 =c_{ij}(t)\int_{\mathcal F_{\rm small}}\bigl(T_z K_{\rm small}(W,\Theta',g,g_1)\bigr)(u)\,d\nu_{\rm small}(u)
\]
with the right-hand integral nonzero for at least one such choice, and \(\iota(r_{ij,t})c_{ij}(\mu_0)=c_{ij}(t)\) (so \(r_{ij,\mu_0}=1\)). Moreover, there exist a common discriminant twelfth-root period \(\delta_{\rm per}\in\C\) of \(\mathcal O_{K_d}\), an unkilled two-torsion point \(t_0\in\C\) for \(\Lambda(\mathcal O_{K_d},\delta_{\rm per})\), and for all \(0\le i,j<m\) standard bottom pairs \((P_{ij},P'_{ij})\) on \(G\) at \((\mu_0,i,j,\mathfrak m_E)\) over \((\mathcal O_{K_d},\delta_{\rm per},t_0,Q_{\rm lev})\) (Definition~\ref{B8.21ca}(3)) with \(\iota(z_A(i,j))=\mathrm{norm}(P_{ij})\) and \(\iota(z_A'(i,j))=\mathrm{norm}(P'_{ij})\). For every \(0\le i,j<m\), \(z_T(i,j,s)=1\) for all \(s\in\Z\), while for \(1\le s\le h_{ij}\) (\(h_{ij}=\lfloor(r_{i+1,j+1}+r_{i,j})/2\rfloor\)), \((z_{T,0}(i,j,s),z_{T,1}(i,j,s))\) and \(g_T(i,j,s,\cdot)\) satisfy the toric-ratio valuation relation of Lemma~\ref{B8.21a} at \(\mu_0\) (for \(\nu_{\mu_0,i,j}\) and prime-twist characters \(\delta_0=\delta\) at \(p_{\rm de}\) and \(\delta_1=\delta'\) at \(p'_{\rm de}\), agreeing with \(z_T,z_{T,c},g_T(w)\) of Lemma~\ref{B8.21a}), and for rational ideal-class partial zeta values \(z(C)=\zeta_F(1-2s,C)\in\Q\) (\(C\in\mathrm{Cl}(\mathcal O_F)\)) of the totally real field \(F\) at \(1-2s\) (as in Definition~\ref{B1.14k} and Lemma~\ref{B8.21a}, with weight \(\kappa_{i,j,s}(\tau)=\max(0,e^{\rm gr}_{\tau,j+1}+c_{\tau,m-i}-s)\) and \(M_{i,j,s}=2sf+2\sum_{\tau\in\mathcal T}\kappa_{i,j,s}(\tau)\)), \(N_0(i,j,s,l)\) is given at every prime \(l\) by
\[
 N_0(i,j,s,l)=\Bigl\lceil\frac{\beta_{\rm vp}(d,s,\kappa_{i,j,s},\mathfrak f_F,z,l)}{12}\Bigr\rceil_+
 =\bigl\lceil c_{E,l}(z)+n_l(d,s,\kappa_{i,j,s},\mathfrak f_F)+M_{i,j,s}\,a_0(\mathfrak f_F,l)\bigr\rceil_+\in\N
\]
(from the partial-zeta, ramification/conductor, and \(\mathfrak f_F\)-valuation exponents of Definition~\ref{B1.14k}, with \(N_0(i,j,s,l)=0\) when \(l\in\N\) is not prime), while for \(s\notin[1,h_{ij}]\) one has \(z_{T,0}(i,j,s)=z_{T,1}(i,j,s)=1\), \(g_T(i,j,s,\cdot)=0\), and \(N_0(i,j,s,\cdot)=0\). Globally, at each finite place \(w\) of \(K\), \(c(w)=1\) if \(p_{\rm de}\in\mathfrak p_w\) and \(c(w)=0\) otherwise; \(S_d=\{l\in\N:l\text{ is prime and }\exists\,i,j,s,\ N_0(i,j,s,l)\ne0\}\) and \(S_e=\{l\in\N:l\text{ is prime and }\exists\,w,\ \mu(w)\ne\mu_0\land l\in\mathfrak p_w\}\); \(d_\ell\) is a grid integrality allowance on \(G\) at constant \(C_F\) (Definition~\ref{B8.21ca}(4)); and \(B_w'=B_w\) is the gap-\(1\) bottom weight bound of Definition~\ref{B8.21ca}(4) (henceforth denoted \(B_w\)).
\end{enumerate}
\end{enumerate}
\end{definition}

\begin{definition}[Grid modulus allowances and evaluation budgets]\label{B8.21e}
For a grid \(G\) over \(F\), recall from Definition~\ref{B8.21c} the finite set \(S_0=\{v_0\}\cup S_1\) of grid places of \(F\), the prime-power absolute norms \(N\mathfrak q>1\) (\(q\in S_0\), also written \(q\) inside logarithms and \(w\)-adic valuations), the grid conductor exponents \(c_{v_0}=1\) and \(c_q=2\) (\(q\in S_0\setminus\{v_0\}\)), \(Q_{\rm lev}=\prod_{q\in S_0}(N\mathfrak q)^{c_q}\), and \(c_S=\log Q_{\rm lev}\). Put \(R_*=k+m-1=r\), \(n_{\rm rk}=k+m=R_*+1\), \(d_s=k+m+5=R_*+6\), and \(N_*=k+m(t+d_s+1)=k+m(t+k+m+6)\). For an estimate constant \(C_1\in\R\) and scale weight \(b=\exp(L)>0\), define the prime-power modulus budget
\[
 P_{\rm pf}=P_{\rm pf}(C_1,b)=C_1f^2(k+m)^6\log\bigl(b(k+m)g_0g_1p_{\rm de}p'_{\rm de}|D_E|\bigr),
\]
and define the \emph{grid modulus allowances} \(e_q=e_q(C_1,b)\in\N\) (\(q\in S_0\)) by
\[
 e_{v_0}(C_1,b)=\bigl\lfloor P_{\rm pf}(C_1,b)/\log N\mathfrak q_{v_0}\bigr\rfloor_+ +1,\qquad
 e_q(C_1,b)=2\quad(q\in S_0\setminus\{v_0\}),
\]
where \(\lfloor y\rfloor_+=\max(0,\lfloor y\rfloor)\in\N\); thus \(c_q\le e_q(C_1,b)\) for all \(q\in S_0\), and put \(\Sigma_0=\Sigma_0(C_1,b)=\sum_{q\in S_0}e_q(C_1,b)\log N\mathfrak q\).

Enumerate \(\mathcal T\simeq\{0,\dots,f-1\}\) and define the integer cloud coordinates \(\mathrm{cl}(\tau,i)=\lfloor v_{\tau,i-2}\rfloor\in\Z\) for \(0\le\tau<f\) and \(2\le i\le k\) (from the Harish--Chandra tuple \(v_\tau\in\Q^{k-1}\) of Definition~\ref{B8.21b}) and \(\mathrm{cl}(\tau,i)=0\) otherwise, with integer cloud spread
\[
 L_c=\max_{0\le\tau<f}\max\bigl(0,\;\mathrm{cl}(\tau,2)-\mathrm{cl}(\tau,k)\bigr)\in\N.
\]
Writing \(\lceil y\rceil_+=\max(0,\lceil y\rceil)\in\N\), define the initial top-node rank \(n_p\in\N_{>0}\) (equal to \(n_m=n'_0=\operatorname{rank}(X_0)\) of Definition~\ref{B8.21c}), the integer gap bound \(\gamma_{\rm gap}(C_1)\in\N\), and the dominant-weight spread factor \(R_p=R_p(C_1)\ge1\) by
\[
 \begin{gathered}
 n_p=(R_*+1)+m(R_*+5)=k+m+m(k+m+4),\\
 \gamma_{\rm gap}(C_1)=\max\Bigl(\bigl\lceil C_1N_*^3+C_1n_{\rm rk}^7+C_1n_{\rm rk}^3\sqrt{L_c}\bigr\rceil_+,\;
 2N_*^3+2d_s+3,\;\bigl\lceil 48n_{\rm rk}^2\sqrt m\bigr\rceil_+\Bigr),\\
 R_p(C_1)=(m+1)\bigl(\gamma_{\rm gap}(C_1)+1\bigr)+L_c+n_p .
 \end{gathered}
\]
Recall the place weights \(\ell_w=\log N\mathfrak p_w/[K:\Q]\) and toric budget \(B_0=C_1fb(\log b+\log x)\) of Lemma~\ref{B8.21a} and the gap-\(1\) bottom weight bound \(B_w\ge1\) of Definition~\ref{B8.21ca}. Define
\[
 E_M=4fn_pR_p\log(n_pR_p),\quad
 H_{I,0}=C_1fmN_*bL+C_1m^2N_*^2\Sigma_0,\quad
 B_F=C_1fbm^2(N_*+N_*^2)(L+\Sigma_0).
\]
\end{definition}

\begin{lemma}[Combination of the eight finite losses]\label{B8.21fa}
Let \(K\) be the coefficient field with embedding \(\iota:K\hookrightarrow\C\), let \(m\ge1\), and let \(L_{\rm form}\subset K^n\) be the full form lattice of Definition~\ref{B8.21d}(3), with dual \(M=L_{\rm form}^{\vee}\). For a finite place \(w\), say that a complex scalar \(z\) has loss at most \(a\) if \(z=\iota(z_K)\) for some \(z_K\in K\) such that \(z_K=0\) or \(v_w(z_K)\ge-a\). For a coefficient lattice \(\mathcal L\) in a \(K\)-vector space, say that a complex coefficient vector has loss at most \(a\) in \(\mathcal L\) if it is the image of a vector \(u\) for which some \(s\in K^{\times}\) satisfies \(v_w(s)\le a\) and \(su\in\mathcal L\).

Choose an \(\mathcal O_K\)-generating set of \(L_{\rm form}\). Suppose the pairings of \(y\in K^n\) with its generators, after applying \(\iota\), have finite expansions of the form
\[
 s_{\rm nor}\sum_{\beta}c_{\beta}
       \sum_{\gamma}\gamma_{\beta,\gamma}
       \sum_{\eta}d_{\beta,\eta}
       \left\langle
       T_{m-1}\cdots T_0\Psi_0,
       R_{m-1}\cdots R_0F_0
       \right\rangle.
\]
Here \(\beta\) runs over the supported finite local terms, \(\gamma\) over their finite class sets, and \(\eta\) over a finite coset expansion. The factors \(c_\beta\), \(\gamma_{\beta,\gamma}\) and \(d_{\beta,\eta}\) are respectively the product of local coefficients, the inverse stabilizer factor, and the product of local datum values. The operators \(T_j\) and \(R_j\), their coefficient lattices and their initial vectors depend on the term indices. They are respectively the successive top and row kernel maps in that expansion. The final pairing is the contraction between the two terminal coefficient spaces.

Let \(\mathcal B\ge0\), and suppose there are eight nonnegative functions
\(\delta_S,\delta_C,\delta_D,\delta_G,\delta_T,\delta_R,\delta_F,\delta_P\), each supported on finitely many finite places, with the following properties. Their weighted sums satisfy
\[
 \begin{array}{c|c}
 \text{loss function}&\displaystyle\sum_w\ell_w\delta(w)\\ \hline
 \delta_S,\delta_C,\delta_D,\delta_G&m\mathcal B\\
 \delta_T,\delta_R,\delta_F,\delta_P&\mathcal B
 \end{array}
\]
where the right column is an upper bound for each function in its row. At every place under consideration, uniformly in all the term indices, \(s_{\rm nor},c_\beta,d_{\beta,\eta},\gamma_{\beta,\gamma}\) have losses at most \(\delta_S(w),\delta_C(w),\delta_D(w),\delta_G(w)\), respectively. The initial top vector \(\Psi_0\) has loss zero, and \(F_0\) has loss at most \(\delta_F(w)\). Each \(T_j\) sends a vector of loss \(a\) to one of loss at most \(a+\delta_T(w)\) in the next coefficient lattice, and each \(R_j\) has the analogous property with \(\delta_R(w)\). If terminal vectors have losses \(a,a'\), their contraction has loss at most \(a+a'+\delta_P(w)\).

Put
\[
 \delta(w)=\delta_S(w)+\delta_C(w)+\delta_D(w)+\delta_G(w)
           +m\delta_T(w)+m\delta_R(w)+\delta_F(w)+\delta_P(w).
\]
Then \(\delta\) is nonnegative and finitely supported,
\[
 \sum_w\ell_w\delta(w)\le(6m+2)\mathcal B\le8m\mathcal B,
\]
and for every such \(w\) there is \(s_w\in K^{\times}\) with
\[
 v_w(s_w)\le\delta(w),\qquad s_wy\in M.
\]
\end{lemma}
\begin{proof}[Outline]
Track losses over \(m\) top and row steps, add the contraction and scalar losses, and apply the nonarchimedean triangle inequality to pairings with \(L_{\rm form}\). Local duality with a uniformizer gives \(s_w\), and summing the eight bounds gives the estimate. The complete proof is in \S\ref{prf:B8.21fa}.
\end{proof}

\begin{lemma}[The finite layer-sum denominator estimate]\label{B8.21f}
Assume \(\mathscr P\). There is \(C_{F0}\in\R\) such that for every \(C_F\ge C_{F0}\) there is \(C_0\in\R\), and for every \(\varepsilon\in\R\) and \(C_1\ge C_0\) there is \(x_0\in\R\) with the following property. Fix an arithmetic datum \(D\) and size datum \(L_{\rm sz}\) at scale \(s\) of exponent \(\varepsilon\) with \(x(s)\ge x_0\), a real realization \(R\), and a level-\(2\) face \(F_{\rm fc}\) such that \((L_{\rm sz},R,F_{\rm fc})\) satisfies the Hida row-setting conditions of Definition~\ref{B8.21d}(1), together with a grid \(G\) at \((1,C_F,b)\), a Weil system \(W\) on \(G\), and a complete choice \(a\) of ambient dimension \(n\) on \((G,W)\) over \(R\) satisfying the evaluation-field condition of Definition~\ref{B8.21d}(2). Put \(\sigma(G)=\sum_{q\in S_0}d_{\rm pl}(q)\log N\mathfrak q\), where \(d_{\rm pl}(q)\) is the prescribed grid depth of Definition~\ref{B8.21c}. Then there exists a nonnegative real-valued denominator-loss function \(\delta_{\rm sum}\) on the finite places of \(K\), vanishing outside a finite set \(T_{\rm sum}\) of exceptional finite places of \(K\) and satisfying
\[
 2\sum_{w\in T_{\rm sum}}\ell_w\delta_{\rm sum}(w)\le C_1fbm^2N_*^2\bigl(L+\sigma(G)\bigr),
\]
such that for every Hermitian test \(\mathcal O_K\)-lattice \((M,(X'_v)_{v\mid\infty})\) in \(K^n\) built from \(a\) as in Definition~\ref{B8.21d}(3), every \(t\in T\) and finite place \(w\) of \(K\) such that every \(\mathfrak P\in Q_t\) either lies over a place in \(S_0\) or has rational residue prime different from that of \(w\), every source element \(x\in M_{t,{\rm lv}}(F_{\rm fc})\), and every vector \(y=(y_1,\dots,y_n)\in K^n\) satisfying \(\iota(y_i)=v_{\rm row}(a,t,\mathrm{form}_{F_{\rm fc}}(t,x))_i\) for all \(1\le i\le n\), there exists a nonzero scalar \(s_0\in K^\times\) with
\[
 v_w(s_0)\le\delta_{\rm sum}(w)\qquad\text{and}\qquad s_0y\in M .
\]
In particular, at every finite place \(w\notin T_{\rm sum}\) one has \(\delta_{\rm sum}(w)=0\), so \(v_w(s_0)\le0\) and \(y\in M_w\).
\end{lemma}

\begin{proof}[Outline]
The local expansion theorem and Definition~\ref{B8.21c} put each class pairing in the form of Lemma~\ref{B8.21fa}. Applying the ten local factor bounds at \(C_1/16\) under \(\mathscr P\) and linearity of \(B_1\) then yields \(\delta_{\rm sum}\) and \(s_0\). The complete proof is in \S\ref{prf:B8.21f}.
\end{proof}

\begin{lemma}[The finite sum estimates]\label{B8.21g}
Assume \(\mathscr P\). There is \(C_{F0}\in\R\) such that for each \(C_F\ge C_{F0}\), in the closed sense of Definition~\ref{B8.39} on the family class of Definition~\ref{B8.24} under the division bound of Lemma~\ref{B8.40a}, there is \(C_0\in\R\) with the following property. For every \(C_1\ge C_0\), eventually at every member point \(p\in\mathcal P\) and every fitted evaluation datum \(d\) at \((p,C_F)\) with unnormalized evaluation matrices \(V_t^0\), there exist a finite set \(T_F\) of finite places of \(K\) and integers \(d_F(w)\) vanishing outside \(T_F\) such that
\[
 V_{\mu(w)}^0(M_{\mu(w),w})\subseteq\mathfrak p_w^{-d_F(w)}M_w\quad(\text{all finite }w),\qquad
 2\sum_{w\in T_F}\ell_w d_F(w)\le B_F .
\]
Moreover, for any family of roots of unity \(\zeta_t\in K\) (\(t\in T\)), replacing \(V_t^0\) by \(\zeta_t^{-1}V_t^0\) preserves both properties with the same \(T_F\) and \(d_F\).
\end{lemma}

\begin{proof}[Outline]
Apply Lemma~\ref{B8.21f} at \(t=\mu(w)\) via place-avoidance and Definition~\ref{B8.21d}(1) with \(d_F=\lfloor\delta_{\rm sum}\rfloor\). Localization gives the inclusion, monotonicity of \(\Sigma_0\) gives \(B_F\), and integrality of \(\zeta_t^{-1}\) preserves both. The complete proof is in \S\ref{prf:B8.21g}.
\end{proof}

\begin{lemma}[Height through a direct quotient metric]\label{B8.21ha}
Let \(K\) be a totally complex number field, put \(n_\infty=[K:\Q]/2\), and write \(A_v\) for the image of a matrix \(A\) under the embedding at an infinite place \(v\). Let \(H\in M_r(K)\) have positive-definite Hermitian images \(H_v\), and let \(P\in M_{q\times r}(K)\) have rank \(q\). Suppose that the invertible source frames \(X_v\) and the direct quotient frames \(X_{Q,v}\) satisfy
\[
 (X_v^{-1})^\dagger X_v^{-1}=H_v,
 \qquad X_{Q,v}=(P_vH_v^{-1}P_v^\dagger)^{1/2}.
\]
For invertible target frames \(X'_v\) and a matrix \(B\in M_{n\times q}(K)\), define
\[
 h_{X_Q,X'}(B)=\frac1{n_\infty}\sum_{v\mid\infty}
       \log\|(X'_v)^{-1}B_vX_{Q,v}\|^2,
\]
and define \(h_{X,X'}(BP)\) by replacing \(B_vX_{Q,v}\) with \(B_vP_vX_v\). Here the norm is the Euclidean operator norm, \(\dagger\) denotes conjugate transpose, and \(\log0=0\) in the height definition. Then
\[
 h_{X_Q,X'}(B)=h_{X,X'}(BP).
\]
\end{lemma}
\begin{proof}
The source-frame identity gives \(H_v^{-1}=X_vX_v^\dagger\). Set
\[
 U_v=X_{Q,v}^{-1}P_vX_v.
\]
Then \(U_vU_v^\dagger=1\). Hence \(U_v\) is a coisometry and \(U_v^\dagger\) is an isometry. For a matrix \(A\) with \(q\) columns, the inequalities
\[
 \|AU_v\|\le\|A\|,\qquad
 \|A\|=\|AU_vU_v^\dagger\|\le\|AU_v\|
\]
give equality. Apply this to \(A=(X'_v)^{-1}B_vX_{Q,v}\), square the norms, take logarithms, and average over the infinite places.
\end{proof}

\begin{lemma}[The base evaluation estimate]\label{B8.21h}
Assume \(\mathscr P\). There is \(C_{F0}\in\R\) such that for each \(C_F\ge C_{F0}\), in the closed sense of Definition~\ref{B8.39} on the family class of Definition~\ref{B8.24} under the division bound of Lemma~\ref{B8.40a}, there is \(C_0\in\R\) with the following property. For every \(C_1\ge C_0\), eventually at every member point \(p\in\mathcal P\) and every fitted evaluation datum \(d\) at \((p,C_F)\) with base matrix \(V_{\mu_0}^0\), every matrix \(V_b\in M_{n\times d_Q(\mu_0)}(K)\) satisfying \(V_bP_{\mu_0}=V_{\mu_0}^0\) obeys
\[
 h(V_b)\le H_{I,0},
\]
where \(n_\infty(K)=[K:\Q]/2\) and \(h(V_b)=\frac1{n_\infty(K)}\sum_{v\mid\infty}\log\|V_b\|_v^2\) (with \(\log 0=0\) in this height definition) is the operator-norm height from \(K^{d_Q(\mu_0)}\) equipped with the direct quotient Hermitian frames \((X_Q(\mu_0)_v)_{v\mid\infty}\) of Definition~\ref{B8.21b} to \((K^n,(X'_v)_{v\mid\infty})\).
\end{lemma}

\begin{proof}[Outline]
Above the Hida and archimedean cutoffs, the archimedean layer-sum and grid-budget bounds give \(h_{X,X'}(V)\le H_{I,0}\). Taking \(X_v=H_V(\mu_0)_v^{-1/2}\) and \(V=V_{\mu_0}^0\), Lemma~\ref{B8.21ha} transfers this bound through \(P_{\mu_0}\) to \(h(V_b)\). The complete proof is in \S\ref{prf:B8.21h}.
\end{proof}

\begin{lemma}[The frame estimate for the evaluation lattice]\label{B8.21i}
For every \(C_F\ge0\), in the closed sense of Definition~\ref{B8.39} on the family class of Definition~\ref{B8.24} under the division bound of Lemma~\ref{B8.40a}, for every \(C_1\ge1\), eventually at every member point \(p\in\mathcal P\) and every fitted evaluation datum \(d\) at \((p,C_F)\), the \(\mathcal O_K\)-submodule \(M\subset K^n\) is finitely generated and nonzero, \(X'_v\) is invertible (\(\det X'_v\ne0\)) at every infinite place \(v\) of \(K\), and with respect to the standard basis of \(K^n\) and the frames \((X'_v)_{v\mid\infty}\),
\[
 \mu_{\max}(M)\le E_M=4fn_pR_p\log(n_pR_p).
\]
\end{lemma}

\begin{proof}[Outline]
Identifies \((M,(X'_v)_{v\mid\infty})\) as the direct-sum test lattice of rank \(n_p\) via \ref{B8.21d}(3). Bounding the dominant weight spread by \(R_p(C_1)\) via \ref{B8.21e} gives nondegeneracy of \(X'_v\), finite generation of \(M\), and \(\mu_{\max}(M)\le E_M\). The complete proof is in \S\ref{prf:B8.21i}.
\end{proof}

\begin{lemma}[The valuation budget of a square tower]\label{B8.21ja}
Let \(K\) be a number field with involution, let \(w\) be a finite place of \(K\), and write \(v=v_w\), with the convention \(v(0)=0\). Let \(e,p,q,e_0,h_0\) be nonnegative integers such that
\[
 e_0+p=e+q,\qquad e_0\le h_0,\qquad h=h_0+p+q.
\]
Let \(R,\zeta,\gamma,A,A'\in K\), let \(t_0,\ldots,t_{p+q-1}\in K\), and let \(\rho_s\in K\) for \(1-e_0\le s\le h\). Suppose that \(A,\gamma,\zeta,t_k,\rho_s\) are nonzero, that \(v(\zeta)=0\), and that
\[
 R\prod_{k=0}^{q-1}\rho_{e+k+1}
 =\zeta\left(\prod_{k=0}^{p+q-1}t_k\right)\gamma
       \prod_{k=0}^{p-1}\rho_{e+q-k}.
\]
Let \(a_s,n_s\ge0\) for \(1\le s\le h\), let \(\eta\ge0\), and let \(d,B\in\R\). Assume
\[
 \begin{gathered}
 v(\rho_s)\le a_s+\eta,\qquad
 -v(\rho_s)\le a_s+n_s\qquad(1\le s\le h),\\
 v(\rho_s)=v(\rho_{1-s})\qquad(1-e_0\le s\le0),\qquad
 |v(t_k)|\le B,\\
 v(\gamma)\le
 v\left(A\prod_{s=1-e_0}^{h_0}\rho_s\right)+v(A')+6d.
 \end{gathered}
\]
Then \(R\ne0\) and
\[
 v(R)\le |v(A)|+|v(A')|
       +4\sum_{s=1}^h a_s+5\sum_{s=1}^h n_s
       +2h\eta+6d+(p+q)B.
\]
\end{lemma}
\begin{proof}[Outline]
Check that all tower indices lie in \([1,h]\), deduce \(R\ne0\), take valuations, and substitute the \(v(\gamma)\) bound. Reflection and the \(v(\rho_s)\) bounds give \(4A_1+5N_1+2h\eta\), and replacing \(v(A),v(A')\) by their absolute values finishes. The complete proof is in \S\ref{prf:B8.21ja}.
\end{proof}

\begin{lemma}[The square valuation estimate]\label{B8.21j}
Assume \(\mathscr P,\mathscr J\), the Hecke--Tate functional equation, and Theorem~\ref{B8.11}. There is \(C_{F0}\in\R\) such that for each \(C_F\ge C_{F0}\), in the closed sense of Definition~\ref{B8.39} on the family class of Definition~\ref{B8.24} under the division bound of Lemma~\ref{B8.40a}, there is \(C_0\in\R\) with the following property. For every \(C_1\ge C_0\), eventually at every member point \(p\in\mathcal P\) and every fitted evaluation datum \(d\) at \((p,C_F)\), every element of \(S_d\cup S_e\) is prime, \(c_q\le e_q\) for all \(q\in S_0\), and for all \(0\le i,j<m\) and every finite place \(w\) of \(K\), putting \(y_{ij}=m-1-j\) (so \(0\le y_{ij}\) and \(y_{ij}+1\le m\)), one has \(h_{ij}\le N_*\), \(r_{ij,\mu_0}=1\), \(r_{ij,\mu(w)}\ne0\),
\[
 \begin{aligned}
 v_w(r_{ij,\mu(w)})\le{}&
 |v_w(z_A(i,j))|+|v_w(z_A'(i,j))|
 +4\sum_{s=1}^{h_{ij}}\bigl(|v_w(z_T(i,j,s))+g_T(i,j,s,w)|+|v_w(z_{T,c(w)}(i,j,s))+g_T(i,j,s,w)|\bigr)\\
 &+5\sum_{s=1}^{h_{ij}}\sum_{l\in S_d}N_0(i,j,s,l)v_w(l)
 +2h_{ij}\sum_{l\in S_e}\frac{v_w(l)}{l^2}
 +6d_\ell(w)+2y_{ij}\sum_{q\in S_0}c_qv_w(q),
 \end{aligned}
\]
and, for every integer \(1\le s\le h_{ij}\), \(N_0(i,j,s,l)\ge0\) for all \(l\in\N\) (vanishing outside \(S_d\)) with \(\sum_{l\in S_d}N_0(i,j,s,l)\log l\le B_0\).
\end{lemma}

\begin{proof}[Outline]
For \(\mu(w)\ne\mu_0\), form \(\rho_s\) from Hecke-function ratios via reflection and toric bounds, compare bottom constants, and descend the seesaw tower. Lemma~\ref{B8.21ja} yields the valuation bound, and Definition~\ref{B8.21e} gives \(c_q\le e_q\). The complete proof is in \S\ref{prf:B8.21j}.
\end{proof}

\begin{lemma}[An explicit prime-denominator allowance]\label{B8.21ka}
Let \(f,N,Q\ge1\) and \(d\ge0\) be integers, let \(B,L\ge1\), and let \(S\subset\N\) be finite. Use \(\log0=0\) throughout this lemma. Suppose that \(a_p\ge0\) for every \(p\in\N\), that \(a_p=0\) outside \(S\), and that
\[
 \sum_{p\in S}a_p\log p\le L,\qquad \log d\le2L.
\]
Put
\[
 M=8^{4fN}(8fN)!,\qquad
 S'=\bigl(\{p:p\mid M\}\cup S\cup\{d\}\cup\{p:p\mid Q\}\bigr)
       \cap\{p:p\text{ is prime}\}.
\]
For prime \(p\), write \(v_p\) for the exponent of \(p\) in a positive integer and set
\[
 e_p=v_p(M)+4fNB\left(a_p+\frac{\mathbf1_{p=d}}4+v_p(Q)\right)
       +4N^2v_p(Q).
\]
Set \(e_p=0\) when \(p\) is not prime. Then \(e_p\ge0\), \(e_p=0\) outside \(S'\), and
\[
 \begin{split}
 \sum_{p\in S'}e_p\log p
 &\le36fN\log(2fN)+6fNBL+4fNB\log Q+4N^2\log Q\\
 &\le I(18;f,N,B,L,Q)\\
 &=36fN\log(2fN)+36fNBL+72fN^2B\log Q.
 \end{split}
\]
\end{lemma}
\begin{proof}[Outline]
Bound \(\log M\le 18r\log r\) for \(r=2fN\) via \((4r)!\le(4r)^{4r}\), rewrite \(\log M\) and \(\log Q\) by prime factorization, and insert the bounds on \(a_p\) and \(\log d\). Elementary bounds using \(fB,N\ge1\) then give \(I(18;f,N,B,L,Q)\). The complete proof is in \S\ref{prf:B8.21ka}.
\end{proof}

\begin{lemma}[The coordinate denominator bound]\label{B8.21k}
There is an absolute constant \(C_{\rm int}>0\) (one may take \(C_{\rm int}=18\)) with the following property. Let \(D\) be an arithmetic datum with \(f=[F:\Q]\), let \(Q,N\ge1\) be integers such that \(Q\) is an admissible level modulus for \(D\) in the sense of Definition~\ref{B8.21ca}(1), and let \(B,L\ge1\) be real numbers with \(\log\operatorname{rd}(E)\le L\). Then there exist a finite set \(S'\) of rational primes and a real-valued exponent function \(\delta_{\rm vp}:\N\to\R_{\ge0}\) vanishing outside \(S'\) and satisfying \(\sum_{p\in S'}\delta_{\rm vp}(p)\log p\le I(C_{\rm int};f,N,B,L,Q)\) such that:
\begin{enumerate}
\item \(\delta_{\rm vp}\) satisfies the \(p\)-adic coordinate integrality bound of Definition~\ref{B8.21ca}(4) for \((D,N,B,Q)\): for every discriminant twelfth-root period \(\delta\in\C\) of \(\mathcal O_{K_d}\) with CM lattice \(\Lambda=\Lambda(\mathcal O_{K_d},\delta)\), unkilled two-torsion point \(t_0\in\C\) for \(\Lambda\), rational prime \(p\), level \(Q'\in\N\) that is \(p\)-compatible with \(Q\) (\(v_p(Q')\le v_p(Q)\)), and \((N,B)\)-small definite CM theta pair \(P\) at \((\mathcal O_{K_d},\delta,t_0,Q')\), every class in \(\mathrm{Cls}_D(d_P,J_P)\) has a representative \(h\) integral over \(p\), every \(h\in G_{\mathbb A}(d_P)\) satisfies \(|\Gamma_D(J_P,h)|\ne0\) and \(v_p(|\Gamma_D(J_P,h)|)\le\delta_{\rm vp}(p)\), and for every \(h\) integral over \(p\) and \(T\in\mathrm{Tup}_D(a_P,d_P,n_P)\) one has \(c_P(h,T)\in\overline{\Q}\) and \(|\lambda(c_P(h,T))|_p\le p^{\delta_{\rm vp}(p)}\) for every ring homomorphism \(\lambda:\overline{\Q}\to\C_p\); and
\item \((S',\delta_{\rm vp})\) satisfies the norm valuation bound of Definition~\ref{B8.21ca}(4) for \((D,N,B,Q)\) with budget \(I(C_{\rm int};f,N,B,L,Q)\): for every discriminant twelfth-root period \(\delta\in\C\) of \(\mathcal O_{K_d}\) with \(\Lambda=\Lambda(\mathcal O_{K_d},\delta)\), unkilled two-torsion point \(t_0\in\C\) for \(\Lambda\), totally complex number field \(K\) with \(\sigma:K\hookrightarrow\C\), index set \(\mathcal I\), families of \((N,B)\)-small definite CM theta pairs \(P_\alpha,P'_\alpha\) (\(\alpha\in\mathcal I\)) at \((\mathcal O_{K_d},\delta,t_0,Q)\) each possessing at least one nonzero coordinate, and elements \(z_{A,\alpha},z'_{A,\alpha}\in K\) with \(\sigma(z_{A,\alpha})=\mathrm{norm}(P_\alpha)\) and \(\sigma(z'_{A,\alpha})=\mathrm{norm}(P'_\alpha)\), setting \(d_\ell(w)=\sum_{p\in S'}\delta_{\rm vp}(p)v_w(p)\ge0\) gives \(z_{A,\alpha}\ne0\), \(z'_{A,\alpha}\ne0\), \(v_w(z_{A,\alpha})\ge-3d_\ell(w)\), and \(v_w(z'_{A,\alpha})\ge-3d_\ell(w)\) at every finite place \(w\) of \(K\), and \(\sum_{w\in U}\ell_w d_\ell(w)\le I(C_{\rm int};f,N,B,L,Q)\) on every finite set \(U\) of finite places of \(K\).
\end{enumerate}
In particular, for any grid \(G\) at \((c_{\rm gap},c_{\rm Dep},b)\) whose level \(Q_{\rm lev}\) is admissible for \(D\) with \(1\le\log b\), \(\log\operatorname{rd}(E)\le\log b\), and \(C\ge C_{\rm int}\), there exists a grid integrality allowance \(d_\ell\) on \(G\) at constant \(C\) in the sense of Definition~\ref{B8.21ca}(4).
\end{lemma}

\begin{proof}[Outline]
Apply Lemma~\ref{B8.21ka} to the genus-one CM derivation-denominator exponents to set \(\delta_{\rm vp}\) with \(C_{\rm int}=18\). The CM jet coordinate bound gives (1), expanding \(\mathrm{norm}(P)\) by the ultrametric inequality gives (2), and specializing finishes. The complete proof is in \S\ref{prf:B8.21k}.
\end{proof}

\begin{lemma}[The lattice norm bound]\label{B8.21l}
Assume \(\mathscr P\), the Hecke--Tate functional equation, and Theorem~\ref{B8.11}. There is \(C_{F0}\in\R\) such that for each \(C_F\ge C_{F0}\), in the closed sense of Definition~\ref{B8.39} on the family class of Definition~\ref{B8.24} under the division bound of Lemma~\ref{B8.40a}, there is \(C_0\in\R\) with the following property. For every \(C_1\ge C_0\), eventually at every member point \(p\in\mathcal P\) and every fitted evaluation datum \(d\) at \((p,C_F)\), one has \(d_\ell(w)\ge0\) at every finite place \(w\) of \(K\), and for every finite set \(U\) of finite places of \(K\), all \(0\le i,j<m\), and each \(z=z_A(i,j),z_A'(i,j)\),
\[
 \sum_{w\in U}\ell_w d_\ell(w)\le D_0,\qquad
 z\ne0,\qquad v_w(z)\ge-3d_\ell(w)\quad(\text{all finite }w),
\]
where \(D_0=2C_1fN_*\bigl(\log(2fN_*)+B_wL+2N_*B_wc_S\bigr)\).
\end{lemma}

\begin{proof}[Outline]
Shows standard bottom pairs \((P_{ij},P'_{ij})\) are \((N_*,B_w)\)-small with a nonzero coordinate via \(\mathscr P\), the Hecke-Tate functional equation, and \ref{B8.11}. Applies the norm valuation bound of \(d_\ell\) in \ref{B8.21ca}(4) and \ref{B8.21d}(3). The complete proof is in \S\ref{prf:B8.21l}.
\end{proof}

\begin{lemma}[The two Galois means]\label{B8.21ma}
There is an absolute constant \(C_A>0\) such that every complex lattice of discriminant one has area at most \(C_A\). Fix such a constant and put
\[
 C_{\rm gal}=18+\frac32\bigl(\log^+(2\pi^2)+\log^+C_A+2\bigr).
\]
Let \(D\) be an arithmetic datum, let \((\mathfrak c,\delta)\) be a curve datum with lattice \(\Lambda\), and let \(t_0\) be an unkilled two-torsion point. Suppose \(L\ge1\), \(\log\operatorname{rd}(E)\le L\), and
\[
 \left|\Re\frac{L'(1,\chi_d)}{L(1,\chi_d)}\right|\le L.
\]
Let \(M\subset\C\) be any number field containing \(g_2(\Lambda),g_3(\Lambda),\wp_\Lambda(t_0)\). For each \(\rho:M\hookrightarrow\C\), let \(\Lambda_\rho\) be the conjugate lattice and choose an unkilled two-torsion point \(t_\rho\) on it. Write \(A_\rho=\operatorname{area}(\Lambda_\rho)\) and let \(\nu_\rho=\nu(\Lambda_\rho,t_\rho)\) be the Gaussian weighted mean square of Definition~\ref{B8.21ca}(1). Then
\[
 \frac1{[M:\Q]}\sum_\rho\log^+(A_\rho^{-1})
       \le(\log^+C_A+2)L\le C_{\rm gal}L,
 \qquad
 \frac1{[M:\Q]}\sum_\rho\log^+\nu_\rho\le C_{\rm gal}L.
\]
\end{lemma}
\begin{proof}[Outline]
Bound \(y^6|\Delta(\tau)|\) on a fundamental domain for \(C_A\), conjugate the curve datum, and average the period identity for \(\log^+(A_\rho^{-1})\). The genus-one odd-theta estimate and theta normalization bound \(\log^+\nu_\rho\), and averaging gives \(C_{\rm gal}L\). The complete proof is in \S\ref{prf:B8.21ma}.
\end{proof}

\begin{lemma}[The archimedean norm estimate]\label{B8.21m}
\begin{enumerate}
\item There is an absolute constant \(C_{\rm arch}>0\) with the following property. Let \(D\) be an arithmetic datum with \(f=[F:\Q]\), let \((\mathfrak c,\delta)\) be a curve datum over \(D\) with CM lattice \(\Lambda=\Lambda(\mathfrak c,\delta)\), let \(t_0\in\C\) be an unkilled two-torsion point for \(\Lambda\) (Definition~\ref{B8.21ca}(1)), and let \(Q,N\ge1\) be integers and \(B,L\ge1\) real numbers satisfying \(\log\operatorname{rd}(E)\le L\) and \(\bigl|\Re(L'(1,\chi_d)/L(1,\chi_d))\bigr|\le L\). For any totally complex number field \(K\) with \(\sigma:K\hookrightarrow\C\), any families of \((N,B)\)-small definite CM theta pairs \(P_\alpha,P'_\alpha\) (\(\alpha\in\mathcal I\)) at \((\mathfrak c,\delta,t_0,Q)\), and any \(z_{A,\alpha},z'_{A,\alpha}\in K\) satisfying \(\sigma(z_{A,\alpha})=\mathrm{norm}(P_\alpha)\) and \(\sigma(z'_{A,\alpha})=\mathrm{norm}(P'_\alpha)\), each \(z=z_{A,\alpha},z'_{A,\alpha}\) satisfies
\[
 \operatorname{mean}_v\log^+|z|_v\le
 C_{\rm arch}f(2N)B\bigl(L+\log(2NB)\bigr)
 +C_{\rm arch}fN^2\bigl(L+\log(2fNB)\bigr)
 +C_{\rm arch}N^2\log Q ,
\]
where \(\operatorname{mean}_v\varphi(v)=\frac2{[K:\Q]}\sum_{v\mid\infty}\varphi(v)\) is the mean over infinite places of \(K\) and \(|z|_v\) denotes the standard complex absolute value at \(v\).
\item Assume \(\mathscr P\). There is \(C_{F0}\in\R\) such that for each \(C_F\ge C_{F0}\), in the closed sense of Definition~\ref{B8.39} on the family class of Definition~\ref{B8.24} under the division bound of Lemma~\ref{B8.40a}, there is \(C_0\in\R\) such that for every \(C_1\ge C_0\), eventually at every member point \(p\in\mathcal P\) and every fitted evaluation datum \(d\) at \((p,C_F)\), every \(0\le i,j<m\) and each \(z=z_A(i,j),z_A'(i,j)\) satisfies \(\operatorname{mean}_v\log^+|z|_v\le A_0\), where
\[
 A_0=C_1fN_*\bigl[2B_w\bigl(L+\log(2N_*B_w)\bigr)+N_*\bigl(L+\log(2fN_*B_w)\bigr)\bigr]+C_1N_*^2c_S .
\]
\end{enumerate}
\end{lemma}

\begin{proof}[Outline]
For (1), pass to a Galois field \(M\), use weighted Cauchy--Schwarz to reduce to averaging \(\log^+S_\rho\), apply class, minor, and conjugate coordinate bounds, and average via Lemma~\ref{B8.21ma}. For (2), combine the family logarithmic-derivative bound with (1). The complete proof is in \S\ref{prf:B8.21m}.
\end{proof}

\begin{lemma}[Existence of a fitted evaluation datum]\label{B8.21n}
Assume \(\mathscr P\), the Hecke--Tate functional equation, Theorem~\ref{B8.11}, and the ordinary lattice-norm transport and partial-zeta denominator clauses of \(\mathscr J\). There is \(C_{F0}\in\R\) such that for each \(C_F\ge C_{F0}\), in the closed sense of Definition~\ref{B8.39} on the family class of Definition~\ref{B8.24} under the division bound of Lemma~\ref{B8.40a}, eventually at every member point \(p\in\mathcal P\), every grid \(G\) fitting \(p\) at \(C_F\) and Weil system \(W\) on \(G\) over \(p\)'s real realization admitting a complete choice that satisfies both the evaluation-field condition and the square-field condition of Lemma~\ref{B8.20} extends to a fitted evaluation datum \(d\) at \((p,C_F)\) with underlying grid \(G\), realizing the higher toric data of Lemma~\ref{B8.21a} and the bottom weight bound \(B_w\) of Definition~\ref{B8.21ca}.
\end{lemma}

\begin{proof}[Outline]
Derives rationality of row sums and seesaw ratios from \(\mathscr P\), the Hecke-Tate functional equation, and \ref{B8.11}. Combines toric and denominator witnesses, bottom pairs, and \(d_\ell\) from \ref{B8.21k} with \(M\) and \(V_t^0\) to satisfy \ref{B8.21d}(3). The complete proof is in \S\ref{prf:B8.21n}.
\end{proof}

\begin{lemma}[Comparison map and base kernel on a fitted grid]\label{B8.21o}
Assume \(\mathscr P\) and Theorems~\ref{B8.11} and \ref{B8.15}. There is \(C_{F0}\in\R\) such that for each \(C_F\ge C_{F0}\), in the closed sense of Definition~\ref{B8.39} on the family class of Definition~\ref{B8.24} under the division bound of Lemma~\ref{B8.40a}, eventually at every member point \(p\in\mathcal P\) and every fitted evaluation datum \(d\) at \((p,C_F)\) with unnormalized matrices \(V_t^0\) and seesaw products \(r_t=\prod_{i,j}r_{ij,t}\), one has \(r_{\mu_0}=1\) and \(r_{\mu(w)}\ne0\) for every finite place \(w\) of \(K\), \(RR_w=\iota_{\mu(w)}P_{\mu(w)}M_{\mu(w),w}\) at every finite place \(w\) of \(K\), \(\iota_{\mu_0}P_{\mu_0}\) maps \(K^{n_V(\mu_0)}\) onto \(RR\otimes_{\mathcal O_K}K\), \(\ker V_{\mu_0}^0=\ker P_{\mu_0}\), and there exist roots of unity \(\zeta_t\in K\) (\(t\in T\)) with \(\zeta_{\mu_0}=1\) and a \(K\)-linear map \(\Psi:RR\otimes_{\mathcal O_K}K\to M\otimes_{\mathcal O_K}K\) satisfying
\[
 \Psi\iota_tP_t=r_t^{-1}(\zeta_t^{-1}V_t^0)\qquad(t=\mu_0\text{ or }t=\mu(w)).
\]
\end{lemma}

\begin{proof}[Outline]
Gets surjectivity of \(\iota_{\mu_0}P_{\mu_0}\) and gluing from link data, and the test expansion from \ref{B8.11} and \ref{B8.21d}(3). Row separation in \ref{B8.15} gives \(\kappa\ne0\), yielding injective \(\Psi\) and \(\ker V_{\mu_0}^0=\ker P_{\mu_0}\). The complete proof is in \S\ref{prf:B8.21o}.
\end{proof}

\begin{lemma}[Fitted-grid size, weight, and modulus bounds]\label{B8.21p}
For every \(C_F\ge0\), in the closed sense of Definition~\ref{B8.39} on the family class of Definition~\ref{B8.24} under the division bound of Lemma~\ref{B8.40a}, for every \(C_1\ge42\), eventually at every member point \(p\in\mathcal P\) and every fitted evaluation datum \(d\) at \((p,C_F)\),
\[
 \begin{gathered}
 1\le x,\quad 1\le L,\quad L+\log x\le3L,\quad
 2\le m,\quad m+3\le k,\quad m^4\le f\le k,\\
 mfN_*^2\le b,\quad 80R_*^4\le N_*,\quad
 1\le n_p\le2mR_*,\quad 1\le R_p\le C_1mN_*^3b,\quad \log(n_pR_p)\le4L,\\
 \Sigma_0\le2P_{\rm pf},\quad c_S\le45fL,\quad
 1\le B_w\le C_1b,\quad \log B_w\le C_1L,
 \end{gathered}
\]
and
\[
 100000\,C_1^3fm^2N_*^2(R_*N_*L+P_{\rm pf})\le x^{C_H}.
\]
\end{lemma}

\begin{proof}[Outline]
Derives rank, gap, and weight bounds from the family windows and \ref{B8.21e}, using \(b=\exp(x^\varepsilon)\) to control \(n_p,R_p\), and \(\log(n_pR_p)\). Sums prime logarithms over \(S_0\) for \(c_S,\Sigma_0\) and bounds the polynomial prefactor by \(x^{C_H}\). The complete proof is in \S\ref{prf:B8.21p}.
\end{proof}

In Proposition~\ref{prop:p8c2-fitted-evaluation}, use the member-point maps, bottom weight bound, and evaluation budgets of Definitions~\ref{B8.21b}, \ref{B8.21ca}, and \ref{B8.21e}, and fix any estimate constant \(C_1\ge C_0\) above a joint threshold \(C_0\ge42\) depending on the eligible family (with the point cutoff \(x_0\) allowed to depend on \(C_1\)).

\begin{proposition}[Evaluation estimates on a fitted grid]\label{prop:p8c2-fitted-evaluation}
Use the hypotheses and fitted choice of Theorem~\ref{B8.22}. There are a Hermitian evaluation lattice \(M\), evaluation matrices \(V_t\), a map \(\Psi:RR\otimes K\to M\otimes K\), and integers \(d_F(w)\) supported on a finite set \(T_F\) with the following simultaneous properties. Put \(r_t=\prod_{i,j}r_{ij,t}\) for the product of the layer-pair factors, with \(r_{\mu_0}=1,\ r_{\mu(w)}\ne0\). After a unit normalization,
\begin{equation}\label{eq:p8c2-psi-comparison}
 \Psi\iota_tP_t=r_t^{-1}V_t\quad(t=\mu_0\text{ or }t=\mu(w)),\qquad
 \ker V_{\mu_0}=\ker P_{\mu_0}.
\end{equation}
Also \(RR_w=\iota_{\mu(w)}P_{\mu(w)}M_{\mu(w),w}\) and \(\iota_{\mu_0}P_{\mu_0}\) is onto \(RR\otimes K\). With direct quotient norms,
\begin{equation}\label{eq:p8c2-evaluation-estimates}
 \begin{gathered}
 V_{\mu(w)}(M_{\mu(w),w})\subseteq\mathfrak p_w^{-d_F(w)}M_w,\qquad
 2\sum_w\ell_w d_F(w)\le B_F,\\
 h(\Psi\iota_{\mu_0})\le H_{I,0},\qquad \mu_{\max}(M)\le E_M ,
 \end{gathered}
\end{equation}
where \(h\) is the mean logarithm of the squared operator norm at infinity.

Use the half-ranks \(h_{ij}\) of Lemma~\ref{B8.20}. Write \(z_A,z_A'\) for the bottom numbers, \(z_T,z_{T,c},g_T\) for the higher data of Lemma~\ref{B8.21a}, and \(c(w)\) for the selected twist. Then
\begin{equation}\label{eq:p8c2-one-pair}
 \begin{aligned}
 v_w(r_{ij,\mu(w)})\le{}&
 |v_w(z_A)|+|v_w(z_A')|\\
 &+4\sum_{s=1}^{h_{ij}}\bigl(|v_w(z_T)+g_T(w)|+|v_w(z_{T,c(w)})+g_T(w)|\bigr)\\
 &+5\sum_{s=1}^{h_{ij}}\sum_{l\in S_d}N_0(i,j,s,l)v_w(l)
 +2h_{ij}\sum_{l\in S_e}\frac{v_w(l)}{l^2}\\
 &+6d_\ell(w)+2y_{ij}\sum_{q\in S_0}c_qv_w(q).
 \end{aligned}
\end{equation}
Here \(S_d,S_e\) are the finite rational prime sets for the count and degree terms, \(d_\ell\ge0\) is the local denominator allowance, and \(S_0,c_q\le e_q\) are the grid primes and modulus exponents. The indices satisfy \(y_{ij}+1\le m\) and \(h_{ij}\le N_*\). The prime count satisfies
\begin{equation}\label{eq:p8c2-count-budget}
 \sum_{l\in S_d}N_0(i,j,s,l)\log l\le B_0
 \qquad(1\le s\le h_{ij}),\qquad N_0(i,j,s,l)\ge0 .
\end{equation}
For bottom weight bound \(B_w\) and logarithmic level allowance \(c_S\) put \(l_1=\log(2fN_*),\ l_2=\log(2N_*B_w),\ l_3=\log(2fN_*B_w)\), and
\[
 \begin{aligned}
 D_0&=2C_1fN_*(l_1+B_wL+2N_*B_wc_S),\\
 A_0&=C_1fN_*[2B_w(L+l_2)+N_*(L+l_3)]+C_1N_*^2c_S .
 \end{aligned}
\]
For every finite set \(U\) of finite places and each \(z=z_A,z_A'\),
\begin{equation}\label{eq:p8c2-bottom-budget}
 \sum_{w\in U}\ell_wd_\ell(w)\le D_0,\qquad
 z\ne0,\quad v_w(z)\ge-3d_\ell(w),\qquad
 \operatorname{mean}_v\log^+|z|_v\le A_0 ,
\end{equation}
where the last mean is over infinite embeddings.
\end{proposition}

\begin{proof}
Under the hypotheses of Theorem~\ref{B8.22}, let \(C_{F0}'\) be the maximum of \(0\) and the grid thresholds in Lemma~\ref{B8.21a} and Lemmas~\ref{B8.21g}--\ref{B8.21h}, \ref{B8.21j}, and \ref{B8.21l}--\ref{B8.21o} (so \(C_{F0}'\ge0\), covering Lemmas~\ref{B8.21i} and \ref{B8.21p}). Fix \(C_F\ge C_{F0}'\), take \(\varepsilon_1>0\) to be the minimum of the exponent bounds of Lemma~\ref{B8.21a} and the closed Lemmas~\ref{B8.21g}--\ref{B8.21j} and \ref{B8.21l}--\ref{B8.21p}, and for \(0<\varepsilon\le\varepsilon_1\) combine their family threshold records \(K\). For an eligible family \(F_m\) meeting \(K\) at \(C_H\) and satisfying the division bound of Lemma~\ref{B8.40a}, let \(C_0\ge42\) be the maximum of \(42\) and the estimate-constant thresholds of Lemma~\ref{B8.21a}, Lemmas~\ref{B8.21g}--\ref{B8.21h}, \ref{B8.21j}, \ref{B8.21l}, and \ref{B8.21m}(2) (so every \(C_1\ge C_0\) also satisfies \(C_1\ge1\) in Lemma~\ref{B8.21i} and \(C_1\ge42\) in Lemma~\ref{B8.21p}); for any \(C_1\ge C_0\), take \(x_0>F_m.x_1\) above the maximum of the point cutoffs at \(C_1\) of Lemma~\ref{B8.21a}, Lemmas~\ref{B8.21g}--\ref{B8.21j}, and Lemmas~\ref{B8.21l}--\ref{B8.21p}.

Fix an indexed point \(p\in\mathcal P\) at \(x\ge x_0\) (so \(x>F_m.x_1\)) carrying a fitted grid \((G,W)\) at \(C_F\) in the sense of Definitions~\ref{B8.21b}--\ref{B8.21d}, with the evaluation budgets of Definition~\ref{B8.21e}. The hypotheses of Theorem~\ref{B8.22} include \(\mathscr P\), \(\mathscr J\), the Hecke--Tate functional equation, and Theorems~\ref{B8.11} and \ref{B8.15}. Lemma~\ref{B8.21n} therefore supplies a fitted evaluation datum \(d\) over \((p,C_F)\) with ambient dimension \(n\), Hermitian frames \((X'_v)_{v\mid\infty}\), evaluation submodule \(M\subset K^n\), matrices \(V_t^0\), seesaw factors \(r_{ij,t}\in K\), denominator function \(d_\ell\), prime sets \(S_d,S_e\), prime-count functions \(N_0(i,j,s,l)\), bottom numbers \(z_A(i,j),z_A'(i,j)\in K\), higher toric ratios \(z_T(i,j,s),z_{T,c}(i,j,s)\in K\) with integer shifts \(g_T(i,j,s,w)\), twist selector \(c(w)\), and bottom weight bound \(B_w\).

Apply Lemma~\ref{B8.21o} to \(d\). It gives the local gluing equalities \(RR_w=\iota_{\mu(w)}P_{\mu(w)}M_{\mu(w),w}\), surjectivity of \(\iota_{\mu_0}P_{\mu_0}\) onto \(RR\otimes_{\mathcal O_K}K\), the base kernel equality \(\ker V_{\mu_0}^0=\ker P_{\mu_0}\), roots of unity \(\zeta_t\in K\) (\(t\in T\)) with \(\zeta_{\mu_0}=1\), and a linear map \(\Psi:RR\otimes_{\mathcal O_K}K\to M\otimes_{\mathcal O_K}K\) such that
\[
 \Psi\iota_tP_t=r_t^{-1}(\zeta_t^{-1}V_t^0)\qquad(t=\mu_0\text{ or }t=\mu(w)),
 \qquad r_t=\prod_{i,j}r_{ij,t}.
\]
Define the unit-normalized evaluation matrices by \(V_t=\zeta_t^{-1}V_t^0\). Since \(\zeta_{\mu_0}=1\), one has \(V_{\mu_0}=V_{\mu_0}^0\), whence \(\ker V_{\mu_0}=\ker P_{\mu_0}\). Lemma~\ref{B8.21j} gives \(r_{ij,\mu_0}=1\) and \(r_{ij,\mu(w)}\ne0\) for all \(0\le i,j<m\) and every finite place \(w\) of \(K\), so \(r_{\mu_0}=1\) and \(r_{\mu(w)}\ne0\). This proves \eqref{eq:p8c2-psi-comparison}.

The finite sum estimates of Lemma~\ref{B8.21g} yield a finite set \(T_F\) of finite places of \(K\) and integers \(d_F(w)\) vanishing outside \(T_F\) such that \(V_{\mu(w)}^0(M_{\mu(w),w})\subseteq\mathfrak p_w^{-d_F(w)}M_w\) for all finite \(w\) and \(2\sum_w\ell_w d_F(w)\le B_F\). By the root-of-unity invariance assertion of Lemma~\ref{B8.21g}, replacing \(V_{\mu(w)}^0\) by \(V_{\mu(w)}=\zeta_{\mu(w)}^{-1}V_{\mu(w)}^0\) preserves the local lattice inclusion at every \(w\) with the same \(T_F\) and \(d_F\). At the base index \(t=\mu_0\), \(r_{\mu_0}=1\), \(V_{\mu_0}=V_{\mu_0}^0\), and \eqref{eq:p8c2-psi-comparison} give \((\Psi\iota_{\mu_0})P_{\mu_0}=V_{\mu_0}^0\). Viewing \(M\otimes_{\mathcal O_K}K\subseteq K^n\) and applying the base evaluation estimate of Lemma~\ref{B8.21h} to \(V_b=\Psi\iota_{\mu_0}\in M_{n\times d_Q(\mu_0)}(K)\) gives \(h(\Psi\iota_{\mu_0})\le H_{I,0}\). The frame estimate for the evaluation lattice in Lemma~\ref{B8.21i} shows that \(M\) is a nonzero finitely generated Hermitian lattice with nondegenerate Hermitian frames (\(\det X'_v\ne0\)) at infinity and \(\mu_{\max}(M)\le E_M\). This establishes all four inequalities in \eqref{eq:p8c2-evaluation-estimates}.

Since \(d\) realizes the higher toric data \(z_T,z_{T,c},g_T\) of Lemma~\ref{B8.21a} by Definition~\ref{B8.21d}(3) and Lemma~\ref{B8.21n}, the square valuation estimate of Lemma~\ref{B8.21j} gives primality of the elements of \(S_d\) and \(S_e\), the index bounds \(y_{ij}=m-1-j\) (hence \(y_{ij}+1\le m\)) and \(h_{ij}\le N_*\), the modulus comparison \(c_q\le e_q\) for every \(q\in S_0\), the valuation bound \eqref{eq:p8c2-one-pair} at every finite place \(w\) of \(K\), and the prime-count bound \eqref{eq:p8c2-count-budget} for \(1\le s\le h_{ij}\).

For the bottom numbers \(z_A=z_A(i,j)\) and \(z_A'=z_A'(i,j)\), with \(B_w\) equal to the bottom weight bound of Definition~\ref{B8.21ca} by Definition~\ref{B8.21d}(3), apply the fitted-grid lattice norm bound of Lemma~\ref{B8.21l}. This gives \(d_\ell(w)\ge0\), \(\sum_{w\in U}\ell_w d_\ell(w)\le D_0\) for every finite set \(U\) of finite places of \(K\), \(z\ne0\), and \(v_w(z)\ge-3d_\ell(w)\) for \(z=z_A,z_A'\). Applying the archimedean norm estimate of Lemma~\ref{B8.21m}(2) to the same bottom pairs with the same parameter substitution gives \(\operatorname{mean}_v\log^+|z|_v\le A_0\) for \(z=z_A,z_A'\), proving \eqref{eq:p8c2-bottom-budget}.

Finally, we check that every height and denominator loss in \eqref{eq:p8c2-evaluation-estimates}--\eqref{eq:p8c2-bottom-budget}, as well as their combined contribution to the slope bound \eqref{eq:p8c2-fun}, is bounded by \(x^{C_H}b\) and is linear in \(b=\exp(x^\varepsilon)\). Because \(C_1\ge42\), Lemma~\ref{B8.21p} gives \(1\le L=x^\varepsilon\), \(L+\log x\le3L\), \(2\le m\), \(1\le f\le k\), \(R_*\le N_*\), \(mfN_*^2\le b\), \(\Sigma_0\le2P_{\rm pf}\), \(c_S\le45fL\), \(1\le B_w\le C_1b\), \(\log B_w\le C_1L\), \(1\le n_p\le2mR_*\), \(1\le R_p\le C_1mN_*^3b\), and \(\log(n_pR_p)\le4L\). Put \(U=C_1^3fm^2R_*N_*^3bL\) and \(V=C_1^3fm^2N_*^2P_{\rm pf}b\). Both \(U\) and \(V\) are products of \(b\) with polynomial-in-\(x\) factors, since \(L=\log b=x^\varepsilon\), \(P_{\rm pf}\) depends on \(b\) through \(\log b=L\) alone, and all rank, degree, and auxiliary-prime logarithms grow polynomially in \(x\). No inversion of \(\Psi\iota_{\mu_0}\) occurs: \(h(\Psi\iota_{\mu_0})\) is bounded forward on the direct quotient metric by \(H_{I,0}=C_1fmN_*bL+C_1m^2N_*^2\Sigma_0\le U+2V\) (with \(b\ge1\)), and the source-frame shift term \(3C_1fmN_*^2L\) used later in Theorem~\ref{B8.22} satisfies \(3C_1fmN_*^2L\le3U\). Likewise \(R_p\) enters \(E_M=4fn_pR_p\log(n_pR_p)\) to the first power while \(\log(n_pR_p)\le4L\), giving \(E_M\le32U\); the finite evaluation budget satisfies \(B_F\le2U+4V\); the toric budget of Lemma~\ref{B8.21a} satisfies \(B_0\le3C_1fbL\le3U\); and \(B_w\le C_1b\) enters \(D_0\) and \(A_0\) to the first power multiplied by \(L\), \(c_S\le45fL\), \(l_1\le1+L\) (from \(2fN_*\le mfN_*^2\le b\)), and \(l_2,l_3\le1+(C_1+1)L\), giving \(D_0\le186C_1^2f^2N_*^2bL\) and \(A_0\le57C_1^3fN_*^2bL\). Setting \(B=\max(A_0+6D_0,B_0)\le1176C_1^3f^2N_*^2bL\) and \(G_*=m^2\bigl(2(2+8N_*)B+10N_*B_0+2N_*+6D_0+4m\Sigma_0\bigr)\) (the combined budget for \(\sum_w\ell_w v_w(r_{\mu(w)})\) obtained in Theorem~\ref{B8.22} from \eqref{eq:p8c2-one-pair}--\eqref{eq:p8c2-bottom-budget} and Lemma~\ref{B8.21a}), the bounds \(f\le k\le R_*\) and \(2(2+8N_*)\le20N_*\) give \(2G_*\le49336U+16V\). Every individual loss and the total budget \(E_M+H_{I,0}+3C_1fmN_*^2L+2G_*+B_F\) are therefore bounded by \(100000C_1^3fm^2N_*^2(R_*N_*L+P_{\rm pf})b\), which is at most \(x^{C_H}b\) by the final inequality of Lemma~\ref{B8.21p}.
\end{proof}

\begin{theorem}[The constant-function Fourier construction satisfies the functional bound]\label{B8.22}
Assume \(\mathscr P,\mathscr J\), the Hecke--Tate functional equation and Theorems~\ref{B8.11} and \ref{B8.15}. Assume the following closed assertion, in the sense of Definition~\ref{B8.39}, on the family class of Definition~\ref{B8.24}. Eventually at every member point in \(\mathcal P\) of Definition~\ref{def:p8c3-admissible}, for all \(C\in\R\), every layer \(j\), block \(p\) and infinite place \(v\), the level-one base eigenvalue satisfies \(\log\lambda^{(1)}_{j,\mu_0,p,v}\ge L_{6,j}\), where \(L_{6,j}\) is the hook-length lower bound of Theorem~\ref{B5.07}.

There is \(C_{F0}'\in\R\) such that for each \(C_F\ge C_{F0}'\) there is \(\varepsilon_1>0\), and for each \(0<\varepsilon\le\varepsilon_1\) there are size thresholds \(K\). Let \(F_m\) belong to the family class of Definition~\ref{B8.24}, let its constants at \(C_H\) meet the thresholds \(K\) of Definition~\ref{B8.39}, and assume the division bound of Lemma~\ref{B8.40a}: for some \(A,x_1\), every point at \(x\ge x_1\) satisfies \(\eta_j\le A D_0L(B_0+N_c)\) in every lift layer \(j\). There is \(x_0\) after which every member point in \(\mathcal P\) carrying a fitted grid at \(C_F\) as specified below satisfies, at level zero,
\[
 \forall C\in\R,\qquad \mu_{\max,RR}\le x^{C_H}b-\langle\log\kappa\rangle_\infty .
\]
The fitted grid must carry a Weil system realizing the given theta action and a complete choice satisfying the coefficient-field condition of Theorem~\ref{B8.17} and the square-field condition of Lemma~\ref{B8.20}. Its finite-group series must satisfy Definition~\ref{def:p8c2-series}.
\end{theorem}

\begin{proof}[Outline]
Glues inclusions from \ref{prop:p8c2-fitted-evaluation} for injective \(\Psi\) and bounds \(h(\Psi)\) via \ref{B5.07}. Sums finite losses via \ref{B8.21a}, then applies \ref{B8.21p} to bound \(\mu_{\max,RR}+\langle\log\kappa\rangle_\infty\) by \(x^{C_H}b\). The complete proof is in \S\ref{prf:B8.22}.
\end{proof}

\section{The functional row on the chosen members}

\begin{theorem}[The functional condition along members]\label{B8.23}
Assume discrete-series coverage, Harish--Chandra realization, rapid decay of cusp forms, Casselman--Wallach globalization, the ray-class-field input and Theorem~\ref{B8.22}. There is \(C_{F0}\in\R\) with the following property. For \(C_F\ge C_{F0}\), take an allowed local Langlands system \(h_L\) with auxiliary extension \(S\). Fix a geometry condition \(\Gamma\). For every norm-choice extension \(N\), assume each datum satisfying \(N\) admits an extension with list at \(C_F\). Adjoin to \(N\) the grid extension condition at \(C_F\) of Lemma~\ref{B8.20} and assume the open conditions of Theorem~\ref{B8.08}.

A condition \(R\) eventually holds for every constructed member over the permitted base frames. On good members satisfying \(R\), a row condition \(E\) eventually holds for every linked bundle. For a compatible selection on \(\mathcal A\) for this row, assume the closed hypotheses on \(W_3(\mathcal A)\) and the division bound of Lemma~\ref{B8.40a}. There is \(C_0\ge0\) such that for each \(C\ge C_0\) the selected bundles eventually satisfy the slope inequality \eqref{eq:p8c2-fun} at parameters \((C,C_H)\), with point threshold allowed to depend on \(C\). The class \(W_3(\mathcal A)\) requires \(W_1(\mathcal A)\), a level-\(3\) height choice and eventual bounds on \(T_u,D_{\Gamma,u}\). For \(D_{\Gamma,u}\) some \(C',x_0\) give \(\log D_{\Gamma,u}\le C'fN(m)\log(fN(m))\) on every member at \(x\ge x_0\) and every word \(u\).
\end{theorem}

\begin{proof}[Outline]
Transport and intersect the divisor and geometry data to select fitted grids tied to the original bundles. Realization and the family conditions of \ref{B8.01k} let \ref{B8.22} apply at its threshold, and the ties carry the inequality back with \(C_0=0\). The complete proof is in \S\ref{prf:B8.23}.
\end{proof}

Use Theorem~\ref{B8.18} in the member supply of Theorem~\ref{B8.07}. Apply Theorem~\ref{B8.03} at \(c_{\rm Dep}=C_F=\max(C_{F0},C_{F1})\) with the grid and functional rows, the linked \(2,3\)-bound of Theorem~\ref{B8.01} and the covolume row of Theorem~\ref{B8.34}.



% ---- p8c3 ----
\chapter{Covolume bounds}\label{ch:B8-covolume}

Let \(\Lambda\) be a rank-\(r\) lattice over a number field \(K\) with \(s_K\) infinite places. In the archimedean frame \(X\), write \(\operatorname{ht}(X,\Lambda)\) for its \(\Z\)-covolume multiplied by \(2^{s_Kr}|\operatorname{disc}K|^{-r/2}\). Put
\[
 \operatorname{lcv}(X,\Lambda)=\frac{2}{[K:\Q]}\log\operatorname{ht}(X,\Lambda).
\]
For a nonzero map \(\varphi\), let \(h(\varphi)\) be the mean over the infinite places of \(\log\|\varphi\|_v^2\).
\begin{equation}\label{eq:p8c3-map-height}
 \|\varphi\|_v\le e^a\quad(v\mid\infty)
 \quad\Longrightarrow\quad h(\varphi)\le2a .
\end{equation}

At a point \(p\) over \(x\), let \(E\) be the host field and put \(f=[E:\Q]/2\), \(L_x=x^\varepsilon\), \(b=e^{L_x}\) and \(E_p=fmbL_x\log\log x\). At layer \(0\le i<m\), write \(j=i+1\). For the layer weights \(w_{j,\rho}\), set
\begin{equation}
 M_{p,i}=(k+j-1-(m-j))\log\Omega\sum_\rho w_{j,\rho},
 \qquad \ell_{3,i}(C,C_H)=M_{p,i}+CE_p .
 \tag{8.24a}\label{eq:p8c3-main}
\end{equation}
At fixed \(p,i,C\) the right side is independent of \(C_H\). For the host field \(E\) and the Kronecker character \(\chi_d\) of \(\Q(\sqrt{-d})\), the period of \ref{B5.07} satisfies
\begin{equation}\label{eq:p8c3-period}
 \log\Omega=\gamma+\log(2\pi/\operatorname{rd}(E))-L_E,\qquad
 L_E=\Re\frac{L'(1,\chi_d)}{L(1,\chi_d)} .
\end{equation}
Here \(\operatorname{rd}(E)=|\operatorname{disc}E|^{1/[E:\Q]}\).

\section{The estimate and its geometric inputs}

Let \(\mathcal U\) be the class of all families. For a family \(F_m\), write \(x_1\) for its lower threshold and \(I_x\) for its index set at \(x\). Let \(\mathcal F_{\rm cov}\) consist of the families in \(W_3(\mathcal U)\) satisfying clauses 1--3 of Definition~\ref{B8.01k} and the following wall bound. Beyond one family cutoff, at every index the prime sum of \ref{B5.22b}, evaluated at the point's own weights, is at most \(x^{-9}\). We use admissibility in the sense of Definition~\ref{def:p8c3-admissible}.

\begin{definition}[Admissible points for the covolume bound]\label{def:p8c3-admissible}
Let \(p\) be a point of the family, \(K\) its realisation field and \(r=k+m-1\). The next three paragraphs define the class \(\mathcal P\) of member points used for the lower-weight estimate. An admissible point for the covolume bound is a point of \(\mathcal P\) satisfying also the two arithmetic requirements below.

At each step of the theta chain, including the last step, the intertwiners and the maps defined by integration against the theta kernel exist. For every tuple the resulting theta maps agree under evaluation, both on the base and on each layer. At the embeddings extending the prescribed base embedding, evaluation at the chosen vector is injective. At the other embeddings the conjugate evaluation is realised by a unit vector whenever the conjugate weight space is identified with the layer space. On each core cusp space the finite level group acts over the core field, preserves integral sections, and under evaluation acts by the realised action multiplied by the prescribed character. In every layer the extremal vector pairs to the prescribed product of minors, and the one-map comparison holds, including the comparison over an extension containing the image of \(K\) under the prime restrictions of \ref{B8.55}.

The primitive and divisor rows satisfy the identities of \ref{B8.08}. Distinct finite indices in the reference row have distinct transported types at every layer. The level-one row and its lifts have the prescribed Hecke decompositions: the inert operators commute, split over the specified field and commute with the finite group action, with distinct indexed inert places. The frames at the two auxiliary primes have the prescribed bases. Both global row lattices agree with their indefinite realisations, the rational lifts span the indicated spaces, and the theta maps and column pairings agree under these identifications. The two level characters come from the prescribed residue-field additive characters under the group identification and satisfy the Weil character relation. At the reference tuple and base character, the selected row subspace has dimension \(D_0\). The rational layer maps into the exported row spaces induce the prescribed maps on coordinates and identify the layer lattices and global row lattice. Every finite row type satisfies the local type conditions of \ref{B8.08}. Every place passing the single-index test for the reference numerical data at \((0,0)\) belongs to the twisting set.

The point comes from compatible spectral and indefinite realisations with the prescribed normalisation and recursion. The core data, base pairing, tame types and series agree under transport. The reference subgroup corresponds to the prescribed level subgroup, and its inverse character gives the two local level characters after scalar extension. At each twist index and layer, on either side, every theta datum transforming by the level character belongs to the allowed set. At each place with nonreference tuple the twist character takes values in the prescribed coefficient field. The embedded realisation field satisfies the ray-level condition at \(g_0g_1p_{\rm de}p'_{\rm de}\), contains the \(g_a(g_a-1)\)-th roots of unity for \(a=0,1\), and contains the ray fields at the tuple levels. The point satisfies the twist and basis conditions of clause 7 of \ref{B8.01k} at its own twist data.

For the first arithmetic requirement let \(F\subset\C\) be the field generated over \(\Q\) by the image of the coefficient ring in the prescribed complex embedding. Put \(H=19fr^2L_x+198f^2L_x\). Require
\[
 \log[F:\Q]\le H,\qquad \log\operatorname{rd}(F)\le H .
\]
For the second, at each \(g=g_0,g_1\) there must be a label datum of \ref{B7.03} with the prescribed Teichmüller map. For every \(P\in\F_g[X]\), the specified series element in \(\C[\GL_r(\F_g)]\) must equal the sum of the central primitive idempotents of the irreducible characters with label \(P\) for that same datum.
\end{definition}

\begin{definition}[The family covolume bound]\label{B8.24}
For each eligible \(F_m,C_H\) in the closed range of \ref{B8.39} on \(\mathcal F_{\rm cov}\), there exist \(C_0\ge0\) and \(x_0\in\R\) such that for every \(x>x_1\) with \(x\ge x_0\), every \(\iota\in I_x\), every admissible point over \((x,\iota)\), every \(C\ge C_0\) and every \(0\le i<m\),
\begin{equation}
 \frac{\operatorname{lcv}(X_{W,i},\Lambda_i)}{D_{\rm norm}}\le M_{p,i}+CE_p .
 \tag{8.24b}
\end{equation}
Here \(X_{W,i}\) and \(\Lambda_i\) are the frame and lattice of layer \(i\) in the reference indefinite realisation at \(p\). The normalising dimension is \(D_{\rm norm}=\dim X_E\), with the convention \(a/0=0\). 
\end{definition}

Write \(r=k+m-1\) and \(n_V=r+t\), where the chain ranks are \(N(j)=k+j(t+1)\). The precision parameter at the point is \(P_*=C_P(frL_x)^{e_P}\), where \(C_P\ge1\) and \(e_P\ge12\) are the chosen precision constants. The layer signature is \((P,Q)=(k+j-1,m-j)\). Use one CM embedding \(\tau\) per conjugate pair. Call a deep cover \(Y\) of the layer good with parameters \((e,o)\), for \(e\in\N\) and \(o\in\R\), under the following conditions. Its level extends the base level by the coprime factor \(15^e\), with auxiliary level parameter one and complementary conjugate ideal factors. The level sections respect the covering map, every geometric frame lifts, and the weight-zero comparison section is nowhere zero. The finite adelic subgroup is the principal congruence subgroup of level \(M'=g_0g_1 15^e\) in the lattice stabiliser. The cover's coefficient field \(F'\) is relatively algebraically closed in the component function fields and satisfies \(\log[F':\Q]+\log\operatorname{rd}(F')\le o\). The intermediate coefficient extension used to form the cover is normal over the core coefficient field. At each embedding extending the prescribed base embedding, the components have their analytic uniformisations, the uniformised frames come by base change from test rings, and the characters of dominant weights are trivial on the component groups acting on the domains. The adelic frame parametrisation agrees with uniformisation. Its distinguished minors have absolute value one. In rank \(r\) the component Gram matrix is the negative transpose of the layer Gram matrix and its polarising scalar is the negative of the host scalar, of prescribed absolute value at every embedding. One embedding of the layer lattice identifies it with the uniformising lattice and gives compatible orthogonal frames on all components at every such embedding. At a nonterminal layer, a good cover without displayed parameters means one with \(e=\lfloor3L_x/\log3\rfloor_{\N}\) and \(o=391(frL_x)^2\).

Use the row weight \(\lambda\) and the anchored weights \(u_W,u_M\) of the quotient construction \eqref{eq:p8c3-quotient}. Write \(\eta_H\) for the Hodge weight, with entry zero in each of the first \(P\) columns and \(-1\) in each of the last \(Q\) columns at every embedding. In zero-based column indices the anchor is \(a_\tau=\lambda_{\tau,k-2}\). The central shift \(c_0=c_0(D,j)\) is the least multiple of \(8\) at least
\[
 \max\!\left(r+\binom{fPQ+1}{2}+2,\ 8fr+4\right).
\]
Write \(G_s(s)\) for the full row weight at shift \(s\). When a row such as \(u_M\) occurs as the argument of a section lattice or its slope, give it scalar exponent zero. The notation \(\mu_Y\) denotes the slope of the indicated section lattice in the squared-norm convention.

\begin{definition}[Section lattices and flat degree]\label{def:p8c3-section-laws}
Let \(r_Y(G)\) be the rank over the coefficient field of the integral cusp lattice of weight \(G\) on \(Y\). A weight has a regular section lattice if its Petersson norms admit Gram frames and its integral cusp sections form a finitely generated module over the coefficient integers spanning the cusp space. On a component this condition uses the sections supported on that component and their component norms.

Deep goodness means that the full row weights at shifts \(0,-c_0,c_0\) and the anchored weight \(u_M\) have regular section lattices, with \(r_Y(G_s(-c_0))>0\). The set of components is finite and nonempty, and the lattice of weight \(u_M\) is regular on each component.

The flat-degree law requires that on every component, for every CM embedding \(\tau\), the flat line at \(\tau\) have an invertible section with constant norm on the frames at each infinite place and constant valuation on the frames at each finite place after subtraction of the local different defect. The corrected finite valuations have finite support. Every such section has degree
\[
 \frac{P-Q}{2}\log\Omega+Q\bigl(\log d-\log\operatorname{rd}(E)\bigr).
\]
\end{definition}

\begin{lemma}[The flat line in the row metric]\label{B8.24a}
For data satisfying the flat-degree law on a nonempty component system of signature \((P,Q)\), the first Chern value at \(\tau\) is
\begin{equation}\label{eq:p8c3-flat-c1}
 c_{1,\tau}=-\tfrac{P-Q}{2}\log\Omega-Q(\log d-\log\operatorname{rd}(E)).
\end{equation}
For a good cover with the deep goodness, flatness, anchor and family size conditions, one has
\[
 \mu_Y(G_s(-c_0))-\mu_Y(u_M)
 \ge-M_{p,i}-1183E_p .
\]
\end{lemma}
\begin{proof}[Outline]
The CM eigenlattice comparison of \ref{B5.07x} and the local different give the first Chern value. The integral twist preserves rank, and comparison with \(M_{p,i}\) and bounds on the different, twist and bad primes give the slope loss \(1183E_p\). The complete proof is in \S\ref{prf:B8.24a}.
\end{proof}

\begin{definition}[The assertions in the Hodge-bundle hypothesis]\label{B8.25l}
Write \(\mathcal H_P\) for the conjunction of the following assertions. Each has its own universally quantified data. For a Hermitian line \(\overline L\) on an integral projective arithmetic variety of generic dimension \(D\), write \(I(\overline L)\) for its arithmetic self-intersection of order \(D+1\).

\begin{enumerate}
\item For a unitary datum of positive rank, a level, its moduli open and its coefficient field, there is a model over that field, smooth of the moduli dimension, with simple normal crossings boundary and the prescribed logarithmic Kodaira--Spencer isomorphism on the designated interior.

\item If \(\overline L\) is nef on fibres and its metric is semipositive, then \(I(\overline L)\le\widehat{\operatorname{vol}}(\overline L)\), with the normalisation of Theorem~\ref{B8.27}.

\item If \(\overline L\) is arithmetically nef, then \(\widehat{\operatorname{vol}}(\overline L)=I(\overline L)\).

\item Arithmetic heights and the geometric and arithmetic intersection functionals exist for finite families of Hermitian lines. The arithmetic functional obeys the rational-section recursion, is invariant under isometry, and is multilinear in tensor powers. Dominant generically finite pullback multiplies it by the generic degree. Adding a smooth conjugation-invariant function \(u\) to one metric weight adds one half of the integral of \(u\) against the product of the remaining Chern forms, summed over the embeddings. In the constant log-norm shift convention this changes \(I(\overline L)\) by \((D+1)t\deg(L_{\Q}^{D})\) for a shift of size \(t\).

\item For an abelian scheme of relative dimension \(g\) over an affine base, \(H^1_{\rm dR}\) is finite locally free of rank \(2g\). Its classes are primitive for the group law, cup product identifies its second exterior power with \(H^2_{\rm dR}\), and the Hodge submodule is a direct summand of rank \(g\). The Hodge sequence commutes with base change. A lift over a nilpotent divided-power thickening identifies the crystal with the de Rham module of the lift. The polarisation pairing is zero on the Hodge submodule and is perfect for a principal polarisation.

\item For a nonzero integral modulus of a number field there is a finite abelian extension unramified away from the modulus. The Frobenius elements away from it generate the Galois group. On coprime ideals the common kernel of the characters is the principal ray subgroup, with totally positive generators congruent to one at the modulus.

\item A principally polarised abelian variety over a finite extension of \(\Q_p\) has a semiabelian model after a finite extension.

\item Let \(G\) be the special unitary group of a nondegenerate Hermitian space of rank at least two over a CM extension \(E/F\). For a finite set \(S\) of places such that \(\prod_{v\in S}G(F_v)\) is noncompact, \(G(F)\) is dense in the adeles of \(G\) away from \(S\). Here \(G\) is simply connected and absolutely almost simple, as required for strong approximation \cite[Th.~7.12]{platonov-rapinchuk-1994}.

\item For \(g\) prime and \(N>0\), the squarefree parameters of \(\GL_N(\F_g)\) in Theorem~\ref{B7.31} give distinct generic irreducible characters. The degree belonging to a parameter times the order of its torus is \(\prod_{i=1}^N(g^i-1)\).

\item For a bundle with a good metric on the open of a smooth compactification, the Chern integral is the intersection number of the extended classes. On its full flag bundle let \(\Delta\) be the root Vandermonde and let \(v\) be the fibre point class. If the alternation of a root polynomial \(Q\) is \(\Delta R\), with \(R\) symmetric in each factor, then \(Q\) and \(Rv\) have the same top intersection on the flag bundle.

\item Nonsingular Hermitian forms of equal rank over a CM field are isometric if their signatures agree at every archimedean place and the quotient of their determinants is a norm from the CM field to its real subfield.

\item For \(\delta,x\in F^\times\), where \(F\) is a number field, \(u^2-\delta v^2=x\) has a solution in \(F\) if it has a solution in each real and finite completion.

\item For semiabelian schemes over a normal Noetherian base, a homomorphism over a dense open extends uniquely over the base \cite[Prop.~3.3.1.5]{lan-2013}.

\item Over a field of characteristic \(p>0\), the induced module for a dominant integral \(\GL_N\)-weight \(\lambda\) is nonzero and simple if \(\lambda_i-\lambda_j+j-i\le p\) for every \(i<j\).

\item A line inducing a principal polarisation on an abelian variety over an algebraically closed field has a one-dimensional space of sections.

\item If \(A\) is an abelian variety of dimension \(g\) over an algebraically closed field and \(n\) is invertible in that field, then \(\#A[n]=n^{2g}\).

\item Over a nilpotent divided-power thickening on which \(p\) is nilpotent, the de Rham module of an abelian lift is identified with the evaluation of the crystal. Given such a lift, an admissible lift of the Hodge filtration is realised by an abelian lift. A homomorphism on the closed fibre lifts uniquely when its action on the crystal preserves the two lifted filtrations.

\item Over the \(p\)-adically complete valuation ring of the algebraically closed characteristic-zero field in Lemma~\ref{B8.25a}, with \(p\) a nonunit, the contravariant prismatic module of a \(p\)-divisible group is free of rank its height. With \(\xi\) the generator of the kernel of Fontaine's map determined by the chosen roots of unity, the Frobenius image contains \(\varphi(\xi)M\), and reduction modulo \(\xi\) gives the functorial Hodge sequence. On the constant and multiplicative height-one groups the Frobenii are respectively \(\varphi\) and \(\varphi(\xi)\varphi\). Frobenius-compatible linear maps from the multiplicative module to the constant module correspond to homomorphisms from the constant group to the multiplicative group.

\item For a smooth integral projective variety over an algebraically closed field of characteristic zero and a nef and big line \(N\), \(H^i(X,K_X\otimes N)=0\) for \(i>0\).

\item On a projective \(D\)-dimensional scheme over an algebraically closed field, a top product of nef lines is nonnegative. On an integral such scheme a nef line \(L\) is big exactly when \(L^D>0\). The section count for a nef line satisfies \(h^0(X,L^m)=(L^D)m^D/D!+O(m^{D-1})\).

\item For a compact complex manifold admitting a positive line, sections prescribed on the transverse finite common zero set of \(n\) sections on an \(n\)-fold extend with the bound \eqref{eq:p8c3-extension}, under the norm and curvature hypotheses stated there.

\item For an arithmetically nef Hermitian line and a fixed Hermitian twist, the unit ball of sections is eventually finite and its logarithmic cardinality has the asymptotic \eqref{eq:p8c3-nef-count}.

\item A geometrically integral projective \(D\)-fold in characteristic zero has a Chow form over its field of definition, unique up to a nonzero scalar. It vanishes on the \((D+1)\)-tuples of hyperplanes meeting the variety in a common point over the algebraic closure. Its degree in each block is the degree of the variety.

\item Arithmetic cycle heights satisfy the rational-section recursion. If an integral arithmetic variety maps generically finitely onto its image of the same dimension, the self-intersection of the pulled-back line is the generic degree times the height of the image cycle.

\item On a smooth projective complex variety analytification identifies the algebraic and holomorphic sections of a line bundle.

\item Let \((X,\Delta)\) be a normal integral projective complex pair with \(\Delta\) effective and \(K_X+\Delta\) \(\Q\)-Cartier. If \(M\) is a line numerically equivalent to \(K_X+\Delta+A\) with \(A\) nef and big, the multiplier ideal \(\mathcal J\) has cosupport the non-klt locus and \(H^i(X,M\otimes\mathcal J)=0\) for \(i>0\). For the second assertion suppose the pair is klt, fix a closed point \(x\), and put \(n=\dim X\ge1\). If \(A^j\cdot Z>c_j^j\) for every irreducible \(j\)-fold \(Z\) through \(x\), with \(c_j>0\) and \(\sum_{j=1}^n j/c_j\le1\), there is an effective rational Cartier divisor \(B\equiv A\) such that \((X,\Delta+B)\) is log canonical near \(x\) and has \(x\) as its isolated non-klt locus there. In both assertions \(A\) is a positive rational multiple of a nef and big line.

\item On a smooth integral projective variety in characteristic zero, a nonzero coherent ideal becomes effective Cartier with simple normal crossings support after a sequence of smooth blow-ups which is an isomorphism off its support.

\item Let \(R\) be Noetherian and complete for an ideal \(I\). Compatible morphisms modulo all \(I^{j+1}\) from a proper \(R\)-scheme to a separated finite-type \(R\)-scheme algebraize uniquely. A compatible system of proper schemes over \(R/I^{j+1}\) with compatible invertible sheaves whose closed-fibre member is ample algebraizes to a proper \(R\)-scheme with a relatively ample invertible sheaf.
\end{enumerate}
\end{definition}

\begin{lemma}[A generator on each component]\label{B8.24b}
Let \(Z\) be a component of a good cover over its coefficient field \(F'\). On the moduli open of \(Z\), let \(A\) be the universal abelian scheme. For an embedding \(\tau\) in the CM type, write \(V_\tau=(\operatorname{Lie}A)_\tau\) and put
\[
             {\cal L}_\tau=\det V_\tau\otimes(\det V_{\bar\tau})^{-1}.
\]
There is an \(F'\)-rational section of \({\cal L}_\tau\) which generates the line everywhere on the open. At each infinite place its invariant norm is constant on the frames of \(Z\).
\end{lemma}
\begin{proof}[Outline]
The polarized Hodge sequence identifies the line with a de Rham determinant, and the level congruence gives trivial determinant monodromy. Deligne's regular-singular comparison algebraizes the horizontal generator, which descends to \(F'\) and has constant Hodge norm. The complete proof is in \S\ref{prf:B8.24b}.
\end{proof}

\begin{lemma}[The projected Hodge lattice and the two integral Lie frames]\label{B8.24c}
Let \(R\) be a discrete valuation ring of mixed characteristic \((0,\ell)\), with fraction field \(K\) containing the images of the embeddings of \(E\). Write \(v\) for its additive valuation. Suppose that the integral Hodge sequence, in homology, is
\[
             0\longrightarrow L^\vee\longrightarrow H
                        \longrightarrow L\longrightarrow0 ,
\]
where \(L\) is the Lie lattice, the action on its dual is conjugated, and \(H\) is locally free of rank \(r\) over \(R\otimes_{\Z}\mathcal O_E\). For an embedding \(\sigma:E\to K\), put
\[
 L[\sigma]=L\cap(L\otimes_RK)_\sigma,\qquad
 H(\sigma)=\operatorname{pr}_\sigma H,\qquad
 h_\sigma=\dim_K(L\otimes_RK)_\sigma .
\]
The Hodge sequence identifies \(\det H(\tau)_K\) with
\(\det L[\tau]_K\otimes\det L[\bar\tau]_K^{-1}\).

Complete \(R\). In the factor of \(\mathcal O_E\otimes R\) containing \(\tau\), let \(S\) be the embeddings with the same residual embedding as \(\tau\). Choose a uniformiser \(\pi\) of the corresponding completion of \(E\) over its maximal unramified subfield, and put \(a_\sigma=\sigma(\pi)\) for \(\sigma\in S\). In the above identification the determinant of a pair of integral Lie frames generates the sublattice of \(\det H(\tau)\) of valuation
\[
                     \sum_{\sigma\in S\setminus\{\tau\}}
                            h_\sigma\,v(a_\sigma-a_\tau).
\]
In particular, a section whose coefficient in a generator of \(\det H(\tau)\) is \(b\) has valuation \(v(b)\) minus this sum in the Lie frames.
\end{lemma}
\begin{proof}[Outline]
Splits the unramified part over completed \(R\) so \(H(\tau)\simeq H/TH\) (\(T=t-a_\tau\)). Lifting a basis of \(L[\tau]\) and adjoining \(\operatorname{pr}_\tau F\simeq L[\bar\tau]^\vee\) gives \(\operatorname{coker}(T|L/L[\tau])\), whose length is the stated sum. The complete proof is in \S\ref{prf:B8.24c}.
\end{proof}

\begin{lemma}[Formal biextension germs and the integral Hodge bundle on a normal model]\label{B8.24dc}
Let \(Z\) be a component of signature \((P,Q)\), with \(r=P+Q\ge1\), of a good cover over its coefficient field \(F'\), and let \(O\subset F'\) be a discrete valuation ring of mixed characteristic \((0,\ell)\) with fraction field \(F'\) and uniformiser \(\varpi\mid\ell\). Let \(\pi_X:X\to\operatorname{Spec}O\) be an integral, normal, flat, projective scheme over \(O\) with function field \(\Omega\) and generic point \(\xi:\operatorname{Spec}\Omega\to X\), whose generic fibre \(X_{F'}\) contains the moduli open \(U\) of \(Z\) as a dense open subscheme. Suppose \(X\) carries a semiabelian scheme \(G\to X\) of relative dimension \(d=\tfrac12[E:\Q]r\) with identity section \(e:X\to G\), an action of \(\mathcal O_E\) on \(G\), and a symmetric cubical line bundle \(L_{\rm PP}\) on \(G\) rigidified along \(e\), whose Mumford biextension \(\Lambda(L_{\rm PP})=m^*L_{\rm PP}\otimes p_1^*L_{\rm PP}^{-1}\otimes p_2^*L_{\rm PP}^{-1}\) on \(G\times_X G\) (where \(m,p_1,p_2:G\times_X G\to G\) are multiplication and the two projections, rigidified along \(e\times\mathrm{id}_G\) and \(\mathrm{id}_G\times e\)) satisfies the following conditions:
\begin{enumerate}
\item[(a)] for every \(a\in\mathcal O_E\), there is a bi-rigidified isomorphism \(\operatorname{ad}_a:(a\times\mathrm{id}_G)^*\Lambda(L_{\rm PP})\xrightarrow{\sim}(\mathrm{id}_G\times\bar a)^*\Lambda(L_{\rm PP})\) on \(G\times_X G\);
\item[(b)] \(L_{\rm PP}\) is isomorphic as a rigidified line bundle to \((\mathrm{id}_G,w_0)^*\Lambda(L_{\rm PP})\) for an element \(w_0\in\mathcal O_E\) with \(w_0+\bar w_0=1\);
\item[(c)] \(U\) equals the abelian locus of \(G_{F'}\) in \(X_{F'}\), and \((G_U,\mathcal O_E\text{-action},\Lambda(L_{\rm PP})_U)\) agrees with the universal principally polarised abelian scheme on \(U\) with its \(\mathcal O_E\)-action and Poincar\'e biextension;
\item[(d)] at every geometric point \(\bar x:\operatorname{Spec}\bar k\to X\) of residue characteristic \(\ell\), with Chevalley torus extension \(0\to T\to G_{\bar x}\xrightarrow{q_{\bar x}}B\to0\) of an abelian variety \(B\) by a torus \(T\), there is a rigidified line bundle \(M\) on \(B\) inducing a principal polarisation \(\lambda_B:B\xrightarrow{\sim}B^\vee\) such that \(L_{\rm PP}|_{G_{\bar x}}\simeq q_{\bar x}^*M\).
\end{enumerate}
Write \(\mathcal L=\operatorname{Lie}(G/X)\) and \(V=\mathcal L_\xi\). Then the following statements hold:
\begin{enumerate}
\item For every integer \(n\ge0\) and every nonempty open subscheme \(W\subseteq X\), pullback along the structure morphism \(\pi_n:G_W^{\times_W n}\to W\) is a bijection \(\Gamma(W,\mathcal O_W)\xrightarrow{\sim}\Gamma(G_W^{\times_W n},\mathcal O_{G_W^{\times_W n}})\).
\item Write \(\mathbf x=(x_1,\ldots,x_d)\) and \(\mathbf y=(y_1,\ldots,y_d)\) (and \(\mathbf x',\mathbf x'',\mathbf y',\mathbf y''\)) for \(d\)-tuples of formal variables. Over a commutative ring \(R\), a \(d\)-dimensional formal biextension germ \(B\) consists of a commutative formal group law \(\mu_B(\mathbf x,\mathbf y)\in R[[\mathbf x,\mathbf y]]^d\) with formal \(\mathcal O_E\)-action \(a\mapsto[a]_B(\mathbf x)\in R[[\mathbf x]]^d\) of linear part \(\rho_B(a)=D([a]_B)(0)\in\operatorname{Mat}_d(R)\), together with bi-rigidified partial-multiplication, symmetry and Rosati-adjointness series \(\varphi_1(\mathbf x',\mathbf x'',\mathbf y)\), \(\varphi_2(\mathbf x,\mathbf y',\mathbf y'')\), \(\mathrm{sy}(\mathbf x,\mathbf y)\) and \(\mathrm{ad}_B(a)(\mathbf x,\mathbf y)\) (\(a\in\mathcal O_E\)) satisfying the symmetric commutative biextension cocycle laws and
\[
 \begin{aligned}
 \varphi_1([a]_B\mathbf x',[a]_B\mathbf x'',\mathbf y)\,\mathrm{ad}_B(a)(\mu_B(\mathbf x',\mathbf x''),\mathbf y)
 &=\mathrm{ad}_B(a)(\mathbf x',\mathbf y)\,\mathrm{ad}_B(a)(\mathbf x'',\mathbf y)\,\varphi_1(\mathbf x',\mathbf x'',[\bar a]_B\mathbf y),\\
 \varphi_2([a]_B\mathbf x,\mathbf y',\mathbf y'')\,\mathrm{ad}_B(a)(\mathbf x,\mu_B(\mathbf y',\mathbf y''))
 &=\mathrm{ad}_B(a)(\mathbf x,\mathbf y')\,\mathrm{ad}_B(a)(\mathbf x,\mathbf y'')\,\varphi_2(\mathbf x,[\bar a]_B\mathbf y',[\bar a]_B\mathbf y'').
 \end{aligned}
\]
There exists a system \(S\) of \(d\)-dimensional formal biextension germs \(B_q\) over \(\mathcal O_{X,q}\) for all points \(q\in X\), obtained from local relative parameters of \(G\) along \(e\) and bi-rigidified local trivialisations of \(\Lambda(L_{\rm PP})\), equipped with compatible formal biextension isomorphisms \((\varphi,\psi,\lambda)\) (where \(\psi=\varphi^{-1}\)) along every specialisation \(q'\rightsquigarrow q\) in \(X\).
\item For any bi-rigidified series \(f(\mathbf x,\mathbf y)\in 1+(\mathbf x)(\mathbf y)R[[\mathbf x,\mathbf y]]\), write \(\operatorname{bil}(f)_{1,1}\in\operatorname{Mat}_d(R)\) for its bilinear \((1,1)\)-coefficient matrix in \(f(\mathbf x,\mathbf y)\equiv 1+\mathbf x^{\mathsf t}\operatorname{bil}(f)_{1,1}\mathbf y\pmod{(\mathbf x)^2+(\mathbf y)^2}\). For any formal biextension germ \(B\) over \(R\) coming by scalar extension from \(S\), setting
\[
 \rho_H(B)(a)(v,\eta)
 =\bigl(\rho_B(a)v,\;\rho_B(\bar a)^{\mathsf t}\eta-\operatorname{bil}(\mathrm{ad}_B(a))_{1,1}^{\mathsf t}v\bigr)
 \qquad(a\in\mathcal O_E,\ v,\eta\in R^d)
\]
defines an \(\mathcal O_E\)-module structure on \(R^d\times R^d\), and every formal biextension isomorphism \((\varphi,\psi,\lambda):B\xrightarrow{\sim}B'\) induces an \(\mathcal O_E\)-linear isomorphism
\[
 \operatorname{transH}(\varphi,\psi,\lambda)(v,\eta)
 =\bigl(D\varphi(0)v,\;D\psi(0)^{\mathsf t}(\eta+\operatorname{bil}(\lambda)_{1,1}^{\mathsf t}v)\bigr)
\]
from the module of \(B'\) to that of \(B\). Gluing the stalk modules of \(S\) inside \(V\times V^\vee\) defines a coherent \(\mathcal O_E\)-module \(\mathcal H\) on \(X\) sitting in a locally split \(\mathcal O_X\)-exact sequence
\[
 0\longrightarrow\mathcal L^\vee(\mathrm{conj})\longrightarrow\mathcal H\longrightarrow\mathcal L\longrightarrow0,
\]
where \(\mathcal O_E\) acts on \(\mathcal L^\vee(\mathrm{conj})\) through conjugation, and at every point \(q\in X\) the stalk \(\mathcal H_q\) is free of rank \(r=P+Q\) over \(\mathcal O_E\otimes_{\Z}\mathcal O_{X,q}\).
\end{enumerate}
\end{lemma}
\begin{proof}
For (1), let \(W\subseteq X\) be nonempty and open. Since \(G\to X\) is smooth with geometrically integral fibres and \(W\) is normal and integral with dense open \(W\cap U\), \(G_W^{\times_W n}\) is normal and integral with dense open \(G_{W\cap U}^{\times_{W\cap U}n}\). Over \(W\cap U\), \(G_{W\cap U}^{\times_{W\cap U}n}\to W\cap U\) is proper and flat of finite presentation with geometrically reduced and connected fibres, so \(\Gamma(W\cap U,\mathcal O_{W\cap U})\xrightarrow{\sim}\Gamma(G_{W\cap U}^{\times_{W\cap U}n},\mathcal O)\). For any \(f\in\Gamma(G_W^{\times_W n},\mathcal O)\), pulling back along the \(n\)-fold identity section \(e_n:W\to G_W^{\times_W n}\) gives \(g=e_n^*f\in\Gamma(W,\mathcal O_W)\) such that \(f-\pi_n^*g\) vanishes along \(e_n\), hence vanishes on \(G_{W\cap U}^{\times_{W\cap U}n}\), and therefore vanishes on the integral scheme \(G_W^{\times_W n}\). Since \(e_n^*\circ\pi_n^*=\mathrm{id}\), \(\pi_n^*\) is bijective. In particular, on every open \(W\subseteq X\) any bi-rigidified automorphism of a line bundle on \(G_W\times_W G_W\) equals \(1\), so \(\operatorname{ad}_a\) in (a) is unique.

For (2) and the gluing in (3), smoothness of \(G\to X\) of relative dimension \(d\) supplies, on an affine open neighbourhood \(\operatorname{Spec}R\subseteq X\) of any \(q\in X\), regular relative parameters \(\mathbf t=(t_1,\ldots,t_d)\) of \(G_R\) along \(e_R\) identifying \(\widehat{\mathcal O}_{G_R,e_R}\simeq R[[\mathbf t]]\), together with a local trivialisation \(\tau_0\) of \(\Lambda(L_{\rm PP})\) near \((e(q),e(q))\) which becomes bi-rigidified after multiplying \(\tau_0(g_1,g_2)\) by \(\tau_0(g_1,e)^{-1}\tau_0(e,g_2)^{-1}\tau_0(e,e)\). Expanding the group law, \(\mathcal O_E\)-action and Mumford partial-multiplication, symmetry and Rosati-adjointness isomorphisms in \((\mathbf t,\tau_0)\) produces a formal biextension germ \(B_q\) over \(\mathcal O_{X,q}\); the cocycle identities hold because the corresponding bi-rigidified line-bundle isomorphisms agree over \(U\) and hence globally by (1). Along a specialisation \(q'\rightsquigarrow q\), comparing two local choices \((\mathbf t,\tau_0)\) and \((\mathbf t',\tau'_0)\) expresses \(\mathbf t=\varphi(\mathbf t')\), \(\mathbf t'=\psi(\mathbf t)=\varphi^{-1}(\mathbf t)\) and \(\tau'_0(\mathbf x,\mathbf y)=\lambda(\mathbf x,\mathbf y)\,\tau_0(\varphi(\mathbf x),\varphi(\mathbf y))\), satisfying
\[
 \begin{aligned}
 \varphi'_{1}(\mathbf x',\mathbf x'',\mathbf y)\,\lambda(\mu_{B'}(\mathbf x',\mathbf x''),\mathbf y)
 &=\lambda(\mathbf x',\mathbf y)\,\lambda(\mathbf x'',\mathbf y)\,\varphi_1(\varphi(\mathbf x'),\varphi(\mathbf x''),\varphi(\mathbf y)),\\
 \varphi'_{2}(\mathbf x,\mathbf y',\mathbf y'')\,\lambda(\mathbf x,\mu_{B'}(\mathbf y',\mathbf y''))
 &=\lambda(\mathbf x,\mathbf y')\,\lambda(\mathbf x,\mathbf y'')\,\varphi_2(\varphi(\mathbf x),\varphi(\mathbf y'),\varphi(\mathbf y'')),\\
 \mathrm{ad}_{B'}(a)(\mathbf x,\mathbf y)\,\lambda(\mathbf x,[\bar a]_{B'}\mathbf y)
 &=\lambda([a]_{B'}\mathbf x,\mathbf y)\,\mathrm{ad}_B(a)(\varphi(\mathbf x),\varphi(\mathbf y)).
 \end{aligned}
\]
By uniqueness of bi-rigidified isomorphisms in (1), \(\operatorname{ad}_{a+b}\) is obtained from \(\operatorname{ad}_a\otimes\operatorname{ad}_b\) via \(\varphi_1,\varphi_2\) (whose series lie in \(1+(\mathbf x')(\mathbf x'')(\mathbf y)\) and \(1+(\mathbf x)(\mathbf y')(\mathbf y'')\) and contribute zero to the \((1,1)\)-jet) \(\operatorname{ad}_1=\mathrm{id}\), and by commutativity of \(\mathcal O_E\), \(\operatorname{ad}_{ab}=\operatorname{ad}_{ba}=((\mathrm{id}_G\times\bar a)^*\operatorname{ad}_b)\circ((b\times\mathrm{id}_G)^*\operatorname{ad}_a)\), whose \((1,1)\)-matrix is \(\rho_B(b)^{\mathsf t}\operatorname{bil}(\mathrm{ad}_B(a))_{1,1}+\operatorname{bil}(\mathrm{ad}_B(b))_{1,1}\rho_B(\bar a)\). Taking transposes shows that \(a\mapsto\rho_H(B_q)(a)\) is a ring homomorphism, and this persists under any scalar extension \(\mathcal O_{X,q}\to R\). Taking the \((1,1)\)-coefficient matrix of the third transition identity above and transposing gives \(\operatorname{transH}(\varphi,\psi,\lambda)\circ\rho_H(B')(a)=\rho_H(B)(a)\circ\operatorname{transH}(\varphi,\psi,\lambda)\). Since \(\operatorname{transH}(\varphi,\psi,\lambda)\) acts by \(D\varphi(0)\) on the quotient \(R^d\) and by \(D\psi(0)^{\mathsf t}=(D\varphi(0)^{-1})^{\mathsf t}\) on the submodule \(\{0\}\times R^d\), the stalk lattices glue inside \(V\times V^\vee\) to a coherent \(\mathcal O_E\)-module \(\mathcal H\) in a locally split \(\mathcal O_X\)-exact sequence \(0\to\mathcal L^\vee(\mathrm{conj})\to\mathcal H\to\mathcal L\to0\).

It remains to prove that \(\mathcal H_q\) is free of rank \(r\) over \(\mathcal O_E\otimes_{\Z}\mathcal O_{X,q}\) at every \(q\in X\). Because \(F'\subset\Omega\) is the coefficient field of a good cover, it contains the images of all \([E:\Q]\) embeddings of \(E\), so every \(\sigma:E\hookrightarrow\Omega\) factors through \(F'\), and the signature \((P,Q)\) gives \(\dim_\Omega V_\sigma+\dim_\Omega V_{\bar\sigma}=P+Q=r\), whence \(\dim_\Omega(V\times V^\vee)_\sigma=r\) for every \(\sigma\). If \(\operatorname{char}\kappa(q)=0\), then \(\varpi\in\mathcal O_{X,q}^\times\), so \(F'=O[1/\varpi]\subseteq\mathcal O_{X,q}\) splits \(\mathcal O_E\otimes_{\Z}\mathcal O_{X,q}=E\otimes_{\Q}\mathcal O_{X,q}\simeq\prod_{\sigma:E\hookrightarrow F'}\mathcal O_{X,q}\) into a finite product of copies of the local domain \(\mathcal O_{X,q}\), on each of which the finite projective summand of \(\mathcal H_q\) has generic rank \(r\) and is therefore free of rank \(r\).

Suppose now that \(\operatorname{char}\kappa(q)=\ell\), and choose a geometric point \(\bar x:\operatorname{Spec}\bar k\to X\) at \(q\) with \(\bar k\) algebraically closed and Chevalley extension \(0\to T\to G_{\bar x}\xrightarrow{q_{\bar x}}B\to0\) as in (d). Any group-scheme homomorphism from the torus \(T\simeq\mathbb G_{m,\bar k}^{d_T}\) to the abelian variety \(B\) is trivial, so \(T\) is \(\mathcal O_E\)-stable and \(\mathcal O_E\) descends to \(B\). Since \(L_{\rm PP}|_{G_{\bar x}}\simeq q_{\bar x}^*M\), we have \(\Lambda(L_{\rm PP})|_{G_{\bar x}\times G_{\bar x}}\simeq(q_{\bar x}\times q_{\bar x})^*\Lambda(M)\), where \(\Lambda(M)=(\mathrm{id}_B\times\lambda_B)^*\mathcal P\) is the Poincar\'e biextension of the principal polarisation \(\lambda_B\). Along the fibres of the \((T\times T)\)-torsor \(q_{\bar x}\times q_{\bar x}:G_{\bar x}\times G_{\bar x}\to B\times B\), any invertible function on \(T\times T\) is a scalar times a pair of characters of \(T\); connectedness of \(B\times B\) makes those characters constant, and the rigidifications along \(e_{\bar x}\times\mathrm{id}_{G_{\bar x}}\) and \(\mathrm{id}_{G_{\bar x}}\times e_{\bar x}\) force both characters to be trivial. By faithfully flat descent along \(q_{\bar x}\times q_{\bar x}\), \(\operatorname{ad}_a|_{G_{\bar x}\times G_{\bar x}}=(q_{\bar x}\times q_{\bar x})^*\operatorname{ad}_{B,a}\) for a unique bi-rigidified Rosati isomorphism \(\operatorname{ad}_{B,a}:(a_B\times\mathrm{id}_B)^*\Lambda(M)\xrightarrow{\sim}(\mathrm{id}_B\times\bar a_B)^*\Lambda(M)\). Choosing formal parameters on \(\widehat G_{\bar x}\) adapted to \(0\to T\to G_{\bar x}\xrightarrow{q_{\bar x}}B\to0\) and pulling back a bi-rigidified trivialisation of \(\Lambda(M)\) from \(\widehat B\times\widehat B\), the matrix \(\operatorname{bil}(\mathrm{ad}_{B_{\bar x}}(a))_{1,1}^{\mathsf t}\) vanishes on \(\operatorname{Lie}(T)\) and takes values in \(\operatorname{Lie}(B)^\vee\subseteq\operatorname{Lie}(G_{\bar x})^\vee\). Consequently \(W_1=\operatorname{Lie}(T)\times\{0\}\) and \(W_2=\operatorname{Lie}(G_{\bar x})\times\operatorname{Lie}(B)^\vee\) are \((\mathcal O_E\otimes_{\Z}\bar k)\)-submodules of \(\mathcal H_{\bar x}\) with \(W_1\simeq\operatorname{Lie}(T)\), \(\mathcal H_{\bar x}/W_2\simeq\operatorname{Lie}(T)^\vee(\mathrm{conj})\) and \(W_2/W_1\simeq\operatorname{Lie}(B)\times\operatorname{Lie}(B)^\vee\). Restricting \(\Lambda(M)\) to \(B\times\operatorname{Inf}^1_0(B)\) constructs the universal vector extension \(0\to\operatorname{Lie}(B)^\vee\to E(B)\to B\to0\) \cite[Ch.~I]{mazur-messing-1974}; representing a tangent vector in \(\operatorname{Lie}(E(B))\) over \(\bar k[\varepsilon]/(\varepsilon^2)\) by \((\mathbf x_B=\varepsilon v,\;1+\varepsilon\eta^{\mathsf t}\mathbf y_B\bmod(\mathbf y_B)^2)\) relative to the formal trivialisation of \(\Lambda(M)\), functorial transport by \(a\in\mathcal O_E\) pulls \(1+\varepsilon\eta^{\mathsf t}\mathbf y_B\) back along \(\bar a_B\) and multiplies by \(\mathrm{ad}_{B,a}(\varepsilon v,\mathbf y_B)^{-1}\equiv1-\varepsilon(\operatorname{bil}(\mathrm{ad}_{B,a})_{1,1}^{\mathsf t}v)^{\mathsf t}\mathbf y_B\), giving an \((\mathcal O_E\otimes_{\Z}\bar k)\)-linear isomorphism \(W_2/W_1\xrightarrow{\sim}\operatorname{Lie}(E(B))\).

Let \(W(\bar k)\) be the Witt ring of \(\bar k\). Decomposing \(\mathcal O_E\otimes_{\Z}\Z_\ell=\prod_{\mathfrak l\mid\ell}\mathcal O_{E,\mathfrak l}\), splitting the maximal unramified subextensions over \(W(\bar k)\) and adjoining uniformisers of Eisenstein polynomials decomposes \(\mathcal O_E\otimes_{\Z}W(\bar k)=\prod_{j=1}^m R_j\) into complete discrete valuation rings \(R_j=\hat e_j(\mathcal O_E\otimes_{\Z}W(\bar k))\) with uniformisers \(\pi_j\) satisfying \(\pi_j^{e_j}\in\ell R_j^\times\), reducing modulo \(\ell\) to \(\mathcal O_E\otimes_{\Z}\bar k=\prod_{j=1}^m A_j\) with \(A_j=R_j/\ell R_j\) and idempotents \(e_j=\hat e_j\bmod\ell\). Let \(D^\vee=\mathbb D(B[\ell^\infty])(W(\bar k))^\vee\) be the \(W(\bar k)\)-dual of the contravariant Dieudonn\'e crystal of \(B[\ell^\infty]\); it is finite free over \(W(\bar k)\) and \(D^\vee/\ell D^\vee\simeq\operatorname{Lie}(E(B))\) compatibly with \(\mathcal O_E\) \cite[Prop.~2.5.2]{berthelot-breen-messing-1982}. For each \(j\), \(\hat e_j D^\vee\) is finitely generated over \(R_j\) and \(W(\bar k)\)-torsion-free, hence \(\pi_j\)-torsion-free because \(\pi_j^{e_j}\in\ell R_j^\times\), so \(\hat e_j D^\vee\) is free over \(R_j\) of some finite rank \(r_{B,j}\) and \(e_j(W_2/W_1)\simeq A_j^{r_{B,j}}\). The cocharacter lattice \(X_*(T)\) is finitely generated and torsion-free over the Dedekind domain \(\mathcal O_E\), hence projective of some rank \(r_T\), so \(\operatorname{Lie}(T)\simeq X_*(T)\otimes_{\Z}\bar k\) and \(\operatorname{Lie}(T)^\vee(\mathrm{conj})\) are free of rank \(r_T\) over \(\mathcal O_E\otimes_{\Z}\bar k\). For each \(j\), freeness of \(e_j(W_2/W_1)\) and \(e_j(\mathcal H_{\bar x}/W_2)\) over \(A_j\) splits the filtration on \(e_j\mathcal H_{\bar x}\) and gives \(e_j\mathcal H_{\bar x}\simeq A_j^{n_j}\) with \(n_j=r_{B,j}+2r_T\). Let \(\mathcal O_{X,q}^{\rm sh}\) be the strict henselisation of \(\mathcal O_{X,q}\), a normal Henselian local domain with separably closed residue field \(k_s=\kappa(q)^{\rm sep}\hookrightarrow\bar k\) and fraction field \(\Omega^{\rm sh}\supset\Omega\). Since \(\mathcal O_E\otimes_{\Z}\F_\ell\) has finite (hence separable) residue fields, the idempotents \(e_1,\ldots,e_m\) already lie in \(\mathcal O_E\otimes_{\Z}k_s\) and lift to orthogonal idempotents \(\tilde e_1,\ldots,\tilde e_m\in\mathcal O_E\otimes_{\Z}\mathcal O_{X,q}^{\rm sh}\) summing to \(1\). Put \(d_j=\dim_{\bar k}A_j\ge1\). The direct summands \(\tilde e_j(\mathcal O_E\otimes_{\Z}\mathcal O_{X,q}^{\rm sh})\) and \(\tilde e_j\mathcal H_q^{\rm sh}\) are projective, hence free, over \(\mathcal O_{X,q}^{\rm sh}\), of closed-fibre ranks over \(k_s\) equal to their \(\bar k\)-dimensions \(d_j\) and \(n_j d_j\). Over \(\Omega^{\rm sh}\), \(\tilde e_j(\mathcal O_E\otimes_{\Z}\Omega^{\rm sh})\) is spanned by \(d_j\) embeddings \(E\hookrightarrow F'\subset\Omega\), each of eigenrank \(r\) on \(V\times V^\vee\), so \(\tilde e_j\mathcal H_q^{\rm sh}\) has generic rank \(rd_j\). Thus \(n_j d_j=rd_j\) and \(n_j=r\) for all \(j\), making \(\mathcal H_{\bar x}\) free of rank \(r\) over \(\mathcal O_E\otimes_{\Z}\bar k\). By faithful flatness of \(k_s\hookrightarrow\bar k\), the closed fibre \(\mathcal H_q^{\rm sh}\otimes_{\mathcal O_{X,q}^{\rm sh}}k_s\) is flat, and hence free of rank \(r\), over the finite product of local Artin rings \(\mathcal O_E\otimes_{\Z}k_s\). Lifting a basis of \(\mathcal H_q^{\rm sh}\otimes_{\mathcal O_{X,q}^{\rm sh}}k_s\) to \(\mathcal H_q^{\rm sh}\) gives an \((\mathcal O_E\otimes_{\Z}\mathcal O_{X,q}^{\rm sh})\)-linear map \((\mathcal O_E\otimes_{\Z}\mathcal O_{X,q}^{\rm sh})^r\to\mathcal H_q^{\rm sh}\) which is surjective by Nakayama's lemma, hence an isomorphism between finite free \(\mathcal O_{X,q}^{\rm sh}\)-modules of equal rank \([E:\Q]r\). Faithfully flat descent along \(\mathcal O_{X,q}\to\mathcal O_{X,q}^{\rm sh}\) makes \(\mathcal H_q\) finite projective of constant rank \(r\) over the Noetherian semilocal ring \(\mathcal O_E\otimes_{\Z}\mathcal O_{X,q}\), and hence free of rank \(r\).
\end{proof}

\begin{lemma}[Formal-biextension conjugate sequence and degree zero on special-fibre curves]\label{B8.24dd}
Keep the hypotheses and notation of Lemma~\ref{B8.24dc}. Let \(k'\) be any field extension of \(k_0=O/\varpi O\), and let \(\Gamma\subset X_{k'}\) be a complete integral curve with normalisation \(\nu:\widetilde\Gamma\to\Gamma\). Then every embedding \(\sigma:E\hookrightarrow\Omega\) maps \(\mathcal O_E\) into the image of \(O\) in \(\Omega\), and the invertible sheaf \(\mathscr D_\sigma=\det(\mathcal H\otimes_{\mathcal O_E,\sigma}\mathcal O_X)\) satisfies \(\deg_{\widetilde\Gamma}(\nu^*\mathscr D_\sigma)=0\).
\end{lemma}
\begin{proof}[Outline]
Glues formal-biextension maps from \([\ell]_B\) via \ref{B8.24dc} into a conjugate exact sequence on \(\widetilde\Gamma\). Comparing degrees with the Hodge sequence around each Frobenius cycle of idempotents forces \(\deg_{\widetilde\Gamma}(\nu^*\mathscr D_\sigma)=0\). The complete proof is in \S\ref{prf:B8.24dd}.
\end{proof}

\begin{lemma}[Finite-place order of the flat section]\label{B8.24d}
Let \(s\) be the horizontal generator of Lemma~\ref{B8.24b} on a component \(Z\) over \(F'\). Fix a finite place \(w\) of \(F'\), and normalise \(v_w\) to have value one on a uniformiser of \(F'_w\). There exists \(\rho_w\in\Q\) such that for every integral frame \(y\) on \(Z\) over a test trait of ramification index \(e\) over \(w\),
\[
                            v_y(s(y))=e\rho_w .
\]
Here \(v_y\) has value one on a trait uniformiser. The numbers \(\rho_w\) vanish outside a finite set. For the different-corrected line the finite coefficient at \(w\mid\ell\) is
\[
 \rho_w-Q\,v_w(\ell)
       \left(\frac{v_\ell(|\operatorname{disc}E|)}{[E:\Q]}
                          -{\bf1}_{\ell=d}\right).
\]
\end{lemma}
\begin{proof}[Outline]
Extends \(\mathcal L_\tau\) of \ref{B8.24b} to a normal model via \ref{B8.24dc}, showing \(\operatorname{div}_X(s)\) is vertical. Applies \ref{B8.24dd}, Stein factorization, and \ref{B8.24c} to obtain \(v_y(s(y))=e\rho_w\) and subtract the different defect. The complete proof is in \S\ref{prf:B8.24d}.
\end{proof}

\begin{theorem}[Geometric inputs for the covolume estimate]\label{B8.25}
Assume characteristic-zero resolution and principalisation, invariant Poincar\'e--Lelong, the arithmetic Kleiman inequality \eqref{eq:p8c3-kleiman} and the two Hilbert--Samuel bounds \eqref{eq:p8c3-hs-upper} for sections on the model and its normalisation, under the hypotheses of \ref{B8.25a}. Assume also the arithmetic volume bound of \ref{B8.27}, Fargues exactness, Zarhin's retraction and nef Riemann--Roch. Assume the two compactification statements of \ref{B8.25a}, with the unitary model over its coefficient field and the Cartier divisor specified by the chosen fan function. Assume semiabelian Lie-bundle seminegativity, point separation, the arithmetic degree formula, the assertions \(\mathcal H_P\) of Definition~\ref{B8.25l} and Grothendieck deformation. Then \ref{B8.24} holds.
\end{theorem}
\begin{proof}[Outline]
Builds the section injection, proves deep goodness via \ref{B8.25b}, \ref{B8.25c}, and flat degree from \ref{B8.24b}, \ref{B8.24d}. Combines \ref{B8.25e}, \ref{B8.26}, \ref{B8.28}, \ref{B8.29}, \ref{B8.25d}, \ref{B8.30}, \ref{B8.31}, \ref{B8.32}, and \ref{B8.33}. The complete proof is in \S\ref{prf:B8.25}.
\end{proof}

\begin{lemma}[Sections on the integral model and its normalisation]\label{B8.25h}
Let \(y=(y_{ki})\) be homogeneous coordinates in a field \(\Omega\), in finitely many nonempty blocks indexed by \(k\), with a nonzero coordinate in each block. Let \(P\) be the kernel of the homomorphism
\[
 \Z[X_{ki}]\longrightarrow\Omega[T_k:k],\qquad X_{ki}\longmapsto y_{ki}T_k .
\]
Suppose \(P\cap\Z=0\) and every element of \(\Omega\) is a quotient of evaluations at \(y\) of integral multihomogeneous forms of the same multidegree. Let \(X\) be the closed subscheme defined by \(P\) in the product of projective spaces over \(\Z\), let \(d=\dim X_{\Q}\), and let \(\nu:X^\nu\to X\) be its normalisation. For an integral multidegree \(a=(a_k)\), write \(L=\mathcal O_X(a)\).

For a choice \(c=(c_k)\) with all \(y_{k c_k}\ne0\), put
\[
 R_c=\Z[y_{ki}/y_{k c_k}:k,i],\qquad m_c=\prod_k y_{k c_k}^{a_k},
\]
and let \(\widetilde R_c\) be the integral closure of \(R_c\) in \(\Omega\). For \(r\ge0\) put
\[
 S_r=\{s\in\Omega:s/m_c^r\in R_c\ \text{for every }c\},\qquad
 \widetilde S_r=\{s\in\Omega:s/m_c^r\in\widetilde R_c\ \text{for every }c\}.
\]
These are the global sections of \(L^r\) on \(X\) and of \(\nu^*L^r\) on \(X^\nu\), respectively. Give \(L\) a smooth conjugation-invariant metric specified on the complex cone by a positive function \(h\) with
\(h((t_k z_{ki}))=\prod_k|t_k|^{a_k}h(z)\).
The point norm of a form \(g\) of multidegree \(ra\) is \(|g(z)|/h(z)^r\). If \(\widehat h^0\) denotes the logarithm of the number of sections of sup norm at most one, then
\[
 0\le\widehat h^0(\widetilde S_r)-\widehat h^0(S_r)=o(r^{d+1}).
\]
Both counts, divided by \(r^{d+1}/(d+1)!\), tend to the arithmetic volume of \(\overline L=(L,\|\cdot\|)\).

Suppose also that \(P_0\subseteq P\) and that the Levi form of \(\log h\) is nonnegative at points of the complex cone of \(P\) on all vectors tangent to the equations of \(P_0\). There is a projective birational model with smooth generic fibre on which the pulled-back metric is smooth and semipositive.
\end{lemma}
\begin{proof}[Outline]
Glue the localised chart rings and use the conductor to compare section counts with a fixed twist via Moriwaki's convergence theorem. Stein factorisation and the projection formula preserve sections on the resolution, and cone lifts give semipositivity. The complete proof is in \S\ref{prf:B8.25h}.
\end{proof}

\begin{lemma}[Detecting a divisor on a completed arithmetic stalk]\label{B8.25i}
Let \(x\) be a closed point of an integral normal scheme of finite type and flat over \(\Z\). Put \(A=\mathcal O_{X,x}\), \(R=\widehat A\), and let \(p\) be the residue characteristic. Fix a dense open \(U\). Test on local homomorphisms \(v:R\to\mathcal O_{\overline{\Q}_p}\) whose generic point maps into \(U\). Such tests exist. If \(a,s\in R\), \(a\ne0\), and \(|v(s)|\le |v(a)|\) on all these tests, then \(a\mid s\) in \(R\). A rational function regular on \(U\) and bounded by one on all tests belongs to \(A\).
\end{lemma}
\begin{proof}[Outline]
Prime avoidance and completion yield one-dimensional tests. If \(s\notin aR\), localise \(R[a/s]\) to get a test with generic point in the prescribed open and \(|v(a)|<|v(s)|\), then use divisibility and faithful flatness to show that bounded rational functions are regular. The complete proof is in \S\ref{prf:B8.25i}.
\end{proof}

\begin{lemma}[The Fourier coordinate at a boundary point]\label{B8.25j}
On a component of the unitary moduli space with its finite theta level, write \(N\) for the theta modulus and \(rN\) for the torsion level used to evaluate the sections. Here \(N=m^2\), with \(m\) an odd prime power at least three, \(r\ge2\), and \(rN\) is invertible on the base. The theta data include a primitive \((rN)^2\)-th root of unity \(\zeta\). Take the normal projective model on which the theta-null tuple and the nonzero finite Fourier tuples have a generating entry locally. Let \(x\) be a closed point, \(A=\mathcal O_{X,x}\), \(R=\widehat A\), and \(X_*\) the period lattice of the completed split degeneration. For \(\alpha\in X_*\), let \(s_\alpha\in R\) be the coefficient of its cubical period section in an integral frame of its line. Two such frames change \(s_\alpha\) by a unit.

For every \(c\in X_*/NX_*\) there are a representative \(\mu_c\) and a nonzero rational function \(e_c\), the quotient of a finite Fourier sum by a theta null, such that
\[
 e_c^N\in A,\qquad e_c^N R=s_{\mu_c}R,\qquad
                    s_\alpha\in s_{\mu_c}R\quad(\alpha\bmod NX_*=c).
\]
\end{lemma}
\begin{proof}[Outline]
The level pairing selects classes on tests from \ref{B8.25i}, and \ref{B8.25k}, theta translates and torsion orthogonality give sharp Fourier bounds. The generating quotient defines \(e_c\), whose \(N\)-th power is regular and divides the class coefficients, one up to a unit. The complete proof is in \S\ref{prf:B8.25j}.
\end{proof}

\begin{lemma}[Order of the chosen quotient on a degeneration trait]\label{B8.25k}
Use the model, closed point \(x\), and choices of Lemma~\ref{B8.25j}. Let a degeneration trait mean the spectrum of a complete discrete valuation ring \(V\) of fraction characteristic zero, with a generic point in the moduli open, the given polarisation and level structure, and the semiabelian degeneration of that point. Suppose its centre on the model specialises to \(x\) and its generic point lifts to the chosen theta open. Write \(X_V\) for its period lattice, \(b_V\) for its monodromy pairing, and normalise the order by \(\operatorname{ord}_V(\varpi)=1\). There is a surjection \(\pi_V:X_*\to X_V\). The choices in Lemma~\ref{B8.25j} satisfy
\[
 \operatorname{ord}_V(e_c^{\,2N})=b_V(\pi_V\mu_c,\pi_V\mu_c)
          \le b_V(\pi_V\alpha,\pi_V\alpha)
          \qquad(\alpha\bmod NX_*=c).
\]
On a trait centred at \(x\) the period map is an isomorphism. In the unitary setting the resulting maps between the period lattices are \(\mathcal O_E\)-linear for the action on characters.
\end{lemma}
\begin{proof}[Outline]
Restrict torus characters over an extension of ramification index one to get a period-lattice surjection. The biextension law and tame inertia compare orders, \ref{B8.25j} gives the order formula and inequality, and positivity forces injectivity. The complete proof is in \S\ref{prf:B8.25k}.
\end{proof}

\begin{lemma}[Finitely many cones for fixed minimum representatives]\label{B8.25m}
Let \(X\) be the lattice of a cusp, \(N\ge3\) the residue modulus, and \(\mu_c\in c\) a representative for each \(c\in X/NX\). Use the smooth polarised fan on the balanced forms on \(X\). Suppose an integral positive definite balanced form \(b_0\) takes its minimum on \(c\) at \(\mu_c\) for every \(c\). There is a finite list of cones in this fixed lattice whose union contains every integral positive semidefinite balanced form with the same minimum representatives.
\end{lemma}
\begin{proof}[Outline]
Minimal evaluation tuples bound the minimisers in each cone with positive definite generator sum. Fan finiteness and the spanning representatives give finitely many cones in the fixed lattice, and closure puts the limit of \((b_0+jB)/j\) in one of them. The complete proof is in \S\ref{prf:B8.25m}.
\end{proof}

\begin{lemma}[Generation of a boundary ideal from approximation on the base]\label{B8.25n}
Let \(\pi:Y\to X\) be a proper birational morphism of integral normal schemes flat and of finite type over \(\Z\). Let \(D\) be an effective Cartier divisor on \(Y\), and let \(U\subset Y\setminus |D|\) be a dense open. At a closed point \(x\in Y\) of residue characteristic \(p\), a test means a local map from \(\widehat{\mathcal O}_{Y,x}\) to \(\mathcal O_{\overline{\Q}_p}\) with generic point in \(U\). Write \(|a(y)|\) for the absolute value of a germ on a test \(y\).

For each closed \(x\) put \(z=\pi(x)\), \(A=\mathcal O_{X,z}\), and let \(\mathfrak m\) be its maximal ideal. Suppose that there is a test \(y_0\) at \(x\) with the following property. For each \(n\ge1\) there is \(F_n\in A\) such that
\[
 |F_n(y_0)|=|s_x(y_0)|,\qquad
 |F_n(y)|\le \max\bigl(|s_{x'}(y)|,\rho_y^n\bigr)
\]
for every test \(y\) at every closed \(x'\) over \(z\). Here \(s_{x'}\) is a local equation of \(D\), and \(\rho_y\) is the maximum of the absolute values on \(y\) of a finite generating set of \(\mathfrak m\). Then the ideal
\[
 J=\ker\bigl(\mathcal O_X\longrightarrow\pi_*\mathcal O_D\bigr)
\]
satisfies \(J\mathcal O_Y=\mathcal O_Y(-D)\).
\end{lemma}
\begin{proof}[Outline]
Take a limit of the approximants in the completed base ring and use \ref{B8.25i} to get divisibility along the fibre, with a unit quotient at the chosen point. Properness and flat base change put the limit in the completed ideal, and faithful flatness gives equality. The complete proof is in \S\ref{prf:B8.25n}.
\end{proof}

\begin{lemma}[The ideal norm on a degeneration trait]\label{B8.25o}
Let \(X_1\to X\) be a finite dominant morphism of integral normal Noetherian schemes with function fields \(L/K\), where \(L/K\) is separable of degree \(h\). For a coherent ideal \(I\) on \(X_1\) there is an ideal \(J\) on \(X\) with the following property. For a trait \(\operatorname{Spec}V\to X\) with \(V\) a complete discrete valuation ring, there is an extension to a complete discrete valuation ring \(V'\) and \(h\) lifts \(y_j:\operatorname{Spec}V'\to X_1\), counted with multiplicity, such that
\[
                  JV'=\prod_{j=1}^h y_j^*I\,\mathcal O_{\operatorname{Spec}V'}.
\]
Here an ideal pulled back to a trait denotes its image in the valuation ring.
\end{lemma}
\begin{proof}[Outline]
Generate the norm ideal by coefficients of products over the field embeddings, and extend the trait using the algebra generated by the conjugate rings. Over the extended valuation ring, coefficient ideals multiply, and a tuple of ideal generators gives the branch product. The complete proof is in \S\ref{prf:B8.25o}.
\end{proof}

\begin{lemma}[The geometric and arithmetic premises]\label{B8.25a}
The fifteen premises of \ref{B8.25} hold for the application with any \(O\) of \ref{B9.01}. The compactification and Hodge-bundle statements are quantified over their own moduli data before the choice of family and level.
\end{lemma}
\begin{proof}[Outline]
Use theta valuations and a normalised blow-up to identify the prescribed toroidal model. Combine resolution, nef Hilbert--Samuel on the model of \ref{B8.25h}, the volume bound \ref{B8.27}, Hodge and deformation comparisons, and the inputs of \ref{B8.25l}. The complete proof is in \S\ref{prf:B8.25a}.
\end{proof}

\begin{lemma}[One size exponent for every layer]\label{B8.25b}
Fix eligible \(F_m,C_H\). For the family size exponent \(C_t\) of Definition~\ref{B8.01k}, put \(A=22+|C_t|\) and \(B=A+2\). Beyond a threshold, at every admissible point and layer, with \(\kappa=\sum_{\tau,a}|\lambda_{\tau,a}|\),
\begin{equation}\label{eq:p8c3-size-main}
 1\le b\le E_p\le x^Bb,\qquad |M_{p,i}|\le x^Bb,\qquad
 \max(b,|M_{p,i}|)\le P_*\kappa .
\end{equation}
At every nonterminal layer a quotient weight \eqref{eq:p8c3-quotient} exists, and every such weight satisfies \eqref{eq:p8c3-quotient-size}.
\end{lemma}
\begin{proof}[Outline]
Host bounds and row sums give a common size exponent, and summing floored adjacent differences from the zero anchor constructs the quotient weight. Gap and spacing bounds control the cross-wall Vandermonde loss and give \(z\le1/2\). The complete proof is in \S\ref{prf:B8.25b}.
\end{proof}

\begin{lemma}[The field and level of a deep cover]\label{B8.25c}
Assume the ray-class-field assertion of \ref{B8.25a}. For each eligible \(F_m,C_H\) there is a threshold beyond which every admissible point and every \(i+1<m\) admit a good cover with \(M'=g_0g_1 15^e\), \(e=\lfloor3L_x/\log3\rfloor\), and coefficient field \(F'\) satisfying
\[
 \log[F':\Q]+\log\operatorname{rd}(F')\le391(frL_x)^2 .
\]
\end{lemma}
\begin{proof}[Outline]
Splitting at \(3\) and \(5\) supplies complementary level ideals, and the ray class field supplies roots of unity and Weber values for compatible frames over the compositum. The ray class sequence and conductor-discriminant formula bound the field size. The complete proof is in \S\ref{prf:B8.25c}.
\end{proof}

\begin{lemma}[The core and realised row metrics]\label{B8.25d}
Assume \(\mathcal H_P\) of \ref{B8.25}. For fixed eligible \(F_m,C_H\), eventually at every admissible point and layer the core Petersson norm has a Gram frame. Its cusp lattice is finitely generated and spans the cusp space over the core coefficient field. The realised norm has a Gram frame, its ambient integral lattice \(\Lambda_{\rm all}\) is finitely generated, and the good lattice \(\Lambda_{\rm good}\) of \eqref{eq:p8c3-bad-lattices} is nonzero. At each embedding of the realisation field, the realised norm on the image of the cusp space is \(c>0\) times the Petersson norm at the restricted embedding, with \(|\log c|\le C_{\rm meas}E_p\) for an absolute \(C_{\rm meas}\). There is a family exponent \(A_1\ge0\) with \(\mu_{\max}(\Lambda_{\rm all})\le x^{A_1}b\), where \(\mu_{\max}\) is the maximum sublattice slope in the squared-norm convention.
\end{lemma}
\begin{proof}[Outline]
Uniformisation, lattice regularity and row realisation give the two frames. Tamagawa measures and Weyl dimensions bound their norm ratio using \ref{B8.25b}, and rescaling and Galois descent transfer the Petersson slope bound. The complete proof is in \S\ref{prf:B8.25d}.
\end{proof}


\section{Arithmetic volume and component heights}

For signature \((p_{\rm sig},q_{\rm sig})\) put \(n=p_{\rm sig}+q_{\rm sig}\) and let \(\Psi\) be the CM type. Write \(\|v\|_1\) and \(\|v\|_\infty\) for the sum and maximum of the absolute weight entries, and \(\operatorname{spr}(v)\) for the sum of the row spreads. A gap \(g>0\) requires strict decrease in each row and \(|v_{\tau,a}-v_{\tau,b}|\ge g\) for \(a\ge p_{\rm sig}\) and \(a\ne b\). Define
\begin{equation}
 B(v,g)=4|\Psi|n^2(1+\|v\|_\infty/g),\quad
 \mathcal E_s(v,g)=(\log B(v,g)+s)\|v\|_1+q_{\rm sig}\operatorname{spr}(v),
 \qquad \ell(D)=\log\operatorname{rd}(E)+|L_E|.
 \tag{8.26a}\label{eq:p8c3-error}
\end{equation}
For a unitary moduli datum over a coefficient field \(F\), a valuation family on its algebra of rules assigns to each height-one prime ideal \(\lambda\) of \(\mathcal O_F\) a set of valuations over \(\lambda\). The record valuation family assigns to \(\lambda\) the valuations \(f\mapsto v(f(y))/e\) obtained by evaluating rules at algebraic integral frames \(y\) on complete discrete valuation traits over \(\lambda\) of ramification index \(e\), in the sense of Definition~\ref{B8.28g}. The height valuation family assigns to \(\lambda\), above a rational prime \(\ell\), the valuations (normalised to extend \(v_\lambda\)) obtained from Hodge--Tate frames over traits above \(\lambda\) that trivialise the Lie eigenspaces through the integral Hodge--Tate map of the \(\ell\)-divisible group over \(\mathcal O_C\) for an algebraically closed \(\ell\)-adically complete valued extension \(C\) of the trait fraction field. For \(B=B(v,g)\) put
\begin{equation}\label{eq:p8c3-mixed-places}
 S_B=\{\ell\text{ prime}:\ell\le B\text{ or }\ell\mid2|\operatorname{disc}E|M'\},\quad
 h_E(\ell)=\frac1{\ell-1}+\frac{40v_\ell(|\operatorname{disc}E|)}{[E:\Q]},
 \quad c_2(B)=\sum_{\ell\in S_B}(h_E(\ell)+\mathbf1_{\ell=2})\log\ell ,
\end{equation}
and form the patched valuation family using the height valuation family at primes \(\lambda\) of \(\mathcal O_F\) above \(S_B\) and the record valuation family elsewhere. On a component \(Z\) over \(F\), let \(n_j(t)\) count the degree-\(j\) cusp sections integral for this patched valuation family whose norm at every good embedding and frame on \(Z\) is at most \(e^{jt}\) in the host scalar scale \(\xi_\R(D)=|\phi_0(\xi_0)|\), where \(\xi_0\) is the host skew-Hermitian element and \(\phi_0:E\hookrightarrow\C\) is the distinguished host embedding. Let \(D_{\rm fl}=|\Psi|\binom n2\) be the flag dimension and \(\deg_Z(v)>0\) the top self-intersection degree on \(Z\) of the line of weight \(v\). Put
\[
 \operatorname{vol}_Z(t)=\limsup_{k\to\infty}
 \frac{D_{\rm fl}!\log n_{8k}(t)}{[F:\Q](8k)^{D_{\rm fl}+1}},
 \qquad
 \nu_Z(v;B)=\inf_{t\in\R}\bigl(\operatorname{vol}_Z(t)/\deg_Z(v)-t\bigr).
\]

\begin{lemma}[Horizontal curve heights and arithmetic self-intersection]\label{B8.26p}
Let \(X\) be the integral projective model cut out by birational multihomogeneous coordinates \(y\) of generic dimension \(D_A\ge0\) with horizontal prime ideal \(P\) as in Lemma~\ref{B8.25h}, equipped with an arithmetic intersection system \(II\) satisfying item~4 of Definition~\ref{B8.25l}. Let \(A\) be a smooth conjugation-invariant Hermitian line on \(X\) of generic degree \(d_A\) and top arithmetic self-intersection \(I=II(P;A^{D_A+1})\). For a horizontal arithmetic curve \(W\) on \(X\) of complex generic degree \(\deg_{\C}W>0\), write \(h_A(W)=II(W;A)/\deg_{\C}W\) for its normalised height. Assume characteristic-zero resolution, the arithmetic Kleiman inequality \eqref{eq:p8c3-kleiman} and Theorem~\ref{B8.27}.
\begin{enumerate}
\item Suppose \(A\) is nef on geometric fibres and its cone metric has nonnegative Levi form at complex cone points of \(P\) along the cone-tangent directions of an ideal \(P_0\subseteq P\). If \(h_A(W)\ge a\) for every horizontal arithmetic curve \(W\) on \(X\), then
\begin{equation}\label{eq:p8c3-constant-shift}
 I\ge (D_A+1)a d_A .
\end{equation}
\item Suppose \(A\) has nonnegative multidegree generated by integral coordinate monomials of point norm at most one and its cone metric has nonnegative Levi form at complex cone points of \(P\) along the cone-tangent directions of an ideal \(P_0\subseteq P\). If horizontal arithmetic curves \(W\) on \(X\) satisfying \(h_A(W)\le a\) form a Zariski-dense collection in \(X_{\Q}\), then
\[
 I\le (D_A+1)a d_A .
\]
\end{enumerate}
\end{lemma}
\begin{proof}[Outline]
For (1), shifts the metric by \(-a\), resolves via \ref{B8.25h}, and applies arithmetic Kleiman. For (2), assuming \(I>(D_A+1)ad_A\), shifts the metric and applies \ref{B8.27} for a small section whose restriction to a dense curve with \(h_A(W)\le a\) gives a contradiction. The complete proof is in \S\ref{prf:B8.26p}.
\end{proof}

\begin{lemma}[Intersection system, semiabelian extension, Hodge--Tate annihilation, Hopf duality and polarised abelian quotients]\label{B8.26q}
Assume the first eight premises of Theorem~\ref{B8.25}.
\begin{enumerate}
\item Characteristic-zero resolution, principalisation and conjugation-invariant Poincar\'e--Lelong yield unrestricted Poincar\'e--Lelong on horizontal multihomogeneous arithmetic varieties together with an arithmetic intersection system \(II\) satisfying item~4 of Definition~\ref{B8.25l} (in particular, symmetric in smooth Hermitian lines, multilinear under tensor products, invariant under isometry, compatible with dominant generically finite pullback and smooth metric shifts, and satisfying the rational-section divisor recursion).
\item Let \(S\) be a locally Noetherian normal scheme, \(U\subset S\) a dense open subscheme and \(G,G'\) semiabelian schemes over \(S\). Every group-scheme homomorphism \(f:G_U\to G'_U\) extends uniquely to a group-scheme homomorphism \(F:G\to G'\).
\item Let \(p\) be prime and \((C,v)\) an algebraically closed valued field of characteristic zero with multiplicative valuation \(v:C\to\R_{\ge0}\) satisfying \(v(p)<1\) and \(p\)-adically complete valuation ring \(\mathcal O_C=\{x\in C:v(x)\le1\}\). For every \(p\)-divisible group \(G\) over \(\mathcal O_C\), the linearised Hodge--Tate map \(\alpha_G:T_pG\otimes_{\Z_p}\mathcal O_C\to\omega_{G^\vee}\) has kernel equal to the annihilator under the Tate pairing of \(\ker(d\log:T_pG^\vee\otimes_{\Z_p}\mathcal O_C\to\omega_G)\), and every \(\kappa\in\mathcal O_C\) satisfying \(v(\kappa)\le v(4)\) for \(p=2\), or \(v(\kappa)^{p-1}\le v(p)\) for \(p\ne2\), satisfies \(\kappa\,\omega_{G^\vee}\subseteq\operatorname{im}(\alpha_G)\).
\item Every proper smooth \(\C\)-scheme decomposes into pairwise disjoint clopen subschemes of constant relative dimension on which compatible holomorphic local sections of a line bundle algebraize and glue uniquely. Over the valuation ring \(R\) of \((C,v)\), for every finite abelian group \(N\) and character isomorphism \(H\simeq\operatorname{Hom}(N,C^\times)\) with \(\#H=\#N\), the Fourier transform maps the \(R\)-dual lattice of any \(R\)-Hopf order \(\mathcal A\subset C^N\) onto a dual \(R\)-Hopf order \(\mathcal A^\vee\subset C^H\) satisfying \((\mathcal A^\vee)^\vee=\mathcal A\) and \((cc')^2=v(\#N)^{\#N}\) for determinant valuations \(c,c'\), while restriction to a subgroup \(K\subseteq N\) and pullback along \(N\to N/K\) yield subgroup and quotient \(R\)-Hopf orders.
\item Let \(R_0\) be a commutative ring, \(A\) an abelian scheme over \(S_0=\operatorname{Spec}R_0\) equipped with a relatively ample line bundle \(L\), and \(G=\operatorname{Spec}B\hookrightarrow A\) a finite locally free subgroup scheme over \(S_0\). There exist an abelian scheme \(A'\) over \(S_0\), a relatively ample line bundle \(L'\) on \(A'\) and a finite, locally free, surjective group-scheme homomorphism \(q:A\to A'\) such that for every \(R_0\)-algebra \(S'\) and every \(S'\)-point \(z:\operatorname{Spec}S'\to A\), one has \(q\circ z=1_{A'}\) if and only if \(z\) factors through \(G\hookrightarrow A\). When \(R_0=R\) is the valuation ring of \((C,v)\) (over which every finite flat subgroup scheme is finite free), \(A'\) embeds as a direct-summand group-scheme retract of a principally polarised abelian scheme over \(S_0\).
\end{enumerate}
\end{lemma}
\begin{proof}[Outline]
Proves (1) by a quadratic twist and divisor induction, (2) by DVR trait lifting and torsion density, and (3) from \ref{B8.25} and Fargues's theorem. Derives (4) from projective GAGA and Fourier duality, and (5) from finite flat quotients, norm bundles, and Zarhin's retraction. The complete proof is in \S\ref{prf:B8.26q}.
\end{proof}

\begin{proposition}[Existence of good root towers]\label{B8.26u}
Assume the first eight premises of \ref{B8.25}. Let \(E\) be an abelian CM host and let \(Y\) be a deep cover as in Proposition~\ref{B8.26}, of split index \(p_{\rm sig}=k+j-1\) and rank \(n=p_{\rm sig}+q_{\rm sig}\), where \(p_{\rm sig},q_{\rm sig}\ge1\). Then there is a good root tower \(\mathscr T\) over \(E\), with model parameter \(4\), in the sense of Definition~\ref{B8.26r}, of split index \(p_{\rm sig}\) and rank \(n\), whose component set \(\operatorname{Comp}\) is the set of algebraic components of \(Y\).
\end{proposition}
\begin{proof}[Outline]
Lemmas~\ref{B8.26p}, \ref{B8.26q}, and \ref{B8.24b} build split Hodge models and determinant generators, with Lemma~\ref{B8.26s} giving flag heights. Degree, wall, Frobenius, flat-height, and Stokes bounds verify Definition~\ref{B8.26r}. The complete proof is in \S\ref{prf:B8.26u}.
\end{proof}

\begin{proposition}[The deep-cover mean estimate]\label{B8.26}
Assume the first eight premises of \ref{B8.25}. There is \(C_{\rm deep}\ge0\) such that for every datum \(D\), every core theta chain \(T\) satisfying the chain laws and \(m^2\le k\), and every \(1\le j<m\), the following holds. Let \(Y\) be a deep cover good with parameters \((e,o)\), for any \(e\in\N\) and \(o\in\R\). Suppose at most fifteen distinct primes divide its level \(M'\). For each integral weight \(v\) with gap \(g>0\) there are real \(s_Z\) such that
\begin{equation}
 s_Z\le\nu_Z(v;B(v,g)),\qquad
 \frac1{\#\operatorname{Comp}(Y)}\sum_Zs_Z\ge-C_{\rm deep}\mathcal E_{\ell(D)}(v,g).
 \tag{8.26b}\label{eq:p8c3-deep}
\end{equation}
\end{proposition}
\begin{proof}[Outline]
Build a global PEL model from Theorem~\ref{B8.25} via Lemmas~\ref{B8.26p} and~\ref{B8.26q}, attaching its good root tower by Proposition~\ref{B8.26u}. Compare model slopes with \(\nu_Z(v;B)\) via a trace basis and Lemma~\ref{B8.26q}(3), bound \(c_2(B)\), and average. The complete proof is in \S\ref{prf:B8.26}.
\end{proof}

For the tower estimates write \(r=p_{\rm sig}\), \(s=q_{\rm sig}\), \(n=r+s\), \(f=|\Psi|\) and \(\lambda_N=\log(N\operatorname{rd}(E))+|L_E|\). Assume the first eight premises of Theorem~\ref{B8.25}, and assume \(s\ge1\) and \(3s+1\le n\), so \(n\ge4\) and \(s<r\). Relative to the arithmetic intersection system \(II\) of Lemma~\ref{B8.26q}(1), let \(\mathscr T\) be a good root tower of split index \(r\) over the abelian CM host in the sense of Definition~\ref{B8.26r} below, satisfying \(\#\mathcal B_{\mathscr T}\le4fn^2\), with \(d=f\binom n2\), \(e=frs\), \(L_v\), \(D\), \(h_{\rm fl}\), \(P_{N,Z}\), \(P_N\), \(V_Z\) and \(V\) as defined there.

\begin{definition}[Root models, good root towers and flag functionals]\label{B8.26r}
Fix a finite nonempty CM type \(\Psi\) of cardinality \(f=|\Psi|\) and integers \(1\le r<n\) with \(s=n-r\), \(d=f\binom n2\) and \(e=frs\). Let \(\R[X_{\tau,i}]\) (\(\tau\in\Psi\), \(0\le i<n\)) be the real polynomial ring in the formal Chern roots. For \(\sigma:\Psi\to S_n\) write \(X^{\operatorname{perm}(\sigma)}=\prod_{\tau\in\Psi}\prod_{i=0}^{n-1}X_{\tau,i}^{\sigma_\tau(i)}\), and define the linear top-degree flag functional
\[
 D(F)=\sum_{\sigma:\Psi\to S_n}\Bigl(\prod_{\tau\in\Psi}\operatorname{sgn}(\sigma_\tau)\Bigr)[X^{\operatorname{perm}(\sigma)}]F ,
\]
which vanishes on homogeneous polynomials of degree different from \(d\). For a real weight \(v=(v_{\tau,i})\) put \(L_v=-\sum_{\tau,i}v_{\tau,i}X_{\tau,i}\) and \(\delta_{\rm fl}(v)=D(L_v^d)\). Call \(v\) dominant if \(v_{\tau,a}\ge v_{\tau,b}\) whenever \(0\le a\le b<n\), and regular dominant if \(v_{\tau,a}>v_{\tau,b}\) whenever \(0\le a<b<n\). Set \(p_{\tau,k}=\sum_{i=0}^{n-1}X_{\tau,i}^k\), \(p_{\tau,k}^{<r}=\sum_{i<r}X_{\tau,i}^k\), and for \(F\in\R[X_{\tau,i}]\) define \(\operatorname{Sh}_r(\tau,F)=D(\sum_{i<r}\partial_{\tau,i}F)\) and \(\operatorname{Tot}(\tau,F)=D(\sum_{i=0}^{n-1}\partial_{\tau,i}F)\).

Embed \(\prod_{\tau\in\Psi}\operatorname{Fl}(\Z^n)\) as a horizontal multihomogeneous arithmetic variety of arithmetic dimension \(d+1\) in the product of Pl\"ucker projective spaces \(\prod_{\tau\in\Psi}\prod_{j=1}^{n-1}\mathbb P^{\binom{n}{n-j}-1}_{\Z}\). On the \(\tau\)-factor let \(\mathcal O^n=E_{\tau,0}\supset E_{\tau,1}\supset\cdots\supset E_{\tau,n}=0\) (\(\operatorname{rk}E_{\tau,j}=n-j\)) be the universal decreasing flag equipped with the subquotient metrics induced by the standard Hermitian inner product on \(\C^n\), and put \(\overline Q_{\tau,i}=(\overline{E_{\tau,i}/E_{\tau,i+1}})^\vee\), so that \(\bigotimes_{i=j}^{n-1}\overline Q_{\tau,i}\simeq(\overline{\det E_{\tau,j}})^\vee\) pulls back the Pl\"ucker Fubini--Study Hermitian line \(\overline{\mathcal O_{\operatorname{Gr}(n-j,n)}(1)}\) on \(\operatorname{Gr}(n-j,n)_{\Z}\). Relative to an arithmetic intersection system \(II\) satisfying item~4 of Definition~\ref{B8.25l}, define the linear arithmetic flag height functional \(h_{\rm fl}:\R[X_{\tau,i}]\to\R\) to vanish on homogeneous polynomials of degree different from \(d+1\) and to map each monomial \(\prod_{m=1}^{d+1}X_{\tau_m,i_m}\) of degree \(d+1\) to \(II\bigl(\prod_{\tau\in\Psi}\operatorname{Fl}(\Z^n);\overline Q_{\tau_1,i_1},\ldots,\overline Q_{\tau_{d+1},i_{d+1}}\bigr)\).

Relative to \(II\), a root model \(M\) for \((\Psi,n)\) consists of a finite block set \(K\) with finite coordinate sets \((\iota_k)_{k\in K}\), a function field \(\Omega\) with birational multihomogeneous coordinates \(y:\bigsqcup_{k\in K}\iota_k\to\Omega\) cutting out a horizontal prime ideal \(\mathfrak p_M\) of arithmetic dimension \(d+1\), multidegrees \(x_{\rm deg}(\tau,i)\in\Z^K\), smooth conjugation-invariant Hermitian root lines \(x(t,(\tau,i))\) (\(t\in\R\)) of multidegree \(x_{\rm deg}(\tau,i)\), an integer \(k_A>0\), a smooth conjugation-invariant auxiliary Hermitian line \(A\), and an integer degree factor \(\delta>0\) equal to the product of the coefficient-field degree and generic flag-cover degree. Evaluating \(II\) on \(\mathfrak p_M\) along the smoothed root lines perturbed by rational multiples of \(A\), dividing by \(\delta\) and passing to the zero-perturbation limit defines a linear limiting intersection functional \(P_M:\R[X_{\tau,i}]\to\R\).

A root tower \(\mathscr T\) of split index \(r\) consists of a finite nonempty component set \(\operatorname{Comp}\), a finite bad-prime set \(\mathcal B_{\mathscr T}\subset\N\) with \(2\in\mathcal B_{\mathscr T}\), a rational weight \(\nu\), and root models \(M(N,Z)\) for each cutoff \(N\in\N\) and \(Z\in\operatorname{Comp}\). Write \(P_{N,Z}=P_{M(N,Z)}\), \(P_N=\sum_{Z\in\operatorname{Comp}}P_{N,Z}\), \(V_Z\) for the geometric component volume on \(Z\), \(V=\sum_{Z\in\operatorname{Comp}}V_Z\), and when \(\delta_{\rm fl}(v)\ne0\) define the model slope
\[
 a_Z(v;N)=\frac{P_{N,Z}(L_v^{d+1})}{(d+1)V_Z\delta_{\rm fl}(v)}.
\]

For an abelian CM host \(E\) with \(\lambda_N=\log(N\operatorname{rd}(E))+|L_E|\), call \(\mathscr T\) a good root tower over \(E\), with model parameter \(4\), if it is equipped with a Frobenius path map \((\ell,\tau,h)\mapsto\operatorname{Fr}_{\ell,h}(\tau)\in\Psi\sqcup\Psi\) (\(\ell,h\in\N\), \(\tau\in\Psi\)) satisfying the following six laws.
\begin{enumerate}
\item Every \(V_Z>0\) takes a common value across \(Z\in\operatorname{Comp}\), and the geometric top intersection functional on \(M(N,Z)\) equals \(V_Z D\).
\item For every dominant weight \(v\) one has \(h_{\rm fl}(L_v^{d+1})\ge0\), with \(h_{\rm fl}(L_v^{d+1})=0\) whenever \(v_{\tau,a}=v_{\tau,a+1}\) for some \(\tau\in\Psi\) and \(0\le a<n-1\); moreover \(h_{\rm fl}(p_{\tau,1}Q)=0\) for all \(\tau\in\Psi\) and all \(Q\) homogeneous of degree \(d\).
\item There exist tower-dependent constants \(C\ge0\) and \(N_1\in\N\) with \(N_1\ge2\) and real sequences \((\alpha_M),(\beta_M)\) such that for every integer \(M\ge N_1\) and every monomial \(m\) of degree \(d+1\),
\[
 \Bigl|P_M(m)-V\Bigl(h_{\rm fl}(m)-(\alpha_M-\beta_M)\sum_{\tau\in\Psi}\operatorname{Sh}_r(\tau,m)-\beta_M\sum_{\tau\in\Psi}\operatorname{Tot}(\tau,m)\Bigr)\Bigr|\le\frac{C\log M}{M}.
\]
\item Whether \(\operatorname{Fr}_{\ell,h}(\tau)\) lies in the left or right copy of \(\Psi\) is independent of \(\tau\in\Psi\). Write \(P_{\ell,\tau,h+1,m}=\sum_{i<r}X_{\sigma,i}^m\) if \(\operatorname{Fr}_{\ell,h}(\tau)\) is the left copy of \(\sigma\in\Psi\), and \(P_{\ell,\tau,h+1,m}=\sum_{i\ge r}(-X_{\sigma,i})^m\) if it is the right copy of \(\sigma\). For all integers \(2\le N\le M\), \(Z\in\operatorname{Comp}\), \(\tau\in\Psi\), \(1\le k\le n\) and \(Q\) homogeneous of degree \(d+1-k\),
\[
 P_{N,Z}(p_{\tau,k}Q)=P_{M,Z}(p_{\tau,k}Q)+V_Z\sum_{\substack{N<\ell\le M\text{ prime}\\\ell\notin\mathcal B_{\mathscr T}}}\delta_\ell(\tau,k,Q),
\]
where
\[
 \delta_\ell(\tau,k,Q)=k\log\ell\left(\sum_{h\ge0}\ell^{-k(h+1)}D\bigl(P_{\ell,\tau,h+1,k-1}Q\bigr)-\frac{D(p_{\tau,k-1}^{<r}Q)}{\ell-1}\right).
\]
\item For \(N\ge2\), \(Z\in\operatorname{Comp}\), \(\tau\in\Psi\) and \(Q\) homogeneous of degree \(d\), \(P_{N,Z}(p_{\tau,1}Q)=h_{N,Z,\tau}V_Z D(Q)\) with \(|h_{N,Z,\tau}|\le4n\lambda_N\) whenever \(N\ge\max(128,\#\mathcal B_{\mathscr T})\).
\item When \(s<r\), putting \(\mu=\sum_{\tau,i\ge r}X_{\tau,i}\), \(J=\prod_{\tau}\bigl(\prod_{0\le a<b<r}(X_{\tau,b}-X_{\tau,a})\prod_{r\le a<b<n}(X_{\tau,b}-X_{\tau,a})\bigr)\) and \(F_*=\mu^{e+1}J\), for each integer \(N\ge\max(128,4fn^2,\#\mathcal B_{\mathscr T})\) and \(Z\in\operatorname{Comp}\), \(M(N,Z)\) admits a base Grassmannian PEL model of generic dimension \(e\) (as in Lemma~\ref{B8.25h}) with an auxiliary ample Hermitian line \(A\) of normalised horizontal curve height at least \(-K_A\) (\(K_A\ge0\)) such that: (i) its interior carries an abelian scheme of relative dimension \(fn\) with \(\mathcal O_E\)-action of signature \((r,s)\) over \(\Psi\) whose adapted determinant metric on \(\mu\) (given at good primes \(\ell\le N\) by the canonical-subgroup and Hodge--Tate frames of Lemma~\ref{B8.26q}(3)) bounds the labelled fibre degree at closed interior horizontal points above by \(fsX-ft\log\operatorname{rd}(E)+x+\tfrac12ft\log(4\pi)\) for some \(0\le t\le s\), where \(X=2\sum_{\ell\le N\text{ or }\ell\in\mathcal B_{\mathscr T}}(\log\ell)/(\ell-1)+2\log\operatorname{rd}(E)\) and the saturated Lie eigenlattice Arakelov degree sums \(x\) (of the rank-\(s\) eigenlattices over \(\overline\Psi\)) and \(y\) (of the rank-\(r\) eigenlattices over \(\Psi\)) satisfy \(|x-y|\le c_{\rm cm}fn\lambda_N\) for an absolute constant \(c_{\rm cm}\ge0\) independent of \(E,\mathscr T,N\) and \(Z\); (ii) complex uniformisation presents each complex sheet over a fundamental set \(\mathcal F\) of finite positive Haar measure for an arithmetic unitary group \(\Gamma\) acting on an \(\mathcal O_E\)-stable Hermitian lattice \(\Lambda\) of signature \((r,s)\) (equipped with an integral trace pairing of modulus at most \(\|v\|\|w\|\) in the positive majorant of first minimum \(\lambda_1\)) together with a nonnegative boundary cusp weight \(U\le fs\bigl(2\log^+(1/\lambda_1)+\log(2fs\operatorname{rd}(E))\bigr)\) such that adding \(\tfrac12 U\) dominates the boundary logarithmic singularity of the adapted determinant metric; and (iii) for all integers \(1\le a\le b\) and any real \(e_0\) lower-bounding the normalised height \(h_{\rm reg}\) of the regularised perturbed line on interior and boundary horizontal curves, \(M(N,Z)\) carries smoothed cleared-denominator bundles \(\overline M_{N,a,b,t'}\) (\(t'>0\)) of denominator \(b\) on \(\mu+(a/b)A\) which for sufficiently small \(t'>0\) are smooth, conjugation-invariant, cone-semipositive along the defining ideal of the base model and fibrewise nef of normalised horizontal curve height at least \(be_0\) on every horizontal arithmetic curve, and whose top arithmetic self-intersection \(I_{t'}\) and positive generic degree \(D_{t'}=b^e\deg_{\rm gen}(\mu+(a/b)A)\) satisfy that \(I_{t'}/((e+1)bD_{t'})\) has Stokes limit as \(t'\downarrow0\) and \(a/b\to0\) bounded above by \(P_{N,Z}(F_*)/((e+1)V_ZD(\mu^eJ))+\tfrac12\operatorname{mean}U\).
\end{enumerate}
\end{definition}

\begin{lemma}[Weyl degree formula, power-sum transition decomposition and flag height positivity]\label{B8.26s}
Adopt the notation of Definition~\ref{B8.26r} for a good root tower \(\mathscr T\), and write \(S(F)=\sum_{\tau\in\Psi}\operatorname{Sh}_r(\tau,F)\) and \(T(F)=\sum_{\tau\in\Psi}\operatorname{Tot}(\tau,F)\).
\begin{enumerate}
\item For every real weight \(v\),
\[
 \delta_{\rm fl}(v)=d!\prod_{\tau\in\Psi}\prod_{0\le a<b<n}\frac{v_{\tau,a}-v_{\tau,b}}{b-a},
\]
which is strictly positive if \(v\) is regular dominant and nonnegative if \(v\) is dominant.
\item Every homogeneous \(F\in\R[X_{\tau,i}]\) of degree \(d+1\) decomposes as \(F=\sum_{\tau\in\Psi}\sum_{k=1}^n p_{\tau,k}Q_{\tau,k}\) with \(Q_{\tau,k}\) homogeneous of degree \(d+1-k\). For \(k\ge1\) and \(\sigma,\tau\in\Psi\), \(\operatorname{Sh}_r(\sigma,p_{\tau,k}Q)=kD(p_{\tau,k-1}^{<r}Q)\) if \(\sigma=\tau\) and \(0\) otherwise, while \(\operatorname{Tot}(\sigma,p_{\tau,k}Q)=nD(Q)\) if \(\sigma=\tau\) and \(k=1\), and \(0\) otherwise. For each \(\tau_0\in\Psi\) there exist \(Q_1\) homogeneous of degree \(d\) with \(D(Q_1)=1\) and \(Q_2\) homogeneous of degree \(d-1\) with \(D(p_{\tau_0,1}^{<r}Q_2)=1\).
\item Let \(\operatorname{rk}_{\rm Lie}:\Psi\sqcup\Psi\to\R\) take the value \(r\) on the left copy and \(s\) on the right copy, fix \(\tau_0\in\Psi\), and put
\[
 \kappa_\ell=\frac{\log\ell}{n}\sum_{h\ge0}\ell^{-(h+1)}\operatorname{rk}_{\rm Lie}(\operatorname{Fr}_{\ell,h}(\tau_0)).
\]
Set \(\rho_\ell(\tau,1,Q)=0\) and, for \(k\ge2\),
\[
 \rho_\ell(\tau,k,Q)=k\log\ell\sum_{h\ge1}\ell^{-hk}D(P_{\ell,\tau,h,k-1}Q).
\]
Then for every prime \(\ell\ge2\), \(\tau\in\Psi\), \(k\ge1\) and \(Q\in\R[X_{\tau,i}]\),
\[
 \delta_\ell(\tau,k,Q)=-\frac{\log\ell}{\ell-1}S(p_{\tau,k}Q)+\kappa_\ell T(p_{\tau,k}Q)+\rho_\ell(\tau,k,Q),
\]
and \(|\rho_\ell(\tau,k,Q)|\le c(\log\ell)/\ell^2\) for some \(c\ge0\) independent of \(\ell\), so the sum over primes \(\ell>N\), \(\ell\notin\mathcal B_{\mathscr T}\) converges absolutely.
\item Assuming the arithmetic Kleiman inequality \eqref{eq:p8c3-kleiman} (and without assuming clause~2 of Definition~\ref{B8.26r}), the arithmetic flag height functional \(h_{\rm fl}\) defined on \(\prod_{\tau\in\Psi}\operatorname{Fl}(\Z^n)\) in Definition~\ref{B8.26r} satisfies clause~2 of Definition~\ref{B8.26r}.
\end{enumerate}
\end{lemma}
\begin{proof}[Outline]
Derives (1) by \(S_n\)-alternation, (2) from flag-variety cohomology and symmetric invariance, and (3) from \ref{B8.26r} and geometric tail bounds. Proves (4) by writing \(L_v\) in Pl\"ucker pullbacks, using arithmetic Kleiman for nonnegativity and partial-flag pullback on walls. The complete proof is in \S\ref{prf:B8.26s}.
\end{proof}

\begin{lemma}[The Hodge test class]\label{B8.26a}
Put \(\mu=\sum_{\tau,i\ge r}X_{\tau i}\) and let \(J\) be the product of \(X_{\tau b}-X_{\tau a}\) over \(a<b\) in each signature block. Set \(F_*=\mu^{e+1}J\) and \(H_*=(e+1)D(\mu^eJ)>0\). There is an absolute \(C_h\ge0\) such that for integers \(N\ge\max(128,4fn^2,\#\mathcal B_{\mathscr T})\),
\[
 P_N(F_*)\ge-VH_*(C_h+2)fs\lambda_N .
\]
\end{lemma}
\begin{proof}[Outline]
Regularize the Hodge metric (\ref{B8.26r}) by \(U/2\), using \ref{B8.26q}(3) and Bost's inequality for horizontal height lower bounds. Bounding \(\operatorname{mean}U\) by isotropic-vector counting and taking the Stokes limit in \ref{B8.26p}(1) yields the estimate. The complete proof is in \S\ref{prf:B8.26a}.
\end{proof}

\begin{lemma}[The two flag integrals for the Hodge test class]\label{B8.26d}
Use the test class \(F_*=\mu^{e+1}J\) and the normalisation \(H_*=(e+1)D(\mu^eJ)\) of Lemma~\ref{B8.26a}. Put \(a=rs\), so that \(e=fa\), where \(f\) is the number of embeddings. Write
\[
 F_*=\sum_{\tau,j}p_{\tau j}Q_{\tau j},
 \qquad p_{\tau j}=\sum_{i=0}^{n-1}X_{\tau i}^{j},
 \qquad 1\le j\le n,
\]
with homogeneous summands. Fix \(2\le k\le n\). At an output embedding \(\sigma\), let \(P\) be the power sum of degree \(k-1\) on the smaller signature block, with either overall sign. Then
\[
 \frac{k|D(PQ_{\tau k})|}{H_*}\le
 \begin{cases}
 \displaystyle\frac{2^{k-1}}{a+1}\binom{a+k}{k},&\sigma=\tau,\\[5pt]
 \displaystyle\frac{2^{k-1}}{k!}\frac{a!}{(a-k+1)!},&\sigma\ne\tau .
 \end{cases}
\]
\end{lemma}
\begin{proof}[Outline]
Use alternation and the Vandermonde to extract a single power-sum coefficient. The Frobenius formula restricts the Schur expansion to hooks at extraction, and tableau counts and multinomial ratios give the bounds for equal and distinct embeddings. The complete proof is in \S\ref{prf:B8.26d}.
\end{proof}

\begin{lemma}[The two cutoff coefficients]\label{B8.26b}
For \(N\) in the range of \ref{B8.26a}, put \(C_0=C_h+2\). Let \(F\) be homogeneous of degree \(d+1\) and choose a decomposition \(F=\sum_{\tau,k}p_{\tau k}Q_{\tau k}\), where \(p_{\tau k}=\sum_i X_{\tau i}^k\), \(1\le k\le n\), and \(Q_{\tau k}\) is homogeneous of degree \(d+1-k\). Set
\[
 S(F)=D\Bigl(\sum_{\tau,i<r}\partial_{\tau i}F\Bigr),\qquad
 T(F)=D\Bigl(\sum_{\tau,i}\partial_{\tau i}F\Bigr).
\]
Write \(h_{\rm fl}\) for the arithmetic flag height functional. The cutoff decomposition
\begin{equation}\label{eq:p8c3-cutoff}
 P_N(F)/V=h_{\rm fl}(F)-AS(F)-bT(F)+R_N(F)
\end{equation}
has coefficients
\[
 A\ge-K\lambda_N,\qquad |rA+nb|\le4n\lambda_N,\qquad K=2(C_0+4)+3 .
\]
\end{lemma}
\begin{proof}[Outline]
Combines \ref{B8.26r} and \ref{B8.26s} as \(M\to\infty\) on test classes to prove the decomposition and bound \(|rA+nb|\). Evaluates it on \(F_*\), using \ref{B8.26p}, \ref{B8.26d}, \ref{B8.26a} to bound \(h_{\rm fl}(F_*)\), \(R_N(F_*)\), \(P_N(F_*)\), and hence \(A\) below. The complete proof is in \S\ref{prf:B8.26b}.
\end{proof}

\begin{proposition}[The mean of the model slopes]\label{B8.26c}
For a real weight \(v\) with gap \(g>0\) and an integer \(N\ge\max(128,\#\mathcal B_{\mathscr T})\) with \(B(v,g)<N+1\), the model slopes satisfy
\begin{equation}
 \frac1{\#\operatorname{Comp}}\sum_Za_Z(v;N)
 \ge-C_T[\lambda_N\|v\|_1+s\operatorname{spr}(v)],
 \qquad C_T=2C_0+17 .
 \tag{8.26c}\label{eq:p8c3-model-mean}
\end{equation}
\end{proposition}
\begin{proof}[Outline]
Applies \ref{B8.26s} and \ref{B8.26b} to \(L_v^{d+1}\), bounding the main term via \(A\) and \(rA+nb\). Bounds \(R_N(L_v^{d+1})\) by Weyl alternation, perturbed Vandermonde reduction, and prime-tail summation, then averages over components using common volumes from \ref{B8.26r}. The complete proof is in \S\ref{prf:B8.26c}.
\end{proof}

\begin{theorem}[Arithmetic volume in the semipositive case]\label{B8.27}
Use the integral projective model \(X\), homogeneous ideal \(P\), and ordinary section lattice \(S_r\) of Lemma~\ref{B8.25h}. Write \(d=\dim X_{\Q}\). Let \(II\) satisfy the arithmetic intersection axioms and let \(A\) be the smooth conjugation-invariant Hermitian line with multidegree \(a\) and cone metric \(h\). Assume that the line is nef on the geometric fibres, including the generic fibre, and that the Levi form of \(\log h\) at the complex points of the cone of \(P\) is nonnegative on vectors tangent to \(P_0\), where \(P_0\subseteq P\). Then
\begin{equation}
 II(P;A^{d+1})\le\limsup_{r\to\infty}
 \frac{\log\#\{s\in S_r:\|s\|_{\sup}\le1\}}
 {r^{d+1}/(d+1)!}.
 \tag{8.27a}\label{eq:p8c3-moriwaki}
\end{equation}
\end{theorem}
\begin{proof}[Outline]
The modification of \ref{B8.25h} gives a smooth semipositive metric and preserves fibrewise nefness. Moriwaki's volume inequality bounds the intersection, which descends by the projection formula, and \ref{B8.25h} identifies the volume with the leading chart count. The complete proof is in \S\ref{prf:B8.27}.
\end{proof}

At a nonterminal layer the product of full flags has dimension \(d_{\rm fl}=f\binom r2\). Put \(T_*=\lceil P_*^2d_{\rm fl}^2 k/m\rceil_{\N}\) and \(\mathcal B=\binom{T_*-1}{d_{\rm fl}}\). Keep the anchored weights \(u_W,u_M\) and write \(\kappa=\|\lambda\|_1\). A quotient weight \(u_1\) has \(u_{1,\tau,k-2}=0\) and, for \(b'=a+1\),
\begin{equation}\label{eq:p8c3-quotient}
 u_{1,\tau,a}-u_{1,\tau,b'}
 =\left\lfloor\frac{u_{W,\tau,a}-u_{W,\tau,b'}
 +\begin{cases}-c_0-(2n-2),&b'=p_{\rm sig},\\2,&b'\ne p_{\rm sig}\end{cases}}
 {T_*+d_{\rm fl}+1}\right\rfloor .
\end{equation}
Write the remainders in these divisions as adjacent differences of the row anchored at zero in position \(k-2\). Let \(u_2\) be \(u_1\) plus that row. For \(u'=u_1-3n\eta_H\) and \(V(u)=\prod_{\tau,a<b'}(u_{\tau,a}-u_{\tau,b'})\), set
\begin{equation}
 \rho=1-V(u')/V(u_1),\qquad z=7(d_{\rm fl}^{\,2}/T_*+d_{\rm fl}\rho).
 \tag{8.28a}\label{eq:p8c3-z}
\end{equation}
By \ref{B8.25b},
\begin{equation}
 T_*\ge3d_{\rm fl}+5,\qquad 0\le\rho<1,\qquad 0<z\le\tfrac12 .
 \tag{8.28b}\label{eq:p8c3-quotient-size}
\end{equation}

\begin{lemma}[Height of a polynomial determinant]\label{B8.25p}
Let \(K\) be a number field. Let \(E\) be a Hermitian \(K\)-space of dimension \(N\ge1\), and let \(w_1,\ldots,w_n\) span \(E\). Suppose that \(\|w_i\|_\sigma\le\Lambda\) at every embedding \(\sigma:K\hookrightarrow\C\), where \(\Lambda\ge1\). For a nonzero multihomogeneous polynomial \(G\) on a product of copies of \(E\), define \(h_0(G)\) as follows. At a finite place take the logarithm of the maximum absolute coefficient after substituting \(\sum_i t_iw_i\) in each block, with separate variables for each block. At an embedding take the mean of \(\log|G|\) on the product of the Hermitian unit spheres with their invariant probability measures. Sum with the usual product-formula weights, normalised by \([K:\Q]\). Write \(e(G)\) for the sum of the block degrees. Then
\[
 h_0(G)\ge-e(G)\left(\log(n\Lambda)+\frac{1+\log n}{2}\right),
 \qquad h_0(G_1G_2)=h_0(G_1)+h_0(G_2).
\]

Let \(M\) be a full Hermitian lattice of rank \(v\) over \(\mathcal O_K\), and fix a \(K\)-basis of its ambient space. Let \(P_1,\dots,P_v\) be polynomial vectors on the same product of spaces. Suppose their coefficients after the finite-place substitutions lie in \(M\), and that \(\|P_j\|_\sigma\le L_j\) on every product of unit spheres, with \(L_j>0\). If \(D=\det(P_1,\ldots,P_v)\) in the chosen basis is nonzero and multihomogeneous, then the normalised Arakelov degree satisfies
\[
             \widehat{\deg}M\ \ge\ h_0(D)-\sum_{j=1}^v\log L_j .
\]
\end{lemma}

\begin{proof}[Outline]
Spherical submean, Cauchy's coefficient estimate and the product formula bound the height below, and finite coefficient norms and spherical means give additivity. Exterior products and Hadamard's inequality turn the determinant height into a lattice degree bound. The complete proof is in \S\ref{prf:B8.25p}.
\end{proof}

\begin{proposition}[Component degrees, ranks and section degrees]\label{B8.25e}
Under the geometric hypotheses of \ref{B8.25}, beyond a threshold for the fixed eligible family, take any admissible point, nonterminal layer, good cover \(Y\) and quotient \(u_1\). Write \(h_Z\) and \(\Delta_Z\) for the Petersson height and the degree of the image of \(Z\) in weight \(u_1\). Set \(\Delta'_Z=\deg_Z(u')=v_{0,Z}V(u')\). Then
\begin{equation}\label{eq:p8c3-shift-degree}
 0<\Delta'_Z\le\Delta_Z\le\Delta'_Z/(1-\rho).
\end{equation}
The \(\Delta'_Z\) have a common value \(\Delta_*>0\). Let \(r_Z\) and \(d_Z\) be the rank and Arakelov degree of the integral cusp lattice on \(Z\) in weight \(u_M\), with its Petersson metric in the squared-norm convention. Write \(\mu_Y(u_M)\) for the slope on the whole cover. Then
\begin{equation}
 (1-z)r_Z\le\mathcal B\Delta_Z\le r_Z.
 \tag{8.29a}\label{eq:p8c3-shares}
\end{equation}
For each cover and quotient weight there exist \(A>0\) and \(R\ge0\) such that for every component \(Z\),
\begin{equation}
 \begin{split}
 d_Z&\ge\mathcal B\Delta_Z
 \bigl(T_*h_Z-3(d_{\rm fl}+1)|h_Z|-A\log(2+|h_Z|)-R\bigr),\\
 A\log\max(A,1)+R&\le C_{\rm HS}E_p ,
 \end{split}
 \tag{8.29b}\label{eq:p8c3-hs}
\end{equation}
where \(C_{\rm HS}\) is uniform in the point, layer and cover. Moreover
\begin{equation}
 2\frac{\sum_Zd_Z}{\sum_Zr_Z}\le\mu_Y(u_M).
 \tag{8.29c}\label{eq:p8c3-degree}
\end{equation}
\end{proposition}
\begin{proof}[Outline]
Top Chern proportionality and moving degrees compare component degrees, while nef section counts and jets bound ranks. Factor the cut determinant by the Chow form, use \ref{B8.25p} and the size bound of \ref{B8.25b}, then compare the summed lattice degrees with the cover slope. The complete proof is in \S\ref{prf:B8.25e}.
\end{proof}

\begin{lemma}[Short generators of the integral lattice]\label{B8.25f}
Let \(K\) be a number field of degree \(a\), \(V\) a finite-dimensional \(K\)-space and \(L\subset V\) a finitely generated \(\mathcal O_K\)-module. For each embedding \(\sigma\) let \(N_\sigma\) be a nonnegative seminorm with \(N_\sigma(cx)=|\sigma c|N_\sigma(x)\). Suppose \(S\subset L\) spans \(V\) over \(K\) and \(N_\sigma(x)\le H\) for \(x\in S\) and all \(\sigma\). With \(v=\dim_K V\), there are at most \(2av\) generators of \(L\) over \(\Z\), each of seminorm at most \(\max(1,a^4v\operatorname{rd}(K)H)\) at every embedding.
\end{lemma}
\begin{proof}[Outline]
Apply transference to the inverse different to find short independent integers, then multiply them by a short \(K\)-basis to form a rational basis. Reduce integral basis coordinates modulo integers and bound the resulting generators by the seminorm triangle inequality. The complete proof is in \S\ref{prf:B8.25f}.
\end{proof}

\begin{lemma}[The short isotropic periods]\label{B8.25q}
Let \(E\) be a CM field of degree \(2f\) and root discriminant \(R_E\). Let \(V\) be a complex Hermitian space on which \(\mathcal O_E\) acts, with the adjoint of \(a\) acting as \(\bar a\). Write \(V_\sigma\) for its summand at the embedding \(\sigma\). Let \(L\subset V\) be a full \(\mathcal O_E\)-stable lattice, principally polarised by the imaginary part \(B\) of the Hermitian form. Let \(\lambda\) be the least length of a nonzero vector of \(L\). Set \(\Omega=2f^{3/2}R_E\), and take \(\theta>0\) with \(\Omega\theta^2\le1\).

There is a saturated \(\mathcal O_E\)-stable isotropic lattice \(I\subset L\). Put \(W=I\otimes_\Z\R\) and define \(u(x)\in W\) by \(\langle u(x),w\rangle_\R=B(x,w)\) for \(w\in W\). Then
\[
 \begin{split}
 x\in L\setminus I&\ \Longrightarrow\ {\rm dist}(x,W)\ge\theta,\\
 x\in L,\quad u(x)\ne0&\ \Longrightarrow\ \|u(x)\|\ge(\Omega\theta)^{-1}.
 \end{split}
\]
One may take \(I=0\) for \(\lambda\ge\theta\). For \(\lambda<\theta\), \(I\) contains every shortest vector. Its real rank is \(2f\rho\), with \(\rho\le\dim_\C V_\sigma\) for every \(\sigma\). If \(x_1,\ldots,x_t\in L\) and the \(u(ax_j)\), \(a\in\mathcal O_E\), span \(W\) over \(\R\), then, in this latter case,
\[
                     \frac{2f}{R_E^2\lambda^2}
                        \le\frac12\sum_{j=1}^t\|u(x_j)\|^2 .
\]
\end{lemma}

\begin{proof}[Outline]
Use short integral scalars from \ref{B8.25f} to adjoin nearby lattice vectors while preserving isotropy and the projection bound. Pair a shortest period with the spanning vectors and use the norm bound for the inverse different. The complete proof is in \S\ref{prf:B8.25q}.
\end{proof}

\begin{lemma}[Curvature and interpolation on the flag chart]\label{B8.25r}
Let \(0<p<r\), and take the product, over \(f\) embeddings, of the unitary flag domains of signature \((p,r-p)\). At an orthogonal frame \(x\) use lower unitriangular coordinates \(n=1+z\). Give the coordinates the Euclidean norm obtained by summing the squared moduli of all their entries. Let \(w\) be the metric potential of weight \(\lambda\). Assume that the adjacent gaps \(\lambda_i-\lambda_{i+1}\) within each signature block, and the gap \(\lambda_{r-1}-\lambda_0\), are at least \(T\) at every embedding. Indices start at zero.

Write \(d(z)\,dz\) for invariant measure in this chart, normalised by \(d(0)=1\), and \(R_x\) for minus the sum of the logarithms of the absolute Gram determinants of the first \(p\) columns. Put \(b_*=4fr^2+1\) and \(U=w-b_*R_x-\log d\). If \(T\ge b_*+2r\), \(U\) is plurisubharmonic and
\[
       \partial\bar\partial U(z)[v,\bar v]
       \ge \tfrac12(T-b_*-2r)\|v\|^2
       \qquad \bigl(\|z\|\le 1/(100r^2)\bigr).
\]
Here and below the Levi form is written as a Hermitian form.

Suppose the ball of this radius embeds in the arithmetic quotient. Put \(\rho=1/(200r^2)\) and \(D=fr(r-1)/2\). Choose a fixed smooth \(\chi:[0,\infty)\to[0,1]\) equal to one on \([0,1/4]\) and zero on \([1,\infty)\). For \(A\ge1\) bounding \(|\chi'|,|\chi''|\), suppose
\[
      T\ge b_*+2r+2\{D(13A+20)+A^2\}/\rho^2 .
\]
There is a holomorphic seed on the chart, of value one at its centre and vanishing at every other translate of the centre in the chart, for which
\[
 \int |F|^2e^{-w+b_*R_x}\,d\mu
 \le e^{C(\|\lambda\|_1+(fr)^8)}\,\mu(B(0,\rho)).
\]
The constant \(C\) depends only on the fixed cutoff.
\end{lemma}

\begin{proof}[Outline]
Compute the Levi form from Gram minors and transport it by reorthogonalising the frame. A weighted \(\bar\partial\) estimate on the pseudoconvex chart corrects a cutoff, with the singular limiting weight forcing the correction to vanish at orbit points. The complete proof is in \S\ref{prf:B8.25r}.
\end{proof}

\begin{lemma}[The kernel sum for the peak section]\label{B8.25s}
In the flag domain of Lemma~\ref{B8.25r} let \(K_x(y)=\exp(-b_*R_x(y))\), with \(R_x\) evaluated in a frame centred at \(x\). Write \(B_x\) for the chart ball of radius \(1/(100r^2)\). If its translates under the arithmetic group \(\Gamma\) are disjoint on the flag domain, there is an absolute \(C\) such that
\[
              \sup_y\sum_{\gamma\in\Gamma}K_x(\gamma y)
                              \le e^{C(fr)^8}.
\]
For a holomorphic seed \(F\) in the weight line with finite integral
\(\int\|F\|^2/K_x\,d\mu\), its periodisation with a unitary character has squared integral over a fundamental domain at most
\[
                         e^{C(fr)^8}\int\|F\|^2/K_x\,d\mu .
\]
\end{lemma}

\begin{proof}[Outline]
Use the chart bounds of \ref{B8.25r} to compare the kernel on disjoint orbit balls, then integrate using the matrix ball and compact flag measures. Weighted Cauchy--Schwarz and unfolding give the bound for the periodised seed. The complete proof is in \S\ref{prf:B8.25s}.
\end{proof}

\begin{lemma}[The zero Fourier coefficient of the theta derivative]\label{B8.25t}
Consider a nonterminal layer of signature \((P,Q)\). At the base level take an even number \(N\) of theta columns. At every CM embedding take a product of minors of the first \(j\) rows of the differential variables, with the minors of size \(j\) occurring to powers \(a_{\tau,j}-a_{\tau,j+1}\). Here the row degrees are decreasing and the degree following the last row is zero. Suppose the differential is harmonic and that \(a_{\tau,P}>0\) for some embedding in the CM type. Every theta derivative with this differential and finite level characteristics is cuspidal.
\end{lemma}
\begin{proof}[Outline]
In a symplectic basis adapted to a cusp, each zero-frequency term is constant in a common differential direction across all columns. The full-row minor annihilates these terms, and absolute convergence permits termwise averaging. The complete proof is in \S\ref{prf:B8.25t}.
\end{proof}

\begin{definition}[Polarised CM period tori and admissible Gaussian parameters]\label{B8.25ga}
Let \(E\) be a CM field of degree \(2f=[E:\Q]\) and root discriminant \(R_E=\operatorname{rd}(E)\). A polarised CM period torus \(\mathcal T=(V,L_0,H,\varepsilon)\) of integer rank \(r\ge1\) over \(E\) consists of the following data.
\begin{enumerate}
\item A multiplicity function \(\tau\mapsto d_\tau\in\N\) on \(\operatorname{Hom}(E,\C)\) satisfying \(d_\tau+d_{\bar\tau}=r\) for every \(\tau\), giving the complex space \(V=\bigoplus_\tau\C^{d_\tau}\) of dimension \(g_A=fr\) with its standard positive-definite Hermitian form \(H\) and coordinatewise \(\mathcal O_E\)-action \((\alpha\cdot v)_\tau=\tau(\alpha)v_\tau\).
\item A full \(\mathcal O_E\)-stable \(\Z\)-lattice \(L_0\subset V\) on which \(\Im H\) is integer-valued and unimodular, so that \(x\mapsto\Im H(x,\cdot)\) maps \(L_0\) isomorphically onto \(\operatorname{Hom}_{\Z}(L_0,\Z)\).
\item A quadratic refinement \(\varepsilon:L_0\to\{\pm1\}\) satisfying \(\varepsilon(x+y)=\varepsilon(x)\varepsilon(y)\cos(\pi\Im H(x,y))\) for all \(x,y\in L_0\), and \(\varepsilon(x)=1\) whenever \(\Im H(\alpha x,x)=0\) for all \(\alpha\in\mathcal O_E\).
\end{enumerate}
Write \(\lambda_1(\mathcal T)=\inf\{\|x\|:0\ne x\in L_0\}\) and, for \(\theta>0\), define the truncated phase \(\phi_\theta(\mathcal T)=2f/(R_E^2\lambda_1(\mathcal T)^2)\) if \(\lambda_1(\mathcal T)<\theta\) and \(\phi_\theta(\mathcal T)=0\) otherwise.

For integers \(r,q,m,R,N_c\ge1\) and real \(R_0>0\), put \(d_*=2fm\), \(\Omega_*=2f^{3/2}R_E\), \(M_1=90m^2d_*^2R_E\),
\[
 \theta(q,R_0)=\sqrt{\frac{fR_0^2R_E^{21}}{q}}>0,\qquad
 b_0(q,R_0)=\frac{q}{R_0^2R_E^{28}},\qquad
 t_0(q,R_0)=qb_0(q,R_0),
\]
and define the archimedean theta-constant bound \(\mathcal P(E,r,q,m,R,N_c)\) by setting \(p_1=1+fr\log(1+6\sqrt{frq})\), \(p_2=\log(2(3d_*)^{d_*})\), \(p_3=r(p_1+p_2+d_*)\), \(A=R^4rd_*/2\), \(K_*=4R_E^7d_*^2\) and
\[
 \mathcal P(E,r,q,m,R,N_c)=2R^4\bigl(p_3+\log N_c+2fr^2\log R\bigr)+A\log(AK_*).
\]
For integers \(D,q,m\ge1\) and real \(R_0>0\), call \((q,R_0,m)\) an admissible Gaussian parameter tuple at differential degree \(D\) on \(\mathcal T\) if \(fr\ge2\), \(\min_\tau d_\tau\le m\), \(m^4\le f\), \(f\le R_E\), \(f^2\le R_E\), \(R_E\ge16\), \(R_0^2\ge\max(270,90(fr+D))\) and
\[
 \begin{gathered}
 d_*^2\Omega_*\theta(q,R_0)^2<1,\qquad b_0(q,R_0)\ge1,\qquad M_1\le q/2,\qquad D\le q,\\
 4D(7+\log(qD))+d_*+1\le t_0(q,R_0),\\
 fr\log2+2d_*+3+(D+1)(r\log r+3r+3)\le t_0(q,R_0).
 \end{gathered}
\]
For integers \(r,D,q,m\ge1\) and real \(R_0>0\), call \((E,r,D,q,R_0,m)\) uniformly admissible if \((q,R_0,m)\) is an admissible Gaussian parameter tuple at differential degree \(D\) on every polarised CM period torus \(\mathcal T\) of rank \(r\) over \(E\) satisfying \(\min_\tau d_\tau\le m\).
\end{definition}

\begin{lemma}[Parameter admissibility and Appell--Humbert reduction of level-\(R\) theta rules]\label{B8.25gb}
Keep the notation of Definition~\ref{B8.25ga}.
\begin{enumerate}
\item Suppose integers \(D,q,m,r\ge1\) satisfy \(fr\ge2\), \(m^4\le f\), \(f^2\le R_E\), \(R_E\ge16\) and \(90R_E^{28}(fr+D)Dr\le q\). Put \(R_0=\sqrt{90(fr+D)}>0\). Then \((E,r,D,q,R_0,m)\) is uniformly admissible. If in addition real \(L\ge2\) satisfies \(m\le r\), \(\log R_E\le2L\), \(\log f\le2L\), \(\log r\le L\) and \(\log q\le70L\), then \(P_0=\mathcal P(E,r,q,m,5,1)\) satisfies \(0\le P_0\le200000(frL)^2\).
\item Let \((X_0,F_c)\) be a unitary moduli datum over \(E\) of signature \((p_{\rm sig},q_{\rm sig})\) with \(r=p_{\rm sig}+q_{\rm sig}\ge1\), level modulus \(M\), good complex embedding \(\sigma:F_c\to\C\) and analytic uniformisation \(U\). Let \(R>1\) be an odd integer with \(R^4\mid M\), let \((x_1,\ldots,x_r)\) span the level-\(R^4\) residue module over \(\mathcal O_E\), and let \(f_0\) be a single-term level-\(R\) theta rule for \((x_1,\ldots,x_r)\). Let \(D,q,m\ge1\) be integers and let \(R_0>0\). At every component \(k\) of \(U\) and domain frame \(y\) on \(k\), write \(\lambda_1(y)\) for the length of the shortest nonzero lattice vector in the positive majorant at \(y\), set \(\Phi=2f/(R_E^2\lambda_1(y)^2)\) if \(\lambda_1(y)<\theta(q,R_0)\) and \(\Phi=0\) otherwise, and write \(\|f_0^2(y)\|\) for the invariant point norm of \(f_0^2\) in section weight \((0,2r)\) (where the second entry of a section weight \((u,k)\) counts theta-column half-weights, i.e. twice the central/scalar parameter of Chapter~\ref{ch:B8-covolume}). Right-multiplication by an upper-triangular Borel matrix \(u_{\rm B}\) produces a domain frame \(x=yu_{\rm B}\) on \(k\) whose distinguished Gram minors all have modulus \(1\), satisfying \(\lambda_1(x)=\lambda_1(y)\) and \(\|f_0^2(y)\|=|f_0(x)|^2\). At \(x\), the complex abelian fibre of dimension \(fr\) carries a period map \(\operatorname{Pm}:\Lambda_k\to\C^{fr}\), a positive-definite Appell--Humbert Hermitian form \(H_m\) on \(\C^{fr}\) whose imaginary part on \(\operatorname{Pm}(\Lambda_k)\) is the integer-valued Riemann form, a \(\C\)-linear isometry \(\psi:(\C^{fr},H_m)\xrightarrow{\sim}(V,H)\) onto a polarised CM period torus \(\mathcal T=(V,L_0,H,\varepsilon)\) of rank \(r\) over \(E\) with \(L_0=\psi(\operatorname{Pm}(\Lambda_k))\), \(\min_\tau d_\tau=\min(p_{\rm sig},q_{\rm sig})\) and \(\phi_{\theta(q,R_0)}(\mathcal T)=\Phi\), an entire principal theta function \(\theta_0:\C^{fr}\to\C\) of unit \(L^2\) norm against \(e^{-\pi\Re H_m(z,z)}(\operatorname{vol}\mathcal F_0)^{-1}\,dz\) on a fundamental domain \(\mathcal F_0\) for \(\operatorname{Pm}(\Lambda_k)\), and vectors \(z_1,\ldots,z_r\in\C^{fr}\) such that \(R^2\psi(z_1),\ldots,R^2\psi(z_r)\) lie in \(L_0\) and span \(L_0/RL_0\) over \(\mathcal O_E\), and
\[
 |f_0(x)|=\prod_{j=1}^r|\theta_0(z_j)|\exp\bigl(-\tfrac\pi2\Re H_m(z_j,z_j)\bigr).
\]
If \((E,r,D,q,R_0,m)\) is uniformly admissible and \(\min(p_{\rm sig},q_{\rm sig})\le m\), then
\[
 \|f_0^2(y)\|\le\exp\bigl((\mathcal P(E,r,q,m,R,1)-\Phi)/R^4\bigr).
\]
\end{enumerate}
\end{lemma}
\begin{proof}[Outline]
Proves (1) by elementary parameter inequalities. For (2), Gram--Schmidt reduces to a unit-minor frame on a polarised CM period torus, bounding theta null-values directly for \(\Phi=0\) or via Fourier expansion along the isotropic sublattice of \ref{B8.25q} for \(\Phi>0\). The complete proof is in \S\ref{prf:B8.25gb}.
\end{proof}

\begin{lemma}[Component Bergman lower bounds from supremum ratios and fundamental-set peaks]\label{B8.25gc}
Let \((X_0,F_c)\) be a unitary moduli datum of rank \(r\) at a good complex embedding \(\sigma:F_c\to\C\) with analytic uniformisation \(U\), adelic quotient \(Q_W\) with probability measure \(\mu_{Q,1}\) and frame map \(\pi_\sigma\), algebraic component \(Z\), class set \(\mathcal C_Z=\{q\in Q_W:\pi_\sigma(q)\text{ lies over }Z\}\) with conditioned measure \(\mu_Z=\mu_{Q,1}(\mathcal C_Z)^{-1}\mu_{Q,1}|_{\mathcal C_Z}\) (zero if \(\mu_{Q,1}(\mathcal C_Z)=0\)), and domain frame \(y=(k_0,y_0)\) over \(Z\). For a section weight with associated model weight \(\lambda=(\lambda_k)_k\), write \(\|P(z)\|_\lambda\) for the invariant point norm and \(\|P\|_{2,Z}=\bigl(\int_{Q_W}\|P(\pi_\sigma(q))\|_\lambda^2\,d\mu_Z(q)\bigr)^{1/2}\), and let \(G(y)\in[0,\infty]\) be the supremum of \(\|P(y)\|_\lambda^2/\|P\|_{2,Z}^2\) over holomorphic cuspidal families \(P=(P_k)_k\) of weight \(\lambda\), with the convention \(c/0=+\infty\) for \(c>0\) and \(0/0=0\).
\begin{enumerate}
\item Suppose \(\pi_\sigma(q)\) lies in the signature domain for every \(q\in Q_W\). Let \(P\) be a holomorphic cuspidal family of weight \(\lambda\) with \(P_{k_0}(y_0)\ne0\), and let \(e>0\). If \(e\,\|P(z)\|_\lambda\le\|P(y)\|_\lambda\) for every domain frame \(z\) on \(U\), then \(G(y)\ge e^2\).
\item Write \(\operatorname{comp}(z)=k\) for the analytic component of a frame \(z=(k,z_0)\). Assume the uniformisation satisfies the level-geometry laws, meaning that \(\pi_\sigma(Q_W)\) lies in the signature domains, all domain frames on a given analytic component \(k\) lie over the same algebraic component, extending a holomorphic cusp form of weight \(\lambda_{k_0}\) on the signature domain of \(k_0\) by \(0\) outside that domain and on \(k\ne k_0\) defines a holomorphic cuspidal family of weight \(\lambda\), each analytic component \(k\) admits an orthogonal frame \(x_{0,k}\) and measurable fundamental set \(\mathcal F_k\) of finite invariant chart measure \(\mu_F(\mathcal F_k)<\infty\), and there is a finite constant \(c_{\rm vol}>0\) such that \(\mu_{Q,1}(\{q\in Q_W:\operatorname{comp}(\pi_\sigma(q))=k_0\})=c_{\rm vol}\mu_F(\mathcal F_{k_0})\) and \(\int_{Q_W}\|P'(\pi_\sigma(q))\|_\lambda^2\,d\mu_{Q,1}(q)=c_{\rm vol}\int_{\mathcal F_{k_0}}\|P_0(x_{0,k_0}n(z))\|_{\lambda_{k_0}}^2\,d\mu_F(z)\) for every such zero-extended family \(P'\). Let \(\beta\ge0\). Suppose that for every orthogonal frame \(x_0\) on \(k_0\) and every measurable fundamental set \(\mathcal F\) in the lower-unitriangular chart \(z\mapsto x_0n(z)\) on \(k_0\) with invariant chart measure \(\mu_F\), there exists a holomorphic cusp form \(P_0\) of weight \(\lambda_{k_0}\) on \(k_0\) such that \(P_0(y_0)\ne0\) and
\[
 \beta\int_{\mathcal F}\|P_0(x_0n(z))\|_{\lambda_{k_0}}^2\,d\mu_F(z)
 \le\|P_0(y_0)\|_{\lambda_{k_0}}^2\,\mu_F(\mathcal F).
\]
Then \(G(y)\ge\beta\).
\end{enumerate}
\end{lemma}
\begin{proof}[Outline]
For (1), since \(\mu_Z\) has mass at most \(1\), \(\|P\|_{2,Z}\le\sup_z\|P(z)\|_\lambda\le e^{-1}\|P(y)\|_\lambda\). For (2), extends \(P_0\) by zero off \(k_0\) and unfolds \(\mu_{Q,1}\) into \(\mathcal F_k\) to bound \(\|P'\|_{2,Z}^2\) by the \(\mathcal F_{k_0}\)-average. The complete proof is in \S\ref{prf:B8.25gc}.
\end{proof}

\begin{lemma}[Uniform bound for cuspidal theta derivatives]\label{B8.25gca}
Fix a layer of the theta chain, with host field \(E\), Hermitian space \(W\) of rank \(r\), level modulus \(M\), coefficient field \(F_c\), analytic uniformisations and adelic frame maps. Assume the level-geometry laws: the moduli datum represents this layer with its distinguished host root and CM type; the auxiliary level factor is one; at every good complex embedding the uniformisation and adelic frame map give the level-\(M\) arithmetic frame quotients of \(W\); the distinguished Gram minors of adelic frames have absolute value one; the stabiliser characters of dominant theta weights with an even number of columns are trivial; and the Hermitian space on every analytic component is a model of \(W\).

Put
\[
 f=[E:\Q]/2,\qquad R_E=\operatorname{rd}(E),\qquad
 \xi_{\R}=|\varphi_0(\xi_0)|>0,
\]
where \(\varphi_0\) and \(\xi_0\) are the distinguished complex embedding and polarising scalar of the host. Fix a good embedding \(\sigma:F_c\hookrightarrow\C\).

Let \((a,N)\) be a theta weight: the row entries \(a_{\tau,i}\) are nonnegative integers, nonincreasing in each row, and \(N\) is even. Write
\[
 D=\sum_{\tau,i}a_{\tau,i}.
\]
Let \(P\in\operatorname{harm}(a,N)\), with the integral coefficient condition of Definition~\ref{B8.28g}(c). For a tuple \(\mathbf x\) of \(N\) lattice characteristics, let \(s=s_{\mathbf x,P}\) be the theta rule formed from \(P\) using the canonical theta connection. Assume that \(s\) belongs to the cuspidal section space of weight \((a,N)\) at every good complex embedding.

Let \(L,\kappa,\rho\in\R\) satisfy
\[
 \begin{gathered}
 L\ge1,\qquad r\le N,\qquad
 \log M\le30L,\qquad
 \log R_E\le3L,\qquad \log(fr)\le L,\\
 D\le\kappa,\qquad frN\le2\kappa,\qquad
 \log\kappa\le20L,\qquad
 \sum_{\nu\in\operatorname{supp}P}|\sigma(P_\nu)|\le\rho.
 \end{gathered}
\]
Then every domain frame \(z\) at \(\sigma\) satisfies
\[
 (2\pi\xi_{\R})^{-D/2}\|s(z)\|_{(a,N)}
 \le \rho\exp\!\bigl(9000(frL)^2\kappa\bigr).
\]
\end{lemma}
\begin{proof}[Outline]
When \(S=\sup_z\|s(z)\|_{(a,N)}>0\), the spread bound \(D+frN\le 3\kappa\) yields an orthonormal frame \(x\) with \(S\le 2|s(x)|\) and \(\lambda_1(x)^2\ge e^{-80L}\). The differentiated theta estimate at \(x\) and parameter bounds give the result. The complete proof is in \S\ref{prf:B8.25gca}.
\end{proof}

\begin{lemma}[Lower bound for a family of block theta derivatives]\label{B8.25gcb}
Retain the layer, full level-geometry laws, good embedding \(\sigma\), and notation \(f,r,R_E,\xi_{\R}\) of Lemma~\ref{B8.25gca}, writing \(q\) for the level modulus and \((p_{\rm sig},q_{\rm sig})\) for the signature. In particular, every uniformised algebraic frame is obtained by base change from a frame over a test algebra, and the polarising scalar \(\xi_{0,k}\) on each analytic component \(k\) satisfies \(|\tau(\xi_{0,k})|=\xi_{\R}\) for every \(\tau:E\hookrightarrow\C\). Let \((a,N)\) be a dominant theta weight with \(N\) even. Let \(P\in\operatorname{harm}(a,N)\), with the integral coefficient condition of Definition~\ref{B8.28g}(c), be exactly a product of powers of the prescribed block determinants, with pairwise disjoint sets of columns for distinct blocks, and assume
\[
 \deg P=D=\sum_{\tau,i}a_{\tau,i}.
\]
Suppose \(\alpha\in\mathcal O_E\) satisfies \(\alpha\bar\alpha=q\), the half-level ideal is \(\mathfrak b=(\bar\alpha)\), and \((E,r,D,q,R_0,m)\) is uniformly admissible in the sense of Definition~\ref{B8.25ga}, with \(q_{\rm sig}\le m\). Write \(\Lambda\) for the \(\mathcal O_E\)-lattice of the level datum. For each of the \(q^{frN}\) tuples \(u\in(\Lambda/\mathfrak b\Lambda)^N\), let \(s_u\) be the theta rule obtained from \(P\) and representatives of \(u\) using the canonical theta connection. Assume that every \(s_u\) belongs to the section space of weight \((a,N)\).

Set
\[
 \begin{gathered}
 d_*=2fm,\qquad K_6=6(2d_*)^{d_*/2},\qquad
 C_\chi=8100m^4d_*^4R_E^7,\\
 A_0=\frac12\left\{\log\pi+
 N\left(\log24+\frac{d_*}{2}\log(2d_*)\right)
 +(r+2)\log2\right\},\qquad
 c_0=\frac{\pi NC_\chi}{q}.
 \end{gathered}
\]
At a domain frame \(y\), let \(\lambda_1(y)>0\) be the first minimum, put \(\phi(y)=2f/(R_E^2\lambda_1(y)^2)\), and let \(\Phi_D(y)\) be the truncated phase with threshold \(\theta(q,R_0)\) of Definition~\ref{B8.25ga}. Then, with
\[
 v_u(y)=(2\pi\xi_{\R})^{-D/2}\|s_u(y)\|_{(a,N)},\qquad
 J_y=4^{-N}K_6^{-N}2^{-(r+2)D}
       \exp\{-2\pi NC_\chi\Phi_D(y)/q\},
\]
one has
\[
 J_y>0,\qquad
 q^D\sum_u v_u(y)^2\ge\pi^{-D}J_y.
\]
For all real \(A,c\) with \(A_0\le A\) and \(c_0\le c\), there is a tuple \(u\) such that
\[
 v_u(y)\ge
 \exp\{-AD-c\phi(y)-(D+frN)\log q\}>0.
\]
\end{lemma}
\begin{proof}[Outline]
An orthogonal frame yields a polarised \(\mathcal O_E\)-torus where Definition~\ref{B8.25ga} and a Gaussian torus estimate lower-bound the sum over translates. Matching translates to characteristics and taking a maximum gives both bounds. The complete proof is in \S\ref{prf:B8.25gcb}.
\end{proof}

\begin{lemma}[The level-five numerator and the Bergman bounds]\label{B8.25g}
Fix eligible \(F_m,C_H\) and assume the packages \(\mathcal H_P,\mathcal H_H\), deep goodness and flat-degree conditions of Theorem~\ref{B8.25} and Definition~\ref{def:p8c3-section-laws}, together with the size and good-cover conditions of Lemmas~\ref{B8.25b} and \ref{B8.25c}. There are an absolute constant \(C_{\rm peak}\ge0\) and a family threshold with the following property. At any admissible point beyond the threshold take a nonterminal layer of signature \((p_{\rm sig},q_{\rm sig})\), a good cover of level \(M\) and a quotient weight \(u_1\) of \eqref{eq:p8c3-quotient}. Write \(r=p_{\rm sig}+q_{\rm sig}\), \(f=[E:\Q]/2\), \(R_E=\operatorname{rd}(E)\), \(Q_0=g_0g_1\) and \(T=T_*\). There are \(P_0,B'_0\ge0\) with
\begin{equation}\label{eq:p8c3-berg-budget}
 TP_0\le200000E_p,\qquad TB'_0\le835C_{\rm peak}E_p .
\end{equation}
For every component \(Z\) of the cover, good embedding \(\sigma\) of the cover field and domain frame \(y\) over \(Z\) at \(\sigma\), let \(\lambda_1\) be the length of the shortest nonzero lattice vector in the positive majorant at \(y\), put \(\phi=2f/(R_E^2\lambda_1^2)\), and let \(G(y)\in[0,\infty]\) be the component Bergman supremum of Lemma~\ref{B8.25gc} for the section weight \((u_1,0)\). For every single-term level-five theta rule \(f_0\) from an \(r\)-tuple spanning the level-\(5^4\) residue module over \(\mathcal O_E\) (as in Lemma~\ref{B8.25gb}(2) with \(R=5\)), let \(a=\|f_0^2(y)\|\) be the invariant point norm of \(f_0^2\) at \(y\). Everywhere \(a\le\exp(P_0/625)\). For \(M\lambda_1^2<b\) one has \(a\le\exp((P_0-\phi)/625)\) and \(G(y)\ge\exp(-\phi/625)\). For \(M\lambda_1^2\ge b\) one has \(G(y)\ge\exp(-B'_0)\).
\end{lemma}
\begin{proof}[Outline]
Lemmas~\ref{B8.25b}, \ref{B8.25c}, and \ref{B8.25gb} give the budget and level-five bounds. For \(M\lambda_1^2<b\), Lemmas~\ref{B8.25gca}, \ref{B8.25gcb}, and \ref{B8.24b} feed Lemma~\ref{B8.25gc}; for \(M\lambda_1^2\ge b\), a seed Poincar\'e series does. The complete proof is in \S\ref{prf:B8.25g}.
\end{proof}

\begin{lemma}[Order of a product of level-five theta constants]\label{B8.28a}
Let \(V\) be a complete discrete valuation ring with fraction field of characteristic zero. Write \(v\) for its additive valuation, with \(v(\varpi)=1\), and extend \(v\) to an algebraic closure. Take an integral frame of the polarised semiabelian degeneration, with full level \(5^4\) on its generic fibre. Let \(X\) be its period lattice and let \(Q:X\to\Z_{\ge0}\) be the order of the cubical period. Its polarisation
\[
              B(x,y)=Q(x+y)-Q(x)-Q(y)
\]
is symmetric and satisfies \(B(ax,y)=B(x,\bar a y)\) for \(a\in\mathcal O_E\). On the real space of symmetric forms with this property, let \(\phi\) be the infimum of \(\ell(B)\) over the real linear functionals \(\ell\) which are integral on integral forms and strictly positive on every nonzero positive semidefinite form.

Let \(f_0\) be the product of the level-five theta constants for a tuple spanning the full level module over \(\mathcal O_E\). Then
\[
                             v(f_0)\ge \phi/625 ,
\]
with \(v(0)=+\infty\).
\end{lemma}
\begin{proof}[Outline]
Bound each theta valuation below by the minimum cubical period order on its label coset divided by \(625\). The minimizing periods span over \(E\) and define a positive integral functional whose value on the period form is at least \(\phi\). The complete proof is in \S\ref{prf:B8.28a}.
\end{proof}

\begin{lemma}[The prime sum in the reference-weight budget]\label{B8.28b}
Use the nonterminal layer and quotient weight of \ref{B8.28}. Put \(L=L_x\), \(T=T_*\), \(r=k+m-1\), and let \(S\) be the width in \ref{B8.25b}. Define
\[
 B_0=\lfloor 8b/T\rfloor,\quad J=2B_0+1,\quad
 U=B_0+1,\quad V=f\{k(S+2r)/T+2q_{\rm sig}B_0\},\quad A=fmU .
\]
Let \(D_E=|\operatorname{disc}E|\) and \(\delta_\ell=v_\ell(D_E)/(2f)\). For a prime \(\ell>J\) put \(e_\ell=0\). For primes \(\ell\le J\) use the first applicable line of the following formula, with \(h_\ell=\log J/\log\ell\):
\[
e_\ell=
\begin{cases}
 A(1+\delta_\ell/2)+V,&\ell\mid Q_S,\\
 (frU+V)/(\ell-1),&\ell\in\{q,d\},\\
 \{A(4+h_\ell/2)+V\}/(\ell-1)
      +f(m+1)r^2(h_\ell+1)+1,&m\ell<k,\\
 (6A+V)/(\ell-1)+1,&\text{otherwise}.
\end{cases}
\]
Let \(N\) be the even number of columns used for the theta weight in \ref{B8.25g}, \(M'\) the covering level, and \(\mathcal B\) the primes dividing \(2dD_EM'\). Set
\[
 a(\ell)=e_\ell+{\bf1}_{\ell\in\{g_0,g_1\}}
              +{\bf1}_{\ell\in\mathcal B}
                    \bigl(1+\tfrac N2fr\,\delta_\ell\bigr).
\]
These nonnegative weights are supported on
\(\{\ell\text{ prime}:\ell\le J\}\cup\{g_0,g_1\}\cup\mathcal B\).
There is a family constant \(C_R\) such that, beyond one threshold,
\[
                         T\sum_\ell a(\ell)\log\ell\le C_R E_p .
\]
\end{lemma}
\begin{proof}[Outline]
Apply elementary prime sums with the width bound \ref{B8.25b}. Bound the extra prime contributions using the level bounds and \ref{B8.01k}, and the different contribution using the theta column count in \ref{B8.25g}. The complete proof is in \S\ref{prf:B8.28b}.
\end{proof}

\begin{lemma}[A short period basis on a trait]\label{B8.28c}
On a split semiabelian uniformisation of a trait, let \(X\) be the period lattice, \(s=\operatorname{rank}_{\Z}X\), and \(Q_X\) the order of its cubical period, with valuation one on a uniformiser. Write \(B_X\) for the polarisation of \(Q_X\) and \(\phi\) for the phase of Lemma~\ref{B8.28a}. Put \(n_E=[E:\Q]\) and \(\kappa=\operatorname{rd}(E)^2/(2n_E)\). There is a basis \(b_1,\ldots,b_s\) of \(X\) such that
\[
 Q_X(b_i)\le\Phi,\qquad |B_X(b_i,z)|\le 2\sqrt{s}\,Q_X(z)
             \quad(z\in X),\qquad \Phi=\kappa s^3\phi .
\]
In the layer of signature \((P,Q)\) with \(P\ge Q\) one has \(s\le n_E Q\).
\end{lemma}
\begin{proof}[Outline]
Bounds \(s\) from the toric Lie algebra, then uses AM--GM on balanced forms with \ref{B8.28a} to lower-bound \(\lambda_1(X^\vee)\). Banaszczyk's transference inequality bounds \(\lambda_s(X)\), and lifting generators along a saturated flag gives the basis. The complete proof is in \S\ref{prf:B8.28c}.
\end{proof}

\begin{lemma}[Integral jets in the theta expansion]\label{B8.28d}
In the notation of Lemma~\ref{B8.28c}, let \(D\) be a nonnegative integer and let \(M\ge4D+4\) satisfy \(2D\sqrt{s}<M\). On any positive number of copies of the uniformisation, fix an integral frame of the theta line at the identity. After scalar extension, every integral jet of total order \(D\), multiplied by a nonzero scalar of valuation at most \(D^2\Phi/M\), is the jet of a level-\(M\) theta section in the following lattice. A section with Fourier coefficient \(c_z\) at the period index \(z\) belongs to the lattice if
\[
                         v(c_z)\ge Q_X(z)/M
\]
in integral frames on the abelian part. On a product, \(Q_X\) in this definition is the sum of the forms of the factors.
\end{lemma}
\begin{proof}[Outline]
Uses \ref{B8.28c} so degree-\(\le D\) sums \(\chi\) uniquely minimize \(Q_X\) mod \(MX\). Prescribes abelian \(D\)-jets at \(\chi\) and extends by period translation so germs span all \(D\)-jets by Nakayama, then scales by valuation \(\lfloor D^2\Phi\rfloor/M\). The complete proof is in \S\ref{prf:B8.28d}.
\end{proof}

\begin{lemma}[The reference depth outside a finite set for fixed data]\label{B8.28e}
Fix a nonterminal member, a good cover and \(u_1\) as in Proposition~\ref{B8.28}, beyond the threshold of Lemma~\ref{B8.25g}. There is a finite set \(\Sigma\) of primes, depending on these fixed data, such that for \(\ell\notin\Sigma\), every algebraic integral frame on a trait of phase \(\phi\) satisfies
\[
  \inf_{f\ \mathrm{integral\ cusp\ section\ of\ weight}\ u_1}
                 \frac{v(f(y))}{v(\ell)}
                  \ \le\ \frac{\phi}{2500v(\ell)} .
\]
\end{lemma}
\begin{proof}[Outline]
Builds cusp sections via \ref{B8.25g} and \ref{B8.25t}, excluding level, \(c_{\tau i}!\), and denominator primes via \ref{B8.24d} and \ref{B8.24c}. For \(\ell\notin\Sigma\), applies \ref{B8.28d} to a monomial jet and decomposes into theta translates via \ref{B8.25l}. The complete proof is in \S\ref{prf:B8.28e}.
\end{proof}

\begin{lemma}[The period estimate for either half-level translation]\label{B8.28f}
Use the base level \(g_0g_1\). Take \(g=g_0,\ g'=g_1\), or \(g=g_1,\ g'=g_0\). Let \(\alpha\in\mathcal O_E\) be the half-level factor with \(\alpha\bar\alpha=g\). With \(X,Q_X,s,\Phi\) as in Lemma~\ref{B8.28c}, every class \(y+\alpha X\) contains \(z_0\) such that
\[
 \frac{Q_X(z_0)}{gg'}\le\frac{s\Phi}{4g'},\qquad
 Q_X(z_0+g'\alpha z)-Q_X(z_0)\ge g'(g'-1)Q_X(\alpha z)
                  \quad(z\in X).
\]
In a theta expansion supported on \(z_0+g'\alpha X\) whose orders are at least \(Q_X(z)/(gg')\), a coefficient at \(z_0\) attaining this bound at the point gives an evaluation of exactly that order.
\end{lemma}
\begin{proof}[Outline]
Uses \(Q_X(\alpha z)=gQ_X(z)\) and randomized rounding in the short basis \(b_i\) to bound \(\min_{z_0\in y+\alpha X}Q_X(z_0)\) by \(gs\Phi/4\). Minimality yields \(|B_X(z_0,\alpha z)|\le Q_X(\alpha z)\), giving the quadratic growth bound and ultrametric order. The complete proof is in \S\ref{prf:B8.28f}.
\end{proof}

\begin{definition}[Block depths, local depths and reference depths]\label{B8.28g}
Fix a nonterminal layer of signature \((P,Q)\) at the base level \(M=g_0g_1\), with coefficient field \(F_c\). Write \(\bar\tau\) for the conjugate of an embedding \(\tau\) of \(E\), and put \(\mu_\tau=P\) for \(\tau\) in the CM type and \(\mu_\tau=Q\) otherwise. A block is a pair \((\tau,i)\) with \(1\le i\le\mu_\tau\). A theta weight \((a,N)\) consists of integers \(a_{\tau,i}\ge0\), nonincreasing in \(i\), and an even number \(N\) of columns. Its gaps are \(c_{\tau,i}=a_{\tau,i}-a_{\tau,i+1}\), with \(a_{\tau,\mu_\tau+1}=0\), and its degree is
\[
 D=\sum_{\tau,i}a_{\tau,i}=\sum_{\tau,i}i\,c_{\tau,i}.
\]

Let \(\lambda\) be a finite place of \(F_c\) above the prime \(\ell\). A trait over \(\lambda\) is a complete discrete valuation ring \(V\) whose fraction field has characteristic zero, with an embedding \(F_c\to\operatorname{Frac}V\) inducing \(\lambda\). Its valuation \(v\) has value one on a uniformiser, so \(v=e\,v_\lambda\) on \(F_c\), where \(e\) is the ramification index. Let \(y\) be an integral frame on \(V\) and \(\widehat G\) the formal group at the origin of its semiabelian scheme. For each \(\tau\) the frame supplies integral Lie vectors \(u_{\tau,1},\dots,u_{\tau,\mu_\tau}\). They form a basis of a saturated summand of the integral Lie algebra, and \(\mathcal O_E\) acts on that summand through \(\tau\). Let \(C\) be an algebraically closed complete valued field containing \(\operatorname{Frac}V\), with valuation \(v_\ell\) normalised by \(v_\ell(\ell)=1\) and with ring of integers \(\mathcal O_C\). Choose integral formal coordinates of \(\widehat G\), and let \(\exp\) be its exponential over \(C\). A power series in these coordinates is integral if its coefficients lie in \(\mathcal O_C\). Such fields and coordinates exist. The infima below run over all of them.

(a) Let \((\tau,i)\) be a block and \(k\ge0\). The block depth \(\nu_{\tau,i}(k;y)\) is the infimum of
\[
 v_\ell\Bigl(\det\bigl(\partial/\partial t_{ab}\bigr)_{a,b\le i}^{\,k}\;
  F\bigl(\exp(\textstyle\sum_{a}t_{a1}u_{\tau,a}),\dots,\exp(\sum_{a}t_{ai}u_{\tau,a})\bigr)\big|_{t=0}\Bigr),
\]
taken over integral power series \(F\) on \(i\) copies of \(\widehat G\). For exponents \(c=(c_{\tau,i})\) put \(\nu_c(y)=\sum_{\tau,i}\nu_{\tau,i}(c_{\tau,i};y)\).

(b) Let \(\varrho\ge0\). The exponents \(c\) have oscillation at most \(\varrho\) at \(\ell\) if every place \(\lambda\mid\ell\) of \(F_c\) has an integer \(A_\lambda\ge0\) such that \(A_\lambda\le\nu_c(y)\) for all integral frames over \(\lambda\) and \(\nu_c(y)\le A_\lambda+\varrho\) for all algebraic integral frames over \(\lambda\).

(c) The set \(\operatorname{harm}(a,N)\) consists of the polynomials \(P\) with coefficients in \(\mathcal O_{F_c}\), in variables \(X_{\tau,i,b}\) with \(1\le b\le N\), that have the following three properties.
\begin{itemize}
\item Each monomial has degree \(a_{\tau,i}\) in the variables of the row \((\tau,i)\).
\item The polynomial lies in the span of products of leading minors \(\det(X_{\tau,i',b_j})_{i',j\le i}\), as in Lemma~\ref{B8.25t}.
\item It is harmonic. This means that it is annihilated by \(\sum_b\partial^2/\partial X_{\tau,i,b}\partial X_{\bar\tau,i',b}\) for every \(\tau\) in the CM type and all \(i,i'\).
\end{itemize}
Replacing \(\mathcal O_{F_c}\) by a ring \(R\subset\operatorname{Frac}V\) defines \(\operatorname{harm}_R(a,N)\).

(d) Let \(P\) be a polynomial over \(\operatorname{Frac}V\) in the variables \(X_{\tau,i,b}\). Let \(P(D)\) be the operator obtained by substituting for \(X_{\tau,i,b}\) the invariant derivation of \(\widehat G^N\) along \(u_{\tau,i}\) in the \(b\)-th factor. Let \(h_L\) be the logarithm along \(\widehat G\) of the canonical trivialisation of the theta line relative to its integral frame. It is the restriction of the canonical gauge of the formal Poincar\'e biextension along \(x\mapsto(x,w_0x)\), where \(w_0\in\mathcal O_E\) satisfies \(w_0+\bar w_0=1\). For \(k\ge0\) the local depth \(\delta_k(P;y)\) is the infimum of
\[
   v_\ell\Bigl(\bigl[P(D)\bigl(F\,e^{k\sum_bh_L(x_b)}\bigr)\bigr](0)\Bigr),
\]
taken over integral power series \(F\) on \(\widehat G^N\), where \(x_b\) is the point of the \(b\)-th factor. Thus \(\delta_k(P;y)\) is the least normalised order at the origin of \(P(\nabla)s\). Here \(s\) runs over the integral germs of the \(k\)-th power of the theta line on \(N\) copies, and \(\nabla\) is the canonical connection that defines theta derivatives.

(e) A cusp section \(f\) over \(F_c\) is integral if \(v(f(y))\ge0\) for every algebraic integral frame \(y\) on every trait over every finite place. Let \(w\) be a weight and \(c\ge0\). The reference depth \(d^c_\ell(w)\) is the supremum of
\[
     \inf_f\frac{v(f(y))}{v(\ell)}-\frac{c\,\phi}{v(\ell)},
\]
where \(f\) runs over the integral cusp sections of weight \(w\). The supremum is taken over the places \(\lambda\mid\ell\), the traits over \(\lambda\), the algebraic integral frames \(y\) on them, and their phases \(\phi\) (Lemma~\ref{B8.28a}). The same definition applies on a cover with its own coefficient field. The number \(d_q(u_1)\) in the proof of Proposition~\ref{B8.28} is \(d^{1/2500}_q(u_1)\) on the cover.
\end{definition}

\begin{lemma}[Bounds for one block]\label{B8.28h}
Let \(y\) be an integral frame over a place above \(\ell\), let \((\tau,i)\) be a block and let \(k\ge0\). Then
\[
     0\le\nu_{\tau,i}(k;y)\le i\,v_\ell(k!)\le\frac{ik}{\ell-1}.
\]
In particular \(\nu_{\tau,i}(k;y)=0\) for \(k<\ell\).
\end{lemma}
\begin{proof}[Outline]
Rewrites \(\partial/\partial t_{ab}\) on \(\widetilde F\) as commuting invariant derivations preserving integral series, giving \(\nu_{\tau,i}(k;y)\ge0\). Evaluates on the diagonal monomial from a unit \(i\times i\) frame minor and applies Legendre's formula. The complete proof is in \S\ref{prf:B8.28h}.
\end{proof}

\begin{lemma}[Summation and rounding of block depths]\label{B8.28i}
Let \(\lambda\mid\ell\) and let \(c\) be exponents. Suppose every block \(f\) has real numbers \(A_f\le B_f\) such that
\begin{itemize}
\item \(A_f\le\nu_f(c_f;y)\) for all integral frames over \(\lambda\), and
\item \(\nu_f(c_f;y)\le B_f\) for all algebraic integral frames over \(\lambda\).
\end{itemize}
Put \(A_\lambda=\lfloor\max(0,\sum_fA_f)\rfloor\). Then \(A_\lambda\le\nu_c(y)\) for all integral frames over \(\lambda\), and
\[
 \nu_c(y)\le A_\lambda+\sum_f(B_f-A_f)+1
\]
for the algebraic ones. Moreover \(c\) has oscillation at most \(0\) at every prime \(\ell>\max_fc_f\). At every prime it has oscillation at most \(D/(\ell-1)\), where \(D=\sum_{\tau,i}i\,c_{\tau,i}\).
\end{lemma}
\begin{proof}[Outline]
Combining \(\nu_c(y)\ge 0\) from \ref{B8.28h} with \(\sum_f A_f\le\nu_c(y)\) gives the lower bound; \(\sum_f B_f\) and \(\max(0,\sum_f A_f)<A_\lambda+1\) give the algebraic upper bound. Setting \(A_\lambda=0\) and summing \ref{B8.28h} over blocks yields both oscillation bounds. The complete proof is in \S\ref{prf:B8.28i}.
\end{proof}

\begin{lemma}[Clearing denominators away from a place]\label{B8.28j}
Let \(f\) be a cusp section of weight \(w\) over \(F_c\). Let \(m\) be a finitely supported real function on the finite places of \(F_c\) such that \(v(f(y))/e\ge m(w')\) for every algebraic integral frame \(y\) on a trait over \(w'\) of ramification index \(e\). Let \(\lambda\) be a finite place and \(k\) an integer such that \(v(f(y))/e\ge k\) for the algebraic integral frames over \(\lambda\). Then there is \(\sigma\in F_c^\times\) with \(v_\lambda(\sigma)=-k\) such that \(\sigma f\) is integral.
\end{lemma}
\begin{proof}[Outline]
Sets \(n_{w'}=-\lfloor m'(w')\rfloor\) on the finite support \(S\) of \(m'\) and applies the Chinese remainder theorem modulo primes dividing \(L\) to build \(\sigma=\beta/L^j\) with \(v_{w'}(\sigma)=n_{w'}\) on \(S\) and \(\ge0\) outside, making \(\sigma f\) integral. The complete proof is in \S\ref{prf:B8.28j}.
\end{proof}

\begin{lemma}[Isotropic columns and a global differential]\label{B8.28k}
Let \(\lambda\) be a place of \(F_c\) above \(\ell\), and let \(V_\lambda\) be the completion of \(\mathcal O_{F_c}\) at \(\lambda\). An embedding \(\tau\) also denotes the map \(\mathcal O_E\to V_\lambda\) through the coefficient field. Let \((a,N)\) be a theta weight with \(N\ge\sum_{c_{\tau,i}\ne0}2i\).

(a) For every embedding \(\tau\) there is \(\alpha_\tau\in V_\lambda\) with the following property. For every \(n_0\ge0\) there are \(u,t\in\mathcal O_E\) with
\[
 \tau(u)\equiv1,\quad \tau(t)\equiv\alpha_\tau \pmod{\ell^{n_0}V_\lambda},\qquad u\bar u+t\bar t\in\ell^{n_0}\mathcal O_E .
\]

(b) Assign disjoint sets of \(2i\) columns \(b_1,\dots,b_{2i}\) to the blocks \((\tau,i)\) with \(c_{\tau,i}\ne0\). Let \(B_{\tau,i}\) be the \(i\times N\) matrix whose \(j\)-th row has entry \(1\) in column \(b_j\), entry \(\alpha_\tau\) in column \(b_{i+j}\), and zeros elsewhere. The column differential
\[
 P_{\rm col}=\prod_{c_{\tau,i}\ne0}\det\Bigl(\sum_bX_{\tau,i',b}(B_{\tau,i})_{j,b}\Bigr)_{i',j\le i}^{\,c_{\tau,i}}
\]
lies in \(\operatorname{harm}_{V_\lambda}(a,N)\). For every \(n_0\) there are \(x_{j,b}\in\mathcal O_E\) with the following properties.
\begin{itemize}
\item \(x_{j,b}=0\) outside the columns of the block.
\item \(\tau(x_{j,b})\equiv(B_{\tau,i})_{j,b}\) modulo \(\ell^{n_0}V_\lambda\).
\item \(\sum_bx_{j,b}\bar x_{j',b}\in\ell^{n_0}\mathcal O_E\) for all \(j,j'\).
\end{itemize}

(c) For every \(n_0\) there is \(P_1\in\operatorname{harm}(a,N)\) with \(P_1-P_{\rm col}\in\ell^{n_0}\operatorname{harm}_{V_\lambda}(a,N)\). Through the canonical map \(V_\lambda\to V\), the same identity holds on every trait \(V\) over \(\lambda\).
\end{lemma}
\begin{proof}[Outline]
Solves \(u\bar u+t\bar t\equiv0\) locally at \(\mathfrak q\mid\ell\) and glues by Chinese remainder for \(\alpha_\tau\). Checks \(P_{\rm col}\in\operatorname{harm}_{V_\lambda}(a,N)\) via Cauchy--Binet, then approximates its coordinates in a \(\Z\)-basis of \(W_\Z\). The complete proof is in \S\ref{prf:B8.28k}.
\end{proof}

\begin{lemma}[Reference depths on a good cover]\label{B8.28l}
Let \(Y\) be a good cover of the base moduli at a nonterminal layer, as in Lemma~\ref{B8.25c}, with coefficient field \(F'\). For every prime \(\ell\), every theta weight \(w=(a,N)\) and every \(c\ge0\), the reference depth \(d^c_\ell(w)\) computed on \(Y\) is at most \(d^c_\ell(w)\) computed on the base moduli.
\end{lemma}
\begin{proof}[Outline]
Forgets the cover level and restricts coefficients to map any algebraic integral frame \(y'\) on \(Y\) to a base frame \(y\) with the same valuation and phase. Pulling back base integral cusp sections preserves values and integrality, yielding the inequality on infima. The complete proof is in \S\ref{prf:B8.28l}.
\end{proof}

\begin{lemma}[Theta derivatives at a place prime to the level]\label{B8.28m}
Let \((a,N)\) be a theta weight of degree \(D\), let \(P\in\operatorname{harm}(a,N)\), and let \(\lambda\mid\ell\) be a place of \(F_c\) with \(\ell\nmid M\). For characteristics \(x=(x_1,\dots,x_N)\) in the level-\(M\) module, let \(\theta_{x,P}\) be the theta derivative with differential \(P\) and characteristics \(x\). Its value at an integral frame is
\[
 \theta_{x,P}(y)=M^{-D}\bigl[P(\nabla)(s_{x_1}\boxtimes\dots\boxtimes s_{x_N})\bigr](0),
\]
where \(s_{x_b}\) is the section of the \(M^2\)-th power of the theta line attached to \(x_b\). This section is the translate of the integral theta section by a point of order dividing \(M^2\) that lifts \(x_b\), normalised at the frame. Let \(y\) be an algebraic integral frame on a trait \(V\) over \(\lambda\).

(1) Suppose \(A\le\delta_k(P;y')\) for all \(k\) and all algebraic integral frames \(y'\) over \(\lambda\). Then \(v(\theta_{x,P}(y))\ge A\,v(\ell)\) for every \(x\).

(2) Suppose \(y\) has a split uniformisation, with \(s\) and \(\Phi\) as in Lemma~\ref{B8.28c}. If \(4D+4\le M\) and \(2D\sqrt s<M\), then
\[
 \inf_x\frac{v(\theta_{x,P}(y))}{v(\ell)}\le\delta_M(P;y)+\frac{D^2\Phi}{M\,v(\ell)} .
\]
\end{lemma}
\begin{proof}[Outline]
For (1), passes to a split uniformization and checks integrality of \(s_{x_b}\) via \ref{B8.28a} to bound by \(\delta_{M^2}\). For (2), lifts the \(D\)-jet of \(F\) via \ref{B8.28d}, expands in theta translates as in \ref{B8.28e}, and uses the ultrametric inequality. The complete proof is in \S\ref{prf:B8.28m}.
\end{proof}

\begin{proposition}[Reference depth from a differential of controlled local depth]\label{B8.28n}
Let \((a,N)\) be a theta weight of degree \(D\) with \(N\ge2r\) and \(a_{\tau,P}>0\) for some \(\tau\) in the CM type. Let \(c\ge0\) satisfy
\[
 4D+4\le g_0g_1,\qquad 2D\sqrt{2fm}<g_0g_1,\qquad \kappa(2fm)^3D^2\le c\,g_0g_1,
\]
with \(\kappa\) as in Lemma~\ref{B8.28c}. Let \(\ell\notin\{g_0,g_1\}\), let \(\lambda\mid\ell\) be a place of \(F_c\), and let \(\varrho\ge0\). Assume the following.
\begin{itemize}
\item[(i)] There are an integer \(A\ge0\) and a differential \(P_1\in\operatorname{harm}(a,N)\) with \(A\le\delta_k(P_1;y)\le A+\varrho\) for all \(k\) and all algebraic integral frames \(y\) over \(\lambda\).
\item[(ii)] For each characteristic tuple \(x\) there is a finitely supported function \(m_x\) on the finite places such that \(v(\theta_{x,P_1}(y))/e\ge m_x(w)\) at every algebraic integral frame over each finite place \(w\).
\end{itemize}
Then every algebraic integral frame \(y\) over \(\lambda\) with phase \(\phi\) satisfies
\[
 \inf_f\frac{v(f(y))}{v(\ell)}-\frac{c\,\phi}{v(\ell)}\le\varrho ,
\]
where \(f\) runs over the integral cusp sections of weight \((a,N)\). If (i) and (ii) hold at every place above \(\ell\), then \(d^c_\ell(a,N)\le\varrho\).
\end{proposition}
\begin{proof}[Outline]
Splits the uniformization, bounds \(s\) via \ref{B8.28c}, and applies \ref{B8.28m}(2) for an upper bound on \(\theta_{x,P_1}\). Checks cuspidality via \ref{B8.25t}, rescales \(\theta_{x,P_1}\) integrally by \ref{B8.28m}(1) and \ref{B8.28j}, and applies \ref{B8.28g}(e). The complete proof is in \S\ref{prf:B8.28n}.
\end{proof}

\begin{lemma}[Inversion in the ring of approximable coefficients]\label{B8.28ua}
Let \(R\) be a Dedekind domain, let \(V\) be a discrete valuation ring, and let \(\varphi:R\to V\) be a ring homomorphism. Let \(q\in R\) have nonzero, nonunit image in \(V\). Define
\[
 S=\{t\in V:\text{ for every }n\ge0\text{ there exist }a\in R,\ z\in V
       \text{ with }\varphi(a)=t+\varphi(q)^nz\}.
\]
Then \(S\) is a subring of \(V\). If \(u\in S\cap V^\times\), then \(u^{-1}\in S\). Consequently, for \(r\in\N\), if an \(r\times r\) matrix \(M\) has entries in \(S\) and \(\det M\in V^\times\), its inverse has entries in \(S\).
\end{lemma}
\begin{proof}[Outline]
The congruences make \(S\) a subring. For \(u\in S\cap V^\times\), approximate \(u\) by \(\varphi(a)\), use maximality of \(\varphi^{-1}(\mathfrak m_V)\) for \(ab+i=1\), invert \(u\) mod \(\varphi(q)^nV\) via a geometric sum, and use the adjugate for \(M^{-1}\). The complete proof is in \S\ref{prf:B8.28ua}.
\end{proof}

\begin{lemma}[Cancellation of a frame with integral pullback]\label{B8.28ub}
Let \(C\) be a valued field of characteristic zero, with valuation ring \(\mathcal O_C\). For \(h,r,s\in\N\), let \(t=(t_1,\ldots,t_h)\), \(Y=(Y_1,\ldots,Y_r)\) and \(X=(X_1,\ldots,X_s)\). Let
\[
 \mathbf e(t)\in C[[t]]^r,\qquad
 \mathbf c(t)\in C[[t]]^s,\qquad
 \psi(Y)\in\mathcal O_C[[Y]]^s,\qquad
 \rho(X)\in\mathcal O_C[[X]]^r
\]
have zero constant terms and satisfy
\[
 \psi(\mathbf e(t))=\mathbf c(t),\qquad
 \rho(\mathbf c(t))=\mathbf e(t).
\]
Let \(W\in C[[X]]\) have constant term \(1\), and suppose that \(U=W(\psi(Y))\) has coefficients in \(\mathcal O_C\). For every \(F\in\mathcal O_C[[Y]]\) and integer \(k\ge0\), the series
\[
 G(X)=(F U^{-k})(\rho(X))
\]
has coefficients in \(\mathcal O_C\) and satisfies
\[
 G(\mathbf c(t))W(\mathbf c(t))^k=F(\mathbf e(t)).
\]
The equality remains true after applying any polynomial in the partial derivatives with respect to \(t\) and evaluating at \(t=0\).
\end{lemma}
\begin{proof}[Outline]
Recursive comparison in \(UU^{-1}=1\) makes \(U^{-1}\) integral, and substituting \(\rho\) keeps \(G\) integral. The curve identities cancel \(U^{-k}\) against \(W^k\) on \(\mathbf c(t)\), and finite-jet linearity of polynomial derivatives at \(0\) finishes. The complete proof is in \S\ref{prf:B8.28ub}.
\end{proof}

\begin{lemma}[Prime idempotents and compatible approximants]\label{B8.28uc}
Let \(R\) be a Dedekind domain, let \(V\) be a discrete valuation ring, and let \(\varphi:R\to V\) be a ring homomorphism. Let \(q\in R\) have nonzero, nonunit image. There are elements \(\varepsilon_N\in R\), for \(N\ge0\), such that
\[
 \varphi(\varepsilon_N)\in1+\varphi(q)^NV,
 \qquad
 \varphi(a)\in\varphi(q)^MV,\ M\le N
 \ \Longrightarrow\ \varepsilon_Na\in q^MR.
\]
If \(t\in V\) and \(x_N\in R\) satisfy \(\varphi(x_N)-t\in\varphi(q)^NV\), then
\[
 \varepsilon_nx_n-\varepsilon_mx_m\in q^mR\quad(m\le n),
 \qquad
 \varphi(\varepsilon_Nx_N)-t\in\varphi(q)^NV.
\]
\end{lemma}
\begin{proof}[Outline]
Factor \((q)=\mathfrak p^rJ\) at \(\mathfrak p=\varphi^{-1}(\mathfrak m_V)\) and lift \(V\)-divisibility to \(\mathfrak p\)-powers. Comaximality of \(\mathfrak p^{k_N}\) and \(J^N\) gives \(r_N+\varepsilon_N=1\) with \(\varepsilon_N\in J^N\), and telescoping finishes. The complete proof is in \S\ref{prf:B8.28uc}.
\end{proof}

\begin{lemma}[Limiting CM actions preserve integral formal points]\label{B8.28ud}
Use the integral formal chart, \(C\), \(\mathcal O_C\), and \(v_\ell\) of Proposition~\ref{B8.28u}. Write \(d[a]\) for the differential of \([a]\), and write \(\exp_{\widehat G}\) and \(\log_{\widehat G}\) for the mutually inverse formal exponential and logarithm. For each embedding \(\tau\), there is a sequence \(\varepsilon_{\tau,N}\in\mathcal O_E\) with the following property. If \(t\in V\) and \(x_N\in\mathcal O_E\) satisfy \(\tau(x_N)-t\in\ell^NV\), there is \(A\in M_d(\mathcal O_C)\) such that
\[
 v_\ell\bigl(A_{ij}-d[\varepsilon_{\tau,N}x_N]_{ij}\bigr)\ge N,
 \qquad Au_{\tau,a}=t\,u_{\tau,a}
\]
for all \(N,i,j,a\); here \(d=\dim\widehat G\), and the coefficient maps to \(V\) and \(C\) are understood.

Let \(\mathcal B\) be a finite set, and let \(w_b\in(t)C[[t]]^d\), for \(b\in\mathcal B\), where \(t\) now denotes a finite set of formal variables. Suppose that \(\exp_{\widehat G}(w_b)\) has coefficients in \(\mathcal O_C\). If matrices \(A_b\in M_d(C)\) and elements \(a_{N,b}\in\mathcal O_E\) satisfy
\[
 v_\ell\bigl((A_b-d[a_{N,b}])_{ij}\bigr)\ge N
 \quad\text{for all }N,b,i,j,
\]
then \(\exp_{\widehat G}(\sum_bA_bw_b)\) has coefficients in \(\mathcal O_C\).
\end{lemma}
\begin{proof}[Outline]
Lemma~\ref{B8.28uc} at \(q=\ell\) makes \(d[\varepsilon_{\tau,N}x_N]\) entrywise Cauchy with integral limit \(A\) acting as \(t\) on \(\tau\)-eigenvectors. Formal-logarithm additivity makes approximants integral, and polynomiality of each coefficient passes to the limit. The complete proof is in \S\ref{prf:B8.28ud}.
\end{proof}

\begin{lemma}[Integral frame along a limiting isotropic inclusion]\label{B8.28ue}
In the chart of Proposition~\ref{B8.28u}, let \(\mathcal B\) be finite and let \(w_b\in(t)C[[t]]^d\) satisfy
\(\exp_{\widehat G}(w_b)\in\mathcal O_C[[t]]^d\).
Let \(A_{b,j}\in M_d(C)\), for \(b\in\mathcal B\) and \(1\le j\le n\). Suppose there are \(a_{N,b,j}\in\mathcal O_E\) such that
\[
 v_\ell\bigl((A_{b,j}-d[a_{N,b,j}])_{rs}\bigr)\ge N,
 \qquad
 \sum_{j=1}^n\overline{a_{N,b,j}}a_{N,b',j}\in\ell^N\mathcal O_E
\]
for every \(N,b,b',r,s\). Then
\[
 \exp\!\left(\sum_{j=1}^n h_L\!\left(
    \exp_{\widehat G}\Bigl(\sum_bA_{b,j}w_b\Bigr)\right)\right)
\]
is an integral power series with constant term \(1\).
\end{lemma}
\begin{proof}[Outline]
Biextension cocycles make \(h(\exp_{\widehat G}w,\exp_{\widehat G}z)\) biadditive and adjoint modulo integral-unit exponentials, reducing each step to \(U_N\exp(T_N)\). The finite-jet estimate of Lemma~\ref{B8.28ud} on \(\exp(T_N)\) and \(A_{b,j}\) gives integrality. The complete proof is in \S\ref{prf:B8.28ue}.
\end{proof}

\begin{proposition}[Column-differential formal-group matching]\label{B8.28u}
Let \(\lambda\mid\ell\) be a finite place of \(F_c\), \(V\) a trait over \(\lambda\), and \(y\) an integral frame on \(V\) with formal group \(\widehat G\). There exist an algebraically closed complete valued field \(C\) containing \(\operatorname{Frac}V\), with valuation \(v_\ell\) normalised by \(v_\ell(\ell)=1\) and integer ring \(\mathcal O_C\), and integral formal coordinates of \(\widehat G\) over \(V\) from an integral formal biextension chart (meaning \(V\)-coordinates descending the formal group law, the \(\mathcal O_E\)-action with \(\tau\)-eigenvectors \(u_{\tau,i'}\), and the formal Poincar\'e biextension cocycle, symmetry, and adjointness data to \(V\), in which the canonical biextension gauge \(h\) on \(\widehat G\times\widehat G\) and \(h_L(x)=h(x,[w_0]x)\) of Definition~\ref{B8.28g}(d) have zero constant term and coefficients in \(\operatorname{Frac}V\)). Let \((a,N)\) be a theta weight of degree \(D\) with gaps \(c=(c_{\tau,i})\) and \(N\ge\sum_{c_{\tau,i}\ne0}2i\), and let \(P_{\rm col}\in\operatorname{harm}_V(a,N)\) be the image on \(V\) of the column differential of Lemma~\ref{B8.28k}(b) formed from matrices \(B_{\tau,i}\). Put \(\mathcal C=\{(\tau,i,j):c_{\tau,i}\ne0,\ 1\le j\le i\}\). In these \(V\)-coordinates over \(C\), set
\[
 \mathcal D_c(F)=\Bigl(\prod_{c_{\tau,i}\ne0}\det\bigl(\partial/\partial t_{\tau,i;a,j}\bigr)_{a,j\le i}^{\,c_{\tau,i}}\Bigr)
 F\Bigl(\Bigl(\exp\Bigl(\textstyle\sum_{a=1}^i t_{\tau,i;a,j}u_{\tau,a}\Bigr)\Bigr)_{(\tau,i,j)\in\mathcal C}\Bigr)\Big|_{t=0}
\]
for an integral power series \(F\) over \(\mathcal O_C\) on \(\widehat G^{\mathcal C}\), and set
\[
 \mathcal N_k(G)=\bigl[P_{\rm col}(D)\bigl(G\,e^{k\sum_{b=1}^N h_L(x_b)}\bigr)\bigr](0)
\]
for \(k\ge0\) and an integral power series \(G\) over \(\mathcal O_C\) on \(\widehat G^N\). Let \(w_0\in\mathcal O_E\) be as in Definition~\ref{B8.28g}(d).
\begin{itemize}
\item[(i)] For each \(k\ge0\) and each integral power series \(G\) over \(\mathcal O_C\) on \(\widehat G^N\), there exists \(K\ge0\) with the following property. Whenever \(n'\ge0\) and elements \(x_{\alpha,b}\in\mathcal O_E\) (\(\alpha=(\tau,i,j)\in\mathcal C\), \(1\le b\le N\)) satisfy \(v_\ell((B_{\tau,i})_{j,b}-\tau(x_{\alpha,b}))\ge n'\) and \(v_\ell\bigl(\tau'\bigl(w_0\sum_{b=1}^N\bar x_{\alpha,b}x_{\alpha',b}\bigr)\bigr)\ge n'\) for all \(\alpha,\alpha'\in\mathcal C\) with \(\alpha'=(\tau',i',j')\), there is an integral power series \(F\) over \(\mathcal O_C\) on \(\widehat G^{\mathcal C}\) such that \(v_\ell(\mathcal N_k(G)-\mathcal D_c(F))\ge n'-K\).
\item[(ii)] For each \(k\ge0\) and each integral power series \(F\) over \(\mathcal O_C\) on \(\widehat G^{\mathcal C}\), there is an integral power series \(G\) over \(\mathcal O_C\) on \(\widehat G^N\) such that \(\mathcal N_k(G)=\mathcal D_c(F)\).
\end{itemize}
\end{proposition}
\begin{proof}[Outline]
For (i), twist a pullback of \(G\) by a biextension unit and compare exponential jets via \(v_\ell\)-Lipschitz continuity. For (ii), build \(\psi,\rho\) from Lemma~\ref{B8.28k}(b) via Lemmas~\ref{B8.28ua}, \ref{B8.28ud}, and~\ref{B8.28ue}, then apply Lemma~\ref{B8.28ub}. The complete proof is in \S\ref{prf:B8.28u}.
\end{proof}

\begin{lemma}[Local depth of a column differential]\label{B8.28o}
In the setting of Proposition~\ref{B8.28u}, let \(\lambda\mid\ell\) be a finite place of \(F_c\), \(V\) a trait over \(\lambda\), and \(y\) an integral frame on \(V\). Let \((a,N)\) be a theta weight of degree \(D\) with gaps \(c=(c_{\tau,i})\) and \(N\ge\sum_{c_{\tau,i}\ne0}2i\), and let \(P_{\rm col}\in\operatorname{harm}_V(a,N)\) be the image on \(V\) of the column differential of Lemma~\ref{B8.28k}(b). Then \(\delta_k(P_{\rm col};y)=\nu_c(y)\) for every \(k\ge0\).
\end{lemma}
\begin{proof}[Outline]
Reduces \(\delta_k(P_{\rm col};y)\) to integral series in the \(V\)-coordinates of \ref{B8.28u}, then matches \(\inf_F v_\ell(\mathcal D_c(F))\) via \ref{B8.28k}(b) and \ref{B8.28u}. Product testing and monomial expansion decompose this across blocks via \ref{B8.28h}. The complete proof is in \S\ref{prf:B8.28o}.
\end{proof}

\begin{lemma}[Conjugate contractions and denominator powers]\label{B8.28va}
Let \(T\) be a finite set with an involution \(\tau\mapsto\bar\tau\). Let \(N,r_\tau\in\N\), and put
\[
 L=\{(\tau,i):\tau\in T,\ 1\le i\le r_\tau\}.
\]
Let \(K\) be a field of characteristic zero, and let \(P\) be a polynomial over \(K\) in variables \(X_{\ell,b}\), with \(\ell\in L\) and \(1\le b\le N\). Suppose
\[
 \sum_{b=1}^N
 \frac{\partial^2P}
 {\partial X_{(\tau,i),b}\partial X_{(\bar\tau,j),b}}=0
 \qquad
 (\tau\in T,\ 1\le i\le r_\tau,\ 1\le j\le r_{\bar\tau}).
\]
If \(d_{\ell,\ell'}\in K\) vanishes unless
\(\ell=(\tau,i)\) and \(\ell'=(\bar\tau,j)\), then
\[
 \sum_{\ell,\ell'\in L}d_{\ell,\ell'}
       \sum_{b=1}^N
       \frac{\partial^2P}{\partial X_{\ell,b}\partial X_{\ell',b}}=0.
\]

Suppose every monomial of \(P\) has row degree \(a_\ell\in\N\), and set \(D=\sum_\ell a_\ell\). Let \(\Pi\) be a finite set, and let
\(c:T\to\Pi\) and \(s:T\to\{+,-\}\) satisfy
\[
 c(\bar\tau)=c(\tau),\qquad s(\bar\tau)\ne s(\tau).
\]
For a monomial \(X^\mu\) in the support of \(P\), define
\[
 n^\pm_{q,b}(\mu)=
       \sum_{\substack{\ell=(\tau,i)\in L\\c(\tau)=q,\ s(\tau)=\pm}}
            \mu_{\ell,b},
 \qquad
 e(\mu)=\sum_{b=1}^N\sum_{q\in\Pi}
            \min(n^+_{q,b}(\mu),n^-_{q,b}(\mu)).
\]
Then \(e(\mu)\le\lfloor D/2\rfloor\).

Let \(V\) be a discrete valuation ring with fraction field \(K\) and valuation \(v\), let \(e_0\) be a positive integer, and let \(\eta\ge0\) be real. Put
\[
 J_\eta=\{z\in K:v(z)/e_0\ge-\eta\},\qquad v(0)=+\infty.
\]
This is a \(V\)-submodule of \(K\) containing \(1\). For \(m\in\N\), let \(J_\eta^m\) be the \(V\)-submodule generated by products of \(m\) elements of \(J_\eta\), with \(J_\eta^0=V\). Then
\[
 z\in J_\eta^m\quad\Longrightarrow\quad v(z)/e_0\ge-m\eta.
\]
In particular,
\[
 z\in J_\eta^{\lfloor D/2\rfloor}
 \quad\Longrightarrow\quad
 v(z)/e_0\ge-\eta D/2.
\]
\end{lemma}
\begin{proof}[Outline]
Each summand vanishes by nonconjugacy or harmonicity, and \(2\min(A,B)\le A+B\) gives \(2e(\mu)\le D\). The ultrametric inequality makes \(J_\eta\) a \(V\)-submodule containing \(1\), valuation additivity bounds \(J_\eta^m\), and \(\lfloor D/2\rfloor\le D/2\) finishes. The complete proof is in \S\ref{prf:B8.28va}.
\end{proof}

\begin{lemma}[Ordered evaluation from the connection identities]\label{B8.28vb}
Let \(V\) be a discrete valuation ring of characteristic zero with fraction field \(K\), let \(\mathcal F\) be a commutative \(d\)-dimensional formal group law over \(V\), and let \(m\in\N\). Put
\[
 \mathscr A=K[[x_1,\ldots,x_d,z_1,\ldots,z_d]],
\]
let \(\iota f=f(x)\), and let \(\operatorname{ct}\) denote constant coefficient. For \(u\in K^d\), define
\[
 D_uf=\sum_{i,j}u_i\frac{\partial\mathcal F_j}{\partial y_i}(x,0)
                  \frac{\partial f}{\partial x_j},
 \qquad
 \beta((u,\phi),(u',\phi'))=\sum_i(\phi'_iu_i-\phi_i u'_i).
\]
Suppose a family of \(K\)-linear operators \(\nabla_\xi\) on \(\mathscr A\), linear in \(\xi\in K^d\oplus K^d\), satisfies
\[
 [\nabla_\xi,\nabla_\eta]=m\beta(\xi,\eta)\operatorname{id},
 \qquad \nabla_{(0,\phi)}\iota f=0,
\]
and, for every finite sequence \(u_1,\ldots,u_r\),
\[
 \operatorname{ct}\bigl(\nabla_{(u_1,0)}\cdots
       \nabla_{(u_r,0)}\iota f\bigr)
 =\operatorname{ct}(D_{u_1}\cdots D_{u_r}f).
\]

Let \(I,\Pi\) be finite sets, fix an ordering of \(I\), and let
\(c:I\to\Pi\) and \(s:I\to\{+,-\}\). Put \(I_\pm=s^{-1}(\pm)\). For a monomial \(Y^\alpha\), define the ordered operator
\[
 O_\xi(Y^\alpha)
 =\left(\prod_{\ell\in I_+}\nabla_{\xi_\ell}^{\alpha_\ell}\right)
  \left(\prod_{\ell\in I_-}\nabla_{\xi_\ell}^{\alpha_\ell}\right),
\]
where each product follows the fixed ordering, composition acts from right to left, and the minus block acts first. Extend \(O_\xi\) linearly to polynomials. For
\(P=\sum_\mu P_\mu X^\mu\in K[X_{\ell,b}:\ell\in I,1\le b\le N]\),
set \((\mu_b)_\ell=\mu_{\ell,b}\) and
\[
 \mathcal E_\xi(P;g_1,\ldots,g_N)
   =\sum_\mu P_\mu\prod_{b=1}^N
          \operatorname{ct}\bigl(O_\xi(Y^{\mu_b})\iota g_b\bigr).
\]
Say two vector families \(A_\ell,B_\ell\) satisfy opposite-pair vanishing if
\(\beta(A_\ell,B_{\ell'})=0\) whenever
\(c(\ell)\ne c(\ell')\) or \(s(\ell)=s(\ell')\).

(a) Let \(u_\ell,\phi_\ell\in K^d\). Suppose the families
\(((u_\ell,0),(0,\phi_\ell))\) and
\(((0,\phi_\ell),(u_\ell,0))\) satisfy opposite-pair vanishing. Set
\[
 \kappa_{\ell\ell'}=
 \begin{cases}
 m\beta((0,\phi_\ell),(u_{\ell'},0)),
   &c(\ell)=c(\ell'),\ s(\ell)=+,\ s(\ell')=-,\\
 0,&\text{otherwise},
 \end{cases}
\]
and
\[
 \Delta=\sum_{b,\ell,\ell'}\kappa_{\ell\ell'}
              \partial_{X_{\ell,b}}\partial_{X_{\ell',b}}.
\]
If \(\Delta P=0\), then
\[
 \mathcal E_{(u,\phi)}(P;g_1,\ldots,g_N)
   =\left[P(D)\prod_b g_b(x_b)\right]_{x_1=\cdots=x_N=0},
\]
where \(X_{\ell,b}\) acts as \(D_{u_\ell}\) in the \(b\)-th copy.

(b) Let \(J\subset K\) be a \(V\)-submodule containing \(1\).
Choose \(\xi^0_\ell=(u_\ell,\eta_\ell)\) with \(u_\ell\in V^d\),
and \(\psi_\ell\in J^d\). Suppose opposite-pair vanishing holds for
\[
 (\xi^0,\xi^0),\qquad ((0,\psi),\xi^0),\qquad
 (\xi^0,(0,\psi)).
\]
Let \(E_0\in K[[x]]\) have constant term \(1\), and suppose
\[
 (\iota E_0)^{-1}\nabla_{\xi^0_\ell}\bigl((\iota E_0)G\bigr)
       \in V[[x,z]]
 \qquad(G\in V[[x,z]],\ \ell\in I).
\]
If \(P\) has coefficients in \(V\) and every \(\mu\) in its support satisfies
\[
 \sum_{b,q}\min\left(
       \sum_{\substack{c(\ell)=q\\s(\ell)=+}}\mu_{\ell,b},
       \sum_{\substack{c(\ell)=q\\s(\ell)=-}}\mu_{\ell,b}\right)\le e,
\]
then, for \(f_b\in V[[x]]\),
\[
 \mathcal E_{\xi^0+(0,\psi)}
       (P;E_0f_1,\ldots,E_0f_N)\in J^e.
\]
Here \(J^e\) is the \(V\)-submodule generated by products of \(e\) elements of \(J\), and \(J^0=V\).
\end{lemma}
\begin{proof}[Outline]
Commuting vertical operators past the base ones expands each column by \(\Delta_0^r/r!\). For (a), horizontal evaluation and \(\Delta P=0\) leave \(r=0\); for (b), integral contraction bounds and conjugation invariance place the value in \(J^e\). The complete proof is in \S\ref{prf:B8.28vb}.
\end{proof}

\begin{lemma}[Formal connection identities in an integral chart]\label{B8.28vc}
Let \(V\) have positive residue characteristic, and use an integral formal biextension chart and canonical gauge \(h\) as in Proposition~\ref{B8.28u}. Let \(\mathcal F\) be its formal group law. Put \(K=\operatorname{Frac}V\), \(d=\dim\widehat G\), and
\[
 \sigma_i(x)=-\frac{\partial h}{\partial y_i}(x,0),
 \qquad S_{ij}=[x_j]\sigma_i(x),
 \qquad
 H_V=\{(u,\phi-Su):u,\phi\in V^d\}.
\]
For \(m\in\N\), let \(E_0(x)=\exp(mh_L(x))\).
There is a family \(\nabla_\xi\) on \(K[[x,z]]\) satisfying the linearity, commutator, vertical annihilation, and horizontal evaluation hypotheses of Lemma~\ref{B8.28vb}. It also satisfies
\[
 (\iota E_0)^{-1}\nabla_\xi\bigl((\iota E_0)G\bigr)
       \in V[[x,z]]
 \qquad(\xi\in H_V,\ G\in V[[x,z]]).
\]
Fix such a family and use the ordered products of Lemma~\ref{B8.28vb}.
\end{lemma}
\begin{proof}[Outline]
Taking \(\nabla_\xi=\mathcal V_\xi-m\lambda_u\) with commuting vector fields yields the commutator and vertical identities. Intertwining \(\nabla_{(u,0)}\) with \(D_u\) gives horizontal evaluation, and differentiating the gauge laws gives integral conjugation. The complete proof is in \S\ref{prf:B8.28vc}.
\end{proof}

\begin{proposition}[Harmonic valuation lower bound in an integral chart]\label{B8.28v}
For the host field \(E\) of conductor \(Q_Sqd\), put \(D_E=|\operatorname{disc}E|\) and \(\delta_p=v_p(D_E)/[E:\Q]\), and let \((a,N)\) be a theta weight of degree \(D\). Let \(V\) be a trait over a finite place of \(F_c\), with uniformiser \(\varpi\) and valuation \(v\) normalised by \(v(\varpi)=1\), let \(y\) be an integral frame on \(V\), and take integral formal coordinates of \(\widehat G\) over \(V\) from an integral formal biextension chart as in Proposition~\ref{B8.28u}. For every pair of integers \(e>0\) and \(k\ge0\), every polynomial \(R\in\operatorname{harm}_V(a,N)\), and every power series \(F\) on \(\widehat G^N\) with coefficients in \(V\), the value \(z=[R(D)(F\,e^{k\sum_{b=1}^N h_L(x_b)})](0)\) satisfies:
\begin{itemize}
\item[(1)] \(v(z)\ge0\) when \(\operatorname{disc}E\) is a unit in \(V\); and
\item[(2)] \(v(z)/e\ge-(v(p)/e)\,\delta_p\,D/2\) when \(p\) is a prime dividing \(Q_Sqd\) that is not a unit in \(V\).
\end{itemize}
\end{proposition}
\begin{proof}[Outline]
The ultrametric inequality reduces to external products of single-factor series. Idempotent eigenlifts split frame directions into integral and polar parts, and Lemmas~\ref{B8.28va}, \ref{B8.28vb}, and \ref{B8.28vc} yield both valuation bounds. The complete proof is in \S\ref{prf:B8.28v}.
\end{proof}

\begin{lemma}[Chart lower bound and global differential of controlled local depth]\label{B8.28p}
Let \((a,N)\) be a theta weight of degree \(D\) with gaps \(c=(c_{\tau,i})\), and use \(D_E,\delta_p,Q_Sqd\) and the chart data \(V,y,\widehat G,e,k,R,F,z\) of Proposition~\ref{B8.28v}.

(1) For every such \(V,y,e,k,R,F\),
\[
 \frac{v(z)}{e}\ge-\sum_{p\mid Q_Sqd}\frac{v(p)}{e}\,\delta_p\,\frac{D}{2}.
\]

(2) Let \(\lambda\mid\ell\) be a finite place of \(F_c\), and suppose \(N\ge\sum_{c_{\tau,i}\ne0}2i\). Let \(A\le B\) be real numbers such that \(A\le\nu_c(y)\) for all integral frames \(y\) over \(\lambda\) and \(\nu_c(y)\le B\) for all algebraic integral frames \(y\) over \(\lambda\). Then there is \(P_1\in\operatorname{harm}(a,N)\) such that \(A\le\delta_k(P_1;y)\le B\) for all \(k\ge0\) and all algebraic integral frames \(y\) over \(\lambda\). In particular, if \(c\) has oscillation at most \(\varrho\) at \(\ell\), then hypothesis~(i) of Proposition~\ref{B8.28n} holds at \(\lambda\).
\end{lemma}
\begin{proof}[Outline]
For (1), uses \ref{B2.46} and \ref{B8.28v} to sum ramified bounds over \(p\mid Q_Sqd\). For (2), approximates \(P_{\rm col}\) modulo \(\ell^{n_0}\) via \ref{B8.28k}(c) and applies (1), \ref{B8.28u}, \ref{B8.28o}, and \ref{B8.28g}(b) to preserve local depth. The complete proof is in \S\ref{prf:B8.28p}.
\end{proof}

\begin{lemma}[Mixed translates and the chart lower bound]\label{B8.28wa}
Let \(V\) be a complete discrete valuation ring with fraction field \(K\) of characteristic zero and valuation \(v\). Fix an integral frame \(y\), an integral formal chart with frames compatible under truncation, and a split PEL uniformisation with its torsion identification. Let \(L\) be the rigidified theta line and \(s_0\) its principal section in this frame.

For a section \(s\) of \(L^k\) and \(\gamma\in\mathcal O_E\) with
\(\gamma\bar\gamma=C\in\N\), write \(\mathcal P_{\gamma,k}s\) for its pullback by \([\gamma]\), viewed as a section of \(L^{kC}\) through the isomorphism compatible with the rigidification. At odd level \(k\), for a point \(z\) killed by \(k\), write \(\mathcal T_{z,k}s\) for the symmetric theta-group translate of \(s\). Its lift has \(k\)-th power equal to the identity and is compatible with inversion. These normalisations fix both operations.

Let \((a,n)\) be a theta weight of degree \(D\), and put
\[
 \nu_a(\gamma)=\prod_\tau\tau(\gamma)^{\sum_i a_{\tau,i}},
\]
using the coefficient-field maps followed by the map to \(K\).
For a germ \(h\) on \(n\) copies of the formal group, put
\[
 \mathcal H_k(P;h)=
 \left[P(D)\left(h\exp\left(k\sum_{b=1}^n h_L(X_b)\right)\right)\right](0).
\]
The derivations, directions and \(h_L\) are those of Proposition~\ref{B8.28u}. For sections \(S_1,\ldots,S_n\), write
\(\mathcal H_k(P;S_1,\ldots,S_n)\) for this value on the product of their germs in the compatible frames.

Let \(e\ge1\) and \(\rho\in\R\), and suppose that
\[
 v(\mathcal H_k(R;h))/e\ge-\rho
\]
for every \(k\ge0\), every \(R\in\operatorname{harm}_V(a,n)\), and every integral formal series \(h\) on the product chart.
Let \(m,n',A,B\in\N\) and \(\alpha,\beta\in\mathcal O_E\) satisfy
\[
 m\text{ odd},\qquad m=\beta\alpha n',\qquad
 \alpha\bar\alpha=A,\qquad\beta\bar\beta=B,\qquad n'\in V^\times.
\]
For \(1\le b\le n\), let \(z_{1,b},z_{2,b}\) be \(K\)-points satisfying
\[
 [m^2]z_{1,b}=[m^2]z_{2,b}=0,\qquad
 [n'^2\alpha\beta]z_{1,b}=0,\qquad
 [n'^2A\beta]z_{2,b}=0,
\]
and set
\[
 F_b=\mathcal T_{z_{2,b},m^2}\mathcal T_{z_{1,b},m^2}
             \mathcal P_{m,1}s_0.
\]
Then \(m^2=n'^2AB\), and there are sections \(T_b\) of \(L^{n'^2A}\), regular on the integral semiabelian model, such that
\[
 F_b=\mathcal P_{\beta,n'^2A}T_b.
\]
Each \(T_b\) has a unique integral formal germ \(f_b\).
For every \(P\in\operatorname{harm}(a,n)\), its image in \(K[X]\) lies in the indicated integral harmonic module and
\[
 \nu_a(\alpha)m^{-D}\mathcal H_{m^2}(P;F_1,\ldots,F_n)
 =n'^{-D}\mathcal H_{n'^2A}(P;T_1,\ldots,T_n).
\]
Consequently the valuation of the left side, divided by \(e\), is at least \(-\rho\).
\end{lemma}
\begin{proof}[Outline]
Build \(T_b\) from \(s_0\) by four pullbacks and translates satisfying \(\mathcal P_{\beta,n'^2A}T_b=F_b\); these preserve the integral theta lattice, giving integral germs \(f_b\). The pullback chain rule, \(\nu_a(m)=m^D\), and the chart bound conclude. The complete proof is in \S\ref{prf:B8.28wa}.
\end{proof}

\begin{proposition}[Split-uniformisation lower bound for theta derivatives]\label{B8.28w}
At the base level \(M=g_0g_1\), let \(\pi_0,\pi_1\in\mathcal O_E\) be the half-level generators satisfying \(\pi_0\bar\pi_0=g_0\) and \(\pi_1\bar\pi_1=g_1\), with half-level ideal \(\mathfrak b=(\pi_0\pi_1)\). Let \((a,N)\) be a theta weight of degree \(D\), and for \(\beta\in\mathcal O_E\) put
\[
 \nu_a(\beta)=\prod_{\tau:E\to\C}\tau(\beta)^{\sum_{i=1}^{\mu_\tau}a_{\tau,i}}\in\mathcal O_{F_c},
\]
where \(\tau:\mathcal O_E\to\mathcal O_{F_c}\) is the coefficient-field map. Write \(\theta_{x,P}\) for the theta derivative of Lemma~\ref{B8.28m}, and use \(Q_Sqd\) and \(\delta_p\) as in Proposition~\ref{B8.28v}. Let \(V'\) be a trait of ramification index \(e'\) over a finite place \(w\) of \(F_c\) with residue characteristic \(\ell_0\), on which an algebraic integral frame \(y'\) has a split uniformisation and rational \(M^2\)-torsion. Put \(\alpha=1\) if \(\ell_0\nmid M\), \(\alpha=\bar\pi_0\) if \(\ell_0=g_0\), and \(\alpha=\bar\pi_1\) if \(\ell_0=g_1\). Then every \(P\in\operatorname{harm}(a,N)\) and every characteristic tuple \(x\) satisfy
\[
 \frac{v'\bigl(\nu_a(\alpha)\,\theta_{x,P}(y')\bigr)}{e'}
 \ge -\sum_{p\mid Q_Sqd}\frac{v'(p)}{e'}\,\delta_p\,\frac{D}{2}.
\]
\end{proposition}
\begin{proof}[Outline]
Factor \(M=\beta\alpha n'\) by whether \(\ell_0\mid M\) and write \(s_{x_b}\) as a \(\mathfrak b^2\)-torsion translate of \(\mathcal P_{M,1}s_0\). Lemma~\ref{B8.28wa} with Lemma~\ref{B8.28p}(1), Proposition~\ref{B8.28u}, and Lemma~\ref{B8.28m} gives the bound. The complete proof is in \S\ref{prf:B8.28w}.
\end{proof}

\begin{lemma}[A finitely supported lower bound for theta derivatives]\label{B8.28q}
At the base level \(M=g_0g_1\), let \(\pi_0,\pi_1\in\mathcal O_E\) be the half-level generators satisfying \(\pi_0\bar\pi_0=g_0\) and \(\pi_1\bar\pi_1=g_1\), with half-level ideal \(\mathfrak b=(\pi_0\pi_1)\). For a theta weight \((a,N)\) of degree \(D\) and \(\beta\in\mathcal O_E\), set
\[
 \nu_a(\beta)=\prod_{\tau:E\to\C}\tau(\beta)^{\sum_{i=1}^{\mu_\tau}a_{\tau,i}}\in\mathcal O_{F_c},
 \qquad \nu=\nu_a(\bar\pi_0\bar\pi_1),
\]
where \(\tau:\mathcal O_E\to\mathcal O_{F_c}\) is the coefficient-field map, and write \(\theta_{x,P}\) for the theta derivative as in Lemma~\ref{B8.28m}. With \(Q_Sqd\) and \(\delta_p\) as in Lemma~\ref{B8.28p}, assume \(\gcd(M,Q_Sqd)=1\).
For each finite place \(w\) of \(F_c\), put \(\rho(w)=\sum_{p\mid Q_Sqd}v_w(p)\,\delta_p\,D/2\) and \(m(w)=-(\rho(w)+v_w(\nu))\).
\begin{itemize}
\item[(1)] The function \(m\) has finite support on the finite places of \(F_c\).
\item[(2)] For every \(P\in\operatorname{harm}(a,N)\), every characteristic tuple \(x\), every finite place \(w\) of \(F_c\), and every algebraic integral frame \(y\) on a trait \(V\) of ramification index \(e\) over \(w\), one has \(v(\theta_{x,P}(y))/e\ge m(w)\).
\item[(3)] Whenever \(w\mid g_0\) or \(w\mid g_1\), one has \(m(w)=-v_w(\nu)\).
\end{itemize}
\end{lemma}
\begin{proof}[Outline]
Gets (1) and (3) from finite support of \(v_w(\nu),v_w(p)\) and \(\gcd(M,Q_Sqd)=1\). For (2), extends \(V\) to a trait with split uniformization and rational \(M^2\)-torsion, applies \ref{B8.28w}, and uses \(v_w(\nu_a(\alpha))\le v_w(\nu)\). The complete proof is in \S\ref{prf:B8.28q}.
\end{proof}

In the host and nonterminal-layer setting of Definition~\ref{B8.28g} and Lemma~\ref{B8.28b} (with \ref{B2.46}), write \(\eta=\sqrt{-d}\in\mathcal O_E\), let \(\eta_0=\phi_0(\eta)\in\C\) be its image under the distinguished host embedding \(\phi_0:E\hookrightarrow\C\), and put
\[
 K=\Q(\eta),\qquad
 \Psi=\{\tau:E\hookrightarrow\C:\tau(\eta)=\eta_0\},\qquad
 2f=[E:\Q],
\]
so that \(\Psi\) is the CM type of the host and layer. For a signature \((P,Q)\), put \(r=P+Q\). As in Definition~\ref{B8.28g}, the Lie multiplicity is \(\mu_\tau=P\) for \(\tau\in\Psi\) and \(\mu_\tau=Q\) for \(\tau\notin\Psi\). A block \((\tau,i)\) has integer size \(1\le i\le\mu_\tau\), and its exponent is a nonnegative integer. All block depths below use the normalization \(v_\ell(\ell)=1\).

\begin{lemma}[Oscillation at split places]\label{B8.28xe}
Use the host and nonterminal-layer data of Lemma~\ref{B8.28b}, with
\[
 (P,Q)=(k+j-1,m-j),\qquad 1\le j<m,\qquad k\ge2.
\]
Call \(\beta=(\tau,i)\) big when \(\tau\in\Psi\) and \(i\ge k-1\), and small otherwise. Let \(c=(c_\beta)\) be a tuple of nonnegative integer exponents, and let \(U,V\in\R\) satisfy
\[
 \sum_{i=k-1}^{P}c_{\tau,i}\le U\quad(\tau\in\Psi),
 \qquad
 \sum_{\beta=(\tau,i)\ \mathrm{small}}i c_\beta\le V.
\]
Let \(\ell\) be a prime.

(a) Let \(\delta\ge0\), and assume the following split-place condition. For every algebraically closed field \(C\) of characteristic zero equipped with a real-valued nonarchimedean valuation, with valuation ring \(O\) containing \(\Z_\ell\) and with \(\ell\) a nonunit, and every injective map
\[
 \tau_O:\operatorname{Hom}_{\Q}(E,\C)
       \longrightarrow\operatorname{Hom}_{\Z}(\mathcal O_E,O)
\]
satisfying
\[
 \tau(\eta)=\tau'(\eta)
 \quad\Longleftrightarrow\quad
 \tau_O(\tau)(\eta)=\tau_O(\tau')(\eta),
\]
where \(\eta\in\mathcal O_E\) is the host root, so that \(\eta^2=-d\), there is a decomposition
\[
 \Z_\ell\otimes_{\Z}\mathcal O_E
     \simeq\prod_{s=1}^{u}(R_s\times\bar R_s)
\]
into primitive factors paired by conjugation, oriented so that the embeddings through the \(R_s\) are precisely \(\tau_O(\Psi)\), and those through the \(\bar R_s\) are precisely \(\tau_O(\Psi^c)\). Here an embedding through \(R_s\) means a map \(\sigma:\mathcal O_E\to O\) whose extension to \(\Z_\ell\otimes\mathcal O_E\) factors through \(R_s\); each embedding belongs to exactly one factor. Put
\[
 n_s=\operatorname{rank}_{\Z_\ell}R_s,
 \qquad
 \delta_s=\frac{v_\ell(\operatorname{disc}(R_s/\Z_\ell))}{n_s}.
\]
The discriminant here is the determinant of the trace-pairing matrix
\(\bigl(\operatorname{Tr}_{R_s/\Z_\ell}(e_a e_b)\bigr)_{a,b}\)
for a \(\Z_\ell\)-basis \((e_a)\) of \(R_s\), and require that this decomposition satisfy
\[
 \sum_{s=1}^{u} n_s=f,\qquad \sum_{s=1}^{u} n_s\delta_s=f\delta.
\]
Then, for every place \(\lambda\mid\ell\) of \(F_c\), there is an integer \(N_\lambda\ge0\) such that
\[
 N_\lambda\le\nu_c(y)
\]
for every integral frame over \(\lambda\), and
\[
 \nu_c(y)\le N_\lambda+fmU(1+\delta/2)+V
\]
for every algebraic integral frame over \(\lambda\). Thus \(c\) has oscillation at most \(fmU(1+\delta/2)+V\) at \(\ell\).

(b) If \(\ell\nmid2\operatorname{disc}E\) and \(\ell\) splits in \(K\), then for every place \(\lambda\mid\ell\) of \(F_c\), one has \(0\le\nu_c(y)\) at every integral frame over \(\lambda\) and
\[
 \nu_c(y)\le\frac{fQU+V}{\ell-1}
              \le\frac{fmU+V}{\ell-1}
\]
at every algebraic integral frame over \(\lambda\); in particular, \(c\) has oscillation at most \((fQU+V)/(\ell-1)\le(fmU+V)/(\ell-1)\) at \(\ell\).
\end{lemma}
\begin{proof}[Outline]
Lemma~\ref{B8.28h} gives \(N_\lambda=0\) and bounds small blocks by \(V/(\ell-1)\). In (a), split-place row bounds \(t_\tau\) with \(\sum t_\tau\le fQ(1+\delta/2)\) control big blocks; in (b), big-block bounds \(c_\beta Q/(\ell-1)\) sum with the small ones. The complete proof is in \S\ref{prf:B8.28xe}.
\end{proof}

\begin{lemma}[Frobenius orbits and character minors]\label{B8.28xfa}
Use the setting of Definition~\ref{B8.28g}. Let \(p\) be a prime unramified in \(E\), let \(y\) be an integral frame over a trait above \(p\), and choose an integral formal chart \(G\) and an algebraically closed valued extension \(C\) as in that definition. Write \(O\) for its valuation ring, \(|\cdot|\) for its multiplicative valuation, and
\[
 v(x)=\frac{\log|x|}{\log|p|}\quad(x\ne0),\qquad
 v(0)=+\infty.
\]
Thus \(v(p)=1\). Put \(\mu_\sigma=P\) for \(\sigma\in\Psi\) and
\(\mu_\sigma=Q\) otherwise. For each embedding \(\tau\) there are an integer \(h>0\) and distinct embeddings \(\tau_s\), indexed by \(s\in\Z/h\Z\), such that \(\tau_0=\tau\) and
\[
 \tau_{s+1}(a)\equiv\tau_s(a)^p\pmod{\mathfrak m_O}
 \qquad(a\in\mathcal O_E).
\]
Here embeddings into \(C\) are obtained from the coefficient embedding of the trait. If \(p\) splits in \(K\), membership of \(\tau_s\) in \(\Psi\) is constant along this orbit. If \(p\) is inert in \(K\), membership changes at every step.

Suppose also that \(p\ne2\) and \(y\) is algebraic. Let
\(u_{\tau,1},\ldots,u_{\tau,\mu_\tau}\) be its full integral
\(\tau\)-eigenframe and set
\[
 R_\tau=\sum_{m\ge0}p^{-m-1}\mu_{\overline{\tau_{m+1}}}.
\]
For every \(1\le i\le\mu_\tau\), there are integral formal characters
\(\chi_1,\ldots,\chi_i\) of \(G\), meaning
\[
 \chi_b(0)=1,\qquad
 \chi_b(G(X,Y))=\chi_b(X)\chi_b(Y),
\]
for which
\[
 v\!\left(\det\bigl(d\chi_b(u_{\tau,a})\bigr)_{a,b=1}^i\right)
 \le R_\tau.
\]
\end{lemma}
\begin{proof}[Outline]
Lifting Frobenius orbits modulo \(\mathfrak m_O\) gives \(\tau_s\) and its \(\Psi\)-behavior. The tangent sequence and Hodge--Tate datum give characters with a maximal minor bounded by \(R_\tau\); expanding along \(i\) rows selects \(\chi_1,\ldots,\chi_i\). The complete proof is in \S\ref{prf:B8.28xfa}.
\end{proof}

\begin{lemma}[Block depth at an odd unramified prime]\label{B8.28xf}
In the setting of Definition~\ref{B8.28g} with signature \((P,Q)\) and CM type \(\Psi\), let \(\ell\nmid2\operatorname{disc}E\) be a prime. Let \(\tau\in\Psi\), let \(1\le i\le P\) be an integer, and let \(n\ge0\) be an integer. For every algebraic integral frame \(y\) over a place above \(\ell\),
\[
 0\le\nu_{\tau,i}(n;y)\le
 \begin{cases}
 \displaystyle\frac{nQ}{\ell-1},
    &\ell\text{ splits in }K,\\[6pt]
 \displaystyle n\left(\frac{Q}{\ell-1}
              +\frac{(P-Q)\ell}{\ell^2-1}\right),
    &\ell\text{ is inert in }K.
 \end{cases}
\]
The lower bound holds for every integral frame. The expression \(P-Q\) in the second upper bound is subtraction in \(\R\).
\end{lemma}
\begin{proof}[Outline]
Applying \(\mathcal D^n\) to \(\prod_{b=1}^i\chi_b(X^{(b)})\) from Lemma~\ref{B8.28xfa} gives \(\nu_{\tau,i}(n;y)\le nR_\tau\). Summing \(R_\tau\) in the split and inert cases yields the upper bounds, while Lemma~\ref{B8.28h} gives nonnegativity. The complete proof is in \S\ref{prf:B8.28xf}.
\end{proof}

\begin{lemma}[Typed formal factors]\label{B8.28xga}
Let \(p\) be a prime unramified in \(E\), and let \(y,G,C,O,v\) be as in the first paragraph of Lemma~\ref{B8.28xfa}. For an \(i\)-tuple
\(u=(u_1,\ldots,u_i)\), write
\(\operatorname{depth}_G(\det D_u^{\,n})\) for the infimum of the valuations obtained by applying
\(\det(D_{u_a}^{(b)})^n\) to integral series on \(i\) copies of \(G\) and evaluating at the identity.

Fix \(\tau\) and its Frobenius orbit
\((\tau_s)_{s\in\Z/h\Z}\).
There are a commutative integral formal group \(G'\) of dimension \(d'\), with logarithm \(\log_{G'}\) having identity linear term, a type function
\[
 t:\{1,\ldots,d'\}\to\Z/h\Z,
\]
and a linear projection \(\Pi\) on Lie coordinates with the following properties.
For every \(\zeta\in C\) satisfying \(\zeta^{p^h-1}=1\), there is an integral formal endomorphism \([\zeta]\) of \(G'\), with zero constant term, such that
\[
 (\log_{G'}[\zeta](X))_e
   =\zeta^{p^{\widetilde{t(e)}}}(\log_{G'}X)_e,
\]
where \(\widetilde{t(e)}\in\{0,\ldots,h-1\}\) represents \(t(e)\). Moreover,
\[
 \#\{e:t(e)=s\}=\mu_{\tau_s}.
\]
For every flag step in the \(\tau_s\)-eigenspace, its projected frame
\(u'_a=\Pi u_{\tau_s,a}\) is integral, has a unit maximal minor, is supported in coordinates of type \(s\), and satisfies, for every integer \(n\ge0\),
\[
 \operatorname{depth}_{G'}(\det D_{u'}^{\,n})
    =\operatorname{depth}_{G}(\det D_u^{\,n}).
\]
These assertions hold for every integral frame.
\end{lemma}
\begin{proof}[Outline]
Take the formal factor \(G'\) for the orbit component of \(\mathcal O_E\otimes\Z_p\), with Teichm\"uller multiplicities \(\mu_{\tau_s}\). Eigencharacter coordinates give the logarithm law, and integral projection gives \(\Pi\) preserving saturation and depths. The complete proof is in \S\ref{prf:B8.28xga}.
\end{proof}

\begin{lemma}[Lower depth on an even Frobenius orbit]\label{B8.28xgb}
In the setting of Lemma~\ref{B8.28xg}, let \(L=\log_\ell n\).
For every integral frame \(y\) and every chart used to define its block depth,
\[
 \operatorname{depth}_{G}
     (\det D_{u_{\tau,1},\ldots,u_{\tau,i}}^{\,n})
 \ge
 \left(i-Q\left(\frac32+\frac L2\right)\right)
       \frac{n\ell}{\ell^2-1}-i^2(L+1).
\]
\end{lemma}
\begin{proof}[Outline]
By Lemma~\ref{B8.28xga}, alternating CM types make \(h\) even, \(u'\) saturated in type zero, and odd types of rank \(Q\). The even-type formal-group depth bound gives a correction \(c\le Q(3/2+L/2)\), and Lemma~\ref{B8.28xga} transfers the bound back to \(G\). The complete proof is in \S\ref{prf:B8.28xgb}.
\end{proof}

\begin{lemma}[Fine block oscillation at an odd inert prime]\label{B8.28xg}
In the setting of Definition~\ref{B8.28g} with signature \((P,Q)\) and CM type \(\Psi\), put \(r=P+Q\), let \(\ell\nmid2\operatorname{disc}E\) be a prime inert in \(K\), and fix a place \(\lambda\mid\ell\) of \(F_c\). Let \(\tau\in\Psi\), let \(1\le i\le P\) be an integer, and let \(n\ge1\) be an integer. Put
\[
 L=\log_\ell n=\frac{\log n}{\log\ell},\qquad
 W=\frac{n(r-2Q-i)\ell}{\ell^2-1}
       +\frac{nQ(3+L/2)}{\ell-1}
       +i^2(L+1).
\]
There are real numbers \(a\le b\), depending on \(\lambda,\tau,i,n\), such that
\[
 a\le\nu_{\tau,i}(n;y)
\]
for every integral frame over \(\lambda\), and
\[
 \nu_{\tau,i}(n;y)\le b
\]
for every algebraic integral frame over \(\lambda\), with \(b-a\le W\). The subtraction in \(r-2Q-i\) is taken in \(\R\).
\end{lemma}
\begin{proof}[Outline]
Set \(a\) and \(b\) to the frame-independent lower and upper bounds from Lemma~\ref{B8.28xgb} and the inert case of Lemma~\ref{B8.28xf}. Direct algebraic comparison of \(s=\ell/(\ell^2-1)\le t=1/(\ell-1)\) verifies \(a\le b\) and \(b-a\le W\). The complete proof is in \S\ref{prf:B8.28xg}.
\end{proof}

\begin{lemma}[Divided-power determinant content]\label{B8.28xha}
For a prime \(p\) and integers \(i,n\ge0\), put
\[
 H_{i,n}(Z)=\det(Z_{ab})_{a,b=1}^i{}^{\!n}
    =\sum_\alpha c_\alpha Z^\alpha\in\Z[Z_{ab}],
 \qquad
 \alpha!=\prod_{a,b}\alpha_{ab}!.
\]
For any characteristic-zero valued field \((C,v)\) with \(v(p)=1\), define its divided-power content by
\[
 \operatorname{cont}_v(H_{i,n})
   =\min_{c_\alpha\ne0}v(c_\alpha\alpha!).
\]
This is the same finite number
\[
 \kappa_{p,i,n}
   =\min_{c_\alpha\ne0}v_p(c_\alpha\alpha!)
\]
for every such field and valuation.
If \(G\) is a commutative integral formal group and \(u\) is an integral
\(i\)-frame with a unit maximal minor, then
\[
 \operatorname{depth}_G(\det D_u^{\,n})\le\kappa_{p,i,n}.
\]
\end{lemma}
\begin{proof}[Outline]
Nonempty support and \(v(m)=v_p(m)\) for \(m\in\Z\setminus\{0\}\) give \(\operatorname{cont}_v(H_{i,n})=\kappa_{p,i,n}\). Dualizing a unit minor of \(u\) gives a test monomial whose derivative at \(0\) is \(c_\alpha\alpha!\), bounding the depth. The complete proof is in \S\ref{prf:B8.28xha}.
\end{proof}

\begin{lemma}[Two radii for a typed determinant]\label{B8.28xhb}
Let \(p\) be prime, let \(C\) be algebraically closed of characteristic zero, and let \(v\) be an additive valuation with \(v(p)=1\), with valuation ring \(O\). Let \(G\) be a commutative integral formal group of dimension \(d\) with a logarithm having identity linear term. Suppose \(h>0\), a function
\(t:\{1,\ldots,d\}\to\Z/h\Z\), and integral formal endomorphisms
\([\zeta]\) satisfy the root-of-unity formula of Lemma~\ref{B8.28xga}.
Let \(u_1,\ldots,u_i\) be an integral frame with a unit maximal minor, supported in type \(\tau\). If at most \(j\) coordinates have type \(\tau+1\), then for every integer \(n\ge0\),
\[
 \kappa_{p,i,n}\le
 \operatorname{depth}_{G}(\det D_u^{\,n})
       +\frac{2in}{p(p-1)}+\frac{2jn}{p-1}.
\]
\end{lemma}
\begin{proof}[Outline]
An adapted typical-curve basis has at most \(j\) curves with \(L_{r(a),1}\ne0\), and its recursion improves valuations on the rest. Degree induction assigns two weighted-order radii, and applying \(H_{i,n}(\partial)\) with Lemma~\ref{B8.28xha} yields the bound. The complete proof is in \S\ref{prf:B8.28xhb}.
\end{proof}

\begin{lemma}[Block oscillation at an inert prime]\label{B8.28xh}
In the setting of Definition~\ref{B8.28g} with signature \((P,Q)\) and CM type \(\Psi\), let \(\ell\) be a prime unramified in \(E\) and inert in \(K\), and fix a place \(\lambda\mid\ell\) of \(F_c\). Let \(\tau\in\Psi\), let \(1\le i\le P\) be an integer, and let \(n\ge0\) be an integer. There are real numbers \(a\le b\), depending on \(\lambda,\tau,i,n\), with
\[
 a\le\nu_{\tau,i}(n;y)
\]
for every integral frame over \(\lambda\), and
\[
 \nu_{\tau,i}(n;y)\le b
\]
for every algebraic integral frame over \(\lambda\), such that
\[
 b-a\le\frac{2in}{\ell(\ell-1)}+\frac{2Qn}{\ell-1}.
\]
If \(k,m\) are nonnegative integers satisfying
\[
 i+1\le k+m,\qquad k\le m\ell,\qquad Q\le m,
\]
the same endpoints satisfy
\[
 b-a\le\frac{5mn}{\ell-1}.
\]
For \(n=0\), one may take \(a=b=0\).
\end{lemma}
\begin{proof}[Outline]
Lemma~\ref{B8.28h} handles \(n=0\); otherwise set \(b=\kappa_{\ell,i,n}\) and \(a=b-W_0\). Lemma~\ref{B8.28xha} gives the upper bound, Lemmas~\ref{B8.28xga} and \ref{B8.28xhb} (\(j=Q\)) give the lower bound, and the size bounds yield \(5mn/(\ell-1)\). The complete proof is in \S\ref{prf:B8.28xh}.
\end{proof}

\begin{proposition}[The four local formal-group oscillation bounds]\label{B8.28x}
Use a nonterminal layer of signature \((P,Q)=(k+j-1,m-j)\) with \(1\le j<m\) and \(k\ge2\), and the quantities \(r,J,U,V,A,\delta_\ell,h_\ell\) of Lemma~\ref{B8.28b}. Call a block \((\sigma,i)\) big if \(\sigma\) lies in the CM type and \(i\ge k-1\), and small otherwise. Let \(c=(c_{\sigma,i})\) be a gap tuple satisfying
\[
 \sum_{i=k-1}^P c_{\tau,i}\le U \quad(\tau\ \text{in the CM type}),\qquad
 \sum_{\text{small }(\sigma,i)}i\,c_{\sigma,i}\le V,\qquad
 \max_{\sigma,i}c_{\sigma,i}\le J.
\]
Put \(K=\Q(\sqrt{-d})\). Then at every prime \(\ell\le J\):
\begin{itemize}
\item[(1)] if \(\ell\mid Q_S\) and \(\Z_\ell\otimes_{\Z}\mathcal O_E\) decomposes into conjugate pairs \((R_s,\bar R_s)\) of primitive factors whose \(\Z_\ell\)-ranks \(n_s\) sum to \(f\) and whose local discriminant valuations \(v_\ell(\operatorname{disc}(R_s/\Z_\ell))\) sum to \(f\delta_\ell\), then \(c\) has oscillation at most \(A(1+\delta_\ell/2)+V\) at \(\ell\);
\item[(2)] if \(\ell\nmid 2\operatorname{disc}E\) splits in \(K\), then \(\nu_\beta(c_\beta;y)\le c_\beta Q/(\ell-1)\) for every big block \(\beta\) and every algebraic integral frame \(y\) over a place above \(\ell\);
\item[(3)] if \(\ell\nmid 2\operatorname{disc}E\) is inert in \(K\), then at each place \(\lambda\mid\ell\) every big block \(\beta\) with \(c_\beta\ge1\) admits \(A_\beta\le B_\beta\) as in Lemma~\ref{B8.28i} with
\[
 \sum_{\beta\ \mathrm{big},\,c_\beta\ge1}(B_\beta-A_\beta)
 \le \frac{A(4+h_\ell/2)}{\ell-1}+f(m+1)r^2(h_\ell+1);
\]
\item[(4)] if \(\ell\nmid 2\operatorname{disc}E\) is inert in \(K\) and \(m\ell\ge k\), then at each place \(\lambda\mid\ell\) every big block \(\beta\) admits \(A_\beta\le B_\beta\) as in Lemma~\ref{B8.28i} with \(B_\beta-A_\beta\le 5mc_\beta/(\ell-1)\).
\end{itemize}
\end{proposition}
\begin{proof}[Outline]
Part~(1) orients conjugate pairs by an idempotent and Lemma~\ref{B8.28xe}. Parts~(2)--(4) use Lemmas~\ref{B8.28xf}, \ref{B8.28xg}, \ref{B8.28xh}, and \ref{B8.28h}, linking to theta differentials via Propositions~\ref{B8.28u}, \ref{B8.28v}, \ref{B8.28w}. The complete proof is in \S\ref{prf:B8.28x}.
\end{proof}

\begin{lemma}[Exponent bounds and the four oscillation regimes]\label{B8.28r}
Use a nonterminal layer of signature \((P,Q)=(k+j-1,m-j)\) with \(1\le j<m\) and \(k\ge2\), the quantities \(T,r,S,B_0,J,U,V,A,\delta_\ell,h_\ell,e_\ell\) of Lemma~\ref{B8.28b}, and the theta weight \((a,N)\) constructed from \(u_1\) in Lemma~\ref{B8.25g}, with gaps \(c=(c_{\tau,i})\) and degree \(D\). Call a block \((\sigma,i)\) big if \(\sigma\) lies in the CM type and \(i\ge k-1\), and small otherwise.

(1) Suppose that, with zero-based column indices on \(u_1\), each row \(u_{1,\tau}\) is nonincreasing and satisfies \(u_{1,\tau,k-2}=0\), \(|u_{1,\tau,s}|\le B_0\) for \(0\le s<r\), \(u_{1,\tau,P-1}-u_{1,\tau,P}\ge N+1\), and \(\sum_{0\le s<k-2}u_{1,\tau,s}\le k(S+2r)/T\). Then
\[
 \sum_{i=k-1}^P c_{\tau,i}\le U \quad(\tau\ \text{in the CM type}),\qquad
 \sum_{\text{small }(\sigma,i)}i\,c_{\sigma,i}\le V,\qquad
 \max_{\sigma,i}c_{\sigma,i}\le J,
\]
and \(D\le frU+V\).

(2) Assume the bounds of (1), and assume that \(2\mid Q_S\) and every prime \(\ell\mid Q_S\) splits the quadratic subfield \(K=\Q(\sqrt{-d})\) over \(\Z_\ell\). Then every prime \(\ell\mid Q_S\) decomposes \(\Z_\ell\otimes_{\Z}\mathcal O_E\) into conjugate pairs \((R_s,\bar R_s)\) of primitive factors whose \(\Z_\ell\)-ranks \(n_s=\operatorname{rk}_{\Z_\ell}(R_s)\) sum to \(f\) and whose rank-weighted local discriminant depths \(n_s\cdot(v_\ell(\operatorname{disc}(R_s/\Z_\ell))/n_s)=v_\ell(\operatorname{disc}(R_s/\Z_\ell))\) sum to \(f\delta_\ell\), and every prime \(\ell\nmid Q_S\) outside \(\{q,d\}\) is odd and does not divide \(\operatorname{disc}E\). Moreover, \(c\) has oscillation at most \(e_\ell\) at every prime \(\ell\).
\end{lemma}
\begin{proof}[Outline]
For (1), telescopes \ref{B8.25g} and sums by parts to bound big and small gaps and \(D\). For (2), splits \(\Z_\ell\otimes_{\Z}\mathcal O_E\) via \ref{B2.46} and \(K\)-idempotents, then combines \ref{B8.28h}, \ref{B8.28i}, and \ref{B8.28x} across all prime regimes. The complete proof is in \S\ref{prf:B8.28r}.
\end{proof}

\begin{lemma}[Harmonic minors on two column sets]\label{B8.28ya}
Let \(\mathcal E=\operatorname{Hom}_{\Q}(E,\C)\), with conjugation \(\tau\mapsto\bar\tau\), and use the multiplicities \(\mu_\tau\) of Definition~\ref{B8.28g}. Let \(a_{\tau,i}\), \(1\le i\le\mu_\tau\), be nonnegative integers that are nonincreasing in \(i\). Put \(a_{\tau,\mu_\tau+1}=0\) and
\[
 c_{\tau,i}=a_{\tau,i}-a_{\tau,i+1},\qquad
 D=\sum_{\tau\in\mathcal E}\sum_{i=1}^{\mu_\tau}a_{\tau,i}.
\]
Let \(S\subset\mathcal E\) contain exactly one member of each conjugate pair. Choose \(n_1,n_2\in\N\) such that \(\mu_\tau\le n_1,n_2\) for every \(\tau\), and put \(N=n_1+n_2\). In the variables \(X_{\tau,i,b}\), \(1\le b\le N\), define
\[
 \Delta_{\tau,j}^{(1)}
 =\det(X_{\tau,i,b})_{1\le i,b\le j},
 \qquad
 \Delta_{\tau,j}^{(2)}
 =\det(X_{\tau,i,n_1+b})_{1\le i,b\le j},
\]
and
\[
 P_S=
 \prod_{\tau\in S}\prod_{j=1}^{\mu_\tau}
       (\Delta_{\tau,j}^{(1)})^{c_{\tau,j}}
 \prod_{\tau\notin S}\prod_{j=1}^{\mu_\tau}
       (\Delta_{\tau,j}^{(2)})^{c_{\tau,j}}.
\]
Then \(P_S\) has integer coefficients and belongs to the harmonic polynomial space \(\operatorname{harm}(a,N)\) of Definition~\ref{B8.28g}(c). It has degree \(D\) and row degree \(a_{\tau,i}\). Its displayed factorisation is preserved by every change of coefficient ring.
\end{lemma}
\begin{proof}[Outline]
Telescoping \(c_{\tau,j}\) gives row degrees \(a_{\tau,i}\), degree \(D\), and row invariance of Definition~\ref{B8.28g}(c). Disjoint columns for conjugate types kill the mixed second derivatives, and determinants commute with ring maps. The complete proof is in \S\ref{prf:B8.28ya}.
\end{proof}

\begin{lemma}[A theta family with a controlled harmonic value]\label{B8.28yb}
Use the base level \(M=gg'\), where \(g,g'\) are distinct odd primes, and half-level generators \(\pi,\pi'\) satisfying
\[
 \pi\bar\pi=g,\qquad \pi'\bar\pi'=g',\qquad
 (\pi)+(\bar\pi)=(\pi')+(\bar\pi')=\mathcal O_E.
\]
Assume that \(g\) splits completely in \(E\). Let \(y\) be an algebraic integral frame over a trait \(V\) of residue characteristic \(g\), with an integral chart compatible under truncation and a split PEL uniformisation with its torsion identification. Write \(X\) for the period lattice and \(Q_X\) for its cubical-period form. Let
\(b_1,\ldots,b_s\) be a basis such that
\[
 Q_X(b_j)\le\Phi,\qquad
 |Q_X(b_j+z)-Q_X(b_j)-Q_X(z)|
       \le2\sqrt s\,Q_X(z)\quad(z\in X),\qquad
 \Phi\ge0.
\]
Let \(F\) be an algebraically closed valued extension of
\(\operatorname{Frac}V\), of characteristic zero, whose additive valuation
\[
 v_F:F\to\Q\cup\{+\infty\}
\]
extends \(v\), whose valuation ring \(\mathcal O_F\) is \(g\)-adically complete, and over which the PEL Tate module is identified with the fixed PEL datum, including the Weil pairings.

Let \((a,N)\) be a theta weight of degree \(D\), and write
\(N=n_1+n_2\ge1\), with \(\mu_\tau\le n_1,n_2\) for all \(\tau\). Put
\[
 S=\{\tau:v_F(\tau(\bar\pi))>0\},
\]
using the coefficient-field maps followed by
\(\operatorname{Frac}V\to F\).
Let \(P_S\) be the polynomial of Lemma~\ref{B8.28ya}, with these two column sets. Assume
\[
 D<g-1,\qquad 4D+4\le g',\qquad 2D\sqrt s<g'.
\]
For a level-\(k\) theta family \(s\) on the product uniformisation, write
\(s\in\mathscr T_k^+\) when its coefficient at every \(\chi\in X^N\), measured in integral frames on the abelian part, has valuation at least
\[
 \frac{Q_N(\chi)}k,\qquad
 Q_N(\chi)=\sum_{b=1}^N Q_X(\chi_b).
\]
This is the theta lattice of Lemma~\ref{B8.28d}.

There exist a family \(s\in\mathscr T_{g'}^+\), its normalised pullback
\(s_2\) by \([\bar\pi]\) of level \(M\), and a finite-order theta-group translate \(t\) of \(s_2\) at a point of \(\ker[\pi]^N\), whose germ \(h\) in the product integral frame satisfies
\[
 \rho:=v_F(\mathcal H_M(P_S;h))\in\Q,
 \qquad
 \rho\le\frac{v(g)D}{g-1}
             +\frac{n_1s\Phi}{4g'}+\frac{2D^2\Phi}{g'}.
\]
Here a finite-order theta-group translate means the action obtained from translation and a nonzero scalar for which some positive iterate acts as the identity on theta families. The notation \(\mathcal H_M\) is that of Lemma~\ref{B8.28wa}.
\end{lemma}
\begin{proof}[Outline]
Type separators and chart characters bound the minor term by \(v(g)D/(g-1)\). The point construction with Lemmas~\ref{B8.28f} and \ref{B8.28d} bounds the period and jet terms, and external products with Lemma~\ref{B8.28ya} pass the bound to \(\mathcal H_M(P_S;h)\). The complete proof is in \S\ref{prf:B8.28yb}.
\end{proof}

\begin{lemma}[Integral expansion and the theta normaliser]\label{B8.28yc}
Retain the field, frame, split uniformisation and half-level data of Lemma~\ref{B8.28yb}, with the base half-level ideal
\(\mathfrak b=(\pi\pi')\), and set
\[
 \alpha=\bar\pi\bar\pi',\qquad \alpha\bar\alpha=M.
\]
Let \(s\in\mathscr T_{g'}^+\), let \(s_2\) be its normalised
\([\bar\pi]\)-pullback, and let \(h\) be the compatible product germ of a finite-order translate of \(s_2\) at a point killed by \([\pi]\) in every factor.
Then there are a finite number \(q\) of scalars \(a_i\in\mathcal O_F\), a principal theta family \(\vartheta\in\mathscr T_1^+\) whose coefficient at period index \(0\) is the base change of an integral section not divisible by a uniformiser of \(V\), and germs \(h_i\) such that
\[
 h=\sum_{i=1}^q a_i h_i.
\]
Each \(h_i\) is the compatible germ of a finite-order translate of the normalised \([\alpha]\)-pullback of \(\vartheta\), at a point killed by
\([\bar\alpha]=[\pi\pi']\) in every factor.
For each \(i\) there are a characteristic tuple \(x_i\) and a scalar
\(u_i\in\mathcal O_F^\times\) such that, for every
\(P\in\operatorname{harm}(a,N)\),
\[
 \mathcal H_M(P;h_i)
       =u_i\,\nu_a(\alpha)\,\theta_{x_i,P}(y).
\]
In particular, if \(\rho=v_F(\mathcal H_M(P;h))<\infty\), some \(i\) satisfies
\[
 v(\theta_{x_i,P}(y))+v(\nu_a(\alpha))\le\rho.
\]
\end{lemma}
\begin{proof}[Outline]
Integral generation expands \(h=\sum_i a_i h_i\) (\(a_i\in\mathcal O_F\)) via translates of the \([\alpha]\)-pullback of \(\vartheta\). Frame splitting and the pullback chain rule give \(u_i\nu_a(\alpha)\theta_{x_i,P}(y)\), and the ultrametric inequality concludes. The complete proof is in \S\ref{prf:B8.28yc}.
\end{proof}

\begin{proposition}[Split-uniformisation upper bound at the level primes]\label{B8.28y}
Use the base level \(M=g_0g_1\), whose distinct odd primes \(g_0,g_1\) split completely in \(E\), and the half-level generators \(\pi_0,\pi_1\in\mathcal O_E\) satisfying \(\mathfrak b+\bar{\mathfrak b}=\mathcal O_E\) as in Proposition~\ref{B8.28w}. Let \((a,N)\) be a theta weight of degree \(D\) with \(N\ge2r\), where \(r=P+Q\), and put \(\nu=\nu_a(\bar\pi_0\bar\pi_1)\). Take \((g,g')=(g_0,g_1)\) or \((g_1,g_0)\), and let \(V'\) be a trait over a place \(\lambda\mid g\) of \(F_c\) on which an algebraic integral frame \(y'\) has a split uniformisation with period lattice of rank \(s\) and short basis as in Lemma~\ref{B8.28c} of parameter \(\Phi'\). If \(D<g-1\), \(4D+4\le g'\), and \(2D\sqrt s<g'\), then there exist \(P_1\in\operatorname{harm}(a,N)\) and a characteristic tuple \(x\) satisfying
\[
 v'\bigl(\theta_{x,P_1}(y')\bigr)+v'(\nu)
 \le\frac{v'(g)\,D}{g-1}+\frac{rs\Phi'}{4g'}+\frac{2D^2\Phi'}{g'}.
\]
\end{proposition}
\begin{proof}[Outline]
Over an algebraically closed \(g\)-adically complete extension \(F\) with a PEL Tate isometry, set \(n_1=r\) and form \(P_1=P_S\) by Lemma~\ref{B8.28ya}. Lemma~\ref{B8.28yb} gives a germ \(h\) with bounded \(\rho\), and Lemma~\ref{B8.28yc} yields \(x\). The complete proof is in \S\ref{prf:B8.28y}.
\end{proof}

\begin{proposition}[Reference depth at the level primes]\label{B8.28s}
Let \((a,N)\) be a theta weight of degree \(D\) with \(N\ge2r\) and \(a_{\tau,P}>0\) for every \(\tau\) in the CM type. Let \(c\ge0\) satisfy, for both \(g\in\{g_0,g_1\}\),
\[
 4D+4\le g,\qquad 2D\sqrt{2fm}<g,\qquad
 \kappa(2fm)^3\Bigl(\frac{r(2fm)}{4}+2D^2\Bigr)\le c\,g,
\]
with \(\kappa\) as in Lemma~\ref{B8.28c}. Assume \(\gcd(g_0g_1,Q_Sqd)=1\). Then \(d^c_g(a,N)\le 1\) for \(g=g_0,g_1\).
\end{proposition}
\begin{proof}[Outline]
Passes to a split-uniformization extension \(V'\) and applies \ref{B8.28c} and \ref{B8.28y} to get a cusp section \(\theta_{x,P_1}\) via \ref{B8.25t}. Clears denominators using \ref{B8.28q} and \ref{B8.28j}, then takes the supremum in \ref{B8.28g}(e). The complete proof is in \S\ref{prf:B8.28s}.
\end{proof}

\begin{proposition}[Weight-change inequality for reference depths]\label{B8.28z}
Fix eligible \(F_m,C_H\) and assume the packages \(\mathcal H_P,\mathcal H_H\), deep goodness, and flat-degree conditions of Theorem~\ref{B8.25} and Definition~\ref{def:p8c3-section-laws}, together with the size and good-cover conditions of Lemmas~\ref{B8.25b} and~\ref{B8.25c}. At an admissible point beyond the common family threshold, at a nonterminal layer take a good cover \(Y\) of level \(M'\) as in Lemma~\ref{B8.25c}, a quotient weight \(u_1\), and the theta weight \((a,N)\) of Lemma~\ref{B8.25g}. Let \(\mathcal B\) be the set of primes dividing \(2dD_EM'\) and put \(\delta_\ell=v_\ell(D_E)/(2f)\) as in Lemma~\ref{B8.28b}. Then on \(Y\), for every prime \(\ell\),
\[
 d^{1/2500}_\ell(u_1)\le d^{1/2500}_\ell(a,N)+{\bf1}_{\ell\in\mathcal B}\bigl(1+\tfrac N2fr\,\delta_\ell\bigr).
\]
\end{proposition}
\begin{proof}[Outline]
Shifts \((a,N)\) to \(u_1\) via determinant and flat sections from \ref{B8.25g}, \ref{B8.24b}, and \ref{B8.24d}. Clears denominators via \ref{B8.28j}, with zero shift off \(\mathcal B\) and the different shift via \ref{B8.24c} on \(\mathcal B\). The complete proof is in \S\ref{prf:B8.28z}.
\end{proof}

\begin{proposition}[Combining the primes and changing to the quotient weight]\label{B8.28t}
Fix eligible \(F_m,C_H\) and assume the packages \(\mathcal H_P,\mathcal H_H\), deep goodness, and flat-degree conditions of Theorem~\ref{B8.25} and Definition~\ref{def:p8c3-section-laws}, together with the size and good-cover conditions of Lemmas~\ref{B8.25b} and~\ref{B8.25c}. At an admissible point beyond the common threshold of Lemmas~\ref{B8.25b}, \ref{B8.25g}, and~\ref{B8.28b} and Proposition~\ref{B8.28z}, use a nonterminal layer, good cover \(Y\) of level \(M'\) as in Lemma~\ref{B8.25c}, quotient weight \(u_1\), theta weight \((a,N)\) of Lemma~\ref{B8.25g} of degree \(D\), and \(L,T,J,e_\ell,a(\ell)\), bad-prime set \(\mathcal B\), and constant \(C_R\) of Lemma~\ref{B8.28b}.

(1) Every prime \(\ell\) satisfies
\[
 d^{1/2500}_\ell(a,N)\le e_\ell+{\bf1}_{\ell\in\{g_0,g_1\}}
\]
both on the base moduli and on \(Y\).

(2) Put \(\mathcal S_{\rm ref}=\{\ell\text{ prime}:\ell\le J\}\cup\{g_0,g_1\}\cup\mathcal B\). Then \(a(\ell)\ge0\) for every prime \(\ell\), \(a(\ell)=0\) for \(\ell\notin\mathcal S_{\rm ref}\), \(d^{1/2500}_\ell(u_1)\le a(\ell)\) for every prime \(\ell\) on \(Y\), and \(T\sum_{\ell\in\mathcal S_{\rm ref}}a(\ell)\log\ell\le C_RE_p\).
\end{proposition}
\begin{proof}[Outline]
Bounds \(d_\ell^{1/2500}(a,N)\) by \ref{B8.28s} at \(g_0,g_1\) and by \ref{B8.28r}, \ref{B8.28p}(2), \ref{B8.28q}, and \ref{B8.28n} elsewhere using \ref{B8.25g} and \ref{B2.46}. Transfers to \(Y\) via \ref{B8.28l}, then combines \ref{B8.28z} and \ref{B8.28b} for (2). The complete proof is in \S\ref{prf:B8.28t}.
\end{proof}

\begin{proposition}[Heights of the covering components]\label{B8.28}
Assume the self-intersection and section-injection comparisons and the packages \(\mathcal H_X,\mathcal H_P,\mathcal H_H\) of \ref{B8.25}, Grothendieck deformation, \ref{B8.26}, and uniform deep goodness, size, cover degree and flat degree. Fix eligible \(F_m,C_H\). There exist \(C_{\rm ht}\ge0\) and \(x_0\) such that for every admissible point with \(x\ge x_0\), every nonterminal layer, good cover \(Y\) and quotient weight \(u_1\), there are \(c_s,c_c,X\ge0\) and real \(s_Z\) satisfying
\begin{equation}
 h_Z+c_s\ge0,\qquad (1-\rho)\max(0,s_Z+c_s)\le h_Z+c_s+c_c,\qquad
 \frac1{\#\operatorname{Comp}(Y)}\sum_Zs_Z\ge-X
 \tag{8.28c}\label{eq:p8c3-height-comp}
\end{equation}
and
\begin{equation}
 T_*X,\ T_*c_c\le C_{\rm ht}E_p,\qquad
 c_s\bigl(6d_{\rm fl}+10+T_*(z+2\rho)\bigr)\le C_{\rm ht}E_p .
 \tag{8.28d}\label{eq:p8c3-height-budget}
\end{equation}
\end{proposition}
\begin{proof}[Outline]
Builds short generators via \ref{B8.25f} for \(h_Z+c_s\ge0\) and a carrying section \(\theta_Z\) via \ref{B8.25g}, \ref{B8.28a}, and \ref{B8.28t}(2). Maps small sections of \(u'\) from \ref{B8.26} into \(u_1\) and compares volumes via \ref{B8.25} and \ref{B8.25a}. The complete proof is in \S\ref{prf:B8.28}.
\end{proof}

\section{From component averages to the core lattice}

\begin{proposition}[Cover and core slopes]\label{B8.29}
Under the uniform goodness, size and component estimates of \ref{B8.25e} and \ref{B8.28},
\begin{equation}
 \mu_Y(u_M)\ge-C_UE_p,\qquad C_U=2(3C_{\rm ht}+C_{\rm HS}+4).
 \tag{8.29d}\label{eq:p8c3-cover-slope}
\end{equation}
Under the geometric hypotheses of \ref{B8.25} and the family conditions above, the core type lattice satisfies, eventually at every admissible point and layer,
\begin{equation}
 \mu_{\rm core}(p,i)\ge-M_{p,i}-C_{\rm core}E_p ,
 \tag{8.29e}\label{eq:p8c3-core-slope}
\end{equation}
with
\begin{equation}
 C_{\rm core}=\max\{C_{\rm term},\,2(C_F+C_U)+9\},
 \tag{8.29f}\label{eq:p8c3-core-constant}
\end{equation}
where \(C_F=1183\) and \(C_{\rm term}\) is the terminal constant below.
\end{proposition}
\begin{proof}[Outline]
Turn the mean height of \ref{B8.28} into a cover slope by degree and rank bounds. Descend via \ref{B8.24a} and \ref{B8.25b}, using a translation determinant to inject into zero-shift lattices at nonterminal layers and the definite Gram comparison at the terminal layer. The complete proof is in \S\ref{prf:B8.29}.
\end{proof}

\section{Removing the bad part}

In the rational realisation let \(\Lambda_{\rm type}\subseteq\Lambda_{\rm all}\) be the type and ambient lattices. For the good projector \(e\) set
\begin{equation}
 \Lambda_{\rm good}=\Lambda_{\rm type}\cap\operatorname{im}e,\qquad
 \Lambda_{\rm bad}=\Lambda_{\rm type}\cap\ker e,\qquad
 \delta=\exp(-x^{e_z/7}),\quad e_z=\varepsilon/J,\qquad J\ge1\text{ the family exponent parameter}.
 \tag{8.30a}\label{eq:p8c3-bad-lattices}
\end{equation}

\begin{proposition}[Rank of the bad lattice]\label{B8.30}
Assume the geometric and section comparisons of \ref{B8.25}, the cuspidal Green parametrisation and the uniform Gram, cover, size, flat degree and share conditions above. For each eligible \(F_m,C_H\), eventually at every admissible point and layer,
\begin{equation}
 \operatorname{rk}\Lambda_{\rm bad}\le3\delta\operatorname{rk}\Lambda_{\rm type}.
 \tag{8.30b}\label{eq:p8c3-bad-rank}
\end{equation}
\end{proposition}
\begin{proof}[Outline]
Use \ref{B7.50} and \ref{B7.51} to bound bad label shares. Orbit separation from \ref{B8.29} embeds nonterminal cover spaces in regular copies, and terminal summands already supply such copies, so projector traces bound the common good rank loss. The complete proof is in \S\ref{prf:B8.30}.
\end{proof}

\begin{lemma}[Slope loss under orthogonal truncation]\label{B8.31}
Over a totally complex number field \(K\), let \(V\) have Hermitian norms \(N_\sigma\) with a Gram frame. Take a finitely generated \(\mathcal O_K\)-lattice \(M\subset V\), \(0\ne L\subseteq M\), and an idempotent \(e\) with
\begin{equation}
 N_\sigma(ex+y-ey)^2=N_\sigma(ex)^2+N_\sigma(y-ey)^2
 \qquad(\sigma,x,y).
 \tag{8.31a}\label{eq:p8c3-ortho}
\end{equation}
Suppose \(g\ge1\) is integral, \(ge(L)\subseteq L\), \(0\le\delta\le1/12\), and
\begin{equation}
 \operatorname{rk}(L\cap\ker e)\le3\delta\operatorname{rk}L,\qquad
 \mu_{\max}(M),\ -\mu(L),\ 2\log g\le\mathcal Q .
 \tag{8.31b}\label{eq:p8c3-cut-premises}
\end{equation}
Then
\begin{equation}
 \mu(L\cap\operatorname{im}e)\ge\mu(L)-12\delta\mathcal Q .
 \tag{8.31c}\label{eq:p8c3-cut-conclusion}
\end{equation}
\end{lemma}
\begin{proof}[Outline]
With \(L_g=L\cap\operatorname{im}e,L_b=L\cap\ker e\), \(gL\subseteq L_g\oplus L_b\) gives \(\operatorname{rk}L=d_g+d_b\) and \(1-e\) embeds \(L/(L_g\oplus L_b)\) in \(g^{-1}L_b/L_b\). Orthogonality and \(d_b\le 4\delta d_g\) bound \(\mu(L)-\mu(L_g)\) by \(12\delta\mathcal{Q}\). The complete proof is in \S\ref{prf:B8.31}.
\end{proof}

\section{The layer inequality and the common constant}

Give the lattices the rational-realisation norm. In the recursive comparison \(d_W(i)=D_{\rm norm}>0\) by the primitive row equality. In the selected basis the frame comparison carries the good norm to the recursive row norm with factor \(1/\sqrt{n_P}\le1\), where \(n_P\) is the positive row multiplicity. The primitive lattice inclusion gives the determinant comparison
\begin{equation}
 \Lambda_{\rm all}\text{ finitely generated},\qquad
 \Lambda_{\rm good}\ne0,\qquad
 \frac{\operatorname{lcv}(X_{W,i},\Lambda_i)}{D_{\rm norm}}\le-\mu(\Lambda_{\rm good}).
 \tag{8.32a}\label{eq:p8c3-gram}
\end{equation}
Let \(G\) be the finite level group acting on the realised row lattice, and write its good projector as \(e=\sum_{h\in G}a_hh\). The tied coefficients satisfy
\[
 \sum_{h\in G}a_ha_{h^{-1}k}=a_k,\qquad
 \sigma(a_{h^{-1}})=\overline{\sigma(a_h)},\qquad |G|a_h\in\mathcal O_K .
\]
The action preserves the lattice and the inner product. The first two identities give idempotence and self-adjointness, hence \eqref{eq:p8c3-ortho}. Set \(g=|G|\), which clears the denominators on \(\Lambda_{\rm type}\). The group-order bound gives \(2\log g\le b\).

By \ref{B8.25b} and \ref{B8.25d}, eventually there are family exponents \(A,B\) with
\begin{equation}
 \mu_{\max}(\Lambda_{\rm all})\le x^Ab,\qquad
 1\le b\le E_p\le x^Bb,\qquad |M_{p,i}|\le x^Bb .
 \tag{8.32b}\label{eq:p8c3-sizes}
\end{equation}
Let \(K\) be the realisation field. The type lattice \(\Lambda_{\rm type}\) is the \(\mathcal O_K\)-span of the integral row forms lying in the type subspace. Under the semilinear identification these forms are exactly the images of the core type lattice. Field extension preserves the slope in the normalised core norm. The norm change in \ref{B8.25d} costs at most \(2C_{\rm meas}E_p\) in slope. With \ref{B8.29} this gives
\begin{equation}
 \mu(\Lambda_{\rm type})\ge\mu_{\rm core}-C_1E_p,\qquad
 \mu_{\rm core}\ge-M_{p,i}-C_2E_p ,
 \tag{8.32c}\label{eq:p8c3-core}
\end{equation}
with \(C_1=2C_{\rm meas}\) and \(C_2=C_{\rm core}\).

\begin{lemma}[Covolume at one layer]\label{B8.32}
At \(x>0\) assume the Gram, recursive, projector, bad-rank, maximum-slope, full size and two core bounds above, with \(C_1,C_2\ge0\). If
\begin{equation}
 12e^{-x^{e_z/7}}(x^A+x^B+C_1+C_2+1)\le1 ,
 \tag{8.32d}\label{eq:p8c3-loss}
\end{equation}
then
\begin{equation}
 \frac{\operatorname{lcv}(X_{W,i},\Lambda_i)}{D_{\rm norm}}
 \le M_{p,i}+(C_1+C_2+1)E_p .
 \tag{8.32e}\label{eq:p8c3-layer}
\end{equation}
\end{lemma}
\begin{proof}
Put \(H=x^A+x^B+C_1+C_2+1\). By \eqref{eq:p8c3-loss} and \(H\ge1\), \(\delta\le1/12\). Equations \eqref{eq:p8c3-sizes}, \eqref{eq:p8c3-core} and \(2\log g\le b\) bound the three quantities in \eqref{eq:p8c3-cut-premises} by \(HE_p\). Apply \ref{B8.31} to \(\Lambda_{\rm type}\subseteq\Lambda_{\rm all}\) with \(\mathcal Q=HE_p\). The cut costs at most \(12\delta HE_p\le E_p\). Apply \eqref{eq:p8c3-gram} and \eqref{eq:p8c3-core}.
\end{proof}

\begin{proposition}[One constant for the family]\label{B8.33}
If the Gram, recursive, truncation, bad-rank, maximum-slope, size and two core bounds hold eventually on the eligible families, then \ref{B8.24} holds with
\begin{equation}
 C_0=C_1+C_2+1 .
 \tag{8.33a}\label{eq:p8c3-family-constant}
\end{equation}
\end{proposition}
\begin{proof}
Fix \(F_m,C_H\) and \(A,B,C_1,C_2\). Since \(e_z>0\), \(e^{-x^{e_z/7}}x^a\to0\) for \(a=A,B,0\). Take the maximum of the input thresholds and a threshold for \eqref{eq:p8c3-loss}. Apply \ref{B8.32} at all admissible points and layers beyond it. The same threshold works for \(C\ge C_0\) by \(E_p\ge0\).
\end{proof}

\section{Transfer to the selected indefinite realisation}

\begin{lemma}[Size of the field generated by the coefficient list]\label{B8.34a}
Let \(E\) be the abelian CM host, of degree \(2f\), with distinguished quadratic subfield \(K=\Q(\sqrt{-d})\) of discriminant \(-d\). Let \(F\subset\C\) be the field generated by the coefficient list of the constructed member in its specified embedding. Suppose \(f\ge8\), \(r\ge62\), \(L\ge4\), and
\[
 \log f\le L,\quad \log\operatorname{rd}(E)\le2L,\quad
 \log g_i\le9L+1\ (i=0,1),\quad
 \log p_{\rm de},\ \log p'_{\rm de}\le L .
\]
Then \(\log[F:\Q]\) and \(\log\operatorname{rd}(F)\) are each at most
\(19fr^2L+198f^2L\).
\end{lemma}
\begin{proof}[Outline]
Place the coefficient field in a compositum of torsion, ray class, cyclotomic and Kummer fields. Use \ref{B6.05f}(2) for the CM class field, bound degrees and differents, and combine the bounds by the compositum inequalities of \ref{B8.25c}. The complete proof is in \S\ref{prf:B8.34a}.
\end{proof}

\begin{theorem}[Covolume bounds on the linked bundles]\label{B8.34}
Assume the discrete-series covering, Harish--Chandra module theory, cusp rapid decay, Casselman--Wallach and ray-class-field statements for the linked bundles, and \ref{B8.24}. Let \(h_L\) give local Langlands parameters, \(N\) be a norm-choice extension, and \(S,\Gamma\) be as in \ref{B8.08}. Assume a height-counting extension with its law for each datum admitted by \(N\). Assume the twelve open conditions of \ref{B8.08} for the parameters determined by \(N,S,\Gamma\). There is a condition \(R\) on members which eventually holds for every constructed member over the permitted base frames. There is also a condition \(\mathcal E\) on bundles such that
\begin{equation}
 \text{eventually every bundle linked to a good member satisfying }R\text{ satisfies }\mathcal E.
 \tag{8.34a}\label{eq:p8c3-linked}
\end{equation}
The same \(R\) and \(\mathcal E\) work for every family class \(\mathcal A\) and every compatible selection \(\sigma\) of such linked bundles satisfying \(\mathcal E\). For each such selection the following estimate is closed on \(W_3(\mathcal A)\) in the sense of \ref{B8.39}. For every eligible \(F_m,C_H\) there is \(C_0\ge0\) such that for each \(C\ge C_0\) there is \(x_0\) for which, at every \(x>x_1\) with \(x\ge x_0\), every selected index \(\iota\in I_x\) and every layer,
\begin{equation}
 \frac{\operatorname{lcv}(X_{W,i},\Lambda_i)}{D_{\rm norm}}\le M_{p,i}+CE_p
 \tag{8.34b}\label{eq:p8c3-selected}
\end{equation}
in the selected indefinite realisation.
\end{theorem}
\begin{proof}[Outline]
Intersect the open rows of \ref{B8.08} to obtain a common witness and identify the selected coefficient field. Apply the degree and root discriminant bounds of \ref{B8.34a}, then use \ref{B8.24} on the selected indefinite realization. The complete proof is in \S\ref{prf:B8.34}.
\end{proof}

The grid-strengthened height-counting extension of \ref{B8.01} gives the covolume row of \ref{B8.03} by \ref{B8.34}. On the selected family, \ref{B8.40a} bounds the local-error sum by \eqref{eq:B8-divinN} and bounds the weighted lengths of the lifts modulo their cuts. Proposition~\ref{B8.38b} bounds the sum of the indices at conjugate twist places by \eqref{eq:B8-paired-index}. Its error terms are nonnegative by \eqref{eq:B8-paired-positive}, and their weighted sum together with the cut exponents satisfies \eqref{eq:B8-paired-error} in every layer. Together with the Gram bound and the pairing-index and prime-square bounds these give the inputs of \ref{B5.08x}. Use \(l_{3,i}=M_{p,i}+CE_p\) from \eqref{eq:p8c3-selected} in \ref{B5.09}, with the passage budget of \ref{B5.09a}. Use \ref{B5.10} to identify the lattice and its metric and pairing. Its numerical bound and the identities \(H^*-\bar\kappa=h_u\) and \(2F_*=\Fo\) give the wordwise inequality
\[
 \mu(RR^\vee)+h_u-\Fo\le S_{\rm div}-S_a\log T_u+E_{\rm div}
\]
of \ref{hyp:p5c4-div}.






% ---- p8c4 ----
\chapter{The definite side and local errors}\label{ch:B8-definite-local}

Use the data of \ref{B8.02}. Write \(\mathcal L\) for the size datum, \(K\) for the coefficient field of the member, \(L=x^\varepsilon\), \(b=e^L\), and \(a=xb\). The definite calculation bounds point evaluation on the divided columns. The local calculation bounds the indices lost in projecting to the good Hecke eigenspaces.

\section{The definite column condition}

\begin{definition}[Numerical conditions on a word]\label{def:B8-word-numbers}
Use \(N=N_m\), \(t=t_{\rm ch}\), \(L_a=\log(a\,\mathrm{rd}E)\), and \(L'=\log(b\,\mathrm{rd}E\,g_0g_1(k+m-1+t))\). Write \(w''=b'/(x|A|)\) for the thin weight floor. Use the row entries \(\nu_{u,j,\rho,i}\) of \ref{B4.03}, the weights of \ref{B4.02b}, and \(\log\Omega\) of \ref{B5.07}. The numerical conditions with constants \(N_c,\lambda_0,c_0,c_e,c_h,S_w,C_1\) are as follows:
\[
 \begin{gathered}
 \log\mathrm{rd}E\le L,\quad \log x\le L,\quad
 N_c\ge0,\quad |\log\Omega|\le L_a,\quad 0\le\lambda_0,\quad m\lambda_0\le b,\\
 N^2\le b,\quad m^3N^2L\le b,\quad m\eta NL\le b,\quad kN\le b,\\
 kS_c\le w'',\quad m(t+1)\eta\le w'',\quad Nw''\le b,\quad
 b'N\le a,\quad mkb'\le a,\\
 1\le c_0\le c_e,\quad 0\le c_h\le c_0,\quad c_hL'\le c_0L_a,\quad
 2\le C_1,\quad 2S_w\le C_1S_c,\quad C_1+1\le c_0 .
 \end{gathered}
\]
For each \(1\le j\le m\) and type \(\rho\) require
\[
 \begin{gathered}
 \tfrac12(k+(2j-1)(t+1))\le w_j(\rho)\le 2b,\qquad
 w_j(\rho)+S_w\le a\,\mathrm{rd}E,\\
 a_s+w_j(0,s)\le a\,\mathrm{rd}E\quad(s\in TT),\\
 \nu_{u,j,\rho,i}\ge0\quad(i<N_{j-1}),\qquad
 \sum_{k\le i<N_{j-1}}\nu_{u,j,\rho,i}\le m(t+1)\eta .
 \end{gathered}
\]
Let \(U_i\) denote the initial column coordinates of the word at \(\rho\), and \(U_i^0\) the reference coordinates. Require \(U_i>v_i>U_{i+1}\) for \(0\le i<k-1\). At an unmarked type require
\[
 U_{k-1}^0+(k-1)/2\ge0,\qquad U_0^0-(k-1)/2\le S_w .
\]
At a marked type \((0,s)\) require \(a_s,b_s\ge0\),
\(v_{k-2}+(k-2)/2\ge0\), and \(v_0-(k-2)/2\le S_w\).
\end{definition}

\begin{definition}[Definite rows for the selected bundle]\label{B8.35}
Fix \(D,O,\mathbf b,c_{\rm sz},C,C_H\). For every choice of level-one primitive linkage, separation, covolume, surjectivity, isometry and passage data, form the reference and selected sides and the definite face \(D_c\). Take \(\lambda_0,c_0,c_e,S_w,C_1\in\R\) satisfying the wordwise numerical conditions of Lemma~\ref{B5.08b} for every word, with the member's \(N_c\), \(C_0=c_0\) and \(C_h=c_{\rm sz}\). Let \(d(z)\) be the dimension of column \(z\) of the definite face. Quantify over every family of \(\mathcal O_K\)-submodules \(U_z\subset K^{d(z)}\) and every division datum \(I\) that agrees, after transport to these sides and face, with the division datum of the selected bundle.

The division ties identify the length and ranks \(m,r,t\), finite group and its idempotent coefficients, tuple labels, infinite places and finite-place selectors. Write \(e_G\) for the group isomorphism and \(P_0,P_s\) for the reference and selected subgroups. Require \(|P_0|=|P_s|(g_0g_1)^{kf}\). The map \(e_G\) identifies \(P_0\) with the level subgroup. For every \(u\in P_0\), twist index \(i\), and coefficient embedding \(b_0\in\{0,1\}\), there is an element \(z\) of the common auxiliary coefficient field \(K_{\rm aux}\) whose image in \(K\) is \(\psi_0(u)^{-1}\) and whose image in the \((i,b_0)\) local coefficient field is \(\psi_{U,i,b_0}(e_Gu)\). In the residue-field matrix models the reference subgroup and its character agree with their level models. On the selected side put \(s=\max(\lfloor m/2\rfloor-1,0)\). In zero-based indices, read entry \((i,j)\) of the level model at \((\sigma(i),\sigma(j))\) of the selected matrix, where \(\sigma(i)=k+i\) for \(i<s\), \(\sigma(i)=i-s\) for \(s\le i<s+k\), and \(\sigma(i)=i\) thereafter, with values taken modulo \(r\). This identifies both the subgroup and its character.

The row maps identify the Hermitian lattices, forms, idempotents and integral tuple modules and commute with projection to the common row space. The reference layer maps identify the layer forms and lattices as well as the zero-summand spans. They commute with projected theta. The theta frames are bases of their complex spans and identify the integral layer lattices. The two sides have the same \(I_j\) and \(j_w\). Their local errors \(\eta_J\) agree on the twist set and vanish off it. The assignment \(u\mapsto w_U(u)\) from words to columns in the division datum is bijective. For each word \(u\), the column isomorphism \(e_u\) from \(K^{d(w_U(u))}\) to the ambient space of the divided lattice carries \(U_{w_U(u)}\) onto the original integral column module and preserves the infinite-place forms. The pairings \(B\) agree under the row and column maps.

The definite column condition requires the following lattice, factorization and norm bounds for every such choice and every word \(u\). Form the numerical table \(R_u\) from the rows of \(u\), padded to length \(N_m\), its level \(T_u\), and the remainders of Lemma~\ref{B5.08c}. Put \(w=w_U(u)\), \(R=R_u\), \(T=T_u\), \(N=N_m\), \(L_a=\log(a\,\operatorname{rd}E)\), and \(d_0=d(w)\). Let \(N_u^0\) denote the input submodule of word \(u\) before division, as in \ref{B4.39}. Require \(d_0>0\) and that \(e_u^{-1}(N_u^0)\) be a full lattice in \(K^{d_0}\).

Let \(J\) be the layers and \(\mathrm{Pl}\) the types. For the degrees \(\nu_{\ell,\tau,i}\) put \(n_{\ell,\tau}=\sum_i\nu_{\ell,\tau,i}\) and \(p_\ell=|\mathrm{Pl}|g_\ell^0\). The height table has volume parameter \(V\), log periods \(\omega_{\tau,v}\), log theta norm factors \(\vartheta_v\), block bounds \(M_{\ell,v}\), and corrections \(h_\ell,h'_\ell\). The bars denote means over infinite places.
\begin{align}
 \bar\omega_\tau&=\log\Omega,\qquad \bar\vartheta\le C_\theta L,\qquad g_\ast\le C_8L,\notag\\
 M_{\ell,v}&\le C_2fm^2N^2(g_\ast+\log T)+C_efN^2g_\ast,\notag\\
 h_\ell,h'_\ell&\le C_4fN^2(L+\log T+g_\ast),\qquad 0\le g_\ell^0\le N^2,\notag\\
 \log(8+32p_\ell V)&\le c_5L_a,\qquad
 \log(2n_{\ell,\tau}+2)\le c_6L_a .\label{eq:B8-kernel-table}
\end{align}
Write \(S,\eta,b',a_\tau,TT,\log\Omega\) for the size, loss, auxiliary weight, top weights, top set and mean log period of \(R\). The terminal table, with leading weights \(\lambda_\tau\), gaps \(D_{0,\tau},D_{1,\tau}\) and evaluation allowance \(e_{\rm ev}\), satisfies
\begin{align}
 A&\le-\sum_\tau n_\tau\log T+\sum_\tau(2c_\tau^*+\xi_\tau^*)+e_{\rm ev},\notag\\
 c_\tau^*&\le(N-1)D_{1,\tau}\log(2a),\qquad
 \xi_\tau^*\le\tfrac12 N(N-1)\log(2a),\notag\\
 n_\tau&=\lambda_\tau-D_{0,\tau},\quad D_{1,\tau}\le3D_{0,\tau}+S+\eta,\quad
 D_{0,\tau}\le S+m\eta+\mathbf1_{TT}(\tau)b_\tau,\notag\\
 b_\tau&\le2b'\quad(\tau\in TT),\qquad
 \lambda_\tau=a_\tau-\tfrac{N-k}{2}\quad(\tau\in TT),\qquad
 \lambda_\tau\ge S-\tfrac{N-1}{2}\quad(\tau\notin TT),\notag\\
 0&\le\log T\le C_TmL,\qquad e_{\rm ev}\le C_{\rm ev}fN^2L_a .
 \label{eq:B8-terminal-table}
\end{align}
The bounding constants are nonnegative, and
\begin{equation}\label{eq:B8-table-constants}
 C_h\ge\max\{2(16+C_T+C_{\rm ev}),\
 C_\theta+c_5+2c_6+1+2C_4+2C_eC_8,\
 2(C_2+C_4)(C_T+C_8)\}.
\end{equation}
Use \(\log0=0\) and \(\log z=\log|z|\) for real \(z\ne0\). Put
\begin{align}
 s_{\ell,v}
 &=\sum_\tau\left[\sum_i(\nu_{\ell,\tau,i}\log\nu_{\ell,\tau,i}-\nu_{\ell,\tau,i})
 -n_{\ell,\tau}\omega_{\tau,v}+1+2g_\ell^0\log(2n_{\ell,\tau}+2)\right]\notag\\
 &\quad+p_\ell\{\vartheta_v+\log(8+32p_\ell V)\}+2M_{\ell,v}+h_\ell+h'_\ell .
 \label{eq:B8-height-size}
\end{align}

Choose a matrix chain \(K^{d_0}\xrightarrow{K_0}\cdots\xrightarrow{K_{s-1}}K^{d_s}\) for the ordinary theta steps and the complementary step. Write \(f_0\) for the initial inclusion of the column into the ramified form space. At the infinite place \(v\), use frames \(X_{\alpha,v}\in\GL_{d_\alpha}(\C)\), taking the given Hermitian frame of column \(w\) at the initial stage. Assign each layer to a step by a map \(\iota:J\to\{0,\ldots,s-1\}\) and a terminal map \(P_{\rm end}:K^{d_s}\to(K^{n_0})^\vee\), where \(n_0\) is the dimension of the reference base form. Write \(P_{w,\mu_0}(\phi,F)\) for the point-evaluation pairing of column \(w\) with the reference vector \(F\). Require
\begin{align}
 P_{w,\mu_0}(\phi,F)&=P_{\rm end}(K_{s-1}\cdots K_0\phi)(F),\label{eq:B8-column-factor}\\
 (Y_v^{-1})^*Y_v^{-1}&=v(H_{\mu_0}),\qquad Y_v\in\GL_{n_0}(\C).\label{eq:B8-definite-frame}
\end{align}
The coefficient maps \(\Phi_{\alpha p,v}\) to normed complex spaces and the continuous linear transitions \(T_{\alpha pq,v}\) satisfy
\begin{equation}\label{eq:B8-coefficient-rules}
\begin{gathered}
 0<\gamma_{\alpha p}\le1,\quad
 \sum_p\gamma_{\alpha p}\|\Phi_{\alpha p,v}(z)\|^2=\|X_{\alpha,v}^{-1}z\|_2^2,\qquad
 t_{\max}(\alpha,v)\ge0,\\
 \Phi_{\alpha+1,q,v}(v(K_\alpha)z)
 =\sum_p\gamma_{\alpha p}\epsilon_{\alpha pq,v}T_{\alpha pq,v}\Phi_{\alpha p,v}(z),
 \qquad |\epsilon_{\alpha pq,v}|=1,\quad
 \|T_{\alpha pq,v}\|^2\le t_{\max}(\alpha,v).
\end{gathered}
\end{equation}
If \(a_\alpha\) counts coefficient groups and \(h^{\rm cl}_\alpha\) bounds their sizes, write \(B_{\alpha,v}=a_\alpha a_{\alpha+1}h^{\rm cl}_\alpha h^{\rm cl}_{\alpha+1}t_{\max}(\alpha,v)\) for the squared transition bound. The required bounds are
\begin{align}
 \|P_{\rm end}\|_{X_{s,v}\to Y_v^\vee}^2&\le e^A,\label{eq:B8-terminal-bound}\\
 \log B_{\alpha,v}&\le\sum_{\iota(\ell)=\alpha}s_{\ell,v}.
 \label{eq:B8-layer-bound}
\end{align}
\end{definition}

\begin{definition}[The size conditions for the height rows]\label{B8.36a}
For a member with a choice of the window and weight parameters of \ref{B3.53}, let \(\mathcal H(c)\) denote the following conditions. With \(D_1=\Re(L'(1,\chi_d)/L(1,\chi_d))\), require
\[
 \begin{gathered}
 \log\mathrm{rd}E\le L,\qquad |D_1|\le L,\qquad \log(4x)\le L,\qquad
 \log g_i\le cL+1\quad(i=0,1),\\
 4N_m^3\le\eta,\quad t_{\rm ch}\le\eta,\quad
 m(t_{\rm ch}+1)\eta\le w'',\quad kS_c\le w'',\quad
 b_s\le2b'\quad(s\in TT).
 \end{gathered}
\]
The spacing \(\Delta\) of the interior word coordinates must be a nonnegative integer. At each thick type \(\tau\) one requires \(w_j(\tau)\le b+\eta+1\) for \(1\le j\le m\).
\end{definition}

\begin{lemma}[Size conditions for the height table]\label{B8.35a}
For a member with its window size datum, the height size conditions at \(Y_0\in\mathbb R\) require an admissible choice of the window parameters on that same datum, as in \ref{B8.05}, and the following inequalities. Write \(N=N_m\), \(t=t_{\rm ch}\), \(r_d=\operatorname{rd}E\), and let \(D_1\) be the quadratic logarithmic derivative in \ref{B5.07}. Then
\[
\begin{gathered}
 \log r_d\le L,\qquad |D_1|\le L,\qquad \log(4x)\le L,\qquad
 \log g_i\le Y_0L+1\quad(i=0,1),\\
 4N^3\le\eta,\qquad t\le\eta,\qquad
 m(t+1)\eta\le w'',\qquad kS_c\le w'',\qquad
 b_s\le2b'\quad(s\in TT).
\end{gathered}
\]
The spacing \(\delta\) must be a natural number, and \(w_{j\rho}\le b+\eta+1\) is required at every thick \(\rho\) and \(1\le j\le m\).
For the selected family these conditions hold at \(Y_0=C_{YL}\) for all sufficiently large indexed \(x\), uniformly in its members.
\end{lemma}

\begin{proof}[Outline]
Use the parameters of \ref{B8.05}, host bounds from \ref{B2.56} and \ref{B2.55}, and prime bounds from \ref{B2.51}. Apply \ref{B8.59} to node sizes and \ref{B4.02b} to \(w''\), then use its ceiling, the box in \ref{B3.50}, and the endpoint in \ref{B3.47}. The complete proof is in \S\ref{prf:B8.35a}.
\end{proof}

\begin{definition}[Column comparisons and kernel size bounds]\label{def:B8-column-comparisons}
Use the coefficient embedding \(\iota:K\hookrightarrow\C\). The real normalizations in Theorem~\ref{B8.09}(5) identify the scaled Haar measure of each signed unitary matrix group with the measure on \(\U(P,Q)\) obtained by ordering its positive and negative coordinates. At each type \(\tau\), the specified discrete series exists with the prescribed formal degree for this measure, and every discrete series with these parameters has that formal degree. For every Hermitian space of the specified signatures the assigned adelic volume is the measure of a fundamental domain for the rational unitary group. The symmetric column forms have the following level-zero property. For each word \(u\) there is an integer \(h_u>0\) such that every value of its level-zero character belongs to \(\iota(K)\) and has \(h_u\)-th power one.

For the point-evaluation comparison, identify the finite row group with the product of the residue-field matrix groups. Identify the selected subgroup and character with the selected level and its character, and the reference pair with the subgroup and character formed from the column additive characters. At a type prime these additive characters are nontrivial. At places over the same \(g_i\) their values on integers agree. The layer characters there are distinct and, in the residue-field coordinates, are the powers \(2j\), \(1\le j\le m\), of one faithful character of \(\F_{g_i}^{\times}\). Under the group identification the column character set consists of the absolutely irreducible \(K\)-valued characters \(\chi\) for which
\[
 \sum_g e_0(g)\,\iota(\chi(g))=\iota(\chi(1)),
\]
where \(e_0(g)\) is the coefficient of \(g\) in the finite-type idempotent.

The column Harish--Chandra parameters agree with those of the member at every word and type. Their row parameters and the two endpoints of the initial column are independent of the word. At each step \(i<m\) there are integral determinant shifts and nonnegative shifts of the polynomial weights. Adding the latter to each coordinate of the natural polynomial weight gives the prescribed jet degree. The current determinant shift plus the character shift on the sharp side equals the nonnegative shift, and the next determinant shift is the negative of the divided-side character shift. These choices work for every tuple, and the adjacent-weight jet belongs to the specified jet space over \(K\).

For every tuple, every intermediate index at \(1\le j\le m-2\), and every selected auxiliary place \(v\) away from the type primes, the local transformation must give the same image of the layer lattice as a matrix of the following form. Write \(N_j\) for the old rank, \(c\) for the number of added coordinates, and \(q_v\) for the residue cardinality. The old diagonal block is the identity, the new diagonal entries are \(q_v^{e_h}\) with \(-j\le e_h\le0\), and the only possible off-diagonal entries form the vector \(q_v^{e_1}z\) above the first new diagonal entry, with \((z,0)\) in the layer lattice. At each word there must also be a complex linear isomorphism from its row space to the complex row model. For a reference tuple with empty auxiliary set, every projected row-coordinate vector has a lift to a level-equivariant form with the prescribed restriction. The isomorphism takes its row image to the vector of values of that form. On the selected subspace the specified column evaluation itself is level-equivariant and has these same values.

Use the column weights for the kernel. Its ranks, host fields, CM type, conductor and primes must agree with those of the member, with the given unramified abelian extension, prime congruences and conjugation-compatible coefficient embedding. A single bijection of archimedean types identifies the marked positions and, for every word, the row parameters and word coordinates. The global and local characters have the prescribed infinity types and Weil normalizations, and the selector chooses the specified place over each base place. The analytic data use the discriminant twelfth root of the ideal lattice, its ideal character and generators of the period ideals over a common number field. The reference tuple gives the trivial character and the sharp level. The kernel uses \(m-2\) intermediate steps, with natural-number subtraction, and the jets for each word at adjacent ranks are defined over the prescribed ambient coefficient field at the selected places and level.

The \emph{kernel size bounds at \((Y,U)\)} require
\[
  2fk<g_i-1,\qquad 2fN_m+1<4a,\qquad
  4a<T_u,\qquad \log T_u\le UmL .
\]
At every word, layer \(0\le j\le m\) and type, the difference between the largest and smallest true weights is at most \(2a\). For each word the numerical rows must also satisfy
\[
 \max\{\log\operatorname{rd}E,\log q,\log d,\log(4x),|D_1|\}\le L,\qquad
 (kb)^8\le g_i,\quad \log g_i\le YL+1,\quad fN_m\le b .
\]
Every degree in a layer row, up to its unpadded length \(N_{j-1}\), lies in \([0,3a]\). At each type, every difference between a coordinate of the base tuple \(b\) and a coordinate of the layer tuple \(e\) has absolute value at most \(\eta\); here \(b\) denotes the tuple of Harish--Chandra coordinates. All reference \(v\)-coordinates and all interior word \(U\)-coordinates lie in \([0,S_c]\). At a marked position require \(b'\le b_s\le2b'\). At an unmarked position, for \(k\ge2\), the reference endpoints satisfy
\[
 U^0_{k-1}=v_{k-2}-\tfrac12,\qquad S_c\le U^0_0-v_0,\qquad U^0_0\le3S_c .
\]
\end{definition}

\begin{theorem}[Columns from point evaluation]\label{B8.36}
For \(C_Y,C_T\in\R\) with \(C_T\ge0\), there is a threshold \(C_0\) uniform in \(D,O,c_{\rm sz},\mathbf b,C,C_H\). Assume the real normalizations and symmetric column forms of Theorem~\ref{B8.09}(5), with the point-evaluation identifications and compatible kernel weights of Definition~\ref{def:B8-column-comparisons}. Assume that for all \(C_Y',C_T'\) the five conditions
\begin{equation}\label{eq:B8-five-size}
 \text{the height size conditions of Lemma~\ref{B8.35a} at }C_Y',\quad
 X'=\exp(x^{\varepsilon/3})\ge8,\quad x\ge2,\quad
 \forall u\ (4a<T_u),\quad \forall u\ (\log T_u\le C_T'mL)
\end{equation}
imply the kernel size bounds of Definition~\ref{def:B8-column-comparisons} at \((C_Y',C_T')\). Suppose they hold at \(C_Y,C_T\), and
\begin{equation}\label{eq:B8-divided-lattice}
 N'_u=I_uN_u^0
\end{equation}
is a full lattice for every word, with \(\operatorname{rank}\nu_u>0\). If \(c_{\rm sz}\ge C_0\), the selected bundle satisfies the definite column condition of Definition~\ref{B8.35} at \(C,C_H\).
\end{theorem}
\begin{proof}[Outline]
Factor column periods in Hermitian frames. Use \ref{B5.07c} for evaluation, Cauchy–Schwarz for kernels, \ref{B5.07x} for mean volume, and \ref{def:B8-column-comparisons} for the threshold, and transport the lattice obtained from a multiple of \(\mathrm{NuLat}_u\). The complete proof is in \S\ref{prf:B8.36}.
\end{proof}

\section{A common bound along the selected family}

\begin{theorem}[The definite condition along members]\label{B8.37}
Fix any local Langlands witness \(h_L\) and any member parameters \(p\) for that witness. For every family class \(\mathcal A\) and every selection linked to \(p\), the definite column condition of Definition~\ref{B8.35} is closed on \(W_3(\mathcal A)\) in the sense of Definition~\ref{B8.39}. Write \(x_1\) for the family's selection cutoff and \(I_x\) for its indices at \(x\), as in Definition~\ref{B8.02}. For every \(C\ge0\) the condition holds eventually for all \(x>x_1\) and \(i\in I_x\). Assume the real second-layer identifications, the selected-family design of \ref{B8.01} with constants \(3,C_H'\), and, beyond one cutoff, \(b^{m/2}\le T_u\) and \(\log T_u\le C_{Tu}mL\) for every member and word. These hypotheses imply eventual fullness of \(N'_u\) and positivity of \(\operatorname{rank}\nu_u\) for every word.
\end{theorem}
\begin{proof}[Outline]
Linkage supplies evaluation and size, and \ref{B5.05b} gives the conditions of \ref{B5.05a}. Clear denominators for fullness, use the lower rank bound for positivity, and apply \ref{B8.36} with size and level bounds from \ref{B8.35a} and \ref{B8.51}. The complete proof is in \S\ref{prf:B8.37}.
\end{proof}

\begin{definition}[Member points for the reference estimates]\label{def:B8-reference-points}
Use the class \(\mathcal P\) defined by the three member-point conditions in Definition~\ref{def:p8c3-admissible}, before its two additional arithmetic requirements. In the reference coordinates these conditions have the following description. The indefinite realization has the layer lattices, inclusions and theta maps of the primitive construction. At good coefficient embeddings the evaluation at the chosen vector is injective; at the other embeddings the conjugate evaluation identity holds. In every layer the action has a geometric realization, an extremal vector pairs to the prescribed delta vector, and the leaf map has the prescribed values and phase. The primitive row types are separated and satisfy the level-one conditions. The frames and the division maps identify the forms, lattices, projectors and pairings of the two realizations as in Definition~\ref{B8.35}. The group map respects the subgroup characters, and the dimension of the reference projected space is \(D_0\). There are linear maps from the arithmetic layer spaces to the exported spaces which agree in layer coordinates and give the prescribed maps into the glued lattice. All finite types are good and every single-index place is twisted.

These identifications must come from the same recursion and spectral realization. The core and the realization use the same base prime and tame types. The exported series agree, the reference character and the local level character have the common coefficient representatives of Definition~\ref{B8.35}, and the exported sets contain every datum transforming by the prescribed level character. At a nonreference place the additive character takes its level-zero values in the embedded coefficient field. The auxiliary field contains the required ray values and the roots of unity of order dividing \(g_i(g_i-1)\). Each tuple field contains the ray field of its level, and the twists agree with the chosen norm.

Let \(\mathcal R\) consist of points satisfying the following extra requirements for the reference calculation. Every single-index place is twisted. Form the reference realization using the numerical data of the point for the two runs. Its core obeys the theta identities with the canonical levels, positive signature rank in every layer and the prescribed half-error estimates. The coefficient-field conditions hold at both recursion primes and the small-field bound holds for the member. The level ideal is the product of the designated prime ideals. For each residue characteristic of a twisted place there is a prime of the intermediate field below every selected prime for that place and tuple. Both determinant characters take values in the designated coefficient field. In this realization the series, inverse reference character and level eigendata agree with the exports. The additive characters agree with the norm choice and have coefficient-field values at the nonreference places. The small field has the prescribed valuations, ray values and roots of unity; the tame types and the ray fields of all tuple levels agree with the recursion.
\end{definition}

\begin{theorem}[One constant for both realisations]\label{B8.38}
Assume the discrete-series, Harish--Chandra, rapid-decay and Casselman--Wallach inputs and the surjectivity and lower-weight rows. At the linked slot assume the primitive, divisor, goodness, geometric compatibility, third-layer covolume, functional, multiplicity, \(\kappa\) and defect rows. For a class \(\mathcal A\) and raw selection satisfying the slot member condition, the following assertion is closed on \(W_3(\mathcal A)\). For every \(F_m,C_H\) in the closed assertion there is \(C\ge0\) such that the following bounds hold eventually at every selected member above \(x_1\). On the indefinite side every single-index place is twisted. The reference covolume bound, the isometry, multiplicity and passage identities, and the \(\kappa\)-bounds hold. Under the division estimate \eqref{eq:B8-divinN} one also has the paired-index and functional bounds. If each layer has the prescribed single leaf map with its fixed values and phase, the selected-side lower weight bound holds. On the definite side the restricted bundle satisfies Definition~\ref{B8.35} at the family's size constant and the same \(C,C_H\).
\end{theorem}

\begin{proof}[Outline]
Use Miller–Schmid decay and \ref{B8.38a} to extend the intertwiners. Use \ref{B8.02} to transport the data for \ref{B8.24}, carry over \ref{B8.38b} by linkage, and apply \ref{B8.39} on \(W_3(\mathcal A)\) to the eight rows, including definiteness from \ref{B8.37}. The complete proof is in \S\ref{prf:B8.38}.
\end{proof}

\begin{proposition}[The paired indices in the reference row]\label{B8.38b}
In the reference realization let \(D_R>0\) be the common rank and \(\mathrm{Tw}\) the twist set. Write \(T_{j,t_w}\) for the selected type lattice, \(\theta_{j,t_w}\) for its layer map, and \(A_{j,t_w}=\theta_{j,t_w}(T_{j,t_w})\) for its image in \(K M_j\). Use the signed index of the two commensurable lattices as in \ref{B5.08a}. Write \(p_w\) for the residue characteristic, \(e_w\) for the ramification index, and \(\ell_w=\log N\mathfrak p_w/[K:\Q]\). Put
\[
 i_{jw}=[M_j:A_{j,t_w}]_w+[M_j:A_{j,t_w}]_{\bar w},\qquad
 c_w=2e_w/p_w^2,\qquad
 R_w=\left(D_R^{-1}\sum_j i_{jw}-2e_w\sum_j F_{0,j,w}-c_w\right)_+ .
\]
Here indices are positive for proper sublattices and \(F_{0,j,w}\) is from \ref{B4.02}. Use the stored error \(R_w/(e_wm)\), so \(e'_{jw}=R_w/m\) contributes to the index. Let \(J_{0,j,w}\) be the least normalized exponent carrying the lift into its cut, \(B_0=fmbL\), and \(N_c\) the size budget of \ref{B8.40}. Under the two-run, rationality, geometry, rank, twist, share, bad-place, exponent and page hypotheses, with the quantitative inputs below, the family division estimate \eqref{eq:B8-divinN} gives the following bounds eventually, uniformly for \(C\) above a common threshold:
\begin{align}
 D_R^{-1}\sum_j i_{jw}&\le\sum_j(2e_wF_{0,j,w}+e'_{jw})+c_w\quad(w\in\mathrm{Tw}),\label{eq:B8-paired-index}\\
 e'_{jw}&\ge0,\label{eq:B8-paired-positive}\\
 \sum_{w\in\mathrm{Tw}}\ell_w(e'_{jw}+e_wJ_{0,j,w})
 &\le C L(B_0+N_c)\quad(1\le j\le m).\label{eq:B8-paired-error}
\end{align}
\end{proposition}

\begin{proof}[Outline]
Apply the two-run remainder bound \ref{B4.52} and the projection estimate \ref{B4.54a} to the lattice in \ref{B5.08a}. Pair conjugate indices, subtract the base terms of \ref{B4.02}, and bound the averaged positive excess and cut exponents by weighted prime sums. The complete proof is in \S\ref{prf:B8.38b}.
\end{proof}

\begin{proposition}[Sobolev bounds for unitary intertwiners]\label{B8.38a}
Let \(r,s\) be nonnegative integers, \(G=\U(r,s)\) and \(K=\U(r)\times\U(s)\). Let \((\sigma,H)\) be irreducible strongly continuous unitary, with finite vectors \(V\) and smooth vectors \(H^\infty\). Let \((\pi,E)\) be smooth complex Fr\'echet with continuous derived operators and defining seminorms \(p_i\) satisfying
\begin{equation}\label{eq:B8-moderate-growth}
 p_i(\pi(g)z)\le h(g)^{N_i}b_i(z),\qquad h(g)=1+\operatorname{tr}(g^*g)
\end{equation}
for continuous seminorms \(b_i\) and \(N_i\in\N\). A linear map \(T_0:H\to E\) intertwining \(K,\mathfrak g=\operatorname{Lie}(G)\) on \(V\) has a unique continuous equivariant extension from \(V\) to \(H^\infty\), with closed image. For \(T_0|_V\ne0\) and each finite list \(S\) of derived words, a continuous seminorm satisfies
\begin{equation}\label{eq:B8-sobolev}
 \sum_{D\in S}\|d\sigma(D)w\|\le q_S(Tw)\qquad(w\in H^\infty).
\end{equation}
The empty word acts by the identity. The map on \(H^\infty\) extends algebraically to \(H\).
\end{proposition}
\begin{proof}[Outline]
Schwartz integration and Casselman–Wallach globalization extend the map. Hahn–Banach bounds coefficients on compact Casimir eigenspaces, and the quotient topology and Casimir Sobolev estimate give Hilbert bounds for derived vectors, with closed image by completeness. The complete proof is in \S\ref{prf:B8.38a}.
\end{proof}

\section{The quantifiers and the local error}

\begin{definition}[Closed assertions on family classes]\label{B8.39}
A threshold datum \(\mathcal K\) has a real pre-threshold and three real functions of lists \(\mathcal K_t,\mathcal K_{\rm mid},\mathcal K_{\rm last}\). Write \(\mathbf o\) for the family's outer list of constants and \(\mathbf m(C_H)\) for its middle list. Say that the family meets the thresholds at \(C_H\) if every preliminary entry is at least \(\mathcal K_{\rm pre}\), every middle entry is at least \(\mathcal K_{\rm mid}(\mathbf o)\), and the following last-entry bound holds. Each last entry is at least \(\mathcal K_{\rm last}(\mathbf o\mathbin{\|}\mathbf m(C_H))\), where \(\|\) denotes concatenation. Define
\[
 P\text{ is closed on }\mathcal C\ \Longleftrightarrow\
 \exists\varepsilon_1>0,\ \forall\,0<\varepsilon\le\varepsilon_1,\ \exists\mathcal K,\
 \forall\, F_m\text{ of parameter }\varepsilon,\ C_H\in\R,\quad
 F_m\in\mathcal C\text{ and }F_m\text{ meets }\mathcal K\text{ at }C_H\ \Longrightarrow\ P(F_m,C_H).
\]
For the closed local-error assertion, take \(P(F_m,C_H)\) to be the family estimate in Definition~\ref{B8.40}. Conjunction uses the smaller cutoff and the componentwise maximum of the thresholds.
\end{definition}

\begin{definition}[The local-error estimate]\label{B8.40}
At index \(i\), let \(O_i\) be the local coefficient ring and \(\ell_i\) its residue characteristic. Write \(\lambda_i=\operatorname{length}_{O_i}/\operatorname{length}_{O_i}(O_i/\ell_iO_i)\), with infinite length and division by zero assigned value zero. For the lattice image \(M_i\) and the eigenspace sums \(V_Z,V_{Z^c}\) of \(Z=Z_{\rm good}(i)\) and its complement, put
\[
 U_i=V_Z\cap(M_i+V_{Z^c}),\quad L_i=M_i\cap V_Z,\quad
 \eta_j(w)=\sum_{i\in I_{j,w}}\lambda_i(U_i/L_i),\quad
 \eta_j(\mathcal L)=\sum_{\substack{w\in S_{\mathcal L}\\\mu(w)\ne\mu_0}}
 \ell_we_w\eta_j(w).
\]
Here \(S_{\mathcal L}\) is the finite support of the size datum, \(I_{j,w}\) is the index set of layer \(j\) at \(w\), and \(\ell_w=\log N\mathfrak p_w/[K:\Q]\).
The finiteness hypothesis requires the lengths \(s_{\ell,e}(E)=\operatorname{length}_{\Z_\ell}\Sha_{\ell,e}(E)\) of Definition~\ref{B6.01}, taken with the member's local data, to be finite for every prime \(\ell\), including \(2\), and every \(1\le e<t\). Under this hypothesis put \(s_\ell=\sum_{e=1}^{t-1}s_{\ell,e}(E)\) at each prime
and \(v_\ell(D)=(\log\ell)s_\ell\), with \(v_\ell(D)=0\) elsewhere. Write \(d_{j,0}=\dim_\C Z_j(\mu_0)\). With \(n=f/2\) and \(T_2\) the host's thick set, the size budget is \(N_c=|T_2|kbL^2/n^2\). The point condition is
\begin{equation}\label{eq:B8-linkN}
 (v_\ell(D))_{\ell\ge0}\text{ summable}\quad\Longrightarrow\quad
 \eta_j(\mathcal L)\le d_{j,0}\left[C\{fmbL^2+(\log x)^2N_c\}+\sum_{\ell\ge0}v_\ell(D)\right].
\end{equation}
The family local-error estimate requires constants \(C,x_0\in\R\) such that \eqref{eq:B8-linkN} holds for every \(x\ge x_0\), every \(i\in I_x\) whose member satisfies the finiteness hypothesis, and every \(1\le j\le m\).
\end{definition}

\begin{lemma}[The sum over twisted places]\label{B8.40b}
Use the place bijection \(\iota\) from the arithmetic coordinates to the row coordinates and the type selector of Lemma~\ref{B4.36a}. Number the arithmetic layers from \(1\), and write \(a(j)\) for the arithmetic number of row layer \(j\). Assume the division data identify the local quotients at each selected type in the two coordinates and that \(\iota^{-1}(\mathrm{Tw})\) lies in the finite support \(S_{\mathcal L}\). If \(a_{jv}\) is the sum of the ordinary local lengths of the lift quotients at the selected type at \(v\), then
\[
 \sum_{v\in\mathrm{Tw}}\ell_v a_{jv}\le\eta_{a(j)}(\mathcal L).
\]
\end{lemma}
\begin{proof}[Outline]
For \(w=\iota^{-1}v\), both coordinates share the prime, weight, and index \(e_w\), so \ref{B8.40} gives \(a_{jv}=e_w\eta_{a(j)}(w)\). Since \ref{B4.36a} picks nonreference type on \(\mathrm{Tw}\), reindexing gives a subsum of \(\eta_{a(j)}(\mathcal{L})\) with nonnegative terms. The complete proof is in \S\ref{prf:B8.40b}.
\end{proof}

\begin{lemma}[From the window bound to the division input]\label{B8.40a}\label{B8.40y}
For the host field \(E\) and the member's number \(t\) of scalar blocks, the window bound of Definition~\ref{B6.01} at budget \(B\) gives the finiteness hypothesis of Definition~\ref{B8.40}, summability, and \(\sum_\ell v_\ell(D)\le B\). Suppose \(\varepsilon>0\) and the family local-error estimate holds with constant \(C_N\). Suppose also that eventually every member satisfies the window bound at budget \(\sqrt b\) and that each \(d_{j,0}\) equals the dimension \(D_0\) of the reference target. Then eventually
\begin{equation}\label{eq:B8-divinN}
 \eta_j(\mathcal L)\le D_0 b_N,\qquad
 b_N=(|C_N|+1)L(B_0+N_c),\qquad B_0=fmbL .
\end{equation}
Thus one constant and one cutoff give the division estimate in every layer of every member beyond that cutoff. In the tied reference coordinates the lengths of lifts modulo their cuts satisfy \(\sum_{v\in\mathrm{Tw}}\ell_v a_{jv}\le D_0b_N\).
\end{lemma}
\begin{proof}[Outline]
Apply \ref{B6.01} to finite sums and use the family local-error estimate. Transfer the bound to twisted lift lengths by \ref{B8.40b} and the division tie, and use surjectivity in \ref{B8.02} to identify the layer image with the reference target. The complete proof is in \S\ref{prf:B8.40y}.
\end{proof}

\section{The primes \(2\) and \(3\)}

\begin{theorem}[Local cohomology]\label{B8.41p}
Let \(F_v/\Q_p\) be finite and \(M\) a finite continuous \(G_{F_v}\)-module killed by \(\ell^a\), with \(\ell\) prime and \(a>0\). Cohomology is finite through degree \(2\) and zero above it. Cup product pairs \(H^i(F_v,M)\) perfectly with \(H^{2-i}(F_v,\operatorname{Hom}(M,\mu_{\ell^a}))\) into \(H^2(F_v,\mu_{\ell^a})\simeq\Z/\ell^a\). Also
\begin{equation}\label{eq:B8-euler}
 |H^1(F_v,M)|=|H^0(F_v,M)||H^2(F_v,M)|
 \begin{cases}|M|^{[F_v:\Q_\ell]},&p=\ell,\\1,&p\ne\ell.\end{cases}
\end{equation}
The local finiteness theorem and Tate duality, together with the Euler characteristic formula, apply also to decomposition groups identified with local Galois groups \cite[Chapter~VII, \S\S2--3]{neukirch-2008}.
\end{theorem}

\begin{definition}[Inputs for the small-prime calculation]\label{B8.41h}
The inputs \(\mathscr H_2\) include Fontaine's crystalline period comparison and fundamental sequence \cite{fontaine-1994}, global Artin reciprocity \cite[Chapter~VIII]{neukirch-2008}, and \ref{B8.41p} at \(p=\ell\). For \(K/\Q_\ell\) finite let \(H_f^1(K,V)\) be the kernel of the period map to \(H^1(K,B_{\rm cris}\otimes V)\). Use its inverse image for a stable lattice \(T\subset V\) and its image for \(V/T\). A cocycle is finite when its class is finite. For \(V\) de Rham, the Bloch--Kato inputs \cite[\S3]{bloch-kato-1990} give
\begin{equation}\label{eq:B8-BK-dim}
 \dim_{\Q_\ell} H_f^1(K,V)=
 \dim_{\Q_\ell} H^0(K,V)+\dim_{\Q_\ell}D_{\rm dR}(V)/\mathrm{Fil}^0,
\end{equation}
and the finite conditions are exact annihilators for rational Tate duals and for an integral lattice and its divisible Tate dual. Over finite \(K/\Q_\ell\), prismatic \(F\)-crystals realize crystalline lattices, and equivariant maps lift uniquely \cite[Theorems~1.2 and~5.6]{bhatt-scholze-2023}. In particular scalar multiplication lifts to the coefficient-ring action.
\end{definition}

Let \(\mathfrak c_{i,\pi}\subset O_i\) be the evaluation at \(\pi\) of the Hecke ideal of operators vanishing off \(Z(i)\). Put \(v_{cc}(i,\pi)=\lambda_i(O_i/\mathfrak c_{i,\pi})\). Let \(s_\pi(\rho_i)\) be the normalized Selmer length of \(\rho_i(\pi)\chi_{i,1}^{-1}\) with divisible coefficients, finite above \(\ell_i\), unrestricted at \(S_i\cup Q_i\) away from \(\ell_i\), and unramified elsewhere. Write \(h_E\) for the class number and normalize \(v_\ell(\ell)=1\). Set
\begin{equation}\label{eq:B8-Bd}
 \mathcal B_{i,\pi}(\rho;B)\ \Longleftrightarrow\
 v_{cc}(i,\pi)\le s_\pi(\rho_i)+B+\sum_{1\le e<t}s_{\ell_i,e}(E)+v_{\ell_i}(h_E)+\ell_i^{-2}.
\end{equation}
Use \(r\) for the rank of \(\rho_i\), \(t\) for the number of scalar blocks, \(n_V=r+t\), \(n=f/2\), \(N=Q_Sqd\) for the host conductor modulus, and \(s_0(i)\) for the depth of the selected cyclotomic inertia element. Put \(H_B=b+S_c+\eta_{\rm sz}+t+2\).

For the small-prime estimate, take data \(X\) giving both the full Galois representation and the rank-\(r\) representation \(\rho_i(\pi)\) at each Hecke point \(\pi\). Impose the following conditions. At an inert index \(i\) the column set is empty. For every place \(w\), the local length allowance at each spectral point \(\pi\) equals the row allowance at \((i,w)\). The row allowance is the long-orbit or the other inert allowance according to \eqref{eq:B8-prime-table-depth}. At split indices it uses the canonical row and column declaration, chosen by the common embedding sign, with the canonical column length in \eqref{eq:B8-split-place-cost}. Its integral width is \(\eta_X=\lfloor\eta_{\rm sz}\rfloor\). The invariant and character exponents \(c_0,c_B\) are given below, with \(j(z)=\lfloor\log_{\ell_i} z\rfloor\) for \(z>0\), \(j(0)=0\), \(h=\lfloor2b\rfloor\) and \(a_0=4nn_V\),
\[
 c_0=5+j(2(h+1))+a_0[4+2j(2(h+1))],\qquad
 c_B=3+j(2h)+t[3+a_0(4+2j(2h))]+4t+3 .
\]
Take \((C_2,C_B,C_C,C_A,C'_B,C_K)=(C_2,320,25,23,32,32)\). Require the norm-choice conditions of Theorem~\ref{B8.42}, admissibility and the polynomial parameter bounds. Require \(L\ge8\), a host splitting cutoff of at least \(60\), the finite Selmer lengths and bounds in \eqref{eq:B8-killing} for \(\rho\), and the finiteness hypothesis of Definition~\ref{B8.40}. The size bounds used below give \(H_B\le2b\).

\begin{theorem}[A family bound at \(2\) and \(3\)]\label{B8.41}
Assume \(\mathscr H_2\). For every \(\varepsilon\in\R\) and every family \(F_m\) of parameter \(\varepsilon\) there are \(C,x_0\in\R\) with the following property. For each \(x\ge x_0\) with \(x>x_1\) and each \(a\in I_x\), use the norm choice of the selected extension at \((x,a)\). For every primewise Wach system, Fontaine datum and comparison, and every \(X,C_2\) satisfying the preceding small-prime hypotheses, one has
\begin{equation}\label{eq:B8-Bd23}
 \mathcal B_{i,\pi}(\rho;Cfmb\log x+Cx^CL^2)
\end{equation}
at every twisted \(i\) with \(\ell_i=2\) or with \(\ell_i=3\) and \(3\nmid Q_S\) for this member, for all \(\pi\in Z_{\rm good}(i)\). This gives the linked closed bound on \(W_3(\mathcal A)\).
\end{theorem}
\begin{proof}[Outline]
Insert the cyclotomic exponents of \ref{B8.41a} into the split congruence bound of \ref{B8.41b}, using inertia valuations for the character and projector terms. Use degree sums, place counts, the tame bound, and parameter bounds to fit the family allowance. The complete proof is in \S\ref{prf:B8.41}.
\end{proof}

\begin{lemma}[The Wach lattice at the lower cyclotomic roots]\label{B8.41ca}
Let \(\ell\) be prime, let \(K_0/\Q_\ell\) be finite unramified with residue field \(k'\), and put \(K'=K_0(\mu_{\ell^s})\), where \(s\ge1\) and \(s\ge2\) if \(\ell=2\). Let \(L\) be a finite extension containing \(K'\), with integer ring \(O\). Put
\[
 A=O[[Z]],\qquad X=(1+Z)^{\ell^{s-1}}-1,\qquad
 Q=\Phi_{\ell^s}(1+Z).
\]
Let \(T\) be a finite free \(O\)-lattice with continuous \(G_{K'}\)-action such that \(T[1/\ell]\) is crystalline. Its prismatic realization gives a free \(A\)-module \(N_i(T)\) of rank \(\operatorname{rank}_O T\) for each embedding \(i:W(k')\hookrightarrow O\). These modules are functorial in equivariant lattice maps. Their completed étale base changes are the components of the cyclotomic \((\varphi,\Gamma)\)-module of \(T\).

Here \(\Gamma=\Gal(K'(\mu_{\ell^\infty})/K')\). Let \(\gamma\in\Gamma\) have cyclotomic value \(a\) with \(v_\ell(a-1)=s\), where \(v_\ell(\ell)=1\). For each \(\zeta^{\ell^{s-1}}=1\), put \(t_\zeta=\log((1+Z)/\zeta)\). Over \(L[[t_\zeta]]\) the component \(N_i(T)\) has a basis fixed by \(\gamma\). In particular,
\[
 (\gamma-1)N_i(T)\subset XN_i(T),\qquad
 (\gamma-a^j)X^jN_i(T)\subset X^{j+1}N_i(T)\quad(j\ge0).
\]
These inclusions also hold on the image of any equivariant lattice map.
\end{lemma}
\begin{proof}[Outline]
Use \ref{B8.41h} to evaluate the crystal on the cyclotomic prism and split it into free components. Construct a \(\gamma\)-fixed basis at \(1\), transport it by Frobenius to the lower roots, and use Weierstrass division to obtain the \(X\)-adic inclusions. The complete proof is in \S\ref{prf:B8.41ca}.
\end{proof}

\begin{lemma}[Integral comparison of two Frobenius lattices]\label{B8.41cb}
Use the cyclotomic extension, coefficient ring \(O\) and components of Lemma~\ref{B8.41ca}. Set \(A_i=O[[Z]]\), \(A_*=\prod_i A_i\), and let \(B\) be the \(\ell\)-adic completion of \(A_*[1/Z]\). Put \(X=(1+Z)^{\ell^{s-1}}-1\) and \(Q=\varphi(X)/X\). Let \(T\) be a crystalline \(O\)-lattice and \(\mathcal D\) its étale module over \(B\). Suppose \(L,L'\subset\mathcal D\) are componentwise finite free, both base change to \(\mathcal D\), and their linearized Frobenius maps are isomorphisms after inverting \(Q\). Then there is \(a\ge0\) such that
\[
 X^aL\subset L',\qquad X^aL'\subset L.
\]
If both lattices carry the same ambient semilinear \(\gamma\), acting as the identity modulo \(Z\) on each lattice, and \(u=\chi_{\rm cyc}(\gamma)\) has depth \(s\) as in Lemma~\ref{B8.41ca}, then \(L=L'\).
\end{lemma}
\begin{proof}[Outline]
Use the étale comparison of \ref{B8.41ca} and Witt ring integrality to clear both transition matrices with a power of \(X\). Cyclotomic equivariance excludes poles at zero, and Frobenius propagates regularity to every root of \(X\). The complete proof is in \S\ref{prf:B8.41cb}.
\end{proof}

\begin{lemma}[Annihilators on a cyclotomic component]\label{B8.41c}
Let \(\ell\) be prime, \(N>0\), and \(4\mid N\) at \(\ell=2\). Set \(K'=\Q_\ell(\mu_{\operatorname{lcm}(\ell,N)})\), \(s=\max(1,v_\ell(N))\), and let \(O\) be the integers of a finite coefficient field containing \(K'\). Use the Wach components over \(K'\) of \ref{B8.41}. Their construction for the crystalline realization is given in Lemma~\ref{B8.41ca}, with \(K_0\) the maximal unramified subfield of \(K'\). The roots of unity of order prime to \(\ell\) generate an unramified extension, so \(K'=K_0(\mu_{\ell^s})\). Each unramified embedding has ring \(A=O[[Z]]\), with components in Frobenius order. Set
\[
 \varphi(Z)=(1+Z)^\ell-1,\quad X=(1+Z)^{\ell^{s-1}}-1,\quad
 Q=\sum_{j=0}^{\ell-1}(1+X)^j,\qquad \varphi(X)=XQ.
\]
Let \(\gamma\) have cyclotomic value \(u\) of depth \(s\). For a surjection \(f:T_1\to T_2\) of crystalline lattices with source jumps in \([-M_i,h_i]\) at component \(i\), \(M_i,h_i\in\N\), write \(N_{ji}\) for the Wach components and \(P_i=fN_{1i}\). For \(H=\max_i(M_i+h_i)\), \(w=2Hs+v_\ell((2H)!)\), one has
\begin{equation}\label{eq:B8-wach-coker}
 X^HN_{2i}\subset P_i,\qquad \ell^wN_{2i}\subset P_i.
\end{equation}
\end{lemma}
\begin{proof}[Outline]
Use the Wach comparisons of \ref{B8.41ca} and \ref{B8.41}, and identify the saturated image with the target by \ref{B8.41cb}. Frobenius descent gives the \(X^H\) bound, and interpolation in the cyclotomic action gives the \(\ell^w\) bound. The complete proof is in \S\ref{prf:B8.41c}.
\end{proof}

\begin{lemma}[Lifting a scalar cocycle at cyclotomic depth \(s\)]\label{B8.41aa}
Let \(K'=K_0(\mu_{\ell^s})\) as in Lemma~\ref{B8.41ca}. For \(e\ge0\), put \(b=s+v_\ell(e+1)\). If \(\bar z\) is a continuous cocycle of \(G_{K'}\) with values in \((\Q_\ell/\Z_\ell)(-e)\), there is a continuous cocycle \(z\) with values in \(\Q_\ell(-e)\) whose reduction is \(\ell^b\bar z\) at every group element.
\end{lemma}
\begin{proof}[Outline]
The cyclotomic valuation shows that \(\ell^b\) kills the dual invariants, so local duality in \ref{B8.41p} gives the same bound for the lifting obstruction. Adjust the rational lift by a coboundary to obtain pointwise agreement. The complete proof is in \S\ref{prf:B8.41aa}.
\end{proof}

\begin{lemma}[Convergence of the cyclotomic primitive]\label{B8.41ab}
Let \(O\) be the integers of a finite extension \(L/\Q_\ell\), with \(|\ell|=\ell^{-1}\). Take \(s\ge1\), with \(s\ge2\) for \(\ell=2\), and \(u\in\Z_\ell^\times\) with \(v_\ell(u-1)=s\). On \(L[[Z]]\) let \(\gamma\) and \(\varphi\) act by the substitutions \((1+Z)^u-1\) and \((1+Z)^\ell-1\), respectively. Put
\[
 X=(1+Z)^{\ell^{s-1}}-1,\qquad Q=\varphi(X)/X,\qquad t_0=\log(1+Z).
\]
On a finite cyclic set of components indexed by \(i\), suppose the matrices of two commuting semilinear operators are
\[
 A_i=\begin{pmatrix}G_i&z_i\\0&1\end{pmatrix},
 \qquad
 B_i=\begin{pmatrix}P_i/Q^m&y_i/\Lambda\\0&1\end{pmatrix},
 \qquad
 \Lambda=\prod_{j\in J}\varphi^j(Q),
\]
where \(J\) is a finite multiset of nonnegative integers, \(m\ge0\), and \(G_i,P_i,z_i,y_i\) have coefficients in \(O[[Z]]\). Assume \(A_i(0)=1\), \(B_i(0)\) is invertible over \(L\), and
\(A_i\gamma(B_i)=B_i\varphi(A_{i-1})\).
The solution with \(w_i(0)=0\) of
\[
                       G_i\gamma(w_i)-w_i=z_i
\]
is unique, and \(t_0^{m+|J|}w_i\) converges on \(|Z|<1\).
\end{lemma}
\begin{proof}[Outline]
Solve for the invariant matrix coefficient by coefficient, bounding the divisors \(1-u^j\) to get convergence on a small disc. Clear the Frobenius poles with a power of \(\log(1+Z)/X\) and iterate the contracting Frobenius substitution. The complete proof is in \S\ref{prf:B8.41ab}.
\end{proof}

\begin{lemma}[Three cyclotomic denominators]\label{B8.41a}
Use \(\mathscr H_2\) and the comparisons of \ref{B8.41}. Let \(K\subset K'\subset K_{\rm coeff}\) be local fields with \(K',s,N\) as in \ref{B8.41c}, and \(O=\mathcal O_{K_{\rm coeff}}\). Put \(j(z)=\lfloor\log_\ell z\rfloor\) for \(z>0\), \(j(0)=0\). For \(h,a,t\in\N\), set
\[
 D_h=s+j(2h)+j(N)+1,\quad E_h=32(2hs+1)+D_h,\quad
 I_{h,a}=20(a+1)D_h,\quad T_{h,a,t}=40(t+1)(a+1)D_h.
\]
For a surjection of crystalline lattices with source jumps in \([-h,h]\), \(\ell^{E_h}\) takes finite target classes into the image of finite source classes. For a crystalline lattice \(T\) with jumps in a set of at most \(a\) integers in \([-h,h]\) and zero absent, \(\ell^{I_{h,a}}\) kills \(H^0(K,T[1/\ell]/T)\).

Suppose \(t\le h\), \(Y\) is crystalline with stable lattice \(Y_0\) and jumps in such a set allowing zero, and \(p:Y\to K_{\rm coeff}\) is invariant linear with \(p(Y_0)\subset O\) and \(p(y_0)=1\) for \(y_0\in Y_0\). For \(0\le e<t\) let \(\phi:Y\to(K_{\rm coeff}/O)(-e)\) be \(O\)-linear and equivariant on \(Y_0\cap\ker p\). Then the cocycle \(\ell^{c_B}\phi(gy_0-y_0)\) is finite for \(c_B=T_{h,a,t}\). On the divisible character write \(Z_D,F_D,B_D,Z_{\rm ur}\) for cocycles, the finite condition, coboundaries and cocycles zero on inertia. Then \(Z_D\cap F_D=B_D\) for \(e>0\), and \(Z_D\cap F_D=Z_D\cap(Z_{\rm ur}+B_D)\) for \(e=0\). With \(\lambda_O\) the length normalized by the length of \(O/\ell O\),
\begin{equation}\label{eq:B8-tbk}
 \lambda_O(\{z\in Z_D:\ell^{c_B}z\in F_D\}/(Z_D\cap F_D))
 \le([K:\Q_\ell]+1)c_B .
\end{equation}
\end{lemma}
\begin{proof}[Outline]
Use \ref{B8.41ab} and \ref{B8.41cb} to compare closed Wach pairs with finite classes, then correct the lifts from \ref{B8.41c} by a polynomial in \(\gamma\). Separate the weights to kill invariants, and use \ref{B8.41aa}, corestriction and duality for the character bound. The complete proof is in \S\ref{prf:B8.41a}.
\end{proof}

\begin{lemma}[The split congruence quotient]\label{B8.41b}
In the setup of \ref{B8.41}, let \(i\) be twisted at \(2\), at \(3\) outside the small set, or at an odd small prime with split embedding types. There is a choice of column places \(A\) and conjugate row places \(A'=\{\bar u:u\in A\}\) such that the following holds for every \(\pi\in Z_{\rm good}(i)\). Put \(\ell=\ell_i\), \(K_i=\operatorname{Frac}(O_i)\), \(E_u\) for the completion of \(E\) at \(u\), \(d_u=[E_u:\Q_\ell]\), and \(W\) for the normalized rank-\(r\) lattice of \(\rho_i(\pi)\chi_{i,1}^{-1}\). Let
\[
 h_W(u)=\lambda_i H^0(E_u,W\otimes_{O_i}(K_i/O_i)),\quad
 \delta=\sum_{1\le e<t}(s_0(i)+v_\ell(e)),\quad
 h_0=\max_{u\in A,\,1\le e<t}\lambda_i H^0(E_u,(K_i/O_i)(-e)),
\]
with empty maximum zero. Use the bounding exponents
\(e'=E_{\lfloor H_B\rfloor}\), \(c_0=I_{\lfloor2b\rfloor+1,4nn_V}\), \(c_B=T_{\lfloor2b\rfloor,4nn_V,t}\) of \ref{B8.41a}. Then \(h_W(u)\le rc_0\) and
\begin{align}
 v_{cc}(i,\pi)\le{}&s_\pi+\sum_{1\le e<t}s_{\ell,e}(E)+v_\ell(h_E)
 +t\sum_{v\in Q_i}\lambda_i(O_i/(\zeta_{i,v}-1))\notag\\
 &+\sum_{u\in A}\bigl(e'+d_um(2e'+\delta)+2\delta+h_W(u)+rc_0\bigr)\notag\\
 &+tc_B\sum_{w\in A'}(d_w+2)+4n_V(|A|+1)\delta+t|A|h_0 .
 \label{eq:B8-split-quotient}
\end{align}
\end{lemma}
\begin{proof}[Outline]
Use \ref{B8.41a} for cyclotomic bounds on column invariants, finite lifts, and character blocks. Bound the relation determinant by elementary divisors and sum the lengths from the local images of Galois difference cocycles and their common global kernel. The complete proof is in \S\ref{prf:B8.41b}.
\end{proof}

\section{Cutting to the good part}

\begin{lemma}[Degree-one Chebotarev]\label{B8.42p}
Let \(F\) be a number field with algebraic closure \(\Omega\), \(U\triangleleft G_F=\Gal(\Omega/F)\) open, \(g\in G_F\), and \(A\) a finite prime set of \(F\). There is \(v\notin A\) of residue degree and ramification index one over \(\Q\) with a prolongation to \(\Omega\) whose inertia lies in \(U\) and a Frobenius representative \(\sigma\) with \(\sigma U=gU\). Existence follows from Chebotarev \cite[\S1A]{thorner-2019}. Primes of residue degree at least two have Dirichlet density zero and finitely many primes ramify over \(\Q\); conjugating the prolongation specifies the representative in \(G_F/U\).
\end{lemma}

\begin{lemma}[The two type primes and a residual characteristic]\label{B8.42pa}
Use \(g_0,g_1\) and their disjoint sets \(B(g_k,r)\) of \ref{B7.50e}. Fix an index \(i\) of nonempty good spectrum and write \(\ell=\ell_i\). A residual realization is a rank-\(r\) representation over the residue field of \(O_i\) with open kernel. For every spectral point \(x\), every split \(v\) outside \(\{v\mid\ell\}\cup S_i\cup Q_i\), every prolongation of \(v\) and every Frobenius representative there, its Frobenius polynomial times \(\prod_{c=1}^t(U-\bar\chi_{i,c}(\mathrm{Frob}_v))\) equals the reduced full Hecke polynomial at \(x,v\).

Assume the label and residual tame eigenvalue identifications, split Frobenius density in finite quotients, and residual Brauer--Nesbitt. The depth-zero labels of realized parameter lists have degrees equal to their local dimensions. For every \(x\in Z_{\rm good}(i)\) and every parameter list \(\Sigma\) for \(x\) of good signature, assume there is an integral \(\rho\) realizing \(\Sigma\) and a character \(\xi\) such that \(\rho\otimes\xi\) realizes \(x\) in rank \(r\). Require \(\xi=1\) on inertia above each \(g_k\ne\ell\). Then each residual realization is absolutely irreducible. At every \(v\mid g_k\ne\ell\) its polynomial on the prescribed tame generator is nonzero at the value of \(\bar\chi_{i,1}\) there.
\end{lemma}
\begin{proof}[Outline]
Use squarefreeness for multiplicity one and \ref{B7.51} to exclude multiple parts and proper stable subspaces by divisor degrees and tame Frobenius. Test inertia by roots, then cancel common character factors and transfer by split Frobenius density and Brauer–Nesbitt. The complete proof is in \S\ref{prf:B8.42pa}.
\end{proof}

A good cut keeps the indices with nonempty good spectrum and spectral sets \(X_i\) with \(Z_{\rm good}(i)\subset X_i\) and \(Z(i)\cap X_i\subset Z_{\rm good}(i)\), where the inclusions allow equality. An integral Hecke-commuting \(\alpha_i\) acts by \(M_0\), the order of the finite row group of Definition~\ref{B8.35}, on good eigenspaces and by zero outside \(X_i\). For every lift into \(X_i\) and every \(w\mid g_k\), let \(a_w\) be the dimension of its local parameter. Require zero monodromy and characteristic polynomial \((U-1)^{a_w}\) on every wild inertia element. For every tame generator, each root \(z\) of its inertia polynomial of multiplicity at least the threshold must satisfy \(z^{Nw-1}=1\). The threshold is \(2\) if \(p_w\) splits completely in \(E\) and either \(a_w=1\) or \(16a_w^2<p_w\). It is \(1\) otherwise.

The four local compatibility conditions specify the paired characters, the depth-zero type with its eigenspan, the inertia polynomial and the local series. The Galois hypotheses for the chosen norm include the first provision, \(Z\)-multiplicity, residual irreducibility, tame eigenvalues and the weight and block dictionaries. They include unramifiedness, depth zero, base change, weight use, localization uniqueness, coefficient largeness, and the norm and same-twist ties. For every parameter list on a good row the triple of good signature supplies a twist place \(w'\) of nonbase type and a coefficient embedding \(\iota_a\) such that \(\iota_a(z)\in\mathfrak m_{O_0(a)}\) exactly when \(v_{w'}(z)>0\), for every \(z\in\mathcal O_{K_{\rm aux}}\). It supplies a nonzero type-idempotent vector with the prescribed Hecke eigenvalues for \(1\le k\le r\) at split places outside its excluded set. Its finite-type labels are good for the characters and their inverses.

\begin{theorem}[Good cuts and restoration of the family bound]\label{B8.42}
Assume \ref{B8.42p}, the decomposition-group inputs of \ref{B8.41p}, and the preceding Galois hypotheses and four local compatibility conditions. There is a good cut \(c\) on which the four local compatibilities persist. At its induced norm choice the cut satisfies the point and admissibility conditions, the first two provisions, residual irreducibility and tame eigenvalue conditions, the weight, block and de Rham jump conditions, and the unramifiedness, depth-zero, base-change and weight identities. Localization is unique, the coefficients satisfy the required largeness and realization conditions, and the norm, good triples, finite types and common twist agree with the cut data. If also \(\mathscr H_2\) holds, \(F_m\) is twist-good, and the Galois hypotheses and four local compatibilities hold at the chosen norm for every \(x>x_1\) and \(i\in I_x\), then \(F_m\) satisfies the family local-error estimate of Definition~\ref{B8.40}.
\end{theorem}
\begin{proof}[Outline]
Cut by the good character idempotent, using \ref{B9.55a} for depth zero and \ref{B8.42pa} with local duality for residual cohomology vanishing. Sum local costs on the cut and use the integral projector and \ref{B4.52} to absorb the transfer error into the family constant. The complete proof is in \S\ref{prf:B8.42}.
\end{proof}

\begin{lemma}[The selected twist and the multiplicity sum]\label{B8.42a}
On a twist-good member with norm-choice facts, take the tied representations and let \(\nu_\ell=\min_{g\in G_E}v_\ell(\varepsilon_\ell(g)^{tr}-1)\), with \(\varepsilon_\ell\) cyclotomic. For each twisted \(i\) and good \(\pi\) the Selmer length is finite and
\begin{equation}\label{eq:B8-killing}
 s_\pi\le\nu_\ell+\ell^{-2},\qquad
 t\sum_{v\in Q_i}\lambda_i(O_i/(\zeta_{i,v}-1))\le\ell^{-2}.
\end{equation}
For every finite prime set \(S\),
\begin{equation}\label{eq:B8-nu-sum}
 \sum_{\ell\in S}\nu_\ell\log\ell\le\log Q_S+\log(tr).
\end{equation}
Under theta scalars, the \(Z\)-multiplicity identity, internal isotypic decomposition and localization uniqueness, every layer \(j\) and twisted place \(w\) satisfy
\begin{equation}\label{eq:B8-mult-rank}
 \sum_{i\in I_{j,w}}\sum_{\pi\in Z_{\rm good}(i)}
 \mathrm{mult}(i,\pi)\le\dim Z_j(\mu_w).
\end{equation}
\end{lemma}
\begin{proof}[Outline]
Cancel common characters and use \ref{B8.42p} for trace equality and \ref{B8.42pa}, matrix density, and trace pairing for integral conjugacy and the twist. Bound Selmer and tame costs by signature and conductor roots, valuations by cyclotomic image, and multiplicity by theta. The complete proof is in \S\ref{prf:B8.42a}.
\end{proof}

\begin{lemma}[The Frobenius remainder at a cut]\label{B8.42ba}
Let \(\ell\) be odd, \(q_0=\ell-1\), \(A=O[[Z]]\) for the integers \(O\) of a finite extension of \(\Q_\ell\), and
\[
 \varphi(Z)=(1+Z)^\ell-1=ZQ.
\]
Index the components by \(i\in\Z/d\Z\), with \(d>0\). Let \(P_i\subset N_i\), where \(N_i\) is a free \(A\)-lattice, and suppose \(Z^{H_*}N_i\subset P_i\) for an integer \(H_*\ge0\). Suppose a cyclotomic generator acts semilinearly by \(\gamma(Z)=(1+Z)^u-1\), trivially modulo \(Z\) on both \(N_i\) and \(P_i\). Here \(u\bmod\ell\) has order \(q_0\) and \(v_\ell(u^{q_0}-1)=1\).

Suppose linearized Frobenius is injective on \(\varphi^*N_{i-1}\) in the ambient rational module. For nonnegative integers \(M_i,h_i\), write \(H_i=M_i+h_i\), and assume
\[
 Q^{M_i}\varphi N_{i-1}\subset N_i,\qquad
 Q^{h_i}P_i\subset A\varphi P_{i-1}.
\]
Let \(c\) be a component, \(J\ge0\), and \(g,h\ge0\) with \(H_*\le2h\) and \(H_{c+k}\le2g\) for \(1\le k\le J\). Put \(C_h=\lceil1+\log_\ell(2h+1)\rceil\). Then \(\ell^{\rho_c}N_c\subset P_c\) for an integer \(\rho_c\ge0\) satisfying
\[
 \rho_c\le C_h+\frac{\ell}{q_0^2}
 \left((2g+C_h+3)\frac{q_0}{q_0-1}+\frac{2h}{q_0^J}\right).
\]
\end{lemma}
\begin{proof}[Outline]
Use the cyclotomic action to trade powers of \(Z\) for powers of \(\ell\), counting the valuations of \(u^k-1\). Expand powers of \(Q\) and use Frobenius to bound each least \(Z\)-exponent in terms of its successor, then iterate the bound over \(J\) steps. The complete proof is in \S\ref{prf:B8.42ba}.
\end{proof}

\begin{lemma}[Simultaneous lifting at two downward positions]\label{B8.42bb}
Let \(O\) be a complete discrete valuation ring of residue characteristic \(\ell\), with uniformizer \(\varpi\). Put \(A=O[[z]]\). On an interval of length \(v\), let \(p_j\), \(0\le j<v\), be monic distinguished polynomials of positive degree, pairwise relatively prime over \(\operatorname{Frac}(O)\). Multiplication of a \(p_j\) by a unit of \(A\) is allowed. For integers \(H_j\ge0\), put \(\Lambda=\prod_j p_j^{H_j}\). Mark two positions \(j_0,j_1\) and put \(D_i=p_{j_i}^{b_i}\), \(b_i\ge0\), for \(i=0,1\).

For a finite free \(O\)-module \(V\), choose saturated submodules \(V_0,V_1\), and write \(V_A=A\otimes_O V\). Define
\[
 {\cal N}(V)=\{\Lambda x:x\in V_A,\ x\in (V_i)_A+D_iV_A\ (i=0,1)\}.
\]
Define \({\cal I}(V)\) by allowing the residue of \(x\) at each marked position, modulo \((V_i)_A+D_iV_A\), to be the image of a constant vector after inverting \(\varpi\). Thus in this definition \(x\in V_A\), and the two rational constant vectors may differ.

Give a second finite free module \(W\) the same data, with its own saturated submodules. Suppose \(F:W\to V\) is linear and \(\ell^{r_0}V\subset F(W)\), where \(r_0\le2a_0\). Then for every \(y\in{\cal I}(V)\) there are \(y_1\in{\cal I}(W)\) and \(\nu\in{\cal N}(V)\) with
\[
 \ell^{2a_0} y=F_A(y_1)+\ell^{2a_0}\nu .
\]
\end{lemma}
\begin{proof}[Outline]
Distinguished division gives integral constant residues, and Smith coordinates express their difference as a sum of filtration vectors and a vector divisible by \(\varpi^\alpha\). Interpolate the residues and lift both constant vectors after scaling by \(\ell^{2a_0}\). The complete proof is in \S\ref{prf:B8.42bb}.
\end{proof}

\begin{lemma}[The numerical factor at a ramified prime]\label{B8.42bc}
Let \(\ell\) be odd and let \(e,s,n\) be positive integers with \(es=\ell-1\). Let \(h,S,\eta,r,m\) be nonnegative integers and \(1\le h'\le h+S+1\). Set \(\log_\ell 0=0\) in the following formulas. Write \(L(k)=k+v_\ell(k!)\), and put
\[
 \begin{split}
 x&=h'/s,\quad R=L(2\lceil2x\rceil-1),\quad D=L(\lceil4x\rceil-1),\\
 Y&=2R+4D+e\bigl(L(S)+\lfloor\log_\ell h'\rfloor\bigr)
       +\max\{2D,\lceil\eta+1/2+\log_\ell h'\rceil\},\\
 B&=eh/\ell+eS+e\log_\ell h+1+\eta .
 \end{split}
\]
Let \(J\subset[-h,2h]\cap\Z\) exclude zero, and set
\[
 E'=2L(4h)+\sum_{j\in J}(1+v_\ell(j))+L(h).
\]
Let \(C\) be a quotient of finite lattice cohomology over a coefficient discrete valuation ring \(O\) above \(\ell\). Suppose \(C\) is generated by at most \(2(ne+1)r\) elements and has a submodule \(C_0\) generated by at most \(ne(e+1)m\) elements, with \(\ell^YC\subset C_0\) and \(\ell^{E'}C_0=0\). Then its length, divided by \(\operatorname{length}_O(O/\ell O)\), is at most
\[
 160\bigl((ne+1)r B+ne(e+1)m(h+1)\bigr).
\]
\end{lemma}
\begin{proof}[Outline]
Count factorial valuations and bound the logarithmic terms to obtain \(Y\le80B\) and \(E'\le18h\). Apply the generator bounds to \(C/C_0\) and \(C_0\) and add their normalized lengths. The complete proof is in \S\ref{prf:B8.42bc}.
\end{proof}

\begin{lemma}[The coefficient field and the crystalline representations at \(q,d\)]\label{B8.42bd}
Let \(\ell\in\{q,d\}\), let \(w\mid\ell\) in the host \(E\), and put \(K=E_w\). If \(K_0\) is the maximal unramified subfield of \(K/\Q_\ell\), then
\[
 K=K_0(\sqrt{-\ell}),\qquad
 K\subset K_0(\mu_\ell),\qquad
 [K_0(\mu_\ell):K]=(\ell-1)/2 .
\]
At an index \(i\) of residue characteristic \(\ell\) on the good cut of Theorem~\ref{B8.42}, write \(R_x\) for the full representation at \(x\in X_i\), and put \(V_x=R_x\otimes\chi_{i,1}^{-1}\). At a good point \(\pi\), put \(W_\pi=\rho_i(\pi)\otimes\chi_{i,1}^{-1}\). Over \(K\), every \(V_x\) and both \(W_\pi\) and \(W_\pi^\vee(1)\) are crystalline. Every subrepresentation of a finite direct sum of the \(V_x^\vee(1)\) is crystalline. Its jumps at an embedding \(\tau\) belong to the union of the sets
\[
             \{-1-a:a\text{ is a jump of }V_x\text{ at }\tau\}
\]
over the summands in that sum.
\end{lemma}
\begin{proof}[Outline]
Use the quadratic Gauss sum of \ref{B2.46} and the cyclotomic Eisenstein polynomial to compute the local field and index. Apply \ref{B8.42} and dimension additivity for crystalline constituents, then use the filtration and Tate twist convention of \ref{B9.04} for the jumps. The complete proof is in \S\ref{prf:B8.42bd}.
\end{proof}

\begin{definition}[Cyclic carries and position data]\label{B8.42bm}
Let \(\ell\) be prime and let \(d',t\ge1\) be integers. For \(h:\Z/d'\Z\to\N\) and \(i\in\Z/d'\Z\), the cyclic carry of \(h\) at \(i\) is
\[
 \sigma_h(i)=\Bigl\lfloor\sum_{j\ge1}h(i+j)\,\ell^{-j}\Bigr\rfloor .
\]
Position data at \(\ell\) of width \(t\) consist of maps \(h,w,\eta':\Z/d'\Z\to\N\), signs \(\varsigma_i\in\{1,-1\}\) and finite multisets \(U_i\) of integers, \(i\in\Z/d'\Z\), with the following six properties. Put \(\sigma_i=\sigma_h(i)\). Let the low window \(W_i\) be \([1,\eta'_i]\) if \(\varsigma_i=1\) and \([h_i-\eta'_i,h_i-1]\) if \(\varsigma_i=-1\). Put \(\bar x=x\) if \(\varsigma_i=1\) and \(\bar x=h_i-x\) if \(\varsigma_i=-1\).
(P1) Every \(u\in U_i\) satisfies \(1\le u\le h_i-1\).
(P2) \(1\le w_i\) and \(w_i+t\le h_i\).
(P3) \(w_i+t-1\le\eta'_i\) if \(\varsigma_i=1\), and \(h_i-\eta'_i\le w_i\) if \(\varsigma_i=-1\).
(P4) \(\eta'_i+4\sigma_i+6<\ell\).
(P5) \(|x-y|>\sigma_i+2\) for every \(x\in U_i\cap W_i\) and every integer \(y\in[w_i,w_i+t-1]\).
(P6) \(\bar x\not\equiv m\pmod\ell\) for every \(x\in U_i\) outside \(W_i\) and every integer \(m\in[-2\sigma_i-2,\ \eta'_i+2\sigma_i+3]\).
\end{definition}

\begin{lemma}[Position data at one embedding]\label{B8.42bn}
Let \(\ell\) be prime, \(t\ge1\) an integer, \(\varsigma\in\{1,-1\}\), and let \(F_l,F_c\) be finite sets of integers. Put \(c_t=(t-1)/2\) and \(\delta(z)=\varsigma(z-c_t)\). Let \(\eta,S,B,W_0\) be real numbers and \(\sigma,\sigma_P\) integers such that
\[
 t\le\eta,\quad B\ge0,\quad S\ge0,\quad 2B+\tfrac t2+3\le W_0,\quad \sigma\le\min(2B,\sigma_P),\quad 2\eta+8B+7<\ell .
\]
Suppose every \(z\in F_l\) satisfies \(2B+t/2+3\le|z-c_t|\le\eta\), every \(z\in F_c\) satisfies \(W_0\le\delta(z)\le W_0+S+\eta+1\), and \(\ell\) divides no integer in
\begin{equation}\label{eq:B8-pos-window}
 [W_0-\eta-2\sigma_P-3,\ W_0+S+2\eta+2\sigma_P+4].
\end{equation}
Let \(m_l\) be the least element of \(F_l\cup\{0\}\) and \(M_l\) the largest element of \(F_l\cup\{t-1\}\). Define \(m_a,M_a\) in the same way with \(F_l\cup F_c\) in place of \(F_l\). Put \(\eta'=M_l-m_l+1\), \(U=\{z+w:z\in F_l\cup F_c\}\) and
\[
 (w,h)=(1-m_l,\ M_a+2-m_l)\ \text{ if }\varsigma=1,\qquad
 (w,h)=(1-m_a,\ M_l+2-m_a)\ \text{ if }\varsigma=-1 .
\]
Then \(h\le W_0+S+2\eta+3\), the multiplicity of \(b+w\) in \(U\) is \(1\) for \(b\in F_l\cup F_c\) and \(0\) for other integers \(b\), and \((h,w,\eta',\varsigma,U)\) satisfies (P1)--(P6) of Definition~\ref{B8.42bm} at one index, with \(\sigma\) in place of \(\sigma_i\).
\end{lemma}
\begin{proof}[Outline]
Bounds \(m_l, M_l,\eta'\) from the \(F_l\) inequalities to verify (P4), (P5), and the multiplicities. Treating \(\varsigma=\pm1\) separately checks (P1)--(P3) and the bound on \(h\), while any congruence violating (P6) places an \(\ell\)-multiple in the excluded interval. The complete proof is in \S\ref{prf:B8.42bn}.
\end{proof}

\begin{definition}[Admissible cocycles and the large-prime statement]\label{B8.42bo}
Let \(\ell\ge5\) be prime and \(K/\Q_\ell\) finite unramified of degree \(d'\). Let \(L'\) be a finite extension of \(\Q_\ell\) containing the images of all embeddings of \(K(\mu_\ell)\), with integers \(O\) and uniformizer \(\varpi\). Number the embeddings \(\tau_i:K\to L'\), \(i\in\Z/d'\Z\), in Frobenius order as in Lemma~\ref{B8.41c}. For a continuous representation \(V\) of \(G_K\) over \(L'\), write \(m_V(\tau,a)\) for the multiplicity of the jump \(a\) at \(\tau\) in the convention of Definition~\ref{B9.04}. Write \(\chi_\ell\) for the cyclotomic character.

Fix \(t\ge1\) and position data \((h,w,\eta',\varsigma,U)\) at \(\ell\) of width \(t\). Let \(V\) be a continuous crystalline representation with \(m_V(\tau_i,a)=\operatorname{mult}_{U_i}(a)\) for all \(i,a\), let \(T_0\subset V\) be a stable \(O\)-lattice, and let \(\chi_1,\dots,\chi_t:G_K\to O^\times\) be continuous crystalline characters with \(\chi_c=\chi_1\chi_\ell^{1-c}\) and with the single jump \(w_i+c-1\) at \(\tau_i\). For \(g\in G_K\) put
\[
 P_g(U)=\det(U-g\mid V)\prod_{c=1}^t(U-\chi_c(g)),\qquad
 \tilde m_i(a)=\operatorname{mult}_{U_i}(a)+\mathbf 1_{[w_i,w_i+t-1]}(a).
\]
A full representation for these data is a continuous crystalline representation \(V'\) of \(G_K\) over \(L'\) with \(m_{V'}(\tau_i,a)=\tilde m_i(a)\) for all \(i,a\), such that \(\det(U-g\mid V')-P_g(U)\) has coefficients in \(\varpi O\) for every \(g\).

Let \(L\) be a stable \(O\)-lattice in \(V\) or in one of the characters \(\chi_c\), let \(\mathcal M\) be a finite \(O\)-module, and let \(z:G_K\to L\otimes_O\mathcal M\) be a continuous cocycle for the action of \(g\) by \(\chi_1(g)^{-1}g\). For \(a\ge1\) with \(\varpi^a\mathcal M=0\), let \(\mathcal E_{z,a}=(L\otimes\mathcal M)\oplus O/\varpi^a\) with
\[
 g(v,s)=(gv+s\chi_1(g)z(g),\ \chi_1(g)s).
\]
The cocycle \(z\) is admissible if there is an integer \(a_0\ge1\) such that for every integer \(a\ge a_0\) with \(\varpi^a\mathcal M=0\) the \(O[G_K]\)-module \(\mathcal E_{z,a}\) is isomorphic to \(L_1/L_2\) for stable \(O\)-lattices \(L_2\subset L_1\) in a finite direct sum of full representations.

Statement (T) asserts, for all such data and every finite \(O\)-module \(\mathcal M\), the following. (T1) A cocycle \(z:G_K\to T_0\otimes\mathcal M\) is admissible exactly when \(z=\sum_j\xi_j\otimes m_j+z_0\), where \(m_j\in\mathcal M\), \(z_0\) is a coboundary, and the \(\xi_j\) are cocycles of \(\chi_1^{-1}V\) with values in \(T_0\) whose classes are finite in the sense of Definition~\ref{B8.41h}. (T2) For \(2\le c\le t\) and a stable lattice \(L\) of \(\chi_c\), a cocycle with values in \(L\otimes\mathcal M\) is admissible exactly when it is a coboundary. (T3) For a stable lattice \(L\) of \(\chi_1\), a cocycle with values in \(L\otimes\mathcal M\) is admissible exactly when its restriction to the inertia group of \(K\) is a coboundary there.
\end{definition}

\begin{definition}[Admissible finite modules and cohomology classes]\label{B8.42boa}
Fix data \((K,L',O,\varpi,\tau_i,t,h,w,\eta',\varsigma,U,V,T_0,\chi_1,\dots,\chi_t)\) at a prime \(\ell\ge5\) as in Definition~\ref{B8.42bo}. A finite \(O[G_K]\)-module \(N\) is \emph{admissible} for these data if there exist a finite direct sum \(\mathcal V\) of full representations, stable \(O\)-lattices \(L_2\subset L_1\subset\mathcal V\), and an \(O[G_K]\)-equivariant surjection \(p:L_1\twoheadrightarrow N\) with \(\ker p=L_2\). Let \(W\) be \(T_0\) or \(O(\chi_c)\) for \(1\le c\le t\), write \(W'=\chi_1^{-1}W\), and let \(\mathcal M\) be a finite \(O\)-module. Since \(O\) is a complete discrete valuation ring and \(\varpi\) is a uniformizer, there exists an integer \(a\ge1\) with \(\varpi^a\mathcal M=0\). Define \(\operatorname{Adm}(W,\mathcal M)\subset H^1(K,W'\otimes_O\mathcal M)\) to be the subset of cohomology classes represented by continuous cocycles \(z:G_K\to W\otimes_O\mathcal M\) (for the action \(g\mapsto\chi_1(g)^{-1}g\)) for which there exists \(a_0\ge1\) such that \(\mathcal E_{z,a}\) is an admissible finite \(O[G_K]\)-module for every integer \(a\ge a_0\) with \(\varpi^a\mathcal M=0\).
\end{definition}

\begin{definition}[Cyclotomic scalar data and shifted kernels]\label{B8.42bod}
Let \(\ell\ge5\) be prime and let \(k\) be a field of characteristic \(\ell\). Put \(A=k[[X]]\) and \(B=k((X))\). Scalar Frobenius and cyclotomic operators are commuting \(k\)-algebra endomorphisms \(\varphi,\gamma:B\to B\) preserving \(A\), with
\[
 \varphi(X)=uX^\ell,\qquad
 \gamma(X)=X(1-\kappa X),\qquad
 u\in A^\times,\quad \kappa\in A,\quad\kappa(0)=1.
\]
To specify the cyclotomic condition, introduce a second variable \(\pi\), put \(A_\pi=k[[\pi]]\), \(B_\pi=k((\pi))\), and set
\[
 \varphi_\pi(f)(\pi)=f(\pi^\ell),\qquad
 \gamma_\pi(f)(\pi)=f((1+\pi)^{1+\ell}-1).
\]
For \(z\in\Z_\ell\), the series \((1+\pi)^z\) means \(\sum_{j\ge0}\overline{\binom zj}\pi^j\), with coefficients reduced from \(\Z_\ell\) to \(k\). Let \(\Delta=\{z\in\Z_\ell^\times:z^{\ell-1}=1\}\); it acts on \(B_\pi\) by the substitutions \(\delta_z(\pi)=(1+\pi)^z-1\). The scalar data are cyclotomic if there is a \(k\)-algebra embedding \(\iota:B\hookrightarrow B_\pi\) such that
\[
 \iota(A)=A_\pi^\Delta,\quad
 \iota(X)=v\pi^{\ell-1}\ (v\in A_\pi^\times),\quad
 \iota\varphi=\varphi_\pi\iota,\quad
 \iota\gamma(a)=\gamma_\pi\iota(a)\quad(a\in A)
\]
(hence \(\iota(B)=B_\pi^\Delta\) and \(\iota\gamma=\gamma_\pi\iota\) on \(B\)). For a finite-dimensional \(B\)-space \(E\), an \(A\)-lattice is a finitely generated \(A\)-submodule \(M\subset E\) with \(BM=E\). For \(a\in\Z\), write \(X^aM=\{X^am:m\in M\}\). For a \(\varphi\)-semilinear map \(f:E\to E\), put \(f^*M=A f(M)\), the \(A\)-span of its image. For a \(k\)-linear map \(\Psi:E\to E\), put
\[
 K_a(M,\Psi)=X^aM\cap\ker\Psi.
\]
For \(b\in B\setminus\{0\}\) and \(y\in B_\pi\setminus\{0\}\), let \(v_X(b)\) and \(v_\pi(y)\) be their least Laurent exponents (with \(v_X(0)=v_\pi(0)=+\infty\)), and write \(M_n(R)\) for the ring of \(n\times n\) matrices over a ring \(R\). Frobenius contraction refers to the equality \(v_X(\varphi(b))=\ell v_X(b)\), or equivalently to multiplication of the \(X\)-adic order by \(\ell\).
\end{definition}
\begin{proof}[Outline]
Approximating \(z,w\) by nonnegative integers verifies the \(\Delta\)-action on \(B_\pi\). Clearing denominators by powers of \(\iota(X)\) gives \(\iota(B)=B_\pi^\Delta\) and \(\gamma\)-intertwining; factoring \(b=X^ma\) gives \(v_X(\varphi(b))=\ell m\). The complete proof is in \S\ref{prf:B8.42bod}.
\end{proof}

\begin{lemma}[The exact \(\psi\)-shift]\label{B8.42boe}
For cyclotomic scalar data as in Definition~\ref{B8.42bod}, every \(y\in B_\pi\) has a unique expansion
\[
 y=\sum_{j=0}^{\ell-1}(1+\pi)^j\varphi_\pi(y_j),\qquad y_j\in B_\pi.
\]
Let \(\psi_\pi(y)=y_0\). There is a \(k\)-linear map \(\psi:B\to B\), determined by \(\iota\psi=\psi_\pi\iota\), with
\[
 \psi(\varphi(b_1)b_2)=b_1\psi(b_2)\quad(b_1,b_2\in B),\qquad
 \psi\varphi=1,\qquad \psi\gamma=\gamma\psi,
 \qquad\psi(X^aA)=A\quad(2-\ell\le a\le1).
\]
Let \(E\) be a finite-dimensional \(B\)-space, let \(f:E\to E\) be \(\varphi\)-semilinear with \(Bf(E)=E\), and let \(g_0,g:E\to E\) be \(\gamma\)-semilinear maps satisfying \(fg_0=gf\). There is a \(k\)-linear map \(\Psi:E\to E\) such that
\[
 \Psi(bf(x))=\psi(b)x\quad(b\in B,\ x\in E),\qquad \Psi g=g_0\Psi.
\]
For every \(A\)-lattice \(M_0\subset E\), put \(P=f^*M_0\). Then
\[
 \Psi(X^aP)=M_0\qquad(2-\ell\le a\le1).
\]
If \((g_0-1)M_0\subset XM_0\), then
\begin{equation}\label{eq:B8-psi-shift-exact}
 (g-1)K_a(P,\Psi)=K_{a+1}(P,\Psi)
 \qquad(2-\ell\le a\le0).
\end{equation}
\end{lemma}
\begin{proof}[Outline]
Changing basis to \((1+\pi)^j\) over \(\varphi_\pi(B_\pi)\) defines \(\psi_\pi\), which descends to \(\psi\) and coordinatewise \(\Psi\). Solving \((1+\pi)^jH\gamma_\pi(y)-y\) recursively for \(1\le j<\ell\) and descending to \(\ker\psi\) proves the exact shift. The complete proof is in \S\ref{prf:B8.42boe}.
\end{proof}

\begin{definition}[Residue models and their extension classes]\label{B8.42bof}
Retain the scalar data of Definition~\ref{B8.42bod}, let \(d'\ge1\) be an integer, and put \(I=\Z/d'\Z\), \(A_I=\prod_{i\in I}k[[X]]\), and \(B_I=\prod_{i\in I}k((X))\). A cyclotomic module over \(B_I\) is a family of finite-dimensional \(B\)-spaces \(D_i\), with \(\varphi\)-semilinear maps \(\varphi_i:D_{i-1}\to D_i\) whose images span \(D_i\), and bijective \(\gamma\)-semilinear maps \(\gamma_i:D_i\to D_i\), satisfying
\[
 \varphi_i\gamma_{i-1}=\gamma_i\varphi_i;
\]
a morphism \((D_i)\to(D'_i)\) is a family of \(B\)-linear maps \(f_i:D_i\to D'_i\) intertwining \(\varphi_i\) and \(\gamma_i\). A lattice in this module is a family \(N_i\subset D_i\) of \(A\)-lattices, regarded as the \(A_I\)-module \(\prod_iN_i\). Trivial cyclotomic action modulo \(X\) means \((\gamma_i-1)N_i\subset XN_i\) for every \(i\).

Let \(K/\Q_\ell\) be a finite extension and let \(\mathcal C\) be a class of finite-dimensional continuous \(k[G_K]\)-representations closed under subquotients and finite direct sums and containing \(k\). A residue realization \(D\) is a \(k\)-linear exact functor from \(\mathcal C\) to these cyclotomic modules, preserving dimension at each component, together with elements \(\mathbf1_i\in D(k)_i\setminus\{0\}\) satisfying
\[
 \varphi_i(\mathbf1_{i-1})=\mathbf1_i,\qquad
 \gamma_i(\mathbf1_i)=\mathbf1_i.
\]
It has the following fixed-vector property: if \(x_i\in D(U)_i\) satisfy \(\varphi_i(x_{i-1})=x_i\) and \(\gamma_i(x_i)=x_i\), then there is a \(G_K\)-equivariant map \(v:k\to U\) with \(D(v)_i(\mathbf1_i)=x_i\).

Fix a collection \(\mathscr B\) of pairs \((U_b,N^b)\), where \(U_b\in\mathcal C\) and \(N^b\) is a lattice in \(D(U_b)\). A model lattice of \(U\in\mathcal C\) is a lattice obtained as follows. Choose \((U_b,N^b)\in\mathscr B\), an injective equivariant map \(j:U_1\hookrightarrow U_b\), and a surjective equivariant map \(p:U_1\twoheadrightarrow U\). Set
\begin{equation}\label{eq:B8-residue-model}
 N_i=D(p)_i\bigl(D(j)_i^{-1}(N_i^b)\bigr).
\end{equation}
In the setting of Definition~\ref{B8.42bo} with \(k=O/\varpi O\), choose \(\gamma\in\Gal(K(\mu_{\ell^\infty})/K(\mu_\ell))\) with \(\chi_\ell(\gamma)=1+\ell\), identify \(A_\pi=k[[\pi]]\) with \(O[[Z]]/\varpi O[[Z]]\) via \(\pi=\bar Z\), and embed \(A=k[[X]]\) into \(A_\pi\) by \(\iota(X)=-\prod_{z\in\Delta}((1+\pi)^z-1)\), which gives cyclotomic scalar data with \(u=1\). Let \(T\) be a stable \(O\)-lattice in a crystalline \(L'\)-representation with \(\overline T=T/\varpi T\in\mathcal C\). Let \(D_{\rm et}(T)\) be the étale cyclotomic module of \(T\) over \(K(\mu_\ell)\) with Wach lattice \(N(T)\) as in Lemmas~\ref{B8.41ca}--\ref{B8.41cb} (at \(s=1\)), on which \(\Delta\) acts via the Teichmüller lift \(\omega\circ\bar\chi_\ell:\Gal(K(\mu_\ell)/K)\xrightarrow{\sim}\Delta\) of the modulo-\(\ell\) cyclotomic character. A residue presentation of \(D\) on \(T\) is a family of surjective additive maps \(\jmath_i:D_{\rm et}(T)^\Delta\to D(\overline T)_i\) compatible with Frobenius, \(\gamma\), and scalar multiplication through \(O[[Z]]^\Delta\to A_\pi^\Delta=\iota(A)\) (at the \(i\)-th component) such that its reduced base lattice
\[
 \overline N_i(T)=
 \{\jmath_i(x):x\in N(T),\ \delta x=x\text{ for every }\delta\in\Delta\}
\]
spans \(D(\overline T)_i\) over \(B\); for the canonical modulo-\(\varpi\) realization \(D(U)_i\) obtained by taking \(\Delta\)-invariants of the étale \((\varphi,\Gamma)\)-module of \(U\) over \(\prod_{i\in I}k((\pi))\), reduction modulo \(\varpi\) supplies \(\jmath_i\).

Let \(U\in\mathcal C\). A cocycle is a continuous function \(z:G_K\to U\) with \(z(gh)=z(g)+g z(h)\). It defines the extension \(E_z=U\oplus k\) with action
\[
 g(u,a)=(gu+a z(g),a).
\]
Such a cocycle has a model if \(E_z\in\mathcal C\) and \(E_z\) has a model lattice of the form \eqref{eq:B8-residue-model}. Coboundaries are the cocycles \(z(g)=gu-u\), with \(u\in U\).

When \(E_z\in\mathcal C\), to define the Herr extension class of \(z\), choose \(e_i\in D(E_z)_i\) mapping to \(\mathbf1_i\). Via the injection \(D(U)_i\hookrightarrow D(E_z)_i\), put
\[
 y_i=\varphi_i(e_{i-1})-e_i,\qquad
 t_i=\gamma_i(e_i)-e_i.
\]
They satisfy \((\gamma_i-1)y_i=\varphi_i(t_{i-1})-t_i\). The Herr extension class is the class of \((y_i,t_i)_i\) in
\[
 \frac{\{(y_i,t_i)_i:(\gamma_i-1)y_i=\varphi_i(t_{i-1})-t_i
                         \text{ for all }i\}}
 {\{(\varphi_i(x_{i-1})-x_i,\ \gamma_i(x_i)-x_i)_i:x_i\in D(U)_i\}}.
\]
For an integer multiset \(J_i\) and an integer shift \(w_i\), define its reduced jump multiplicity at \(\alpha\in k\) by
\[
 r_i(J_i,w_i;\alpha)=\#\{u\in J_i:\overline{u-w_i}=\alpha\},
\]
with multiplicities counted in \(J_i\). The reduced jump multiset is \(\{\overline{u-w_i}:u\in J_i\}\).
\end{definition}
\begin{proof}[Outline]
Exactness of \(D\) and Noetherianity of \(A\) show that model and reduced lattices are finitely generated over \(A\) and span over \(B\). The cocycle law gives the \(E_z\)-action, and \(\varphi_i,\gamma_i\)-commutation makes the Herr class well-defined. The complete proof is in \S\ref{prf:B8.42bof}.
\end{proof}

\begin{lemma}[The geometric carry bound]\label{B8.42boba}
Let \(d'\ge1\), \(\ell\ge2\) and \(H\ge0\) be integers. Let
\(h:\Z/d'\Z\to\N\) satisfy \(h_i\le H\) for every \(i\), and put
\[
 \sigma_i=\left\lfloor\sum_{j\ge1}h_{i+j}\ell^{-j}\right\rfloor.
\]
Then \(\sigma_i\le H/(\ell-1)\) for every \(i\).
\end{lemma}
\begin{proof}
For every \(j\ge1\),
\(0\le h_{i+j}\ell^{-j}\le H\ell^{-j}\). The geometric series on the
right converges and has sum \(H/(\ell-1)\). Comparison gives convergence
of the series defining \(\sigma_i\) and
\[
 \sigma_i\le\sum_{j\ge1}h_{i+j}\ell^{-j}
              \le\sum_{j\ge1}H\ell^{-j}=\frac{H}{\ell-1}.
\]
\end{proof}

\begin{lemma}[Excluding false roots in a downward window]\label{B8.42bobb}
Let \(\ell,d',t,h,w,\eta',\varsigma,U\) be as in Definition~\ref{B8.42bm}, and assume (P1)--(P4) and (P6). Put
\(J_i=U_i+\{w_i,w_i+1,\ldots,w_i+t-1\}\), with multiset addition. If
\(\varsigma_i=-1\), \(u\in J_i\), and \(a\) is an integer with
\(0\le a\le h_i-w_i\), then
\[
 \ell\mid a-(u-w_i)\quad\Longrightarrow\quad a=u-w_i.
\]
\end{lemma}
\begin{proof}[Outline]
Set \(m=h_i-w_i-a\in[0,\eta'_i]\) by (P3), so \(h_i-u\equiv m\pmod\ell\). Condition (P6) rules out \(u\in U_i\) outside the low window; otherwise (P1)--(P4) give \(\lvert(h_i-u)-m\rvert\le\eta'_i<\ell\), forcing \(h_i-u=m\). The complete proof is in \S\ref{prf:B8.42bobb}.
\end{proof}

\begin{lemma}[Transporting independent admissible residue classes]\label{B8.42bobc}
Retain the data of Definition~\ref{B8.42bo}, let \(k=O/\varpi O\), and
put \(\overline W=(\chi_1^{-1}T_0)/\varpi(\chi_1^{-1}T_0)\). Suppose
the residue realization of Definition~\ref{B8.42bof} carries every
\(\chi_1^{-1}\)-twisted admissible finite module to a module with a
model lattice. Assume also that admissibility of a residue class is
equivalent to admissibility of its extension at exponent \(1\).
Given \(s\) linearly independent classes in
\(\operatorname{Adm}(T_0,k)\), one can choose continuous cocycles
\(z_1,\ldots,z_s:G_K\to\overline W\) whose extensions have model
lattices and whose classes modulo coboundaries are linearly
independent. Their Herr extension classes are linearly independent.
\end{lemma}
\begin{proof}[Outline]
Quotienting and twisting identifies coefficients with \(\overline W\), giving modelled cocycles and preserving independence. A zero Herr class yields \(\varphi_i\)- and \(\gamma_i\)-fixed lifts splitting the extension, so its coefficients vanish. The complete proof is in \S\ref{prf:B8.42bobc}.
\end{proof}

\begin{lemma}[Character quotients and the two residual vanishings]\label{B8.42bobd}
Let \(k\) be a finite field of characteristic \(\ell\), let
\(T\) be a finite-dimensional continuous \(k[G_K]\)-representation,
and let \(\omega=\bar\chi_\ell\). Suppose that, for
each \(e\in\{-1,0\}\), every finite-dimensional representation
\(A\), every injective equivariant map \(j:A\hookrightarrow T\),
and every equivariant map \(p:A\to k(\omega^{-e})\), one has \(p=0\).
Then
\[
 \operatorname{Hom}_{G_K}(k,T)=0,
 \qquad \operatorname{Hom}_{G_K}(T,k(\omega))=0.
\]
Under local Tate duality these imply \(H^0(K,T)=H^2(K,T)=0\).
\end{lemma}
\begin{proof}
A nonzero equivariant map \(k\to T\) is injective. Taking \(A=k\),
\(e=0\), and \(p=\operatorname{id}_k\) contradicts the hypothesis.
For the second assertion take \(A=T\), \(j=\operatorname{id}_T\),
and \(e=-1\). The first vanishing says \(T^{G_K}=0\). Local Tate
duality identifies the dual of \(H^2(K,T)\) with
\(H^0(K,T^\vee(1))\), and the latter is the second displayed Hom
space. This proves both cohomological vanishings.
\end{proof}

\begin{lemma}[Unramified extensions and the residue dimension one test]\label{B8.42bobe}
Retain the setting of Definition~\ref{B8.42bo}. Assume the crystalline
sum and unramified twist rules, closure of admissible modules under
subquotients and finite direct sums, and admissibility of
\(O(\chi_1)\otimes_O X\) for every finite \(O\)-module \(X\).
Assume also that a class is admissible whenever its extension is
admissible at one exponent annihilating its coefficient module.
Then
\[
 H^1_{\rm ur}(K,X)\subset\operatorname{Adm}(O(\chi_1),X)
 \qquad\text{for every finite }X.
\]
Suppose, in addition, that these admissible classes form submodules,
are functorial in \(X\), satisfy pullback along injections as in
Lemma~\ref{B8.42bob}(a), and have residue dimension at most one.
Then
\[
 \operatorname{Adm}(O(\chi_1),X)=H^1_{\rm ur}(K,X),\qquad
 \dim_k\operatorname{Adm}(O(\chi_1),k)=1.
\]
\end{lemma}
\begin{proof}[Outline]
A lattice in \(\chi_1\mu\oplus\chi_1\) for \(\mu=(1+q)^{\nu}\) pushes out to show every unramified class is admissible. Subtracting the unramified part and applying Nakayama's lemma to any nonzero remainder contradicts the residue dimension bound. The complete proof is in \S\ref{prf:B8.42bobe}.
\end{proof}

\begin{lemma}[The carry recurrence and the standard residue counts]\label{B8.42bobf}
Let the position data satisfy (P1)--(P6) of Definition~\ref{B8.42bm}. Put
\[
 d_i=\sum_{u<w_i}\operatorname{mult}_{U_i}(u),\qquad
 \rho_i(\alpha)=r_i(U_i,w_i;\alpha)\quad(\alpha\in k),
\]
where \(k\) has characteristic \(\ell\) and \(r_i\) is defined in Definition~\ref{B8.42bof}. Then
\[
 \sigma_i+1\le\ell(\sigma_{i-1}+1)-h_i.
\]
No \(u\in U_i\) is congruent modulo \(\ell\) to an integer in
\[
 [w_i-\sigma_i-2,\ w_i+t+1+\sigma_i].
\]
In particular, \(\rho_i(\bar a)=0\) for \(0\le a\le\sigma_i+1\). The two block counts are
\[
 \sum_{j=0}^{w_i-1}\rho_i(\overline{-w_i+j})=d_i
       \quad(\varsigma_i=1),
 \qquad
 d_i+\sum_{a=0}^{h_i-w_i}\rho_i(\bar a)=|U_i|
       \quad(\varsigma_i=-1).
\]
\end{lemma}
\begin{proof}[Outline]
The identity \(\ell S_{i-1}=h_i+S_i\) for \(S_i=\sum_{j\ge1}h_{i+j}\ell^{-j}\) gives the carry bound. Conditions (P4)--(P6) exclude cloud jumps and turn low-window congruences modulo \(\ell\) into equalities, proving the excluded band and block counts. The complete proof is in \S\ref{prf:B8.42bobf}.
\end{proof}

\begin{lemma}[A common model for the residual Herr representatives]\label{B8.42bobg}
Assume \(\mathscr H_2\) and the data of Definition~\ref{B8.42bo}. Put \(k=O/\varpi O\) and
\(\overline W=(\chi_1^{-1}T_0)/\varpi(\chi_1^{-1}T_0)\), and use the canonical residue realization of Definition~\ref{B8.42bof}. Take the base pairs to be the reductions of stable lattices in finite sums of \(\chi_1^{-1}\)-twisted full representations. There is a model lattice \(N^{\max}\) of \(D(\overline W)\) containing every model lattice of \(\overline W\), including its reduced Wach lattice \(N^{\rm std}\).

Put \(P_i^{\max}=A\varphi_i(N_{i-1}^{\max})\). Every admissible residue class has a Herr representative \((b_i,g_i)_i\) satisfying
\[
 (\gamma_i-1)b_i=\varphi_i(g_{i-1})-g_i,\qquad
 g_i\in XN_i^{\max},
\]
and
\[
 b_i\in X^{-w_i}N_i^{\max}\quad(\varsigma_i=1),
 \qquad
 b_i\in X^{-(h_i-w_i)}P_i^{\max}\quad(\varsigma_i=-1).
\]
The common lattice satisfies
\[
 X^{h_i-w_i}N_i^{\max}\subset P_i^{\max}
       \subset X^{-w_i}N_i^{\max},\qquad
 (\gamma_i-1)N_i^{\max}\subset XN_i^{\max}.
\]
At a downward index, with \(\nabla_i=X^{-1}(\gamma_i-1)\), there is \(f_i\in k[Y]\) such that
\[
 f_i(0)\ne0,\qquad
 \nabla_i f_i(\nabla_i)N_i^{\max}
       \subset X^{h_i-w_i+1}N_i^{\max}.
\]
An independent family of admissible residue classes gives an independent family of these Herr classes.
\end{lemma}
\begin{proof}[Outline]
The largest-model theorem on \(N^{\rm std}\) gives \(N^{\max}\), and extension-model comparison places all Herr pairs inside it. Height, \(\gamma\), and downward polynomial relations pass via a base presentation; Lemma~\ref{B8.42bobc} gives independence. The complete proof is in \S\ref{prf:B8.42bobg}.
\end{proof}

\begin{lemma}[The character obstruction]\label{B8.42bobh}
Assume \(\mathscr H_2\) and the data of Definition~\ref{B8.42bo}. Put \(k=O/\varpi O\) and
\(\overline W=(\chi_1^{-1}T_0)/\varpi(\chi_1^{-1}T_0)\) and \(\omega=\bar\chi_\ell\). For every integer \(-1\le e<t\), every finite-dimensional continuous \(k[G_K]\)-representation \(E\), every injective equivariant map \(j:E\hookrightarrow\overline W\), and every equivariant map \(p:E\to k(\omega^{-e})\), one has \(p=0\).
\end{lemma}
\begin{proof}[Outline]
If \(p\ne0\), pulling \(N^{\rm std}\) back by \(D(j)\) and pushing forward by \(D(p)\) yields a model \(N^q\). Carry comparison with \(C^{(e)}\) and subquotient multiplicity bounds place a jump of \(U_0\) in the excluded band of Lemma~\ref{B8.42bobf}. The complete proof is in \S\ref{prf:B8.42bobh}.
\end{proof}

\begin{lemma}[The residue character bounds]\label{B8.42bobi}
Assume \(\mathscr H_2\) and the data of Definition~\ref{B8.42bo}, and put \(k=O/\varpi O\). Then
\[
 \operatorname{Adm}(O(\chi_c),k)=0\quad(2\le c\le t),
 \qquad \dim_k\operatorname{Adm}(O(\chi_1),k)=1.
\]
\end{lemma}
\begin{proof}[Outline]
For \(2\le c\le t\), Definition~\ref{B8.42bod} and (P3)--(P4) clear poles on \(C^{(c-1)}\) and Frobenius contraction yields a boundary. For \(c=1\), constant-pair reduction bounds the dimension by \(1\); Lemma~\ref{B8.42bobe} gives a nonzero class. The complete proof is in \S\ref{prf:B8.42bobi}.
\end{proof}

\begin{lemma}[The \(\psi\)-shift dimension estimate and admissible cohomology classes]\label{B8.42bob}
Assume \(\mathscr H_2\). Fix data \((K,L',O,\varpi,\tau_i,t,h,w,\eta',\varsigma,U,V,T_0,\chi_1,\dots,\chi_t)\) at a prime \(\ell\ge5\) as in Definition~\ref{B8.42bo}. At each embedding \(\tau_i:K\to L'\) (\(i\in\Z/d'\Z\)), the Hodge--Tate jumps \(u\in U_i\) of \(V\) and \(w_i+c-1\) of \(\chi_c\) (\(1\le c\le t\)) lie in \([1,h_i-1]\), with no upper bound of the form \(h_i<\ell\) on the spread \(h_i\); relative to \(\ell\), the position data \((h,w,\eta',\varsigma,U)\) satisfy conditions (P1)--(P6) of Definition~\ref{B8.42bm} in terms of the cyclic carries \(\sigma_i=\sigma_h(i)=\lfloor\sum_{j\ge1}h_{i+j}\ell^{-j}\rfloor\), the low windows \(W_i\) (\([1,\eta'_i]\) if \(\varsigma_i=1\) and \([h_i-\eta'_i,h_i-1]\) if \(\varsigma_i=-1\)), and the reflected jumps \(\bar x\) (\(x\) if \(\varsigma_i=1\) and \(h_i-x\) if \(\varsigma_i=-1\)):
\[
 \eta'_i+4\sigma_i+6<\ell,\qquad
 |x-y|>\sigma_i+2\quad(x\in U_i\cap W_i,\ y\in[w_i,w_i+t-1]),
\]
and \(\bar x\not\equiv m\pmod\ell\) for every \(x\in U_i\setminus W_i\) and every integer \(m\in[-2\sigma_i-2,\ \eta'_i+2\sigma_i+3]\). Let \(k=O/\varpi O\) be the residue field, put \(r=\dim_{L'}V\), \(W_0=\chi_1^{-1}T_0\), and
\[
 d_-=\sum_{i\in\Z/d'\Z}\sum_{u<w_i}\operatorname{mult}_{U_i}(u),
\]
the total multiplicity of negative Hodge--Tate jumps of \(\chi_1^{-1}V\) across all embeddings \(\tau_i\). For a finite \(O\)-module \(\mathcal M\) with trivial \(G_K\)-action, write \(H^1_{\rm ur}(K,\mathcal M)=\ker\bigl(H^1(K,\mathcal M)\to H^1(I_K,\mathcal M)\bigr)\), where \(I_K\subset G_K\) is the inertia subgroup. Then:
\begin{enumerate}
\item[(a)] Admissible finite \(O[G_K]\)-modules are closed under submodules, quotients, and finite direct sums, and \(T_0\otimes_O\mathcal M\) and \(O(\chi_c)\otimes_O\mathcal M\) (\(1\le c\le t\)) are admissible for every finite \(O\)-module \(\mathcal M\). For \(W=T_0\) or \(W=O(\chi_c)\) (\(1\le c\le t\)) and \(W'=\chi_1^{-1}W\), the subset \(\operatorname{Adm}(W,\mathcal M)\subset H^1(K,W'\otimes_O\mathcal M)\) coincides with the set of classes represented by cocycles \(z\) such that \(\mathcal E_{z,a}\) is admissible for some (equivalently, every) integer \(a\ge1\) with \(\varpi^a\mathcal M=0\); it is an \(O\)-submodule, functorial in \(\mathcal M\), and satisfies \(i_*^{-1}(\operatorname{Adm}(W,\mathcal M'))=\operatorname{Adm}(W,\mathcal M)\) for every injective \(O\)-linear map \(i:\mathcal M\hookrightarrow\mathcal M'\). A cocycle \(z:G_K\to W\otimes_O\mathcal M\) is admissible in the sense of Definition~\ref{B8.42bo} if and only if its cohomology class belongs to \(\operatorname{Adm}(W,\mathcal M)\).
\item[(b)] Over the residue field \(k\), one has
\[
 \dim_k\operatorname{Adm}(T_0,k)\le d_-,\qquad
 \dim_k\operatorname{Adm}(O(\chi_c),k)=0\quad(2\le c\le t),\qquad
 \dim_k\operatorname{Adm}(O(\chi_1),k)=1 .
\]
\item[(c)] Furthermore, \(H^0(K,W_0/\varpi W_0)=H^2(K,W_0/\varpi W_0)=0\), the \(O\)-module \(H^1(K,W_0)\) is finite free of rank \(d'r\) with \(H^1_f(K,W_0)\) an \(O\)-direct summand of rank \(d_-\), the natural map \(H^1(K,W_0)\otimes_O\mathcal M\to H^1(K,W_0\otimes_O\mathcal M)\) is an isomorphism for every finite \(O\)-module \(\mathcal M\), and
\begin{align*}
 \operatorname{Adm}(T_0,\mathcal M)&=\operatorname{im}\bigl(H^1_f(K,W_0)\otimes_O\mathcal M\to H^1(K,W_0\otimes_O\mathcal M)\bigr),\\
 \operatorname{Adm}(O(\chi_c),\mathcal M)&=\{0\}\quad(2\le c\le t),\qquad
 \operatorname{Adm}(O(\chi_1),\mathcal M)=H^1_{\rm ur}(K,\mathcal M).
\end{align*}
\end{enumerate}
\end{lemma}
\begin{proof}[Outline]
Part~(a) uses lattice presentations; (b) combines modulo-\(\varpi\) Wach models with Lemmas~\ref{B8.42boe}, \ref{B8.42bobf}, \ref{B8.42bobg}, and \ref{B8.42bobi}. Part~(c) uses Lemmas~\ref{B8.42bobh}, \ref{B8.42bobd}, and \ref{B8.42bobe} with devissage. The complete proof is in \S\ref{prf:B8.42bob}.
\end{proof}

\begin{remark}[Range of Hodge--Tate weights relative to \(\ell\) in the global application]\label{B8.42boc}
In Lemma~\ref{B8.42bp}, Lemma~\ref{B8.42bob} and Proposition~\ref{B8.42bo1} are applied at primes \(\ell>Z\) (with \(\ell\ge11\)) outside the exceptional prime divisors of \eqref{eq:B8-exception-intervals}. There the shifted jump spread satisfies \(h_i\le b+S_c+2\eta_{\rm sz}+3\), so for \(Z<\ell\le b+S_c+2\eta_{\rm sz}+3\) the Hodge--Tate weight spread can exceed \(\ell\), i.e. lie outside the Fontaine--Laffaille range (while (P4) at \(i-1\) gives \(\lfloor h_i/\ell\rfloor\le\sigma_{i-1}\le(\ell-7)/4\), hence \(h_i<\ell(\ell-3)/4\)); nonetheless \(\sigma_i\le a_*=\lfloor2b/Z\rfloor\) and \(\eta'_i\le2\eta_{\rm sz}+1\) satisfy \(\eta'_i+4\sigma_i+6\le2\eta_{\rm sz}+8b/Z+7<Z<\ell\), and excluding prime divisors of \eqref{eq:B8-exception-intervals} enforces (P6).
\end{remark}

\begin{proposition}[Admissible cocycles at large primes]\label{B8.42bo1}
Assume \(\mathscr H_2\). Then statement (T) of Definition~\ref{B8.42bo} holds for all data of that definition.
\end{proposition}
\begin{proof}[Outline]
Identify twisted-action and twisted-coefficient cocycles, applying Lemma~\ref{B8.42bob}(a),(c) to \(W_0=\chi_1^{-1}T_0\) for (T1) via Definition~\ref{B8.41h}. Transport along \(O(\chi_c)\simeq L_c\) and apply Lemma~\ref{B8.42bob}(a),(c) to obtain (T2) and (T3). The complete proof is in \S\ref{prf:B8.42bo1}.
\end{proof}

\begin{lemma}[Zero local loss off the exceptional primes]\label{B8.42bp}
Work on the twist-good cut family of Theorem~\ref{B8.42}. At a member put \(n=f/2\) and \(\eta=\eta_{\rm sz}\), let \(N(m)\) be the terminal rank, and set
\[
 Z=\frac{mb}{nt^2(N(m)+t)L^2+kn^2(S_c+t+L)},\qquad
 a_*=\lfloor2b/Z\rfloor .
\]
The base sizes are \(b\), \(w''(n)\) and \(b_T+\lfloor b'/x\rfloor\) at top positions \(T\), where \(w''(n),b_T,b'\) are the ordinary lower endpoint, top-position size and auxiliary scale. A prime \(\ell>Z\) is exceptional if it divides an integer in one of the intervals
\begin{equation}\label{eq:B8-exception-intervals}
 [W-\eta-2a_*-3,\ W+S_c+2\eta+2a_*+4]
\end{equation}
with \(W\) a base size.

(a) There is a family cutoff beyond which every prime \(\ell>Z\) satisfies \(\ell\ge11\), \(t+1<\ell\) and \(\ell\notin\{q,d\}\), lies outside the small set, and is unramified in \(E\). If moreover \(\ell\) is not exceptional, then for every layer \(1\le j\le m\), every \(d'\ge1\), and every choice of signs \(\varsigma_i\) and weight types at \(i\in\Z/d'\Z\), there are position data \((h,w,\eta',\varsigma,U)\) at \(\ell\) of width \(t\) with \(\operatorname{mult}_{U_i}(b+w_i)=\mathbf 1_{J_i}(b)\) for all \(i\) and \(b\). Here \(J_i\) is the jump set of layer \(j\) for the weight type at \(i\) and the sign \(\varsigma_i\), described in the proof.

(b) Assume \(\mathscr H_2\) and statement (T) of Definition~\ref{B8.42bo}. Beyond the cutoff of (a), let \(i\) be a twisted index whose residue characteristic \(\ell=\ell_i>Z\) is not exceptional and lies outside \(\{g_0,g_1\}\). Let \(\pi\in Z_{\rm good}(i)\) and \(w\mid\ell\). Let \(L'\) be a finite extension of the fraction field \(K_i\) of \(O_i\) containing the images of all embeddings of \(E_w(\mu_\ell)\), with integers \(O\) and uniformizer \(\varpi\). For \(x\in X_i\) write \(R_x\) for the full representation at \(x\) and \(V_x=R_x\otimes\chi_{i,1}^{-1}\), as in Lemma~\ref{B8.42bd}. For an \(O_i[G_E]\)-lattice \(M\) of rank \(k_M\) with action \(\mu\), a cut cocycle of \(M\) at \(w\) is a map \(G_{E_w}\to M\otimes K_i/O_i\) of the form \(g\mapsto\phi(gv_0-v_0)\). Here \(x_1,\dots,x_{a'}\in X_i\), \(v_0\in\bigoplus_jV_{x_j}\), \(\mathcal Y\subset\bigoplus_jV_{x_j}\) is a finitely generated \(O_i\)-submodule stable under \(G_E\) and containing \(gv_0-v_0\) for all \(g\in G_E\), the map \(\phi:\mathcal Y\to M\otimes K_i/O_i\) is \(O_i\)-linear with \(\phi(gy)=\mu(g)\phi(y)\), and \(G_{E_w}\) acts through the decomposition embedding. Let \(M\) be the normalized rank-\(r\) lattice \(W_\pi\) of \(\rho_i(\pi)\chi_{i,1}^{-1}\), or the character \(\chi_{i,c}\chi_{i,1}^{-1}\) with \(1\le c\le t\). Let \(\kappa\) be a cut cocycle of \(M\) at \(w\) and \(\bar\kappa\) its image in \(M\otimes L'/O\). Then \(\bar\kappa\) is a finite cocycle in the sense of Definition~\ref{B8.41h} if \(M=W_\pi\), a coboundary on \(G_{E_w}\) if \(c\ge2\), and a coboundary on the inertia group of \(E_w\) if \(c=1\).
\end{lemma}
\begin{proof}[Outline]
For (a), combines family size bounds with \ref{B2.46} and \ref{B8.42bn} to build position data. For (b), twists by a Lubin--Tate character (\ref{B8.42bd}), extends congruences via \ref{B8.42pa} and \ref{B8.42p}, proves admissibility by Artin--Rees, and applies statement (T). The complete proof is in \S\ref{prf:B8.42bp}.
\end{proof}

\begin{definition}[The annihilator statement at a ramified prime]\label{B8.42bq}
Let \(\ell\) be odd, \(K_0/\Q_\ell\) finite unramified of degree \(n_K\), and \(K\subset K_0(\mu_\ell)\) an extension of \(K_0\) of degree \(e\). Put \(s=(\ell-1)/e\). Let \(L'\) be a finite extension of \(\Q_\ell\) containing the images of all embeddings of \(K_0(\mu_\ell)\), with integers \(O\). For a stable \(O\)-lattice \(T\) in a continuous representation of \(G_K\) over \(L'\), let \(Z_f(T)\) be the \(O\)-module of continuous cocycles \(G_K\to T\) whose classes are finite in the sense of Definition~\ref{B8.41h}, and \(B(T)\) the module of coboundaries.

Admissible data at \(K\) consist of the following. Crystalline lattices \(T_1,T_2\) and an equivariant surjection \(F:T_1\to T_2\), with \(r=\operatorname{rank}T_2\) and \(I_F=F(Z_f(T_1))+B(T_2)\). Integers \(S,\eta,m,h\ge0\). For each embedding \(\tau:K\to L'\) a sign \(\varsigma_\tau\) and an integer \(\beta_\tau>\eta\). They are required to satisfy three conditions. The jumps of \(T_1\) at \(\tau\) lie in \([-h,h]\), and in \([\beta_\tau,\beta_\tau+S]\cup[-\eta,\eta]\) if \(\varsigma_\tau=1\), in \([-\beta_\tau-S,-\beta_\tau]\cup[-\eta,\eta]\) if \(\varsigma_\tau=-1\). At each \(\tau\), the lattice \(T_2\) has at most \(m\) jumps in \([-\eta,\eta]\), counted with multiplicity. If \(e\ne2\), the sign \(\varsigma_\tau\) depends only on the restriction of \(\tau\) to \(K_0\).

Statement (A) at \(K\) asserts the following for all admissible data at \(K\). There are an integer \(h'\) with \(1\le h'\le h+S+1\) and a set \(\Sigma_0\) of at most \(n_Ke(e+1)m\) elements of \(Z_f(T_2)\) with
\[
 \ell^Y Z_f(T_2)\subset O\Sigma_0+I_F ,
\]
where \(Y\) is the integer of Lemma~\ref{B8.42bc} for \((\ell,e,s,h',S,\eta)\).
\end{definition}

\begin{lemma}[Normalization, \(\Delta\)-invariant ring relations, and component splitting]\label{B8.42bqa}
Assume \(\mathscr H_2\). Let \(\ell,K_0,K,n_K,e,s,L',O\) be as in Definition~\ref{B8.42bq}, put \(K'=K_0(\mu_\ell)\), \(\Delta=\Gal(K'/K)\), and \(\Gamma=\Gal(K'(\mu_{\ell^\infty})/K')\), and fix a topological generator \(\gamma\in\Gamma\) with \(v_\ell(\chi_\ell(\gamma)-1)=1\).
\begin{enumerate}
\item[(a)] The degree \(e=[K:K_0]\) divides \(\ell-1\), so \(es=\ell-1\) and \(s\ge1\). For any admissible data \((T_1,T_2,F,r,S,\eta,m,h,\varsigma_\tau,\beta_\tau)\) at \(K\), every jump of \(T_2\) at an embedding \(\tau:K\to L'\) is a jump of \(T_1\) at \(\tau\), and there exist integers \(0\le\eta_1\le\eta\), \(1\le h'\le h+S+1\), and \(\beta'_\tau\) (\(\tau:K\to L'\)) with \(\eta_1<\beta'_\tau\) and \(\beta'_\tau+S\le h'\) such that at every \(\tau\), the jumps of both \(T_1\) and \(T_2\) lie in \([\beta'_\tau,\beta'_\tau+S]\cup[-\eta_1,\eta_1]\) when \(\varsigma_\tau=1\) and in \([-\beta'_\tau-S,-\beta'_\tau]\cup[-\eta_1,\eta_1]\) when \(\varsigma_\tau=-1\) (hence in \([-h',h']\)), and \(T_2\) has at most \(m\) jumps in \([-\eta_1,\eta_1]\), counted with multiplicity.
\item[(b)] Let \(A_*=\prod_{c\in\Z/n_K\Z}O[[Z]]\) be the level-one Wach ring over \(K'\). There exist a \(\Delta\)-invariant parameter \(\pi'=(\pi'_c)_{c\in\Z/n_K\Z}\in A_*^\Delta\) and degree-\(e\) polynomials \(E_c\in O[U]\) (\(c\in\Z/n_K\Z\)) such that \(A_*^\Delta=\prod_{c\in\Z/n_K\Z}O[[\pi'_c]]\), \(\pi'\) is a unit of \(A_*\) times \(Z^s\), the cyclotomic polynomial \(Q=\Phi_\ell(1+Z)\) is a unit of \(A_*\) times \(E(\pi')=(E_c(\pi'_c))_c\), \(\varphi(\pi')\) is a unit times \(\pi'E(\pi')^s\), and \((\gamma-1)\pi'\) is a unit times \(\pi'E(\pi')\).
\item[(c)] For \(h'\ge1\), put \(M=\lceil 4h'/s\rceil\), \(L(k)=k+v_\ell(k!)\), \(\rho_A=L(2\lceil 2eh'/(\ell-1)\rceil-1)\), and \(\sigma_A=L(M-1)\). Let \(\mathscr E\) be the \(\ell\)-adic completion of \(A_*[Z^{-1}]\). For a crystalline \(O\)-lattice \(T\) with jumps in \([-h',h']\), let \(D(T)\) be its étale \((\varphi,\Gamma)\)-module over \(\mathscr E\), with Wach lattice \(N(T)\). Let \(\epsilon_c\) be the idempotent of the \(c\)-th factor of \(A_*\), and put \(E_\pi=E(\pi')\) and \(N_c(T)^\Delta=\epsilon_cN(T)^\Delta\). Define
\[
 \begin{split}
 \Phi_c(T)&=A_*\text{-span}\{\varphi z:z\in\epsilon_{c-1}N(T)\},\\
 \Lambda_c^P(T)&=\{w\in\epsilon_cD(T)^\Delta:E_\pi^{h'}w\in\Phi_c(T)\},\\
 U_k(T)&=\pi'^kN(T)^\Delta,\\
 \mathcal P_c^{(k)}(T)&=
 \frac{\{w\in\Lambda_c^P(T):(\gamma-1)w\in U_k(T)\}}
 {\pi'^kN_c(T)^\Delta}\qquad(k\ge1).
 \end{split}
\]
Here \(N_c(T)^\Delta\subset\Lambda_c^P(T)\), and the denominator is contained in the numerator. For every equivariant map \(f:T\to T'\) of crystalline \(O\)-lattices with jumps in \([-h',h']\), write \(f_*:\mathcal P_c^{(k)}(T)\to\mathcal P_c^{(k)}(T')\) for the map induced by \(D(f)\). A \emph{Herr pair} is a pair \((w,z)\in D(T)^\Delta\times D(T)^\Delta\) with \((\gamma-1)w=(\varphi-1)z\). Under the étale comparison \(D(T)\subset(W(C^\flat)\otimes_{\Z_\ell}T)^{G_{K'(\mu_{\ell^\infty})}}\) of Lemma~\ref{B8.41cb}, choosing \(\tilde b\in W(C^\flat)\otimes_{\Z_\ell}T\) with \((\varphi-1)\tilde b=w\) and letting \(g\mapsto z_g\in D(T)^\Delta\) be the continuous \(1\)-cocycle on \(G_K\) factoring through \(\Delta\times\Gamma\) with \(z_g=0\) for \(g\in\Delta\) and \(z_\gamma=z\) defines a cocycle \(g\mapsto z_g-(g-1)\tilde b\) in \(T\); its cohomology class \(\operatorname{cl}_{K,T}(w,z)\in H^1(K,T)\) is \(O\)-linear, natural in \(T\), and vanishes on \(((\varphi-1)v,(\gamma-1)v)\) (\(v\in D(T)^\Delta\)). The map \(\varphi-1:U_M(T)\to U_M(T)\) is bijective. Every numerator vector \(w\) for \(\mathcal P_c^{(M)}(T)\) therefore has a unique \(z\in U_M(T)\) completing it to a Herr pair, and that pair has a class in \(H_f^1(K,T)\), the image of \(Z_f(T)\) in \(H^1(K,T)\). This defines an \(O\)-linear map
\[
 \operatorname{cl}_c^H:\mathcal P_c^{(M)}(T)\longrightarrow H_f^1(K,T),
 \qquad [w]\longmapsto\operatorname{cl}_{K,T}(w,z),
\]
which is natural under equivariant maps \(f:T\to T'\) of such crystalline lattices (\(f_*\circ\operatorname{cl}_c^H=\operatorname{cl}_c^H\circ f_*\)). Let \(\mathcal J_c(T)\) be the \(O\)-span of all \(\kappa\in Z_f(T)\) for which there are such \(w,z\) with \([\kappa]=\operatorname{cl}_{K,T}(w,z)\). Then
\[
 \ell^{\rho_A+\sigma_A}Z_f(T)\subset\sum_{c\in\Z/n_K\Z}\mathcal J_c(T)+B(T).
\]
\end{enumerate}
\end{lemma}
\begin{proof}[Outline]
Tower degree and jump additivity give (a), and tame descent gives (b). For (c), Frobenius contraction on \(U_M(T)\) defines \(\operatorname{cl}_c^H\), and \(\Delta\)-averaging with Lagrange interpolation and \(\gamma\)-shifts decomposes \(\ell^{\rho_A+\sigma_A}Z_f(T)\). The complete proof is in \S\ref{prf:B8.42bqa}.
\end{proof}

\begin{definition}[Coordinate lattices and evaluation maps]\label{B8.42bqd}
Use the local data of Lemma~\ref{B8.42bqa}, fix an integer \(h'\ge1\) with \(M=\lceil4h'/s\rceil\), fix \(c\in\Z/n_K\Z\), and let \(T\) be a crystalline \(O\)-lattice over \(K\) with jumps in \([-h',h']\). Put \(A^\Delta=A_*^\Delta\), \(E_\pi=E(\pi')\), and \(V=T[1/\ell]\). Fix compatible primitive roots \((\zeta_{\ell^n})_{n\ge1}\) as in the proof of Lemma~\ref{B8.41ca} (at depth \(1\)), so \(\zeta_\ell\in K'\) identifies \(W(k')[[Z]]/(Q)\) with \(\mathcal O_{K'}\) via \(1+Z\mapsto\zeta_\ell\), and let \(B_{\rm dR}=\operatorname{Frac}(B_{\rm dR}^+)\) be Fontaine's de Rham period field, where \(B_{\rm dR}^+\) is the \(\ker\theta\)-adic completion of \(W(\mathcal O^\flat)[1/\ell]\) with uniformizer \(t_{\rm cris}=\log[\epsilon]\). Put \(\mathscr V_T=V\otimes_{\Q_\ell}B_{\rm dR}\); then \(D_{\rm dR}(V)=\mathscr V_T^{G_K}=\bigoplus_{\tau:K\to L'}D_{\rm dR}(V)_\tau\subset\mathscr V_T\), with \(\tau\)-component projections \(\mathrm{pr}_\tau:D_{\rm dR}(V)\to D_{\rm dR}(V)_\tau\). Let \(\iota_c:K_0\to L'\) be the embedding indexing the component \(c\), and let \(\epsilon_c\) be its idempotent in \(A_*\). An embedding \(\tau:K\to L'\) is \emph{above \(c\)} if \(\tau|_{K_0}=\iota_c\). An \emph{evaluation pair above \(c\)} is a pair \((\zeta,\tau)\) for which an embedding \(\widetilde\tau:K'\to L'\) satisfies
\[
 \widetilde\tau|_{K_0}=\iota_c,\qquad
 \widetilde\tau|_K=\tau,\qquad
 \widetilde\tau(\zeta_\ell)=\zeta.
\]
For \(a=(a_i)_i\in A_*\), write
\[
 \operatorname{ev}_{c,\zeta}(a)(t)=a_c(\zeta\exp(t)-1)\in L'[[t]],
\]
with coefficientwise \(\ell\)-adic evaluation. Under the level-one prismatic comparison of Lemmas~\ref{B8.41ca}--\ref{B8.41cb} (\(1+Z\mapsto[\epsilon^{1/\ell}]=\zeta_\ell\exp(t_{\rm cris}/\ell)\) in \(B_{\rm dR}^+\)), completed along \(Z=\zeta\exp(t)-1\) via \(\widetilde\tau\) with \(t\) tracking \(t_{\rm cris}/\ell\), each \(z\in N(T)\) has a Laurent expansion \(\operatorname{ev}^{N}_{c,\zeta,T}(z)\in\mathscr V_T((t))\) (lying in \(D_{\rm dR}(V)((t))\) when \(z\in N(T)^\Delta\)). Write \([t^k]F\) for the coefficient of \(t^k\) in a Laurent series \(F\) (with \([t]=[t^1]\)). For \(z\in D(T)\), if there is a nonnegative integer \(k\) with \(Q^k\varphi z\in N(T)\), take the least such \(k\) and set
\[
 v_{c,\zeta,T}(z)=
 \bigl([t]\operatorname{ev}_{c,\zeta}(Q)\bigr)^{-k}
 [t^k]\operatorname{ev}^{N}_{c,\zeta,T}(Q^k\varphi z).
\]
Set this value equal to zero when there is no such \(k\). Define the decreasing filtration
\[
 \mathscr F_\tau^a(T)=
 \operatorname{Fil}^aD_{\rm dR}(V)_\tau
 +\sum_{\tau'\ne\tau}D_{\rm dR}(V)_{\tau'}\subset\mathscr V_T
 \qquad(a\in\Z).
\]
The residue module and its scalar extension are
\[
 L_c(T)=\frac{N_{c-1}(T)^\Delta}
 {N_{c-1}(T)^\Delta\cap\pi'N(T)},\qquad
 \mathscr B_c(T)=\mathscr E\otimes_O L_c(T).
\]
For \(z\in N_{c-1}(T)^\Delta\), write \(\bar z\) for its image in \(L_c(T)\). For an \(O\)-submodule \(S_0\subset L_c(T)\), set
\[
 \begin{split}
 R_c(S_0)&=\operatorname{im}\bigl(A^\Delta\otimes_O S_0\longrightarrow\mathscr B_c(T),
                 \ a\otimes s\longmapsto\epsilon_ca\otimes s\bigr),\\
 \Lambda_c^L(T;S_0)&=\{x\in\mathscr B_c(T):E_\pi^{h'}x\in R_c(S_0)\},\qquad
 \Lambda_c^L(T)=\Lambda_c^L(T;L_c(T)).
 \end{split}
\]
Multiplication by \(A^\Delta\) acts on the first tensor factor, as does \(\gamma_L=\gamma_{\mathscr E}\otimes1\), where \(\gamma_{\mathscr E}\) is the \(\gamma\)-action on \(\mathscr E\). Let \(\Phi_c(T)\) and \(\Lambda_c^P(T)\) be the modules in Lemma~\ref{B8.42bqa}(c), and put
\[
 \begin{split}
 P_c^\varphi(T)&=\Phi_c(T)\cap\epsilon_cD(T)^\Delta,\\
 J_c^L(T)&=\pi'^ME_\pi^{3h'}R_c(L_c(T)),\qquad
 J_c^P(T)=\pi'^ME_\pi^{3h'}P_c^\varphi(T),\\
 \overline\Lambda_c^L(T)&=\Lambda_c^L(T)/(\Lambda_c^L(T)\cap J_c^L(T)),\\
 \overline\Lambda_c^P(T)&=\Lambda_c^P(T)/(\Lambda_c^P(T)\cap J_c^P(T)).
 \end{split}
\]
Write \(q_L\) and \(q_P\) for these quotient maps. An equivariant map \(f:T\to T'\) of crystalline \(O\)-lattices induces \(f_P:D(T)\to D(T')\), its reduction \(f_0:L_c(T)\to L_c(T')\), and \(f_L=1\otimes f_0:\mathscr B_c(T)\to\mathscr B_c(T')\).
\end{definition}
\begin{proof}[Outline]
Since \(v_\ell(\zeta-1)>0\), \(\operatorname{ev}_{c,\zeta}\) converges with \([t]\operatorname{ev}_{c,\zeta}(Q)\ne0\), making \(v_{c,\zeta,T}\) well-defined. Bilinearity defines the coordinate modules and quotients, and equivariance of \(f_P\) induces \(f_0\) and \(f_L\). The complete proof is in \S\ref{prf:B8.42bqd}.
\end{proof}

\begin{definition}[Model lattices, Frobenius-horizontal frames, and slot spaces]\label{B8.42bqb}
Use Definition~\ref{B8.42bqd}, with \(h'\ge1\), \(M=\lceil4h'/s\rceil\), and \(\sigma_A=L(M-1)\).
\begin{enumerate}
\item[(a)] For a crystalline \(O\)-lattice \(T\) over \(K\) with jumps in \([-h',h']\), a \emph{representation} of \(x\in\mathscr B_c(T)\) consists of finitely many \(a_j\in A^\Delta\) and \(z_j\in N_{c-1}(T)^\Delta\) such that
\[
 E_\pi^{h'}x=\sum_j\epsilon_ca_j\otimes\bar z_j.
\]
Its numerator expansion at \((\zeta,\tau)\) is
\[
 G_{c,\zeta}(a,z;t)=\sum_j\operatorname{ev}_{c,\zeta}(a_j)(t)v_{c,\zeta,T}(z_j).
\]
The \emph{model lattice} \(\mathcal M_c(T)\) is the \(O\)-span of the elements admitting such a representation for which
\[
 [t^i]G_{c,\zeta}(a,z;t)\in\mathscr F_\tau^{h'-i}(T)
 \quad\text{for every evaluation pair above \(c\) and every \(i\in\Z\)}.
\]
Coefficients with negative index in a power series are zero.

\item[(b)] A \emph{Frobenius-horizontal frame} is a family of \(O\)-linear maps
\[
 \alpha_T:\Lambda_c^L(T)\longrightarrow\overline\Lambda_c^P(T),\qquad
 \beta_T:\Lambda_c^P(T)\longrightarrow\overline\Lambda_c^L(T),
\]
for crystalline lattices \(T\) with jumps in \([-h',h']\), subject to the following conditions. They kill \(\Lambda_c^L(T)\cap J_c^L(T)\) and \(\Lambda_c^P(T)\cap J_c^P(T)\), respectively. Their induced quotient maps satisfy
\[
 \bar\beta_T\bar\alpha_T=\ell^{2\sigma_A}\operatorname{id},\qquad
 \bar\alpha_T\bar\beta_T=\ell^{2\sigma_A}\operatorname{id}.
\]
They satisfy the following representative compatibilities whenever the arguments lie in the indicated domains. If \(\alpha_T(x)=q_P(y)\), then
\[
 \alpha_T(ax)=q_P(ay)\quad(a\in A^\Delta),\qquad
 \alpha_T(\gamma_Lx)=q_P(\gamma y).
\]
If \(\beta_T(y)=q_L(x)\), then
\[
 \beta_T(ay)=q_L(ax)\quad(a\in A^\Delta),\qquad
 \beta_T(\gamma y)=q_L(\gamma_Lx).
\]
For every equivariant map \(f:T\to T'\) of crystalline \(O\)-lattices with jumps in \([-h',h']\), these identities are natural under \(f_L,f_P\): the first equality \(\alpha_T(x)=q_P(y)\) implies \(\alpha_{T'}(f_Lx)=q_P(f_Py)\), and the second equality \(\beta_T(y)=q_L(x)\) implies \(\beta_{T'}(f_Py)=q_L(f_Lx)\). Finally, whenever \(\alpha_T(x)=q_P(y)\), \((a_j,z_j)_j\) is a representation of \(x\), and \(k\ge0\) satisfies \(Q^kE_\pi^{h'}y\in N(T)\), they satisfy
\[
 [t^i]\operatorname{ev}^{N}_{c,\zeta,T}(Q^kE_\pi^{h'}y)
 =\ell^{\sigma_A}[t^i]\bigl(\operatorname{ev}_{c,\zeta}(Q^k)G_{c,\zeta}(a,z;t)\bigr)
 \qquad(i<k+4h')
\]
for every evaluation pair above \(c\).

\item[(c)] For an \(O\)-submodule \(S_0\subset L_c(T)\), define
\[
 \begin{split}
 \mathcal I_c^L(T;S_0)&=\{x\in\Lambda_c^L(T;S_0):
                         (\gamma_L-1)x\in\pi'\mathcal M_c(T)\},\\
 \mathcal R_c(T;S_0)&=\frac{\mathcal I_c^L(T;S_0)}
 {\mathcal I_c^L(T;S_0)\cap\pi'\mathcal M_c(T)},\qquad
 \mathcal H_c(T)=\mathcal R_c(T;L_c(T)).
 \end{split}
\]
Given a crystalline \(O\)-lattice \(T_2\) with jumps in \([-h',h']\) and a chosen \(O\)-submodule \(L_r\subset L_c(T_2)\), the \emph{residual module} is \(\mathcal I_{\rm rest}=\mathcal R_c(T_2;L_r)\).

\item[(d)] Fix the normalized windows \((S,\eta_1,\varsigma_\tau,\beta'_\tau)\) of Lemma~\ref{B8.42bqa}(a); call \(\tau\) (and any evaluation pair \((\zeta,\tau)\)) \emph{upward} if \(\varsigma_\tau=1\) and \emph{downward} if \(\varsigma_\tau=-1\). A scalar \(f\in\mathscr E\) is an \emph{admissible slot coefficient} if there exist \(a,g\in A^\Delta\) with
\[
 \epsilon_ca=a,\qquad E_\pi^{h'}f=a,\qquad
 E_\pi^{h'}(\gamma f-f)=\pi'g,
\]
and, at every evaluation pair above \(c\), the following coefficient conditions hold:
\[
 \begin{array}{ll}
 \varsigma_\tau=1:&[t^i]\operatorname{ev}_{c,\zeta}(a)=0
                    \quad(0\le i<h'-\beta'_\tau),\\
 \varsigma_\tau=-1:&\text{there is }b_{\zeta,\tau}\in L'\text{ such that}
 \\
 & [t^i]\operatorname{ev}_{c,\zeta}(a)
   =b_{\zeta,\tau}[t^i]\operatorname{ev}_{c,\zeta}(E_\pi^{h'})
       \quad(0\le i<h'+\beta'_\tau+S).
 \end{array}
\]
An \emph{elementary slot tensor} for \(T\) is \(f\otimes\bar z\), where \(f\) is such a coefficient, \(z\in N_{c-1}(T)^\Delta\), and
\[
 v_{c,\zeta,T}(z)\in\mathscr F_\tau^{\beta'_\tau}(T)
 \quad\text{whenever \((\zeta,\tau)\) is above \(c\) and \(\varsigma_\tau=1\)}.
\]
Let \(\mathcal T_c(T)\subset\Lambda_c^L(T)\) be their \(O\)-span. When \(\mathcal T_c(T)\subset\mathcal I_c^L(T;L_c(T))\), write \(\mathcal C_{\rm lift}(T,c)\) for its image in \(\mathcal H_c(T)\).
\end{enumerate}
\end{definition}

\begin{lemma}[Model comparison and invariant coordinates]\label{B8.42bqe}
Assume \(\mathscr H_2\), and let \(T\) be crystalline with jumps in \([-h',h']\). Use Definitions~\ref{B8.42bqd} and~\ref{B8.42bqb}.
\begin{enumerate}
\item[(a)] The module \(\mathcal M_c(T)\) is stable under \(A^\Delta\) and \(\gamma_L\), and
\[
 E_\pi^{h'}R_c(L_c(T))\subset\mathcal M_c(T)\subset\Lambda_c^L(T),
 \qquad J_c^L(T)\subset\pi'^M\mathcal M_c(T).
\]
It is \(\ell\)-saturated in \(\Lambda_c^L(T)\): if \(x\in\Lambda_c^L(T)\), \(a\ge0\), and \(\ell^ax\in\mathcal M_c(T)\), then \(x\in\mathcal M_c(T)\). The quotient
\[
 \mathcal Q_c(T)=\Lambda_c^L(T)/(\Lambda_c^L(T)\cap\pi'\mathcal M_c(T))
\]
is finite free over \(O\).
\item[(b)] Suppose \(x\in\Lambda_c^L(T)\) and \((\gamma_L-1)x\in\mathcal M_c(T)\). For a representation \((a_j,z_j)_j\) of \(x\), let \(f(t)\in\mathscr V_T((t))\) satisfy
\[
 \operatorname{ev}_{c,\zeta}(E_\pi^{h'})f(t)=G_{c,\zeta}(a,z;t).
\]
Then \([t^i]f(t)\in\mathscr F_\tau^{-i}(T)\) for every evaluation pair \((\zeta,\tau)\) above \(c\) and every \(i\ne0\).
\item[(c)] For any Frobenius-horizontal frame, \(\alpha_T(\mathcal M_c(T))\) has representatives in \(N_c(T)^\Delta\), and \(\beta_T(N_c(T)^\Delta)\) has representatives in \(\mathcal M_c(T)\). For \(1\le j\le M\), \(x\in\Lambda_c^L(T)\), and \(y\in\Lambda_c^P(T)\),
\[
 \begin{split}
 \alpha_T(x)=q_P(y),\ x\in\pi'^j\mathcal M_c(T)
      &\Longrightarrow y\in\pi'^jN_c(T)^\Delta,\\
 \beta_T(y)=q_L(x),\ y\in\pi'^jN_c(T)^\Delta
      &\Longrightarrow x\in\pi'^j\mathcal M_c(T).
 \end{split}
\]
Moreover, for \(x\in\Lambda_c^L(T)\) and \(y\in\Lambda_c^P(T)\), \(\alpha_T(x)=q_P(y)\) with \(x\in\mathcal I_c^L(T;L_c(T))\) implies \((\gamma-1)y\in\pi'N_c(T)^\Delta\), and \(\beta_T(y)=q_L(x)\) with \((\gamma-1)y\in\pi'N_c(T)^\Delta\) implies \(x\in\mathcal I_c^L(T;L_c(T))\). Consequently \(\alpha_{1,T}([x])=[y]\) (when \(\alpha_T(x)=q_P(y)\)) and \(\beta_{1,T}([y])=[x]\) (when \(\beta_T(y)=q_L(x)\)) define natural \(O\)-linear maps
\[
 \alpha_{1,T}:\mathcal H_c(T)\longrightarrow\mathcal P_c^{(1)}(T),\qquad
 \beta_{1,T}:\mathcal P_c^{(1)}(T)\longrightarrow\mathcal H_c(T),
\]
with \(\alpha_{1,T}\beta_{1,T}=\ell^{2\sigma_A}\operatorname{id}\) and \(\beta_{1,T}\alpha_{1,T}=\ell^{2\sigma_A}\operatorname{id}\).
\end{enumerate}
\end{lemma}
\begin{proof}[Outline]
Filtration stability, \(E_\pi\)-division, and saturation give (a). Induction on Laurent exponents of \(f(ut)-f(t)\) proves (b), and \ref{B8.42bqb}(b) with \(\gamma\)-commutation matches \(\pi'\)-powers to define \(\alpha_{1,T}\) and \(\beta_{1,T}\). The complete proof is in \S\ref{prf:B8.42bqe}.
\end{proof}

\begin{lemma}[Division at an integral evaluation root]\label{B8.42bqca}
Use the local ring data of Lemma~\ref{B8.42bqa}(b) and the evaluation maps of Definition~\ref{B8.42bqd}. Thus \(O\) is the ring of integers of the finite extension \(L'/\Q_\ell\), \(E_c\) is the polynomial defining \(E_\pi\) on component \(c\), and each evaluation uses a primitive \(\ell\)-th root of unity. Assume \(e=2\), put \(s=(\ell-1)/2\), and let \(r\in O\) satisfy \(\ell\mid r^2\). For every \(a\in A^\Delta\), every component \(c\), and every integer \(d\ge0\), there are \(p_0,\ldots,p_{d-1}\in O\) and \(a_d\in A^\Delta\) such that, at every evaluation pair above \(c\),
\[
 \operatorname{ev}_{c,\zeta}(a)
 =\sum_{j<d}p_j\bigl(\operatorname{ev}_{c,\zeta}(\pi')-r\bigr)^j
  +\bigl(\operatorname{ev}_{c,\zeta}(\pi')-r\bigr)^d
       \operatorname{ev}_{c,\zeta}(a_d).
\]
If \(r=[t^0]\operatorname{ev}_{c,\zeta}(\pi')\), then
\(\operatorname{ev}_{c,\zeta}(\pi')-r=tU(t)\) with \(U(0)\ne0\). Consequently, vanishing through degree \(d-1\) at one such evaluation is equivalent to \(p_0=\cdots=p_{d-1}=0\), and then holds at every evaluation with the same value of \(r\).
\end{lemma}
\begin{proof}[Outline]
Repeated division of \(F(\pi')\) by \(\pi'-r\), convergent in \(O\) since \(\ell\mid r^2\), gives the expansion. The simple zero of \(\Phi_\ell(\zeta\exp(t))\) at \(t=0\) yields \(\operatorname{ev}_{c,\zeta}(\pi')-r=tU(t)\) (\(U(0)\ne 0\)), proving the last claim. The complete proof is in \S\ref{prf:B8.42bqca}.
\end{proof}

\begin{lemma}[The quadratic root separation and its denominator cost]\label{B8.42bqcb}
Continue with \(e=2\). Let \(r_U,r_D\in O\) be distinct evaluation values of \(\pi'\) above a component \(c\), with \(\ell\mid r_U^2\) and \(\ell\mid r_D^2\), and put \(\delta_0=r_D-r_U\). Then
\[
 E_c(X)=a_2(X-r_U)(X-r_D),\qquad
 a_2\in O^\times,\qquad \delta_0^2\mid\ell.
\]
In particular \(v_\ell(\delta_0)\le1/2\). Let \(b,B\) be integers with \(1\le B\le b\), put \(d=b-B+1\), and let \(x\in L'\) satisfy
\(o_0=x\binom b{B-1}\delta_0^d\in O\).
If an integer \(y\ge0\) satisfies
\[
 d+2v_\ell\binom b{B-1}\le2y,
\]
then \(\ell^yx\in O\). More precisely, write
\(v=v_\ell\binom b{B-1}\) and \(\binom b{B-1}=\ell^v c'\), where \(c'\in O^\times\), and choose \(r_0\in O\) with \(\ell=\delta_0^2r_0\). Then
\[
 \ell^yx=o_0(c')^{-1}\delta_0^{2(y-v)-d}r_0^{y-v}.
\]
\end{lemma}
\begin{proof}[Outline]
Factoring \(E_c\) at \(r_U,r_D\) and comparing constant and linear terms in \(Q=w'E_\pi\) proves \(a_2\in O^\times\) and \(\delta_0^2\mid\ell\). Substituting \(o_0\) and \(\ell=\delta_0^2r_0\) under \(2(y-v)\ge d\) verifies the integral formula for \(\ell^y x\). The complete proof is in \S\ref{prf:B8.42bqcb}.
\end{proof}

\begin{proposition}[Frames, slot lifting, and the component bound]\label{B8.42bqc}
Assume \(\mathscr H_2\). In the setting of Definitions~\ref{B8.42bqd} and~\ref{B8.42bqb} and Lemma~\ref{B8.42bqe}, fix normalized data \((T_1,T_2,F,S,\eta_1,m,h',\varsigma_\tau,\beta'_\tau)\) as in Lemma~\ref{B8.42bqa}(a), with the sign condition of Definition~\ref{B8.42bq}, and fix \(c\in\Z/n_K\Z\). Put
\[
 \kappa_A=e\bigl(L(S)+\lfloor\log_\ell h'\rfloor\bigr),\qquad
 y_c(\eta_1)=\left\lceil\eta_1+\tfrac12+\log_\ell h'\right\rceil,\qquad
 y_1=\rho_A+\max\bigl(2\sigma_A,y_c(\eta_1)\bigr).
\]
Define the finite set of integers
\[
 \begin{split}
 \Omega_c={}&\bigcup_{\substack{\tau\text{ above }c\\\varsigma_\tau=1}}
       \{i\in\Z:-\beta'_\tau-S\le i<-\beta'_\tau\}\\
 &\mathbin{\cup}\bigcup_{\substack{\tau\text{ above }c\\\varsigma_\tau=-1}}
       \{i\in\Z:\beta'_\tau\le i<\beta'_\tau+S\},\qquad
 d_c=\prod_{i\in\Omega_c}(1-\chi_\ell(\gamma)^i)\in O.
 \end{split}
\]
Then the following statements hold.
\begin{enumerate}
\item[(a)] A Frobenius-horizontal frame exists. For \(j=1,2\), the reduction
\(r_c:\mathcal P_c^{(M)}(T_j)\to\mathcal P_c^{(1)}(T_j)\) induced by numerator inclusion \(U_M(T_j)\subset U_1(T_j)\) is injective and its image contains \(\ell^{\sigma_A}\mathcal P_c^{(1)}(T_j)\). There is a unique \(O\)-linear map \(\operatorname{lift}_c\) with \(r_c\operatorname{lift}_c=\ell^{\sigma_A}\operatorname{id}\). It satisfies \(\operatorname{lift}_cr_c=\ell^{\sigma_A}\operatorname{id}\) and commutes with \(F_*\).

\item[(b)] For any such frame there are an \(O\)-submodule \(L_r\subset L_c(T_2)\), an integer \(k\ge0\), and vectors \(z_1,\ldots,z_k\in N_{c-1}(T_2)^\Delta\) such that
\[
 \operatorname{rank}_O L_r\le em,\qquad
 O\bar z_1+\cdots+O\bar z_k+L_r=L_c(T_2),
\]
and \(v_{c,\zeta,T_2}(z_j)\in\mathscr F_\tau^{\beta'_\tau}(T_2)\) at every upward evaluation pair above \(c\). The module \(L_r\) is finite free. For every \(w\in\mathcal P_c^{(1)}(T_2)\), there are admissible slot coefficients \(f_1,\ldots,f_k\) and \(x_r\in\mathcal I_c^L(T_2;L_r)\) such that each \(f_j\otimes\bar z_j\) belongs to \(\mathcal I_c^L(T_2;L_c(T_2))\) and
\[
 \beta_{1,T_2}(d_cw)=\sum_{j=1}^k[f_j\otimes\bar z_j]+[x_r]
       \quad\text{in }\mathcal H_c(T_2).
\]
Here square brackets denote the quotient by \(\pi'\mathcal M_c(T_2)\).

\item[(c)] For every \(O\)-submodule \(S_0\subset L_c(T_2)\), \(\mathcal R_c(T_2;S_0)\) is finite free and
\[
 \operatorname{rank}_O\mathcal R_c(T_2;S_0)
       \le(e+1)\operatorname{rank}_O S_0.
\]
In particular, \(\operatorname{rank}_O\mathcal I_{\rm rest}\le e(e+1)m\) for the \(L_r\) of (b). Inclusion of the numerators induces an injection \(\mathcal R_c(T_2;S_0)\hookrightarrow\mathcal H_c(T_2)\).

\item[(d)] There is \(a_c\in O\) such that \(a_cd_c=\ell^{\kappa_A}\). Every elementary slot tensor for \(T_2\) belongs to \(\mathcal I_c^L(T_2;L_c(T_2))\). For any Frobenius-horizontal frame, if \(f\otimes\bar z\) is such a tensor and \(y\in\Lambda_c^P(T_2)\) satisfies \(\alpha_{T_2}(f\otimes\bar z)=q_P(y)\), then
\[
 (\gamma-1)y\in\pi'N_c(T_2)^\Delta,\qquad
 \ell^{y_1}y-F_P(y_1')\in\pi'N_c(T_2)^\Delta
\]
for some \(y_1'\in\Lambda_c^P(T_1)\) with \((\gamma-1)y_1'\in\pi'N_c(T_1)^\Delta\). This holds when all signs above \(c\) are downward, when they are all upward, and when \(e=2\). Equivalently,
\[
 \ell^{y_1}\alpha_{1,T_2}\bigl(\mathcal C_{\rm lift}(T_2,c)\bigr)
       \subset F_*\bigl(\mathcal P_c^{(1)}(T_1)\bigr).
\]

\item[(e)] There is a finite subset \(\Sigma_c\subset Z_f(T_2)\), of cardinality at most \(e(e+1)m\), such that
\[
 \ell^{\sigma_A+2\sigma_A+\kappa_A+y_1}\mathcal J_c(T_2)
       \subset I_F+O\Sigma_c.
\]
\end{enumerate}
\end{proposition}

\begin{proof}[Outline]
Backward Frobenius and graded \(\gamma-1\) lifting give (a), filtration subspaces give (b), and Lemma~\ref{B8.42bqe}(a) gives (c). Wach lifting with Lemma~\ref{B8.42bqe}(c), Lemma~\ref{B8.42bqca}, and Lemma~\ref{B8.42bqcb} proves (d), and Herr pairs yield (e). The complete proof is in \S\ref{prf:B8.42bqc}.
\end{proof}

\begin{proposition}[The annihilator statement at \(q\) and \(d\)]\label{B8.42bq1}
Assume \(\mathscr H_2\). Then statement (A) of Definition~\ref{B8.42bq} holds for every local field \(K\) of that definition, and in particular at every completion \(K=E_w\) of the host \(E\) above \(q\) or \(d\).
\end{proposition}
\begin{proof}[Outline]
Normalize data and split \(Z_f(T_2)\) into \(\mathcal J_c(T_2)\) via Lemma~\ref{B8.42bqa} (with Lemma~\ref{B8.42bd} at \(q,d\)), applying Proposition~\ref{B8.42bqc}(e) at each \(c\) for \(\Sigma_0\). Comparing exponents with Lemma~\ref{B8.42bc} via \(\eta_1\le\eta\) gives (A). The complete proof is in \S\ref{prf:B8.42bq1}.
\end{proof}

\begin{lemma}[The local lifting bound at a ramified prime]\label{B8.42br}
Assume \(\mathscr H_2\) and statement (A) of Definition~\ref{B8.42bq} at \(K\). For admissible data at \(K\), put \(B=eh/\ell+eS+e\log_\ell h+1+\eta\), with \(\log_\ell0=0\). Then
\[
 \lambda_O\bigl(Z_f(T_2)/I_F\bigr)\le160\bigl((n_Ke+1)rB+n_Ke(e+1)m(h+1)\bigr),
\]
where \(\lambda_O\) is the length normalized by the length of \(O/\ell O\).
\end{lemma}
\begin{proof}[Outline]
Bounds generators of \(Z_f(T_2)/B(T_2)\) by \(2(n_Ke+1)r\) via \ref{B8.41p}, Nakayama's lemma, and Bloch--Kato dimension. Lifts finite classes at exponent \(E'\) over \(K_0(\mu_\ell)\) via \ref{B8.41c}, \ref{B8.41a}, and corestriction, then applies (A) and \ref{B8.42bc}. The complete proof is in \S\ref{prf:B8.42br}.
\end{proof}

\begin{definition}[The paired-character statement at the type primes]\label{B8.42bs}
Let \(C_T\ge40\) be the type threshold and \(C'\ge8C_T\) the paired-character constant chosen in the proof of Theorem~\ref{B8.42}. At a member of the family, let \(i\) be an index with \(\ell=\ell_i=g_k\), \(k\in\{0,1\}\). Statement (P) at the member asserts the following for every such \(i\). Let \((R_x)_{x\in X_i}\) be full representations realizing the points of \(X_i\), and \(\rho=(\rho_x)\) rank-\(r\) representations with \(R_\pi\simeq\rho_\pi\oplus\bigoplus_c\chi_{i,c}\) for \(\pi\in Z(i)\). For an integer \(\eta_0\ge0\) and \(\pi\in Z_{\rm good}(i)\) put \(s_\pi=s_\pi(\rho)\), \(\delta_X=\sum_{0<e<t}(s_0(i)+v_\ell(e))\) and
\begin{multline*}
 R_i(\pi)=4fn_VC'n_V^2(n_V+\eta_0+t)+2fr^2+s_\pi+\sum_{0<e<t}s_{\ell,e}(E)+v_\ell(h_E)\\
 +4n_V\delta_X+t\sum_{v\in Q_i}\lambda_i(O_i/(\zeta_{i,v}-1)) .
\end{multline*} Let \(\eta_0,s,S\in\N\). Let \(\mathrm{sg}\) assign a sign to each place \(w\mid\ell\) of \(E\), with opposite signs at distinct conjugate places, and let \(U_w\subset\Z\) be finite. Suppose the following three conditions hold.
(i) \(\eta_0+t+1<s\), \(\log g_{1-k}\le2\log g_k\), and \(C_T((s+S)(n_V+1)^2+(n_V+1)^3)<g_k\).
(ii) For every \(w\mid\ell\), \(|U_w|\le n_V\), every jump at \(w\) of every \(R_x\chi_{i,1}^{-1}\) lies in \(U_w\), and every \(a\in U_w\) satisfies \(a\notin\{-1,t\}\) together with one of \(0\le a<t\), \(|a-(t-1)/2|\le\eta_0+1\), or \(0<\mathrm{sg}(w)(a-(t-1)/2)\le s+S-(t-1)/2\).
(iii) The lengths \(s_\pi\), \(\pi\in Z_{\rm good}(i)\), and \(s_{\ell,e}(E)\), \(0<e<t\), are finite.
Then
\[
 \lambda_i(U_i/L_i)\le\sum_{\pi\in Z_{\rm good}(i)}\mathrm{mult}(i,\pi)R_i(\pi),
\]
with \(U_i,L_i\) as in Definition~\ref{B8.40}.
\end{definition}

\begin{lemma}[Congruence length from a family of endomorphisms]\label{B8.42bs1a}
Fix an index \(i\), and let \(M_i\) be the integral Hecke lattice of Definition~\ref{B8.40}. Put \(K_i=\operatorname{Frac}(O_i)\), and write \(V_x\subset M_i\otimes_{O_i}K_i\) for the eigenspace on which every \(f\in T_i\subset O_i^{X_i}\) acts by \(f(x)\). For a subset \(\mathcal S\subset X_i\), put
\[
 V_{\mathcal S}=\bigoplus_{x\in\mathcal S}V_x,\quad
 U_{\mathcal S}=(M_i+V_{\mathcal S^c})\cap V_{\mathcal S},\quad
 L_{\mathcal S}=M_i\cap V_{\mathcal S}.
\]
Extend endomorphisms of \(M_i\) to its scalar extension when applying them to \(V_x\). Let \(\mathcal A\) be a family of \(O_i\)-linear endomorphisms of \(M_i\) commuting with \(T_i\), and suppose every member of \(\mathcal A\) vanishes on \(V_x\) for \(x\notin\mathcal S\). For every \(\pi\in\mathcal S\), suppose that \(a_\pi\in O_i\setminus\{0\}\) and that some \(\alpha_\pi\in\mathcal A\) acts by \(a_\pi\) on \(V_\pi\). Then
\[
 \lambda_i(U_{\mathcal S}/L_{\mathcal S})
 \le\sum_{\pi\in\mathcal S}\dim_{K_i}V_\pi\,
          \lambda_i(O_i/a_\pi O_i).
\]
\end{lemma}
\begin{proof}[Outline]
Each \(\alpha\in\mathcal A\) preserves \(V_x\) and maps \(U_{\mathcal S}\) into \(L_{\mathcal S}\). Filtering by \(F_j=\bigoplus_{a\le j}V_{\pi_a}\) gives \(a_{\pi_j}p_j(U_j)\subset p_j(L_j)\), and summing exact-sequence lengths proves the bound. The complete proof is in \S\ref{prf:B8.42bs1a}.
\end{proof}

\begin{lemma}[An endomorphism supported on the good spectrum]\label{B8.42bs1b}
Use the Hecke algebra and eigenspaces of Lemma~\ref{B8.42bs1a}. Suppose \(\ell_i=g_k\), and put
\[
 M_0=|\GL_r(\F_{g_0})|^f|\GL_r(\F_{g_1})|^f.
\]
Let \(\pi\in Z_{\rm good}(i)\), and let \(Y\subset X_i\) contain \(\pi\), with \(Y\cap Z(i)\subset Z_{\rm good}(i)\). Let
\(\beta\in\operatorname{End}_{O_i}(M_i)\) commute with the Hecke algebra, vanish on \(V_x\) for \(x\notin Y\), and act by a nonzero \(b\in O_i\) on \(V_\pi\). Suppose
\[
 \lambda_i(O_i/bO_i)\le v_{g_k}(M_0).
\]
Define
\[
 \mathfrak c_{Y,\pi}
 =\{f(\pi):f\in T_i,\ f(x)=0\text{ for all }x\in Y\setminus Z(i)\},
 \qquad v_{Y,\pi}=\lambda_i(O_i/\mathfrak c_{Y,\pi}).
\]
There are \(g\in T_i\) and \(c\in O_i\setminus\{0\}\) such that
\(\mathfrak c_{Y,\pi}=cO_i\), \(g(\pi)=c\), and \(g\) vanishes on
\(Y\setminus Z(i)\). The endomorphism \(\alpha=g\beta\) commutes with the Hecke algebra, vanishes on every \(V_x\) for \(x\notin Z_{\rm good}(i)\), and acts by \(a=cb\ne0\) on \(V_\pi\). Moreover,
\[
 \lambda_i(O_i/aO_i)\le v_{Y,\pi}+v_{g_k}(M_0).
\]
\end{lemma}
\begin{proof}[Outline]
A product of separating elements in \(T_i\) shows \(\mathfrak c_{Y,\pi}=cO_i\ne 0\), supplying \(g\in T_i\). Then \(\alpha=g\beta\) vanishes outside \(Z_{\rm good}(i)\), acts by \(cb\) on \(V_\pi\), and length additivity bounds \(\lambda_i(O_i/cbO_i)\). The complete proof is in \S\ref{prf:B8.42bs1b}.
\end{proof}

\begin{lemma}[Length after absorption of finite pushouts]\label{B8.42bs1c}
Assume the local cohomological inputs of Definition~\ref{B8.41h}. Let \(K/\Q_\ell\) and \(L/\Q_\ell\) be finite, let \(O=\mathcal O_L\), and let \(X,U\) be continuous finite-dimensional de Rham \(L\)-representations of \(G_K\). Write \(U'=U^\vee(1)\). Fix dual stable lattices \(L_U\subset U\), \(L_{U'}\subset U'\), put \(r_U=\operatorname{rank}_O L_U\), and normalize length by \(\lambda_O(O/\ell O)=1\).

For each stable lattice \(L_X\subset X\) and equivariant \(O\)-linear map \(\phi:L_X\to U/L_U\), push forward the finite integral cocycles of \(L_X\). Let \(S\subset Z^1(K,U/L_U)\) be the \(O\)-span of all these pushouts, and let \(Z_f^1(K,U/L_U)\) denote the cocycles whose classes are finite. Let \(e\ge0\) be an integer and assume the absorption condition
\(\ell^eS\subset Z_f^1(K,U/L_U)\). Then
\[
 \lambda_O\bigl((S+Z_f^1)/Z_f^1\bigr)
 \le e\bigl(\operatorname{rank}_O H_f^1(K,L_{U'})+r_U\bigr).
\]
\end{lemma}
\begin{proof}[Outline]
Local Tate duality and Definition~\ref{B8.41h} inject \((S+Z_f^1)/Z_f^1\) into \(\operatorname{Hom}_O(A/\ell^e A,L/O)\) for \(A=H_f^1(K,L_{U'})\). Bounding \(A[\varpi]\) via \(H^0(K,L_{U'}/\varpi L_{U'})\) and the structure theorem over \(O\) give the length bound. The complete proof is in \S\ref{prf:B8.42bs1c}.
\end{proof}

\begin{lemma}[The local lifting input at a type prime]\label{B8.42bs1d}
Assume \(\mathscr H_2\), let \(\ell\) be prime, and let \(C_T\ge40\) be an integer. Fix an algebraic closure \(\Omega/\Q_\ell\), a finite Galois intermediate extension \(K'/\Q_\ell\), and a finite extension \(L/\Q_\ell\) contained in \(\Omega\), with ring of integers \(O\). Let \(K'_0\) be the maximal unramified subfield of \(K'\), and let \(\iota:K'_0\to B_{\rm cris}\) be its embedding in the chosen period datum. Assume that its image is stable under \(G=G_{\Q_\ell}\). Put
\[
 H'=G_{K'},\qquad
 I=\{g\in G:g\iota(c)=\iota(c)\text{ for every }c\in K'_0\}.
\]
Choose a Frobenius element \(F\) for this period datum, and a finite set \(\Sigma\) of characters \(I\to L^\times\) stable under conjugation by \(F\). Let \(T_1,T_2\) be stable \(O\)-lattices, write \(V_j=T_j\otimes_O L\), and let \(f:T_1\twoheadrightarrow T_2\) be equivariant. Suppose \(V_1,V_2\) are crystalline over \(K'\), and that
\[
 \chi|_{I\cap H'}=1,\quad \gcd(\operatorname{ord}\chi,\ell)=1
       \quad(\chi\in\Sigma),\qquad
 D_{{\rm cris},K'}(V_j)=\bigoplus_{\chi\in\Sigma}
                  D_{{\rm cris},K'}(V_j)[\chi]
       \quad(j=1,2),
\]
where the bracket denotes the \(I\)-eigenspace. Let \(N,\eta',H\in\N\) and \(J\subset\Z\) satisfy
\[
 1\in\Sigma,\quad |\Sigma|\le N,\quad N\ge2,\quad
 0\in J\subset[-\eta',H],\quad |J|\le N+1,\quad
 1\le\eta'<H,\quad C_T(HN^2+N^3)<\ell.
\]
Suppose all Hodge--Tate jumps of both representations, at every embedding, belong to \(J\). Let \(r_{\rm orb}\) be the maximum length of a conjugation orbit of a ratio \(\chi\psi^{-1}\), for \(\chi,\psi\in\Sigma\). Then
\[
 4r_{\rm orb}\le N^2,\qquad
 \ell^{C_T(Nr_{\rm orb}+\eta'(N+r_{\rm orb}))}Z_f^1(T_2)
       \subset f_*Z_f^1(T_1)+B^1(T_2).
\]
Here \(Z_f^1\) consists of integral cocycles whose rational class is finite, and \(B^1\) consists of integral coboundaries.
\end{lemma}
\begin{proof}[Outline]
Character averaging separates inertia eigenspaces under \(\mathscr H_2\). Extending coefficients to contain the degree-\((\ell-1)!\) unramified extension, tame descent builds good crystalline lattices, applies the lifting and orbit bounds, and descends. The complete proof is in \S\ref{prf:B8.42bs1d}.
\end{proof}

\begin{lemma}[Passing the lifting bound to a family]\label{B8.42bs1e}
Assume \(\mathscr H_2\). Let \((V_x)_{x\in A}\) be a family of continuous \(L\)-representations of \(G_{\Q_\ell}\) of a common finite dimension, crystalline over one extension \(K'\), with the extension and period stability of Lemma~\ref{B8.42bs1d}. Suppose the family has a common Frobenius-stable character set \(\Sigma\) and jump set \(J\), with the inertia and jump conditions of that lemma for every member. Retain its integers \(C_T,N,\eta',H\). Put
\[
 r'=\lfloor N^2/4\rfloor,\qquad
 e_{\rm fam}=C_T\bigl(Nr'+\eta'(N+r')\bigr).
\]
Let \(W_1,W_2\) be subrepresentations of finite direct sums of members of the family. For every pair of stable lattices \(L_j\subset W_j\) and equivariant surjection \(f:L_1\twoheadrightarrow L_2\),
\[
 \ell^{e_{\rm fam}}Z_f^1(L_2)
       \subset f_*Z_f^1(L_1)+B^1(L_2).
\]
\end{lemma}
\begin{proof}[Outline]
Fontaine comparison transfers crystallinity over \(K'\), the \(\Sigma\)-eigenspace sum, and jumps in \(J\) to \(W_1,W_2\). Lemma~\ref{B8.42bs1d} then applies to \(L_1,L_2\), and \(r_{\rm orb}\le r'\) gives the exponent \(e_{\rm fam}\). The complete proof is in \S\ref{prf:B8.42bs1e}.
\end{proof}

\begin{lemma}[The local cut-cocycle bound from either family]\label{B8.42bs1f}
Assume \(\mathscr H_2\), and retain the type-locus data in Proposition~\ref{B8.42bs1}. Fix \(w\mid\ell\) with \(E_w=\Q_\ell\), put \(V_x=R_x\chi_{i,1}^{-1}\), and choose \(\pi\in X_s\cap Z(i)\). Let \(U\) be either the normalized rank-\(r\) constituent \(\rho_\pi\chi_{i,1}^{-1}\), with \(k=r\), or the character constituent \(\chi_{i,c}\chi_{i,1}^{-1}\), with \(k=1\) and \(1\le c\le t\). The semisimple decomposition of \(V_\pi\) embeds \(U\) into it. Fix a stable lattice \(L_U\subset U\). Suppose the family \((V_x)_{x\in X_s}\) is crystalline over a common extension and has a common finite character set satisfying the extension, period-stability and inertia conditions of Lemma~\ref{B8.42bs1d}. Suppose that every lattice surjection in subrepresentations of finite sums of the family, or every such surjection for its Tate-dual family, lifts finite cohomology after multiplication by \(\ell^e\). Let \(S_{w,U}\) be the span of the cut cocycles in \(U/L_U\), with all spectral points in \(X_s\). Then
\[
 \lambda_i\bigl((S_{w,U}+Z_f^1)/Z_f^1\bigr)\le2ke.
\]
\end{lemma}
\begin{proof}[Outline]
Fontaine comparison puts cut cocycles into finite pushouts from a local span \(X\). The direct or dual lifting hypothesis and the Bloch--Kato dimension formula give the bound \(2ke\) on finite collections, extending to the full span via direct sums. The complete proof is in \S\ref{prf:B8.42bs1f}.
\end{proof}

\begin{proposition}[The paired-character estimate at the type primes]\label{B8.42bs1}
Assume \(\mathscr H_2\). For every integer \(C_T\ge40\) and every real \(C'\ge8C_T\), statement (P) of Definition~\ref{B8.42bs} with type threshold \(C_T\) and paired-character constant \(C'\) holds at every member satisfying the norm-choice conditions and four local compatibility conditions of Theorem~\ref{B8.42}. Because any fixed integer \(C_T\ge40\) (such as \(C_T=40\)) and \(C'\ge8C_T\) depend on \(\mathscr H_2\) alone, they may be chosen before the cuts in Theorem~\ref{B8.42}, and statement (P) then holds at every member above \(x_1\) of the twist-good cut family of Theorem~\ref{B8.42}.
\end{proposition}
\begin{proof}[Outline]
Lemmas~\ref{B8.42bs1c}, \ref{B8.42bs1d}, \ref{B8.41a}, \ref{B8.42bs1e}, and \ref{B8.42bs1f} bound cut cocycles on type loci. Global reciprocity then bounds \(v_{cc}\), and Lemmas~\ref{B8.42bs1a} and \ref{B8.42bs1b} bound \(\lambda_i(U_i/L_i)\). The complete proof is in \S\ref{prf:B8.42bs1}.
\end{proof}

\begin{lemma}[The local error at the type primes]\label{B8.42bt}
Consider a member in the small-prime setup preceding Theorem~\ref{B8.41}, with hub width \(\eta_X\) and type threshold \(40\le C_T\le b\), at which statement (P) of Definition~\ref{B8.42bs} holds. Let \(i\) be a twisted index with \(\ell_i=g_k\). Then \(\lambda_i(U_i/L_i)\le\sum_{\pi\in Z_{\rm good}(i)}\mathrm{mult}(i,\pi)R_i(\pi)\) with \(\eta_0=\eta_X\).
\end{lemma}
\begin{proof}[Outline]
If \(Z_{\rm good}(i)\) is empty, \(U_i=L_i=0\) by \ref{B8.40}. Otherwise the parameter split bound, the type-prime inequality \(g_{1-k}\le2g_k\), and the finiteness hypothesis of \ref{B8.40} verify conditions (i)--(iii) of \ref{B8.42bs}, and statement (P) gives the bound. The complete proof is in \S\ref{prf:B8.42bt}.
\end{proof}

\begin{definition}[The coefficient-chain statement for words with every letter thick]\label{B8.42bu}
Let \(\ell\ge11\) be prime and \(q_\ell=\ell-1\). Let \(K/\Q_\ell\) be unramified of even degree \(d'\), and let \(L'\) be a finite extension of \(\Q_\ell\) containing a primitive \(\ell\)-th root of unity and the images of all embeddings of \(K(\mu_\ell)\), with integers \(O\). Number the embeddings \(\tau_i\), \(i\in\Z/d'\Z\), in Frobenius order. Let \(F:T_1\to T_2\) be an equivariant surjection of crystalline \(O\)-lattices, \(r=\operatorname{rank}T_2\), and use \(Z_f\), \(B\) and \(I_F\) as in Definition~\ref{B8.42bq}. Let \(p_0\in\{0,1\}\) and \(S,\delta,m,\beta_i\ge0\) be integers with the following properties. The jumps of \(T_1\) at \(\tau_i\) lie in \([-\beta_i-S,-\beta_i]\cup[-\delta,\delta]\) when \(i+p_0\) is even, for the representative of \(i\) in \([0,d')\), and in \([\beta_i,\beta_i+S]\cup[-\delta,\delta]\) otherwise. At each \(\tau_i\), the lattice \(T_2\) has at most \(m\) jumps in \([-\delta,\delta]\). With
\[
 h=\max_i\beta_i+S,\quad H=h+\delta,\quad \lambda=1+\log_\ell(2H),\quad A=\frac{H(2\ell-1)}{q_\ell^2}+3+\lambda,
\]
with \(\log_\ell0=0\),
one has \(h/2\le\beta_i\) and \(2A+\delta+1\le\beta_i\) for every \(i\). Put
\[
 \sigma=\frac{\lfloor h/q_\ell\rfloor+1+\delta}{q_\ell}+1,\qquad \sigma_2=\frac{2H}{q_\ell^2}+1 .
\]
Statement (R) asserts the following for all such data. There are an integer \(a_1\le A+1\), real numbers \(s_1,E_1,e_2\) and an integer \(e\ge0\) with
\[
 s_1\le\frac{2a_1+2\delta+S}{q_\ell}+2\lambda,\quad E_1\le\frac{3a_1+S+2\delta}{q_\ell}+3\lambda,\quad
 e_2\le\frac h{\ell q_\ell}+E_1+2\sigma_2+2S+5+\lambda,
\]
and \(e\le s_1+2\sigma+e_2+S+E_1+\sigma_2+1\). There are also a finite free \(O\)-module \(F_0\) of rank at most \((d'+1)r\), an \(O\)-linear map \(\jmath:F_0\to Z_f(T_2)\) surjective modulo \(B(T_2)\), a chain \(F_0\supset F_1\supset\cdots\supset F_k\), and a set \(\mathcal G\) of at most \(4d'\) indices below \(k\) with \(k-|\mathcal G|\le3d'/2\). For \(i\in\mathcal G\), \(F_i\) lies in the sum of \(F_{i+1}\) and the span of at most \(m\) of its elements. For \(i<k\) outside \(\mathcal G\), the length of \(F_i/(F_i\cap(F_{i+1}+\jmath^{-1}(I_F)))\), normalized by that of \(O/\ell O\), is at most \(rS+2m((H+1)/q_\ell+1)\). Finally \(\ell^e\jmath(F_k)\subset I_F\).
\end{definition}

\begin{lemma}[Coordinate-subset basis refinement for a saturated submodule]\label{B8.42bua}
Let \(O\) be a discrete valuation ring, \(M\) a finite free \(O\)-module with finite basis \((b_i)_{i\in I}\), and \(S\subset I\) a subset. Let \(P\subset\sum_{i\in S}Ob_i\) be an \(O\)-saturated submodule of \(M\), meaning that whenever \(a\in O\setminus\{0\}\), \(x\in M\), and \(ax\in P\), one has \(x\in P\). Then there exists an \(O\)-basis \((b'_i)_{i\in I}\) of \(M\) such that \(b'_i=b_i\) for every \(i\notin S\), \(\sum_{i\in S}Ob'_i=\sum_{i\in S}Ob_i\), and \(P=\sum_{i\in I,\,b'_i\in P}Ob'_i\).
\end{lemma}
\begin{proof}[Outline]
Restrict to \(Q=\sum_{i\in S}Ob_i\), where \(P\) is \(O\)-saturated; over the DVR \(O\), \(P\) is a direct summand of \(Q\) with adapted basis \((c_j)_{j\in S}\). Replacing \((b_i)_{i\in S}\) by \((c_i)_{i\in S}\) while fixing \(b_i\) for \(i\notin S\) gives the basis of \(M\). The complete proof is in \S\ref{prf:B8.42bua}.
\end{proof}

\begin{definition}[Components and coefficient chains for thick words]\label{B8.42buc}
Use the data of Definition~\ref{B8.42bu}, and put \(n=d'\), \(q=\ell-1\), and
\[
 J_-=\{c\in\Z/n\Z:c+p_0\text{ is even}
       \text{ for the representative of }c\text{ in }[0,n)\}.
\]
Thus \(|J_-|=n/2\). Put \(\Delta=\Gal(K(\mu_\ell)/K)\), and choose a topological generator \(\gamma\) of \(\Gal(K(\mu_{\ell^\infty})/K(\mu_\ell))\). Write \(\mathcal A=\prod_{i\in\Z/n\Z}O[[Z]]\), and let \(e_i\) be its \(i\)-th idempotent. For a crystalline \(O\)-lattice \(T\), use the Wach lattice \(N(T)\subset\mathcal D(T)\) and its associated étale module \(\mathcal D(T)\) as in Lemmas~\ref{B8.41ca}--\ref{B8.41c} (at \(s=1\)), and let \(\operatorname{cmp}_T:N(T)\to D_{\rm cris}(T[1/\ell])=(B_{\rm cris}\otimes_{\Q_\ell}T[1/\ell])^{G_K}\) be the specialization map modulo \(Z\) in the proof of Lemma~\ref{B8.41ca} composed with scalar extension to \(L'\). Write \(\varphi\) for Frobenius on these modules, and \(F_D\) for the map induced by \(F\). Choose the invariant parameter \(\pi_0\) at unramified level one, and put \(Q_0=\pi_0+\ell\). The parameter is normalized by \(\mathcal A^\Delta=\prod_iO[[\pi_0]]\); the cyclotomic polynomial \(\Phi_\ell(1+Z)\) is a unit times \(Q_0\).

For \(k\ge0\), put \(P_k=\prod_{j=0}^{k-1}\varphi^j(Q_0)\), with \(P_0=1\). Thus \(P_k^{a_0}x\in N(T)\) for some \(k,a_0\ge0\) means that \(x\) has finite poles along these Frobenius translates. Define
\[
\begin{aligned}
 \mathcal D_i^\Delta(T)&=\{x\in\mathcal D(T):e_ix=x,
                  \ \varepsilon x=x\text{ for every }\varepsilon\in\Delta\},\\
 \mathcal N_i(T)&=N(T)\cap\mathcal D_i^\Delta(T),\qquad
 \mathcal U^j(T)=\pi_0^jN(T)^\Delta\quad(j\ge0),\\
 \mathcal R_i^{(j)}(T)&=\mathcal A\text{-span}\{
                 \varphi^jx:x\in\mathcal N_i(T)\},\\
 \mathcal V_i(T)&=\{x\in\mathcal D_i^\Delta(T):
       (\gamma-1)x\in\mathcal U^1(T),
       \text{ and }P_k^{a_0}x\in N(T)\text{ for some }k,a_0\ge0\},\\
 \mathcal W_i(T)&=\{x\in\mathcal V_i(T):
                         Q_0^\delta x\in\mathcal R_{i-1}^{(1)}(T)\},\\
 \mathcal C_i(T)&=\operatorname{cmp}_T(\mathcal N_i(T)),\\
 \mathcal C_i(F)&=\{\operatorname{cmp}_{T_2}(F_Dx):
       x\in\mathcal N_i(T_1),\ F_Dx\in N(T_2)\}.
\end{aligned}
\]
A Herr cocycle of \(T\) is a pair \((z,v)\in\mathcal D(T)^\Delta\times\mathcal D(T)^\Delta\) satisfying \((\gamma-1)z=(\varphi-1)v\). Write \(\mathcal H(T)\) for their \(O\)-module and put \(\operatorname{cl}_T=\operatorname{cl}_{K,T}:\mathcal H(T)\to H^1(K,T)\) for the natural integral Herr class map of Lemma~\ref{B8.42bqa}(c) (with \(e=1\)). Its kernel consists of the pairs \(((\varphi-1)t,(\gamma-1)t)\) with \(t\in\mathcal D(T)^\Delta\). The map \(F_{\mathcal H}:\mathcal H(T_1)\to\mathcal H(T_2)\) applies \(F_D\) to both coordinates.

Set
\[
 a=\lceil A\rceil,\qquad M=\lfloor h/q\rfloor+1,\qquad
 \widehat\sigma=\lceil\sigma\rceil,\qquad
 C_{\rm len}=rS+2m\left(\frac{H+1}{q}+1\right),
\]
and define the real numbers
\[
\begin{aligned}
 B_1&=\frac{2a+2\delta+S}{q}+\lambda,\qquad
 E_G=\frac{3a+S+2\delta}{q}+\lambda,\\
 B_{\rm red}&=B_1+2\sigma+2,\\
 B_{\rm vec}&=\frac{h}{\ell q}+\frac{S+1}{q}+S+E_G
                           +2\sigma_2+6+\lambda,\\
 B_{\rm int}&=S+E_G+\sigma_2+\lambda.
\end{aligned}
\]
Here \(A,h,H,\lambda,\sigma,\sigma_2\) have the meanings fixed in Definition~\ref{B8.42bu}.

For submodules \(U'\subset U\) of a finite free \(O\)-module, a \emph{generator step} means that \(U/U'\) is generated by at most \(m\) elements. A \emph{length step} means that
\(\lambda_O(U/U')\le C_{\rm len}\), where
\(\lambda_O=\operatorname{length}_O/\operatorname{length}_O(O/\ell O)\).
Equivalently, in the \emph{obstruction-map formulation}, a step \(U'\subset U\) is a generator step (respectively, a length step) if there is an \(O\)-linear map \(U\to Q\) with kernel \(U'\) such that \(Q\) is generated over \(O\) by at most \(m\) elements (respectively, \(\lambda_O(Q)\le C_{\rm len}\)). A chain with at most \(a'\) generator steps and \(b'\) length steps is a finite descending chain whose steps can be assigned these two types with the indicated counts.
\end{definition}
\begin{proof}
The obstruction-map formulation of a step gives an injection \(U/U'\hookrightarrow Q\). If \(Q\) is generated by at most \(m\) elements, choose a surjection \(O^m\to Q\). The inverse image of \(U/U'\) is a submodule of \(O^m\), hence is free of rank at most \(m\), and maps onto \(U/U'\). If \(\lambda_O(Q)\le C_{\rm len}\), the injection gives the same bound for \(U/U'\). Conversely, take \(Q=U/U'\) and the quotient map. Thus the two formulations of coefficient chains agree.
\end{proof}

\begin{lemma}[Cocycle reduction on the negative components]\label{B8.42bud}
Assume \(\mathscr H_2\) and the data of Definitions~\ref{B8.42bu} and~\ref{B8.42buc}. There is an integer \(u_1\ge0\) with \(u_1\le B_1\) such that
\[
 s_{\rm red}=u_1+2\widehat\sigma\le B_{\rm red}
\]
has the following property. For every finite free \(O\)-module \(P\) and \(O\)-linear map \(g:P\to Z_f(T_2)\), there are a submodule \(P_2\subset P\), a chain from \(P\) to \(P_2\) with at most \(n\) generator steps and no length steps, and \(O\)-linear maps
\[
 w_c:P_2\longrightarrow\mathcal D(T_2)\qquad(c\in J_-)
\]
with \(w_c(P_2)\subset\mathcal W_c(T_2)\). Write \(Z^1(K,T_2)\) for the \(O\)-module of continuous cocycles \(G_K\to T_2\). For each \(f\in P_2\) and \(c\in J_-\), one can choose an integral cocycle \(\kappa_c(f)\in Z^1(K,T_2)\) and a Herr cocycle \(X_c(f)=(z_c(f),v_c(f))\) such that
\[
 v_c(f)\in\mathcal U^M(T_2),\qquad
 \operatorname{cl}_{T_2}(X_c(f))=[\kappa_c(f)],
\]
for a sign \(\epsilon(f)\in\{1,-1\}\),
\[
 \ell^{s_{\rm red}}g(f)-\epsilon(f)\sum_{c\in J_-}\kappa_c(f)\in B(T_2),
\]
and whenever \(e_*\ge0\) is an integer and \(u\in\mathcal V_c(T_1)\) satisfies \(\ell^{e_*}w_c(f)-F_Du\in\mathcal U^1(T_2)\), there exists \(Y\in\mathcal H(T_1)\) with \(\operatorname{cl}_{T_1}(Y)\in H_f^1(K,T_1)\) and \(\ell^{e_*}X_c(f)-F_{\mathcal H}Y\in\mathcal U^M(T_2)\times\mathcal U^M(T_2)\).
\end{lemma}
\begin{proof}[Outline]
From extension Herr data, a \(D\)-polynomial correction and uniformizer truncation with \(n\) kernel steps give \(P_2\) and \(w_c\). Bijectivity of \(\varphi-1\) on \(\mathcal U^M(T_2)\) and truncation then construct \(X_c\), \(\kappa_c\), and the lift \(Y\). The complete proof is in \S\ref{prf:B8.42bud}.
\end{proof}

\begin{lemma}[A recursive letter and its interpolation cost]\label{B8.42buea}
Assume \(\mathscr H_2\) and the data of Definitions~\ref{B8.42bu} and~\ref{B8.42buc}. For \(i\notin J_-\), put
\[
 E_i=\left\lfloor\frac{3a+S+2\delta-1}{q}+\log_\ell H\right\rfloor,
 \qquad s_2=\lceil\sigma_2\rceil.
\]
Then \(E_i\ge0\) and
\[
 E_i\le\frac{3A+S+2\delta}{q}+\lambda\le E_G.
\]
For every \(y\in\mathcal N_{i-2}(T_2)\) such that
\(d=\varphi^2\operatorname{cmp}_{T_2}(y)\in\operatorname{Fil}^{\beta_i}_{\rm cris}\), there is \(z\in\mathcal N_i(T_1)\) with
\[
 F_Dz\in N(T_2),\qquad
 \operatorname{cmp}_{T_2}(F_Dz)
       =\ell^{E_i+s_2+\beta_{i-1}+S-\beta_i}d.
\]
The exponent on the right is an integer and may be negative.
\end{lemma}
\begin{proof}[Outline]
Two-Frobenius truncation and Laurent expansions build \(x\in\mathcal N_i(T_2)\) with comparison value \(\ell^{s_2+\beta_{i-1}+S-\beta_i}d\). Lifting \(P_I(D_2)x\) and \(P_J(D_2)x\) and combining them via polynomial interpolation at exponent \(E_i\) gives \(z\). The complete proof is in \S\ref{prf:B8.42buea}.
\end{proof}

\begin{lemma}[First Frobenius correction of the chosen coordinates]\label{B8.42bue}
Assume \(\mathscr H_2\) and the data of Definitions~\ref{B8.42bu} and~\ref{B8.42buc}. There is an integer \(e_{\rm vec}\ge0\) with \(e_{\rm vec}\le B_{\rm vec}\) such that, for every \(c\in J_-\), every finite free \(O\)-module \(P\), every submodule \(U\subset P\), and every \(O\)-linear map \(w:U\to\mathcal D(T_2)\) with \(w(U)\subset\mathcal W_c(T_2)\), there are a submodule \(U'\subset U\), an \(O\)-linear map \(y_c:U'\to\mathcal D(T_1)\), and a chain from \(U\) to \(U'\) with at most three generator steps and one length step, such that
\[
 y_c(f)\in\mathcal V_c(T_1),\qquad
 r_c(f):=\ell^{e_{\rm vec}}w(f)-F_Dy_c(f)\in\mathcal N_c(T_2)
 \quad(f\in U').
\]
\end{lemma}
\begin{proof}[Outline]
Constant-coefficient reduction links \(D_f(0)\) to \(x_f\), and Lemma~\ref{B8.42buea} with truncation builds source lifts \(v_d\). A determinant quotient pullback and two filtration kernels yield \(U'\) and \(y_c\) with residual in \(\mathcal N_c(T_2)\). The complete proof is in \S\ref{prf:B8.42bue}.
\end{proof}

\begin{lemma}[Restricting a coefficient module through a quotient]\label{B8.42bufa}
Let \(\ell\) be a prime, let \(O\) be the ring of integers of a finite extension of \(\Q_\ell\), and normalize length by
\(\lambda_O(Q)=\operatorname{length}_O(Q)/\operatorname{length}_O(O/\ell O)\).
Let \(a,m\ge0\) be integers with \(a\le m\), let \(C\ge0\) be real, let
\(K\subset X\subset N\) be \(O\)-modules, and let \(s_1,\ldots,s_a\in N\). Suppose
\[
 \lambda_O\bigl(N/(X+\textstyle\sum_{j=1}^aOs_j)\bigr)\le C.
\]
Let \(U\) be an \(O\)-module. For an \(O\)-linear map
\(\psi:U\to N/K\), put
\[
 U_1=\psi^{-1}\bigl((X+\textstyle\sum_jOs_j)/K\bigr),
 \qquad U_2=\psi^{-1}(X/K).
\]
Then \(U\supset U_1\supset U_2\) has one length step of normalized length at most \(C\) and one generator step with at most \(m\) generators.
\end{lemma}
\begin{proof}
The map induced by \(\psi\) embeds \(U/U_1\) in
\(N/(X+\sum_jOs_j)\), proving the length estimate. It embeds \(U_1/U_2\) in
\((X+\sum_jOs_j)/X\), which is generated by the images of the \(s_j\).
A submodule of a module generated by \(a\) elements over a discrete valuation ring is generated by at most \(a\) elements: lift it to a submodule of \(O^a\), take a basis of that submodule, and project. This proves the generator estimate.
\end{proof}

\begin{lemma}[Explicit reconstruction after three crystalline comparisons]\label{B8.42bufb}
Use the modules of Definition~\ref{B8.42buc} and their level-one Wach properties. Explicitly, specialization kills precisely
\(\pi_0\mathcal N_i(T_j)\) in each \(\mathcal N_i(T_j)\), commutes with maps of lattices, and satisfies
\[
 \operatorname{cmp}_{T_j}(Q_0^k\varphi z)
       =\ell^k\varphi\operatorname{cmp}_{T_j}(z)
\]
whenever \(z\in\mathcal N_{i-1}(T_j)\), \(k\ge0\), and
\(Q_0^k\varphi z\in\mathcal N_i(T_j)\).
Also \(Q_0(0)=\ell\) and
\((\gamma-1)\mathcal N_i(T_j)\subset\pi_0\mathcal N_i(T_j)\).

Let \(c\) be a component and let \(B\ge0\) be an integer. Suppose
\(Q_0^B\varphi\mathcal N_{c-1}(T_j)\subset\mathcal N_c(T_j)\) for \(j=1,2\).
Let \(r\in\mathcal N_c(T_2)\), let
\(d_1,d_2,d_3\in D_{\rm cris}(T_2[1/\ell])\), and let
\(s_1,s_2,b_0,h_0,e,t\) be nonnegative integers and \(k\in\Z\) satisfying
\[
\begin{aligned}
 \ell^{s_1}\operatorname{cmp}_{T_2}(r)&=\ell^B\varphi d_1,\\
 \ell^{b_0}d_1&=\varphi d_2,\\
 \ell^{s_2}d_2&=\ell^{h_0}\varphi d_3,\\
 e+h_0&=t+k+b_0.
\end{aligned}
\]
Suppose \(z\in\mathcal N_{c-1}(T_1)\) satisfies
\[
 F_Dz\in N(T_2),\qquad
 \operatorname{cmp}_{T_2}(F_Dz)=\ell^k\varphi^2d_3.
\]
Then
\[
 x=\ell^tQ_0^B\varphi z\in\mathcal V_c(T_1),\qquad
 \ell^{e+s_1+s_2}r-F_Dx\in\mathcal U^1(T_2).
\]
\end{lemma}
\begin{proof}[Outline]
Height bounds give \(x\in\mathcal V_c(T_1)\) and \(F_Dx\in\mathcal N_c(T_2)\). Weighted specialization and the identities equate specializations of \(F_Dx\) and \(\ell^{e+s_1+s_2}r\), putting the difference in \(\pi_0\mathcal N_c(T_2)\subset\mathcal U^1(T_2)\). The complete proof is in \S\ref{prf:B8.42bufb}.
\end{proof}

\begin{lemma}[Integral residual correction and compatible Herr classes]\label{B8.42buf}
Assume \(\mathscr H_2\) and the data of Definitions~\ref{B8.42bu} and~\ref{B8.42buc}. Put
\[
 e_{\rm int}=S+\lceil E_G+\sigma_2\rceil.
\]
Then \(e_{\rm int}\le B_{\rm int}\). For every \(c\in J_-\), every finite free \(O\)-module \(P\), every submodule \(U\subset P\), and every \(O\)-linear map \(r:U\to\mathcal D(T_2)\) with \(r(U)\subset\mathcal N_c(T_2)\), there are a submodule \(U''\subset U\) and a chain from \(U\) to \(U''\) with at most three generator steps and two length steps such that, for each \(f\in U''\), there is an element \(x_f\in\mathcal V_c(T_1)\) satisfying
\[
 \ell^{e_{\rm int}}r(f)-F_Dx_f\in\mathcal U^1(T_2).
\]

When \(U=U'\) and \(r=r_c\) are obtained in Lemma~\ref{B8.42bue} from \(w=w_c|_V\) on any submodule \(V\subset P_2\) of Lemma~\ref{B8.42bud}, set \(u_f=\ell^{e_{\rm int}}y_c(f)+x_f\) for \(f\in U''\). Then
\[
 \ell^{e_{\rm vec}+e_{\rm int}}w_c(f)-F_Du_f\in\mathcal U^1(T_2),
\]
and there is a Herr cocycle \(Y_c(f)\in\mathcal H(T_1)\) such that
\[
 \operatorname{cl}_{T_1}(Y_c(f))\in H_f^1(K,T_1),\qquad
 \ell^{e_{\rm vec}+e_{\rm int}}X_c(f)-F_{\mathcal H}Y_c(f)
                  \in\mathcal U^M(T_2)\times\mathcal U^M(T_2).
\]
Consequently,
\[
 \ell^{e_{\rm vec}+e_{\rm int}}\kappa_c(f)\in I_F.
\]
\end{lemma}
\begin{proof}[Outline]
Using Lemma~\ref{B8.42bufa} twice and a filtration kernel gives \(U''\), where Lemmas~\ref{B8.42buea} and \ref{B8.42bufb} build \(x_f\). Lemma~\ref{B8.42bud} lifts to \(Y_c(f)\), and bijectivity of \(\varphi-1\) makes the error a coboundary. The complete proof is in \S\ref{prf:B8.42buf}.
\end{proof}

\begin{lemma}[A basis adapted to a saturated submodule]\label{B8.42bu1a}
Let \(\ell\) be prime, let \(L'/\Q_\ell\) be finite, let \(O\) be its ring of integers, let \(M\) be a finite free \(O\)-module, and let \((b_i)_{i\in I}\) be an \(O\)-basis with \(I\) finite. A submodule \(P\subset M\) is called saturated if
\[
 a\in O\setminus\{0\},\quad x\in M,\quad ax\in P
       \quad\Longrightarrow\quad x\in P.
\]
If \(P\) is saturated, there are an \(O\)-basis \((c_i)_{i\in I}\) of \(M\) and a subset \(J\subset I\) such that
\[
 P=\bigoplus_{i\in J}Oc_i.
\]
The change of basis and its inverse have entries in \(O\). Thus both matrices have entries of nonnegative valuation, and their determinants have valuation zero, for the valuation normalized by \(v_\ell(\ell)=1\).
\end{lemma}
\begin{proof}[Outline]
Smith normal form gives a basis \((c_i)\) of \(M\) and \(p_j=a_jc_{\alpha(j)}\) spanning \(P\) with \(a_j\ne 0\). Saturation gives \(c_{\alpha(j)}\in P\), so \(J=\alpha(J_0)\) spans \(P\), and the \(O\)-transition matrices are inverses with unit determinants. The complete proof is in \S\ref{prf:B8.42bu1a}.
\end{proof}

\begin{lemma}[Refining a basis inside a prescribed coordinate submodule]\label{B8.42bu1b}
Let \(O,M,I,(b_i)\) be as in Lemma~\ref{B8.42bu1a}, let \(S\subset I\), and set \(Q=\bigoplus_{i\in S}Ob_i\). If \(P\subset Q\) is saturated in \(M\), there is an \(O\)-basis \((b'_i)_{i\in I}\) of \(M\) such that
\[
 b'_i=b_i\quad(i\notin S),\qquad
 \bigoplus_{i\in S}Ob'_i=Q,\qquad
 P=\bigoplus_{\{i:b'_i\in P\}}Ob'_i.
\]
The change of basis acts as the identity on
\(\bigoplus_{i\notin S}Ob_i\). Its matrix and inverse have integral entries and unit determinant.
\end{lemma}
\begin{proof}[Outline]
Saturation in \(M\) makes \(P\) saturated in \(Q\), so Lemma~\ref{B8.42bu1a} gives an adapted basis \((c_i)_{i\in S}\) of \(Q\). Keeping \(b_i\) for \(i\notin S\) yields a basis of \(M\) with an integral, unit-determinant block-diagonal transition matrix. The complete proof is in \S\ref{prf:B8.42bu1b}.
\end{proof}

\begin{lemma}[A Wach basis adapted to the central jumps]\label{B8.42bu1c}
Assume \(\mathscr H_2\) and the data of Definitions~\ref{B8.42bu} and~\ref{B8.42buc}. Let \(i\in J_-\), put \(h_i=\beta_i+S\), and let \(e_1,\ldots,e_r\) be an \(O[[\pi_0]]\)-basis of \(\mathcal N_{i-1}(T_2)\). There are a matrix \(B\in\GL_r(O)\), an integer \(m'\le m\), and integers \(-\delta\le u_j\le\delta\), \(1\le j\le m'\), such that
\[
 e'_j=\sum_{k=1}^r B_{jk}e_k
\]
is an \(O[[\pi_0]]\)-basis and
\[
 \iota\varphi\operatorname{cmp}_{T_2}(e'_j)
       \in\operatorname{Fil}^{u_j}D_{\rm dR}\quad(1\le j\le m'),
 \qquad
 \sum_{u=-\delta}^{\delta}\operatorname{mult}_{\tau_i}(u)(h_i+u)
       \le\sum_{j=1}^{m'}(h_i+u_j).
\]
The basis change and its inverse have integral entries, so their determinant valuations are zero.
\end{lemma}
\begin{proof}[Outline]
Pulling back \(\operatorname{Fil}^u D_{\rm dR}\) via the basis \(w_k=\iota\varphi\operatorname{cmp}_{T_2}(e_k)\) gives saturated submodules \(F^u\subset O^r\). Iterating Lemma~\ref{B8.42bu1b} produces \(B\in\GL_r(O)\) and \(u_j\) matching the jump multiplicities. The complete proof is in \S\ref{prf:B8.42bu1c}.
\end{proof}

\begin{proposition}[The coefficient-chain statement for thick words]\label{B8.42bu1}
Assume \(\mathscr H_2\). Then statement (R) of Definition~\ref{B8.42bu} holds for all data of that definition.
\end{proposition}
\begin{proof}[Outline]
Assemble Wach and Herr packages from \(\mathscr H_2\) with Lemma~\ref{B8.42bu1c} and Lemma~\ref{B8.42buea}. Presenting \(H_f^1(K,T_2)\) and applying Lemma~\ref{B8.42bud}, Lemma~\ref{B8.42bue}, and Lemma~\ref{B8.42buf} over \(J_-\) proves (R). The complete proof is in \S\ref{prf:B8.42bu1}.
\end{proof}

\begin{lemma}[Separated-cluster lifting and colength bound at unramified primes]\label{B8.42bv}
Assume \(\mathscr H_2\). Let \(\ell\ge11\) be prime and \(q_\ell=\ell-1\). Let \(K/\Q_\ell\) be finite unramified of degree \(d_u\ge1\), and let \(L'\) be a finite extension of \(\Q_\ell\) containing the images of all embeddings of \(K(\mu_\ell)\), with integers \(O\). Number the embeddings \(\tau_a:K\to L'\), \(a\in\Z/d_u\Z\), in Frobenius order. Let \(F:T_1\to T_2\) be an equivariant surjection of crystalline \(O\)-lattices over \(K\) with \(r=\operatorname{rank}_O T_2\), use \(Z_f\), \(B\), and \(I_F\) as in Definition~\ref{B8.42bq}, and write \(\lambda_O\) for length normalized by \(\operatorname{length}_O(O/\ell O)\). Let \(b_*,w_*,S,\delta,J\ge0\) be integers with \(1\le J<d_u\), and for each \(a\in\Z/d_u\Z\) let \(\varsigma_a\in\{1,-1\}\) and \(s_a\in\N\) be such that the jumps of \(T_1\) at \(\tau_a\) lie in \([s_a,s_a+S]\cup[-\delta,\delta]\) when \(\varsigma_a=1\) and in \([-s_a-S,-s_a]\cup[-\delta,\delta]\) when \(\varsigma_a=-1\). Call \(a\in\Z/d_u\Z\) \emph{thick} if \(b_*\le s_a\le b_*+\delta\) and \emph{nonthick} if \(s_a\le w_*\). Assume every \(a\in\Z/d_u\Z\) is thick or nonthick, and that the thick positions form singletons or adjacent pairs separated from the next cluster along \(\Z/d_u\Z\) by at least \(J\) consecutive nonthick positions. Put \(g=w_*+S+\delta\), \(g_*=\max(1,g)\), and \(h=b_*+\delta+S\) if a thick position exists, \(h=g\) otherwise, with \(\log_\ell0=0\).

There is a threshold \(C_0\ge1\), depending on \(\mathscr H_2\) alone, such that whenever \(C_2\ge C_0\) and \(C_2(J+1)(w_*+S+\delta+1)\le b_*\), there exists an integer \(e_{\rm cl}\ge0\) with \(e_{\rm cl}\le800\{h_1/q_\ell^{J+1}+g_*(3J+1)+\log_\ell h_1\}\), where \(h_1=\max(h,1)\), satisfying \(\ell^{e_{\rm cl}}Z_f(T_2)\subset I_F\) and
\begin{equation}\label{eq:B8-cluster-col}
 \lambda_O\bigl(Z_f(T_2)/I_F\bigr)\le C_2(d_u+1)r\bigl\{h/q_\ell^{J+1}+(g+1)(J+1)+\log_\ell h\bigr\}.
\end{equation}
\end{lemma}
\begin{proof}[Outline]
Cuts after thick clusters with Lemma~\ref{B8.42ba} bound remainders, and shift killers with Lemma~\ref{B8.42bb} give \(\ell^{e_{\rm cl}}Z_f(T_2)\subset I_F\). Bounding \(Z_f(T_2)/I_F\) by \((d_u+1)r\) generators then yields the colength estimate. The complete proof is in \S\ref{prf:B8.42bv}.
\end{proof}

\begin{lemma}[Unrestricted and short-interval coordinate lifting at unramified primes]\label{B8.42bw}
Assume \(\mathscr H_2\). Let \(\ell\ge11\), \(q_\ell=\ell-1\), \(K/\Q_\ell\) unramified of degree \(d_u\ge1\), \(L'\), \(O\), \(\tau_a\) (\(a\in\Z/d_u\Z\)), \(F:T_1\to T_2\), \(r=\operatorname{rank}_O T_2\), \(Z_f\), \(B\), \(I_F\), and \(\lambda_O\) be as in Lemma~\ref{B8.42bv}.
\begin{enumerate}
\item[(a)] Suppose the jumps of \(T_1\) at \(\tau_a\) lie in \([-M_a,h_a]\) with \(M_a,h_a\in\N\), and set \(H'=\max_{a\in\Z/d_u\Z}(M_a+h_a)\). Then there is an integer \(e'\ge0\) satisfying \(\ell^{e'}Z_f(T_2)\subset I_F\) and
\begin{equation}\label{eq:B8-unrestricted-exp}
 e'\le23\{2+H'/\ell+\log_\ell(H'+1)\}.
\end{equation}
\item[(b)] Let \(b_*,w_*\ge1\) and \(S,\delta,m\ge0\) be integers with \(20w_*,20\delta,20S\le b_*\). For each \(a\in\Z/d_u\Z\), let \(\varsigma_a\in\{1,-1\}\) and \(s_a>\delta\) be integers with \(\varsigma_{a+1}=-\varsigma_a\) such that the jumps of \(T_1\) at \(\tau_a\) lie in \([s_a,s_a+S]\cup[-\delta,\delta]\) when \(\varsigma_a=1\) and in \([-s_a-S,-s_a]\cup[-\delta,\delta]\) when \(\varsigma_a=-1\), and \(T_2\) has at most \(m\) jumps in \([-\delta,\delta]\) counted with multiplicity. Assume every \(a\in\Z/d_u\Z\) is either thick (\(b_*\le s_a\le b_*+\delta\)) or nonthick (\(s_a\le w_*\)), at least one \(a\in\Z/d_u\Z\) is nonthick, and no three consecutive positions in \(\Z/d_u\Z\) are thick. Put \(h=b_*+\delta+S\) if a thick position exists and \(h=w_*+S\) otherwise, set \(U_0=h/\ell^2+(w_*+S+\delta)/\ell+\log_\ell(2h)+1\), and let \(I=1\) when some nonthick position \(a\) adjacent to a thick position satisfies
\begin{equation}\label{eq:B8-local-band}
 |\ell s_a-b_*|\le320\{\ell\delta+\ell S+b_*/\ell+s_a+\ell\log_\ell(2h)\},
\end{equation}
and \(I=0\) otherwise. Then, for every integer \(e'\ge0\) with \(\ell^{e'}Z_f(T_2)\subset I_F\) (by (a), applied with \((M_a,h_a)=(\delta,s_a+S)\) when \(\varsigma_a=1\) and \((s_a+S,\delta)\) when \(\varsigma_a=-1\), such an \(e'\) exists and satisfies \eqref{eq:B8-unrestricted-exp}),
\begin{equation}\label{eq:B8-short-col}
 \lambda_O\bigl(Z_f(T_2)/I_F\bigr)\le320(d_u+1)r\bigl\{h/\ell^2+(w_*+\delta)/\ell+S+1+\log_\ell(2h)+\delta I\bigr\}+320d_ume'.
\end{equation}
\end{enumerate}
\end{lemma}
\begin{proof}[Outline]
Prove (a) from two level-one Wach lifts and the shifted-weight kernel valuation. For (b), use the unrestricted lift if \(b_*\) is small; otherwise cut before nonthick positions into intervals of length at most \(3\), split central jumps via (a), and apply arc-series lifts. The complete proof is in \S\ref{prf:B8.42bw}.
\end{proof}

\begin{lemma}[The non-place part of the row bound]\label{B8.42bxa}
Retain the cut setting of Lemma~\ref{B8.42bx}(b). Let \(\mathcal R\) and \(\mathcal C\) be the row and column sets, and put \(d_u=[E_u:\Q_\ell]\). Write \(h_W(u)\) and \(h_0\) for the normalized invariant lengths of Lemma~\ref{B8.41b}, and rename its projector valuation \(\delta_{\rm pr}\). Let \(c\ge1\), and suppose
\[
 c_B\le ct n n_VL,\qquad h_0\le ct,\qquad
 c_0(u)\le cn n_VL,\qquad h_W(u)\le crn n_VL
 \quad(u\in\mathcal C).
\]
Put
\[
\begin{aligned}
 D_0={}&tc_B\sum_{w\in\mathcal R}(d_w+2)
       +4n_V(|\mathcal C|+1)\delta_{\rm pr}
       +t|\mathcal C|h_0,\\
 D_{\rm sp}={}&D_0+\sum_{u\in\mathcal C}
       \bigl(d_um\delta_{\rm pr}+2\delta_{\rm pr}+h_W(u)+rc_0(u)\bigr),\\
 B_c={}&cfn t^2n_VL.
\end{aligned}
\]
Then \(D_0\le32B_c\) and \(D_{\rm sp}\le48B_c\). In particular, for \(c=120\), both are at most \(7200fn t^2n_VL^2\).
\end{lemma}
\begin{proof}[Outline]
The degree sum \(2f\) bounds \(|\mathcal C|\le2f\) and \(\sum_{w\in\mathcal R}(d_w+2)\le6f\). With \(m,r\le t\) and \(\delta_{\rm pr}\le2t\), termwise bounds in \(B_c\) give \(D_0\le32B_c\), \(D_{\rm sp}\le48B_c\), and \(c=120\). The complete proof is in \S\ref{prf:B8.42bxa}.
\end{proof}

\begin{lemma}[Inert word dichotomy and final-row congruence passage]\label{B8.42bx}
Assume \(\mathscr H_2\), and work on the twist-good cut family of Theorem~\ref{B8.42} with the norm-choice facts and four ties. Fix \(C_2\ge\max(1,C_0)\), where \(C_0\) is the threshold of Lemma~\ref{B8.42bv}. At a member satisfying the finiteness hypothesis of Definition~\ref{B8.40}, use the size width \(S_c\), error width \(\eta=\eta_{\rm sz}\), \(n=f/2\), \(n_V=r+t\), \(H_B=b+S_c+\eta+t+2\), thick density \(\delta_2=|T_2|/n>0\), progression-window length \(L_w\), order \(o_\ell\) of \(\ell\) modulo \(2^{e+2}\) up to sign for the host tower exponent \(e\), and base sizes \(b\), \(w''(n)\), \(b_T+\lfloor b'/x\rfloor\) (\(T\in TT\)). Put
\[
 Y=\frac{kL}{m\delta_2},\quad
 \eta_3=\eta+t+2,\quad
 w_B=\max\bigl\{w''(n)+\eta,\ \max_{T\in TT}(b_T+b'/x+\eta)\bigr\},\quad
 w_P=w_B+\tfrac{t+1}{2}.
\]
Eventually along the family, for every member and every twisted index \(i\) whose residue characteristic \(\ell=\ell_i\ge5\) does not divide \(Q_Sqdg_0g_1\), the following hold:
\begin{enumerate}
\item[(a)] Suppose \(\ell\) is inert in \(\Q(\sqrt{-d})\). Then \(\ell\ge11\), and for every place \(u\mid\ell\) of \(E\) the completion \(K=E_u\) is unramified of degree \(d_u\) over \(\Q_\ell\). Set \(T_0=(t+1)/2\), \(b_*=\lceil b\rceil-T_0\), \(S=\lfloor S_c\rfloor\), \(\delta=\lfloor\eta\rfloor+t+2\), and \(J=\max(1,\lceil\log_{\ell-1}Y\rceil)\). At every \(\pi\in Z_{\rm good}(i)\), over a finite splitting extension \(L'\) of \(\Q_\ell\) containing the images of all embeddings of \(E_u(\mu_\ell)\), with integers \(O\), the local comparison supplies an equivariant surjection \(F:T_1\to T_2\) of crystalline \(O\)-lattices over \(E_u\) onto the rank-\(r\) Tate-dual target lattice \(T_2\), together with an integer \(w_*\in\N\) and, for each \(a\in\Z/d_u\Z\), a sign \(\varsigma_a\in\{1,-1\}\) and an integer \(s_a\in\N\) satisfying the jump hypotheses of Lemma~\ref{B8.42bv}, such that \(T_2\) has at most \(m\) jumps in \([-\delta,\delta]\) whenever \(s_a>\delta\), every \(a\in\Z/d_u\Z\) is either thick (\(b_*\le s_a\le b_*+\delta\)) or nonthick (\(s_a\le w_*\), with \(|s_a-y|\le T_0\) for some \(y\in[w''(n),w_B]\)), and with \((M_a,h_a)=(\delta,\max(\delta,s_a+S))\) for \(\varsigma_a=1\) and \((\max(\delta,s_a+S),\delta)\) for \(\varsigma_a=-1\),
\begin{equation}\label{eq:B8-word-box}
 M_a+h_a\le2H_B,\quad b_*+\delta+S\le H_B,\quad w_*\le w_P,\quad S\le S_c,\quad\delta\le\eta_3.
\end{equation}
If \(\ell<Y\) and \(o_\ell>4L_w\), then \(2J+2\le L_w\), \(1\le J<d_u\),
\begin{equation}\label{eq:B8-long-word-condition}
 C_2(J+1)(w_*+S+\delta+1)\le b_*,
\end{equation}
and the thick positions in \(\Z/d_u\Z\) form singletons or adjacent pairs separated by at least \(J\) nonthick positions. Otherwise \(\varsigma_{a+1}=-\varsigma_a\) for all \(a\in\Z/d_u\Z\), and either \(d_u\) is even and every \(a\in\Z/d_u\Z\) is thick with \(\tfrac12(\max_{a'}s_{a'}+S)\le s_a\) and \(2A+\delta+1\le s_a\) (where \(A\) is given by Definition~\ref{B8.42bu} for \(\beta_a=s_a\) and \(d'=d_u\)), or
\begin{equation}\label{eq:B8-short-word-condition}
 1\le b_*,w_*,\quad 20w_*,20\delta,20S\le b_*,\qquad s_a>\delta\quad(\text{all }a),
\end{equation}
at least one position in \(\Z/d_u\Z\) is nonthick, and no three consecutive positions in \(\Z/d_u\Z\) are thick.
\item[(b)] Let \(L_{i,u}\in\Q_{\ge0}\) (\(u\mid\ell\)) be place allowances such that when \(\ell\) splits in \(\Q(\sqrt{-d})\), \(L_{i,u}=(2d_um+1)e_A\) on the canonical column set of \(i\) and \(0\) outside it, where \(e_A=\lceil23\{2+2\lfloor H_B\rfloor/\ell+\log_\ell(2\lfloor H_B\rfloor+1)\}\rceil\), and when \(\ell\) is inert in \(\Q(\sqrt{-d})\), \(L_{i,u}\) bounds \(\lambda_O(Z_f(T_2)/I_F)\) at \(u\) across all \(\pi\in Z_{\rm good}(i)\) (with \(L_{i,u}=0\) if \(Z_{\rm good}(i)=\varnothing\)). Then for every \(\pi\in Z_{\rm good}(i)\),
\[
 \mathcal B_{i,\pi}\Bigl(\rho;\ 7200 fn t^2 n_V L^2+\sum_{u\mid\ell}L_{i,u}\Bigr).
\]
\end{enumerate}
\end{lemma}
\begin{proof}[Outline]
Local de Rham and Fontaine comparisons build a common cyclic word, and scale estimates prove the inert dichotomy in (a). For (b), character-block and invariant estimates combine with Lemma~\ref{B8.42bxa} to bound the non-place congruence terms. The complete proof is in \S\ref{prf:B8.42bx}.
\end{proof}

\begin{lemma}[Local costs on the cut]\label{B8.42b}
Assume \(\mathscr H_2\). On the twist-good cut family of \ref{B8.42} with the norm-choice facts and four ties, there is one family constant \(C>0\) such that eventually, at every member satisfying the finiteness hypothesis of Definition~\ref{B8.40}, with nonnegative costs \(A_\ell\) specified below, such that
\begin{align}
 \eta_{j,\ell}&\le d_{j,0}\log\ell\bigl(A_\ell+\nu_\ell+2\ell^{-2}
 +\sum_{1\le e<t}s_{\ell,e}(E)+v_\ell(h_E)+v_\ell(M_0)\bigr),\label{eq:B8-cost-prime}\\
 6+\log(Q_Str)&+\log h_E+\log M_0+\sum_{\ell\in S}A_\ell\log\ell
 \le2000(C+1)^3 fmbL^2\qquad(S\text{ any finite set of primes}).\label{eq:B8-cost-sum}
\end{align}
Here \(\eta_{j,\ell}\) is the part of \(\eta_j(\mathcal L)\) above \(\ell\). In applying \eqref{eq:B8-cost-sum} to the family estimate of Definition~\ref{B8.40}, take \(S\) to be the characteristics of the twisted support.
\end{lemma}
\begin{proof}[Outline]
Bound \(\eta_{j,\ell}\) by cases via Theorem~\ref{B8.41}, Lemmas~\ref{B8.42bp}, \ref{B8.42br}, \ref{B8.42bt}, \ref{B8.42bv}, \ref{B8.42bw}, \ref{B8.42bx}, and Propositions~\ref{B8.42bq1}, \ref{B8.42bu1}. Summing \(A_\ell\log\ell\) via integral comparisons gives the global bound. The complete proof is in \S\ref{prf:B8.42b}.
\end{proof}

\section{The closed local condition in the witness}

\begin{lemma}[Artin action on the ratio of toric sums]\label{B8.43i}
Use the host \(E\supset K=\Q(\sqrt{-d})\) and the toric data of \ref{B1.14k}(3). Thus \(E/E^+\) is unramified at finite places, the CM embeddings restrict to one embedding of \(K\), \(s,\kappa_i\ge1\), and \(\chi\) has the prescribed base conductor \(\mathfrak f\). Let \(Q\) be a finite set of split primes prime to \(\mathfrak f\). Write \(T(q)\) for the class term at \(q\in\operatorname{Pic}(O_Q)\), and \(T_0\ne0\) for the sum at the base level with trivial Picard twist. Fix an algebraic closure of \(K\) in \(\C\). For every character \(\psi\) of the finite group \(\operatorname{Pic}(O_Q)\), the number
\[
             R_\psi=\frac{\sum_q\psi(q)^{-1}T(q)}{T_0}
\]
belongs to this algebraic closure. Suppose \(g\) fixes the embedded CM-type values and the finite character values \(\chi_f(y)|y|_{E,f}^{s}\). Choose \(t\in\mathbb A_{K,f}^{\times}\) with \(\operatorname{Art}_K(t)^{-1}=g|_{K^{\rm ab}}\), and write \([t]\) for its image in \(\operatorname{Pic}(O_Q)\). Then
\[
                      gR_\psi=(g\psi)([t])^{-1}R_{g\psi}.
\]
\end{lemma}
\begin{proof}[Outline]
Apply \ref{B1.14k}(3), proved in \ref{B1.15a}, to the twisted and base sums with a common period. Translate the Picard classes by \([t]^{-1}\) in the transport formula, then divide to cancel the common period and transport factors. The complete proof is in \S\ref{prf:B8.43i}.
\end{proof}

\begin{lemma}[The denominator of the finite Fourier expansion]\label{B8.43j}
Let \(L\) be the completed finite tensor lattice of a self-dual source and target pair in the Weil representation. Write \(\langle\, ,\,\rangle\) for its alternating pairing and \(\psi\) for the additive character, normalized so that \(L\) is its own annihilator. At a collection of split places let \(e\) select one factor and \(\bar e\) the conjugate factor. Put \(L_+=eL\), \(L_-=\bar eL\). Let \(M\ge1\), and let \(R\subset\C\) contain the \(M\)-th roots of unity. Let \(\Phi\) be an \(R\)-valued function supported on \(L_-\) and constant modulo \(ML_-\). Write \(f_\Phi\) for its vector in the lattice model and \(f_0\) for the vacuum. There is an integer \(K\), fixed by the original tensor lattice, such that \(M^K f_\Phi\) belongs to the \(R\)-span of the Heisenberg translates of \(f_0\) by \(M^{-1}L_+\). In particular, at a place over \(\ell\), for \(M=M_0d'\) with \(\ell\nmid d'\), the power \(\ell^{K\operatorname{ord}_\ell(M_0)}\) suffices over the local ring containing the values of \(\Phi\) and the roots of unity.
\end{lemma}
\begin{proof}[Outline]
Project integral lattice generators to obtain the Fourier characters. Self-duality and character orthogonality express \(M^K f_\Phi\) as a sum of vacuum translates, and the prime-to-\(\ell\) part of \(M^K\) is a local unit. The complete proof is in \S\ref{prf:B8.43j}.
\end{proof}

\begin{lemma}[Two denominators independent of the auxiliary level]\label{B8.43k}
Fix a prime \(\ell\). Let \(J\) be a complex linear map into a polynomial space, and let \(L_0\) be a finitely generated additive lattice in its domain. There is \(c_1\ge0\) such that, in every embedded number field containing the coefficients of \(J(L_0)\), all coefficients of \(\ell^{c_1}J(L_0)\) are integral at every place over \(\ell\).

For the level factor, let \(E\) be the CM field, \(\mathcal L\) the integral marking lattice of \(\Z\)-rank \(s\), \(r\) the target rank, \(n\) the source rank and \(D\) the degree of the harmonic polynomial. Put \(M=M_0d'\) and \(N=Mu\), where \(M_0,d',u\ge1\) and \(\ell\nmid d'u\). Let \(\mathfrak b\) be the marking ideal with \(M\mathcal O_E\subseteq\mathfrak b\). For a lattice \(\Lambda\subset E^n\) satisfying \(N\mathcal O_E^n\subseteq\Lambda\cap\mathcal O_E^n\), put
\[
 I=[\Lambda\cap\mathcal O_E^n:N\mathcal O_E^n],\qquad
 A=\mathcal L/\mathfrak b\mathcal L,\qquad
 H=I^r |A|^n(2|\operatorname{disc}E|N)^D .
\]
Then \(\operatorname{ord}_\ell H\le c_3\), where
\[
 c_3=([E:\Q]nr+sn)\operatorname{ord}_\ell M_0
                 +D\operatorname{ord}_\ell(2|\operatorname{disc}E|M_0).
\]
\end{lemma}
\begin{proof}[Outline]
Clear the \(\ell\)-denominators of the polynomial coefficients on additive generators of \(L_0\). Bound the level denominator using \(I\mid N^{[E:\mathbb Q]n}\), \(|A|\mid M^s\), and the equal \(\ell\)-valuations of \(M,N,M_0\). The complete proof is in \S\ref{prf:B8.43k}.
\end{proof}

\begin{lemma}[The different in descent to the coefficient lattice]\label{B8.43l}
Let \(O\subset\C\) be a discrete valuation ring, let \(\ell\) be a rational prime in its maximal ideal, and put \(e=\operatorname{ord}_O(\ell)\). Let \(K\subset L\subset\operatorname{Frac}(O)\) be number fields. Denote the places selected by \(O\) on \(K,L\) by \(v,w\). Use integer-valued valuations on both fields. For \(y\in L\), suppose there is \(b\ge0\) such that
\[
 \operatorname{Tr}_{L/K}(ay)\in\mathcal O_{K,v}
 \quad\text{for every } a\in\mathcal O_L
 \text{ with }a\in (w')^b
 \text{ at every }w'\mid v,\ w'\ne w .
\]
Here \(w'\) in an ideal power denotes the prime ideal. Then \(\ell^{e^2}y\in\mathcal O_{L,w}\). If \(L_w/K_v\) is unramified, \(y\in\mathcal O_{L,w}\).
\end{lemma}
\begin{proof}[Outline]
Use the Chinese remainder theorem and the completed trace pairing to place \(y\) in the local inverse different. Bound its exponent using the trace from the maximal unramified subfield, and use the perfect residue trace pairing in the unramified case. The complete proof is in \S\ref{prf:B8.43l}.
\end{proof}

\begin{lemma}[The short theta value on an integral frame]\label{B8.43m}
Use the lattice and marking of Lemma~\ref{B8.43k}. Write \(n\) for the number of theta columns, \(r\) for the target rank, \(A=\mathcal L/\mathfrak b\mathcal L\), and \(\Delta=\operatorname{disc}E\). Suppose that \(\Lambda\) is self-dual, agrees with the standard lattice at the primes dividing \(M\), and that \(N\) is an odd multiple of \(M\) with \(N\mathcal O_E^n\subseteq\Lambda\cap\mathcal O_E^n\). Fix a place of the coefficient field \(K\). Let \(w:A^n\to K\) and the coefficients of a homogeneous harmonic polynomial \(P\) of degree \(D\) be integral there. On every integral frame of good reduction over that place the short theta value \(v\), with class coefficient one, satisfies
\[
       I^r |A|^n (2\Delta N)^D v\ \in V ,
       \qquad I=[\Lambda\cap\mathcal O_E^n:N\mathcal O_E^n].
\]
Here \(V\) is the complete discrete valuation ring of the test frame. Good reduction means that its generic abelian variety with its endomorphisms and polarisation extends to an abelian scheme over \(V\).
\end{lemma}

\begin{proof}[Outline]
Rationalizes \(N^2\)-torsion and checks integrality of the germs \(F_{qzb}\). Adjusts the logarithmic gauge by a balanced symmetric form (killed by harmonic \(P\)) via a trace-dual basis, then proves by induction that \(2\Delta\) clears each covariant derivative. The complete proof is in \S\ref{prf:B8.43m}.
\end{proof}

\begin{lemma}[Testing a cusp section at a degenerating frame]\label{B8.43n}
Let \(M\) be the level of a uniformised unitary moduli space, let \(\lambda\mid\ell\) be a coefficient place, and suppose that \(M/\ell^{\operatorname{ord}_\ell M}\ge3\). Let \(f\) be a cusp section of nonnegative dominant weight and nonnegative theta degree, cuspidal at every designated complex embedding, and let \(a\) be in the coefficient field. If \(af\) is integral at every good-reduction integral frame over \(\lambda\) on complete discrete valuation rings with finite residue field, it is integral at every integral frame on those rings over \(\lambda\).
\end{lemma}

\begin{proof}[Outline]
Excludes definite degeneration and lifts a split degeneration to a toroidal arc giving \(h(t)/(\pi t)^d\). Kills negative \(t\)-powers via Cauchy's formula on the complex arc, and proves \(\pi^d\mid h/t^d\) by specializing over \(V[[t]]_{(\pi)}\) to finite-residue traits. The complete proof is in \S\ref{prf:B8.43n}.
\end{proof}

\begin{proposition}[One clearing power on the ray covers]\label{B8.43o}
Fix a standard layer, its pinned base frame, its lowest-weight tensor and its coefficient map. There is a phase \(u\), independent of the cover, with the following property. Fix a rational prime \(\ell\), a finite source translate \(s\), and a finitely generated additive lattice \(L_0\) of coefficient vectors. On all the covers and for all the twists, places and integral tests specified in Theorem~\ref{B8.43h}, the reading multiplied by \(\ell^c u\) belongs to the indicated integral cusp-section lattice, for a single integer \(c\ge0\). On a good rational ray cover the phase-normalised reading is rational for every rational test and rational coefficient vector, also after the allowed changes of splitting.
\end{proposition}

\begin{proof}[Outline]
Fixes \(u\) from the base pin, expands tests via \ref{B8.43k} and \ref{B8.43j}, flips half-lattices, and moves the point to split primes by unitary reflections. Forms integral cusp sections via \ref{B8.43m}, \ref{B8.25t}, and \ref{B8.43n}, then deduces rationality. The complete proof is in \S\ref{prf:B8.43o}.
\end{proof}

\begin{lemma}[Integral descent of an invariant reading]\label{B8.43p}
Let \(K\subset L\) be the base and cover coefficient fields of a uniformised tower cover of base level \(M_0\), pinned at the nonnegative dominant weight and nonnegative theta degree in question. Let \(F\) be a rational cover cusp reading invariant under the base compact subgroup. Fix a rational prime \(\ell\) with \(M_0/\ell^{\operatorname{ord}_\ell M_0}\ge3\), and let \(S\) be a set of places of \(L\) over \(\ell\) where \(F\) is integral. If every \(w\in S\) is unramified over \(K\), then \(F\) is a finite sum of base cusp readings integral at the places below \(S\), with coefficients in \(L\) integral at \(S\). At a single possibly ramified place \(w\mid\ell\) where \(F\) is integral, if \(K,L\) are in the fraction field of the fixed valuation ring \(O\) of Lemma~\ref{B8.43l} and \(\operatorname{ord}_O(\ell)=e\), the same conclusion holds at \(S=\{w\}\) for \(\ell^{e^2}F\).
\end{lemma}

\begin{proof}[Outline]
Étale descent and flatness write \(F\) in a local integral base cusp basis (\ref{B8.43n}). Lifting integral frames gives integral traces so \ref{B8.43l} bounds the coefficients at each \(w\in S\), and the Chinese remainder theorem glues the decompositions. The complete proof is in \S\ref{prf:B8.43p}.
\end{proof}

\begin{lemma}[Transport of the theta generator by a Hecke isogeny]\label{B8.43q}
Let \(B\) be a marked geometric fibre of the unitary moduli space in characteristic zero, with principal rigidified theta line \(L\). Let \(P\) be a maximal ideal of \(\mathcal O_E\) above the rational prime \(p\), with \(P+\bar P=\mathcal O_E\) and \(p\) prime to the full level. Write \(g=\dim B\) and \(q_P=\mathrm NP\). For an integer order \(i\ge0\), the kernels are the \(\mathcal O_E\)-stable subgroups \(C\) satisfying
\[
 B[(p)P^{-1}]\subseteq C\subseteq B[(p)\bar P],\qquad
 e_{p^2}|_{C\times C}=1,\qquad
 |C|=p^{2g},\qquad [B[p]:C\cap B[p]]=q_P^i ,
\]
where \(B[J]\) is annihilated by \(J\) and \(e_{p^2}\) is the Weil pairing of the principal polarisation. There are finitely many such kernels. For each of them there is a unique quotient framed point with the prescribed transported marking. If \(\Psi:B\to B'=B/C\) is the quotient, its tangent basis is \(p^{-1}d\Psi\) of the source basis. For a generator \(s\) of \(H^0(B,L)\), its theta generator \(s'\) is uniquely determined by
\[
           \Psi^*s'=\sum_{c\in C/(C\cap B[p])}
                    t_c^*\bigl(\beta^{-1}[p]^*s\bigr).
\]
In this formula \([p]^*L\) is identified with \(L^{p^2}\) by rigidification, \(\beta:\Psi^*L'\simeq L^{p^2}\) is rigidified, \(L'\) is the marked principal line on \(B'\), and \(t_c\) denotes translation by any representative. The summands on the right are sections of \(\Psi^*L'\), with its natural kernel linearisation.
\end{lemma}

\begin{proof}[Outline]
Descends the polarisation and marking to \(B/C\) from isotropy of \(C\subset B[p^2]\) and identifies \(\Psi^*L'\simeq L^{p^2}\) via biextensions. Descends the \(C\)-invariant trace sum by faithful flatness, and proves nonvanishing by \(B[p]\)-character orthogonality. The complete proof is in \S\ref{prf:B8.43q}.
\end{proof}

\begin{lemma}[Hecke frames and integral values away from \(p\)]\label{B8.43r}
Let \(\lambda\nmid p\) be a place of the coefficient field. Start with an algebraic integral test frame over a complete discrete valuation ring \(V\) above \(\lambda\) with finite residue field. In an algebraic closure of its fraction field take any Hecke target of Lemma~\ref{B8.43q}. After finite valued extension it is obtained from an integral target frame by a triangular frame change with unit diagonal, including a unit on the theta line. If the source has good reduction the target frame may be taken of good reduction. A section integral on the algebraic test frames at \(\lambda\) has value integral over \(V\) at this target.
\end{lemma}

\begin{proof}[Outline]
Extends \(\Psi,\Psi'\) to semiabelian models so \(p^{-1}d\Psi\) identifies integral Lie lattices. Proves regularity of the trace in \ref{B8.43q} via Fourier bounds (\(v(a)\ge0\)), then uses a reverse trace along \([p]\Psi'\) and orthogonality to get \(v(a)=0\). The complete proof is in \S\ref{prf:B8.43r}.
\end{proof}

\begin{proposition}[The Hecke transform on sections and cusps]\label{B8.43s}
Use the kernels of Lemma~\ref{B8.43q}. For a marked unitary space of positive rank \(r\), \(0\le i\le r\), nonnegative even theta degree \(n\), scalar \(b\) in the coefficient field and section \(f\) of weight \(a\) and theta degree \(n\), there is a unique section of the same weight and degree whose value at a geometric frame \(y\) is
\[
       (T_{P,i,b} f)(y)=b\,q_P^{-in/2}\sum_C f(y_C).
\]
Here \(y_C\) is the transported frame of Lemma~\ref{B8.43q}. In the uniformised layers this operator preserves the cusp subspace for dominant weight. More generally, whenever a section transform of dominant weight is given by this kernel formula, it preserves the cusp subspace. At \(\lambda\nmid p\) the displayed operator preserves integral sections if \(b\) is integral at \(\lambda\).
\end{proposition}

\begin{proof}[Outline]
Realizes the kernels of \ref{B8.43q} as a finite étale correspondence, builds the section by trace, and checks integrality via \ref{B8.43r}. Bounds \(|t_C|\le q_P^i\) by Cauchy--Schwarz and dominates the transform by finitely many rational isometries for cuspidality. The complete proof is in \S\ref{prf:B8.43s}.
\end{proof}

\begin{proposition}[The Hecke reading in the pinned frame]\label{B8.43t}
Fix the designated embedding \(\iota\) of a pinned layer \(j\) and put \(n=n_V\), its even source rank. At an admissible split Hecke prime \(P\mid p\), take the element \(\alpha\) of order \(i\) in the split matrix model and a set \(\mathcal R\) of representatives of \(K\alpha K/K\), where \(K\) is the member's compact level. Write \(\mu_j(P)\) for the reference unramified character at a uniformiser of \(P\). For a section \(f\) of the layer weight let \(R_f(h)\) be its evaluation in the frame at \(h\), multiplied by the layer's Lie normalisation, its pinned component constant and its prescribed adelic multiplier. If \(f'\) is the transform of Proposition~\ref{B8.43s}, then
\[
            R_{f'}(h)=\iota(b)\mu_j(P)^{-i}
                                  \sum_{a\in\mathcal R}R_f(ha).
\]
\end{proposition}

\begin{proof}[Outline]
Bijects the kernels of \ref{B8.43q} with \(K\alpha K/K\), matching frames up to unipotent Lie changes. Compares the Weil action \(\omega(1,a)f_0\) with the trace in \ref{B8.43q} at a rational period vector as in \ref{B8.43o} to find \(t^n\), then sums over \ref{B8.43s}. The complete proof is in \S\ref{prf:B8.43t}.
\end{proof}

\begin{theorem}[Series laws for the third-layer construction]\label{B8.43h}
Use the third-layer nonendoscopic inputs and the geometric fields of \(P\) in \ref{B8.01}. Over a nonarchimedean local field, Bernstein's unitary irreducibility theorem is used for the normalized parabolic induction of two smooth irreducible unitarizable blocks \cite[\S8.2]{bernstein-1984}. For two smooth irreducibles with disjoint unramified twist classes of cuspidal support, the normalized product is irreducible by \cite[Proposition~8.5]{zelevinsky-1980}. In this case any two segments from opposite multisegments have disjoint cuspidal lines, so they are unlinked. The local and Galois compatibilities include \ref{B8.42p} and the residual Brauer--Nesbitt input of \ref{B8.42pa}. At each embedding \(\tau\) over a de Rham place \(w\) of a rank-\(a\) representation \(V\), the geometric jump clause is \(\sum_i\dim\operatorname{gr}^iD_{\rm dR}(V)_\tau=a\).

For each pinned standard layer the phase-normalised series is rational on every good rational ray cover, including for the permitted determinant twist and change of splitting in Proposition~\ref{B8.43o}. For a fixed standard layer with its pinned frame, extremal tensor and coefficient map, the integral assertion uses a phase \(u\in\C\), \(|u|=1\). For each prime \(\ell\), finite source translate \(s\) and finitely generated additive lattice \(L_0\) of coefficient vectors, there is an integer \(c\ge0\). It works on every uniformized tower cover with squarefree auxiliary level prime to \(\ell Q_Sq\,\mathrm N\mathfrak c_K\), pinned at every dominant weight of the prescribed degree. The cover field must contain the ray field of the full level in the designated embedding, the coefficients of \(L_0\), and the normalized finite values of the reference splitting character of weight \(-1/2\). Let \(S\) be any set of its places over \(\ell\). For every level test with coefficient function integral at \(S\), every \(v\in L_0\), and every determinant twist \(\xi\) defined over the cover field on norm-one idèles and trivial there on rational and auxiliary-congruence elements, the function \(g\mapsto\ell^c u\,\xi(\det g_f)\Theta(s,g;v)\) is the reading of a cusp section integral at \(S\). Here \(\Theta\) uses the indicated level test and coefficient map. The descent assertion concerns a cover integral cusp form invariant under the base compact level. If the cover coefficient field is unramified at every place of \(S\) over the base coefficient field, the form is a finite linear combination of base cusp forms integral at the places below \(S\), with coefficients in the cover field integral at \(S\). At each split Hecke prime outside the level, the transported frames and their values satisfy Lemmas~\ref{B8.43q} and~\ref{B8.43r}, the latter at coefficient places away from that prime. Proposition~\ref{B8.43s} gives the section transform and preserves cuspidality. At the designated member embedding its reading is the normalised double-coset sum of Proposition~\ref{B8.43t}. For a nonzero base toric sum \(T_0\), the split ray-level character sums divided by \(T_0\) are algebraic, with the action of Lemma~\ref{B8.43i}. In particular, for a ray character pulled back to a Picard character \(\psi\) and an integer \(h\), the action is \(gR_{\psi^h}=(g\psi)([t])^{-h}R_{(g\psi)^h}\) whenever \(g,t\) satisfy the hypotheses of that lemma. This follows by applying its formula to \(\psi^h\), since \(g\) commutes with integer powers.
\end{theorem}
\begin{proof}[Outline]
Verifies the hypotheses of \ref{B8.43o} for rational and integral series readings and uses \ref{B8.43p} for unramified descent. Obtains the Hecke laws from \ref{B8.43q}, \ref{B8.43r}, \ref{B8.43s}, and \ref{B8.43t}, and the toric action from \ref{B8.43i}. The complete proof is in \S\ref{prf:B8.43h}.
\end{proof}

\begin{theorem}[The closed family local-error estimate]\label{B8.43}
Assume \(\mathscr H_2\), the third-layer and geometric hypotheses of Theorem~\ref{B8.43h}, depth-zero support, and the decomposition-group inputs of \ref{B8.41p}. For a local Langlands witness \(h_L\) and \(c_{\rm Dep}\in\R\), take the list-extension parameters \(p=p_0(h_L,c_{\rm Dep})\) of \ref{B8.01}, with the default geometric closure. For these parameters, the linked closed estimate at \(2,3\) implies that the family local-error estimate of Definition~\ref{B8.40} is closed on \(W_3(\mathcal A)\) for every linked selection.
\end{theorem}
\begin{proof}[Outline]
Verifies the hypotheses of \ref{B8.42} from the third-layer, series, and depth-zero data via \ref{B8.42pa}, then applies \ref{B8.42} and \ref{B8.41} on each cut family. Strengthens and projects to the selection of \ref{B8.04} using \ref{B8.03} and \ref{B8.43h}. The complete proof is in \S\ref{prf:B8.43}.
\end{proof}



% ---- p8c5 ----
\chapter{Primes in short intervals, masses and counts}\label{ch:p8c5}

Let \(E\) be the CM field of the arithmetic datum \(D\), with the Artin map of \ref{B2.01}. Put \(F=E^+\) and \(f=[F:\Q]\). Write \(m,k,r\) for the chain length and the starting and spectral ranks, respectively. Let \(T_u\) be the level prime at word \(u\). Put \(L=x^\varepsilon\) and \(b=e^L\), with \(0<\varepsilon<1/52\).

\section{Uniform estimates for Dirichlet \(L\)-functions}

For primitive \(\chi\bmod N\), put \(\psi_0(\chi,u)=\sum_{n\le u}\Lambda(n)\chi(n)-{\bf1}_{\chi=1}u\). Write \(m_\chi(\rho)\) for zero multiplicity. The sets \(Z_\chi(T)\) and \(Z_\chi[t]\) consist of distinct zeros in the open critical strip with \(|\Im\rho|<T\) and \(|\Im\rho-t|\le1\), respectively.

\begin{definition}[The uniform analytic input]\label{B8.44}
The uniform analytic input consists of Theorems~\ref{B8.45} and~\ref{B8.46} with positive constants \(C_1,c,C_7,C_{11}\) independent of the arithmetic datum \(D\).
\end{definition}

\begin{theorem}[Density near \(1\) and logarithmic zero regions]\label{B8.45}
There are \(C_1,c>0\) such that, for \(Q\ge3\), every finite set of triples \((N,\chi,\rho)\) with primitive \(\chi\bmod N\), \(1\le N\le Q\), and
\[
 L(\rho,\chi)=0,\qquad \Re\rho\ge399/400,\qquad |\Im\rho|\le Q
\]
has total multiplicity at most \(C_1Q^{1/4}\). For primitive \(\chi\bmod N\), \(N\ge3\), the region \(\Re s\ge1-c/\log(\max(N,N|\Im s|))\) contains at most one zero. Such a zero is real and simple and \(\chi\) is real. Every open-strip zero of \(\zeta\) has \(\Re\rho\le1-c/\log(|\Im\rho|+3)\).
\end{theorem}

Montgomery's density estimate gives \(O(Q^{1/8})\) occupied cells \((N,\chi,\lfloor\Im\rho\rfloor)\) at this threshold \cite[Ch.~12]{montgomery-1971}. By \ref{B8.46} their total multiplicity is \(O(Q^{1/8}\log Q)\), hence \(O(Q^{1/4})\). The two regions use a common constant in Davenport's zero-free regions \cite[\S\S13--14]{davenport-1980}.

\begin{theorem}[Explicit formula and unit-window zero count]\label{B8.46}
There are \(C_7,C_{11}>0\) such that for every primitive \(\chi\bmod N\), \(N\ge1\), and \(2\le T\le u\), the set \(Z_\chi(T)\) is finite and
\[
 \left|\psi_0(\chi,u)+\sum_{\rho\in Z_\chi(T)}m_\chi(\rho)\frac{u^\rho}{\rho}\right|
 \le C_7\left(\frac{u\log^2(NuT)}{T}+\log(Nu)
 +\sum_{\rho\in Z_\chi(1)}\frac{m_\chi(\rho)}{|\rho|}\right).
 \tag{8.46.1}\label{eq:p8c5-explicit}
\]
For each \(t\in\R\), \(Z_\chi[t]\) is finite and
\[
 \sum_{\rho\in Z_\chi[t]}m_\chi(\rho)\le C_{11}\log\bigl(N(|t|+2)\bigr).
 \tag{8.46.2}\label{eq:p8c5-window}
\]
\end{theorem}

The first bound replaces Davenport's signed low-zero correction by absolute values. The conductor-one formula and the unit-window count use the corresponding estimates for zeta and primitive characters. \cite[\S\S15--17 and \S19, (9)]{davenport-1980}

Let \(q,d\) and \(Q_S\) be the two primes and the prime-to-\(qd\) factor of the Artin modulus in \ref{B2.01}. Write \(N_E=Q_Sqd\) for that modulus and put \(\lambda_E=\max(1,\log N_E)\). Let \(s_D\) be the odd part of \(Q_S\). A primitive Dirichlet character \(\chi\) belongs to \(D\) if there is a character \(\omega:\Gal(E/\Q)\to\C^\times\) such that \(\chi(a)=\omega(\operatorname{Art}_D(a))\) for every positive integer \(a\) coprime to \(N_E\). For \(y\ge1\) define
\[
 g_\chi(y)=\sum_{n\le y}\Lambda(n)\chi(n)(y/n-1)
 +\begin{cases}yL'(1,\chi)/L(1,\chi),&\chi\ne1,\\
 -y(\log y-1-\gamma),&\chi=1.
 \end{cases}
 \tag{8.47.1}\label{eq:p8c5-centred-sum}
\]
Put \(X_1(Y)=e^{Y^{\varepsilon/5}}\), \(X_2(Y)=e^{Y^{\varepsilon/5}(\log Y)^2}\). Let \(\mathcal Z(H)\) be the set of real primitive characters of conductor at most \(H\) for which \(L(s,\chi)\ne0\) throughout \(\Re s>399/400,\ |\Im s|\le H\). The zero-free box hypothesis for a family of exponent \(J\) requires that for all sufficiently large \(x\) and every member \(i\in I_x\) there be a \(Y\in[x^{1/J},x]\) with the following properties. For every real character \(\chi\) of conductor \(N\) belonging to the datum at \((x,i)\),
\[
 \begin{split}
 d\mid N&\ \Longrightarrow\ \chi\in\mathcal Z(8s_DX_1(Y)X_2(Y)),\\
 q\mid N,\ d\nmid N&\ \Longrightarrow\ \chi\in\mathcal Z(8s_DX_1(Y)).
 \end{split}
 \tag{8.47.2}\label{eq:p8c5-boxes}
\]
At this same \(Y\), put \(\chi_q=(-q/\cdot)\), \(\chi_d=(-d/\cdot)\), and \(\chi_8=(2/\cdot)\). For \(\eta=1,\chi_8\), every primitive character agreeing with \(\eta\chi_q\) on integers coprime to \(8qd\) must belong to \(\mathcal Z(8X_1(Y))\). Every primitive character agreeing there with \(\eta\chi_d\) or \(\eta\chi_d\chi_q\) must belong to \(\mathcal Z(8X_1(Y)X_2(Y))\). The hypothesis on the auxiliary primes requires that for all sufficiently large \(x\) and every \(i\in I_x\) there be a \(Y\in[x^{1/J},x]\), possibly different from the one used for the boxes, such that
\[
 \begin{gathered}
 q\equiv3\bmod8,\quad X_1(Y)^{1/2}<q\le X_1(Y),\qquad
 d\equiv7\bmod16,\quad X_2(Y)^{1/2}<d\le X_2(Y),\\
 (-d/q)=1,\qquad (-d/p)=1\quad(p\le Y^{\varepsilon/5},\ p\text{ prime}),\qquad
 \log N_E\le2Y^{\varepsilon/5}(\log Y)^2 ,
 \end{gathered}
 \tag{8.47.3}\label{eq:p8c5-qd}
\]
where the symbols are Kronecker symbols.

\begin{proposition}[The character-sum estimate on a family]\label{B8.47}
Assume the uniform analytic input of Definition~\ref{B8.44}. Let \(F_m\) be a family with starting scale \(x_1>1\) and fixed family constant \(C_o\) satisfying \(o\le C_o\) and \(Q_S\le C_o|A|\) for all \(x>x_1\) and \(i\in I_x\) (as built into the family structure for the host group \(A\), character order \(o\ge1\), and \(2\)-exponent \(e\in\N\) of \ref{B2.01}, where \(|A|=2^e o\) and \(2^{e+2}\mid Q_S\)). Suppose \(F_m\) satisfies the zero-free box hypothesis, including the six bounds following \eqref{eq:p8c5-boxes}, and the auxiliary-prime hypothesis \eqref{eq:p8c5-qd}. There are \(c_6>0\) and \(x_0>x_1\) such that for every \(x\ge x_0\), every member \(i\in I_x\), and every character \(\chi\) belonging to the datum at \((x,i)\),
\[
 |g_\chi(y)|\le y e^{-c_6x^{\varepsilon/8}}
 \qquad(y\ge e^{x^{\varepsilon/3}}/3).
 \tag{8.47.4}\label{eq:p8c5-thin}
\]
In the closed assertion of Definition~\ref{B8.39} one also assumes that some \(C_5>0\) works for all sufficiently large \(x\) and all \(i\in I_x\) as follows. Every primitive \(\chi\) agreeing with \(\chi_q\), \(\chi_d\), or \(\chi_q\chi_d\) on integers coprime to \(8qd\) satisfies \( |g_\chi(x)|\le x^{1/2+\varepsilon}\) and \( |L'(1,\chi)/L(1,\chi)|\le C_5\log x\).
\end{proposition}

\begin{proof}[Outline]
Combine the zero-free boxes and \ref{B8.45} with reflection and unit-window counts to bound the reciprocal zero weights. Apply the explicit formula at \(T=e^{\sqrt{\log u}}\), then use partial summation and the conductor bound to estimate the tail integral. The complete proof is in \S\ref{prf:B8.47}.
\end{proof}

\section{Standard masses and hyperplane volumes}

For a CM field \(K\) put \(h_M(z,w)=\sum_a\overline{z_a}w_a\), \(\Lambda_M=\mathcal O_K^M\), \(f_K=[K^+:\Q]\), and \(v_K=|\operatorname{disc}K|^{1/(2f_K)}/2\). Use the involution induced by complex conjugation in every embedding \(K\hookrightarrow\C\). The finite unramifiedness hypothesis says that for every maximal ideal \(\mathfrak P\subset\mathcal O_K\) some \(a\in\mathcal O_K\) satisfies \(a-\bar a\notin\mathfrak P\). Equivalently, \(K/K^+\) is unramified at every finite place. The mass and Euler factor are
\[
 \operatorname{mass}(h,\Lambda)
 =\sum_{[\Lambda']\in\operatorname{gen}(\Lambda)}|\operatorname{Aut}(h,\Lambda')|^{-1},
 \tag{8.48.1}\label{eq:p8c5-mass}
\]
\[
 L_\eta(K,M)=\prod_{v<\infty}(1-\eta_v^M q_v^{-M})^{-1},
 \tag{8.48.2}\label{eq:p8c5-euler}
\]
where \(v\) runs over the finite places of \(K^+\), \(q_v\) is its residue cardinality, and \(\eta_v=1\) at split places and \(-1\) at inert places. Set \(M_j=\operatorname{mass}(h_j,\mathcal O_E^j)\).

\begin{theorem}[The adjacent-rank mass formula]\label{B8.48}
For a CM field \(K\) unramified over \(K^+\) at every finite place and every \(M\ge3\),
\[
 \frac{\operatorname{mass}(h_M,\Lambda_M)}{\operatorname{mass}(h_{M-1},\Lambda_{M-1})}
 =\frac{\bigl((M-1)!(v_K/\pi)^M\bigr)^{f_K}}{(2v_K)^{f_K/2}}L_\eta(K,M).
 \tag{8.48.3}\label{eq:p8c5-mass-ratio}
\]
\end{theorem}

Apply Gan, Hanke and Yu's mass formula to the unramified standard lattice \cite[(4.2) and Propositions 4.4--4.5]{gan-hanke-yu-2001}. Its adjacent-rank quotient is \(2^{-f_K}|L(1-M,\eta^M)|\). The functional equation, with finite conductor one and \(|\operatorname{disc}K^+|=(2v_K)^{f_K}\), puts the quotient in the form \eqref{eq:p8c5-mass-ratio}.

The quotient-volume hypothesis on \(D\) is the following. For each Hermitian space \(W\) of compatible constant signature \((P,Q)\), take Haar measure on \(\U(W)(\mathbb A_F)\) with the compact-dual normalization at infinity and finite lattice stabilizer of measure one. Such a measure exists. For every such measure the action of \(\U(W)(F)\) admits a fundamental domain, and every fundamental domain has the prescribed volume \(\operatorname{vol}_D(W)\).

\begin{theorem}[The volume ratio for a unit hyperplane]\label{B8.49}
Assume \ref{B8.13} and the quotient-volume hypothesis on \(D\). Let \(W\) have rank \(M\ge3\) and compatible constant signature \((P,Q)\). For \(s=1\) if \(P\ge1\), or \(s=-1\) if \(Q\ge1\), there is a vector of norm \(s\) with complement \(W'\) of signature \((P-1,Q)\) or \((P,Q-1)\), respectively. With \(v_E=|\operatorname{disc}E|^{1/(2f)}/2\),
\[
 \frac{\operatorname{vol}_D(W')}{\operatorname{vol}_D(W)}
 =\left(\frac{(2\pi)^M}{(M-1)!}\right)^f
 \frac{(2v_E)^{f(1/2-M)}}{L_\eta(E,M)}.
 \tag{8.49.1}\label{eq:p8c5-volume-ratio}
\]

\end{theorem}

\begin{proof}[Outline]
Choose a unit coordinate vector in the standard signature form. Compare the one-vector Gram mean of \ref{B8.13} with its unfolding along the vector stabilizer, computing the real fibre by differentiation and the finite densities by residue norm counts and lifting. The complete proof is in \S\ref{prf:B8.49}.
\end{proof}

\begin{proposition}[The definite genus mass]\label{B8.50}
For every datum \(D\), let \(W\) have compatible signature \((M,0)\) at every real place, with \(M\ge1\). If \(h\) is its Hermitian form and \(\Lambda_0\) its self-dual base lattice, then \(\operatorname{mass}(h,\Lambda_0)=M_M\).
\end{proposition}

\begin{proof}[Outline]
Checking local equivalence at inert, split, and real places lets Landherr's theorem give a rational isometry with \(h_M\). Jacobowitz's local classification puts the transported self-dual lattice in the standard genus, and the isometry preserves automorphism orders. The complete proof is in \S\ref{prf:B8.50}.
\end{proof}

\begin{lemma}[Degree of a separated row]\label{B8.50a}
Use the family's compact-dual measure and signature \((p,q)=(p_Y,q_Y)\), \(p,q\ge1\), \(r=p+q\). Put \(F_s=\prod_{a=0}^{s-1}a!\). For the decreasing Harish--Chandra parameter \(\lambda\) at column \(\rho\), put \(\Delta(\lambda)=\prod_{a<b}(\lambda_a-\lambda_b)\), \(V(\lambda)=\Delta(\lambda)/F_r\), and \(H=\max_{a,b}|\lambda_a-\lambda_b|\). For \(e^{-r}\le t<1/2\), assume that opposite-sign differences have absolute value at least \(8192\min(p,q)(pq)^3\log r/t^2\). Then the formal degree and lowest-type dimension satisfy
\[
 (1-t)V(\lambda)\le d_{\Sigma,\rho}\le4d_\tau(3pqH)^{pq}.
 \tag{8.50.1}\label{eq:p8c5-degree}
\]
\end{lemma}

\begin{proof}[Outline]
The coordinate increment formula keeps the degree proportional to the Vandermonde along a ray of increasing gaps. Compare the Gaussian limit from Schur orthogonality and the minor bounds with the compact Weyl dimension formula to fix the constant and bound the degree. The complete proof is in \S\ref{prf:B8.50a}.
\end{proof}

\begin{lemma}[Volume at the starting signature]\label{B8.51b}
Assume the adjacent-rank formulas of \ref{B8.48} and \ref{B8.49}, and the definite rank-two volume formula of \ref{B8.50}. Let \(W\) be an admissible Hermitian space of constant signature \((P,Q)\) at the real places, with \(P\ge2\). Then \(\operatorname{vol}_D(W)=M_{P+Q}\). In particular \(\operatorname{vol}_D(Y_0)=M_r\) for the starting row space.
\end{lemma}
\begin{proof}[Outline]
Take successive orthogonal complements of unit vectors using \ref{B8.49}, removing negative directions first. By \ref{B8.48} each step preserves the volume-to-mass ratio, which is one by the definite rank-two formula of \ref{B8.50}. The complete proof is in \S\ref{prf:B8.51b}.
\end{proof}

\begin{proposition}[The scale of \(T_u\)]\label{B8.51}
Assume \ref{B8.49} and the rank-two case of \ref{B8.50} for every datum. For every member condition and bundle association, the following choice suffices for the bounds on the level primes:
\[
 c_{T,\min}(y)=\max(11,16+3\max(5,y+1)).
 \tag{8.51.1}\label{eq:p8c5-level-constant}
\]
On the family class \(W_1(\mathcal A)\) of \S\ref{ch:p8c1}, use the quantifiers of Definition~\ref{B8.39}. If \(C_{Tu}\ge c_{T,\min}(C_{YL})\), then for all sufficiently large \(x\), every member \(i\in I_x\) and every word \(u\) satisfy
\[
 b^{m/2}\le T_u,\qquad \log T_u\le C_{Tu}mL .
 \tag{8.51.2}\label{eq:p8c5-level-scale}
\]
Here \(C_{YL}\) is the family constant for the type primes.
\end{proposition}

\begin{proof}[Outline]
Insert the starting volume from \ref{B8.51b} into the product formula. Above the cutoff of \ref{B5.11c}, use the product and denominator bounds of \ref{B5.11} to bound \(\log T_u\) on both sides. The complete proof is in \S\ref{prf:B8.51}.
\end{proof}

\begin{lemma}[Selecting a prime at the word scale]\label{B8.51a}
Fix \(0<\varepsilon<1/52\) and \(c_s,C_Y>0\). For sufficiently large \(x\), let \(d\ge100\) be prime with \(d\equiv3\bmod4\), and take a positive multiple \(N\) of \(d\) with \(\log N\le3x^{\varepsilon/5}(\log x)^2\). Let \(g_0,g_1,m\) be positive integers with \(\log g_i\le C_Yx^\varepsilon\) and \(m\ge x^{1/3}\). For \(1\le V\le x^{12}\) and \(\log S\ge c_smx^\varepsilon\), there is a prime \(S<T\le(1+1/V)S\) with
\[
 T=\pi_T\bar\pi_T,\qquad
 \pi_T\equiv1\bmod M\mathcal O_K,\qquad
 M=N(g_0g_1)^2,\quad K=\Q(\sqrt{-d}).
\]
\end{lemma}

\begin{proof}[Outline]
Bound \(\log(4dM^2)\) and \(\log V\) by \(c\sqrt{\log S}\) and apply \ref{B2.50} with width \(1/V\). The resulting quadratic representation and congruences give an integral \(\pi_T\) of norm \(T\) congruent to \(1\) modulo \(M\mathcal O_K\). The complete proof is in \S\ref{prf:B8.51a}.
\end{proof}

\section{Two counting statements on the same selected data}

Write \(\mathcal D_3=W_3(\mathcal A)\) for the selected family class of \S\ref{ch:p8c1}, and \(\ell_{\rm out}\) for the family's list of outer constants in Definition~\ref{B8.39}.

\begin{lemma}[The labelled space before counting]\label{B8.52a}
Fix a point with the label realization of \ref{B8.08}. Let \(X\) denote that realization, including its unit lowest-type vectors, level subgroup, Haar measures and finite level projector \(e\). Write \(\sigma\) for the product representation at infinity, \(w\) for the product unit vector, and \(T_\phi\) for the real convolution of \ref{B8.03a}. Assume that the rank is at least \(1801\) and the adjacent auxiliary weight gaps are at least \(16r^{19}\), where \(r\) is the rank.

A label \(t\) specifies the finite type at every place over both distinguished primes. For a set \(S\) of such labels, let \(e_S\) be \(e\) multiplied, in the finite convolution algebra, by the sum of the type projectors for \(t\in S\). Put
\[
 V_S=\operatorname{im}(R_{e_S}T_\phi),\qquad d_t=\dim_\C V_{\{t\}} .
\]
Then \(V_S\) is finite dimensional and is identified by evaluation at \(w\) with the image of \(R_{e_S}\) on the continuous \(G_\infty\)-intertwiners from \(\sigma\) into \(L^2(X)\). There is a full index set with \(d_t\) indices labelled \(t\). Its good indices can be identified, preserving the label, with the selected indices of the point. For every \(S\),
\[
 \dim_\C V_S=\sum_{t\in S}d_t .
\]
\end{lemma}
\begin{proof}[Outline]
Coefficient integrability from \ref{B8.02y} permits \ref{B8.03a}, and Schur orthogonality identifies the image by evaluation at the cyclic vector. The orthogonal projectors of \ref{B7.44} split it by label, and \ref{B8.08} identifies the good basis. The complete proof is in \S\ref{prf:B8.52a}.
\end{proof}

\begin{theorem}[Counting constants from the label condition]\label{B8.52}
Assume \ref{B8.13} and the label row of \ref{B8.08} for \(K_P,Q,Q',L_k\), enlargement \(R\), and condition \(Q'\wedge R\). There are a constant \(c_{\rm Ic}>0\) and, for each finite real list \(\ell\), a tuple \(c_{30}(\ell)=(C_{\rm res},n_0,A,C_b,C,c,c_b)\). At \(\ell=\ell_{\rm out}\) its entries satisfy
\[
 0\le C,\qquad 2(C_{YL}+1)\le A.
 \tag{8.52.1}\label{eq:p8c5-count-outer}
\]
On \(\mathcal D_3\), with the quantifiers of Definition~\ref{B8.39}, these constants give the estimate of Lemma~\ref{B8.53} on the specified framed Hermitian realization whenever \(\max_v\min(p_v,q_v)\ge1\). Here the frames identify the base lattice and its constant signature. The real measures have the compact-dual normalization, quotient volumes agree in equal rank, and in definite signature the finite class decomposition has sum of reciprocal stabilizer orders equal to the standard mass. The level-character idempotents in this assertion have the form \(e=\mathbf1_H\chi^{-1}/\mu_f(H)\), where \(H\) is a compact open subgroup of the lattice stabilizer with unipotent reduction at \(g_\circ\) and \(\chi\) is a unitary character of \(H\).

At each sufficiently large point there is an arithmetic row system \(A\) on the selected base data with a common nontrivial additive character at each distinguished prime. It satisfies the join conditions for the selected real parameters and for the split level prime and finite projectors. The selected indices inject into the full label index set of \ref{B8.52a}, with image exactly the good indices. For any subset \(S'\) of label pairs, \(A(S')\) is the projected intertwiner space of that realization. Its dimension is the number of full indices with label in \(S'\). The total finite idempotent \(e'\) has \(D_A^\star(e')>0\). For \(k_0=0,1\) and a subset \(S\) of types at \(k_0\), put \(\operatorname{cyl}_{k_0}S=\{t:t_{k_0}\in S\}\), with idempotent \(e_{\operatorname{cyl}_{k_0}S}\). Its main term is nonnegative and
\[
 \left|\dim A(\operatorname{cyl}_{k_0}S)-D_A^\star(e_{\operatorname{cyl}_{k_0}S})\right|
 \le e^{-c_{\rm Ic}rf\log r}D_A^\star(e_{\operatorname{cyl}_{k_0}S}).
 \tag{8.52.2}\label{eq:p8c5-cylinder-count}
\]
\end{theorem}

\begin{proof}[Outline]
Use \ref{B8.02} and \ref{B8.10} for the gap test of \ref{B8.02y} and the size bounds, combining thresholds by \ref{B8.39}. Form cylinder projectors using \ref{B7.44} on the labelled space of \ref{B8.52a}, lift them, and use \ref{B8.53} for the constants and counts. The complete proof is in \S\ref{prf:B8.52}.
\end{proof}

Let \(I\) be a rank-\(n\) Hermitian datum with self-dual lattice over \(E/F\) unramified at finite places. Put \(\mathrm{rd}=|\operatorname{disc}E|^{1/(2f)}\). Write \(U\) for its unitary group, \(K_f\) for the finite lattice stabilizer and \(K(g)\) for its congruence kernel modulo \(g\). Use product Haar measure on \(U(\mathbb A_F)\) with finite factor \(\mu_f\), and let \(\mathscr F\) be a fundamental domain for \(U(F)\), of volume \(V\). Write \(L^2=L^2(U(F)\backslash U(\mathbb A_F))\). At each real place take a discrete series \(\sigma_v\) with parameter \(\lambda_v\) and unit lowest-type vector \(w_v\). Let \(\sigma=\widehat\bigotimes_v\sigma_v\) on \(H_\sigma\), \(w=\bigotimes_vw_v\), \(d_\infty=\prod_v d(\sigma_v)\) in the prescribed measures, and \(G_\infty=\prod_{v\mid\infty}U(F_v)\). For right convolution \(R_e\) by \(e\), put
\[
 D_{\rm count}=\dim_\C\bigl(R_e\circ\operatorname{Hom}^{\rm cont}_{G_\infty}(H_\sigma,L^2)\bigr),
 \qquad D^\star=d_\infty V\Re e(1).
 \tag{8.53.1}\label{eq:p8c5-count-main}
\]

\begin{lemma}[Counting idempotents]\label{B8.53}
Assume the finite-volume and nonsingular Gram-density assertions of \ref{B8.13}. For each \(A\in\R\) there are \(n_0\in\N,\ C_0\ge0\) such that for any additional level predicate \(P\), \(C\ge C_0\), and \(c=c_b=1/16\), the following holds. Take
\[
 \begin{gathered}
 n\ge n_0,\quad b\ge2,\quad\log\log\mathrm{rd}\le n,\quad
 |\lambda_{v,a}|\le b^2,\quad(p_v,q_v)=(p,q),\quad q^2\le n,\\
 g\in\N,\quad1\le g\le b^A,\qquad 4n<g_\circ,\quad g_\circ\text{ prime}.
 \end{gathered}
 \tag{8.53.2}\label{eq:p8c5-count-size}
\]
Require \(g_\circ\) to split completely and unramified in \(E\). The locally constant, self-adjoint convolution idempotent \(e\) is of positive type, supported in \(K_f\), bi-invariant under \(K(g)\), and has unipotent reduction on its support modulo \(g_\circ\). Assume \(P(I,\mu_f,g_\circ,e)\). Each \(\lambda_v\) decreases strictly in \((n-1)/2+\Z\), with absolute opposite-sign differences at least \(h\), where
\[
 Cn^{19}f^5\log(b\,\mathrm{rd})\le h,\qquad
 \log f\le\sqrt n,\qquad \log\log b\le c_bnf\log n .
 \tag{8.53.3}\label{eq:p8c5-count-gap}
\]
Then
\[
 |D_{\rm count}-D^\star|\le e^{-cnf\log n}D^\star .
 \tag{8.53.4}\label{eq:p8c5-count-bound}
\]
\end{lemma}

For \(x\in U(\mathbb A_F)\), \(\Lambda_x\) is the lattice cut out by the finite translate \(x_f\), with Euclidean norms in the real signature frames at \(x_\infty\). The size of a nonzero rank-\(a\) lattice \(B\), with induced or quotient metrics, is
\[
 \operatorname{size}(B)=
 \left(\frac{2^{fa}\operatorname{covol}(B)}{|\operatorname{disc}E|^{a/2}}\right)^{1/(2fa)} .
 \tag{8.53.19}\label{eq:p8c5-size}
\]
Write \(\operatorname{least}(x)\) for the infimum of nonzero sublattice sizes and \(\tau_1(x)=\min(1/2,\operatorname{least}(x))\). The largest size of a lattice is the supremum of its nonzero saturated quotient sizes. In the Harder--Narasimhan filtration \(G_{\le t}\) is the sum of sublattices of largest size at most \(t\). For \(y\in U(F_v)\) let \(d_v(y)=(\sum_a(\log s_a(y))^2)^{1/2}\) in terms of singular values \(s_a(y)\). Set
\[
 H=nf\log n,\quad N=n^2f,\quad X=nf\,\mathrm{rd}\,gb,\quad
 \eta=(10^6n^8f^2)^{-1},\quad \tau_*=X^{-11q}.
\]
For \(q\ge1\), call \(\operatorname{least}(x)\ge\tau_*\) the moderate region and its complement the deep region.

\begin{proof}[Outline]
Unfold the projector trace, with only the identity surviving in definite signature by the congruence. In indefinite signature use \ref{B8.53a} for decay and packing, \ref{B8.53b} for flag densities and deep volume, and \ref{B8.53e} for the deep trace. The complete proof is in \S\ref{prf:B8.53}.
\end{proof}

\begin{lemma}[Uniqueness of the lowest-type germ]\label{B8.53a-jet}
Let \(S\) be the lowest type of a unitary discrete series of \(U(p,q)\), \(n=p+q\), with compact group \(K=U(p)\times U(q)\). Let the decreasing parameter \(\lambda\) have opposite-sign gap at least \(g>2n\). Let \(\mathfrak p_\C\) be the span of opposite-sign matrix units and \(\Lambda_S\) the highest weight of \(S\). Let \(V_-\subset\mathfrak p_\C\otimes S\) be the sum of isotypic components of highest weight \(\Lambda_S+e_a-e_b\) for \(b<a\) of opposite signs. The homogeneous Schmid equation kills the first derivative contracted against \(V_-\). Two analytic \(\operatorname{End}(S)\)-valued germs at \(1\), with lowest-type left and right covariance, agree if they have the same value and both satisfy this equation.
\end{lemma}

\begin{proof}[Outline]
Induct on noncompact derivatives, using covariance and commutators to make the first possible nonzero derivative symmetric. The differentiated Schmid equation removes the lowering components, and the compact Casimir gap forces that derivative to vanish. The complete proof is in \S\ref{prf:B8.53a-jet}.
\end{proof}

\begin{lemma}[Local decay and arithmetic balls]\label{B8.53a}
At a real place of \ref{B8.53} with \(q=\min(p_v,q_v)\ge1\) and \(h\ge120nq^{3/2}\), the lowest-type compression of \(\sigma_v(y)\) has norm at most \(e^{-(h/8)\min(d^2,d/\sqrt{2q})}\), \(d=d_v(y)\). For \(n\ge1\), differences of distinct rational operators preserving \(\Lambda_x\) have geometric mean of real operator norms at least \(\tau_1(x)^2\). The arithmetic ball in \(x^{-1}U(F)x\) with finite component in \(K_f\) and radii \(d_v\le R_v\), for \(R_v\ge0\), satisfies \eqref{eq:p8c5-ball}.
\end{lemma}

\begin{proof}[Outline]
Use \ref{B8.53a-jet} to identify the compression with a compact average, then bound its minor weights by Cauchy–Binet. For the arithmetic ball, determinants give separation and rescaling the real radii gives the packing bound. The complete proof is in \S\ref{prf:B8.53a}.
\end{proof}

\begin{lemma}[Two volume estimates]\label{B8.53b}
Use the quotient of \ref{B8.53}, put \(r=\mathrm{rd}\), and let \(c_v\mu_v\) be Haar measure with Hilbert--Schmidt density one at the identity. Write \(P_0=\prod_v c_v\), \(k_R=\mu_f(K(g))\). Assume \ref{B8.13}, \(n=p+q\ge8,\ q^2\le n,\ Z\ge2,\ r\le Z,\ 2fr\le Z^2,\ g\ge1\). Then
\[
 P_0V\le Z^N\mu_f(K_f)\le Z^Ng^{2N}k_R.
 \tag{8.53.24}\label{eq:p8c5-absolute-volume}
\]
If also \(\log n\ge15,\ \log r\le e^n,\ \log f\le\sqrt n\), the probability measure of filtration flags of type \((j_1,\ldots,j_a)\), \(j_i\ge1,\ 1\le J=\sum j_i\le q\), is bounded by
\[
 \kappa\prod_\lambda\lambda^{m_\lambda}\prod_i 2fj_i\,\frac{dt_i}{t_i}
 \qquad(0<t_1<\cdots<t_a<1/2),
 \tag{8.53.25}\label{eq:p8c5-flag-measure}
\]
with \eqref{eq:p8c5-layer-scales} and \(\kappa\le e^{-JH/2}\). In particular \(\Pr(\operatorname{least}<1/2)\le e^{-H/3}\).
\end{lemma}

\begin{proof}[Outline]
Iterate the unit-hyperplane calculation of \ref{B8.49} and bound the rank-two class count by Minkowski's theorem. Unfold isotropic flags and bound the resulting Levi volumes, zeta factors and compact volume ratios to estimate the flag density. The complete proof is in \S\ref{prf:B8.53b}.
\end{proof}

\begin{lemma}[The deep-region exponent]\label{B8.53c}
Let \(n\ge8,\ f,q\ge1,\ q^2\le n\) be integers and \(r,g\ge1,\ b\ge2\). Put \(X=nfrgb,\ N=n^2f,\ k=8qN\). For \(0<\tau\le X^{-11q}\), define
\[
 \rho=\frac{n^2\sqrt f\,gr}{4}(2\tau)^{1/q},\quad
 B_j=\sqrt n\,b^2+\tfrac32n^{3/2}+3j,\quad B=B_{k+N},\qquad
 U=(2\rho\sqrt f B)^{2k}(6B)^Ne^f(17\sqrt n/\tau^2)^{2N}.
\]
Then \(\rho\le\log2\) and \(Ug^{2N}X^N\le e^{-N}\).
\end{lemma}

\begin{proof}
Put \(Y=X^q\). Since \(X\ge16\), \((2\tau)^{1/q}X^{11}\le2\) and \(n^2\sqrt f\,gr\le X^2\), one has \(\rho\le1/2\). The rank bounds give \(B\le30X^3\) and \(n^2fgrB/2\le15X^5\le X^6\). Thus \((2\rho\sqrt f B)^{16q}\le Y^{96}(2\tau)^{16}\). With \(Bng^2X\le Y^{10}\) and \(\tau Y^{11}\le1\),
\[
 (Ug^{2N}X^N)^{1/N}\le2^{16}\cdot6\cdot289\,e\,Y^{96}\tau^{12}Bng^2X
 \le2^{16}\cdot6\cdot289\,e\,Y^{-26}\le e^{-1}.
\]
\end{proof}

\begin{lemma}[A finite family of cuspidal isometries]\label{B8.53d}
Use the data of \ref{B8.53} with \(n_0\ge8,\ C\ge\max(0,8\cdot10^{14}(A+4))\) and \(q\ge1\). Let \(T_i:H_\sigma\to L^2\), \(1\le i\le D\), be isometric intertwiners with cuspidal images. Let \(V_j\) be the span of derivatives through order \(j\) of lowest-type tensors, taken one place at a time. Suppose \(T_i u\) has continuous right \(K(g)\)-invariant representative \(F_i(u)\) on these spaces, with right derivative \(F_i(d\sigma(Z)u)\) in direction \(Z\). Put \(F(u,y)=(F_i(u)(y))_i\). At \(\tau=\operatorname{least}(x)<X^{-11q}\), use \ref{B8.53c} with \(r=\mathrm{rd}\). Then
\[
 \|F(w,x)\|_2^2\le UDP_0/k_R .
 \tag{8.53.31}\label{eq:p8c5-deep-family}
\]
\end{lemma}

\begin{proof}[Outline]
Casimir derivative bounds and local cube integration with the packing estimate of \ref{B8.53a} bound evaluation. A gap in the small flag gives a short cover of a unipotent orbit, where vanishing constant terms allow repeated integration of derivatives. The complete proof is in \S\ref{prf:B8.53d}.
\end{proof}

\begin{corollary}[Mass of the counting projector]\label{B8.53e}
Under the indefinite hypotheses of \ref{B8.53}, the counted space is finite dimensional. Its projector has continuous kernel \eqref{eq:p8c5-kernel} and diagonal \eqref{eq:p8c5-trace}, and
\[
 \int_{\rm deep}K(x,x)\,dx\le e^{-N}D_{\rm count}.
 \tag{8.53.33}\label{eq:p8c5-deep-integral}
\]
\end{corollary}

\begin{proof}[Outline]
Smooth averaging and coefficient decay give continuous cuspidal vectors. Use \ref{B8.53d}, \ref{B8.53c} and the volume bound for the deep estimate, then use Bessel's inequality on the moderate region for finiteness and unfold to identify the kernel. The complete proof is in \S\ref{prf:B8.53e}.
\end{proof}

\section{Exact class counts}

For the upper unitriangular subgroup \(U_k(\F_g)\) put
\[
 I_k(g)=[\GL_k(\F_g):U_k(\F_g)],\qquad
 I_{01}=\prod_{\nu=0}^1 I_k(g_\nu)^f.
 \tag{8.54.1}
\]
A rational prime \(p\) splits completely in \(F\) precisely when \(F\) has \(f=[F:\Q]\) primes above \(p\), each with residue field \(\F_p\).

\begin{theorem}[The class number at constructed points]\label{B8.54}
Use the quantifiers of \ref{B8.39} on \(W_1(\mathcal A)\) for every member condition and bundle association. At each constructed point \(x>x_1\), \(i\in I_x\), suppose both \(g_0\) and \(g_1\) split completely in \(F\) and
\[
 \exists\nu\in\{0,1\},\qquad 2fk<g_\nu-1 .
 \tag{8.54.2}\label{eq:p8c5-stabilizer-prime}
\]
For every word \(u\), \(T_u\equiv1\bmod Q_Sqd\) and \(T_u\) splits completely in \(E\) and \(F\). Write \(\mathcal C_{u,0}\) for the initial column class set of the word \(u\). Its cardinality is
\[
 \#\mathcal C_{u,0}=M_k I_{01}(T_u^k-1)^f.
 \tag{8.54.3}
\]
\end{theorem}

\begin{proof}[Outline]
Use \ref{B8.51a} for splitting and compute the level index from the unitriangular and vector stabilizers. The cyclotomic degree bound forces the genus automorphism actions on level cosets to be free, so their orbit count is the index times the mass of \ref{B8.50}. The complete proof is in \S\ref{prf:B8.54}.
\end{proof}

\section{Bounded real rows and reciprocity data}

\begin{lemma}[Descent of a complex spanning relation]\label{B8.55a}
Let \(K\) be a field, \(\sigma:K\hookrightarrow\C\) an embedding, \(M\) a \(K\)-vector space, \(X\) a complex vector space, \(C\subseteq M\) a \(K\)-subspace, and \(L:M\to X\) a \(\sigma\)-semilinear map. Suppose that for every finite \(K\)-linearly independent family \(b_1,\ldots,b_m\) in \(C\), the family \(L(b_1),\ldots,L(b_m)\) is linearly independent over \(\C\). For \(S\subseteq C\) and \(f\in C\),
\[
 L(f)\in\operatorname{span}_{\C}L(S)
 \quad\Longrightarrow\quad
 f\in\operatorname{span}_K S.
\]
\end{lemma}

\begin{proof}
Choose a finite subset \(S_0\subseteq S\) whose images span \(L(f)\), and choose a \(K\)-basis \(b_1,\ldots,b_m\) of \(\operatorname{span}_K S_0\) from \(S_0\). Semilinearity gives
\(L(f)\in\operatorname{span}_{\C}\{L(b_1),\ldots,L(b_m)\}\).
If \(f\notin\operatorname{span}_K S\), then \(f,b_1,\ldots,b_m\) is \(K\)-linearly independent. All its members lie in \(C\), so the hypothesis makes \(L(f),L(b_1),\ldots,L(b_m)\) linearly independent over \(\C\), a contradiction.
\end{proof}

\begin{lemma}[Independence of comparison values]\label{B8.55b}
Fix the moduli, uniformization and comparison data at a selected layer of Proposition~\ref{B8.55}. Let \(K\subseteq\C\) be its second comparison field and \(C_K\) its space of algebraic cusp sections of the specified weight and degree. For an algebraic section \(s\), let \(R(s)(k,x)\) be its value in the complex frame \(x\) of component \(k\), extended by zero outside the frame domain. If the comparison assigns the component and frame \((\alpha(g),x_g)\) to an adelic point \(g\), put
\[
 E(s)(g)=R(s)(\alpha(g),x_g).
\]
For every finite \(K\)-linearly independent family \(s_1,\ldots,s_m\) in \(C_K\), the functions \(E(s_1),\ldots,E(s_m)\) are linearly independent over \(\C\).
\end{lemma}

\begin{proof}[Outline]
Applying \(K\)-linear functionals \(\lambda:\C\to K\) to \(\sum z_iR(s_i)=0\) and descending from complex points gives \(\sum\lambda(z_i)s_i=0\), proving \(\C\)-independence of \(R(s_i)\). Equivariance extends vanishing from comparison frames to all frames. The complete proof is in \S\ref{prf:B8.55b}.
\end{proof}

\begin{proposition}[Bounded rows of the indefinite realization]\label{B8.55}
Assume the following four inputs of the real theta construction.
\begin{enumerate}
\item For \(p,q\ge0\) with \(p+q\ge1\), use the determinant double cover
\[
 \widetilde U(p,q)=\{(g,z):g\in U(p,q),\ z\in\C^\times,
                         \ z^2=\det g\},
\]
For every realization of this cover as a topological group, with a chosen Haar measure and compatible exponential map, use the following assertion. Put \(I=\{1,\ldots,p\}\sqcup\{1,\ldots,q\}\), and let \(\prec_c\) be the compact-positive relation of the parameter convention. An allowed parameter is an injective map \(\lambda:I\to\Q\) such that \(a\prec_c b\) implies \(\lambda_b<\lambda_a\), with \(2\lambda_a\in\Z\) and \(\lambda_a-\lambda_b\in\Z\) for all \(a,b\).

There is a family \(\Pi_\lambda\) of nonzero irreducible unitary discrete-series representations indexed by all allowed parameters, and every such discrete-series representation is unitarily equivalent to some \(\Pi_\lambda\). Here a discrete-series representation includes an injective complex-linear realization of its Lie-algebra module whose image is exactly the compact-finite vectors and which intertwines the derivatives along the specified exponential map. It also includes a positive formal degree \(d_\Pi\); all squared matrix coefficients below are integrable and satisfy
\[
 \int_{\widetilde U(p,q)}
       |\langle y,\Pi(g)x\rangle|^2\,dg
       =\frac{\|x\|^2\|y\|^2}{d_\Pi}
 \qquad(x,y\in H_\Pi).
\]
Use the Blattner weight \(b(\lambda)\), the highest-weight cone \(\mathcal C_\lambda\), the rational parameter scalar \(q_{p,q}(\lambda)\), and the rational rank scalar \(r_{p+q}\) of the real theta parameter convention. The \(b(\lambda)\)-highest-weight space of \(\Pi_\lambda\)'s compact-finite module has dimension one, its Casimir acts by \(q_{p,q}(\lambda)-r_{p+q}\), and every highest weight \(\mu\) of that module belongs to \(\mathcal C_\lambda\).

\item Use the Harish--Chandra module realization, admissibility and irreducibility conventions of the real theta construction. For every such cover and irreducible unitary representation, all compact-finite vectors are smooth, there is a Harish--Chandra module realization, and every such realization is admissible and irreducible. Two irreducible unitary representations with isomorphic realized complex Lie-algebra modules are unitarily equivalent. Conversely, a unitary representation with an admissible irreducible Harish--Chandra module realization is irreducible.

There is also the following analytic assertion for the matrix unitary groups. For every \(n\ge0\), sign vector \(\epsilon\in\{\pm1\}^n\), and nonzero complex Hilbert space carrying a continuous irreducible unitary representation \(\sigma\) of
\[
 G_\epsilon=\{g\in\GL_n(\C):g^*J_\epsilon g=J_\epsilon\},
 \qquad J_\epsilon=\operatorname{diag}(\epsilon),
\]
the orbit map \(g\mapsto\sigma(g)v\) is real analytic for every vector finite under the standard maximal compact subgroup.

\item For every arithmetic datum and Hermitian space \(V\), choose the compatible infinite embeddings and an exponential family agreeing with matrix exponentiation at each infinite place. Let \(H_V\) and \(\overline H_V\) denote the adelic height and quotient height in this construction. Suppose \(\varphi\) is a continuous function on \(U(V)(\mathbb A_F)\), invariant on the left by \(U(V)(F)\), and smooth in the infinite directions. Require an open subgroup \(U\) of the full adelic group such that \(\varphi(gk)=\varphi(g)\) whenever \(k\in U\) belongs to the finite-adelic subgroup. Suppose that the integral of \(\varphi\) along every rational unipotent radical defined by an isotropic frame vanishes, for every translate and every compatible Haar measure and fundamental domain. For a word \(X=(X_1,\ldots,X_s)\) in the infinite Lie algebra, let \(D_X\) be its prescribed word derivative. Assume uniform moderate growth in the precise form
\[
 \exists N\in\N\quad\forall X\quad\exists C_X\in\R
 \quad\forall g,\qquad
       |D_X\varphi(g)|\le C_X H_V(g)^N.
\]
Then
\[
 \forall A\in\N\quad\exists C_A\in\R\quad\forall g,
 \qquad |\varphi(g)|\,\overline H_V(g)^A\le C_A.
\]
The rapid-decay reference used here is Miller and Schmid, \emph{On the rapid decay of cuspidal automorphic forms} \cite{miller-schmid-2012}.

\item Let \(G_\epsilon\) be as above, with Lie algebra \(\mathfrak g_\epsilon\) and standard maximal compact subgroup \(K_\epsilon\). Let \(\sigma\) be a continuous unitary representation on a complex Hilbert space \(H\) whose closed invariant subspaces are \(0\) and \(H\). Write \(H_{K_\epsilon}\) and \(H^\infty\) for its compact-finite and smooth vectors. Let \(E\) be a complete Hausdorff complex vector space whose topology is generated by a countable family of seminorms \((p_i)\). Suppose \(\pi(g)\) and \(d\pi(X)\) are continuous complex-linear operators on \(E\), with \(\pi(1)=1\), \(\pi(gh)=\pi(g)\pi(h)\), and continuous orbit maps. Require \(d\pi(X)v\) to be the differentiated action along the prescribed exponential in direction \(X\), for every \(X\in\mathfrak g_\epsilon\) and \(v\in E\).
Require the growth condition
\[
 \forall i\quad\exists q_i,N_i\quad\forall g\in G_\epsilon,
 \ v\in E,\qquad
 p_i(\pi(g)v)\le
 \left(1+\sum_{a,b}|g_{ab}|^2\right)^{N_i}q_i(v),
\]
where \(q_i\) is a continuous seminorm and \(N_i\in\N\). For a finite list \(\mathcal S\) of Lie-algebra words, put
\[
 Q_{\mathcal S}(w)=
 \sum_{(X_1,\ldots,X_s)\in\mathcal S}
       \|d\sigma(X_1)\cdots d\sigma(X_s)w\|,
\]
with the empty word acting as the identity.

Let \(T_0:H\to E\) be complex-linear and intertwine the \(K_\epsilon\)-action and the derived \(\mathfrak g_\epsilon\)-action on \(H_{K_\epsilon}\). There is a complex-linear map \(T:H\to E\) agreeing with \(T_0\) on \(H_{K_\epsilon}\) and intertwining \(G_\epsilon\) on \(H^\infty\), such that
\[
 \forall i\quad\exists\mathcal S,C\quad
 \forall w\in H^\infty,\qquad
       p_i(Tw)\le C Q_{\mathcal S}(w).
\]
Every other map with these properties agrees with \(T\) on \(H^\infty\), and \(T(H^\infty)\) is closed in \(E\). If \(T_0v\ne0\) for some \(v\in H_{K_\epsilon}\), then for every finite list \(\mathcal S\) there is a continuous seminorm \(q_{\mathcal S}\) on \(E\) such that
\[
 Q_{\mathcal S}(w)\le q_{\mathcal S}(Tw)
       \qquad(w\in H^\infty).
\]
The globalization reference used here is Wallach, \emph{Real Reductive Groups II}, Chapter~11 \cite[Ch.~11]{wallach-1992}.
\end{enumerate}
For every member condition and bundle association the following hold at each selected point.

For every \(j<m\) there are smooth unit extremal vectors \(u_\rho\) of the specified weights in the lowest types and a unit vector \(w\) in the paired space such that the Fock-model pairing of \(\bigotimes_\rho u_\rho\) with \(w\) is the prescribed product of minors \(\Delta_j\). Let \(K_{\rm c}\) be the coefficient field of the selected point, with its fixed embedding \(\iota:K_{\rm c}\hookrightarrow\C\). At the reference tuple \(t=\mu_0\), the comparison uses the field \(\iota(K_{\rm c})\). The component coefficients \(c_j\) and phase \(\phi_j\) have modulus one. There are a central level element \(z\) and a root of unity \(\zeta_t\in\iota(K_{\rm c})\). For each allowed input \(x\), its cuspidal theta class has, for every admissible representative \(f\) and every \(h\in G_{\Q}(W)\backslash G_{\mathbb A}(W)\), where \(W\) is the layer-\(j\) target Hermitian space of the moduli datum on which the theta function is evaluated, the identity
\[
 \phi_j\Theta_{j,t}(x)(h)
 =\zeta_t c_j(\alpha(h))\,a_j\,J_j(z\cdot f)(h).
\]
Here \(\Theta_{j,t}(x)\) is the adelic theta function of the input, \(\alpha(h)\) is the component of its comparison frame at \(h\), and \(J_j(f)(h)\) is the adelic evaluation of the section of \(f\) in that frame before multiplication by \(c_j\). If \(A_j\) is the sum of the section's weight exponents over all embeddings and frame coordinates, then \(a_j=(2\pi\xi)^{-A_j/2}\), where \(\xi=(\sqrt d\,\mathrm N\mathfrak c_K)^{-1}\) and \(\mathfrak c_K\) is the host curve ideal. Thus the component-normalized section is \(\varphi_j(c_j,f)(h)=c_j(\alpha(h))J_j(f)(h)\). Every layer lift \((g,k_\ell)\) preserves the cusp space and satisfies \(\varphi_j(c_j,gf)(h)=\chi_{\det}(g)\varphi_j(c_j,f)(hk_\ell)\), where \(\chi_{\det}\) is the product of the designated determinant characters. For each reference input there are finitely many \(a_\nu\in K_{\rm c}\) and algebraic cusp forms \(f_\nu\) of the specified weight such that
\[
 \phi_j\mathcal T_j(x)(h)=a_j\sum_\nu\iota(a_\nu)\,\varphi_j(c_j,f_\nu)(h)
 \tag{8.55.1}\label{eq:p8c5-theta-extension}
\]
for every \(h\), where \(\mathcal T_j\) uses the reference vector. At every weight and degree the base forms span over the comparison field. Scalar extension holds over each common coefficient algebra and the level action preserves coefficients.

For a tuple \(t\), write \(Q_t\) for its prescribed set of primes of \(E\). The prime compatibility condition requires that any two primes in the same \(Q_t\) over one rational prime have the same contraction to \(\mathcal O_{\Q(\sqrt{-d})}\), and that no prime in any \(Q_t\) divide \(Q_Sq\,\mathrm N(\mathfrak c_K)\). Under this condition one comparison at layer \(j\), over a field containing \(\iota(K_{\rm c})\), realizes the prescribed lattices, level actions and normalized theta class sums for all tuples. The theta identity holds for every tuple with the same central element \(z\) and phase \(\phi_j\), and with a root of unity \(\zeta_t\in\iota(K_{\rm c})\) depending on the tuple. Its component factor is the prescribed unit multiplier for sections at layer \(j+1\), taken at the embedding of the base field induced by the enlarged comparison field. The scalar-extension and level identities hold over the new field.
\end{proposition}

\begin{proof}[Proof (outline)]
The selected member has
\[
 n_V\ge8fr+4.
 \tag{8.55.2}\label{eq:p8c5-real-rank}
\]
The four stated inputs supply the smooth unitary realization, its lowest types and the required uniform rapid decay. Carry the unit-vector pairing to the Fock model. The pure-tensor map is injective on the paired data, so the image of \(\Delta_j\) is nonzero. With the other factors fixed, the map on the \(\rho\)-th lowest type is compact-equivariant and nonzero, hence injective. The product of minors has the prescribed torus character there, so its preimage \(u_\rho\) has that weight.

Fix a layer \(j\) and the reference tuple \(\mu_0\). Put \(n=n_V\) and let
\[
 K=\iota(K_{\rm c})\subseteq\C,
 \qquad \sigma:K\hookrightarrow\C
\]
be the inclusion. This is the second comparison field at the reference tuple. Let \(C_K\) be the \(K\)-space of algebraic cusp sections of the prescribed weight, and let \(\mathcal S\subseteq C_K\) consist of the normalized reference theta class sums belonging to allowed rational inputs. The normalized class-sum construction gives a member of \(C_K\) for each such input. We will prove \(C_K=\operatorname{span}_K\mathcal S\).

Use the readouts \(R\) and \(E\) of Lemma~\ref{B8.55b}. Choose the canonical covariant derivatives on the theta line. Let \(M=g_0g_1\) be the odd base level, let \(\mathfrak m\) be its chosen half-level ideal, and let \(L_W\) be the target lattice in the moduli datum. Choose an \(n\)-tuple \(x_1\) spanning \(L_W/\overline{\mathfrak m}L_W\) over \(\mathcal O_E\), whose point indicator gives the prescribed finite reference test. Write \(w\) for the finite Fourier transform of that indicator, using the level pairing between \(L_W/\overline{\mathfrak m}L_W\) and \(L_W/\mathfrak mL_W\). The normalized reference data also give an integer \(B>0\), divisible by \(2\) and \(M\), containing the bad primes. In the standard definite source frame the Gram matrix is \(1\) and the lattice is \(\Lambda_0=\mathcal O_E^n\).

The short-theta spanning theorem applies to this uniformization with this one fixed \(x_1\). Here are its hypotheses at the selected datum: the target rank is \(r\ge1\); the auxiliary level parameter is \(1\); the component Gram matrices and imaginary scalar agree; the arithmetic groups contain the principal congruence groups of level \(M\); the weight is dominant and \(n\) is even; and the corresponding arithmetic automorphy character is trivial. Moreover \([E:\Q]=2f\), so its numerical requirement is
\[
 4[E:\Q]r+4=8fr+4\le n.
\]
For \(s\in C_K\), its conclusion supplies finitely many algebraic sections \(s_\ell\) over \(K\) and \(b_\ell\in\C\) such that
\[
 R(s)=\sum_\ell b_\ell R(s_\ell).
\]
Each \(s_\ell\) is given on reduced, finitely generated \(K\)-algebras by a short theta value with finite weight \(w\), a self-dual lattice \(\Lambda_\ell\) equal to \(\Lambda_0\) at the primes dividing \(B\) and \(M\), an admissible odd level \(N_\ell\), and a polynomial of the prescribed harmonic highest weight. These are the algebraic short values formed by applying the canonical covariant derivatives to the finite sum of torsion translates of the theta section, with the lattice-index, finite-level and degree normalizations. The short-value identity is natural in the test algebra and the frame. The theorem initially permits any generating tuple of the conjugate half-level quotient. A rational unitary change of variables transfers those generators to \(x_1\); covariance of the short values expresses each resulting section in the \(K\)-span of the sections for \(x_1\). Applying the \(\sigma\)-semilinear readout gives the displayed complex spanning relation.

For each of these chosen short sections the split-place move gives a finite adelic unitary element \(h_\ell\), supported at split places away from \(B\), a new level \(N'_\ell\), a new harmonic polynomial \(P'_\ell\), and a scalar \(z_\ell\in K\). The new lattice is the complex conjugate of \(h_\ell\Lambda_0\); the lattice \(h_\ell\Lambda_0\) is self-dual, standard at the bad primes, and supported at split places. The new level is admissible for that lattice. At every framed point over a \(K\)-algebra, whenever the original short value is \(v\), there is a new short value \(v'\) with
\[
 v=\sigma_S(z_\ell)v',
\]
where \(\sigma_S:K\to S\) is the structure map. Indeed the rational unitary move is congruent to the identity modulo \(M\), so it preserves \(w\); pushing the polynomial along it preserves its harmonic weight and lattice conditions. Naturality of the short values descends the comparison scalar to \(K\). Applying this at each complex comparison frame gives a function \(t_\ell(g)\), the new short value there, such that
\[
 E(s_\ell)(g)=\sigma(z_\ell)t_\ell(g).
\]
This identity also covers \(z_\ell=0\).

Fix one of these terms and write \(h,P',t\) for its data. The polynomial coefficient map over \(K\) has image the prescribed harmonic highest-weight module, so \(P'\) is the image of a coefficient vector \(v\) over \(K\). Let \(\mathcal K_{h,v}(g)\) be the oscillator theta kernel with source point \(h\), coefficient vector \(v\), and the chosen lowest-type tensor \(\bigotimes_\rho u_\rho\). Let \(c(k)\) be the provisional unit component constants, and let \(\mu(g)\) be the unitary comparison multiplier, so that \(J_j(s)(g)=\mu(g)E(s)(g)\). Put
\[
 d_a=\sum_{\tau,l}a_{\tau,l},\qquad
 \ell=(2\pi\xi)^{-d_a/2},\qquad
 A(g)=c(\alpha(g))\ell\mu(g).
\]
Here \(a_{\tau,l}\) are the section's weight exponents. The Fock coefficient comparison supplies a unit complex phase \(\eta\), and the one-class theta comparison supplies a sign \(\delta\in\{1,-1\}\). Define
\[
 C_{\rm cmp}=\eta\delta^{d_a}(-1)^{d_a}
                 \bigl(\sqrt{2\pi\xi}\bigr)^{-d_a}.
\]
If \(\epsilon_h\in K\) is the normalized algebraic class scalar for \(h\Lambda_0\), the kernel comparison is the pointwise identity
\[
 \mathcal K_{h,v}(g)
 =C_{\rm cmp}\,c(\alpha(g))\mu(g)\sigma(\epsilon_h)t(g)
 =b_h A(g)t(g),\qquad
 b_h=C_{\rm cmp}\ell^{-1}\sigma(\epsilon_h).
\]
The class-scalar formula identifies \(\sigma(\epsilon_h)\), after multiplication by nonzero rational lattice indices, with a unit character value times a positive square root of an index. Hence \(\epsilon_h\ne0\). Also \(\ell>0\), \(|\eta|=|c(k)|=|\mu(g)|=1\), so \(b_h\ne0\) and \(A(g)\ne0\) for every \(g\).

The function \(\mathcal K_{h,v}\) belongs to the complex span of the theta lifts of the rational definite automorphic inputs at \(\mu_0\), with the same lowest-type tensor. To see the reduction, choose finitely many source double-coset representatives. The finite theta test transforms under the source level group by the inverse prescribed character. Interpolation on these finitely many classes produces a definite automorphic form whose theta lift is proportional to \(\mathcal K_{h,v}\), with nonzero proportionality constant. The rational class functions span the complex definite automorphic realization, and the theta lift is complex linear. This proves the asserted span.

For each rational input \(x\) in this span, the class presentation supplies an allowed class-coefficient function and hence \(s_x\in\mathcal S\). The reference class identity has the form
\[
 \Theta(x)(g)=u A(g)E(s_x)(g),\qquad |u|=1,
\]
with one \(u\) independent of \(x\) and \(g\). Therefore, if
\(\mathcal K_{h,v}=\sum_\nu q_\nu\Theta(x_\nu)\), \(q_\nu\in\C\), cancellation of the nonzero function \(A\) gives
\[
 t=\frac{u}{b_h}\sum_\nu q_\nu E(s_{x_\nu})
 \in\operatorname{span}_{\C}E(\mathcal S).
\]
Multiplying by \(\sigma(z_\ell)\) and using the raw spanning relation now gives
\(E(s)\in\operatorname{span}_{\C}E(\mathcal S)\).

All spans in the preceding analytic comparison are over \(\C\). For the descent, \(E\) is \(\sigma\)-semilinear on algebraic section rules, \(\mathcal S\subseteq C_K\), and Lemma~\ref{B8.55b} preserves every finite independent family from \(K\) to \(\C\). Lemma~\ref{B8.55a} consequently gives
\[
 s\in\operatorname{span}_{K}\mathcal S,
 \qquad K=\iota(K_{\rm c}).
\]
This proves \(C_K=\operatorname{span}_K\mathcal S\), the spanning assertion of the theorem. The expansion \eqref{eq:p8c5-theta-extension} is the reverse relation: the normalized class-sum construction above places each reference theta class in \(C_K\), and transporting its coefficients through \(\iota:K_{\rm c}\simeq K\) gives coefficients in \(K_{\rm c}\). The factor \(\ell=a_j\) is the factor in extremal-vector evaluation and in the core Petersson norm. The coefficient rules for the sections commute with scalar extension on common charts, and their level action contributes precisely the product of the designated determinant characters.

Finally the reference phase normalization replaces the provisional component constants and phase by the prescribed layer-\(j+1\) component multiplier and \(\phi_j\). The selected tuple characters already have the prescribed values at the used primes and the prescribed formulas away from their twist supports. Under the stated prime compatibility condition, the normalized comparison data can be constructed for all tuples over one field containing \(K\). The class identities hold for every normalized family, and the reference class sums span \(C_K\); thus the reference comparison extends to this family. The resulting comparison has one central level element and one phase for all tuples, the stated tuple roots of unity, and the same scalar-extension and level identities.
The complete proof is \Lean{GlobalTheory.\allowbreak{}Fc.\allowbreak{}R1.\allowbreak{}row\_\allowbreak{}hRB\_\allowbreak{}univ\_\allowbreak{}alt}.
\end{proof}

\begin{lemma}[The unramified CM model and its inert parameter]\label{B8.56f}
Let \(d>0\) be squarefree, \(d\equiv3\bmod4\), and let a complex period lattice be stable under \(\omega=(1+\sqrt{-d})/2\). For \(p\ge5\), \(p\nmid d\), fix a field isomorphism \(\C\simeq\C_p\). Put \(W=W(\overline{\F}_p)\), \(K^{\rm ur}=\operatorname{Frac}W\). In the homothety class of the given lattice there are short Weierstrass invariants in \(W\) of discriminant one. If \(p\) is inert in \(\Q(\sqrt{-d})\), and \(f\) is the logarithm of this good model with tangent one, there are \(\epsilon\in W^\times\) and mutually inverse integral parameters \(v,h\) with
\[
 f(v(T))=\epsilon\ell(T),\qquad
 \ell(h(T))=\epsilon^{-1}f(T),\qquad
 \ell(T)=\sum_{k\ge0}p^{-k}T^{p^{2k}} .
\]
\end{lemma}

\begin{proof}[Outline]
Use CM reciprocity and \ref{B8.42p} for unramifiedness and the auxiliary kernel choice, then use abscissa products for the valuation divisibilities. At an inert prime reduce the logarithm to type \(p-a\tau\) and apply Hazewinkel's lemma. The complete proof is in \S\ref{prf:B8.56f}.
\end{proof}

\begin{proposition}[Reciprocity supplies the twisting and kernel conditions]\label{B8.56}
Use the weak local correspondence of Theorem~\ref{B9.46} in the extension parameters of \((8.6)\). Take a bundle association for which each associated reciprocity datum is identified with the compatible datum of a constructed point. The following assertions hold for the good members of that construction.

Use the family quantifiers of Definition~\ref{B8.39} on \(\mathcal D_3\). Beyond one family threshold, at every \(x>x_1\) and every member \(i\in I_x\), choose levels and local lifts on the prescribed fillings for all words and twists. At the used twists and permitted places outside the twist primes these choices have the comparisons of Definition~\ref{B3.54k} and the following property:
\[
 \text{the local systems satisfy the stabilizer and complementary-orbit bounds of Theorem~\ref{B4.16} with initial saving one.}
 \tag{8.56.1}\label{eq:p8c5-twist}
\]
At every member beyond the selection threshold with \(x>x_1\), adjacent class scalars have norm one at every level-free \(w\), for every word and tuple. For each word \(u\) and used twist \(\mu\), the adjacent pair maps preserve rational row spaces at such places. On the reference chain and on each used twisted chain, the lift sends the complex span of the embedded rational class functions at layer \(j\) into the corresponding span at layer \(j+1\). Each of these chains has trivial initial stabilizers.

For the selected entries at this member there are coordinate scalars, orbit phases and signs \((e_w,\beta_w,\sigma_w)_w\). Write \(p\) for the residue characteristic of \(w\). A place is level-free for \((u,\mu)\) when \(p\ne T_u\), \(p\notin\{g_0,g_1\}\) and \(p\notin Q'(\mu)\). At each such place the following assertions hold.
\begin{itemize}
 \item If \(p\) is in the inert class with \(h_P\) of even order, or is in the signal class, every set of orbit representatives satisfies the letter conditions of Proposition~\ref{B4.01b} for these data. For every \(j<m\) with \(N_jN_{j+1}\) even, all class-pair entries have the all-place expansion of Proposition~\ref{B4.11a}.
 \item At an inert-class prime and every layer of even rank product, the entries have the one-position expansion when \(h_P=0\), and the all-real orbit expansion when \(h_P\ne0\) has odd order.
 \item If \(p\ne d\), the minor evaluations of Definition~\ref{B3.11} have integral coefficients at every class pair and every layer of even rank product.
\end{itemize}
At every level-free place in the discriminant class, each entry at \(j<m\) has the form
\[
 \Theta=\exp\left(-S\sum_t u_t\sum_{a,b}X_{(0,t),a,b}X_{(1,t),a,b}\right)\Theta_0,
 \quad |S|_p\le p^{1/2},\quad |u_t|_p\le1,\quad
 \left|\left(\prod_\alpha\nu_\alpha!\right)[X^\nu]\Theta_0\right|_p\le1
 \tag{8.56.2}\label{eq:p8c5-ramified}
\]
for every multi-index \(\nu\).

If \(5\le p\le X''\) is in the inert class, \(h_P\ne0\), and the orbit length is even, then for every word \(u\), every \(j<m\), every selected first-sheet orbit representative \((0,t)\), and every series with the prescribed expansion on that orbit, all coefficients of its minor evaluation at zero have norm at most \(p^{-\max(g_0(t,j+1),0)}\) for every column choice. The gain \(g_0\) and the minor evaluation are those of \ref{B3.48} and \ref{B4.10a}.
\end{proposition}

\begin{proof}[Outline]
Use reciprocity \ref{B1.15b}, the norm choices of \ref{B9.46} and the model of \ref{B8.56f}. Combine the chain of \ref{B8.56c} with the cyclic bound of \ref{B4.01b}, then use \ref{B8.56a}, \ref{B8.56b}, \ref{B8.56d} and the phase defects for expansions and jet bounds. The complete proof is in \S\ref{prf:B8.56}.
\end{proof}

\begin{lemma}[The pairing of the CM sigma series]\label{B8.56e}
Let \(K=\Q(\sqrt{-d})\) be the host quadratic field, \(\mathfrak c\) a nonzero curve ideal of norm prime to \(2d\), and \(\Lambda=\delta\mathfrak c\) of discriminant one. Let \(p\ge5\) be inert in \(K\) with \(p\nmid\mathrm N\mathfrak c\), and fix \(\lambda:\overline\Q\hookrightarrow\C_p\). Put \(W=W(\overline\F_p)\) and \(K_0=\operatorname{Frac}W\). Let \(\sigma\) be Witt Frobenius. Use an integral good Weierstrass model in the prescribed homothety class with logarithm \(l'(0)=1\), and write \(m_a=l^{-1}(al)\), \(t\oplus u=l^{-1}(l(t)+l(u))\).

For a period \(\omega\), let \(\eta_\Lambda(\omega)\) be the increment of Weierstrass zeta. For a positively oriented basis \(\omega_1,\omega_2\), put
\[
 s_2=\frac{\eta_\Lambda(\omega_1)\bar\omega_2-\eta_\Lambda(\omega_2)\bar\omega_1}
 {\omega_1\bar\omega_2-\omega_2\bar\omega_1},\qquad
 \Sigma_\Lambda(z)=\sigma_\Lambda(z)\exp(-s_2z^2/2).
\]
For model invariants \(e^4\lambda(g_2),e^6\lambda(g_3)\), with \(g_2,g_3\) those of \(\Lambda\), let \(\Sigma(z)=e^{-1}(\lambda_*\Sigma_\Lambda)(ez)\). Its coefficients are independent of the extension of \(\lambda\) used in the transfer. Write
\[
 P_E(t)=-\log(\Sigma(l(t))/t),\qquad B_E(t,u)=P_E(t\oplus u)-P_E(t)-P_E(u).
\]
Here \(A\sim B\) means \(A-B\) has zero constant term and \(\exp(A-B)\) is integral. Then \(P_E\in K_0[[t]]\), \(B_E\) is biadditive modulo \(\sim\), and \(m_a\) is integral for \(a\in W^{\sigma^2}\), with
\[
 P_E(m_a(t))\sim a\sigma(a)P_E(t),\qquad
 B_E(m_a(t),u)\sim B_E(t,m_{\sigma(a)}(u)).
\]
Let \(h\) be an integral parameter isomorphism with \(\ell(h(t))=e_0l(t)\), where \(\ell\) is as in \ref{B3.04} and \(e_0\in W^\times\). Put \(v=h^{-1}\), \(\epsilon=e_0^{-1}\), \(F(X,Y)=\ell^{-1}(\ell(X)+\ell(Y))\), and
\[
 P(T)=-\log\frac{\Sigma(\epsilon\ell(T))}{\epsilon T},\qquad
 b(X,Y)=B_E(l^{-1}(\epsilon X),l^{-1}(\epsilon Y)).
\]
Then \(b\in XYK_0[[X,Y]]\) satisfies \ref{B3.05}, and \(b^\ell=B_E(v(X),v(Y))\sim B_P\), where \(B_P(X,Y)=P(F(X,Y))-P(X)-P(Y)\). The norm and adjoint laws hold for \(P,B_P\) with \([a]=\ell^{-1}(a\ell)\).
\end{lemma}

\begin{proof}[Outline]
The sigma addition and isogeny formulas give biadditivity and the CM norm law modulo integral exponentials. Continuity and polarization give the scalar and adjoint laws, which transfer through the height-two parameter to the pairing of \ref{B3.05}. The complete proof is in \S\ref{prf:B8.56e}.
\end{proof}

\begin{lemma}[The ordinary gain on a standard orbit series]\label{B8.56d}
Assume \eqref{eq:p8c5-entry-shape} with \(d_{\ell1}\ge5,\ d_{\ell2}\ge4\). Let \(w\) have \(5\le p\le X''\), residue class zero, \(h_P\ne0\) and even orbit length \(d_0=n_P\ge2\). Take a word \(u\), a layer \(j<m\), and a selected representative \((0,t)\) of a first-sheet orbit at \(p\). On the orbit of \((0,t)\) use the standard letter series below with scalar \(|\xi|_p=1\) and sign \(-1\). For every \(\Theta\) with this expansion, each coefficient of the column-minor evaluation of \ref{B4.10a} has norm at most \(p^{-\max(0,g_0(t,j+1))}\).
\end{lemma}

\begin{proof}[Outline]
Undo the trace-dual change by integral substitution and remove the torus average modulo integral exponentials via \ref{B3.05}. Use \ref{B3.05b} for exponent zero and \ref{B3.26} for \(g_A\), then compare the caps and digit errors in \ref{B3.13} and \ref{B3.47}. The complete proof is in \S\ref{prf:B8.56d}.
\end{proof}

\begin{lemma}[The corrected Taylor series in an orthonormal basis]\label{B8.56a}
A complex period lattice stable under \((1+\sqrt{-d})/2\), \(d\) squarefree and \(d\equiv3\bmod4\), has for every \(p\ge5\) prime to \(d\) and isomorphism \(\C\simeq\C_p\) a homothetic short Weierstrass model over the completed unramified extension with integral coefficients and unit discriminant.

Take inert \(p\ge5\) in \(K=\Q(\sqrt{-d})\), a theta datum for a nonzero ideal \(\mathfrak c\) of norm prime to \(2d\), and \(\Lambda=\delta\mathfrak c\) of discriminant one. Its character \(\chi\) is multiplicative on norm-one ideals with \(\chi((z))=\prod_\tau\tau(z)^{-1}\) on unitary \(z\), for the prescribed CM type of \(E\). Fix \(\lambda:\overline\Q\hookrightarrow\C_p\). Let \(L,L'\) be self-dual of even ranks. Let \(F\) be a finite linear combination with coefficients in \(\mathcal O_{\C_p}\) of the transferred Taylor series of translates of the sections in Lemma~\ref{B3.35d}. For each summand the two lattices agree with \(L,L'\) after localization at \(p\), and some positive integer prime to \(p\) carries its shift into \(P(L,L')\).

Use any integral Weierstrass model \(W\) of discriminant one obtained from the transferred period equation by a Weierstrass change of variables. Write \(l_W\) for its logarithm of tangent one and \(e\) for the inverse of the leading scale of the change, so that \(e\,l_W\) is the period differential coordinate. Put \(A(T)=\sigma_\Lambda^*(e\,l_W(T))/(eT)\), with \(\sigma_\Lambda^*=\Sigma_\Lambda\) transferred by \(\lambda\). Let \(B_1,\ldots,B_r\) be a \(K\)-linear basis localized at \(p\) of \(\mathfrak a_0\operatorname{Hom}(L,L')\), where \(\mathfrak a_0=\mathfrak b^{-1}\) as in Lemma~\ref{B3.35d}. Thus the \(B_i\) form an \(\mathcal O_{K_p}\)-basis of the completed module. Require it to be orthonormal for \(\operatorname{Tr}_{E/K}\sum_{a,b}x_{ab}\overline{y_{ab}}\). Then
\[
 F\left(\left(e\sum_i\lambda(\tau(B_{i,ba}))l_W(T_i)\right)_{\tau,b,a}\right)
 \prod_i A(T_i)^{-1}\in\mathcal O_{\C_p}[[T_1,\ldots,T_r]] .
 \tag{8.56.5}\label{eq:p8c5-orthogonal-taylor}
\]
\end{lemma}

\begin{proof}[Outline]
Use \ref{B8.56f} and the sigma product for prime-to-\(p\) translates, with \ref{B8.56e} for unitary basis changes. The four-copy identity of \ref{B4.22c} gives integrality coefficientwise, and \ref{B3.35d} and \ref{B3.35a-arithmetic} handle the endpoint moves. The complete proof is in \S\ref{prf:B8.56a}.
\end{proof}

\begin{lemma}[The change to paired variables]\label{B8.56b}
Use the height-two phase \(P\) and pairing \(b\) of \ref{B8.56e}, or any phase with the same biadditivity, norm and adjoint laws and a pairing satisfying \ref{B3.05} with \(b^\ell\sim B_P\). Let \(f_e,f'_e\in(W^{\sigma^2})^r\) for a finite index set \(e\). For \(\langle u,v\rangle=\sum_i\sigma(u_i)v_i\), assume
\(\langle f_e,f_h\rangle=\langle f'_e,f'_h\rangle=0,\ \langle f_e,f'_h\rangle=\delta_{eh}\).
Put \(L_i=\sum_e(f_{ei}X_e+f'_{ei}Y_e)\), \(S=\sum_e b(X_e,Y_e)\). If \(H(T)=\Theta((\ell(T_i))_i)\exp(\sum_iP(T_i))\) is integral, then \(\Theta(L)\exp S\) is integral after logarithmic substitution. The phase \(S\) is even modulo the equivalence of \ref{B3.04}.
\end{lemma}

\begin{proof}[Outline]
Expand the phase of the formal sum using the norm, adjoint and biadditivity laws. The Gram relations leave only the matched cross-pairings, giving integrality after substitution, and adjointness for \(-1\) gives evenness. The complete proof is in \S\ref{prf:B8.56b}.
\end{proof}

\begin{proposition}[The kernel on the selected chain]\label{B8.56c}
Fix a constructed arithmetic column, a word \(u\), a used twist \(\mu\), and a coefficient place \(w\) of residue characteristic \(p\). Assume the permission condition of Theorem~\ref{B3.45}, \(p\notin Q'(\mu)\), and \(2r_G\le k\). Use the prescribed filling and the initial stabilizer condition of Theorem~\ref{B4.16}. The curve with its prescribed period lattice and embedding has the good model and parameter of Lemma~\ref{B8.56f}.

There are levels and a chain on that filling for which the arithmetic kernel satisfies Definition~\ref{B3.05a}. One embedding extending the coefficient place, one family of target readings and one family of normalized entry series give the class expansions at every layer.

On these same levels and chain there are semilinear comparisons on the scalar extensions of the rational lift spaces. They agree with the initial comparison, intertwine adjacent lifts and transfer terminal local valuation bounds to the terminal rational lattice with the same exponent. On the selected levels and chain the system has \(d_G(0)=0\) and \(d_G(j)=v_p(D_{\Gamma,u})\) for \(j>0\), with \(D_{\Gamma,u}\) the common stabilizer divisor.
\end{proposition}

\begin{proof}[Outline]
Use the good model \ref{B8.56f}, finite tests \ref{B3.35e} and paired variables \ref{B8.56b} to build the chart series. Transport contractions by \ref{B3.10a} and sum with stabilizer weights to intertwine lifts, then use \ref{B4.13a} to transfer the terminal bound. The complete proof is in \S\ref{prf:B8.56c}.
\end{proof}


% ---- p8c6 ----
\chapter{The chain length}\label{ch:p8c6}

Fix a scale $s=(x,\varepsilon,J)\in\R^3$ with $x>1$,
$0<\varepsilon<1/52$ and $J\ge1$. Put
\[
 L=x^\varepsilon,\qquad \delta_1=x^{-\varepsilon/8},\qquad
 \alpha=\sqrt{x\delta_1}=x^{1/2-\varepsilon/16},\qquad
 \beta=x^{9/2+10\varepsilon}.
\]

\phantomsection\label{B8.57}
For $y\in\R$ let $\lceil y\rceil_+=\max(0,\lceil y\rceil)$ and set
$e(y)=2\lceil y/2\rceil_+$. The prescribed chain length (the number of lifts) is $m_*(s)=e(\alpha)$.
Theorem~\ref{B8.59} uses the same rounding formula to choose $m$
and the starting rank $k$ on a prescribed frame.

For $y>0$ we have $y/2\le\lceil y/2\rceil_+<y/2+1$,
so $y\le e(y)<y+2$. Every even natural number at least $y$
is $2a$ with $a\ge\lceil y/2\rceil_+$. Thus $m_*(s)$ is the
least even natural number at least $\alpha$, with
$\alpha\le m_*(s)<\alpha+2$.

\phantomsection\label{B8.58}
Fix the family constants $C$. Write $c_W$ for the coefficient multiplying $r_*$ in the formula for $w_C$ below, before division by $8$, and $c_T$ for the exponent of $x$ in the formula for $t_C$. The family also prescribes a bound $c_o$ on the index of the cyclic subgroup generated by the distinguished element of the frame. For $n\in\N$ put $f=2n$ and define
\[
 \begin{aligned}
 k_*&=e(\beta),& r_*&=k_*+m_*-1,\\
 w_C&=8\left\lceil\frac{c_Wr_*}{8}
       \left(f+\frac{f^2(k_*-1)^3L^2}{m_*}\right)\right\rceil_+,
 &t_C&=\lceil x^{c_T}\rceil_+w_C+8fr_*+4-r_* .
 \end{aligned}
\]
Subtraction in $\N$ is truncated at zero.
For a frame $\varphi$ as in \ref{B2.01}, let $A$ be its finite abelian group and $a_0$ its distinguished element. Put $n=|A|$ and $o_\varphi=[A:\langle a_0\rangle]$. Rank data on $\varphi$ with starting rank $k$, chain length $m$ and auxiliary rank $t$ satisfy the family's rank selection condition at scale $s$ if and only if
\[
 x^{2+8\varepsilon}\le n\le2x^{2+8\varepsilon},\qquad
 k=k_*,\quad m=m_*,\quad t=t_C,\quad o_\varphi\le c_o.
\]
The prescribed triple $(k_*,m_*,t_C)$ is admissible for frames of order $n$ if and only if $t_C$ is odd and
\[
 \begin{gathered}
 2\le m_*,\quad m_*+3\le k_*,\quad r_*+5\le t_C,\quad
 3m_*<t_C,\\
 m_*\le f\le k_*,\qquad
 m_*^2\le k_*+m_*(t_C+1).
 \end{gathered}
\]
When the triple is admissible, adjoining it to any frame of order $n$ gives rank data on that frame. Both $k_*$ and $m_*$ are even by the definition of $e$.

\begin{theorem}[Ranks at a window point]\label{B8.59}
For each $0<\varepsilon<1/52$ one may take $x_0=16$.
For every scale $s$ with this $\varepsilon$ and $x\ge x_0$, and every
arithmetic frame $\varphi$ of order $n$ satisfying
$x^{2+8\varepsilon}\le n\le2x^{2+8\varepsilon}$, there are rank data $(k,m,t)$ on $\varphi$ such that
\[
 \alpha\le m\le2\alpha,\qquad
 \beta\le k\le2\beta,\qquad
 r+t=8fr+4+TW,
\]
where $f=2n$, $r=k+m-1$, and
\[
 T_O=T=\lceil x^{377}\rceil_+,\qquad
 n_W=W=8\left\lceil\frac{10r}{8}
       \left(f+\frac{f^2(k-1)^3L^2}{m}\right)\right\rceil_+.
\]

\end{theorem}
\begin{proof}[Outline]
Use \ref{B8.57} and the order window to bound \(m,k\), then set \(t=8fr+4+TW-r\) to obtain the identity for \(r+t\) and \(t\ge15r+4\). Even \(m,k\) make \(r,t\) odd, and the lower bound on \(t\) gives the remaining rank inequalities for \(\varphi\). The complete proof is in \S\ref{prf:B8.59}.
\end{proof}


\part{Galois representations and local Langlands in the archive}

% ---- p9c1 ----
\chapter{Galois representations attached to polarized cuspidal representations}\label{ch:galois-polarized}
\section{The global theorem}

Let $E/F$ be a CM extension with conjugation $c$. Fix an algebraic closure $\Omega$ of $E$, a prime $\ell$, and an isomorphism $\iota:\bar{\Q}_\ell\simeq\C$. Put $G_E=\Gal(\Omega/E)$. A polarized input consists of a cuspidal $\pi$ on $\GL_a(\mathbf A_E)$, $a\ge1$, with distinct infinitesimal weights in $(a-1)/2+\Z$, and a Hecke character $\chi$ of $F$ with $\chi_v(-1)=(-1)^a$ at each real place and
\[
 \pi^c\simeq(\pi\otimes(\chi\circ N_{E/F})^{-1})^\vee.
\]
For a finite place $w$ write $E_w$ for its completion, $k(w)$ for its residue field, and $q_w=\#k(w)$. For a Satake multiset $B$ at $w$ put $A_a(w,B)=q_w^{(a-1)/2}B$ and $P_{w,B}(X)=\prod_{\alpha\in A_a(w,B)}(X-\alpha)$. We use inverse arithmetic Frobenius in the Satake formulas and transport coefficients by $\iota$. A prolongation of $w$ to $\Omega$ specifies the restriction $r_w$ of a representation $r$ to the local Galois group. Write $W_w$, $I_w$ and $P_w$ for the Weil group, inertia and wild inertia. The degree on $W_w$ sends arithmetic Frobenius to $1$ and has kernel $I_w$. Galois representations are continuous and semisimple unless specified otherwise.

\begin{definition}[Global and local inputs]\label{B9.01}
Let $h$ denote Theorem~\ref{B9.02} and $L$ the local existence theorem \ref{B9.46}. The three inputs, in order, are $(h,L,h)$. Theorem~\ref{B6.04} uses the global theorem to prove the window bound of Definition~\ref{B6.03} with $C_t=401$ for every primitive Dirichlet character.
\end{definition}

\begin{theorem}[Global theorem]\label{B9.02}
For each polarized input there exist $r:G_E\to\GL_a(\bar{\Q}_\ell)$, an integer $w_{\rm gt}$ and integer multisets $H_\tau$ as in \ref{B9.04}. These data satisfy \ref{B9.05} at every $w\nmid\ell$ and \ref{B9.06} at every $w\mid\ell$, for every prolongation and every map of \ref{B9.47}.
\end{theorem}
\begin{proof}
Apply \ref{B9.03}.
\end{proof}

\begin{proposition}[Assembly]\label{B9.03}
Theorems \ref{B9.08} and \ref{B9.39} imply \ref{B9.02}.
\end{proposition}
\begin{proof}
The geometric seed gives \ref{B9.28}, and the approximations of \ref{B9.31} give the odd-rank limit of \ref{B9.35}. The stable-range theorem gives the two lifts of \ref{B9.41}. Propositions \ref{B9.42}--\ref{B9.45} recover the even-rank data by subtracting the added lines. The spherical theorem \ref{B9.38} gives weak attachment over CM fields. Apply \ref{B9.07}, with the tame places restricted as in \ref{B9.05}.
\end{proof}

\begin{definition}[Weights]\label{B9.04}
For an embedding $\tau:E\hookrightarrow\bar{\Q}_\ell$, form $H_\tau$ by replacing each infinitesimal weight $z$ of $\pi$ at $\iota\circ\tau$ by $(a-1)/2-z$, with its multiplicity. The weight $w_{\rm gt}$ satisfies $H_{\tau c}=w_{\rm gt}-H_\tau$.
\end{definition}

For a cuspidal $\Pi$ of rank $a\ge2$ over a CM field $E$, cohomological for an irreducible algebraic coefficient representation and satisfying $\Pi^c\simeq\Pi^\vee$, let $R_\ell(\Pi)$ be the representation of Shin and Chenevier--Harris \cite{shin-2011,chenevier-harris-2013}. Caraiani proves that at each finite place $w$,
\[
 \iota\,\mathrm{WD}(R_\ell(\Pi)|_{G_{E_w}})^{\mathrm{F\text{-}ss}}
 \simeq \mathcal L_{a,E_w}(\Pi_w),
\]
where the geometric normalization is
\[
 \mathcal L_{a,E_w}(\Pi_w)
 =\operatorname{rec}_{\rm cl}\!\left(\Pi_w^\vee\otimes|\det|_w^{(1-a)/2}\right),
\]
with $\operatorname{rec}_{\rm cl}$ the correspondence compatible with local $L$- and $\epsilon$-factors. The isomorphism includes monodromy at every finite place \cite[Thm.~1.1]{caraiani-2014,caraiani-2012}. Theorem~\ref{B9.02} allows general polarizations and uses the following local tests.

\begin{definition}[Away from $\ell$]\label{B9.05}
Write a Weil--Deligne parameter $z$ as $(\rho_z,N_z)$, with Weil action $\rho_z$ and monodromy $N_z$. For $w\nmid\ell$, a prolongation and a map $R$ of \ref{B9.47}, put $x=\mathrm{WD}_\iota(r_w)$ and $y=R(\pi_w\otimes|\det|_w^{(1-a)/2})$. Here $\mathrm{WD}_\iota$ denotes the Frobenius-semisimplified Weil--Deligne class with its monodromy, transported by $\iota$.

The toral test requires characteristic polynomial $(X-1)^a$ on every Weil commutator acting on $y$. The tame test requires this polynomial on $P_w$. For each lift $\gamma\in I_w$ of a topological generator of $I_w/P_w$, it also requires $z^{q_w-1}=1$ for every repeated root $z$ of $\det(X-\rho_y(\gamma))$. The full tame test adds a continuous isomorphism $E_w\simeq\Q_p$, where $p$ is the residue characteristic, and the condition $a=1$ or $p>16a^2$. The local comparison of $x$ with $y$ requires that they have the same rank. If $N_y=0$, it requires equality of classes under the toral test, and $N_x=0$ with equality of all inertia characteristic polynomials under the full tame test.
\end{definition}

\begin{definition}[Above $\ell$]\label{B9.06}
At $w\mid\ell$, take any prolongation. Use any period-ring data $P,A_{\inf}$ admitting the following identifications. The crystalline and de Rham rings of $P$ admit $G_{\Q_\ell}$-equivariant ring isomorphisms with $B_{\rm cris}$ and $B_{\rm dR}$ that preserve the images of $\Z_\ell$. The crystalline isomorphism respects Frobenius and the map from $A_{\inf}$. The de Rham isomorphism carries each filtration step for every finite extension of $\Q_\ell$ to the corresponding step of $B_{\rm dR}$. For $z\in\bar{\Q}_\ell$, the image of the coefficient $z$ lies in $B_{\rm dR}^+$ and reduces to $z$ under $\theta$. The distinguished crystalline period maps in $B_{\rm dR}$ to $\log([\varepsilon])$ for the chosen system $\varepsilon$ of $\ell$-power roots of unity. Under these identifications the map $B_{\rm cris}\to B_{\rm dR}$ agrees with the natural map on the uncompleted divided-power envelope and sends $A_{\rm cris}$ into $B_{\rm dR}^+$. In each finite filtration quotient of $B_{\rm dR}^+$, a single power of $\ell$ carries the image of $A_{\rm cris}$ into the image of $A_{\inf}$.
 Take any period bridge, that is, a comparison map inducing the specified local embedding. For a map $R$ of \ref{B9.47} put $y=R(\pi_w\otimes|\det|_w^{(1-a)/2})$. If $\pi_w$ is spherical, $r_w$ is crystalline and de Rham with jumps $H_\tau$ at every induced embedding $\tau$. If $N_y=0$ and either the toral or the full tame test of \ref{B9.05} holds, $r_w$ is potentially semistable and de Rham with those jumps. Its Frobenius-semisimplified period parameter $x$ satisfies the local comparison with $y$ in Definition~\ref{B9.05}.
\end{definition}

\begin{lemma}[Parallel infinity exponent and polarized weights]\label{B9.07a}
For every polarized input $(\pi,\chi)$ of rank $a\ge1$, there is an integer $b$ such that $\chi_v(x)=x^b$ for all $x\in\R_{>0}$ at every real place $v$ of $F$. With $w_{\rm gt}=a-1-b$, the integer multisets $H_\tau$ of Definition~\ref{B9.04} satisfy $H_{\tau c}=w_{\rm gt}-H_\tau$ for every $\tau:E\hookrightarrow\bar{\Q}_\ell$.
\end{lemma}
\begin{proof}[Outline]
Polarization pairs infinitesimal weights above \(v\) by \(z\mapsto b_v-z\), making \(b_v\in\Z\). Because \(\chi\) kills a finite-index subgroup of totally positive units, Dirichlet's unit theorem forces all \(b_v=b\), giving \(H_{\tau c}=a-1-b-H_\tau\) in \ref{B9.04}. The complete proof is in \S\ref{prf:B9.07a}.
\end{proof}

\begin{proposition}[Local clauses]\label{B9.07}
Assume \ref{B9.38} over CM fields. Every polarized input has common witnesses to \ref{B9.04}--\ref{B9.06}. In both local clauses the wild-inertia and repeated-root tests give the tame conclusion for every rank $a\ge1$ when $E_w\simeq\Q_p$ continuously and $N_y=0$. The toral clause applies over every $E_w$.
\end{proposition}
\begin{proof}[Outline]
Takes \(w_{\rm gt},H_\tau\) of \ref{B9.04} from \ref{B9.07a} and \(r\) from the spherical theorem. Uses genericity for the toral clause; for the tame clause over \(E_w\simeq\Q_p\), kills wild inertia and compares inertia polynomials and jumps after an unramified lift. The complete proof is in \S\ref{prf:B9.07}.
\end{proof}

\section{Geometric seeds in odd rank}

Let $a$ be odd, $x^2=-d$ in $E$ with $d\in\Z_{>0}$, $-d\in(\Q_\ell^\times)^2$, and $[F:\Q]$ even. Assume every nonsplit finite place is unramified over $F$ and lies away from $\ell$. Places are separated if their restrictions to $F$ are distinct. A split place over $p$ is usable if $-d\in(\Q_p^\times)^2$.

Take separated $w_0,v'\mid\ell$ inducing the same place on $\Q(x)$. Let $U$ be a definite-unitary automorphic representation of rank $a$, unramified at nonsplit places off a finite nonsplit set $T$, and spherical at $v'$ and at split places separated from $w_0$ and every member of a finite set $S$. At $w_0$ assume $U$ contains an elliptic type given by a nontrivial character of a torsion-free compact open $J$ through a quotient $J\twoheadrightarrow\Z/\ell^m$, with $m\ge1$, that forces cuspidality. For a dominant weight $\lambda=(\lambda_i)_{0\le i<a}$ put $J(\lambda)=\{i-\lambda_i:0\le i<a\}$. A finite set is lacunary if disjoint nonempty subsets have different means.

\begin{theorem}[Geometric seed]\label{B9.08}
Suppose $J(\lambda)$ is lacunary for a weight of $U$ at an embedding over $v'$. There are a dimension $d_\rho$ and a representation $\rho:G_E\to\GL_{d_\rho}(\bar{\Q}_\ell)$ with a lacunary subset of size $a$ in an actual jump multiset at $v'$. At usable $w\nmid\ell$ off a finite set, inertia is trivial and $P_{w,B}$ annihilates inverse Frobenius for every Satake multiset $B$ of $U$ there and every prolongation.

For a specified usable spherical $w\nmid\ell$ one can require the same local conclusion at $w$. At a spherical $w\mid\ell$ separated from $w_0,v'$ and agreeing with $v'$ on $\Q(x)$, one can require crystallinity with every prescribed jump at every induced embedding. Alternatively, one can require the crystalline divisibility of \ref{B9.19} for all $R=A_a(w,B)$ and all compatible period data. Each choice gives its own witness to the first paragraph.
\end{theorem}
\begin{proof}[Outline]
Twist the semisimplified eigenspace of \ref{B9.15} by the Hecke character in \ref{B9.09}. Use \ref{B9.17} for annihilators and \ref{B9.18a} and \ref{B9.19} for Hodge degrees, then shift the jumps and scale the roots by the twist. The complete proof is in \S\ref{prf:B9.08}.
\end{proof}

\begin{lemma}[Extension of the unitary eigensystem to similitudes]\label{B9.09a}
Use \(E/F,x,a,U,w_0\) of \ref{B9.08}, and put \(I=\Q(x)\). On \(E^a\) take the standard definite Hermitian form. Its group of finite adelic similitudes with rational adelic multiplier is
\[
 H=\{(g,s)\in\GL_a(\mathbf A_{E,f})\times\mathbf A_{\Q,f}^{\times}
             :{}^t\bar g g=s\,1_a\}.
\]
Write \(\nu(g,s)=s\), \(H^1=\ker\nu\), and let \(\Gamma\) be the subgroup of rational points with positive rational multiplier. For a finite idèle \(y\) of \(E\), put
\[
 \zeta(y)=
 \bigl(\overline{N_{E/I}y}\,1_a,N_{E/\Q}y\bigr).
\]
This lies in the centre of \(H\), since the product of \(N_{E/I}y\) and its conjugate is \(N_{E/\Q}y\).

For each infinite place \(v\) of \(E\), choose the embedding \(\tau_v\) used to describe the weight of \(U\). Write \(\lambda_v=(\lambda_{v,i})_{0\le i<a}\) for its highest weight, \(W_v\) for the corresponding algebraic representation of \(\GL_a(\C)\), and \(W=\bigotimes_v W_v\). Define
\[
 b_v=\sum_i|\lambda_{v,i}|,\qquad
 b=\sum_v\sum_i(\lambda_{v,i}+b_v),\qquad
 \xi(g,s)=s^b\bigotimes_v W_v^\vee(\tau_v g)
 \quad ((g,s)\in\Gamma).
\]
Choose a nonzero algebraic form in the weight space of \(U\) with the prescribed type at \(w_0\) and the prescribed spherical eigenvalues. There is a nonzero form \(\Phi:H\to W^\vee\) with left transformation law \(\xi\) and the same type and spherical eigenvalues for the multiplier-one operators. There is an algebraic Hecke character \(\kappa\) of \(E\) such that
\[
                 \Phi(h\zeta(y))=\kappa_f(y)^{-1}\Phi(h).
\]
Its exponent at an embedding \(\tau:E\hookrightarrow\C\) is
\[
 m_\tau=b-
       \sum_{\{v:\,\tau_v|_I=\bar\tau|_I\}}\sum_i\lambda_{v,i}.
\]
\end{lemma}
\begin{proof}[Outline]
Extend the unitary form by zero outside its rational and level translates, preserving the multiplier-one eigenvalues. Take a common eigenvector in the finite span of central translates and read the character's infinity exponents from the rational central action. The complete proof is in \S\ref{prf:B9.09a}.
\end{proof}

\begin{lemma}[The ambient PEL datum and the standard definite summand]\label{B9.09ab}
Under the field hypotheses of \ref{B9.08}, let \(X_0\) be a finite set of rational primes with \(-d\in(\Q_q^\times)^2\) for every \(q\in X_0\), let \(v_0\) be a distinguished real place of \(F\), and let \(\omega\in\mathcal O_I\subseteq\mathcal O_E\) generate the ring of integers of \(I=\Q(x)\) over \(\Z\).
\begin{enumerate}
\item There exist an \(E/F\)-Hermitian form \(h_{\rm PEL}\) on \(E^a\) of signature \((1,a-1)\) at \(v_0\) and \((0,a)\) at every real place \(v\ne v_0\) of \(F\), an \(\mathcal O_E\)-lattice \(\Lambda\subset E^a\), and a positive integer \(e\) prime to \(\prod_{q\in X_0}q\) such that the alternating form \(e_\Lambda(u,v)=\operatorname{Tr}_{E/\Q}(x\,h_{\rm PEL}(u,v))\) takes values in \(\Z\), satisfies \(e\,\Lambda^\vee\subseteq\Lambda\), and at every rational prime \(q\nmid e\) with \(-d\in(\Q_q^\times)^2\) (including every \(q\in X_0\)) the form \(e_\Lambda\) is perfect over \(\Z_q\) and the minimal polynomial of \(\omega\) has distinct roots modulo \(q\).
\item For every finite set \(X_1\) of rational primes, there is a rational prime \(p_0\notin X_1\) with \(p_0\nmid e\) and \(-d\in(\Q_{p_0}^\times)^2\). Specializing \((E^a,\mathcal O_E,\Lambda,e_\Lambda,\omega)\) at \(p_0\) gives a local PEL datum satisfying the hypotheses of Proposition~\ref{B9.10a} with \(m=1\) and \(n=a\), local similitude group \(G_{p_0}\), compact open stabilizer \(K_{p_0}\subset G_{p_0}\) of \(\Lambda\otimes\Z_{p_0}\), and an away-from-\(p_0\) trivialization identifying the away-from-\(p_0\) PEL similitude group with the away-from-\(p_0\) factor \(H^{p_0}\) of \(H\) in Lemma~\ref{B9.09a}.
\item There is \(\mu\in F^\times\) with \(\mu_{v_0}>0\) and \(\mu_v<0\) for every real \(v\ne v_0\) such that \(h_{\rm std}=\mu\,1_a\) on \(E^a\) has signature \((a,0)\) at \(v_0\), signature \((0,a)\) at every real \(v\ne v_0\), the same finite local invariants as \(h_{\rm PEL}\) at every finite place of \(F\), and rational positive-multiplier similitude group \(\Gamma=\{(g,s)\in\GL_a(E)\times\Q_{>0}:{}^t\bar g g=s\,1_a\}\).
\end{enumerate}
\end{lemma}
\begin{proof}[Outline]
Builds \(h_{\rm std}=\mu 1_a\) by Chebotarev's theorem and global norm reciprocity, and \(h_{\rm PEL}\) of matching finite invariants by weak approximation and Definition~\ref{B1.15ac}. Forms unimodular lattices at \(X_0\) and picks \(p_0\) split in \(I\) via Chebotarev. The complete proof is in \S\ref{prf:B9.09ab}.
\end{proof}

\begin{lemma}[The auxiliary prime and regular Iwahori operator]\label{B9.09ac}
Keep the notation of Lemma~\ref{B9.09ab}, let \(\Phi\) be the similitude form of Lemma~\ref{B9.09a}, let \(H,C\) be finite sets of usable places as in \ref{B9.09}, and fix the model prime \(p_0\) away from \(\ell\), the residue characteristics of \(H\cup C\), the bad level places of \(\Phi\), and a finite set outside which a chosen argument \(g_0\) with \(\Phi(g_0)\ne0\) is integral and unimodular. There exist a rational prime \(r_0\notin\{p_0,\ell\}\) outside the residue characteristics of \(H\cup C\), the bad level places of \(\Phi\), and the prime factors of \(e\), together with a split place \(W\mid r_0\) of \(E\), such that for every principal level \(N\ge3\) prime to \(p_0\) with \(r_0\mid N\) whose congruence subgroup \(K(N)\subset H^{p_0}\) stabilizes \(\Phi\) away from \(w_0\), the away-from-\(p_0\) Hecke algebra at level \(K(N)\) contains an idempotent \(e_{\rm typ}\) and elements \(u,u'\) supported at \(W\) over the unit of \(K(N)\) away from \(W\) with the following properties:
\begin{enumerate}
\item \(e_{\rm typ}\) acts as the identity on the transported form of \(\Phi\), \(u\) commutes with \(e_{\rm typ}\) and with every deck correspondence and spherical Hecke operator attached to places of residue characteristic different from \(r_0\), and \(u'ue_{\rm typ}=e_{\rm typ}\);
\item for every integer \(n\ge1\) and every word polynomial \(b\) in those deck and spherical operators, every element in the away-from-\(p_0\) support of \(b(T)u^ne_{\rm typ}\) has distinct eigenvalue valuations at \(W\), and therefore has a separable characteristic polynomial of degree \(a\).
\end{enumerate}
\end{lemma}
\begin{proof}[Outline]
Chebotarev's theorem chooses \(r_0\) and \(W\mid r_0\) to define \(e_{\rm typ}\) and the Iwahori operator \(u\). Bernstein's presentation gives \(u'\) and support in \(I_W t_W^n I_W\), where unit minors give \(a\) distinct Newton-polygon slopes. The complete proof is in \S\ref{prf:B9.09ac}.
\end{proof}

\begin{proposition}[Geometric plans]\label{B9.09}
Under the field and local representation hypotheses of \ref{B9.08}, there is a plan with a tower of fibre powers and an algebraic character $\kappa$. Its Hecke operators $T_g$ have a prescribed eigensystem $(\theta_g)_g$ over $\bar{\Q}_\ell$. On words applied to a spectral test $f$ the Hecke operators commute. A word polynomial is a finite linear combination of products of the operators, evaluated at $\theta$ by replacing each operator with its scalar. For every word polynomial $s$ with $s(\theta)\ne0$, some word polynomial $a_0$ gives a regular test $F=a_0sf$ with nonzero trace $D(F)$ on the direct sum of the definite similitude form spaces at the transported weight. The geometric orbital terms with a marked root and reducible polynomial vanish, as do the definite orbital terms outside the regular irreducible locus.

The plan satisfies the Hodge condition at $v'$. At a Hodge place the selected correspondences descend to a common smooth proper model and $\kappa$ is unramified. At each embedding $\tau$ over the place the transported weights give the degrees $j+m_\tau$ of \ref{B9.18a}. Here $j$ ranges over the prescribed jumps and $m_\tau$ is the infinity exponent of $\kappa$. Off a finite set of usable $w\nmid\ell$, $\kappa$ is unramified and the correspondence roots are $\kappa(w)^{-1}A_a(w,B)$ for every $B$. One may prescribe this root identity at a usable spherical $w\nmid\ell$. At $w\mid\ell$ as in \ref{B9.08}, one may prescribe either the Hodge condition or the root identity at $w$.
\end{proposition}
\begin{proof}[Outline]
Pairs the PEL datum, form \(\Phi\), model prime \(p_0\), and auxiliary Iwahori operator \(u\) to place the eigensystem in the definite trace. Newton's identities on the eigenspace projector give \(D(F)\ne0\), and constant-term vanishing at \(w_0\) kills reducible terms. The complete proof is in \S\ref{prf:B9.09}.
\end{proof}

\begin{definition}[PEL objects]\label{B9.10}
Fix the base datum over $R$ with CM field $F_{\rm PEL}$, order $\mathcal O$, and maps $t_i:\mathcal O_{F_{\rm PEL}}\to R$ taking the distinguished element $\omega$ to a common scalar $\rho$. A PEL object is an abelian scheme $A/S$ with an $\mathcal O$-action. For every $R$-algebra $R'$ and map $x:\operatorname{Spec}R'\to S$ over $R$, each tangent vector at the identity of $x^*A$ on which $\omega$ acts by the image of $\rho$ must be a sum of vectors on which $\mathcal O_{F_{\rm PEL}}$ acts through the respective maps $t_i$ followed by $R\to R'$. At each geometric point the $t_i$-parts have the prescribed Lie multiplicities.

The polarization is given by a symmetric line bundle $P$ on $A\times_S A$, trivializable on the axes and relatively ample on the diagonal. For every order element $b$ it satisfies $(b\times1)^*P\simeq(1\times b^*)^*P$, where $b^*$ is the involution of the order. The level is an $\mathcal O$-linear map $\eta:\Lambda/N\Lambda\to A[N](S)$ from the datum lattice modulo $N$. At every geometric point it must be the level-$N$ evaluation of a Tate module isomorphism compatible with the specified Weil pairings. The universal object classifies these objects by pullback up to the prescribed isomorphism.
\end{definition}

\begin{proposition}[The integral PEL moduli problem]\label{B9.10a}
Let \(p\) be a prime. Write \(K_0\) for the CM field of the PEL datum, \(\mathcal O\) for the order with involution, \(\Lambda\) for its lattice, and \(e_\Lambda\) for the alternating form on \(\Lambda\). Suppose the form is perfect over \(\Z_p\) and
\[
 \mathcal O\otimes\Z_p\simeq M_m(\mathcal O_{K_0}\otimes\Z_p).
\]
Write \(\operatorname{rank}_{\Z}\Lambda=mn[K_0:\Q]\). Let \(\omega\in\mathcal O_{K_0}\) generate an imaginary quadratic order whose defining polynomial has distinct roots modulo \(p\). Use the polarization and the complex signature specified by the datum in \ref{B9.10}.

Let \(R\) be a Noetherian \(\Z_{(p)}\)-algebra. On the chosen side of \(\omega\), let \(t_1,\ldots,t_b:\mathcal O_{K_0}\to R\) be the structural maps. They take \(\omega\) to the same root. For \(i\ne j\), the differences \(t_i(c)-t_j(c)\), as \(c\) varies, generate \(R\), and the reductions of \(t_i,t_j\) at each geometric point of characteristic \(p\) have distinct kernels in \(\mathcal O_{K_0}\). On this side the prescribed Lie multiplicity at each \(t_i\) is \(m\).

For \(N\ge3\) prime to \(p\), use the objects of \ref{B9.10} with the following integral conditions. The kernel of the polarization homomorphism defined by the bundle on \(A\times A\) is killed by an integer prime to \(p\). At a geometric point the level is induced by a pairing-compatible isomorphism of the prime-to-\(p\) Tate modules. Impose the Lie condition of \ref{B9.10} after every scalar extension, including extensions to rings with nilpotents. This functor is represented by a smooth quasi-projective \(R\)-scheme \(X_N\) of relative dimension \(b(n-1)\).

If \(R\) is a domain of characteristic zero, \(X_N\) is projective under either of the following conditions:
\[
  2b<[K_0:\Q],\qquad
  \text{or}\qquad
  \mathcal O\otimes\Q\text{ is a division algebra and } n<2m.
\]
\end{proposition}
\begin{proof}[Outline]
Rigidifies over the Siegel scheme with a level-\(M\) frame, represents PEL data via bounded homomorphism graphs, and descends for \(N\ge3\). Proves smoothness by Dieudonn\'e freeness and Grothendieck--Messing, and projectivity via the valuative criterion ruling out a torus. The complete proof is in \S\ref{prf:B9.10a}.
\end{proof}

\begin{theorem}[Representability]\label{B9.11}
Over a Noetherian $\Q$-algebra, the PEL datum of \ref{B9.10a} with its structural maps and level $N\ge3$ prime to $p$ represents the moduli problem of \ref{B9.10} by a scheme.
\end{theorem}
\begin{proof}[Outline]
Apply \ref{B9.10a} in characteristic zero. Restrict the full Tate frame away from \(p\), and reconstruct it by adjoining dual bases at \(p\) while preserving the level sections, identifying the two moduli functors. The complete proof is in \S\ref{prf:B9.11}.
\end{proof}

\begin{proposition}[Fibre powers]\label{B9.12}
Over a characteristic-zero field, compatible compact PEL data give smooth projective fibre powers of every nonnegative exponent with finite étale first Hecke projections of prescribed constant degree.
\end{proposition}
\begin{proof}[Outline]
Use \ref{B9.10a} and \ref{B9.11} for the base and universal abelian scheme, and use exterior products of relatively ample bundles on its fibre powers. The Hecke projections come from finite étale level forgetting and the factorwise isogeny. The complete proof is in \S\ref{prf:B9.12}.
\end{proof}

\begin{proposition}[Integral models and the two Hecke projections at a selected place]\label{B9.12a}
Let \(w\) be a place selected in \ref{B9.09}, \(p\) its residue characteristic, \(R=\mathcal O_{E,(w)}\), and \(q=|k(w)|\). Identify \(K_0\) with \(E\) by the embedding in the plan and write \(\mathfrak w\) for the inverse image of \(w\). The lower model at \(w\) and its abelian fibre powers are smooth projective over \(R\).

At a root place and for \(0\le j\le a\), the lower Hecke correspondence on the moduli scheme has an integral model
\[
 X\xleftarrow{s_j}Y_j\xrightarrow{t_j}X.
\]
Both arrows are finite, \(Y_j\) is flat over \(R\), and both arrows are étale on the generic fibre. The universal isogeny has polarization multiplier \(q^j\). On the chosen side its generic kernel is a \(j\)-dimensional \(k(w)\)-subspace of \(A[\mathfrak w]\). On a fibre power the first projection is finite, is étale after inverting \(p\), and then has degree \(\#\operatorname{Gr}_j(k(w)^a)\). At a Hodge place the selected operators admit correspondences on one smooth projective lower model with finite étale first projections over \(R\).
\end{proposition}
\begin{proof}[Outline]
Use the chart of \ref{B9.09}, the model of \ref{B9.10a}, and fibre powers as in \ref{B9.12}. Parametrize root-place kernels by Grassmannians and use the tangent bound to obtain finite projections on normalization, with étale subgroup indices at Hodge places for \ref{B9.19}. The complete proof is in \S\ref{prf:B9.12a}.
\end{proof}


\begin{definition}[Total cohomology]\label{B9.13}
For a fibre power $Z/E$, use geometric pullback on the total rational cohomology $\bigoplus_j\Q_\ell\otimes_{\Z_\ell}\varprojlim_mH^j_{\mathrm{\acute et}}(Z_\Omega,\Z/\ell^m)$ \cite{milne-1980}.
\end{definition}

\begin{lemma}[Correspondences]\label{B9.14}
For a correspondence $Z\xleftarrow{s}Y\xrightarrow{t}Z$ with $s$ finite étale of constant degree, set $T=\operatorname{Tr}_s t^*$, where $\operatorname{Tr}_s$ is the finite étale trace. If a Galois cover $b$ of $Z$ splits $s$ into branches $\sigma_j$ viewed as maps to $Y$, then $b^*Tu=\sum_j(t\circ\sigma_j)^*u$ on cohomology.
\end{lemma}
This is finite étale trace base change with rational $\ell$-adic coefficients \cite{milne-1980}.

\begin{proposition}[Eigenspaces]\label{B9.15}
In \ref{B9.13} tensored with $\bar{\Q}_\ell$, the common eigenspace for the correspondences and eigenvalues of \ref{B9.09} is Galois-stable with finite-dimensional continuous semisimplification.
\end{proposition}
\begin{proof}[Outline]
Decompose the smooth projective geometric fibre into finitely many integral components and apply \ref{B9.26} for finite dimension, with continuity from inverse-limit cohomology. Correspondences over \(E\) commute with Galois, and a composition flag gives the semisimplification. The complete proof is in \S\ref{prf:B9.15}.
\end{proof}

\begin{lemma}[Ordinary points]\label{B9.16}
At an ordinary point of the integral PEL model at $w$, let $T_j$ be the branch multiset indexed by $j$-planes of $H=k(w)^a$. For residue Frobenius $F_*$ there are multisets $A_j$ with $A_0=T_0$, $A_a=0$ and $T_{j+1}=F_*A_j+q_w^{j+1}A_{j+1}$ for $0\le j<a$. Consequently
\[
 \sum_{j=0}^a(-1)^jq_w^{\binom j2}F_*^{a-j}T_j=0.
\]
Let $Y$ be the plan's integral PEL model, and let $V$ be a valuation ring over its base with algebraically closed fraction field and residue characteristic $p$ nonzero in $V$. If $y:\operatorname{Spec}V\to Y$ sends the closed point to a generic point of the special fibre, the reduced branches satisfy the same identity in the free abelian group on residue-field points.
\end{lemma}
\begin{proof}[Outline]
Use deformations on the lower model of \ref{B9.12a} to get generic connected height one. Planes containing the reduction kernel line yield Frobenius, and maps into the line count transverse lifts, giving a telescoping recurrence. The complete proof is in \S\ref{prf:B9.16}.
\end{proof}

\begin{proposition}[Annihilator]\label{B9.17}
At $w\nmid\ell$ the root identity of \ref{B9.09} gives trivial inertia and annihilator $P_{w,B}$ at inverse Frobenius on the piece of \ref{B9.15} twisted by $\kappa_\ell$.
\end{proposition}
\begin{proof}[Outline]
Descend the deck invariants of \ref{B9.09} by normalized trace on the model of \ref{B9.12a}. Use \ref{B9.14} and \ref{B9.27} to carry the branch relation of \ref{B9.16} to generic cohomology, then divide by correspondence degrees, normalize and twist. The complete proof is in \S\ref{prf:B9.17}.
\end{proof}

\begin{theorem}[Holomorphic Lefschetz]\label{B9.18}
Let $Z/\C$ be smooth projective of dimension $D$ and $Z\xleftarrow{c_1}Y\xrightarrow{c_2}Z$ a finite étale correspondence with finitely many transverse fixed points. At each fixed point $y$, put $A_y=dc_2(dc_1)^{-1}$ on the tangent space of $Z$ at $c_1(y)=c_2(y)$. For $j\ge0$,
\[
 \sum_{b=0}^D(-1)^b\operatorname{Tr}(\operatorname{Tr}_{c_1}c_2^*
 \mid H^b(Z,\Omega_Z^j))
 =\sum_y\frac{[T^j]\det(1+TA_y)}{\det(1-A_y)}.
\]
\end{theorem}
This is the Atiyah--Bott formula restricted to the indicated bundles and correspondences \cite{atiyah-bott-1968}.

\begin{lemma}[The orbital weight of a marked eigenline]\label{B9.18b}
Use the odd rank \(a\), the single indefinite place, and the regular test \(F\) of \ref{B9.09}. Give a complex base fixed point of characteristic polynomial \(P\) weight \(c(P,\mu)\), where \(\mu\) is the eigenvalue on its positive line at the indefinite place. Use the same finite-level and stabilizer weights as for the Hecke trace. The sum of these weights is the sum over the regular irreducible definite classes with polynomial \(P\), with their definite orbital weights multiplied by
\[
                         \sum_{P(\mu)=0}c(P,\mu).
\]
The roots are taken at the chosen complex embedding.
\end{lemma}
\begin{proof}[Outline]
Uniformization identifies fixed points with similitudes preserving a positive line; reducible terms vanish by \ref{B9.09}. For irreducible \(P\), global reciprocity in \(E[\gamma]/B^+\) twists the form to match definite orbital weights of \ref{B9.09ab} for each root \(\mu\). The complete proof is in \S\ref{prf:B9.18b}.
\end{proof}

\begin{lemma}[Hodge degrees of a spectral test]\label{B9.18a}
For a dominant-weight plan of \ref{B9.09}, the fibre exponent $k$ can be chosen uniformly for the conjugates of the weight. For each complex embedding $\tau$ over a Hodge place choose the automorphism $\sigma\in\operatorname{Aut}(\C)$ provided by the plan's weight transport. For each prescribed jump $j$ there are a datum embedding indexed by $i_0$ and an index $i$ in its Hodge exponent vector $h$ with $h_i=j+m_\tau$. On the fibre power $Z$ over $\sigma\circ\tau$, some $H^b(Z,\Omega_Z^{h_i})$ contains a nonzero common eigenvector for the correspondences with eigenvalues $\sigma(\iota(\theta_g))$.
\end{lemma}
\begin{proof}[Outline]
Use the regular test of \ref{B9.09} and interpolate determinants after canceling the fibre denominator. Expand Schur alternants and apply \ref{B9.18b} and \ref{B9.18} to compute the Hodge trace, then isolate the common eigenkernel. The complete proof is in \S\ref{prf:B9.18a}.
\end{proof}

\begin{proposition}[Crystalline data]\label{B9.19}
At $w\mid\ell$, a surjective good-model cover as in \ref{B9.17} gives crystallinity of the semisimplified piece $\rho$ of \ref{B9.15}. If the correspondences descend to one model with the prescribed multiplicities, a common Hodge eigenclass of degree $h$ gives jump $h$ at the induced embedding. Under the root identity of \ref{B9.09}, write $R$ for the correspondence root multiset before twisting. The piece is crystalline, and for each compatible period datum, prolongation and bridge there are a Weil element $g$ of degree $-1$ and an integer $m\ge0$ such that the Weil polynomial of $\rho$ at $g$ divides $(\prod_{\alpha\in R}(X-\alpha))^m$.
\end{proposition}
\begin{proof}[Outline]
Apply smooth proper comparison to the lower model of \ref{B9.12a} using descent from \ref{B9.17}, and retain the degree from \ref{B9.18a} by strict filtered comparison. The branch relation of \ref{B9.16} bounds the Frobenius spectrum after degree normalization. The complete proof is in \S\ref{prf:B9.19}.
\end{proof}

\section{Cohomological inputs}

Write $H^j(Z,i)=H^j_{\mathrm{\acute et}}(Z,\Z/b(i))$ for the Tate twist $i$ with $b\ge1$ invertible on the base. The notation includes simplicial presheaves.

\begin{definition}[Core]\label{B9.20}
C2--C5 and D1 are \ref{B9.21}--\ref{B9.25}, respectively.
\end{definition}

\begin{theorem}[Affine curves, C2]\label{B9.21}
For a finite-type Dedekind $k$-algebra $A$ over an algebraically closed field $k$, $H^j(\operatorname{Spec}A,i)=0$ for $j\ge2$.
\end{theorem}
This is affine-curve cohomological dimension for Tate coefficients \cite{milne-1980}.

\begin{proposition}[Generic restriction, C3]\label{B9.22}
For an absolutely integrally closed local domain $T$ with fraction field $K$ and an étale $T[t]$-algebra $B$, restriction from $\operatorname{Spec}B$ to $\operatorname{Spec}(K\otimes_T B)$ is an isomorphism in every degree and twist.
\end{proposition}
\begin{proof}[Outline]
Use smooth local acyclicity on strict henselian stalks to compare the Tate twist with its derived pushforward from the geometric generic fibre. Taking global cohomology gives the restriction isomorphism. The complete proof is in \S\ref{prf:B9.22}.
\end{proof}

\begin{theorem}[Affine lines, C4]\label{B9.23}
Projection from the affine line over any ring induces isomorphisms in every degree and twist.
\end{theorem}
This is étale homotopy invariance for invertible Tate coefficients \cite{artin-grothendieck-verdier-1973}.

\begin{theorem}[Laurent lines, C5]\label{B9.24}
For strictly henselian local $A$, the Laurent line $\operatorname{Spec}A[t,t^{-1}]$ has degree-one group $\Z/b$ additively and vanishing in degrees at least two, in every twist.
\end{theorem}
This is the Laurent-line formula over a strictly henselian base \cite{milne-1980}.

\begin{proposition}[Projective bundles, D1]\label{B9.25}
For a rank-$N$ bundle $V$ on a simplicial presheaf $\mathcal E$, let $\pi:\mathbf P(V)\to\mathcal E$ be its bundle of lines and $L$ its tautological line. If $b$ is invertible on schemes with zero-simplices, the map $(u_j)_j\mapsto\sum_j\pi^*u_jc_1(L)^j$ identifies $\bigoplus_{0\le j<N}H^{n-2j}(\mathcal E,i-j)$ with $H^n(\mathbf P(V),i)$.
\end{proposition}
\begin{proof}[Outline]
For the trivial bundle, apply the scheme projective bundle formula to the map of complexes given by Chern-power cup products. Apply this in each simplicial degree with compatible Chern cocycles, then compare the bounded-below total complexes. The complete proof is in \S\ref{prf:B9.25}.
\end{proof}

\begin{theorem}[Duality]\label{B9.26}
For smooth projective integral $Z$ over an algebraically closed field $k$ with $\ell$ invertible in $k$, rational $\ell$-adic cohomology is finite-dimensional in each degree and vanishes for $j>2\dim Z$. The cup product followed by trace in degree $2\dim Z$ pairs complementary degrees perfectly after Tate trivialization.
\end{theorem}
This is smooth projective Poincaré duality with rational $\ell$-adic coefficients \cite{milne-1980}.

\begin{theorem}[Specialization]\label{B9.27}
Let $Y$ be smooth proper over a ring $O$, and let $V$ be a valuation ring over $O$ with algebraically closed fraction and residue fields. Restriction from $Y_V=Y\times_O\operatorname{Spec}V$ to either fibre is an isomorphism for moduli invertible in the residue field. Rational specialization at such a prime $\ell$ respects pullbacks and compatible automorphisms.
\end{theorem}
This is smooth proper specialization over $V$ \cite[Exp.~XVI, {\S}2]{artin-grothendieck-verdier-1973}.


\section{Rank extraction and congruence limits}

For a dominant weight $\lambda=(\lambda_0,\ldots,\lambda_{a-1})$, write $j_i=i-\lambda_i$. The ordered jumps have the stronger gaps if $j_{i+1}-j_0>a(j_i-j_0)$ for $0\le i<a-1$. A multiset of $a$ nonzero roots is in general position if the following polynomial is nonzero in every enumeration.
\begin{equation}\label{eq:p9c1-general-position}
 \prod_{i\ne j}\prod_{h=1}^a(\alpha_i^h-\alpha_j^h)
 \prod_{j\ne i,\ j\ne k}(\alpha_i\alpha_k-\alpha_j^2)
 \prod_{\substack{\sum_i m_i=a,\ m_i\in\Z_{\ge0}\\(m_i)\ne(1,\ldots,1)}}
 \left(\prod_i\alpha_i^{m_i}-\prod_i\alpha_i\right).
\end{equation}
For a spherical Satake multiset $B$ the test may be applied to either $B$ or $A_a(w,B)$, since each factor is homogeneous.

\begin{proposition}[Seed tests]\label{B9.28}
In the odd-rank field setting of \ref{B9.08}, fix separated split $w_0,v'\mid\ell$, a reference representation $U$ and finite sets $S,Y$ of places. There is a finite set $Z_0$ of usable places away from $\ell$ separated from $w_0,S$ with the following property. Suppose $U'$ is unramified at nonsplit places off $T$, has elliptic type at $w_0$, agrees with the weights of $U$ off the pair, and is spherical at $v'$ and at split places separated from $w_0,S$. Suppose it has stronger gaps at an embedding inducing $v'$ and is in general position on $Z_0$.

There is a continuous semisimple $r:G_E\to\GL_a(\bar{\Q}_\ell)$ and there are a dense open set $V\subset G_E$ and a finite set $X$ of places with the following properties. At each usable $w\nmid\ell$ separated from $S$ and outside $X$, for each Satake multiset $B$ there, $r$ is unramified and the transported polynomial at $s^{-1}$ is $P_{w,B}$ for every arithmetic Frobenius lift $s\in V$. At every usable $w\nmid\ell$ in $Y$ separated from $S$, a general-position $B$ gives unramifiedness and this polynomial at every inverse arithmetic Frobenius lift. At spherical $w\mid\ell$ in $Y$ separated from $w_0,v',S$, $r$ is crystalline and de Rham with the prescribed jumps. For each general-position $B$ there, its unramified Weil polynomial at $\sigma\in W_{E_w}$ has roots $\alpha^{-\deg\sigma}$ for $\alpha\in A_a(w,B)$. The place statements hold for every prolongation and the period statements for every compatible datum and bridge.
\end{proposition}
\begin{proof}[Outline]
Use the gaps and split coordinates to apply \ref{B9.08}, normalize the determinant by \ref{B9.50}, and extract rank by \ref{B9.29}. Use degree-one Chebotarev in \ref{B8.42p} to compare the witnesses and \ref{B9.30} to recover multiplicities and local data. The complete proof is in \S\ref{prf:B9.28}.
\end{proof}

\begin{proposition}[Rank extraction]\label{B9.29}
Let $\rho$ be a continuous semisimple $\bar{\Q}_\ell$-representation of $G_E$ which is Hodge--Tate at a place above $\ell$. Suppose an actual jump multiset there contains a lacunary subset of size $a\ge1$. Suppose that on a dense set $\rho(g)$ has a monic degree-$a$ annihilator $Q_g$ with $(-1)^aQ_g(0)=\eta(g)$ for a continuous character $\eta$. There is a continuous semisimple rank-$a$ representation $r$ of determinant $\eta$, with dense open separable locus and $r\hookrightarrow\rho\hookrightarrow r^{\oplus m}$ for some $m$.
\end{proposition}
\begin{proof}[Outline]
Pass the minimal-polynomial bound to the Lie algebra and use Sen's theorem and lacunarity to obtain scalar corners. Take one standard module per matrix factor and use the dense separable locus to recover characteristic polynomials and the determinant. The complete proof is in \S\ref{prf:B9.29}.
\end{proof}

\begin{lemma}[Multiplicities]\label{B9.30}
Let an $a$-by-$a$ matrix $M$ be annihilated by the root polynomial of a size-$a$ multiset $S$, and suppose $\det M=\prod_{\alpha\in S}\alpha$. If $S$ is the unique size-$a$ multiset supported in $\operatorname{supp}S$ with this product, then $M$ has spectrum $S$.
\end{lemma}
\begin{proof}[Outline]
The annihilator, determinant, and uniqueness in \ref{B9.28} identify the spectrum. Above \(\ell\), crystallinity passes to witnesses: the jump witness forces all \(a\) distinct jumps, while \ref{B9.19}, determinant comparison, and unramifiedness give the Weil polynomial. The complete proof is in \S\ref{prf:B9.30}.
\end{proof}

\begin{lemma}[Approximation with one varying level]\label{B9.30a}
Let $E/F$ be CM, $G$ the definite unitary group of rank $a\ge1$, and $G_\infty$ its archimedean factor. Fix $\iota:\bar{\Q}_\ell\simeq\C$, a split $w_0\mid\ell$, an algebraic weight $V$, and a compact open level $K$ away from the place below $w_0$, embedded in $G(\mathbf A_{F,f})$ with component $1$ there. Let $D$ be a finite set of split places separated from $w_0$ where $K$ has hyperspecial factors. For each $w\in D$ take a nonzero polynomial $\Phi_w\in\C[X_1,\ldots,X_a]$. A spherical system is good at $w$ if $\Phi_w$ is nonzero on every ordering of its Satake roots.

Let $f(v,g)$ be an algebraic weight form fixed by $K$, linear for $v\in V$, left invariant under $G(F)$ and smooth at finite places, with $f(v,gx)=f(xv,g)$ for $x\in G_\infty$. For every $\epsilon>0$ there is a finite sum $F$ of $K$-fixed forms of weight $V$ from representations good throughout $D$ and a scalar function $e$ such that
\begin{equation}\label{eq:p9c1-generality-relative}
 f(v,g)-F(v,g)=e(g)f(v,g),\qquad |\iota^{-1}e(g)|_\ell\le\epsilon
 \quad(v\in V,\ g\in G(\mathbf A_F)).
\end{equation}
\end{lemma}
\begin{proof}[Outline]
Use spherical Hecke polynomials of small Plancherel norm and deep-level separation to bound the exceptional span. On the finite \(\ell\)-group, use the determinant estimate on translates of an annihilating row and iterate on the pointwise error. The complete proof is in \S\ref{prf:B9.30a}.
\end{proof}

\begin{proposition}[Hecke approximation]\label{B9.31}
Let $E/F$ be CM, $a\ge1$ odd, $\ell$ prime, and $\iota:\bar{\Q}_\ell\simeq\C$. Let $T$ be a finite nonsplit set and $U$ an automorphic representation of the definite unitary group of rank $a$, unramified at nonsplit places off $T$. Fix split $w_0,v'\mid\ell$ above distinct places of $F$. There are compact open subgroups $K_w$ at split $w$ separated from $w_0$. At each split $w$ separated from both places, $U_w^{K_w}\ne0$, and $K_w=\GL_a(\mathcal O_{E_w})$ if $U_w$ is spherical.

For finite sets $Z,W$ of split spherical places of $U$ separated from both places, with $W$ above $\ell$, there exist nonzero complex $c_{w,n}$ for $(w,n)\in W\times\Z$. For any index set $J$, let $f_j$ be finitely supported local operators at split $w_j\nmid\ell$, with $|\iota^{-1} f_j(g)|_\ell\le1$ on their coefficients, each acting by a scalar on the $K_{w_j}$-invariants of each smooth irreducible representation. Write $\theta_j$ for its scalar on $U$, transported by $\iota^{-1}$. Augment the tuple by
\[
 \theta_{w,n}=\iota^{-1}\left(c_{w,n}\sum_{\alpha\in A_a(w,B_w)}\alpha^n\right),
 \qquad (w,n)\in W\times\Z,
\]
where $B_w$ is the Satake multiset of $U$. The augmented tuple satisfies \ref{B9.32} with the augmented tuples of representations $U'$ for the same operators and constants. Each $U'$ is unramified at nonsplit places off $T$, has nonzero $K_w$-invariants at split places separated from the pair, and has the weights of $U$ at embeddings inducing places separated from the pair. It has elliptic type at $w_0$, is spherical at $v'$ with the stronger gaps at some embedding over $v'$, and is in general position at every place of $Z\cup W$. If every $f_j$ is rational-valued, a single finite extension $L/\Q_\ell$ contains all $\theta_j$ for $j\in J$.
\end{proposition}
\begin{proof}[Outline]
Scale general-position and Laurent operators (\ref{B9.30a}) to contract the lattice norm, and apply \ref{B9.30a}. Composing polynomial weight-shift and congruence elliptic-type steps gives the target tuples; finite Galois orbits bound the rational eigenvalue field. The complete proof is in \S\ref{prf:B9.31}.
\end{proof}

\begin{definition}[Polynomial approximation]\label{B9.32}
For $\vartheta\in k^J$ and $\mathcal S\subset k^J$ over an ultrametric field $k$, we say that $\vartheta$ admits polynomial approximation by $\mathcal S$ if for each $\delta>0$ there are a finite set $s\subset\mathcal S$ and weights $c_\eta\in k$ such that
\[
 \left|P(\vartheta)-\sum_{\eta\in s}c_\eta P(\eta)\right|\le\delta
 \qquad(P\in\Z[X_j:j\in J]).
\]
\end{definition}

\begin{proposition}[Trace limit]\label{B9.33}
Let $x:J\to\Gamma$ have dense range in a compact group and $\vartheta:J\to L$ with $[L:\Q_\ell]<\infty$. If the samples of a family of continuous rank-$a$ traces over $\bar{\Q}_\ell$ give \ref{B9.32}, then $\vartheta$ extends to the trace of a continuous semisimple rank-$a$ representation, with \ref{B9.32} on $\Gamma$ for the same family.
\end{proposition}
\begin{proof}[Outline]
Applying \ref{B9.32} to \(X_i-X_j\) gives Cauchy convergence in \(L\), extending \(\vartheta\) continuously to \(\Gamma\) with \ref{B9.32}. Testing \(X_1-a\) at the identity and taking limits of cyclic and degree-\((a+1)\) alternating trace identities lets \ref{B9.34} apply. The complete proof is in \S\ref{prf:B9.33}.
\end{proof}

\begin{theorem}[Pseudocharacters]\label{B9.34}
Over an algebraically closed field of characteristic zero, every dimension-$a$ group pseudocharacter is the trace of a semisimple rank-$a$ representation, unique up to isomorphism. Over a topological coefficient field this representation is continuous if its trace is continuous.
\end{theorem}
The trace realization is the group case of Taylor--Rouquier \cite{taylor-1991,rouquier-1996}. For continuity, choose a basis of the semisimple image algebra from group elements. The trace pairing is nondegenerate, since on each matrix factor it is the ordinary trace pairing times the positive multiplicity of its standard module. Its dual basis expresses the matrix coordinates as linear combinations of translates of the trace.

\begin{proposition}[Removing gaps]\label{B9.35}
Under the odd-rank field hypotheses of \ref{B9.08}, take split $w_0,v'\mid\ell$ separated over $F$, and automorphic $U$ unramified at nonsplit places off the finite nonsplit set $T$. There is a continuous semisimple $r:G_E\to\GL_a(\bar{\Q}_\ell)$. At every usable spherical $w\nmid\ell$ it is unramified and has transported polynomial $P_{w,B}$ at every inverse arithmetic Frobenius lift. At every spherical $w\mid\ell$ separated from the pair it is crystalline and de Rham with the jumps prescribed by the weight of $U$. For each such $w\mid\ell$, each Satake multiset $B$, compatible period datum, prolongation, bridge and Weil element $\sigma$, there are seeds with the common jumps and the joint approximation of \ref{B9.36}, with
\[
 S_n=\iota^{-1}\left(\sum_{\alpha\in A_a(w,B)}\alpha^{-n\deg\sigma}\right).
\]
\end{proposition}
\begin{proof}[Outline]
Approximate by \ref{B9.31} and \ref{B9.28} on dense samples from \ref{B8.42p}, then pass \ref{B9.32} to the limit and use \ref{B9.33} to realize the trace. Use \ref{B9.34} for inertia, Newton's identities for Frobenius, and \ref{B9.36} above \(\ell\). The complete proof is in \S\ref{prf:B9.35}.
\end{proof}

\begin{lemma}[The integral action and finite de Rham congruences]\label{B9.35aa}
Let \(K/\Q_\ell\) be finite, with residue field \(k\) and uniformizer \(\pi\). Write \(C\) for its completed algebraic closure and \(G_K=\Gal(\bar K/K)\). Choose compatible \(\ell\)-power roots \(\pi^\flat\) of \(\pi\). Put
\[
 \mathfrak S=W(k)[[u]],\qquad A_{\inf}=W(\mathcal O_{C^\flat}),
 \qquad \mathcal W=W(C^\flat).
\]
Embed \(\mathfrak S\) into \(A_{\inf}\) by \(u\mapsto[\pi^\flat]\). Let \(E\in W(k)[u]\) be the monic Eisenstein polynomial of \(\pi\), and use \(E\) also for its image, which generates the kernel of \(\theta:A_{\inf}\to\mathcal O_C\). Write \(\phi\) for Witt Frobenius, with \(\phi(u)=u^\ell\).

Let \(\Lambda\) be a finite free \(\Z_\ell\)-lattice of rank \(d\) with continuous action \(\rho\) of \(G_K\), and suppose \(\Lambda[1/\ell]\) is crystalline. Choose a basis of its Breuil--Kisin module and a basis of \(\Lambda\). Let \(T\in\GL_d(\mathcal W)\) have as columns the images of the first basis in the second under étale comparison. Put
\[
 \varepsilon_g=T^{-1}\rho(g)g(T).
\]
Then \(\varepsilon_g-I_d\in uM_d(A_{\inf})\) for every \(g\in G_K\). For each \(j\ge1\) there are \(k_j\ge0\) and \(B\in M_d(\mathfrak S)\) such that
\[
 B\equiv\ell^{k_j}I_d\pmod{E\mathfrak S},\qquad
 \phi(\varepsilon_g)g(B)\equiv B\pmod{E^jA_{\inf}}
 \quad(g\in G_K).
\]
\end{lemma}
\begin{proof}[Outline]
Use the crystalline lattice theorem and Tate twists to get the identity action matrix modulo \(u\). On the Frobenius pullback, normalize the constant isocrystal comparison to the identity modulo \(E\) and clear denominators at each finite order. The complete proof is in \S\ref{prf:B9.35aa}.
\end{proof}

\begin{lemma}[A common denominator at fixed order]\label{B9.35ab}
Keep \(K,\pi^\flat,\mathfrak S,A_{\inf},E,\theta\) from Lemma~\ref{B9.35aa}. For each \(J\ge1\) there is an integer \(c_J\ge0\) with the following property. Let \(d\) be any rank and let \(P_h\in M_d(A_{\inf})\) for \(h\in G_K\). Suppose that for some \(k\ge0\) a matrix \(B\in M_d(\mathfrak S)\) satisfies
\[
 B\equiv\ell^kI_d\pmod E,\qquad
 P_hh(B)\equiv B\pmod{E^J A_{\inf}}\quad(h\in G_K).
\]
Then a matrix over \(\mathfrak S\) satisfies both congruences with \(c_J\) in place of \(k\). The exponent is chosen before \(d\), the matrices \(P_h\), and \(k\).
\end{lemma}
\begin{proof}[Outline]
Use the cyclotomic periods of \ref{B9.35aa} to choose \(g\) with \(\lambda=\theta(g(E)/E)\) of infinite order. Induct on the powers of \(E\), bounding the denominator at step \(i\) by the valuation of \(\lambda^i-1\). The complete proof is in \S\ref{prf:B9.35ab}.
\end{proof}

\begin{lemma}[Uniform denominator for a Weil element]\label{B9.35a}
Assume the joint approximation and the representation and period hypotheses of \ref{B9.36}, and fix its Weil element. The restriction $r_w$ is crystalline. There are integral matrix lattices for the limit and the finite layers of seeds, with projection matrices $\Pi_m$ to the limit and identity columns $e,e_m$. Their Weil operators have integral numerators $V,V_m$ with a common denominator $s\in A_{\rm cris}$ invertible in $B_{\rm cris}$. For each precision $L$ some $m\ge L$ satisfies
\begin{equation}\label{eq:p9c1-weil-intertwine}
 V\Pi_m\equiv\Pi_m V_m,\qquad \Pi_m e_m\equiv e\pmod{\ell^L}.
\end{equation}
Here the operator on invariant columns represents the action at the fixed Weil element after extension to the crystalline period ring.
\end{lemma}
\begin{proof}[Outline]
Use \ref{B9.34} and \ref{B9.36} to project seed lattices to the limit and \ref{B9.35aa} and \ref{B9.35ab} to pass de Rham congruences to it. Iterate Frobenius to construct crystalline bases and compare Weil operators with a common denominator using \ref{B9.06}. The complete proof is in \S\ref{prf:B9.35a}.
\end{proof}

\begin{proposition}[Crystalline limits]\label{B9.36}
Let $r:G_E\to\GL_a(\bar{\Q}_\ell)$ be continuous and semisimple, with $a\ge1$. Let $C$ be a set of continuous representations of the same group and rank. Fix a prolongation at $w\mid\ell$, compatible period data and a bridge for which every $\rho\in C$ is crystalline with the same jumps $J_\tau$ at every induced embedding $\tau$. Fix $\sigma\in W_{E_w}$ and put
\[
 s_j(\rho)=\iota^{-1}\operatorname{tr}\big(\mathrm{WD}_\iota(\rho_w)(\sigma)^j\big)\qquad(j\ge0).
\]
Suppose $S_j\in\bar{\Q}_\ell$ are given and for every $N\ge0$ and $\epsilon>0$ there are a finite $H\subset C$ and weights $b_\rho\in\bar{\Q}_\ell$ with
\begin{equation}\label{eq:p9c1-joint}
 \begin{aligned}
 \left|\operatorname{tr}r(g)-\sum_{\rho\in H}b_\rho\operatorname{tr}\rho(g)\right|_\ell
 &\le\epsilon &&(g\in G_E),\\
 \left|S_j-\sum_{\rho\in H}b_\rho s_j(\rho)\right|_\ell
 &\le\epsilon &&(0\le j\le N).
 \end{aligned}
\end{equation}
Then $r_w$ is crystalline and de Rham with jumps $J_\tau$ at each induced embedding, and its Weil trace at $\sigma^j$ is $\iota(S_j)$ for every $j\ge0$.
\end{proposition}
\begin{proof}[Outline]
Use \ref{B9.35a} for crystallinity and the element of \ref{B9.35ab} to encode jumps as powers of an infinite-order scalar. Recover the jumps from limiting power sums by Newton's identities and pass the weighted Weil traces to the limit using \ref{B9.35a}. The complete proof is in \S\ref{prf:B9.36}.
\end{proof}

\section{Descent to general linear groups}

\begin{lemma}[Unitary character with prescribed infinity type]\label{B9.36c}
Let $E/F$ be CM and unramified at finite places, with quadratic idèle class character $\eta_{E/F}$. Given a finite set $S$ of finite places of $E$ and odd integers $e_v$ at its complex places, there is a unitary Hecke character $\mu$ of $E$ restricting to $\eta_{E/F}$ on the idèles of $F$, unramified on $S$, with
\[
 \mu_w=(-1)^{\operatorname{ord}_w}\quad(w\text{ nonsplit finite}),\qquad
 \mu_v(z)=(z/|z|)^{e_v}.
\]
The exponents are relative to the chosen embeddings at the complex places.
\end{lemma}
\begin{proof}[Outline]
Prescribe the infinity and quadratic characters on the two idèle subgroups, using the auxiliary congruence and odd parity for agreement. Extend by divisibility of the unit circle and read the nonsplit values from units and a uniformizer. The complete proof is in \S\ref{prf:B9.36c}.
\end{proof}

\begin{lemma}[The terms of the split Hecke test]\label{B9.36b}
Let $E/F$ and $n$ be as in \ref{B6.04b}. Let $\pi$ be a regular algebraic conjugate self-dual cusp form on $\GL_n(\mathbf A_E)$, spherical at every nonsplit finite place. Choose $\mu$ as in \ref{B6.04b}. There exist a weight $\lambda$ and a section $\Sigma$ of finite places over $F$, with a finite set $Z\subset\Sigma$ of split places containing representatives of all pairs where $\pi$ or $\mu$ ramifies. For $w\notin Z\cup\bar Z$, let $B(w)$ be the Satake multiset of $\pi_w$. At $w\in Z$ there are tests $f_w$ with $\operatorname{tr}\pi_w(f_w)\ne0$, one with regular support.

On a finite set $T\subset\Sigma$ of split places disjoint from $Z$ take spherical Hecke polynomials $e_{1,w},e_{2,w}$. For any family of multisets $X$ put
\[
 v(X)=\prod_{w\in T}e_{1,w}(X(w))e_{2,w}(X(\bar w)^{-1}),
\]
where inversion is entrywise. Let $f_k$ use the fixed tests at $Z$ and the $k$-th power of the product of the two Hecke polynomials at each place in $T$. Use the unit elsewhere. Denote by $S_i(k)$ the stable terms of \ref{B6.04b} before their signs. There are finite lists, chosen before $T$ and the polynomials, with the following properties.

The normalized unitary trace is $\sum_h A_h v(B_h)^k$, where $h$ ranges over automorphic summands unramified at every nonsplit place and $B_h$ is their split Satake family. On a proper block with $a_i b_i>0$ the list consists of pairs of cuspidal lists of ranks $a_i,b_i$, with Satake unions $X_1,X_2$. At a good place $w$ with uniformizer $\varpi_w$ put $X(w)=\mu_w(\varpi_w)^{b_i}X_1(w)\sqcup\mu_w(\varpi_w)^{a_i}X_2(w)$. If $|v(X)|<\rho$ for every entry, with $\rho>0$, then
\begin{equation}\label{eq:p9c1-proper-block}
 |S_i(k)|\le C(1+k)^N\rho^k\qquad(k\ge0).
\end{equation}
For $(n,0)$ the list consists of cuspidal lists of total rank $n$ omitting $\pi$. Under the same strict inequality on that list,
\begin{equation}\label{eq:p9c1-end-block}
 |S_{n,0}(k)-Pv(B)^k|\le C(1+k)^N\rho^k,\qquad
 P=c\prod_{w\in Z}\operatorname{tr}\pi_w(f_w),\qquad c\ne0.
\end{equation}
Here $c$ is fixed before the varying tests, whereas $C,N$ may depend on the tests and $\rho$. The $(0,n)$ term equals $S_{n,0}(k)$. The competing rank-$n$ lists, including the twisted unions from proper blocks, have regular infinitesimal characters and inversion-invariant Satake unions at all but finitely many nonsplit places.
\end{lemma}
\begin{proof}[Outline]
Use the local matching in \ref{B6.04b} and finiteness in \ref{B6.04u} to form the Hecke lists. Spectral separation and the trace estimates of \ref{B6.04c} bound competitors, leaving the same cuspidal contribution at both ends. Nonsplit tests enforce inversion invariance. The complete proof is in \S\ref{prf:B9.36b}.
\end{proof}

\begin{lemma}[Weak descent]\label{B9.36a}
Suppose $E/F$ is CM and unramified at finite places. A regular algebraic conjugate self-dual cuspidal $\pi$ of positive rank, spherical at every nonsplit finite place, has a definite-unitary automorphic match at almost every split place, unramified at every nonsplit finite place.
\end{lemma}
\begin{proof}[Outline]
Choose \(\mu\) by \ref{B9.36c} and use split Hecke tests on the lists of \ref{B9.36b} to isolate the nonzero end contributions in \ref{B6.04b}. The Jacquet–Shalika bound makes the factors of a matching list unitary, and Rankin–Selberg orthogonality forces one to be \(\pi\). The complete proof is in \S\ref{prf:B9.36a}.
\end{proof}

\begin{proposition}[Unitary descent]\label{B9.37}
For $\chi=\chi_{E/F}^a$, assume $E$ contains an imaginary quadratic subfield and $[F:\Q]$ is even. Suppose every nonsplit finite place lies away from $\ell$ and is unramified over $F$. Assume $\pi$ is spherical at these places. Given a finite set $Z$, there is a definite-unitary automorphic $U$ matching $\pi$ at split places, with weights giving its infinitesimal characters. At nonsplit places $U$ is unramified off a finite nonsplit set $T$ disjoint from $Z$.
\end{proposition}
\begin{proof}[Outline]
Take the weak match by \ref{B9.36a} and apply Labesse’s base change. The Jacquet–Shalika bound forces every Speh length to be one, and the split Rankin–Selberg argument of \ref{B9.36a} identifies the transfer with \(\pi\), giving full split and infinitesimal matching. The complete proof is in \S\ref{prf:B9.37}.
\end{proof}

\begin{lemma}[A ramified place for quadratic base change]\label{B9.37a}
Let $M/K$ be quadratic with $K$ totally complex and let $\pi$ be cuspidal of rank $a\ge1$ over $K$. If $M/K$ ramifies over a spherical place of $\pi$, there is a cuspidal lift. At $w\mid v$ it has multiset $\{\alpha^{f(w/v)}:\alpha\in B_v\}$ whenever $B_v$ exists, with residue degree $f(w/v)$. Infinitesimal characters pull back by restriction of embeddings. Local components agree under continuous completion isomorphisms extending $K\to M$.
\end{lemma}
\begin{proof}
Twisting by the quadratic character changes the unramified cuspidal support at the designated place, so $\pi$ has no quadratic self-twist. Prime-degree cyclic base change gives the cuspidal lift, and the local norm raises unramified parameters to the residue degree \cite[Ch.~III, Thms.~4.2 and 5.1]{arthur-clozel-1989}.
\end{proof}

\begin{proposition}[Spherical theorem]\label{B9.38}
Call a continuous semisimple $r:G_E\to\GL_a(\bar{\Q}_\ell)$ weakly attached to $\pi$ if outside a finite set of $w$ there is a Satake multiset $B_w$ such that inertia acts trivially and every inverse arithmetic Frobenius has polynomial $P_{w,B_w}$ after applying $\iota$. Every polarized input has a weakly attached $r$. For every such $r$ and every Satake multiset $B$ at a spherical $w\nmid\ell$, inertia acts trivially and the inverse-Frobenius polynomial is $P_{w,B}$. At a spherical $w\mid\ell$, $r_w$ is crystalline and de Rham with the prescribed jumps. For every compatible period datum and bridge its Weil polynomial at $h\in W_{E_w}$ is $\prod_{\alpha\in A_a(w,B)}(X-\alpha^{-\deg h})$ for every Satake multiset $B$ there.
\end{proposition}
\begin{proof}[Outline]
For quadratic polarization over split fields, applies \ref{B9.37}, \ref{B9.35}, and \ref{B9.36} in odd rank and the even-rank case, descending via \ref{B9.37a} and index-two patching. Reduces general polarization by an algebraic Hecke twist from \ref{B9.07a}. The complete proof is in \S\ref{prf:B9.38}.
\end{proof}

\begin{lemma}[Exponents of unlinked characters]\label{B9.38a}
Let $(\pi,\chi)$ be a polarized input of rank $a$, $K=E_w$, $\nu_K=|\cdot|_K$, and $s=(1-a)/2$. Suppose the cuspidal support of $\sigma^*=\pi_w\otimes|\det|_K^s$ consists of unlinked smooth characters $\xi_1,\ldots,\xi_a$. For characters, unlinking means $\xi_i\ne\xi_j\nu_K$ for every $i,j$. With the weight of \ref{B9.04} there are real numbers $e_j$ such that
\begin{equation}\label{eq:p9c1-unlinked-exponents}
 |e_j|<\tfrac12,\qquad
 |\xi_j(z)|=\nu_K(z)^{e_j-w_{\rm gt}/2}\quad(z\in K^\times).
\end{equation}
\end{lemma}
\begin{proof}[Outline]
The central-character exponent \(t=(a-1-w_{\rm gt})/2\) makes \(\sigma^*|\det|^{w_{\rm gt}/2}\) unitarizable and generic. In the generic unitary dual, blocks have twists \(|e|<1/2\); unlinking forces length one, characters force rank one, and removing the twist gives the bound. The complete proof is in \S\ref{prf:B9.38a}.
\end{proof}

\begin{lemma}[The support polynomial in the local functional equation]\label{B9.38e}
Let \(k\) be a nonarchimedean local field of characteristic zero, \(q\) its residue cardinality, and \(\varpi\) a uniformizer. Let \(\tau\) be an irreducible generic representation of \(\GL_n(k)\), \(S\) its multiset of supercuspidal support for normalized induction, and \(\lambda\) a smooth character of \(k^\times\). Define
\[
 D_k(S,\lambda;T)=
 \prod_{\substack{\mu\in S,\ \operatorname{rank}\mu=1\\
                   (\mu\lambda)|_{\mathcal O_k^\times}=1}}
              \bigl(1-(\mu\lambda)(\varpi)T\bigr).
\]
For any nontrivial additive character of \(k\), the standard local \(\gamma\)-factor of \(\tau\otimes(\lambda\circ\det)\), in the variable \(T=q^{-s}\), has the form
\[
                   cT^d\,\frac{D_k(S,\lambda;T)}
                                        {D_k(S,\lambda;qT)}
 \qquad(c\in\C^\times,\ d\in\Z).
\]
\end{lemma}
\begin{proof}[Outline]
Multiply the local gamma factors over the twisted support. Higher-rank supercuspidals and ramified characters give epsilon monomials, and Tate's functional equation gives the quotient factors for unramified characters. The complete proof is in \S\ref{prf:B9.38e}.
\end{proof}

\begin{lemma}[A spherical lift fixing a prescribed Weil element]\label{B9.38b}
Assume the hypotheses of \ref{B9.38a} and let $r$ be weakly attached to $\pi$. Fix a prolongation of $w$ to $\Omega$ and the reciprocity map $b_K:W_K\to K^\times$ of \ref{B9.50}. For each $g\in W_K$ of nonzero degree there exist a finite CM extension $M/E$ in $\Omega$, fixed pointwise by $g$, and a polarized cuspidal lift $\Pi$ of $\pi$. Then $r|_{G_M}$ is weakly attached to $\Pi$. At the place $u$ determined by the prolongation put $L=M_u$. The component $\Pi_u$ has a Satake multiset $B$ such that for every $h\in W_L\subset W_K$,
\[
 \prod_{\alpha\in A_a(u,B)}(X-\alpha^{-\deg_L h})
 =\prod_{j=1}^a(X-\xi_j(b_K(h))).
\]
In particular the formula applies to $g$.
\end{lemma}
\begin{proof}[Outline]
Use reciprocity and Grunwald–Wang for a CM tower that unramifies the characters and excludes self-twists by trace nonvanishing. Compare functional equations by \ref{B9.38e} for support multiplicities and use the window of \ref{B9.38a} for spherical induction. The complete proof is in \S\ref{prf:B9.38b}.
\end{proof}

\begin{proposition}[Comparison on unlinked character support]\label{B9.38c}
With $r,\sigma^*,\xi_j$ as in \ref{B9.38b}, put $Y=\bigoplus_{j=1}^a(\xi_j\circ b_K,0)$. If $w\nmid\ell$, then $\mathrm{WD}_\iota(r_w)=Y$. If $w\mid\ell$, then $r_w$ is potentially semistable and de Rham with the jumps of \ref{B9.04} at every induced embedding, and its Frobenius-semisimplified period parameter for every compatible datum and bridge equals $Y$. In both cases $R(\sigma^*)=Y$ for every map $R$ of \ref{B9.47}.
\end{proposition}
\begin{proof}[Outline]
Use \ref{B9.38b} and \ref{B9.38} to compare Weil polynomials and descend zero monodromy and the jumps. Recover inertia traces by Cayley–Hamilton and identify the Weil action by Brauer–Nesbitt and the parameter by \ref{B9.50} and \ref{B9.47}. The complete proof is in \S\ref{prf:B9.38c}.
\end{proof}

\begin{lemma}[An unramified extension for depth-zero support]\label{B9.38d}
Let $(\pi,\chi),r,w$ be as in \ref{B9.38} and fix a prolongation. Suppose $K=E_w\simeq\Q_p$ continuously. Suppose $\pi_w$ has first principal congruence invariants and unlinked cuspidal support, with its terms of rank at least two pairwise inequivalent up to unramified twist. Given any $R$ of \ref{B9.47}, there is a CM tower $M/E$ with a polarized cuspidal lift $\Pi$ weakly attached to $r|_{G_M}$. Its completion $L=M_u$ at the prolongation is unramified of degree $f=a!$ over $K$. The normalized component $\Pi_u^*=\Pi_u|\det|_L^{(1-a)/2}$ has unlinked character support. There is an upper map $R'$ of \ref{B9.47} such that for every $t\in I_K=I_L$,
\[
 \det(X-\rho_{R'(\Pi_u^*)}(t))
 =\det(X-\rho_{R(\pi_w|\det|_K^{(1-a)/2})}(t)).
\]
The infinitesimal character of $\Pi$ at an embedding of $M$ is that of $\pi$ at its restriction to $E$.
\end{lemma}
\begin{proof}[Outline]
Use a cyclic tower as in \ref{B9.38b} and split Green orbits by Shintani matching with the labels of \ref{B7.31} and \ref{B9.51}. Compute inertia by \ref{B9.49} and \ref{B9.47} and compare with the lifted characters by \ref{B9.46}. The complete proof is in \S\ref{prf:B9.38d}.
\end{proof}

\section{Two theta lifts and even rank}

\begin{proposition}[Stable range]\label{B9.39}
Let $E/F$ be CM unramified at finite places with an imaginary quadratic subfield, and let its inert places have odd residue characteristic and norm greater than three. Choose a complex embedding of $E$ above each real place of $F$ and an $F$-embedding $E\hookrightarrow F_v$ for each $v$ in a designated finite set $P$ of split finite places. Use an additive character trivial on $F$ and equal to $\exp(-2\pi i\operatorname{Tr}(x))$ at infinity. Require it to be unramified at $P$ and the inert places. Equip the definite Hermitian source $X$ with matrix $1_n$, $n\ge1$, and a self-dual lattice. Choose compatible archimedean Cartan data and split bases at $P$. There is a bound $d_0$ for definite skew-Hermitian targets $W$ with matrix $\delta1_d$, where $\bar\delta=-\delta$ and $\delta\ne0$. Equip the target with a lattice self-dual for the trace pairing and choose compatible archimedean Cartan data and split bases at $P$.

For $d\ge d_0$ use splitting characters $\chi_X,\chi_W$ restricting to the quadratic powers $n,d$ with infinity exponents $m_X,m_W$. A definite-unitary automorphic $U$ with $2\lambda_{v,k}\ge d+m_W(v)$ for $0\le k<n$ has a lift $U'$ of weights
\[
 2\lambda'_{v,k}=n+m_X(v)+
 \begin{cases}
 2\lambda_{v,k}-d-m_W(v)& k<n,\\
 0& n\le k<d.
 \end{cases}
\]
For every finite place $v$ and every $F$-embedding $t:E\hookrightarrow F_v$, the components $\tau,\tau'$ of $U,U'$ at the place determined by $t$ admit a nonzero pairing
$\beta:\mathcal S(M_{n,d}(F_v))\times\tau\to\tau'$.
It is linear in the first slot and conjugate-linear in the second, with $\beta(\omega(g,h)\Phi,\tau(g)u)=\tau'(h)\beta(\Phi,u)$ for the split Weil action $\omega$.
\end{proposition}
\begin{proof}[Outline]
Use the doubling embedding and determinant bounds to control height-layer volumes and get integrability at sufficiently high powers. Pad the partition and restore the determinant shift, apply \ref{B9.40} for nonvanishing, and project to obtain the split local pairings. The complete proof is in \S\ref{prf:B9.39}.
\end{proof}

\begin{lemma}[Integral criterion]\label{B9.40}
In \ref{B9.39}, write $[G]=\U(X)(F)\backslash\U(X)(\mathbf A_F)$. Let $C_\Phi$ be the Weil matrix coefficient of a polynomial times finite Schwartz vector $\Phi$, and $M_\varphi$ the right-translation matrix coefficient of a continuous form $\varphi$ on $[G]$. Suppose the doubling height satisfies $\int_{\U(X)(\mathbf A_F)}h^t<\infty$ for $t\ge t_0$, with $t_0<d/2$. Then $I=\int C_\Phi\overline{M_\varphi}$ converges absolutely and vanishes for a zero lift. If $d\ge2n$ and intertwiners from one finite-dimensional irreducible archimedean module take a common vector to nonzero polynomial $f_0$ and nonzero $\varphi$ in irreducible automorphic $U$, some finite Schwartz data give a nonzero lift.
\end{lemma}
\begin{proof}[Outline]
Apply \ref{B8.12} and unfold to a matrix-coefficient integral, using the height bound to reach the theta point. Use translated finite vacua and Schur orthogonality so the positive integral dominates the height tail. The complete proof is in \S\ref{prf:B9.40}.
\end{proof}

\begin{lemma}[Split spherical parameters of a lift]\label{B9.40a}
At a split place $w$, suppose $U,U'$ of ranks $n,n+t$ have the local pairing of \ref{B9.39}, with $t>n$ and unramified unitary splitting characters. Write $x=\chi_X(w)$, $y=\chi_W(w)$, $q=q_w$. A Satake multiset $B$ of $U$ gives the Satake multiset
\[
 (xy^{-1})B\sqcup\{xq^{(t-1)/2-j}:0\le j<t\}
\]
of $U'$ at the same place.
\end{lemma}
\begin{proof}[Outline]
Filter the split Schwartz model by rank and use strict contraction to force the pairing onto the full-rank orbit. Apply Frobenius reciprocity at its stabilizer and take the contragredient to obtain the twisted Satake multiset and the determinant-character string. The complete proof is in \S\ref{prf:B9.40a}.
\end{proof}

\begin{proposition}[Two lifts]\label{B9.41}
Under the field conditions of \ref{B9.39}, with $[F:\Q]$ even and inert $v\nmid\mathfrak D_{F/\Q}$, take definite-unitary automorphic $U$ of positive even rank $n$ with weights $\lambda_\tau$ at every embedding and a split place $w_0$. For some odd $t$ there are two lifts $U_i$ of rank $n+t$ and algebraic characters $\kappa,\psi_{ij}$ with disjoint $\psi$-families and $\kappa,\psi_{1j}$ unramified at $w_0$. At split $w$ where $\kappa$ and the $i$-th family are unramified, each Satake multiset $B$ of $U$ admits $B_i$ with
\[
 A_{n+t}(w,B_i)
 =\kappa(w)^{-1}A_n(w,B)\sqcup\{\psi_{ij}(w)\}_{0\le j<t}.
\]
For infinity exponents $m_\tau,p_{ij,\tau}$, the jumps of $U_i$ are $(J(\lambda_\tau)+m_\tau)\sqcup\{-p_{ij,\tau}\}_j$.
\end{proposition}
\begin{proof}[Outline]
Choose odd excess rank by \ref{B9.39} and an unramified character by \ref{B9.36c}. Use two pairs of powers with the same quotient in \ref{B9.39}, rescale \ref{B9.40a} to get the two added families, and use infinity exponents to separate them and compute jumps. The complete proof is in \S\ref{prf:B9.41}.
\end{proof}

\begin{lemma}[Frobenius identity]\label{B9.42}
Let $Q$ be split places with dense inverse Frobenius away from $\ell$ after every finite removal. Suppose continuous equal-rank $r_i$ and continuous characters $\chi_{ij}$ are cofinitely unramified there, with inverse-Frobenius roots of $r_i$ a common multiset plus its own family's values. Then for all $g\in G_E$,
\[
 \det(X-r_1(g))\prod_j(X-\chi_{2j}(g))
 =\det(X-r_2(g))\prod_j(X-\chi_{1j}(g)).
\]
\end{lemma}
\begin{proof}
At each inverse Frobenius in the dense set both sides multiply the common polynomial by both line products. Apply coefficientwise continuity.
\end{proof}

\begin{lemma}[Line subtraction]\label{B9.43}
Under \ref{B9.42} with semisimple $r_i$ and disjoint families, $r_1\simeq s\oplus\bigoplus_j\chi_{1j}$. For line roots $c_j$ and jumps $m_{j,\tau}$, multisets $S\sqcup\{c_j\}_j$ and $J_\tau\sqcup\{m_{j,\tau}\}_j$ for $r_1$ become $S,J_\tau$ for $s$. At $w\mid\ell$ suppose $r_1$ and the lines are crystalline and their Weil polynomials at each $g$ use their respective roots raised to $-\deg g$. Then $s$ is crystalline with polynomial $\prod_{\alpha\in S}(X-\alpha^{-\deg g})$.
\end{lemma}
\begin{proof}
Brauer--Nesbitt identifies the enlarged sums in \ref{B9.42}. On the second side the first-family lines project trivially to the second-family lines, so on the first side they occur in $r_1$ with at least their listed multiplicities. Take their semisimple complement $s$. A direct summand of a crystalline representation is crystalline. Cancel the line factors in the Frobenius and Weil polynomials and subtract their jumps by additivity.
\end{proof}

\begin{proposition}[One place]\label{B9.44}
Assume the density of \ref{B9.42} and odd-rank spherical compatibility on $Q$. Suppose there are two lifts of rank-$n$ $U$ as in \ref{B9.41} at $w_0\in Q$. There is a rank-$n$ representation $r$ weakly compatible on $Q$. At all but finitely many $w\in Q$ away from $\ell$, it is unramified and has polynomial $P_{w,B}$ for every Satake multiset $B$ there. For every Satake multiset $B$ of $U$ at $w_0$ the following local assertions hold. If $w_0\nmid\ell$, $r$ is unramified there with polynomial $P_{w_0,B}$. If $w_0\mid\ell$, it is crystalline de Rham with jumps $J(\lambda_{\iota\circ\tau})$ at induced $\tau$ and the Weil polynomials of \ref{B9.43} for $S=A_n(w_0,B)$ and every compatible period datum.
\end{proposition}
\begin{proof}[Outline]
Apply odd-rank compatibility to both lifts, using injectivity of character realization to keep line families disjoint and cofinite Satake unions for \ref{B9.42}. Combining \ref{B9.41} at \(w_0\), \ref{B9.43} subtraction, and \(\kappa_\ell\)-twisting gives \(r\) and compatibility. The complete proof is in \S\ref{prf:B9.44}.
\end{proof}

\begin{proposition}[All places]\label{B9.45}
With \ref{B9.44} available at every place of $Q$, weakly compatible representations exist and all have its local data at every spherical place in $Q$.
\end{proposition}
\begin{proof}
Take the witness of \ref{B9.44} at the target. Remove the exceptional sets of the witness and the given weakly compatible representation. Also remove the finitely many split places where $U$ has no Satake multiset. Use the same multiset for both representations at every remaining place. Density and Brauer--Nesbitt identify them.
\end{proof}



% ---- p9c2 ----
\chapter{A weak local Langlands correspondence}\label{ch:weak-local-langlands}

Let \(K=E_w\) be a finite completion of a number field \(E\), with residue field \(\F_q\) of characteristic \(p\). Fix a valuation ring \(A\) over \(w\) in an algebraic closure \(\Omega\) of \(E\). The Weil group \(W_K\) has inertia \(I_K\) and wild inertia \(P_K\). The degree map \(\deg:W_K\to\Z\) sends arithmetic Frobenius to \(1\). Write \(\mathcal I(K)\) for the smooth irreducible complex classes \(x\) of \(\GL_{d(x)}(K)\), \(d(x)\ge0\). Cuspidal means supercuspidal of positive rank. In rank \(n\), put \(K_0=\GL_n(\mathcal O_K)\) and \(K_1=\ker(K_0\to\GL_n(\F_q))\).

The complex Weil--Deligne classes \(V=(\rho_V,N_V)\) in \(\mathcal W(K)\) have elementwise semisimple Weil action and open inertia kernel. The nilpotent \(N_V\) satisfies \(\rho_V(\sigma)N_V\rho_V(\sigma)^{-1}=q^{\deg\sigma}N_V\). Irreducibility refers to \(\rho_V\), and tameness to trivial \(P_K\)-action.

Put \(\nu_K=|\cdot|_K\) and \(\nu_W=q^{\deg}\). Twists of \(x\) use the determinant. Write \(V(s)=V\otimes\nu_W^s\) for \(s\in\R\). The self-dual \(\operatorname{Sp}_m\) has one monodromy block and Weil characters \(\nu_W^{j-(m-1)/2}\), \(0\le j<m\).

\section{The existence theorem and its laws}

\begin{theorem}[Existence]\label{B9.46}
There is a map \(\operatorname{rec}_T:\mathcal I(K)\to\mathcal W(K)\) satisfying Definition~\ref{B9.47}.
\end{theorem}

\begin{proof}[Outline]
Extend \ref{B9.53} by \ref{B9.54}. The cuspidal product gives \ref{B9.47}, block operations give \ref{B9.48}, indecomposable blocks recover good segments, and the compact picture and exactness of invariants reduce tameness and the type law to cuspidal support. The complete proof is in \S\ref{prf:B9.46}.
\end{proof}

\begin{definition}[The weak correspondence]\label{B9.47}
A map \(R:\mathcal I(K)\to\mathcal W(K)\) is a \emph{weak local Langlands correspondence} if there exists a local reciprocity map \(a:W_K\to K^\times\) as in Theorem~\ref{B9.50} for which the following laws hold. For every \(x\), we have \(\dim R(x)=d(x)\) and \(R(x\nu_K^s)=R(x)(s)\) for all \(s\in\R\). In rank one, \(R(\chi)=(\chi\circ a,0)\) for every smooth character \(\chi\) of \(K^\times\). For the Langlands multisegment \(\mathfrak m\) of \(x\),
\[
 R(x)=\bigoplus_{[u,m]\in\mathfrak m}R(u)((m-1)/2)\otimes\operatorname{Sp}_m.
 \tag{2}\label{eq:p9c2-segments}
\]
Cuspidal parameters are irreducible. For every \(x\), \(|\det\rho_{R(x)}|=|\omega_x\circ a|\). For good \(x\) and smooth \(\chi\),
\[
 R(x\otimes\chi)=R(x)\otimes(\chi\circ a),\qquad
 \det\rho_{R(x)}=\omega_x\circ a,\qquad R(x^\vee)=R(x)^\vee.
\]
If \(R(x)=R(y)\) and \(y\) is good, then \(x=y\). Good classes with nonzero \(K_1\)-invariants have tame parameters. The map satisfies the type law of \ref{B9.49}.
\end{definition}

The full correspondence combines Scholze's cuspidal bijection \cite[Thm.~1.2(c--d)]{scholze-2013} with Rodier's multisegment extension \cite[\S4.4]{rodier-1982}. 

\begin{definition}[Good cuspidal classes and good support]\label{B9.48}
Depth zero means nonzero \(K_1\)-invariants. Let \(\mathcal C_K\subseteq\mathcal I(K)\) consist of the rank-one classes and, when \(K\) is isomorphic to \(\Q_p\) as a topological field, all classes \(u_0\otimes(\chi\circ\det)\) with \(u_0\) of depth zero and \(\chi\) a smooth character of \(K^\times\). A class is \emph{good}, or has \emph{good cuspidal support}, if every term of its cuspidal support belongs to \(\mathcal C_K\). For a cuspidal class \(u\), this means \(u\in\mathcal C_K\), since its support is \(\{u\}\).
\end{definition}

\begin{definition}[The depth-zero type law]\label{B9.49}
Fix a continuous isomorphism \(i:\Q_p\simeq K\) and \(n\ge1\). Take \(E'/E\) in \(\Omega\) with place \(w'\) induced by \(A\), \(p\notin(w')^2\) and \(Nw'=p^n\). Let \(a'\) be reciprocity at \(w'\) and \(\mathrm{red}:\mathcal O_{E'}\twoheadrightarrow k_n\) reduction with kernel \(w'\), where \(k_n\) is the residue field of order \(p^n\). Take a character \(\theta:k_n^\times\to\C^\times\) of Frobenius orbit length \(n\) and a cuspidal irreducible \(\tau\) of \(\GL_n(\F_p)\) with
\[
 \operatorname{tr}\tau(g)=(-1)^{n-1}\sum_{j=0}^{n-1}\theta(z^{p^j})
 \tag{3}
\]
when \(g\) has irreducible characteristic polynomial with root \(z\in k_n^\times\). If \(\pi\in\mathcal I(K)\) of rank \(n\) contains the inflation of \(\tau\) via \(i\) on \(\GL_n(\Z_p)\), then \(N_{R(\pi)}=0\). For \(t\in I_{E'_{w'}}\), let \(t_E\) be its image in \(I_K\). For \(b\in\mathcal O_{E'}\) with \(|a'(t)-b|_{w'}<1\) and \(z=\mathrm{red}(b)\in k_n^\times\),
\[
 \det(X-\rho_{R(\pi)}(t_E))=\prod_{j=0}^{n-1}(X-\theta(z^{p^j})).
 \tag{4}
\]
\end{definition}

Formula (3) is Green's elliptic character formula \cite{green-1955}. The law uses the cuspidal case over \(\Q_p\) of the depth-zero inertia formula \cite[Addendum (A.1)]{silberger-zink-2008}.

\section{Cuspidal parameters}

\begin{theorem}[Local reciprocity]\label{B9.50}
There is a unique surjective homomorphism \(a:W_K\to K^\times\), continuous on \(I_K\), with \(\operatorname{ord}_K(a(\sigma))=-\deg\sigma\) for all \(\sigma\). Every smooth Weil character factors as \(\chi\circ a\) with \(\chi\) smooth. For every global Artin map \(\mathrm{Art}_E:\mathbf A_E^\times\to\Gal(\Omega/E)^{\mathrm{ab}}\) and \(\sigma\in W_K\),
\[
 \mathrm{Art}_E(\operatorname{loc}_w(a(\sigma)))=\bar\sigma,
\]
where \(\operatorname{loc}_w\) is the local id\`ele inclusion and \(\bar\sigma\) is the abelianized image.
\end{theorem}

Local reciprocity gives existence in this Frobenius convention \cite[Ch.~XIII, \S4]{serre-1979}, with local--global compatibility as in \cite[Ch.~VI, \S5]{neukirch-1999}. For \(u\ne1\) in \(K^\times\), choose an open subgroup of finite index excluding \(u\). Characters of the finite quotient by this subgroup separate its elements. Choose one nontrivial on \(u\) and extend it from the single place \(w\) to a finite-order idèle class character by Grunwald--Wang \cite[Thm.~2.2]{wang-2015}. Take the finite abelian extension cut out by this character. In that extension the global Artin image of \(\operatorname{loc}_w(u)\) is the nontrivial local Artin image of \(u\). Thus \(\mathrm{Art}_E\circ\operatorname{loc}_w\) is injective, and the Artin identity gives uniqueness.

\begin{theorem}[Parameters on the good cuspidal classes]\label{B9.51}
Fix \(a\). Over \(K\simeq\Q_p\) continuously, each depth-zero cuspidal \(u\) has a unique irreducible tame parameter \(V_u\) of dimension \(d(u)\) satisfying the type law for \(u\) and \(\det\rho_{V_u}=\omega_u\circ a\). The map \(u\mapsto V_u\) is injective. The formulas
\[
 r_0(\chi)=(\chi\circ a,0),\qquad
 r_0(u_0\otimes\chi)=V_{u_0}\otimes(\chi\circ a)
\]
define an assignment on good cuspidals that commutes with smooth twists. Here \(\chi\) is smooth. In the second formula \(K\simeq\Q_p\) continuously and \(u_0\) is depth-zero cuspidal.
\end{theorem}

\begin{proof}[Outline]
Compact induction and Green's formula \ref{B7.31} give the type and central character. In \ref{B9.49}, reciprocity gives the inertia orbit, its Frobenius determinant fixes the cycle, Newton's identities give the type law, and the label and central character track twists. The complete proof is in \S\ref{prf:B9.51}.
\end{proof}

\begin{lemma}[Auxiliary irreducible parameters]\label{B9.52}
For \(n\ge2\) there is an irreducible \(J_n\) of dimension \(n\) with unitary determinant. If \(K\simeq\Q_p\), one can require \(J_n\ne V\otimes\xi\) for every tame parameter \(V\) and smooth Weil character \(\xi\).
\end{lemma}

\begin{proof}[Outline]
Reciprocity on principal units and the trace pairing give a full Frobenius orbit of inertia characters. A determinant-one Frobenius cycle gives a finite-image irreducible parameter, and its nonscalar wild inertia action rules out all character twists of tame parameters. The complete proof is in \S\ref{prf:B9.52}.
\end{proof}

\begin{proposition}[The assignment on all cuspidal classes]\label{B9.53}
Fix \(J_n\) as in \ref{B9.52}. For cuspidal \(u\) of rank \(n\), define \(e(u)\in\R\) by \(|\omega_u|=\nu_K^{ne(u)}\). The assignment
\[
 r(u)=
 \begin{cases}
 r_0(u),&u\in\mathcal C_K,\\
 J_{d(u)}(e(u)),&\text{otherwise}
 \end{cases}
 \tag{6}
\]
satisfies the cuspidal laws of Definition~\ref{B9.47}. In particular, \(r(u)=r(v)\) with \(v\) good forces \(u=v\).
\end{proposition}

\begin{proof}[Outline]
Real twisting and \ref{B9.52} give the second branch its laws, and wild inertia separates the branches. On good classes, \ref{B9.51} gives the higher-rank laws and injectivity, inversion gives duality, and reciprocity gives rank-one tameness and injectivity. The complete proof is in \S\ref{prf:B9.53}.
\end{proof}

\section{Multisegments and uniqueness}

\begin{theorem}[Langlands extension]\label{B9.54}
Each \(x\) has a unique Langlands multisegment. A segment \([u,m]\) consists of \(u\nu_K^j\), \(0\le j<m\), with \(u\) cuspidal and \(m\ge1\). Formula~\eqref{eq:p9c2-segments}, using \(r(u)\) for each first term, extends any cuspidal assignment \(r\). The empty multisegment gives zero.
\end{theorem}

In Rodier's parametrization \cite[Thm.~3]{rodier-1982}, take the unique irreducible quotient for each segment by normalized induction in increasing exponent order. Order segments so none precedes a later one, where on an integer twist line \([a,b]\) precedes \([c,d]\) if \(a<c\le b+1\) and \(b<d\). The normalized induction of the segment quotients has unique irreducible quotient \(x\). Zelevinsky gives the submodule form \cite[Thm.~6.1]{zelevinsky-1980}. Permuting the segments permutes the direct summands. A cuspidal \(u\) has multisegment \(\{[u,1]\}\), so \(\operatorname{Sp}_1=\mathbf1\) gives \(R(u)=r(u)\).

\begin{theorem}[Uniqueness on good representations]\label{B9.55}
Let \(R,R'\) satisfy \ref{B9.47}. Then \(R(x)=R'(x)\) for every good \(x\).
\end{theorem}

\begin{proof}
By \ref{B9.50} the maps agree on characters. Over \(K\simeq\Q_p\) continuously they agree on depth-zero cuspidals by \ref{B9.51} and on their twists by \ref{B9.47}. Apply (2) to the first terms of good \(x\).
\end{proof}

A cuspidal support \(s\) is unlinked if \(u\ne v\nu_K\) for every \(u,v\in s\). Two terms are inertially equivalent if one is an unramified character twist of the other.

\begin{lemma}[Support under the local tests]\label{B9.55a}
Let \(x\) be generic with \(N_{R(x)}=0\) for a map \(R\) as in \ref{B9.47}. If every Weil commutator has polynomial \((X-1)^{d(x)}\) on \(R(x)\), then \(x\) has unlinked rank-one support. If \(K\simeq\Q_p\) continuously and \(P_K\) acts trivially on \(R(x)\), then \(x\) has unlinked depth-zero support and \(x^{K_1}\ne0\). In this second case, assume also that for each lift \(\gamma\in I_K\) of a topological generator of \(I_K/P_K\), every repeated root \(z\) of \(\det(X-\rho_{R(x)}(\gamma))\) satisfies \(z^{q-1}=1\). Then the support terms of rank at least two are pairwise inertially inequivalent. Both support conclusions give goodness by \ref{B9.48}.
\end{lemma}

\begin{proof}[Outline]
Genericity and zero monodromy give singleton segments, and semisimplicity makes the commutator test force rank one. The finite orbit injection uses \ref{B9.51} and \ref{B9.47} to give depth-zero preimages, and equal inertia polynomials make the repeated-root test force rank one. The complete proof is in \S\ref{prf:B9.55a}.
\end{proof}

In \ref{B9.07} put \(x=\pi_w\nu_K^{(1-d(\pi_w))/2}\) with \(N_{R(x)}=0\). Whenever \ref{B9.55a} gives good support, \ref{B9.55} gives \(R'(x)=R(x)\) for any weak correspondence \(R'\). For unlinked character support, \ref{B9.38c} gives equality with the Galois parameter. Over \(K\simeq\Q_p\) continuously, the wild-inertia and repeated-root tests give \(K_1\)-invariants and unlinked support with pairwise inertially inequivalent terms of rank at least two. The normalizing twist preserves these properties. Apply \ref{B9.38c} on the unramified lift of \ref{B9.38d}. Restriction keeps the monodromy operator and inertia group, so the Galois parameter downstairs has zero monodromy and the inertia polynomials of \(R(x)\).

For any \(x\in\mathcal I(K)\), suppose the inertia polynomials of \(R(x)\) agree with those of a sum of \(d(x)\) Weil characters. For \(s,t\in W_K\), the commutator \([s,t]\) belongs to \(I_K\), and every one-dimensional character takes value \(1\) at \([s,t]\). By \eqref{eq:p9c2-support} and semisimplicity, each commutator acts trivially on each cuspidal parameter. Irreducibility gives rank-one support, so \ref{B9.55} applies.



% ---- appendix ----
\part{Appendix: complete proofs}\label{part:appendix}

This appendix gives the proofs that the main text summarises. Each is headed by the statement it proves.

\section*{Proofs for Chapter~\ref{ch:p1c1}: Five statements and the critical line}

\subsection*{Proof of \ref{B1.14a}}\label{prf:B1.14a}\phantomsection
\begin{proof}
For rational $t>r$ choose $c$ with $v(c)=t$ by taking a root of a
power of $\ell$. The $h$th elementary symmetric function $S_h$ of
$(cH_i)^{12}$ has weight $12hw$ for the full modular group.
Translation by one eliminates fractional powers of $q$. The
modular function $S_h/\Delta^{hw}$ is holomorphic on the upper
half-plane and has pole order at most $hw$ at the cusp. Its
Laurent coefficients lie in $R$, since $\Delta/q$ has integral
coefficients and constant term one. Cancel the negative powers
in succession by subtracting multiples of powers of $j$. The
multipliers lie in $R$ since $j$ has leading coefficient one and
integral coefficients. The remaining holomorphic function on
$X(1)$ is its constant coefficient. Thus $S_h/\Delta^{hw}\in R[j]$.

At $\tau$, the numbers $(cH_i(\tau)/\omega^w)^{12}$ are roots
of a monic polynomial over $R$. A root of negative valuation
would make the leading term the unique term of least valuation.
Thus $v(H_i(\tau)/\omega^w)\ge-t$. Let $t$ decrease to $r$.
\end{proof}

\subsection*{Proof of \ref{B1.14b}}\label{prf:B1.14b}\phantomsection
\begin{proof}
Put $X(P)=\wp_\Lambda(P)/\Omega^2$. By homothety take
$\Lambda=\Lambda_\tau$ and $\omega=\Omega/p$. Since the curve has
CM, $j(\tau)$ is an algebraic integer \cite{shimura-1971}.
Write $P=(a\tau+b)/N$ with $0\le a,b<N$ and
$\zeta=\exp(2\pi i_{\C}b/N)$. In the partial fractions for $\wp$
expand $z/(1-z)^2=\sum_{h\ge1}hz^h$. Then
$(G_2+\wp_{\Lambda_\tau}(P))/p^2$ has constant term zero for $a\ne0$
and $\zeta/(1-\zeta)^2$ for $a=0$. Its coefficient at $q^{u/N}$,
$u>0$, is
\begin{equation}\label{eq:p1c1-torsion-coefficients}
 \sum_{h\mid u}h
 \left(\mathbf1_{u/h\equiv a\ (N)}\zeta^h+
 \mathbf1_{u/h\equiv-a\ (N)}\zeta^{-h}\right).
\end{equation}
When $a=0$ one has $\zeta\ne1$ and
$N/(1-\zeta)=-\sum_{h=0}^{N-1}h\zeta^h$. Thus all coefficients
are integral at $\ell$ if $\ell\nmid N$. The series $G_2/p^2$
has constant term $-1/12$ and coefficient $2\sum_{h\mid j}h$ at
$q^j$, $j>0$. The families $12\wp(P)/p^2$ and
$(\wp(P)-\wp(P'))/p^2$ are permuted in weight two by changes of
period basis at a common level prime to $\ell$. Lemma~\ref{B1.14a}
therefore gives
\begin{equation}\label{eq:p1c1-torsion-integrality}
 12X(P)\in R,\qquad X(P)-X(P')\in R
\end{equation}
for all such nonzero points $P,P'$.

For a nonzero multiplier $\alpha$ of $\Lambda$ put $h=\alpha\bar\alpha$.
Cancel the poles in
$\alpha^2\wp_\Lambda(\alpha z)-\sum_{T\in\ker\alpha}\wp_\Lambda(z+T)$.
The resulting constant is $-\sum_{T\ne0}\wp_\Lambda(T)$ by comparison
at $z=0$. Integrate using $\zeta_\Lambda'=-\wp_\Lambda$ and take the
increment by a period $u$. The coefficients of $\bar u$ cancel
and the coefficients of $u$ give
\[
 (\alpha^2-h)s_2(\Lambda)=\sum_{0\ne T\in\ker\alpha}\wp_\Lambda(T).
\]
When $\ell\nmid h$ the kernel points are killed by $h$, so
\begin{equation}\label{eq:p1c1-isogeny-integrality}
 (\alpha^2-h)U-(\alpha^2-1)X(P)
 =\sum_{T\ne0}\bigl(X(T)-X(P)\bigr)\in R.
\end{equation}
In each case $\alpha^2-h\ne0$, which gives algebraicity of $U$.

For $\ell\ne d$ take $\alpha=w_0$ to get $-2dU+(d+1)X(P)\in R$.
This suffices outside $6d$. At $2$, $v(X(P))\ge-2$ and
$v(d+1)\ge2$ for $d\equiv3\pmod4$, while $v(d+1)=1$ for
$d\equiv1\pmod4$. At $3$ the same equation suffices if $3\mid d+1$.
Otherwise $d\equiv1\pmod3$. Use $\alpha_\pm=1\pm w_0$ of norm
$1+d$. Multiply each relation by $\alpha_\mp^2-1$, subtract and
multiply by $w_0$ to get $-4d^2U\in R$. At $\ell=d$ use
$\alpha=1+w_0$. Now $12$ and $1+d$ are units and
$\alpha^2-(1+d)=2w_0(1+w_0)$ has valuation $1/2$.
\end{proof}

\subsection*{Proof of \ref{B1.14c}}\label{prf:B1.14c}\phantomsection
\begin{proof}
Put $\theta=(-1+w)/2$ if $d\equiv3\pmod4$ and $\theta=w$
otherwise. Then $\tau=\phi_i(\theta)$ is independent of $i$ and
$\Lambda=\Z\tau+\Z$. Set $D=-(\theta-\bar\theta)^2$, equal to
$d$ or $4d$, and $h=v_\ell(D)$. Write $\sigma_i=\phi_i|_F$.

Choose positive integers $A,B$ prime to $\ell$ with
$AY\subset\mathcal O_E$ and $B\ell^a\mathcal O_E\subset Y$.
The inclusions hold locally above $\ell$ by the conductor
exponents, and away from $\ell$ by clearing lattice denominators.
Put
\[
 L=\{(u,v')\in F^2:u\theta+v'\in ADY\},\qquad
 m_0=\ell+1,\qquad N=m_0ADB\ell^a .
\]
The coordinates of $z=(u\theta+v')/D\in AY$ are
\[
 u=D\frac{z-\bar z}{\theta-\bar\theta},\qquad
 v'=D\frac{\theta\bar z-\bar\theta z}{\theta-\bar\theta}.
\]
Both coordinates and their squares divided by $D$ are integral.
Together with $B\ell^a\mathcal O_E\subset Y$ this gives
\begin{equation}\label{eq:p1c1-lattice-coordinates}
 N\mathcal O_F^2\subseteq L\subseteq\mathcal O_F^2,\qquad
 v(\sigma_i(u)),v(\sigma_i(v'))\ge h/2\quad((u,v')\in L).
\end{equation}
Choose the nonzero vector $y_0=(N/m_0,N/m_0)\bmod N$, killed by $m_0$.
\leanref{GlobalTheory.\allowbreak{}RoadE.\allowbreak{}passl}{Development.\allowbreak{}AbelianVarietyTheory.\allowbreak{}Proofs.\allowbreak{}DivisorProofs.\allowbreak{}MemPassScaled}{42--116}

For a lattice $L'$ between $N\mathcal O_F^2$ and $\mathcal O_F^2$,
let $V_{L'}((b_i'),(a_i'))$ be $c_k$ times the sum of
$\prod_i(\sigma_i(u)b_i'+\sigma_i(v')a_i')^{-k}$ over
$(L'\setminus0)/\mathcal O_F^\times$, with the same continuation at
$k=2$. For a multi-index $\nu$ put $|\nu|=\sum_i\nu_i$,
$W_\nu=n+2|\nu|$, and
\[
 D_\nu(L';t,C,T)=
 \left.\partial_z^\nu V_{L'}\bigl((t+z_i(Ct+T))_i,(1+Cz_i)_i\bigr)
 \right|_{z=0}.
\]
If $ADt'=u\theta+v'$, the deformed $i$th coordinate of $t'$,
multiplied by $AD$, is
$\sigma_i(u)(\tau+z_i\bar\tau)+\sigma_i(v')(1+z_i)$.
The substitution respects the real units, so
\begin{equation}\label{eq:p1c1-deformed-substitution}
 \partial^\kappa g(0)=(AD)^n D_\kappa(L;\tau,1,\bar\tau-\tau).
\end{equation}
For $k=2$ the regularized sums give the additional factor
$(AD)^{2f\sigma}$, equal to one at zero. The Fourier expansion justifies continuation and differentiation at zero.
\leanref{GlobalTheory.\allowbreak{}HilbFam.\allowbreak{}iterD\_\allowbreak{}defSum\_\allowbreak{}eq\_\allowbreak{}jet}{Development.\allowbreak{}AbelianVarietyTheory.\allowbreak{}Proofs.\allowbreak{}DivisorProofs.\allowbreak{}DeformedLatticeSumDerivativesEqEisensteinJets}{21--53}

For $\alpha\in\mathrm{GL}_2(\Z/N)$ let $L_\alpha$ be the inverse
image of $(L\bmod N)\alpha$ in $\mathcal O_F^2$, and set
$y_\alpha=y_0\alpha$. Since $L_\alpha=L\widetilde\alpha+
N\mathcal O_F^2$ for an integral lift $\widetilde\alpha$, it satisfies
\eqref{eq:p1c1-lattice-coordinates}. Write $y_\alpha=(a',b')$
and $P_\alpha(t)=(a't+b')/N\bmod\Lambda_t$. This point is nonzero
and killed by $m_0$. Put
\[
 C_\alpha(t)=\frac{G_2(t)+\wp_{\Lambda_t}(P_\alpha(t))}{p^2},
 \qquad H_{\alpha,\nu}(t)=p^{-W_\nu}
 D_\nu(L_\alpha;t,p^2C_\alpha(t),-p).
\]
For $\gamma=\left(\begin{smallmatrix}A'&B'\\C'&D'\end{smallmatrix}\right)
\in\mathrm{SL}_2(\Z)$ put $j_\gamma=C't+D'$. The transformation of
$G_2$ and homogeneity of $\wp$ give
$C_\alpha(\gamma t)=j_\gamma^2 C_{\alpha\bar\gamma}(t)-C'j_\gamma/p$
\cite[Ch.~2]{shimura-1971}. Multiplying the coordinates at $\gamma t$
by $j_\gamma$ multiplies the lattice row by $\gamma$ and replaces
$(C,T)$ by $(C-TC'j_\gamma,Tj_\gamma^2)$. For
$(C,T)=(p^2C_\alpha(\gamma t),-p)$ this is
$j_\gamma^2(p^2C_{\alpha\bar\gamma}(t),-p)$. The coordinates contribute
$j_\gamma^n$ and differentiation contributes $j_\gamma^{2|\nu|}$.
Thus the family is permuted in weight $W_\nu$.

Fix $L'=L_\alpha$. Let $\mathfrak b$ be its first-coordinate image
and $0\times\mathfrak a^\vee$ its vertical kernel, where
$\mathfrak a^\vee=\mathfrak D_F^{-1}\mathfrak a^{-1}$ and
$\mathfrak D_F$ is the different. Choose $b_j\in\mathfrak b$,
$c_j\in\mathfrak b^{-1}$ with $\sum_j b_jc_j=1$ and lifts
$(b_j,t_j)\in L'$. Then $x=\sum_j c_jt_j$ satisfies $(u,ux)\in L'$
for all $u\in\mathfrak b$, and
$L'=\{(u,ux+v'):u\in\mathfrak b,\ v'\in\mathfrak a^\vee\}$.
Put $W_{\mathfrak a}=|D_F|^{1/2}\operatorname{Nm}\mathfrak a$.
For $x=0$ the expansion of $V_{L'}((t_i),1)/(p^n W_{\mathfrak a}^k)$
is $\sum_\xi A_\xi\exp(p\sum_i\sigma_i(\xi)t_i)$. General $x$ shifts
each $t_i$ by $\sigma_i(x)$. The coefficients are
\begin{equation}\label{eq:p1c1-hilbert-coefficients}
 A_0=2^{-f}z([\mathfrak a]^{-1}),\qquad
 A_\xi=\sum_{\substack{u\in\mathfrak a,\ b\in\mathfrak b,\ ub=\xi}}
 \left(\frac{|N_{F/\Q}u|}{\operatorname{Nm}\mathfrak a}\right)^{k-1}
 \quad(\xi\gg0),
\end{equation}
with zero coefficients elsewhere. Here $\xi\gg0$ means totally
positive, and the pairs are taken modulo
$(u,b)\mapsto(\epsilon^{-1}u,\epsilon b)$ for
$\epsilon\in\mathcal O_F^\times$. The integral ideal
$(u)\mathfrak a^{-1}$ runs through divisors of
$(\xi)(\mathfrak a\mathfrak b)^{-1}$, so this is a finite sum
of powers of ideal norms.

For $b\ne0$, apply Poisson summation on $\mathfrak a^\vee$, with
trace dual $\mathfrak a$. For even $k$, the real Fourier transform
of $(t+z)^{-k}$ at $u$ vanishes for $u\Im z<0$ and for $u\Im z>0$
equals
\[
 \operatorname{sgn}(\Im z)\frac{(-p)^k u^{k-1}}{(k-1)!}\exp(puz).
\]
Closing the contour in the half-plane containing $-z$ in the
second case gives the residue formula. Use $z=\sigma_i(b)t_i$ at
each embedding. The resulting support is $\xi=ub\gg0$. The inverse
covolume is $W_{\mathfrak a}$, and multiplication of the transforms
with $c_k$ gives the nonconstant terms. The row $b=0$ sums over
nonzero elements of $\mathfrak a^\vee$ modulo units and contributes
the partial zeta function of $C=[\mathfrak D_F\mathfrak a]$ at $k$.
The partial zeta functional equation with trivial sign and zero
additive shift identifies
$|D_F|^{z/2}\Gamma_{\R}(z)^f\zeta_F(z,C)$ at $z=k$ with the
corresponding value at $1-k$ for $C'=[\mathfrak a]^{-1}$, since
$CC'=[\mathfrak D_F]$ \cite{chapdelaine-2009}. Substituting
$\Gamma(s)=(s-1)!$ and
$\Gamma(1/2-s)=(-4)^s s!\sqrt\pi/(2s)!$ gives $A_0$.

For $k=2$ first apply Poisson summation with the regularizing factor
for $\Re\sigma>0$. The zero frequency from the rows $b\ne0$ is a
regular factor times
\[
 \zeta_F(1+2\sigma,[\mathfrak b]^{-1})\,c(\sigma)^f,\qquad
 c(\sigma)=\int_\R(t+i_{\C})^{-2}|t+i_{\C}|^{-2\sigma}\,dt .
\]
The integral is holomorphic near zero and $c(0)=0$ by the
antiderivative $-(t+i_{\C})^{-1}$. The partial zeta function has
at most a simple pole, so $f\ge2$ makes this contribution vanish
at zero.

For the nonzero frequencies, fix compact sets of $\sigma$ in a
neighborhood of zero and of $(t_i)$ in the product of upper
half-planes. Fix a derivative order in the real coordinates of
$(t_i)$. In the Fourier integrals rescale each real integration
variable by $|\sigma_i(b)|$. The imaginary parts of the rescaled
parameters stay away from zero. Repeated integration by parts in
each real variable gives, for sufficiently large $B$, a bound
\[
 C_B\prod_i
 \frac{(|\sigma_i(b)|+|\sigma_i(b)|^{-1})^m}
 {(1+|\sigma_i(ub)|)^B}
\]
for every derivative of the fixed order, with $m$ independent of
$B$. After rescaling, differentiation of the integrand
in the parameters gives integrable functions with any fixed
number of integrable real-variable derivatives, uniformly on
the compact sets. The scale factors in $b$ are bounded by the
displayed numerator. Integration by parts gives inverse powers
for $|\sigma_i(ub)|\ge1$, and the absolute integral bounds the
remaining frequencies.

Choose one $b$ per unit class with
$|\sigma_i(b)|\le C_1|\sigma_j(b)|$ for all $i,j$, by moving
its logarithms into a bounded fundamental domain for the unit
lattice after subtracting their mean. Replace $u$ by the inverse
unit multiple when changing $b$. This preserves $ub$ and permutes
the nonzero points of $\mathfrak a$. Write
$S(b)=\sum_i|\sigma_i(b)|$ and $S(u)=\sum_i|\sigma_i(u)|$.
Lattice discreteness bounds $S(b)$ below by a positive constant,
so $|\sigma_i(b)|+|\sigma_i(b)|^{-1}\ll S(b)$.
Also $\prod_i(1+|\sigma_i(ub)|)\ge\sum_i|\sigma_i(ub)|
\ge S(b)S(u)/(fC_1)$. With $B\ge mf+f+1$ the bound is at most
a constant times $S(b)^{mf-B}S(u)^{-B}$. Each factor is summable
over its lattice by the $O(R^f)$ count of points with $S\le R$.
This gives normal convergence for the nonzero frequencies and
their derivatives. Evaluating at $\sigma=0$ gives
\eqref{eq:p1c1-hilbert-coefficients} and its holomorphic derivatives.
\leanref{GlobalTheory.\allowbreak{}HeckeTwo.\allowbreak{}heckeWeightTwo\_\allowbreak{}of\_\allowbreak{}fe}{Development.\allowbreak{}AbelianVarietyTheory.\allowbreak{}Proofs.\allowbreak{}DivisorProofs.\allowbreak{}HeckeWeightTwoContinuationOfFunctionalEquation}{22--63}
\leanref{GlobalTheory.\allowbreak{}HilbertEis.\allowbreak{}eisExp}{Development.\allowbreak{}AbelianVarietyTheory.\allowbreak{}Proofs.\allowbreak{}DivisorProofs.\allowbreak{}HilbertEisensteinExpansionOfLipschitzHeckeZeta}{61--99}

Put $b_0=h+a$, $d_F'=v_\ell(|D_F|)$ and $\rho=2c+2a+h/2$.
One has $d_F'/f\le c$ since conductor primes are unramified in
$E$ and $|D_F|^2\mid |D_E|$. Since $N\in\mathfrak a^\vee$,
$Nu\in\mathfrak D_F^{-1}$ for $u\in\mathfrak a$. As $F/\Q$ is
Galois, the exponents of the different are constant above each
rational prime. Taking norms gives $\mathfrak D_F^f=(|D_F|)$.
Using \eqref{eq:p1c1-lattice-coordinates} for $b\in\mathfrak b$ gives
\begin{equation}\label{eq:p1c1-frequency-bounds}
 v(\sigma_i(\xi))\ge-b_0-d_F'/f+h/2\ge-\rho
 \quad(\xi\ne0,\ A_\xi\ne0).
\end{equation}
Also $N\mathfrak D_F\mathfrak a$ is integral, so
$v_\ell(\operatorname{Nm}\mathfrak a)\ge-fb_0-d_F'$ and
$v(W_{\mathfrak a}^k)\ge-sd_F'-nb_0$.

Write $r(k,j,t)=(-1)^{j-t}\binom jt(k+t)_{j-t}$ for $0\le t\le j$,
with $(x)_u=x(x+1)\cdots(x+u-1)$ and $(x)_0=1$. Replacing $C$ by
$C+v'$ in $D_\nu$ gives a sum over $\mu\le\nu$ with coefficients
$\prod_i r(k,\nu_i,\mu_i)(v')^{|\nu|-|\mu|}$. In one coordinate
use $z\mapsto z/(1+v'z)$ and the factor $(1+v'z)^{-k}$. The order
$t$ term contributes $j!/t!$ times the coefficient of $z^{j-t}$
in $(1+v'z)^{-k-t}$, which is the stated multiplier after
multiplication by $j!/t!$.

Starting with direction $(0,-p)$, define
\[
 B_{\mu,m}=\sum_{N\operatorname{Tr}_{F/\Q}\xi=m}
 A_\xi\exp(p\operatorname{Tr}_{F/\Q}(\xi x))
 \prod_i\sigma_i(\xi)^{\mu_i}.
\]
The sum is over the support of \eqref{eq:p1c1-hilbert-coefficients},
with exponent zero giving the factor one. It is finite since a
fixed trace bounds all totally positive coordinates.
The indices are integers by $N\mathfrak a\mathfrak b
\subseteq\mathfrak D_F^{-1}$, and the phase is a root of unity.
Differentiating the exponential in direction $-p$ gives
$-p^2\sigma_i(\xi)$ at each differentiation. The change of
direction with $v'=p^2C_\alpha$ gives
\[
 H_{\alpha,\nu}=
 W_{\mathfrak a}^k
 \sum_{\mu\le\nu}\left(\prod_i r(k,\nu_i,\mu_i)(-1)^{\mu_i}\right)
 C_\alpha^{|\nu|-|\mu|}\sum_{m\ge0}B_{\mu,m}q^{m/N}.
\]
Equations \eqref{eq:p1c1-hilbert-coefficients} and
\eqref{eq:p1c1-frequency-bounds} give $v(B_{\mu,m})\ge-e-|\mu|\rho$.
The coefficients of $C_\alpha$ are integral by
\eqref{eq:p1c1-torsion-coefficients} since $m_0$ is prime to
$\ell$, even when $\ell\mid N$. Each coefficient of $H_{\alpha,\nu}$
therefore has valuation at least $-R_\nu$, where
\begin{equation}\label{eq:p1c1-modular-coefficient-bound}
 R_\nu=e+s d_F'+nb_0+|\nu|\rho .
\end{equation}
\leanref{GlobalTheory.\allowbreak{}HilbFam.\allowbreak{}sz\_\allowbreak{}hfun\_\allowbreak{}coef}{Development.\allowbreak{}AbelianVarietyTheory.\allowbreak{}Proofs.\allowbreak{}DivisorProofs.\allowbreak{}HilbertEisensteinCoefficientValuationBound}{226--361}
\leanref{GlobalTheory.\allowbreak{}HilbFam.\allowbreak{}isFam\_\allowbreak{}hfun\_\allowbreak{}GL}{Development.\allowbreak{}AbelianVarietyTheory.\allowbreak{}Proofs.\allowbreak{}DivisorProofs.\allowbreak{}HilbertFamilyWeierstrassTwist}{27--35}

At the CM point $\omega=\Omega/p$ satisfies
$\omega^{12}=\Delta(\tau)$ and $j(\tau)$ is integral.
Lemma~\ref{B1.14a} gives $v(H_{1,\nu}(\tau)/\omega^{W_\nu})\ge-R_\nu$.
Put $V=\pi/\Im\tau$, so $V(\bar\tau-\tau)=-p$. The change of
direction now gives
\begin{equation}\label{eq:p1c1-jet-identity}
 \frac{D_\kappa(L;\tau,1,\bar\tau-\tau)V^{K_\kappa}}
 {p^{W_\kappa}\omega^{W_\kappa}}
 =
 \sum_{\nu\le\kappa}\left(\prod_i r(k,\kappa_i,\nu_i)\right)
 \left(\frac{V/p^2-C_1(\tau)}{\omega^2}\right)^{K_\kappa-|\nu|}
 \frac{H_{1,\nu}(\tau)}{\omega^{W_\nu}}.
\end{equation}
The parenthesis is $-(s_2(\Lambda)+\wp_\Lambda(P_1(\tau)))/\Omega^2$.
Lemma~\ref{B1.14b} bounds its valuation below by $-\rho$.
At $\ell=d$ one has $h=1$. At $\ell=2$ one has $c\ge1/2$,
since $2$ ramifies in the Galois field $E$ and ramification
index $e_0\ge2$ gives different exponent at least $e_0-1$
\cite[Ch.~III, \S~6]{serre-1979}. Thus $\rho\ge1$ for
$d\equiv3\pmod4$ and $\rho\ge2$ for $d\equiv1\pmod4$, where
$h=2$. At the other primes $\rho\ge0$ suffices.
Hence the left side has valuation at least $-R_\kappa$.
\leanref{GlobalTheory.\allowbreak{}HilbFam.\allowbreak{}sz\_\allowbreak{}jet\_\allowbreak{}conj}{Development.\allowbreak{}AbelianVarietyTheory.\allowbreak{}Proofs.\allowbreak{}DivisorProofs.\allowbreak{}SzJetConj}{24--57}

Let $\delta_\ell=\tfrac12\mathbf1_{\ell=d}$ and choose algebraic
$\gamma_\ell$ with valuation $c+\delta_\ell$. Take
$a_i=\Omega/\ell^a$ and $b_i=2\pi\gamma_\ell/(\sqrt d\,a_i)$.
These give the first two bounds of
\eqref{eq:p1c1-coordinate-estimates}. By
\eqref{eq:p1c1-deformed-substitution}, its normalized derivative
is the left side of \eqref{eq:p1c1-jet-identity} multiplied by
\[
 (AD\ell^a)^n
 \left(\frac{2\Im\tau}{\sqrt d}\gamma_\ell\ell^{2a}\right)^{K_\kappa}.
\]
The ratio $2\Im\tau/\sqrt d$ is $1$ or $2$ according as $d$
is $3$ or $1\bmod4$, and its valuation is $h/2-\delta_\ell$.
The denominator exponent is at most
\[
 R_\kappa-n(h+a)-K_\kappa(h/2+c+2a)
 = e+s d_F'+K_\kappa c\le e+(n+K_\kappa)c .
\]
The coordinates depend on the lattice and period root and work
for every $\kappa$.
\end{proof}

\subsection*{Proof of \ref{B1.15a}}\label{prf:B1.15a}\phantomsection
\begin{proof}
Siegel--Klingen gives rational $z(C)$ before the conductors,
weights or local embeddings are chosen \cite{klingen-1962,deligne-ribet-1980}.
For $A\ge2$ prime to $\ell$, apply Deligne--Ribet at weight $k$
to the ordinary class indicator on ideals prime to $\ell$ and
the ideal $(A)$ \cite{deligne-ribet-1980}. The indicator takes values in $\Z_\ell$. It is constant on ray classes modulo $(\ell)$ and invariant under translation by $(A)$. Thus $A^n-1$ times the partial
zeta value with the primes above $\ell$ removed is integral at
$\ell$, since $\operatorname{Nm}(A)=A^f$. For each omitted
$\mathfrak p$, the removed Euler factor is
$1-(\operatorname{Nm}\mathfrak p)^{k-1}S$, with $S$ translation
by $[\mathfrak p]^{-1}$ on classes. If $S^j=1$, its inverse is
the finite geometric sum divided by
$1-(\operatorname{Nm}\mathfrak p)^{(k-1)j}$, an $\ell$-adic unit.
Restore the factors and divide by $2^f$ to obtain the first
inequality of \eqref{eq:p1c1-class-denominator}.

For odd $\ell$, a primitive root modulo $\ell^2$ attains the
minimum. If $\ell-1\mid n$, the valuation is $1+v_\ell(n)$
by binomial expansion along powers of $\ell$. Otherwise it is
zero already modulo $\ell$. For odd $A$ and even $n$,
\[
 v_2(A^n-1)=v_2(A-1)+v_2(A+1)+v_2(n)-1 .
\]
The minimum is attained at $A=3$. In particular $e_\ell=0$ if
$\ell>n+1$, and if $\ell-1$ is a proper divisor of $n$ then
$\ell\le n/2+1$. The product of primes up to $u$ is at most
$4^u$. For $u=2m$ the primes above $m$ divide
$\binom{2m}{m}$, and for $u=2m+1$ those above $m+1$ divide
$\binom{2m+1}{m}$. Both binomial coefficients are at most $4^m$,
so induction gives the bound. Separating a possible prime $n+1$
now gives
\[
 \sum_\ell e_\ell\log\ell
 \le \log n+(f+2)\log2+\log(n+1)+(n/2+1)\log4
 \le n\log(n+1)\qquad(n\ge16).
\]
The last step uses $2f\le n$ and $\log(n+1)\ge4\log2$.
\leanref{GlobalTheory.\allowbreak{}ZetaDenom.\allowbreak{}zetaDenominators}{Development.\allowbreak{}AbelianVarietyTheory.\allowbreak{}Proofs.\allowbreak{}DivisorProofs.\allowbreak{}PartialZetaDenominatorsCMFieldOfSiegelKlingenDel}{146--164}

For the toric bound choose $y$ equal to one above $\ell$ and
use Lemma~\ref{B1.14c}. Put
\[
 J_\kappa=\partial^\kappa g(0)\prod_i
 \frac{(b_i/a_i)^{\kappa_i}}{a_i^k},\qquad
 A_i=a_i/\Omega,\qquad B_i=\frac{2\pi}{a_i b_i\sqrt d}.
\]
Differentiation gives
$\partial^\kappa g(0)=c_k(-1)^{K_\kappa}\prod_i(k)_{\kappa_i} S_y$.
At $k=2$ the regularized holomorphic derivatives have factors $(k+\sigma)_{\kappa_i}$. The continued Fourier expansion gives their evaluation at zero. Since $(k)_j=\Gamma(k+j)/(k-1)!$,
\begin{equation}\label{eq:p1c1-derivative-value}
 x_\Omega(S_y)=(-1)^{K_\kappa}J_\kappa
 \prod_i B_i^{\kappa_i}A_i^{k+2\kappa_i}.
\end{equation}
If $c_\ell=0$, then $\ell\ne d$, since $d$ ramifies in the
quadratic subfield and the conductor primes are unramified in $E$.
The coordinate estimates give denominator exponent $e_\ell+Ma_\ell$.
For $c_\ell>0$, dividing $J_\kappa$ by $\prod_i\kappa_i!$ gives
denominator exponent at most $e_\ell+nc_\ell+K_\kappa(c_\ell+1)$,
since $\sum_i v_\ell(\kappa_i!)=\sum_{i,j\ge1}
\lfloor\kappa_i/\ell^j\rfloor\le K_\kappa$, where $i$ ranges over
the CM type. Restore the integral factorials. The $B_i$ contribute
at most $K_\kappa(c_\ell+\tfrac12\mathbf1_{\ell=d})$ and the $A_i$
at most $Ma_\ell$. This proves $v(x_\Omega(S_y))\ge-b_\ell$.

Every order class has a representative equal to one above $\ell$.
Weak approximation first moves these components into the open local
order units, which can then cancel them. At a split conductor place
an order unit is a diagonal real-field unit times a pair congruent
to one modulo the conductor. The character is trivial on both
factors by its real-field restriction and conductor. Away from the
conductor it is unramified. Thus $\chi_f$ and its Picard twists are trivial on $\widehat O_Q^\times$. A principal
change of representative acts on $Y_y$ by multiplication. The
infinity character and the product formula give the inverse of
the resulting scaling of $S_y$ in the factor
$(\chi_f(y)|y|_{E,f}^s)^{-1}$.

Let $H_y$ be the class order and write $y^{H_y}=t v'$ with
$t\in E^\times$ and $v'\in\widehat O_Q^\times$. For
$u=(\chi_f(y)|y|_{E,f}^s)^{-1}$ one has
$u^{H_y}=\chi_\infty(t)|N_{E/\Q}t|^s$. Since $t$ is a unit above
$\ell$, each embedded conjugate of $t$ is an $\ell$-unit. Thus $u$
is an algebraic $\ell$-unit, and the normalized class term
$u x_\Omega(S_y)$ has the same bound. Picard character values are
roots of unity. The ultrametric inequality for the finite twisted
sum and the twelfth power give \eqref{eq:p1c1-toric-denominator}.
The argument applies at every embedding above $\ell$.
\leanref{GlobalTheory.\allowbreak{}ToricDenomRow.\allowbreak{}toricDenominator}{Development.\allowbreak{}AbelianVarietyTheory.\allowbreak{}Proofs.\allowbreak{}DivisorProofs.\allowbreak{}ToricDenominatorOfZetaDenominators}{39--96}

For transport, choose a basis of the integer lattice of $K$ and
a root $\Omega_0^{12}=D(\Lambda_K)$. The elliptic main theorem of complex multiplication gives the conjugate algebraic elliptic curve and a differential that determines the period scale $\Omega_1$. In Proposition~\ref{B6.05f}(1) the arithmetic Artin symbol of $t$ is $\operatorname{Art}_K(t)^{-1}$. This is $g|_{K^{\rm ab}}$ by the choice of $t$. The theorem of complex multiplication therefore sends the lattice $\mathfrak a$ to $t^{-1}\mathfrak a$ and the torsion marking $z$ to $t^{-1}z$ on that lattice \cite[Ch.~6]{shimura-1971}.
\leanref{GlobalTheory.\allowbreak{}LatVal.\allowbreak{}exists\_\allowbreak{}globalArtinMap}{Development.\allowbreak{}AbelianVarietyTheory.\allowbreak{}Proofs.\allowbreak{}DivisorProofs.\allowbreak{}PeriodExistenceCm}{32--79}
\leanref{GlobalTheory.\allowbreak{}LatVal.\allowbreak{}periodEx\_\allowbreak{}of\_\allowbreak{}mainTheorem}{Development.\allowbreak{}AbelianVarietyTheory.\allowbreak{}Proofs.\allowbreak{}DivisorProofs.\allowbreak{}PeriodExistenceCm}{81--103}
 The scales work at every weight.
The Fourier calculation uses $[F:\Q]\ge2$. In the twisted
assertion, if $[F:\Q]=1$, then $F=\Q$ and finite unramifiedness
gives $|D_E|=|D_F|^2=1$, which is impossible for an imaginary
quadratic field.

Choose a common finite level $N$ after clearing the denominators
of the two coordinates of $Y_y$. The construction of $H_{\alpha,\nu}$
uses the resulting full lattice and a nonzero torsion marking.
Put $g\zeta_N=\zeta_N^a$ and $d_a=\operatorname{diag}(1,a)$
modulo $N$. In the coefficient
\eqref{eq:p1c1-torsion-coefficients}, $g$ preserves the row's
first coordinate and multiplies the second by $a$. In the Hilbert
expansion the constant term and ideal norms are rational.
The phase $\exp(p\operatorname{Tr}_{F/\Q}(\xi x))$ goes to the
phase with $x$ replaced by $ax$. The lattice substitution
$(b,bx+v')\mapsto(b,abx+av')$ modulo level leaves the first-coordinate
ideal and vertical kernel unchanged. The equality $g\phi_i=\phi_{\pi(i)}$
permutes the factors $\sigma_i(\xi)^{\mu_i}$, with new order
$\mu'_{\pi(i)}=\mu_i$. Thus $g$ sends the coefficients of the family at $\alpha$ to those of the conjugate family at $\alpha d_a$.

Let $A$ be the torsion marking matrix from CM reciprocity, and
write $\Lambda_j=c_j(\Z\tau_j+\Z)$ for the two period lattices.
Put $p=2\pi i$ and $\omega_j=c_j\Omega_j/p$. Indeed, lattice homothety gives $g_2(c\Lambda)=c^{-4}g_2(\Lambda)$ and $g_3(c\Lambda)=c^{-6}g_3(\Lambda)$. Hence $D(\Lambda_j)=c_j^{-12}p^{12}\Delta(\tau_j)$. CM reciprocity carries $g_2(\Lambda_0)/\Omega_0^4$ and $g_3(\Lambda_0)/\Omega_0^6$ to the corresponding invariants at $\Lambda_1$. Their discriminant at $\Lambda_0$ is $D(\Lambda_0)/\Omega_0^{12}=1$, so $D(\Lambda_1)/\Omega_1^{12}=g(1)=1$. Hence $\Omega_j^{12}=c_j^{-12}p^{12}\Delta(\tau_j)$. Thus $\omega_j^{12}=\Delta(\tau_j)$ at both lattices. For a nonzero torsion point represented by $z$ on $\Z\tau_j+\Z$, homothety gives
\[
 \frac{\wp_{\Lambda_j}(c_jz)}{\Omega_j^2}
 =\frac{\wp_{\Z\tau_j+\Z}(z)/p^2}{\omega_j^2},\qquad
 \frac{\wp'_{\Lambda_j}(c_jz)}{\Omega_j^3}
 =\frac{\wp'_{\Z\tau_j+\Z}(z)/p^3}{\omega_j^3}.
\]
\leanref{GlobalTheory.\allowbreak{}TransportProof.\allowbreak{}disc\_\allowbreak{}j\_\allowbreak{}of\_\allowbreak{}torsTransport}{Development.\allowbreak{}AbelianVarietyTheory.\allowbreak{}Proofs.\allowbreak{}DivisorProofs.\allowbreak{}HilbertClassFieldLatticeTransport}{54--73}
\leanref{GlobalTheory.\allowbreak{}WpFam.\allowbreak{}Transp.\allowbreak{}main}{Development.\allowbreak{}AbelianVarietyTheory.\allowbreak{}Proofs.\allowbreak{}DivisorProofs.\allowbreak{}EisensteinWeierstrassTorsionConjTransport}{20--50}
CM reciprocity transports $E_4(\tau_j)/\omega_j^4$ and
$E_6(\tau_j)/\omega_j^6$, with their usual weight normalizations,
and transports the normalized $\wp,\wp'$ values at a torsion row
$x$ to those at $xA$. To pass from the action on coefficients to
the action on values, use the two basis torsion rows. Their
pairs of $\wp,\wp'$ values, normalized in weights two and three,
determine the ordered basis on the elliptic curve. Thus they
separate the labels $\alpha$. In the family with conjugate
coefficients the label giving the transported values of these
four functions is $\beta=\alpha A d_a^{-1}$, since
$\beta d_a=\alpha A$. The same label works for every normalized
form. Adjoin the four separating functions to the family. Use the elementary symmetric functions of linear combinations with independent indeterminates to form the polynomial of its values over the labels. The homogeneous coefficients are
full-level modular forms, so their normalized values are
expressed in the normalized $E_4,E_6$. Conjugating the
coefficients and specializing the separating functions identifies
the factor at $\alpha$ with that at $\beta$ in the conjugate family.
Consequently the CM value at $\alpha$ transports to the
conjugate value at $\alpha A$. For the correction in
\eqref{eq:p1c1-jet-identity}, $\wp$ transports by the same marking.
The isogeny identity of Lemma~\ref{B1.14b} for a nonreal multiplier
in $K$ expresses $s_2$ as a sum of $\wp$ values on its kernel
divided by a nonzero coefficient in $K$. Thus the normalized
correction transports as well.

At the initial lattice use $c_0=1$, and at the inverse ideal
choose the basis so that $c_1$ is a positive rational. Put
$t_E$ for the image of $t$ in the finite id\`eles of $E$.
On the finite completions the marking multiplies
$y\widehat O_Q$ by $t_E^{-1}$. In the basis $(\tau_1,1)$ its
coordinates therefore describe $c_1^{-1}Y_{yt_E^{-1}}$.
Multiplying an element by $c_1^{-1}$ multiplies the undeformed
summand of weight $(k,\kappa')$ by $c_1^{kf}$, since at each
embedding the two exponents of this real scalar sum to $-k$.
Write $n_0>0$ for the common integer clearing the two lattice coordinates. Put $h=K_\kappa$, $M=kf+2h$, and $q_j=2\Im\tau_j/\sqrt d$. Let $X_j$ be the right side of \eqref{eq:p1c1-jet-identity} at the marking $\tau_j$, evaluated with its coordinate lattice and torsion row. Its total period denominator is $\omega_j^M$. Each term of derivative order $\nu$ has total weight $2(h-|\nu|)+(kf+2|\nu|)=M$. Its integer coefficient is preserved by permutation of the embeddings. The transport of the correction of weight two and of the modular family of weight $kf+2|\nu|$ therefore gives $gX_0=X_1$. No action on $p$ is needed in this equality.

For the coordinate lattice at $j$, denote its sum before clearing denominators by $S^{\rm coord}_j$. Thus $S^{\rm coord}_0=S_y$, and $S^{\rm coord}_1$ is the sum for $c_1^{-1}Y_{yt_E^{-1}}$ with the permuted derivative orders. Substitution into the left side of \eqref{eq:p1c1-jet-identity} gives
\[
 x_{\Omega_j}(S^{\rm coord}_j)=(-1)^h n_0^{kf}q_j^h c_j^M X_j .
\]
Clearing the coordinates contributes $n_0^{kf}$ and converting the derivative contributes $(-1)^h q_j^h$. Since $p\omega_j=c_j\Omega_j$, replacing the denominator $(p\omega_j)^M$ by $\Omega_j^M$ multiplies the value by $c_j^M$. On the continued sum the change $t\mapsto c_1^{-1}t$ contributes $c_1^{kf+2f\sigma}$ before evaluation at $\sigma=0$. Thus $S^{\rm coord}_1=c_1^{kf}S'_{yt_E^{-1}}$. Dividing by $c_1^{kf}$ leaves $c_1^{M-kf}=c_1^{2h}$. Since $c_0=1$ and $q_0$ is rational, the two identities are
\leanref{GlobalTheory.\allowbreak{}LatConjEll.\allowbreak{}sum\_\allowbreak{}transport}{Development.\allowbreak{}AbelianVarietyTheory.\allowbreak{}Proofs.\allowbreak{}DivisorProofs.\allowbreak{}LatticeValueConjugationOfEllipticCMTheorems}{140--165}

\[
 \begin{split}
 g x_{\Omega_0}(S_y)
 &=(-1)^{K_\kappa} n_0^{kf}
 (2\Im\tau_0/\sqrt d)^{K_\kappa} X_1,\\
 x_{\Omega_1}(S'_{yt_E^{-1}})
 &=(-1)^{K_\kappa} n_0^{kf}
 (2c_1^2\Im\tau_1/\sqrt d)^{K_\kappa} X_1.
 \end{split}
\]
To see the rationality used here, write the integral element $\delta$ with $\delta^2=-d$ as $a\tau_0+b$ in the embedded integer lattice. Taking imaginary parts gives $q_0=2/|a|$. The positive determinant taking a basis of $\mathcal O_K$ to a basis of the inverse ideal is its fractional ideal norm. Computing the same determinant by areas gives
\[
 A=|t|_{K,f}=\frac{c_1^2\Im\tau_1}{\Im\tau_0}.
\]
Multiplication of the first displayed value by $A^{K_\kappa}$ gives the second, also when $X_1=0$. Hence $g x_{\Omega_0}(S_y)=A^{-K_\kappa}x_{\Omega_1}(S'_{yt_E^{-1}})$. Permuting the derivative orders preserves their sum $K_\kappa$ and the product $\prod_i\Gamma(k+\kappa_i)$, hence also $C_0$.
\leanref{GlobalTheory.\allowbreak{}LatConjEll.\allowbreak{}value\_\allowbreak{}eq}{Development.\allowbreak{}AbelianVarietyTheory.\allowbreak{}Proofs.\allowbreak{}DivisorProofs.\allowbreak{}LatticeValueConjugationOfEllipticCMTheorems}{166--203}
\leanref{GlobalTheory.\allowbreak{}LatConjEll.\allowbreak{}latConj}{Development.\allowbreak{}AbelianVarietyTheory.\allowbreak{}Proofs.\allowbreak{}DivisorProofs.\allowbreak{}LatticeValueConjugationOfEllipticCMTheorems}{204--396}
\leanref{GlobalTheory.\allowbreak{}LatConjEll.\allowbreak{}covol}{Development.\allowbreak{}AbelianVarietyTheory.\allowbreak{}Proofs.\allowbreak{}DivisorProofs.\allowbreak{}LatticeValueConjugationOfEllipticCMTheorems}{55--114}
\leanref{GlobalTheory.\allowbreak{}LatConjEll.\allowbreak{}ratio\_\allowbreak{}rat}{Development.\allowbreak{}AbelianVarietyTheory.\allowbreak{}Proofs.\allowbreak{}DivisorProofs.\allowbreak{}LatticeValueConjugationOfEllipticCMTheorems}{115--139}

\leanref{GlobalTheory.\allowbreak{}T.\allowbreak{}latticeValueIsAlgebraic\_\allowbreak{}of\_\allowbreak{}ellipticCM\_\allowbreak{}theorems}{Development.\allowbreak{}AbelianVarietyTheory.\allowbreak{}Proofs.\allowbreak{}DivisorProofs.\allowbreak{}LatticeValueIsAlgebraicOfEllipticCMTheoremsStatement}{9--11}
\leanref{GlobalTheory.\allowbreak{}T.\allowbreak{}latticeValueConjugation\_\allowbreak{}of\_\allowbreak{}ellipticCM\_\allowbreak{}theorems}{Development.\allowbreak{}AbelianVarietyTheory.\allowbreak{}Proofs.\allowbreak{}DivisorProofs.\allowbreak{}LatticeValueConjugationOfEllipticCMTheoremsStatement}{9--11}

The normalized modular form values and correction are algebraic.
The class order calculation shows that the finite class factors
are algebraic as well. In the twisted sum, after transport the
term from $y$ carries
$[\chi'_f(y)|y|_{E,f}^s]^{-1}\psi'([y])^{-1}$.
With $y'=yt_E^{-1}$ this is the factor at $y'$ times
$[l_\chi\psi'([t])]^{-1}$. Multiplication by $t_E^{-1}$ permutes
the order classes, so summing gives \eqref{eq:p1c1-twisted-transport}.
Both periods come from the elliptic lattices and are chosen before
the extension, type, conductor or twist.
\leanref{GlobalTheory.\allowbreak{}TwistedTransportProof.\allowbreak{}twistedTransport\_\allowbreak{}of\_\allowbreak{}latClaims}{Development.\allowbreak{}AbelianVarietyTheory.\allowbreak{}Proofs.\allowbreak{}DivisorProofs.\allowbreak{}TwistedTransportOfLatClaims}{189--203}

For ordinary transport the Hilbert class field theorem gives
$K(j(\Lambda_K))$ with one conjugate $j$ for each ideal class
\cite{shimura-1971}. Choose $t_g$ by CM reciprocity and use its
inverse ideal for $\Lambda_g$. Restriction of embeddings of $L$
over $K$ to the Hilbert class field has fibres of cardinality
$[L:K(j(\Lambda_K))]$. Set the Picard character to one and take
twelfth powers in the transport identity to obtain the first
line of \eqref{eq:p1c1-ordinary-transport}. Since
$|t_g|_{E,f}=A^{[E:K]}=A^f$ and $\chi'_f$ is unitary, one has
$|l_\chi|=A^{sf}$. The absolute value therefore gives the factor
$A^{-12(K_\kappa+sf)}=A^{-6M}$. For equal weights, the periods and
$A^{-K_\kappa}$ cancel when dividing the two normalized values.
A field automorphism carries a nonzero denominator to a nonzero
value, which proves the ratio assertion.
\leanref{GlobalTheory.\allowbreak{}TransportProof.\allowbreak{}transport\_\allowbreak{}of}{Development.\allowbreak{}AbelianVarietyTheory.\allowbreak{}Proofs.\allowbreak{}DivisorProofs.\allowbreak{}HilbertClassFieldLatticeTransport}{160--192}

For reflection take $C_{F,0}=C'_0=x_0=0$. Fix a recursive size datum of \ref{B8.19}, an admissible level-character system of \ref{B8.20}, and a grid of \ref{B8.21} as in Definition~\ref{B1.14k}(5). The host conditions of \ref{B8.19} supply the CM type, and the splitting conditions and conductor exponents in \ref{B8.19}--\ref{B8.21} give the base conductor conditions on $\chi_0=\chi_{0,i,j}$. For each $t$, write $\xi=\xi_t$. By Definition~\ref{B1.14k}(5)(a), $\xi$ is trivial on $\mathbb A_F^\times$, has finite order, and is tame, and the coefficient-field condition supplies all its required local values. By Definition~\ref{B1.14k}(5)(b), at each place $v$ where $\xi$ ramifies and $\chi_0$ is unramified, the residue field has prime order $q_v\mid Q_n(t)$ with $v\in W_t$ and $d_v\ge1$, so the root condition supplies all $q_v$th roots of unity in $K$, and these places occur in conjugate pairs under $E/F$. By Definition~\ref{B1.14k}(5)(c), at each place where both $\xi$ and $\chi_0$ ramify, the residue field has prime order, the place is unramified over its rational prime, and $\chi_0$ has conductor exponent two.

By Definition~\ref{B1.14k}(5)(a) and~(d), $\chi_0$ and $\chi_0\xi$ share the infinity exponents $m_\tau=\kappa_\tau(h)+h>h\ge e_0$, where $h$ is half the sum of the grid ranks at $(i,j)$ and $(i+1,j+1)$. Thus $|m_\tau|+s\ge 1+m_\tau-e_0>0$ and $|m_\tau|+1-s>0$, so the archimedean gamma factors at $s$ and $1-s$ are finite and nonzero throughout $1-e_0\le s\le0$.

Apply the local epsilon-factor formulas \cite{tate-1967}.
At an unramified twist place, with reference conductor exponent
$a_v$, different exponent $d_v$ and uniformizer $\varpi_v$, the
multiplier is $\xi_v(\varpi_v^{a_v+d_v})$. At a place where the
twist has conductor one and $\chi_0$ is unramified, the
conductor-one sum gives $\chi_0(\varpi_v)$ times a finite-order
character value, a power of $q_v$ and a nontrivial Gauss sum.
Its additive character takes values in the stipulated roots of
unity. For the residue character, still denoted $\xi$, additive
orthogonality gives $G(\xi)G(\xi^{-1})=\xi(-1)q_v$. Hence $G(\xi)$
is a unit away from $q_v$. At a place where both ramify, write
units modulo $\varpi_v^2$ as residue lifts times $1+\varpi_v x$.
The sum over $x$ selects a single residue unit, since $\chi_0$
has conductor exponent two. The tame twist is constant on
the principal units, so the quotient is its value at the
selected unit times the uniformizer factor.

The places where the twist ramifies and the reference character is unramified occur in conjugate pairs.
The two reference factors $\chi_0(\varpi_v)$ in each pair are
inverse by triviality on the real-field id\`eles. The remaining
product is $\iota_D(u)$ for a unit away from $Q_n(t)$. In the
two global Hecke--Tate equations the infinite factors agree.
Triviality on real-field id\`eles identifies each inverse
character with its conjugate under $E/F$, and conjugation
permutes the Euler factors of its Hecke function. Cross-multiply
the two functional equations and cancel the common nonzero
gamma factors to obtain \eqref{eq:p1c1-grid-reflection}.
\leanref{GlobalTheory.\allowbreak{}SqTop.\allowbreak{}refl\_\allowbreak{}square}{Development.\allowbreak{}AbelianVarietyTheory.\allowbreak{}Proofs.\allowbreak{}DivisorProofs.\allowbreak{}ReflectionOfFunctionalEquation}{156--250}
\leanref{GlobalTheory.\allowbreak{}SqTop.\allowbreak{}refl\_\allowbreak{}of\_\allowbreak{}functionalEquation}{Development.\allowbreak{}AbelianVarietyTheory.\allowbreak{}Proofs.\allowbreak{}DivisorProofs.\allowbreak{}ReflectionOfFunctionalEquation}{252--262}
\end{proof}

\subsection*{Proof of \ref{B1.15ba}}\label{prf:B1.15ba}\phantomsection
\begin{proof}
A simple module $M$ over a $\C$-algebra of at most countable complex dimension has scalar commutant. Indeed a nonscalar endomorphism $A$ has $A-z$ invertible for all $z\in\C$ by Schur's lemma. For $m\ne0$, the vectors $(A-z)^{-1}m$ are $\C$-linearly independent, since clearing denominators in a finite relation gives a nonzero polynomial in $A$ killing $m$, and its linear factors are invertible; this contradicts countable dimension.

For a linearly independent tuple in such a simple module $M$, its orbit spans $M^b$, where $b$ is the tuple length. A proper span would admit a nonzero module map $M^b\to M$ killing the tuple, by semisimplicity, and the scalar commutant would give a linear dependence. Thus the acting algebra prescribes the images of any independent finite family. In $\bigotimes_{i=1}^r M_i$, apply this in each slot to isolate a nonzero pure tensor from a finite sum and send it to arbitrary pure tensors; for $r=0$ the empty product $\C$ is simple. In the restricted tensor product, any nonzero vector and any target vector lie in $\bigotimes_{v\in S}M_v$ for a common finite set $S\supseteq S_0$ by adjoining the nonzero vectors $s_v$, so simplicity follows from the finite case on $S$.
\end{proof}

\subsection*{Proof of \ref{B1.15b}}\label{prf:B1.15b}\phantomsection
\begin{proof}
For the polarized pair take $r$, the weight datum $w_{\rm gt}$
and jump multisets $H_\tau$ from the field of $O$ in
\ref{B9.02}. On a cuspidal representation of $\GL_n$ over the
totally complex field, inverse transpose inverts the unramified
Satake parameters. Cuspidal rigidity identifies this action
with the contragredient \cite{jacquet-shalika-1981}.
\leanref{PositivityTheory.\allowbreak{}JMult1.\allowbreak{}gelfandKazhdanDual\_\allowbreak{}of\_\allowbreak{}rigidity}{Development.\allowbreak{}WeilBound.\allowbreak{}CoreStatements.\allowbreak{}TemperedRepresentationTheory.\allowbreak{}MultiplicityOneRigidity.\allowbreak{}GelfandKazhdanDualOfRigidity}{72--78}
For the second representation field use the dual local--global
field of $O$ on conjugate self-dual regular algebraic cuspidal
representations of rank $a\ge2$, with polarization $\chi_{E/F}^a$.
Fix $r,w_{\rm gt},H_\tau$ before the place $w\mid\ell$,
prolongation, period datum and bridge. For each correspondence of Definition~\ref{B9.47}, suppose its normalized local parameter passes the inertia test of Definition~\ref{B1.15ac}. Wild inertia acts trivially since its image is finite and all its elements are unipotent. When $d_w=2$, complete splitting gives $E_w\simeq\Q_\ell$, and the repeated-root condition gives the full tame test of Definition~\ref{B9.05}. When $d_w=1$, every eigenvalue of a tame generator has order dividing $q_w-1$. Its image is semisimple of finite order, so its $q_w$th power equals itself. Frobenius therefore centralizes the inertia image, which gives the toral test. In both cases Definition~\ref{B9.06} gives the stated comparison.

The shared local fields use Bernstein--Zelevinsky finiteness,
Bernstein's unitary product theorem and Zelevinsky's polynomial
ring on segments \cite{bernstein-zelevinsky-1977,bernstein-1984,zelevinsky-1980}.
Use \ref{B8.12}--\ref{B8.13} for the global Weil and doubling
formulas. For the unramified local integral, the spherical section
restricted to the doubling embedding is constant on each $K$
double coset. The zonal coefficient has the same sums over its
right $K$ cosets as the Borel spherical function with parameter
$\alpha$. Thus with $\operatorname{vol}K=1$ the integral has the
same double-coset sums as that spherical calculation. In the
length generating series the variable is $X=q^{-s-a/2}$. By the unramified local doubling calculation of Piatetski-Shapiro and Rallis \cite[Part~A, \S3, Theorem~3.1]{gelbart-piatetski-shapiro-rallis-1987} and Lapid and Rallis \cite[Proposition~3]{lapid-rallis-2005} for the split unitary group $\GL_a$, in terms of $X$ the resulting rational function is the quotient of the standard local $L$-factor $L(\alpha,s+1/2)L(\alpha^{-1},s+1/2)$ by the unramified degenerate principal-series normalizing factor $\prod_{j=1}^a(1-q^{-2s-j})^{-1}$, namely
\[
 \frac{\prod_{j=1}^a(1-q^{a-j}X^2)}
 {\displaystyle\prod_{i=1}^a
 (1-\alpha_i q^{(a-1)/2}X)
 (1-\alpha_i^{-1}q^{(a-1)/2}X)} .
\]
The double-coset series is absolutely summable for $|X|<\delta$
for some $\delta>0$ depending on the local parameters. The
zonal coefficient is constant on each double coset, so its
absolute contribution is the absolute value of its coset sum
times the absolute section value. Hence the same summability
bound applies to the integral. Substitute $X=q^{-s-a/2}$ and take
$s_0=-\log_q\delta-a/2$ to obtain the asserted two Satake factors
and finite product.

The Harish--Chandra results give the fixed-level finiteness,
discrete series and formal degree, and square-integrable Schur
orthogonality in the forms stated above. Use \ref{B8.38a} for
the smooth globalization and the seminorm bound for the
intertwiner. For cuspidal decay at fixed open level, use reduction theory. In a cusp the average over the unipotent radical is zero.
The difference between the value and this average is bounded
by the size of the reduced unipotent shifts times a finite
sum of first-derivative bounds on a fixed compact thickening.
The shift size is bounded by a negative power of the quotient
height. Differentiation preserves the level and cuspidality,
so iteration improves any fixed moderate-growth exponent
to the prescribed decay order. The bounded-height region uses
the estimate of order zero. This gives finitely many derivative
words and a constant at each level and pair of growth and decay
orders, the finite-seminorm form of cuspidal rapid decay used
here \cite[Theorem~A]{miller-schmid-2012}. Taking the largest
word length gives the stated Sobolev bound.

For a smooth compactly supported test $f$, the derivatives of
$R(f)u$ are convolutions of $u$ against derivatives of $f$, with
support contained in $\operatorname{supp} f$. The quotient
kernel for such a test is the sum of its translates by rational
points. The number of contributing translates on the support
is bounded by a power of the group height, uniformly in the
evaluation point. By Cauchy--Schwarz on the quotient, these
convolutions have moderate-growth bounds with constants
$C_l\|u\|_2$ and a common exponent determined by the support.
The convolution has a fixed open level and is cuspidal.
Apply the finite-word decay estimate with decay order zero
to obtain $|R(f)u(g)|\le C_f\|u\|_2$ at every $g$. For any
finite orthonormal system $(u_i)$, Bessel's inequality for
evaluation of the convolution gives
$\sum_i|R(f)u_i(g)|^2\le C_f^2$. Integration over the
finite-volume quotient gives a bound independent of the
orthonormal system, hence $R(f)$ is Hilbert--Schmidt.
Dixmier--Malliavin factorization, with the first support in
the prescribed identity neighborhood, writes $f$ as a finite
sum of convolutions \cite{dixmier-malliavin-1978}. At the finite
places use normalized averaging at a common open level of the
test. Thus $R(f)$ is a finite sum of products of two
Hilbert--Schmidt operators and is trace class.

Flath's theorem gives restricted tensor factorization of the
admissible irreducible adelic representation, with spherical
factors almost everywhere \cite{flath-1979}. Strong approximation
applies to the special unitary group in the stated noncompact
range \cite{platonov-rapinchuk-1994}. Theorem~\ref{B7.31} gives
the regular Green field of $P$ from \ref{B7.05}, and Green's
classification supplies the cuspidal and series entries in
$F$ \cite{green-1955}. Use Clozel for cohomologicality and purity
in the regular algebraic range \cite{clozel-1990}. Apply
\ref{B6.04a} with finite unramifiedness and rank $n\ge1$ to
obtain the single definite base change at all places. The
geometric-lemma filtration and the segment classification
give the unlinked products and standard cosocle. Use
depth-zero support and Jacquet injectivity for the
first-congruence fields. The generic unitary dual gives the
complementary products, and the strict Satake bounds and
cuspidal rigidity are the Jacquet--Shalika results in the
unitary cuspidal case \cite{jacquet-shalika-1981}.

For the finite Galois prime-existence field, apply \ref{B8.42p}
to the kernel of restriction to the extension and a lift of
the specified automorphism. The resulting prime has residue
degree and ramification index one over $\Q$. Its Frobenius
restricts to the automorphism and its inertia acts trivially on
the extension. For the second field use the prescribed open
normal subgroup and finite avoidance set in \ref{B8.42p}.

Artin reciprocity gives the $\ell$-adic and ray-character fields.
Above $\ell$, de Rham comparison for algebraic characters and
potential semistability give the Weil--Deligne comparison.
For every compatible local component $\xi$ and every admissible
correspondence $R$, local reciprocity gives the character parameter
of $\xi$ with zero monodromy. The rank-one law of \ref{B9.47}
identifies it with $R(\xi)$. For the integral crystalline lattices,
use the prismatic comparison of \ref{B8.41h} and its functoriality
on lattice maps. The CM curve clause uses Gross's $A(d)$ over
$\Q(j(\mathcal O_K))$ \cite[Ch.~5]{gross-1980}. Its minimal
discriminant ideal gives the stated local equations and good
reduction. For the Weber field use
the elliptic main theorem of complex multiplication at modulus
$(n)$ \cite[Ch.~6]{shimura-1971}. This is the principal-modulus
specialization of the ray class field theorem. The two formulas
in \ref{B8.13} give the convergent Siegel--Weil fields, the second
with its volume and Gram factors. The Schr\"odinger formula on
the Fock vector gives the indicated special value and
\ref{B8.12} gives continuation.

For the nonendoscopic hypotheses, combine the local and global results above with the segment classification, the depth-zero results, local converse and the existence of the local correspondence in Definition~\ref{B9.01}. The cohomology factor uses the local and
global finiteness and Euler characteristics, Poitou--Tate,
local duality and unramified orthogonality, Greenberg--Wiles
and local Bloch--Kato. Attach Lemma~\ref{B1.15a} and the
Hecke--Tate functional equation \cite{tate-1967}.
The Weil action and the discrete-series results, together with trace class, the two factorizations, the integral formulas and partial Hecke continuation, give the analytic hypotheses in Definition~\ref{B1.15ac}. For its
de Rham count, fix a compatible embedding inducing $w$ and
take its de Rham component over a local coefficient field
containing the image of the completion at $w$. De Rham
comparison gives dimension $a$ for that component. Its
filtration is exhaustive and separated, so the dimensions
of its graded pieces sum to $a$. These dimensions are the
jump multiplicities. For the remaining local cohomology
fields use \ref{B8.41p}. Fontaine's $B_{\rm cris}$ theorem, global reciprocity, local Bloch--Kato, Tate duality, local Euler characteristic and finiteness, and Bhatt--Scholze integral comparison give $\mathscr H_2$ as in Definition~\ref{B8.41h}.
\leanref{UpperTheory.\allowbreak{}LiteratureTheorems.\allowbreak{}literature\_\allowbreak{}nonexistence}{Development.\allowbreak{}AutomorphicArithmetic.\allowbreak{}ThetaCorrespondence.\allowbreak{}ThetaPairLifts.\allowbreak{}ThirdPlaceSeries.\allowbreak{}EveryCharacterOfLiteratureOfVWin}{639--644}

For the prime bound in $F$, apply the general progression
estimate of Thorner--Zaman for $q>2$ and $\gcd(a,q)=1$.
Their relative error is less than one at a sufficiently large
fixed power of $q$, and the main term is positive
\cite[\S1.2]{thorner-2024}. For $q=1,2$, take $p=3$ and enlarge the
absolute constant. Use the CM theorem for $j$
\cite{shimura-1971}, finite-field cuspidal mirabolic restriction
in the stated Kirillov model and Weyl branching.
The convergent Rallis formula of \ref{B8.13} uses theta
convergence, doubling and the first convergent Siegel--Weil
formula. Take the Hermitian classification and Hasse principle,
and quadratic class field theory. The Borel--Casselman
isomorphism gives the Iwahori projection for admissible $\GL_n$
over the nonarchimedean local field \cite{borel-1976,casselman-1980}.
Radford's theorem gives the polynomial shuffle ring on the
Lyndon words in characteristic zero \cite{radford-1979}.
Hecke continuation and nonvanishing at one give the analytic
Hecke fields \cite{tate-1967}.

For \eqref{eq:p1c1-log-derivative}, a primitive $\chi$ of
conductor $q>1$ is nonprincipal. Its $L$-function is entire
and nonzero on $\Re z\ge1$. Use the concentric disks about $2$
of radii $7/2$ and $15/4$. On the larger disk the bound for
$\Re z\ge-7/4$ gives $|L(z,\chi)|\ll q^4$ uniformly, while
the Euler product gives $|L(2,\chi)|\ge1/2$. Write $D$ for
the smaller closed disk and $N_D$ for the zero count there
with multiplicities. Jensen's formula gives
$N_D\le C_1\log(2q)$. The logarithmic derivative formula on
nested disks gives
\[
 \left|\frac{L'}{L}(1,\chi)
 -\sum_{\rho\in D}\frac{\operatorname{ord}_\rho L}{1-\rho}\right|
 \le C_1\log(2q),
\]
with $C_1$ enlarged if necessary. Divide by the Blaschke factors for the zeros in $D$, with respect to the larger disk. The quotient is zero-free in the smaller disk. Its logarithmic modulus has the same upper bound on the larger boundary, and its modulus at the centre is at least $|L(2,\chi)|$. The interior derivative bound for a holomorphic
logarithm gives the estimate at $1$. The derivatives of the
denominators of the Blaschke factors contribute $O(N_D)$
since $1$ is in a fixed interior disk. Every zero in the
sum of \eqref{eq:p1c1-log-derivative} belongs to $D$.
Each other zero in $D$ has $\Re\rho\le0$ or $|\Im\rho|>1$,
so $|1-\rho|\ge1$. Removing its contribution costs at most
$N_D$ in total. This proves the bound with $C=2C_1$.

For the real unitary fields use the discrete-series lowest
types and the holomorphic chamber. The closure of an admissible
$(\mathfrak g,K)$-submodule of a unitary $\mathrm U(p,q)$-representation
is invariant and has that submodule as its compact-finite
space. An isomorphism of the compact-finite modules of two
irreducible unitary representations extends, after positive
scalar normalization, to a unitary intertwiner
\cite{wallach-1992}. These supply the admissible closure and
infinitesimal equivalence fields. Cuspidal boundedness for
the form and its invariant derivatives follows by the
decay estimate above.
\leanref{GlobalTheory.\allowbreak{}gkGlobalisation\_\allowbreak{}holds}{Development.\allowbreak{}AbelianVarietyTheory.\allowbreak{}Proofs.\allowbreak{}DivisorProofs.\allowbreak{}InfinitesimalEquivalenceUnitaryGroupPQ}{737--738}

Use the two real metaplectic realizations and the adelic Weil
representation for the CM unitary pair, with Stone--von Neumann
and general linear Howe duality. For the split local member the
quadratic extension, nontrivial unitary additive character and
the restriction of the splitting character give Kudla's action
with the Levi multiplier in Definition~\ref{B1.15ac}
\cite{kudla-1994}. Conjugating translation by this action gives
$H(w)\omega(g)=\omega(g)H(wg)$. In the Levi formula the square of
the norm multiplier is the Jacobian of $x\mapsto xa$, which
cancels in the integral. The Fourier operators are isometric
with the self-dual measure in their construction. Every additive
Haar measure is a scalar multiple of that measure, so the
isometry holds for each of them. For rational theta invariance on $\operatorname{Hom}_E(V,V')=V^*\otimes_E V'$ (where $V=E^N$ and $V'=E^{N'}$ carry the standard forms with identity Gram matrix and their standard lattices $\mathcal O_E^N$, $\mathcal O_E^{N'}$, $NN'$ is even, and $(\gamma,\gamma')\in\U(V)\times\U(V')$ acts by $x\mapsto\gamma'x\gamma^{-1}$), global Poisson summation gives $\Theta(\omega(\gamma,\gamma')\varphi)=\Theta(\varphi)$. Since $E$ is a CM field, these standard forms are definite at every archimedean place, so the archimedean components of rational $\gamma,\gamma'$ lie in compact unitary groups, which preserve the Gaussian line. When $\varphi=\varphi_{\infty,0}\otimes\varphi_f$ uses the archimedean Gaussian $\varphi_{\infty,0}$, Kudla's splitting \cite[Theorem~3.1 and Proposition~4.3]{kudla-1994} acts on $\varphi_{\infty,0}$ at each CM-type embedding $\tau$ by the partner splitting-character exponents $\epsilon'n'_\tau/2$ on $\det\gamma_\tau$ and $\epsilon n_\tau/2$ on $\det\gamma'_\tau$, multiplied by the metaplectic Gaussian half-determinant character of $x\mapsto\gamma'_\tau x\gamma_\tau^{-1}$, which contributes $-N'/2$ on $\det\gamma_\tau$ (right multiplication by $\gamma_\tau^{-1}$ across $N'$ rows) and $+N/2$ on $\det\gamma'_\tau$ (left multiplication by $\gamma'_\tau$ across $N$ columns), matching the vacuum weights in Proposition~\ref{B9.39}. Thus $\omega_\infty(\gamma,\gamma')\varphi_{\infty,0}=\bigl(\prod_\tau\tau(\det\gamma)^{a_\tau}\tau(\det\gamma')^{a'_\tau}\bigr)\varphi_{\infty,0}$ with $2a_\tau=\epsilon'n'_\tau-N'$ and $2a'_\tau=\epsilon n_\tau+N$. Rewriting $\Theta(\varphi_{\infty,0}\otimes\omega_f(\gamma,\gamma')\varphi_f)=\Theta(\omega_\infty(\gamma^{-1},\gamma'^{-1})\varphi_{\infty,0}\otimes\varphi_f)$ pulls back the finite test by $(\gamma,\gamma')$ and multiplies by $\prod_\tau\tau(\det\gamma)^{-a_\tau}\tau(\det\gamma')^{-a'_\tau}$.
\leanref{GlobalTheory.\allowbreak{}Published.\allowbreak{}ThetaRationalInvariance}{Development.\allowbreak{}AbelianVarietyTheory.\allowbreak{}PublishedTheorems.\allowbreak{}GlobalTheta.\allowbreak{}ThetaRationalInvarianceHolds}{10--21}
\leanref{Formalism.\allowbreak{}MainNotions.\allowbreak{}StandardLatticesGL3.\allowbreak{}GA}{Formalism.\allowbreak{}SharpRestriction.\allowbreak{}GlobalSchurBounds.\allowbreak{}StandardLatticesGL3}{374--377}
\leanref{Formalism.\allowbreak{}MainNotions.\allowbreak{}StandardLatticesGL3.\allowbreak{}stdL}{Formalism.\allowbreak{}SharpRestriction.\allowbreak{}GlobalSchurBounds.\allowbreak{}StandardLatticesGL3}{395--397}
Use the uniqueness of the Hermitian discrete-series realization,
the one-orbit ideal invariant and the Weil representation for
the good models.

For the Siegel-set estimate, integrate in Iwasawa coordinates
$g=pak$ with $p$ in the modulus-one parabolic. In these coordinates
Haar measure has the factor $\delta_P(t)^{-1}$ and positive
torus measure $\prod_i dt_i/t_i$. Bound the test by $h(t)$ and
integrate over the compact parabolic part and $K$. Their
measures and the Haar normalizations give one constant $C$
for all $f,h$. The reduction-theory height comparison on this
set gives $C',A$ uniformly for $g\in\mathfrak S$ and rational
$\gamma$. Applying cuspidal decay in the quotient height and
using this comparison gives the asserted decay on the set.

At an infinite place $w$ let $\mathcal P_w$ be the polynomial
oscillator module for the real unitary pair with the specified
Weil splitting, with action $\omega_w$. Divide by the span
of $d\omega_w(X)f$ for source Lie algebra elements $X$ and
$\omega_w(k)f-f$ for elements of the source maximal compact.
Call the quotient target module $D_w$. On the oscillator
double cover this is the theta module for the genuine
one-dimensional character corresponding under the splitting
to the trivial source character. The pair is reductive with
compatible Cartan involutions. If $D_w\ne0$, the character
occurs as a source quotient, and Howe gives an admissible
finite-length target module with irreducible head
\cite[Theorems~1A and~2.1]{howe-1989}.

In polynomial coordinates the harmonics for a compact group
are the simultaneous kernels of the constant-coefficient
quadratic differential operators in its symplectic centralizer.
Take joint harmonics for the two compact groups in the prescribed
source compact character. This character has a single compact
type. Howe's Lemmas~3.3 and~4.1(c) give the corresponding
target type $\tau_w$ \cite{howe-1989}. Equation~(4.5) there gives
generation of the source isotypic space by the joint type
under the target enveloping algebra, and generation of the
target $\tau_w$-isotypic space by that type under the source
enveloping algebra. On the source-character quotient the
latter algebra acts by scalars. Thus $D_w[\tau_w]$ is one
copy of $\tau_w$ and generates $D_w$.

For the product of real groups the polynomial space is the
algebraic tensor product of the $\mathcal P_w$. Quotienting
by the source relations in all slots gives $D=\bigotimes_w D_w$.
If some $D_w=0$, the assertions hold for $D=0$. Otherwise the
tensor of the generating compact types occurs once and generates
$D$. Projection to this isotypic space preserves compact-stable
subspaces. Each proper target submodule has zero component of
this type, as does the sum of two proper submodules. Each local composition factor is a subquotient of the countable-dimensional
polynomial space $\mathcal P_w$, so their external tensor products
are simple by Lemma~\ref{B1.15ba}. Tensoring the finite composition
filtrations therefore bounds the length.
The Fock map onto the polynomial test space identifies stable
subspaces containing the source relations with submodules of
$D$. Pullback preserves strict inclusions and sums. This gives
the chain and proper-sum assertions.
\leanref{GlobalTheory.\allowbreak{}howeDualityReal\_\allowbreak{}holds}{Development.\allowbreak{}AbelianVarietyTheory.\allowbreak{}Proofs.\allowbreak{}DivisorProofs.\allowbreak{}HoweDualityRealHolds}{18--48}
\end{proof}

\subsection*{Proof of \ref{B1.15}}\label{prf:B1.15}\phantomsection
\begin{proof}
Fix $h\in\mathcal R$. Theorem~\ref{B6.04} gives $\mathcal V$.
Write $h_L$ for the conclusion of Theorem~\ref{B9.46}. For every number field, every finite place and every prolongation to an algebraic closure, there exists a local correspondence satisfying the laws of Definition~\ref{B9.47}. Put $O=(h,h_L,h)$ in \ref{B9.01}.
Proposition~\ref{B1.15b} gives $P,F$. Theorem~\ref{B7.37}
applied to the regular-character field gives $\mathcal G$.

Apply Theorem~\ref{B8.11} with the doubling continuation and
definite Siegel--Weil value of \ref{B8.12}. Its local quotient
maps supply the matched tests. In the definite comparison
let $I_y,J_y$ be the two integrals at a match $y$, and $c,\lambda$
the pairing constants. For $D_y=I_y-cJ_y$, the proportional
pairings give $\langle D_y,J_z\rangle=0$ and
$\langle D_y,D_z\rangle=\lambda\langle J_y,J_z\rangle$.
Choose $J_y\ne0$. If $v'=\operatorname{vol}[\U(V')]$ for
the primed source $V'$, the reverse lifts against the primed
test at $y$ are constant with values $0$ and $\|J_y\|^2/v'$.
The squared-norm identity of \ref{B8.11d} gives
$0=\lambda\|J_y\|^4/v'$. Thus $\lambda=0$ and all $D_z$ vanish
by continuity. A nonzero unprimed test and its match give
$c\ne0$. Induction on signatures in \ref{B8.11} supplies the
full integral identity.

Apply \ref{B8.15} with $F$, the vacuum generation of
\ref{B8.14} and the first formula of \ref{B8.13}, in the
range $a<b\le2a$ and $a+\operatorname{Witt}(B)<b$.
Lemmas~\ref{B8.15b} and~\ref{B8.15c} give a column chain with nonzero bottom period. By the fixed-depth test calculation of Lemma~\ref{B8.15e} and the two exchanges of Lemma~\ref{B8.15f},
every nonzero row $t$ and prescribed nonzero bottom vector $w$
have a nonzero period on an allowed column chain. Local
Bessel functional dimensions are at most one, so the period
factors as $A_w(\rho,\xi)c_{t,\Psi,\Phi}$ with real test pair
$(\rho,\xi)$ and column tests $(\Psi,\Phi)$. In a nonzero
row space choose the real pair using one nonzero row, and for
each nonzero row choose a column with $c_{t,\Psi,\Phi}\ne0$.
In a zero row space any real pair works. Theorem~\ref{B8.18}
uses this separation with the integral identity and local
matching to supply the tests in \ref{B8.01}.

Apply Theorem~\ref{B8.25} with the fifteen proofs of
Lemma~\ref{B8.25a}, available before choosing the Galois and
local correspondence data in $O$. The Siegel model uses one
polarized fan for every odd level $N\ge3$ of a fixed genus $g\ge1$.
For the unitary model use a polarized fan without self-intersection.
Its function $\phi$ gives the weighted boundary $BD$, with
Cartier ideal $(\prod_i t_i^{d\phi(\gamma_i)})$ in the boundary
coordinates $t_i$ for the primitive ray vectors $\gamma_i$.
Here $d,k\ge1$ are chosen with $d$ clearing the denominators
of $\phi$. If $\det\omega$ is a line bundle, then
$(\det\omega)^k\otimes\mathcal O(-BD)$ is ample.
Lemmas~\ref{B8.25b} and~\ref{B8.25c} give the common size
exponent and the nonterminal covers. Theorem~\ref{B8.26c}
supplies $C_T$ uniform in the field and tower. Take
$C_{\rm deep}=C_T+120$ in \ref{B8.26}. The resulting bound
$h_X$ of \ref{B8.24} supplies $C_0\ge0,x_0$ for each eligible
$F_m,C_H$, with the same $x_0$ valid for every $C\ge C_0$
at every eligible point and layer. The transfer of
\ref{B8.34} uses the functional row of \ref{B8.23} on the same
extension.
\leanref{UpperTheory.\allowbreak{}LiteratureTheorems.\allowbreak{}published\_\allowbreak{}results\_\allowbreak{}arithmetic}{Development.\allowbreak{}AutomorphicArithmetic.\allowbreak{}ThetaCorrespondence.\allowbreak{}ThetaPairLifts.\allowbreak{}ThirdPlaceSeries.\allowbreak{}EveryCharacterOfLiteratureOfVWin}{781--809}
\leanref{SelfIntersectionVolumeEstimates.\allowbreak{}topCovolSlim2\_\allowbreak{}use\_\allowbreak{}alt\_\allowbreak{}alt}{Development.\allowbreak{}AutomorphicArithmetic.\allowbreak{}ThetaCorrespondence.\allowbreak{}ThetaPairLifts.\allowbreak{}ThirdPlaceSeries.\allowbreak{}TopCovolSlim2Use}{16--22}

For each $c\in\R$, use the parameters $p_0(h_L,c)$ of Theorem~\ref{B8.01}. At each point the norm choices carry the rationality law for the period maps and a grid of parameters $(1,c,b)$. The coefficient field contains the grid's prescribed layer-character values and roots of unity, and its abelian character list agrees with the two sides of the twisting data. Use the crystalline lattice realization in $\mathscr H_2$ supplied by $P$. On the cyclotomic extensions in the small-prime calculation, Lemma~\ref{B8.41ca} gives the functorial Wach lattices and the integral action at the lower roots used in Lemma~\ref{B8.41c}. Apply Theorem~\ref{B8.41} with these components. This gives the linked
bound at twisted indices above $2$ and also above $3$ when $3$
is outside the small-prime set of the point's host.
Theorem~\ref{B8.01} gives an element of $\mathcal H$ from
this bound, $P,F$, the two theta conclusions, $h_X$ and $h_L$.
Apply Theorem~\ref{B1.16}.
\end{proof}

\subsection*{Proof of \ref{B1.17}}\label{prf:B1.17}\phantomsection
\begin{proof}
Take $h\in\mathcal H$. Apply Theorem~\ref{B7.54} with Green's assertion $\mathcal G$ and the family class, selection and estimates supplied by $h$.  This gives the connected-family assertion used in Theorem~\ref{B1.18}, with chain exponents $(e_D,e_L)=(2,0)$.
Indeed, fix $\lambda>1$ and a request $A$ for which hosts exist. Consider the families $F_m$ with $\lambda$-windows that satisfy $A$ at every point indexed by $i\in I_{F_m}(x)$ for $x>x_1(F_m)$. Suppose each such family has a threshold $x_0$ such that, for all $x\ge x_0$ and all $i\in I_{F_m}(x)$, the host field $E$, truncation index $t$ and weight $b$ at that point satisfy Definition~\ref{B6.01} with bound $\sqrt b$. Thus for every choice of local data the lengths $s_{\ell,e}(E)$ are finite for $1\le e<t$, and every finite prime sum in \eqref{eq:p6c1-total-bound} is at most $\sqrt b$. Then for every sufficiently small $\varepsilon>0$ there is a family of parameter $\varepsilon$ in this class and a choice of chain bounds valid at every sufficiently large window point, as in Theorem~\ref{B5.33}.
\leanref{GlobalTheory.Fc.forall\_goal\_of\_shape\_of\_hyps}{Development.\allowbreak{}AutomorphicArithmetic.\allowbreak{}ThetaCorrespondence.\allowbreak{}ThetaPairLifts.\allowbreak{}ThirdPlaceSeries.\allowbreak{}ShapeOf}{31--32}

For every primitive $\psi$, Theorem~\ref{B2.61} uses
$\mathcal V$ to supply the required hosts with $\lambda=3/2$.
\leanref{UpperTheory.\allowbreak{}Top.\allowbreak{}supplyChi}{Development.\allowbreak{}AutomorphicArithmetic.\allowbreak{}CharacterBoxes.\allowbreak{}SupplyChi}{8--16}
Apply Theorem~\ref{B1.18}.
\end{proof}

\section*{Proofs for Chapter~\ref{ch:arithmetic-windows}: How arithmetic windows enter the zero-free criterion}

\subsection*{Proof of \ref{B1.18}}\label{prf:B1.18}\phantomsection
\begin{proof}
Fix a primitive $\psi$ and take $A=A_\psi$ in the selection principle.
At an indexed scale $(x,\varepsilon,J)$ write $t$ for its auxiliary rank and $TT$ for its set of top positions. 
The first clause of $A_\psi$ gives the arithmetic requirement of Definition~\ref{B2.58} for the host and $TT$ at the cyclic level $\mathcal L_\psi$, with $C=4$.
Let $E$ be the number field of this host and put $b=e^{x^\varepsilon}$. For every algebraic closure
of $E$ and every compatible family of local data from \ref{B6.01},
the request also requires finite
$\Z_\ell$-lengths of $\Sha_{\ell,e}(E)$ for every prime $\ell$
and $1\leq e<t$, with
$\sum_{\ell\in S}\log\ell\sum_{1\leq e<t}
\operatorname{length}_{\Z_\ell}\Sha_{\ell,e}(E)\leq\sqrt b$
for every finite set $S$ of rational primes.

For a family in $\mathcal C_{\lambda,A}$, each indexed scale has
coordinates $x,\varepsilon$ when $x\geq x_1(\mathcal F)+1$.
Thus $x_1(\mathcal F)+1$ serves both as the threshold in $E_{\mathcal L_\psi}(\mathcal F)$, with $C=4$, and as the threshold $x_V$ for the length bound.
Family selection now gives (2), since each selected family belongs to $\mathcal C_{\lambda,A_\psi}$ and therefore has $\lambda$-windows.

To prove (1), fix $e_D,e_L\in\R$, $\varepsilon>0$, and an arithmetic parameter family $\mathcal F$ at $\varepsilon$ satisfying the two hypotheses of (1).
For all sufficiently large $x\in W(\mathcal F)$, the tied chain gives
an index $i$ with the chain hypothesis of \ref{B2.20} at $D_{x,i}$.
The condition $E_{\mathcal L_\psi}(\mathcal F)$ gives the arithmetic
hypothesis at the same index. Thus \ref{B2.20} gives (1).
Apply \ref{B1.19} and \ref{B1.22}.
\end{proof}

\subsection*{Proof of \ref{B1.20}}\label{prf:B1.20}\phantomsection
\begin{proof}
Choose $e_D,e_L,\varepsilon_1$ from (2).
Fix a kept character and $s$ with $\sigma=\Re s>1/2$. Set
\[
\begin{gathered}
b_0=\tfrac12\min(\sigma-\tfrac12,\tfrac12),\qquad
\theta=\tfrac12+b_0,\qquad D=e_D+e_L+7,\\
\varepsilon=\min\left(\tfrac{\varepsilon_1}{2},
                      \tfrac{b_0}{|D|+1}\right).
\end{gathered}
\]
Then $0<b_0\leq1/4$, $1/2<\theta<\min(\sigma,1)$,
$0<\varepsilon<\varepsilon_1$, and $D\varepsilon\leq b_0$.
Take the family from (2) and the threshold $x_1$ in (1).
On its allowed scales above $\max(x_1,1)$,
\[
G_m^{\mathcal L}(y)\leq y^{1/2+D\varepsilon}\leq y^\theta.
\]
Choosing each family window beyond $\max(Y_0,x_1,1)$ gives
$\mathcal H_\lambda(\{y\in[1,\infty):G_m^{\mathcal L}(y)\leq y^\theta\})$.
The assumed window implication gives nonvanishing throughout $\Re s>\theta$, which contains the chosen point since $\theta<\sigma$.
\end{proof}

\subsection*{Proof of \ref{B1.23}}\label{prf:B1.23}\phantomsection
\begin{proof}
Index $F$ by $\mathcal X_{\mathcal L}$, with its characters
and moduli, weight $1$ at modulus one and $w_{\mathcal L}(\chi)>0$
elsewhere. A primitive principal character has modulus $1$, so
admissibility holds.

The $n=0$ term vanishes since $\Lambda(0)=0$.
On the remaining support, $1\leq n\leq\lfloor y\rfloor_+$ gives
$n\leq y$, hence $\max(y/n-1,0)=y/n-1$. For $y<1$ the support
is empty. At modulus one, $\chi(n)=1$ for all $n$, and the centred
term is $G_z(y)$. For each nonprincipal $\chi$, both definitions give
$S_\chi(y)+yL'(1,\chi)/L(1,\chi)$, with $S_\chi$ as in \ref{B1.26}.
The nonvanishing assertions have identical indices.
\end{proof}

\subsection*{Proof of \ref{B1.29}}\label{prf:B1.29}\phantomsection
\begin{proof}
Given $\eta>0$ and $U_0\in\R$, use $\mathrm{WH}(F;\theta,\delta)$
to choose $Y\geq\max(e^{U_0},1)$ and its envelope $R$.
Put $U=\log Y$ and $R_U(u)=R(e^u)$.
Then $U\geq U_0$, and $\exp$ maps $I_U$ bijectively onto $J_Y$,
since $e^{(1+\delta)U}=Y^{1+\delta}$.
The envelope inequalities follow.
Apply Folland's change-of-variables theorem
\cite[Theorem~2.47(a)]{folland-1999} to $u\mapsto e^u$ and
$y\mapsto R(y)y^{-\theta-1}$ on $J_Y$, extended by zero outside
$J_Y$. This specialization gives integrability and
\[
\begin{aligned}
\int_{I_U}R(e^u)e^{-\theta u}\,du
&=\int_{I_U}R(e^u)(e^u)^{-\theta-1}e^u\,du\\
&=\int_{J_Y}R(y)y^{-\theta-1}\,dy
\leq Y^\eta=e^{\eta U}.
\end{aligned}
\]
\end{proof}

\section*{Proofs for Chapter~\ref{ch:zeros-become-poles}: Zeros become poles}

\subsection*{Proof of \ref{B1.30}}\label{prf:B1.30}\phantomsection
\begin{proof}
By \ref{B1.38}, \(g\) is Borel measurable and
\(|g(u)|\le(\sum_i a_iC_{\chi_i})e^{2u}\) for \(u\ge0\).
By \ref{B1.45}, \(2\Re G_{\chi_i}=G_{\chi_i}+G_{\chi_i^{-1}}\).
Multiply by \(e^{-su}\) and integrate termwise for \(\Re s>2\)
by \ref{B1.37}.
Both characters satisfy \ref{B1.34}, since
\(\chi_i^{-1}=1\iff\chi_i=1\).
Fix \(z\) with \(\Re z>0\). By \ref{B1.34} and finiteness of \(I\), there
is a punctured neighbourhood of \(z\) on which
\(L(w,\chi_i)L(w,\chi_i^{-1})\ne0\) for every \(i\in I\).
By \ref{B1.45}, this neighbourhood misses \(Z_F\).
If \(\Re s>0\) and \(s\notin Z_F\), each summand is analytic at \(s\)
by \ref{B1.34} and \ref{B1.45}. By definition \(Z_F\) lies in
\(\Re s>0\), and \ref{B1.46} gives the residues.
\end{proof}

\subsection*{Proof of \ref{B1.34}}\label{prf:B1.34}\phantomsection
\begin{proof}
For \(\chi\ne1\), apply \ref{B1.36} with \(b=0\) to the entire
function \(f=L(\,\cdot\,,\chi)\), which has \(f(1)\ne0\). Use
\ref{B1.41} and \ref{B1.37} for the integral.

If \(\chi=1\), then \(N=1\). The entire function \(f(s)=(s-1)\zeta(s)\) for
\(s\ne1\) has \(f(1)=1\) and \(f'(1)=\gamma\) by the Laurent expansion
\cite[Section~25.2]{olver-2010}. On \(\Re s>2\), the Euler product gives
\(\zeta(s)\ne0\), so \(f'/f=\zeta'/\zeta+1/(s-1)\).
With \(b=1\), \ref{B1.36} gives the numerator
\(\zeta'/\zeta+1/(s-1)-s\gamma+1-s\) of \ref{B1.43}.
Apply \ref{B1.37} for the integral. The factor \(s-1\) preserves
zero orders away from \(1\). At \(1\), \(f(1)\) and the assigned \(\zeta(1)\) are
both nonzero.
\end{proof}

\subsection*{Proof of \ref{B1.36}}\label{prf:B1.36}\phantomsection
\begin{proof}
Since \(N(1)=0\), \(T(s)=N(s)/(s-1)\) extends analytically with
\(T(1)=N'(1)\). Near \(1\) use \(\Phi_f(s)=-T(s)/s\).
Away from \(0,1\) use the stated quotient.
The identity theorem and \(f(1)\ne0\) give isolated zeros.
At a zero \(\rho\ne0,1\), \ref{B1.35} gives
\(D(s)=m/(s-\rho)+h(s)\), with \(m=\operatorname{ord}_\rho f\).
For \(Q(s)=-1/[s(s-1)]\),
\[
 \begin{split}
 \Phi_f(s)={}&\frac{mQ(\rho)}{s-\rho}
 +m\frac{Q(s)-Q(\rho)}{s-\rho}\\
 &+Q(s)\bigl(h(s)-sD(1)+b(1-s)\bigr).
 \end{split}
\]
The second term extends with value \(mQ'(\rho)\), and the third is
analytic. The residue is \(mQ(\rho)=m/[\rho(1-\rho)]\).
If \(\rho\ne0\) and \(f(\rho)\ne0\), take \(m=0\) and \(h_\rho=\Phi_f\).
\end{proof}

\subsection*{Proof of \ref{B1.44}}\label{prf:B1.44}\phantomsection
\begin{proof}
Choose \(a_*\in(0,1]\) representing \(a\), so \(H_a(s)=\zeta(s,a_*)\)
for \(s\ne1\) \cite[Section~25.11]{olver-2010}.
The identity \(\Gamma(\bar z)=\overline{\Gamma(z)}\) and \(\pi>0\)
give conjugation for \(\Gamma_{\R}\) and its entire reciprocal
\(J=1/\Gamma_{\R}\) \cite[Section~5.2]{olver-2010}.
Since \(J(1)=1\) and \(H_a\) has residue \(1\),
\(R_a(s)=H_a(s)-J(s)/(s-1)\) extends across \(1\).
In Mathlib's normalization,
\(\lim_{s\to1,\,s\ne1}(H_a(s)-J(s)/(s-1))=H_a(1)\).
Hence the extension has \(R_a(1)=H_a(1)\).
The Hurwitz series gives \(R_a(\bar s)=\overline{R_a(s)}\) on
\(\Re s>1\). The identity theorem extends this equality to \(\C\).
Add the polar term away from \(1\) and use the extension value at \(1\).
At \(s=-2n\) with \(n\in\Z_{\ge0}\), both sides of the gamma
conjugation identity have the assigned value zero.
\end{proof}

\subsection*{Proof of \ref{B1.45}}\label{prf:B1.45}\phantomsection
\begin{proof}
For modulus \(N\), the definition gives
\[
 L(s,\chi)=N^{-s}\sum_{j\in\Z/N\Z}\chi(j)H_{j/N}(s).
\]
This is Apostol's finite Hurwitz expansion away from \(1\)
\cite[Chapter~12, (8), p.~255]{apostol-1976},
extended to \(1\) by the values in \ref{B1.44}.
For \(N=1\), the sole class \(j=0\) contributes \(H_0=\zeta\).
Conjugate the sum using \ref{B1.44} and
\(\overline{\chi(j)}=\chi^{-1}(j)\).
For nonprincipal \(\chi\), differentiate
\(L(s,\chi^{-1})=\overline{L(\bar s,\chi)}\) at \(1\)
to obtain the identity for \(D(1)\).

The real weights \(\Lambda(n)(y/n-1)_+\) and the identity
\(\overline{\chi(n)}=\chi^{-1}(n)\) give conjugation for the ramps.
The principal centring term is real. The identity for \(D(1)\) gives
conjugation for the nonprincipal centring terms.
\end{proof}

\section*{Proofs for Chapter~\ref{ch:p2c1}: Counting primes class by class}

\subsection*{Proof of \ref{B2.02}}\label{prf:B2.02}\phantomsection
\begin{proof}
The map $\{\pm1\}\to\Z/2\Z$ displayed above is an isomorphism.
Read the two signs in the order $d,q$ and write the last coordinate
additively. The inverse reads the signs in the order $q,d$.
The formulas for the distinguished elements follow by substitution.
Under $T$, the condition $\sigma_d=1$ defining $\Psi_d$ becomes $v_d=0$.
Composition with an isomorphism is a bijection of character groups,
with inverse composition with $T^{-1}$.
Evaluating $\theta_\alpha=\alpha\circ T$ at $\tau_1(s)$ for $s\in A_0$ gives the assertion about trivial restrictions. Evaluating it at $\iota_d$ gives the correspondence for $X_+$.
\end{proof}

\subsection*{Proof of \ref{B2.03}}\label{prf:B2.03}\phantomsection
\begin{proof}
The divisibility $2^{e+2}\mid Q_S$ gives $2\in S$.
Also $q\in S$ and $d\nmid Q_S$, while $q\ne d$ because their
residues modulo $8$ differ. Thus $d\notin S$.
If $p$ is prime and $p\notin S\cup\{d\}$, then $p\nmid\mathfrak m$.
The local unit subgroup defining
inertia is trivial, so $I_p=\{0\}$.

Suppose $p\in S$. Since $d\mid\mathfrak m^{(p)}$, a unit defining
inertia is $1$ modulo $d$. Its $d$-sign is therefore $1$.
A unit defining Frobenius is $p$ modulo $d$, and its $d$-sign
is $(p/d)$. For odd $p$, quadratic reciprocity and
$d\equiv3\pmod4$ give
\[
 \left(\frac{p}{d}\right)=\left(\frac{-d}{p}\right)=1.
\]
This uses the splitting condition on $Q_S$ when $p\mid Q_S$,
and the host condition when $p=q$. For $p=2$ the same conclusion
follows from $d\equiv7\pmod8$. Thus $F_p\in\Phi$ and
$I_p\subseteq\Phi$.

At $d$, the local units are those congruent to $1$ modulo $Q_Sq$, so their Artin images lie in $\{1,\iota_d\}$.
The Chinese remainder theorem supplies such a unit which is
$-1$ modulo $d$. Its $d$-sign is $(-1/d)=-1$, so its image is $\iota_d$. Lemma~\ref{B2.02} now gives
$I_d=\{0,\epsilon_d\}$. Addition by $\epsilon_d$ exchanges
$\Phi$ and its complement.
\end{proof}

\subsection*{Proof of \ref{B2.06}}\label{prf:B2.06}\phantomsection
\begin{proof}
For $n>N$ and $\rho\in\Phi$ the positive part
$[a/n-w(\rho)]_+$ is zero. Expand the von Mangoldt sum
into prime powers, using $\Lambda(p^i)=\log p$, and interchange
the finite sums. Lemma~\ref{B2.05} makes the contribution
of a prime power $p^i$ equal to
\[
 \frac{2\log p}{|I_p|}
 \sum_{\substack{g\in iF_p+I_p\\g\notin\Phi}}
 [a/p^i-w(r_t(g))]_+ .
\]
For $p\in S$, the coset $iF_p+I_p$ is contained in $\Phi$,
so the inner sum is empty. For $p\notin S\cup\{d\}$, inertia
is trivial and the contribution is
\[
 2\log p\,
 \mathbf1_{\,iF_p\notin\Phi}
 [a/p^i-w(r_t(iF_p))]_+ .
\]
If $F_p\in\Phi$, all its multiples lie in $\Phi$.
If $F_p\notin\Phi$, its first binary coordinate is $1$,
so $iF_p\notin\Phi$ precisely when $i$ is odd.
At $d$, inertia has order $2$, which gives the second
line of the definition of $\ASK_N$.

Since $p^i\le N$ implies $p,i\le N$, the rectangular cutoff
retains every nonzero term. When $d>N$, every $d$-term vanishes.
\end{proof}

\subsection*{Proof of \ref{B2.07}}\label{prf:B2.07}\phantomsection
\begin{proof}
For the trivial character every inertia average equals $1$, so $v_{\mathbf1}(n)=1$.
Fix $g\notin\Phi$ and put $W=w(r_t(g))$. Since $r_t(g)\in\Phi$,
we have $W>0$ and $a/W\le N$. Then
\[
 \sum_{n=1}^N\Lambda(n)v_\alpha(n)[a/n-W]_+
 = WG_\alpha(a/W)+R_\alpha(W),
\]
where
\[
 R_\alpha(W)=
 \begin{cases}
 a\log a-a\log W-(1+\gamma)a,&\alpha=\mathbf1,\\
 -a\ell_\alpha,&\alpha\ne\mathbf1.
 \end{cases}
\]
Indeed the terms with $n>a/W$ vanish, and multiplication of
the defining sum for $G_\alpha(a/W)$ by $W$ gives the displayed
formula.

The half-group character sum is
\[
 \sum_{g\notin\Phi}\overline{\alpha(g)}
 =\begin{cases}
 f,&\alpha=\mathbf1,\\
 -f,&\alpha=\chi_d,\\
 0,&\alpha\notin\{\mathbf1,\chi_d\}.
 \end{cases}
\]
Since $G\setminus\Phi=\epsilon_d+\Phi$, the character sum is
$\overline{\alpha(\epsilon_d)}\sum_{h\in\Phi}\overline{\alpha(h)}$,
which vanishes when $\alpha|_\Phi$ is nontrivial.
The two characters trivial on $\Phi$ are $\mathbf1$ and
$\chi_d$, whose values on the other coset are $1$ and $-1$.

Substitution in Lemma~\ref{B2.06} gives $\Delta(a,t,w)$ and the remainders. The coset sum leaves $a\ell_{\chi_d}$ from the
nontrivial characters. In the trivial remainder, change variables by
$\rho=r_t(g)$ to obtain
$a\log a-(1+\gamma)a-(a/f)\sum_{\rho\in\Phi}\log w(\rho)$.
Together these give $\Main(a,w)$.
\end{proof}

\subsection*{Proof of \ref{B2.08}}\label{prf:B2.08}\phantomsection
\begin{proof}
The identity $-e_1(s+u-t)=e_1(t)+\tau(-s)+\tau(-u)$ gives
\[
 \sum_{u\in A_0}\alpha(-e_1(s+u-t))
 =\alpha(e_1(t))\alpha(\tau(-s))
 \sum_{u\in A_0}\alpha(\tau(-u)).
\]
For $\alpha\in P$ the last two factors are $1$ and $|A_0|$.
Otherwise the restriction to $A_0$ is a nontrivial character.
Translation by an element on which it differs from $1$ shows
that its sum is zero. Multiply the identity for each $\alpha$ by
$bG_\alpha(y)/f$, then sum over $s\in T_2$ and $\alpha\in\widehat G$.
The surviving coefficient is $|A_0|b|T_2|/f=|A_0|b\delta/2$.
\end{proof}

\subsection*{Proof of \ref{B2.09}}\label{prf:B2.09}\phantomsection
\begin{proof}
The condition on $a/b$ gives $0<(x-1)b\le a\le(x+1)b$.
The inverse image of the thick position $((0,1),s+u)$
under $r_t$ is $e_1(s+u-t)$.
Since $\overline{\alpha(g)}=\alpha(-g)$, replacing
$w(r_t(g))G_\alpha(a/w(r_t(g)))$ by $bG_\alpha(a/b)$
at those positions, and by zero at the other positions,
gives $\Delta_{\rm thick}(u,t;a/b)$.

At a thick position, \eqref{eq:b2-lip} bounds the change by
\[
 \frac{3a}{b}|w(\rho)-b|
 \le3(x+1)(\eta+1).
\]
At a position $\rho\in\Phi\setminus\mathcal T_u$, we have
\[
 \frac{X'}3w(\rho)
 \le\frac23 b'X'
 \le\frac23xb
 \le(x-1)b
 \le a.
\]
Thus $a/w(\rho)\ge X'/3$, and \eqref{eq:b2-thin} gives
\[
 \left|w(\rho)G_\alpha(a/w(\rho))\right|
 \le a\theta\le(x+1)b\theta.
\]
There are $f$ positions and $|\widehat G|=|G|=2f$
characters. Divide the sum of these bounds by $f$ and
use $|\alpha(g)|=1$. The resulting bound on $|\Delta(a,t,w)-\Delta_{\rm thick}(u,t;a/b)|$ gives the inequality of real parts.
\end{proof}

\subsection*{Proof of \ref{B2.10}}\label{prf:B2.10}\phantomsection
\begin{proof}
Write $\Psi(z)=\sum_{1\le n\le z}\Lambda(n)$.
Rosser and Schoenfeld prove $\Psi(z)<1.03883z$ for $z>0$
\cite[Theorem~12]{rosser-schoenfeld-1962}.
For $z>0$ this implies $\Psi(z)\le3z$, which also holds at $z=0$ since $\Psi(0)=0$.
This bound gives \eqref{eq:b2-lip} for every $D$ and $\ell$.

Fix $a>0$. For $\alpha\ne\mathbf1$ and $v>0$, multiplication of the
definition of $G_\alpha$ by $v$ gives
\[
 vG_\alpha(a/v)=
 \sum_{n\ge1}\Lambda(n)v_\alpha(n)[a/n-v]_+ +a\ell_\alpha.
\]
The constant $a\ell_\alpha$ cancels
in the difference for $v,v'$. For $n\le a/\min(v,v')$, each
positive part changes by at most $|v-v'|$. All other summands vanish.
Summing the changes with $|v_\alpha(n)|\le1$ gives at most $\Psi(a/\min(v,v'))|v-v'|$. The Chebyshev bound gives \eqref{eq:b2-lip}.

For the trivial character, take $0<v\le v'$ and let $B$
be the decrease of $\sum_n\Lambda(n)[a/n-v]_+$ when $v$
is replaced by $v'$. Then
\[
 0\le B\le3(a/v)(v'-v),\qquad
 0\le a\log(v'/v)\le(a/v)(v'-v).
\]
The difference $vG_z(a/v)-v'G_z(a/v')$ equals
$B-a\log(v'/v)$. Both nonnegative quantities are bounded
by $3(a/v)(v'-v)$, so their difference has absolute
value at most that bound. Interchanging $v,v'$ covers
the other order.
\leanref{Formalism.\allowbreak{}LFunctionChainParameters.\allowbreak{}ThickSetMeanBounds.\allowbreak{}AN.\allowbreak{}psi\_\allowbreak{}le\_\allowbreak{}three}{Formalism.\allowbreak{}SizeWindowBrackets.\allowbreak{}ChainSlopeWindows.\allowbreak{}LipschitzWeightAndMinimum}{668--691}
\leanref{Formalism.\allowbreak{}LFunctionChainParameters.\allowbreak{}ThickSetMeanBounds.\allowbreak{}lipschitz\_\allowbreak{}weight\_\allowbreak{}all}{Formalism.\allowbreak{}SizeWindowBrackets.\allowbreak{}ChainSlopeWindows.\allowbreak{}LipschitzWeightAndMinimum}{957--965}

For each $t\in TT$ the hypotheses imply $a_t>0$.
Theorem~\ref{B2.07} and Lemma~\ref{B2.09} therefore give,
for every $u\in A_0$ and $0\le j<m$,
\[
 \Re\Delta_{\rm thick}(u,t;a_t/b)
 \le \ASK_N(a_t,t,w_{u,j})-\Re\Main(a_t,w_{u,j})+\frac{2b}{x}.
\]
The two cost inequalities bound each error contribution by $b/x$.
After summation over $t,j$, the assumed bound on the gap yields
\[
 \sum_{t\in TT}\sum_{j=0}^{m-1}
 \Re\Delta_{\rm thick}(u,t;a_t/b)
 \le\mathscr E+\frac{2Q}{x}.
\]
Sum over $u\in A_0$. By Lemma~\ref{B2.08}, the left side becomes
\[
 |A_0|\,m b\,\frac{\delta}{2}
 \sum_{t\in TT}\Re\sum_{\alpha\in P}
 \alpha(e_1(t))G_\alpha(a_t/b).
\]
Divide by $|A_0|$. Multiplying the assumed lower bound for the sum
over $t\in TT$ by $mb\delta/2\ge0$ gives
$Q(\delta/2)Y_*\le\mathscr E+2Q/x$.
\end{proof}

\section*{Proofs for Chapter~\ref{ch:p2c2}: Reading off the target sum}

\subsection*{Proof of \ref{B2.12}}\label{prf:B2.12}\phantomsection
\begin{proof}
For \(\alpha=1\), sum \(G_z(y_t)\ge G_z(x)-C_d\) from \ref{B2.11} and use \(C_d\le R\).
For \(\alpha\ne1\) in \(K_p\), \(\alpha(e_1(t))=\alpha(0,t)\). Put \(w=w_c(\alpha)\) and \(S_\alpha=\sum_{t\in TT}\alpha(0,t)\). Then \(|S_\alpha-wn_T|\le\rho n_T\) by hypothesis, with equality to zero on the left when \(\alpha|_{\{0\}\times A}=1\). We have
\[
 \begin{split}
 \sum_{t\in TT}\alpha(0,t)G_\alpha(y_t)-wn_TG_\alpha(x)
 ={}&\sum_{t\in TT}\alpha(0,t)(G_\alpha(y_t)-G_\alpha(x))\\
 &+(S_\alpha-wn_T)G_\alpha(x).
 \end{split}
\]
By \ref{B2.11}, its absolute value is at most
\(n_T\{\log(x+1)+C_d+K_o+\rho x(\log x+C_d+K_o)\}\le n_TR\).
Taking real parts gives the bound.

For \(\alpha\in P\setminus K_p\), \ref{B2.11} gives
\(|G_\alpha(y_t)|\le |G_\alpha(x)|+\log(x+1)+C_d+L\le |G_\alpha(x)|+R\).
Since \(|\alpha(e_1(t))|=1\), each real part is at least \(v_\alpha-R\). For \(K_p\subseteq P\), sum \(v_\alpha\) separately over \(K_p\) and \(P\setminus K_p\) to obtain the identity.
\end{proof}

\subsection*{Proof of \ref{B2.13}}\label{prf:B2.13}\phantomsection
\begin{proof}
Fix \(u\in U\). By \ref{B5.43}, \(\ASK_{\mathscr D}(u)-\Main_{\mathscr D}(u)\le E_8(u)\), with \(E_8\) defined in \ref{B5.42}. The bijection \(e_T\) gives \(|\mathcal T|=n_T\). Apply \ref{B5.45} with the assumed bounds on \(B_0,N_c\) to get
\[
 E_8(u)\le x^{C_{\rm err}\varepsilon}
 \bigl[n_Ta_E(1+m^2/k)+m(B_0+N_c)\bigr]
 \le x^{C_{\rm err}\varepsilon}\mathcal B.
\]
By \ref{B5.46}, the weight and constant identifications and \(f=|\Phi|\) give the arithmetic main sum after reindexing by \(e_T,e_P\). Subtract it from the assumed expression for \(\ASK_{\mathscr D}(u)\). The result is \(\operatorname{gap}(\operatorname{tr}(u))\) by \ref{B2.10}.
\end{proof}

\subsection*{Proof of \ref{B2.14}}\label{prf:B2.14}\phantomsection
\begin{proof}
Each local factor in \ref{B2.04} is an average of complex numbers of modulus one. Hence their product satisfies \(|V_\alpha(n)|\le1\) for \(n\ge1\). The positive part \([y/n-1]_+\) changes by at most \(|y-z|/n\) between \(y\) and \(z\), and both positive parts vanish for \(n>\max(y,z)\). Thus for \(\alpha\ne1\) and \(y,z\ge1\),
\[
 |G_\alpha(y)-G_\alpha(z)|
 \le |y-z|\left(\sum_{n\le\max(y,z)}\frac{\Lambda(n)}n+|\ell_\alpha|\right).
\]
Apply \ref{lem:p2c2-mertens} with \(r=\lfloor\max(y,z)\rfloor\) to obtain \eqref{eq:p2c2-deriv} with \(C_d\ge\log4+4\).
On \((r,r+1)\) for an integer \(r\ge1\), differentiation of \ref{B1.26} gives
\[
 G_z'(y)=\sum_{n\le r}\frac{\Lambda(n)}n-\log y+\gamma .
\]
By \ref{lem:p2c2-mertens} and \(1\le y/r\le2\), its absolute value is at most \(C_d=\log4+4+\log2+|\gamma|\). Each new summand vanishes at the integer where it enters, so \(G_z\) is continuous there. Partition \([\min(y,z),\max(y,z)]\) at its interior integers. On each piece the mean value theorem bounds the absolute increment by \(C_d\) times the length. Summing gives \eqref{eq:p2c2-zeta-lip}.
\leanref{Formalism.\allowbreak{}LFunctionChainParameters.\allowbreak{}ThickSetMeanBounds.\allowbreak{}derivative\_\allowbreak{}integral\_\allowbreak{}all}{Formalism.\allowbreak{}SizeWindowBrackets.\allowbreak{}ChainSlopeWindows.\allowbreak{}LipschitzWeightAndMinimum}{966--982}

Put \(A_1=5+C_d+|K_o|+|M'|(4+C_d+|K_o|)+4|C_E|\) and \(R=A_1x^{1/4+\varepsilon}\).
Choose \(x_1\ge3\) so that for \(x\ge x_1\),
\begin{equation}\label{eq:p2c2-arithmetic-threshold}
 26\le x^{3\varepsilon/8},\qquad |M'|A_1\le x^{1/4}.
\end{equation}
Use \(y_t=a_t/b\), \(\rho=|M'|x^{-1+\varepsilon}\), and the bounds \(|K_o|\) and \(L=|C_E|\log x\) in \ref{B2.12}. They satisfy the character hypotheses by \eqref{eq:p2c2-chars}. The inequality \(\log u\le u-1\) for \(u>0\) gives \(\log x\le4x^{1/4}\) at \(u=x^{1/4}\) and \(\log2\le1\). Hence \(\log(x+1)\le5x^{1/4}\), and
\[
 \rho x(\log x+C_d+|K_o|)
 \le |M'|(4+C_d+|K_o|)x^{1/4+\varepsilon}.
\]
Since \(x^\varepsilon\ge1\), these bounds give the three requirements on \(R\) in \ref{B2.12}.

Put \(G_*=G_m^A(x)-\Sigma_A(x)\). Sum \ref{B2.12} over \(P\) and use its identity for \(\sum_{\alpha\in P}v_\alpha\). By \eqref{eq:p2c2-chars} and \eqref{eq:p2c2-arithmetic-threshold}, \(|P|R\le x^{1/2+\varepsilon}\). In particular,
\begin{equation}\label{eq:p2c2-top-lower}
 n_T(G_*-4x^{1/2+\varepsilon})
 \le\sum_{t\in TT}\Re\sum_{\alpha\in P}\alpha(e_1(t))G_\alpha(a_t/b).
\end{equation}
Apply \ref{B2.10} with \eqref{eq:p2c2-top-lower} and the assumed error \(x^{C\varepsilon}\mathcal B\). This gives
\begin{equation}\label{eq:p2c2-balance}
 Q\frac{\delta_2}{2}(G_*-4x^{1/2+\varepsilon})
 \le x^{C\varepsilon}\mathcal B+\frac{2Q}{x}.
\end{equation}
In \ref{B5.55} take \(G_x=G_*\) and \(E=x^{C\varepsilon}\mathcal B\). The size conditions follow from \eqref{eq:p2c2-sizes} and \(T_2\subseteq A\), and the threshold condition is \eqref{eq:p2c2-arithmetic-threshold}. Thus \(G_*\le x^{1/2+(C+3)\varepsilon}\).
\end{proof}

\subsection*{Proof of \ref{B2.15}}\label{prf:B2.15}\phantomsection
\begin{proof}
Put \(M_H=Q_Sqd\) and \(r_H=T\circ\operatorname{art}\), with \(T\) as in \ref{B2.02}. On units modulo \(M_H\), \(\theta_\alpha\circ\operatorname{art}=\alpha\circ r_H\). The Dirichlet character given by this pullback has primitive representative \(\chi_\alpha\).

Let \(f_\alpha\mid M_H\) be its conductor. For a prime \(p\), inertia \(I_p\) is the image under \(r_H\) of the kernel of reduction of units modulo the prime-to-\(p\) part of \(M_H\). Thus \(\alpha|_{I_p}=1\) precisely when the pullback factors through this part, or equivalently when \(p\nmid f_\alpha\). In that case the Frobenius unit is congruent to \(p\) modulo \(f_\alpha\), so for \(i\ge1\),
\[
 \frac1{|I_p|}\sum_{h\in I_p}\alpha(iF_p+h)=\chi_\alpha(p)^i.
\]
When \(p\mid f_\alpha\), translation by an \(h_0\in I_p\) with \(\alpha(h_0)\ne1\) multiplies the sum by \(\alpha(h_0)\), so the sum is zero, as is \(\chi_\alpha(p)\). Multiplication over prime powers gives \(V_\alpha(n)=\chi_\alpha(n)\) for \(n\ge1\). Compatibility supplies the logarithmic derivative for \(\alpha\ne1\).

Equal primitive representatives induce equal characters on units modulo \(M_H\), so Artin surjectivity gives equal host characters. The trivial character reduces to the primitive trivial character. By injectivity and the character isomorphism of \ref{B2.02}, \(\chi_\alpha=1\) exactly when \(\alpha=1\). In that case both sums equal \(G_z\). Otherwise, substitute the local values and the logarithmic derivative into \ref{B2.04}. In \ref{B1.26}, \([y/n-1]_+\) equals \(y/n-1\) for \(1\le n\le y\) and vanishes for \(n>y\), so the sums agree.
\end{proof}

\subsection*{Proof of \ref{B2.17}}\label{prf:B2.17}\phantomsection
\begin{proof}
Each \(\chi\in\mathcal K\) has a representative in \(X_+\) by \ref{B2.58}, unique by \ref{B2.15}. Thus \(K_H\subseteq P\) and each member takes value \(1\) on \(\epsilon_d\). The trivial character belongs to \(K_H\). By \ref{B2.16} and \ref{B2.58}, \([A:A_0]=o\) and \(|P|=4o\).

Put \(N_\beta=|TT\cap\beta|\) for \(\beta\in A/A_0\), and \(h=x^{-1+\varepsilon}\). The shares of \ref{B2.58} give
\[
 |N_\beta-\lambda_\beta n_T|\le hn_T,\qquad
 \lambda_0=\frac4{3o+1},\quad
 \lambda_\beta=\frac3{3o+1}\quad(\beta\ne0).
\]
For \(\alpha\in P\) nontrivial on \(A\), \(\bar\alpha(t+A_0)=\alpha(0,t)\) defines a nontrivial character of \(A/A_0\). Translation by a coset on which \(\bar\alpha\) differs from \(1\) shows that \(\sum_\beta\bar\alpha(\beta)=0\). Hence
\[
 \sum_\beta\bar\alpha(\beta)\lambda_\beta
 =\frac{3\sum_\beta\bar\alpha(\beta)+1}{3o+1}=c,
\]
and
\[
 \left|\sum_{t\in TT}\alpha(0,t)-cn_T\right|
 =\left|\sum_\beta\bar\alpha(\beta)(N_\beta-\lambda_\beta n_T)\right|
 \le ohn_T\le4ohn_T .
\]
For retained \(\alpha\ne1\), \ref{B2.15} identifies \(|\ell_\alpha|\) with a summand of \(K_o\). For \(\alpha\in P\setminus K_H\), apply the logarithmic derivative bound of \ref{B2.58} using \eqref{eq:p2c2-bounded} and \ref{B2.15} to get \(|\ell_\alpha|\le C\log x\).
\end{proof}

\subsection*{Proof of \ref{B2.18}}\label{prf:B2.18}\phantomsection
\begin{proof}
By \ref{B2.02}, \ref{B2.15}, and \ref{B2.58}, the map \(\alpha\mapsto\chi_\alpha\) gives a bijection \(K_H\simeq\mathcal K\). A retained \(\alpha\) has \(\alpha(\epsilon_d)=1\) by \ref{B2.58}. If it is trivial on \(A\), it factors through \(E_0\), where this condition leaves \(1\) and \(\chi_q\). By \ref{B2.02}, the host character \(\theta_{\chi_q}\) is the projection \(((\sigma_q,\sigma_d),a)\mapsto\sigma_q\).

The trivial character contributes \(G_z(y)\) with weight \(1\). When \(\chi_{\rm sp}\) is designated, \ref{B2.58} gives \(\pi_H(\theta_{\chi_q})=\chi_{\rm sp}\), so its contribution also has weight \(1\). In the absence of a designated character the same condition puts \(\chi_q\) outside \(K_H\). Every other retained \(\alpha\) is nontrivial on \(A\), so \(w_c(\alpha)=c\). By injectivity in \ref{B2.15}, its primitive representative is distinct from \(1\) and from the designated character when one is present. Hence \(w_\Lambda(\chi_\alpha)=c\). Apply \ref{B2.15} and reindex by the bijection. The cyclic level has \(\mathcal K=\mathcal X\) and no designated character.
\end{proof}

\subsection*{Proof of \ref{B2.20}}\label{prf:B2.20}\phantomsection
\begin{proof}
Put \(C'=\max(e_D+e_L,0)+3\). Take the threshold of \ref{B2.14} for \(C'\), \(c=(3o+1)^{-1}\), \(M'=4o\), \(C_E=C\), and \(K_o=\sum_{\chi\in\mathcal K}|\kappa_\chi|\). Here \(C'\ge0\) and \(|c|\le1\). Since \(\varepsilon>0\) and \(\log x/x^\varepsilon\to0\), choose \(x_1\ge3\) above this threshold so that for all \(x\ge x_1\),
\begin{equation}\label{eq:p2c2-host-threshold}
 \log x\le x^\varepsilon,\qquad
 200(1+K)^2(1+c_{HH})(1+e_K)\le x^\varepsilon,\qquad
 1+4o\le x^\varepsilon .
\end{equation}
Take a datum and compatible \(\ell\) from \ref{def:p2c2-core-at-host}. The window conditions of \ref{B5.45} give \(e_{\rm Div},e_{\rm Loss},c_{HH}(\mathcal C),e_K\ge0\). Thus \(e_D,e_L\ge0\) and \(C'=e_D+e_L+3\). Since \(c_{HH}(\mathcal C)\le c_{HH}\), \eqref{eq:p2c2-host-threshold} gives the large-\(x\) condition of \ref{B2.13} for \(\mathcal C\). For each \(s\in A_0\), choose \(u\in U\) with \(\operatorname{tr}(u)=s\). Apply \ref{B2.13} and use \(x\ge1\) and \(\mathcal B\ge0\) to get
\[
 \operatorname{gap}(s)
 \le x^{(e_{\rm Div}+e_{\rm Loss}+3)\varepsilon}\mathcal B
 \le x^{C'\varepsilon}\mathcal B .
\]
The same function \(a:A\to\R\) serves every translation, as required in \ref{B2.10}.

Use the local conditions of \ref{B2.03} and the size and weight conditions of \ref{def:p2c2-core-at-host}. Apply \ref{B2.14} with \(K_p=K_H\), using \ref{B2.17} for the character conditions \eqref{eq:p2c2-chars}. Rewrite its left side by \ref{B2.18} to get
\[
 G_{m,\Lambda}(x)\le x^{1/2+(e_D+e_L+6)\varepsilon}+\Sigma_A(x).
\]
By \ref{B2.19}, \(\Sigma_A(x)\le4o\,x^{1/2+\varepsilon}\). Since \(e_D+e_L\ge0\), the right side is at most \((1+4o)x^{1/2+(e_D+e_L+6)\varepsilon}\). The last bound in \eqref{eq:p2c2-host-threshold} gives \(G_{m,\Lambda}(x)\le x^{1/2+(e_D+e_L+7)\varepsilon}\).
\end{proof}

\section*{Proofs for Chapter~\ref{ch:p2c3}: Choosing the auxiliary primes \(q\) and \(d\)}

\subsection*{Proof of \ref{B2.23}}\label{prf:B2.23}\phantomsection
\begin{proof}
Because \(8\mid m_0\), we have \(m_0\geq8\) and \(\log m_0\geq2\). Let \(c_0^{\rm ZF}>0\) be the classical zero-free constant for Dirichlet \(L\)-functions, so that every modulus \(q\geq2\) has at most one real zero \(\beta>1-c_0^{\rm ZF}/\log q\) among all nonprincipal characters modulo \(q\), and put \(c_Z=\min\{c_0^{\rm ZF},1/20\}>0\). Call \((\chi_1,\beta_1)\) an exceptional pair modulo \(m_0\) if \(\chi_1\bmod m_0\) satisfies \(\chi_1^2=1\), \(\chi_1\ne1\), \(L(\beta_1,\chi_1)=0\), and \(1-1/(8\log m_0)<\beta_1<1\). Select the optional exceptional pair modulo \(m_0\) by the two-case rule:
\begin{enumerate}
\item if there exists an exceptional pair \((\chi_1,\beta_1)\) satisfying \(1-c_Z/\log m_0<\beta_1<1\), choose that pair;
\item otherwise, if at least one exceptional pair modulo \(m_0\) exists (so every such pair satisfies \(1-1/(8\log m_0)<\beta_1\leq1-c_Z/\log m_0\)), choose one such pair \((\chi_1,\beta_1)\); if none exists, choose no exceptional pair.
\end{enumerate}
By construction, choosing no exceptional pair occurs precisely when no exceptional pair modulo \(m_0\) exists, and whenever an exceptional zero lies in \((1-c_Z/\log m_0,1)\) it is the chosen pair. Set \(\eta(m_0)=\min\{1,(1-\beta_1)\log m_0\}\) when a pair \((\chi_1,\beta_1)\) is chosen and \(\eta(m_0)=1\) otherwise; Siegel's theorem gives an absolute constant \(c_{\rm S}>0\) with \(\eta(m_0)\geq c_{\rm S}m_0^{-1}\) for all \(m_0\geq8\).

Put \(\psi(y,\chi)=\sum_{n\leq y}\Lambda(n)\chi(n)\) and \(\mathrm{main}(y,\chi)=\mathbf1_{\chi=\chi_0}y-\mathbf1_{\chi=\chi_1}y^{\beta_1}/\beta_1\) (dropping the second term when no exceptional pair is chosen). Split \(\sum_{\chi\bmod m_0}|\psi(y,\chi)-\mathrm{main}(y,\chi)|\) at an absolute threshold \(\log m_0\leq c_A\sqrt{\log y}\) with \(0<c_A\leq c_Z\). In the near range \(\log m_0\leq c_A\sqrt{\log y}\), the classical zero-free region modulo \(m_0\) bounds \(|\psi(y,\chi)-\mathrm{main}(y,\chi)|\ll y\,e^{-c'\sqrt{\log y}}\) for every \(\chi\bmod m_0\): in case~(1) the sole zero in \((1-c_Z/\log m_0,1)\) is subtracted in \(\mathrm{main}(y,\chi_1)\), while in case~(2) there is no zero in \((1-c_Z/\log m_0,1)\) and the subtracted term satisfies \(y^{\beta_1}/\beta_1\leq2y\,e^{-(1-\beta_1)\log y}\leq2y\,e^{-c_Z\log y/\log m_0}\leq2y\,e^{-\sqrt{\log y}/2}\) since \(\beta_1>1-1/(8\log m_0)\geq1/2\) and \(\log m_0\leq c_Z\sqrt{\log y}\). Summing over the \(\varphi(m_0)\leq m_0\leq e^{c_A\sqrt{\log y}}\) characters and inserting \(1\leq c_{\rm S}^{-1}m_0\eta(m_0)\) gives error \(O\bigl(y\,\eta(m_0)e^{-c_1\sqrt{\log y}}\bigr)\) after choosing \(c_A\) small enough against the zero-free constant.

In the far range \(\log y\geq C_0\log m_0\log(\log m_0+2)\) with \(\log m_0>c_A\sqrt{\log y}\), apply the smoothed explicit formula at height \(T=m_0^{12}\) and difference it over \([y,(1+\delta)y]\) with \(\delta=m_0^{-5}\) to obtain, for each \(\chi\bmod m_0\),
\[
 \bigl|\psi(y,\chi)-\mathrm{main}(y,\chi)\bigr|
 \ll yE/\delta + y\sum_{\rho\in Z(\chi)}'\bigl(m(\rho)-\mathbf1_{(\chi,\rho)=(\chi_1,\beta_1)}\bigr)y^{\Re\rho-1}
 + (\delta y+1)\log(2y) + \mathbf1_{\chi=\chi_1}y^{1-\kappa/2},
\]
where \(Z(\chi)\) consists of the zeros of \(L(s,\chi)\) in \(\Re\rho\geq1-\kappa/2\), \(|\Im\rho|\leq T\), and \(E\ll(m_0T)^{C_1}y^{-\kappa/2}+T^{-1}\log(2y)\log^2(m_0T)\). Summing the non-zero-sum terms over \(\chi\bmod m_0\) gives \(O(y\,m_0^{-2})\ll y\,\eta(m_0)e^{-c_2\log y/\log m_0}\) once \(C_0\) is large enough, using \(m_0^{-1}\ll\eta(m_0)\) and \(m_0^{-1}\leq e^{-c_A^2\log y/\log m_0}\). For the remaining zero sum \(\sum_\chi W(\chi)\), combine the logarithmic zero-density bound \(N(\sigma,T,m_0)\ll(m_0T)^{a(1-\sigma)}(\log(m_0T))^B\) with the lossy Deuring--Heilbronn repulsion of Theorem~\ref{B2.22} (which introduces a factor \(\eta(m_0)(\log(m_0T))^N\) when \(1-\beta_1\) is small). Writing \(e^{-c_3\log y/\log m_0}=e^{-(c_3/2)\log y/\log m_0}\cdot e^{-(c_3/2)\log y/\log m_0}\), the hypothesis \(\log y\geq C_0\log m_0\log(\log m_0+2)\) with \(C_0\) sufficiently large gives \((\log m_0)^{B+N+1}e^{-(c_3/2)\log y/\log m_0}\leq1\), absorbing all powers of \(\log m_0\) and yielding \(\sum_\chi W(\chi)\ll\eta(m_0)e^{-(c_3/2)\log y/\log m_0}\).

Gluing the near and far bounds across \(y\in[\sqrt t,t]\) gives \(\sum_{\chi\bmod m_0}|\psi(y,\chi)-\mathrm{main}(y,\chi)|\ll y\,\eta(m_0)\bigl(e^{-c'\log y/\log m_0}+e^{-c'\sqrt{\log y}}\bigr)\) on \([\sqrt t,t]\). Character orthogonality modulo \(m_0\) (using \(\chi_1(a_0)^{-1}=\chi_1(a_0)\in\{\pm1\}\) for \((a_0,m_0)=1\)) and Abel partial summation against \(1/\log u\), together with \(H_{a_0}(t)\gg\eta(m_0)\operatorname{li}_2(t)/\varphi(m_0)\), yield the stated relative bound on \(\pi(t;m_0,a_0)\) (compare also \cite[Theorem~1.1]{thorner-2024} in the full Linnik range \(t\geq m_0^{c_2}\) under the same two-case selection rule).
\leanref{Formalism.\allowbreak{}HermitianLattices.\allowbreak{}Ext.\allowbreak{}PrimesInProgressions.\allowbreak{}psi\_\allowbreak{}relative\_\allowbreak{}near\_\allowbreak{}range}{Formalism.\allowbreak{}GL3MeanSquares.\allowbreak{}GenusMeanAsymptotics.\allowbreak{}JacobiCharacterSums.\allowbreak{}PNTRangeA}{161--240}
\leanref{Formalism.\allowbreak{}HermitianLattices.\allowbreak{}Ext.\allowbreak{}PrimesInProgressions.\allowbreak{}psi\_\allowbreak{}relative\_\allowbreak{}far\_\allowbreak{}range}{Formalism.\allowbreak{}GL3MeanSquares.\allowbreak{}GenusMeanAsymptotics.\allowbreak{}JacobiCharacterSums.\allowbreak{}PNTBAsm}{15--167}
\leanref{Formalism.\allowbreak{}HermitianLattices.\allowbreak{}Ext.\allowbreak{}PrimesInProgressions.\allowbreak{}glue\_\allowbreak{}ranges}{Formalism.\allowbreak{}GL3MeanSquares.\allowbreak{}GenusMeanAsymptotics.\allowbreak{}JacobiCharacterSums.\allowbreak{}PNTGlue}{22--57}
\leanref{Formalism.\allowbreak{}HermitianLattices.\allowbreak{}Ext.\allowbreak{}PrimesInProgressions.\allowbreak{}prime\_\allowbreak{}count\_\allowbreak{}of\_\allowbreak{}psi}{Formalism.\allowbreak{}GL3MeanSquares.\allowbreak{}GenusMeanAsymptotics.\allowbreak{}JacobiCharacterSums.\allowbreak{}PrimeCountOfPsi}{12--92}
\end{proof}

\subsection*{Proof of \ref{B2.23a}}\label{prf:B2.23a}\phantomsection
\begin{proof}
Use McCurley's zero-free region for the product of all Dirichlet \(L\)-functions at one modulus, in the form of \cite[(1.2)]{benli-2026}. Take \(k=3m\), so that the result applies also when \(m=1,2\). For \(t\in\R\),
\[
 \max\{k,k|t|,10\}\leq \bigl(m(|t|+2)\bigr)^4.
\]
Hence \(c=1/50\) puts the stated closed region inside McCurley's open region at modulus \(k\). There the product has at most one zero, a simple real zero of a real nonprincipal character. Induction from modulus \(m\) to modulus \(k\) is injective on characters. It multiplies each \(L\)-function by Euler factors that are nonzero for \(\Re s>0\), so it preserves the zeros in the stated region and their multiplicities. The conclusion follows on restricting the induced characters to modulus \(m\).
\end{proof}

\subsection*{Proof of \ref{B2.25}}\label{prf:B2.25}\phantomsection
\begin{proof}
Apply the second assertion of Theorem~\ref{B2.24} with \(\nu=1\).
Decrease its constant to obtain \(0<c_*\leq1\) such that
\(1-\beta_1\geq c_*/m_0\) for every exceptional pair modulo \(m_0\).

For \(\log U\geq4\), the corrected main term from Theorem~\ref{B2.23} satisfies
\[
 H_{a_0}(U)\geq
 \frac{c_*U}{2m_0^2\log U}.
 \tag{1}\label{eq:aux-count-main}
\]
Integration over \([U/2,U]\) gives
\(\operatorname{li}_2(U)\geq U/(2\log U)\).
For a selected exceptional zero, \(\beta_1>1/2\) gives \(U^{\beta_1}\geq\sqrt U>2\).
This proves \eqref{eq:aux-count-main} when the pair is absent or \(\chi_1(a_0)=-1\).
If \(\chi_1(a_0)=1\), the gap
\(\operatorname{li}_2(U)-\operatorname{li}_2(U^{\beta_1})\)
is bounded below by \((U-U^{\beta_1})/\log U\).
The inequality
\(1-e^{-v}\geq\tfrac12\min\{v,1\}\) for \(v\geq0\), together with
\(1-\beta_1\geq c_*/m_0\), gives
\(U-U^{\beta_1}\geq c_*U/(2m_0)\).
Finally \(\varphi(m_0)\leq m_0\), which proves
\eqref{eq:aux-count-main}.

Write \(L=\log Y\) and \(u=\log U\).
Admissibility gives
\[
 \log m_0\leq\frac{4u}{L^2},
 \qquad
 Y^{\alpha/6}\leq u\leq Y^{\alpha/5}L^2.
 \tag{2}\label{eq:aux-count-scales}
\]
Thus
\(\log(\log m_0+2)=O_\alpha(L)\) and
\(u/\log m_0\geq L^2/4\).
For sufficiently large \(Y\), these inequalities place \(U\) in the range of Theorem~\ref{B2.23} and make its relative error at most \(1/2\).
Consequently \(\pi(U;m_0,a_0)\geq\tfrac12H_{a_0}(U)\).

Uniformly in the allowed \(u\),
\[
 \frac1u
 \log\!\left(\frac{c_*U}{2m_0^2\log U}\right)
 \geq
 1-\frac8{L^2}-\frac{\log u+O(1)}u
 \longrightarrow1.
\]
Increase \(Y_0\) so that \(H_{a_0}(U)\geq4U^{9/10}\).
Removing primes at most \(\sqrt U\) removes at most \(\sqrt U\) integers.
Hence
\[
 |P(m_0,a_0,U)|
 \geq2U^{9/10}-U^{1/2}
 \geq U^{9/10}.
\]
\end{proof}

\subsection*{Proof of \ref{B2.26}}\label{prf:B2.26}\phantomsection
\begin{proof}
Put \(L=\log Y\), \(u=\log U\), \(D=m_0Y^A\), and \(\ell=\log D\).
Choose the optional exceptional pair of Theorem~\ref{B2.23} for every modulus \(m_0s\).
For an allowed \(s\), denote by
\(\lambda_s(v)=(v/s)\) the quadratic character on
\((\Z/s\Z)^\times\).

Call \(s\) exceptional if the selected pair
\((\chi_s,\beta_s)\) modulo \(m_0s\) satisfies
\[
 1-\beta_s<\frac{c_{\mathrm P}}{\ell},
 \qquad
 \chi_s(v)=\lambda_s(v)
 \quad
 \bigl(v\equiv1\pmod{m_0},\ (v,s)=1\bigr).
 \tag{4}\label{eq:aux-aligned-exception}
\]
There is at most one such \(s\).
Indeed, Theorem~\ref{B2.24} makes two corresponding exceptional characters agree on integers coprime to the product of their moduli.
Their restrictions in \eqref{eq:aux-aligned-exception} would therefore make
\(\lambda_{s_1}\) and \(\lambda_{s_2}\) agree on integers congruent to \(1\) modulo \(m_0\) and coprime to \(s_1s_2\).
If \(s_1\ne s_2\), choose a prime dividing exactly one of them.
At that prime choose a quadratic nonresidue, and at every other prime dividing \(s_1s_2\) choose \(1\).
The Chinese remainder theorem, also imposing \(v\equiv1\pmod{m_0}\), gives an integer on which the two Jacobi symbols differ.
Take \(n_*\) to be that \(s\), or \(0\) if none exists.

Fix \(s\ne n_*\), and apply Theorem~\ref{B2.23} modulo \(m_0s\).
Let
\[
 \mathcal B_s=
 \{b\in(\Z/m_0s\Z)^\times:b\equiv a_0\pmod{m_0}\}.
\]
The Chinese remainder theorem identifies \(\mathcal B_s\) with
\((\Z/s\Z)^\times\).
Therefore
\[
 \sum_{b\in\mathcal B_s}\lambda_s(b)=0.
 \tag{5}\label{eq:aux-primary-cancel}
\]
For the exceptional contribution, character orthogonality on the kernel of reduction modulo \(m_0\) gives
\[
 \sum_{b\in\mathcal B_s}\lambda_s(b)\chi_s(b)=0
\]
unless the restriction of \(\chi_s\) to that kernel is \(\lambda_s\).
In this aligned case, \(s\ne n_*\) gives
\(\beta_s\leq1-c_{\mathrm P}/\ell\).
Since \(2\leq U^{\beta_s}\leq U\), the bounds
\(\operatorname{li}_2(U^{\beta_s})\leq2U^{\beta_s}\) and
\(\operatorname{li}_2(U)\geq U/(2u)\) give
\[
 \frac{|\operatorname{li}_2(U^{\beta_s})|}
      {\operatorname{li}_2(U)}
 \leq r,\qquad
 r=4u\exp(-c_{\mathrm P}u/\ell).
 \tag{6}\label{eq:aux-exception-ratio}
\]

Put
\(k=A_0(e^{-cu/\log(m_0s)}+e^{-\sqrt{cu}})\).
Write \(H_b\) for the corrected main term in class \(b\), and put
\(A_U=\operatorname{li}_2(U)/\varphi(m_0s)\).
These main terms are nonnegative.
In the aligned case, \(r\leq1/2\) gives \(H_b\geq(1-r)A_U\), while the exceptional term in each class has magnitude at most \(rA_U\).
Its signed sum is therefore bounded by
\(r(1-r)^{-1}\sum_bH_b\leq2r\sum_bH_b\).
The signed sum vanishes when there is no exceptional pair and, by orthogonality, in the unaligned case.
For \(|k|\leq1/2\), the relative estimate gives
\[
 \sum_{b\in\mathcal B_s}H_b
 \leq2\sum_{b\in\mathcal B_s}\pi(U;m_0s,b)
 \leq2(|P|+\sqrt U),
\]
where \(P=P(m_0,a_0,U)\).
The summed errors contribute at most \(|k|\sum_bH_b\), and the primary terms cancel by \eqref{eq:aux-primary-cancel}.
Primes dividing \(s\), and the passage from primes at most \(U\) to \(P\), are covered by the primes at most \(\sqrt U\).
Here \(s\leq Y^A<\sqrt U\) for large \(Y\).
It follows that
\[
 \left|\sum_{q\in P}\lambda_s(q)\right|
 \leq
 (4r+2|k|)(|P|+\sqrt U)+\sqrt U.
 \tag{7}\label{eq:aux-signed-relative}
\]

Since \(s\leq Y^A\), we have \(\log(m_0s)\leq\ell=\log m_0+AL\).
For large \(Y\), admissibility gives \(\log m_0\leq4u/L^2\) and \(AL\leq4u/L^2\), so \(u/\ell\geq L^2/8\).
Also \(\log u\leq(\alpha/5+2)L\) and \(u\geq Y^{\alpha/6}\).
These bounds place every \(m_0s\) in the range of the progression theorem and yield, after increasing \(Y_0\),
\[
 4r+2|k|\leq\tfrac14Y^{-B},\qquad
 U^{-2/5}\leq\tfrac14Y^{-B}.
\]
Lemma~\ref{B2.25} gives \(\sqrt U\leq U^{-2/5}|P|\).
Substitution in \eqref{eq:aux-signed-relative} proves
\eqref{eq:aux-kernel-cancel}.
\end{proof}

\subsection*{Proof of \ref{B2.27}}\label{prf:B2.27}\phantomsection
\begin{proof}
For a positive integer \(v\), let
\[
 r(v)=\prod_{\substack{p\mid v\\
                        v_p(v)\text{ odd}\\p\nmid m_0}}p.
\]
The product contains each prime at most once and excludes every divisor of \(m_0\), including \(2\).
Quadratic reciprocity gives, for \(1\leq n,n'\leq N\),
\[
 \left|
 \sum_{q\in P}\left(\frac nq\right)\left(\frac{n'}q\right)
 \right|
 =
 \left|\sum_{q\in P}\left(\frac q{r(nn')}\right)\right|.
 \tag{8}\label{eq:aux-reciprocity-pair}
\]
The reciprocity sign and the factors contributed by primes dividing \(m_0\) are constant on the progression.
The factor at \(2\) is constant because \(8\mid m_0\).
The condition \(q>N^2\) prevents a candidate prime from dividing \(nn'\).

Expand the square and use \(|\kappa(n)|\leq1\).
The resulting upper bound is
\[
 \sum_{n,n'\leq N}c_y(n)c_y(n')
 \left|
 \sum_{q\in P}\left(\frac nq\right)\left(\frac{n'}q\right)
 \right|.
 \tag{9}\label{eq:aux-moment-expansion}
\]
The coefficients vanish outside prime powers.
If \(n=p^a\), its kernel after removing primes dividing \(m_0\) is \(p\) when \(a\) is odd and \(p\nmid m_0\), and is \(1\) otherwise.

Pairs with \(r(nn')=1\) either have the same prime base or have both entries in the set counted by \(R_0\).
Their total coefficient is at most
\[
 \sum_p s_p(y,N)^2+R_0(y,m_0,N)^2.
\]
For a fixed kernel \(r>1\), two distinct nontrivial kernels \(p,p'\) can produce \(r\) only when \(r=pp'\).
These pairs have total weight at most \(2B^2\).
For prime \(r\), the two kernel orders are \((r,1)\) and \((1,r)\).
Together they contribute at most \(2s_r(y,N)R_0(y,m_0,N)\leq2BR_0(y,m_0,N)\).
Thus the exceptional nontrivial kernel has total coefficient at most
\(2B^2+2BR_0\).
If \(n_*\leq1\), there is no exceptional nontrivial kernel.

The total coefficient of all remaining pairs is at most
\((\sum_{n\leq N}c_y(n))^2\).
Their character sums have magnitude at most \(\eta|P|\), so their contribution to \eqref{eq:aux-moment-expansion} is at most
\(\eta|P|(\sum_{n\leq N}c_y(n))^2\).
Combining these bounds with \eqref{eq:aux-moment-expansion} proves the assertion.
\end{proof}

\subsection*{Proof of \ref{B2.28}}\label{prf:B2.28}\phantomsection
\begin{proof}
For a prime \(p\),
\[
 s_p(y,N)\leq
 \log p\sum_{a\geq1}\min\{1,y/p^a\}.
\]
If \(p\leq y\), split the sum at \(p^a\leq y\).
The initial part contributes at most \(\log y\), and the geometric tail contributes at most \(2\log p\).
If \(p>y\), the whole sum is at most
\(2y\log p/p\leq2\log y\), since \((\log t)/t\) decreases for \(t\geq3\).
This proves the first bound.

To sum the squares, first consider \(p\leq\sqrt y\).
Their contribution is \(O(\sqrt y(\log y)^2)=O(y\log y)\).
For \(\sqrt y<p\leq y\), the bound \(s_p\leq3\log y\) and
\(\log y\leq2\log p\) give \(s_p^2\ll(\log y)\log p\).
The elementary Chebyshev bound
\(\sum_{p\leq t}\log p\leq t\log4\) handles this range.
For \(p>y\), we have \(s_p^2\leq4y^2(\log p)^2/p^2\).
Partial summation of the same Chebyshev bound gives
\[
 \sum_{p>y}\frac{(\log p)^2}{p^2}\ll\frac{\log y}{y}.
\]
This proves the second bound in \eqref{eq:aux-prime-power-bounds}.

For the even powers, write \(p^{2a}=(p^a)^2\).
Their contribution is bounded by
\[
 \sum_{v\geq1}\Lambda(v)\min\{1,y/v^2\}\ll\sqrt y.
\]
The last estimate follows by splitting at \(v=\sqrt y\) and applying
\(\sum_{v\leq t}\Lambda(v)\ll t\).
There are at most \(w\) prime bases \(p\mid m_0\) with \(p\leq N\).
Each contributes at most \(3\log y\).
This proves \eqref{eq:aux-R0-bound}.

Finally, \(c_y(n)\leq y\Lambda(n)/n\), and the elementary Mertens upper bound
\(\sum_{n\leq Z}\Lambda(n)/n\ll\log Z\) proves
\eqref{eq:aux-total-coefficient}.
\end{proof}

\subsection*{Proof of \ref{B2.29}}\label{prf:B2.29}\phantomsection
\begin{proof}
Apply Lemma~\ref{B2.26} with its parameters \(A=800K\) and \(B=K+1\).
It supplies the cancellation hypothesis of Lemma~\ref{B2.27} with
\(\eta=Y^{-(K+1)}\), since \(\lfloor Z\rfloor_+^2\leq Y^{800K}\).
For large \(Y\), admissibility gives
\(\log U\geq Y^{\alpha/6}\geq1600K\log Y\).
Therefore \(\sqrt U\geq Z^2\), and every candidate \(q\) exceeds
\(\lfloor Z\rfloor_+^2\).

Put \(w=Y^{\alpha/5}\).
By \eqref{eq:aux-modulus-gap}, every prime divisor \(p\mid m_0\) with
\(p\leq\lfloor Z\rfloor_+\) is at most \(w\).
Since \(\alpha\leq1\) and \(y\geq Y\), we have \(w\leq y^{1/5}\).
For large \(Y\), the estimates of Lemma~\ref{B2.28} consequently give
\(R_0\leq5C_1\sqrt y\) and \(B=C_2\log y\).
Thus \(R_0^2+2B^2+2BR_0\ll y\log y\) with an absolute implied constant.

For the remaining term, use
\(\log Z=400K\log Y\) and \((400K\log Y)^2\leq Y\).
Then
\[
 \begin{split}
 \eta\left(\sum_{n\leq Z}c_y(n)\right)^2
 &\leq
 C_4^2Y^{-(K+1)}y^2(400K\log Y)^2\\
 &\leq C_4^2y^2Y^{-K}
 \leq C_4^2y
 \leq C_4^2y\log y.
 \end{split}
\]
Insert these estimates and
\(\sum_p s_p^2\leq C_3y\log y\) into Lemma~\ref{B2.27}.
Take \(C_{\mathrm m}\) to be the sum of the absolute constants in the resulting bounds.
\end{proof}

\subsection*{Proof of \ref{B2.30}}\label{prf:B2.30}\phantomsection
\begin{proof}
For fixed \(Z,q,\kappa\), the function
\(S_q(y)=S(Z,y,q,\kappa)\) is continuous and piecewise linear.
Its right derivative is
\[
 D_q(v)=
 \sum_{\substack{1\leq n\leq Z\\n>v}}
 \frac{\Lambda(n)\kappa(n)}n\left(\frac nq\right).
\]
For \(v>0\),
\[
 vD_q(v)=
 S\bigl(Z,v,q,\kappa\,1_{\{n>v\}}\bigr).
 \tag{14}\label{eq:aux-derivative-twist}
\]
The coefficient function on the right has modulus at most \(1\).
The full quantifier over \(\kappa\) in the mean-square hypothesis therefore gives
\[
 \sum_{q\in P}|D_q(v)|^2
 \leq C\frac{\log v}{v}|P|
 \qquad(Y\leq v\leq X).
 \tag{15}\label{eq:aux-derivative-moment}
\]

Consider \(T\leq W\leq2T\) with \([T,W]\subseteq[Y,X]\).
Let \(B_T\) consist of candidates with a witness in \([T,W]\).
For every \(q\in P\), choose \(y_q\in[T,W]\) and take it to be a witness when \(q\in B_T\).
The fundamental theorem of calculus and Cauchy--Schwarz give
\[
 |S_q(y_q)|^2
 \leq2|S_q(T)|^2+
 2T\int_T^W|D_q(v)|^2\,dv.
\]
After summing and using \eqref{eq:aux-derivative-moment},
\[
 \sum_{q\in P}|S_q(y_q)|^2
 \leq4CT\log(2T)|P|.
 \tag{16}\label{eq:aux-block-maximal}
\]
For \(q\in B_T\), the chosen witness satisfies \(|S_q(y_q)|^2>T^{1+2\tau}/4\).
Hence
\[
 |B_T|
 \leq16C\log(2T)T^{-2\tau}|P|
 \leq\frac{32C}{\tau}T^{-\tau}|P|,
 \tag{17}\label{eq:aux-block-bad}
\]
where \(T\geq2\) and \(\log(2T)\leq2T^\tau/\tau\).

If \(X<Y\), the witness hypothesis gives \(B=\varnothing\).
Otherwise, for each integer \(i\geq0\) with \(T_i=2^iY\leq X\), put
\(W_i=\min\{2T_i,X\}\).
Let \(B_i\) consist of the \(q\in B\) with a witness in \([T_i,W_i]\).
These intervals cover \([Y,X]\), so \(|B|\leq\sum_i|B_i|\).
Summing \eqref{eq:aux-block-bad} gives
\[
 |B|
 \leq
 \frac{32C}{\tau}
 \frac{Y^{-\tau}}{1-2^{-\tau}}|P|.
\]
Choose \(Y\) so that the coefficient is at most \(\eta\).
\end{proof}

\subsection*{Proof of \ref{B2.31}}\label{prf:B2.31}\phantomsection
\begin{proof}
Set \(s_0=399/400\). Let \(Q_0\geq20\) be an absolute threshold chosen below; for \(3\leq Q<Q_0\), there are finitely many characters of conductor at most \(Q_0\), which are absorbed into \(C_{\mathrm z}\). Assume \(Q\geq Q_0\). Put \(X=\lceil Q\rceil\), \(Y=X^4\), \(J=8(\lfloor\log_2 X\rfloor+1)\), \(N=2^J\geq X^8\), and \(K=YN=Y2^J\). Remove the at most one character of conductor \(1\) from \(\mathcal E\), leaving a subset \(\mathcal E'\) of nonprincipal primitive characters. For each \(\chi\in\mathcal E'\) of conductor \(q_\chi\leq Q\), fix a zero \(\rho_\chi\) of \(L(s,\chi)\) with \(s_0<\Re\rho_\chi<1\) and \(|\Im\rho_\chi|\leq Q\).

For each \(\chi\in\mathcal E'\), define the mollified Dirichlet polynomial coefficients \(c_n=\sum_{d\mid n,\,d\leq Y,\,n/d\leq N}\mu(d)\). Expanding the product of the partial sum of \(L(\rho_\chi,\chi)\) up to \(N\) and the M\"obius mollifier up to \(Y\) gives the mollifier identity
\[
 \left(\sum_{n=1}^N\chi(n)n^{-\rho_\chi}\right)
 \left(\sum_{d=1}^Y\mu(d)\chi(d)d^{-\rho_\chi}\right)
 =1+\sum_{Y<n\leq YN}c_n\chi(n)n^{-\rho_\chi}.
\]
Because \(L(\rho_\chi,\chi)=0\) and partial sums of \(\chi\) are bounded by \(q_\chi\), partial summation gives
\[
 \left|\sum_{n=1}^N\chi(n)n^{-\rho_\chi}\right|
 \leq 2(q_\chi+1)\left(1+\frac{|\rho_\chi|}{\Re\rho_\chi}\right)N^{-\Re\rho_\chi}
 \leq 2(2X)(3X)(X^8)^{-15/16}
 \leq \frac{3}{XY}\leq\frac{1}{2Y},
\]
using \(X\geq20\) and \(\Re\rho_\chi>15/16\). Since the mollifier sum has modulus at most \(Y\), their product has modulus at most \(1/2\), and therefore \(\bigl|\sum_{Y<n\leq YN}c_n\chi(n)n^{-\rho_\chi}\bigr|\geq1/2\). Splitting \((Y,YN]\) into the \(J\) dyadic blocks \((U_j,2U_j]\) with \(U_j=Y2^j\) (\(0\leq j<J\)), each \(\chi\in\mathcal E'\) selects an index \(j\) with
\[
 \left|\sum_{U_j<n\leq 2U_j}c_n\chi(n)n^{-\rho_\chi}\right|\geq\frac{1}{2J}.
\]

Fix \(0\leq j<J\), put \(U=U_j\), and let \(T\subseteq\mathcal E'\) be the set of characters assigned to block \(j\). Write \(b_n=c_n n^{-s_0}\) and \(v_\chi(n)=\chi(n)n^{-(\rho_\chi-s_0)}\). The divisor bound \(|c_n|\leq d(n)\leq C_1 n^{1/400}\) gives
\[
 \sum_{U<n\leq 2U}|b_n|^2\leq C_1^2 U^{-99/100}\leq C_1^2 U^{-197/200}.
\]
For distinct \(\chi,\chi'\in T\), the product \(\chi\overline{\chi'}\) is a nonprincipal character modulo \(q_\chi q_{\chi'}\leq Q^2\), so its partial sums are bounded by \(q_\chi q_{\chi'}\leq Q^2\). Putting \(w=(\rho_\chi-s_0)+\overline{\rho_{\chi'}-s_0}\) (which has \(\Re w>0\) and \(|w|\leq 1+2Q\)), summation by parts yields
\[
 \left|\sum_{U<n\leq 2U}v_\chi(n)\overline{v_{\chi'}(n)}\right|
 \leq 2q_\chi q_{\chi'}(1+|w|)\leq 8Q^3,
\]
while the diagonal sum \(\chi=\chi'\) is at most \(U\). By the Hal\'asz--Montgomery duality inequality,
\[
 \frac{|T|}{4J^2}
 \leq\sum_{\chi\in T}\left|\sum_{U<n\leq 2U}b_n v_\chi(n)\right|^2
 \leq\left(\sum_{U<n\leq 2U}|b_n|^2\right)\bigl(U+8Q^3|T|\bigr)
 \leq C_1^2 U^{3/200}+8Q^3 C_1^2 U^{-197/200}|T|.
\]
Since \(U\geq Y\geq Q^4\), we have \(U^{-197/200}\leq Q^{-197/50}\), so \(8Q^3 C_1^2 U^{-197/200}\leq 8C_1^2 Q^{-47/50}\leq 1/(8J^2)\) once \(Q_0\) is chosen sufficiently large (as \(J\leq 32\log Q\)). Subtracting \(|T|/(8J^2)\) gives \(|T|\leq 8J^2 C_1^2 U^{3/200}\leq 8J^2 C_1^2 K^{3/200}\). Summing over \(0\leq j<J\) and restoring the conductor-\(1\) character gives
\[
 |\mathcal E|\leq 1+8J^3 C_1^2 K^{3/200}.
\]
Because \(J\leq 32\log Q\) and \(K=Y2^J\ll Q^{12}\), we have \(J^3 K^{3/200}\ll(\log Q)^3 Q^{9/50}\ll Q^{1/4}\), yielding \(|\mathcal E|\leq C_{\mathrm z}Q^{1/4}\).
\leanref{Formalism.\allowbreak{}HermitianLattices.\allowbreak{}Ext.\allowbreak{}PrimitiveCharacterCasts.\allowbreak{}ZD.\allowbreak{}mollifier\_\allowbreak{}identity}{Formalism.\allowbreak{}GL3MeanSquares.\allowbreak{}GenusMeanAsymptotics.\allowbreak{}JacobiCharacterSums.\allowbreak{}CountZMoll}{79--132}
\leanref{Formalism.\allowbreak{}HermitianLattices.\allowbreak{}Ext.\allowbreak{}PrimitiveCharacterCasts.\allowbreak{}ZD.\allowbreak{}zero\_\allowbreak{}detection}{Formalism.\allowbreak{}GL3MeanSquares.\allowbreak{}GenusMeanAsymptotics.\allowbreak{}JacobiCharacterSums.\allowbreak{}CountZDetect}{43--163}
\leanref{Formalism.\allowbreak{}HermitianLattices.\allowbreak{}Ext.\allowbreak{}PrimitiveCharacterCasts.\allowbreak{}ZD.\allowbreak{}coefficient\_\allowbreak{}square\_\allowbreak{}bound}{Formalism.\allowbreak{}GL3MeanSquares.\allowbreak{}GenusMeanAsymptotics.\allowbreak{}JacobiCharacterSums.\allowbreak{}CountZMoll}{133--192}
\leanref{Formalism.\allowbreak{}HermitianLattices.\allowbreak{}Ext.\allowbreak{}PrimitiveCharacterCasts.\allowbreak{}ZD.\allowbreak{}pair\_\allowbreak{}sum\_\allowbreak{}bound}{Formalism.\allowbreak{}GL3MeanSquares.\allowbreak{}GenusMeanAsymptotics.\allowbreak{}JacobiCharacterSums.\allowbreak{}CountZPair}{98--102}
\leanref{Formalism.\allowbreak{}HermitianLattices.\allowbreak{}Ext.\allowbreak{}PrimitiveCharacterCasts.\allowbreak{}ZD.\allowbreak{}pair\_\allowbreak{}twist\_\allowbreak{}bound}{Formalism.\allowbreak{}GL3MeanSquares.\allowbreak{}GenusMeanAsymptotics.\allowbreak{}JacobiCharacterSums.\allowbreak{}CountZTwist}{32--74}
\leanref{Formalism.\allowbreak{}HermitianLattices.\allowbreak{}Ext.\allowbreak{}PrimitiveCharacterCasts.\allowbreak{}ZD.\allowbreak{}duality\_\allowbreak{}inequality}{Formalism.\allowbreak{}GL3MeanSquares.\allowbreak{}GenusMeanAsymptotics.\allowbreak{}JacobiCharacterSums.\allowbreak{}CountZDual}{9--125}
\leanref{Formalism.\allowbreak{}HermitianLattices.\allowbreak{}Ext.\allowbreak{}PrimitiveCharacterCasts.\allowbreak{}ZD.\allowbreak{}blockCountOf\_\allowbreak{}statement}{Formalism.\allowbreak{}GL3MeanSquares.\allowbreak{}GenusMeanAsymptotics.\allowbreak{}JacobiCharacterSums.\allowbreak{}CountZBlock}{23--167}
\leanref{Formalism.\allowbreak{}HermitianLattices.\allowbreak{}Ext.\allowbreak{}PrimitiveCharacterCasts.\allowbreak{}ZD.\allowbreak{}numeric\_\allowbreak{}constants}{Formalism.\allowbreak{}GL3MeanSquares.\allowbreak{}GenusMeanAsymptotics.\allowbreak{}JacobiCharacterSums.\allowbreak{}CountZNum}{10--135}
\leanref{Formalism.\allowbreak{}HermitianLattices.\allowbreak{}Ext.\allowbreak{}PrimitiveCharacterCasts.\allowbreak{}ZD.\allowbreak{}final\_\allowbreak{}bound\_\allowbreak{}of\_\allowbreak{}assembly}{Formalism.\allowbreak{}GL3MeanSquares.\allowbreak{}GenusMeanAsymptotics.\allowbreak{}JacobiCharacterSums.\allowbreak{}CountZFinal}{10--130}
\end{proof}

\subsection*{Proof of \ref{B2.33}}\label{prf:B2.33}\phantomsection
\begin{proof}
Under the Chinese remainder isomorphism
\[
 (\Z/fq\Z)^\times
 \simeq(\Z/f\Z)^\times\times(\Z/q\Z)^\times,
\]
take the product of \(\xi\) with the nontrivial quadratic character modulo \(q\).
A character of coprime product level has conductor equal to the product of the conductors of its two factors.
Both factors here are primitive, so the conductor is \(fq\).
The quadratic factor also makes the product nonprincipal.
Extending the characters by zero gives \eqref{eq:aux-primitive-twist} on all integers.

For uniqueness, induce both characters to level \(NN'L\).
They agree on every unit of that ring by the assumed equality.
Their induced characters at level \(NN'\) therefore agree as well.
A character induced from both levels \(N\) and \(N'\) factors through
\(\gcd(N,N')\).
Primitivity gives \(N\leq\gcd(N,N')\) and \(N'\leq\gcd(N,N')\).
Hence \(N=N'\), and injectivity of induction gives \(\chi=\chi'\).
\end{proof}

\subsection*{Proof of \ref{B2.34}}\label{prf:B2.34}\phantomsection
\begin{proof}
Write
\[
 x=Y^K,\quad V=x^{400},\quad
 \theta=x^{-2},\quad h=V\theta,\quad
 \delta=\frac{\alpha}{1600K}.
\]
In particular \(0<\delta\leq1/800\).
Introduce
\[
 \mathcal S_\chi(t)
 =
 \sum_{1\leq n\leq t}
 \Lambda(n)\chi(n)\left(\frac tn-1\right).
\]

For every \(0<\delta\leq1/800\), there is \(A_\delta\geq0\) such that, for \(Q\geq2\), every \(\chi\in\mathcal Z(Q)\) obeys
\[
 \begin{split}
 \bigl|\mathcal S_\chi(V+V\theta)
       -\mathcal S_\chi(V)-V\theta D_\chi\bigr|
 \leq A_\delta\log Q\bigl(
    (\log Q)\theta V^{1-1/400+\delta}
    +V^2/Q\bigr)
 \end{split}
 \tag{19}\label{eq:aux-contour-difference}
\]
for \(V\geq1\) and \(0<\theta\leq1\).

Every \(\chi\in\mathcal Z(Q)\) is a nonprincipal Dirichlet character modulo \(N\) with \(1\leq N\leq Q\) and \(L(w,\chi)\ne0\) on the rectangle \(\Re w>399/400\), \(|\Im w|\leq Q\); here \(\chi\) need not be primitive. On \(\Re s=2\), the integrated Perron kernel for \(\mathcal S_\chi(t)\) is
\[
 \frac{-L'(s,\chi)}{L(s,\chi)}
       \frac{t^s}{s(s-1)}.
\]
Move the segment \(|\Im s|\leq Q/2\) to \(\sigma_0=1-1/400+\delta\). Because \(\chi\) is nonprincipal, \(L(s,\chi)\) is entire, so the only pole crossed is the simple pole of \((s-1)^{-1}\) at \(s=1\) with residue \(tD_\chi\), where \(D_\chi=-L'(1,\chi)/L(1,\chi)\).

Suppose \(\sigma_0\leq\Re s<11/10\) and \(|\Im s|\leq Q-1\). To bound \(L'/L(s,\chi)\) without assuming \(\chi\) primitive, note first that orthogonality modulo \(N\) gives \(\bigl|\sum_{1\leq n\leq u}\chi(n)\bigr|\leq N\) for every \(u\geq1\), so partial summation of \(\sum_{n\geq1}\chi(n)n^{-z}\) on \(\Re z\geq1/2\) yields the convexity bound \(|L(z,\chi)|\leq 9N(1+|z|)\) directly for \(L(z,\chi)\) \cite[Ch.~12]{davenport-2000}. Second, for \(\Re z>1\), complete multiplicativity of \(\chi\) gives \(L(z,\chi)^{-1}=\sum_{n=1}^\infty\mu(n)\chi(n)n^{-z}\) and \(|L(z,\chi)|^{-1}\leq\Re z/(\Re z-1)\). Put \(z_0=21/20+i\Im s\). Then \(|L(z_0,\chi)|\geq1/21\), and every \(z\) in the closed disk \(\overline D(z_0,11/20)\subset\{\Re z\geq1/2\}\) satisfies \(1+|z|\leq2(|\Im s|+2)\), whence \(|L(z,\chi)/L(z_0,\chi)|\leq 378\,N(|\Im s|+2)\). Since \(\sigma_0>399/400>26/27\) and \(L(s,\chi)\ne0\), Landau's Borel--Carath\'eodory local zero-expansion lemma on the concentric disks \(\overline D(z_0,11/60)\subset\overline D(z_0,11/20)\) \cite[\S12]{davenport-2000}, \cite[Lemma~6.3 and Theorem~11.4]{montgomery-vaughan-2007} expresses \(L'/L(s,\chi)\), up to an error \(O\bigl(\log(N(|\Im s|+2))\bigr)=O(\log Q)\), as \(\sum_{\rho\in Z}m(\rho)/(s-\rho)\), where \(Z\) is the set of zeros of \(L(z,\chi)\) in \(\overline D(z_0,11/60)\), with multiplicities \(m(\rho)\) satisfying \(\sum_{\rho\in Z}m(\rho)\ll\log(N(|\Im s|+2))\ll\log Q\) (as \(N(|\Im s|+2)\leq Q(Q+1)\leq Q^3\) for \(Q\geq2\)). (Equivalently, if \(\chi\) is induced by a primitive character \(\chi^*\bmod q^*\) with \(q^*\mid N\), then \(L(z,\chi)=L(z,\chi^*)\prod_{p\mid N,\,p\nmid q^*}(1-\chi^*(p)p^{-z})\); the extra Euler factors have all their zeros on \(\Re z=0\), disjoint from \(\overline D(z_0,11/20)\), and contribute \(\sum_{p\mid N,\,p\nmid q^*}\chi^*(p)(\log p)/(p^s-\chi^*(p))=O(\log(2N))=O(\log Q)\) to \(L'/L(s,\chi)\) on \(\Re s\geq1/2\) \cite[Ch.~9]{davenport-2000}, \cite[Lemma~6.1]{montgomery-vaughan-2007}.) Every \(\rho\in Z\) satisfies \(|\Im\rho|\leq Q-1+11/60<Q\), so \(\chi\in\mathcal Z(Q)\) gives \(\Re\rho\leq399/400\) and \(|s-\rho|\geq\Re(s-\rho)\geq\delta\). For \(\Re s\geq11/10\), the absolutely convergent Dirichlet series \(-\sum_{n\geq1}\Lambda(n)\chi(n)n^{-s}\) gives \(|L'/L(s,\chi)|\leq 1/(\Re s-1)+O(1)=O(1)=O(\log Q)\). Thus
\[
 \left|\frac{L'}L(s,\chi)\right|
 \ll_\delta\log Q
 \quad
 \bigl(\Re s\geq\sigma_0,\ |\Im s|\leq Q-1\bigr).
\]
For the difference of the two left-edge integrals, use
\[
 \bigl|(V+V\theta)^s-V^s\bigr|
 \ll \theta|s|V^{\sigma_0}
 \qquad(\Re s=\sigma_0).
\]
After \(|s|\) cancels, use
\(\int_{-Q/2}^{Q/2}|\sigma_0-1+it|^{-1}\,dt\ll\log Q\), since
\(1-\sigma_0\geq1/800\).
The horizontal segments at \(|\Im s|=Q/2\leq Q-1\), and the tails of the original line at real part \(2\), contribute \(O_\delta((\log Q)V^2/Q)\).
These estimates give \eqref{eq:aux-contour-difference}.
\leanref{Formalism.\allowbreak{}HermitianLattices.\allowbreak{}Ext.\allowbreak{}BoxTail.\allowbreak{}tail\_\allowbreak{}box\_\allowbreak{}negLD}{Formalism.\allowbreak{}GL3MeanSquares.\allowbreak{}GenusMeanAsymptotics.\allowbreak{}JacobiCharacterSums.\allowbreak{}TailBox}{35--120}
\leanref{Formalism.\allowbreak{}HermitianLattices.\allowbreak{}Ext.\allowbreak{}BoxTail.\allowbreak{}tail\_\allowbreak{}contour}{Formalism.\allowbreak{}GL3MeanSquares.\allowbreak{}GenusMeanAsymptotics.\allowbreak{}JacobiCharacterSums.\allowbreak{}TailContour}{14--136}

There is also a direct finite-sum estimate
\[
 \left|
 \mathcal S_\chi(V+h)-\mathcal S_\chi(V)
 -h\sum_{n\leq V}\frac{\Lambda(n)\chi(n)}n
 \right|
 \leq
 h(h+1)\frac{\log(V+h)}V.
 \tag{20}\label{eq:aux-finite-difference}
\]
Indeed, terms with \(n\leq V\) cancel exactly.
There are at most \(h+1\) integers in \(V<n\leq V+h\).
Each corresponding term satisfies
\[
 \left|\Lambda(n)\chi(n)
       \left(\frac{V+h}{n}-1\right)\right|
 \leq h\,\frac{\log(V+h)}V.
\]
Put \(\Delta=\mathcal S_\chi(V+h)-\mathcal S_\chi(V)\) and
\(A_V=\sum_{n\leq V}\Lambda(n)\chi(n)/n\).
Since \(hT_\chi(V)=(\Delta-hA_V)-(\Delta-hD_\chi)\), the triangle inequality and
\eqref{eq:aux-contour-difference}--\eqref{eq:aux-finite-difference}, followed by division by \(h\), give
\[
 |T_\chi(V)|
 \leq
 A_\delta(\log Q)^2V^{-1/400+\delta}
 +A_\delta(\log Q)\frac{V}{Q\theta}
 +(h+1)\frac{\log(V+h)}V.
 \tag{21}\label{eq:aux-tail-general}
\]

Multiply by \(x\).
Since \(V=x^{400}\), the first term becomes
\[
 A_\delta(\log Q)^2x^{400\delta}
 \leq A_\delta Y^{\alpha/2}Y^{\alpha/4}
 =A_\delta Y^{3\alpha/4}.
\]
The second becomes \(A_\delta(\log Q)x^{403}/Q=o(1)\), uniformly because
\(Q\geq\exp(Y^{\alpha/6})\).
The third is \(O(x^{-1}\log x+x^{-399}\log x)\).
Their sum is at most \(Y^\alpha\) for sufficiently large \(Y\).
\end{proof}

\subsection*{Proof of \ref{B2.35}}\label{prf:B2.35}\phantomsection
\begin{proof}
For the first assertion, apply Lemma~\ref{B2.34} with \(K=1\).
It gives \(|T_\chi(Y^{400})|\leq1\), since \(\alpha\leq1\).
Consequently
\[
 \left|\frac{L'}L(1,\chi)\right|
 \leq1+\sum_{n\leq Y^{400}}\frac{\Lambda(n)}n
 \leq C_{\mathrm L}\log Y
\]
with an absolute constant.

For the explicit coefficient, there are absolute constants
\(C_0,\ell_0\), with \(\ell_0\geq1\), such that
\[
 \left|\frac{L'}L(1,\chi)\right|
 \leq1000\log\log Q+C_0
 \qquad(\log Q\geq\ell_0,\ \chi\in\mathcal Z(Q)).
 \tag{24}\label{eq:aux-kappa-loglog}
\]
Put \(\ell=\log Q\) and choose
\[
 x=\ell^{5/2},\qquad V=x^{400}=\ell^{1000},
 \qquad\theta=x^{-2},\qquad\delta=1/4000.
\]
Use \eqref{eq:aux-tail-general} with this fixed \(\delta\).
Its first two terms are bounded respectively by constant multiples of
\[
 \ell^2V^{-1/400+1/4000}=\ell^{-1/4},
 \qquad
 \ell\,\frac{Vx^2}{Q}=\frac{\ell^{1006}}{e^\ell}.
\]
The finite-difference term is bounded as well.
Thus \(|T_\chi(V)|\leq C_1\), for an absolute \(C_1\) and all sufficiently large \(\ell\).

Mertens's upper bound gives
\[
 \sum_{n\leq V}\frac{\Lambda(n)}n\leq\log V+2
 \qquad(V\geq1).
\]
It follows that
\(|D_\chi|\leq C_1+\log V+2=1000\log\ell+C_1+2\), which proves
\eqref{eq:aux-kappa-loglog}.
\leanref{Deduction.\allowbreak{}mertens\_\allowbreak{}upper}{Development.\allowbreak{}AutomorphicArithmetic.\allowbreak{}DedekindLFunctions.\allowbreak{}LandauTheorems.\allowbreak{}mertensUpper}{103--113}

Finally,
\[
 \log\log Q
 \leq
 \log A_1+\frac{\varepsilon}{5}\log Y+2\log\log Y.
\]
The coefficient of \(\log Y\) obtained from
\eqref{eq:aux-kappa-loglog} is \(200\varepsilon<50/13\).
Since \(4-50/13=2/13>0\), the terms
\(2000\log\log Y+1000\log A_1+C_0\) are absorbed by the remaining multiple of \(\log Y\).
This gives \eqref{eq:aux-kappa-four}.
\end{proof}

\subsection*{Proof of \ref{B2.36}}\label{prf:B2.36}\phantomsection
\begin{proof}
Put \(Z=(Y^K)^{400}\).
For \(0<y\leq Z\), splitting at \(n=y\) gives
\[
 g_\chi(y)
 =
 -\sum_{n\leq Z}\Lambda(n)\chi(n)\min\{1,y/n\}
 -yT_\chi(Z).
 \tag{25}\label{eq:aux-finite-splitting}
\]
Indeed, for \(n\leq y\),
\(y/n-1=y/n-\min\{1,y/n\}\), whereas for \(y<n\leq Z\), the two terms on the right cancel.

The first sum in \eqref{eq:aux-finite-splitting} is \(S(Z,y,q,\kappa)\).
Lemma~\ref{B2.34} gives
\(|yT_\chi(Z)|\leq|Y^KT_\chi(Z)|\leq Y^\alpha\).
For \(Y\geq16\), \(\alpha\leq1/4\), and \(y\geq Y\), we have
\(Y^\alpha\leq\tfrac12y^{1/2+\tau}\).
The triangle inequality proves the desired bound for \(g_\chi(y)\) at every point of the window.
The absolute constant in \eqref{eq:aux-kappa-C} gives the logarithmic derivative bound.
\end{proof}

\subsection*{Proof of \ref{B2.37}}\label{prf:B2.37}\phantomsection
\begin{proof}
Set \(Q=FU\) and \(Z=(Y^K)^{400}\).
The inequalities \(1\leq F\leq\sqrt U\) give
\[
 Y^{\alpha/6}\leq\log Q
 \leq\tfrac32\log U
 \leq\tfrac32Y^{\alpha/5}(\log Y)^2
 \leq Y^{\alpha/4}
\]
for sufficiently large \(Y\).
Lemma~\ref{B2.25} gives \(|P|\geq U^{9/10}>0\).

Define \(B_1\subseteq B\) to consist of primes admitting a primitive character of conductor \(fq\), with \(f\leq F\) and \(fq\leq Q\), whose \(L\)-function has a zero in \(\mathcal R(Q)\).
By Lemma~\ref{B2.32},
\(|B_1|\leq C_{\mathrm b}Q^{1/4}\leq C_{\mathrm b}U^{3/8}\).
Since \(3/8<9/10\), increasing \(Y_0\) gives
\[
 |B_1|\leq\tfrac12\eta|P|.
 \tag{26}\label{eq:aux-stage-zero-bad}
\]

For each \(\xi\in\Xi\), let \(B_2(\xi)\subseteq B\) consist of primes for which
\(|S(Z,y,q,\xi)|>\tfrac12y^{1/2+\tau}\) at some \(y\in[Y,Y^K]\).
Theorem~\ref{B2.29} supplies its mean-square hypothesis.
Apply Theorem~\ref{B2.30} with proportion \(\eta/(2(R+1))\).
It gives
\[
 |B_2(\xi)|\leq
 \frac{\eta}{2(R+1)}|P|,
 \qquad
 \left|\bigcup_{\xi\in\Xi}B_2(\xi)\right|
 \leq\tfrac12\eta|P|.
 \tag{27}\label{eq:aux-stage-sum-bad}
\]

We prove
\[
 B\subseteq B_1\cup\bigcup_{\xi\in\Xi}B_2(\xi).
 \tag{28}\label{eq:aux-stage-cover}
\]
Take \(q\in B\) outside the right-hand side, and let \(\xi,L\) witness its defining failure.
With \(f=\operatorname{cond}(\xi)\), we have \(q>\sqrt U\geq F\geq f\), so \(q\nmid f\).
Also \(q\equiv a_0\pmod{m_0}\), where \(8\mid m_0\) and \((a_0,m_0)=1\), so \(q\) is odd.
Lemma~\ref{B2.33} constructs a nonprincipal primitive
\(\chi\) of conductor \(fq\) with
\(\chi(n)=\xi(n)(n/q)\) for every \(n\).
Its conductor is at most \(FU\).
Since \(q\notin B_1\), \(L(s,\chi)\) has no zero in \(\mathcal R(Q)\). Thus \(\chi\in\mathcal Z(Q)\).
Since \(q\notin B_2(\xi)\), its finite sum obeys the required bound throughout the window.
Lemma~\ref{B2.36} gives its box.

Any primitive representative of the same function at level \(L\) agrees with \(\chi\) on all integers coprime to \(L\).
The uniqueness assertion of Lemma~\ref{B2.33} identifies it with \(\chi\).
Thus every primitive representative at level \(L\) belongs to the box, contrary to the choice of \(\xi,L\). This proves \eqref{eq:aux-stage-cover}.
Combine \eqref{eq:aux-stage-zero-bad} and
\eqref{eq:aux-stage-sum-bad}.
\end{proof}

\subsection*{Proof of \ref{B2.38}}\label{prf:B2.38}\phantomsection
\begin{proof}
For each character \(\xi\) modulo \(N_0\), choose its primitive representative \(\xi^*\).
Its conductor divides \(N_0\), and
\(\xi^*(n)=\xi(n)\) whenever \((n,N_0)=1\).
The set \(\Xi\) of these representatives has at most
\(R=\#\{\text{Dirichlet characters modulo }N_0\}\) members.

Use \(\alpha=\varepsilon\), \(U=e^z\), \(m_0=8\), \(a_0=3\), and \(F=N_0\) in Theorem~\ref{B2.37}.
For large \(Y\), we have
\(\tfrac14(\log Y)^2\log8\leq z\),
\(Y^{\varepsilon/6}\leq z\leq z(\log Y)^2\), and
\(N_0\leq e^{z/2}\).
Thus the progression is admitted and the conductor bound holds.
The sole prime divisor of \(m_0=8\) is \(2\leq z\), so
\eqref{eq:aux-modulus-gap} holds.

If \(q\) fails the first selection condition, some character \(\xi\) modulo \(N_0\) has a \(q\)-twist with a primitive representative at level \(N_0q\) outside \(\mathcal B(N_0e^z,C_1,\tau,Y,Y^K)\).
On integers coprime to this level the values of \(\xi\) and \(\xi^*\) agree.
The same representative therefore witnesses a failure for the member
\(\xi^*\in\Xi\) in Theorem~\ref{B2.37}.
Apply Theorem~\ref{B2.37} to this bad set.
\end{proof}

\subsection*{Proof of \ref{B2.39}}\label{prf:B2.39}\phantomsection
\begin{proof}
Write \(m_0=16\mathfrak m(z,q)\).
The integer \(\mathfrak m(z,q)\) is odd.
By the Chinese remainder theorem there is a reduced residue class \(a_0\) modulo \(m_0\) specified by
\(a_0\equiv7\pmod{16}\) and
\(a_0\equiv-1\pmod{\mathfrak m(z,q)}\).
Hence
\[
 P_2(z,z',q)=P(m_0,a_0,e^{z'}).
 \tag{30}\label{eq:aux-second-progression}
\]
Use the parameters
\[
 \alpha=\frac65\varepsilon,\qquad
 U=e^{z'},\qquad F=N_0e^z
 \tag{31}\label{eq:aux-second-parameters}
\]
in Theorem~\ref{B2.37}.

The elementary primorial bound gives
\[
 \prod_{\substack{3\leq p\leq z\\p\text{ prime}}}p
 \leq4^{\lfloor z\rfloor_+}.
\]
Since \(q\leq e^z\),
\[
 \log m_0\leq\log16+z+z\log4\leq4z
 \tag{32}\label{eq:aux-second-log-modulus}
\]
for large \(Y\).
It follows that
\(\tfrac14(\log Y)^2\log m_0\leq z(\log Y)^2\leq z'=\log U\).
For the other two admission bounds, observe that
\(Y^{\alpha/6}=Y^{\varepsilon/5}=z\) and
\(Y^{\alpha/5}=zY^{\varepsilon/25}\).
The lower bound \(z'\geq z(\log Y)^2\) gives
\(Y^{\alpha/6}\leq\log U\).
The upper bound on \(z'\), together with \(K^2\leq Y^{\varepsilon/25}\), gives
\(\log U\leq Y^{\alpha/5}(\log Y)^2\).
Thus the progression in \eqref{eq:aux-second-progression} is admitted.

Every prime factor of \(m_0\) other than \(q\) is at most \(z\) for large \(Y\).
The small factors are at most \(Y^{\alpha/5}\).
For \(q\), use \(q>e^{z/2}\) and \(z\geq800K\log Y\) to obtain
\[
 q>(Y^K)^{400}.
 \tag{33}\label{eq:aux-q-beyond-truncation}
\]
This proves the prime-divisor condition \eqref{eq:aux-modulus-gap}.

For large \(Y\), every \(q\in P_1(z)\) exceeds \(N_0\).
For each character \(\xi\) modulo \(N_0\) and each \(a\in\{0,1\}\), choose a primitive character \(\xi_{q,a}^*\) representing
\(n\mapsto\xi(n)(n/q)^a\) on integers coprime to \(N_0\).
For \(a=0\), use the primitive representative of \(\xi\).
For \(a=1\), Lemma~\ref{B2.33} gives the twist of that representative by \(q\).
In both cases the conductor is at most \(N_0q\leq F\).
The resulting family has at most
\(R=2\#\{\text{Dirichlet characters modulo }N_0\}\) members, uniformly in \(q\).

Finally, \(\log F=\log N_0+z\leq z'/2\) for large \(Y\), because
\(z'\geq z(\log Y)^2\).
Hence \(1\leq F\leq\sqrt U\).
The parameter \(\alpha=6\varepsilon/5\) is at most \(1/4\).

The box height in Theorem~\ref{B2.37} is \(FU=N_0e^{z+z'}\).
If \(d\) fails the second selection condition for \(\xi,a\), the same primitive character witnesses failure in Theorem~\ref{B2.37} for \(\xi_{q,a}^*\) at level \(L=N_0qd\), since the two twists agree on integers coprime to \(L\).
Apply Theorem~\ref{B2.37} to this bad set.
\end{proof}

\subsection*{Proof of \ref{B2.40}}\label{prf:B2.40}\phantomsection
\begin{proof}
Apply Lemmas~\ref{B2.38} and \ref{B2.39}, and take the larger of their constants.
A larger constant preserves the logarithmic derivative bound because
\(\log Y\geq0\).
Restriction from \([Y,Y^K]\) to \([Y,X]\) preserves every box.

Let \(G_1\) be the complement in \(P_1(z)\) of the candidates failing the first box.
For each \(q\in G_1\), let \(G_2(q)\) be the complement in its own candidate set \(P_2(z,z',q)\) of the candidates failing the second box.
The two bad-set estimates give the stated cardinality inequalities.
The candidate sets are nonempty by the same lemmas.
\end{proof}

\subsection*{Proof of \ref{B2.41}}\label{prf:B2.41}\phantomsection
\begin{proof}
Use Theorem~\ref{B2.40} with \(\eta=1/2\).
Its first good set has cardinality at least half the positive cardinality of \(P_1(z)\), so it is nonempty.
Choose \(q\) from it.
Choose \(d\) from the second good set associated to \(q\), which is nonempty by the same cardinality bound.
Membership in \(P_1(z)\) and \(P_2(z,z',q)\) gives \eqref{eq:aux-pair-rows}.
Membership in the corresponding good sets gives both box bounds.

For large \(Y\), we have \(e^{z/2}>N_0\) and \(z'\geq2z\).
Hence \(q>N_0\) and \(d>e^{z'/2}\geq e^z\geq q\).
\end{proof}

\subsection*{Proof of \ref{B2.43}}\label{prf:B2.43}\phantomsection
\begin{proof}
Apply Theorem~\ref{B2.41} with \(N_0=8M\), \(\tau=\varepsilon/2\),
\(K=A\), \(X=Y^A\), and \(z'=z'_0\).
It gives primes \(q,d\) satisfying the window conditions and \(8M<q<d\).
Thus neither prime divides \(M\), and \((M,qd)=1\).

Let \(f:\N\to\C\) agree, on integers coprime to \(8M\), with a character
\(\xi_8\) modulo \(8M\).
Let \(\chi\) be a primitive representative at level \(2Mqd\) of
\(n\mapsto f(n)(n/q)^i(n/d)^j\), where
\((i,j)\in\{0,1\}^2\setminus\{(0,0)\}\).
We claim that
\[
 \chi\in\mathcal B(H_j,C,\varepsilon/2,Y,Y^A)
 \tag{37}\label{eq:aux-final-box}
\]
holds, with the constant \(C\) from Theorem~\ref{B2.41}.

If \(j=1\), restrict the agreement defining \(\chi\) to integers coprime to
\(8Mqd\).
On these integers \(f=\xi_8\), so \(\chi\) represents the twist in the second selection condition with \(a=i\).
Its height is \(8Me^{z+z'_0}=H_1\).

If \(j=0\), then \(i=1\).
Choose the primitive representative \(\xi_8^*\) of \(\xi_8\), with conductor
\(f_0\mid8M\).
Since \(q>8M\), Lemma~\ref{B2.33} gives a primitive character of conductor
\(f_0q\) whose values are \(\xi_8^*(n)(n/q)\).
It agrees with \(\chi\) on integers coprime to \(2Mqd\).
The uniqueness part of Lemma~\ref{B2.33} therefore identifies the characters.
In particular, \(\chi\) represents \(\xi_8(n)(n/q)\) at level \(8Mq\).
The first selection condition gives \eqref{eq:aux-final-box} at height \(8Me^z=H_0\).

For the first family in Definition~\ref{B2.42}, induce the given character
\(\xi\) modulo \(M\) to level \(8M\).
Induction preserves its values on integers coprime to \(8M\).
The claim supplies \eqref{eq:aux-final-box}.
In particular, \(\chi\) is nonprincipal, so its centred sum is the function \(g_\chi\) of Definition~\ref{B2.21}.
\leanref{QH.\allowbreak{}pairEx\_\allowbreak{}g}{Development.\allowbreak{}AutomorphicArithmetic.\allowbreak{}MeanSquareBounds.\allowbreak{}pairEx}{61--78}
Its box yields
\(|g_\chi(y)|\leq y^{1/2+\varepsilon/2}\) for \(Y\leq y\leq Y^A\).

For \(j=0,1\), we have \(\log H_j=\log(8M)+z+jz'_0\).
For \(Y\geq\exp(1+\log(8M))\),
\[
 Y^{\varepsilon/6}\leq\log H_j
 \leq3Y^{\varepsilon/5}(\log Y)^2.
\]
By \eqref{eq:aux-final-box}, \(\chi\in\mathcal Z(H_j)\), so \eqref{eq:aux-kappa-four} applies with \(Q=H_j\) and \(A_1=3\).
It gives
\(\left|L'/L(1,\chi)\right|\leq4\log Y\leq4\log y\).
This proves \eqref{eq:aux-pair-good-bounds}.

Finally, let \(\rho\) be a real primitive character of conductor dividing \(8M\).
Induce it to level \(8M\) and apply the same claim with \(f=\rho\).
By \eqref{eq:aux-final-box}, \(L(s,\chi)\ne0\) throughout \(\mathcal R(H_j)\), as required by \eqref{eq:aux-pair-good-boxes}.
Take \(Y_1=\max\{Y_{\rm pair},Y_4,\exp(1+\log(8M))\}\), where
\(Y_{\rm pair}\) is the threshold in Theorem~\ref{B2.41} for the parameters
used above and \(Y_4\) is the threshold in \eqref{eq:aux-kappa-four} for \(A_1=3\).
\end{proof}

\section*{Proofs for Chapter~\ref{ch:arithmetic-window-host}: Building the arithmetic host of a window}

\subsection*{Proof of \ref{B2.44}}\label{prf:B2.44}\phantomsection
\begin{proof}
The surjection $A^{(e)}\to H_e$ has kernel $\{1\}\times C_e$, where
$C_e=\{b\in B:(1,b)\in A^{(e)}\}$. Thus $|A^{(e)}|=2^e|C_e|$.
Reduction gives $C_{e+1}\subset C_e$, so $|C_e|$ stabilizes to $c$.
The projection onto $B$ has kernel $A_0$ embedded in the cyclic group
$H_e$, and $|B|=|\Xi|$.

Let $\eta(a)\in A^{(e)}$ be the joint image, and let
$w:\{\pm1\}\to A^{(e)}$ take a sign to the joint image of that integer.
The Artin map is
\begin{equation}\label{eq:artin-with-sign}
 \operatorname{Art}(a)=
 \left(\left(\frac aq\right),\left(\frac ad\right),
       \eta(a)w\left(\left(\frac aq\right)\right)\right).
\end{equation}
Its kernel has both signs $1$ and $\eta(a)=1$, which gives the
congruence and $\psi(a)=1$. Primitivity of $\psi$ gives $M\mid Q_S$
because $\psi$ factors through $Q_S$.
\end{proof}

\subsection*{Proof of \ref{B2.46}}\label{prf:B2.46}\phantomsection
\begin{proof}
Apply \ref{B2.44} with $e_0\geq\max(2,4M,64o_\psi)$.
Take $E$ to be the fixed field of the kernel of \eqref{eq:artin-with-sign}
inside $\Q(\zeta_{Q_Sqd})$. For each unit residue modulo $Q_S$,
the Chinese remainder theorem realizes all four sign pairs, so the
Artin map surjects onto $\{\pm1\}^2\times A$.

The third coordinate projects to $[a]\in H_e$ because $w(\pm1)$ projects
to $1$. For $a\equiv1\bmod Q_Sq$, the Artin image lies in $\{1,\epsilon_d\}$. For $a\equiv1\bmod Q_Sd$ the third coordinate is
$w((a/q))$ and the image squares to $1$. For each unit $a$ modulo $N=Q_Sqd$, the automorphism $\zeta_N\mapsto\zeta_N^a$ restricts on $E$ to the element with coordinates $\operatorname{Art}(a)$, as required in Definition~\ref{B2.01}.
For $\ell=q,d$, the quadratic Gauss sum
$\tau_\ell=\sum_{u\bmod\ell}(u/\ell)\zeta_\ell^u$ satisfies
$\tau_\ell^2=-\ell$ and $\sigma_a(\tau_\ell)=(a/\ell)\tau_\ell$
\cite[Theorems~9.5 and~9.7]{montgomery-vaughan-2007}.
The Artin kernel therefore fixes both sums, and the two signs cut
out $\Q(\sqrt{-q})$ and $K=\Q(\sqrt{-d})$.
Complex conjugation maps to $((-1,-1),1)$. These data give the raw
host of \ref{B2.01}.

The local calculation of \ref{B2.03} gives trivial inertia outside
$Q_Sqd$ and inertia generated by $\epsilon_d$ at $d$. Since $\epsilon_d$ acts nontrivially on $K$, the inertia at $d$ in
$E/K$ is trivial. Thus every prime factor of $\mathfrak D_{E/K}$
lies above a rational prime dividing $Q_Sq$. These primes split in
$K$ by \ref{B2.03}. For each conjugate pair of primes of $K$ above
them, choose one member and put into $\mathfrak b$ every prime
factor of $\mathfrak D_{E/K}$ above that member, with its
multiplicity. Complex conjugation preserves the different and
interchanges the two members of each pair. Every element of
$\Gal(E/K)$ preserves the different and fixes each prime of $K$.
Consequently
\[
 \mathfrak D_{E/K}=\mathfrak b\,\overline{\mathfrak b},\qquad
 \mathfrak b+\overline{\mathfrak b}=\mathcal O_E,\qquad
 \sigma(\mathfrak b)=\mathfrak b\quad(\sigma\in\Gal(E/K)).
\]


Fix a prime $\ell\mid Q_S$. If $2^{e+2}\mid Q_S/\ell$, minimality of
$Q_S$ gives a unit $1\bmod Q_S/\ell$ with nontrivial joint image.
Otherwise $\ell=2$ and $v_2(Q_S)=e+2$. The unit $1+Q_S/2$ is
$1\bmod Q_S/2$ and $1+2^{e+1}\bmod 2^{e+2}$, so its $H_e$-coordinate
is nontrivial. Lift the witness to a unit equal to $1\bmod qd$.
The lift is $1\bmod Q_Sqd/\ell$ and its Artin image is $((1,1),z)$
with $z\ne1$. The bounds follow from
$Q_S\leq2^{e+2}M\leq4Mn\leq n^2$ and
$o_\phi^{64}=o_\psi^{64}\leq2^{64o_\psi}\leq2^e$.
Every odd prime $\ell\mid Q_S$ divides $M$, so $4\ell\leq B_0$.

For $\xi\in\Xi$ let $\xi_A$ be its pullback to $A$. The character
$\theta_\xi((s_q,s_d),z)=\xi_A(z\,w(s_q)^{-1})$ is trivial on $A_0$
and takes value $1$ on $\epsilon_d$. Substitution into
\eqref{eq:artin-with-sign} gives the primitive representative
$\xi^{\mathrm{prim}}$. On a unit equal to $1\bmod Q_Sd$ and nonsquare
modulo $q$, the retained characters take value $1$ and $\theta_q$ takes value $-1$. Apply \ref{B2.44} and \ref{B2.45} for the kernel and
twist descriptions.
\end{proof}

\subsection*{Proof of \ref{B2.47}}\label{prf:B2.47}\phantomsection
\begin{proof}
Let $\theta\in\mathcal B=(X_+\setminus K_H(\mathcal X))\cup X_-$ and let $\chi$ be a primitive Dirichlet representative of $\theta$.
Since every element of $X_+\cup X_-$ is trivial on $A_0$, Lemma~\ref{B2.45} gives $\xi\in\Xi=\langle\psi\rangle$ and $i,j\in\{0,1\}$ such that $\theta(\epsilon_d)=(-1)^j$ and
\[
 \theta(\operatorname{Art}(a))
   =\xi(a)\left(\frac aq\right)^i\left(\frac ad\right)^j
 \qquad\bigl((a,Q_Sqd)=1\bigr).
\]

We claim that $(i,j)\ne(0,0)$.
Suppose $(i,j)=(0,0)$.
If $\theta\in X_-$, then $\theta(\epsilon_d)=-1$, contradicting $\theta(\epsilon_d)=(-1)^0=1$.
If $\theta\in X_+\setminus K_H(\mathcal X)$, then $\theta(\operatorname{Art}(a))=\xi(a)$ for every integer $a$ coprime to $Q_Sqd$.
Since $M\mid Q_S$, every integer coprime to $Q_Sqd$ is coprime to $M$, so the primitive representative $\xi^{\mathrm{prim}}\in\mathcal X$ of $\xi$ agrees with $\xi$ and hence with $\theta\circ\operatorname{Art}$ on integers coprime to $Q_Sqd$.
Thus $\theta$ is represented by $\xi^{\mathrm{prim}}\in\mathcal X$, contradicting $\theta\notin K_H(\mathcal X)$.
Hence $(i,j)\ne(0,0)$.

Let $a$ be an integer with $(a,2Mqd)=1$.
Then $(a,q)=(a,d)=1$, and since every prime divisor of $Q_S$ divides $2M$, we also have $(a,Q_S)=1$, so $(a,Q_Sqd)=1$.
Combining the representative identity $\chi(a)=\theta(\operatorname{Art}(a))$ with \eqref{eq:host-twist-form} yields
\[
 \chi(a)=\xi(a)\left(\frac aq\right)^i\left(\frac ad\right)^j=f_{\xi,i,j}(a)
 \qquad\bigl((a,2Mqd)=1\bigr).
\]
Therefore $\chi$ is a primitive representative of $f_{\xi,i,j}$ at level $2Mqd$.
Since $(q,d)$ is a good pair for modulus $M$ on $[Y,x]$ with exponent $\varepsilon$ and $Y\leq x$, \eqref{eq:aux-pair-good-bounds} of Definition~\ref{B2.42} applies at $y=x$ and gives
\[
 |g_\chi(x)|\leq x^{1/2+\varepsilon/2}\leq x^{1/2+\varepsilon},
 \qquad
 \left|\frac{L'}{L}(1,\chi)\right|\leq 4\log x,
\]
using $x\geq1$ and $\varepsilon>0$ in the first inequality.
\end{proof}

\subsection*{Proof of \ref{B2.48}}\label{prf:B2.48}\phantomsection
\begin{proof}
Set \(P=2+8\varepsilon\) and \(\lambda=c^{1/P}\). For \(v\ge1\) take \(a(v)=\lceil\log_2(v^P)\rceil\), and take \(a(v)=0\) for \(v<1\). Then \(v^P\le2^{a(v)}\le2v^P\) for \(v\ge1\). Apply Definition~\ref{B6.03} to this page function and exponent \(401\), obtaining \(J\ge1\) and \(X_0\). Let \(Y_1\) be the cutoff of \ref{B2.43} for this \(J\) and \(M=\operatorname{cond}\psi\).

Given \(y\), choose \(X\ge X_0\) with \(Y=X^{1/J}\ge\max(1,y,Y_1,\lambda^3,\lambda^{400})\). By \ref{B2.43} there is a good pair on \([Y,Y^J]=[Y,X]\). Its prime conditions at \(Y\) give the window pair required by \ref{B6.03}. Keeping this pair fixed, choose \(W\ge Y\) with \(W^2\le X\) as in that definition. Put \(Z=\lambda W\ge y\).

For \(Z\le x\le Z^{3/2}\), use \(v=x/\lambda\) and \(e=a(v)\). Lemma~\ref{B2.48a} gives \(W\le v\le x\le W^2\le X\) and the required bounds on \(c2^e\). Thus \(x^{1/J}\le Y\le x\), and restricting the good pair to \([Y,x]\) gives \eqref{eq:window-datum}. Every field with the prescribed kernel at \((e,q,d)\) satisfies the upstream bound at \(v\). All its summands are nonnegative. The cutoff and budget inequalities of Lemma~\ref{B2.48a} therefore give \eqref{eq:window-vv} for the same field and all the local data of Definition~\ref{B6.01}.
\end{proof}

\subsection*{Proof of \ref{B2.49a}}\label{prf:B2.49a}\phantomsection
\begin{proof}
Write $L_\chi(s)=\sum_{n\geq1}a_\chi(n)n^{-s}$ and
$P_\chi(s)=\sum_{n\geq1}\widehat\lambda_\chi(n)n^{-s}$.
For $\Re s>1$ the convolution gives $L'_\chi=-P_\chi L_\chi$.
Integration from $+\infty$, with $a_\chi(1)=1$ and
$|\widehat\lambda_\chi|\leq2\Lambda$, gives
$|L_\chi(\sigma+it)|\geq\zeta(\sigma)^{-2}$.
Put $r=\kappa|G|$. By \eqref{eq:host-family-counts} and character
orthogonality, $E_\chi(y)=\sum_{n\leq y}a_\chi(n)-\mathbf1_{\chi=1}ry$
satisfies $|E_\chi(y)|\leq10D\sqrt y$ for $y\geq1$. Partial summation gives
\[
 L_\chi(s)=\mathbf1_{\chi=1}\frac{rs}{s-1}
        +s\int_1^\infty E_\chi(y)y^{-s-1}\,dy.
\]
Hence $F_\chi$ extends to $\Re s>1/2$ and
$|F_\chi(s)|\leq300D(1+|s|)^2$ for $\Re s\geq11/20$.
On the disc of radius $7/10$ about $z_t=5/4+it$, the logarithm of the
ratio of the maximum modulus to $|F_\chi(z_t)|$ is $O(\ell(t))$.
Apply \cite[Lemmas~6.1 and~6.3]{montgomery-vaughan-2007} to
$z\mapsto F_\chi(z_t+7z/10)$ with $r=5/7$ and $R=6/7$ in those lemmas.
For $|s-z_t|\leq1/2$ away from the zeros,
\begin{equation}\label{eq:host-local-logderivative}
 \frac{F'_\chi}{F_\chi}(s)=
 \sum_{|\rho-z_t|\leq3/5}\frac{m_\rho}{s-\rho}+O(\ell(t)),
 \qquad \sum_{|\rho-z_t|\leq3/5}m_\rho\ll\ell(t),
\end{equation}
where $m_\rho$ is the multiplicity. Since
$P_\chi(s)=\mathbf1_{\chi=1}/(s-1)-F'_\chi(s)/F_\chi(s)$ and all zeros
have real part at most $1$, each zero contributes a nonpositive real
part to $P_\chi(s)$ when $\Re s>1$.
\leanref{GlobalTheory.\allowbreak{}HkLocal.\allowbreak{}logDeriv\_\allowbreak{}local}{Development.\allowbreak{}AbelianVarietyTheory.\allowbreak{}Proofs.\allowbreak{}GlobalInterpolation.\allowbreak{}RayClassLFunctionLogDerivLocalFormula}{21--157}

Write a zero as $1-b+it$, with $0\leq b\leq v\leq1/400$ and $v>0$.
At $\sigma=1+6v$, the inequality $3+4\cos u+\cos 2u\geq0$ gives
$3\Re P_1(\sigma)+4\Re P_\chi(\sigma+it)+\Re P_{\chi^2}(\sigma+2it)\geq0$.
If $\chi^2\ne1$, both shifted pole terms vanish. If $|t|\geq30v$,
their real parts before weighting are at most $1/(156v)$. In either
case the indicated zero and \eqref{eq:host-local-logderivative} give
\[
 4/(7v)\leq3/(6v)+5/(156v)+O(\ell(t)).
\]
For $\chi^2=1$ and $0<|t|<30v$, sum $P_\xi(\sigma)$ over
$\xi\in\langle\chi\rangle$ at $\sigma=1+60v$. Orthogonality makes the
coefficients nonnegative. The conjugate pair contributes
$2(60v+b)/((60v+b)^2+t^2)\geq120/(4621v)$ against one pole of size
$1/(60v)$. Both comparisons give $1/v\ll\ell(t)$, which excludes
these zeros for $v=c_0/\ell(t)$ and sufficiently small $c_0$.
\leanref{GlobalTheory.\allowbreak{}Hk.\allowbreak{}ZeroFree.\allowbreak{}zero\_\allowbreak{}free\_\allowbreak{}region}{Development.\allowbreak{}AbelianVarietyTheory.\allowbreak{}Proofs.\allowbreak{}GlobalInterpolation.\allowbreak{}RayClassLFunctionZeroFreeRegion}{286--312}

If the strip contains two real zeros counted with multiplicity, sum
$P_\xi(\sigma)$ at $\sigma=1+2c_0/\ell(0)$ over the subgroup generated
by their characters. It has at most four elements and the sum has
nonnegative coefficients. By \eqref{eq:host-local-logderivative} the
zeros contribute at least $2\ell(0)/(3c_0)$ against one pole
$\ell(0)/(2c_0)$ and an error $O(\ell(0))$. This is impossible for
small $c_0$.
\leanref{GlobalTheory.\allowbreak{}Hk.\allowbreak{}RealZero.\allowbreak{}main}{Development.\allowbreak{}AbelianVarietyTheory.\allowbreak{}Proofs.\allowbreak{}GlobalInterpolation.\allowbreak{}RayClassLFunctionUniqueSimpleExceptionalZero}{135--259}

Write $H_\chi(t)=\sum_n a_\chi(n)e^{-nt}$. The partial summation formula
gives $F_1(1)=r\geq D^{-1}$. For $n\geq1$, nonnegativity and
\eqref{eq:host-family-counts} give $a_g(n)\leq\kappa n+10\sqrt n\leq14n$.
For a nontrivial real $\chi$, \eqref{eq:host-family-counts} and this
bound give
\[
 |H_1(t)-r/t|+|H_\chi(t)|\ll D^{11}\min(t,1/t).
\]
For $t\geq1$, summing the coefficient bound against $e^{-nt}$ gives
$|H_\chi(t)|\ll De^{-t}$. Together with the bound at zero, this extends
$\int_0^\infty H_\chi(t)t^{s-1}\,dt=\Gamma(s)F_\chi(s)$
from $\Re s>1$ to $s=1$. Hence $F_\chi(1)=\int_0^\infty H_\chi(t)\,dt$.

Put $J(y)=\int_0^\infty e^{-y/w-w}\,dw/w$. For $0<y\leq1$,
expand the product and use $a_1*a_\chi\geq0$ and
$(a_1*a_\chi)(1)=1$ to get
\[
 \begin{aligned}
 J(1)\leq J(y)&\leq\sum_{n\geq1}(a_1*a_\chi)(n)J(ny)\\
 &=\int_0^\infty H_1(y/w)H_\chi(w)\,\frac{dw}{w}
 =\frac r y F_\chi(1)+O(D^{22}\sqrt y).
 \end{aligned}
\]
The inequality $e^{-z}\ll z^{-3}$ gives $J(z)\ll z^{-3}$, so
$\sum_n a_1(n)n^{-3}<\infty$ gives absolute convergence. In the error
integral use $\min(y/w,w/y)\leq\sqrt{y/w}$ and the integrability of
$\min(w,1/w)w^{-3/2}$. Take $\sqrt y=\delta D^{-22}$ for a sufficiently
small absolute $\delta>0$. Since $r\leq4D$, this gives
$|F_\chi(1)|\gg D^{-45}$.
\leanref{GlobalTheory.\allowbreak{}Hk.\allowbreak{}RealChar.\allowbreak{}norm\_\allowbreak{}F\_\allowbreak{}one\_\allowbreak{}ge}{Development.\allowbreak{}AbelianVarietyTheory.\allowbreak{}Proofs.\allowbreak{}GlobalInterpolation.\allowbreak{}NormRayClassFamilyFOneGeOfRealCharacter}{343--423}

The Cauchy integral formula with radius
$1/10$ and the bound for $F_\chi$ give $|F'_\chi|\leq18000D$ on
$|s-1|\leq1/4$. For either a trivial or a nontrivial real character,
$|F_\chi(1)-F_\chi(\beta)|\leq18000D(1-\beta)$ gives
$1-\beta\gg D^{-46}$.
\leanref{GlobalTheory.\allowbreak{}Hk.\allowbreak{}LFam.\allowbreak{}real\_\allowbreak{}zero\_\allowbreak{}far\_\allowbreak{}from\_\allowbreak{}one}{Development.\allowbreak{}AbelianVarietyTheory.\allowbreak{}Proofs.\allowbreak{}GlobalInterpolation.\allowbreak{}ExistsRayClassLFunctionRealZeroFarFromOne}{47--77}
\end{proof}

\subsection*{Proof of \ref{B2.49b}}\label{prf:B2.49b}\phantomsection
\begin{proof}
Put $w=\sqrt{\log u}$ and $\alpha=c_0/32$. For $\alpha w\leq1$,
Chebyshev's bound gives $|\Psi_\chi(u)|\leq2\sum_{n\leq u}\Lambda(n)\ll u$
\cite[Theorem~2.4]{montgomery-vaughan-2007}, and any $\beta\in Z_\chi$
has $\beta>1/2$. Otherwise take $T=e^{\alpha w}$, with $2\leq T\leq u$
and $\ell(T)\leq3w$. Truncated Perron inversion on $\Re s=1+1/\log u$
\cite[Theorem~5.2 and Corollary~5.3]{montgomery-vaughan-2007}
costs $O(u\log^2u/T+\log u)$. For $u/2<n<2u$, sum the kernel error
$O(\log u)\min(1,u/(T|n-u|))$, with the minimum taken as $1$ at $n=u$.
For the other $n$, use $|\log(u/n)|\geq\log2$ and
$\sum_n\Lambda(n)n^{-1-1/\log u}\ll\log u$ by partial summation
from Chebyshev's bound.
\leanref{GlobalTheory.\allowbreak{}PerronCoeff.\allowbreak{}truncated\_\allowbreak{}perron\_\allowbreak{}of\_\allowbreak{}le\_\allowbreak{}two\_\allowbreak{}vonMangoldt}{Development.\allowbreak{}AbelianVarietyTheory.\allowbreak{}Proofs.\allowbreak{}GlobalInterpolation.\allowbreak{}TruncatedPerronOfLeTwoVonMangoldt}{9--158}

In $1-c_0/(2\ell(t))\leq\Re s\leq3$ at height $t$, away from
$Z_\chi$, we have
\[
 \frac{F'_\chi}{F_\chi}(s)-\sum_{\beta\in Z_\chi}\frac1{s-\beta}
 = O_{c_0}(\ell(t)^2).
\]
In \eqref{eq:host-local-logderivative} there are $O(\ell(t))$ zeros
counted with multiplicity. For each nonexceptional zero in that sum,
$\ell(\Im\rho)\leq4\ell(t)/3$ gives a horizontal distance
$\gg c_0/\ell(t)$ from $s$. If $|\beta-z_t|>3/5$ and
$|s-z_t|\leq1/2$, then $|s-\beta|>1/10$. For $\Re s\geq7/4$ the
absolutely convergent series for $P_\chi$ gives the bound directly.
\leanref{GlobalTheory.\allowbreak{}Hk.\allowbreak{}LogDerivHalf.\allowbreak{}logDeriv\_\allowbreak{}sub\_\allowbreak{}exceptional\_\allowbreak{}le}{Development.\allowbreak{}AbelianVarietyTheory.\allowbreak{}Proofs.\allowbreak{}GlobalInterpolation.\allowbreak{}LogDerivSubExceptionalLe}{92--261}

Take $b=c_0/(8\ell(T))$. Shift the integral of $P_\chi(s)u^s/s$ to
the line $\sigma=1-2b$ or $1-4b$ at distance at least $b$ from
$Z_\chi$. The residues give $\mathbf1_{\chi=1}u$ and $-u^\beta/\beta$
for $\beta>\sigma$. A zero to the left contributes at most
$2u^{1-2b}$ in absolute value. On the new sides
$|P_\chi|\ll_{c_0}\ell(T)^2+1/b$. The horizontal and vertical integrals
cost $O_{c_0}(w^2u/T)$ and $O_{c_0}(Tw^2u^{1-2b})$, respectively.
Since $2b\log u\geq c_0w/12\geq2\log T$, the total error is
$O_{c_0}(uw^4e^{-\alpha w})=O_{c_0}(ue^{-\alpha w/2})$.
\end{proof}

\subsection*{Proof of \ref{B2.49c}}\label{prf:B2.49c}\phantomsection
\begin{proof}
Use the $c_0$ of \ref{B2.49a} in \ref{B2.49b}. Put
$L=\sqrt{\log X}$, $U=X(1+h)$ and $\Delta=c_9D^{-a}$. Choose
$c=\min\{1,c_7/(2(a+2)),c_7c_9/(24C)\}$. Then $cL\geq\log D\geq1$,
so $X\geq\mathrm e$ and both endpoints satisfy $\log D\leq\sqrt{\log u}$.
By orthogonality,
$\sum_\chi\Psi_\chi(u)=|G|\sum_{1\leq n\leq u}\lambda_0(n)$.
Subtract \eqref{eq:host-endpoint-weight} at the endpoints and write
$W=\sum_{X<n\leq U}\lambda_0(n)$. Since $U\leq2X$,
\[
 |G|W=hX-B+R,\qquad |R|\leq3CDX e^{-c_7L}.
\]
Here $B=0$ in the absence of an exceptional pair. If the pair is present,
\[
 B=\int_X^U t^{\beta-1}\,dt\leq hX e^{-\Delta\log X}\leq hX(1-\Delta/2)
\]
since $0<\Delta\leq1$ and $\log X\geq1$. The remainder satisfies
\[
 \frac{|R|}{hX\Delta}\leq\frac{3C}{c_9}e^{((a+2)c-c_7)L}
 \leq\frac{3C}{c_9}e^{-c_7L/2}\leq\frac14 .
\]
The last step uses $e^{c_7L/2}\geq c_7L/2\geq c_7/(2c)\geq12C/c_9$.
Hence $|G|W\geq hX\Delta/4$. Use $|G|\leq D$ and take
$A=a+1$, $\eta=c_9/4$.
\end{proof}

\subsection*{Proof of \ref{B2.50}}\label{prf:B2.50}\phantomsection
\begin{proof}
Let $K=\Q(\sqrt{-d})$, $G$ its ray class group modulo $N\mathcal O_K$,
and $D=4dN^2$. For ideals prime to $N$ of norm $n$ in class $g$,
let $a_g(n)$ count the ideals and $\lambda_g(n)$ sum their ideal von
Mangoldt weights. The kernel of the map to the ordinary class group
is a quotient of $(\mathcal O_K/N\mathcal O_K)^\times$, so $|G|\leq h_KN^2$,
where $h_K$ is the class number. Minkowski's bound
\cite[Chapter~I, Theorem~6.3]{neukirch-1999} gives an ideal of norm
at most $\sqrt d$ in each ordinary class. At norm $n$ there are at most
$\tau(n)\leq n$ ideals by prime ideal factorization, where $\tau$
counts positive divisors. Thus $h_K\leq d$ and $|G|\leq D$.

Choose an integral ideal $J$ prime to $N$ in class $-g$.
Ideals $I$ in class $g$ correspond to generators $z$ of $IJ$ with
$z\equiv1\bmod N$. Such a generator is unique since $N\geq3$ and
$\mathcal O_K^\times=\{\pm1\}$. The possible $z$ form a translate of
$NJ$ in $J$. In a reduced basis of $NJ$, write
\[
 \frac{|z|^2}{\mathrm{N}J}
 =\alpha(m+\theta_1+\ell(k+\theta_2))^2+\gamma(k+\theta_2)^2
 \qquad(m,k\in\Z),
\]
where $\theta_1,\theta_2$ describe the translate. For the shortest
basis vector $w_1$, $\alpha=|w_1|^2/\mathrm{N}J\geq N^2$ since
$w_1/N\in J$. Choose the second vector with projection
$|\ell|\leq1/2$ on the first. Its length is at least $|w_1|$, hence
$\gamma\geq3\alpha/4$. The covolume gives $\alpha\gamma=dN^4/4$, so
$\kappa=\pi/\sqrt{\alpha\gamma}=2\pi/(\sqrt d\,N^2)$ lies in $[D^{-1},4]$.
Each horizontal chord costs at most one in rounding. The sum of chord
lengths differs from its integral by at most the total variation
$4\sqrt{y/\alpha}$. The area is $\kappa y$ and the count error is at
most $2\sqrt{y/\gamma}+3+4\sqrt{y/\alpha}\leq10\sqrt y$.

For the Gaussian sum use Poisson summation twice
\cite[Theorem~D.3]{montgomery-vaughan-2007}. At scale $s$, the
one-dimensional Gaussian sum divided by its integral differs from
$1$ by at most $2\sum_{j\geq1}e^{-\pi^2j^2/s}\leq s^2$, by
$e^{-z}\leq2/z^2$. First sum over $m$, uniformly in the shift
$\theta_1+\ell(k+\theta_2)$, then over $k$. The relative error is at
most $(\alpha t)^2+(\gamma t)^2+(\alpha\gamma t^2)^2$.
For $0<t\leq1$ the absolute error is at most
$3\pi(\alpha\gamma)^2t\leq D^{10}t$. This gives
\eqref{eq:host-family-counts}.

The counts give $a_g(0)=0$ and $a_g(1)=\mathbf1_{g=0}$. Factorization
gives the logarithmic convolution and $\sum_g\lambda_g(n)\leq2\Lambda(n)$
since $[K:\Q]=2$. For real $\chi$, the coefficient at an ideal $I$
prime to $N$ in $L_1L_\chi$ is
$\prod_{\mathfrak p^e\parallel I}\sum_{j=0}^e\chi(\mathfrak p)^j$.
It is nonnegative and at least one on squares. Summing over ideals
of norm $n$ gives $(a_1*a_\chi)(n)\geq0$. The $a_1(n)$ ideals of
norm $n$ prime to $N$ have distinct squares of norm $n^2$, each with
coefficient at least one. Hence $(a_1*a_\chi)(n^2)\geq a_1(n)$.
By \ref{B2.49c}, $W=\sum_{X<n\leq U}\lambda_0(n)\geq\eta hX/D^A$,
where $U=X(1+h)$.
\leanref{GlobalTheory.\allowbreak{}Hk.\allowbreak{}rayCountTheta}{Development.\allowbreak{}AbelianVarietyTheory.\allowbreak{}Proofs.\allowbreak{}GlobalInterpolation.\allowbreak{}RayCountThetaFibreCount}{53--122}
\leanref{GlobalTheory.\allowbreak{}Hk.\allowbreak{}exists\_\allowbreak{}lFam\_\allowbreak{}ray}{Development.\allowbreak{}AbelianVarietyTheory.\allowbreak{}Proofs.\allowbreak{}GlobalInterpolation.\allowbreak{}ExistsLFamRay}{73--129}

If $p$ is prime and $\lambda_0(p)>0$, an ideal of norm $p$ has a
generator $z=(u-v)/2+v(1+\sqrt{-d})/2\equiv1\bmod N$. Its norm
and congruence give \eqref{eq:norm-one-coordinates}. Each such prime
contributes at most $2\log U$, and the other positive weights lie
on higher prime powers. Thus
\[
 W\leq2\log U\,|\mathcal P(d,N;X,U)|+2\sum_{j\geq2}\vartheta(U^{1/j}),
\]
where $\vartheta(y)=\sum_{p\leq y}\log p$. By Chebyshev's bound
\cite[Theorem~2.4]{montgomery-vaughan-2007}, the last term is
$O(\sqrt U+U^{1/3}\log U)=O(\sqrt X)$. For $v=\sqrt{\log X}$,
\[
 hX/(D^A\sqrt X)\geq\exp(v^2/2-(A+1)cv),\qquad v\geq\log4/c .
\]
Shrink $c$ so that the higher powers cost at most $\eta hX/(2D^A)$.
Since $\log U\leq2\log X$, take $c'=\eta/8$.
\end{proof}

\subsection*{Proof of \ref{B2.51}}\label{prf:B2.51}\phantomsection
\begin{proof}
Take $V=(kb)^8$ and $D=4dN^2$. Since $d\mid N$,
$\log D\leq\log4+9x^{\varepsilon/5}(\log x)^2$ and
$\log V\geq8x^\varepsilon$. Thus $\log D=o(\sqrt{\log V})$ and
$A\log D+\log\log V-\log c'\leq(\log V)/10$ for large $x$.
Apply \ref{B2.50} with $h=1$ to obtain
\begin{equation}\label{eq:norm-candidate-count}
 |\mathcal P(d,N;V,2V)|\geq V^{9/10}.
\end{equation}
\leanref{UpperTheory.\allowbreak{}T3.\allowbreak{}normPrimesN}{Development.\allowbreak{}AutomorphicArithmetic.\allowbreak{}ShortNorms.\allowbreak{}normPrimesN}{108--158}
By \eqref{eq:host-rank-ranges}, $r\leq x^5$ for large $x$.
Fix $g_*\in[V,2V]\cap\Z$. For each $j$, four distinct exceptional
prime powers would divide $g_*^j-1$ and have product greater than
$g_*^j$. Hence $|\mathcal B_r(g_*)|\leq3r+1$.

For composite $\ell$, membership $\ell\in\mathcal B_r(g)$ forces
$g=\ell$. For prime $\ell$ and $1\leq j\leq r$, take the least $a$
with $\ell^a>V^{j/4}$. For $g\in[V,2V]$, the inequality $g^j<\ell^{4v_\ell(g^j-1)}$ gives $v_\ell(g^j-1)\geq a$, hence $g^j\equiv1\bmod\ell^a$. For odd $\ell$ there are at most $j$ roots
by cyclicity. For $\ell=2$, the cyclic subgroup $\langle5\rangle$
has index at most two, so there are at most $2j$. Each residue class
contains at most $1+V/\ell^a\leq1+V^{3/4}$ candidates. For each
$\ell\in\mathcal B_r(g_*)$, sum over $1\leq j\leq r$ and include
$g=\ell$. At most
\begin{equation}\label{eq:bb-exclusion-count}
 (3r+1)\bigl(1+2r^2(V^{3/4}+1)\bigr)
\end{equation}
integers $g\in[V,2V]$ have $\mathcal B_r(g)\cap\mathcal B_r(g_*)\ne\varnothing$.
Since $20r^3<V^{3/20}$ for large $x$, \eqref{eq:bb-exclusion-count}
is less than $V^{9/10}$. Choose $g_0\in\mathcal P(d,N;V,2V)$.
By \eqref{eq:norm-candidate-count}, some $g_1\in\mathcal P(d,N;V,2V)$
lies outside the exclusion set for $g_*=g_0$. Disjointness gives
$g_0\ne g_1$.

Let $A\geq1$ be the exponent of \ref{B2.49}. Choose $p,p',q_0$
successively in the classes $1$ modulo $15N$, $15Np$, and $Npp'$.
They are pairwise distinct. Since $N\geq100$,
\[
 p\leq N^{2A},\qquad p'\leq N^{(2A+2)A},\qquad
 q_0\leq N^{(1+2A+(2A+2)A)A}.
\]
Take $C_8=(1+2A+(2A+2)A)A$. For large $x$,
$C_8\log N\leq3C_8x^{\varepsilon/5}(\log x)^2<x^\varepsilon$.
Thus $p,p',q_0<b\leq(kb)^8\leq g_i$.
\leanref{GlobalTheory.\allowbreak{}choice\_\allowbreak{}count}{Development.\allowbreak{}AbelianVarietyTheory.\allowbreak{}Proofs.\allowbreak{}FoundationProofs.\allowbreak{}DivLeRpowThreeQuarters}{221--222}
\end{proof}

\subsection*{Proof of \ref{B2.52}}\label{prf:B2.52}\phantomsection
\begin{proof}
The primitive conductor $N_\chi$ divides $Q_Sqd$. For a real
character, the local component at an odd prime power is trivial on the
principal units, which have odd order. At a power of $2$ the component
factors through modulus $8$, since every unit congruent to $1\bmod8$
is a square. The Chinese remainder theorem separates the components
at $q$ and $d$, each either trivial or Legendre. The remaining
component has a primitive real representative $\rho$ with
$N_\rho\mid8Q_{\mathrm{odd}}$. Thus on units modulo $Q_Sqd$,
\[
 \chi(a)=\rho(a)\left(\frac aq\right)^i\left(\frac ad\right)^j
 \qquad(i,j\in\{0,1\}).
\]
The product of these primitive characters at coprime conductors is
primitive at $N_\rho q^i d^j$, so \ref{B2.33} gives
$N_\chi=N_\rho q^i d^j$. If $q$ or $d$ divides $N_\chi$, the twist
is nonzero and $j=\mathbf1_{d\mid N_\chi}$.
Every integer coprime to $2Mqd$ is coprime to $Q_Sqd$ by the hypothesis on the prime divisors of $Q_S$. The twist identity therefore holds on the units used in Definition~\ref{B2.42}. Since $N_\rho\mid8M$, the boxes in \ref{B2.42}
apply. The prime sizes there give the conductor bound.

The nontrivial real host characters of conductor prime to $qd$ have
conductor at most $8\lfloor B\rfloor$.
Nonvanishing at $1$ gives this finite family a common zero-free
interval to the left of $1$
\cite[Theorem~4.9]{montgomery-vaughan-2007}. Take the smaller of its
width and $1/400$ to cover the other real host characters.

Let $\chi$ be any primitive host character and put
$\lambda=\max(1,\log N_\chi)$. Combining the classical zero-free
region with the preceding exclusion gives the following bound for
each nontrivial zero $\rho$ of $L(s,\chi)$
\cite[\S\S~6.1 and~11.1]{montgomery-vaughan-2007}.
\begin{equation}\label{eq:thin-zero-free-region}
 \Re\rho\leq1-\frac{c}{\lambda+\log(|\Im\rho|+3)},
\end{equation}
where $c>0$ depends on $B$. By the functional equation, apply
\eqref{eq:thin-zero-free-region} to $1-\rho$ for $\overline\chi$.
For $|\Im\rho|\leq2$ this gives $|\rho|\gg_B1/\lambda$.
The zero count in a unit interval at height $h$ is
$O(\lambda+\log(|h|+3))$ with multiplicity
\cite[Theorems~10.13 and~10.17]{montgomery-vaughan-2007}.
The low zeros contribute $O_B(\lambda^2)$. Add the other unit-height
counts with weights $1/j$ to obtain
\[
 \sum_{|\Im\rho|\leq T}\frac{m_\rho}{|\rho|}
 \ll_B(\lambda+\log T)^2\qquad(T\geq1),
\]
where $m_\rho$ denotes multiplicity.

Put $\psi_0(\chi,u)=\sum_{n\leq u}\Lambda(n)\chi(n)-\mathbf1_{\chi=1}u$
and
\[
 Z_\chi(T)=\{\rho:0<\Re\rho<1,\ L(\rho,\chi)=0,\ |\Im\rho|<T\}.
\]
For $2\leq T\leq u$ the truncated explicit formula gives
\[
 \left|\psi_0(\chi,u)+\sum_{\rho\in Z_\chi(T)}m_\rho\frac{u^\rho}{\rho}\right|
 \ll \frac{u\log^2(N_\chi uT)}{T}+\log(N_\chi u)
      +\sum_{\rho\in Z_\chi(1)}\frac{m_\rho}{|\rho|}.
\]
For $u\geq\mathrm e$, put $v=\log u$, $T=\exp(\sqrt v)$ and
$D_v=\exp(-c'v/(\lambda+\sqrt v))$, with $c'=\min(c/3,1)$.
Then $2\leq T\leq u$ and \eqref{eq:thin-zero-free-region} gives
$|u^\rho|\leq uD_v$ for $\rho\in Z_\chi(T)$. The weighted zero
count bounds the zero sum by $O_B(u(v+\lambda)^2D_v)$ and
the last term of the error by $O_B(\lambda^2)$. Since $T^{-1}\leq D_v$
and $uD_v\geq1$, we obtain
\begin{equation}\label{eq:thin-primitive-weight}
 |\psi_0(\chi,u)|\ll_B u(v+\lambda)^2D_v.
\end{equation}

This bound implies absolute convergence of
$I_\chi=\int_1^\infty\psi_0(\chi,u)u^{-2}\,du$. For real $s>1$,
partial summation gives
\[
 \int_1^\infty\psi_0(\chi,u)u^{-s-1}\,du
 =-\frac{L'(s,\chi)}{sL(s,\chi)}-\frac{\mathbf1_{\chi=1}}{s-1}.
\]
As $s\downarrow1$, dominated convergence and the Laurent expansion
$\zeta(s)=(s-1)^{-1}+\gamma+O(s-1)$ give
$I_\chi=-L'(1,\chi)/L(1,\chi)$ for $\chi\ne1$ and
$I_1=-1-\gamma$. For $y\geq1$, termwise integration on $[1,y]$ gives
\[
 y\int_1^y\psi_0(\chi,u)u^{-2}\,du
 =\sum_{n\leq y}\Lambda(n)\chi(n)(y/n-1)
  -\mathbf1_{\chi=1}y\log y.
\]
Subtract $yI_\chi$ and use \ref{B1.26} to obtain
$G_\chi(y)=-y\int_y^\infty\psi_0(\chi,u)u^{-2}\,du$.

Since $N_\chi\mid Q_Sqd$, we have
$\lambda\leq1+2x^{\varepsilon/5}(\log x)^2$.
For $v\geq V=\log y\geq x^{\varepsilon/3}/2$ and large $x$, this gives
$\lambda\leq v^{61/100}$. Thus
$(v+\lambda)^2D_v\leq4v^2e^{-(c'/2)v^{39/100}}$.
Substituting $u=\mathrm e^v$ in \eqref{eq:thin-primitive-weight}
bounds $|G_\chi(y)|/y$ by
\[
 4C e^{-(c'/4)V^{39/100}}\int_0^\infty v^2e^{-(c'/4)v^{39/100}}\,dv
\]
with $C$ depending on $B$. The integral is finite and
$V^{39/100}\gg x^{13\varepsilon/100}$, with $13/100>1/8$.
This proves \eqref{eq:thin-character-bound} with a threshold depending
on $B,\varepsilon$.
\leanref{QH.\allowbreak{}psi\_\allowbreak{}thin}{Development.\allowbreak{}AutomorphicArithmetic.\allowbreak{}MeanSquareBounds.\allowbreak{}PsiThin}{385--422}
\leanref{QH.\allowbreak{}g\_\allowbreak{}integral}{Development.\allowbreak{}AutomorphicArithmetic.\allowbreak{}MeanSquareBounds.\allowbreak{}GIntegral}{229--234}
\leanref{QH.\allowbreak{}int\_\allowbreak{}thin}{Development.\allowbreak{}AutomorphicArithmetic.\allowbreak{}MeanSquareBounds.\allowbreak{}IntThin}{169--263}
\end{proof}

\subsection*{Proof of \ref{B2.53}}\label{prf:B2.53}\phantomsection
\begin{proof}
Let $N\geq1$ and $L\geq2$. One sphere of constant squared norm
in $\{0,\ldots,d-1\}^S$ contains at least $d^S/(Sd^2)$ points.
Map them to integers by $v\mapsto\sum_{\nu=0}^{S-1}v_\nu(4Ld)^\nu$.
In a putative relation \eqref{eq:no-short-relation}, each coefficient
has absolute value less than half the base. Successive reduction
modulo the base makes every coefficient zero. The three vectors
would lie on a line, which meets the sphere in at most two points.

Choose the parameters so that $4Ls<N$ and the cardinality is at least
\begin{equation}\label{eq:integer-sphere-bound}
 \frac NL\exp(-32\sqrt{\log N\log L}).
\end{equation}
If this is at most $1$, take $\{0\}$. Otherwise put $v=\log N$,
$w=\log L$, $\sigma=\sqrt{vw}$, $S=\lfloor\sigma/w\rfloor$,
$a=(v-w-3)/S$ and $d=\lfloor e^a/(4L)\rfloor$.
Then $S\geq32$, $a\geq7+w$ and $e^a/(8L)\leq d\leq e^a/(4L)$.
Since $N=e^{aS}Le^3$, we have $4L(4Ld)^S\leq N$.
The sphere count gives \eqref{eq:integer-sphere-bound} because
$e^3(8L)^SS e^{2a}\leq e^{32\sigma}$. This follows from
$Sw\leq\sigma$, $a\leq2\sigma$ and $3+S(7+w)+S+2a\leq32\sigma$.

Take $N=n_0$, $E=\exp(32\sqrt{\log n_0\log L})$, and $Q=n_0/(LE)$.
The bounds $\log n_0=O(\log x)$ and $\log L=O(\log\log x)$ give
$4C_oLE\leq x^\varepsilon$ eventually, hence $Q\geq4x^{-\varepsilon}n$.
Since $x^{-\varepsilon}n\geq1$ and $E\geq2$, deleting $0$ leaves at
least $Q-1\geq\max(x^{-\varepsilon}n,Q/E)$ points by
\eqref{eq:integer-sphere-bound}. This gives \eqref{eq:thick-size-bounds}
with $C_T=64$. The remaining $s$ are positive and $4Ls<n_0$,
so $4L|T_2|\leq n_0$.
\leanref{Formalism.\allowbreak{}FrobeniusOrbits.\allowbreak{}hostThickExists}{Formalism.\allowbreak{}FrobeniusOrbits.\allowbreak{}FrobeniusOrbitsProofs.\allowbreak{}ThickExists}{181--335}
\end{proof}

\subsection*{Proof of \ref{B2.54}}\label{prf:B2.54}\phantomsection
\begin{proof}
Give $\{0,1\}^A$ the product of the Bernoulli measures with parameters
$p_a=\delta_1$ on $A_0$ and $p_a=3\delta_1/4$ elsewhere. Let $T$ be the
set of coordinates equal to $1$. For each coset $C=\gamma+A_0$, put
$Y_C=|T\cap C|-\mathbb E|T\cap C|$ and $D=\sqrt{2\delta_1n}$. Then
\[
 \begin{gathered}
 \mu=\mathbb E|T|=\frac{3o_\phi+1}{4o_\phi}\delta_1n,\qquad
 \mathbb E|T\cap C|=\lambda_\gamma\mu,\\
 \mathbb E\sum_CY_C^2=\sum_ap_a(1-p_a)\leq D^2/2.
 \end{gathered}
\]
Condition on $t\in T$. The other indicators retain their product law
with parameters in $[3\delta_1/4,\delta_1]$.
For large $x$, $(M_0+1)\ell=O_\varepsilon(x^{\varepsilon/8}\log x)
\leq x^{\varepsilon/4}\leq N_{\rm orb}$, so a run contains $M_0+1$
distinct full tiles. A tile away from $t$ is marked with probability
at least
\[
 q_*=((1-\delta_1)^u)^{3R_\varepsilon+1}
 \left(\lfloor u/2\rfloor\frac{3\delta_1}{4}
                 (1-\delta_1)^{u-1}\right)^{3R_\varepsilon}.
\]
At least $M_0$ tiles in a run have disjoint coordinates away from $t$,
so the run fails with probability at most $(1-q_*)^{M_0}$.
A short-window failure needs at least $K_\varepsilon$ further selected
points and a long-window failure needs at least $N_1-1$. The union
bound over windows and runs gives conditional orbit failure probability
at most
\begin{equation}\label{eq:orbit-failure-probability}
 N_{\rm orb}\left[
 \binom{L_2}{K_\varepsilon}\delta_1^{K_\varepsilon}
 +(1-q_*)^{M_0}
 +\sum_{j=N_1-1}^{G'}\binom{G'}{j}\delta_1^j(1-\delta_1)^{G'-j}\right],
\end{equation}
where the sum is the probability of at least $N_1-1$ successes in $G'$ independent trials with success probability $\delta_1$. The two window bounds
use at most $L_2$ or $G'$ distinct coordinates, each with conditional
selection probability at most $\delta_1$ away from $t$.

Write $\delta=\delta_1\leq1/100$. Since $u=\lceil1/\delta\rceil$,
\[
 (1-\delta)^u\geq\exp\left(-\frac{1+\delta}{1-\delta}\right)>\frac3{10},
 \qquad \delta\lfloor u/2\rfloor\geq\frac{1-\delta}{2}.
\]
An occupied-cell factor is at least $(3/4)(99/200)(3/10)>1/10$.
Thus $q_*\geq10^{-(6R_\varepsilon+1)}$ and
$(1-q_*)^{M_0}\leq x^{-(H_\varepsilon+3)}$.
For $a=N_1-1$, a union bound over $a$-element subsets gives the binomial
tail bound $\binom{G'}a\delta_1^a\leq(\mathrm e G'\delta_1/a)^a$.
The definition of $N_1$ gives $x^{-(H_\varepsilon+3)}$ as an upper
bound for this term as well. Multiplying either bound by
$N_{\rm orb}\leq n\leq2x^{H_\varepsilon}$ gives $2x^{-3}$.

Put $C=\max(C_8,2)$. Since every prime at most $Y^{\varepsilon/5}$ satisfies $(-d/p)_K=1$, we have $p>Y^{\varepsilon/5}\geq x^{\varepsilon/(5C)}$ for inert-class
primes. For $x\geq16$, \eqref{eq:host-rank-ranges} gives
$\log_p((kx^{26})^4)\leq620C/\varepsilon$, so
$L_2\leq2(620C/\varepsilon+1)$. Since
$\delta_1^{K_\varepsilon}\leq x^{-(H_\varepsilon+2)}$, the first term
of \eqref{eq:orbit-failure-probability}, including $N_{\rm orb}$, is
$O_{\varepsilon,C_8}(x^{-2})$. The conditional probability is thus
at most $\beta=C_bx^{-2}$ for suitable $C_b$.

Let $B(T)$ be the left side of \eqref{eq:bad-top-budget} with $TT=T$,
and put $m_B=\beta\delta_1n\sum_{p\in\mathcal P_s}\log p/p$.
The conditional bound gives $\mathbb EB(T)\leq m_B$. If $m_B>0$,
then $\mathbb E(\sum_CY_C^2/D^2+B(T)/(4m_B))\leq3/4$.
Some $T$ has all $|Y_C|\leq D$ and $B(T)\leq4m_B$. If $m_B=0$,
then $B(T)=0$ almost surely and the second-moment bound gives the
same choice. For this choice,
\[
 \bigl||T|-\mu\bigr|\leq o_\phi D,\qquad
 \left||T\cap C|-\lambda_\gamma|T|\right|
 \leq (1+\lambda_\gamma o_\phi)D.
\]
By \eqref{eq:host-rank-ranges}, $\delta_1n\geq x^{2+63\varepsilon/8}$.
Since $o_\phi$ is bounded, these inequalities give the size bounds
and \eqref{eq:top-shares} for $TT=T$ at large $x$. The inequality $B(T)\leq4m_B$ gives \eqref{eq:bad-top-budget} for this top set.
\leanref{UpperTheory.\allowbreak{}BernoulliSumsFrobeniusOrbits.\allowbreak{}not\_\allowbreak{}good\_\allowbreak{}le}{Development.\allowbreak{}AutomorphicArithmetic.\allowbreak{}ThickFrames.\allowbreak{}NotGoodLe}{84--139}
\leanref{UpperTheory.\allowbreak{}BernoulliSumsFrobeniusOrbits.\allowbreak{}select}{Development.\allowbreak{}AutomorphicArithmetic.\allowbreak{}ThickFrames.\allowbreak{}select}{70--175}
\end{proof}

\subsection*{Proof of \ref{B2.55}}\label{prf:B2.55}\phantomsection
\begin{proof}
By \ref{B2.42}, the pair has the required prime ranges and
$-d\equiv1$ modulo $q$ and every odd prime $p\leq Y^{\varepsilon/5}$.
The congruence $d\equiv7\pmod{16}$ gives Kronecker value $1$ at $2$.
Put $z=Y^{\varepsilon/5}$. By \eqref{eq:host-rank-ranges},
$n\leq2x^{9/4}$ and $\log Q_S\leq2C_0+\frac94\log x$.
Using $\log x\leq C_0\log Y$, for large $x$ we have $\log Y\geq2$,
$z\geq820C_0\log Y$ and $\log Q_S\leq z((\log Y)^2-1)$.
Add $\log q\leq z$ and $\log d\leq z(\log Y)^2$ for the pair-row bound.

For the pointwise bounds use the twists of the trivial character in
\ref{B2.42} and $x^{1/2+\varepsilon/2}\leq x^{1/2+\varepsilon}$.
For the six quadratic twists use the base characters $1,\chi_8$ with
conductors $f=1,8$. The product of the primitive factors at the
coprime conductors is real and has conductor
$fq^id^j\leq8X_1(Y)X_2(Y)^j$. On units modulo $8qd$, quadratic
reciprocity identifies $(-q/\,\cdot\,)$ and $(-d/\,\cdot\,)$ with the
two sign characters. By \ref{B2.33}, every primitive representative
in either family agrees with
the corresponding primitive twist. In \ref{B2.42}, take the base
modulus $M$ for the bounds and $8M$ for the boxes. The respective
twist identities hold on units modulo $2Mqd$ and $8Mqd$. Restrict
the box to height $8X_1(Y)X_2(Y)^j$. The host data include
$q,d>4tr+1$.
\end{proof}

\subsection*{Proof of \ref{B2.56}}\label{prf:B2.56}\phantomsection
\begin{proof}
Take the thickness constant $C_{T_2}$ from Theorem~\ref{B2.53}, the exponent $C_1\geq1$ from Theorem~\ref{B2.51} (using Theorems~\ref{B2.49} and \ref{B2.50} for its progression-prime and short-interval norm-prime inputs), $c_6=1$ from Theorem~\ref{B2.52}, and $C_b>0$ from Theorem~\ref{B2.54}.
Set $C_8=\max(C_0,C_1)$, and take $x_0$ to be the maximum of the scale thresholds of Theorems~\ref{B2.51}--\ref{B2.54}, Lemma~\ref{B2.55}, the threshold for $4\log x\leq x^\varepsilon$, and the threshold for the prime-size lower bounds below.

Put $N=Q_Sqd$ and $z=Y^{\varepsilon/5}$.
By \eqref{eq:host-rank-ranges} and \ref{B6.61}, we have $r\leq x^5$ and $t\leq x^{400}$ for large $x$.
The lower bound $\log q>z/2$ from Definition~\ref{B2.42}, together with $z\geq x^{\varepsilon/(5C_0)}$ and $q\mid d+1$, yields $d\geq100$ and $q,d>4tr+1$ for large $x$.
To bound $\log N$ before choosing the final level primes, form an auxiliary choice record over $H_0$: take two primes greater than $Q_S$ and two larger primes $1\bmod15$ as the level set, set $\mathfrak c_K=\mathcal O_K$, $TT=T_2=\varnothing$, $L_w=2$, and $C_{T_2}=0$.
Admissibility of $R$ for $H_0$ at splitting bound $B_0=\lfloor z\rfloor$ yields an initial host $D_0$ extending $H_0$.
Applying Lemma~\ref{B2.55} to $D_0$ using the good-pair hypothesis for $(q,d)$ on $[Y,x]$ gives $\log N\leq2z(\log Y)^2\leq2x^{\varepsilon/5}(\log x)^2$.
Since $d\equiv7\pmod{16}$, we have $d\equiv3\pmod4$.
Theorem~\ref{B2.51} now supplies primes $g_0,g_1,p,p',q_0$ satisfying the size bounds, norm-prime conditions for $(d,N)$, $\mathcal B_r$-disjointness, congruences $p\equiv p'\equiv1\pmod{15N}$ and $q_0\equiv1\pmod N$, and upper bounds $p,p',q_0\leq N^{C_1}\leq N^{C_8}$.

Set $L_w=10\lceil405\log x\rceil+2\lceil30/\varepsilon\rceil+22$, so $L_w\leq(4084+60/\varepsilon)\log x$ for $\log x\geq1$.
Theorem~\ref{B2.53} with $C_o=\max(C_0,2)$ gives a thick set $T_2$ satisfying \eqref{eq:thick-size-bounds} and \eqref{eq:no-short-relation} at $L=L_w$.
By \eqref{eq:norm-one-coordinates} and $d\equiv3\pmod4$, the norm primes satisfy \eqref{eq:level-norm-and-congruence} and $g_i\equiv1\pmod N$.
All five primes $g_0,g_1,p,p',q_0$ are $1\bmod N$, hence coprime to $N$; they split completely in $\Q(\zeta_N)$, and therefore in $E$ and $K$.
Form the choice record with level set $\{g_0,g_1,p,p',q_0\}$, $\mathfrak c_K=\mathcal O_K$, $TT=\varnothing$, $T_2$, $L_w$, and $C_{T_2}$.
Admissibility of $R$ for $H_0$ at splitting bound $\lfloor z\rfloor\geq\lfloor x^{\varepsilon_z/5}\rfloor$ yields a host $D$ extending $H_0$.

Apply Lemma~\ref{B2.55} to $D$ using the good-pair hypothesis (Definition~\ref{B2.42}) on $[Y,x]$.
This supplies:
\begin{enumerate}
 \item the size, splitting, and modulus conditions at $Y$;
 \item the pointwise square-root bounds $|g_\chi(x)|\leq x^{1/2+\varepsilon}$ and $|(L'/L)(1,\chi)|\leq4\log x$ (constant $C_5=4$) for every primitive representative $\chi$ at level $8qd$ of the three sign characters $(\,\cdot\,/q)$, $(\,\cdot\,/d)$, and $(\,\cdot\,/q)(\,\cdot\,/d)$ (obtained from \eqref{eq:aux-pair-good-bounds} with the principal character $\xi=1$ modulo $M$);
 \item the reality, conductor, and zero-free box conditions in $\mathcal R(8X_1(Y)X_2(Y)^j)$ for the six quadratic twists (obtained from \eqref{eq:aux-pair-good-boxes}).
\end{enumerate}
Apply Theorem~\ref{B2.54} to $D$ to obtain a top set $TT\subset A$ satisfying the size bounds, the bad-orbit bound \eqref{eq:bad-top-budget} with constant $C_b$, and the coset share bound \eqref{eq:top-shares} with $\rho=x^{-1+\varepsilon}$.
Form $H$ by replacing the top set of $D$ with $TT$.

The divisibility $Q_S\mid\operatorname{lcm}(2^{e+2},M)$ implies $Q_{\mathrm{odd}}\mid M$ and that every prime divisor of $Q_S$ divides $2M$.
By the first part of Theorem~\ref{B2.52} applied to the good-pair boxes \eqref{eq:aux-pair-good-boxes} of Definition~\ref{B2.42}, every primitive real host character $\chi$ with $q\mid N_\chi$ or $d\mid N_\chi$ satisfies $N_\chi\leq H_{Q_{\mathrm{odd}}}(j)$ and $L(s,\chi)\neq0$ on $\mathcal R(H_M(j))$, hence on $\mathcal R(H_{Q_{\mathrm{odd}}}(j))$ and in particular has no real zero greater than $399/400$.
Since $Q_{\mathrm{odd}}\leq M\leq C_0$ and $\log N\leq2x^{\varepsilon/5}(\log x)^2$, the second part of Theorem~\ref{B2.52} gives the thin-position bound \eqref{eq:thin-character-bound} with $c_6=1$.
The logarithmic-derivative bound of Lemma~\ref{B2.55} for $\chi_d=(\,\cdot\,/d)$ gives $|\Re(L'/L)(1,\chi_d)|\leq4\log x\leq x^\varepsilon$ for $x\geq x_0$.
Since $\operatorname{lcm}(15,N)\mid15N$, the congruences $p,p'\equiv1\pmod{\operatorname{lcm}(15,N)}$ hold as well.
Thus $H$ and $q_0$ satisfy all completed-host conditions with constants $C_8,4,c_6,C_b$.
\end{proof}

\subsection*{Proof of \ref{B2.57}}\label{prf:B2.57}\phantomsection
\begin{proof}
Choose ranks on $\phi$ by \ref{B8.59}. Fix a choice record and $B_0$
as in admissibility. The field rows give the square root of $-q$
with its Galois action. The specified CM type consists of the
embeddings taking the chosen square root of $-d$ to
$i\sqrt d$. It has $2|A|>0$ embeddings, so choose one.
Since $K=\Q(\sqrt{-d})$ and an automorphism acts on $\sqrt{-d}$ by its $d$-sign, $\Gal(E/K)$ corresponds to the subgroup $\{((s_q,1),z):s_q\in\{\pm1\},\ z\in A\}$. The arithmetic datum supplies a
$\Gal(E/K)$-stable ideal $\mathfrak b$ with
\[
 \mathfrak D_{E/K}=\mathfrak b\,\overline{\mathfrak b},\qquad
 \mathfrak b+\overline{\mathfrak b}=\mathcal O_E.
\]
For the datum of \ref{B2.46}, take the ideal from the choice of one
prime in each conjugate pair of primes of $K$. Adjoin these choices and set the splitting
bound to $B_0$.
\leanref{HostParam.\allowbreak{}Seam.\allowbreak{}exists\_\allowbreak{}host}{Development.\allowbreak{}AutomorphicArithmetic.\allowbreak{}CoprimeRanks.\allowbreak{}SecondSetNoRelation}{100--105}
\end{proof}

\subsection*{Proof of \ref{B2.60}}\label{prf:B2.60}\phantomsection
\begin{proof}
Fix $0<\varepsilon<1/52$.
Take $c\geq1$ and $e_0$ from Theorem~\ref{B2.46}.
Apply Theorem~\ref{B2.48} with the cohomological window hypothesis of Definition~\ref{B6.03} and the unconditional good-pair existence theorem \ref{B2.43} to obtain $J\geq1$ and the family of windows $[Z,Z^{3/2}]$.
Choose $C_0\geq\max(1,J,o_\psi,8M)$.
From the completion hypothesis (discharged by Theorem~\ref{B2.56}) at $(\varepsilon,C_0)$ obtain $C_{T_2},C_8,c_6,C_b$ and threshold $x_{0,A}$.
Take the rank-supply threshold $x_{0,B}$ and the threshold $x_{0,T}$ for $t\leq x^{400}$ from \ref{B6.61}.
Increase $x_1\geq\max(x_{0,A},x_{0,B},x_{0,T})$ so that $x>c2^{e_0}$ and $Y^{\varepsilon/5}\geq\max(5,4M)$ whenever $x\geq x_1$ and $Y\geq x^{1/J}$.

Let $\mathcal W$ consist of the real $x$ admitting $Y,q,d,e$ as in \eqref{eq:window-datum} for which \eqref{eq:window-vv} holds for every field with the prescribed $\psi$-kernel at $(e,q,d)$.
By Theorem~\ref{B2.48}, $\mathcal W$ satisfies \eqref{eq:window-set-property} with $\lambda=3/2$.
Fix $x\geq x_1$ in $\mathcal W$ and its datum $(Y,q,d,e)$.
Since $(q,d)$ is a good pair for $(M,\varepsilon)$ on $[Y,x]$ in the sense of Definition~\ref{B2.42}, it satisfies \eqref{eq:host-prime-hyp}, $(M,qd)=1$, and $(-d/p)=1$ for every odd prime $p\leq Y^{\varepsilon/5}$.
The lower bound $\log q>Y^{\varepsilon/5}/2\geq5/2$ gives $q\neq3$, and $c2^e\geq x^{2+8\varepsilon}\geq x>c2^{e_0}$ gives $e\geq e_0$.

Apply Theorem~\ref{B2.46} with $B_0=\lfloor Y^{\varepsilon/5}\rfloor\geq4M$ to obtain a frame $\phi$ with $n=c2^e$ and $o_\phi=o_\psi$, and an arithmetic datum $H_0$ on $\phi$ with primes $q,d$, satisfying the field conditions at $B_0$, \eqref{eq:field-conductor-bounds}, $M\mid Q_S$, \eqref{eq:cyclic-artin-kernel}, the twist representation of Lemma~\ref{B2.45}, and the retained-character and $\theta_q$ clauses of Definition~\ref{B2.58}.
The rank-supply hypothesis gives ranks $R$ on $\phi$ satisfying the rank conditions at $s=(x,\varepsilon,J)$ and admissible for $H_0$.
Since $Q_S\leq4Mn\leq C_0n$ and $M\leq C_0$, the completion hypothesis at $(Y,C_0)$ produces a host $H$ extending $H_0$ and a prime $q_0$ satisfying the completed-host conditions with constants $C_8,4,c_6,C_b$ (including the pointwise bounds for $\chi_q,\chi_d,\chi_q\chi_d$ with $C_5=4$ from Lemma~\ref{B2.55} and the real-character zero-free boxes and thin-position bound from Theorem~\ref{B2.52}) together with the top-set shares \eqref{eq:top-shares} at $\rho=x^{-1+\varepsilon}$.

From $Q_S\mid\operatorname{lcm}(2^{e+2},M)$, every prime divisor of $Q_S$ divides $2M$.
Because $M\mid Q_S$, the twist representation of Lemma~\ref{B2.45} holds on $H$, and $(q,d)$ is a good pair for $(M,\varepsilon)$ on $[Y,x]$, Theorem~\ref{B2.47} gives $|g_\chi(x)|\leq x^{1/2+\varepsilon}$ and $|(L'/L)(1,\chi)|\leq4\log x$ for every $\theta\in\mathcal B_H(\mathcal X)$ and every primitive representative $\chi$ of $\theta$.
Together with the frame index, retained-character, $\theta_q$, and top-set share conditions, this establishes the arithmetic requirement of Definition~\ref{B2.58} for $\mathcal L_\psi$ with $C=4$.

Restriction from $\Q(\zeta_{Q_Sqd})$ to $E$ is the Artin map, so \eqref{eq:cyclic-artin-kernel} shows that $E$ has the prescribed $\psi$-kernel at $(e,q,d)$.
The window bound \eqref{eq:window-vv} therefore yields $V_E(\lfloor x^{400}\rfloor)\leq\sqrt b$.
Since $t\leq x^{400}$, restricting the nonnegative sum in \eqref{eq:vv-budget} gives $V_E(t)\leq\sqrt b$.
Taking $C_5=4$ and $x_0=x_1$ completes the proof of the first statement; applying Theorem~\ref{B2.56} discharges the completion hypothesis.
\end{proof}

\subsection*{Proof of \ref{B2.61}}\label{prf:B2.61}\phantomsection
\begin{proof}
Apply Theorem~\ref{B2.60} with the rank supply of Theorem~\ref{B2.57}.
In Theorem~\ref{B2.60}, the good pairs $(q,d)$ on $[Y,x]$ are supplied unconditionally by Theorem~\ref{B2.43} via Theorem~\ref{B2.48}, and the completion hypothesis is discharged by Theorem~\ref{B2.56}.
The resulting completed hosts at $C_5=4$ inherit the pointwise square-root bounds for $\chi_q$, $\chi_d$, and $\chi_q\chi_d$ from Lemma~\ref{B2.55}, the real-character zero-free boxes and thin-position bound from Theorem~\ref{B2.52}, and the square-root bounds for every character in $\mathcal B_H(\mathcal X)$ from Theorem~\ref{B2.47}.
\end{proof}

\section*{Proofs for Chapter~\ref{ch:p3c1}: The local kernel and the differential operator}

\subsection*{Proof of \ref{B3.05b}}\label{prf:B3.05b}\phantomsection
\begin{proof}
Leibniz's rule and the chain rule preserve integrality of derivatives at zero under products and zero-constant substitutions. The derivatives of $\log(1+T)$ are $(-1)^{n-1}(n-1)!$ at order $n\ge1$. Those of $\ell$ are integral since $(p^{2h})!$ is divisible by $p^h$. The same holds for $\ell^{-1}$ by differentiation of $\ell\circ\ell^{-1}=T$ and induction, since $\ell'(0)=1$. Thus for $A\simeq0$, the identity $A=\log(\exp(A^\ell))\circ\ell^{-1}$ gives integral derivatives.

Differentiating the first additivity defect in $X,X'$ and setting $X'=0$ gives $\partial_X^2b$. The other defect gives $\partial_Y^2b$. For the remaining mixed coefficient, use the conjugation defect with $\eta$ the Teichm\"uller lift of an element of $\F_{p^2}\setminus\F_p$. Then $(\eta-\sigma\eta)[XY]b\in W$ and $\eta-\sigma\eta$ is a unit.
\end{proof}

\subsection*{Proof of \ref{B3.06}}\label{prf:B3.06}\phantomsection
\begin{proof}
Put $F=\exp\Psi$ and $\varphi Q=(\sigma Q)(X^p)$. By Legendre's formula \cite[(1.7)]{krattenthaler-2015}, exponential and logarithm identify $p(X)W[[X]]$ with $1+p(X)W[[X]]$, where $(X)$ is the zero-constant ideal. The right side is therefore equivalent to $\varphi F/F^p\in1+p(X)W[[X]]$. Integral $F$ gives this by reduction modulo $p$. Conversely, suppose it holds and the terms of $F$ below degree $h\ge1$ are integral. Divide $F$ by their sum $T$ to obtain $H=F/T$. Since $\varphi T/T^p$ lies in the same group, so does $\varphi H/H^p$. Its degree-$h$ part is $-pH_h$, since $\varphi(H-1)$ starts in degree $ph$. Thus $F_h=H_h$ is integral.
\end{proof}

\subsection*{Proof of \ref{B3.08}}\label{prf:B3.08}\phantomsection
\begin{proof}
Write $P\equiv Q$ if $P-Q$ has zero constant term and $\exp(P-Q)$ is integral over $W$. Set $\beta=b^\ell$ and $F(X,T)=\ell^{-1}(\ell(X)+\ell(T))$. The relation $\ell(T)-\ell(T^{p^2})/p=T$ gives $F\in\Z_p[[X,T]]$ by Hazewinkel's functional equation lemma with $q=p^2$ and $s_1=p^{-1}$ \cite[Lemma 2.3, (2.5)]{hazewinkel-1979}. Let $\delta f=f(F(X,T))-f(X)-f(T)$, with $d_1,d_2$ the corresponding defects in each variable. The pairing laws give $d_i\beta\equiv0$ for $i=1,2$.

For a primitive $(p^2-1)$st root $\zeta\in W$, average $\beta(\zeta^k X,\zeta^{-pk}Y)$ over $0\le k<p^2-1$ to obtain $B$. Since $\ell(\zeta T)=\zeta\ell(T)$, the conjugation law makes each summand equivalent to $\beta$. The criterion \ref{B3.06} allows division by $p^2-1$ in this average and gives $B\equiv\beta$. The equivariance $F(\zeta X,\zeta T)=\zeta F(X,T)$ gives $d_iB\equiv0$. The average has support
\begin{equation}\label{eq:p3c1:phase-support}
 [X^iY^l]B\ne0\ \Longrightarrow\ i,l\ge1,\qquad i\equiv pl\pmod{p^2-1}.
\end{equation}
For integral zero-constant $P',D$ and a series $C$ with integral derivative,
\begin{equation}\label{eq:p3c1:taylor}
 \bigl(C(P'+pD)-C(P')\bigr)/p\quad\text{is integral}.
\end{equation}
Indeed, the $D^k$ term from $[T^n]C$ has scalar $n[T^n]C\binom{n-1}{k-1}p^{k-1}/k$ for $k\ge1$, and $v_p(k)\le k-1$.

Put $h(T)=\ell(T^p)/p$. Set $P=F(X,T)^p$ and $P'=F(X^p,T^p)$. Then $P-P'\in p\Z_p[[X,T]]$. Since $\ell'$ is integral, \eqref{eq:p3c1:taylor} gives $\delta h=(\ell(P)-\ell(P'))/p$ integral. Also $\delta\ell=0$. If $f(0)=0$ and $\delta f$ is integral, then
\begin{equation}\label{eq:p3c1:one-coboundary}
 f=a h+b_0\ell+R,\qquad a\in W,\quad b_0\in K_0,\quad R\in W[[T]].
\end{equation}
Subtract $[T]f$ times $\ell$. If the remainder is integral below degree $n\ge2$ and has coefficient $u$ there, its mixed defect terms give $\binom ni u\in W$ for $0<i<n$. The base-$p$ digit expansion of $(1+X)^n$ modulo $p$ gives a unit intermediate coefficient when $n$ is different from every power of $p$. For $n=p^{k+1}$, the identity $\binom{p^{k+1}}{p^k}=p\binom{p^{k+1}-1}{p^k-1}$ gives $pu\in W$ by the same expansion. At degree $p^{2j+1}$ subtract $p^{j+1}u h$, and at degree $p^{2j+2}$ subtract $p^{j+1}u\ell$. The earlier coefficients change by integral amounts. The increments belong to $p^jW$, so the multipliers converge and the limiting remainder is integral.

Differentiate $\exp(d_1B)$ in $T$ at zero, where its value is one. Using $\partial_2F(X,0)=1/\ell'(X)$ gives $\partial_XB-\ell'(X)[X]B$ integral. Also $\delta([X]B)$ is integral by the $X$-linear term of $\exp(d_2B)$. By \eqref{eq:p3c1:phase-support}, $[Y^{p^{2k}}][X]B=0$ for all $k$. Thus in \eqref{eq:p3c1:one-coboundary} for $[X]B$ one has $b_0/p^k\in W$ for all $k$, so $b_0=0$. Exchanging the variables gives $c,c'\in W$ with
\begin{equation}\label{eq:p3c1:phase-derivatives}
 \partial_XB-c\ell'(X)h(Y),\qquad
 \partial_YB-c'\ell'(Y)h(X)\quad\text{integral}.
\end{equation}
For any variable tuple $\mathbf X$, put $D_\sigma Q=((\sigma Q)(\mathbf X^p)-pQ)/p$ and set
\[
 J=D_\sigma B-\sigma(c)h(X)\ell(Y)-\sigma(c')\ell(X)h(Y).
\]
By \ref{B3.06}, $D_\sigma(d_1B)$ is integral. The commutator $d_1(D_\sigma B)-D_\sigma(d_1B)$ equals
$((\sigma B)(P,Y^p)-(\sigma B)(P',Y^p))/p$.
Apply \eqref{eq:p3c1:taylor} to $(\sigma B)(U,Y)-\sigma(c)\ell(U)h(Y)$ in $U$. Modulo integral series the commutator is $\sigma(c)\delta h\,h(Y^p)=\sigma(c)\delta h\,(\ell(Y)-Y)$. Subtracting the first defects of the two products in $J$ leaves an integral series since $\delta h$ is integral and $\delta\ell=0$. The second defect is integral by the exchanged calculation.

For $Q\in XYK_0[[X,Y]]$ with integral defects,
\begin{equation}\label{eq:p3c1:two-coboundaries}
 Q\in\gamma\ell(X)\ell(Y)+W[[X,Y]]\qquad(\gamma\in K_0).
\end{equation}
Apply \eqref{eq:p3c1:one-coboundary} to $[Y^m]Q$ and $[X^i]Q$. With $h_i,\ell_i$ the coefficients of $h,\ell$, we obtain
\[
 A_mh_i+B_m\ell_i-A'_ih_m-B'_i\ell_m\in W,\qquad A_m,A'_i\in W.
\]
Fix $i=p^{2k+1}$ and take $m=p^{2l}$. The first term has valuation at least $-k-1$, so $B'_i/p^l$ has the same bound for all $l$. Hence $B'_i=0$. Fix $m$ and vary $k$ to get $A_m=0$. Taking $i=p^{2k}$ gives $B_m=0$ when $\ell_m=0$. For $m=p^{2l}$, subtract $p^{-l}$ times the comparison at $m=1$. The valuation of $(B_m-p^{-l}B_1)/p^k$ is bounded below uniformly in $k$, so $B_m=p^{-l}B_1$. This proves \eqref{eq:p3c1:two-coboundaries}.

The congruence in \eqref{eq:p3c1:phase-support} is preserved by $D_\sigma$ and holds for the two products in $J$. Thus $[X^{p^{2k}}Y^{p^{2l}}]J=0$, and \eqref{eq:p3c1:two-coboundaries} gives $\gamma/p^{k+l}\in W$ for all $k,l$. Hence $J$ is integral. For $m=p^{2l+1}$, the coefficient of $XY^m$ in $B$ is congruent to $c/p^{l+1}$ by \eqref{eq:p3c1:phase-derivatives} and to $-\sigma(c')/p^{l+1}$ by integrality of $J$. Varying $l$ gives $c=-\sigma(c')$. Exchanging variables gives $c'=-\sigma(c)$ and $\sigma^2c=c$. Reindexing gives $A(X^p,Y^p)=p(h(X)\ell(Y)-A(Y,X))$. Consequently $D_\sigma(B-G_c)=J$ and $\beta\equiv G_c$ by \ref{B3.06}.

The coefficient of $T^{p^{2b}}$ in $\pi^{-1}\ell(\pi T)$ is $p^{\sum_{a<b}p^{2a}-b}$. The series lies in $T\Z[[T]]$ and reduces to $T+T^{p^2}$ modulo $p$. Its inverse $\mu$ therefore has integral coefficients and reduces to $\sum_{l\ge0}(-1)^lT^{p^{2l}}$. The identity $\ell(\pi\mu(T))=\pi T$ transports phase equivalences by substitution of $\pi\mu(L)$ in the logarithmic parameters. Also
\begin{equation}\label{eq:p3c1:scaled-phase}
 G_c(\pi X,\pi Y)=\frac{\pi^{p+1}}p(cXY^p-\sigma(c)X^pY)
 +\pi^{p+1}R(X,Y),\qquad R\in(X,Y)R_\pi[[X,Y]].
\end{equation}
For an $A$ term with $l\ge1$, the exponent of $p$ after factoring $\pi^{p+1}$ is $\sum_{a<i}(p^{2a}-1)+\sum_{a<l}(p^{2a+1}-1)-1\ge p-2$.

Reindexing the reduction of $\mu(L)$ gives
\[
 \mu(L^X_{r,a,b})\equiv\sum_{s\ {\rm even}}\sigma^s(\omega_r^*)S_{s,a,b},
 \qquad
 \mu(L^Y_{r,a,b})\equiv\sum_{s\ {\rm odd}}\sigma^s(\omega_r)S_{s,a,b}\pmod p .
\]
Trace duality makes the even conjugate matrices inverse transposes. Hence for $s$ even and $t$ odd,
$\sum_r\sigma^s(\omega_r^*)\sigma^{t+1}(\omega_r)=\mathbf1_{s=t+1}$,
with cyclic indices. Applying $\sigma$ gives the other parity. Frobenius on the $p$th powers in \eqref{eq:p3c1:scaled-phase} therefore gives $H_u$ modulo $p$ in the term with prefactor $\pi^{p+1}/p$. The resulting phase error belongs to $\pi^{p+1}(Z)R_\pi[[Z]]$. If $i!=p^a m_i$ with $p\nmid m_i$, Legendre's formula gives $(p-1)a\le i$, so $\pi^{(p+1)i}/i!=m_i^{-1}\pi^{(p+1)i-(p^2-1)a}\in R_\pi$. Exponentiating the error gives an integral unit, whose product with the units from phase equivalence gives $V$.
\end{proof}

\subsection*{Proof of \ref{B3.09}}\label{prf:B3.09}\phantomsection
\begin{proof}
If $V=\varnothing$, the generator phase in Definition~\ref{B3.05} is a constant equivalent to the empty sum. By Definition~\ref{B3.04} that constant is zero. The generator therefore equals its integral cofactor, so take $U_\varnothing=\Theta_A\in O_c$. Otherwise $p\ge5$. In the coordinate tuple $\mathcal L=(L^X,L^Y)$, the integral cofactor becomes $U(\pi\mathcal L)=U^\ell(\pi\mu(\mathcal L))$. Put $q=\pi^{p+1}$. At a single place \ref{B3.08} gives the series as $V_0\exp(-qH_u/p)$ with $V_0$ integral. Take $U_i=V_0(-q)^i\exp(qG_v)/i!$.

For the generator apply \ref{B3.08} to the pairing phase at each place. Put $\Delta=\sum_{v,e}b(X_{v,e},Y_{v,e})-\Phi_A$. The unit $E=\exp(\Delta^\ell)$ has coefficients in $W$. Since $\ell(\pi\mu(T))=\pi T$, the factor $\exp(\Delta(\pi\mathcal L))=E(\pi\mu(\mathcal L))$ belongs to $1+(Z)R'[[Z]]$. Absorb it into the integral cofactor $V_A$. The series is $V_A\prod_v\exp(-qH_{u_v}/p)$, so take $U_{\mathbf i}=V_A\prod_v((-q)^{i_v}\exp(qG_v)/i_v!)$. The factorial bound in \ref{B3.08} gives integrality, and the zero constant terms of $H'_v$ give the displayed degree bound.
\end{proof}

\subsection*{Proof of \ref{B3.10b}}\label{prf:B3.10b}\phantomsection
\begin{proof}
For every product \(Q\) of flag minors of the same heights as \(H\),
\([H,Q]=[H,H]Q(1)\). Indeed a nonzero pairing requires equality of the degrees in each row. Row one must then occur in every nonempty minor. After deleting it, row two must occur in every remaining nonempty minor. Continuing gives the initial row sets in every column. In that case the ordering signs give both \(Q(1)\) and \(Q/H\). The result extends by linearity to the span of these products. Cauchy--Binet shows that this span is stable under row substitutions.

First take \(z\) to be an \(s\times s\) matrix of independent variables. Reverse its columns to obtain \(R\) and work in the fraction field. For each \(h\), factor \(R\) out of \(C_h(Y,z)\). The last \(h\) columns left inside the determinant are the last \(h\) unit columns in reverse order. Thus, with \(\epsilon=\prod_k(-1)^{h_k(h_k-1)/2}\),
\[
 \prod_k\det C_{h_k}(Y,z)=\epsilon(\det R)^c H(R^{-1}Y).
\]
Substitution by \(R^{-1}\) in the first argument of the factorial pairing is adjoint to substitution by \(R^{-\mathsf t}\) in the second. This follows by comparing coefficients in \(\exp(\sum_{i,j}Y_{ij}Z_{ij})\). The normalized pairing is consequently \(\epsilon(\det R)^cQ_I(R^{-\mathsf t})\). Multiplying \([E_{I_k}\mid z_{:,1:h_k}]\) by \(R^{-1}\) and expanding along the last \(h_k\) columns gives precisely its \(k\)-th factor, with sign \((-1)^{h_k(h_k-1)/2}\). Both sides are polynomials in the original entries. The identity therefore survives every specialization. Unused columns of \(z\) may be filled arbitrarily.
\end{proof}

\subsection*{Proof of \ref{B3.10a}}\label{prf:B3.10a}\phantomsection
\begin{proof}
Choose basis matrices $A_v$ of $L_v$. Lift to $L'_v$ a basis of a complement to $L''_v/pL'_v$, then a basis of that subspace with lifts in $L''_v$. This gives a basis matrix $C_v$ of $L'_v$. Multiplying its first $i=i_v$ columns by $p$ gives a basis matrix $B_v$ of $L''_v$. Thus $C_v=B_vD_i$ with $D_i=\operatorname{diag}(p^{-1}I_i,I_{N_{j+1}-i})$. Take $\phi=(A_v,B_v)_v$.

In one block $(v,s)$ let $A,B,C$ denote the embedded matrices, and recall from Definition~\ref{B3.01a} the shift $d_l=d_{l,v,s}$, partition $\kappa=\kappa_{l,v,s}$, highest-weight minor product $H_\kappa$, unframed Weyl lattice $\Delta_O(\kappa)$, and tuple Weyl weight lattice $H_l(L)$ transported by $Y\mapsto E^{-1}Y$ with scalar $(\det E)^{-d_l}$. For a permitted arithmetic model of initial rank $k\ge5$, Proposition~\ref{B3.52b} bounds the index of the Weyl lattice in the tableau lattice used for the edge tests.
\leanref{Formalism.\allowbreak{}L2.\allowbreak{}PairingIndexBounds.\allowbreak{}Ent.\allowbreak{}delta\_\allowbreak{}zero\_\allowbreak{}eq\_\allowbreak{}twisted\_\allowbreak{}delta}{Formalism.\allowbreak{}LocalToGlobal.\allowbreak{}PairingIndexBounds.\allowbreak{}RowActionCoefficients.\allowbreak{}AdmissibleInterpolation.\allowbreak{}Bridge}{45--63}
\leanref{Formalism.\allowbreak{}L2.\allowbreak{}PairingIndexBounds.\allowbreak{}Ent.\allowbreak{}branching\_\allowbreak{}source\_\allowbreak{}identity}{Formalism.\allowbreak{}LocalToGlobal.\allowbreak{}PairingIndexBounds.\allowbreak{}RowActionCoefficients.\allowbreak{}AdmissibleInterpolation.\allowbreak{}BranchingSourceIdentity}{48--59}
\leanref{Formalism.\allowbreak{}L2.\allowbreak{}PairingIndexBounds.\allowbreak{}Ent.\allowbreak{}branching\_\allowbreak{}target\_\allowbreak{}identity}{Formalism.\allowbreak{}LocalToGlobal.\allowbreak{}PairingIndexBounds.\allowbreak{}RowActionCoefficients.\allowbreak{}AdmissibleInterpolation.\allowbreak{}BranchingSourceIdentity}{60--71}

Put $n=N_j$, $n'=N_{j+1}$, $w=w_{v,j+1}(s)$, and $c=\nu_{j+1,s,1}$. Write $h_k$ for the $c$ target column heights; they satisfy $h_k\le n$ since the final $t+1$ rows at level $j+1$ have length zero, and $\sum_{k=0}^{c-1}h_k=\sum_{a=1}^{n'}\nu_{j+1,s,a}=b$. For a square matrix $M$ let $C_h(M,x)$ take its first $\dim M-h$ columns and then the first $h$ columns of $x$. Write $x_{a,h}=\tau_{v,s,h,a}$ with matrix indices shifted to start at one, and let $x_0$ consist of the first $n$ columns. Put $\mu=(\nu_{j+1,s,a})_{a=1}^n$ and $e=\mu_n$. Use the source polynomials
\[
 S^+(x,Y)=(\det x_0)^e H_{\mu-e\mathbf1}(x_0^{\mathsf t}Y),\qquad
 S^-(x,Y)=\prod_k\det C_{h_k}(Y,x).
\]
The even source shape is $\mu-e\mathbf1$. The odd shape has heights $n-h_k$ and rows $c-\mu_{n+1-a}$. An elementary row substitution $Y\mapsto EY$ transfers to $x\mapsto E^{\mathsf t}x$ in $S^+$ and $x\mapsto E^{-1}x$ in $S^-$, since $\det E=1$. Thus the auxiliary coefficient spans are stable under divided powers. Evaluation at the identity columns for $S^+$ and the reversed identity columns for $S^-$ gives the corresponding highest-weight polynomial up to sign. Hence these coefficients span a module containing the Weyl generators. Under transport $Y\mapsto A^{-1}Y$, the free columns become $A^{-\mathsf t}x$ with scalar $(\det A)^e$ in the even case, and $Ax$ with scalar $(\det A)^{-c}$ in the odd case.

Compute the symbol $R_G(d)$ of Definition~\ref{B3.05aa} on $S^+$ and $S^-$, writing $Y$ now for the $n'\times n'$ target variable matrix in the block. At an even position, substituting $D=(dY)_{:,1:n}$ and multiplying by $(\det D)^e$ combines with $(\det x_0)^e$ to supply the $e$ height-$n$ minors, so $N_-=1$ in this block and $R_{S^+}(d)=H_{\nu_{j+1,s}}(z_0^{\mathsf t}Y)$ with $z=d^{\mathsf t}x$ and $z_0$ its first $n'$ columns.

At an odd position, with $T=(\operatorname{adj}Y)_{\mathrm{last}\ n\ \mathrm{rows}}d^{\mathsf t}$, the left inverse $\Pi_j$ in Definition~\ref{B3.05aa} evaluates the symbol of $S^-$ via the column identity
\[
 \det T\,(\det Y)^{n'-h}\det C_h(\operatorname{adj}T,x)
 =(\det T)^{n-h}(\det Y)^{n'-1}\det C_h(Y,d^{\mathsf t}x).
\]
Multiplying the first column matrix by $T$ and the second by $\operatorname{adj}Y$ gives initial columns of scalar identity matrices, and expansion along those columns leaves the same minor in the last $h$ rows and first $h$ columns of $(\operatorname{adj}Y)d^{\mathsf t}x$. Put $u=\det Y$, $v=\det T$, $P=\prod_k\det C_{h_k}(\operatorname{adj}T,x)$ and $Q=\prod_k\det C_{h_k}(Y,d^{\mathsf t}x)$. Multiplying the column identities over $k$ gives $v^cu^{cn'-b}P=v^{cn-b}u^{c(n'-1)}Q$. Since $u$ and $v$ are nonzero polynomials, in their localization this reads $Q=u^{c-b}v^{c-cn+b}P=(N_+/N_-)P$ in the block, so $N_+P=N_-Q$ in the polynomial ring and every linear left inverse $\Pi_j$ to multiplication by $N_-$ gives $Q$ on $N_+P$, also coefficientwise in the source parameters.
\leanref{Formalism.\allowbreak{}L2.\allowbreak{}G.\allowbreak{}Pend.\allowbreak{}polDel}{Formalism.\allowbreak{}LocalToGlobal.\allowbreak{}GroupLattices.\allowbreak{}PendingLemmas.\allowbreak{}DetYtMatNeZero}{100--102}
\leanref{Formalism.\allowbreak{}L2.\allowbreak{}G.\allowbreak{}Pend.\allowbreak{}col\_\allowbreak{}id}{Formalism.\allowbreak{}LocalToGlobal.\allowbreak{}GroupLattices.\allowbreak{}PendingLemmas.\allowbreak{}CardFilterValLt}{73--101}
\leanref{Formalism.\allowbreak{}L2.\allowbreak{}G.\allowbreak{}Pend.\allowbreak{}cols\_\allowbreak{}id}{Formalism.\allowbreak{}LocalToGlobal.\allowbreak{}GroupLattices.\allowbreak{}PendingLemmas.\allowbreak{}CardFilterValLt}{102--113}

In Definition~\ref{B3.05aa}(3), the displayed symbols $H_{\nu_{j+1,s}}(z_0^{\mathsf t}Y)$ and $Q=\prod_k\det C_{h_k}(Y,d^{\mathsf t}x)$ already have homogeneous $d$-degree $\sum_k h_{v,s,k}=\sum_{a=1}^{n'}\nu_{j+1,s,a}$ in each block $(v,s)$, so on the source generators and their linear span the degree filter in $\mathcal D_j(G,f)$ is automatic. The arithmetic substitution $B_\phi$ on $Z$ is $Z\mapsto A^{\mathsf t}ZB^{-\mathsf t}$ in an even block and $Z\mapsto A^{-1}ZB$ in an odd block; comparing coefficients in $\sum_{a,a'}d_{a,a'}Z_{a,a'}$ gives the transposed action $d\mapsto AdB^{-1}$ and $d\mapsto A^{-\mathsf t}dB^{\mathsf t}$, respectively, preserving each block degree.

For a target partition $\kappa$, use the pairing $[P,Q]=\sum_\alpha\alpha![Y^\alpha]P[Y^\alpha]Q$ and put
\[
 C_\kappa=\sum_\alpha\alpha!\bigl([Y^\alpha]H_\kappa\bigr)^2\in\Z_{>0}.
\]
Use the image of this integer in the coefficient field. Row substitution by $E$ is adjoint to substitution by $E^{\mathsf t}$. For a product $Q$ of initial-column minors with selected rows and shape $\kappa$, $[H_\kappa,Q]/C_\kappa=Q(I)$. Indeed, equality of row weights forces row $1$ to appear in every column, row $2$ in every column of height at least two, and so on. These tests equal $H_\kappa$ with the row-order sign. Multilinearity gives the same identity for row substitutions in $Q$.

At an even block the normalized pairing of $H_\kappa(z_0^{\mathsf t}Y)$ with $Q$ is therefore $Q(z_0)$. Apply Lemma~\ref{B3.10b} with an independent matrix of free columns before specializing it to $d^{\mathsf t}x$. At an odd block $\kappa$ has rows $c-\nu_{j+1,s,n'+1-a}$. Reverse the first $n'$ columns of $z$ to form $R$. Factoring $R$ from each $C_{h_k}(Y,z)$ and applying adjointness gives $(\det R)^c Q(R^{-\mathsf t})$ with the column reversal signs. If $I_k$ lists the rows of the $k$th test minor, write $E_{I_k}$ for its unit columns. The minor of $R^{-\mathsf t}$ on rows $I_k$ and the first $n'-h_k$ columns, multiplied by $\det R$, is the determinant with first columns $E_{I_k}$ and last $h_k$ columns of $R$. This follows by multiplication by $R^{-1}$ and expansion along the last columns. Undoing the reversal gives the determinant with columns $E_{I_k}$ followed by the first $h_k$ columns of $z$. Polynomiality extends the formula to singular $R$.

For the framed test $Q(C^{\mathsf t}Y)$, move $C$ to the output by adjointness. After the source transport and the transposed frame action, the free columns become
\[
 C^{\mathsf t}(AdB^{-1})^{\mathsf t}A^{-\mathsf t}x=D_i d^{\mathsf t}x,\qquad
 C^{-1}(A^{-\mathsf t}dB^{\mathsf t})^{\mathsf t}Ax=D_i^{-1}d^{\mathsf t}x.
\]
At an odd block the substitution $Y\mapsto CY$ also contributes $(\det C)^c$. Repeated rows give zero tests. For the other tests let $J_k$ be the selected rows in an even block and their complementary rows in an odd block. At an even position retain the order of the rows in the test minor. At an odd position take the complementary rows in increasing order and include the sign from expansion along the unit columns. Write $q_k$ for the number of rows in $J_k$ among the first $i$ and $M=\sum_k\min(i,h_k)$. The minor $D_k$ of $d^{\mathsf t}x$ on rows $J_k$ and its first $h_k$ columns is the transpose of the corresponding column determinant in \ref{B3.10}. Since $q_k\le\min(i,h_k)$, the normalized pairings before weight twists are
\[
 (\det A)^e p^{-M}\prod_k p^{\min(i,h_k)-q_k}D_k
 \quad\text{and}\quad
 \pm(\det A)^{-c}(\det C)^c\prod_k p^{q_k}D_k
\]
at even and odd positions, respectively.
\leanref{Formalism.\allowbreak{}L2.\allowbreak{}PairingIndexBounds.\allowbreak{}Ent.\allowbreak{}MiddleFrameLemmas.\allowbreak{}blkX}{Formalism.\allowbreak{}LocalToGlobal.\allowbreak{}PairingIndexBounds.\allowbreak{}RowActionCoefficients.\allowbreak{}AdmissibleInterpolation.\allowbreak{}ProdItePow}{372--415}

Let $\widetilde y$ be the asserted output on the universal transported source product, with its target twist $\prod_{v,s}(\det C)^{-d_1}$ divided out. Denote a product of framed tests by $Q_I^C$. First evaluate on $f_d=\exp(\sum_z d_zZ_z)$, where $z$ indexes the coordinates. Contraction with $f_d$ evaluates the symbol at $d$. Multiply the block identities, including the source twists $(\det A)^{-d_0}$. The factorial and target pairing both factor over disjoint blocks. This gives
\begin{equation}\label{eq:p3c1:frame-pairing}
 \frac{[\widetilde y,Q_I^C]}{\prod_{v,s}C_\kappa}
 =\pm U\,\mathcal D_{P_{J,\mathbf i}}(f),\qquad
 U=\left(\prod_v p^{N'_{v,j+1}(i_v)}\right)D(L,L'')
 \prod_{v,s}(\det A)^{-d_0}(\det C)^{d_1}U_s,
\end{equation}
for $f=f_d$, where $J$ is the column choice given by the $J_k$, and $U_s=(\det A)^ep^{-M}$ or $(\det A)^{-c}(\det C)^c$ according to parity. For general $f=\sum_\alpha f_\alpha Z^\alpha$, apply the functional $d^\alpha\mapsto\alpha! f_\alpha$ on the finite block degrees of the identity, coefficientwise in $\tau$. Indeed, for any matrix $T$, comparison of the two expansions of $\exp(d^{\mathsf t}TZ)$ gives $\alpha![Z^\alpha]f(TZ)=\sum_\gamma\gamma! f_\gamma[d^\gamma](T^{\mathsf t}d)^\alpha$. Thus this functional moves the transposed action back to the stated substitution on the series, with every factorial retained as in \ref{B3.11}. On the left $f_d(B_\phi Z)=\exp(\sum_z(B_\phi^{\mathsf t}d)_zZ_z)$. The block degrees are preserved, so ordinary contraction recovers $\mathcal D_j(G,f(B_\phi Z))$ on the given generators. On the right it gives the indicated test. The symbolic outputs lie in the $K_c$-span of products of the target Weyl spaces by their highest-weight formulas for invertible free-column matrices. The formulas are polynomial, so the outputs at singular matrices and the symbol coefficients also lie in that span. Hence contraction preserves it.
\leanref{Formalism.\allowbreak{}L2.\allowbreak{}PairingIndexBounds.\allowbreak{}Ent.\allowbreak{}xMid\_\allowbreak{}of}{Formalism.\allowbreak{}LocalToGlobal.\allowbreak{}PairingIndexBounds.\allowbreak{}RowActionCoefficients.\allowbreak{}AdmissibleInterpolation.\allowbreak{}ProdItePow}{504--545}

Put $a=\operatorname{ord}_p\det A_v$ and $a'=\operatorname{ord}_p\det B_v$. Write $\lambda_s^{[i]}=\sum_{k=1}^i\lambda_{v,s,k}$. The embeddings preserve the determinant valuations since $p$ maps to $p$ and units and their inverses map to integral elements. Thus $C$ has determinant valuation $a'-i$, and $D(L,L'')$ contributes $w(a'-a)$ at an even position and its negative at an odd position. At an even position $d_1=-w$ and $e=d_0+w$. Since $i\le r\le N_0$, counting boxes in the first $i$ rows gives $M=\lambda_s^{[i]}+iw$. At an odd position negative reversal gives $d_1=d_0$ and $c=w-d_0$. With the terms $\lambda_s^{[i]}$ and $iw$ of $N'_{v,j+1}(i)$ included, the contributions to $v_p(U)$ are
\begin{equation}\label{eq:p3c1:frame-unit}
 \begin{aligned}
 \lambda_s^{[i]}+w(a'-a)-d_0a+d_1(a'-i)+ea-M&=0,\\
 iw-w(a'-a)-d_0a+d_1(a'-i)-ca+c(a'-i)&=0.
 \end{aligned}
\end{equation}
\leanref{Formalism.\allowbreak{}L2.\allowbreak{}PairingIndexBounds.\allowbreak{}Ent.\allowbreak{}xUnit\_\allowbreak{}of}{Formalism.\allowbreak{}LocalToGlobal.\allowbreak{}PairingIndexBounds.\allowbreak{}RowActionCoefficients.\allowbreak{}AdmissibleInterpolation.\allowbreak{}SumIccEqRange}{233--309}

The standard bideterminants of the Schur lattice and its duality with the Weyl lattice are the straight-shape cases of \cite[Theorem II.2.16 and Proposition II.4.1]{akin-1982}. For a divided-power word $u$,
$[uH_\kappa,Q]/C_\kappa=(u^{\mathsf t}Q)(I)$,
where transposition reverses the word and transposes the substitutions. This is integral for integral $Q$. Conversely, restrict $Q$ to lower unitriangular matrices. Its column degrees determine the diagonal exponents, so this restriction reads the lower triangular monomial coefficients. Produce the lower unitriangular matrix from $I$ by row additions in decreasing order of source row index. Each source row is still a unit row at the time of addition. Thus every coefficient of the restriction is evaluation of a divided-power word at $I$.

Order exponents lexicographically by column and then by row. For a minor on increasing rows $i_1<\cdots<i_h$, the leading monomial is $\prod_{r=1}^hY_{i_r,r}$ with coefficient one. Its product over the columns of a standard tableau determines the tableau because each tableau row is ordered. These leading monomials are lower triangular. Successive subtraction of the leading terms therefore gives an integral Schur expansion whenever the normalized pairings are integral. The same argument with zero pairings gives nondegeneracy on the common characteristic-zero span. Hence the field dual of an integral Weyl basis for $[-,-]/C_\kappa$ consists of integral combinations of the tests.

Transport the Weyl basis by $Y\mapsto C^{-1}Y$ and the dual basis by $Y\mapsto C^{\mathsf t}Y$. Their products are dual under the product normalization. By \eqref{eq:p3c1:frame-pairing} and \eqref{eq:p3c1:frame-unit}, all pairings with framed tests have valuation at least $g$ coefficientwise in $\tau$. The source coefficient spans and the integral embedding of $O$ give the same bound for every $G\in H_j(L)$. Each basis coordinate of $\widetilde y$ has this bound by integral duality. Restoring the target twists puts the basis vectors in $H_{j+1}(L')$ and gives the stated sum.
\leanref{Formalism.\allowbreak{}L2.\allowbreak{}PairingIndexBounds.\allowbreak{}Ent.\allowbreak{}xPair}{Formalism.\allowbreak{}LocalToGlobal.\allowbreak{}PairingIndexBounds.\allowbreak{}RowActionCoefficients.\allowbreak{}AdmissibleInterpolation.\allowbreak{}NuTAntitone}{36--100}
\leanref{Formalism.\allowbreak{}L2.\allowbreak{}PairingIndexBounds.\allowbreak{}Ent.\allowbreak{}dual\_\allowbreak{}bases\_\allowbreak{}holds}{Formalism.\allowbreak{}LocalToGlobal.\allowbreak{}PairingIndexBounds.\allowbreak{}RowActionCoefficients.\allowbreak{}AdmissibleInterpolation.\allowbreak{}PolyKTorsionFree}{44--125}
\leanref{Formalism.\allowbreak{}L2.\allowbreak{}PairingIndexBounds.\allowbreak{}Ent.\allowbreak{}tgt\_\allowbreak{}bidet}{Formalism.\allowbreak{}LocalToGlobal.\allowbreak{}PairingIndexBounds.\allowbreak{}RowActionCoefficients.\allowbreak{}AdmissibleInterpolation.\allowbreak{}TgtBidet}{33--75}
\end{proof}

\section*{Proofs for Chapter~\ref{ch:b3-entry-estimates}: Entry estimates: factorials against denominators}

\subsection*{Proof of \ref{B3.15}}\label{prf:B3.15}\phantomsection
\begin{proof}
For a monomial \(z^m\), the contributions from multiplying by \(z_x\)
on the series side and differentiating on the polynomial side both
vanish when \(m_x=0\). When \(m_x>0\), their equality is
\[
 m!\,[z^{m-\mathbf1_x}]F
 =m_x(m-\mathbf1_x)!\,[z^{m-\mathbf1_x}]F.
\]
By linearity, multiplication by \(L_k\) is adjoint to
\(\partial_{b_k}=\sum_x b_{k,x}\partial/\partial z_x\).
Expanding the monomials of \(P(z+\sum_kt_kb_k)\) gives
\[
 \prod_k\partial_{b_k}^{\beta_k}P
   =\left(\prod_k\beta_k!\right)D_b^{[\beta]}P.
\]
The multinomial coefficients are integers, so this identity holds in
every characteristic. Apply the adjoint relation repeatedly to
obtain \eqref{eq:b3-hasse-pair}.

The expansion takes place over \(\mathcal O_K\) when \(P\) and the
directions are integral. For homogeneous \(P\), every term has total
\((z,t)\)-degree \(d\). If \(|\beta|\le d\), every monomial in the
coefficient of \(t^\beta\) has \(z\)-degree \(d-|\beta|\).
If \(|\beta|>d\), the coefficient is zero.
\end{proof}

\subsection*{Proof of \ref{B3.16}}\label{prf:B3.16}\phantomsection
\begin{proof}
If \(I_c\) repeats a label, \(\Delta_c=0\). For injective \(I_c\), a
direction at \(s_c\) with label \(b\in\operatorname{im}I_c\) adds
\(t_kM_c\xi_k\) to the unique column labeled \(b\). Every other
direction leaves the matrix unchanged. The determinant has degree at
most one in the parameters attached to any single column. If two
directions have the same row vector, a term using both replacement
columns vanishes by alternation.

The Hasse product formula is
\begin{equation}\label{eq:b3-hasse-product}
 D_b^{[\beta]}\prod_c\Delta_c
   =\sum_{\beta^{(1)}+\cdots+\beta^{(C)}=\beta}
        \prod_cD_b^{[\beta^{(c)}]}\Delta_c.
\end{equation}
In a nonzero summand, directions from \(\mathcal K_0\) go to factors
at \(s_0\). With a common row vector, each such factor receives at most
one direction. With labels in \(J\), it receives at most one per label
in \(\operatorname{im}I_c\cap J\). Summing the respective bounds over
these factors proves \eqref{eq:b3-hits-cap} and
\eqref{eq:b3-hits-columns}.
\end{proof}

\subsection*{Proof of \ref{B3.17}}\label{prf:B3.17}\phantomsection
\begin{proof}
If \(E=0\), then \(D=0\). Both families are empty and \(h\) has
zero constant term, so \(Z_{\ge p}=Y=0\) suffices.
Assume \(E\ge1\) and replace each \(S_{s,a,b}\) by its sum over
\(0\le l<E\). The omitted terms have degree at least \(p^{2E}\).
The linear substitution \eqref{eq:b3-top-low-sub} preserves this
degree bound, so the coefficients of \(h\) through degree \(D\)
agree with those of the finite substitution.

Work modulo \(p\mathcal O_K\), where Frobenius is additive. Since
\(p\) is odd, the \(p\)-th power preserves the signs \((-1)^l\).
For fixed \(v,a,b\), the substituted product
\(S_{v,a,b}^pS_{v+1,a,b}\) becomes \((Z_1+T_1)(Z_2+T_2)\), where
\begin{align*}
 Z_1&=\sum_{l<E}(-1)^ly_{v-2l,a,b}^{p^{2l+1}},\\
 T_1&=\sum_{l<E}(-1)^l
       \bigl(\mathbf1_{v-2l\in T}\xi_{v-2l,a}t_{v-2l,b}
                    \bigr)^{p^{2l+1}},\\
 Z_2&=\sum_{l'<E}(-1)^{l'}y_{v+1-2l',a,b}^{p^{2l'}},\\
 T_2&=\sum_{l'<E}(-1)^{l'}
       \bigl(\mathbf1_{v+1-2l'\in T}
                   \xi_{v+1-2l',a}t_{v+1-2l',b}
                    \bigr)^{p^{2l'}}.
\end{align*}
Write \(Z_2=y_{v+1,a,b}+Z_2'\), where every monomial of \(Z_2'\)
has remaining degree at least \(p^2\). Then
\begin{equation}\label{eq:b3-sort-four}
 (Z_1+T_1)(Z_2+T_2)
   =T_1T_2+T_1y_{v+1,a,b}
      +Z_1(Z_2+T_2)+T_1Z_2'.
\end{equation}

Multiply by \(u_v\) and sum over \(v,a,b\). Expanding \(T_1T_2\)
gives the first sum of \eqref{eq:b3-four-parts}, with the exponents
and coefficients in \eqref{eq:b3-da} and \eqref{eq:b3-aa}.
In the second term of \eqref{eq:b3-sort-four}, change variables
from \(v\) to \(s=v-2l\). Summing over \(a\) gives
\[
 (-1)^lu_{s+2l}\,
 t_{s,b}^{p^{2l+1}}
 \sum_a\xi_{s,a}^{p^{2l+1}}y_{s+2l+1,a,b},
\]
when \(s\in T\), and zero otherwise. Summing over \(s\in S\),
\(0\le l<E\) and \(0\le b<r_1\) gives the sum indexed by \(\kappa=(s,l,b)\).

Define \(Z_{\ge p}\) by summing
\(u_v[Z_1(Z_2+T_2)+T_1Z_2']\) over \(v,a,b\). Its monomials have
\(y\)-degree at least \(p\) by the bounds for \(Z_1\) and \(Z_2'\).
Subtract the two indexed sums and \(Z_{\ge p}\) from the finite
substitution and divide by \(p\). The congruence modulo
\(p\mathcal O_K\) gives an integral quotient \(Y\), which completes
\eqref{eq:b3-four-parts}.
\end{proof}

\subsection*{Proof of \ref{B3.18}}\label{prf:B3.18}\phantomsection
\begin{proof}
The polynomial \(t^eQ_0\) is homogeneous of degree \(D\).
Thus \eqref{eq:b3-class-recall} expresses its pairing with \(g\) as
a finite sum of pairings with \(p^{-i}h^iU_i\), \(0\le i\le D\).
In these pairings we may use the displayed expression for \(h\),
because multiplication by power series preserves agreement through
degree \(D\).

Expand \(h^i\) by the multinomial formula. A summand chooses
\(m_\phi\) occurrences of each pure-top term,
\(\beta_\kappa\) occurrences of each term with a linear remaining
factor, \(i_2\) occurrences of \(Z_{\ge p}\), and \(i_3\) occurrences
of \(pY\). Put
\[
 A=\sum_\phi m_\phi,\qquad B=\sum_\kappa\beta_\kappa,
 \qquad i=A+B+i_2+i_3.
\]
Its multinomial coefficient and the exponent of its extracted top
monomial are
\[
 M=\frac{i!}
        {(\prod_\phi m_\phi!)(\prod_\kappa\beta_\kappa!)i_2!i_3!},
 \qquad
 q=\sum_\phi m_\phi d^A_\phi+\sum_\kappa\beta_\kappa d^B_\kappa.
\]
Use \ref{B3.15} to move the product of linear factors to the polynomial
side. This contributes \(\prod_\kappa\beta_\kappa!\) and replaces
\(Q_0\) by \(D_b^{[\beta]}Q_0\), with \(t^e\) unchanged. The resulting
pairing is a sum over \(n\in\operatorname{supp}D_b^{[\beta]}Q_0\)
with factorial factor \(e!n!\). A nonzero coefficient requires \(q\le e\).

Apart from integral scalar factors and a coefficient of the integral
Hasse derivative, such a summand has the form
\begin{equation}\label{eq:b3-term-shape}
 p^{-i+i_3}M
 \left(\prod_\kappa\beta_\kappa!\right)e!n!\,
 [t^{e-q}y^n](Z_{\ge p}^{\,i_2}Y^{i_3}U_i).
\end{equation}
If it is nonzero, every scalar factor selected with positive
multiplicity is nonzero. The minimum remaining degree of
\(Z_{\ge p}^{\,i_2}\) gives \(pi_2\le|n|\).

The multinomial product is divisible by the factorial of the total
number \(B\) of linear factors. Indeed,
\[
 M\prod_\kappa\beta_\kappa!
   =B!\,
      \frac{i!}{(\prod_\phi m_\phi!)B!i_2!i_3!},
\]
where the quotient is an integer multinomial coefficient. Hence
\(\left|M\prod_\kappa\beta_\kappa!\right|\le|B!|\).
The coefficient in \eqref{eq:b3-term-shape} has norm at most one
and \(|p^{-i+i_3}|=p^{A+B+i_2}\). For each nonzero summand \(w\)
this gives \(|w|\le p^{A+B+i_2}|e!|\,|B!|\,|n!|\) with its own data.
Choose a summand of maximal norm. The ultrametric inequality gives
\(|J|\le |w|\), so its data satisfy \eqref{eq:b3-term-factorials}.
\end{proof}

\subsection*{Proof of \ref{B3.20}}\label{prf:B3.20}\phantomsection
\begin{proof}
Put
\[
 c_p=\frac1{p-1},\qquad
 A=\sum_s\alpha_s,\qquad L=\sum_sL_s,\qquad
 I=\sum_{s\in S}\sum_{l\in E}i_{s,l}.
\]
By \ref{B3.19}, monotonicity of \(\ell_p\), and
\eqref{eq:b3-star-local},
\begin{align}
 \sum_bv_p(e_{s,b}!)
   &\ge c_p\alpha_s-r_1\ell_p(\alpha_s),
       \label{eq:b3-star-top}\\
 \sum_{a,b}v_p(a_{s,a,b}!)
   &\ge c_p\sum_{a,b}a_{s,a,b}
             -r_0r_1\ell_p(L_s).
       \label{eq:b3-star-rem}
\end{align}
For the middle family,
\begin{equation}\label{eq:b3-star-logsum}
 \sum_{l\in E}\ell_p(i_{s,l})\le\ell_p(\alpha_s)^2.
\end{equation}
For \(\alpha_s=0\), all the \(i_{s,l}\) vanish. A positive \(i_{s,l}\)
implies \(p^{2l+1}\le\alpha_s\), hence \(l<\ell_p(\alpha_s)\).
There are at most \(\ell_p(\alpha_s)\) such indices, and each
\(i_{s,l}\le\alpha_s\). Another application of \ref{B3.19} gives
\[
 \sum_{l\in E}v_p(i_{s,l}!)
   \ge c_p\sum_{l\in E}i_{s,l}-\ell_p(\alpha_s)^2.
\]
Outside \(T\), this inequality has zero on both sides. Summing the
three bounds over \(s\) and using \eqref{eq:b3-star-total} gives
\begin{equation}\label{eq:b3-three-factorial-sum}
 \mathcal F\ge c_p(A+L)-\delta.
\end{equation}
Since \(c_p-\rho_p=\gamma_p\), it follows that
\[
 \mathcal F-\rho_p(A+L)-(i_0+I+i_2)
 \ge \gamma_p A-i_0-I-\delta+\gamma_p L-i_2.
\]
By \eqref{eq:b3-star-total} and \(1/p\le\gamma_p\), we have
\(i_2\le L/p\le\gamma_p L\).
\end{proof}

\subsection*{Proof of \ref{B3.21}}\label{prf:B3.21}\phantomsection
\begin{proof}
Transposing the linear substitution \eqref{eq:b3-top-low-sub}
in the differential pairing gives
\begin{equation}\label{eq:b3-entry-transpose}
 \mathcal J
  =\left\langle t^eQ_0(y),
       F\bigl(y+\mathbf1_T\xi t\bigr)\right\rangle.
\end{equation}
A derivative in \(t_{s,b}\) is carried by the chain rule to
\(\sum_a\xi_{s,a}\partial/\partial z_{s,a,b}\) when \(s\in T\).
Top exponents outside \(T\) vanish. A derivative in \(y_{s,a,b}\)
becomes \(\partial/\partial z_{s,a,b}\). Iterating and evaluating
at the origin gives \eqref{eq:b3-entry-transpose}.

The substituted \(F\) belongs to \(\mathcal C_p(h)\).
Choose \(E\) with \(D<p^{2E}\) and apply \ref{B3.17}.
Each determinant contributes its height to the degree of \(Q_0\),
so \(Q_0\) is homogeneous of degree \(\sum_sL_s\).
Apply \ref{B3.18} to \eqref{eq:b3-entry-transpose}.

For a nonzero entry, the resulting \(m,\beta,i_2,n\) satisfy
\eqref{eq:b3-entry-support} by the nonzero scalar conditions.
The other conclusions give \eqref{eq:b3-entry-budget},
\eqref{eq:b3-entry-hasse-nonzero}, and
\begin{equation}\label{eq:b3-entry-factorial-value}
 v_p(\mathcal J)\ge
 v_p(e!)+v_p(B!)+v_p(n!)
       -\left(\sum_\phi m_\phi+B+i_2\right),
 \qquad B=\sum_\kappa\beta_\kappa.
\end{equation}
Every integer has norm at most one. For an integer \(u\) with
\(p\nmid u\), a Bézout identity \(au+bp=1\) gives \(|u|=1\) since
\(|p|<1\). Writing each factorial factor \(m\) in
\eqref{eq:b3-entry-factorial-value} as \(p^ku\) with \(p\nmid u\)
gives \(v_p(m)=k=\operatorname{ord}_p(m)\), where \(\operatorname{ord}_p\)
denotes the arithmetic valuation.

For each source \(s\), set \(i_{s,l}=\sum_{b<r_1}\beta_{(s,l,b)}\)
for \(l<E\) and \(i_{s,l}=0\) for \(l\ge E\). The top budget gives
\(p^{2l+1}i_{s,l}\le\alpha_s\). Since \(\sum_{s,l}i_{s,l}=B\),
\[
 \prod_{s,l}i_{s,l}!\ \mid\ B!,
 \qquad
 v_p(B!)\ge\sum_{s,l}v_p(i_{s,l}!).
\]
A monomial \(y^n\) of \(D_b^{[\beta]}Q_0\) comes from a monomial
\(y^{n_0}\) of \(Q_0\) with \(n\le n_0\) and \(|n|+B=|n_0|\).
Every monomial of \(Q_0\) has degree \(L_s\) in the variables at \(s\),
so
\[
 \sum_{a,b}n_{s,a,b}\le L_s,\qquad
 |n|+B=\sum_sL_s.
\]
Also \(pi_2\le|n|\le\sum_sL_s\). Apply \ref{B3.20} with
\(a_{s,a,b}=n_{s,a,b}\), \(i_0=\sum_\phi m_\phi\), and the grouped
array \(i_{s,l}\). Its conclusion and
\eqref{eq:b3-entry-factorial-value} give \eqref{eq:b3-entry-star}.
\end{proof}

\subsection*{Proof of \ref{B3.22}}\label{prf:B3.22}\phantomsection
\begin{proof}
The positive integer \(a-a'\) satisfies \(s+(a-a')=s'\), so
minimality gives \(d\le a-a'\le a\).

For \(s\in T\), write \(x_l=\alpha_sp^{-(2l+1)}\), with \(d=2l+1\).
At an excluded distance, \eqref{eq:b3-charge-excluded} gives
\(0\le x_l-\operatorname{ch}_{s,l}-i_{s,l}\). At an unexcluded
distance, \eqref{eq:b3-charge-support} and \eqref{eq:b3-charge-free} give
\[
 [x_l-C_{s+d}]_+
 \le x_l-i_{s,l}
 =x_l-\operatorname{ch}_{s,l}-i_{s,l}.
\]
The series of positive parts is bounded termwise by the convergent
series \(\sum_lx_l=\gamma_p\alpha_s\). The other two arrays have finite
support. Summation therefore gives
\[
 \sum_{\substack{d\ge1,\ d\ {\rm odd}\\d\notin D_s}}
      [\alpha_sp^{-d}-C_{s+d}]_+
 \le\gamma_p\alpha_s
      -\sum_l\operatorname{ch}_{s,l}-\sum_li_{s,l}.
\]
Sum over \(s\in T\) and subtract \(\delta\).
\end{proof}

\subsection*{Proof of \ref{B3.23}}\label{prf:B3.23}\phantomsection
\begin{proof}
The bound holds for \(\mathcal J=0\). For a nonzero entry take the
multiplicities of \ref{B3.21}. Set \(i_{s,l}=\sum_b\beta_{(s,l,b)}\)
for \(l<E\) and \(i_{s,l}=0\) for \(l\ge E\), as in that proof.
For fixed \(s,l\) with \(l<E\), the directions share the target
\(s+2l+1\) and the row vector \((\xi_{s,a}^{p^{2l+1}})_a\).
By \eqref{eq:b3-entry-hasse-nonzero} and \eqref{eq:b3-hits-cap},
\(i_{s,l}\le C_{s+2l+1}\). The top budget gives
\(p^{2l+1}i_{s,l}\le\alpha_s\), so for every \(l\ge0\)
\begin{equation}\label{eq:b3-arrival-source-min}
 i_{s,l}\le\min(\alpha_sp^{-(2l+1)},C_{s+2l+1}).
\end{equation}

Assign each occurrence of a pure-top term to an excluded distance.
For \(\phi=(v,l,l',a,b)\), put
\[
 s_1=v-2l,\qquad s_2=v+1-2l',\qquad
 a_1=2l+1,\qquad a_2=2l'.
\]
Then \(s_1+a_1=s_2+a_2\). The endpoints have opposite parity, and
\(a_1,a_2\) are unequal. Choose as source the endpoint with larger
\(a_i\), and as target the other endpoint. Let \(d_\phi\) be the
first positive distance from the source to the target.
By \ref{B3.22}, \(d_\phi\le\max(a_1,a_2)\). The distance is odd.
If \(m_\phi>0\), both endpoints belong to \(T\), so
\(d_\phi\in D_{\operatorname{source}(\phi)}\). The pure-top monomial
consumes at least \(p^{d_\phi}\) units of top degree at its chosen
source in column \(b\).

For \(d=2l+1\), let
\[
 \operatorname{ch}_{s,l}
  =\sum_{\substack{\phi:\ \operatorname{source}(\phi)=s\\
                             d_\phi=2l+1}}m_\phi.
\]
These counts have finite support, and
\[
 \sum_{s\in T,l}\operatorname{ch}_{s,l}=\sum_\phi m_\phi,
 \qquad
 \operatorname{ch}_{s,l}=0\quad\text{if }2l+1\notin D_s.
\]
For fixed \(s,l\), sum the coordinatewise top budget
\eqref{eq:b3-entry-budget} over the columns at \(s\).
The charged \(A\)-terms and the \(B\)-terms of distance \(2l+1\)
consume at least \(p^{2l+1}(\operatorname{ch}_{s,l}+i_{s,l})\)
units of top degree. Consequently
\begin{equation}\label{eq:b3-combined-charge-budget}
 \operatorname{ch}_{s,l}+i_{s,l}
       \le\alpha_sp^{-(2l+1)}.
\end{equation}
Apply the scalar part of \ref{B3.22} with
\[
 V=
 \gamma_p\sum_{s\in T}\alpha_s
       -\sum_\phi m_\phi-\sum_\kappa\beta_\kappa-\delta.
\]
By \eqref{eq:b3-combined-charge-budget} and
\eqref{eq:b3-arrival-source-min}, this gives \(V\ge g_A\).
The common estimate \eqref{eq:b3-entry-star} proves the theorem.
\end{proof}

\subsection*{Proof of \ref{B3.24}}\label{prf:B3.24}\phantomsection
\begin{proof}
For a nonzero entry take the data of \ref{B3.21}, with \(i_{s,l}\)
extended by zero as above, and put
\[
 A=\sum_\phi m_\phi,\qquad
 B_{\mathrm{ev}}=\sum_{s\in T_{\mathrm{ev}}}\sum_{l\ge0}i_{s,l},
 \qquad
 B_{\mathrm{odd}}=\sum_{s\in T_{\mathrm{odd}}}\sum_{l\ge0}i_{s,l}.
\]
Each pure-top term consumes at least one unit of top degree at its
even endpoint. Each even-source linear-factor term also consumes at
least one unit at an even position. Summing the top budget
\eqref{eq:b3-entry-budget} over the even coordinates gives
\begin{equation}\label{eq:b3-even-payment}
 A+B_{\mathrm{ev}}\le A_{\mathrm{ev}}.
\end{equation}
For odd sources, \eqref{eq:b3-arrival-source-min} gives
\[
 B_{\mathrm{odd}}
 \le\sum_{s\in T_{\mathrm{odd}}}
       \sum_{l\ge0}\min(\alpha_sp^{-(2l+1)},C_{s+2l+1}).
\]
Substitute both inequalities in the common lower bound
\[
 v_p(\mathcal J)\ge
 \gamma_p\sum_{s\in T}\alpha_s
       -A-B_{\mathrm{ev}}-B_{\mathrm{odd}}-\delta+\rho_pD.
\]
This gives \eqref{eq:b3-entry-a2}. The norm bound also holds for
a zero entry.
\end{proof}

\subsection*{Proof of \ref{B3.25}}\label{prf:B3.25}\phantomsection
\begin{proof}
For a nonzero entry take the multiplicities of \ref{B3.21}.
Let \(\mathcal P\) consist of the even-source coordinates \((s,b)\)
with \(b\notin J\) and the odd-source coordinates with \(b\in J\).
The endpoints of a pure-top term have opposite parity and share a
column label, so exactly one belongs to \(\mathcal P\). Each occurrence
consumes at least one unit of top degree there, as does a \(B\)-term
whose source coordinate belongs to \(\mathcal P\). Summing the top
budget over \(\mathcal P\) gives
\begin{equation}\label{eq:b3-top-column-payment}
 \sum_\phi m_\phi+
 \sum_{\kappa:\ (\operatorname{source}(\kappa),
                         \operatorname{col}(\kappa))\in\mathcal P}
                     \beta_\kappa
 \le\sum_{(s,b)\in\mathcal P}e_{s,b}.
\end{equation}

The even-source \(B\)-terms left outside this payment have labels
in \(J\) and odd targets. Group their directions by target.
By the nonvanishing in \eqref{eq:b3-entry-hasse-nonzero},
\eqref{eq:b3-hits-columns} applies at each such target with label set \(J\). Summing the bounds gives
\[
 \sum_{\substack{\kappa:\ \operatorname{source}(\kappa)\ {\rm even}\\
                              \operatorname{col}(\kappa)\in J}}
       \beta_\kappa
 \le
 \sum_{c:\ s_c\ {\rm odd}}\#(\operatorname{im}I_c\cap J).
\]
Together with \eqref{eq:b3-top-column-payment}, this gives
\begin{equation}\label{eq:b3-full-column-payment}
 \sum_\phi m_\phi+
 \sum_{\kappa:\ \operatorname{source}(\kappa)\ {\rm even}}
       \beta_\kappa
 \le N_J.
\end{equation}

For odd-source \(B\)-terms use \eqref{eq:b3-arrival-source-min}.
Add \(N_J\) to the common estimate \eqref{eq:b3-entry-star} and use
\eqref{eq:b3-full-column-payment} to obtain
\[
 N_J+v_p(\mathcal J)\ge
 \gamma_p\sum_{s\in T}\alpha_s
 -\sum_{s\in T_{\mathrm{odd}}}
       \sum_{l\ge0}\min(\alpha_sp^{-(2l+1)},C_{s+2l+1})
 -\delta+\rho_pD.
\]
This is \(v_p(p^{N_J}\mathcal J)\ge g_B+\rho_pD\). The norm
bound also holds for a zero entry.
\end{proof}

\subsection*{Proof of \ref{B3.26}}\label{prf:B3.26}\phantomsection
\begin{proof}
Let \(F\in\mathcal C_p(H)\). For \(A(\tau)\in K[\tau]\), put
\[
 \|A\|_{\mathrm G}=\max_\eta|[\tau^\eta]A|,
 \qquad \|0\|_{\mathrm G}=0.
\]
To prove multiplicativity, divide each of two nonzero polynomials
by a coefficient of maximal norm. Each normalized polynomial is
integral and has nonzero reduction modulo the maximal ideal of
\(\mathcal O_K\). Their reduced product is nonzero because the residue
polynomial ring is a domain. The normalized product therefore has
Gauss norm one. Rescaling proves multiplicativity.

The Gauss norm extends to the rational function field \(K(\tau)\).
Write \(Y_{s,i,a}\) for its coordinate functions. For every
\(A\in K[\tau]\), evaluation at \(Y\) gives
\begin{equation}\label{eq:b3-gauss-generic-eval}
 |A(Y)|=\|A\|_{\mathrm G}.
\end{equation}
Taking \(A\) constant gives the isometric inclusion of \(K\).
Taking \(A=\tau_{s,i,a}\) gives \(|Y_{s,i,a}|=1\).

Evaluate the frame variables in \(\mathcal P\) at \(Y\).
Separate the height-one factors at marked positions. Each has
the form \(\sum_{a<r_0}Y_{s,0,a}z_{s,a,b}\). Choose
\(\xi_{s,a}=Y_{s,0,a}\) and let \(e_{s,b}\) count the separated
factors at \(s\) with label \(b\). If \(r_0=0\), the vectors are
empty and there are no height-one factors. Leave the remaining
determinants in \(Q_0\). Thus \(\mathcal P(Y,z)=P_{\mathrm{ent}}(z)\)
of \eqref{eq:b3-entry-polynomial}, with integral frame and direction
coordinates and counts \eqref{eq:b3-gauss-alpha} and
\eqref{eq:b3-gauss-counts}. The isometric extension preserves the
expansion-class identities for \(F\).

Evaluation in the finite frame-jet sum gives
\(\mathcal J_F(Y)=\langle P_{\mathrm{ent}},F\rangle\).
By \eqref{eq:b3-gauss-generic-eval},
\begin{equation}\label{eq:b3-coeff-to-entry}
 |[\tau^\eta]\mathcal J_F|
 \le\|\mathcal J_F\|_{\mathrm G}
 =|\langle P_{\mathrm{ent}},F\rangle|.
\end{equation}
Apply \ref{B3.23} and \ref{B3.24} to this entry to get
\begin{align}
 |[\tau^\eta]\mathcal J_F|
       &\le p^{-(g_A+\rho_pD)},\label{eq:b3-scaled-a1}\\
 |[\tau^\eta]\mathcal J_F|
       &\le p^{-(g_B-A_{\mathrm{ev}}+\rho_pD)}.
          \label{eq:b3-scaled-a2}
\end{align}
\leanref{Formalism.\allowbreak{}InvariantFactorPaths.\allowbreak{}G7.\allowbreak{}EGaussK.\allowbreak{}reduce}{Formalism.\allowbreak{}ThetaLogCoordinates.\allowbreak{}UniformThetaSubstitution.\allowbreak{}FrobeniusLiftNorms.\allowbreak{}EGaussK}{38--101}

For the prescribed-column bound, apply \ref{B3.25} with \(J=J_i\).
A separated height-one factor contributes \(1\) to \(N_{J_i}\) if
its position is even and its label lies outside \(J_i\), or if its
position is odd and its label lies in \(J_i\). Since \(i\ge1\), these
are exactly its contributions to \(M_i\). Each remaining determinant
at an odd position contributes \(q_{s,c}(i)\) to both. At an even
position, a remaining determinant adds
\(\min(h_{s,c},i)-q_{s,c}(i)\) to \(M_i\). This is nonnegative
because \(q_{s,c}(i)\le h_{s,c}\) and \(q_{s,c}(i)\le |J_i|\le i\).
Hence
\begin{equation}\label{eq:b3-edge-payment-comparison}
 N_{J_i}\le M_i.
\end{equation}
Since \(|p|\le1\), \ref{B3.25} and \eqref{eq:b3-coeff-to-entry} give
\begin{equation}\label{eq:b3-scaled-b}
 |p^{M_i}[\tau^\eta]\mathcal J_F|
 \le p^{-(g_B+\rho_pD)}.
\end{equation}

Take \(F(z)=\Theta(\pi z)\). The product \(\mathcal P(\tau,z)\)
is homogeneous of \(z\)-degree \(D\), since each determinant
contributes its height. Substitution multiplies \([z^n]\Theta\)
by \(\pi^{|n|}\), and the frame pairing uses terms of degree \(D\).
Together with \(\pi^{p^2-1}=p\) this gives
\begin{align}
 \mathcal J_{\Theta(\pi z)}&=\pi^D\mathcal J_\Theta,
       \label{eq:b3-frame-scaling}\\
 |\pi^D|&=p^{-\rho_pD}.\label{eq:b3-pi-degree}
\end{align}
Dividing each scaled bound by \(|\pi^D|=p^{-\rho_pD}\) proves
the three asserted estimates.
\leanref{Formalism.\allowbreak{}InvariantFactorPaths.\allowbreak{}G7.\allowbreak{}EGaussK.\allowbreak{}descend}{Formalism.\allowbreak{}ThetaLogCoordinates.\allowbreak{}UniformThetaSubstitution.\allowbreak{}FrobeniusLiftNorms.\allowbreak{}EGaussK}{132--144}
\end{proof}

\section*{Proofs for Chapter~\ref{ch:b3-one-lift}: One lift on a lattice and the Hecke commutation}

\subsection*{Proof of \ref{B3.27}}\label{prf:B3.27}\phantomsection
\begin{proof}
Write \(\nu_{s,a}=\nu_{j+1,s,a}\) for the target row lengths of Definition~\ref{B3.02}. The column heights \(h_{s,c}=\#\{a:\nu_{s,a}>c\}\) are at most \(N_j\), since the last \(t+1\) rows vanish. At a marked position delete the \(\alpha_{v,s}=\nu_{s,1}-\nu_{s,2}\) height-one columns. The reduced diagram has \(\nu_{s,2}\) columns there and \(\nu_{s,1}\) elsewhere, with size \(L_s=\sum_a\nu_{s,a}-\alpha_{v,s}\). This is the diagram of Theorem~\ref{B3.26}. The contraction sets variables outside \(v\) to zero, and the injection of local auxiliary variables into the full set identifies the local coefficients and gives zero off their support.

Here \(W_v\subseteq O\) gives \(e_O>0\). Complete \(K_c\) for \(|z|=p^{-v(z)}\). Equip the algebraic closure of the completion with the extending spectral norm and choose \(\pi^{p^2-1}=p\) there. For the integral linear forms \(L\) in the local substitution, the inverse \(\mu\) of \(\pi^{-1}\ell(\pi T)\) from Theorem~\ref{B3.08} gives \(U(\pi L)=U^\ell(\pi\mu(L))\), an integral series. The scaled phase is \(\exp(-\pi^{p+1}H_u/p)V\), where \(H_u=\sum_{s,a,b}u_sS_{s,a,b}^{p}S_{s+1,a,b}\) is the formal series of Definition~\ref{B3.07}. The series \(V\) and the scalars \((-\pi^{p+1})^d/d!\) are integral. Thus the scaled product is the sum of \(p^{-d}H_u^d\) times the integral series \((-\pi^{p+1})^d VU(\pi L)/d!\). Theorem~\ref{B3.26} gives the norm bound \(p^{-A_v}\) after contraction in the original variables.
\end{proof}

\subsection*{Proof of \ref{B3.29}}\label{prf:B3.29}\phantomsection
\begin{proof}
Definition~\ref{B3.05a} gives \(\Theta=\exp(-\Phi)U_0\). For \(\Delta=\Phi-\sum_v\Phi_v\), phase equivalence makes \(\exp(\Delta^\ell)\) an integral unit of constant term one, so take \(U_1=\exp(-\Delta)U_0\). If there are no positions, the contraction is the integral constant coefficient. Otherwise expand \(U_1^\ell=\sum_\beta a_\beta X^\beta\) and substitute the inverse \(\mathfrak e=\ell^{-1}\) coordinatewise. Each \(a_\beta\in O_c\), and the corresponding term of \(U_1\) is \(a_\beta\prod_v\mathfrak e(Y_v)^{\beta_v}\). Substitution and contraction factor over positions in this product. Only \(|\beta|\le\deg_Z P_{I,\mathbf i}\) can contribute, since \(\mathfrak e(0)=0\). The factor \(\mathfrak e(Y_v)^{\beta_v}\) becomes \(Y_v^{\beta_v}\) after logarithmic substitution. By \eqref{eq:p3c3-choice-exponents} and Theorem~\ref{B3.28}, the contracted factor at \(v\) has coefficient bound \(q_v\). Multiplication adds the bounds and finite addition preserves them.
\end{proof}

\subsection*{Proof of \ref{B3.30}}\label{prf:B3.30}\phantomsection
\begin{proof}
For \(f\in V_j\) and a target lattice pair \(L'\), Definition~\ref{B3.05a} gives over \(K_c\)
\[
 (K_jf)(L')=\sum_{L\in\mathcal R_j}
 \frac{c_K}{|\Gamma_{j,L}|}\operatorname{ker}_{j,L,L'}(f(L))
 \qquad(j+1<m),
 \tag{3.30.4}\label{eq:p3c3-class-expansion}
\]
where \(\mathcal R_j\) contains one representative of each lattice orbit in each component. At the terminal step it gives
\[
 (K_jf)(L')=\sum_x\frac{c_K}{w_x}\sum_y
 \operatorname{ker}_{j,L_{x,y},L'}(f(L_{x,y})).
 \tag{3.30.5}\label{eq:p3c3-terminal-expansion}
\]
Both sums are finite. The terminal index sets and their associated weights and source pairs are fixed for all \(f,L'\).

For each source \(L\) and each simultaneous lower neighbour \(L''\) of \(L'\), Definition~\ref{B3.05a} supplies a frame from \(L\) to \(L''\) in which the bound of Lemma~\ref{B3.10a} holds for every series and every real lower bound. Apply it to the framed theta series with source value \(f(L)\in H_j(L)\), where \(f\in\Lambda_j\). The value \(\operatorname{ker}_{j,L,L''}(f(L))\), multiplied by \(\prod_v p^{N'_{v,j+1}(i_v)}\), is a sum of elements of \(H_{j+1}(L')\) with coefficients of bound \(g\). For each class denominator \(w\), \(p^{d_G(j)}/w\in O_c\). Consequently the value at \(L'\) of the normalized class sum is
\[
 y=\sum_\alpha a_\alpha z_\alpha,\qquad
 z_\alpha\in H_{j+1}(L'),\qquad
 \nu_c(a_\alpha)\ge h=g-d_G(j).
 \tag{3.30.7}\label{eq:p3c3-class-coefficients}
\]
In a nonzero \(K\)-linear relation \(\sum b_\alpha z_\alpha=0\), choose \(b_\gamma\ne0\) of least \(\varpi\)-order. Eliminating \(z_\gamma\) preserves the coefficient bound because the ratios \(b_\alpha/b_\gamma\) lie in \(O_c\cap K=\iota(O)\). Repeat until the vectors are independent. Since \(y\) is defined over \(K\), an invertible coordinate minor gives \(a_\alpha\in K\). For \(e_O>0\), \(\operatorname{ord}_\varpi(a_\alpha)\ge\lceil e_Oh\rceil\) by \eqref{eq:p3c3-coeff-valuation}. For \(e_O=0\) the same condition gives \(a_\alpha\in\iota(O)\). Thus \(y\in\iota(\varpi)^{\lceil e_Oh\rceil}H_{j+1}(L')\).

The lift sends \(V_j\) into \(V_{j+1}\) by admissibility. A level element gives an index-preserving bijection between the lower neighbours of a lattice and those of its translate. An upper endpoint of an edge with full lower endpoint is full. Thus the edge sums preserve equivariance and support, and the fibre bounds give \eqref{eq:p3c3-class-bound}.
\end{proof}

\subsection*{Proof of \ref{B3.32}}\label{prf:B3.32}\phantomsection
\begin{proof}
Take \(i_v=0,\ q_v=A_v\) by Theorem~\ref{B3.27} and use Lemma~\ref{B3.31} with the empty list. In Definition~\ref{B3.13}, for large odd distances \(d\) one has \(\alpha_{v,s}p^{-d}<1\). The tail in each class modulo \(n'_v\) then vanishes if the capacity is positive and is geometric if the capacity is zero. Hence \(A_v\in\Q\), so Theorem~\ref{B3.29} gives the coefficient bound. For \(\mathbf i=\mathbf0\), Definition~\ref{B3.05a} gives the coefficient bound \(g=0\) at every lattice pair, frame and column choice. Lemma~\ref{B3.30} then gives \eqref{eq:p3c3-integral-kernel}.
\end{proof}

\subsection*{Proof of \ref{B3.34}}\label{prf:B3.34}\phantomsection
\begin{proof}
For the lift bound take the list \([v]\) in Lemma~\ref{B3.31}, with \(i_v=i,\ q_v=B_v\) and \(i_{v'}=0,\ q_{v'}=A_{v'}\) elsewhere. For integrality choose an adapted basis on an index-\(i\) edge \(A\to_i B\) at level \(b\), so
\[
 A_v=gW_v^N,\qquad B_v=gDW_v^N,\qquad
 D=\operatorname{diag}(p^{-1}I_i,I_{N-i}),\qquad N=N_b.
 \tag{3.34.3}
\]
At one embedding let \(\kappa^*\) be the actual weight, \(e=\kappa_N^*\), and \(\nu_a=\kappa_a^*-e\).  The untwisted integral Weyl module is spanned by divided-power words on the highest-weight product of minors of shape \(\nu\). A divided power of order \(n\) in direction \(E_{ab}\) takes \([z^n]\) after \(Y\mapsto(I+zE_{ab})Y\). It shifts the total degree in the first \(i\) rows by \(n(\mathbf1_{b\le i}-\mathbf1_{a\le i})\), so each nonzero word is homogeneous in those rows. Row substitution expands each minor into minors of the same height. A height-\(h\) minor uses at most \(\min(i,h)\) of the first \(i\) rows. Thus each nonzero word \(F\) has degree \(d\le M_i=\sum_{a=1}^i\nu_a\) in the first \(i\) rows, so \(p^{M_i}F(DY)=p^{M_i-d}F(Y)\).

Let \(G\) be the image of \(g\) at this embedding. A fibre element at \(A\) has the form \(F=(\det G)^{-e}F_0\), where the polynomial \(H(Y)=F_0(GY)\) lies in the untwisted module. The preceding degree bound puts \(p^{M_i}H(DY)=p^{M_i}F_0(GDY)\) in the untwisted module, so \(p^{M_i}F_0\) satisfies the row condition at \(B\). Its determinant twist is \((\det(GD))^{-e}p^{M_i}F_0=p^{M_i+ie}F\), where \(M_i+ie=\sum_{a\le i}\kappa_a^*\). Tensoring over the embeddings gives \(p^{\mathcal A_i(\kappa)}H_b(A)\subseteq H_b(B)\). Sum over the edges, with equivariance and support as in Lemma~\ref{B3.30}. For \(b=j+1,\ i\le r\), the even terms use the initial entries and every odd row entry is at least \(-w_{v,j+1}(s)\). Thus
\[
 \mathcal A_i(\kappa)\le
 \sum_{s\ {\rm even}}\sum_{a=1}^i\lambda_{v,s,a}
 +i\sum_{s\ {\rm odd}}w_{v,j+1}(s)=N'_{v,j+1}(i).
 \tag{3.34.5}
\]
This proves integrality for the stated normalization.
\end{proof}

\subsection*{Proof of \ref{B3.35a}}\label{prf:B3.35a}\phantomsection
\begin{proof}
In an abstract system the fourth item of Definition~\ref{B3.05a} gives the units and the additive maps with the stated quantifiers. For a constructed arithmetic system take the signed determinant lengths and the record of the two components and all entries outside \(v\). Lemma~\ref{B3.35a-arithmetic} gives \eqref{eq:p3c3-indicator} for these data. Its formula \eqref{eq:p3c3-two-units} chooses the units before the fibre, and \eqref{eq:p3c3-chart-action} identifies the chart value with the value obtained from a representative.
\end{proof}

\subsection*{Proof of \ref{B3.35d}}\label{prf:B3.35d}\phantomsection
\begin{proof}
For any three full lattices of the same rank,
\(\mathfrak d(U,W)=\mathfrak d(U,V)\mathfrak d(V,W)\), by taking determinants over each local ring. Hermitian duality sends the determinant ideal to its inverse conjugate. Thus the relative determinant of self-dual lattices has norm one, and its character value is nonzero.

The trace dual of \(\mathfrak b^{-1}\) for the Hermitian trace pairing is
\(\mathfrak D_{E/K_0}^{-1}\bar{\mathfrak b}=\mathfrak b^{-1}\). In dual bases of \(A,B\) this calculation applies to each matrix entry. On the curve, division by the area makes an oriented period basis have alternating pairing one. Hence \(P(A,B)\) is self-dual. On \(Q\) the operators \(W_y\) obey the additive group law. If \(h\in P\cap Q\), then \(W_{y+h}F=W_yW_hF=W_yF\) for a \(P\)-fixed section.

Use the complex move identities \eqref{eq:B4.22c-source} and \eqref{eq:B4.22c-target}. Their hypotheses are self-duality and even ranks, as here. Their proof uses the normalized averages on complex period lattices before making the split-prime moment estimate. To see explicitly how the choice of ideal character enters, denote the Weil family's character by \(\chi_*\), its sections by \(e^*_{A,B}\) and its determinant factor by \(C_*(A,B)\). On the standard pair the theta space is a line and the Weil section is nonzero. Write the chosen standard section as \(a_0\) times that section. Transport first the standard source to \(A\), then the standard target to \(B\). Linearity of both averages gives
\[
 e_{A,B}=a_0\,\frac{C(A,B)}{C_*(A,B)}e^*_{A,B}.
\]
The quotient makes sense even if \(a_0=0\). In a source move \(A\to A_1\) the ratio of the two \(C\)'s formed with \(\chi\) is \(\chi(\mathfrak d(A,A_1))^{n'/2}\). In a target move \(B\to B_1\) it is \(\chi(\mathfrak d(B,B_1))^{-n/2}\). Dividing the corresponding move identities by the new \(C\) therefore gives \(I_{P,Q}\) in both cases. Finally, a rational unitary change contributes a principal ideal generated by its determinant. On such generators both characters have value \(\prod_{\tau\in\Psi}\tau(\det g)^{-1}\). Thus the ratio \(C/C_*\) is invariant under the changes at both ends. Applying the Weil family's pullback identity \eqref{eq:B4.22c-equiv} gives the asserted formula for \(e\).
\end{proof}

\subsection*{Proof of \ref{B3.35e}}\label{prf:B3.35e}\phantomsection
\begin{proof}
An element of the determinant class group of the construction has a finite idèle representative. Strip off the finitely many uniformizer powers of that representative. The remaining integral idèle is in the kernel of the class character. At each stripped-off uniformizer the character is the unramified value of the defining Hecke character, which belongs to \(H\) by the choice of coefficient field. Their product with the respective integer exponents therefore belongs to \(H^\times\). In particular the class character to power \(b_l\) at each chosen representative belongs to \(H^\times\).

For \(l=0\) the power is zero. For \(l>0\) transport \(\operatorname{diag}(x,1,\ldots,1)\) from the standard Hermitian form to the level-\(l\) form. This gives a finite unitary point \(Y\) of determinant \(x\). In the chosen component write \(Y=\gamma y_C k\) with \(\gamma\) rational and \(k\) in the level. The determinant twist of this chart at \(Y\) is \(\chi_I(x)^{b_l}\). Its value on \(k\) is one. Its value on \(\gamma\) is the inverse of the \(b_l\)-th power of the type norm of \(\det\gamma\), whose factors are among the specified embeddings into \(H\). Multiplication by the class value at \(y_C\) proves the assertion.
\leanref{GlobalTheory.\allowbreak{}TwV.\allowbreak{}Xc.\allowbreak{}ChiPow.\allowbreak{}tw\_\allowbreak{}ch\_\allowbreak{}rat}{Development.\allowbreak{}AbelianVarietyTheory.\allowbreak{}Proofs.\allowbreak{}SatakeParameters.\allowbreak{}KernelSiteNormalisingCharPowValuesInCoeffField}{79--99}

For the test, carry the lattice, its prescribed residue maps and its local frames by \(y\), and do the same at the target using \(y'\). A frame at a marked prime is a basis inducing the prescribed coordinates on its finite quotient. Lifts of a basis of the quotient give such frames. At primes of the word level take the indicator that \(Z\) maps the source lattice into the target lattice and preserves the prescribed residue map. Multiply it by the following factors at the type and twist primes. At a matrix step read \(Z\) in the two local frames and take the finite Fourier sum of the local test vector over its free matrix entries. At the penultimate step append the column in the specified target coset before evaluating the local vector. At the last step evaluate the test vector on the projected matrix and its finite Fourier factor on the omitted row against the specified source coset. Definition~\ref{B3.35f} gives the finite summation, its volume factor and the last-row Fourier character in the coordinates of these frames. In particular the coset belongs to the target in the second case and to the source in the third.

Two lifts of the same frame differ by a transformation acting trivially on the finite quotient on which the local vector is read. In the Fourier sum this transformation permutes the summation classes and preserves the additive character. In the appended column the difference lies in its defining congruence lattice. In the last row its pairing with the Fourier variable is integral. These observations prove independence of the lifts and the congruence assertion. They also show that equality of the two adelic points at the test places gives equality of the functions, including the residue-map indicator. The values are algebraic: the finite sums use values of the finite characters and rational volume factors. Their denominators are supported on the prescribed exceptional primes.

A rational change carries each decorated point and its frames to the changed point. Reading \(Z\) in those frames replaces it by \((g')^{-1}Zg\). For a right level change, first keep the target fixed. In a matrix step the change of source frame permutes the free entries in the finite sum. The determinant character in the Weil action contributes the source factor of \(\mu_{j+1}(\det k)\), and the eigenvalue of the chosen source vector contributes the inverse factor of \(\lambda_j(k)\). In the appended-column test the same change acts on the matrix columns with the appended column fixed. For a target change the appended column is carried within its prescribed coset. In the projected-row test the change acts on the row and by the inverse action on its Fourier variable. Reindexing the finite sum preserves their pairing. At either endpoint the local determinant factors multiply to the step character on its determinant, by the definition of the step character as the untwisted character times the finite twist. The level-vector factors multiply to the level character. The target factors have the inverse determinant and direct level character. Multiplication gives the second displayed identity, also when the test value is zero.
\leanref{GlobalTheory.\allowbreak{}TwV.\allowbreak{}levAt\_\allowbreak{}mat}{Development.\allowbreak{}AbelianVarietyTheory.\allowbreak{}Proofs.\allowbreak{}SatakeParameters.\allowbreak{}FiniteDatumLevelEquivarianceAtMatrixSteps}{18--77}
\leanref{GlobalTheory.\allowbreak{}AdLevCpl.\allowbreak{}levAt\_\allowbreak{}coupled}{Development.\allowbreak{}AbelianVarietyTheory.\allowbreak{}Proofs.\allowbreak{}SatakeParameters.\allowbreak{}AdelicDatumLevelEquivarianceFirstCoupledStep}{13--37}
\leanref{GlobalTheory.\allowbreak{}TwTop.\allowbreak{}levAt\_\allowbreak{}top}{Development.\allowbreak{}AbelianVarietyTheory.\allowbreak{}Proofs.\allowbreak{}SatakeParameters.\allowbreak{}LastStepAdelicDatumLevelEquivariance}{42--115}
\end{proof}

\subsection*{Proof of \ref{B3.35g}}\label{prf:B3.35g}\phantomsection
\begin{proof}
Write \(M_i\) for the period module at pair \(i\), identifying its elements with their analytic images. Let \(A_i\) have localizations \(c^{-1}M_i\) on \(S\), \(cM_i\) on \(\bar S\), and \(M_i\) elsewhere. Interchange \(c\) and \(c^{-1}\) to define \(B_i\). For the moment suppress \(i\). Choose sets \(X,T,Y\) of representatives of
\[
 A/(A\cap M),\qquad (M\cap A)/(M\cap A\cap B),\qquad B/(B\cap M).
\]
Write \(B_H(x,y)=\Im H(x,y)\) for the alternating form and \(W_x\) for the translations in \eqref{eq:p3c3-translation}. Fourier transformation in the middle quotient with counting measure of total mass one gives the operator
\[
 F\longmapsto\frac1{|T|}
 \sum_{x\in X,\ t\in T,\ y\in Y}
 \xi(t-x)\exp(2\pi i B_H(t,y))\,W_xW_yF.
\]
Indeed the Fourier coefficient at \(y\) is the sum over \(t\) divided by \(|T|\). Its inverse Fourier term acts by \(W_y\), and translation of that test coset acts by \(W_x\). The three quotients are finite because all their lattices are commensurable.

Here are the three changes of representatives on an \(M\)-fixed section. Adding \(h\in M\cap A\cap B\) to \(t\) preserves \(\xi\) by its congruence at \(S\), and \(B_H(h,y)\) is integral since \(h,y\in B\). Adding \(h\in B\cap M\) to \(y\) preserves the character since \(t,h\in M\), and \(W_{y+h}F=W_yW_hF=W_yF\). If \(x\) is replaced by \(x+h\) with \(h\in A\cap M\), replace \(t\) by \(t+h\) modulo the middle denominator. The argument \(t-x\) stays fixed. The additional Fourier factor is \(\exp(2\pi i B_H(h,y))\). Moving \(W_h\) past \(W_y\) gives its inverse, and then \(W_hF=F\).

Fix a reference pair \(0\). At \(P,\bar P\) choose exponents large enough to carry every source lattice into the reference source and conversely, and likewise for every target. The Chinese remainder theorem gives \(g\in\mathcal O_E\) in these prime powers with \(g\equiv1\bmod c\) and \(g\equiv1\bmod Q\) at each \(Q\in S\) prime to \(c\). Put \(b=g^2\). For a homomorphism \(z\) from the \(i\)-th source to the \(i\)-th target, \(b z(v)=g z(gv)\) carries the reference source into the reference target. The same calculation in the opposite direction shows that \(b\) carries the respective period modules into one another. Localizing proves the same inclusions for \(A_i,B_i\) and their indicated intersections.

For each of the three quotients, denote the numerator by \(U_i\) and the denominator by \(V_i\). The local descriptions give \(cU_i\subseteq V_i\). Multiplication by \(b\) defines maps in both directions between the reference quotient and the \(i\)-th quotient. Their composite is multiplication by \(b^2\), which acts as the identity since \(b^2-1\in c\mathcal O_E\). Thus the actual elements \(bx,bt,by\) obtained from one choice of reference representatives serve at every \(i\). They have the same analytic images, the same character factors, and the same test values by Lemma~\ref{B3.35e}. The resulting finite sum is one complex linear operator on the ambient function space. The representative-independence proved above identifies its restriction to each fixed section with the prescribed test there.

For the coefficient assertion, first consider the elliptic factor of the standard section. The discriminant normalization gives an algebraic Weierstrass equation: the CM \(j\)-invariant is algebraic, and the discriminant fixes the scale of \(g_2,g_3\). The coordinate differential is the invariant differential of this equation. The reduced odd theta section with derivative one has algebraic Taylor coefficients, by \cite[Corollary~2.2]{bannai-kobayashi-2010}. The divisor of this section is the origin, hence symmetric and defined over the coefficient field. Its translate by a rational period point has algebraic coefficients by the translation theorem for a symmetric divisor, after adjoining the torsion coordinates. An algebraic local equation for the translated divisor has algebraic Taylor coefficients in the invariant differential, so restoring this local equation also preserves algebraicity. This is the elliptic case of \cite[Theorem~2.9]{bannai-kobayashi-2010}. Translating to the prescribed nonzero half-period and dividing by its nonzero value gives the normalized even section. Further torsion translations compose up to roots of unity.

For the one-entry matrix section use the fourth-power normalization in Lemma~\ref{B4.22e} with its initial rational isometry \(g_0\). Its scalar \(\kappa_0\) is algebraic by the defining root equations. The entries of the isometry in CM embedding coordinates are algebraic: the trace matrix, the change to a Hermitian orthogonal basis and the four-square matrices are over number fields. Translating the four factors by the same rational period point gives a finite sum of translates of the elliptic product. Commensurability makes each translating point torsion modulo the elliptic product lattice. The normalized averaging coefficients are inverse square roots of positive integers times roots of unity. It follows that the fourth tensor power of the translated one-entry section has algebraic coefficients. If \(a_\gamma\ne0\) is any coefficient of that section, the tensor coefficients at \((\gamma,\gamma,\gamma,\gamma)\) and \((\nu,\gamma,\gamma,\gamma)\) give \(a_\gamma^4\) and \(a_\nu a_\gamma^3\). Hence every \(a_\nu\) is algebraic. The zero section satisfies the same assertion.

Tensor over the matrix entries and perform the two endpoint averages of Lemma~\ref{B3.35d}. The determinant character values in those averages are algebraic. Indeed for a norm-one ideal \(\mathfrak a\), principalize \(\mathfrak a^h=(z)\); then \(\mathfrak a^{2h}=(z/\bar z)\), to which the principal-ideal character formula applies. Finally every \(x,y,t\) in the finite sum is a rational period point. Its translation phases are roots of unity, and the values of \(\xi\) are algebraic by Definition~\ref{B3.35f}. The output is a finite algebraic combination of entire torsion translates with algebraic Taylor coefficients. Taylor expansion, followed coefficientwise by the prescribed embedding of \(\overline{\Q}\), therefore carries every rational linear relation between these normalized outputs to the same relation on their series.
\end{proof}

\subsection*{Proof of \ref{B3.35a-arithmetic}}\label{prf:B3.35a-arithmetic}\phantomsection
\begin{proof}
Write \(E/F\) for the host CM extension and \(K_0\) for its imaginary quadratic subfield. Write \(K\hookrightarrow\C_p\) for the local coefficient embedding, compatible with \(\lambda:\overline{\Q}\hookrightarrow\C_p\) on the global coefficient field. Fix a component at each endpoint. At level \(l\) let \(M_l,J_l\) be the base lattice and form, \(C_l\) the change to standard orthogonal coordinates, and \(a_*\) the level parameter. The point \(y\) of a tuple \(L\) has its component representative away from the moving places. At a moving split place its first matrix, in the chart basis, is a basis matrix of \(L_v\). Its other matrix is determined by the unitary equation. With \(\#\) denoting the \(J_l\)-dual, the kernel lattice and determinant idèle are
\[
 M(y)=E^{N_l}\cap y\widehat M_l,\qquad
 A(y)=C_l\bigl(M(y)+a_*M(y)^\#\bigr),\qquad d_y=\det y .
 \tag{3.35.8}\label{eq:p3c3-chart-lattice}
\]
Indeed completion gives the corresponding translate of the base lattice, and taking the dual commutes with a unitary translate. The base sum in parentheses is the standard self-dual companion in orthogonal coordinates. Thus \(A(y)\) is self-dual, and its first completion at \(v\) is the image of \(L_v\) under the chart matrix. At level zero use the initial chart and at positive levels the target chart of the preceding lift.
\leanref{GlobalTheory.\allowbreak{}pairOK\_\allowbreak{}allSites}{Development.\allowbreak{}AbelianVarietyTheory.\allowbreak{}Proofs.\allowbreak{}SatakeParameters.\allowbreak{}KernelSiteLatticeTuplePairsAdmissible}{23--48}

Fix the place \(v\) and the layer \(j\). Write \(n=N_j,\ n'=N_{j+1}\), \(2\eta=[\operatorname{Frac}W_v:\Q_p]\), and put
\[
 \Delta_w(l)=\sum_{s\ {\rm even}}w_{w,l}(s)
                  -\sum_{s\ {\rm odd}}w_{w,l}(s),\qquad
 \Delta=\Delta_v(j+1).
 \tag{3.35.9}\label{eq:p3c3-signed-weights}
\]
Let \(\pi\) be the norm-one idèle equal to \((p,p^{-1})\) at the moving pair \(P,\bar P\) and one elsewhere. Put
\(M_\pi=\lambda(\mu_{j+1}(\pi))\) and \(C_\pi=\lambda(\chi_I(\pi))\).
The lift \((p,1)\) under \(x\mapsto x/\bar x\) is supported above \(p\). There the untwisted step character is unramified, since \(p\) differs from the type primes. Its value is the product of its uniformizer values to their respective orders. Those values and the value of the finite twist belong to the kernel coefficient field included in the global coefficient field. Hence \(M_\pi\in K^\times\). By Lemma~\ref{B3.35e}, \(C_\pi^{b_j}\) and \(C_\pi^{b_{j+1}}\) also belong to \(K\) if both levels have components.
\leanref{GlobalTheory.\allowbreak{}TwV.\allowbreak{}Xc.\allowbreak{}MuLL.\allowbreak{}mu\_\allowbreak{}mem\_\allowbreak{}LL}{Development.\allowbreak{}AbelianVarietyTheory.\allowbreak{}Proofs.\allowbreak{}SatakeParameters.\allowbreak{}StepCharacterUniformizingIdeleMemCoeffField}{91--176}

To compute the valuations, raise \(\pi\) to a positive power \(m_0\) such that its first ideal factor to that power is principal, say \((y_0)\). Then \(\pi^{m_0}=z k\) with \(z=y_0/\bar y_0\) and \(k\) an integral norm-one idèle. On the compact units the step character has finite image. On \(z\) it is \(\prod_\tau\tau(z)^{e_\tau}\), where the filling gives
\(e_{w,s}=w_{w,j+1}(s)+b_j+n'/2\).
The valuations at even and odd embeddings have opposite signs. The constant shift cancels because there are equally many of each. Only \(v\) contributes, with first valuation \(m_0\). Thus \(v_p(M_\pi)=\Delta\).

For a norm-one ideal \(\mathfrak a\), write it as \(\mathfrak b_1/\bar{\mathfrak b}_1\) by choosing one prime from each conjugate pair in its support. Principalizing a positive power of \(\mathfrak b_1\) gives \(\mathfrak a^h=(z_1)\), \(z_1\bar z_1=1\). Hence \(\chi(\mathfrak a)^h=\prod_{\tau\in\Psi}\tau(z_1)^{-1}\). At an inert prime of \(K_0\) a decomposition-group element inducing Frobenius on \(K_0\) interchanges \(\Psi\) and its conjugate. It preserves the valuation and carries this product to its inverse. The product has valuation zero. This applies to the determinant-ideal value of \(\pi\). Consequently
\[
 a=p^{-\Delta}M_\pi C_\pi^{b_j},\qquad
 a'=p^\Delta M_\pi^{-1} C_\pi^{-b_{j+1}}
 \tag{3.35.11}\label{eq:p3c3-two-units}
\]
are units in the local coefficient field. Their definition involves neither component nor lattice. If either component set is empty take \(a=a'=1\) and all maps zero.
\leanref{GlobalTheory.\allowbreak{}TwV.\allowbreak{}Xc.\allowbreak{}lambdaCharacter\_\allowbreak{}norm\_\allowbreak{}eq\_\allowbreak{}p\_\allowbreak{}pow\_\allowbreak{}of\_\allowbreak{}pos}{Development.\allowbreak{}AbelianVarietyTheory.\allowbreak{}Proofs.\allowbreak{}SatakeParameters.\allowbreak{}NormLambdaCharacterNormOneIdeleEqPPow}{19--85}
\leanref{GlobalTheory.\allowbreak{}TwV.\allowbreak{}Xc.\allowbreak{}two\_\allowbreak{}units}{Development.\allowbreak{}AbelianVarietyTheory.\allowbreak{}Proofs.\allowbreak{}SatakeParameters.\allowbreak{}TwoUnitsAtPlaceUniformizingIdele}{38--82}

Fix the two components and every entry away from \(v\). This is the record \(\omega\). Consider a nonempty finite family of pairs in this fibre and designate one pair as \(0\). Denote its members' kernel lattices by \(A_i,B_i\), and their source and target lengths by \(s_i,s'_i\). In the fixed chart matrices the first completion of the matrix lattice is
\(D_i=G'_i\operatorname{Mat}_{n'\times n}(W_v)G_i^{-1}\), where \(G_i,G'_i\) are bases of the two tuple entries. Its indicator in these coordinates is precisely the containment test. The fixed ideal and chart factors cancel in relative volumes. The \(n\) columns see \(G'_i\) on the left and the \(n'\) rows see \(G_i^{-1}\) on the right. With \(d_i=ns'_i-n's_i\), the length convention gives
\[
 \operatorname{vol}(D_i)/\operatorname{vol}(D_0)
                =p^{2\eta(d_i-d_0)} .
 \tag{3.35.12}\label{eq:p3c3-volume}
\]

Here is the finite Fourier calculation that fixes the sign of this power. Put \(P_i=P(A_i,B_i)\), \(P_-=\bigcap_iP_i,\ P_+=\sum_iP_i\). By self-duality \(P_+^\#=P_-\). The alternating pairing on \(P_+/P_-\) is perfect. The only primary parts are the two over \(P,\bar P\). Conjugation interchanges these primes, so each part is isotropic and the second is the character dual of the first. Write the first as \(V\). The image of \(P_i\) in the two halves is \(D_i\oplus D_i^\perp\), writing \(D_i\) also for its finite image. A local generator of the curve ideal identifies these first images with the preceding matrix lattice quotients.

On functions on \(V\) use \(T_af(t)=f(t+a)\) and \(M_bf(t)=b(t)f(t)\). On sections the corresponding operations are
\[
 W_xF(Z)=\alpha(x)F(Z+x)
          \exp\bigl(-\pi H(Z,x)-\pi H(x,x)/2\bigr),
 \tag{3.35.14}\label{eq:p3c3-translation}
\]
where \(H(Z,x)=\sum_{\tau,a,b}Z_{\tau,ba}\overline{x_{\tau,ba}}/
\operatorname{area}(\delta\mathfrak c)\) and
\(\alpha(x)=\exp(\pi i\sqrt d\,\Re H(x,x)/2)\).
Their commutator is \(b(a)=\exp(-2\pi i\,\Im H(a,b))\) in period coordinates.
To choose lifts in each half with its additive translation law, choose an integer of \(E\) in \(P\) coprime to its conjugate, taking \((\varrho)=P^h\) for a positive multiple \(h\) of the ideal class order. For sufficiently large \(J\) the two norm-one multiples
\((\bar\varrho/\varrho)^{\pm2J}P_0\) contain the respective local halves. Intersect these multiples with each \(P_i\). The intersections sum to \(P_i\): this can be tested locally, where one of the two multipliers enlarges the required lattice and the two lattices agree away from the moving pair. The semicharacter law on each multiple is transported from \(P_0\), since the multiplier preserves \(H\).

Let \(\theta_i=C(A_i,B_i)^{-1}e_{A_i,B_i}\). There is a linear map \(S\) sending \(T_aM_b\mathbf1_{D_0}\) to \(W_aW_b\theta_0\). Indeed shifts by \(D_0\) and modulations by \(D_0^\perp\) fix the generating vectors. Commuting a shift past a modulation gives the same scalar on each side. On each coset of \(D_0\), the distinct characters of \(D_0\) are independent and span the functions there. Distinct cosets have disjoint support. This accounts for all relations. Put \(J_i=D_i\cap D_0\), \(m_i=[D_i:J_i]\), \(n_i=[D_0:J_i]\). Sum first over the characters of \(D_0/J_i\); this gives \(n_i\mathbf1_{J_i}\) on \(D_0\). Then sum over translations by \(D_i/J_i\); this gives \(n_i\mathbf1_{D_i}\). The normalized move divides by \(\sqrt{m_i n_i}\), since the second-half index is \(n_i\). Apply Lemma~\ref{B3.35d} first to the source and then to the target, including the intermediate pair in the family. The two square roots of relative volume multiply. Thus
\[
 S(\mathbf1_{D_i})=
       \sqrt{\operatorname{vol}(D_i)/\operatorname{vol}(D_0)}\,\theta_i .
 \tag{3.35.15}\label{eq:p3c3-move}
\]
\leanref{Formalism.\allowbreak{}L2.\allowbreak{}PairingIndexBounds.\allowbreak{}Mv.\allowbreak{}MovesU\_\allowbreak{}of}{Formalism.\allowbreak{}LocalToGlobal.\allowbreak{}PairingIndexBounds.\allowbreak{}ModelLatticeMoves.\allowbreak{}ModelIndexPairings.\allowbreak{}MovesU}{100--234}

The finite test applies one linear operator to these sections on the whole finite family. To see this, for each \(P_i\) enlarge its local lattice by \(c^{-1}\) at \(S\) and shrink it by \(c\) at \(\bar S\), obtaining \(A\); let \(B\) use the opposite changes. Choose representatives \(X,T,Y\) of \(A/(A\cap P_i)\), \((P_i\cap A)/(P_i\cap A\cap B)\), \(B/(B\cap P_i)\), respectively. The operator is
\[
 {\cal T}_\xi F=|T|^{-1}
 \sum_{x\in X,\,t\in T,\,y\in Y}
 \xi(t-x)e^{2\pi i\,\Im H(t,y)}W_xW_y F .
 \tag{3.35.16}\label{eq:p3c3-finite-theta}
\]
Use period preimages in the argument of \(\xi\). Lemma~\ref{B3.35e} makes \(\xi\) invariant under \(A\cap B\). Changing \(t\) by a period in the denominator of its quotient leaves the additive character unchanged. Adding \(h\in B\cap P_i\) to \(y\) preserves the Fourier character because \(t,h\in P_i\), and \(W_{y+h}F=W_yW_hF=W_yF\) on a \(P_i\)-fixed section. For \(x\mapsto x+h\), \(h\in A\cap P_i\), reindex \(t\mapsto t+h\); its extra Fourier factor cancels the commutator of \(W_h\) with \(W_y\).

The moving primes are outside \(S,\bar S\) and do not divide \(c\). By the Chinese remainder theorem choose \(g\equiv1\bmod c\mathcal O_E\) clearing the finitely many lattice comparisons in both directions. Then \(g^2\) clears the Hom-lattice comparisons, one factor at either endpoint. These three finite quotients are killed by \(c\). Multiplication by \(g^2\) in the two directions induces inverse maps since \(g^4\equiv1\bmod c\). By Lemma~\ref{B3.35g}, multiply each of the three reference sets by this same \(g^2\) in the rational period module and then embed it into the analytic space. This gives three sets of actual points common to all the operators in \eqref{eq:p3c3-finite-theta}. Lemma~\ref{B3.35e} gives the same \(\xi\) on them. If \(H_i={\cal T}_\xi e_{A_i,B_i}\), every finite rational relation \(\sum_i q_i\mathbf1_{D_i}=0\) consequently gives
\[
       \sum_i q_i\,p^{\eta(d_i-d_0)} C(A_i,B_i)^{-1}H_i=0 .
 \tag{3.35.17}\label{eq:p3c3-complex-relation}
\]
\leanref{GlobalTheory.\allowbreak{}KerCommon.\allowbreak{}kerTh\_\allowbreak{}common\_\allowbreak{}op}{Development.\allowbreak{}AbelianVarietyTheory.\allowbreak{}Proofs.\allowbreak{}SatakeParameters.\allowbreak{}KernelFunctionsCommonLinearOperatorOffSplitPrime}{128--175}
\leanref{GlobalTheory.\allowbreak{}KerRel.\allowbreak{}kerTh\_\allowbreak{}relation\_\allowbreak{}of\_\allowbreak{}indicator}{Development.\allowbreak{}AbelianVarietyTheory.\allowbreak{}Proofs.\allowbreak{}SatakeParameters.\allowbreak{}KernelThetaRelationOfIndicatorRelationSplit}{76--122}

Lemma~\ref{B3.35g} gives algebraic Taylor coefficients for each \(H_i\) and algebraicity of the nonzero character factors \(C(A_i,B_i)\). Also \(p^{\eta(d_i-d_0)}\) is rational. Taking each Taylor coefficient of \eqref{eq:p3c3-complex-relation} therefore gives an equality in \(\overline{\Q}\) to which the same embedding \(\lambda\) applies for every \(i\).
\leanref{GlobalTheory.\allowbreak{}AlgTranslate.\allowbreak{}coefIn\_\allowbreak{}translate}{Development.\allowbreak{}AbelianVarietyTheory.\allowbreak{}Proofs.\allowbreak{}SatakeParameters.\allowbreak{}CoefInTranslate}{53--82}
\leanref{GlobalTheory.\allowbreak{}kerTh\_\allowbreak{}entire\_\allowbreak{}algCoef}{Development.\allowbreak{}AbelianVarietyTheory.\allowbreak{}Proofs.\allowbreak{}SatakeParameters.\allowbreak{}KerThEntireAlgCoef}{22--39}

In the fixed coordinates the series is obtained by Taylor expansion and \(\lambda\), followed by the chart substitution and restriction, from
\[
 p^{\sum_w\Delta_w(j+1)(s'_{i,w}-s_{i,w})}
 \frac{\mu_{j+1}(d_{y'_i})}{\mu_{j+1}(d_{y_i})}
 c(A_i)^{-b_j}c(B_i)^{b_{j+1}}\, C(A_i,B_i)^{-1}H_i .
 \tag{3.35.18}\label{eq:p3c3-chart-series}
\]
The two chart matrices in the substitution are constant on the fibre. At level \(l\), \(c(A(y))\) equals \(\chi_I(d_y)\) times a constant depending on the fixed rational change to standard coordinates. If the source and target lengths increase by \(h,h'\), the determinant idèles change by \(\pi^{-h},\pi^{-h'}\) times integral units supported at the moving pair. Both characters are trivial on those units. The ratio of the scalars preceding \(C(A_i,B_i)^{-1}H_i\) is therefore
\[
 p^{\Delta(h'-h)}M_\pi^{h-h'}C_\pi^{b_jh-b_{j+1}h'}=a^h(a')^{h'} .
\]
Set \(k_i=a^{s_i}(a')^{s'_i}p^{-\eta d_i}\). Applying the common substitutions to \eqref{eq:p3c3-complex-relation} now gives \(\sum_iq_i k_i^{-1}\Theta_i^F=0\). Thus \(\mathbf1_{D_i}\mapsto k_i^{-1}\Theta_i^F\) defines a rational linear map on the span of the fibre indicators. Extend it by zero on a vector-space complement. This is \(\Xi_\omega\) in \eqref{eq:p3c3-indicator}.

Finally let a source group element have chart matrices \(G_w\), and write the changed point as \(g y k\), with \(g\) rational unitary and \(k\) in the level. Lemma~\ref{B3.35e} gives the multiplier \(\mu_{j+1}(\det k)\lambda_j(k)^{-1}\) and the substitution \(Z\mapsto Zg\) in standard coordinates. The source idèle denominator in \eqref{eq:p3c3-chart-series} cancels \(\mu_{j+1}(\det k)\). The principal-ideal law cancels the constant \(b_j+n'/2\) in the type at \(g\). The remaining exponents are exactly the weights \(w_{w,j+1}(s)\). The lengths change by \(-\operatorname{ord}_w\det G_w\). Hence, writing \(\tau_{w,s}\) for the indicated embedding and \(S_G\) for the source-row substitution by \(\tau_{w,s}(G_w)^{\mathsf T}\) at even \(s\) and by \(\tau_{w,s}(G_w)^{-1}\) at odd \(s\), we obtain
\[
 \Theta^F_{\gamma L,L'}=\kappa(\gamma)^{-1} S_G\Theta^F_{L,L'},\qquad
 \kappa(\gamma)=\chi_{\rm lev}(\gamma)
 \prod_{w,s}\tau_{w,s}(\det G_w)^{(-1)^s w_{w,j+1}(s)}
 \prod_w p^{-\Delta_w(j+1)\operatorname{ord}_w\det G_w}.
 \tag{3.35.19}\label{eq:p3c3-chart-action}
\]
Here \(\chi_{\rm lev}(\gamma)\) is the local coefficient realization of \(\lambda_j(k)\). For a stabilizer this formula gives the stabilizer identity. Applying it to any element taking the chosen representative to \(L\) recovers the original series at \(L\). A lower neighbour of index \(i\) multiplies the basis determinant by a factor of order \(i\), so its signed length decreases by \(i\). The other lengths and \(\omega\) are unchanged.
\leanref{GlobalTheory.\allowbreak{}XiGen.\allowbreak{}xiEx\_\allowbreak{}theta\_\allowbreak{}general}{Development.\allowbreak{}AbelianVarietyTheory.\allowbreak{}Proofs.\allowbreak{}SatakeParameters.\allowbreak{}XiExThetaGeneral}{86--304}
\leanref{GlobalTheory.\allowbreak{}TwV.\allowbreak{}site\_\allowbreak{}law\_\allowbreak{}xiIn}{Development.\allowbreak{}AbelianVarietyTheory.\allowbreak{}Proofs.\allowbreak{}SatakeParameters.\allowbreak{}KernelSiteLawAndFixedCoordinateStabiliserLaw}{20--48}
\end{proof}

\subsection*{Proof of \ref{B3.35c}}\label{prf:B3.35c}\phantomsection
\begin{proof}
For representatives \(A,B\) in a component the left contribution of their orbits is \(\sum_{\gamma:A\to_h\gamma B}\phi(A,\gamma B)/(w(A)w(B))\), since each \(\gamma B\) occurs \(w(B)\) times. Inversion sends the finite indexing set to \(\gamma^{-1}A\to_h B\). Simultaneous invariance identifies the summands. Summing over orbits and components gives \eqref{eq:p3c3-orbit-adjoint}.

Apply it coefficientwise to \(\operatorname{ker}_{j,B,L'}(f(A))\). For each edge \(A\to_h B\), the edge calculation in Theorem~\ref{B3.34} gives \(p^{\mathcal A_h}H_j(A)\subseteq H_j(B)\), so \(f(A)\in K H_j(B)\). For a source polynomial \(G\), put \(\rho(\gamma)G=\chi_{\rm lev}(\gamma)(\gamma G)\), where \(\gamma G\) is the level-group action on polynomials in Definition~\ref{B3.05a}. Then \(f(\gamma A)=\rho(\gamma)f(A)\) and Definition~\ref{B3.05a} gives
\(\operatorname{ker}_{j,\gamma B,L'}(\rho(\gamma)f(A))=\operatorname{ker}_{j,B,L'}(f(A))\).
Linearity sums the input at \(B\) to \((T_{v,h}^{(j)}f)(B)\). This proves \eqref{eq:p3c3-hecke-adjoint} for \eqref{eq:p3c3-class-expansion}. At the terminal step it is the terminal neighbour clause of Definition~\ref{B3.05a}. For the arithmetic class data, \(y\) again has one value and \(x\) runs over the source representatives with their stabilizer orders.
\end{proof}

\subsection*{Proof of \ref{B3.36}}\label{prf:B3.36}\phantomsection
\begin{proof}
If \(xA\not\subseteq B\), both sides of \eqref{eq:p3c3-neighbour-count} vanish. Otherwise let \(d=\dim_{\F_q}\ker(\bar x:A/pA\to B/pB)\). The lower neighbours correspond to codimension-\(i\) subspaces of \((B/pB)/\operatorname{im}\bar x\) of dimension \(d+g\). Multiplication by \(p\) identifies the upper neighbours in the \(\ell\)-term with \((i-\ell)\)-dimensional subspaces of \(\ker\bar x\). Thus the count reduces to
\[
 {d+g\brack i}_q=\sum_{\ell=0}^{\min(i,g)}
 q^{(i-\ell)(g-\ell)}{g\brack\ell}_q{d\brack i-\ell}_q.
 \tag{3.36.5}\label{eq:p3c3-gaussian-convolution}
\]
To count \(i\)-dimensional subspaces of \(D\oplus G\), with dimensions \(d,g\), fix the intersection \(J\) with \(G\) and the image \(H\) in \(D\). For \(\dim J=\ell\), the subspace modulo \(J\) is the graph of a map \(H\to G/J\). There are \(q^{(i-\ell)(g-\ell)}\) such maps for each \(J,H\).

For the lift take \(g=t+1,\ N=N_j,\ \eta=\eta_v,\ \Delta=\Delta_v(j+1)\). A target lower edge of index \(i\) decreases \(s'\) by \(i\) and \(d=Ns'-N's\) by \(Ni\). A source upper edge of index \(h\) increases \(s\) by \(h\) and decreases \(d\) by \((N+g)h\). Other lengths and \(\omega\) are fixed by Lemma~\ref{B3.35a}. With \(b,G\) fixed in Lemma~\ref{B3.35b}, the respective multipliers of \(J_{L,L'}(b;G)\) are
\[
 C_A(i)=p^{\Delta i}(a')^{-i}p^{\eta Ni},\qquad
 C_C(h)=p^{\Delta h}a^hp^{\eta(N+g)h}.
 \tag{3.36.6}\label{eq:p3c3-edge-multipliers}
\]
\leanref{Formalism.\allowbreak{}MainNotions.\allowbreak{}FormalGroupSystem.\allowbreak{}Pf.\allowbreak{}kernel\_\allowbreak{}target\_\allowbreak{}identity}{Formalism.\allowbreak{}SharpRestriction.\allowbreak{}GrowthBounds.\allowbreak{}MultisetValuation.\allowbreak{}IntegralSpanCoordinates}{80--126}
\leanref{Formalism.\allowbreak{}MainNotions.\allowbreak{}FormalGroupSystem.\allowbreak{}Pf.\allowbreak{}kernel\_\allowbreak{}source\_\allowbreak{}identity}{Formalism.\allowbreak{}SharpRestriction.\allowbreak{}GrowthBounds.\allowbreak{}MultisetValuation.\allowbreak{}IntegralSpanCoordinates}{127--174}

Apply \(J_{L,L'}(\,\cdot\,;G)\) at the original pair to \eqref{eq:p3c3-neighbour-count} as an identity of indicator sums. For the target sum, multiply by \(C_A(i)\) to obtain the actual kernel values. For each source sum, divide the actual kernel values by \(C_C(i-\ell)\) to obtain the evaluation at the original pair. The target neighbour sum of kernels is therefore the sum of the source neighbour sums with coefficients \(C_A(i)C_C(i-\ell)^{-1}q^{(i-\ell)(g-\ell)}{g\brack\ell}_q\).
\leanref{Formalism.\allowbreak{}MainNotions.\allowbreak{}FormalGroupSystem.\allowbreak{}Pf.\allowbreak{}ker\_\allowbreak{}dn\_\allowbreak{}sum}{Formalism.\allowbreak{}SharpRestriction.\allowbreak{}GrowthBounds.\allowbreak{}MultisetValuation.\allowbreak{}CoupledSumOfEntries}{38--64}

Use \(G=f(L_{x,y})\) and sum over the class data with weights \(c_K/w_x\). The target sum gives \(T_{v,i}^{(j+1)}K_j f\). By Lemma~\ref{B3.35c}, the source neighbour sums act on \(f\) to give \(K_jT_{v,i-\ell}^{(j)}f\). Theorem~\ref{B3.34} ensures that \(T_{v,h}^{(j)}\) preserves \(V_j\) for \(h\le N_j\), and the sum is empty for \(h>N_j\). The resulting equality descends to \(K\).
\leanref{Formalism.\allowbreak{}MainNotions.\allowbreak{}FormalGroupSystem.\allowbreak{}Pf.\allowbreak{}hk\_\allowbreak{}point}{Formalism.\allowbreak{}SharpRestriction.\allowbreak{}GrowthBounds.\allowbreak{}MultisetValuation.\allowbreak{}HeckeConditionOfAdm}{15--74}

Put \(Q=p^\eta\). Then
\[
 {g\brack\ell}_{Q^2}=Q^{\ell(g-\ell)}
 e_\ell(Q^{g-1},Q^{g-3},\ldots,Q^{1-g}).
 \tag{3.36.7}\label{eq:p3c3-gaussian-elementary}
\]
Indeed the elementary generating polynomial at \(g\) is the one at \(g-1\) with \(z\) replaced by \(Qz\), times \(1+Q^{1-g}z\). After the factor \(Q^{\ell(g-\ell)}\), this gives the recursion with multipliers \(Q^{2\ell}\) and \(1\) in \eqref{eq:p3c3-gaussian-convolution} for dimensions \(g-1\) and \(1\). Here \(g=t+1\), so the arguments are \(\beta_{v,a}\). In the normalized kernel coefficient the exponent of \(Q\) is
\[
 -i(N+g-i)+Ni-(N+g)(i-\ell)+2(i-\ell)(g-\ell)+\ell(g-\ell)
 =-(i-\ell)(N-i+\ell).
\]
This is the source normalization. The remaining factors are \(((aa')^{-1})^i(ap^\Delta)^\ell e_\ell((\beta_{v,a})_a)\). Thus \eqref{eq:p3c3-one-step} holds with \(u_{j+1}=(aa')^{-1}\) and \(\mathfrak a_{j+1}=ap^\Delta\). Both sides are supported on lattice tuples.
\end{proof}

\subsection*{Proof of \ref{B3.38}}\label{prf:B3.38}\phantomsection
\begin{proof}
Fix \(v\) and suppress it in subscripts. Take \(u_l,\mathfrak a_l\) from Lemma~\ref{B3.36}. Set \(U_0=1\) and \(\mathfrak m_l=\mathfrak a_l/U_{l-1}\). Since the \(U_l\) are units, \eqref{eq:p3c3-prefix-unit} follows. Set \(\widehat{\mathcal E}_i^{(b)}=U_b^{-i}\mathcal E_i^{(b)}\) and \(\widehat{\mathcal E}^{(b)}(z)=\sum_{i\ge0}\widehat{\mathcal E}_i^{(b)}z^i\). After normalization the scalar at \(i,\ell\) in \eqref{eq:p3c3-one-step} is
\(U_{j+1}^{-i}u_{j+1}^i\mathfrak a_{j+1}^\ell U_j^{i-\ell}=\mathfrak m_{j+1}^\ell\).
Thus the \(j\)th lift inserts \(\prod_{a=0}^t(1+\mathfrak m_{j+1}\beta_a z)\) in the generating polynomial. Iterate and multiply by the reciprocal product \(\sum_{\ell\ge0}(-1)^\ell h_\ell(\mathcal F_b)z^\ell\). The coefficient of \(z^r\) gives
\[
 \left(\sum_{\ell=0}^r(-1)^\ell h_\ell(\mathcal F_b)
 \widehat{\mathcal E}_{r-\ell}^{(b)}\right)K^{[b]}
 =K^{[b]}\mathcal E_r^{(0)}.
 \tag{3.38.8}\label{eq:p3c3-reciprocal-product}
\]
On \(V_b\) the index-zero operator is the identity, so
\(\widetilde T_i^{\,u,(b)}=p^{N'_b(i)+\eta i(N_b-i)}\widehat{\mathcal E}_i^{(b)}\)
also holds at zero. Since
\(\ell(N_b-2r+\ell)+(r-\ell)(N_b-r+\ell)=r(N_b-r)\),
the sum on the left of \eqref{eq:p3c3-reciprocal-product} multiplied by \(p^{N'_b(r)+\eta r(N_b-r)}\) is \(Z_b\). Inserting \(\mathcal E_r^{(0)}=p^{-\eta r(k-r)-\mathcal A_v}y_v\) gives exponent
\[
 N'_b(r)-\mathcal A_v+\eta rb(t+1)=P_v^{\Z}(b).
\]

For \(1\le l\le b\), Definition~\ref{B3.01} gives \(w_l(s)\ge w_1(s)\) at even \(s\) and \(w_l(s)\le w_b(s)\) at odd \(s\). Since \(v_p(\beta_a)\ge-\eta t\), every argument of \(h_\ell(\mathcal F_b)\) has valuation at least \(\sum_{s\ {\rm even}}w_1(s)-\sum_{s\ {\rm odd}}w_b(s)-\eta t\). Each monomial has degree \(\ell\). Substituting this bound in \eqref{eq:p3c3-output-polynomial} cancels the odd weight sum and gives
\[
 v_p(c_{b,\ell})\ge
 \sum_{s\ {\rm even}}\sum_{a=r-\ell+1}^{r}(\lambda_{s,a}+w_1(s))
 +\eta\ell(N_b-2r+\ell-t).
 \tag{3.38.14}
\]
The double sum is nonnegative by dominance and \(\lambda_{s,k}+w_1(s)\ge0\). Also \(N_b-2r+\ell-t=(k-2r)+(b-1)t+b+\ell\ge0\). At \(\ell=r\), the term at \(a=1\) at each marked even position is at least \(\lambda_{s,1}-\lambda_{s,2}=\alpha_{v,s}\). This proves \eqref{eq:p3c3-prefix-coefficients}.

For \(b=j+1\) and \(\ell<r\), Theorem~\ref{B3.34} gives the bound \(B_v+R_j(v)\) on the Hecke term. Its coefficient is integral and \(U_b\) is a unit. At \(\ell=r\) the Hecke term is the identity and the scalar adds at least \(E_v\) to the bound \(B_v-E_v+R_j(v)\) of Theorem~\ref{B3.33}. All terms therefore lie in the lattice of exponent \(B_v+R_j(v)\).
\end{proof}

\section*{Proofs for Chapter~\ref{ch:chain-induction}: Induction along the chain}

\subsection*{Proof of \ref{B3.39}}\label{prf:B3.39}\phantomsection
\begin{proof}
For the empty list, \eqref{eq:chain-multiple-commutation} is the identity. For \(L=(v,L')\), pass the scalar from the induction hypothesis through \(Z_{v,b}\) and apply \eqref{eq:chain-prefix-relation} to \(Y_{L'}F\in K\Lambda_0\). The two powers of \(p\) multiply to the power with exponent \(\sum_{v\in L}P_v^{\Z}(b)\).

For the bound, expand the product of the \(Z_{v,j+1}\). A summand is a scalar \(a\) times the ordered product of \(\widetilde T_{v,i_v,j+1}\) on \(K_jF\), where \(0\le i_v\le r\) on \(L\) and \(i_v=0\) elsewhere. The inverse powers of \(U_v(j+1)\) are units. By the assumed coefficient bounds, each \(v\in L\) with \(i_v=0\) contributes \(E_v\) to the bound for \(a\), and the other factors are integral. Thus \(\nu_K(a)\ge H\), where \(H=\sum_{v\in L,\ i_v=0}E_v\).

Put \(q'_v=q_v-E_v\) if \(v\in L\) and \(i_v=0\), and \(q'_v=q_v\) otherwise. For \(v\in L\), the inequality \(q_v\le B_v(j+1)\) gives \(q'_v\le B_v(j+1)-E_v\) when \(i_v=0\) and \(q'_v\le B_v(j+1)\) when \(i_v>0\). Outside \(L\), use \(\mathcal L_{v,j}(q_v)\). Thus the \(q'_v\) satisfy (3.29.3) of \ref{B3.29}. Applying the real-exponent argument of Lemma~\ref{B3.31} to the assumed coefficient estimate and coefficient-to-lattice implication gives the operator bound \(\sum_v q'_v-d_j=\sum_v q_v-H-d_j\). Multiplication by \(a\) restores \(H\), since
\[
 \lceil e(x+y)\rceil\le\lceil ex\rceil+\lceil ey\rceil.
\]
The target lattice is closed under addition.
\end{proof}

\subsection*{Proof of \ref{B3.41}}\label{prf:B3.41}\phantomsection
\begin{proof}
At \(b=0\) the assertion is \(F\in\Lambda_0\). Suppose it holds at \(b<m\) for every integral input. Put \(C=C_b\) and let \(D\) be its complement in \(L\). Write \(H_b=\sum_{j<b}\beta_j\). Commutativity gives \(Y_L^{b+1}=Y_CY_L^bY_D\). Since \(Y_DF\in\Lambda_0\), the induction hypothesis gives \(K^{[b]}W=\varpi^{\lceil eH_b\rceil}F'\) with \(W=Y_L^bY_DF\) and \(F'\in\Lambda_b\). Each \(y_v\) preserves \(\Lambda_0\), so \(W\in\Lambda_0\).

Set \(N=\sum_{v\in C}P_v^{\Z}(b+1)\). By \ref{B3.39} and linearity,
\[
 K^{[b+1]}Y_L^{b+1}F
 =p^{-N}Z_{C,b+1}K^{[b+1]}W
 =p^{-N}\varpi^{\lceil eH_b\rceil}Z_{C,b+1}K_bF'.
\]
Use \eqref{eq:chain-multiple-bound} with \(q_v=G^+_v(b+1)\) on \(C\) and \(q_v=G^0_v(b+1)\) elsewhere. Put \(Q=\sum_v q_v-d_b\). Since \(p\) is a unit times \(\varpi^e\), the right-hand side lies in
\[
 \varpi^{\lceil eH_b\rceil+\lceil eQ\rceil-eN}\Lambda_{b+1}.
\]
The exponent is at least \(\lceil e(H_b+Q-N)\rceil=\lceil e(H_b+\beta_b)\rceil\), which gives the bound at \(b+1\).
\end{proof}

\subsection*{Proof of \ref{B3.42}}\label{prf:B3.42}\phantomsection
\begin{proof}
Take a generator in \ref{B3.40} with \(C\subseteq V\) and \(F_0\in\Lambda_0\). Enumerate \(C\) without repetition and apply \ref{B3.41} with
\[
 C_j=\{v\in C:G^0_v(j+1)\le G^+_v(j+1)-P^*_v(j+1)\}.
\]
Let \(g=\sum_{j=0}^{m-1}\beta_j\) be the resulting lower bound for \(K^{[m]}Y_C^mF_0\).
On \(C_j\), the contribution to \(\beta_j\) is at least the maximum in \eqref{eq:chain-final-budget} because \(P_v\le P^*_v\). On \(C\setminus C_j\) and outside \(V\), both contributions equal \(G^0_v(j+1)\). At each \(v\in V\setminus C\), the excess in \eqref{eq:chain-final-budget} over all layers is the sum of \([G^+_v(b)-P^*_v(b)-G^0_v(b)]_+\) over \(b\). The bound on this sum gives
\begin{equation}\label{eq:chain-real-comparison}
 T\le g+\sum_{v\in V\setminus C}S_v.
\end{equation}

Put \(n=\sum_{v\in V\setminus C}\lceil eS_v\rceil\). Since \(e\ge0\), \eqref{eq:chain-real-comparison} gives \(eT\le\lceil eg\rceil+n\). The right side is integral, hence \(\lceil eT\rceil\le\lceil eg\rceil+n\). By \ref{B3.41} and linearity, the lifted generator lies in \(\varpi^{\lceil eg\rceil+n}\Lambda_m\subseteq\varpi^{\lceil eT\rceil}\Lambda_m\). Linearity extends the inclusion to the \(O\)-span.
\end{proof}

\subsection*{Proof of \ref{B3.43}}\label{prf:B3.43}\phantomsection
\begin{proof}
The summands of \(S\) are nonnegative. Dominance gives \(\alpha_s\ge0\). The lower bound on \(w(1,s)\) and monotonicity in \(j\) give \(\lambda_{s,i}+w(j,s)\ge0\) for \(1\le i\le k_1+1\) and \(1\le j\le m\). Thus the capacities in \ref{B3.02} are nonnegative. The positive parts in \ref{B3.13} and the minima in \ref{B3.14} are nonnegative. With \(\delta_j\) the digit loss of \ref{B3.03}, this gives
\[
 g_A(j)\ge-\delta_j,\qquad
 g_B(j)\le\gamma_p\sum_{s\in T}\alpha_s-\delta_j.
\]
The inner sum in the price formula of \ref{B3.37} has \(r\) terms, so
\[
 P(j)=\sum_{s\ {\rm odd}}\sum_{i=k_1+2-r}^{k_1+1}
 (w(j,s)+\lambda_{s,i})+\frac{n'rj(t+1)}2\ge0.
\]
Hence \([g_B(j)-P(j)-g_A(j)]_+\le\gamma_p\sum_{s\in T}\alpha_s\). Sum over \(1\le j\le m\).
\end{proof}

\subsection*{Proof of \ref{B3.45}}\label{prf:B3.45}\phantomsection
\begin{proof}
For each \(v\in\mathcal P\), the local data include a DVR \(W_v\) with maximal ideal \((p)\) and a finite ring map \(j_v:W_v\to O\). Its kernel is contained in \((p)\), and \(W_v/\ker j_v\subseteq O\) is integral. The lying-over theorem of Cohen and Seidenberg says that every prime of the smaller ring in an integral extension is the contraction of a prime of the larger ring \cite{cohen-1946}. Applied to the image of \((p)\) in \(W_v/\ker j_v\), it gives a prime of \(O\) containing \(p\). Thus \(p\) is a nonunit, and \((p)=(\varpi^e)\) forces \(e>0\).

The kernel field \(K_c\) contains \(\operatorname{Frac}W(\overline{\F}_p)\), where \(W\) denotes Witt vectors, so \(p\ne0\) in \(K_c\) and in \(K\). For \(e>0\), the map \(\iota\) is injective, since a nonzero kernel would contain a power of \(\varpi\) and therefore kill \(p\). The valuation ring \(O_c\subseteq K_c\) contains the image of \(z\) or \(z^{-1}\) for every \(z\in K^\times\). By Definition~\ref{B3.12}, its inverse image in \(K\) is \(\iota(O)\). Thus each \(z\in K^\times\) belongs to \(\iota(O)\) or is the inverse of a nonzero element of \(\iota(O)\), and \(K=\operatorname{Frac}(O)\) via \(\iota\).

The cited results give the first nine assertions. In Theorem~\ref{B3.38}, the required inclusion \(K_j(K\Lambda_j)\subseteq K\Lambda_{j+1}\) for \(0\le j<m\) is part of Definition~\ref{B3.05a}.

By Definition~\ref{B3.05a}, the input operators \(y_v\) commute pairwise on \(M_0\), and each maps \(\Lambda_0\) into itself. For \(e>0\), apply \ref{B3.42} with the sequences from \ref{B3.38}. Use the joint coefficient estimate of \ref{B3.29} and the coefficient-to-lattice implication of Lemma~\ref{B3.30}.

If \(e=0\), then \(\mathcal P=\varnothing\) and \(I(S,\varnothing)=\Lambda_0\). Since \(\lceil eh\rceil=0\) for every \(h\), the basic lift bound of \ref{B3.32} gives \(K_j\Lambda_j\subseteq\Lambda_{j+1}\) at each step. Thus \(K^{[m]}\Lambda_0\subseteq\Lambda_m\), which is \eqref{eq:chain-final-bound} at \(e=0\).

The local bounds of \ref{B3.27} and the external bound \eqref{eq:chain-final-bound} give \ref{B3.44}. Lemma~\ref{B3.43} gives the asserted upper bound for the enhancement budget. Apply the system assertion at \(p=\operatorname{char}k(w)\) to the given system at \((u,w)\).
\end{proof}

\section*{Proofs for Chapter~\ref{ch:global-frobenius-orbit-parameters}: Global Frobenius-orbit parameters}

\subsection*{Proof of \ref{B3.50}}\label{prf:B3.50}\phantomsection
\begin{proof}
Permission gives the inert class and \(n_p=\operatorname{ord}(h)\). The map \(\beta_h\) is constant on each \(h\)-coset and takes its value there. Its fibre over \(t_0\) is therefore the image of \(\iota_{t_0}\) in Definition~\ref{def:p3c5-orbit-model}. Write \(w_j(i)\) for the local shift, \(c_{\rm loc}\le c\) for the local row bound, and \(\lambda_{i,\ell}\) for the base weights, numbered from one. The filling identification of Definition~\ref{def:p3c5-orbit-model} gives the branch rank \(r\) and
\[
 z_e(\iota(i))=(0,\iota((i+e)\bmod n_p)),\qquad
 w_j(i)=B_j(0,\iota(i)).
 \tag{3a}\label{eq:p3c5-tied-orbit}
\]
Local marks correspond to \(T(t_0)\) and odd marks to \(T(t_0)\setminus T^V(t_0)\). By \ref{B3.37},
\[
 P_v(j)=\sum_{\substack{0\le i<n_p\\ i\ {\rm odd}}}
 \left(rw_j(i)+\sum_{\ell=k-r+1}^{k}\lambda_{i,\ell}\right)
 +\frac{n_prj(t+1)}2. \tag{3}\label{eq:p3c5-exchange}
\]
For \(r=1\), the marked summand at \(s=\iota(i)\) is \(B_j(0,s)-b_s\). At unmarked positions \(w_j(i)\le w''+\eta\) by \ref{B3.47} and \(\lambda_{i,k}\le c\) by \ref{B3.01}. Since \(j(t+1)/2\le N_m\), this gives the large-prime bound.

In \eqref{eq:p3c5-size-clauses}, \(x^\theta(\log x)\Lambda/k\ge0\), so \(4R\le x^{-2}\le1\). With \eqref{eq:p3c5-top-window} this gives
\[
 b'\le b_s\le\frac54 b'\qquad(s\in TT).
 \tag{3b}\label{eq:p3c5-bottom-window}
\]
For \(r=r_G<k\), the last \(r\) rows omit the first and each has weight at most \(c_{\rm loc}\). At a mark, \(w_j(i)\le b_s+b'/x+\eta\le2b'+\eta\) by \eqref{eq:p3c5-bottom-window} and \(x\ge2\). The unmarked shift is at most \(w''+\eta\). Thus \eqref{eq:p3c5-exchange} is at most \(r_G(2N_v(t_0)b'+n_pM_{\mathrm{Pr}})\), which gives the small-prime bound.
\end{proof}

\subsection*{Proof of \ref{B3.52b}}\label{prf:B3.52b}\phantomsection
\begin{proof}
Initial-column minors are invariant under the specified column operations by alternation. Row substitutions commute with these operations, so \(D_\lambda\subseteq A_\lambda\).

Order monomials lexicographically by column and then by row, with earlier variables first. For a nonzero \(P\in A_\lambda\), let \(m_{ij}\) be the exponents of its leading monomial. Homogeneity gives \(\sum_i m_{ij}=\nu_j\). For adjacent columns \(j,j+1\), write \(x_i=Y_{i,j}\), \(y_i=Y_{i,j+1}\). Retain the component with the leading monomial's total degree in each pair \((x_i,y_i)\) and its degrees in each other variable. This component is invariant under \(y\mapsto y+tx\), which preserves those degrees.

In that component successively take the coefficient of \(x_a^{K}y_a^e\) in row \(a\), where \(K\) is the largest exponent of \(x_a\) and \(K+e\) is the fixed pair degree. At this step the polynomial has the form \(\sum_{i=0}^K x_a^iy_a^{e+K-i}P_i\), with \(P_K\ne0\) and \(P_i\) independent of \(x_a,y_a\). After \(y\mapsto y+tx\), factor out \((y_a+tx_a)^e\). If the transformed polynomial has \(t\)-degree at most \(s\), then \(e\le s\). The remaining factor has \(t\)-degree at most \(s-e\) and \(y_a\)-degree at most \(K\). Evaluating it at \(x_a=1,\ y_a=-t\) gives the transformed \(P_K\), of \(t\)-degree at most \(s-e+K\). The initial bound is \(s=0\). The successive pairs \((K,e)\) are \((m_{a,j},m_{a,j+1})\), since the degrees outside the two columns are fixed. Iterating the degree bound gives
\[
 \sum_{i\le a}m_{i,j+1}\le\sum_{i<a}m_{i,j}.
\]
Make a tableau of shape \(\nu\) by placing \(m_{ij}\) copies of \(i\) in row \(j\) in increasing order. The inequality ensures strict increase down each column. For each tableau column take the minor on its indicated rows and the first matrix columns. Its diagonal term is leading with coefficient one, so the product has the same leading monomial as \(P\). Subtract the leading coefficient of \(P\) times this product and induct on the leading monomial. This gives spanning over \(O\), and the distinct leading monomials give independence.

The resulting basis is the straight-shape case of the standard basis of Akin, Buchsbaum and Weyman \cite[Theorem II.2.16]{akin-1982}, with shape \(\nu^{\mathsf t}\) in their notation, where a single row gives an exterior power. The normalized pairing of \ref{B3.10a} gives \(D_\lambda\) the same rank as this Schur lattice. Thus \(V=FA_\lambda=FD_\lambda\).
\leanref{Formalism.\allowbreak{}GradedWeylDivisibility.\allowbreak{}leadInterlaced}{Formalism.\allowbreak{}GradedWeylDivisibility.\allowbreak{}FinalC}{36--36}
\leanref{Formalism.\allowbreak{}GradedWeylDivisibility.\allowbreak{}semistandard\_\allowbreak{}basis}{Formalism.\allowbreak{}GradedWeylDivisibility.\allowbreak{}FockSpaceTableaux.\allowbreak{}SSBasis}{10--67}
\leanref{Formalism.\allowbreak{}L2.\allowbreak{}G.\allowbreak{}Model.\allowbreak{}sw\_\allowbreak{}weyl\_\allowbreak{}eq}{Formalism.\allowbreak{}LocalToGlobal.\allowbreak{}GroupLattices.\allowbreak{}FockModels.\allowbreak{}SwRowGenMonomial}{177--198}

For \(i\ne j\) let \(e_{ij}\) be the first coefficient of the substitution adding row \(i\) to row \(j\), and let \(e_{ii}\) be the degree operator in row \(i\). The coefficient of order \(a\) is \(e_{ij}^{(a)}=e_{ij}^a/a!\). These operators preserve \(A_\lambda\) by column invariance and \(D_\lambda\) by its generating words. Each tableau product uses \(l\) minors. A given row occurs at most once in each minor. Its vector of row degrees therefore has coordinates in \([0,l]\). Row weight \(\nu\) occurs on the line \(Fh_\lambda\). Row \(0\) must occur in every minor and row \(1\) in every minor of size at least two. Each subsequent row \(i\) must occur in every minor of size at least \(i+1\). Let \(\xi\) project to row weight \(\nu\) and evaluate at the identity matrix. Then \(\xi(A_\lambda)\subseteq O\) and \(\xi(h_\lambda)=1\).

Put \(\mathscr H=\{T\in\operatorname{End}_F(V):T(A_\lambda)\subseteq D_\lambda\}\) and \(u(y)=\xi(y)h_\lambda\). Then \(u\in\mathscr H\) has weight zero and trace one. On endomorphisms use \(E_{ij}T=e_{ij}T-Te_{ij}\). The off-diagonal divided powers preserve \(\mathscr H\) by
\[
 E_{ij}^{(a)}T=\sum_{b=0}^{a}(-1)^b e_{ij}^{(a-b)}T e_{ij}^{(b)}.
\]
Endomorphism weights are differences of row weights of \(V\), so each coordinate lies in
\[
 [-l,l].
 \tag{11c}\label{eq:p3c5-adjoint-range}
\]

For a row index \(0\le c'<k\) and \(S\subseteq\{0,\ldots,c'-1\}\), define operators on \(\operatorname{End}_F(V)\) by
\[
 W_0(S)=1,\qquad W_{q+1}(S)=\sum_{i\in S} E_{c'i}W_q(S)E_{ic'} .
\]
For \(S_{c'}=\{0,\ldots,c'-1\}\) put
\[
 P_{c'}=\sum_{q=0}^l b_qW_q(S_{c'}),\qquad
 b_q=\frac{(-1)^q c'!}{(q+c')!q!}.
\]
Suppose \(z\) has weight zero and \(E_{ab}z=0\) for \(a,b<c'\). Write \(W_q=W_q(S_{c'})\). Each \(W_q\) preserves weight and commutes with \(E_{ab}\) for \(a,b<c'\). In the recursion the commutators \([E_{ab},E_{c'i}]=-\delta_{ia}E_{c'b}\) and \([E_{ab},E_{ic'}]=\delta_{bi}E_{ac'}\) cancel by commutation with \(W_q\) at the preceding step.

For \(j<c'\), let \(y\) be a product of \(s\) operators \(E_{ic'}\) with \(i<c'\) applied to \(z\). Then
\[
 E_{c'c'}y=-sy,\qquad
 \sum_{i<c'}E_{ji}E_{ic'}y=(s+c')E_{jc'}y .
 \tag{11d}\label{eq:p3c5-column-identities}
\]
For the second identity at \(s=0\), commute \(E_{ji}\) to the right and use \(E_{ji}z=0\). Inserting \(E_{bc'}\) adds one copy of \(E_{jc'}E_{bc'}y\), since the column operators commute and \([E_{ji},E_{bc'}]=\delta_{ib}E_{jc'}\). Commuting \(E_{jc'}\) through the recursion gives
\[
 E_{jc'}W_qy
 =W_q E_{jc'}y+q(2s+c'+q)W_{q-1}E_{jc'}y ,
 \tag{11e}\label{eq:p3c5-projector-commutator}
\]
with zero second term at \(q=0\). In the step from \(q\) to \(q+1\), the recursion contributes \(q(2s+2+c'+q)\). The commutator \([E_{jc'},E_{c'i}]=E_{ji}-\delta_{ij}E_{c'c'}\) contributes \((s+c')+(s+1)\) by \eqref{eq:p3c5-column-identities}. The sum is \((q+1)(2s+c'+q+1)\).

At \(s=0\) the terms of \(E_{jc'}P_{c'}z\) cancel because \(b_{q+1}(q+1)(c'+q+1)=-b_q\). The terminal term \(W_lE_{jc'}z\) vanishes by \eqref{eq:p3c5-adjoint-range}, since its rightmost \(l+1\) column operators would give coordinate \(-l-1\) at \(c'\). For \(v=P_{c'}z\), put \(e=E_{jc'}\) and \(f_-=E_{c'j}\). Then \(ev=[e,f_-]v=0\), and
\[
 ef_-^a v=-a(a-1)f_-^{a-1}v\qquad(a\ge1).
\]
If \(f_-v\ne0\), take the last nonzero \(f_-^av\), which exists by \eqref{eq:p3c5-adjoint-range}. The identity at \(a+1\) gives a contradiction. Thus the new row operators also kill \(v\), and its diagonal operators vanish by weight zero. For \(q>0\), each summand of \(W_qz\) has trace zero by its outer commutator \(E_{c'i}\). Hence \(P_{c'}\) preserves trace.
\leanref{Formalism.\allowbreak{}MainNotions.\allowbreak{}G.\allowbreak{}GLWeightTheory.\allowbreak{}raising\_\allowbreak{}casimirPower\_\allowbreak{}commutator}{Formalism.\allowbreak{}SharpRestriction.\allowbreak{}GrowthBounds.\allowbreak{}RootWeights.\allowbreak{}CommutatorTripleProductLeibniz}{159--201}
\leanref{Formalism.\allowbreak{}MainNotions.\allowbreak{}G.\allowbreak{}GLWeightTheory.\allowbreak{}extremal\_\allowbreak{}projector\_\allowbreak{}properties}{Formalism.\allowbreak{}SharpRestriction.\allowbreak{}GrowthBounds.\allowbreak{}RootWeights.\allowbreak{}CommutatorTripleProductLeibniz}{259--303}

Put \(d=\lfloor\log_p(l+k-1)\rfloor\) and \(C_{q,S}=c'!W_q(S)/((q+c')!q!)\). The operations \(U\mapsto E_{c'i}UE_{ic'}\) commute for distinct \(i<c'\). Expand the \(q\)-th power after adjoining \(i\notin S\) and convert both outside powers to divided powers. This gives
\[
 C_{q,S\cup\{i\}}
 =\sum_{a=0}^q \binom{q+c'}{a}^{-1}
 E_{c'i}^{(a)} C_{q-a,S} E_{ic'}^{(a)}.
 \tag{11f}\label{eq:p3c5-projector-denom}
\]
The scalar is \(\binom qa(a!)^2(q-a+c')!(q-a)!/((q+c')!q!)\), which equals \(\binom{q+c'}a^{-1}\). For \(0\le a\le N=q+c'\) and \(0\le q\le l\),
\[
 v_p\binom Na=\sum_{j\ge1}
 \left(\left\lfloor\frac{N}{p^j}\right\rfloor-
 \left\lfloor\frac{a}{p^j}\right\rfloor-
 \left\lfloor\frac{N-a}{p^j}\right\rfloor\right)\le d .
 \tag{11g}\label{eq:p3c5-binomial-value}
\]
Each term is zero or one and vanishes for \(p^j>N\). The prime-to-\(p\) part of the binomial coefficient is a unit in \(O\), so \(p^d/\binom Na\in O\). For \(S=\varnothing\), \(C_{0,S}=1\) and \(C_{q,S}=0\) for \(q>0\). Induction on \(|S|\) in \eqref{eq:p3c5-projector-denom}, using divided-power stability, gives
\[
 p^{|S|d}C_{q,S}\mathscr H\subseteq\mathscr H
 \quad(0\le q\le l),\qquad
 p^{c'd}P_{c'}\mathscr H\subseteq\mathscr H .
 \tag{11h}\label{eq:p3c5-projector-integral}
\]
\leanref{Formalism.\allowbreak{}MainNotions.\allowbreak{}G.\allowbreak{}GLWeightTheory.\allowbreak{}casimirPowerScaled\_\allowbreak{}insert\_\allowbreak{}dividedPower}{Formalism.\allowbreak{}SharpRestriction.\allowbreak{}GrowthBounds.\allowbreak{}RootWeights.\allowbreak{}CasimirPowerScaledInsertDividedPower}{86--99}
\leanref{Formalism.\allowbreak{}MainNotions.\allowbreak{}G.\allowbreak{}GLWeightTheory.\allowbreak{}extremalProjector\_\allowbreak{}denominator\_\allowbreak{}bound}{Formalism.\allowbreak{}SharpRestriction.\allowbreak{}GrowthBounds.\allowbreak{}RootWeights.\allowbreak{}CasimirPowerScaledInsertDividedPower}{217--235}

Apply \(P_0,\ldots,P_{k-1}\) to \(u\) in order. The result commutes with every \(e_{ij}\). It acts by a scalar on \(Fh_\lambda\) and by the same scalar on its divided-power words. These span \(V\), so trace one makes the result \(1_V/\dim_F V\). By \eqref{eq:p3c5-projector-integral}, multiplication by \(p^{d\sum_{c'=0}^{k-1}c'}\) puts it in \(\mathscr H\). Multiplying by \(\dim_F V\) gives \(p^{d_p(\lambda)}1_V\in\mathscr H\), which proves \eqref{eq:p3c5-factor-index}.
\leanref{Formalism.\allowbreak{}MainNotions.\allowbreak{}G.\allowbreak{}GLWeightTheory.\allowbreak{}restA\_\allowbreak{}of\_\allowbreak{}linf\_\allowbreak{}integrality}{Formalism.\allowbreak{}SharpRestriction.\allowbreak{}GrowthBounds.\allowbreak{}RootWeights.\allowbreak{}LinfAdjointLeEllWtHighest}{51--88}
\leanref{Formalism.\allowbreak{}L2.\allowbreak{}G.\allowbreak{}Model.\allowbreak{}sandwichO}{Formalism.\allowbreak{}LocalToGlobal.\allowbreak{}GroupLattices.\allowbreak{}FockModels.\allowbreak{}Sandwich}{125--145}

For a frame matrix \(g\), use \(Y\mapsto g^{-1}Y\) and the twist \((\det g)^{-\lambda_{k-1}}\) as in \ref{B3.10a}. Both commute with scalar multiplication. The Weyl lattice is also independent of the frame. Over the DVR, integral invertible matrices are generated by elementary shears and diagonal units. The shears preserve \(D_\lambda\) by divided-power stability, and diagonal units preserve its row-homogeneous generators. Thus the framed inclusions apply to the intersections defining the two factors. Negative reversal at an odd position preserves \(l\).

Use \(R=2n\delta_1x^\varepsilon\Lambda/k\) as in \ref{B3.50}. The parameter choice gives, for \(s\in TT\),
\[
 0<4S\le\Lambda b'/k,\qquad
 4R+C_\rho(\log x)\Lambda/k\le x^{-2},\qquad b_s\le(1+R)b',
\]
where \(C_\rho>0\). The first inequality gives \(\Lambda>0\), so the second gives \(R\le x^{-2}/4\le1/4\). Hence at a top position the spread \(a_s+b_s\) is bounded using \eqref{eq:p3c5-top-window} and \(b_s\le5b'/4\).

At an off-top embedding \(\tau\), let \((u_i)\) and \((v_i)\) be the upper and lower interlacing rows of its Harish--Chandra parameter. The rank shift gives \(l=u_0-u_{k-1}-(k-1)\). The cloud and extreme-coordinate clauses give \(u_0-v_0\le2S\), \(0\le v_i\le S\), and \(u_{k-1}=v_{k-2}-1/2\). Thus
\[
 l\le
 \begin{cases}
 \frac65 a_\ast+\frac54 b',&\text{at a top position},\\
 3S+\frac12-(k-1),&\text{at an off-top position},
 \end{cases}
 \qquad l\le3a_\ast,
 \tag{11i}\label{eq:p3c5-weight-spread}
\]
where \(kS\le w''\le b'\le a_\ast\) and \(k\ge5\). Since \(\Delta\ge N_mm^{10}\) and \(S=4k\Delta\), we have \(S\ge1\) and \(k\le kS\le w''\le a_\ast\). Hence \(l+k-1\le4a_\ast-1\le a_\ast^2\), and
\[
 d_p(\lambda)\le\binom{k}{2}\,2\log_p a_\ast\le k^2\log_p a_\ast .
 \tag{11j}\label{eq:p3c5-factor-exponent}
\]
\leanref{Formalism.\allowbreak{}BasicNotions.\allowbreak{}ParameterChoice'.\allowbreak{}bT\_\allowbreak{}le'}{Formalism.\allowbreak{}LocalRingCardinality.\allowbreak{}ColumnWeightSpreadOffTop}{85--102}
\leanref{Formalism.\allowbreak{}BasicNotions.\allowbreak{}ParameterChoice'.\allowbreak{}k\_\allowbreak{}le\_\allowbreak{}a}{Formalism.\allowbreak{}LocalRingCardinality.\allowbreak{}ColumnWeightSpreadOffTop}{77--81}
\leanref{Formalism.\allowbreak{}BasicNotions.\allowbreak{}columnWeight\_\allowbreak{}spread\_\allowbreak{}le}{Formalism.\allowbreak{}LocalRingCardinality.\allowbreak{}ColumnWeightSpreadOffTop}{182--195}
\leanref{Formalism.\allowbreak{}BasicNotions.\allowbreak{}gat\_\allowbreak{}ellWt\_\allowbreak{}lamAt\_\allowbreak{}le}{Formalism.\allowbreak{}LocalRingCardinality.\allowbreak{}ColumnWeightSpreadOffTop}{241--247}

At a full lattice tuple the fibres are spanned by products in disjoint variables over the pairs \((v,s)\). Apply \eqref{eq:p3c5-factor-index} to each factor. By \eqref{eq:p3c5-factor-exponent} the sum of the exponents over the \(f\) pairs is at most \(fk^2\log_p a_\ast\le e_p\). The remaining integral power of \(p\) preserves the Weyl fibre. The inclusion \(D_\lambda\subseteq A_\lambda\) gives the reverse fibre inclusion. Both function lattices use the support and equivariance conditions of \ref{B3.01}, so the fibre inclusions give \eqref{eq:p3c5-total-index} pointwise.
\end{proof}

\subsection*{Proof of \ref{B3.53}}\label{prf:B3.53}\phantomsection
\begin{proof}
By \ref{B3.02}, at a mark \(\alpha_i=a_{\iota(i)}-\lambda_{i,2}\) and \(C_i=w_j(i)+\lambda_{i,2}\). Elsewhere \(\alpha_i=0\) and \(C_i=w_j(i)+\lambda_{i,1}\). These row weights have absolute value at most \(c_{\rm loc}\le c\), and \(w_j(i)\ge w_1(i)\ge c_{\rm loc}\) by \ref{B3.01}. Thus \eqref{eq:p3c5-tied-orbit} gives, with local indices modulo \(n_p\),
\[
 \bar\alpha_{\iota(i)}\le\alpha_i\quad(i\text{ marked}),\qquad
 0\le C_{i+e}\le\bar C_{\iota(i),e}\quad(e\ge0).
 \tag{13}\label{eq:p3c5-capacity}
\]
Even marks correspond to \(T^V(t_0)\). The local exclusions \(D_i\) of \ref{B3.13} are the least positive distances to opposite-parity marks. The cycle bijection and \eqref{eq:p3c5-tied-orbit} give
\[
 e\in D_i,\quad e\text{ odd}\quad\Longrightarrow\quad
 e\in\mathcal E_p(\iota(i)).
 \tag{14}\label{eq:p3c5-exclusions}
\]

Put \(A=2a_\ast\) and \(z=\log_p A=z_p\). The rank conditions give \(m\le k\), \(t\ge1\), and \(fk^2\ge3\), so \(N_m\le k(t+2)\le3kt\le t^3fk^3\). By \eqref{eq:p3c5-size-clauses},
\[
 N_m\le b\le a_\ast,\qquad c\le4w''/5,\qquad \eta\le w''/2 .
 \tag{15a}\label{eq:p3c5-size-aux}
\]
At a mark of the represented first-sheet orbit the shift is at most \(5b'/4+b'/x+\eta\le11b'/4\) by \ref{B3.47} and \eqref{eq:p3c5-bottom-window}. At the other positions of that orbit it is at most \(w''+\eta\le3b'/2\). After removing the single-box columns, each row has at most \(w_j(i)+c_{\rm loc}\) boxes. By \eqref{eq:p3c5-top-window}, \eqref{eq:p3c5-size-aux}, and \(b'\le a_\ast\),
\[
 \alpha_i\le A,\qquad
 L_i\le N_j(w_j(i)+c_{\rm loc})\le a_\ast(4a_\ast)=A^2,\qquad
 N_{j-1},N_j\le N_m .
 \tag{15}\label{eq:p3c5-local-size}
\]
The digit lengths in \ref{B3.03} satisfy \(\ell_p(\alpha_i)\le1+z\le2z\) and \(\ell_p(L_i)\le1+2z\le3z\). Hence
\[
 \delta_v\le n_p(2N_mz+3N_m^2z)+4N_v(t_0)z^2
 \le5n_pN_m^2z+4N_v(t_0)z^2\le\bar\delta(t_0).
 \tag{16}\label{eq:p3c5-digit-error}
\]
By \eqref{eq:p3c5-capacity} and \eqref{eq:p3c5-summable}, the positive-part series are summable and each global summand is bounded by its local counterpart. By \eqref{eq:p3c5-exclusions} each retained ordinary summand is retained locally. For the enhanced gain, the capacities in \eqref{eq:p3c5-capacity} permit \eqref{eq:p3c5-enhanced}. Compare its positive parts at odd marks with those of \ref{B3.14}, and at even marks use \(\gamma_p\bar\alpha_{\iota(i)}\le\gamma_p\alpha_i\). Subtract the errors using \eqref{eq:p3c5-digit-error}. Proposition~\ref{B3.50} gives the price bound.
\end{proof}

\subsection*{Proof of \ref{B3.54}}\label{prf:B3.54}\phantomsection
\begin{proof}
The logarithm bound gives \(p>2a_\ast\).  The rank and parameter conditions give \(m\ge2\), \(m+3\le k\), \(N_mm^{10}\le\Delta\), and \(S=4k\Delta\). In particular,
\[
 S\ge100,\qquad \kappa_j\le N_m\le S,\qquad
 kS\le w''\le b'/x\le b'\le a_\ast/2 .
 \tag{17a}\label{eq:p3c5-large-scales}
\]
The fillings are nonnegative by \ref{B3.47} and \eqref{eq:p3c5-top-window}, so every capacity is at least \(4S-\kappa_j\ge3S\). Since \(a_\ast\ge2kS\ge2c\) and \(z_p\ge0\),
\[
 0\le\bar\alpha_s\le6a_\ast/5\quad(s\in T(t_0)),\qquad
 \bar C_{s,e}\ge1,\qquad \bar\delta(t_0)\ge0 .
 \tag{18}\label{eq:p3c5-large-capacity}
\]
For \(e\ge1\), \(\bar\alpha_s/p^e\le6a_\ast/(5p)<1\), so \(g_0=-\bar\delta\le0\). The minima in \ref{B3.49} are nonnegative. As \(a_\ast<p/2\) and \(3p^2/5\le p^2-1\) for \(p\ge2\),
\[
 g_U(t_0,j)\le\gamma_p\sum_{s\in T(t_0)}\bar\alpha_s\le N_v(t_0).
 \tag{19}\label{eq:p3c5-large-enhancement}
\]
The coset has at most \(n_p\) elements by \ref{B3.46}. At a mark \(B_j(0,s)-b_s\ge b'/x-S\ge0\), and \(K_{\mathrm{Pr}},M_{\mathrm{Pr}}\ge1\). Thus the large-prime branch of \eqref{eq:p3c5-price} is at least \(n_p\ge N_v(t_0)\). The small-prime branch is at least \(N_v(t_0)\) by \(r_G\ge1\) and \(2b'+c\ge1\). This proves \eqref{eq:p3c5-large-residue}, and \eqref{eq:p3c5-gain} gives zero counted gain.
\end{proof}

\subsection*{Proof of \ref{B3.55}}\label{prf:B3.55}\phantomsection
\begin{proof}
At step \(j-1\) and position \(v\), the ordinary and zero bounds give the coefficient estimate of \ref{B3.12} with exponent \(\max\{g_{A,v}(j),0\}\), hence with every smaller exponent \(q\).
\leanref{Formalism.\allowbreak{}MainNotions.\allowbreak{}FormalGroupSystem.\allowbreak{}lowAt\_\allowbreak{}max}{Formalism.\allowbreak{}SharpRestriction.\allowbreak{}GrowthBounds.\allowbreak{}FormalGroupSystem}{718--725}

Put \(B=v(V')\). In \ref{B3.42} set the gain and price inputs \(G^+,P^*\) to zero off \(B\). If \(z_p\ge1\), use \((G^+,P^*)=(g_U,\Pi_p)\) on \(B\) and the ordinary input \(G^0=[g_0]_+\) on \(v(\mathcal R_p)\), with the corresponding orbit and layer arguments. By \eqref{eq:p3c5-small-residue}, these meet the gain and price bounds and \(G^0\le\max\{g_{A,v},0\}\). If \(z_p<1\), use \((G^+,P^*)=(\min\{g_{B,v},P_v\},P_v)\) on \(B\) and \(G^0=0\) on \(v(\mathcal R_p)\). The gain and price bounds hold, and the represented positions in \eqref{eq:p3c5-local-final} contribute zero by \eqref{eq:p3c5-large-residue}. Off \(v(\mathcal R_p)\) take \(G^0(v,j)=F_n(v,j)\) in both cases. This gives the required coefficient bound at every position.

Set \(S(v)=\sum_{j=1}^m[G^+(v,j)-P^*(v,j)-G^0(v,j)]_+\). Injectivity on \(\mathcal R_p\) makes \(y_{V'}F\) a generator of \(I(S,B)\) in \ref{B3.40}, with all positions of \(B\) used. Theorem~\ref{B3.42} gives the layer contribution \(\max\{G^+-P^*,G^0\}\) on \(B\) and \(G^0\) elsewhere. On \(v(\mathcal R_p)\) these are the two sums of \eqref{eq:p3c5-local-final}. The complementary positions and the stabilizer subtraction give at least \(\mathrm{low}\) by \eqref{eq:p3c5-complement}. Multiplication by \(e_{\mathcal O}\) and taking ceilings preserve the bound. Splitting \(\mathcal R_p\setminus V'\) into \(V_s\setminus V'\) and \(\mathcal R_p\setminus V_s\) gives the two ordinary sums of \ref{B4.09}.
\end{proof}

\subsection*{Proof of \ref{B3.56}}\label{prf:B3.56}\phantomsection
\begin{proof}
Fix the twist, the comparison and the permitted place. The first clause gives a system on precisely this filling, with the chosen levels and chain. It gives coefficient bounds on every unrepresented position and their sum satisfies \eqref{eq:p3c5-complement} after subtraction of the stabilizer losses. Apply \ref{B3.45} to this admissible system and its kernel. It supplies the ordinary coefficient bound and the conditional chain inequality. The zero bound at every position is part of the admissibility conditions of \ref{B3.05a}. The coefficient ring embeds in its fraction field. These are the local hypotheses of \ref{B3.55}. The two numerical hypotheses in the statement cover its two ranges of $p$, so that theorem gives the asserted inequality. Splitting its ordinary sum at the chosen subset gives the form in \ref{B4.09}.
\end{proof}

\section*{Proofs for Chapter~\ref{ch:layer-records}: What each layer records}

\subsection*{Proof of \ref{B4.01d}}\label{prf:B4.01d}\phantomsection
\begin{proof}
By \ref{B3.46}, the period is $\operatorname{ord}(h_p)$ off $\mathcal S$
and $2\operatorname{ord}(2h_p)$ on $\mathcal S$, where a return takes an
even step. In the filling of \ref{B3.47} put $b'=xb/X'$.
At a top position $t$, the weight is at least $b_t+b'/x$, with $b_t\ge b'>0$. The lower bounds at thick and thin positions are $b$ and $b'/(x|A|)$, respectively. Since $k\ge m+3$, each orbit capacity $(k-1)w_j(z)$
is nonnegative. The feasible set is nonempty because it contains zero. Its defining inequalities are closed and $0\le x_{s,i}\le p^{-i}d_s$ bounds each of its finitely many coordinates. The continuous total flow therefore attains its maximum, including when the index set is empty. Replacing $z$ by $P_p^az$ replaces $d_s$ by $d_{s+a}$. The permutation $x'_{s,i}=x_{s+a,i}$ carries the old feasible set onto the new one. It preserves total flow, and the permutation of the capacities preserves their sum.
\end{proof}

\subsection*{Proof of \ref{B4.01a}}\label{prf:B4.01a}\phantomsection
\begin{proof}
Write $Lx$ for the vector of loads in \ref{B4.01} and set
$|x|=\sum_{s,i}x_{s,i}$. Let $M$ be the maximum at capacity $d$, attained
by Lemma~\ref{B4.01d}. Apply finite-dimensional linear programming duality to $Lx\le d$, $-x\le0$, with objective $\mathbf1^{\mathsf T}x$. The primal is feasible and has finite maximum $M$. The dual has nonnegative multipliers $\tilde y,z$ with $L^{\mathsf T}\tilde y-z=\mathbf1$ and minimum $d\cdot\tilde y=M$ \cite[Cor.~7.1g]{schrijver-1986}. Eliminating $z$ gives precisely $L^{\mathsf T}\tilde y\ge\mathbf1$.
Set $y_s=\min(1,\tilde y_s)$. Each edge inequality involving a capped
coordinate retains a summand at least one. Thus $y$ is feasible, and
$d\cdot y\le M$ since $d\ge0$. For any feasible $x,y$,
\[
 d\cdot y-|x|
 =\sum_s y_s(d_s-(Lx)_s)
 +\sum_{s,i}x_{s,i}(p^iy_s+y_{s+i}-1)\ge0.
\]
This proves equality of the optima. At equality every summand vanishes.
If $y_s\ge2/p$, capacity at $s$ is tight and $p^iy_s\ge2$ on each outgoing
edge. Hence $d_s=\sum_i x_{s-i,i}$. Summing over such vertices gives at
most $|x|$, and summing all loads gives $p|x|\le\sum_s d_s$.

For the comparison, take a dual minimizer $y$ at $W$ and set
$S=\sum_s w_s$. Then $F(W)=\sum_s W_s(\gamma-y_s)$. With
$(a)_+=\max(a,0)$, testing $y$ at $W+e$ gives
\[
 F(W+e)\ge F(W)-\sum_s e_s(y_s-\gamma)_+\ge F(W)-8S/p.
\]
Indeed, $(y_s-\gamma)_+\le1/p$ where $y_s<2/p$, since $\gamma\ge1/p$.
On the other vertices it is at most one, and the demand estimate gives
a total $w_s$ of at most $S/p$ there. The capacities $D$ are nonnegative,
and every flow at $D$ is feasible at $W+e$. Hence
\[
 F(W+e)\le F(D)+2\gamma\sum_{s\in T}w_s\le F(D)+3S/p,
\]
since $2p^2\le3(p^2-1)$.
\end{proof}

\subsection*{Proof of \ref{B4.01c}}\label{prf:B4.01c}\phantomsection
\begin{proof}
Solve $L(g)=Z$ by total degree from
\[
 g_i=Z_i-\sum_{h\ge1}p^{-h}g_{P^hi}^{p^{2h}}.
\]
In each degree the sum uses finitely many $h$ and smaller-degree
coefficients. The identity linear part gives a two-sided inverse.
Inductively the contribution at $h$ has weight at least
$p^{2h}\rho-h\ge\rho$, since $(p^{2h}-1)/(p^2-1)\ge h$.

Use the shift in each copy of the letters. For either coboundary
$B=\Delta_X\beta,\Delta_Y\beta$, put $V=(\exp B)\circ L$.
The integral series $V-1$ has zero constant term, so
$(V-1)\circ g=\exp B-1$ has weight at least $\rho$ in each term.
Splitting a letter commutes with simultaneous torus scaling on its
copies. Hence an exponent $n$ with nonzero coefficient in $\exp B-1$
satisfies
\[
 p^r-1\ \bigm|\ \sum_i n_i(-1)^{s_i}p^{s_i},
\]
where $s_i$ is the letter of variable $i$. Because $r$ is even, reduction
modulo $p+1$ gives $(p+1)\mid |n|$. The coefficients of $\exp B$ lie in
$K_0$, so the positive weights of $\exp B-1$ lie in $q_*\Z$ and are at
least $q_*$. In the logarithm of $\exp B$, the part of order $h\ge1$
has weight at least $hq_*-v_p(h)\ge q_*$, since
$h\ge p^{v_p(h)}\ge1+(p-1)v_p(h)$. Thus every term of $B$ has weight
at least $q_*$.

Split a parity block of a monomial $cZ^n$ between two copies so that each has positive total degree. Assign degree $k_i$ of each variable to
the first copy. Its coboundary coefficient is
$c\prod_i\binom{n_i}{k_i}$ since the copy exponents recover $n$ and the
subtracted series have zero split coefficient. If the block uses at
least two variables, separating whole variables gives factor one. For
a single exponent $m$ with base-$p$ digits $m_h$, use
\[
 (1+T)^m\equiv\prod_h(1+T^{p^h})^{m_h}\pmod p.
\]
If $m$ differs from every power of $p$, choose $m_h\ne0$ with $p^h<m$.
The factor $\binom{m}{p^h}$ has residue $m_h$. Both splits preserve
degree and valuation, so the monomial has weight at least $q_*$.

In the remaining case each block consists of one variable to a power
of $p$. The bidegree condition gives a power $p^h$ with $h>0$.
Splitting off $p^{h-1}$ costs one in valuation because
$\binom{p^h}{p^{h-1}}=p\binom{p^h-1}{p^{h-1}-1}$ and the second
factor is a unit by the digit expansion. This gives weight at least
$q_*-1$. For the pair on opposite-parity letters $s,t$, torus invariance
gives $p^{s+u}\equiv p^{t+v}\pmod{p^r-1}$. The residues of
$1,p,\ldots,p^{r-1}$ are distinct and $p^r\equiv1$, so
$s+u=t+v\bmod r$. In particular $u-v$ is odd.
\end{proof}

\subsection*{Proof of \ref{B4.01b}}\label{prf:B4.01b}\phantomsection
\begin{proof}
Normalize $v_p(p)=1$ and put $q_*=1/(p-1)$. Let $g=\ell^{-1}$ in the
matrix coordinates. Then
\[
 \Theta(\xi Z)=(U\circ g)(Z)\exp\Phi(Z),\qquad
 \Phi(Z)=\epsilon\sum_{a,b}\beta(Z_{\bullet,ab}).
\]
Use the base weight $v_p(c)+\rho|n|$ of \ref{B4.01c} for $cZ^n$.

Index rows and matrix entries from zero. Let $Y=(Y_{t,k,a})$ be
independent variables and equip $\C_p(Y)$ with the Gauss norm of
\ref{B3.26}. Mark the letters of $T$ and call all rows except their first
rows ordinary. At target $t$ take labels $(s,j)$ with $j$ a signed odd
integer, $s=t-j\in T$, and
\[
 j>0\ \Longrightarrow\ p^j\le A_s,\qquad
 j<0\ \Longrightarrow\ p^{-j}\le|\nu_t|.
\]
Keep distinct integers $j$ even when they agree modulo $r$. Since
$|j|\le\log_p a$, there are at most $2\lfloor\log_p a\rfloor+1$ labels
at $t$. Enumerate them from zero in increasing order of $j$. Label $i$
uses row $e_i=i+\mathbf1_T(t)<N_0$ by the rank hypothesis.

Choose an integer $M\ge0$ at least all the $|j|$ in the lists. Let
$y_{t,k}$ have coordinates $Y_{t,k,a}$ except on marked first rows, where
$y_{t,0,a}=Y_{t,0,a}^{p^M}$. At target $t$ write $(s_i,j_i)$ for label $i$.
Put $L_i=p^{M+j_i}$ and set
\[
 \lambda_i(x)=\sum_a Y_{s_i,0,a}^{L_i}x_a,\qquad
 G_{ih}=\lambda_i(y_{t,e_h}).
\]
Reduce the coefficients of $\det G$ modulo the maximal ideal of
$\mathcal O_{\C_p}$, then specialize the target variables $Y_{t,k,a}$
to $\delta_{ka}$. Since $j_i$ is odd and $r$ even,
$s_i\ne t$, so the specialization gives
$\det(Y_{s_i,0,e_h}^{L_i})_{i,h}$. A permutation monomial records the
pair $(s_i,L_i)$ chosen at each distinct coordinate $e_h$. These pairs
are distinct, so the diagonal monomial has coefficient one. Thus
$\det G$ has Gauss norm one, also for the empty matrix, and the
adjugate gives an integral inverse.
\leanref{Formalism.\allowbreak{}DworkPowerSeriesWeights.\allowbreak{}labelsExist\_\allowbreak{}statement}{Formalism.\allowbreak{}GaussJetBounds.\allowbreak{}GaussJetProofs.\allowbreak{}Labels}{57--128}
\leanref{Formalism.\allowbreak{}DworkPowerSeriesWeights.\allowbreak{}generic\_\allowbreak{}minor\_\allowbreak{}norm}{Formalism.\allowbreak{}GaussJetBounds.\allowbreak{}GaussJetProofs.\allowbreak{}Generic}{78--124}

Set $v_i=\sum_h(G^{-1})_{hi}y_{t,e_h}$, so $\lambda_j(v_i)=\delta_{ji}$.
Keep each marked first row. Set $x_{t,e_i}=v_i$ and use
$x_{t,k}=y_{t,k}-\sum_i\lambda_i(y_{t,k})v_i$ on unreserved ordinary
rows. Thus $\lambda_i(x_{t,k})=\delta_{e_i k}$ on ordinary rows.
Recover the reserved rows by $y_{t,e_h}=\sum_iG_{ih}v_i$ and the others
by adding back the subtracted sums. Write these integral relations as
$y_{t,h}=\sum_k C^t_{hk}x_{t,k}$. At a marked target $C^t_{00}=1$
and $C^t_{0k}=C^t_{k0}=0$ for $k>0$.
\leanref{Formalism.\allowbreak{}DworkPowerSeriesWeights.\allowbreak{}frameA\_\allowbreak{}letter}{Formalism.\allowbreak{}GaussJetBounds.\allowbreak{}GaussJetProofs.\allowbreak{}FrameA}{48--195}

Put $K_x(V)=\Theta(\xi Z(V))$ with
$Z_{t,ab}(V)=\sum_k x_{t,k,a}V_{t,k,b}$. By the chain rule, direction
$y_{t,h}$ in column $b$ is represented by
$\sum_k C^t_{hk}\partial_{V_{t,k,b}}$. With commuting differential
symbols $D_{t,k,b}$ and a column choice $I$ indexed from zero, a minor
of height $h$ has symbol
\[
 \det\left(\sum_k C^t_{ik}D_{t,k,I(j)}\right)_{i,j<h}
 =\sum_{k_0,\ldots,k_{h-1}}
 \det(C^t_{i,k_j})_{i,j<h}\prod_{j<h}D_{t,k_j,I(j)} .
\]
Repeated $k_j$ give equal columns in the coefficient determinant.
At a marked letter each surviving minor uses row zero once by the form
of $C^t$. A symbol exponent $\kappa$ has degree $|\nu_t|$ at $t$.
Each row is used at most once per diagram column. At a marked letter
row zero has degree $A_t$, and ordinary rows occur in the $B_t$ columns
of height at least two. Thus the ordinary degree at $t$ is $D_t$ and
each labelled row has degree at most $c_t$ over all matrix columns.
\leanref{Formalism.\allowbreak{}DworkPowerSeriesWeights.\allowbreak{}op\_\allowbreak{}prod\_\allowbreak{}shape}{Formalism.\allowbreak{}GaussJetBounds.\allowbreak{}GaussJetProofs.\allowbreak{}ThmAOp}{244--298}

Write $a_\kappa$ for the coefficient of $D^\kappa$ in the transformed
minor product. By \ref{B3.11}, the evaluation of the scaled kernel at
the frame $\tau=y$ is
\[
 \sum_\kappa a_\kappa\,\kappa![V^\kappa]K_x(V),
 \qquad \kappa!=\prod_{t,k,b}\kappa_{t,k,b}!.
\]
In a product of series, exponents contributing to $[V^\kappa]$ are at
most $\kappa$ coordinatewise.
Integral linear substitution preserves base weight bounds by homogeneity.
By \ref{B4.01c}, a phase term of base weight below $q_*$ has weight at
least $q_*-1$ and takes the form $cZ_s^{p^u}Z_t^{p^v}$ in one matrix
entry, with $t=s+u-v$ and $u-v$ odd. Let $\mathcal O$ be the Gauss
valuation ring. The $p$-power map is a ring endomorphism of
$(\mathcal O/p\mathcal O)[V]$. Its $u$-fold iterate gives
\[
 \left(\sum_k d_k V_k\right)^{p^u}
 \equiv\sum_k d_k^{p^u}V_k^{p^u}\pmod{p\mathcal O[V]}
\]
for $d_k\in\mathcal O$. A term that splits
either power among several rows gains at least one in valuation and
has base weight at least $q_*$. For an unsplit term on rows $k,k'$ in
a common matrix column, the sum over the original row coordinate $a$
gives the multiplier $S=\sum_a x_{s,k,a}^{p^u}x_{t,k',a}^{p^v}$.

Orient a mixed pair with marked first row $(s,0)$ and ordinary row
$(t,k')$. Put $j=u-v$. The support bounds give $p^u\le A_s$ and
$p^v\le|\nu_t|$, hence $p^j\le A_s$ for $j>0$ and $p^{-j}\le|\nu_t|$
for $j<0$. Thus $(s,j)$ is a label $i$ at $t$. Since
$x_{s,0,a}^{p^u}=(Y_{s,0,a}^{L_i})^{p^v}$ by $M+u=M+j+v$, apply the
same congruence to the integral elements $Y_{s,0,a}^{L_i}x_{t,k',a}$
and sum over $a$ to obtain
\[
 S\equiv\lambda_i(x_{t,k'})^{p^v}\pmod{p\mathcal O}.
\]
For $k'\ne e_i$ we have $\lambda_i(x_{t,k'})=0$, so $S\in p\mathcal O$.
The extra unit of valuation gives base weight at least $q_*$.
\leanref{Formalism.\allowbreak{}DworkPowerSeriesWeights.\allowbreak{}pure\_\allowbreak{}cross\_\allowbreak{}term\_\allowbreak{}weight}{Formalism.\allowbreak{}GaussJetBounds.\allowbreak{}GaussJetProofs.\allowbreak{}PureCrossTermWeight}{163--260}

Take a dual minimizer $\sigma_t$ at capacities $D_t$ by \ref{B4.01a}.
For every positive odd $e$, its residue $e_0$ modulo $r$ is odd with
$1\le e_0<r$, so
$p^e\sigma_t+\sigma_{t+e}\ge p^{e_0}\sigma_t+\sigma_{t+e_0}\ge1$.
For marked $s$, let $b_s$ be the first positive odd $e<h_s$ for which
$A_sp^{-e}\le c_{s+e}$, or $h_s$ if there is none. Set
$\delta_s=p^{-b_s}$ with $p^{-\infty}=0$. Enlarge the weight of every
variable on row $(t,k)$ from $\rho$ to $\rho+d_{t,k}$, where
\[
 d_{t,k}=\begin{cases}
 \delta_t,&k=0,\ t\in T,\\
 \sigma_t,&\text{ordinary row with no label},\\
 \max\{\sigma_t,1-p^j\delta_s\},&\text{row labelled }(s,j),\ j>0,\\
 \max\{\sigma_t,p^j\},&\text{row labelled }(s,j),\ j<0 .
 \end{cases}
\]
The increments are nonnegative. For an unsplit pair on two ordinary
rows orient the powers so that $e=u-v>0$. The increase is at least
$p^v(p^e\sigma_s+\sigma_{s+e})\ge1$. For two marked first rows use the
same orientation. The target is marked, so $b_s\le h_s\le e$ and
$p^u\delta_s\ge p^{u-e}\ge1$. For a mixed pair with marked source, the
congruence for $S$ leaves the labelled partner to consider. For $j>0$
the increase is at least $p^v(p^j\delta_s+1-p^j\delta_s)=p^v$.
For $j<0$ the partner gives at least $p^v p^j=p^u$. Thus every phase
term of index $\le\kappa$ has weight at least $q_*$.
\leanref{Formalism.\allowbreak{}DworkPowerSeriesWeights.\allowbreak{}dualRadiiInequalities}{Formalism.\allowbreak{}GaussJetBounds.\allowbreak{}GaussJetProofs.\allowbreak{}OneLetterRadiusIdentity}{124--207}

The phase has zero constant term. Terms of index $\le\kappa$ in its
$h$-th power have weight at least $hq_*$. Legendre's formula
$v_p(h!)=\sum_{d\ge1}\lfloor h/p^d\rfloor\le hq_*$ then gives
nonnegative exponential weight through $\kappa$
\cite[Ch.~I, \S2, Ex.~14]{koblitz-1984}. Write $U=\sum_n u_nZ^n$.
By \ref{B4.01c}, each term contributed by $u_ng^n$ has base weight at
least $v_p(u_n)+\rho|n|\ge0$. The integral linear substitution and
the nonnegative row increments preserve this bound. Each coordinate of
$\kappa$ is at most $a$, so \ref{B3.19} gives digit loss at most
$\log_p a+1$ per factorial. The factorial-weighted coefficient
has valuation at least
\[
 \sum_{t,k,b}\kappa_{t,k,b}(\gamma-d_{t,k})
 -rN_0N_1(\log_p a+1),\qquad \gamma=q_*-\rho .
\]

Before the extra charges on labelled rows, the ordinary rows contribute
$\sum_t D_t(\gamma-\sigma_t)=F(D)$ by \ref{B4.01a}. Fix a marked $s$
and put $a_e=A_sp^{-e}$. The difference $\gamma A_s-H_s$ is
\[
 \sum_{e<h_s}\min(a_e,c_{s+e})+
 \begin{cases} a_{h_s},&h_s<\infty,\\0,&h_s=\infty.\end{cases}
\]
A positive label $e$ costs at most $c_{s+e}(1-p^e\delta_s)_+$
above the dual charge. For $e<b_s$ this is bounded by
$c_{s+e}=\min(a_e,c_{s+e})$. For $e\ge b_s$ it vanishes. The first-row
charge $A_s\delta_s$ is zero at $b_s=\infty$. Otherwise it equals
$a_{b_s}$, the minimum at $b_s<h_s$ or the endpoint at $b_s=h_s$.
The remaining amount is at least $H_s$. A negative label $-e$ costs
at most $p^{-e}c_{s-e}$ and satisfies $p^e\le|\nu_{s-e}|$, so
$c_{s-e}\le M_s$. These labels cost at most $\gamma M_s$ in total.
\leanref{Formalism.\allowbreak{}DworkPowerSeriesWeights.\allowbreak{}rowRadiusAssembly}{Formalism.\allowbreak{}GaussJetBounds.\allowbreak{}GaussJetProofs.\allowbreak{}CapacityLeReachCapacity}{278--507}

The integral symbol coefficients preserve the valuation bound in the
differential evaluation. The substitution of $p^M$-th powers for marked
first-row variables preserves the Gauss norm because it sends distinct
frame monomials to distinct monomials. The refined bound follows from
$J_s\le H_s$ term by term and \ref{B4.01a}.

For one letter use the variables $Y_{k,a}$ for the frame and substitute
$Z_{ab}=\sum_kY_{k,a}V_{k,b}$, with all $V$-variables of weight $\gamma$.
Apply \ref{B4.01c} at period two before identifying the letters.
An exceptional pair has degree $p^u+p^v\ge p+1$, so increasing the
variable weight from $\rho$ to $\gamma$ adds at least
$(p+1)(\gamma-\rho)=1$. Substitution by an integral linear form and its
unit multiple preserves degree and the resulting weight bound $q_*$.
The factor $\epsilon/2$ has valuation zero. The same substitution bound
gives $U\circ g$ nonnegative weight at $\gamma$. The exponential argument
gives valuation at least $-\gamma|\kappa|$ before factorials. Since
$|\kappa|=|\nu|$ and all $N_0N_1$ coordinates are at most $a$,
\ref{B3.19} gives $(q_*-\gamma)|\nu|-N_0N_1(\log_p a+1)$ after factorials.
Here $q_*-\gamma=\rho$, and the Gauss norm bounds each frame coefficient.
\leanref{Formalism.\allowbreak{}DworkPowerSeriesWeights.\allowbreak{}oneLetterR}{Formalism.\allowbreak{}GaussJetBounds.\allowbreak{}OneLetterRadius}{23--26}

The minor product is homogeneous of degree equal to the box count,
so scaling the kernel variables by $\xi$ multiplies the evaluation by
a unit by \ref{B3.11}.
\end{proof}

\subsection*{Proof of \ref{B4.02a}}\label{prf:B4.02a}\phantomsection
\begin{proof}
The long row is obtained from \(z_i+e_l\) by subtracting
\((k+m-2+t-2i)/2+(p-(m-j))/2\). This is \(p-1-i+(t+1)/2\). The short row uses the remaining \(m-j\) entries with the sign displayed above. These rational expressions are integers since \(t\) is odd. In the construction of the row space the prescribed integer coordinates are the integers with these rational values, so the formulas apply to the row space used for the period.

The entries of \(z\) decrease by at least one. This follows within the \(v\)'s by interlacing, at their last entry by \(v_{k-2}>U_{k-1}>-b_0\), and within the \(b\)'s by adding the two half-gap inequalities. Thus \(\kappa^{\rm big}\) is weakly decreasing. Its last entry is \(e_l-b_l-(t+1)/2\ge1\). When \(j<m\), the negative of the first short entry is \(b_j-e_l-(t+1)/2\ge1\). The increment in the negative short tuple is \(b_{j+c+1}-b_{j+c}-1\ge0\).

For the degree in a column, subtract \((N_l+N_{l+1}-1-2i)/2\) from the source parameter plus \(e_l\). The source parameter is \(U_i\) for \(i<k\), and is \(-e_q+(t-2d)/2\) for \(i=k+q(t+1)+d\). Substitution gives the two formulas. The head decreases weakly because \(U_i-U_{i+1}\ge1\). Adding the upper half-gap at \(q\) to the lower half-gap at \(q+1\) gives
\[
 e_{q+1}-e_q\ge t+3 .
\]
Hence successive tail blocks decrease by at least two, and the last block, if present, has value at least two. The last head entry minus the first tail entry is
\(U_{k-1}+e_0-1-t/2\ge1\), since \(U_{k-1}\ge-b_0+1/2\). At \(l=0\) this same expression is the last head entry. This proves positivity, including the first layer where the tail is empty.
\end{proof}

\subsection*{Proof of \ref{B4.02b}}\label{prf:B4.02b}\phantomsection
\begin{proof}
Use the integer base at each position in the weight construction of \ref{B8.05}. The adjustment of the base to the unit character adds a residue in \([0,|\mu(E)|)\) at the designated thick position. The same adjusted base is used for all layers and words. The construction adds \(qh_0\) for the \(q\)-th auxiliary coordinate and a further \(h_0/2\) for the extension coordinate. At a mark the integer endpoints before addition of \(1/2\) are \(a_s+k/2-1\) and \(-b_s-k/2\). Elsewhere they are \(4\delta(k-2)+4k\delta\) and \(-1\). This gives the asserted tuples.

For an interior coordinate the gaps to the two adjacent \(v\)'s are \(1/2+\xi_i\delta\) and \(4\delta-1/2-\xi_i\delta\), both positive. At a mark \(a_s\ge a\), \(b_s\ge b'\), and eventually \(S_c\le a\) and \(\lfloor b'/x\rfloor\ge k\). Indeed \(S_c,k,n\) grow at most polynomially by \ref{B8.02ac} and \ref{B8.59}, whereas \(b'/x=\exp(x^\varepsilon-x^{\varepsilon/3})\). These inequalities give \(U_0>v_0\) and \(U_{k-1}>-B_\rho\). Off the marks \(B_\rho\ge1\) gives the latter inequality. Finally \(t\le x^{401}\) at large selected points by the rank formula of \ref{B8.59}, and \(c_h\ge402\) in \ref{B8.05}, so \(h_0\ge t+3\). Substitution into \ref{B4.02a} gives the tail degrees.
\end{proof}

\section*{Proofs for Chapter~\ref{ch:local-to-column}: From local lifts to bounds on whole columns}

\subsection*{Proof of \ref{B4.07c}}\label{prf:B4.07c}\phantomsection
\begin{proof}
The certificate proves the prime conditions, and the assigned comparison and model are independent of its choice.
\leanref{Formalism.\allowbreak{}BasicNotions.\allowbreak{}GIdx}{Formalism.\allowbreak{}LocalRingCardinality.\allowbreak{}RhoDetPrimeEqMuTildeMulChiPow}{68--70}

For \(z\in V\), the tied comparison at a component \(j\) and full lattice tuple \(\mathbf L\) uses a polynomial \(Q_{z,j,\mathbf L}\) over \(K\). To construct it, use basis matrices of \(\mathbf L\) at \(p\) and the representative of \(j\) away from \(p\). Write the resulting adelic point as \(\gamma y_d\kappa\), with \(y_d\) a class representative, \(\gamma\) rational unitary and \(\kappa\) in the level group. Take the polynomial \(Q_d\) over \(K\) for \(z\) at \(d\). Form the block matrix of \(\gamma\) through the specified embeddings \(E\hookrightarrow K\) at the positions and their Frobenius translates, where \(E\) is the host field. Act by this matrix on \(Q_d\) and multiply by the element \(r_\kappa\in K\) whose complex image is the level character \(\chi(\kappa)\). Evaluation in the transported frame then gives \(\chi(\kappa)\) times the value at \(d\). Two presentations of the same tuple give adelic points related by a level element of character one, so the transformation law gives the same reading. The comparison applies \(i\) to the coefficients of \(Q_{z,j,\mathbf L}\) and gives zero outside full lattice tuples.
\leanref{Formalism.\allowbreak{}L2.\allowbreak{}G.\allowbreak{}Cmp.\allowbreak{}IsVal}{Formalism.\allowbreak{}LocalToGlobal.\allowbreak{}GroupLattices.\allowbreak{}ComparisonMaps.\allowbreak{}CmpLApply}{18--26}
\leanref{Formalism.\allowbreak{}L2.\allowbreak{}G.\allowbreak{}Cmp.\allowbreak{}cmpC}{Formalism.\allowbreak{}LocalToGlobal.\allowbreak{}GroupLattices.\allowbreak{}ComparisonMaps.\allowbreak{}CmpLApply}{30--32}
\leanref{Formalism.\allowbreak{}L2.\allowbreak{}G.\allowbreak{}Cmp.\allowbreak{}CmpTied.\allowbreak{}valRat}{Formalism.\allowbreak{}LocalToGlobal.\allowbreak{}GroupLattices.\allowbreak{}ComparisonMaps.\allowbreak{}NpDescFOfValRatF}{18--22}

An edge \(S\to L_v\) has \(pL_v\subseteq S\subseteq L_v\) with quotient of residue dimension \(r\). Its possibilities are finite by reduction modulo \(pL_v\). The neighbour is
\[
 (y_vH)(j,\mathbf L)=p^{\mathcal A_v}\sum_{S\to L_v}H(j,\mathbf L[v\leftarrow S]),
 \tag{4.6b}\label{eq:p4c2-neighbour}
\]
where \(\mathcal A_v\) sums the first \(r\) actual weights over the translates at \(v\). This formula preserves coefficients in \(i(K)\). To see integrality, choose a basis of \(L_v\) whose last \(k-r\) vectors lift a basis of \(S/pL_v\). With \(\omega^2=p\), the first \(r\) basis vectors of \(S\) are divided by \(\omega\) at each of the two half-steps from \(S\) to \(L_v\). For a weight \(\lambda\), put \(\nu_i=\lambda_i-\lambda_{k-1}\). A tableau product of this shape has degree \(d\le\sum_{i<r}\nu_i\) in the first \(r\) rows, since a minor of height \(h\) uses at most \(\min(h,r)\) of them. The Weyl generators lie in the same tableau lattice. In the enlarged frame a homogeneous generator changes by \(\omega^{-(d+r\lambda_{k-1})}\), including its determinant twist. The normalisation leaves the integral factor \(\omega^{\sum_{i<r}\nu_i-d}\). In the tensor product, with \(H_S,H_a,H_{L_v}\) the successive fibres, this gives
\[
 \omega^{\mathcal A_v}H_S\subseteq H_a,\qquad
 \omega^{\mathcal A_v}H_a\subseteq H_{L_v}.
 \tag{4.6c}\label{eq:p4c2-edge-scaling}
\]
Level transport permutes edges and carries fibre lattices by unit character factors. Thus \(y_v\) preserves both \(M\) and the tableau lattice \(M_{\rm na}\) of \ref{B3.52b}.
\leanref{Formalism.\allowbreak{}L2.\allowbreak{}G.\allowbreak{}Model.\allowbreak{}weyl\_\allowbreak{}scale}{Formalism.\allowbreak{}LocalToGlobal.\allowbreak{}GroupLattices.\allowbreak{}FockModels.\allowbreak{}OperatorsPreserveLattices}{76--127}
\leanref{Formalism.\allowbreak{}L2.\allowbreak{}G.\allowbreak{}Model.\allowbreak{}operators\_\allowbreak{}preserve\_\allowbreak{}lattices}{Formalism.\allowbreak{}LocalToGlobal.\allowbreak{}GroupLattices.\allowbreak{}FockModels.\allowbreak{}OperatorsPreserveLattices}{334--344}

If \(H=\sum_i t_i c(z_i)\) has coefficients in \(i(K)\), choose a \(K\)-linear retraction \(\pi:\mathcal K\to K\) of \(i\). For any coefficient evaluation \(e\), with \(e(c(z_i))=i(a_i)\), one has
\[
 e\!\left(c\!\left(\sum_i\pi(t_i)z_i\right)\right)=i\!\left(\sum_i\pi(t_i)a_i\right)=i\pi(e(H))=e(H).
\]
Evaluations separate functions. Apply this to \(H=y_vc(z)\), which lies in \(\mathcal K M\) and has rational coefficients by \eqref{eq:p4c2-neighbour}. Injectivity gives a unique \(K\)-linear \(y_v^K\) intertwining with \(y_v\).
\leanref{Formalism.\allowbreak{}L2.\allowbreak{}G.\allowbreak{}Cmp.\allowbreak{}yvRat\_\allowbreak{}of\_\allowbreak{}valRat}{Formalism.\allowbreak{}LocalToGlobal.\allowbreak{}GroupLattices.\allowbreak{}ComparisonMaps.\allowbreak{}YvRatOfValRat}{269--280}

Replace \(y_v\) by \(y_v^K\) in \(Y_B\) to form \(Y_B^K\). The module
\[
 C_w=\left(\sum_{B\subseteq W_p}Y_B^K(p^{n(B)}L_w)\right)\cap p^{b_0}L_w\cap V^r
\]
extends to \(N_P\). Indeed, \(A\to\mathcal O\) is local by \(i^{-1}(\mathcal O)=A\) and flat by torsion-freeness over a DVR. The first two modules in the intersection extend by intertwining and the input comparison. Flatness commutes with the kernels of the difference map for the first intersection and the map to \(V/V^r\) for the second. Inverting a uniformiser of \(A\) in \(\mathcal O\) gives \(\mathcal K\), and the basis independence identifies the resulting scalar extension with its image in \(\mathscr P\). The same kernel argument gives \(\mathcal O c(R_w)=M\cap V^{\rm ram}\).

For any finite free \(A\)-submodule of \(V\), extend an \(A\)-basis of the submodule to a \(K\)-basis of \(V\). Membership of a rational vector's image in the extended submodule forces its complementary coordinates to vanish and the others to lie in \(i^{-1}(\mathcal O)=A\). Apply this to \(C_w\), which is finite free as a submodule of \(p^{b_0}L_w\), to obtain \(J_w=C_w\). The upper inclusion follows from the second and third terms of the intersection. The lower inclusion follows from the empty subset's contribution \(p^aL_w\).
\leanref{GlobalTheory.\allowbreak{}EX.\allowbreak{}OB.\allowbreak{}mem\_\allowbreak{}span\_\allowbreak{}of\_\allowbreak{}semilinear}{Development.\allowbreak{}AbelianVarietyTheory.\allowbreak{}Proofs.\allowbreak{}GlobalInterpolation.\allowbreak{}MemSpanOfSemilinear}{16--78}
\leanref{Formalism.\allowbreak{}BasicNotions.\allowbreak{}FBC.\allowbreak{}span\_\allowbreak{}inf\_\allowbreak{}span\_\allowbreak{}tower}{Formalism.\allowbreak{}LocalRingCardinality.\allowbreak{}Lemmas}{368--467}
\end{proof}

\subsection*{Proof of \ref{B4.07b}}\label{prf:B4.07b}\phantomsection
\begin{proof}
The half-step inclusions \eqref{eq:p4c2-edge-scaling} for the fibres of \ref{B3.52a} give \(H_L\subseteq R_a\) on outgoing edges and \(H_L^\vee\subseteq R'_a\) on incoming edges. The intersections give the first two pairs of inclusions. For \(z=\omega^{-A_v}h\in R_a\) and a functional \(\phi=\omega^{-A_v}\psi\) in the defining intersection for \(P_L^\vee\), one has \(\phi(p^{A_v}z)=\psi(h)\in\mathcal O\). This proves the last inclusion and gives integrality of each summand of \(y_v\) on \(L^+_{B,v}\).

For the orbit function, every contributing edge is a transport of \(S\to L'\). The terms with \(L'\ne L_v\) lie in the target fibre multiplied by \(p^E\). At a transport of \(\mathbf L\), triviality of its stabiliser makes the transport of the specified edge unique. The remaining term is the orbit function of \(p^{A_v}\xi\). The difference divided by \(p^E\) is equivariant and fibrewise integral.
\end{proof}

\subsection*{Proof of \ref{B4.07a}}\label{prf:B4.07a}\phantomsection
\begin{proof}
If \(W=\varnothing\), Definition~\ref{B4.07} gives \(N_P=R\), and both budgets are zero by \ref{B3.52}. Suppose \(W\ne\varnothing\). Set \(D=D_p,\ q=p^D,\ S=\max_{v\in W}S_v\) and \(s=m\sum_{v\in W}\kappa_v\). For \(A=M,M_{\rm na}\) write
\[
 I(A)=\sum_{B\subseteq W}p^{\sum_{v\notin B}S_v}
                   \Bigl(\prod_{v\in B}y_v^m\Bigr)A .
\]
Here and below complements of subsets of positions are taken in \(W\). The summand for \(B=\varnothing\) shows that \(I(A)\) has full rank in \(A\).

Use the fibres of Lemma~\ref{B4.07b}. At a target of local fibre rank \(d\), choose a Smith basis \(b_j\) of \(P_L\) for the span of the tested vectors \(p^{\mathcal A_v}F_i\). Write its generators as \(a_jb_j\), allowing \(a_j=0\). The forms and combinations in \ref{B3.52a} give integral matrices with
\(C\operatorname{diag}(a_j)X=qI_{h_0}\), where \(h_0\ge(1-\kappa_v)d\). Discard the terms with \(a_j\nmid q\). Each discarded entry belongs to \(q\mathfrak m_{\mathcal O}\). The remaining product divided by \(q\) is therefore an integral matrix congruent to \(I_{h_0}\) in the residue field. Its rank over the fraction field is \(h_0\). Thus at least \(h_0\) indices satisfy \(a_j\mid q\), and for each of them \(qb_j\) is an integral combination of the tested vectors.

For \(\varnothing\ne B\subseteq W\), use these bases on \(B\) and the base fibre bases elsewhere. If a tensor \(b\) uses a passing index at \(v\in B\), insert the integral combination for that index at its source edges. At every other edge, \(F_i\in\omega^{c_*}R_{a'}\) and \(p^{\mathcal A_v}R_{a'}\subseteq P_{L'}\). Since \(\omega^2=p\) and \(c_*\ge4(D+1)\), the other contributions lie in \(p^{D+1}L_B\). Trivial stabilizers allow Lemma~\ref{B4.07b} to be applied to each orbit function. With its notation this gives
\[
 q\delta_{i_0,\mathbf L}(b)\equiv
 y_v\sum_i u_i\delta_{i_0,\mathbf L[S_i]}(\xi_i)
 \pmod{p^{D+1}L_B}.
 \tag{4.6g}\label{eq:p4c2-good-products}
\]
Here \(u_i\) are the combination coefficients and \(\xi_i\) has \(F_i\) in the replaced factor. Let \(C_B\) consist of the orbit functions of products failing at every \(v\in B\). At each representative the number of such products is at most its fibre rank times \(\prod_{v\in B}\kappa_v\). Summing over representatives gives
\[
 |C_B|\le d_M\prod_{v\in B}\kappa_v .
 \tag{4.6h}\label{eq:p4c2-bad-products}
\]
Set \(T_B=\sum_{v\in B}y_vL^+_{B,v}+\mathcal O C_B\), where \(\mathcal O C_B\) denotes the span. All basis products give \(qL_B\subseteq T_B+pqL_B\). The finitely generated module \((qL_B+T_B)/T_B\) equals its \(p\)-multiple. Nakayama's lemma gives \(qL_B\subseteq T_B\), in particular
\[
 p^{D+1}L_B\subseteq T_B .
 \tag{4.6i}\label{eq:p4c2-stage-input}
\]

Order \(W\), fix \(0\le\beta<|W|\), and put \(h=2(D+1)\),
\(d_v'=2q_v+h\), \(K_v=m d_v'\), \(K_\beta=\sum_v K_v+2\beta S\).
A stage consists of \(B\subseteq W\) of size \(\beta+1\) and integers \(0\le j_v<m\) for \(v\in B\). Put \(P=\{u<\max B:u\notin B\}\),
\[
 Y=\prod_{u\in P}y_u^m\prod_{v\in B}y_v^{j_v},
 \qquad a=\sum_{u\in P}K_u+\sum_{v\in B}(q_v+j_vd_v').
\]
Let \(C^{(\beta)}\) be the union of \(YC_B\) over these stages. Since \(P_L\subseteq H_L\) and all the \(y_v\) preserve \(M_{\rm na}\), this is a subset of \(M_{\rm na}\). We prove
\[
 \omega^{K_\beta-a}YL_B
 \subseteq I(M_{\rm na})+\mathcal O C^{(\beta)} .
 \tag{4.6j}\label{eq:p4c2-stage}
\]
Use descending induction on \(m|P|+\sum_{v\in B}j_v\), which is bounded by \(m|W|\). At an active position the unused part of its allowance is
\[
 K_v-(q_v+j_vd_v')=(m-1-j_v)d_v'+q_v+h.
\]
Thus \(K_\beta-a\ge h\). For \(z\in L_B\), write
\(\omega^h z=\sum_{v\in B}y_vz_v+c\) using \eqref{eq:p4c2-stage-input}, with \(z_v\in L^+_{B,v}\) and \(c\in\mathcal O C_B\). Multiply by \(\omega^{K_\beta-a-h}Y\). The term from \(c\) belongs to the claimed span. Fix one of the remaining terms.

If \(j_v+1<m\), increase \(j_v\) by one. By \eqref{eq:p4c2-fibres}, multiplication by \(\omega^{2q_v}\) takes \(U\) into \(P_L\), via \(H_L\). Apply the induction hypothesis to \(z'=\omega^{2q_v}z_v\). The new operator is \(Yy_v\) and the new cost is \(a'=a+2q_v+h\). Thus \(K_\beta-a'+2q_v=K_\beta-a-h\), the exponent required for this term.

If \(j_v=m-1\) and there is a position after \(\max B\), take the first such position \(w\). Replace \(v\) by \(w\) in \(B\), with \(j_w=0\). The completed set becomes \(P\cup\{v\}\). At \(v\), multiplication by \(\omega^{q_v}\) takes \(U\) into \(H_L\). At \(w\), multiplication by \(\omega^{q_w}\) takes \(H_L\) into \(P_L\). Hence \(z'=\omega^{q_v+q_w}z_v\) lies in the new stage lattice. Completing \(v\) uses its remaining \(q_v+h\), and entering \(w\) uses \(q_w\). The new cost is \(a'=a+h+q_v+q_w\), so \(K_\beta-a'+q_v+q_w=K_\beta-a-h\). Commutation gives the new operator \(Yy_v\). In both transitions the induction index increases by exactly one.

Finally let \(j_v=m-1\) and \(\max B=\max W\). Put \(B'=B\setminus\{v\}\) and \(n'=\sum_{u\in B'}S_u\). The vector
\(\prod_{u\in B'}y_u^{j_u}(\omega^{q_v}z_v)\) belongs to \(M_{\rm na}\). The summand of \(I(M_{\rm na})\) indexed by \(W\setminus B'\) therefore contains
\(\omega^{2n'+q_v}Yy_vz_v\). Since \(n'\le\beta S\), the \(2\beta S\) in \(K_\beta\) and the remaining \(q_v+h\) at \(v\) give \(K_\beta-a-h\ge 2n'+q_v\). This proves the terminal case and finishes the induction.

Start with the first \(\beta+1\) positions, all at exponent zero. The product of the inclusions \(\omega^{q_v}H_L\subseteq P_L\) pays its initial cost \(a\). For fixed \(B\) there are \(m^{\beta+1}\) exponent choices. Expanding \((\sum_v\kappa_v)^{\beta+1}\) counts each product on distinct indices \((\beta+1)!\) times. Hence
\[
 |C^{(\beta)}|\le d_Mm^{\beta+1}
       \sum_{|B|=\beta+1}\prod_{v\in B}\kappa_v
 \le d_M\frac{s^{\beta+1}}{(\beta+1)!}.
 \tag{4.6k}\label{eq:p4c2-stage-count}
\]
By \ref{B3.52}, \(K_\beta/2+b_p=\tau_p+\beta S\). Multiply \eqref{eq:p4c2-stage} at the initial stage by \(p^{b_p}\). The sandwich \eqref{eq:p4c2-sandwich} and linearity give an inclusion of \(p^{\tau_p+\beta S}M\) in \(I(M)\) plus a span in \(M\) with at most \(|C^{(\beta)}|\) generators.

Write \(a_j\) for the uniformizer exponents of \(I(M)\subseteq M\). The image of \(p^{\tau_p+\beta S}M\) in \(M/I(M)\) needs one generator for each \(a_j>e(\tau_p+\beta S)\). A submodule of a module with \(b\) generators over a DVR needs at most \(b\) generators, by lifting to a submodule of \(\mathcal O^b\). Thus \eqref{eq:p4c2-stage-count} bounds this count. For \(\beta\ge|W|\) it is zero by the empty summand of \(I(M)\). Intersecting \(I(M)\) with \(p^{\tau_p}M\) changes the exponents to \(\max(a_j,e\tau_p)\). For \(S>0\), sum the excess over successive intervals of width \(eS\). It is at most \(ed_M S\sum_{\beta\ge0}s^{\beta+1}\le2ed_M Ss\), since \(s\le1/2\). For \(S=0\) the empty summand gives \(I(M)=M\). Saturation gives \(R\cap p^{\tau_p}M=p^{\tau_p}R\), and the inclusion of \(R\) induces an injection \(R/N_P\hookrightarrow M/(I(M)\cap p^{\tau_p}M)\). Now \(\mathrm{io}_p=2Ss\) by \ref{B3.52}, which gives both assertions.
\end{proof}

\subsection*{Proof of \ref{B4.09}}\label{prf:B4.09}\phantomsection
\begin{proof}
The matching system gives the coefficient and terminal estimates of \ref{B3.45} and \ref{B3.42}. At every position and every step \(0\le i<m\), its differentiated coefficients satisfy the bound of \ref{B3.12} with exponent zero by the final sentence of item~7 of Definition~\ref{B3.05a}. The filling and position map carry the tie and \eqref{eq:p4c2-small-prime} to the twist. Apply \ref{B3.55} with \(\mathcal O\hookrightarrow\mathcal K\), lower exponent \(B^0_{u,w}(\sigma)\) and \eqref{eq:p4c2-complement}--\eqref{eq:p4c2-large-prime}, and split the complementary sum at \(V\).
\end{proof}

\subsection*{Proof of \ref{B4.10}}\label{prf:B4.10}\phantomsection
\begin{proof}
For each \(t\in W_p\), the ceiling gives
\[
 \sum_j H_{t,j}\le\sum_j A_{t,j}+S_t.
 \tag{4.18}\label{eq:p4c2-price-ceiling}
\]
Apply \eqref{eq:p4c2-hecke} with \(V=W_p\), \(V'=B\) and multiply by \(p^{n(B)}\), which adds \(n(B)\) to the local exponent. The bound suffices since
\[
 \sum_j g^{\rm c}_{p,j}\le n(B)+\sum_j
 \left(\sum_{t\in B}H_{t,j}+\sum_{t\in R_p\setminus B}A_{t,j}\right).
 \tag{4.19}\label{eq:p4c2-generators}
\]
This follows from \eqref{eq:p4c2-price-ceiling} on omitted representatives, and also holds when the gain vanishes. The bound extends to the span by \(\mathcal O\)-linearity. The filling, position map and ramified submodule in \eqref{eq:p4c2-np} are common.
\end{proof}

\subsection*{Proof of \ref{B4.10b}}\label{prf:B4.10b}\phantomsection
\begin{proof}
First use frames with entries in the embedded coefficient field. Such frames exist simultaneously at the primes over \(p\). Lift a basis of each localization by approximation in the host field. The condition of being a basis is open, since the change matrix must have integral entries and unit determinant. The Hermitian dual basis gives the companion frame. At a level-free place the marking conditions at the level and at the row primes impose no further change on these frames.

In one frame let \(D_\nu\) be the integral span of the products of flag minors of shape \(\nu\), and take the tensor product over the types. Its integral dual under the normalized polynomial pairing is the Weyl lattice. Indeed, the semistandard products are an integral basis of \(D_\nu\), and the divided-power construction and the evaluation at the identity in \ref{B3.10a} give the dual lattice. Consequently a functional lies in that Weyl lattice exactly when its values on all the flag-minor products are integral. The determinant power is part of the frame change on both sides of this test. In the source its removal and the inclusion of the full-height columns give the source polynomial of shape \(\nu\). In the target the same shape is read in rank \(N_{i+1}\). Changing a lattice basis acts on these two integral lattices by inverse dual automorphisms. Thus the test gives membership in the prescribed lattice in any local basis.
\leanref{GlobalTheory.\allowbreak{}LatSCore.\allowbreak{}prSC\_\allowbreak{}int\_\allowbreak{}iff}{Development.\allowbreak{}AbelianVarietyTheory.\allowbreak{}Proofs.\allowbreak{}IwasawaDecomposition.\allowbreak{}FormsLatticeAtClassIffTwistedPolynomial}{20--33}
\leanref{GlobalTheory.\allowbreak{}Mem.\allowbreak{}IntDualT.\allowbreak{}mem\_\allowbreak{}locW\_\allowbreak{}iff\_\allowbreak{}prTC\_\allowbreak{}le\_\allowbreak{}one}{Development.\allowbreak{}AbelianVarietyTheory.\allowbreak{}Proofs.\allowbreak{}IwasawaDecomposition.\allowbreak{}MemLocalisedWeylLatticeIffTargetPairingIntegral}{21--38}

Here is the rationality detail for the target test. A rational representative, evaluated on a rational matrix in a rational frame, has value in \(\iota_w(K)\). Write the evaluation polynomial as a functional on the universal product of minors. Choose an \(\iota_w(K)\)-linear retraction \(\pi:\C_p\to\iota_w(K)\), and replace the value of the functional on each monomial by its image under \(\pi\). Extend the resulting assignment \(\C_p\)-linearly on monomials. The two functionals give the same evaluation polynomial on matrices over \(\iota_w(K)\), since the universal product has integral coefficients. This field is infinite, so the evaluation polynomials agree identically. Their difference vanishes on each flag-minor test. To see this, polarize the columns of the universal product, or apply Cauchy--Binet with indeterminate columns and read the coefficients in the semistandard basis. The diagram fits in the rank, so all the columns used in this operation exist. The reading of the original functional therefore equals its image under \(\pi\). In the adapted frames of the kernel one may further multiply the frame at each type by a local unit. Homogeneity multiplies each reading by products of powers of those units and of their determinants. Its absolute value is unchanged.
\leanref{GlobalTheory.\allowbreak{}KerVal.\allowbreak{}canM\_\allowbreak{}ker}{Development.\allowbreak{}AbelianVarietyTheory.\allowbreak{}Proofs.\allowbreak{}IwasawaDecomposition.\allowbreak{}TargetPairingSizesInCoefficientFieldValueGroup}{62--75}
\leanref{GlobalTheory.\allowbreak{}Pdv.\allowbreak{}KerScale.\allowbreak{}kerEntHyp\_\allowbreak{}scB}{Development.\allowbreak{}AbelianVarietyTheory.\allowbreak{}Proofs.\allowbreak{}SatakeParameters.\allowbreak{}TrueColumnChoicesKernelDatumAdaptedBases}{160--212}

It remains to pass from values to forms. In representatives the rational space consists of the tuples of rational weight polynomials satisfying the transformation equations for forms. Indeed, transport from representatives determines a form, and two transports to the same point differ by a stabilizer. The equations at a representative are precisely invariance under that stabilizer, with its prescribed character and weight action. The section polynomial is determined by its evaluation in the frame model. Hence rational representatives for a tuple already in the complex form space give a rational form. The integral form lattice imposes the local Weyl conditions on these values. To localize this description, take a rational tuple satisfying the conditions at \(w\). It has denominators at finitely many other places, since there are finitely many representatives and each value lattice is full. Choose an integer of \(\mathcal O_K\) in sufficiently high powers of these other prime ideals and congruent to \(1\) modulo \(\mathfrak p_w\). Its multiple satisfies all the integral conditions, and this integer is invertible in \(A\). Conversely, division of an integral tuple by an element outside \(\mathfrak p_w\) preserves the local conditions. Applying the frame test at every representative proves both equivalences.
\end{proof}

\subsection*{Proof of \ref{B4.10a}}\label{prf:B4.10a}\phantomsection
\begin{proof}
Write \(K_i\) for the lift and \(h_c\) for the stabilizer order at a source class \(c\). Let \(x\) be an integral source vector. By \ref{B4.10b}, its value at \(c\) is represented by a rational functional \(v_c\) with
\(\lvert v_c(\Delta_J)\rvert_p\leq1\) on every source product of minors \(\Delta_J\).

Expand each minor of \(X_zZ_z\) by Cauchy--Binet before applying the differential operator. The resulting \(P_I\) lies in the span of the \(\Delta_J\). Fixed-shape straightening gives the semistandard products as an integral basis of this span of products over \(\Z\) \cite[Theorem~3.2(1)]{deconcini-1980}. In a diagonal monomial order the leading monomial of a semistandard product has coefficient one. It records the multiset in every tableau row, which determines that row in its prescribed order. Thus the leading monomials are distinct. In decreasing order, subtract the leading coefficient times its basis product. The entry hypothesis bounds this coefficient by \(p^{-B}\), and subtraction preserves the bound on the remaining polynomial since the basis products have integral coefficients. The argument in the product monomial order works for all types together. We obtain
\[
 P_I=\sum_J a_{IJ}\Delta_J,\qquad |a_{IJ}|_p\leq p^{-B}.
\]

By rationality of the pair map, its value from \(c\) to \(c'\) has a rational target representative \(v_{cc'}\). Its \(I\)-reading is
\[
 t_{cc'I}=\epsilon_{cc'}v_c(P_I)
          =\epsilon_{cc'}\sum_J a_{IJ}v_c(\Delta_J).
\]
This is the defining differential evaluation of the kernel in the two frames. The pair scalar \(\epsilon_{cc'}\) has absolute value one. The ultrametric inequality gives \(|t_{cc'I}|_p\leq p^{-B}\). Choose a uniformizer \(\pi_w\in K\) and set \(a=\pi_w^{\lceil e_wB\rceil}\). By \ref{B4.10b}, a nonzero \(t_{cc'I}\) has norm \(p^{-v_w(b)/e_w}\) for some \(b\in K^\times\). Hence \(v_w(b)\geq\lceil e_wB\rceil\). This proves \(|t_{cc'I}/\iota_w(a)|_p\leq1\), including zero readings.

Let \(n\) be a positive common multiple of the \(h_c\), and put \(k=v_p(n)\). The scalar \(p^k/h_c=(n/h_c)/(n/p^k)\) lies in \(\mathcal O_{K,(w)}\). At \(c'\) form the rational value
\[
 z_{c'}=\sum_c\frac{p^k}{a h_c}v_{cc'}.
\]
Every reading is integral. These values evaluate to \(\iota(p^k/a)K_i(x)\), since the class weight is \(1/h_c\). The lift preserves the complex space of forms. The target implication of \ref{B4.10b} gives a vector \(z\) in the localized target lattice with exactly this evaluation. The lift of \(x\) is represented by \((a/p^k)z\), and
\[
 v_w(a/p^k)=\lceil e_wB\rceil-e_wk\geq e_w(B-k).
\]
Take \(n=D\Gamma_u\) by \ref{B4.03}. When \(i=0\) and \(\sigma=1\), all \(h_c=1\), so take \(n=1\). These choices give \eqref{eq:p4c2-entry-lift}. At \(\mu_0\) the stage isomorphisms identify the two lifts and carry the integral lattices onto each other. Extending those isomorphisms to \(\mathcal O_{K,(w)}\) transports the same factor \(a/p^k\).
\end{proof}

\subsection*{Proof of \ref{B4.11}}\label{prf:B4.11}\phantomsection
\begin{proof}
In (i) the place is level-free. On the signal and even-order branches use the bound without a neighbour operator on kind \(A\) and \ref{B4.11a}(2) on kind \(BC\). Kind \(A\) is the first sheet at even-order inert primes, and all signal orbits have kind \(BC\). Add by \ref{B4.11a}(1). The full base sum is \(F_j(w)\), and \(I^*_{j,p}\) covers \eqref{eq:p4c2-orbit-error} and the positive part of the kind-\(A\) average. The entry exponent is at least
\[
 \sum_{t\in R_p}A_{t,j}+F_j(w)-I^*_{j,p}.
 \tag{4.25}\label{eq:p4c2-even-entry}
\]
On the odd-order branch use \ref{B4.11a}(5), with the large-error constants of \ref{B4.03}, to obtain \(F_j(w)-I^*_{j,p}\). Both branches give, by \eqref{eq:p4c2-entry-lift} and the error identity, the step exponent
\[
 g_{0,S}(p,j)+F_j(w)-E^{\rm ev}_{u,j}(w)+\sigma\mathbf1_{\{i=0\}}\delta_{u,w},
 \qquad j=i+1.
 \tag{4.27}\label{eq:p4c2-orbit-step}
\]

In (ii) use \ref{B4.11a}(4). In (iii) the host class gives \(p=d\) and admissibility gives \(g_0,g_1\ne d\). The entry bound of \ref{B4.11a}(3) is at least \(F_j(w)-D^*_j\) by the maximum over words in \ref{B4.03}. Apply \eqref{eq:p4c2-entry-lift} in both cases and add the step bounds in (i)--(iii) by \ref{B4.12}. In (iv) use the entry exponent zero of admissibility and \eqref{eq:p4c2-entry-lift}, with \(Q'(\mu_0)=\varnothing\) at the reference index.
\end{proof}

\subsection*{Proof of \ref{B4.16a}}\label{prf:B4.16a}\phantomsection
\begin{proof}
For (1), setting the variables outside a block to zero eliminates the other phases because their constant terms vanish. For addition, expand the integral multiplier into monomials and factor each monomial by blocks. Inverse logarithmic substitution gives a product of kernels with integral monomial multipliers in logarithmic coordinates. Differential evaluation factors by \ref{B3.11}. At each order the sum is finite, so the exponents add.

For (2) put \(z_s=P^sr\), \(w_s=w_j(z_s)\), \(W_s=(k-1)w_s\). Let \(T\) be the marked indices and \(A_s,B_s\) the first two rows of \(\nu_s\). By \eqref{eq:p4c2-shape}, the refined cyclic estimate applies at size \(a_\ast^2\) with base \(F_{\mathrm{Orb},j}(p,r)\). Indeed \(p\le a_\ast\), the rank inequality holds at \(N_i\), and the capacities are \(W_s+e_s\) off \(T\), \(W_s-2w_s+e_s\) on \(T\), with \(0\le e_s\le4w_s\).

At a signal prime, \(r=(0,t_0)\) and \(z_s=(s\bmod2,t_0+sh_P)\). As \(n_P=2\operatorname{ord}(2h_P)\), the marks correspond to \(T(t_0)\) of \ref{B3.48}. An odd positive distance \(e\) from \((0,t)\) ends on sheet one, so no positive odd distance reaches another mark and none is excluded. For the capacity \(c_{s+e}\) of \ref{B4.01b}, \eqref{eq:p4c2-shape} gives \(c_{s+e}\le B_j(z_{s+e})+c_G=\bar C_{t,e}\) and \(A_s-B_s\ge\bar\alpha_t\). Hence
\[
 [\bar\alpha_t p^{-e}-\bar C_{t,e}]_+
 \le[(A_s-B_s)p^{-e}-c_{s+e}]_+\le A_sp^{-e}.
\]
Sum over odd \(e>0\) and the marks. The marked terms \(J_s\) of \ref{B4.01b} cover \(g_{0,r,j}\) because the digit subtraction is nonnegative and each \(J_s\ge0\). On the second inert sheet \(g_{0,r,j}=0\).

The reverse-partner maximum is at most \(b+\eta+1+3S_0\), and \(C_6\le c_{e,6}\). Thus the first two terms of \eqref{eq:p4c2-orbit-error} cover these losses. For \(l=\log_pa_\ast\ge1\), one has \(C(2l+1)\le(1+c_{\mathrm{dl},1})l\) by \eqref{eq:p4c2-constants}. The rank and linear digit terms cover the error \(Cn_PN_iN_{i+1}(2l+1)\) on the orbit, since their logarithms are at least \(l\) and \(N_iN_{i+1}\le N_m^2\).

For (3) the quadratic phase pairs the two sheets at \(t\) with scalar \(-Su_t\) of norm at most \(p^{1/2}\). Its factorial-weighted coefficient for matching powers \(h,h\) at a matrix entry is \(h!(-Su_t)^h\), with zero for unequal powers. Put \(\Delta_t=\sum_{a,b}\partial_{(0,t),ab}\partial_{(1,t),ab}\). The contraction of order \(n_t\) has operator \((-Su_t)^{n_t}\Delta_t^{n_t}/n_t!\). The operator after removing the scalar is integral, by the multinomial expansion with divided derivatives on one side and ordinary derivatives on the other. For diagrams \(\alpha,\lambda\) on the two sheets, with first rows \(a_1,b_1\), at most \(|\alpha|-a_1\) contractions use lower rows on the \(\alpha\) side. A first-row contraction differentiates in that frame-row direction on the other sheet. Each of the \(b_1\) minors is affine in this translation, since two columns replaced by the same vector give zero. At most \(\min(a_1,b_1)\) first-row contractions survive. In particular \(n_t\le|\alpha|\) off the top and \(n_t\le|\alpha|-a_1+\alpha_2+b_1\) on the top. These are the contraction bounds of \ref{B4.03}. The residual series has integral factorial-weighted coefficients, which proves (3).

For (4), \eqref{eq:p4c2-shape} gives \(|\nu_z|\ge(k-1)w_j(z)\) everywhere. On the top this uses both top inequalities. Thus the base term in \(B_C\) covers \(F_{\rm inert}\) at \(n_P=1\). Since \(C(2l+1)\le c_{\rm dig}l\), the digit term of \(C^*_{j,p}\) covers the subtraction in \(B_C\). Apply the one-letter bound at size \(a_\ast^2\) and add by (1).

For (5) use optimal dual coordinates \(0\le y_s\le1\) of \ref{B4.01a} on each kind-\(BC\) orbit. With \(W_s=(k-1)w_s\), these give
\[
 F(W)=\sum_s W_s(\gamma(p)-y_s),\qquad
 p^e y_s+y_{s+e}\ge1\quad(1\le e<n_P,\ e\text{ odd}).
\]
Setting \(y'_s=\max(y_s,p^{-3})\) enforces all positive odd distances, since every omitted distance is at least three. This costs at most \(p^{-3}\sum_s W_s\). These orbits have no top positions and \(0\le|\nu_s|-W_s\le4w_s\). For \(H=\{s:y_s\ge2/p\}\), \ref{B4.01a} gives \(\sum_{s\in H}W_s\le p^{-1}\sum_s W_s\). Since \(\gamma(p)\ge1/p\),
\[
 \sum_s(|\nu_s|-W_s)[y'_s-\gamma(p)]_+
 \le \frac4p\sum_{s\notin H}w_s+4\sum_{s\in H}w_s
 \le\frac8p\sum_s w_s.
 \tag{4.26a}\label{eq:p4c2-odd-weight}
\]
On kind \(A\) take \(y'_s=1/(p+1)\), which satisfies the odd-distance constraints and gives \(\gamma(p)-y'_s\ge0\).

Apply \ref{B4.01c} at period \(2n_P\) to the phases in \eqref{eq:p4c2-odd-product}. Give each letter at \(z\) weight \(\rho(p)+y'_z\) in addition to the coefficient valuation. A term of base weight at least \(q_\ast=1/(p-1)\) keeps that bound. Each remaining term is an exceptional pair of base weight at least \(q_\ast-1\). With powers \(p^\alpha,p^\beta\), \(\alpha>\beta\), and \(z\) carrying the larger power, its partner projects to \(P^ez\) for \(e=\alpha-\beta>0\) odd. The weight increase \(p^\beta(p^ey'_z+y'_{P^ez})\) is at least one, also if the positions coincide after projection. Identifying the letters preserves weights. The factors \(\epsilon_r/2\) are units, so the resulting phase \(\Phi\) has weight at least \(q_\ast\) and zero constant term.

Write \(g=\ell^{-1}\). Each coordinate of \(g\) has at least the base weight \(\rho(p)\) by \ref{B4.01c}, so \(U\circ g\) has nonnegative weight. Also \(v_p(n!)=\sum_{r\ge1}\lfloor n/p^r\rfloor\le nq_\ast\). Hence \(\Theta(\xi Z)=(U\circ g)(Z)\exp\Phi(Z)\) has nonnegative weight. For \(0\le n\le a_\ast^2\), the factorial formula using base-\(p\) digits gives
\[
 v_p(n!)\ge nq_\ast-(2l+1),\qquad l=\log_pa_\ast,
 \tag{4.26b}\label{eq:p4c2-factorial}
\]
since the digit sum of \(n>0\) is at most \((p-1)(1+\log_p(a_\ast^2))\). For a differential index \(\kappa\), homogeneity gives \(\sum_{a,b}\kappa_{z,ab}=|\nu_z|\). Applying \eqref{eq:p4c2-factorial} to each coordinate and using \(q_\ast-\rho(p)=\gamma(p)\) gives
\[
 v_p(\kappa![Z^\kappa]\Theta)\ge
 \sum_z|\nu_z|(\gamma(p)-y'_z)-fN_iN_{i+1}(2l+1).
\]
Frame coefficients are integral combinations of these by \ref{B3.11}. Sum the dual values and \eqref{eq:p4c2-odd-weight} over kind-\(BC\) orbits. Each orbit contributes its base value once, equivalently its position sum divided by \(n_P\). Since \(2l+1\le3l\), the result is \(L_{\rm odd}\). The positive kind-\(A\) average in \(I^*_{j,p}\) covers that part of \(F_{\rm inert}\). The cubic weight loss is in \(X^*_{j,p}\), and \(c_{e,6}\ge12\) covers the \(8/p\) weight loss in \eqref{eq:p4c2-orbit-error} and the logarithmic loss in \(X^*_{j,p}\).
\end{proof}

\subsection*{Proof of \ref{B4.13a}}\label{prf:B4.13a}\phantomsection
\begin{proof}
For the arithmetic assertion let \(L_j^{\rm loc}\) be the lattice in the local charts and \(\jmath:K\hookrightarrow\mathcal K\) the coefficient embedding. At each \(1\le j\le m\) choose a finite set of full lattice tuples meeting every level-group orbit in each component, as in Definition~\ref{def:p4c2-charts}. Choose frames at its representatives over the host field \(E\) by simultaneous local approximation of basis matrices. After clearing denominators this follows from the Chinese remainder theorem, since a sufficiently close matrix stays in the given right \(\mathrm{GL}_k(W_v)\) coset, with \(W_v\) the local lattice ring at that position. The block embeddings and rational level characters give the stabiliser equations over \(K\). In each rational frame take the invariant space for the stabiliser action with its character, and intersect with the framed integral lattice over \(\mathcal O_{K,(w)}\). Clearing denominators expresses this intersection as a kernel in that lattice. Flat base change to \(\mathcal O\) preserves the kernel. Evaluation at the finitely many representatives gives
\[
 L_j^{\rm loc}=\operatorname{span}_{\mathcal O}
 \{F\in L_j^{\rm loc}: F\text{ has rational readings}\}.
 \tag{4.31a}\label{eq:p4c2-rational-readings}
\]
For an empty class set this is the zero module.
\leanref{Formalism.\allowbreak{}MainNotions.\allowbreak{}G.\allowbreak{}CmpJ.\allowbreak{}StepIn.\allowbreak{}desc}{Formalism.\allowbreak{}SharpRestriction.\allowbreak{}GrowthBounds.\allowbreak{}ChartComparison.\allowbreak{}CmpJMemSpan}{58--85}

Write \(V_j^{\rm rat}\) for the rational space and \(\Lambda^{\rm rat}_{j,(w)}\) for its localised lattice. At every positive level, evaluation in the rational embedding blocks gives a semilinear map \(c_j\) with rational polynomial readings. If \(c_j(x)=0\), these readings vanish at every class representative. The rational form therefore evaluates to zero in every frame, so injectivity of the form embedding gives \(x=0\).
\leanref{Formalism.\allowbreak{}MainNotions.\allowbreak{}G.\allowbreak{}CmpJ.\allowbreak{}stepsMu}{Formalism.\allowbreak{}SharpRestriction.\allowbreak{}GrowthBounds.\allowbreak{}ChartComparison}{128--147}
 At the designated place \(w\), the chart lattice test in Definition~\ref{def:p4c2-charts} applies. The normalising factors for the determinant twist are units there, so integral local readings characterise membership in \(\Lambda^{\rm rat}_{j,(w)}\). For a \(K\)-basis \(b_i\), the images \(c_j(b_i)\) are independent over \(\mathcal K\). Indeed a nonzero relation, after applying to the reading equations a \(K\)-linear functional on \(\mathcal K\) nonzero on one of its coefficients, would give a nonzero rational relation. The images lie in \(\mathcal K L_j^{\rm loc}\), by scaling into the integral test. The dimension hypothesis gives \(\dim_K V_j^{\rm rat}=\operatorname{rk}_{\mathcal O}L_j^{\rm loc}\) for \(1\le j\le m\), so the independent images span over \(\mathcal K\). For \(F=\sum_i a_i c_j(b_i)\) with rational readings, apply a \(K\)-linear retraction \(\pi:\mathcal K\to K\) of \(\jmath\) to each reading. Separation gives \(F=c_j(\sum_i\pi(a_i)b_i)\). For integral \(F\), the test puts this preimage in \(\Lambda^{\rm rat}_{j,(w)}\). Together with \eqref{eq:p4c2-rational-readings} this yields
\[
 L_j^{\rm loc}=\operatorname{span}_{\mathcal O}c_j(\Lambda^{\rm rat}_{j,(w)}).
 \tag{4.31b}\label{eq:p4c2-chart-lattice}
\]
At level zero this equality is the input comparison.
\leanref{Formalism.\allowbreak{}MainNotions.\allowbreak{}G.\allowbreak{}CmpJ.\allowbreak{}StepsMu.\allowbreak{}cmp\_\allowbreak{}int}{Formalism.\allowbreak{}SharpRestriction.\allowbreak{}GrowthBounds.\allowbreak{}ChartComparison.\allowbreak{}ComparisonZeroWeightBase}{201--204}

At \(j=m\) take a basis of \(\Lambda^{\rm rat}_{m,(w)}\). A bound \(s\) on \(c_m(y)\) gives \(v_{\mathcal O}(\jmath(a_i))\ge\lceil v_{\mathcal O}(p)s\rceil\) on the nonzero coordinates \(a_i\) of \(y\), by \eqref{eq:p4c2-chart-lattice}. Since \(a_i^{e_w}/p^{v_w(a_i)}\) is a \(w\)-unit, its image is a unit and \(e_wv_{\mathcal O}(\jmath(a_i))=v_{\mathcal O}(p)v_w(a_i)\). Thus \(v_w(a_i)\ge e_ws\). Taking a coordinate of least valuation as the scalar, or zero for \(y=0\), proves the arithmetic descent.
\leanref{Formalism.\allowbreak{}MainNotions.\allowbreak{}G.\allowbreak{}CmpJ.\allowbreak{}CmpUpMu.\allowbreak{}val\_\allowbreak{}of}{Formalism.\allowbreak{}SharpRestriction.\allowbreak{}GrowthBounds.\allowbreak{}ChartComparison.\allowbreak{}ComparisonZeroWeightBase}{313--332}

For the general implication, start at \(y_0=e^\mu_0x\) and take rational representatives \(y_{j+1}\) of the successive lifts. Intertwining gives \(c_{j+1}(y_{j+1})=K_j^{\rm loc}c_j(y_j)\), so \(y_m\) represents the complex column and \(c_m(y_m)\) is the local output. Apply the assumed terminal lattice bound to \(y_m\). At \(\mu_0\) the inverse reference isomorphism at \(m\) carries this representative and the localised lattice to the reference chain.
\leanref{Formalism.\allowbreak{}MainNotions.\allowbreak{}G.\allowbreak{}CmpJ.\allowbreak{}chainMu\_\allowbreak{}cmp\_\allowbreak{}K}{Formalism.\allowbreak{}SharpRestriction.\allowbreak{}GrowthBounds.\allowbreak{}ChartComparison.\allowbreak{}ComparisonZeroWeightBase}{115--120}
\leanref{GlobalTheory.\allowbreak{}Proof.\allowbreak{}LatticeValuationBound.\allowbreak{}exists\_\allowbreak{}chain}{Development.\allowbreak{}AbelianVarietyTheory.\allowbreak{}Proofs.\allowbreak{}SatakeParameters.\allowbreak{}LatticeValuationBound}{27--40}
\leanref{GlobalTheory.\allowbreak{}Proof.\allowbreak{}ColumnWeightComparison.\allowbreak{}colGe0\_\allowbreak{}of\_\allowbreak{}colGeTw}{Development.\allowbreak{}AbelianVarietyTheory.\allowbreak{}Proofs.\allowbreak{}SatakeParameters.\allowbreak{}ColumnWeightComparison}{25--35}
\end{proof}

\subsection*{Proof of \ref{B4.14}}\label{prf:B4.14}\phantomsection
\begin{proof}
The ordinary bound is \ref{B4.13}, and the zero gain follows from \ref{B4.08}. The local conclusion follows from \ref{B4.09}--\ref{B4.10}. If \(x\in\mathcal N_{u,w}\) through a permitted model, the common filling and position map carry the tie to the twist. Use \eqref{eq:p4c2-local-np} on the common lattice and descend by \eqref{eq:p4c2-descend-terminal}. The counted gain dominates the served gain by \ref{B4.08}.

Otherwise \(W_p=\varnothing\) and \(x\in S^{\rm ram}_{u,w}\). Under \eqref{eq:p4c2-zero-gain} use \eqref{eq:p4c2-ordinary-column}, using the base-term identity from \ref{B3.54k}(3). In the remaining cases \(p\le X''\) is nonexceptional and signal or nonzero even-order inert. Formula \eqref{eq:p4c2-counted-gain} gives \(g^{\rm s}_{p,j}=g_{0,S}(p,j)\). Apply \eqref{eq:p4c2-ordinary-gain}.
\end{proof}

\subsection*{Proof of \ref{B4.15}}\label{prf:B4.15}\phantomsection
\begin{proof}
Take \(x\ge3\) and put \(U=mh\eta\). By \eqref{eq:p3c5-bottom-window}, \(b_T(t)\le5b'/4\). In the row formula of \ref{B4.03}, interlacing bounds the initial word weights at indices \(1,\ldots,k-1\) by \(S_0\), and at index zero off the top by \(3S_0\), since the first base weight is at most \(2S_0\) above the interlacing weight. Thus, for \(w=w_j(\tau)\), the first \(k\) rows obey \(\nu_i\le S_0+w\) for \(1\le i<k\) and \(\nu_0\le3S_0+w\) off the top. The later rows at \(k+\ell h+a\), \(0\le a<h,\ \ell\le j-2\), equal \(w_j(\tau)-w_{\ell+1}(\tau)-(j-\ell-1)h\). The windows have width at most \(\eta+1\), and \(h\ge1\), so each such row is at most \(\eta\) and their sum is at most \(2U\). With \(A=\sum_i\nu_i\), this gives
\[
 A-\nu_0+\nu_1\le k(S_0+w)+2U,\qquad
 A\le(k+2)S_0+kw+2U\quad\text{off the top}.
 \tag{4.37a}\label{eq:p4c2-size-rows}
\]
The sheet-zero weight on \(TT\) is at most \(b_T+b'/x+\eta\le23b'/12\). The first row on sheet one is at most \(3S_0+b+\eta+1\) at thick positions and \(3S_0+z+\eta\) at thin positions. In \eqref{eq:p4c2-size-rows} these give contraction bounds \(23kb'/12+5z+2b\) and \(23kb'/12+6z\), respectively. They are at most \(4kb'+2b\) and \(4kb'\) since \(z\le b'/3\) and \(k\ge5\). Off \(TT\), sheet zero is thin and \(\mathrm{ct}_{u,j,t}\le(k+2)S_0+2kz+2U\le4kb'\). Sum over \(n=f/2\) positions to obtain \eqref{eq:p4c2-size-contractions}. The stabiliser size assumption gives \eqref{eq:p4c2-size-denominator} for \(C\ge C_D\).
\leanref{PI.\allowbreak{}SzCt.\allowbreak{}UnitaryGroupMassFormula.\allowbreak{}ct\_\allowbreak{}le}{Formalism.\allowbreak{}PeriodIntegrals.\allowbreak{}UnitaryGroupMassFormula.\allowbreak{}SizeCount.\allowbreak{}SzCtOf}{9--83}

For the prime sum put \(c_1=12+5|c_{e,6}|\) and
\[
 \mathcal R=N_m^2(\log(8nN_ma_\ast)+\log Y)
 +|c_{\mathrm{dl},1}|N_m^2\log(2a_\ast)+2|c_{\mathrm{dl},2}|(\log(2a_\ast))^2.
 \tag{4.41}\label{eq:p4c2-size-remainder}
\]
Here \(b=e^{x^\varepsilon}\), \(a_\ast=xb\), \(X'=e^{x^{\varepsilon/3}}\), \(b'=a_\ast/X'\) and \(Y=xnX'\). The size bounds give \(n\le x^3\), \(N_m^2\le e^{x^{\varepsilon/3}}\), so \(Y\mathcal R\le C_\varepsilon x^{4+2\varepsilon}e^{2x^{\varepsilon/3}}\) for a family constant at large \(x\). Consequently
\[
 b'\le b,\qquad 0\le Y\mathcal R\le b,\qquad \log Y\ge1
 \tag{4.42}\label{eq:p4c2-size-scale}
\]
eventually. The weight windows then give \(0\le w_j(z')\le5b\) at every position \(z'\), and \(S_0\le b\).

Subtract the kind-\(A\) positive part from \(I^*_{j,p}\). By \ref{B4.03} the rest is \eqref{eq:p4c2-orbit-error} plus \(X^*_{j,p}\). The first two terms of \eqref{eq:p4c2-orbit-error} are bounded by \(c_1fb/p\), since \(\gamma(p)\le2/p\) and \(\#TT_{BC},\#BC\le f\). After multiplication by \(\log p\), the other terms are bounded by \(f\mathcal R\) using \(n_P\le2n,\ p\le Y,\ N_i\le N_m\) and \(1/\log p\le2\). Hence
\[
 \log p\,(I^*_{j,p}-[\overline F_{A,j}(p)]_+)
 \le c_1fb\,\frac{\log p}{p}+f\mathcal R+\log p\,X^*_{j,p}.
 \tag{4.43}\label{eq:p4c2-size-single-prime}
\]

For primes supporting \(X^*_{j,p}\), the inert shift \(h_P\ne0\) has odd order. Let \(2^{e+2}\) be the modulus of the real cyclotomic factor in the host construction and put \(q_e=2^e\). The projection of the second Frobenius component to \((\Z/2^{e+2}\Z)^\times/\{\pm1\}\) is trivial because that quotient is a \(2\)-group and the component has odd order. Arithmetic Frobenius acts on the cyclotomic quotient through the class of \(p\), so \(p\equiv\pm1\bmod 4q_e\). The lower bound on \(n\) gives \(n\ge x^{2+\theta}\), and \(n\le o\,2^e\) with \(o\le x^\theta\) gives \(q_e\ge x^2\). Also \(k\le x^{9/2+10\varepsilon+\theta}\le x^5\). For every finite set \(S'\) of positive integers in the two classes,
\[
 \sum_{a\in S'}\frac{\log a}{a^3}\le\frac{\log q_e}{q_e^3}.
 \tag{4.44}\label{eq:p4c2-progression}
\]
Indeed \(1\) contributes zero, and for \(a=4q_er\pm1\), \(r\ge1\), the bounds \(\log a\le2(a/q_e)\log q_e\) and \(a\ge3q_er\) give \(\log a/a^3\le\log q_e/(4q_e^3r^2)\). Sum both signs with \(\sum_{r\ge1}r^{-2}\le2\). Using \(k\le q_e^{5/2}\) and \(\log q_e\le2\sqrt{q_e}\), the cubic term of the weighted \(X^*_{j,p}\) sum is at most \(10fb\). Its logarithmic term is at most \(|c_{e,6}|fb\) by \(\#S\le Y\) and \(YN_m^2\log a_\ast\le b\). Therefore
\[
 \sum_{p\in S}\log p\,X^*_{j,p}\le(10+|c_{e,6}|)fb\log Y.
 \tag{4.45}\label{eq:p4c2-size-odd}
\]
Rosser and Schoenfeld give \(\sum_{p\le Y}\log(p)/p<\log Y\) for \(Y>1\) \cite[(3.24)]{rosser-schoenfeld-1962}. Nonnegativity gives the same upper bound for a subset of the primes. Sum \eqref{eq:p4c2-size-single-prime} using \(\#S\,\mathcal R\le b\) and \eqref{eq:p4c2-size-odd}. Take \(C=\max\{C_D,c_1+1+10+|c_{e,6}|,0\}\) and \(x_0\) above the size thresholds and \(x_D\).
\end{proof}

\subsection*{Proof of \ref{B4.16}}\label{prf:B4.16}\phantomsection
\begin{proof}
Combine the closed conditions by \eqref{eq:p4c2-closed}. For the numerical comparison take \(\varepsilon_1=1\) and zero staged thresholds. The lower bound on \(k\) in \ref{B4.15} gives \(k\ge x\) for \(x\ge1\), hence
\[
 r_G=1+\left\lceil8(4+8\varepsilon)/\varepsilon\right\rceil<32/\varepsilon+66\le k
 \tag{4.48}\label{eq:p4c2-rank}
\]
for \(x\ge x_n=\max\{x_1+1,2,32/\varepsilon+66\}\). With \(K_{\delta,1}\ge5,\ K_{\delta,2}\ge4,\ K_{\rm Pr}\ge1\), the family sizes give the hypotheses of \ref{B3.53}. The parameter choice and nonnegative constants give those of \ref{B3.54}.
\leanref{Formalism.\allowbreak{}MainNotions.\allowbreak{}G.\allowbreak{}rG\_\allowbreak{}lt\_\allowbreak{}k}{Formalism.\allowbreak{}SharpRestriction.\allowbreak{}GrowthBounds.\allowbreak{}PlaneCountTheoremWithConstants}{1798--1826}
\leanref{Formalism.\allowbreak{}MainNotions.\allowbreak{}G.\allowbreak{}numCmpF\_\allowbreak{}fixed}{Formalism.\allowbreak{}SharpRestriction.\allowbreak{}GrowthBounds.\allowbreak{}PlaneCountTheoremWithConstants}{1828--1831}

Let \(x_b,x_c\) be the thresholds of the two clauses and \(C_s,x_a\) the choices of \ref{B4.15}. For the splitting exponent \(e_z>0\) put
\[
 x_0=\max\{x_a,x_b,x_c,x_n,3^{5/e_z}\}.
 \tag{4.49}\label{eq:p4c2-service-threshold}
\]
For \(x\ge x_0\) with \(x>x_1\), the small-prime splitting condition applies to every odd prime up to \(\lfloor x^{e_z/5}\rfloor\ge3\). Tower divisibility gives \(2\mid Q_S\), so \(2\) belongs to the small-prime set. At \(3\) the splitting condition gives \(({-d}/3)=1\). The host class rule excludes both primes from the inert and signal classes. At signal primes \(n_P=2\operatorname{ord}(2h_P)\), and at even-order inert primes \(n_P=\operatorname{ord}(h_P)\), so in both cases \(p\ge5\) and \(n_P\) is even and at least two. The level primes \(g_0,g_1\) split and \(d\) is the discriminant prime, so \(g_0,g_1\ne d\).

Form \(Y'\) and \(W\) from the first clause. Under the error hypotheses of \ref{B4.03} it supplies the matching systems and \eqref{eq:p4c2-complement} on all its specified twists and places. The filling and representative map give \eqref{eq:p4c2-small-prime} by \ref{B3.53}, and \ref{B3.54} gives \eqref{eq:p4c2-large-prime}. Apply \ref{B4.09}--\ref{B4.10}. The terminal hypothesis of \ref{B4.14} applies to every matching admissible system and kernel, even at places with \(p\in Q'(\mu)\). There the formal-group theorem of \ref{B3.45} gives the terminal estimate of \ref{B3.42}.
\leanref{Formalism.\allowbreak{}MainNotions.\allowbreak{}Pk.\allowbreak{}SysTw}{Formalism.\allowbreak{}SharpRestriction.\allowbreak{}SmoothDesignExistence.\allowbreak{}SumOverPrimesHypotheses}{116--130}
\leanref{Formalism.\allowbreak{}MainNotions.\allowbreak{}WeilJacquetShalikaLaw.\allowbreak{}local\_\allowbreak{}theorem\_\allowbreak{}of\_\allowbreak{}statement}{Formalism.\allowbreak{}SharpRestriction.\allowbreak{}GrowthBounds.\allowbreak{}WeilJacquetShalikaLaw}{371--375}

The second clause gives \ref{B4.10a} at \(\sigma=1\), and \ref{B4.01b} supplies the three letter estimates at \((C_6,C)=(11,1)\). The fixed constants \(c_{e,6}=12,\ c_{\rm dig}=3,\ c_{\mathrm{dl},1}=5,\ c_{\mathrm{dl},2}\ge0\) satisfy \eqref{eq:p4c2-constants}. Thus the orbit and kernel data satisfy Definition~\ref{def:p4c2-orbit}. The restriction and addition laws and the orbit estimates follow from \ref{B4.11a}. Apply \ref{B4.11}--\ref{B4.12} and then \ref{B4.14} to \(Y'\). The filling and position map are common to \(Y\) and \(Y'\), as is the span of the rational ramified subspace. The base terms also agree. Thus the three conditions of \ref{B3.54k} and membership in the enhanced input set of \ref{B4.07} carry over to \(Y'\). The errors use the same global data. Hence \eqref{eq:p4c2-enhanced-column} and \eqref{eq:p4c2-ordinary-column}, together with \ref{B4.15}, give \eqref{eq:p4c2-servf} with \(C=C_s\) and threshold \(x_0\).
\leanref{Formalism.\allowbreak{}MainNotions.\allowbreak{}Pk.\allowbreak{}serv\_\allowbreak{}re}{Formalism.\allowbreak{}SharpRestriction.\allowbreak{}SmoothDesignExistence.\allowbreak{}SumOverPrimesHypotheses}{98--107}

For a prescribed selection, the first closed condition supplies the same levels, chains and comparisons. For \(x_c\) take the threshold assigned by the selection to the family \(\mathcal F\in\mathscr C\). For \(x\ge x_c,\ x>x_1,\ i\in I_x\) the selected member's kernel has the required properties.
\end{proof}

\section*{Proofs for Chapter~\ref{ch:columns-to-periods}: From columns to periods}

\subsection*{Proof of \ref{B4.18a}}\label{prf:B4.18a}\phantomsection
\begin{proof}
In the assumed reference and twisted evaluations, the integral interlacing
map gives the lowered representative in local class coordinates.
Reference frame changes multiply it by inverse finite-order
character values in \(\iota(K)\), which are roots of unity in \(K\) and hence local units at every finite place. The terminal lattice at
one finite place spans the terminal \(K\)-space, so these representatives
extend linearly. Injectivity of the point embedding identifies the
extension with the local representatives at all places. Intersect the
local point lattices.

Lift an \(\mathcal O\)-integral expression for \(H\) on the row generators to
their rational lifts. At each allowed place \(w\), the lifted row polynomial \(r_c\) lies in
the \(w\)-local dual of the point lattice for the Weyl pairing, so
\(b_c=\langle x_c,r_c\rangle\) is integral at \(w\). The Petersson class weight
is \(1/h_c\). The point stabilizer injects into the terminal
stabilizer, whose order divides the chain denominator. The reference chain of word \(u\) has denominator \(\Delta_u\), so the common-denominator condition gives \(h_c\mid\Delta_u\) on the twists as well. Thus
\(\Delta_u y=\sum_c(\Delta_u/h_c)b_c\) is integral at \(w\).
\end{proof}

\subsection*{Proof of \ref{B4.19a}}\label{prf:B4.19a}\phantomsection
\begin{proof}
Let \(e\) identify the complex input space with the ramified space of initial reference class functions. The rational ramified space is realized as a subspace of \(V^0_{u,0}\). Its embedding in the complex input space sends \(\phi\) to \(e^{-1}(\jmath^0_{u,0}(j\phi))\). Applying \(e\) cancels the inverse, so its initial reference class function is \(\jmath^0_{u,0}(j\phi)\).

For the reference tuple the period on a complex input \(\xi\) is defined by applying the reference column to \(e\xi\), evaluating and lowering its output, and pairing against the inverse reference test identification applied to \(H\). For a different tuple its definition applies the twisted column to \(E_\mu e\xi\) and uses the inverse twisted test identification on \(H\). Substitution of the preceding expression for \(\xi\) gives the displayed formula in each case. On an embedded rational terminal vector the same evaluation, lowering and test identification define \(\mathcal T_u^\mu\) in \ref{B4.18}.
\end{proof}

\subsection*{Proof of \ref{B4.19}}\label{prf:B4.19}\phantomsection
\begin{proof}
By \ref{B4.06} the column bound supplies \(a\in K\) and \(z\in(\Lambda^\mu_{u,m})_{(w)}\) with
\[
 C^\mu E_\mu\jmath^0_{u,0}(j\phi)=\jmath^\mu_{u,m}(az),
 \qquad a=0\quad\text{or}\quad v_w(a)\geq e_ws .
\]
Here and at the reference tuple we use the identifications of \ref{B4.19a}. Since \(H\) is admissible and \(p_w\notin Q'(\mu)\), \ref{B4.18} applied to \(z,H\) gives \(y_0\in K\) with \(\iota(y_0)=\mathcal T_u^\mu(z,H)\) and \(v_w(y_0)\geq-e_w\delta_{u,w}\) if \(y_0\ne0\). Evaluation and lowering are linear, and the pairing with the fixed test is linear in the terminal variable. Thus \ref{B4.19a} gives
\[
 P_u^\mu(\phi,H)=\mathcal T_u^\mu(az,H)=\iota(a y_0).
\]
Take \(y=ay_0\). When \(y\ne0\), add the valuations of its two factors.
\end{proof}

\subsection*{Proof of \ref{B4.21}}\label{prf:B4.21}\phantomsection
\begin{proof}
The stabilizer condition on the design gives a cutoff \(x_s\) beyond which every prime divisor of every \(\Delta_u\) is at most \(2fN_m+1\). The level condition gives a cutoff \(x_T\) beyond which \(b^{m/2}\le T_u\) for every word. Both cutoffs are uniform in the member index at \(x\). Take \(x_0=\max(x_s,x_T)\). If \(x\ge x_0\) and \(x>x_1\), the member satisfies the size condition of Lemma~\ref{B4.21b}. Hence \(T_u>2fN_m+1\). Since \(T_u\) is prime, the stabilizer condition excludes \(T_u\mid\Delta_u\). In the definition of closedness in \ref{B4.16}, take the \(\varepsilon\)-cutoff to be \(1\) and all four size thresholds to be zero. The preceding choice of \(x_0\) then works for each family in the design class.
\end{proof}

\subsection*{Proof of \ref{B4.21a}}\label{prf:B4.21a}\phantomsection
\begin{proof}
Each vector of \(V^r\) has a nonzero integer multiple in \(L^0\).
Since \(\mathcal O\) is Noetherian, \(L^0\cap V^r\) is finitely
generated, so \(M_{H,u}\) is full. Clearing denominators outside
\(\mathfrak p_w\) gives \(j((M_{H,u})_{(w)})=S^{\mathrm{ram}}_{u,w}\).
At a selected permitted place \ref{B4.07c} gives
\[
 p_w^{a+b}(M_{H,u})_{(w)}\subseteq B_w
 \subseteq (M_{H,u})_{(w)}
\]
for the nonnegative integers \(a,b\) of that lemma. The first
inclusion comes from the empty Hecke product. The second uses
\(N_P\subseteq M\) and the coordinate test.

At \(w\mid T_u\), frame changes in a neighbor sum act by
substitution in rational matrices, determinant powers and the
rational level character. Thus \(U_v\) preserves rational
coefficients. Rescaling markings permutes the neighbor incidences,
so \(U_v\) preserves \(V^r\), the common kernel of the diamond sums.
In the down-up calculation of \ref{B5.05a}, the coefficients at a
marked vector \(e\) are \(T_u^{k-1}\) at \(e\), zero at the other
nonzero vectors on its line and \(T_u^{k-2}\) off the line.
On the common kernel the line sums and the unrestricted trace
vanish. Thus \(D_vU_v=T_u^{k-1}I\), where \(D_v\) is the down
operator. It follows that \(U_v\) is invertible on \(V^r\).
Under the complex embedding, \(N_{\lambda,u,w}=G S^{\mathrm{ram}}_{u,w}\)
with \(G=T_u^{a_0}U^m\) a rational automorphism of \(V^r\).
Clearing denominators of \(G\) and \(G^{-1}\) on generators of
\(L^0\cap V^r\) gives an integer \(n_w>0\) with
\(n_w S^{\mathrm{ram}}_{u,w}\subseteq G S^{\mathrm{ram}}_{u,w}
\subseteq n_w^{-1} S^{\mathrm{ram}}_{u,w}\). For an empty product take
\(G=1\).
\leanref{GlobalTheory.\allowbreak{}LevelHecke.\allowbreak{}uop\_\allowbreak{}rat}{Development.\allowbreak{}AbelianVarietyTheory.\allowbreak{}Proofs.\allowbreak{}IwasawaDecomposition.\allowbreak{}UpwardOperatorPreservesRationality}{147--183}
\leanref{GlobalTheory.\allowbreak{}LevelCommensurable.\allowbreak{}commensurable}{Development.\allowbreak{}AbelianVarietyTheory.\allowbreak{}Proofs.\allowbreak{}IwasawaDecomposition.\allowbreak{}RamifiedFormsCommensurableOfUStableInjective}{36--93}

The exceptional set \(S\) of selected permitted places and places
above \(T_u\) is finite, since nonempty selection requires \(p_w<X'\)
by \ref{B3.51}. At its places the preceding inclusions give
\[
 n_w(M_{H,u})_{(w)}\subseteq B_w
 \subseteq n_w^{-1}(M_{H,u})_{(w)}
\]
for positive integers \(n_w\). With \(n=\prod_{w\in S}n_w\) this also
gives \(nM_{H,u}\subseteq N_u\subseteq n^{-1}M_{H,u}\).
Each \(B_w\) in the display is full and finitely generated over
\(\mathcal O_{(w)}\). The \(\mathcal O\)-span of a set of its local
generators is a full lattice in \(V_u\) with localization \(B_w\).
Use these lattices at the exceptional places and \(M_{H,u}\)
elsewhere in \ref{B4.35} to obtain the intersection and all its
localizations.
\leanref{GlobalTheory.\allowbreak{}EX.\allowbreak{}Glue.\allowbreak{}exists\_\allowbreak{}glue}{Development.\allowbreak{}AbelianVarietyTheory.\allowbreak{}Proofs.\allowbreak{}GlobalInterpolation.\allowbreak{}MemLocSub}{69--114}
\leanref{GlobalTheory.\allowbreak{}EX.\allowbreak{}Glue.\allowbreak{}inLoc\_\allowbreak{}iInf}{Development.\allowbreak{}AbelianVarietyTheory.\allowbreak{}Proofs.\allowbreak{}GlobalInterpolation.\allowbreak{}MemLocSub}{51--68}

The normalized local colength is
\[
 \ell_{u,w}=\frac{1}{e_w}
 \operatorname{length}_{\mathcal O_{(w)}}
 \left(\bigl(M_{H,u}/(N_u\cap M_{H,u})\bigr)_{(w)}\right).
\]
The quotient is killed by \(n\). Away from \(T_u\) the denominator
lattice localizes to \((N_u)_{(w)}\). At \(T_u\) the intersection
in the formula gives the denominator.
\leanref{GlobalTheory.\allowbreak{}EX.\allowbreak{}Led.\allowbreak{}setLat}{Development.\allowbreak{}AbelianVarietyTheory.\allowbreak{}Proofs.\allowbreak{}GlobalInterpolation.\allowbreak{}ExistsLedPart}{187--191}
\leanref{GlobalTheory.\allowbreak{}EX.\allowbreak{}Led.\allowbreak{}exists\_\allowbreak{}ledPart}{Development.\allowbreak{}AbelianVarietyTheory.\allowbreak{}Proofs.\allowbreak{}GlobalInterpolation.\allowbreak{}ExistsLedPart}{195--226}
\end{proof}

\subsection*{Proof of \ref{B4.22}}\label{prf:B4.22}\phantomsection
\begin{proof}
The level-prime description of \((N_u)_{(w)}\) gives \(j_u\phi\in N_{\lambda,u,w}\).
The formula in \ref{B4.03}, with \(m\) even and \(t\) odd, gives
\(n_0=mT^*_{u,w}\in\N\). At \(\mu_0\), take \(c=T_u^{n_0}\)
in \ref{B4.22a} and divide its terminal representative by \(c\).
On another used tuple the twisted bound gives exponent \(-n_0\).
Apply \ref{B4.19} at \(\mu_w\) with \(\delta_{u,w}=0\).
At \(w\mid T_u\), the identities of \ref{B4.20} give zero raw gain and zero base term at each lift, with loss \(T^*_{u,w}\). The screened gain
also vanishes. Summing over the \(m\) lifts gives \(S_{u,w}=-n_0\).
\end{proof}

\subsection*{Proof of \ref{B4.22d}}\label{prf:B4.22d}\phantomsection
\begin{proof}
Each \(b_\nu\) has norm at most one by continuity from \(\N^r\). For a finite expansion \(f=\sum_\nu a_\nu b_\nu\), its coefficient \(a_\nu\) is \((\Delta_1^{\nu_1}\cdots\Delta_r^{\nu_r} f)(0)\), where \(\Delta_j f(c)=f(c+e_j)-f(c)\). Hence \(\|f\|_\infty=\max_\nu|a_\nu|_p\). These finite expansions are dense. In one variable this is Mahler's interpolation theorem \cite{mahler-1958}; it applies in particular to the \(\Q_p\)-valued indicators of residue disks. Products approximate the indicators of residue boxes in \(r\) variables. Locally constant functions, finite linear combinations of these indicators with coefficients in \(\C_p\), approximate every continuous function by compactness. The rule \(f\mapsto\sum_\nu a_\nu A_\nu\) now extends uniquely by completeness.

Choose \(k\) so large that on each box modulo \(p^k\) the oscillation of every \(g_i\) is less than \(\eta\). Take its representative \(c\in\{0,\ldots,p^k-1\}^r\) and put \(d_c=\mu(\mathbf1_{c+p^k\Z_p^r})\). The assertion follows by applying \(\mu\) to the difference between \(g_i\) and the corresponding step function. For \(r=0\) there is one box.
\end{proof}

\subsection*{Proof of \ref{B4.22e}}\label{prf:B4.22e}\phantomsection
\begin{proof}
Use Lebesgue measure in \(H\)-orthonormal complex coordinates. The integral of \(e^{-\pi H(z,z)}\) over the period space is one. A lattice self-dual for \(B=\Im H\) has covolume one for this measure: in real orthonormal coordinates the matrix of \(B\) has determinant one. The trace dual of \(\mathfrak b^{-1}\) is \(\mathfrak D_{E/K_d}^{-1}\bar{\mathfrak b}=\mathfrak b^{-1}\); the area factor gives alternating pairing one on an oriented basis of \(\Lambda\). Hence \(P_1\) is self-dual.

Direct multiplication gives \(U_xU_y=e^{-\pi i B(x,y)}U_{x+y}\). On each indicated lattice the semicharacter makes \(W_xW_y=W_{x+y}\). For self-dual \(L,L'\) the commutator induces a perfect pairing of \(L/(L\cap L')\) with \(L'/(L\cap L')\). Summing it shows \(I_{L',L}I_{L,L'}=1\) on \(L\)-fixed sections. Change of variable in the Gaussian norm on a common period domain identifies the reverse average with the adjoint, for the norm with period-domain measure normalized to mass one. The averages are therefore isometries. The fixed holomorphic space for \(\Lambda\) is a line: division by its theta section gives an elliptic function with at most one simple pole, which the residue theorem excludes. Applying this in one variable at a time proves the product statement. The inverse averages transport this line to every commensurable self-dual period lattice.

Diagonalize the positive trace form over \(K_d\). For a positive rational diagonal coefficient \(A/B_0\), write \(AB_0\) as the norm of an integral quaternion by the four-square theorem. Its left multiplication matrix divided by \(B_0\) has transpose times itself equal to \((A/B_0)I_4\). This gives an isometry \(g_0\) on four copies. Write \(P_0=g_0P_1^4\), \(P=\Lambda^{4f}\). At a split prime \(l\) choose one factor \(\mathfrak l\), a generator \(\pi_l\) of \(\mathfrak l^h\), with \(h\) the class number, and a root \(r_l^{2h}=\pi_l/\bar\pi_l\). Let \(y_l\) be the determinant order at \(\mathfrak l\) of a local unitary map taking \(P_0\) to \(P\). Such a map sends a basis on the first half to a basis there and acts by the inverse adjoint on the second half. The order is independent of the bases. Put \(\kappa_0=\prod_l r_l^{-y_l}\); the two lattices agree away from finitely many primes. The fourth-power normalization chooses \(\Theta\) so that
\[
 \Theta^{\otimes4}(z)=\kappa_0(I_{P,P_0}\theta_\Lambda^{\otimes4f})(g_0z).
\]
The fourfold line is the tensor product of four one-entry lines, so this chooses a fourth root.

At \(p\) the maps taking \(P_0\) to \(P\) form an open coset. Choose a rational unitary \(u\) in it with no eigenvalue one. Indeed multiplying a local map by norm-one scalars arbitrarily close to one avoids eigenvalues \(1,-1\). In the Cayley chart the resulting open set pulls back to an open set in the vector space of skew-Hermitian matrices, where rational points are dense. Set \(g=ug_0\), \(Q=uP_0\). Since \(\det u\) has orders \(y_p,-y_p\) at the factors over \(p\), the defining equation of \(r_p\), and the fact that \(r_l\) is a \(p\)-unit for \(l\ne p\), give
\begin{equation}\label{eq:B4.22c-kappa}
 |\jmath(\kappa_0)|_p^2|\jmath(\det u)|_p=1 .
\end{equation}

We compute the pullback factor. Let \(v\) be unitary on a period space, let \(L,vL\) be commensurable, and assume \(A=1-v^{-1}\) invertible. Put \(L_0=L\cap vL\), \(J=[L:L_0]=[vL:L_0]\), \(D=\det A\), \(P_A=A^{-1}-1/2\). Pullback by \(v^{-1}\) and \(I_{L,vL}\) carry the same theta line isomorphically to the \(vL\)-line. Write their ratio as \(c_v\). Define
\[
 K(w,z)=\sum_{x\in L}W_x^{(w)} e^{\pi H(w,z)},\qquad
 \phi(x)=\alpha(x)e^{-\pi H(P_Ax,x)}.
\]
For fixed \(z\), \(K(\cdot,z)\) belongs to the \(L\)-line by reindexing. The sum and the following integrals converge absolutely: on bounded sets of \(w,z\) the terms are bounded by a constant times \(e^{-c\|x\|^2+C\|x\|}\) with \(c>0\).

Apply the two maps in the first variable, set \(w=z\), multiply by \(e^{-\pi H(z,z)}\) and integrate over a fundamental domain for \(L_0\). In the term of the average with representative \(y\in vL\), the summand indexed by \(x\in L\) is a constant times \(e^{2\pi i B(y+x,z)}\). Since \(B(y+x,L_0)\subseteq\Z\), its integral vanishes unless \(y+x=0\). Choose zero as the representative of the trivial coset. The surviving term has coefficient one. Thus the integral of the average is \(J/\sqrt J\).

For the pullback the summand indexed by \(x\) is, by expanding the Hermitian form,
\[
 \alpha(x)\,(U_x^{(w)} e^{\pi H(w,z)})\big|_{w=v^{-1}z}
 e^{-\pi H(z,z)}
 =\phi(x)\exp\{-\pi H(A(z-A^{-1}x),z-A^{-1}x)\}.
\]
For \(h\in L_0\), \(Ah=h-v^{-1}h\in L\). The semicharacter law and \(\alpha(v^{-1}h)=\alpha(h)\) give \(\phi(Ah)=1\). The cross factor in \(\phi(x+Ah)\) is \(e^{-2\pi iB(A^{-1}Ah,x)}=1\), using symmetry of the polarization on \(L\). Thus the sum over \(x\) in one coset modulo \(AL_0\) unfolds to the integral over the whole space. In an orthonormal eigenbasis of \(v\), each eigenvalue \(\lambda\) of \(A\) has positive real part, and \(2\pi\int_0^\infty e^{-\pi\lambda r^2}r\,dr=\lambda^{-1}\). The integral is \(D^{-1}\). Since \([L:AL_0]=|D|^2J\), comparison gives
\[
 c_v=\frac{|D|}{D}G,\qquad
 G=\frac{1}{\sqrt{|V|}}\sum_{x\in V}\phi(x),\qquad V=L/AL_0 .
\]

For completeness, \(A^{-1}+(A^{-1})^*=1\) makes \(P_A\) skew-adjoint. Substituting in the preceding expression for \(\phi\) and using the semicharacter law on \(L\) gives polarization \(e^{-2\pi i B(A^{-1}x,y)}\). This is trivial for every \(y\in L\) exactly when \(A^{-1}x\in L\). In that case \(v^{-1}A^{-1}x=A^{-1}x-x\in L\), so \(A^{-1}x\in L_0\). This proves nondegeneracy on \(V\). Also \(\phi(nx)=\phi(x)^{n^2}\), since \(\alpha(x)=\pm1\) on \(L\). The finite quadratic-module Gauss formula gives \(G^8=1\) \cite[Proposition~2.11 and Lemmas~3.6--3.8]{stromberg-2013}.

There is also a factor from composing averages. For \(L_1,L_2,L_3\) the composition is a scalar \(\mu\) times \(I_{L_1,L_3}\). Project the composition back to the \(L_1\)-line by the inverse average. Commutator orthogonality retains the pairs \(z\in L_3,\ y\in L_2\) with \(z+y\in L_1\). They form a finite group modulo \((t,h-t)\), \(t\in L_3\cap L_2,\ h\in L_2\cap L_1\). Expanding \(W_zW_y\) gives their phase
\[
 \alpha(z)\alpha(y)\alpha(z+y)^{-1}e^{-\pi i B(z,y)}.
\]
The group laws on the three lattices make this expression independent of representatives; its polarization is \(e^{2\pi i B(y,z')}\) and it is homogeneous quadratic. The averaging constants are positive. If its restriction to the radical were nontrivial, translation there would annihilate the sum. The composition is invertible, so the phase descends to the nondegenerate quotient. Its normalized sum has eighth power one by the same Gauss formula. Since the composition is an isometry, the positive factor relating \(\mu\) to that normalized sum is one.

Apply this with \((L_1,L_2,L_3)=(P,P_0,Q)\) and \(v=u\) acting on \(P_0\). The scalar \(c\) relating the two normalized standard sections is \(c_u\mu\). From \(A^*=-uA\) in dimension \(4f\) we get \(\bar D/D=\det u\). Consequently
\[
 c^8=(|D|/D)^8=(\det u)^4.
\]
In particular \(c\) is algebraic. Taking \(p\)-adic absolute values gives \(|\jmath(c)|_p^2=|\jmath(\det u)|_p\). Thus \(e=\kappa_0c\) has \(|\jmath(e)|_p=1\) by \eqref{eq:B4.22c-kappa}.

For the endpoint factors we need the sign of a repeated Gauss sum. For every integer \(m_0>0\) there is \(\lambda\in\mathcal O_E\) with \(\lambda\bar\lambda=N\in\Z\) and \(N\equiv-1\bmod m_0\). In the host take
\(\lambda=(a+b\sqrt{-q})(c_1+d_1\sqrt{-d})\). At an odd prime \(l\mid m_0\) choose \(t\in\{q,d\}\) prime to \(l\). The sets of squares and of elements \(-1-ty^2\) each have \((l+1)/2\) elements in \(\F_l\), so \(x^2+ty^2=-1\) has a solution. One partial derivative is a unit, giving a lift to every power of \(l\); use \(1\) in the other norm factor. At \(2\), \(q\equiv3\bmod8\), so \(-5/q\equiv1\bmod8\) has a square root in \(\Z_2\). Equivalently one lifts an odd square root from modulus \(2^k\) to \(2^{k+1}\), \(k\ge3\), by choosing between \(y\) and \(y+2^{k-1}\). Use \(x=2\) in \(x^2+qy^2=-1\). The Chinese remainder theorem on the four integer coefficients proves the claim.
\leanref{Formalism.\allowbreak{}L2.\allowbreak{}PairingIndexBounds.\allowbreak{}Eq.\allowbreak{}J.\allowbreak{}Neg.\allowbreak{}rep\_\allowbreak{}int}{Formalism.\allowbreak{}LocalToGlobal.\allowbreak{}PairingIndexBounds.\allowbreak{}SplitPairEquations.\allowbreak{}IntegralSquareRoots.\allowbreak{}HenselQuad}{137--149}

In one row take \(m_0=4|V|\). Multiplication by \(\lambda\) and by \(\bar\lambda\) preserves both period lattices and commutes with the endpoint action and \(A\). It permutes \(V\), since their product is \(-1\) on \(V\). The integer \(N\) is the same at every embedding, hence \(H(\lambda x,\lambda y)=NH(x,y)\) and \(\phi(\lambda x)=\phi(x)^N\). The polarization at \((x,x)\) equals \(\phi(x)^2\) and has order dividing \(|V|\). Thus \(\phi(\lambda x)=\overline{\phi(x)}\). The sum \(G\) is real, so \(G^8=1\) gives \(G^2=1\).

For a source change the \(b\) rows give identical one-row factors. On one row the determinant of the action on period coordinates is \(\prod_{\tau}\tau(\overline{\det\gamma})\). The row dimension \(af\) is even, whence \(c_v^2=\bar D/D=\prod_\tau\tau(\overline{\det\gamma})\). Taking \(b\) rows gives the claimed exponent \(b/2\). For a target change use \(a\) columns; the determinant on one column is \(\prod_\tau\tau(\det\gamma')\). If the action has eigenvalue one, factor it by a rational norm-one scalar chosen so that neither factor has eigenvalue one. The triangle of averages in this factorization is a product of identical row triangles, or identical column triangles. On each triangle choose \(m_0\) divisible by the orders of its finite phases and of its finite group. The same \(\lambda\) permutes that group and raises the displayed triangle phase to \(N\), hence conjugates it. Its scalar is real of absolute value one. An even number of copies has scalar one, so the factors multiply also in this case. For zero rank the empty product equals one.
\end{proof}

\subsection*{Proof of \ref{B4.22c}}\label{prf:B4.22c}\phantomsection
\begin{proof}
For commuting directional derivatives \(D_j\) in the scaled directions of the statement, the coefficient at \(Z^\nu\) in
\(\prod_j\exp(D_j\log(1+Z_j))F(0)\) is \(A_\nu(F;v)\), after transport by \(\jmath\). By \ref{B4.22d} a bound of one on these coefficients gives a measure with transform
\[
 \int_{\Z_p^r}\prod_j(1+Z_j)^{c_j}\,d\mu(c).
\]
Expressing powers in falling factorials shows that the power moments give the iterated derivatives. Products of such transforms correspond to pushing product measures forward by addition. Integral \(p\)-adic linear substitutions in the directions correspond to pushforward by the transpose on the measure variables. Both operations preserve the upper bound one on the norm.

Apply \ref{B4.22e} to the fourth tensor power. If \(nt\in P_1\), \(n>0\), \(p\nmid n\), put \(v=g(t,t,t,t)\) and \(J=[Q:P\cap Q]\). Commuting \(U_v\) past the terms of the average gives
\[
 (U_t\Theta)^{\otimes4}(z)=
 \frac{e}{\sqrt J}\sum_{x\in Q/(P\cap Q)}
 \alpha(x)e^{-\pi iB(v,x)}
 \prod_{\ell=1}^{4f}(U_{v_\ell+x_\ell}\theta_\Lambda)((gz)_\ell).
\]
Here \(nJ(v_\ell+x_\ell)\in\Lambda\). The phases are roots of unity since \(nv,x\in Q\), and \(e/\sqrt J\) is a \(p\)-adic unit. Multiplication by \(J\) puts each transformed period direction in \(P\), and \(1/J\in\Z_p\). The elliptic hypothesis and the product and substitution formulas give integral coefficients for every list in the four blocks. Use independent variables in the blocks, with the same list and index \(\nu\) in each. The coefficient at \((\nu,\nu,\nu,\nu)\) is \(A_\nu(U_t\Theta;v_\bullet)^4\). Hence the one-entry coefficients are integral. Taking products over the entries proves the assertion for the standard matrix vector.

For clarity, write \(P(L,L')=\Lambda(L,L')\), let \(L_0,L'_0\) be standard, and form the vector \(F_{L,L'}\) by averaging the standard vector first to \(P(L,L'_0)\), then to \(P(L,L')\). Its normalization is
\[
 e_{L,L'}=
 \chi(\det(L_0,L))^{b/2}\chi(\det(L'_0,L'))^{-a/2}F_{L,L'}.
\]
The relative determinant here is the determinant ideal of one lattice in a basis of the other. The character is multiplicative and
\(\chi((z))=\prod_\tau\tau(z)^{-1}\) for \(z\bar z=1\).

We record why lattice moves compose with this normalization. For moves with the other lattice standard, the triangle scalar is one by the paired Gauss calculation in \ref{B4.22e}. For moves in different variables consider the four period lattices modulo their common intersection. At every finite place where the lattices differ, the two target lattices admit the same two complementary isotropic halves. At a split place of \(E/E^+\) take the two factors exchanged by conjugation. In the host of \ref{B2.01}, the prime \(d\) satisfies \(d\equiv7\pmod{16}\), so the defining polynomial \(X^2-X+(1+d)/4\) of \(\mathcal O_{K_d}\) splits into two distinct linear factors modulo \(2\). Consequently \(2\) splits in \(K_d\), and therefore every prime of \(E^+\) above \(2\) splits in \(E/E^+\). Every inert place of \(E/E^+\) thus has odd residue characteristic. At an inert place the extension \(E_v/E^+_v\) is unramified of odd residue characteristic, and the lattices are self-dual of even rank. Reduction modulo the maximal ideal gives a nondegenerate Hermitian space of even dimension over a finite quadratic extension, which is hyperbolic. A hyperbolic pair lifts integrally: lift an isotropic basis vector to a primitive vector, adjust its Hermitian norm to zero by Hensel's lemma (as the residue characteristic is odd and the Hermitian form is unimodular), choose a second vector pairing to one with the first, and subtract a multiple of the first using surjectivity of the trace map \(\mathcal O_{E_v}\to\mathcal O_{E^+_v}\) for the unramified extension so that the second vector is also isotropic. Splitting off the hyperbolic plane inductively proves that the local space is hyperbolic. In odd residue characteristic, the two self-dual lattice norms are points of the Bruhat--Tits building of the unitary group, and a common apartment splits them simultaneously in a single Witt basis by the norm description of \cite[\S\S2.8--2.12]{bruhat-tits-1987}. Projection onto one vector of each hyperbolic pair and its complementary projection preserve both lattices. For the finite quotients it suffices to use these projections modulo a common denominator carrying either source lattice into the other. The Chinese remainder theorem gives the projections to that precision, with \(e^2=e\) and \(e^*=1-e\) on the quotients.

This reduces the finite Heisenberg calculation to a lattice \(A\) in one half and its annihilator \(A^\perp\) in the other. In the function realization, translations by \(A\) fix \(\mathbf1_A\), and characters from \(A^\perp\) fix it. For a change \(A\to D\), summation of the characters cuts \(\mathbf1_A\) down to \(\mathbf1_{A\cap D}\) with multiplicity \([A:A\cap D]\). Summation over \(D/(A\cap D)\) then gives \(\mathbf1_D\). The normalized average consequently takes \(\mathbf1_A\) to the positive multiple \(([A:A\cap D]/[D:A\cap D])^{1/2}\mathbf1_D\). The two paths round the square have the same factor by multiplicativity of relative volumes. On each isotropic half the Weyl operators multiply as translations: their Hermitian cross terms vanish on an exact half, and the congruence projections give the same operators on the finite quotients. Thus the square commutes. Interchanging successive source and target moves now transports the standard triangle identities to arbitrary endpoints.
\leanref{Formalism.\allowbreak{}L2.\allowbreak{}Ker.\allowbreak{}Cd.\allowbreak{}PairingIndexBounds.\allowbreak{}Eq.\allowbreak{}SquareSplitPairs.\allowbreak{}sqSplit\_\allowbreak{}of}{Formalism.\allowbreak{}LocalToGlobal.\allowbreak{}KernelLattices.\allowbreak{}ColumnData.\allowbreak{}ColumnDataEquations.\allowbreak{}SqLMain}{21--145}
\leanref{Formalism.\allowbreak{}L2.\allowbreak{}Ker.\allowbreak{}Cd.\allowbreak{}PairingIndexBounds.\allowbreak{}Eq.\allowbreak{}SqG.\allowbreak{}eqP\_\allowbreak{}car}{Formalism.\allowbreak{}LocalToGlobal.\allowbreak{}KernelLattices.\allowbreak{}ColumnData.\allowbreak{}ColumnDataEquations.\allowbreak{}SplitPairOfGen}{55--59}

Pullback preserves \(H,\alpha\) and carries each summation quotient bijectively onto the moved quotient. The two factors computed in \ref{B4.22e} are the changes of the character normalization on moving the standard endpoints. This gives, with \(P_{\gamma,\gamma'}F(z)=F(\gamma'^*z\gamma)\),
\begin{align}
 P_{\gamma,\gamma'}e_{L,L'}&=e_{\gamma L,\gamma'L'},\label{eq:B4.22c-equiv}\\
 e_{L,L''}&=\chi(\det(L',L''))^{-a/2}I_{P(L,L'),P(L,L'')}e_{L,L'},\label{eq:B4.22c-target}\\
 e_{L_1,L'}&=\chi(\det(L,L_1))^{b/2}I_{P(L,L'),P(L_1,L')}e_{L,L'}.\label{eq:B4.22c-source}
\end{align}

Finally move \(L,L'\) by rational unitary maps so that both agree with their standard lattices above \(p\). The basis and dual basis at the split factors give the local maps, and Cayley approximation as in \ref{B4.22e} gives the rational maps. For the moved pair both averaging indices are prime to \(p\). Its relative determinant ideals have valuation zero above \(p\). For any such norm-one ideal \(\mathfrak a\), choose \(h>0\) with \(\mathfrak a^h=(z)\). Then \(\mathfrak a^{2h}=(z/\bar z)\), so the principal-ideal character formula makes \(\chi(\mathfrak a)\) a \(p\)-adic unit. Expanding the two averages and the prime-to-\(p\) translation gives a finite sum of translates of the standard vector with integral coefficients. Each translating point has prime-to-\(p\) denominator, since the original shift and the two finite quotients do. The same denominator clearing puts the directions in the standard period lattice. The standard bound and integral substitution apply term by term. Equation \eqref{eq:B4.22c-equiv} transports both translations and directions back to \(L,L'\). The measure assertion follows from \ref{B4.22d}.
\end{proof}

\subsection*{Proof of \ref{B4.22f}}\label{prf:B4.22f}\phantomsection
\begin{proof}
At the split prime the connected ordinary formal group has a multiplicative parameter. Normalize it as follows. On the connected \(p^n\)-torsion, the multiplicative coordinate and the pairing with the period \(1\) both give primitive characters into \(\mu_{p^n}\). For the pairing, the denominator-\(p^n\) periods modulo periods are the two cyclic split factors in \(\mathcal O_{K_d}/p^n\). The alternating area pairing is perfect, and the projection of \(1\) onto the complementary factor is a generator. The congruence on \(\bar v\) selects the connected factor. The powers relating the two characters are units modulo \(p^n\), compatible under multiplication by \(p\); their inverse limit is a unit in \(\Z_p\). Changing the parameter by this unit power makes the coordinate of \(y\) exactly \(\zeta_y\). In this chart an elliptic function regular on the origin disk and with bounded expansion satisfies
\[
 \widehat{F(\,\cdot+y)}=S_{\zeta_y}\widehat F .
\]
Indeed the parameter identifies the formal group law with multiplication of \(1+Z\). Substitution first gives the formula on its convergent coordinate functions, hence on rational functions wherever the denominator is nonzero, and then on their regular germs. We compute the extra factor for the theta section.

Put \(f=\widehat\theta\). Both \(f\) and its inverse have integral coefficients. Here are the details of the coefficient bound. For a nonzero \(\alpha\in\mathcal O_{K_d}\) put \(m=\alpha\bar\alpha\) and
\[
 E_\alpha(z)=\theta(\alpha z)/\theta(z)^m .
\]
The period multipliers cancel, so \(E_\alpha\) is elliptic with \(E_\alpha(0)=1\). Its possible poles are at the nonzero half-period, and its possible zeros lie over that half-period under \(\alpha\). Neither meets the origin disk: reduction commutes with \(\alpha\), and reduction is injective on the prime-to-\(p\) torsion \cite[Ch.~VII, Proposition~3.1]{silverman-2009}. Both \(E_\alpha\) and its reciprocal have bounded expansions. Explicitly, clear their poles by products of \(X-X(P)\) in good Weierstrass coordinates. For such a pole \(P\), \(X(P)\) is integral, and in a formal parameter \(T\), \(T^2(X-X(P))\) is an integral unit series. The numerator is polynomial in the affine coordinates; multiplication by a power of \(T\) bounds its coefficients. Regularity removes the negative powers. The same argument works after the integral change to the multiplicative parameter. The Gauss norm on bounded series is multiplicative, by taking the supremum of the multiplicative Gauss norms on closed disks of radii tending to one. The two series are inverse with constant term one, so both have norm one. Choose \(\alpha\equiv1\) at the chosen prime and \(0\) at the other. Write \(a=\jmath(\alpha)\) and \([a](Z)=(1+Z)^a-1\). In
\[
 f([a](Z))=f(Z)^m\widehat E_\alpha(Z)
\]
the coefficient of the unknown \(j\)-th coefficient of \(f\), for \(j>0\), is \(a^j-m\), a unit. Starting with \(f(0)=1\) proves integrality degree by degree. Inversion preserves it.
\leanref{GlobalTheory.\allowbreak{}MomOne.\allowbreak{}ThetaQuot.\allowbreak{}theta\_\allowbreak{}quot}{Development.\allowbreak{}AbelianVarietyTheory.\allowbreak{}Proofs.\allowbreak{}JetHarmonics.\allowbreak{}ThetaNormQuotientEllipticUnitSeries}{71--140}
\leanref{GlobalTheory.\allowbreak{}MomOne.\allowbreak{}theta\_\allowbreak{}series\_\allowbreak{}integral}{Development.\allowbreak{}AbelianVarietyTheory.\allowbreak{}Proofs.\allowbreak{}JetHarmonics.\allowbreak{}CmEvenThetaSeriesChartCoeffIntegral}{15--77}

First take \(t=0\). Put \(g=\widehat{W_y\theta}\), \(h=S_{\zeta_y}f\), and \(R=g/h\). The denominator is a unit series, since \(S_{\zeta_y}(f^{-1})\) is its inverse. Also \(g(0)\ne0\): the point of \(y\) has \(p\)-power order, whereas a zero of \(\theta\) represents \(1/2\).

For every positive \(s_1\) prime to \(p\), choose \(\alpha\in s_1\mathcal O_{K_d}\) congruent to \(1\) and \(0\) at the two primes modulo \(p^{n+1}\), respectively. Then \(m=\alpha\bar\alpha\ge2\), \(p\mid m\), \(s_1\mid m\), and \(a=\jmath(\alpha)\) is a unit congruent to \(1\bmod p^n\). We have \(W_{\alpha y}\theta=W_y\theta\). \((\alpha-1)y\) is a period: the chosen component of \(\alpha-1\) and the other component of \(v\) clear its denominator. Moreover \(y\overline{(\alpha-1)y}\) is integral, since \(v\bar v\) is divisible by \(p^{2n}\). In the principal lattice \(H(z,w)=2z\bar w/\sqrt d\), \(\alpha(z)=e^{\pi i z\bar z}\); inserting the period in the composition law for \(W\) therefore gives factor one. In coordinates this last assertion uses
\(\Re h+\Im h/\sqrt d\in\Z\) for \(h\in\mathcal O_{K_d}=\Z+\Z(1+i\sqrt d)/2\), with either sign of the imaginary term.

The two factors in the definition of \(W_y\), including its constant semicharacter, give the following identity of germs:
\[
 (W_{\alpha y}\theta)(\alpha z)=(W_y\theta)(z)^m E_\alpha(z+y).
\]
Indeed \(H(\alpha z,\alpha y)=mH(z,y)\) and \(\alpha(\alpha y)=\alpha(y)^m\), where in the latter equality \(\alpha(\,\cdot\,)\) denotes the semicharacter. Thus
\[
 g([a](Z))=g(Z)^m S_{\zeta_y}\widehat E_\alpha(Z).
\]
The same equation holds for \(h\). Apply \(S_{\zeta_y}\) to the equation for \(f\) and use \(\zeta_y^a=\zeta_y\), since \(a\equiv1\bmod p^n\). Division gives \(R([a](Z))=R(Z)^m\). Its constant coefficient \(c\ne0\) satisfies \(c^{m-1}=1\). If \(r_jZ^j\) is its first nonconstant term, comparison gives \(a^jr_j=mc^{m-1}r_j\), impossible since \(a\) is a unit and \(p\mid m\). Hence \(R=c\). Take \(s_1=1\) to obtain one exponent \(m_1\), then \(s_1=m_1-1\) to obtain \(m_2\) divisible by \(m_1-1\). Now \(c^{m_2}=c^{m_2-1}=1\), so \(c=1\). This determines the scalar for the chosen section.

For the given \(t\), write \(F=U_t\theta/\theta\). The function \(F(sz)\) is elliptic because \(st\) is a period. Its poles avoid the origin disk since \(s\) is a unit there and the denominator vanishes at the nonzero half-period. Apply elliptic translation to \(F(sz)\) at \(y/s\). The coordinate at \(y/s\) has \(s\)-th power \(\zeta_y\). Substitution by \([s]\) is invertible with inverse \([1/s]\), so cancellation of this substitution gives \(\widehat{F(\,\cdot+y)}=S_{\zeta_y}\widehat F\). Multiply by the identity for \(W_y\theta\), using \(W_y(F\theta)(z)=F(z+y)W_y\theta(z)\). This proves the claim for \(U_t\theta\), including when its value at the origin vanishes.
\end{proof}

\subsection*{Proof of \ref{B4.22g}}\label{prf:B4.22g}\phantomsection
\begin{proof}
The host hypotheses give \(d\) prime, \(d\equiv7\bmod16\), and \(p\) odd split with \(p\nmid N\mathfrak c\). A basis as stated exists. Choose bases locally at both primes, lift them simultaneously modulo those primes in \(M\), and apply Nakayama's lemma. Independence follows by extension to either completion. Multiplication by a common prime-to-\(p\) denominator puts the dual basis in the dual module.

We specify the normalization and the norm bound. Write \(h=N\mathfrak c\). The isogeny from the lattice \(h\mathcal O_{K_d}\) to \(\mathfrak c\) has degree \(h\) prime to \(p\); it identifies integral invariant differentials up to a unit. The discriminant has weight twelve, so \(|\jmath(\delta/\delta_0)|_p=1\). For the principal chart of \ref{B4.22f}, \(\Omega_0=\omega\jmath(\delta_0)\) is a unit by the same weight-twelve formula on a good model. The ideal-curve period is
\(\Omega=h\Omega_0\jmath(\delta/\delta_0)\). Each differential generator is a unit at the chosen place. Indeed the split coordinate lattice is integral and contains a vector with unit coordinate at the specified type, by taking the corresponding coordinate of a split basis and its trace dual. Thus \(|u_\tau|_p=1\).
\leanref{GlobalTheory.\allowbreak{}MomN.\allowbreak{}PeriodData.\allowbreak{}laws\_\allowbreak{}and\_\allowbreak{}chart}{Development.\allowbreak{}AbelianVarietyTheory.\allowbreak{}Proofs.\allowbreak{}GlobalInterpolation.\allowbreak{}LawsAndChart}{12--36}
\leanref{GlobalTheory.\allowbreak{}MomDelta.\allowbreak{}delta\_\allowbreak{}ratio}{Development.\allowbreak{}AbelianVarietyTheory.\allowbreak{}Proofs.\allowbreak{}GlobalInterpolation.\allowbreak{}DiscOneScaleRatioIdealRingAlgebraicUnit}{101--143}
\leanref{GlobalTheory.\allowbreak{}EigenGenerator.\allowbreak{}norm\_\allowbreak{}generator\_\allowbreak{}eq\_\allowbreak{}one}{Development.\allowbreak{}AbelianVarietyTheory.\allowbreak{}Proofs.\allowbreak{}SatakeParameters.\allowbreak{}NormPadicEmbeddingLeOneOfMemInverseHalfDifferent}{202--222}

For the elliptic input to \ref{B4.22c}, pull \(\theta_\Lambda\) back by \(z\mapsto\delta h z\). The pulled-back Hermitian form is \(h\) times the principal form. Division by \(\theta(z)^h\) gives an elliptic function of value one at zero. Its zeros and poles have order dividing \(2h\) as points of the elliptic curve and avoid the origin disk. The denominator calculation of \ref{B4.22f}, applied also to its reciprocal, gives norm one. The ordinary expansion of the unshifted ideal-curve theta is therefore integral. Directions in \(\delta\mathfrak c\) substitute \(c/h\), \(c\in\mathfrak c\), integral at the chosen split prime.

We prove that prime-to-\(p\) translation preserves the Gauss norm on any finite linear combination of translates of \(\theta^h\). Let \(E/\mathcal O_K\) be the good Weierstrass model of the principal curve \(\C/\mathcal O_{K_d}\) over a finite extension \(K/\Q_p\) containing the relevant prime-to-\(p\) torsion points and coefficients, with smooth irreducible special fibre \(E_k\) and reduced origin \(0_k\in E_k\). Let \(v_\eta\) be the discrete valuation on \(K(E)\) at the generic point \(\eta\) of \(E_k\), normalized to extend \(v_p\) on \(K\), and put \(|f|_\eta=p^{-v_\eta(f)}\). Let \(f\in K(E)^\times\) be a rational function whose horizontal polar divisor avoids \(0_k\). Choose \(c\in K^\times\) with \(v_\eta(c^{-1}f)=0\) and set \(f_0=c^{-1}f\). Because \(f_0\) has neither vertical zero nor vertical pole along \(E_k\) and no horizontal pole through \(0_k\), it belongs to the local ring \(\mathcal O_{E,0_k}\) and reduces to a nonzero rational function \(\bar f_0\in k(E_k)\) regular at \(0_k\). In the completed local ring \(\widehat{\mathcal O}_{E,0_k}\simeq\mathcal O_K[[T]]\), and hence also in the normalized multiplicative parameter \(Z\) (which is an integral unit power series in \(T\)), the expansion \(\widehat f_0(Z)\in\mathcal O_K[[Z]]\) reduces modulo \(\mathfrak m_K\) to the formal expansion of \(\bar f_0\) at the smooth point \(0_k\), which is nonzero. Thus \(\|\widehat f_0\|_{\mathrm{Gauss}}=1\) and hence
\[
 \|\widehat f\|_{\mathrm{Gauss}}=|c|_p=|f|_\eta.
\]
Moreover, for every \(P\in E(K)\), translation \(\tau_P\) extends to an \(\mathcal O_K\)-automorphism of the abelian scheme \(E\), so it fixes \(\eta\) and satisfies \(|f\circ\tau_P|_\eta=|f|_\eta\). The same holds for negation \([-1]\).

First take \(t\) with \(p\nmid s\), \(st\in\mathcal O_{K_d}\), and \(t\notin\tfrac12+\mathcal O_{K_d}\), and define \(F_t(z)=(U_t\theta)(sz)/\theta(sz)\). For any period \(\lambda\in\mathcal O_{K_d}\), the Heisenberg commutation law gives \(W_{s\lambda}(U_t\theta)=e^{-2\pi iB(\lambda,st)}U_t(W_{s\lambda}\theta)=U_t\theta\) since \(st\in\mathcal O_{K_d}\), so numerator and denominator share the same period multiplier and \(F_t\) is elliptic. By the definition of \(U_t\), direct cancellation yields \(F_t(z)F_{-t}(z+t/s)=1\), and evenness of \(\theta\) gives \(F_{-t}(-z)=F_t(z)\). Evaluating \(|\cdot|_\eta\) on these two identities gives \(1=|F_t|_\eta|F_{-t}(\,\cdot+t/s)|_\eta=|F_t|_\eta|F_{-t}|_\eta=|F_t|_\eta^2\), whence \(|F_t|_\eta=1\).

Now let \(h\ge1\) and let \(G=\sum_j c_j U_{t_j}(\theta^h)\) be a nonzero finite linear combination of prime-to-\(p\) translates of \(\theta^h\) with coefficients in \(\jmath^{-1}(K)\) (the case \(G=0\) being immediate), of Hermitian form \(hH\). For a prime-to-\(p\) shift \(t\) with \(t\notin\tfrac12+\mathcal O_{K_d}\), choose a positive integer \(s\) prime to \(p\) such that \(st\in\mathcal O_{K_d}\) and \(st_j\in\mathcal O_{K_d}\) for every \(j\). Set
\[
 F(z)=\frac{G(sz)}{\theta(sz)^h},\qquad
 F^{(t)}(z)=\frac{(U_t G)(sz)}{\theta(sz)^h}.
\]
Both \(F\) and \(F^{(t)}\) are elliptic by the same period-multiplier check. Since \(G\) and \(U_t G\) are holomorphic, every pole of \(F\) and of \(F^{(t)}\) is a zero of \(\theta(sz)\), that is, a point with \(sz\in\tfrac12+\mathcal O_{K_d}\). Because \(p\) is odd and \(p\nmid s\), these are nonzero prime-to-\(p\) torsion points, and reduction modulo \(p\) is injective on prime-to-\(p\) torsion \cite[Ch.~VII, Proposition~3.1]{silverman-2009}; hence no pole of \(F\) or \(F^{(t)}\) reduces to \(0_k\). The Gauss-norm identification therefore applies to both \(F\) and \(F^{(t)}\). Furthermore, expanding \(U_t G\) and \(F_t\) gives the identity in \(K(E)\)
\[
 F^{(t)}(z)=F(z+t/s)\,F_t(z)^h.
\]
Taking \(|\cdot|_\eta\) and using translation invariance and \(|F_t|_\eta=1\) gives \(|F^{(t)}|_\eta=|F|_\eta\). Since \(\widehat\theta(Z)\) is an integral unit series by \ref{B4.22f} and \([s](Z)=(1+Z)^s-1\) is an automorphism of \(\mathcal O_K[[Z]]\) for \(s\in\Z_p^\times\), multiplying by \(\widehat\theta([s](Z))^h\) and inverting \([s]\) yields
\[
 \|\widehat{U_t G}\|_{\mathrm{Gauss}}
 =\|\widehat{F^{(t)}}\|_{\mathrm{Gauss}}
 =|F^{(t)}|_\eta
 =|F|_\eta
 =\|\widehat F\|_{\mathrm{Gauss}}
 =\|\widehat G\|_{\mathrm{Gauss}}.
\]
If instead \(t\in\tfrac12+\mathcal O_{K_d}\), write \(t=t_1+t_1\) with \(t_1=t/2\). Because \(p\) is odd, \(t_1\) has prime-to-\(p\) order and \(t_1\notin\tfrac12+\mathcal O_{K_d}\). Since \(B(t_1,t_1)=0\), we have \(U_t G=U_{t_1}(U_{t_1}G)\), and applying the preceding equality twice gives \(\|\widehat{U_t G}\|_{\mathrm{Gauss}}=\|\widehat G\|_{\mathrm{Gauss}}\).

For the ideal curve use the finite theta space on the sublattice \(\delta h\mathcal O_{K_d}\). Its sections in this frame are spanned by translates of \(\theta^h\) of prime-to-\(p\) order. To see spanning, diagonalize one of the two cyclic translation groups of order \(h\); translations in the other cyclic group permute its characters transitively by the perfect commutator. Each character line has dimension one by descent to a principal lattice. The invariant span of the translates of the nonzero section is therefore the entire space. Applying the preceding norm equality in that frame proves the prime-to-\(p\) bound for the ideal curve. Thus \ref{B4.22c} applies to \(U_t e\), also for the companion lattices.
\leanref{GlobalTheory.\allowbreak{}MomOne.\allowbreak{}TranslateNorm.\allowbreak{}bd\_\allowbreak{}translate}{Development.\allowbreak{}AbelianVarietyTheory.\allowbreak{}Proofs.\allowbreak{}JetHarmonics.\allowbreak{}CpSeriesBoundedOneTorsionTranslateCmTheta}{72--154}
\leanref{GlobalTheory.\allowbreak{}MomN.\allowbreak{}latInt\_\allowbreak{}th0\_\allowbreak{}of\_\allowbreak{}split\_\allowbreak{}prime}{Development.\allowbreak{}AbelianVarietyTheory.\allowbreak{}Proofs.\allowbreak{}GlobalInterpolation.\allowbreak{}EvenThetaFunctionCmLatticeIntegralSplitPrime}{23--143}
\leanref{GlobalTheory.\allowbreak{}MomN.\allowbreak{}latInt\_\allowbreak{}opU\_\allowbreak{}of\_\allowbreak{}coprime}{Development.\allowbreak{}AbelianVarietyTheory.\allowbreak{}Proofs.\allowbreak{}GlobalInterpolation.\allowbreak{}CmThetaTranslatesPrimeToQIntegral}{202--237}

Choose an orthogonal frame in the rational period module whose squared lengths for the \(K_d\)-valued Hermitian form are positive integers \(a_j\) prime to \(p\). Locally a basis and its dual on the two split factors give a vector of unit length. Lift such a vector, and project to its orthogonal complement by dividing by its unit length. This preserves the localized module. Repetition and clearing prime-to-\(p\) denominators give the frame. In this frame the section \(e\) is a finite linear combination \(\sum_k c_k\prod_j U_{t_{kj}}(\theta^{a_j})\) of products of prime-to-\(p\) translates of \(\theta^{a_j}\), by applying the spanning argument just given in each factor. The coefficients \(c_k\) may be taken algebraic for the given algebraic jets: a basis chosen among the translates has an invertible minor of Taylor coefficients in algebraic scaled directions, and solving with that minor produces algebraic \(c_k\). Applying \(U_t\) shifts the factors, and the composition phases are constants independent of \(y\).

Each factor \(U_{t_{kj}}(\theta^{a_j})\) has integral transform by the preceding norm equality, so the finite algebraic linear combination has bounded binomial moments in the orthogonal frame, with bound \(\max_k|\jmath(c_k)|_p\) (no integrality of the coefficients \(c_k\) is needed at this step). By \ref{B4.22d} this determines a bounded \(\C_p\)-valued measure \(\mu\) on \(\Z_p^r\) of some finite norm \(N\). The change from the orthogonal frame to \((\beta_j)\) is invertible over the chosen \(\Z_p\), since both are bases locally and their trace-pairing denominators are units there. Push \(\mu\) forward by the dual coordinate change to obtain a bounded measure \(\mu_t\). Its binomial moments are those of \(U_t e\) in the directions \((\beta_j)\), which have \(p\)-adic norm at most one by Lemma~\ref{B4.22c}; uniqueness in Lemma~\ref{B4.22d} therefore gives the sharp norm bound \(\|\mu_t\|\le1\).

We check that this change of frame also transports the Weyl operator. At torsion level \(p^n\) write an allowed shift as \(y=N_0z/p^n\), with \(z\in M\) divisible by \(p^{2n}\) on the other half. For \(y\in N_0vM\) one can take \(n=1\), since \(pv\) is the square of the level factor that is a unit in the chosen component. Choose \(\epsilon\in N_0\mathcal O_{K_d}\) congruent to \(1,0\) modulo the two prime powers \(p^{2n}\). Reduce the chosen coordinates of \(N_0z\) in the basis \(\beta_j\) modulo \(p^n\), obtaining integers \(b_j\), and put \(y'=\epsilon\sum b_j\beta_j/p^n\). The difference \(l=y-y'\) lies in \(N_0M\) and is divisible by \(p^n\) on the other half locally. Indeed these assertions follow at the chosen half by the coordinate congruences, at the other half by divisibility by \(p^{2n}\) before division, and away from \(p\) by \(N_0\mid\epsilon\). The period \(l\) fixes \(U_t e\), since \(N_0t\) is a period. The Hermitian cocycle of \(l,y'\) is integral: on the two factors the possible \(p^n\) denominator of \(y'\) pairs with the \(p^n\) divisibility of \(l\), and its other component already contains \(p^n\). Thus \(W_yU_t e=W_{y'}U_t e\), and \([y,x]=[y',x]\) on periods.
\leanref{GlobalTheory.\allowbreak{}MomN.\allowbreak{}goodLift\_\allowbreak{}eq}{Development.\allowbreak{}AbelianVarietyTheory.\allowbreak{}Proofs.\allowbreak{}GlobalInterpolation.\allowbreak{}GoodLiftEqIdempotentLift}{63--333}

Multiply this lift, if necessary, by a common multiple of the prime-to-\(p\) frame lengths and translate denominators congruent to \(1\bmod p^n\). The same calculation leaves its operator and bracket unchanged. Each orthogonal coordinate now satisfies the hypotheses of \ref{B4.22f}, with an integer \(s\) clearing all the translates in that coordinate. The transform of each product is therefore shifted by its torsion coordinates \(\zeta_j\); the finite sum is shifted by the same substitution. On a transform,
\[
 \prod_j(\zeta_j(1+Z_j))^{c_j}
 =\prod_j\zeta_j^{\,c_j\bmod p^n}\prod_j(1+Z_j)^{c_j}.
\]
Here the identity first holds for nonnegative integers and extends by continuity. Orthogonality gives the first factor in the new coordinates as \(\jmath([y,x(c)])\): the Hermitian pairing is the sum of the coordinate pairings, and the dual change of variables preserves this sum. This identifies the translated measure as that character times \(\mu_t\).

Finally a directional derivative in the rational matrix direction \(Y\) reads the linear function
\(\Omega^{-1}\sum_jc_j\jmath(\operatorname{Tr}_{E/K_d}\sum_{kl}Y_{kl}\overline{(\beta_j)_{kl}})\), using the chosen embedding of \(K_d\). Extension in the direction variable to individual type and matrix coordinates gives the entries of \(x(c)^\dagger/\Omega\) in those embeddings. Multiplication by the differential generators and repetition for \(\nu\) give \(M_{x(c)}(\nu)\). Since \(W_t=\alpha(t)U_t\), the one remaining scalar is \(\jmath(\alpha(t))\), a root of unity independent of \(y,\nu\). This proves \eqref{eq:B4.22b-shared}.
\leanref{GlobalTheory.\allowbreak{}MomOne.\allowbreak{}Translates.\allowbreak{}moments\_\allowbreak{}of\_\allowbreak{}translates}{Development.\allowbreak{}AbelianVarietyTheory.\allowbreak{}Proofs.\allowbreak{}JetHarmonics.\allowbreak{}CmLatticeTranslationOperatorEqPhaseMulNegComp}{288--368}
\leanref{GlobalTheory.\allowbreak{}MomT.\allowbreak{}jetMoments\_\allowbreak{}of\_\allowbreak{}unitPeriod}{Development.\allowbreak{}AbelianVarietyTheory.\allowbreak{}Proofs.\allowbreak{}GlobalInterpolation.\allowbreak{}JetMomentsOfUnitPeriod}{19--143}
\end{proof}

\subsection*{Proof of \ref{B4.22h}}\label{prf:B4.22h}\phantomsection
\begin{proof}
All index computations can be made on the primary factors. At \(P\mid\varpi\), \(M'_P\subseteq M_P\) imposes one residue equation per source basis vector. Thus \(M/I\) injects into the product of \(f\) copies of \(\F_p^a\). At \(\bar P\) the inclusion is reversed, and away from \(p\) the lattices agree. The perfect alternating pairing between the relative quotients of the two self-dual period lattices gives \([M':I]=[M:I]\).

For \(b\in I\), the condition \(rb\in Z_0\) at \(P\mid\varpi\) is \(b\in\varpi M_P\), since \(c_Ag\bar\varpi\) is a unit there. At row primes it holds for every \(b\in I\), by the factor \(g\) and integrality of \(b\). Away from the selected level primes \(\varpi\) is a unit. Thus the preimage of \(Z_0\) in \(I\) is \(\varpi M\); the inclusion \(\varpi M\subseteq I\) also follows directly at the selected primes. Further, \(g\varpi M\subseteq M\cap B\). Taking the indices in
\[
 g\varpi M\subseteq\varpi M\subseteq I\subseteq M
\]
gives the asserted divisibility, since \([\varpi M:g\varpi M]\) is prime to \(p\).

At primes away from \(p\), every class of \(B/(B\cap M)\) has a representative in \(pB\). These representatives annihilate \(rI\), because \(prI\subseteq B\) and \(B\) is integral. At \(p\) the quotient is killed by \(p\); its classes are represented on the enlarged conjugate half. Pairing there with \(rI_P\) is pairing with \(I_P\), since \(r\) is a unit at \(P\). Its annihilator is the dual of \(I_P=M'_P\), hence \(M'_{\bar P}\) modulo \(M_{\bar P}\). The containment \(M'_{\bar P}\subseteq p^{-1}M_{\bar P}=B_{\bar P}\) follows by dualizing \(pM_P\subseteq M'_P\). This identifies the \(p\)-part with \(M'/I\). The Chinese remainder theorem realizes the local representatives by elements of \(M'\cap B\); at primes away from \(p\) prescribe the zero class. The two primary parts have coprime orders, so their sum is direct.
\end{proof}

\subsection*{Proof of \ref{B4.22b}}\label{prf:B4.22b}\phantomsection
\begin{proof}
Use representatives standard at the primes over \(\varpi\). Write \(R(x)\) for the requirement that at each such \(P\), \(xN_P\subseteq N'_P\) and \(l'_P x=l_P\) on \(N_P\). Set \(\Xi=\mathbf1_R\Xi_0\), \(H_S=\operatorname{Hom}_S(N,N')\), and \(M=M(N,N')\). For a positive integer \(h\) let \(B(h)\) have localizations \(hM_Q\) on \(S\), \(h^{-1}M_Q\) on \(\bar S\), and \(M_Q\) elsewhere. Let \(A(h)\) interchange the first two prescriptions. Write \(B_0=B(pg)\). The curve and half-different factors of \(M\) are units at the test primes and above \(p\), by the prime exclusions in the constructed family.

A class \(x\in (c_Apg)^{-1}H_S/(pgH_S)\) (in particular any test class \(x\in c_A^{-1}H_S/(pgH_S)\)) has a lift \(t_x\in A(c_Apg)\) with \(t_x-x\in pgH_S\), by the Chinese remainder theorem after clearing a common denominator. Its theta section is
\begin{equation}\label{eq:B4.22b-theta-class}
 \Theta_x=\frac1n W_{-t_x}\sum_{y\in B_0/(M\cap B_0)}W_y e_{N,N'},
 \qquad n=[M\cap A(pg):M\cap A(pg)\cap B_0].
\end{equation}
The finite Fourier transform of the indicator of zero on the quotient of order \(n\) is the constant \(1/n\). Translation of the indicator gives \(W_{-t_x}\). Two allowed lifts differ by an element of \(B_0\); the integral period cocycle on the half-lattices and reindexing the sum give the same section. The natural map into \(M/(M\cap B_0)\) gives \(n\mid[M:M\cap B_0]\). On every contributing class \(c_Ag t_x\in M\): at \(P\mid\varpi\) use integrality of \(x\) and its congruence with \(t_x\); at the other primes use the defining localizations of \(A(c_Apg)\).
\leanref{GlobalTheory.\allowbreak{}MomT.\allowbreak{}ThShape.\allowbreak{}thInd\_\allowbreak{}shape}{Development.\allowbreak{}AbelianVarietyTheory.\allowbreak{}Proofs.\allowbreak{}GlobalInterpolation.\allowbreak{}ThetaMovedCosetIndicatorVacuumShape}{284--396}

Write
\[
 c_\chi(N,N')=\chi(\det(L_0,N))^{b/2}\chi(\det(L'_0,N'))^{-a/2}
\]
for the Weil-splitting normalization of \ref{B4.22c}, and put \(\theta_{N,N'}=c_\chi(N,N')^{-1}e_{N,N'}\) for the unnormalized theta vector \(F_{N,N'}\) of \ref{B4.22c}. Dividing \eqref{eq:B4.22c-target} and \eqref{eq:B4.22c-source} by \(c_\chi\) cancels the relative character powers by multiplicativity of relative determinants and of \(\chi\), yielding
\[
 \theta_{N,N''}=I_{P(N,N'),P(N,N'')}\theta_{N,N'},\qquad
 \theta_{N_1,N'}=I_{P(N,N'),P(N_1,N')}\theta_{N,N'}.
\]
Set \(\widetilde\Theta_{N,N',x}=c_\chi(N,N')^{-1}\Theta_x\), obtained from \eqref{eq:B4.22b-theta-class} by replacing \(e_{N,N'}\) with \(\theta_{N,N'}\). Let \(J_\nu\in K\) be the rational jet sum determined by
\[
 \sigma_0(J_\nu)=\sum_{x}\sigma_0(\Xi(x))\,\partial_c^\nu(\widetilde\Theta_{N,N',x})(0)
 =c_\chi(N,N')^{-1}\sum_{x}\sigma_0(\Xi(x))\,\partial_c^\nu(\Theta_x)(0),
\]
where \(\sigma_0:K\hookrightarrow\C\) is the complex embedding and the index \(\nu\) lies in the rational-jet range stipulated in the statement. The point kernel \(k(d,d')\) equals the rational jet operator attached to \((J_\nu)_\nu\), scaled by the step-character prefactor
\[
 \operatorname{scal}_i(N,\mathrm{fr};N',\mathrm{fr}')
 =\mu_{\mathrm{inv},i}(N,\mathrm{fr})\,\mu_{\mathrm{inv},i}(N',\mathrm{fr}')^{-1},
\]
where
\[
 \mu_{\mathrm{inv},i}(N,\mathrm{fr})
 =\mu_{\mathrm{off},i}(N)\prod_{Q\in S_{\rm row}}m_{I,i}(Q,\det\mathrm{fr}_Q),
 \qquad
 \mu_{\mathrm{off},i}(N)
 =\prod_{\substack{v\notin S_{\rm row}|_{E^+}\\ v\text{ split in }E/E^+}}m_{U,i}(\operatorname{sel}(v))^{\operatorname{ord}_{\operatorname{sel}(v)}\det(L_0,N)}.
\]
Here \(E^+\) is the maximal totally real subfield of \(E\), \(S_{\rm row}|_{E^+}\) denotes the places of \(E^+\) below \(S_{\rm row}\) (above places \(v\) not split in \(E/E^+\), self-duality of \(L_0\) and \(N\) forces \(\operatorname{ord}_Q\det(L_0,N)=0\)), \(\operatorname{sel}(v)\) is the chosen prime of \(E\) over a split place \(v\), \(m_{U,i}(Q)\in K^\times\) is the unramified step-character value at \(Q\notin S_{\rm row}\) (of infinity type \(e_{i,\tau}\)), and \(m_{I,i}(Q,\cdot)\) is the local character at \(Q\in S_{\rm row}\). At standard representatives above \(p\), both \(c_\chi(N,N')\) and \(\operatorname{scal}_i(N,\mathrm{fr};N',\mathrm{fr}')\) are units at \(w\): the relative determinant ideals have order zero above \(p\), every unramified character value away from \(p\) is a \(w\)-unit, and the local frame factors at the row primes are \(w\)-units. We determine \((\iota_w J_\nu)\) on any finite set of rational indices.
\leanref{GlobalTheory.\allowbreak{}ColPeriod.\allowbreak{}kerOf\_\allowbreak{}thDistC}{Development.\allowbreak{}AbelianVarietyTheory.\allowbreak{}Proofs.\allowbreak{}SatakeParameters.\allowbreak{}ValuationLeOneOfMemZmu}{198--221}
\leanref{GlobalTheory.\allowbreak{}CharVal.\allowbreak{}charValAtT}{Development.\allowbreak{}AbelianVarietyTheory.\allowbreak{}Proofs.\allowbreak{}GlobalInterpolation.\allowbreak{}CharacterValueAtT}{865--913}

At an unseen place put \(v=\varpi/\bar\varpi\). Localizing at \(S,\bar S\) and their complement gives
\begin{equation}\label{eq:B4.22b-half-period}
 gvM\subseteq B_0,\qquad gv^{-1}B_0\subseteq M,\qquad gpB_0\subseteq M .
\end{equation}
For example at \(P\mid\varpi\), \(v=p/\bar\varpi^2\) with \(\bar\varpi\) a unit. The conjugate place reverses the factors, and \(p,\varpi,\bar\varpi\) are units at the row primes.

Represent the image of \(pB_0\) modulo \(M\) by \(R_2\). Each \(z-t_x\), \(z\in R_2\), has denominator dividing \(c_Ag\) modulo \(M\). Choose \(N_1\) prime to \(p\) divisible by \(c_Ag\) and all the integers \(N_0\) of \ref{B4.22g} for these translates. Represent the image of \(N_1vM\) modulo \(M\) by \(R_1\). In \(B_0/(B_0\cap M)\) these two images form a direct sum. The first is killed by \(g\), the second by \(p\) because \(pv=\varpi^2\), and \(N_1gB_0\subseteq N_1vM\) by \eqref{eq:B4.22b-half-period}. Bézout for \(N_1g,p\) proves that they generate. Multiplication by \(N_1\) preserves the image of \(vM\) modulo \(M\), and \(vM\) is self-dual since \(v\bar v=1\). The perfect alternating pairing on the two relative quotients gives
\[
 |R_1|=[vM:vM\cap M]=[M:vM\cap M].
\]

For fixed \(z\) apply \ref{B4.22g} with \(t=z-t_x\) to every \(y\in R_1\). In \eqref{eq:B4.22b-theta-class} commute the period operators to this order. The operators at \(z,y\) commute on \(B_0\); moving the displacement \(-t_x\) past \(y\) contributes \([y,-t_x]\). Thus at a matrix table point \(b\in M\) the \(y\)-dependent phase is \([y,b-t_x]\). It is a character of the indicated finite image: \(N_1(b-t_x)\in M\); a representative in \(N_1vM\cap M\) belongs to \(N_1M\), since \(pvM\subseteq M\) and \(N_1\) is prime to \(p\). For a character nontrivial at a class \(y_0\), translation by \(y_0\) multiplies its sum by a scalar different from one, so its sum is zero. The trivial character contributes \(|R_1|\). This is a single orthogonality test at each table point, common to all the jet coordinates. In the latter case \(N_1(b-t_x)\) pairs integrally with \(vM\), hence lies in \(vM\). Locally at \(P\mid\varpi\) it follows that \(b-t_x\in pM_P\), so \(R(b)\) holds. The denominator is absorbed because
\[
 n\mid[M:M\cap B_0]\mid
 [M:M\cap vM]\,[M\cap vM:M\cap vM\cap B_0].
\]
The second index is prime to \(p\), since its quotient is killed by \(g\). Hence \(|R_1|/n\) is integral at \(p\). More explicitly, in using \ref{B4.22d} take one approximation for all pairs \((y,\nu)\) for this fixed \(z\). Orthogonality then applies to each approximate table point. The original tolerance may be made arbitrarily small compared with \(|1/n|_p\). Summing the resulting finite approximations over \(z\) and the test classes proves the unseen row of the table below.
\leanref{GlobalTheory.\allowbreak{}MomT.\allowbreak{}CosetBar.\allowbreak{}clo\_\allowbreak{}class}{Development.\allowbreak{}AbelianVarietyTheory.\allowbreak{}Proofs.\allowbreak{}GlobalInterpolation.\allowbreak{}RelIndexNeZeroOfSmul}{207--248}
\leanref{GlobalTheory.\allowbreak{}MomT.\allowbreak{}CosetBar.\allowbreak{}moments\_\allowbreak{}thInd}{Development.\allowbreak{}AbelianVarietyTheory.\allowbreak{}Proofs.\allowbreak{}GlobalInterpolation.\allowbreak{}RelIndexNeZeroOfSmul}{251--458}

At a seen place put \(D'=L_{d'}\), \(N''=D'+\varpi(D')^\vee\), \(M''=M(N,N'')\), and \(I=M\cap M''\). At \(P\mid\varpi\), \(N''_P=\ker l'_P\); at the conjugate it is the dual of that kernel, and elsewhere it is \(N'\). These descriptions give self-duality and the inclusions \(\varpi N'\subseteq N''\), \(\bar\varpi N''\subseteq N'\). Write \(F=fa\). For standard local source basis vectors \(e_j\) consider
\[
 M\longrightarrow\prod_{P\mid\varpi} k(P)^a,\qquad
 b\longmapsto\bigl(l'_P(be_j)\bigr)_{P,j}.
\]
Its kernel is contained in \(I\). The assertion at \(P\mid\varpi\) follows by testing a basis; away from those primes \(\varpi b\in M''\) and \(\varpi\) is invertible. The target has order \(p^F\) by \ref{B8.54}. Taking the index of the kernel and then that of \(I\) proves
\[
 h_1=[M:I]\mid p^F .
\]
The alternating pairing gives a perfect duality between \(M/I\) and \(M''/I\), since the annihilators of \(M\) and \(M''\) are themselves. Thus \([M'':I]=h_1=p^\ell\), \(0\le\ell\le F\).
\leanref{GlobalTheory.\allowbreak{}MomT.\allowbreak{}MomentCompanion.\allowbreak{}companion\_\allowbreak{}gen}{Development.\allowbreak{}AbelianVarietyTheory.\allowbreak{}Proofs.\allowbreak{}GlobalInterpolation.\allowbreak{}MomentCompanion}{7--58}
\leanref{GlobalTheory.\allowbreak{}MomT.\allowbreak{}IdxBound.\allowbreak{}relIndex\_\allowbreak{}per\_\allowbreak{}dvd}{Development.\allowbreak{}AbelianVarietyTheory.\allowbreak{}Proofs.\allowbreak{}GlobalInterpolation.\allowbreak{}RelIndexPerDvd}{91--118}

Put \(r_0=c_Ag\bar\varpi\), \(Z_0=pgH_S\), and let \(J\) represent \(r_0I/(r_0I\cap Z_0)\). Addition of its elements preserves the test, by the row periodicity and \(I_P N_P\subseteq\ker l'_P\). For \(b\in M\), \(r_0b\in Z_0\) if and only if \(b\in\varpi M\). At \(P\mid\varpi\) cancel the units \(c_A,g,\bar\varpi\); at a row prime \(b\) is already integral and \(p\) is a unit. Division by \(\varpi\) preserves the other lattice conditions. Since \(\varpi M\subseteq I\), this gives \(|J|=[I:\varpi M]\). Also \(g\varpi M\subseteq M\cap B_0\) by the local definitions. Lemma~\ref{B4.22h}, applied with the row primes and the marked primes above \(\varpi\), identifies the refinement quotient and gives
\begin{equation}\label{eq:B4.22b-refinement}
 n\mid[M:g\varpi M]=h_1|J|[\varpi M:g\varpi M],
 \qquad v_p(n)\le\ell+v_p(|J|).
\end{equation}

Group the indicators in translates of \(J\) and use the lifts \(t_x+j\). Each \(W_{-j}\) fixes \(e_{N,N'}\). The sum of the characters produced by commuting it through \(W_y\) is \(|J|\) on \(B'=\{y\in B_0:B(r_0 I,y)\subseteq\Z\}\) and zero elsewhere. By the annihilator assertion of Lemma~\ref{B4.22h} with \(M'=M''\) and \(B=B_0\), the surviving classes modulo \(M\) consist of \(h_1\) classes of \(p\)-power order and the image of \(pB_0\) of order prime to \(p\). Thus the sum over the first factor is a complete companion-lattice average for each fixed element of the second. The first subgroup pairs integrally with \(r_0I\) by self-duality of \(M''\). The second does so because \(pr_0I\subseteq B_0\) and \(B_0\) is integral. Conversely, for \(y\in B'\), use \(r_0\bar\varpi M''\subseteq r_0 I\) and move scalars across the pairing to get \(\varpi^2c_Ag y\in M''\). Also \(g(\bar\varpi/\varpi)y\in M\) by \eqref{eq:B4.22b-half-period}, giving \(\bar\varpi c_Ag y\in M''\). Since \(\varpi^2,\bar\varpi\) are coprime, \(c_Ag y\in M''\). Bézout for \(c_Ag,p\) gives the required decomposition. The images intersect trivially: the first is killed by \(p\) since \(\bar\varpi M''\subseteq M\), the second by \(g\). Finally \(g\varpi M''\subseteq B_0\) and \(\bar\varpi M''\subseteq M\); Bézout for \(g\varpi,\bar\varpi\) shows that \(M''\cap B_0\) represents all of \(M''/I\).

By \eqref{eq:B4.22c-target} the sum over \(M''/I\) contributes \(\sqrt{h_1}\,\kappa e_{N,N''}\), where
\[
 \kappa=\chi(\det(N',N''))^{a/2}
       =\left(\prod_\tau\tau(\bar\varpi/\varpi)\right)^{a/2}.
\]
The local determinant ideal has order one at \(P\mid\varpi\), minus one at its conjugate, and zero elsewhere. At a seen embedding \(|\jmath(\kappa)|_p^2=p^F\).

For \(z\in pB_0\) the remaining translate \(t_1=z-t_x\) satisfies \(m_*:=c_Ag\varpi t_1\in M''\). Write \(A_1c_Ag\bar\varpi+B_1\varpi=1\) in \(\mathcal O_E\) and put \(t'=B_1\varpi t_1\). Then \(c_Ag\,t'\in M''\). Use \ref{B4.22g} for this \(t'\) and the companion pair, and write \(A_2N_0+B_2p=1\). Direct substitution gives
\[
 t_1=t'+N_0(\bar\varpi/\varpi)A_2A_1m_*+
                 B_2\bar\varpi^2 A_1m_* .
\]
The last term is a period and the middle term is an allowed displacement for the companion measure. The Heisenberg composition phases are roots of unity. The outer coefficient is therefore \(b_z=(|J|/n)\sqrt{h_1}\,\kappa\) times a root of unity. By \eqref{eq:B4.22b-refinement},
\[
 |\jmath(p^F b_z)|_p^2
 =p^{-2F}\left|\frac{|J|h_1}{n}\right|_p^2
   |h_1|_p^{-1}|\jmath(\kappa)|_p^2
 \le p^{-2F+\ell+F}\le1 .
\]
A period \(b\in M''\) sends \(N_{\bar P}\) into \(N''_{\bar P}\); taking adjoints gives \(b^\dagger(\ker l'_P)\subseteq N_P\). Applying \ref{B4.22d} and summing the finite class contributions with their integral test coefficients proves that, on every finite set of rational indices, \((\iota_wJ_\nu)_\nu\) lies in the closure of the integral span of
\begin{equation}\label{eq:B4.22b-tables}
 \begin{array}{c|c|l}
 &\text{table}&\text{support}\\ \hline
 \text{seen}&(p^{-fa}M_b(\nu))_\nu&
       b^\dagger(\ker l'_P)\subseteq N_P\quad(P\mid\varpi)\\
 \text{unseen}&(M_b(\nu))_\nu&R(b).
 \end{array}
\end{equation}
\leanref{GlobalTheory.\allowbreak{}MomT.\allowbreak{}UnitTr.\allowbreak{}moments\_\allowbreak{}unit\_\allowbreak{}translate}{Development.\allowbreak{}AbelianVarietyTheory.\allowbreak{}Proofs.\allowbreak{}GlobalInterpolation.\allowbreak{}MomentsTranslateRationalPointUnitDenominator}{71--160}
\leanref{GlobalTheory.\allowbreak{}MomT.\allowbreak{}CosetSum.\allowbreak{}moments\_\allowbreak{}coset\_\allowbreak{}sum}{Development.\allowbreak{}AbelianVarietyTheory.\allowbreak{}Proofs.\allowbreak{}GlobalInterpolation.\allowbreak{}MomentsCosetSum}{363--492}
\leanref{GlobalTheory.\allowbreak{}MomT.\allowbreak{}type\_\allowbreak{}p}{Development.\allowbreak{}AbelianVarietyTheory.\allowbreak{}Proofs.\allowbreak{}GlobalInterpolation.\allowbreak{}MemSuppBarOfNe}{267--326}
\leanref{GlobalTheory.\allowbreak{}MomT.\allowbreak{}type\_\allowbreak{}bar}{Development.\allowbreak{}AbelianVarietyTheory.\allowbreak{}Proofs.\allowbreak{}GlobalInterpolation.\allowbreak{}MemSuppBarOfNe}{248--266}

We pass from tables to the polynomial value lattice and record the determinant-twist compatibility across the step. Write \(\lambda=\tilde\lambda_i\) and \(\lambda'=\tilde\lambda_{i+1}\) for the source and target point weights in zero-based coordinates, where \(\tilde\lambda_{j,\tau,r}=\lambda^u_{j,\tau,r}-js\) (matching \ref{B4.21a} at \(j=0\); in the one-based indexing of \ref{B4.22a} the coordinate index is shifted by one). Put \(\sigma_\tau=e_{i,\tau}-a/2\) and \(\sigma'_\tau=e_{i,\tau}-b/2\), where \(e_{i,\tau}\in\Z\) is the common chain parameter of the adjacent point weights. Both \(\sigma_\tau\) and \(\sigma'_\tau\) are integers because \(a,b\) are even. The point weights satisfy
\[
 \lambda'_{\tau,r}+\sigma_\tau=\lambda_{\tau,r}+\sigma'_\tau\quad(0\le r<a),\qquad
 \lambda'_{\tau,r}=-\sigma_\tau\quad(a\le r<b).
\]
Hence the target determinant shift (the last coordinate of \(\lambda'_\tau\)) is \(h'_\tau=\lambda'_{\tau,b-1}=-\sigma_\tau\), agreeing with the chain determinant shift \(s_{i,\tau}=-h'_\tau=\sigma_\tau\) of \ref{B4.22a} and with the zero-based last-coordinate convention \(b_v=-\lambda^u_{0,\tau_w(v),k-1}\) of \ref{B4.21a}. Furthermore \(\sigma_\tau-\sigma'_\tau=(b-a)/2=s\), and writing \(h_\tau=\lambda_{\tau,a-1}\) for the source determinant shift, dominance of \(\lambda'_\tau\) gives
\[
 n_\tau:=h_\tau+\sigma'_\tau=\lambda'_{\tau,a-1}+\sigma_\tau\ge\lambda'_{\tau,a}+\sigma_\tau=0,
\]
which equals the last coordinate of the step partition \(\nu^u_{i+1,\tau}\).
\leanref{GlobalTheory.\allowbreak{}ColPeriod.\allowbreak{}chainWeights}{Development.\allowbreak{}AbelianVarietyTheory.\allowbreak{}Proofs.\allowbreak{}SatakeParameters.\allowbreak{}SubHalfRankInt}{17--70}
\leanref{GlobalTheory.\allowbreak{}JetPoint.\allowbreak{}detShift\_\allowbreak{}eq}{Development.\allowbreak{}AbelianVarietyTheory.\allowbreak{}Proofs.\allowbreak{}GlobalInterpolation.\allowbreak{}ExistsMulOfApprox}{115--119}
\leanref{GlobalTheory.\allowbreak{}BarStd.\allowbreak{}sub\_\allowbreak{}eq\_\allowbreak{}half}{Development.\allowbreak{}AbelianVarietyTheory.\allowbreak{}Proofs.\allowbreak{}GlobalInterpolation.\allowbreak{}AdmissibleAdjointRationalPointPadicMatrix}{369--385}

On the polynomial part \(P\) of a source value first multiply by \(\prod_\tau\det(X_\tau)^{n_\tau}\). Substitute the adjoint matrix in the first \(a\) columns in target coordinates. Equivalently, make this substitution with matrix indeterminates \(Y\) and write the result as \(\sum_\nu Q_\nu(X)Y^\nu\). A table at \(b\) evaluates \(Y_\tau\) at \(u_\tau\iota_w(\tau(b^\dagger))\). All these matrices are integral in the indicated frames: at the seen half this is the first support condition, and at the unseen half \(b\) carries the source marking kernel into the target marking kernel, so its adjoint is integral on their duals.

Each \(Q_\nu\) has the target polynomial weight \((\lambda'_{\tau,r}-h'_\tau)_{0\le r<b}\): multiplying by \(\det(X_\tau)^{n_\tau}\) adds \(n_\tau\) full columns to the source diagram, and padding its minors to rank \(b\) produces the target diagram with its determinant shift \(h'_\tau\) removed, since \((\lambda_{\tau,r}-h_\tau)+n_\tau=\lambda'_{\tau,r}+\sigma_\tau=\lambda'_{\tau,r}-h'_\tau\) for \(0\le r<a\) and \(\lambda'_{\tau,r}-h'_\tau=0\) for \(a\le r<b\). Matrix substitution preserves the span of these minors by Cauchy--Binet. A linear functional annihilating the target span thus annihilates \(\sum Q_\nu Y^\nu\) at every rational matrix \(Y\), and consequently annihilates each coefficient because \(K\) is infinite. A source Weyl vector belongs to the dual Weyl lattice, so its polynomial part evaluated in a local source basis \(B_\tau\) equals \(\prod_\tau\det(B_\tau)^{-h_\tau}\) times a polynomial with integral coefficients (that is, has integral coefficients after scaling by \(\prod_\tau\det(B_\tau)^{h_\tau}\)). Multiplying by \(\prod_\tau\det(X_\tau)^{n_\tau}\) therefore introduces the source determinant factor \(\prod_\tau\det(B_\tau)^{n_\tau-h_\tau}=\prod_\tau\det(B_\tau)^{\sigma'_\tau}\). Fix a target monomial in a local target basis \(B'_\tau\). Write \(d_\nu\) for its coefficient in the framed \(Q_\nu\), and let \(D\) be the finite set of indices occurring in the substitution. Equation~\eqref{eq:B4.22b-tables} gives, simultaneously on \(D\), sums \(J'_\nu=\kappa\sum_j c_j M_{b_j}(\nu)\) approaching the embedded jets, with \(|c_j|_p\le1\), with all \(b_j\) in the indicated support, and with \(\kappa=p^{-fa}\) or \(1\) on the respective side. After removing \(\prod_\tau\det(B_\tau)^{\sigma'_\tau}\) and \(\kappa\), the coefficient of each point evaluation is integral because the adjoint matrix maps the target frame into the source frame. Their integral combination has the same property. The error before division is bounded by \(\max_{\nu\in D}|d_\nu|_p\max_{\nu\in D}|\iota_wJ_\nu-J'_\nu|_p\). The divisor is fixed and nonzero, so closedness of the valuation ring gives the assertion for the actual coefficient. The output is rational, so its integral coefficients lie in \(\mathcal O_{K,(w)}\). Restoring the target determinant twist \(\prod_\tau\det(B'_\tau)^{h'_\tau}=\prod_\tau\det(B'_\tau)^{-\sigma_\tau}\) proves membership in the target dual Weyl lattice after removing \(\kappa\) times the total frame factor \(\prod_\tau\det(B_\tau)^{\sigma'_\tau}\det(B'_\tau)^{-\sigma_\tau}\).

For every positive root \(j<k\) of \(\mathrm{GL}_b\) the pairing with \(\lambda'_\tau+\rho\) is \(\lambda'_{\tau,j}-\lambda'_{\tau,k}+k-j<p\) by the size hypothesis. Over an algebraic closure of the residue field, the Jantzen sum formula for the Weyl module has no terms: for each positive root the indexing integers satisfy \(0<mp<\langle\lambda'_\tau+\rho,\alpha^\vee\rangle<p\) \cite[II, 8.19]{jantzen-2003}. Its radical is zero, so the Weyl and induced modules of this highest weight are simple. The canonical map between them is nonzero: the highest bideterminant takes value one at the identity, also after reduction. It is therefore onto, and surjectivity descends to the residue field by faithful scalar extension. The integral cokernel is finitely generated and has zero reduction, so it vanishes by Nakayama's lemma. Apply this at every type, tensor the resulting equalities, and restore the determinant twists in the chosen frames.
\leanref{GlobalTheory.\allowbreak{}JetPoint.\allowbreak{}closureIntegral}{Development.\allowbreak{}AbelianVarietyTheory.\allowbreak{}Proofs.\allowbreak{}GlobalInterpolation.\allowbreak{}ExistsMulOfApprox}{313--399}
\leanref{GlobalTheory.\allowbreak{}KerT.\allowbreak{}localization\_\allowbreak{}delAtr\_\allowbreak{}eq\_\allowbreak{}localization\_\allowbreak{}nabAtr}{Development.\allowbreak{}AbelianVarietyTheory.\allowbreak{}Proofs.\allowbreak{}GlobalInterpolation.\allowbreak{}LocPDelAtrEqLocPNabAtr}{24--36}

Evaluate the combined frame factor \(\prod_\tau\det(B_\tau)^{\sigma'_\tau}\det(B'_\tau)^{-\sigma_\tau}\) on each half. At the seen half the source basis is standard (\(B_\tau=I_a\)) and the target marking kernel has \(\det(B'_\tau)=p\) up to a \(w\)-unit. Since \(h'_\tau=-\sigma_\tau\), the frame factor together with the table scale \(p^{-fa}\) is \(p^{\sum_\tau h'_\tau-fa}\) up to a \(w\)-unit. Multiplying by the \(w\)-unit prefactor \(\operatorname{scal}_i(N,\mathrm{fr};N',\mathrm{fr}')\) proves the seen inclusion. At the unseen half both dual marking frames have determinant of valuation \(-e_w\), so the frame factor has valuation \(e_w\sum_\tau(\sigma_\tau-\sigma'_\tau)=e_wfs\). Multiplying by the \(w\)-unit prefactor \(\operatorname{scal}_i(N,\mathrm{fr};N',\mathrm{fr}')\) proves the unseen inclusion. For arbitrary representatives, choose rational unitary changes \(\gamma\in\U(E^a)\), \(\gamma'\in\U(E^b)\) moving \(N,N'\) to standard representatives above \(p\) as in \ref{B4.22e}. The total kernel scalar separates into source and target factors
\[
 \operatorname{scal}_i(N,\mathrm{fr};N',\mathrm{fr}')\,c_\chi(N,N')^{-1}
 =\rho_i(N,\mathrm{fr})\,\rho'_i(N',\mathrm{fr}'),
\]
where
\[
 \rho_i(N,\mathrm{fr})=\mu_{\mathrm{inv},i}(N,\mathrm{fr})\,\chi(\det(L_0,N))^{-b/2},\qquad
 \rho'_i(N',\mathrm{fr}')=\mu_{\mathrm{inv},i}(N',\mathrm{fr}')^{-1}\,\chi(\det(L'_0,N'))^{a/2}.
\]
By the infinity type \(e_{i,\tau}\) of the step character (with principal-ideal law matching \(\chi((z))=\prod_\tau\tau(z)^{-1}\) of \ref{B4.22c}, so that \(\mu_{\mathrm{inv},i}(\gamma N,\gamma\mathrm{fr})/\mu_{\mathrm{inv},i}(N,\mathrm{fr})=\prod_\tau\tau(\det\gamma)^{-e_{i,\tau}}\)), on a determinant change of order one at \(P\mid\varpi\) and minus one at \(\bar P\) the factor \(\rho_i\) changes by valuation \(\sigma'_\tau\) times that of \(\chi(P\bar P^{-1})\) and \(\rho'_i\) by \(-\sigma_\tau\) times that of \(\chi(P\bar P^{-1})\), where \(\tau\in\Sigma\) is the unique type corresponding to \(P\) (at a seen place \(w\)) or \(\bar P\) (at an unseen place \(w\)). Since \(v_w(\chi(P\bar P^{-1}))\) has opposite sign to \(v_w(\tau(\det\gamma))\), these changes cancel the determinant-twist exponents \(\sigma'_\tau\) and \(-\sigma_\tau=h'_\tau\) in the frame factor and in the \(\GL\)-action on \(V_\lambda,V_{\lambda'}\). Pullback carries the test to \(x\mapsto\Xi_0(\gamma'^{-1}x\gamma)\); \eqref{eq:B4.22c-equiv} transports the normalized theta section, and the integral frame actions transport the two value lattices equivariantly.

For the neighbor identity fix a rational map \(x\) and one \(P\mid\varpi\). A lower source neighbor \(d_0\uparrow_P d\) has underlying lattice \(N_0\) with \((N_0)_P=\ker l_P\), \((N_0)_{\bar P}=(\ker l_P)^\vee\), and \((N_0)_Q=N_Q\) for \(Q\notin\{P,\bar P\}\). Its marking, if it contributes, is forced to be \(z\mapsto l'_P(xz)\) on \(\ker l_P\). If a target upper neighbor \(d'\uparrow_P d'_1\) contributes, this functional takes value one at \(pz\) for \(l_P(z)=1\), by the defining marking of \(d'_1\). It therefore gives a contributing lower source neighbor \(d_0\uparrow_P d\). Conversely, present the move from a contributing lower source neighbor \(d_0\uparrow_P d\) by adjoining \(p^{-1}z_0\) to \((N_0)_P\), where \(z_0\in(N_0)_P\) has marking one on \(d_0\). Then \(xz_0\in N'_P\) has marking one on \(d'\). Adjoining \(p^{-1}xz_0\) to \(N'_P\) gives the underlying lattice \(N'_1\) of the target upper neighbor \(d'_1\) (with \((N'_1)_{\bar P}\) its Hermitian dual and \((N'_1)_Q=N'_Q\) for \(Q\notin\{P,\bar P\}\)). One can use a global lift by clearing a denominator outside \(P\) and inverting it modulo \(P\). Two possible lifts differ by \(pN'_P\), so \(d'_1\) is unique. Writing \(\phi_T(d,d';x)\) for the indicator of \(R(x)\), this proves
\[
 \sum_{d'\uparrow_P d'_1}\phi_T(d,d'_1;x)=\sum_{d_0\uparrow_P d}\phi_T(d_0,d';x)
\]
for every rational map \(x\). Every \(x\) with a nonzero term on either side satisfies \(x\in\operatorname{Hom}_P(N,N'_1)\cap\operatorname{Hom}_P(N_0,N')\subseteq p^{-1}\operatorname{Hom}_P(N,N')\). Off \(P,\bar P\) the decorations and frames are fixed.

We relate the unnormalized theta sections on both sides to the common base pair \((N,N')\). Put \(M=M(N,N')\), \(M_0=M(N_0,N')\), \(M_1=M(N,N'_1)\), and locally write \(B,B_0,B_1\) for their level-\(pg\) half-lattices \(B(pg)\) (and \(A_0,A_1\) for the corresponding \(A\)-lattices of \(M_0,M_1\)). For any lattice pair \((L,L')\) of ranks \((a,b)\), the local index \([B(L,L')_{Q'}:M(L,L')_{Q'}\cap B(L,L')_{Q'}]\) is \(1\) for \(Q'\notin\bar S\) and \([M(L,L')_{\bar Q}:pgM(L,L')_{\bar Q}]=[M(L,L')_Q:pgM(L,L')_Q]\) at \(Q'=\bar Q\) for \(Q\in S\); the product over \(Q\in S\) depends only on \(a,b,S,pg\) and equals \(n=[M\cap A(pg):M\cap A(pg)\cap B]\). Hence \([B(L,L'):M(L,L')\cap B(L,L')]=n\), and the unnormalized class section built from \(\theta_{L,L'}\) is
\[
 \widetilde\Theta_{L,L',x}=n^{-1/2}W_{-t_x}I_{M(L,L'),B(L,L')}\theta_{L,L'}.
\]
For \((N,N'_1)\), we have \(\theta_{N,N'_1}=I_{M,M_1}\theta_{N,N'}\). The lattices \(M,M_1,B_1,B\) agree away from \(S\cup\bar S\) and split along the complementary isotropic halves at \(S\) and \(\bar S\), so the isotropic-half composition law of \ref{B4.22c} gives \(I_{M_1,B_1}I_{M,M_1}=I_{B,B_1}I_{M,B}\) with composition scalar \(1\). Locally at \(P\), \(B_P=pM_P\subseteq pM_{1,P}=B_{1,P}\) has index \([M_{1,P}:M_P]=[(N'_1)_P:N'_P]^a=p^a\); at \(\bar P\), \(B_{1,\bar P}=p^{-1}M_{1,\bar P}\subseteq p^{-1}M_{\bar P}=B_{\bar P}\); elsewhere \(B\) and \(B_1\) agree. Thus \([B_1:B\cap B_1]=p^a\), and by the Chinese remainder theorem each coset of \(B_1/(B\cap B_1)\) has a representative \(u\in B_1\) with \(u_{\bar Q}\in c_ApgM_{1,\bar Q}\) for every \(Q\in S\). For a test class \(x\in c_A^{-1}\operatorname{Hom}_S(N,N'_1)\), choose a lift \(t_x\in A_1(c_Apg)\) of \(x\) modulo the finer lattice \(B_S=pg\operatorname{Hom}_S(N,N')\subseteq B_{1,S}\). Because \(t_x\in A_1(c_Apg)\) and \(u\in B_1\) satisfies \(u_{\bar Q}\in c_ApgM_{1,\bar Q}\) for all \(Q\in S\), we have \(\frac{\sqrt d}{2}H(t_x,u)\in\mathcal O_{K_d}\), whence \(W_{-t_x}W_u=W_{-(t_x-u)}\) as in \ref{B4.22f}. Moreover, \(x_P,u_P\in M_{1,P}\subseteq p^{-1}M_P=A(c_Apg)_P\) and \(M_{1,\bar P}\subseteq M_{\bar P}\) (and likewise \(M_{0,P}\subseteq p^{-1}M_P\) and \(M_{0,\bar P}\subseteq M_{\bar P}\) below), so \(t_x-u\) lies in \(A(c_Apg)\); since \(t_x-x\in B_S\), it lifts \(x-u\) modulo \(B_S\), which runs over the \(p^a\) cosets of \(B_S\) inside \(x+B_{1,S}\). Therefore
\[
 \widetilde\Theta_{N,N'_1,x}
 =p^{-a/2}\sum_{\substack{x'\equiv x\bmod B_{1,S}\\ x'\bmod B_S}}\widetilde\Theta_{N,N',x'}.
\]
By the same calculation for \((N_0,N')\), we have \(\theta_{N_0,N'}=I_{M,M_0}\theta_{N,N'}\), \(I_{M_0,B_0}I_{M,M_0}=I_{B,B_0}I_{M,B}\), and \([B_0:B\cap B_0]=[M_{0,P}:M_P]=[N_P:(N_0)_P]^b=p^b\), giving
\[
 \widetilde\Theta_{N_0,N',x}
 =p^{-b/2}\sum_{\substack{x'\equiv x\bmod B_{0,S}\\ x'\bmod B_S}}\widetilde\Theta_{N,N',x'}.
\]
The functions \(\phi_T(d,d'_1;\cdot)\Xi_0\) and \(\phi_T(d_0,d';\cdot)\Xi_0\) are invariant under \(B_{1,S}\) and \(B_{0,S}\) respectively: at \(P\), \(p\operatorname{Hom}_P(N,N'_1)\) sends \(N_P\) into \(p(N'_1)_P\subseteq\ker l'_{1,P}\) and \(p\operatorname{Hom}_P(N_0,N')\) sends \((N_0)_P\) into \(pN'_P\subseteq\ker l'_P\), while at \(S_{\rm row}\) the lattices are unchanged. Summing over the test quotients unfolds both sides onto the common quotient \(p^{-1}c_A^{-1}\operatorname{Hom}_S(N,N')/B_S\):
\begin{align*}
 \sum_{d'\uparrow_P d'_1}\sum_{x\bmod B_{1,S}}
 \phi_T(d,d'_1;x)\Xi_0(x)\widetilde\Theta_{N,N'_1,x}
 &=p^{-a/2}\sum_{x'\bmod B_S}
 \left(\sum_{d'\uparrow_P d'_1}\phi_T(d,d'_1;x')\right)
 \Xi_0(x')\widetilde\Theta_{N,N',x'},\\
 \sum_{d_0\uparrow_P d}\sum_{x\bmod B_{0,S}}
 \phi_T(d_0,d';x)\Xi_0(x)\widetilde\Theta_{N_0,N',x}
 &=p^{-b/2}\sum_{x'\bmod B_S}
 \left(\sum_{d_0\uparrow_P d}\phi_T(d_0,d';x')\right)
 \Xi_0(x')\widetilde\Theta_{N,N',x'}.
\end{align*}
By the indicator equality above, the left-hand sum equals \(p^{(b-a)/2}=p^s\) times the right-hand sum.

It remains to compare the step-character prefactors \(\operatorname{scal}_i\) on the two sides. Since \(P,\bar P\notin S_{\rm row}\) and all lattices and frames away from \(P,\bar P\) are fixed, only the place \(v_0=P\cap\mathcal O_{E^+}\) changes in \(\mu_{\mathrm{inv},i}\). Set \(\varepsilon=1\) if \(\operatorname{sel}(v_0)=P\) and \(\varepsilon=-1\) if \(\operatorname{sel}(v_0)=\bar P\), and define the split step-character value \(\mu_P=m_{U,i}(\operatorname{sel}(v_0))^\varepsilon\in K^\times\). Replacing \(N\) by \(N_0\) decreases the \(P\)-localization by index \(p\) and increases the \(\bar P\)-localization by index \(p\), so \(\operatorname{ord}_{\operatorname{sel}(v_0)}\det(L_0,N_0)=\operatorname{ord}_{\operatorname{sel}(v_0)}\det(L_0,N)+\varepsilon\) and
\[
 \mu_{\mathrm{inv},i}(N_0,\mathrm{fr})=\mu_P\,\mu_{\mathrm{inv},i}(N,\mathrm{fr}).
\]
Replacing \(N'\) by \(N'_1\) increases the \(P\)-localization by index \(p\) and decreases the \(\bar P\)-localization by index \(p\), so \(\operatorname{ord}_{\operatorname{sel}(v_0)}\det(L'_0,N'_1)=\operatorname{ord}_{\operatorname{sel}(v_0)}\det(L'_0,N')-\varepsilon\) and
\[
 \mu_{\mathrm{inv},i}(N'_1,\mathrm{fr}')^{-1}=\mu_P\,\mu_{\mathrm{inv},i}(N',\mathrm{fr}')^{-1}.
\]
Every summand on both sides therefore carries the same prefactor \(\mu_P\,\operatorname{scal}_i(N,\mathrm{fr};N',\mathrm{fr}')\), independent of the ranks \(a,b\). Applying the linear rational jet operator from \(V_\lambda\) to \(V_{\lambda'}\) and multiplying by this common scalar yields the neighbor identity.
\leanref{GlobalTheory.\allowbreak{}KerT.\allowbreak{}UK.\allowbreak{}sum\_\allowbreak{}phiT}{Development.\allowbreak{}AbelianVarietyTheory.\allowbreak{}Proofs.\allowbreak{}GlobalInterpolation.\allowbreak{}UkTermOfMoveStep}{374--391}
\leanref{GlobalTheory.\allowbreak{}KerT.\allowbreak{}UK.\allowbreak{}ukTerm\_\allowbreak{}of\_\allowbreak{}moveStep}{Development.\allowbreak{}AbelianVarietyTheory.\allowbreak{}Proofs.\allowbreak{}GlobalInterpolation.\allowbreak{}UkTermOfMoveStep}{491--534}
\end{proof}

\subsection*{Proof of \ref{B4.22a}}\label{prf:B4.22a}\phantomsection
\begin{proof}
Put \(p=T_u\), \(s=(t+1)/2\), and let \(V\) be the selected set of level indices from \ref{B4.21a}. Write \(r=|V|\), \(K_i\) for the \(i\)-th lift, \(C_0=1\), and \(C_{i+1}=K_iC_i\). Let \(S_i\) be the complex image of the localized form lattice and \(U_i\) the ordered product of the upward operators at \(V\) on stage \(i\).

The class stabilizers have orders prime to \(p\). Indeed a stabilizer acts faithfully on \(E^{N_i}\). An element of order \(p\) acting on this space as a rational vector space would have a primitive \(p\)-th root of unity as an eigenvalue. Its minimal polynomial would have degree at least \(p-1\), whereas the rational dimension is \(2fN_i\le2fN_m<p-1\). Cauchy's theorem gives the assertion. Thus the class weights \(1/|\operatorname{Stab}(c)|\) preserve the local bounds of \ref{B4.22b}.

At a seen place \ref{B4.22b} puts the class sum, multiplied by \(p^{fN_i-\sum_\tau h'_\tau}\), in the target value lattice with every marking kernel inserted. At an index in \(V\), the marked kernel of an upper neighbor is the ordinary lattice at the lower endpoint. Summing over the upper neighbors therefore erases that index. At a remaining index choose a basis in which the marking reads the first coordinate modulo \(p\); the kernel basis change is \(\operatorname{diag}(p,1,\dots,1)\). Multiplying by \(p^{\tilde\lambda_1}\) erases that index, where \(\tilde\lambda_1\) is the first target point weight. Indeed if \(h'\) is its last weight, a generator of the polynomial part has degree \(\beta_1\le\tilde\lambda_1-h'\) in the first row. Each of its \(\tilde\lambda_1-h'\) diagram columns uses that row at most once. This bound also holds on the Weyl generators obtained by integral row substitutions and coefficient extraction. The net power after the determinant twist is \(\tilde\lambda_1-h'-\beta_1\ge0\).
\leanref{GlobalTheory.\allowbreak{}LvT.\allowbreak{}pow\_\allowbreak{}smul\_\allowbreak{}mem\_\allowbreak{}valB\_\allowbreak{}erase}{Development.\allowbreak{}AbelianVarietyTheory.\allowbreak{}Proofs.\allowbreak{}GlobalInterpolation.\allowbreak{}LevelPrimePowSmulMemValueLatticeErase}{534--622}

Use one-based coordinates and write \(\tilde\lambda_{j,\tau,k}=\lambda^u_{j,\tau,k}-js\). At stage \(i+1\), \(-h'_\tau=s_{i,\tau}\) is its chain determinant shift. The total exponent just paid is
\[
 fN_i+\sum_\tau s_{i,\tau}
       +\sum_{v\notin V}\tilde\lambda_{i+1,\tau_w(v),1}.
\]
Each \(s_{i,\tau}\) is bounded above by \(W_{u,\tau}\) of \ref{B4.03} since it is minus the last point weight. For \(v\notin V\), \(\tau_w(v)\) is a non-top type, so the first point weight is also bounded above by \(W_{u,\tau_w(v)}\). The map \(v\mapsto\tau_w(v)\) is bijective. Hence we may use \(d_0=fN_m+2\sum_\tau W_{u,\tau}\) at every stage. At an unseen place \(V=\varnothing\); the factor \(p^{fs}\) supplied by \ref{B4.22b} is integral, and we take \(d_0=0\). The lift and the upward operators preserve the space of forms by their equivariant class-sum definitions. Testing at representatives now gives
\begin{equation}\label{eq:B4.22a-step}
 p^{d_0}U_{i+1}K_iS_i\subseteq S_{i+1}.
\end{equation}

For one level index the last identity of \ref{B4.22b} passes to the class sum with its stabilizer weights. An incident-pair orbit through \((d,e)\) has weight
\([\operatorname{Stab}(d):\operatorname{Stab}(d,e)]/|\operatorname{Stab}(d)|=1/|\operatorname{Stab}(d,e)|\);
counting from \(e\) gives the same weight. Thus \(U_{i+1,v}K_i=p^sK_iU_{i,v}\). Move the factors of the ordered product through in order to obtain
\[
 U_{i+1}K_i=p^{sr}K_iU_i,\qquad U_iC_i=p^{sri}C_iU_0 .
\]
For \(z\in S_0\) put \(b_0=0\) and \(b_{i+1}=b_i+d_0+sr(i+1)\). If \(p^{b_i}C_iU_0^iz\in S_i\), apply \eqref{eq:B4.22a-step} and the second identity with \(i+1\). This gives \(p^{b_{i+1}}C_{i+1}U_0^{i+1}z\in S_{i+1}\). Hence \(b_m=md_0+sr\,m(m+1)/2\). At a seen place \(r=n_T\), and \ref{B4.03} gives \(b_m=mT^*_{u,w}\). At an unseen place both sides are zero.

By \ref{B4.21a} the input is \(p^{a_0}U_0^m z\) after embedding, for a localized ramified vector with image \(z\in S_0\). Here \(a_0\) is minus \(m\) times the sum of the selected last initial weights; the weight signs give \(a_0\ge0\). The scalar \(cp^{a_0-b_m}\) lies in \(\mathcal O_{K,(w)}\) by the hypothesis on \(v_w(c)\). Multiplying the terminal inclusion by this scalar proves the reference assertion.

At a used twist with \(p_w\notin Q'(\mu)\), the level-zero identification preserves the localized base lattice. On corresponding decorated representatives the auxiliary decoration at stage zero is unique, so the frame action and class value are identified. Scaling markings commutes with this correspondence, which therefore carries the kernels defining the ramified subspace to one another. Upper neighbors correspond by the same residue vectors of marking one, so each upward sum, and then its ordered product and \(m\)-th power, intertwine. The twisted input thus has the same presentation with the same \(a_0\) and the corresponding twisted \(z\). Applying the iteration on the twisted chain puts \(p^{b_m}\) times its column in the localized terminal lattice. The scalar \(p^{-b_m}\) has valuation \(-e_wmT^*_{u,w}\), giving exactly the bound of \ref{B4.06}.
\leanref{GlobalTheory.\allowbreak{}LvT.\allowbreak{}ColSharp.\allowbreak{}chain\_\allowbreak{}into}{Development.\allowbreak{}AbelianVarietyTheory.\allowbreak{}Proofs.\allowbreak{}GlobalInterpolation.\allowbreak{}LevelChainSharpColumnValuationOfStability}{91--115}
\leanref{GlobalTheory.\allowbreak{}LvT.\allowbreak{}TrueStep0.\allowbreak{}step0Tie\_\allowbreak{}trueCols}{Development.\allowbreak{}AbelianVarietyTheory.\allowbreak{}Proofs.\allowbreak{}GlobalInterpolation.\allowbreak{}TwistedChainInitialLatticeCompatibleTrueColumns}{167--174}
\end{proof}

\subsection*{Proof of \ref{B4.23a}}\label{prf:B4.23a}\phantomsection
\begin{proof}
Fix \(w\mid p=g_i\ne T_u\), a word and a used tuple. At step
\(j\) use representatives standard above \(p\) with lattices
\(N,N'\). Let \(S\) be the summation prime set and \(S_i\)
its primes over \(\pi_i\). The constructed test \(\Xi\) has
\(|\Xi|_w\le1\), is supported in \(\operatorname{Hom}_{S_i}(N,N')\)
and is invariant under
\[
 p\operatorname{Hom}_{S_i}(N,N')\cap
 c_B\operatorname{Hom}_{S\setminus S_i}(N,N').
\]
Here \(c_B\) is its period modulus with \(p\mid c_B\), and
\(c_A\) its support denominator. Write \(c_Ac_B=p^n c'\)
with \(p\nmid c'\). Both \(c_B\) and \(pc'\) divide \(c_Ac_B\)
and their period groups preserve the test. At \(S_i\) both
are divisible by \(p\), and elsewhere in \(S\) the factor
\(p\) is a unit and both moduli are divisible locally by
\(c_B\). Refining either sum to modulus \(c_Ac_B\) partitions
the cosets with the same test value and preserves the normalized
theta sum. Rationality and injectivity of \(\iota\) give equality
in \(K\). At modulus \(pc'\), integral support on \(S_i\)
and removal of the unit \(p^n\) elsewhere give
\(x\in(c')^{-1}\operatorname{Hom}_S(N,N')\) for a nonzero
summand.

Write \(M\) for the split period lattice and put
\(d_0=fN_{j-1}N_j\). The prime-to-\(p\) part of the average
has representative set \(S_1\); its shifts \(t_k\) satisfy
\(m_0 t_k\in M\) for one positive integer \(m_0\) prime to \(p\), and its phase
coefficients \(r_k\) have norm at most one. For the \(p\)-part
let \(S_a,S_b\) represent \((\bar\pi_i/\pi_i)M/\bar\pi_iM\)
and \(M/\bar\pi_iM\), respectively. These dual quotients have
cardinality \(q=p^{d_0}\), with one residue coordinate per
entry per type. On \(S_b\) take \(\Psi=|S_1|^{-1}\mathbf1_0\).
For \(a\in S_a\) its Fourier coefficient is
\[
 \widehat\Psi(a)=q^{-1}\sum_{b\in S_b}\Psi(b)e^{-2\pi iB_H(b,a)}
 =(|S_1|q)^{-1}.
\]
Write \(G=\sum_{a\in S_a}\widehat\Psi(a)W_a e_{N,N'}\).
The test-class function is the sum of the \(r_k U_{t_k}G\).
In the following estimate \(\iota_p:\C\simeq\C_p\) extends
the split embedding at \(w\), \(M_{\mathrm{fr}}\) is the integral
matrix of direction coefficients in the period frame, and
\(A_{\mathrm{fr}}=\iota_p(\bar\pi_i)M_{\mathrm{fr}}\). Denote the
Taylor series of a function \(F\), read through \(\iota_p\)
and evaluated in directions with coefficient matrix \(B\),
by \(\mathcal S(F;B)\). For a finite tuple of variables \(Z\), define the ring of Hurwitz-integral power series
\[
 \mathcal O_{\C_p}\langle\langle Z\rangle\rangle
 =\left\{\sum_\alpha c_\alpha Z^\alpha\in\C_p[[Z]] :
 \alpha!\,c_\alpha\in\mathcal O_{\C_p}\text{ for all }\alpha\right\},
\]
where \(\alpha!=\prod_s\alpha_s!\). This subring contains \(\mathcal O_{\C_p}[[Z]]\) and is closed under \(\mathcal O_{\C_p}\)-linear combinations, products (since the \(Z^\gamma/\gamma!\)-coefficient of a product of \(\sum a_\alpha Z^\alpha/\alpha!\) and \(\sum b_\beta Z^\beta/\beta!\) is \(\sum_{\alpha+\beta=\gamma}\binom{\gamma}{\alpha}a_\alpha b_\beta\)), and \(\mathcal O_{\C_p}\)-linear substitutions of the variables (since for \(\ell_s\in\mathcal O_{\C_p}\) the multinomial expansion gives \((\sum_s\ell_s W_s)^n/n!=\sum_{|\beta|=n}(\prod_s\ell_s^{\beta_s})W^\beta/\beta!\), so each \(Z^\alpha/\alpha!\) maps into \(\mathcal O_{\C_p}\langle\langle W\rangle\rangle\) and each degree receives finitely many terms). The local estimate is
\begin{equation}\label{eq:B4.23a-split}
 p^{d_0}\mathcal S\!\left(\sum_k r_k U_{t_k}G;
 A_{\mathrm{fr}}\right)
 \in\mathcal O_{\C_p}\langle\langle Z\rangle\rangle .
\end{equation}
Here \(Z\) denotes the variables of the chosen directions.
\leanref{Formalism.\allowbreak{}L2.\allowbreak{}AlgTh.\allowbreak{}split\_\allowbreak{}bound\_\allowbreak{}of\_\allowbreak{}xi\_\allowbreak{}lemmas}{Formalism.\allowbreak{}LocalToGlobal.\allowbreak{}AlgebraicTheory.\allowbreak{}AlgebraicTheoryProofs.\allowbreak{}XiConn}{464--496}

We prove the étale translation estimate: if a shift \(t\) satisfies \(N_0\pi_0 t\in M\) and \(N_0 p t\in M\) for a positive integer \(N_0\) with \(p\nmid N_0\) and a CM integer \(\pi_0\in\mathcal O_{K_d}\) satisfying \(\pi_0\bar\pi_0=p\) and \(|\iota_p(\pi_0)|_p=1\), then \(\mathcal S(U_t e_{N,N'};B)\in\mathcal O_{\C_p}\langle\langle Z\rangle\rangle\) for every direction matrix \(B\) over \(\C_p\) with entries of \(p\)-adic norm at most \(p^{-1}\). Because \(p=g_i\) is a fixed auxiliary prime of the host and \(2fN_m+1<p\) (with \(f\ge1\) and \(N_m\ge2\)), we have \(p>5\), \(p\nmid d\), \(p\nmid\operatorname{discr}(E)\), and \(p\) splits unramified in \(K_d\).

First reduce from \(U_t e_{N,N'}\) to one-variable elliptic theta translates. Write \(L_0,L'_0\) for the standard lattices of ranks \(N_{j-1},N_j\) and \(M_{\mathrm{std}}\) for the standard period lattice of \((L_0,L'_0)\). Since \(N,N'\) agree with \(L_0,L'_0\) locally at \(p\), there is a positive integer \(M_1\) prime to \(p\) such that \(M_1 N\subseteq L_0\), \(M_1 L_0\subseteq N\), \(M_1 N'\subseteq L'_0\), and \(M_1 L'_0\subseteq N'\); hence \(M_1^2 M\subseteq M_{\mathrm{std}}\) and \(M_1^2 M_{\mathrm{std}}\subseteq M\). As in the proof of \ref{B4.22c}, averaging \(\Theta^{\otimes N_{j-1}N_j}\) from \(M_{\mathrm{std}}\) to \(M\) (with prime-to-\(p\) indices and \(p\)-adic unit character factors) and applying \(U_t\) expresses \(U_t e_{N,N'}\) as a finite linear combination of translates \(U_{t+s_\rho}\Theta^{\otimes N_{j-1}N_j}\) with coefficients of \(p\)-adic norm at most \(1\), where each \(s_\rho\) satisfies \(M_1^2 s_\rho\in M_{\mathrm{std}}\). Put \(M_2=M_1^2 N_0\), still prime to \(p\). Then \(M_2\pi_0(t+s_\rho)\in M_{\mathrm{std}}\) and \(M_2 p(t+s_\rho)\in M_{\mathrm{std}}\) for every \(\rho\), so each matrix entry \(t_{\rho,i_1 i_2}\) (\(1\le i_1\le N_{j-1}\), \(1\le i_2\le N_j\)) satisfies \(M_2\pi_0 t_{\rho,i_1 i_2}\in P_1\) and \(M_2 p t_{\rho,i_1 i_2}\in P_1\), where \(P_1\) is the one-entry period lattice. Applying \ref{B4.22e} to the fourth power of entry \((i_1,i_2)\) yields
\[
 (U_{t_{\rho,i_1 i_2}}\Theta)^{\otimes 4}(z)
 =\frac{e}{\sqrt J}\sum_{x\in Q/(P\cap Q)}\alpha(x)e^{-\pi i B_H(v,x)}
 \prod_{\nu=1}^{4f}(U_{v_\nu+x_\nu}\theta_\Lambda)\bigl((gz)_\nu\bigr),
\]
where \(P=\Lambda^{4f}\), \(Q=g(P_1^4)\), \(J=[Q:P\cap Q]\) is prime to \(p\), \(|\iota_p(e)|_p=1\), and \(v=g(t_{\rho,i_1 i_2},t_{\rho,i_1 i_2},t_{\rho,i_1 i_2},t_{\rho,i_1 i_2})\). Since \(g\) is \(\C\)-linear, it commutes with scalar multiplication by \(\pi_0\), so \(M_2\pi_0(v+x)\in Q\). Putting \(N_1=M_2 J\) (still prime to \(p\)) and using \(JQ\subseteq P=\Lambda^{4f}\) gives \(N_1\pi_0(v_\nu+x_\nu)\in\Lambda\) for every \(\nu\) and \(x\). Let \(G_g\) be the \(4f\times 4f\) matrix of \(g\) in the standard complex coordinates of \(\C^{4f}\). Because the auxiliary curve ideal is trivial, \(P_1=\{(\delta\tau(a))_{\tau\in\Sigma}:a\in\mathcal O_E\mathfrak b^{-1}\}\) and \(P=(\delta\mathcal O_{K_d})^{4f}\). Canceling \(\delta\ne0\) from \(JQ\subseteq P\) shows that for every \(a=(a_1,a_2,a_3,a_4)\in\mathcal O_E^4\subseteq(\mathcal O_E\mathfrak b^{-1})^4\) the vector \(J G_g((\tau(a_q))_{q,\tau})\) has entries in \(\varphi_0(\mathcal O_{K_d})\), where \(\varphi_0:K_d\hookrightarrow\C\) is the common restriction of all \(\tau\in\Sigma\); hence \(\iota_p(G_g((\tau(a_q))_{q,\tau}))\in\mathcal O_{\C_p}^{4f}\) because \(p\nmid J\). Choose a \(\Z\)-basis \((e_1,\dots,e_{2f})\) of \(\mathcal O_E\). The \(2f\times 2f\) matrix \(V=(\iota_p(\sigma(e_\kappa)))_{\sigma\in\Sigma\cup\bar\Sigma,\,1\le\kappa\le 2f}\) has entries in \(\mathcal O_{\C_p}\) and satisfies \((\det V)^2=\operatorname{discr}(E)\); since \(p\nmid\operatorname{discr}(E)\), we have \(V\in\GL_{2f}(\mathcal O_{\C_p})\). Each standard basis vector of \(\C_p^{4f}\) is therefore an \(\mathcal O_{\C_p}\)-linear combination of the embedded vectors \((\iota_p(\tau(e_\kappa))\delta_{q,q_0})_{q,\tau}\), so every entry of \(\iota_p(G_g)\) lies in \(\mathcal O_{\C_p}\). For the \(f\)-row submatrix \(B_{i_1 i_2}\) of \(B\) at entry \((i_1,i_2)\), the block-diagonal matrix \(\widetilde B=\operatorname{diag}(B_{i_1 i_2},B_{i_1 i_2},B_{i_1 i_2},B_{i_1 i_2})\) has entries of \(p\)-adic norm at most \(p^{-1}\), so \(\iota_p(G_g)\widetilde B\) has entries of \(p\)-adic norm at most \(p^{-1}=|N_1 p|_p\) and factors as \((N_1 p)L\) with \(L\) having entries in \(\mathcal O_{\C_p}\).

Next establish the one-variable Hurwitz estimate on the elliptic curve \(E_\Lambda=\C/\Lambda\), where \(\Lambda=\delta\mathcal O_{K_d}\) has the discriminant normalization \(\delta^{12}=\Delta(\mathcal O_{K_d})\) of \ref{B5.07x} (so \(\Delta(\Lambda)=1\)) and \(\theta=\theta_\Lambda\) is its normalized even theta section with \(\theta(0)=1\) and simple zero at the nonzero half-period \(P_0=\delta/2\bmod\Lambda\in E_\Lambda[2]\setminus\{0\}\). Fix \(\tau\in\C\) with \(N_1\pi_0\tau\in\Lambda\), and put \(\beta=N_1\pi_0\in\mathcal O_{K_d}\) and \(m_\beta=\beta\bar\beta=N_1^2 p\). Since \(E_\Lambda\) has complex multiplication by \(\mathcal O_{K_d}\), its \(j\)-invariant \(j(E_\Lambda)=1728 g_2(\Lambda)^3\) is an algebraic integer. As \(p\ge5\), the integers \(2\) and \(3\) are units in \(\Z_p\); from \(1728 g_2(\Lambda)^3=j(E_\Lambda)\) and \(27 g_3(\Lambda)^2=g_2(\Lambda)^3-1\), both \(\iota_p(g_2(\Lambda))\) and \(\iota_p(g_3(\Lambda))\) are algebraic and lie in \(\mathcal O_{\C_p}\). Hence, after transporting coefficients via \(\iota_p\), the short Weierstrass equation
\[
 Y^2=X^3-\frac{\iota_p(g_2(\Lambda))}{4}X-\frac{\iota_p(g_3(\Lambda))}{4},
\]
parametrized analytically over \(\C\) by \((X,Y)=(\wp(z),\wp'(z)/2)\), has coefficients in the ring of integers \(\mathcal O_F\) of a finite extension \(F/\Q_p\) inside \(\C_p\) and unit discriminant \(\Delta(\Lambda)=1\), defining a smooth proper elliptic curve \(\mathcal E/\mathcal O_F\); via \((X,Y)=(\wp(z),\wp'(z)/2)\) and \(\iota_p\), we identify \(E_\Lambda(\C)\) with \(\mathcal E(\C_p)\) and elliptic functions on \(E_\Lambda\) with rational functions on \(\mathcal E_{\C_p}\). Write \(T=-X/Y\) (corresponding analytically to \(-2\wp(z)/\wp'(z)\)) for the formal-group parameter at the origin, and for a rational function \(R\) on \(\mathcal E_{\C_p}\) regular at \(0\), write \(\widehat R(T)\in\C_p[[T]]\) for its formal \(T\)-expansion at \(0\). The invariant differential \(\frac{dX}{2Y}\) equals \(dz\) analytically and \(\omega(T)\,dT\) formally, with \(\omega(T)\in 1+T\mathcal O_F[[T]]\) \cite[Ch.~IV, \S1]{silverman-2009}, so \(z=\log_{\widehat E}(T)=\int_0^T\omega(T)\,dT\) and its compositional inverse is \(T(z)=\exp_{\widehat E}(z)\).

For any \(\alpha\in\mathcal O_{K_d}\), multiplication by \(\alpha\) acts on \(T\) by \([\alpha](T)=\exp_{\widehat E}(\iota_p(\alpha)\log_{\widehat E}(T))\). For every \(n\in\Z\), the \(n\)-fold formal group sum \([n](T)=\exp_{\widehat E}(n\log_{\widehat E}(T))\) lies in \(\mathcal O_F[[T]]\); since each coefficient of \(\exp_{\widehat E}(a\log_{\widehat E}(T))\) is a polynomial in \(a\) and \(\iota_p(\alpha)\in\Z_p\) (because \(p\) splits in \(K_d\)), density of \(\Z\) in \(\Z_p\) gives \([\alpha](T)\in\iota_p(\alpha)T+T^2\mathcal O_F[[T]]\). Because \(\mathcal E/\mathcal O_F\) is the smooth Néron model of its generic fiber, \([\alpha]\) extends uniquely to an endomorphism of the group scheme \(\mathcal E/\mathcal O_F\), so \(T([\alpha](Q))=[\alpha](T(Q))\) for every point \(Q\) in the origin residue disk \(\widehat E(\mathfrak m_{\C_p})\). Since \(|\iota_p(2\beta)|_p=|2N_1|_p|\iota_p(\pi_0)|_p=1\), the ultrametric inequality gives \(|[2\beta](t)|_p=|t|_p\) for all \(t\in\mathfrak m_{\C_p}\), whence
\[
 E_\Lambda[2\beta]\cap\widehat E(\mathfrak m_{\C_p})=\{0\}.
\]

Enlarge the finite extension \(F/\Q_p\) inside \(\C_p\) so that all points of \(E_\Lambda[2\beta]\) are rational over \(F\). Consider the elliptic functions \(E_\beta(z)=\theta(\beta z)/\theta(z)^{m_\beta}\) and \(G_0(z)=E_\beta(z)^{-1}=\theta(z)^{m_\beta}/\theta(\beta z)\), which satisfy \(E_\beta(0)=G_0(0)=1\) (the period multipliers cancel because \(H(\beta z,\beta l)=m_\beta H(z,l)\) and \(\alpha(\beta l)=\alpha(l)^{m_\beta}\) for \(l\in\Lambda\)). On \(E_\Lambda\),
\[
 \operatorname{div}(G_0)=m_\beta[P_0]-\sum_{Q\in[\beta]^{-1}(P_0)}[Q],
\]
whose support is contained in \(E_\Lambda[2\beta]\setminus\{0\}\) (as \([\beta](0)=0\ne P_0\) and \([2](P_0)=0\)), hence is disjoint from \(\widehat E(\mathfrak m_{\C_p})\); both \(E_\beta\) and \(G_0\) are defined over \(F\). As in the proofs of \ref{B4.22f} and \ref{B4.22g}, let \(\pi_F\) be a uniformizer of \(\mathcal O_F\) and \(\eta_s\) the generic point of the smooth special fiber \(\mathcal E_s\): if \(R\in F(\mathcal E)^\times\) is regular at \(0\) and its generic-fiber poles on \(\mathcal E_F\setminus\{0\}\) are cleared by a monic polynomial \(P(X)\in\mathcal O_F[X]\) of degree \(N_P\) (in particular whenever all poles lie in \(\mathcal E(F)\setminus\widehat E(\mathfrak m_{\C_p})\), since \(Q\notin\widehat E(\mathfrak m_{\C_p})\) implies \(|X(Q)|_p\le 1\), so \(X(Q)\in\mathcal O_F\)), then \(v_{\eta_s}(P(X))=0\) and \(P(X)R=\pi_F^{v_{\eta_s}(R)}U\) with \(U=A(X)+YB(X)\in\mathcal O_F[X,Y]\setminus\pi_F\mathcal O_F[X,Y]\) of pole order at most \(2N_P\) at \(0\) (so \(\deg A\le N_P\) and \(2\deg B+3\le 2N_P\)); since \(T^2X(T),-T^3Y(T)\in 1+T\mathcal O_F[[T]]\) \cite[Ch.~IV, \S1]{silverman-2009}, multiplying by \(T^{2N_P}\) gives \(T^{2N_P}P(X(T))\in\mathcal O_F[[T]]^\times\) and \(T^{2N_P}U(X(T),Y(T))\in\mathcal O_F[[T]]\setminus\pi_F\mathcal O_F[[T]]\), whence \(\widehat R(T)\in\pi_F^{v_{\eta_s}(R)}(\mathcal O_F[[T]]\setminus\pi_F\mathcal O_F[[T]])\). Applying this to \(E_\beta\) and \(G_0=E_\beta^{-1}\), since \(\widehat E_\beta(0)=\widehat G_0(0)=1\) forces \(v_{\eta_s}(E_\beta)\le 0\) and \(v_{\eta_s}(G_0)\le 0\) while \(v_{\eta_s}(E_\beta)+v_{\eta_s}(G_0)=0\), we obtain \(v_{\eta_s}(E_\beta)=v_{\eta_s}(G_0)=0\) and \(\widehat E_\beta(T),\widehat G_0(T)\in\mathcal O_F[[T]]^\times\). Write \(f_0(T)\in\C_p[[T]]\) for the \(T\)-expansion of \(\theta\) transported by \(\iota_p\). From \(\theta(\beta z)=\theta(z)^{m_\beta}E_\beta(z)\) we obtain \(f_0([\beta](T))=f_0(T)^{m_\beta}\widehat E_\beta(T)\). Writing \(f_0(T)=\sum_{k=0}^\infty a_k T^k\) with \(a_0=1\) and comparing \(T^k\)-coefficients for \(k\ge1\) yields \((\iota_p(\beta)^k-m_\beta)a_k\in\mathcal O_{\C_p}\) by induction on \(k\), because \([\beta](T)\in\iota_p(\beta)T+T^2\mathcal O_F[[T]]\) and \(\widehat E_\beta(T)\in 1+T\mathcal O_F[[T]]\). Since \(|\iota_p(\beta)|_p=1\) and \(p\mid m_\beta\), the factor \(\iota_p(\beta)^k-m_\beta\) is a unit in \(\Z_p\), so \(f_0(T)\in 1+T\mathcal O_{\C_p}[[T]]\subseteq\mathcal O_{\C_p}[[T]]^\times\) and hence \(f_0([\beta](T))\in\mathcal O_{\C_p}[[T]]^\times\).

Now establish the Gauss-norm bound on the étale \(N_1\pi_0\)-torsion disk via multiplication by \(\beta=N_1\pi_0\). Write \(H_\tau(T)\in\C_p[[T]]\) for the \(T\)-expansion of \(U_\tau\theta\) transported by \(\iota_p\), and define \(G_\tau(z)=(U_\tau\theta)(z)^{m_\beta}/\theta(\beta z)\). Because \(\beta\tau\in\Lambda\), the transformation law \(W_{\beta\tau}\theta=\theta\) evaluated at \(\beta z\), together with \(H(\beta z,\beta\tau)=m_\beta H(z,\tau)\) and \(H(\beta\tau,\beta\tau)=m_\beta H(\tau,\tau)\), gives
\[
 \theta(\beta z)
 =\alpha(\beta\tau)\,\theta(\beta z+\beta\tau)
 \exp\!\left(-\pi m_\beta H(z,\tau)-\tfrac{\pi}{2}m_\beta H(\tau,\tau)\right).
\]
Dividing \((U_\tau\theta)(z)^{m_\beta}=\theta(z+\tau)^{m_\beta}\exp(-\pi m_\beta H(z,\tau)-\frac{\pi}{2}m_\beta H(\tau,\tau))\) by \(\theta(\beta z)\) cancels the exponential factors and yields
\[
 G_\tau(z)=\alpha(\beta\tau)^{-1}\frac{\theta(z+\tau)^{m_\beta}}{\theta(\beta(z+\tau))}
 =\alpha(\beta\tau)^{-1}G_0(z+\tau),
\]
where \(\alpha(\beta\tau)=\pm1\). Put \(P_\tau=\tau\bmod\Lambda\in E_\Lambda[\beta]\subseteq\mathcal E(\mathcal O_F)\). Translation by \(P_\tau\) is an automorphism of the smooth group scheme \(\mathcal E/\mathcal O_F\), so it preserves the valuation at \(\eta_s\) and gives \(v_{\eta_s}(G_\tau)=v_{\eta_s}(G_0)=0\). Since \([\beta](P_\tau)=0\), the divisor of \(G_\tau\) on \(E_\Lambda\) is
\[
 \operatorname{div}(G_\tau)=m_\beta[P_0-P_\tau]-\sum_{Q\in[\beta]^{-1}(P_0)}[Q],
\]
where \(P_0-P_\tau\in E_\Lambda[2\beta]\) because \([2\beta](P_0-P_\tau)=[2\beta](P_0)-2[\beta](P_\tau)=0\). Passing to \(T\)-expansions gives \(H_\tau(T)^{m_\beta}=\widehat G_\tau(T)\,f_0([\beta](T))\). We treat the two cases for \(P_\tau\):
\begin{itemize}
\item If \(P_\tau\ne P_0\), then \(P_0-P_\tau\in E_\Lambda[2\beta]\setminus\{0\}\), so neither \(P_0-P_\tau\) nor any point of \([\beta]^{-1}(P_0)\) lies in \(\widehat E(\mathfrak m_{\C_p})\). Thus both \(G_\tau\) and \(G_\tau^{-1}\) have valuation \(0\) at \(\eta_s\), are regular on \(\widehat E(\mathfrak m_{\C_p})\), and have all horizontal poles outside \(\widehat E(\mathfrak m_{\C_p})\). Applying the generic-point criterion above to both \(G_\tau\) and \(G_\tau^{-1}\) (whose poles all lie in \(E_\Lambda[2\beta]\setminus\{0\}\subseteq\mathcal E(F)\setminus\widehat E(\mathfrak m_{\C_p})\)) shows that \(\widehat G_\tau(T)\) and \(\widehat G_\tau(T)^{-1}\) lie in \(\mathcal O_F[[T]]\), so \(\widehat G_\tau(T)\in\mathcal O_{\C_p}[[T]]^\times\). Set \(R_\tau(T)=H_\tau(T)\) and \(C(T)=\widehat G_\tau(T)\,f_0([\beta](T))\in\mathcal O_{\C_p}[[T]]^\times\).
\item If \(P_\tau=P_0\), then \(0\notin[\beta]^{-1}(P_0)\), so \(G_\tau\) has a zero of exact order \(m_\beta\) at \(0\) and poles at \([\beta]^{-1}(P_0)\), all outside \(\widehat E(\mathfrak m_{\C_p})\). Consider the rational function \(\widetilde G_\tau=T^{-m_\beta}G_\tau\) on \(\mathcal E\). On \(\mathcal E/\mathcal O_F\), the Weierstrass coordinates \(X,Y\) are nonzero on \(\mathcal E_s\), so \(v_{\eta_s}(T)=0\) and hence \(v_{\eta_s}(\widetilde G_\tau)=0\). Inside \(\widehat E(\mathfrak m_{\C_p})\), \(T=-X/Y\) has a single simple zero at \(0\) and no other zeros or poles (since \(0<|T(Q)|_p<1\) for \(Q\in\widehat E(\mathfrak m_{\C_p})\setminus\{0\}\)). Therefore both \(\widetilde G_\tau\) and \(\widetilde G_\tau^{-1}\) have valuation \(0\) at \(\eta_s\), are regular on \(\widehat E(\mathfrak m_{\C_p})\), and have all horizontal poles outside \(\widehat E(\mathfrak m_{\C_p})\). The poles of \(\widetilde G_\tau=(-Y/X)^{m_\beta}G_\tau\) lie in \([\beta]^{-1}(P_0)\subseteq\mathcal E(F)\setminus\widehat E(\mathfrak m_{\C_p})\) and at \(X=0\) (cleared by \(X^{m_\beta}\prod_{Q\in[\beta]^{-1}(P_0)}(X-X(Q))\in\mathcal O_F[X]\)), and the poles of \(\widetilde G_\tau^{-1}=(-X/Y)^{m_\beta}G_\tau^{-1}\) lie in \(E_\Lambda[2]\setminus\{0\}\subseteq\mathcal E(F)\setminus\widehat E(\mathfrak m_{\C_p})\) (cleared by \(\prod_{Q\in E_\Lambda[2]\setminus\{0\}}(X-X(Q))^{m_\beta}\in\mathcal O_F[X]\)); applying the generic-point criterion above to \(\widetilde G_\tau\) and \(\widetilde G_\tau^{-1}\) therefore gives \(\widehat{\widetilde G_\tau}(T)\in\mathcal O_F[[T]]^\times\subseteq\mathcal O_{\C_p}[[T]]^\times\). Because \(\theta\) has a simple zero at \(P_0\) and \(T(z)=z+O(z^5)\), the series \(R_\tau(T)=H_\tau(T)/T\) lies in \(\C_p[[T]]\) and satisfies \(R_\tau(T)^{m_\beta}=C(T):=\widehat{\widetilde G_\tau}(T)\,f_0([\beta](T))\in\mathcal O_{\C_p}[[T]]^\times\).
\end{itemize}
In both cases \(R_\tau(T)^{m_\beta}=C(T)\in\mathcal O_{\C_p}[[T]]^\times\), and from \(R_\tau(0)^{m_\beta}=C(0)\in\mathcal O_{\C_p}^\times\) we obtain \(R_\tau(0)\in\mathcal O_{\C_p}^\times\).

We now pass from the formal-group parameter \(T\) to the scaled additive coordinate \(z=N_1 p\,\tilde z\) and extract the \(m_\beta\)-th root in \(\mathcal O_{\C_p}\langle\langle\tilde z\rangle\rangle\). Define the derivation \(D=\omega(T)^{-1}\frac{d}{dT}\) on \(\mathcal O_{\C_p}[[T]]\).
\leanref{Formalism.\allowbreak{}L2.\allowbreak{}Th1.\allowbreak{}A.\allowbreak{}D}{Formalism.\allowbreak{}LocalToGlobal.\allowbreak{}TorsionIntegrality.\allowbreak{}WeierstrassTwoTorsion.\allowbreak{}CompInv}{57--58}
Under \(z=\log_{\widehat E}(T)\), \(D\) corresponds to \(\frac{d}{dz}\). Because \(\omega(T)^{-1}\in 1+T\mathcal O_{\C_p}[[T]]\), induction on \(k\ge0\) shows \(D^k(T)\in\mathcal O_{\C_p}[[T]]\), so \((D^k T)(0)\in\mathcal O_{\C_p}\). The compositional inverse \(T(z)=\exp_{\widehat E}(z)\) therefore satisfies
\[
 T(z)=\sum_{k=1}^\infty\frac{(D^k T)(0)}{k!}\,z^k.
\]
By Legendre's formula, for every \(k\ge1\) and \(p\ge5\),
\[
 v_p(k!)=\sum_{r\ge1}\left\lfloor\frac{k}{p^r}\right\rfloor<\frac{k}{p-1}<k,
\]
so \(v_p(k!)\le k-1\) and \(p^k/k!\in p\Z_p\). Substituting \(z=N_1 p\,\tilde z\) yields
\[
 T(N_1 p\,\tilde z)
 =\sum_{k=1}^\infty (D^k T)(0)\,N_1^k\frac{p^k}{k!}\,\tilde z^k
 \in p\,\tilde z\,\mathcal O_{\C_p}[[\tilde z]].
\]
Substituting \(T=T(N_1 p\,\tilde z)\) into \(R_\tau(T)^{m_\beta}=C(T)\) gives
\[
 R_\tau\bigl(T(N_1 p\,\tilde z)\bigr)^{m_\beta}
 =C(0)\bigl(1+p\,y(\tilde z)\bigr),\qquad
 y(\tilde z)\in\tilde z\,\mathcal O_{\C_p}[[\tilde z]].
\]
Because \(\C_p\) has characteristic \(0\) and \(R_\tau(0)\in\mathcal O_{\C_p}^\times\) satisfies \(R_\tau(0)^{m_\beta}=C(0)\), uniqueness of the formal \(m_\beta\)-th root with given nonzero constant term gives
\[
 R_\tau\bigl(T(N_1 p\,\tilde z)\bigr)
 =R_\tau(0)\sum_{k=0}^\infty\left(k!\binom{1/m_\beta}{k}p^k\right)\frac{y(\tilde z)^k}{k!}.
\]
Since \(m_\beta=N_1^2 p\) with \(p\nmid N_1\), we have
\[
 k!\binom{1/m_\beta}{k}p^k
 =\prod_{s=0}^{k-1}\bigl(N_1^{-2}-sp\bigr)\in\Z_p\subseteq\mathcal O_{\C_p}
\]
for every \(k\ge0\). Write \(y(\tilde z)=\sum_{r\ge1}b_r\tilde z^r\) with \(b_r\in\mathcal O_{\C_p}\). By the multinomial expansion, for every \(k\ge0\) and \(n\ge k\), the coefficient of \(\tilde z^n\) in \(y(\tilde z)^k/k!\) multiplied by \(n!\) equals
\[
 \sum_{\substack{(k_r)_{r\ge1}\\ \sum k_r=k,\;\sum r k_r=n}}
 \frac{n!}{k!}\cdot\frac{k!}{\prod_{r\ge1}k_r!}\prod_{r\ge1}b_r^{k_r}\in\mathcal O_{\C_p},
\]
because \(n=\sum r k_r\ge\sum k_r=k\) makes both \(n!/k!\) and \(k!/\prod k_r!\) integers; for \(n<k\) the coefficient vanishes. Hence \(y(\tilde z)^k/k!\in\mathcal O_{\C_p}\langle\langle\tilde z\rangle\rangle\) for all \(k\ge0\). Since each \(\tilde z^n\)-coefficient receives a finite sum over \(0\le k\le n\), it follows that \(R_\tau(T(N_1 p\,\tilde z))\in\mathcal O_{\C_p}\langle\langle\tilde z\rangle\rangle\). When \(P_\tau\ne P_0\), this gives \(\iota_p((U_\tau\theta)(N_1 p\,\tilde z))=R_\tau(T(N_1 p\,\tilde z))\in\mathcal O_{\C_p}\langle\langle\tilde z\rangle\rangle\); when \(P_\tau=P_0\), multiplying by \(T(N_1 p\,\tilde z)\in\mathcal O_{\C_p}[[\tilde z]]\subseteq\mathcal O_{\C_p}\langle\langle\tilde z\rangle\rangle\) gives \(\iota_p((U_\tau\theta)(N_1 p\,\tilde z))=T(N_1 p\,\tilde z)\,R_\tau(T(N_1 p\,\tilde z))\in\mathcal O_{\C_p}\langle\langle\tilde z\rangle\rangle\).

To return to \(U_t e_{N,N'}\), take four independent copies \(Z^{(1)},Z^{(2)},Z^{(3)},Z^{(4)}\) of the variable tuple \(Z\) and substitute the \(\mathcal O_{\C_p}\)-linear forms \((L(Z^{(1)},Z^{(2)},Z^{(3)},Z^{(4)}))_\nu\) (from \(\iota_p(G_g)\widetilde B=(N_1 p)L\)) into the one-variable series \(\iota_p((U_{v_\nu+x_\nu}\theta_\Lambda)(N_1 p\,\tilde z))\). Because \(|\iota_p(e/\sqrt J)|_p=1\), \(\alpha(x)=\pm1\), and \(e^{-\pi i B_H(v,x)}\) is a root of unity, closure of the Hurwitz ring under \(\mathcal O_{\C_p}\)-linear substitutions, products over \(\nu=1,\dots,4f\), and \(\mathcal O_{\C_p}\)-linear combinations over \(x\in Q/(P\cap Q)\) yields
\[
 \prod_{q=1}^4\mathcal S(U_{t_{\rho,i_1 i_2}}\Theta;B_{i_1 i_2})(Z^{(q)})
 =\mathcal S\!\left((U_{t_{\rho,i_1 i_2}}\Theta)^{\otimes 4};\widetilde B\right)
 \in\mathcal O_{\C_p}\langle\langle Z^{(1)},Z^{(2)},Z^{(3)},Z^{(4)}\rangle\rangle .
\]
Write \(\mathcal S(U_{t_{\rho,i_1 i_2}}\Theta;B_{i_1 i_2})(Z)=\sum_\alpha c_\alpha Z^\alpha\). Reading the multi-index \((\alpha,\alpha,\alpha,\alpha)\) gives \((\alpha,\alpha,\alpha,\alpha)!\,c_\alpha^4=(\alpha!\,c_\alpha)^4\in\mathcal O_{\C_p}\), hence \(\alpha!\,c_\alpha\in\mathcal O_{\C_p}\) for every \(\alpha\), so \(\mathcal S(U_{t_{\rho,i_1 i_2}}\Theta;B_{i_1 i_2})\in\mathcal O_{\C_p}\langle\langle Z\rangle\rangle\). Taking products over all \(N_{j-1}N_j\) matrix entries \((i_1,i_2)\) and summing over \(\rho\) with coefficients of \(p\)-adic norm at most \(1\) proves \(\mathcal S(U_t e_{N,N'};B)\in\mathcal O_{\C_p}\langle\langle Z\rangle\rangle\), completing the étale translation estimate.

If \(|\iota_p(\bar\pi_i)|_p<1\), then
\(|\iota_p(\bar\pi_i)|_p=p^{-1}\) and \(\iota_p(\pi_i)\) is
a unit, since \(\pi_i\bar\pi_i=p\) and the prime splits
unramified in \(K_d\). Thus \(A_{\mathrm{fr}}\) has entries
of norm at most \(p^{-1}\). For \(a\in S_a\) one has
\(\pi_i a\in M\) and \(p a\in M\), so \(m_0\pi_i\) and \(m_0 p\)
kill \(t_k+a\) modulo \(M\). Expand \(U_{t_k}W_a\) using the
translation law. The phase is a root of unity and the
coefficient \(q r_k\widehat\Psi(a)\) has norm at most one.
The étale estimate on the shifted series gives
\eqref{eq:B4.23a-split}.

If \(|\iota_p(\bar\pi_i)|_p=1\), set \(s=\pi_i/\bar\pi_i\),
so \(|\iota_p(s)|_p=p^{-1}\). The scalar change-of-frame
formula gives
\[
 e_{N,N'}(sz)=c\,q^{-1/2}\sum_{a\in S_a}W_a e_{N,N'}(z),
 \qquad c^8=(\bar\pi_i/\pi_i)^{4d_0}.
\]
It is the Cayley Gauss comparison of \ref{B4.22e} for the
scalar isometry \(s\), with average denominator \(\sqrt q\).
Expand \(\widehat\Psi\) and exchange the sums in \(a,b\). Then commute \(W_b\) through the \(a\)-sum. The character produced
in that commutation is \(e^{-2\pi iB_H(b,a)}\), since \(W_b\)
fixes \(e_{N,N'}\) for \(b\in S_b\subset M\).
Transport this \(W_b\) by the isometry \(s\). The resulting
formula is
\[
 G(z)=D^{-1}\sum_{b\in S_b}\Psi(b)\alpha(b)
 (U_{sb}e_{N,N'})(sz),\qquad
 D=c\,q^{1/2},\qquad D^8=\bar\pi_i^{8d_0}.
\]
Thus \(D\) is a unit. After \(U_{t_k}\) the translates in
this sum are \(s(t_k+b)\), with the additional unit phase
\(e^{-\pi i B_H(st_k,sb)}\). They are killed modulo \(M\)
by \(m_0\bar\pi_i\) and \(m_0 p\), since
\(\bar\pi_i s=\pi_i\) and \(ps=\pi_i^2\).
Apply the same étale estimate with the matrix
\(\iota_p(s)A_{\mathrm{fr}}\), whose entries have norm at
most \(p^{-1}\). In this direction the series in \eqref{eq:B4.23a-split} lies in \(\mathcal O_{\C_p}\langle\langle Z\rangle\rangle\) before multiplication by \(p^{d_0}\).

In the diagonal differential frame \(A_{\mathrm{fr}}\) gives
diagonal entries \(\iota_p(\bar\pi_i)c_{D,\tau}\), with the
generators read through the comparison. Extracting the
coefficient at multi-index \(\alpha\) and multiplying by
\(\alpha!\) gives the jet with the factor
\(\prod_\tau\tau(\bar\pi_i)^{A_{u,j,\tau}}\) when the
degree at \(\tau\) is \(A_{u,j,\tau}\). The normalizing
character on the standard pair is a unit, and integrality
of \(\Xi\) permits the weighted sum. Thus
\begin{equation}\label{eq:B4.23a-jet}
 d_j J_\alpha\in\mathcal O_{(w)},\qquad
 d_j=p^{fN_{j-1}N_j}\prod_\tau
 \tau(\bar\pi_i)^{A_{u,j,\tau}} .
\end{equation}
\leanref{Formalism.\allowbreak{}L2.\allowbreak{}AlgTh.\allowbreak{}jetAtSplit\_\allowbreak{}of}{Formalism.\allowbreak{}LocalToGlobal.\allowbreak{}AlgebraicTheory.\allowbreak{}AlgebraicTheoryProofs.\allowbreak{}CardTwoSidedQuotient}{349--602}

On a tensor Weyl generator perform the polynomial substitution
of \ref{B4.22b} with its determinant factors and the matrix
entries as indeterminates. Cauchy--Binet expands each minor
in target minors whose column heights are the source column heights
together with the full-height source columns added by the determinant
power; by weight compatibility these are the column heights of the target partition \(\kappa_\tau\) at \(\tau\), of total box count \(|\kappa_\tau|=A_{u,j,\tau}\). Straightening by Pl\"ucker relations expresses each differential coefficient in the standard bideterminant basis of the target dual Weyl lattice \(\nabla_{\mathcal O_{(w)}}(\kappa_\tau)\) \cite[Theorem~II.2.16]{akin-1982}. To pass from \(\nabla_{\mathcal O_{(w)}}(\kappa_\tau)\) to the target Weyl lattice \(\Delta_{\mathcal O_{(w)}}(\kappa_\tau)\), recall from Lemma~\ref{B3.10a} \cite[Proposition~II.4.1]{akin-1982} that if \((e_r)\) is an \(\mathcal O_{(w)}\)-basis of \(\Delta_{\mathcal O_{(w)}}(\kappa_\tau)\) and \((d_r)\) is its dual basis over \(K\) under the factorial pairing \([P,Q]=\sum_\alpha\alpha![Y^\alpha]P[Y^\alpha]Q\), then \((C_{\kappa_\tau}d_r)\) is an \(\mathcal O_{(w)}\)-basis of \(\nabla_{\mathcal O_{(w)}}(\kappa_\tau)\), where \(C_{\kappa_\tau}=[H_{\kappa_\tau},H_{\kappa_\tau}]\). Every \(Q\in\nabla_{\mathcal O_{(w)}}(\kappa_\tau)\) has coefficients in \(\mathcal O_{(w)}\), so \([Q,C_{\kappa_\tau}d_r]\in\mathcal O_{(w)}\) and \(C_{\kappa_\tau}Q=\sum_r[Q,C_{\kappa_\tau}d_r]e_r\in\Delta_{\mathcal O_{(w)}}(\kappa_\tau)\). By the repeated column-addition Capelli identity in Lemma~\ref{B5.07y} \cite{howe-umeda-1991}, \(C_{\kappa_\tau}=\operatorname{Hook}(\kappa_\tau)\) is the product of the hook lengths of the diagram of \(\kappa_\tau\). Every hook of \(\kappa_\tau\) is a subset of its \(|\kappa_\tau|=A_{u,j,\tau}\le D_{u,j}\) boxes, and \(2D_{u,j}+2<p\) gives \(D_{u,j}+1<p\), so every hook length of \(\kappa_\tau\) is positive and strictly less than \(p\) (likewise \(\kappa_{\tau,1}-\kappa_{\tau,N_j}+N_j-1\le A_{u,j,\tau}+N_m-1\le D_{u,j}+N_m-1<p\) places \(\kappa_\tau\) inside the fundamental \(p\)-alcove where Jantzen's sum formula gives \(\Delta_{\mathcal O_{(w)}}(\kappa_\tau)=\nabla_{\mathcal O_{(w)}}(\kappa_\tau)\) \cite[Part~II, Corollary~5.6]{jantzen-2003}). Thus \(C_{\kappa_\tau}\in\mathcal O_{(w)}^\times\) and \(\nabla_{\mathcal O_{(w)}}(\kappa_\tau)=\Delta_{\mathcal O_{(w)}}(\kappa_\tau)\), so each differential coefficient belongs to the target Weyl lattice \(\Delta_{\mathcal O_{(w)}}(\kappa_\tau)\).
Since the jet degree at \(\tau\) is \(A_{u,j,\tau}\), \eqref{eq:B4.23a-jet} gives integral jet multipliers at \(w\) after scaling by \(d_j\). By linearity and equivariance the scaled operator preserves the local lattice condition.

Class weights are units by the stabilizer calculation of
\ref{B4.22a}. In the last two kernels use the same decoration
\(Y\) on the intermediate lattice \(N\). Such a decoration is a coset \(y+L_+\). At a row prime \(P\), write \(q_P=|k(P)|\). At the distinguished generator primes the local conditions are \(y\in q_P^{-1}N_P\) and \((L_+)_P=N_P\). At the other row primes they are \(y\in q_P^{-M}N_P\) and \((L_+)_P=q_P^M N_P\), where \(M\) is the row modulus exponent. These restrictions are imposed at the row primes alone. Multiply their bounds before summing
over \(Y\). Put \(d_*=\prod_{j=1}^m d_j\). The row tie
lifts \(H\), and the pairing tie gives the rational period
\(y\). In terminal class coordinates, with unscaled column
values \(\psi_c\) and row values \(H_c\), the class sum is
\[
 d_*y=\sum_c |\operatorname{Stab}(c)|^{-1}
 \langle\operatorname{br}(d_*\psi_c),H_c\rangle .
\]
The interlacing weights have the same determinant shift,
so integral branching sends the source Weyl lattice into
the target Weyl lattice. The level factors introduced by branching are integral. Since \(\pi_T\bar\pi_T=T_u\ne g_i\), both are units at \(w\).
The scaled column lies in the source lattice, and \(H_c\)
lies in the pairing dual of the target lattice. Every summand
is integral. Hence, for \(y\ne0\),
\[
 \frac{v_w(y)}{e_w}\ge-\sum_j(\nu_{i,u,j,w}+fN_{j-1}N_j)
 \ge-\sum_j\nu_{i,u,j,w}-\mathcal E_G .
\]
\leanref{GlobalTheory.\allowbreak{}ColPeriod.\allowbreak{}valuation\_\allowbreak{}PtilE\_\allowbreak{}le\_\allowbreak{}one}{Development.\allowbreak{}AbelianVarietyTheory.\allowbreak{}Proofs.\allowbreak{}SatakeParameters.\allowbreak{}NotDvdCardOfInjective}{268--305}

For the budget put \(A=f(N_m^2+m^2\eta)\).
By \ref{B4.02}--\ref{B4.03} the base terms cancel and
\(F_j-E_{u,j}=-\nu_{i,u,j,w}-c_g A-\delta_{u,w}\).
Since \(t\le N_m\), one has \(\mathcal E_G\le2mA\), so
\(mc_g A+m\delta_{u,w}\ge C_1\mathcal E_G\).
Substitution gives \(B_{u,w}\le-\sum_j\nu_{i,u,j,w}-C_1\mathcal E_G\).

For the family size assertion the middle threshold gives
\(c_g\ge2\). Put \(A'=kb\ge10\), \(N=N_m\ge2\).
The choice of fixed primes gives a size parameter \(Y\) with
\((A')^8\le Y\le g_i\) and
\(N^4(Q_Sqd)^{C_Y}\le Y^{1/4}\), with \(C_Y>0\) and \(Q_Sqd\ge1\). Thus \(N^{16}\le Y\).
For \(D_{u,j}\), the rows are decreasing, so use the zeroth
entry of \ref{B4.03}. It is at most the first input weight
plus \(w_j(\tau)\) minus \((2N_{j-1}+t)/2\).
Write \(A_*=xb\). The parameter choice gives
\(m(t+1)\eta\le w''\) and \(kS_c\le w''\), where \(S_c\) is the cloud-coordinate scale of the member, with
\(w''\le b'\). Since \(m(t+1)\ge2\) and \(k\ge5\), it follows
that \(2\eta\le b'\) and \(5S_c\le b'\).
At a top type the first input weight is \(a_T+(k-1)/2\)
with \(a_T\le6A_*/5\), and \(w_j\le b_T+b'/x+\eta\)
with \(b_T\le2b'\). Using \(b'\le A_*/2\) gives an entry
at most \(6A_*/5+7b'/2\le3A_*\).
Off the top the parameter choice gives \(u_0-v_0\le2S_c\), and the first cloud coordinate satisfies \(v_0\le S_c\).
Thus \(u_0\le3S_c\).
The thick and thin windows give \(w_j\le b+\eta+1\) and
\(w_j\le b'+\eta\), respectively. With the bounds on \(S_c,\eta\) and \(1\le b\le A_*\), the row entry is
again at most \(3A_*\). Summing over \(N_{j-1}\) coordinates
and \(f\) types and using \(x,f\le k\) gives
\(D_{u,j}\le3(A')^2N\).
\leanref{GlobalTheory.\allowbreak{}CoreHt.\allowbreak{}RowsRange.\allowbreak{}nuRow\_\allowbreak{}zero\_\allowbreak{}le\_\allowbreak{}three\_\allowbreak{}mul}{Development.\allowbreak{}AbelianVarietyTheory.\allowbreak{}Proofs.\allowbreak{}DecoratedLatticeProofs.\allowbreak{}NuRowZeroLeThreeMul}{34--91}

Each size expression is at most \(6(A')^2N+2\). For
\(N\le(A')^5\) this is at most \(6(A')^7+2<(A')^8\).
For \(N>(A')^5\) it is at most \(6N^2+2<N^{16}\).
These bounds are uniform in the member index.
\leanref{GlobalTheory.\allowbreak{}AtG.\allowbreak{}KernelSizes.\allowbreak{}sizeAt\_\allowbreak{}of\_\allowbreak{}params}{Development.\allowbreak{}AbelianVarietyTheory.\allowbreak{}Proofs.\allowbreak{}SatakeParameters.\allowbreak{}KernelDegreesSmallAgainstTypePrimes}{115--182}
\end{proof}

\subsection*{Proof of \ref{B4.25}}\label{prf:B4.25}\phantomsection
\begin{proof}
By \ref{B3.51}, \(G_{p,j}\ge0\). The screened gain vanishes on \(\mathcal D\). The raw gain vanishes outside the finite support of the size data, and also when \(p_w\) lies in neither the inert nor the signal class. It also vanishes at \(g_0,g_1\).
In the inert and signal classes, the primes dividing
\(Q_Sqd\,g_0g_1\) are fixed primes. In the remaining
cases the filling identity of \ref{B4.04} identifies the raw gain with the counted gain in \eqref{eq:p4c2-counted-gain}. Outside \(\mathcal D\) the inert case of odd order greater than one is excluded by hypothesis, so \eqref{eq:p4c2-served-gain} identifies \(G_{p_w,j}\) with the same counted gain.
\end{proof}

\subsection*{Proof of \ref{B4.26}}\label{prf:B4.26}\phantomsection
\begin{proof}
The denominator \(1\) gives \(\phi\in(N_u)_{(w)}\).
For nonempty selection take the permitted model and use
the equivalence to put \(c_{u,w}(j_u\phi)\) in \(N_{P,u,w}\). For empty selection the assumed inclusion
puts \(\phi\) in \((M_{H,u})_{(w)}\), so its image lies in
\((\Lambda^0_{u,0})_{(w)}\cap V^r=S^{\mathrm{ram}}_{u,w}\).
Proposition~\ref{B4.21a} supplies these local properties on
the arithmetic member. In the first case the comparison image belongs to the enhanced lattice. In the second case the selection is empty and the initial vector belongs to the localized ramified lattice. Either condition admits the input to \eqref{eq:p4c2-enhanced-column}. By \ref{B4.25}, the selected column has bound \(s^{\rm enh}_{u,w}\ge S_{u,w}+\delta_{u,w}\). The common stabilizer denominator in \ref{B4.03} allows us to apply \ref{B4.19}, which subtracts \(\delta_{u,w}\) from this exponent.
\end{proof}

\subsection*{Proof of \ref{B4.30}}\label{prf:B4.30}\phantomsection
\begin{proof}
Take the maximum of the given cutoffs and \(3^{5/e_z}\), where \(J\) is the family scale exponent and \(e_z=\varepsilon/J>0\). To check the permitted-model hypothesis of \ref{B4.26}, fix \(p=p_w\ne T_u\) with \(W_p\ne\varnothing\).
By \ref{B3.51}, \(p\) is an inert candidate and
\(\operatorname{ord}(h_p)\ge x^{\varepsilon/4}\).
An odd order would divide \(o\) because the image in the
cyclic \(2\)-group would vanish. But
\(o\le x^{\varepsilon/100}<x^{\varepsilon/4}\).
The inert class makes \(p\) prime to the conductor.
At each \(P\mid p\), a Frobenius has the same image in the
abelian host group as the representative defining \(h_p\).
Its second component is the image of \(h_p\) under the frame
embedding and has that even order. The inert class excludes \(2\).
The host splitting condition applies to primes up to at least \(\lfloor x^{e_z/5}\rfloor\ge3\).
The host splitting condition therefore forces \(3\) to split
in \(K_d\), incompatible with the inert class of a candidate.
Host level primes are split. These facts and \(p\ne T_u\)
give the conditions of \ref{B4.07}.

For an arithmetic member, \ref{B4.21a} gives both local
equivalences for the same \(N_u\) and its inclusion in
\((M_{H,u})_{(w)}\) away from \(T_u\). For a general member
these are the corresponding clauses of its family data.
The error identities \ref{B4.02}--\ref{B4.04} apply. Together with the agreement of the local fillings and the ramified-space and base-term identities of \ref{B4.14}, they give the error and compatibility hypotheses of \ref{B4.29}. Lemma~\ref{B4.20} gives the required size identities at every level prime. An inert-class prime \(p\) with \(h_p\) of odd order greater than one meets the screening criterion of \ref{B4.04}. Every level prime is split, so \(p\) is distinct from every \(T_u\) and belongs to the discarded set \(\mathcal D\). With filling and period linearity apply \ref{B4.29}.

For closedness use the finite conjunction rule of \ref{B8.39}.
Once the common threshold admits a family, the eighth input
supplies the witness \(C_1\) for that family. Use it in the
common-cutoff argument to obtain both eventual assertions.
\end{proof}

\subsection*{Proof of \ref{B4.31}}\label{prf:B4.31}\phantomsection
\begin{proof}
The design gives the logarithmic stabilizer bound of Definition~\ref{def:p4c3-design} on the entire class. By \ref{B4.16} and the two closed selection hypotheses, the following assertion is closed there: for some \(C\ge0\), the size estimates of \ref{B4.15} and both column estimates of \ref{B4.14} at normalization \(1\) hold beyond one family cutoff.
Lemma~\ref{B4.21} gives closed eventual coprimality.
Together with the four point and period hypotheses, the pair of sharp level-prime estimates and the common fixed-prime bound, these give the eight closed hypotheses of \ref{B4.30} on \(\mathcal D_3\). Theorem~\ref{B4.30} therefore gives closedness of the eventual rounded period bound. Taking the quantifiers in \eqref{eq:p4c2-closed} in order gives \(\varepsilon_1\) and then \(\mathcal K\). The arithmetic
inputs for the constructed family come from the closed rows
of \ref{B8.09} and the fixed-prime row of \ref{B8.06}.
\end{proof}

\section*{Proofs for Chapter~\ref{ch:hecke-row-gluing}: Gluing row lattices across Hecke characters}

\subsection*{Proof of \ref{B4.36a}}\label{prf:B4.36a}\phantomsection
\begin{proof}
Take $U=\{a:\chi_a=\xi_w\text{ for some }w\in\mathrm{Tw}\}$ and $TU=\{\tau_a:a\in U\}$. For each $w\in\mathrm{Tw}$, choose a label $a$ with $\chi_a=\xi_w$. Then $\tau_a\in TU$. Take $t_w=e^{-1}(\tau_a)$. By \ref{B4.33}, $\theta(t_w)=a$, so $\chi_{\theta(t_w)}=\xi_w$. Off $\mathrm{Tw}$ take $t_w=t_0$. 

The valuation of $a\in\mathcal O$ at the conjugate arithmetic place equals that of $\bar a$ at the original place. Thus the place correspondence respects conjugation and $\mathrm{Tw}$ is stable. Off $\mathrm{Tw}$ both types are $t_0$. On $\mathrm{Tw}$ we have $p_{\bar w}=p_w$ and $\psi_{\bar w}=\psi_w$, so $\xi_{\bar w}=\xi_w$. On local units above the selected auxiliary prime, $\psi_w$ is trivial and $\xi_w$ has order $15$. The surjection $a\mapsto a_f$ and the embedding $\sigma$ make the character map in \ref{B4.33} injective, with $\chi_{\theta(t_0)}=1$. Hence $t_{\bar w}=t_w\ne t_0$ on $\mathrm{Tw}$.

In either system, \ref{B4.36} gives $RR_{(z)}=(L_{t_z})_{(z)}$. For $z=w,\bar w$ the type $t_z$ is $t_w$.
\end{proof}

\subsection*{Proof of \ref{B4.38}}\label{prf:B4.38}\phantomsection
\begin{proof}
Fix $H\in M_t$. By the assumed integral decomposition at $\theta(t)$, write $J_tH=\sum_{y\in S}\sigma(b_y)y$ with $S\subseteq\widetilde M_{\theta(t)}\cap V_0(\theta(t))$ finite and $b_y\in\mathcal O$.
For $x\in N$, choose the period representative $a_y\in\mathcal V_{w,s}$ for each $y\in S$. For the fixed column input $j_NE_ux$, let $z$ be its image under the reference column chain followed by evaluation and lowering. Let $e$ be the complex-linear identification from the period test space to the ambient row space. Let $\ell_z$ be the complex-linear functional on the period test space obtained by fixing $z$ in the first argument of the terminal period pairing. The row functional is $\ell_z\circ e^{-1}$. The twisted chain has the same formula with its own $z$ and $e$. Thus the period is complex-linear in the row in both cases. The bridge identity gives
\[
 \sigma(P_{u,t}(x,H))=\sum_{y\in S}\sigma(b_y)\sigma(a_y)
 =\sigma\left(\sum_{y\in S}b_ya_y\right).
\]
Injectivity of $\sigma$ gives $P_{u,t}(x,H)=\sum_{y\in S}b_ya_y\in\mathcal V_{w,s}$. Since $\phi\mapsto P_{u,t}(\phi,H)$ is $\mathcal O$-linear, the bound extends to $\operatorname{span}_{\mathcal O}N$.

At $w$, take $t=t_w$. Apply the first part with $N=\mathcal N_u$ and $s=s_{u,w}$. For $x\in\mathcal N_u$, we have $E_ux\in N_u$. Since $\theta(t_w)=\mu_w$, Theorem~\ref{B4.29} supplies the required period representatives. The span of $\mathcal N_u$ is $\mathcal N_u$.
\end{proof}

\section*{Proofs for Chapter~\ref{ch:pairing-after-division}: The pairing after division}

\subsection*{Proof of \ref{B4.43}}\label{prf:B4.43}\phantomsection
\begin{proof}
For $x\in(\iota\pi M)_{(w)}$, choose
$t\in\mathcal O\setminus\mathfrak p_w$ and $F\in M$ with
$tx=\iota\pi F$. Then $P(F)=\zeta\,t\,b(x)$. If $b(x)\ne0$, this is
nonzero, and $v_w(\zeta)=v_w(t)=0$ gives
$v_w(b(x))=v_w(P(F))\ge c$.

Take $x\in R$. The first assertion gives
$v_w(b(x))\ge0$ at every finite place when $b(x)\ne0$.
The principal fractional ideal $(b(x))$ therefore has nonnegative
exponent at every prime, so it is an integral ideal and
$b(x)\in\mathcal O$. The value zero also belongs to $\mathcal O$.
\end{proof}

\subsection*{Proof of \ref{B4.45}}\label{prf:B4.45}\phantomsection
\begin{proof}
Apply Theorem~\ref{B4.44} with $R=R^c$, $B=B^c_i$, $M_w=M_{t_w}$,
$\pi_w=A_w$, $\iota_w=\mathrm{id}$ and $P_w=P_{i,t_w}$.

Put $N=E_u^{-1}(N_u)$.
Proposition~\ref{B4.38} gives the bound on $N\times M_{t_w}$ with
$\sigma_w=s_{u,w}$ viewed in $\R$. By Lemma~\ref{B4.41}, $N^c_u$
is the shift of $N$ by $\tau_w=s_{u,w}=\lceil\sigma_w\rceil$.
The first assertion gives $B^c_u(N^c_u)\subset(R^c)^\vee$.
For $\phi\in N'_u$ and $r\in RR$, take $x\in N^c_u$ and $y\in R^c$
mapping to $\phi$ and $r$. Then $B_u(\phi)(r)=B^c_u(x)(y)\in\mathcal O$.
\end{proof}

\section*{Proofs for Chapter~\ref{ch:p4c6}: Slope change and the data for step 5}

\subsection*{Proof of \ref{B4.53}}\label{prf:B4.53}\phantomsection
\begin{proof}
Form \(P_u\) from the primes at most \(X''\) and the four exceptional primes.
Adjoin the prime divisors of \(M_u\). Outside \(P_u\) the raw gain
and loss vanish, so the gain after dropping primes vanishes as well.
The size datum gives \(F_j(w)=0\) for \(p_w>X''\) with
\(p_w\notin\{g_0,g_1,d\}\). Substitution in the net budget of
\ref{B4.05} gives \(S_{u,w}=0\).
By \ref{B4.51}, \ref{B4.52} and \(e_w\ge1\), the error is at most
\(2\sum_{p\in P_u}\log p\).

Put \(P(n)=\prod_{p\le n}p\). For \(k\ge1\), the primes in
\((k+1,2k+1]\) divide \(\binom{2k+1}{k}\le4^k\), since the binomial
sum \(2^{2k+1}\) contains two equal central terms. Thus
\(P(2k+1)\le P(k+1)4^k\). Since \(P(2k)=P(2k-1)\) for \(k\ge2\),
induction from \(n=0,1,2\) gives \(P(n)\le4^n\).
Take \(n=\lfloor X''\rfloor\), use \(\prod_{p\mid M_u}p\mid M_u\)
and include the four exceptional logarithms.
\end{proof}

\subsection*{Proof of \ref{B4.54a}}\label{prf:B4.54a}\phantomsection
\begin{proof}
Choose \(a\in K^\times\) with \(v(a)=E\) and put \(B'=B\cap a^{-1}\Theta\).
In an elementary-divisor basis of \(B\) let \(b_i\) be the exponents
of \(\Theta\). Take \(V\) spanned by the basis vectors with \(b_i>E\).
By \eqref{eq:p4c6-theta-cut} and \((\Theta_1)_V\subseteq\Theta_V\),
\begin{equation}\label{eq:p4c6-theta-remainder}
 \operatorname{length}(B/B')=\sum_i(b_i-E)_+
 =[B_V:\Theta_V]_v-E\dim V\le r .
\end{equation}

The operator \(H=\sum_{u\in U}\psi(u)^{-1}\rho(u)\) sends a stable
lattice into its intersection with \(W\) and acts by \(h\) on \(W\).
For \(g\notin U\), the specified \(u\) gives
\(\psi(u)H\rho(g)|_W=\psi(g^{-1}ug)H\rho(g)|_W\), so \(H\rho(g)|_W=0\).
For left coset representatives \(g_i\) this yields
\[
 H\rho(g_i^{-1}g_j)|_W=h\delta_{ij}\operatorname{id}_W .
\]
Thus the translates \(\rho(g_i)W\) are in direct sum and \(s\dim W\le n\).
Put \(Y=B_W/B'_W\). Summing translates and applying the operators
\(H\rho(g_i^{-1})\) give maps \(Y^s\to B/B'\to Y^s\) whose composite
is multiplication by \(h\). Both maps descend by stability of \(B,B'\).
Hence \(s\operatorname{length}(hY)\le r\). The quotient \(Y/hY\) is
killed by \(h\) and has at most \(\dim W\) generators. Since
\(B_Z/B'_Z\hookrightarrow Y\),
\begin{equation}\label{eq:p4c6-theta-eigenspace}
 s[B_Z:B'_Z]_v\le s[B_W:B'_W]_v\le r+n v(h).
\end{equation}

For \(z\in B'_Z\), one has \(az/\sigma\in(\Theta_K)_W\), so
\(haz/\sigma\in\Theta_0\cap V_0\) is fixed by \(e\).
Thus \((ha/\sigma)B'_Z\subseteq e\Theta_0\subseteq e\Theta_J\).
Scaling adds \(d_Z(E-v(\sigma)+v(h))\) to the index.
Apply \eqref{eq:p4c6-theta-eigenspace},
\eqref{eq:p4c6-theta-comparison} and \(X_3\subseteq A\).
\end{proof}

\subsection*{Proof of \ref{B4.55}}\label{prf:B4.55}\phantomsection
\begin{proof}
The shift is full by \ref{B4.40}. Apply \ref{B4.45} with
\(\tau_w\le\lceil\sigma_w\rceil\) for the inclusion and \ref{B4.50}
for the slope difference.

The positive-dimensional reference base quotient is spanned by its
spectral subspaces, so \(\mathcal I\ne\varnothing\). The coefficient
bound of \ref{B5.07a} and \eqref{eq:p4c6-height-inputs} give
\(\|K_i\|_v^2\le\exp(\sum_{j:i(j)=i}k_{j,v})\), with \(K_i=0\) when
its coefficient bound is zero. For \(B\ne0\), \ref{B5.07} gives
\(\log\|B\|_v^2\le t_v\), where
\[
 t_v=h_P+\sum_j k_{j,v}-\log\kappa_v
 -\min_{\alpha\in\mathcal I}\sum_j\log\lambda_{j,\alpha,v}.
\]
Apply the terminal and kernel table estimates
\eqref{eq:B5-terminal-table-bound} and \eqref{eq:B5-kernel-table-bound}
of \ref{B5.07a} to the ranks and weights in \(R\).
They give \(h_P\le H_P\) and \(\langle k_{j,v}\rangle_v\le H_{B,j}\).
By \eqref{eq:p4c6-selected-spectrum} and \eqref{eq:p4c6-height-inputs},
\[
 \min_{\alpha\in\mathcal I}\sum_j\log\lambda_{j,\alpha,v}
 \ge\sum_jl_{6,j}\ge\sum_jH_{X,j}.
\]
Averaging gives \(\bar t:=\langle t_v\rangle_v\le H-\bar\kappa\).
Take \(c_v=\exp(t_v+H-\bar\kappa-\bar t)\) for \(B\ne0\) and
\(c_v=\exp(H-\bar\kappa)\) for \(B=0\).

By \ref{B5.08x}, \(\mu_0\le B_*\). The row budgets give
\(0\le\delta\le1/3\), so passage at \(U=h_{\mathrm r}-\bar\kappa\) gives
\[
 \mu_C-\bar\kappa\le\mu_0+L_{\rm pass},\qquad
 L_{\rm pass}=3\delta(\mu_0+h_{\mathrm r}+28\lambda_{\mathrm r})
 +28\lambda_{\mathrm r}.
\]
By \eqref{eq:p4c6-row-budgets}, \(L_{\rm pass}\le
1+28\lambda_{\mathrm r}\le1+C_0m\lambda_0\), since \(m\ge1\) and
\(\lambda_0\ge0\). The same budgets give
\(\mu_0\le\sum_{j,\rho}\beta_j w_{j,\rho}\log\Omega+2F_*+E_b\).
Apply \ref{B5.08} with \(\widetilde\mu_C=\mu_C-\bar\kappa\).
\end{proof}

\subsection*{Proof of \ref{B4.55b}}\label{prf:B4.55b}\phantomsection
\begin{proof}
For the fixed scale write \(j\) for the arithmetic layer number in
the comparison, and let \(p\) be the residue characteristic of
\(\iota^{-1}y\). At \(g_0,g_1\), Lemma~\ref{B4.02a} identifies the rows as
\(\kappa^{\rm big}=a^+\) and \(\kappa^{\rm small}=-a^-\).
At each position use the auxiliary entries \(b_i,e_i\) of its weight
record, numbered from zero, and put \(t=t_{\rm ch}\).
The last index of \(a^+\) is \(k+j-2\), where the spectral
parameter is \(-b_{j-1}\). Substitution in the two rows gives
\[
 a^+_{\rm last}=e_{j-1}-b_{j-1}-(t+1)/2,\qquad
 a^-_0=b_j-e_{j-1}-(t+1)/2\quad(j<m).
\]
The gaps \(e_{j-1}-b_{j-1}\) and \(b_j-e_{j-1}\) are at least
\((t+3)/2\). The spectral parameter decreases by at least \(1\)
at each step of the \(a^+\) row, while its index term increases by \(1\).
The \(a^-\) increment is \(b_{j+c+1}-b_{j+c}-1\) from entry \(c\)
to entry \(c+1\), nonnegative by the gaps between auxiliary entries.
Thus \(a^+\) is nonincreasing and \(a^-\) is nondecreasing. Every
entry of either row is at least \(1\). After reindexing the embeddings,
the sum of absolute weights is \(\sum a^++\sum a^-\), and both base
terms are minus half this Fock degree. At \(d\) both equal
\(-(k-1)/2\) times the sum of the minima of the two weights.

For the other primes, use the position map
\[
 (\beta,t')\longmapsto e_\tau^{-1}\bigl(\tau_{\beta+1}(t')\bigr),
 \qquad \beta\in\{0,1\},
\]
where \(e_\tau\) is the bijection from embedding indices to the union of the two branches in \ref{B3.52}. Each branch is injective. Their images partition the embedding indices
under \(e_\tau^{-1}\).
For an inert or signal prime within the cutoff, \(p\) is coprime to
the conductor. By \ref{B3.46}, the second component of the Artin symbol of \(p\), expressed in the additive coordinates of the positions, gives the same Frobenius displacement \(h_p\) on both sides. Hence both sides have the same position sequence and period
\(n_p\). By \ref{B3.52}, corresponding positions have equal weights,
so the capacity vectors of \ref{B4.01} agree entry by entry.
The fixed-class sums with denominator \(p^2-1\) agree, as do the cyclic
orbit values and averages. Both entries vanish in the
remaining prime classes or above the cutoff outside the exceptional
primes. These cases exhaust the formula for \(F_j(w)\) in \ref{B4.02}, proving \eqref{eq:p4c6-base-tie}.

Enlarge both place sums to \(\mathrm{Tw}\cup\iota(S)\). At a place outside \(\mathrm{Tw}\), untwistedness gives \(F_{0,j,y}=0\) for every layer. At a place outside \(\iota(S)\), arithmetic support gives \(F_{a(j)}(\iota^{-1}y)=0\). Thus \eqref{eq:p4c6-base-tie} puts every place with nonzero base term in both sets. In particular the whole fibre of a type prime contributes whenever its Fock degree is positive. The place
bijection preserves residue norms and ramification indices, and
the weights in \ref{B5.02} are \(2\ell_we_w\), twice those in \(F_*\).
Sum \eqref{eq:p4c6-base-tie} using the layer bijection.
\end{proof}

\subsection*{Proof of \ref{B4.55a}}\label{prf:B4.55a}\phantomsection
\begin{proof}
Put \(l=\log M\). In the ambient frames the change of mean slope
from \(P_i\) lies in \([-2\log h_i,0]\) by inclusion and \ref{B4.46}.
The same interval bounds the change of maximum slope, since multiplication
by \(h_i\) sends submodules of \(P_i\) into \(h_iP_i\).
Frame scaling adds \(\log h_i\). Hence
\begin{equation}\label{eq:p4c6-frame-comparison}
 |\mu(R_i)-\mu(P_i)|\le l,\qquad
 |\mu_{\max}(R_i)-\mu_{\max}(P_i)|\le l .
\end{equation}

Write \(P_{i,s}\) for the cut in a nonzero ambient isotypic space.
Scalar endomorphisms give absolute irreducibility by scalar extension
of endomorphisms and semisimplicity. Over \(\overline K\) the eigenvalues
of \(g\in G\) are roots of unity. Their sum is the \(K\)-valued trace
\(\chi_s(g)\), hence belongs to \(\mathcal O\). The central projector
formula \cite[Section 2.6, Theorem 8]{serre-1977} has numerator
coefficients \(\dim(s)\chi_s(g^{-1})\) and denominator \(M\). Thus
\begin{equation}\label{eq:p4c6-type-cuts}
 MP_i\subseteq\bigoplus_s P_{i,s}\subseteq P_i ,
\end{equation}
with the sum orthogonal by unitarity.

On a copy of \(s\), the translates of the base character line span
the copy. Project onto the selected character space and choose
\(n_s\) translates whose projections form a basis. The resulting
projected translates \(A_i\) give the map \((x_i)_i\mapsto\sum_i A_ix_i\)
from \(n_s\) copies of the base cut space to the selected cut space.
It is an isomorphism on each copy of \(s\) and hence on the isotypic
component. The reverse map projects the inverse translates onto the
base character space. At infinity it is the adjoint, so it is also
an isomorphism. Each operator uses at most one central projector
and two subgroup averages. The subgroup characters have integral
coefficients and the subgroup orders divide \(M\). Multiplication by
\(M^3\) gives integral maps between \(P_{0,s}^{\oplus n_s}\) and
\(P_{1,s}\) in both directions.

Each \(A_i\) has norm at most one, and at every infinite place
\[
 \left\|\sum_{i=1}^{n_s}A_ix_i\right\|^2
 \le n_s\sum_{i=1}^{n_s}\|A_ix_i\|^2
 \le n_s\sum_{i=1}^{n_s}\|x_i\|^2 .
\]
The adjoint has the same operator norm. The regular representation
contains each irreducible with multiplicity its dimension
\cite[Section 2.4, Proposition 5]{serre-1977}, so
\(n_s\le\dim(s)\le\sqrt M\). Either integral map \(f_s\) therefore satisfies
\begin{equation}\label{eq:p4c6-transfer-norm}
 \log\|f_s\|_v^2\le 6l+\log n_s\le\tfrac{13}{2}l .
\end{equation}

For an injective lattice map with \(\log\|f\|_v^2\le c\) at every
infinite place, the image slope is at least the source slope minus
\(c\). In rank \(r\) its real volume grows by at most
\(\exp(n_Krc/2)\). A same-rank target has slope at least that of the
image by the index formula. Put \(r_s=\operatorname{rk}_{\mathcal O}P_{0,s}\),
\(d_s=r_s\mu(P_{0,s})\), \(R=\sum_s r_s>0\) and \(W=\sum_s n_sr_s\).
Use the forward norm bound when adding the degrees of the orthogonal
type cuts. Then apply \eqref{eq:p4c6-type-cuts} and \ref{B4.46}. This gives
\begin{equation}\label{eq:p4c6-transfer-degrees}
 R\mu(P_0)\le\sum_s d_s+2lR,\qquad
 W\mu(P_1)\ge\sum_s n_sd_s-\tfrac{17}{2}lW .
\end{equation}

An orthogonal sum \(\bigoplus_i Q_i\) has maximum slope
\(\max_i\mu_{\max}(Q_i)\) over its nonzero summands. Indeed, filter
a submodule \(L\) by \(L_j=L\cap\bigoplus_{i\le j}Q_i\). Projection
embeds \(L_j/L_{j-1}\) in \(Q_j\) with norm at most one for the
quotient metrics. Degree additivity and the slope inequality
\cite[Propositions 5.5(2) and 5.6]{fonseca-2019} bound the slope of
\(L\) by the maximum, which the summands also bound from below.
Applying this to the targets of the reverse maps and using
\eqref{eq:p4c6-transfer-norm} gives
\(\mu_{\max}(P_{1,s})\le\mu_{\max}(P_0)+13l/2\).
By \eqref{eq:p4c6-type-cuts} and \ref{B4.46},
\(\mu_{\max}(P_1)\le\mu_{\max}(P_0)+17l/2\). Thus
\eqref{eq:p4c6-frame-comparison} gives
\(\mu_{\max}(R_1)\le U+21l/2\le U+14\lambda\).

Set \(A=U+l\) and \(z_s=Ar_s-d_s\). By \eqref{eq:p4c6-frame-comparison}
and \(P_{0,s}\subseteq P_0\) these are nonnegative. By
\eqref{eq:p4c6-transfer-degrees}, \(\sum_s z_s/R\le A-\mu(P_0)+2l\).
The multiplicity bounds \eqref{eq:p4c6-multiplicities} give
\begin{equation}\label{eq:p4c6-weighted-deficit}
 \frac{\sum_s n_sz_s}{W}
 \le\frac{1+\delta}{1-\delta}\frac{\sum_s z_s}{R}
 \le(1+3\delta)\frac{\sum_s z_s}{R}.
\end{equation}
Substitute \(\sum_s n_sd_s=AW-\sum_s n_sz_s\) in
\eqref{eq:p4c6-transfer-degrees} and apply
\eqref{eq:p4c6-weighted-deficit} and \eqref{eq:p4c6-frame-comparison}.
This gives
\[
 -\mu(R_1)\le-\mu(R_0)+\tfrac{25}{2}l+
 3\delta\bigl(U-\mu(R_0)+4l\bigr).
\]
The slope changes sign under duality by duality of the metrized
determinant line \cite[Definition 5.4]{fonseca-2019}. The inequality \(0\le l\le\lambda\) gives the row dual bound.
\end{proof}

\subsection*{Proof of \ref{B4.56}}\label{prf:B4.56}\phantomsection
\begin{proof}
By \ref{B4.05}, \(s_{u,w}=0\) off \(S\). Thus the
product over all places in \ref{B4.40} is \(I_u\) of \ref{B4.39},
and the coordinate shift is \(I_u\mathcal N_u\). Its image under
\(E_u\) is \(N'_u=I_uN_u\). Cancel \(I_u\) using \ref{B4.48} to get
\(E_u(\mathcal N_u)=N_u\).

At each finite place \(w\), apply the integral decomposition of \ref{B4.37} to \(J_{t_w}F\) at the label \(\mu_w\). The period bound of \ref{B4.29} puts the value on each summand in \(\mathcal V_{w,s_{u,w}}\) under the coefficient embedding. This module is closed under integral linear combinations, so the period identity of \ref{B4.38} gives \(P_{u,t_w}(\phi,F)\in\mathcal V_{w,s_{u,w}}\) for \(\phi\in\mathcal N_u\). Apply \ref{B4.55}.
The pairing identity transports the coordinate inclusion to
\(B_u(N'_u)\subseteq RR^\vee\). If \(B_u\ne0\), the isometries give
\(h(B_u)\le\langle\log c_v\rangle_v=H_u-\bar\kappa=h_u\).
The dual row isometry gives \(\mu_C=\mu(RR^\vee)\). Substitute
\eqref{eq:p4c6-package-comparisons} in
\(\mu_C-\bar\kappa+H_u-2F_*\le R_{*,u}\).
\end{proof}

\section*{Proofs for Chapter~\ref{ch:lattice-after-division}: The lattice after division}

\subsection*{Proof of \ref{B5.05a}}\label{prf:B5.05a}\phantomsection
\begin{proof}
Let a diamond representative \(\delta\) fix the class of \(y\). Write \(y=\gamma y\delta\kappa\) with \(\gamma\) rational and \(\kappa\) in the level group. The diamonds normalize the level. Their marked multipliers lie in \(\F_T^\times\), so \((\delta\kappa)^H\) is in the level and \(\gamma^H\) is in the class stabilizer. That stabilizer is trivial by \ref{B8.54}. At either \(g_i\), the diamond is the identity locally and \(y^{-1}\gamma y\) is in the local level. Its reduction lies in a unipotent \(g_i\)-group and the congruence kernel is pro-\(g_i\). Thus the finite order of \(\gamma\) is a power of each of the two distinct primes \(g_0,g_1\). Hence \(\gamma=1\) and \(\delta=\kappa^{-1}\). Each marked multiplier of \(\delta\) is one. The diamond action is free.
\leanref{Formalism.\allowbreak{}L2.\allowbreak{}GlobalTheoryTheory.\allowbreak{}LocalLatticeBases.\allowbreak{}S0.\allowbreak{}fixed\_\allowbreak{}class\_\allowbreak{}mem\_\allowbreak{}level\_\allowbreak{}general}{Formalism.\allowbreak{}LocalToGlobal.\allowbreak{}GlobalTheory.\allowbreak{}GL6.\allowbreak{}DiagonalSlopes.\allowbreak{}EqOneOfCoprimeTorsionTwoPrimes}{572--621}
\leanref{Formalism.\allowbreak{}L2.\allowbreak{}GlobalTheoryTheory.\allowbreak{}LocalLatticeBases.\allowbreak{}S0.\allowbreak{}diamond\_\allowbreak{}free\_\allowbreak{}on\_\allowbreak{}classes}{Formalism.\allowbreak{}LocalToGlobal.\allowbreak{}GlobalTheory.\allowbreak{}GL6.\allowbreak{}DiagonalSlopes.\allowbreak{}EqOneOfCoprimeTorsionTwoPrimes}{671--699}

Diamond transport identifies the orthogonal coefficient slices on each orbit isometrically. The sum in one coordinate has kernel on \(\Z^H\) with basis \(e_i-e_H\), \(i<H\), and Gram determinant \(H\). The sum to \(\Z\) splits by \(1\mapsto e_H\). Tensor the \(f\) splittings with the coefficient lattice. The common kernel has rank fraction \((1-1/H)^f\) and its coordinate Gram determinant is \(H^{f(H-1)^{f-1}}\). The coefficient lattice contributes equally per rank to the two slopes. This gives the mean formula, while the induced metric gives the maximum-slope inequality.
\leanref{Formalism.\allowbreak{}L2.\allowbreak{}GlobalTheoryTheory.\allowbreak{}LocalLatticeBases.\allowbreak{}S0.\allowbreak{}free\_\allowbreak{}rank\_\allowbreak{}model\_\allowbreak{}of\_\allowbreak{}diamonds}{Formalism.\allowbreak{}LocalToGlobal.\allowbreak{}GlobalTheory.\allowbreak{}GL6.\allowbreak{}DiagonalSlopes.\allowbreak{}GroupActionMulSingleEqDilation}{287--301}

At \(w\mid T\) let \(S_0=(\Lambda_{u,0})_{(w)}\) and let \(V^{\rm ram}\) be the common diamond kernel. Put \(R=S_0\cap V^{\rm ram}\). For an active position \(v\), write \(\tau=\tau_w(v)\) and let \(\lambda=\lambda_{u,\tau}\) be the rank-\(k\) highest weight. Put \(b_v=-\lambda_k\ge0\). Let \(B_v\) use the frame of the kernel of the marked residue functional. The down operator \(D_v\) sums transports to that kernel. The up operator \(U_v\) sums from the upper neighbors formed by adjoining \(T^{-1}z\) with \(z\) of marking one. The extended marking recovers the original lattice as its kernel. Integral frame changes preserve the determinant-twisted Weyl lattice, and the level character is a unit by \ref{B8.09}(5). Hence \(D_vS_0\subseteq B_v\) and \(U_vB_v\subseteq S_0\).
\leanref{PI.\allowbreak{}IntegralAtBasePointLattice.\allowbreak{}integral\_\allowbreak{}at\_\allowbreak{}t0\_\allowbreak{}lattice}{Formalism.\allowbreak{}PeriodIntegrals.\allowbreak{}OrbitalBoundsConclusion.\allowbreak{}IntegralAtT0Lattice}{6--8}

The adapted frame change is \(\operatorname{diag}(T,1,\ldots,1)\). The multiplier on a torus weight \(\beta\) is \(T^{-\beta_1}\), with \(\beta_1\ge-b_v\). Thus \(T^{b_v}S_0\subseteq B_v\), with length \(e_w(b_v+\beta_1)\) per weight basis vector. By permutation symmetry the mean of \(\beta_1\) is \(k^{-1}\sum_i\lambda_i\). Consequently
\begin{equation}\label{eq:B5-level-weight-length}
 \operatorname{length}_w(B_v/T^{b_v}S_0)
 =e_wd_M\left(b_v+k^{-1}\sum_i\lambda_i\right).
\end{equation}
\leanref{Formalism.\allowbreak{}L2.\allowbreak{}GlobalTheoryTheory.\allowbreak{}LemB.\allowbreak{}cls\_\allowbreak{}len\_\allowbreak{}D}{Formalism.\allowbreak{}LocalToGlobal.\allowbreak{}GlobalTheory.\allowbreak{}CoordinateLengths.\allowbreak{}PairingAdjointness}{81--137}

In \(\F_T^k\) let \(e\ne0\) be the marked vector. Paths in \(D_vU_v\) with endpoint \(y\ne0\) are indexed by \(g(e)=g(y)=1\) for a linear functional \(g\). The count is \(T^{k-2}\) off the line of \(e\) and \(T^{k-1}\) at \(e\). It is zero at the other points of that line. With the first condition dropped the count is \(T^{k-1}\) for every \(y\ne0\). Transports at a fixed endpoint agree up to the level. If \(\mathrm{Nm}_v\) is the diamond sum and \(\mathrm{Tr}_v\) the unrestricted sum divided by \(T^{k-1}\), this gives
\begin{equation}\label{eq:B5-level-du}
 D_vU_v=T^{k-1}I-T^{k-2}\mathrm{Nm}_v+T^{k-2}\mathrm{Tr}_v .
\end{equation}
Rescaling the markings permutes the incidence sets, so both operators preserve \(V^{\rm ram}\). Since \(\mathrm{Tr}_v\mathrm{Nm}_v=H\mathrm{Tr}_v\), the trace vanishes on \(V^{\rm ram}\). There \(D_vU_v=T^{k-1}I\) and, by finite dimensionality, \(U_vD_v=T^{k-1}I\).
\leanref{Formalism.\allowbreak{}L2.\allowbreak{}GlobalTheoryTheory.\allowbreak{}UnramifiedTrace.\allowbreak{}du\_\allowbreak{}formula\_\allowbreak{}of\_\allowbreak{}fibre\_\allowbreak{}count}{Formalism.\allowbreak{}LocalToGlobal.\allowbreak{}GlobalTheory.\allowbreak{}UnramifiedTrace.\allowbreak{}DuFormulaOfFibreCount}{15--42}
\leanref{Formalism.\allowbreak{}L2.\allowbreak{}GlobalTheoryTheory.\allowbreak{}LocalLatticeBases.\allowbreak{}S0.\allowbreak{}hecke\_\allowbreak{}ramified\_\allowbreak{}zero\_\allowbreak{}of\_\allowbreak{}level\_\allowbreak{}zero}{Formalism.\allowbreak{}LocalToGlobal.\allowbreak{}GlobalTheory.\allowbreak{}GL6.\allowbreak{}DiagonalSlopes.\allowbreak{}NaturalTranslateAdelicFunction}{473--491}

Put \(R_B=B_v\cap V^{\rm ram}\). Then \(T^{k-1}R\subseteq U_vR_B\subseteq R\), and \(U_vR_B/T^{b_v}U_vR\) is an image of \(R_B/T^{b_v}R\). The latter injects into \(B_v/T^{b_v}S_0\). At stage zero the leading entry is at most \(a_{\max}=\max_{t\in TT}a_t\), and the other entries have absolute value at most \(W_{u,\tau}\) from \ref{B4.03}. The selected sizes give
\begin{equation}\label{eq:B5-level-sizes}
 a_{\max}\le2a,\qquad k\le N_m\le b',\qquad b_v\le2b',\qquad W_{u,\tau}\le3b'.
\end{equation}
The last bound uses \(2b'\) for the top-type entries and \(b'\) for the shift \(i(t_{\rm ch}+1)/2\). Add the bound \(e_w(k-1)d\) for \(R/U_vR_B\) to \eqref{eq:B5-level-weight-length}. It yields
\begin{equation}\label{eq:B5-level-one-step}
 \operatorname{length}_w(R/T^{b_v}U_vR)
 \le e_wd_M(a_{\max}/k+2W_{u,\tau}+k+b_v)
 \le9e_wd_M(a/k+b').
\end{equation}
\leanref{Formalism.\allowbreak{}MainNotions.\allowbreak{}product\_\allowbreak{}first\_\allowbreak{}zeroth\_\allowbreak{}parts.\allowbreak{}twist\_\allowbreak{}le}{Formalism.\allowbreak{}SharpRestriction.\allowbreak{}GlobalSchurBounds.\allowbreak{}FiniteModelOfAdmissibleTransports}{1706--1722}
\leanref{Formalism.\allowbreak{}MainNotions.\allowbreak{}J.\allowbreak{}Wt\_\allowbreak{}top\_\allowbreak{}le}{Formalism.\allowbreak{}SharpRestriction.\allowbreak{}AdmissibleParameterJoins.\allowbreak{}DIVfromGLo}{112--138}
\leanref{Formalism.\allowbreak{}MainNotions.\allowbreak{}J.\allowbreak{}fifthLevelSizes\_\allowbreak{}of\_\allowbreak{}weights}{Formalism.\allowbreak{}SharpRestriction.\allowbreak{}AdmissibleParameterJoins.\allowbreak{}DIVfromGLo}{211--248}

The congruence lattice is \((\prod_v T^{b_v}U_v)^mR\). For injective \(A,B\) preserving \(R\), the colengths add because \(AR/ABR\simeq R/BR\). The active \(v\) map injectively via \(\tau_w\) into the top positions, so there are at most \(n_T\) factors per round. At an unseen place the product is the identity. Thus the length is at most \(9e_wd_Mmn_T(a/k+b')\). The level size gives \(H>2f\), so \(d/d_M\ge1-f/H\ge1/2\). Multiply by \(2\ell_w/d\) and sum using \(\sum_{w\in F}\ell_we_w\le\log T\) from \ref{B4.52}.
\leanref{Formalism.\allowbreak{}MainNotions.\allowbreak{}J.\allowbreak{}card\_\allowbreak{}VV\_\allowbreak{}le}{Formalism.\allowbreak{}SharpRestriction.\allowbreak{}AdmissibleParameterJoins.\allowbreak{}DIVfromGLo}{185--202}
\end{proof}

\subsection*{Proof of \ref{B5.06}}\label{prf:B5.06}\phantomsection
\begin{proof}
Definition~\ref{B5.03} gives \(s_{uw}=\lfloor e_w\sg(u,w)\rfloor\). Thus \(R_u=\sum_{w\in S}(c_w/e_w)\{e_w\sg(u,w)\}\ge0\), where braces denote fractional part and \(c_w/e_w=2\ell_w>0\). The rank identity gives \(0\le\delta\le2f/T_u\). By \ref{B5.02} and \ref{B5.04},
\[
 \mu(N'_u)=\mu(M_{H,u})+J_u-R_u
 -\sum_{w\in S}c_w(\ell^{\rm col}_{uw}/d-\tau_w).
\]
Off \(S_T\) the parenthesis is at most \((\mathrm{io}_w+\delta\tau_w)/(1-\delta)\). The total sum is at most \(P_u+4\mathrm{IO}+4fb^2/T_u\). The mean formula of \ref{B5.05a} gives the lower bound. For the maximum slope, \ref{B5.05} bounds the increase from \(M_{H,u}\) by \(\sum_wc_w(\sigma'_w-\tau_w)=J_u-R_u\le J_u\).

For the selected family put \(X'=\exp(x^{\varepsilon/3})\) and let \(\tau_p\) be the common value over \(p\). A nonempty selection in \ref{B3.52} requires \(p<X'\). If \(N_{pz}\) counts top positions over a selected \(z\), its cut terms cost at most \(CN_{pz}a/p\) each, and its rounding unit costs at most \(N_{pz}a/p\) by \ref{B5.20a}. Use \(\sum_zN_{pz}\le n_T\) and the denominator allowance \(Cfk^2\log a/\log p\) to obtain
\begin{equation}\label{eq:B5-tau-prime}
 (\log p)\tau_p\le Cfk^2\log a+(2C+1)mn_Ta\frac{\log p}{p}.
\end{equation}
Rosser and Schoenfeld prove \(\sum_{p\le Y}(\log p)/p<\log Y\) for \(Y>1\) \cite[(3.24)]{rosser-schoenfeld-1962}. In particular, \(2\log Y\) bounds the sum. Group by rational prime using \ref{B4.52}. For \(Z=\exp(x^{\varepsilon/2})\), the bounds \(C,k,2\log a,X'\le Z,\ f\le2Z\) of \ref{B8.02} give \(2Cfk^2\log a\,\lceil X'\rceil\le Z^8\) eventually. Thus
\begin{equation}\label{eq:B5-tau-family}
 \mathcal T\le4(2C+1)mn_Ta\log X'+b^{8x^{-\varepsilon/2}}.
\end{equation}
\leanref{Formalism.\allowbreak{}MainNotions.\allowbreak{}G.\allowbreak{}TopChoiceBounds.\allowbreak{}tau\_\allowbreak{}sum\_\allowbreak{}family\_\allowbreak{}of\_\allowbreak{}mertens}{Formalism.\allowbreak{}SharpRestriction.\allowbreak{}GrowthBounds.\allowbreak{}TopChoiceBounds}{4503--4505}
The right side is \(O(x^8b)\) by \(m,n_T\le2x^3,\ \log X'\le x\), and is eventually at most \(b^2\).

The level sizes \(16f\le b,\ 4\le b,\ b^3\le T_u\) bound \(2f/T_u\) by \(1/2\) and \(4fb^2/T_u\) by \(1/4\). Using \(\log(T_u-1)\le2\sqrt{T_u}\) also bounds \(f\log(T_u-1)/(T_u-2)\) by \(1/4\). The level fibre has empty selection since \(X'\le T_u\). Use \ref{B5.05a} there and \ref{B5.20a} elsewhere for the colengths and containments. By \ref{B4.51} and \ref{B4.54},
\[
 0\le R_u=2\sum_{w\in S}\ell_w\{e_w\sg(u,w)\}
 \le 2(\log4+3+C_T+C_{\rm stab})X'' ,
\]
where \(C_T,C_{\rm stab}\) are the level and stabilizer constants of \ref{B4.54}. Thus \(c_1\ge\max\{36,2(\log4+3+C_T+C_{\rm stab})\}\) suffices.
\end{proof}

\subsection*{Proof of \ref{B5.07}}\label{prf:B5.07}\phantomsection
\begin{proof}
At \(\mu_0\) the comparison matrix is the identity in the image basis.
\leanref{Formalism.\allowbreak{}IndefR.\allowbreak{}Cx.\allowbreak{}CxData.\allowbreak{}rowMat\_\allowbreak{}base}{Formalism.\allowbreak{}LocalToGlobal.\allowbreak{}IndefiniteReal.\allowbreak{}CxLemmas}{352--354}

At each infinite place \(v\), the comparison yields \eqref{eq:B5-chain-metric}. For \(v\ne v_0\), let \(P_p\) be the projection onto \(Q_p\) along the other rational spectral summands, extended at \(v\). Distinct types have disjoint joint Hecke spectra. The spherical operators commute with one another and with their adjoints in the quotient metric, so these summands are orthogonal. In the reference coordinates the constructed row metric is
\[
 \langle\iota z,\iota y\rangle_v=\sum_{p\in I}\left(\kappa_v\prod_{j=1}^m\lambda_{jpv}\right)\langle P_pz,P_py\rangle_v .
\]
The weights are positive, and the projections sum to the identity. For \(z\in Q_p\) the only surviving term is the one indexed by \(p\), with \(\langle z,P_py\rangle_v=\langle z,y\rangle_v\). This proves \eqref{eq:B5-chain-metric} also when a layer contributes the empty-part weight \(e^{L_{6,j}}\) of \ref{B5.07g}: that very weight occurs in the displayed definition of the row metric.
\leanref{Formalism.\allowbreak{}L2.\allowbreak{}DIV.\allowbreak{}weightMet}{Formalism.\allowbreak{}LocalToGlobal.\allowbreak{}Divisibility.\allowbreak{}Frames}{49--53}
\leanref{Formalism.\allowbreak{}IndefR.\allowbreak{}Cx.\allowbreak{}ComplexDataConditionsBaseLevel.\allowbreak{}partsOrth}{Formalism.\allowbreak{}LocalToGlobal.\allowbreak{}IndefiniteReal.\allowbreak{}PackageFromPrimitiveStatements}{66--73}
 At \(v_0\), let \(P_b\) be the level subgroup for the side of the row under consideration. The row Gram matrix is the Petersson Gram matrix divided by \(|P_b|\), and its undivided Petersson ratio on \(Q_p\) is \(|P_b|c_{pv_0}\).
\leanref{Formalism.\allowbreak{}IndefR.\allowbreak{}Cx.\allowbreak{}rowXE\_\allowbreak{}part}{Formalism.\allowbreak{}LocalToGlobal.\allowbreak{}IndefiniteReal.\allowbreak{}CxLemmas}{280--316}
\leanref{Formalism.\allowbreak{}IndefR.\allowbreak{}Cx.\allowbreak{}CxData.\allowbreak{}scal\_\allowbreak{}rowsR}{Formalism.\allowbreak{}LocalToGlobal.\allowbreak{}IndefiniteReal.\allowbreak{}CxCoreRG}{60--88}
\leanref{Formalism.\allowbreak{}IndefR.\allowbreak{}Cx.\allowbreak{}CxData.\allowbreak{}colsOfR}{Formalism.\allowbreak{}LocalToGlobal.\allowbreak{}IndefiniteReal.\allowbreak{}CxCoreRG}{151--170}

Equation \eqref{eq:B5-chain-metric} makes \(Q_p\) a \(c_{pv}\)-eigenspace of \(\iota^*\iota\). Thus the squared singular values of \(\iota^\vee\) are at least \(a_v=\min_p c_{pv}\). The dual \(q^\vee\) is isometric, so
\begin{equation}\label{eq:B5-chain-norm}
 a_v\|B_u\|_v^2\le\|(\iota q)^\vee B_u\|_v^2
 \le\|P_u^{\rm term}\|_v^2\prod_i\|K_i\|_v^2 .
\end{equation}
Take logarithms and average. For the selected realization apply \ref{B5.07a} and \ref{B5.07b}. The constructed point assigns to each word the terminal allowance plus the sum of the kernel allowances, less the sum of the spectral lower bounds and the embedding mean \([K:\Q]^{-1}\sum_\sigma\log\kappa_{[\sigma]}\). The family takes this assignment from the point at every index above its cutoff. Since each infinite place has exactly two embeddings, this last term equals \(\langle\log\kappa_v\rangle_\infty\).
\leanref{GlobalTheory.\allowbreak{}EX.\allowbreak{}OB.\allowbreak{}LedPart.\allowbreak{}ht}{Development.\allowbreak{}AbelianVarietyTheory.\allowbreak{}Definitions.\allowbreak{}GlobalTwisting.\allowbreak{}SetUnipotentTranslate}{379--380}
\leanref{GlobalTheory.\allowbreak{}EX.\allowbreak{}OB.\allowbreak{}LedPart.\allowbreak{}pt}{Development.\allowbreak{}AbelianVarietyTheory.\allowbreak{}Definitions.\allowbreak{}GlobalTwisting.\allowbreak{}SetUnipotentTranslate}{381--402}
\leanref{GlobalTheory.\allowbreak{}EX.\allowbreak{}famOfKits}{Development.\allowbreak{}AbelianVarietyTheory.\allowbreak{}Proofs.\allowbreak{}GlobalInterpolation.\allowbreak{}ExistsFamPOfKits}{56--158}
\end{proof}

\subsection*{Proof of \ref{B5.07c}}\label{prf:B5.07c}\phantomsection
\begin{proof}
Since \(A_i\le D_{0,\tau}\), one has \(\sigma_i\in I_i\). For \(v\in I_i\), the number \(v-i\) lies strictly below \(\sigma_j-j\) when \(j<i\), and strictly above it when \(j>i\). Excluding \(\sigma_i\) ensures \(c_\tau\ne0\).

Let \(Q\) extract the coefficient of \(\prod_iX_{N,i}^{B_i}\) with last-column degree zero and project to the \(\sigma\)-module. The unscaled map is \(P_\tau=c_\tau\kappa_\tau Q\). In \(\xi\), the required degree of \(X_{N,n}\) forces the last-row terms of all \(Z_n\) to use column \(n\). Descend in \(i\) to get \(Q\xi=h_\sigma\). Row repetition makes \(\xi\) invariant under the upper unipotents of \(\GL_n\), and its row degrees \(\sum_{j\ge i}A_j+\sum_{j>i}B_j\) give \(\sigma_i\) after restoring the shift. By multiplicity-one interlacing \cite[\S8.1]{goodman-wallach-2009}, \(\xi\) spans the highest line of the \(\sigma\)-summand. In unit highest-vector frames \(\|\xi\|^2=\kappa_\tau\), so adjunction and equivariance give \(Q^*h_\sigma=\xi/\kappa_\tau\) and \(P_\tau P_\tau^*=c_\tau^2\kappa_\tau I\). Tensoring yields
\begin{equation}\label{eq:B5-branch-adjoint}
 P_tP_t^*=\prod_\tau(t_\tau\bar t_\tau)^{-n_\tau}c_\tau^2\kappa_\tau I,
 \qquad P_t=\left(\prod_\tau t_\tau^{-n_\tau}\right)\bigotimes_\tau P_\tau .
\end{equation}
The adjoint image attains the squared norm. At a coefficient embedding \(s_0\), \(t_\tau=s_0\tau(\bar\varpi)\) for the generator \(\varpi\bar\varpi=T_u\) of \ref{B4.03}. CM conjugation therefore gives the log squared norm
\begin{equation}\label{eq:B5-branch-log}
 -\sum_\tau(\lambda_{\tau,1}-D_{0,\tau})\log T_u+
 \sum_\tau(2\log|c_\tau|+\log\kappa_\tau).
\end{equation}
\leanref{GlobalTheory.\allowbreak{}BranchNorm.\allowbreak{}br\_\allowbreak{}tr}{Development.\allowbreak{}AbelianVarietyTheory.\allowbreak{}Proofs.\allowbreak{}DivisorProofs.\allowbreak{}LastMapNorm}{192--213}

For a class with stabilizer order \(h_c\), the contribution is \(h_c^{-1}b_\sigma(P_t\psi_c,F_c)\). The bilinear form is \(b_\sigma(z,y)=z^{\mathsf T}G_\sigma y\) for a real symmetric positive Gram matrix. Cauchy--Schwarz on \(\bar z,y\) bounds it by \(\|z\|\|y\|\). Apply \eqref{eq:B5-branch-adjoint} and Cauchy--Schwarz in the class sum with weights \(h_c^{-1}\).
\end{proof}

\subsection*{Proof of \ref{B5.07d}}\label{prf:B5.07d}\phantomsection
\begin{proof}
Write \(J_V,J_W\) for the Gram matrices and \(X^*=\bar X^{\,t}\). In the adelic model on \(r\)-by-\(n\) matrices the alternating form and the Heisenberg operators are
\[
 [X,Y]=\operatorname{Tr}_{E/F}
  \bigl(\delta\,\operatorname{tr}(J_V^{-1}Y^*J_WX)\bigr),
 \qquad
 D(Y)f(X)=\psi([X,Y]/2)f(X+Y).
\]
Here \(\bar\delta=-\delta\ne0\) is the scalar in the Weil datum and \(\psi\) is its standard additive character of \(\mathbb A_F/F\). Evaluation at zero in this model is the theta functional. For a rational matrix \(Y\), the rational transformation law of a theta vector gives
\[
 D(Y)f(X)=\psi([X,Y])f(X).
\]
Write \(\omega\) for the Weil action. It satisfies
\(\omega(h,g)D(Y)=D(gYh^{-1})\omega(h,g)\).

Choose a nonzero rational isotropic column \(\xi\) in the first nonzero member of the flag defining the parabolic. Put
\[
 z(t)=1+t\xi\xi^*J_W,\qquad \bar t=-t,
 \qquad Z_\xi=\{z(t)\}.
\]
Since \(\xi^*J_W\xi=0\), multiplication in \(Z_\xi\) adds the parameters. Its rational quotient is a copy of \(\mathbb A_F/F\). For a theta vector \(f\) and an adelic row \(a\) of length \(n\), put
\[
 C_f(a)=\int_{[Z_\xi]}(\omega(1,z)f)(\xi a)\,dz ,
 \qquad [Z_\xi]=Z_\xi(F)\backslash Z_\xi(\mathbb A_F).
\]
Use Haar probability on this compact quotient. The integrand is periodic in \(a\) modulo \(E^n\), since the alternating form vanishes on pairs \(\xi a,\xi b\). It descends in \(z\) by rational invariance of theta evaluation and \(z\xi=\xi\). These functions are continuous and bounded on the two compact quotients. We shall prove \(C_f(a)=C_f(0)\).

Let \(\mathcal A(f)=\int_{\mathbb A_E^n/E^n} f(\xi a)\,da\), with Haar probability. Translation of \(a\) gives \(\mathcal A(D(\xi b)f)=\mathcal A(f)\). The characters
\(a\mapsto\psi([\xi a,Y])\), for rational \(Y\), exhaust the characters of this quotient. Indeed \(Y\mapsto m=\xi^*J_WY\) is onto \(E^n\), and the trace pairing with \(\delta J_V^{-1}\) identifies these rows with the dual of \(\mathbb A_E^n/E^n\). Thus, up to the choice of sign of the frequency, the Fourier coefficients of \(C_f\) are
\[
 \int_{[Z_\xi]}\mathcal A(D(Y)\omega(1,z)f)\,dz .
\]

Suppose \(m\ne0\). The element \(N=mJ_V^{-1}m^*\) belongs to \(F^\times\), since \(J_V^{-1}\) is positive definite at a real place. Choose \(s\in\mathbb A_F\) supported at a split finite place such that \(\psi(\delta^2Ns)\ne1\). This is possible because the local component of the standard character is nontrivial. Set \(t=\delta s\) and \(b=tm\). Then
\[
 z(t)Y=Y+\xi b,\qquad [\xi b,Y]/2
       =\tfrac12\operatorname{Tr}_{E/F}(\delta tN)=\delta^2Ns .
\]

We give the local action on \(\mathcal A\). In the split Schr\"odinger model, \(z(t)\) acts by the linear substitution \(X\mapsto z(t)^{-1}X\) on one half of the matrix space, with a positive modulus factor \(c_t\). The determinant character is one since \(\det z(t)=1\). The line operators from the opposite half act by the characters of the row \(\xi^*J_WX\), using the perfect pairing with the coefficient space. A functional invariant under these operators depends only on the restriction of a local test to \(\xi^*J_WX=0\). To see this, take a locally constant compactly supported test which vanishes on that subspace. Its support has a finite partition into compact open sets on each of which some such character has constant value different from one. Invariance kills the restriction to each set. On the subspace itself \(z(t)^{-1}X=X-t\xi(\xi^*J_WX)=X\). Applying this argument to the local test with the other factors fixed gives
\(\mathcal A(\omega(1,z(t))f)=c_t\mathcal A(f)\), with \(c_t>0\).

The Heisenberg multiplication law now gives, for every theta vector \(f'\),
\[
 \mathcal A(D(Y)\omega(1,z(t))f')
   =\psi([\xi b,Y]/2)c_t\,\mathcal A(D(Y)f').
\]
Indeed one commutes \(\omega(1,z(t))\) across \(D(Y)\), writes
\(D(\xi b)D(Y)=\psi([\xi b,Y]/2)D(Y+\xi b)\), and uses the invariance of \(\mathcal A\) under \(D(\xi b)\). The multiplier differs from one. Its absolute value is \(c_t\), and if \(c_t=1\) its phase is \(\psi(\delta^2Ns)\ne1\). Translation of \(z\) by \(z(t)\) leaves the integral defining the Fourier coefficient unchanged. The coefficient therefore vanishes. When \(m=0\) the frequency is trivial. Fourier uniqueness on the compact additive quotient and continuity prove the asserted constancy. Since \(z\xi=\xi\), covariance also gives
\[
 \int_{[Z_\xi]}(\omega(1,z)D(\xi a)f)(0)\,dz=C_f(0)
 \qquad\text{for every }a .
\]
\leanref{GlobalTheory.\allowbreak{}TransvConst.\allowbreak{}exists\_\allowbreak{}transv}{Development.\allowbreak{}AbelianVarietyTheory.\allowbreak{}Proofs.\allowbreak{}IwasawaDecomposition.\allowbreak{}ConstantTermHeisOAlong}{113--159}
\leanref{GlobalTheory.\allowbreak{}SplitTransvection.\allowbreak{}transvection\_\allowbreak{}pos}{Development.\allowbreak{}AbelianVarietyTheory.\allowbreak{}Proofs.\allowbreak{}IwasawaDecomposition.\allowbreak{}SplitPlaceTransvectionActsPositivelyInvariantFor}{251--286}
\leanref{GlobalTheory.\allowbreak{}TransvConst.\allowbreak{}exists\_\allowbreak{}eigen}{Development.\allowbreak{}AbelianVarietyTheory.\allowbreak{}Proofs.\allowbreak{}IwasawaDecomposition.\allowbreak{}ConstantTermHeisOAlong}{229--242}
\leanref{GlobalTheory.\allowbreak{}TransvConst.\allowbreak{}constantTerm\_\allowbreak{}heisO\_\allowbreak{}along}{Development.\allowbreak{}AbelianVarietyTheory.\allowbreak{}Proofs.\allowbreak{}IwasawaDecomposition.\allowbreak{}ConstantTermHeisOAlong}{245--346}

Fix a translate \((h,g)\) of the kernel. In the marked split frame take the local column of \(g^{-1}\xi\) and divide by a coordinate of maximal absolute value. The resulting column \(v\) is integral and \(\bar v\ne0\) in the residue field \(k_w\). Let \(A\) be the first \(r\) rows of the local test variable. On the support, \(\bar A\bar v\ne0\). Fix a nontrivial additive character \(\psi_k\) of \(k_w\). For each \(\ell\in k_w^{r\,\vee}\) take the Heisenberg modulation with multiplier
\(\psi_k(\ell(\bar A\bar v))\). These are line operators: lift \(\ell\) in the first \(r\) dual coordinates, extend it by zero, and scale by the inverse of the coordinate used to make \(v\) primitive and by the conductor of the residue pairing. The split trace pairing realizes precisely these multipliers. Their sum on the test is zero, since the sum of a nontrivial character on \(k_w^{r\,\vee}\) is zero. Transporting the operators by \((h,g)\) gives \(D(\xi a_\ell)\) on the translated vector, with the coefficient rows transported by \(h^{-1}\). These finite-place operators commute with the prescribed real polynomial operators. Taking the transvection integral of their sum gives \(|k_w|^r C_f(0)=0\) for this translated kernel.
\leanref{GlobalTheory.\allowbreak{}LevelSupport.\allowbreak{}exists\_\allowbreak{}sum\_\allowbreak{}zero}{Development.\allowbreak{}AbelianVarietyTheory.\allowbreak{}Proofs.\allowbreak{}DivisorProofs.\allowbreak{}LevelFockVectorHeisenbergAverageIsotropicEqZero}{45--225}

The definite quotient is compact. We may integrate the vanishing transvection integral over it by Fubini. Finally, \(Z_\xi\) belongs to the unipotent radical \(N\) of the parabolic and centralizes it: each \(u\in N\) fixes \(\xi\), so \(u(\xi\xi^*J_W)u^{-1}=\xi\xi^*J_W\). Average the constant term on \([N]\) also over \([Z_\xi]\). Translation invariance on \([N]\) gives the original constant term. Integrating first over \([Z_\xi]\) gives zero at every remaining translate. This proves cuspidality.
\end{proof}

\subsection*{Proof of \ref{B5.07e}}\label{prf:B5.07e}\phantomsection
\begin{proof}
Write \(Q=Q_w\). Let \(L\) be the lattice in the symplectic space of the dual pair and put \(I=[gL:gL\cap L]\). In the lattice model the intertwiner on the vacuum is the sum of the Heisenberg translates indexed by \(gL/(gL\cap L)\), divided by the signed square root of \(I\) prescribed by the splitting. In its pairing with the original vacuum only the zero coset survives. Indeed, for a representative \(y\notin L\), self-duality supplies \(a\in L\) whose commutator character with \(y\) differs from one. Moving the lattice operator for \(a\) across the pairing multiplies that summand by this character and leaves the vacuum pairing unchanged. The zero coset contributes the original squared norm. Thus division by the value at the identity removes exactly that norm.
\leanref{GlobalTheory.\allowbreak{}LatCoef.\allowbreak{}coef\_\allowbreak{}of\_\allowbreak{}local\_\allowbreak{}vacuum}{Development.\allowbreak{}AbelianVarietyTheory.\allowbreak{}Proofs.\allowbreak{}TwistedOperators.\allowbreak{}WeilLatticeIntertwinerCoefficientAtPlace}{104--152}

At a split place choose one prime \(P\mid w\) and use its \(\GL_r\) chart. For Cartan exponents \(d_1,\ldots,d_r\), the two halves of \(L\) have moved indices \(Q^{n\sum_i\max(d_i,0)}\) and \(Q^{n\sum_i\max(-d_i,0)}\). The splitting uses the positive square root in this chart and multiplies by the local determinant character. At an inert place the exponents occur in opposite pairs, with central zero in odd rank. Put \(D=\sum_i\max(d_i,0)\). There \(I=(Q^2)^{nD}\). The lattice splitting on a conjugation-stable simple residue quotient is the negative square root of its order, hence takes \(Q^2\) to \(-Q\). This is also the sign required by \(\chi_P=(-1)^n\) on the uniformizer. Consequently the two coefficients on a Cartan representative \(a_d\) are
\[
 \phi_w^\circ(a_d)=
 \begin{cases}
   \chi_P^{\,\sum_i d_i}Q^{-n\sum_i|d_i|/2},&w\text{ split},\\
   ((-1)^n Q^{-n})^{D},&w\text{ inert}.
 \end{cases}
\]
Both have value one on \(K_w\).
\leanref{GlobalTheory.\allowbreak{}LamInert.\allowbreak{}lam\_\allowbreak{}qV\_\allowbreak{}inert}{Development.\allowbreak{}AbelianVarietyTheory.\allowbreak{}Proofs.\allowbreak{}TwistedOperators.\allowbreak{}IndexNormalizationInertNormEqNeg}{16--49}

We compare with the spherical doubling section normalized to one on the doubled hyperspecial. In the split chart its restriction to the graph has radial variable \(X=Q^{-s-r/2}\) raised to \(\sum_i|d_i|\), with the unramified determinant factor. In the inert chart the determinant of the change from the diagonal lattice to the integral graph frame has valuation \(D\): in a hyperbolic Cartan pair of exponents \(d,-d\) its integral graph basis contributes \(|d|\), and a zero exponent contributes zero. The unramified inducing law on this frame gives \((\chi_P Q^{-(2s+r)})^D\). Taking \(s=t/2\) gives the two preceding coefficients. This checks both the value at the identity and the restriction on every double coset.
\leanref{GlobalTheory.\allowbreak{}UGraph.\allowbreak{}sphericalDoublingSection\_\allowbreak{}graph\_\allowbreak{}diagonal\_\allowbreak{}cartan}{Development.\allowbreak{}AbelianVarietyTheory.\allowbreak{}Proofs.\allowbreak{}TwistedOperators.\allowbreak{}SphericalDoublingSectionGraphDiagonalEq}{156--205}
\leanref{GlobalTheory.\allowbreak{}ZetaGL.\allowbreak{}zeta\_\allowbreak{}sum\_\allowbreak{}GL}{Development.\allowbreak{}AbelianVarietyTheory.\allowbreak{}Proofs.\allowbreak{}TwistedOperators.\allowbreak{}UnramifiedZetaSumGLCosets}{58--272}

Apply the unramified doubling calculation to the conjugate vacuum and the normalized zonal spherical function of \(F_1\). The extension is split or unramified quadratic, both splitting characters are unramified off the marked places, the spherical section has the normalization just computed, and the two spherical vectors pair to one. With the inverse twist its denominator is
\[
 \prod_{a=0}^{r-1}L_w(2s+r-a,\chi_V^{-1}|_{F_w^\times}\eta_w^a),
\]
where \(\eta_w\) is the quadratic character of \(E/F\). Since \(\chi_V|_{F_w^\times}=\eta_w^n\), the substitution \(h=t+r-a=n-a\) gives \(d_w\), including the sign at every inert place. The numerator is the standard base-change factor at \(s+1/2\). For both split and inert unramified places this is the identity in \cite[\S4.2.1, pp.~48--49, display immediately before (53)]{eischen-harris-li-skinner-2020}; their (53) defines the denominator. Their local rank is our \(r\), their unramified inducing character is \(\chi_V^{-1}\), and the section and vector pairing have value one in the normalizations above. In split coordinates the numerator uses the two multisets \(\chi_P^{-1}B_P\) and \(\chi_P B_P^{-1}\). In inert coordinates it uses \((-1)^n B_P\) and \(NP=Q^2\).

On each double coset the sum of \(R(g)F_1\) over right \(K_w\)-cosets equals the Hecke eigenvalue times \(F_1\). Applying \(\ell\) gives the same sum as for the normalized zonal function times \(\ell(F_1)\), even when \(\ell(F_1)=0\). Conjugate the resulting identity. At a split place unitarity gives \(\overline{B_P}=B_P^{-1}=B_{\bar P}\) as multisets, and \(\chi_{\bar P}=\chi_P^{-1}\). At an inert place \(\overline{B_P}=B_P\). Thus conjugation preserves the displayed local product, giving \eqref{eq:B5-spherical-cosets}.

To justify evaluation at \(s=t/2\), first use the trivial spherical representation and trivial inducing character. Group the positive Cartan terms by \(\sum_i|d_i|\) in the split case and by \(2D\) in the inert case. This gives a power series in \(Q^{-s}\) whose coefficients are nonnegative. Near zero the spherical integral converges and the cited formula identifies its Taylor series with the rational Euler quotient. For the trivial representation the standard Satake parameters have modulus at most \((NP)^{(r-1)/2}\). Thus every possible pole of the quotient satisfies \(\Re s\le(r-2)/2\); the reciprocals of the normalizing factors are polynomials. The Taylor series consequently converges on a disk containing \(Q^{-t/2}\) in its interior. Its value there is the sum of the absolute vacuum weights. It dominates the series with a unitary twist and a normalized unitary spherical function throughout a neighborhood of that point. The identity on the initial half-plane therefore extends to \(s=t/2\) with the integral represented by its convergent coset sum. In the split case its Satake exponents have absolute value at most \((r-1)/2\) in base \(Q\). In the inert case the base-change bound is the same in base \(Q^2\). Since \(t\ge r+5\), every Euler argument in the resulting denominator is at most \(Q^{-2}\) in modulus. There are at most \(2r\) such arguments and \(r\) numerator factors of absolute value at most \(1+Q^{-2}\). Hence
\[
 1\le\sum_{gK_w}|\phi_w^\circ(g)|
       \le (1+Q^{-2})^r(1-Q^{-2})^{-2r}.
\]
Boundedness of \(\ell(R(g)F_1)\) justifies the double-coset summation. The excess of this majorant over one is \(O_r(Q^{-2})\), summable over the places of \(F\).
\leanref{GlobalTheory.\allowbreak{}ZetaSumU.\allowbreak{}hasSum\_\allowbreak{}zetaSum\_\allowbreak{}inert}{Development.\allowbreak{}AbelianVarietyTheory.\allowbreak{}Proofs.\allowbreak{}TwistedOperators.\allowbreak{}HasSumUnramifiedZetaSumUnitaryInert}{139--221}
\leanref{GlobalTheory.\allowbreak{}FinZeta.\allowbreak{}AbsConv.\allowbreak{}local\_\allowbreak{}place}{Development.\allowbreak{}AbelianVarietyTheory.\allowbreak{}Proofs.\allowbreak{}TwistedOperators.\allowbreak{}LocalZetaSumsAbsolutelyConvergentUniform}{81--116}
\end{proof}

\subsection*{Proof of \ref{B5.07y}}\label{prf:B5.07y}\phantomsection
\begin{proof}
Write \(F=E^+\), \(f=[F:\Q]\), \(G=\U(W)\) and \(H=\U(V)\), with \(\dim_E W=r\) and \(\dim_E V=n=r+t\). For an \(F\)-group \(J\) put \([J]=J(F)\backslash J(\mathbb A_F)\). Write \(S_f\) for the places below the marked primes. They are split, with distinct restrictions to \(F\). Use the following product Haar measures, taking counting measure on rational points in forming the quotients.

On \(G(F_w)\) give the lattice stabilizer volume one for \(w\notin S_f\). For \(w\in S_f\) give
\[
 K_w^1=\ker\bigl(\GL_r(\mathcal O_w)\longrightarrow\GL_r(k_w)\bigr)
\]
volume one. Thus \(\GL_r(\mathcal O_w)\) has volume \(|\GL_r(k_w)|\) at a marked place. At a real label \(\rho\) use \(dg_\rho=V_r\,dg_{\rho,\mathrm{cd}}\), where \(dg_{\rho,\mathrm{cd}}\) has the invariant density transferred from Haar probability on the compact dual \(\U(r)\), and
\[
                 V_r=\prod_{a=0}^{r-1}\frac{2\pi^{a+1}}{a!}.
\]
Indeed the invariant density in independent skew-Hermitian coordinates has a factor two for each complex off-diagonal coordinate. The specified Lie algebra measure divides by these factors. Projection of the compact group to its last column then has the ordinary sphere density and fiber the same coordinate measure on \(\U(r-1)\). Multiplication of the sphere areas gives \(V_r\). On \(H\) use probability at infinity and give the finite definite level volume one. On the additive symplectic space use rational covolume one, with the vacuum normalized to unit norm. At a marked split place the additive measure in the integral matrix chart gives \(\operatorname{Mat}_{n,r}(\mathcal O_w)\) volume one. These conventions fix the quotient integrals and the metrics defining the adjoint.
\leanref{Formalism.\allowbreak{}IndefR.\allowbreak{}Lev.\allowbreak{}IsMeasLevW}{Formalism.\allowbreak{}LocalToGlobal.\allowbreak{}IndefiniteReal.\allowbreak{}SubgroupCharExtensionOneApply}{122--132}
\leanref{Formalism.\allowbreak{}IndefR.\allowbreak{}Links.\allowbreak{}RealDat.\allowbreak{}models\_\allowbreak{}νW}{Formalism.\allowbreak{}LocalToGlobal.\allowbreak{}IndefiniteReal.\allowbreak{}LinksChar}{223--225}
\leanref{ReciprocityTheory.\allowbreak{}ChartRestriction.\allowbreak{}DS.\allowbreak{}volU}{ReciprocityTheory.\allowbreak{}ReductiveGroups.\allowbreak{}DiscreteSeries.\allowbreak{}FormalDegree.\allowbreak{}VolUFactProdMulLe}{14--15}
\leanref{GlobalTheory.\allowbreak{}LowRow.\allowbreak{}CompactDualScale.\allowbreak{}isMinf\_\allowbreak{}eq\_\allowbreak{}volU\_\allowbreak{}smul\_\allowbreak{}cd}{Development.\allowbreak{}AbelianVarietyTheory.\allowbreak{}Proofs.\allowbreak{}DivisorProofs.\allowbreak{}StandardMeasureUnitaryPQEqVolumeSmulCompactDual}{29--52}

Use inner products conjugate linear in the first variable. For a definite class form \(x\), extend its values equivariantly to \(x(h)\). Let \(\Phi(v,b)\) be the reference Schwartz vector with lowest-type argument \(v\) and definite coefficient \(b\) of \ref{B8.55}. Define the full lift on lowest-type vectors by
\begin{equation}\label{eq:B5-full-lift}
 (Tx)(v)(g)=\int_{[H]}\Theta(h,g;\Phi(v,x(h)))\,dh .
\end{equation}
The marked test has invertible upper block, so Lemma~\ref{B5.07d} gives cuspidality at every translate. The level is principal and its character is trivial. At \(\rho\), put \(p=p_\rho,\ q=q_\rho\). In the \(r\)-by-\(n\) Fock array let \(M_s^+\) be the minor in the first \(s\) rows and columns, and \(M_s^-\) the minor in rows \(p+1,\ldots,p+s\) and the last \(s\) columns. On the paired unit extremal vectors of \ref{B8.55} the polynomial is
\[
 \Delta_\rho=
 \prod_{s=1}^{p}(M_s^+)^{a^+_s-a^+_{s+1}}
 \prod_{s=1}^{q}(M_s^-)^{a^-_s-a^-_{s+1}},
\]
with parts numbered from one and a zero appended. Each mixed lowering operator \(\sum_c\partial_{ac}\partial_{bc}\), \(a\le p<b\), kills it: \(p+q=r\le n\), so the two sets of columns are disjoint. Compact equivariance gives the lowering equations on the lowest type. These equations give the holomorphic lowest-weight module in the cuspidal realization. More explicitly, on words in raising operators the invariant pairing is determined by its value on the lowest type: move a raising operator to the opposite side as a lowering operator and commute it through the word. The commutators are compact operators and the last lowering operator kills the lowest vector. Induction gives the pairing on every pair of words, hence the same null space as in the prescribed unitary lowest-weight module. This gives the intertwiner on its irreducible quotient, and therefore the multiplicity vector in \(Y_\sigma\) with evaluation \eqref{eq:B5-full-lift}.
\leanref{GlobalTheory.\allowbreak{}LowRow.\allowbreak{}CuspLift.\allowbreak{}liftIsCuspC}{Development.\allowbreak{}AbelianVarietyTheory.\allowbreak{}Proofs.\allowbreak{}DivisorProofs.\allowbreak{}AdelicLiftMemCuspFormMultiplicitySpace}{22--50}
\leanref{GlobalTheory.\allowbreak{}LowRow.\allowbreak{}exists\_\allowbreak{}isLift\_\allowbreak{}lowC}{Development.\allowbreak{}AbelianVarietyTheory.\allowbreak{}Proofs.\allowbreak{}DivisorProofs.\allowbreak{}AdelCOfInnerProductFormula}{142--156}

Let \(d_\rho\) be the dimension of the lowest compact type, \(N_K=\prod_\rho d_\rho\), and \(d_c\) the dimension of the definite coefficient space. Choose unit extremal vectors \(u_\rho\), put \(u=\bigotimes_\rho u_\rho\) and \(F_z=z(u)\). Schur's lemma on the common discrete series gives a scalar Petersson pairing on the lowest type. Its trace is the multiplicity pairing, and evaluation at \(u\) gives
\begin{equation}\label{eq:B5-petersson-evaluation}
 b(z,z')=\int_{[G]}\overline{F_z(g)}F_{z'}(g)\,dg
         =N_K^{-1}\langle z,z'\rangle .
\end{equation}
In particular evaluation at \(u\), as a function on \([G]\), is injective. Define \(\Psi_z\) by
\[
 \langle\Psi_z(h),b\rangle
  =\int_{[G]}\Theta(h,g;\Phi(u,b))\overline{F_z(g)}\,dg .
\]
The transformation law of the Schwartz vector makes \(\Psi_z\) a definite class form. For representatives \(h_\alpha\) with stabilizer orders \(s_\alpha\), the quotient integral of an invariant pairing on \([H]\) is its class sum with weights \(1/s_\alpha\). Applying \eqref{eq:B5-petersson-evaluation} to \(z,Tx\) therefore gives \(T^*z=N_K\Psi_z\). With an orthonormal basis \(b_1,\ldots,b_{d_c}\) of the coefficient space,
\begin{equation}\label{eq:B5-adjoint-class-pairing}
 \langle T^*z,T^*z'\rangle
  =N_K^2\sum_{a=1}^{d_c}\int_{[H]}
    \langle\Psi_z(h),b_a\rangle
      \overline{\langle\Psi_{z'}(h),b_a\rangle}\,dh .
\end{equation}
\leanref{GlobalTheory.\allowbreak{}AdelicScalar.\allowbreak{}pet\_\allowbreak{}eq}{Development.\allowbreak{}AbelianVarietyTheory.\allowbreak{}Proofs.\allowbreak{}DivisorProofs.\allowbreak{}FormOfClassSums}{75--78}
\leanref{GlobalTheory.\allowbreak{}LowRow.\allowbreak{}form\_\allowbreak{}key\_\allowbreak{}lowC}{Development.\allowbreak{}AbelianVarietyTheory.\allowbreak{}Proofs.\allowbreak{}DivisorProofs.\allowbreak{}AdelCOfInnerProductFormula}{160--197}

Fix \(z\in Y_\sigma(c)\); take \(z'\) arbitrary in \(Y_\sigma\). Double the two theta kernels in \eqref{eq:B5-adjoint-class-pairing}. On the two restrictions of the doubled space the test is \(\Phi_1(X_L)\overline{\Phi_2(-X_R)}\). The splitting contributes \(\chi_V(\det g_2)\) on the second copy of \(G\), which is the multiplier in the doubling integral with character \(\chi_V^{-1}\). The Siegel exponent is \(s_0+r/2=n/2\), hence \(s_0=t/2\). The definite space has Witt index zero, the lattice frames are self-dual, and the doubled polynomial test is compact-finite. The restriction of \(\chi_V\) is the quadratic character to power \(n\). Since \(n>2r\), the definite identity of \ref{B8.12} replaces the theta integral over \([H]\) by \(\operatorname{vol}[H]\) times the Eisenstein value. Here \(s_0>r/2\), in the half-plane of absolute unfolding. In the orbit decomposition each orbit with nonzero isotropic intersection integrates a proper parabolic constant term of a cusp form. The open orbit has representatives \((g,1)\); its section there is the Weil matrix coefficient, and the other factor is the conjugate right-translation pairing of \(F_z,F_{z'}\). Absolute convergence permits integration in this order.
\leanref{GlobalTheory.\allowbreak{}Core.\allowbreak{}Rk.\allowbreak{}IP.\allowbreak{}innerProduct\_\allowbreak{}convergent}{Development.\allowbreak{}AbelianVarietyTheory.\allowbreak{}Proofs.\allowbreak{}DecoratedLatticeProofs.\allowbreak{}UnitaryThetaInnerProductFormulaConvergentRange}{36--156}

At infinity, the invariant-pairing calculation on raising operators gives
\begin{equation}\label{eq:B5-fock-discrete-coefficient}
 \langle\Phi_\rho,\omega_\rho(g)\Phi_\rho\rangle
 =\|\Phi_\rho\|^2\langle u_\rho,\pi_\rho(g)u_\rho\rangle
\end{equation}
for paired unit extremals, where \(\pi_\rho\) is the prescribed discrete series. The identity on the generated modules extends to their unitary completions. The commuting compact definite action makes the corresponding sesquilinear form on its irreducible coefficient space invariant, so its value on each \(b_a,b_a\) equals the value on the paired unit extremal. For a partition \(\beta\) supported on the first \(s\) rows let \(P_\beta\) be its product of principal minors. In the Fock metric the adjoint of multiplication by the \(s\)-row determinant is \(\det(\partial)\). The determinant Capelli identity on this highest vector gives
\begin{equation}\label{eq:B5-fock-capelli}
 \det(\partial)(\det(X)P_\beta)=
           \prod_{i=1}^{s}(\beta_i+1+s-i)\,P_\beta .
\end{equation}
Indeed the upper row operators give determinants with a repeated row, and the diagonal row degrees after multiplication are \(\beta_i+1\). This is the column-determinant identity of \cite{howe-umeda-1991} applied to the \(s\)-by-\(s\) corner. The same product is the quotient of the two hook products on adding a column of height \(s\), by the factorial formula of \ref{B5.07}. Adding columns in order of increasing height gives the squared norm of each minor product. Since the two corners use disjoint variables, their norms multiply. Thus
\[
 H_\rho:=\|\Phi_\rho\|^2=
 \operatorname{Hook}(a^+_{\sigma,\rho})\operatorname{Hook}(a^-_{\sigma,\rho}).
\]
The real part of the open-orbit integral contributes \(\prod_\rho H_\rho J_\rho\), where
\[
 J_\rho=\int_{\U(p_\rho,q_\rho)}
           |\langle u_\rho,\pi_\rho(g)u_\rho\rangle|^2\,dg_\rho
       =V_r\prod_{a<b}\frac{b-a}{|\Lambda_{\rho,a}-\Lambda_{\rho,b}|}.
\]
The second equality is Schur orthogonality and Harish-Chandra's discrete-series formal-degree formula in compact-dual measure \cite{harish-chandra-1966}. Its hypotheses hold: the parameter \(\Lambda_\rho\) is strictly decreasing and \(\Lambda_{\rho,a}-(r-1)/2\) is integral. These are the ordered parameters of the prescribed holomorphic discrete series; a unit vector has fourth power of its norm equal to one in the orthogonality formula.
\leanref{GlobalTheory.\allowbreak{}FockCoefReal.\allowbreak{}fock\_\allowbreak{}coef\_\allowbreak{}real}{Development.\allowbreak{}AbelianVarietyTheory.\allowbreak{}Proofs.\allowbreak{}DecoratedLatticeProofs.\allowbreak{}JointHarmonicsLowestKTypeCoefficientReal}{41--81}
\leanref{GlobalTheory.\allowbreak{}LowRowFock.\allowbreak{}fock\_\allowbreak{}delta}{Development.\allowbreak{}AbelianVarietyTheory.\allowbreak{}Proofs.\allowbreak{}DivisorProofs.\allowbreak{}FockFormMinorProductEqProdHookProducts}{189--203}
\leanref{GlobalTheory.\allowbreak{}LowRow.\allowbreak{}CompactDualCoef.\allowbreak{}inv\_\allowbreak{}integral\_\allowbreak{}coef\_\allowbreak{}sq}{Development.\allowbreak{}AbelianVarietyTheory.\allowbreak{}Proofs.\allowbreak{}DivisorProofs.\allowbreak{}InvIntegralSqDiscreteSeriesCoefficientCompactDua}{90--125}

At \(w\in S_f\), write \(Q=|k_w|\) and \(X=\binom{A}{B}\). At principal level the support is \(A\equiv I_r,\ B\equiv0\) modulo the maximal ideal. Its additive volume and the squared norm of the test are \(m_w=Q^{-nr}\). If both \(X\) and \(Xg^{-1}\) meet the support, then \(A,Ag^{-1}\in K_w^1\), hence \(g\in K_w^1\). On \(K_w^1\) the test has trivial total level character: the Weil determinant multiplier cancels the inverse splitting factor in the test. The cusp evaluation is fixed by \(K_w^1\). Thus the marked integral contributes \(m_w\operatorname{vol}_G(K_w^1)=m_w\). More generally a residue subgroup \(H_w\) has inverse image of volume \(|H_w|\) in this measure; the level-normalized adjoint divides by the residue subgroup order. Here that order is one.

The definite level at \(w\) reduces to the matrices
\(\left(\begin{smallmatrix}I_r&C\\0&D\end{smallmatrix}\right)\) with \(D\in\GL_t(k_w)\). Reduction of the full lattice stabilizer surjects onto \(\GL_n(k_w)\): a lift of an invertible residue matrix is integral invertible, and its inverse adjoint determines the other split component. Its index \(i_w\) satisfies
\begin{equation}\label{eq:B5-marked-index}
 i_wm_w
 =\frac{|\GL_n(k_w)|}{Q^{rt+nr}|\GL_t(k_w)|}
 =\prod_{h=t+1}^{n}(1-Q^{-h})=d_w^{-1}.
\end{equation}
The middle equality follows by counting ordered independent columns and using \(n^2-t^2=rt+nr\). Let \(M\) be the mass of the full standard definite lattice genus, the sum of the reciprocal automorphism orders. Refining the level replaces a class with stabilizer \(A_0\) by its orbits on the level cosets. The sum of their inverse stabilizer orders is the number of cosets divided by \(|A_0|\). Consequently \(\operatorname{vol}[H]=M\prod_{w\in S_f}i_w\).
\leanref{GlobalTheory.\allowbreak{}ThetaOrth.\allowbreak{}orth}{Development.\allowbreak{}AbelianVarietyTheory.\allowbreak{}Proofs.\allowbreak{}TwistedOperators.\allowbreak{}ThetaTranslateOrthogonalOfSupportMovedOff}{184--252}
\leanref{GlobalTheory.\allowbreak{}BaseScalar.\allowbreak{}adjLev\_\allowbreak{}triv}{Development.\allowbreak{}AbelianVarietyTheory.\allowbreak{}Proofs.\allowbreak{}DivisorProofs.\allowbreak{}AdelicScalarBase}{60--66}
\leanref{GlobalTheory.\allowbreak{}LevelIndex.\allowbreak{}levelIndexNormC}{Development.\allowbreak{}AbelianVarietyTheory.\allowbreak{}Proofs.\allowbreak{}TwistedOperators.\allowbreak{}ThetaLevelIndexMulSchwartzNormFormula}{24--69}

At \(w\notin S_f\) all nonzero evaluations of a vector in \(Y_\sigma(c)\) have the Satake multisets of \(c\). Indeed evaluation commutes with each finite sum of translates, and the Satake transform identifies a spherical character with its Weyl orbit. For the second evaluation take the functional \(\ell(F)=\int_{[G]}F\overline{F_{z'}}\). Cauchy--Schwarz and right invariance bound it uniformly on translates of \(F_z\). Lemma~\ref{B5.07e} applies also after translation at other places, since those translations commute. Its absolute majorant justifies passing from finite sets of places to the restricted product. Put
\[
 \mathcal E=\prod_{P:\,P|_F\notin S_f}\prod_{\beta\in B_P}
       (1-\chi_P^{-1}\beta(NP)^{-(1+t)/2})^{-1},
 \qquad d_O=\prod_{w\notin S_f}d_w,\qquad d_S=\prod_{w\in S_f}d_w .
\]
The unmarked integral contributes \(\mathcal E/d_O\). Combining it with \eqref{eq:B5-marked-index} and using \eqref{eq:B5-petersson-evaluation} once on the remaining pairing gives
\begin{equation}\label{eq:B5-full-pairing}
 \langle T^*z,T^*z'\rangle
 =M d_c N_K\left(\prod_\rho H_\rho J_\rho\right)
       \frac{\mathcal E}{d_Od_S}\langle z,z'\rangle .
\end{equation}
One factor \(N_K\) survives from the two adjoints, and \(d_c\) counts the terms in \eqref{eq:B5-adjoint-class-pairing}. For \(z'=z\ne0\) the coefficient is real and nonnegative. Testing against arbitrary \(z'\) proves \(TT^*z=l_cz\), with
\begin{equation}\label{eq:B5-full-scalar}
 l_c=M d_c\prod_\rho(d_\rho H_\rho J_\rho)
                  \frac{|\mathcal E|}{d_Od_S}.
\end{equation}
\leanref{GlobalTheory.\allowbreak{}FinZeta.\allowbreak{}finiteZetaL}{Development.\allowbreak{}AbelianVarietyTheory.\allowbreak{}Proofs.\allowbreak{}TwistedOperators.\allowbreak{}FiniteZetaIntegralThetaLiftEqPartialLValue}{177--232}
\leanref{GlobalTheory.\allowbreak{}LowRow.\allowbreak{}scalar\_\allowbreak{}of\_\allowbreak{}form\_\allowbreak{}key\_\allowbreak{}lowC}{Development.\allowbreak{}AbelianVarietyTheory.\allowbreak{}Proofs.\allowbreak{}DivisorProofs.\allowbreak{}AdelCOfInnerProductFormula}{214--243}

For the size estimate, the normalized spherical functions of the unitary constituents are bounded by one. The bounded spherical-function criterion \cite[\S4.7]{macdonald-1971} puts the real Satake exponents in the convex hull of the Weyl translates of the half-modulus exponent. In the standard base-change coordinates this gives \(|\beta|\le (NP)^{(r-1)/2}\); in the inert chart the radius is measured in \(NP=Q_w^2\). As \(t\ge r+5\), each Euler argument is at most \((NP)^{-2}\) in modulus. For a number field \(L\), grouping the places over each rational prime gives \(\sum_P(NP)^{-2}\le [L:\Q]\sum_{a\ge2}a^{-2}\le[L:\Q]\). From \(\log|1-z|\le |z|\) and \(-\log(1-y)\le2|y|\) for real \(|y|\le1/4\), we obtain nonzero convergent products and
\begin{equation}\label{eq:B5-euler-bounds}
 \log|\mathcal E|\ge-2rf,\qquad \log d_O,\ \log d_S\le4rf .
\end{equation}

The identity class gives \(M\ge1/(w_E^nn!)\), where \(w_E=|\mu(E)|\). Indeed an integral automorphism of the standard definite lattice has entries of absolute value at most one at every embedding. A nonzero entry has absolute value one everywhere by its integral norm and is a root of unity: the monic polynomials of all its powers have bounded integral coefficients. The unit column lengths then give one nonzero entry in each column. Since \(\varphi(w_E)\le2f\) and the prime-power formula gives \(\varphi(h)^2\ge h/2\), we have \(w_E\le8f^2\). Hence
\begin{equation}\label{eq:B5-mass-bounds}
 -\log M\le n\log 2+\log(n!)+4nf .
\end{equation}
\leanref{GlobalTheory.\allowbreak{}LowRow.\allowbreak{}StdMass.\allowbreak{}massLower\_\allowbreak{}of\_\allowbreak{}stdV}{Development.\allowbreak{}AbelianVarietyTheory.\allowbreak{}Proofs.\allowbreak{}DivisorProofs.\allowbreak{}UnitaryMassLowerBoundOfOrthonormalLatticeFrames}{148--171}

The entries of \(\Lambda_\rho\) are the cloud entries and negatives of the auxiliary entries in decreasing order. Adding \(e_j\) puts the former in \([0,3b]\) and the latter in \([-b,b]\). Thus each factor \(|\Lambda_{\rho,a}-\Lambda_{\rho,b}|\) in the numerator of \(J_\rho^{-1}\) is at most \(4b\), and each denominator factor \(b-a\), for the two coordinate indices \(a<b\), is a positive integer. Also \(V_r^{-1}\le r^{r^2}\) by bounding each \(a!\le r^r\) in the volume product. In the selected size datum \(b\ge1\) and \(L'=\log(b\,\mathrm{rd}(E)g_0g_1n)\), where \(\mathrm{rd}(E)\ge1\). The bounds \(k\ge5\) and \(g_0,g_1\ge(kb)^8\) give \(L'\ge5.5\), \(\log n\le L'\) and \(\log(4b)\le L'\). In particular,
\begin{equation}\label{eq:B5-formal-degree-bound}
 J_\rho>0,\qquad
 \log J_\rho^{-1}\le r^2\log r+\tfrac12r(r-1)\log(4b)
                     \le 2+5r^2L'.
\end{equation}
As \(2r<n\), \(f\ge1\) and \(n\ge8\),
\(n(\log2+\log n)+f(2+5r^2L')\le2fn^2L'\).
Equations \eqref{eq:B5-euler-bounds}--\eqref{eq:B5-formal-degree-bound}, together with \(d_c,d_\rho\ge1\), give
\[
 l_c>0,\qquad
 \log l_c\ge\sum_\rho\log H_\rho-2fn^2L'-4nf-10rf
            \ge\sum_\rho\log H_\rho-8fn^2L'.
\]
The embedding permutes the real labels, preserving their signatures. The minus-hook factors are positive integers. Hence \(l_c\ge e^{D_j}\).
\leanref{GlobalTheory.\allowbreak{}LowRow.\allowbreak{}lam\_\allowbreak{}par\_\allowbreak{}width}{Development.\allowbreak{}AbelianVarietyTheory.\allowbreak{}Proofs.\allowbreak{}DivisorProofs.\allowbreak{}LiftModelParameterDecreasingHalfIntegralWidth}{118--163}
\leanref{GlobalTheory.\allowbreak{}LowRow.\allowbreak{}scalar\_\allowbreak{}bound\_\allowbreak{}lowC}{Development.\allowbreak{}AbelianVarietyTheory.\allowbreak{}Proofs.\allowbreak{}DivisorProofs.\allowbreak{}AdelCOfInnerProductFormula}{52--87}

The spherical Hecke operators commute and their adjoints belong to the same algebra by inversion of double cosets. Their joint eigenspaces therefore span \(Y_\sigma\) orthogonally. It follows that \(TT^*\ge e^{D_j}I\). Let \(\Gamma\) be the finite level group, with action \(\rho_\sigma\), and \(\mathcal X_j\) its retained characters. The coefficients
\[
 e_g=|\Gamma|^{-1}\sum_{\chi\in\mathcal X_j}\chi(1)\chi(g^{-1})
\]
give the orthogonal projector \(e=\sum_g\sigma(e_g)\rho_\sigma(g)\) by character orthogonality. Applying it to \eqref{eq:B5-full-lift} gives the required projected evaluation, so injectivity in \eqref{eq:B5-petersson-evaluation} gives \(S=eT\). For \(y\in\operatorname{im}e\),
\(\langle y,SS^*y\rangle=\langle y,TT^*y\rangle\ge e^{D_j}\|y\|^2\).
If \(S^*Sx=\lambda x\) with \(\lambda>0\), then \(Sx\ne0\) is an eigenvector of \(SS^*\) in \(\operatorname{im}e\) with the same eigenvalue. The spectral decomposition of \(S^*S\) on \((\ker S)^\perp\) gives the claimed source inequality.
\leanref{GlobalTheory.\allowbreak{}LowForm.\allowbreak{}parts\_\allowbreak{}scalar\_\allowbreak{}six}{Development.\allowbreak{}AbelianVarietyTheory.\allowbreak{}Proofs.\allowbreak{}DivisorProofs.\allowbreak{}RationalLiftLowerBoundOffBasePlace}{178--193}
\leanref{GlobalTheory.\allowbreak{}LowForm.\allowbreak{}lower\_\allowbreak{}of\_\allowbreak{}parts}{Development.\allowbreak{}AbelianVarietyTheory.\allowbreak{}Proofs.\allowbreak{}DivisorProofs.\allowbreak{}ExistsFullLift}{158--173}
\end{proof}

\subsection*{Proof of \ref{B5.07a}}\label{prf:B5.07a}\phantomsection
\begin{proof}
Put \(U=S_c+m\eta\) and \(I_\rho=\mathbf1_{TT}(\rho)\), and extend \(a_\rho\) by zero off \(TT\). By \eqref{eq:B8-terminal-table}, \(D_{0,\rho}\le U+I_\rho b_\rho\), with \(b_\rho\le2b'\) at a mark. Since \(m\ge1\) and \(\eta\ge0\), \(S_c+\eta\le U\). Hence \(D_{1,\rho}\le3D_{0,\rho}+S_c+\eta\le4U+6I_\rho b'\). At a mark,
\[
 -\lambda_\rho=-a_\rho+(N-k)/2\le-a_\rho+(N-1)/2 .
\]
Elsewhere \(-\lambda_\rho\le (N-1)/2-S_c\). Adding \(D_{0,\rho}\) and using \(n_\rho=\lambda_\rho-D_{0,\rho}\) gives
\begin{equation}\label{eq:B5-terminal-weight-bounds}
 D_{0,\rho}\le U+2I_\rho b',\qquad
 D_{1,\rho}\le4U+6I_\rho b',\qquad
 -n_\rho\le-I_\rho a_\rho+(N-1)/2+U+2I_\rho b'.
\end{equation}

Write \(T=T_u\). Since \(m\le N\) and \(L\le L_a\), the table gives \(0\le\log T\le C_TNL_a\). Also \(a=xe^L\), \(\operatorname{rd}E\ge1\) and \(L_a\ge1\), so \(\log(2a)\le2L_a\). Keep \(-I_\rho a_\rho\log T\) at its exact value. The remaining exponent term is bounded by \(C_TNL_a((N-1)/2+U+2I_\rho b')\). The other two entries of the terminal table satisfy
\[
 2c_\rho^*\le 4NL_a(4U+6I_\rho b'),\qquad
 \xi_\rho^*\le N^2L_a .
\]
Sum over \(|\mathcal P|\le f\) types, of which \(n_T\) are marked, and add \(e_{\rm ev}\le C_{\rm ev}fN^2L_a\). This gives
\begin{equation}\label{eq:B5-terminal-table-bound}
 A+S_a\log T\le L_a\bigl((2C_T+24)n_Tb'N+
 (C_T+16)fUN+(C_T/2+1+C_{\rm ev})fN^2\bigr).
\end{equation}
The first bound in \eqref{eq:B8-table-constants} dominates each coefficient. The right side is therefore at most \(H_{P,u}+S_a\log T\). In the selected table the evaluation allowance is zero: \(\varpi\bar\varpi=T\) in \ref{B5.07c}, and \eqref{eq:B8-terminal-period} gives the squared norm bound in the stated frames. For a general table the same norm bound is \eqref{eq:B8-terminal-bound}.

For a layer \(j\), take the mean in \eqref{eq:B8-height-size}. By \ref{B5.07x} the mean of \(\omega_{\rho,v}\) is \(\log\Omega\). The term \(-n_{j,\rho}\log\Omega\) is therefore exactly \(-\sum_i\nu_{uj\rho i}\log\Omega\). Together with the factorial terms this gives \(\sum_\rho K_{uj\rho}\) with the normalization of \ref{B5.07}. The constant and degree-log terms cost at most \(fN^2(1+2c_6)L_a\). Since \(p_j=|\mathcal P|g_j^0\le fN^2\), the two terms multiplied by \(p_j\) cost at most \(fN^2(C_\theta+c_5)L_a\).

For the remaining terms, insert \(g_\ast\le C_8L\) and \(\log T\le C_TmL\) in \eqref{eq:B8-kernel-table}. The term \(2C_efN^2g_\ast\) costs at most \(2C_eC_8 fN^2L_a\). In the two corrections \(h_j,h'_j\), the unmultiplied \(L\) costs at most \(2C_4fN^2L_a\). The terms containing \(g_\ast+\log T\) are at most \(2(C_2+C_4)(C_T+C_8)fm^3N^2L\), since \(m\ge1\). Thus
\begin{equation}\label{eq:B5-kernel-table-bound}
 \langle s_{jv}\rangle_\infty-\sum_\rho K_{uj\rho}
 \le fN^2\bigl((C_\theta+c_5+2c_6+1+2C_4+2C_eC_8)L_a
 +2(C_2+C_4)(C_T+C_8)m^3L\bigr).
\end{equation}
The last two bounds in \eqref{eq:B8-table-constants} give \(\langle s_{jv}\rangle_\infty\le H_{K,uj}\).

Finally, apply weighted Cauchy--Schwarz to the sum over input indices in \eqref{eq:B8-coefficient-rules}. At stage \(i\) the sum of the weights is at most the number \(a_i h_i^{\rm cl}\) of indices, since every weight is at most one. Summing the resulting squared estimates with the output weights costs at most \(a_{i+1}h_{i+1}^{\rm cl}\). The two frame identities consequently give \(\|K_i\|_v^2\le a_i h_i^{\rm cl}a_{i+1}h_{i+1}^{\rm cl}t_{\max}(i,v)\). If the product vanishes, \(K_i=0\). Otherwise exponentiate \eqref{eq:B8-layer-bound}.
\end{proof}

\subsection*{Proof of \ref{B5.07fa}}\label{prf:B5.07fa}\phantomsection
\begin{proof}
The determinant ideal of \(L_a^h\) is \(\mathfrak c_a/\overline{\mathfrak c_a}\), with \(\mathfrak c_a\) integral and prime to its conjugate. Indeed, at a chosen split prime \(P\) write \(e_1,\ldots,e_n\) for the elementary-divisor exponents relative to \(L_0\), with \(\varpi^{e_i}\mathcal O_{E,P}\) as the convention. At \(\bar P\) the exponents are the negatives, up to order, because the second matrix is the inverse adjoint. Put \(s_P=\sum_i e_i\) and \(Q_P=\mathrm N P\). Taking the positive part of the determinant divisor gives
\[
 \mathrm N\mathfrak c_a=\prod_P Q_P^{|s_P|}.
\]
The contribution to \([L_a^h:L_a^h\cap L_0]\) at \(P\) is \(Q_P^{\sum_i\max(-e_i,0)}\), and at \(\bar P\) it is \(Q_P^{\sum_i\max(e_i,0)}\). Outside the moving pairs the two lattices agree. Conjugation gives an additive isomorphism on the intersection quotients, so
\[
 I_a=\prod_P Q_P^{r\sum_i|e_i|},\qquad Z_a=m_a I_a.
\]
The second equality also follows directly by multiplying the indices in the tower from \(N_aL_0\) to \(L_a\).

Here is the normalization of the class coefficient in \ref{B8.55}. Write \(B(\mathfrak c_a)\) for the algebraic determinant value with its positive norm factor. The coefficient is
\[
 \epsilon_a=\frac{t_a}{B(\mathfrak c_a)},\qquad
 t_a\in\Q_{>0},\qquad t_a^2 I_a=\mathrm N\mathfrak c_a.
\]
In particular the rational factor in this definition can be computed before choosing a coefficient embedding. The ranks in \ref{B8.58}--\ref{B8.59} give \(r\) odd. Since \(|s_P|\equiv\sum_i|e_i|\pmod 2\), the unique positive choice is
\[
 t_a=\prod_P Q_P^{\left(|s_P|-r\sum_i|e_i|\right)/2}.
\]
\leanref{Formalism.\allowbreak{}L2.\allowbreak{}Pj.\allowbreak{}EpIs}{Formalism.\allowbreak{}LocalToGlobal.\allowbreak{}ProjectiveLattices.\allowbreak{}ThClass}{106--115}
\leanref{GlobalTheory.\allowbreak{}ClassTieCompare.\allowbreak{}exists\_\allowbreak{}epIs\_\allowbreak{}iff}{Development.\allowbreak{}AbelianVarietyTheory.\allowbreak{}Proofs.\allowbreak{}UnitaryExistence.\allowbreak{}ExistsFirstOfSecond}{69--92}
\leanref{GlobalTheory.\allowbreak{}ClassScalarCompare.\allowbreak{}stdIndex\_\allowbreak{}conjLat}{Development.\allowbreak{}AbelianVarietyTheory.\allowbreak{}Proofs.\allowbreak{}UnitaryExistence.\allowbreak{}NotIsOfFinOrderUnitQuot}{139--156}

We compute the image of \(B(\mathfrak c_a)\). At the reference tuple the twist is trivial and the splitting character is \(\mu_0\eta^{-1}\). The moving primes avoid the bad modulus, so the local characters there are unramified. On the two primes of a moving pair their product character has trivial restriction to the diagonal local field. Thus a factor with \(s_P<0\) on the chosen half can equally be evaluated with exponent \(-s_P\) on the opposite half. The ideal-character value on \(\mathfrak c_a\) is consequently \((\mu_0\eta^{-1})(d_a)\).

For the change of embedding, let \(\pi_P\) denote the idèle of a uniformizer on the chosen half, equal to one off \(P\). The conjugate of \(\mu_0(d_a)\) is specified by its coefficient-field value, since \(d_a\) has infinite part one. For \(\eta\), the coefficient-field value \(b_P\) at \(\pi_P\) has base reading \(\eta(\pi_P)/\sqrt{Q_P}\) and reading \(\eta_\sigma(\pi_P)/\sqrt{Q_P}\) at \(\sigma\). Hence the algebraic product
\[
 \prod_P b_P^{-s_P} Q_P^{(|s_P|-s_P)/2}
\]
has reading \(\eta_\sigma(d_a)^{-1}\prod_P Q_P^{|s_P|/2}\). Multiplying by the coefficient of \(\mu_0(d_a)\) gives \(B(\mathfrak c_a)\) at the base embedding, by its defining ideal-character formula. Injectivity at the base now identifies the algebraic elements. At \(\widetilde\sigma\) this gives
\[
 \widetilde\sigma\bigl(B(\mathfrak c_a)\bigr)
       =\chi_{W,\sigma}(d_a)\sqrt{\mathrm N\mathfrak c_a}.
\]
Since \(\sqrt{\mathrm N\mathfrak c_a}=t_a\sqrt{I_a}\), substitution gives
\[
 \widetilde\sigma(\epsilon_a)\frac{Z_a}{m_a}
 =\frac{t_a I_a}{\chi_{W,\sigma}(d_a)\sqrt{\mathrm N\mathfrak c_a}}
 =\chi_{W,\sigma}(d_a)^{-1}\sqrt{I_a}.
\]
\leanref{GlobalTheory.\allowbreak{}GalOne.\allowbreak{}ClassConj.\allowbreak{}coefficientField\_\allowbreak{}reading\_\allowbreak{}at\_\allowbreak{}embedding\_\allowbreak{}prod}{Development.\allowbreak{}AbelianVarietyTheory.\allowbreak{}Proofs.\allowbreak{}DivisorProofs.\allowbreak{}EmbeddingConjugateOscillatorSixth}{148--182}
\leanref{GlobalTheory.\allowbreak{}GalOne.\allowbreak{}ClassConj.\allowbreak{}exists\_\allowbreak{}ideal\_\allowbreak{}conj}{Development.\allowbreak{}AbelianVarietyTheory.\allowbreak{}Proofs.\allowbreak{}DivisorProofs.\allowbreak{}EmbeddingConjugateOscillatorSixth}{185--295}
\leanref{GlobalTheory.\allowbreak{}ClassScalarValue.\allowbreak{}zCard\_\allowbreak{}eq}{Development.\allowbreak{}AbelianVarietyTheory.\allowbreak{}Proofs.\allowbreak{}DivisorProofs.\allowbreak{}OscillatorModelValuesThird}{293--299}
\leanref{GlobalTheory.\allowbreak{}ClassScalarValue.\allowbreak{}value\_\allowbreak{}calc}{Development.\allowbreak{}AbelianVarietyTheory.\allowbreak{}Proofs.\allowbreak{}DivisorProofs.\allowbreak{}OscillatorModelValuesThird}{317--334}
\end{proof}

\subsection*{Proof of \ref{B5.07f}}\label{prf:B5.07f}\phantomsection
\begin{proof}
Write \(u_\sigma\) for the prescribed unit extremal tensor and \(\operatorname{ev}_\sigma(y)(h)=y(u_\sigma)(h)\) for evaluation in the target multiplicity space. If \(f\) is an algebraic cusp section, denote its rational coordinate column by \([f]\). At the framed point belonging to \(h\), write \(R_\sigma(f,h)\) for the value of the section in its algebraic frame, \(C(h)\) for the component, and \(\mu_\sigma(h)\) for the adelic multiplier taking that framed value to the adelic section. If \(c_\sigma(C)\) denotes the chosen component coefficient, the realization of \([f]\) is specified by
\[
 \operatorname{ev}_\sigma(J_\sigma[f])(h)
 =a_j c_\sigma(C(h))\mu_\sigma(h)R_\sigma(f,h),\qquad
 a_j=(2\pi\xi)^{-A_j/2}.
\]
Here \(\xi=(\sqrt d\,\mathrm N\mathfrak c_K)^{-1}\) and \(A_j\) is the total degree in \ref{B8.55}. In particular, the function called \(\varphi_j(c_j,f)\) in \eqref{eq:p8c5-theta-extension} acquires the factor \(a_j\) upon extremal evaluation of the rational column. This formula on the rational cusp frame extends by semilinearity to its coordinate space.
\leanref{Formalism.\allowbreak{}L2.\allowbreak{}Pj.\allowbreak{}CoreJ.\allowbreak{}phiJ}{Formalism.\allowbreak{}LocalToGlobal.\allowbreak{}ProjectiveLattices.\allowbreak{}FamilyUniformisationBasePointOrthogonalBis}{132--135}
\leanref{Formalism.\allowbreak{}IndefR.\allowbreak{}FWT.\allowbreak{}FWTie.\allowbreak{}evaluates}{Formalism.\allowbreak{}LocalToGlobal.\allowbreak{}IndefiniteReal.\allowbreak{}FWExist}{645--650}

Let \(\Gamma\) be the finite level group and let \(\rho_\sigma\) denote its action by right translation. Write \(e_g\in K\) for the coefficients of the retained central idempotent in \ref{B5.07y}, so that \(e_\sigma=\sum_{g\in\Gamma}\sigma(e_g)\rho_\sigma(g)\). Write \(\delta:\Gamma\to K^\times\) for the character in the level action of \ref{B8.55}. Choose definite class representatives. For \(x\in X\), take the algebraic theta section of its list of values on those representatives. By \ref{B8.55} it has an expansion \(f_x=\sum_\nu a_\nu f_\nu\) after base change from cusp sections, with \(a_\nu\) read in \(K\) through the reference field identification. Use this expansion of the class section itself, before the central translation in the reference identity. Denote that translation by \(z_0\) and the corresponding root of unity, pulled back to \(K\), by \(\zeta_0\). Put
\begin{equation}\label{eq:B5-rational-lift-formula}
 \theta x=\sum_{g\in\Gamma} e_g\delta(g)^{-1}
           \sum_\nu\zeta_0a_\nu[\,g(z_0 f_\nu)\,].
\end{equation}
The level action preserves the rational cusp space. At the base embedding its determinant multiplier cancels \(\delta(g)^{-1}\). By the reference class identity of \ref{B8.55}, with phase one and with the factor \(a_j\) in the displayed evaluation formula, the inner sum then gives the unprojected theta function. The outer sum projects it. Hence
\(\operatorname{ev}_0(J_0\theta x)=\operatorname{ev}_0(S_0j_0x)\).
Evaluation is injective by \eqref{eq:B5-petersson-evaluation}, and the realized cusp frame is independent. Thus the column in \eqref{eq:B5-rational-lift-formula} is uniquely determined by \(x\). Additivity and \(K\)-linearity follow by applying \(J_0\) to the two sides of the corresponding identities, using its semilinearity and that of \(j_0\).
\leanref{GlobalTheory.\allowbreak{}GalLift.\allowbreak{}Assembly.\allowbreak{}ratOf}{Development.\allowbreak{}AbelianVarietyTheory.\allowbreak{}Proofs.\allowbreak{}DivisorProofs.\allowbreak{}GalLiftOfOneMap}{127--131}
\leanref{GlobalTheory.\allowbreak{}ConjRatLiftFive.\allowbreak{}base\_\allowbreak{}lift5}{Development.\allowbreak{}AbelianVarietyTheory.\allowbreak{}Proofs.\allowbreak{}DivisorProofs.\allowbreak{}ConjugatesOfRationalThetaLiftFromCharacterisatio}{46--68}

Fix \(\sigma\) over the distinguished embedding of \(K_d\). The map from \(K\) to the reference comparison field is an isomorphism onto that field, so composition of its inverse with \(\sigma\) gives an embedding \(\widetilde\sigma\) of the entire chart. All subsequent evaluations of algebraic sections use this embedding.

First consider the finite test. Let \(M\) be the covering modulus, \(A\) the finite group of \(n\)-tuples of arguments of the level indicator \(\varphi\), and \(A^\vee\) its dual under the level pairing. At the reference tuple the twisting set is empty, so \(\varphi\) takes values in \(\{0,1\}\). For a primitive \(M\)-th root \(\zeta_M\), use the Fourier weights
\[
 w(q)=\sum_{b\in A}\varphi(b)\zeta_M^{\langle b,q\rangle},
 \qquad q\in A^\vee.
\]
Choose \(d_\sigma\) prime to \(M\) such that \(\widetilde\sigma(\zeta_M)=\zeta_M^{d_\sigma}\) in the reference coordinates. Then \(\widetilde\sigma(w(q))=w(d_\sigma q)\). The sign \(\varepsilon\) converting the finite adelic Heisenberg action to the action on the theta chart is fixed by the choice of the two actions, hence is the same for every class. The required change of marking is multiplication by \(n_\sigma=\varepsilon d_\sigma\). The base modulus divides \(M\). At every marked prime the scalar matrix \(n_\sigma I_r\) is therefore invertible. Under the level-group identification with the product of the marked general linear groups these matrices give an element \(z\in Z(\Gamma)\). Its action on the marked half of the level lattice is multiplication by \(n_\sigma\) modulo the base level.
\leanref{GlobalTheory.\allowbreak{}GalOne.\allowbreak{}conj\_\allowbreak{}coeff\_\allowbreak{}eq\_\allowbreak{}dilate\_\allowbreak{}ofFive}{Development.\allowbreak{}AbelianVarietyTheory.\allowbreak{}Proofs.\allowbreak{}DivisorProofs.\allowbreak{}CoefMapNormEquivariantBasisOfFive}{2428--2455}
\leanref{GlobalTheory.\allowbreak{}GalCentral.\allowbreak{}central\_\allowbreak{}lev\_\allowbreak{}ofFive}{Development.\allowbreak{}AbelianVarietyTheory.\allowbreak{}Proofs.\allowbreak{}DivisorProofs.\allowbreak{}CoefMapNormEquivariantBasisOfFive}{2483--2506}

Fix a framed point \(h\), let \(L\) be its self-dual target lattice, and let \(v\) be a unit vacuum for \(L_0\otimes L\). For a class \(a\) let \(v_a\) be the average of the Heisenberg translates of \(v\) by \(L_a\otimes L\), weighted by the lattice character. This is the projection onto the vacuum line for \(L_a\otimes L\). Its normalization can be seen in the split Schr\"odinger model. On one half write \(D_0,D_a\) for the old and new position lattices on the tensor space, \(D_\cap=D_0\cap D_a\), \(m=[D_a:D_\cap]\), and \(n'=[D_0:D_\cap]\). Averaging the modulations for the new dual half restricts \(\mathbf1_{D_0}\) to \(D_\cap\). Averaging its translations over \(D_a/D_\cap\) gives \(m^{-1}\mathbf1_{D_a}\). The unitary change of variables sends the old indicator to \(\sqrt{n'/m}\,\mathbf1_{D_a}\) times the inverse splitting character on the chosen determinant half. Thus its multiplier relative to the projection is the character times \(\sqrt{mn'}\). The calculation applies to successive disjoint split pairs. Localizing \(L\) gives \(r\) free coordinates, so the product of the \(mn'\) is \(I_a\) of \ref{B5.07fa}. Consequently the analytic change of vacuum has coefficient \(\chi_{W,\sigma}(d_a)^{-1}\sqrt{I_a}\) against \(v_a\). By \ref{B5.07fa} it is exactly the embedded class coefficient times \(Z_a/m_a\).
\leanref{GlobalTheory.\allowbreak{}OneMapRef.\allowbreak{}SplitProj.\allowbreak{}step}{Development.\allowbreak{}AbelianVarietyTheory.\allowbreak{}Proofs.\allowbreak{}IwasawaDecomposition.\allowbreak{}SplitFinitePointMovesVacuumToProjection}{115--210}
\leanref{GlobalTheory.\allowbreak{}T.\allowbreak{}splitFinitePoint\_\allowbreak{}moves\_\allowbreak{}vacuum\_\allowbreak{}to\_\allowbreak{}projection}{Development.\allowbreak{}AbelianVarietyTheory.\allowbreak{}Proofs.\allowbreak{}IwasawaDecomposition.\allowbreak{}SplitFinitePointMovesVacuumToProjectionStatement}{16--18}

We give the marking calculation, since a change of a Heisenberg lift can introduce a scalar. For \(q\in A^\vee\) choose a displacement \(Y_a(q)\) for the signed marking. Row by row it is congruent to the marked representative divided by \(M\), modulo \(L\). The lifts are taken with the clearing condition that a sufficiently high power of the marked half-ideal carries them into \(2N_a L\). Such lifts exist by the Chinese remainder theorem on the two coprime halves of the level. Equivalently, one prescribes the marking on its supporting half and divisibility by \(2N_a\) off that half. Take \(N_a\) divisible by \(M\) and take a second clearing integer \(k_a\), prime to \(M\), carrying \(L_0\) into \(L_a\). These choices are possible because the moving places avoid the level.

We show that these displacements give the transported translates with their characters. Write \(H\) for the completed tensor lattice of the moved vacuum, \(X'\) for the displacement transported from the reference marking, and \(Y'\) for the displacement in the new marking, both in finite adelic coordinates with the same sign. It suffices to consider \(q\) with nonzero Fourier weight. Its coordinates are annihilated by an integer supported at the bad primes of the test, disjoint from the moving pairs. Indeed, the indicator is constant on cosets of the kernel of reduction to its test modulus, so summing a nontrivial character on that kernel gives zero. At the moving pairs \(q\) is therefore integral as a displacement. At the other places the two markings give the same coset modulo the original integral lattice. Clearing that lattice into the moved lattice gives
\[
 k_a(X'-Y')\in H,\qquad k_a M Y'\in H .
\]
Let \(\mathfrak b\) be the marked half-ideal. For \(l\) large enough for both root systems, choose \(s\in\mathfrak b^l\) and \(t\) with \(\bar t\in\mathfrak b^l\) and \(s+2k_a^2t=1\). This is possible since the level is odd, \(k_a\) is prime to it and its two halves are coprime. By the clearing conditions on the old and new roots,
\[
 sX'=2b_1,\quad sY'=2b_2,\quad
 \bar t X'=2M b_3,\quad k_a\bar t Y'=2M b_4,\qquad b_i\in H.
\]
Put \(b=b_1-b_2+k_at\,k_a(X'-Y')\). Then \(b\in H\) and \(X'-Y'=2b\). If \(\omega\) is the alternating form in the Heisenberg character, sesquilinearity gives
\[
 \omega(Y',b)=\omega(2b_2,b)
            +\omega(k_a M Y',\,2(k_a b_3-b_4)).
\]
Both terms belong to the kernel of the additive character by integrality on \(H\). The doubled lattice \(2H\) acts trivially on the moved vacuum. The Heisenberg multiplication law now gives identical actions for \(X'\) and \(Y'\). This also proves independence of the admissible choices of roots. Replacing \(q\) by \(d_\sigma q\) gives admissible roots for the marking multiplied by \(d_\sigma\).
\leanref{GlobalTheory.\allowbreak{}OneMapRef.\allowbreak{}Term.\allowbreak{}heis\_\allowbreak{}congr}{Development.\allowbreak{}AbelianVarietyTheory.\allowbreak{}Proofs.\allowbreak{}IwasawaDecomposition.\allowbreak{}OmegaFormSignIndependent}{64--78}
\leanref{GlobalTheory.\allowbreak{}OneMapRef.\allowbreak{}Term.\allowbreak{}heis\_\allowbreak{}root}{Development.\allowbreak{}AbelianVarietyTheory.\allowbreak{}Proofs.\allowbreak{}IwasawaDecomposition.\allowbreak{}OmegaFormSignIndependent}{304--390}
\leanref{GlobalTheory.\allowbreak{}OneMapRef.\allowbreak{}Term.\allowbreak{}heis\_\allowbreak{}root\_\allowbreak{}zak}{Development.\allowbreak{}AbelianVarietyTheory.\allowbreak{}Proofs.\allowbreak{}IwasawaDecomposition.\allowbreak{}OmegaFormSignIndependent}{412--460}

We next compare the polynomial applied to this finite vector. In the complex definite coefficient space at \(\sigma\), consider the map obtained by pairing with the analytic extremal tensor, and the map obtained by embedding the algebraic polynomial coefficients and relabelling their variables. Put both maps into the same Fock space by applying the polynomials to the standard vacuum. Call the maps \(A\) and \(B\), respectively. The relabelling keeps the column index. With \(p=k+j-1\), \(q_-=m-j\), a row numbered \(i\) from zero goes to row \(q_-+i\) on the positive side and \(q_--1-i\) on the conjugate side in the framed derivative coordinates. Thus the prescribed product of minors is the same polynomial in these coordinates for the two maps. The algebraic highest vector has norm one after embedding: its Gram value in the standard definite frame is the rational number one. On the analytic side the vector giving the minor polynomial has norm one by \ref{B8.55}. Changing the definite orthonormal basis acts unitarily, including the vacuum character. We obtain unit vectors \(w_A,w_B\) with \(Aw_A=Bw_B\ne0\). Nonvanishing follows also by taking the leading monomial of the minor product in the polynomial realization of the Fock space.

We check that the comparison on these vectors holds on the whole coefficient space. For a rational unitary element \(\gamma\), the analytic map intertwines its action with the oscillator action on the output. For the algebraic map, substituting the matrix of \(\gamma\) in the source columns of the minor polynomial gives the same oscillator action times a scalar: the scalar is the embedded coefficient-line multiplier divided by the vacuum character of the orthonormal change of basis. Thus, writing \(\rho\) for the coefficient action, there are an operator \(T_\gamma\) and a nonzero scalar \(b_\gamma\) with
\[
 A\rho(\gamma)=T_\gamma A,\qquad
 B\rho(\gamma)=b_\gamma T_\gamma B.
\]
The rational unitary group is dense in the product of the compact real groups. Indeed the Cayley transform gives rational coordinates near the identity from the skew-adjoint matrices, and weak approximation there gives a dense subset of a neighborhood of the identity. The compact unitary groups are connected. Hence the irreducible coefficient representation remains irreducible on the rational group. Both maps are injective, since their kernels are invariant and they have a common nonzero value. Likewise \(A^{-1}(\operatorname{im}B)\) is invariant and contains \(w_A\). It is the full space, and equality of dimensions gives \(\operatorname{im}A=\operatorname{im}B\). The automorphism \(V=B^{-1}A\) satisfies
\[
 V\rho(\gamma)V^{-1}=b_\gamma^{-1}\rho(\gamma).
\]
For every \(\gamma\) of nonzero trace, taking the trace gives \(b_\gamma=1\). These elements generate the rational group. To see this, the trace is a nonzero real analytic function on the connected compact group, since its value at the identity is the dimension. Its nonvanishing set is open dense. For any rational \(\gamma\), density gives a rational \(\beta\) with both \(\beta\) and \(\gamma\beta\) in this set. Schur's lemma gives \(V=\alpha\,1\). As \(\alpha w_A=w_B\), \(|\alpha|=1\). Put \(p_\sigma=\alpha^{-1}\). This phase works for every class value and every input. Finally each derivative letter in the passage from Fock to Lie coordinates contributes \(1/\sqrt{2\pi\xi}\). The coefficient polynomials are homogeneous of degree \(A_j\), as follows by counting the row degrees of their generating minor product and applying linear substitutions. After multiplication by \(p_\sigma\) the analytic differential polynomial is therefore \(a_j\) times the embedded algebraic one.
\leanref{GlobalTheory.\allowbreak{}AdTh.\allowbreak{}PhaseMap.\allowbreak{}proportional}{Development.\allowbreak{}AbelianVarietyTheory.\allowbreak{}Proofs.\allowbreak{}DivisorProofs.\allowbreak{}OscillatorModelExtensionFirst}{31--76}
\leanref{GlobalTheory.\allowbreak{}GalOne.\allowbreak{}PhaseAt.\allowbreak{}io\_\allowbreak{}eq\_\allowbreak{}phase\_\allowbreak{}smul\_\allowbreak{}at\_\allowbreak{}ofFive}{Development.\allowbreak{}AbelianVarietyTheory.\allowbreak{}Proofs.\allowbreak{}DivisorProofs.\allowbreak{}CoefMapNormEquivariantBasisOfFive}{827--996}
\leanref{GlobalTheory.\allowbreak{}OneMapRef.\allowbreak{}IoCmp.\allowbreak{}two\_\allowbreak{}div\_\allowbreak{}fockScale}{Development.\allowbreak{}AbelianVarietyTheory.\allowbreak{}Proofs.\allowbreak{}IwasawaDecomposition.\allowbreak{}OmegaFormRationality}{78--80}
\leanref{GlobalTheory.\allowbreak{}GalOne.\allowbreak{}IoCmp.\allowbreak{}ioAt\_\allowbreak{}of\_\allowbreak{}conj\_\allowbreak{}ofFive}{Development.\allowbreak{}AbelianVarietyTheory.\allowbreak{}Proofs.\allowbreak{}DivisorProofs.\allowbreak{}CoefMapNormEquivariantBasisOfFive}{1033--1143}

We can now write the two evaluations at \(h\). Let \(\mathcal D_h(P,v')\) mean that the theta series of the test vector \(v'\) is differentiated by \(P\) in the framed Lie directions and evaluated at zero. Denote the embedded, relabelled polynomial of \(x\) at class \(a\) by \(P_{a,\sigma}\), the common representative level volume by \(V_{\rm lev}\), its stabilizer order at \(a\) by \(s_a\), and Heisenberg translation by \(D(Y)\). The common finite sum is
\[
 \mathcal C(h,x)=
 \sum_a\frac{V_{\rm lev}}{s_a}
 \frac{\widetilde\sigma(\epsilon_a)Z_a}{m_a|A^\vee|}
 \mathcal D_h\!\left(P_{a,\sigma},
       \sum_{q\in A^\vee}\widetilde\sigma(w(q))
                   D(Y_a(d_\sigma q))v_a\right).
\]
For the analytic lift, finite Fourier inversion expresses the level test as the sum of translates with base weights \(w(q)/|A^\vee|\). The equality \(\widetilde\sigma(w(q))=w(d_\sigma q)\) and the permutation \(q\mapsto d_\sigma q\) give exactly the vector displayed above. The vacuum projection calculation supplies its class coefficient. Integration over each definite class contributes \(V_{\rm lev}/s_a\). On the algebraic side the central translation \(z\) pulls the marking back to \(n_\sigma\) times the original marking, so that the embedded Fourier weights and the roots \(Y_a(d_\sigma q)\) occur together. These operations commute with evaluating a polynomial differential operator on a finite sum.

For the remaining frame scalars let \(\beta_h\) be the scalar expressing the product of the degree-one theta generators of the framed point in the theta function of \(v\). Both are sections of the same principal theta lines, each of dimension one. Use the same generators on the covering level by pullback. Let \(c'_\sigma(C)\) denote the component coefficient of the analytic vacuum in this framing after removing \(\mu_\sigma(h)\). More explicitly, at one point of each component compare the transported unit vacuum with the theta generator normalized to have mean squared norm one. Their quotient, divided by the unitary multiplier \(\mu_\sigma(h)\), gives \(c'_\sigma(C)\), of modulus one. At another framed point of the component the real change of frame transports the Gaussian and its arguments together. The finite change between the two transported target lattices carries the self-dual vacuum by the quotient of their splitting multipliers. Dividing by \(\mu_\sigma\) cancels that quotient. A rational change of representative fixes the theta series by change of summation variable. This gives the same \(c'_\sigma(C)\) in every such chart. In particular the scalar of the analytic vacuum relative to \(v\) at \(h\) is \(c'_\sigma(C(h))\mu_\sigma(h)\beta_h\).
\leanref{GlobalTheory.\allowbreak{}OneMapRef.\allowbreak{}FrameConst.\allowbreak{}exists\_\allowbreak{}frConst}{Development.\allowbreak{}AbelianVarietyTheory.\allowbreak{}Proofs.\allowbreak{}IwasawaDecomposition.\allowbreak{}OmegaFormScalarExtension}{246--316}
\leanref{GlobalTheory.\allowbreak{}GalOne.\allowbreak{}FrameConst.\allowbreak{}exists\_\allowbreak{}frAtEmb\_\allowbreak{}ofFive}{Development.\allowbreak{}AbelianVarietyTheory.\allowbreak{}Proofs.\allowbreak{}DivisorProofs.\allowbreak{}CoefMapNormEquivariantBasisOfFive}{516--536}

The covering section is constant at the reference tuple. At the distinguished embedding its value at a marked point is the prescribed root of unity times the twist character of a determinant root. The reference twist is trivial. The marked charts cover the complex components, and evaluation on them identifies the section with the constant section of that value. This equality descends to the comparison field: scalar extension to \(\C\) is injective on sections, and evaluation at the chosen coefficient-field point gives the constant in that field. If its base value has order \(m_0\), its coefficient-field value has \(m_0\)-th power one by the same injectivity. Denote its reading under \(\widetilde\sigma\) by \(c_0\). Thus \(|c_0|=1\).

Write \(T_\sigma\) for the unprojected analytic lift. Substitution in the differentiated theta series and in the evaluation of the translated algebraic class section gives, respectively,
\[
 \begin{aligned}
 p_\sigma\,\operatorname{ev}_\sigma(T_\sigma j_\sigma x)(h)
   &=c'_\sigma(C(h))\mu_\sigma(h)\beta_h\,a_j\,\mathcal C(h,x),\\
 \sum_\nu\sigma(a_\nu)R_\sigma(zf_\nu,h)
   &=c_0\beta_h\,\mathcal C(h,x).
 \end{aligned}
\]
For the second equality the expansion of \(f_x\) is an equality of algebraic sections, so it can first be translated and then evaluated over \(\widetilde\sigma\). Multiplying the first equality by \(c_0\) and the second by \(a_j c'_\sigma(C(h))\mu_\sigma(h)\) identifies the two sides even at a zero of the evaluated section.
\leanref{GlobalTheory.\allowbreak{}OneMapRefCompare.\allowbreak{}cover\_\allowbreak{}const}{Development.\allowbreak{}AbelianVarietyTheory.\allowbreak{}Proofs.\allowbreak{}IwasawaDecomposition.\allowbreak{}OneMapTypingsEquiv}{180--206}
\leanref{GlobalTheory.\allowbreak{}GalOne.\allowbreak{}adelic\_\allowbreak{}value\_\allowbreak{}ofFive}{Development.\allowbreak{}AbelianVarietyTheory.\allowbreak{}Proofs.\allowbreak{}DivisorProofs.\allowbreak{}CoefMapNormEquivariantBasisOfFive}{2533--2588}
\leanref{GlobalTheory.\allowbreak{}GalOne.\allowbreak{}geometric\_\allowbreak{}value\_\allowbreak{}ofFive}{Development.\allowbreak{}AbelianVarietyTheory.\allowbreak{}Proofs.\allowbreak{}DivisorProofs.\allowbreak{}CoefMapNormEquivariantBasisOfFive}{2591--2648}
\leanref{GlobalTheory.\allowbreak{}GalOne.\allowbreak{}identity\_\allowbreak{}of\_\allowbreak{}frConst\_\allowbreak{}ofFive}{Development.\allowbreak{}AbelianVarietyTheory.\allowbreak{}Proofs.\allowbreak{}DivisorProofs.\allowbreak{}CoefMapNormEquivariantBasisOfFive}{2651--2732}

Put \(u(h)=c'_\sigma(C(h))/c_\sigma(C(h))\). This function descends to the components of the quotient at principal level and has modulus one. Multiplication by \(u\) preserves the real representation and the principal level. It also preserves cuspidality: on the orbit of a unipotent radical the component is constant, by approximation of its finite adelic coordinates by rational ones and invariance under the real group. Multiplication by \(\bar u\) is its inverse. The Petersson integral in \eqref{eq:B5-petersson-evaluation} shows that it induces a unitary operator \(U_\sigma\) on the multiplicity space. Write \(J'=U_\sigma J_\sigma\) and \(\kappa=p_\sigma c_0\). The two evaluations just proved say
\[
 \kappa\,\operatorname{ev}_\sigma(T_\sigma j_\sigma x)
       =\sum_\nu\sigma(a_\nu)\operatorname{ev}_\sigma(J'[zf_\nu]).
\]

In this frame the level law is
\[
 J'[gf]=\sigma(\delta(g))\,\rho_\sigma(g)J'[f].
\]
Indeed pullback under \(g\) identifies the framed fibre at \(h\) with the fibre at \(hk_g\), where \(k_g\) is a level lift. If their degree-one generators differ by \(\lambda\), the section value of degree \(n\) changes by \(\lambda^n\). Comparing the two vacuum readings gives \(c'_\sigma(C(h))=c'_\sigma(C(hk_g))\lambda^n\): the finite vacuum multiplier is \(\mu_\sigma(k_g)\) and \(\mu_\sigma(hk_g)=\mu_\sigma(h)\mu_\sigma(k_g)\). One can compare at a nonzero value of the generator since its mean squared norm is one. The generator factors thus cancel in the framed pullback. The remaining level coefficient is the embedded determinant character in \ref{B8.55}. The common factor \(a_j\) multiplies both evaluations. Injectivity of extremal evaluation proves the level law on vectors.
\leanref{GlobalTheory.\allowbreak{}GalOne.\allowbreak{}Equivar.\allowbreak{}eqAt\_\allowbreak{}of\_\allowbreak{}frAtEmb\_\allowbreak{}ofFive}{Development.\allowbreak{}AbelianVarietyTheory.\allowbreak{}Proofs.\allowbreak{}DivisorProofs.\allowbreak{}CoefMapNormEquivariantBasisOfFive}{1244--1323}
\leanref{GlobalTheory.\allowbreak{}CompPhase.\allowbreak{}exists\_\allowbreak{}isCompPhase\_\allowbreak{}five}{Development.\allowbreak{}AbelianVarietyTheory.\allowbreak{}Proofs.\allowbreak{}DivisorProofs.\allowbreak{}CentralShiftPhaseConjReadingsFiveClauses}{38--46}

Set \(Z_x=\sum_\nu\sigma(a_\nu)J'[f_\nu]\). Apply \(e_\sigma\) to the evaluation identity. By \ref{B5.07y} its analytic side is the evaluation of \(\kappa S_\sigma j_\sigma x\). The level law and injectivity give the first equality below. Applying the same level law directly to \eqref{eq:B5-rational-lift-formula} gives the second:
\[
 \kappa S_\sigma j_\sigma x
       =\sigma(\delta(z))\,e_\sigma\rho_\sigma(z)Z_x,\qquad
 J'\theta x=\sigma(\zeta_0\delta(z_0))\,e_\sigma\rho_\sigma(z_0)Z_x.
\]
Both translations commute with the projector. Multiplying the first equality by \(\rho_\sigma(z_0z^{-1})\) and substituting in the second proves \eqref{eq:B5-conjugate-lift} with
\[
 g_\sigma=\rho_\sigma(z_0z^{-1}),\qquad
 \xi_\sigma=\sigma(\zeta_0)\,
             \frac{\sigma(\delta(z_0))}{\sigma(\delta(z))}\,p_\sigma c_0 .
\]
All four scalar factors are independent of \(x\). The two determinant values have order dividing \(|\Gamma|\), and \(\zeta_0\) has finite order. Every complex embedding sends these values to the unit circle. The norm computations for \(p_\sigma\) and \(c_0\) therefore give \(|\xi_\sigma|=1\). Right translation preserves the quotient measure, so \(g_\sigma\) is isometric. At the distinguished embedding the defining base equality itself allows \(U_0=g_0=1\) and \(\xi_0=1\).

For the metric comparison, the squared norm of \(J_\sigma[f]\) is, by \eqref{eq:B5-petersson-evaluation}, the integral of \(|c_\sigma\mu_\sigma R_\sigma(f)|^2\) times \(a_j^2\prod_\rho\dim \tau_\rho\), where \(\tau_\rho\) is the lowest compact type. The core Petersson norm on the same section also includes the factor \(a_j\) on lengths. Thus on \(d'\) independent columns the squared Gram determinants on both sides of the core-to-realization comparison contain \(a_j^{2d'}\). Their ratio is the measure and lowest-type ratio of \ref{B8.25d}. In the height calculation the rational target carries this realized metric, so the equality just proved gives \(\|J_\sigma\theta x\|^2=\|S_\sigma j_\sigma x\|^2\) with the metric used for the singular values in \ref{B5.07y}.
\leanref{Formalism.\allowbreak{}L2.\allowbreak{}Pj.\allowbreak{}l2C}{Formalism.\allowbreak{}LocalToGlobal.\allowbreak{}ProjectiveLattices.\allowbreak{}L2DefiniteCuspForms}{157--160}
\leanref{GlobalTheory.\allowbreak{}GalLift.\allowbreak{}Assembly.\allowbreak{}conj\_\allowbreak{}abstract}{Development.\allowbreak{}AbelianVarietyTheory.\allowbreak{}Proofs.\allowbreak{}DivisorProofs.\allowbreak{}GalLiftOfOneMap}{84--103}
\leanref{GlobalTheory.\allowbreak{}ConjRatLiftFive.\allowbreak{}conj\_\allowbreak{}lift5}{Development.\allowbreak{}AbelianVarietyTheory.\allowbreak{}Proofs.\allowbreak{}DivisorProofs.\allowbreak{}ConjugatesOfRationalThetaLiftFromCharacterisatio}{71--139}
\end{proof}

\subsection*{Proof of \ref{B5.07g}}\label{prf:B5.07g}\phantomsection
\begin{proof}
We first describe how to find \(Z_{jp}\) over \(K\). The Hecke identification in the primitive hypothesis of \ref{B8.38} has the following content. The base lift intertwines the operators on class forms with a commuting family on its target. In that family \(W_{jp}\) is the sum of the joint generalized eigenspaces with the prescribed layer type of \(p\). A layer type prescribes a Hecke polynomial at each good split prime and a scalar for each spherical double coset at each good inert prime. Here the good primes are taken off the marked places of the reference tuple.

Decompose the rational class space into the simultaneous generalized eigenspaces of its split and inert operators. These operators commute and split over the coefficient field of the selected realization. The rational subspace of forms whose base images belong to \(\operatorname{im}S_0^*\) is stable under them. Indeed, the kernel of the lift is stable under the intertwining operators, and their adjoints belong to the same family. Each component of \(Bx\) therefore still has base image in \(\operatorname{im}S_0^*\). Intertwining sends each component into the target component with the same eigenvalues at the base embedding. On \(\operatorname{im}S_0^*=(\ker S_0)^\perp\) the lift is injective. Thus \(S_0j_0Bx\in W_{jp}\) precisely when \(Bx\) lies in the sum of the rational components with the prescribed restrictions. This gives \(Z_{jp}\) by taking an inverse image over \(K\). Its basis can be found by linear equations: for an operator \(A\) and eigenvalue \(a\), the generalized eigenspace is \(\ker(A-aI)^{\dim_K X}\). In an intersection of such kernels one retains a new equation whenever it reduces the dimension. The resulting finite intersection gives the same space.
\leanref{Formalism.\allowbreak{}IndefR.\allowbreak{}Cx.\allowbreak{}HeckeIntertwiningLifts}{Formalism.\allowbreak{}LocalToGlobal.\allowbreak{}IndefiniteReal.\allowbreak{}PrimitiveStatementsWithPassage}{48--57}
\leanref{Formalism.\allowbreak{}IndefR.\allowbreak{}Cx.\allowbreak{}part\_\allowbreak{}iff\_\allowbreak{}lifts3}{Formalism.\allowbreak{}LocalToGlobal.\allowbreak{}IndefiniteReal.\allowbreak{}LiftCoordinatesInGoodPartFamily}{114--128}

Fix an embedding \(\tau\) agreeing with the base embedding on \(K_d\). Fix \(p\) before choosing \(x\). The layer transfer assigns to \(p\) one list of split Hecke polynomials and inert double-coset values at the good places. By the reference Hecke identification of \ref{B8.38}, every summand contributing a rational vector to \(Z_{jp}\) has this restriction. Thus it suffices to recover the spherical generator values from this one list. The recovery works on the entire sum of the rational components. At a split prime with residue cardinality \(q\), the coefficient of degree \(h\) of the prescribed polynomial is
\((-1)^h q^{h(h-1)/2}a_h\), where \(a_h\) is the eigenvalue of the \(h\)-th elementary double coset. Thus the polynomial determines every \(a_h\). At an inert prime the type specifies the double-coset values individually. At the base embedding the spherical operators are normal: their adjoints are given by inversion in the commutative spherical algebra. Hence their generalized eigenspaces are eigenspaces. Each resulting equation on a class form is an equality of finite sums of translates in its rational coefficient frame. Applying \(\tau\) gives the equation with eigenvalue \(\tau(a_h)\), or the corresponding embedded inert value.
\leanref{GlobalTheory.\allowbreak{}ZPartEig.\allowbreak{}points\_\allowbreak{}eig}{Development.\allowbreak{}AbelianVarietyTheory.\allowbreak{}Proofs.\allowbreak{}TwistedOperators.\allowbreak{}ThetaLiftPointsPartEigenformsCoefficientField}{198--208}
\leanref{GlobalTheory.\allowbreak{}ConjScalar.\allowbreak{}eigF\_\allowbreak{}conj}{Development.\allowbreak{}AbelianVarietyTheory.\allowbreak{}Proofs.\allowbreak{}TwistedOperators.\allowbreak{}SemilinearInjectiveOfIndep}{53--61}

We show that the vectors \(z=S_\tau j_\tau Bx\) with \(x\in Z_{jp}\) have one target spherical system whenever one of them is nonzero. At a good place the test is the self-dual lattice vacuum, the splitting characters are unramified, and the retained projector at the marked places commutes with these spherical operators. Matching the actions of two operators on the vacuum gives their intertwining relation on the lift: insert the finite coset sum in the theta integral and change the source integration variable by each representative. This replaces the source double coset by its inverse.

At an inert place the two lattices are unimodular over an unramified quadratic extension, and the standard additive character has conductor the integers of the fixed field. The residue characteristic is odd since the host decomposition condition excludes conjugation from the stabilizer of a prime above \(2\). The integral Heisenberg translations fix the vacuum and its translates span the lattice model. A target spherical operator applied to the vacuum is fixed by the target integral group. By \ref{B8.14} it is a finite linear combination of source translates of the vacuum. It is fixed by the source integral group as well, since the two actions commute. Averaging this combination over the source integral group gives a source spherical operator with the same value on the vacuum. Its inverse-double-coset action on \(j_\tau Bx\) has the common source eigenvalue just computed. This determines the target eigenvalue for all \(x\).
\leanref{GlobalTheory.\allowbreak{}InertLoc.\allowbreak{}exists\_\allowbreak{}locPair}{Development.\allowbreak{}AbelianVarietyTheory.\allowbreak{}Proofs.\allowbreak{}SatakeParameters.\allowbreak{}ExistsUnramifiedLocalPairOfInertWeilDatum}{30--100}
\leanref{GlobalTheory.\allowbreak{}ConjScalar.\allowbreak{}LiftPart.\allowbreak{}loc\_\allowbreak{}inert}{Development.\allowbreak{}AbelianVarietyTheory.\allowbreak{}Proofs.\allowbreak{}TwistedOperators.\allowbreak{}ConjugateDatumLiftJointEigenspaceIntoOnePart}{135--151}

At a split place write \(n=r+t\) for the source rank and \(r\) for the target rank. Let \(q\) be the residue cardinality, and let \(s\) and \(w\) be the source and target splitting characters evaluated on the idèle which is a uniformizer at the chosen prime and one elsewhere. Write \(T_h\) for the target elementary double coset with \(h\) uniformizer entries, and \(\alpha_k\) for the embedded source eigenvalue. The split vacuum calculation gives, for \(0\le k\le r\),
\[
 \alpha_k z=\sum_{h=0}^k A_{kh}T_hz,\qquad
 A_{kh}=q^{(r-h)(k-h)}{t\brack k-h}_q
             (wq^{r/2})^{-k}(sq^{n/2})^h .
\]
Here \({t\brack i}_q\) counts the \(i\)-planes in the complementary residue space of dimension \(t\), with value zero for \(i>t\). For a fixed intersection of dimension \(h\) with the target block, choose the image of dimension \(k-h\) in the complement. Its lifts are graphs of maps into the target block modulo the intersection, giving \(q^{(r-h)(k-h)}\) possibilities. The two determinant multipliers in the split Weil action give the last two factors. Since \(A_{kk}=(sq^{n/2}/(wq^{r/2}))^k\ne0\), induction on \(k\) gives \(T_kz=\beta_kz\) with \(\beta_k\) independent of \(x\). The elementary cosets and the inverse determinant coset generate the split spherical algebra; the determinant coset is invertible. Together with the inert calculation this gives the desired target system \(c\).
\leanref{GlobalTheory.\allowbreak{}Hk3.\allowbreak{}SplitMatch.\allowbreak{}match\_\allowbreak{}split}{Development.\allowbreak{}AbelianVarietyTheory.\allowbreak{}Proofs.\allowbreak{}SatakeParameters.\allowbreak{}ThetaVacuumHeckeIdentitySplitPlace}{148--187}
\leanref{GlobalTheory.\allowbreak{}ConjScalar.\allowbreak{}LiftPart.\allowbreak{}loc\_\allowbreak{}split}{Development.\allowbreak{}AbelianVarietyTheory.\allowbreak{}Proofs.\allowbreak{}TwistedOperators.\allowbreak{}ConjugateDatumLiftJointEigenspaceIntoOnePart}{154--189}

Put \(u=j_\tau Bx\), and write \(S_\tau=eT\) as in \ref{B5.07y}. If some \(z\ne0\), the space \(Y_\tau(c)\) is nonzero. Lemma~\ref{B5.07y} gives \(TT^*z=l_cz\), with \(l_c>0\) independent of \(x\). Since \(ez=z\),
\[
 S_\tau S_\tau^*S_\tau u=eTT^*z=l_c S_\tau u.
\]
The vector \(S_\tau^*S_\tau u-l_cu\) belongs to \(\ker S_\tau\). It also belongs to its orthogonal complement by the hypothesis on \(u\). It is therefore zero. For every \(x'\in Z_j\), adjunction gives
\[
 \langle S_\tau j_\tau Bx,S_\tau j_\tau Bx'\rangle
       =l_c\langle j_\tau Bx,j_\tau Bx'\rangle .
\]
If every \(z\) vanishes, the same orthogonality implies \(j_\tau Bx=0\) for every \(x\in Z_{jp}\).
\leanref{GlobalTheory.\allowbreak{}ConjScalar.\allowbreak{}adjoint\_\allowbreak{}apply\_\allowbreak{}of\_\allowbreak{}orthogonal\_\allowbreak{}ker}{Development.\allowbreak{}AbelianVarietyTheory.\allowbreak{}Proofs.\allowbreak{}TwistedOperators.\allowbreak{}AdjointApplyOfOrthogonalKer}{17--30}

Use \(\tau=\sigma\) if its restriction to \(K_d\) agrees with the base embedding. The two unitary maps and the common unit phase in \eqref{eq:B5-conjugate-lift} preserve the target inner product. The source inner product is the class sum with matrix \(\sigma(H)\). Dividing both pairings by \(N\) therefore gives
\(\langle y_v(x),y_v(x')\rangle=l_cM_v(x,x')\).
For the other restriction to \(K_d\), take \(\tau=\bar\sigma\). On the same two rational vectors the target pairing at \(\sigma\) is the conjugate of the pairing at \(\bar\sigma\), by the conjugate evaluation in the reference realization of \ref{def:B8-reference-points}. The displayed class sum has the same conjugation law. Conjugating the equality preserves the real scalar \(l_c\).
\leanref{GlobalTheory.\allowbreak{}ConjScalar.\allowbreak{}readings\_\allowbreak{}scale\_\allowbreak{}core}{Development.\allowbreak{}AbelianVarietyTheory.\allowbreak{}Proofs.\allowbreak{}DivisorProofs.\allowbreak{}RationalLiftReadingsScaleClassSumEigenspaces}{21--59}
\leanref{GlobalTheory.\allowbreak{}ConjScalarCoord.\allowbreak{}liftMat\_\allowbreak{}isRatLift\_\allowbreak{}inner}{Development.\allowbreak{}AbelianVarietyTheory.\allowbreak{}Proofs.\allowbreak{}TwistedOperators.\allowbreak{}LiftWeightComplexPointEqTwist}{68--79}

The diagonal \(M_v(x,x)\) is the squared norm of \(j_\tau Bx\) divided by \(N\), for either choice of \(\tau\). If it vanishes on every member of a \(K\)-basis of \(Z_{jp}\), semilinearity gives \(j_\tau Bx=0\) throughout \(Z_{jp}\); both pairings then vanish against every \(x'\). Otherwise the first basis member \(b\) with nonzero diagonal has \(M_v(b,b)>0\). Orthogonality to the kernel gives \(S_\tau j_\tau Bb\ne0\), so the computation with \(c\) applies. Evaluating at \((b,b)\) identifies the quotient in the statement with \(l_c\). It also proves uniqueness of the scalar in the paired row. The scalar characterization in the selected realization uses this unique value on an occurring part and the exponential of the fixed layer allowance on an empty part.
\leanref{GlobalTheory.\allowbreak{}LowOff.\allowbreak{}partAt\_\allowbreak{}iff}{Development.\allowbreak{}AbelianVarietyTheory.\allowbreak{}Proofs.\allowbreak{}DivisorProofs.\allowbreak{}LogLiftEigenvalueLowerBoundOtherPlaces}{38--43}
\leanref{Formalism.\allowbreak{}IndefR.\allowbreak{}Links.\allowbreak{}RealDat3.\allowbreak{}IsScal}{Formalism.\allowbreak{}LocalToGlobal.\allowbreak{}IndefiniteReal.\allowbreak{}ZeroPartEigenvalueConvention}{217--224}
\end{proof}

\subsection*{Proof of \ref{B5.07b}}\label{prf:B5.07b}\phantomsection
\begin{proof}
At the scale of \ref{B5.07} the sixth entry is
\(L_{6,j}=H_j^+-C_{\rm sz}fn_V^2L'\). The exponent of \ref{B5.07y} is \(D_j=H_j^+-8fn_V^2L'\). Since \(C_{\rm sz}\ge8\) and \(L'\ge0\), \(L_{6,j}\le D_j\). At the distinguished place take \(0\ne w\in W_{jp}\) if this space occurs. Its assigned scalar \(\lambda>0\) satisfies \(S_0S_0^*w=\lambda w\), so \(\|S_0^*w\|^2=\lambda\|w\|^2\). As \(S_0^*w\perp\ker S_0\), the source estimate of \ref{B5.07y}, with the stipulated reference primes and sizes, gives
\[
 e^{D_j}\lambda\|w\|^2\le \|S_0S_0^*w\|^2
                         =\lambda^2\|w\|^2.
\]
Thus \(\lambda\ge e^{D_j}\). On an empty space the assigned value is \(e^{L_{6,j}}\).
\leanref{GlobalTheory.\allowbreak{}LowRow.\allowbreak{}l6m\_\allowbreak{}le\_\allowbreak{}log\_\allowbreak{}lam3\_\allowbreak{}base\_\allowbreak{}place}{Development.\allowbreak{}AbelianVarietyTheory.\allowbreak{}Proofs.\allowbreak{}DivisorProofs.\allowbreak{}LogLiftEigenvalueLowerBoundBasePlace}{13--21}

Write \(K_0=\{a\in K:\tau(\bar a)=\overline{\tau(a)}\text{ for every }\tau:K\hookrightarrow\C\}\). Let \(H\) be the class Gram matrix in the chosen class basis, \(V\) the inclusion of the retained block, and \(H_1=V^*HV\). We check the metric hypotheses on the primitive reference row. Definition~\ref{def:p4c6-definite-blocks} gives \(H\) over \(K_0\), positive at every infinite place, and group matrices over \(K_0\). The character projector and the level-character average have entries over \(K_0\): their character values are sums of roots of unity, on which every embedding commutes with conjugation. The image of their product therefore has a basis over \(K_0\), obtained by taking independent columns. The block inclusion uses such a basis. At each \(\tau\), an injective minor of \(V\) stays nonzero, and
\[
 \tau(H_1)=\tau(V)^*\,\tau(H)\,\tau(V)>0.
\]
For the conjugate of a representative embedding this follows as well by conjugating its positive Gram matrix.

For clarity, the quotient-frame assertion on this row can also be checked before transporting its orthogonal section. The ordinary and inert operators of the primitive reference row have matrices over \(K_0\), commute, split over \(K\), and have their adjoints in the specified families for the class pairing. Hence their complex realizations are commuting normal operators at each embedding. If \(T'\) is the specified adjoint of \(T\), and their eigenvalues on a common \(K\)-eigenvector are \(\alpha,\alpha'\), adjunction at the distinguished embedding gives \(\alpha'=\bar\alpha\). At any other embedding it gives \(\tau(\alpha')=\overline{\tau(\alpha)}\), since the embedded vector has positive norm. Thus \(\alpha\in K_0\). The selected sum of joint eigenspaces in block coordinates consequently has a basis over \(K_0\). Indeed each eigenspace is an intersection of kernels over \(K_0\), and taking a sum or passing to the block coordinates preserves this property. Finitely many kernel equations suffice, by dimension. The reference coordinate basis is chosen over \(K_0\) whenever such a basis exists.

In these coordinates the section matrix \(E_1\) has entries in \(K_0\). Put \(G_1=E_1^*H_1E_1\). The orthogonal quotient projection satisfies \(P=G_1^{-1}E_1^*H_1\), because subtracting \(E_1Px\) from \(x\) must give a vector orthogonal to each column of \(E_1\). In particular \(PE_1=1\). For any embedding put \(A=\tau(H_1)\), \(C=\tau(E_1)\) and \(G'=\tau(G_1)\). Since the entries of \(H_1,E_1\) lie in \(K_0\), one has \(G'=C^*AC\). The columns of \(C\) are independent by a nonzero minor, and \(A>0\), so \(G'>0\). The embedded quotient map is \(G'^{-1}C^*A\); its kernel is the \(A\)-orthogonal complement of the columns of \(C\). A lift of \(y\) is uniquely \(Cy+z\) with \(z\) in this kernel, and its squared norm is \(y^*G'y+z^*Az\). The infimum is therefore \(y^*G'y\) at this embedding. Its orthonormal coordinate frame \(X_v\) therefore satisfies
\((X_v^{-1})^*X_v^{-1}=\tau(G_1)\) at the representative of \(v\). This proves the asserted frame compatibility for the selected realization. The ordinary and inert rationality and adjoint conditions here are precisely the two conditions on the primitive Hecke action used with the reference eigenspace selection in \ref{B8.08}.
\leanref{Formalism.\allowbreak{}L2.\allowbreak{}DIV.\allowbreak{}BlockDat.\allowbreak{}HV\_\allowbreak{}pos}{Formalism.\allowbreak{}LocalToGlobal.\allowbreak{}Divisibility.\allowbreak{}RowLatLemmas}{38--40}
\leanref{Formalism.\allowbreak{}IndefR.\allowbreak{}Cx.\allowbreak{}okOp\_\allowbreak{}opAll}{Formalism.\allowbreak{}LocalToGlobal.\allowbreak{}IndefiniteReal.\allowbreak{}MajAtEqZeroOfMemQpOfNe}{18--36}
\leanref{Formalism.\allowbreak{}IndefR.\allowbreak{}Cx.\allowbreak{}proviso\_\allowbreak{}of\_\allowbreak{}lift\_\allowbreak{}eigenvalues}{Formalism.\allowbreak{}LocalToGlobal.\allowbreak{}IndefiniteReal.\allowbreak{}MajAtEqZeroOfMemQpOfNe}{45--55}
\leanref{GlobalTheory.\allowbreak{}LowRow.\allowbreak{}baseLinks\_\allowbreak{}of\_\allowbreak{}prims}{Development.\allowbreak{}AbelianVarietyTheory.\allowbreak{}Proofs.\allowbreak{}DivisorProofs.\allowbreak{}SplitPartBaseLinksOfPrimitiveRows}{15--29}

Here is the transport argument using the quotient-frame hypothesis directly. It applies in any of the rational coordinates in the statement. Write \(\tau'(a)=\overline{\tau(\bar a)}\), \(C=\tau(E_1)\), \(D=\tau'(E_1)\), \(A=\tau(H_1)\), \(G'=\tau(G_1)\), and \(P'=\tau(P)\). The identities over \(K\) give
\[
 P'C=1,\qquad D^*A=G'P',\qquad D^*AC=G'.
\]
In particular \(P'\) is surjective. Since \(A>0\), \(P'A^{-1}P'^*>0\). Its inverse is the quotient Gram matrix at this embedding: the lift
\[
 y\longmapsto A^{-1}P'^*(P'A^{-1}P'^*)^{-1}y
\]
is a right inverse and is \(A\)-orthogonal to \(\ker P'\). Comparing with the quotient frame gives
\[
 G'=(P'A^{-1}P'^*)^{-1}.
\]
For the other embedding representing the same place this identity follows by complex conjugation. In this calculation the adjoint is taken on the complex quotient map \(P'\).

Taking the adjoint of \(D^*A=G'P'\) gives \(D=A^{-1}P'^*G'\), since \(A\) and \(G'\) are Hermitian. It follows that \(P'D=1\) and \(D^*AD=G'\). Apply this result at \(\tau'\). The matrices of \(H_1,G_1\) there are still \(A,G'\), since they have entries in \(K_0\), and twisting \(\tau'\) again gives \(\tau\). Thus \(C^*AC=G'\). Together with \(D^*AC=G'\) and its adjoint this gives
\[
 (C-D)^*A(C-D)=0.
\]
Positivity on each column implies \(C=D\). This proves \(\tau(E_1^*)=\tau(E_1)^*\) at every embedding. The identity \(P=G_1^{-1}E_1^*H_1\) then gives \(\tau(P^*)=\tau(P)^*\) as well. For the full inclusion \(B=VE_1\) we obtain \(\tau(B^*)=\tau(B)^*\).
\leanref{Formalism.\allowbreak{}L2.\allowbreak{}DIV.\allowbreak{}BlockDat.\allowbreak{}XQ\_\allowbreak{}quot}{Formalism.\allowbreak{}LocalToGlobal.\allowbreak{}Divisibility.\allowbreak{}RowLatLemmas}{67--69}
\leanref{GlobalTheory.\allowbreak{}QuotRational.\allowbreak{}entries\_\allowbreak{}mem\_\allowbreak{}cmMax}{Development.\allowbreak{}AbelianVarietyTheory.\allowbreak{}Proofs.\allowbreak{}DivisorProofs.\allowbreak{}MapConjTransposeTwist}{70--113}
\leanref{GlobalTheory.\allowbreak{}QuotRational.\allowbreak{}fE\_\allowbreak{}mem\_\allowbreak{}cmMax}{Development.\allowbreak{}AbelianVarietyTheory.\allowbreak{}Proofs.\allowbreak{}DivisorProofs.\allowbreak{}MapConjTransposeTwist}{114--120}

Fix \(\sigma\) over the distinguished embedding of \(K_d\). The complex class space is \(\C\otimes_{K,\sigma}X\). Indeed, at each representative the allowed values are the image of the average of its finite stabilizer action, with the level character included. This rational idempotent extends to the corresponding complex idempotent, and the values at the representatives determine the class form. Now choose a \(K\)-basis of \(\ker\theta\) and extend it to a basis of \(X\). The images under \(\theta\) of the additional basis vectors are independent. Their coordinate columns remain independent after embedding, by a nonzero minor, and the realized target frame is independent. Applying \ref{B5.07f} to this basis proves that \(\ker S_\sigma\) is spanned by \(j_\sigma(\ker\theta)\), including when the rank is zero or full.
\leanref{GlobalTheory.\allowbreak{}ConjSpan.\allowbreak{}span\_\allowbreak{}conjForms}{Development.\allowbreak{}AbelianVarietyTheory.\allowbreak{}Proofs.\allowbreak{}DivisorProofs.\allowbreak{}SpanEqTopOfTie}{126--134}
\leanref{GlobalTheory.\allowbreak{}ker\_\allowbreak{}le\_\allowbreak{}span\_\allowbreak{}image\_\allowbreak{}ker\_\allowbreak{}of\_\allowbreak{}semilinear}{Development.\allowbreak{}AbelianVarietyTheory.\allowbreak{}Proofs.\allowbreak{}DivisorProofs.\allowbreak{}KerLeSpanImageKerOfSemilinear}{6--56}

Write \(C_\sigma(z,z')=\overline{\sigma(z)}^{\,t}\sigma(H)\sigma(z')\). This equals the realized class pairing by the stabilizer-weighted formula. For \(u\in\ker\theta\), \ref{B5.07f} gives \(S_0j_0u=0\). The reference layer subspace of rational class forms is the inverse image, under class realization, of the adjoint range of the phase-normalized lift. In block coordinates one takes its inverse image under the block inclusion \(V\); the columns of \(E_1\) are a basis of that inverse image. Thus \(Bx=VE_1x\) belongs to the reference layer subspace for every \(x\). To read its defining range in the base realization, multiplication of a lift by a unit phase multiplies its adjoint by the conjugate phase, and transport by the unitary or conjugate-unitary class identifications carries the adjoint range onto the transported adjoint range. Undoing these identifications gives \(j_0Bx\in\operatorname{im}S_0^*\).
\leanref{GlobalTheory.\allowbreak{}ConjScalarOff.\allowbreak{}jV\_\allowbreak{}mem\_\allowbreak{}range}{Development.\allowbreak{}AbelianVarietyTheory.\allowbreak{}Proofs.\allowbreak{}TwistedOperators.\allowbreak{}LiftScalarIdentityOffFixedPlaceBaseLevel}{32--36} Since \(S_0j_0u=0\), we have \(0=\langle j_0Bx,j_0u\rangle=C_{\sigma_0}(Bx,u)=\overline{\sigma_0(x)}^{\,t}\sigma_0(B)^*\sigma_0(H)\sigma_0(u)\) for every \(x\in K^{d_j}\) and every \(u\in\ker\theta\). Using \(\sigma_0(B)^*=\sigma_0(B^*)\) and letting \(x\) run through the standard basis of \(K^{d_j}\) gives \(\sigma_0(B^*Hu)=0\), whence \(B^*Hu=0\) over \(K\) by injectivity of \(\sigma_0\). At any embedding \(\sigma:K\hookrightarrow\C\) over the distinguished embedding of \(K_d\), applying \(\sigma\) and \(\sigma(B^*)=\sigma(B)^*\) yields \(\langle j_\sigma Bx,j_\sigma u\rangle=C_\sigma(Bx,u)=\overline{\sigma(x)}^{\,t}\sigma(B)^*\sigma(H)\sigma(u)=\overline{\sigma(x)}^{\,t}\sigma(B^*Hu)=0\) for all \(u\in\ker\theta\). Because \(\ker S_\sigma\) is spanned over \(\C\) by \(j_\sigma(\ker\theta)\), it follows that \(j_\sigma Bx\perp\ker S_\sigma\). Applying Lemma~\ref{B5.07y} at \(\sigma\), followed by the isometries of Lemma~\ref{B5.07f}, gives
\begin{equation}\label{eq:B5-rational-source-norm}
 e^{D_j}\Re C_\sigma(Bx,Bx)\le\|S_\sigma j_\sigma Bx\|^2
                               =\|J_\sigma\theta Bx\|^2.
\end{equation}
If \(\sigma\) induces the other embedding of \(K_d\), then \(\bar\sigma\) agrees with the distinguished embedding on \(K_d\), so \(j_{\bar\sigma}Bx\perp\ker S_{\bar\sigma}\) and \eqref{eq:B5-rational-source-norm} holds at \(\bar\sigma\). The two target pairings on rational columns are conjugate by Definition~\ref{def:B8-reference-points}, and the formula for \(C_\sigma\) gives the same conjugation law on the source. The diagonal inequality thus holds at \(\sigma\).
\leanref{GlobalTheory.\allowbreak{}LowForm.\allowbreak{}lower\_\allowbreak{}overK\_\allowbreak{}six}{Development.\allowbreak{}AbelianVarietyTheory.\allowbreak{}Proofs.\allowbreak{}DivisorProofs.\allowbreak{}RationalLiftLowerBoundOffBasePlace}{208--237}

For \(v\ne v_0\), divide by the order \(N\) of the reference level subgroup. In the notation of \ref{B5.07g} the source form is \(M_v=N^{-1}C_\sigma(B\,\cdot,B\,\cdot)\), and \(y_v(x)=N^{-1/2}J_\sigma\theta Bx\). Hence
\begin{equation}\label{eq:B5-off-base-gram}
 e^{D_j}\Re M_v(x,x)\le\|y_v(x)\|^2 .
\end{equation}
We have checked the kernel-orthogonality hypothesis of \ref{B5.07g} at every required embedding. On an occurring part take its basis vector \(b\) with \(M_v(b,b)>0\). The paired-row identity gives \(\|y_v(b)\|^2=\lambda_{jpv}M_v(b,b)\), so division gives \(\lambda_{jpv}\ge e^{D_j}\). On an empty part the specified scalar is \(e^{L_{6,j}}\). Summing \(L_{6,j}\le\log\lambda_{jpv}\) at fixed \(p,v\) and taking the minimum over the nonempty finite set \(I\) proves the second inequality. Replacing the zero-based layer index by \(j-1\) identifies the plus-hook sum and its subtraction term with \(\sum_jH_{X,j}\) in \ref{B5.08}.
\leanref{GlobalTheory.\allowbreak{}LowForm.\allowbreak{}lowForm\_\allowbreak{}of\_\allowbreak{}emb\_\allowbreak{}six}{Development.\allowbreak{}AbelianVarietyTheory.\allowbreak{}Proofs.\allowbreak{}DivisorProofs.\allowbreak{}RationalLiftLowerBoundOffBasePlace}{238--261}
\leanref{GlobalTheory.\allowbreak{}LowOff.\allowbreak{}l6m\_\allowbreak{}le\_\allowbreak{}log\_\allowbreak{}lam3\_\allowbreak{}other\_\allowbreak{}places}{Development.\allowbreak{}AbelianVarietyTheory.\allowbreak{}Proofs.\allowbreak{}DivisorProofs.\allowbreak{}LogLiftEigenvalueLowerBoundOtherPlaces}{46--54}
\leanref{GlobalTheory.\allowbreak{}LowMainTerm.\allowbreak{}sum\_\allowbreak{}logHook\_\allowbreak{}eq}{Development.\allowbreak{}AbelianVarietyTheory.\allowbreak{}Proofs.\allowbreak{}DivisorProofs.\allowbreak{}HookLogSumSubSizeTermLeLowerBound}{78--90}
\end{proof}

\subsection*{Proof of \ref{B5.07x}}\label{prf:B5.07x}\phantomsection
\begin{proof}
Logarithmic Chowla--Selberg gives the class mean of \(\log(4\pi\Im z_C|\eta(z_C)|^4)\) as \(\gamma+\log(2\pi/\sqrt d)-D_1\), with \(z_C\) the ideal-class points and \(\eta\) the Dedekind eta function \cite{selberg-1967}. In discriminant-one normalization the argument of the logarithm is \(\operatorname{area}/\pi\). Take a common extension containing the Hilbert class field of \(K_d\). By the elliptic CM theorem of \ref{B1.15a}, embeddings fixing \(K_d\) traverse the ideal classes with equal multiplicities. Embeddings inducing its other embedding are obtained by conjugation, which preserves area. For a model of unit discriminant \(u\), scaling by \(c\) with \(c^{12}=u\) gives discriminant one and adds \(\frac16\log|\sigma u|\) to the log area. Its mean is zero by the product formula. Each embedding of \(M\) extends with the same multiplicity, so the resulting normalized mean descends to \(M\).

Let \(\mathfrak b\bar{\mathfrak b}=\mathfrak D_{E/K_d}\) for the host ideal \(\mathfrak b\). The differential ideal is the extension of \(\tau(\mathfrak b^{-1}\mathfrak D_{E/K_d})\). With \(B=N_{E/\Q}(\mathfrak b^{-1}\mathfrak D_{E/K_d})\), one has \(B^2=N_{E/\Q}\mathfrak D_{E/K_d}\) and
\[
 \frac{\log|N_{M/\Q}(s_\tau)|}{[M:\Q]}=\frac{\log B}{[E:\Q]} .
\]
The identity \(|\operatorname{disc}E|=d^{[E:K_d]}B^2\) shows that subtracting this doubled log norm replaces \(\sqrt d\) by \(\mathrm{rd}E=2v_E\) in the class mean. Multiplication by \(-\sum_i\nu_{uj\rho i}\) gives the period term in \ref{B5.07}.
\end{proof}

\subsection*{Proof of \ref{B5.08d}}\label{prf:B5.08d}\phantomsection
\begin{proof}
By Lemma~\ref{B4.02b}, \(e_q=B_\rho+(q+1/2)h_0\), where \(h_0>0\) and \(\eta=mh_0\). In particular
\[
 e_{m-1}=B_\rho+\eta-h_0/2 .
\]
At a mark, \(B_\rho=b_s+\lfloor b'/x\rfloor\le 3b'\) by \eqref{eq:p3c5-size-clauses} and \(x\ge1\). At a thin position \(B_\rho=\lceil w''\rceil\le w''+1\). At a thick position \(B_\rho=\lceil b\rceil+d_\rho\), with \(0\le d_\rho<|\mu(E)|\), including the designated position at which the weight is adjusted.

Let \(h=|\mu(E)|\). Since \(\Q(\zeta_h)\subseteq E\), Euler's totient satisfies \(\varphi(h)\le [E:\Q]=2f\). The prime-power formula gives \(\varphi(h)^2/h\ge1/2\). Indeed, the factor at \(p^a\) is \(p^{a-2}(p-1)^2\), which is at least one except at \(2^1\), where it is \(1/2\). Thus \(h\le8f^2\). By the rank bounds of \ref{B8.59} and the size threshold of \ref{B8.35a}, eventually \(f\le x^3\) and \(h_0\ge x^{1300}\). Hence \(h\le h_0/2\). The three upper bounds for \(e_{m-1}\) are consequently \(3b'+\eta\), \(w''+\eta\), and \(b+1+\eta\), respectively. By \ref{B4.02b}, \(b'/b=x\exp(-x^{\varepsilon/3})\), \(w''/b\le\exp(-x^{\varepsilon/3})\), and \(\eta\) grows at most polynomially in \(x\) at fixed family parameters. Since \(b=\exp(x^\varepsilon)\), each bound is at most \(2b\) eventually.
\end{proof}

\subsection*{Proof of \ref{B5.08b}}\label{prf:B5.08b}\phantomsection
\begin{proof}
Take the integer spacing and \(h_0\) of \ref{B4.02b}. By \ref{B8.59} the ranks and \(n\) are polynomially bounded, and by \ref{B4.02b} so are \(S_c,\eta\). For fixed \(d>0\), \(d\log x=o(x^{\varepsilon/3})\) and \(x^{\varepsilon/3}+d\log x=o(L)\). Thus \(X'\), \(b\), and \(w''\) eventually exceed each fixed power of \(x\). This gives the four rank and \(\eta\) bounds with right side \(b\). Since \(N,mk\le X'\), the identities \(b'=a/X'\), \(w''=b/(nX')\) give the other three product bounds \(b'N,mkb'\le a\) and \(Nw''\le b\). The inequalities \(kS_c,B\le w''\) follow also from \ref{B8.35a}.

The finite row group has \(2f\) general linear factors of rank \(k+m-1\), over the residue fields at the two split type primes. Their orders are at most \(g_i^{(k+m-1)^2}\), by counting all matrices. Hence \(0\le m\log M\le2mf(k+m-1)^2\log\max_i g_i\), which is polynomially bounded by \ref{B8.35a}. The formula for \(N_c\) is nonnegative.

By \ref{B8.35a}, \(\log r_d,\log x,|D_1|\le L\). From \ref{B5.07},
\[
 |\log\Omega|\le \gamma+\log(2\pi)+\log r_d+L\le L_a
\]
as soon as \(x\ge16\). Here \(L_a=\log x+L+\log r_d\) and \(r_d\ge1\). Also \(\log n_V\le L\) eventually, so
\(L_p\le(3+2C_{YL})L+2\). For example \(C_0\ge\max\{28,C_1+1,(5+2C_{YL})C_h\}\) suffices for the constant bounds.

Since \(s_j\le N\le w''\), all three lower endpoints of \ref{B3.47} give \(w_{j\rho}\ge s_j\). At each position \(w_{j\rho}=e_{j-1}(\rho)\) by \ref{B4.02b}. Lemma~\ref{B5.08d} gives the upper bound, including the root-of-unity adjustment at the designated thick position. Since \(S_c\le b\) and \(x\ge3+|C_c|\) eventually, \(w_{j\rho}+S_w\le a\). At a mark the first box of \ref{B5.22a} gives \(a_s\le a+b\). The field \(E\) contains the quadratic field of discriminant \(-d\), so the discriminant identity in a tower gives \(r_d\ge\sqrt d\ge2\). Consequently \(a_s+w_{j\rho}\le a+3b\le A_*\) for \(x\ge3\).

Positivity and interlacing are in \ref{B4.02a}. A tail block of \ref{B4.02b} has \(t+1\) copies of a number at most \((j-1-q)h_0\le\eta\). This proves the stated sum bound, including \(j=1\). Off the marks the bottom endpoint is \(-1/2\) and the shifted top endpoint is at most \(2S_c\), by \ref{B4.02b}. At a mark the two box weights are positive, \(v_{k-2}=0\), and \(v_0\le S_c\). These give every endpoint condition with the specified \(S_w\).
\end{proof}

\subsection*{Proof of \ref{B5.08c}}\label{prf:B5.08c}\phantomsection
\begin{proof}
Write \(w=w_{j\rho}\), \(W=w-s_j\), \(d_j=(j-1)t\), and \(p=k+j-1\). Both \(s_j\) and \(d_j\) belong to \([0,N]\), and \(W\) is a nonnegative integer. Remove the rank shifts from the two original parameters by setting
\[
 c_i=U_i-(k-1)/2+i\quad(i<k),\qquad
 d_i'=v_i-(k-2)/2+i\quad(i<k-1).
\]
These are integers since \(k\) is even. By \ref{B4.02a} the head of the column and the first \(k-1\) hook rows are exactly
\[
 y_i:=\nu_i=W+c_i,\qquad r_i:=\alpha_i=W+d_i'+d_j .
\]
The shifts differ by half a unit in the original coordinates. Hence \(U_i>v_i>U_{i+1}\) gives \(c_i\ge d_i'\ge c_{i+1}\).

We use the following factorial bound for these particular \(r_i\):
\[
 \log H_{j\rho}\ge\sum_{i<p}\log(\alpha_i!)\ge\sum_{i<k-1}\varphi(r_i).
\]
Indeed, the hook of a cell at column \(h\), numbered from zero, has length at least \(\alpha_i-h\). The product of these lower bounds along row \(i\) is \(\alpha_i!\). To identify the cell product with the gamma quotient in \ref{B5.07}, delete the first column of a diagram with \(p\) positive rows. The cell product and the gamma quotient both lose the factor \(\prod_{i<p}(\alpha_i+p-1-i)\). The pairwise differences in the denominator are unchanged. Removing a trailing zero row leaves the quotient unchanged, since the new denominator factors cancel the additional numerator factors. Induction from the empty diagram proves the identity. Finally \(e^r\ge r^r/r!\) gives \(\varphi(r)\le\log(r!)\) for each positive integer \(r\), with equality at zero under the convention \(\varphi(0)=0\). The remaining \(j\) factorials have nonnegative logs.

For nonnegative integers \(y,r\) with \(y\le A_*\) we shall use
\[
 \varphi(y)-\varphi(r)\le |y-r|L_a .
\]
For \(y,r\ge1\), integrate \(\log s\) when \(r\le y\), and use monotonicity when \(r>y\). For \(r=0<y\), use \(\varphi(y)\le yL_a\). For \(y=0<r\), use \(-\varphi(r)\le1\le rL_a\). The tail degrees have sum at most \(B\le w''\le b\le A_*\), so their entire contribution to the \(\varphi\)-sum is at most \(BL_a\).

Suppose first that \(\rho\) is unmarked. The endpoint bounds and interlacing give \(0\le c_{k-1}\le c_i\le c_0\le S_w\). Pair \(y_i\) with \(r_i\) for \(0\le i<k-1\). Since \(W+S_w\le A_*\),
\[
 \sum_{i<k-1}\{\varphi(y_i)-\varphi(r_i)\}
 \le L_a\sum_{i<k-1}(c_i-d_i'+d_j)
 \le [S_w+(k-1)d_j]L_a .
\]
The last inequality uses \(\sum_{i<k-1}(c_i-d_i')\le c_0-c_{k-1}\).
The remaining head entry contributes at most \((W+S_w)L_a\). Thus
\[
 e_{j\rho}\le[W+2S_w+(k-1)d_j+B]L_a .
\]

At a mark the endpoints are \(c_0=a_s\), \(c_{k-1}=-b_s\). The endpoint conditions on \(v\) give
\(0\le d_{k-2}'\le d_0'\le S_w\). Pair \(y_i\) with \(r_i\) for \(1\le i\le k-2\). Now \(d_{i-1}'\ge c_i\ge d_i'\), so
\[
 \sum_{i=1}^{k-2}(c_i-d_i')\le d_0'-d_{k-2}'\le S_w .
\]
The paired cost is at most \([S_w+(k-2)d_j]L_a\). The first head entry contributes at most \(\varphi(a_s)+WL_a\), since \(a_s+W\le A_*\). The last is \(W-b_s\in[0,W]\), and contributes at most \(WL_a\). The unpaired hook term costs \(-\varphi(r_0)\le1\). Hence
\[
 e_{j\rho}\le[2W+S_w+(k-2)d_j+B]L_a+1 .
\]
In both cases use \(W\le w\), \(kd_j\le N^2\), \(2S_w\le C_1S_c\), \(C_1\ge2\), and \(1\le N^2L_a\). These give the first claimed bound since \(C_0\ge C_1+1\).

For the degree sum write \(b_{\rm tail}=\sum_{k\le i<N_{j-1}}\nu_i\). Before replacing \(W\) by \(w\), the discrepancy from \(I_\rho a_\rho+kW\) is exactly
\[
 \begin{cases}
 \displaystyle\sum_{i<k}c_i+b_{\rm tail},&I_\rho=0,\\[2mm]
 \displaystyle-b_s+\sum_{i=1}^{k-2}c_i+b_{\rm tail},&I_\rho=1.
 \end{cases}
\]
The first lies in \([0,kS_w+B]\), and the second in \([-W,(k-2)S_w+B]\), since \(W-b_s\ge0\). Restoring \(w\) subtracts \(ks_j\). Therefore
\[
 |\Delta_{j\rho}|\le W+kS_w+B+kN.
\]
The desired bound follows from \ref{B5.08b} and \(I_\rho kb'\ge0\). Finally \(\varphi(0)=0\) proves the assertion about extension of the tuple.
\end{proof}

\subsection*{Proof of \ref{B5.08}}\label{prf:B5.08}\phantomsection
\begin{proof}
For each pair \(j,\rho\), the row contribution, the subtracted hook and the term from the reference divisor have sum at most
\[
 I_\rho\{\varphi(a_\rho)-a_\rho\log\Omega\}
 +e_{j\rho}+\{(\beta_j-k)w_{j\rho}-\Delta_{j\rho}\}\log\Omega .
\]
This follows by substituting the degree identity before estimating the term with \(\log\Omega\), whose sign need not be fixed. For the arithmetic data the two remainders and the required bounds are those of \ref{B5.08c}. Each marked position occurs once in every layer. Adding the terminal term \(-\sum_{\rho\in\mathcal U}a_\rho\log T\) gives exactly \(\mathcal D(T)\).

Put \(Z=n_TaL_a\), \(V=mL(B_0+N_c)\). We have \(L\le L_a\le3L\), \(V\ge fm^2bL^2\), \(L\ge1\). There are at most \(mf\) pairs. Per pair the error \(e_{j\rho}\) is at most \(5C_0bL_a\), because \(w_{j\rho}\le2b\) and \(S_c,n_*^2,m(t_c+1)\eta\le b\). Thus its total is at most \(15C_0V\). The \(\beta_j-k\) terms cost at most \(2fm^2bL_a\le6V\). In the degree discrepancy the unmarked part of the bound is at most \(5C_0b\) per pair, using \(kn_*,kS_c\le b\). Its cost is at most \(15C_0V\). The additional marked allowance costs at most \(C_0mn_Tkb'L_a\le C_0 Z\).

The terminal error is at most \(C_0 Z+9C_0V\), since \(b'n_*\le a\) and \(m\eta n_*,S_cn_*,n_*^2\le b\). Here \(S_c\le w''\) and \(n_*w''\le b\). The kernel and hook allowances are bounded by
\[
 mC_hfn_*^2(L_a+m^3L)\le4C_0V,\qquad
 mC_hfn_V^2L_p\le3C_0V .
\]
For the first use \(n_*^2,m^3n_*^2L\le b\), and for the second use \(n_V\le n_*\), \(C_hL_p\le C_0L_a\). Finally \(E_b\le c_eV\) and \(L_{\rm pass}\le1+C_0m\lambda_0\le2C_0V\). Adding all terms gives at most \(2C_0Z+55c_eV\) beyond \(\mathcal D(T)\), which proves the displayed inequality. In the arithmetic application \ref{B5.08x} and \eqref{eq:p4c6-row-budgets} give the required reference bound with \(E_b=E_{l3}+L_G+\sum_j(E_{1,j}+E_{2,j})+E_3\).
\end{proof}

\subsection*{Proof of \ref{B5.08a}}\label{prf:B5.08a}\phantomsection
\begin{proof}
Fix \(t=t_w=t_{\bar w}\) and take bases over the semilocal ring at \(w,\bar w\). Let \(a\) be the determinant of the change from the \(L_{t_0}\) image basis to that of \(L_t\). By elementary divisors \(v_z(a)=[L_{t_0}:R_0]_z\) for \(z=w,\bar w\). The corresponding layer determinants satisfy \(v_z(b_j)=[M_j:A_{jt}]_z-[M_j:A_{j,t_0}]_z\). If \(h_{jt}\) is the layer Gram determinant, the inverse Gram gives \([T_{jt}:T_{jt}^\sharp]_w=-v_w(h_{jt})\).

Let \(D_{\mathrm{row}}(t,v)\) and \(D_j(t,v)\) be the squared determinant ratios for \(\iota_t,\theta_{jt}\). For \(t=t_0\) or \(v=v_0\), the metric identities in \ref{def:p4c6-core} give factors \(\kappa_v\prod_j\lambda_{jt\alpha v}\) and \(\lambda_{jt\alpha v}\) on the row and layer blocks indexed by \(\alpha\). Each such block has the same dimension on the row and on each layer. Multiplying their factors with these dimensions as exponents gives
\[
 D_{\mathrm{row}}(t,v)=\kappa_v^{D_C}\prod_jD_j(t,v)
 \qquad(t=t_0\ \text{or}\ v=v_0),
\]
where \(D_C\) is the row rank.
\leanref{Formalism.\allowbreak{}L2.\allowbreak{}Indef.\allowbreak{}IndefRows.\allowbreak{}detId}{Formalism.\allowbreak{}LocalToGlobal.\allowbreak{}IndefiniteForms.\allowbreak{}IndefUse}{196--206}

Divide the identities for \(t_w,t_0\) at \(v_0\). The common \(\kappa_{v_0}^{D_C}\) cancels. The Gram multipliers for the image basis changes are \(\bar a a\) and \(\bar b_j b_j\). By injectivity of the embedding,
\[
 \bar a a\,G_{t_0,w}\prod_jh_{jt}
 =G_{t,w}\prod_j(\bar b_j b_j h_{j,t_0}).
\]
Taking valuations gives the formula, since \(v_w(\bar b)=v_{\bar w}(b)\), also at a self-conjugate place where the row index occurs twice.
\end{proof}

\subsection*{Proof of \ref{B5.08x}}\label{prf:B5.08x}\phantomsection
\begin{proof}
By \ref{B5.08a} and \(\ell_w=\ell_{\bar w}\),
\[
 \operatorname{lcv}(R_0)-\operatorname{lcv}(L_{t_0})
 =2\sum_w\ell_w[L_{t_0}:R_0]_w
 =\sum_w\ell_w(C_w(t_w)-C_w(t_0)).
\]
For the base quotient map \(P\) and test Gram matrix \(H\) put \(g=\det(PH^{-1}P^*)^{-1}\in K^\times\). At \(\sigma\), write \(d_\sigma=\det(P_\sigma H_\sigma^{-1}P_\sigma^*)^{-1}\), with complex adjoint. For \(\sigma'(a)=\overline{\sigma(\bar a)}\) the entries of the test metric give \(H_{\sigma'}=H_\sigma\). The two row frames have mixed Gram determinant \(\sigma(g)^{-1}\). Cauchy--Schwarz on their exterior products yields \(d_\sigma d_{\sigma'}\le|\sigma(g)|^2\). Average over the permutation of embeddings and use the product formula. A basis change of determinant \(c\) adds \(v_w(c)+v_{\bar w}(c)\) to the local Gram valuation and the corresponding weighted term to the covolume. This gives \(\operatorname{lcv}(\Lambda_{t_0})\le\sum_w\ell_wv_w(G_{t_0,w})\).

The algebraic layer metrics and the product formula give
\(\operatorname{lcv}(T_{j,t_0})=-\sum_w\ell_w[T_{j,t_0}:T_{j,t_0}^\sharp]_w\).
Thus the layer determinant ratio satisfies
\[
 \langle\log D_j(t_0,v)\rangle_\infty=\operatorname{lcv}(M_j)+
 \sum_w\ell_w\bigl(2[M_j:A_{j,t_0}]_w+[T_{j,t_0}:T_{j,t_0}^\sharp]_w\bigr).
\]
The row determinant identity at \(t_0\) adds \(D_C\bar\kappa_0\) and the layer log ratios to \(\operatorname{lcv}(\Lambda_{t_0})\). Adding the index formula for \(R_0\) cancels the contributions at \(t_0\), so
\[
 \mu(R_0^\vee)\le\bar\kappa_0+\sum_j\operatorname{lcv}(M_j)/D_C
 +D_C^{-1}\sum_{w\in S_y}\ell_wC_w(t_w)
\]
for a support set \(S_y\).

Write \(i_{jw}=[M_j:A_{j,t_w}]_w+[M_j:A_{j,t_w}]_{\bar w}\) and \(q_{jw}=[T_{j,t_w}:T_{j,t_w}^\sharp]_w\). The Gram term is at most \(L_G\) by \eqref{eq:p4c6-gram-inputs}. The combined layer contribution vanishes where the selected type \(t_w\) equals the base type \(t_0\). On the places with \(t_w\ne t_0\), \ref{B8.38b} and \eqref{eq:p4c6-gram-inputs} give
\[
\begin{aligned}
 D_C^{-1}\sum_ji_{jw}&\le\sum_j(2e_wF^0_{jw}+\mathrm{err}_{jw})+2e_w/p_w^2,\\
 D_C^{-1}q_{jw}&\le e_wJ^0_{jw}+\eta_{jw}/D_C,\\
 \sum_{w:\,t_w\ne t_0}\ell_w(\mathrm{err}_{jw}+e_wJ^0_{jw})&\le E_{1,j},&
 \sum_{w:\,t_w\ne t_0}\ell_w\eta_{jw}/D_C&\le E_{2,j}.
\end{aligned}
\]
Here \(F^0,J^0\) are the base and pairing-index terms and \(\mathrm{err},\eta\) their allowances. Sum with \(F_*=\sum_j\sum_{w:\,t_w\ne t_0}\ell_we_wF^0_{jw}\).

In the construction let \(M\) be the level group order and \(s>0\) the integral denominator of \ref{B4.55}. The adjoint of the quotient map embeds \(\Lambda_t^\sharp\) in the test dual. Multiplication by \(Ms\) sends the test dual into the test lattice, hence \(Ms\,\Lambda_t^\sharp\subseteq\Lambda_t\) on projecting back. In the semilocal basis \(Ms\,G^{-1}\) is integral at \(w\) for the quotient Gram matrix \(G\). Its determinant gives \(v_w(G_{t,w})\le D_Cv_w(Ms)\), so \(L_G=\lambda_0+\log s\) suffices by \ref{B4.52} when \(\log M\le\lambda_0\).

On a twisted layer let \(J_w\) be the least nonnegative containment exponent between the cut lattices in \ref{B8.38b} and \(a_{jw}\) the length of the lift lattice modulo its cut. Containment and the two denominators give
\[
 q_{jw}\le a_{jw}+D_C(J_w+v_w(d_{j,t_w})+v_w(s)),\qquad J_w\le e_wJ^0_{jw}.
\]
Use \ref{B8.40a} with \(D_0=D_C\) for \(\sum_w\ell_wa_{jw}/D_C\le b_j\), with \(b_j=b_N\) its cut allowance. Since \(d_{j,t_w}\mid d_g\), \ref{B4.52} gives \(E_{2,j}=b_j+\log d_g+\log s\). For \(E_3\), the sum of \(2\log p/p^2\) through \(256\) is less than \(0.979\). Partial summation with \(\sum_{p\le y}\log p<1.01624y\) \cite[(3.32)]{rosser-schoenfeld-1962} bounds the tail by \(4.06496/256\), so \(E_3=1\) suffices.
\end{proof}

\subsection*{Proof of \ref{B5.09}}\label{prf:B5.09}\phantomsection
\begin{proof}
The passage cost is at most \(1+28\lambda\le1+C_0m\lambda_0\). Apply \ref{B5.08} with \(\widetilde\mu_C=\mu_C-\bar\kappa\). The \(z_v\) of \ref{B5.07} have mean at most \(H^*-\bar\kappa\). Add the difference to every \(z_v\) and exponentiate to get positive squared norm bounds of exact mean log \(H^*-\bar\kappa\). For \(B_u=0\) use the constant bound \(\exp(H^*-\bar\kappa)\). Apply \ref{B4.55} with \((\delta,Q_{\rm pass},h_{\rm r},\lambda_{\rm r})=(\delta_0,Q_0,H_0,\lambda)\). In the arithmetic coordinates \ref{B4.38} allows \(\sigma_w=s_{uw}\) from \ref{B5.03}.
\end{proof}

\subsection*{Proof of \ref{B5.09a}}\label{prf:B5.09a}\phantomsection
\begin{proof}
Under the row group isomorphism and coefficient embedding, a retained \(\chi\) is \(\prod_w\chi_w\). Each \(\chi_w\) is the character induced from the prescribed cuspidal datum on its block parabolic. Reordering the local labels conjugates the subgroup character by an inner automorphism and preserves the normalized twisted trace. Thus \ref{B7.47b} and \ref{B7.48} give the trace of the local subgroup character projector as
\[
 n_w=\sum_{j=0}^k c_j\gamma_q(k-j),\qquad
 \gamma_q(0)=1,\quad\gamma_q(s)=q^s-q^{s-1}\ (s>0),
\]
where \(q=q_w\) and \(c_j\) counts subsets of cuspidal blocks of total length \(j\). A subset is determined by its union of positions, so \(c_j\le\binom rj\). The empty subset contributes \(q^k-q^{k-1}\). The other terms divided by \(q^k\) sum to at most \((1+q^{-1})^r-1\le2r/q\). Hence
\begin{equation}\label{eq:B5-local-multiplicities}
 1-q_w^{-1}\le n_w/q_w^k\le1+(2r-1)/q_w.
\end{equation}
Over the \(2f\) places, \(q_w\ge g\) and \(4fr/g\le1\) give a product between \(1-2f/g\) and \(\exp(4fr/g)\le1+8fr/g\). The row budget with \(2F_*=\Fo\) from \ref{B4.55b} gives \(B_*\le B\).

The sizes in \ref{B5.08}, with \(L_a\le4P\), give \(A_3\le16P^4b\). For \(F_{jw}\) in \ref{B4.02} use the weight bound \(W=4P^5b\). The type-prime and ramified terms are nonpositive. In the inert term, \(1/(p^2-1)\le1/p\) and \(p/(p^2-1)\le2/p\). The orbit length division counts each orbit once by \ref{B4.01}. Each orbit costs at most \(2(k-1)W/p\) per unit length, with at most \(P\) orbit labels and a factor two in the inert sum. Thus \(F_{jw}\le16P^7b/p\), with positive terms confined to \(p\le X''\). Group by \ref{B4.52} and use \cite[(3.24)]{rosser-schoenfeld-1962} in the form used in \ref{B5.06}. The cost per layer is at most \(64P^7b\log(X''+1)\). Since \(m\le P,\ \log(X''+1)\le4P\), this yields \(\Fo\le256P^9b\).

For \(n=|A|\), \(T_2\subseteq\{0,\ldots,n-1\}\) gives
\[
 N_c=\frac{|T_2|kbL^2}{n^2}\le P^3b,\qquad
 B_0=fmbL\le P^3b,\qquad L(B_0+N_c)\le2P^4b.
\]
The level size identification is \(M=\prod_w|\GL_r(k_w)|\) over the \(2f\) selected residue fields, whose sizes are at most \(2g_0\). Thus \(\log M\le2fr^2\log(2g_0)\). Use \(2r^2\le b\) and the selected prime bound \(\log(2g_0)\le2L^2\) to get \(\log M\le2mL(B_0+N_c)\le4P^5b\).
\leanref{Formalism.\allowbreak{}L2.\allowbreak{}DIV.\allowbreak{}Arith4S.\allowbreak{}Nc\_\allowbreak{}le\_\allowbreak{}rec}{Formalism.\allowbreak{}LocalToGlobal.\allowbreak{}Divisibility.\allowbreak{}EigenvalueSumLeClosedOf}{1743--1783}
\leanref{Formalism.\allowbreak{}L2.\allowbreak{}DIV.\allowbreak{}Cons.\allowbreak{}log\_\allowbreak{}card\_\allowbreak{}le\_\allowbreak{}unitE}{Formalism.\allowbreak{}LocalToGlobal.\allowbreak{}Divisibility.\allowbreak{}DivisibilityConstruction.\allowbreak{}RBLam}{62--115}

\(B\le(272+2c_e)P^9b\). Since \(P\) is polynomially bounded in \(x\), eventually \(B\le b^2,\ x^{C_H}\le b,\ 28\log M\le b^2\). This gives \(Q_0\le3b^2\). The prime size \(g\ge(kb)^8\) gives \(24frQ_0\le72frb^2\le b^8\le g\) eventually.
\end{proof}

\subsection*{Proof of \ref{B5.10}}\label{prf:B5.10}\phantomsection
\begin{proof}
Fix a member beyond the common cutoff and a word \(u\). Use the bijections of positions and layers in \ref{B8.35}. They carry the marks to \(\{0\}\times TT\) and preserve \(a_t\) at every mark. Padding a column adds terms \(\varphi(0)-0\log\Omega=0\). Lemma~\ref{B4.02a} identifies its positive row partition with \(\kappa^{\rm big}\). Thus the terminal, kernel and hook entries in \ref{B5.08} are respectively \(H_{P,u},H_{K,uj},L_{6,j}\) of \ref{B5.07}. With \(\bar\kappa=\langle\log\kappa_v\rangle_\infty\), the defining formula for \(h_u\) gives \(H^*-\bar\kappa=h_u\). The fixed scale and the untwisted reference row give \(2F_*=\Fo\) by \ref{B4.55b}.

Apply \ref{B5.07a} to the definite tables and \ref{B5.07b} to the reference spectrum. These give the transition bounds and the spectral subtraction in \ref{B5.07}. The paired indices of \ref{B8.38b}, the cut estimate of \ref{B8.40a} and the reference Gram bounds give the bound for \(\mu_0\) in \ref{B5.08x}. The numerical row bounds are \ref{B5.08b} and \ref{B5.08c}; the numerical passage bounds on \(\delta_0,Q_0\) and \(\lambda\) are those of \ref{B5.09a}, while the passage inequality for \(\mu(RR^\vee)-\bar\kappa\) follows from \ref{B4.55a}, \ref{B8.23} and \ref{B5.09a}. Theorem~\ref{B5.09} therefore applies to the coordinate pairing \(B_u^{\rm c}\) with \(T=T_u\). Its image on the shifted coordinate lattice lies in the dual coordinate row lattice. Its squared operator norms are bounded by positive \(c_{u,v}\) with
\[
 \langle\log c_{u,v}\rangle_\infty=H^*-\bar\kappa.
\]
Its slope inequality, using the right hand side of \ref{B5.08}, is
\[
 \mu_C+h_u-\Fo\le\mathcal D(T_u)
 +(50C_0+10)n_TaL_a+60c_emL(B_0+N_c).
\]
Here the row isometry of \ref{B4.56} identifies the two dual row slopes. Since \(B_0,N_c\ge0\), \(x\ge1\) and \(L=x^\varepsilon\le x^{e_{\rm Div}\varepsilon}\), the last two terms are at most \(E_{\rm div}\).

Let \(E_R,E_u\) be the rational row and column isometries. For arbitrary rational coordinate vectors \(\phi,F\) one has
\[
 B_u(E_u\phi)(E_RF)=B_u^{\rm c}(\phi)(F).
\]
Both isometries carry the specified lattices onto the actual lattices. Evaluating on lattice vectors gives \(B_u(N'_u)\subseteq RR^\vee\). After scalar extension to \(v\), take the supremum of the absolute value over the two unit balls. Both isometries are onto, so \(\|B_u\|_v=\|B_u^{\rm c}\|_v\). If \(B_u\ne0\), some matrix entry over \(K\) is nonzero at every embedding. Taking logarithms and averaging gives \(h(B_u)\le\langle\log c_{u,v}\rangle_\infty=h_u\).

The support assertion follows from \ref{B4.52a} with the stabilizer denominator of the same word. The frame assertion follows from \ref{B5.10a} applied with the chosen \(C_H\) to these full positive-rank lattices. Lemma~\ref{B5.13} bounds the loss with constant \(c_E\); Lemma~\ref{B5.14} bounds the difference of the adjusted heights with constant \(20+6/\varepsilon\). Both are at most \(c_{\rm Div}\). The fullness and complex period identifications in the hypotheses supply the remaining clauses of Definition~\ref{B5.10b}.
\end{proof}

\subsection*{Proof of \ref{B5.10a}}\label{prf:B5.10a}\phantomsection
\begin{proof}
Put \(q_\kappa=\langle\log\kappa_v\rangle_\infty\). By \ref{B8.23} and \ref{B4.55a}, \(\mu_{\max}(RR)+q_\kappa\le x^{a_{\rm row}}b+14\lambda_0\), with \(\lambda_0\le4b\). Take \(F_0=\max(a_{\rm row},1)+1\), which bounds the right side by \(x^{F_0}b\) for \(x\ge57\).

The rank windows give \(P=1+f+k+m+t_{\rm ch}+N_m+L\le x^M\) eventually. In the maximum-slope bound of \ref{B5.06}, the \(\mu_{\max}(M_u)\) term is at most \(x^{4M+1}b\) eventually by \ref{B5.12}. For \(\Ga\), the retained terms of \ref{B3.51}--\ref{B3.52} give \(\sum_j\gainD_j(w)-\tau_w\le4mn_Ta/p_w\) through \(X''\), with zero gain above \(X''\) by \ref{B4.52a}. Grouping by \ref{B4.52} and applying \cite[(3.24)]{rosser-schoenfeld-1962} gives
\[
 \Ga\le16mn_Ta\log(X''+1)\le64P^4b
\]
using \(n_Tx\le2k\) and \(\log(X''+1)\le4P\). Add \(\Fo\le256P^9b\) and \(-\Ee_u\le c_EmB_0\). Include \(c_1mB_0\), with \(mB_0\le P^4b\). The arithmetic part is eventually at most \(x^{9M+2}b\).

For the height, \eqref{eq:B8-five-size} gives \(S_a\log T_u\le2|C_{Tu}|P^4b\). The terminal size coefficient before \(L_a\) is at most \(4P^3b\), hence \(-H_{P,u}\le(2|C_{Tu}|+16|C_{\rm sz}|)P^9b\). For the row entries, \ref{B8.35}--\ref{B8.37} give \(|\nu|\le M_\nu=11P^5b,\ \log M_\nu\le16P,\ |\log\Omega|\le7P\). A single entry costs at most \(M_\nu(23P+1)+1\le300P^6b\) in the negative kernel sum. Thus \(-K_{uj\rho}\le300P^7b\).

Fix a type \(\rho\) and a layer \(1\le j\le m\). The plus hook has \(k+j-1\) rows, so its last gamma argument is \(1+\kappa^{\rm big}_{j\rho,k+j-2}\). The last entry of the long row is \(e_{j-1}-b_{j-1}-(t_{\rm ch}+1)/2\). Hence the gamma argument is \(e_{j-1}-b_{j-1}-(t_{\rm ch}-1)/2\ge2\) by \ref{B4.03a}. Dominance gives the same lower bound on the earlier arguments, and their upper bound is \(Y=10P^5b\). Thus each factorial log is at most \(Y\log Y\le160P^6b\). There are at most \(P\) numerator terms and the denominator factors are at least one, so \(\log\operatorname{Hook}(a^+_{j\rho})\le160P^7b\).

Summing over at most \(P^2\) layer-type pairs and using \(fN_m^2L_a\le4P^4,\ fm^3N_m^2L\le P^7,\ fn_V^2L'\le6P^5\) per layer gives
\[
 -h_u-q_\kappa\le(500+2|C_{Tu}|+30|C_{\rm sz}|)P^9b\le x^{9M+1}b
\]
eventually. Add the bounds for \(\mu(N'_u),-h_u-q_\kappa\) and \(\mu_{\max}(RR)+q_\kappa\). The choice of \(C_H\) absorbs the fixed number of contributions for \(x\ge5\). Average with weights \(d_u/\sum_{u'}d_{u'}\).
\end{proof}

\subsection*{Proof of \ref{B5.11c}}\label{prf:B5.11c}\phantomsection
\begin{proof}
The bounds on \(f=2n,k,m\) follow from \eqref{eq:host-rank-ranges} since \(9/2+10\varepsilon<5\) and \(m^2L\le4x^{1+\varepsilon}\) eventually. The discriminant bound is \ref{B8.35a}; the quadratic subfield gives \(2v_E\ge\sqrt d\), with \(d\) tending to infinity by \ref{B2.55}. Apply \ref{B2.51} for the lower prime bound and \ref{B8.35a} with \(L\ge1\) for the upper bound.

Use the exact coordinates of \ref{B4.02b}. Off the marks \(U_0-v_0=S_c+1/2\) and \(U_{k-1}=-1/2\). At a mark the spread is \(a_s+b_s+k-1\), with \(a\le a_s\le a+b\) and \(b'\le b_s\le2b'\) by \ref{B5.22a} and \ref{B8.35a}. Here \(b/a=1/x\), \(b'/a=1/X'\), and \(k,S_c\) grow polynomially while \(b'\) exceeds every fixed power. These give all the endpoint inequalities, including the span bound in both cases. The base \(b_0\) is at least \(w''-1\) by the three base formulas. Thus \(b_0/x\ge\max(S_c,h_0/2)\) eventually. Since \(e_q=b_q+h_0/2\) and \(b_q\ge b_0\), this proves the relative bound. Apply Lemma~\ref{B5.08d} to the same base and step size to obtain \(e_{m-1}\le2b\). Finally \(b_{q+1}-b_q=h_0\) and \(e_q=b_{q+1}-h_0/2\). The node threshold in the proof of \ref{B8.35a} gives \(h_0\ge x^{1300}\), while \(r\le x^5\), \(f\le x^3\), \(L\le x\) by \ref{B8.59}. This gives the asserted gap.
\end{proof}

\subsection*{Proof of \ref{B5.11}}\label{prf:B5.11}\phantomsection
\begin{proof}
The prime interval gives \(\log T_u=\log X_u-2\log\Gamma_u/(fk)+\epsilon_u\) with \(0\le\epsilon_u\le\log(1+1/(fk))\le1\). This loses at most \(2\) in the lower bound and adds at most \(1\) to the upper bound. The common denominator cancels in differences, and \(|\epsilon_u-\epsilon_{u'}|\le1\). These errors fit in \(c_{\rm Count}\). In particular \(-S_a\log T_u\le-S_a\lambda_5\), as used in \ref{B5.36}.

For the family, take \(C_0=C_Y=\max(5,C_{YL}+1)\), with \(C_{YL}\) the type-prime size constant. Take \(x\ge44^4\) with \(x>x_1\) beyond the cutoff in \ref{B5.11c}. Apply that lemma to the same word and type, numbering the coordinates from one in the following calculation. Put $r=k+m-1$, and let $\Delta$ be the spacing of the two choices of a word coordinate. The scalar conditions are $f,k\le x^{C_0}$, $m^2L\le k$, $(kb)^8\le g_i$, $\log g_i\le C_YL$ and the bounds on $m,k,v_E$ below. The consecutive auxiliary weights have gaps at least $r^{19}f^5L/C_0$. Off the top the leading gap is between $4k\Delta$ and $8k\Delta$. On the top the leading coordinate is at most $2a$. The remaining coordinate conditions are given in the next four sentences. To obtain them, write the interior word coordinates using the two choices in \ref{B8.02c}; their common interlacing row has successive spacing $4\Delta$ and lies in $[0,S_c]$, where $S_c=4k\Delta$. The first and last coordinates agree with those of the reference word. On a marked type they equal $a_s+(k-1)/2$ and $-b_s-(k-1)/2$. The top window \eqref{eq:p3c5-top-window}, $b_s\le2b'$, $k\le b$, and $100b\le a$ give the asserted span and endpoint bounds there, using $S_c\le b'\le b$. Off the top the first gap is in $[S_c,2S_c]$ and the last coordinate is $v_{k-1}-1/2$, giving span at most $3S_c+1/2\le2a$. The rank window of \ref{B8.01k} gives $44\le m\le x$ eventually. The prime window $Y\le g_i\le2Y$ with $(kb)^8\le Y$ and $\log Y\le C_{YL}L$ gives $\log g_i\le C_YL$ since $L\ge1$. The auxiliary-weight spacing follows from $b_{j+1,\rho}-b_{j\rho}\ge\eta/m$ and $r^{19}f^5L\le\eta/m$ in the selected window.
\leanref{Formalism.\allowbreak{}MainNotions.\allowbreak{}GlobalTheoryTheory.\allowbreak{}JDB.\allowbreak{}db\_\allowbreak{}cloud}{Formalism.\allowbreak{}SharpRestriction.\allowbreak{}GlobalSchurBounds.\allowbreak{}JointsL}{738--786}
\leanref{Formalism.\allowbreak{}MainNotions.\allowbreak{}GlobalTheoryTheory.\allowbreak{}JDB.\allowbreak{}dataB\_\allowbreak{}of}{Formalism.\allowbreak{}SharpRestriction.\allowbreak{}GlobalSchurBounds.\allowbreak{}JointsL}{1085--1091}
\leanref{Formalism.\allowbreak{}MainNotions.\allowbreak{}GlobalTheoryTheory.\allowbreak{}JDB.\allowbreak{}dataA\_\allowbreak{}fam}{Formalism.\allowbreak{}SharpRestriction.\allowbreak{}GlobalSchurBounds.\allowbreak{}GlobalRankLatticeProofs.\allowbreak{}FamilyPartFiveOfExplicitConstants}{303--315}
At each \(\rho\) the interlacing entries \(v_i\) have spacing \(v_i-v_{i+1}=4\Delta\), \(\Delta\ge1\), and lie in \([0,4k\Delta]\). A middle shifted column entry is \(v_{i-1}-1/2\) or \(v_{i-1}-1/2-\Delta\). At a nontop place the last entry is \(v_{k-1}-1/2\); at a top place it is at most \(-4k\Delta\) and at least \(-a\). The leading coordinate exceeds \(v_1\) by at least \(4k\Delta\), and the column span is at most \(2a\). The row sizes include \(44\le m\le\min(k,x),\ \log(2k)\le L,\ \pi\le v_E,\ \log(2v_E)\le L\). For the negative row entries \(-b_{j\rho}\), use
\[
 1\le b_{j\rho}<w_{j\rho}\le(1+1/x)b_{j\rho},\qquad w_{j\rho}\le2b.
\]
The auxiliary entries are separated by the intervening weights, with \(b_{1\rho}/x\ge4k\Delta\).

Write \(z_a\) for the shifted column coordinates and \(\Delta_\rho(u)=\prod_{a<b}(z_a-z_b)\). The first-entry gaps change by at most \(\Delta\) and are at least \(4k\Delta\). For the other pairs the change is at most \(2\Delta\) and the gaps are at least \(3(b-a)\Delta\). By \(|\log y-\log y'|\le|y-y'|/\min(y,y')\), the change in \(\log\Delta_\rho\) is at most \(1/4+\frac23 k(1+\log k)\). The Weyl denominator is common to the words, so \(|\log X_u-\log X_{u'}|\le2\log k\).

Set \(r=k+m-1,\ F_n=\prod_{i=0}^{n-1}i!\). For the row parameter \(\lambda_\rho\), \ref{B8.50a} at \(t=1/3\ge e^{-r}\) gives \(d_{\Sigma,\rho}\ge(2/3)\Delta(\lambda_\rho)/F_r\), where \(\Delta(\lambda_\rho)\) is its Vandermonde. The required gap is at most \(9C_3^+r^9\), and the supply is \(r^{19}f^5L/C_0\). Thus \(L\ge9C_3^+C_0\) suffices. The \(2k-3\) column gaps incident to the first or last entry are at most \(8a\); for interior pairs \(z_i-z_l\le v_{i-1}-v_l\). Each cross gap at \(-b_{j\rho}\) is at least \(b_{j\rho}\); unused row gaps are at least one. It follows that
\[
 \Delta_\rho(u)\le(8a)^{2k-3}\Delta(v),\qquad
 \Delta(\lambda_\rho)\ge\Delta(v)\prod_j b_{j\rho}^{k-1}.
\]
The first power costs at most \(5kL\) once \(2(\log8+1+\log x)\le L\). Also \(k\log w_{j\rho}-(k-1)\log b_{j\rho}\le k/x+2L\). With \(m\le\min(k,x)\), this yields
\begin{equation}\label{eq:B5-dimension-lower}
 S_{\rm dim}-k\sum_{\rho,j}\log w_{j\rho}\ge-f(1+k+7kL),\qquad
 S_{\rm dim}=\sum_\rho\log\frac{d_{\Sigma,\rho}F_r}{\Delta_\rho(u)}.
\end{equation}

The volume identity \ref{B8.51b} and adjacent-rank mass ratio \ref{B8.48} cancel the Weyl factorials. With \(\ell_E=\log(v_E/\pi)\) their product is
\begin{equation}\label{eq:B5-count-product}
\begin{aligned}
 fk\log X_u={}&f\left((m-1)k+\frac{m(m-1)}2\right)\ell_E
 -\frac{f(m-1)}2\log(2v_E)+S_L\\
 &+f\left(km+\frac{m(m-1)}2\right)\ell_g+fS_g+S_{\rm dim},
\end{aligned}
\end{equation}
where \(S_L=\sum_{n=k+1}^r\log L_\eta(E,n)\) and \(S_g=\sum_{i=0}^1\sum_{n=k+1}^r\log(1-g_i^{-n})\). The Euler bounds give \(|S_L|\le fm,\ -2\le S_g\le0\). Subtract \(fk(\mathcal W+m\ell_g)\). By \eqref{eq:B5-dimension-lower} the dimension term costs \(f(1+k+7kL)\). The unmatched linear term in \(\ell_E\) costs at most \(fkL\), the discriminant term costs at most \(f(m-1)L/2\), and the two Euler sums cost at most \(fm+2f\). The remaining quadratic terms in \(\ell_E,\ell_g\) are nonnegative. Hence the difference is at least \(-f(8kL+(m-1)L/2+m+3+k)\). Take the family cutoff so that \(L\ge2\). Since \(m\le k\) and \(k\ge3\), \(m+3+k\le3k\le 3kL/2\); the displayed cost is at most \(11fkL\).

For the upper bound, \(v_i-v_l\le z_i-z_{l+1}\) gives \(\Delta(v)\le\Delta_\rho(u)\), since the unused column gaps are at least one. The compact Weyl product in \ref{B8.50a} has at most \(2km\) factors left after the common-shift interlacing entries, each at most \(H+2k\) for largest row gap \(H\le4e^L\). Its denominator \(F_pF_q\) is at least \(F_k\) in signature \((p,q)\), and \(pq\le km,\ r\le2k\). Hence
\[
 \log\frac{d_{\Sigma,\rho}}{\dim V(u,\rho)}
 \le\log4+2km\log(H+2k)+km\log(3pqrH)\le km(12+10L).
\]
Using \(\log(F_r/F_k)\le2mkL\) gives \(S_{\rm dim}\le fkm(12+12L)\). In \eqref{eq:B5-count-product}, \(\ell_E\le L,\ \ell_g\le2C_YL\) give \(\log X_u\le(15+3C_Y)mL\) for \(L\ge9,\ k\ge5\).
\leanref{Formalism.\allowbreak{}MainNotions.\allowbreak{}StandardLattice.\allowbreak{}part\_\allowbreak{}five\_\allowbreak{}family\_\allowbreak{}of\_\allowbreak{}mass\_\allowbreak{}volume}{Formalism.\allowbreak{}SharpRestriction.\allowbreak{}GlobalSchurBounds.\allowbreak{}LinkMin}{318--423}

The lower bound and \(\ell_g\ge16L,\ \mathcal W\ge0\) give \(\log T_u^\star\ge16mL-11L-2\ge mL\) eventually. For this choice of primes, \ref{B8.05} gives \(m\ge x^{1/3},\ 1\le fk\le x^{12},\ g_i\le b^9,\ \log(Q_Sqd)\le2x^{\varepsilon/5}(\log x)^2\). Lemma~\ref{B8.51a} applies with scale \(T_u^\star\) and interval parameter \(fk\), and its congruence yields the splitting of \ref{B8.54}. Take \(C\ge\max\{11,16+3C_Y,2\}\) in \eqref{eq:B5-count-torus}, and take \(x_0\) above these cutoffs and the family cutoff \(x_1\).
\end{proof}

\subsection*{Proof of \ref{B5.11b}}\label{prf:B5.11b}\phantomsection
\begin{proof}
Fix \(\rho\) and spacing \(\Delta>0\). In the shifted coordinates the first entry is fixed. The middle entries are \(v_i-1/2-\delta_i\Delta\) with \(v_i-v_l=4(l-i)\Delta,\ \delta_i\in\{0,1\}\). At a nontop type the last entry is in the list with bit zero. At a top type it is fixed at \(-b_{\rm end}\), where \(b_{\rm end}=b_v+(k-1)/2\) for \(b_v\) in \ref{B5.05a}. For \(J=k-1\) regularly spaced entries at a nontop type, and \(J=k-2\) at a top type, the ratio to the unperturbed Vandermonde is
\[
 R(\delta)=\prod_{i<l\le J}\left(1+\frac{\delta_l-\delta_i}{4(l-i)}\right).
\]
For \(|z|\le1/4\), \(|\log(1+z)-z|\le2z^2/3\), so the quadratic errors total at most \(J\pi^2/144\). At distance \(d\), the linear sum telescopes to two boundary sums with difference bounded by \(\min(d,J-d)\). With \(q=\lceil J/2\rceil\), four times the absolute linear contribution is at most
\[
 q+\sum_{d=q+1}^J(J/d-1)\le2q-J+J\log(J/q)\le1+0.7J .
\]
Thus \(|\log R|\le(J+1)/4\). The first-to-middle gaps are at least \(4k\Delta\) and change by at most \(\Delta\). Their total contribution is less than \(1/3\). The first-to-last gap is fixed. At a top type the remaining edge changes cost at most \((k-2)\Delta/b_{\rm end}\). Since \(b_{\rm end}\ge b_1/2,\ b_1\ge C_\ell mk\Delta,\ C_\ell\ge40,\ m\ge2\), one has \(4k\Delta/b_{\rm end}\le1/10\). Combining the two word bounds for \(R\) with these edges gives
\begin{equation}\label{eq:B5-vandermonde-ratio}
 |\log(\Delta_\rho(u')/\Delta_\rho(u))|\le k/2+1\le3k/4\qquad(k\ge4).
\end{equation}
\leanref{Formalism.\allowbreak{}FourierDecayAlt.\allowbreak{}P.\allowbreak{}vandermondeRatioHarmonic}{Formalism.\allowbreak{}HermitianCurvature.\allowbreak{}HermitianCurvatureProofs.\allowbreak{}VandermondeRatioHarmonic}{379--447}
\leanref{Formalism.\allowbreak{}FourierDecayAlt.\allowbreak{}P.\allowbreak{}vandermondeWords}{Formalism.\allowbreak{}HermitianCurvature.\allowbreak{}HermitianCurvatureProofs.\allowbreak{}SumInvSqLeSub}{738--750}

Put \(g=g_0g_1,\ r=k+m-1\) and define \(t_{u\rho}>0\) by
\[
 t_{u\rho}^{\,k}=(v_E/\pi)^{(m-1)k+m(m-1)/2}(2v_E)^{-(m-1)/2}
 g^{km+m(m-1)/2}\frac{d_{\Sigma,\rho}F_r}{\Delta_\rho(u)}.
\]
Then \(t_{u\rho}/t_{u'\rho}\le e^{3/4}\). Write \(d_u^S\) for the level-zero dimension and \(A_u=\prod_\rho(t_{u\rho}/T_u)^k\). The stabilizer condition of \ref{B8.54} follows from \(g_i\ge(kb)^8\), so \(d_u^S=(T_u^k-1)^fd_u^\star\). Equation \eqref{eq:B5-count-product} gives
\[
 D^\star/d_u^S=E_u A_u,\qquad
 E_u=\frac{\prod_{j=k+1}^{r}L_\eta(E,j)
 \prod_{i=0}^1\prod_{j=k+1}^r(1-g_i^{-j})^f}{(1-T_u^{-k})^f}.
\]
For \(j\ge k+1\ge6\), \(\zeta(j)-1\le2^{1-j}\) by integral comparison. The \(L\)-terms cost at most \(2f/2^k\) in absolute log by \ref{B8.48}. Each type-prime product and the denominator cost at most \(2f/2^k\) by \(-\log(1-y)\le2y\). Hence \(|\log E_u|\le z=f2^{3-k}\).
\leanref{Formalism.\allowbreak{}MainNotions.\allowbreak{}BridgeLemma.\allowbreak{}degree\_\allowbreak{}ratio\_\allowbreak{}product\_\allowbreak{}formula}{Formalism.\allowbreak{}SharpRestriction.\allowbreak{}CuspSpanNormBounds.\allowbreak{}BruhatCellsKronecker.\allowbreak{}DegreeRatioFormulaStatement}{294--329}

The spectral row gives \(D^\star/(C_0\Gamma^2)\le d_u^S\le C_0D^\star/\Gamma^2\) with \(\Gamma\) common. Thus \(A_{u'}/A_u\le e^{2z}C_0^2\), and
\[
 (T_u/T_{u'})^{kf}
 =\frac{A_{u'}}{A_u}\prod_\rho(t_{u\rho}/t_{u'\rho})^k
 \le e^{3kf/4+2z}C_0^2\le(C_0e^{7/4})^{kf}.
\]
Here \(2z\le kf,\ 2\le kf\). The Vandermonde bound follows from \eqref{eq:B5-vandermonde-ratio}.
\end{proof}

\subsection*{Proof of \ref{B5.11a}}\label{prf:B5.11a}\phantomsection
\begin{proof}
For the rank-\(k\) class lattice \(A\subseteq A^\vee\) over \(E\), the normalized degree is \(-\log[A^\vee:A]/[E:\Q]\) in the Hermitian dual metric. The local factors give \([A^\vee:A]=T_u^{2f}\), so the slope is \(-l/k\). Isometric scalar extension preserves it.

Let \(h_\nu\) be the squared norm of the highest Fock vector and \(c_\nu=h_\nu/\prod_i\nu_i!\). Write \(\mu_\nu\) for minus the log Gram determinant per rank of the standard Weyl lattice. In representation dimension \(d'\) the determinant character is \(\det^{d'\sum_i\lambda_i/k}\) by permutation symmetry of weights. Tensor over \(\rho\), then sum the orthogonal, equal-rank classes to get
\[
 \mu(M_u)=\sum_\rho[\mu_\nu-(d_\rho/k+t_\rho)l].
\]
For an integral basis matrix \(V\) of the dual polynomial lattice in monomial coordinates, its Gram matrix is \(V^{\mathsf T}\operatorname{diag}(m!/h_\nu)V\), where \(m!=\prod_{a,i}m_{ai}!\) and \(\sum_a m_{ai}=\nu_i\). The multinomial bound gives \(m!/h_\nu\ge k^{-d_\rho}/c_\nu\). A nonzero maximal minor of \(V\) is an integer. Cauchy--Binet and duality give \(\mu_\nu\ge-d_\rho\log k-\log c_\nu\). In the hook formula of \ref{B5.07} the denominator is an integer at least one. Dividing by \(\prod_i\nu_i!\) leaves at most \(k(k-1)/2\) factorial factors of size at most \(s_\rho+k\). 

For the maximum slope use \(e=[E:\Q]\) and the trace metric on \(A^\vee\) of Euclidean rank \(ek\). For a nonzero Euclidean dual vector \(z\), testing against \(\mathcal O_Ey,\ y\in A^\vee\), gives \(h(z,y)\in\mathfrak D_E^{-1}\), with \(\mathfrak D_E\) the absolute different. For \(\delta\in\mathfrak D_E\), \(\bar\delta z\) pairs integrally with \(A\) and lies in \(A^\vee\). Testing against this vector gives \(\mathfrak D_E^2 h(z,z)\subseteq\mathcal O_E\). Taking norms and applying the arithmetic-geometric mean to the positive conjugates gives \(\operatorname{Tr}h(z,z)\ge e/\mathrm{rd}(E)^2\). Banaszczyk's transference bound \(\lambda_n(\Lambda)\lambda_1(\Lambda^*)\le n\) for Euclidean rank \(n\) \cite{banaszczyk-1993} gives \(ek\) real independent vectors in the trace lattice of squared length at most \(ek^2(2v_E)^2\). Choose \(k\) independent \(E\)-vectors \(y_a\) among them.

Pair with \(y_a\) to form the substitution matrix on each column. Its Gram operator is bounded by its trace. The polynomial substitution in degree \(d_\rho\) therefore gives a squared norm multiplier at \(v\) at most \((\sum_a h_{v,\rho}(y_a,y_a))^{d_\rho}\) in the normalized monomial metric. The trace-length bounds and concavity of \(\log\) give mean log at most \(d_\rho\log(k^3(2v_E)^2)\) for this factor.

Multiply the polynomial coordinates by the determinant factor of exponent \(t_\rho\) in each column. The coordinates of the class lattice then lie in a common fractional ideal \(I\). At each finite place its generator comes from the product of the local column determinant lattices to powers \(t_\rho\). At an infinite place the squared inverse determinant factor of the column frames is a positive multiplier \(D_v\). In the doubled slope convention the determinant degree of the \(\rho\)-column lattice is \(k\,\mu(A)=-l\). Combining the ideal norm and the frame determinants therefore gives
\[
 -\frac{2\log N(I)}{[K:\Q]}+\langle\log D_v\rangle_\infty
 \le-\sum_\rho t_\rho l .
\]
In local column bases the ideal is the product of the determinants of the basis matrices to powers \(t_\rho\), and at infinity \(D_v\) is the squared norm of the corresponding inverse frame determinant raised to these powers. The two log sums together are \(\sum_\rho kt_\rho\mu(A)\) for the determinant lattices. This comparison uses the signed powers before passing to \(|t_\rho|\).
\leanref{Formalism.\allowbreak{}L2.\allowbreak{}GlobalTheoryTheory.\allowbreak{}TraceIntegrality.\allowbreak{}classMapT\_\allowbreak{}of}{Formalism.\allowbreak{}LocalToGlobal.\allowbreak{}GlobalTheory.\allowbreak{}GL4.\allowbreak{}PacketOfFrameDegree}{72--105}

For a rank-\(s\) sublattice, divide the coordinates locally by a generator of \(I\) and take a nonzero maximal coordinate minor over \(K\). Its finite valuations, measured in a local basis of the sublattice, are bounded below by \(s v_w(I)\). At infinity Cauchy--Binet bounds the Gram determinant below by the squared modulus of that minor times the product of the monomial weight lower bounds. The product formula combines these estimates to bound the slope of the coordinate image by \(\sum_\rho(d_\rho\log k+\log c_\nu)-2\log N(I)/[K:\Q]\). The squared norm bound for the substitution adds its mean log and \(\langle\log D_v\rangle_\infty\) to the source slope bound. Apply the determinant comparison and \(-t_\rho l\le |t_\rho|l\). As \(v_E\ge2\), \(4k^4v_E^2\le(kv_E)^4\). This gives the asserted bound on each class. Integral orthogonal projections, with their kernel-image rank filtration on a sublattice, give the same bound for the sum of classes.
\end{proof}

\subsection*{Proof of \ref{B5.12}}\label{prf:B5.12}\phantomsection
\begin{proof}
The first assertion follows from \ref{B5.11a}. Put \(C_n=\max(3,c_{\rm Count},C_{YL}+1),\ c_L=3\) and \(l_T=\log T_u\le c_{\rm Count}mL\). Increase the cutoff for \(x^{14}\le X',\ 7\log x\le L,\ m\ge c_L+c_{\rm Count}+6C_n+9\). The size rows give \(\log(xbv_E)\le c_LL,\ N_mm^4\le b'\). With \(|D_1|\le L\) these give the numerical conditions with log argument \(c_LL\).

At a top column the first entry is \(\alpha_{u\rho}\in[a,2a]\) and the negative of the last is in \([b',2b')\). The interior entries are in \([0,S_c]\). Other columns have entries in \([0,C_nS_c]\). With \(B=b'+S_c\), the decomposition of \ref{B5.11a} gives
\[
 \sum_\rho d_\rho\le2n_Ta+O_{C_n}(fkB),\quad
 \sum_\rho|t_\rho|\ll_{C_n} fB,\quad
 \sum_{\rho,i}\lambda_i\le\sum_{\rho\ {\rm top}}\alpha_{u\rho}
 +O_{C_n}(fkS_c).
\]
Since \(S_c\le b'\le b,\ k\le x^5\), \(\log(s_\rho+k)\ll_{C_n} c_LL\). Put \(S_1=\sum_{\rho\ {\rm top}}\alpha_{u\rho}l_T/k\). Substitution in \ref{B5.11a} gives constants \(C_5,C_6\) depending on \(C_n\) with
\[
\begin{aligned}
 \mu(M_u)&\ge-S_1-C_5(n_Ta\log k+fkB\log k+fS_cl_T+fk^2c_LL),\\
 \mu_{\max}(M_u)&\le C_6(n_Ta\log(kv_E)+fB(k\log(kv_E)+l_T)+fk^2c_LL).
\end{aligned}
\]
\leanref{Formalism.\allowbreak{}MainNotions.\allowbreak{}flipTwoThree.\allowbreak{}nums\_\allowbreak{}wts\_\allowbreak{}of\_\allowbreak{}fam}{Formalism.\allowbreak{}SharpRestriction.\allowbreak{}GlobalSchurBounds.\allowbreak{}GlobalRankLatticeProofs.\allowbreak{}SevenLogLeLOfPowerBound}{698--747}
\leanref{Formalism.\allowbreak{}MainNotions.\allowbreak{}flipTwoThree.\allowbreak{}numsL\_\allowbreak{}of\_\allowbreak{}fam}{Formalism.\allowbreak{}SharpRestriction.\allowbreak{}GlobalSchurBounds.\allowbreak{}GlobalRankLatticeProofs.\allowbreak{}SlopesMFOfCountFam}{17--23}

Here \(B\le2a/X'\), \(f=O(x^3)\), \(k\le x^5\), and \(n_T\ge1\), so both \(fkB\) and \(fk^2\) are at most \(n_Ta\) eventually. With \(m\le k,\ L\le L_a,\ \log k\le5L_a,\ \log(kv_E)\le6L_a\), the cost of \(S_1\) in units of \(n_TaL_a\) is at most \(2c_{\rm Count}m/k\). The terms multiplied by \(C_5,C_6\) cost at most \(10+c_{\rm Count}+c_L\) and \(12+c_{\rm Count}+c_L\), respectively. Taking \(c_{\rm Slope}=2c_{\rm Count}+\max(C_5,C_6,0)(12+c_{\rm Count}+c_L)\) suffices.
\end{proof}

\subsection*{Proof of \ref{B5.13}}\label{prf:B5.13}\phantomsection
\begin{proof}
Write \(v_{i,u}=v_{g_i}(D_{\Gamma,u})\). Since the loss is supported on \(S\), the fixed-prime formula in \ref{B4.03}, \(2F_g(j)=-D_j^{\rm Fock}\) from \ref{B4.02}, the weights of \ref{B4.52}, and clause 2 of \ref{B5.12k} give
\[
 \sum_{\substack{w\in S\\p_w=g_i}}c_w E_{ujw}
 =\bigl(2v_{i,u}-D_j^{\rm Fock}+2H_g+A_{uj}\bigr)\log g_i.
\]
Subtract \(S_a\log g_i\). Since \(v_{i,u},H_g\ge0\), clauses 3 and 4 bound the absolute difference by \(C_gB_0\). On the level fibre \(E_{ujw}=E^T_{u,w}\ge0\), so the absolute sum there is at most \(C_TB_0\). On the other fibres group by rational prime. The stabilizer term costs at most \(2\sum_p v_p(D_{\Gamma,u})\log p=2\log D_{\Gamma,u}\). The discriminant term costs at most \(2\log d\,|D_j^*|\), and the remaining terms cost at most \(\sum_{p\in P}2\log p\,|J^*_{j,p}|\) over eligible primes meeting \(S\). Clauses 6--8 bound their sum by \((C_{\rm stab}+C_d+C_I)B_0\). Add the level cost for the off-type bound and the two fixed-prime differences for the first inequality. Sum over \(j\) for \(\Ee_u\).
\end{proof}

\subsection*{Proof of \ref{B5.14}}\label{prf:B5.14}\phantomsection
\begin{proof}
Add \(S_a\log T_u\) to the height formula in \ref{B5.07}. The remaining terminal allowance is independent of the word. The hook, the two layer allowances and the mean metric correction are also common. Subtraction therefore gives the displayed identity.

By \ref{B4.02a}, every unpadded degree is at least one. By the coordinate formula there, changing the word changes a head degree by the change in the corresponding \(U_i\), and all tail degrees remain fixed. The endpoints are common to the words. At each interior head coordinate the two choices differ by at most \(\delta\le S_c\) and the degrees are at most \(w_{j\rho}+S_c\). Write \(M=6x^6b\). For \(1\le y,y'\le M\), integrating the derivative \(\log s-\log\Omega\) gives
\[
 \left|(\varphi(y)-y\log\Omega)-(\varphi(y')-y'\log\Omega)\right|
 \le |y-y'|(\log M+|\log\Omega|).
\]
The stated bounds on \(D_1\) and \(\log\operatorname{rd}E\), with \ref{B5.07}, give \(|\log\Omega|\le6+2L\). Since \(\log x\le L/\varepsilon\) and \(L\ge1\), the last factor is at most \((14+6/\varepsilon)L\). There are at most \(k\) varying coordinates per layer and position. The resulting bound is \((14+6/\varepsilon)mfkS_cL\le(14+6/\varepsilon)B_0\). This implies the claimed inequality for each member and hence for each translated member with its own pair of words.
\end{proof}

\section*{Proofs for Chapter~\ref{ch:prime-sums-net-gain}: Prime sums and net gain}

\subsection*{Proof of \ref{B5.15}}\label{prf:B5.15}\phantomsection
\begin{proof}
Group by rational prime. The model uses the full fiber on the
geometric support, and
\(\sum_{w\mid p}e_w\log N\mathfrak p_w/[K:\Q]=\log p\) by
\(N\mathfrak p_w=p^{f_w}\) and \(\sum_{w\mid p}e_wf_w=[K:\Q]\)
\cite[Chap.~I, Prop.~8.2]{neukirch-1999}. The values agree on the
fiber by \ref{B3.52}. Equation~\eqref{eq:p3c5-top-window} gives
\(a_t\le B\) since \(a_t\) is integral. Partners have weight
\(\ge w''\ge1\) by \ref{B5.17a} and \ref{B5.21a}, so all ordinary
terms above \(B\) vanish.

For the cutoff, a positive term has \(p^i<a_t/w_j(\rho)\le N\),
using \(d\) in place of \(p\) for a ramified term. The ordinary primes
in the arithmetic sum are those of the signal and inert classes, and
their exponents are odd. At \(d\), the coset \(iF_d+I_d\) has one
element outside \(\Phi\). In either case the partner
\(r_t(g)=c+g+((0,0),t)\) has first sign bit zero and hence a
unique matching place by \ref{B5.21a}. At an ordinary exponent \(i\)
its position has \(A\)-coordinate \(t+ih_p\) and the signal indicator
as its remaining bit. Thus the model filling and the presentation
read the same lift-\((j+1)\) weight there. The ramified terms read the
weight at the matching place of the ramified coset. Substitution in
the arithmetic sum gives the model sum.
\leanref{Formalism.\allowbreak{}MainNotions.\allowbreak{}Lnk.\allowbreak{}askIsF}{Formalism.\allowbreak{}SharpRestriction.\allowbreak{}GlobalSchurBounds.\allowbreak{}FrameCodeOfAdditive}{546--556}
\end{proof}

\subsection*{Proof of \ref{B5.17}}\label{prf:B5.17}\phantomsection
\begin{proof}
Geometric series give
\[
 A_{jtp}\le a_t p/(p^2-1)\le2a/p,\qquad
 D_{jt}\le a_t\log d/(d-1)\le2a\log d/d . \tag{5.17.2a}\label{5.17.2a}
\]
These give \(E_{\rm out},E_1,E_2,E_d\) in \eqref{5.16.3}. Above
\(X''\) any contribution comes from the first power. Its contribution is at
most \(w''\vartheta\), and it vanishes for \(p\ge X''(1+\vartheta)\).
The intervening interval has at most \(X''\vartheta+1\) integers.
This gives \(E_h\). For \(E_{\rm fl},E_{\rm cp}\) use \(|S|\le X''\),
\(|TT_p|\le n_T\), and \(1/\log p\le2\).

For \(Y\ge1\) one has
\[
 \Theta(Y)=\sum_{p\le Y}\log p\le Y\log4,\qquad
 \sum_{p\le Y}\frac{\log p}{p}\le\log Y+\log4 .
 \tag{5.17.6}\label{5.17.6}
\]
The first bound follows by induction using \(\binom{2r}{r}\le2^{2r}\)
and \(\binom{2r+1}{r}\le2^{2r}\) to count primes above \(r\) and
\(r+1\), respectively. For \(n=\lfloor Y\rfloor\), multiply the
second sum by \(n\) and use
\(\sum_{p\le n}\lfloor n/p\rfloor\log p+\Theta(n)
\le\log(n!)+\Theta(n)\).
The first bound gives \(E_{\rm pr}\). For \(E_{\rm or}\) use the
second bound, \(p^{-L_1(p)}\le1/(K_0p)\) below \(p_0\), and
\[
 \sum_{n\ge p_0}\frac{\log n}{n^2}
 \le\frac{2(\log(p_0-1)+1)}{p_0-1}, \tag{5.17.7}\label{5.17.7}
\]
by integrating \((\log y)/y^2\) from \(p_0-1\).
\end{proof}

\subsection*{Proof of \ref{B5.18}}\label{prf:B5.18}\phantomsection
\begin{proof}
For odd \(s>1\) and \(0\le j\le e\),
\[
 h_e(s)^{2^j}=1\iff s\equiv\pm1\pmod{2^{e+2-j}} .
 \tag{5.18.3}\label{5.18.3}
\]
For \(j\ge1\), successive differences of squares give
\(v_2(s^{2^j}-1)=v_2(s-1)+v_2(s+1)+j-1\), and one of \(s\pm1\)
has valuation one. For \(j=0\) use the quotient definition. With
\(r=\min(e,\lfloor\log_2J_0\rfloor)\) and \(M=2^{e+2-r}\), the
small-order primes have the form \(qM\pm1\) with \(q\ge1\).
Since \(p\ge qM/2\), their sum is at most
\(4\log Y_0(1+\log Y_0)/M\). Use \(2^e\le J_0M\).
\leanref{Formalism.\allowbreak{}RankGapBounds.\allowbreak{}LocalQuadraticParts.\allowbreak{}ordersSumY\_\allowbreak{}of}{Formalism.\allowbreak{}RankGapBounds.\allowbreak{}RankGapProofs.\allowbreak{}PmModulusFacts}{252--287}
\leanref{Formalism.\allowbreak{}RankGapBounds.\allowbreak{}LocalQuadraticParts.\allowbreak{}powTwoPowEqOneIff\_\allowbreak{}holds}{Formalism.\allowbreak{}RankGapBounds.\allowbreak{}RankGapProofs.\allowbreak{}FourNotDvdOddSqAddOne}{132--138}

In the band \(2i\le\log(2X')\le2\log X'\), and taking \(i\)-th
roots puts \(p\) in the interval \(y_i[1-\rho,1+\rho]\), where
\(y_i=(X')^{1/i}\). The interval has at most \(2\rho y_i+1\) integers. Each
prime contributes at most \(\log(2X')\min(2/y_i,1/p_0)\).
Summation proves the band bound.
\end{proof}

\subsection*{Proof of \ref{B5.20}}\label{prf:B5.20}\phantomsection
\begin{proof}
For \(m=0\), use \(\tau_p\le0\) on \(S\) from \eqref{5.16.1} and
\eqref{5.19.2} with \(m=0\). For \(m>0\) put \(E=x^\nu\),
\(S_x=e^{x^{\varepsilon/4}}\), \(B=e^{z/2}\), and \(D_0=100K_sE^4x^4\).
Thus
\[
 X'\ge D_0,\quad B\ge D_0X'S_x,\quad b=B^2,\quad
 x\ge100,\quad 0\le\sigma_x\le1 . \tag{5.20.8}\label{5.20.8}
\]
In particular \(\log x,\nu\log x\le u\le z,\ z^2\le x\), and
\[
 \log a\le2z,\qquad \log(2a)\le x . \tag{5.20.9}\label{5.20.9}
\]
The bounds for \(E_{\rm fl}/Q,E_{\rm cp}/Q,E_{\rm pr,low}/Q\)
from \eqref{5.17.5} are respectively
\[
 12K_sE^4x^3X'/b,\quad 6K_sE^2x^2X'S_x/b,\quad
 3K_sE^2x^2/X'. \tag{5.20.10}\label{5.20.10}
\]
Each is at most \(\sigma_x\). For the upper price part,
\(X'Q_2/b-C\) is at most the sum of
\(2K_sX'S_x/b,\ K_sz^2/x^2,\ K_sX'Ex/b,\ X'/B\).
Each is at most \(\sigma_x/4\), so
\[
 E_{\rm pr,high}/Q\le2C\log4+3\sigma_x .
 \tag{5.20.12}\label{5.20.12}
\]
In \(E_{\rm or}\) the three factors inside braces in
\eqref{5.17.5} use \(\mathcal L/k\le z^2/x^2\),
\(K_0^{-1}\le x^{-2}\), \(\log X'+\log4\le3z\), and a tail at
most \(4z/x^2\). This gives \(E_{\rm or}/Q\le20K_sz^4/x\).
In \ref{B5.18} use \(J_0,\lfloor\log X'\rfloor\le z\),
\(\log(2^{e_2}\max(Y,X''))\le4z\), and \(2^{e_2}\ge x^2\).
Then
\[
 E_{\rm ord}/Q\le64K_sz^3/x,\qquad E_{\rm band}/Q\le40z^2/x .
 \tag{5.20.13}\label{5.20.13}
\]
The remaining four bounds are
\[
 E_{\rm out}/Q\le4z/x,\quad E_2/Q\le2x^{-8},\quad
 E_h/Q\le16z/x,\quad E_d/Q\le4z/x .
 \tag{5.20.14}\label{5.20.14}
\]
Here the first two use \ref{B5.19}. For the last two use
\(\log(2X'')\le4z\), \(a\vartheta^2,w''\vartheta\le b/x\), and
\(\log d/d\le2\log x/x^2\). Add the ten bounds by \(\sigma_x\) to
\eqref{5.20.12}. The index bound follows from
\eqref{5.19.3} and \(amn_Tx^{-9}=x^{-8}Q\).
Finally \(x^{4\varepsilon}=o(x)\), \(\log x=o(u)\), and
\(u+x^{\varepsilon/4}=o(z)\).
\end{proof}

\subsection*{Proof of \ref{B5.20d}}\label{prf:B5.20d}\phantomsection
\begin{proof}
With \(d_I=\det H_I\), each minor is
\(d_I^{p^{z_j}}\det(T_{jli}^{p^{z_j-u_{j,l}}})\). If \(d_{I_0}=0\),
the initial factor vanishes. Otherwise row-reduce the first
block to \(I_r\). On a competing plane a replacement minor is
nonzero and its factor vanishes because
\(d^{\sum_j p^{z_j}(Q_j-1)}=d^{p^{z_b}-p^{z_0}}=1\) for \(d\in F^\times\).
On the coordinate plane each replacement vanishes and the
\(\Delta_j\) have unit diagonal terms. Lift cyclic positions
by a common multiple of \(n'\) when necessary. A Witt lift
of \(L^h\) has order at least \(h\) on the competing planes.
\end{proof}

\subsection*{Proof of \ref{B5.20f}}\label{prf:B5.20f}\phantomsection
\begin{proof}
In the numerator of \(\binom{u+p^hZ}{i}\) at least
\(\lfloor i/p^\ell\rfloor\) factors have both coefficients in
\(p^\ell\Z_p\) for \(1\le\ell\le h\). Thus the polynomial is integral.
In Smith coordinates at \(u\), the Gauss norm on \(|Z|\le p^{-E}\)
of the determinant is at most \(p^{-E\beta_u}\), by the row bounds.
Multiplicativity of the Gauss norm and the roots \(\zeta\) of the
determinant give
\[
 E\beta_u\le v(b)+\sum_\zeta\min(v(u-\zeta),h+E).
\]
At heights \(\ell-1<y\le\ell\), the integers with
\(v(u-\zeta)\ge y\) lie in one residue class modulo \(p^\ell\).
Integration to height \(h\) costs at most \(H/(p-1)+h\) per root,
and above \(h\) at most one point contributes, at cost \(E\).
Use \(s\le dJ\). Removing the integral scalar subtracts its
valuation from the finite diagonal valuations.
\end{proof}

\subsection*{Proof of \ref{B5.20m}}\label{prf:B5.20m}\phantomsection
\begin{proof}
Write \(P_k(b)=\binom{b+k-1}{k-1}\). The total count is \(R=\prod_i P_k(b_i)\). Indeed, tensoring an irreducible with one symmetric power adds a horizontal strip with multiplicity one. Dualizing gives the removal rule. Restriction to the subgroup fixing one coordinate gives one copy of each interlacing row. These two dimension identities also follow by grouping semistandard tableaux, respectively by the horizontal strip in the Pieri rule and by the boxes occupied by the largest letter in the branching rule. Evaluating dimensions after the last restriction gives precisely the labels in \ref{B5.20l}. Negative entries are handled by a common determinant twist before restriction.

Keep two lists at each stage. The first consists of the degrees of the vector letters already used, in reverse chronological order, and the second is the corresponding list for the covectors. At a restriction step both lists stay fixed. For the current row form two decreasing lists \(u^+,u^-\) of its positive parts and of the absolute values of its negative parts. Extend lists by zeros. If \(B\) is the list of degrees of the corresponding sign, then for every \(h\ge0\)
\[
 \sum_{j>h}u_j\le \sum_{j>h} B_j.                                      \tag{1}
\]
For example, on adding a strip of size \(b\) the sum of positive parts increases by at most \(b\), and interlacing gives \(u^+_{j+1,\mathrm{new}}\le u^+_{j,\mathrm{old}}\). These are the inequalities for \(h=0\) and \(h\ge1\) after \(b\) is prepended to \(B\). Negative parts decrease coordinatewise. Reversing and negating gives the covector step. In a restriction both new lists are coordinatewise bounded by the old ones. This proves (1) by induction. It also proves that a deficit \(X\) in (1) persists, with \(h\) increased by one each time a subsequent signed letter of the corresponding sign is used.

For either sign let \(B_1,\ldots,B_v\) be its final list of degrees. Inequality (1) at \(h=v\) makes all subsequent \(u_j\) zero. Put
\[
 F(y)=\sum_{d=1}^q\log(1+y/d),\qquad
 T_h=\sum_{j=h}^{v}(u_j-B_j).
\]
Concavity and summation by parts give
\[
 \sum_j\{F(u_j)-F(B_j)\}
 \le \sum_{h=1}^{v}\{F'(B_h)-F'(B_{h-1})\}T_h,\qquad F'(B_0):=0.
                                                                    \tag{2}
\]
Every \(T_h\le0\). Since \(B_h\ge k\ge q\) and \(B_{h-1}\ge4B_h\) when \(h>1\),
\[
 F'(B_h)\ge q/(2B_h),\qquad F'(B_{h-1})\le F'(B_h)/2.
\]
Consequently a deficit \(X\) beginning at the entry \(B_h\) costs at least \(qX/(4B_h)\) in (2).

For a failed positive- or negative-part test, the loss of mass creates a deficit of at least \(\xi_i\) at the start of the appropriate list. The list at that point is nonempty, since mass was lost. Its first degree is an earlier signed degree and is at most \(4v_i/p\). Its position in the final list is the one specified by persistence in (1). The cost is at least \(qL/(16k)\). For a failed test at coordinate \(r\), omit the already restored positions from the row. In the positive degree list the corresponding entry is \(b_{o(i)}\). Apply (1) to the tail strictly after that entry and then include the entry itself. The deficit is at least \(b_{o(i)}/3\), so the cost is at least \(q/12\).

For the final row there are at most \(t\) positive and \(t\) negative entries, by interlacing from zero. In its Weyl dimension formula reserve the middle \(K-2t=q\) zero coordinates. Pairing one exterior entry of absolute value \(y\) with these coordinates contributes at most \(\exp F(y)\), since their distances from it are at least \(1,\ldots,q\) in increasing order. The remaining pairs lie among \(2t\) coordinates. Each factor is at most \(1+2at\), since the spread is at most the sum of all signed-letter degrees. Also \(P_k(b)\ge\exp F(b)\). Applying (2) to both signs therefore bounds the number of labels on one discarded path by
\[
 R\exp\left(9t^2\log(1+2at)-\min(qL/(16k),q/12)\right).
\]
Here \(\binom{2t}{2}\le9t^2\) suffices.

At each stage record the first \(t\) positive parts and the last \(t\) negative parts. These determine the row, including the shorter rows at restriction steps. Each recorded integer lies between \(0\) and \(2at\). The initial row is fixed, so there are at most \((1+2at)^{2t^2}\) paths. Multiply the preceding bound by this number and use \(q\ge k/2\).
\end{proof}

\subsection*{Proof of \ref{B5.20ha}}\label{prf:B5.20ha}\phantomsection
\begin{proof}
Consider one coordinate, with power \(P>1\), digit \(d\), and quotient \(b\). Splitting the factorial into consecutive intervals of length \(P\) gives
\[
 \frac{(d+Pb)!}{d!\,b!\,(P!)^b}
 =
 \prod_{h=0}^{b-1}\prod_{j=1}^{P-1}\left(1+\frac Pj h\right)
 \prod_{j=1}^{d}\left(1+\frac Pj b\right).
\]
Indeed the multiples \(P,2P,\ldots,bP\) contribute \(P^b b!\). For \(1\le j<P\), \(v_p(j)<v_p(P)\). Every displayed factor belongs to \(1+p\Z_p\). The factor for \(P=1\) is one.

Write \(t_j=P/(pj)\in\Z_p\). Replace \(p\) by an indeterminate \(T\) in the factors \(1+pt_jh\) and \(1+pt_jb\), keeping \(t_j\) fixed. Fix a sign \(\epsilon\in\{1,-1\}\). The coefficient of \(T^a\) in \(\prod_{j<P}(1+Tt_jh)^\epsilon\), with \(j\) starting at one, is a polynomial \(V_a(h)\) of degree at most \(a\) with coefficients in \(\Z_p\). For the negative sign use the formal geometric series. In particular \(V_0=1\).

Let \(H_e(b)\) be the coefficient of \(T^e\) in the product over \(0\le h<b\). Then \(H_0=1\), \(H_e(0)=0\) for \(e>0\), and
\[
 H_e(b+1)-H_e(b)=\sum_{a=1}^{e}V_a(b)H_{e-a}(b).
\]
By induction \(H_e\) is an integer-valued polynomial on \(\N\) of degree at most \(2e\). Indeed the right side has degree at most \(2e-1\), and summation from \(0\) to \(b-1\) takes \(\binom hj\) to \(\binom b{j+1}\). Integrality of its values also follows directly from the product. Multiplying by the last \(d\) factors preserves the degree bound \(2e\) for the coefficient of \(T^e\).

In the product over \(I\), the indices contributing to \(T^e\) sum to \(e\), so the total degree of that coefficient is at most \(2e\), independently of \(|I|\). At \(T=p\) the terms with \(e\ge c\) vanish modulo \(p^c\). The remaining polynomial takes values in \(\Z_p\) on \(\N^I\). Its coefficients in the product binomial basis are its iterated forward differences at zero and belong to \(\Z_p\). Its degree is at most \(2(c-1)\le4c\).
\end{proof}

\subsection*{Proof of \ref{B5.20h}}\label{prf:B5.20h}\phantomsection
\begin{proof}
Fix an incoming edge \(e\), write \(K=k-r\), \(c'=c_*+o\), and
normalize \(v(p)=1\), \(v(\omega)=1/2\). A source family on \(e\) has target tests for its primitive
combinations and satisfies the bounds
\[
\begin{split}
 T&=C_tN\left(CN_1^2\log(a_0)\frac{a_0}{pk}
 +a_0p^{-L_1}+kp^{G'}\right),\\
 t_S&\le T+o,\quad t_T\le T+Nc_t+o,\quad
 D_e\le C_DNk^2\log_p a_0,\\
 n_e&\ge (1-C_rNr^2k/Y)(1-kc_t/a_0)^N
 \dim V\,p^{-n'rK}.
\end{split}\tag{5.20.16}\label{eq:p5c2-plane-edge}
\]
It suffices to treat \(CNr^2k<Y\).
\leanref{Formalism.\allowbreak{}MainNotions.\allowbreak{}Pf14.\allowbreak{}EdgeStep}{Formalism.\allowbreak{}SharpRestriction.\allowbreak{}GrowthBounds.\allowbreak{}ScaleEdges.\allowbreak{}ManyEdges}{68--84}
\leanref{Formalism.\allowbreak{}MainNotions.\allowbreak{}Pf14.\allowbreak{}edgeStep\_\allowbreak{}of}{Formalism.\allowbreak{}SharpRestriction.\allowbreak{}GrowthBounds.\allowbreak{}ScaleEdges.\allowbreak{}InitialLevelMonoWeights}{30--180}

Number the base-weight rows from \(1\). At a mark split off
\(d_s=\lambda_{s,1}-\lambda_{s,2}\). Let
\(B_k(d)=\binom{d+k-1}{k-1}\), \(P=\prod_{s\in S}B_k(d_s)\), and
\(B\) be the product of the residual dimensions. The Weyl formula
\cite[Lecture~24]{fulton-1991} gives at each mark the ratio
\[
 \prod_{j=2}^k\frac{d_s+j-1}{d_s+\lambda_{s,2}-\lambda_{s,j}+j-1}
 \ge1-kc_t/a_0.
\]
Here \(d_s\ge a_0\) and the residual spread is at most \(c_t\).
Negated reversal at odd marks preserves dimension. Hence
\(PB\ge(1-kc_t/a_0)^N\dim V\).
\leanref{Formalism.\allowbreak{}MainNotions.\allowbreak{}Pf14.\allowbreak{}top\_\allowbreak{}count}{Formalism.\allowbreak{}SharpRestriction.\allowbreak{}GrowthBounds.\allowbreak{}ScaleEdges.\allowbreak{}RemainingWeightVectorCount}{30--81}
\leanref{Formalism.\allowbreak{}MainNotions.\allowbreak{}TestSideCounting.\allowbreak{}card\_\allowbreak{}all}{Formalism.\allowbreak{}SharpRestriction.\allowbreak{}GrowthBounds.\allowbreak{}DualFockLinear.\allowbreak{}TestSideCounting}{321--397}

Lift the marks periodically to \(\Z\). In each marked tile take the
middle \(r\) of its \(3r\) even marks, \(f_{j,0}<\cdots<f_{j,r-1}=z_j\).
If there are \(J_f\) marked tiles in a period, extend by
\(f_{j+J_f,i}=f_{j,i}+n'\). Segment \(j\) contains the marks in
\([f_{j-1,0},z_j]\), with the \(r\) marks of frame \(j-1\) as its
left frame. The gaps in a marked tile and \(L_2<u_1\) separate its
middle frame from the adjacent marks. Thus every letter after the
left frame has \(r\) even predecessors in the segment. The oldest
is more than \(L_1\) steps back, since otherwise a window of length
\(L_2\) would contain those \(r\) marks. The full-tile condition
bounds each segment length by \(G'\), including the segment crossing
the incomplete tile, and each has at most \(N_s=\min(N_1,2N)\) letters.
\leanref{Formalism.\allowbreak{}MainNotions.\allowbreak{}Pf14.\allowbreak{}frames\_\allowbreak{}of}{Formalism.\allowbreak{}SharpRestriction.\allowbreak{}GrowthBounds.\allowbreak{}ScaleEdges.\allowbreak{}FramesOf}{25--57}

At a letter use \(r\) plane and \(K\) complementary variables.
Use plane variables on the left frame. The normal variables are the
complementary ones at even letters and the plane ones at odd letters.
The segment substitution raises each variable at \(x\) to
\(M_x=p^{z_j-x}\) at its position modulo \(n'\). For a set of letters
\(I\), let \(v_I,l_I\) be the least and greatest exponents occurring
under a generic conjugate of \(\operatorname{diag}(T,1,\ldots,1)\)
on \(I\), with inverse transpose on covectors. Use \(+\infty,-\infty\)
for zero. Set \(Q_*=p^{G'}\),
\(\alpha_j=\mathbf1_{j=0}+rK(p^{z_{j+1}-z_j}-1)\), and
\[
 m_*=(N_1+1)((2Q_*+3)c'+kQ_*+c'\max_j\alpha_j+3c').
\]
For a non-left-frame letter \(x\) let \(y\) be its oldest of \(r\)
even predecessors, \(b_x=d_{y\bmod n'}/M_y\), and
\(m_x=m_*+\lfloor b_x/(k^2Y)\rfloor\). Span products \(A B'\) of a
plane and a complementary polynomial with \(v_{<x}(A)\ge b_x/2\).
At even \(x\) require \(l_{\le x}(B')+m_x\le v_{<x}(A)\), and at odd
\(x\) require \(l_{<x}(B')+m_x\le v_{\le x}(A)\). Impose the normal
degree bound
\[
 M_x^{-1}\left(C_{\mathcal L}N_s^2\log(a_0)\frac{a_0}{pk}
 +\epsilon_xa_0p^{-L_1}\right),
 \qquad \epsilon_x=10\ \text{at even }x,\quad \epsilon_x=1\ \text{at odd }x .
 \tag{5.20.16a}\label{eq:p5c2-symbol-bound}
\]
Here \(C_{\mathcal L}\) is absolute. A symbol is a sum of products of
the substituted segment polynomials and residual monomials with
exponents below \(M_x\) at the non-left-frame occurrence. Its
multidegree is what remains after reserving the fixed determinant
factor below. The left frame uses no normal variables, and the
residual monomial contributes at most \(kQ_*\) per row. Thus
\eqref{eq:p5c2-symbol-bound} gives the bound \(T\).
\leanref{Formalism.\allowbreak{}MainNotions.\allowbreak{}Pf14.\allowbreak{}symbolSpaces}{Formalism.\allowbreak{}SharpRestriction.\allowbreak{}GrowthBounds.\allowbreak{}ScaleEdges.\allowbreak{}SymbolRealisableBound}{10--22}
\leanref{Formalism.\allowbreak{}MainNotions.\allowbreak{}SelectionLemmas.\allowbreak{}Pf.\allowbreak{}gen\_\allowbreak{}bound}{Formalism.\allowbreak{}SharpRestriction.\allowbreak{}GrowthBounds.\allowbreak{}SegmentDegree.\allowbreak{}StrictMonoIntExists}{524--588}

For the count put
\[
 R_j=\left\lfloor\log_p\left\lfloor
 \frac{a_0}{8Yk^2N_sp^{z_j-f_{j-1,0}}}\right\rfloor_+\right\rfloor,\quad
 M^c_{j,x}=p^{z_j-x+R_j},\quad Q_j=p^{z_{j+1}+R_{j+1}-z_j-R_j}.
\]
Use zero for the rounded logarithm of zero. The shifted endpoints
are the maxima of \(z_j\) and
\(f_{j-1,0}+\lfloor\log_p(a_0/(8Yk^2N_s))\rfloor\), so increase and
have total increment \(n'\) per period. In particular
\(\prod_jQ_j=p^{n'}\). In addition to the fixed lifting factor in \eqref{eq:p5c2-fixed-factor}, reserve a determinant for each non-left-frame occurrence \(x\). In the unpowered segment variables it is the determinant of the last \(r\) preceding vector letters on their plane coordinates, raised to \(s_x=\lfloor b_x/(k^2Y)\rfloor\). Apply the unshifted segment substitution to these determinants and multiply over occurrences. At a mark \(y\) their total degree is \(\sum_{j,x}s_x\sum_{v}\mathbf1_{v\equiv y\bmod n'}p^{z_j-v}\), where \(v\) runs through those \(r\) predecessors in segment \(j\). Let \(a'_y\) be the marked degree after subtracting both reserves. The second reserve is used when converting the coarse symbols back to the unshifted prefix spaces. Its reduction is a product of nonzero determinants. At each non-left-frame occurrence
sum over coordinate residues \(0\le u_c<M^c_{j,x}\) with
\(\sum_cu_c\le a'_x\) and \(\sum_cu_c\equiv a'_x\bmod M^c_{j,x}\).
The remaining degree is \(h=(a'_x-\sum_cu_c)/M^c_{j,x}\).
At a right-frame slot choose \(0\le e_{ji}<rQ_j\) in the class of
\(h\) modulo \(Q_j\). The left-frame degree of the next segment
is \((h-e_{ji})/Q_j\).
\leanref{Formalism.\allowbreak{}MainNotions.\allowbreak{}SelectionLemmas.\allowbreak{}Pf.\allowbreak{}Dat.\allowbreak{}Qc}{Formalism.\allowbreak{}SharpRestriction.\allowbreak{}GrowthBounds.\allowbreak{}SegmentDegree.\allowbreak{}Coarse}{43--45}
\leanref{Formalism.\allowbreak{}MainNotions.\allowbreak{}SelectionLemmas.\allowbreak{}Pf.\allowbreak{}Dat.\allowbreak{}bfam}{Formalism.\allowbreak{}SharpRestriction.\allowbreak{}GrowthBounds.\allowbreak{}SegmentDegree.\allowbreak{}RestCountOf}{43--51}

For fixed digits use \ref{B5.20l}. Before the right frame assign degrees \(b_i=h_i\) to the interior letters and \(b_i=[B_w-m_0]_+\) to the left-frame letters. At the successive right-frame letters perform the \(r\) restriction steps. Equivalently, at the \(j\)-th such step adjoin \(j\) leading entries \(b_*=\sum b_i\) to the restricted row. This gives a vector strip of uniquely determined degree. On the first \(r\) coordinates use the highest tableau and on the last \(K\) use the freely chosen label of \ref{B5.20l}.
The path polynomials are products of the exterior columns
described below, with monomial labels in the boundary coordinates.
Set the designated complementary variables to zero at the
interior letters and the designated non-plane coordinates at the
frame letters. On the resulting polynomials, a column's diagonal
term puts its entries in height order. In the product the exponent
differences at each letter recover the strips of the path, and
the remaining frame exponents recover its labels. The leading
coefficients are units, so these terms give a triangular
coefficient matrix over the residue field. On the left frame
multiply by the determinant of the preceding frame to power
\(m_0=kN_s+1\) and by the compensating first-coordinate monomials
of exponents \(h_{\rm left}-B_w\), where
\[
 B_w=\left\lfloor a'_x/M^c_{j,x}\right\rfloor-(k+r+1).
\]
The left-frame degree bound is \(h_{\rm left}\ge B_w\).
These factors are common within the degree class and their
reductions are nonzero, so independence is preserved. For products
over segments, the exterior-column tensors are independent for
fixed classes. Distinct residue arrays are recovered before
quotienting by each \(M^c_{j,x}\). On a shared frame the first \(r\)
exponents of the right-frame polynomial are below \(Q_j\).
Their residues recover the right-frame part and the quotients
recover the left-frame part. The other \(K\) coordinates of the
shared frame belong to the right-frame part. This gives
independence across the boundary choices.
\leanref{Formalism.\allowbreak{}MainNotions.\allowbreak{}SelectionLemmas.\allowbreak{}Pf.\allowbreak{}restInd\_\allowbreak{}of}{Formalism.\allowbreak{}SharpRestriction.\allowbreak{}GrowthBounds.\allowbreak{}SegmentDegree.\allowbreak{}ScaleSCoeffZero}{162--188}

Apply \ref{B5.20m} with \(t\) the number of segment letters, \(v_i=a_0/M^c_{j,i}\), and \(L=C_{\mathcal L}N_s^2\log a_0\). The left frame and the last \(r\) letters before the right frame are vector letters by the choice of the middle marks of a marked tile. Earlier letters of the same parity are at least two positions back. The divided degrees before the right frame lie between \((2/5)a_0/M^c_{j,i}\) and \(2a_0/M^c_{j,i}\): subtracting the reserved degree leaves at least \(9a_0/10\) before division, and the digit sum and left-frame reserve after division cost at most \(k+r+2+m_0\). The definition of \(R_j\) gives \(a_0/M^c_{j,i}\ge 8Yk^2N_s\) when the rounded logarithm is positive. When it is zero, use \(M^c_{j,i}\le p^{G'}\) and the last size hypothesis instead. These bounds give \(b_i\ge k\), the growth factor \(p^2/5\), and the earlier-degree bound \(4v_i/p\). Since \(t\le N_s\), enlarging \(C\) gives \(K-2t\ge k/2\). At a covector the plane degree removed is at most the total decrease of positive parts, since the first \(r\) entries are nonnegative. At a vector, interlacing bounds the increase of positive parts in the complementary row by the preceding plane entry at position \(r\). This entry is at most \(2b_{o(i)}\): interlace back across the intervening vector strips and sum vector degrees through the oldest one, using the geometric growth. The decrease in absolute negative parts is at most \(\xi_i\). The retained paths thus have complementary increase at most \(2b_{o(i)}+\xi_i\) at a vector and plane decrease at most \(\xi_i\) at a covector, as required. Finally \(1+2a_0t\le a_0^2\) and \(\min(L,k)\ge1088N_s^2\log a_0\). The exceptional fraction in \ref{B5.20m} is at most \(a_0^{-N_s^2}\). The unrestricted labelled count there supplies exactly the asserted product of binomials. For a shared slot the path determines a strip degree \(d\le2b_*\) by interlacing: after deletion of one leading coordinate the change in the shorter row sum lies between its former smallest and largest entries, both in \([-b_*,b_*]\). The reserve satisfies \(4Yb_*\le Q_j\). Indeed, with \(D=z_j-f_{j,0}\), summing the distinct powers from letters before the right frame gives \(b_*\le4a_0p^{-R_j-D-L_1-1}\). The definition of the next rounded logarithm gives \(p^{R_{j+1}}>a_0/(8pYk^2N_sp^{z_{j+1}-f_{j,0}})\). Dividing by \(Q_j\) gives \(4Yb_*/Q_j\le128Y^2k^2N_sp^{-L_1}\le1\). Apply \ref{B5.20n} with \(b=b_*\). The exponent array from that lemma is the boundary monomial label, and \(e_{ji}=d+\sum v_i\). Its contribution is at least \(Q_j^{r-1}(1-r/(2Y))\) per slot.
Thus the comparison factor for a shared row, after restoring
the divided scale and moving its left-frame count to the
preceding segment, is
\[
 B_k([B_w-m_0]_+)\,Q_j^{r-1}(1-r/(2Y))(M^c_{j,x})^{k-1}Q_j^K ,
\]
where \(B_w\) uses segment \(j+1\). For an interior row of reduced degree \(d=a'_x\) and scale \(M=M^c_{j,x}\), the sum needed here is exactly
\[
 \sum_{\substack{0\le u_c<M,\ |u|\le d\\ |u|\equiv d\pmod M}} B_k((d-|u|)/M)=B_k(d).
\]
Indeed, coordinatewise Euclidean division gives a bijection from the exponent vectors of degree \(d\) to the pairs counted on the left. At a shared slot only the count of residues is needed. There are \(M^{k-1}\), since the last digit is determined by the other digits and the congruence, and \(kM\le a'_x\) makes every such choice admissible. For fixed residue rows the allowed boundary totals are chosen separately at each right frame. The lower bound at the next left frame uses \(B_w\), independent of the preceding choice of total. Thus summing the product over boundary totals gives the product of the segment sums. Summing next over the residue rows factors into the row sums above.

Here are the two dimension comparisons. Write \(d_x\) for the original degree. In the fixed factor \eqref{eq:p5c2-fixed-factor}, a determinant contributes degree at most \(Q_*\) in a given row. There are at most \(N\) predecessor determinants and a row occurs in at most one frame determinant. Since \(\alpha_j\le1+rkQ_*\), the charge of this factor is at most \(Q_*N((2Q_*+3)c'+kQ_*)+Q_*(2+rkQ_*)c'\), which is at most \(a_0/(kY)\) by the last size hypothesis, on enlarging the absolute constant. For the other reserve, a row at \(y\) among the predecessors contributes \(s_xp^{z_j-y}\le2a_0/(k^2Y)\), since the oldest predecessor is at or before \(y\). There are at most \(N_s\) such occurrences per segment and at most two segments containing the row. Hence \(a'_x\ge d_x-5a_0/(kY)\). In the case under consideration \(Y>CNr^2k\), so in particular \(a'_x\ge9a_0/10\). Each factor of the binomial product for \(B_k(a'_x)/B_k(d_x)\) is at least \(1-5/(kY)\), and their product is at least \(1-5/Y\).

For a shared row put \(M'=M^c_{j,x}Q_j\), its scale in the next segment. From the definition of the rounded logarithm, \(M'\le\max(p^{G'},a_0/(8Yk^2N_s))=a_0/(8Yk^2N_s)\), using the last size hypothesis for the equality. Thus \(a'_x/M'\ge(36/5)Yk^2N_s\). For \(b=[B_w-m_0]_+\), the floor and the two subtractions give \(a'_x-M'b\le(k+r+2+m_0)M'\le4kN_sM'\le a'_x/(kY)\). For \(1\le i<k\) it follows that \(M'(b+i)/(a'_x+i)\ge1-1/(kY)\). Multiplication gives \((M')^{k-1}B_k(b)\ge(1-1/Y)B_k(a'_x)\). The scale restoration has
\(r-1+K=k-1\) at each of the \(r\) shared slots. Put \(A_0=1-a_0^{-N_s^2}\). The exceptional fraction in the segment count of \ref{B5.20m}, with the parameters just checked, is at most \(a_0^{-N_s^2}\) in each of the \(J_f\) segments. Removing \(Q_j^{rK}\) from each segment's scale factor and multiplying gives retained dimension at least
\[
 \frac{P}{p^{n'rK}}A_0^{J_f}
 \bigl((1-5/Y)(1-1/Y)(1-r/(2Y))\bigr)^N .
 \tag{5.20.16b}\label{eq:p5c2-symbol-count}
\]
Indeed \(N_s\ge1\) and the size hypotheses give \(a_0\ge32Y\ge1\), so \(a_0^{N_s^2}\ge32Y\) and \(1-1/(32Y)\le A_0\le1\).
\leanref{Formalism.\allowbreak{}MainNotions.\allowbreak{}SelectionLemmas.\allowbreak{}Pf.\allowbreak{}Dat.\allowbreak{}firstSumFactor}{Formalism.\allowbreak{}SharpRestriction.\allowbreak{}GrowthBounds.\allowbreak{}SegmentDegree.\allowbreak{}FamilyDegreesAdmissible}{43--43}
\leanref{Formalism.\allowbreak{}MainNotions.\allowbreak{}SelectionLemmas.\allowbreak{}Pf.\allowbreak{}first\_\allowbreak{}sum\_\allowbreak{}factor\_\allowbreak{}bounds}{Formalism.\allowbreak{}SharpRestriction.\allowbreak{}GrowthBounds.\allowbreak{}SegmentDegree.\allowbreak{}FamilyDegreesAdmissible}{139--150} Since \(J_f\le N\)
Since \(J_f\le N\), Bernoulli's inequality gives \(A_0^{J_f}\ge1-N/(32Y)\). For factors \(1-u_i\) in \([0,1]\), induction gives \(\prod_i(1-u_i)\ge1-\sum_i u_i\). Applying this directly to all the factors in \eqref{eq:p5c2-symbol-count} bounds their product below by \(1-N(1/32+5+1+r/2)/Y\ge1-Nr^2k/Y\), since \(k>r\ge2\).
\leanref{Formalism.\allowbreak{}MainNotions.\allowbreak{}SelectionLemmas.\allowbreak{}Pf.\allowbreak{}restSeg\_\allowbreak{}of}{Formalism.\allowbreak{}SharpRestriction.\allowbreak{}GrowthBounds.\allowbreak{}SegmentDegree.\allowbreak{}GoodPathWeight}{108--139}
\leanref{Formalism.\allowbreak{}MainNotions.\allowbreak{}SelectionLemmas.\allowbreak{}Pf.\allowbreak{}secondPartClosed}{Formalism.\allowbreak{}SharpRestriction.\allowbreak{}GrowthBounds.\allowbreak{}SegmentDegree}{33--36}
\leanref{Formalism.\allowbreak{}MainNotions.\allowbreak{}SelectionLemmas.\allowbreak{}Pf.\allowbreak{}restSum\_\allowbreak{}of}{Formalism.\allowbreak{}SharpRestriction.\allowbreak{}GrowthBounds.\allowbreak{}SegmentDegree.\allowbreak{}SumTotalLoss}{101--249}
\leanref{Formalism.\allowbreak{}MainNotions.\allowbreak{}SelectionLemmas.\allowbreak{}Pf.\allowbreak{}sum\_\allowbreak{}total\_\allowbreak{}loss}{Formalism.\allowbreak{}SharpRestriction.\allowbreak{}GrowthBounds.\allowbreak{}SegmentDegree.\allowbreak{}SumTotalLoss}{11--54}

Choose a Witt matrix \(g\) for the edge. Write \(g_x\) for its
\(x\)-th Frobenius conjugate. In adapted coordinates the edge
substitution \(s_g\) scales the normal entries by \(\omega\).
In the original coordinates its row substitutions are
\(X_x\operatorname{diag}(I_r,\omega I_K)g_x^{-1}\) at even
letters and \(X_x\operatorname{diag}(\omega I_r,I_K)g_x^{\mathsf t}\)
at odd letters. Let \(\Delta_j\) be the first \(r\)-column minor
of frame \(j\), with row powers \(p^{z_j-x}\). Let \(D_x\) be the
first \(r\)-column determinant of the \(r\) preceding even rows
in reverse order, with row powers \(p^{x-y}\). With
\(R_0=\prod_j\Delta_j^{\alpha_j}\) the fixed factor is
\[
 U=R_0^{c'}\prod_j\Delta_j^{c'}\prod_{x\in S}D_x^{(2Q_*+3)c'+kQ_*}.
 \tag{5.20.16c}\label{eq:p5c2-fixed-factor}
\]
It has normal degree zero and nonzero reduction, so multiplication
by its reduction is injective. Choose symbols giving a basis of
the reduced span in \eqref{eq:p5c2-symbol-count}. Their coefficient
matrix has a unit maximal minor.
\leanref{Formalism.\allowbreak{}MainNotions.\allowbreak{}SelectionLemmas.\allowbreak{}Pf.\allowbreak{}restRank\_\allowbreak{}of}{Formalism.\allowbreak{}SharpRestriction.\allowbreak{}GrowthBounds.\allowbreak{}SegmentDegree.\allowbreak{}Coarse}{128--160}
\leanref{Formalism.\allowbreak{}MainNotions.\allowbreak{}SelectionLemmas.\allowbreak{}Pf.\allowbreak{}restSel}{Formalism.\allowbreak{}SharpRestriction.\allowbreak{}GrowthBounds.\allowbreak{}SegmentDegree.\allowbreak{}SelIsComplRestrictSupport}{106--144}
\leanref{Formalism.\allowbreak{}MainNotions.\allowbreak{}Pf14.\allowbreak{}SPtrue}{Formalism.\allowbreak{}SharpRestriction.\allowbreak{}GrowthBounds.\allowbreak{}ScaleEdges.\allowbreak{}SymbolRealisableBound}{43--47}

A column path on segment \(j\) starts at zero weight, with strictly increasing letter positions ending at \(z_j\), and satisfies \(N_s+r\le K\) by \(N_1+3r\le k\). Add a horizontal strip
at an even letter and remove one at an odd letter. At weight
\(\mu\), column \(c\in\Z\) has height \(\#\{i:\mu_i\ge c\}\).
Build the exterior column polynomial \(K_c(g;J)\) over \(O[X]\) using the \(g\)-twisted rows \(v^{(g)}_x=X_xg_x\) (even \(x\)) and \(w^{(g)}_x=X_xg_x^{-\mathsf t}\) (odd \(x\)) as in Lemma~\ref{B5.20he}(1): at the first change use power one, at each subsequent change \(x\) raise the exterior coordinates \(K_I\) to the polynomial power \(p^{x-y}\) (where \(y\) is the preceding change), apply exterior multiplication for rising positive or falling nonpositive columns and contraction for the other changes (writing \(u=v^{(g)}_x\) at even \(x\) and \(u=w^{(g)}_x\) at odd \(x\)),
\[
 \sum_{b\in I}(-1)^{\#\{i\in I:i>b\}}u_bK_{I\setminus b},
 \qquad
 \sum_{b\notin I}(-1)^{\#\{i\in I:i>b\}}u_bK_{I\cup b},
\]
and after the last change at \(y\) raise to \(p^{z_j-y}\). By Lemma~\ref{B5.20he}(1), each \(K_c(g;J)\) is integral of multidegree \(p^{z_j-x}\) at every used letter \(x\), satisfies \(\operatorname{ord}_{g'}(K_c(g;J))\ge\mathrm{oc}(c)\) for every \(g'\in\GL_k(W(\mathbf F_{p^{n'}}))\), and agrees under \(s_g\) modulo \(\omega^{\mathrm{oc}(c)+1}O[X]\) with \(\omega^{\mathrm{oc}(c)}\) times the entrywise segment substitution of the separated column \(K^{\rm sep}_c(J)\), where \(\mathrm{oc}(c)=\sum_{x\text{ costly}}p^{z_j-x}\le2Q_*\). The product of these entrywise-substituted leading exterior coordinates is the segment symbol.
\leanref{Formalism.\allowbreak{}MainNotions.\allowbreak{}G.\allowbreak{}KKstate}{Formalism.\allowbreak{}SharpRestriction.\allowbreak{}GrowthBounds.\allowbreak{}ObjPol}{538--552}
\leanref{Formalism.\allowbreak{}MainNotions.\allowbreak{}G.\allowbreak{}symbolGeneratorData.\allowbreak{}key\_\allowbreak{}column\_\allowbreak{}leading\_\allowbreak{}bound}{Formalism.\allowbreak{}SharpRestriction.\allowbreak{}GrowthBounds.\allowbreak{}PathWeight.\allowbreak{}LetterWeightPos}{283--338}

At each mark \(x\in S\) (and in particular at each non-left-frame letter of a segment), the \(r\) preceding even marks \(y_1>\cdots>y_r\) in reverse order define \(D_x(g)\) (so \(s_g(D_x(g))=D_x\)), the bordered minor \(B_{x,b}(g)\) (at even \(x\), \(r<b\le k\)), and the adjugate-pairing factor \(P_{x,b}(g)\) (at odd \(x\), \(1\le b\le r\)) of Lemma~\ref{B5.20he}(2). By Lemma~\ref{B5.20he}(2), both \(B_{x,b}(g)\) and \(P_{x,b}(g)\) belong to \(O[X]\), satisfy \(\operatorname{ord}_{g'}(B_{x,b}(g))\ge1\) and \(\operatorname{ord}_{g'}(P_{x,b}(g))\ge1\) for every \(g'\in\GL_k(W(\mathbf F_{p^{n'}}))\), and satisfy \(s_g(B_{x,b}(g))\equiv\omega s_g(D_x(g))X_{x,b}\) and \(s_g(P_{x,b}(g))\equiv\omega s_g(D_x(g))X_{x,b}\pmod{\omega^2O[X]}\), with every error term having normal degree equal to its \(\omega\)-order (at least \(p\)). The reserve \(D_x^{kQ_*}\) pays for the normal entries of the residual monomial, each lifted by its corresponding factor \(B_{x,b}(g)\) or \(P_{x,b}(g)\).
\leanref{Formalism.\allowbreak{}MainNotions.\allowbreak{}G.\allowbreak{}WittFrameLocaliserOrders.\allowbreak{}firstKeyMinor}{Formalism.\allowbreak{}SharpRestriction.\allowbreak{}GrowthBounds.\allowbreak{}WittFrameLocaliserOrders}{3153--3158}
\leanref{Formalism.\allowbreak{}MainNotions.\allowbreak{}G.\allowbreak{}WittFrameLocaliserOrders.\allowbreak{}pairingKey}{Formalism.\allowbreak{}SharpRestriction.\allowbreak{}GrowthBounds.\allowbreak{}WittFrameLocaliserOrders}{3159--3162}

Use \(L^{c'}\) of \ref{B5.20d} in place of \(R_0^{c'}\), with unshifted frame powers. Define a term \(S\) at level \(h\) to be a product of an explicit power \(\omega^m\), leading exterior column factors of orders \(a_c=\mathrm{oc}(c)\), \(q\) leading bordered or pairing factors \(\omega s_g(D_x(g))X_{x,b}\), integral nonnormal polynomials (including the reserve factors of \(D_x\) and \(\Delta_j\)), and \(R_0^{c'}\), with \(h=m+\sum_c a_c+q\). Its raw lift \(F^0\) replaces each leading column by \(K_c(g;J)\), each leading bordered or pairing factor by \(B_{x,b}(g)\) or \(P_{x,b}(g)\), each nonnormal variable by its \(g\)-twisted coordinate (which has order \(\ge0\) under every \(s_{g'}\)), and \(R_0^{c'}\) by \(L^{c'}\) (which has order \(\ge c'\) on every competing plane by Lemma~\ref{B5.20d}). By Lemma~\ref{B5.20he}(4), the raw lift \(F^0\) satisfies \(\operatorname{ord}_{g'}(F^0)\ge m+\sum_c a_c+q+c'=h+c'\) on every competing plane \(g'\). For a normal component of degree \(b\), the initial term has \(m=0\) and \(h=\sum_c a_c+q=b\), where \(q\) counts the normal entries of the residual monomial. Keep with each term at level \(b+d\) at least \((2Q_*+3)c'-(2Q_*+1)d\) copies of \(D_x\) at each row and \(c'-d\) copies of each \(\Delta_j\), after accounting for the initial residual monomial. Expand \(s_g(F^0)-S\). Terms of order at least \(b+c'\) already meet the precision bound. Every other term has an additional rise \(d'\ge1\). For an exterior column \(K_c(g;J)\) of order \(a=\mathrm{oc}(c)\le2Q_*\), every monomial of normal degree \(t\) in \(s_g(K_c(g;J))\) has coefficient divisible by \(\omega^t\) because \(K_c(g;J)\in O[X]\), and after subtracting its leading part its coefficients are also divisible by \(\omega^{a+1}\). Write such a monomial as \(\omega^z X^\nu H\) with \(H\) integral in the nonnormal coordinates, \(|\nu|=t\), and \(z=\max(t,a+1)\), absorbing the remaining coefficient into \(H\). Its rise is \(e=z-a\ge1\), and \(t\le a+e\le(2Q_*+1)e\). For a bordered or pairing factor \(B_{x,b}(g)\) or \(P_{x,b}(g)\) of order \(1\), Lemma~\ref{B5.20he}(2) gives rise \(e=p^{x-y_a}-1\ge p-1\) and normal degree \(t=e+1\le2e\le(2Q_*+1)e\). For a localizer error of rise \(e\), choosing replacement minors in the product of \ref{B5.20d} loses one \(\Delta_j\) for each replaced copy at frame \(j\); each replacement minor contributes at least one to the order after substitution and each normal entry contributes one power of \(\omega\), so both the deficit of \(\Delta_j\) and the normal degree are at most \(e\).

Expand the product of the columns, bordered and pairing factors, and \(c'\) localizers, subtracting the product of their leading terms. Leave unchanged column, bordered, and pairing factors in their leading form. In each remaining summand the rises add. At total rise \(e\) there are at most \((2Q_*+1)e\) exposed normal variables (with explicit \(\omega\)-power at least the number of exposed variables) and at most \(e\) missing factors of each \(\Delta_j\). Pair each exposed \(\omega X_{x,b}\) with one reserve copy of \(D_x\) to form a leading bordered or pairing factor \(\omega s_g(D_x(g))X_{x,b}\), keeping any remaining explicit \(\omega\)-power. Set aside \((2Q_*+1)e\) copies of each \(D_x\) and \(e\) copies of each \(\Delta_j\) (the latter restoring \(R_0^{c'}\)); unused copies multiply the nonnormal coefficient. The resulting child terms have level \(h+e\), preserve every row multidegree, and retain the two prescribed reserves for level \(h+e\).

Put \(b_*=b+c'\) for this recursion. If \(h\ge b_*\), the zero lift gives the required agreement modulo \(\omega^{b_*}\). Otherwise let \(F^0\) be the raw lift of a term \(S\) at level \(h\). The expansion just obtained gives \(s_gF^0-S=\sum_\nu S_\nu+R\), with \(R\in\omega^{b_*}O[X]\) and \(h< h_\nu<b_*\). By descending induction in \(h\), choose lifts \(F_\nu\) of the children to precision \(b_*\) and take \(F=F^0-\sum_\nu F_\nu\). Then \(s_gF-S=R-\sum_\nu(s_gF_\nu-S_\nu)\) has the required order. By Lemma~\ref{B5.20he}(4), \(F^0\) has order at least \(h+c'\ge b_*\) on every competing plane, and every subtracted lift \(F_\nu\) has order at least \(b_*\) there by induction. All operations take place over \(O\).

This also treats a target symbol \(P_b=(ZU+p^{c'}Y')_b\) with
\(Z\) in the symbol space and \(Y'\) integral. The term
\(\omega^bp^{c'}Y'_b\) already belongs to \(\omega^{b+c'}O[X]\).
Projecting to the row multidegree commutes with the substitution.
The integral lift \(F_b\) therefore satisfies
\[
 s_gF_b-\omega^bP_b\in\omega^{b+c'}O[X],\qquad
 s_{g'}F_b\in\omega^{b+c'}O[X]\quad(g'\text{ on a competing edge}).
 \tag{5.20.16d}\label{eq:p5c2-clean-symbol}
\]
Sum \(\omega^{-b}F_b\) for \(b\le\lfloor T\rfloor\).
\leanref{Formalism.\allowbreak{}MainNotions.\allowbreak{}G.\allowbreak{}WittFrameLocaliserOrders.\allowbreak{}clean\_\allowbreak{}symbol\_\allowbreak{}order}{Formalism.\allowbreak{}SharpRestriction.\allowbreak{}GrowthBounds.\allowbreak{}WittFrameLocaliserOrders}{1247--1300}
\leanref{Formalism.\allowbreak{}MainNotions.\allowbreak{}G.\allowbreak{}symbolGeneratorData.\allowbreak{}core}{Formalism.\allowbreak{}SharpRestriction.\allowbreak{}GrowthBounds.\allowbreak{}PathWeight.\allowbreak{}FixedSymbolFactorIntegral}{148--207}
\leanref{Formalism.\allowbreak{}MainNotions.\allowbreak{}Pf14.\allowbreak{}exists\_\allowbreak{}pol\_\allowbreak{}family\_\allowbreak{}of\_\allowbreak{}mem\_\allowbreak{}gen\_\allowbreak{}of\_\allowbreak{}vsDeg\_\allowbreak{}le}{Formalism.\allowbreak{}SharpRestriction.\allowbreak{}GrowthBounds.\allowbreak{}ScaleEdges.\allowbreak{}ExistsPolFamilyOfMemGenOfVsDegLe}{140--252}
\leanref{Formalism.\allowbreak{}MainNotions.\allowbreak{}Pf14.\allowbreak{}flatSrc}{Formalism.\allowbreak{}SharpRestriction.\allowbreak{}GrowthBounds.\allowbreak{}ScaleEdges.\allowbreak{}FlatSrc}{60--111}
\leanref{Formalism.\allowbreak{}MainNotions.\allowbreak{}Pf14.\allowbreak{}edgeFamOf}{Formalism.\allowbreak{}SharpRestriction.\allowbreak{}GrowthBounds.\allowbreak{}ScaleEdges.\allowbreak{}ScaleZpowMono}{48--118}

Multiply the lifts by the following \(B\) residual tableaux. Number rows and columns from one, and let \(\nu\) be the partition of the residual active weight after subtracting its last entry. At a vector mark take the standard tableaux whose leading monomial has exponent \(\nu_1\) at matrix entry \((1,1)\). At a covector mark require that the entire first matrix row have exponent zero. Off the marks take the full tableau basis. Thus the dimension factor at a mark is the Weyl dimension in rank \(k-1\) of \((\lambda_{s,2},\ldots,\lambda_{s,k})\).
The additional order and depth per row are bounded by the sum of
its first \(r\) residual entries minus \(r\) times its last entry,
at most \(rc_t\) at marks and \(2rc_b\) elsewhere. Reduce a proposed relation in the adapted coordinates. For each residual tableau tuple \(i\), let \(Q_i\) be its coefficient polynomial in the flat letters, including the common factor \(U\). These polynomials have the prescribed flat degree in each factor. The relation in the class variables substitutes the first-column entries for the flat letters at a vector mark, and the signed initial minors omitting one row for those letters at a covector mark. Expand each \(Q_i\) in monomials. We recover both the flat monomial and the residual tableau from the leading exponent of their product, using the column-then-row order of \ref{B3.52b}.

At a vector mark the residual first column has total degree \(\nu_1\), all at entry \((1,1)\). Subtract this known column from the first column of the product to recover the exponents of all the flat letters. At a covector mark write \(a_c\) for the exponent of the minor omitting row \(c\). The leading term of that minor uses the diagonal entry \((t,t)\) precisely when \(t<c\), for \(1\le t<k\). The residual tableau contributes zero on this diagonal. Indeed its entries omit the letter one, and strict increase down each tableau column makes the entry at depth \(t\) at least \(t+1\). The diagonal exponents of the product therefore give the tails \(\sum_{c>t}a_c\). Adjoining the prescribed total degree and the empty tail recovers every \(a_c\) by subtraction. Off the marks the flat degree is zero. At every factor, subtracting the recovered flat exponent leaves the leading exponent of the residual tableau, which determines its rows by \ref{B3.52b}. Different factors have disjoint matrix variables. All these leading coefficients are units, including the cofactor signs. The largest leading monomial of a nonzero term in the expanded relation consequently cannot cancel. Hence every \(Q_i\) vanishes over the residue field.
\leanref{Formalism.\allowbreak{}MainNotions.\allowbreak{}TestSideCounting.\allowbreak{}key\_\allowbreak{}inj}{Formalism.\allowbreak{}SharpRestriction.\allowbreak{}GrowthBounds.\allowbreak{}DualFockLinear.\allowbreak{}DgFalseDiag}{168--237}
\leanref{Formalism.\allowbreak{}MainNotions.\allowbreak{}TestSideCounting.\allowbreak{}bigInd}{Formalism.\allowbreak{}SharpRestriction.\allowbreak{}GrowthBounds.\allowbreak{}DualFockLinear.\allowbreak{}ThAllMonomial}{154--158}
Cancel the common
nonzero reduction of \(U\) and use the unit coefficient minor
of the symbols. All coefficients in the original relation
vanish modulo \(\omega\). This proves primitivity of a combination
with a unit coefficient and the claimed source rank.
\leanref{Formalism.\allowbreak{}MainNotions.\allowbreak{}TestSideCounting.\allowbreak{}primStep\_\allowbreak{}of}{Formalism.\allowbreak{}SharpRestriction.\allowbreak{}GrowthBounds.\allowbreak{}DualFockLinear.\allowbreak{}TheoremIIClosedOfTestBound}{13--135}

For the target side write \(\lambda=\phi+\eta\), where \(\phi\)
is the split-off flat weight, zero off the marks. Let \(\Psi_\xi\)
identify functionals with vectors by the inverse normalized
factorial form at weight \(\xi\) as in \ref{B3.10a}. Write
\[
 k_\xi=\binom{k}{2}\sum_{s:l_s>0}\lfloor\log_p(l_s+k-1)\rfloor,
 \qquad l_s=\xi^{\rm act}_{s,1}-\xi^{\rm act}_{s,k}
\]
on the marked factors. Transposing the divided raising operators
shows that the image is the divided-power lattice. By
\ref{B3.52b},
\(p^{k_\xi}R_{e^\vee}^{\xi}\subseteq\Psi_\xi((R_e^\xi)')
\subseteq R_{e^\vee}^{\xi}\) for the reversed dual edge \(e^\vee\).
Put \(B_0=Nk^2\log_p a_0\). Since \(l_s+k\le a_0^2\),
\(\kappa=k_\lambda+k_\phi\le4B_0\).
\leanref{Formalism.\allowbreak{}InvariantFactorPaths.\allowbreak{}INST.\allowbreak{}Top.\allowbreak{}sandwichLocal}{Formalism.\allowbreak{}ThetaLogCoordinates.\allowbreak{}SandwichModelDomains.\allowbreak{}FinalW}{50--52}
\leanref{Formalism.\allowbreak{}L2.\allowbreak{}G.\allowbreak{}Model.\allowbreak{}testSide\_\allowbreak{}of\_\allowbreak{}core}{Formalism.\allowbreak{}LocalToGlobal.\allowbreak{}GroupLattices.\allowbreak{}FockModels.\allowbreak{}TestSide}{36--99}
\leanref{Formalism.\allowbreak{}MainNotions.\allowbreak{}TestSideCounting.\allowbreak{}sizeStep}{Formalism.\allowbreak{}SharpRestriction.\allowbreak{}GrowthBounds.\allowbreak{}DualFockLinear.\allowbreak{}LogSizeBoundLe}{94--167}

For \(\beta\in R_{e^\vee}^\eta\), multiplication \(m_\beta\)
defines \(L_\beta=p^{k_\lambda}\Psi_\lambda^{-1}m_\beta\Psi_\phi\).
Its transpose for the normalized edge pairings is \(\rho_\beta\),
given by \(T(\rho_\beta x)=p^{A_\eta}(L_\beta T)(x)\).
Choose an integral frame of the lattice of the reversed edge at each translate, as in \ref{B3.52a}. For a dominant active weight \(\xi\), remove the determinant twist \(\xi_k\). In the remaining tableau the number of columns of height \(h<k\) is \(\xi_h-\xi_{h+1}\). Both this number and the removed twist are additive in \(\lambda=\phi+\eta\), also after negated reversal. In a tableau product of minors for \(\lambda\), assign \(\phi_h-\phi_{h+1}\) columns of each height to the first factor and the rest to the second, using the active weights at that translate. Each subtableau retains its column order. Its product of minors is an integral section of the assigned weight after restoring its determinant twist. By the tableau basis of \ref{B3.52b} this factors a spanning set. Taking products over translates gives the assertion for the tensor lattice. The edge factors are \(\omega^{-A_\xi}\) in \ref{B3.52a}; the sums of the first \(r\) active entries satisfy \(A_\lambda=A_\phi+A_\eta\), so this factorization gives the span on \(R_{e^\vee}^\lambda\) itself.

For each flat factor \(y\) just used, \(T=\Psi_\phi^{-1}(p^{k_\phi}y)\) belongs to the flat test lattice by the preceding inclusions. Then \(L_\beta T=p^\kappa\Psi_\lambda^{-1}(y\beta)\). As \(\Psi_\lambda(R'_e)\subseteq R_{e^\vee}^\lambda\), these tests span a lattice containing \(p^\kappa R'_e\). If every \(\rho_\beta x\) belonged to \(\omega^{2\kappa+1}R_e^\phi\), its normalized edge pairing with every integral flat test would belong to \(\omega^{2\kappa+1}O\). By the transpose identity and additivity of \(A_\xi\), the same would hold for the normalized pairing of \(x\) with \(p^\kappa R'_e\). Since \(p^\kappa=\omega^{2\kappa}\), perfection would imply \(x\in\omega R_e^\lambda\). For primitive \(x\) some integral component therefore has exact order \(a\le2\kappa\).
\leanref{Formalism.\allowbreak{}MainNotions.\allowbreak{}TestSideCounting.\allowbreak{}ontoAll}{Formalism.\allowbreak{}SharpRestriction.\allowbreak{}GrowthBounds.\allowbreak{}DualFockLinear.\allowbreak{}RestCloudBOfTwo}{23--84}
\leanref{Formalism.\allowbreak{}MainNotions.\allowbreak{}TestSideCounting.\allowbreak{}rhoAtG\_\allowbreak{}of\_\allowbreak{}pair}{Formalism.\allowbreak{}SharpRestriction.\allowbreak{}GrowthBounds.\allowbreak{}DualFockLinear.\allowbreak{}RhoAtGOfPair}{74--221}
\leanref{Formalism.\allowbreak{}MainNotions.\allowbreak{}TestSideCounting.\allowbreak{}rhoStep\_\allowbreak{}of\_\allowbreak{}top}{Formalism.\allowbreak{}SharpRestriction.\allowbreak{}GrowthBounds.\allowbreak{}DualFockLinear.\allowbreak{}RestCloudBOfTop}{115--133}

The flat polynomial \(G\) of \(y\) in edge coordinates has integral
coefficients and a unit coefficient. The residual partition
contributes words in \(X_{s,a}\partial_{s,b}\) of length at most
\(c_t\) per mark. This follows by differentiating the products
of its at most \(c_t\) columns. A row substitution conjugates
each first-order operator to a linear combination of such
operators. The normal-degree bound becomes \(\lfloor T\rfloor+Nc_t\).
\leanref{Formalism.\allowbreak{}MainNotions.\allowbreak{}TestSideCounting.\allowbreak{}restND\_\allowbreak{}of\_\allowbreak{}adm}{Formalism.\allowbreak{}SharpRestriction.\allowbreak{}GrowthBounds.\allowbreak{}DualFockLinear.\allowbreak{}RestNDOfAdm}{14--81}

For a test polynomial \(Z\), put \(v_d=v(\prod_s d_s!)\) and
\(k_d=\sum_{s:d_s>0} k\lfloor\log_p d_s\rfloor\). The integral
normalized pairing is
\[
 \frac{p^{k_d}}{\prod_s d_s!}
 \sum_\alpha\alpha!\,[X^\alpha]G\,[X^\alpha]Z .
 \tag{5.20.16e}\label{eq:p5c2-flat-pairing}
\]
The valuation of a multinomial coefficient in row \(s\) is at
most \(k\lfloor\log_p d_s\rfloor\), by successive binomial choices
and \ref{B3.19}. For one digit split \(\alpha=u+Pb\), where
\(P\) is a power of \(p\) and \(0\le u<P\), the unit
\(\gamma(u,b)=(u+Pb)!/(u!\,b!\,(P!)^b)\) has valuation zero
by the same formula for factorial valuations. Separate the
plane and normal variables. For a polynomial \(H\) let
\(v_{\rm F}(H)=\min_\alpha(v(H_\alpha)+v(\alpha!)/2)\).
Decompose the normal part by its prefix isotypic projections.
One component has valuation no greater than the original
minimum because these projections sum to the identity. At a
normal index attaining this minimum select its plane coefficient
\(P_c\). If \(m\) minimizes \(v_{\rm F}(P_c)\), divide by
\([X^m]P_c\,\omega^{v(m!)}\) to normalize it to valuation zero.
\leanref{Formalism.\allowbreak{}MainNotions.\allowbreak{}Proposition5Numerics.\allowbreak{}pick\_\allowbreak{}core}{Formalism.\allowbreak{}SharpRestriction.\allowbreak{}GrowthBounds.\allowbreak{}DiagonalSupport.\allowbreak{}PickCore}{23--166}

\paragraph{Projection with denominators.}
On the multihomogeneous space of a signed alphabet, write
\(d_{\rm tot}=\sum_i d_i\) for the sum of the divided plane degrees.
For a prefix \(S_0\), the condition on the least signed weight
for a conjugate of \(\operatorname{diag}(T,1,\ldots,1)\) selects
the isotypic components with last highest-weight entry at least
the prescribed threshold. The projection onto this space loses
at most \(e_C=r^2(1+\log_p(2d_{\rm tot}+r))\) in \(v_{\rm F}\).

To see the denominator bound, choose a lattice basis \(b_i\) of
the irreducible representation \(L_\lambda\). The invariant
projector \(P_{\rm inv}\) on \(\operatorname{Hom}(L_\lambda,V)\)
gives
\[
 e_\lambda(v)=\dim L_\lambda\sum_i
 P_{\rm inv}(b_i^\vee\otimes v)(b_i).
\]
On an equivariant copy of \(L_\lambda\) the invariant projection
is normalized trace on endomorphisms. On a different isotypic
summand the Hom space has zero invariants. This proves the
formula. Use the lattice for \(v_{\rm F}\), whose monomial
generators have coefficient scale \(\omega^{-v(\alpha!)}\).
It is preserved by divided root operators. For example, moving
\(h\) powers between two coordinates has squared normalized
coefficient \(\binom{a}{h}\binom{b+h}{h}\) on exponents \(a,b\).
For the Hom lattice take its intersection on the irreducible
factor and the corresponding dual lattice. The invariant
projection denominator is bounded by \(\binom r2\) times the
largest \(\lfloor\log_p M\rfloor\) for the root strings, with
\(M\le2d_{\rm tot}+r\). In the extremal-projector expansion each
positive root contributes a divided raising and lowering
product and the inverse of a binomial coefficient with upper
argument at most \(M\). Its valuation cost is at most
\(\lfloor\log_p M\rfloor\) by \ref{B3.19}. Evaluation against
the lattice basis and multiplication by \(\dim L_\lambda\)
introduce no further denominator. The sum of the selected
\(e_\lambda\) has the same bound by the ultrametric inequality.

The factorial form is contravariant for transposition of the
signed operators, so the isotypic projections are self-adjoint.
For disjoint sets of letters the actions commute. For nested
sets, the small-set projector commutes with the small-set
action and with the action on the complement in the large
set. Hence it preserves the large-set isotypic spaces and
commutes with their projectors. Each product of \(h\) such
projections therefore loses at most \(he_C\) and preserves
the pairing against polynomials already in every selected
subspace.
\leanref{Formalism.\allowbreak{}MainNotions.\allowbreak{}G.\allowbreak{}GLWeightTheory.\allowbreak{}isotypic\_\allowbreak{}projector\_\allowbreak{}trace\_\allowbreak{}formula}{Formalism.\allowbreak{}SharpRestriction.\allowbreak{}GrowthBounds.\allowbreak{}RootWeights.\allowbreak{}IsotypicProjectorTraceFormula}{21--101}
\leanref{Formalism.\allowbreak{}MainNotions.\allowbreak{}G.\allowbreak{}GLWeightTheory.\allowbreak{}invariant\_\allowbreak{}denominator\_\allowbreak{}sharp\_\allowbreak{}of}{Formalism.\allowbreak{}SharpRestriction.\allowbreak{}GrowthBounds.\allowbreak{}RootWeights.\allowbreak{}InvariantProjectionExtremalFormula}{92--103}
\leanref{Formalism.\allowbreak{}MainNotions.\allowbreak{}PlaneGram.\allowbreak{}PROJ}{Formalism.\allowbreak{}SharpRestriction.\allowbreak{}GrowthBounds.\allowbreak{}MonomialOrderSegments.\allowbreak{}LiftBoundOfHypotheses}{168--179}
\leanref{Formalism.\allowbreak{}MainNotions.\allowbreak{}PlaneGram.\allowbreak{}prodProj}{Formalism.\allowbreak{}SharpRestriction.\allowbreak{}GrowthBounds.\allowbreak{}MonomialOrderSegments.\allowbreak{}ProjKeepsLetter}{56--113}

Number the right-frame slots from \(1\) to \(r\), and let \(I_j\) be the interior letters of segment \(j\), excluding both frames. The link indexed by \((j,m)\), \(1\le m<r\), is \(I_j\cup\{(j,h):m<h\le r\}\). The letter \((j,h)\) here denotes the shared boundary letter in the original plane alphabet. The links are nested for fixed \(j\) and disjoint for different \(j\). In the enlarged alphabet the corresponding set consists of the powered letters of the preceding frame, the ordinary letters of \(I_j\), and the surplus letters of the indicated right-frame suffix. A chain prefix consists of all the preceding powered frame letters together with the indicated initial part of \(I_j\). Put
\(D_l=J_f(r-1)e_C\), \(n_0=v(\prod_i d_i!)/2\), and \(k'=\sum_i r(1+\log_p(d_i+r))\). On the original plane alphabet fix the low digits \(d_{xa}\) of the selected class. In row \(x\), let \(M_x=p^{z_j-x}\) at its non-left-frame occurrence. Its factorial unit on plane exponent \(\alpha\) is \(u(\alpha)=\prod_{x,a}\gamma(d_{xa},\alpha_{xa})\), with power \(M_x\) in each factor. For an exponent \(m\) minimizing the normalized \(P_c\), put \(t=v(m!)/2\), and let \(Y\) be the link projection of \(X^m\). The preceding projector bound gives \(v_{\rm F}(Y)\ge t-D_l\). Write \(\rho_\beta x=\omega^a y\) for the component of exact order selected above, using the exact symbol representatives for the source combination. Apply Lemma~\ref{B5.20hb} to the flat polynomial of \(y\), with the selected normal projector and normal index. It gives \(\langle P_c,Y\rangle=[X^m]P_c\,m!\) of valuation \(t\), and gives the chain bound on every boundary digit component of this same \(P_c\). The transfer to the lifted source representatives uses the precision of \eqref{eq:p5c2-clean-symbol} after the test has been constructed.

Take an integral binomial approximation \(g\) to \(u^{-1}\) modulo \(p^{c'}\) from \ref{B5.20ha}. Write \(D_g\) for its diagonal action. For all polynomials \(A,B\), the factorial pairing satisfies \(v(\langle A,B\rangle)\ge v_{\rm F}(A)+v_{\rm F}(B)\), term by term. Thus \(\langle P_c,(D_uD_g-I)Y\rangle\) has valuation at least \(c'+t-D_l>t\). Put \(a=\lceil n_0+D_l-t\rceil\). Since every monomial factorial in these row degrees divides \(\prod_i d_i!\), the coefficients of \(p^aD_gY\) are integral. Its pairing against \(P_c\) with multiplier \(u\) has valuation \(a+t\le n_0+D_l+1\). A binomial of degree \(d\) in a coordinate exponent acts by \(X^d\partial_X^{[d]}\). In any basis each summand of its conjugate lowers the least signed first-coordinate weight by at most \(d\). Thus every link retains its threshold minus \(6c'\), using \(4c'\le6c'\) in \ref{B5.20ha}.

Let \(T_0=p^{\lceil n_0+D_l-t\rceil}D_gY\). It is integral and has every original link bound with loss \(6c'\), as just proved. Lemma~\ref{B5.20hc} gives an integral allowable preimage \(Z\), with the same bounds on the enlarged links. Denote recombination by \(\Phi\). Thus \(\Lambda(Z)=\langle P_c,D_u\Phi Z\rangle=\langle P_c,D_uT_0\rangle\) has valuation at most \(n_0+D_l+1\). For an exponent array \(\zeta\) on this alphabet, let \(\bar\zeta\) be its array on the original letters. It keeps the interior exponents and assigns exponent \(\rho+P_t\beta\) to a boundary variable with remainder exponent \(\rho\) and quotient exponent \(\beta\). The separate right-frame surplus contributes zero to \(\bar\zeta\). At a boundary \(t\) denote the remainder and quotient
rows of \(\zeta\) by \(\rho_t,\beta_t\). The diagonal unit is
\[
 \Gamma(\zeta)=u(\bar\zeta)
 \prod_t\prod_{a=1}^r\gamma(\rho_{t,a},\beta_{t,a}),
\]
with \(u\) as above and with power \(P_t\) in the second product. On a fixed remainder class apply \ref{B5.20ha} once with the original row powers and digits, and once with the boundary powers and remainders. Both applications hold for either sign. Substitution of \(\rho+P_t\beta\) into the first approximation preserves total degree, and the resulting polynomial is integer-valued, so its binomial coefficients in the free exponents are integral by forward differences. The class selection acts on the remainder letters alone. These letters are outside every chain prefix and enlarged link, so the selection commutes with their group actions.
The remaining factorial constants depend on row degrees and
remainder rows. On the allowed low remainder exponents the
binomial expansions of these units and their inverses give
integral multipliers \(g_+,g_-\) with
\(v(g_+-\Gamma),v(g_--\Gamma^{-1})\ge c'\).
A binomial in a quotient exponent is a divided diagonal
derivative. There are two factors for each sign. The original factorial unit evaluated on the collapsed exponents and the product of quotient--remainder units each have a binomial approximation of total degree at most \(6c'\). Their product costs \(12c'\). The error still has valuation at least \(c'\), since both exact factors are units and both approximations are integral. A binomial in an affine exponent \(u+Pv\) of degree \(d\) has an integral binomial expansion in \(v\) of degree at most \(d\), by finite differences. On a fixed low-digit class each of the two diagonal multipliers therefore lowers a prefix order by at most \(12c'\). In the enlarged alphabet the resulting chain bound is its original threshold minus its assigned slack and \(12c'\), and the link bound is its original threshold minus \(30c'\). For the quotient letters from the preceding frame the slack used before the projection is at least \(r d_I/P+13rc'\) for each boundary power \(P\), where \(d_I\) is the total plane degree of that preceding interior. Summing the \(r\) quotient loads costs at most the assigned slack minus \(12c'\), since \(r\ge1\).

On the enlarged alphabet, within each segment a prefix contains
its left-frame powered letters and an initial part of its
interior. These sets are nested there and disjoint across
segments. They are also contained in, or disjoint from, the
enlarged link sets of their segment. The product \(P_{\rm pre}\)
of their projections preserves the link conditions and commutes
with diagonal derivatives on the remainder letters outside
these prefixes. Write \(D_m=n_c e_C\), where \(n_c\) counts the
prefixes used. The operator is
\[
 \operatorname{diag}(g_-)\,P_{\rm pre}\,
 \operatorname{diag}(g_+).
\]
The remainder degrees are preserved because the projections
commute with those coordinate operators. Each prefix action
preserves every letter degree, including the surplus right-frame
degree. Hence the output still obeys the digit and row constraints.
The projection imposes the chain thresholds, and the derivative
slack preserves the required link thresholds after the last
diagonal multiplication. Multiplying by \(p^{\lceil D_m+k'/2\rceil}\)
clears the remaining denominator in the plane test. To compare its value, take \(\theta_{ji}=H_{ji}-s_i/8\) and \(\sigma_{ji}=3s_i/8\) in Lemma~\ref{B5.20hd}. We check the slack at a cut \(x_i\) of segment \(j\). Put \(D=z_{j-1}-f_{j-1,0}\) and \(P=p^{z_j-z_{j-1}}\). Each interior position of segment \(j-1\) is at distance at least \(D+L_1+1\) from its endpoint and has divided degree at most \(2a_0\) times the inverse power at that distance. Summing distinct distances gives \(d_{I_{j-1}}\le3a_0p^{-D-L_1-1}\). The oldest even predecessor in segment \(j\) is at or after \(f_{j-1,0}\), so \(a_0p^{-D}/P\le b_{x_i}\). The spacing bound and \(Y\ge r\) imply \(48rYk^2\le p^{L_1+1}\), whence \(r d_{I_{j-1}}/P\le s_i/8\), using \(b_{x_i}/(Yk^2)\le2s_i\). The last size bound and \(b_{x_i}\ge a_0/p^{G'}\) give \(s_i\ge100r^2c'\) on enlarging the absolute constant. This proves the required slack. The row digit approximations above give the two unit congruences, and Lemma~\ref{B5.20hb} gives the source chain bounds. Write \(\Omega=D_{g_-}P_{\rm pre}D_{g_+}\). Lemma~\ref{B5.20hd} now gives \(v(\Lambda(\Omega Z)-\Lambda(Z))\ge c'+n_0-2k'-D_m>n_0+D_l+1\), since \(k'\ge1\) and the size cutoff gives \(c'>4(D_m+D_l+3k')\). Thus the two values have equal valuation. With \(N_2=\lceil D_m+k'/2\rceil\), the integral test \(p^{N_2}\Omega Z\) has value of valuation at most \(n_0+D_l+1+N_2\le n_0+D_l+D_m+3k'\). Here the denominator bound can also be read directly on each block: the label and surplus factorials are fixed, and the variation of the sum of the free letter factorial valuations is at most \(k'\) by the row multinomial bound. Hence a loss \(D_m\) in \(v_{\rm F}\) costs at most \(D_m+k'/2\) in coefficient valuation. This is the plane contribution used in \eqref{eq:p5c2-plane-cost}. Straightening
in the signed column basis expresses the projected component as
products with those prefix bounds. Recombining letters and
multiplying by the chosen normal monomial therefore puts each
normal homogeneous component in the dual-edge symbol space
with its fixed factor, up to the \(p^{c'}\) remainder allowed
in \eqref{eq:p5c2-clean-symbol}. The normal degree stays at most
\(\lfloor T\rfloor+Nc_t\).
\leanref{Formalism.\allowbreak{}MainNotions.\allowbreak{}PlaneGram.\allowbreak{}TwistData.\allowbreak{}gaP}{Formalism.\allowbreak{}SharpRestriction.\allowbreak{}GrowthBounds.\allowbreak{}MonomialOrderSegments.\allowbreak{}CollapseExponent}{57--59}
\leanref{Formalism.\allowbreak{}MainNotions.\allowbreak{}PlaneGram.\allowbreak{}stepDEF\_\allowbreak{}of}{Formalism.\allowbreak{}SharpRestriction.\allowbreak{}GrowthBounds.\allowbreak{}MonomialOrderSegments.\allowbreak{}OmegaValue}{108--120}
\leanref{Formalism.\allowbreak{}MainNotions.\allowbreak{}PlaneGram.\allowbreak{}plane\_\allowbreak{}gram\_\allowbreak{}theorem\_\allowbreak{}of\_\allowbreak{}five\_\allowbreak{}lemmas}{Formalism.\allowbreak{}SharpRestriction.\allowbreak{}GrowthBounds.\allowbreak{}MonomialOrderSegments.\allowbreak{}PlaneGramTheoremOfLemmas}{13--17}

Write \(b=\log_p a_0\). There are at most \(N\) original plane letters, each of degree at most \(2a_0\). Both the number of chain cuts and the number of links are at most \(N\): a segment with \(l\) interior letters contributes \(l+1\) cuts and \(r-1\) links, and the interiors together with one frame per segment use each mark once. Since \(p\le a_0\), \(4N\le a_0\) and \(r\le k\le a_0\), the logarithmic cost \(1+\log_p(2\sum d_i+r)\) is at most \(4b\). This gives \(D_m,D_l,k'\le4B_0\); summing with coefficients \(1,1,3\) gives \(20B_0\). For the normal selection, the joint normal projection costs at most \(2N k^2(3b)\): there are at most \(2N\) cuts and each root-string upper argument is at most \(a_0^3\). The two further costs are half the valuations of the plane and normal row multinomials. They are at most \(Nr(3b)/2\) and \(NK(3b)/2\), by \ref{B3.19}. Their sum with the projection cost is at most \(9B_0\). More explicitly, if \(b_p,b_n\) are the two quotient exponent arrays of the selected unit monomial, the fixed digit factor times \(b_p!b_n!\) has the valuation of the original monomial factorial, at most \(v_d\). Normalizing the minimizing plane coefficient subtracts half the plane row-factorial valuation \(n_0\); the excess is bounded by the two multinomial costs just computed. Hence
\[
 D_m+D_l+3k'\le20B_0,\qquad k_d\le20B_0,\qquad
 v(\text{normal selection scalar})\le v_d-n_0+9B_0 .
 \tag{5.20.16f}\label{eq:p5c2-plane-cost}
\]
The last scalar restores the factorials of the normal variables
and the normalization of \(P_c\). The main term before the common
factorial normalization has valuation at most \(v_d+29B_0\).
The multiplier errors and the discarded plane remainder have
valuation at least \(v_d+c'-20B_0-k_d\). Choose \(C\) so that
\(c'>4(D_m+D_l+3k')\) and \(c'>(9+3\cdot20)B_0\).
The main term determines the valuation in
\eqref{eq:p5c2-flat-pairing}, which is at most \(49B_0\). Lift each normal
component by \eqref{eq:p5c2-clean-symbol} and divide by its
\(\omega\)-degree. The multinomial bound puts the resulting
pairing error in \(\omega^{c'}O\). We also have
\(c'>2\lceil49B_0\rceil\), so the lifted flat test has valuation
at most \(D_F=\lceil49B_0\rceil\), depth at most \(T+Nc_t\), and
order \(c'\) on competing incoming edges.
\leanref{Formalism.\allowbreak{}MainNotions.\allowbreak{}Proposition5Numerics.\allowbreak{}value\_\allowbreak{}core}{Formalism.\allowbreak{}SharpRestriction.\allowbreak{}GrowthBounds.\allowbreak{}DiagonalSupport.\allowbreak{}Proposition5OfThreeParts}{61--124}
\leanref{Formalism.\allowbreak{}MainNotions.\allowbreak{}Proposition5Numerics.\allowbreak{}sizeStep'}{Formalism.\allowbreak{}SharpRestriction.\allowbreak{}GrowthBounds.\allowbreak{}DiagonalSupport.\allowbreak{}SizeStepAlt}{11--145}
\leanref{Formalism.\allowbreak{}MainNotions.\allowbreak{}Proposition5Numerics.\allowbreak{}numStep\_\allowbreak{}of}{Formalism.\allowbreak{}SharpRestriction.\allowbreak{}GrowthBounds.\allowbreak{}DiagonalSupport.\allowbreak{}LeOfOneLeMul}{182--231}
\leanref{Formalism.\allowbreak{}MainNotions.\allowbreak{}TestSideCounting.\allowbreak{}restFlat\_\allowbreak{}of}{Formalism.\allowbreak{}SharpRestriction.\allowbreak{}GrowthBounds.\allowbreak{}DualFockLinear.\allowbreak{}FockSumRight}{47--140}

Pull back by \(L_\beta\). The factor \(\omega^a\) leaves the valuation at most
\(\kappa+D_F\le54B_0\). The residual factors
cost at most \(o\) in depth and order, so \(c'-o=c_*\).
A source approximation error in \(\omega^{c_*}O\) preserves the
value bound since \(c_*>2D_e\).
\leanref{Formalism.\allowbreak{}MainNotions.\allowbreak{}TestSideCounting.\allowbreak{}cloudBAt\_\allowbreak{}of\_\allowbreak{}rho}{Formalism.\allowbreak{}SharpRestriction.\allowbreak{}GrowthBounds.\allowbreak{}DualFockLinear.\allowbreak{}CloudBAtOfRho}{87--298}
\leanref{Formalism.\allowbreak{}MainNotions.\allowbreak{}TestSideCounting.\allowbreak{}test\_\allowbreak{}bound\_\allowbreak{}of\_\allowbreak{}two\_\allowbreak{}halves}{Formalism.\allowbreak{}SharpRestriction.\allowbreak{}GrowthBounds.\allowbreak{}DualFockLinear.\allowbreak{}TestSideExplicitIdentification}{149--199}

For \(q=p^{n'}\) use the \(q^{rK}\) incoming edges of
\ref{B3.52a} given by the kernels of
\((u,v)\mapsto\bar u+A\bar v\), \(A\in M_{r,K}(\mathbf F_q)\).
The residue vectors \((-Ae_j,e_j)\) recover \(A\), so the edges
are distinct. A lattice basis at the target and its Frobenius
conjugates transport the construction to this family. The
size bound gives \(Nkc_tY\le a_0\), hence
\((1-kc_t/a_0)^N\ge1-1/Y\). Summing \eqref{eq:p5c2-plane-edge}
and applying \ref{B5.20e} proves the result with \(C\) enlarged
to absorb \(C_t,C_D,C_r+1\).
\leanref{Formalism.\allowbreak{}MainNotions.\allowbreak{}Pf14.\allowbreak{}manyEdges}{Formalism.\allowbreak{}SharpRestriction.\allowbreak{}GrowthBounds.\allowbreak{}ScaleEdges.\allowbreak{}ColSpanMap}{166--211}
\leanref{Formalism.\allowbreak{}MainNotions.\allowbreak{}Pf14.\allowbreak{}prvE\_\allowbreak{}of\_\allowbreak{}gl\_\allowbreak{}operators\_\allowbreak{}form\_\allowbreak{}representation}{Formalism.\allowbreak{}SharpRestriction.\allowbreak{}GrowthBounds.\allowbreak{}ScaleEdges.\allowbreak{}TheoremIIOfEdgeBoundConstants}{9--74}
\end{proof}

\subsection*{Proof of \ref{B5.20hb}}\label{prf:B5.20hb}\phantomsection
\begin{proof}
First take one product of the segment symbols \(A_j B'_j\) in \eqref{eq:p5c2-symbol-bound}, one term of its low monomial, and one monomial of the fixed determinant factor \eqref{eq:p5c2-fixed-factor}. All calculations in this paragraph precede the residual differential operators. In each non-left-frame occurrence the substituted segment exponents are multiples of \(M_x\). The left-frame occurrence of a shared row uses \(M_xP_t\). Thus the first residue selection either kills the entire product or leaves a plane factor times a normal factor. On the plane side the result is a fixed monomial \(X^b\) times the product of the renamed \(A_j\), with each left-frame variable raised to its boundary power. The exponent \(b\) comes from the quotient of the low monomial times the fixed determinant monomial. It vanishes on interior plane covectors: the low exponents were below \(M_x\), and the determinants use even plane rows. On the normal side there are no left-frame variables. Selecting the two row multidegrees preserves this factorization. The normal projection followed by coefficient extraction now multiplies the plane factor by a scalar.

Suppose that scalar is nonzero. On each normal prefix its type occurs in the product of the \(B'_j\) after row-degree selection. Only \(B'_j\) uses the normal letters of that prefix. Selection of row degrees commutes with the prefix action. Consequently \(k_{ju}\) is at most the greatest signed weight of \(B'_j\) on that prefix. Indeed, in an occurring irreducible a highest vector attains the first highest-weight entry, and a bound on the greatest weight in every frame holds throughout the module generated by the polynomial. The inequalities defining \(A_jB'_j\) therefore give \(v_{<x_u}(A_j)\ge H_{ju}\). For the term involving \(u-1\), the cut after \(x_{u-1}\) is the cut before \(x_u\). An intervening odd letter adds no even predecessor, so the oldest even predecessor and the reserve \(s_u\) are the same at these two letters.
\leanref{Formalism.\allowbreak{}MainNotions.\allowbreak{}PairMembershipBlocks.\allowbreak{}restN}{Formalism.\allowbreak{}SharpRestriction.\allowbreak{}GrowthBounds.\allowbreak{}PlaneCoefficients.\allowbreak{}CompDEqHomComp}{17--43}
\leanref{Formalism.\allowbreak{}MainNotions.\allowbreak{}PairMembershipBlocks.\allowbreak{}restTh}{Formalism.\allowbreak{}SharpRestriction.\allowbreak{}GrowthBounds.\allowbreak{}PlaneCoefficients.\allowbreak{}ValSCongrSupp}{151--203}

For a chain cut, keep separate variables for the left and right occurrences of each shared row. Use a quotient letter at the left occurrence and a temporary label letter at the right occurrence. On this alphabet the cut acts on exactly the original prefix of \(A_j\). The other segment factors use disjoint letters there. Projecting to the prescribed degrees after recombination commutes with this action on the separated alphabet: the degrees tested are sums of whole letter degrees with coefficients \(1\) or \(P_t\). Expand next in the temporary label monomials. Extracting any one of their coefficients preserves the cut since no temporary label lies in it. If such a monomial has exponent \(h_t\) at a shared row, adding \(b_t\) selects the residue of \(h_t+b_t\) and multiplies the quotient component by the monomial of exponent \(\lfloor(h_t+b_t)/P_t\rfloor\). The other new factor on the cut is the interior part of \(X^b\). Both have nonnegative least weight there, since their variables on the cut are vector variables. Summing the terms with the prescribed residues proves the chain bound \(H_{ju}\) for every digit component of the undifferentiated product.
\leanref{Formalism.\allowbreak{}MainNotions.\allowbreak{}TraceSpectra.\allowbreak{}restChain\_\allowbreak{}of}{Formalism.\allowbreak{}SharpRestriction.\allowbreak{}GrowthBounds.\allowbreak{}AevalValuation.\allowbreak{}MemSupportedOf}{280--359}

For a link, decompose \(A_j\) by its individual letter degrees \(\delta_y\). This preserves its prefix bound. Removing the left frame from the action can lower its least weight by at most \(\sum_{y\ {\rm in\ left\ frame}}\delta_y\). To see this, apply the inverse twist on those rows to the twisted prefix polynomial. Its entries on a vector of degree \(\delta_y\) have powers at least \(-\delta_y\). The removed rows are now outside the link, so their power substitutions preserve the resulting lower bound. The other segment factors involve, on this link, at most shared vector letters. Their least weights there are nonnegative.

Write \(d^{\rm pl}_y\) for the selected physical plane degree, and \(b_y^0\) for the row degree of \(X^b\). In a nonzero component,
\[
 b_y^0+P_y\delta_y\le d^{\rm pl}_y
 \qquad(y\text{ in the left frame}).
\]
Indeed the remaining factors have nonnegative degree in that row. The total loss is consequently at most
\[
 L_j=\sum_{y\ {\rm in\ left\ frame}}
              (d^{\rm pl}_y-b_y^0)_+/P_y .
\]
Put \(g=z_j-z_{j-1}\). In the first digit split the nonfixed degree at row \(y\), before division by \(p^{z_{j-1}-y}\), is at most \(2a_0\). Each summand in \(L_j\) is therefore at most \(2a_0p^{-g-(z_{j-1}-y)}\). The distances \(z_{j-1}-y\) are distinct nonnegative integers. As \(p\ge 3\), summation gives \(L_j\le3a_0p^{-g}\). At a right-frame cut the oldest of the preceding \(r\) even letters lies at least \(L_1+1\) places beyond \(z_{j-1}\): the first \(r\) interior letters after that frame are even, with the gap furnished by the marked tile. Thus \(a_0p^{-g}\le b_{x_u}p^{-L_1-1}\). The spacing hypothesis gives \(48Yk^2\le p^{L_1+1}\). Also \(b_{x_u}\ge a_0/p^{G'}\), so the last size hypothesis gives \(b_{x_u}/(Yk^2)\ge1\), whence \(b_{x_u}/(Yk^2)\le2s_u\). It follows that \(L_j\le s_u/8\). This proves the link bound \(H_{ju}-s_u/8\) before differentiation.
\leanref{Formalism.\allowbreak{}MainNotions.\allowbreak{}PairMembershipBlocks.\allowbreak{}restLink\_\allowbreak{}of}{Formalism.\allowbreak{}SharpRestriction.\allowbreak{}GrowthBounds.\allowbreak{}PlaneCoefficients.\allowbreak{}SupportedOfSupport}{77--199}
\leanref{Formalism.\allowbreak{}MainNotions.\allowbreak{}PairMembershipBlocks.\allowbreak{}arithCondition\_\allowbreak{}of\_\allowbreak{}one\_\allowbreak{}le}{Formalism.\allowbreak{}SharpRestriction.\allowbreak{}GrowthBounds.\allowbreak{}PlaneCoefficients.\allowbreak{}FortyEightBoundArith}{51--214}

We propagate these two bounds through a residual word. Keep the assertions simultaneously for every first digit block, every normal type and every extracted coefficient. Charge one unit to cut \(u\) for a differentiation or multiplication at row \(s\) if that physical row occurs in the segment at or before \(x_u\), and zero otherwise. The following digit calculation gives this charge at both levels of digit selection. In one coordinate write a polynomial as \(\sum_{h=0}^{P-1}X^h f_h(X^P)\), with quotient variable \(V\). Multiplication takes \(f_h\) to the next label, multiplying by \(V\) on a carry. Differentiation on a term with \(h>0\) gives \((h+PV\partial_V)f_h\) at label \(h-1\). At \(h=0\) it gives \(P\partial_V f_0\) at label \(P-1\). The Euler operator \(X\partial_X\) acts on a fixed label by \(h+PV\partial_V\). These formulas include \(P=1\). A carry or a borrow changes the input quotient degree by one in exactly the affected letter, so the simultaneous assertion supplies the requisite input block.

In any frame, multiplication or differentiation shifts the least signed weight down by at most one on a cut containing its letter, and by zero outside it. A conjugate of \(V_a\partial_{V_a}\) is a sum of \(V_b\partial_{V_c}\); the two coordinate weights have difference at least \(-1\) for either sign of the letter. Thus the Euler term also costs at most one. This proves the charge for plane operators, on the original links and, by the same digit calculation, on every quotient prefix. A quotient letter from the preceding frame is charged at its left occurrence.

For a normal operator, insert the finite sum of all joint normal projections before acting. On a prefix containing its letter, a multiplication, differentiation or Euler term raises the greatest signed weight by at most one, by the same coordinate calculation. Outside the prefix it preserves that bound. An input type with first entry \(k\) can therefore contribute to an output type with first entry \(k'\) only if \(k'\le k+1\) on a touched prefix and \(k'\le k\) on the others. Each active normal entry in the maximum defining \(H_{ju}\) uses a prefix ending at or before \(x_u\). The change in that maximum is at most the assigned charge. After projecting, each output coefficient is a scalar linear combination of plane coefficients of the contributing input types. Least-weight bounds pass to such combinations. This proves the induction step for normal operators as well.
\leanref{Formalism.\allowbreak{}MainNotions.\allowbreak{}PairMembershipBlocks.\allowbreak{}oneOpStep\_\allowbreak{}of}{Formalism.\allowbreak{}SharpRestriction.\allowbreak{}GrowthBounds.\allowbreak{}PlaneCoefficients.\allowbreak{}RestOneOpOf}{106--131}
\leanref{Formalism.\allowbreak{}MainNotions.\allowbreak{}LetterSubstitution.\allowbreak{}noOpStep}{Formalism.\allowbreak{}SharpRestriction.\allowbreak{}GrowthBounds.\allowbreak{}MonomialSupport.\allowbreak{}Lam1SPderivAndXMulLeAddIndicator}{38--146}

Each \(X_{s,a}\partial_{s,b}\) incurs at most two charges. In the segment construction of Lemma~\ref{B5.20h}, the entire segment, including both frames, lies in a window of length \(G'\), so it has at most \(N_1\) physical rows; there are at most \(c_t\) occurrences per marked row in a residual word. For \(C\ge16\) the last size bound of Lemma~\ref{B5.20h} gives
\[
 16N_1c_tYk^2p^{G'}\le a_0 ,
\]
using \(N_1\le rN\) and \(c'\ge rNc_t\). Hence \(2N_1c_t\le s_u/8\). Subtracting this from the chain bound gives the first assertion. Subtracting it from the link bound gives at least \(H_{ju}-s_u/4\), which implies the second assertion.

To apply the induction, write the primitive combination in the exact symbol representatives as \(\sum_{i,b}c_{ib}x_{ib}\), with \(i\) the flat symbol and \(b\) the residual tableau. For each summand the residual contraction in the proof of Lemma~\ref{B5.20h} gives a linear combination of the words just considered applied to symbol \(i\). Changing to the edge coordinates conjugates a row matrix unit to a linear combination of matrix units in the same row. The two bounds are linear subspace conditions over the fraction field. They therefore hold for the sum with coefficients \(c_{ib}\), for \(\omega^{-a}\rho_\beta\sum c_{ib}x_{ib}\) when the component has exact order \(a\), and after the nonzero scalar division defining \(P_c\). In particular the second assertion puts \(P_c\) in the image of every link projector. Moving the self-adjoint projectors across the factorial form now gives the stated coefficient identity. At a minimizing monomial \(v([X^{m_0}]P_c)=-v(m_0!)/2\), which gives its valuation.
\leanref{Formalism.\allowbreak{}MainNotions.\allowbreak{}TestSideCounting.\allowbreak{}cloudBAt\_\allowbreak{}of\_\allowbreak{}rho}{Formalism.\allowbreak{}SharpRestriction.\allowbreak{}GrowthBounds.\allowbreak{}DualFockLinear.\allowbreak{}CloudBAtOfRho}{87--298}
\leanref{Formalism.\allowbreak{}MainNotions.\allowbreak{}TestSideWords.\allowbreak{}memAt\_\allowbreak{}of\_\allowbreak{}wordConditionAt}{Formalism.\allowbreak{}SharpRestriction.\allowbreak{}GrowthBounds.\allowbreak{}TwistedSubstitution.\allowbreak{}MemRed}{40--62}
\end{proof}

\subsection*{Proof of \ref{B5.20hc}}\label{prf:B5.20hc}\phantomsection
\begin{proof}
Order a segment's interior first and append its right frame in chronological order \(r,\ldots,1\). Thus each of its original links is an initial segment ending before the last letter. We give the integral basis argument used to separate the variables.

For this ordered word use paths of horizontal strips as in Definition~\ref{B5.20l}, in rank \(r\), with a standard terminal tableau and no restriction steps. Form a polynomial for each path and tableau by the exterior multiplication and contraction formulas in Lemma~\ref{B5.20h}, with power one throughout. At a cut let \(\mu\) be the path row and \(h(c)=\#\{a:\mu_a\ge c\}\) its height in column \(c\). Under the prefix action the state of a positive column is in \(\bigwedge^{h(c)}V\); a nonpositive column uses \(\bigwedge^{r-h(c)}V^\vee\). Later operations are linear in that state with coefficients in the suffix variables. The least possible first-coordinate weights of these two exterior powers are respectively
\[
 {\bf1}_{h(c)=r},\qquad -{\bf1}_{h(c)<r}.
\]
The sum over columns is \(\mu_r\), by counting integer columns on either side of zero. Each path polynomial consequently satisfies every lower bound read from its path rows.

These polynomials form an integral basis. Here are the independence and count arguments. They work over every residue field. Under a generic matrix change of frame all weights at the end of a path are dominated by its terminal row \(\mu\): in each exterior power the sum of the largest \(a\) coordinate weights is attained by the initial index set, or by its complementary final set for a dual column. Attaining \(\mu\) forces that choice in every column. By exterior change of basis its coefficient is the canonical path polynomial with these index sets times the product of the corresponding flag minors of the generic matrix. For a negative column the complementary-minor identity changes the inverse minor to a flag minor divided by the determinant, up to sign. At fixed \(\mu\) the determinant power is common. The flag products for distinct standard terminal tableaux are independent: under row-by-row lexicographic order on the entries of the generic matrix, the diagonal monomial is the unique leading monomial of each flag minor and hence of each flag product, with leading coefficient one, and its exponents give the entries in each tableau row, whose weak ordering recovers the row.

The canonical path polynomials at fixed \(\mu\) are independent even over the field of the generic matrix entries. Inspect the last letter first, ordering its coordinate exponents lexicographically from \(r\) down to \(1\) for a vector and from \(1\) up to \(r\) for a covector. If the preceding row is \(\lambda\), the leading coefficient in this letter is the preceding canonical polynomial, up to sign, at the monomial with exponents \(|\mu_a-\lambda_a|\). Indeed a changing column ends on an initial set or its complementary final set. In the exterior sum the extreme admissible coordinate is the index of the changed row. Erasing it in a wedge or adjoining it in a contraction leaves the preceding canonical index. The term has coefficient \(1\) or \(-1\). Multiplication over columns gives the indicated exponents. Since the direction of the strip is fixed, these exponents recover \(\lambda\). Induction from the empty word proves independence. In a relation among all the labelled paths take a maximal terminal row in dominance order and extract that weight after the generic change of frame. Only this row contributes. Independence of the canonical polynomials, followed by independence of the flag products, kills all its coefficients. This proves independence over the residue field.
\leanref{Formalism.\allowbreak{}InvariantFactorPaths.\allowbreak{}GelfandTsetlinBases.\allowbreak{}pathProductLinearIndependent}{Formalism.\allowbreak{}ThetaLogCoordinates.\allowbreak{}JacobiPathWeights.\allowbreak{}JacobiMinorProofs.\allowbreak{}IsWeightedHomogeneousZeroStripPathWord}{752--890}
\leanref{Formalism.\allowbreak{}InvariantFactorPaths.\allowbreak{}GelfandTsetlinBases.\allowbreak{}Top.\allowbreak{}psiPath\_\allowbreak{}basis}{Formalism.\allowbreak{}ThetaLogCoordinates.\allowbreak{}JacobiPathWeights.\allowbreak{}PathBasisTheorem}{23--26}

By the horizontal-strip dimension identity used in Lemma~\ref{B5.20m}, summing terminal tableau counts gives the product of the dimensions of the symmetric powers of the individual letters. For a dual letter the identity follows by dualizing, and shifting all entries of a row by the same integer tensors with a determinant character of dimension one. This count equals the number of ordinary monomials of the prescribed multidegree. The square coefficient matrix is therefore invertible modulo the maximal ideal and has unit determinant over the coefficient DVR.

At one cut the paths with \(\mu_r\ge\theta\) span the required subspace over the fraction field. They lie in it by the exterior-power computation. Their number is its dimension: Pieri decomposition on the prefix gives one copy along each prefix path, the specified types are precisely those with \(\mu_r\ge\theta\), and continuing along unrestricted suffix letters multiplies the dimension by their symmetric-power dimensions. Since the full family is an integral basis, an integral vector satisfying the cut has zero coefficients at all other paths and integral coefficients at the retained ones. Applying this to each cut in the same basis gives the simultaneous assertion. The alphabets comprising an interior and its right frame are disjoint across segments, so their tensor product gives the corresponding expansion for \(T\).

Initially write every boundary variable as \(R_t\). In each column that uses at least one interior letter of its segment, replace a boundary occurrence at \(h>1\) by \(S_{j,h}\). Keep slot \(1\) on \(R_{j,1}\). Leave columns without interior occurrences on the label letters. Each change recombines to the original column. A column is multilinear in its used letters, since a horizontal strip changes each column at most once per letter. There are at most \(d_{I_j}\) columns using the interior. Hence each surplus degree is at most \(d_{I_j}\).

Consider the cut for \(m\) on the interior and surplus letters \(h>m\). For a changed column this action pulls back to the original prefix action, so the column retains its exterior-power lower bound. An unchanged column is independent of the letters in the new cut. Its old lower bound is at most zero: a positive column using no interior letter could have received at most \(r-m<r\) vector occurrences by the cut, and a nonpositive column has lower bound at most zero at every height. The column product therefore retains the full path lower bound at every cut.

Finally expand in all \(R\)-monomials. Their coefficients separately satisfy the bounds on the interiors and surplus letters. Replace \(R_t^h\) coordinatewise by \(R_t^{h\bmod P_t}W_t^{\lfloor h/P_t\rfloor}\). This preserves integrality and recombination. In the enlarged link the additional \(W\)'s are vectors and contribute nonnegative least weight in every frame. This proves the assertion after summing the expansion.
\end{proof}

\subsection*{Proof of \ref{B5.20hd}}\label{prf:B5.20hd}\phantomsection
\begin{proof}
Fix an allowable exponent \(\zeta\). Its block fixes all letter degrees and all individual coordinates on the label and surplus letters. Denote its surplus exponent array by \(M\), extended by zero on the interiors of the original alphabet. Recombination has exponent \(\alpha+M\). Prefix projections preserve the block: they preserve each letter degree and act only on interior and quotient letters. Thus the proof can be made separately on each block.

Let \(\partial^{[M]}\) be the product of the divided coordinate derivatives, and set \(A=D_{u^{-1}}\partial^{[M]}D_uF\). The coefficient tested at \(\alpha\) is
\[
 A_\alpha
 =u(\alpha)^{-1}u(\alpha+M)F_{\alpha+M}
                   \prod_i\binom{\alpha_i+M_i}{M_i}.
\]
Write \(C_B=\prod_t(P_t!)^{|\beta_t|}\), where \(|\beta_t|\) is the fixed quotient degree of slot \(t\). Factoring the boundary factorials coordinatewise gives
\[
 \alpha!\,M!=C_B\,\zeta!\,\Gamma(\zeta)/u(\alpha),\qquad
 \Lambda(X^\zeta)=C_B\,\zeta!\,\Gamma(\zeta)A_\alpha .
\]
In particular \(C_B\) is constant on the block.
\leanref{Formalism.\allowbreak{}MainNotions.\allowbreak{}PlaneGram.\allowbreak{}TwistData.\allowbreak{}coeff\_\allowbreak{}GM}{Formalism.\allowbreak{}SharpRestriction.\allowbreak{}GrowthBounds.\allowbreak{}MonomialOrderSegments.\allowbreak{}SubstitutedExponentApplyLabel}{108--125}
\leanref{Formalism.\allowbreak{}MainNotions.\allowbreak{}PlaneGram.\allowbreak{}TwistData.\allowbreak{}Lam\_\allowbreak{}monomial\_\allowbreak{}blk}{Formalism.\allowbreak{}SharpRestriction.\allowbreak{}GrowthBounds.\allowbreak{}MonomialOrderSegments.\allowbreak{}SubstitutedExponentApplyLabel}{141--158}

We compute the chain orders of the digit component of \(A\) at \(\rho\). Every coefficient of \(F\) has valuation at least \(-n_0\), because its monomial factorial divides the product of the row factorials. Perform the calculation modulo coefficients of valuation at least \(c'-n_0\). Both \(D_u\) and \(D_{u^{-1}}\) may then be replaced by their integral binomial approximations of degree at most \(6c'\) from Lemma~\ref{B5.20ha}. Its hypotheses hold for the first row split, whose digits lie between zero and \(M_x-1\), with \(p\ge13\). On any fixed boundary class, substitution of \(\rho_t+P_t\beta_t\) preserves polynomial degree and integer-valuedness. The expansion in binomials of the free exponents is integral by forward differences. A factor \(\binom{\beta_a}{h}\) acts by \(W_a^h\partial_{W_a}^{[h]}\). In any frame its least signed weight shift is at least \(-h\), since the coordinate weight differences are at least \(-1\). Each of the two approximations therefore costs at most \(6c'\) at every cut. All exact units, approximate multipliers and divided derivatives in this computation preserve coefficient lower bounds.

For the intermediate derivative consider a single shared row with power \(P\), input label \(d\), and derivative multi-index \(m\). In one coordinate define
\[
 a_\ell=[T^m](1+T)^d\bigl((1+T)^P-1\bigr)^\ell .
\]
The binomial theorem gives
\[
 \binom{d+Pb}{m}=\sum_{\ell=0}^{m}a_\ell\binom b\ell.
\]
Here \(a_\ell=0\) if \(d+P\ell<m\). Since \((1+T)^P-1=T^P+pH(T)\) with \(H\in\Z[T]\), at most \(\lfloor m/P\rfloor\) of the \(\ell\) factors can contribute \(T^P\). Hence \(v(a_\ell)\ge(\ell-\lfloor m/P\rfloor)_+\). Multiply these expansions over the coordinates of the row. Terms with \(\sum_a(\ell_a-\lfloor m_a/P\rfloor)_+\ge c'\) contribute to the discarded coefficient error. Every remaining nonzero term has
\[
 |\ell|\le |m|/P+c',\qquad d_a+P\ell_a-m_a\ge0.
\]
On the quotient polynomial its operator is a scalar times
\[
 W^{\,\lfloor(d+P\ell-m)/P\rfloor}\partial_W^{[\ell]} .
\]
Indeed on \(W^b\) its coefficient is \(\prod_a\binom{b_a}{\ell_a}\), its output quotient is \(b-\ell+\lfloor(d+P\ell-m)/P\rfloor\), and its output label is \(d-m\bmod P\). The multiplier monomial uses a vector letter and has nonnegative least weight. Thus the loss is at most \(|m|/P+c'\) on a cut containing this quotient letter and zero on other cuts. This calculation treats every input and output label, including borrows.
\leanref{Formalism.\allowbreak{}MainNotions.\allowbreak{}G.\allowbreak{}GLWeightTheory.\allowbreak{}derivative\_\allowbreak{}digit\_\allowbreak{}expansion}{Formalism.\allowbreak{}SharpRestriction.\allowbreak{}GrowthBounds.\allowbreak{}RootWeights.\allowbreak{}BinomialExpansionCoeffZeroOfDeg}{110--173}

The quotient letters on a chain cut are exactly those of the preceding frame. By allowability each \(|M_t|\le d_{I_{j-1}}\) there. The hypothesis on \(\sigma_{ji}\) gives
\[
 12c'+\sum_{t\ {\rm in\ frame}\ j-1}(|M_t|/P_t+c')
 \le \sigma_{ji}+(12-12r)c'\le\sigma_{ji}.
\]
It follows that in the digit component of \(A\) the retained polynomial \(Q\) satisfies every bound \(\theta_{ji}-\sigma_{ji}\); each coefficient of the difference has valuation at least \(c'-n_0\). Multiply \(Q\) by the fixed label and surplus monomial and select the block letter degrees. Call the result \(Q_B\). Both operations commute with all the prefix actions. Consequently \(P_{\rm pre}Q_B=Q_B\).

For every \(\zeta'\) in this block, with original collapsed exponent \(\alpha'\), put
\[
 e_{\zeta'}=g_-(\zeta')\Gamma(\zeta')A_{\alpha'}-[X^{\zeta'}]Q_B .
\]
The coefficient formula gives \(v(A_{\alpha'})\ge-n_0\), and \(g_-\Gamma\equiv1\bmod p^{c'}\). Hence \(v(e_{\zeta'})\ge c'-n_0\). For \(W\) supported on the block we obtain the exact equality
\[
 \Lambda(D_{g_-}W)
 = C_B\langle Q_B,W\rangle+
       C_B\sum_{\zeta'}W_{\zeta'}e_{\zeta'}\,\zeta'!.
\]
The first terms agree for \(W=P_{\rm pre}X^\zeta\) and \(W=X^\zeta\), by self-adjointness.

For either choice of \(W\), the projector bound, also valid for the identity with loss zero, gives
\[
 v(W_{\zeta'})\ge\tfrac12v(\zeta!)-D_m-\tfrac12v(\zeta'!).
\]
We record the factorials needed to use this estimate. For any original exponent \(b\) in the prescribed degrees,
\(v(b!)\ge2n_0-k'\). In a row of degree \(d_i\), successively choosing its coordinate exponents uses at most \(r\) binomial coefficients with upper argument at most \(d_i\). Each costs at most \(\lfloor\log_p d_i\rfloor\) for \(d_i>0\), by the carry formula from Legendre's identity in Lemma~\ref{B3.19}. A zero row costs zero. Splitting \(b=\alpha'+M\) costs at most another \(k'\) by the same bound on each coordinate binomial. The block factorial identity therefore gives
\[
 v(C_B)+v(\zeta'!)=v(\alpha'!)+v(M!)\ge 2n_0-2k'.
\]
It also holds for \(\zeta\). Average these two inequalities and add the bounds for \(W_{\zeta'}\) and \(e_{\zeta'}\). Each summand of the error has valuation at least \(c'+n_0-2k'-D_m\). Subtraction proves this bound for the difference of the two evaluations with \(D_{g_-}\).

Finally \(D_{g_+}X^\zeta=g_+(\zeta)X^\zeta\). The additional error is
\((g_+(\zeta)g_-(\zeta)-1)\Lambda(X^\zeta)\). Its first factor has valuation at least \(c'\). The second has valuation at least \(\tfrac12v((\alpha+M)!)\ge n_0-k'/2\), by \(v_{\rm F}(F)=0\) and the unit property of \(u\). This proves the claimed precision on a monomial. Taking the finite linear combination with the integral coefficients of \(Z\) proves the result.
\end{proof}

\subsection*{Proof of \ref{B5.20t}}\label{prf:B5.20t}\phantomsection
\begin{proof}
Use the invariant projection on a space of linear maps, first extracting its zero-weight part. For lattices \(L_1,L_2\) preserved by divided roots, the lattice of maps from \(L_1\) to \(L_2\) is preserved by the divided adjoint action. Its \(j\)-th operator sends \(T\) to
\[
 \sum_{b=0}^j(-1)^b E_{uv}^{(j-b)}T E_{uv}^{(b)}.
\]
Weight-zero extraction is integral when the lattices are graded. In the applications below all weights of the space of maps have entries in \([-l,l]\), since both its arguments have weights from \(A\). Thus the row recursion \eqref{eq:p3c5-projector-commutator} terminates at \(q=l\). By \eqref{eq:p3c5-projector-denom}, adjoining one of the \(j\) earlier indices in row \(j\), numbered from zero, costs at most \(\lfloor\log_p(l+h-1)\rfloor\). Summing \(j\) gives \(d\), as in \eqref{eq:p3c5-projector-integral}. These identities use the matrix-unit commutators and the weight interval, so they apply to each of the spaces of maps used here. Write \(I(T)\) for the resulting invariant map. On endomorphisms of an irreducible module \(U\) it is \(\operatorname{tr}(T)1_U/\dim U\), since the recursion preserves trace. It commutes with composition by intertwiners, by the same recursion.

Choose one irreducible copy \(U\subset A\) of type \(\lambda\), and a basis \(b_j\) of \(U\cap L\), with dual basis \(b_j^\vee\). For \(v\in A\) form
\[
 (\dim U)\sum_j I\bigl(b_j^\vee\otimes v\bigr)(b_j).
\]
On a copy of \(U\) this gives \(v\) by the trace formula. On any other irreducible type it is zero since there is no nonzero intertwiner from \(U\) to that type. Complete reducibility applies over \(\Q_p\) to the split semisimple Lie algebra with the center acting diagonally. Thus this is the type projector. For \(v\in L\), the maps inside the sum take \(U\cap L\) into \(L\). The denominator calculation gives the first assertion, including sums of types.

On \(A_\lambda\) take the map extracting weight \(\lambda\). It takes \(L\cap A_\lambda\) into \(\Delta_\lambda\), since vectors of this weight in that component are highest. Apply the same recursion to maps from \(L\cap A_\lambda\) to \(\Delta_\lambda\). Both lattices are divided-root stable and weight graded. On each irreducible copy the map being averaged has trace one. Its average is the inclusion divided by \(\dim U\). Multiplication by \(p^d\dim U\) gives the second inclusion. In rank one both assertions follow directly from integral weight extraction.
\end{proof}

\subsection*{Proof of \ref{B5.20u}}\label{prf:B5.20u}\phantomsection
\begin{proof}
A cap on a group of coordinates says that the first-coordinate weight of every translate is at most its threshold. Indeed, the orbit span is the smallest stable subspace containing the polynomial. In a representation of the general linear group the largest first-coordinate weight of an irreducible summand is the first entry of its highest weight.

Restriction commutes with the actions on the proper prefixes and with the action of \(\mathfrak{gl}(I)\) on the whole prefix. To see the cap on the first \(m\) letters, differentiate a polynomial of last degree \(d\) in \(W\). The span of all its first derivatives, as the polynomial varies in a full cap, stays in that cap. The derivatives transform among themselves and each preserves or lowers first-coordinate weight. Differentiating \(d\) times shows that every coefficient in \(W\) satisfies the preceding cap, since the full action at last degree zero is the preceding action. This proves \(RY_d\subseteq X_d\). With no finite full cap each element of \(X_d\) lifts to itself. The same holds in degree zero.

Let \(C\) be the least finite full cap. Write \(A_{ab}=E^m_{ab}\) for the action on the earlier letters and \(D_W=\sum_bW_b\partial_{W_b}\). For \(J=I\) or \(J=\{1,\ldots,r\}\), and \(a\in J\), put
\[
 U_a^J=W_a(C-D_W)-\sum_{b\in J}W_bA_{ab}.
\]
Restriction intertwines the full \(U_a\) with \(U_a^I\) for \(a\in I\). It also intertwines \(\partial_{W_a}\) and the matrix units with both indices in \(I\).

The full \(U_a\) sends \(Y_j\) into \(Y_{j+1}\). At an earlier cut, \(A_{ab}\) is the sum of the action at that cut and the action on the remaining earlier letters, both of which preserve its cap. For the full cut put \(J_{ab}=A_{ab}+W_a\partial_{W_b}\). Then
\[
 U_a=CW_a-\sum_bW_bJ_{ab},\qquad [J_{uv},U_a]=\delta_{va}U_u .
\]
Consequently the span of the \(U_aY_j\) is stable under the full group. Only \(U_1\) can raise first-coordinate weight. On the weight-\(C\) boundary, \(J_{11}=C\) and \(J_{1b}=0\) for \(b>1\), since the latter operators raise that weight. Thus \(U_1\) vanishes on the boundary. The span satisfies the full cap.

We prove surjectivity by induction on \(j\). Write \(d'>0\) for the new degree. Derivatives \(\partial_{W_a}\), \(a\in I\), take \(X_{d'}\) into \(X_{d'-1}\) by the derivative argument above for the smaller group. By induction and the intertwining formula,
\[
 \Omega X_{d'}\subseteq RY_{d'},\qquad
 \Omega=\sum_{a\in I}U_a^I\partial_{W_a}
 =(C-d'+1)d'-\sum_{a,b\in I}W_b\partial_{W_a}A_{ab}.
\]
Here is the kernel calculation. The operator commutes with \(\mathfrak{gl}(I)\), since the two indices in its sum are contracted. A nonzero kernel on \(X_{d'}\) therefore contains a highest vector \(H=\sum_\beta W^\beta v_\beta\) of weight \(\theta\) with \(\theta_1\le C\). Choose a nonzero coefficient maximizing \(\sum_{a\in I}a\beta_a\), and put \(\lambda_a=\theta_a-\beta_a\). For \(a<b\), the equation \((A_{ab}+W_a\partial_{W_b})H=0\) gives
\[
 A_{ab}v_\gamma=-(\gamma_b+1)v_{\gamma-\epsilon_a+\epsilon_b}
\]
when \(\gamma_a>0\), and gives zero when \(\gamma_a=0\). This implies \(\lambda_a\ge\theta_b\). Indeed, with \(k=\beta_b\), start at \(\gamma=\beta+k(\epsilon_a-\epsilon_b)\). Applying \(A_{ab}^k\) gives a nonzero multiple of \(v_\beta\), and one further application gives zero by maximality. In the finite-dimensional \(\mathfrak{sl}_2\) for \(a,b\), the resulting highest vector has weight \(\lambda_a-\lambda_b\) and lies \(k\) steps up a root string. Hence \(\lambda_a-\lambda_b\ge k\).

At \(W^\beta\), the diagonal summands \(W_a\partial_{W_a}A_{aa}H\) contribute \(\sum_a\beta_a\lambda_a v_\beta\). For \(a>b\) the required coefficient of \(H\) has larger \(\sum a\beta_a\) and vanishes. For \(a<b\) the highest-vector equation gives \(-(\beta_a+1)\beta_b v_\beta\). Thus the coefficient of \(\Omega H\) is \(v_\beta\) times
\[
 (C-d'+1)d'-\sum_a\beta_a\lambda_a
                  +\sum_{a<b}(\beta_a+1)\beta_b .
\]
Successive use of \(\lambda_a\ge\lambda_{a+1}+\beta_{a+1}\) and \(C\ge\lambda_1+\beta_1\) gives \(C-\lambda_a\ge\sum_{j\le a}\beta_j\). Substitution and \(\sum_a\beta_a=d'\) bound the displayed integer below by \(d'+\sum_{a<b}\beta_b>0\). This contradicts \(\Omega H=0\). Hence \(\Omega\) is injective on \(X_{d'}\). Its image lies in \(RY_{d'}\subseteq X_{d'}\), and finite dimension gives equality.
\leanref{Formalism.\allowbreak{}InvariantFactorPaths.\allowbreak{}L4.\allowbreak{}omega\_\allowbreak{}injective\_\allowbreak{}holds}{Formalism.\allowbreak{}ThetaLogCoordinates.\allowbreak{}FockResolutionOnto.\allowbreak{}FockMonomialProofs.\allowbreak{}OmegaInjectiveHolds}{673--726}

We now keep the thresholds fixed and vary the last degree over all \(j\). Let \(d_0\) be the sum of the earlier degrees. If \(C\ge S\), the full cap is automatic at the prescribed degree \(d\) and an integral element lifts to itself. Suppose \(C<S\). Set
\[
 X=\bigoplus_{j\ge0}X_j,\qquad L_X=X\cap L,\qquad
 L_R=\bigoplus_{j\ge0} R(Y_j\cap L).
\]
The preceding kernel calculation also bounds the possible degrees. On a highest vector in \(X_j\), its last nonzero leading coefficient as chosen above satisfies \(\lambda_{r-1}+j\le C\). Since \(\lambda_{r-1}\) is an earlier-letter weight, it is at least \(-d_0\). Thus \(X_j=0\) for \(j>C+d_0\).

On \(X\) take matrix indices \(0,1,\ldots,r-1\) and use the operators
\[
 E_{00}=C-D_W,\quad E_{0b}=\partial_{W_b},\quad
 E_{a0}=U_a^I,\quad E_{ab}=A_{ab}+W_a\partial_{W_b}\quad(a,b\in I).
\]
They preserve \(X\) by rational surjectivity and the intertwining just proved. They satisfy the matrix-unit relations. In particular \([\partial_{W_b},U_a^I]=\delta_{ab}(C-D_W)-E_{ab}\), the \(U_a^I\) commute by expansion using \([A_{ab},A_{ij}]=\delta_{bi}A_{aj}-\delta_{ja}A_{ib}\), and \([E_{ij},U_a^I]=\delta_{ja}U_i^I\). The weights are read off on monomials. Their zeroth entry is \(C-j\), and the others are earlier signed weights plus the corresponding \(W\)-exponents. Each lies in \([-d_0,C+2d_0]\) on the nonzero spaces. For the zeroth entry this uses \(C\ge-d_0\) whenever \(X\ne0\). The coordinate differences are at most \(C+3d_0\le4S\).

Both \(L_X\) and \(L_R\) are stable under the divided roots. For the derivatives this follows from the binomial formula for a divided derivative. For \(E_{ab}\), \(a\ne b\) in \(I\), it follows by taking coefficients of an elementary shear. For \(U_a^J\) one can check all divided powers at once. Put \(q=1+tW_a\). Multiply by \(q^C\) and make the following substitutions in a polynomial:
\[
 W\longmapsto W/q,\qquad
 V_b\longmapsto V_b-tW_bV_a/q\ (b\in J),\qquad
 Z_a\longmapsto Z_a+t\sum_{b\in J}W_bZ_b .
\]
Here \(V\) runs through the earlier vectors, \(Z\) through the earlier covectors, and their undisplayed coordinates stay fixed. These substitutions compose by addition of \(t\), including the prefactor, and their derivative at zero is \(U_a^J\). The coefficient of \(t^h\) is therefore \((U_a^J)^h/h!\). All coefficients are integral, also for negative integral \(C\), since \(q\) has constant term one. Taking \(J\) full and then putting \(W_r=0\) gives the formula for \(J=I\). Applying a full divided power to an integral lift thus gives an integral lift of the restricted divided power. Weight extraction on \(L_X\) is coefficient extraction by the monomial weights.
\leanref{Formalism.\allowbreak{}InvariantFactorPaths.\allowbreak{}L4.\allowbreak{}raising\_\allowbreak{}divided\_\allowbreak{}integral}{Formalism.\allowbreak{}ThetaLogCoordinates.\allowbreak{}FockResolutionOnto.\allowbreak{}FockMonomialProofs.\allowbreak{}DerivativeDividedIntegral}{434--438}
\leanref{Formalism.\allowbreak{}InvariantFactorPaths.\allowbreak{}L4.\allowbreak{}restrictedLattice\_\allowbreak{}divided\_\allowbreak{}power}{Formalism.\allowbreak{}ThetaLogCoordinates.\allowbreak{}FockResolutionOnto.\allowbreak{}FockMonomialProofs.\allowbreak{}RestrictedLatticeDividedPower}{52--78}

Apply \ref{B5.20t} to \(X,L_X\) with \(l=4S\). For \(x\in L_X\), the component \(x_\lambda\) of each type satisfies \(p^b x_\lambda\in L_X\). A second application puts \(p^{2b}x_\lambda\) in the divided-root span of integral highest vectors for the displayed action. Each such vector is annihilated by every \(E_{0b}=\partial_{W_b}\). In characteristic zero it is independent of \(W\), so belongs to \(X_0\cap L\) and lifts to itself. Divided-root stability puts the whole span in \(L_R\). Sum over \(\lambda\) and take the last-degree-\(d\) component. This gives the asserted integral lift.
\leanref{Formalism.\allowbreak{}InvariantFactorPaths.\allowbreak{}L4.\allowbreak{}sat\_\allowbreak{}main}{Formalism.\allowbreak{}ThetaLogCoordinates.\allowbreak{}FockResolutionOnto.\allowbreak{}FockMonomialProofs.\allowbreak{}RestrictedSaturationHolds}{49--88}
\end{proof}

\subsection*{Proof of \ref{B5.20v}}\label{prf:B5.20v}\phantomsection
\begin{proof}
Write \(x=W_r\). For a cut containing \(W\), multiplication by \(x^e\) raises the maximum first-coordinate weight over all translates by exactly \(e\). To see this, take a generic frame. The largest weight of the translated polynomial is attained on a nonempty Zariski open set, and the coefficient of \(W_1\) in the translated \(x\) is nonzero on another such set. On their intersection the top weight terms multiply to a nonzero polynomial. For an earlier cut \(x\) is fixed. Division by \(x^e\) therefore gives exactly the stated shifts, preserves coefficient integrality, and subtracts \(e\epsilon_r\) from the weight. It commutes with the raising operators of \(\mathfrak{gl}_{r-1}\).

After division, \ref{B5.20u} describes the rational image before the weight and highest-vector conditions. Full weight extraction commutes with restriction. On the fixed weight, projection onto highest vectors of the smaller group is projection onto the isotypic component whose highest weight equals the restriction of that weight. If the restricted weight is not dominant, the projection is zero. This description follows because the common kernel of the raising operators in an irreducible module is its highest line. Restriction is equivariant for the smaller group, and the prefix conditions are stable under it. These projections therefore commute with restriction and preserve the caps. This proves \(\Q_pG_e=x^eH\).

For an integral \(y\in H\), first lift \(p^{2e_2}y\) by \ref{B5.20u}. The total degree after division is \(S-e\le S\). Extracting the specified full weight from the lift is integral. For the smaller-group type projection use \ref{B5.20t}. In any coordinate the weights of the signed polynomial space lie between minus the sum of covector degrees and the sum of vector degrees, an interval of width at most \(S\). The coefficient lattice is weight graded and stable under divided elementary shears. Thus \(p^{e_1}\) clears the projection. Multiplying the result by \(x^e\) gives an element of \(M_{\ge e}\) with initial part \(p^{e_1+2e_2}x^ey\).

For the partner projection, first keep just the monomials with \(x\)-exponent \(e\) and strip off \(x^e\). On the resulting space impose the full-group type caps for prefix lengths \(0,\ldots,m\), combining coincident caps by their minimum. At length \(m\) include the bound inherited from each full cut. On all letters impose the smaller-group type cap, extract the full weight, and impose the smaller-group highest-vector projection. These projections commute. A type projector for an earlier full group commutes with its matrix units and with the matrix units on the remaining letters, hence with the later group action. The smaller group preserves the exponent-zero condition on \(x\). The factorial pairing interchanges \(E_{ab}\) and \(E_{ba}\), so distinct isotypic components are orthogonal, and weight extraction is orthogonal on monomials. Restoring \(x^e\) multiplies the factorial form on the retained monomials by a common scalar. The resulting map is consequently the required orthogonal projection.

At lengths zero and one there is a single full-group type: respectively the trivial type and a symmetric power of the standard module or its dual. Thus the two projections are integral. Lengths \(2,\ldots,m\) cost at most \((m-1)e_1\) altogether by \ref{B5.20t}. The smaller-group cap and the highest-vector projection cost at most \(e_1\) each. This gives \(D=(m+1)e_1\). For \(m=0\) the specified full weight singles out at most one monomial in \(W\). The caps and the highest-vector equations either retain that line or kill it. A retained monomial is its own integral lift and coefficient extraction is its orthogonal projector.
\leanref{Formalism.\allowbreak{}InvariantFactorPaths.\allowbreak{}L4.\allowbreak{}highest\_\allowbreak{}vector\_\allowbreak{}projection\_\allowbreak{}holds}{Formalism.\allowbreak{}ThetaLogCoordinates.\allowbreak{}FockResolutionOnto.\allowbreak{}FockMonomialProofs.\allowbreak{}HighestVectorProjectionPos}{252--514}
\end{proof}

\subsection*{Proof of \ref{B5.20p}}\label{prf:B5.20p}\phantomsection
\begin{proof}
Let \(u\) be primitive in \(G_e\). Select a coefficient \(a=u_\alpha\) of least valuation. Then \(a^{-1}u\in V_e\cap P_{\Z_p}\). If \(v_p(a)>c\), the equality
\(u=(a/p^c)(p^ca^{-1}u)\) contradicts primitivity. The vector \(p^{c+D}y_\alpha\) belongs to \(G_e\), and
\[
 v_p([u,p^{c+D}y_\alpha]_0)
 =c+D+v_p(a)-\sum_i v_p\binom{d_i}{\alpha_{i1},\ldots,\alpha_{ir}}
 \le D+2c .
\]
Apply this to every source basis vector in a Smith reduction of the Gram matrix over \(\Q_p\) by unimodular integral row and column operations. A zero diagonal would annihilate that row. A diagonal of valuation greater than \(D+2c\) would take every integral partner to a value of greater valuation. Hence every diagonal is nonzero with the stated bound, and taking their product proves the assertion.
\end{proof}

\subsection*{Proof of \ref{B5.20q}}\label{prf:B5.20q}\phantomsection
\begin{proof}
Each \(M_{\ge e+1}\) is saturated in \(M_{\ge e}\), as it is the kernel of coefficient extraction. Starting at \(e=o\) and extending bases downwards gives the asserted basis. In particular the \(g_b\) of order \(e\) form an integral basis of \(G_e\).

First take monomial inputs with coefficient one. The two substitutions have integer coefficients. In their expansions write the two powers of \(\eta\) as \(n-S+v,n-S+v'\), before applying contractions. The power of \(W_\rho\) in each term is at most \(v\), respectively \(v'\). For the term with \(l\) contractions move \(E^l\) across the form. If \(a,a'\le l\) contractions on the two sides use \(\partial_{Y_\rho}\), matching monomials require \(v-a=v'-a'=b\ge0\). The \(n\)-dependent factors, apart from constants, are
\[
 \binom{o+n+r+z-1-l}{o-l}
 (n-S+v)_{\underline a}(n-S+v')_{\underline {a'}}
 \frac{(n-S+b)!}{(n-S)!},
\]
where \(x_{\underline j}=x(x-1)\cdots(x-j+1)\). The displayed expression is polynomial of degree at most
\(o-l+a+a'+b=o-l+a'+v\le o+S\). The binomial in front follows by cancelling \(\binom{o+n+r+z-1}{l}\) against the prefactor. This proves polynomiality over \(\Q\) on monomials and hence over \(\Q_p\) on the given basis.

We compute the top term first on real monomials, where the factorial form is positive. On bidegree \((o,n)\) the spaces \(E^j\ker F\), with the kernel taken in degree \((o-j,n-j)\), are mutually orthogonal and span. Indeed \(F\) is adjoint to \(E\), so the orthogonal complement of \(E\) times the preceding bidegree is \(\ker F\). Iterate this equality. From \([F,E]=r+\deg_W+\deg_Y\), the eigenvalue of \(T_z\) on \(E^j\ker F\) is
\[
 \sum_l(-1)^l\binom jl
 \frac{\binom{o+n+r-j-1}{l}}{\binom{o+n+r+z-1}{l}}
 =\frac{(z+j)_{\underline j}}{(o+n+r+z-1)_{\underline j}}.
\]
The last identity is inclusion-exclusion for choosing \(j\) elements from a set of size \(o+n+r+z-1\) while avoiding a specified subset of size \(o+n+r-j-1\).

For a monomial of exact order \(e\), its substituted polynomial is \(E^e V_n\), where \(V_n=\eta^{n-S}V\), \(V\) is independent of \(n\), and \(V\) contains no \(W_\rho\). The largest possible \(\eta\)-degree in \(V\) is \(S-e\). Its coefficient there is the original monomial divided by \(W_\rho^e\). Thus the squared norm of \(V_n\), divided by \((n-S)!\,n^{S-e}\), tends to the factorial of this latter monomial. The same normalized squared norm of \(FV_n\) stays bounded: the \(W_\rho\)-derivative vanishes and every other \(Y\)-derivative acts on a bounded degree. On every nonharmonic component of \(V_n\) the positive eigenvalue of \(EF\) is at least \(n-o\). Its total normalized squared norm is consequently \(O(1/n)\).

On harmonic level \(j\) of the reduced bidegree, the eigenvalue of \(F^eT_z E^e\) is
\[
 (j+e)_{\underline e}
 (o+n+r-e-j-1)_{\underline e}
 \frac{(z+j+e)_{\underline{j+e}}}
      {(o+n+r+z-1)_{\underline{j+e}}}.
\]
It tends to \((e!)^2\binom{z+e}{e}\) for \(j=0\) and is \(O(n^{-j})\) for \(j>0\). The preceding norm estimate and
\(\binom{o+n+r+z-1}{o}/n^o\longrightarrow1/o!\)
now give, for equal orders \(e\), the coefficient at degree \(S+o-e\):
\[
 \frac{p^u e!\binom{z+e}{e}}{o!\,S!}\,[g_b,g_{b'}]_0 .
\]
For distinct monomials the corresponding factorial pairing is zero. More generally, in the norm defined by \(T_z\), a monomial of order \(e\), after this normalization, has all its limiting mass in harmonic level \(e\) of the full bidegree. For two different orders, orthogonality and Cauchy--Schwarz therefore make the mixed degree strictly less than
\(\big((S+o-e)+(S+o-e')\big)/2\). These are coefficient identities and vanishings of rational polynomials for pairs of monomials. Bilinearity transfers them to arbitrary \(\Q_p\)-coefficients. In a basis vector of order \(e\), higher-order monomials do not affect the coefficient just calculated.

In the determinant expansion the sum of the half-degree bounds along any permutation equals \(s_A\). Equality requires preservation of each order block. The coefficient at \(s_A\) is therefore the product of the determinants of the displayed blocks. It is nonzero under the hypothesis.
\end{proof}

\subsection*{Proof of \ref{B5.20r}}\label{prf:B5.20r}\phantomsection
\begin{proof}
Evaluation at \(Y_\rho=1,\ Y_c=0\) for \(c\ne\rho\) recovers \(f\). The upper root operators with both indices different from \(\rho\) act on \(f[n]\) by acting on \(f\) before substitution. Those with second index \(\rho\) annihilate the substituted generators and \(Y_\rho\), by the vector and covector substitution formulas. Hence \(f[n]\) is a nonzero highest vector of weight \(\lambda=w-n\epsilon_\rho\). Put \(a=n'-n\). The vector \(Y_\rho^a f[n]\) equals \(f[n']\) and has weight \(\lambda'=\lambda-a\epsilon_\rho\).

In the space of products of the prefix module with degree-\(a\) polynomials in \(Y\), the largest possible weight is \(\lambda'\), with a one-dimensional weight space. Indeed each factor has a one-dimensional highest weight space, and every other weight is lower in the dominance order. Let \(\mathcal C=\sum_{i,j}E_{ij}E_{ji}\) for the prefix including \(Y\). On the irreducible of highest weight \(\nu\) its scalar is \(c(\nu)=\sum_j\nu_j(\nu_j+r+1-2j)\), by applying the commutators to a highest vector. If \(\nu<\lambda'\) has the same total weight, summation by parts against the strictly decreasing sequence
\(\lambda'_j+\nu_j+r+1-2j\) gives \(c(\nu)<c(\lambda')\). The polynomial \(P(\mathcal C)=\prod_{\nu<\lambda'}(\mathcal C-c(\nu))\), with one factor for each other irreducible type present, maps into the module generated by \(f[n']\) and is nonzero on its highest line.

Use the map \(v\mapsto P(\mathcal C)(Y_\rho^a v)\). It commutes with all raising operators and shifts every weight by \(-a\epsilon_\rho\). A nonzero kernel on the irreducible prefix module would contain a highest vector: take a maximal weight in the kernel. The unique highest line maps nontrivially. Extend the injection coefficientwise on the disjoint suffix monomials. For a cut before \(Y\), the new highest vector satisfies its threshold by the substitution, which is equivariant on those letters, and the prefix action preserves that cut space. For a cut including \(Y\), multiplication by a covector polynomial preserves an upper first-coordinate bound: every weight in its orbit has first coordinate at most zero. The prefix Casimir preserves such a cut space because it commutes with the action on the whole larger prefix and preserves its weight spaces. The image thus lies in \(Z_{n'}\).
\end{proof}

\subsection*{Proof of \ref{B5.20s}}\label{prf:B5.20s}\phantomsection
\begin{proof}
Use polynomial representatives as in \ref{B3.52b}, subtracting the last entry from every entry before forming the column degrees. Put
\[
 \beta_j=c+\lambda_{j+1},\quad \alpha_j=Q+\beta_j,\quad
 N=Q+2c,\quad \alpha'_j=N-\alpha_{h+1-j}\qquad(1\le j\le h).
\]
The middle coefficient spaces have shapes \(\alpha,\alpha'\). We need a pairing \(\ell\) with
\[
 \ell(gv,gt)=\det(g)^N\ell(v,t).
\]
Here is its integral normalization. Put \(e=\alpha_h,\ \nu=\alpha-e\mathbf1\). Divide the first input by \(\det^e\). For the second, on a generic \(h\)-square matrix \(Z\) substitute
\(Z_{ij}\mapsto\operatorname{adj}(Z)_{h+1-j,i}\), then divide by \(\det(Z)^{|\alpha'|-(N-e)}\).
Jacobi's complementary minor identity sends an initial minor of height \(d\) to the complementary initial minor times \(\pm\det^{d-1}\), with the indicated reversal of columns. The full product has \(N-e\) columns, so the displayed power divides it over \(O\). Complementing the columns again recovers every generator of shape \(\nu\), up to sign and the prescribed full determinant columns. Thus both maps onto shape \(\nu\) are integral surjections. Under \(g\) in the second input the complement substitution transforms by \(\det(g)^{N-e}\) and inverse transpose. The first input contributes \(\det(g)^e\).

On shape \(\nu\) take the contravariant form normalized to have value one on the highest vector. Its dual of the dual Weyl lattice is the Weyl lattice: transposing a divided root operator interchanges its two indices, and evaluation on the highest line is the integral highest coefficient. Equivalently, the pairings with all divided-power words are the coefficients of the successive elementary row substitutions. The integral dual Weyl lattice consists of the sections with all those coefficients integral, by the tableau basis in \ref{B3.52b}. If \(d=\binom h2\lfloor\log_p(\nu_1+h-1)\rfloor\), the inclusions in \ref{B3.52b} give an integral form after multiplication by \(p^d\), whose image in the dual contains \(p^d\) times that dual. The hypotheses there hold since \(h\ge5\) and \(\nu_1\le2c\). Smith reduction and rescaling a partner by a unit and, when necessary, by a power of the uniformizer give full-rank pairs \(u_b,u'_a\) for \(\ell\) with
\(\ell(u_b,u'_a)=\delta_{ab}p^{k_b}\), \(k_b\le d\). For a Smith entry of half-integral valuation use the ceiling of its valuation. Since \(d\) is integral this preserves the bound. Lift the pairs through the two surjections.
\leanref{Formalism.\allowbreak{}MainNotions.\allowbreak{}flipOneThree.\allowbreak{}Pf.\allowbreak{}tauP}{Formalism.\allowbreak{}SharpRestriction.\allowbreak{}GrowthBounds.\allowbreak{}HermitianTensor.\allowbreak{}TauP}{303--482}
\leanref{Formalism.\allowbreak{}MainNotions.\allowbreak{}flipOneThree.\allowbreak{}Pf.\allowbreak{}pb\_\allowbreak{}smith\_\allowbreak{}lam}{Formalism.\allowbreak{}SharpRestriction.\allowbreak{}GrowthBounds.\allowbreak{}HermitianTensor.\allowbreak{}PbSmithLam}{27--69}

On a generic \(k\)-square matrix \(X\), let \(D(Z)=\operatorname{diag}(1,Z,1)\). Define a bilinear map on sections by
\[
 \mathcal B(v,t)(X)=
 \ell_{Z,Z'}\bigl(v(XD(Z)),t(XD(Z'))\bigr).
\]
The subscript means to pair coefficients of the two polynomials in \(X\) and convolve their \(X\)-exponents. The coefficients in \(Z,Z'\) have the middle column degrees of \(v,t\), respectively. Addition of an earlier middle column to a later one preserves these coefficients by the defining column invariance. They are therefore integral elements of the two middle lattices whenever \(v,t\) are integral. Left multiplication on \(X\) commutes with both substitutions. The determinant covariance of \(\ell\) makes the middle degrees of the output all equal to \(N\). Adding the first column to a middle column amounts, in either input, to adding multiples of its first column to its middle columns. The last column has degree zero. These checks give the defining column invariance and degrees of the upper flat space when \(v,t\) have weights \(\lambda,\delta^\vee\), with \(\delta^\vee_j=-\delta_{k+1-j}\).
\leanref{Formalism.\allowbreak{}MainNotions.\allowbreak{}flipOneThree.\allowbreak{}Pf.\allowbreak{}lamOf}{Formalism.\allowbreak{}SharpRestriction.\allowbreak{}GrowthBounds.\allowbreak{}HermitianTensor.\allowbreak{}CoeffSumMonomialLinear}{57--59}
\leanref{Formalism.\allowbreak{}MainNotions.\allowbreak{}flipOneThree.\allowbreak{}Pf.\allowbreak{}lamMemP}{Formalism.\allowbreak{}SharpRestriction.\allowbreak{}GrowthBounds.\allowbreak{}HermitianTensor.\allowbreak{}SubCoeffP}{10--57}

Write \(u_b=\det^Q v_b\), which is possible integrally since every tableau of shape \(\alpha\) has \(Q\) full columns available for removal. For each middle initial minor in \(v_b\) adjoin the first row and column, shifting the middle row indices up by one. If the product has \(d_0\) columns, multiply by \(X_{11}^{2c-d_0}\). This exponent is nonnegative since \(\beta_1\le2c\). Extend linearly and denote the lift by \(s_b\). On \(X=D(Z)\) it restricts to \(v_b(Z)\). The same construction on \(u'_a\), whose first row has length at most \(2c\), gives \(t_a\) restricting to \(u'_a\). These have weights \(\delta,\delta^\vee\); their first-row degrees are \(2c\) and their last-row degrees zero. Set \(U_b F=Fs_b\), \(L_a v=\mathcal B(v,t_a)\).
\leanref{Formalism.\allowbreak{}MainNotions.\allowbreak{}flipOneThree.\allowbreak{}Pf.\allowbreak{}liftP}{Formalism.\allowbreak{}SharpRestriction.\allowbreak{}GrowthBounds.\allowbreak{}HermitianTensor.\allowbreak{}InitialMinorLiftFirstEntry}{284--315}

For flat \(F\), \(F(XD(Z))=\det(Z)^Q F(X)\): the height-one columns see the first column, and each height-\((k-1)\) minor sees all the middle columns. Thus \(L_a U_bF=F H_{ab}\) with \(H_{ab}\) independent of \(F\). To identify \(H_{ab}\), take \(F\) to be the \(Q\)-th power of the initial principal \((k-1)\)-minor. The output has the maximum possible degree in the first row and degree zero in the last. In the two-column-height tableau description of a flat space this line is spanned by the highest tableau. Indeed all height-one factors must use row one, and all height-\((k-1)\) factors must omit row \(k\). Evaluation at \(X=I\) gives its scalar
\(\ell(\det^Qv_b,u'_a)=\delta_{ab}p^{k_b}\).
Cancel the nonzero initial minor power. The remaining highest tableau is the one of flat weight \([2c,2c]\), denoted by \(z_0\). Hence \(iF=z_0F\) has the asserted composition identity on every \(F\).

In an adapted frame, the normalized \(z_0\) is primitive since its tableau polynomial has a unit diagonal coefficient. If \(z_0F\) belongs to the integral upper lattice, choose the least \(j\ge0\) with \(\omega^jF\) integral in the lower frame. For \(j>0\), reduction of \(z_0(\omega^jF)\) is zero. Both factors have nonzero reduction as polynomials, a contradiction. In using coefficientwise reduction here, the integral tableau lattices are saturated in the polynomial coefficient lattices by their unit leading monomials. Undoing the common determinant twists gives saturation in the selected edge lattice. The resulting retraction is required in this one lattice.
\end{proof}

\subsection*{Proof of \ref{B5.20w}}\label{prf:B5.20w}\phantomsection
\begin{proof}
Fix a tie \(i\), and let \(j\) be the length of its prefix including the covector. The highest vectors on this prefix are the polynomials \(f[n_i]\) of \ref{B5.20q}. Indeed, evaluation at the last coordinate covector recovers \(f\). Conversely, on the chart \(Y_r\ne0\), the upper elementary substitutions with indices \(a,r\) take a covector to \(Y_re_r^\vee\). Invariance under those substitutions and homogeneity in \(Y\) recover exactly the formula in \ref{B5.20q}. Evaluation preserves integrality, and the inverse formula is integral since \(n_i\) is at least the preceding degree. A cap on the earlier letters is preserved by each simultaneous change of their frame, and evaluation recovers the same cap. Thus this is also an identification of the integral highest-vector lattices, weight by weight. The new highest weight is \(w-n_i\epsilon_r\).

Use left and right integral Smith bases, with substituted highest vectors \(h_a,h'_a\), and diagonal entries \(d_a\). Write \(B_{\mathrm F}\) for the unnormalized factorial pairing on the \(j\) letters. With \(o\) the marked vector degree and \(S\) the preceding total degree, the common scalar is
\[
 c_i=\frac{p^{\lceil\kappa\rceil}\binom{o+n_i+r+z_i-1}{o}}
                 {(n_i-S)!\,S!\prod_{\nu<j}d_\nu!},
\]
where \(d_\nu\) here denotes a letter degree. Thus \(c_iB_{\mathrm F}(\Phi_i h_a,h'_b)=\delta_{ab}d_a\) within each weight. Also \(c_iB_{\mathrm F}(h_a,h_b)\) is integral. In a monomial pairing the normalized factorial denominators are row multinomials. By the factorial formula of \ref{B3.19} their total valuation is at most \(jr\lfloor\log_p X\rfloor\): at each power \(p^\alpha\le X\) the difference of floors in one row is at most \(r-1\). Replacing \(n_i!\) by \((n_i-S)!S!\) contributes the integer \(\binom{n_i}{S}\). The other binomial in \(c_i\) is integral.

Define \(\pi_i h_a=h_a\) if \(v_p(d_a)\le E\), and zero otherwise. On highest vectors define a second map by
\[
 A_i h_b=
 \sum_{\substack{a:\ {\rm wt}(h_a)={\rm wt}(h_b)\\v_p(d_a)\le E}}
       \frac{p^{\lfloor E\rfloor}}{d_a}\,
       c_iB_{\mathrm F}(h_a,h_b)\,h'_a .
\]
These coefficients are integral. Extend both maps equivariantly on the rational prefix space satisfying its caps. This extension exists because in each isotypic component the highest vectors are its multiplicity space. Precompose with the commuting type projections for the caps on the prefix, and act coefficientwise on the suffix monomials. On \(Y\) these preliminary projections fix every prefix coefficient. Each extended map preserves \(Y\): it takes its values in the earlier cap spaces, and for any later cut it commutes with that cut's group action and preserves weights. In particular it carries the largest stable space with the stated weight bound into itself.

There is a uniform denominator for \(A_i\). Projection onto the prefix cap space costs at most \(ja_0\), by \ref{B5.20t} applied to the distinct positive prefix lengths. It commutes with the action on the \(j\) letters, so takes a highest vector to a highest vector. The map just defined takes integral highest vectors in the cap space to integral polynomials. On an arbitrary integral prefix coefficient, project first onto each type, at cost \(a_0\), and apply the divided-root generation assertion of \ref{B5.20t}, at cost another \(a_0\). The map is equivariant, and divided roots preserve integral coefficients. This proves
\[
 p^{K_i}A_i M\subseteq M,\qquad K_i=(j+2)a_0\le\kappa .
\]
All applications have the required weight bound since a coordinate on any prefix ranges from minus its covector degree to its vector degree.

The Smith identities on highest vectors give
\[
 B_0(\Phi_i v,A_i v')=p^{\lfloor E\rfloor}B_0(\pi_i v,v')
\]
for \(v\) satisfying the caps through tie \(i\) and \(v'\in Y\). To extend the identity, move each root operator on one argument to its transpose on the other. The two forms are contravariant and agree on the highest multiplicity spaces, so agree on the modules they generate. Distinct suffix monomials are orthogonal and matching monomials multiply both sides by the same factorial. Also \(\pi_i\) commutes with every later \(\Phi_j\), since the later pair is disjoint from its prefix.

In \(Y\), the codimension of the fixed space of \(\pi_i\) is \(\sum_w b_{iw}z_{iw}\). To see this, split the rational prefix module into the copies generated by the \(h_a\). Projection onto any one copy extends coefficientwise and preserves the later caps by equivariance, as above. Its image on \(Y\) is the intersection of those caps with the product of the copy and the suffix polynomial space. Copies of the same type have equal image dimension: the isomorphism carrying one highest vector to the other and its inverse both extend coefficientwise and preserve every later cap. Their common dimension is \(z_{iw}\). This argument also gives \(\dim Y=\sum_w(\operatorname{rank}M_w)z_{iw}\) for each fixed tie.

Let \(W\) be the intersection in \(Y\) of all these fixed spaces. Its codimension is at most the sum of their codimensions. The lattice \(W\cap M\) is saturated in \(M\). A primitive \(v\) in it has a unit polynomial coefficient, since division by \(p\) preserves every rational condition. Pair \(v\) with the monomial of that coefficient. The normalized factorial on this monomial is an inverse product of multinomials, so the pairing has valuation at most zero. Project the monomial onto all the caps and multiply by \(p^{K_0}\), \(K_0=(L+1)a_0\). This yields \(y\in M\) with \(v_p B_0(v,y)\le K_0\). The projections are self-adjoint and fix \(v\).

Take \(y'=(p^{K_1}A_1)\cdots(p^{K_t}A_t)y\in M\). Applying the pairing identity in order \(1,\ldots,t\) gives
\[
 B(v,y')=p^{\sum_i(K_i+\lfloor E\rfloor)}B_0(v,y).
\]
At step \(i\), the later tie operators preserve the earlier prefix caps and commute with \(\pi_i\), which fixes \(v\). Hence the valuation is at most \(\lfloor t(E+\kappa)+\kappa\rfloor\). If the number of Smith entries above this threshold exceeded \(\operatorname{codim}_Y W\), their rational left Smith span would meet \(W\) nontrivially. Scale an element of the intersection to be primitive in \(M\). Every partner in \(M\) would then have pairing of valuation above the threshold, contrary to the construction. This proves the bound.
\end{proof}

\subsection*{Proof of \ref{B5.20x}}\label{prf:B5.20x}\phantomsection
\begin{proof}
Let \((u_a)\) represent the residue line over \(\F_{p^{n'}}\). Modulo the uniformizer, the edge substitution at \(s_j\) sends \(X_{ja}\) to \(u_a^{p^{s_j}}\) times the surviving coordinate of that row. For \(u\ne0\),
\[
 \prod_j(u^{p^{s_j}})^{h_j}
 =\prod_j\frac{u^{p^{s_j+G_j}}}{u^{p^{s_j}}}=1,
\]
since \(u^{p^{n'}}=u\) at the wraparound. If \(u_0=0\), the first factor of \(R\) vanishes. Otherwise a competing line has \(u_a\ne0\) for some \(a>0\). Its \(a\)-th difference vanishes by the product identity. At the distinguished line the nonzero coordinate is normalized to \(u_0=1\), which gives the indicated monomial with coefficient one. Taking the \(c\)-th power gives the divisibility and the exact multidegrees.
\end{proof}

\subsection*{Proof of \ref{B5.20ia}}\label{prf:B5.20ia}\phantomsection
\begin{proof}
There are at most \(k\) digits in \(D_i\), so \(D_i\le kQ\). Since successive odd distances differ by at least two,
\[
 P_i C_i\le a_1/p^{d_i}
       +2b'\sum_{\nu<i}p^{d_\nu-d_i}
 \le a_1/p^{d_i}+2b'/p .
\]
In the first range insert this inequality in the bound on \(e_i\), using \(P_iH_i\le b_i\le2b'\). The excess of \(P_i e_i\) over the minimum is at most
\[
 8b'/p^2+L'a_1/(pk)+2L'b'/(pk).
\]
In the other range the minimum is at least \(a_1/p^{d_i}\). The excess is at most \(2b'/p+2L'b'/k\). Thus in either range an upper bound is \(4b'/p+L'a_1/(pk)+4L'b'/k\).

Finally \(C_s\le a_1/p^2+2s b'/p\), since \(G\ge2\) and \(P_i\ge p\). The terminal contribution is at most
\[
 a_1/p^2+2s b'/p+2s L'a_1/(pk)+8s a_1/p^3+2L'b'/k .
\]
Adding the \(s\) interior bounds and their digits proves the inequality coefficient by coefficient, using \(p\ge5\).
\end{proof}

\subsection*{Proof of \ref{B5.20ib}}\label{prf:B5.20ib}\phantomsection
\begin{proof}
By the polynomial calculation in \ref{B5.20q}, each entry has degree at most \(S+o\le2S\). Evaluate it at \(n=S,S+1,\ldots,S+2S\). The two input polynomials \(f_b[n]\) and \(f_{b'}[n]\) have integral coefficients, since the substitutions in \ref{B5.20q} have integral coefficients and \(n\ge S\).

For the term of contraction order \(l\), set \(D=o+n+r+z-1\). Cancellation gives
\[
 \frac{\binom D o}{\binom D l}
   =\frac{\binom{D-l}{o-l}}{\binom ol}.
\]
The denominator costs at most \(\ell_X\) in valuation, by \ref{B3.19} with \(l\le o\le S\). The size of \(D\) does not enter this estimate.

On a monomial with tied exponents \(\alpha,\beta\in\N^r\), the contribution of a multi-index \(a\) with \(|a|=l\) to \(F^l/(l!)^2\) is zero unless \(a\le\alpha,\beta\). Otherwise its coefficient is
\[
 \frac{\prod_{h=1}^r\binom{\alpha_h}{a_h}\binom{\beta_h}{a_h}}
      {\binom{l}{a_1,\ldots,a_r}}.
\]
This follows by expanding the power of the sum of the commuting contractions and writing each falling factorial as \(a_h!\) times a binomial. The denominator is a product of at most \(r-1\) binomials with upper argument at most \(l\). Its cost is at most \(r\ell_X\). Multiplication by \(E^l\) has integral coefficients and restores the row degrees.

There are now \(m+2\) rows, of degrees \(d_1,\ldots,d_{m+1},n\). On two matching monomials the factorial pairing divided by \(n!\prod_i d_i!\) is the product of the inverses of their row multinomials. Every row degree is at most \(3S\le pX\). By successive binomial choices and \ref{B3.19}, each row costs at most \(r\ell_X\). Replacing \(n!\) by \((n-S)!S!\) multiplies the value by \(\binom nS\). The total denominator cost is at most \(((m+2)r+r+1)\ell_X\), so the scaled entry is integral at every evaluation point.

If \(a(n)\) is this entry, its coefficient of \(\binom{n-S}{j}\) is \(\Delta^j a(S)\). This is an integral linear combination of the first \(j+1\) values just computed. Finally
\[
 \binom{n-S}{j}=\sum_{h=0}^j\binom{-S}{j-h}\binom nh ,
\]
and \(\binom{-S}{j-h}\in\Z\).
\end{proof}

\subsection*{Proof of \ref{B5.20ic}}\label{prf:B5.20ic}\phantomsection
\begin{proof}
For one variable \(x\), the operator of index \(n\) takes the coefficient of \(t^n\) after \(x\mapsto(1+t)x\). On the reduced word \(y\) this is \(x^n\partial_x^{[n]}\), where the divided derivative has monomial coefficient \(\binom{\mu_x}{n}\). At a prefix containing that letter, every conjugate of the operator is a sum of products of \(n\) coordinate multiplications and \(n\) derivatives on the same letter. The signed first-coordinate weight increases by at most \(n\). Indeed on a vector letter each multiplication has weight zero or one and each derivative has weight zero or minus one. On a covector letter the signs reverse. A prefix omitting the letter commutes with the operation. Applying this argument in every basis gives the stable prefix bounds.

If \(x\) occurs in an internal tie polynomial, write \(H_h=A+B\), where \(A\) is its unique summand containing \(x\). The coefficient of \(t^l\) in the transformed \(H_h^{j_h}\) is
\[
 \binom{j_h}{l}H_h^{j_h-l}A^l .
\]
For an ordinary coordinate \(A=W_{ha}Y_{ha}\). Multiplying the reduced word by \(A^l\) adds \(l\) to the two reduced degrees and uses at most \(l\) in any prefix bound. Under any prefix change of basis each vector-covector product has greatest first-coordinate weight at most one, including when only its vector occurs in the prefix. For a scalar coordinate \(A=Z_h S'_h=H_h-E_h\). Expanding \(A^l\) gives terms with \(d\le l\) factors of \(E_h\), level \(j_h-d\), and cost at most \(d\) by the same argument. The factor \(Z_h^{z_h}\) contributes integral binomial coefficients. The other ties use disjoint variables. In extracting \(t^n\), the orders falling on all factors sum to \(n\). This proves the one-coordinate assertion, and composition in the different coordinates proves the displayed inclusion.

For saturation, express a polynomial in the field span as
\(\sum_{j,\alpha}(\prod_h H_h^{j_h}Z_h^{z_h})y_{j,\alpha}\xi^\alpha\), collecting equal indices. Set all \(Z_h=1\), then substitute \(S'_h-E_h\) for \(S'_h\) at every tie. These substitutions preserve integral coefficients and take \(H_h\) to \(S'_h\), also at the terminal tie where \(E_h\) involves \(\xi\). The coefficient of \(\prod_h(S'_h)^{j_h}\xi^\alpha\) is \(y_{j,\alpha}\). Hence all these reduced polynomials are integral. They already satisfy the prefix and degree conditions in the field span, so the displayed expression belongs to \(L(b)\).
\end{proof}

\subsection*{Proof of \ref{B5.20i}}\label{prf:B5.20i}\phantomsection
\begin{proof}
Put \(r=k-1\), \(\ell=1+\log_p(8a_1+r)\), \(Q_p=p^{G'}\),
and \(q_p'=p^{n'}\). The number of edges of rank one into
a target is \(E_q=((q_p')^k-1)/(q_p'-1)\), by counting
nonzero residue forms modulo scalars. At a source count
residue lines. For flat weight \((A_s,0,\ldots,0,-B_s)\) at the
marks and zero elsewhere, let \(f(A,B)\) be the least nonnegative
loss allowing \((1-f(A,B))\dim V/E_q\) pairs on each edge at
depth \(t_*\), order \(c_4\), and common diagonal valuation at
most \(D_1\). Index the budgets by the gaps between successive
even marks. Set
\[
\begin{gathered}
 \kappa_0=2Nk\log_p a_1,\quad
 D_j=C_A Y^2N_1^2 k^2\ell+2\kappa_0,\quad
 D_0=\sum_jD_j+4\kappa_0,\quad
 u_0=\lceil(k+1)\ell\rceil,\quad D_1=D_0+4Nu_0,\\
 t_0(A,B)=J_{\rm odd}(A,B)+C_tN
 (a_1/p^2+b'/p+L'a_1/(pk)+L'b'/k+kQ_p+Q_pc_4),\qquad
 t_*=\lceil t_0(A,B)+NJ_0\rceil_+ .
\end{gathered}\tag{5.20.19a}\label{eq:p5c2-rank-one-budgets}
\]
Here \(J_{\rm odd}(A,B)\) uses the \(A\)-weight at the left even
endpoint and the \(B\)-weights at interior odd marks.
\leanref{Formalism.\allowbreak{}MainNotions.\allowbreak{}flipOneThree.\allowbreak{}Pf.\allowbreak{}Dat.\allowbreak{}kF}{Formalism.\allowbreak{}SharpRestriction.\allowbreak{}GrowthBounds.\allowbreak{}ColumnSpanDensity.\allowbreak{}RhoLowerBoundLargeConductor}{138--140}

On the tensor of the two full symmetric powers at each mark let
\(g=\kappa_P+\sum_j\kappa_{T,j}\), where \(\kappa_{T,j}\) is the fraction
of rank in the direct sum of the kept block lattices on segment \(j\)
lost by retaining Smith factors of valuation at most \(D_j\). Use
loss zero for rank zero. Define
\[
 \kappa_P=C_\kappa N\left(e^{-L'/8}+
 \frac{kQ_p(m_0+N_1k)}{b'}+\frac1Y+\frac{kE_0}{a_1}\right),
 \tag{5.20.19b}\label{eq:p5c2-rank-one-margin}
\]
where \(m_0=C_m kN_1Q_pc_4+12c_4\) and
\[
 E_0=C_{E1}(a_1/p^2+b'/p+L'N_1a_1/(pk)+L'b'/k)
 +C_{E2}k^2N_1Q_p^2c_4 .
\]
The constants here are absolute. Equation~\eqref{eq:p5c2-rank-one-size}
gives \(e^{-L'/8}\le1/(NY)\),
\(k^3N_1Q_p^2c_4/b'\le2N_1/(HNY)\), and
\(k^2N_1Q_p/b'\le2N_1/(HNY)\). The other contributions of \(kE_0/a_1\)
are bounded by a constant times \(2k/p^2+L'(N_1+1)/p\), which is
\(O(N_1/Y)\). Thus \(\kappa_P\le C_sNN_1/Y\).
\leanref{Formalism.\allowbreak{}MainNotions.\allowbreak{}flipOneThree.\allowbreak{}Pf.\allowbreak{}Dat.\allowbreak{}kx}{Formalism.\allowbreak{}SharpRestriction.\allowbreak{}GrowthBounds.\allowbreak{}ColumnSpanDensity.\allowbreak{}RhoLowerBoundLargeConductor}{142--144}
\leanref{Formalism.\allowbreak{}MainNotions.\allowbreak{}G.\allowbreak{}kappaMarginBound}{Formalism.\allowbreak{}SharpRestriction.\allowbreak{}GrowthBounds.\allowbreak{}ObjSeg}{434--435}

Consider one gap of length \(G\), with inner distances \(d_i\)
for \(0\le i<s\), weights \(a_i,b_i\), and right endpoint
weights \(a_+,b_+\). Set \(P_i=p^{G-d_i}\) and
\(\sigma=\#\{i:p^{d_i}<a_1/b'\}\). In each row of \(k\) exponents
fix digits \(u_{ic},v_{ic}\in[0,P_i)\) with
\[
 H_i=(b_i-\sum_cu_{ic})/P_i\in\N,\qquad
 K_i=(a_i-\sum_cv_{ic})/P_i\in\N.
\]
The blocks in this digit class have nonnegative integer
parameters \(o_i,w_i,h_i,o_+\). Their signed capacity is given by
\[
 C'_0=\left\lfloor\frac{a_1-E_0}{p^G}\right\rfloor\in\Z,\quad
 C_0^{\rm cap}=C'_0,\quad C_{i+1}^{\rm cap}=C_i^{\rm cap}+H_i-e_i,\qquad
 e_i=\begin{cases}w_i+h_i&i<\sigma,\\ o_i&i\ge\sigma.\end{cases}
\]
Write \(C_i^{\rm cap}=C_i\) in the segment formulas. For \(i<\sigma\)
keep \(o_i=0,\ w_i\le H_i,\ h_i\le K_i\) and
\[
 w_i+h_i\le(1+4/p^2)H_i+L'C_i/k,\qquad h_i\le2L'C'_0/k .
\]
For \(i\ge\sigma\) keep \(w_i=h_i=0,\ o_i\le H_i\) and
\(o_i\le C_i+L'H_i/k\). At the terminal letter keep
\(o_+\le C_s+(L'/k+4/p^2)\sum_i2a_1/P_i+L'b_+/k\).
The \(r\)-variable letters alternate between vectors and covectors,
starting with a vector. The degrees are
\((w_i,K_i-h_i)\) for \(i<\sigma\) and \((o_i,K_i)\) for \(i\ge\sigma\),
followed by \(o_+\). Adjoin a covector \(\xi\) of degree \(b_+\).
\leanref{Formalism.\allowbreak{}MainNotions.\allowbreak{}G.\allowbreak{}segData}{Formalism.\allowbreak{}SharpRestriction.\allowbreak{}GrowthBounds.\allowbreak{}ObjSeg}{319--331}
\leanref{Formalism.\allowbreak{}MainNotions.\allowbreak{}G.\allowbreak{}Kept}{Formalism.\allowbreak{}SharpRestriction.\allowbreak{}GrowthBounds.\allowbreak{}ObjSeg}{161--173}

Internal ties pair the vector and covector at \(i\ge\sigma\) with shift
\(z_i=H_i-o_i\). The terminal tie pairs the final vector and \(\xi\)
with shift \(z_+=a_+-o_+\). A level chooses a count \(j_i\) at each
internal tie and \(j'\) at the terminal tie, at most the lesser
degree in its pair, and subtracts it from both degrees there.
Let \(Y_{\rm pol}\) be the subspace on the first \(2s+1\) reduced
letters cut out by the prefix bounds
\(\lfloor C_i-m_0\rfloor\) through the vector at \(i\ge\sigma\) and
\(\lfloor C_s-m_0\rfloor\) on the full prefix. At each prefix take
the largest stable subspace in the allowed signed first-coordinate
weight span. It is obtained by intersecting successively with
inverse images under the prefix operators
\(\sum_{\rm vec}X_{la}\partial_{lb}-\sum_{\rm covec}X_{lb}\partial_{la}\).
The level lattice is the sum of \(Y_{\Z}\xi^\mu\) over
\(|\mu|=b_+-j'\), where \(Y_{\Z}=Y_{\rm pol}\cap\Z_p[X]\).
\leanref{Formalism.\allowbreak{}MainNotions.\allowbreak{}G.\allowbreak{}blockData}{Formalism.\allowbreak{}SharpRestriction.\allowbreak{}GrowthBounds.\allowbreak{}ObjSeg}{207--210}

For a tied pair \(w,e\), let \(E\) multiply by \(\sum_cw_ce_c\)
and \(F=\sum_c\partial_{w_c}\partial_{e_c}\) be its adjoint for
\(\langle X^\alpha,X^\beta\rangle=\delta_{\alpha\beta}\alpha!\).
At reduced degrees \(o,n\) and shift \(z\) put
\[
 \Phi_z=\sum_{l=0}^{\min(o,n)}
 \frac{(-1)^lE^lF^l}{(l!)^2\binom{o+n+r+z-1}{l}}.
 \tag{5.20.19c}\label{eq:p5c2-tie-operator}
\]
With a fresh scalar pair \(Z,S'\) for the tie put
\(\widehat E=E+ZS'\), \(\widehat F=F+\partial_Z\partial_{S'}\).
The extended harmonic projector on \(Z^zy\) has this formula
with hats. Indeed \(\widehat F^l(Z^zy)=Z^zF^ly\) and
\([\widehat F,\widehat E]=r+1+\deg_{(w,Z)}+\deg_{(e,S')}\)
cancels consecutive terms under \(\widehat F\). Specializing
\(Z=1,S'=0\) gives \(\Phi_z\). Let \(T\) compose the internal
operators and \(d_F\) be the product of the reduced degree
factorials, including that of \(\xi\). The preterminal form
is \(B(u,v)=\langle Tu,v\rangle/d_F\). The internal ties are
disjoint, so the operators commute and \(B\) is symmetric.
\leanref{Formalism.\allowbreak{}GM.\allowbreak{}tiePhi}{Formalism.\allowbreak{}MahlerMeasure.\allowbreak{}LevelSum}{21--25}
\leanref{Formalism.\allowbreak{}GM.\allowbreak{}M.\allowbreak{}tieP}{Formalism.\allowbreak{}MahlerMeasure.\allowbreak{}TieVars}{18--22}

For a lattice with a form write \(g_D\) for the greatest
number of integral pairs with matrix \(p^DI\) and \(b_D\)
for the complementary rank. For integral pairs \(x_a,y_a\)
in \(Y_{\Z}\) and \(m=b_+-j'\), use
\(x_a\xi^\mu,(m!/\mu!)y_a\xi^\mu\) for every \(|\mu|=m\).
These are orthogonal across distinct \(\mu\) and preserve
the diagonal scalar. If \(d_\xi=\binom{m+r-1}{r-1}\) and
\(r_l=d_\xi\dim Y_{\rm pol}\), the bad rank before the
terminal operation is at most \(d_\xi b_D(Y_{\Z})\).
\leanref{Formalism.\allowbreak{}MainNotions.\allowbreak{}G.\allowbreak{}SymmetricFormBlocks.\allowbreak{}plusAW\_\allowbreak{}of\_\allowbreak{}blockOK}{Formalism.\allowbreak{}SharpRestriction.\allowbreak{}GrowthBounds.\allowbreak{}MonomialFactorials.\allowbreak{}PlusAWICondition}{63--146}

\paragraph{Internal Gram average.}
Fix digits, block, and level. The allowed \(a_i\) form one progression of step \(P_i\) per
coordinate, so the active covector degrees \(n_i=K_i-j_i\)
range over a product of integer intervals of lengths \(Y_i\). The block and level
indices are constant there. Indeed \(kP_i\le b'\),
\(b_i\le5b'/2\), and \(a_i\ge a_1>25b'\) imply \(H_i\le K_i\).
For \(i<\sigma\), \(P_i h_i\le a_1/2\) by the upper bound on
\(h_i\) and \(2L'/k\le1/2\). For \(i\ge\sigma\) use
\(j_i\le o_i\le H_i\). The terminal index bound uses \(o_+,b_+\)
alone.
\leanref{Formalism.\allowbreak{}MainNotions.\allowbreak{}G.\allowbreak{}SymmetricFormBlocks.\allowbreak{}idx\_\allowbreak{}sub}{Formalism.\allowbreak{}SharpRestriction.\allowbreak{}GrowthBounds.\allowbreak{}MonomialFactorials.\allowbreak{}DigitsOKOfCompat}{88--188}

Let \(K^{(i)}\) replace \(K_i\) by its largest compatible
value and let \(S_i\) be the degree of the prefix through its
vector letter at the upper endpoint of the grid. If \(n_{0i}\)
is the least reduced covector degree, then
\[
 P_i n_{0i}\ge a_1-kP_i-\tfrac52 b'.
\]
Here the digit sum is at most \(kP_i\), and \(P_i j_i\le b_i\).
The prefix and window bounds are
\[
 P_iS_i\le\tfrac52 b'+(2s+1)2a_1/p,\qquad
 P_iY_i\ge a_1/(2kY).
\]
For the first bound use \(H_i\) for the current vector and the
smaller divided degrees for earlier letters, since
\(P_i/P_h\le1/p\) for \(h<i\). The second follows by counting
the integers of the residue class in the upper weight interval
and allowing one incomplete step at each end. As \(kP_i\le b'\),
\(2s+1\le k\), \(p\ge k^2Y^2\), and \(kY^2b'\le a_1/2\), these
bounds give \(S_i\le n_{0i}\) and \(S_i/Y_i\le7/Y\). At the upper
endpoint, \(K_i\le2a_1/P_i\), \(o_i,w_i\le5b'/(2P_i)\), and
\(o_+\le C_s+(L'/k+4/p^2)\sum_i2a_1/P_i+L'b_+/k\).
Inserting the signed capacity and using \(P_i\ge p\) bounds
the total reduced degree by \(8a_1+r\). For each internal tie
the pair degree plus \(r+z_i+1\) is at most
\((2a_1+5b'/2)/P_i+r+1\le8a_1+r\).
Also \(Y_i\le 2a_1+1\) and \(2S_i+1\le8a_1+r\) by the prefix
bound. These inequalities hold throughout the grid by
monotonicity of the degrees.
\leanref{Formalism.\allowbreak{}CountingSumsFrames.\allowbreak{}TheoremMV7}{Formalism.\allowbreak{}PolynomialDimension.\allowbreak{}ChainOrderLetters}{14--36}
\leanref{Formalism.\allowbreak{}MainNotions.\allowbreak{}G.\allowbreak{}SymmetricFormBlocks.\allowbreak{}lower\_\allowbreak{}bound\_\allowbreak{}initial\_\allowbreak{}multiplicity}{Formalism.\allowbreak{}SharpRestriction.\allowbreak{}GrowthBounds.\allowbreak{}MonomialFactorials.\allowbreak{}LowerBoundInitialMultiplicity}{23--75}

Put \(\kappa_S=3(s+2)r^2\ell\), let \(t\) be the number of
internal ties, and set
\[
 D_{1r}=t(Y+C_M)\kappa_S+C_M\kappa_S,\qquad D'=\lfloor D_{1r}\rfloor
\]
with \(C_M=1\).
At tie \(i\) use the lattice in \ref{B5.20o} on the prefix ending immediately before the varying covector. The weight \(w\) is a full torus weight and is highest for the group on the first \(r-1\) coordinates. Fix all other degrees. This prefix lattice is independent of \(n_i\). Put \(S=\sum d_j\le S_i\), \(o=d_{m+1}\), \(z=z_i\), and \(u=\lceil\kappa_S\rceil\). The matrix for the window is the matrix of \ref{B5.20q}, including its binomial multiplier and its \((n_i-S)!\) normalization. Apply \ref{B5.20v} to this prefix, whose last letter is a vector. For \(m>0\) it gives \(c=e_1+2e_2\) and \(D=(m+1)e_1\) in \ref{B5.20p}, with \(e_1,e_2\) as in \ref{B5.20v}. Take for each partner the orthogonal projection of its monomial. For \(m=0\) take \(c=D=0\). With these bounds, \ref{B5.20p} gives \(v_p(\det H_e)\le d_e\kappa_4\), where \(\kappa_4=(s+4)r^2\ell\): both logarithms are at most \(\ell\) since \(2S+1\le8a_1+r\), \(p\ge3\), and \(m+7\le2(s+4)\). For the scalar \(\beta_e=p^u e!\binom{z+e}{e}/(o!S!)\) in \ref{B5.20q}, the factorial formula in \ref{B3.19} gives \[v_p(\beta_e)+\frac{S+o-e}{p-1}=u+v_p\binom{z+e}{e}+\frac{s_p(S)+s_p(o)-s_p(e)}{p-1}\le\kappa_S+3\ell,\] where \(s_p\) denotes digit sum. Indeed \(S,o,z+e\le8a_1+r\), the binomial costs at most \(\lfloor\log_p(8a_1+r)\rfloor\), and the ceiling of \(\kappa_S\) costs less than one. Multiplying along the order blocks proves \(v_p(b_A)+s_A/(p-1)\le d(\kappa_S+3\ell+\kappa_4)\le2d\kappa_S\). In particular the leading determinant is nonzero. Apply \ref{B5.20ib} with \(X=8a_1+r\). Here \(2S+1\le2S_i+1\le X\). There are at most \(2s+1\) reduced letters before the terminal covector, so \(m\le2s\) and \(((m+2)r+r+1)\ell\le3(s+2)r^2\ell=\kappa_S\le u\). Thus the matrix has integral binomial coefficients. In \ref{B5.20f} take \(H=Y_i\), \(J=2S_i\), and \(h=1+\lfloor\log_p X\rfloor\). The grid bounds give \(H,J+1\le X<p^h\), and \(h\le\ell\). With the leading-coefficient bound just proved, the Smith window estimate at
threshold \(Y\kappa_S\) thus gives the coefficient
\[
 \frac{2S_i}{Y_i}+
 \frac{2\kappa_S+2S_i\ell/Y_i}{Y\kappa_S}
 \le\frac{30+C_M}{Y}.
\]
Apply \ref{B5.20w} at each grid point with \(L=2s+1\), \(X=\lfloor8a_1+r\rfloor\), \(E=Y\kappa_S\), and \(\kappa=\kappa_S\). The caps end at the marked vectors or the final letter. The preceding degree bounds give \(n_i\ge n_{0i}\ge S_i\). The total degree is at most \(8a_1+r\), so \(S_{\rm tot}+r-1\le p(8a_1+r)\) and the exponent \(a_0\) in the lemma is at most \(r^2\ell/2\). Both denominator inequalities follow from \(\kappa_S=3(s+2)r^2\ell\). The threshold in the conclusion is exactly \(D'\). On a fibre in coordinate \(i\) fix the same prefix basis and the same nonzero representative for each nonzero weight space. By \ref{B5.20r}, its suffix dimension is at most its value at the top of the fibre. Multiply the single-variable count by that value and sum over weights. The dimension identity in the proof of \ref{B5.20w} gives \(\dim Y_{\rm pol}(K^{(i)})\) at this endpoint. The step in \(n_i\) is one because the step in \(a_i\) is \(P_i\). Summing over the other coordinates gives
\[
 \sum_{K\in\mathcal C} b_{D'}(Y_{\Z}(K))
 \le\frac{30+C_M}{Y}\sum_{\text{active }i}\sum_{K\in\mathcal C}
 \dim Y_{\rm pol}(K^{(i)})
 \tag{5.20.19f}\label{eq:p5c2-internal-mean}
\]
for the digit progression \(\mathcal C\).
\leanref{Formalism.\allowbreak{}InvariantFactorPaths.\allowbreak{}MC.\allowbreak{}theoremMXB\_\allowbreak{}of\_\allowbreak{}point}{Formalism.\allowbreak{}ThetaLogCoordinates.\allowbreak{}PartnerMeanCoordinates.\allowbreak{}PartnerMeanProofs.\allowbreak{}TheoremMXBOfPoint}{7--230}
\leanref{Formalism.\allowbreak{}InvariantFactorPaths.\allowbreak{}MC.\allowbreak{}theoremMV7}{Formalism.\allowbreak{}ThetaLogCoordinates.\allowbreak{}PartnerMeanCoordinates.\allowbreak{}PartnerMeanProofs.\allowbreak{}GramHypothesisOfLogBound}{188--250}
\leanref{Formalism.\allowbreak{}InvariantFactorPaths.\allowbreak{}MC.\allowbreak{}Top.\allowbreak{}theoremMV7}{Formalism.\allowbreak{}ThetaLogCoordinates.\allowbreak{}PartnerMeanCoordinates.\allowbreak{}LatticeTheoremEps28Valid}{7--8}
\leanref{Formalism.\allowbreak{}MainNotions.\allowbreak{}G.\allowbreak{}CountingStatementReduction.\allowbreak{}alternative\_\allowbreak{}mean\_\allowbreak{}value\_\allowbreak{}theorem}{Formalism.\allowbreak{}SharpRestriction.\allowbreak{}GrowthBounds.\allowbreak{}CountingStatementReduction}{19--23}
\leanref{Formalism.\allowbreak{}MainNotions.\allowbreak{}G.\allowbreak{}SymmetricFormBlocks.\allowbreak{}restTM\_\allowbreak{}closed}{Formalism.\allowbreak{}SharpRestriction.\allowbreak{}GrowthBounds.\allowbreak{}MonomialFactorials.\allowbreak{}EpsilonBoundMConstant}{228--236}

\paragraph{Terminal Gram average.}
Diagonalize the symmetric preterminal form over \(\Z_p\).
A smallest-valuation off-diagonal entry can be put on the
diagonal by summing the two basis vectors, since \(2\) is a unit.
Subtract integral multiples of the resulting vector to clear
its row and repeat. Diagonal entries of valuation \(\le D'\)
give \(g_{D'}\), by Smith reduction. For the terminal reduced
degrees \(o,n\) put \(J=\min(o,n)\), \(q_0=o+n+r-1\).
Multiplying the terminal form by
\(\sigma_0=p^{x_0}\binom{q_0+z}{J}\) gives the matrix entries
\[
 \frac{p^{x_0}}{d_F}\sum_{\lambda=0}^J(-1)^\lambda
 \binom{q_0+z-\lambda}{J-\lambda}
 \frac{\langle E^\lambda F^\lambda Tu,v\rangle}
 {(\lambda!)^2\binom J\lambda}.
 \tag{5.20.19d}\label{eq:p5c2-terminal-gram}
\]
This uses
\(\binom{x}{J}\binom J\lambda=\binom{x}{\lambda}\binom{x-\lambda}{J-\lambda}\).
With \(n_L=2s+1,\ a_{\rm big}=(3s+2)a_1\),
\(\ell_p=1+\lfloor\log_p(8a_{\rm big}+r)\rfloor\), take
\(x_0=((t+1)(r+1)+(n_L+1)r+2(t+1))\ell_p\).
In divided contractions the coordinate binomial factors are
integral and the multinomial denominator costs at most
\((r-1)\ell_p\). The tie and terminal binomial denominators
cost at most \(\ell_p\) each. Comparing coordinate factorials
to row factorials costs at most \(((n_L+1)r+2(t+1))\ell_p\).
These are the two contributions to \(x_0\). Vandermonde's
identity writes \eqref{eq:p5c2-terminal-gram} integrally in
\(\binom zi\), with coefficient \(p^{x_0}B\) at \(i=J\).
\leanref{Formalism.\allowbreak{}MainNotions.\allowbreak{}G.\allowbreak{}GramWindowAlt.\allowbreak{}phi\_\allowbreak{}pair}{Formalism.\allowbreak{}SharpRestriction.\allowbreak{}GrowthBounds.\allowbreak{}FockGramWindow.\allowbreak{}PhiPair}{18--46}
\leanref{Formalism.\allowbreak{}MainNotions.\allowbreak{}G.\allowbreak{}GramWindowAlt.\allowbreak{}winLvl}{Formalism.\allowbreak{}SharpRestriction.\allowbreak{}GrowthBounds.\allowbreak{}FockGramWindow.\allowbreak{}WinLvl}{77--97}

On the preceding good directions, the leading determinant
coefficient is the product of the scaled preceding entries
divided by \((J!)^d\) in rank \(d\). Since
\(J/(p-1)-v_p(J!)\le\ell_p\), apply \ref{B5.20f} with
\(L\le x_0+D'+\ell_p\) and \(E_*=\lfloor YD_{1r}\rfloor+1\).
The terminal window has \(H_a\ge a_1/(kY)\), while \(J\le5b'/2\).
Thus \(J/H_a\le1/(2Y)\). With \(h=\lceil\log_p(a_1+1)\rceil\),
\(H_a,J+1\le p^h\), and \(x_0+\ell_p,h\le\kappa_S\le D_{1r}\).
The fraction discarded on the good directions is at most
\(1/(2Y)+(2+1/(2Y))/Y\le5/Y\). For a retained entry \(d_i\),
scale its partner by \(p^{E_*-1}\sigma_0/d_i\), which is
integral and gives the terminal diagonal \(p^{E_*-1}\).
Thus
\[
 \frac1{H_a}\sum_{a_+}b_{\lfloor YD_{1r}\rfloor}(L,l,a_+)
 \le d_\xi b_{D'}(Y_{\Z})+5r_l/Y .
 \tag{5.20.19e}\label{eq:p5c2-terminal-loss}
\]
The first term bounds the entire rank if the preceding good
span has rank zero.
\leanref{Formalism.\allowbreak{}MainNotions.\allowbreak{}G.\allowbreak{}GramWindowAlt.\allowbreak{}restWinNum\_\allowbreak{}proved}{Formalism.\allowbreak{}SharpRestriction.\allowbreak{}GrowthBounds.\allowbreak{}FockGramWindow.\allowbreak{}InitialPointNumericBound}{183--235}
\leanref{Formalism.\allowbreak{}MainNotions.\allowbreak{}G.\allowbreak{}GramWindowAlt.\allowbreak{}bad\_\allowbreak{}le\_\allowbreak{}smith}{Formalism.\allowbreak{}SharpRestriction.\allowbreak{}GrowthBounds.\allowbreak{}FockGramWindow.\allowbreak{}GramWindow2Proved}{31--133}
\leanref{Formalism.\allowbreak{}MainNotions.\allowbreak{}G.\allowbreak{}SymmetricFormBlocks.\allowbreak{}restPlus\_\allowbreak{}of}{Formalism.\allowbreak{}SharpRestriction.\allowbreak{}GrowthBounds.\allowbreak{}MonomialFactorials.\allowbreak{}RestPlusOf}{53--60}

\paragraph{From levels to block lattices.}
At each level \(l\), the ordinary generators are
\(\prod_h\widehat E_h^{j_h}Z_h^{z_h}u\) with \(u\) in the level
lattice, including the terminal tie. Define \(\Theta_l u\) by
harmonically projecting each pair after multiplication by
\(Z_h^{z_h}\) and before applying \(\widehat E_h^{j_h}\).
With \(L_p=\lfloor\log_p(8a_{\rm big}+r)\rfloor\) and
\(U=(t+1)\lceil(r+1)(1+\log_p(8a_{\rm big}+r))\rceil\),
\(p^U\Theta_lu\) belongs to the integral span of ordinary
generators. The contraction coefficient for
\(F^{h'}/(h'!)^2\) at \(|\nu|=h'\) is
\[
 \frac{\prod_c\binom{\alpha_c}{\nu_c}\binom{\beta_c}{\nu_c}}
 {\binom{h'}{\nu_1,\ldots,\nu_r}} .
\]
The denominator costs at most \((r-1)L_p\) and the projector
binomial at most \(L_p\). A contraction lowers the degrees
and increases the level count after the compensating
multiplication. At a prefix containing both letters it
commutes with the action. At the prefix ending at the vector
letter each derivative lowers its signed first-coordinate
weight or leaves it fixed. The derivative span is stable by
the commutator with \(X_a\partial_b\). Thus the new
generators satisfy the prefix constraints of their levels.
\leanref{Formalism.\allowbreak{}MainNotions.\allowbreak{}G.\allowbreak{}SymmetricFormBlocks.\allowbreak{}restLL\_\allowbreak{}proved}{Formalism.\allowbreak{}SharpRestriction.\allowbreak{}GrowthBounds.\allowbreak{}MonomialFactorials.\allowbreak{}BlockValuationZpOfValuationBound}{190--195}

The images of different levels are orthogonal. Move a power
of \(\widehat E\) at a differing tie across the pairing as
\(\widehat F\) until a contraction meets a harmonic factor.
At level \(l\) the factorial pairing on lifted vectors is the
level form times
\[
 a_l=d_F\prod_h(j_h!)^2
 \binom{o_h+n_h+r+z_h+j_h}{j_h}z_h!,
\]
where \(o_h,n_h\) are reduced degrees. This follows by
iterating the commutator on the harmonic space and evaluating
the scalar coordinates. The factorials are nonzero.
Put \(M=\max_l v(a_l)\), \(b_l^{\,0}=\lfloor YD_{1r}\rfloor\).
For level pairs \(u_i,v_i\) of diagonal \(p^{b_l^{\,0}}\) use
\(p^U\Theta_lv_i\) and \((p^M/a_l)p^U\Theta_lu_i\).
Their common factorial diagonal is \(p^e\) with
\(e=b_l^{\,0}+2U+M\).
\leanref{Formalism.\allowbreak{}MainNotions.\allowbreak{}G.\allowbreak{}SymmetricFormBlocks.\allowbreak{}Th\_\allowbreak{}fock}{Formalism.\allowbreak{}SharpRestriction.\allowbreak{}GrowthBounds.\allowbreak{}MonomialFactorials.\allowbreak{}GoodCountMem}{101--132}
\leanref{Formalism.\allowbreak{}MainNotions.\allowbreak{}G.\allowbreak{}SymmetricFormBlocks.\allowbreak{}Th\_\allowbreak{}orth}{Formalism.\allowbreak{}SharpRestriction.\allowbreak{}GrowthBounds.\allowbreak{}MonomialFactorials.\allowbreak{}GoodCountMem}{133--151}

The field span of ordinary generators has dimension at most
the sum of level ranks. Conversely each \(\Theta_l\) is
injective on the level space. Multiplication by \(\prod_h\widehat E_h^{j_h}\) is injective in the polynomial ring. Remove these factors and set \(Z_h=1,S'_h=0\) at each tie. By \eqref{eq:p5c2-tie-operator} the specialized map is the product of the operators \(\Phi_{z_h}\). For one pair of reduced degrees \(o,n\), decompose into the harmonic levels \(E^a\ker F\), \(0\le a\le\min(o,n)\), as in \ref{B5.20q}. On this level the eigenvalue is \(\binom{z+a}{a}/\binom{o+n+r+z-1}{a}\). For \(a=0\) it is one. For \(a>0\) both binomial coefficients are positive integers since \(z\ge0\), \(r\ge1\) and \(o,n\ge a\). This computation over the rationals remains valid after extension to \(\Q_p\). Each specialized operator is invertible, including at the terminal tie, and the operators act on disjoint pairs. A zero image before specialization therefore implies zero input.
The distinct level images are independent by the distinct
harmonic degrees in the disjoint tied pairs. They lie in the
block span by the denominator bound above. This gives the
level-dimension sum bound by the block rank.
\leanref{Formalism.\allowbreak{}MainNotions.\allowbreak{}G.\allowbreak{}SymmetricFormBlocks.\allowbreak{}blkZ\_\allowbreak{}gen}{Formalism.\allowbreak{}SharpRestriction.\allowbreak{}GrowthBounds.\allowbreak{}MonomialFactorials.\allowbreak{}DegFactNeZero}{58--127}

Put \(s_b=\prod_{i<\sigma}(H_i-w_i)!\,h_i!\).
The block lattice in this span is defined coefficientwise
by \(2v(c_\mu)+v(\mu!)+v(s_b)\ge0\), with form
\(s_b\langle\ ,\ \rangle\). For a monomial of an ordinary
generator,
\[
 M-v(\mu!)\le((2s+2)r+3(t+1))L_p .
 \tag{5.20.19g}\label{eq:p5c2-level-cost}
\]
Assign the two scalar coordinates to their original letters.
The total degrees per letter are fixed at every level.
The factorials in \(a_l\), except its \(t+1\) binomials,
divide the product of these total-degree factorials.
Comparing to \(\mu!\) costs at most \(L_p\) per coordinate.
Use \(L_p\le\ell\) and \(r\ge4\) to get
\(e\le b_l^{\,0}+L_*+v(\mu!)\) with
\(L_*=10(s+1)(r+1)\ell+\kappa_0\).
\leanref{Formalism.\allowbreak{}MainNotions.\allowbreak{}G.\allowbreak{}SymmetricFormBlocks.\allowbreak{}val\_\allowbreak{}gen}{Formalism.\allowbreak{}SharpRestriction.\allowbreak{}GrowthBounds.\allowbreak{}MonomialFactorials.\allowbreak{}FinPosEqCast}{199--238}

With \(f_0=\lfloor b_l^{\,0}+L_*\rfloor\) scale both families
by \(\omega^{f_0-e-v(s_b)}\). They belong to the block lattice
and have a diagonal of valuation \(f_0\). Saturating the
first family preserves rank. A primitive vector in its
saturation has some coordinate of valuation \(\le0\) in the
original family, so the opposite vector detects it at
valuation \(\le f_0\). Smith reduction gives at least that
many good directions in the block. Increasing \(C_A\) makes
\(f_0\le D_j\) since \(YD_{1r}=O(Y^2(s+2)^2r^2\ell)\).
\leanref{Formalism.\allowbreak{}MainNotions.\allowbreak{}G.\allowbreak{}SymmetricFormBlocks.\allowbreak{}block\_\allowbreak{}dual\_\allowbreak{}family\_\allowbreak{}bound}{Formalism.\allowbreak{}SharpRestriction.\allowbreak{}GrowthBounds.\allowbreak{}MonomialFactorials.\allowbreak{}BlockBoundOfDualFamilies}{105--214}

For a digit class the full dimension in the \(k\)-variable
rows is
\[
 m_{\mathcal C}(K)=P(C'_0)P(b_+)\prod_i P(H_i)P(K_i).
\]
The sum of all level ranks over the kept blocks in this class
is at most \(m_{\mathcal C}(K)\). First use the independent
level maps into each block. Then embed the block spaces into
the full digit-class space. Their first-coordinate degrees
and contracted-pair degrees distinguish the block labels.
The prefix capacities are the bounds on these first-coordinate
degrees in the successive interlacing steps. Counting without
those bounds gives precisely the displayed product, since
each full row is a symmetric power on \(k\) variables. 

Changing one coordinate from \(x=K_i\) to its class maximum
\(y\) changes one factor. The residue sum is at most \(kP_i\le b'\),
so \(xP_i\ge a_1-b'\ge3a_1/4\). The enlarged \(a\)-window
gives \(r(y-x)P_i\le 3(3a_1/4)\) by its width
\(a/(kY)+J_b+c_b\) and the size inequalities. Hence \(x\ge1\),
\(r(y-x)\le3x\), and
\[
 \frac{P(y)}{P(x)}=\prod_{j=1}^{r}\frac{y+j}{x+j}
 \le(y/x)^r\le e^3 .
\]
Multiply \eqref{eq:p5c2-internal-mean} by \(d_\xi\) and sum
\eqref{eq:p5c2-terminal-loss} over levels and classes.
The internal loss is at most \((30+C_M)s e^3/Y\) times the
summed class dimension and the terminal loss at most \(5/Y\)
times that dimension. The comparison with the full
segment rank costs at most \(2e^3\) in the weighted count.
It follows by summing the class dimensions over digits and
using the kept-block comparison below, with the degree
ratio bound for variation inside each progression. Thus
each segment contributes at most \(C_A(s+1)N_1/Y\), for
\(C_A\ge2((30+C_M)e^3+5)e^3\). For fixed \(B\), the product
average restricts to the uniform average of the segment's
interior \(A\)-weights and its right endpoint. Sum using
\(\sum_j(s_j+1)\le N\) and add \(\kappa_P\). This proves
\eqref{eq:p5c2-rank-one-mean}.
\leanref{Formalism.\allowbreak{}MainNotions.\allowbreak{}G.\allowbreak{}CountingStatementReduction.\allowbreak{}mCl}{Formalism.\allowbreak{}SharpRestriction.\allowbreak{}GrowthBounds.\allowbreak{}CountingStatementReduction}{2946--2950}
\leanref{Formalism.\allowbreak{}MainNotions.\allowbreak{}G.\allowbreak{}SymmetricFormBlocks.\allowbreak{}claim2}{Formalism.\allowbreak{}SharpRestriction.\allowbreak{}GrowthBounds.\allowbreak{}MonomialFactorials.\allowbreak{}SegmentWeightedCountingBoundOfAllRemainingBounds}{40--54}
\leanref{Formalism.\allowbreak{}MainNotions.\allowbreak{}G.\allowbreak{}SymmetricFormBlocks.\allowbreak{}ratioCondition\_\allowbreak{}of\_\allowbreak{}segFacts}{Formalism.\allowbreak{}SharpRestriction.\allowbreak{}GrowthBounds.\allowbreak{}MonomialFactorials.\allowbreak{}RatioConditionOfSegFacts}{74--168}
\leanref{Formalism.\allowbreak{}MainNotions.\allowbreak{}G.\allowbreak{}SymmetricFormBlocks.\allowbreak{}segment\_\allowbreak{}weighted\_\allowbreak{}counting\_\allowbreak{}bound\_\allowbreak{}of\_\allowbreak{}all\_\allowbreak{}remaining\_\allowbreak{}bounds}{Formalism.\allowbreak{}SharpRestriction.\allowbreak{}GrowthBounds.\allowbreak{}MonomialFactorials.\allowbreak{}SegmentWeightedCountingBoundOfAllRemainingBounds}{134--173}
\leanref{Formalism.\allowbreak{}MainNotions.\allowbreak{}G.\allowbreak{}GramWindowAlt.\allowbreak{}theorem\_\allowbreak{}A\_\allowbreak{}closed\_\allowbreak{}form\_\allowbreak{}holds}{Formalism.\allowbreak{}SharpRestriction.\allowbreak{}GrowthBounds.\allowbreak{}FockGramWindow}{24--25}
\leanref{Formalism.\allowbreak{}MainNotions.\allowbreak{}flipOneThree.\allowbreak{}Pf.\allowbreak{}segment\_\allowbreak{}mean\_\allowbreak{}kappa\_\allowbreak{}bound}{Formalism.\allowbreak{}SharpRestriction.\allowbreak{}GrowthBounds.\allowbreak{}ColumnSpanDensity.\allowbreak{}LocalHypothesisOfGlobal}{143--177}

\paragraph{Boundary rank.}
Let \(r_j\) be the total rank of the kept block lattices on
segment \(j\). The ranks satisfy
\[
 \prod_jr_j\ge (1-\kappa_P)\prod_{s\in S}P(A_s)P(B_s)/E_q.
\]
In a fixed digit class count interlacing paths of the
\(k\)-row chain starting at a single row of degree
\(C'_0-\lceil m_0\rceil\). At the end retain the
\(r\)-row terminal weight, with multiplicity its Weyl
dimension. To compare with the \(r\)-row blocks, strip
the first row at each step. If the first-row length before
a vector-covector pair is \(c_i\) and afterwards \(c_{i+1}\),
the removed degree is \(c_i+H_i-c_{i+1}\).
In the early pairs split it into the increase of the sum of
the remaining rows and the part removed on the covector
step. These are \(w_i,h_i\). In the late pairs the removed
degree gives \(o_i\) with the residual contraction count
recorded at its tie. At the terminal pair use the first-row
shift and the terminal contraction count. The remaining
rows, these counts, and the boundary row recover the first
row by its degree at every step. More explicitly, adjoin a terminal index \(0\le j'\le\lfloor L'b_+/k\rfloor\). Write \(|t|=\sum_a t_a\) for the sum of a row and \(t^-=(t_2,\ldots,t_k)\) for its tail in this paragraph. At an inner pair let \(t,v,t'\) be the rows before the vector, after the vector, and after the covector. Set \(w=|v^-|-|t^-|\), \(h=v_1-t'_1\). Interlacing shows \(0\le w\le H_i\), \(0\le h\le K_i\). The two tail degrees are \(w,K_i-h\), and \(w+h=t_1+H_i-t'_1\). At an early pair the block records \((w_i,h_i)=(w,h)\). At a late pair it records \(o_i=w+h\) and its level is \(j_i=h\), so the reduced degrees are again \(w,K_i-h\). If \(t\) is the last full row and \(\lambda\) the terminal interlacing row, set \(d=|\lambda|-|t^-|\ge0\), \(o_+=d+j'\). Adjoining \(\lambda\) to the tail chain gives the last vector strip of degree \(d\). The output includes a tableau of shape \(\lambda\) with its determinant twist, and a monomial of degree \(b_+-j'\) in \(\xi\). These are precisely the two multiplicities. From the reduced path recover each full row by appending its missing first entry, equal to the prescribed signed degree sum to that point minus the tail sum. The final row gives \(\lambda\), and the terminal level gives \(j'\). This proves injectivity. Along the path \(C_i=c_i+\lceil m_0\rceil\), by the capacity recursion and its initial value. The largest entry of the tail after the vector is \(v_2\le t_1=c_i=\lfloor C_i-m_0\rfloor\). At the last prefix \(\lambda_1\le c_s\). These are the prefix caps of the reduced path.

For the unrestricted path count, Pieri interlacing counts
the product of the symmetric-power dimensions at its
letters. For the paired-chain count put \(A=C'_0-\lceil m_0\rceil\) and use the signed degrees \(A,H_0,K_0,\ldots,H_{s-1},K_{s-1}\), with signs \(+,+,-,\ldots,+,-\). After the last signed step choose an interlacing \(r\)-row \(\lambda\), weighted by \(\dim V_\lambda\). Number the inner pairs from zero. Let \(c_i\) be the first coordinate just before pair \(i\), \(c_s\) the last first coordinate, \(R_i\) the second coordinate just after the vector of pair \(i\), \(u_i=c_i+H_i-c_{i+1}-R_i\), and \(n_i=-\lambda_{r-s+i+1}\) in one-based row indices. The two rejection tests are \[u_i+4(H_i+r)/p^2>L'\max(c_i,H_i)/k,\qquad n_i<K_i(1-L'/k).\] We check that a pair passing both tests gives a kept block. Write \(a_c=C'_0\), \(\theta=L'/k\) and \(A=a_c-\lceil m_0\rceil\). The digit bounds give \((b'-kP_i)/P_i\le H_i\le 2b'/P_i\). Successive odd distances differ by at least two, so \(\sum_{h<i} H_h\le4H_i/p^2\), using \(kQ_p\le b'/4\) and the geometric series. The size bounds give \(E_0/a_1,\ (m_0+1)Q_p/a_1,\ kQ_p/b'\le\theta/2\) on increasing the fixed thresholds. Indeed the first four terms of \(E_0/a_1\) are bounded by \(C_{E1}(2/p^2+\theta(N_1+1)/p)\), using \(b'/a_1<1/p\). For the remaining terms use \(Q_p^4 c_4/b'\le2/(Hk^3NY)\), \(m_0=(C_m kN_1Q_p+12)c_4\), \(p\ge k^2\), and \(HN\mathcal L^2\le k\). In particular \(\lceil m_0\rceil/a_c\le2\theta\). At an early position \(p^{d_i}b'<a(1-\rho_{\rm loc})\), and \(b_i/b'\le1+\Delta/b'+3\mathcal L/k\). At a late position \(p^{d_i}b'>a(1+\rho_{\rm loc})\). Using \(a_1\le a+1\), \(1/a\le\theta\), and \(\rho_{\rm loc}\ge HN_1\theta+4\Delta/b'\), these give \(H_i\le(1-10N_1\theta)a_c\) early and \(a_c\le(1-10N_1\theta)H_i\) late, for sufficiently large \(H\).

Write \(O_i=w+h\) in the preceding tail construction. The second entry before a vector is at most the total of the preceding vector degrees \(\sum_{h<i}H_h\), since the initial second entry is zero and covector steps decrease entries. Hence \(R_i\le w+4H_i/p^2\). Also \(R_i\le c_i\) by interlacing and \(c_i\le A+\sum_{h<i}H_h\). Early, this gives \(c_i\le2a_c\). At each preceding early step the first acceptance test gives \(O_i\le H_i+4H_i/p^2+2\theta a_c\). The recursion then gives \(c_i\ge A-i(4/p^2+2\theta)a_c\ge(1-6N_1\theta)a_c\ge H_i\), using \(4/p^2\le\theta\). The maximum in the test is therefore \(c_i\). We obtain \(O_i\le(1+4/p^2)H_i+\theta c_i\) and \(h=O_i-w\le\theta c_i\le2\theta a_c\). Late, \(c_i\le a_c+4H_i/p^2\le(1-\theta)H_i\). Thus \(O_i\le c_i+\theta H_i\le H_i\). These are the two inner block tests. For the terminal test the increase in positive parts from the last tail to \(\lambda\) is at most \(c_s\), by telescoping the interlacing inequalities. The decrease in absolute negative parts is at most \(\sum_i K_i-\sum_i n_i\), since the covectors introduced at most \(\sum_i K_i\) negative boxes. The second acceptance test bounds this by \(\theta\sum_iK_i\). Hence \(o_+\le c_s+\theta\sum_i K_i+\theta b_+\), which is within the stated terminal bound since \(K_i\le2a_1/P_i\). To estimate the rejected weight put \(C=c_i+r\), \(B=(H_i+r)(1+4/p^2)\), \(R=R_i+r\), \(v=u_i+4(H_i+r)/p^2\). Then \[R(c_{i+1}+r)=CB-(C-R)(B-R)-Rv.\] The interlacing bounds give \(0\le R\le\min(C,B)\), \(v\ge0\), \(\lambda_1\le c_s\), and \(0\le n_i\le K_i(1+4/p^2)\). For \(\tau=L'/(4k)\le1/16\), failure gives \(v>2\tau\max(C,B)\): here \(H_i,K_i,A\ge k^3\), \(p\ge k^2\), and \(\max(C,B)\le2\max(c_i,H_i)\). If \(R\ge\min(C,B)/2\), then \(Rv\ge\tau CB\). Otherwise \((C-R)(B-R)\ge CB/4\). A failed step therefore costs a factor \(1-\tau\) in the product. A failed second test gives the same factor in the bound for \(n_i+r\) after the additive \(r\) is absorbed. Indeed, with \(\eta=k^{-2}\), \(r/K_i\le\eta\), and \(1-L'/k+\eta\le (1-\tau)e^{2\eta}\) for \(L'\ge1\). At all other steps \((H_i+r)(1+4/p^2)\le e^{2\eta}H_i\) and \(n_i+r\le e^{2\eta}K_i\). Multiply the inequalities for \(R(c_{i+1}+r)\); the intermediate \(c_i+r\) cancel. Using \(A+r\le e^\eta A\) and raising the result to power \(r-1=k-2\) gives the loss \(\exp(2-L'/5)\), since \((k-2)(4s+1)/k^2\le2\) and \((k-2)/(4k)\ge1/5\) for \(k\ge10\). To pass to Weyl dimensions use \(y^{r-1}/(r-1)!\le\binom{y+r-1}{r-1}\le(y+r)^{r-1}/(r-1)!\). There are at most \((8a_1)^{4(N_1+1)^2}\) unlabelled pairs: each intermediate row and the last interlacing row is determined by its first \(s+1\) and last \(s\) entries, since all intervening ones vanish by interlacing. For clarity, let \(D_r(y)=\binom{y+r-1}{r-1}\). The Weyl product for a decreasing row \(\lambda\) is at most \(\prod_aD_r(|\lambda_a|)\). In its factor with indices \(a<b\), bound \(1+(\lambda_a-\lambda_b)/(b-a)\) by \((1+\max(\lambda_a,0)/(b-a))(1+\max(-\lambda_b,0)/(b-a))\), then multiply. Repeated interlacing bounds the positive entries after the first by \(R_{s-1},\ldots,R_0\) in order: a covector step decreases each entry, and at a vector step each entry after the first is at most the preceding entry of the previous row. The last \(s\) entries have absolute values \(n_0,\ldots,n_{s-1}\), and all intervening entries vanish. The same interlacing in the last \(s-1-i\) covector steps bounds \(n_i\) by \(\sum_{h\le i}K_h\le K_i(1+4/p^2)\), by the geometric series for the divided covector degrees. Thus the product just estimated bounds the rejected Weyl dimension by \(e^{2-L'/5}D_r(A)\prod_iD_r(H_i)D_r(K_i)\). Each \(D_r(y)\le P(y)\), so the ratio is relative to the unrestricted \(k\)-row dimension. The additional factor \((8a_1+k)^{(2N_1+2)^2}\ge1\) may also be included in this upper bound. Thus the logarithm of the total exceptional fraction is at most \(2-L'/5+64N_1^2\log a_1\le-L'/8\) for \(C_0\ge880\), using \(k\le a_1\), \(a_1\ge9\), and \(\log a_1\ge1\). The terminal cutoff retains at least
\(1-e^{-L'/3}\). For example its count sums the homogeneous
dimensions on \(r\) variables for degrees
\(b_+-j'\) with \(j'\le\lfloor L'b_+/k\rfloor\); the binomial
sum gives \(P(b_+)-P(b_+-\lfloor L'b_+/k\rfloor-1)\).
The reserve at the start retains at least
\(1-r\lceil m_0\rceil/C'_0\) by the same dimension ratio.
Summing over the digits restores \(P(a_i)P(b_i)\), since each
monomial has a unique quotient and remainder row.
For a shared even mark the starting degree of the next
segment is \(\lfloor(a_1-E_0)/p^{G}\rfloor\), compared with
the undivided \(A\)-weight in the preceding endpoint.
The binomial product compares these dimensions with
scale \(p^{-Gr}\) and losses bounded by a constant times
\(kQ_p(m_0+N_1k)/b'+1/Y+kE_0/a_1\). Here the first term
allows the reserved and rounded degrees, the second
allows the width \(a_1/(kY)\), and the last allows the
subtraction of \(E_0\). Each boundary row is used once in
this comparison. Multiplication over segments uses
\(\sum_jG_j=n'\) and
\(p^{-n'r}E_q=\sum_{i=0}^{r}p^{-n'i}\ge1\).
The sum of the segment losses is at most \(\kappa_P\) by
\eqref{eq:p5c2-rank-one-margin}. If that sum is at least one
the desired lower bound follows from nonnegativity.
\leanref{GlobalTheory.\allowbreak{}EX.\allowbreak{}RangeOne.\allowbreak{}tensor\_\allowbreak{}of\_\allowbreak{}pairLift}{Development.\allowbreak{}AbelianVarietyTheory.\allowbreak{}Proofs.\allowbreak{}GlobalInterpolation.\allowbreak{}TensorOfPairLift}{583--595}

\paragraph{Paired lifting.}
For a tuple \(\beta\) of kept blocks, map each segment into
the symbol coordinates. Insert its common low digits and
raise the divided coordinates to the powers \(P_i\). Insert
the prescribed scalar coordinates of its ties.
Here are the two symbol maps on a segment block. Assign each ordinary coarse coordinate to the corresponding nonzero coordinate of its physical row. At an internal or terminal tie assign \(Z\) to coordinate zero of the vector row and \(S'\) to coordinate zero of the covector row, where the vector row here is the physical row carrying the coarse vector and its fixed shift. Thus \(Z\) supplies the shift \(H_i-o_i\), or \(a_+-o_+\) at the end. At an early pair the unused coordinate-zero slots carry the fixed quotients \(H_i-w_i,h_i\). Insert the prescribed digits, using power \(P_i\) at interior rows and power one and digit zero at the endpoint. Denote the resulting fine exponent of a coarse monomial \(X^\mu\) by \(a(\mu)\).

For a physical row \(x\) let \(d_x\) be its fine degree and \(e_x\) the sum of the coarse exponents in its assigned variable slots, excluding fixed quotients. This sum is constant on the block. Indeed \(W_aY_a\) and \(ZS'\) have the same degree in each physical row, and lowering a pair of reduced degrees by one is compensated by one tie factor. Put \(s_b=\prod_{i<\sigma}(H_i-w_i)!h_i!\), \(N_b=s_b\prod_x e_x!\), and \(t_b=\omega^{v_p(N_b)}\). Write \(J_b(\mu)\) in this paragraph for the product of the digit units in the variable slots. By \ref{B5.20ha} with \(c=c_4\), choose \(g\equiv J_b^{-1}\pmod{p^{c_4}}\) of total binomial degree at most \(4c_4\). Coordinates of power one contribute one, so the indices used by \(g\) are all interior coordinates, including the internal scalar coordinates. The maps before substitution are \(u\mapsto t_bD_gu\) and \(v\mapsto t_bv\), respectively.

For a block-lattice coefficient \(u_\mu\), the defining inequality gives \(2v(t_bu_\mu)\ge v(N_b)-v(s_b\mu!)\ge0\), since \(\mu!\mid\prod_x e_x!\). Thus \(t_bu\) has integral coefficients and lies in the field span of the ordinary generators. Apply saturation in \ref{B5.20ic}, then its derivative bound with \(b_J=\lfloor C_i-m_0(C_m)\rfloor\) at each prescribed cut. It puts \(t_bD_gu\) in the integral span with reserve \(m_0(C_m)-4c_4\), which is at least \(m_0(13)\). The second map has the full original reserve. This proves the two simultaneous symbol-membership assertions used below.
We carry out the rank-one lift explicitly using Lemma~\ref{B5.20he} with plane rank \(1\) and complementary rank \(r=k-1\), numbering the \(k\) coordinates from \(0\) to \(r\) so that coordinate \(0\) is the distinguished line. Fix a Witt matrix \(g\in\GL_k(W(\mathbf F_{p^{n'}}))\) for the edge \(e\). At each mark \(x\), the twisted vector row \(v^{(g)}_x=X^{\rm vec}_xg_x\) has nonnormal coordinate \(0\) and normal coordinates \(1,\ldots,r\), while the twisted covector row \(w^{(g)}_x=X^{\rm cov}_xg_x^{-\mathsf t}\) has normal coordinate \(0\) and nonnormal coordinates \(1,\ldots,r\). For the depth bound, apply \ref{B5.20ia} with \(Q=Q_p\) and \(a_S\) the weight at the left even endpoint: \(J_b+c_b\le3\mathcal L b'/k\) (since \(L'\le2\mathcal L\)) gives \(b'\le b_i,b_+\le b'(1+\rho_{\rm loc})<2b'\), and \(1/a<\rho_{\rm loc}\) by \eqref{eq:p5c2-rank-one-size} for large \(H\). If \(p^{d_i}<a_1/b'\), the excluded odd-power interval forces \(p^{d_i}<(a/b')(1-\rho_{\rm loc})\), so \(p^{d_i}b_i\le a_S\); otherwise \(a_1/p^{d_i}\le b'\le b_i\). Summing \ref{B5.20ia} over segments using \(\sum_j(s_j^{\rm int}+1)=N\) (where \(s_j^{\rm int}\) counts interior marks of segment \(j\)) bounds the initial block normal degree \(h_0\) by \(t_0\) for \(C_t\ge30\).

For the one-sided lift put \(K=13kN_1Q_p\), \(m_0(v)=vkN_1Q_pc_4+12c_4\), and \(m'_0=m_0(13)=12c_4+c_4K\in\N\); taking \(C_m\ge17\) gives \(m_0(C_m)-4c_4\ge m'_0\). Because the segment symbol maps have disjoint alphabets (each containing its interior pairs and its right endpoint), at every even start \(s_j\) (the right endpoint of segment \(j-1\)) the preceding terminal shift \(Z_{+,j-1}^{a_{s_j}-o_{+,j-1}}=(X^{\rm vec}_{s_j,0})^{a_{s_j}-o_{+,j-1}}\) has exponent \(a_{s_j}-o_{+,j-1}\ge a_1-E_0\ge C'_{0,j}p^{G_j}\) (for \(C_{E1}\ge3\), using \(\sum_i P_i^{-1}\le p/(p^2-1)\), \(G_{j-1}\ge2\), and \(b_{+,j-1}<2b'\) in the kept-block bound for \(o_{+,j-1}\)), supplying \((C'_{0,j}-m'_0)p^{G_j}\) powers of \(X^{\rm vec}_{s_j,0}\) for segment \(j\)'s column starts, \(m'_0p^{G_j}\) reserve powers of \(X^{\rm vec}_{s_j,0}\) at \(s_j\), and \((a_{s_j}-o_{+,j-1})-C'_{0,j}p^{G_j}\ge0\) leftover powers. On each segment \(j\) (from \(s_j\) to \(s_{j+1}=s_j+G_j\), with interior marks \(x_i=s_j+d_i\), \(0\le i<s\), and powers \(P_i=p^{G_j-d_i}\)), expand each reduced polynomial \(y\) from either symbol map (which has integral coefficients and satisfies the prefix bounds \(\lfloor C_i-m'_0\rfloor=C_i-m'_0\) at \(i\ge\sigma\) and at \(s\)) in the integral rank-\(r\) exterior-column path basis from the proof of Lemma~\ref{B5.20hc} (applied symmetrically to upper bounds \(\kappa_1\le C_i-m'_0\), since the greatest first-coordinate weight of a column is \(\mathbf1_{h(c)>0}\) for \(c>0\) and \(-\mathbf1_{h(c)=0}\) for \(c\le0\), summing to \(\kappa_1\)). At each early mark \(x_i\) (\(i<\sigma\)), lift the \(h_i\) fixed covector quotient units by \(h_i\) two-letter columns \(P_{x_i,0}(g)^{P_i}\) on \((s_j^{\rm vec},x_i^{\rm cov})\) (positions \(s_j<x_i\), order \(P_i\le Q_p\), leading term \(\omega^{P_i}X_{s_j,0}^{p^{G_j}}(X^{\rm cov}_{x_i,0})^{P_i}\) modulo \(\omega^{P_i+1}O[X]\) by Lemma~\ref{B5.20he}(1)). Since \(h_l\le2\theta C'_0\) gives \(\sum_{l<\sigma}h_l\le2N_1\theta C'_0\), the early inequality \(\sum_{h\le i}w_h\le(1+4/p^2)H_i\le(1-9N_1\theta)C'_0\le C'_0-m'_0-\sum_{l<\sigma}h_l\) (from \(H_i\le(1-10N_1\theta)C'_0\) and \(m'_0/C'_0\le2\theta\)) together with the prefix bounds \(C_i-m'_0\) at \(i\ge\sigma\) and at \(s\) ensures that every complementary path \(\kappa\) occurring in \(y\) joins with the nondecreasing rank-one plane path \(\nu\) (starting at \(C'_0-m'_0-\sum_{l<\sigma}h_l\) at \(s_j^{\rm vec}\), gaining \(H_i-w_i\) for \(i<\sigma\) or \(H_i-o_i=z_i\) for \(i\ge\sigma\) at \(x_i^{\rm vec}\), and changing by \(0\) at every \(x_i^{\rm cov}\) and at \(s_{j+1}^{\rm vec}\), so \(\nu=C_i-m'_0\) before \(x_i^{\rm vec}\) for \(i\ge\sigma\) and at \(s_{j+1}^{\rm vec}\)) to give a rank-\(k\) path \(\mu\) on the \(2s+2\le2N_1+2\le r\) letters \((s_j^{\rm vec};x_0^{\rm vec},x_0^{\rm cov},\ldots,x_{s-1}^{\rm vec},x_{s-1}^{\rm cov};s_{j+1}^{\rm vec})\). The only equal-position letter pairs are \((x_i^{\rm vec},x_i^{\rm cov})\): every step at \(x_i^{\rm vec}\) follows letters strictly before \(x_i\), while at \(x_i^{\rm cov}\) the plane entry \(\nu\ge0\) is unchanged, so any column \(c\in\Z\) used at \(x_i^{\rm cov}\) still has \(\mu_0=\nu\ge c\) after \(x_i^{\rm cov}\) (for \(c>0\) because \(c\le\kappa_1^{\rm before}\le\nu\), and for \(c\le0\) because \(\nu\ge0\ge c\)), hence height \(\ge1\) after \(x_i^{\rm cov}\) and height \(\ge2>1\) before it; thus \(x_i^{\rm cov}\) is never costly. Lemma~\ref{B5.20he}(1) therefore applies to every active column \(c\) of \(\mu\) with \(z=s_{j+1}\): \(K_c(g;J)\in O[X]\) has degree \(p^{G_j}\), \(P_i\), or \(1\) at each used letter at \(s_j\), \(x_i\), or \(s_{j+1}\) (and \(0\) at unused letters), satisfies \(\operatorname{ord}_{g'}(K_c(g;J))\ge\mathrm{oc}(c)\) for every \(g'\in\GL_k(W(\mathbf F_{p^{n'}}))\), agrees modulo \(\omega^{\mathrm{oc}(c)+1}O[X]\) under \(s_g\) with \(\omega^{\mathrm{oc}(c)}\) times the entrywise-powered separated column, and has \(\mathrm{oc}(c)\le2p^{G_j}\le2Q_p\); taking terminal index \(J=\{0\}\cup J_\kappa\) for \(c>0\) and \(J=J_\kappa\) for \(c\le0\) (where \(J_\kappa\subseteq\{1,\ldots,r\}\) is the terminal index in \(\kappa\)) gives \(K^{\rm sep}_c(\{0\}\cup J_\kappa)=X^{\rm vec}_{a_0,0}K^\kappa_c(J_\kappa)\) for \(c>0\) (with \(a_0\in\{s_j,x_i\}\) the step where \(\nu\) crosses \(c\)) and \(K^{\rm sep}_c(J_\kappa)=K^\kappa_c(J_\kappa)\) for \(c\le0\), with no sign change since \(0<1,\ldots,r\) and coordinate \(0\) vanishes in all subsequent updates.

Lift the remaining factors of each generator on segment \(j\) as follows. Each internal tie \(\widehat E_i^{j_i}\) (\(i\ge\sigma\)) lifts as the product of \(j_i\) tie factors \(E_{x_i}^{P_i}\), and the terminal tie \(\widehat E_+^{j'}\) lifts as \(j'\) factors \(E_{s_{j+1}}\); by Lemma~\ref{B5.20he}(3) (with \(P=P_i\) and \(P=1\)), each \(E_{x_i}^{P_i}\) has exact order \(P_i\) under every \(s_{g'}\) and agrees modulo \(\omega^{P_i+2}O[X]\) with \(\omega^{P_i}\operatorname{pow}_{s_{j+1}}(\widehat E_i)\) under \(s_g\), while each \(E_{s_{j+1}}\) has exact order \(1\) and zero error under \(s_g\). At each interior mark \(x_i\), lift each normal digit \(u_{ic}<P_i\) (\(1\le c\le r\)) by \(B_{x_i,c}(g)^{u_{ic}}\) and \(v_{i,0}<P_i\) by \(P_{x_i,0}(g)^{v_{i,0}}\), where \(D_{x_i}(g)=(v^{(g)}_{s_j,0})^{p^{d_i}}\), \(B_{x,c}(g)\) is the \(2\times2\) minor on columns \(\{0,c\}\), and \(P_{x,0}(g)=\sum_{a=0}^r(v^{(g)}_{s_j,a})^{p^{x-s_j}}w^{(g)}_{x,a}\) are the rank-one bordered and pairing factors of Lemma~\ref{B5.20he}(2). Each has order at least \(1\) under every \(s_{g'}\) and leading term \(\omega X_{s_j,0}^{p^{d_i}}X^{\rm vec}_{x_i,c}\) or \(\omega X_{s_j,0}^{p^{d_i}}X^{\rm cov}_{x_i,0}\) modulo \(\omega^2O[X]\) under \(s_g\). Across all \(k\) normal coordinates and \(s\le N_1\) interior marks, lifting the initial normal digits consumes at most \(kN_1p^{G_j}\) degrees of \(X_{s_j,0}\) (at most \(kN_1\) units of \(p^{G_j}\)). Since the early fixed vector quotients \((X^{\rm vec}_{x_i,0})^{P_i(H_i-w_i)}\) (\(i<\sigma\)) and the shifts \(Z_i^{z_i}\) (\(i\ge\sigma\)) are already supplied by the plane steps of \(\nu\) at \(x_i^{\rm vec}\), and \(Z_{+,j-1}^{a_{s_j}-o_{+,j-1}}\) supplies segment \(j\)'s \(C'_{0,j}p^{G_j}\) start and reserve powers together with the \((a_{s_j}-o_{+,j-1})-C'_{0,j}p^{G_j}\) leftover powers of \(X_{s_j,0}\), across all segments we lift \(X_{s_{j_0},0}^{c_4}\prod_j X_{s_j,0}^{c_4(k-1)(p^{G_j}-1)}\) (carved from the reserve powers of \(X_{s_j,0}\)) by the rank-one localizer \(R(g)^{c_4}\) of \ref{B5.20x} in the twisted rows \(v^{(g)}_{s_j}\), which has order at least \(c_4\) on every competing line, leading term \(X_{s_{j_0},0}^{c_4}\prod_j X_{s_j,0}^{c_4(k-1)(p^{G_j}-1)}\) under \(s_g\), and degree \(c_4((k-1)(p^{G_j}-1)+\mathbf1_{j=j_0})\le kc_4p^{G_j}\) at start \(s_j\) (at most \(kc_4\) units of \(p^{G_j}\)), while the remaining nonnormal digits, the terminal monomial \(\xi^\alpha\), and all remaining reserve or leftover powers of \(X_{s_j,0}\) lift by substituting \(v^{(g)}_{x,0}\) and \(w^{(g)}_{x,c}\) (\(c\ge1\)), which are integral of order \(\ge0\) under every \(s_{g'}\) and exact under \(s_g\). Together the localizer and initial normal digits consume at most \(kc_4+kN_1\le2kN_1Q_pc_4\) units of \(p^{G_j}\) at start \(s_j\), leaving at least \(11kN_1Q_pc_4+12c_4\) units of \(p^{G_j}\) in \(X_{s_j,0}\) for the recursion.

We now lift any term \(S\) at level \(h=h_0+d\) (\(0\le d<c_4\)) by descending induction in \(h\) to precision \(h_*=h_0+c_4\), maintaining at least \(11kN_1Q_p(c_4-d)+12c_4\) reserve units of \(p^{G_j}\) in \(X_{s_j,0}\) at every start \(s_j\) beyond the localizer degree. Here \(h\) is the sum of the explicit \(\omega\)-power, column orders \(\mathrm{oc}(c)\) (and \(P_i\) for the early two-letter columns), individual tie orders \(P_i\) and \(1\), and number of bordered or pairing factors. Let \(F^0\) be the raw lift of \(S\). By Lemma~\ref{B5.20he}(4), \(F^0\) satisfies \(\operatorname{ord}_{g'}(F^0)\ge h+c_4\ge h_*\) on every competing line \(g'\). Expand \(s_g(F^0)-S\), leaving unchanged column, tie, and bordered or pairing factors in their leading form, and discard terms in \(\omega^{h_*}O[X]\). In every remaining summand, rises add across factors: an exterior column error of rise \(e_{\rm col}\ge1\) has normal degree \(t\le\mathrm{oc}(c)+e_{\rm col}\le(2Q_p+1)e_{\rm col}\); an internal tie factor error of rise \(e_{\rm tie}\ge2\) has normal degree \(P_i\le Q_p\le Q_pe_{\rm tie}\); a bordered or pairing error of rise \(e_{\rm bor}=p^{x-s_j}-1\ge1\) has normal degree \(e_{\rm bor}+1\le2e_{\rm bor}\); and a localizer error of rise \(e_{\rm loc}\ge1\) has normal degree \(e_{\rm loc}\) and loses at most \(e_{\rm loc}\) powers of \(X_{s_j,0}\) at each start \(s_j\). Thus a child summand of total rise \(e\ge1\) has at most \((2Q_p+1)e\) exposed normal variables (with available \(\omega\)-power at least the number of exposed variables) and loses at most \((2Q_p+1)e\) powers of \(X_{s_j,0}\) at each start \(s_j\). Pair each exposed normal variable (\(X^{\rm vec}_{x,b}\) for \(1\le b\le r\), or \(X^{\rm cov}_{x,0}\)) at \(x\in(s_j,s_{j+1}]\) (using the preceding start \(s_{j-1}\) when \(x=s_j\)) with \(s_g(D_x(g))=(X^{\rm vec}_{s_j,0})^{p^{x-s_j}}\) from the reserve powers of \(X^{\rm vec}_{s_j,0}\) to form a leading factor \(\omega(X^{\rm vec}_{s_j,0})^{p^{x-s_j}}X^{\rm vec}_{x,b}\) or \(\omega(X^{\rm vec}_{s_j,0})^{p^{x-s_j}}X^{\rm cov}_{x,0}\) (preserving every row multidegree, lifted at the next step by \(B_{x,b}(g)\) or \(P_{x,0}(g)\)), and restore the localizer leading monomial \(X_{s_{j_0},0}^{c_4}\prod_j X_{s_j,0}^{c_4(k-1)(p^{G_j}-1)}\). Together these consume at most \((2Q_p+1)e\) units of \(p^{G_j}\) for the bordered factors plus at most \((2Q_p+1)e\) powers of \(X_{s_j,0}\) to restore lost start degrees, altogether at most \((4Q_p+2)e\le11kN_1Q_pe\) units of \(p^{G_j}\) at each start \(s_j\). Subtracting the inductively constructed lifts \(F_\nu\) of all child terms gives an integral lift \(F=F^0-\sum_\nu F_\nu\) of the exact multidegree satisfying \(s_g(F)-S\in\omega^{h_0+c_4}O[X]\) and \(s_{g'}(F)\in\omega^{h_0+c_4}O[X]\) on every competing line \(g'\). Applying this at \(d=0\) to \(S=\omega^{h_0}S_{\rm sym}\) (for the assembled symbol \(S_{\rm sym}\)) and dividing \(F\) by \(\omega^{h_0}\) yields an integral vector in the distinguished edge coordinates agreeing with \(S_{\rm sym}\) modulo \(\omega^{c_4}\) and having order at least \(c_4\) on every competing edge of the same source. Apply the identical construction on the reversed dual edge \(e^\vee\), at depth \(t_1=\lfloor t_0(A,B)\rfloor_+\). The normalized dual identification turns the second lift into a target functional and takes its competing-source conditions to the competing-incoming conditions.

A fine exponent determines its block. Reduction modulo \(P_i\) gives both digit rows, and hence \(H_i,K_i\). At an early pair the two coordinate-zero quotients are \(H_i-w_i,h_i\). After removing these fixed quotients and all the digits, division by the row powers recovers every variable exponent, since distinct variable slots go to distinct coordinates. At a late pair recover \(o_i\) as the sum of the ordinary vector exponents plus the exponent of \(S'_i\). Each factor \(\widehat E_i\) contributes one to this sum, and the fixed shift \(Z_i^{z_i}\) contributes zero. The same sum at the terminal tie recovers \(o_+\). Thus distinct blocks have disjoint fine supports and their factorial pairing vanishes exactly.

Fix a segment block and use its row degrees and maps above. Put \(L_j=\sum_x\lfloor\log_p(\max(1,d_x))\rfloor\) and \(z_j=p^{kL_j}/\prod_x d_x!\). There is a positive rational number \(C_b\), constant on the block, such that \(a(\mu)!=C_b\,\mu!\,J_b(\mu)\). Indeed a variable slot of digit \(d\) and power \(P\) contributes \(d!(P!)^{\mu_v}\) after its coarse factorial and digit unit are removed. The sum of the exponents of \(P!\) in a physical row is \(e_x\), and every unused slot has a fixed fine exponent. Set \(\kappa_j=z_jp^{v(N_b)}C_b/s_b\). For two block coefficients the quantity \(s_b\mu!u_\mu v_\mu\) is integral, by adding their two lattice inequalities. The factorial identity and \(gJ_b\equiv1\pmod{p^{c_4}}\) give the normalized block pairing \(\kappa_j(s_b\langle u,v\rangle+\omega^{c_4}y)\) for some \(y\in O\).

To bound \(\kappa_j\), take the level-zero exponent \(\mu_0\) placing each ordinary letter degree on its first ordinary coordinate, each shift on \(Z\), and zero on \(S'\). This exponent has the required row degrees, whether or not its monomial meets the prefix cuts. Let \(M_f=\prod_xd_x!/a(\mu_0)!\) and \(M_c=\prod_x e_x!/\mu_0!\). Since the digit factor is a unit, \(v(\kappa_j)=kL_j-v(M_f)+v(M_c)\). Each fine multinomial has at most \(k\) slots. Each coarse multinomial at \(\mu_0\) has at most two nonzero slots and \(e_x\le d_x\). By \ref{B3.19}, \(0\le v(M_f)\le(k-1)L_j\) and \(0\le v(M_c)\le L_j\). Hence \(0\ne\kappa_j\in O\) and \(v(\kappa_j)\le(k+1)L_j\).

The segment alphabets are disjoint: each contains its interior pairs and its right endpoint. The factorial form and its row normalization therefore multiply. Set \(\kappa_\beta=\prod_j\kappa_j\). There are two rows per mark and each degree is at most \(a_1^2\), so \(\sum_jL_j\le4N\log_p a_1\) and \(v(\kappa_\beta)\le4\kappa_0\). Finally write the distinguished substitutions of the two lifts as \(P,Q\), with integral symbols \(S,T\). The difference between the normalized pairings is the sum of the pairings of \(P-S,T\) and \(P,Q-T\). Each row normalization is integral on integral polynomials because its monomial denominator is a multinomial of valuation at most \((k-1)\lfloor\log_p(\max(1,d_x))\rfloor\). Both differences therefore contribute elements of \(\omega^{c_4}O\). Under the normalized dual-edge identification this gives
\[
 [F_{\beta,u},G_{\beta',v}]
 \equiv\mathbf1_{\{\beta=\beta'\}}\kappa_\beta
 \prod_j\langle u_j,v_j\rangle_j
 \pmod{\omega^{c_4}O},\qquad
 0\ne\kappa_\beta\in O,\quad v(\kappa_\beta)\le4\kappa_0 .
 \tag{5.20.19h}\label{eq:p5c2-pair-lift}
\]
\leanref{Formalism.\allowbreak{}MainNotions.\allowbreak{}ColumnSystems.\allowbreak{}Pf.\allowbreak{}srcHalf1\_\allowbreak{}of}{Formalism.\allowbreak{}SharpRestriction.\allowbreak{}GrowthBounds.\allowbreak{}CoarseMultiplicity.\allowbreak{}PairedLiftAtOfIdentity}{85--162}
\leanref{Formalism.\allowbreak{}MainNotions.\allowbreak{}ColumnSystems.\allowbreak{}Pf.\allowbreak{}paired\_\allowbreak{}lift\_\allowbreak{}at\_\allowbreak{}of\_\allowbreak{}identity}{Formalism.\allowbreak{}SharpRestriction.\allowbreak{}GrowthBounds.\allowbreak{}CoarseMultiplicity.\allowbreak{}PairedLiftAtOfIdentity}{163--235}
\leanref{Formalism.\allowbreak{}MainNotions.\allowbreak{}ColumnSystems.\allowbreak{}Pf.\allowbreak{}paired\_\allowbreak{}lifts\_\allowbreak{}of\_\allowbreak{}identity}{Formalism.\allowbreak{}SharpRestriction.\allowbreak{}GrowthBounds.\allowbreak{}CoarseMultiplicity.\allowbreak{}PairedLiftAtOfIdentity}{236--242}
\leanref{Formalism.\allowbreak{}MainNotions.\allowbreak{}ColumnSystems.\allowbreak{}Pf.\allowbreak{}restForm\_\allowbreak{}of\_\allowbreak{}oplus}{Formalism.\allowbreak{}SharpRestriction.\allowbreak{}GrowthBounds.\allowbreak{}CoarseMultiplicity.\allowbreak{}BFFinal}{26--27}

Take the retained Smith pairs on each segment. The matrix on
their lifted tuples is diagonal modulo \(\omega^{c_4}\), with
diagonal valuations at most \(D_0\). Since \(4D_0<c_4\), division
of each column by its diagonal gives an integral matrix
congruent to the identity modulo \(\omega\). Let \(z\) be a
diagonal of greatest valuation. After multiplying the
columns by the integral ratios of \(z\) to their diagonals,
use the inverse unit matrix to obtain exact matrix \(zI\).
The number \(n\) of pairs is the product of the retained
segment ranks. The boundary rank comparison and
\(\prod_j(1-\kappa_{T,j})\ge1-\sum_j\kappa_{T,j}\) give
\[
 n\ge(1-g(A,B))\prod_sP(A_s)P(B_s)/E_q .
 \tag{5.20.19i}\label{eq:p5c2-tensor-rank}
\]
For \(g>1\) this follows from \(n\ge0\).
\leanref{Formalism.\allowbreak{}MainNotions.\allowbreak{}ColumnSystems.\allowbreak{}Pf.\allowbreak{}liftVecAt\_\allowbreak{}of}{Formalism.\allowbreak{}SharpRestriction.\allowbreak{}GrowthBounds.\allowbreak{}BinomialRatio.\allowbreak{}PairLift}{102--153}
\leanref{Formalism.\allowbreak{}MainNotions.\allowbreak{}ColumnSystems.\allowbreak{}Pf.\allowbreak{}factors\_\allowbreak{}of\_\allowbreak{}perturb}{Formalism.\allowbreak{}SharpRestriction.\allowbreak{}GrowthBounds.\allowbreak{}BinomialRatio.\allowbreak{}FactorsOfPerturb}{72--144}
\leanref{Formalism.\allowbreak{}MainNotions.\allowbreak{}ColumnSystems.\allowbreak{}Pf.\allowbreak{}extended\_\allowbreak{}statement\_\allowbreak{}of\_\allowbreak{}parts}{Formalism.\allowbreak{}SharpRestriction.\allowbreak{}GrowthBounds.\allowbreak{}BinomialRatio.\allowbreak{}ExtendedStatementOfParts}{93--152}

\paragraph{Passing to harmonic summands.}
At a mark let \(h=\sum_iX_iY_i\) on bidegree \((a,b)\) in two
rows of \(k\) variables. Exchange the roles of the rows at odd marks.
The quotient by \(h\) in this bidegree is the flat dual Weyl
lattice. Map \(X_i\) to entry \(i\) of the first column of a matrix
and \(Y_i\) to the signed minor on its first \(k-1\) columns
omitting row \(i\). Cofactor expansion gives the relation.
Straightening gives the standard tableau basis with columns
of heights one and \(k-1\), together with its ordinary
integral coefficient image.
\leanref{Formalism.\allowbreak{}MainNotions.\allowbreak{}flipOneThree.\allowbreak{}Pf.\allowbreak{}PsiXY}{Formalism.\allowbreak{}SharpRestriction.\allowbreak{}GrowthBounds.\allowbreak{}ColumnSpanDensity.\allowbreak{}DeltaPsiH}{12--15}

Let \(E^{[l]}\) multiply by \(h^l/l!\), and
\(F^{[l]}=(\sum_i\partial_{X_i}\partial_{Y_i})^l/l!\).
The harmonic and level projectors are
\[
 H_{a,b}=\sum_{l=0}^{\min(a,b)}
 \frac{(-1)^lE^{[l]}F^{[l]}}{\binom{a+b+k-2}{l}},
 \qquad
 P_j=\frac{E^{[j]}H_{a-j,b-j}F^{[j]}}{\binom{a+b+k-j-1}{j}} .
 \tag{5.20.19j}\label{eq:p5c2-harmonic-projections}
\]
The commutator \([F,h]=k+\deg_X+\deg_Y\) gives
\(FH_{a,b}=0,\ H_{a,b}\equiv1\bmod h\). On a harmonic
\(v\) of bidegree \((a-i,b-i)\) it gives
\(F(h^iv)=i(a+b+k-i-1)h^{i-1}v\). Hence \(P_j(h^iv)=\delta_{ij}h^iv\).
Reducing modulo \(h\) and iterating proves \(\sum_jP_j=1\)
and gives the harmonic dimension \(d(a,b)\).

With \(L=\lfloor\log_p(a+b+k)\rfloor\), the two divided
operators preserve the lattice with basis \(X^\alpha
Y^\beta/\alpha!\). The projectors cost at most \(2L\)
there by the binomial bounds. The inclusions between
that lattice times \(a!\) and the ordinary lattice cost
at most \((k-1)L\) by the multinomial bound. Thus
\(p^{(k+1)L}\) clears the projectors in ordinary
coordinates, including in every dual edge basis since \(h,F\)
are invariant. Use \(i_j[f]=h^jH_{a-j,b-j}f\) and
\(r_jx=[P_jx/h^j]\). Division by \(h^j\) on polynomial
multiples preserves integrality because its leading
coefficient is a unit. The tensor maps satisfy \(\sum_j i_jr_j=1\). By
\(a+b+1\le8a_1\), their denominators in both directions
are at most \(p^u\) for \(u=Nu_0\).
The summand dimension is \(d_j^{\rm harm}=\prod_s d(A_s-j_s,B_s-j_s)\).
\leanref{Formalism.\allowbreak{}MainNotions.\allowbreak{}flipOneThree.\allowbreak{}Pf.\allowbreak{}oneTop\_\allowbreak{}abs}{Formalism.\allowbreak{}SharpRestriction.\allowbreak{}GrowthBounds.\allowbreak{}ColumnSpanDensity.\allowbreak{}HarProjEqZeroIff}{87--225}

Let \(A_H=\sum_{s\ {\rm even}}A_s+\sum_{s\ {\rm odd}}B_s\),
\(A_j=A_1\) on the summand, and \(d_j=A_H-A_j\ge0\).
Use \(r'_j=\omega^{d_j}r_j,\ i'_j=\omega^{-d_j}i_j\).
For the tensor pairs \(x_a,y_b\) use in the summand
\(x_{ja}=p^ur'_jx_a\) and \(y_{jb}=p^{2u+d_j}y_b\circ i'_j\).
A lattice exponent \(-A_H\), or \(-A_H+c_4\) on a competing
edge, becomes \(-A_j\), or \(-A_j+c_4\), respectively.
Indeed \(p^{d_j}\omega^{-d_j}=\omega^{d_j}\) on the
dual side. Depth \(t_1\) persists because \(d_j\ge0\).
\leanref{Formalism.\allowbreak{}MainNotions.\allowbreak{}flipOneThree.\allowbreak{}Pf.\allowbreak{}lattice\_\allowbreak{}maps\_\allowbreak{}of\_\allowbreak{}functor\_\allowbreak{}comparison}{Formalism.\allowbreak{}SharpRestriction.\allowbreak{}GrowthBounds.\allowbreak{}ColumnSpanDensity.\allowbreak{}LatticeMapsOfFunctorComparison}{100--220}

For \(M_j=(y_{jb}(p^{A_j}x_{ja}))_{b,a}\), cancellation
of the \(\omega\)-powers and \(A_j+d_j=A_H\) give
\(\sum_jM_j=p^{3u}zI_n\). Apply \ref{B5.20k}. It gives
\(n_j\) pairs in each summand, \(\sum_jn_j\ge n\), with
diagonal valuation at most \(D_0+3u\le D_1\).
\leanref{Formalism.\allowbreak{}MainNotions.\allowbreak{}flipOneThree.\allowbreak{}Pf.\allowbreak{}rho\_\allowbreak{}sum\_\allowbreak{}bound\_\allowbreak{}of\_\allowbreak{}maps\_\allowbreak{}splitting}{Formalism.\allowbreak{}SharpRestriction.\allowbreak{}GrowthBounds.\allowbreak{}ColumnSpanDensity.\allowbreak{}RhoSumBoundOfMapsSplitting}{88--211}

Transport these pairs to the \(E_q\) outgoing edges of
one source. After division by block diagonals their matrix
is the identity modulo \(\omega\) since \(2D_1<c_4\).
Hence \(E_qn_j\le d_j^{\rm harm}\). If all \(j_s\le J_b\),
each term \(\min(B_y,A_x/p^d)\) in \(J_{\rm odd}\) decreases
under the shift by at most \(J_0\), also across the cyclic
boundary. Hence \(t_1\le t_*(A-j,B-j)\) and
\(E_qn_j\le(1-f(A-j,B-j))d_j^{\rm harm}\). Since
\(\sum_jd_j^{\rm harm}=\prod_sP(A_s)P(B_s)\) by
telescoping, \eqref{eq:p5c2-tensor-rank} gives the
pointwise inequality of \ref{B5.20g}.
\leanref{Formalism.\allowbreak{}MainNotions.\allowbreak{}flipOneThree.\allowbreak{}Pf.\allowbreak{}keyIneq\_\allowbreak{}of}{Formalism.\allowbreak{}SharpRestriction.\allowbreak{}GrowthBounds.\allowbreak{}ColumnSpanDensity.\allowbreak{}FactorizationBoundMonoLattices}{169--285}
\leanref{Formalism.\allowbreak{}MainNotions.\allowbreak{}flipOneThree.\allowbreak{}Pf.\allowbreak{}rho\_\allowbreak{}upper\_\allowbreak{}bound\_\allowbreak{}of\_\allowbreak{}edge\_\allowbreak{}count}{Formalism.\allowbreak{}SharpRestriction.\allowbreak{}GrowthBounds.\allowbreak{}ColumnSpanDensity.\allowbreak{}RhoUpperBoundOfEdgeCount}{106--205}

\paragraph{Middle rows.}
Write \(\eta=12Nkc_b/b'\). When the lower shift leaves the
box, rank zero suffices. Otherwise let \(V_-,V_+\) have the
flat extremes \((A_s-c_b,-B_s+c_b)\) and
\((A_s+c_b,-B_s-c_b)\), with respective marked dimensions
\(d_-,d_+\). Let \(d_{\rm mid}\) be the product of the
middle-row dimensions in the general dominant weight
\(\lambda\). The Weyl formula gives
\[
 d_{\rm mid}d_-\le\dim V_{\rm marked}\le E_cd_{\rm mid}d_-,
 \qquad d_+\le E_cd_-,\qquad
 E_c=(1+6c_b/b')^{(2k-3)N}.
\]
At each mark the varying root factors are the \(2k-3\) factors
meeting an extreme row. Each varies by at most \(6c_b/b'\)
of its lower value. Since \(E_c\le e^\eta\), \(\eta\le1/9\),
we have \(E_c\le2\) and \((1-\eta)E_c\le1\).

Apply \ref{B5.20s} at each mark with \(c=c_b\), \(P=A_s-c_b\), \(Q=B_s-c_b\). On the lower-shift domain both degrees are nonnegative. The size hypotheses with the stated large constants give \(k\ge10\), \(2c_b\le8a_1\). At odd marks use reversed dual coordinates. Tensor the maps, indexing a middle direction by a tuple of local Smith indices. An off-diagonal tuple has a zero factor, and the diagonal exponents add. The maps on the marked lattices are \(i:V_-\to V_+\),
\(U_b:V_-\to V_\lambda\), and \(L_a:V_\lambda\to V_+\).
Here \(i\) multiplies by the extra extreme columns,
and the \(U_b\) insert the middle tableau indexed by \(b\).
Use the corresponding dual middle basis for \(L_a\),
clearing its divided-power denominator. In column
coordinates the paired middle basis contracts to
\[
 L_aU_b=\delta_{ab}p^{k_b}i,\qquad
 0\le k_b\le k_{\rm mid}=N\lceil r^2\ell\rceil .
\]
For the frame calculation let \(g\) carry the standard frame to the selected frame and let \(D_t=\operatorname{diag}(t^{-1},1,\ldots,1)\). An inserted section of determinant twist \(-b\) and first-row degree \(a+b\) is transported as \(\det(g)^b g^{-1}v\), with matrix actions on the rows. Checking \(t^a\) times this section in the \(gD_t\) frame gives the scalar \((\det(g)t^{-1})^{-b}t^a\det(g)^b t^{-(a+b)}=1\). For the two lifted middle factors take \(a=b=c_b\); for the multiplier of \(i\) take \(a=b=2c_b\). Use \(t=1,\omega,\omega^2\) at source, selected edge and target. Hence the multipliers for \(i,U_b\) are integral at the source and in the normalized selected edge; the source lattice is contained in every normalized outgoing lattice for these dominant residual weights. Multiplication therefore preserves the competing-edge order as well as the source depth. Saturation at the selected edge follows from \ref{B5.20s}, giving an integral retraction there. For a target functional \(g_+\) take \(p^d g_+\circ L_a\), where \(d=A_1(V_+)-A_1(V_\lambda)=Nc_b\). The residual factor for \(L_a\) is integral at the selected normalized edge, and multiplying it by \(\omega^{2d}\) puts it in the target vertex lattice. This vertex lattice lies in each normalized incoming residual lattice: in an adapted frame every weight of that factor is between its two extremes, which are \(c_b,-c_b\) at each mark. Bilinearity in the common frame now gives target depth \(t\), selected-edge integrality, and competing-edge order \(c_2-2d\) from an upper test of order \(c_2\). Finally \([v,p^dg_+\circ L_a]_\lambda=[L_av,g_+]_+\) by the definition of the normalized pairing.

Take \(n_-,n_+\) pairs with scalar diagonals \(z_-,z_+\)
at the lower and upper flat weights. Push the lower
source family by \(i\) and pull its test family back by
the retraction, with the normalization powers, to the
upper edge lattice. Let \(M\) be the matrix pairing the
pushed family with the upper test family. In Smith bases
retain the \(h\) entries dividing \(z_-z_+\), and scale
them to that scalar. The \(s=n_--h\) other rows have
\(M_{\rm bad}\in z_-z_+\mathfrak m\). Pair these \(s\)
source directions, together with the \(n_+\) upper
source directions, against the retracted tests and the
upper tests. The block matrix has diagonals \(z_-I_s,z_+I\)
and off-diagonal blocks \(M_{\rm bad},B\) with \(B\)
integral. Eliminating the first diagonal gives a Schur
complement \(z_+I-BM_{\rm bad}/z_-\), which is \(z_+\)
times a unit matrix. Multiplication by \(z_-z_+\) clears
the inverse of the block matrix. The \(s+n_+\) resulting
directions have the upper source conditions and tests in
the distinguished upper edge dual. On the \(E_q\) outgoing
edges their normalized matrix reduces to the identity
because \(4D_1<c_4\). Thus \(s+n_+\le d_+/E_q\) and hence
\(h\ge n_-+n_+-d_+/E_q\). Apply every \(U_b\) to each
retained source and the corresponding \(L_a\) to its test.
The displayed composition identity, with the additional
factor \(p^{k_{\rm mid}-k_a}\) on the test, gives
\(d_{\rm mid}h\) pairs of valuation at most \(2D_1+k_{\rm mid}\).
\leanref{Formalism.\allowbreak{}MainNotions.\allowbreak{}flipOneThree.\allowbreak{}Pf.\allowbreak{}cloud\_\allowbreak{}count}{Formalism.\allowbreak{}SharpRestriction.\allowbreak{}GrowthBounds.\allowbreak{}ColumnSpanDensity.\allowbreak{}CloudCount}{30--186}
\leanref{Formalism.\allowbreak{}MainNotions.\allowbreak{}flipOneThree.\allowbreak{}Pf.\allowbreak{}cloud\_\allowbreak{}prop\_\allowbreak{}of\_\allowbreak{}maps\_\allowbreak{}retraction\_\allowbreak{}basis}{Formalism.\allowbreak{}SharpRestriction.\allowbreak{}GrowthBounds.\allowbreak{}ColumnSpanDensity.\allowbreak{}CloudPropOfMapsRetractionBasis}{36--54}
\leanref{Formalism.\allowbreak{}MainNotions.\allowbreak{}flipOneThree.\allowbreak{}Pf.\allowbreak{}latRetract}{Formalism.\allowbreak{}SharpRestriction.\allowbreak{}GrowthBounds.\allowbreak{}ColumnSpanDensity.\allowbreak{}CloudPropOfMapsBasis}{8--101}

The Weyl inequalities now give loss at most
\(f_-+2f(A+c_b,B+c_b)+\eta\). The denominator is at most \(D\)
by \eqref{eq:p5c2-rank-one-size} and \eqref{eq:p5c2-rank-one-budgets},
since \(N_1^2\le k\) and there are at most \(N\) segments.
Off-mark factors cost at most \(2n'c_b\) in depth and
order. The marked pullbacks leave order \(c_4-2Nc_b-2n'c_b=c_*\).
The upper shift in \(J_{\rm odd}\) costs at most \(Nc_b\).
Substitution in \(t_*(A+c_b,B+c_b)+2n'c_b\) gives
\eqref{eq:p5c2-rank-one-depth} after enlarging \(C\).
Apply \ref{B5.20e}.
\leanref{Formalism.\allowbreak{}MainNotions.\allowbreak{}flipOneThree.\allowbreak{}Pf.\allowbreak{}midLP\_\allowbreak{}of}{Formalism.\allowbreak{}SharpRestriction.\allowbreak{}GrowthBounds.\allowbreak{}ColumnSpanDensity.\allowbreak{}LatticeCountPrimeMonoLevel}{183--276}
\end{proof}

\subsection*{Proof of \ref{B5.20aa}}\label{prf:B5.20aa}\phantomsection
\begin{proof}
Write \(N=|T|\), \(n'=n_p\), and let \(T_{\rm ev}\) be the even marks. For \(s\in T\) put
\[
 M_s=\sum_{\substack{h\ge1\\h\ {\rm odd}}}
       \min\{a_s/p^h,W_j(z_h(s))\},
 \qquad \gamma_p=\frac{p}{p^2-1}.
\]
By \ref{B5.17a} the capacities \(W_j\) are at least \(w''\ge1\) at these scales. Both series are bounded termwise by the geometric series with terms \(a_s/p^h\). The identity \([u-w]_++\min(u,w)=u\) gives \(A_s+M_s=\gamma_pa_s\).

The numerator and capacity in \ref{B3.49} are \(\bar\alpha_s=a_s-2c\) and \(\bar C_{s,h}=W_j(z_h(s))-\kappa_j+c\). The size conditions of \ref{def:p3c5-size} give \(0\le\kappa_j\le N_m\le S_c\) and \(2c\le a_s\) eventually. Thus \(0\le\bar C_{s,h}\le W_j(z_h(s))+c\). At an odd mark, replacing \(a_s,W_j(z_h(s))\) by \(\bar\alpha_s,\bar C_{s,h}\) decreases any positive-part summand by at most \(2c/p^h+c\le3c\). A decrease requires a positive original summand, hence \(p^h<a_s\le2a\) by \eqref{eq:p3c5-top-window}. There are at most \(z\) such positive integers \(h\). At an even mark the corresponding decrease in \(A_s+M_s\) is \(2c\gamma_p\le 3cz\), since \(p<X'=a/b'\le a\). Summing the two estimates and using the definition of \(g_U\) gives
\[
 \sum_{s\in T}A_s+\sum_{s\in T_{\rm ev}}M_s
 \le g_U(t_0,j)+\bar\delta(t_0)+3Ncz .
 \tag{5.20.20}\label{eq:p5c2-capacity-perturbation}
\]
\leanref{Formalism.\allowbreak{}MainNotions.\allowbreak{}AuxBlock8.\allowbreak{}Pf.\allowbreak{}St.\allowbreak{}gU\_\allowbreak{}lower}{Formalism.\allowbreak{}SharpRestriction.\allowbreak{}GrowthBounds.\allowbreak{}LogarithmicSums.\allowbreak{}PartnerCountOfCentral}{308--344}

Put \(G_1=C_{31}L/\delta_1\), \(U=40n'\delta_1\), \(N_1=N_{\rm win}\), \(q_p=p^{G'}\), and \(z_a=\log_p a\). By \ref{B5.20j},
\[
 N\ge1,\qquad n'\le NG_1,\qquad U\le40NC_{31}L,\qquad N_1\le C_NL .
\]
All coefficients below are the nonnegative family coefficients of \ref{B3.50} and \ref{B3.52a}. We use one common bound
\[
 Y_0=\exp(x^{23\varepsilon/48}).
\]
For sufficiently large \(x\), each of \(G_1,U,n',N_1,k,N_m,L,z_a,x^{13},c,\eta,q_p\) and each fixed coefficient is at most \(Y_0\), uniformly in \(p<X'\). Indeed \(n'\le n\), and \ref{def:p3c5-size} bounds \(n,k,m,t\) by fixed powers of \(x\). Also \(z_a\le\log_2 a\), \(c=16k\Delta\), \(\log\Delta\le x^{\varepsilon/4}/2\), and \(\log\eta\le x^{\varepsilon/4}\). Finally \eqref{eq:p3c5-window} gives \(G'=O_\varepsilon(x^{\varepsilon/8}L)\). Thus
\[
 \log q_p\le G'x^{\varepsilon/3}
       =O_\varepsilon(x^{22\varepsilon/48}L)
       \le x^{23\varepsilon/48}
\]
eventually. Each fixed power of \(Y_0\) is \(o(b')\), since
\(\log b'=L+x^\varepsilon-x^{\varepsilon/3}\).

First suppose \(p<p_0\). Write \(r=r_G\) and put
\[
 M=w''+\eta+c+N_m,\qquad
 D'=K_{\rm DII}(N_1^3+r^5+k^2)z_a .
\]
We have \(kq_p,D'\le b'\) by the preceding bounds. Also \(G_1M\le b'\) eventually. Indeed \(G_1w''/b'=C_{31}L/(xn\delta_1)\to0\), and \(G_1(\eta+c+N_m)\le3Y_0^2=o(b')\). In particular \(n'M\le Nb'\). With \(c_0=2b'+2c\le3b'\), the small-prime formulas of \ref{B3.50} and \ref{B3.52a} now give
\[
 \begin{aligned}
 P_{v,j}&\le 6rNb',\\
 \lceil t_v\rceil_+
   &\le K_{\rm tII}rN(H_-+4b')+2rNb'+1,\\
 D_p&\le40NC_{31}Lb'+2rNb'+120rNC_{31}Lb'+1 .
 \end{aligned}
\]
In the last line the three terms before rounding bound \(UD'\), \(2rn'c\), and \(rUc_0\), respectively. Since \(N,r,b',L\ge1\), adding these lines and the unit in \eqref{eq:p5c2-selected-comparison} gives
\[
 P_{v,j}+\lceil t_v\rceil_++D_p+1
 \le N(C_*r b'L+C_*LH_-)
\]
as soon as \(C_*\ge K_{\rm tII}r\) and \(C_*\ge13+4K_{\rm tII}+160C_{31}\). Now use \eqref{eq:p5c2-capacity-perturbation} and \(M_s\ge0\).

Suppose next that \(p_0\le p<X'\). By \ref{B5.20j} the cycle has even length and contains an even mark. For each odd mark \(y\), let \(h\) be the positive backward distance to the preceding even mark \(s\). Then \(h\) is odd and \(z_h(s)=(0,y)\) in orbit coordinates. The pair \((s,h)\) determines \(y\), so different odd marks use different summands of \(\sum_{s\in T_{\rm ev}}M_s\). By \ref{B3.47},
\[
 \min(b_y,a_s/p^h)\le\min(W_j(0,y),a_s/p^h),
 \qquad B_j(0,y)-b_y\le b'/x+\eta .
\]
The first inequality bounds the contribution of \(y\) to \(J_{\rm odd}\) in \(t_v\). The second bounds its contribution to the marked part of \(P_{v,j}\). Their sum over the odd marks is at most \(\sum_{s\in T_{\rm ev}}M_s+N(b'/x+\eta)\).
\leanref{Formalism.\allowbreak{}MainNotions.\allowbreak{}AuxBlock8.\allowbreak{}Pf.\allowbreak{}St.\allowbreak{}kinkHH}{Formalism.\allowbreak{}SharpRestriction.\allowbreak{}GrowthBounds.\allowbreak{}LogarithmicSums.\allowbreak{}PartnerCountOfCentral}{369--439}

To bound the remaining terms, put
\[
 E=K_{\rm DI}(x^{13})^2(N_1^5+k^3)z_a .
\]
By \ref{B3.52a}, \(D_p=D_I\le UE+1\) and
\[
 c_6=8(D_I+1)+K_2Ukz_a+4Uc+2n'c .
\]
In the price, \(K_{\rm Pr}n'w''\le NC_*w''L/\delta_1\) for \(C_*\ge K_{\rm Pr}C_{31}\). In the depth, \(K_1NN_1H_+\le NC_*LH_+\) for \(C_*\ge K_1 C_N\). After dividing by \(N\), the sum of the remaining costs is at most
\[
 K_{\rm Pr}G_1(\eta+c+N_m)
 +K_2(q_pc_6+kq_p+kz_a+c)+2G_1c+40C_{31}LE+3 .
\]
Here \(n'/N\le G_1\) and \(U/N\le40C_{31}L\). The final \(3\) covers the two ceiling errors and the unit in the asserted inequality, since \(N\ge1\).

Write \(Y=Y_0\). Replacing the nonnegative factors by their upper bounds gives
\[
 E\le Y^4(Y^5+Y^3),\qquad
 D_I\le d(Y):=Y^5(Y^5+Y^3)+1,\qquad
 c_6\le h(Y):=8(d(Y)+1)+Y^4+6Y^2 .
\]
Consequently the preceding cost per mark is at most
\[
 \begin{aligned}
 &3Y^3+Y\{Yh(Y)+2Y^2+Y\}+2Y^2
       +40Y^2Y^4(Y^5+Y^3)+3\\
 &\quad=8Y^{12}+40Y^{11}+8Y^{10}+40Y^9+Y^6
             +6Y^4+5Y^3+19Y^2+3 .
 \end{aligned}
\]
The coefficients sum to \(130\), so for \(Y\ge130\) this is at most \(130Y^{12}\le Y^{13}\). Since \(13x^{23\varepsilon/48}\le x^{\varepsilon/2}\) eventually, \(Y^{13}\le b_0\). This gives coefficient one for \(b_0\). Combining the price and depth bounds yields
\[
 P_{v,j}+\lceil t_v\rceil_++D_p+1
 \le \sum_{s\in T_{\rm ev}}M_s+
 N\bigl(b'/x+\eta+C_*w''L/\delta_1+b_0+C_*LH_+\bigr).
\]
Substitution in \eqref{eq:p5c2-capacity-perturbation} proves \eqref{eq:p5c2-selected-comparison}, since \(3cz\le C_*(c+\eta)z\).
\leanref{Formalism.\allowbreak{}MainNotions.\allowbreak{}AuxBlock8.\allowbreak{}Pf.\allowbreak{}St.\allowbreak{}B.\allowbreak{}hboB\_\allowbreak{}arith}{Formalism.\allowbreak{}SharpRestriction.\allowbreak{}GrowthBounds.\allowbreak{}LogarithmicSums.\allowbreak{}StarNumB}{98--130}
\leanref{Formalism.\allowbreak{}MainNotions.\allowbreak{}AuxBlock8.\allowbreak{}Pf.\allowbreak{}St.\allowbreak{}B.\allowbreak{}facts\_\allowbreak{}ev}{Formalism.\allowbreak{}SharpRestriction.\allowbreak{}GrowthBounds.\allowbreak{}LogarithmicSums.\allowbreak{}StarNumB}{159--356}
\end{proof}

\subsection*{Proof of \ref{B5.20a}}\label{prf:B5.20a}\phantomsection
\begin{proof}
The duties of \ref{B8.02} identify the cyclic coordinates of a selected representative with those of the twist. The rows are dominant and their unspecified entries have absolute value at most \(\lfloor c\rfloor_+\), where \(c=4S_c\). Below \(p_0\) use \ref{B5.20b}. Above or at \(p_0\) use the attained minimum of \ref{B3.52a}. The coefficient-ring map of the design preserves the Frobenius frames and the chosen square root of \(p\). The tableau generators in each frame therefore extend to generators of the corresponding lattice over the model ring \(O\), including the intermediate edge lattice. Source vectors and integral tests extend by scalars, and their pairing matrix extends to the same scalar identity over \(O\). This preserves the rank and the inclusions defining the test.

Below \(p_0\), \ref{B3.51} bounds the share by \(C_{\rm frac}N_vr_G^2k/x^{26}\). Since \(r_G\) is fixed and \(k\le x^6\) eventually by \ref{def:p3c5-size}, this is at most \(N_vx^{-12}\). Above or at \(p_0\), \ref{B5.20c} gives the mean bound \(C_mN_vN_{\rm win}/x^{13}\), and \ref{B5.20j} gives \(C_mN_{\rm win}\le x\) eventually. The selection rule \eqref{eq:p3c5-selection} includes all good representatives of a candidate with \(\Sigma_p\le x^{-\varepsilon}\).
\leanref{Formalism.\allowbreak{}L2.\allowbreak{}G.\allowbreak{}Model.\allowbreak{}lpAt\_\allowbreak{}dom}{Formalism.\allowbreak{}LocalToGlobal.\allowbreak{}GroupLattices.\allowbreak{}FockModels.\allowbreak{}LPBase}{256--384}
\leanref{Formalism.\allowbreak{}MainNotions.\allowbreak{}G.\allowbreak{}Joints.\allowbreak{}linear\_\allowbreak{}programming\_\allowbreak{}input\_\allowbreak{}of\_\allowbreak{}hypotheses}{Formalism.\allowbreak{}SharpRestriction.\allowbreak{}GrowthBounds.\allowbreak{}Joints}{2040--2127}

Fix \(u,w\) with \(p=p(w)\ne T_u\). For nonempty selection, \ref{B4.30} gives a permitted place, and the orbit identification gives distinct positions for the selected representatives. The definite level stabilizers are finite. Each embeds into pro-groups at two distinct primes by the two unipotent level conditions, so each has order one. The neighbor sums at distinct positions commute and are integral by \ref{B4.07c}. Put \(E_p=\lceil C_e f k^2\log_p a\rceil_+\). By \ref{B3.52b} the level-zero lattices satisfy \(p^{E_p}M_{\rm na}\subseteq M\subseteq M_{\rm na}\). The intersection \(R=M\cap V^{\rm ram}\) is saturated in \(M\). We have \(m\ge2\) and \(m\sum_{v\in V_p}\kappa_v=\Sigma_p\le x^{-\varepsilon}\le1/2\) at the selected scales. These are the lattice and share hypotheses of \ref{B4.07a}.
\leanref{Formalism.\allowbreak{}L2.\allowbreak{}G.\allowbreak{}gatEF\_\allowbreak{}of\_\allowbreak{}tied}{Formalism.\allowbreak{}LocalToGlobal.\allowbreak{}GroupLattices.\allowbreak{}LinkGConcreteNoKer}{20--27}
\leanref{GlobalTheory.\allowbreak{}Mem.\allowbreak{}Top.\allowbreak{}LayOfCore.\allowbreak{}ledOK}{Development.\allowbreak{}AbelianVarietyTheory.\allowbreak{}Proofs.\allowbreak{}IwasawaDecomposition.\allowbreak{}ExistsLayerParametersAdmissible}{65--77}

Write \(A=\mathcal O_{K,(w)}\) and \(e_O=v_O(p)\), with \(v_O\) normalized by a uniformizer of \(O\). Under the initial inclusion and the comparison of \ref{B4.07c} and \ref{B4.21a}, the extensions of \((M_{H,u})_{(w)}\) and \((N_u)_{(w)}\) are \(R\) and the enhanced lattice \(N_P\). Indeed localization commutes with intersection, and \((N_u)_{(w)}\subseteq(M_{H,u})_{(w)}\) off \(T_u\). If \(\pi\) is a uniformizer of \(A\), then \(p\) is a unit times \(\pi^{e_w}\), so \(v_O(\pi)=e_O/e_w\). Flat extension carries each composition factor of the \(A\)-quotient to \(O/\pi O\). Hence
\[
 \operatorname{length}_O(R/N_P)
 =\frac{e_O}{e_w}\operatorname{length}_A
       ((M_{H,u})_{(w)}/(N_u)_{(w)})
 =e_O\ell_{u,w}.
\]
Dividing the bound of \ref{B4.07a} by \(e_O\) gives
\(\ell_{u,w}\le d_M(\tau_w+\mathrm{io}_w)\). For containment, put \(b_\tau=\lceil\tau_p\rceil_+\). By \eqref{eq:p4c2-np} the image of \(p^{-b_\tau}x'\) lies in \(M\) for every \(x'\in(N_u)_{(w)}\). The coordinate test of \ref{B4.07c} puts \(p^{-b_\tau}x'\) in \((M_{H,u})_{(w)}\), and \(e_wb_\tau\ge\lfloor e_w\tau_w\rfloor_+\) by \ref{B3.52}. For empty selection the two lattices coincide by \ref{B4.21a}, and both budgets vanish by \ref{B3.52}.
\leanref{Formalism.\allowbreak{}BasicNotions.\allowbreak{}LatCmp.\allowbreak{}NP\_\allowbreak{}eq\_\allowbreak{}span}{Formalism.\allowbreak{}LocalRingCardinality.\allowbreak{}Lemmas}{990--1026}
\leanref{Formalism.\allowbreak{}BasicNotions.\allowbreak{}LatCmp.\allowbreak{}length\_\allowbreak{}quotient\_\allowbreak{}of\_\allowbreak{}nonempty}{Formalism.\allowbreak{}LocalRingCardinality.\allowbreak{}Lemmas}{1061--1117}
\leanref{Formalism.\allowbreak{}MainNotions.\allowbreak{}G.\allowbreak{}nuServed\_\allowbreak{}of}{Formalism.\allowbreak{}SharpRestriction.\allowbreak{}GrowthBounds.\allowbreak{}PlaneCountTheoremWithConstants}{664--709}

For the upper bounds on gains, \(\gamma_p=p/(p^2-1)\le4/(3p)\) and \(0\le\bar\alpha_s\le6a/5\) by \eqref{eq:p3c5-top-window} and \ref{def:p3c5-size}. Both \(g_0(v,j)\) and \(g_U(v,j)\) are at most \(\gamma_p\sum_{s\in T(v)}\bar\alpha_s\) by \ref{B3.48} and \ref{B3.49}. Here \(T(v)\) denotes the fibre of the representative of \(v\). The subtracted minima and the logarithmic error are nonnegative. The price in \ref{B3.50} is also nonnegative, since at a mark \(B_j(0,s)-b_s\ge b'/x-\kappa_j\ge0\) by \ref{B3.47} and \(N_m\le S_c\le w''\le b'/x\). Each maximum in \eqref{eq:p3c5-gain} is therefore between zero and \(8N_va/(5p)\). Summing over the disjoint fibres gives the asserted bound \(4n_Ta/p\). For \(S_v\), \ref{B3.52} takes the ceiling of a sum bounded by \(8mN_va/(5p)\). If \(mN_v=0\) this ceiling is zero. Otherwise selection implies \(p<X'\) and \(a/p\ge b'\ge500\) eventually, so \(8mN_va/(5p)+1\le2mN_va/p\).
\leanref{Formalism.\allowbreak{}MainNotions.\allowbreak{}AuxBlock8.\allowbreak{}Pf.\allowbreak{}orbit\_\allowbreak{}max\_\allowbreak{}le}{Formalism.\allowbreak{}SharpRestriction.\allowbreak{}GrowthBounds.\allowbreak{}LogarithmicSums.\allowbreak{}NotSingularOfCentral}{469--483}
\leanref{Formalism.\allowbreak{}MainNotions.\allowbreak{}G.\allowbreak{}svBound\_\allowbreak{}of}{Formalism.\allowbreak{}SharpRestriction.\allowbreak{}GrowthBounds.\allowbreak{}PlaneCountTheoremWithConstants}{409--546}

Fix a signal or inert-class prime \(p\le X''\). Let \(g^{\rm count}_{p,j}\) be the counted gain of \ref{B3.51}. We first bound the total logarithmic error. Each fibre of \(\beta_{\Delta_p}\) in \(H\) has \(\operatorname{ord}(\Delta_p)\) elements by \ref{B3.46}. Thus \(\sum_{v\in\mathcal R_p}n_p\) is \(n\) in the inert class and \(2n\) in the signal class. Both are at most \(f\). Also \(\sum_{v\in\mathcal R_p}N_v=n_T\le f\). Write \(z_a=\log_p a\). Since \(p\le X''=a/w''\le a\) and \(a\ge2\), we have \(z_a\ge1\) and \(\log_p(2a)\le2z_a\). By \ref{B3.47},
\[
 \sum_{v\in\mathcal R_p}\bar\delta(v)
 \le 2K_{\delta,1} fN_m^2z_a+4K_{\delta,2} fz_a^2 .
\]
Moreover \(E_p\le(C_e+1)fk^2z_a\) by \ref{B3.52a}. Take \(C_*\) in \ref{B5.20aa} also at least \(2K_{\delta,1}+C_e+1\) and \(4K_{\delta,2}\). Since \(k\le N_m\), \eqref{5.17.1} with \(N_0=N_m\) gives, for every final coefficient \(C\ge C_*\),
\[
 E_p+\sum_{v\in\mathcal R_p}\bar\delta(v)\le F_p .
\]
\leanref{Formalism.\allowbreak{}MainNotions.\allowbreak{}AuxBlock8.\allowbreak{}Pf.\allowbreak{}sum\_\allowbreak{}dl\_\allowbreak{}add\_\allowbreak{}ep\_\allowbreak{}le}{Formalism.\allowbreak{}SharpRestriction.\allowbreak{}GrowthBounds.\allowbreak{}LogarithmicSums.\allowbreak{}NotSingularOfCentral}{326--350}

Suppose \(V_p\ne\varnothing\). By \ref{B3.51} this is an inert-class candidate with \(p<X'\) and \(\operatorname{ord}(h_p)\ge x^{\varepsilon/4}>1\). Thus the maximum branch of \eqref{eq:p3c5-gain} applies. At a selected orbit its contribution is
\[
 \max\{g_U(v,j)-P_{v,j},\,g_0(v,j),\,0\}
 \ge g_U(v,j)-P_{v,j},
\]
where \(P_{v,j}\) is the global price defined in \ref{B5.20aa}. Each unselected orbit contributes \([g_0(v,j)]_+\ge0\). Hence
\[
 \sum_{v\in V_p}\bigl(g_U(v,j)-P_{v,j}\bigr)
 \le g^{\rm count}_{p,j}.
\]
By \ref{B3.52}, \(TT_p=\coprod_{v\in V_p}T(v)\). Put \(d_v=\lceil t_v\rceil_++D_p+1\). Summing \eqref{eq:p5c2-selected-comparison} over that partition, for the fixed lift \(j\), gives
\[
 \sum_{s\in TT_p}(A_{jsp}-\ell_p^{\rm int})
 \le g^{\rm count}_{p,j}
      +\sum_{v\in V_p}\bar\delta(v)-\sum_{v\in V_p}d_v .
\]
Here \(A_{jsp}=A_s\): the partner in \eqref{5.17.2} has weight \(W_j(z_h(s))\) by \ref{B5.17a} and \ref{B3.52}. All terms beyond the finite arithmetic cutoff vanish, since \(p^h>a_s\) gives \(a_s/p^h<1\le W_j(z_h(s))\).

There are \(m\) lifts. The depth costs sum to \(m\sum_{v\in V_p}d_v=\tau_p-E_p\) by \ref{B3.52}. Subtracting \(mF_p\) after summing over the lifts gives
\[
 \begin{aligned}
 \sum_{j=1}^m\sum_{s\in TT_p}(A_{jsp}-\ell_p^{\rm int})-mF_p
 &\le \sum_{j=1}^m g^{\rm count}_{p,j}-\tau_p
      +E_p+m\sum_{v\in V_p}\bar\delta(v)-mF_p\\
 &\le \sum_{j=1}^m g^{\rm count}_{p,j}-\tau_p .
 \end{aligned}
\]
The last line uses \(m\ge1\), \(E_p\ge0\), and
\(\sum_{v\in V_p}\bar\delta(v)+E_p\le F_p\).
\leanref{Formalism.\allowbreak{}MainNotions.\allowbreak{}AuxBlock8.\allowbreak{}Pf.\allowbreak{}sum\_\allowbreak{}TTP\_\allowbreak{}served}{Formalism.\allowbreak{}SharpRestriction.\allowbreak{}GrowthBounds.\allowbreak{}LogarithmicSums.\allowbreak{}NotSingularOfCentral}{601--609}
\leanref{Formalism.\allowbreak{}MainNotions.\allowbreak{}AuxBlock8.\allowbreak{}Pf.\allowbreak{}gain\_\allowbreak{}bound\_\allowbreak{}at\_\allowbreak{}of\_\allowbreak{}star}{Formalism.\allowbreak{}SharpRestriction.\allowbreak{}GrowthBounds.\allowbreak{}LogarithmicSums.\allowbreak{}NotSingularOfCentral}{619--729}

Take \(C=4C_*+1\) in \eqref{5.17.1} and \eqref{5.20.1}. At a candidate, \(C(S_c+k+\eta)\ge C_*(c+\eta)\). For \(p<p_0=k^2x^{26}\) we have \(pk^2x^{78}\le(kx^{26})^4\). Taking base-\(p\) logarithms and ceilings gives \(L_1(p)\le L_1^{\rm loc}\), so \(ap^{-L_1^{\rm loc}}\le ap^{-L_1(p)}\). In the other branch,
\[
 b'/x+\eta+C_*w''L/\delta_1+b_0
 \le C\{b'/x+\eta+S_c+(w''+t)L/\delta_1\}+b_0 .
\]
These comparisons and \(\Lambda=\mathcal L\) from \ref{B3.52} give \(\ell_p^{\rm int}\le\ell_p\) with the stated \(\ell_p\).

It remains to consider \(V_p=\varnothing\). Then \(\tau_p=0\). At an inert-class prime \(p<X'\), \(TT_p\) is empty by \ref{B3.52}, so nonnegativity of \(F_p\) and of the counted gain proves the inequality. In the other cases \(TT_p=TT\). At a signal prime every odd partner is on sheet one, so none of the odd exponents is excluded in \ref{B3.48}. At \(p\ge X'\), for \(h\ge1\) we have
\[
 a_s/p^h\le\bigl(a+a/(kx^{13})\bigr)/X'
              \le b'+b'/x .
\]
Every excluded partner is a mark, where \(W_j\ge b'+b'/x\) by \ref{B3.47}. Its summand in \(A_{jsp}\) is therefore zero. In particular, if \(p\) is inert-class with \(h_p=0\), every odd partner is the mark \((0,s)\), and both the prime sum and the counted gain vanish. In the remaining cases apply the positive-part comparison in the proof of \ref{B5.20aa} on the unexcluded odd exponents. For each orbit it gives
\[
 \sum_{s\in T(v)}(A_{jsp}-3c\log_p(2a))
 \le g_0(v,j)+\bar\delta(v)
 \le [g_0(v,j)]_++\bar\delta(v).
\]
Sum over all orbits. The positive parts sum to \(g^{\rm count}_{p,j}\) by \eqref{eq:p3c5-gain}, and the errors sum to at most \(F_p\). Since \(\ell_p\ge C_p\ge3c\log_p(2a)\), summing over the \(m\) lifts proves the same inequality with the stated losses.

Finally \ref{B3.52} identifies the counted gain with the gain on the supported fibre. At a served prime \ref{B4.04} leaves this gain in the sum, so \(g_{jp}=g^{\rm count}_{p,j}\). Taking \(r_p=F_p\) gives \eqref{5.16.1}.
\leanref{Formalism.\allowbreak{}MainNotions.\allowbreak{}G.\allowbreak{}ask\_\allowbreak{}eq\_\allowbreak{}askP}{Formalism.\allowbreak{}SharpRestriction.\allowbreak{}GrowthBounds.\allowbreak{}PlaneCountTheoremWithConstants}{223--245}
\leanref{Formalism.\allowbreak{}MainNotions.\allowbreak{}G.\allowbreak{}lo\_\allowbreak{}le\_\allowbreak{}loG}{Formalism.\allowbreak{}SharpRestriction.\allowbreak{}GrowthBounds.\allowbreak{}PlaneCountTheoremWithConstants}{102--206}
\leanref{Formalism.\allowbreak{}L2.\allowbreak{}GlobalTheoryTheory.\allowbreak{}Loss.\allowbreak{}loss\_\allowbreak{}bound\_\allowbreak{}at}{Formalism.\allowbreak{}LocalToGlobal.\allowbreak{}GlobalTheory.\allowbreak{}PowFourLtCeilThresholdPrime}{766--817}
\end{proof}

\subsection*{Proof of \ref{B5.20y}}\label{prf:B5.20y}\phantomsection
\begin{proof}
Take \(T_0,T_2,C_*\) from \ref{B5.20i}. Fix \(h_G=2+8\varepsilon\), \(r_G=\lceil8(h_G+2)/\varepsilon\rceil+1\), as in \ref{B3.51}. In \eqref{eq:p3c5-window} we have \(u_1\delta_1\le2\) and \(M_0\le(h_G+3)10^{6r_G+1}\log x+1\). Thus for \(\log x\ge1\) that formula gives \(N_{\rm win}\le A_\varepsilon\log x\), where one may take
\[
 A_\varepsilon=2+\max\!\left\{
 12\exp(1)(6r_G+1)\big((h_G+3)10^{6r_G+1}+6\big),
 (h_G+3)/\log 2\right\}.
\]
Require every middle entry to be at least \(\max(1,T_0)\). For a concatenated outer and middle list \(\ell\), take the last threshold to be
\[
 \max\bigl(\{A_\varepsilon\}\cup
     \{\,2T_2(s),\,2C_*(s):s\in\ell,\ s\ge T_0\,\}\bigr).
\]
The pre-threshold and the unused threshold function can be zero. The entry \(C_0\) belongs to the middle list at \(CH\), so the last threshold bounds \(2T_2(C_0)\) and \(2C_*(C_0)\). Applied to the depth, denominator, band and two loss coefficients of \ref{B5.20i}, these are exactly the lower bounds needed in \eqref{eq:p5c2-family-thresholds}. Applied to \(C_N\), the same test gives the window bound. All these choices precede the choice of a family. For a family passing the tests, the cutoff \(x\ge\max(\exp(1),x_1+1)\) ensures that the scale in the displayed window formula is the indexed scale \(x\).
\leanref{GlobalTheory.\allowbreak{}EX.\allowbreak{}RangeOne.\allowbreak{}rangeOne\_\allowbreak{}abs}{Development.\allowbreak{}AbelianVarietyTheory.\allowbreak{}Proofs.\allowbreak{}GlobalInterpolation.\allowbreak{}TensorOfPairLift}{637--640}
\end{proof}

\subsection*{Proof of \ref{B5.20j}}\label{prf:B5.20j}\phantomsection
\begin{proof}
Use \(k\le x^6\), \(p\ge2\), \(u_1=\lceil x^{\varepsilon/8}\rceil\)
in \(L_2\) and \eqref{eq:p3c5-window} for
\(G'=O_\varepsilon(x^{\varepsilon/8}\log x)\) and the stated
bound on \(N_1\). An odd orbit order would have trivial image
in the \(2\)-group and divide \(o\le x^{\varepsilon/100}\),
contrary to the candidate condition \(n'\ge x^{\varepsilon/4}\).
The run condition on full tiles in the good orbit gives even
gaps at most \(G'\), so \(n'\le NG'\) and there are at least
two gaps by \(3G'\le n'\). Cover the cycle with
\(\lceil n'/G'\rceil\) windows and use \(N_1\le30G'\delta_1\)
for the upper bound on \(N\).
\end{proof}

\subsection*{Proof of \ref{B5.20b}}\label{prf:B5.20b}\phantomsection
\begin{proof}
Use \ref{B5.20h} with \(Y=x^{26}\), \(a_0=a-c\),
\(c_b=\lfloor c\rfloor_+\), \(c_t=\lfloor c_0\rfloor_+\),
\(D=D_{II}\), \(c_*=8(D+1)\), and \(c=4S_c\).
By \ref{B5.20j}, \(N\ge x^{\varepsilon/8}/(C_{31}\log x)\),
so \(N_1^2\le rN\). Let \(C_{\rm rank}\) be the family coefficient in the rank requirement \(C_{\rm rank}N_1^2\log a\le k\). Since \(k\ge x^{9/2}\), both this requirement and \(N_1+3r\le k\) hold eventually.
The host split-prime threshold forces candidate \(p\ge13\).
Since \(N\le x^3,\ p<p_0\le k^2x^{26}+1\),
\[
 256k^2NY^2\le k^2x^{78}\le(kx^{26})^4/p\le p^{L_1^{\rm loc}-1}.
\]
The family thresholds bound \(C\) above by each of the following coefficients: the share coefficient \(C_{\rm frac}\), the depth coefficient in \(t_{II}\), the coefficient of \(N_1^2\log a\) in \(\mathcal L\), the denominator coefficient \(K_{\rm DII}\), the coefficient of \(k^3r^2Np^{5G'}(c_*+rNc_0+2rn'c)Y\) in the weight-size requirement, and \(C_{\rm rank}\).
Use \(N\le40n'\delta_1\) in \(D_{II}\).
Since \(D=O(x^{16}b')\) and
\(\log(p^{G'})=O_\varepsilon(x^{\varepsilon/8}\log^2x)\),
\(k^3r^2Np^{5G'}(c_*+rNc_0+2rn'c)Y=o(a)\) by \(a/b'=X'\).
At a mark, \(a_0+c_b\le a\le a_s\). By \eqref{eq:p3c5-top-window}, \(a_s\le6a/5\le2a_0-c_b\) eventually, since \(c=o(a)\). By \eqref{eq:p3c5-bottom-window}, \(0\le b_s\le5b'/4\), so the integer \(b_s+2c_b\) is at most \(\lfloor2b'+2c\rfloor=c_t\). Put \(z=\log_p a\) and \(U=40n'\delta_1\). The denominator hypothesis of \ref{B5.20h} follows from
\[
 CNk^2\log_p a_0\le K_{\rm DII}U(N_1^3+r^5+k^2)z\le D_{II}.
\]
For the depth, \(CN_1^2\log a_0\le\mathcal L\), \(a_0\le a\), \(c_t\le c_0\), and \(c_b\le c\). The expression in \eqref{eq:p5c2-plane-depth} is therefore at most
\[
 K_{\rm tII}rN\{\mathcal La/(pk)+ap^{-L_1^{\rm loc}}+c_0+kq_p\}+2rn'c=t_{II}.
\]
Thus the integer depth \(\lceil t_{II}\rceil_+\) is admissible. The witness has rank at least \((1-\min\{1,CNr^2k/Y\})\dim V\), which is at least the required rank because \(C\le C_{\rm frac}\).
\leanref{Formalism.\allowbreak{}MainNotions.\allowbreak{}G.\allowbreak{}Joints.\allowbreak{}thirteen\_\allowbreak{}le}{Formalism.\allowbreak{}SharpRestriction.\allowbreak{}GrowthBounds.\allowbreak{}SizesAtRecord}{163--183}
\end{proof}

\subsection*{Proof of \ref{B5.20c}}\label{prf:B5.20c}\phantomsection
\begin{proof}
Use \ref{B5.20j} and \(k^2Y^2\le p<X'\). Take
\(\Delta=b_{\rm top}(x)-b'\), \(c_b=\lfloor4S_c\rfloor_+\),
\(Y_b=x^\varepsilon/20\), \(D=D_I\), and take \(C_0\) to be the family coefficient in \(\mathcal L=C_0N_1^2\log a\), as in \ref{B3.51}.
Let \(C_{\rm I}\) be the family coefficient in the rank-one size requirement \(C_{\rm I}N\mathcal L^2\le k\), and let \(C_\rho\) be the coefficient of \(C_N(\log x)\mathcal L/k\) in the band width \(\rho_B\) of \ref{B3.51}. The family thresholds give
\[
 C_0\ge T_0,\qquad H=T_2(C_0)\le\min(C_{\rm I},K_{\rm DI},C_\rho/2),
 \qquad 2C_*(C_0)\le K_1,\quad C_*(C_0)\le K_2 .
 \tag{5.20.22}\label{eq:p5c2-family-thresholds}
\]
Use \(N\le x^3,\ k\ge x^{9/2}\) and
\(\mathcal L=O_\varepsilon(x^\varepsilon\log^2 x)\) for
\(2N_1+2\le k-1,\ 2Y_bN\mathcal L<k,\ HN\mathcal L^2\le k\).
The width of \(I_b\) is at least \(Y_bN\mathcal L b'/k\)
by \(N\le40n'\delta_1,\ n'\le n\). By \ref{B3.51},
\(\Delta\le\rho_B b'/4\) and the local band half-width is
at most \(\rho_B=o(1)\). The candidate band exclusion thus applies.
The remaining elementary size inequalities use \(X'=e^{x^{\varepsilon/3}}\)
and \(D_I\le x^{49}\), with
\(HNY^2(N_1^5+k^3)\log_p a\le D_I\). In particular
\[
 Hk^3Nq_p^4[8(D_I+1)+(N+n')c_b]Y=o(b').
\]
The logarithm of the left side is bounded by
\(O(\log x)+O(x^{23\varepsilon/48})+O(x^{\varepsilon/4})\)
from the bounds on \(q_p,c_b\), while \(\log b'\sim x^\varepsilon\).
Both integer windows are nonempty eventually. This gives
\eqref{eq:p5c2-rank-one-size}.

With the notation of \ref{B5.20i}, let \(m'=k-1\).
For \(0\le i\le J_b\) and \(u\ge a_1\),
\(P(u)/P(u+i)\ge1-m'J_0/a_1\). The binomial product gives
\[
 \frac{d(u,v)}{P(u)P(v+i)}\ge
 \frac{P(v)}{P(v+i)}-\frac{P(v)}{P(v+i+1)} .
\]
For \(v\in I_b+\tau\), \(\tau=0,c_b\), one has
\(v+J_b+1+m'\le2b'\) eventually, so
\(P(v)/P(v+J_b+1)\le e^{-L'/3}\).
Summing gives the density in \ref{B5.20g} with
\(\delta=(1-m'J_0/a_1)(1-e^{-L'/3})\).
For the interval sizes \(w_a,w_b\),
\[
 N(J_b+c_b)/w_b\le6/Y_b,\qquad
 N(J_b+c_b)/w_a\le6N\mathcal L Y/X',\qquad
 Nm'J_0/a_1\le NL'/X'.
\]
These tend to zero, as does \(Ne^{-L'/3}\). Thus the
cardinality inequality of \ref{B5.20g} holds eventually.
The downward boundary probability is at most
\(Nc_b(1/w_a+1/w_b)\). This and \(12Nkc_b/b'\) are at
most \(1/Y\) eventually, by the exponential size of \(b'\).
Therefore
\[
 \mathbb E_0(f_-+2f(A+c_b,B+c_b)+12Nkc_b/b')
 \le6(C_A+C_s)NN_1/Y+2/Y .
 \tag{5.20.23}\label{eq:p5c2-original-mean}
\]

Teichm\"uller expansions identify the source with
\(W(\mathbf F_{p^{n'}})\). The coefficient DVR has
unramified subring with residue field \(\mathbf F_{p^{2n'}}\),
and separating even and odd powers of \(\omega\) in its
expansion gives the full ring by adjoining \(\omega^2=p\).
Extend the smaller residue-field identification to the
larger one and carry \(\omega\) to its counterpart. This
transports the lattices and pairings, so the displayed
majorant holds for every local model and dominant weight.
\leanref{Formalism.\allowbreak{}MainNotions.\allowbreak{}flipOneThree.\allowbreak{}Pf.\allowbreak{}restOIso}{Formalism.\allowbreak{}SharpRestriction.\allowbreak{}GrowthBounds.\allowbreak{}ColumnSpanDensity.\allowbreak{}RestOIso}{8--12}
\leanref{Formalism.\allowbreak{}MainNotions.\allowbreak{}flipOneThree.\allowbreak{}Pf.\allowbreak{}theorem\_\allowbreak{}I\_\allowbreak{}of\_\allowbreak{}intermediate\_\allowbreak{}bounds}{Formalism.\allowbreak{}SharpRestriction.\allowbreak{}GrowthBounds.\allowbreak{}ColumnSpanDensity.\allowbreak{}TheoremIOfIntermediateBounds}{28--77}

Let \(C=C_*(C_0)\) be the coefficient of \eqref{eq:p5c2-rank-one-depth}, put \(U=40n'\delta_1\), and write \(z=\log_p a\). Since \(p<X'=a/b'\), the sum of the four terms in the first parenthesis of that depth bound is at most \(2(a/p^2+\mathcal La/(pk))\). Its contribution is bounded by the corresponding term of \(t_I\) because \(2C\le K_1\) by \eqref{eq:p5c2-family-thresholds}. In the second parenthesis, \(N\le U\) by \ref{B5.20j}, \(C\le K_2\), and \(c_b\le c=4S_c\). Thus
\[
 c_*+CNkz+4Nc_b+2n'c_b
 \le8(D_I+1)+K_2Ukz+4Uc+2n'c=c_6 .
\]
Multiplication by \(q_p=p^{G'}\) preserves this inequality. The terms \(kq_p+kz+c_b\) are at most \(kq_p+kz+c\), the outer coefficient \(CN\) is at most \(K_2N\), and the last term \(2n'c_b\) is at most \(2n'c\). Each odd mark uses the same least positive distance back to an even mark in both formulas, by the cyclic order in \ref{B5.20j}. This gives the required depth \(\lceil t_I\rceil_+\) with denominator \(D_I\) and order \(8(D_I+1)\). For each pair of weights the transported fraction is now admissible for every model and dominant weight in \ref{B3.52a}, so it bounds the least passing fraction from above. Average this pointwise inequality over the original box. Since \(NN_1\ge1\), \eqref{eq:p5c2-original-mean} gives the assertion with \(C_m=6(C_A+C_s)+2\), or any larger family constant.
\leanref{Formalism.\allowbreak{}MainNotions.\allowbreak{}G.\allowbreak{}Joints.\allowbreak{}tI\_\allowbreak{}le\_\allowbreak{}tIf}{Formalism.\allowbreak{}SharpRestriction.\allowbreak{}GrowthBounds.\allowbreak{}Joints}{1161--1266}
\end{proof}

\subsection*{Proof of \ref{B5.21}}\label{prf:B5.21}\phantomsection
\begin{proof}
The family sizes give \eqref{5.20.2} with \(C=\max(C_G,0)\)
and eventually \(b_0\le b^{1/2}\). Increase the threshold for
\eqref{5.20.3}, \(k^2x^{26}\le X'\), and \(n_T>0\).
\leanref{Formalism.\allowbreak{}L2.\allowbreak{}GlobalTheoryTheory.\allowbreak{}Loss.\allowbreak{}windows\_\allowbreak{}ev}{Formalism.\allowbreak{}LocalToGlobal.\allowbreak{}GlobalTheory.\allowbreak{}PowFourLtCeilThresholdPrime}{1555--1595}

Use \(K_0=k^2x^{78}\), \(L_1(p)=1+\lceil\log_pK_0\rceil\),
and \(\vartheta=(kx^{13})^{-1}\). Lemma~\ref{B5.17a} gives
the ordinary side conditions. At \(d\) the exponents are
positive and the weights nonnegative.
\leanref{Formalism.\allowbreak{}L2.\allowbreak{}GlobalTheoryTheory.\allowbreak{}Loss.\allowbreak{}side\_\allowbreak{}at}{Formalism.\allowbreak{}LocalToGlobal.\allowbreak{}GlobalTheory.\allowbreak{}PowFourLtCeilThresholdPrime}{1199--1209}

Use \(\rho=\rho_B,\ Y=X'(1+C_r\rho_B)\) from \ref{B4.04}.
The word primes split. By the discard criterion of \ref{B4.04}, a prime \(p\in P_s\) is discarded if \(p>X''\), or if it is of inert class and either \(h_p\ne0\) has odd order or \(p<Y\) satisfies the small-order, band, or excess-share test of \ref{B4.04}. At each such prime the gain and deduction vanish. Served primes satisfy \(p\le X''\).
\leanref{Formalism.\allowbreak{}L2.\allowbreak{}GlobalTheoryTheory.\allowbreak{}Loss.\allowbreak{}nonserved\_\allowbreak{}zero}{Formalism.\allowbreak{}LocalToGlobal.\allowbreak{}GlobalTheory.\allowbreak{}PowFourLtCeilThresholdPrime}{687--698}

Use \ref{B5.20a} for the served-prime comparisons. In the
small-order branch of \(P_1\) the order of \([p]\) divides
\(\operatorname{ord}(h_p)<x^{\varepsilon/4}\). In particular
an inert-class \(h_p\ne0\) of odd order has trivial \(2\)-group
image. The band branch satisfies the second alternative of
\ref{B5.18}, so that lemma applies with \(Y_0=\max(Y,X'')\).
The remaining discarded candidates have \(\Sigma_p>\theta\).

Signal primes and primes \(p\ge X'\) have \(TT_p=TT\).
Every other served prime is an inert candidate below \(X'\)
with eligible share. Its \(TT_p\) consists of the positions on good representatives.
Every omitted position therefore lies on a bad orbit, even when
there are no good representatives. This gives the required \(R_{\rm out}\).
If \(V_p\) is nonempty, \(\mathrm{io}_p=2\Sigma_p\max_v S_v\)
by \ref{B3.52}. Apply \(S_v\le2mn_Ta/p\). If \(V_p\) is empty,
the index vanishes. Hence \eqref{5.19.2} holds.
\leanref{Formalism.\allowbreak{}L2.\allowbreak{}GlobalTheoryTheory.\allowbreak{}Loss.\allowbreak{}unserved\_\allowbreak{}bad}{Formalism.\allowbreak{}LocalToGlobal.\allowbreak{}GlobalTheory.\allowbreak{}PowFourLtCeilThresholdPrime}{722--764}
\leanref{Formalism.\allowbreak{}L2.\allowbreak{}GlobalTheoryTheory.\allowbreak{}Loss7.\allowbreak{}loss\_\allowbreak{}setting\_\allowbreak{}family}{Formalism.\allowbreak{}LocalToGlobal.\allowbreak{}GlobalTheory.\allowbreak{}LossRequirementsOfSmallShares}{41--80}

The index is supported on full candidate fibers below \(X'\).
The fiber identity of \ref{B5.15} gives
\(\mathrm{IO}_i=\sum_{p\in\mathcal P_c}\log p\,\mathrm{io}_{i,p}\).
\leanref{Formalism.\allowbreak{}L2.\allowbreak{}GlobalTheoryTheory.\allowbreak{}Loss.\allowbreak{}IO\_\allowbreak{}at}{Formalism.\allowbreak{}LocalToGlobal.\allowbreak{}GlobalTheory.\allowbreak{}PowFourLtCeilThresholdPrime}{253--268}

By \ref{B5.21a}, \(\lceil2a\rceil\) is an eligible cutoff.
For every eligible cutoff, \ref{B5.15} identifies the ordinary
weights and the ramified coset.
\leanref{Formalism.\allowbreak{}MainNotions.\allowbreak{}Lnk.\allowbreak{}wPart\_\allowbreak{}eq}{Formalism.\allowbreak{}SharpRestriction.\allowbreak{}GlobalSchurBounds.\allowbreak{}FrameCodeOfAdditive}{464--478}
\leanref{Formalism.\allowbreak{}MainNotions.\allowbreak{}Lnk.\allowbreak{}reads\_\allowbreak{}w}{Formalism.\allowbreak{}SharpRestriction.\allowbreak{}GlobalSchurBounds.\allowbreak{}FrameCodeOfAdditive}{223--235}
\leanref{Formalism.\allowbreak{}MainNotions.\allowbreak{}ZetaZeroFreeRegion.\allowbreak{}Br.\allowbreak{}isInert\_\allowbreak{}iff}{Formalism.\allowbreak{}SharpRestriction.\allowbreak{}GlobalSchurBounds.\allowbreak{}RamificationTwoAndConjugationAtQAndD}{371--381}
\leanref{Formalism.\allowbreak{}MainNotions.\allowbreak{}Lnk.\allowbreak{}askIsF}{Formalism.\allowbreak{}SharpRestriction.\allowbreak{}GlobalSchurBounds.\allowbreak{}FrameCodeOfAdditive}{546--556}
Apply \ref{B5.20} with \(c_{\rm Loss}=3K_s+13\).
\end{proof}

\subsection*{Proof of \ref{B5.21a}}\label{prf:B5.21a}\phantomsection
\begin{proof}
The bijection \(e_\tau\) labels the indices by the host subgroup
\(\Psi\). The position \(T(e_\tau(\tau))\) of \ref{B2.02} has
first sign bit zero and recovers the host element from its
signs and \(A\)-coordinate. Its image has size \(f=2n=|\Phi|\).
The bounds of \ref{B3.47} give \(\sqrt b\le w''\le b'\),
\(t_0\eta_{\rm sz}\le w''/x\), and \(t_0\ge1\). At a mark,
\(b'\le b_s\le(1+v)b'\) with
\(v=2n\delta_1x^\varepsilon\mathcal L/k\) and \(4v\le x^{-2}\).
Its weight interval \([b_s+b'/x,b_s+b'/x+\eta_{\rm sz}]\)
is contained in \([1,2b']\), as is the thin interval
\([w'',w''+\eta_{\rm sz}]\). The thick labels
\(((0,1),s_2+u_0)\) have the indicated thick interval.
\end{proof}

\subsection*{Proof of \ref{B5.22a}}\label{prf:B5.22a}\phantomsection
\begin{proof}
Take \(\varepsilon_0=1/53\). For each \(0<\varepsilon\le\varepsilon_0\) choose the numerical threshold \(\mathcal K_W(\varepsilon)\) of \ref{B5.20y}. Now fix any family and any real \(CH\) passing this threshold. Lemmas~\ref{B5.20j} and \ref{B5.20c} apply to this family. Take the maximum of their point cutoffs and the family starting scale plus one. The averaging below is over trial weights in the box of this fixed point.
Eventually the two interval widths are at least one by
\(k\le x^6\), \(a=xe^{x^\varepsilon}\), \(b'=ae^{-x^{\varepsilon/3}}\).
Give the integer box its uniform law. Below \(p_0\), the family bound \(C_{\rm frac}r_G^2k\le x^{14}\) gives
\(\kappa_v\le N_vx^{-12}\) pointwise by \ref{B3.51}.
Above \(p_0\), \ref{B5.20c} and \ref{B5.20j} give
\(\mathbb E\kappa_v\le C_mN_vN_1/x^{13}\le N_vx^{-12}\).
\leanref{GlobalTheory.\allowbreak{}EX.\allowbreak{}RangeOne.\allowbreak{}rangeOne\_\allowbreak{}abs}{Development.\allowbreak{}AbelianVarietyTheory.\allowbreak{}Proofs.\allowbreak{}GlobalInterpolation.\allowbreak{}TensorOfPairLift}{637--640}
\leanref{Formalism.\allowbreak{}MainNotions.\allowbreak{}G.\allowbreak{}Joints.\allowbreak{}type\_\allowbreak{}one\_\allowbreak{}hypothesis\_\allowbreak{}of\_\allowbreak{}segments}{Formalism.\allowbreak{}SharpRestriction.\allowbreak{}GrowthBounds.\allowbreak{}Joints}{1283--1325}
\leanref{Formalism.\allowbreak{}MainNotions.\allowbreak{}AuxBlock8.\allowbreak{}Pf.\allowbreak{}restHypI\_\allowbreak{}ev}{Formalism.\allowbreak{}SharpRestriction.\allowbreak{}GrowthBounds.\allowbreak{}LogarithmicSums.\allowbreak{}MeanNum}{416--660}

The orbit coordinate map is injective on a period, so its induced
law is uniform. The good orbits are disjoint, so
\(\mathbb E\Sigma_p\le mn_Tx^{-12}\). For \(u=\log X'\) put
\(I_p=\mathbf1_{\{\Sigma_p>x^{-\varepsilon}\}}\) and
\(Z=\sum_{p\in\mathcal P_c}(\log p/p)(I_p+\Sigma_p)\).
Using \(I_p\le x^\varepsilon\Sigma_p\) and \eqref{5.17.6} gives
\[
 4\mathbb EZ\le4(x^\varepsilon+1)mn_Tx^{-12}u(2+u)
 \le48x^{-19/2+(9+2/100+2/3)\varepsilon}\le x^{-9}
\]
eventually, by \(m\le x^{1/2+\varepsilon/100}\) and
\(n_T\le2x^{2+8\varepsilon+\varepsilon/100}\).
Choose a pair with \(Z\le\mathbb EZ\) and use \(R_W\le4Z\).
\end{proof}

\subsection*{Proof of \ref{B5.22b}}\label{prf:B5.22b}\phantomsection
\begin{proof}
At a window point beyond the thresholds choose a model \(i_0\in I(x)\)
and carry its weights along the top bijection to \(\mathcal T(x)\).
Extend by zero. This pair satisfies \eqref{eq:p5c2-common-box} and
is realized by \(i_0\). Every realization satisfies the share
bound by hypothesis. For the constructed family, Proposition~\ref{B8.02a} gives \(R_W(i)\le x^{-9}\) at the model's own weights for every \(x>x_1\) and every \(i\in I(x)\).
\leanref{Formalism.\allowbreak{}MainNotions.\allowbreak{}FamCore.\allowbreak{}coverU\_\allowbreak{}hyp\_\allowbreak{}self}{Formalism.\allowbreak{}SharpRestriction.\allowbreak{}FamCore}{207--212}
\end{proof}

\section*{Proofs for Chapter~\ref{ch:p5c3}: A slope bound for families of maps}

\subsection*{Proof of \ref{B5.26}}\label{prf:B5.26}\phantomsection
\begin{proof}
Put \(T=\operatorname{tr}Q_Z\) and \(a_i=\|P_i-I\|_{\mathrm{HS}}\). For an orthogonal projection \(E\) and an operator \(M\) on the same space, Cauchy--Schwarz gives
\[
 |\langle E,M\rangle_{\mathrm{HS}}|\le\|E\|_{\mathrm{HS}}\|M\|_{\mathrm{HS}}.
\]
Since \(E=E^*=E^2\), this yields
\begin{equation}\label{eq:p5c3-trace-projection}
 |\operatorname{tr}(EM)|\le\sqrt{\operatorname{rank}E}\,\|M\|_{\mathrm{HS}} .
\end{equation}
With \(E=I\), \(M=P_i-I\), the same inner product inequality for the sequences \((\sqrt{d_i})_i\) and \((a_i)_i\) and \eqref{eq:p5c3-frame-source} give
\begin{equation}\label{eq:p5c3-positive-trace}
 T=\sum_i\operatorname{tr}P_i
 \ge R-\sum_i\sqrt{d_i}\,a_i
 \ge R-\sqrt{R\sum_i a_i^2}
 \ge (1-\sqrt e)R\ge R/2>0.
\end{equation}
Let \(\Pi_W\) be the orthogonal projection onto \(W=(\sum_i\widetilde B_i V'_i)^\perp\). Put \(T_W=\operatorname{tr}(\Pi_WQ_Z)\). Let \(E_i\) project onto \((V'_i)^\perp\), of dimension \(k_i\). Since \(\Pi_W\widetilde B_i\) vanishes on \(V'_i\), we have \(\Pi_W\widetilde B_i=\Pi_W\widetilde B_iE_i\) and
\[
 \operatorname{tr}(\Pi_W c_i\widetilde B_i\widetilde B_i^*)
 =c_i\|\Pi_W\widetilde B_iE_i\|_{\mathrm{HS}}^2
 \le c_i\|\widetilde B_iE_i\|_{\mathrm{HS}}^2
 =\operatorname{tr}(E_iP_iE_i).
\]
By \eqref{eq:p5c3-trace-projection},
\[
 \operatorname{tr}(E_iP_iE_i)
 =k_i+\operatorname{tr}(E_i(P_i-I))
 \le k_i+\sqrt{k_i}\,a_i .
\]
The subspaces satisfy \(\sum_i k_i\le zR\). Apply Cauchy--Schwarz to \((\sqrt{k_i})_i\) and \((a_i)_i\) and use \eqref{eq:p5c3-frame-source} to obtain
\begin{equation}\label{eq:p5c3-trace-above}
 T_W\le\sum_i k_i+\sqrt{\sum_i k_i}\sqrt{\sum_i a_i^2}
 \le R(z+\sqrt{ze}).
\end{equation}
By \eqref{eq:p5c3-frame-target}, the operator \(A=Q_Z-(T/D_Z)1\) satisfies
\[
 \|A\|_{\mathrm{HS}}^2=\operatorname{tr}(Q_Z^2)-T^2/D_Z
 \le\eta T^2/D_Z .
\]
Put \(q=\dim_\C W/D_Z\). Apply \eqref{eq:p5c3-trace-projection} to \(\Pi_W,A\) to get \(T_W/T\ge q-\sqrt{\eta q}\). By \eqref{eq:p5c3-positive-trace} and \eqref{eq:p5c3-trace-above},
\[
 q-\sqrt{\eta q}\le 2(z+\sqrt{ze}).
\]
For \(z\ge1/4\), use \(q\le1\). For \(0\le z<1/4\), use \(\sqrt{\eta q}\le(\eta+q)/2\) and \(\sqrt{ze}\le\sqrt e\) to obtain \(q\le\eta+4z+4\sqrt e\).
\end{proof}

\subsection*{Proof of \ref{B5.30}}\label{prf:B5.30}\phantomsection
\begin{proof}
By \eqref{eq:p5c3-positive-trace}, \(D_Z>0\) and hence \(D_C>0\). For \(\tau\ge0\), define
\[
 \theta_i(\tau)=\frac{s-\tau+h_i-c_K}{2},\qquad
 N_i^\tau=(N_i)_{\ge\theta_i(\tau)},\qquad
 x_{ij}=2\lambda_j(N_i)+c_K-h_i\quad(1\le j\le n_Kd_i).
\]
A rank-zero source contributes an empty sum. By \ref{B5.24}, the sum of the adjusted slopes is \(n_K\sum_i d_i(\mu(N_i)-h_i)\ge n_KRs\) and each is at most \(s+U\). The count in \ref{B5.25}, divided by \(n_K\) using \ref{B5.24}, gives
\begin{equation}\label{eq:p5c3-source-corank}
 \sum_i(d_i-r(N_i^\tau))\le Rz_U(\tau).
\end{equation}
Let \(S^\tau=(C^\vee)_{\ge(s-\tau-c_K)/2}\). For a nonzero image \(B_iN_i^\tau\), \ref{B5.24} and \ref{B5.23} give
\[
 \mu_{\min}(B_iN_i^\tau)\ge
 \mu_{\min}(N_i^\tau)-h(B_i)\ge s-\tau.
\]
By \ref{B5.24} the image lies in \(S^\tau\). Thus \(M^\tau=\sum_i B_iN_i^\tau\subseteq S^\tau\). Apply \ref{B5.27} with \eqref{eq:p5c3-source-corank} and use \(D_Z\le D_C/(1-\delta)\) to obtain
\begin{equation}\label{eq:p5c3-target-corank}
 D_C-r(S^\tau)\le D_C-r(M^\tau)\le D_Cw(\tau),
 \qquad w(\tau)=\frac{\eta+4z_U(\tau)+4\sqrt e}{1-\delta}.
\end{equation}
By \ref{B5.24}, the adjusted target slopes \(y_j=2\lambda_j(C^\vee)+c_K\) have mean \(\mu(C^\vee)\) and are at least \(s-H\). By \eqref{eq:p5c3-target-corank}, at most \(n_KD_Cw(\tau)\) lie below \(s-\tau\). The function \(w\) is continuous on \([0,H]\), including when \(U=0\). Apply \ref{B5.28} and integrate \(z_U\) to get
\[
 s-\mu(C^\vee)\le\int_0^H w(\tau)\,d\tau
 =E_{\mathrm{fr}}(U,H;\delta,\eta,e).
\]
The second inequality follows since \(\log\#I\ge0\).
\end{proof}

\subsection*{Proof of \ref{B5.31}}\label{prf:B5.31}\phantomsection
\begin{proof}
By \eqref{eq:p5c3-code-input},
\begin{equation}\label{eq:p5c3-large-total}
 \mathcal R\ge \#\mathcal W\,\frac{D_*}{C_0\Gamma^2}\ge\frac{D_*}{\sigma}.
\end{equation}
Put \(a=\sqrt{e_0}\), \(t=\sqrt{d_e}\) and
\[
 Q_f=\sum_uc_uA_uA_u^*,\qquad Q_Z=\sum_uc_u A_uJ_uJ_u^*A_u^*.
\]
Complete the columns of each \(J_u\) to an orthonormal basis. In these bases \(J_u^*MJ_v\) is a submatrix of \(M\), so its squared Hilbert--Schmidt norm is bounded by that of \(M\). Apply this to \(c_u\widetilde B_u^*\widetilde B_u-I=J_u^*(K_{uu}-I)J_u\) to get a source defect sum at most \(e_0\mathcal R\) by \eqref{eq:p5c3-diagonal}. By \eqref{eq:p5c3-rank-loss}, \(R\ge(1-d_e)\mathcal R\), so the sum is at most \(2e_0R\). By \eqref{eq:p5c3-small-input}, \(2e_0\le\sigma^2/2\le1/800\).

By \eqref{eq:p5c3-diagonal}, \(\|K_{uu}\|_{\mathrm{HS}}\le(1+a)\sqrt{f_u}\). Expand the trace and use \eqref{eq:p5c3-cross} to obtain
\[
 \operatorname{tr}(Q_f^2)=\sum_{u,v}\|K_{uv}\|_{\mathrm{HS}}^2
 \le (1+\eta_0)\frac{\mathcal R^2}{D_*}+(1+a)^2\mathcal R .
\]
Here \(1+\eta_0\ge0\). By \eqref{eq:p5c3-large-total},
\begin{equation}\label{eq:p5c3-compressed-square}
 \operatorname{tr}(Q_f^2)\le\frac{\mathcal R^2}{D_*}
 \bigl(1+\eta_0+(1+a)^2\sigma\bigr).
\end{equation}
The restricted cross matrices are \(J_u^*K_{uv}J_v\), so contraction in each pair gives \(\operatorname{tr}(Q_Z^2)\le\operatorname{tr}(Q_f^2)\).

By \eqref{eq:p5c3-diagonal} and \eqref{eq:p5c3-trace-projection}, \(\operatorname{tr}Q_f=\sum_u\operatorname{tr}K_{uu}\ge(1-a)\mathcal R\). For the projection \(E_u=I-J_uJ_u^*\), of rank \(f_u-d_u\), the same estimates give
\begin{align*}
 \operatorname{tr}Q_f-\operatorname{tr}Q_Z
 &=\sum_u\operatorname{tr}(K_{uu}E_u)\\
 &\le \sum_u(f_u-d_u)+a\sum_u\sqrt{f_u(f_u-d_u)}\\
 &\le (d_e+a\sqrt{d_e})\mathcal R
\end{align*}
by Cauchy--Schwarz for \((\sqrt{f_u})_u\) and \((\sqrt{f_u-d_u})_u\) and \eqref{eq:p5c3-rank-loss}. Consequently
\begin{equation}\label{eq:p5c3-compressed-trace}
 \operatorname{tr}Q_Z\ge\beta\mathcal R,\qquad
 \beta=1-a-t^2-ta\ge57/80,
\end{equation}
since \(a\le1/40\) and \(t\le1/2\).

For the coefficient in \eqref{eq:p5c3-compressed-square}, we have
\[
 (1+\eta_1)\bigl(1+\eta_0+(1+a)^2\sigma\bigr)
 \le (1+5\sigma)(1-a)^2.
\]
Indeed \(1+\eta_1\ge0\) by \eqref{eq:p5c3-rank-loss}. The left side increases with \(\eta_0,\eta_1,a\), and the right side decreases with \(a\). At \(\eta_0=\eta_1=a=\sigma/2\) the difference of right and left is
\[
 \sigma\left(2-\frac{13}{2}\sigma+\frac12\sigma^2-\frac18\sigma^3\right)\ge0
 \qquad(0<\sigma\le1/20).
\]
Put \(P=1+5\sigma+3t\) and \(B=1-a\). Write \(v=t(t+a)\). Then \(P\le11/4\) and \(t+a\le21/40\), so \(2P(t+a)\le231/80<117/40\le3B\). Hence
\[
 P\beta^2-(1+5\sigma)B^2
 =tB\bigl(3B-2P(t+a)\bigr)+Pv^2\ge0.
\]
Combining \eqref{eq:p5c3-rank-loss}, \eqref{eq:p5c3-compressed-square} and \eqref{eq:p5c3-compressed-trace} gives
\[
 D_Z\operatorname{tr}(Q_Z^2)\le P\beta^2\mathcal R^2
 \le P(\operatorname{tr}Q_Z)^2.
\]
Thus \eqref{eq:p5c3-frame-source}--\eqref{eq:p5c3-frame-target} hold with the stated constants. The definition of \(s\) gives the mean hypothesis of \ref{B5.29} with equality. By \ref{B5.24}, \(U\ge0\), and the definition of \(U\) gives the required upper bound on each adjusted maximum slope. Apply \ref{B5.30}.
\end{proof}

\subsection*{Proof of \ref{B5.31a}}\label{prf:B5.31a}\phantomsection
\begin{proof}
For each word \(u\) let \(K_2\) be the compact open section level, \(F_{N,u}\) a fundamental domain for the unipotent lattice, \(H_u\) the compact real stabilizer and \(V_{\pi_u}\) its nonzero coefficient space. With the Haar measures of the selected realization the unfolding constant is
\[
 \alpha_u=\frac{\operatorname{vol}(K_2)\operatorname{vol}(F_{N,u})\operatorname{vol}(H_u)}{\dim_{\C}V_{\pi_u}}>0.
\]
The compact open subgroup has positive finite Haar measure, the lattice has positive finite covolume, and the compact group has positive finite total measure. Thus every block can be divided by \(\sqrt{\alpha_u\alpha_v}\).
\leanref{Formalism.\allowbreak{}L2.\allowbreak{}SG.\allowbreak{}PoincareInHH.\allowbreak{}PoincareSetting.\allowbreak{}c0X}{Formalism.\allowbreak{}LocalToGlobal.\allowbreak{}SiegelLatticeTheory.\allowbreak{}HaarTransversals.\allowbreak{}PoincareInHH}{257--258}
\leanref{Formalism.\allowbreak{}L2.\allowbreak{}SG.\allowbreak{}TermOp.\allowbreak{}IdBlocks.\allowbreak{}idDiagBlocks\_\allowbreak{}of}{Formalism.\allowbreak{}LocalToGlobal.\allowbreak{}SiegelLatticeTheory.\allowbreak{}ClosedLocalOrbits.\allowbreak{}GaugeLatticeShifts.\allowbreak{}IdBlocks}{785--800}

For \(u\ne v\), \ref{B5.11b} gives both orders of the level and Vandermonde comparisons with \(A=3e^{7/4}\). By Lemma~\ref{B8.02d}, these comparisons and the source conditions imply all the counting conditions of Definition~\ref{B8.02c} for the pair \((u,v)\) with this value of \(A\). With the size thresholds of \ref{B8.10}, the ties identify the column dimensions with \(f_u,f_v\) and the main term with \(D_*=D^\star\). Apply \ref{B8.10b} to the full normalized Gram operator. Its right hand side is \((1+b^{-6})\Delta_{u,v}\), where the dimension and main-term identifications just made give \(\Delta_{u,v}=f_uf_v/D_*\). This is \eqref{eq:p5c3-cross} with \(\eta_0=b^{-6}\). The bound is uniform because the differing coordinate used in \ref{B8.10e} can be chosen for every pair of distinct words and the threshold in \ref{B8.10b} is common to those pairs.

Fix \(u\) and use \(r_0=(m-2)/2\). For a term in the unfolding of \ref{B8.10}, let \(p_\ell\) be the order-\(\ell\) leading pairing minor of the two flags, and let \(T_u\) be the level prime. The Gram invariant is
\[
 G=p_{r_0+1}\overline{p_{r_0+1}}-T_u^2p_{r_0}\overline{p_{r_0}},
\]
where the bar is CM conjugation. The full normalized block is \(K_0+K''\), with \(K_0\) the sum over the identity cell and the \(G\ne0\) terms, and \(K''\) the sum over the remaining \(G=0\) terms.

Let \(E/F_{\mathrm{real}}\) be the host CM extension. For the fixed flag \(F\) and the transported flag \(F'\), write \(\mathfrak n_F,\mathfrak n_{F'}\) for their unipotent Lie algebras and \(\ell_F,\ell_{F'}\) for the simple-root functionals of \ref{B8.10y}. Relevance means
\[
 \operatorname{Tr}_{E/F_{\mathrm{real}}}\ell_F(Y)
 =\operatorname{Tr}_{E/F_{\mathrm{real}}}\ell_{F'}(Y)
 \qquad(Y\in\mathfrak n_F\cap\mathfrak n_{F'}).
\]
If the equality fails at \(Y\), set \(\Delta=(\ell_F-\ell_{F'})(Y)+\overline{(\ell_F-\ell_{F'})(Y)}\ne0\) and \(s=g_0g_1/(2f\Delta)\). The real image of \(\exp(sY)\) belongs to both unipotent groups. The character ratio at this element is
\[
 \exp\left(\frac{2\pi i}{g_0g_1}
 \sum_{\tau}\tau(s\Delta)\right)=\exp(\pi i)=-1,
\]
where \(\tau\) runs over the \(f\) embeddings of the real field. The finite support and the global character identity make the ratio trivial on the arithmetic subgroup of the intersection used in the orbital integral. In the quotient average over the intersection, translation by \(\exp(sY)\) preserves the measure and multiplies the character integral by \(-1\). Thus the orbital integral vanishes.

Let \(h_\rho\) represent a finite column class \(\rho\) with row lattice \(\Lambda_\rho\). Let \(x_u\) be the fixed level translate and \(K_2\) the subgroup supporting the finite section. For a term \(t\), write \(t_f\) for the finite component of its rational representative. The term meets \((\rho,\rho')\) exactly when its ordered row lattices are \((\Lambda_\rho,\Lambda_{\rho'})\) and
\begin{equation}\label{eq:p5c3-meeting}
 (h_\rho x_u)^{-1}t_f(h_{\rho'}x_u)\in K_2.
\end{equation}
In the identity cell the rational stabilizer element is \(n\eta\), with \(n\) in the unipotent radical and \(\eta\) in the rational Levi. Write \(U_u^{\rm lev}\) for the column level, the Levi subgroup cut out by membership in \(K_2\) after conjugation by \(x_u\). The level decomposition projects \eqref{eq:p5c3-meeting} to \(h_\rho^{-1}\eta_fh_{\rho'}\in U_u^{\rm lev}\). Thus the two representatives define the same double coset and \(\rho=\rho'\). Neatness of the column level forces \(\eta=1\). The finite condition then implies that \(n\) belongs to the unipotent stabilizer of the row lattice. The ratio of any two such representatives lies in this stabilizer, so they represent the same term in the unipotent quotient. The unit representative supplies a meeting term. The finite support makes every nonmeeting contribution zero.
\leanref{Formalism.\allowbreak{}L2.\allowbreak{}SG.\allowbreak{}TermOp.\allowbreak{}IdBlocks.\allowbreak{}Meets}{Formalism.\allowbreak{}LocalToGlobal.\allowbreak{}SiegelLatticeTheory.\allowbreak{}ClosedLocalOrbits.\allowbreak{}GaugeLatticeShifts.\allowbreak{}IdBlocks}{717--721}
\leanref{Formalism.\allowbreak{}L2.\allowbreak{}SG.\allowbreak{}Pf.\allowbreak{}IdMeet.\allowbreak{}idMeetKeyC\_\allowbreak{}holds}{Formalism.\allowbreak{}LocalToGlobal.\allowbreak{}SiegelLatticeTheory.\allowbreak{}ClosedLocalOrbits.\allowbreak{}GaugeLatticeShifts.\allowbreak{}IdMeetKeyPf}{136--178}
\leanref{Formalism.\allowbreak{}L2.\allowbreak{}SG.\allowbreak{}Pf.\allowbreak{}IdMeet.\allowbreak{}idMeetOneC\_\allowbreak{}of}{Formalism.\allowbreak{}LocalToGlobal.\allowbreak{}SiegelLatticeTheory.\allowbreak{}ClosedLocalOrbits.\allowbreak{}GaugeLatticeShifts.\allowbreak{}IdMeetOnePf}{61--90}

In the meeting block put \(k_\rho=(h_\rho x_u)^{-1}n_f(h_\rho x_u)\). The level decomposition gives \(\chi_2(k_\rho)=\psi_f(n_f)\), where \(\chi_2\) is the level character and \(\psi_f\) is the finite unipotent character. By unitarity, \(\overline{\psi_\infty(n_\infty^{-1})}=\psi_\infty(n_\infty)\). The global character identity on rational points gives
\[
 \chi_2(k_\rho)\overline{\psi_\infty(n_\infty^{-1})}
 =\psi_f(n_f)\psi_\infty(n_\infty)=1.
\]

For \(h\) in the finite adelic Levi, let \(F(h)\) be the covolume of its real unipotent lattice \(\Lambda(h)\). Passing to \(\Lambda(g)\cap\Lambda(gh)\) multiplies each covolume by the index of this intersection in the corresponding lattice. The indices are the relative indices in the finite unipotent level subgroups. These subgroups for \(g\) and \(gh\) are simultaneous conjugates by \(g\) of the subgroups for \(1\) and \(h\). Conjugation preserves both indices, so
\[
 \frac{F(gh)}{F(g)}=\frac{F(h)}{F(1)}.
\]
Right multiplication by the column level leaves the unipotent subgroup unchanged. Left multiplication by a rational Levi representative conjugates the real lattice by its compact component, which preserves the unipotent Haar measure. Thus \(F\) is constant on column double cosets and has finite range. The positive character \(F(h)/F(1)\) is one, since the powers of each of its values also belong to its finite range. Every meeting block therefore has the same unipotent volume \(\operatorname{vol}_N=F(1)\).

Let \(\pi_u\) be the irreducible coefficient representation and \(\operatorname{vol}_H\) the compact-group Haar volume. The real averaging operator after phase cancellation commutes with \(\pi_u\) and is scalar. For the coefficient vector \(W\), \(\operatorname{tr}(W\otimes W^*)=\|W\|^2\). With the quotient measure \(\nu\) the normalization is \(\int\|W\|^2\,d\nu=1\). The Weil integration formula thus gives trace \(F(1)\operatorname{vol}_H\). Multiplication by the finite volume \(\operatorname{vol}(K_2)\) gives the scalar \(c_{0,X}=\operatorname{vol}(K_2)\operatorname{vol}_N\operatorname{vol}_H/\dim\pi_u\) of the unnormalized block. The level and unfolding ties give the normalized identity block \(\delta_{\rho\rho'}I_{\pi_u}\). Consequently \(K_0-I=K_{\rm gen}+K_{\rm ng}\), where \(K_{\rm gen}\) and \(K_{\rm ng}\) are the generic and nongeneric nonzero-\(G\) sums.
\leanref{Formalism.\allowbreak{}L2.\allowbreak{}SG.\allowbreak{}Pf.\allowbreak{}IdMeet.\allowbreak{}idMeetPhaseC\_\allowbreak{}of}{Formalism.\allowbreak{}LocalToGlobal.\allowbreak{}SiegelLatticeTheory.\allowbreak{}ClosedLocalOrbits.\allowbreak{}GaugeLatticeShifts.\allowbreak{}IdMeetPhasePf}{14--42}
\leanref{Formalism.\allowbreak{}L2.\allowbreak{}SG.\allowbreak{}TermOp.\allowbreak{}IdBlocks.\allowbreak{}idDiagBlocks\_\allowbreak{}of}{Formalism.\allowbreak{}LocalToGlobal.\allowbreak{}SiegelLatticeTheory.\allowbreak{}ClosedLocalOrbits.\allowbreak{}GaugeLatticeShifts.\allowbreak{}IdBlocks}{761--821}

The adjacent weight gaps in \ref{B8.02c} imply the pair gaps of \ref{B8.10y}. The same source conditions give \((kb)^8\le g_i\le b^{C_g}\). They also give \(T_u\le b^{C_Tm}\) and \((T_u,g_0g_1)=1\). In the normal form of \ref{B8.10y}, let \(X,X'\) be the spans of the isotropic vectors of the two frames over \(E\). That proposition gives \(p_j\ne0\) at \(j=r_0-\dim_E(X\cap X')\), with \(1\le j\le r_0\) for each relevant term of \(K''\).

For a rank-\(j\) label \(\delta\), use the leading minors \(\delta_\ell=p_\ell\) and the quantities \(q_\ell,H_1,\mathfrak d,D_\delta,S_\delta\) of \ref{B8.10y}. Here \(\mathfrak d\) is the extended different of the real subfield. Write \(N\) for the absolute ideal norm and \(\operatorname{rd}(E)\) for the root discriminant. Since \(E/F_{\mathrm{real}}\) is unramified at finite places, \(N(\mathfrak d)=\operatorname{rd}(E)^{2f}\). With \(\tau\) in the CM type,
\[
 0<D_\delta=\prod_\tau\bigl(|\tau(\delta_j)|\operatorname{rd}(E)^{-2q_j}\bigr)^2
 \le(2H_1)^{4r_0 f}S_\delta^2.
\]
Let \(A(j,\delta)\) and \(M(j,\delta)\) be the amplitude and counting factors of \ref{B8.10y}, with the constants \(C_s,C_p\) there. The bound for one vector in \eqref{eq:p8c1-type-factors} supplies \(D_\delta^{-(k-1-j)/2}S_\delta^{-C_sk}\). The sum of relative stabilizer weights supplies \(D_\delta^{k/2}S_\delta^{C_pk}\). Thus the label factor of \(AM\) satisfies
\begin{equation}\label{eq:p5c3-label-factor}
 D_\delta^{(j+1)/2}S_\delta^{(C_p-C_s)k}
 \le (2H_1)^{2r_0 f(j+1)}S_\delta^{-3r_0}.
\end{equation}
Since \(S_\delta\ge1\), it suffices that \((C_p-C_s)k+j+1\le-3r_0\). This follows from \(C_s\ge C_p+4\) and the bounds \(j\le r_0\) and \(k\ge r_0+3\).

The pairing minors are algebraic integers because the relevant vectors of both frames lie in the integral row lattice. Put \(R_\ell=\max(1,\operatorname{rd}(E)^{2q_\ell}(2H_1)^{2r_0})\). For a finite set of values \(\xi\) of one minor,
\begin{equation}\label{eq:p5c3-minor-sum}
 \sum_\xi\prod_\tau\max(1,|\tau(\xi)|/R_\ell)^{-3}
 \le (64R_\ell)^{2f}.
\end{equation}
Partition the real and imaginary coordinates of the embedding vector into intervals of length \(1/2\). A box contains at most one \(\xi\), since two distinct integers in the box would have \(0<|N_{E/\Q}(\xi-\xi')|<1\). The weight in one complex coordinate is bounded by the product of the weights \(\max(1,|z|/R_\ell)^{-3/2}\) of its two real coordinates \(z\). In each real coordinate, the sum of the suprema over the half-unit intervals is at most
\[
 2+4\int_0^\infty\max(1,t/R_\ell)^{-3/2}\,dt
 =2+12R_\ell\le64R_\ell .
\]

There are at most \(2^j\) choices for the nonzero minor positions, and their values determine the label. Since \(j\le r_0\), \(S_\delta^{-3r_0}\) is bounded by the product of the weights in \eqref{eq:p5c3-minor-sum} over the nonzero positions. Thus for every finite set \(\mathcal D_j\) of occurring rank-\(j\) labels,
\begin{equation}\label{eq:p5c3-label-sum}
 \sum_{\delta\in\mathcal D_j}S_\delta^{-3r_0}
 \le 2^j\prod_{\ell=1}^j(64R_\ell)^{2f}
 \le \exp(C_2m^2 f Lr_0).
\end{equation}
For the last inequality, use \(q_\ell\le r_0^2\) and the fact that \(H_1\) is a fixed power of \(b^m\). By \ref{B8.10y}, \(\log\operatorname{rd}(E)\le L/C_\ell\). Thus \(\log(64R_\ell)=O(m^2L)\) with a uniform constant, and \(C_2\) can be chosen uniformly.

By \eqref{eq:p5c3-label-factor} and \eqref{eq:p5c3-label-sum}, the product \(AM\) satisfies
\begin{equation}\label{eq:p5c3-amplitude-count}
 \sum_{\delta\in\mathcal D_j} A(j,\delta)M(j,\delta)\le a_j,
 \qquad a_j=\exp(Cm^6fL)(C/k)^{j(k-1-j)f}.
\end{equation}
Here \(C\ge1\) enlarges the constants of \ref{B8.10y}. The height factors there, \(H(C)^{Cr_0^4}\) and \(H(C_h)^{C_hr_0^4}\) with \(H(c)=b^{cm^2 f}\), have logarithms \(O(m^6fL)\) before the enlargement. The extra \(2H_1\) factor and the label-sum bound have the same upper order.

Let \(v_t\) be the orbital vector of term \(t\). For a set of terms \(S\), use the operator \(\operatorname{part}(S)\) of \ref{B8.10} and put
\[
 w_i(S)=\sum_{\substack{t\in S\\\iota(\mathrm{row}(t))=i}}v_t,
\]
where \(\iota\) gives the finite-type index of a row. For finite \(S\), expand the Hilbert--Schmidt square as a sum of inner products over the column blocks. In each pair of terms, the finite character coefficients pair to the entry \(G_{ij}\) of the finite-type Gram matrix, where \(i,j\) are their finite-type indices. The remaining factors pair to the orbital-vector inner product. Grouping the expansion by \(i,j\) gives
\[
 \|\operatorname{part}(S)\|_{\mathrm{HS}}^2
 =\Re\sum_{i,j}G_{ij}\langle w_i(S),w_j(S)\rangle .
\]
Absolute summability gives this identity for the full cell. Apply \ref{B8.10f} to bound it by \(P_g\Gamma_1\sum_i\|w_i(S)\|^2\), where \(P_g=\prod_{i=0}^1(1+3/g_i)^f\) and \(\Gamma_1\) is the finite-type normalization. Within a fixed type \(i\), each \(v_t\) is supported on its ordered row-lattice pair, so different rows give zero cross terms.

In the singular cell, bound the same-row inner products by products of norms. By \ref{B8.10y} and \eqref{eq:p5c3-amplitude-count}, the sum of these norm products over any finite set of pairs of ranks \(j,j'\) is at most \(\operatorname{mass}_u\dim(\pi_u)a_ja_{j'}\). Here \(\operatorname{mass}_u\) is the sum of the inverse row-lattice stabilizer orders. If \(n_\Lambda\) is the number of column classes over a row lattice \(\Lambda\) with stabilizer order \(a_\Lambda\), the finite level gives at least \(\Gamma_1\) distinct right cosets over \(\Lambda\). At most \(a_\Lambda\) of these can lie in one double coset, since the row stabilizer identifies the right cosets over \(\Lambda\). Thus \(\Gamma_1/a_\Lambda\le n_\Lambda\). Neatness gives a contribution of \(\dim\pi_u\) from each column class, so
\[
 \Gamma_1\operatorname{mass}_u\dim\pi_u
 \le \dim\pi_u\sum_\Lambda n_\Lambda=f_u.
\]
Since \(f\le k\) and \(g_i\ge(kb)^8\), \(P_g\le2\) for large \(x\). Sum over the two ranks and take the supremum of the finite nonnegative pair subsums. Absolute convergence gives
\begin{equation}\label{eq:p5c3-singular}
 \|K''\|_{\mathrm{HS}}^2\le 2f_u\left(\sum_{j=1}^{r_0}a_j\right)^2
 \le f_u e^{-kf\log k}.
\end{equation}
For the last inequality, \(j(k-1-j)\ge k-2\) in this range, so for large \(k\) the sum is at most \(k\exp(Cm^6fL)(C/k)^{(k-2)f}\). With \(C_\ell\ge8C\) and \(C_\ell m^6 L\le k\log k\), the squared height factor is at most \(\exp(kf\log k/4)\). For large \(k\) the squared power is at most \(\exp(-3kf\log k/2)\), and \(2k^2\le\exp(kf\log k/4)\).

For \((u,u)\), the level and Vandermonde factors agree, so \ref{B8.02d} gives the counting hypotheses with \(A=1\). Put \(T=m^5fL\) and \(\Delta_{uu}=f_u^2/D^\star\). Set \(\ell=e^{-kf\log k/3}\). In a same-row pair from a nonzero-\(G\) cell, \(G(t)\ne0\) gives linear independence of the \(2r_0+2\) vectors of the fixed frame and the frame of \(t\). The pair space lies in their orthogonal complement and has dimension at most \((k+2r_0+1)-(2r_0+2)=k-1\). Thus the pair has corank one or corank at least two.

Apply the norm expansion and \ref{B8.10f} to the two nonzero-\(G\) cells with the signed same-row inner products. The higher-corank and nonscalar loop bounds of \ref{B8.10a}, together with the signed corank-one estimates of \ref{B8.10d}, give
\[
 \begin{aligned}
 \|K_{\rm gen}\|_{\rm HS}^2 &\le P_g(e^{C_BT}+2\ell)\Delta_{uu},\\
 \|K_{\rm ng}\|_{\rm HS}^2 &\le P_g(\ell+2^{-(k-2r_0-5)}e^{C_4T})\Delta_{uu}.
 \end{aligned}
\]
Using \(\ell\le1\) and \(k\ge2r_0+5\), their sum is at most \(P_g(e^{C_BT}+3+e^{C_4T})\Delta_{uu}\). For large \(x\), \(10P_g\le e^T\). Since \(C_B,C_4\ge0\) and \(C_\Gamma\ge1+\max(C_B,C_4)\), twice the sum is at most \(\Gamma\Delta_{uu}\). Hence \(\|K_0-I\|_{\mathrm{HS}}^2\le\Gamma\Delta_{uu}\) by the squared triangle bound \(\|M+N\|_{\mathrm{HS}}^2\le2(\|M\|_{\mathrm{HS}}^2+\|N\|_{\mathrm{HS}}^2)\). Apply this bound to \(K_0-I\) and \(K''\). With \eqref{eq:p5c3-singular} and the column-size bound \(\Gamma f_u/D^\star\le C_0/\Gamma\) of \ref{B8.02c}, the full squared defect is at most \(e_0f_u\).

In \ref{B5.10} let \(V\) be the common self-isometry of the retained row space and \(U_u\) the source isometries. Write \(a_u\ne0\) for the scalars. Let \(\jmath\) denote the antilinear isometric embedding of the retained row space into the full target. For an orthonormal basis \((q_a)_a\) of the retained rows, extend \((\jmath(Vq_a))_a\) to an orthonormal basis of the full target. Let \((e_l)_l\) and \((e'_j)_j\) be orthonormal bases of the unrestricted and restricted sources, respectively, and let \(\iota_u\) be the isometric source inclusion. Set \((J_u)_{lj}=\langle e_l,\iota_u U_u e'_j\rangle\). Parseval gives
\[
 (J_u^*J_u)_{jj'}
 =\langle\iota_u U_u e'_j,\iota_u U_u e'_{j'}\rangle
 =\delta_{jj'}.
\]
Let \(P_u^{\rm per}\) be the unrestricted period operator and \(A_u\) its Bessel matrix. The retained entries of \(A_uJ_u\) are \(\langle\jmath(Vq_a),P_u^{\rm per}\iota_u U_u e'_j\rangle\), with the inner product linear in the second argument. The antilinearity of \(\jmath\) makes these pairings linear in \(q_a\). If \(P_E,P_u\) are the changes from orthonormal to arithmetic coordinates before applying the isometries, the Bessel identity therefore gives the retained block as
\[
 P_E^{\mathsf T}B_{u,v}(a_u^{-1}P_u).
\]
Conjugate all matrices for the conjugate embedding. The target change is common, and the invertible source change includes \(a_u^{-1}\).

With \(c_u=\alpha_u^{-1}\), the normalized matrices are \(\sqrt{c_uc_v}A_u^*A_v\), up to entrywise conjugation, which preserves the squared norms.

In \ref{B8.52} take the entire type set at either distinguished prime. Its cylinder contains every label pair and its dimension is \(D_{\rm all}=D_Z\). By \eqref{eq:green-cylinder}, its idempotent is the full level idempotent \(e\). For the selected realization, the row main term \(D^\star\) equals the spectral main term \(D_A^\star(e)\) of the full level idempotent. Thus \eqref{eq:p8c5-count-main} gives
\[
 D_A^\star(e)=D^\star=d_\infty\operatorname{vol}(\mathscr F)\Re e(1)>0,
\]
with the degree factor \(d_\infty\) and fundamental domain \(\mathscr F\) of \ref{B8.53}. By the cylinder estimate of \ref{B8.52} with \(c_{\rm Ic}=1/16\) from \eqref{eq:p8c5-count-constants},
\[
 D_Z-D^\star\le e^{-rf\log r/16}D^\star=\eta_1D^\star .
\]

The inclusion in \ref{B8.52} identifies the retained rows with the indices good at both primes. Every omitted index belongs to at least one of the two bad cylinders. For these cylinders \ref{B8.52} supplies the dimension and relative-count conditions. Equation \eqref{eq:green-cylinder} gives the main term for a finite set \(S\) of bad tuples at the distinguished prime \(p_j\),
\[
 D_j(S)=D^\star\sum_{t\in S}\prod_{v\mid p_j}w_{j,v}(t_v),
\]
where \(w_{j,v}\) is the local weight of this row system at \(v\). With the cuspidal parametrization of Definition~\ref{B7.06}, Theorem~\ref{B7.53} gives the local laws and size bounds for this row system. Apply \ref{B7.52} to obtain \eqref{eq:p5c3-row-defect} for \(\delta=1-D_{\rm rows}/D_Z\).

In the level model of \ref{B8.37}, the commensurability and ideal twist in \ref{B5.04} preserve the ramified column rank. By \ref{B5.05a}, \(d_u=f_u(1-1/(T_u-1))^f\). By \ref{B8.51}, \(T_u\ge b^{m/2}\ge b\ge3\) for large \(x\). For \(0\le p\le1\),
\begin{equation}\label{eq:p5c3-geometric}
 1-(1-p)^f=p\sum_{j=0}^{f-1}(1-p)^j\le fp.
\end{equation}
At \(p=1/(T_u-1)\), \eqref{eq:p5c3-geometric} gives \(f_u-d_u\le f_uf/(T_u-1)\le2ff_u/b\). Summing over \(u\) proves the rank-loss bound.
\end{proof}

\subsection*{Proof of \ref{B5.32}}\label{prf:B5.32}\phantomsection
\begin{proof}
Take \(x\) above the common cutoffs and \(20^{1/C_H}\). By \ref{B5.31}, \(\eta=\bar\eta\) and \(e=2e_0\le\bar e\). The numerical hypothesis gives
\[
 (\bar\eta+4\sqrt{\bar e})x^{C_H}
 =5+(3\sqrt{d_e}+4\sqrt{e_B})x^{C_H}\le6.
\]
Arakelov slope duality \cite[\S2.4]{bost-chen-2013} and the height clause of \ref{B5.10} give
\[
 s-\mu_{\min}(RR^\vee)=s+\mu_{\max}(RR)\le x^{C_H}b
\]
at each translation. Apply \ref{B5.31} with this \(H\) and use \((\eta+4\sqrt e)H\le6b\) in \eqref{eq:p5c3-frame-error}.
\end{proof}

\subsection*{Proof of \ref{B5.32b}}\label{prf:B5.32b}\phantomsection
\begin{proof}
Take the maximum of the class-number and column cutoffs, \(3\), and the family starting scale plus one. The type primes split in the real host field by the host splitting conditions of \ref{B2.56}. The selected norm-prime congruence of \ref{B8.51a} gives \(T_u\equiv1\bmod Q_Sqd\), hence splitting in the cyclotomic host and its real subfield. The type-prime size bound of \ref{B2.51} gives \(g_0\ge(kb)^8\). Since \(b\ge2\) and \(1\le f\le k\), it follows that \(g_0-1>2fk\). The freeness argument of \ref{B8.54} applies to every level-zero class.

Let \(c\) run through those classes, and choose a representative \(y_c\). By the symmetric column identification of \ref{B8.09}, a lowest form is an equivariant function with coefficient space \(\bigotimes_\rho V(u,\rho)\), where \(\rho\) runs over the archimedean embeddings of the CM type. Its values at \(y_c\) determine it. Conversely prescribe a coefficient vector at each \(y_c\) and extend by the rational action and the level character. Two expressions for the same adelic point differ by a stabilizer at \(y_c\); this stabilizer is trivial. Evaluation therefore identifies the lowest form space with one copy of the coefficient space for each class. The archimedean weights here are the column weights: passing from the shifted coordinates to the Weyl dimension product subtracts a common shift, which leaves all its differences unchanged. Hence
\[
 f_u=\#\mathrm{cls}_u(0)\prod_\rho\dim V(u,\rho).
\]
Formula \((8.54.3)\) of \ref{B8.54} now gives the asserted rank.
\leanref{Formalism.\allowbreak{}MainNotions.\allowbreak{}Ext.\allowbreak{}Lw.\allowbreak{}finrank\_\allowbreak{}lowest\_\allowbreak{}forms\_\allowbreak{}space}{Formalism.\allowbreak{}SharpRestriction.\allowbreak{}GlobalSchurBounds.\allowbreak{}GlobalRankLatticeProofs.\allowbreak{}UnitTraceSumBounds.\allowbreak{}RankConditionsOfLayerTwo}{230--256}
\leanref{Formalism.\allowbreak{}MainNotions.\allowbreak{}Ext.\allowbreak{}Lw.\allowbreak{}card\_\allowbreak{}ssIdx\_\allowbreak{}eq}{Formalism.\allowbreak{}SharpRestriction.\allowbreak{}GlobalSchurBounds.\allowbreak{}GlobalRankLatticeProofs.\allowbreak{}UnitTraceSumBounds.\allowbreak{}RankConditionsOfLayerTwo}{257--282}
\leanref{Formalism.\allowbreak{}MainNotions.\allowbreak{}Ext.\allowbreak{}Lw.\allowbreak{}rankTiesF\_\allowbreak{}of\_\allowbreak{}C2S\_\allowbreak{}alt}{Development.\allowbreak{}AutomorphicArithmetic.\allowbreak{}ThetaCorrespondence.\allowbreak{}ThetaPairLifts.\allowbreak{}ThirdPlaceSeries.\allowbreak{}ClosedRankTiesFC2S}{12--19}

Put \(N=fk\). The lower endpoint of the word window is \(T_u^\star=(D^\star/(d_u^\star\Gamma^2))^{1/N}\), and the selected prime lies between \(T_u^\star\) and \((1+1/N)T_u^\star\). Substitution gives
\begin{equation}\label{eq:p5c3-word-dimension}
 \frac{f_u\Gamma^2}{D^\star}
 =\left(\frac{T_u}{T_u^\star}\right)^N(1-T_u^{-k})^f .
\end{equation}
The first factor lies in \([1,(1+1/N)^N]\subseteq[1,\exp(1)]\). The second is at most one and at least \(1-fT_u^{-k}\ge1-k2^{-k}\ge1/2\), since \(T_u\ge2\) and \(k\ge3\). These bounds are uniform in the word.
\leanref{Formalism.\allowbreak{}MainNotions.\allowbreak{}LinkB.\allowbreak{}rank\_\allowbreak{}formula\_\allowbreak{}of\_\allowbreak{}part\_\allowbreak{}four}{Formalism.\allowbreak{}SharpRestriction.\allowbreak{}GlobalSchurBounds.\allowbreak{}TopDensityNGeSixteen}{102--103}
\leanref{Formalism.\allowbreak{}MainNotions.\allowbreak{}proviso\_\allowbreak{}e}{Formalism.\allowbreak{}SharpRestriction.\allowbreak{}GlobalSchurBounds.\allowbreak{}MainSup}{196--245}
\end{proof}

\subsection*{Proof of \ref{B5.32a}}\label{prf:B5.32a}\phantomsection
\begin{proof}
Fix a sufficiently small positive \(\varepsilon\). In \ref{B8.05} take the maximum of the requested threshold and the thresholds for the class-number row, \ref{B5.31a} and \ref{B5.10a}, with the outer-list substitution prescribed in \ref{B8.05}. Thus the outer constants are fixed before the common middle value \(C_F=C_H=C_\Gamma\ge1\) is chosen, and the final thresholds are evaluated on the resulting lists. The selection supplies the symmetric column of \ref{B8.09} on these same points.

Let \(L=x^\varepsilon\) and \(\Gamma=\exp(C_\Gamma m^5fL)\). By \eqref{eq:host-rank-ranges} and \(f=2n\),
\[
 m^5fL\le128x^{9/2+9\varepsilon},\qquad
 k\ge x^{9/2+10\varepsilon}.
\]
At fixed parameters, \(\log(3\Gamma^2x^{C_H})=o(k)\) uniformly at the constructed points. Since \(k\to\infty\), this logarithm is eventually at most \((k-3)\log2\). Exponentiating gives \(3\Gamma^2x^{C_H}\le2^{k-3}\). The right side is an integer, hence also bounds the ceiling. This proves the logarithmic bound required in \ref{B5.45}.
\leanref{GlobalTheory.\allowbreak{}EX.\allowbreak{}szRow\_\allowbreak{}card\_\allowbreak{}words\_\allowbreak{}le}{Development.\allowbreak{}AbelianVarietyTheory.\allowbreak{}Proofs.\allowbreak{}GlobalInterpolation.\allowbreak{}SzRowCardWordsLe}{75--104} The code in \ref{B8.05} injects exactly that many indices into the binary words of length \(k-2\) and uses the same indices at each translation. Lemma~\ref{B5.32b} gives the rank bound on this family. Since \([1/2,\exp(1)]\subset[1/3,3]\), the word count and rank bound give both inequalities in \eqref{eq:p5c3-code-input} with \(C_0=3\).

For the family satisfying the additional hypotheses of \ref{B5.31a}, put \(y=x^{\varepsilon/8}\) and \(E=e^{-y}\). Then \(L=y^8\). The displayed rank bounds give \(k\ge8\log\Gamma\) eventually. With \(r=k+m-1\ge k\), the formulas of \ref{B5.31a} yield
\[
 \eta_0,\eta_1\le E,\qquad
 e_0=2C_0/\Gamma+2e^{-kf\log k}
 \le(2C_0+2)e^{-C_\Gamma y^8}\le E^2.
\]
Its source rank-loss bound \(2f/b\) is at most \(E^2\) eventually, since \(f\) grows polynomially and \(b=e^{y^8}\). Take \(d_e=E^2,\ e_B=2E^2\). Eventually \(Ex^{C_H}\le1/18\), so \eqref{eq:p5c3-small-input} holds, \(d_e,e_B\le1/4\), and
\[
 (3\sqrt{d_e}+4\sqrt{e_B})x^{C_H}
 \le(3+4\sqrt2)/18<1.
\]
These are the common bounds required in \ref{B5.32}.
\end{proof}

\section*{Proofs for Chapter~\ref{ch:p5c4}: The pointwise bound and its eight error terms}

\subsection*{Proof of \ref{B5.33}}\label{prf:B5.33}\phantomsection
\begin{proof}
Take \(e_K\) from the size bounds and \(c_{\rm Count},c_{\rm Slope},c_1,c_{\rm Loss}\ge0\) from the family estimates. Use a common constant \(C_d\) for \ref{B5.10}, \ref{B5.13} and \ref{B5.14}. The budgets are nonnegative, so take
\[
 c_{\rm Div}=c_E=c_H=C_d^+=\max(C_d,0),\qquad
 K=C_d^++c_{\rm Count}+c_{\rm Slope}+c_1+c_{\rm Loss}+6 .
\]
Take \(x_0\) above the common thresholds and at least \(\max(x_1+1,3)\), where \(x_1\) is the family threshold. At \(x\ge x_0\), choose \(i_0,a_v,b_v,i_1\) by the loss hypothesis. The covering in \ref{B5.22} from \(i_1\) supplies a model \(i(u_0)\) at each \(u_0\in A\) realising the same pair. Use the frame and global data of \(i_0\).

At fixed \(x\) the family fixes \(\mathrm{rd}(E)\) and \(d\). Thus \(v=\mathrm{rd}(E)/2\) and the quadratic logarithmic derivative \(L_pL\) are common, so \(\log\Omega\) is common. The shared scale \(a\) gives the same \(L_a=\log(a\,\mathrm{rd}(E))\), and the family fixes \(g_0,g_1\) and hence \(\ell_g\). The common ranks, scales and thick-label count give the same \(B_0,N_c\). Identify each top set with \(TT\) by the integer labels. The weight at label \(t\) is the same value of \(a_v\) in each selected model. The sum in the division estimate of \ref{B5.10} therefore becomes
\[
 \sum_{t\in TT}
 \bigl[m(a_t\log a_t-a_t-a_t\log\Omega)-a_t\log T_u\bigr]
 = S_{\rm div}-S_a\log T_u .
\]
Its error becomes \(E_{\rm div}\), which gives \ref{hyp:p5c4-div}. The triangle inequality sums the \(m\) layer bounds of \ref{B5.13} to give \ref{hyp:p5c4-E}. The remaining estimates \ref{hyp:p5c4-H}--\ref{hyp:p5c4-N} follow from \ref{B5.14}, \ref{B5.11}, \ref{B5.12} and \ref{B5.06}.

Put \(\varepsilon_z=\varepsilon/J_{\mathcal F}>0\), with \(J_{\mathcal F}\ge1\) the family scale parameter. By \ref{B7.52}, \(0\le\delta_{i(u_0)}\le\bar\delta=\exp(-x^{\varepsilon_z/7})<1/2\). Use this same \(\delta_{i(u_0)}\) in the frame inequality, with \(\eta_{i(u_0)},e_{i(u_0)}\). By \ref{B5.10} and target duality,
\[
 s-\mu_{\min}(RR^\vee)=s+\mu_{\max}(RR)\le x^{C_H}b .
\]
Thus \(\mathcal H=x^{C_H}b>0\) gives \ref{hyp:p5c4-frame}--\ref{hyp:p5c4-parameters}. Apply Lemma~\ref{B5.21a} with the scale signs and lattice hypotheses to obtain \ref{hyp:p5c4-setting}.

Admissibility gives \(a\le a_t\le a+a/(kx^{13})\le2a\) and \(0\le a_t/b-x\le1/(kx^{12})\le1\), since \(k\ge1\) and \(x\ge3\). Lemma~\ref{B5.21a} gives the thick and other node bounds and \(w_{u_0j\rho}\ge1\) on \(\Phi\). Hence \(a_t/w_{u_0j\rho}\le N\).

The family fixes \(q,d,Q_S\) and the Artin map on units in the common frame. By \ref{B2.01}--\ref{B2.02}, reduction to the prime-to-\(p\) part of \(Q_Sqd\) gives the same inertia subgroup and Frobenius coset as for the datum \(D_{\rm ar}\) of \(i_0\). Use the top labels and the place bijection of \ref{B5.21a}, which identifies the weights at lift \(j+1\) with \(w_{u_0j}\) on \(\Phi\). Extend \(w_{u_0j}\) by \(1\) off \(\Phi\). Then
\[
 \ASK(u_0)=\sum_{t\in TT}\sum_{j=0}^{m-1}
 \ASK(D_{\rm ar},N,a_t,t,w_{u_0j}).
\]
In \ref{B2.04} each contributing \(g\notin\Phi\) has its partner \(\rho\) in \(\Phi\), for ordinary primes and for \(d\). A positive summand requires \(p^i<a_t/w_{u_0j\rho}\le N\), which implies \(p,i\le N\). The common-weight loss hypothesis gives \ref{hyp:p5c4-loss}.

The size and rank budgets, with \(f=2n\) and \(|S_2|=n_2\), give
\[
 L=x^\varepsilon,\qquad B_0=2nmbx^\varepsilon,\qquad
 N_c=\frac{n_2kbx^\varepsilon}{n(nx^{-\varepsilon})}.
\]
By \eqref{eq:p5c4-top-density}, \(n_T\ge1\). The seven coefficients in \ref{B5.45} lie in \([0,K]\), and \eqref{eq:p5c4-frame-budget} gives
\[
 (\bar\eta+4\sqrt{\bar e})x^{C_H}b
 \le6b\le Kx^\varepsilon n_Ta .
\]
We have \(b'X'=a=xb\). Enlarge \(x_0\) so that \(x^\varepsilon\) exceeds both \(\log x\) and \(200(1+K)^2(1+c_{HH})(1+e_K)\), as required in \ref{B5.45}.

Apply \ref{B2.03} to the base host with \(S=\{q\}\cup\{p:p\text{ prime},\ p\mid Q_S\}\) for the local conditions of \ref{B5.48} and \ref{B5.49}. Apply Lemma~\ref{B5.33a} to \(D_{\rm ar}\) using the hypotheses at \(i_0\). Use Euler's constant \(\gamma\) and \(L_pL=\Re\ell(\chi_C)\). The character data and the displayed sum for \(\ASK\) then give \ref{B5.46} via the top and place bijections. Lemma~\ref{B5.22} gives \(e_L=0\) for its family.
\end{proof}

\subsection*{Proof of \ref{B5.33a}}\label{prf:B5.33a}\phantomsection
\begin{proof}
Pull \(\xi\) back by the change of coordinates \(T\) of \ref{B2.02}. Both the model representative \(\chi_\xi\) and the primitive reduction of the pulled-back Artin character evaluate \(\xi\) on Frobenius at primes unramified in the host field \(E\). Induce both to the product of their conductors and \(|\operatorname{disc}E|\). Dirichlet's theorem \cite[\S4]{davenport-1980} gives a prime in each unit class at this modulus. The induced characters agree, so their primitive reductions agree. By Lemma~\ref{B2.15}, \(\chi_\xi(n)=\prod_{p\mid n}V_\xi(p^{v_p(n)})\) for every \(n\ge1\).

For \(\xi\ne\mathbf1\), set \(\ell(\xi)=L'(1,\chi_\xi)/L(1,\chi_\xi)\). Choose \(\ell(\mathbf1)=0\), since Definition~\ref{B2.04} sets \(G_{\mathbf1}=G_z\) independently of this value. Any primitive character with the prescribed values equals \(\chi_\xi\) by Lemma~\ref{B2.33}. For \(n\) coprime to \(Q_Sqd\), the equality at unramified primes gives \(\chi_\xi(n)=\xi(T(\operatorname{Art}(n)))\) by multiplicativity. This is the membership condition in \ref{B8.47}. By \ref{B2.15}, \(G_\xi=g_{\chi_\xi}\), with \(g_1=G_z\) for the primitive trivial character. This gives the thin-character bound.

On integers coprime to \(Q_Sqd\), \ref{B2.46} and \ref{B2.01} give the Artin signs \((n/q)\) and \((n/d)\). Since \(2\mid Q_S\) and \(q,d\equiv3\pmod4\), quadratic reciprocity \cite[Chap.~5]{ireland-rosen-1990} gives
\[
 \left(\frac{-p}{n}\right)
 =(-1)^{(n-1)/2}\left(\frac{p}{n}\right)
 =\left(\frac{n}{p}\right)\qquad(p=q,d).
\]
The product sign has the product of these values. For each sign let \(\chi^*\) be the primitive reduction of its character modulo \(8qd\). On the units modulo \((8qd)(Q_Sqd)\) the reciprocity identity gives \(\chi^*=\chi_\xi\), so Lemma~\ref{B2.33} identifies the primitive characters. Thus \eqref{eq:p5c4-qd-input} bounds \(G_\xi(x)\) and \(\ell(\xi)\). The \(d\)-sign cuts out \(\Q(\sqrt{-d})\) by \ref{B2.01}, so its representative is the Kronecker character \(\chi_d\) of discriminant \(-d\). The constant \(D_1=\Re(L'/L)(1,\chi_d)\) in \ref{B5.07x} gives \(L_pL=D_1=\Re\ell(\chi_C)\).

For the hosts of Theorem~\ref{B2.59}, let \(M\) be the target conductor and fix \(K=J_{\rm win}\) from Theorem~\ref{B2.48}. For sufficiently large \(Y\), Theorem~\ref{B2.43} selects \(q,d\) on \([Y,Y^K]\). In the proof of Theorem~\ref{B2.48}, the same pair determines the window containing \(x\), with estimates restricted to \([Y,x]\).

Put \(z=Y^{\varepsilon/5}\), \(z'=z(\log Y)^2\). In the two stages of \ref{B2.38} and \ref{B2.39}, the candidate \(p\) lies in \((\sqrt{U_{\rm p}},U_{\rm p}]\), with \((U_{\rm p},F)=(e^z,8M)\) and \((e^{z'},8Me^z)\), respectively, where \(F\) bounds the conductors of the primitive bases. Apply Theorem~\ref{B2.37} in those progressions with \(\tau=\varepsilon/2\) and discarded proportion \(1/2\). The maximal estimate of Theorem~\ref{B2.30} discards at most a quarter of each candidate set after summing over the bases. Lemmas~\ref{B2.25} and~\ref{B2.32} bound the proportion discarded for zeros by \(C_{\rm b}U_{\rm p}^{3/8-9/10}\le1/4\) for large \(Y\), with \(C_{\rm b}\ge0\). For each surviving prime, Lemma~\ref{B2.36} gives the window bound.
\leanref{QH.\allowbreak{}stage}{Development.\allowbreak{}AutomorphicArithmetic.\allowbreak{}MeanSquareBounds.\allowbreak{}StagePow}{41--171}

Choose \(q\) from the surviving first set, then \(d\) from the surviving second set for that \(q\). The first stage uses the primitive reductions of characters \(\xi_8\) modulo \(8M\) as bases. The second uses the primitive reductions of \(\xi_8(n)(n/q)^i\) at level \(8Mq\), \(i=0,1\). Thus the first stage controls \((i,j)=(1,0)\) and the second controls \((0,1),(1,1)\), where \(i,j\) are the exponents of the \(q\)- and \(d\)-factors. A character \(\xi\) modulo \(M\) lifts to \(\xi_8\) by reduction modulo \(M\). On integers coprime to \(8Mqd\) the corresponding twists have the same values, so Lemma~\ref{B2.33} transfers the stage estimates to the primitive representatives at level \(2Mqd\).

For each stage \(H_j=FU_{\rm p}=8Me^{z+jz'}\). A primitive base \(\kappa\) has conductor \(f_\kappa\le F\), and the odd candidate prime \(p\) exceeds \(F\). Lemma~\ref{B2.33} gives a nonprincipal primitive twist of conductor \(f_\kappa p\le H_j\). The exclusion in Lemma~\ref{B2.32} gives nonvanishing for \(\Re s>399/400\) and \(|\Im s|\le H_j\). Since \(\log(8M)\le z\) and \(\log Y\ge1\) for large \(Y\),
\[
 Y^{\varepsilon/6}\le\log H_j\le3Y^{\varepsilon/5}(\log Y)^2 .
\]
Here \(0<\varepsilon<1/52\), so \(200\varepsilon<50/13<4\). In the proof of \eqref{eq:aux-kappa-four} the remaining \(2000\log\log Y+O(1)\) is eventually at most \((4-200\varepsilon)\log Y\). Thus the refined assertion of Lemma~\ref{B2.35} applies with \(A_1=3\) and gives \(|(L'/L)(1,\chi)|\le4\log Y\).
\leanref{QH.\allowbreak{}pairEx}{Development.\allowbreak{}AutomorphicArithmetic.\allowbreak{}MeanSquareBounds.\allowbreak{}pairEx}{112--169}

Take the principal character modulo \(M\) in Definition~\ref{B2.42}. If \((a,2Mqd)=1\), then \(a\) is a unit modulo \(M\) and \(8qd\). The principal character equals \(1\) there, and reciprocity gives \(\chi^*(a)=(a/q)^i(a/d)^j\) for the exponent pair of the sign. Hence \(\chi^*\) is a primitive representative at level \(2Mqd\). Evaluating at \(Y\le x\le Y^K\) gives the bounds of \ref{B2.55}
\[
 |g_{\chi^*}(x)|\le x^{1/2+\varepsilon/2}\le x^{1/2+\varepsilon},
 \qquad
 |(L'/L)(1,\chi^*)|\le4\log Y\le4\log x .
\]
\leanref{qdq\_\allowbreak{}core}{Development.\allowbreak{}AutomorphicArithmetic.\allowbreak{}CharacterBoxes.\allowbreak{}PairRows}{264--297}
\leanref{qdq\_\allowbreak{}of\_\allowbreak{}pair}{Development.\allowbreak{}AutomorphicArithmetic.\allowbreak{}CharacterBoxes.\allowbreak{}PairRows}{299--321}
\end{proof}

\subsection*{Proof of \ref{B5.36}}\label{prf:B5.36}\phantomsection
\begin{proof}
By \ref{hyp:p5c4-N}, \(\mu_u-h_u\ge\mu_{M,u}+\Ga-\Ee_u+\Fo-h_u-\mathrm{Lo}_u\). By \ref{hyp:p5c4-div} and \ref{hyp:p5c4-count}, since \(S_a\ge0\),
\[
 \Fo-h_u\ge\mu_C-S_{\rm div}+S_a\log T_u-E_{\rm div}
 \ge\mu_C-S_{\rm div}+S_a\lambda_5(u_0)-E_{\rm div}.
\]
Use \(\Ee_u\le m(S_a\ell_g+c_EB_0)\) and \(\mu_{M,u}\ge-c_{\rm Slope}n_TaL_a(1+m/k)\) from \ref{hyp:p5c4-E} and \ref{hyp:p5c4-M}. This gives \(\mu_u-h_u\ge\Ga+\mu_C-K_5(u_0)-\mathrm{Lo}_u\). Take the \(d\)-weighted mean.
\end{proof}

\subsection*{Proof of \ref{B5.38}}\label{prf:B5.38}\phantomsection
\begin{proof}
Subtract the heights in the upper bound of \ref{hyp:p5c4-N} for \(u\) and in the lower bound for \(u'\). Their difference cancels \(\Ga+\Fo\). By \ref{hyp:p5c4-M} and \ref{B5.37},
\begin{align*}
 (\mu_u^+-h_u)-(\mu_{u'}-h_{u'})
 &\le\mu_{M,u}^+-\mu_{M,u'}+
 (-\Ee_u-h_u)-(-\Ee_{u'}-h_{u'})\\
 &\quad+c_1mB_0+\mathrm{Lo}_{u'}
 \le K_3+\mathrm{Lo}_{u'} .
\end{align*}
Average over \(u'\) with weights \(d_{u'}\) and maximise over \(u\). For the lower bound, each \(\mu_u-h_u\le\mu_u^+-h_u\) by \ref{hyp:p5c4-setting}, so their weighted mean is at most \(\max_u(\mu_u^+-h_u)\).
\end{proof}

\subsection*{Proof of \ref{B5.40}}\label{prf:B5.40}\phantomsection
\begin{proof}
The setting and the bounds on height and frame parameters give \ref{hyp:p5c4-frame}. With \ref{B5.36} this gives
\[
 \Ga+\mu_C-K_5-J\le s\le\mu_C+\log|W(u_0)|
 +\frac{4\vartheta(U,\mathcal H)+(\eta+4\sqrt e)\mathcal H}{1-\delta}.
\]
Cancel \(\mu_C\) and use \(K_5(u_0)=\Main(u_0)+K_0\) by \ref{B5.35}. By \ref{B5.38} and \ref{hyp:p5c4-HH}, \(0\le U\le K_3+J\) and \(0<\mathcal H\le Xb\). Thus \ref{B5.39} and \ref{hyp:p5c4-parameters} give
\[
 4\vartheta(U,\mathcal H)\le4\log(1+X)(K_3+J+b),\qquad
 (\eta+4\sqrt e)\mathcal H\le(\bar\eta+4\sqrt{\bar e})Xb .
\]
The numerator is nonnegative and \(1/(1-\delta)\le1/(1-\bar\delta)\).
\end{proof}

\section*{Proofs for Chapter~\ref{ch:p5c5}: Four terms and the square-root bound}

\subsection*{Proof of \ref{B5.44}}\label{prf:B5.44}\phantomsection
\begin{proof}
Put \(r=m^2/k\), \(A_1=n_Ta\), \(Y_0=y_Dy_Ly\), and
\[
 E=A_1(1+r)+m(B_0+N_c),\qquad P=Y_0E.
\]
We have \(P\ge E\ge A_1\ge a\ge1\) and \(Y_0\ge y,y_D,y_L\).
In \(R_1\), use \(S_a(1+r)\le2E\) and \(L_a\le2Y_0\).
In \(R_4\) and \(R_6\), use \(m/k\le1\) and \(mb\le a\), respectively.
The definitions give
\[
\begin{aligned}
 R_1&\le4KP,& R_2&\le3KP,& R_3&\le KP,\\
 R_4&\le4KP,& R_5&\le P,& R_6&\le KP.
\end{aligned}
\]
Since \(b'\le a/k\), we have \(n_Tm^2(a/k+b')\le2A_1r\le2E\).
The middle term of \(J_0\) is therefore at most \(4K^2P\).
The other terms satisfy \(c_1P_0+c_1X''+1\le(K^2+K+1)P\), so
\begin{equation}\label{eq:J0-bound}
 J_0\le(5K^2+K+1)P.
\end{equation}
The slope term in \(\bar U\) is at most \(6KP\). The first two
summands of \(V_0\) contribute at most \(5KP\) together, and
\(S_ac_{\rm Count}\ell_k\le2KgP\).
Also \(c_1mB_0\le KP\) and \(b\le P\).
Together with \eqref{eq:J0-bound}, these give
\[
 \bar U+b\le uP,\qquad u=5K^2+13K+2Kg+2.
\]
Since \(\tau\le KyA_1\le KP\),
\[
 R_8\le d_I\{4hy\,uP+KP\}\le2\{4hy\,uP+KP\}.
\]
Using \(y\ge1\) gives
\[
 \sum_i R_i\le (5K^2+14K+2)P+2\{4hy\,uP+KP\}
 \le (5K^2+16K+2+8hu)yP.
\]
Since \(h,g\ge1\),
\[
\begin{aligned}
 5K^2+16K+2+8hu
 &\le hg(45K^2+136K+18)\\
 &\le200(1+K)^2hg\le y.
\end{aligned}
\]
Thus \(\sum_iR_i\le y^2P=y_Dy_Ly^3E\).
\end{proof}

\subsection*{Proof of \ref{B5.45}}\label{prf:B5.45}\phantomsection
\begin{proof}
Apply Theorem~\ref{B5.44} with
\[
\begin{gathered}
 y=x^\varepsilon,\quad y_D=x^{e_{\rm Div}\varepsilon},\quad
 y_L=x^{e_{\rm Loss}\varepsilon},\quad
 g=1+e_K,\quad h=1+c_{HH},\\
 \ell_W=\log\#W(u_0),\quad \ell_k=\log k,\quad
 \ell_H=\log(1+x^{c_{HH}}),\\
 d_I=(1-\bar\delta)^{-1},\qquad
 \tau=(\bar\eta+4\sqrt{\bar e})x^{c_{HH}}b.
\end{gathered}
\]
The signs of the exponents give \(y,y_D,y_L,h,g\ge1\).
Since \(\varepsilon\le1\), we have \(m\le x\), so \(mb\le xb\le a\).
Also \(m\le xm\le k\), \(1\le b\le xb\le a\), and
\(0\le S_a\le2n_Ta\) by summation. The logarithm bounds are
\[
\begin{aligned}
 \log\#W(u_0)&\le(k-3)\log2\le k\le a,\\
 0\le\log k&\le e_K\log x\le(1+e_K)y,\\
 0\le\log(1+x^{c_{HH}})&\le\log2+c_{HH}\log x\le(1+c_{HH})y.
\end{aligned}
\]
The last line uses \(1+x^{c_{HH}}\le2x^{c_{HH}}\),
\(\log2\le1\), and \(\log x\le y\). Also \(0<d_I\le2\) by the
bound on \(\bar\delta\).
The terms \(R_1,\ldots,R_8\) sum to \(E_8(u_0)\) by
Definition~\ref{B5.42}, and \(y_Dy_Ly^3=x^{C_{\rm err}\varepsilon}\).
\end{proof}

\subsection*{Proof of \ref{B5.54}}\label{prf:B5.54}\phantomsection
\begin{proof}
Since \(x^\varepsilon x^{-\varepsilon}=1\), the left side equals
\[
 (mb)\frac{n_2k}{n^2}x^{2\varepsilon}.
\]
Since \(k\ge0\), we have \(n_2k\le nk\). Apply the upper bound on \(k\)
and then the density bound to obtain
\[
 \frac{n_2k}{n^2}x^{2\varepsilon}
 \le n x^{1/2+\varepsilon}
 \le2n_Tx^{1/2+9\varepsilon/8}\le2n_TY.
\]
Multiply by \(mb\ge0\). For \(Q>0\), the normalized term is
\[
 \frac{m}{Q}\,\frac{n_2kbx^\varepsilon}{n(nx^{-\varepsilon})}
 =\frac{n_2}{n_T}\frac{k}{n^2}x^{2\varepsilon}.
\]
\end{proof}

\subsection*{Proof of \ref{B5.55}}\label{prf:B5.55}\phantomsection
\begin{proof}
The lower bounds imply \(n_T,m,k>0\), hence \(Q>0\).
Put \(P=x^{C\varepsilon}\), \(X=x^{1/2+\varepsilon}\), and
\(Y=x^{1/2+13\varepsilon/8}\).
Then \(P,Y\ge1\) and \(X\le Y\).
Lemmas~\ref{B5.51}--\ref{B5.54} give
\[
\begin{aligned}
 \mathcal B_x(a_E)
 &=n_Ta_E+n_Ta_E\frac{m^2}{k}+mB_*+mN_*\\
 &\le(2+2+4+2)YQ=10YQ.
\end{aligned}
\]
Thus \(E/Q\le10PY\). Multiply the assumed inequality by \(2/Q\)
to obtain
\[
 \delta_2(G_x-4X)\le20PY+\frac4x\le22PY,
\]
where \(x\ge3\) and \(PY\ge1\). Since \(0<\delta_2^{-1}\le x^\varepsilon\)
and \(X\le PYx^\varepsilon\), we get
\[
 G_x\le4X+22PYx^\varepsilon
 \le26x^{1/2+(C+21/8)\varepsilon}
 \le x^{1/2+(C+3)\varepsilon}.
\]
The last inequality uses \(26\le x^{3\varepsilon/8}\).
\end{proof}

\section*{Proofs for Chapter~\ref{ch:window-inputs}: The window bound and its inputs}

\subsection*{Proof of \ref{B6.04}}\label{prf:B6.04}\phantomsection
\begin{proof}
Use Proposition~\ref{B6.05} to assemble the automorphic data from the attachment \(\mathcal R\) and the transfer proved in Lemma~\ref{B6.04a}. Remark~\ref{B6.05ca} specifies the published results used in the deductions from these data. For its descent step use Lemma~\ref{B9.36a} in ranks two through four, with the split set containing the places over the chosen prime and with spherical data outside it. These are the hypotheses on the descent input of Proposition~\ref{B6.06}. That proposition gives the assertion of Definition~\ref{B6.03}, with exponent \(401\).
\end{proof}

\subsection*{Proof of \ref{B6.04y}}\label{prf:B6.04y}\phantomsection
\begin{proof}
Regroup the diagonal of the Hecke kernel on the finite double coset space by rational conjugacy. The regular support condition makes each term separable. The trace on an automorphic summand is $d_\lambda$ times the trace on its finite fixed vectors. Compact support bounds the coefficient denominators, and the unit-circle condition bounds the conjugates of the coefficients by the binomial coefficients in rank $n$. These bounds leave finitely many polynomials.

Exclude the test places and primes dividing coefficient denominators. Also exclude primes dividing the determinants or discriminants. At the remaining places the order generated by the regular norm is finite \'etale over the local integer ring. Its factors are unramified integer rings, and every full stable lattice is a product of their fractional ideals. At a split place the centralizer translates any such lattice to the standard one. At an inert place use the trace description of the forms and lattices proved in \ref{B6.04e}. On a field fixed by the involution, the trace dual of the integer ring is itself. Self-duality for coefficient $c$ thus requires the ideal exponent $j$ to satisfy $2j=-v(c)$. This leaves one contributing norm class, since unit norms are surjective. On exchanged factors the two ideal exponents have prescribed sum, and the norm-one group acts transitively on pairs with that sum. The same lattice conditions apply to an integral twisted representative by the unit coefficient construction in the proof of \ref{B6.04e}. Thus in each case the stable unit integral has one contributing integral conjugacy orbit. The stabilizer is the integral centralizer, so its contribution is one under the volume conventions. Classes without an integral orbit contribute zero.

For a fixed elliptic $P$, the algebra $E[X]/P$ has involution given by conjugation on $E$ and inversion of the generator. On each field factor an invariant form is $\operatorname{Tr}(c\,x\bar y)$ with $c$ fixed by the involution and specified modulo norms. Indeed, pairing with $1$ defines the trace functional and the adjoint relation determines the form. Let $u_j$ be the product of the local determinant signs on factor $j$. Quadratic norm reciprocity realizes the prescribed signs with positive infinity type exactly when all $u_j=1$. The orthogonal sum is then locally isometric to the standard form, hence globally isometric to it by the local-global theorem for Hermitian forms. The same theorem on each factor gives uniqueness of the rational tuple \cite{scharlau-1985}. Summing the sign over $P_1\mid P$ gives $\prod_j(1+u_j)$, so the surviving tuples acquire multiplicity $2^{r(P)}$.

To compare masses, take a matrix $N$ over $E$ with characteristic polynomial $P$. Conjugate inversion preserves $P$, and separability implies that $(N^{-1})^*$ and $N$ are conjugate, where ${}^*$ denotes conjugate transpose. Thus there is an invertible $X$ with $NXN^*=X$. Put $\epsilon=(-1)^{n-1}$ and choose $c\in E^\times$ with $c/\bar c$ outside the spectrum of $-\epsilon NX^*X^{-1}$. The possible ratios form an infinite set. Then
\[
 M=cX+\epsilon\bar cNX^*,\qquad NM^*=\epsilon M,\qquad
 \delta=MJ_n^t .
\]
The choice of $c$ makes $M$ invertible. Since $J_n^t=\epsilon J_n$ and $J_nJ_n^t=1$, $\delta\theta_J(\delta)=\epsilon M(M^*)^{-1}=N$. The twisted centralizer consists of $x$ with $xMx^*=M$. Such $x$ commute with $N$, and the adjoint on $E[N]$ sends scalars to their conjugates and $N$ to $N^{-1}$. Thus the centralizer is the same norm-one torus as for a unitary representative. This identifies rational points and maximal compact subgroups, so the unitary mass is $2^{r(P)}\kappa(P)$. For coprime $P_1,P_2$ the tori and their mass factors split. The geometric sum can therefore be written over ordered factorizations $P=P_1P_2$ with summand
\[
 \kappa(P_1)\kappa(P_2)h_\lambda(P)
 \prod_{w\in\Sigma}O_w^{P_1}(P;f_w).
\]
\leanref{GlobalTheory.\allowbreak{}PreStabU.\allowbreak{}fixedP}{Development.\allowbreak{}WeilBound.\allowbreak{}CoreStatements.\allowbreak{}AutomorphicKernels.\allowbreak{}IsobaricBaseChange.\allowbreak{}IntegralCoefficientBaseChange.\allowbreak{}UnitaryBaseChangeCharacters.\allowbreak{}UnitaryBaseChangeLinks.\allowbreak{}UnitaryGeometricSideSeparableTest}{357--435}

Inverting the two root multisets gives $R/\bar R=(-1)^{ab}D$ for the global resultant. At a split pair $\Delta_w$ is the product of $\mu_u(R)|R|_u^{1/2}$ over the pair, by the restriction of $\mu$ to the base. At an inert place this expression is the stated factor. Write $R_v$ for the image of $R$ at the chosen embedding over $v$. The product formulas give
\begin{equation}\label{eq:p6c1-resultant-product}
 \prod_{w\in\Sigma}\Delta_w
 \prod_{v\mid\infty}(R_v/|R_v|)^{e_v}|R_v|=1 .
\end{equation}
Expand the Weyl numerator at $v$ along the first $a$ rows. The minor in columns $j_0<\cdots<j_{a-1}$ has exponents $\lambda_{v,j_i}+n-1-j_i$, hence highest weights $\lambda_{v,j_i}+b+i-j_i$ after division by the block denominator. The complementary minor has the formula with $a$ in place of $b$. Put $h_v=(e_v+1)/2$. The cross denominator in the conjugated character gives
\[
 \frac{(R_v/|R_v|)^{e_v}|R_v|}{\overline R_v}
 =((-1)^{ab}D_v)^{h_v}.
\]
This factor subtracts $bh_v$ and $ah_v$ from the respective highest weights in the conjugated weight characters. In infinitesimal parameters the shifts are $-be_v/2$ and $-ae_v/2$. The local sign is $(-1)^{\sum_{j\in I_v}j+a(a-1)/2+abh_v}$. Multiply over $v$ and use \eqref{eq:p6c1-orbital-match} and \eqref{eq:p6c1-resultant-product}. An empty block gives the full minor with sign $+1$.
\leanref{UpperTheory.\allowbreak{}ConvolutionTheory.\allowbreak{}G2.\allowbreak{}termId}{Development.\allowbreak{}AutomorphicArithmetic.\allowbreak{}Foundations.\allowbreak{}Convolution.\allowbreak{}WeightRankUnits.\allowbreak{}TermId}{24--41}
\leanref{UpperTheory.\allowbreak{}ConvolutionTheory.\allowbreak{}G2.\allowbreak{}realProdS}{Development.\allowbreak{}AutomorphicArithmetic.\allowbreak{}Foundations.\allowbreak{}Convolution.\allowbreak{}WeightRankUnits.\allowbreak{}RealProdS}{31--59}
\leanref{UpperTheory.\allowbreak{}ConvolutionTheory.\allowbreak{}G2.\allowbreak{}realSigned}{Development.\allowbreak{}AutomorphicArithmetic.\allowbreak{}Foundations.\allowbreak{}Convolution.\allowbreak{}WeightRankUnits.\allowbreak{}RealSigned}{9--154}
\end{proof}

\subsection*{Proof of \ref{B6.04ea}}\label{prf:B6.04ea}\phantomsection
\begin{proof}
Beuzart-Plessis's linearized Jacquet--Rallis theorem equates the orbital integrals of the integral matrix indicators for a relatively regular $F$-matrix of size $r+1$ and its match on the standard Hermitian space over $L$ \cite[Thm.~1]{beuzart-plessis-2021}. Matching is conjugacy by $\GL_r(L)$ embedded in the upper left block. The linear integral vanishes when the standard Hermitian orbit is absent. The integral subgroups of $\GL_r(F)$ and $U_r(F)$ have volume one. The linear integrand at $Y$ carries $\eta$ of the determinant with rows $e^*,e^*Y,\ldots,e^*Y^r$, where $e^*$ is the last coordinate row. Here $r$ indexes the acting group.
\leanref{GlobalTheory.\allowbreak{}JR.\allowbreak{}unitMatch\_\allowbreak{}all}{Development.\allowbreak{}AbelianVarietyTheory.\allowbreak{}Proofs.\allowbreak{}InertiaEigenvalues.\allowbreak{}JacquetRallisUnitMatchAllRanks}{9--11}

Put $u=cw'\sigma w'$, and let $A$ be multiplication by $g$. In the power basis take the linear matrix with blocks $(A,u,\operatorname{Tr}_{B/L},0)$. On $(B,h_c)$ use $(A,w',h_c(\,\cdot\,,w'),0)$. Multiplication by $c\sigma w'$ identifies the marked columns and rows. The moment matrix $(\operatorname{Tr}_{B/L}(ug^{i+j}))_{i,j<r}$ is invertible, so the column is cyclic and the row is cocyclic. The moments determine the form on the basis $g^iw'$, so the standard match exists exactly when $(B,h_c)$ is standard.

The linear cosets correspond to $\sigma$-stable lattices. Choose $z\in O$ with $z-\bar z$ a unit, so $O=O_F+zO_F$. For $m\in M$, both $b=(m-\sigma m)/(z-\bar z)$ and $m-zb$ are fixed by $\sigma$, which gives descent to $O_F$. The matrix blocks impose $gM\subset M$ and $u\in M$. The dual condition is $1\in M^*$. Subtract earlier rows from $e^*Y^j$ to obtain the trace rows $\operatorname{Tr}(g^i\,\cdot\,)$ on the first block. Expansion along $e^*$ gives $s_g(M)$.

On the Hermitian side every unimodular lattice has an orthonormal basis. In its reduction some vector has nonzero length, since polarization with residual scalar multiples and the nondegenerate finite-field trace pairing would otherwise make the reduced form zero. Lift such a vector and normalize its length by a unit norm. Unit norms are surjective by the finite-field norm and successive lifting with the residue trace. Split off the resulting line and iterate. This works also in residue characteristic two. The integral cosets then correspond to the self-dual lattices for which $gM\subset M$ and $w'\in M$. Self-duality gives the covector condition. Each contributes one. If $(B,h_c)$ has no standard match, both the lattice count and the linear integral vanish. This proves \eqref{eq:p6c1-marked}.
\end{proof}

\subsection*{Proof of \ref{B6.04eb}}\label{prf:B6.04eb}\phantomsection
\begin{proof}
Put $V=\operatorname{Hom}_L(B_2,B_1)$ and $H=\operatorname{Hom}_O(M_2,M_1)$. Lifting a basis of $M_2$ writes $M=\{(x,y):y\in M_2,\ x-f(y)\in M_1\}$ for a unique $f\bmod H$ in $V/H$. Stability says $D(f)=g_1f-fg_2\in H$. Over a splitting field the eigenvalues of $D$ are $\alpha-\beta$ for the roots of the two polynomials, so $\det D=R$. The parameter group $D^{-1}H/H\simeq H/DH$ has cardinality $q^{2v(R)}$ by elementary divisors. In the self-adjoint case $\tau(f)=\sigma_1f\sigma_2$ defines a semilinear involution on the parameter group. With $z$ as in \ref{B6.04ea}, each class is uniquely $a+zb$ with $a,b$ fixed, by $b=(f-\tau f)/(z-\bar z)$ and $a=f-zb$. The invariant lattices correspond to the fixed classes, whose cardinality is the square root of the total.
\end{proof}

\subsection*{Proof of \ref{B6.04ec}}\label{prf:B6.04ec}\phantomsection
\begin{proof}
A stable full lattice forces its generator to be integral by Cayley--Hamilton. If either generator is nonintegral, both sides vanish. For an integral generator in a block or the product, let $t$ be a uniformizer of $F$ and choose $m\ge0$ with $t^mD_B\subset O[g]\subset D_B$. For a stable lattice $M$ set $I_M=D_BM$. Then
\begin{equation}\label{eq:p6c1-conductor-sandwich}
 t^m I_M\subset M\subset I_M .
\end{equation}
If $M$ is self-dual for $h_c$, the $h_c$-dual $I_M^\#$ satisfies $t^{2m}I_M\subset I_M^\#\subset I_M$. On each factor the dual exponent is a fixed integer minus the exponent on the conjugate factor. Thus a fixed field has bounded exponent and an exchanged pair has bounded sum of exponents. Norm-one elements translate the exponents on a pair by any $j,-j$. This gives finitely many $I_M$ modulo $T_B$, and the finite quotient in \eqref{eq:p6c1-conductor-sandwich} gives finitely many lattices above each. The argument also applies to an integral norm-one generator.

To remove the mark, let $\rho_B(k)$ count cosets $uT_B^0$ with $u^{-1}t^k\in D_B$. A norm-one element has valuation zero on a fixed field, and every pair $j,-j$ occurs on exchanged fields. Let $b_i$ be the normalized valuation of $t$ in one field of the $i$th exchanged pair. Then
\begin{equation}\label{eq:p6c1-rho}
 \rho_B(k)=\prod_{i\ {\rm exchanged}}(2kb_i+1)\qquad(k\ge0).
\end{equation}
For fixed $c$, choose $b\ge0$ integral with $t^bD_B\subset M\subset t^{-b}D_B$ for all the finitely many orbit representatives in $X_B(c)$. Let $A_M(k)$ count the translates of $M$ containing $t^k$. The translates are indexed by $T_B/(T_B)_M$, and the membership condition is $u^{-1}t^k\in M$. For $k\ge b$ this gives
\begin{equation}\label{eq:p6c1-translates}
 [T_B^0:(T_B)_M]\rho_B(k-b)
 \le A_M(k)\le [T_B^0:(T_B)_M]\rho_B(k+b).
\end{equation}
The ratios $\rho_B(k\pm b)/\rho_B(k)$ tend to one. Sum and use \eqref{eq:p6c1-marked} with $w'=t^k$ to get
\begin{equation}\label{eq:p6c1-marked-limit}
 \lim_{k\to\infty}\frac{G_B(ct^{2k})}{\rho_B(k)}=N_B(c).
\end{equation}
\leanref{GlobalTheory.\allowbreak{}UnitLemma.\allowbreak{}unitaryLimit}{Development.\allowbreak{}AbelianVarietyTheory.\allowbreak{}Proofs.\allowbreak{}InertiaEigenvalues.\allowbreak{}UnitaryLatticeCountRatioTendsto}{16--81}

Let $S$ index the fields fixed by $\sigma$. On such a field $H$ the extension $H/H^\sigma$ is unramified because the residue involution acts nontrivially on the residue field of $L$. If $\#k_H=q^{2d}$, the $q^d$-power map restricts to the $q$-power map on $k_L$, so $d$ is odd. Let $\pi_i$ be a uniformizer of $H^\sigma$ on factor $i$ and one elsewhere. Then $\eta(N_{B/L}\pi_i)=(-1)^d=-1$. The $\pi_J=\prod_{i\in J}\pi_i$ for $J\subset S$ represent $C_B$. On fixed fields the norms have even valuation and exhaust the units. On an exchanged pair every fixed invertible element is a norm.

For $I\subset S$ define $D_I f(y)=\sum_{J\subset I}(-1)^{|J|}f(\pi_Jy)$ and put
\[
 H_B(I,k)=\sum_M s_g(M)D_I1_M(t^{2k})D_{S\setminus I}1_{M^*}(1),
\]
over full $g$- and $\sigma$-stable lattices. Each expanded term bounds $M$ below by $O[g]$ times an invertible vector and above by the trace dual of another such lattice, so the sums are finite. For $J'\subset S\setminus I$, the change $M\mapsto\pi_{J'}M$ moves $\pi_{J'}$ from the dual test to the vector test and multiplies the basis sign by $(-1)^{|J'|}$. It follows that
\begin{equation}\label{eq:p6c1-alternating-count}
 H_B(I,k)=\sum_{J\subset S}(-1)^{|I\cap J|}G_B(\pi_Jt^{2k}).
\end{equation}
By \eqref{eq:p6c1-marked-limit}, the limit after division by $\rho_B(k)$ is the signed sum of $N_B(\pi_J)$.
\leanref{GlobalTheory.\allowbreak{}UnitLemma.\allowbreak{}Desc.\allowbreak{}HH\_\allowbreak{}eq\_\allowbreak{}alternating\_\allowbreak{}GG}{Development.\allowbreak{}AbelianVarietyTheory.\allowbreak{}Proofs.\allowbreak{}InertiaEigenvalues.\allowbreak{}SignedInvolutionStableLatticeSumEqAlternatingSum}{37--170}

Write $S_i$ for the fixed factors of $B_i$. Use $I=S_1$ in $B$ and in $B_1$, and $I=\varnothing$ in $B_2$. For a product lattice put $M_1=M\cap B_1$ and $M_2=\operatorname{pr}_2M$. These are full stable lattices in the two blocks. A basis of $M_1$ and lifts of a basis of $M_2$ give
\begin{equation}\label{eq:p6c1-basis-sign}
 s_g(M)=\eta(R)s_{g_1}(M_1)s_{g_2}(M_2).
\end{equation}
Indeed the remainder map from polynomials of degree $<a+b$ to the two polynomial quotients has determinant $\pm R$, where $a,b$ are the degrees of $P_{g_1},P_{g_2}$. Use the basis $1,\ldots,X^{a-1},P_{g_1},\ldots,X^{b-1}P_{g_1}$, for which the second diagonal block is multiplication by $P_{g_1}$ modulo $P_{g_2}$. The trace evaluations in the lattice basis then have the two block trace determinants on the diagonal.

Index the involution orbits of field factors by $j$ and let $D_j$ be the integer ring product on orbit $j$. The exponents of $D_BM$ on an exchanged pair agree by $\sigma$-stability. Round each orbit exponent down to a multiple of the valuation of $t$. By \eqref{eq:p6c1-conductor-sandwich} there are integers $p_j$ with
\begin{equation}\label{eq:p6c1-coordinate-sandwich}
 \bigoplus_j t^{p_j+N}D_j\ \subset M\ \subset\ \bigoplus_j t^{p_j}D_j,
 \qquad N=m+1.
\end{equation}
The restricted tuples give the same inclusions for $M_1,M_2$. Each tuple has at most $A$ lattices, where $A$ is the number of subgroups of $D_B/t^ND_B$.

For a constant $C$ independent of $k$, a nonzero product of differences requires
\begin{equation}\label{eq:p6c1-lattice-window}
 -C\le p_j\le2k+C,\qquad
 p_j\ge2k-C\ (j\in I),\qquad
 p_j\le C\ (j\in S\setminus I).
\end{equation}
Some $\pi_Jt^{2k}$ belongs to $M$ and some $\pi_{J'}$ belongs to $M^*$. The upper inclusion and the dual of the lower inclusion bound these two memberships. The trace dual of each $D_j$ is a fixed fractional ideal, and the $\pi_J$ form a finite set. If $p_j<2k-C$ for $j\in I$, the lower inclusion puts each $j$-component of $\pi_Jt^{2k}$ in $M$, so toggling $j$ cancels the vector difference. If $p_j>C$ for $j\in S\setminus I$, the dual of the upper inclusion does the same for the covector difference. Choose $C$ to cover the bounds in both projections as well.

For $k>2C$ consider a tuple in \eqref{eq:p6c1-lattice-window} with $C\le p_j\le2k-C$ on exchanged pairs. With $C$ increased for the fixed trace duals and multipliers, \eqref{eq:p6c1-coordinate-sandwich} puts all vector components outside $S_1$ in $M$ and all covector components outside $S_2$ in $M^*$. The vector test then uses $M_1$, and the covector test uses the trace dual of $M_2$. The same discards in each block give the product of the separate tests. This applies wherever either test product is nonzero, because the separate tests impose \eqref{eq:p6c1-lattice-window} on the restricted tuples.

If $h$ is the number of exchanged pairs and $h>0$, at most $C'(1+k)^{h-1}$ tuples can contribute to the difference. The fixed coordinates have boundedly many choices, and one exchanged coordinate lies within bounded distance of an endpoint. Each tuple has at most $A$ lattices and each test product has absolute value at most $2^{|S_1|+|S_2|}$. By \ref{B6.04eb} each pair of block lattices has $q^{v(R)}$ invariant lifts. With \eqref{eq:p6c1-basis-sign} this gives
\begin{equation}\label{eq:p6c1-boundary-count}
 \left|H_B(S_1,k)-(-q)^{v(R)}H_{B_1}(S_1,k)H_{B_2}(\varnothing,k)\right|
 \le 2^{1+|S_1|+|S_2|}AC'(1+k)^{h-1}.
\end{equation}
For $h=0$ the two tests agree for $k>2C$, so the left side is zero. In either case division by $\rho_B(k)=\rho_{B_1}(k)\rho_{B_2}(k)$ makes the difference tend to zero.
\leanref{GlobalTheory.\allowbreak{}UnitLemma.\allowbreak{}Desc.\allowbreak{}term\_\allowbreak{}window}{Development.\allowbreak{}AbelianVarietyTheory.\allowbreak{}Proofs.\allowbreak{}InertiaEigenvalues.\allowbreak{}AlternatingDifferenceTermWindow}{217--329}
\leanref{GlobalTheory.\allowbreak{}UnitLemma.\allowbreak{}Desc.\allowbreak{}hh\_\allowbreak{}prod\_\allowbreak{}tendsto}{Development.\allowbreak{}AbelianVarietyTheory.\allowbreak{}Proofs.\allowbreak{}InertiaEigenvalues.\allowbreak{}TendstoSignedLatticeSumProductSubFactors}{63--202}

For $c=(c_1,c_2)$ represented by $\pi_J$, the product sign in \eqref{eq:p6c1-alternating-count} is $\eta(N_{B_1/L}c_1)$. In a common basis $\det h_{c_1}=N_{B_1/L}(c_1)\det h_1$. A contributing class on $B_1$ has a self-dual lattice, so its Gram determinant has even valuation. The limits for $B_1$ and $B$ are therefore $\eta(\det h_1)S(B_1)$ and $\eta(\det h_1)K(B_1,B_2)$, with $h_1$ on $B_1$ in both. The limit for $B_2$ is $S(B_2)$. Cancel the common sign.
\leanref{GlobalTheory.\allowbreak{}UnitLemma.\allowbreak{}LieEnd.\allowbreak{}unitLemmaLie\_\allowbreak{}of}{Development.\allowbreak{}AbelianVarietyTheory.\allowbreak{}Proofs.\allowbreak{}InertiaEigenvalues.\allowbreak{}UnitLemmaSelfAdjointLieOfCountLimitDescent}{137--198}
\end{proof}

\subsection*{Proof of \ref{B6.04ed}}\label{prf:B6.04ed}\phantomsection
\begin{proof}
Nonintegral generators give zero as in \ref{B6.04ec}. For integral generators set $g=(g_1,g_2)$ and $R_0=O[g]$ in $B=B_1\times B_2$. Since $g^{-1}=\sigma g$ is integral, multiplying its monic equation by the power of $g$ one less than its degree gives $g^{-1}\in R_0$. Hence $\sigma R_0=R_0$. Induct on $\dim_L B$ from the zero algebra.

If $R_0$ has a proper nonzero $\sigma$-fixed idempotent $e=f(g)$ with $f\in O[X]$, every stable lattice splits into its $e$ and $1-e$ parts. The two parts are orthogonal for $h_c$. The parameters split as a product, as do the norm-one groups and stabilizers. The Gram signs multiply. The two parts have coprime reductions of their characteristic polynomials since $f$ takes values one and zero on their respective reduced roots. Thus the cross resultants between parts are units. The counts multiply and the resultant valuations add. Apply the induction hypothesis to both parts.
\leanref{GlobalTheory.\allowbreak{}UnitLemma.\allowbreak{}OrbitAll.\allowbreak{}identity\_\allowbreak{}integral}{Development.\allowbreak{}AbelianVarietyTheory.\allowbreak{}Proofs.\allowbreak{}DivisorProofs.\allowbreak{}UnitLemmaUnitaryGroupEveryResidueOfLie}{93--144}

The Artinian reduction $R_0/tR_0$ is a product of local rings. Its idempotents lift by Hensel's lemma for $X^2-X$. Conjugation permutes the lifted local factors. Without a proper nonzero fixed idempotent, there are either two exchanged rings or one local ring.
\leanref{GlobalTheory.\allowbreak{}UnitLemma.\allowbreak{}OrbitAll.\allowbreak{}one\_\allowbreak{}orbit}{Development.\allowbreak{}AbelianVarietyTheory.\allowbreak{}Proofs.\allowbreak{}DivisorProofs.\allowbreak{}UnitLemmaUnitaryGroupEveryResidueOfLie}{66--92}

In the exchanged case choose $e$ with $e+\sigma e=1$. On each block and the product the norm-one group projects isomorphically to the invertible group on the $e$ part, with inverse $u\mapsto(u,\sigma(u)^{-1})$. Every fixed invertible parameter is a norm. With $c=1$ and the second part transported by $\sigma$, a self-dual stable lattice is $(M,M^*)$, since the halves pair by the ordinary trace. Conversely any full $eg$-stable lattice gives such a pair, since $(eg)^{-1}$ belongs to the order it generates. The stabilizer indices are those for the action on the first part.

Write $v_+,v_-$ for the within-part cross-resultant valuations of the two blocks. The mixed resultants are units by the idempotent polynomial. The first count in \ref{B6.04eb} on the $e$ part gives
\begin{equation}\label{eq:p6c1-exchanged-count}
 N_B(1)=q^{2v_+}N_{B_1}(1)N_{B_2}(1).
\end{equation}
For the weighted count, fix a pair of lattice-orbit representatives with stabilizer $H$. Their lifts form a finite $H$-set. A lift orbit has weight equal to the pair's weight times $[H:H_M]$, where $H_M$ stabilizes the lift. Summing gives the pair's weight times the fibre size. Counting on the other part gives \eqref{eq:p6c1-exchanged-count} with $v_-$ instead. If the product count vanishes, so does $N_B(1)$. Otherwise $v_+=v_-$ and the self-dual lattice on the first block gives Gram sign one. The required multiplier is $q^{2v_+}=(-q)^{v_++v_-}$.
\leanref{GlobalTheory.\allowbreak{}UnitLemma.\allowbreak{}Swap.\allowbreak{}identity\_\allowbreak{}of\_\allowbreak{}swap}{Development.\allowbreak{}AbelianVarietyTheory.\allowbreak{}Proofs.\allowbreak{}InertiaEigenvalues.\allowbreak{}InvolutiveAlgebraProductIdentity}{979--1103}

Suppose $R_0$ is local with residue field $k'$ of size $q^{2d}$, generated over $k_L$ by the residue of $g$. The residue involution is the $q^d$-power map, so $d$ is odd. The order contains the integers $O'$ of the unramified extension $L'/L$ of degree $d$. Lift a generator of $k'^\times$ to a root of $X^{q^{2d}-1}-1$ and adjoin it. Its minimal polynomial on the residue field lifts as a factor with unit derivative. By uniqueness, conjugation acts on the lifted roots by the $q^d$-power map. Let $s\in O'^\times$ lift the residue of $g$ in this group of roots, so $s\bar s=1$. The fixed field $F'$ has residue cardinality $q'=q^d$.

The inclusion $O'\subset R_0$ gives $L'$-scalars on the algebras. The image of $g-s$ lies in the maximal ideal of the Artinian local ring $R_0/tR_0$, so it is nilpotent. The projected orders give full $O'$-lattices with the same nilpotence property. Thus the characteristic polynomials $P_i'$ over $L'$ reduce to $(X-\widetilde s)^{\deg P_i'}$.
\leanref{GlobalTheory.\allowbreak{}UnitLemma.\allowbreak{}UnrAny.\allowbreak{}exists\_\allowbreak{}unr\_\allowbreak{}embedding\_\allowbreak{}any}{Development.\allowbreak{}AbelianVarietyTheory.\allowbreak{}Proofs.\allowbreak{}DivisorProofs.\allowbreak{}ExistsUnramifiedEmbeddingOfInvolutionAlgebra}{52--220}

Choose $x\in O'$ with residue outside $k_{F'}$ and put $y=\bar x$. Let $\xi=sx/y$. The elements $x,y,x-y,\xi$ are units. Also $\xi\bar\xi=1$ and $\widetilde\xi\ne\widetilde s$. Put
\[
 Y_i=\frac{yg_i-x\xi}{g_i-\xi},\qquad Y=(Y_1,Y_2).
\]
For $P=P_1'P_2'$ and for each block polynomial, $P(\xi)$ is a unit. Dividing $P(X)-P(\xi)$ by $X-\xi$ shows that the denominator is invertible in the order. Conjugating the fraction and using $\sigma_i g_i=g_i^{-1}$ gives $\sigma_i Y_i=Y_i$. For the reverse inclusion let $m=\deg P$ and form
\[
 H(Z)=(Z-y)^m P\!\left(\frac{\xi Z-x\xi}{Z-y}\right)\in O'[Z].
\]
The relation $g(Y-y)=\xi Y-x\xi$ gives $H(Y)=0$, while $H(y)=(\xi(y-x))^m$ is a unit. Division of $H(Z)-H(y)$ by $Z-y$ gives $(Y-y)^{-1}\in O'[Y]$. Thus the orders generated by $g$ and $Y$ agree on the product and on the blocks. The invertible substitution preserves distinct roots.

Over $O'$ both generators have the same stable lattices. The forms and weights also agree. Fix stable full lattices in the blocks, for example the projected orders. Their extension lattices are the same for both generators. By \ref{B6.04eb} their count is $(q'^2)^v$ for the respective cross-resultant valuation $v$, so the valuations agree. Apply \ref{B6.04ec} over $L'/F'$ to $Y_1,Y_2$.
\leanref{GlobalTheory.\allowbreak{}UnitLemma.\allowbreak{}CayleyAny.\allowbreak{}identity\_\allowbreak{}of\_\allowbreak{}residually\_\allowbreak{}unipotent\_\allowbreak{}any}{Development.\allowbreak{}AbelianVarietyTheory.\allowbreak{}Proofs.\allowbreak{}DivisorProofs.\allowbreak{}UnitLemmaResiduallyUnipotentCayleyTransform}{107--224}

Every $g$-stable $O$-lattice carries $O'$. Its Hermitian duals over $O$ and $O'$ agree by trace transitivity and trace self-duality of $O'/O$. The norm-one groups and integral units are the same. For a Gram matrix $H'$ of size $r'$ over $L'$ and an integral basis $(b_j)_j$ of $O'/O$, the determinant over $L$ is
\[
 \det\bigl(\operatorname{Tr}_{L'/L}(b_j\bar b_k)\bigr)^{r'}N_{L'/L}(\det H').
\]
The first determinant is a unit and the norm multiplies valuation by the odd integer $d$. Restriction of scalars gives $P_{g_i}=\prod_{\tau\in\Gal(L'/L)}\tau P_i'$. For $\tau\ne\nu$, the reductions of $\tau P_1'$ and $\nu P_2'$ have distinct roots, since $\widetilde s$ generates $k'/k_L$. Degree zero contributes resultant one. The $d$ diagonal resultants have the same valuation $v'$, so the exponent over $L$ is $dv'$. Since $(-q)^{dv'}=(-q')^{v'}$, the multipliers agree.
\leanref{GlobalTheory.\allowbreak{}UnitLemma.\allowbreak{}UnrAny.\allowbreak{}unitLemma\_\allowbreak{}local\_\allowbreak{}of\_\allowbreak{}lie\_\allowbreak{}any}{Development.\allowbreak{}AbelianVarietyTheory.\allowbreak{}Proofs.\allowbreak{}DivisorProofs.\allowbreak{}UnitLemmaLocalAnyResidueCharOfLieAlgebra}{58--143}
\leanref{GlobalTheory.\allowbreak{}UnitLemma.\allowbreak{}Unr.\allowbreak{}ResVal.\allowbreak{}val\_\allowbreak{}resultant\_\allowbreak{}unr}{Development.\allowbreak{}AbelianVarietyTheory.\allowbreak{}Proofs.\allowbreak{}InertiaEigenvalues.\allowbreak{}ValResultantEqDegreeMulUnramifiedExt}{91--218}
\end{proof}

\subsection*{Proof of \ref{B6.04c}}\label{prf:B6.04c}\phantomsection
\begin{proof}
In rank $d$ put $p(j)=\tau j(2d+1-j)$. The cut requires $\deg_x W\le p(\dim W)$ for every rational subspace $W$ and $\deg_x E^d\ge-\tau$. Here degree is the negative logarithm of the adelic norm of a rational basis wedge measured at $x$. Write $\mathcal C_\tau$ for the product of the cuts and $K_h(x,y)=\sum_\gamma h(x^{-1}\gamma y)$ for the periodized kernel. If the logarithmic support size is at most $r$ in each factor, then
\begin{equation}\label{eq:p6c1-cut-trace}
 \mathcal E(h)=\int_{\mathcal C_\tau}K_h(x,\theta x)\,dx
 \qquad\bigl(\tau\ge C_0(1+r)\bigr)
\end{equation}
for a field- and rank-dependent $C_0$.

 Adelic reduction gives a Siegel decomposition $\gamma x=b\,a(t)k_0$ in each rank \cite{borel-1963}. Here $b$ lies in a compact upper triangular set. The matrix $a(t)$ is positive diagonal and $k_0$ lies in the maximal compact. The simple-root inequalities put $a(t)^{-1}ba(t)$ in a fixed compact. Taking $m_i=\lfloor t_i\rfloor$ gives cells $a(m)\Omega_0$ for a compact $\Omega_0$. There the coordinate line degrees are $[E:\Q]m_i+O(1)$ and the total degree is $[E:\Q]\sum_i m_i+O(1)$. These formulas follow from the action of $a(m)$ on the coordinate wedges and bounded operator norms for $o,o^{-1}$ with $o\in\Omega_0$. The cut bounds each $m_i$ above by $O(1+\tau)$ and bounds the sum below by $-O(1+\tau)$. Subtract the other upper bounds to bound each $m_i$ below. The cut is closed by continuity of each degree, and it lies in the images of finitely many compact cells.

For any representative $h_0$ of a point in a cell, rational changes of basis transport the subspaces whose degrees are measured. A rational line has degree at the identity bounded above, since a nonzero coordinate of its basis vector has product of absolute values one. Write $L(g)$ for the logarithmic size of $g$ defined in the statement. Comparison by exterior powers gives $\deg_{h_0}W\le C(1+L(h_0))$ for every rational line $W$. Applying the determinant estimate also to the inverse gives $\deg_{h_0}E^d\ge-C(1+L(h_0))$. The same coordinate and subtraction argument bounds $\|m\|_\infty$ by $O(1+L(h_0))$. For a class $x_d$ represented by $g$ in rank $d$, put $\rho_d(x_d)=\inf_{\gamma\in\GL_d(E)}L(\gamma g)$. Taking the infimum gives the first bound below. The representative $a(m)o$ and the cut bounds give the second, with $H(x)=1+\rho_a(x_a)+\rho_b(x_b)$.
\begin{equation}\label{eq:p6c1-reduction}
 1+\|m\|_\infty\ll 1+\rho_d(x_d),
 \qquad H(x)\ll1+\tau\quad(x\in\mathcal C_\tau).
\end{equation}

Unfold the compact cut by rational twisted classes. In one factor let $g$ represent a class and put $A=g\theta(g)$ and $\operatorname{rad}_g W=W\cap W^{\perp_g}$. On the support of the orbital integrand,
\begin{equation}\label{eq:p6c1-degree-bounds}
 \deg_x(AW)\ge\deg_x W-D,\qquad
 |\deg_x E^d|\le D,\qquad
 \deg_x W\le\deg_x\operatorname{rad}_gW+D\quad(AW=W),
 \qquad D=C_d(1+r).
\end{equation}
For $j=\dim W$ and a basis wedge $v$ of $W$, use
$\|(\bigwedge^j A)v\|_x\le\|\bigwedge^j(x^{-1}Ax)\|\|v\|_x$ at each place, take negative logarithms and sum. The two factors of $x^{-1}Ax$ are $x^{-1}g\theta(x)$ and its $\theta$-image, so their logarithmic norms are bounded by the support size. The determinant of $x^{-1}g\theta(x)$ bounds the total degree by the product formula for $\det g$. For stable $W$, the pairing with orthogonality $W^{\perp_g}$ induces a nondegenerate pairing on the quotient by the radical. Its determinant on a rational basis has product of absolute values one. The determinant bound in quotient norms gives an upper bound $D$ for the quotient degree. Multiplicativity of subspace and quotient covolumes gives the last inequality.

At fixed $x$ the subspace degrees are bounded above by the exterior-power comparison with the identity. Put $b(W)=\deg_x W-p(\dim W)$ and $s=\sup_W b(W)$. If $s>0$, take $W$ within $\delta$ of the supremum, for $0<\delta<\min(s,\tau-D)$ and $\tau>D$. If $AW\ne W$, put $j=\dim W$ and $h_0=j-\dim(W\cap AW)>0$. Projecting wedge bases modulo the intersection gives the degree inequality for intersection and sum. Since $2p(j)-p(j-h_0)-p(j+h_0)=2\tau h_0^2$, it gives
\begin{equation}\label{eq:p6c1-degree-intersection}
 b(W\cap AW)+b(W+AW)\ge2b(W)-D+2\tau h_0^2>2s .
\end{equation}
Thus $W$ is stable. Its radical is stable and isotropic. A drop in dimension from $W$ to its radical decreases $p$ by at least $\tau$, and \eqref{eq:p6c1-degree-bounds} would give $b(\operatorname{rad}_gW)>s$. Thus a violation is detected on a stable isotropic subspace. The total degree condition holds by \eqref{eq:p6c1-degree-bounds}. A regular elliptic norm has no nonzero stable isotropic subspace, so the cut leaves its whole orbital integral.
\leanref{PositivityTheory.\allowbreak{}Rg.\allowbreak{}BcB.\allowbreak{}D1.\allowbreak{}onlyIsotropic}{Development.\allowbreak{}WeilBound.\allowbreak{}RepresentationTheoryGL.\allowbreak{}ArchimedeanBaseChange.\allowbreak{}ArchimedeanBaseChangeProof.\allowbreak{}IsotropicMahlerBounds.\allowbreak{}LogAbsoluteWedges.\allowbreak{}OnlyIsotropic}{60}
\leanref{PositivityTheory.\allowbreak{}Rg.\allowbreak{}BcB.\allowbreak{}D1.\allowbreak{}ellTerm}{Development.\allowbreak{}WeilBound.\allowbreak{}RepresentationTheoryGL.\allowbreak{}ArchimedeanBaseChange.\allowbreak{}ArchimedeanBaseChangeProof.\allowbreak{}IsotropicMahlerBounds.\allowbreak{}LogAbsoluteWedges.\allowbreak{}OnlyIsotropic}{61}

For a regular nonelliptic norm, the stable isotropic subspaces $W_e$ correspond to signed selections $e$ of the paired irreducible norm factors, with a zero coordinate for an unselected pair. Order the selections by inclusion of signed coordinates. The meet keeps common signs. The signed join keeps the union except at conflicting coordinates, where it keeps neither sign. Consequently
\[
 W_{e\wedge e'}=W_e\cap W_{e'},\qquad
 W_{e\vee e'}=\operatorname{rad}_g(W_e+W_{e'}).
\]
Indeed, the pairing is nondegenerate between the two factors of a pair and pairs distinct pairs trivially. Put $U=W_e+W_{e'}$ and $y_e=p(\dim W_e)-\deg_x W_e$. The degree inequality for $W_e,W_{e'}$ bounds the sum of the degrees of the intersection and $U$ below by the sum of the original degrees. Concavity gives $p(\dim(W_e\cap W_{e'}))+p(\dim U)\le p(\dim W_e)+p(\dim W_{e'})$. If the radical has smaller dimension than $U$, then $p(\dim U)-p(\dim\operatorname{rad}_gU)\ge\tau$, which absorbs the error $D$ in replacing $\deg_xU$ by $\deg_x\operatorname{rad}_gU$. In equal dimension the radical equals $U$. Hence $y_0=0$ and $y_{e\wedge e'}+y_{e\vee e'}\le y_e+y_{e'}$.

Let $T$ be the adelic twisted centralizer and $\Gamma$ its rational subgroup. Write $a(t)$ for the logarithmic coordinates of $t\in T$ on the paired factors. Translation by $t$ changes $\deg_xW_e$ by $e\cdot a(t)$. The dot product is additive across signed meet and join, including at conflicting coordinates. For $b(e)=e\cdot a(t)-y_e$ the preceding inequality gives
\begin{equation}\label{eq:p6c1-signed-meet}
 b(e\wedge e')+b(e\vee e')\ge b(e)+b(e').
\end{equation}

For every real vector $y$ indexed by signed selections, define $\gamma(y,z)$ as the sum of $(-1)^{|C|}$ over chains $C$ of nonzero selections with $e\cdot z>y_e$ for each $e\in C$. Include the empty chain. Define
\[
 V(y)=\int_{\Gamma\backslash T}\gamma(y,a(t))\,dt.
\]
For the actual cut bounds, $\gamma$ is the indicator of $e\cdot z\le y_e$ for all $e$. If there is a violation, take a maximizer $e_*$ of $b$ with $b(e_*)>0$. By \eqref{eq:p6c1-signed-meet}, every violator $e$ has $b(e\wedge e_*)\ge b(e)>0$. Toggle $e_*$ on a chain contained in its down-set. On any other chain take the first member $e$ outside that down-set and toggle $e\wedge e_*$. Earlier members lie below the meet and later ones above it, so this is a sign-reversing involution. When all inequalities hold, the empty chain is the sole contribution.

For arbitrary $y$ with $\max_e|y_e|\le M$, the integrand is bounded by the number of chains and vanishes unless $|z_i|\le2M$ for every coordinate. For example, if $z_i>2M$, drop a minus sign at $i$ on all violating selections that have one. The resulting selections are nonzero violators, since their dot products increase by more than $2M$. This order-preserving map lies below the identity. Next add a plus sign at $i$ to any selection still zero there. This map lies above the identity and also preserves violations. The two order homotopies leave the alternating chain sum unchanged and give the subposet of violators with plus at $i$. It has minimum the selection consisting of that sign alone, so toggling the minimum cancels the sum. For $z_i<-2M$ interchange the signs. The inverse image of the box under $a$ is compact modulo $\Gamma$. On paired factors the coordinates bound the id\`ele norms and the norm-one id\`ele class groups are compact fibres. The self-paired factors have compact norm-one torus quotients. Thus $V(y)$ is absolutely defined and bounded for bounded $y$. Changing one coordinate $y_e$ changes the chain sum by terms from chains containing $e$, so the difference depends on that change and the coordinates comparable with $e$.

With no paired factors the weight is constant, and the unipotent integral vanishes at a complex place with a root off the unit circle. It uses the flag from the two opposite collections of eigenlines off the circle. With paired factors, use two infinite places. Such places exist because an imaginary quadratic field ramifies over $\Q$ at a finite prime. The local degrees give the cut weight as $V(p_1(x_1)+p_2(x_2)+z_0)$, with $p_i$ the local vectors and $z_0$ the fixed remaining terms. Subtract $V(p_1(x_1))$ by replacing the coordinates one at a time. In the difference at $e$, every relevant subspace is contained in $W_e$ or contains $W_e$ and lies in $W_e^{\perp_g}$. The unipotent radical for $W_e\subseteq W_e^{\perp_g}$ acts trivially on their determinant lines. Thus the difference is invariant along the unipotent integration at the first place and its integral vanishes. The remaining weight $V(p_1(x_1))$ is constant at the second place, and the same vanishing applies there. Regular orbit maps are proper modulo the centralizer on the test support. On compact representatives the local degree vectors are bounded, so the box estimate for $V$ gives absolute integrability in these decompositions. This proves \eqref{eq:p6c1-cut-trace} in each factor.
\leanref{PositivityTheory.\allowbreak{}Rg.\allowbreak{}BcB.\allowbreak{}CV.\allowbreak{}pf\_\allowbreak{}ClassVanish}{Development.\allowbreak{}WeilBound.\allowbreak{}RepresentationTheoryGL.\allowbreak{}ArchimedeanBaseChange.\allowbreak{}ArchimedeanBaseChangeProof.\allowbreak{}QuotientVanishing.\allowbreak{}LocalQuotientIntegrals.\allowbreak{}PClassVanish}{43--136}
\leanref{PositivityTheory.\allowbreak{}Rg.\allowbreak{}BcB.\allowbreak{}CV.\allowbreak{}ellFormula2J}{Development.\allowbreak{}WeilBound.\allowbreak{}RepresentationTheoryGL.\allowbreak{}ArchimedeanBaseChange.\allowbreak{}ArchimedeanBaseChangeProof.\allowbreak{}QuotientVanishing.\allowbreak{}LocalQuotientIntegrals.\allowbreak{}EllFormula2J}{12}

Let $u_i(x)$ and $v_i(y)$ be the periodized rows and columns of the end kernels, with $M$ absorbed in the columns. Smoothing both variables by mass-one approximate identities gives $L^2$ convergence of the rows and columns by continuity and compact support. Boundedness of $V$ and continuity of the kernel give $K_h(x,\theta x)=\sum_i\langle u_i(x),Vv_i(\theta x)\rangle$. For compact $\Omega$, put $n_\Omega(x)=\#\{\gamma:x^{-1}\gamma x\in\Omega\}$ in the product rational group. Unfolding a squared row or column norm pairs translates whose difference is conjugate into a product of supports, with one inverted. Choose $\Omega$ containing those products, pulled back by $\theta$ for the columns. Each overlap integral is uniformly bounded by the fixed kernel norms, so Cauchy--Schwarz gives
\begin{equation}\label{eq:p6c1-diagonal-count}
 |K_h(x,\theta x)|\le C_1 n_\Omega(x)\|V\|.
\end{equation}
On a cell $a(m)\Omega_0$, each element counted by $n_\Omega(x)$ gives a distinct rational translate of $x$ in $a(m)\Omega_0\Omega$. Unfolding bounds the integral of the count on the quotient cell by $\operatorname{vol}(\Omega_0\Omega)$. By \eqref{eq:p6c1-reduction}, for $L=2(a+b)>0$ the integral of $H^{-2L}n_\Omega$ is bounded by a constant times the product of the sums of $(1+\|m\|_\infty)^{-L}$ over $\Z^a$ and $\Z^b$. Both converge, with one term in rank zero. When $a+b=0$ the count is constant. Hence
\begin{equation}\label{eq:p6c1-integrated-count}
 \int_{\mathcal C_\tau}n_\Omega(x)\,dx\ll(1+\tau)^{2L}.
\end{equation}
Apply \eqref{eq:p6c1-cut-trace}, \eqref{eq:p6c1-diagonal-count} and \eqref{eq:p6c1-integrated-count} with $\tau=C_0(1+r_0+kr_1)$.
\end{proof}

\subsection*{Proof of \ref{B6.04l}}\label{prf:B6.04l}\phantomsection
\begin{proof}
Take the exterior algebra of the complexified Hermitian matrices tensored with the two weight coefficient factors. Write $\Gamma$ for exterior parity and $S_s$ for multiplication by $s^j$ in degree $j$. Let $T$ be the involution on the coefficient factors. The matrix kernel $k$ acts on right compact-equivariant sections. The scalar trace is $\phi(y)=\operatorname{tr}(k(y)T)$, with the exterior signs in the kernel formulation below. Compact equivariance and commutation of $T$ with the compact action give the compact conjugation invariance.

In a trace-orthonormal Hermitian basis write $\epsilon_i,\iota_i$ for exterior multiplication and contraction. They satisfy $\epsilon_i\iota_j+\iota_j\epsilon_i=\delta_{ij}$. The convolution kernel is obtained by integrating the even wave solution against a smooth even compactly supported time function and applying the central differential operators. The wave construction gives finite propagation. On a joint eigenfunction the wave equation is scalar, so the time integral acts by the integral of the time function against $\cosh(t\sqrt{c})$, where $c$ is the eigenvalue of the wave operator. The central operators multiply this scalar by their eigenvalues. Finite sums and convolutions have the corresponding sums and products of values.

For the unipotent integral, use the field of Hermitian matrices
\[
 H(x)=P_{x^{-1}W_1}+P_{x^{-1}W_2}-1,
\]
where $P$ is standard orthogonal projection. It is invariant on the left under the flag stabilizer and equivariant on the right under the compact group. The complement relations give $H(d)=-H(1)$. On sections averaged along the flag radical, the wave operator and the central differential operators commute with multiplication by $\epsilon(H(x))$. In flag coordinates the central operators on left radical-invariant sections act by their block-diagonal parts, and the field is constant along the left parabolic directions. The wave commutation then follows from uniqueness for the wave equation. Integrate the resulting kernel identity along the radical. For $B=\int_N k(ud)\,du$ it reads
\[
 \epsilon(H(1))B=B\epsilon(H(d))=-B\epsilon(H(1)).
\]
The averaging uses Haar invariance on $N$, and the integrals have compact support there. Put $C=\epsilon(H(1))$ and $D=\iota(H(1))$. Then $CD+DC=\operatorname{tr}(H(1)^2)1$, and this scalar is nonzero since $W_1\ne0$. As $T$ commutes with $C$, cyclicity of trace gives $\operatorname{tr}((CD+DC)BT)=0$. This proves the unipotent vanishing.
\leanref{PositivityTheory.\allowbreak{}RgBcC.\allowbreak{}VC.\allowbreak{}vcKer\_\allowbreak{}of}{Development.\allowbreak{}WeilBound.\allowbreak{}RepresentationTheoryGL.\allowbreak{}ArchimedeanBaseChange.\allowbreak{}ArchimedeanBaseChangeProof.\allowbreak{}ArchimedeanGelfandKernels.\allowbreak{}LeftEquivariantFunctions.\allowbreak{}GelfandFlags.\allowbreak{}UnitaryMemGLs}{37--97}

For the elliptic integral, use the positive Hermitian cone with polar coordinate $v=xx^*$. Write $F(x)$ for the action on the exterior and weight coefficients in this frame. For $v=xx^*$ and $w=yy^*$, put
\[
 K(v,w)=c_r\det(w)^{-r}F(x)k(x^{-1}d(y^*)^{-1})\Gamma T F(y)^{-1},
\]
where $c_r$ is the polar Haar constant. The expression descends to the two cone variables by compact equivariance. The cochain identity for the multiplier kernel, in these coordinates, is
\[
 \sum_i\epsilon_i\partial_{1,i}K+\sum_i(\partial_{2,i}K)\epsilon_i=0,\qquad
 \Gamma K=K\Gamma .
\]
The identity comes from commutation of the central wave multiplier with the exterior differential. In the polar change of variables the frame derivative in direction $xX$, for $X=X^*$, is the cone derivative in direction $2xXx^*$. In the second variable the density contributes $-r\operatorname{tr}(w^{-1}dw)\det(w)^{-r}$, which is the density term in the second cochain derivative.
\leanref{PositivityTheory.\allowbreak{}RgBcC.\allowbreak{}LfSt.\allowbreak{}kChain\_\allowbreak{}holds}{Development.\allowbreak{}WeilBound.\allowbreak{}RepresentationTheoryGL.\allowbreak{}ArchimedeanBaseChange.\allowbreak{}ArchimedeanBaseChangeProof.\allowbreak{}ArchimedeanGelfandKernels.\allowbreak{}LeftEquivariantFunctions.\allowbreak{}GLMatrixDerivatives.\allowbreak{}KChainHolds}{7--8}

For $0\le s\le1$ integrate $\operatorname{tr}(\Gamma S_s K(sv+(1-s)1,v))$ over the cone. Its derivative in $s$ is the integral of the divergence with $i$th component
\[
 \operatorname{tr}\bigl(\Gamma\epsilon_i S_s\iota(v-1)K(sv+(1-s)1,v)\bigr),
 \qquad
 \iota(v-1)=\sum_i(v-1)_i\iota_i .
\]
Indeed $S_s'=\sum_i\epsilon_i S_s\iota_i$, and in the first-variable derivative use $s\epsilon_i S_s=S_s\epsilon_i$ and $\epsilon_i\iota(v-1)=(v-1)_i-\iota(v-1)\epsilon_i$. In the second-variable derivative, moving $\epsilon_i$ to the right under the trace contributes the parity sign. The remaining terms cancel by the cochain identity. The interpolated kernels have common compact support in the cone when the $z_i^2$ are distinct. In polar coordinates, boundedness of $x^{-1}d(y^*)^{-1}$ and its inverse, with $xx^*=s yy^*+(1-s)1$, bounds $y$ and $y^{-1}$ uniformly for $0\le s\le1$. Hence the divergence integrates to zero. At $s=1$ polar integration gives the twisted integral. At $s=0$ the projection $S_0$ onto degree zero gives the coefficient eigenfunctions $x\mapsto F(x)^{-1}S_0u$ with parameter $\inf(l)$. The convolution thus acts there by $\operatorname{val}(\inf(l))$. The trace of $F(d)T$ on degree zero uses the flip of the two coefficient factors. The identity $\operatorname{tr}((A\otimes B)\mathrm{flip})=\operatorname{tr}(AB)$ gives $\overline{\operatorname{ch}_l(z_1^2,\ldots,z_r^2)}$. Rank zero uses the scalar kernel on the trivial group. Fubini gives the two-block version.
\leanref{PositivityTheory.\allowbreak{}RgBcC.\allowbreak{}LfSt.\allowbreak{}lAsm}{Development.\allowbreak{}WeilBound.\allowbreak{}RepresentationTheoryGL.\allowbreak{}ArchimedeanBaseChange.\allowbreak{}ArchimedeanBaseChangeProof.\allowbreak{}ArchimedeanGelfandKernels.\allowbreak{}LeftEquivariantFunctions.\allowbreak{}GLMatrixDerivatives.\allowbreak{}TwMemGLs}{61--122}
\end{proof}

\subsection*{Proof of \ref{B6.04m}}\label{prf:B6.04m}\phantomsection
\begin{proof}
Fix one elliptic norm polynomial $P$ of degree $r$ and a tuple of local classes contributing to its stable integral. Let $L_i$ be its irreducible field factors and $K_i$ their fixed fields for conjugate inversion. The elliptic condition makes $L_i/K_i$ CM. Relative to one rational twisted representative, the form description in \ref{B6.04e} identifies rational classes with $\prod_i K_i^\times/N_{L_i/K_i} L_i^\times$. The same description at each finite place gives the local norm classes, with the split factors automatic. The local tuple has the unit class relative to the representative almost everywhere by the finite \'etale argument in \ref{B6.04b}. For each $i$, the finite norm symbols therefore prescribe one product of signs at the real places of $K_i$. By quadratic norm reciprocity, every choice of real signs with that product is realized. The Hasse norm theorem identifies two global coefficients with the same signs at all places \cite{neukirch-1999}. Thus the number of rational classes in the fibre is $\prod_i 2^{[K_i:\Q]-1}=2^{r|\Sigma_\infty|-r(P)}$.

Each such class has mass $2^{r(P)}\kappa(P)$. Its weight character and its orbital integrals at the prescribed finite classes are fixed. Apply this count to both blocks and then sum over the local tuples. The regular finite place ensures that the contributing norm polynomials are coprime, since their product becomes separable after extending scalars to that completion. 

For the second equality, decompose the adelic group into the product of the infinite and finite groups. Its Haar measure is a nonzero constant times the product of the chosen measures, by uniqueness of Haar measure. For each elliptic class the archimedean centralizer is compact. Use probability measure there in forming the quotient, so the archimedean orbital integrals are the group integrals of \ref{B6.04l}. On the finite group, disintegrate first by the adelic centralizer and then by its maximal compact subgroup. Each double coset in $\Gamma\backslash C/C_c$ has volume $1/\#(\Gamma\cap C_c)$, since $C$ is abelian. The resulting volume is the centralizer mass. At a split pair the two local group integrals combine by convolution on one factor after applying the twist to the second. At an inert place there is one local orbital integral. The product of local integrals is finite after removing the unit factors. This gives the rational-class summand times $c_{a,b}\prod_v\nu_v$. Twist translation preserves the centralizers, so the constant is common to the classes.
\end{proof}

\subsection*{Proof of \ref{B6.04da}}\label{prf:B6.04da}\phantomsection
\begin{proof}
At \(w\in T'\) the block level is the product \(L\) of the two Iwahori subgroups. The matrix indices there are the assignments of the chamber coordinates to the two blocks. The fixed kernel is diagonal in these indices. Each admissible diagonal entry is a weighted sum of characteristic functions of \(L\)-double cosets, and an inadmissible entry is zero. At \(w\in S_0\) the indices are coset representatives for the split test level. Each entry of the fixed matrix is bi-invariant under the chosen block level \(L\). Indeed, the level embeds in the Levi part of the split test level and the determinant twist is trivial on it. Left and right translation by \(L\) in the constant-term integral therefore leave the entry unchanged.

In both cases the matrix of every variable finite kernel at this position is \(P_{ij}=\delta_{ij}e_L\). The same holds for its adjoint, since \(e_L^*=e_L\), with \(f^*(x)=\overline{f(x^{-1})}\). If \(M\) is the matrix of the fixed kernel, Haar averaging gives, entry by entry,
\[
 (M*P)_{ij}=M_{ij}*e_L=M_{ij}
             =e_L*M_{ij}=(P*M)_{ij}.
\]
This calculation permits arbitrary bi-\(L\)-invariant entries of \(M\).
\leanref{PositivityTheory.\allowbreak{}Rg.\allowbreak{}BcB.\allowbreak{}F.\allowbreak{}LocComm.\allowbreak{}beta\_\allowbreak{}m}{Development.\allowbreak{}WeilBound.\allowbreak{}RepresentationTheoryGL.\allowbreak{}ArchimedeanBaseChange.\allowbreak{}ArchimedeanBaseChangeProof.\allowbreak{}SatakeExponents.\allowbreak{}UniformizerValuations.\allowbreak{}AdelicHeckeBlocks.\allowbreak{}PLocComm}{337--348}

At a variable split place the block level is
\(K=\GL_a(\mathcal O_w)\times\GL_b(\mathcal O_w)\). The transferred test is bi-\(K\)-invariant by the compact average in \eqref{eq:p6c1-split-match} and triviality of the twist on integral units. Its adjoint is again bi-\(K\)-invariant. Transposition in both factors preserves every \(K\)-double coset: elementary divisors give a diagonal representative in each factor. Transposition preserves the Haar measure normalized by \(\operatorname{vol}(K)=1\) and reverses convolution. It fixes each bi-\(K\)-invariant function, including the convolution of two such functions. Thus the two convolutions in opposite orders agree. At this place the finite factor of \(\beta_s\) is the level idempotent. At all remaining finite positions the variable factors are level idempotents. Products over distinct places commute by factorization of convolution.
\leanref{PositivityTheory.\allowbreak{}Rg.\allowbreak{}BcB.\allowbreak{}F.\allowbreak{}sphComm}{Development.\allowbreak{}WeilBound.\allowbreak{}RepresentationTheoryGL.\allowbreak{}ArchimedeanBaseChange.\allowbreak{}ArchimedeanBaseChangeProof.\allowbreak{}SatakeExponents.\allowbreak{}UniformizerValuations.\allowbreak{}AdelicHeckeBlocks.\allowbreak{}IsCompactMaximalPairPlace}{61--71}
\leanref{PositivityTheory.\allowbreak{}Rg.\allowbreak{}BcB.\allowbreak{}F.\allowbreak{}locComm}{Development.\allowbreak{}WeilBound.\allowbreak{}RepresentationTheoryGL.\allowbreak{}ArchimedeanBaseChange.\allowbreak{}ArchimedeanBaseChangeProof.\allowbreak{}SatakeExponents.\allowbreak{}UniformizerValuations.\allowbreak{}AdelicHeckeBlocks.\allowbreak{}PLocComm}{355--376}

At infinity, composing two wave multipliers convolves their even time functions and composes their central differential operators. The time convolutions commute, as do the central operators. The adjoint of a wave multiplier is again a sum of wave multipliers: inversion and adjunction take the central word to its formal adjoint, and the undifferentiated real even wave kernel is self-adjoint. These identities on compactly supported smooth sections identify the convolution kernels by testing against approximate identities. They give commutation in either order with every variable multiplier and its adjoint.
\leanref{GlobalTheory.\allowbreak{}CompanionTheory.\allowbreak{}WaveMult.\allowbreak{}asmMult\_\allowbreak{}main}{Development.\allowbreak{}AutomorphicSpectrum.\allowbreak{}Proofs.\allowbreak{}BlockDecomposition.\allowbreak{}AsmMultMain}{25--75}

On \(V\) the variable finite operators act by their spherical scalars on each constituent summand. The adjoints act by the conjugate scalars. The operator \(u\) preserves each summand. We have thus proved in every coordinate that two generators commute and that each generator commutes with the adjoint of every generator. Taking adjoints gives the identities for two adjoints. Expanding products gives commutativity of \(\mathcal A\). The same expansion, using the identities with the fixed operator in each coordinate, proves the two assertions about \(u\) and \(\beta_s\).
\leanref{PositivityTheory.\allowbreak{}Rg.\allowbreak{}BcB.\allowbreak{}F.\allowbreak{}SetupAlgPf.\allowbreak{}algCR}{Development.\allowbreak{}WeilBound.\allowbreak{}RepresentationTheoryGL.\allowbreak{}ArchimedeanBaseChange.\allowbreak{}ArchimedeanBaseChangeProof.\allowbreak{}SatakeExponents.\allowbreak{}UniformizerValuations.\allowbreak{}SetupAlgPf}{95--103}
\end{proof}

\subsection*{Proof of \ref{B6.04db}}\label{prf:B6.04db}\phantomsection
\begin{proof}
Take a composition series of \(M\), which has finite length by \cite[\S2]{bernstein-zelevinsky-1977}. In the first member containing the given vector, its image in the successive quotient \(\tau\) is nonzero and satisfies the same equations. Let \(N\) be the upper unitriangular group. The image \(y\) of this vector in the Jacquet module \(\tau_N\) is nonzero by the Iwahori projection lemma \cite[Lem.~4.7]{borel-1976}. Here the subgroup is the Iwahori for \(B\), the representation is smooth admissible, and the coefficients are complex.

The right cosets in \(I_jt_xI_j\) have representatives \(ut_x\), with \(u\) running through \((N\cap I_j)/t_x(N\cap I_j)t_x^{-1}\). There are \(\delta_B(t_x)^{-1}\) such cosets, by taking the index in each upper triangular entry. In the coinvariants all these translates have the same image. Hence the action of \(A_x\) there is the action of \(t_x\) in the normalized Jacquet module \(r_B\tau\). The vector \(y\) is fixed by the integral diagonal torus. Every integer tuple is a difference of two strictly decreasing tuples, so its torus character is
\[
 \chi(\operatorname{diag}(u_i\varpi^{e_i}))=\prod_i\xi_i^{e_i},
 \qquad u_i\in O_L^\times .
\]
This character also occurs as a quotient of \(r_B\tau\). Indeed the integral-torus invariants form a finite-dimensional module for the commuting uniformizer translations. Its simultaneous generalized \(\chi\)-space is nonzero and has a simple quotient of character \(\chi\). Projection onto the invariants by compact averaging and then onto this generalized eigenspace gives the quotient on \(r_B\tau\).

Normalized Frobenius reciprocity embeds \(\tau\) in \(\operatorname{Ind}_B^{\GL_j(L)}\chi\). Write \(\chi_i=|\cdot|_L^{s_i}\chi_i^u\) with \(\chi_i^u\) unitary. Since \(|\varpi|_L<1\), the displayed order gives \(s_1\le\cdots\le s_j\). The spherical Langlands classification, in the subrepresentation convention obtained by taking contragredients of the standard quotient, gives a unique irreducible subrepresentation of this induction. It is spherical with Satake multiset \(\{\xi_i\}\). This applies also to equal real exponents, for which the corresponding block is induced from unitary unramified characters of a general linear group \cite[Ch.~XI, \S2]{borel-wallach-2000}.

Compact averaging lifts a nonzero \(K_j\)-fixed vector of \(\tau\) first to the member of the composition series and hence to \(M\). By \(\GL_j(L)=PK_j\), evaluation at the identity identifies
\[
 M^{K_j}\simeq\bigotimes_i\sigma_i^{K_{j_i}}.
\]
Thus each factor is spherical. In normalized induction the spherical Hecke character on this line is obtained by concatenating the Satake parameters of the factors: the normalized constant term sends a symmetric Laurent polynomial to the same polynomial evaluated on the block variables. Its image in \(\tau^{K_j}\) has the same character. The elementary symmetric polynomials determine the asserted multiset.
\end{proof}

\subsection*{Proof of \ref{B6.04u}}\label{prf:B6.04u}\phantomsection
\begin{proof}
Apply the adelic counting theorem of M\"uller \cite[Theorem~1.4, (1.17)--(1.18)]{mueller-2024} to the restriction of scalars of \(\GL_d/E\). This group satisfies the analytic hypothesis of that theorem, as in its Corollary~1.3. Its adelic quotient is taken modulo the positive real split center. Fix one of the compact types. On this type a bounded infinitesimal parameter bounds the absolute value of the Casimir eigenvalue, since the Casimir acts by a fixed polynomial in the parameter. Both the cuspidal and the residual counting functions are finite at this bound. Every representation in question contributes at least one dimension. Sum over the finite type list. There are finitely many compositions of integers at most \(n\), so taking products proves the assertion for tuples. The finite lists used for the discrete leaves in the trace comparison are covered by this statement.
\end{proof}

\subsection*{Proof of \ref{B6.04dc}}\label{prf:B6.04dc}\phantomsection
\begin{proof}
We first compute the integral of a single double-coset indicator. For the general linear group use the involution \(\theta(g)=J_r(\bar g^{-1})^tJ_r^{-1}\) of \ref{B6.04b}. On a diagonal it sends \(d_i\) to \(\bar d_{i^*}^{-1}\), with reversal within the block.

Here is the diagonalisation used in the computation. If \(z\) is strictly decreasing, every element of \(I_rt_zI_r\) is \(I_r\)-conjugate to a diagonal of valuations \(z\). Indeed, an initial conjugation puts the element in the form \(t_z u\), \(u\in I_r\). Divide by \(\varpi^{z_{r-1}}\). For \(r>1\) the reduction is diagonal with a simple unit in the last position and zeros elsewhere. Hensel factorisation gives a unit eigenvalue and an invariant complementary summand. Write the unit eigenvector as \((v_0,\ldots,v_{r-2},1)\). The row equations give \(v_i\in\varpi^{z_i-z_{r-1}}\mathcal O_L\). In the quotient by this line the matrix, before the division, has entries
\[
 \varpi^{z_i}\bigl(u_{ij}-\varpi^{z_{r-1}-z_i}v_i u_{r-1,j}\bigr)
 \qquad(i,j<r-1).
\]
The matrix in parentheses is upper triangular modulo \(\varpi\) with unit diagonal, since \(u_{r-1,j}\in\varpi\mathcal O_L\). Induction diagonalises the quotient. Lifting its basis to the invariant complement preserves the reduction of that basis. This gives the required Iwahori change of basis. Rank one starts the induction.

For twisted conjugation take \(\delta\in I_rt_zI_r\). Its norm belongs to \(I_rt_{z-z^*}I_r\). To see the double-coset product here, factor the middle Iwahori element in upper unipotent, diagonal and lower unipotent factors, using Gaussian elimination with unit pivots. Conjugation by the first dominant translation moves the upper factor into the Iwahori, and inverse conjugation by the second moves the lower factor into the Iwahori. The diagonal factor commutes with both translations. The entries of \(z-z^*\) are strictly decreasing. Diagonalise the norm by the preceding argument, and use the same change of basis for twisted conjugation of \(\delta\). If the resulting norm is \(D\), then \(D\delta=\delta\theta(D)\). The ordered valuations of \(D\) and \(\theta(D)\) agree and are distinct. The entry equations therefore make \(\delta\) diagonal. Its valuations are \(z\): reduction with the standard lattice chain distinguishes the diagonal Iwahori double cosets, since the relative positions of that chain and its translate record the translation of each one-dimensional successive quotient.

In the unitary group the ordinary diagonalisation can be made isometric. In the ordered integral eigenbasis, invariance of the Hermitian form pairs an eigenline only with the line whose eigenvalue is its conjugate inverse. The pairings at the reverse positions are units by unimodularity. Rescale one vector in each pair to give the prescribed Witt pairing. The norm on units of an unramified quadratic extension is onto, so one may also normalise the middle vector in odd rank. These rescalings preserve the Iwahori flag.

Fix monic \(Q_1,Q_2\) as in \eqref{eq:p6c1-orbital-match}. On the unitary support their product has roots \(\lambda_i\in L^\times\) with \(v(\lambda_i)=h_i\) and \(\lambda_{i^*}=\bar\lambda_i^{-1}\). Thus the roots in \(Q_1\) determine exactly one admissible \(A\). Conversely, any such list gives a diagonal unitary representative. For a general linear block with exponent vector \(z\), a diagonal twisted representative must solve
\[
 d_i/\bar d_{i^*}=\lambda_i,\qquad v(d_i)=z_i.
\]
Choose \(d_{i^*}\) of the prescribed valuation and solve for \(d_i\) on each pair. The valuation condition holds because \(v(\lambda_i)=z_i-z_{i^*}\). At a middle position Hilbert's theorem 90 solves \(d/\bar d=\lambda\), and multiplication by a power of \(\varpi\) sets the valuation. This constructs both blocks. If the prescribed roots and valuations do not occur, neither side meets the test support.

There is one contributing orbit on each side when these conditions hold. For the unitary group this follows from the ordered isometric eigenbases. In a twisted block any two diagonal solutions of the same valuations differ by units \(e_i\) with \(e_i=\bar e_{i^*}\). A diagonal twisted change of basis by units accounts for both entries of each pair. At the middle position it accounts for \(e_i\) by surjectivity of the norm on units. In particular two conjugates in the support are conjugate there by the Iwahori. For a representative in the support the set of conjugators contributing to the orbital integral is consequently \(CI\), where \(C\) is its centralizer, twisted on the general linear side. In diagonal coordinates \(C\cap I\) is its maximal compact subgroup: commutation with the regular diagonal norm first makes the centralizer diagonal, and the remaining condition pairs its entries by conjugate inversion. Thus the orbital integral of the indicator is \(\operatorname{vol}(I)/\operatorname{vol}(C\cap I)=\operatorname{vol}(I)\). The normalized indicator has integral one. This also gives one for the rank-zero factor.

On \(\ker Q_1\) the Hermitian form is an orthogonal sum of the selected hyperbolic planes and, if selected, the integral middle line. Its determinant has valuation zero in this basis. The sign in the unitary orbital sum is therefore \(+1\). Across blocks all root valuations differ, so
\[
 v(\operatorname{Res}(Q_1,Q_2))
 =\sum_{i\in A,j\in B}v(\lambda_i-\lambda_j)=v_A(x).
\]
The factor \(\Delta_w\) of \ref{B6.04b} is precisely \((-\sqrt q)^{-v_A(x)}\). This proves the identity, including its vanishing cases.
\leanref{PositivityTheory.\allowbreak{}Rg.\allowbreak{}Bc.\allowbreak{}ArchbcB.\allowbreak{}B.\allowbreak{}M4.\allowbreak{}asmL\_\allowbreak{}A}{Development.\allowbreak{}WeilBound.\allowbreak{}RepresentationTheoryGL.\allowbreak{}ArchimedeanBaseChange.\allowbreak{}ArchimedeanBaseChangeProof.\allowbreak{}IwahoriFactorisation.\allowbreak{}TwistedConjugationNorms.\allowbreak{}Development}{802--868}
\leanref{PositivityTheory.\allowbreak{}Rg.\allowbreak{}Bc.\allowbreak{}ArchbcB.\allowbreak{}B.\allowbreak{}M4.\allowbreak{}asmR\_\allowbreak{}A}{Development.\allowbreak{}WeilBound.\allowbreak{}RepresentationTheoryGL.\allowbreak{}ArchimedeanBaseChange.\allowbreak{}ArchimedeanBaseChangeProof.\allowbreak{}IwahoriFactorisation.\allowbreak{}TwistedConjugationNorms.\allowbreak{}Development}{997--1030}

For the powers, two translations in the same chamber have additive length. Explicitly, the right cosets in \(ItI\) are indexed by the upper unipotent part \(U_0\) of \(I\) modulo \(tU_0t^{-1}\). For translations \(t,t'\) in that chamber the inclusions
\(tt'U_0(tt')^{-1}\subseteq tU_0t^{-1}\subseteq U_0\)
give a product of transversals, the second transversal inside \(tU_0t^{-1}\) obtained by conjugating the one for \(t'\). The triangular factorisation used above gives this coset description also for the unitary Iwahori: uniqueness of the triangular factors in Witt coordinates makes them fixed under the unitary involution. Thus the convolution of the two indicators divided by Iwahori volume on each factor is the similarly normalized indicator for \(tt'\). For fixed \(A\), restriction of exponents is additive, and \(\min(h_i,h_j)\) always selects the same index on this chamber. Hence \(v_A(x+x')=v_A(x)+v_A(x')\). Expand the power over ordered tuples of translations and apply the single-translation identity to their sum. The assignment \(A\) is made once for that sum. This gives the asserted order of summation and convolution.
\end{proof}

\subsection*{Proof of \ref{B6.04d}}\label{prf:B6.04d}\phantomsection
\begin{proof}
Fix the exceptional places $S_0\cup J$ and the archimedean data of \ref{B6.04da} from $\pi$ and $S$. Choose the level $K_w$ of the final induced representation to be $\GL_n(\mathcal O_w)$ away from these places and their conjugates. At each exceptional place choose a principal congruence subgroup deep enough that the diagonal Levi blocks of its intersection with each standard block parabolic belong to the prescribed pair level and that the two determinant twists are trivial there. Include the conjugate pair level at a split place. Such a depth exists by openness and the finite number of blocks. The two archimedean weights of \ref{B6.04b} fix the compact types. These choices precede $T$.

For the separating tests, now fix $T$ and adjoin an auxiliary good inert place to obtain $T'$. Use the pair of Iwahori subgroups at every place of $T'$, the prescribed levels at $S_0\cup J$, and maximal compact subgroups elsewhere. At these levels there are finitely many possible discrete inducing data for either block. Indeed the ranks partition the block rank and the infinitesimal multisets partition the prescribed block multiset at every embedding. For a vector in a parabolic induction fixed by the test level $U$, evaluation on representatives of $P\backslash G/U$ gives fixed vectors for open subgroups of the Levi. There are finitely many representatives by Iwasawa decomposition and compactness. Restriction of the prescribed compact types to the compact Levi likewise gives a finite list. For each resulting rank, level and compact type, the specified infinitesimal characters give a finite-codimensional ideal in the enveloping-algebra center. At finite level the adelic quotient has finitely many arithmetic components by ideal class finiteness. Harish-Chandra's finiteness theorem on each component gives finite dimension for these data, with the split central character fixed \cite[Thm.~8.4]{borel-2007}. This bounds the number of discrete representations that occur. Their local inductions have finite length and finite-dimensional fixed spaces. We use this list for the chosen $T'$ throughout the interpolation and the choice of the power $m$. Changing the test coefficients changes none of its defining levels or types.
\leanref{UpperTheory.\allowbreak{}HeckeLevelFrameExistence.\allowbreak{}CF.\allowbreak{}cls\_\allowbreak{}fin}{Development.\allowbreak{}AutomorphicArithmetic.\allowbreak{}Foundations.\allowbreak{}BaseChangeIdentities.\allowbreak{}FrameMatching.\allowbreak{}HeckeEigenvalueFiniteness.\allowbreak{}ClsFin}{31--46}

At each inert testing place use the Weyl representative $(\beta_1,\ldots,\beta_e)$ of the spherical parameter with $|\beta_1|\ge\cdots\ge|\beta_e|\ge1$, where $e=\lfloor n/2\rfloor$. Permutations and inversions in the relative Weyl group give this order. The spherical representation embeds in the unramified principal series in subrepresentation order by the spherical Langlands classification \cite[Ch.~XI, \S2]{borel-wallach-2000}. Frobenius reciprocity therefore gives the ordered character as a quotient of its normalized Jacquet module. Its generalized character space for the integral-torus invariants is finite-dimensional and contains an eigenline. The Iwahori--Jacquet isomorphism carries this line back to a common translation eigenline, with the double-coset index accounted for as in Lemma~\ref{B6.04db} \cite[Lem.~4.7]{borel-1976}. Use for $\alpha$ its character on the comparison exponent lattice. Its ordered coordinates are $(\beta_1,\ldots,\beta_e,\beta_e^{-1},\ldots,\beta_1^{-1})$, with $1$ inserted in the middle for odd $n$.
\leanref{PositivityTheory.\allowbreak{}Rg.\allowbreak{}BcB.\allowbreak{}F.\allowbreak{}Eig.\allowbreak{}eigInert}{Development.\allowbreak{}WeilBound.\allowbreak{}RepresentationTheoryGL.\allowbreak{}ArchimedeanBaseChange.\allowbreak{}ArchimedeanBaseChangeProof.\allowbreak{}SatakeExponents.\allowbreak{}UniformizerValuations.\allowbreak{}InertEigenvalues.\allowbreak{}EigInert}{8--8}
\leanref{PositivityTheory.\allowbreak{}Rg.\allowbreak{}BcB.\allowbreak{}F.\allowbreak{}FX.\allowbreak{}coneEigOf}{Development.\allowbreak{}WeilBound.\allowbreak{}RepresentationTheoryGL.\allowbreak{}ArchimedeanBaseChange.\allowbreak{}ArchimedeanBaseChangeProof.\allowbreak{}SatakeExponents.\allowbreak{}UniformizerValuations.\allowbreak{}StrictJoinCandidates.\allowbreak{}IsHeckeFsum}{241--273} For each block and each allowed assignment of chamber coordinates to the two factors, include the transfer multiplier when pulling the candidate translation characters back to the common exponent lattice. These form a finite set of characters. Regard them as points of a complex torus. For each point $\gamma\ne\alpha$, choose a coordinate $i(\gamma)$ on which they differ. The Laurent polynomial
\[
 L(Z)=\prod_{\gamma\ne\alpha}
 \frac{Z_{i(\gamma)}-\gamma_{i(\gamma)}}{\alpha_{i(\gamma)}-\gamma_{i(\gamma)}}
\]
is one at $\alpha$ and zero at the other points. Multiply by $Z^v/\alpha^v$ with $v$ deep enough in the positive chamber that the support has consecutive gaps at least two. Its coefficients give the local test. They fix the target translation eigenvector. A nonzero value on a transferred candidate for any allowed assignment forces equality with $\alpha$ for that character.

At each place of $S_0$ take the target local representation at the chosen level and list the irreducible subquotients of the finite list of candidate inductions. The inductions have finite length. For a compact open level, the nonzero fixed spaces of the irreducibles are simple modules of its Hecke algebra. Isomorphic simple modules give the same irreducible representation generated by the fixed vectors. The density theorem applied to the finite list of distinct simple modules gives an element acting by a rank-one projection on the target module and zero on every other module in the list. Normalize the test by the level volume for the chosen Haar measure. If all subquotients of a candidate are different from the target, exactness of fixed vectors gives a fixed-space composition filtration on which this operator strictly lowers the filtration. Its power of degree the fixed-space dimension is zero.

Choose a bound $m_0$ for the dimensions of candidate Iwahori fixed spaces at all the inert testing places, including the auxiliary place. For each $w\in S_0$ choose a bound $m_w$ for the dimensions of the split fixed spaces. Take $m=m_0+\sup_{w\in S_0}m_w+1$, with empty supremum zero. The tensor of the local target vectors gives a nonzero $v$ in the $\pi$ summand with $uv=v$. Each variable generator acts on this summand by a scalar, and its adjoint acts by the conjugate scalar. Their action on $\C v$ defines a unital $*$-character $p$ of $\mathcal A$. The reference generator of Lemma~\ref{B6.04da} satisfies $p(a_0)=1$.
\leanref{UpperTheory.\allowbreak{}HeckeLevelFrameExistence.\allowbreak{}frameEx\_\allowbreak{}of}{Development.\allowbreak{}AutomorphicArithmetic.\allowbreak{}Foundations.\allowbreak{}BaseChangeIdentities.\allowbreak{}FrameMatching.\allowbreak{}HeckeEigenvalueFiniteness.\allowbreak{}KUEqKLOfNotMemBad}{116--230}

At a split spherical place where $\mu_w$ is unramified, the normalized constant-term map followed by the determinant twist preserves convolution. At the other split places transfer the complete test by \eqref{eq:p6c1-split-match}, with the factor $\mu_w((-1)^{ab})$ evaluated once. At spherical inert places the unit matches by \ref{B6.04e}.

At a bad inert place let $K(j)$ be the principal congruence subgroup of level $j\ge1$ in $K_w$. Put $L=E_w$ and let $q_0$ be the base residue cardinality. Take a base-field element $t$ of valuation $j$. Write the two norm generators as $1+tg_i$ with polynomials $Q_i$ of degrees $a,b$. Their resultants satisfy
\begin{equation}\label{eq:p6c1-congruence-resultant}
 v(\operatorname{Res}(P_{g_1},P_{g_2}))=v(\operatorname{Res}(Q_1,Q_2))-jab.
\end{equation}
The relevant lattice conditions are stability under $g_i$ and under $(g_1,g_2)$. For nonintegral generators the counts vanish. For integral generators choose $x\in O$ with $x-\bar x$ a unit and put
\[
 Y_i=\frac{x\bar x\,g_i}{x-\bar x+xtg_i}.
\]
The norm-one relation gives $\sigma_i(g_i)=-g_i/(1+tg_i)$ and hence $\sigma_i(Y_i)=Y_i$. The denominator is a unit in the integral order. The inverse transformation has unit denominator $x\bar x-xtY_i$, so the orders coincide on the blocks and the product. The difference of the images of roots $u',v'$ is
\[
 (u'-v')\,\frac{x\bar x(x-\bar x)}
 {(x-\bar x+xtu')(x-\bar x+xtv')}.
\]
The second factor is a unit. Thus the lattice counts and the cross-resultant valuation are preserved. Apply \eqref{eq:p6c1-lie-count} and \eqref{eq:p6c1-congruence-resultant}. With $q_w=q_0^2$, the transfer of $1_{K(j)}$ is $(-\sqrt{q_w})^{-jab}$ times the indicator of integral pairs with twisted norms congruent to one modulo $\varpi^j$, for a uniformizer $\varpi$. The level projection and its transfer carry the index $[K_w:K(j)]$.

Apply Lemma~\ref{B6.04dc} at each inert translation place, including the auxiliary place. Its ordered positive coordinates $l_i$ have opposites $-l_i$ and a middle zero in odd rank. Assign the pairs indexed by $A$ to block $a$ and the rest to block $b$. The allowed assignments have $|A|=\lfloor a/2\rfloor$, $|A^c|=\lfloor b/2\rfloor$ and $(a\bmod2)+(b\bmod2)=n\bmod2$. Other assignments contribute zero. The coefficient of each allowed term is multiplied by
\[
 [K_a\times K_b:I_a\times I_b](-\sqrt{q_w})^{-v_{\rm Res}},
\]
where $K_a,K_b$ are the integral subgroups, $I_a,I_b$ the Iwahori subgroups and $v_{\rm Res}$ the cross-resultant valuation. Each difference of two roots of distinct valuations has valuation the smaller one. Across two opposite pairs the four differences thus contribute $-2\max(l_i,l_j)$, and a middle coordinate against a pair contributes $-l_i$. On the chamber this is linear in the assigned coordinates, which gives the character multiplier used in the interpolation. The distinct valuations give separability of the product norm polynomial on the auxiliary test support. At this place the fixed kernel is diagonal in the assignments. Its $m$th convolution power uses products of translations in the same chamber, so the exponents add after normalization by Iwahori volume. A variable finite letter and its adjoint both have the diagonal level idempotent at this position, as in Lemma~\ref{B6.04da}.
\leanref{GlobalTheory.\allowbreak{}CongMatch.\allowbreak{}match\_\allowbreak{}principal\_\allowbreak{}congruence\_\allowbreak{}inert}{Development.\allowbreak{}AbelianVarietyTheory.\allowbreak{}Proofs.\allowbreak{}DivisorProofs.\allowbreak{}CongruenceSubgroupIndicatorWeight}{41--66}
\leanref{PositivityTheory.\allowbreak{}Rg.\allowbreak{}BcB.\allowbreak{}F.\allowbreak{}B3.\allowbreak{}locMatch}{Development.\allowbreak{}WeilBound.\allowbreak{}RepresentationTheoryGL.\allowbreak{}ArchimedeanBaseChange.\allowbreak{}ArchimedeanBaseChangeProof.\allowbreak{}SatakeExponents.\allowbreak{}UniformizerValuations.\allowbreak{}WordFrameBounds.\allowbreak{}PLocMatch}{143--208}

To change twists take $u_r=J_r$ in odd rank and an imaginary scalar times $J_r$ in even rank. Then $\theta_J=\operatorname{Ad}(u_r)\theta_0$ and $u_r\theta_0(u_r)=1$. The map $\delta\mapsto\delta u_r$ preserves the norm and identifies the twisted classes. Evaluate the finite tests at $xu_r^{-1}$ in each block. At a split pair the two evaluations give the convolution in the split orbital integral, with the twist on the second component.

At infinity apply \ref{B6.04l} for the two weights, with fixed smooth wave multipliers whose product of values at the target parameters is one. The central operators of the word multiply the target value, and adjoints give the conjugated values. The resulting kernels have the very cuspidal vanishing of \ref{B6.04c}. The auxiliary place gives the separability required by \ref{B6.04m} for both norm polynomials together. The rational-to-stable factor there is $2^{n|\Sigma_\infty|}$, and the adelic factor is $c_{a,b}$. Thus the signed term of \ref{B6.04b} contributes the block coefficient $s_I/(c_{a,b}2^{n|\Sigma_\infty|})$.
\leanref{PositivityTheory.\allowbreak{}Rg.\allowbreak{}BcB.\allowbreak{}F.\allowbreak{}F0.\allowbreak{}chainEx\_\allowbreak{}of}{Development.\allowbreak{}WeilBound.\allowbreak{}RepresentationTheoryGL.\allowbreak{}ArchimedeanBaseChange.\allowbreak{}ArchimedeanBaseChangeProof.\allowbreak{}SatakeExponents.\allowbreak{}UniformizerValuations.\allowbreak{}HyperbolicEntryBounds.\allowbreak{}InfinitesimalWeightZero}{42--88}
\leanref{UpperTheory.\allowbreak{}BaseChangeIdentities.\allowbreak{}FrameWordIdentity.\allowbreak{}traceword\_\allowbreak{}pf}{Development.\allowbreak{}AutomorphicArithmetic.\allowbreak{}Foundations.\allowbreak{}BaseChangeIdentities.\allowbreak{}FrameMatching.\allowbreak{}HeckeWordBounds.\allowbreak{}PTraceword}{19--60}
\leanref{UpperTheory.\allowbreak{}BaseChangeIdentities.\allowbreak{}FrameWordIdentity.\allowbreak{}wordId}{Development.\allowbreak{}AutomorphicArithmetic.\allowbreak{}Foundations.\allowbreak{}BaseChangeIdentities.\allowbreak{}FrameMatching.\allowbreak{}HeckeWordBounds.\allowbreak{}FrameBound}{30--40}

Write $g=\sum_j c_j w_j$, where $w_j$ is a word in the variable generators and their adjoints, with the empty word allowed. On the unitary side $u$ commutes with every letter by Lemma~\ref{B6.04da}. In block $s$ prefix each $w_j$ by $m$ copies of $\beta_s$. The resulting operator is $D_sR_s(w_j)$. Thus for indices $j_1,\ldots,j_k$ the exceptional finite positions contain $mk$ forward copies of the fixed kernel. Apply \eqref{eq:p6c1-two-block-trace} with the local matches above to this ordered product. Multiply by $c_{j_1}\cdots c_{j_k}$ and sum. Distributivity gives the middle operator $(D_sR_s(g))^k$ between the two fixed end convolutions $T_{1,s},T_{2,s}$.
\leanref{PositivityTheory.\allowbreak{}Rg.\allowbreak{}BcB.\allowbreak{}F.\allowbreak{}B3.\allowbreak{}bw}{Development.\allowbreak{}WeilBound.\allowbreak{}RepresentationTheoryGL.\allowbreak{}ArchimedeanBaseChange.\allowbreak{}ArchimedeanBaseChangeProof.\allowbreak{}SatakeExponents.\allowbreak{}UniformizerValuations.\allowbreak{}WordFrameBounds.\allowbreak{}WordIdAt}{62--71}
The concatenations have length at most $k\ell_0$, where $\ell_0$ bounds the original word lengths including the repetitions of the fixed test. Matrix multiplication, applied to the entries and inverse entries, bounds the logarithmic support by $r_0+kr_1$. The summed kernel retains the common level. It also satisfies the auxiliary regularity and archimedean vanishing conditions. Apply \eqref{eq:p6c1-growth} and absorb the fixed block coefficient into $C_s$.
\end{proof}

\subsection*{Proof of \ref{B6.04aa}}\label{prf:B6.04aa}\phantomsection
\begin{proof}
At an embedding $\tau:F\hookrightarrow\C$, the infinitesimal multiset of $\pi$ is
\[
 H(\lambda_\tau)=\bigl\{\lambda_{\tau,k}+\tfrac{n-1}{2}-k:0\le k<n\bigr\}\subset\tfrac{n-1}{2}+\Z.
\]
Its entries are distinct because $\lambda_{\tau,k}$ is weakly decreasing. By Lemma~\ref{B6.04a}, in unnormalized parabolic induction the $i$-th component embeds in the induction of $\rho_i|\det|^{(a_i-1)((d_i-1)/2-j)}$ ($0\le j<d_i$); subtracting the half-modulus exponent $a_i((d_i-1)/2-j)$ on block $j$, it embeds in the normalized parabolic induction of $\rho_i|\det|^{j-(d_i-1)/2}$ ($0\le j<d_i$), so its infinitesimal multiset at $\tau$ is the union of the shifts $m(\rho_i,\tau)+j-(d_i-1)/2$, equivalently $\bigcup_{j=0}^{d_i-1}(m(\rho_i,\tau)+(d_i-1)/2-j)$. Since this union is a submultiset of $H(\lambda_\tau)$, the entries of $m(\rho_i,\tau)$ are distinct and lie in
\[
 \tfrac{n-1}{2}-\tfrac{d_i-1}{2}+\Z=\tfrac{n-d_i}{2}+\Z.
\]
Let $\eta=\eta_{F/F^+}$ be the quadratic Hecke character of $F/F^+$. Since $F/F^+$ is CM, $\eta_v(-1)=-1$ at every real place $v$ of $F^+$, while $\eta\circ N_{F/F^+}=1$. Thus $\chi=\eta^{a_i}$ satisfies $\chi_v(-1)=(-1)^{a_i}$ and $\chi\circ N_{F/F^+}=1$. Let $\varepsilon_\ell:G_F\to\bar{\Q}_\ell^\times$ be the $\ell$-adic cyclotomic character, which is the realization of $|\cdot|_F$ in Proposition~\ref{B6.05c}(9); it sends inverse arithmetic Frobenius at $w\nmid\ell$ to $q_w^{-1}$ and has filtration jump $-1$ at every embedding $\tau:F\hookrightarrow\bar{\Q}_\ell$.

First suppose $n-a_i-d_i+1$ is even. Then $(n-d_i)/2\equiv(a_i-1)/2\pmod1$, so the distinct weights of $\rho_i$ lie in $(a_i-1)/2+\Z$. Since $\rho_i$ is conjugate self-dual, $(\rho_i,\chi)$ is a polarized input for Theorem~\ref{B9.02}. Let $r_i:G_F\to\GL_{a_i}(\bar{\Q}_\ell)$ be the resulting representation, and for $0\le j<d_i$ set
\[
 m_j=\frac{n-a_i-d_i+1}{2}+j\in\Z,\qquad r_{i,j}=r_i\otimes\varepsilon_\ell^{-m_j}.
\]
Because $(1-a_i)/2-m_j=(d_i-1)/2-j+(1-n)/2$, twisting the local comparisons of Definitions~\ref{B9.05}--\ref{B9.06} for $r_i$ by $\varepsilon_\ell^{-m_j}$ yields assertions (1) and (2): at $w\nmid\ell$ the Weil--Deligne parameter is twisted by $|\cdot|_w^{-m_j}$, multiplying Frobenius eigenvalues by $q_w^{m_j}$ to give $q_w^{(a_i-1)/2+m_j}B_w(\rho_i)=q_w^{(n-d_i)/2+j}B_w(\rho_i)$, and at $w\mid\ell$ each filtration jump $(a_i-1)/2-z$ increases by $m_j$ to $(n-d_i)/2+j-z$.

Now suppose $n-a_i-d_i+1$ is odd. Then $(n-d_i)/2\equiv a_i/2\pmod1$, so the weights of $\rho_i$ lie in $a_i/2+\Z$. Fix a CM type $\Sigma$ of $F/F^+$, and write $\epsilon(\tau_0)=1$ if $\tau_0\in\Sigma$ and $\epsilon(\tau_0)=-1$ if $\tau_0\notin\Sigma$. Given any finite set $S$ of finite places of $F$, Lemma~\ref{B9.36c} with $e_v=1$ at every complex place provides a unitary Hecke character $\mu$ of $F$ unramified on $S$, restricting to $\eta$ on $\mathbf A_{F^+}^\times$, with $\mu_v(z)=z/|z|$ relative to $\tau_0\in\Sigma$. Since $\eta\circ N_{F/F^+}=1$, we have $\mu(x)\mu(c(x))=\eta(N_{F/F^+}(x))=1$, so $\mu^c=\mu^{-1}$.

Put $\tilde\rho_i=\rho_i\otimes(\mu\circ\det)$. Then $\tilde\rho_i$ is cuspidal on $\GL_{a_i}(\mathbf A_F)$ and satisfies
\[
 \tilde\rho_i^c\simeq\rho_i^\vee\otimes(\mu^{-1}\circ\det)=\tilde\rho_i^\vee=(\tilde\rho_i\otimes(\chi\circ N_{F/F^+})^{-1})^\vee.
\]
At each embedding $\tau_0:F\hookrightarrow\C$, the infinitesimal weights of $\tilde\rho_i$ are $z+\tfrac12\epsilon(\tau_0)$ for $z\in m(\rho_i,\tau_0)$; they are distinct and lie in $(a_i-1)/2+\Z$. Thus $(\tilde\rho_i,\chi)$ is a polarized input for Theorem~\ref{B9.02}. Let $\tilde r_i:G_F\to\GL_{a_i}(\bar{\Q}_\ell)$ be the attached representation.

Consider the algebraic Hecke character $\psi_0=\mu|\cdot|_F^{1/2}$ of $F$. At $\tau_0\in\Sigma$, $\psi_{0,\tau_0}(z)=(z/|z|)(z\bar z)^{1/2}=z$, so $\psi_0$ has integral infinity type $n_{\tau_0}=1$ for $\tau_0\in\Sigma$ and $n_{\tau_0}=0$ for $\tau_0\notin\Sigma$. By Proposition~\ref{B6.05c}(9), $\psi_0$ admits a continuous $\ell$-adic realization $\psi_\mu:G_F\to\bar{\Q}_\ell^\times$ which is de Rham and potentially crystalline at every $w\mid\ell$ (crystalline when $\mu_w$ is unramified), has filtration jump $-n_{\iota\circ\tau}$ at $\tau:F\hookrightarrow\bar{\Q}_\ell$, and has Weil--Deligne character $(\mu_w|\cdot|_w^{1/2})\circ\operatorname{rec}_w$ with zero monodromy. For $0\le j<d_i$ set
\[
 m_j=\frac{n-a_i-d_i}{2}+j\in\Z,\qquad r_{i,j}=\tilde r_i\otimes\psi_\mu^{-1}\otimes\varepsilon_\ell^{-m_j}.
\]
Because $n-a_i-d_i$ is even, $m_j$ is an integer, and $-a_i/2-m_j=(d_i-1)/2-j+(1-n)/2$. For every finite place $w\nmid\ell$ at which $\mu_w$ is unramified (in particular every $w\in S$ with $w\nmid\ell$), the character twist $\mu_w|\cdot|_w^{1/2+m_j}$ is unramified and hence trivial on inertia $I_w$ and on Weil commutators; thus zero monodromy and the toral and full tame tests of Definition~\ref{B9.05} are preserved between $\tilde\rho_{i,w}\otimes|\det|_w^{(1-a_i)/2}$ and $(\rho_{i,w}|\det|_w^{(d_i-1)/2-j})\otimes|\det|_w^{(1-n)/2}$, and untwisting $\tilde r_i$ by $\psi_\mu^{-1}\otimes\varepsilon_\ell^{-m_j}$ twists the Weil--Deligne parameter by $\mu_w^{-1}|\cdot|_w^{-1/2-m_j}$, matching
\[
 \bigl(\tilde\rho_{i,w}\otimes|\det|_w^{(1-a_i)/2}\bigr)\otimes(\mu_w^{-1}|\det|_w^{-1/2-m_j})
 =\bigl(\rho_{i,w}|\det|_w^{(d_i-1)/2-j}\bigr)\otimes|\det|_w^{(1-n)/2}.
\]
When $\rho_{i,w}$ is also spherical at such a place $w\nmid\ell$ where $\mu_w$ is unramified, inverse arithmetic Frobenius acts on $r_{i,j}$ with eigenvalues $q_w^{(a_i-1)/2}\mu_w(\varpi_w)B_w(\rho_i)\cdot\mu_w(\varpi_w)^{-1}q_w^{1/2}\cdot q_w^{m_j}=q_w^{(n-d_i)/2+j}B_w(\rho_i)$. At every $w\mid\ell$ in $S$, $\mu_w$ is unramified, so $\tilde\rho_{i,w}$ is spherical if and only if $\rho_{i,w}$ is spherical, the character $\psi_\mu|_{G_{F_w}}$ is crystalline (preserving crystallinity, potential semistability, zero monodromy, and the period comparison of Definition~\ref{B9.06}), and the filtration jumps of $r_{i,j}$ at $\tau:F\hookrightarrow\bar{\Q}_\ell$ are
\[
 \left(\frac{a_i-1}{2}-z-\frac{\epsilon(\iota\circ\tau)}{2}\right)+n_{\iota\circ\tau}+m_j
 =\frac{n-d_i}{2}+j-z,
\]
since $n_{\iota\circ\tau}-\epsilon(\iota\circ\tau)/2=1/2$ for both $\iota\circ\tau\in\Sigma$ and $\iota\circ\tau\notin\Sigma$.

For any two choices $\mu,\mu'$ corresponding to finite sets $S,S'$, the resulting continuous semisimple representations $r_{i,j}$ and $r'_{i,j}$ have identical Frobenius characteristic polynomials at every finite place $w\nmid\ell$ outside the finite ramification locus of $\rho_i,\mu,\mu'$. By Chebotarev density \cite[Ch.~VII, \S13]{neukirch-1999} and Brauer--Nesbitt semisimplicity, $r_{i,j}\simeq r'_{i,j}$. Choosing $S$ to contain any given finite place $w$ together with all places above $\ell$ therefore establishes (1) and (2) at every finite place $w$ for one fixed representation $r_{i,j}$.

Finally, summing over all $(i,j)$ gives $r_\pi$ of dimension $\sum_i a_id_i=n$. At any finite place $w$ where $\pi$ is unramified, the normalized parabolic induction $\Pi_w=\bigtimes_i\bigtimes_{j=0}^{d_i-1}\rho_{i,w}|\det|_w^{j-(d_i-1)/2}$ has the spherical subquotient $P_w$ of Lemma~\ref{B6.04a}. Exactness of $\GL_n(\mathcal O_w)$-invariants gives $\Pi_w^{\GL_n(\mathcal O_w)}\ne0$, and evaluation at the identity via the Iwasawa decomposition $\GL_n(F_w)=P(F_w)\GL_n(\mathcal O_w)$ identifies $\Pi_w^{\GL_n(\mathcal O_w)}$ with $\bigotimes_{i,j}\rho_{i,w}^{\GL_{a_i}(\mathcal O_w)}$. Hence every $\rho_{i,w}$ is spherical, and $B_w(P_w)=\bigsqcup_i\bigsqcup_{j=0}^{d_i-1}q_w^{(d_i-1)/2-j}B_w(\rho_i)=\bigsqcup_i\bigsqcup_{j=0}^{d_i-1}q_w^{j-(d_i-1)/2}B_w(\rho_i)$. By (2), $r_\pi$ is crystalline at $w$ when $w\mid\ell$; by (1), when $w\nmid\ell$ it is unramified at $w$ with inverse Frobenius eigenvalue multiset $\bigsqcup_{i,j}q_w^{(n-d_i)/2+j}B_w(\rho_i)=q_w^{(n-1)/2}B_w(P_w)$. At each $\tau:F\hookrightarrow\bar{\Q}_\ell$, the union of the multisets $(n-1)/2-(z+(d_i-1)/2-j)$ equals $(n-1)/2-H(\lambda_{\iota\circ\tau})=\{k-\lambda_{\iota\circ\tau,k}:0\le k<n\}$.
\end{proof}

\subsection*{Proof of \ref{B6.04g}}\label{prf:B6.04g}\phantomsection
\begin{proof}
Fix \(s\). For \(a\in\mathcal A\), Lemma~\ref{B6.04da} gives \(D_sR_s(a)=R_s(a)D_s\). Apply the same equality to \(a^*\) and take adjoints. This gives \(D_s^*R_s(a)=R_s(a)D_s^*\). Hence \(B_s=D_s^*D_s\) commutes with \(R_s(a)\). Since \(B_s\) is positive and \(R_s(\mathcal A)\) is closed under adjoints, the algebraic unital \(*\)-algebra they generate is commutative. Its operator-norm closure is \(\mathcal C_s\). A contractive character on the algebraic algebra extends uniquely to this closure.
\leanref{PositivityTheory.\allowbreak{}Rg.\allowbreak{}BcB.\allowbreak{}F.\allowbreak{}ExponentialIntegralMeasures.\allowbreak{}adjoin\_\allowbreak{}comm}{Development.\allowbreak{}WeilBound.\allowbreak{}RepresentationTheoryGL.\allowbreak{}ArchimedeanBaseChange.\allowbreak{}ArchimedeanBaseChangeProof.\allowbreak{}SatakeExponents.\allowbreak{}UniformizerValuations.\allowbreak{}PeakTraceBoundHolds}{187--215}

Suppose that for each \(s\) every character \(\eta\) with \(\eta\circ R_s=p\) vanishes on \(B_s\). Put \(K_0=1+\sum_s\|R_s(a_0)\|\) and \(\delta=|p(a_0)|/(2K_0)>0\). In the character space of \(\mathcal C_s\) the set \(Y_s=\{\eta:\eta(B_s)\ge\delta^2\}\) is compact. For each \(\eta\in Y_s\) choose \(a\) with \(\eta(R_s(a))\ne p(a)\). A scalar multiple \(d\) of \(a-p(a)1\) satisfies \(|\eta(R_s(d))|>1\), an open condition. Take a finite subcover for each \(s\) and collect the resulting \(d\)'s in a finite list. Put \(h=\sum_d d^*d\). Then \(p(h)=0\), each \(R_s(h)\) is positive, and \(\eta(R_s(h))>1\) on \(Y_s\).

Choose \(M>1+\max_s\|R_s(h)\|\) and set \(b_L=(1-h/M)^L\). We have \(p(b_L)=1\). Apply the Gelfand norm formula to \(R_s(b_L)^*B_sR_s(b_L)\) in the concrete commutative \(C^*\)-algebra just constructed \cite[Thm.~2.1.10]{murphy-1990}. A zero block contributes norm zero. On \(Y_s\) we have \(0\le1-\eta(R_s(h))/M\le1-1/M\). Off \(Y_s\) the same quantity lies in \([0,1]\) and \(\eta(B_s)<\delta^2\). Thus
\begin{equation}\label{eq:p6c1-peak-norm}
 \|D_sR_s(b_L)\|^2
 \le\max\{\delta^2,\|B_s\|(1-1/M)^{2L}\}.
\end{equation}
There are finitely many blocks, so one \(L\) gives \(\|D_sR_s(b_L)\|\le\delta\) for all \(s\). Put \(g=a_0b_L/p(a_0)\). Moving \(R_s(a_0)\) past \(D_s\) gives \(\|D_sR_s(g)\|\le1/2\). On the unitary side \(A=u\rho(g)\) satisfies \(Av=v\). Applying \eqref{eq:p6c1-trace-power} to this \(g\) and summing the bounds gives
\begin{equation}\label{eq:p6c1-peak-trace}
 |\operatorname{tr}(A^k)|\le C(1+k)^N2^{-k}\qquad(k\ge1).
\end{equation}
Let \(\xi_1,\ldots,\xi_r\) be the distinct eigenvalues of \(A\), with algebraic multiplicities \(m_i\). Triangularizing \(A\) gives
\(\operatorname{tr}(A^{k+j})=\sum_i\xi_i^j m_i\xi_i^k\) for \(0\le j<r\). All \(r\) traces tend to zero. The matrix \((\xi_i^j)_{j,i}\) is an invertible Vandermonde matrix, so every \(m_i\xi_i^k\) tends to zero. The eigenvalue \(1\) has positive multiplicity since \(v\ne0\). This contradiction proves the assertion.
\end{proof}

\subsection*{Proof of \ref{B6.04fa}}\label{prf:B6.04fa}\phantomsection
\begin{proof}
Write \(Y\) for the rational quotient and \(\Gamma\) for the rational group. For a compact set \(\Omega\) in the adelic group put
\(n_\Omega(x)=\#\{\gamma\in\Gamma:x^{-1}\gamma x\in\Omega\}\).
Write \(k_i(x)\) for the \(L^2(Y)\) row, including the finite matrix indices, which represents evaluation of the \(i\)-th component of \(A_0h\) at \(x\). Periodizing the compact reference kernel gives these rows. Unfolding their squared norms pairs two translates of its support \(C\). An overlap requires \(x^{-1}\gamma x\in CC^{-1}\); each overlap integral is bounded by the squared \(L^2\) norm of an entry by Cauchy--Schwarz. Summing over the finite matrix indices gives \(\|k_i(x)\|^2\le c\,n_{CC^{-1}}(x)\). Translation continuity of the kernel in \(L^2\) gives continuous rows.

For the two ranks \(a,b\) take \(M=2(a+b)>0\). On a reduction cell with representatives in \(t\Omega_0\), the integral of \(n_\Omega\) over the image of the cell is at most \(\operatorname{vol}(\Omega_0\Omega)\). Indeed, each element counted at a representative \(x\) gives a distinct translate in \(t\Omega_0\Omega\), and the sum of its indicator unfolds. By \eqref{eq:p6c1-reduction}, \(H^{-2M}\) on the cell with indices \(l\in\Z^a,l'\in\Z^b\) is bounded by a constant times \((1+\|l\|_\infty)^{-M}(1+\|l'\|_\infty)^{-M}\). Both sums converge, with a single term for a zero rank. Consequently
\[
 \sum_i\int_Y H(x)^{-2M}\|k_i(x)\|^2\,dx<\infty .
\]
\leanref{PositivityTheory.\allowbreak{}Rg.\allowbreak{}BcB.\allowbreak{}E1.\allowbreak{}vecNorm}{Development.\allowbreak{}WeilBound.\allowbreak{}RepresentationTheoryGL.\allowbreak{}ArchimedeanBaseChange.\allowbreak{}ArchimedeanBaseChangeProof.\allowbreak{}ThetaCellCovers.\allowbreak{}KernelRowsColumnsSquareIntegrable}{51}
\leanref{PositivityTheory.\allowbreak{}Rg.\allowbreak{}BcB.\allowbreak{}E1.\allowbreak{}wtCount}{Development.\allowbreak{}WeilBound.\allowbreak{}RepresentationTheoryGL.\allowbreak{}ArchimedeanBaseChange.\allowbreak{}ArchimedeanBaseChangeProof.\allowbreak{}ThetaCellCovers.\allowbreak{}KernelDiagonalIntegrableHolds}{23}

Let \(U\) be a neighborhood of \(\chi\) in the character space on which \(|\eta(A_0)|>|\chi(A_0)|/2\). The representation of \(\mathcal C_s\) is faithful by construction. A continuous function supported in \(U\) and equal to one at \(\chi\) therefore has nonzero operator image by the Gelfand representation theorem. Choose a unit vector in its range. The cyclic subspace generated by that vector has the form \(J L^2(\nu)\), where \(J\) is an isometry and \(\nu\) is a probability measure carried by \(U\). Indeed, the vector state on continuous functions gives \(\nu\), and sending a continuous function to its action on the vector gives the isometry by the definition of the state. Under \(J\) every \(z\in\mathcal C_s\) acts by multiplication by \(\eta(z)\).

Use the convention that the inner product is linear in its second argument. The rows \(J^*k_i(x)\) give measurable functions \(F_i(x,\eta)=\overline{(J^*k_i(x))(\eta)}\), with
\[
 (A_0Jh)_i(x)=\int F_i(x,\eta)h(\eta)\,d\nu(\eta).
\]
The weighted row bound and Fubini give \(H^{-M}F(\,\cdot\,,\eta)\in L^2(Y)\) almost everywhere. These functions are nonzero on a set of positive \(\nu\)-measure, since \(A_0J\) is bounded below by \(|\chi(A_0)|/2\). For the identity, each variable generator and \(X\), denote the corresponding operator by \(z\). The product \(zA_0\) has a continuous compactly supported matrix convolution kernel, also when \(z\) is scalar. Apply
\((zA_0)A_0Jh=A_0J(\eta(zA_0)h)\) to a countable dense set of \(h\in L^2(\nu)\). Testing on compact subsets of \(Y\) and using Fubini gives
\((zA_0)F=\eta(zA_0)F\) almost everywhere in both variables. The list of equations is countable. Take one \(\eta\) where they all hold and the weighted vector is nonzero. The equation with \(z=1\) gives the continuous representative \(f=A_0F/\eta(A_0)\): convolution of a locally integrable function with the continuous compact reference kernel is continuous. The other equations now hold everywhere by continuity. Convolution associativity and \(\eta(A_0)\ne0\) give \(zf=\eta(z)f\). Taking \(U\) successively within distance \(1/j\) of the prescribed values on the first \(j\) generators and on \(X\) gives the sequence.
\leanref{PositivityTheory.\allowbreak{}Rg.\allowbreak{}BcB.\allowbreak{}E2.\allowbreak{}oneStep\_\allowbreak{}holds}{Development.\allowbreak{}WeilBound.\allowbreak{}RepresentationTheoryGL.\allowbreak{}ArchimedeanBaseChange.\allowbreak{}ArchimedeanBaseChangeProof.\allowbreak{}EigenvalueSeparation.\allowbreak{}AlmostEverywhereEigen.\allowbreak{}OrthogonalKeepingHolds}{35}
\end{proof}

\subsection*{Proof of \ref{B6.04f}}\label{prf:B6.04f}\phantomsection
\begin{proof}
Work in the Hilbert space \(H_s\) of \ref{B6.04da}. Let \(A_0=R_s(a_0)\) and \(X=(\beta_s^m)^*\beta_s^m\). By \ref{B6.04g}, \(\chi(A_0)=1\) and \(\chi(X)>0\). Apply Langlands' spectral decomposition to \(\operatorname{Res}_{E/\Q}(\GL_a\times\GL_b)\), omitting any rank-zero factor. This is a connected reductive group over \(\Q\); its rational adelic quotient with Haar measure is the one defining \(H_s\). Tensor the decomposition with the fixed finite matrix multiplicity. Its fibres are normalized inductions from discrete representations on the norm-one adelic Levi, extended trivially on the real split centre and then twisted by \(\lambda\in i\mathfrak a_M^*\). The split centre of the ambient group is included in this parameter since the quotient defining \(H_s\) retains it. Right convolutions act fibrewise. These are the decomposition and equivariance assertions of \cite[Thm.~7.2]{arthur-2005}.

Enumerate the generators in \ref{B6.04da}, including the differentiated reference operators, and include \(A_0,X\) at every stage. For the first \(j\) operators, the open joint spectral neighbourhood of radius \(1/j\) about their values at \(\chi\) has nonzero spectral projection. Indeed a continuous function of these coordinates supported in the neighbourhood and taking value one at \(\chi\) is a nonzero element of the concrete commutative \(C^*\)-algebra. The projection is decomposable, so it is nonzero on some induction fibre. For \(j\ge3\), \(A_0\) is invertible on its range there. This puts the range in the finite level and compact types of the reference kernel. That space is finite-dimensional by admissibility. Simultaneous diagonalisation there gives a unit joint eigenvector \(F_j\) with eigenvalues in the stated neighbourhood. Write \(a_j\) for its \(A_0\)-eigenvalue. In particular \(a_j\to1\).

Write \(A_0=WP\) with \(P\) the finite level projection and \(W\) the infinite convolution. Then \(PF_j=F_j\) and \(WF_j=a_jF_j\). If \(t\) is a finite generator, its finite convolution has eigenvalue \(a_j^{-1}\) times the eigenvalue of \(R_s(t)\). The infinite kernel is self-adjoint at the reference operator by the real even wave function of \ref{B6.04l}. The finite convolution for \(X\) therefore has eigenvalue \(a_j^{-2m}\) times its full eigenvalue. A central differential operator applied to \(W\) gives its infinitesimal-character value times \(a_j\).

We record how the latter values bound the inducing data. In rank \(r>0\), for either complex central action, write \(D\) for the upper triangular matrix with entries \(\xi_i+(r-1)/2\) on the diagonal and \(-1\) above it, where \(\{\xi_i\}\) is the infinitesimal multiset. For \(P_\xi(z)=\det(z-D)\) the matrix determinant lemma gives
\[
 \frac{P_\xi(z-1)}{P_\xi(z)}
 =1-\sum_{k\ge0}({\bf1}^tD^k{\bf1})z^{-k-1}.
\]
Indeed \(D+{\bf1}{\bf1}^t\) is lower triangular with each diagonal entry increased by one. In the identity obtained by clearing \(P_\xi(z)\), the coefficient of \(z^{r-j-1}\) determines its \(j\)-th nonleading coefficient from the preceding ones: the two multipliers of this coefficient are \(r-j\) and \({\bf1}^t{\bf1}=r\). This works through \(j=r\). Thus convergence of the finitely many central values gives convergence of the characteristic polynomials and of the root multisets. Their limits are the shifted block weights by \ref{B6.04b} and \ref{B6.04l}.
\leanref{PositivityTheory.\allowbreak{}Rg.\allowbreak{}BcB.\allowbreak{}F.\allowbreak{}Rd.\allowbreak{}B.\allowbreak{}infLim}{Development.\allowbreak{}WeilBound.\allowbreak{}RepresentationTheoryGL.\allowbreak{}ArchimedeanBaseChange.\allowbreak{}ArchimedeanBaseChangeProof.\allowbreak{}SatakeExponents.\allowbreak{}UniformizerValuations.\allowbreak{}FrameReading.\allowbreak{}GelfandValueContinuity.\allowbreak{}InfLim}{77--116}

Pass to a subsequence with fixed Levi ranks. Each discrete factor of positive rank \(r\) has a norm twist \(|\det|_{\mathbb A_E}^{z_j}\) with \(\Re z_j=0\), by the domain of integration in the spectral decomposition. Its normalized discrete representation is trivial on its real split centre. Hence the sum of the imaginary parts of its infinitesimal parameters over the embeddings of \(E\) is zero. After twisting, the sum is \(r[E:\Q]\Im z_j\). Its summands occur in the bounded block multisets; changing from normalized to unnormalized induction adds fixed real shifts. This bounds \(z_j\), and then bounds the parameters of the normalized discrete factor.

The inducing levels and compact types belong to finite lists. At finite level \(L_0\) this follows by evaluating an induced section on the finite set \(P_M\backslash G/L_0\) at the finite places, using the compact model and maximal level almost everywhere. At infinity restrict the prescribed compact types to the compact Levi. With these levels and types fixed, there are finitely many normalized discrete factors in the bounded parameter set. The finite adelic level has finitely many arithmetic components. The nonsplit central directions modulo arithmetic units are compact after the real split centre is removed, so their bounded characters form a finite set. On the semisimple components the Casimir formula bounds the homogeneous-bundle Laplacian on the indicated compact types. The arithmetic discrete-spectrum bound of \cite[Thm.~1.5.2]{ji-1997} applies to each such bundle; passage to a torsion-free subgroup of finite index gives the same finiteness for the original component. This proves the assertion also for residual discrete factors. In rank one the central-character argument suffices.

Take a subsequence on which these discrete factors are fixed and all \(z_j\) converge. In the compact model the level-and-type space is now one fixed finite-dimensional space: norm twists are trivial on the compact subgroups. Its unit vectors \(F_j\) have a convergent subsequence with unit limit \(F\). Each finite kernel in question is a finite sum of coset operators at the chosen level. Its matrix entries in this model are continuous in the twists, since each translate introduces a fixed positive norm raised to the twist parameter. Thus all the finite eigenvalue equations pass to \(F\), with the two ratios in the statement.

Each limiting twist is purely imaginary. Every limiting parameter of the twisted factor is real, since it occurs in the limiting block multiset up to a real shift. Summing its imaginary parts gives \(r[E:\Q]\Im z=0\). Hence every limiting twist is zero. The fixed discrete factors give \(P_a,P_b\) with the required infinitesimal characters. Finally the eigenvalue of \((b_s^\dagger)^{*m}*b_s^{*m}\) on \(F\) is positive by \ref{B6.04g}, so \(b_s^{*m}F\ne0\).
\end{proof}

\subsection*{Proof of \ref{B6.05a}}\label{prf:B6.05a}\phantomsection
\begin{proof}
For (1) and (2) use the finiteness theorem and the geometric lemma for smooth representations, with both functors normalized \cite[\S\S2--3, in particular Lemma~2.12]{bernstein-zelevinsky-1977}. In the integral for the duality pairing, the two half-modulus factors give the modulus in the integration formula for \(P\backslash G_n\). Additivity along the filtration gives the multiplicity formula. The segment subrepresentation, its Jacquet module and the polynomial generators are given in \cite[\S\S3,6--7]{zelevinsky-1980}. Assertion (4) is Bernstein's unitary product theorem \cite[\S8.2]{bernstein-1984}.

For (5), take the Zelevinsky multisegments of \(\pi_1,\pi_2\). Every term of a segment is a member of the corresponding cuspidal support. If two segments from different multisegments were linked, then either they would share a term, or the last term of one would be adjacent to the first term of the other on a cuspidal line. These alternatives give respectively \(j=0\) and \(j=\pm1\) in the excluded isomorphism. Zelevinsky's product formula for multisegments with no linked pair across the two factors gives the irreducible indexed by their union \cite[Prop.~8.5]{zelevinsky-1980}. In particular the hypothesis holds in \ref{B1.15ac} and \ref{B8.43h} when the two supports occupy disjoint unramified twist classes.

The segment quotient and the generic product description are \cite[\S9, in particular Thm.~9.7]{zelevinsky-1980}. The tempered classification uses the square-integrable factors in unitary normalization \cite[\S2]{tadic-1986}. For (8), the standard ordering gives the Langlands quotient. The classification gives uniqueness and exhaustion. The constituent for the sum occurs in the product by the multisegment multiplication formula, equivalently by passing to the dual parametrization in the Grothendieck ring \cite[\S\S6--9]{zelevinsky-1980}.
\end{proof}

\subsection*{Proof of \ref{B6.05b}}\label{prf:B6.05b}\phantomsection
\begin{proof}
In (1), \(J\) is the inverse image of the unipotent radical of the indicated parabolic in the finite general linear group. To apply the depth-zero Jacquet isomorphism, test both sides against the representations of \(M(\mathcal O_K)/J_M\) induced from cuspidals of its finite Levi subgroups. These induced representations contain every irreducible of that finite group. For a test induced from a smaller Levi \(M_0\), finite-group Frobenius reciprocity identifies its multiplicity space on the left with the space for the cuspidal type on the refined parahoric. Apply the Moy--Prasad isomorphism for that type in \(G_n\) and in \(M\). Transitivity of Jacquet restriction identifies the two resulting spaces on \(M_0\). Both maps are induced by the coinvariant projection, and modulus characters are trivial on the compact Levi subgroups. This gives (1) \cite[Prop.~6.7]{moy-prasad-1996}.

For (2), choose an ordering of the support for which its exterior tensor product occurs in a Jacquet module of \(V\). Such an ordering exists by cuspidal support and the geometric lemma of \ref{B6.05a}. Tensoring the given fixed vectors gives a nonzero \(J_M\)-invariant. Fixed vectors under a compact subgroup are an exact functor over \(\C\), by averaging. Thus the Jacquet module has a nonzero \(J_M\)-invariant. Its lift under (1) is fixed by \(K_n^1\subset J\).

The converse theorem in (3) is \cite[Conj.~1.1 and Thm.~1.3]{jacquet-liu-2018}. To check the application, take a generic test representation in rank \(m\). Each supercuspidal in its support has rank smaller than \(m+1\). The cuspidal pair formula gives both local \(L\)-factors equal to one, also for the dual pair. The sign from the central character of the test is common to the two Whittaker equations. Applying the hypothesis to every twist \(\tau|\det|^z\) therefore gives equality of the gamma factors as functions of \(s\) in rank \(m\). We use the pair functional equation and multiplicativity under normalized induction. \cite[\S\S2,3,8]{jacquet-piatetskii-shapiro-shalika-1983}

For a unitary character \(\alpha\) of \(K^\times\), let \(R(\alpha)\) be the ratio of the two gamma factors with test \(\alpha\), in the field of rational functions of \(q^{-s}\). Fix another unitary character \(\xi\). The normalized products of \(m\) unitary characters are irreducible and generic by \ref{B6.05a}(4) and the Whittaker model for induction \cite[Thm.~9.7]{zelevinsky-1980}. Testing with \(m\) copies of \(\xi\), and then replacing one copy by \(\alpha\), gives \(R(\xi)^m=1\) and \(R(\alpha)R(\xi)^{m-1}=1\). Hence \(R(\alpha)\) has one common value for all unitary characters.

Write \(\omega_i\) for the central character of \(\pi_i\). For every \(c\in K^\times\) of sufficiently negative valuation there is a unitary character \(\alpha_c\) of sufficiently large conductor with \(R(\alpha_c)=\omega_2(c)/\omega_1(c)\). This is the highly ramified twist formula and the character construction in \cite[Prop.~2.6 and proof of Cor.~2.7]{jiang-nien-stevens-2015}. In conductor-zero normalization one prescribes \(\alpha_c(1+x)=\psi(cx)\) for \(v(x)\ge\lfloor -v(c)/2\rfloor+1\) and extends to the units and then unitarily to \(K^\times\). For a scaled additive character the change-of-character factor in the pair functional equation gives the same ratio with \(c\) measured against that character. Given \(z\in K^\times\), choose \(c\) so that both \(c\) and \(cz\) are in this range. Constancy of the ratio gives \(\omega_1(z)=\omega_2(z)\). The common value of \(R\) is consequently one.

Finally take any generic irreducible \(\tau\) of rank \(r\le\lfloor n/2\rfloor\). If \(r=m\), the equality is already known. Otherwise choose a unitary character \(\alpha\) with conductor greater than the conductors of all rank-one members of the cuspidal support of \(\tau\). Its unramified twists are disjoint from those members. Appending \(m-r\) copies of \(\alpha\) by normalized induction gives an irreducible generic test in rank \(m\), by \ref{B6.05a}(4),(5) and the generic induction criterion. Multiplicativity gives the desired equality for \(\tau\) after cancellation of the nonzero rank-one factors. This verifies every twist and the central-character condition.

For (4), apply Brauer induction to \(1_G\), and multiply by the virtual character on the left. The tensor identity reduces the assertion to degree-zero virtual characters of elementary subgroups, allowing twists by linear characters \cite[Ch.~10]{serre-1977}. An elementary group is a product of a cyclic group and a finite \(p\)-group. Its irreducible characters are monomial. Modulo the subgroup generated by inductions of linear-character differences, each monomial \(\operatorname{Ind}_H^E\lambda\) equals \([E:H]1_E\). Indeed one first replaces \(\lambda\) by \(1_H\). In a subnormal chain from \(H\) to \(E\) with cyclic prime-order factors, the induction of the trivial character at each step is the sum of all characters of the cyclic quotient. Each summand is congruent to the trivial character. Transitivity of induction proves the assertion along the chain. Hence each degree-zero combination vanishes in the quotient.

Assertion (5) is Theorem~\ref{B7.31}. For (6), apply Chevalley's congruence theorem to the finitely generated subgroup \(\mathcal O_E^\times\subset E^\times\). Its multiplicative congruence modulo \(a\) for an integral unit is exactly \(u-1\in a\mathcal O_E\) \cite[pp.~36--37]{chevalley-1951}.
\end{proof}

\subsection*{Proof of \ref{B6.05c}}\label{prf:B6.05c}\phantomsection
\begin{proof}
For (1)--(3) apply cuspidal rapid decrease on Siegel sets and Harish-Chandra finiteness with the given compact types and central eigensystem \cite{harish-chandra-1968}\cite[Thm.~8.4]{borel-2007}. A finite set of Siegel sets covers the quotient at each fixed finite level. A bound there by each inverse power of height gives the stated bound with the infimum over rational translates. In (3) the absolute value has already been made constant in the split central direction. The cuspidal discrete decomposition gives the asserted constituents. Assertion (4) is the isobaric uniqueness theorem and the strict local bound of Jacquet and Shalika \cite{jacquet-shalika-1981}. For (5), the transform and the contragredient have reciprocal Satake roots at every common spherical place. Apply (4) to the two singleton lists and transport the evaluation pairing. Item (6) applies Theorem~\ref{B9.02} to the polarized cuspidal representation in its statement. Item (7) is Clozel's purity lemma and the cohomology calculation for regular algebraic type \cite[Lem.~4.9 and \S3.5]{clozel-1990}.

For (8) use global and local reciprocity with the inverse arithmetic convention \cite[Ch.~VI]{neukirch-1999}. Dividing by the positive archimedean factors and the finite ray unit groups gives the indicated narrow ray group. Characters separate the elements of this finite abelian group. In (9), locally at \(\ell\), the character on an open subgroup of the units is the algebraic character with the prescribed exponents. The remaining smooth factor has finite inertia. The rank-one period comparison for locally algebraic characters gives the asserted jumps and Weil character. Over the extension killing the finite inertia factor the comparison is crystalline, and this already holds over \(E_v\) when the local smooth factor is unramified \cite[Ch.~III, Appendix]{serre-1998}\cite{fontaine-1994}. Item (10) is Berger's \(p\)-adic monodromy theorem \cite{berger-2002}.
\end{proof}

\subsection*{Proof of \ref{B6.05d}}\label{prf:B6.05d}\phantomsection
\begin{proof}
Local finiteness, local duality and the local Euler characteristic are the finite-module theorems of \cite[Ch.~I, \S2]{milne-2006}, also in the completion formulation of \ref{B8.41p}. To see (2), both modules are unramified because \(v\nmid\ell\). Their unramified classes inflate from the procyclic residue Galois group, of \(\ell\)-cohomological dimension one, so their products vanish. For either module \(B\), the unramified subgroup is the cokernel of \({\rm Frob}_v-1\) on \(B\), of order \(|H^0(F_v,B)|\). By (1) and local duality, the product of the two orders is \(|H^1_v(A)|\). This proves equality with both annihilators \cite[Thm.~I.2.6]{milne-2006}.

For (3) apply the global finiteness and Euler characteristic theorems for a finite module with \(S\) containing the places dividing its order \cite[Ch.~I, \S\S4--5]{milne-2006}. Each complex place contributes \(|A|\) to \(a_\infty\). For the normalized sequence use \ref{B6.09}, with the uniformizer Kummer class first in the cup fixing the invariant. For (4) let \(B=\bigoplus_{v\in S} H^1_v(A)\), \(D\) be the image of global localization in \(B\), and \(L=\bigoplus L_v\). By \ref{B6.09} the dual localization image is \(D^\perp\). The identity
\[
 |D\cap L|/|D^\perp\cap L^\perp|=|D|\,|L|/|B|
\]
follows by taking the annihilator of \(D+L\). Including the kernels of localization gives the ratio
\[
 \frac{h_1\,|L|}{|B|\,|\Sha^1_S(A^\dagger)|}.
\]
Exactness in degree two gives
\(h_2=|\Sha^1_S(A^\dagger)|\prod_{v\in S}|H^2_v(A)|/|H^0(F,A^\dagger)|\).
Substitution of (1) and (3), with
\(\sum_{v\mid\ell}[F_v:\Q_\ell]=[F:\Q]\), gives (4). Assertion (5) is Bloch--Kato's local dimension and duality theorem with its lattice and divisible-quotient conditions \cite[\S3]{bloch-kato-1990}.

In the strict condition of \ref{B6.29} the local submodule at \(\ell\) is zero, so its dual is all of \(H^1_\ell(A^\dagger)\) at finite coefficient level. In \ref{B6.10} the modulus \(H=f^{(\ell')}\ell^m\) contains \(\ell\), and the tested class is locally zero at \(p\mid H\). At the auxiliary \(q\nmid H\) both coefficients are unramified, so (2) applies there.
\end{proof}

\subsection*{Proof of \ref{B6.05e}}\label{prf:B6.05e}\phantomsection
\begin{proof}
The bound on divided-power products in (1) makes the thickening locally divided-power nilpotent. Apply Grothendieck's deformation theorem for abelian schemes \cite[Thm.~5.4.3]{grothendieck-1974}. The equivalence there identifies the lift with the filtered crystal. Essential surjectivity gives the prescribed summand, and full faithfulness gives the assertion for a specified homomorphism. Dualizing the cohomological filtration gives the homology formulation.

The de Rham calculation and base change in (2) are \cite[Prop.~2.5.2]{berthelot-breen-messing-1982}. Primitivity also follows by pulling back along addition: on each of the two coordinate axes the pullback is the original class, and degree one on the product is the direct sum of the two degree-one groups. The Poincaré pairing identifies the homology of the dual with the dual of the homology of \(A\), with the two Hodge filtrations annihilating one another. Pullback by the homomorphism of a line bundle gives isotropy, and a principal polarization is an isomorphism. This proves the pairing assertion and uses the duality and biduality of \cite[Ch.~I, \S1]{faltings-chai-1990}. For the comparison in (3), use smooth coordinates locally on the lift. In the divided-power envelope of the diagonal the relative coordinate differences have divided powers \(u^{[j]}\). In each coordinate the augmentation to the relative de Rham complex is a resolution: contraction in positive degree sends \(u^{[j]}du\) to \(u^{[j+1]}\), since \(d u^{[j+1]}=u^{[j]}du\). Tensor these contractions in the coordinate directions and use étale descent. A finite affine cover and its intersections compute the comparison for the separated quasi-compact lift.

For (4), de Rham comparison identifies integration with the map from integral singular homology to its complexification. Evaluation on invariant differentials gives the differential of uniformization. Functoriality gives the endomorphism assertion. The first Chern class on the torus is the integral alternating form of its Appell--Humbert bundle. The comparison for \(d\log\) multiplies this class by \(2\pi i\), with the sign determined by the order used for the alternating pairing \cite[Chs.~2--3]{birkenhake-lange-2004}. For (5), shrink to a simply connected parameter neighborhood and continue a basis of integral homology in the local system. Integration of relative invariant differentials on this basis gives the holomorphic lattice vectors. The homology maps of endomorphisms are locally constant in that basis. In the resulting flat frame, the coefficients of a relative de Rham class are its periods. Crystalline transport in characteristic zero is the Taylor isomorphism of the Gauss--Manin connection \cite[\S\S1--2]{katz-oda-1968}. A constant cycle is horizontal, so pairing commutes with each Taylor coefficient. All terms of degree greater than \(s\) evaluate to zero. For (6) apply the good reduction criterion and its potential and complex multiplication forms \cite[Thms.~1--3]{serre-tate-1968}.
\end{proof}

\subsection*{Proof of \ref{B6.05f}}\label{prf:B6.05f}\phantomsection
\begin{proof}
Items (1) and (2) are the idèlic main theorem of complex multiplication and its \(j\)-invariant specialization for the maximal order \cite[Ch.~5]{shimura-1971}. In (2) two ideal lattices are homothetic precisely when their ideal classes agree. The ray-class rationality theorem and the integral difference theorem apply to (3) by linearity in the rational locally constant function on the ray group \cite{klingen-1962,deligne-ribet-1980}. In the difference theorem the smoothing ideal is prime to the modulus, so translation preserves the set of ideals in the series. For (4) apply the partial-zeta functional equation to the class and its dual lattice class \cite{chapdelaine-2009}. For (5) multiply the local functional equations of \cite[\S\S3--4]{tate-1967}. A complex place of angular exponent \(n_\tau\) contributes \((-i)^{|n_\tau|}\). The finite-place factors are those for the same global additive character. The nonzero angular exponent excludes a pole from a pure norm character.
\end{proof}

\subsection*{Proof of \ref{B6.05}}\label{prf:B6.05}\phantomsection
\begin{proof}
Take the polarized Galois representation from \(\mathcal R\) and the definite-unitary transfer from \ref{B6.04a}. The local representation, finite-group and unit assertions of \(h_N\) are \ref{B6.05a} and \ref{B6.05b}. Proposition~\ref{B6.05c} supplies the automorphic and character assertions, using \(\mathcal R\) in item (6). The transfer proved in \ref{B6.04a} supplies the definite-unitary entry in the assembled data. The cohomology assertions are \ref{B6.05d}. The abelian-scheme and special-value assertions are \ref{B6.05e} and \ref{B6.05f}.

For the weak descent used in \ref{B6.06}, let \(E/E^+\) be CM and unramified at finite places, and let \(S\) consist of split places and contain the places above the specified rational prime. A regular algebraic conjugate self-dual cuspidal representation spherical outside \(S\) is spherical at every nonsplit finite place. Lemma~\ref{B9.36a} therefore supplies the definite-unitary match, including its unramified components at those places, in each required rank \(2,3,4\). Proposition~\ref{B6.06} now gives \ref{B6.03}.
\end{proof}

\subsection*{Proof of \ref{B6.06a}}\label{prf:B6.06a}\phantomsection
\begin{proof}
Fix the infinitesimal character of \(\pi\) at every real place. In matrix coordinates these characters specify the scalars of the central Capelli operators: at a place \(v\) their eigenvalues are the values of one polynomial in the spectral parameter. The coefficients generate the center of the enveloping algebra for \(\mathfrak{gl}_a\). The kernel \(J\) of the resulting character of the archimedean enveloping-algebra center has codimension one and annihilates \(\pi\).

Choose a compact open subgroup \(U_0\subset U\). Let \(V\) be the space of \(\xi\). For each intertwiner \(A:V\to\pi^U\) and each adelic point \(g\), take the linear functional \(v\mapsto A(v)(g)\). Their common kernel \(V_0\) is \(K_\infty\)-stable by equivariance. On \(V/V_0\) these functionals span the dual. Their matrix coefficients are continuous since \(A(v)\) is an automorphic form and \(A(\xi(k)v)(g)=A(v)(gk)\). Thus \(V/V_0\) is a continuous finite-dimensional compact-group representation. Every \(A\) factors through it. In particular there is a single finite set of continuous compact types containing all its constituents.

At level \(U_0\) there are finitely many arithmetic quotients, by finiteness of the adelic class set for \(\GL_a\) (equivalently, finiteness of the ideal class group at maximal level and finite index upon shrinking compact open level). On each quotient apply Harish-Chandra's finiteness theorem with the ideal \(J\) and the projector onto this finite set of compact types \cite[Thm.~8.4]{borel-2007}. The split central direction is taken with the central action of \(\pi\). The resulting space is finite dimensional. Its intersection with \(\pi\) contains \(A(V)\) for every \(A\).
\end{proof}

\subsection*{Proof of \ref{B6.06b}}\label{prf:B6.06b}\phantomsection
\begin{proof}
The highest weight of \(\Xi\), shifted by the half-sum of the positive roots, is regular. Thus the matching in (1) gives the regular algebraic infinitesimal character for the dual coefficients. Apply the archimedean cohomology calculation of \cite[\S3.5]{clozel-1990}. Item (2) is \cite[Thm.~3.13]{clozel-1990}, including the permutation of embeddings.

For (3), choose in each \(\pi\in S\) a nonzero cohomology class at infinity and a nonzero vector in \(\pi_f^U\). Their tensors belong to distinct direct summands of cuspidal cohomology. The cuspidal cohomology injection preserves their independence \cite[Lem.~3.15(i),(ii)]{clozel-1990}. At finite level its target is the sum of the ordinary cohomologies of the arithmetic components with coefficients in \(\Xi\). Indeed, decompose \(Y\) into \(G(E)\)-orbits and apply Shapiro's lemma to the module of functions on each orbit. This is the group cohomology on the right, including for compact open levels with torsion stabilizers since the coefficients have characteristic zero.
\end{proof}

\subsection*{Proof of \ref{B6.06c}}\label{prf:B6.06c}\phantomsection
\begin{proof}
Given a finite extension \(D/F\) and \(v\), the Galois closure of \(D\) contains finitely many quadratic fields. Choose a member split at \(v\) outside that finite list. It is linearly disjoint from \(D\). Thus the family satisfies the generality hypothesis with empty exceptional set in the prime-degree patching lemma of \cite[\S1]{sorensen-2020}. The other two hypotheses of that lemma are the stated invariance and compatibility. Its conclusion identifies the restriction at every member. The cofinite conclusion suffices here.
\end{proof}

\subsection*{Proof of \ref{B6.06d}}\label{prf:B6.06d}\phantomsection
\begin{proof}
For (1) apply cuspidal strong multiplicity one \cite{jacquet-shalika-1981}. At spherical places inverse transpose and contragredience both invert the Satake roots, so multiplicity one identifies the two representations and gives the nondegenerate bilinear pairing. Use \ref{B6.06a} for admissibility. Twist by a power of the adelic norm to make the split central character unitary. Write \(J\) for the annihilator of the twisted form in the centre of the enveloping algebra. It has finite codimension because the centre generates a finite-dimensional space on the form. Choose the compact idempotent for the finitely many archimedean types in its compact translates. Cuspidal rapid decrease gives square integrability modulo the split centre. Project onto the irreducible summands of the cuspidal discrete spectrum. Every projection has the same finite level, is fixed by the compact idempotent and is killed by \(J\), since the projections intertwine right translation and its differential. At this level the adelic quotient has finitely many arithmetic components. Harish-Chandra finiteness on each component gives a finite-dimensional space for this idempotent and ideal \cite[Thm.~8.4]{borel-2007}. The nonzero projections are mutually orthogonal in that space, so finitely many sum to the form. Undo the twist. For the newform in (3), adelisation gives the open level stabiliser and cuspidality. Its lowest-weight vector has one rotation weight, together with its reflected weight for the full compact group, and its infinitesimal character is the lowest-weight infinitesimal character. A norm twist preserves these finiteness conditions.
\leanref{UpperTheory.\allowbreak{}V.\allowbreak{}Af.\allowbreak{}cuspSemisimpleReal\_\allowbreak{}dB}{Development.\allowbreak{}AutomorphicArithmetic.\allowbreak{}AutomorphicVarieties.\allowbreak{}AutomorphicForms.\allowbreak{}TemperedLocalComponents.\allowbreak{}PCuspSSDB}{58--64}

Items (2), (4) and (5) use cyclic base change \cite[Ch.~III, Thms.~4.2 and~5.1]{arthur-clozel-1989}. In (4) and (5) strong multiplicity one turns cofinite invariance into invariance of the upper representation. The fibre characters are the characters trivial on idèlic norms. Restriction of a local unramified parameter takes Frobenius to its residue-degree power. At a split place the completions agree. At infinity restriction is restriction of the Weil parameter. A norm twist used to normalize the weights shifts both multisets equally and cancels on twisting back. For (3), adelize the lowest-weight newvector and apply the norm twist with the indicated infinitesimal character. The spherical Hecke polynomial of the newvector at \(p\nmid M\) is the polynomial displayed above, with one factor \(p^{1/2}\) from algebraic normalization and \(p^j\) from the twist \cite[Ch.~III]{jacquet-langlands-1970}. For (6), prescribe invariant \(1/2\) at each real place and zero at each finite place. They sum to zero by even degree. The quaternion ramification theorem gives the algebra, and a maximal order gives the integral splittings \cite[Thm.~14.6.1]{voight-2021}.
\end{proof}

\subsection*{Proof of \ref{B6.06e}}\label{prf:B6.06e}\phantomsection
\begin{proof}
Take the quadratic CM fields admitting cuspidal base change of \(\pi\). For a given finite \(v\mid q\), choose primes \(p\equiv-1\bmod q\) for odd \(q\), or \(p\equiv-1\bmod 8\) for \(q=2\). Omit \(2,q\), the primes ramified in \(F\), and the primes below nonspherical components. Dirichlet's theorem supplies infinitely many \(p\). Then \(-p\) is a square in \(\Q_q\), and \(L=F(\sqrt{-p})\) is split at \(v\) and ramified above \(p\). Its character cannot fix \(\pi\): at a place above \(p\) the spherical support is unramified and the quadratic twist has ramified support. By \ref{B6.06d}(2) the lift is cuspidal. These fields include infinitely many distinct members since each is ramified at its chosen prime.

The lift preserves regularity and pulls back the polarization. If necessary multiply the polarizing character on \(F\) by the quadratic character of \(L/F\) to give the sign required in \ref{B9.02}. Its norm pullback is unchanged and its sign changes at every real place. Apply \ref{B9.02} over these CM fields. On a pairwise compositum the two restrictions have the same Frobenius polynomials almost everywhere: for a place of residue degree \(f\) over \(F\) both multisets are \((q_v^{1/2}B_v)^f\), with powers taken root by root. The same computation gives invariance under conjugation. The restrictions are semisimple by normality and finite index. Chebotarev and equality of semisimple characters give the isomorphisms in \ref{B6.06c}. Patch to \(r\).

For each specified place and prolongation choose a split member outside the finite exceptions of the patching lemma. Its completion equals \(F_v\), its local automorphic component is \(\pi_v\), and its decomposition group on the chosen prolongation is the entire downstairs decomposition group. Thus the comparisons of \ref{B9.05} and \ref{B9.06} at a toral unlinked or spherical component apply to \(r\) itself. For the ratio in (2), a cuspidal representation becomes unitary after a global norm twist. In rank two an unlinked principal-series component in that normalization has exponents of absolute value less than \(1/2\), by the unitary classification \cite{tadic-1986}. Moduli of the Weil roots therefore have ratio less than \(q_v\), unchanged by the common normalization. At a spherical cusp component genericity excludes a linked pair. Finally the weights pull back under embeddings by \ref{B6.06d}(2). On one CM member \ref{B9.04} reflects the weights at conjugate embeddings over the same real place. They both pull back the downstairs multiset, so the same integer works over \(F\).
\end{proof}

\subsection*{Proof of \ref{B6.06f}}\label{prf:B6.06f}\phantomsection
\begin{proof}
Take a tower from \(K\) to \(L\) with cyclic prime-degree steps. All its fields contain the imaginary quadratic field \(K\), and all are CM since \(L/\Q\) is abelian. In a step \(M/D\) write \(W_D=W|_{G_D}\). Every \(W_D\) is irreducible because \(W|_{G_L}\) is. Suppose an attachment is known over \(M\). Its Satake multisets are invariant at conjugate good places since the attached representation extends to \(G_D\). By \ref{B6.06d}(4) choose a cuspidal descent \(\pi_D\). Its infinitesimal weights are regular algebraic.

We explain why this descent has a Galois representation to compare with \(W_D\). For any two-dimensional representation the exterior product is a perfect pairing with multiplier its determinant. Conjugation over \(D^+\) preserves \(W_D\) up to isomorphism because \(W\) is defined over \(\Q\). Over \(M\), the attachment and its conjugate dual, with the determinant and algebraic norm correction, consequently have equal good Satake multisets. In detail, if \(d=\det W\) and \(B\) is the Satake pair at a good place \(u\) of \(M\), then
\[
 B=\bigl(q_u^{-1}\iota d({\rm Frob}_u^{-1})\bigr)B^{-1}.
\]
The scalar is the value of the pullback of a Hecke character of \(D^+\): take the character of \(d\) by reciprocity and multiply by the idèlic norm, with both characters restricted by norm from \(\Q\). Thus \(\pi_D^c\) and the correspondingly twisted \(\pi_D^\vee\) have the same cuspidal base change. The fibre assertion of \ref{B6.06d}(4) identifies their discrepancy with a character of \(\Gal(M/D)\). This character extends to \(\Gal(M/D^+)\): the latter is abelian and complex conjugation splits its quotient of order two. Choose the extension trivial on this conjugation. After multiplying the polarizing character by this finite character or its inverse according to the direction of the comparison, strong multiplicity one gives the polarization of \(\pi_D\). The real signs agree since the original multiplier is restricted from \(\Q\) and complex conjugation is the same element at every real place in this abelian extension. Multiplying, if necessary, by the character of \(D/D^+\) makes all signs positive, as required in rank two, without changing the norm pullback. Thus \ref{B9.02} attaches a continuous semisimple \(R_D\) to \(\pi_D\).

At almost every \(u\mid v\) of \(M/D\), the Frobenius roots of \(R_D|_{G_M}\) and \(W|_{G_M}\) agree: taking \(f(u/v)\)-th powers commutes with the factor \(q_v^{1/2}\). Normality gives semisimplicity upon restriction, so Chebotarev identifies the restrictions. The space
\(\operatorname{Hom}_{G_M}(R_D,W_D)\) is one-dimensional by irreducibility. The action \(T\mapsto W_D(g)T R_D(g)^{-1}\) on this line defines a character \(\delta\) of \(\Gal(M/D)\). A nonzero \(T\) is invertible and gives \(W_D\simeq R_D\otimes\delta\). Twist \(\pi_D\) by the finite Hecke character with these Galois Frobenius values. Its attachment is \(W_D\). Repeating the step gives the specified attachment over \(K\).

At the last step \(K/\Q\), let \(\omega_\pi\) be the central character of this attachment. At good primes the determinant identity gives \(\omega_\pi|\cdot|^{-1}\) finite order, by reciprocity and character multiplicity one. Hence the two infinitesimal weights sum to \(1\). They are distinct and belong to \(\tfrac12+\Z\), so their difference is a nonzero even integer. This checks the gap in \ref{B6.06d}(5). Choose a rational cuspidal descent matching at the split primes. Its rank-two central character polarizes it, and \(\Q\) has just one real place, so \ref{B6.06e} supplies \(R_\Q\).

For each finite Galois extension \(D/K\), its normal closure over \(\Q\) contains \(K\). Chebotarev in this closure realizes every conjugacy class needed for \(D/K\) by Frobenius elements over rational primes split in \(K\): lift the chosen element to the subgroup fixing \(K\), and choose a prolongation with that representative. One may exclude finitely many primes. Thus Frobenius comparison at the split primes gives \(R_\Q|_{G_K}\simeq W|_{G_K}\). Apply the same one-dimensional Hom argument with quotient \(\Gal(K/\Q)\). The resulting trivial or quadratic twist of the rational descent attaches to \(W\).
\end{proof}

\subsection*{Proof of \ref{B6.06}}\label{prf:B6.06}\phantomsection
\begin{proof}
The coefficient, conjugation and counting assertions are \ref{B6.06b}. The totally real attachment is \ref{B6.06e}, which uses the patching lemma \ref{B6.06c}. The remaining real and cyclic assertions are \ref{B6.06a} and \ref{B6.06d}. In particular \ref{B6.06d}(5) supplies the split-prime descent entry of \(\mathscr H\), and \ref{B6.06f} derives descent for the specified Galois representation when that conclusion is needed.

Apply \ref{B6.15} to \(h_N\) for the small-twist bound. The common marking and the two attachments in the first count are supplied by \ref{B6.52}(1)--(3). With the vanishings of \ref{B6.43} these give \ref{B6.30}. The rank-four input of \ref{B6.49} gives the second and third counts in \ref{B6.31}. For the fifth class use \ref{B6.34}, with the Gaussian lifting assertion \ref{B6.51} supplied by \ref{B6.52}(4). These are the four counts in \ref{B6.07}, which gives the window bound.
\end{proof}

\subsection*{Proof of \ref{B6.07}}\label{prf:B6.07}\phantomsection
\begin{proof}
Apply \ref{B6.27} at $C_t=401$ with the small-twist bound and the local finiteness and duality inputs from $h_N$. The generator bounds of \ref{B6.18} follow from \ref{B6.17} and Ferrero--Washington vanishing of $\mu_\ell$ for abelian number fields \cite{ferrero-washington-1979}. On this window \ref{B6.39} gives the odd-prime plus bound from the supplied counts using \ref{B6.32} and \ref{B6.35}, with the class Z count supplied by \ref{B6.36}. Poitou--Tate duality and unramified orthogonality give the odd-prime minus bound of \ref{B6.13}. Combine these with the $2$-primary bound of \ref{B6.27} by \ref{B6.14}.
\end{proof}

\section*{Proofs for Chapter~\ref{ch:vwin-three-length-bounds}: Three length bounds}

\subsection*{Proof of \ref{B6.09a}}\label{prf:B6.09a}\phantomsection
\begin{proof}
Von Staudt--Clausen and $B_1=-1/2$ clear the coefficients of $B_j(X)$
with $\mathrm{primorial}(j+1)$ \cite[p.~56]{washington-1997}.
Since $f\mid N$, this gives integrality. At a level raised by a prime
$p$, Bernoulli distribution over all lifts gives the original sum.
If $p\nmid f$ and $p$ is new to the level, the omitted multiples of $p$
contribute $\chi(p)p^{j-1}$ times that sum. Iteration gives the product.
The functional equation and Euler product at $j$ give primitive
nonvanishing in the stated parity \cite[Ch.~4]{washington-1997}, and the
extra factors are nonzero since $|\chi(p)p^{j-1}|>1$.

The Fourier series gives $|B_j(x)|\le4j!/(2\pi)^j$ on $[0,1]$
\cite[p.~267]{apostol-1976}. For odd $h=2a+1$ the primes in $(a+1,h]$
divide $\binom{2a+1}{a}\le2^{2a}$. Induction gives
$\mathrm{primorial}(h)\le4^h$ since even $h>2$ adds no prime.
Use both bounds at $m=f,N$.
\end{proof}

\subsection*{Proof of \ref{B6.10}}\label{prf:B6.10}\phantomsection
\begin{proof}
Since $|\Gal(E/\Q)|\in O^\times$, inflation--restriction identifies
the twisted global cohomology with the invariants after restriction.
The finite quotients of inertia and decomposition groups have order
dividing $n$, so the same argument identifies the local kernels.
Untwisting takes the invariants to the $\theta$-part.

Let $\chi=\widetilde\theta^{\rm prim}$ have conductor $f$ and put $j=k+1$.
The bound at conductor $f$ suffices since $A_{N,j}(\widetilde\theta)=(N/f)A_{f,j}(\chi)$.
Inflate $\theta$ to $G_\Q$. Write $W=K_O/O$ and
$\alpha=\theta^{-1}\varepsilon_\ell^{-k}$. Complex conjugation $\delta$
acts by $-1$, so each class has a unique cocycle with $c(\delta)=0$.
At one lift over $\ell$, write $c|_{D_\ell}=(\alpha-1)m_\ell$ and choose
$a_0=\pi^h$ killing the finite image of $c$ and $m_\ell$. At the
$\sigma$-conjugate lift use $\alpha(\sigma)m_\ell-c(\sigma)\in W[a_0]$.
For $p\mid f$, $p\ne\ell$, primitivity and the prime-to-$\ell$ order
give $g\in I_p$ with $\alpha(g)-1$ a unit. The inertia cobounding vector
is $(\alpha(g)-1)^{-1}c(g)\in W[a_0]$. Subtract its coboundary on $D_p$.
The remaining values lie in $W^{I_p}=0$. Thus the normalized representatives
lie in the union of the modules $T(a_0)$ of normalized $W[a_0]$-valued
cocycles unramified outside $f\ell$ and coboundaries in $W[a_0]$ on every
$D_p$ for $p\mid f\ell$.

First suppose $\alpha(g_0)-\varepsilon_\ell(g_0)$ is a unit.
Take $M_0=\ell^m$ with $m\ge\max(1,v_\ell(f))$, and $H=f^{(\ell')}M_0$,
where $f^{(\ell')}$ is the prime-to-$\ell$ part of $f$. Write $[t]_M$
for least nonnegative residues. Let $\sigma_t$ lift the $t$-power action
on $\mu_H$ and $1\le b<H$ represent $g_0$ there. Put
\[
h_{M,b}(t)=\lfloor b[t]_M/M\rfloor-\lfloor b[-t]_M/M\rfloor,\qquad
z=\sum_{t\in(\Z/H\Z)^\times}h_{H,b}(t)\alpha(\sigma_t)^{-1}.
\tag{6.3}
\]
Set $C_{M,j}(u)=M^{j-1}B_j(u/M)$ and $d_j=2j\mathrm{primorial}(j+1)$.
For $u=[t]_H$, $a=\lfloor bu/H\rfloor$ and $v=bu-Ha$,
we have $h_{H,b}(t)=2a-(b-1)$. Expand the first two terms
$v^j/H-jv^{j-1}/2$ of $C_{H,j}(v)$. The lower terms have a factor $H$
after clearing denominators, so
\[
d_j\bigl(b^kh_{H,b}(t)u^k+2(C_{H,j}(v)-b^jC_{H,j}(u))/j\bigr)\in H\Z.
\]
Sum with weights $\chi_H(bt)$ and reindex the term at $v$. Since
$\alpha(\sigma_t)^{-1}\equiv u^k\chi(t)\bmod M_0O$, this gives
\[
d_j\bigl(b^k\chi(b)z-2(\chi(b)b^j-1)B_{j,\chi_H}/j\bigr)\in M_0O.
\tag{6.4}
\]
The choice of $g_0$ makes $\chi(b)b^j-1$ a unit. The level formula in
\ref{p6c2:bernoulli} preserves the valuation of $B_{j,\chi}$ here,
since the single possible extra factor $1-\chi(\ell)\ell^{j-1}$ is
$1\bmod\ell O$. Choose $m$ with $|M_0|_\ell<|a_0d_jB_{j,\chi}/j|_\ell$.
Then
\[
|z|_\ell=|B_{j,\chi}/j|_\ell,\qquad a_0z\mid M_0,\qquad
z\mid A_{f,j}(\chi).
\tag{6.5}\label{p6c2:z}
\]

Let $M>2$.
For a degree-one prime $P'\nmid M$ of $\Q(\mu_M)$, put $p=\mathrm NP'$.
Reduction identifies $\mu_M$ with its image in $\F_p^\times$,
so $M\mid p-1$. Let $\eta$ lift the $D$-th residue power for
$D=(p-1)/M$, and let $\sigma_{M,t}$ act by $t$-th powers on $\mu_M$.
For $1\le i\le M-2$, the characters $\eta,\eta^i,\eta^{i+1}$ are
nontrivial. The Gauss-sum identity
$J_i=J(\eta,\eta^i)=G(\eta)G(\eta^i)/G(\eta^{i+1})$ gives
$J_i\bar J_i=p$, since each Gauss sum has squared modulus $p$.
Modulo $\sigma_{M,t}^{-1}P'$ the sum is
$\sum_{x\in\F_p}x^{D[t]_M}(1-x)^{D[it]_M}$.
The powers of $x$ have positive degree below $2(p-1)$. The term of
degree $p-1$ occurs with coefficient prime to $p$ exactly when
$[t]_M+[it]_M\ge M$. Since $p$ is unramified, this proves
\[
v_{\sigma_{M,t}^{-1}P'}(J_i)=1-C_i(t),\qquad
C_i(t)=\lfloor([t]_M+[it]_M)/M\rfloor .
\tag{6.5a}\label{p6c2:jacobi}
\]
We may choose such a prime in any ideal class and avoid any finite set.
Chebotarev in the Hilbert class field gives positive density in the
class, and primes of residue degree greater than one over $\Q$ have
density zero \cite[Ch.~VII, Theorem~13.4]{neukirch-1999}.

Choose a place $A_0$ of $\overline\Q$ above a prime $r\equiv1\bmod H$.
Let $d$ be squarefree with prime factors $q\equiv1\bmod H$,
$q\ne r$, and $\ell\nmid\mathrm{ord}_q(r)$. Put $M=Hd$ and let $P$
lie below $A_0$ at level $M$. Choose $P'$ as above with $PP'=(y)$,
and put $J=\prod_{i=1}^{b-1}J_i$. The carries telescope to
$\lfloor b[t]_M/M\rfloor$, so $\bar J/J$ has valuation $h_{M,b}(t)$
at $\sigma_{M,t}^{-1}P'$. Hence
$X_d=\prod_t(\sigma_{M,t}^{-1}y)^{h_{M,b}(t)}/(\bar J/J)$ satisfies
\[
(X_d)=\prod_t(\sigma_{M,t}^{-1}P)^{h_{M,b}(t)},\qquad X_d\bar X_d=1.
\tag{6.5b}\label{p6c2:x-ideal}
\]
The second identity follows from $h_{M,b}(-t)=-h_{M,b}(t)$.

For a new prime $q$ of the same kind, let $\phi\in D_q$ lift arithmetic
Frobenius. It acts by $q$-th powers on $\mu_M$. Write $s=s(d,q)$ for the
residue degree over $r$ from level $M$ to $qM$, which divides
$\mathrm{ord}_q(r)$. Since $b<H<q$, the sum of the first floor in
$h_{qM,b}$ over all $q$ lifts $[t]_M+iM$ is
\[
\left\lfloor b[t]_M/M\right\rfloor+
\sum_{i=0}^{q-1}\left\lfloor ib/q\right\rfloor .
\]
$ib\bmod q$ runs over all residues. Remove the lift divisible
by $q$, whose first floor is $\lfloor b[q^{-1}t]_M/M\rfloor$,
and subtract the corresponding sum for $-t$. The sum over unit lifts
is $h(t)-h(q^{-1}t)$ for $h=h_{M,b}$. By \eqref{p6c2:x-ideal},
$\mathrm{Norm}(X_{qd})/(X_d/\phi^{-1}X_d)^s$ is a unit of modulus $1$
at every embedding. Kronecker's theorem makes it a root of unity,
whose order in $\Q(\mu_M)$ divides $2M$.

Choose $P',y$ at level $qM$ and take norms
to level $M$. The upper prime over $r$ has norm $P^s$, so the
expression for $X_d$ using these norms equals $X_d^s$ up to a
$2M$-th root of unity. At a prime above $q$, total ramification makes
the residue $F_t$ of $\sigma_{qM,\tilde t}^{-1}y$ independent of the
lift, and $F_t=F_{qt}^q$. Thus the upper principal factor reduces to
$\prod_tF_t^{h(t)-h(q^{-1}t)}=\prod_tF_{qt}^{(q-1)h(t)}$, the factor for
$\phi^{-1}\mathrm{Norm}(y)$. The two residue characters on $\F_p^\times$
satisfy $\eta_{qM}^q=\eta_M$. At $q$, $\eta_{qM}(u)$ and
$\phi^{-1}(\eta_M(u))$ have the same $q$-th power, so they agree after
reduction. The Jacobi sums and their conjugates therefore agree term
by term after applying $\phi^{-1}$ at level $M$. Their denominators
are units at $q$ by $J_i\bar J_i=p$. Hence
$X_{qd}\equiv\zeta'(\phi^{-1}X_d)^s$ at each prime over $q$,
where $(\zeta')^{2qM}=1$.

Let $u_0$ be the least nonnegative integer representing the action
of $g_0$ on $\mu_{M_0}$ and put $P_0(v)=g_0(v)/v^{u_0}$.
For $x_d=P_0(X_d)^{2f^{(\ell')}d}$ the root factors vanish.
Thus $\mathrm{Norm}(x_{qd})=(x_d/\phi^{-1}x_d)^S$ and
$x_{qd}\equiv(\phi^{-1}x_d)^S$ above $q$, with $S=qs(d,q)$.
Write $L$ for the list of primes of $d$ and $L'=(q,L)$.
Take $\tau_p$ generating the action on $\mu_p$ and acting trivially
on the other coprime cyclotomic factors. Let
$\mathcal D_p=\sum_{i=0}^{p-2}i\tau_p^i$ and apply
$\prod_{p\in L}\mathcal D_p$ multiplicatively to $x_d,x_{qd}$ to obtain
$y_L,w_L$, respectively. For the relative norm $N_q$,
$(\tau_q-1)\mathcal D_q=q-1-N_q$ gives
\[
\tau_q(y_{L'})/y_{L'}=w_L^{q-1}(\phi^{-1}y_L/y_L)^S.
\]
Put $U=G_{\Q(\mu_H)}$. For every $g\in U$, $g(y_L)/y_L$ has an
$M_0$-th root in $\Q(\mu_{Hd})$. For the empty list the quotient is
$1$. For each $p\in L$, put $p$ first and use induction on the
remaining list. The displayed identity gives the root
$w^{(p-1)/M_0}z_{\phi^{-1}}^S$, where $w$ is the derivative omitting
$p$ and $z_{\phi^{-1}}$ is the root at the lower level.
The Chinese remainder theorem gives generators for the relative
cyclotomic group from these $\tau_p$. Products extend the assertion
because $z_g\,g(z_h)$ is a root for $gh$ in the same field.

Choose $Y_L^{M_0}=y_L$ and roots $z_g$ as above. For primitive
$\zeta\in\mu_{M_0}$ and a generator $w$ of $W[M_0]$ over $O$, define
$\kappa_L:U\to\Z/M_0\Z$ and $\nu_L:U\to W[M_0]$ by
\[
\zeta^{\kappa_L(g)}=\frac{gP_0(Y_L)}{P_0(Y_L)P_0(z_g)},\qquad
\nu_L(g)=\sum_{t\in(\Z/H\Z)^\times}
\beta(\sigma_t)\kappa_L(\sigma_t^{-1}g\sigma_t)w,\quad
\beta=\varepsilon_\ell\alpha^{-1}.
\]
The root $z_g\,g(z_h)$ gives additivity since $P_0$ commutes with the
action on the cyclotomic layer and kills $\mu_{M_0}$. Reindexing
gives $\nu_L(\sigma g\sigma^{-1})=\beta(\sigma)\nu_L(g)$.
At $p\nmid Hrd$, inertia fixes the cyclotomic layer, so take $z_g=1$
for $g\in I_p$. Here $P_0(Y_L)$ is a unit by \eqref{p6c2:x-ideal}.
The resulting root of unity reduces to $1$, and reduction is
injective on $\mu_{M_0}$.

By total tame ramification over the level-$H$ field at $r$,
choose $\rho^{M_0}=r$ and $\tau_r\in I_{A_0}$ with $\tau_r\rho=\zeta\rho$. Normalize $v_{A_0}(r)=1$. For
$K(t)=v_{A_0}(\sigma_tP_0(y_\varnothing))$, the element
$\sigma_tP_0(Y_\varnothing)/\rho^{K(t)}$ is a unit. Its quotient under
$\tau_r$ belongs to $\mu_{M_0}$ and reduces to $1$. Thus
$t\kappa_\varnothing(\sigma_t^{-1}\tau_r\sigma_t)\equiv K(t)\bmod M_0$.
If $d_0$ represents $g_0$ on $\mu_H$, then $P_0$ sends a valuation
function $k(t)$ to $k(td_0)-u_0k(t)$. Reindex the sum with weights
$\alpha(\sigma_t)^{-1}$ to obtain the multiplier $\alpha(g_0)-u_0$
modulo $M_0O$. This occurs once in $x_1$ and once in $P_0(y_\varnothing)$.

At $q$, take $\tau_q$ in inertia and
$z_{\tau_q}=w_L^{(q-1)/M_0}z_{\phi^{-1}}^S$. Write $B=P_0(Y_L)$.
The residue formula for $w_L$ gives
$P_0(w_L)^{(q-1)/M_0}\equiv(\phi^{-1}B)^{S(q-1)}$.
The quotient for $\kappa_{L'}(\tau_q)$ therefore reduces to
$(\phi^{-1}B/(BP_0(z_{\phi^{-1}})))^S$, since inertia fixes residues
and $B\equiv(\phi^{-1}B)^q$. Use injectivity of reduction on
$\mu_{M_0}$ and apply the same computation at the conjugate places
to obtain
\[
\nu_\varnothing(\tau_r)=2f^{(\ell')}(\alpha(g_0)-u_0)^2zw,\qquad
\nu_{L'}(\tau_q)=-q\,s(d,q)\nu_L(\Frob_q).
\tag{6.6}\label{p6c2:inertia}
\]
The coefficients $2f^{(\ell')}(\alpha(g_0)-u_0)^2$ and $-q\,s(d,q)$
are units, as is $\beta(g_0)-1$. Equivariance extends $\nu_L$ to the
$\beta$-cocycle
\[
\widetilde\nu_L(g)=-(\beta(g_0)-1)^{-1}\nu_L(gg_0g^{-1}g_0^{-1}).
\]
The extension remains unramified outside $Hrd$ since the commutator
of $g_0$ with inertia there is a product of two such inertia elements.

Put $I_L=\langle\nu_L(U)\rangle_O$, $(A_L)=\mathrm{Ann}(I_L)$, and let
$S_L\subset T(a_0)$ consist of cocycles zero on decomposition groups
over $r$ and the primes in $L$. Since $\operatorname{Ann}(w)=(M_0)$,
\eqref{p6c2:inertia} gives
\[
\operatorname{Ann}(\nu_\varnothing(\tau_r))=(M_0/z),\qquad
a_0\mid M_0/z\mid A_\varnothing\mid M_0 .
\]
Suppose $a_0\mid A_L$ and $0\ne c\in S_L$. The twists on $U$ are trivial
modulo $M_0$. If $c|_U=0$, it descends to an abelian quotient, where
commutation with $\delta$ and $c(\delta)=0$ give $c=0$.
Let $V\subset U$ fix $\rho$. If $c|_V=0$, conjugation acts on $c(U)$
by both $\alpha$ and, through the radical quotient, $\varepsilon_\ell$.
Their unit difference at $g_0$ gives $c|_U=0$. The same argument for
$\nu_L\bmod\pi I_L$ uses $\beta(\delta)-\varepsilon_\ell(\delta)=2$
and Nakayama's lemma. Here $I_L\ne0$ since $a_0$ is a nonunit and
$a_0\mid A_L$. If $I_L=0$, the generation condition on $\nu_L$ is automatic.

The kernel of $c|_V$ and the kernel of $\nu_L|_V\bmod\pi I_L$ are
proper open subgroups of $V$. The subgroup fixing $\mu_{H\ell}$ is
also proper. The degree-$\ell$ extension $\Q(\mu_{H\ell})/\Q(\mu_H)$
is unramified at $r$, whereas $\Q(\mu_H,\rho)/\Q(\mu_H)$ is totally
ramified there. Three proper subgroups covering a group must all have
index two, and the first kernel has odd index. Their complement
in $V$ is thus nonempty and open. Chebotarev gives a prime $q\nmid Hrd$
with a Frobenius lift $\phi$ in this complement. We have
$c(\phi)\ne0$ and $\nu_L(\phi)$ generates $I_L$. Also $r$ is an
$M_0$-th power modulo $q$ and $v_\ell(q-1)=m$, so
$\ell\nmid\mathrm{ord}_q(r)$.

By \eqref{p6c2:inertia}, the annihilator of $y=\nu_{L'}(\tau_q)$ is
$(A_L)$. Hence $a'=A_{L'}=A_Le$ for some $e\in O$, with
$a_0\mid a'\mid M_0$. Use a generator of $W[a']$ to identify it with
$B'=O/(a')$, in which $y$ generates the ideal $(e)$. The
$R$-linear forms $\lambda:B'\to R$, for $R=\Z/M_0\Z$, separate points.
Each cyclic summand $\Z/\ell^s\Z$ of the finite $R$-module
embeds in $R$ by $1\mapsto\ell^{m-s}$.

Cup $\widetilde\nu_{L'}$ with $c'\in S_L$ using multiplication
followed by $\lambda$. The twists give coefficients $R(1)$.
Local triviality of $c'$ kills the cups at $r,L$ and $p\mid H$.
For $p\ne\ell$, the unramified subgroups for a finite unramified
$\ell$-primary module and its Tate dual are exact annihilators.
Thus the other cups away from $q$ vanish
by \ref{B6.05d}(2) and \cite[Theorem~I.2.6]{milne-2006}. By \ref{B6.09}, the invariant at
$q$ vanishes as well, so the local cup is zero.

Since $q\equiv1\bmod H$, $D_q$ acts trivially on the finite coefficient
modules. For an unramified $R$-valued character $u$,
$u(h)c'(\phi)=u(\phi)c'(h)$ on $D_q$ by agreement on inertia and on
$\phi$. Thus $g\mapsto\lambda(\widetilde\nu_{L'}(g)c'(\phi))$ cups
trivially with every such $u$. Exact orthogonality makes this
homomorphism unramified, so $\lambda(yc'(\phi))=0$ for all $\lambda$.
It follows that $yc'(\phi)=0$. Evaluation at $\phi$ maps $S_L$ to
$W[e]$ with kernel $S_{L'}$, since $c'$ is unramified at $q$ and
vanishing on $D_q$ passes to all conjugate lifts. Therefore
\[
\mathrm{length}_O S_L\le\mathrm{length}_O S_{L'}+v_O(e).
\tag{6.7}\label{p6c2:descent}
\]
For the initial step use the same pairing at $r$ with
$\widetilde\nu_\varnothing$, $c'\in T(a_0)$ and
$A_\varnothing=(M_0/z)e_0$. This gives
$\mathrm{length}_O(T(a_0)/S_\varnothing)\le v_O(e_0)$.
At each subsequent step $v_O(e)>0$ since $c(\phi)\ne0$.
The bound $A_L\mid M_0$ forces termination at $S_{L_*}=0$.
Summing \eqref{p6c2:descent} gives
\[
\mathrm{length}_O T(a_0)\le
v_O(e_0)+v_O(A_{L_*})-v_O(A_\varnothing)
=v_O(A_{L_*})-v_O(M_0/z)\le v_O(z).
\]
By \eqref{p6c2:z}, the increasing images of $T(\pi^h)$ in $C_\theta$
have lengths bounded by $v_O(A_{f,j}(\chi))$. They stabilize.

In the remaining case $\alpha\equiv\varepsilon_\ell\bmod\pi$,
the prime-to-$\ell$ order of $\theta$ makes it unramified off $\ell$.
Normalize at $t\in I_\ell$ acting by $-1$ on $\mu_\ell$. Then
$c|_{D_\ell}=0$ since $\alpha(t)-1$ is a unit. For a $W[\pi]$-valued
cocycle, each $\F_\ell$-linear form on $W[\pi]$ gives a Kummer class
of $\pm\ell^s$. Local triviality gives $\ell\mid s$, so the class
vanishes. Its cocycle vanishes too by normalization at $t$.
The forms separate points, and induction on the power of $\pi$
killing the cocycle gives $C_\theta=0$.
\end{proof}

\subsection*{Proof of \ref{B6.10a}}\label{prf:B6.10a}\phantomsection
\begin{proof}
For $N\le2$ the class group is trivial. Otherwise take a degree-one
prime $P\nmid N$ in $c$ as above, and put $p=\mathrm NP$, $P_t=\tau_t^{-1}P$.
For $i\not\equiv0,-1\bmod N$, \eqref{p6c2:jacobi} gives
$\prod_tP_t^{C_i(t)}=(p/J_{[i]_N})$ at level $N$. For
$i\equiv0,-1$ the product is $(1),(p)$, respectively.
The identity $NC_i(t)=[it]_N+[t]_N-[(i+1)t]_N$ shows that
$\sum_{0\le i<a}C_i(t)=\lfloor a[t]_N/N\rfloor$. Apply $\tau_a^{-1}$
to the product over $0\le i<a$. The exponent at $P_b$ becomes $j_b$.
Since $\prod_bP_b=(p)$, the fractional ideal with exponents $E_b$
is principal.

As $K$ is totally complex, $h$ factors through the class group by
inverse Artin reciprocity. That map sends $\sigma_b^{-1}u\sigma_b$
to the $\tau_b^{-1}$-translate of the class for $u$
\cite[Appendix, \S3]{washington-1997}. The class identity gives the
homomorphism identity.
\end{proof}

\subsection*{Proof of \ref{B6.12}}\label{prf:B6.12}\phantomsection
\begin{proof}
Put $H=N\ell$, $b=2H+1$, $j=k+1$ and $\Gamma=\Gal(E/\Q)$.
Write $\sigma_t$ for the cyclotomic action of $t$ on $E$. Set
\[
\Theta_j=-\sum_{t\in(\Z/H\Z)^\times}C_{H,j}([t]_H)[\sigma_t]/j,\qquad
Y=(1-b^j[\sigma_b])\Theta_j.
\]
For $L=H\ell^m$, Bernoulli distribution gives
$\sum_{i=0}^{L/H-1}C_{L,j}(u+iH)=C_{H,j}(u)$.
Since $\ell\mid H$, every lift of a unit class modulo $H$ is a unit
modulo $L$. Thus the same formula at level $L$ gives $Y$ after
summing over fibres. For $v=[b^{-1}t]_L$ and $a_t=(bv-[t]_L)/L$,
expansion of the first two Bernoulli terms gives
\[
-\frac{C_{L,j}([t]_L)-b^jC_{L,j}(v)}j
\equiv -[t]_L^k((b-1)/2-a_t)\pmod{(L/d_j)\Z}.
\]
Hence the image $z$ of $Y$ lies in $\Z_\ell[\Gamma]$. For
$v_\ell(L)\ge u+v_\ell(d_j)+1$, the integral sum of the right-hand
coefficients at level $L$ represents $z$ modulo $\ell^u$.

For $x\in\Sha_{\ell,k}(E)[\ell^u]$, choose a cocycle $\xi$ representing
$x$. Enlarge $L$ until the twist is trivial on the finite image of
$\xi$ and on local cobounding vectors at $\ell$. Thus
$h=\xi|_{G_{\Q(\mu_L)}}$ is an unramified homomorphism.
Let $U_\infty=\ker(\varepsilon_\ell|_{G_E})$ and
$f=\operatorname{res}_{U_\infty}(zx)$. For
$g\in U_\infty\cap G_{\Q(\mu_H)}$, $g$ and its conjugates fix
$\mu_L$. The cyclotomic weight $[t]_L^k$ cancels the action of the
$t$-power lift on the restricted cocycle values. For such a lift $\sigma_t$,
\[
f(g)=-\sum_t\bigl((b-1)/2-a_t\bigr)h(\sigma_t^{-1}g\sigma_t)=0
\]
by \ref{B6.10a}. Since $g^{\varphi(H)}\in G_{\Q(\mu_H)}$ for every
$g\in U_\infty$, we have $\varphi(H)f=0$.
For odd $\ell$ the cyclotomic image of $G_E$ is open and procyclic
in $\Z_\ell^\times$. A topological generator $\gamma$ has infinite
order, so multiplication by $\gamma^{-k}-1$ is surjective on
$\Q_\ell/\Z_\ell$. Subtracting the coboundary that matches a
cocycle on $\gamma$ kills it by continuity. Inflation--restriction
therefore makes restriction to $U_\infty$ injective, and
$z\varphi(H)x=0$.

Let $c$ be complex conjugation. In
$\Z_\ell[\Gamma]/(c-(-1)^j)$ the adjugate of multiplication by $z$
gives $\Delta\varphi(H)\Sha^-_{\ell,k}(E)=0$, where $\Delta$ is
its determinant. Put $d_Y=Hj(j+1)!$ and $u_Y=2b^jH^jj!$.
Clearing coefficients of $B_j$ gives
$H\mathrm{primorial}(j+1)C_{H,j}(u)\in\Z$. Since
$\mathrm{primorial}(j+1)\mid(j+1)!$, we have $d_YY\in\Z[\Gamma]$.
The $n/2$ eigenvalues $(1-b^j\theta(b))(-B_{j,\theta}/j)$, with
$\theta$ pulled back from $\Gamma$ to level $H$ in the stated parity,
are nonzero by \ref{p6c2:bernoulli}. The Fourier bound for $B_j(x)$
gives $|B_{j,\theta}|\le4H^jj!/(2\pi)^j$, so each eigenvalue has
absolute value at most $u_Y$. Thus $D=d_Y^{n/2}\Delta\in\Z$ and
\[
0<|D|\le(d_Yu_Y)^{n/2},\qquad
\log(d_Yu_YH)\le30j\log(k\ell N).
\tag{6.9a}\label{p6c2:det}
\]
Put $T=3kH$. Then $2,H,j\le T$. Also $b\le T^3$ and
$j+1\le T^2$. These give $d_Yu_YH\le T^{7j+6}\le T^{10j}$, and
$\log T\le3\log(kH)$. Taking $r=v_\ell(D\varphi(H))$ gives
\eqref{p6c2:exponent}.

If $z_0$ represents a class killed by $\ell$, write $\ell z_0=dm$,
where $d$ is the coboundary map. Take $\ell m_1=m$. Then
$z_0-dm_1$ has values killed by $\ell$. Restrict to $E'=E(\mu_\ell)$,
where the residual twist is trivial, to obtain an $\F_\ell$-valued
homomorphism unramified outside the places $S$ over $\ell$.
Choose one lift per class. Equal restricted homomorphisms give
equal original classes by restriction and corestriction, since
$[E':E]$ is prime to $\ell$. Kummer theory identifies the group of
these unramified-outside-$S$ homomorphisms with
\[
V_S=\{a\in E'^\times:v(a)\equiv0\bmod\ell\ (v\notin S)\}/E'^{\times\ell}.
\]
Its quotient by the $S$-units modulo $\ell$-th powers is
$\mathrm{Cl}(\mathcal O_{E',S})[\ell]$. For $m'=[E':\Q]$, the unit
rank is $r_1+r_2-1+|S|$, with $(r_1,r_2)$ the signature and
$|S|\le m'$. The roots of unity contribute at most $\ell$.
Hence $|V_S|\le\ell^{2m'}h(E')$.

Each of the $m'$ field characters has conductor dividing $H$.
The conductor--discriminant formula gives
$\mathfrak D=|\operatorname{disc}E'|\le H^{m'}$
\cite[Ch.~VII, Theorem~11.9]{neukirch-1999}. By Minkowski's bound
each ideal class has a representative of norm at most
$Z=\sqrt{\mathfrak D}$ \cite[Ch.~I, \S5]{neukirch-1999}.
The ideal count at each norm is bounded by the $m'$-fold divisor
function. Sum over the first $m'-1$ factors using
$\sum_{a\le Z}1/a\le1+\log Z$ to obtain
\[
\log h(E')\le\log Z+(m'-1)\log(1+\log Z)\le3m'\log H ,
\tag{6.9b}\label{p6c2:class}
\]
since $m'\le H$. With $m'\le\ell n$ and $\log\ell\ge\log2>5/8$,
the $\ell$-torsion dimension is at most $8\ell n\log H$.
For $P=\Sha^-_{\ell,k}(E)$ and each $i\ge0$,
$P[\ell^{i+1}]/P[\ell^i]\hookrightarrow P[\ell]$ by
$[x]\mapsto\ell^ix$. Sum these bounds up to the exponent in
\eqref{p6c2:exponent} to obtain \eqref{p6c2:at-degree}.

The Galois group of the stated compositum embeds in a $2$-group
times $\Gal(K_M/\Q)$, which proves the assertion about odd divisors.
\end{proof}

\subsection*{Proof of \ref{B6.14}}\label{prf:B6.14}\phantomsection
\begin{proof}
Fix $\varepsilon,a$ as in \ref{B6.03} and put $m=\mathrm{cond}\psi$.
The page and height conditions for $A_1=4$ follow from
\[
4\le x^2\le2^{a(x)}\le2x^{2+8\varepsilon}\le x^4,\qquad
2\le t(x):=\lfloor x^{C_t}\rfloor\le x^{C_t}\qquad(x\ge2).
\tag{6.13}\label{p6c2:page}
\]
Take the joint window constants $J_0>0,C,P,X_1$, put $J=\max(J_0,1)$
and $U=X^{1/J}$. For sufficiently large $X\ge X_1$, \ref{B6.02} gives
$q,d>\max(2X^{C_t}+1,P,m)$ for every window pair at $U$, and
\[
\log(qd)\le U^{\varepsilon/5}(1+\log^2U)
\le U^{\varepsilon/4}\le X^{\varepsilon/(4J_0)} .
\tag{6.14}\label{p6c2:pair}
\]
Thus $\gcd(qd,m)=1$ and the joint window gives
$U\le X^{1/J_0}\le Y$ with $Y^2\le X$.

Fix $x\in[Y,Y^2]$ and $E$ as in \ref{B6.03}. In the common cyclotomic
layer the four conditions of \ref{B6.02} fix, respectively,
$\Q(\mu_{2^{a(x)+2}})^+$, the field $K_\psi$ cut out by $\psi$,
and the two quadratic character fields. Since $q,d\equiv3\bmod4$,
these last fields are $\Q(\sqrt{-q})$ and $\Q(\sqrt{-d})$.
The intersection of the fixing subgroups corresponds to their
compositum, which is abelian and CM. The same intersection describes
the field at level $\operatorname{lcm}(2^{a(x)+2},m,q,d)$, so
$n\le F\le2^{a(x)+2}mqd$. By \eqref{p6c2:page} and \eqref{p6c2:pair},
for sufficiently large $U$ and all $x$ in the window,
\[
\log(Ft)\le3\log2+\log m+(2+8\varepsilon+C_t)\log x+\log(qd)
\le2x^{\varepsilon/4}.
\tag{6.15}\label{p6c2:log}
\]
Put $W=x^{\varepsilon/4}$, $b=x^{4\varepsilon/5}$ and $o=[K_\psi:\Q]$.
By \eqref{p6c2:minus} the negative sum is at most $B_-(240,o,n,F,t)$.
By \eqref{p6c2:minus-exp} and \eqref{p6c2:log} this is at most
$Ke^{8W}$ for $K=3+240\sum_{\ell\mid o,\ \ell\ {\rm odd\ prime}}\ell^3$.
Combine \eqref{eq:B6-two-sum} and \eqref{eq:b6-pluswindow} with this
bound in \ref{B6.08}. The resulting total is at most
\[
e^b\log2+Ke^{8W}+(Ft)^C+Ce^b
\le(\log2+K+1+|C|)e^{(9+2|C|)b}\le e^{x^\varepsilon/2}.
\]
The last inequality holds uniformly for $x\ge U$ as $U\to\infty$,
since $x^\varepsilon/b=x^{\varepsilon/5}$.
\end{proof}

\section*{Proofs for Chapter~\ref{ch:B6-cyclotomic-window}: The cyclotomic tower and a common window}

\subsection*{Proof of \ref{B6.15}}\label{prf:B6.15}\phantomsection
\begin{proof}
Write $A=(\Q_\ell/\Z_\ell)(-k)$ and let $\lambda:M\to P$ be localization, with $M=H^2(G_{L,S_\ell},\Z_\ell(k+1))$ and $P=\prod_{v\mid\ell}H^2(L_v,\Z_\ell(k+1))$. Write $\operatorname{len}$ for $\Z_\ell$-length. The reductions in degree two identify the finite-coefficient localization map with $M/\ell^r M\to P/\ell^r P$. The degree-two reduction maps are surjective with kernels the indicated multiples, and $M$ and $P$ are finitely generated.

Since $k\ne0$, the local cyclotomic images give finite $H^0(L_v,A)$. Choose $\ell^a$ killing all these invariants. Finite local Tate duality gives
$\#(P/\ell^rP)=\prod_{v\mid\ell}\#H^0(L_v,A[\ell^r])$. The right side stabilizes, so $P$ is finite. Choose $\ell^b$ killing $P$. For $r\ge a+b+1$ the finite degree-two localization kernel is
$\ker\lambda/(\ker\lambda\cap\ell^r M)$. By \ref{B6.09} it has the same order as the degree-one localization kernel for $A[\ell^r]$.

The latter kernel injects into $\Sha_{\ell,k}(L)$ for $r\ge a$. Indeed, a representing cocycle that becomes a global coboundary $(g-1)t$ in $A$ is locally $(g-1)t_v$ for $t_v\in A[\ell^r]$. Then $t-t_v\in H^0(L_v,A)\subset A[\ell^r]$, so $t\in A[\ell^r]$. The images of the finite-level kernels increase and exhaust Sha. A continuous representative and the finitely many local coboundaries all occur at a finite level. The filtration $\ker\lambda\cap\ell^r M$ is separated. Its quotient orders therefore give $\operatorname{len}\ker\lambda=s_\ell(L,k)$. The finite-level cokernel has order $\#H^0(L,A[\ell^r])$ by \ref{B6.09}. For large $r$ this gives $\operatorname{len}\operatorname{coker}\lambda=\operatorname{len}H^0(L,A)$. Thus
\begin{equation}\label{eq:B6-H2-local-correction}
 \begin{gathered}
 V_\ell(L,k)+\operatorname{len}H^0(L,A)
 =s_\ell(L,k)+\sum_{v\mid\ell}\operatorname{len}H^0(L_v,A),\\
 s_\ell(L,k)\le V_\ell(L,k)
 \le s_\ell(L,k)+\sum_{v\mid\ell}\operatorname{len}H^0(L_v,A).
 \end{gathered}
\end{equation}

If $u\in\Z_\ell^\times$, $c=v_\ell(u-1)\ge1$ is finite, and $c\ge2$ for $\ell=2$, then
\begin{equation}\label{eq:B6-lte}
 v_\ell(u^e-1)=c+v_\ell(e)\qquad(e\ge1).
\end{equation}
Write $e=\ell^j e_0$ with $\ell\nmid e_0$. In the binomial expansion of $u^{e_0}-1$ the linear term has valuation $c$ and all higher terms have larger valuation. For a unit $1+t$ of the indicated level, $\ell t$ is the unique term of least valuation in $(1+t)^\ell-1$. Iterating $j$ times proves the formula.

The local cyclotomic image has index $d_v\le[L_v:\Q_\ell]\le n$ and contains $(1+\ell^2)^{d_v}$. By \eqref{eq:B6-lte}, its invariants on $A$ have length at most $2+v_\ell(d_vk)$. Hence
\begin{equation}\label{eq:B6-local-H0}
 \sum_{v\mid\ell}\operatorname{len}H^0(L_v,A)
 \le n(2+\lfloor\log_\ell n\rfloor+k).
\end{equation}
Multiplication by $\log\ell$ bounds this term by $n(3+k)\log(n+1)$ since $\ell\le n+1$. Apply \eqref{eq:B6-small-Sha} in \eqref{eq:B6-H2-local-correction} and increase the exponent to absorb this term. A module of length $v<\infty$ has order $\ell^v$.
\end{proof}

\subsection*{Proof of \ref{B6.16}}\label{prf:B6.16}\phantomsection
\begin{proof}
Write $Y=F^\times$ additively and $\partial:Y\to I=\operatorname{Div}(F)$ for the divisor map. Put $g=\sigma^{-1}$ and $\delta=g-a$, so $(g\mathfrak d)(P)=\mathfrak d(\sigma P)$. Set
\[
 U=\partial^{-1}(\delta I),\qquad W=U/(\delta Y+Y_{\mathrm{tors}}),
 \qquad C=I/(\partial Y+\delta I),\qquad
 E=\frac{\ker\partial}{\delta(\ker\partial)+Y_{\mathrm{tors}}}.
\]
Since $g^h=1$ and $a>1$, $\delta$ is injective on $I$. The snake lemma gives $0\to E\to W\to\ker(\delta|\mathrm{Cl}(F))\to0$. The last group has order $\#C$, so
\begin{equation}\label{eq:B6-W-factor}
 \#W=\#E\,\#C.
\end{equation}
On complex places $g$ has $r_2(L)$ cycles of length $h$. The logarithmic unit lattice omits the constant direction, of eigenvalue $1-a$ for $\delta$. Taking determinants on this lattice gives
\begin{equation}\label{eq:B6-unit-det}
 \#E\,(a-1)=(a^h-1)^{r_2(L)}.
\end{equation}
Also $\mathcal E_{\nu,k}=\operatorname{Hom}(C,\Z/\ell^\nu\Z)$ has order at most $\ell^{v_\ell(\#C)}$. Formula \eqref{eq:B6-lte} gives $v_\ell(a-1)=c+v_\ell(k)$ and $v_\ell(a^h-1)=\nu+v_\ell(k)$.

In additive notation $\mathcal N_a=\sum_{j=0}^{h-1}a^jg^{-j}$, so $\mathcal N_a\delta=\delta\mathcal N_a=(1-a^h)g$. Thus
\[
 f:W\longrightarrow F^\times/(F^{\times\ell^\nu}\mu(F)),\qquad
 [\alpha]\longmapsto[\mathcal N_a(\alpha)^{\ell^b}]
\]
is well defined, with $\mu(F)$ the group of roots of unity. For $[u]\in\ker f$, write $\ell^b\mathcal N_a u=\ell^\nu y$ modulo torsion. Apply $\delta$ and cancel $\ell^\nu$ modulo torsion. This puts $\ell^b(1-a^h)gu/\ell^\nu$ in $\delta Y+Y_{\rm tors}$. Since $g$ is invertible on $W$, the kernel is killed by $\ell^b(a^h-1)/\ell^\nu$, of valuation $b+v_\ell(k)$. Representatives of distinct image classes satisfy the counting hypothesis, so $\#\operatorname{im}f\le B$.

For $[u]\in W[\ell]$ choose $y$ with $\ell u=\delta y$ modulo torsion. If $y-y'=\ell z$ modulo torsion for two such choices, then $u-u'=\delta z$ modulo torsion. Thus the classes of $y$ modulo $\ell Y+Y_{\mathrm{tors}}$ distinguish the classes in $W[\ell]$. Because $u\in U$, choose $D_u\in I$ with $\partial u=\delta D_u$. Taking divisors gives $\delta(\partial y-\ell D_u)=0$, so $\partial y=\ell D_u$. Also $a\equiv1\bmod\ell$ implies invariance of the class of $y$ under $g$. In multiplicative notation, for a representative $\alpha$ of this class and a generator $\rho$ of $\mu(F)$, write
\[
 \sigma(\alpha)=\alpha v^\ell\rho^r,\qquad 0\le r<\ell,\qquad
 (N_{F/L}v)^\ell=N_{F/L}(\rho)^{-r}.
\]
There are at most $\ell^2$ pairs $(r,N_{F/L}v)$. For equal pairs Hilbert 90 applied to the ratio of the $v$'s gives a representative in $L^\times$ for the ratio of the $\alpha$'s modulo $\ell$th powers. The action of $\sigma$ on $\zeta$ has order $h$ by \eqref{eq:B6-lte}, so $F=L(\zeta)$ is unramified away from $\ell$. The descended ratios have valuations divisible by $\ell$ at primes away from $\ell$. Valuations above $\ell$ contribute at most $[L:\Q]$ dimensions modulo $\ell$. Classes with all valuations divisible by $\ell$ map to $\mathrm{Cl}(L)[\ell]$, with kernel the units modulo $\ell$th powers, of dimension at most $[L:\Q]$. This gives $\#W[\ell]\le\ell^{g(L)}$.

Put $K=(\ker f)[\ell^\infty]$. It is killed by $\ell^{b+v_\ell(k)}$, and $K[\ell]\subseteq W[\ell]$ gives at most $g(L)$ cyclic factors. Thus $\#K\le\ell^{(b+v_\ell(k))g(L)}$ and
\[
 \ell^{v_\ell(\#W)}
 \le\#\operatorname{im}f\,\ell^{(b+v_\ell(k))g(L)}
 \le B\,\ell^{(b+v_\ell(k))g(L)}.
\]
Equations \eqref{eq:B6-W-factor} and \eqref{eq:B6-unit-det} give \eqref{eq:B6-finite-class-count}.

To bound the weighted norms, first assume
\begin{equation}\label{eq:B6-room}
 2k+1<\ell^c,\qquad 2k+1<r_2(L).
\end{equation}
Put $G_r=\GL_r(\mathcal O_L)$ and let $T_1,T_2$ be the torsion orders of $H_i(G_{4k+6},\Z)$ for $i=2k+1,2k$, respectively. They are finite by \ref{B6.16c}. The counting hypothesis holds with
\begin{equation}\label{eq:B6-few}
 b=2k+(2k+1)(c+2k+1),\qquad B=(\ell^\nu)^{r_2(L)}T_1T_2.
\end{equation}
Put $q_0=\ell^\nu$, $A_\nu=\Z/q_0\Z$, $\kappa=1+\ell^c$ and $M_*=2k+1$. Define
\[
 P(T)=\prod_{\substack{0\le j<M_*\\j\ne k}}(T-\kappa^j),\qquad
 Q(T)=(T-a)P(T),\qquad p_{k,c}=P(a).
\]
For $0<d<M_*$, formula \eqref{eq:B6-lte} gives $v_\ell(\kappa^d-1)=c+v_\ell(d)<c+M_*$. Apply this to each factor of $P(a)$. In $A_\nu$,
\begin{equation}\label{eq:B6-pe}
 p_{k,c}\mid\ell^{M_*(c+M_*)}.
\end{equation}

Let $C_{q_0}=\langle\zeta\rangle$, $D=C_{q_0}\times C_{q_0}$ and $d_t(s)=(s,s^t)$ for $0\le t\le M_*$. Over $A_\nu$ the bar cycle $z$ is the sum of
$[\zeta^{i_1}\mid\zeta\mid\cdots\mid\zeta^{i_k}\mid\zeta]$ over $i_1,\ldots,i_k\bmod q_0$. Translation of the indices cancels its boundary modulo $q_0$. The tensor product of the two cyclic periodic resolutions has $2k+1$ basis vectors in degree $2k$ and zero differentials over $A_\nu$. Take a primitive integral relation among lifts of the $2k+2$ vectors representing $(d_t)_*[z]$. On reduction,
\[
 \sum_t n_t(d_t)_*z=\partial\xi,\qquad n_{t_0}\in A_\nu^\times
\]
for some $t_0$.

Identify $D$ with $A_\nu^2$ by exponents and put $\pi_t(i,j)=i+tj$. We construct $\eta\in\Z[A_\nu^2]$ with
\[
 (\pi_t)_*\eta=0\ (t\ne t_0),\qquad
 (\pi_{t_0})_*\eta=\sum_i[T^i]Q(T)[\kappa^i].
\]
For $J=(X^{\ell^c}-1)\subset\Z[X]$, the substitution $S(f)(X)=f(X^\kappa)$ acts by $\kappa^j$ on $J^j/J^{j+1}$. Indeed $S(f)\equiv f\bmod J$ and
\[
 S(X^{\ell^c}-1)=(X^{\ell^c}-1)\sum_{i=0}^{\kappa-1}X^{i\ell^c}.
\]
Apply the factors of $Q(S)$ in order of increasing exponent to obtain $Q(S)X\in J^{M_*}$. Under $X\mapsto[1]\in\Z[A_\nu]$, the desired pushforward is therefore a multiple of $([\ell^c]-[0])^{M_*}$. For $t\ne t_0$, condition \eqref{eq:B6-room} gives $t_0-t\mid\ell^c$ in $A_\nu$. Geometric summation gives $[t_0]-[t]\mid[\ell^c]-[0]$. Divide the desired pushforward by the product of these $M_*$ factors and lift the quotient by $i\mapsto(i,0)$. Multiply this lift by $\prod_{t\ne t_0}([(0,1)]-[(t,0)])$ to obtain $\eta$.
\leanref{UpperTheory.\allowbreak{}V.\allowbreak{}ZK.\allowbreak{}sEta}{Development.\allowbreak{}AutomorphicArithmetic.\allowbreak{}AutomorphicVarieties.\allowbreak{}ModularityArgument.\allowbreak{}ClassCounting.\allowbreak{}ExpStep}{116--182}

For $\partial\alpha=(g-a)\mathfrak d$, put $\beta=\prod_i\sigma^{-i}(\alpha)^{[T^i]P(T)}$. Then $\partial\beta=Q(g)\mathfrak d$. On $F^D$ use the diagonal representation $\rho_0$ with character $(s_1,s_2)\mapsto\zeta^{is_1+js_2}$ at coordinate $(i,j)$. It preserves $\mathcal O_F^D$ and the lattice whose coordinate $(i,j)$ is the ideal of divisor $\eta(i,j)\mathfrak d$.

Construct an $L$-linear isomorphism $G_t$ from the standard to the second lattice in scalar blocks $\pi_t=x$. For $t\ne t_0$ each block has total target divisor zero. For $t=t_0$ apply $\sigma^{-i}$ to coordinates with $\pi_{t_0}=\kappa^i$, taking $i$ modulo $h$. The scalar becomes $\zeta$ and the divisors transform by $g^i$. Leave other coordinates fixed. The exponent-$1$ block now has total divisor $\partial\beta$ and all other blocks have divisor zero. Steinitz classification gives an isomorphism in each block. Multiplication by a unit on an ideal summand adjusts the determinant to $\beta$ on the distinguished block and $1$ elsewhere. Undo the coordinate conjugations. The resulting $G_t$ commutes with $\rho_0d_t$.
\leanref{UpperTheory.\allowbreak{}V.\allowbreak{}ZK.\allowbreak{}Cl.\allowbreak{}cBetaDiv}{Development.\allowbreak{}AutomorphicArithmetic.\allowbreak{}AutomorphicVarieties.\allowbreak{}ModularityArgument.\allowbreak{}ClassCounting.\allowbreak{}LatticeBlockClasses.\allowbreak{}CBetaDiv}{33--62}
\leanref{UpperTheory.\allowbreak{}V.\allowbreak{}ZK.\allowbreak{}Cl.\allowbreak{}cGtEq}{Development.\allowbreak{}AutomorphicArithmetic.\allowbreak{}AutomorphicVarieties.\allowbreak{}ModularityArgument.\allowbreak{}ClassCounting.\allowbreak{}LatticeBlockClasses.\allowbreak{}GtEqZeta}{54--106}

Let $H$ be the stabilizer of the standard lattice. Put $G_*=G_{t_0}$, $\rho_1=G_*^{-1}\rho_0G_*$, and $b_t=G_*^{-1}G_t$. Both representations take values in $H$, and $b_t\in H$ satisfies $\operatorname{conj}_{b_t^{-1}}\rho_1d_t=\rho_0d_t$, where $\operatorname{conj}_b(v)=bvb^{-1}$. The bar prism $p_b$ inserts $b$ with alternating signs and conjugates subsequent entries by $b^{-1}$. Its boundary satisfies $\partial p_b+p_b\partial=(\operatorname{conj}_{b^{-1}})_*-1$. Hence
\[
 y=(\rho_1)_*\xi-(\rho_0)_*\xi+\sum_t n_t p_{b_t}((\rho_1d_t)_*z)
\]
is a cycle, since the boundary of the sum is $(\rho_0)_*\partial\xi-(\rho_1)_*\partial\xi$.

Choose $d\in\mathcal O_L\setminus\{0\}$ to clear denominators of the maps, $\alpha$ and their inverses. Let $G_d$ be the group of $L$-linear automorphisms preserving $\mathcal O_F[1/d]^D$, with inclusion $\iota:H\hookrightarrow G_d$. For a homogeneous $n$-simplex $(v_0,\ldots,v_n)$ in a group bar complex, take the oriented staircase chain for $\Delta^n\times\Delta^2$ with its three vertex rows $v$, $vb$ and $vbb'$. The terms for $0\le i\le j\le n$ are
\[
 (-1)^{i+j}[v_0,\ldots,v_i,\ v_i b,\ldots,v_j b,\
 v_jbb',\ldots,v_nbb'] .
\]
Interior faces cancel in pairs. The boundary over $\partial\Delta^n$ is the same construction on the boundary of the input. The other faces are the prisms between rows, with the long edge opposite in sign to the two consecutive edges. In inhomogeneous coordinates the row $vb$ has entries conjugated by $b^{-1}$. Thus on any cycle $z'$,
\[
 p_{bb'}z'-p_bz'-p_{b'}(\operatorname{conj}_{b^{-1}})_*z'
 \in\operatorname{im}\partial .
\]
Use this with $b=G_*$, $b'=b_t$ and $Z_t=(\rho_0d_t)_*z$. The prism boundary identity also gives
$(\rho_1)_*\xi-(\rho_0)_*\xi\equiv p_{G_*}(\rho_0)_*(\partial\xi)$ modulo boundaries. Summing over $t$ yields
\[
 \iota_*y\equiv\sum_t n_t p_{G_t}Z_t\pmod{\operatorname{im}\partial}.
\]
The standard lattice is projective over $\mathcal O_L$. Adjoin a complementary summand with trivial action to obtain a free lattice. Push $n_{t_0}^{-1}y$ into the resulting general linear group. Call this degree-$2k+1$ cycle $y_\alpha$.
\leanref{UpperTheory.\allowbreak{}V.\allowbreak{}ZK.\allowbreak{}Cl.\allowbreak{}cChain}{Development.\allowbreak{}AutomorphicArithmetic.\allowbreak{}AutomorphicVarieties.\allowbreak{}ModularityArgument.\allowbreak{}ClassCounting.\allowbreak{}LatticeBlockClasses.\allowbreak{}CChain}{8--63}
\leanref{UpperTheory.\allowbreak{}V.\allowbreak{}ZK.\allowbreak{}Cl.\allowbreak{}cJoin}{Development.\allowbreak{}AutomorphicArithmetic.\allowbreak{}AutomorphicVarieties.\allowbreak{}ModularityArgument.\allowbreak{}ClassCounting.\allowbreak{}LatticeBlockClasses.\allowbreak{}CJoin}{5--187}

Take a prime $P$ of $F$ unramified of residue degree one over $L$, away from $\ell$ and $d$. Put $K=k(P)$ and $d_K=\#K-1$. The reduction $\bar\zeta$ has order $q_0$. Choose a generator $g$ of $K^\times$ with $g^{d_K/q_0}=\bar\zeta$. For $v\in K^\times$ let $L_K(v)$ be its logarithm to base $g$ in $[0,d_K)$ and put $\operatorname{ind}(v)=L_K(v)\bmod q_0$. The function
\[
 U_K(v,w)=\frac{L_K(v)+L_K(w)-L_K(vw)}{d_K}\pmod{q_0}
\]
is a symmetric two-cocycle by cancellation of logarithms. In evaluating $U_K^{\smile k}\smile\operatorname{ind}$ on a scalar prism, adjacent insertions within each $U_K$ pair cancel by symmetry. The last insertion contributes $\operatorname{ind}(v)$ for the inserted scalar $v$. For scalar exponent $x$ each pair contributes
\[
 \sum_{i\bmod q_0}U_K(\bar\zeta^{xi},\bar\zeta^x)
 =q_0L_K(\bar\zeta^x)/d_K=x\quad\text{in }A_\nu.
\]

For a finite-dimensional $K$-space $V$, choose a nonzero vector on each line. For a tuple $(g_1,\ldots,g_{2k+1})$ and initial line $\lambda_0$, set $\lambda_i=g_i^{-1}\lambda_{i-1}$. Record the scalar $u_i$ of $g_i$ on the chosen vectors from $\lambda_i$ to $\lambda_{i-1}$. Summing the cocycle on the recorded tuples over $\lambda_0$ defines $\theta_{P,V}$. Multiplication of recorded scalars agrees with the bar faces. If the vectors on the successive lines are rescaled by $d_i$, the new scalars are $u_i'=d_{i-1}^{-1}u_i d_i$. At each cut $i$, insert $d_i$ between the old prefix and the new suffix, with sign $(-1)^i$. Faces at successive cuts cancel by $u_i d_i=d_{i-1}u_i'$. The two endpoint faces give the difference of the records, so the two cocycles differ by a coboundary. In particular their evaluations on cycles agree and are invariant under isomorphism.

For representations $\rho:G\to\GL(V)$ and $\sigma:G\to\GL(W)$ and $z'\in Z_{2k+1}(G;A_\nu)$,
\[
 \theta_{P,V\oplus W}((\rho\oplus\sigma)_*z')
 =\theta_{P,V}(\rho_*z')+\theta_{P,W}(\sigma_*z').
\]
Indeed, the projection of lines outside $W$ to $V$ is equivariant for the block diagonal action, and each line of $V$ has $(\#K)^{\dim W}\equiv1\bmod q_0$ lifts. Choose vectors on the lifts by projection to the chosen vectors in $V$. The lines in $W$ contribute once each. A trivial summand contributes zero.

For a scalar action of exponent $x$, the prism multiplication identity makes the evaluation a homomorphism in the inserted matrix $M$, since conjugation fixes the scalar cycle. It factors through the determinant. Elementary transvections generate $\mathrm{SL}(V)$ and are commutators with diagonal matrices, since $\#K>2$. Replace $M$ by $\operatorname{diag}(\det M,1,\ldots,1)$. Additivity and the scalar computation give $x^k\operatorname{ind}(\det M)$.
\leanref{UpperTheory.\allowbreak{}V.\allowbreak{}ZK.\allowbreak{}sCoc}{Development.\allowbreak{}AutomorphicArithmetic.\allowbreak{}AutomorphicVarieties.\allowbreak{}ModularityArgument.\allowbreak{}ClassCounting.\allowbreak{}FProd}{101--150}
\leanref{UpperTheory.\allowbreak{}V.\allowbreak{}ZK.\allowbreak{}sTrans}{Development.\allowbreak{}AutomorphicArithmetic.\allowbreak{}AutomorphicVarieties.\allowbreak{}ModularityArgument.\allowbreak{}ClassCounting.\allowbreak{}STrans}{411--413}
\leanref{UpperTheory.\allowbreak{}V.\allowbreak{}ZK.\allowbreak{}sThetaOf}{Development.\allowbreak{}AutomorphicArithmetic.\allowbreak{}AutomorphicVarieties.\allowbreak{}ModularityArgument.\allowbreak{}ClassCounting.\allowbreak{}SThetaOf}{192--273}

If $r:\mathcal O_F\to K$ is reduction at $P$, the maps $r\circ\sigma^j$ identify $\mathcal O_F\otimes_{\mathcal O_L}K$ with $\prod_{j=0}^{h-1}K$ by the Chinese remainder theorem at the split prime below $P$. Reduce $(\rho_0d_t,G_t)$ in block coordinates. For $t=t_0$ the distinguished block in component $j$ has exponent $\kappa^j$ and determinant $\overline{\sigma^j(\beta)}$. All other determinants are $1$. Write $\theta_P$ for evaluation after adjoining the trivial complement. The chain congruence and $n_{t_0}^{-1}$ give
\[
 \theta_P(y_\alpha)=\sum_{j=0}^{h-1}a^j\operatorname{ind}(\sigma^j(\beta))
 =P(a)\sum_{j=0}^{h-1}a^j\operatorname{ind}(\sigma^j(\alpha))
 =p_{k,c}\operatorname{ind}(\mathcal N_a(\alpha)).
\]
Here $\operatorname{ind}$ acts after reduction at $P$. Expand $\beta$. Since $a^h\equiv1\bmod q_0$, shifting $j$ by $i$ in the coefficient of $[T^i]P(T)$ gives the middle equality.
\leanref{UpperTheory.\allowbreak{}V.\allowbreak{}ZK.\allowbreak{}Cl.\allowbreak{}cBetaVal}{Development.\allowbreak{}AutomorphicArithmetic.\allowbreak{}AutomorphicVarieties.\allowbreak{}ModularityArgument.\allowbreak{}ClassCounting.\allowbreak{}LatticeBlockClasses.\allowbreak{}CBetaVal}{31--122}

For a finite set of $\alpha$'s let $\Phi=(\theta_P)_P$ over the common good primes. If $\ell^{2k}\Phi(y_\alpha)=\ell^{2k}\Phi(y_{\alpha'})$, then by \eqref{eq:B6-pe}
\[
 \mathcal N_a(\alpha)^{\ell^b}\mathcal N_a(\alpha')^{-\ell^b}
\]
is a $q_0$th power at every such prime. Its residue roots all lie in $k(P)$, and reduction is injective on $\mu_{q_0}$. Hence Frobenius fixes a root in the Kummer extension and acts trivially. Primes of $F$ of residue degree greater than one over $L$ have density zero by the bound on their norms by powers of the underlying prime norms. Chebotarev gives the global $q_0$th power relation \cite[Chap.~VII, \S13]{neukirch-1999}.
\leanref{LateralTheory.\allowbreak{}Q.\allowbreak{}SDetectCut.\allowbreak{}goal}{Development.\allowbreak{}EtaleCohomology.\allowbreak{}CyclotomicDiscriminants.\allowbreak{}SDetectCut}{101--142}

Stabilize the cycles to a common rank $r\ge4k+6$. Additivity preserves $\Phi$ under stabilization. The Bockstein for $0\to\Z\xrightarrow{q_0}\Z\to A_\nu\to0$ has image in $H_{2k}(G_r,\Z)[q_0]$. This subgroup has at most $T_2$ elements by \ref{B6.16a}, since $2(2k)+4\le4k+6$. Differences within a fibre lift to integral homology in degree $2k+1$. For such a lift $x$, let $D_\ell$ repeat a matrix in $\ell$ diagonal blocks and let $S$ stabilize to the same rank by an identity block. Set
\[
 z_0(x)=x,\qquad z_{t+1}(x)=(D_\ell)_*z_t(x)-\ell^{t+2}S_*z_t(x).
\]
At each step the value of $\Phi$ after reduction is multiplied by $\ell(1-\ell^{t+1})$. The second factor is a unit in $A_\nu$. Thus $\ell^{2k}\Phi(x\bmod q_0)$ is the evaluation of an integer multiple of $z_{2k}(x)$.

Put $d=2k+1$ and $r'=\ell^{2k}r$. By $d<r_2(L)$ in \eqref{eq:B6-room}, the stable comparison gives an $N\ge r'$ at which every real degree-$d$ cohomology class is represented on cycles by an invariant form on the product of complex symmetric spaces. Stabilization from $r'$ to $N$ is injective on real homology by \ref{B6.16a}. Borel's stable description gives primitive trace forms at the complex places and degree-one determinant logarithms \cite[\S\S10--12]{borel-1974}. These are additive under block sum and compatible with stabilization. A product of $j$ primitives pulls back under $D_\ell$ to $\ell^j$ times the product. Each degree-$d$ evaluation at $N$ is of the form
\[
 \sum_{w\mid\infty}c_w\tau_w+\sum_{j=2}^{d}a_j,
\]
where $\tau_w$ is the degree-$d$ primitive at $w$ and $a_j$ is a sum of products of $j$ primitives. The recurrence multiplies weight $j$ by $\prod_{t=0}^{2k-1}(\ell^j-\ell^{t+2})$, which vanishes for $2\le j\le d$. If the $\tau_w$ vanish on a real linear combination of recurrence images, its stabilization pairs to zero with every real cohomology class. It therefore vanishes already at $r'$. Thus the integral image has rank at most $r_2(L)$, the real rank of $K_{2k+1}(\mathcal O_L)$ in this case \cite[Prop.~12.2]{borel-1974}.
\leanref{PositivityTheory.\allowbreak{}KB.\allowbreak{}sBStarOdd\_\allowbreak{}of}{Development.\allowbreak{}AutomorphicArithmetic.\allowbreak{}AutomorphicVarieties.\allowbreak{}LiteratureResults.\allowbreak{}KTheoryStability.\allowbreak{}ReductionBoundsTorsion.\allowbreak{}OddSpanRank.\allowbreak{}SBStarOddOf}{13--78}

At $r'$ the torsion in degree $2k+1$ has order $T_1$ by \ref{B6.16a}. In each Bockstein fibre choose $y_0$. The integral recurrence images of lifts of $y_\alpha-y_0$ generate a group of rank at most $r_2(L)$ and torsion order at most $T_1$. Its image under reduction and $\Phi$ contains every $\ell^{2k}(\Phi(y_\alpha)-\Phi(y_0))$ in the fibre, and has at most $q_0^{r_2(L)}T_1$ elements. There are at most $T_2$ fibres. Separation by $\ell^{2k}\Phi$ proves \eqref{eq:B6-few}.
\leanref{UpperTheory.\allowbreak{}V.\allowbreak{}ZK.\allowbreak{}zSeq}{Development.\allowbreak{}AutomorphicArithmetic.\allowbreak{}AutomorphicVarieties.\allowbreak{}ModularityArgument.\allowbreak{}ClassCounting.\allowbreak{}HasVal}{56--60}
\leanref{UpperTheory.\allowbreak{}V.\allowbreak{}ZK.\allowbreak{}BOdd}{Development.\allowbreak{}AutomorphicArithmetic.\allowbreak{}AutomorphicVarieties.\allowbreak{}ModularityArgument.\allowbreak{}ClassCounting.\allowbreak{}LemmaKOfSteps}{11--15}
\leanref{UpperTheory.\allowbreak{}V.\allowbreak{}ZK.\allowbreak{}selmer\_\allowbreak{}upper\_\allowbreak{}bound\_\allowbreak{}of\_\allowbreak{}inputs}{Development.\allowbreak{}AutomorphicArithmetic.\allowbreak{}AutomorphicVarieties.\allowbreak{}ModularityArgument.\allowbreak{}ClassCounting.\allowbreak{}SUpperOf}{21--62}
\leanref{UpperTheory.\allowbreak{}V.\allowbreak{}ZK.\allowbreak{}sUpperOf}{Development.\allowbreak{}AutomorphicArithmetic.\allowbreak{}AutomorphicVarieties.\allowbreak{}ModularityArgument.\allowbreak{}ClassCounting.\allowbreak{}SUpperOf}{64--66}
\leanref{UpperTheory.\allowbreak{}V.\allowbreak{}ZK.\allowbreak{}bOdd}{Development.\allowbreak{}AutomorphicArithmetic.\allowbreak{}AutomorphicVarieties.\allowbreak{}ModularityArgument.\allowbreak{}ClassCounting.\allowbreak{}BOdd}{3--3}
\leanref{LateralTheory.\allowbreak{}Q.\allowbreak{}SFewOfCut.\allowbreak{}sFewOf}{Development.\allowbreak{}EtaleCohomology.\allowbreak{}CyclotomicDiscriminants.\allowbreak{}SFewOfCut}{153--211}

Substituting \eqref{eq:B6-few} into \eqref{eq:B6-finite-class-count} gives
\begin{equation}\label{eq:B6-E-bound}
 \#\mathcal E_{\nu,k}\le T_1T_2\,\ell^{(b+v_\ell(k))g(L)+c+v_\ell(k)}.
\end{equation}
Take $c$ maximal with $\mu_{\ell^c}\subset L$. The cyclotomic image is closed in $1+\ell\Z_\ell$ for odd $\ell$ and in $1+4\Z_2$ for $\ell=2$. Its logarithm is a closed additive subgroup, hence an ideal of $\Z_\ell$. Maximality of $c$ identifies this ideal as $\ell^c\Z_\ell$, so the image is $1+\ell^c\Z_\ell$. At $F=L(\mu_{\ell^\nu})$, $\nu\ge c$, choose $\sigma$ acting on $\zeta$ by $1+\ell^c$. Its order is $\ell^{\nu-c}$ by \eqref{eq:B6-lte}.

Put $H_\nu=G_{L(\mu_{\ell^\nu})}$ and $H_\infty=\bigcap_\nu H_\nu\subset G_L$. By \ref{B6.20} and its restriction description, Sha corresponds to equivariant $A$-valued characters of $H_\infty$ unramified everywhere and split above $\ell$. For such a character $\phi$, choose open normal $N\subset G_L$ with $H_\infty\cap N\subset\ker\phi$. The quotient $H_\infty N/N$ extends $\phi$ to a character $\widetilde\phi$ of $H_\infty N$, which contains $H_\nu$ for large $\nu$ by compactness. Inertia away from $\ell$ lies in $H_\infty$. For splitting, the elements in the conjugates of the finitely many decomposition groups over $\ell$ on which $\widetilde\phi$ is nonzero form a compact subset of $H_\infty N$ disjoint from $H_\infty$. For equivariance, the pairs $(\gamma,u)\in G_L\times H_\infty N$ for which
\[
 \widetilde\phi(\gamma u\gamma^{-1})\ne\chi_\ell(\gamma)^{-k}\widetilde\phi(u)
\]
form a compact set disjoint from $G_L\times H_\infty$. The conditions are locally constant, and the conjugates of each decomposition group are parametrized by the compact group $G_L$. Increasing $\nu$ excludes these bad sets from $H_\nu$ and $G_L\times H_\nu$. The finite image of $\widetilde\phi$ is killed by $\ell^\nu$ for large enough $\nu$.

For a finite set of Sha classes, choose a common such $\nu\ge\max(c,2)$ and identify $A[\ell^\nu]$ with $A_\nu$. At this layer take $\varphi(P)=\widetilde\phi(\Frob_P)$. Global reciprocity makes the sum on every principal divisor zero. Equivariance under $\sigma^{-1}$ gives $\varphi(\sigma^{-1}P)=a\varphi(P)$. Class field theory recovers the characters of $H_\nu$ from their prime values, and restriction recovers the original characters. Bound \eqref{eq:B6-E-bound} thus holds for every finite subset of Sha. Its length $s$ satisfies
\begin{equation}\label{eq:B6-Sha-count}
 \ell^s\le T_1T_2\,\ell^{(b+v_\ell(k))g(L)+c+v_\ell(k)}.
\end{equation}
\leanref{UpperTheory.\allowbreak{}V.\allowbreak{}Kum.\allowbreak{}xdualExt}{Development.\allowbreak{}AutomorphicArithmetic.\allowbreak{}AutomorphicVarieties.\allowbreak{}GaloisLibrary.\allowbreak{}KummerTheory.\allowbreak{}XdualExt}{82--117}
\leanref{UpperTheory.\allowbreak{}V.\allowbreak{}Kum.\allowbreak{}eqExt}{Development.\allowbreak{}AutomorphicArithmetic.\allowbreak{}AutomorphicVarieties.\allowbreak{}GaloisLibrary.\allowbreak{}KummerTheory.\allowbreak{}CycZConjGal}{132--177}
\leanref{PositivityTheory.\allowbreak{}QuaternionShaBounds.\allowbreak{}ShaAssumptions.\allowbreak{}good\_\allowbreak{}of}{Development.\allowbreak{}AutomorphicArithmetic.\allowbreak{}AutomorphicVarieties.\allowbreak{}ModularityArgument.\allowbreak{}ClassCounting.\allowbreak{}RingEquivalences.\allowbreak{}OrderOfOneAddPrimePowZmodPrimePow}{390--410}
\leanref{PositivityTheory.\allowbreak{}QuaternionShaBounds.\allowbreak{}ShaAssumptions.\allowbreak{}card\_\allowbreak{}le}{Development.\allowbreak{}AutomorphicArithmetic.\allowbreak{}AutomorphicVarieties.\allowbreak{}ModularityArgument.\allowbreak{}ClassCounting.\allowbreak{}RingEquivalences.\allowbreak{}OrderOfOneAddPrimePowZmodPrimePow}{412--448}
\leanref{PositivityTheory.\allowbreak{}QuaternionShaBounds.\allowbreak{}ShaAsm\_\allowbreak{}holds}{Development.\allowbreak{}AutomorphicArithmetic.\allowbreak{}AutomorphicVarieties.\allowbreak{}ModularityArgument.\allowbreak{}ClassCounting.\allowbreak{}RingEquivalences.\allowbreak{}OrderOfOneAddPrimePowZmodPrimePow}{454--486}

To arrange \eqref{eq:B6-room}, take the layer $L'/L$ of index $m=\lceil\log_\ell(4k+4)\rceil$. Write $\delta=[L':L]=\ell^m$, $n=[L:\Q]$, $n'=\delta n$, and let $c=c(L')$ be the maximal root level in $L'$. The layer is generated over $L$ by a primitive $\ell^{c(L)+m}$th root, so $c\ge c(L)+m$. The derivative of $X^{\ell^{c(L)+m}}-1$ bounds the relative different by divisibility into $(\ell^{c(L)+m})$. Use $\ell\le n+1$ and the degree of $\Q(\mu_{\ell^c})$ to obtain
\begin{equation}\label{eq:B6-promotion}
 2k+1<\min(\ell^c,r_2(L')),\qquad
 \delta\le(n+1)(4k+4),\qquad \ell^c\le2n',\qquad
 |D_{L'}|\le |D_L|^\delta(\ell^c)^{n'}.
\end{equation}
With $X=X_\ell(L)$ and $T=u^k\tau$ from \ref{B6.20}, divisibility of $T^{\ell^m}-1$ by $T-1$ gives a surjection $X/(T^{\ell^m}-1)X\to X/(T-1)X$. Thus $s_\ell(L,k)\le s_\ell(L',k)$.
\leanref{LateralTheory.\allowbreak{}Q.\allowbreak{}SLayerCut.\allowbreak{}goal\_\allowbreak{}of}{Development.\allowbreak{}EtaleCohomology.\allowbreak{}CyclotomicDiscriminants.\allowbreak{}SLayerCut}{55--107}

Put $v=2kn'$ and $H'=\log(2|D_{L'}|)$. By Minkowski's bound each class has an integral ideal of norm at most $2^{n'/2}|D_{L'}|^{1/2}$ \cite[Ch.~I, \S5, Ex.~3]{neukirch-1999}. Indeed the Minkowski factor in degree $n'$ with $r_2$ complex pairs is $(4/\pi)^{r_2}n'!/(n')^{n'}\le2^{n'/2}$. The number of ideals of norm at most $B$ is at most $B^2\zeta_{L'}(2)$, and $\zeta_{L'}(2)\le\zeta(2)^{n'}<2^{n'}$ by the Euler factors. Hence $h_{L'}\le4^{n'}|D_{L'}|$ and
\begin{equation}\label{eq:B6-class-gens-bound}
 g(L')\le2n'+4n'H'+2.
\end{equation}
With $c,\log\ell,v_\ell(k)\le v$ and $b\le v+2v^2$, use \eqref{eq:B6-class-gens-bound} and \eqref{eq:B6-Sha-count}. Bounding both torsion orders by \eqref{eq:B6-homology-torsion} at rank $4k+6$ gives
\[
 s_\ell(L',k)\log\ell\le2\exp(v^{24}H')+v^9H'\le\exp(v^{26}H').
\]
By \eqref{eq:B6-promotion}, $v\le(2kn)^5$ and $H'\le(2kn)^{11}\log(2|D_L|)$. This gives \eqref{eq:B6-small-Sha} with $C=141$.
\end{proof}

\subsection*{Proof of \ref{B6.16a}}\label{prf:B6.16a}\phantomsection
\begin{proof}
The Bass stable rank of $\mathcal O_L$ is at most two. For a unimodular row $(a_1,\ldots,a_t,b)$ with $t\ge2$, if all the $a_j$ vanish, add the unit $b$ to the first entry. Otherwise take $a_j\ne0$ and $h\ne j$. Choose $u$ by the Chinese remainder theorem to be $1$ at primes dividing both $a_j$ and $a_h$, and $0$ at the other primes dividing $a_j$. A prime containing the first $t$ entries after replacing $a_h$ by $a_h+ub$ would divide $a_j$, hence $a_h$ and $b$, contrary to unimodularity.

For $i\ge1$, van der Kallen's integral constant-coefficient theorem gives an isomorphism from rank $r$ to rank $r+1$ when $r\ge1+\max\{2i,i+\operatorname{sr}(\mathcal O_L)\}$ \cite[Theorem~4.11(i) and Corollary~4.12(ii)]{kallen-1980}. Take the full linear group in his notation. The preceding stable-rank calculation gives $\operatorname{sr}(\mathcal O_L)\le2$, and $2i+4\ge1+\max\{2i,i+2\}$. In degree zero stabilization is the identity on $\Z$. Compose the consecutive maps and use \ref{B6.16c} for finiteness.
\end{proof}

\subsection*{Proof of \ref{B6.16b}}\label{prf:B6.16b}\phantomsection
\begin{proof}
For $N=0$ take tuples of one point, with $0^0=1$. For $N\ge1$ let $P$ be the space of tuples $h=(h_w)$ of Hermitian forms at the complex places cut out by $h[v]=\sum_w v_w^*h_wv_w\ge1$ for $0\ne v\in\mathcal O_L^N$. The standard form scaled up lies in $P$ by discreteness. In an integral basis the real Gram matrix $B$ of $h\in P$ is positive semidefinite by density of rational directions. Write $B+\delta\,1=S^tS$, $\delta>0$. Minkowski's theorem for a box gives $0\ne z\in\Z^d$ with $\|Sz\|_\infty\le2|\det S|^{1/d}$ \cite[Ch.~I, Thm.~4.4]{neukirch-1999}. Here the lattice is $S\Z^d$, of covolume $|\det S|$, and the centered box of radius $2|\det S|^{1/d}$ has volume $4^d|\det S|$. Thus $1\le4d|\det S|^{2/d}$ and $(4d)^d\det B\ge1$ on letting $\delta\to0$. All $h_w$ are positive definite. The integral basis has real covolume $2^{-n/2}\sqrt D$ in the complex embeddings. Each Hermitian determinant occurs twice in the real determinant, so
\begin{equation}\label{eq:B6-real-gram}
 2^d\det B=D^N\Bigl(\prod_w\det h_w\Bigr)^2,\qquad
 \prod_w\det h_w\ge I^{-1}.
\end{equation}
The inequality uses $D\ge1$ and $d\ge2$.

Let $X_p$ consist of ordered $(p+1)$-tuples, with repetitions, of vertices of $P$ minimizing a common linear functional positive on every nonzero positive semidefinite tuple. Such a functional has compact sublevels on $P$ by its positive minimum on the unit section of the semidefinite cone. Its minimizing face has a vertex. Near each $h\in P$ the forms have a uniform positive lower eigenvalue, so locally finitely many lattice inequalities can become active. At a vertex the active evaluations span the dual space, since an annihilated direction would allow perturbations of both signs.

For a finite cycle choose a witnessing functional for each tuple and let $K$ be their convex hull. Fix $h_0\in P$. A vertex minimizing a functional of $K$ lies in a witness sublevel $f(h)\le f(h_0)$, by the convex combination defining that functional. These finitely many compact sublevels contain finitely many vertices by local finiteness. Their minimizing cones intersect $K$ in a finite closed convex cover. The ordered nerve of this cover lies in $X_\bullet$ and contains the cycle.

For a finite closed convex cover of a nonempty compact convex $K$, the ordered nerve has exact augmented chains. Induct on the number of realizable subsets of indices, namely subsets whose cover members meet at a point of $K$. If all members meeting $K$ share a point, prepending one fixed index gives a contraction, also for repeated vertices. Otherwise choose a maximal realizable subset and a member meeting $K$ but disjoint from that subset's intersection in $K$. Strictly separate the two compact convex sets by a hyperplane. The two closed halves of $K$ and their common section are nonempty. Each loses at least one realizable subset. The ordered nerves of the halves have union the full nerve and intersection the nerve of the section. The segment joining the witnesses in the two halves lies in every member of the tuple. The short exact sequence of augmented chain complexes for this union and intersection, with the first map given by the two inclusions with opposite signs, proves exactness by induction. Thus the cycle in $X_\bullet$ bounds. A minimizing vertex gives surjectivity of the augmentation, and change of basis preserves the tuple condition.

At the average of a tuple the active vectors are its common minimal vectors. If they spanned a proper $L$-subspace, choose a nonzero $L$-linear functional $\lambda$ vanishing there. Subtracting a small multiple of the semidefinite tuple $(|\lambda_w|^2)_w$ preserves the active equalities and stays in $P$ by local finiteness. It decreases the common functional, a contradiction. Take $N$ independent common minima as columns of $M$. Put $H_w=M_w^*h_wM_w$ at a vertex and $i=|N_{L/\Q}(\det M)|$. The diagonal entries of $H_w$ are at most one. Hadamard's inequality follows by Gram--Schmidt, since the volume of the columns is the product of their successive orthogonal projection lengths, bounded by the product of their lengths. It gives $\det H_w\le1$ and, by \eqref{eq:B6-real-gram},
\[
 i/I\le\prod_w\det H_w=i\prod_w\det h_w\le1.
\]
Thus $1\le i\le I$ and every $\det H_w\ge I^{-1}$. The matrix $C=iM^{-1}$ is integral because $\mathcal O_L^N/M\mathcal O_L^N$ has order $i$. For fixed $i$, equal residues of $C,C'$ modulo $i$ make $M'M^{-1}=1+M'(C-C')/i$ and its inverse integral. This gives at most $\sum_{i=1}^I i^{nN^2}\le I^{1+dN}$ left $G$-classes of $M$.

Fix a representative $M$. The trace bound $N$ gives least eigenvalue at least $I^{-1}N^{1-N}$ for each $H_w$. If $v$ is minimal at such a vertex, $\|M_w^{-1}v_w\|_2^2\le IN^{N-1}$. Each conjugate coordinate of $Cv$ therefore has absolute value at most $R$. Distinct algebraic integers are separated by at least one in the maximum embedding norm, by the field norm. Packing radius-$1/2$ balls bounds the count per coordinate by $(2R+1)^n$. As $v\mapsto Cv$ is injective, there are at most $(2R+1)^d$ possible minimal vectors. At most $dN$ independent active equations determine a vertex. Allowing an empty slot for padding gives at most $\bigl((2R+1)^d+1\bigr)^{dN}$ lists. Count the $p+1$ entries to obtain $U_p$. With $B_0=2Dd$, one has $I\le B_0^{2d}$ and $R\le B_0^{5d}$. The resulting exponent is at most
\[
 2d(1+dN)+\bigl((5d+2)d+1\bigr)dN(p+1)\le12(p+1)d^3N.
\]

A stabilizer fixes the first form and is finite, since its columns are lattice vectors of bounded length. If its kernel modulo $3$ is nontrivial, take an element of prime order $q$ in it, of the form $1+3^aY$ with $a\ge1$ maximal. For $q\ne3$ its $q$th power differs from $1$ by $q3^aY$ modulo $3^{a+1}$. For $q=3$ its cube differs from $1$ by $3^{a+1}Y$ modulo $3^{a+2}$. Both contradict its order. Reduction embeds the stabilizer into the set of $3^{nN^2}$ residue matrices.
\end{proof}

\subsection*{Proof of \ref{B6.16c}}\label{prf:B6.16c}\phantomsection
\begin{proof}
Put $C_p=\Z[X_p]$, $K_0=\Z$ and $K_{p+1}=\ker(C_p\twoheadrightarrow K_p)$, using exactness of \ref{B6.16b}. For $p<m$ and $j=m-p$, Shapiro's lemma identifies $H_j(G,C_p)$ with the sum of $H_j(S,\Z)$ over the orbit stabilizers $S$. Transfer kills each summand by $|S|$. Its bar cycles lie in a free group of rank $|S|^j$, so the summand has order at most $|S|^{|S|^j}$. Since $U_p\le A$, this gives $\#H_j(G,C_p)\le s^{As^j}$. In the sequence
\[
 H_j(G,C_p)\longrightarrow H_j(G,K_p)\longrightarrow H_{j-1}(G,K_{p+1})
\]
the torsion in the middle has at most $\#H_j(G,C_p)$ elements over each torsion value on the right.

The group $(K_m)_G$ is the cokernel of $(C_{m+1})_G\to(C_m)_G$. The target is free on at most $A$ orbits and each column has $\ell^1$-norm at most $m+2$ from the alternating deletion sum. If the matrix has rank $t$, its torsion order divides each nonzero $t$-rowed minor by Smith normal form. Hadamard's inequality bounds that minor by $(m+2)^t$. Multiplication along the $m$ exact sequences gives $(m+2)^A s^{A\sum_{j=1}^m s^j}\le(m+2)^A s^{mAs^m}$.

Put $u=kn$, $H=\log(2D)$ and $a=12(m+1)d^3N\log(2Dd)$. By \eqref{eq:B6-orbit-bound}, $A\le e^a$. Using $\log(m+2)\le m+1$ and $\log s\le s$ gives
\[
 \log T\le A(2m+1)s^{m+1}
 \le\exp\!\bigl(a+2(m+1)nN^2+2m\bigr).
\]
Here $d,N\le10u$, $m+1\le4u$ and $\log(2Dd)\le16uH$. Hence $12d^3N\log(2Dd)+2nN^2\le1\,920\,300u^5H$ and $2m\le12u^6H$. The exponent is at most $7\,681\,212u^6H\le(2u)^{24}H$.
\end{proof}

\subsection*{Proof of \ref{B6.18}}\label{prf:B6.18}\phantomsection
\begin{proof}
By \ref{B6.17} and class field theory, the unramified mod-$\ell$ character groups of the finite layers have order at most $\ell^r$. Their restriction images increase and exhaust the infinite-layer group by the finite-quotient extension argument of \ref{B6.16}, with inertia at $\ell$ in place of decomposition in the compactness argument. Compact Nakayama gives at most $r$ generators for $\mathrm{XX}$ and its quotient $\mathrm{XP}$. Take $R_\ell(A)=\max(1,r)$.

In the tower of $K_M(i)$, allow ramification above $q,d$. Tameness gives at most one mod-$2$ inertia dimension per place. In the tower over $\Q(i)$, an odd $p\equiv3\bmod4$ has Frobenius $p^2\in1+4\Z_2$ and decomposition index $2^{v_2(p+1)-1}$. For $q,d$ these indices are $2,4$, so after base change there are at most $6n_M$ places in all.

At one quadratic adjunction take index-two Galois subgroups $U\subset V$ and choose $\sigma\in V\setminus U$. Conjugation by $\sigma^{-1}$ induces an involution $\iota$ on characters of $U$, since $\sigma^2\in U$. Over $\mathbf F_2$, $(\iota-1)^2=0$. An invariant character $\phi$ extends exactly when $\phi(\sigma^2)=0$. The obstruction has at most two values and any extending character has two choices on $\sigma$, both preserving ramification restrictions because the excluded inertia groups lie in $U$. Thus the invariant dimension is at most the character dimension on $V$. Since $\operatorname{im}(\iota-1)\subseteq\ker(\iota-1)$, the total dimension on $U$ is at most twice the invariant dimension. Two adjunctions give $4(g+6n_M)$ for the unramified subspace and hence for both duals.
\end{proof}

\subsection*{Proof of \ref{B6.19}}\label{prf:B6.19}\phantomsection
\begin{proof}
Restriction preserves Sha. If $E'=E$, restriction is the identity. For a quadratic extension inflation--restriction identifies the kernel on $H^1$ with $\ker(1+\sigma:A\to A)/(\sigma-1)A$, where $A=((\Q_2/\Z_2)(-k))^{G_{E'}}$ and $\sigma$ generates the quotient. This subquotient of a locally cyclic group is killed by $2$ and has length at most one.

Adjoin $\mu_{2^{a+2}}$ to $K_{M,q,d}$ to get $L_m$ with $m\le a+2$. Under a common embedding $E\subset L_m$ has index dividing two, by the index of the real cyclotomic subfield. Apply \eqref{eq:B6-quadratic-loss}. The three quadratic adjunctions give the degree bound, and their conductors $4,q,d$ give the embedding.
\end{proof}

\subsection*{Proof of \ref{B6.20}}\label{prf:B6.20}\phantomsection
\begin{proof}
Write $\Gamma_m=\Gal(L_\infty/L_m)$ and $A=(\Q_\ell/\Z_\ell)(-k)$. The generator minus the identity acts on $A$ by the surjective scalar $u^{-k\ell^m}-1$. Hence $H^1(\Gamma_m,A)=0$. Also $H^2(\Gamma_m,A)=0$ by cohomological dimension one. Restriction identifies global classes with equivariant characters of $G_{L_\infty}$. Inertia away from $\ell$ lies in this subgroup. Above $\ell$, a local class whose character vanishes on the kernel in a decomposition group descends to an open subgroup of $\Gamma_m$, where $H^1$ again vanishes by the nonzero scalar. Thus Sha identifies with $(\mathcal D_\ell(L_0)(-k))^{\Gamma_m}$. Duality preserves length and turns invariants into coinvariants for $(u^k\tau)^{\ell^m}$.
\end{proof}

\subsection*{Proof of \ref{B6.21}}\label{prf:B6.21}\phantomsection
\begin{proof}
Lift $\tau^{\ell^j}$ to a matrix $S$ on $\Z_\ell^R$ along a surjection to $X$. For $t=u^{k\ell^j}$, the cokernel of $tS-1$ surjects onto the one defining $w_{\tau,u}(k,j)$. Its length is the valuation of the determinant by Smith normal form, hence $\sum_i v_\ell(tz_i-1)$ for the eigenvalues $z_i$ of $S$ with multiplicity.

For $\ell$-adic units $A,B$ and algebraic $z$, $|A-B|_\ell\le\max(|Az-1|_\ell,|Bz-1|_\ell)$. Use the ultrametric inequality when $|z|_\ell\ge1$ and the fact that both right-hand norms are $1$ when $|z|_\ell<1$. Formula \eqref{eq:B6-lte} with exponent $\ell^j|a-b|$ gives
\[
 |(u^{\ell^j})^a-(u^{\ell^j})^b|_\ell
 =\ell^{-(j+c)}|a-b|_\ell\qquad(a,b\in\N).
\]
For $a=b$ both sides are zero. The balls in $\Z_\ell$ centred at $a\in\N$ with radii $r_a=\ell^{j+c}|u^{a\ell^j}z_i-1|_\ell$ are nested because $|a-b|_\ell\le\max(r_a,r_b)$. Compactness gives $s_i$ in their intersection. Take $a=k$ to obtain \eqref{eq:B6-centres-fixed}.

Put $g=u^k\tau$ and $h=\lfloor\log_\ell R\rfloor+1$. The case $R=0$ has $X=0$. Otherwise $g-1$ is nilpotent modulo $\ell$ of index at most $R$. The binomial theorem gives $(g^{\ell^h}-1)X\subseteq\ell X$ and $(g^{\ell^{h+1}}-1)X\subseteq\ell^2X$. If $(w-1)X\subseteq\ell^2 X$, the factor $1+\cdots+w^{\ell-1}$ in $w^\ell-1$ has image $\ell X$ by Nakayama, since it agrees with $\ell$ modulo $\ell^2X$. Hence $(w-1)X/(w^\ell-1)X$ is a quotient of $X/\ell X$ and has length at most $R$. Iterate from $j_0=h+2$ and use divisibility for $m\le j_0$ to get
\begin{equation}\label{eq:B6-growth}
 w_{\tau,u}(k,m)\le w_{\tau,u}(k,j_0)+R(m-j_0)_+.
\end{equation}
Apply \eqref{eq:B6-centres-fixed} at $j_0$. Under the distance condition each valuation is at most $b$.
\end{proof}

\subsection*{Proof of \ref{B6.24}}\label{prf:B6.24}\phantomsection
\begin{proof}
Set $U=\log X$, $B=C_t/\eta$ and $T=\{\log k/C_t:k\in\mathcal E,\ k\ge1\}$. Seek $U/(2^{R+2}B^R)\le y\le U/2$ with
\begin{equation}\label{eq:B6-log-window}
 a\in T,\quad y\le u\le2y,\quad a\le u\quad\Longrightarrow\quad Ba\le u.
\end{equation}
If $T$ misses $(U/(2B),U]$, take $y=U/2$. Otherwise remove a point in this interval and apply induction on $R$ with upper endpoint $U/(2B)$. The new bound for $2y$ excludes the removed point and the new lower bound is $U/(2^{R+2}B^R)$. Take $Y=e^y$ and $u=\log x$. For $k=0$ the conclusion is automatic.
\end{proof}

\subsection*{Proof of \ref{B6.25}}\label{prf:B6.25}\phantomsection
\begin{proof}
Since $k\le Z$, outside $\mathcal E$ use \eqref{eq:B6-first-bound}. Inside $\mathcal E$, \eqref{eq:B6-good} gives $k\le x^\eta$, so the exponent in \eqref{eq:B6-small-H2} at $L_{j_0}$ is at most $(2n_0)^{C_K}c_0x^{3\varepsilon/4}$. Divide by $\log\ell>1/2$ and use \eqref{eq:B6-second-bound}. Together the cases give
\begin{equation}\label{eq:B6-point-before-logs}
 s_\ell(L_m,k)\le R(m+j_0+c+\lfloor\log_\ell Z\rfloor)
 +2\exp\bigl((2n_0)^{C_K}c_0x^{3\varepsilon/4}\bigr).
\end{equation}
Since $\log X\le J\log x$, we have $\lfloor\log_\ell Z\rfloor\le2C_tJ\log x$. Take $C=R(A_0+j_0+c_1+2C_tJ)+1+(2n_0)^{C_K}c_0$.

For $x\ge2$ and arbitrary $C_t,C\in\R$, the product $x^{C_t}P_{C,\varepsilon}(x)$ is at most
\[
 3\exp\!\left((|C_t|+|C|)\log x+|C|+|C|x^{3\varepsilon/4}\right).
\]
The logarithm of this bound is $o(x^{4\varepsilon/5})$.
\end{proof}

\subsection*{Proof of \ref{B6.26}}\label{prf:B6.26}\phantomsection
\begin{proof}
At $2$ take $L_0=K_{M,q,d}$ and $j=j_0(2,R_0(M))$. For constants $n_0,c_0$ depending on $M,j$,
\begin{equation}\label{eq:B6-layer-disc}
 [L_j:\Q]\le n_0,\qquad
 \log(2|D_{L_j}|)\le c_0w\quad(w\ge1,\ \log(qd)\le w).
\end{equation}
Take $n_0=8n_M2^j$ by \ref{B6.19}. Write $L_j=L_0(\mu_{2^r})$ with $r$ minimal. Then $2^r\le2n_0$ and $L_j$ embeds into a cyclotomic field of conductor dividing $2^r(4N_Mqd)$. The derivative of $X^h-1$ bounds the root discriminant of a field generated by an $h$th root by $h$. The tower formula passes this bound to subfields. Thus $\log(2|D_{L_j}|)\le\log2+n_0\log(8n_0N_Mqd)$, which proves \eqref{eq:B6-layer-disc}.

Take $\mathcal E_2$ at $Z=\lfloor X^{C_t}\rfloor$ by \ref{B6.18}, \ref{B6.22} and \ref{B6.23}. Here $\log(qd)\le x^{\varepsilon/4}$ and $c\le2^{c-1}\le[L_0:\Q]\le8n_M$. Use \eqref{eq:B6-layer-disc} with $w=x^{\varepsilon/4}$. By \ref{B6.19} and \eqref{eq:B6-admitted}, each $E$ admits a layer with $m\le a(x)+2\le(2|A_1|+2)(1+\log x)$ and $s_2(E,k)\le s_2(L_m,k)+1$. Apply \ref{B6.25} with these uniform bounds.

At an odd prime choose one realization and take $\mathcal E_\ell$ at the same cutoff. Apply \ref{B6.25} with $c_1=c$, $n_0=[L_{j_\ell}:\Q]$, $c_0=\log(2|D_{L_{j_\ell}}|)$ and $A_0=j_\ell$. Layer isomorphism transports the bound with this set to every realization. In both cases there are at most $x^{C_t}$ terms and each is at most $P_{C,\varepsilon}(x)$, with the quadratic cost included at $2$. Take $x_0$ from \eqref{eq:B6-final-absorption}.
\end{proof}

\section*{Proofs for Chapter~\ref{ch:b6-classwise}: Summing the plus component by classes}

\subsection*{Proof of \ref{B6.28a}}\label{prf:B6.28a}\phantomsection
\begin{proof}
The host quotient embeds in the product of the dyadic quotient by $\{\pm1\}$ with the image of $\chi$ and the two square-class groups. The dyadic quotient has order $2^a$. Restriction to the $r$-factor for odd $r\mid f$ has image in the factor from $\chi$. For $r=q,d$ it has image in a square-class group. The dyadic order divides $2^{c_2(\theta)}$. If the global order is odd, the factors at $2,q,d$ are trivial and the remaining local conductors divide the corresponding parts of $f$. The modulus defining $f_0$ divides $4o$, and $2^{n_0}\le\ell^{f_0}-1$.

Largeness puts $\mathcal B_\chi$ below $q,d$. The factors at $2,q,d$ survive in $\theta_P$ with $2$-power order, so odd order of $\theta_P$ forces them to be trivial. Outside $\mathcal B_\chi$, the prime $\ell$ is prime to $2of$. If $\ell-1\mid2k$, the bounds $5\le\ell\le2k+1<q,d$ also put $\ell$ outside the prime divisors of $N$.
\end{proof}

\subsection*{Proof of \ref{B6.30aa}}\label{prf:B6.30aa}\phantomsection
\begin{proof}
Choose a simple \(H\)-submodule. Its \(G\)-translates span the representation, so the restriction is semisimple and \(G\) permutes its isotypic summands transitively. All summands have the same dimension and multiplicity. Since \(n\) is prime, a reducible restriction is scalar or is a sum of \(n\) distinct lines permuted transitively by \(G\). The element \(u\) excludes the scalar case. In the second case these lines are the eigenlines of \(u\). Each space in the ordinary flag is a sum of a subset of them. An element of \(D\) permutes the lines and preserves each member of this complete flag, so it fixes every line.

Let \(J\) be the kernel of the permutation action. The characters on the \(n\) lines over \(J\) have integral unit values. For instance, intersect a \(G\)-stable lattice with each line, after extending the coefficient field to define the lines. The sum of these intersections is a \(G\)-stable lattice. Reducing its characteristic polynomials on \(J\) identifies the multiset of the reduced line characters with the restrictions of the residual characters of \(G\), by uniqueness of semisimplification. Each of the latter is invariant under conjugation by \(G\). The reduced line characters are conjugate to one another by transitivity, hence all equal on \(J\). Since \(D\subseteq J\), the residual characteristic polynomial at \(s\) would have only one distinct root. Its roots include \(\bar\lambda_a(s)\) and \(\bar\lambda_b(s)\), a contradiction.

For the last assertion take the normal core \(U_0\) of \(U\). Apply the first assertion to \(H\cap U_0\) and \(u^{[G:U_0]}\). Irreducibility on \(H\cap U_0\) implies irreducibility on \(H\cap U\).
\end{proof}

\subsection*{Proof of \ref{B6.30ac}}\label{prf:B6.30ac}\phantomsection
\begin{proof}
Write \(H=G_K\) and \(\Psi=\varepsilon_\ell^k\vartheta\). Fix a place above \(\ell\). Its rational inertia belongs to \(H\), because \(\ell\nmid D_K\), and maps onto \(\Z_\ell^\times\) by \(\varepsilon_\ell\). The local test of \eqref{eq:p6c5-link-input-tests} gives \(x\in H\) in the decomposition group with \(\bar\Psi(x)^2\ne1\).

In Type I take the normal core \(U\) of the given open subgroup in \(H\). Choose \(\sigma\) in inertia of cyclotomic value \(1+\ell\), and take \(u=\sigma^m\) with \(m=[H:U]\operatorname{ord}\vartheta\). Then \(u\in U\), and its three eigenvalues \((1+\ell)^{km},1,(1+\ell)^{-km}\) are distinct by the ordinary inertia condition. The restriction to \(U\) is semisimple: the translates under \(H\) of a simple \(U\)-submodule span the irreducible \(H\)-module and are again simple \(U\)-submodules. Its isotypic summands are permuted transitively. In dimension three a reducible restriction thus has either scalar \(U\)-action or three lines permuted transitively. Scalar action is excluded by \(u\). The lines in the second case are precisely the eigenlines of \(u\). The ordinary flag is made up of successive sums of these lines, and \(x\) preserves it. Its permutation therefore fixes each line. A transitive subgroup of the permutations of three letters contains a three-cycle; choose \(g\in H\) inducing one. Then \(gx^j\) has trace zero for every \(j\ge0\). Reducing the trace in the marked lattice gives
\[
 A B^j+1+(A B^j)^{-1}=0,\qquad
 A=\bar\Psi(g),\quad B=\bar\Psi(x).
\]
For \(j=0,1,2\) the three distinct elements \(A,AB,AB^2\) would all be roots of \(X^2+X+1\). This contradiction proves the first assertion.

In Type II put \(N=H\cap\ker\varepsilon_\ell\). On the chosen inertia group \(I\) the corner has trace \(\Psi+\Psi^{-1}\): subtract the trivial middle character from the three prescribed inertia eigenvalues. The cyclotomic surjection gives \(H=IN\). Choose \(u\in I\) with \(\Psi(u)\) of infinite order, using a power of the same \(\sigma\). If \(N\) acts by scalars then the projective image of \(H\) is abelian. It centralizes the projective image of \(u\), whose eigenvalue ratio has infinite order. That centralizer preserves each eigenline, contradicting irreducibility on \(H\). Otherwise a reducible restriction to \(N\) has either a unique invariant line or two distinct isotypic lines, by the same socle argument. A unique line is \(H\)-invariant. In the remaining case an element of \(I\) must interchange the two lines, since \(H=IN\). If \(s\) does so, all \(su^{2j}\) interchange them and have trace zero. The inertia trace formula would give \(\Psi(s)^2\Psi(u)^{4j}=-1\) for every \(j\), contrary to the choice of \(u\).
\end{proof}

\subsection*{Proof of \ref{B6.30ab}}\label{prf:B6.30ab}\phantomsection
\begin{proof}
In Type I, \ref{B6.30ac} verifies the open-subgroup hypothesis of \ref{B6.52}(2). It gives a regular algebraic conjugate self-dual cuspidal attachment \(\pi\) over \(K\) for \(r|_{G_K}\varepsilon_\ell^{-1}\). We check irreducibility on \(N=G_{K(\mu_{\ell^\infty})}\). As in \ref{B6.30ac}, restriction to the normal subgroup \(N\) is semisimple. Three permuted lines would give an open line stabilizer in \(G_K\), contrary to that lemma. Here the stabilizer is open because the eigenline action is continuous and its orbit finite. If \(N\) acted by scalars, the projective image of \(G_K\) would be abelian. It would preserve the finite set of eigenlines of an inertia element with three distinct eigenvalues, again contradicting open-subgroup irreducibility. These exhaust the reducible possibilities in prime dimension. A scalar twist preserves this irreducibility.

In Type II the attachment of \(W\varepsilon_\ell^{-k}\) is cuspidal on \(\GL_2/\Q\) by \ref{B6.52}(3), of weight \(2k+1\ge3\). Its quadratic base change to \(K\) is cuspidal by \ref{B6.06}: the exceptional case gives a self-twist by \(\eta_K\), whose intertwiner would commute with \(W|_{G_K}\) and hence be scalar, whereas on the other coset it would intertwine a matrix with its negative. The attachment after base change has the restricted good Frobenius polynomials. Its weights remain distinct, and the exterior-product pairing of \(W\), with \(\det W=\eta_K\), gives essential conjugate self-duality. More explicitly, for complex conjugation \(c\) and \(J=\left(\begin{smallmatrix}0&1\\-1&0\end{smallmatrix}\right)\), the matrix \(P=JW(c)^{-1}\) satisfies \(W(h)^tP W(chc^{-1})=P\) on \(G_K\). It is symmetric since \(W(c)^2=1\) and \(\det W(c)=-1\). Twisting adds the corresponding scalar to the multiplier. Irreducibility on \(N\) is \ref{B6.30ac}.

In each case Chebotarev and Brauer--Nesbitt identify the attached Galois representation with the indicated scalar twist, using the good Frobenius polynomials. On \(G_K\) the adjoint is conjugation on the full matrix algebra, so scalar twists cancel. At \(c\) its action is \(Z\mapsto-P^{-1}Z^tP\) for the polarising form; another form with the same multiplier differs by a scalar by irreducibility. This identifies the polarised adjoint in the assertion with the one in \cite[Thm.~6.1]{thorne-2022}, as in the proof of \ref{B6.53}. The field \(K\) is CM with \(K^+=\Q\), the rank is at least two, the attachment is regular algebraic cuspidal and polarised, and its cyclotomic restriction has just been proved irreducible. That theorem gives the stated vanishing.
\end{proof}

\subsection*{Proof of \ref{B6.30ad}}\label{prf:B6.30ad}\phantomsection
\begin{proof}
Let \(v_0\) be the distinguished split place. Identify its group with \(\GL_3(H)\), where \(H=F_{v_0}\). Write \(k_0\) for its residue field, \(q_0=|k_0|\), \(K^0=\GL_3(\mathcal O_H)\), \(K^1=\ker(K^0\to\GL_3(k_0))\), and \(U\) for the upper unitriangular group. The level at \(v_0\) is \(K^1\). Fix a nontrivial additive character \(\psi_0\) of \(k_0\). For \(a,b\in k_0^\times\) set
\[
 \chi_{a,b}(u)=\psi_0(au_{12}+bu_{23})\qquad(u\in U(k_0)),
\]
using indices \(1,2,3\) in this formula. On \(K^1\)-invariants take
\[
 e=\frac{1}{q_0^3}\sum_{u\in U(k_0)}
                \left(\sum_{a,b\in k_0^\times}\chi_{a,b}(u)^{-1}\right)u .
\]
Each inner coefficient is an integer. It is the product of two sums, each equal to \(q_0-1\) or \(-1\) according as its argument is zero or nonzero. Also \(v_0\nmid\ell\), so \(q_0\) is a unit in \(O\). The lifts of \(u\) in \(K^0\) give the same operator on \(K^1\)-invariants. Character orthogonality makes \(e\) the sum of mutually orthogonal projectors. The Hecke operators and weight characters used to define \(B_0/J\) act away from \(v_0\). Thus \(e\) acts on the integral residual summand, commutes with this algebra and specializes under \ref{B6.48d}.
\leanref{UpperTheory.\allowbreak{}V.\allowbreak{}Z.\allowbreak{}Desc.\allowbreak{}V0.\allowbreak{}eOp}{Development.\allowbreak{}AutomorphicArithmetic.\allowbreak{}AutomorphicVarieties.\allowbreak{}ModularityArgument.\allowbreak{}CuspidalDescent.\allowbreak{}GLDenseRanges.\allowbreak{}InertiaPartMaps.\allowbreak{}MapInj}{86--169}

We compute the rank of \(e\) on \(\tau^{K^1}\) for a generic irreducible smooth complex representation \(\tau\) of \(\GL_3(H)\) with \(\tau^{K^1}\ne0\). Extend \(\psi_0\) to an additive character of \(H\) of conductor \(\mathfrak m_H\), and lift \(a,b\) to units. This gives a nondegenerate character \(\chi\) of \(U(H)\). Take a Whittaker functional \(l\) for \(\chi\). Conjugation by a diagonal matrix takes any choice of nondegenerate Whittaker character to this one. The functions \(W_z(g)=l(\tau(g)z)\) on the mirabolic subgroup, consisting of matrices with last row \((0,0,1)\), separate the vectors \(z\) by \cite[Thm.~4.11]{bernstein-zelevinsky-1977}.

First choose \(z\in\tau^{K^1}\) and \(g\) with \(W_z(g)\ne0\). By Iwasawa decomposition and normality of \(K^1\) in \(K^0\), we can absorb the compact factor into \(z\) and the upper unipotent factor into a nonzero scalar. We have \(l(\tau(d)z)\ne0\) with \(d=\operatorname{diag}(d_1,d_2,d_3)\). Comparing the actions of \(1+sE_{i,i+1}\) and its \(d\)-conjugate for \(s\in\mathfrak m_H\) gives \(v_H(d_i)\ge v_H(d_{i+1})\). Average \(\tau(d)z\) over \(K^1\). Its value under \(l\) is unchanged. Indeed write each element of \(K^1\) as upper unipotent times diagonal times lower unipotent, with all factors in \(K^1\). The character on the first factor is trivial. The diagonal commutes with \(d\), and \(d^{-1}\) conjugates the lower factor into \(K^1\) by the valuation inequalities. Now average with character \(\chi^{-1}\) over \(U(\mathcal O_H)\). This preserves the value under \(l\) and produces a nonzero \(\chi_{a,b}\)-vector in \(\tau^{K^1}\).

On this character space \(l\) is injective. To see this, suppose \(W_z(p)\ne0\) for a mirabolic matrix \(p\). Work up from the last row. If rows below row \(i\) are upper unitriangular, compare right translation by \(1+sE_{j,i+1}\), \(s\in\mathcal O_H,\ j\le i\), with left translation by its \(p\)-conjugate. The conjugate is upper unipotent. Its only possibly nonzero simple-root entry is in position \(i,i+1\), with value \(sp_{ij}\). The character identity gives
\[
       \psi(t_i s(p_{ij}-\delta_{ij}))=1
              \quad(s\in\mathcal O_H),\qquad (t_1,t_2)=(a,b),
\]
where \(\psi\) denotes the lifted additive character and \(t_i\) the chosen unit lift. Hence \(p_{ij}-\delta_{ij}\in\mathfrak m_H\) for \(j\le i\). Right multiplication by an element of \(K^1\) makes this row upper unitriangular while preserving the rows below: on the first \(i\) columns use a matrix congruent to the identity to carry \((p_{i1},\ldots,p_{ii})\) to \((0,\ldots,0,1)\). Repetition reduces \(p\) to an upper unipotent without changing \(W_z(p)\). Therefore \(l(z)=0\) would give \(W_z=0\) on the mirabolic. This proves injectivity. Each of the \((q_0-1)^2\) characters contributes exactly one dimension.
\leanref{UpperTheory.\allowbreak{}V.\allowbreak{}Z.\allowbreak{}Desc.\allowbreak{}V0.\allowbreak{}isoNe}{Development.\allowbreak{}AutomorphicArithmetic.\allowbreak{}AutomorphicVarieties.\allowbreak{}ModularityArgument.\allowbreak{}CuspidalDescent.\allowbreak{}GLDenseRanges.\allowbreak{}InertiaPartMaps.\allowbreak{}ValOrd}{280--392}
\leanref{UpperTheory.\allowbreak{}V.\allowbreak{}Z.\allowbreak{}Desc.\allowbreak{}V0.\allowbreak{}isoLe}{Development.\allowbreak{}AutomorphicArithmetic.\allowbreak{}AutomorphicVarieties.\allowbreak{}ModularityArgument.\allowbreak{}CuspidalDescent.\allowbreak{}GLDenseRanges.\allowbreak{}InertiaPartMaps.\allowbreak{}TvApply}{379--393}
\leanref{UpperTheory.\allowbreak{}V.\allowbreak{}Z.\allowbreak{}Desc.\allowbreak{}V0.\allowbreak{}locE}{Development.\allowbreak{}AutomorphicArithmetic.\allowbreak{}AutomorphicVarieties.\allowbreak{}ModularityArgument.\allowbreak{}CuspidalDescent.\allowbreak{}GLDenseRanges.\allowbreak{}InertiaPartMaps.\allowbreak{}RankEqOneOfInj}{22--44}

We next count the other factors in the weight fibre. Fix \(\iota:\overline{\Q}_\ell\simeq\C\). The generalized \(x\)-space before projection is nonzero. Indeed \(M\) is finite and faithful over \(B_0/J\), and Nakayama's lemma gives a nonzero summand at \(x\) after specialization. Control in \ref{B6.48d} identifies it with the corresponding classical space. Decompose that space into definite automorphic constituents, counting occurrences. The good spherical operators on each contributing constituent have the prescribed eigenvalues. In \ref{B6.04a} and \ref{B6.04aa} the semisimple sum of the attached representations of its cuspidal factors therefore has the determinant at \(x\), by Frobenius density. Irreducibility leaves one cuspidal factor of rank three with multiplicity one. Strong multiplicity one identifies the same base change \(\pi_x\) for every contributing constituent, as in \ref{B6.48a}.

At the infinite places the algebraic type selects the prescribed irreducible representation of the compact group. At a split finite place the component is determined by \(\pi_x\). At a nonsplit finite place the level is hyperspecial, and spherical base change determines the spherical representation: its Satake parameters are recovered up to inversion and permutation, the spherical Weyl group. There is therefore one possible unitary representation with these data. Its automorphic multiplicity is one by \cite[Thm.~5.4]{labesse-2011}. Here \(L/F\) is unramified at finite places, \([F:\Q]=d>3C_0+7\), the unitary group is quasi-split at the inert finite places, and its compact real groups give discrete series. Its automorphic representations are cuspidal by anisotropy modulo centre. The base change is cohomological and conjugate self-dual. For Labesse's sign condition \((*)\), the conjugate character of the prescribed compact representation, with probability Haar measure, has trace one on that representation and zero on the others. This gives the constant occurrence count needed here. The archive's trace comparison uses the weaker conclusion that this count depends only on the infinite weight.
\leanref{UpperTheory.\allowbreak{}V.\allowbreak{}Z.\allowbreak{}Mult.\allowbreak{}mult\_\allowbreak{}free}{Development.\allowbreak{}AutomorphicArithmetic.\allowbreak{}AutomorphicVarieties.\allowbreak{}ModularityArgument.\allowbreak{}MultiplicityOne.\allowbreak{}CompOfFree}{15--19}
\leanref{UpperTheory.\allowbreak{}V.\allowbreak{}Z.\allowbreak{}Desc.\allowbreak{}V0.\allowbreak{}count\_\allowbreak{}e\_\allowbreak{}of\_\allowbreak{}sph\_\allowbreak{}same\_\allowbreak{}at\_\allowbreak{}same\_\allowbreak{}bc\_\allowbreak{}dims}{Development.\allowbreak{}AutomorphicArithmetic.\allowbreak{}AutomorphicVarieties.\allowbreak{}ModularityArgument.\allowbreak{}ThresholdArguments.\allowbreak{}DenseThresholdCases.\allowbreak{}TCountE}{147--192}

At \(v\mid\ell\) the prescribed ordinary refinement contributes one dimension, including at \(x\). The generalized chamber eigenspace injects into the Jacquet module: a vector in the kernel is killed by averaging over a sufficiently large upper unipotent subgroup, and translation turns this average into a power of a strictly dominant chamber operator up to a nonzero scalar. That operator is invertible on its generalized eigenspace. Take a torus character occurring in the image. Frobenius reciprocity gives principal-series support. The geometric lemma bounds the multiplicity of any ordered tuple in the Jacquet module by the number of permutations of its inducing tuple giving that character \cite[Lem.~2.12]{bernstein-zelevinsky-1977}. The chamber eigenvalues for the first \(j\) uniformizers give successive partial products, hence determine each smooth uniformizer value in its assigned position, with the factors \(q_v^{-i}\) of \ref{B6.48a} removed. These values have distinct \(\ell\)-adic valuations. In the notation of that lemma the value for position \(i\) has valuation
\(\sum_{\tau}h_{i,\tau}v_\ell(\tau(\varpi_v))\), since the Galois character has unit value on a reciprocity lift. This also follows by taking successive quotients of the weight factors in \ref{B6.48d}. The labelled weights increase strictly at each embedding. Thus at most one permutation contributes to the generalized eigenspace, and its nonvanishing gives dimension one.
\leanref{UpperTheory.\allowbreak{}V.\allowbreak{}Z.\allowbreak{}J.\allowbreak{}D.\allowbreak{}dense\_\allowbreak{}threshold\_\allowbreak{}order\_\allowbreak{}one}{Development.\allowbreak{}AutomorphicArithmetic.\allowbreak{}AutomorphicVarieties.\allowbreak{}ModularityArgument.\allowbreak{}ThresholdArguments.\allowbreak{}DenseThresholdCases.\allowbreak{}ExistsEig}{218--322}

Let \(S_t\) be the tame type places, \(K_t\) the product of their maximal compact groups, and \(V_t\) the prescribed tame coefficient representation in the convention that the fixed vectors are in \(V_t\otimes\bigotimes_{u\in S_t}\pi_{x,u}\). Set
\[
                 b(x)=\dim_\C
          \left(V_t\otimes\bigotimes_{u\in S_t}\pi_{x,u}\right)^{K_t}.
\]
The tame type in the comparison selects zero monodromy and the fixed multiset of inertia characters extending to the Weil group. Indeed a nonzero form in the constituent supplies a nonzero tuple transforming by this coefficient type: apply a coefficient functional nonzero on one of its values. This is the defining local test of the chosen tame type. Each Weil commutator lies in inertia and all the extending characters are one on it. Thus the Frobenius-semisimple Weil action is a sum of characters. The local representation is generic, since \(\pi_x\) is cuspidal. With zero monodromy the segments in the local Langlands classification are singletons, and genericity makes their induction irreducible. Its restriction to the maximal compact, by Iwasawa decomposition, is induction from the integral Borel of the restrictions of these characters to units. Local reciprocity identifies those restrictions with the prescribed inertia characters. Their order can be chosen the same for \(x\) and \(y\), since the irreducible inductions permit permutation. The resulting compact representations are isomorphic, so \(b(x)=b(y)\).
\leanref{UpperTheory.\allowbreak{}V.\allowbreak{}Hi.\allowbreak{}partUR2Type\_\allowbreak{}alt}{Development.\allowbreak{}AutomorphicArithmetic.\allowbreak{}AutomorphicVarieties.\allowbreak{}ModularityArgument.\allowbreak{}ThresholdArguments.\allowbreak{}DenseThresholdCases.\allowbreak{}AutomorphicTensorInductionOfLocalPlaces}{76--113}
\leanref{UpperTheory.\allowbreak{}V.\allowbreak{}Z.\allowbreak{}J.\allowbreak{}D.\allowbreak{}resEq2\_\allowbreak{}alt}{Development.\allowbreak{}AutomorphicArithmetic.\allowbreak{}AutomorphicVarieties.\allowbreak{}ModularityArgument.\allowbreak{}ThresholdArguments.\allowbreak{}DenseThresholdCases.\allowbreak{}ResEq2}{27--47}

The unprojected space gives \(b(x)>0\) and \(\pi_{x,v_0}^{K^1}\ne0\). Its component at \(v_0\) is generic for the same cuspidality reason. All the remaining spherical factors are lines. Taking the tensor product of the spaces just computed gives
\[
                    \mu_e(x)=b(x)(q_0-1)^2.
\]
The same calculation at \(y\) proves the assertion.
\end{proof}

\subsection*{Proof of \ref{B6.30ae}}\label{prf:B6.30ae}\phantomsection
\begin{proof}
Write \(T=B_0/J\), \(P=eM\), and \(T'\) for the reduced image of \(T\) acting on \(P\). The module \(P\) is finite free over \(\Lambda\), being a direct summand of a finite free module over a local ring. The weight ring is a power series ring over \(O\) on the free local unit coordinates. An inertia derivation as in the statement kills \(\Lambda\): those coordinates are values of the ordered characters up to fixed scalar factors. Differentiation of power series at an integral arithmetic point is determined by the coordinate derivatives. Indeed after finite coefficient extension the ideal of that weight is generated by the coordinate differences.

Fix the algebraic weight of \(x\). The characters with pairwise distinct finite unit parts at each ordinary place detect \(T\). To see this, all finite torus characters at that algebraic weight detect \(\Lambda\) by Fourier inversion on the finite quotients, after translating by the algebraic character. The separated ones still detect it: multiply by the product, for each pair of coordinates at each place, of the differences of their universal finite-part values on one free generator. This product is nonzero in the power series domain and vanishes at every omitted character. In a separated fibre the classical operators act by scalars on each simultaneous constituent by \ref{B6.48c} and \ref{B6.48d}. Applying separation to the matrix entries on \(M\) proves the assertion. In particular \(T\) is reduced. Its image on \(P\) is already reduced, as it is detected on the projected summands of the same fibres.

There is \(b\in T\) with \(x(b)\ne0\) such that every detecting character \(y\) with \(y(b)\ne0\) is irreducible. Choose nine group elements whose images at \(x\) form a matrix basis, and take the determinant of their matrix of trace pairings. At \(x\) it is nonzero. At a detecting character the classical semisimple representation has this same specialized trace pairing, so nonvanishing implies full matrix span. By Lemma~\ref{B6.30ad} every such \(y\), and \(x\) itself, belongs to the support of \(P\). If \(a\) belongs to the kernel of \(T\to T'\), then \(ba\) vanishes at every detecting character, hence \(ba=0\). Evaluation and the product rule at \(x\) give \(x(a)=D(a)=0\) for any derivation \(D\) at \(x\). Thus we can work on \(T'\).
\leanref{UpperTheory.\allowbreak{}V.\allowbreak{}J.\allowbreak{}irrOpenW\_\allowbreak{}alt}{Development.\allowbreak{}AutomorphicArithmetic.\allowbreak{}AutomorphicVarieties.\allowbreak{}ModularityArgument.\allowbreak{}ThresholdArguments.\allowbreak{}DenseThresholdCases.\allowbreak{}KeepIrrOpenW}{15--31}
\leanref{UpperTheory.\allowbreak{}V.\allowbreak{}Z.\allowbreak{}Desc.\allowbreak{}V0.\allowbreak{}unrAlgE}{Development.\allowbreak{}AutomorphicArithmetic.\allowbreak{}AutomorphicVarieties.\allowbreak{}ModularityArgument.\allowbreak{}CuspidalDescent.\allowbreak{}GLDenseRanges.\allowbreak{}InertiaPartMaps.\allowbreak{}UnrAlgE}{5--32}

Put \(K_\Lambda=\operatorname{Frac}\Lambda\). The finite torsion-free reduced \(\Lambda\)-algebra \(T'\) embeds into a product of finite separable fields over \(K_\Lambda\). Choose \(t\in T'\) separating its geometric generic characters and also the characters in the weight fibre of \(x\). This can be done by taking an integral linear combination of a finite set of module generators: each pair of distinct characters rules out a proper linear hyperplane, and there are finitely many pairs. Write the characteristic polynomial of \(t\) on \(P\) as \(\prod_i f_i^{m_i}\), with the \(f_i\) distinct, monic and irreducible over \(K_\Lambda\). They belong to \(\Lambda[X]\) by integrality and normality of \(\Lambda\). Faithfulness ensures that every generic factor occurs. The polynomial \(f=\prod_i f_i\) kills \(t\) in \(T'\).

Let \(r_i\) be the multiplicity of \(x(t)\) in \(f_i\) after specialization at the weight of \(x\). Separation in that fibre and simultaneous primary decomposition give
\[
                       \mu_e(x)=\sum_i m_i r_i.
\]
For each \(i\) with \(r_i>0\), the branch indexed by \(f_i\) passes through \(x\). Indeed \(\Lambda[t]\simeq\Lambda[X]/(f)\), the branch algebra is integral over \(\Lambda[X]/(f_i)\), and lying over gives a character on that branch at the specified root. Separation in the fibre identifies it with \(x\). On this branch choose a detecting \(y\) with \(b(y)\ne0\), \(f_i'(y(t))\ne0\), and \(f_j(y(t))\ne0\) for \(j\ne i\), where coefficients are evaluated at the weight of \(y\). These conditions define a nonempty open on the branch: \(b\) is nonzero at \(x\), and the generic roots are separate and simple. The detecting characters with \(b\ne0\) meet it. Indeed, after multiplication by \(b\), their vanishing test is the test on \(T\) above. To isolate a branch, multiply also by an element in every other minimal prime and outside the chosen one. Such an element exists by incomparability and finiteness of the minimal primes. Equivalently one may perform this test on \(T[1/b]=T'[1/b]\).

At \(y\) the generalized \(y\)-space is contained in the primary space for \(t\) at \(y(t)\), of dimension \(m_i\). Lemma~\ref{B6.30ad} gives \(\mu_e(x)\le m_i\). Since \(\mu_e(x)>0\), the displayed sum forces exactly one \(r_i=1\) and the others zero. Thus \(x(f'(t))\ne0\). Differentiating \(f(t)=0\) gives \(D(t)=0\).

For any \(a\in T'\) there is \(H_a\in\Lambda[X]\) satisfying \(f'(t)a=H_a(t)\). Indeed, over a splitting field write \(t_s,a_s\) for their values at the geometric generic characters and take
\[
                        H_a(X)=\sum_s a_s\prod_{u\ne s}(X-t_u).
\]
Its coefficients are Galois invariant and integral over \(\Lambda\), so lie in \(\Lambda\). Evaluation at the generic characters proves the identity by torsion-freeness. Differentiation now gives \(x(f'(t))D(a)=0\), proving the lemma.
\end{proof}

\subsection*{Proof of \ref{B6.30af}}\label{prf:B6.30af}\phantomsection
\begin{proof}
Put \(\mathfrak q=\ker\xi\), \(d=[F:\Q]\), \(D=1+3d\). We give the dimension and reducedness arguments separately.

The map \(\phi\) is onto. The traces and ordinary characters generate \(B\) topologically. For the traces, the difference of the matrices at the marking and its polarized inverse is diagonal with residually distinct entries. Its spectral projectors are polynomials in it over the trace coefficient ring, by simple-root lifting. Inserting the projectors in traces of words recovers the products of entries around cycles. In rank three the polarisation identifies an entry with the reflected transpose entry at the polarized inverse group element. A product of entries of total weight zero is a product of cycles after using these identifications: pair weights \(1,-1\) and \(2,-2\), and group the remaining weights in triples \(1,1,-2\) or \(-1,-1,2\). A weight-one entry can occupy either \(01\) or \(12\) by reflection, and likewise for weight minus one. This gives the asserted generators. The link is onto on the matrix group and on the ordinary local groups, and its twist multiplies the generators by coefficient units. The image is closed by completeness and compactness of the determinant ring. The component propagation in \ref{B6.51a}, first for rank two and then across the rank-three block locus, gives \(\phi(J)\subset P\) at every prime \(P\) off the three indicated loci. Every prime below \(\mathfrak q\) is such a prime.
\leanref{UpperTheory.\allowbreak{}V.\allowbreak{}Pr.\allowbreak{}phiEx}{Development.\allowbreak{}AutomorphicArithmetic.\allowbreak{}AutomorphicVarieties.\allowbreak{}CurveDVRSpecialization.\allowbreak{}PxIsCont}{231--288}
\leanref{UpperTheory.\allowbreak{}V.\allowbreak{}J.\allowbreak{}absGood}{Development.\allowbreak{}AutomorphicArithmetic.\allowbreak{}AutomorphicVarieties.\allowbreak{}ModularityArgument.\allowbreak{}LinkedUnramifiedChains.\allowbreak{}DerivativeDivisibility.\allowbreak{}AbsGood}{10--29}

Every component of \(B_0/J\) has dimension \(D\). It is a finite torsion-free reduced algebra over the power series weight ring in \(3d\) variables over \(O\), by \ref{B6.48c} and the detection argument of Lemma~\ref{B6.30ae}. Its minimal primes contract to zero by torsion-freeness, and finite integral extension preserves dimension. Off the triangular loci the matrix fibre over \(B\) has dimension at most one. For example, invert nonzero homogeneous entries of opposite signs. Powers of every homogeneous generator then become, up to a weight-zero factor, integral over the rational function field of one nonzero homogeneous coordinate. This gives the fibre bound also for the complete local charts, by the dimension formula for a finite-type algebra over a complete local ring, applied before completion as in \ref{B6.51a}. Consequently a branch through \(\mathfrak q\) in the completed residual chart has dimension at most \(D+1\).

Here are the local heights. Complete at the reduction of the integral point, allowing a finite residue extension. The evaluation extends to this completion, and its kernel \(Q\) has one-dimensional quotient: the image lies in a finite coefficient DVR and contains \(O\). The completed chart has a presentation over a coefficient DVR with \(t\) variables and \(r\) equations, \(t-r=D\), by the ungraded presentation in \ref{B6.51a}. The preimage of \(Q\) in the power series ring has height \(t\). Thus \(\operatorname{ht}Q\ge t-r=D\). The branch bound gives \(\operatorname{ht}Q\le D\), since one can append the closed point to a chain ending at \(Q\). In the regular local ring at the preimage of \(Q\), the ideal of the \(r\) equations therefore has height \(r\). The equations form a regular sequence. The resulting local complete intersection is Cohen--Macaulay and has no embedded associated primes. This property descends to \(A_{\mathfrak q}\) by the faithfully flat local map to the localized completion. For the height calculation downstairs the formal fibre at \(Q\) has dimension zero: the quotient of the residual local ring by \(\mathfrak q\) is a finite order over \(O\), already complete. In particular the height at an integral evaluation is at least \(D\). Also, every minimal branch below \(\mathfrak q\) has contraction of dimension \(D\) in \(B\). Indeed flat going down lifts that branch to a minimal branch of the completion. The presentation gives dimension at least \(1+t-r=D+1\) there, and the fibre bound gives the opposite inequality \(1+\dim(B/(P\cap B))\).
\leanref{UpperTheory.\allowbreak{}V.\allowbreak{}J.\allowbreak{}htUp}{Development.\allowbreak{}AutomorphicArithmetic.\allowbreak{}AutomorphicVarieties.\allowbreak{}ModularityArgument.\allowbreak{}LinkedUnramifiedChains.\allowbreak{}ScalarIrreducibility.\allowbreak{}HeightLeOfDim}{45--69}
\leanref{UpperTheory.\allowbreak{}V.\allowbreak{}J.\allowbreak{}htLow}{Development.\allowbreak{}AutomorphicArithmetic.\allowbreak{}AutomorphicVarieties.\allowbreak{}ModularityArgument.\allowbreak{}LinkedUnramifiedChains.\allowbreak{}ScalarIrreducibility.\allowbreak{}HtOfSum}{34--58}
\leanref{UpperTheory.\allowbreak{}V.\allowbreak{}J.\allowbreak{}compDim}{Development.\allowbreak{}AutomorphicArithmetic.\allowbreak{}AutomorphicVarieties.\allowbreak{}ModularityArgument.\allowbreak{}LinkedUnramifiedChains.\allowbreak{}ScalarIrreducibility.\allowbreak{}CompDim}{62--88}

Fix such a minimal prime \(P\). By surjectivity and the dimension equality, \(\phi^{-1}(P\cap B)\) is a minimal prime of \(J\). Use at each \(v\mid\ell\) the auxiliary labelled weights \(h_{i,\tau}=i\) except at one embedding \(\tau_v\), where \(h_{i,\tau_v}=2i\), \(0\le i\le2\). The separated characters at this weight are dense on the chosen coefficient branch by the detection and branch-isolation argument in Lemma~\ref{B6.30ae}. There is a point above one of them on \(P\), off all three excluded loci. Indeed invert the product of one surviving upper, lower and odd-weight entry on \(P\). The resulting finite-type dominant map to the coefficient branch has image containing a nonempty open by Chevalley's theorem. Choose a separated point in that open and lift it over \(\overline{\Q}_\ell\).

The lift can be taken with all off-diagonal entries of positive valuation. Under the grading, rescaling by an element of valuation \(u\) changes the valuation of a weight-\(j\) coordinate by \(ju\). For the finitely many homogeneous generators these are simultaneous strict lower bounds on \(u\) for positive \(j\) and upper bounds for negative \(j\). Every lower bound is strictly below every upper bound: multiply suitable positive powers of the two coordinates to get weight zero. This product reduces to zero at the diagonal residual point, and hence has positive valuation. Choose a rational \(u\) in the interval and extend coefficients. The same bounds hold on all off-diagonal entries by homogeneous generation with the residual diagonal reduction. Minimal components are preserved by the grading torus, since it is connected and the components form a finite set. Denote the resulting evaluation by \(z\).
\leanref{UpperTheory.\allowbreak{}V.\allowbreak{}J.\allowbreak{}genPt}{Development.\allowbreak{}AutomorphicArithmetic.\allowbreak{}AutomorphicVarieties.\allowbreak{}ModularityArgument.\allowbreak{}LinkedUnramifiedChains.\allowbreak{}ScalarIrreducibility.\allowbreak{}ChevPrime}{41--105}
\leanref{UpperTheory.\allowbreak{}V.\allowbreak{}J.\allowbreak{}wtEx}{Development.\allowbreak{}AutomorphicArithmetic.\allowbreak{}AutomorphicVarieties.\allowbreak{}ModularityArgument.\allowbreak{}LinkedUnramifiedChains.\allowbreak{}ScalarIrreducibility.\allowbreak{}PlaceFree}{131--166}

We show that the local ring at \(z\) is regular. The marked representation there is irreducible. An invariant subspace is a sum of lines of the diagonal marking difference, whose eigenvalues are distinct. Polarisation makes the reflected annihilator invariant as well. A proper invariant coordinate subspace and its reflected annihilator impose one of the two triangular shapes or split off the middle coordinate. All three were excluded. Apply \ref{B6.48a} to the twisted determinant and its ordered flag. Use the two choices of distinguished place in \ref{B6.51a}; the same twisted matrices and ordered characters give arithmetic evaluations for both, and component propagation puts each on its Hecke quotient. The two attachments agree at the good primes, so strong multiplicity one identifies them and gives the witness at all places. At either distinguished place the underlying matrices and the fixed twist are unramified by the choice of that place.

At \(v\mid\ell\) the local degree exceeds five by \eqref{eq:p6c5-linked-layer}. Choose a sufficiently deep norm-one principal unit of infinite order, using the exponential of a nonzero trace-zero element. Its reciprocity lift has cyclotomic value one. After taking a power to kill the finite characters, the ordered values on that lift are \(\tau_v(u)^{-i}\), since the contributions with weight \(i\) multiply to a power of the norm. Their ratios have infinite order. Hence the flag is unlinked. The residual local marking separates the three residual characters. Lemma~\ref{B6.30aa} gives irreducibility on the cyclotomic kernel and on its open subgroups. The rank-three coset argument of \ref{B6.51a} now gives density of elements with separable polynomial on that kernel.

These facts give the determinant tests in \ref{B6.54}. Indeed, on the cyclotomic kernel a twisted adjoint cocycle is an ordinary adjoint cocycle. If all its trace tests on the separable elements vanish, continuity and \ref{B6.53a} make it a coboundary there. After subtraction of that boundary, its values on the quotient lie in the scalar centralizer. The quotient is abelian and acts on scalars by the nontrivial cyclotomic character. Commuting an arbitrary element with one of cyclotomic value different from one shows that this scalar cocycle is a boundary. Finitely many tests suffice because the continuous cocycle space is finite dimensional: on an open subgroup acting trivially on a residual lattice, cocycles into the lattice give homomorphisms to the semidirect product with a pro-\(\ell\) congruence group. Its source has finitely generated maximal pro-\(\ell\) quotient by the finite character condition of \ref{B6.51frame}. Values on generators and finitely many coset representatives determine the cocycle, also after clearing denominators. Thus the common kernel of the tests is cut out by finitely many of them.

A fixed-inertia derivation at \(z\) differentiates the matrices to an ordinary cocycle by multiplication on the right by the inverse matrix. The lifted flags give the local correcting coboundaries and the inertia diagonal is zero. The fixed twist has zero derivative. Full matrix span, the witness, the distinguished inertia condition and the determinant tests just proved give all the hypotheses of \ref{B6.54}. Hence the matrix derivative is \([X,r(g)]\) for one \(X\). The marking difference is diagonal and has distinct entries, so its off-diagonal derivative forces \(X\) to be diagonal. Polarisation then says that \(X\) plus its reflected transpose commutes with the representation. Irreducibility makes this sum scalar. Modulo scalars \(X\) is therefore a multiple of \(\operatorname{diag}(1,0,-1)\). A conjugation derivative fixes the ordered characters: their values are determined by the flag projectors of a local element with distinct residual eigenvalues. Thus the fixed-inertia tangent space has dimension at most one. The three characters at \(v\) contribute at most \(3d_v\) inertia derivatives by the generators in \ref{B6.51frame}, applied to their logarithmic derivatives. The full tangent dimension is at most \(1+3d=D\). The residue field of the characteristic-zero local ring is separable algebraic over \(\operatorname{Frac}O\), so its embedding dimension is bounded by this tangent dimension. Its height is at least \(D\) by the presentation above. It is regular, hence a domain.
\leanref{UpperTheory.\allowbreak{}V.\allowbreak{}J.\allowbreak{}fixOf\_\allowbreak{}alt}{Development.\allowbreak{}AutomorphicArithmetic.\allowbreak{}AutomorphicVarieties.\allowbreak{}ModularityArgument.\allowbreak{}ThresholdArguments.\allowbreak{}DenseThresholdCases.\allowbreak{}BFixOf}{316--342}
\leanref{UpperTheory.\allowbreak{}V.\allowbreak{}J.\allowbreak{}fixLe}{Development.\allowbreak{}AutomorphicArithmetic.\allowbreak{}AutomorphicVarieties.\allowbreak{}ModularityArgument.\allowbreak{}LinkedUnramifiedChains.\allowbreak{}ScalarIrreducibility.\allowbreak{}FixLe}{8--43}

Localizing further at \(P\) shows that \(A_P\) is a field. We have proved this for every minimal prime below \(\mathfrak q\). If \(\phi(a)\) survived in \(A_{\mathfrak q}\) for some \(a\in J\), the nonzero cyclic module it generates would have an associated prime. This is an associated prime of \(A_{\mathfrak q}\), hence a minimal prime by the complete-intersection argument. At that prime the local ring is a field and \(\phi(a)\) belongs to its zero ideal by component propagation. This contradicts survival of the cyclic module at its associated prime. Clearing the localization denominator gives the stated \(s\).
\end{proof}

\subsection*{Proof of \ref{B6.30ag}}\label{prf:B6.30ag}\phantomsection
\begin{proof}
Use the integral flag basis at \(v\), and write \(t(g)\) for its upper triangular matrix. Such a basis is obtained by intersecting the flag with the lattice; the successive quotients are free over the DVR. With \(D_v\) the local group, put
\[
                   c(g)=t(g)_{i,i+1}/t(g)_{i+1,i+1}.
\]
The adjacent two-dimensional subquotient twisted by \(\lambda_{v,i+1}^{-1}\) is an extension of \(1\) by \(\varepsilon_\ell\) with cocycle \(c\). Over a finite extension on which \(t\) is crystalline, the subquotient and its twisting line are crystalline. Thus the restricted Kummer class is a unit class by \cite[Ex.~3.9]{bloch-kato-1990}. Restriction multiplies the valuation by the ramification index, so the original rational class is also a unit class. Its cup product with the unramified additive character \(u\) sending arithmetic Frobenius to one is zero by local reciprocity. Consequently there is a continuous cochain \(\beta\) such that
\[
       \beta(gh)=\beta(g)+\varepsilon_\ell(g)\beta(h)-c(g)u(h).
\]
Compactness permits multiplying \(u,\beta\) by the same scalar to make them integral, with \(u\ne0\).
\leanref{UpperTheory.\allowbreak{}V.\allowbreak{}Z.\allowbreak{}J.\allowbreak{}stepF}{Development.\allowbreak{}AutomorphicArithmetic.\allowbreak{}AutomorphicVarieties.\allowbreak{}ModularityArgument.\allowbreak{}LinkedUnramifiedChains.\allowbreak{}CrystallineCardinalities.\allowbreak{}TwoApply}{122--155}
\leanref{UpperTheory.\allowbreak{}V.\allowbreak{}Z.\allowbreak{}J.\allowbreak{}cupUnr}{Development.\allowbreak{}AutomorphicArithmetic.\allowbreak{}AutomorphicVarieties.\allowbreak{}ModularityArgument.\allowbreak{}LinkedUnramifiedChains.\allowbreak{}CrystallineCardinalities.\allowbreak{}CupUnr}{101--134}

Put \(O_m=O/\ell^mO\), \(q=|O_m|\), \(m\ge1\). At each ordinary place \(w\), let \(Z_w^\triangle\) consist of the continuous upper triangular adjoint cocycles over \(O_m\) with zero diagonal on inertia. For a finite local module \(V\), \ref{B6.05d}(1) and \(|B^1|=|V|/|H^0|\) give
\[
                 |Z^1(D_w,V)|=|V|^{1+d_w}|H^0(D_w,V^\dagger)|.
\]
The local action is continuous with finite image. Upper triangular matrices have rank six, and the diagonal is a trivial quotient for conjugation. Requiring its three characters to vanish on inertia costs at most \(q^{3d_w}\), by the \(d_w\) inertia generators of \ref{B6.51frame}. Hence \(|Z_w^\triangle|\ge q^{6+3d_w}\).

At \(v\) first use upper triangular matrices with the two indicated diagonal entries equal. Their rank is five, and their \(i,i+1\) entry is an equivariant surjection onto \(O_m(1)\): the extra conjugation term is a multiple of the difference of the two diagonal entries. The dual invariants thus have at least \(q\) elements. There are two independent diagonal inertia characters to kill. The same formula gives at least \(q^{6+3d_v}\) cocycles with the two diagonal entries equal on the whole local group.

There is an additional direction. In indices \(0,1,2\) use the matrix cochain
\[
 Z(g)=
 \begin{cases}
 u(g)E_{00}+\beta(g)E_{01},&i=0,\\
 -u(g)E_{22}+\beta(g)E_{12},&i=1 .
 \end{cases}
\]
Modulo the line \(O_m E_{02}\) it is a cocycle. Indeed conjugation on an adjacent entry \(j,j+1\) multiplies that entry by \(\lambda_j/\lambda_{j+1}\) and subtracts
\((t_{j,j+1}/t_{j+1,j+1})(Z_{jj}-Z_{j+1,j+1})\). This gives the equation for \(\beta\). The obstruction to lifting is in \(H^2(D_v,O_m(\lambda_0/\lambda_2))\). At the local marking element the three diagonal values are distinct modulo the maximal ideal. Since one adjacent ratio equals \(\varepsilon_\ell\), the outer ratio differs from \(\varepsilon_\ell\) there by a unit difference. Local duality in \ref{B6.05d}(1) makes the obstruction group zero. Correcting the \(02\) entry therefore gives a cocycle in \(Z_v^\triangle\). Its indicated diagonal difference at some \(g\) is the image of a fixed \(c_0=u(g)\ne0\). Multiples realizing distinct elements of \(c_0O_m\) give disjoint translates of the equal-diagonal subgroup. With \(C=|O/c_0O|\) this proves \(C|Z_v^\triangle|\ge q^{7+3d_v}\), since \(|c_0O_m|\ge q/C\).

Let \(Z_w\) include arbitrary local coboundaries, after returning to the marked basis. Adding these gives at least \(q^3\) times as many cocycles as in \(Z_w^\triangle\). Indeed a matrix whose coboundary is upper triangular is itself upper triangular: conjugate the separating local element to diagonal form by an upper triangular matrix, and use the unit differences on the three entries below the diagonal. The map from \(Z_w^\triangle\times M_3(O_m)\) onto \(Z_w\) thus has kernel of order at most \(q^6\).

Let \(Z_{\rm glob}\) be the polarized adjoint cocycles on the frame group satisfying these local conditions. Proposition~\ref{B6.05d}(4) gives
\[
             q^9\prod_w |Z_w|
             \ \le\ |Z_{\rm glob}|\,q^{3d}q^{9\,\#\{w\mid\ell\}} .
\]
To check the factors, the global dual invariants are zero: on the sign kernel the residual matrix composition characters are powers of \(\bar\Psi\), all distinct from \(\bar\varepsilon_\ell\) by \ref{B6.51frame}. A nonzero dual invariant at any depth would give one on a residual composition factor. At a real place the invariants identify with alternating three-by-three matrices and have order \(q^3\), since \(2\) is a unit. At the additional finite ramification places take the classes inflated from the quotient by the prescribed inertia kernel. The remaining inertia has order prime to \(\ell\); averaging followed by the procyclic Frobenius calculation gives a local \(H^1\) of the same order as \(H^0\). These classes descend to the frame quotient because the killed inertia groups normally generate the kernel. Passing from classes to cocycles gives the displayed factors of \(q^9\); the dual Selmer group has order at least one.

Normalize by imposing \(z(c)=0\) at the distinguished conjugation and by preserving the marking difference at \(g_0\). This costs at most \(q^8\). The equation at \(c^2=1\) makes \(z(c)\) symmetric under reflection in the antidiagonal, leaving six entries. When these vanish, the derivative of \(r(g_0)-r(cg_0^{-1}c)\) is the difference of a matrix and its reflected transpose. Its off-diagonal vanishing imposes the two entries \(01,10\). Write \(Z_m\) for the normalized subgroup. Multiplying the local bounds, using \(d=\sum_w d_w\), and cancelling gives
\[
                              q^2\le C|Z_m|.
\]
A normalized cocycle gives a derivation into \(O_m\) by applying the universal property to the first-order lift with matrix derivative \(z(g)r(g)\). The local coboundary corrects the flag basis, and the vanishing diagonal on inertia fixes its characters there. This construction is injective on \(Z_m\): the matrix derivatives recover \(z\) on the sign kernel, and \(z(c)=0\) recovers it on the other coset.

Let \(I=\ker\xi\). These fixed-inertia derivations are homomorphisms from the quotient of \(I/I^2\) by the classes of the inertia character values minus their evaluations. This is a finite \(O\)-module in the sense of finite generation, since \(R\) is Noetherian and \(R/I=O\). If its free rank were at most one, the structure theorem over a DVR would bound the number of homomorphisms to \(O_m\) by \(C' q\), with \(C'\) the order of its torsion. This contradicts \(q^2\le CC'q\) as \(m\) grows. Two dual basis vectors of its free part give the desired derivations.
\end{proof}

\subsection*{Proof of \ref{B6.30ah}}\label{prf:B6.30ah}\phantomsection
\begin{proof}
Take \(M_0=\ell N|D_K|\) for the ramification modulus of the link. It is a positive multiple of \(\ell N\) containing the ramification of the quadratic sign. The level \(N\) is even, and parity gives \(\Psi(c)=1\). Use the field \(L\), frame and seed constructed in \ref{B6.52} from the two tests \eqref{eq:p6c5-link-input-tests}. Its degree bounds are \eqref{eq:p6c5-linked-layer}. Changing from the original marking to the marking of the frame amounts to integral change of basis: the marked matrix difference has distinct residual eigenvalues, so its ordered spectral projectors lift, and the polarization pairs opposite eigenlines. Choose paired generators and normalize the middle generator, enlarging coefficients if necessary. Restriction now gives an integral point \(\xi:R\to O\) of the rank-three frame. Use the two definite comparisons of \ref{B6.51a}, enlarging \(O\) to contain the coefficients of the point and the local embeddings. The rank-three twist in both is \(\varepsilon_\ell^{-1}\).
\leanref{UpperTheory.\allowbreak{}Gal.\allowbreak{}projectedFormsSystemPart.\allowbreak{}openEx2E\_\allowbreak{}alt}{Development.\allowbreak{}AutomorphicArithmetic.\allowbreak{}AutomorphicVarieties.\allowbreak{}ModularityArgument.\allowbreak{}ThresholdArguments.\allowbreak{}DenseThresholdCases.\allowbreak{}UnitaryRankTwoFlags}{62--115}

Suppose a level flag over \(K\) has a link. The local alternative gives potential crystallinity. Transport the prolongation to one of the framed prolongations above \(\ell\) and restrict to its decomposition group over \(L\). Conjugation on the sign kernel transports the flag by change of basis. On the other coset it takes the reversed dual flag, so a relation \(\kappa_i=\varepsilon_\ell\kappa_j\) becomes \(\kappa_j^{-1}=\varepsilon_\ell\kappa_i^{-1}\) after transport. The two restricted flags, the transported flag and the flag of \(\xi\), have the same multiset of composition characters by Jordan--Hölder. Thus the framed flag is potentially crystalline there and two of its characters satisfy \(\lambda_a=\varepsilon_\ell\lambda_b\). On inertia its characters are \(\Psi^{1-a}\). Evaluate on an inertia element of infinite cyclotomic order after raising it to a power killing \(\vartheta\). This gives \(k(b-a)=1\), so \(k=1\) and \(b=a+1\). Lemma~\ref{B6.30ag} supplies two independent fixed-inertia derivations \(D_1,D_2:R\to O\).

Write \(A,B,\phi\) as in Lemma~\ref{B6.30af}. The prime of \(\xi\) on \(A\) is off both triangular loci and the middle block locus because \(G_L\) is open in \(G_K\) and the restriction there is irreducible. The determinant point \(x=\xi\circ\phi\), with \(B\) included in \(R\), is irreducible for the same reason. It is arithmetic: after the twist its inertia exponents are \(k(1-i)-1\), so in the locally algebraic convention of \ref{B6.48a} its labelled weights \(1-k(1-i)\) strictly increase. The remaining inertia factors are finite.

For \(a\in J\) take \(s\) as in Lemma~\ref{B6.30af}. Evaluation of \(s\phi(a)=0\) gives \(x(a)=0\). For either \(D=D_1,D_2\), differentiating the same equality gives
\[
                     0=\xi(s)D(\phi(a))+x(a)D(s),
                     \qquad D(\phi(a))=0 .
\]
The pullback of \(D\) to \(B_0\) therefore descends to a derivation on \(B_0/J\). It kills ordinary inertia because \(\phi\) sends each ordinary character to the framed one times the fixed scalar twist. Lemma~\ref{B6.30ae} kills this derivation. Surjectivity of \(\phi\) gives \(D(B)=0\).

Choose monomials \(p,q\) of weights \(1,-1\) with \(\xi(p)\xi(q)\ne0\). Indeed there are surviving entries of both signs off the triangular loci and one of odd weight off the block locus. If the other sign occurs only with absolute weight two, multiply that entry by the surviving odd entry. For a derivation \(E\) killing \(B\), the equality \(E(p)=0\) implies \(E(q)=0\) by differentiating \(pq\in B\). For an entry \(a\) of weight \(j>0\), differentiate \(aq^j\in B\); for \(j<0\) differentiate \(ap^{-j}\in B\). Thus \(E\) kills all matrix entries. The flag characters belong to \(B\), and the entries and characters topologically generate the complete ring. The derivation is continuous, since the derivative of a product of \(m\) elements of the maximal ideal takes values in \(\mathfrak m_O^{m-1}\). Hence \(E=0\).

Apply this to \(E=D_2(p)D_1-D_1(p)D_2\). It kills \(p\), so it is zero. Independence forces \(D_1(p)=D_2(p)=0\), whence \(D_1=0\) by the same argument. This contradiction excludes the link. Conjugating within \(G_K\) gives the conclusion at every prolongation of the place.
\end{proof}

\subsection*{Proof of \ref{B6.30a}}\label{prf:B6.30a}\phantomsection
\begin{proof}
Factor a characteristic-zero point of the marked complete ring through a finite coefficient DVR and retain the marking and its ordinary local lift. In Type I apply Lemma~\ref{B6.30aa} to $G=H=G_K$ and inertia at $\ell$. Since $\ell\nmid D_K$, rational inertia lies in $G_K$ and maps onto $\Z_\ell^\times$ under the cyclotomic character. The inertia characters are $\Psi,1,\Psi^{-1}$, with $\Psi=\varepsilon_\ell^k\vartheta$. A power of an inertia lift of $1+\ell$ kills $\vartheta$ and separates the three eigenvalues at every positive power. On inertia $\bar\vartheta^2=1$, and $\ell-1\nmid2k$, so some inertia element has $\bar\Psi^2\ne1$. The residual semisimplification is the sum of the three prescribed characters. The lemma gives irreducibility on every open subgroup of $G_K$, as required in \ref{B6.52}(2). That theorem supplies the cuspidal attachment with witness after the twist by $\varepsilon_\ell^{-1}$. Away from $\ell$ the prescribed inertia type has level dividing $N$. At each place above $\ell$ take the ordered ordinary flag. Its characters on inertia are $\Psi,1,\Psi^{-1}$, so its weights are $(k,0,-k)$ and its finite parts factor through the action on $\mu_N$. We check the local alternative in Lemma~\ref{B6.30ah}. In the distinguished local problem of Definition~\ref{B6.40m}, Lemma~\ref{B6.40n}(3) shows that the marked point is crystalline on the open subgroup $G_{K_1}$, hence potentially crystalline. For the conjugate place use the reversed dual flag; taking duals and changing basis preserve crystallinity. This gives admissibility with the local alternative. Lemma~\ref{B6.30ah} now gives an unlinked level flag, including for the adjacent weights at $k=1$.\par
\leanref{UpperTheory.\allowbreak{}V.\allowbreak{}Loc.\allowbreak{}rtPoints}{Development.\allowbreak{}AutomorphicArithmetic.\allowbreak{}AutomorphicVarieties.\allowbreak{}PadicHodgeTheory.\allowbreak{}RtPoints}{11--15}
\leanref{UpperTheory.\allowbreak{}V.\allowbreak{}U3.\allowbreak{}admissible\_\allowbreak{}rank\_\allowbreak{}three\_\allowbreak{}of\_\allowbreak{}type\_\allowbreak{}one\_\allowbreak{}of\_\allowbreak{}automorphic\_\allowbreak{}point}{Development.\allowbreak{}AutomorphicArithmetic.\allowbreak{}AutomorphicVarieties.\allowbreak{}ModularityArgument.\allowbreak{}ThresholdArguments.\allowbreak{}DenseThresholdCases.\allowbreak{}ExistsFlagCharactersOfUpperTriangular}{163--214}
 This puts the point in $\mathcal A_3$. Let $L$ be the auxiliary CM field of the linked frame. Its group is open in $G_K$, so the representation remains irreducible on $G_L$. At an ordinary place the witness of Definition~\ref{def:b6-cuspidal-witness} and the local comparison give an element of the cyclotomic kernel with all eigenvalue ratios of infinite order: use the positive Weil degree argument in Lemma~\ref{B6.48k}, with weights $k,0,-k$, and take a positive power to enter $G_L$. On the linked flag a local element with $\bar\Psi^2\ne1$ is supplied by \ref{B6.51frame} and the frame of \ref{B6.52}. Its residual flag characters, up to the common twist, are $\bar\Psi,1,\bar\Psi^{-1}$. Lemma~\ref{B6.30aa} with $G=G_L$ and $H=G_{L(\mu_{\ell^\infty})}$ gives irreducibility there and on every open subgroup of $H$. For dense separability one may now use the rank-three coset argument in \ref{B6.51a}, with this same separating element. For a permitted tangent, the finite inertia quotient away from $\ell$ and the ordinary finite-condition comparison \eqref{eq:p6c5-local-finite} at $\ell$ give the finite adjoint class. Apply \ref{B6.53} on the linked frame and descend the coboundary by averaging over the finite orbit of a cobounding element.

In Type II the middle line on $G_K$ is trivial and the corner extends to a representation $W$ with determinant $\eta_K$, the quadratic character of $K$, and trace zero at complex conjugation. Its ordered weights are $(k,-k)$. By \ref{B6.52}, $W\varepsilon_\ell^{-k}$ has a weight-$(2k+1)$ newform attachment and $W$ has the corresponding witness. To apply \ref{B6.53}, first restrict the irreducible corner from $G_K$ to the auxiliary $G_L$, normal in $G_K$. Its determinant there is one. The linked marking $t\in G_L$ has $\bar\Psi(t)^2\ne1$ by \ref{B6.51frame}, so the residual roots at $t$ are distinct. A reducible restriction would therefore be the sum of two distinct lines interchanged by some $g\in G_K$. Both $W(g)$ and $W(gt)$ would have trace zero. Reduction of their traces gives $\bar\Psi(g)^2=-1=\bar\Psi(gt)^2$, contradicting $\bar\Psi(t)^2\ne1$. Thus the restriction is irreducible and solvable base change gives the cuspidal attachment on $L$. The rank-two ordinary comparison over $K$ gives a separating element in the cyclotomic kernel by the witness argument of Lemma~\ref{B6.48k}, with weights $k,-k$; take a power into $G_L$. The local test $\bar\Psi(s)^2\ne1$ of \ref{B6.51frame} distinguishes the two residual characters $\bar\Psi,\bar\Psi^{-1}$ at $s$. The scalar twist on the linked frame preserves these tests and irreducibility. Apply Lemma~\ref{B6.30aa} with $G=G_L$ and $H=G_{L(\mu_{\ell^\infty})}$. This checks the cyclotomic restriction in \ref{B6.53} on the field where that theorem gives block rigidity. In the full adjoint the two off-block summands are $W$ and $W\otimes\eta_K$. Their finite classes vanish by \ref{B6.43}. For the restriction over $K$ itself, put $I=I_\ell\subset G_K$. The trace on $I$ is $\Psi+\Psi^{-1}$. If the restriction to $G_K\cap\ker\varepsilon_\ell$ were reducible, the action of the abelian quotient on its commuting endomorphisms would give a nontrivial self-twist $\chi$ through the cyclotomic character: the commuting endomorphisms have dimension greater than one, the conjugation actions commute, and their simultaneous eigenspaces give the intertwiners. Equivalently, in rank two one may use Lemma~\ref{B6.30aa} directly with the separating local element just constructed over $K$ and the inertia element on which $\bar\Psi^2\ne1$. In the self-twist argument $\chi^2=1$ by determinants. Inertia surjects onto the cyclotomic image, so choose $s\in I$ with $\chi(s)\ne1$. On this coset trace zero forces $\Psi(s)^4=1$. Choose $u\in I$ above $1+\ell$; since $k>0$ and $\vartheta$ is finite, $\Psi(u)$ has infinite order. The same trace identity at $su^2$ gives $\Psi(u)^8=1$, a contradiction. The remaining diagonal summand is $C_\ell(\eta_K)$, whose finite Selmer group vanishes by Lemma~\ref{B6.40p}(3).
\end{proof}

\subsection*{Proof of \ref{B6.30}}\label{prf:B6.30}\phantomsection
\begin{proof}
Choose distinct primes $r_1,r_2>30fp_K+1$ and put $K=\Q(\sqrt{-p_K})$. Its discriminant is $D_K=-p_K$, and $r_i\nmid30N|D_K|$. The characters have prime-to-$\ell$ order and $\theta_\ell^2=1$ by \ref{B6.28a}. Inertia at $p_K$ fixes $\bar\varepsilon_\ell$ and $\theta_P$. It acts nontrivially on $K$. Parity, $2\mid N$, $p_K\nmid N$ and $\ell\nmid p_Kh_*$ give the remaining datum conditions of Definition~\ref{B6.40m} and \ref{B6.40}, and the standing hypotheses $\mathscr H$ (\ref{B6.05}, \ref{B6.09}) together with \ref{B6.15} supply the cohomological inputs of Definition~\ref{B6.40m}. Corollary \ref{B6.30a} supplies its representation hypotheses, and \ref{B6.41} gives the trace estimate.

For each occurring character kernel $J$ in $(\Z/N\Z)^\times$, take the integers $D_{J,1},D_{J,2}$ of \ref{B6.42} at $r_1,r_2$. Apply the two-prime sum (6.42.3--4) in the proof of \ref{B6.42} to the occurring kernels. It gives finiteness and, with the constants $c_K,A_K$ there, bounds the sum by $2c_K(\log r_1+\log r_2)N[N(k+1)]^{A_K}$, which is at most a fixed power of $Nt$.
\end{proof}

\subsection*{Proof of \ref{B6.30b}}\label{prf:B6.30b}\phantomsection
\begin{proof}
Take \(k=2j\), \(t=2j+1\), \(a=2\), and \(O_\ell=\Z_\ell\). The trivial character belongs to \(X_{\ell,\Z_\ell}(2j)\) by \ref{B6.28}. Choose \(q\equiv3\bmod 8\) and \(d\equiv7\bmod16\), both larger than \(2t+1\), the fixed cutoffs in \ref{B6.28}, and the two primes in \ref{B6.30}. They can be chosen with \(q,d\ll_\chi t\). Indeed, apply \ref{B2.23} with modulus \(16\) to the reduced residue classes \(3,7\). For this fixed modulus its error is \(o(X/\log X)\), and any exceptional term is \(O(\operatorname{li}_2(X^{\beta_1}))=o(X/\log X)\). Subtraction at \(2X\) and \(X\) gives a prime in each class in \((X,2X]\) for all large \(X\). Take \(X\) to be the maximum of \(2t+1\) and a sufficiently large fixed cutoff. Thus \(N=\operatorname{lcm}(16,f)qd\ll_\chi t^2\).

Fix an odd prime \(\ell\). Put \(D=(\Q_\ell/\Z_\ell)(-2j)\), and let \(G_S\) be the Galois group over \(\Q\) unramified at finite primes outside \(S=\{\ell\}\). The strict group of \ref{B6.29} for the trivial character is
\[
 \ker\bigl(H^1(G_S,D)\longrightarrow H^1(\Q_\ell,D)\bigr).
\]
To see this, inertia at \(p\ne\ell\) acts trivially on \(D\), so a cocycle unramified there vanishes on inertia and descends through the unramified quotient. At a prime \(p\mid N\), \(p\ne\ell\), its unramified local class belongs to \(D/(p^{-2j}-1)D\). The scalar is a nonzero element of \(\Z_\ell\), and multiplication by it on \(D\) is onto. This accounts for every extra decomposition-group condition in \ref{B6.29}.

For \(u\ge1\), take \(D_u=D[\ell^u]\). Its Cartier dual is \(T_u=(\Z/\ell^u)(2j+1)\). The restriction \(H^0(G_S,D_u)\to H^0(\Q_\ell,D_u)\) is an isomorphism. Both groups take invariants under the same cyclotomic image, since \(\Q_\ell(\mu_{\ell^u})/\Q_\ell\) has Galois group \((\Z/\ell^u)^\times\). The first four terms of \ref{B6.09}, with ground field \(\Q\), odd \(\ell\) and \(S=\{\ell\}\), therefore identify the dual of \(H^2(G_S,T_u)\) with the kernel in degree one for \(D_u\). Passing to the direct limit on the discrete side and the inverse limit on the compact side gives
\[
 \mathcal S_{\ell,2j,1}^{\vee}
   \simeq H^2_{\mathrm{\acute et}}\bigl(\Z[1/\ell],\Z_\ell(2j+1)\bigr).
\]
Here finite-coefficient cohomology agrees with the cohomology of \(G_S\) by \cite[II, Proposition~2.9]{milne-2006}, applied to the locally constant \(\ell\)-power torsion sheaves on \(\operatorname{Spec}\Z[1/\ell]\). The finite degree-one groups satisfy the Mittag--Leffler condition. Weibel's calculation uses the Rost--Voevodsky norm-residue isomorphism \(K^M_i(F)/\ell\simeq H^i(F,\mu_\ell^{\otimes i})\) for fields \(F\) of characteristic different from \(\ell\), through the motivic-to-\'etale comparison in the motivic spectral sequence. For \(\operatorname{Spec}\Z[1/\ell]\) the \(\ell\)-cohomological dimension is two since \(\ell\) is odd. In positive even degree \(4j\), Theorem~8.2 therefore identifies the group \(H^2_{\mathrm{\acute et}}(\Z[1/\ell],\Z_\ell(2j+1))\) on the right with \(K_{4j}(\Z)[\ell^\infty]\). \cite[VI, Thms.~4.1, 4.2 and 8.2]{weibel-2013}

Consequently its logarithmic order is \(a_{\ell,2j,1}\log\ell\). For \(\ell\notin B\) with \(\ell-1\mid4j\), we have \(\ell\ge5\) and \ref{B6.36} applies. Take for \(S\) in \ref{B6.30} the remaining primes outside \(B\) dividing the order of the finite group \(K_{4j}(\Z)\). Retaining the trivial-character summand gives at most \((Nt)^{C_{\mathrm I}}\ll_\chi t^{3C_{\mathrm I}}\), as asserted.
\end{proof}

\subsection*{Proof of \ref{B6.31a}}\label{prf:B6.31a}\phantomsection
\begin{proof}
Let $z$ be the cocycle and take a finite $\ell$-adic field $K_0$ containing the values of $\eta$. Choose a countable family of finite extensions $K_i$ of $\Q_\ell$ covering $C_\ell$. Such a family exists since there are finitely many extensions of each bounded degree in a fixed algebraic closure. The cocycle identity makes $\{g:z(g)\in K_0K_i\}$ a subgroup, and continuity makes it closed. By Baire's theorem one of these subgroups is open. Adjoin the values of $z$ at finitely many coset representatives to its coefficient field $K_0K_i$. The cocycle identity now puts all values in the resulting finite extension $E$. On any subgroup where $\lambda$ is nontrivial, a local cobounding scalar is $z(g)/(\lambda(g)-1)$ for any $g$ with $\lambda(g)\ne1$. If $\lambda$ is trivial there, the local cocycle is zero and the scalar can be taken to be zero. Thus the witnesses may be taken in $E$.

Let $O_E$ be the ring of integers of $E$. The strict group for $(E/O_E)(\lambda)$ with the same local conditions has finite length. Indeed, restrict to a finite abelian field $L$ containing $\mu_\ell$ and trivializing $\eta$. The image has length at most $e_{O_E}s_{\ell,e}(L)$, finite by \eqref{eq:B6-H2-local-correction} and \ref{B6.15}. The restriction kernel argument of \ref{B6.32a} bounds its kernel by a finite multiple of the length of the global invariants. These have finite length since $\lambda$ is nontrivial on inertia at $\ell$.

Choose a uniformizer $\pi$ of $O_E$ and $\pi^a$ killing the strict group. For each $m\ge1$, reducing $\pi^{-m-a}z$ modulo $O_E$ gives a class in that group. Multiplication by $\pi^a$ makes it zero, so some $u_m\in E$ satisfies
\[
 z(g)-(\lambda(g)-1)u_m\in\pi^m O_E\qquad(g\in G_\Q).
\]
For $g_*$ with $\lambda(g_*)\ne1$, evaluation gives $u_m\to z(g_*)/(\lambda(g_*)-1)$. The congruences show that this limit cobounds $z$.
\end{proof}

\subsection*{Proof of \ref{B6.31d}}\label{prf:B6.31d}\phantomsection
\begin{proof}
Take the exponents of \eqref{eq:b6-sp4-levels}.
Use $(n',p')=(n_{\mathrm{II}},1)$ in II and $(n_{\mathrm{III}},p_1)$ in III. Enlarge $O$ to contain the $2^{n_{\mathrm{II}}+n_{\mathrm{III}}}$-th roots of unity as well as the values of the auxiliary characters. Take a primitive odd character $\xi_2$ modulo $2^{n'}$. In II put $\xi_0=1$. In III take $\xi_0$ modulo $p_1$ of order $\operatorname{ord}\theta$, which divides $o_{\rm odd}\mid p_1-1$ by \ref{B6.28a}. In particular $\xi_0$ is even. Put $\xi=\xi_0\xi_2$, $\Psi=\theta\varepsilon_\ell^k$ and $\alpha=\varepsilon_\ell\omega^{-1}\xi$. The distinguished lift is $\rho_0=\operatorname{diag}(\alpha\Psi,\Psi,\alpha,1)$.

Let $T=\{\ell\}\cup\{v\ne\ell:\kappa_v\ne1\}$. In II, every $v\in T$ divides $\ell f$ by \ref{B6.28a}, and $2^s$ with $s=c_2(\theta)+1+v_2(o)$ kills the $2$-parts of $\theta$ and $\xi_0$. In III, $T=\{\ell\}$ and $s=0$. Since $\ell\ne p_1$ in III, \eqref{eq:b6-sp4-levels} gives
\begin{equation}\label{eq:b6-sp4-level-bounds}
 c_2(\theta)<n',\qquad 2^sv^2\le2^{n'},\qquad (p',v)=1
 \quad(v\in T).
\end{equation}

For $v\in T$ take arithmetic Frobenius lifts $\sigma_v$, with $\varepsilon_\ell(\sigma_\ell)=1$ by adjustment in inertia. At $v\ne\ell$ the nontrivial $\ell$-power character at $v$ factors through the prime-to-$v$ part of the local unit quotient, so $\ell\mid v-1$. Thus $\bar\varepsilon_\ell(\sigma_v)=1$ throughout $T$. We have
\begin{equation}\label{eq:b6-sp4-separation}
 \overline{\xi\Psi^i}(\sigma_v)\ne1
 \qquad(v\in T,\ i\in\Z).
\end{equation}
Indeed, let $u$ be an odd common multiple of the odd parts of the orders of $\theta$ and $\xi_0$. An equality to one, raised to $2^su$, gives $\xi_2(v)^{2^su}=1$ by injectivity of reduction on $2$-power roots of unity. The primitive character $\xi_2$ is faithful on the units congruent to $1\bmod4$. Since $v^{2^{s+1}u}\equiv1\bmod4$, we would have $2^{n'}\mid v^{2^{s+1}u}-1$. But \eqref{eq:B6-lte} applied to $v^2$, followed by \eqref{eq:b6-sp4-level-bounds}, gives
\[
 v_2(v^{2^{s+1}u}-1)=s+v_2(v^2-1)<n'.
\]

The lifts of $\bar\xi$, $\bar\Psi\bar\xi^{-1}$ and $\bar\Psi\bar\xi$ have respective odd-modulus and dyadic factors
\[
 (\xi_0,\xi_2),\qquad
 (\omega^k\theta_1\xi_0^{-1},\theta_2\xi_2^{-1}),\qquad
 (\omega^k\theta_1\xi_0,\theta_2\xi_2).
\]
The products are odd since $\Psi$ is even. Evaluation at $1+2^{n'-1}$ gives primitivity of the dyadic factors by \eqref{eq:b6-sp4-level-bounds}. The primitive characters inducing the odd-modulus factors have conductors $f'\mid p'\ell\operatorname{cond}(\theta_1)$ and orders dividing $\ell^{f_0}-1$. For $\xi_0$ this follows from its triviality in II and $\theta_1=\theta$ in III. By \eqref{eq:b6-sp4-levels}, with $\operatorname{cond}(\theta_1)\mid f$ in III, we have $f'4^{n_0}\le2^{n'}$. Also $4\mid\ell^{f_0}-1$ by the definition of $f_0$. At $v\in T$ a local coboundary vanishes on inertia in the character kernel. At other places the local condition kills the intersection of inertia with the kernel of the residual diagonal action, of index prime to $\ell$. The product criterion \ref{B6.59} gives residual Selmer vanishing for the three characters and their inverses.

In each class $\theta_P$ is nontrivial and has conductor and order prime to $\ell$. Choose $g$ fixing $\mu_\ell$ with $\bar\theta_P(g)\ne1$. The prime-to-$\ell$ part of the local character at $\ell$ factors through $\mu_\ell$, so $\bar\Psi(g)=\bar\theta_P(g)$. As $\xi$ is odd and $\Psi$ even, either $g$ or its product with complex conjugation gives $g_0$ with $\bar\xi(g_0)\ne1$ and $\bar\Psi(g_0)\ne1$. Complete the good datum $\mathcal D$ with $S_{\mathcal D}=\{p:p\mid2\ell Np'\}$ and
\[
 \phi=\lim_{m\to\infty}\sigma_\ell^{\,\ell^{m!}},\qquad
 s_0\in I_\ell\cap\ker\theta\cap\ker\omega,\qquad
 \varepsilon_\ell(s_0)\text{ of infinite order}.
\]
The powers defining $\phi$ stabilize in each finite quotient at the prime-to-$\ell$ part of the image of $\sigma_\ell$. For $s_0$, raise an inertia preimage of $1+\ell$ to $\operatorname{ord}(\theta)(\ell-1)$. The character $\xi$ is unramified at $\ell$, so $\xi(s_0)=1$.
\leanref{UpperTheory.\allowbreak{}V.\allowbreak{}X.\allowbreak{}xvan}{Development.\allowbreak{}AutomorphicArithmetic.\allowbreak{}AutomorphicVarieties.\allowbreak{}LevelCohomology.\allowbreak{}HcXvan}{19--41}
\end{proof}

\subsection*{Proof of \ref{B6.31e}}\label{prf:B6.31e}\phantomsection
\begin{proof}
With the alternating form $J$ of antidiagonal $(1,1,-1,-1)$, a matrix $X\in M_4$ lies in $\mathfrak{sp}_4$ precisely when $X_{ij}=-s_is_jX_{3-j,3-i}$ for $s=(1,1,-1,-1)$. The $10$-dimensional space $\mathfrak{sp}_4$ decomposes under $\operatorname{Ad}r_0$ into root spaces of bidegrees:
\begin{itemize}
\item first degree zero: $(0,0)$ of dimension $2$ (diagonal $X_{00}=-X_{33}$, $X_{11}=-X_{22}$, character $1$), $(0,1)$ of dimension $1$ ($X_{02}=X_{13}$, character $\Psi$), and $(0,-1)$ of dimension $1$ ($X_{20}=X_{31}$, character $\Psi^{-1}$);
\item nonzero first degree (six $1$-dimensional cross-block spaces): $(1,0)$ ($X_{01}=-X_{23}$, character $\alpha$), $(-1,0)$ ($X_{10}=-X_{32}$, character $\alpha^{-1}$), $(1,1)$ ($X_{03}$, character $\alpha\Psi=\nu$), $(-1,-1)$ ($X_{30}$, character $\nu^{-1}$), $(-1,1)$ ($X_{12}$, character $\Psi/\alpha$), and $(1,-1)$ ($X_{21}$, character $\alpha/\Psi$).
\end{itemize}
The four strictly upper-triangular roots ($i<j$) have bidegrees $(0,1),(1,0),(-1,1),(1,1)$, and the four strictly lower-triangular roots ($i>j$) are their negatives.

For (1), fix $v\in T\setminus\{\ell\}$ and put $N_v=I_v\cap\ker\theta$. Since $D_v/I_v\simeq\widehat\Z$ is abelian, every commutator $[g,h]$ in $D_v$ lies in $I_v$. Because $\theta:D_v\to O^\times$ is a homomorphism into an abelian group, $[g,h]\in\ker\theta$ as well. Hence $[D_v,D_v]\subseteq N_v$, so $D_v/N_v$ is abelian, topologically generated by $\sigma_v$ and an inertia generator $\tau_v$ of order $m_v=\#\theta(I_v)$ with $\sigma_v\tau_v=\tau_v\sigma_v$. Since $\ell\mid m_v$, $\zeta=\theta(\tau_v)\ne1$, and because $\alpha$ is unramified at $v$, $r_0(\tau_v)=\operatorname{diag}(\zeta,\zeta,1,1)$.

Let $U_v^\wedge=\widehat{U_v[1/\ell]}_{\ker x_{0,v}}$, and consider a lift $b:D_v/N_v\to\mathrm{GSp}_4(A)$ of $r_0|_{D_v}$ over a local Artinian $E$-algebra $A$ with residue field $E$ satisfying $\det(XI_4-b(\tau_v))=(X-\zeta)^2(X-1)^2$. Since $b(\tau_v)^{m_v}=I_4$ and $X^{m_v}-1$ is separable over $E$, writing $X^{m_v}-1=(X-\zeta)(X-1)Q(X)$ shows that $Q(b(\tau_v))$ is invertible over $A$, whence $(b(\tau_v)-\zeta I_4)(b(\tau_v)-I_4)=0$. As $\zeta-1\in E^\times$, $A^4=W_\zeta\oplus W_1$ splits into the free rank-$2$ eigenspaces of $b(\tau_v)$ for $\zeta$ and $1$, lifting $Ee_0\oplus Ee_1$ and $Ee_2\oplus Ee_3$. Since $b(\tau_v)$ has similitude $\nu(\tau_v)=\zeta\ne0,1$, both $W_\zeta$ and $W_1$ are isotropic for $J$, hence complementary Lagrangian subspaces of $(A^4,J)$. Because $b(\sigma_v)$ commutes with $b(\tau_v)$, it preserves $W_\zeta$ and $W_1$; conversely, any complementary Lagrangian pair $(W_\zeta,W_1)$ lifting $(Ee_0\oplus Ee_1,Ee_2\oplus Ee_3)$ together with any invertible $M_\zeta\in\GL(W_\zeta)$ lifting $\operatorname{diag}(\nu(\sigma_v),\Psi(\sigma_v))$ uniquely determines $b(\tau_v)$ and $b(\sigma_v)$ (the action on $W_1\simeq W_\zeta^\vee$ being forced by similitude $\nu(\sigma_v)$). Near the base point, $W_\zeta$ is the graph of a matrix $C$ in positions $(2,0),(2,1),(3,0),(3,1)$ satisfying the single Lagrangian condition $C_{20}=C_{31}$ (three unconstrained parameters of bidegrees $(0,-1),(1,-1),(-1,-1)$), $W_1$ is the graph of a matrix $D$ in positions $(0,2),(0,3),(1,2),(1,3)$ with $D_{02}=D_{13}$ (three unconstrained parameters of bidegrees $(0,1),(1,1),(-1,1)$), and $M_\zeta$ supplies four unconstrained parameters of bidegrees $(0,0),(0,0),(1,0),(-1,0)$. Thus $U_v^\wedge/J_vU_v^\wedge$ is formally smooth over $E$ of dimension $6+4=10$, with bigraded coordinate counts $(1,1,2)$ in $((0,-1),(0,1),(0,0))$ and $1$ in each of the six cross-block bidegrees.

Since $U_v^\wedge/J_vU_v^\wedge\simeq E[[t_1,\ldots,t_{10}]]$ is a domain, $I_v$ is the unique minimal prime of $\operatorname{sat}_\ell(\sqrt{J_v})$ contained in $\ker x_{0,v}$, and $B_v=U_v/I_v$ embeds in $\widehat{B_v[1/\ell]}_{\ker x_{0,v}}\simeq E[[t_1,\ldots,t_{10}]]$ by Krull's intersection theorem. Because $J_v$ is invariant under $\mathrm{Sp}_4$-conjugation, the preimage of $I_v\cdot(U_v\widehat\otimes_O O[[z]])$ under universal formal $\mathrm{Sp}_4$-conjugation congruent to $I_4$ is a prime ideal containing $\operatorname{sat}_\ell(\sqrt{J_v})$ and contained in $\ker x_{0,v}$, hence contains $I_v$; this gives conjugation invariance and a bigraded homogeneous presentation. For each cross-block root $\chi=\alpha^{j_1}\Psi^{j_2}$ ($j_1\ne0$), \eqref{eq:b6-sp4-separation} gives $\bar\chi(\sigma_v)\ne1$, while $\bar\varepsilon_\ell(\sigma_v)=1$ since $\ell\mid v-1$; thus $H^0(D_v,E(\chi))=H^2(D_v,E(\chi))=0$, so $H^1(D_v,E(\chi))=0$ by local Euler characteristic and $Z^1(D_v,E(\chi))=B^1(D_v,E(\chi))$ is $1$-dimensional.

For (2), let $U_\ell^\wedge=\widehat{U_\ell[1/\ell]}_{\ker x_{0,\ell}}$ and consider an $A$-point $b:D_\ell\to\mathrm{GSp}_4(A)$ of $U_\ell^\wedge/J_\ell U_\ell^\wedge$ over a local Artinian $E$-algebra $A$ with residue field $E$. Since $k\ge2$ and $\varepsilon_\ell(s_0)$ has infinite order, the eigenvalues $d_i=\varepsilon_\ell(s_0)^{w_i}$ ($w=(k+1,k,1,0)$) of $r_0(s_0)$ are pairwise distinct in $E^\times$. The operators $e_i=P_i/\prod_{j\ne i}(d_i-d_j)$ are orthogonal rank-$1$ idempotents summing to $I_4$ with $e_i^tJe_j=0$ for $i+j\ne3$, giving a $D_\ell$-stable isotropic full flag in $A^4$ with inertia characters $(\delta_0,\delta_1,\delta_2,\delta_3)$. Choosing a symplectic basis adapted to this flag congruent to the standard basis modulo $\mathfrak m_A$ expresses $b=hb_{\mathrm{tri}}h^{-1}$ with $b_{\mathrm{tri}}:D_\ell\to B_{\mathrm{GSp}_4}(A)$ upper triangular of similitude $\nu$ and inertia characters $(\delta_i)$, and $h\in\mathrm{Sp}_4(A)\cap(I_4+\mathfrak m_AM_4)$ unique up to right multiplication by $B_{\mathrm{Sp}_4}(A)\cap(I_4+\mathfrak m_AM_4)$.

We check that upper-triangular lifts $b_{\mathrm{tri}}$ with fixed similitude $\nu$ and fixed inertia characters $(\delta_i)$ are formally smooth over $E$. The diagonal characters satisfy $\lambda_0\lambda_3=\lambda_1\lambda_2=\nu$ and $\lambda_i|_{I_\ell}=\delta_i|_{I_\ell}$, so they are parameterised freely by two unramified scalars $\lambda_0(\sigma_\ell),\lambda_1(\sigma_\ell)$ in bidegree $(0,0)$. The four strictly upper roots are $\chi\in\{\Psi,\alpha\Psi,\alpha,\Psi/\alpha\}$, with cyclotomic weights $k,k+1,1,k-1\ge1$ (as $k\ge2$), so $H^0(D_\ell,E(\chi))=0$ for all four. For $H^2(D_\ell,E(\chi))\simeq H^0(D_\ell,E(\varepsilon_\ell\chi^{-1}))^\vee$: for $\chi=\Psi$ or $\alpha\Psi$, the cyclotomic weight of $\varepsilon_\ell\chi^{-1}$ is $1-k\le-1$ or $-k\le-2$; for $\chi=\alpha=\varepsilon_\ell\omega^{-1}\xi$, $(\varepsilon_\ell\alpha^{-1})(\sigma_\ell)=\xi(\sigma_\ell)^{-1}\ne1$ by \eqref{eq:b6-sp4-separation}; for $\chi=\Psi/\alpha$, either $k\ge3$ (cyclotomic weight $2-k\le-1$) or $k=2$ and $(\varepsilon_\ell(\Psi/\alpha)^{-1})(\sigma_\ell)=(\xi\Psi^{-1})(\sigma_\ell)\ne1$ by \eqref{eq:b6-sp4-separation}. Thus $H^2(D_\ell,E(\chi))=0$ for all four strictly upper roots, and local Euler characteristic gives $\dim_E H^1(D_\ell,E(\chi))=1$ and $\dim_E Z^1(D_\ell,E(\chi))=2$. Lifting along the central series of the nilradical is unobstructed at each step. Adjoining the four strictly lower-triangular framing directions $h\in\mathrm{Sp}_4/B_{\mathrm{Sp}_4}$ (each contributing $\dim_E B^1(D_\ell,E(\chi^{-1}))=1$ since $H^0(D_\ell,E(\chi^{-1}))=0$) shows that $U_\ell^\wedge/J_\ell U_\ell^\wedge$ is formally smooth over $E$ of dimension $2+4\cdot2+4\cdot1=14$, with the stated bidegree counts. Just as in (1), $I_\ell$ is therefore a single prime ideal invariant under $\mathrm{Sp}_4$-conjugation congruent to $I_4$, and $B_\ell=U_\ell/I_\ell$ is a complete Noetherian local $O$-domain with a bigraded homogeneous presentation.

For (3), consider the framed presentation of \ref{B6.58} over $B=B_\ell\widehat\otimes_O\widehat\bigotimes_{v\in T\setminus\{\ell\}}B_v$ in any cross-block bidegree $j=(j_1,j_2)$ with $j_1\ne0$. By the residual Selmer clause of Definition~\ref{B6.31b} and \eqref{eq:b6-sp4-separation} at $\sigma_v$ ($v\in T$), we have $H^1_{\mathrm{fr},j}=0$; since $\bar\alpha^{j_1}\bar\Psi^{j_2}$ is odd and differs from $\bar\varepsilon_\ell$ at $\sigma_\ell$, Poitou--Tate duality \ref{B6.09} and the local and global Euler characteristics give $\dim\operatorname{Ob}_j-\dim H^1_{\mathrm{fr},j}=0$, hence $\operatorname{Ob}_j=0$. Thus there are no relative variables or relative relations in cross-block bidegrees; the cross-block coordinates of the ambient formal chart are the $9$ from $B_\ell$ and the $6$ from each $B_v$ ($v\in T\setminus\{\ell\}$). At $\ell$, $r(\phi)=b_\ell(\phi)$ depends only on $B_\ell$, and along each of the six cross-block coboundary directions $\delta X$, the corresponding entry of $b_\ell(\phi)-r_0(\phi)$ has derivative $(r_0(\phi)_{ii}-r_0(\phi)_{jj})X_{ij}$, a unit by \eqref{eq:b6-sp4-separation}. At each $v\in T\setminus\{\ell\}$, when the variables of $B_\ell$ are set to zero, we have $z|_{D_\ell}=0$, $z|_{D_v}=\beta_v+\delta b_v\in B^1(D_v,E(\chi))$ ($v\in T\setminus\{\ell\}$), and $z|_{I_q\cap\ker\bar r}=0$ ($q\notin T$); if $z\ne0$, scaling $z$ to take primitive values in $O$ keeps each local cobounding scalar $u_v=z(\sigma_v)/(\chi(\sigma_v)-1)$ integral in $O$ by \eqref{eq:b6-sp4-separation}, so reducing modulo $\mathfrak m_O$ produces a nonzero element of $H^1_{\mathrm{fr},j}=0$, a contradiction. Thus $z=0$, so for the frame $h_v=I_4-\epsilon b_v$ of \ref{B6.58} the intertwining relation $z|_{D_v}=\beta_v+\delta b_v$ with $\beta_v=\delta X_v\in B^1(D_v,V_j)$ gives $b_v=-X_v$; thus the six cross-block entries of $h_v-I_4$ have derivative $1$ along the six cross-block coordinates of $B_v$. The Jacobian of the $6+6(|T|-1)$ cross-block equations with respect to these $6+6(|T|-1)$ directions is therefore block triangular with invertible diagonal blocks. By the formal implicit function theorem over $E$, imposing these equations eliminates all $6$ cross-block coordinates of each $B_v$ ($v\in T\setminus\{\ell\}$) and $6$ of the $9$ cross-block coordinates of $B_\ell$, leaving $(1,2,2)$ at $\ell$ and $(1,1,2)$ at each $v\in T\setminus\{\ell\}$ in first-degree-zero bidegrees, together with the three remaining cross-block coordinates $u:(1,0)$, $v:(-1,1)$, and $w:(1,1)$. A monomial $u^av^bw^c$ has first degree $a-b+c=0$ precisely when $b=a+c$, i.e.\ $u^av^bw^c=(uv)^a(wv)^c$; hence the first-degree-zero formal subring of $E[[u,v,w]]$ is $E[[uv,wv]]$, generated by the algebraically independent coordinates $uv$ of weight $1$ and $wv$ of weight $2$.
\end{proof}

\subsection*{Proof of \ref{B6.31}}\label{prf:B6.31}\phantomsection
\begin{proof}
Choose a good datum $\mathcal D$ by Lemma~\ref{B6.31d}. Corollary~\ref{B6.29a} permits its finite coefficient extension, since the normalized length is unchanged. Use the notation of Definitions~\ref{B6.31b} and~\ref{B6.31c}. Write $L(\mathcal D)$ for the length of the strict group with coefficients in $A_O(\Psi^{-1})$, with decomposition conditions at $S_{\mathcal D}$ and vanishing on finite inertia. At each added prime Frobenius acts by a root of unity times $p^{-k}$, so Frobenius minus the identity is surjective on the divisible coefficient module. The procyclic cohomology formula for the discrete torsion module and the unramified quotient generated topologically by Frobenius \cite[Prop.~1.7.7(i)]{neukirch-2008} gives zero unramified local cohomology there. Hence $L_{\ell,k,\theta}\le L(\mathcal D)$. The strict group has finite length by the finite-coefficient argument in \ref{B6.31a}.

Since $\Psi^{-1}(g_0)-1$ is a unit, every class $x$ in the strict group for $\mathcal D$ has a unique representative $z_x$ with $z_x(g_0)=0$. These representatives are additive in $x$. Choose $b_0\ne0$ killing the group. At each $v\in T$ choose $s_v\in I_v$ with $\Psi(s_v)\ne1$. Such an element exists by the infinite cyclotomic image at $\ell$ and by $\kappa_v\ne1$ elsewhere. Evaluation at $s_v$ shows that $b_0(\Psi^{-1}(s_v)-1)$ kills the local cobounding element. Hence $b=b_0\prod_{v\in T}(\Psi^{-1}(s_v)-1)$ puts the cocycles and local witnesses in $A_O[b]\simeq O/(b)$.

Use zero-based matrix indices and put $X_-=E_{20}+E_{31}$, where $E_{ij}$ is a matrix unit. Over $O\oplus\epsilon O/(b)$ with $\epsilon^2=0$ define
\[
 \rho_x(g)=(1+\epsilon z_x(g)X_-)\rho_0(g).
\]
Conjugation by $\rho_0$ acts on $X_-$ by $\Psi^{-1}$, so the cocycle identity gives the representation law. Use the local domains $B_\ell$ and $B_v$ ($v\in T\setminus\{\ell\}$) constructed in Lemma~\ref{B6.31e}. At each $v\in T$, lift the local witness to $O$; since the defining ideal of $B_v$ is invariant under $\mathrm{Sp}_4$-conjugation congruent to the identity by Lemma~\ref{B6.31e}(1)--(2), its infinitesimal conjugation gives an integral derivation on the local variables killing the equations of $B_v$. Lift the values on the new weight-$-1$ variables to $O$ and set the other values to zero. The derivation kills the equations modulo $b$ because it reduces to the deformation with its local witnesses. Its values on representation entries modulo $b$ are additive in $x$. The $(2,0)$ value is $\rho_{0,00}(g)z_x(g)$, with unit multiplier. The entry values therefore give the additive injection required in \ref{B6.58}, which yields $L(\mathcal D)\le\operatorname{length}_O C_{-1}(F)$ for the weight-$-1$ relation list $F$.

In the framed presentation of \ref{B6.58}, give entry $ij$ bidegree $(l_i-l_j,m_i-m_j)$, where $l=(1,0,1,0)$ and $m=(1,1,0,0)$. At $\ell$ use entries $(0,1),(0,3),(2,1),(1,0),(3,0),(1,2)$ of $\rho(\phi)$, with the frame at $\ell$ fixed to the identity. The derivative along the corresponding root vector is $\rho_{0,jj}(\phi)-\rho_{0,ii}(\phi)$. These differences are units by \eqref{eq:b6-sp4-separation}, since $\phi$ and $\sigma_\ell$ have the same residual diagonal. At the other $v\in T$ use the same six entries of the frame at $v$, with derivative $1$ in each corresponding direction. The six equations at each place therefore eliminate those directions.

The residual adjoint action in bidegree $j=(j_1,j_2)$ is scalar with character $\bar\alpha^{j_1}\bar\Psi^{j_2}$. The framed Euler characteristic gives
\[
 \dim\operatorname{Ob}_j-\dim H^1_{\rm fr,j}
 =2\,\mathbf1_{j=(0,0)}+\mathbf1_{j=(0,1)}+\mathbf1_{j=(0,-1)}.
\]
Here $H^1_{\rm fr,j}$ is the space of relative tangent classes with the local frame conditions. In the Euler characteristic the local terms cancel the contributions at the finite places, and the remaining term is the dimension fixed by complex conjugation. The weights of first degree zero contribute $2$ in degree $(0,0)$ and $1$ in each corner degree. In the cross-block roots the character is odd. The dual invariant term vanishes there by \eqref{eq:b6-sp4-separation} at $\sigma_\ell$, whose cyclotomic value is one, and in the even degrees it vanishes by evaluation at complex conjugation. In each cross-block root the relative tangent dimension is zero by the residual Selmer vanishing above.
\leanref{UpperTheory.\allowbreak{}V.\allowbreak{}GradeZeroInterpolatingElement.\allowbreak{}nums}{Development.\allowbreak{}AutomorphicArithmetic.\allowbreak{}AutomorphicVarieties.\allowbreak{}SymplecticGSp4.\allowbreak{}NumsZero}{47--168}

Write $b_T=|T|-1$ and let $a_-,a_+,a_0$ be the relative variable counts in bidegrees $(0,-1),(0,1),(0,0)$. By Lemma~\ref{B6.31e}(3), after eliminating the $6+6b_T$ cross-block directions, the local variables at $\ell$ contribute $(1,2,2)$ and each further place $v\in T\setminus\{\ell\}$ contributes $(1,1,2)$. Thus
\begin{equation}\label{eq:b6-sp4-counts}
 r=1+b_T+a_-,\qquad r'=2+b_T+a_+,\qquad m=2+2b_T+a_0.
\end{equation}
Here $m$ denotes the number of degree-zero variables. At each further place take the $02,20,00,11$ entries of its frame minus the identity, and take the $02,20$ entries of the marking matrix. Adding these relations to the obstruction counts gives $(r+1,m,r')$ in the three degrees. The three local coordinates with nonzero first degree have bidegrees $u:(1,0)$, $v:(-1,1)$ and $w:(1,1)$. A first-degree-zero monomial in these variables has exponent of $v$ equal to the sum of the exponents of $u,w$. Hence the first-degree-zero chart adds $uv$ of weight $1$ and $wv$ of weight $2$. In weights $(-2,-1,0,1,2)$ its coordinate counts are $(0,r,m,r'+1,1)$ and its relation counts are $(0,r+1,m,r',0)$. Lemma \ref{B6.31a} kills the negative-weight tangents modulo coboundaries, and unramified reciprocity kills the weight-zero tangents since $\operatorname{Cl}(\Q)=1$ and the strict local conditions remove the ramified directions \cite[Ch.~VI, Prop.~6.9]{neukirch-1999}. The marking removes the coboundary directions. Hence the nonpositive relation linear terms span the corresponding cotangent spaces over the fraction field.
\leanref{UpperTheory.\allowbreak{}V.\allowbreak{}GradeZeroInterpolatingElement.\allowbreak{}exModel}{Development.\allowbreak{}AutomorphicArithmetic.\allowbreak{}AutomorphicVarieties.\allowbreak{}SymplecticGSp4.\allowbreak{}ExModel}{263--268,304--332,425--453}
\leanref{UpperTheory.\allowbreak{}V.\allowbreak{}GradeZeroInterpolatingElement.\allowbreak{}aSmooth}{Development.\allowbreak{}AutomorphicArithmetic.\allowbreak{}AutomorphicVarieties.\allowbreak{}SymplecticGSp4.\allowbreak{}ASmooth}{1049--1135}

Let $x$ be a two-sided point of the first-degree-zero quotient. In the bigraded algebra generated by the homogeneous pieces of the representation model, projection onto first degree zero is linear over the first-degree-zero subalgebra. It splits the inclusion, so the fibre after base change by $x$ is nonzero. This algebra is of finite type over its first-degree-zero part, so the fibre has a $C_\ell$-point, say $\widetilde x$. By \ref{B6.49} the bidegree-zero image is finite over $O$. At $\widetilde x$, every nonconstant monomial of bidegree zero in the centered coordinates has absolute value less than one. Its value is integral by finiteness. If a nonzero value had integral inverse, reversing a monic equation for the inverse would give a polynomial in the monomial with constant term one that vanishes at $\widetilde x$. That polynomial is a unit in the degree-zero formal algebra.

There are therefore torus scalars making all nonzero centered coordinate values have absolute value less than one. In valuation terms, one seeks a linear functional on the two weights whose addition to the coordinate valuations makes them all positive. A failure of these strict inequalities would give a nonzero nonnegative rational relation among the weights with nonpositive sum of valuations. Clearing denominators would give a bidegree-zero monomial of nonpositive valuation. Choose a rational solution and torus scalars with those valuations. The scalars and the finitely many coordinate values lie in a finite extension of the coefficient field. Evaluation of the formal series in its complete ring of integers gives an allowed point $\rho$. Torus scaling preserves the two-sided condition and the bidegree-zero values.
\leanref{UpperTheory.\allowbreak{}V.\allowbreak{}GradeZeroInterpolatingElement.\allowbreak{}ptLift}{Development.\allowbreak{}AutomorphicArithmetic.\allowbreak{}AutomorphicVarieties.\allowbreak{}SymplecticGSp4.\allowbreak{}PtLift}{124--198}
\leanref{UpperTheory.\allowbreak{}V.\allowbreak{}GradeZeroInterpolatingElement.\allowbreak{}ptScale}{Development.\allowbreak{}AutomorphicArithmetic.\allowbreak{}AutomorphicVarieties.\allowbreak{}SymplecticGSp4.\allowbreak{}PtScale}{45--68}
\leanref{UpperTheory.\allowbreak{}V.\allowbreak{}Alg.\allowbreak{}points\_\allowbreak{}finite}{Development.\allowbreak{}AutomorphicArithmetic.\allowbreak{}AutomorphicVarieties.\allowbreak{}CommutativeAlgebra.\allowbreak{}PointsFinite}{310--379}

Write $i\to j$ if $\rho(g)_{ij}\ne0$ for some $g$. The span of the representation contains all four diagonal projections. Indeed, $\rho(\phi)$ is block diagonal in the blocks $\{0,2\},\{1,3\}$, and their characteristic polynomials are coprime by residual separation. Polynomial projections therefore give the two block idempotents. The marking matrix
\[
 M_\rho=\rho(g_0)+(\alpha\Psi)(g_0)\rho(g_0^{-1})
\]
is diagonal on the $\{0,2\}$ block, with reduced difference $(\bar\alpha(g_0)-1)(\bar\Psi(g_0)-1)$. It separates the two projections there. Symplectic adjunction in the span gives the projections at the reflected indices $i\mapsto3-i$. Multiplication by the projections turns a nonzero entry into a matrix unit. Thus $\to$ is transitive, and symplectic adjunction sends $i\to j$ to $3-j\to3-i$.

A mutually linked cross-block pair cannot form a component of size two. For such a pair, transitivity confines the multiplication law on its two-by-two subquotient to its two indices. The ordinary flag induces a stable local line in that subquotient. Orient the pair as $(i,j)$ so a vector on that line is $(a,1)$ with $|a|_\ell\le1$. In the integral model write $\rho(g)_{ij}=b h(g)$, where $b$ generates the nonzero ideal of the $ij$ entries. Dividing the multiplication law by $b$ and reducing gives the cocycle $\bar h(g)/\bar\rho_{0,jj}(g)$ for the residual ratio $\bar\rho_{0,ii}/\bar\rho_{0,jj}$. The stable line gives
\[
 b h(g)=a\bigl(\rho(g)_{jj}-\rho(g)_{ii}+a\rho(g)_{ji}\bigr)
 \qquad(g\in D_\ell).
\]
Compare this identity at $g$ and $\sigma_\ell$ by cross multiplication. After cancellation of $b$ and reduction, the last term in parentheses vanishes and the diagonal difference at $\sigma_\ell$ is a unit. This makes the residual cocycle a local coboundary at $\ell$. At $v\in T\setminus\{\ell\}$ the local first cohomology of this root is zero by local duality and the Euler characteristic, since \eqref{eq:b6-sp4-separation} and $\bar\varepsilon_\ell(\sigma_v)=1$ kill invariants and dual invariants. Elsewhere the cocycle vanishes on the prescribed inertia kernels. The residual Selmer vanishing makes it a global coboundary. Its value at $\phi$ is zero, and the cross-block ratio at $\phi$ differs from one. The cobounding scalar is therefore zero, contrary to the choice of $b$.

A mutual component of size at least three meets its symplectic reflection, and their union is all four indices. It then gives full matrix span. If each mutual component is a singleton, the support admits a triangular order. In the remaining case both within-block pairs are mutually linked by reflection, and transitivity permits cross-block entries in at most one direction. Full matrix span gives irreducibility. In the remaining case the points are block triangular.
\leanref{UpperTheory.\allowbreak{}V.\allowbreak{}GradeZeroInterpolatingElement.\allowbreak{}eii}{Development.\allowbreak{}AutomorphicArithmetic.\allowbreak{}AutomorphicVarieties.\allowbreak{}SymplecticGSp4.\allowbreak{}EiiUnit}{121--151,169--186}
\leanref{UpperTheory.\allowbreak{}V.\allowbreak{}GradeZeroInterpolatingElement.\allowbreak{}exists\_\allowbreak{}third\_\allowbreak{}index\_\allowbreak{}linked\_\allowbreak{}pair}{Development.\allowbreak{}AutomorphicArithmetic.\allowbreak{}AutomorphicVarieties.\allowbreak{}SymplecticGSp4.\allowbreak{}ExistsThirdIndexLinkedPair}{190--348}
\leanref{UpperTheory.\allowbreak{}V.\allowbreak{}GradeZeroInterpolatingElement.\allowbreak{}tri}{Development.\allowbreak{}AutomorphicArithmetic.\allowbreak{}AutomorphicVarieties.\allowbreak{}SymplecticGSp4.\allowbreak{}TriComb}{7--34}

In the triangular case compare the diagonal characters in the triangular order with those in the prescribed local ordinary flag. Equality of traces and separation at $s_0$ match the two lists by a permutation. The cross-block residual separation at $\sigma_\ell$ makes it preserve $\{0,2\}$ and $\{1,3\}$. The symplectic pairing, evaluated at $s_0$, makes it commute with $i\mapsto3-i$. The possibilities are the identity and $(0\,2)(1\,3)$. In either order, a cover $i\to j$ of the support order gives a two-dimensional extension, with cocycle $\rho(g)_{ij}/\rho(g)_{jj}$. If it runs against the prescribed order of the weights, its character is a negative cyclotomic power times a finite character, by the inertia comparison. The local ordinary flag splits that extension at $\ell$, and the prescribed finite inertia quotients give local coboundaries elsewhere. By \ref{B6.31a} it splits globally. The diagonal projections in the span force the cobounding scalar to vanish, since a coboundary would make the $ij$ entry a fixed multiple of the diagonal difference on the entire span. Transitivity expresses every strict support relation as a chain of covers. Hence the negative-weight entries vanish in the first order and the positive-weight entries vanish in the second. The respective $(0,0)$ values at $s_0$ are $(\alpha\Psi)(s_0)$ and $\alpha(s_0)$, which differ since $\Psi(s_0)=\varepsilon_\ell(s_0)^k$ has infinite order. This entry has weight zero, so a zero-weight fibre has at most one strict sign in the triangular case.
\leanref{UpperTheory.\allowbreak{}V.\allowbreak{}GradeZeroInterpolatingElement.\allowbreak{}diagChar}{Development.\allowbreak{}AutomorphicArithmetic.\allowbreak{}AutomorphicVarieties.\allowbreak{}SymplecticGSp4.\allowbreak{}DiagChar}{126--173}
\leanref{UpperTheory.\allowbreak{}V.\allowbreak{}GradeZeroInterpolatingElement.\allowbreak{}diagonal\_\allowbreak{}character\_\allowbreak{}order\_\allowbreak{}lt\_\allowbreak{}of\_\allowbreak{}stable\_\allowbreak{}flag}{Development.\allowbreak{}AutomorphicArithmetic.\allowbreak{}AutomorphicVarieties.\allowbreak{}SymplecticGSp4.\allowbreak{}ContentBoundLevel}{161--248}

In the remaining reducible case take the block-diagonal part of $\rho(g)$ for every $g$. All matrices are block triangular in the same direction, so taking their block-diagonal parts preserves multiplication. This point is the limit under the first torus action in the direction that kills the cross-block entries, which have first degrees of one sign. It retains the entries in positions $(0,2)$ and $(2,0)$ and obeys the point equations. Its corner span contains both diagonal projections and, by multiplying with them, both off-diagonal units in the $\{0,2\}$ block. This contradicts the block assertion of \ref{B6.49}. Thus two-sided points are irreducible.
\leanref{UpperTheory.\allowbreak{}V.\allowbreak{}GradeZeroInterpolatingElement.\allowbreak{}twoIrr}{Development.\allowbreak{}AutomorphicArithmetic.\allowbreak{}AutomorphicVarieties.\allowbreak{}SymplecticGSp4.\allowbreak{}TwoIrr}{42--66}

Use the isomorphism $\iota:C_\ell\simeq\C$ supplied by \ref{B6.49} for $\mathcal D$. Each irreducible allowed point $\rho$ is continuous. For every imaginary quadratic field $K$ such that $\rho|_{G_K}$ is irreducible, Theorem~\ref{B6.49} gives $\rho|_{G_K}$ a cuspidal attachment with witness in the sense of Definition~\ref{def:b6-cuspidal-witness}, for this same $\iota$. Its multiplier is $\varepsilon_\ell^{k+1}\theta\xi\omega^{-1}$. Put $\mathcal N=\ell2^{n'}Np'$ and $w=(k+1,k,1,0)$. Away from $\ell$, inertia is diagonal on successive quotients with characters factoring through $\mu_{\mathcal N}$. At $\ell$, let $\lambda_i$ for $i=0,1,2,3$ be the ordered ordinary flag characters. Their restrictions to inertia agree with the corresponding diagonal characters of $\rho_0$, and their reductions agree with those characters on $D_\ell$. In particular the inertial characters are cyclotomic powers of weights $w_i$ times finite characters factoring through $\mu_{\mathcal N}$.

For every open subgroup $U\le D_\ell$ containing $\sigma_\ell$ and all $i,j$ we have
\begin{equation}\label{eq:b6-sp4-unlinking}
 \lambda_i|_U\ne(\lambda_j\varepsilon_\ell)|_U .
\end{equation}
Since $U$ has finite index in the profinite group $D_\ell$, some positive power $s_0^m$ belongs to $U$. On $s_0^m$ the four characters are the powers of $\varepsilon_\ell(s_0^m)$ with exponents $w_i$. This cyclotomic value has infinite order. An equality on $U$ would therefore give $w_i=w_j+1$. This excludes $i=j$ and puts $i,j$ in different sets of the pair $\{0,2\},\{1,3\}$, since the two weights within either set differ by $k\ge2$. At $\sigma_\ell$ both the cyclotomic and Teichmüller values are one. The ratio of the two residual diagonal characters from different sets is, up to inversion, $\overline{\xi\Psi^a}(\sigma_\ell)$ for some $a\in\{-1,0,1\}$. This differs from one by \eqref{eq:b6-sp4-separation}, whereas reduction of the alleged equality on $U$ makes the ratio one. Taking $U=D_\ell$ proves the unlinking required for level $(\mathcal N;w)$.
\leanref{UpperTheory.\allowbreak{}V.\allowbreak{}Loc.\allowbreak{}spPtL}{Development.\allowbreak{}AutomorphicArithmetic.\allowbreak{}AutomorphicVarieties.\allowbreak{}PadicHodgeTheory.\allowbreak{}SplAdjUnit}{264--331}
\leanref{UpperTheory.\allowbreak{}V.\allowbreak{}GradeZeroInterpolatingElement.\allowbreak{}level}{Development.\allowbreak{}AutomorphicArithmetic.\allowbreak{}AutomorphicVarieties.\allowbreak{}SymplecticGSp4.\allowbreak{}LevBasis}{116--189}

To lift a tangent at $x$, start with a point $\widetilde x$ above $x$. Torus scaling is conjugation, so the irreducibility of the scaled point gives irreducibility at $\widetilde x$. Choose entries $u=\rho(g)_{01}$ and $v=\rho(h)_{10}$ in the bigraded algebra whose values at $\widetilde x$ are nonzero. Such entries exist by full matrix span. Their first degrees are $1,-1$, so $uv$ belongs to the first-degree-zero algebra and has unit value at the dual-number point. After localizing at $uv$, both $u$ and $v$ are units. Multiplication by their powers moves each homogeneous element to first degree zero. In particular the localization is a Laurent polynomial algebra in $u$ over its first-degree-zero part. Give $u$ its value at $\widetilde x$ and extend the dual-number point to this algebra. The torus scalars used to form the allowed base point also make the scalar parts of the lifted centered coordinates topologically nilpotent. The finitely many coordinate values lie in a finite coefficient field with dual numbers, where the formal series evaluate. This gives an allowed lift.

By \ref{B6.48}, the lifted tangent acts on the representation matrices by a commutator with $X\in\mathfrak{sp}_4(C_\ell)$. The block condition at $\phi$ forces $X$ to be block diagonal, and the two marked corners of $[X,M_\rho]$ force $X_{02}=X_{20}=0$ by the distinct marked diagonal values. The symplectic Lie algebra relation kills $X_{13},X_{31}$. Its diagonal then acts on bidegree $(a,b)$ by a linear function of $(a,b)$. On the first-degree-zero algebra the tangent has the form $y\mapsto cj x(y)$ on elements $y$ of weight $j$. This gives the tangent dimension required in \ref{B6.58}.
\leanref{UpperTheory.\allowbreak{}V.\allowbreak{}GradeZeroInterpolatingElement.\allowbreak{}m4}{Development.\allowbreak{}AutomorphicArithmetic.\allowbreak{}AutomorphicVarieties.\allowbreak{}SymplecticGSp4.\allowbreak{}ExistsLiftOfDerivationOnIdeal}{131--174,208--237}
\leanref{UpperTheory.\allowbreak{}V.\allowbreak{}GradeZeroInterpolatingElement.\allowbreak{}rigidDiag}{Development.\allowbreak{}AutomorphicArithmetic.\allowbreak{}AutomorphicVarieties.\allowbreak{}SymplecticGSp4.\allowbreak{}RigidDiag}{99--172}

An intertwiner $U$ of two allowed irreducible points is block diagonal by the disjoint residual spectra at $\phi$. On the $\{0,2\}$ block the marking matrices are diagonal, and their two corresponding cross differences are units by the same reduced difference computed above. Intertwining the marking matrices gives $u_{02}=u_{20}=0$. The representations have the same multiplier. If $J$ is the symplectic form, $U^tJU$ and $J$ give two forms with that multiplier on the source representation. Schur's lemma gives $U^tJU=cJ$ with $c\ne0$. Since $U$ is diagonal and invertible on $\{0,2\}$, the symplectic equation forces $u_{13}=u_{31}=0$ and $u_{00}u_{33}=u_{11}u_{22}$. Hence $U$ acts by the marked torus on the entries, which determine the points. In first degree zero this is the graded torus of \ref{B6.58}.
\leanref{UpperTheory.\allowbreak{}V.\allowbreak{}GradeZeroInterpolatingElement.\allowbreak{}conjTorus}{Development.\allowbreak{}AutomorphicArithmetic.\allowbreak{}AutomorphicVarieties.\allowbreak{}SymplecticGSp4.\allowbreak{}ConjTorus}{208--245}

Take a prime $p_0\mid3\ell Np'+1$. Then $p_0\ge5$ and $p_0\nmid\mathcal N$. Let $\sigma$ be arithmetic Frobenius at $p_0$. The trace of $\rho_0$ there is $p_0^{k+1}z_1+p_0^kz_2+p_0z_3+1$ with $z_i^h=1$ for $h=(\ell-1)2^{n'}(2o)$, by \eqref{eq:b6-host-orders} and the orders of $\omega$ and $\xi$. Apply \ref{B6.50} with roots $\iota(z_i)$ to singleton families. The sum of the weight-zero lifts of the four diagonal entries at $\sigma$ minus this distinguished trace vanishes at $\rho_0$ and is nonzero at every two-sided point. Apply the orbit estimate of \ref{B6.58}. Take one representative per orbit. Their representations form a pairwise nonisomorphic family $\Xi$. Since $|\ell|_\ell=\ell^{-1}$, the orbit estimate gives
\begin{equation}\label{eq:b6-sp4-length}
 \frac{L(\mathcal D)}{e_O}\log\ell\le
 2\sum_{\rho\in\Xi}
 -\log|\operatorname{tr}\rho(\sigma)-\operatorname{tr}\rho_0(\sigma)|_\ell .
\end{equation}
Apply \ref{B6.50} to $\Xi$ with the same $\iota$. Its nonzero integer $D$ satisfies $v_\ell(D)\log\ell\le\log|D|$. Thus the bound is $2ch[\mathcal N(k+1)]^A\log p_0$ with the constants $c,A$ of that theorem. Equations \eqref{eq:b6-sp4-levels} and \eqref{eq:b6-host-orders}, together with $2^{c_2(\theta)}\mid N$, $p'\le p_1$ and $p_0\le3\ell Np'+1$, bound this by a fixed power of $\ell Nk$. Corollary \ref{B6.29a} returns the bound to the original ring.
\end{proof}

\subsection*{Proof of \ref{B6.32a}}\label{prf:B6.32a}\phantomsection
\begin{proof}
Put $G=G_\Q$, $U=G_{L_1}$ and $M=M_{\ell,k,\theta}$ with scalar action $\lambda$. Inflation--restriction for the closed normal subgroup $U$ \cite[Prop.~1.6.7]{neukirch-2008} represents the restriction kernel by cocycles $Z$ vanishing on $U$. Evaluation on $G/U$ embeds $Z$ into $M^n$. Since the quotient is abelian,
\[
 (\lambda(g)-1)z(h)=(\lambda(h)-1)z(g).
\]
Thus multiplication by any $\lambda(g)-1$ sends these cocycles to coboundaries. If $a$ generates the ideal of these differences, the restriction kernel is a quotient of $Z/aZ$. Lift generators of any finitely generated submodule of $Z/aZ$ and let $W$ be their span in $Z$. It has finite length, and its image has length at most $\operatorname{length}_O(W/aW)=\operatorname{length}_O W[a]$. Consequently $\operatorname{length}_O(Z/aZ)\le\operatorname{length}_O Z[a]$, and $Z[a]$ embeds into $(M^G)^n$.

Lift $1+\ell^s$ to inertia under $\varepsilon_\ell$, with $s=\max\{v_\ell(N),1\}$. This lift lies in $\ker\theta$, and its scalar minus $1$ has valuation $s+v_\ell(k)$ by \eqref{eq:B6-lte}. If $\ell\nmid N$ and $\ell-1\nmid k$, a lift of a residual cyclotomic generator gives a unit scalar difference. This bounds the inertia invariants. The restriction image lies in $O\otimes_{\Z_\ell}\Sha_{\ell,k}(L_1)$ since the local conditions persist and the twist becomes $\varepsilon_\ell^{-k}$. This gives \eqref{eq:b6-transfer}, with $\nu_\ell(N,k)\le N+k\le2Nk$.
\end{proof}

\subsection*{Proof of \ref{B6.34}}\label{prf:B6.34}\phantomsection
\begin{proof}
Put $K=\Q(i)$, $\Psi=\theta\varepsilon_\ell$, and $c=c_2(\theta)$. Let $\eta_4$ be the quadratic character of $K$, let $\mathbf k=O/\mathfrak m_O$ be the residue field of $O$, and set
\[
 S=\{v:v\mid\ell N\},\qquad
 S_b=\{v\mid N:v\ne\ell,\ \kappa_v\ne1\},\qquad
 T_b=\{\ell\}\cup S_b .
\]
The odd local orders divide $2o$ by \ref{B6.28a}, and the primes in $T_b$ divide $\ell f$. The class V threshold gives
\begin{equation}\label{eq:b6-v-threshold}
 8ov^2<2^c\quad(v\in T_b),\qquad c\ge4,\qquad
 \ell f\,4^{n_0}\le2^{c-1}.
\end{equation}
For each $v\in T_b$ choose an arithmetic Frobenius lift $\tau\in D_v$, with $\varepsilon_\ell(\tau)=1$ at $\ell$. At each $v\in S_b$, the nontrivial $\ell$-power part $\kappa_v$ forces $v\equiv1\bmod\ell$; since $\ell$ is odd, every $v\in T_b$ is odd, so $\eta_4$ is unramified at $v$ and $\eta_4(I_v)=1$. Put $x_v=\tau^2$. If $\bar\Psi(x_v)^2=1$, raising the equality to $o$ kills the odd-modulus factors, whose orders divide $2o$ and which occur to the fourth power. Injectivity of reduction on dyadic roots then gives $\theta_2(v)^{4o}=1$. Faithfulness on units congruent to $1\bmod4$ would give $v_2(v^{4o}-1)\ge c$. But \eqref{eq:B6-lte} and \eqref{eq:b6-v-threshold} give
\[
 v_2(v^{4o}-1)=v_2(v^2-1)+1+v_2(o)<c .
\]
Hence
\begin{equation}\label{eq:b6-v-split}
 \eta_4(x_v)=\bar\varepsilon_\ell(x_v)=1,\qquad
 \bar\Psi(x_v)^2\ne1 .
\end{equation}
\leanref{UpperTheory.\allowbreak{}V.\allowbreak{}UnitaryThreeInertia.\allowbreak{}hostData\_\allowbreak{}of\_\allowbreak{}host}{Development.\allowbreak{}AutomorphicArithmetic.\allowbreak{}AutomorphicVarieties.\allowbreak{}UnitaryThreeDeformations.\allowbreak{}HostDataOf}{9--34}
\leanref{UpperTheory.\allowbreak{}V.\allowbreak{}UnitaryThreeInertia.\allowbreak{}exists\_\allowbreak{}sigma}{Development.\allowbreak{}AutomorphicArithmetic.\allowbreak{}AutomorphicVarieties.\allowbreak{}UnitaryThreeDeformations.\allowbreak{}DyadicCharacterPowerNeOne}{202--243}

Let $G=\Gamma_{\Q,S}$ be the Galois quotient unramified outside $S$, and let $F$ be the Gaussian frame on $G$ described before \ref{B6.51}, with distinguished lift $r_0$ having sign $\eta_4$ and matrix part $\rho_0=\operatorname{diag}(\Psi,1,\Psi^{-1})$, and write $\bar r$ for the reduction of $r_0$. For each $v\in S\setminus\{\ell\}$ put
\[
 N_v=I_v\cap\ker\eta_4\cap\ker\bar\Psi\le D_v,
\]
where $\bar\Psi:G\to\mathbf k^\times$ is the residual character. At $v\in S\setminus T_b$ we have $\kappa_v=1$, so $\bar\Psi(s)=1$ on $I_v$ implies $\Psi(s)=1$, and therefore $r_0$ kills $N_v$. Global finiteness of first cohomology for finite $\F_\ell$-representations of $G$ gives finiteness of the set of continuous $\F_\ell$-characters on every open subgroup of $G$. Hence the unrestricted lifting problem for $\bar r$ on $G$ admits a universal complete Noetherian local $O$-algebra with universal lift $u_{\mathrm{unr}}$. Fix $\iota:C_\ell\simeq\C$ and choose a marking $g_1$ from \ref{B6.51} and \ref{B6.52}. At each $v\in T_b$, Hensel's lemma and \eqref{eq:b6-v-split} give orthogonal spectral projectors $P_{v,0},P_{v,1},P_{v,2}$ for the matrix of $u_{\mathrm{unr}}(x_v)$ reducing to $E_{00},E_{11},E_{22}$, expressing the local condition of $F$ at $v$ by the closed ideal of relations $P_{v,j}u_{\mathrm{unr}}(g)_{\mathrm{mat}}P_{v,i}=0$ ($i<j$ at $\ell$ and $i\ne j$ at $v\in S_b$) on $D_v\cap\ker\eta_4$ and $P_{v,i}(u_{\mathrm{unr}}(s)_{\mathrm{mat}}-\Psi(s)^{1-i})P_{v,i}=0$ on $I_v\cap\ker\eta_4$. Cutting out the closed ideal defined by the marking, triviality on $N_v$ for $v\in S\setminus T_b$, and these projector relations at $T_b$ yields the universal complete Noetherian local $O$-algebra $R$ with universal lift $u:G\to\GL_3(R)\rtimes\{1,-1\}$. Since $r_0$ satisfies all frame conditions, $R$ is nontrivial and admits a base homomorphism $R\to O$ specializing $u$ to $r_0$.
\leanref{UpperTheory.\allowbreak{}V.\allowbreak{}UnitaryThreeInertia.\allowbreak{}exists\_\allowbreak{}hostFrame}{Development.\allowbreak{}AutomorphicArithmetic.\allowbreak{}AutomorphicVarieties.\allowbreak{}UnitaryThreeDeformations.\allowbreak{}HostEx}{55--114}
\leanref{UpperTheory.\allowbreak{}V.\allowbreak{}UnitaryThreeInertia.\allowbreak{}hostFrame\_\allowbreak{}ok}{Development.\allowbreak{}AutomorphicArithmetic.\allowbreak{}AutomorphicVarieties.\allowbreak{}UnitaryThreeDeformations.\allowbreak{}HostOk}{239--330}
\leanref{UpperTheory.\allowbreak{}V.\allowbreak{}UnitaryThreeInertia.\allowbreak{}hostFrame\_\allowbreak{}phi}{Development.\allowbreak{}AutomorphicArithmetic.\allowbreak{}AutomorphicVarieties.\allowbreak{}UnitaryThreeDeformations.\allowbreak{}HostPhi}{182--216}
\leanref{UpperTheory.\allowbreak{}V.\allowbreak{}UnitaryThreeInertia.\allowbreak{}exists\_\allowbreak{}isUniv}{Development.\allowbreak{}AutomorphicArithmetic.\allowbreak{}AutomorphicVarieties.\allowbreak{}UnitaryThreeDeformations.\allowbreak{}ExistsIsUniv}{497--512}
\leanref{UpperTheory.\allowbreak{}V.\allowbreak{}UnitaryThreeInertia.\allowbreak{}isPt\_\allowbreak{}rz}{Development.\allowbreak{}AutomorphicArithmetic.\allowbreak{}AutomorphicVarieties.\allowbreak{}UnitaryThreeDeformations.\allowbreak{}MapDiagonalRepMatrixEq}{35--105}

By \ref{B6.51}, the weight-zero algebra $R_0$ of $u$ is finite over $O$. A Type I point has an irreducible matrix restriction $r$ to $G_K$, with attachment and witness for $r\varepsilon_\ell^{-1}$. Its integral ordinary flag has inertia characters $\Psi,1,\Psi^{-1}$ of weights $(1,0,-1)$ and the corresponding residual characters on the full ordinary group. At $x_\ell$ the cyclotomic value is one, and the three residual values $\bar\Psi(x_\ell),1,\bar\Psi(x_\ell)^{-1}$ are distinct. This separates the flag characters and excludes an equality between two of them after a single cyclotomic twist. For equal indices a nontrivial cyclotomic inertia value excludes that equality. Away from $\ell$, the cyclotomic character is unramified; on $I_v\cap G_K$ the inertia action is conjugate to $\operatorname{diag}(\theta,1,\theta^{-1})$, directly from the frame condition at $v\in S_b$ and by averaging $r(s)\rho_0(s)^{-1}$ over the prime-to-$\ell$ quotient $(I_v\cap G_K)/N_v$ at $v\in S\setminus T_b$ (and unramifiedness outside $S$). These characters factor through $\mu_N$. Thus $r$ belongs to the rank-three family of \ref{B6.41}.
\leanref{UpperTheory.\allowbreak{}V.\allowbreak{}UnitaryThreeInertia.\allowbreak{}levelAway\_\allowbreak{}of\_\allowbreak{}pt}{Development.\allowbreak{}AutomorphicArithmetic.\allowbreak{}AutomorphicVarieties.\allowbreak{}UnitaryThreeDeformations.\allowbreak{}LevelAway}{176--233}

Pass to the linked CM frame of \ref{B6.51a}. The determinant law and the transported ordinary characters give the arithmetic point with its prescribed flag and automorphic witness. The separated flags and local comparison for the attachment give the infinite-order test at each ordinary place. Take a Weil element $s$ of positive degree $d_0$ in the kernel of the cyclotomic character. Local reciprocity gives one by taking the class of a rational $\ell$-power in the local multiplicative group, with the direction chosen for positive degree. If $\kappa_i$ are the flag characters of the automorphically normalized representation with weights $m_i$, the toral support characters have moduli $|\iota(\kappa_i(s))|q_v^{m_i d_0}$ there, where $q_v$ is the residue cardinality. By \eqref{eq:p6c5-toral-support} the ratio of any two of these moduli lies strictly between $q_v^{-d_0}$ and $q_v^{d_0}$. If $\kappa_i(s)/\kappa_j(s)$ were a root of unity, then $|m_i-m_j|<1$, impossible for distinct integer weights. The common twist cancels from these ratios. Taking a positive power to enter a linked local subgroup preserves the test.
\leanref{UpperTheory.\allowbreak{}V.\allowbreak{}Z.\allowbreak{}J.\allowbreak{}hostUnlE\_\allowbreak{}alt}{Development.\allowbreak{}AutomorphicArithmetic.\allowbreak{}AutomorphicVarieties.\allowbreak{}ModularityArgument.\allowbreak{}ThresholdArguments.\allowbreak{}DenseThresholdCases.\allowbreak{}EHostUnl}{109--137}
\leanref{UpperTheory.\allowbreak{}V.\allowbreak{}Z.\allowbreak{}Win.\allowbreak{}RR.\allowbreak{}free\_\allowbreak{}core\_\allowbreak{}alt}{Development.\allowbreak{}AutomorphicArithmetic.\allowbreak{}AutomorphicVarieties.\allowbreak{}ModularityArgument.\allowbreak{}ThresholdArguments.\allowbreak{}DenseThresholdCases.\allowbreak{}RRFree}{62--116}

Write $\Gamma'$ for the cyclotomic kernel of the ordinary frame of \ref{B6.53}. The rank-three irreducibility test applies to the intersection of $\Gamma'$ with the normal core of any open subgroup of the matrix group domain. A positive power of the local element just constructed lies in this intersection and has three distinct eigenvalues with pairwise ratios of infinite order. By the prime-rank restriction criterion, a reducible restriction to this normal subgroup would be scalar or a sum of three lines permuted transitively by the full group. The distinct eigenvalues exclude the scalar case. In the line case the ordinary flag is a flag of sums of these eigenlines, so the local group fixes the lines individually. On the kernel of the permutation action, the three line characters reduce, up to a permutation, to the restrictions of the three global residual characters. This follows by comparing characteristic polynomials and their semisimplifications. Transitivity and conjugation invariance of the global residual characters would make all three reduced line characters equal on the local group, contrary to the transported residual separation. Thus the intersection is irreducible, and in particular $\Gamma'$ acts irreducibly.

Each open coset in $\Gamma'$ contains an element with separable characteristic polynomial. Take a coset $yH$ with $H$ the intersection of $\Gamma'$ with an open subgroup. Diagonalize a power of the local test element in $H$ and write its matrix as $D=\operatorname{diag}(\lambda_1,\lambda_2,\lambda_3)$. Every finite-index open subgroup used to form $H$ acts irreducibly on intersection with $\Gamma'$. The same holds for the preimage of the identity component of the Zariski closure of the image of $H$, since this preimage is closed and of finite index in $H$. Irreducibility and density in that component give an element $b\in H$ whose matrix and inverse have all entries nonzero. In the eigenbasis of $D$, the diagonal entries of $r(y)D^m r(b)$ and of its inverse are nonzero for some $m$. Each such entry is an exponential sum in the $\lambda_i^m$ or $\lambda_i^{-m}$ with a nonzero coefficient. The product of these sums is nonzero for some $m$, by Vandermonde after grouping the bases and using that the ratios have infinite order. In rank three the inverse diagonal entries give the principal minors of size two. For a matrix $A$ with all these principal minors and diagonal entries nonzero, the discriminant of $AD^{m'}$ is also a nonzero exponential sum. Order the $\lambda_i$ by a real linear functional on their multiplicative group modulo torsion, with values $s_1>s_2>s_3$. In the discriminant the highest term comes from the square of the trace times the square of the second elementary coefficient, at weight $4s_1+2s_2$. Its coefficient is $(A_{11}\det A_{\{1,2\},\{1,2\}})^2$. All other discriminant terms have smaller weight. Vandermonde gives a power $m'$ with nonzero discriminant. The resulting matrix is still in the image of the coset.
\leanref{UpperTheory.\allowbreak{}V.\allowbreak{}Th.\allowbreak{}nice\_\allowbreak{}zero\_\allowbreak{}of\_\allowbreak{}rank\_\allowbreak{}two\_\allowbreak{}or\_\allowbreak{}three\_\allowbreak{}irreducible\_\allowbreak{}free}{Development.\allowbreak{}AutomorphicArithmetic.\allowbreak{}AutomorphicVarieties.\allowbreak{}ModularityArgument.\allowbreak{}ThresholdArguments.\allowbreak{}DenseThresholdCases.\allowbreak{}TNiceSmall}{24--43}
\leanref{UpperTheory.\allowbreak{}V.\allowbreak{}Pr.\allowbreak{}regCoset}{Development.\allowbreak{}AutomorphicArithmetic.\allowbreak{}AutomorphicVarieties.\allowbreak{}CurveDVRSpecialization.\allowbreak{}RegCoset}{71--162}

In the ordinary basis at $\ell$, a first-order lift is upper triangular with fixed inertia diagonals. The comparison \eqref{eq:p6c5-local-finite}, applied successively to the positive-weight Hom spaces on the subgroup killing the finite inertial characters, puts the adjoint cocycle in the finite local group. Away from $\ell$ it is an inertia coboundary by averaging over the finite inertia quotient. The flag weights and the cyclotomic-kernel properties give the hypotheses of \ref{B6.53}. The cocycle is a coboundary on the linked frame and hence on the original group by averaging over the finite quotient. Indeed a restriction coboundary can first be subtracted on the whole group. The remaining cocycle factors through the quotient with values in the subgroup invariants, and division by the quotient order gives a coboundary there.

Choose a prime divisor $p\equiv3\bmod4$ of $12\ell N-1$, which exists since this integer is $3\bmod4$. Then $p\nmid6\ell N$ and $7\le p\le12\ell N$. Take $\sigma$ to be inverse Frobenius at the Gaussian prime of norm $p^2$. The distinguished trace is $\iota^{-1}((p^2)^{-1}z+1+p^2z^{-1})$ with $z=\iota(\theta(\sigma))$ and $z^{\varphi(N)}=1$. Apply \ref{B6.41} to singletons using the indexed family of all $\varphi(N)$-th roots of unity. Every Type I trace difference at $\sigma$ is nonzero.

Let $V_j$ be the weight-$j$ residual adjoint space over $\mathbf k$. For every $v\in S\setminus\{\ell\}$ (including $v\in S_b$), let $Z^1_{\mathcal L,v}$ be the space of continuous $1$-cocycles on $D_v$ with values in $V_j$ vanishing on $N_v=I_v\cap\ker\eta_4\cap\ker\bar\Psi$. At $v=\ell$, take $Z^1_{\mathcal L,\ell}=B^1(D_\ell,V_j)$ for $j<0$, $Z^1_{\mathcal L,\ell}=Z^1_{\mathrm{ur}}(D_\ell,V_0)$ (cocycles vanishing on $I_\ell$) for $j=0$, and $Z^1_{\mathcal L,\ell}=Z^1(D_\ell,V_j)$ for $j>0$. Since $N_v$ acts trivially on $V_j$ and $I_\ell$ acts trivially on $V_0$ (as $\eta_4(I_\ell)=1$), every global coboundary in $B^1(G,V_j)$ restricts into $Z^1_{\mathcal L,v}$ for all $v\in S$. Let $Z^1_{\mathcal L,j}\subseteq Z^1(G,V_j)$ be the space of global cocycles restricting into $Z^1_{\mathcal L,v}$ at every $v\in S$, and let $H^1_{\mathcal L,j}=Z^1_{\mathcal L,j}/B^1(G,V_j)$. The obstruction space $\operatorname{Ob}_j$ consists of pairs $(c_2,(u_v)_{v\in S})$ with $c_2\in Z^2(G,V_j)$ and $du_v=c_2|_{D_v}$, modulo $(da,(a|_{D_v})_v)$ for global $1$-cochains $a$ and $(0,(z_v)_v)$ for $z_v\in Z^1_{\mathcal L,v}$.

For each $v\in T_b$ and $j\ne0$, the element $x_v\in D_v$ of \eqref{eq:b6-v-split} satisfies $\eta_4(x_v)=\bar\varepsilon_\ell(x_v)=1$ and $\bar\Psi(x_v)^2\ne1$. On $V_j$ (which is zero for $|j|>2$), $x_v$ acts by the scalar $\bar\Psi(x_v)^j\ne1$, and on the Tate dual $V_j^\vee(1)$ it acts by $\bar\Psi(x_v)^{-j}\ne1$. Hence $H^0(D_v,V_j)=0$ and, by local Tate duality, $H^2(D_v,V_j)=0$. For $v\in S_b$ (where $v\ne\ell$), the local Euler characteristic formula gives $h^1(D_v,V_j)=h^0(D_v,V_j)+h^2(D_v,V_j)=0$, so $H^1(D_v,V_j)=0$. In particular, at $v\in S_b$ we have $Z^1_{\mathcal L,v}=Z^1(D_v,V_j)=B^1(D_v,V_j)$ for $j\ne0$; every cocycle in $Z^1_{\mathcal L,v}$ is therefore a coboundary on $D_v$ and produces a first-order lift locally conjugate to $\bar r$, matching the diagonalizability and inertia-diagonal condition of $F$. At $v\in S_b$ and $j=0$, since $v$ is odd we have $\eta_4(I_v)=1$, so $I_v$ acts trivially on the diagonal space $V_0$; because $[I_v:N_v]$ divides $2(\#\mathbf k-1)$ and is prime to $\ell$, vanishing on $N_v$ is equivalent to vanishing on $I_v$, giving $Z^1_{\mathcal L,v}=Z^1_{\mathrm{ur}}(D_v,V_0)$, which again matches the weight-zero frame condition at $v\in S_b$.
\leanref{UpperTheory.\allowbreak{}V.\allowbreak{}UnitaryThreeInertia.\allowbreak{}hostFrame\_\allowbreak{}locVan}{Development.\allowbreak{}AutomorphicArithmetic.\allowbreak{}AutomorphicVarieties.\allowbreak{}UnitaryThreeDeformations.\allowbreak{}HostLocVan}{155--201}

The primitive factor inducing $(\omega\theta_1)^2$ has odd conductor $f'\mid\ell f$, since the odd factors outside $f$ have square one by \ref{B6.28a}. Its order divides $\ell^{f_0}-1$. Squaring $\theta_2$ halves its order on the cyclic units congruent to $1\bmod4$, on which $\eta_4$ is trivial. Thus $\theta_2^2\eta_4$ has conductor $2^{c-1}$. Its product with the odd-modulus factor is odd and reduces to $\eta_4\bar\Psi^2$. The bounds in \eqref{eq:b6-v-threshold} give $f'4^{n_0}\le2^{c-1}$, and $4\mid\ell^{f_0}-1$ by the definition of $f_0$. The subgroup $\ker(\eta_4,\bar\Psi)\le G_\Q$ has index dividing $2(\#\mathbf k-1)$, prime to $\ell$. For every $v\in S\setminus\{\ell\}$, a cocycle in $H^1_{\mathcal L,-2}$ vanishes on $N_v=I_v\cap\ker(\eta_4,\bar\Psi)$, and at $v=\ell$ it is a local coboundary, hence vanishes on inertia in the character kernel. The inverse-character assertion of \ref{B6.59} gives $H^1_{\mathcal L,-2}=0$.
\leanref{UpperTheory.\allowbreak{}V.\allowbreak{}X.\allowbreak{}dirData}{Development.\allowbreak{}AutomorphicArithmetic.\allowbreak{}AutomorphicVarieties.\allowbreak{}LevelCohomology.\allowbreak{}HcDirData}{29--110}
\leanref{UpperTheory.\allowbreak{}V.\allowbreak{}X.\allowbreak{}xpp}{Development.\allowbreak{}AutomorphicArithmetic.\allowbreak{}AutomorphicVarieties.\allowbreak{}LevelCohomology.\allowbreak{}Xpp}{11--85}
\leanref{UpperTheory.\allowbreak{}V.\allowbreak{}X.\allowbreak{}xppOf}{Development.\allowbreak{}AutomorphicArithmetic.\allowbreak{}AutomorphicVarieties.\allowbreak{}LevelCohomology.\allowbreak{}HcXppOf}{60--99}

In nonzero weight $x_\ell$ kills global dual invariants by \eqref{eq:b6-v-split}. In weight zero take an element fixing $i$ with $\bar\varepsilon_\ell\ne1$, which exists since $\Q(i)$ and $\Q(\mu_\ell)$ are disjoint. It acts trivially on $V_0$ and kills the dual invariants. Thus Poitou--Tate in \ref{B6.09} gives surjectivity in degree two onto the local groups. The obstruction exact sequence then gives
\[
 e_j:=\dim H^1_{\mathcal L,j}-\dim\operatorname{Ob}_j
 =h^1(G,V_j)-h^2(G,V_j)
  +\sum_{v\in S}\bigl(h^2(D_v,V_j)-\dim Z^1(D_v,V_j)
                         +\dim Z^1_{\mathcal L,v}\bigr).
\]
Write $d_j=\dim_{\mathbf k} V_j$ and $d_j^+=\dim_{\mathbf k} V_j^{c_\infty=1}$, where $c_\infty$ is complex conjugation. The global Euler characteristic gives $h^1-h^2=h^0+d_j-d_j^+$. For every $v\in S\setminus\{\ell\}$ (including $v\in S_b$), the homomorphism $(\eta_4,\bar\Psi):I_v\to\{\pm1\}\times\mathbf k^\times$ has kernel $N_v$, so $[I_v:N_v]$ divides $2(\#\mathbf k-1)$ and is prime to $\ell$, while $N_v$ acts trivially on $V_j$. Restriction $Z^1_{\mathcal L,v}\to Z^1(I_v,V_j)$ maps into $B^1(I_v,V_j)$ because $H^1(I_v/N_v,V_j)=0$, and surjects onto $B^1(I_v,V_j)$ because $B^1(D_v,V_j)\subseteq Z^1_{\mathcal L,v}$; its image has dimension $d_j-\dim V_j^{I_v}$. Its kernel is $Z^1_{\mathrm{ur}}(D_v,V_j)=Z^1(D_v/I_v,V_j^{I_v})$, which evaluation at a Frobenius generator of $D_v/I_v$ identifies with $V_j^{I_v}$, of dimension $\dim V_j^{I_v}$. Hence $\dim Z^1_{\mathcal L,v}=d_j$. The local Euler formula gives $\dim Z^1(D_v,V_j)=d_j+h^2(D_v,V_j)$, so these terms cancel.

At $\ell$ the local Euler formula gives $\dim Z^1(D_\ell,V_j)=2d_j+h^2(D_\ell,V_j)$. In nonzero weight $x_\ell$ gives $h^0(G,V_j)=h^0(D_\ell,V_j)=h^2(D_\ell,V_j)=0$. The negative local condition then has dimension $d_j$ and the positive condition has dimension $2d_j$. In weight zero the local condition has dimension $d_0$ because $V_0$ is unramified at $\ell$. The action on $V_0$ factors through $\eta_4$, so $V_0^G=V_0^{c_\infty=1}$ and the remaining terms cancel. We obtain
\begin{equation}\label{eq:b6-v-euler}
 e_j=
 \begin{cases}
 -d_j^+,&j<0,\\
 0,&j=0,\\
 d_j-d_j^+,&j>0 .
 \end{cases}
\end{equation}
Here $(d_j,d_j^+)=(3,1)$ for $j=0$, $(2,1)$ for $|j|=1$, $(1,0)$ for $|j|=2$, and $(0,0)$ for $|j|\ge3$, so $e_{-1}=-1$, $e_1=e_2=1$, and $e_j=0$ for all other $j$. The local and global Euler formulas used here are those of \cite[I, Theorems~2.8 and~5.1]{milne-2006}.
\leanref{UpperTheory.\allowbreak{}selmerComplex\_\allowbreak{}finite\_\allowbreak{}and\_\allowbreak{}euler\_\allowbreak{}characteristic}{Development.\allowbreak{}AutomorphicArithmetic.\allowbreak{}AutomorphicVarieties.\allowbreak{}UnitaryThreeDeformations.\allowbreak{}SelmerComplexFiniteAndEulerCharacteristic}{174--184}
\leanref{UpperTheory.\allowbreak{}V.\allowbreak{}UnitaryThreeInertia.\allowbreak{}hostFrame\_\allowbreak{}count}{Development.\allowbreak{}AutomorphicArithmetic.\allowbreak{}AutomorphicVarieties.\allowbreak{}UnitaryThreeDeformations.\allowbreak{}HostCount}{288--352}
\leanref{UpperTheory.\allowbreak{}finrank\_\allowbreak{}gradedPiece\_\allowbreak{}and\_\allowbreak{}complexConjugation\_\allowbreak{}invariants}{Development.\allowbreak{}AutomorphicArithmetic.\allowbreak{}AutomorphicVarieties.\allowbreak{}UnitaryThreeDeformations.\allowbreak{}FinrankGradedPieceAndComplexConjugationInvariant}{90--205}
\leanref{UpperTheory.\allowbreak{}global\_\allowbreak{}cohomology\_\allowbreak{}finite\_\allowbreak{}set}{Development.\allowbreak{}AutomorphicArithmetic.\allowbreak{}AutomorphicVarieties.\allowbreak{}UnitaryThreeDeformations.\allowbreak{}GlobalCohomologyFiniteSet}{34--65}
\leanref{UpperTheory.\allowbreak{}global\_\allowbreak{}euler\_\allowbreak{}characteristic\_\allowbreak{}rationals\_\allowbreak{}odd\_\allowbreak{}prime}{Development.\allowbreak{}AutomorphicArithmetic.\allowbreak{}AutomorphicVarieties.\allowbreak{}UnitaryThreeDeformations.\allowbreak{}GlobalEulerCharacteristicRationalsOddPrime}{15--53}
\leanref{UpperTheory.\allowbreak{}local\_\allowbreak{}euler\_\allowbreak{}characteristic\_\allowbreak{}rationals}{Development.\allowbreak{}AutomorphicArithmetic.\allowbreak{}AutomorphicVarieties.\allowbreak{}UnitaryThreeDeformations.\allowbreak{}LocalEulerCharacteristicRationals}{32--79}

Choose a minimal adapted homogeneous presentation $\pi_W:P=O[[X_1,\ldots,X_n]]\twoheadrightarrow R$ whose coordinates map to centered matrix entries $Y_i=u(g_i)_{p_iq_i}-\rho_0(g_i)_{p_iq_i}$ of weights $w_i=w(p_i)-w(q_i)\in\{-2,-1,0,1,2\}$ (with $(w(0),w(1),w(2))=(1,0,-1)$), whose kernel $J=\ker\pi_W$ is homogeneous, and for which $\pi_W$ maps the weight-$j$ power series $P_j$ onto the weight-$j$ piece $R_j$ of $R$ under diagonal torus conjugation. The base homomorphism $R\to O$ specializing $u$ to $r_0$ kills every $Y_i$, so its composition with $\pi_W$ is the constant-term map $x_0:P\to O$. Put $n_j=\#\{i:w_i=j\}$. For each weight $j$, every cocycle in $Z^1_{\mathcal L,j}$ differs by a global coboundary from a normalized cocycle, meaning one vanishing at $c_\infty$ and having zero $(0,1)$ and $(1,0)$ entries in its variation of $r(g_1)-r(c_\infty g_1^{-1}c_\infty)$. Thus the normalized subspace $Z^1_{\mathcal L,j,\mathrm{norm}}\subseteq Z^1_{\mathcal L,j}$ surjects onto $H^1_{\mathcal L,j}$. By the local condition matching above, each cocycle in $Z^1_{\mathcal L,j,\mathrm{norm}}$ defines a dual-number lift of the marked frame over $\mathbf k[\epsilon]/(\epsilon^2)$, hence factors through $R$ and vanishes if it kills the weight-$j$ adapted coordinates. Evaluation on those coordinates therefore injects $Z^1_{\mathcal L,j,\mathrm{norm}}$ into $\mathbf k^{n_j}$, giving $\dim_{\mathbf k}H^1_{\mathcal L,j}\le n_j$.

We construct homogeneous generators of $J$ directly from the unframed obstruction spaces $\operatorname{Ob}_j$. By the Artin--Rees lemma in the Noetherian local ring $P$, choose $k_0\ge1$ such that $\mathfrak m_P^{k_0}\cap J\subseteq\mathfrak m_P J$. Let $\mathfrak t\subseteq\mathfrak m_P^{k_0}$ be the homogeneous truncation ideal of level $k_0$ so that $P/\mathfrak t$ is finite and $\Z$-graded, and put $\bar J=J\cdot(P/\mathfrak t)$. For each weight $j$, let $K_j\subseteq\bar J$ be the homogeneous ideal of $P/\mathfrak t$ generated by $\mathfrak m_{P/\mathfrak t}\bar J$ and all homogeneous components of $\bar J$ of weight different from $j$. The quotient $q_j:(P/\mathfrak t)/K_j\twoheadrightarrow(P/\mathfrak t)/\bar J$ is a small extension whose kernel is concentrated in weight $j$ and annihilated by the maximal ideal. Pushing out $q_j$ by a functional $\lambda\in\ker(q_j)^\vee$ gives a small $O$-algebra extension $E_\lambda\twoheadrightarrow R/\pi_W(\mathfrak t)$ with kernel $\mathbf k\epsilon$ in weight $j$, while $q_j$ is an isomorphism in all weights different from $j$. Lifting the matrix entries of the pushforward $\bar u$ homogeneously in $(P/\mathfrak t)/K_j$ (preserving the polarizing involution since $2\in O^\times$) produces a $2$-cocycle $c_2\in Z^2(G,V_j)$. Locally, $c_2|_{D_v}=du_v$ with local frame compatibility modulo $Z^1_{\mathcal L,v}$: for $v\in S\setminus T_b$, $\bar u|_{D_v}$ factors through $D_v/N_v$, where $H^2(D_v/N_v,V_j)=0$ since $\ell\nmid[I_v:N_v]$ and $D_v/I_v\simeq\widehat{\Z}$; for $v\in T_b$ and $j\ne0$, $H^2(D_v,V_j)=0$ gives $u_v$, with local frame compatibility by $H^1(D_v\cap\ker\eta_4,V_j)=0$ at $v\in S_b$, unchanged entries in weights $\le0$ at $v=\ell$ when $j>0$, and unique homogeneous lifting in weights $\ge0$ (extended across the index-at-most-two sign quotient) at $v=\ell$ when $j<0$; for $v\in T_b$ and $j=0$, in the homogeneous Henselian spectral-projector basis of $\bar u(x_v)$ the off-diagonal entries on $D_v\cap\ker\eta_4$ lie in nonzero weights where $q_0$ is an isomorphism, while the diagonal ratios to $\rho_0$ factor through $(D_v\cap\ker\eta_4)/I_v\simeq\widehat{\Z}$ and lift across $q_0$ compatibly with the polarizing involution, extending across the sign quotient to $D_v$. If the resulting class in $\operatorname{Ob}_j$ vanishes, global correction and normalization at $c_\infty$ and $g_1$ yield a marked frame lift over $E_\lambda$, hence by universality of $u$ a local $O$-algebra map $R\to E_\lambda$ lifting $R\twoheadrightarrow R/\pi_W(\mathfrak t)$; its pullback to $P$ differs from the canonical quotient $P\to E_\lambda$ by an $O$-derivation into $\mathbf k\epsilon$, which kills $\mathfrak m_P^2\supseteq J$ (by $x_0(J)=0$ and minimality of $\pi_W$), forcing $\lambda=0$ and hence $\dim_{\mathbf k}\ker(q_j)\le\dim_{\mathbf k}\operatorname{Ob}_j$. Hence $\bar J/\mathfrak m_{P/\mathfrak t}\bar J$ is generated in weight $j$ by $\dim_{\mathbf k}\operatorname{Ob}_j$ homogeneous elements. Taking weight-$j$ components of lifts in $J$ (using homogeneity of $J$ and $\mathfrak t$) produces $\dim_{\mathbf k}\operatorname{Ob}_j$ homogeneous power series $f_{j,i}\in J\cap P_j$ whose generated ideal $I_0\subseteq J$ satisfies $J\subseteq I_0+\mathfrak m_P J+(\mathfrak t\cap J)\subseteq I_0+\mathfrak m_P J$. Nakayama's lemma yields $J=I_0$.
\leanref{UpperTheory.\allowbreak{}V.\allowbreak{}UnitaryThreeInertia.\allowbreak{}hasOrigin}{Development.\allowbreak{}AutomorphicArithmetic.\allowbreak{}AutomorphicVarieties.\allowbreak{}UnitaryThreeDeformations.\allowbreak{}Origin}{73--110}
\leanref{UpperTheory.\allowbreak{}V.\allowbreak{}UnitaryThreeInertia.\allowbreak{}finrank\_\allowbreak{}H1L\_\allowbreak{}le}{Development.\allowbreak{}AutomorphicArithmetic.\allowbreak{}AutomorphicVarieties.\allowbreak{}UnitaryThreeDeformations.\allowbreak{}CommAlgFunctorCoboundaryZeroMemCocyclesOne}{370--422}
\leanref{UpperTheory.\allowbreak{}V.\allowbreak{}UnitaryThreeInertia.\allowbreak{}exists\_\allowbreak{}relations}{Development.\allowbreak{}AutomorphicArithmetic.\allowbreak{}AutomorphicVarieties.\allowbreak{}UnitaryThreeDeformations.\allowbreak{}ExistsRelations}{121--258}

Combining $\dim_{\mathbf k}H^1_{\mathcal L,j}\le n_j$ with \eqref{eq:b6-v-euler} gives
\[
 \dim_{\mathbf k}\operatorname{Ob}_j+e_j\le n_j \qquad(j\in\Z).
\]
Since $w_i\in\{-2,-1,0,1,2\}$, we have $n_j=0$ and hence $\dim_{\mathbf k}\operatorname{Ob}_j=0$ for $|j|>2$. For a diagonal point over $O[\epsilon]/(\epsilon^2)$ or $C_\ell$, the ratios to the distinguished characters on $G_K$ are unramified, hence trivial by reciprocity and $\operatorname{Cl}(\Q(i))=1$. This identifies the unique point where all nonzero weights vanish with $r_0$ and kills normalized $O$-valued tangents of weight zero. For $j=-n<0$ the scalar lines have characters $\eta_4^s\theta^{-n}\varepsilon_\ell^{-n}$ with $s=0$ or $1$. Lemma \ref{B6.31a} kills these $O$-valued tangents modulo coboundaries. The marking removes the remaining coboundaries in nonpositive weights, so every $O$-valued derivation at $x_0$ supported in weight $j\le0$ and killing $J$ is zero. By duality, the linear terms of the generators in weights $-2,-1,0$ span the corresponding rational cotangent spaces.
\leanref{UpperTheory.\allowbreak{}V.\allowbreak{}UnitaryThreeInertia.\allowbreak{}frame\_\allowbreak{}class\_\allowbreak{}one}{Development.\allowbreak{}AutomorphicArithmetic.\allowbreak{}AutomorphicVarieties.\allowbreak{}UnitaryThreeDeformations.\allowbreak{}FrameClassOne}{166--269}
\leanref{UpperTheory.\allowbreak{}V.\allowbreak{}UnitaryThreeInertia.\allowbreak{}eq\_\allowbreak{}rz\_\allowbreak{}of\_\allowbreak{}diag}{Development.\allowbreak{}AutomorphicArithmetic.\allowbreak{}AutomorphicVarieties.\allowbreak{}UnitaryThreeDeformations.\allowbreak{}UnitEqOneOfPowEqOneOfResidueOne}{165--209}
\leanref{UpperTheory.\allowbreak{}V.\allowbreak{}UnitaryThreeInertia.\allowbreak{}tangent\_\allowbreak{}coord\_\allowbreak{}eq\_\allowbreak{}zero\_\allowbreak{}of\_\allowbreak{}nonpos\_\allowbreak{}weight}{Development.\allowbreak{}AutomorphicArithmetic.\allowbreak{}AutomorphicVarieties.\allowbreak{}UnitaryThreeDeformations.\allowbreak{}TangentCoordEqZeroOfNonposWeight}{70--121}

For $0\ne b\in O$, let $\mathcal S_F(b)\subseteq H^1(G,(O/(b))(\Psi^{-1}))$ denote the classes on $G$ that are coboundaries on $D_v$ for $v\in T_b$ and vanish on $N_v$ for $v\in S\setminus T_b$. The unit $\Psi^{-1}(g_1)-1$ gives each class a unique representative $z$ with $z(g_1)=0$. On $N_v$ for $v\in S\setminus T_b$ we have $\Psi=1$, so subtracting the normalizing coboundary preserves the zero values. On each $D_v$ for $v\in T_b$, subtract its scalar from the local cobounding element. Since $\Psi(c_\infty)=1$ and $\ell$ is odd, $z(c_\infty)=0$, which also gives $z(c_\infty g_1^{-1}c_\infty)=0$.

Put $T_-=E_{10}-E_{21}$ in zero-based indices. Then $(1+\epsilon z T_-)\rho_0$ is the square-zero lift, with the original sign on the full group. The root line spanned by $T_-$ has action $\Psi^{-1}$, so this lift obeys the multiplication and polarization equations. The local witnesses give its local conjugations, and the two marking values give its marking equations. For an adapted coordinate $Y_i=u(g_i)_{p_iq_i}-\rho_0(g_i)_{p_iq_i}$ its value under the point derivation is
\[
 d_z(Y_i)=(T_-)_{p_iq_i}\rho_0(g_i)_{q_iq_i}\,z(g_i)\pmod b ,
\]
where matrix parts are used in the entry notation. This assignment is additive and supported in weight $-1$. It kills the linear terms of $J$ for every $b$. Lifting its coordinate values gives an integral derivation supported in weight $-1$. Conversely the values on the coordinates determine the square-zero representation, and its $(1,0)$ entries recover $z$ because the diagonal factors are units. This gives the additive injection required for the cotangent comparison (6.58.2) of \ref{B6.58}. If $n_{-1}=0$, it gives $\mathcal S_F(b)=0$ for every $b$. In this case take $V=0$ and no representatives.
\leanref{UpperTheory.\allowbreak{}V.\allowbreak{}UnitaryThreeInertia.\allowbreak{}mprime\_\allowbreak{}coord}{Development.\allowbreak{}AutomorphicArithmetic.\allowbreak{}AutomorphicVarieties.\allowbreak{}UnitaryThreeDeformations.\allowbreak{}PrecompleteOfPowLe}{385--427}

Suppose $n_{-1}>0$. At a marked point write $r(g)$ for its matrix part when taking entries. The normalization and the polarization give
\[
 r(g_1)-r(c_\infty g_1^{-1}c_\infty)
             =\operatorname{diag}(\zeta,0,-\zeta),
 \qquad
 \bar\zeta=\bar\Psi(g_1)-\bar\Psi(g_1)^{-1}\ne0.
\]
Thus every invariant subspace on the sign kernel is spanned by coordinate axes. Polarization pairs the relation that entry $ij$ vanishes throughout the sign kernel with the corresponding relation for $(2-j,2-i)$. In particular the vanished entries $10,21$ occur together, as do $01,12$. The coordinate subspaces of dimension one or two now give the upper-triangular case if $10,20$ vanish and the lower-triangular case if $01,02$ vanish. The remaining reducible pattern has the middle line and the corner on $\{0,2\}$, with both corner directions present. The corner is irreducible by the two distinct marked diagonal values and the presence of its two off-diagonal entries. A two-sided point excludes both triangular cases.
\leanref{UpperTheory.\allowbreak{}V.\allowbreak{}U3.\allowbreak{}shape\_\allowbreak{}cases}{Development.\allowbreak{}AutomorphicArithmetic.\allowbreak{}AutomorphicVarieties.\allowbreak{}UnitaryThreeGroups.\allowbreak{}ExistsCoordOfNotMatIrr}{45--115}

In the corner case, pass to the integral model. Entry $(2,1)$ vanishes on the sign kernel and then on the other coset by multiplication by the normalized complex conjugation, whose matrix part is the identity. Choose a nonzero $(2,0)$ entry $u_*$ of least valuation and write $r(g)_{20}=u_*b(g)$. Matrix multiplication, with the polarizing involution on the other coset, gives
\[
 \bar b(gh)=\bar b(g)\bar\Psi(h)+\eta_4(g)\bar\Psi(g)^{-1}\bar b(h).
\]
Thus $\bar b(g)\bar\Psi(g)^{-1}$ is a cocycle in weight $-2$, zero at complex conjugation. Outside $T_b$ the subgroups $N_v$ kill it. At $v\in T_b$ let $H\equiv1\bmod\mathfrak m$ carry the ordinary upper-triangular basis (which is diagonal for $v\in S_b$) to the given basis. For $g$ in the local sign kernel, comparison of the $0$ and $2$ entries of the first column of $r(g)H=H r^{\mathrm{tri}}(g)$ gives
\[
 u_* b(g)H_{00}^{\,2}=H_{20}V(g),\qquad
 V(g)=\sum_{j=0}^2 r(g)_{0j}H_{j0}-H_{00}r(g)_{22}.
\]
Here $\overline{V(g)}=\bar\Psi(g)-\bar\Psi(g)^{-1}$. Take $s=x_v$. Multiply the identities at $g$ and $s$ by $V(s)$ and $V(g)$ respectively, subtract, cancel $u_*$ in the integral domain and reduce. Since $\overline{V(s)}\ne0$ by \eqref{eq:b6-v-split}, this gives
\[
 \bar b(g)=w_v\bigl(\bar\Psi(g)-\bar\Psi(g)^{-1}\bigr),
 \qquad w_v=\bar b(s)/\overline{V(s)}.
\]
The cocycle is a coboundary on the local sign kernel, hence on the decomposition group by averaging over the quotient of order at most two. In particular it vanishes on $N_v$ for $v\in S_b$ and is a coboundary on $D_\ell$. Residual Selmer vanishing $H^1_{\mathcal L,-2}=0$ gives a global coboundary. Its value at complex conjugation forces the cobounding scalar to vanish since the character there is $-1$. This contradicts the nonzero reduction at the entry chosen for $u_*$. Every two-sided point is therefore Type I.
\leanref{UpperTheory.\allowbreak{}V.\allowbreak{}UnitaryThreeInertia.\allowbreak{}typeI\_\allowbreak{}of\_\allowbreak{}twoSigned}{Development.\allowbreak{}AutomorphicArithmetic.\allowbreak{}AutomorphicVarieties.\allowbreak{}UnitaryThreeDeformations.\allowbreak{}TypeIOfTwoSigned}{476--580}

At these points rigidity expresses a tangent by a commutator. Polarization at complex conjugation and the $g_1$ marking reduce it to the diagonal torus direction, which acts on weight $j$ by a scalar times $j$. An intertwiner on $G_K$ between two normalized Type I points intertwines their diagonal marking matrices. Their corresponding diagonal values have the same reductions, while the three reductions in each matrix are distinct. Thus the intertwiner is diagonal. It transports the polarizations to proportional forms by irreducibility, so its diagonal action is that of the marked torus up to scalar. Hence distinct orbits give different isomorphism classes for a fixed embedding of $K$.
\leanref{UpperTheory.\allowbreak{}V.\allowbreak{}UnitaryThreeInertia.\allowbreak{}tangent\_\allowbreak{}eq\_\allowbreak{}orbit}{Development.\allowbreak{}AutomorphicArithmetic.\allowbreak{}AutomorphicVarieties.\allowbreak{}UnitaryThreeDeformations.\allowbreak{}ContinuousCompOfFilteredContinuity}{262--400}
\leanref{UpperTheory.\allowbreak{}V.\allowbreak{}UnitaryThreeInertia.\allowbreak{}conj\_\allowbreak{}of\_\allowbreak{}matIso}{Development.\allowbreak{}AutomorphicArithmetic.\allowbreak{}AutomorphicVarieties.\allowbreak{}UnitaryThreeDeformations.\allowbreak{}NormTwoEqOne}{32--56}

Now set $r=n_{-1}\ge1$, $s=n_{-2}$, $m=n_0$, $r'=n_1-1$, and $s'=n_2-1$; here $r',s'\ge0$ because $n_1\ge\dim_{\mathbf k}\operatorname{Ob}_1+1\ge1$ and $n_2\ge\dim_{\mathbf k}\operatorname{Ob}_2+1\ge1$. Let $A=\bigoplus_{j\in\Z}P_j\subset P$ be the graded subalgebra of finite sums of homogeneous power series, with base point $x_0|_A:A\to O$ and canonical graded formal chart $A\hookrightarrow(\operatorname{Frac}O)[[X_1,\ldots,X_n]]$. Its variable counts in weights $(-2,-1,0,1,2)$ are $(s,r,m,r'+1,s'+1)$, with no other weights. Padding the homogeneous generators $f_{j,i}\in J\cap P_j$ with zeros to match $\dim_{\mathbf k}\operatorname{Ob}_j\le n_j-e_j$ produces homogeneous equations $G_1,\ldots,G_s$ in weight $-2$, $F_1,\ldots,F_{r+1}$ in weight $-1$, $Z_1,\ldots,Z_m$ in weight $0$, $E_1,\ldots,E_{r'}$ in weight $1$, and $E'_1,\ldots,E'_{s'}$ in weight $2$ generating $I=\ker(A\twoheadrightarrow\bigoplus_{j\in\Z}R_j)$. The image of $A_0$ in $A/I$ is $R_0$, finite over $O$; the unique fixed $C_\ell$-point is $x_0|_A$; every two-sided $C_\ell$-point is Type I with tangent space of dimension at most one; and a lift $a_F\in A_0\cap\ker x_0$ of $\operatorname{tr}u(\sigma)-\operatorname{tr}r_0(\sigma)\in R_0$ is nonzero at every two-sided $C_\ell$-point. Apply the cotangent comparison (6.58.2) (via the weight-$-1$ injection above) and the orbit estimate (6.58.3) of \ref{B6.58}, with the normalized valuation $e_Ov_\ell$, to $A$, $I$, and $a_F$. In both cases for $n_{-1}$ there are $V\in\N$ and Type I orbit representatives with matrix restrictions $\rho_i$ to $G_K$ such that
\begin{equation}\label{eq:b6-frame}
 \operatorname{length}_O\mathcal S_F(b)\le V
 \le2e_O\sum_i v_\ell\bigl(\operatorname{tr}\rho_i(\sigma)-
                         \operatorname{tr}\rho_0(\sigma)\bigr)
 \qquad(0\ne b\in O).
\end{equation}
\leanref{UpperTheory.\allowbreak{}V.\allowbreak{}UnitaryThreeInertia.\allowbreak{}Abs.\allowbreak{}coord\_\allowbreak{}count}{Development.\allowbreak{}AutomorphicArithmetic.\allowbreak{}AutomorphicVarieties.\allowbreak{}UnitaryThreeDeformations.\allowbreak{}AbsCoord}{68--280}
\leanref{UpperTheory.\allowbreak{}norm\_\allowbreak{}bound\_\allowbreak{}of\_\allowbreak{}integral\_\allowbreak{}values}{Development.\allowbreak{}AutomorphicArithmetic.\allowbreak{}AutomorphicVarieties.\allowbreak{}UnitaryThreeDeformations.\allowbreak{}NormBoundOfIntegralValues}{43--65}
\leanref{UpperTheory.\allowbreak{}V.\allowbreak{}UnitaryThreeInertia.\allowbreak{}Frame.\allowbreak{}ShaO}{Development.\allowbreak{}AutomorphicArithmetic.\allowbreak{}AutomorphicVarieties.\allowbreak{}UnitaryThreeDeformations.\allowbreak{}IsCharZeroPt}{95--98}
\leanref{UpperTheory.\allowbreak{}V.\allowbreak{}UnitaryThreeInertia.\allowbreak{}abstract\_\allowbreak{}count}{Development.\allowbreak{}AutomorphicArithmetic.\allowbreak{}AutomorphicVarieties.\allowbreak{}UnitaryThreeDeformations.\allowbreak{}AbstractCount}{25--165}

A class in $\mathcal S_{\ell,1,\theta}[b]$ for $b=\ell^m$ has a representative in $A_O[b]$. Subtract the coboundary of a $b$-division point of the witness for $bz$. At $v\in T_b$, evaluation at $x_v$ shows that the local witness is killed by $b$, since $\Psi(x_v)-1$ is a unit by \eqref{eq:b6-v-split}. Outside $S$ we have $\Psi|_{I_p}=1$, so the representative vanishes on $I_p$ and descends to $G$; on $N_v$ for $v\in S\setminus T_b$ we likewise have $\Psi=1$, so the representative vanishes there. Thus $\mathcal S_F(b)$ maps under inflation and $A_O[b]\hookrightarrow A_O$ onto a subgroup containing $\mathcal S_{\ell,1,\theta}[b]$. Continuous cocycles into $A_O$ have finite image, so \eqref{eq:b6-frame} gives $L_{\ell,1,\theta}\le V$.

Group the representatives by the embedding of $K$ into the algebraic closure used for their witnesses. For each group apply \ref{B6.41} with the restrictions at the index $z=\iota(\theta(\sigma))$ and empty sets at the other roots of unity. The restrictions at that index are pairwise nonisomorphic. With the constants $c_K,A_K$ of \ref{B6.41}, the sum of valuations times $\log\ell$ is at most $c_K\varphi(N)(2N)^{A_K}\log p$ per group, and there are at most two groups. Divide \eqref{eq:b6-frame} by $e_O$ and use $p\le12\ell N$.
\end{proof}

\subsection*{Proof of \ref{B6.35}}\label{prf:B6.35}\phantomsection
\begin{proof}
By \eqref{eq:b6-host-orders}, $\operatorname{ord}\theta<2oR_1(\ell,\theta)\le32o^2(\ell f)^2\ell^{2F_0}$. Let $L$ be the compositum of the field cut out by $\theta$ with $\Q(\mu_\ell)$. It is abelian and totally imaginary. With $M_\chi=\max\mathcal B_\chi$,
\[
 n=[L:\Q]\le\ell\operatorname{ord}\theta\le
 n_\chi:=32o^2(M_\chi f)^2M_\chi^{2F_0+1}.
\]
Set $X=\ell N$. Then $L\subseteq F=\Q(\mu_X)$. Discriminant transitivity gives the lower bound \cite[Ch.~III, Cor.~2.10]{neukirch-1999}. For the upper bound write $Y^X-1=\Phi_X(Y)G(Y)$ and evaluate its derivative at a primitive root $\zeta_X$. Thus $\Phi_X'(\zeta_X)$ divides $X\zeta_X^{-1}$ in $\Z[\zeta_X]$. The absolute discriminant of this power basis is $|N_{F/\Q}\Phi_X'(\zeta_X)|$, and the field discriminant divides it. Together these give $|D_L|^{[F:L]}\le |D_F|\le X^{[F:\Q]}$. Hence $\log(2|D_L|)\le(n_\chi+1)\log X$. With $C_K$ from \ref{B6.15}, equations \eqref{eq:B6-H2-local-correction} and \eqref{eq:B6-small-H2} give
\[
 s_{\ell,1}(L)\log\ell\le
 \log\#H^2(G_{L,S_\ell},\Z_\ell(2))\le X^A,\qquad
 A=(2n_\chi)^{\lceil C_K\rceil}(n_\chi+1).
\]
Apply \eqref{eq:b6-transfer}. The normalized character length times $\log\ell$ is at most $X^A+n\nu_\ell(N,1)\log\ell\le X^A+2n_\chi X$. A fixed power of $X$ bounds the last expression.
\end{proof}

\subsection*{Proof of \ref{B6.36}}\label{prf:B6.36}\phantomsection
\begin{proof}
Write $\alpha=\varepsilon_\ell^{-k}$. It suffices to kill the $\ell$-torsion because continuous cocycles have finite image. If $\ell z$ is the coboundary of $m$, subtract the coboundary of an $\ell$-division point of $m$ to put $z$ in $A_O[\ell]$.

If $\ell-1\mid k$, then $z$ is a homomorphism unramified outside $\ell$. Kronecker--Weber \cite[Ch.~V, Thm.~1.10]{neukirch-1999} puts its field in a cyclotomic field, and inertia at each $p\ne\ell$ removes the $p$-part of its conductor. Both $z$ and $\alpha$ factor through $\Gal(\Q(\mu_{\ell^\infty})/\Q)$. The decomposition group at $\ell$ surjects onto this group, so the local coboundary formula holds globally.

Otherwise $k=(\ell-1)u/2$ with $u$ odd and $\ell\equiv1\bmod4$. The residual character cuts out $K=\Q(\sqrt\ell)$ of discriminant $\ell$. Normalize $z(\delta)=0$ at $\delta\in I_\ell$ with $\alpha(\delta)-1$ a unit. On $G_K$ the cocycle is a homomorphism unramified away from $\ell$. The local coboundary witness at $\ell$ is killed by $\ell$, by evaluation at an inertia element with unit scalar difference. Hence $z$ vanishes on the decomposition groups over $\ell$ in $G_K$. The real places contribute zero by odd order, so $z|_{G_K}$ factors through $\operatorname{Cl}(K)$ by the small Hilbert class field theorem \cite[Ch.~VI, Prop.~6.9]{neukirch-1999}.

Minkowski's ideal-class bound \cite[Ch.~I, \S5, Ex.~3]{neukirch-1999} gives an integral ideal of norm at most $\sqrt\ell/2$ in each class: the degree is two, both embeddings are real, and the discriminant is $\ell$. In a quadratic field the number of ideals of norm $n$ is at most the divisor count $d(n)$ by prime-power factorization. With $T=\sqrt\ell$,
\begin{equation}\label{eq:b6-z-class-number}
 \#\operatorname{Cl}(K)\le\sum_{n\le T}d(n)
 =\sum_{a\le T}\left\lfloor T/a\right\rfloor
 \le T(1+\log T)<T^2=\ell .
\end{equation}
Consequently $z|_{G_K}=0$, and the normalization on the other coset gives $z=0$. In class Z the range follows from \ref{B6.28a} and parity.
\end{proof}

\subsection*{Proof of \ref{B6.36a}}\label{prf:B6.36a}\phantomsection
\begin{proof}
Put $\Delta=\Gal(E/\Q)$ and extend coefficients to $O$. For a cyclic generator $T$ of order $m$, $\prod_{\zeta^m=1}(T-\zeta)=0$. For endomorphisms $A,B$, the sequence exact at its first two terms
\[
 0\longrightarrow\ker B\longrightarrow\ker(AB)
      \xrightarrow{B}\ker A
\]
gives $\operatorname{length}\ker(AB)\le\operatorname{length}\ker B+\operatorname{length}\ker A$. Iterate over the linear factors and the cyclic factors of $\Delta$. The conjugation sign selects the characters in $X_{\ell,O}(k)$. The sum of their eigenspace lengths bounds the length.

Twist the $\theta$-eigenspace to invariants with coefficients $W=M_{\ell,k,\theta}$. Transgression embeds the global restriction cokernel in $H^2(\Delta,W^{G_E})$ \cite[Prop.~1.6.7]{neukirch-2008}. For a finite abelian group $B$ acting by scalars $\lambda(g)$ on a discrete $O$-module $U$ with finite-length cyclic submodules, each $\lambda(g)-1$ kills cohomology because $g$ acts as the identity there \cite[Prop.~1.6.3]{neukirch-2008}. If $(a)$ is the ideal of the scalar differences and $Z^i$ is the module of $i$-cocycles, $H^i(B,U)$ is a quotient of $Z^i/aZ^i$. The finite-submodule argument of \ref{B6.32a} bounds its length by that of $Z^i[a]\subseteq(U^B)^{|B|^i}$. Applied to the global term in degree two, this bounds the loss by $e_On^2\nu_\ell(N,k)$.

For a global class downstairs whose restriction upstairs is strict, localize above each relevant rational prime. The local values lie in the local restriction kernels, and their joint kernel imposes the downstairs conditions. Since the places over each rational prime are conjugate, one representative suffices. At $\ell$ the decomposition quotient has order at most $n$ and the invariant length is at most $e_O\nu_\ell(N,k)$ by \ref{B6.32a}. The scalar bound in degree one gives loss at most $e_On\nu_\ell(N,k)$. At $p\ne\ell$ the inertia quotient is trivial for $p\nmid N$. For $p\mid N$, its order divides $n$. Its first cohomology is killed by $n$ \cite[Prop.~1.6.1]{neukirch-2008}. Divisibility allows representatives in $W[n]$, so evaluation at quotient elements bounds the length by $e_Onv_\ell(n)$.

The added decomposition condition at $p\mid N$, $p\ne\ell$, costs at most the length of $H^1(D_p/I_p,W^{I_p})$. For unramified $\theta$, Frobenius acts by a root of unity times $p^{-k}$, so Frobenius minus the identity is surjective on $W$. For ramified $\theta$, a nontrivial inertia value $\zeta$ has $\zeta^n=1$, so $\zeta-1$ divides $n$ and $W^{I_p}\subseteq W[n]$. The formula $H^1(D_p/I_p,W^{I_p})=W^{I_p}/(\mathrm{Frob}_p-1)W^{I_p}$ for the discrete torsion module \cite[Prop.~1.7.7(i)]{neukirch-2008} bounds the cost by $e_Ov_\ell(n)$. These terms give \eqref{eq:b6-loss} using \ref{B6.32a}, $n,\omega(N),v_\ell(n)\le Nk$ and $Nk\ge16$. If $\ell\nmid n$, averaging kills the finite-quotient cohomology. The added conditions cost zero by the Frobenius calculation and $W[n]=0$.
\end{proof}

\subsection*{Proof of \ref{B6.37}}\label{prf:B6.37}\phantomsection
\begin{proof}
Choose $O$ as in \ref{B6.36a}, and put $M=Nt$. Among the characters occurring for $1\le k<t$, those with $\operatorname{ord}\theta_P$ odd have conductor dividing $f$ by \ref{B6.28a}. There are at most $f$ distinct such characters. Each contributes at most $B+M^{C_4}$ before multiplication by $\log\ell$, by \ref{B6.32}. The other characters contribute finite weighted terms of size at most $M^{2b}$ by \ref{B6.31}, \ref{B6.34} and \ref{B6.35}, where $b=\max\{C_{\mathrm{II}},C_{\mathrm V},C_{\mathrm{Vs}}\}$ and $\ell\le N$. There are at most $N$ terms per twist. Equations \eqref{eq:b6-redall} and \eqref{eq:b6-loss} give finiteness and bound the weighted sum by
\[
 f\log\ell\,(B+M^{C_4})+M^{2b+1}+M^7\log\ell .
\]
Choose $C_{\ell,j}$ to absorb these terms.
\end{proof}

\subsection*{Proof of \ref{B6.38}}\label{prf:B6.38}\phantomsection
\begin{proof}
By \ref{B6.28a}, $[E:\Q]\mid2^{a+2}o$, so \ref{B6.36a} applies with zero loss. Use \ref{B6.30} with its two primes to sum class I for each twist. Every other prime in the sum has $\ell\le2t-1$ by \ref{B6.28a}. There are at most $N$ characters per twist and prime. Class Z contributes zero by \ref{B6.36}, and classes II and III contribute at most $(\ell Nk)^{\max(C_{\mathrm{II}},C_{\mathrm{III}})}$ per character by \ref{B6.31}. These finite bounds sum to at most a fixed power of $Nt$.
\end{proof}

\subsection*{Proof of \ref{B6.39}}\label{prf:B6.39}\phantomsection
\begin{proof}
Apply \ref{B6.37} to $S\cap\mathcal B_\chi$ and \ref{B6.38} to its complement. The fixed number of summands and the coefficients of $B$ are absorbed into $C$. Singletons give finiteness at every odd prime.

Take $X_0\ge2^J$. On each $K$-window, $2\le x\le X$, so admissibility and the prime bounds give largeness for $t(x)$. The odd-layer conditions in the definition of the window supply $B=D_x$ in \eqref{eq:b6-point}. The character group of $E$ contains $\chi$ and the quadratic characters of conductors $q,d$. It also contains a faithful character of the dyadic quotient by $\{\pm1\}$, of conductor $2^{a(x)+2}$. These conductors have least common multiple $N$. Since $E$ is realizable in $\Q(\mu_N)$, it follows that $\operatorname{cond}(E)=N$. Now \eqref{eq:b6-point} gives \eqref{eq:b6-pluswindow}. For the specified tower data, Theorem \ref{B6.27} supplies the window and the bound at the prime $2$. The floor choice lies between $2$ and $x^{401}$ for every $x\ge2$.
\end{proof}

\section*{Proofs for Chapter~\ref{ch:vwin-automorphic}: Automorphic counts and deformation rigidity}

\subsection*{Proof of \ref{B6.40n}}\label{prf:B6.40n}\phantomsection
\begin{proof}
Write \(U^\wedge=\widehat{U[1/\ell]}_{\mathfrak p_0}\) for the completion of \(U[1/\ell]\) at \(\mathfrak p_0=(\ker x_{0,\ell})U[1/\ell]\); it represents framed deformations of \(r_0|_{D_\ell}:D_\ell\to\mathcal G_3(E)\) over local Artinian \(E\)-algebras with residue field \(E\). Let \(R^{\mathrm{cris}}_{r_0}=U^\wedge/I_{\mathrm{cris}}\) be the quotient parameterising those deformations whose restriction to \(G_{K_1}\) is crystalline.

We first show that \(R^{\mathrm{cris}}_{r_0}\) is formally smooth over \(E\) of dimension \(12\) with graded tangent dimensions \((1,2,3,4,2)\). Over \(E\), the adjoint representation \(\operatorname{ad}r_0=M_3(E)\) decomposes into five weight spaces \(V_j\) (\(j=b-a\in\{-2,-1,0,1,2\}\)) under diagonal conjugation:
\[
 V_2=E(\Psi^2\eta),\quad
 V_1=E(\Psi)\oplus E(\Psi\eta),\quad
 V_0=E(1)\oplus E(\eta)^{\oplus 2},\quad
 V_{-1}=E(\Psi^{-1})\oplus E(\Psi^{-1}\eta),\quad
 V_{-2}=E(\Psi^{-2}\eta).
\]
A first-order deformation \((1+\epsilon z)r_0|_{D_\ell}\) over \(E[\epsilon]/(\epsilon^2)\) is crystalline on \(G_{K_1}\) precisely when \(z|_{G_{K_1}}\in Z^1_f(G_{K_1},\operatorname{ad}r_0)\). Because \([D_\ell:G_{K_1}]\) is finite and \(\operatorname{char}E=0\), Galois descent on \(B_{\mathrm{cris}}\)-cohomology identifies these cocycles with \(Z^1_f(D_\ell,\operatorname{ad}r_0)\). For each \(D_\ell\)-representation \(V\) that is crystalline on \(G_{K_1}\), Bloch--Kato's dimension formula \cite[Proposition~3.8]{bloch-kato-1990} and \(\dim_E B^1(D_\ell,V)=\dim_E V-\dim_E H^0(D_\ell,V)\) give
\[
 \dim_E Z^1_f(D_\ell,V)
 =\dim_E H^1_f(D_\ell,V)+\dim_E V-\dim_E H^0(D_\ell,V)
 =\dim_E V+\dim_E\bigl(D_{\mathrm{dR}}(V)/\mathrm{Fil}^0 D_{\mathrm{dR}}(V)\bigr).
\]
In the convention that \(\Q_\ell(1)\) has filtration jump \(-1\), a character of cyclotomic weight \(m\) contributes \(1\) to \(\dim_E(D_{\mathrm{dR}}/\mathrm{Fil}^0)\) if \(m>0\) and \(0\) if \(m\le0\). The Hodge--Tate weights of \(V_{-2},V_{-1},V_0,V_1,V_2\) are \(-2k\); \(-k,-k\); \(0,0,0\); \(k,k\); and \(2k\), respectively. Hence \(\dim_E Z^1_f(D_\ell,V_j)\) equals
\[
 1+0=1,\quad 2+0=2,\quad 3+0=3,\quad 2+2=4,\quad 1+1=2
\]
for \(j=-2,-1,0,1,2\), summing to \(12\). This holds for every \(k\ge1\) (including resonant \(k=1\) with \(\Psi|_{D_\ell}=\varepsilon_\ell\) or \((\eta\Psi)|_{D_\ell}=\varepsilon_\ell\)) and for both split and inert \(\ell\).

To prove that \(R^{\mathrm{cris}}_{r_0}\) is formally smooth over \(E\), let \(A'\twoheadrightarrow A\) be a small extension of local Artinian \(E\)-algebras with residue field \(E\), and let \(r:D_\ell\to\mathcal G_3(A)\) be a deformation of \(r_0|_{D_\ell}\) crystalline on \(G_{K_1}\). Let \(\Delta=\operatorname{Gal}(K_1/K_v)\) (of order \(e\in\{1,2\}\)) and let \(K_{1,0}=K_v\) be the maximal unramified subfield of \(K_1\). By the equivalence of categories between crystalline \(G_{K_1}\)-representations and weakly admissible filtered \(\varphi\)-modules \cite[Thm.~A and \S4.1]{colmez-fontaine-2000}, \(r|_{G_{K_v}}\) corresponds to a free \(K_v\otimes_{\Q_\ell}A\)-module \(D_A\) of rank \(3\) with a \(\Delta\)-action, a \(\Delta\)-equivariant \(\sigma\)-semilinear Frobenius \(\varphi_A:D_A\to D_A\), and a \(\Delta\)-stable Hodge filtration \(\mathrm{Fil}^\bullet_A\) on \((D_A)_{K_1}\) with jumps \(-k,0,k\).
\begin{itemize}
\item If \(\ell\) splits in \(K\) (\(K_v=\Q_\ell\)), lifting \(r\) as an unframed representation amounts to lifting the free \(A\)-module \(D_A\) with its finite \(\Delta\)-action (unobstructed since \(\operatorname{char}E=0\)), lifting the invertible \(\Delta\)-equivariant operator \(\varphi_A\) to \(\varphi_{A'}\), and lifting the full flag \(\mathrm{Fil}^\bullet_A\) in the rank-\(3\) free \(A\)-module \(((D_A)_{K_1})^\Delta\) to a full flag \(\mathrm{Fil}^\bullet_{A'}\) in \(((D_{A'})_{K_1})^\Delta\). There is no compatibility equation between \(\varphi_{A'}\) and \(\mathrm{Fil}^\bullet_{A'}\).
\item If \(\ell\) is inert in \(K\) (\([K_v:\Q_\ell]=2\)), fix a finite extension \(E'/E\) containing \(K_v\) so that \(K_v\otimes_{\Q_\ell}E'=E'\times E'\), and to check formal smoothness of \(R^{\mathrm{cris}}_{r_0}\otimes_E E'\) over \(E'\), take a small extension \(A'\twoheadrightarrow A\) of local Artinian \(E'\)-algebras with residue field \(E'\), writing \(D_A=D_{1,A}\oplus D_{2,A}\) accordingly. Choose \(\phi_0\in D_\ell\setminus G_{K_v}\); because \(\operatorname{Gal}(K_1/\Q_\ell)\) has exponent dividing \(2\), we have \(\phi_0^2\in G_{K_1}\). Setting \(P=r(\phi_0)_{\mathrm{mat}}\), the group relations \(P\theta(r(g)_{\mathrm{mat}})P^{-1}=r(\phi_0 g\phi_0^{-1})_{\mathrm{mat}}\) (\(g\in G_{K_v}\)) and \(P\theta(P)=r(\phi_0^2)_{\mathrm{mat}}\) say that \(f=PJ:V^\vee\xrightarrow{\sim}V^{\phi_0}\) is a \(G_{K_v}\)-isomorphism satisfying \((f^{-1})^{\phi_0}\circ r(\phi_0^2)_{\mathrm{mat}}=(f^\vee)^{-1}\). Under \(D_{\mathrm{cris}}\), since \(\phi_0^2\in G_{K_1}\) acts trivially on \(D_A\), this corresponds to a \(\Delta\)-equivariant isomorphism \(F=(F_1,F_2):D_A^\vee\xrightarrow{\sim}D_A^{\phi_0}\) of filtered \(\varphi\)-modules satisfying \(F^{\phi_0}=F^\vee\), i.e.\ \(F_2=F_1^\vee\) for an unconstrained \(\Delta\)-equivariant isomorphism \(F_1:D_{1,A}^\vee\xrightarrow{\sim}D_{2,A}\). Writing \(\varphi_A=(\varphi_{1,A},\varphi_{2,A})\) with \(\varphi_{1,A}:D_{1,A}\to D_{2,A}\) and \(\varphi_{2,A}:D_{2,A}\to D_{1,A}\), compatibility \(F\circ(\varphi_A^\vee)^{-1}=\varphi_A^{\phi_0}\circ F\) is equivalent to the single formula \(\varphi_{2,A}=F_1^\vee(\varphi_{1,A}^\vee)^{-1}F_1^{-1}\), and compatibility with the filtration says that \(\mathrm{Fil}^\bullet_{2,A}\) is the orthogonal complement of \(\mathrm{Fil}^\bullet_{1,A}\) under \(F_1\). Thus lifting \((D_A,\varphi_A,\mathrm{Fil}^\bullet_A,F)\) to \(A'\) amounts to lifting \(D_{1,A}\) with its \(\Delta\)-action, lifting \(F_1\), lifting \(\varphi_{1,A}\) arbitrarily (which defines \(\varphi_{2,A'}\)), and lifting the \(\Delta\)-stable flag \(\mathrm{Fil}^\bullet_{1,A}\) in \(((D_{1,A})_{K_1})^\Delta\) arbitrarily (which defines \(\mathrm{Fil}^\bullet_{2,A'}\)).
\end{itemize}
In both cases weak admissibility persists automatically across the small extension \(A'\twoheadrightarrow A\): since the steps of \(\mathrm{Fil}^\bullet_{A'}\) are free over \(A'\), we have a short exact sequence of filtered \(\varphi\)-modules \(0\to D_0\otimes_E\ker(A'\to A)\to D_{A'}\to D_A\to0\) (with \(E'\) in place of \(E\) when \(\ell\) is inert), and weak admissibility is closed under extensions (for any \(\varphi\)-stable \(K_v\)-subspace \(M\subseteq D_{A'}\), setting \(M_1=M\cap(D_0\otimes\ker)\) and \(M_2=\operatorname{im}(M\to D_A)\) gives \(t_H(M)\le t_H(M_1)+t_H(M_2)\le t_N(M_1)+t_N(M_2)=t_N(M)\), with equality for \(M=D_{A'}\)). Thus the unframed crystalline \(\mathcal G_3\)-representation lifts to \(A'\), and since a basis of a free \(A\)-module lifts to any free \(A'\)-module lifting it, the framing lifts as well. Hence \(R^{\mathrm{cris}}_{r_0}\) is formally smooth over \(E\) of dimension \(12\) (descending from \(E'\) to \(E\) by faithful flatness when \(\ell\) is inert), and therefore \(R^{\mathrm{cris}}_{r_0}\simeq E[[t_1,\ldots,t_{12}]]\).

We next compare \(I\cdot U^\wedge\) with \(I_{\mathrm{cris}}\). First, let \(J_{\mathrm{cris}}\subset U\) be the closed ideal cutting out the framed crystalline locus of \(r^{\mathrm{univ}}|_{G_{K_1}}\) with Hodge--Tate weights \((k,0,-k)\) and trivial inertial type, obtained by pulling back the crystalline quotient of \cite[Thm.~1.2.2]{kisin-2010} along \(G_{K_1}\hookrightarrow D_\ell\). At any characteristic-zero point \(y\) of \((U/J_{\mathrm{ord}})[(\ell\delta_1\delta_2)^{-1}]\) (or any finite Artinian \(E'\)-algebra point), the polynomials \(P_i/c_i\) give orthogonal rank-\(1\) idempotents summing to \(I_3\), so \(r_y|_{G_{K_v}}\) preserves a full flag with diagonal characters \(\lambda_0,\lambda_1,\lambda_2\) whose restrictions to \(I_{K_1}\) are \(\varepsilon_\ell^k,1,\varepsilon_\ell^{-k}\). Since \(\delta_1(y)\delta_2(y)\ne0\) and \(\varphi_v\in G_{K_1}\), we have \(\lambda_0(\varphi_v)\ne\varepsilon_\ell(\varphi_v)\lambda_1(\varphi_v)\) and \(\lambda_1(\varphi_v)\ne\varepsilon_\ell(\varphi_v)\lambda_2(\varphi_v)\), while the outer ratio \(\lambda_0/\lambda_2=\varepsilon_\ell^{2k}\) on \(I_{K_1}\) differs from \(\varepsilon_\ell\) since \(2k\ne1\). Lemma~\ref{B6.48h} therefore shows that \(r_y|_{G_{K_1}}\) is crystalline, so \(y\) kills \(J_{\mathrm{cris}}\). Because \(U[1/\ell]\) is a Jacobson ring and \((U/I_{\mathrm{sat}})[1/\ell]\) is reduced with \(\delta_1\delta_2\) a nonzerodivisor, \(J_{\mathrm{cris}}\subseteq I_{\mathrm{sat}}\subseteq I\). In particular, \(I_{\mathrm{cris}}=J_{\mathrm{cris}}U^\wedge\subseteq I\cdot U^\wedge\).

For the reverse inclusion, consider the universal crystalline deformation \(r^{\mathrm{cris}}\) over the regular complete local domain \(R^{\mathrm{cris}}_{r_0}\simeq E[[t_1,\ldots,t_{12}]]\) and the inverse-limit filtered \(\varphi\)-module \((D,\varphi,\mathrm{Fil}^\bullet)=\varprojlim_n D_{\mathrm{cris},K_1}(r^{\mathrm{cris}}\bmod\mathfrak m^n)\) over \(K_v\otimes_{\Q_\ell}R^{\mathrm{cris}}_{r_0}\). Let \(d_v=[K_v:\Q_\ell]\in\{1,2\}\). Modulo \(\mathfrak m_{R^{\mathrm{cris}}_{r_0}}\), the \(K_v\)-linear Frobenius \(\varphi^{d_v}\) has eigenvalues in \(E\) of pairwise distinct \(\ell\)-adic valuations \(-kd_v,0,kd_v\) (corresponding to the distinct weights \(k,0,-k\)). Applying Hensel's lemma on each local factor of the finite product of complete local rings \(K_v\otimes_{\Q_\ell}R^{\mathrm{cris}}_{r_0}\), \(D\) splits uniquely into three \(\varphi^{d_v}\)-stable and \(\Delta\)-stable rank-\(1\) direct summands \(D=L_0\oplus L_1\oplus L_2\) lifting the base eigenspaces \(L_{i,0}\) of slopes \(-kd_v,0,kd_v\); since \(\varphi\) commutes with \(\varphi^{d_v}\) and preserves each \(L_{i,0}\) (with \(\sigma\)-fixed base eigenvalues in \(E\)), uniqueness of the splitting forces \(\varphi(L_i)=L_i\). At the base point, \(\mathrm{Fil}^k_0=(L_{2,0})_{K_1}\) and \(\mathrm{Fil}^0_0=(L_{2,0}\oplus L_{1,0})_{K_1}\), so \((L_{0,0})_{K_1}\cap\mathrm{Fil}^0_0=0\) and \((L_{0,0}\oplus L_{1,0})_{K_1}\cap\mathrm{Fil}^k_0=0\). Transversality persists over \(K_1\otimes_{\Q_\ell}R^{\mathrm{cris}}_{r_0}\), giving \((L_0)_{K_1}\cap\mathrm{Fil}^0=0\) and \((L_0\oplus L_1)_{K_1}\cap\mathrm{Fil}^k=0\) while \((L_0\oplus L_1)_{K_1}\cap\mathrm{Fil}^0\) is a free direct summand of rank \(1\). Thus the induced filtrations on \(L_0\) and \(L_0\oplus L_1\) are flat over \(R^{\mathrm{cris}}_{r_0}\), with Hodge jump \(-k\) and Newton slope \(-k\) on \(L_0\), and Hodge jumps \(-k,0\) and Newton slopes \(-k,0\) on \(L_0\oplus L_1\). Over every Artinian quotient \(R^{\mathrm{cris}}_{r_0}/\mathfrak m^n\), \(L_0\) and \(L_0\oplus L_1\) are successive extensions of the weakly admissible base modules \(L_{0,0}\) and \(L_{0,0}\oplus L_{1,0}\), so \(0\subset L_0\subset L_0\oplus L_1\subset D\) is a flag of weakly admissible filtered \(\varphi\)-modules with \(\Delta\)-action. By \cite[Thm.~A]{colmez-fontaine-2000}, \(r^{\mathrm{cris}}|_{G_{K_v}}\) preserves a full flag over every \(R^{\mathrm{cris}}_{r_0}/\mathfrak m^n\) whose graded inertia characters are \((\Psi,1,\Psi^{-1})\). Hence \(J_{\mathrm{ord}}\) maps to zero in \(R^{\mathrm{cris}}_{r_0}\).

Moreover \(\ell\delta_1\delta_2\) has nonzero image in the domain \(R^{\mathrm{cris}}_{r_0}\): in the diagonal unramified crystalline family \(g\mapsto(\operatorname{diag}(\Psi(g)u(g),1,\Psi(g)^{-1}u(g)^{-1}),\eta(g))\) with \(u\) unramified and \(u(\varphi_v)=1+T\), both \(\lambda_0(\varphi_v)/\lambda_1(\varphi_v)=\Psi(\varphi_v)(1+T)\) and \(\lambda_1(\varphi_v)/\lambda_2(\varphi_v)=\Psi(\varphi_v)(1+T)\) are nonconstant, so \(\delta_1\) and \(\delta_2\) do not vanish identically on \(R^{\mathrm{cris}}_{r_0}\). Because \(R^{\mathrm{cris}}_{r_0}\) is a characteristic-zero domain, its defining prime ideal \(I_{\mathrm{cris}}\subset U^\wedge\) contains \(I_{\mathrm{sat}}U^\wedge\). Combined with \(I_{\mathrm{cris}}\subseteq I\cdot U^\wedge\), and since \(I\cdot U^\wedge\) is the intersection of the minimal primes of \(I_{\mathrm{sat}}U^\wedge\) inside \(\mathfrak m_{U^\wedge}\), it follows that \(I_{\mathrm{cris}}\) is the unique minimal prime of \(I_{\mathrm{sat}}U^\wedge\) inside \(\mathfrak m_{U^\wedge}\) and \(I\cdot U^\wedge=I_{\mathrm{cris}}\).

Because \(I\) is reduced and every minimal prime of \(I\) lies in \(\ker x_{0,\ell}\), the composite \(B=U/I\to(U/I)_{\ker x_{0,\ell}}\to\widehat{B[1/\ell]}_{\ker x_{0,\ell}}=R^{\mathrm{cris}}_{r_0}\) is injective by Krull's intersection theorem. In particular \(I\) is a single prime ideal and \(B\) is a complete Noetherian local \(O\)-domain.

For conjugation invariance, let \(R\) be a complete Noetherian local \(O\)-algebra with residue field \(\mathbf k\) and let \(h\in\GL_3(R)\) satisfy \(h\equiv I_3\bmod\mathfrak m_R\). For \(g\in G_{K_v}\), \((h,1)r(g)(h,1)^{-1}\) has matrix part \(hr(g)_{\mathrm{mat}}h^{-1}\), which conjugates \(P_i\) to \(hP_ih^{-1}\) and leaves \(\det(XI_3-r(s_0)_{\mathrm{mat}})\), \(t_i\), \(\delta_1\), and \(\delta_2\) unchanged. Thus \(J_{\mathrm{ord}}\), \(\sqrt{J_{\mathrm{ord}}}\), and \(I_{\mathrm{sat}}\) are invariant under conjugation by \((h,1)\). On the universal formal conjugation map \(U\to U\widehat\otimes_O O[[z_{ab}]]\) (for \(h=I_3+(z_{ab})\)), the preimage of the prime ideal \(I\cdot(U\widehat\otimes_O O[[z_{ab}]])\) is a prime ideal containing \(I_{\mathrm{sat}}\) and contained in \(\ker x_{0,\ell}\); uniqueness of the minimal prime \(I\) forces this preimage to contain \(I\), proving conjugation invariance of \(I\). Specialising to the diagonal torus \(\operatorname{diag}(t,1,t^{-1})\) shows that the kernel of a minimal graded power series presentation \(O[[X]]\twoheadrightarrow U\twoheadrightarrow B=O[[X]]/(H)\) is homogeneous, and specialising to \(h=I_3+\epsilon Y\) over \(O[\epsilon]/(\epsilon^2)\) gives the derivation killing \(H\). Since the diagonal torus acts semisimply on \(R^{\mathrm{cris}}_{r_0}/\mathfrak m^n\) for every \(n\) with cotangent weights \((1,2,3,4,2)\), lifting a homogeneous basis of \(\mathfrak m/\mathfrak m^2\) to homogeneous elements of \(\mathfrak m\) yields a graded isomorphism \(R^{\mathrm{cris}}_{r_0}\simeq E[[t_1,\ldots,t_{12}]]\) and hence the graded rational formal chart \(\varphi:U\to E[[t_1,\ldots,t_{12}]]\) with \(\ker\varphi=I\). This proves (1) and (2).

Finally, for (3), since \(J_{\mathrm{cris}}\subseteq I\), every characteristic-zero point \(f:B\to C_\ell\) kills \(J_{\mathrm{cris}}\) and \(J_{\mathrm{ord}}\), so \(r_f|_{G_{K_1}}\) is crystalline and (using the idempotents \(P_i/c_i\) over \(C_\ell\)) \(r_f|_{G_{K_v}}\) preserves a full flag with inertia characters \((\Psi,1,\Psi^{-1})\). Any characteristic-zero tangent vector \(B\to C_\ell[\epsilon]/(\epsilon^2)\) at \(f\) likewise kills \(J_{\mathrm{cris}}\), so its deformation is crystalline on \(G_{K_1}\), and finite-index Galois descent places its cocycle in \(Z^1_f(D_\ell,\operatorname{ad}r_f)\).
\end{proof}

\subsection*{Proof of \ref{B6.40p}}\label{prf:B6.40p}\phantomsection
\begin{proof}
For (1), let \(r\) be a two-sided characteristic-zero point of \(S\). Since \(r(c)_{\mathrm{mat}}=I_3\), for every \(g\in G_K\) we have \(r(cg^{-1}c)_{\mathrm{mat}}=J(r(g)_{\mathrm{mat}})^tJ\). Applying this to \(g_1\) shows that \(Z=r(g_1)_{\mathrm{mat}}-r(cg_1^{-1}c)_{\mathrm{mat}}\) satisfies \(JZ^tJ=-Z\), so \(Z_{ab}=-Z_{2-b,2-a}\). Together with the marking equations \(Z_{01}=Z_{10}=0\), this forces \(Z=\operatorname{diag}(\zeta,0,-\zeta)\) with \(\bar\zeta=\bar\Psi(g_1)-\bar\Psi(g_1)^{-1}\ne0\). Hence \(Z\) has three pairwise distinct eigenvalues, so every \(G_K\)-invariant subspace of \(C_\ell^3\) is spanned by coordinate axes. Polarisation \(r(cg^{-1}c)_{\mathrm{mat}}=J(r(g)_{\mathrm{mat}})^tJ\) pairs the vanishing of entry \((a,b)\) throughout \(G_K\) with that of \((2-b,2-a)\); in particular \((1,0),(2,1)\) vanish together and \((0,1),(1,2)\) vanish together. Because \(r\) is two-sided, neither all negative-weight entries nor all positive-weight entries vanish on \(G_K\), excluding upper-triangular and lower-triangular shapes. Any proper nonzero \(G_K\)-invariant coordinate subspace therefore forces the middle coordinate \(1\) to split from \(\{0,2\}\), so all odd-weight entries vanish on \(G_K\), while both \((0,2)\) and \((2,0)\) are nonzero on \(G_K\), making the corner on \(\{0,2\}\) irreducible.

In Type~II, put \(\kappa(g)=r(g)_{\mathrm{mat},11}\) for \(g\in G_K\). Then \(\kappa:G_K\to C_\ell^\times\) is a continuous character satisfying \(|\kappa(g)-1|_\ell<1\). At any finite place \(A\) of \(\overline\Q\) not above \(\ell\), \(r\) kills \(I_A\cap\ker\eta\cap\ker\vartheta\), whose quotient in \(I_A\cap G_K\) has order dividing \(\operatorname{ord}\vartheta\) and hence prime to \(\ell\); since \(r\) and \(r_0\) have the same residual values on this finite quotient, they are conjugate on \(I_A\cap G_K\). At any place \(A\mid\ell\), conjugating to the distinguished place \(A_0\) and applying Lemma~\ref{B6.40n}(3) shows that \(r|_{I_A}\) is conjugate to an upper-triangular representation with diagonal \(r_0|_{I_A}\). Consequently, for every finite place \(A\) and every \(g\in I_A\cap G_K\), \(\det(X-r(g)_{\mathrm{mat}})=(X-\Psi(g))(X-1)(X-\Psi(g)^{-1})\). If \(\bar\Psi(g)\ne1\), then \(|\Psi(g)^{\pm1}-1|_\ell=1\), whereas \(|\kappa(g)-1|_\ell<1\), forcing \(\kappa(g)=1\). If \(\bar\Psi(g)=1\) and \(A\nmid\ell\), then \(g\in\mathcal K_{\mathrm{away}}\), so \(r(g)=1\) and \(\kappa(g)=1\). If \(\bar\Psi(g)=1\) and \(A\mid\ell\), then because \(\ell-1\nmid2k\) and \(\operatorname{comp}_\ell(\vartheta)^2=1\), there exists \(s\in I_A\subset G_K\) with \(\bar\Psi(s)^2\ne1\); then both \(\bar\Psi(s)\ne1\) and \(\bar\Psi(gs)\ne1\), so \(\kappa(s)=\kappa(gs)=1\) and hence \(\kappa(g)=1\). Thus \(\kappa\) is unramified at all finite places of \(K\). Class field theory gives \(\kappa^{h_K}=1\), and since \(\ell\nmid h_K\) and \(|\kappa(g)-1|_\ell<1\), it follows that \(\kappa=1\) on \(G_K\). The same argument applied to \(\det(r|_{G_K})=\det(\operatorname{corner}(r|_{G_K}))\) (which equals \(1\) on every inertia group and satisfies \(|\det r(g)-1|_\ell<1\)) gives \(\det(\operatorname{corner}(r(g)_{\mathrm{mat}}))=1\) for all \(g\in G_K\). For any \(M\in\mathrm{SL}_2(C_\ell)\), \(J_2(M^{-1})^tJ_2=\operatorname{diag}(-1,1)M\operatorname{diag}(-1,1)\), where \(J_2=\bigl(\begin{smallmatrix}0&1\\1&0\end{smallmatrix}\bigr)\). Since \(r(cgc)_{\mathrm{mat}}=J(r(g)_{\mathrm{mat}}^{-1})^tJ\) for \(g\in G_K\), the formula for \(W\) defines a continuous homomorphism \(W:G_\Q\to\GL_2(C_\ell)\) extending the irreducible corner, with \(W(c)=\operatorname{diag}(-1,1)\), \(\operatorname{tr}W(c)=0\), and \(\det W=\eta\).

For (2), suppose \(r,r'\) are two-sided points of \(S\) and \(U\in\GL_3(C_\ell)\) satisfies \(Ur(g)_{\mathrm{mat}}U^{-1}=r'(g)_{\mathrm{mat}}\) for all \(g\in G_K\). Then \(U\) intertwines the diagonal marking matrices \(Z=\operatorname{diag}(\zeta,0,-\zeta)\) and \(Z'=\operatorname{diag}(\zeta',0,-\zeta')\), which therefore have the same eigenvalue set. Since \(\bar\zeta=\bar\zeta'\ne0\) and \(\ell\) is odd, \(\zeta'=\zeta\), so \(U=\operatorname{diag}(u_0,u_1,u_2)\) is diagonal. In Type~I, \(r'(cg^{-1}c)_{\mathrm{mat}}=UJU^t(r'(g)_{\mathrm{mat}})^t(UJU^t)^{-1}\); irreducibility of \(r'|_{G_K}\) forces \(UJU^t\) to be a scalar multiple of \(J\), so \(u_0u_2=u_1^2\). In Type~II, replacing \(u_1\) by \(\sqrt{u_0u_2}\in C_\ell^\times\) preserves \(Ur(g)_{\mathrm{mat}}U^{-1}\) on the block-diagonal matrices \(r(g)_{\mathrm{mat}}\) and achieves \(u_0u_2=u_1^2\). In both cases \(U=u_1\operatorname{diag}(t,1,t^{-1})\) with \(t=u_0/u_1\) (so \(t^{-1}=u_2/u_1\) by \(u_0u_2=u_1^2\)), whose conjugation preserves \(r(c)_{\mathrm{mat}}=I_3\) and scales weight-\(j\) entries by \(t^j\).

For (3), in weights \(j=-2,-1\), every character \(\lambda\) occurring in \(V_j\) has the form \(\lambda=\chi^{-1}\varepsilon_\ell^{-e}\) with \(\chi\) of finite order and \(e\in\{k,2k\}\ge1\). By \cite[Proposition~3.8]{bloch-kato-1990}, \(\dim_{C_\ell}H^1_f(\Q_\ell,C_\ell(\lambda))=\dim_{C_\ell}H^0(\Q_\ell,C_\ell(\lambda))+0=0\) because \(\lambda|_{I_\ell}\) has infinite order. Hence any class in \(H^1_f(\Q,C_\ell(\lambda))\) is a coboundary on all finite inertia groups and on \(D_\ell\), so it vanishes by Lemma~\ref{B6.31a}. In weight \(j=0\), we have \(V_0=C_\ell(1)\oplus C_\ell(\eta)^{\oplus2}\). For the trivial character, \(H^1_f(\Q_\ell,C_\ell)=H^1_{\mathrm{unr}}(\Q_\ell,C_\ell)\), so classes in \(H^1_f(\Q,C_\ell)\) are continuous unramified homomorphisms \(G_\Q\to C_\ell\), which vanish because \(\operatorname{Cl}(\Q)=1\). For \(C_\ell(\eta)\), restriction embeds \(H^1_f(\Q,C_\ell(\eta))\) into \(H^1_f(K,C_\ell)\) (since \([K:\Q]=2\) is invertible in \(C_\ell\)), and classes in \(H^1_f(K,C_\ell)\) are everywhere unramified continuous homomorphisms \(G_K\to C_\ell\), which factor through the finite group \(\operatorname{Cl}(K)\) and hence vanish \cite[Ch.~VI, Prop.~6.9]{neukirch-1999}.
\end{proof}

\subsection*{Proof of \ref{B6.40}}\label{prf:B6.40}\phantomsection
\begin{proof}
By \ref{B6.29} work over the unramified ring generated by \(\vartheta\). Choose the marking \(g_1\) from \ref{B6.52}. Use the marked deformation problem and distinguished local ring \(B=U/I\) of Definition~\ref{B6.40m}, with base lift \(r_0(g)=(\operatorname{diag}(\Psi(g),1,\Psi(g)^{-1}),\eta(g))\). By Lemma~\ref{B6.40n}(1)--(2), \(B=O[[X]]/(H)\) has a homogeneous presentation centred at \(r_0\) and a graded rational formal chart with \(1,2,3,4,2\) coordinates in weights \(-2,-1,0,1,2\), respectively.

Apply the relative presentation of \ref{B6.58} over \(B\) to global lifts vanishing on \(\mathcal K_{\mathrm{away}}\) before imposing the marking equations, with one framed place \(v_0=\ell\) (frame fixed to the identity) and local inertia kernels \(N_q=I_q\cap\ker\bar r_0\) at the finite ramified primes \(q\ne\ell\) dividing \(N|D_K|\). Let \(G_S\) be the maximal Galois quotient unramified outside \(S=\{\ell\}\cup\{q:q\mid N|D_K|\}\). Write \(Y\) for the extra coordinates and \(d_j=\dim_{\mathbf k}H^1_{\mathrm{fr},j}\) for their number in weight \(j\), chosen from matrix entries minus their base values.

We verify that the relative obstruction space \(\operatorname{Ob}_j\) of \ref{B6.58} has dimension \(d_j+\mathbf1_{|j|\le1}\). Let \(\bar V_j\) be the weight-\(j\) summand of \(\operatorname{ad}\bar r_0\) over \(\mathbf k\). Its complex-conjugation eigenspace decomposition is
\[
 \bar V_2=\mathbf k(\bar\Psi^2\eta),\quad
 \bar V_1=\mathbf k(\bar\Psi)\oplus\mathbf k(\bar\Psi\eta),\quad
 \bar V_0=\mathbf k(1)\oplus\mathbf k(\eta)^{\oplus2},\quad
 \bar V_{-1}=\mathbf k(\bar\Psi^{-1})\oplus\mathbf k(\bar\Psi^{-1}\eta),\quad
 \bar V_{-2}=\mathbf k(\bar\Psi^{-2}\eta),
\]
so \(\dim_{\mathbf k}\bar V_j^{c=1}=\mathbf1_{|j|\le1}\). We claim that \(H^0(G_S,\bar V_j^\vee(1))=0\) for every \(j\in\{-2,-1,0,1,2\}\). An even constituent \(\bar\Psi^j\) takes value \(1\) at \(c\), whereas \(\bar\varepsilon_\ell(c)=-1\ne1\). For an odd constituent \(\bar\Psi^j\eta\), since \(K\) is imaginary quadratic there is a prime \(p\mid D_K\); because \((N,D_K)=1\) and \(\ell\nmid D_K\), we have \(p\nmid\ell N\). On \(I_p\), both \(\bar\varepsilon_\ell\) and \(\bar\vartheta\) (hence \(\bar\Psi\)) are trivial, while \(\eta|_{I_p}\) takes the value \(-1\ne1\). Thus no constituent of \(\bar V_j\) equals \(\bar\varepsilon_\ell\), proving \(H^0(G_S,\bar V_j^\vee(1))=0\).

By Poitou--Tate duality \ref{B6.09}, the map \(H^2(G_S,\bar V_j)\to\bigoplus_{v\in S}H^2(D_v,\bar V_j)\) is surjective, so the image of \(\operatorname{Ob}_j\to H^2(G_S,\bar V_j)\) has dimension \(h^2(G_S,\bar V_j)-\sum_{v\in S}h^2(D_v,\bar V_j)\). The kernel of \(\operatorname{Ob}_j\to H^2(G_S,\bar V_j)\) is the cokernel of
\[
 Z^1(G_S,\bar V_j)\longrightarrow H^1(D_\ell,\bar V_j)\oplus\bigoplus_{q\in S\setminus\{\ell\}}Z^1(D_q,\bar V_j)/Z^1_{\mathcal L,q},
\]
where \(Z^1_{\mathcal L,q}=\{e_q\in Z^1(D_q,\bar V_j):e_q|_{N_q}=0\}\). For each \(q\in S\setminus\{\ell\}\), \([I_q:N_q]\) divides \(2\operatorname{ord}\vartheta\) and is prime to \(\ell\), so \(\dim Z^1_{\mathcal L,q}=\dim\bar V_j\) and local Euler characteristic gives \(\dim Z^1(D_q,\bar V_j)-\dim Z^1_{\mathcal L,q}=h^2(D_q,\bar V_j)\). Let \(W\subseteq Z^1(G_S,\bar V_j)\) be the kernel of the displayed map. Every global coboundary \(\delta v\) (\(v\in\bar V_j\)) vanishes on each \(N_q\) and restricts to \(B^1(D_\ell,\bar V_j)\), so \(W\to B^1(D_\ell,\bar V_j)\) is surjective with kernel \(H^1_{\mathrm{fr},j}\); hence \(\dim W=d_j+\dim\bar V_j-h^0(D_\ell,\bar V_j)\). Combining this with the global Euler characteristic \(h^1(G_S,\bar V_j)-h^2(G_S,\bar V_j)-h^0(G_S,\bar V_j)=\dim\bar V_j-\dim\bar V_j^{c=1}\) and the local Euler characteristic \(h^1(D_\ell,\bar V_j)-h^2(D_\ell,\bar V_j)=\dim\bar V_j+h^0(D_\ell,\bar V_j)\) cancels all local terms and gives
\[
 \dim_{\mathbf k}\operatorname{Ob}_j=d_j+\dim_{\mathbf k}\bar V_j^{c=1}=d_j+\mathbf1_{|j|\le1}.
\]

We next count the marking equations. Since \(c^2=1\) and \(\eta(c)=-1\), the matrix \(A=r(c)_{\mathrm{mat}}\) satisfies \(A\theta(A)=I_3\), i.e.\ \(JA^tJ=A\), so \(A_{ab}=A_{2-b,2-a}\). Hence the condition \(A=I_3\) is cut out by \(6\) homogeneous equations \(A_{20}=0\), \(A_{10}=0\), \(A_{00}-1=A_{11}-1=0\), \(A_{01}=0\), \(A_{02}=0\), with weight multiplicities \(1,1,2,1,1\). Together with the two marking equations at \(g_1\) in positions \((1,0)\) (weight \(-1\)) and \((0,1)\) (weight \(1\)), the marking adds \(8\) equations with multiplicities \(1,2,2,2,1\) in weights \(-2,-1,0,1,2\), respectively.

Put \(R=B[[Y]]\), let \(A=\bigoplus_j A_j\subset R\) be the sum of its homogeneous pieces, and restrict the quotient map to \(q:A\to S\), where \(S\) is the marked global ring with universal lift \(r_S\). The local chart of Lemma~\ref{B6.40n}(2) extends by the free coordinates \(Y\) to the rational formal chart for \(A\). With
\[
 s=d_{-2}+1,\quad r=d_{-1}+2,\quad m=d_0+3,\quad
 r'=d_1+3,\quad s'=d_2+1
\]
in the notation of the orbit estimate of \ref{B6.58}, the chart has \(s,r,m,r'+1,s'+1\) coordinates of weights \(-2,-1,0,1,2\), with \(r\ge1\). The kernel \(J\) of \(q\) is generated by the relative and marking equations, whose counts in the five weights are \((s,r+1,m,r',s')\), giving the families \(G,F,Z,E,E'\) of \ref{B6.58}. Let \(x_0:S\to O\) be the base point and \(x=x_0q\). Write \(C_{-1}(F)\) for the cotangent quotient by the weight-\(-1\) equations \(F_i\).
\leanref{UpperTheory.\allowbreak{}V.\allowbreak{}Loc.\allowbreak{}rtChart}{Development.\allowbreak{}AutomorphicArithmetic.\allowbreak{}AutomorphicVarieties.\allowbreak{}PadicHodgeTheory.\allowbreak{}RtChart}{41--74}
\leanref{UpperTheory.\allowbreak{}V.\allowbreak{}U3.\allowbreak{}exists\_\allowbreak{}model}{Development.\allowbreak{}AutomorphicArithmetic.\allowbreak{}AutomorphicVarieties.\allowbreak{}UnitaryThreeGroups.\allowbreak{}ModelEx}{156--178}
\leanref{UpperTheory.\allowbreak{}V.\allowbreak{}U3.\allowbreak{}model\_\allowbreak{}of\_\allowbreak{}pres}{Development.\allowbreak{}AutomorphicArithmetic.\allowbreak{}AutomorphicVarieties.\allowbreak{}UnitaryThreeGroups.\allowbreak{}ModelOfPres}{320--329}

By \ref{B6.40a} and \ref{B6.58}, \(\operatorname{length}_O\mathcal S_{\ell,k,\vartheta}[b]\le\operatorname{length}_O C_{-1}(F)\) for \(b\ne0\). Continuous cocycles into the discrete coefficient module have finite image. Thus every class belongs to some \(\mathcal S_{\ell,k,\vartheta}[\ell^j]\) and
\[
 L_{\ell,k,\vartheta}\le \operatorname{length}_O C_{-1}(F).
 \tag{6.40.2}\label{B6.40-length}
\]

The residue-characteristic assertion of \ref{B6.52} gives finiteness over \(O\) of the subalgebra of \(S\) generated by weight-zero products of matrix entries. The ring \(S\) is Hausdorff and this finite \(O\)-module is compact. The image \(q(A_0)\) lies in it by approximation by homogeneous polynomials, so \(q(A_0)\) is finite over \(O\).

Let \(\xi:A/J\to C_\ell\) be an \(O\)-point. Write \(z_i\in A\) for the images of the centred coordinates and \(j_i\) for their weights. Each nonconstant monomial of weight zero lies in the maximal ideal of \(A_0\). Its image under \(\xi\) has positive valuation since the local algebra \(q(A_0)\) is finite over \(O\). For nonzero coordinate values put \(v_i=v_\ell(\xi(z_i))\). A positive-weight coordinate and a negative-weight coordinate give a weight-zero monomial \(z_i^{-j_h}z_h^{j_i}\). Its positive valuation gives
\[
 -v_i/j_i < v_h/(-j_h) \qquad (j_i>0>j_h).
\]
Choose a rational number above all the lower bounds on the left and below all the upper bounds on the right. Take \(t\in C_\ell^\times\) of that valuation. Now \(t^{j_i}\xi(z_i)\) has positive valuation whenever \(\xi(z_i)\ne0\), including the weight-zero coordinates. These scaled values define evaluation on \(R\) in the integers of a finite extension of \(\operatorname{Frac}O\). To see agreement on \(A\), group each homogeneous series into finitely many monomials times weight-zero series. The weight-zero coordinate ideal is generated by finitely many nonconstant weight-zero monomials. Both evaluations send those monomials to elements of positive valuation, so agreement on polynomials gives agreement on weight-zero series. The evaluation kills the relative and marking equations, so it descends to a local point \(f:S\to C_\ell\) with \((t\cdot\xi)(z)=f(qz)\) for \(z\in A\).

At \(f\circ q\), a tangent derivation \(D\) killing \(J\) gives values \(f(qz_i)+\epsilon D(z_i)\) on the coordinates, with \(\epsilon^2=0\). The map on \(A_0\) with this first-order term factors through \(q(A_0)\). Choose a finite extension \(E'\) containing its values on module generators and the finitely many coordinate values and derivatives. In \(C_\ell[\epsilon]/(\epsilon^2)\) these lie in a complete ring \(O_{E'}\oplus\epsilon\varpi_{E'}^{-a}O_{E'}\) for some \(a\ge0\), where \(O_{E'}\) is the ring of integers and \(\varpi_{E'}\) a uniformiser. The coordinate values reduce to the maximal ideal modulo \(\epsilon\), so formal evaluation converges. It agrees with \(f\circ q+\epsilon D\) on homogeneous series and kills the equations. Thus it defines a lift \(S\to C_\ell[\epsilon]/(\epsilon^2)\). At \(\xi\) use the same \(t\) on tangent vectors, multiplying their values on weight \(j\) by \(t^j\).
\leanref{UpperTheory.\allowbreak{}V.\allowbreak{}U3.\allowbreak{}Model}{Development.\allowbreak{}AutomorphicArithmetic.\allowbreak{}AutomorphicVarieties.\allowbreak{}UnitaryThreeGroups.\allowbreak{}ConjStable}{54--85}
\leanref{UpperTheory.\allowbreak{}V.\allowbreak{}U3.\allowbreak{}model\_\allowbreak{}upper}{Development.\allowbreak{}AutomorphicArithmetic.\allowbreak{}AutomorphicVarieties.\allowbreak{}UnitaryThreeGroups.\allowbreak{}ModelUpper}{103--250}

Put \(r_f=f_*r_S\). Negative weights vanish at \(f\circ q\) precisely when \(r_f\) is upper triangular on \(G_K\), and positive weights vanish precisely when it is lower triangular there. This follows from the weights of the matrix entries and their homogeneous products. The marked diagonal point is \(r_0\), so the only fixed solution is \(x\). This gives the one-sided fibre condition of \ref{B6.58}.

By Lemma~\ref{B6.40n}(3) and finite-quotient averaging on inertia away from \(\ell\), every lifted tangent cocycle at a characteristic-zero point \(f\) lies in the Bloch--Kato finite Selmer group \(Z^1_f(G_\Q,\operatorname{ad}r_f)\). At \(x\), a derivation of weight \(j\le0\) killing \(J\) lifts to a normalised cocycle in \(Z^1_f(G_\Q,V_j)\), since the base matrices are diagonal. By Lemma~\ref{B6.40p}(3), \(H^1_f(\Q,V_j)=0\) for \(j\le0\), so the cocycle is a coboundary. The marking forces its conjugating vector into the Lie algebra generated by \(\operatorname{diag}(1,0,-1)\), which centralises \(r_0\). Thus the cocycle vanishes, as does the derivation, since the matrix entries determine the lift. Duality now gives the required rational spans: a functional on the weight-\(-1\) cotangent space killing the \(F_i\) extends by zero in the other weights to a derivation killing \(J\), and the same argument applies to the \(G_i\) in weight \(-2\) and the \(Z_i\) in weight \(0\).

At a two-sided point, Lemma~\ref{B6.40p}(1) classifies \(r_f\) as Type~I or Type~II. Use adjoint rigidity in Type~I. In Type~II, Lemma~\ref{B6.40p}(1) shows that the middle line is trivial on \(G_K\) and the corner extends to \(W\) of determinant \(\eta\) and \(\operatorname{tr}W(c)=0\). The adjoint decomposes into the polarised adjoint of the corner (killed by block adjoint rigidity), the two off-block summands \(W\) and \(W\otimes\eta\) (killed by \ref{B6.43}, using Lemma~\ref{B6.40n}(3), Lemma~\ref{B6.30ac}, and \ref{B6.52}(3)), and the \((1,1)\) diagonal summand \(C_\ell(\eta)\) (killed by Lemma~\ref{B6.40p}(3)). Each lifted tangent cocycle is therefore a coboundary with conjugating vector in the one-dimensional torus Lie algebra generated by \(\operatorname{diag}(1,0,-1)\). The cocycle determines the coordinate derivatives, so the tangent space at \(\xi\) has dimension at most one.
\leanref{UpperTheory.\allowbreak{}V.\allowbreak{}Z.\allowbreak{}J.\allowbreak{}ptHyp\_\allowbreak{}of\_\allowbreak{}autB\_\allowbreak{}of\_\allowbreak{}lemW\_\allowbreak{}of\_\allowbreak{}isBase}{Development.\allowbreak{}AutomorphicArithmetic.\allowbreak{}AutomorphicVarieties.\allowbreak{}ModularityArgument.\allowbreak{}ThresholdArguments.\allowbreak{}DenseThresholdCases.\allowbreak{}FPtHypU}{15--40}

Type~I admissibility and the Type~II endpoint of \ref{B6.52} put the two-sided points in \(\mathcal A_3\cup\mathcal A_{2+1}\). Take a weight-zero lift \(a\in A\) of \(\operatorname{tr}r_S(\sigma)-T\). It vanishes at \(x\) and is nonzero at the two-sided points, so the orbit estimate of \ref{B6.58} applies. Lemma~\ref{B6.40p}(2) places points with isomorphic \(G_K\)-restrictions in the same \(\mathbb G_m\)-orbit. The orbit representatives give \(X\). Combine their bound with \eqref{B6.40-length}. Over the unramified ring the length and valuation are already normalised, so \ref{B6.29} returns the bound to the original \(O\).
\end{proof}

\subsection*{Proof of \ref{B6.40a}}\label{prf:B6.40a}\phantomsection
\begin{proof}
Put \(\alpha=\Psi^{-1}\). At \(g_1\) and at a test element of the distinguished local group, \(\Psi^2-1\) is a unit. Subtract the unique coboundary making the Selmer cocycle zero at \(g_1\). By uniqueness the normalised cocycle is \(b\)-torsion. Its local form is \(s-\alpha(h)s\), so the local unit test gives \(bs=0\). Identify coefficient \(b\)-torsion with \(O/(b)\) and call the cocycle \(t\). On the prescribed kernels \(\mathcal K_{\mathrm{away}}\), \(\alpha=1\) and the local coboundary vanishes. At complex conjugation \(c\), the identities \(c^2=1\), \(\Psi(c)=1\), and \(2\in O^\times\) give \(t(c)=0\), hence \(t(cg_1^{-1}c)=0\).

In zero-based matrix units, put \(V=e_{10}-e_{21}\). Negative reflected transpose \(X\mapsto-JX^tJ\) fixes \(V\), and conjugation by \(r_0\) acts on \(V\) by \(\alpha\). Hence \((1+\epsilon tV)r_0\) over \(O\oplus\epsilon O/(b)\) (\(\epsilon^2=0\)) satisfies the marked local and global conditions modulo \(b\). Choose a lift \(\widetilde s\in O\) of \(s\). By Lemma~\ref{B6.40n}(1), the defining ideal \(I\) of \(B=O[[X]]/(H)\) is invariant under conjugation by \(I_3+\epsilon\widetilde sV\), so the infinitesimal conjugation \((I_3+\epsilon\widetilde sV,1)r_0|_{D_\ell}(I_3+\epsilon\widetilde sV,1)^{-1}\) over \(O[\epsilon]/(\epsilon^2)\) defines an \(O\)-valued derivation on \(B\) at \(x_{0,\ell}\) supported in weight \(-1\) that kills \(H\). The ambient algebra \(A\) lies in \(O[[X,Y]]/(H)\), where \(Y\) are the extra global variables. Take this local derivation on \(X\), lift the values on \(Y\) in weight \(-1\) to \(O\), and assign value zero to all variables of other weights. The resulting linear-coefficient functional kills \(H\), defines an \(O\)-valued derivation at \(x\) supported in weight \(-1\), and reduces modulo \(b\) to the derivation through \(S\), which kills \(F_i\). The \((1,0)\) entry at \(g\) has derivative \(r_0(g)_{\mathrm{mat},00}\,t(g)=\Psi(g)t(g)\), which recovers \(t(g)\) because \(\Psi(g)\) is a unit.
\end{proof}

\subsection*{Proof of \ref{B6.41-valuation}}\label{prf:B6.41-valuation}\phantomsection
\begin{proof}
Use \(S=S_1\sqcup S_2\) and \(t(s)\) from \ref{B6.41a}. Put \(R=q^k\) and \(F_{u,s}=Rt(s)-z_u-R-R^2z_u^{-1}\). By \eqref{B6.41-traces} these factors are integral and
\[
 |Rt(s)-z_u-R|\le R(\sqrt q+3)+1<R^2
\]
since \(R\ge q\ge7\). Thus \(0<|F_{u,s}|<2R^2\). The product \(D=\prod_{u,s}F_{u,s}\) is fixed by all complex field automorphisms, so \(D\in\Z\setminus\{0\}\) and \(\log|D|\le\#\mathcal O\,\#S(2k\log q+1)\). By \eqref{B6.41-labels} and \(q=p^f\), \(1\le f\le2\), take \(c=5(c_1+c_2|D_K|^{A_2})\) and \(A=A_1+A_2+1\). For each \(u\) the labels of \(X_u\) are distinct. Under \(\iota^{-1}\) their factors are \(q^k\Delta_{u,\rho}\). The multiplier is an \(\ell\)-adic unit and the other factors have nonnegative valuation.
\leanref{UpperTheory.\allowbreak{}GalN.\allowbreak{}V.\allowbreak{}Win.\allowbreak{}frobenius\_\allowbreak{}trace\_\allowbreak{}norm\_\allowbreak{}nonzero\_\allowbreak{}log\_\allowbreak{}bound}{Development.\allowbreak{}AutomorphicArithmetic.\allowbreak{}Foundations.\allowbreak{}GaloisRepresentations.\allowbreak{}AdmissibleResidues.\allowbreak{}FrobeniusTraceCountBoundOfRatEigensystemCount}{329--400}
\end{proof}

\subsection*{Proof of \ref{B6.41b}}\label{prf:B6.41b}\phantomsection
\begin{proof}
Because \(r\) satisfies the witness condition of Definition~\ref{def:p6c5-witness} and every place \(w\mid\ell\) is unlinked or potentially unlinked, every \(w\mid\ell\) is potentially unlinked. Hence at every finite place \(w\) of \(K\), the cuspidal support of \(\pi_w\) consists of characters \(\xi\) of \(K_w^\times\) satisfying \(q_w^{t_\pi-1/2}<|\xi(\varpi_w)|<q_w^{t_\pi+1/2}\) for a single \(t_\pi\in\R\). The ratio of moduli of any two support characters lies strictly between \(q_w^{-1}\) and \(q_w\), so no two belong to a nontrivial Bernstein--Zelevinsky segment. By \ref{B9.47}, the local Langlands parameter of \(\pi_w\) is a sum of characters with zero monodromy, and \(\pi_w\) is an irreducible principal series representation induced from those characters.

Since \(\pi\) is regular algebraic, cuspidal, and polarised, and \(r\) is irreducible, Chebotarev density and Brauer--Nesbitt identify \(r\) with the Galois representation of \ref{B9.02} attached to \(\pi\), so local-global compatibility (\ref{B9.02}, \ref{B9.04}, \ref{B9.06}) applies to \(r\) at every finite place \(w\) of \(K\): for \(w\nmid\ell\) the Frobenius-semisimplified Weil--Deligne representation of \(r|_{G_{K_w}}\) matches the local Langlands parameter of \(\pi_w\) under \(\iota\), and for \(w\mid\ell\) (via the toral zero-monodromy clause of \ref{B9.06}) \(r|_{G_{K_w}}\) is potentially crystalline with matching Weil--Deligne representation and with Hodge--Tate weights at each \(\tau:K\hookrightarrow C_\ell\) inducing \(w\) matching the infinitesimal weights of \(\pi\) at \(\iota\circ\tau\).

For (1), given \(\varphi:K\hookrightarrow\C\), set \(\tau=\iota^{-1}\circ\varphi\) and let \(w\mid\ell\) be the place induced by \(\tau\). On \(G_{K_w}\), \(r\) is de Rham and preserves a flag whose graded pieces have inertia characters \(\varepsilon_\ell^{m_i}\chi_i\) with \(\chi_i\) of finite order. The Hodge--Tate weights of a de Rham representation with a stable flag are those of its graded pieces, unaffected by finite-order twists, so they equal \(m=(m_1,\ldots,m_n)\). Local-global compatibility \ref{B9.04} then identifies the infinitesimal character of \(\pi\) at \(\varphi\) as that determined by \(m\).

For (2), by the restricted tensor product decomposition of \(\pi_{\mathrm{fin}}\) it suffices to show that \(\pi_w\) has a nonzero \(K_w(M)\)-fixed vector at every finite place \(w\), where \(K_w(M)=\ker(\GL_n(\mathcal O_{K_w})\to\GL_n(\mathcal O_{K_w}/M\mathcal O_{K_w}))\).
\begin{itemize}
\item Suppose \(w\nmid\ell\). By the Galois level hypothesis, \(r|_{I_w}\) is diagonal with characters \(\chi_i\) factoring through the action on \(\mu_M\). Since the Weil--Deligne parameter of \(\pi_w\) has zero monodromy and matches that of \(r|_{G_{K_w}}\), the inducing characters \(\xi_i\) of the principal series \(\pi_w\) correspond under local class field theory to \(\chi_i\) on \(\mathcal O_{K_w}^\times\) and therefore vanish on the local norm group of \(K_w(\mu_M)/K_w\). Writing \(M=p^aM'\) with \(p\) the residue characteristic of \(w\) and \(p\nmid M'\), the extension \(K_w(\mu_{M'})/K_w\) is unramified, while \(N_{K_w/\Q_p}(1+p^a\mathcal O_{K_w})\subseteq1+p^a\Z_p\) acts trivially on \(\mu_{p^a}\) under local reciprocity. Thus every \(\xi_i\) is trivial on \((1+M\mathcal O_{K_w})\cap\mathcal O_{K_w}^\times\) (which is \(\mathcal O_{K_w}^\times\) when \(w\nmid M\)). Every element of \(B(K_w)\cap K_w(M)\) has diagonal entries in \((1+M\mathcal O_{K_w})\cap\mathcal O_{K_w}^\times\), on which both \(\delta_B^{1/2}\) and every \(\xi_i\) are trivial; hence the function in \(\pi_w=\operatorname{Ind}_B^{\GL_n}(\xi_1\otimes\cdots\otimes\xi_n)\) supported on the open double coset \(B(K_w)K_w(M)\) with value \(1\) on \(K_w(M)\) is well-defined, right \(K_w(M)\)-invariant, and nonzero, giving \(\pi_w^{K_w(M)}\ne0\).
\item Suppose \(w\mid\ell\). Since \(r|_{G_{K_w}}\) is potentially crystalline and preserves a flag with graded inertia characters \(\varepsilon_\ell^{m_i}\chi_i\), exactness of \(D_{\mathrm{pst}}\) shows that the inertia action on its Weil--Deligne representation has the same characters as \(\bigoplus_i D_{\mathrm{pst}}(\kappa_i)\), where \(\kappa_i|_{I_w}=\varepsilon_\ell^{m_i}\chi_i\). The crystalline factor \(\varepsilon_\ell^{m_i}\) acts trivially on \(D_{\mathrm{pst}}\) inertia, leaving the finite-order characters \(\chi_i|_{I_w}\). If \(\ell\mid M\), each \(\chi_i\) factors through the action on \(\mu_M\), so the same norm calculation as above shows that every inducing character \(\xi_i\) of \(\pi_w\) is trivial on \(1+M\mathcal O_{K_w}\), giving \(\pi_w^{K_w(M)}\ne0\). If \(\ell\nmid M\), then \(K_w(\mu_M)/K_w\) is unramified, so \(I_w\) fixes \(\mu_M\) and \(\chi_i|_{I_w}=1\) for every \(i\); with zero monodromy, \(\pi_w\) is an unramified principal series and therefore has a nonzero \(\GL_n(\mathcal O_{K_w})\)-fixed vector.
\end{itemize}

For (3), if \(w\nmid\ell\) and \(\pi_w\) is spherical with Satake multiset \(B\), then the local Langlands parameter of \(\pi_w\) is unramified with inverse Frobenius eigenvalues \(q_w^{(n-1)/2}B\) and zero monodromy. Local-global compatibility at \(w\nmid\ell\) implies that \(r\) is unramified at \(w\) and that \(\iota\) sends the eigenvalue multiset of \(r(\sigma_w^{-1})\) to \(q_w^{(n-1)/2}B\).

For (a), let \(\eta_{K/K^+}\) be the quadratic idele class character of \(K/K^+\). Its base change \(\eta_{K/K^+}\circ N_{K/K^+}\) to \(K\) is trivial and its value at \(-1\) at every real place of \(K^+\) is \(-1\), so \(\chi=\eta_{K/K^+}^n\) satisfies \(\chi_v(-1)=(-1)^n\) and polarises \(\pi\).

For (b), restricting \(r_0\) from \(G_\Q\) to \(G_{K'}\) preserves the Galois level \((M;m)\) because every inertia and decomposition group of \(K'\) is conjugate into the corresponding subgroup over \(\Q\). At any prolongation \(A\) of a place \(w\mid\ell\) of \(K'\), the flag characters \(\kappa_i\) satisfy \(\kappa_i|_{I_A}=\varepsilon_\ell^{m_i}\chi_i|_{I_A}\) with \(\chi_i\) of finite order; on an open subgroup of \(I_A\) killing \(\chi_j/\chi_i\), an equality \(\kappa_i=\kappa_j\varepsilon_\ell\) would force \(\varepsilon_\ell^{m_i-(m_j+1)}=1\), contrary to \(m_i\ne m_j+1\) and the infinite image of \(\varepsilon_\ell\) on \(I_A\). Conjugating across all prolongations of \(w\) shows that \(w\) is unlinked for \(r\) in the first case. In the second case, choose an imaginary quadratic field \(K_0\) in which \(\ell\) splits and on which \(r_0\) remains irreducible. The decomposition groups of \(K_0\) above \(\ell\) are those of \(\Q\), so \(r_0|_{G_{K_0}}\) is unlinked above \(\ell\), and the quadratic witness condition makes it potentially unlinked there. Hence the flag characters of \(r_0\) at \(\ell\) satisfy \(\kappa_i^a\ne(\kappa_j\varepsilon_\ell)^a\) for all \(i,j\) and all \(a\ge1\). If \(H\) is an open subgroup of a decomposition group \(D\) and a character \(\chi\) of \(D\) satisfies \(\chi^a|_H=1\), then \(\chi^a\) factors through the finite quotient of \(D\) by the normal core of \(H\), so \(\chi^{ad}=1\) for the order \(d\) of that quotient. Applied to \(\chi=\kappa_i/(\kappa_j\varepsilon_\ell)\) and \(H=D\cap G_{K'}\), this shows that every place of \(K'\) above \(\ell\) is potentially unlinked for \(r\). Finally, \(\theta^{-1}\varepsilon_\ell^{1-n}\) is unramified outside \(M\ell\) of cyclotomic weight \(-\mathrm{wt}+1-n\), so by global class field theory there is an algebraic Hecke character \(\chi'\) of \(\Q\) with \(\chi'_p(p)=p^{n-1}\iota(\theta(\sigma_p^{-1}))^{-1}\) for \(p\nmid M\ell\), which we multiply by \(\eta_{K'/\Q}\) if needed to arrange \(\chi'_\infty(-1)=(-1)^n\) while preserving \(\chi'\circ N_{K'/\Q}\). Comparing Satake parameters at conjugate unramified places of \(K'\) via \(r_0(g)^tPr_0(g)=\theta(g)P\) and applying strong multiplicity one \cite{jacquet-shalika-1981} gives \(\pi^\vee\simeq\pi^c\otimes(\chi'\circ N_{K'/\Q}\circ\det)\). Hence \(\pi\) is polarised and satisfies (1)--(3).
\end{proof}

\subsection*{Proof of \ref{B6.41-labels}}\label{prf:B6.41-labels}\phantomsection
\begin{proof}
Use spherical Hecke systems away from the level as CM labels. For each \(\rho\in\mathcal A_3(K,\iota;k,N)\), the twist \(r=\rho\varepsilon_\ell^{-1}\) is irreducible of Galois level \((N;(k-1,-1,-k-1))\), is separated ordinary (so every place above \(\ell\) is unlinked), and has a regular algebraic conjugate self-dual cuspidal attachment \(\pi\) with witness. By Lemma~\ref{B6.41b}(a), \(\pi\) has infinitesimal character determined by \((k-1,-1,-k-1)\) at every complex embedding, possesses a nonzero \(K(N)\)-fixed vector, and has algebraic Satake parameters \(q_vB_v\) matching the inverse Frobenius eigenvalues of \(r\) at every finite place \(v\nmid N\ell\) (since \(\pi_v\) is spherical whenever \(v\nmid N\)). These Frobenius polynomials determine the semisimple representation \(r\) (hence \(\rho\)) up to isomorphism by \ref{B9.02} and Chebotarev density.

For each regular algebraic cuspidal constituent \(\pi\) in this automorphic family, \ref{B6.06b}(2) conjugates its finite Hecke system and permutes the embeddings indexing its infinitesimal weights, keeping it in the family. For each fixed infinitesimal-weight family, choose the algebraic coefficient representation with the negative infinitesimal multisets. Proposition~\ref{B6.06b}(1) gives cohomology with these common coefficients. Since \(K(3N)\subset K(N)\), a nonzero \(K(N)\)-fixed vector is fixed by the neat compact open group \(K(3N)\). Choose one constituent per distinct Hecke system. They are pairwise nonisomorphic, so \ref{B6.06b}(3) bounds every finite selection by the cohomology dimension at level \(3N\), and hence bounds the whole set \(S_1\). Bound the cochain dimension by the cell count at level \(3\) times the covering index and coefficient dimension. Reduction modulo \(3N\) bounds the index polynomially in \(N\), and the Weyl dimension formula bounds the coefficient dimension polynomially in \(k+1\). The same bound holds in any fixed rank \(n\ge2\) for regular algebraic polarised cuspidal automorphic representations over a CM field at level \(3M\) and weight bound \(H\).

Now consider the rational label count in rank \(n\ge2\) for irreducible representations \(W:G_\Q\to\GL_n(C_\ell)\) of Galois level \((M;m)\) satisfying one of the two alternatives of Lemma~\ref{B6.41b}(b), quadratic automorphy and witness, and twisted self-duality \(W(g)^tPW(g)=\theta(g)P\). Choose \(n^2+2\) imaginary quadratic fields of pairwise coprime discriminants and retain \(n^2+1\) unramified at \(p\). Distinct quadratic self-twists occupy independent character lines in the \(n^2\)-dimensional endomorphism representation. The index-two restriction criterion gives an irreducible restriction \(W|_{G_{K'}}\) in at least one retained field \(K'\). By Lemma~\ref{B6.41b}(b), the cuspidal attachment \(\pi\) of \(W|_{G_{K'}}\) over \(K'\) is polarised, has infinitesimal character determined by \(m\), has a nonzero \(K(M)\)-fixed vector, and satisfies Frobenius--Satake compatibility at every place \(v\nmid M\ell\) of \(K'\). Hence \(W|_{G_{K'}}\) receives a CM label from the finite conjugation-stable set of Hecke systems over \(K'\) of size at most \(P(M,H)\). Any irreducible restriction to \(G_{K'}\) has at most two extensions to \(G_\Q\) up to isomorphism, interchanged by the quadratic twist if both occur. The restricted inverse Frobenius roots at a place \(v\mid p\) of \(K'\) are the \(f\)-th powers of the rational roots, with residue degree \(f\le2\), so each multiset admits at most \(2^n\) lifts. Taking two copies of each pair consisting of a restricted Hecke system and a lifted multiset yields a rational label set of size at most \(2^{n+1}(n^2+1)P(M,H)\), on which automorphisms of \(\C\) act on the pair and fix the copy index. For \(\rho\simeq W|_{G_K}\oplus1\in\mathcal A_{2+1}(K,\iota;k,N|D_K|)\), the representation \(W\) is irreducible over \(G_\Q\) of level \((N|D_K|;(k,-k))\) with \(k-(-k)=2k\ge2\ne1\), satisfies \(W(g)^tJ_0W(g)=(\det W(g))J_0\) for \(J_0=\bigl(\begin{smallmatrix}0&1\\-1&0\end{smallmatrix}\bigr)\) with \(\det W\) of cyclotomic weight zero, and satisfies the quadratic automorphy and witness conditions (via quadratic base change of its weight-\((2k+1)\) newform attachment); applying this rank-\(2\) rational count at level \(M=N|D_K|\) gives \(S_2\) and the second bound of \eqref{B6.41-labels}.

The support bounds and purity put the three inverse Frobenius root moduli in \([q^{-1/2},q^{1/2}]\) with product one. The extreme logarithmic triples are the permutations of \((\tfrac12\log q,0,-\tfrac12\log q)\). Convexity bounds the sum of moduli by \(\sqrt q+1+q^{-1/2}\). For the block the two inverse rational roots satisfy the bounds at \(p\) and have product of moduli one. Raise to the residue degree at \(w\) and add the trivial root to obtain the same bound. The inverse spherical operators on the dual integral lattice give integral inverse roots after scaling by the residue norm to \(\max_i m_i\) at weights \((m_i)\). Use \(k-1\) for the twisted rank-three roots, since undoing the twist on inverse Frobenius contributes \(q^{-1}\). In rank two \(p^k\) times an inverse rational root is integral. Raising it to the residue degree gives the scaling of \eqref{B6.41-traces}.
\leanref{UpperTheory.\allowbreak{}GalN.\allowbreak{}V.\allowbreak{}Win.\allowbreak{}N.\allowbreak{}countK\_\allowbreak{}alt}{Development.\allowbreak{}AutomorphicArithmetic.\allowbreak{}AutomorphicVarieties.\allowbreak{}ModularityArgument.\allowbreak{}ThresholdArguments.\allowbreak{}DenseThresholdCases.\allowbreak{}FrobeniusTraceCountBoundOfCmEigensystemCount}{63--125}
\leanref{UpperTheory.\allowbreak{}GalN.\allowbreak{}V.\allowbreak{}Win.\allowbreak{}N.\allowbreak{}frobenius\_\allowbreak{}trace\_\allowbreak{}count\_\allowbreak{}bound\_\allowbreak{}of\_\allowbreak{}cm\_\allowbreak{}eigensystem\_\allowbreak{}count\_\allowbreak{}alt}{Development.\allowbreak{}AutomorphicArithmetic.\allowbreak{}AutomorphicVarieties.\allowbreak{}ModularityArgument.\allowbreak{}ThresholdArguments.\allowbreak{}DenseThresholdCases.\allowbreak{}FrobeniusTraceCountBoundOfCmEigensystemCount}{418--461}
\leanref{UpperTheory.\allowbreak{}GalN.\allowbreak{}V.\allowbreak{}Win.\allowbreak{}N.\allowbreak{}frobenius\_\allowbreak{}trace\_\allowbreak{}count\_\allowbreak{}bound\_\allowbreak{}of\_\allowbreak{}rat\_\allowbreak{}eigensystem\_\allowbreak{}count\_\allowbreak{}alt}{Development.\allowbreak{}AutomorphicArithmetic.\allowbreak{}AutomorphicVarieties.\allowbreak{}ModularityArgument.\allowbreak{}ThresholdArguments.\allowbreak{}DenseThresholdCases.\allowbreak{}FrobeniusTraceCountBoundOfCmEigensystemCount}{348--398}
\leanref{UpperTheory.\allowbreak{}V.\allowbreak{}Z.\allowbreak{}Win.\allowbreak{}A4.\allowbreak{}newAdmS\_\allowbreak{}alt}{Development.\allowbreak{}AutomorphicArithmetic.\allowbreak{}AutomorphicVarieties.\allowbreak{}ModularityArgument.\allowbreak{}ThresholdArguments.\allowbreak{}DenseThresholdCases.\allowbreak{}NewAdmS}{18--25}
\end{proof}

\subsection*{Proof of \ref{B6.42}}\label{prf:B6.42}\phantomsection
\begin{proof}
The datum requires \(K\) imaginary quadratic, so take \(D=1\) in all other fields. For \(w\mid p\) use \ref{B6.41} with \(z_\vartheta=\vartheta(\mathrm Nw)^{-1}\) for complex characters of kernel \(H\). There are \(\varphi([G:H])\) for cyclic \(G/H\) and zero otherwise. A good \(O\)-valued character supplies all orbit values, and the datum conditions are invariant under powers coprime to the character order. At inverse Frobenius, \(\Psi(\sigma^{-1})\) corresponds to \(q^{-k}z_\vartheta\). Sum \ref{B6.40} over the orbit and apply \eqref{B6.41-valuation} to the resulting sets.

For two distinct auxiliary primes with integers \(D,D'\), let \(S_H\) consist of the primes admitting good data for \(H\) in the chosen coefficient rings. Apply the bound from a prime different from \(\ell\) for each \(\ell\in S_H\). Prime factorisation and disjoint character sets give
\[
 \sum_{\ell\in S_H}\log\ell\sum_{\ker\vartheta=H}a_{\ell,k,\vartheta}
 \le2\log|D|+2\log|D'|,
 \qquad
 \sum_{H\ \mathrm{occurring}}\varphi([G:H])\le \#G\le N.
 \tag{6.42.3--4}
\]
\leanref{UpperTheory.\allowbreak{}V.\allowbreak{}U3.\allowbreak{}corGK}{Development.\allowbreak{}AutomorphicArithmetic.\allowbreak{}AutomorphicVarieties.\allowbreak{}UnitaryThreeGroups.\allowbreak{}CorGSumLog}{37--74}
\leanref{UpperTheory.\allowbreak{}V.\allowbreak{}U3.\allowbreak{}countI}{Development.\allowbreak{}AutomorphicArithmetic.\allowbreak{}AutomorphicVarieties.\allowbreak{}UnitaryThreeGroups.\allowbreak{}CountI}{22--137}
\end{proof}

\subsection*{Proof of \ref{B6.43b}}\label{prf:B6.43b}\phantomsection
\begin{proof}
First identify \(f\). The Atkin--Lehner--Li decomposition at weight \(h\) and level \(\Gamma_1(M)\) expresses the cusp space as the direct sum of the images of the new spaces of levels \(M'\mid M\) under \(g(z)\mapsto g(dz)\), \(d\mid M/M'\). Each new space is a sum of distinct eigenlines for the Hecke algebra away from \(M\). The degeneracy maps commute with those Hecke operators. Thus each nonzero component of \(f\) comes from a newform with eigenvalues \(a_p\) for \(p\nmid M\). The lower-level exclusion gives \(M'=M\), and then \(d=1\). Multiplicity one gives a single eigenline at level \(M\). The normalization \(a_1=1\) identifies \(f\) with its normalized primitive eigenform. This is the newform of \cite[\S6.1]{kato-2004}. \cite[Theorem~9.4]{stein-2007}

Put \(F=\Q(a_n:n\ge1)\). This is a number field containing the values of \(\omega\). The restriction of \(\iota^{-1}\) to \(F\) determines a place \(\lambda\mid\ell\) and an embedding \(E=F_\lambda\hookrightarrow\overline{\Q}_\ell\). Its image is \(\Q_\ell(\iota^{-1}a_n:n\ge1)\), since the coefficient field is generated by finitely many of the \(a_n\). Let \(U=V_E(f)\) be Kato's two-dimensional realization. For \(p\nmid M\ell\), his formula on inertia invariants is a formula on all of \(U\), and gives
\[
 \det(1-u\,U(\sigma_p^{-1}))
   =1-\iota^{-1}(a_p)u+\iota^{-1}(\omega(p))p^{2e}u^2 .
\]
Thus \(U\) and \(W(-e)\) have the same characteristic polynomial on inverse arithmetic Frobenius outside a finite set. Chebotarev density in each finite Galois quotient and continuity give equality of traces on \(G_\Q\). Brauer--Nesbitt identifies the semisimplification of \(U\otimes_E\overline{\Q}_\ell\) with \(W(-e)\). The latter is irreducible of dimension two, so the former representation already has a single composition factor of dimension two. Twisting gives
\[
                 W\simeq V\otimes_E\overline{\Q}_\ell,
                 \qquad V=V_E(f)(e).
\]
In particular the twist in Kato's notation has sign \(+e\): at \(\sigma_p^{-1}\), a twist by \(e\) multiplies each root by \(p^{-e}\). \cite[\S6.3 and (14.10.3)]{kato-2004} \cite[Ch.~VII, Theorem~13.4]{neukirch-1999}
\leanref{UpperTheory.\allowbreak{}V.\allowbreak{}Af.\allowbreak{}Newform.\allowbreak{}heckePoly}{Development.\allowbreak{}AutomorphicArithmetic.\allowbreak{}AutomorphicVarieties.\allowbreak{}AutomorphicForms.\allowbreak{}Newform}{41--42}
\leanref{UpperTheory.\allowbreak{}V.\allowbreak{}Lit.\allowbreak{}coefF}{Development.\allowbreak{}AutomorphicArithmetic.\allowbreak{}AutomorphicVarieties.\allowbreak{}LiteratureResults.\allowbreak{}CoefFieldEq}{48--50}
\leanref{UpperTheory.\allowbreak{}V.\allowbreak{}Lit.\allowbreak{}FrobPoly}{Development.\allowbreak{}AutomorphicArithmetic.\allowbreak{}AutomorphicVarieties.\allowbreak{}LiteratureResults.\allowbreak{}CoefFieldEq}{51--57}

Choose a \(G_\Q\)-stable \(\mathcal O_E\)-lattice \(T\subset V\). Such a lattice exists because the compact Galois image has bounded translates of any lattice, and their sum is a stable lattice. Kato's Theorem~14.2(1) gives finiteness over every finite abelian extension for each stable lattice in \(V_E(f)(r)\) when \(1\le r\le h-1\) and \(2r\ne h\). Here \(K/\Q\) has degree two, \(r=e\), \(1\le e\le 2e=h-1\), and \(2e=h-1\ne h\). This includes \(e=1\). Hence \(\operatorname{Sel}(K,T)\) is finite. Kato's theorem includes newforms with complex multiplication. \cite[Theorem~14.2(1), p.~235, and \S15]{kato-2004}

We pass from this finite group to rational coefficients. Put \(A=V/T\). By \cite[\S14.1]{kato-2004}, \(\operatorname{Sel}(K,T)\) is the subgroup of \(H^1(G_K,A)\) defined at each place \(v\) by the image of \(H_f^1(K_v,V)\) in \(H^1(G_{K_v},A)\). Therefore reduction modulo \(T\) maps \(H_f^1(K,V)\) into \(\operatorname{Sel}(K,T)\). Its image is divisible, since \(H_f^1(K,V)\) is an \(E\)-vector space. A finite group has zero divisible subgroup.

Let \(z:G_K\to V\) represent a finite class. For each \(n\ge0\) the class of \(\ell^{-n}z\) reduces to zero in \(H^1(G_K,A)\). Lift a vector giving its coboundary from \(A\) to \(V\), and multiply that lift by \(\ell^n\). The resulting vector \(x_n\in V\) satisfies
\[
             z(g)-(gx_n-x_n)\in\ell^nT
                       \qquad\text{for every }g\in G_K .
\]
The coboundaries form a finite-dimensional \(E\)-subspace of the space of continuous functions \(G_K\to V\) with its supremum norm. This subspace is closed because \(E\) is complete. The displayed uniform convergence puts \(z\) in that subspace. Thus \(H_f^1(K,V)=0\).

The same vanishing holds with coefficients in \(\overline{\Q}_\ell\). To see this for an arbitrary continuous cocycle \(z\) with these coefficients, realize it in the representation
\[
          g\longmapsto
          \begin{pmatrix} V(g)&z(g)\\0&1\end{pmatrix}
                        \quad(g\in G_K).
\]
Its image is contained in the matrices over a finite extension \(E'/E\). Indeed, the finite subextensions of \(\overline{\Q}_\ell/E\) are countable: by Krasner's lemma each is generated by a root of a polynomial over a fixed countable dense subfield of \(E\). The preimages of the matrix groups over these finite extensions are closed subgroups covering \(G_K\). Baire's theorem gives one with interior, hence of finite index. Adjoining the entries of finitely many coset representatives gives \(E'\).

The local conditions descend to \(E'\). Away from \(\ell\), apply an \(E'\)-linear retraction \(\overline{\Q}_\ell\to E'\) to a vector giving the inertia coboundary. At \(v\mid\ell\), apply the same retraction on coefficients, tensored over \(\Q_\ell\) with the identity of \(B_{\mathrm{cris}}\), to the vector giving the crystalline coboundary. Both operations commute with Galois because the matrices of \(V\) are over \(E\). In an \(E\)-basis of \(E'\), the coordinates of \(z\) are now continuous \(V\)-valued cocycles satisfying the finite conditions. Each is a coboundary by the preceding paragraph. Recombining their cobounding vectors proves \(H_f^1(K,W)=0\).

Finally take \(D=W\) or \(D=W\otimes\theta\). Restriction takes \(H_f^1(\Q,D)\) into \(H_f^1(K,W)\), since \(\theta|_{G_K}=1\). At \(w\mid p\ne\ell\) an inertia coboundary restricts to an inertia coboundary on \(I_w\), including when \(p\) ramifies in \(K\). At \(w\mid\ell\) a vector giving the \(B_{\mathrm{cris}}\)-coboundary still gives it on \(G_{K_w}\). The archimedean \(H^1\) with rational coefficients is zero. The two restrictions correspond respectively to the \(+1\) and \(-1\) actions of \(\Gal(K/\Q)\) on classes over \(K\), since twisting multiplies the action of its nonidentity element by \(\theta=-1\). For each choice of \(D\), corestriction after restriction on \(H^1(G_\Q,D)\) is multiplication by \([K:\Q]=2\). Both restrictions have zero image and are injective. This proves the two assertions.
\end{proof}

\subsection*{Proof of \ref{B6.43a}}\label{prf:B6.43a}\phantomsection
\begin{proof}
Put \(\phi=\delta^{-1}\chi^{-1}\) and take the primitive newform \(g\) attached to \(V\phi\varepsilon_\ell^{-k}\). The flag of \(V\phi\) meets the spherical test. Strong multiplicity one identifies the norm-twisted adelisation of \(g\) with the tested attachment \cite{jacquet-shalika-1981}. If \(\ell\mid M_g\), replacing the vector at \(\ell\) by a spherical vector and retaining the other local vectors and the lowest-weight vector at infinity gives the same good Hecke eigenvalues at level \(M_g/\ell\), contrary to primitivity.

Let \(\kappa_0,\kappa_1\) be the flag characters. For a Frobenius lift \(\sigma\) at \(A\), put \(u=(\kappa_0\phi)(\sigma)\varepsilon_\ell(\sigma)^{-k}\). The spherical comparison gives \(b\) in the Satake multiset \(B\) with \((\ell^k b)^{-1}=\iota(u)\). The normalised spherical polynomial of \(g\) has roots \(\ell^kb'\) for \(b'\in B\). Compactness gives \(|u|_\ell=1\), and the determinant formula gives valuation \(2k\) for the other root. Hence \(a_\ell(g)\) is a unit. If the unit root were a root of unity, the other root would have complex modulus \(\ell^{2k}\). The Hecke recurrence at \(\ell^j\) would contradict \eqref{eq:p6c5-coeff}.

The inertia exponent excludes \(\kappa_1=\varepsilon_\ell\). Equality \(\kappa_0=\varepsilon_\ell\) would force \(k=1\) and \(u=\phi(\sigma)\), a root of unity. In a stable lattice over finite coefficients, extend a basis of the saturated intersection with the first line. These matrices satisfy \eqref{eq:p6c5-normal} with \(h=\ell-1\). The source identity holds with \(\xi=\chi^{-1}\). Take \(S\) to contain \(\ell\) and the ramified primes.
\leanref{GlobalTheory.\allowbreak{}Ord2.\allowbreak{}sources\_\allowbreak{}of\_\allowbreak{}spherical}{Development.\allowbreak{}AutomorphicArithmetic.\allowbreak{}AutomorphicVarieties.\allowbreak{}ModularityArgument.\allowbreak{}ThresholdArguments.\allowbreak{}DenseThresholdCases.\allowbreak{}TwistNewform}{17--112}
\leanref{GlobalTheory.\allowbreak{}NewLevel.\allowbreak{}not\_\allowbreak{}dvd\_\allowbreak{}level\_\allowbreak{}of\_\allowbreak{}satake\_\allowbreak{}hecke\_\allowbreak{}poly\_\allowbreak{}cofinite}{Development.\allowbreak{}AutomorphicArithmetic.\allowbreak{}AutomorphicVarieties.\allowbreak{}ModularityArgument.\allowbreak{}ThresholdArguments.\allowbreak{}DenseThresholdCases.\allowbreak{}UnramifiedAlgebraicWeights}{59--133}
\leanref{UpperTheory.\allowbreak{}V.\allowbreak{}Z.\allowbreak{}Kato.\allowbreak{}WeilSphericalRecords.\allowbreak{}frobRoot\_\allowbreak{}alt}{Development.\allowbreak{}AutomorphicArithmetic.\allowbreak{}AutomorphicVarieties.\allowbreak{}ModularityArgument.\allowbreak{}ThresholdArguments.\allowbreak{}DenseThresholdCases.\allowbreak{}FrobRoot}{13--26}
\end{proof}

\subsection*{Proof of \ref{B6.46d}}\label{prf:B6.46d}\phantomsection
\begin{proof}
Cuspidality bounds \(y^{k+1/2}|g'(x+iy)|\). Integrate its square at \(y=X^{-1}\) and apply Parseval for \eqref{eq:p6c5-coeff}, also before depletion. Put \(M=N\operatorname{mod}(\xi)\) and \(b(n)=\xi(n)a_n n^{-k}\) on integers prime to \(M\), zero elsewhere. Partial summation gives
\[
 \sum_{n\le X}|b(n)|^2\ll X,\qquad
 S(1+s)=\sum_n|b(n)|^2 n^{-1-s}\ll s^{-1}\quad(0<s\le1).
 \tag{6.46.17}\label{eq:p6c5-square}
\]
The series \(L_b(z)=L(g'\otimes\xi,z+k)\), with the indicated omission, is holomorphic near the positive real axis by the Mellin integral of the twisted cusp form. It converges absolutely for real \(z>1\) by Cauchy--Schwarz.

At a good prime let \(\alpha,\beta\) be the normalised twisted roots, with \(|\alpha\beta|=1\). The prime-power bound gives \(|\alpha|,|\beta|\le p^{1/2}\) by the generating function \((1-\alpha x)^{-1}(1-\beta x)^{-1}\). Multiplying geometric series gives
\[
 \sum_{r\ge0}|b(p^r)|^2x^r
 =\frac{1-x^2}{\prod_{u,v\in\{\alpha,\beta\}}(1-u\bar v x)} .
\]
At \(x=p^{-z}\) the local logarithm of \(\zeta(z)|L_b(z)|^2S(z)\zeta(2z)\) is \(\sum_{j\ge1}|1+\alpha^j+\beta^j|^2x^j/j\). The zeta factors at omitted primes also have nonnegative logarithms. Absolute convergence gives nonvanishing for \(z>1\) and, for each finite set \(T\) of good primes,
\[
 \exp\!\left(\sum_{p\in T}|1+b(p)|^2p^{-1-s}\right)
 \le 2(1+s^{-1})|L_b(1+s)|^2 S(1+s).
 \tag{6.46.18}\label{eq:p6c5-positivity}
\]
The Petersson adjoint is \(T_p^*=\overline{\omega(p)}T_p\), so \(a_p=\omega(p)\overline{a_p}\). For \(p\equiv-1\bmod M\) we have \(\omega(p)=-1\) by odd weight and \(\xi(p)=\pm1\), so \(|1+b(p)|^2\ge1\). If \(L_b(1)=0\), holomorphy and \eqref{eq:p6c5-square} bound the right side of \eqref{eq:p6c5-positivity} as \(s\downarrow0\). Take the limit on finite sets of primes in this progression and use the divergence of their reciprocal sum \cite{dirichlet-1837}.
\end{proof}

\subsection*{Proof of \ref{B6.48j}}\label{prf:B6.48j}\phantomsection
\begin{proof}
Use the local domain \(B=B_\ell\) of Definition~\ref{B6.31c}. The values \(d_i=\delta_i(s_0)\) are powers of \(\varepsilon_\ell(s_0)\) with distinct exponents \(k+1,k,1,0\). After inverting their differences the projectors onto the four eigenlines of \(b_\ell(s_0)\) are \(P_i/\prod_{j\ne i}(d_i-d_j)\). The local equations say that their successive partial sums define a stable flag and that inertia acts on its \(i\)-th quotient by \(\delta_i\). The differences are nonzero in \(O\), so both characteristic-zero specializations preserve this calculation.

For \(g\in D_\ell\) let \(z_i\in\operatorname{Frac}B\) be the trace of the \(i\)-th projector times \(b_\ell(g)\). It is an eigenvalue of \(b_\ell(g)\), hence integral over \(B\). The finite domain \(B[z_i]\) is local, since a finite algebra over a Henselian local ring is a product of local rings. Its two evaluations are defined by the projector formula and send \(z_i\) to \(\delta_i(g)\) and \(\lambda_i(g)\). The first evaluation has image \(O\), by integrality and integral closedness of \(O\). Thus \(z_i-\delta_i(g)\) belongs to the maximal ideal. Under the second evaluation every element \(a\) of this ideal has image of absolute value less than one. Indeed, if \(|a'|\ge1\) for its image \(a'\), the algebraic number \(1/a'\) is integral over \(\Z_\ell\). Its monic equation, multiplied by a power of \(a'\), would send an element of \(1+aB[z_i]\) to zero. This proves the reduction assertion.

A positive power of \(s_0\) belongs to \(U\). Equality of \(\lambda_i\) and \(\varepsilon_\ell\lambda_j\) there would force their exponents to differ by one. Indices in either one of the two coordinate planes have exponent difference \(k\ge2\). For indices in different planes the residual ratio at \(\sigma_\ell\), or its inverse, is \(\overline{\xi\Psi^b}(\sigma_\ell)\) for some integer \(b\), since \(\varepsilon_\ell(\sigma_\ell)=\omega(\sigma_\ell)=1\). Definition~\ref{B6.31b} excludes equality at this element of \(U\).
\end{proof}

\subsection*{Proof of \ref{B6.48h}}\label{prf:B6.48h}\phantomsection
\begin{proof}
Each \(A\)-line \(A(\kappa_i)\), viewed over \(E\), is an unramified representation twisted by \(\varepsilon_\ell^{w_i}\), hence is crystalline. Induct on the flag. Let \(P\) be the preceding crystalline subrepresentation and let \(L=A(\kappa_j)\) be the next quotient. The extension class belongs to \(H^1(K,V)\) for \(V=\operatorname{Hom}_E(L,P)\). This representation is crystalline. Filtering the \(A\)-lines by powers of the maximal ideal gives character factors \(\bar\kappa_i/\bar\kappa_j\) in \(V\), with \(i<j\). The corresponding factors of \(V^\vee(1)\) have no invariants by hypothesis, so \(H^0(K,V^\vee(1))=0\). The cyclotomic exponents of \(V\) are \(w_i-w_j>0\). Consequently \(\mathrm{Fil}^0D_{\rm dR}(V)=0\), with the convention that \(\Q_\ell(1)\) has filtration jump \(-1\).

The crystalline dimension formula \cite[Proposition~3.8]{bloch-kato-1990}, local Tate duality and the local Euler characteristic formula \cite[Ch.~II, \S5]{serre-2002} now give
\begin{equation}\label{eq:p6c5-local-finite}\tag{6.48.3}
 \dim_E H^1_f(K,V)
 =\dim_E H^0(K,V)+[K:\Q_\ell]\dim_E V
 =\dim_E H^1(K,V).
\end{equation}
Thus the extension class dies over \(B_{\rm cris}\). A splitting there identifies the middle term after extension of scalars with the sum of its two crystalline ends. Taking period invariants gives its full dimension, so the middle term is crystalline.
\end{proof}

\subsection*{Proof of \ref{B6.48k}}\label{prf:B6.48k}\phantomsection
\begin{proof}
Write \(\omega=\bar\varepsilon_\ell\). We record the residual test used in choosing the tower. Its initial CM field has an element \(t\) with \(\omega(t)=1\) at which the four residual diagonal characters have distinct values. In II the values, up to the choice of order, are
\[
 -b,\quad -1,\quad b,\quad 1,\qquad b=\bar\xi(t),\quad b^2\ne1.
\]
Here the element with \(\bar\Psi(t)=-1\) can be multiplied, if necessary, by the element with \(\bar\Psi=1,\ \omega=1,\ \bar\xi^4\ne1\) in the same CM subgroup. In III there are two choices of initial field. The rigidity comparison uses the field cut out, before the auxiliary quadratic conditions, by \(\bar\xi^2=\bar\Psi^{\epsilon}\), where \(\epsilon\in\{1,-1\}\) is the chosen label order. Dividing the ordered characters by the first gives \(1,b^{-1},b^{-2},b^{-3}\) at its marking, with \(b=\bar\xi(t)\). Here \(\omega(t)=1\) and \(b\) is the negative of a nontrivial element of odd order, so its first four positive powers differ from one. The tensor seed in \ref{B6.49} uses a separately chosen field on which \(\xi=\lambda_{K'}\). At its marking \(\omega(t)=1\), \(\xi(t)=-1\), and the four characters divided by the first have values \(1,-1,z^{-1},-z^{-1}\), with \(z=\bar\Psi(t)^\epsilon\) of odd order greater than one. These values are distinct since \(\ell\) is odd. In this construction \(-z\) has order twice the order of \(z\), hence exceeds four. In either construction use the marking belonging to the chosen field. In the tower enlargement a power \(t^{\ell^a}\) belongs to \(G_L\). Raising to this power preserves distinctness in the residue field. In particular the four residual characters restrict distinctly to \(G_L\). Also, a sign character of \(G_L\) factoring through \(\omega\) and permuting the residual characters by multiplication must be trivial. Indeed it takes value one at the marking, so its permutation fixes every index. A common residual twist cancels from this test.
\leanref{UpperTheory.\allowbreak{}V.\allowbreak{}Ws.\allowbreak{}elts2}{Development.\allowbreak{}AutomorphicArithmetic.\allowbreak{}AutomorphicVarieties.\allowbreak{}WeakShimuraTheory.\allowbreak{}FrobeniusFields.\allowbreak{}Elts2}{19--52}
\leanref{UpperTheory.\allowbreak{}V.\allowbreak{}Ws.\allowbreak{}ok2A}{Development.\allowbreak{}AutomorphicArithmetic.\allowbreak{}AutomorphicVarieties.\allowbreak{}WeakShimuraTheory.\allowbreak{}FrobeniusFields.\allowbreak{}OddPowNe}{55--177}
\leanref{UpperTheory.\allowbreak{}V.\allowbreak{}Ws.\allowbreak{}exists\_\allowbreak{}tame\_\allowbreak{}tower\_\allowbreak{}ok\_\allowbreak{}three}{Development.\allowbreak{}AutomorphicArithmetic.\allowbreak{}AutomorphicVarieties.\allowbreak{}WeakShimuraTheory.\allowbreak{}CompactIntegralFields.\allowbreak{}NegSqPowNeOneOfOddOrder}{181--280}

First \(r|_{G_L}\) is irreducible. Otherwise its semisimple restriction has a nonscalar endomorphism. The finite abelian group \(\Gal(L/\Q)\) acts by conjugation on the endomorphisms commuting with \(G_L\), with scalar fixed space by Schur's lemma. A nontrivial eigenspace therefore supplies an intertwiner with a nontrivial character twist of \(r\). This intertwiner is invertible by irreducibility over \(G_\Q\). Choose \(g\) where the twisting character differs from one. The trace of \(r(gh)\) vanishes for all \(h\in G_L\). Reducing gives
\[
 \sum_i\bar\delta_i(g)\bar\delta_i(h)=0\qquad(h\in G_L),
\]
contrary to independence of the distinct characters \(\bar\delta_i|_{G_L}\). This also proves irreducibility over \(K\).

Put \(H=G_{L(\mu_{\ell^\infty})}\). There is an element of \(H\) with simple spectrum. Here is the local argument used to obtain it. At a place of \(K\) over \(\ell\), the tower splits \(K/\Q\) locally. Lemma~\ref{B6.48j} separates the ordered characters from their cyclotomic translates for \(i<j\). For the reverse pairs the distinct inertia weights exclude equality, and the cyclotomic character is nontrivial for the equal-index pairs. Thus the witness of Definition~\ref{def:b6-cuspidal-witness} applies at this place. Choose a Weil element of positive residue degree \(d\) with cyclotomic value one. Such an element is obtained by raising a Frobenius lift to a power and multiplying by inertia, since the cyclotomic image of inertia is open in the local cyclotomic image. Write \(q\) for the residue cardinality. In the ordinary local comparison the moduli of the smooth characters at this element are the moduli of the ordinary characters with their respective algebraic weight factors \(q^{\pm d m_i}\), with the same convention for all \(i\). Their exponents at a uniformiser lie in one open interval of length one by the witness. A root-of-unity ratio of two ordinary eigenvalues would therefore give \(|m_i-m_j|<1\). The weights \(k+1,k,1,0\) are distinct integers. A positive power of the element belongs to \(G_L\) and still separates all four eigenvalues.
\leanref{UpperTheory.\allowbreak{}V.\allowbreak{}Z.\allowbreak{}Win.\allowbreak{}RR.\allowbreak{}free\_\allowbreak{}core\_\allowbreak{}alt}{Development.\allowbreak{}AutomorphicArithmetic.\allowbreak{}AutomorphicVarieties.\allowbreak{}ModularityArgument.\allowbreak{}ThresholdArguments.\allowbreak{}DenseThresholdCases.\allowbreak{}RRFree}{97--116}

Suppose the restriction to \(H\) is reducible. As in Lemma~\ref{B6.30aa}, it is semisimple with transitively permuted isotypic summands. Their multiplicity is one, since repeated copies give repeated eigenvalues at the separating element. The number of summands is therefore two or four. The action on the summands factors through the abelian cyclotomic quotient of \(G_L\). The stabiliser of one summand fixes them all, and the resulting transitive abelian permutation group acts regularly. It has a nontrivial sign character \(\eta\). The stabiliser is closed, being the stabiliser of a subspace, and has finite index, so \(\eta\) is continuous. Multiplication by the sign of the coset on each summand intertwines \(r|_{G_L}\) with its twist by \(\eta\). Since \(\ell\) is odd, the pro-\(\ell\) subgroup of the cyclotomic image lies in \(\ker\eta\). Thus \(\eta\) factors through \(\omega\), and reduction preserves its two possible values. Reduction of the isomorphism of twists shows that \(\bar\eta\) permutes the four residual characters. The marking test gives \(\bar\eta=1\), a contradiction. Finally, a scalar twist preserves all invariant subspaces.
\end{proof}

\subsection*{Proof of \ref{B6.48}}\label{prf:B6.48}\phantomsection
\begin{proof}
Write \(r'=(1+\epsilon z)r\), with \(\epsilon^2=0\). The fixed similitude puts \(z\) in \(\mathfrak{sp}_4\), and multiplication gives
\[
 z(gh)=z(g)+\operatorname{Ad}(r(g))z(h).
 \tag{6.48.2}
\]
Away from \(\ell\), average over the finite inertia quotient through which the lift factors to obtain an inertia coboundary.

Work over a finite coefficient field \(E\) at \(\ell\). The marked inertia eigenvalue differences are units over the dual numbers, so the lifted ordered flag reduces to the prescribed flag. Write \(\chi=(\vartheta,\omega)\) on the local group, let \(\sigma\) be the specified Frobenius, and let \(f\) be the unramified character with values in \(\operatorname{im}\chi\) and \(f(\sigma)=\chi(\sigma)\). On \(U=\ker(\chi f^{-1})\) the inertia kills both finite characters. The graded characters there have cyclotomic weights \(k+1,k,1,0\). Since \(\sigma\in U\), Lemma~\ref{B6.48j} gives \(\kappa_i\ne\varepsilon_\ell\kappa_j\) on \(U\) for \(i<j\), where \(\kappa_i\) are the graded characters of \(r\).

Apply Lemma~\ref{B6.48h} over the fixed field of \(U\), with \(A=E[\epsilon]/(\epsilon^2)\). Its inertia hypothesis holds for the lifted characters, and the reduction of the ordered lifted flag is the flag just considered. The resulting crystalline self-extension splits over \(B_{\rm cris}\) on \(U\). Average a splitting over the finite quotient \(D_\ell/U\) to obtain a \(D_\ell\)-equivariant splitting. The symplectic projection \(Y\mapsto(Y-J^{-1}Y^tJ)/2\), for the matrix \(J\) of the alternating form, is equivariant under conjugation by a similitude and fixes the values of \(z\). Hence \(z\) satisfies the finite local condition.
\leanref{tb\_\allowbreak{}atL}{Development.\allowbreak{}AutomorphicArithmetic.\allowbreak{}AutomorphicVarieties.\allowbreak{}SymplecticGSp4.\allowbreak{}TbTriDiag}{287--374}
\leanref{UpperTheory.\allowbreak{}V.\allowbreak{}Loc.\allowbreak{}pr6}{Development.\allowbreak{}AutomorphicArithmetic.\allowbreak{}AutomorphicVarieties.\allowbreak{}PadicHodgeTheory.\allowbreak{}CrystallineTransferSurjection}{70--98}

Take the fitted tower for the ordinary model, with quadratic field \(K\) and CM compositum \(L\). The residual marking argument in Lemma~\ref{B6.48k} gives irreducibility on \(G_K\), so the attachment hypothesis applies there. That lemma then gives irreducibility of the twisted representation on \(G_{L(\mu_{\ell^\infty})}\). Its image in the level group is the full subgroup \(\Gamma'\) of \ref{B6.53}, because the quotient respects the cyclotomic character. For the symplectic form matrix \(J\), put \(X^\dagger=J^{-1}X^tJ\) and \(P(X)=(X-X^\dagger)/2\). The projection \(P\) commutes with \(\operatorname{Ad}(r(g))\). For the sign \(\eta\) of \(K\) use the polarised adjoint action
\[
 R(g)X=\operatorname{Ad}(r(g))P(X)
 +\eta(g)\operatorname{Ad}(r(g))(X-P(X)).
 \tag{6.48.4}
\]
On \(G_K\) this is the conjugation action. The scalar twist over \(L\) preserves the adjoint. Replacing a complex conjugation \(c_K\) by \(wc_K\), for \(w\in G_K\), replaces its polarising matrix by its right product with \(r(w)^{-1}\). On excluded inertia the adjoint action is trivial, so the local coboundary vanishes there and \(z\) descends to the frame over \(L^+\). Solvable base change transfers the attachment and witness. By \ref{B6.53}, the affine action \(Y\mapsto R(g)Y+z(g)\) has a \(G_{L^+}\)-fixed point. Its \(G_\Q\)-orbit is finite. Project the barycentre to \(\mathfrak{sp}_4\) to obtain \(X\) with \(z(g)=X-\operatorname{Ad}(r(g))X\).
\leanref{UpperTheory.\allowbreak{}V.\allowbreak{}Z.\allowbreak{}Win.\allowbreak{}Dn.\allowbreak{}exists\_\allowbreak{}cm\_\allowbreak{}field\_\allowbreak{}twist\_\allowbreak{}h1}{Development.\allowbreak{}AutomorphicArithmetic.\allowbreak{}AutomorphicVarieties.\allowbreak{}ModularityArgument.\allowbreak{}ThresholdArguments.\allowbreak{}DenseThresholdCases.\allowbreak{}ExistsCmFieldTwistH1}{28--103}
\leanref{UpperTheory.\allowbreak{}V.\allowbreak{}Th.\allowbreak{}adDown}{Development.\allowbreak{}AutomorphicArithmetic.\allowbreak{}AutomorphicVarieties.\allowbreak{}TheoremsMuLinked.\allowbreak{}GaloisRepresentationTypes.\allowbreak{}AdDown}{92--194}
\end{proof}

\subsection*{Proof of \ref{B6.48i}}\label{prf:B6.48i}\phantomsection
\begin{proof}
Set \(\Phi(h\otimes m)=h(u)i(m)\). The elementary symmetric relations in the statement make this well defined. First work at one place and omit \(j\). In the polynomial ring on the ordered roots set
\[
 \Delta=\prod_{a<b}(x_b-x_a),\qquad H=\prod_{a<b}(x_b-Qx_a),
 \qquad \partial h=\Delta^{-1}\sum_{\sigma\in S_n}\operatorname{sgn}(\sigma)\sigma h .
\]
The alternating numerator is divisible by each \(x_b-x_a\), so \(\partial h\) is a symmetric polynomial. These integral identities may first be calculated with indeterminate \(Q\).

With the Iwahori subgroup of volume one, \(iT=\sum_{\sigma}T_\sigma\): the finite flag variety partitions into the Bruhat cells indexed by \(\sigma\). Write \(E_0\) for this sum. In the polynomial representation the Bernstein relation for the adjacent transposition \(s=(a,a+1)\) gives
\[
 T_sh=Q\,s(h)+(Q-1)x_a\frac{h-s(h)}{x_a-x_{a+1}} .
\]
Here the split group has root system of type \(A_{n-1}\), all simple parameters equal \(Q\), and the denominator divides the indicated difference. This is the Bernstein presentation in \cite[\S1.15]{haines-kottwitz-prasad-2010}. Since \(T_sE_0=QE_0\), the output of \(E_0\) is symmetric. In its expansion \(\sum_\sigma c_\sigma\sigma\) over the rational function field this gives \(c_{\tau\sigma}=\tau(c_\sigma)\), by linear independence of the distinct field automorphisms. The longest permutation occurs only in its own reduced word. Its coefficient is the product over inversions of \((x_a-Qx_b)/(x_a-x_b)\), the longest permutation applied to \(H/\Delta\). Hence \(E_0 h=\partial(Hh)\). The calculation gives identities in the Hecke algebra: the expansions of the \(T_\sigma\) in permutations are triangular in Bruhat order with nonzero leading coefficients, and the Bernstein basis is free over the root-polynomial algebra. Polynomial identities so obtained specialize at the given \(Q\).

Let \(A\) run over subsets of the pairs \(a<b\). Put \(a_A=\prod_{(a,b)\in A}x_b\), \(c_A=\prod_{(a,b)\notin A}(-x_a)\). Then
\[
 \sum_A a_A\,\sigma(c_A)=\prod_{a<b}(x_b-x_{\sigma(a)})
 =\begin{cases}\Delta,&\sigma=1,\\0,&\sigma\ne1.\end{cases}
\]
Indeed a nonidentity permutation has some \(\sigma(a)>a\). Define
\(\Psi_0(v)=\sum_A a_A(x)\otimes T(c_A(u)v)\).
This map intertwines the roots. To see this before specialization, take two lists \(x,y\) with the same elementary symmetric functions. The tensor \(\sum_A a_A(x)c_A(y)\) is killed by \(x_i-y_i\): the identity just proved shows this after substituting every permutation of \(x\) for \(y\). These substitutions detect identities. The ordered splitting algebra in \(y\) is free over the polynomial ring in \(x\), by successive monic root adjunctions, and over its fraction field it splits as the product over the permutations. Symmetric functions can be moved through \(T\) since they are central. Specializing the tensor identity proves the required intertwining, including for coincident specialized roots.

On \(V\) the composite \(\Phi\Psi_0\) is \(\sum_A a_A E_0 c_A\). Substituting \(E_0=\sum_\sigma\sigma(H/\Delta)\sigma\) and using the subset identity gives multiplication by \(H\). On \(h\otimes m\) the other composite gives \(\sum_A a_A\partial(Hc_Ah)\otimes m=Hh\otimes m\), by the same identity. Here \(T(p(u)i(m))=\partial(Hp)m\), since an Iwahori simple generator acts by \(Q\) on a spherical vector and the symmetric coefficients act as specified.

For several places take the product of the forward operations and the sum over tuples of subsets in the reverse operation. The actions at distinct places commute, so both composites are multiplication by \(\prod_jH_j\). Multiply the reverse map by the polynomial \(f/\prod_jH_j\), which exists since every factor \(x_{jb}-Q_jx_{ja}\), \(a<b\), occurs in \(f\). All operations commute with coefficients through their action on the common representation. Intertwining gives the same identities after applying the indicated idempotents.
\end{proof}

\subsection*{Proof of \ref{B6.48ca}}\label{prf:B6.48ca}\phantomsection
\begin{proof}
Enlarge the coefficient field to contain the character values and the images of the embeddings. The rank-one local reciprocity comparison gives \(\kappa_i\) filtration jump \(h_{i,\tau}\) at \(\tau\) and smooth Weil character \(\omega_i\circ a\). Indeed, its restriction to an open subgroup of the units is the displayed algebraic character. The rank-one crystalline periods for the embeddings give its indicated jumps. Removing that algebraic factor leaves the displayed smooth character on both units and a uniformiser. This is the locally algebraic character comparison of \cite[\S B.3]{fontaine-2004}, with reciprocity as in \ref{B9.06}.

The inertia action on the Weil--Deligne space factors through a finite group. Over the coefficient field of characteristic zero it splits into its character spaces. The polynomial hypothesis gives each of the \(n\) distinct inertia characters multiplicity one. Each extends to \(W_K\), so its line is \(W_K\)-stable. On the \(i\)-th line the Weil character is \(\omega_i\circ a\). To see the identification at a Frobenius lift \(s\), take traces at \(su\) for inertia elements \(u\) and use linear independence of the inertia characters as functions of \(u\). The monodromy operator commutes with inertia. On each line it is a nilpotent scalar and hence zero.

Twist by \(\kappa_0^{-1}\). The line indexed by zero in the resulting period module has trivial Weil character. As a \((\varphi,N)\)-module with descent data it is the rank-one module of the trivial character: the inertia action is trivial, and the linear Frobenius for one residue-field period acts as one. Descent for the semilinear Frobenius gives a Frobenius-fixed generator. Its Newton number is zero. At every embedding all jumps of the ambient twisted filtration are \(h_{i,\tau}-h_{0,\tau}\ge0\), so the induced filtration on the line has nonnegative Hodge number. Weak admissibility bounds this number above by its Newton number. Thus every labelled jump of the line is zero.

Over a finite Galois extension where the representation is semistable, this filtered line is weakly admissible. The equivalence for filtered \((\varphi,N)\)-modules gives a subrepresentation, and full faithfulness transports the descent data on the inclusion. The line with its jump-zero filtration and trivial descent character gives the trivial subrepresentation of the twist. Here the base is a finite extension of \(\Q_\ell\), and both the ambient module and the line have equal Hodge and Newton numbers, as required in \cite[Thm.~A and \S4.1]{colmez-fontaine-2000}. Untwisting gives a subline of character \(\kappa_0\). Exactness on semistable representations, with the induced and quotient filtrations, removes its Weil character and its one labelled jump at each embedding. The quotient satisfies the same hypotheses for indices \(1,\ldots,n-1\). Induction gives the flag.
\end{proof}

\subsection*{Proof of \ref{B6.48cb}}\label{prf:B6.48cb}\phantomsection
\begin{proof}
Use principal level one at \(v_0\) and a free coefficient lattice \(O^b\) for the tame type \(\tau\). Summing translates of a lattice under the finite image of \(\tau\) makes it invariant, and principal level one at the odd prime \(p\ne\ell\) with \(p\nmid D_L\) makes stabilizers act trivially on coefficient fibres. At a place \(v\) over \(\ell\) the free pro-\(\ell\) part of the local units has rank \([F_v:\Q_\ell]\), giving the power series presentation \(A=\Lambda=O[[X_1,\ldots,X_{n[F:\Q]}]]\) with zero residual values for the variables \(X_m\).

Over any complete Noetherian local \(O\)-algebra \(A'\) with continuous character \(\theta'\), write \(Y'=Y(A',\theta')\). At fixed precision the diagonal character \(\theta'\) factors through a finite quotient, and sufficiently high powers of \(U_z\) contract the lower unipotent variation to the corresponding level. There are finitely many double cosets at that level. On the resulting finite module the factorial powers eventually project to the invertible part. The projectors commute with reduction and give a continuous ordinary projector \(e\).

Choose the double cosets at precision one. Let \(i:Y'\to (A')^t\) be evaluation on their representatives, with the coefficient coordinates included, and \(\sigma:(A')^t\to Y'\) the map obtained by extending values and applying \(e\). The stabilisers act trivially on coefficient fibres at the chosen level, so the extension is \(A'\)-linear. For \(f\in Y'\cap V_j\), evaluate at double cosets at precision \(j+1\), extend these values and apply \(e\). This gives \(f\) modulo \(V_{j+1}\) as a linear combination of projected extensions of coordinate vectors, with coefficients in \(\mathfrak m_{A'}^j\). Each projected extension is congruent modulo \(V_1\) to its approximation by \(\sigma i\). Thus \(f\) agrees modulo \(V_{j+1}\) with \(\sigma(a_j)\) for some \(a_j\in\mathfrak m_{A'}^j (A')^t\). Successive corrections sum in \((A')^t\). With \(j=0\) this proves surjectivity of \(\sigma\), and with general \(j\) it proves \(Y'\cap V_j=\mathfrak m_{A'}^jY'\). In particular \(\sigma i\) is the identity modulo \(\mathfrak m_{A'} Y'\). It is invertible on the finite module \(Y'\) by Nakayama's lemma and the fact that a surjective endomorphism of a Noetherian module is injective. The map \((\sigma i)^{-1}\sigma\) splits \(i\). Hence \(Y'\) is projective and is finite free over the local ring \(A'\).
\leanref{UpperTheory.\allowbreak{}V.\allowbreak{}Hi.\allowbreak{}Frame.\allowbreak{}Ord}{Development.\allowbreak{}AutomorphicArithmetic.\allowbreak{}AutomorphicVarieties.\allowbreak{}HeckeAlgebras.\allowbreak{}ContinuousTorusCharacter}{229--231}
\leanref{UpperTheory.\allowbreak{}V.\allowbreak{}Hi.\allowbreak{}ordFree}{Development.\allowbreak{}AutomorphicArithmetic.\allowbreak{}AutomorphicVarieties.\allowbreak{}HeckeAlgebras.\allowbreak{}OrdinaryHeckeCharacters.\allowbreak{}OrdFree}{441--482}
\leanref{UpperTheory.\allowbreak{}V.\allowbreak{}Hi.\allowbreak{}ordAbs}{Development.\allowbreak{}AutomorphicArithmetic.\allowbreak{}AutomorphicVarieties.\allowbreak{}HeckeAlgebras.\allowbreak{}OrdinaryHeckeCharacters.\allowbreak{}OrdFree}{185--205}
\leanref{UpperTheory.\allowbreak{}V.\allowbreak{}Hi.\allowbreak{}ordStruct}{Development.\allowbreak{}AutomorphicArithmetic.\allowbreak{}AutomorphicVarieties.\allowbreak{}HeckeAlgebras.\allowbreak{}OrdinaryHeckeCharacters.\allowbreak{}OrdFree}{357--375}

Specialising to \((A,\theta)\), the simultaneous primary decomposition for the commuting operators \(T_r\) on \(Y/\mathfrak m_A^a Y\) lifts by idempotents in the finite commutative Hecke algebra over \(A\). Hence \(M(K)\) is a free direct summand over \(A\).
\leanref{UpperTheory.\allowbreak{}V.\allowbreak{}Hi.\allowbreak{}HSet.\allowbreak{}MK}{Development.\allowbreak{}AutomorphicArithmetic.\allowbreak{}AutomorphicVarieties.\allowbreak{}HeckeAlgebras.\allowbreak{}CoefficientRepTrivialOnUnramified}{160--162}
\end{proof}

\subsection*{Proof of \ref{B6.48c}}\label{prf:B6.48c}\phantomsection
\begin{proof}
Let \(H\) be the CM subgroup of the level group and \(c\) its complex conjugation. The restricted-ramification global group and the local groups have finite continuous \(\F_\ell\)-character spaces on all open subgroups by class field theory. Use the continuous deformation ring of the split determinant \(\prod_i\bar\chi_i\) on \(H\) \cite{chenevier-2014}. Tensor it with the local character rings for the ordered characters \(\psi_{vi}\). In these rings fix the prescribed \(\ell\)-primary torsion values in the completed group algebra and use Teichmüller lifts on the prime-to-\(\ell\) part.

Write \(\lambda_i,\operatorname{tr}_D\) for the characteristic coefficients and trace of the universal determinant, and set \([g]^*=\varepsilon_\ell(g)^{1-n}[cg^{-1}c^{-1}]\) on group algebra basis elements. In the completed tensor product impose the relations
\[
 \lambda_i([g]^*)-\lambda_i([g]),\qquad
 \operatorname{tr}_D\!\left([h]\prod_{i=0}^{n-1}([g_i]-\psi_{vi}(g_i))\right),
 \qquad \operatorname{tr}_D([h]([u]-1)).
\]
Here \(g,h\in H\), the \(g_i\) are in the decomposition group of \(v\) in the given flag order, and \(u\) ranges over the prescribed tame inertia kernels. The assumed pair makes the defining ideal proper. The quotient represents ordinary determinant pairs by the universal properties of its factors and relations.

Take the away level \(K\), tame coefficient lattice \(O^b\), weight ring \(A=\Lambda\), compact diagonal character \(\theta\), ordinary module \(Y\), ordinary projector \(e\), precision-one evaluation basis, and finite free direct summand \(M(K)\subseteq Y\) of Lemma~\ref{B6.48cb}. The spherical characteristic polynomials are at good inverse Frobenii. The ordinary operators give the ordered products of the flag characters on uniformisers, with the factors from \(\theta\). Suppose first that \(M(K)\ne0\), and write \(T\subseteq\operatorname{End}_A M(K)\) for the algebra generated by these operators and the weights. Use algebraic weight zero. Among its arithmetic specialisations use those for which, at each \(v\mid\ell\), the finite diagonal characters on \(\mathcal O_{L_v}^\times\) are pairwise distinct. These weights separate \(A\). Indeed the free local unit groups have ranks \([L_v:\Q_\ell]>0\). On their product the weight coordinates vary independently, after a fixed character twist on each factor. All finite characters separate its completed group algebra by Fourier inversion on each finite quotient. On a free generator at each \(v\), multiply a proposed vanishing element by the product of the differences of the universal values of the finite parts for \(i<j\). Each difference is nonzero because it involves two independent coordinates. The product times the element vanishes at all finite-character specialisations: at a collision the multiplier vanishes, and elsewhere the weight is among those selected. The weight ring, after extending scalars for the fixed twists, is a power series domain. Cancellation proves the assertion.

Use the coefficient control isomorphism of \ref{B6.48d}, whose construction uses the ordinary projector and the evaluation basis of Lemma~\ref{B6.48cb}. At a selected weight decompose the classical space into its definite automorphic constituents. On each constituent the ordinary subspace injects into the upper-unipotent Jacquet module in each split factor at \(\ell\). Indeed a vector mapping to zero is a finite sum of unipotent differences. Averaging over a sufficiently large compact unipotent subgroup kills it. On upper-integral-unipotent invariants translation of this average is a nonzero index factor times a power of the strictly dominant ordinary operator, which is invertible on the ordinary part. The torus on the image has the prescribed ordered unit characters. Frobenius reciprocity gives principal-series support, and two characters with distinct unit restrictions cannot differ by an unramified norm twist. The induction is irreducible by \cite[Thm.~4.2]{zelevinsky-1980}. Its Jacquet module has the permutations of the inducing tuple as successive characters, one per permutation, by \cite[Lem.~2.12]{bernstein-zelevinsky-1977}. Precisely one permutation has the specified ordered unit restrictions. Hence the chamber operators act by scalars on this part, including on every copy if the constituent occurs with multiplicity. The good spherical operators also act by scalars there. The selected classical Hecke characters therefore separate \(T\): an operator annihilated by them acts as zero on each selected classical fibre, and every matrix entry in an \(A\)-basis then vanishes by separation of \(A\).

For such a Hecke character \(x\), apply \ref{B6.04a} to its constituent and \ref{B9.02} to the regular algebraic polarized cuspidal inputs of base change, with their prescribed twists. Write \(r_x\) for their semisimple sum. Its good inverse Frobenius polynomials are the spherical Hecke polynomials. At a selected split place above \(\ell\), the local support consists of singleton characters. Indeed its multiset is the support just computed, and a longer segment or a repeated cuspidal block in base change would repeat an inertia character. In the algebraic normalisation write these Weil characters as \(\omega_i\circ a\), with \(a\) the reciprocity map of Lemma~\ref{B6.48ca}. Their parameter has zero monodromy and is trivial on Weil commutators. Each cuspidal factor therefore satisfies the toral test of \ref{B9.05}--\ref{B9.07}. Direct sum and the base-change twists give the stated Weil polynomial for \(r_x\), and give its filtration jumps in the order \(h_{i,\tau}=i\) at weight zero.

Let \(U_{v,j}\) denote the unit operator for \(j\) initial uniformisers and put \(U_{v,0}=1\). On a reciprocity lift of a uniformiser prescribe the value \(U_{v,i+1}/U_{v,i}\) for \(\psi_{vi}\varepsilon_\ell^i\), and on units prescribe the inverse of the \(i\)-th coordinate of \(\theta\). Units and a uniformiser give a character of \(L_v^\times\). Modulo each power of the maximal ideal of \(T\) its image is finite, so local reciprocity and inverse limit give a continuous decomposition-group character. On the Jacquet module the unipotent index of the chamber with \(j\) initial uniformisers is \(q_v^{j(n-j)}\). Including the half-modulus and the algebraic normalisation, the chamber eigenvalue is \(\prod_{i<j}\omega_i(\varpi_v)q_v^{-i}\). Successive quotients give the character formula of Lemma~\ref{B6.48ca} on a uniformiser, and the compact diagonal action gives it on units. That lemma supplies one flag for \(r_x\) with ordered characters \(x\psi_{vi}\). For arbitrary \(g_i\in D_v\), the rightmost factor \(r_x(g_{n-1})-x\psi_{v,n-1}(g_{n-1})\) lowers the last flag space. Each following factor lowers the next space. The product is zero, including after multiplication on the left by any \(r_x(h)\).

Embed \(T\) in the product of the coefficient fields at the selected characters. Its image is closed since \(T\) is finite over the compact weight ring. The coefficients of the product representation's characteristic polynomials take values in this image on good inverse Frobenii. Continuity and Frobenius density give the same conclusion on every group element. Amitsur's formula for the determinant of a sum expresses its polynomial-law coefficients on finite group-algebra sums integrally in the characteristic coefficients of words, hence gives the determinant law over \(T\) itself \cite[Lem.~1.12(ii)]{chenevier-2014}. Its reduction is the prescribed law by the same Frobenius calculation. Polarisation gives the involution identities at each selected character. On a prescribed tame inertia kernel every lifted type character is one, so the inertia difference acts by zero there. The ordered products vanish by the common flag just constructed. Separation of \(T\) now gives all three families of relations in \(T\). For the zero module the annihilator is the unit ideal.
\leanref{UpperTheory.\allowbreak{}V.\allowbreak{}Hi.\allowbreak{}eq\_\allowbreak{}zero\_\allowbreak{}of\_\allowbreak{}vanish\_\allowbreak{}at\_\allowbreak{}good\_\allowbreak{}points\_\allowbreak{}dense\_\allowbreak{}weight}{Development.\allowbreak{}AutomorphicArithmetic.\allowbreak{}AutomorphicVarieties.\allowbreak{}ModularityArgument.\allowbreak{}ThresholdArguments.\allowbreak{}DenseThresholdCases.\allowbreak{}TCoreSep}{37--66}
\leanref{UpperTheory.\allowbreak{}V.\allowbreak{}Hi.\allowbreak{}coreDet\_\allowbreak{}alt}{Development.\allowbreak{}AutomorphicArithmetic.\allowbreak{}AutomorphicVarieties.\allowbreak{}ModularityArgument.\allowbreak{}ThresholdArguments.\allowbreak{}DenseThresholdCases.\allowbreak{}TCoreDet}{291--322}
\leanref{UpperTheory.\allowbreak{}V.\allowbreak{}Z.\allowbreak{}Ord.\allowbreak{}core\_\allowbreak{}hecke\_\allowbreak{}of\_\allowbreak{}automorphic\_\allowbreak{}hecke\_\allowbreak{}dense\_\allowbreak{}weight}{Development.\allowbreak{}AutomorphicArithmetic.\allowbreak{}AutomorphicVarieties.\allowbreak{}ModularityArgument.\allowbreak{}ThresholdArguments.\allowbreak{}DenseThresholdCases.\allowbreak{}OrdCoreH}{20--61}
\leanref{UpperTheory.\allowbreak{}V.\allowbreak{}Hi.\allowbreak{}heckeDeterminant\_\allowbreak{}satisfies\_\allowbreak{}deformation\_\allowbreak{}relations\_\allowbreak{}alt}{Development.\allowbreak{}AutomorphicArithmetic.\allowbreak{}AutomorphicVarieties.\allowbreak{}ModularityArgument.\allowbreak{}ThresholdArguments.\allowbreak{}DenseThresholdCases.\allowbreak{}HeckeDeterminantSatisfiesDeformationRelations}{15--35}
 The coefficients, ordered products and weights generate the Hecke algebra. The induced action of \(B\) has kernel \(J=\operatorname{Ann}_B M(K)\), with \(J=B\) if \(M(K)=0\).
\leanref{UpperTheory.\allowbreak{}V.\allowbreak{}Hi.\allowbreak{}HSet.\allowbreak{}MK}{Development.\allowbreak{}AutomorphicArithmetic.\allowbreak{}AutomorphicVarieties.\allowbreak{}HeckeAlgebras.\allowbreak{}CoefficientRepTrivialOnUnramified}{160--162}
\leanref{UpperTheory.\allowbreak{}V.\allowbreak{}Hi.\allowbreak{}nonempty\_\allowbreak{}unitaryHeckeData}{Development.\allowbreak{}AutomorphicArithmetic.\allowbreak{}AutomorphicVarieties.\allowbreak{}ModularityArgument.\allowbreak{}ThresholdArguments.\allowbreak{}DenseThresholdCases.\allowbreak{}HdEx}{20--27}
\leanref{UpperTheory.\allowbreak{}V.\allowbreak{}Hi.\allowbreak{}exists\_\allowbreak{}hecke\_\allowbreak{}ideal}{Development.\allowbreak{}AutomorphicArithmetic.\allowbreak{}AutomorphicVarieties.\allowbreak{}ModularityArgument.\allowbreak{}ThresholdArguments.\allowbreak{}DenseThresholdCases.\allowbreak{}ExistsHeckeIdeal}{70--83}

At auxiliary place \(j\) of residue cardinality \(Q_j\) take the full integral group, the upper Iwahori and the kernel in that Iwahori of the map from diagonal residues to their \(\ell\)-primary quotient. Write \(M,M_w,M_1\) for their modules with the same away localisation. The last covering group is \(\Delta=\prod_{j,i}\Z/\ell^{a_j}\), where \(a_j=v_\ell(Q_j-1)\ge r\) at depth \(r\) by \ref{B6.56}. The distinguished congruence level makes the action free on each double-coset fibre. Select one sheet over each coset, multiply by its indicator and apply the ordinary projector. The resulting endomorphism \(e_1\) of ordinary functions satisfies \(\sum_{d\in\Delta}d e_1 d^{-1}=1\). Thus
\[
 v\longmapsto\sum_{d\in\Delta}d\otimes e_1(d^{-1}v)
\]
splits the action map \(A[\Delta]\otimes_A M_1\to M_1\) on the ordinary module before localisation. The equivariant residual Hecke and root projections preserve projectivity. On the regular module the norm identifies coinvariants with invariants by sending the class of \(1\) to \(\sum_d d\). The identity extends to projective modules. The invariants give \(M_w\), and the kernel of the norm is generated by \(d-1\).
\leanref{UpperTheory.\allowbreak{}V.\allowbreak{}Hi.\allowbreak{}ordProj}{Development.\allowbreak{}AutomorphicArithmetic.\allowbreak{}AutomorphicVarieties.\allowbreak{}HeckeAlgebras.\allowbreak{}OrdinaryHeckeCharacters.\allowbreak{}OrdProj}{379--500}

For the trace identity take the root operators \(u_{ji}\), diamond generators \(1+y_{ji}\) and multiplier \(f\) of \ref{B6.57}(1). On a classical generalised eigenspace order the Frobenius roots \(\alpha_i\) by the torus character at the place under consideration. The splitting algebra has rank \(b=n!\), by adjoining successive roots in degrees \(n,n-1,\ldots,1\), so \((u_i-\alpha_i)^b=0\) in the local factor. If two roots coincide or \(\alpha_i=Q_j\alpha_a\) with \(i\ne a\), the corresponding factor of \(f\) is a linear combination of commuting nilpotents of exponent \(b\). Its \(2b\)-th power vanishes. In the other case the roots are distinct and monodromy vanishes, since it connects roots of ratio \(Q_j\). The tame generator \(\tau_{r,j}\) acts on line \(i\) by \(1+y_{ji}\), giving
\[
 \prod_{a\ne i}(\phi_{r,j}-\alpha_a)(\tau_{r,j}-1-y_{ji})=0.
\]
For \(E=2b\), evaluation at \(\alpha_a\) shows that \((X-u_a)^{E+1}\) is divisible by \(X-\alpha_a\). Multiplication and trace give the relation of \ref{B6.57}(1). Classical specialisations separate the module, so it holds integrally. At the split auxiliary places take the Iwahori module obtained by the diamond norm and the spherical module at the same away level. The spherical Hecke polynomial supplies the signed elementary symmetric coefficients on spherical vectors. Lemma~\ref{B6.48i} applies to these modules and then to their ordinary residual summands. Its multiplier is exactly the product \(f\) of \ref{B6.57}(1).
\leanref{UpperTheory.\allowbreak{}V.\allowbreak{}Z.\allowbreak{}Win.\allowbreak{}T2.\allowbreak{}level\_\allowbreak{}relation\_\allowbreak{}two\_\allowbreak{}x}{Development.\allowbreak{}AutomorphicArithmetic.\allowbreak{}AutomorphicVarieties.\allowbreak{}ModularityArgument.\allowbreak{}ThresholdArguments.\allowbreak{}DenseThresholdCases.\allowbreak{}ReductionOfPolynomialElementPow}{105--130}
\leanref{UpperTheory.\allowbreak{}V.\allowbreak{}Pa.\allowbreak{}Lev.\allowbreak{}twoOf\_\allowbreak{}alt}{Development.\allowbreak{}AutomorphicArithmetic.\allowbreak{}AutomorphicVarieties.\allowbreak{}ModularityArgument.\allowbreak{}ThresholdArguments.\allowbreak{}DenseThresholdCases.\allowbreak{}LevTwoOf}{103--117}
\leanref{UpperTheory.\allowbreak{}V.\allowbreak{}Th.\allowbreak{}opOfX\_\allowbreak{}alt}{Development.\allowbreak{}AutomorphicArithmetic.\allowbreak{}AutomorphicVarieties.\allowbreak{}ModularityArgument.\allowbreak{}ThresholdArguments.\allowbreak{}DenseThresholdCases.\allowbreak{}OpOpOf}{139--171}
\end{proof}

\subsection*{Proof of \ref{B6.48d}}\label{prf:B6.48d}\phantomsection
\begin{proof}
Work in a finite extension of the coefficient field containing the values of the arithmetic specialization, the images of the local embeddings and the values of \(\alpha\). Such an extension suffices since \(A\) is a power series ring in finitely many variables, their images are topologically nilpotent, and a finite extension of \(\Q_\ell\) is complete. Write \(O'\) for its integer ring and \(\varpi\) for a uniformizer. We use the given embedding into \(\C\) when speaking of classical automorphic forms.

First specialize the coefficients of the scalar-valued ordinary functions, retaining the tame coefficient coordinates. This gives a map
\[
       j:O'\otimes_A Y\longrightarrow Y',
\]
where \(Y'\) is defined by the same ordinary condition with diagonal character \(x\theta\). Use the evaluation coordinates at precision one of Lemma~\ref{B6.48cb}. Denote them by \(e_i\), and denote the ordinary projections of the coordinate extensions by \(w_i\). Coefficient specialization commutes with finite double-coset sums and hence with the ordinary projector, by its factorial-power construction. The specialized \(w_i\) are therefore the projected extensions for \(Y'\). The endomorphism \(f\mapsto\sum_i e_i(f)w_i\) is the identity modulo the maximal ideal on each of \(Y,Y'\), by the precision-one approximation of Lemma~\ref{B6.48cb}. On each module it is an automorphism by Nakayama's lemma and finite generation. The map \(f'\mapsto\sum_i e_i(f')(1\otimes w_i)\) composes with \(j\) in the two orders to give these automorphisms, with scalars extended in the first one. Thus \(j\) is an isomorphism.
\leanref{UpperTheory.\allowbreak{}V.\allowbreak{}Hi.\allowbreak{}ordBC}{Development.\allowbreak{}AutomorphicArithmetic.\allowbreak{}AutomorphicVarieties.\allowbreak{}HeckeAlgebras.\allowbreak{}OrdinaryHeckeCharacters.\allowbreak{}OrdBCMapOrd}{198--233}

Let \(p_0\) be the lowest-weight coordinate functional, fixed by the upper unipotents. For a classical form \(F\) with algebraic coefficients define
\[
                  c(F)(g)=p_0\!\left(W_\lambda(g_\ell)^{-1}F(g)\right)
\]
in each tame coefficient coordinate. Choose the level deep enough that \(\alpha\) is trivial on its diagonal congruence subgroup and the lower unipotent subgroup preserves the coefficient lattice. The right transformation law for \(c(F)\) on the diagonal is then \((\chi_\lambda\alpha)^{-1}\). The formula commutes with the tame Hecke sums. For a positive diagonal \(a\), take upper-unipotent representatives for its double coset. Applying \(p_0\) to each translate gives
\[
        c\!\left(\chi_\lambda(a)^{-1}[a]F\right)=[a]c(F).
\]
In particular this holds for the ordinary operator \(U_z\) of Lemma~\ref{B6.48cb}.

Use the weight lattice in the algebraic representation, with the lowest-weight coordinate first. The upper integral unipotents act integrally and preserve this coordinate. On every other weight the normalized action of \(z\) is divisible by \(\varpi\). Indeed the weight differs from the lowest weight by a nonzero sum of positive roots, and each simple root takes a uniformizer value on the strictly dominant diagonal \(z\). Equivalently, in the integral weight basis all columns of \(\chi_\lambda(z)^{-1}W_\lambda(z)\) except the first vanish modulo \(\varpi\). Thus if the untwisted values \(W_\lambda(g_\ell)^{-1}F(g)\) are divisible by \(\varpi^a\) and their first coordinates by \(\varpi^{a+1}\), application of the normalized \(U_z\) makes all coordinates divisible by \(\varpi^{a+1}\). This follows term by term in the upper-unipotent sum. On the integral ordinary part a positive power of \(U_z\) is the identity at any fixed precision. If \(c(F)=0\), induction on \(a\) now puts \(F\) in every \(\varpi^a\) times the value lattice. Hence \(F=0\).
\leanref{UpperTheory.\allowbreak{}V.\allowbreak{}Hi.\allowbreak{}wtRep}{Development.\allowbreak{}AutomorphicArithmetic.\allowbreak{}AutomorphicVarieties.\allowbreak{}HeckeAlgebras.\allowbreak{}OrdinaryHeckeCharacters.\allowbreak{}ProdZpowConst}{95--260}
\leanref{UpperTheory.\allowbreak{}V.\allowbreak{}Hi.\allowbreak{}ctrlOrd}{Development.\allowbreak{}AutomorphicArithmetic.\allowbreak{}AutomorphicVarieties.\allowbreak{}HeckeAlgebras.\allowbreak{}OrdinaryHeckeCharacters.\allowbreak{}CtrlOrd}{220--372}

For surjectivity, let \(f\in Y'\). At precision one choose a deeper lower-unipotent level at which \(f\) is invariant and which acts trivially on the coefficient lattice modulo \(\varpi\). There are finitely many double cosets at this level. On their representatives prescribe a classical form by giving its untwisted first coordinate the value of \(f\) and its other coordinates zero. Right translation extends the values: the stabilizers act trivially on the coefficient fibres by the level condition of Lemma~\ref{B6.48cb}. The resulting integral form \(G\) satisfies \(c(G)\equiv f\bmod\varpi\) on all cosets. To return to the original lower-unipotent level, apply a power \(U_z^b\), with \(b\) at least the difference of the depths and divisible by a period of \(f\bmod\varpi\). Conjugation by \(z^{-b}\) sends the lower-unipotent subgroup at the original depth into the deeper one. The upper-unipotent coset sum gives invariance at the original depth; the finite diagonal character is already trivial on its congruence subgroup. Thus \(U_z^bG\) is classical at that level and its first-coordinate function is still congruent to \(f\). Applying the ordinary projector there preserves this congruence. We have obtained an integral classical ordinary lift modulo \(\varpi\).

On \(Y'\), divisibility of all values by \(\varpi\) is divisibility in the module: divide the values, and use periodicity one precision higher in checking the ordinary condition. Nakayama's lemma now gives surjectivity from the integral ordinary lattice. This lattice spans the classical ordinary space since, in the finite double-coset coordinates, the ordinary projector is the limit of factorial powers of an integral matrix and selects exactly its unit generalized eigenvalues. Inverting \(\varpi\), extending scalars, and applying the residual Hecke idempotent prove the assertions. This is the fixed-level function-space form of ordinary control; the corresponding specialization of ordinary Hecke algebras is given in \cite[Lemma~2.6.4]{geraghty-2019}.
\leanref{UpperTheory.\allowbreak{}V.\allowbreak{}Hi.\allowbreak{}Sp.\allowbreak{}absSpec}{Development.\allowbreak{}AutomorphicArithmetic.\allowbreak{}AutomorphicVarieties.\allowbreak{}HeckeAlgebras.\allowbreak{}HeckeFixedLevelSpaces.\allowbreak{}FreeMap}{425--588}
\leanref{UpperTheory.\allowbreak{}V.\allowbreak{}Hi.\allowbreak{}partAct0\_\allowbreak{}alt}{Development.\allowbreak{}AutomorphicArithmetic.\allowbreak{}AutomorphicVarieties.\allowbreak{}ModularityArgument.\allowbreak{}ThresholdArguments.\allowbreak{}DenseThresholdCases.\allowbreak{}HeckeOperatorPreservesPart}{308--335}
\end{proof}

\subsection*{Proof of \ref{B6.48a}}\label{prf:B6.48a}\phantomsection
\begin{proof}
Write \(\mathbb T=B/J\) for the Hecke algebra acting faithfully on \(M(K)\). The point \(x\) gives a character of \(\mathbb T\). The module is finite over \(\mathbb T\), and its support is all of \(\operatorname{Spec}\mathbb T\) since its annihilator is zero. Thus the local factor at this character in the characteristic-zero weight fibre of \(M(K)\) is nonzero. In that factor the last nonzero power of the maximal ideal acting on the module is killed by the maximal ideal. This gives a simultaneous eigenvector for the specified character, including its ordinary operators. By \ref{B6.48d} it is a classical form at the prescribed weight, finite diagonal character and away level. Decomposition of the definite automorphic space gives a constituent in which its projection is nonzero.

Apply definite base change \ref{B6.04a} to this constituent and the Galois attachment of \ref{B9.01} to its regular algebraic cuspidal factors, with the twists prescribed by base change. Their semisimple sum has the specified characteristic polynomials at good inverse Frobenii, by the defining spherical action of \ref{B6.48c}. Frobenius density and equality of semisimple characters identify its determinant on the level group with \(D_x\). Pullback to \(G_L\) preserves semisimplicity since \(p\) is surjective. The pullback has the determinant of \(r\), hence is isomorphic to \(r\). Irreducibility leaves one cuspidal factor, with one copy. Denote it by \(\pi_0\). Its infinitesimal character and conjugate duality are those of \ref{B6.04a}.
\leanref{UpperTheory.\allowbreak{}V.\allowbreak{}Z.\allowbreak{}Ord.\allowbreak{}core\_\allowbreak{}hecke\_\allowbreak{}of\_\allowbreak{}automorphic\_\allowbreak{}hecke\_\allowbreak{}dense\_\allowbreak{}weight}{Development.\allowbreak{}AutomorphicArithmetic.\allowbreak{}AutomorphicVarieties.\allowbreak{}ModularityArgument.\allowbreak{}ThresholdArguments.\allowbreak{}DenseThresholdCases.\allowbreak{}OrdCoreH}{20--61}

Suppose now that \(r\) is continuous. At a chosen split place above \(\ell\), put \(K=L_v\), \(q=|\mathcal O_K/\mathfrak m_K|\), \(\nu=|\cdot|_K\), and choose a uniformizer \(\varpi_K\). Let \(\kappa_i\), \(0\le i<n\), be the specialized ordinary characters and let \(h_{i,\tau}\) be their strictly increasing weights at \(\tau:K\hookrightarrow\overline{\Q}_\ell\). On the local Weil group use the reciprocity map \(a\) to \(K^\times\) for which a lift of \(\varpi_K\) has degree \(-1\).

We describe the smooth characters furnished by the same classical eigenvector. Its image in the upper-unipotent Jacquet module is nonzero. Indeed a vector in the kernel of the coinvariant map is a finite sum of unipotent differences. Averaging over a sufficiently large compact unipotent subgroup kills such a sum. On upper-integral-unipotent invariants these averages, followed by translation, are positive powers of the strictly dominant double-coset operator, up to nonzero index factors. That operator has nonzero eigenvalue on our vector. In the Jacquet module each chamber translation acts by its Hecke eigenvalue divided by the unipotent index. The diagonal units act by the inverse finite character. The finite-length torus module therefore has the resulting character as a constituent. More explicitly, localizing at that torus character gives a nonzero finite-length module, which has a quotient of that character. Frobenius reciprocity gives principal-series support for the local automorphic constituent.

Write its support, in the normalization for parabolic induction, as
\(\{\omega_i\nu^{(n-1)/2}:0\le i<n\}\), in the order read by the chamber operators. At diagonal units the \(\omega_i\) are the inverses of the finite parts in \ref{B6.48d}. For the diagonal with \(j\) uniformizers first, the scalar before division by \(\chi_\lambda\) is
\(\prod_{i<j}\omega_i(\varpi_K)q^{-i}\). Indeed the unipotent index is \(q^{j(n-j)}\); its half-density and the common twist \(\nu^{(n-1)/2}\) give \(q^{-j(j-1)/2}\). Taking successive quotients gives the individual values on \(\varpi_K\). Together with the unit values these are precisely the ordinary-character identities
\[
 \kappa_i(s)=
       \iota^{-1}(\omega_i(a(s)))
       \prod_{\tau}\tau(a(s))^{-h_{i,\tau}}
       \qquad(s\in W_K).
\]
Here the algebraic part of the \(i\)-th compact character has exponents \(h_{i,\tau}-i\); the other \(i\) units of weight are the cyclotomic correction in \ref{B6.48c}. Units and a uniformizer generate \(K^\times\), so the identity holds on the whole Weil group.
\leanref{UpperTheory.\allowbreak{}V.\allowbreak{}Z.\allowbreak{}Ord.\allowbreak{}ordJac}{Development.\allowbreak{}AutomorphicArithmetic.\allowbreak{}AutomorphicVarieties.\allowbreak{}ModularityArgument.\allowbreak{}OrdinaryJacquetModules.\allowbreak{}OrdJac}{13--124}
\leanref{UpperTheory.\allowbreak{}V.\allowbreak{}Hi.\allowbreak{}core\_\allowbreak{}eigen\_\allowbreak{}of\_\allowbreak{}weight\_\allowbreak{}char\_\allowbreak{}hecke\_\allowbreak{}level}{Development.\allowbreak{}AutomorphicArithmetic.\allowbreak{}AutomorphicVarieties.\allowbreak{}ModularityArgument.\allowbreak{}ThresholdArguments.\allowbreak{}DenseThresholdCases.\allowbreak{}TCoreEig}{189--215}
\leanref{UpperTheory.\allowbreak{}V.\allowbreak{}Z.\allowbreak{}Ord.\allowbreak{}endOrd\_\allowbreak{}of\_\allowbreak{}printsH}{Development.\allowbreak{}AutomorphicArithmetic.\allowbreak{}AutomorphicVarieties.\allowbreak{}ModularityArgument.\allowbreak{}ThresholdArguments.\allowbreak{}DenseThresholdCases.\allowbreak{}IsoOfDet}{124--150}

The local flag has these characters. To see this from the determinant relations, the trace of every global translate of
\(\prod_i(r(g_i)-\kappa_i(g_i))\) vanishes by \ref{B6.48c}. The global matrices span the full matrix algebra by irreducibility, so nondegeneracy of its trace pairing makes every such product zero. The spans of the images of the successive suffix products give a filtration from zero to the whole space. Its \(i\)-th quotient carries \(\kappa_i\), since left multiplication by \(r(g)-\kappa_i(g)\) takes that span into the preceding one. To see that every quotient is a line, take a Weil lift of a sufficiently deep rational unit \(1+\ell^b\) on which all smooth characters are trivial. The ordinary weights have strictly increasing sums over \(\tau\), so the values of the \(\kappa_i\) there are distinct. The specialized local characteristic polynomial has these \(n\) roots by the diagonal character relation of the ordinary Hecke action. Thus each occurs once. In particular every flag at this prolongation has the same multiset of characters.
\leanref{UpperTheory.\allowbreak{}V.\allowbreak{}Z.\allowbreak{}Ord.\allowbreak{}flagJ\_\allowbreak{}alt}{Development.\allowbreak{}AutomorphicArithmetic.\allowbreak{}AutomorphicVarieties.\allowbreak{}ModularityArgument.\allowbreak{}ThresholdArguments.\allowbreak{}DenseThresholdCases.\allowbreak{}OrdFlagJ}{102--133}

Assume the place is unlinked. For a pair \(i,j\) put
\(\delta=\kappa_i/(\kappa_j\varepsilon_\ell)\) and
\(d_\tau=h_{j,\tau}-h_{i,\tau}-1\). The local character formula for \(\varepsilon_\ell\) gives
\[
 \delta(s)=\iota^{-1}\!\left(
       \frac{\omega_i(a(s))}{\omega_j(a(s))\nu(a(s))}
                 \right)\prod_\tau\tau(a(s))^{d_\tau}.
\]
If \(\omega_i=\omega_j\nu\), evaluate at a Weil lift of \(\ell\in K^\times\). The character \(\delta\) has compact image on the decomposition group, so all its values have \(\ell\)-adic norm one. This gives \(\sum_\tau d_\tau=0\). For \(i<j\) all \(d_\tau\) are nonnegative; for \(i\ge j\) they are all negative. Consequently all are zero. Then \(\delta=1\) on the dense Weil group and by continuity on the decomposition group, contrary to unlinking. Thus the toral support of \(\pi_{0,v}\) is unlinked. Lemma~\ref{B9.38a} now applies to this support. The common power of \(\nu\) relating the two normalizations shifts all exponents by the same amount.

Suppose next that \(\delta^b=1\) for some \(b\ge1\). Evaluate at a Weil lift of a rational integer \(1+\ell^c>1\) sufficiently close to one to kill both smooth characters. This gives \((1+\ell^c)^{b\sum_\tau d_\tau}=1\). The sign argument again gives \(d_\tau=0\) for all \(\tau\). At a lift of \(\varpi_K\) the complex absolute value of \(\iota(\delta)\) is one, whence
\[
                   |\omega_j(\varpi_K)|=q\,|\omega_i(\varpi_K)|.
\]
The two support characters cannot satisfy this equality: their absolute values lie in the same open interval whose endpoints have ratio \(q\), by \ref{B9.38a}. This excludes every positive power in the definition of potential unlinking.
\leanref{UpperTheory.\allowbreak{}V.\allowbreak{}Z.\allowbreak{}Ord.\allowbreak{}exp\_\allowbreak{}torsion\_\allowbreak{}of\_\allowbreak{}literal\_\allowbreak{}automorphy\_\allowbreak{}n}{Development.\allowbreak{}AutomorphicArithmetic.\allowbreak{}AutomorphicVarieties.\allowbreak{}ModularityArgument.\allowbreak{}ThresholdArguments.\allowbreak{}DenseThresholdCases.\allowbreak{}OrdExpTor}{265--306}

Other prolongations over the chosen place are obtained by conjugation. Over the other place of the split pair, polarization takes the flag to the reversed dual flag with the common multiplier \(\varepsilon_\ell^{1-n}\). This follows also by comparing the conjugate and twisted-dual Frobenius polynomials and applying density to the semisimple representations. The common multiplier cancels in character ratios, and inversion interchanges the two indices in a cyclotomic equality. The automorphic supports transform by conjugate duality as well. Thus the conclusions hold at both places. In applying the support argument at a potentially unlinked place, its flag satisfies the unlinking condition by taking the first power.

Away from \(\ell\) and the distinguished place, \(S'=S\) leaves spherical places and the specified split tame types. At a spherical place the support is unramified toral, and a generic spherical constituent has unlinked characters: in the unramified classification its parameter has zero monodromy, and genericity excludes linked singleton segments. At a tame type place, the classical vector realizes the prescribed compact type. The type condition gives zero monodromy and diagonal inertia with the specified Teichmüller characters. Each of these inertia characters extends to the decomposition group, so every Weil commutator acts trivially. Apply \ref{B9.55a} to the algebraically normalized generic constituent. Undoing the common normalization preserves unlinking.
\leanref{UpperTheory.\allowbreak{}V.\allowbreak{}Hi.\allowbreak{}partUR2Type\_\allowbreak{}alt}{Development.\allowbreak{}AutomorphicArithmetic.\allowbreak{}AutomorphicVarieties.\allowbreak{}ModularityArgument.\allowbreak{}ThresholdArguments.\allowbreak{}DenseThresholdCases.\allowbreak{}AutomorphicTensorInductionOfLocalPlaces}{76--113}
\leanref{UpperTheory.\allowbreak{}Gal.\allowbreak{}projectedFormsSystemPart.\allowbreak{}typePlain}{Development.\allowbreak{}AutomorphicArithmetic.\allowbreak{}Foundations.\allowbreak{}GaloisRepresentations.\allowbreak{}PlainTorsionPairs.\allowbreak{}PlainTorOfMono0}{114--131}

Finally, any other cuspidal attachment has the same unramified Satake multisets as \(\pi_0\) at almost every place, by equality with the Frobenius polynomials of \(r\) in the same normalization. Strong multiplicity one identifies the two. The global twist making this one cuspidal representation unitary supplies one value of \(t\) in \eqref{eq:p6c5-toral-support} at all the required places, by \ref{B9.38a}.
\leanref{UpperTheory.\allowbreak{}V.\allowbreak{}Z.\allowbreak{}Ord.\allowbreak{}born\_\allowbreak{}j\_\allowbreak{}of\_\allowbreak{}printed\_\allowbreak{}polarization\_\allowbreak{}window}{Development.\allowbreak{}AutomorphicArithmetic.\allowbreak{}AutomorphicVarieties.\allowbreak{}ModularityArgument.\allowbreak{}ThresholdArguments.\allowbreak{}DenseThresholdCases.\allowbreak{}OrdBornJ}{32--50}
\end{proof}

\subsection*{Proof of \ref{B6.48ac}}\label{prf:B6.48ac}\phantomsection
\begin{proof}
Put \(S=L[[T_1,\ldots,T_h]]\) and let \(T_i\) act by \(y_i\). The power series act because \(B\) is complete and the \(y_i\) belong to its maximal ideal. Lift an \(L\)-basis of \(M/YM\) to \(m_1,\ldots,m_e\in M\). These lifts give a map \(S^e\to M\). The displayed isomorphism identifies its associated graded map for \((T_1,\ldots,T_h)\) and \(Y\) with an isomorphism in every degree. A nonzero vector of series has a first nonzero homogeneous part, so the map is injective. Given \(m\in M\), first represent its residue modulo \(YM\), then its successive remainders in \(Y^aM/Y^{a+1}M\). The resulting series map to \(m\), since \(\bigcap_aY^aM=0\). Thus \(M\) is finite free over \(S\).

Multiplication embeds \(B/\operatorname{Ann}_B M\) in the finite \(S\)-module \(\operatorname{End}_S(M)\), so that ring is finite over \(S\). A minimal support prime \(Z\) is associated to \(M\). If \(Z=\operatorname{Ann}_B(m)\) with \(m\ne0\), freeness over the domain \(S\) gives \(Z\cap S=0\), where contraction is through the action map. Hence \(B/Z\) is integral over a copy of \(S\). Lying over, going up and incomparability give \(\dim B/Z=\dim S=d_0+h\).
\end{proof}

\subsection*{Proof of \ref{B6.48ad}}\label{prf:B6.48ad}\phantomsection
\begin{proof}
In \(\mathbb P^{g-1}_P\), with line coordinate \([v]\), take the closed subscheme \(X\) given by \(Zv=0\) and \(u_iv_j=u_jv_i\). Its closed fibre is \(\mathbb P^{g-1}_{P/\mathfrak m}\). On the chart \(v_c=1\) its equations are
\[
 z_{ic}+\sum_{j\ne c}z_{ij}v_j=0,\qquad
 u_b-u_cv_b=0\quad(b\ne c).
\]
At a point of the closed fibre the ambient local ring is regular. The parameters of \(P\) remain independent in its cotangent space after extending the residue field. Indeed the chart is a polynomial algebra over \(P\), and parameters of \(P\) followed by lifts of parameters of the local ring of the fibre form a regular system of parameters. In the displayed equations the coefficients of the parameters \(z_{ic}\) and \(u_b\), \(b\ne c\), form an identity matrix. The quotient local ring is regular.

Every point of \(X\) specialises to a point of the closed fibre. Its closure is proper over \(\operatorname{Spec}P\), so its image is closed and contains the closed point. Regularity passes to localisations. Thus the local rings of \(X\) are domains and its finitely many irreducible components are disjoint. Each component meets the connected closed fibre. It follows that \(X\) is irreducible. Over the residue field of any prime containing \((Zu,\det Z)\), take the line spanned by \(u\) if \(u\ne0\), and any kernel line of \(Z\) otherwise. Conversely the field-valued incidence equations imply \(Zu=0\) and \(\det Z=0\). The image of \(X\) is therefore \(V(Zu,\det Z)\), which is irreducible.

To bound its dimension, add the \(g^2\) equations \(z_{ij}=\delta_{ij}u_i\) to \((Zu)\). They give \(u_i^2=0\) in the quotient. Its radical is the ideal of the \(g^2+g\) independent parameters. Cutting by \(g^2\) elements decreases dimension by at most \(g^2\), so \(\dim P/(Zu)\le\dim P-g\). The adjugate identity gives \((\det Z)u=0\bmod P_0\), hence \(\det Z\in P_0\). The prime radical found above lies below \(P_0\). Strict inclusion would give dimension greater than \(\dim P-g\) for its quotient, contradicting the bound.
\end{proof}

\subsection*{Proof of \ref{B6.48aad}}\label{prf:B6.48aad}\phantomsection
\begin{proof}
Extend the matrices and characters linearly to the group algebra. The product identity extends to independently chosen group-algebra elements by multilinearity. It holds also for \(r_0\), since its \(i\)-th factor takes the \((i+1)\)-st flag space into the \(i\)-th.

The evaluation vectors of the residual characters span \(F^n\) by linear independence of distinct characters. Choose \(z\in F[D]\) at which their values are pairwise distinct. Put \(a_{b,i}=\chi_{b,i}(z)\). Every difference \(a_{b,i}-a_{b,j}\), \(i\ne j\), is a unit. Taking every slot equal to \(z\) gives the annihilating polynomial \(\prod_i(X-a_{b,i})\) for \(r_b(z)\). Its Lagrange projectors \(P_{b,i}\) are orthogonal and sum to one. The projectors for the two systems agree modulo \(I\). They belong to the given subalgebra: the denominators are scalar matrices in it, and their inverses belong to it by the finite geometric series for a unit with nilpotent difference from a nonzero element of \(F\).

For \(i>j\), multiply the ordered identity on the left by \(P_{b,i}\) and on the right by \(P_{b,j}\). Use \(g\) in slot \(j\) and \(z\) in every other slot. The factors to the left of slot \(j\) act on the left projector by units, and the factors to its right act on the right projector by units. This gives \(P_{b,i}r_b(g)P_{b,j}=0\). The same calculation with \(i=j\) gives
\[
 P_{b,i}r_b(g)P_{b,i}=\chi_{b,i}(g)P_{b,i}.
\]
For \(b=0\) the sums of the first projectors give the original flag: they preserve that flag and act by one on precisely its first graded lines.

Put \(S=\sum_iP_{1,i}P_{0,i}\) and \(Y=1-S\). Then \(Y\in M_n(I)\), \(S^{-1}=1+Y\), and \(SP_{0,i}=P_{1,i}S\). Thus \(S^{-1}r_1(g)S\) preserves the original flag with ordered diagonal \(\chi_{1,i}(g)\). Multiplication by \(r_0(g)^{-1}\) and expansion using \(I^2=0\) give
\[
 S^{-1}r_1(g)Sr_0(g)^{-1}-1=Z(g)-dY(g).
\]
This gives both assertions, including the stated diagonal on every element of \(D\).
\end{proof}

\subsection*{Proof of \ref{B6.48aae}}\label{prf:B6.48aae}\phantomsection
\begin{proof}
At level \(N\) put \(Q_j=\mathrm Nw_{N,j}\), \(a_j=v_\ell(Q_j-1)\) and \(\Delta_N=\prod_{j,i}\Z/\ell^{a_j}\). In the deepest module of \ref{B6.48c} take the ordinary residual Hecke summand and the ordered-root summand on which every \(U_{ji}-\widetilde u_{ji}\) is topologically nilpotent. After extension by \(s\), call this \(M_N\). The sheet splitting and the equivariant projections give projectivity over \(\Lambda[\Delta_N]\). This ring is local because in its residue group algebra the augmentation ideal is nilpotent. Thus \(M_N\) is free over
\[
 \Lambda[\Delta_N]=S/\bigl((1+y_{ji})^{\ell^{a_j}}-1:j,i\bigr).
\]
The norm gives the corresponding Iwahori module at augmentation.
\leanref{UpperTheory.\allowbreak{}V.\allowbreak{}Hi.\allowbreak{}ordProj}{Development.\allowbreak{}AutomorphicArithmetic.\allowbreak{}AutomorphicVarieties.\allowbreak{}HeckeAlgebras.\allowbreak{}OrdinaryHeckeCharacters.\allowbreak{}OrdProj}{379--500}
\leanref{UpperTheory.\allowbreak{}V.\allowbreak{}Pa.\allowbreak{}Lev.\allowbreak{}opToTW}{Development.\allowbreak{}AutomorphicArithmetic.\allowbreak{}AutomorphicVarieties.\allowbreak{}PatchingModules.\allowbreak{}NilpotentLevelLowering.\allowbreak{}OpToTW}{181--286}

Apply Lemma~\ref{B6.48i} before passing to the limit in the auxiliary levels. In its notation let \(V\) be the Iwahori module, \(M\) the spherical module, \(i:M\to V\) inclusion and \(T:V\to M\) the unnormalised reverse coset sum. If \(D_N\) is the ordered splitting algebra, define \(\phi:D_N\otimes M\to V\) by \(p\otimes v\mapsto p(u)i(v)\). The symmetric polynomials in the root operators \(u\) act on spherical vectors by the signed Frobenius coefficients.

At one place put
\[
 D(x)=\prod_{a<b}(x_b-x_a),\qquad
 H_Q(x)=\prod_{a<b}(x_b-Qx_a),\qquad
 \partial h=D^{-1}\sum_\sigma\operatorname{sgn}(\sigma)\sigma h .
\]
For a subset \(A\) of pairs \(a<b\), set \(a_A=\prod_{(a,b)\in A}x_b\), \(c_A=\prod_{(a,b)\notin A}(-x_a)\). Expansion yields
\[
 \sum_A a_A\sigma(c_A)=\prod_{a<b}(x_b-x_{\sigma(a)})
 =\mathbf1_{\sigma=1}D(x),
\]
since a nonidentity permutation has \(\sigma(a)>a\) for some \(a\). In the Bernstein polynomial representation the adjacent operator acts by
\[
 T_s h=Q\,s(h)+(Q-1)x_a(h-s(h))/(x_a-x_{a+1}).
\]
The spherical sum \(E=\sum_\sigma T_\sigma\) obeys \(T_s E=QE=ET_s\), which determines its permutation coefficients from the longest one. Multiplication along a reduced word gives \(E=\sum_\sigma \sigma(H_Q/D)\sigma\), or \(E(h)=\partial(H_Qh)\). The triangular permutation expansion with nonzero leading coefficients identifies this computation with the Hecke operator identity. At all places define
\[
 \psi_0(v)=\sum_{(A_j)}\left(\prod_j a_{A_j}(x_j)\right)
 \otimes T\!\left(\prod_j c_{A_j}(u_j)v\right).
\]
The subset identity gives \(\sum_A a_A E c_A=H_Q\) and \(\sum_A a_A\partial(H_Qc_Ah)=H_Qh\). These compute \(\phi\psi_0\) and \(\psi_0\phi\) at one place. Operators at distinct places commute. Multiply \(\psi_0\) by the polynomial quotient of
\[
 f_N=\prod_j\prod_{i\ne b}(U_{ji}-U_{jb})(U_{ji}-Q_jU_{jb})
\]
by \(\prod_jH_{Q_j}\). Both composites now act by \(f_N\). The maps intertwine the root operators and restrict to the chosen residual ordering. Via the norm they compare \(M_N/(y)M_N\) with the base module tensored with the splitting algebra.
\leanref{UpperTheory.\allowbreak{}V.\allowbreak{}AlgebraTools.\allowbreak{}hkDualEx}{Development.\allowbreak{}AutomorphicArithmetic.\allowbreak{}AutomorphicVarieties.\allowbreak{}AlgebraTools.\allowbreak{}HkDualEx}{14--42}
\leanref{UpperTheory.\allowbreak{}V.\allowbreak{}Pa.\allowbreak{}Lev.\allowbreak{}multiTwo}{Development.\allowbreak{}AutomorphicArithmetic.\allowbreak{}AutomorphicVarieties.\allowbreak{}PatchingModules.\allowbreak{}NilpotentLevelLowering.\allowbreak{}MultiTwo}{174--315}
\leanref{UpperTheory.\allowbreak{}V.\allowbreak{}Pa.\allowbreak{}Lev.\allowbreak{}twoOf\_\allowbreak{}alt}{Development.\allowbreak{}AutomorphicArithmetic.\allowbreak{}AutomorphicVarieties.\allowbreak{}ModularityArgument.\allowbreak{}ThresholdArguments.\allowbreak{}DenseThresholdCases.\allowbreak{}LevTwoOf}{103--117}

Write \(Y=(y_{ji})\), \(\mathfrak l=\mathfrak m_\Lambda\) and \(X=B^{\rm sp}\otimes M_0\), with tensor product over the base of the splitting algebra after extending \(M_0\). Observe that \(f\in R\) equals \(f_N\) with \(Q_j=1\) for all \(j\). If \(a_i\ge b+j\), the diamond relations lie in \(Y^b+\mathfrak l^j\), since for \(1\le m<b\),
\[
 v_\ell\binom{\ell^{a_i}}{m}\ge a_i-v_\ell(m)\ge j .
\]
Hence the truncated modules are free over \(S/(Y^b+\mathfrak l^j)\) at sufficiently high levels. At augmentation, for each \(j\) choose \(e\) with \(\mathfrak m_{B_0}^eM_0\subset\mathfrak l^jM_0\) by finite length modulo \(\mathfrak l^j\). The inverse Frobenii tend to \(s_i\) in the finite quotient determining the polynomials at that precision. The ordered splitting presentations then agree by identification of corresponding roots, and the comparison maps intertwine the actions of \(R\). Their composites reduce to \(f\) since \(Q_i\equiv1\bmod\ell^N\).
\leanref{UpperTheory.\allowbreak{}V.\allowbreak{}Pa.\allowbreak{}Core.\allowbreak{}freeN}{Development.\allowbreak{}AutomorphicArithmetic.\allowbreak{}AutomorphicVarieties.\allowbreak{}PatchingModules.\allowbreak{}AnnihilatorBaseChange.\allowbreak{}FreeNQuot}{29--45}
\leanref{UpperTheory.\allowbreak{}V.\allowbreak{}Pa.\allowbreak{}Core.\allowbreak{}eventually\_\allowbreak{}exists\_\allowbreak{}patchedModule\_\allowbreak{}quotient\_\allowbreak{}maps}{Development.\allowbreak{}AutomorphicArithmetic.\allowbreak{}AutomorphicVarieties.\allowbreak{}PatchingModules.\allowbreak{}AnnihilatorBaseChange.\allowbreak{}EventuallyExistsPatchedModuleQuotientMaps}{193--240}

For \(b\ge1\) set
\[
 L_b=\prod_N M_N\big/\{(v_N):\
 \forall j,\ v_N\in(Y^b+\mathfrak l^j)M_N
 \text{ on a set in the ultrafilter}\}.
\]
Use \(\mathcal B=\widehat{R[1/f]}^{\,Y}\) and \(\mathcal M=\varprojlim_b L_b[1/f]\). At \(b=1\) choose augmentation comparisons at depths tending to infinity in the ultrafilter. The constant-module limit is \(X\), by finiteness of its truncations and completeness. Both comparison composites become invertible after inverting \(f\). At finite precision freeness identifies the \(Y\)-graded pieces with the copies of the augmentation indexed by monomials. Each degree uses finitely many monomials, so the identities pass to \(L_b\) and through localisation. The transitions are surjective with kernels given by powers of \(Y\). Hence
\[
 \mathcal M/Y\mathcal M\simeq\mathcal B\otimes_R X,\qquad
 \operatorname{gr}_Y\mathcal M\simeq(\mathcal M/Y\mathcal M)[y_{ji}].
\]
Lift finite generators of the augmentation successively through the complete tower to obtain finiteness over \(\mathcal B\).
\leanref{UpperTheory.\allowbreak{}V.\allowbreak{}Pa.\allowbreak{}Core.\allowbreak{}pmodEx}{Development.\allowbreak{}AutomorphicArithmetic.\allowbreak{}AutomorphicVarieties.\allowbreak{}PatchingModules.\allowbreak{}AnnihilatorBaseChange.\allowbreak{}PmodEx}{32--96}

Every support prime contains \(\mathfrak a\). On classical eigenspaces where \(f_N\ne0\) the auxiliary projector onto line \(i\) is
\[
 \frac{\prod_{b\ne i}(\phi_j-U_{jb})}{\prod_{b\ne i}(U_{ji}-U_{jb})}.
\]
Local compatibility gives tame character \(1+y_{ji}\) there. Thus at every classical eigencharacter \(f_N\) times each auxiliary translated trace vanishes, with the zero of \(f_N\) accounting for coincident roots and ratios \(Q_j\). The permutation filtration of the Jacquet module has length at most \(n!\). The corresponding power of the operator product kills the generalised eigenspaces, and the classical specialisations separate vectors. Polarisation, ordinary and prescribed inertia relations hold at finite level. For each generator \(a\) of \(\mathfrak a\), finite-quotient agreement of the common-cover lifts and of \(f_N,f\) makes a fixed power of \(fa\) agree to arbitrary maximal-ideal order with finite-level annihilators. The finite-length augmentation comparison and the monomial graded-piece identity pass these approximations through every truncation. After inverting \(f\), \(a\) belongs to the radical of the annihilator.
\leanref{UpperTheory.\allowbreak{}V.\allowbreak{}Z.\allowbreak{}Win.\allowbreak{}T2.\allowbreak{}level\_\allowbreak{}relation\_\allowbreak{}one\_\allowbreak{}two}{Development.\allowbreak{}AutomorphicArithmetic.\allowbreak{}AutomorphicVarieties.\allowbreak{}ModularityArgument.\allowbreak{}ThresholdArguments.\allowbreak{}DenseThresholdCases.\allowbreak{}LocOkLevelModuleOfThetaData}{118--153}
\end{proof}

\subsection*{Proof of \ref{B6.48aa}}\label{prf:B6.48aa}\phantomsection
\begin{proof}
Use the base ordinary module \(M_0\) of \ref{B6.48c}, localised at the residual system. The action sends determinant coefficients and ordered characters to the Hecke coefficients and character operators, with kernel \(J_0\). The initial support hypotheses in (a) and (b) give \(M_0\ne0\).
\leanref{UpperTheory.\allowbreak{}V.\allowbreak{}Hi.\allowbreak{}hida\_\allowbreak{}unitary\_\allowbreak{}of\_\allowbreak{}printed\_\allowbreak{}hida}{Development.\allowbreak{}AutomorphicArithmetic.\allowbreak{}AutomorphicVarieties.\allowbreak{}ModularityArgument.\allowbreak{}ThresholdArguments.\allowbreak{}DenseThresholdCases.\allowbreak{}THidaU}{38--48}

Put \(A=k[[t]]\), \(F=k((t))\). In each block choose a stable lattice and diagonalise the polarising form, with entries \(t^{d_i}b_i\) for units \(b_i\). Integrality of \(\rho(s)\) and \(P\rho(s)P^{-1}\) bounds entry \((i,j)\) below by \(\max(0,d_j-d_i)\) in valuation. Replace \(t\) by \(t^2\), extend \(k\) to take the square roots of the units, and rescale vector \(i\) by \(t^{-d_i}b_i^{-1/2}\), with \(b_i\) read after base change. The form becomes the identity. The entry bound becomes \(2\max(0,d_j-d_i)+d_i-d_j\ge0\), so the matrices remain integral. The same argument applies to one block. These injective DVR extensions preserve the kernels. Intersect the flags with the lattices. The steps are saturated, and in the two-block case the distinct local characters let us choose the flag bases with each column in its assigned block.

Write \(d_rY(s)=r(s)Yr(s)^{-1}-Y\). An ordinary adjoint cocycle transforms under \(c\) by negative transpose and vanishes on the prescribed inertia kernels. At \(v\mid\ell\), subtracting \(d_rY_v\) must make the cocycle upper triangular in the flag basis, with zero diagonal on the preimage of the \(\ell\)-primary torsion of the local abelianisation. Global boundaries use \(Y^t=-Y\). Laurent coefficients have bounded entries, so a power of \(t\) makes them integral and adically continuous. Call matrices even on the diagonal blocks and odd on the off-diagonal blocks. The integral odd cocycles form a finite \(A\)-module. Their values on generators of the maximal pro-\(\ell\) quotient of the open residual kernel and on finitely many coset representatives determine them. The finite-character hypothesis of \ref{B6.56} bounds the number of these generators. Choose integral ordinary odd cocycles \(\kappa_1,\ldots,\kappa_g\) giving a basis of the classes over \(F\). For one block take \(g=0\).

Use the commutative coefficient algebra with basis \(1,e_i,e'_i,\eta\), \(1\le i\le g\), in which \(e_i e'_j=\delta_{ij}\eta\) and all other products of positive-degree elements vanish. Set
\[
 \rho_1(s)=\left(1+\sum_i e_i\kappa_i(s)\right)\rho(s).
\]
For a flag basis \(w_v\) and local witnesses \(Y_{iv}\) for \(\kappa_i\), its flag basis is \((1-\sum_i e_iY_{iv})w_v\). Over a coefficient ring \(X\), let \(V_1(X)\) have even entries in \(X\) and odd entries in \(\sum_i Xe_i\). Let \(V_2(X)\) have even entries in \(\sum_i Xe'_i\) and odd entries in \(X\eta\). Both use conjugation by \(\rho_1\). The trace dual of \(V_1\) has even entries in \(X\eta\) and odd entries in \(\sum_i Xe'_i\), with pairing given by the \(\eta\)-coefficient of trace. Use the cyclotomic twist of conjugation on this dual. On \(H'\) the tests on its cocycles are \(z\mapsto[\eta]\operatorname{tr}(z(s)\rho_1(s))\).

To prove detection, write \(z=\eta p+\sum_b e'_b l_b\) and extend \(K_b(s)=\kappa_b(s)\rho(s)\) and \(L_b(s)=l_b(s)\rho(s)\) linearly on \(F[H']\). Absolute irreducibility and inequivalence give a surjection from this algebra to the product of the full matrix blocks. For \(c_1,c_2\) in its kernel the product rule kills the terms with \(\rho(c_1)\) or \(\rho(c_2)\). Multiply by lifts of the diagonal matrix units selecting opposite-block entries \(i,j\). The trace test gives
\[
 \sum_b\bigl(K_b(c_1)_{ij}L_b(c_2)_{ji}
 +L_b(c_1)_{ij}K_b(c_2)_{ji}\bigr)=0.
\]
For fixed \(c_2\), multiplication on each side by lifts of matrix units shows that this combination kills all relations in the corresponding corner. It descends to a derivation of the product algebra in its corner bimodule. Applying the derivation law to the central idempotents shows that this derivation is inner. The untwisted and twisted parts restrict to distinct eigenspaces for conjugation by an element of \(H\) of nontrivial residual cyclotomic value, which exists since \(\mu_\ell\not\subset L\). Each part is therefore inner on \(H'\). An odd cocycle vanishing on \(H'\) takes values in intertwiners between the two blocks, hence vanishes. Polarised transposition determines the opposite corner. Independence of the \(\kappa_b\) now gives \(L_b(c_2)_{ji}=0\) for every \(b,c_2\). The corner derivation argument gives
\[
 l_b(s)=\varepsilon_\ell(s)\operatorname{Ad}(\rho(s))Y_b-Y_b.
\]
Subtract the twisted boundary of \(\sum_b e'_bY_b\), whose even term is \(\eta\varepsilon_\ell(s)\sum_b[\kappa_b(s),\operatorname{Ad}(\rho(s))Y_b]\). For the remaining even derivative, apply the trace identity to a relation multiplied by a group element. Nondegeneracy of trace on the full blocks makes the derivative vanish on all relations. Its descent is inner. Subtracting this boundary leaves a cocycle on the cyclotomic quotient with values in the block scalars. If \(\bar\varepsilon_\ell(h)\ne1\), its cocycle law gives
\[
 (\bar\varepsilon_\ell(h)-1)z(s)=(\bar\varepsilon_\ell(s)-1)z(h).
\]
Thus it is a boundary. Average a bounding matrix with its polarised image, using \(2\in A^\times\), and clear denominators. This gives detection modulo \(t\)-power torsion. For one block the even derivative argument applies directly.

The intersection of the separable loci of the finitely many curves is dense and open in \(H'\). Continuity gives detection on that intersection. Choose tests decreasing the dimension of the common kernel modulo boundaries, take the union for the curves and call the resulting elements \(s_1,\ldots,s_q\). Split their polynomials and equalise the residue fields by finite injective extensions of the DVRs.
\leanref{UpperTheory.\allowbreak{}V.\allowbreak{}Pa.\allowbreak{}Det.\allowbreak{}relPair}{Development.\allowbreak{}AutomorphicArithmetic.\allowbreak{}AutomorphicVarieties.\allowbreak{}PatchingModules.\allowbreak{}DeterminantPadicFree.\allowbreak{}RelPair}{23--177}
\leanref{UpperTheory.\allowbreak{}V.\allowbreak{}Pa.\allowbreak{}Det.\allowbreak{}ann}{Development.\allowbreak{}AutomorphicArithmetic.\allowbreak{}AutomorphicVarieties.\allowbreak{}PatchingModules.\allowbreak{}DeterminantPadicFree.\allowbreak{}BlockPrinc}{75--161}
\leanref{UpperTheory.\allowbreak{}V.\allowbreak{}Pa.\allowbreak{}Det.\allowbreak{}detAlg}{Development.\allowbreak{}AutomorphicArithmetic.\allowbreak{}AutomorphicVarieties.\allowbreak{}PatchingModules.\allowbreak{}DeterminantPadicFree.\allowbreak{}DetAlg}{71--142}
\leanref{UpperTheory.\allowbreak{}V.\allowbreak{}Pa.\allowbreak{}choice}{Development.\allowbreak{}AutomorphicArithmetic.\allowbreak{}AutomorphicVarieties.\allowbreak{}PatchingModules.\allowbreak{}AutomorphicDetection.\allowbreak{}Choice}{14--41}
\leanref{UpperTheory.\allowbreak{}V.\allowbreak{}Pa.\allowbreak{}Kap}{Development.\allowbreak{}AutomorphicArithmetic.\allowbreak{}AutomorphicVarieties.\allowbreak{}PatchingModules.\allowbreak{}PairedDualFirstOrderDualMatrixSpace}{306--315}
\leanref{UpperTheory.\allowbreak{}V.\allowbreak{}Pa.\allowbreak{}Detects}{Development.\allowbreak{}AutomorphicArithmetic.\allowbreak{}AutomorphicVarieties.\allowbreak{}PatchingModules.\allowbreak{}ThetaCoefficientMapEq}{257--258}

Apply \ref{B6.56} to these elements. The residual action on the kernel of the diagonal characters is unipotent and its congruence kernels are pro-\(\ell\). Put \(d_+=[L^+:\Q]\), \(D_*=n(d_++q)\), and \(\nu=\binom a2+\binom{n-a}2+ga(n-a)\) for first block size \(a\), with \(a=n\) in the one-block case. At auxiliary level \(N\) and depth \(m\) set \(A_m=A/(t^m)\). Let \(W_m\) be the ordinary \(V_1(A_m)\)-cocycles for the pulled-back \(\rho_1\) modulo skew-transpose boundaries. Use the full local cohomology condition at each added place.
\leanref{UpperTheory.\allowbreak{}V.\allowbreak{}Pa.\allowbreak{}SelZ}{Development.\allowbreak{}AutomorphicArithmetic.\allowbreak{}AutomorphicVarieties.\allowbreak{}PatchingModules.\allowbreak{}PairedDualFirstOrderDualMatrixSpace}{240--244}
\leanref{UpperTheory.\allowbreak{}V.\allowbreak{}Pa.\allowbreak{}Sel1}{Development.\allowbreak{}AutomorphicArithmetic.\allowbreak{}AutomorphicVarieties.\allowbreak{}PatchingModules.\allowbreak{}PairedDualFirstOrderDualMatrixSpace}{344--344}
\leanref{UpperTheory.\allowbreak{}V.\allowbreak{}Pa.\allowbreak{}TWGrp.\allowbreak{}B1}{Development.\allowbreak{}AutomorphicArithmetic.\allowbreak{}AutomorphicVarieties.\allowbreak{}PatchingModules.\allowbreak{}ThetaCoefficientMapEq}{54--56}

For the obstruction group use the odd module \(M\) of \(\rho\) over \(A_m\), with \(c\) acting by negative transpose. At \(v\mid\ell\) let \(\pi_v:M\to Q_v\) quotient by the upper triangular odd matrices in the flag basis and take \(J_v=G_v\). At a prescribed tame group use \(\pi_v=1\), \(Q_v=M\), and \(J_v\) its prescribed inertia subgroup. A degree-two representative consists of a global two-cocycle \(b\) and continuous local one-cochains \(a_v\) with \(da_v=\pi_v(b|_{G_v})\). Quotient by the pairs with \(b=dc_1\) and \(a_v-\pi_v(c_1|_{G_v})\) a cocycle restricting to a boundary on \(J_v\). Call the group \(E_m\). For each \(m\), at all sufficiently high \(N\),
\begin{equation}\tag{6.48a.1}\label{eq:p6c5-patch-cohomology}
 \log_{\#k}\#W_m\le mD_*+C,\qquad
 \log_{\#k}\#(E_m/t^eE_m)\le ge+C'\quad(1\le e\le m),
\end{equation}
with \(C,C'\) independent of \(m,N\), by Lemma~\ref{B6.48aaa}.

The local factors used for the first bound have exponents at most
\begin{equation}\tag{6.48a.2}\label{eq:p6c5-patch-local}
 m[n(1+d'_v)+\nu(1+d_v)]+C_v\quad(v\mid\ell),
 \qquad mn+C_j\quad(v\text{ auxiliary}).
\end{equation}
Here \(d'_v\) is the free inertial abelian rank and \(d_v\) the local degree in the strict upper count, with both lists summing to \(d_+\). At an added place the pulled-back representation kills the tame generator \(\tau_{N,j}\). The tame relation, the congruence \(Q_j\equiv1\pmod{\ell^N}\) and the cocycle equation imply that its cocycle value commutes modulo \(t^m\) with \(\rho_1(h_N(\phi_{N,j}))\). At fixed \(m\), this matrix equals \(\rho_1(s_j)\) modulo \(t^m\) for sufficiently high \(N\). Over \(F\) the centralizer in the paired coefficient module has dimension \(n\). In a block eigenbasis the constant term is diagonal, and the simple root differences determine the off-diagonal coefficients in the first-order slots from that diagonal. Equivalently, one can diagonalise the lifted matrix by a first-order change of basis over \(F\), where the root differences are invertible. Apply Smith reduction over \(A\) to the commutator map. Each free kernel direction contributes \(m\) to the length of the truncated kernel. A nonzero elementary divisor \(t^b\) contributes \(\min(m,b)\), which is bounded by its fixed valuation \(b\). This gives \(mn+C_j\) for the inertia term of the full local condition after division by the unramified factor.
\leanref{UpperTheory.\allowbreak{}V.\allowbreak{}Pa.\allowbreak{}h0q}{Development.\allowbreak{}AutomorphicArithmetic.\allowbreak{}AutomorphicVarieties.\allowbreak{}PatchingModules.\allowbreak{}EpsilonEtaMatrices.\allowbreak{}FinrankCentralizerW1Le}{106--146}

At \(D_0\) the twisted intertwiner equations have zero rational kernel, so Smith reduction bounds their truncated invariant lengths. The detecting evaluations bound the dual kernel modulo boundaries by the same argument. At \(\ell\) the diagonal abelian characters contribute \(n(1+d'_v)\) and the strict upper terms contribute at most \(\nu(1+d_v)\) per unit of depth. Finite duality in \ref{B6.09} and the global Euler characteristic \cite[Ch.~II, \S5]{serre-2002} subtract the boundary and flag contribution \(m[(n+\nu)\#\{v\mid\ell\}+\nu d_+]\) from the sum of \eqref{eq:p6c5-patch-local}. The remaining coefficient of \(m\) is \(D_*\).

For the second bound the product defect of a \(V_2\)-lift gives \(b\) and the local defects give \(a_v\). Changing the lifts changes the pair by its defining boundaries. The depth sequence \(m-e,m,e\) bounds \(E_m/t^eE_m\) by the obstruction in depth \(e\) and the degree-two localisation cokernel. The base Selmer term in \ref{B6.09} counts ordinary odd classes. Evaluate cocycles on finitely many generators and include local witnesses. The rows of the equations generate a finite \(A\)-module. Smith reduction bounds the excess length of the truncated solution space over the integral solutions by the sum of the valuations of its nonzero elementary divisors. The integral classes have rank \(g\), giving \(ge+C_0\). At added places the odd centralizer has zero rational kernel since the block spectra are disjoint. At ordinary places the lower graded ratios are nontrivial, and a separating element bounds the global dual odd invariants. The truncated invariant lengths are thus bounded by fixed elementary divisors. By \ref{B6.09} with all ordinary and added places included, the global dual invariants bound the degree-two localisation cokernel. This gives \(C'\).
\leanref{UpperTheory.\allowbreak{}V.\allowbreak{}Pa.\allowbreak{}B2.\allowbreak{}sel0}{Development.\allowbreak{}AutomorphicArithmetic.\allowbreak{}AutomorphicVarieties.\allowbreak{}PatchingModules.\allowbreak{}ArchimedeanH0Cards.\allowbreak{}ImgCard}{340--365}

The \((g+1)\)-st largest elementary divisor of \(E_m\) has length at most \(C'\), by applying \eqref{eq:p6c5-patch-cohomology} at that length. For \(g+1\) obstruction vectors, lift their projection matrix on the first \(g\) cyclic summands to \(A\). A primitive vector in its kernel gives a combination supported in the other summands. Multiplication by \(t^{C'}\) kills it, so this multiple of the primitive combination lifts.
\leanref{UpperTheory.\allowbreak{}V.\allowbreak{}Pa.\allowbreak{}W1\_\allowbreak{}bound\_\allowbreak{}arithmetic\_\allowbreak{}rank\_\allowbreak{}two}{Development.\allowbreak{}AutomorphicArithmetic.\allowbreak{}AutomorphicVarieties.\allowbreak{}PatchingModules.\allowbreak{}EpsilonEtaMatrices.\allowbreak{}W1BoundArithmeticRankTwo}{40--166}
\leanref{UpperTheory.\allowbreak{}V.\allowbreak{}Pa.\allowbreak{}B2.\allowbreak{}obstructionBound\_\allowbreak{}of\_\allowbreak{}duality}{Development.\allowbreak{}AutomorphicArithmetic.\allowbreak{}AutomorphicVarieties.\allowbreak{}PatchingModules.\allowbreak{}ArchimedeanH0Cards.\allowbreak{}ObstructionBoundOfDuality}{145--217}

For an open normal subgroup \(H\) of a profinite group, write \(N_\ell(H)\) for the kernel of its maximal pro-\(\ell\) quotient. At each level in \ref{B6.56}, divide the level group by this kernel for the inverse image of the base subgroup \(H\). Their open pro-\(\ell\) subgroups have uniformly bounded generator numbers. Take surjections \(\pi_N:\Phi\to\Gamma_N\) onto these quotients from a free profinite group on lifts of the common finite quotient, a lift of complex conjugation and enough pro-\(\ell\) generators. Fix a nonprincipal ultrafilter on the levels. In each finite quotient of \(\Phi\), take the image of \(\ker\pi_N\) occurring on a set in the ultrafilter. Intersect their inverse images to obtain a closed normal \(K\), and put \(\Gamma_\infty=\Phi/K\). The marked elements and local maps pass to the limit in finite quotients. The limiting \(\tau_j,\phi_j\) commute by (6.56.1)--(6.56.2) and map to \(1,s_j\) in the base group.

At fixed depth, generator and local witness values belong to finite sets, and a representation with its cocycle is a homomorphism to a finite semidirect product. The simultaneous transfer of cocycles and boundaries at this precision, and the passage through all depths, are given in Lemma~\ref{B6.48aab}. Compactness gives compatible lifts in all depths, retaining a unit coordinate in the primitive combination before multiplying by \(t^{C'}\). Clearing a common denominator and inverting \(t\) gives a nonzero liftable combination of any \(g+1\) tuples of even classes. For \(b\) independent \(V_1(F)\)-classes, Smith reduction on finite generator evaluations modulo boundaries gives at least \((\#k)^{(m-c)b}\) classes at depth \(m\) for fixed \(c\). The first bound of \eqref{eq:p6c5-patch-cohomology} gives \(b\le D_*\). An odd cocycle on the limit has value at each \(\tau_j\) in the odd centralizer of \(\rho(s_j)\), hence zero. The \(\tau_j\) normally generate the kernel to the base in the pro-\(\ell\) quotient. Therefore \([\kappa_j]\) still form a basis of the ordinary odd classes.
\leanref{UpperTheory.\allowbreak{}V.\allowbreak{}Pa.\allowbreak{}cohLim}{Development.\allowbreak{}AutomorphicArithmetic.\allowbreak{}AutomorphicVarieties.\allowbreak{}PatchingModules.\allowbreak{}AutomorphicDetection.\allowbreak{}CohLim}{12--20}
\leanref{UpperTheory.\allowbreak{}V.\allowbreak{}Pa.\allowbreak{}koszulThree\_\allowbreak{}of\_\allowbreak{}det}{Development.\allowbreak{}AutomorphicArithmetic.\allowbreak{}AutomorphicVarieties.\allowbreak{}PatchingModules.\allowbreak{}SeparableMatrixStability.\allowbreak{}KoszulThreeOfDet}{8--50}

Extend \(O\) unramified to the common residue field and adjoin a power series variable \(s\) evaluated at \(t\) on each curve. Lift the curve and shared primes through the finite free extension by lying-over, going-down and going-up, using the extended shared primes for the power series variable. Adjoin the ordered roots at \(s_j\) by successive monic root adjunctions and division by the corresponding linear factors. The local factor \(B^{\rm sp}\) for the prescribed residual order is nonzero and finite flat over the enlarged base. Lift the primes through this extension by the same prime-lifting argument. The ordered roots are integral on the extended curve DVRs.

Put \(d=nd_+\), \(\Lambda=O[[s,w_1,\ldots,w_d]]\) and \(S=\Lambda[[y_{ji}]]\), where \(w_b\) are the weight parameters and there are \(nq\) variables \(y_{ji}\). Use the common complete local ring \(R\) for the continuous determinant on the CM subgroup of \(\Phi\), the ordered characters with their prescribed torsion values, and the independent variables \(s,U_{ji}-\widetilde u_{ji},y_{ji}\), with \(\widetilde u_{ji}\) lifting the residual roots. Pullback of the base data, the adjoined roots and \(y_{ji}\mapsto0\) give an adically surjective map \(\pi:R\to B^{\rm sp}\).

Let \(\mathfrak a\subset R\) be generated by the polarisation relations and the traces of the following words after left multiplication by arbitrary group elements. Use \([k]-1\) for the kernel to the limit, the ordered ordinary products and prescribed inertia differences of (b), and
\[
 \prod_i([\phi_j]-U_{ji}),\qquad
 \prod_{b\ne i}([\phi_j]-U_{jb})([\tau_j]-1-y_{ji}).
\]
Use all lifts of the indicated elements to the common cover. These relations vanish under \(\pi\). For the preimage \(p_i\subset R\) of curve prime \(i\), set \(C_i=\widehat{(R/\mathfrak a)_{p_i/\mathfrak a}}\). Its coefficient DVR is \(O_*=\widehat{O[[s]]_{(\varpi)}}\), with \(\varpi\) a uniformiser of \(O\) and residue field \(F=k((t))\) via \(s\mapsto t\).
\leanref{UpperTheory.\allowbreak{}V.\allowbreak{}Pa.\allowbreak{}Core.\allowbreak{}rrEx}{Development.\allowbreak{}AutomorphicArithmetic.\allowbreak{}AutomorphicVarieties.\allowbreak{}PatchingModules.\allowbreak{}AnnihilatorBaseChange.\allowbreak{}FiniteOpenHomsToFiniteGroup}{225--329}
\leanref{UpperTheory.\allowbreak{}V.\allowbreak{}Pa.\allowbreak{}Core.\allowbreak{}Rel}{Development.\allowbreak{}AutomorphicArithmetic.\allowbreak{}AutomorphicVarieties.\allowbreak{}PatchingModules.\allowbreak{}AnnihilatorBaseChange.\allowbreak{}SpecializationMapCompUniversalSection}{103--117}
\leanref{UpperTheory.\allowbreak{}V.\allowbreak{}Pa.\allowbreak{}ooEx}{Development.\allowbreak{}AutomorphicArithmetic.\allowbreak{}AutomorphicVarieties.\allowbreak{}PatchingModules.\allowbreak{}SelmerScaling.\allowbreak{}OoExKer}{106--196}

Fix \(C=C_i\). Form the finite Cayley--Hamilton algebra \(E\) over \(C\) from the group algebra of the CM limit group by the closed Cayley--Hamilton, ordinary, inertia and auxiliary relations and their prescribed adjoints. The curve representation maps \(E\) onto the full block algebras over \(F\). Lift the diagonal matrix units to orthogonal self-adjoint idempotents \(p_b\) summing to one. At each step average a lift with its adjoint in the complementary corner and use the projector for the Hensel factor at eigenvalue one. The determinant rank of \(p_b\) is one by reduction, so Cayley--Hamilton gives \(p_bzp_b=\operatorname{tr}(p_bzp_b)p_b\). In each block choose a distinguished index \(o(b)\) and lift \(E_{o(b),b}\) to \(x_b\in p_{o(b)}Ep_b\). Rescale by the square root of a scalar unit congruent to one to get \(x_b^*x_b=p_b\), \(x_bx_b^*=p_{o(b)}\).

For two blocks let \(p,p'\) be their distinguished idempotents, \(N=pEp'\) and \(\beta(v,w)=\operatorname{tr}(vw^*)\). The corner identities give \(\beta(v,w)z=v(w^*z)=\beta(w,z)v\). The pairing is symmetric and vanishes in the residue field. Form the finite commutative algebra
\[
 A'=C\oplus N,\qquad (b,v)(b',w)=(bb'+\beta(v,w),bw+b'v).
\]
Read \(x_i z x_j^*\) as a scalar in each diagonal block, as an element of \(N\) in \(pEp'\), and by its adjoint in \(p'Ep\). Insertion of \(\sum_jx_j^*x_j=1\) gives the matrix multiplication law. Denote the group matrices by \(r_u\). For one block take \(N=0\), \(A'=C\). Put \(I_N=(\beta(v,w):v,w\in N)\) and \(I_2=(\varpi)+\mathfrak m_C^2+I_N\).
\leanref{UpperTheory.\allowbreak{}V.\allowbreak{}Pa.\allowbreak{}idem}{Development.\allowbreak{}AutomorphicArithmetic.\allowbreak{}AutomorphicVarieties.\allowbreak{}PatchingModules.\allowbreak{}ClosureConjugation.\allowbreak{}OBLt}{43--66}
\leanref{UpperTheory.\allowbreak{}V.\allowbreak{}Pa.\allowbreak{}gr\_\allowbreak{}E}{Development.\allowbreak{}AutomorphicArithmetic.\allowbreak{}AutomorphicVarieties.\allowbreak{}PatchingModules.\allowbreak{}ClosureConjugation.\allowbreak{}ComplOddPart}{48--122}

Let \(T=(\mathfrak m_C/I_2)^\vee\). An \(O_*\)-derivation \(\alpha\in T\) gives an ordinary even cocycle \(A_\alpha\) by applying \(b+v\mapsto\bar b+\delta\alpha(b)\), \(\delta^2=0\), to \(r_u\) and multiplying the derivative by \(\rho^{-1}\). If \(A_\alpha\) is a boundary, infinitesimal conjugation makes the matrices constant. The determinant coefficients are then constant. In the \(i\)-th diagonal entry of each ordered-product identity at \(v\mid\ell\), specialise the other slots at elements separating \(\psi_{v,j}\) from \(\psi_{v,i}\). All factors except the \(i\)-th are units, so \(\psi_{v,i}\) is constant. Simple-root uniqueness fixes the auxiliary roots, and the projected tame relations fix the \(y_{ji}\). Since \(s\) belongs to \(O_*\), these generators force \(\alpha=0\).

Let \(U\subset T^g\) consist of tuples \((\alpha_j)\) for which a bounded ordinary \(V_2(F)\)-cocycle \(\Psi\) has \([e'_j]\Psi=A_{\alpha_j}\) for every \(j\). The lifting bound on the limit gives \(\operatorname{codim}U\le g\). For a basis element \(\alpha\) of \(W=\bigcap_j\iota_j^{-1}U\), take witnesses \(\Psi_j\) for \(A_\alpha\) in slot \(j\). Then
\[
 y_\alpha=A_\alpha+\sum_j e_j[\eta]\Psi_j
\]
is an ordinary \(V_1(F)\)-cocycle. Its constant component and \(e'_jy_\alpha=\Psi_j\) give the cocycle equation in every degree. For the local condition choose a witness \(Y_v\) for \(A_\alpha\). The primed witness in slot \(b\) for \(\Psi_j\) differs from \(\delta_{jb}Y_v\) by an invariant modulo the upper triangular matrices. The nontrivial graded ratios force its lower entries to vanish successively. Subtract the upper triangular corrections in the lifted flag basis. The \(\eta\)-components of the resulting witnesses give the local \(e_j\)-components for \(y_\alpha\).

The classes of the \(y_\alpha\) and the \(g^2\) classes \(e_j\kappa_i\) are independent together. The products \(e_je_b=0\) give the cocycle and flag equations for \(e_j\kappa_i\). In a relation equal to \(d_{\rho_1}Y\), the constant part and injectivity for \(A_\alpha\) kill the coefficients of the \(y_\alpha\). The constant term of \(Y\) then commutes with the full blocks and is skew-transpose, so it is zero. The \(e_j\)-components of the relation now give linear combinations of the \(\kappa_i\) that are ordinary boundaries, and all their coefficients vanish. The limit bound gives \(\dim W+g^2\le D_*\). Apply \ref{B6.48ab}, and lift its spanning forms \(X_b\) to \(x_b\in\mathfrak m_C\).
\leanref{UpperTheory.\allowbreak{}V.\allowbreak{}Pa.\allowbreak{}lift1}{Development.\allowbreak{}AutomorphicArithmetic.\allowbreak{}AutomorphicVarieties.\allowbreak{}PatchingModules.\allowbreak{}SelmerScaling.\allowbreak{}FirstOrderMatrixMemSquareZero}{279--389}
\leanref{UpperTheory.\allowbreak{}V.\allowbreak{}Pa.\allowbreak{}claim1}{Development.\allowbreak{}AutomorphicArithmetic.\allowbreak{}AutomorphicVarieties.\allowbreak{}PatchingModules.\allowbreak{}SelmerScaling.\allowbreak{}FirstOrderLiftConstantPart}{134--245}

For \(\mu\in(N/\mathfrak m_CN)^\vee\), differentiating the odd entries of \(r_u\) and multiplying by \(\rho^{-1}\) gives an odd cocycle \(K_\mu\). An inner derivative vanishes on the lifted idempotents, since their represented matrices are the constant \(E_{ii}\). An odd bounding matrix commuting with all \(E_{ii}\) is zero. The group algebra spans the distinguished corner, so \(\mu\mapsto[K_\mu]\) is injective. Write its coordinates on \([\kappa_j]\) as \(\mu(m_j)\) by duality. Lifts of the \(m_j\) generate \(N\) by Nakayama's lemma.

For an \(F\)-linear functional \(\lambda:N/I_2N\to F\), write the bilinear difference \(\lambda(bv)-\bar b\lambda(v)\) on the two cotangent spaces as \(\sum_h\alpha_h(b)\mu_h(v)\). With paired slots indexed by \(h\) the map
\[
 b+v\longmapsto
 \bar b+\sum_h e_h\mu_h(v)+\sum_h e'_h\alpha_h(b)+\eta\lambda(v)
\]
is multiplicative. Its mixed products give the bilinear difference and the \(\alpha_h\) kill \(I_N\). Apply the map to \(r_u\) and use first-order odd conjugation to replace each \(K_{\mu_h}\) by \(\sum_j d_{hj}\kappa_j\), where \(d_{hj}=\mu_h(m_j)\). Use \(\kappa_j\) in the unprimed slots and \(\sum_h d_{hj} A_{\alpha_h}\) in the primed slots and keep the constant and \(\eta\)-parts. These changes are transposes for the mixed pairing, hence preserve the product identities, polarisation and inertia relations. Denote the resulting representation by \(r'\), put \(I=(e'_1,\ldots,e'_g,\eta)\) in the slot algebra, and set \(Z(s)=r'(s)\rho_1(s)^{-1}-1\). The constant and unprimed coordinates agree with those of \(\rho_1\), so \(Z\) has entries in \(I\). The block parity leaves primed entries on the diagonal blocks and \(\eta\)-entries on the off-diagonal blocks, precisely \(V_2(F)\). Since \(I^2=0\), multiplication gives \(Z(st)=Z(s)+\rho_1(s)Z(t)\rho_1(s)^{-1}\). The common polarisation identity gives negative transpose under conjugation, using \((1+Z)^{-1}=1-Z\). On the prescribed tame inertia kernels both representations are one.

We check the denominators before passing to cohomology. Write \(A=k[[t]]\), \(F=k((t))\), and let \(R_0\) be the integral coefficient ring before localization and completion at the curve prime. Its curve map to \(A\) is onto, with the power series in the curve coordinate lifted from the coefficient ring. Under the slot map their reductions are those series in \(t\). The kernel of the curve map has finitely many generators since \(R_0\) is Noetherian. Their images lie in the augmentation ideal \(J'\) of the slot algebra, where \((J')^3=0\). Successively expanding modulo \(J',(J')^2,(J')^3\) shows that the image of \(R_0\) lies in the finite \(A\)-algebra generated by these nilpotent images: subtract a lift of the curve value and write the remainder in the kernel generators at each step. Take entries of a finite set of generators of the integral Cayley--Hamilton matrix module. Their images span, over this finite algebra, a finite \(A\)-module containing the slot images of every group matrix. Each coordinate of this module in \(F\) has a common denominator \(t^b\). Moreover the maximal ideal of \(R_0\) maps into \((t)+J'\) in the finite algebra, so its \((m+3)\)-rd power maps into \(t^m\) times that algebra. Adic continuity of the integral group matrices therefore makes all the cleared coordinates continuous modulo every \(t^m\). Conjugation by the fixed change of basis, finite slot sums, multiplication, and substitution \(s\mapsto s^{-1}\) preserve this bounded continuity. The same holds for \(\rho_1(s^{-1})\), whose entries are formed from the integral base matrices and the finitely many chosen bounded cocycles \(\kappa_j\). It follows for \(Z\).

On each ordinary decomposition group apply Lemma~\ref{B6.48aad} with \(r_0=\rho_1\), \(r_1=r'\), and this \(I\). The local characters agree modulo \(I\), their augmentations are pairwise distinct by the curve hypothesis, and the transformed ordered products vanish. Use the subalgebra with the displayed block parities. Its intersection with \(M_n(I)\) is \(V_2(F)\), so the lemma gives a single local witness \(Y_v\in V_2(F)\) for all primed coordinates. In the flag basis \(Z-dY_v\) is upper triangular. Its diagonal vanishes on the prescribed preimage of the local abelian torsion because both systems of flag characters take value one there. This proves the ordinary local condition. Finally \([e'_j]Z=\sum_h d_{hj}A_{\alpha_h}=A_{a_j}\), where the bilinear difference evaluated at \(m_j\) gives \(a_j(b)=\lambda(bm_j)-\bar b\lambda(m_j)\). This bounded ordinary \(Z\) witnesses \((a_j)\in U\), and \ref{B6.48ab} gives
\[
 \lambda\!\left(\sum_{j,b}\widetilde c_{ij,b}x_bm_j\right)
 =\sum_{j,b}c_{ij,b}X_b(a_j)=0
\]
for lifts \(\widetilde c_{ij,b}\in O_*\). Since the functionals separate \(N/I_2N\), the row sums belong to \(I_2N\).
\leanref{UpperTheory.\allowbreak{}V.\allowbreak{}Pa.\allowbreak{}dualEx}{Development.\allowbreak{}AutomorphicArithmetic.\allowbreak{}AutomorphicVarieties.\allowbreak{}PatchingModules.\allowbreak{}SelmerScaling.\allowbreak{}DualEx}{94--135}
\leanref{UpperTheory.\allowbreak{}V.\allowbreak{}Pa.\allowbreak{}claim2}{Development.\allowbreak{}AutomorphicArithmetic.\allowbreak{}AutomorphicVarieties.\allowbreak{}PatchingModules.\allowbreak{}SelmerScaling.\allowbreak{}Claim2}{11--34}

Map \(P=O_*[[X_1,\ldots,X_{D_*},u_1,\ldots,u_g]]\) to \(A'\) by \(X_b\mapsto x_b\), \(u_j\mapsto m_j\). The images span the relative cotangent space because \(I_N\) maps into \(\mathfrak m_{A'}^2\), so completeness gives surjectivity. Express row \(i\) as \(\sum_j w_{ij}m_j\) with \(w_{ij}\in I_2\), and lift \(w_{ij}\) to \(v_{ij}\in(\varpi)+\mathfrak m_P^2\). Put \(z_{ij}=\sum_b\widetilde c_{ij,b}X_b-v_{ij}\) and \(Z=(z_{ij})\). Then \(Zu=0\) in \(A'\). The \(z_{ij}\) give independent relative cotangent vectors by \ref{B6.48ab}, independent also of the \(u_j\). For one block this is a surjection from \(O_*[[X_1,\ldots,X_{D_*}]]\) to \(C\).
\leanref{UpperTheory.\allowbreak{}V.\allowbreak{}Pa.\allowbreak{}pres}{Development.\allowbreak{}AutomorphicArithmetic.\allowbreak{}AutomorphicVarieties.\allowbreak{}PatchingModules.\allowbreak{}SelmerScaling.\allowbreak{}AlCSmul}{12--148}

Apply Lemma~\ref{B6.48aae} to the auxiliary tower, the ultrafilter, \(R\), \(\Lambda\), \(S\), \(\mathfrak a\), \(Y=(y_{ji})\), \(f=\prod_j\prod_{i\ne b}(U_{ji}-U_{jb})^2\), and \(X=B^{\rm sp}\otimes_{B_0[[s]]}(M_0\otimes_O O[[s]])\). On each curve the image of \(f\) is nonzero by separability at \(s_j\). Lemma~\ref{B6.48aae}(3) gives the finite \(\mathcal B\)-module \(\mathcal M=\varprojlim_b L_b[1/f]\) over \(\mathcal B=\widehat{R[1/f]}^{\,Y}\) with \(\mathcal M/Y\mathcal M\simeq\mathcal B\otimes_R X\), \(\operatorname{gr}_Y\mathcal M\simeq(\mathcal M/Y\mathcal M)[y_{ji}]\), and \(\mathfrak a\mathcal B\subseteq\sqrt{\operatorname{Ann}_{\mathcal B}\mathcal M}\).

Complete \(\mathcal B\) locally at the prime over \(p_i\) and write \(\widehat M_i\) for the base change of \(\mathcal M\). The ring before quotienting by \(\mathfrak a\) is the completion of \(R_{p_i}\), since \(Y\subset p_i\) and \(f\notin p_i\). Let \(S_i\) be the image of \(\operatorname{Ann}(\widehat M_i)\) in \(C_i\). Dividing by \(\mathfrak a\) preserves the support dimensions. The weight prime of the curve has completed local ring regular of dimension \(1+d\), since \(\Lambda\) has dimension \(2+d\) and the curve quotient is \(k[[t]]\). The augmentation of \(\widehat M_i\) is finite free over that ring. Indeed, \(M_0\) is free over weight space, the root adjunction is finite free over the enlarged base, and completion at the chosen prime selects a direct summand of the module completed at the weight prime. Lemma~\ref{B6.48aac} identifies each completed graded piece by its multiplication map. Thus the \(nq\) auxiliary parameters form a regular sequence on the nonzero completed module. A system of weight parameters continues it to a regular sequence with finite-length quotient. Apply \ref{B6.48ac} to the completed weight ring, the completed module and its \(nq\) auxiliary parameters. The augmentation is finite free as above, and flat base change preserves the displayed graded-piece isomorphism. Here \(d_0=1+d\) and \(h=nq\), so \(d_0+h=1+D_*\). Since \(\mathfrak a\) lies in the radical of the annihilator, quotienting the ambient ring by \(\mathfrak a\) identifies the prime quotients at the minimal support primes. We obtain
\begin{equation}\tag{6.48a.3}\label{eq:p6c5-patch-support}
 \dim C_i/Z_0=1+D_*\qquad(Z_0\text{ minimal over }S_i).
\end{equation}

Let \(Q\) be a prime of \(C_i\), with preimage \(\widetilde Q\) in the completion before imposing \(\mathfrak a\). Suppose \(Y\subseteq\widetilde Q\) and \(\widetilde Q\cap R=\pi^{-1}(P_Q)\) for a prime \(P_Q\) of \(B^{\rm sp}\). Then
\begin{equation}\tag{6.48a.5}\label{eq:p6c5-patch-fibre}
 \begin{split}
 S_i\subseteq Q
 &\ \Longleftrightarrow\ 
 \widehat M_i\otimes\kappa(\widetilde Q)\ne0\\
 &\ \Longleftrightarrow\ X\otimes_R\kappa(\widetilde Q)\ne0
 \ \Longleftrightarrow\ J_0\subseteq P_Q\cap B_0 .
 \end{split}
\end{equation}
The first equivalence is Nakayama's lemma, the second uses augmentation with \(f\) invertible, and the third uses extension of the residue field of the original ordinary fibre.
\leanref{UpperTheory.\allowbreak{}V.\allowbreak{}Pa.\allowbreak{}Core.\allowbreak{}pmFib}{Development.\allowbreak{}AutomorphicArithmetic.\allowbreak{}AutomorphicVarieties.\allowbreak{}PatchingModules.\allowbreak{}AnnihilatorBaseChange.\allowbreak{}PmFib}{7--32}
\leanref{UpperTheory.\allowbreak{}V.\allowbreak{}Pa.\allowbreak{}Core.\allowbreak{}pmDim}{Development.\allowbreak{}AutomorphicArithmetic.\allowbreak{}AutomorphicVarieties.\allowbreak{}PatchingModules.\allowbreak{}AnnihilatorBaseChange.\allowbreak{}PmDimQuot}{99--150}
\leanref{UpperTheory.\allowbreak{}V.\allowbreak{}Pa.\allowbreak{}s5}{Development.\allowbreak{}AutomorphicArithmetic.\allowbreak{}AutomorphicVarieties.\allowbreak{}PatchingModules.\allowbreak{}AutomorphicDetection.\allowbreak{}AnnBc}{69--86}
\leanref{UpperTheory.\allowbreak{}V.\allowbreak{}Pa.\allowbreak{}core}{Development.\allowbreak{}AutomorphicArithmetic.\allowbreak{}AutomorphicVarieties.\allowbreak{}PatchingModules.\allowbreak{}AutomorphicDetection.\allowbreak{}LiftPt}{76--176}

In (a), \eqref{eq:p6c5-patch-fibre} puts the curve prime in the support. For a minimal support prime \(Z_0\) below it the surjection \(O_*[[X_1,\ldots,X_{D_*}]]\to C_i/Z_0\) has zero kernel by \eqref{eq:p6c5-patch-support}, since a nonzero ideal decreases dimension. Thus \(Z_0=0\). Lift a prime below the curve kernel to the splitting algebra by going-down and to the curve completion by faithful flatness. The lift contains \(S_i\subseteq Z_0\), so \eqref{eq:p6c5-patch-fibre} gives (a).

For (b), stability is detected by a nonzero minor \(\det(\operatorname{tr}_D(g_i h_j))_{i,j\le n^2}\). Indeed full span gives a nonsingular trace pairing, and a nonzero minor forces full span in the semisimple representation. Thus stability persists under generalisation. At a stable prime \(Q\subset C\), \(I_N\not\subseteq Q\), since killing the odd products gives a direct sum determinant. At a prime with nonsemisimple specialised group algebra one has \(\operatorname{Ann}_C N\subseteq Q\). Otherwise the odd corners vanish and the algebra is a quotient of the product of full matrix algebras. A lift of a thick base prime at augmentation has nonsemisimple specialised algebra. The projection to the base group sends \(\tau_j\) to one and \(\phi_j\) to \(s_j\). The auxiliary relations vanish by Cayley--Hamilton and \(y=0\). The closed ordinary and inertia relations map into the base closed ideal by the finite quotients defining closure. After a common residue field extension the specialised algebra therefore surjects onto the nonsemisimple base algebra.
\leanref{UpperTheory.\allowbreak{}V.\allowbreak{}Pa.\allowbreak{}gr\_\allowbreak{}thick}{Development.\allowbreak{}AutomorphicArithmetic.\allowbreak{}AutomorphicVarieties.\allowbreak{}PatchingModules.\allowbreak{}ClosureConjugation.\allowbreak{}GrThick}{77--86}
\leanref{UpperTheory.\allowbreak{}V.\allowbreak{}Pa.\allowbreak{}Core.\allowbreak{}nonssT}{Development.\allowbreak{}AutomorphicArithmetic.\allowbreak{}AutomorphicVarieties.\allowbreak{}PatchingModules.\allowbreak{}AnnihilatorBaseChange.\allowbreak{}NonssT}{69--100}

Let \(Z_0\) be a minimal support prime with stable contraction and \(Q\) a target of either kind. Write \(\zeta\) for the image of \(Z\) in \(C\). The adjugate gives \((\det\zeta)N=0\). Hence \(\det\zeta\in Q\) in the nonsemisimple case by the annihilator condition and in the stable case by \((\det\zeta)I_N=0\). Lift \(Z_0,Q\) to \(A'\) by integrality, then pull back to \(Z',Q'\subset P\). Both contain \((Zu)\). Some \(u_j\notin Z'\), since otherwise \(I_N\subseteq Z_0\). In particular \(g>0\), and the adjugate gives \(\det Z\in Z'\). The target condition gives \(\det Z\in Q'\).

The entries of \(Z,u\) have independent cotangent classes in the regular local presentation. Lemma~\ref{B6.48ad} therefore gives a least prime over \((Zu,\det Z)\). Adjoin a kernel line with chart \(v_c=1\) and equations \(Zv=0,\ u_b=u_cv_b\) for \(b\ne c\). Over the closed point of \(P\) the direction coordinates form an affine space. At each prime of that fibre the ambient ring is regular. The row equations have independent linear parts in column \(c\), and the vector equations have independent parts in the \(u_b\), \(b\ne c\). This follows from relative cotangent independence and \(v_c=1\). The local quotient is regular and is a domain. Its generic residual direction gives a contraction to \(P\) common to the charts, by changing coordinates \(v\mapsto v/v_b\) on their overlap. At each specialised residual direction the regular local quotient has this same generic contraction. Given any prime over \((Zu,\det Z)\), dominate its quotient by a valuation ring. Over its fraction field choose a kernel vector, using \(u\) itself if nonzero. Divide by a coordinate of least valuation to define a point in one chart. Its closed centre is in the residual direction fibre, so the common generic contraction is contained in the given prime. This proves the least-prime assertion.

Adjoin \(Z=\operatorname{diag}(u)\) to the row equations. Then \(u_j^2=0\) for all \(j\), so the radical of the enlarged ideal is generated by \(g^2+g\) independent parameters. The height theorem gives
\begin{equation}\tag{6.48a.4}\label{eq:p6c5-row-dimension}
 \dim P/(Zu)\le\dim P-g .
\end{equation}
By integrality and \eqref{eq:p6c5-patch-support}, \(\dim P/Z'=\dim C/Z_0=\dim P-g\). The least prime over \((Zu,\det Z)\) is contained in \(Z'\). Strict inclusion would make its quotient dimension larger than \(\dim P-g\), contrary to \eqref{eq:p6c5-row-dimension}. Therefore it equals \(Z'\) and is contained in \(Q'\). This gives \(Z_0\subseteq Q\).

Lift the starting stable prime to the curve completion. It is in the support by \eqref{eq:p6c5-patch-fibre}, and a minimal support prime beneath it has stable contraction by the trace minor. At a shared thick prime lift the common fibre prime to \(Q\) in the current completion. The containment just proved puts \(Z_0\) below \(Q\). Before imposing \(\mathfrak a\), contract the preimage of \(Z_0\) to \(z\in\operatorname{Spec}\mathcal B\). The contraction of \(\widetilde Q\) contains \(Y\), and in \(\mathcal B/Y=(R/Y)[1/f]\) it lies under the next curve prime. Hence \(z\) lies under that curve prime as well. Faithful flatness gives a prime above \(z\) in the next completion. Both completed fibres over \(z\) are field extensions of the same nonzero fibre of \(\mathcal M\), so the new prime is in the support. It has the same stable contraction to \(R/\mathfrak a\). Iterate the construction and at the last curve use \(Z_0\subseteq Q\) for the target prime, followed by \eqref{eq:p6c5-patch-fibre}. For \(h=0\) this last argument applies at the starting curve.
\leanref{UpperTheory.\allowbreak{}V.\allowbreak{}Pa.\allowbreak{}Core.\allowbreak{}pmMove}{Development.\allowbreak{}AutomorphicArithmetic.\allowbreak{}AutomorphicVarieties.\allowbreak{}PatchingModules.\allowbreak{}AnnihilatorBaseChange.\allowbreak{}PmMove}{8--24}
\leanref{UpperTheory.\allowbreak{}V.\allowbreak{}Pa.\allowbreak{}presOut0}{Development.\allowbreak{}AutomorphicArithmetic.\allowbreak{}AutomorphicVarieties.\allowbreak{}PatchingModules.\allowbreak{}SelmerScaling.\allowbreak{}PresOut0}{9--69}
\leanref{UpperTheory.\allowbreak{}V.\allowbreak{}Pa.\allowbreak{}presOut1}{Development.\allowbreak{}AutomorphicArithmetic.\allowbreak{}AutomorphicVarieties.\allowbreak{}PatchingModules.\allowbreak{}SelmerScaling.\allowbreak{}CotangentIndependentOfResidues}{74--212}
\leanref{UpperTheory.\allowbreak{}V.\allowbreak{}AlgebraTools.\allowbreak{}under\_\allowbreak{}le\_\allowbreak{}of\_\allowbreak{}det\_\allowbreak{}mem}{Development.\allowbreak{}AutomorphicArithmetic.\allowbreak{}AutomorphicVarieties.\allowbreak{}AlgebraTools.\allowbreak{}CaDetCore}{9--49}
\leanref{UpperTheory.\allowbreak{}V.\allowbreak{}Pa.\allowbreak{}Core.\allowbreak{}chainAbs}{Development.\allowbreak{}AutomorphicArithmetic.\allowbreak{}AutomorphicVarieties.\allowbreak{}PatchingModules.\allowbreak{}AnnihilatorBaseChange.\allowbreak{}ChainAbs}{7--71}
\end{proof}

\subsection*{Proof of \ref{B6.48aaa}}\label{prf:B6.48aaa}\phantomsection
\begin{proof}
First work over \(F\), with the polarising matrix made the identity as in \ref{B6.48aa}. Put \(\chi=\bar\varepsilon_\ell\). The subgroup on which we test is \(H'=H\cap\ker\varepsilon_\ell\), using the full cyclotomic character. In the coefficient algebra of that proof the perfect dual pairing is the \(\eta\)-coefficient of matrix trace. Over \(A_m\), follow it by the \(t^{m-1}\)-coefficient and \(\operatorname{Tr}_{k/\F_\ell}\) to get the finite-coefficient pairing. Opposite matrix units and complementary powers of \(t\) show that it is perfect. The dual action has the additional factor \(\chi\).

Here are the details of detection. Write \(r=\rho\), \(r_1=(1+\sum e_j\kappa_j)r\), and write a twisted dual cocycle as \(z=\eta p+\sum e'_jl_j\). On \(F[H']\) extend the functions
\[
 R(s)=r(s),\qquad K_j(s)=\kappa_j(s)r(s),\qquad L_j(s)=l_j(s)r(s)
\]
linearly. The map \(R\) has image the product \(D\) of the full diagonal block algebras. This follows from absolute irreducibility on \(H'\) and, for two blocks, inequivalence of the two simple modules. Put \(I=\ker R\). The \(K_j,L_j\) satisfy the derivation rule with left and right action through \(R\). Extend also \(T(s)=[\eta](z(s)r_1(s))\) linearly. Multiplying the coefficient matrices and using \(e_ie'_j=\delta_{ij}\eta\) gives
\[
 T(uv)=T(u)R(v)+R(u)T(v)
       +\sum_j\bigl(K_j(u)L_j(v)+L_j(u)K_j(v)\bigr).
\]
If all trace tests vanish, \(\operatorname{tr}T\) vanishes on the group algebra. For \(u,v\in I\), insert between them a lift of \(E_{jj}\) and on the outside a lift of \(E_{ii}\), where \(i,j\) belong to opposite blocks. Terms differentiating either lift vanish because they contain \(R(u)\) or \(R(v)\). Thus, using \(b\) for the cocycle index,
\[
 \sum_b K_b(u)_{ij} L_b(v)_{ji}
       +L_b(u)_{ij}K_b(v)_{ji}=0.
\]
Fix \(v\). Replacing \(u\) by \(a u a'\) for lifts of matrix units makes the same linear combination zero on the entire corresponding corner of \(I\). The combination of corner derivations therefore factors through \(D\). A derivation of a product of two matrix algebras into the off-diagonal bimodule is inner: applying the product rule to the central idempotent of the target block and to its products with the two summands expresses it as commutation with its value on that idempotent.

In ordinary \(H^1(H',\operatorname{Hom}(r_2,r_1))\), the restrictions of the \(K\)- and \(L\)-cocycles in this combination transform under \(H\) by the trivial character and by \(\chi^{-1}\), respectively. This follows by applying the cocycle identity to \(hsh^{-1}\) and discarding its coboundary term. Some \(h\in H\) has \(\chi(h)\ne1\), since \(\mu_\ell\not\subset L\). Each of the two combinations is consequently inner on \(H'\). A bounding matrix in the chosen corner works on \(H\): after subtracting its boundary the cocycle vanishes on \(H'\), and its remaining values are \(H'\)-intertwiners, which are zero. Polarisation gives the opposite corner by transposition. In particular a combination of the \(\kappa_b\) that is inner in this corner is a polarised boundary in both corners. Independence of their classes forces \(L_b(v)_{ji}=0\) for each \(b\). This holds for every \(v\in I\) and every opposite pair. Each \(L_b\) now factors through \(D\) on \(H'\); the same derivation and intertwiner arguments make \(l_b\) a twisted boundary on \(H\).
\leanref{UpperTheory.\allowbreak{}V.\allowbreak{}Pa.\allowbreak{}Det.\allowbreak{}relPair}{Development.\allowbreak{}AutomorphicArithmetic.\allowbreak{}AutomorphicVarieties.\allowbreak{}PatchingModules.\allowbreak{}DeterminantPadicFree.\allowbreak{}RelPair}{23--177}
\leanref{UpperTheory.\allowbreak{}V.\allowbreak{}Pa.\allowbreak{}Det.\allowbreak{}ann}{Development.\allowbreak{}AutomorphicArithmetic.\allowbreak{}AutomorphicVarieties.\allowbreak{}PatchingModules.\allowbreak{}DeterminantPadicFree.\allowbreak{}BlockPrinc}{75--161}

Subtract these boundaries in the paired coefficient algebra. This removes their mixed \(\eta\)-terms as well, since conjugation in the boundary is by \(r_1\). The remaining even derivative on \(H'\) satisfies the ordinary derivation rule. For \(u\in I\) its product with every \(R(v)\) has trace zero, by the test on \(uv\). The trace form on \(D\) is nondegenerate, so the derivative vanishes on \(I\). Every derivation of a full matrix algebra into itself is inner: for a derivation \(\delta\) the matrix \(\sum_i\delta(E_{i1})E_{1i}\) has commutator \(\delta(E_{ab})\) with \(E_{ab}\). Apply this in each block. The remaining cocycle factors through the abelian quotient \(H/H'\) on the block scalars with action \(\chi\). For a fixed \(h\) with \(\chi(h)\ne1\), commutativity gives
\((\chi(h)-1)z(s)=(\chi(s)-1)z(h)\), which supplies its bounding scalar. Averaging the bounding matrix with its polarised image gives a bounding matrix of the required transpose type because \(\ell\) is odd. For one block the proof starts with the even derivative. A power of \(t\) clears all denominators.
\leanref{UpperTheory.\allowbreak{}V.\allowbreak{}Pa.\allowbreak{}Det.\allowbreak{}detAlg}{Development.\allowbreak{}AutomorphicArithmetic.\allowbreak{}AutomorphicVarieties.\allowbreak{}PatchingModules.\allowbreak{}DeterminantPadicFree.\allowbreak{}DetAlg}{71--142}
\leanref{UpperTheory.\allowbreak{}V.\allowbreak{}Pa.\allowbreak{}Det.\allowbreak{}detAsm}{Development.\allowbreak{}AutomorphicArithmetic.\allowbreak{}AutomorphicVarieties.\allowbreak{}PatchingModules.\allowbreak{}DeterminantPadicFree.\allowbreak{}DetAsm}{160--269}

The common separable locus of the curves is open dense in \(H'\), being a finite intersection of open dense sets there. A continuous trace vanishing on it vanishes on \(H'\). The space on which the tests act modulo boundaries is finite-dimensional: values on a finite set of generators of the open pro-\(\ell\) residual group, together with values on coset representatives, determine a cocycle. The reduction to this group is valid for finite truncations since the semidirect product of the residual unipotent action and the additive cocycle module on the residual kernel is an \(\ell\)-group. The finite-character hypothesis in \ref{B6.56} gives finite generation of the pro-\(\ell\) group. Successively choose an evaluation nonzero on the current kernel until that kernel is zero. Take the union over the curves and split the finitely many polynomials by finite DVR extensions as in \ref{B6.48aa}.
\leanref{UpperTheory.\allowbreak{}V.\allowbreak{}Pa.\allowbreak{}choice}{Development.\allowbreak{}AutomorphicArithmetic.\allowbreak{}AutomorphicVarieties.\allowbreak{}PatchingModules.\allowbreak{}AutomorphicDetection.\allowbreak{}Choice}{14--41}

We spell out the integral estimate used below. In generator coordinates, cocycle relations, local conditions with their bounding matrices, and the selected evaluations are linear equations over \(A\). One may take all relations: their rows generate a submodule of a finite free \(A\)-module, hence a finite list of rows suffices. Concretely the coordinates before the relations describe crossed homomorphisms on a free profinite cover. At every depth the crossed homomorphism together with the representation is a homomorphism to a finite semidirect product. Evaluation on a word of the cover gives compatible linear forms at all depths and hence a row over \(A\). Truncation thus imposes exactly the reductions of the rows. A nonzero Smith divisor \(t^b\) introduces a factor of at most \(Q^b\) in the ratio of the truncated kernel to the reductions of integral solutions. Quotienting by boundaries has the same bounded-error property: use the map from bounding matrices to generator values, and project away any local witness variables. Equivalently, take successively the kernel, the projection of the integral kernel, and the image of the boundary map; the indices of these last two lattices in their saturations are fixed. A rational space of classes of dimension \(u\) therefore contributes at most \(Q^{mu+c}\), with \(c\) independent of \(m\).

Apply this observation to the base twisted dual satisfying the trace tests. Its quotient by boundaries has bounded order, by detection. At level \(N\) the dual of the full auxiliary local condition is zero in local \(H^1\). A representative there is a boundary at both the tame generator and Frobenius. The tame generator acts trivially on the pulled-back module, so the cocycle vanishes on it. Its semidirect-product homomorphism consequently descends to the base by the normal-generation assertion of \ref{B6.56}, after the pro-\(\ell\) reduction above. At fixed \(m\) there are finitely many base cocycles. For all sufficiently high \(N\) their evaluations at the images of the auxiliary Frobenii agree simultaneously with their evaluations at the selected \(s_j\). At \(s_j\) the cyclotomic twist is trivial, and a boundary multiplied by \(r_1(s_j)\) has trace zero by cyclic invariance. This gives the base tests. Subtraction of a boundary normalises the value at complex conjugation to zero, by averaging over its group of order two. Descent then gives an injection on the dual classes in question. Denote the resulting uniform length bound by \(c_2\).
\leanref{UpperTheory.\allowbreak{}V.\allowbreak{}Pa.\allowbreak{}dualN}{Development.\allowbreak{}AutomorphicArithmetic.\allowbreak{}AutomorphicVarieties.\allowbreak{}PatchingModules.\allowbreak{}EpsilonEtaMatrices.\allowbreak{}DualN}{10--155}

Write \(\lambda=\log_Q\#\) for the following counts and put
\(\nu=\binom a2+\binom{n-a}2+ga(n-a)\). At \(v\mid\ell\), first count cocycles already upper triangular in the lifted flag. A diagonal entry is additive and vanishes on the prescribed torsion. The local abelianisation gives at most \(m(1+d'_v)\) for its length. Filter the upper entries by distance from the diagonal. After smaller distances are fixed, differences in position \(ij\) are cocycles for \(\psi_i/\psi_j\). For this one-dimensional module, the local Euler characteristic and Tate duality in \ref{B6.05d} give
\[
 \lambda Z^1(G_v,A_m(\psi_i/\psi_j))
 =m(1+d_v)+\lambda H^0(G_v,A_m(\chi\psi_j/\psi_i)).
\]
The last term is bounded by the valuation of one nonzero character difference, since \(\psi_i\ne\chi\psi_j\) for \(i<j\). There is one slot for each within-block upper entry and \(g\) slots for each between-block upper entry. This proves the first estimate in \eqref{eq:p6c5-patch-local}. Allowing a change of flag adds at most the index of the upper triangular matrices in \(V_1(A_m)\), since a boundary from an upper triangular matrix is already among the cocycles counted. In the Greenberg--Wiles local factor this leaves the subtraction \(m(n+\nu)\) at each ordinary place: the upper triangular matrices have length \(m(n+\nu)\), and the local factor for cocycles is their count divided by the order of the full coefficient module.

At an auxiliary place the full local factor has length \(\lambda H^0(G_w,V_1^\vee(1))\), by the Euler characteristic away from \(\ell\). At the chosen Frobenius, over \(F\) an invariant has all its odd primed entries zero: each would intertwine blocks with disjoint spectra. Its remaining \(\eta\)-matrix commutes with the simple-spectrum matrix \(r(s_j)\), so has dimension at most \(n\). The Smith estimate gives \(mn+c_j\). At the distinguished place the hypothesis on twisted intertwiners kills successively the primed and the \(\eta\)-parts over \(F\), giving a bound \(c_0\). At a tame place with unramified condition the local factor is one, since a class is determined by its Frobenius value modulo Frobenius boundaries. Globally a skew invariant for \(V_1\) has zero constant part: this part is block scalar and skew. Its odd coefficients are then intertwiners and vanish. The truncated invariants thus have bounded length \(c_1\). Each real place subtracts \(m\nu\), the length of the skew part for negative transposition, since each within-block unordered pair contributes one slot and each cross pair contributes \(g\).
\leanref{UpperTheory.\allowbreak{}V.\allowbreak{}Pa.\allowbreak{}locV}{Development.\allowbreak{}AutomorphicArithmetic.\allowbreak{}AutomorphicVarieties.\allowbreak{}PatchingModules.\allowbreak{}EpsilonEtaMatrices.\allowbreak{}EntryZero}{204--323}
\leanref{UpperTheory.\allowbreak{}V.\allowbreak{}Pa.\allowbreak{}h0q}{Development.\allowbreak{}AutomorphicArithmetic.\allowbreak{}AutomorphicVarieties.\allowbreak{}PatchingModules.\allowbreak{}EpsilonEtaMatrices.\allowbreak{}FinrankCentralizerW1Le}{106--132}

All modules here are finite \(\ell\)-primary modules, the level set includes the places above \(\ell\), and the pairing at the start identifies the Tate dual used in \ref{B6.05d} and \ref{B6.09}. The Greenberg--Wiles formula, with the dual Selmer factor bounded by \(c_2\), now gives
\[
 \lambda W_m\le c_1+c_2+c_0+\sum_j(mn+c_j)-m\nu d_+
 +\sum_{v\mid\ell}\{m[n(1+d'_v)+\nu(1+d_v)] +c_v-m(n+\nu)\}.
\]
Both local-degree sums are \(d_+\), which proves the first bound.
\leanref{UpperTheory.\allowbreak{}V.\allowbreak{}Pa.\allowbreak{}oneLevel}{Development.\allowbreak{}AutomorphicArithmetic.\allowbreak{}AutomorphicVarieties.\allowbreak{}PatchingModules.\allowbreak{}EpsilonEtaMatrices.\allowbreak{}OneLevel}{236--484}

For the second bound let \(M_e\) be the odd module at depth \(e\). At the ordinary and prescribed tame places use the quotients \(Q_{v,e}\) and the subgroups \(J_v\) in the definition of \(E_e\). Write
\[
 T^1_e=\bigoplus_v H^1(G_v,Q_{v,e})/
           \ker\bigl(H^1(G_v,Q_{v,e})\to H^1(J_v,Q_{v,e})\bigr).
\]
Let \(S_e\) be the kernel of \(H^1(G,M_e)\to T^1_e\), and \(D_e\) the cokernel of \(H^2(G,M_e)\to\bigoplus_v H^2(G_v,Q_{v,e})\), where \(G\) is the global level group. The map from \(E_e\) forgetting the local cochains takes values in the kernel of this degree-two localisation. Its kernel is a quotient of \(T^1_e/\operatorname{im}H^1(G,M_e)\): subtract a bounding global cochain from a representative whose global class vanishes, and read the remaining local cocycles. Hence
\[
 \#E_e\le
 \#S_e\,\#D_e\,
 \frac{\#H^2(G,M_e)}{\#H^1(G,M_e)}
 \frac{\#T^1_e}{\prod_v\#H^2(G_v,Q_{v,e})}.
\]
Put \(u=a(n-a)\). The odd module has rank \(2u\), the real-place invariants have rank \(u\), and the ordinary lower quotient has rank \(u\). The global Euler characteristic bounds the first ratio by \(Q^{-eu d_+}\). At an ordinary place the local Euler characteristic bounds its contribution to the second ratio by \(Q^{eu d_v}\#H^0(G_v,Q_{v,e})\). At a tame place the unramified subgroup has order \(\#H^0\), so its contribution is one. These degree factors cancel. By \ref{B6.09}, \(\#D_e\) is bounded by the order of the global invariants of \(M_e^\vee(1)\): first localise to all the level places, then project to the indicated quotients. Local \(H^2\) surjects onto \(H^2\) of each quotient by local cohomological dimension two. The rational odd dual has no \(H'\)-invariants, by block inequivalence, giving a uniform bound \(Q^{c_d}\). Each ordinary lower quotient has a filtration by nontrivial character ratios. It has no rational invariants, so \(\#H^0(G_v,Q_{v,e})\le Q^{c_h}\) for a common constant.
\leanref{UpperTheory.\allowbreak{}V.\allowbreak{}Pa.\allowbreak{}B2.\allowbreak{}odCount}{Development.\allowbreak{}AutomorphicArithmetic.\allowbreak{}AutomorphicVarieties.\allowbreak{}PatchingModules.\allowbreak{}ArchimedeanH0Cards.\allowbreak{}OdCount}{8--80}

At base level the generator-and-witness count above gives \(\#S_e\le Q^{ge+c_b}\). At auxiliary level evaluate an odd cocycle on the tame generators. These evaluations are invariant under addition of a boundary. Their simultaneous kernel descends to the base, with its old local conditions. Indeed the tame generators act trivially and normally generate the new kernel in the pro-\(\ell\) quotient. Their values commute with the respective Frobenii by the tame relation and \(\mathrm Nw\equiv1\bmod\ell\). The odd centralizer at \(s_j\) is zero over \(F\). The Smith estimate on its commutator therefore bounds the product of all the evaluation image sizes by \(Q^{c_n}\), uniformly for \(e\le m\) once the Frobenii agree to depth \(m\). This proves \(\lambda S_e\le ge+c_b+c_n\) at that level.
\leanref{UpperTheory.\allowbreak{}V.\allowbreak{}Pa.\allowbreak{}B2.\allowbreak{}sel0}{Development.\allowbreak{}AutomorphicArithmetic.\allowbreak{}AutomorphicVarieties.\allowbreak{}PatchingModules.\allowbreak{}ArchimedeanH0Cards.\allowbreak{}ImgCard}{340--365}
\leanref{UpperTheory.\allowbreak{}V.\allowbreak{}Pa.\allowbreak{}B2.\allowbreak{}selN}{Development.\allowbreak{}AutomorphicArithmetic.\allowbreak{}AutomorphicVarieties.\allowbreak{}PatchingModules.\allowbreak{}ArchimedeanH0Cards.\allowbreak{}ExistsCore}{385--454}

Finally use both coefficient sequences with middle term \(A_m\), one with quotient \(A_e\), the other with quotient \(A_{m-e}\). Denote the induced maps \(E_{m-e}\to E_m\to E_e\) by \(i,p\), and the reduction \(E_m\to E_{m-e}\) by \(r\). The inclusion on coefficients for \(i\) is multiplication by \(t^e\), so \(ir=t^e\). We give the two index estimates. If a pair reduces to a boundary under \(p\), lift its bounding global cochain and its allowed local corrections, then subtract their boundary. At an ordinary place a local correction is a boundary and its bounding vector lifts. At a tame place inertia acts trivially and an unramified cocycle lifts by lifting its Frobenius value. The remaining pair has coefficients in \(t^e A_m\). Thus \(\ker p\subseteq\operatorname{im}i\). Moreover \(\#\operatorname{coker}r\le\#D_e\). To see this, lift the global two-class of a pair at depth \(m-e\) to depth \(m\), using global \(\ell\)-cohomological dimension two, and adjust by a cochain so its reduction is the given representative. Lift the local one-cochains. The differences of their differentials from the local images of the lifted global representative are local two-cocycles with coefficients in the kernel \(A_e\). Their tuple modulo global degree-two localisation is the defect in \(D_e\). A zero defect permits subtraction of a global kernel two-cocycle and correction by local kernel one-cochains. The pair then lifts. Changing the lifts changes the defect by a local boundary or the localisation of a global kernel class. The same construction on differences gives the asserted index bound. Consequently
\[
 \#(E_m/t^eE_m)\le \#E_e\,\#D_e
 \le Q^{ge+c_b+c_n+c_h\#\{v\mid\ell\}+2c_d}.
\]
All constants were chosen from the original curve, the fixed local groups and the finite detecting set.
\leanref{UpperTheory.\allowbreak{}V.\allowbreak{}Pa.\allowbreak{}B2.\allowbreak{}exactMid}{Development.\allowbreak{}AutomorphicArithmetic.\allowbreak{}AutomorphicVarieties.\allowbreak{}PatchingModules.\allowbreak{}ArchimedeanH0Cards.\allowbreak{}ExactMid}{100--116}
\leanref{UpperTheory.\allowbreak{}V.\allowbreak{}Pa.\allowbreak{}B2.\allowbreak{}cokerBound}{Development.\allowbreak{}AutomorphicArithmetic.\allowbreak{}AutomorphicVarieties.\allowbreak{}PatchingModules.\allowbreak{}ArchimedeanH0Cards.\allowbreak{}CokerBound}{151--169}
\end{proof}

\subsection*{Proof of \ref{B6.48aab}}\label{prf:B6.48aab}\phantomsection
\begin{proof}
Let \(\Phi\) be the finitely generated common cover, \(K_N\) its level kernels, \(K_\infty\) its limit kernel and \(\mathcal U\) the level ultrafilter. For every finite quotient of \(\Phi\), the image of \(K_N\) equals the image of \(K_\infty\) on a set in \(\mathcal U\). In the construction by intersection this last assertion follows by compactness: the chosen subgroup images are compatible under maps of finite quotients, and each element of a chosen image lifts to the inverse limit of the images. The same finite-quotient construction gives agreement of the marked elements and of the local maps at that precision.

Fix \(m\). There are finitely many continuous maps from the common covering group to any specified finite group that are homomorphisms, since their values on topological generators determine them. This applies also to its open CM subgroup. Encode a cocycle by the homomorphism with values in the semidirect product of the additive matrix module and the finite matrix image. Intersect the kernels of the finitely many candidate homomorphisms, taking the normal core in the cover if necessary. In the resulting finite quotient agreement of kernel images tests factorisation simultaneously for all candidates. At this precision the coefficient action and the local maps also agree. Upper triangularity after subtraction of a local boundary, its diagonal torsion condition, and the polarised transpose equation can all be tested there. Each local bounding matrix ranges over a finite set. This proves simultaneous transfer in both directions for the even, \(V_1\)- and \(V_2\)-cocycles. Boundaries transfer using the same matrix on the common cover. In particular the first bound of \ref{B6.48aaa} holds on the limit at every depth with the same constant.
\leanref{UpperTheory.\allowbreak{}V.\allowbreak{}Pa.\allowbreak{}t1}{Development.\allowbreak{}AutomorphicArithmetic.\allowbreak{}AutomorphicVarieties.\allowbreak{}PatchingModules.\allowbreak{}SeparableMatrixStability.\allowbreak{}LcOfAdic}{122--209}
\leanref{UpperTheory.\allowbreak{}V.\allowbreak{}Pa.\allowbreak{}taylorWiles\_\allowbreak{}first\_\allowbreak{}condition\_\allowbreak{}at\_\allowbreak{}limit\_\allowbreak{}of\_\allowbreak{}group\_\allowbreak{}condition}{Development.\allowbreak{}AutomorphicArithmetic.\allowbreak{}AutomorphicVarieties.\allowbreak{}PatchingModules.\allowbreak{}SeparableMatrixStability.\allowbreak{}TaylorWilesFirstConditionAtLimitOfGroupCondition}{19--42}

Suppose \(b\) rational classes are independent. Clear denominators in the cocycles and their local witnesses. Evaluate on generators, including the columns of the global boundary map. The span of the \(b\) columns in the quotient by the boundary lattice has rank \(b\). Smith form gives a constant \(c\) such that a combination reducing to a boundary modulo \(t^m\) has every coefficient in \(t^{m-c}A\), for \(m\ge c\). Hence \((m-c)b\le mD_*+C_1\) for every such \(m\). Division by \(m\) gives \(b\le D_*\).

At a finite level the product and local flag defects of a lift to \(V_2\) give a pair defining a class in \(E_m\). Indeed the unprescribed slot is the odd \(\eta\)-slot, which is square-zero; associativity is the global two-cocycle equation and the flag equations say its local projection is a differential. Changing that slot by a one-cochain and changing a flag witness changes the pair by its defining boundary. Thus vanishing of the class is equivalent to a lift. The defect is linear in the prescribed primed tuple, since products of two primed slots vanish.

At depth \(m\), the \((g+1)\)-st largest cyclic summand of \(E_m\), if present, has length \(b\le C_2\): take \(e=b\) in \ref{B6.48aaa}. Project the \(g+1\) defects to the largest \(g\) summands, lift their coordinates to a matrix over \(A\), and choose a kernel vector over \(F\). Scale to get integral coordinates with one unit. Multiplication by \(t^{C_2}\) kills the remaining summands of its defect. Take one level at which the bound and both finite transfers hold. Transfer the prescribed tuples there, take this combination and its lift, and transfer the lift back. Equality of the primed components is equality on the common cover, so it survives this operation.
\leanref{UpperTheory.\allowbreak{}V.\allowbreak{}Pa.\allowbreak{}taylorWiles\_\allowbreak{}second\_\allowbreak{}condition\_\allowbreak{}at\_\allowbreak{}limit\_\allowbreak{}of\_\allowbreak{}group\_\allowbreak{}condition}{Development.\allowbreak{}AutomorphicArithmetic.\allowbreak{}AutomorphicVarieties.\allowbreak{}PatchingModules.\allowbreak{}SeparableMatrixStability.\allowbreak{}TaylorWilesSecondConditionAtLimitOfGroupConditio}{7--27}

For the passage through the depths, first clear denominators in the given tuples. For each \(m\) retain the coefficient vector \emph{before} multiplication by \(t^{C_2}\), the lifted semidirect-product homomorphism and its flag witnesses. Their reductions at a fixed smaller depth belong to finite sets. Select a diagonal subsequence with constant reductions at each depth and a fixed unit coordinate modulo \(t\). The compatible homomorphisms define a continuous integral cocycle, since each finite reduction is continuous. Its primed tuple is \(t^{C_2}\) times the limiting combination, by equality at every depth and \(t\)-adic separation. The coefficient with the fixed unit coordinate remains nonzero over \(F\).
\leanref{UpperTheory.\allowbreak{}V.\allowbreak{}Pa.\allowbreak{}secondCohomologyBound\_\allowbreak{}of\_\allowbreak{}boundaryBound}{Development.\allowbreak{}AutomorphicArithmetic.\allowbreak{}AutomorphicVarieties.\allowbreak{}PatchingModules.\allowbreak{}SeparableMatrixStability.\allowbreak{}SecondCohomologyBoundOfBoundaryBound}{44--120}

For odd cohomology, let \(p:\Gamma_\infty\to\Gamma_{\rm base}\) be the map on the pro-\(\ell\) reductions used in the construction. Its kernel is the closed normal subgroup generated by the marked tame elements \(\tau_j\). Here is the finite-quotient check of the reverse inclusion. Let \(V\) be an open normal subgroup of the cover giving a quotient of \(\Gamma_\infty\) that kills all \(\tau_j\). On a set in \(\mathcal U\), the same quotient factors through the level group and kills its tame generators, hence factors through the base by \ref{B6.56}. There are finitely many homomorphisms from the finitely generated base to this finite quotient. Take the one occurring on a set in \(\mathcal U\). Finite-quotient agreement of the maps to the base identifies its composite with the limit projection. Thus every such finite quotient kills \(\ker p\), as required.
\leanref{UpperTheory.\allowbreak{}V.\allowbreak{}Pa.\allowbreak{}ker\_\allowbreak{}oo}{Development.\allowbreak{}AutomorphicArithmetic.\allowbreak{}AutomorphicVarieties.\allowbreak{}PatchingModules.\allowbreak{}SeparableMatrixStability.\allowbreak{}KerOo}{15--80}

In the limit \(\phi_j\tau_j=\tau_j\phi_j\), with \(p(\tau_j)=1\) and \(p(\phi_j)=s_j\). Evaluating an odd cocycle on the two products shows that \(z(\tau_j)\) commutes with \(r(s_j)\). The two blocks have disjoint spectra there, so \(z(\tau_j)=0\). The cocycle homomorphism therefore descends; on the CM subgroup one also uses the polarisation equation for the conjugates by complex conjugation. The local maps to the base identify the local tests on the descended cocycle. Conversely pullback preserves those tests. If an inflated linear combination is a boundary \(d_rY\), the same matrix \(Y\) bounds it on the base by surjectivity. Hence the inflated \([\kappa_j]\) both span and are independent over \(F\).
\leanref{UpperTheory.\allowbreak{}V.\allowbreak{}Pa.\allowbreak{}koszulThree\_\allowbreak{}of\_\allowbreak{}det}{Development.\allowbreak{}AutomorphicArithmetic.\allowbreak{}AutomorphicVarieties.\allowbreak{}PatchingModules.\allowbreak{}SeparableMatrixStability.\allowbreak{}KoszulThreeOfDet}{8--50}
\end{proof}

\subsection*{Proof of \ref{B6.48aac}}\label{prf:B6.48aac}\phantomsection
\begin{proof}
Write \(\mathfrak l=\mathfrak m_\Lambda\) and \(h=nq\). At fixed precisions \(b,j\), the modules \(M_N/(Y^b+\mathfrak l^j)M_N\) are free over \((\Lambda/\mathfrak l^j)[T_1,\ldots,T_h]/(T)^b\) for all sufficiently high \(N\), by Lemma~\ref{B6.48aae}(1). Indeed the binomial coefficient of degree \(i<b\) in the \(N\)-th diamond relation is divisible by \(\ell^j\) once \(v_\ell(\mathrm Nw_N-1)\ge b+j\). Hence, for \(0\le a<b\), multiplication by the monomials of degree \(a\) identifies copies of the augmentation with the degree-\(a\) layer. This map is canonical; a basis at level \(N\) proves it is an isomorphism and does not enter its definition.

Use the quotients \(L_b\) by sequences tending to zero modulo \(Y^b+\mathfrak l^j\) along the ultrafilter for every \(j\). The maps \(L_{a+1}\to L_a\) are surjective, since the same representative sequence serves at both precisions. For \(a\ge1\) their kernel is canonically
\(\bigoplus_{|\alpha|=a}L_1T^\alpha\). For injectivity, a relation maps to zero at every \(\mathfrak l^j\)-precision. The finite-level monomial isomorphism makes each coefficient zero at that precision. For surjectivity, take a representative in the kernel. At each precision \(j\) it belongs to \(Y^a+\mathfrak l^j\) on an ultrafilter set on which that monomial isomorphism holds. Intersect the sets for precisions \(1,\ldots,j\) and with \(N\ge j\). For each \(N\) use the largest applicable \(j\le N\), choosing arbitrary values if there is none. This precision tends to infinity along the ultrafilter. Lift the degree-\(a\) coefficients at these precisions to \(M_N\). Their sequences give the required finitely many elements of \(L_1\).

At augmentation the two comparison maps with \(X\) of Lemma~\ref{B6.48aae}(2) have composites \(f\), at each \(\mathfrak l^j\)-precision. Choose them at precisions tending to infinity in the same manner. The ultrapower modulo convergence of the constant module \(X\) is \(X\): its quotients \(X/\mathfrak l^jX\) are finite, and \(X\) is finite over the complete ring \(\Lambda\). Each sequence selects one compatible residue at each depth. It follows that \(L_1[1/f]\simeq X[1/f]\). Localisation is exact, so the descriptions of all the successive kernels persist after inverting \(f\).
\leanref{UpperTheory.\allowbreak{}V.\allowbreak{}Pa.\allowbreak{}Core.\allowbreak{}eventually\_\allowbreak{}exists\_\allowbreak{}patchedModule\_\allowbreak{}quotient\_\allowbreak{}maps}{Development.\allowbreak{}AutomorphicArithmetic.\allowbreak{}AutomorphicVarieties.\allowbreak{}PatchingModules.\allowbreak{}AnnihilatorBaseChange.\allowbreak{}EventuallyExistsPatchedModuleQuotientMaps}{193--240}

Take \(\mathcal M=\varprojlim L_b[1/f]\). Lift a finite generating set of \(X[1/f]\) to \(\mathcal M\) by surjectivity of the transitions. Successively subtract homogeneous polynomials in \(Y\) times these lifts, using the kernel isomorphisms. The coefficients form convergent series in \(\mathcal B=\widehat{R[1/f]}^{\,Y}\). This proves finite generation. Starting the subtraction in degree \(b\) for an element zero in \(L_b[1/f]\) gives coefficients in \(Y^b\mathcal B\). It follows that the kernel of projection is exactly \(Y^b\mathcal M\), and proves the graded identity before local completion.

Since \(Y\subset p\) and \(f\notin p\), the ideal \(p\mathcal B\) is prime: modulo \(Y\), the completion has quotient \(R[1/f]/Y\). More generally
\(\mathcal B/p^j\mathcal B=R[1/f]/p^jR[1/f]\), since \(Y^j\subset p^j\) and completion is exact for finite modules over the Noetherian ring \(R[1/f]\). Localising these identities and taking the inverse limit identifies \(\widehat{\mathcal B_{p\mathcal B}}\) with \(B'\). This map from \(\mathcal B\) is flat. Tensor the exact sequences for two consecutive powers of \(Y\) to obtain the asserted graded isomorphism over \(B'\).
\leanref{UpperTheory.\allowbreak{}V.\allowbreak{}Pa.\allowbreak{}Core.\allowbreak{}yPrime}{Development.\allowbreak{}AutomorphicArithmetic.\allowbreak{}AutomorphicVarieties.\allowbreak{}PatchingModules.\allowbreak{}AnnihilatorBaseChange.\allowbreak{}YPrime}{23--124}

Put \(p_\Lambda=p\cap\Lambda\). The series variable \(s\) maps to \(t\), so \(\Lambda/p_\Lambda=k[[t]]\). More explicitly, lift the images of the \(d=nd_+\) weight coordinates to series in \(O[[s]]\). Subtracting those series from the coordinates identifies \(p_\Lambda\) with the ideal generated by the resulting \(d\) variables and the uniformiser. Thus \(L=\widehat{\Lambda_{p_\Lambda}}\) is regular of dimension \(1+d\). The augmentation module \(X\) is finite free over \(\Lambda\) by base-level freeness and finite free root adjunction, with the prescribed local factor selected by an idempotent. To see the effect of the second completion, let \(D\) be the image of \(R\) in \(\operatorname{End}_\Lambda(X)\). It is finite over \(\Lambda\). The finite \(L\)-algebra \(D\otimes_\Lambda L\) is the product of its complete local factors above \(p_\Lambda\). Completing \(X\) at \(p\) selects the corresponding direct summand of \(X\otimes_\Lambda L\), or zero if \(p\) is outside its support. Thus \(B'\otimes_R X\) is finite free over \(L\). If \(M'\) has a support prime, this augmentation is nonzero by Nakayama.

Apply \ref{B6.48ac} to \(L\to B'\), the finite nonzero module \(M'\), and the \(h\) elements of \(Y\) in its maximal ideal. Its polynomial graded-map hypothesis is the isomorphism just proved. Each minimal support prime has quotient dimension \((1+d)+h\). Finally \(\mathfrak aB'\subseteq\sqrt{\operatorname{Ann}_{B'}M'}\) by Lemma~\ref{B6.48aae}(3). Passage to \(C\) therefore identifies each minimal support prime and its quotient ring, giving quotient dimension \(1+n(d_++q)\).
\end{proof}

\subsection*{Proof of \ref{B6.48bb}}\label{prf:B6.48bb}\phantomsection
\begin{proof}
To the first list adjoin the minimal primes of each \(B_j\) which contain none of the \(A_i\). Every ideal on the enlarged list has quotient dimension less than \(b\). A prime avoiding this list avoids every \(B_j\), since a prime containing \(B_j\) contains a minimal prime over it.

Cutting \(R/I\) by \(s\) gives dimension at least \(b\). Choose a component of this dimension or greater, with prime \(\mathfrak q_0\). From each ideal on the enlarged list take an element outside \(\mathfrak q_0\), and denote the product of their images in \(D=R/\mathfrak q_0\) by \(f\). This product is nonzero. While \(\dim D>1\), choose a nonzero \(z\in\mathfrak m_D\) outside the finitely many height-one primes containing \(f\). This is possible by prime avoidance. Each minimal prime over \((z)\) has height one and avoids \(f\). Passage to its quotient preserves \(f\ne0\) and decreases the dimension by one, by the dimension theorem for a principal ideal in a complete local domain. The final domain gives \(\mathfrak p\).

In equal characteristic its normalisation is finite, since a complete Noetherian local ring is excellent. The normalisation of this one-dimensional domain is local: it is a finite algebra over a Henselian local ring and has no nontrivial idempotents. It is a complete DVR with finite residue extension \(k'\). A coefficient field and a uniformiser identify it with \(k'[[t]]\).
\end{proof}

\subsection*{Proof of \ref{B6.48bf}}\label{prf:B6.48bf}\phantomsection
\begin{proof}
Let \(k'=k_\Lambda\), let \(K=G_F\cap\ker\bar r_D\), and form the inertia quotient \(\Gamma\) of \(G_F\) by the closed normal subgroup generated by \(I_q\cap K\) at all finite places \(q\nmid\ell\) of \(F\). Because \(\bar r_D\) takes values in \((k'^\times)^4\), which has order prime to \(\ell\), the open normal subgroup \(K\triangleleft G_F\) has index prime to \(\ell\). Let \(V=\mathfrak{sp}_4(k')\) be the \(10\)-dimensional adjoint representation of \(\Gamma\) (infinitesimal symplectic matrices with fixed similitude \(\nu|_{G_F}\)), and at each \(v\mid\ell\) let \(b_v\subseteq V\) be the \(6\)-dimensional Borel subalgebra preserving the residual ordinary flag at \(v\).

We verify \(\operatorname{Hom}_\Gamma(V,k'(\bar\varepsilon_\ell))=0\). In both residual support shapes of Lemma~\ref{B6.48b}, the residual representation \(\bar r\) is lower triangular after ordering the basis as \((0,1,2,3)\) for support in \(\{-\alpha,-\beta,-\psi,-\theta\}\) or as \((0,2,1,3)\) for support in \(\{\beta,-\alpha,-\theta,-\psi\}\). Hence every Jordan--H\"older constituent of \(V\) is one-dimensional, with character either \(1\) (twice, on the diagonal) or one of the eight root ratios \(d_i/d_j\) (\(i\ne j\)). For any diagonal constituent, complex conjugation \(c\in G_F\) satisfies \(1(c)=1\ne-1=\bar\varepsilon_\ell(c)\) since \(\ell\) is odd. For any root constituent \(d_i/d_j\) (\(i\ne j\)), the marking \(g_0\in G_F\) of the fitted tower has residual cyclotomic value \(\bar\varepsilon_\ell(g_0)=1\) and four pairwise distinct eigenvalues \(d_0(g_0),d_1(g_0),d_2(g_0),d_3(g_0)\) (preserved when raising \(g_0\) to an \(\ell\)-power upon passing to a deep layer, as \((d_i/d_j)(g_0)\in k'^\times\) has order prime to \(\ell\)), whence \((d_i/d_j)(g_0)\ne1=\bar\varepsilon_\ell(g_0)\). Thus no constituent of \(V\) equals \(k'(\bar\varepsilon_\ell)\) on \(G_F\), giving \(\operatorname{Hom}_\Gamma(V,k'(\bar\varepsilon_\ell))=0\).

At each real place \(w\mid\infty\) of \(F\), complex conjugation \(c_w\in G_F\) has symplectic similitude \(-1\) and eigenvalues \((1,1,-1,-1)\), so its centralizer in \(\mathfrak{sp}_4(k')\) is isomorphic to \(\mathfrak{gl}_2(k')\) and has dimension \(4\). Because \(\sum_{v\mid\ell}d_v=d\), Lemma~\ref{B6.51framea} (which uses only the cohomological inputs of Definition~\ref{B6.48be}) applies to \(V\) and \(b_v\), yielding finite-dimensional cone cohomology groups \(H_b^i\) (\(0\le i\le2\)) with
\[
 \dim_{k'}H_b^1-\dim_{k'}H_b^2-\dim_{k'}H_b^0
 \ge \sum_{v\mid\ell}d_v\dim_{k'}b_v-\sum_{w\mid\infty}\dim_{k'}V^{c_w=1}
 = 6d-4d=2d .
\]

Let \(R_0\) be the completed universal deformation ring at \(x\) of the marked symplectic representation with fixed similitude \(\nu|_{G_F}\) and ordinary flags at \(v\mid\ell\), before imposing the prescribed \(\ell\)-primary torsion values on the local flag characters. Take a minimal presentation \(S=\Lambda[[X_1,\dots,X_{t_0}]]\twoheadrightarrow R_0\) with kernel \(I_0\) generated by \(r_0\) elements. By the square-zero obstruction calculation in Lemma~\ref{B6.51ab}, pushing \(S/\mathfrak m_S I_0\to S/I_0\) out along any \(k'\)-linear functional on \(I_0/\mathfrak m_S I_0\) yields a class in \(H_b^2\) that vanishes if and only if the functional is zero, so \(r_0\le\dim_{k'}H_b^2\). For the relative tangent space, every cone cocycle in \(H_b^1\) lifts to a first-order marked deformation \((1+\epsilon z)\bar r\) with ordinary flags (conjugating by a suitable first-order symplectic matrix to diagonalize the marking \(r(g_0)\), which is possible because its four residual eigenvalues are distinct). The kernel of the resulting surjection \(\operatorname{Tan}(R_0)\twoheadrightarrow H_b^1\) consists of coboundaries preserving the diagonal marking \(r(g_0)\), hence corresponds to the centralizer of \(\bar r(g_0)\) in \(\mathfrak{sp}_4(k')\) modulo \(H_b^0\). Because \(\bar r(g_0)\) has four distinct eigenvalues, its centralizer in \(\mathfrak{sp}_4(k')\) is the \(2\)-dimensional diagonal torus \(\mathfrak t\supseteq H_b^0\), so the kernel has dimension \(2-\dim_{k'}H_b^0\). Thus \(t_0\ge\dim_{k'}H_b^1+2-\dim_{k'}H_b^0\), and combining with the Euler characteristic bound gives
\[
 t_0-r_0\ge\dim_{k'}H_b^1-\dim_{k'}H_b^2-\dim_{k'}H_b^0+2\ge 2d+2 .
\]

Finally, \(R\) is obtained from \(R_0\) by imposing the prescribed values of the four local flag characters \(\psi_{v,0},\dots,\psi_{v,3}\) on the \(\ell\)-primary torsion subgroup of \(I_v^{\mathrm{ab}}\) at each of the \(s\) places \(v\mid\ell\). By local class field theory, this \(\ell\)-primary torsion subgroup is isomorphic to \(\mu_{\ell^\infty}(F_v)\), which is finite cyclic. Choose a generator \(\gamma_v\) of it. The symplectic pairing fixes \(\psi_{v,0}\psi_{v,3}=\psi_{v,1}\psi_{v,2}=\nu|_{I_v}\), so fixing the two characters \(\psi_{v,0}(\gamma_v)\) and \(\psi_{v,1}(\gamma_v)\) determines all four characters on the \(\ell\)-primary torsion of \(I_v^{\mathrm{ab}}\). Imposing these conditions across all \(s\) places adds at most \(2s\) relations to the presentation of \(R_0\), producing a presentation of \(R\) with \(t=t_0\) variables and \(r\le r_0+2s\) relations satisfying \(t-r\ge 2d+2-2s\).
\end{proof}

\subsection*{Proof of \ref{B6.48b}}\label{prf:B6.48b}\phantomsection
\begin{proof}
Lift minimal primes to the completed local ring \(R\) at \(x\). Write \(u:A\to R\), \(\varphi=uj\), \(b=2d+1-2s\) and \(d_v=d/s\), the common local degree at \(\ell\). By Lemma~\ref{B6.48bf}, \(R\) admits a presentation over \(\Lambda\) with at least \(b+1=2d+2-2s\) more variables than relations. Complete local connectedness \cite[Exp.~XIII, Thm.~2.1]{grothendieck-1968} gives a chain of minimal primes with adjacent \(P,Q\) satisfying
\begin{equation}\tag{6.48b.1}\label{eq:p6c5-four-adjacent}
 \dim R/(P+Q)\ge b+1 .
\end{equation}

Suppose \(JR\subseteq P\). The reducible support loci in the first four alternatives and in \(H\) apart from \(\{\pm\beta,\pm\theta\}\) have dimension at most \(d+20C+1\). Lemma~\ref{B6.48bc} gives the degeneration and the entry for each support. At its limiting flagged point the chosen odd entry and two diagonal entries cost at most \(d+3C\) tangent dimensions, the ten symplectic kernel coordinates at most \(10C\), and the flag entries at most \(6s\). The coefficient parameter contributes one. In the Hecke special fibre Lemma~\ref{B6.48bd} gives \(2d+1-d_v\) parameters at a coincidence on \(U_v\), including the single inertia value used to extend equality to all of \(I_v\). Off the preceding reducible loci, the matrix fibre contributes at most \(2+4s\). The marked projectors separate the entries, and trace products determine the opposite entry from one entry of each pair. Since \(s\le C\) and \(d\ge30(C+1)^2\),
\begin{equation}\tag{6.48b.2}\label{eq:p6c5-four-exclusions}
 d+20C+1<b,\qquad 2d+1-d_v+2+4s<b.
\end{equation}
On the remaining support \(\{\pm\beta,\pm\theta\}\), the nonzero flag entry \((n_v)_{10}\) gives
\[
 (\psi_{v,1}-\psi_{v,0})(\psi_{v,1}-\psi_{v,3})=0
 \quad\text{on }G_v .
\]
The two equality subgroups cover \(G_v\), so one is all of \(G_v\). This support also lies in the coincidence locus.

Apply \ref{B6.48bb} with \(I=P+Q\) and \(s=\varpi\), a uniformiser of \(\Lambda\). For the first list use the reducible support loci bounded by \(d+20C+1\). For the second list use the loci defined by \(J\), \(\varpi\) and coincidence of two specified characters on \(U_v\). Outside the first list their dimensions are bounded by \(2d+1-d_v+2+4s\), where here \(s\) denotes the number of flags. Both bounds are strictly less than \(b\) by \eqref{eq:p6c5-four-exclusions}, and \eqref{eq:p6c5-four-adjacent} gives the dimension hypothesis on \(I\). The remaining block support forces one of the character coincidences just excluded. The lemma gives a local map \(R\to k'[[t]]\) through \(R/(P+Q,\varpi)\) whose generic point avoids all these loci. Compose with \(u\) to obtain \(y\). Its generic span is full, its characters are separated and its constant term is residual. It kills \(J\) because \(JR\subseteq P\). The curve condition applies to \(\varphi^{-1}(Q)\subseteq\ker(yj)\), so \(JR\subseteq Q\). Induction along the chain and contraction give the result.
\end{proof}

\subsection*{Proof of \ref{B6.48bc}}\label{prf:B6.48bc}\phantomsection
\begin{proof}
Index rows and columns from zero, and write \(i^*=3-i\). The marked diagonal matrix has four distinct residual entries. Over any field at a prime in the completion its spectral projectors are the four coordinate projectors. Consequently the matrix algebra spanned by the representation contains \(E_{ij}\) precisely when some matrix has nonzero \(ij\)-entry. The relation of containing \(E_{ij}\) is reflexive and transitive, and is invariant under \(ij\mapsto j^*i^*\) by the similitude inverse formula. Thus the full matrix algebra is the only support with all eight roots. The following table gives the proper supports needed, and the exceptional odd entry to retain in the tangent count. An entry \(ij\) means the single matrix coordinate in that position. A dash means the residual entries in that column of cases exclude the row.
\[
\begin{array}{c|l|ccc}
 & S & 1,2 & 3,4 & 5\\ \hline
1&\beta,-\beta,-\alpha,-\theta,-\psi&12&12&12\\
2&\beta,-\beta,\alpha,\theta,\psi&-&-&-\\
3&\theta,-\theta,\psi,\beta,-\alpha&-&12&-\\
4&\theta,-\theta,-\psi,-\beta,\alpha&-&-&03\\
5&\alpha,-\alpha,-\theta,-\psi,-\beta&01&01&01\\
6&\alpha,-\alpha,\theta,\psi,\beta&-&01&-\\
7&\psi,-\psi,\beta,-\alpha,-\theta&12&12&12\\
8&\psi,-\psi,-\beta,\alpha,\theta&-&-&03\\
9&\beta,-\beta,\theta,-\theta&-&03&\text{exception}
\end{array}
\]
For exhaustiveness one can enumerate with eight bits, one for each entry \(01,02,03,10,12,20,21,30\). Copy each bit to its reflected position and put ones on the diagonal. Discard a choice if \(b_{ij}b_{jk}>b_{ik}\) for some \(i,j,k\). Equivalently, any proper transitive choice has a nonempty proper invariant coordinate subspace (take a terminal class of its directed graph). The algebra preserves its symplectic perpendicular as well. A coordinate line, together with its perpendicular, gives four maximal supports; a coordinate isotropic plane gives four more. The remaining coordinate planes are the perpendicular nondegenerate pair \(03,12\), giving row 9. Writing down the entries preserving each indicated subspace gives the nine sets displayed. A support through \(x\) must contain all its nonzero residual entries. In columns \(1,2\) this requires \(-\psi\) and one of \(\beta,-\alpha\), or both \(-\alpha,-\beta\). In columns \(3,4\) it requires \(-\alpha\), or both \(\beta,-\theta\). In column 5 it requires \(-\psi\), or both \(-\beta,-\theta\). These tests give the dashes.
\leanref{UpperTheory.\allowbreak{}V.\allowbreak{}Ws.\allowbreak{}redList}{Development.\allowbreak{}AutomorphicArithmetic.\allowbreak{}AutomorphicVarieties.\allowbreak{}WeakShimuraTheory.\allowbreak{}CompactIntegralFields.\allowbreak{}RedList}{124--140}

Let \(R_S\) be the completed local quotient cutting out entries outside \(S\). If it is nonzero, its relative tangent space over \(\Lambda\) injects into pairs \((Z,Y_v)\) consisting of the derivative of the representation times its inverse and the derivative of each flag matrix times its inverse. Indeed the values of both derivatives recover the lift to the dual numbers, by the defining property of the marked matrix ring. The representation derivative satisfies the cocycle and infinitesimal similitude equations. We bound this tangent space after diagonal degeneration. For exponents \(m_i\), conjugate the matrices by \(\operatorname{diag}(T^{m_i})\); each entry has weight \(m_i-m_j\). In each of the following choices \(m_i+m_{i^*}\) is constant, so this conjugation preserves the equations. The residual representation and flags have no negative-weight entries. Filter the tangent pairs by least entry weight, simultaneously at every group element and every flag. In the linearised multiplication, form and flag equations, the term of least weight satisfies the same equations at the weight-zero residual data. The support equations pass to that term too. Each graded piece thus injects into the corresponding weight subspace of the tangent equations at the limit. There are finitely many entry weights and the target weight subspaces are independent. Summing gives the desired tangent inequality, over the original finite residue field.
\leanref{UpperTheory.\allowbreak{}V.\allowbreak{}Ws.\allowbreak{}tan\_\allowbreak{}flow}{Development.\allowbreak{}AutomorphicArithmetic.\allowbreak{}AutomorphicVarieties.\allowbreak{}WeakShimuraTheory.\allowbreak{}RedTanFlow}{93--127}
\leanref{UpperTheory.\allowbreak{}V.\allowbreak{}Cu.\allowbreak{}redA\_\allowbreak{}of}{Development.\allowbreak{}AutomorphicArithmetic.\allowbreak{}AutomorphicVarieties.\allowbreak{}ModularityArgument.\allowbreak{}CuspidalTwists.\allowbreak{}PRedA}{57--151}

For cases 1 and 2 use \(m=(0,2,1,3)\) for the first residual shape of \ref{B6.48b} and \(m=(0,1,2,3)\) for the second. Inspection of their allowed entries gives the diagonal representation and identity flags in the limit. For cases 3 and 4 use \(m=(0,1,0,1)\). The surviving flag is \(1+b_v(E_{20}+E_{31})\) with \(b_v=(n_v)_{20}\ne0\). For case 5 use \(m=(0,0,1,1)\). Its flag is \(1+b_v(E_{10}-E_{32})\), \(b_v=(n_v)_{10}\ne0\). The signs follow from \(n_v\) being lower unipotent symplectic. The representation in both limits is diagonal. The local flag condition gives \(\bar\chi_0=\bar\chi_2\) in cases 3 and 4, and \(\bar\chi_0=\bar\chi_1\) in case 5: apply the representation to the first flag column and take its nonzero \(2\)- or \(1\)-coordinate.

At a diagonal representation \(r_0\), each \(Z_{ij}\) is a cocycle for \(\bar\chi_i/\bar\chi_j\), and is zero off the support for \(i\ne j\). At a local group put
\(D=Z+\operatorname{Ad}(r_0)Y_v-Y_v\), with the sign convention for the derivative of the flag matrix just chosen. Conjugating \(D\) by the limiting flag basis makes it upper triangular. The six independent odd positions are \(01,03,21,12,10,30\). Fix to zero the exceptional global odd coordinate in the table and the global diagonal coordinates \(00,11\). We show that each remaining odd cocycle is locally a boundary.

For the identity flag the entries below the diagonal are boundaries, and among \(01,03,12\) at most the indicated exceptional entry belongs to the row. For the flag with \(b=b_v=(n_v)_{20}\), multiplication by its inverse on the left and by the flag on the right gives
\[
 D_{10}=-bD_{12},\qquad D_{30}=-b^2D_{12},\qquad
 D_{21}=2bD_{01}+b^2D_{03}.
\]
These relations also hold for local cohomology classes; entries combined in any equation have the same local character by \(\bar\chi_0=\bar\chi_2\) and similitude. In the first two equations the class of \(12\) vanishes if any of \(12,10,30\) is absent or has been fixed to zero. In the third equation two zero classes force the third to vanish, since \(2b\) and \(b^2\) are units. In rows \(1,3,5,6,7,9\) of the middle column, deleting the displayed coordinate leaves respectively at least one zero among \(12,10,30\) and at least two among \(21,01,03\). This proves the assertion in both cases, including row 9 of case 4.

For the flag with \(b=(n_v)_{10}\) the lower equations give
\[
 D_{21}=D_{30}=0,\qquad
 D_{10}=b^2D_{01}+b(D_{00}-D_{11}).
\]
Here \(D_{00}=D_{11}=0\) in the kernel under consideration. The local characters for \(10,01\) are trivial. In rows \(1,4,5,7,8\) of the last column, deleting the indicated coordinate kills \(03,12\) and at least one of \(01,10\). The displayed equations kill the remaining local classes.

The exceptional coordinate has global cocycle dimension at most \(d+C\), and the two diagonal coordinates contribute at most \(2C\). In the kernel, reflection in the antidiagonal determines the Lie algebra matrix from the ten positions \(i\le j^*\). Each of these contributes at most \(C\): for an odd position use the strict cocycle bound just proved applicable, and for the other positions use the bounds at \(0,\pm\psi\). The flag derivatives are strictly lower and have at most six coordinates per place. Hence the relative tangent dimension is at most \(d+3C+10C+6s\le d+20C\). Lifts of a relative cotangent basis together with a uniformiser generate the maximal ideal by Nakayama's lemma. The height theorem gives \(\dim R_S\le1+d+20C\).
\leanref{UpperTheory.\allowbreak{}V.\allowbreak{}Cu.\allowbreak{}ReductionTangentLimit.\allowbreak{}rank\_\allowbreak{}A}{Development.\allowbreak{}AutomorphicArithmetic.\allowbreak{}AutomorphicVarieties.\allowbreak{}ModularityArgument.\allowbreak{}CuspidalTwists.\allowbreak{}PRedTanA}{160--255}
\leanref{UpperTheory.\allowbreak{}V.\allowbreak{}Ws.\allowbreak{}redDim}{Development.\allowbreak{}AutomorphicArithmetic.\allowbreak{}AutomorphicVarieties.\allowbreak{}WeakShimuraTheory.\allowbreak{}RedDim}{88--124}

In the exception work over the fraction field at the prime. The representation preserves the planes \(03\) and \(12\), each symplectic of dimension two, so both block determinants equal the multiplier \(\nu(s)\). The first flag vector has nonzero projections to both planes: its coordinates \(0,1\) are \(1,(n_v)_{10}\), and the second is a unit at \(x\). For every local element \(s\) each block therefore has eigenvalue \(\psi_{v,0}(s)\) and other eigenvalue \(\nu(s)/\psi_{v,0}(s)=\psi_{v,3}(s)\). Thus \(\psi_{v,1}(s)\) equals one of these two values. The two equality subgroups cover the decomposition group. A group cannot be the union of two proper subgroups, since the product of an element exclusive to each lies in neither. One equality holds on the whole group.
\end{proof}

\subsection*{Proof of \ref{B6.48bd}}\label{prf:B6.48bd}\phantomsection
\begin{proof}
In a prime quotient of the special fibre the inertial characters take values in the principal units and opposite characters are inverse. These follow by reduction of the prescribed inertia datum and the similitude. At a local place \(u\) take \(d_v\) inertia elements lifting generators of the torsion-free pro-\(\ell\) part of the image of inertia in the local abelianisation. Its rank is \(d_v\) by local reciprocity, since the principal units modulo torsion have this rank. These elements detect triviality of an inertial character in a prime quotient: the successive principal-unit quotients give pro-\(\ell\) image, and principal units of a characteristic-\(\ell\) local domain have no torsion. For the latter assertion, an \(\ell\)-power root of unity satisfies \((z-1)^{\ell^a}=0\), and prime-to-\(\ell\) torsion reducing to one is trivial by the geometric-series factorisation.

At the exceptional place \(v\), choose also \(s_0\in I_v\) mapping to a topological generator of \(\varepsilon_\ell(I_v)\). This exists because \(\ell\) is odd and every closed subgroup of \(\Z_\ell^\times\) is procyclic. The residual inertia datum of the fitted frame is trivial there on \(I_v\cap\ker\varepsilon_\ell\), so the latter is contained in \(U_v\). Equality on \(U_v\) and at \(s_0\) therefore gives equality on \(I_v\), by continuity. This is the one extra parameter in the count.
\leanref{UpperTheory.\allowbreak{}exists\_\allowbreak{}pinning\_\allowbreak{}elements\_\allowbreak{}of\_\allowbreak{}locAb}{Development.\allowbreak{}AutomorphicArithmetic.\allowbreak{}AutomorphicVarieties.\allowbreak{}BuildingsResidueChars.\allowbreak{}ExistsPinningElementsOfLocAb}{26--148}
\leanref{UpperTheory.\allowbreak{}exists\_\allowbreak{}mem\_\allowbreak{}pins\_\allowbreak{}insert\_\allowbreak{}ker\_\allowbreak{}of\_\allowbreak{}padic\_\allowbreak{}character}{Development.\allowbreak{}AutomorphicArithmetic.\allowbreak{}AutomorphicVarieties.\allowbreak{}BuildingsResidueChars.\allowbreak{}ExistsMemPinsInsertKerOfPadicCharacter}{6--33}

If \(j\ne i^*\), take \(a=j\); if \(j=i^*\), choose \(a\notin\{i,i^*\}\). The pairs \(a,a^*\) and \(i,i^*\) exhaust the indices. Impose \(\psi_{u,a}-1\) on all the chosen generators at all places, \(\psi_{u,i}-1\) on them for \(u\ne v\), and \(\psi_{v,i}(s_0)-\psi_{v,j}(s_0)\). All belong to the maximal ideal. In any prime containing them and \(q\), the first pair is trivial on all inertia, and the second is trivial away from \(v\). At \(v\) it is trivial by equality with the first pair if \(a=j\). If \(j=i^*\), equality with its inverse gives square one and hence triviality for a principal unit. Condition \eqref{eq:p6c5-potential-finite} now puts the whole coefficient maximal ideal in this prime. The number of elements imposed is \(sd_v+(s-1)d_v+1=2d+1-d_v\). Their ideal in \(B/q\) has maximal radical, proving the first assertion by the height theorem.

For the fibre count first use the finite-type matrix algebra over \(B\) before local completion. Over an algebraic closure of the coefficient prime's residue field, the marked eigenvalues are fixed up to a finite choice. Their projectors are polynomials in the marked matrix. In the full-span locus two tuples with the same character polynomials are conjugate: traces of products agree, and the trace pairing on the full matrix algebra identifies their matrix relations and gives an isomorphism of matrix algebras. Such an isomorphism is inner. With the ordering of the marked eigenvalues fixed the conjugator is diagonal. The two pulled-back symplectic forms differ by a scalar, since their quotient commutes with an irreducible representation. Hence the diagonal conjugator satisfies \(d_0d_3=d_1d_2\), and common scalars act trivially. The representation fibre has dimension at most two.

On the exceptional two-plane locus outside the small supports, both entries of each opposite pair \(03,30\) and \(12,21\) occur. Otherwise the support lies in one of rows \(1,2,3,4\) of \ref{B6.48bc} other than row 9, by inspection of the missing root. In each full two-dimensional block the represented cycles have fixed traces: insert the marked projectors into the trace of products and expand the projectors as polynomials. Specifying a single nonzero \(03\)-entry recovers every \(30\)-entry from such a trace, after which one nonzero \(30\)-entry recovers all \(03\)-entries. The same applies to \(12,21\); the diagonal entries are already traces with one projector. This again gives at most two parameters. For each place a lower unipotent symplectic flag matrix is determined by \(10,20,30,21\), since \(n_{32}=-n_{10}\) and \(n_{31}=n_{20}-n_{21}n_{10}\). This adds at most \(4s\).
\leanref{UpperTheory.\allowbreak{}V.\allowbreak{}Ws.\allowbreak{}fibTorus}{Development.\allowbreak{}AutomorphicArithmetic.\allowbreak{}AutomorphicVarieties.\allowbreak{}WeakShimuraTheory.\allowbreak{}FibreTraces.\allowbreak{}FibTTrace}{120--178}
\leanref{UpperTheory.\allowbreak{}V.\allowbreak{}Cu.\allowbreak{}torusY}{Development.\allowbreak{}AutomorphicArithmetic.\allowbreak{}AutomorphicVarieties.\allowbreak{}ModularityArgument.\allowbreak{}CuspidalTwists.\allowbreak{}PTorusY}{87--129}

For a prime in the completion contract it first to the finite-type matrix algebra. The dimension of the completed prime quotient is at most the local dimension of the contracted quotient at \(x\): the former is a quotient of the completion of the latter, and completion preserves local dimension. The dimension formula for a finite-type domain over the complete local domain \(B/q\) bounds this local dimension by \(\dim B/q\) plus the transcendence degree of the generic fibre. Indeed the base is universally catenary, and the closed-point residue extension at \(x\) is finite. The preceding count bounds that transcendence degree by \(2+4s\). This proves the second assertion.
\end{proof}

\subsection*{Proof of \ref{B6.48ba}}\label{prf:B6.48ba}\phantomsection
\begin{proof}
For \(h\ge3\) put \(L_h(X)=(X\,d/dX)^{h-1}(X/(1-X))\), \(\zeta=e^{2\pi i/H}\), and
\[
 E_{h,a}(z)=\frac{H^h(h-1)!}{(-2\pi i)^h}
 \sum_{(u,v)\equiv(0,a)\bmod H}'(uz+v)^{-h}.
\]
The prime omits \((0,0)\). Under \(\Gamma_0(H)\), the index transforms by multiplication by the bottom right entry. For unit \(a\) the constant term at a cusp with bottom left entry indivisible by \(H\) is zero, since no transformed row has first entry zero. The constant at infinity is \(L_h(\zeta^a)\) by expansion of the second-coordinate sum. Set
\[
 E=\frac{1}{2\tau(\chi^{-1})}
 \sum_{a\bmod H}\chi^{-1}(a)E_{h,a},\qquad
 \tau(\chi^{-1})=\sum_{a\bmod H}\chi^{-1}(a)\zeta^a .
\]
Fourier orthogonality gives the positive coefficients \(s(n)=\sum_{d\mid n}\chi(d)d^{h-1}\), with the factor two from parity. Primitivity and \(\ell\nmid H\) make the Gauss sum a unit.

Write \(c\) for the constant of \(E\). The denominator \((1-X)^h\) is a unit on nontrivial roots of order dividing \(H\). Reduce \(\theta^{-1}\) to \(1\) on the unit indices. For each \(b\bmod D\) sum over its lifts and subtract the one divisible by \(w\). Differentiating the geometric series gives \(\sum_{j=0}^{w-1} L_h(x\zeta^{Dj})=w^h L_h(x^w)\). Using \(\phi(w)=1\) gives
\begin{equation}\tag{6.48b.3}\label{eq:p6c5-eis-constant}
 2\tau(\chi^{-1})c
 \equiv(w^h-1)\sum_{b\bmod D}\phi^{-1}(b)L_h(\zeta^{wb})
 \equiv0\pmod{\mathfrak m_{C_\ell}}.
\end{equation}
The sum uses zero character coefficient at nonunits.
\leanref{UpperTheory.\allowbreak{}V.\allowbreak{}Sd.\allowbreak{}eisChi}{Development.\allowbreak{}AutomorphicArithmetic.\allowbreak{}AutomorphicVarieties.\allowbreak{}IdeleGenerators.\allowbreak{}HolomorphicIrreducibility.\allowbreak{}EisChi}{82--149}
\leanref{UpperTheory.\allowbreak{}V.\allowbreak{}Sd.\allowbreak{}constSmall}{Development.\allowbreak{}AutomorphicArithmetic.\allowbreak{}AutomorphicVarieties.\allowbreak{}IdeleGenerators.\allowbreak{}HolomorphicIrreducibility.\allowbreak{}ConstSmall}{168--225}

To produce an integral form with constant \(1\) at infinity and zero at the other cusps, separate the constants at indices \(1,b\) for each \(b\ne\pm1\) among the units. For \(u(x)=x/(1-x)^2\), the difference \(u(\zeta)-u(\zeta^b)\) has unit numerator \((\zeta-\zeta^b)(1-\zeta^{1+b})\) and unit denominator. The values of \(L_3=x(1+x)/(1-x)^3\) are units at both roots. Differentiation gives \(L_4=u(1+6u)\) and \(L_5=L_3(1+12u)\). In characteristic three use \(L_4\). Otherwise use \(L_3^2\) if its residues differ, and \(L_3L_5\) if they agree, since multiplication by \(1+12u\) separates them. Accordingly choose from the families \(E_{4,a}, E_{3,a}^2,E_{3,a}E_{5,a}\). Subtract the constant at \(b\) times the level-one form \(E_4,E_6,E_4^2\) of the same weight. The product over \(b\ne\pm1\) has even weight \(h_0\), unit constant at \(1\) and vanishing constants at those \(b\).

For \(h-h_0\ge10\), multiply by \(E_{3,a}\) if the difference is odd, and by \(E_{3,a}^2\) if it is even. Fill the remaining even weight by a monomial in \(E_4,E_6\). Let \(F_a\) be the resulting family and \(b_1\) the unit constant of \(F_1\). The Fourier expansions are integral. The character average \(G=(2b_1)^{-1}\sum_a\chi(a^{-1})F_a\) has constant \(1\) at infinity by parity. At other cusps the weight-three factor has zero constant. Conjugation of the root-of-unity coefficients gives the same construction and transformation laws for the conjugate character.
\leanref{UpperTheory.\allowbreak{}V.\allowbreak{}Sd.\allowbreak{}lkSep}{Development.\allowbreak{}AutomorphicArithmetic.\allowbreak{}AutomorphicVarieties.\allowbreak{}IdeleGenerators.\allowbreak{}HolomorphicIrreducibility.\allowbreak{}EulerianNumeratorThree}{103--181}
\leanref{UpperTheory.\allowbreak{}V.\allowbreak{}Sd.\allowbreak{}unitFam}{Development.\allowbreak{}AutomorphicArithmetic.\allowbreak{}AutomorphicVarieties.\allowbreak{}IdeleGenerators.\allowbreak{}HolomorphicIrreducibility.\allowbreak{}UnitFam}{248--264}
\leanref{UpperTheory.\allowbreak{}V.\allowbreak{}Sd.\allowbreak{}unitForm}{Development.\allowbreak{}AutomorphicArithmetic.\allowbreak{}AutomorphicVarieties.\allowbreak{}IdeleGenerators.\allowbreak{}HolomorphicIrreducibility.\allowbreak{}UnitForm}{69--155}

By \eqref{eq:p6c5-eis-constant}, \(f=E-cG\) is cuspidal and integral with first coefficient congruent to one. For \(A_p=\chi(p)p^{h-1}\) the divisor sum satisfies
\[
 s(pn)+A_p\mathbf1_{p\mid n}s(n/p)=(1+A_p)s(n)
\]
also at the level primes. Thus \(f\) gives a nonzero simultaneous residual eigenvector. In the Deligne--Serre lifting lemma \cite[Lem.~6.11]{deligne-serre-1974}, the residual eigenvector determines a maximal ideal of the commuting Hecke algebra on the integral Fourier lattice. A minimal prime below it survives in characteristic zero by torsion-freeness over the coefficient DVR. Its field of fractions gives the lifted system, and its generalised eigenspace contains a common eigenvector. One may work in a space containing the coefficientwise conjugates. The Hecke recursion determines all positive coefficients from the first, so the eigenform can be normalised. Its eigenvalue at \(\ell\) reduces to one.
\leanref{UpperTheory.\allowbreak{}V.\allowbreak{}Sd.\allowbreak{}eigLift}{Development.\allowbreak{}AutomorphicArithmetic.\allowbreak{}AutomorphicVarieties.\allowbreak{}IdeleGenerators.\allowbreak{}HolomorphicIrreducibility.\allowbreak{}EigLift}{158--281}
\leanref{UpperTheory.\allowbreak{}V.\allowbreak{}Z.\allowbreak{}Sym.\allowbreak{}Unl.\allowbreak{}eigLiftC}{Development.\allowbreak{}AutomorphicArithmetic.\allowbreak{}AutomorphicVarieties.\allowbreak{}ModularityArgument.\allowbreak{}SymmetricPowerTransfer.\allowbreak{}UnlinkedHalfIntegral.\allowbreak{}UnlinkedHeckeRoots.\allowbreak{}EigLiftC}{206--332}
\end{proof}

\subsection*{Proof of \ref{B6.49a}}\label{prf:B6.49a}\phantomsection
\begin{proof}
The monoid of exponent vectors of weight zero is the set of lattice points of the nonnegative rational cone cut out by the weight equations. It has finitely many generators. The images of their positive-degree monomials, together with \(\varpi\), generate the maximal ideal of the complete weight-zero ring \(R_0\). Indeed, in the power series presentation the monomials generate its weight-zero subring topologically, and homogeneous relations let the weight-zero component pass to the quotient.

Let \(u\) be one of these monomials. If \(u\notin\sqrt{\varpi R}\), take a prime over \((\varpi)\) avoiding \(u\). Its quotient is a complete local domain, and \(u\) is a nonzero nonunit there. In dimension greater than one, choose a nonzero parameter avoiding the finitely many height-one primes dividing \(u\) and take a minimal prime over that parameter. The resulting domain still has positive dimension and \(u\) survives. Repetition gives a one-dimensional quotient in which \(u\ne0\). Its normalisation at a maximal ideal is a complete DVR with finite residue field, since normalisation of the complete one-dimensional local domain is finite. The composite \(R\to A\) contradicts the monomial hypothesis. Thus some power of each \(u\) belongs to \(\varpi R\). Taking the weight-zero component shows that the same power belongs to \(\varpi R_0\). A uniform power of \(\mathfrak m_{R_0}\) therefore lies in \(\varpi R_0\), and \(R_0/\varpi R_0\) is finite over \(O/(\varpi)\). Lift generators and use successive \(\varpi\)-adic approximation to obtain a finite \(O\)-module. For integer weights a nonconstant monomial of total weight zero contains either a zero-weight factor or factors of both strict signs, which proves the sufficient condition.
\end{proof}

\subsection*{Proof of \ref{B6.49b}}\label{prf:B6.49b}\phantomsection
\begin{proof}
We first compute the full cuspidal trace. At each real place of weight \(k\) take the matrix coefficient of a unit lowest-weight vector in the weight-\(k\) discrete series, multiplied by the formal degree and with the convolution convention giving the projection onto that line. The central signs agree with \(\omega\) since both weight changes are even. Both \(k\)'s exceed two, so these coefficients are integrable modulo the center. The finite test is the convolution for the word at the given level with its central character. It is a finite sum of cosets modulo level and center and is unchanged by convolution with its level idempotent. These are the tests of Jacquet and Langlands \cite[\S16, proof of Thm.~16.1, especially (16.1.7)]{jacquet-langlands-1970}. At each of the two real places the integral of the matrix coefficient over the upper unipotent group, averaged over compact conjugation as in (16.1.7), vanishes. The two distinct vanishing places eliminate terms (iv)--(viii) of their trace formula. The archimedean projectors select the prescribed cusp space; one-dimensional automorphic representations contribute zero.

We use the compact expression for the remaining trace in the same comparison \cite[\S16, (16.1.4)--(16.1.8)]{jacquet-langlands-1970}. The quaternion algebra over \(F\) ramified at both real places and split at every finite place exists since there are two ramified places. Thus every finite double coset, including those at \(v,v'\), acts with its original level and multiplicity. At a real place the corresponding representation modulo the positive center is
\(\operatorname{Sym}^{k-2}\otimes\operatorname{nrd}^{\,1-k/2}\), where \(\operatorname{nrd}\) denotes reduced norm. For eigenvalues \(A,B\) of a quaternion in its complex matrix realization, put \(C_k(A,B)=\sum_{j=0}^{k-2} A^{k-2-j}B^j\). The character is \(C_k(A,B)(AB)^{1-k/2}\).

Here is the finite sum with which the two weights are compared. Let \(U\) be the product of the finite level groups, \(Z_f\) the finite adelic center, and \(D\) the quaternion algebra. Take representatives \(x\) of
\[
 D^\times(F)\backslash D^\times(\mathbf A_{F,f})/(Z_fU).
\]
This is the finite ideal class set at maximal level, modulo central ideals, with a finite refinement at the chosen level. The fiber at \(x\) consists of vectors fixed under the action of its projective stabilizer, with the specified central character included. This stabilizer is finite: the rational projective group is discrete adelically, and the stabilizing subgroups are compact at infinity and at finite places. Its order is a power of two. Indeed its image at \(u\) lies, up to conjugation, in the upper triangular subgroup modulo the residue field of order two, with pro-\(2\) reduction kernel. The projective rational group injects into its local group at \(u\). Thus averaging over the stabilizer introduces an \(\ell\)-integral denominator.

Expand the word into right cosets and take the trace on these fibers. A term on the diagonal comes from a returning coset \(a\) with \(xa=\gamma xz k_0\), where \(\gamma\in D^\times(F)\), \(z\in Z_f\) and \(k_0\in U\). Insert the stabilizer average to take the trace on the invariant fiber. Its summands have the same description with \(\gamma\) ranging through the stabilizer translates. The scalar from the finite character is a root of unity independent of the weight. This proves that the same finite list, with the same coefficients, computes both traces.

Write \(z'_w=z_w^{-1}\). The element \(\gamma\) at \(w\) is conjugate to \(z'_w M_w\), where \(M_w=a_w k_{0,w}^{-1}\) is integral and \(v_w(\det M_w)=b_w\). At either place \(w\mid\ell\) the word representatives have the form
\(\left(\begin{smallmatrix}\varpi_w^{b_w}&x_w\\0&1\end{smallmatrix}\right)\), with \(x_w\in O_w\). This follows by multiplying the representatives \(\left(\begin{smallmatrix}\varpi_w&x\\0&1\end{smallmatrix}\right)\) for a single letter, with \(x\) running through residue classes. Multiplication by \(k_{0,w}^{-1}\) gives \(\operatorname{tr}M_w\equiv1\bmod\ell\). The characteristic polynomial therefore has, by Hensel's lemma, roots \(A_w,B_w\in\Z_\ell\) with \(A_w\) a unit and \(v_\ell(B_w)=b_w\). In particular a returning coset has distinct roots and cannot contribute to the central term. Every noncentral rational quaternion in this algebra is elliptic at both real places.

Put \(P=\prod_w(\mathrm Nw)^{b_w}\) and \(Z=\prod_w(\mathrm Nw)^{v_w(z'_w)}\). The reduced norm of \(\gamma\) is totally positive, and its rational norm is \(Z^2P\). By homogeneity of \(C_h\), the two real character factors in the parallel weight, multiplied by the Hecke normalization \(P^{h/2-1}\) and transported to \(\overline{\Q}_\ell\), give
\[
 \left(\frac{z'_v z'_{v'}}{Z}\right)^{h-2}
 C_h(A_v,B_v) C_h(A_{v'},B_{v'}).
\]
Here \(F_v=F_{v'}=\Q_\ell\). The prefactor is a unit since
\(v_\ell(Z)=v_\ell(z'_v)+v_\ell(z'_{v'})\).
In the mixed weight the prefactor and the \(v'\)-factor remain the same. The \(v\)-factor becomes
\[
 C_{h+2t}(A_v,B_v)\left(\frac{\ell^{b_v}}{A_v B_v}\right)^t .
\]
Indeed the additional determinant power contributes \(((z'_v)^2A_vB_v)^{-t}\), the symmetric power contributes \((z'_v)^{2t}\) in extracting the scalar, and the word contributes \(\ell^{b_vt}\).

Both expressions are integral. Modulo \(\ell^{b_v}\) one has \(C_k(A_v,B_v)\equiv A_v^{k-2}\). The ratio of the resulting unit terms in the two weights is
\[
 \left(A_v^2\frac{\ell^{b_v}}{A_vB_v}\right)^t .
\]
Its base belongs to \(\Z_\ell^\times\). For \(N>0\), raising a unit first to \(\ell-1\) and then \(N\) times to \(\ell\) gives \(1\bmod\ell^N\) by the binomial theorem. The assertion for \(N=0\) is integrality. Summation with the stabilizer denominators proves the trace congruence.
\leanref{UpperTheory.\allowbreak{}Esk.\allowbreak{}s\_\allowbreak{}TF}{PositivityTheory.\allowbreak{}AutomorphicForms.\allowbreak{}PublishedInputs.\allowbreak{}JacquetLanglands.\allowbreak{}TorusKernels.\allowbreak{}FinalTF}{92--92}
\leanref{UpperTheory.\allowbreak{}Esk.\allowbreak{}pf\_\allowbreak{}GTerm}{PositivityTheory.\allowbreak{}AutomorphicForms.\allowbreak{}PublishedInputs.\allowbreak{}JacquetLanglands.\allowbreak{}TorusKernels.\allowbreak{}PfGTerm}{37--139}

For the eigenform assertion, simultaneously triangularize the commuting operators on both spaces over \(\overline{\Q}_\ell\). There are finitely many joint systems, counted on the diagonal with their multiplicities. For each system \(u\) on either space which differs from \(s\) by an element of absolute value at least one in some coordinate, select one such coordinate \(j(u)\). Form the polynomial
\[
 P_s(X)=\frac{X_vX_{v'}}{s_vs_{v'}}
 \prod_u\frac{X_{j(u)}-u_{j(u)}}{s_{j(u)}-u_{j(u)}} .
\]
Its coefficients are integral. For each linear factor the denominator has absolute value at least one, and the constant coefficient is integral by writing \(-u/(s-u)=1-s/(s-u)\) at the selected coordinate. On every excluded system the polynomial is zero. On every system congruent to \(s\) it is one modulo the maximal ideal. For all sufficiently large \(a\), the \(\ell^a\)-th powers of the latter values are one modulo \(\ell^N\): for \(v_\ell(y)>0\), the binomial expansion gives
\(v_\ell((1+y)^\ell-1)\ge\min(1+v_\ell(y),\ell v_\ell(y))\), and there are finitely many values.

Each monomial in \(P_s^{\ell^a}\) contains at least \(\ell^a\) copies of both ordinary operators. Choose also \(\ell^a\ge N\). The trace congruence for the monomials shows that the two multiplicities of systems congruent to \(s\) agree modulo \(\ell^N\). In parallel weight this multiplicity is an integer between \(1\) and \(N<\ell^N\). Hence the mixed space contains a nonzero joint generalized eigenspace for one such system. On that space the commuting nilpotent parts have a common nonzero kernel, giving the required cusp eigenform.
\leanref{UpperTheory.\allowbreak{}Esk.\allowbreak{}finite\_\allowbreak{}system\_\allowbreak{}proof}{PositivityTheory.\allowbreak{}AutomorphicForms.\allowbreak{}PublishedInputs.\allowbreak{}JacquetLanglands.\allowbreak{}TorusKernels.\allowbreak{}FiniteSystemProof}{206--410}
\end{proof}

\subsection*{Proof of \ref{B6.49c}}\label{prf:B6.49c}\phantomsection
\begin{proof}
For a characteristic-zero specialisation \(y:R\to\overline{\Q}_\ell\), every \(x\in\mathfrak m_R\) has \(|y(x)|<1\). Otherwise \(y(x)^{-1}\) is integral over \(\Z_\ell\), and its monic equation, multiplied by a power of \(y(x)\), puts a unit of \(1+\mathfrak m_R\) in \(\ker y\). Every element of \(R\) is the sum of an element of \(O\) and one of \(\mathfrak m_R\), so all values are integral. If \(x_1,\ldots,x_b\) generate \(\mathfrak m_R\), the values on \(\mathfrak m_R^a\) have norm at most \(\max_i|y(x_i)|^a\). Successive approximation by polynomials in the \(x_i\) with coefficients in \(O\) now puts the image in the integers of the finite extension generated by their values.

For the flag, let \(S\) be the local domain at \(\ell\) of Definition~\ref{B6.31c}, with diagonal specialisation \(x:S\to O\). Over its fraction field the four eigenvalues \(d_i=\delta_i(s_0)\) are distinct. The projector equations give the stable flag in their prescribed order and give \(\delta_i\) on inertia. The similitude pairs eigenline \(i\) with eigenline \(3-i\), so this is a symplectic flag. For \(g\in D_\ell\) its \(i\)-th quotient eigenvalue \(a_i(g)\) is integral over \(S\). Put \(b_i=\prod_{j\ne i}(d_i-d_j)\). The projector trace formula puts \(b_i a_i(g)\) in \(S\), with image \(b_i\delta_i(g)\) under \(x\). The finite domain \(S[a_i(g)]\) is local since \(S\) is henselian. At a prime over \(\ker x\), the difference \(a_i(g)-\delta_i(g)\) vanishes, since \(b_i\) has nonzero image. The difference thus belongs to the maximal ideal. Its image in the finite algebra modulo \(\mathfrak m_S\) is nilpotent. The determinant trick gives a monic annihilating polynomial for the difference reducing to a power of \(X\).

Impose these polynomial equations on the quotient lines of the symplectic flag variety over \(S\), together with flag stability and the specified characters on inertia. These are closed conditions. Their common locus contains a generic point over \(\operatorname{Frac}S\), and its projection is closed by properness. Thus over any of the specialisations under consideration there is a flag over an algebraic closure of the fraction field of the coefficient DVR. Its finitely many coordinates lie in a finite extension. Intersect its spaces with the integral lattice. The intersections are saturated direct summands, and extension of a basis of an isotropic flag to a symplectic basis gives an integral flag matrix \(u\). The monic equations give the prescribed residual characters, also in equal characteristic. The inertia identities specialise as identities.
\leanref{UpperTheory.\allowbreak{}V.\allowbreak{}Ws.\allowbreak{}Pp.\allowbreak{}genTri}{Development.\allowbreak{}AutomorphicArithmetic.\allowbreak{}AutomorphicVarieties.\allowbreak{}WeakShimuraTheory.\allowbreak{}SymplecticIwahori.\allowbreak{}GenTri}{81--125}
\leanref{UpperTheory.\allowbreak{}V.\allowbreak{}Ws.\allowbreak{}Pp.\allowbreak{}triSpec}{Development.\allowbreak{}AutomorphicArithmetic.\allowbreak{}AutomorphicVarieties.\allowbreak{}WeakShimuraTheory.\allowbreak{}SymplecticIwahori.\allowbreak{}TriSpec}{6--35}

At \(\sigma_\ell\) the residual eigenvalues on \(A\) are disjoint from those on \(B\), by Definition~\ref{B6.31b}. The first flag line therefore lies in \(A\) after reduction. The second graded line has its eigenvalue in the \(B\)-spectrum. The first flag plane is the sum of the first line and a line in \(B\), by the spectral projectors for the two disjoint spectra. Use the original label order if \(u_{00}\) is a unit. Otherwise \(u_{20}\) is a unit, and use the symplectic signed permutation of order of indices \((2,3,0,1)\). In this case the first-column eigenvector equation at every \(g\in D_\ell\) gives \(\bar\delta_2(g)=\bar\delta_0(g)\), hence \(\bar\Psi(g)=1\). The relabelling therefore respects the residual flag characters. In the chosen order the first line projects nontrivially onto coordinate \(0\). Its annihilator in \(B\) for the alternating form is the line in the first flag plane. That line projects nontrivially onto coordinate \(1\). Thus the first two leading minors of the flag matrix are units. The third is a unit by symplectic duality.

Choose the tower with this label order and its marking \(t\). Hensel lifting gives four rank-one spectral projectors for \(r(t)\), since its four residual eigenvalues are distinct. Apply the projectors to the corresponding basis vectors and normalise the two opposite symplectic pairings. This gives a symplectic eigenbasis congruent to the chosen basis modulo the maximal ideal. When \(r\) preserves \(B\), its projectors preserve \(B\), and this basis change does also. The flag matrix still has unit leading minors. Gaussian elimination writes it as \(n b\), with \(n\) lower unitriangular and \(b\) upper triangular. Both factors are symplectic by uniqueness of the factorisation and the inverse-transpose involution giving the symplectic group. In reduction the first line lies in \(A\) and the first plane splits across \(A,B\). Thus \(n_{10}=n_{30}=n_{21}=0\) in the residue field. The symplectic equations \(n_{32}=-n_{10}\) and \(n_{31}=n_{20}-n_{10}n_{21}\) give the asserted weight condition.
\leanref{UpperTheory.\allowbreak{}V.\allowbreak{}Ws.\allowbreak{}Lab.\allowbreak{}choice}{Development.\allowbreak{}AutomorphicArithmetic.\allowbreak{}AutomorphicVarieties.\allowbreak{}WeakShimuraTheory.\allowbreak{}FlagsAndDiagonals.\allowbreak{}Choice}{37--72}
\leanref{UpperTheory.\allowbreak{}V.\allowbreak{}Ws.\allowbreak{}Lab.\allowbreak{}exNorm}{Development.\allowbreak{}AutomorphicArithmetic.\allowbreak{}AutomorphicVarieties.\allowbreak{}WeakShimuraTheory.\allowbreak{}FlagsAndDiagonals.\allowbreak{}ExNorm}{7--193}

In the irreducible case the span of the group matrices over the algebraic closure is the full matrix algebra, so entries \(02,20,12\) all occur. In the extension case the same argument on \(B\) gives entries \(13,31\). If the entries from \(A\) to \(B\) all vanished, \(A\) would be an invariant complement. Their three possible weights are \(\beta,-\alpha,-\theta\).
\leanref{UpperTheory.\allowbreak{}V.\allowbreak{}Ws.\allowbreak{}Lab.\allowbreak{}starConj}{Development.\allowbreak{}AutomorphicArithmetic.\allowbreak{}AutomorphicVarieties.\allowbreak{}WeakShimuraTheory.\allowbreak{}FlagsAndDiagonals.\allowbreak{}StarConj}{128--173}
\end{proof}

\subsection*{Proof of \ref{B6.49d}}\label{prf:B6.49d}\phantomsection
\begin{proof}
Normalise the valuation of a uniformiser \(\varpi\) to one. For a root \(w\), collect the valuations of the nonzero representation entries of weight \(w\), over all group elements. Include the valuation of the corresponding nonzero flag coordinate from the following table:
\[
 \begin{array}{c|rrrr}
 w&-\alpha&-\beta&-\psi&-\theta\\ \hline
 \text{coordinate}&n_{10}&n_{21}&n_{20}&n_{30}.
 \end{array}
\]
Let \(c_w\) be the minimum of that collection, omitting the inequality for an empty collection. For \(z=(a,b)\) impose \(c_w+w\cdot z\ge0\). These inequalities give integrality for every entry of the conjugated representation and flag. Indeed the conjugating diagonal has exponents \(a+b,b,a,0\); the remaining lower flag entries are given by the two symplectic equations in Lemma~\ref{B6.49c}.

Among feasible real pairs maximise \(b\), then \(a\). The origin is feasible. The entry of weight \(-\psi\) bounds \(b\) above, and at a fixed \(b\) an entry of weight \(\beta,-\alpha\) or \(-\theta\) bounds \(a\) above. These maxima are attained in the closed polyhedron defined by the finitely many inequalities. The selected vertex has half-integral coordinates, since any two independent root normals have determinant of absolute value one or two. Adjoin \(\sqrt{\varpi}\) and use the diagonal with entries \(\varpi^{a+b},\varpi^b,\varpi^a,1\). This preserves the marking, flag characters and similitude law. Write \(S\) for the tight weights. A weight belongs to \(S\) precisely when a representation entry or its listed flag coordinate survives. All initial valuations are positive except possibly the one for \(n_{20}\). Hence \(w\cdot z<0\) for every tight \(w\ne-\psi\), and tightness of \(-\psi\) implies \(b\ge0\).
\leanref{UpperTheory.\allowbreak{}V.\allowbreak{}Ws.\allowbreak{}h7}{Development.\allowbreak{}AutomorphicArithmetic.\allowbreak{}AutomorphicVarieties.\allowbreak{}WeakShimuraTheory.\allowbreak{}CompactIntegralFields.\allowbreak{}RootSystemLinearProgramCases}{99--180}
\leanref{UpperTheory.\allowbreak{}V.\allowbreak{}Ws.\allowbreak{}Lat.\allowbreak{}vertex}{Development.\allowbreak{}AutomorphicArithmetic.\allowbreak{}AutomorphicVarieties.\allowbreak{}WeakShimuraTheory.\allowbreak{}AdicLattices.\allowbreak{}Vertex}{45--111}

We need an arithmetic restriction on \(S\). Suppose \(ij\) has odd weight, namely one of \(\pm\alpha,\pm\beta,\pm\theta\), and
\[
 R_{ix}(g)R_{xj}(h)=0\qquad(x\ne i,j;\ g,h\in G_\Q).
\]
The diagonal entries are \(d_i\), since diagonal conjugation leaves them unchanged. Matrix multiplication makes \(R_{ij}/d_j\) a cocycle for \(\gamma=d_i/d_j\). It vanishes at the marking and on the excluded inertia kernels. If
\[
 R_{ij}(s)=m(d_j(s)-d_i(s))\qquad(s\in D_\ell),
\]
then \(R_{ij}=0\) globally and \(m=0\). To see this, the odd-weight characters in either allowed order are precisely the reductions of the six characters in the last condition of Definition~\ref{B6.31b}. At each \(p\in T\setminus\{\ell\}\), inertia has an \(\ell\)-primary quotient detected by \(\vartheta\); thus \(p\equiv1\bmod\ell\) by the tame inertia relation. At the prescribed Frobenius the residual cyclotomic character is one and \(\gamma\ne1\) by the separation in that definition. Hence both the invariants of \(\gamma\) and those of its Tate dual vanish. Local duality and the local Euler characteristic for the finite \(\ell\)-primary module over \(\Q_p\), \(p\ne\ell\), give \(H^1(\Q_p,\gamma)=0\) \cite[I, Cor.~2.3 and Thm.~2.8]{milne-2006}. Apply the global coboundary condition of Definition~\ref{B6.31b}. For extended residue coefficients, apply it coordinate by coordinate over the original residue field, using the finite span of the values of the locally constant cocycle. Evaluation at the marking, where \(\gamma(t)\ne1\), kills the global cobounding scalar. Evaluation at \(\sigma_\ell\), where the odd-weight character also differs from one, kills \(m\).
\leanref{UpperTheory.\allowbreak{}V.\allowbreak{}Ws.\allowbreak{}Lat.\allowbreak{}oddVan}{Development.\allowbreak{}AutomorphicArithmetic.\allowbreak{}AutomorphicVarieties.\allowbreak{}WeakShimuraTheory.\allowbreak{}AdicLattices.\allowbreak{}OddVan}{55--87}
\leanref{UpperTheory.\allowbreak{}V.\allowbreak{}Ws.\allowbreak{}Lat.\allowbreak{}entVan}{Development.\allowbreak{}AutomorphicArithmetic.\allowbreak{}AutomorphicVarieties.\allowbreak{}WeakShimuraTheory.\allowbreak{}AdicLattices.\allowbreak{}EntVan}{31--64}

For \(s\in D_\ell\) the first column of \(R(s)N\) is \(d_0(s)\) times the first column of \(N\). If \(-\alpha\notin S\), then \(R_{10}=R_{32}=N_{10}=0\). There is no two-step product from \(0\) to \(3\), and the first-column identity gives \(R_{30}=N_{30}(d_0-d_3)\) locally. The preceding argument kills \(R_{30}\) and \(N_{30}\). If \(\beta,\psi\notin S\), then \(R_{12}=R_{13}=0\). There is no two-step product from \(0\) to \(1\), and the same identity gives \(R_{10}=N_{10}(d_0-d_1)\) locally. Thus \(R_{10}=N_{10}=0\). The similitude identity applied to the inverse group element gives \(R_{32}=0\). We have proved
\[
 -\theta\in S\ \Longrightarrow\ -\alpha\in S,\qquad
 -\alpha\in S\ \Longrightarrow\ \beta\in S\ \text{or}\ \psi\in S .
\]
\leanref{UpperTheory.\allowbreak{}V.\allowbreak{}Ws.\allowbreak{}Lat.\allowbreak{}noNeg}{Development.\allowbreak{}AutomorphicArithmetic.\allowbreak{}AutomorphicVarieties.\allowbreak{}WeakShimuraTheory.\allowbreak{}AdicLattices.\allowbreak{}NotTight}{45--101}

At the chosen vertex each of the directions \((1,0),(0,1),(-1,1)\) violates a tight inequality for every sufficiently small positive displacement. Taking dot products with the eight roots gives, respectively,
\[
 S\cap\{-\alpha,\beta,-\theta\}\ne\varnothing,\quad
 S\cap\{-\psi,-\beta,-\theta\}\ne\varnothing,\quad
 S\cap\{-\psi,\alpha,-\beta\}\ne\varnothing.
\]
If \(\beta\notin S\), the first intersection and the two arithmetic implications give \(-\alpha,\psi\in S\). Hence \(a>0>b\), which excludes every weight in the third intersection. Thus \(\beta\in S\) and \(b<a\). This excludes \(-\beta\). If \(-\psi\notin S\), the third intersection gives \(\alpha\in S\), so \(b<a<0\). The second intersection then requires \(-\theta\in S\), which requires \(a+b>0\). Consequently \(-\psi\in S\) and \(a>b\ge0\). The strict sign test excludes \(\alpha,\psi,\theta,-\beta\). This proves \(S\subseteq S_P\).

There is no flag coordinate for \(\beta\), so \(R_{12}\) is nonzero somewhere. Also \(N_{21}=0\), which gives the stated flag support. If \(E(-\psi)\) fails, tightness of \(-\psi\) gives \(N_{20}\ne0\). Row \(2\) of the first-column identity gives \(d_2N_{20}=d_0N_{20}\) on \(D_\ell\). Finally the two-step products for entry \(12\) vanish, since \(R_{02}=R_{13}=0\). A scalar \(m\) as in the statement would kill \(R_{12}\) by the odd-weight argument. This proves the second alternative.
\leanref{UpperTheory.\allowbreak{}V.\allowbreak{}Ws.\allowbreak{}Lat.\allowbreak{}shape}{Development.\allowbreak{}AutomorphicArithmetic.\allowbreak{}AutomorphicVarieties.\allowbreak{}WeakShimuraTheory.\allowbreak{}AdicLattices.\allowbreak{}LowCase}{39--104}
\end{proof}

\subsection*{Proof of \ref{B6.49e}}\label{prf:B6.49e}\phantomsection
\begin{proof}
At a prolongation \(v\) write \(g\) for an element transporting the original place to \(v\). Use the lower unitriangular representative \(N_v\) of the transported flag. The factorisation exists integrally. Modulo the maximal ideal the first two leading minors of \(R(g)N\) are \(d_0(g)\) and \(d_0(g)d_1(g)\), by the support in Lemma~\ref{B6.49d} and \(N_{21}=0\). The third is a unit by the similitude identity. Elimination and reduction therefore commute. The first two columns in the factorisation give
\[
 (N_v)_{21}=0,\qquad
 (N_v)_{20}d_0(g)=R_{20}(g)+d_2(g)N_{20}.
\]
Conjugating the local group transports the flag characters. Their inertia restrictions are the prescribed ones, since the characters come from \(G_\Q\). The excluded inertia acts trivially by the defining inertia conditions of the tower and of the rational representation. The tower marking remains diagonal since it is a power of the original one. This gives the residual point with representation support in \(S_P\) and flag support in \(S_f\).

We use that \(F/\Q\) is normal and \(G_F\) contains a marking \(t_F\) separating the four labels. If \(f=R_{ij}/d_j\) satisfies the cocycle law for \(\gamma=d_i/d_j\) and vanishes on \(G_F\), evaluate it at \(g t_F g^{-1}\in G_F\). The result is \((1-\gamma(t_F))f(g)=0\), so \(f=0\) on \(G_\Q\). On \(S_P\) this applies to entries \(20,31,12\), by matrix multiplication. Thus the first alternative over \(\Q\) retains both weights over \(F\).

In the other alternative \(R_{20}=R_{31}=0\), and the displayed transport identity gives \((N_v)_{20}\ne0\) at every \(v\). The support is as required for \(V\). If \(E(-\alpha)\) failed over \(F\), entry \(10\) would vanish there. It too satisfies the two-term product law, because \(R_{20}=R_{13}=0\); hence it would vanish over \(\Q\). On the original \(D_\ell\), row \(1\) of the first-column identity would give
\[
 R_{12}(s)N_{20}=N_{10}(d_0(s)-d_1(s)).
\]
Since \(d_0=d_2\) there, \(m=N_{10}/N_{20}\) would be the scalar excluded in Lemma~\ref{B6.49d}. This proves \(V\land E(-\alpha)\).
\end{proof}

\subsection*{Proof of \ref{B6.49f}}\label{prf:B6.49f}\phantomsection
\begin{proof}
Write \(O_1\) for the coefficient DVR and \(\pi\) for its uniformiser. The planes \(A,B\) are complementary Lagrangians. Let \(M\) consist of the infinitesimally symplectic maps from \(A\) to \(B\), extended by zero on \(B\). It is a free lattice, its elements have pairwise zero products, and \(G_\Q\) acts by conjugation through \(r\). Let \(P\subseteq M\) consist of the maps strictly lowering the ordinary flag. It is saturated and stable under \(D_\ell\).

Write \(s=\sigma_\ell\). The residual spectra of \(r(s)\) on \(A,B\) are disjoint, so \(s-1\) acting on \(M\) is invertible. The first ordinary line lies in \(A\), and the first ordinary plane is its sum with a line in \(B\): apply the two block spectral projectors of \(r(s)\), whose block characteristic polynomials are coprime over \(O_1\). If \(w\) spans this latter line over the fraction field, the nonzero map \(v\mapsto\langle w,v\rangle w\) belongs to the fraction space of \(P\). It kills the first plane and takes the entire space into that plane. Thus \(\operatorname{rank}P>0\). Complex conjugation \(c\) acts on \(A,B\) by \(-1,1\), respectively: these are its residual signs and the two projectors \((1\pm r(c))/2\) are integral. It acts by \(-1\) on \(M\).
\leanref{UpperTheory.\allowbreak{}V.\allowbreak{}Ws.\allowbreak{}Bl.\allowbreak{}flagVec}{Development.\allowbreak{}AutomorphicArithmetic.\allowbreak{}AutomorphicVarieties.\allowbreak{}WeakShimuraTheory.\allowbreak{}BlockMatrixPlaces.\allowbreak{}FlagVec}{87--176}

Put \(M_n=M/\pi^nM\), \(P_n=P/\pi^nP\). Use unramified cohomology away from \(\ell\) and the image of \(H^1(\Q_\ell,P_n)\) at \(\ell\). The latter map into \(H^1(\Q_\ell,M_n)\) is injective: \(s-1\) is bijective on \(M_n\), on its stable submodule \(P_n\) by finiteness, and on their quotient. Its image has order at least \(|P_n|\) by the local Euler characteristic formula for the finite \(\ell\)-primary module \(P_n\) over \(\Q_\ell\) \cite[I, Thm.~2.8]{milne-2006}. Apply the finite-module Selmer ratio \eqref{eq:p6c5-link-selmer-ratio} over \(\Q\) with the places of ramification and \(\ell\). At \(p\ne\ell\) the unramified group is the group of Frobenius coinvariants of \(M_n^{I_p}\), of the same order as the invariants. The local invariants at \(\ell\) vanish by \(s-1\). The global dual invariants vanish because \(\varepsilon_\ell(s)=1\) and \(s-1\) is surjective; the real invariants vanish because \(\ell\) is odd. Thus there are at least \(|P_n|\) Selmer classes. Every class has a unique representative vanishing at \(c\), since \(c-1=-2\). Each representative vanishes on \(I_p\cap\ker r\) away from \(\ell\), and on \(D_\ell\) differs from a \(P_n\)-valued cocycle by a coboundary.
\leanref{UpperTheory.\allowbreak{}V.\allowbreak{}Ws.\allowbreak{}Bl.\allowbreak{}gwCount}{Development.\allowbreak{}AutomorphicArithmetic.\allowbreak{}AutomorphicVarieties.\allowbreak{}WeakShimuraTheory.\allowbreak{}BlockMatrixPlaces.\allowbreak{}GwCount}{12--76}

Consider all normalised cocycles satisfying the last two conditions. At each depth they form a finite set, by finite-module global cohomological finiteness and finiteness of the module of coboundaries. Their cardinalities are unbounded by the preceding estimate. At every depth there is a cocycle with nonzero reduction modulo \(\pi\). Otherwise at that and all greater depths every cocycle would be divisible by \(\pi\), and division would inject the set at each successive depth into the preceding set. For the local condition in this division, write \(z|_{D_\ell}-dy\in P_{n+1}\). If \(z\) is divisible by \(\pi\), evaluation at \(s\) shows \(y\bmod\pi\in P_1\), since \(s-1\) is invertible on \(M_1/P_1\). Subtract a lift in \(P_{n+1}\) from \(y\), then divide by \(\pi\); saturation preserves the condition. The other conditions divide entrywise. The sets would consequently have bounded cardinality. Taking a branch in the finitely branching tree of these cocycles and their local witnesses gives, by completeness, a nonzero integral cocycle \(z\) and \(y\in M\) with \(z|_{D_\ell}-dy\) valued in \(P\), where \(dy(g)=r(g)y r(g)^{-1}-y\).
\leanref{UpperTheory.\allowbreak{}V.\allowbreak{}Ws.\allowbreak{}Bl.\allowbreak{}cocT}{Development.\allowbreak{}AutomorphicArithmetic.\allowbreak{}AutomorphicVarieties.\allowbreak{}WeakShimuraTheory.\allowbreak{}BlockMatrixPlaces.\allowbreak{}CocT}{11--56}

Replace \(z,y\) by \(\pi z,\pi y\). Put \(r'(g)=(1+z(g))r(g)\) and take the flag \((1-y)u\), where \(u\) was the original flag matrix. The cocycle law and \(M^2=0\) give multiplicativity. The new flag matrix conjugates the local action to the old one multiplied in the ordinary basis by \(1+u^{-1}(z-dy)u\), upper unitriangular there. This preserves its diagonal characters. The new representation preserves the similitude and the prescribed inertia kernels. Its reduction and block diagonal entries agree with the old ones. In particular all characteristic polynomials agree. The action on \(B\) is the similitude twist of the dual action on \(A\); hence it has full matrix span. A complement would be the graph of a map \(Y:A\to B\) and would express \(z\) as a coboundary in \(\operatorname{Hom}(A,B)\). Evaluation at \(c\) gives \(2Y=0\), since \(z(c)=0\). This forces \(Y=0\) and then \(z=0\), a contradiction over the algebraic closure as well.
\leanref{UpperTheory.\allowbreak{}V.\allowbreak{}Ws.\allowbreak{}Bl.\allowbreak{}cocA}{Development.\allowbreak{}AutomorphicArithmetic.\allowbreak{}AutomorphicVarieties.\allowbreak{}WeakShimuraTheory.\allowbreak{}BlockMatrixPlaces.\allowbreak{}TwistMul}{28--62}
\end{proof}

\subsection*{Proof of \ref{B6.49g}}\label{prf:B6.49g}\phantomsection
\begin{proof}
Let \(C\) be the identity component of the Zariski closure of \(r(G)\) in \(\GL_n\), and \(N=r^{-1}(C)\). This is a normal subgroup of finite index. It is closed in the profinite topology since the matrix entries are continuous and the equations for \(C\) are defined over a finite coefficient extension. Hence \(N\) is open. Its image is dense in \(C\): the components are open and closed in the algebraic group. For any open subgroup \(U\le G\), \(r(U\cap N)\) is also dense in \(C\). Indeed, finitely many translates of its closure cover \(C\), and \(C\) is irreducible. Thus an open coset \(\sigma U\) contains elements with image dense in \(r(\sigma)C\).

Write \(A\) for the linear span of \(r(N)\). It is a unital matrix algebra normalized by \(r(G)\). A positive power of \(u\) belongs to \(N\) and still has distinct eigenvalues. Its spectral projectors give all \(E_{ii}\in A\). Consequently an entry \(ij\) occurs in \(A\) precisely when \(E_{ij}\in A\), by multiplying on both sides by projectors. This relation is reflexive and transitive. Group together indices related in both directions. The algebra contains every matrix unit within each class. The span \(J\) of its matrix units between distinct classes is a nilpotent two-sided ideal: a nonzero product follows a strictly ordered chain of classes. The quotient by \(J\) is the product of the full class matrix algebras, so \(J\) is the Jacobson radical and is normalized by \(r(G)\). Its common nullspace is nonzero by nilpotence. Irreducibility makes that nullspace the whole representation. Thus \(J=0\).

Conjugation permutes the central idempotents of the resulting product. It acts transitively on the classes, since a proper union of orbits would span an invariant subspace. All classes therefore have the same size. Singleton classes would make \(r(N)\) commutative. There is one class in ranks two and three. In rank four there is one class or there are two classes of size two.

The linear span of \(r(\sigma)C\) is \(r(\sigma)A\). If \(\sigma\) preserves the classes, this span contains every diagonal matrix unit. Each diagonal coordinate is consequently nonzero somewhere on the coset. Inversion maps the coset to \(r(\sigma^{-1})C\), so the same holds for each inverse diagonal coordinate. If \(\sigma\) exchanges two classes, its coset has zero diagonal blocks and its span contains every off-diagonal block matrix unit. Each required nonvanishing is a nonempty Zariski open condition on the same irreducible coset. Their finite intersection meets \(r(\sigma(U\cap N))\).
\end{proof}

\subsection*{Proof of \ref{B6.49h}}\label{prf:B6.49h}\phantomsection
\begin{proof}
For (1), take \(h\) past the threshold of Lemma~\ref{B6.48ba}. The resulting normalized cuspidal eigenform of weight \(h\), level \(\Gamma_1(Dw)\) and character \(\phi\theta\) is new because its nebentype has primitive conductor \(Dw\). Attach its cuspidal representation of \(\GL_2(\mathbf A_\Q)\) by Proposition~\ref{B6.06d}(3) and its rank-two Galois representation by Proposition~\ref{B6.06e}. Twist from determinant \(D_h^{-1}\) by \(D_h\) to obtain a representation \(\rho\) of determinant \(D_h\). For a two-by-two matrix \(R\) of determinant \(D_h^{-1}\), \(\operatorname{tr}(D_h R)=\operatorname{tr}(R^{-1})\). Thus the good-prime congruences in Lemma~\ref{B6.48ba}, density, and \(\bar\theta=1\) give residual polynomial \((X-1)(X-\bar\Phi_h)\) on \(G_\Q\), with \(\Phi_h=\varepsilon_\ell^{h-1}\phi\). Local compatibility in Proposition~\ref{B6.06e}(2) gives spherical support outside \(Dw\), and at \(p\mid Dw\) the primitive nebentype of conductor \(Dw\) gives unlinked principal-series support with inertia type \(1,\phi\theta\) away from \(\ell\). At \(\ell\) the unit root reduces to one by \(a_\ell\equiv1\) and the positive valuation of the other root. The unramified flag character in the initial normalisation therefore reduces to one on the decomposition group. The twist gives the ordered inertia weights \((h-1,0)\) with residual labels \((\bar\Phi_h,1)\) on the decomposition group. Cuspidality gives irreducibility in the rank-two attachment.
\leanref{UpperTheory.\allowbreak{}V.\allowbreak{}Z.\allowbreak{}Win.\allowbreak{}Birth.\allowbreak{}propEGP\_\allowbreak{}of\_\allowbreak{}printsH}{Development.\allowbreak{}AutomorphicArithmetic.\allowbreak{}AutomorphicVarieties.\allowbreak{}ModularityArgument.\allowbreak{}ThresholdArguments.\allowbreak{}DenseThresholdCases.\allowbreak{}BirthEG}{231--297}
\leanref{UpperTheory.\allowbreak{}V.\allowbreak{}Z.\allowbreak{}Sym.\allowbreak{}W.\allowbreak{}prop\_\allowbreak{}eg\_\allowbreak{}complex\_\allowbreak{}of\_\allowbreak{}printed\_\allowbreak{}hida}{Development.\allowbreak{}AutomorphicArithmetic.\allowbreak{}AutomorphicVarieties.\allowbreak{}ModularityArgument.\allowbreak{}ThresholdArguments.\allowbreak{}DenseThresholdCases.\allowbreak{}WEGCore}{235--303}

For restriction, let \(H\) be open normal in \(G_\Q\). A positive power of an inertia element of infinite cyclotomic order lies in \(H\) and has distinct eigenvalues because \(h>1\). If \(\rho|_H\) fixes a line, its translates are among these two eigenlines. Irreducibility over \(G_\Q\) forces a transitive action on both. Every element interchanging them has trace zero and hence \(\bar\Phi_h=-1\). Multiplying an interchanging element by any element fixing both lines shows that \(\bar\Phi_h=1\) on the stabiliser, which contains \(H\). Thus a nontrivial value of \(\bar\Phi_h\) on \(H\) gives irreducibility there, and a value with \(\bar\Phi_h^2\ne1\) anywhere gives irreducibility on every open subgroup by taking normal cores.
\leanref{UpperTheory.\allowbreak{}V.\allowbreak{}Sd.\allowbreak{}sI}{Development.\allowbreak{}AutomorphicArithmetic.\allowbreak{}AutomorphicVarieties.\allowbreak{}IdeleGenerators.\allowbreak{}HolomorphicIrreducibility.\allowbreak{}NormOneOfFin}{23--56}

For (2), take an algebraic Hecke character of \(K\) of type \((-1,0)\) by extending the character on principal ideles of this type with the odd quadratic character \(\eta_K\) of discriminant \(D_K\) times the inverse norm on rational ideles. At finite places away from \(\ell\) the restrictions to units can be taken of finite order dividing the prime-to-\(\ell\) part of the relevant residue-unit orders, since the discriminant \(D_K\) is prime to \(\ell\), and unramified outside \(D_K\). On transfer from \(G_\Q\) for the \(\ell\)-adic realization \(\kappa_0\), the rational-idelic restriction is \(\eta_K\varepsilon_\ell^{-1}\), whose restriction to \(G_K\) gives \(\kappa_0(g)\kappa_0(cgc^{-1})=\varepsilon_\ell(g)^{-1}\).
\leanref{UpperTheory.\allowbreak{}V.\allowbreak{}Cu2.\allowbreak{}kapEx\_\allowbreak{}alt}{Development.\allowbreak{}AutomorphicArithmetic.\allowbreak{}AutomorphicVarieties.\allowbreak{}ModularityArgument.\allowbreak{}ThresholdArguments.\allowbreak{}DenseThresholdCases.\allowbreak{}MPKapEx}{21--117}
\leanref{UpperTheory.\allowbreak{}V.\allowbreak{}Sd.\allowbreak{}twist\_\allowbreak{}of\_\allowbreak{}printed\_\allowbreak{}results}{Development.\allowbreak{}AutomorphicArithmetic.\allowbreak{}AutomorphicVarieties.\allowbreak{}ModularityArgument.\allowbreak{}ThresholdArguments.\allowbreak{}DenseThresholdCases.\allowbreak{}TTwist}{111--178}
\end{proof}

\subsection*{Proof of \ref{B6.49i}}\label{prf:B6.49i}\phantomsection
\begin{proof}
For (1), let \(c\in\overline k^\times\) be the single residual root of \(\det(X-\overline{\rho(s_0)})\), nonzero since \(\rho(s_0)\in\GL_m(A)\). The characteristic polynomials \(P_1(T)=\det(T-\chi(s_0)\rho(s_0))\) and \(P_0(T)=\det(T-\rho(s_0))\) have coefficients in \(A\) and reductions \((T-\bar\chi(s_0)c)^m\) and \((T-c)^m\) over \(\overline k\). Since \(c\ne0\) and \(\bar\chi(s_0)\ne1\), these reductions are coprime over \(k\), hence \(P_1\) and \(P_0\) are coprime in \(F[T]\). Choose \(U,V\in F[T]\) with \(U(T)P_1(T)+V(T)P_0(T)=1\). Any \(X\in M_m(F)\) satisfying \(\chi(s_0)\rho(s_0)X=X\rho(s_0)\) satisfies \(V(\chi(s_0)\rho(s_0))P_0(\chi(s_0)\rho(s_0))X=XV(\rho(s_0))P_0(\rho(s_0))\), so
\[
 X=U(\chi(s_0)\rho(s_0))P_1(\chi(s_0)\rho(s_0))X+XV(\rho(s_0))P_0(\rho(s_0))=0
\]
by the Cayley--Hamilton theorem.

For (2), the line characters on \(P\) have unit values, since for \(p\in P\) both \(\rho(p)\) and \(\rho(p)^{-1}=\rho(p^{-1})\) belong to \(\GL_m(A)\) and each eigenvalue and its inverse are roots of monic polynomials over \(A\). Reduce them using the fixed extension of the valuation. For \(p\in P\) and \(g\in H\), every residual eigenvalue of the commutator \([p,g]=pgp^{-1}g^{-1}\) is one, since \(\bar\chi_i([p,g])=1\) for every \(i\). On a line \(L_a\) sent by \(g\) to \(L_b=gL_a\), the eigenvalue of \(\rho([p,g])\) on \(L_b\) is the ratio of the character of \(L_b\) at \(p\) to that of \(L_a\) at \(p\), so the two reduced values at \(p\) coincide. By transitivity they all coincide at a common reduced value \(\bar\alpha(p)\). Over \(\overline k\) the characteristic polynomial of \(\overline{\rho(p)}\) is \((X-\bar\alpha(p))^m=\prod_{i=1}^m(X-\bar\chi_i(p))\), giving \(\bar\chi_1(p)=\cdots=\bar\chi_m(p)=\bar\alpha(p)\).
\leanref{UpperTheory.\allowbreak{}V.\allowbreak{}Ws.\allowbreak{}P1.\allowbreak{}irrL}{Development.\allowbreak{}AutomorphicArithmetic.\allowbreak{}AutomorphicVarieties.\allowbreak{}WeakShimuraTheory.\allowbreak{}ProjectiveLineFlags.\allowbreak{}TraceOfCp}{17--75}
\leanref{UpperTheory.\allowbreak{}V.\allowbreak{}Ws.\allowbreak{}h4}{Development.\allowbreak{}AutomorphicArithmetic.\allowbreak{}AutomorphicVarieties.\allowbreak{}WeakShimuraTheory.\allowbreak{}CompactIntegralFields.\allowbreak{}H4}{19--118}
\leanref{UpperTheory.\allowbreak{}V.\allowbreak{}Pr.\allowbreak{}irrNormal}{Development.\allowbreak{}AutomorphicArithmetic.\allowbreak{}AutomorphicVarieties.\allowbreak{}CurveDVRSpecialization.\allowbreak{}IrrNormal}{62--180}
\leanref{UpperTheory.\allowbreak{}V.\allowbreak{}Pr.\allowbreak{}nice\_\allowbreak{}specialization\_\allowbreak{}twist\_\allowbreak{}satisfies\_\allowbreak{}third\_\allowbreak{}condition}{Development.\allowbreak{}AutomorphicArithmetic.\allowbreak{}AutomorphicVarieties.\allowbreak{}CurveDVRSpecialization.\allowbreak{}NiceSpecializationTwistSatisfiesThirdCondition}{164--214}
\end{proof}

\subsection*{Proof of \ref{B6.49}}\label{prf:B6.49}\phantomsection
\begin{proof}
Use the host inequalities and cases of \ref{B6.28}. Write \(e=k\), \(\psi=\vartheta\) as a Galois character, and \(\Psi=\varepsilon_\ell^e\psi\). Take the auxiliary odd character \(\xi\), similitude \(\nu\) and level primes \(S_D\) of Definition~\ref{B6.31b}. A tower uses an imaginary quadratic \(K\), a totally real abelian \(F\) of even degree, \(L=FK\), a marking \(g_0\in G_F\), and prolongations over \(\ell\) inducing the chosen place of \(K\). Write \(d_i\) for the ordered residual diagonal characters. At auxiliary rational set \(W\) the frame kills inertia outside \(W\cup\{\ell\}\) on the kernel of the residual diagonal. At \(\ell\) it records the decomposition group, inertia, the preimage of the \(\ell\)-primary abelian torsion and the flag. Auxiliary split primes are trivial on \(F\), the residual diagonal and \(\mu_\ell\). On a deep layer \(\psi|_{G_F}\) has order prime to \(\ell\).

For a sign \(s=\pm1\) and a finite list \(S_0\), choose a fresh prime \(p\equiv s\bmod 8\prod_{q\in S_0}q\). Its quadratic character has sign \(s\) and splits at the listed primes by quadratic reciprocity and the supplement at \(2\). In II take an even character \(\phi\) split at \(S_D\), then two odd characters \(\eta,\eta'\), adding each previous prime to the splitting list. Let \(\kappa_b\) be the \(\ell\)-primary part of \(\psi\) with its local factor at \(\ell\) removed. Set
\[
 G_F=\{g:\bar\Psi(g)=\phi(g),\ \eta(g)=\eta'(g),\ \kappa_b(g)=1\},
 \qquad G_K=\ker\eta,\qquad G_{k_0}=\ker\phi .
\]
Here \(k_0\) is real quadratic and split at \(2,\ell\). The equality of the odd characters puts conjugation in \(G_F\) and gives even degree. At the conductor of \(\eta\), inertia in \(G_F\) kills \(\eta\) because \(\eta'\) is unramified. Thus \(FK/F\) is unramified at finite places. The root of weight \(\psi\) in \ref{B6.48b} cuts out \(Fk_0/F\). The marking has \(\bar\Psi=-1,\ \bar\xi^2\ne1\) and residual cyclotomic value one. The fourth-power test \(\bar\xi^4\ne1\) holds on elements of the decomposition groups at \(\ell\) and of inertia at \(2\).
\leanref{UpperTheory.\allowbreak{}V.\allowbreak{}Ws.\allowbreak{}quadEx}{Development.\allowbreak{}AutomorphicArithmetic.\allowbreak{}AutomorphicVarieties.\allowbreak{}WeakShimuraTheory.\allowbreak{}FrobeniusFields.\allowbreak{}QuadEx}{75--119}
\leanref{UpperTheory.\allowbreak{}V.\allowbreak{}Ws.\allowbreak{}exists\_\allowbreak{}tame\_\allowbreak{}tower\_\allowbreak{}ok\_\allowbreak{}two}{Development.\allowbreak{}AutomorphicArithmetic.\allowbreak{}AutomorphicVarieties.\allowbreak{}WeakShimuraTheory.\allowbreak{}CompactIntegralFields.\allowbreak{}ExistsTameTowerOkTwo}{22--60}

In III choose an odd quadratic character \(\lambda_{K'}\) split at \(S_D\) and ramified at a fresh prime beyond the conductors. Choose odd \(\eta,\eta'\) by the same rule, split also at this prime, with \(\eta'\) split at the conductor of \(\eta\). Set
\[
 G_F=\{g:\xi(g)=\lambda_{K'}(g),\ \bar\Psi(g)=\bar\psi(g),
 \ \eta(g)=\eta'(g)\},\qquad G_K=\ker\eta .
\]
As in II, \(F\) is totally real of even degree and \(FK/F\) is unramified at finite places. Write \(\omega\) for the Teichmüller lift of \(\bar\varepsilon_\ell\). Here \(\ell\ge5,\ \ell-1\mid2e\), and \(\psi\) has odd order prime to \(\ell\). For the marking multiply an inertia element at \(2\) with \(\xi=-1\), one changing the sign of \(\lambda_{K'}\) at its fresh conductor, and a square detecting the odd part of \(\psi\). The fresh inertia kills the old characters. Their product lies in \(G_F\cap G_K\) with \(\omega=1\). Its relative labels are \(1,-1,z,-z\) up to inversion, with \(z\ne1\) of odd order. Thus a quadratic permutation twist through \(\omega\) fixes every index. The ratio \(-z\) has nontrivial first four powers. At \(\ell\), squares of inertia lie in \(G_F\), and some have nontrivial residual cyclotomic value. These are the residual separation and cyclotomic tests.
\leanref{UpperTheory.\allowbreak{}V.\allowbreak{}Cu.\allowbreak{}existsTensorSeedTower\_\allowbreak{}holds}{Development.\allowbreak{}AutomorphicArithmetic.\allowbreak{}AutomorphicVarieties.\allowbreak{}ModularityArgument.\allowbreak{}CuspidalTwists.\allowbreak{}ComplexConjugationsConjugateCaseThree}{538--721}

In II let \(\omega_0\) be the finite Hecke character of \(k_0\) corresponding to \(\xi|_{G_{k_0}}\). It is odd at both real places and unramified over \(\ell\). The fourth-power test makes it ramified at a place of norm two and gives \(\bar\omega_0(\varpi_v),\bar\omega_0(\varpi_{v'})\ne1\) at the two places \(v,v'\mid\ell\). Let \(E\) be the parallel weight-one Hilbert Eisenstein series for \(1,\omega_0\). Write \(V_w\) for the normalized level-raising operator which sends the coefficient at \(\mathfrak a\) to the coefficient at \(\mathfrak a/\mathfrak p_w\), interpreted as zero for a nonintegral index. Take \((1-V_v)(1-\omega_0(\varpi_{v'})V_{v'})E\). The two roots at each of these unramified places are \(1,\omega_0(\varpi_w)\), so the retained roots are \(\omega_0(\varpi_v)\) and \(1\), respectively. The constant coefficient of \(E\) at infinity is the same on the ideal-class components. In weight one the normalized \(V_v\) translates the component of this constant coefficient with multiplier one. Thus \(1-V_v\) kills it. The coefficient at the unit ideal remains one.
\leanref{UpperTheory.\allowbreak{}V.\allowbreak{}Sd.\allowbreak{}emix}{Development.\allowbreak{}AutomorphicArithmetic.\allowbreak{}AutomorphicVarieties.\allowbreak{}IdeleGenerators.\allowbreak{}HolomorphicIrreducibility.\allowbreak{}NotDvdMul}{270--300}

Choose the multiplier at level one. Since \(\ell\) splits in the real quadratic field \(k_0\), its total Hilbert Hasse invariant has parallel weight \(\ell-1\) and expansion one at every cusp \cite[Thm.~2.1(1),(2)]{goren-2001}. At sufficiently small auxiliary tame level, a sufficiently high power lifts to an integral characteristic-zero form by \cite[Lems.~2.2--2.3]{deo-dimitrov-wiese-2024}. To obtain level one, take a normal full tame level contained in an auxiliary level satisfying condition~(20) of that reference, and pull back the lift. Take the product of its translates by the finite quotient of the full level-one compact group by this normal subgroup. The product is invariant at level one and has trivial character. Its reduction is the corresponding power of the level-one Hasse invariant, since that section is fixed by every translate. Thus its expansion is congruent to one at every cusp. Taking a further power gives parallel weight \(a\) divisible by \(\ell-1\) with \(a+1\ge4\). Multiply the stabilized series by this form.
\leanref{GlobalTheory.\allowbreak{}HilbertThetaSeries.\allowbreak{}hilbertOne}{Development.\allowbreak{}AutomorphicArithmetic.\allowbreak{}AutomorphicVarieties.\allowbreak{}GaussSumFiltrations.\allowbreak{}HasQOfSmallSubOne}{33--76} Use level exponent one above \(\ell\) and the conductor exponents of \(\omega_0\) elsewhere. The product has a unit coefficient at the unit ideal and preserves the residual system. At good places the weight change preserves the Hecke recurrence modulo the maximal ideal because \((\mathrm N\mathfrak p)^a\equiv1\bmod\ell\).

For the Hilbert lifting put \(h=a+1\). Transport the Fourier expansion of the finite-dimensional holomorphic space by \(\iota^{-1}\) into the space of coefficient functions on the nonzero integral ideals of \(k_0\). At index \(\mathfrak a\), \(T_{\mathfrak p}\) takes the coefficient at \(\mathfrak p\mathfrak a\) and, if \(\mathfrak p\mid\mathfrak a\), adds the coefficient at \(\mathfrak a/\mathfrak p\) times \(\chi(\mathfrak p)(\mathrm N\mathfrak p)^{h-1}\), with \(\chi=\omega_0\) extended by zero at the level primes. At level primes this is the \(U_{\mathfrak p}\) formula. All these multipliers are integral, so the commuting operators preserve both the coefficientwise integrality condition and the submodule of functions with every coefficient in the maximal ideal. The product above gives a nonzero simultaneous residual eigenvector in this coefficient space. Apply the Deligne--Serre lifting lemma \cite[Lem.~6.11]{deligne-serre-1974} to this operator family. The Fourier map is injective on the holomorphic space in weight \(h\), so the lifted coefficient eigenvector pulls back to a holomorphic eigenform. Its \(U_v,U_{v'}\) eigenvalues reduce to \(\bar\omega_0(\varpi_v),1\), respectively.
\leanref{UpperTheory.\allowbreak{}V.\allowbreak{}Sd.\allowbreak{}holLift\_\allowbreak{}alt}{Development.\allowbreak{}AutomorphicArithmetic.\allowbreak{}AutomorphicVarieties.\allowbreak{}ModularityArgument.\allowbreak{}ThresholdArguments.\allowbreak{}DenseThresholdCases.\allowbreak{}THolLift}{129--207}

If the lift is Eisenstein with inducing characters \(\eta_1,\eta_2\), the good eigenvalues and the central character imply
\((\bar\eta_1(g)-1)(\bar\eta_2(g)-1)=0\) by reciprocity and Frobenius density. The two kernels cover the group, so one character is trivial. In the ordinary Eisenstein lift both unit \(U\)-roots come from the first inducing character, since the second has the positive norm power. Its residual values at \(v,v'\) are \((\bar\omega_0(\varpi_v),1)\), and the product character gives \((1,\bar\omega_0(\varpi_{v'}))\) for the second. Both characters are nontrivial. Thus the lift is cuspidal.

Take \(h\ge4\) and let \(d_0\) be the dimension of the parallel-weight cusp space at the chosen level and central character. Put \(t_0=2(\ell-1)\ell^{d_0}\). Change the weight at the real embedding assigned to \(v\) from \(h\) to \(h+2t_0\), with the factor \(\ell^{t_0}\) in the normalisation of the operator at \(v\). For a product containing at least \(r\ge1\) ordinary operators at each place over \(\ell\), the normalised traces satisfy
\begin{equation}\tag{6.49.2}\label{eq:p6c5-mixed-traces}
 \operatorname{tr}_{h,h}\equiv
 \operatorname{tr}_{h+2t_0,h}\pmod{\ell^{\min(r,d_0)}} .
\end{equation}
Apply Lemma~\ref{B6.49b} to the two cusp spaces, with \(F=k_0\), \(N=d_0\), \(t=t_0\), and \(\omega=\omega_0\). The field is real quadratic and \(\ell\) splits. The finite character has the parity of \(h\) at each real place, its conductor divides the level, and the norm-two place has positive level exponent. Both exponents above \(\ell\) are one. The chosen \(t_0\) is divisible by \((\ell-1)\ell^{d_0}\). This gives \eqref{eq:p6c5-mixed-traces} for every indicated Hecke word. The parallel cusp eigenform constructed above has integral Hecke eigenvalues with the prescribed reductions, including the units \(\bar\omega_0(\varpi_v)\) and \(1\). The second assertion of the lemma gives a mixed-weight cusp eigenform with both ordinary refinements in the displayed normalization.
\leanref{UpperTheory.\allowbreak{}V.\allowbreak{}Z.\allowbreak{}Bc.\allowbreak{}propMHolZ\_\allowbreak{}alt}{Development.\allowbreak{}AutomorphicArithmetic.\allowbreak{}AutomorphicVarieties.\allowbreak{}ModularityArgument.\allowbreak{}ThresholdArguments.\allowbreak{}DenseThresholdCases.\allowbreak{}MPropMHol}{108--223}
\leanref{UpperTheory.\allowbreak{}Esk.\allowbreak{}pf\_\allowbreak{}TrCong}{PositivityTheory.\allowbreak{}AutomorphicForms.\allowbreak{}PublishedInputs.\allowbreak{}JacquetLanglands.\allowbreak{}TorusKernels.\allowbreak{}PfTrCong}{101--129}
\leanref{UpperTheory.\allowbreak{}Esk.\allowbreak{}pf\_\allowbreak{}HolWC}{PositivityTheory.\allowbreak{}AutomorphicForms.\allowbreak{}PublishedInputs.\allowbreak{}JacquetLanglands.\allowbreak{}TorusKernels.\allowbreak{}PfHolWC}{174--305}

Attach the rank-two representation by Propositions~\ref{B6.06d}(1) and~\ref{B6.06e} and twist by \(\varepsilon_\ell^a\) in that normalisation. The resulting \(r\) on \(G_{k_0}\) has determinant \(\varepsilon_\ell^a\xi\), residual polynomial \((X-1)(X-\bar\xi)\) and weights \((a+t_0,-t_0)\), \((a,0)\) at the two places. The residual orders are \((\bar\xi,1)\), \((1,\bar\xi)\). The squared marking splits the lattice into its two eigenlines. Scale one line so the upper entries are integral with at least one unit. Then the opposite entries lie in the maximal ideal by the residual products. Induce from \(G_{k_0}\). The summands on \(G_{k_0}\) are \(r,r^{g_0}\), and the other coset exchanges them with marked square \(r(g_0^2)\). The weights distinguish the summands on open subgroups. The marked fourth-power test and the incompatible ordinary orders give irreducibility. The similitude is \(\mu=\varepsilon_\ell^a\xi\phi\). On \(G_F\), twist by the unique square root of \(\nu/\mu\) congruent to one, which is a character by Hensel uniqueness. Base change to the compositum with \(k_0\) and quadratic automorphic induction give the attachment, since the conjugate infinity types are distinct \cite{arthur-clozel-1989}. Perform the lattice scaling on \(G_{Fk_0}\): the restriction there is irreducible, the squared marking has distinct residual eigenvalues, and the upper entries attain their maximum absolute value by compactness. Make this maximum one. Write the two conjugate residual matrices there as \(\left(\begin{smallmatrix}1&x\\0&c\end{smallmatrix}\right)\) and \(\left(\begin{smallmatrix}1&y\\0&c\end{smallmatrix}\right)\), where \(c=\bar\xi(g)\). Let \(a,b\) be the two chosen square roots of the eigenvalues of the squared marking, in the order of these two diagonal positions in the rank-two lattice, and put \(\gamma=\bar b/\bar a\). In the marked induction basis the matrix is
\[
 \begin{pmatrix}c&0&0&0\\ D&1&-S&0\\0&0&c&0\\ S&0&-D&1\end{pmatrix},\qquad D=(\gamma x-y)/2,\quad S=(\gamma x+y)/2.
\]
Indeed, on each of the two pairs of exchanged lines take the sum and difference eigenvectors of the marking. In the upper entry of the first summand the change of scale between the pairs is \(b/a\), and in the second summand it is one. Adding and subtracting the resulting entries gives the displayed matrix. The marking itself is diagonal with entries \(b,-a,-b,a\). Multiplication by this diagonal covers the other coset and preserves the support. The common scalar twist reduces to one. Choose \(g\in G_{Fk_0}\) with \(x\ne0\) by the scaling. Either \(D\ne0\), giving entry \(10\), or \(D=0\) and \(S=\gamma x\ne0\), giving entries \(12,30\). Their respective weights are \(-\alpha,\beta,-\theta\).

For the local flags write the two ordinary lines as \((u,1)\) and \((1,w)\), interchanging the two summands where the ordinary orders are reversed. Here \(u\) is integral and \(w\) is in the maximal ideal: the line of residual character \(1\) is the first coordinate line for an upper triangular residual matrix with distinct diagonal characters on the local group, and the line of the other character has unit second coordinate. In the integral sum-and-difference basis write the lower flag matrix as
\[
 \begin{pmatrix}1&0&0&0\\p&1&0&0\\q&s&1&0\\r&t&-p&1\end{pmatrix}.
\]
With \(c_0=b/a\), its entries \((p,q,r,s,t)\) in the two orders are respectively
\[
 \left(c_0u,-1,c_0u,\frac{-2w}{1+c_0uw},\frac{-(1-c_0uw)}{1+c_0uw}\right),\qquad
 \left(-u,1,u,\frac{-2w}{c_0+uw},\frac{c_0-uw}{c_0+uw}\right).
\]
These are obtained by putting the two ordinary eigenvectors in the first two flag spaces and making the flag matrix lower unitriangular. Direct multiplication gives \(t+ps=q\), the symplectic flag equation. Both denominators are units. Thus \(n_{21}=0\) after reduction and \(n_{20}=\pm1\) at every place. This proves \(V\) with the asserted alternative.
\leanref{UpperTheory.\allowbreak{}exists\_\allowbreak{}ribet\_\allowbreak{}lattice}{Development.\allowbreak{}AutomorphicArithmetic.\allowbreak{}AutomorphicVarieties.\allowbreak{}BuildingsResidueChars.\allowbreak{}ExistsRibetLattice}{7--41}
\leanref{UpperTheory.\allowbreak{}mix\_\allowbreak{}flag\_\allowbreak{}exists}{Development.\allowbreak{}AutomorphicArithmetic.\allowbreak{}AutomorphicVarieties.\allowbreak{}BuildingsResidueChars.\allowbreak{}MixFlagExists}{116--144}

\leanref{UpperTheory.\allowbreak{}Esk.\allowbreak{}pf\_\allowbreak{}PropMV\_\allowbreak{}sv\_\allowbreak{}alt}{Development.\allowbreak{}AutomorphicArithmetic.\allowbreak{}AutomorphicVarieties.\allowbreak{}ModularityArgument.\allowbreak{}ThresholdArguments.\allowbreak{}DenseThresholdCases.\allowbreak{}SVRPropMV}{21--152}
\leanref{UpperTheory.\allowbreak{}V.\allowbreak{}Z.\allowbreak{}Win.\allowbreak{}Birth.\allowbreak{}mix\_\allowbreak{}birth\_\allowbreak{}of\_\allowbreak{}prints}{Development.\allowbreak{}AutomorphicArithmetic.\allowbreak{}AutomorphicVarieties.\allowbreak{}ModularityArgument.\allowbreak{}ThresholdArguments.\allowbreak{}DenseThresholdCases.\allowbreak{}BirthMix}{17--120}
\leanref{UpperTheory.\allowbreak{}V.\allowbreak{}Z.\allowbreak{}Bc.\allowbreak{}seedII\_\allowbreak{}of\_\allowbreak{}p\_\allowbreak{}alt}{Development.\allowbreak{}AutomorphicArithmetic.\allowbreak{}AutomorphicVarieties.\allowbreak{}ModularityArgument.\allowbreak{}ThresholdArguments.\allowbreak{}DenseThresholdCases.\allowbreak{}BirthSeedII}{120--148}

In III choose an algebraic Hecke character of \(K'\) of type \((-1,0)\) by Lemma~\ref{B6.49h}(2), with \(\ell\)-adic character \(\kappa_0\); here \(\ell\nmid D_{K'}\) since the fresh conductor prime of \(\lambda_{K'}\) avoids \(\ell\) and is congruent to \(-1\bmod 8\). Choose \(b\ge1\) divisible by the residual image exponent and tame inertia exponents of \(\kappa_0^2\varepsilon_\ell\). Then \(\lambda_0=(\kappa_0^2\varepsilon_\ell)^b\) reduces to one and kills inertia away from \(\ell\). Its type is \((-b,b)\) and conjugation inverts it. Choose \(w\) split in the conductor, residual and cyclotomic fields by Chebotarev \cite[\S1A]{thorner-2019}, take a character \(\theta\) of order \(\ell\) at \(w\), and apply Lemma~\ref{B6.49h}(1) to the odd primitive character \(\phi\) inducing \(\xi\psi^{\pm1}\), with sign given by the order of the labels, at a sufficiently large weight \(h\ge18b+3\) with \(h-1\equiv e\bmod\ell-1\). The resulting ordinary representation \(\rho\) has residual pair \((\bar\Phi_h,1)\) for \(\Phi_h=\varepsilon_\ell^{h-1}\xi\psi^{\pm1}\), with \(\bar\Phi_h^2\ne1\) at the marking, so Lemma~\ref{B6.49h}(1) gives irreducibility on every open subgroup of \(G_\Q\). Scale the eigenlines in the reverse order to II so the lower entry is integral and somewhere a unit. The ordinary line has chart \((1,x_v)\) with \(x_v\) integral by the local separation test.

Choose \(1\le j\le9\) so \(2bj\) avoids the at most eight differences of the two infinitesimal weights at an embedding of \(K'\) and the two at its conjugate, in both orders. Put \(m=bj\), \(\lambda=\lambda_0^j\), \(\tau=\operatorname{Ind}_{G_{K'}}^{G_\Q}\lambda\). In the sum-and-difference basis \(f_0,f_1\), the two diagonal entries on \(G_{K'}\) are \((\lambda+\lambda^{-1})/2\) and the other two are \((\lambda-\lambda^{-1})/2\). On the other coset use the matrix at \(gg_0^{-1}\) times \(\operatorname{diag}(1,-1)\), since \(\lambda(g_0^2)=1\). The matrices preserve \(\operatorname{diag}(1,-1)\), with ordinary lines \((1,s_v)\), \(s_v=\pm1\), of weights \(m,-m\).

Use \(\rho\otimes\tau\) in basis \((e_0f_0,e_0f_1,-e_1f_1,e_1f_0)\), where \(e_if_j=e_i\otimes f_j\) and \(e_0,e_1\) are the marked basis of \(\rho\). The tensor of the alternating and symmetric forms is symplectic. Let \(u\equiv1\) be the square root of \(\varepsilon_\ell/\omega\). For \(\xi\psi\) twist by \(\delta=u^{e+2-h}\). For \(\xi\psi^{-1}\) take \(\delta=u^{e+2-h}\omega^e\psi\). Then \(\delta^2\Phi_h=\nu\) by \(\omega^{h-1}=\omega^e\). The tower enlargement kills the character of conductor \(w\) in \(\det\rho\). The local flag is obtained by applying \(\left(\begin{smallmatrix}1&0\\x_v&1\end{smallmatrix}\right)\otimes1\) to the lower unitriangular matrix with entries \(s_v\) at \(10\), \(-s_v\) at \(32\), and zero at the other off-diagonal positions. Before the common twist the weights are
\[
 (h-1+m,h-1-m,m,-m).
\]
Their successive gaps are at least two by \(h\ge18b+3\) and \(m\le9b\).

Over \(LK'\) the constituents are \(\rho\lambda,\rho\lambda^{-1}\). They have distinct infinity types by the choice of \(j\). Solvable base change and quadratic induction give the cuspidal attachment over \(L\) \cite{arthur-clozel-1989}. An isomorphism of the two constituents on the cyclotomic kernel would give \(\lambda^4=1\) there by determinants. The opposite nonzero weights at the two split places give \(\lambda\) infinite image on that kernel, by taking the quotient of two inertia elements of equal cyclotomic value. The other coset exchanges the constituents. This yields the irreducible symplectic seed. Its residual frame initially allows \(w\). Write the residual rank-two factors as \(\left(\begin{smallmatrix}A&0\\ y&1\end{smallmatrix}\right)\) and \(\operatorname{diag}(1,e)\), where \(e\) is the sign of \(K'\). The second formula follows from \(\bar\lambda=1\) and the action on the two cosets. Omitting the common nonzero scalar, the tensor in the displayed basis is
\[
 \begin{pmatrix} A&0&0&0\\0&Ae&0&0\\0&-ey&e&0\\y&0&0&1\end{pmatrix}.
\]
The local lower flag matrix has \(10\) entry \(s_v=\pm1\), so this support gives \(H\). To see that both entries survive restriction, use \(y(gh)=y(g)A(h)+y(h)\). The lattice has \(y(g)\ne0\) for some \(g\in G_\Q\). At the marking \(t\in G_F\) of the enlarged tower, \(A(t)\ne1\) by distinctness of labels \(0,3\). If \(y(t)=0\), then \(y(gtg^{-1})=y(g)(A(t)-1)/A(g)\ne0\), and \(gtg^{-1}\in G_F\) by normality. Thus one element of \(G_F\) has both \(21\) and \(30\) entries nonzero. These give \(E(-\beta)\) and \(E(-\theta)\).
\leanref{UpperTheory.\allowbreak{}V.\allowbreak{}Cu.\allowbreak{}entOf4}{Development.\allowbreak{}AutomorphicArithmetic.\allowbreak{}AutomorphicVarieties.\allowbreak{}ModularityArgument.\allowbreak{}CuspidalTwists.\allowbreak{}ConjRes}{36--89}

\leanref{UpperTheory.\allowbreak{}V.\allowbreak{}Cu2.\allowbreak{}kapEx\_\allowbreak{}alt}{Development.\allowbreak{}AutomorphicArithmetic.\allowbreak{}AutomorphicVarieties.\allowbreak{}ModularityArgument.\allowbreak{}ThresholdArguments.\allowbreak{}DenseThresholdCases.\allowbreak{}MPKapEx}{21--117}
\leanref{UpperTheory.\allowbreak{}V.\allowbreak{}Cu.\allowbreak{}twEx4}{Development.\allowbreak{}AutomorphicArithmetic.\allowbreak{}AutomorphicVarieties.\allowbreak{}ModularityArgument.\allowbreak{}CuspidalTwists.\allowbreak{}EqOfRd}{68--172}
\leanref{UpperTheory.\allowbreak{}V.\allowbreak{}Cu.\allowbreak{}Dat.\allowbreak{}Q}{Development.\allowbreak{}AutomorphicArithmetic.\allowbreak{}AutomorphicVarieties.\allowbreak{}ModularityArgument.\allowbreak{}CuspidalTwists.\allowbreak{}OverOf4}{51--53}
\leanref{UpperTheory.\allowbreak{}V.\allowbreak{}Cu.\allowbreak{}rank\_\allowbreak{}four\_\allowbreak{}automorphy\_\allowbreak{}of\_\allowbreak{}tensor\_\allowbreak{}of\_\allowbreak{}base\_\allowbreak{}change\_\allowbreak{}of\_\allowbreak{}induction}{Development.\allowbreak{}AutomorphicArithmetic.\allowbreak{}AutomorphicVarieties.\allowbreak{}ModularityArgument.\allowbreak{}ThresholdArguments.\allowbreak{}DenseThresholdCases.\allowbreak{}AutomorphicLiftGsp4OfInputs}{69--180}
\leanref{UpperTheory.\allowbreak{}V.\allowbreak{}Cu.\allowbreak{}fourth\_\allowbreak{}initial\_\allowbreak{}of\_\allowbreak{}rank\_\allowbreak{}four\_\allowbreak{}data}{Development.\allowbreak{}AutomorphicArithmetic.\allowbreak{}AutomorphicVarieties.\allowbreak{}ModularityArgument.\allowbreak{}ThresholdArguments.\allowbreak{}DenseThresholdCases.\allowbreak{}BirthSeedIV}{18--63}

To join residual points put \(A_0=d_1/d_0,\ B_0'=d_2/d_0\), so \(d_3/d_0=A_0B_0'\). A pair of type \(P\) has the first representation and flag supports of \ref{B6.48b} and agrees in entry \(20\) of its matrices and flags. A pair of type \(M\) has the second representation support and agrees in entry \(10\) of matrices and flags. The divided entries \(x(g)=r_{10}(g)/d_0(g)\) and \(y(g)=r_{20}(g)/d_0(g)\) on their respective supports are cocycles of \(A_0,B_0'\). They vanish at \(g_0\). On a local group the flag identity gives \(x=(1-A_0)m_v,\ y=(1-B_0')n_v\), where \(m_v,n_v\) are the corresponding flag entries.

Given finitely many such \(y_j\), choose a finite coefficient field \(k_0'\) containing the residual values and let \(V_0\) be the \(k_0'\)-span of the values of \(y_j\). The values form finite sets by compactness. On the finite \(\F_\ell\)-space \(D'=k_0'\oplus V_0^{\{j\}}\) use
\[
 g(u,(v_j))=\bigl(A_0(g)u,
 (A_0(g)B_0'(g)v_j+A_0(g)u y_j(g))_j\bigr).
\]
The cocycle law for \(y_j\) gives the action law. Put \(m_0=(1,0)\). Let \(S_0\) contain the places over \(S_D\cup\{\ell\}\) and the primes already allowed, and impose zero local class there. Choose infinitely many new rational primes split in the normal closure of the fields cut out by this action, the old residual representation, \(F\) and \(\mu_\ell\), by Chebotarev \cite[\S1A]{thorner-2019}. At a new place allow inertia values in \(\F_\ell m_0\). The action there is trivial and the tame relation with norm \(Q\equiv1\bmod\ell\) allows a nonzero inertia value. The condition has at least \(2|D'|\) elements by varying the unramified part.

For a resulting set \(S\) let \(\mathcal S\) be the Selmer group with these conditions. Finite duality in \ref{B6.09} and the Euler characteristic \cite[Ch.~II, \S5]{serre-2002} give
\begin{equation}\tag{6.49.5}\label{eq:p6c5-gysin-count}
 2^{|S-S_0|}\le |\mathcal S|\, |(D')^\vee|\,|D'|^{[F:\Q]+|S_0|}.
\end{equation}
Each new factor divided by the invariant order \(|D'|\) is at least two. Each old local invariant factor costs at most \(|D'|\), each infinite factor at most \(|D'|\), and the global dual invariant order is at most \(|(D')^\vee|\). Let \(N_0\) be the number of cocycles unramified outside \(S_0\), finite by global finiteness and finiteness of the coboundaries. Set \(|S-S_0|=N_0 |(D')^\vee|\,|D'|^{[F:\Q]+|S_0|}\). Since \(2^b>b\) for \(b\ge0\), the Selmer group contains a class with nonzero inertia at a new place, of value in \(\F_\ell m_0\). Add the underlying rational primes to \(W\). On the still-excluded part of old inertia the action is trivial and the local class is a boundary, so its cocycle value vanishes.
\leanref{UpperTheory.\allowbreak{}V.\allowbreak{}Ws.\allowbreak{}Gys.\allowbreak{}selBig}{Development.\allowbreak{}AutomorphicArithmetic.\allowbreak{}AutomorphicVarieties.\allowbreak{}WeakShimuraTheory.\allowbreak{}RamifiedLineCharacters.\allowbreak{}SelBig}{7--43}
\leanref{UpperTheory.\allowbreak{}V.\allowbreak{}Ws.\allowbreak{}Gys.\allowbreak{}gysCore}{Development.\allowbreak{}AutomorphicArithmetic.\allowbreak{}AutomorphicVarieties.\allowbreak{}WeakShimuraTheory.\allowbreak{}RamifiedLineCharacters.\allowbreak{}GysCore}{9--62}
\leanref{UpperTheory.\allowbreak{}V.\allowbreak{}Ws.\allowbreak{}Gys.\allowbreak{}extRep}{Development.\allowbreak{}AutomorphicArithmetic.\allowbreak{}AutomorphicVarieties.\allowbreak{}WeakShimuraTheory.\allowbreak{}RamifiedLineCharacters.\allowbreak{}ExtRep}{34--193}

Project the cocycle to \(x,z_j\). Then \(x\) is nonzero somewhere on \(\ker A_0\) and
\[
 z_j(gh)=z_j(g)+A_0(g)B_0'(g)z_j(h)+y_j(g)A_0(g)x(h).
\]
Locally \(x=(1-A_0)m_v,\ z_j=(1-A_0B_0')o_{jv}-y_jA_0m_v\). Subtract the coboundary with coordinates
\[
 b=\frac{x(g_0)}{1-A_0(g_0)},\qquad
 c_j=\frac{z_j(g_0)}{1-A_0(g_0)B_0'(g_0)}
\]
to normalise at \(g_0\). The denominators are nonzero, and \(y_j(g_0)=0\). This subtracts \((1-A_0)b\) and \((1-A_0B_0')c_j-y_jA_0b\) from the cocycle coordinates and translates the local witnesses by \(-b,-c_j\).

Suppress \(j\) and put
\[
 U(x,y,t)=
 \begin{pmatrix}1&0&0&0\\x&1&0&0\\y&0&1&0\\t&y&-x&1\end{pmatrix},
 \qquad
 r(g)=U(x(g),y(g),2z(g)-x(g)y(g))\operatorname{diag}(d_i(g)).
\]
If \(y|_{G_v}=(1-B_0')n_v\), take flag matrix \(U(m_v,n_v,2o_v-m_vn_v)\). The product of two \(U\)'s adds the first two parameters and replaces the third by \(t+t'+yx'-xy'\). Diagonal conjugation scales the three parameters by \(A_0,B_0',A_0B_0'\). Thus the equation for \(z\) gives multiplicativity, and the local equations give the prescribed flag characters. The \(U\)'s preserve the form, and \(d_0d_3=d_1d_2\) gives the similitude. Replacing \(z\) by \(xy-z\) exchanges the cross terms in the cocycle equation. The reversed construction therefore preserves given \(10\) data and produces a nonzero \(20\) coordinate.
\leanref{UpperTheory.\allowbreak{}V.\allowbreak{}Ws.\allowbreak{}Gys.\allowbreak{}cupEx}{Development.\allowbreak{}AutomorphicArithmetic.\allowbreak{}AutomorphicVarieties.\allowbreak{}WeakShimuraTheory.\allowbreak{}RamifiedLineCharacters.\allowbreak{}GysinPsOf}{13--31}
\leanref{UpperTheory.\allowbreak{}V.\allowbreak{}Ws.\allowbreak{}Gys.\allowbreak{}ob}{Development.\allowbreak{}AutomorphicArithmetic.\allowbreak{}AutomorphicVarieties.\allowbreak{}WeakShimuraTheory.\allowbreak{}RamifiedLineCharacters.\allowbreak{}SuppMMono}{126--233}

For an irreducible point take its integral pair from Lemma~\ref{B6.49c}. In the block case with full first-block span use the nonsplit companion of Lemma~\ref{B6.49f}. Choose the label order of Lemma~\ref{B6.49c} in constructing the initial tower. Apply Lemma~\ref{B6.49d} over \(\Q\), and let \(p_0\) be the reduction of this lattice with its transported flags over the initial totally real field. Lemma~\ref{B6.49e} applies both here and after enlargement. In III start with the seed's \(10\) data and \(W=\{w\}\). The reversed construction gives \(p_2\) with the same \(10\) data and \(E(-\psi)\). Apply the simultaneous construction to the \(20\) data of \(p_0,p_2\), giving \(p_3,p_4\) with the respective \(20\) data and a common \(10\) coordinate with \(E(-\alpha)\). In II apply it to two copies of the \(20\) data of \(p_0\), giving \(p_3\) with \(E(-\alpha)\). Enlarge \(F\) by the cyclic fields of degree \(\ell\) unramified outside one \(w'\in W\). These are subfields of \(\Q(\mu_{w'})\) and are unramified at \(\ell\). The new tame inertia is killed by restriction. In III its unipotent image has exponent dividing \(\ell\) since the nilpotent part has fourth power zero and \(\ell\ge5\). In II the displayed \(U\)-shape has exponent dividing \(\ell\) also for \(\ell=3\). Transport flags to new prolongations by conjugation and take their lower unitriangular factors modulo upper triangular matrices, which preserves the matched coordinates. Write \(F_1\) for this tame enlargement and \(F_m\) for the totally real abelian cyclotomic layer over \(F_1\) cut out below by \(\varepsilon_\ell^{\ell-1}\equiv1\bmod\ell^{m+1}\). Raising the marking to an \(\ell\)-power puts \(t=g_0^{\ell^a}\) in the enlarged group \(G_{F_m}\) (which is normal of \(\ell\)-power index in \(G_F\), since both \(F_1\) and \(F_m\) are abelian over \(\Q\)) and preserves the residual eigenvalue inequalities because \(A_0,B_0'\) take values in \(k_\Lambda^\times\) of order prime to \(\ell\). In particular \(A_0(t)\ne1\) and \(B_0'(t)\ne1\). Consequently the constructed Selmer cocycle \(x\) (and in III the reversed Selmer cocycle \(y\) of \(p_2\)) remains nonzero after restriction to \(G_{F_m}\): if \(x|_{G_{F_m}}=0\), then for every \(g\in G_F\) both \(t\) and \(gtg^{-1}\) lie in \(G_{F_m}\), so comparing \(x(gt)=x(g)+A_0(g)x(t)=x(g)\) with \(x((gtg^{-1})g)=x(gtg^{-1})+A_0(gtg^{-1})x(g)=A_0(t)x(g)\) gives \((1-A_0(t))x(g)=0\), forcing \(x\equiv0\) on \(G_F\), contrary to its nonzero tame inertia at \(W\). In III the same computation with \(B_0'(t)\ne1\) shows \(y|_{G_{F_m}}\ne0\).

Call the enlarged field \(F_1\). The field fixed by the intersection of \(G_{F_1}\) with the residual diagonal and cyclotomic kernels is abelian and contains \(\mu_\ell\). By \ref{B6.17} one integer \(r\) bounds the unramified mod-\(\ell\) character dimensions over its cyclotomic layers. Let \(g\) bound the number of places over \(S_D\) and \(\ell\) in those layers. Such a bound exists because the Frobenius norm away from \(\ell\) generates an open subgroup of the cyclotomic \(\Z_\ell\)-quotient, and inertia has open image above \(\ell\). Take \(C=r+2+2g\). Restriction of a strict root cocycle to the residual kernel gives an unramified character, with at most one extra dimension for coboundaries. Old tame places cost at most one inertia generator each. The global Euler characteristic adds the negative real eigendimensions, zero for \(0,\pm\psi\) and at most \([F:\Q]\) for an odd root. Local duality and the local Euler characteristic bound the product-local odd cocycle dimension by \([F:\Q]+2g\). These are the bounds of \ref{B6.48b}. On the layers \(\varepsilon_\ell^{\ell-1}\equiv1\bmod\ell^{m+1}\), the \(\ell\)-primary part of \(\psi\) kills the infinite intersection by the initial kernel condition and hence a finite layer by compactness. The degrees are unbounded. Take a layer with \(d\ge30(C+1)^2\) on which \(\psi\) has order prime to \(\ell\).
\leanref{UpperTheory.\allowbreak{}V.\allowbreak{}Ws.\allowbreak{}diesEx}{Development.\allowbreak{}AutomorphicArithmetic.\allowbreak{}AutomorphicVarieties.\allowbreak{}WeakShimuraTheory.\allowbreak{}CompactIntegralFields.\allowbreak{}DiesEx}{112--184}
\leanref{UpperTheory.\allowbreak{}V.\allowbreak{}Ws.\allowbreak{}restrT}{Development.\allowbreak{}AutomorphicArithmetic.\allowbreak{}AutomorphicVarieties.\allowbreak{}WeakShimuraTheory.\allowbreak{}CompactIntegralFields.\allowbreak{}RestrT}{156--272}
\leanref{UpperTheory.\allowbreak{}V.\allowbreak{}Ws.\allowbreak{}x2}{Development.\allowbreak{}AutomorphicArithmetic.\allowbreak{}AutomorphicVarieties.\allowbreak{}WeakShimuraTheory.\allowbreak{}CompactIntegralFields.\allowbreak{}X2}{14--38}
\leanref{UpperTheory.\allowbreak{}V.\allowbreak{}Ws.\allowbreak{}layerEx}{Development.\allowbreak{}AutomorphicArithmetic.\allowbreak{}AutomorphicVarieties.\allowbreak{}WeakShimuraTheory.\allowbreak{}CompactIntegralFields.\allowbreak{}LayerEx}{112--305}

Write \(q_i\) for the restricted points and \(q_{\rm seed}\) for the seed reduction. In III the chain is
\[
 q_0\xrightarrow{P}q_3\xrightarrow{M}q_4\xrightarrow{P}q_2
 \xrightarrow{M}q_{\rm seed}.
 \tag{6.49.3}
\]
Number the alternatives of \ref{B6.48b} in their displayed order. Lemma~\ref{B6.49e} puts \(q_0\) in (1) or (3). In both chains \(q_3\) has the displayed unipotent shape \(U\) and its \(10\) coordinate is nonzero somewhere. Its \(20\) matrix and flag coordinates agree with those of \(q_0\). If its \(20\) matrix coordinate is nonzero somewhere, \(q_3\) satisfies (1). Otherwise the paired \(31\) entry vanishes as well, \(q_0\) must be in (3), and the common flag entries at \(20\) are nonzero at every place. This gives (3) at \(q_3\).

In III the \(M\)-match with \(q_3\) gives \(q_4\) the same nonzero \(10\) matrix entry, and the \(P\)-match with \(q_2\) gives it the nonzero \(20\) matrix entry of the reversed construction. Thus \(q_4\) satisfies (1). At \(q_2\) the \(10\) matrix entry agrees with the seed's zero entry and the \(10\) flag entries agree with its nonzero flag entries. In the displayed \(U\)-matrix the vanishing of \(10\) also kills \(32\). The remaining support is contained in that for \(H\), and the reversed construction gives \(E(-\psi)\). This proves (5) at \(q_2\). The tensor computation gives (5) at the seed with \(E(-\beta),E(-\theta)\).

In II take the two diagonal blocks of the lower block triangular matrices and flags of \(q_3\) in order \((0,1;2,3)\). This projection preserves multiplication and the symplectic adjoint, which exchanges the two diagonal blocks. It retains the \(10\) data and leaves only weight \(-\alpha\) off the diagonal. Write \(X_v\) for the resulting flag, and take \(y_v\) to be the seed's \(20\) flag entry. Keep the projected representation and use the flag \(X_v(1+y_v(E_{20}+E_{31}))\). The right factor is symplectic. In the projected flag basis the local matrices are diagonal, since they are both lower and upper triangular. Their diagonal entries satisfy \(d_0=d_2\) and \(d_1=d_3\) on the local group by the triviality of the root character of weight \(\psi\) there. Hence the right factor commutes with them. The new flag has the old \(10\) entry, has \(20\) entry \(y_v\ne0\), and has support in \(\{-\alpha,-\psi,-\theta\}\). This gives (3) at \(q_4\). Its matrices and the seed matrices have zero \(20\) entries and their \(20\) flags agree. Thus the chain is \(q_0\xrightarrow{P}q_3\xrightarrow{M}q_4\xrightarrow{P}q_{\rm seed}\). The induced matrix calculation gives (3) or (4) at the seed. For (4), at each ordinary place the tower supplies \(g\) in its decomposition group at which the residual character of weight \(\alpha\) has fourth power different from one. This character is \(d_0/d_1\), by the coordinate weights in \ref{B6.48b}. Hence \((d_0/d_1)(g)\ne1\), the required local test.

Each arrow is linear interpolation \((1-X)r+Xr'\) of the matrices and flags. In \(M\) the differences lie in the lower-left block for \((0,1;2,3)\) by agreement of the \(10\) data and the symplectic equations. In \(P\), use \((0,2;1,3)\). Agreement at \(20\) gives agreement at \(31\) by \(r_{00}r_{31}=r_{11}r_{20}\). The differences again occupy the lower-left block, and products of two differences vanish. Expansion of the defining equations gives points over \(k_\Lambda[X]\).
\leanref{UpperTheory.\allowbreak{}V.\allowbreak{}Ws.\allowbreak{}lineP}{Development.\allowbreak{}AutomorphicArithmetic.\allowbreak{}AutomorphicVarieties.\allowbreak{}WeakShimuraTheory.\allowbreak{}CompactIntegralFields.\allowbreak{}AffSmulOne}{127--216}
\leanref{UpperTheory.\allowbreak{}V.\allowbreak{}Ws.\allowbreak{}lineM}{Development.\allowbreak{}AutomorphicArithmetic.\allowbreak{}AutomorphicVarieties.\allowbreak{}WeakShimuraTheory.\allowbreak{}CompactIntegralFields.\allowbreak{}GlobalTheory}{131--199}
\leanref{UpperTheory.\allowbreak{}V.\allowbreak{}Ws.\allowbreak{}deep\_\allowbreak{}tower\_\allowbreak{}residue\_\allowbreak{}chain\_\allowbreak{}exists}{Development.\allowbreak{}AutomorphicArithmetic.\allowbreak{}AutomorphicVarieties.\allowbreak{}ModularityArgument.\allowbreak{}ThresholdArguments.\allowbreak{}DenseThresholdCases.\allowbreak{}Chain2}{13--92}
\leanref{UpperTheory.\allowbreak{}V.\allowbreak{}Cu.\allowbreak{}chainIV\_\allowbreak{}alt}{Development.\allowbreak{}AutomorphicArithmetic.\allowbreak{}AutomorphicVarieties.\allowbreak{}ModularityArgument.\allowbreak{}ThresholdArguments.\allowbreak{}DenseThresholdCases.\allowbreak{}CuPChainIV}{16--82}

Let \(F\) be the totally real field of the fitted final layer and \(K\) its imaginary quadratic field. Choose the common polarising twist \(\kappa\) over their CM compositum \(L=FK\) for both Hecke comparisons. Extend coefficients to contain its values, the auxiliary character values and all images of \(\mathcal O_L\). Use the split levels of \ref{B6.48c}, taking the second distinguished place away from the first rational prime. The residual cyclotomic test gives \(\mu_\ell\not\subset L\). The twisted determinants and characters give maps \(\phi_i:B_i\to B\). For the two ideals \(J_i\) from \ref{B6.48aa}, set \(J=\sum_i\phi_i(J_i)B\).
\leanref{UpperTheory.\allowbreak{}Gal.\allowbreak{}projectedFormsSystemPart.\allowbreak{}exists\_\allowbreak{}two\_\allowbreak{}links\_\allowbreak{}same\_\allowbreak{}kappa\_\allowbreak{}not\_\allowbreak{}dvd\_\allowbreak{}of\_\allowbreak{}fit\_\allowbreak{}deep}{Development.\allowbreak{}AutomorphicArithmetic.\allowbreak{}AutomorphicVarieties.\allowbreak{}ModularityArgument.\allowbreak{}ThresholdArguments.\allowbreak{}DenseThresholdCases.\allowbreak{}LinkTwo}{13--51}

This \(J\) satisfies \eqref{eq:p6c5-potential-finite}. In any prime quotient containing the right side, the weight variables specialise to their residual values. Hecke finiteness over weight space makes the Hecke coefficients residual as well. The twisted flag values and polynomial coefficients from the two maps generate the coefficient ring in the fibre. Hence the prime contains \(\mathfrak m_B\), and a uniform power follows by Noetherianity.
\leanref{UpperTheory.\allowbreak{}V.\allowbreak{}Ws.\allowbreak{}exists\_\allowbreak{}pow\_\allowbreak{}maximalIdeal\_\allowbreak{}le\_\allowbreak{}of\_\allowbreak{}modular\_\allowbreak{}phi}{Development.\allowbreak{}AutomorphicArithmetic.\allowbreak{}AutomorphicVarieties.\allowbreak{}ModularityArgument.\allowbreak{}ThresholdArguments.\allowbreak{}DenseThresholdCases.\allowbreak{}TPmFin}{126--173}

For a curve \(y\) as in \ref{B6.48b}, the maps \(f_i=yj\phi_i\) specialise the determinant and ordered characters to the twisted restrictions. Elements separating the flags on \(U_v\) lift to the local link groups with cyclotomic value one. The residual tests for the good datum give \(\psi_i\ne\psi_j\varepsilon_\ell\) in the required order. The distinguished local test has equal residual eigenvalues and nontrivial residual cyclotomic value, so Lemma~\ref{B6.49i}(1) with \(m=4\) and \(\chi=\varepsilon_\ell|_{D_0}\) gives zero twisted intertwiner space. Full generic span survives the quadratic restriction. Otherwise a sign self-twist would force trace zero on the other coset, which contradicts linear independence of the four restricted residual characters, with unit coefficients from the coset representative.

Each pair of flag characters attains a nontrivial principal unit on \(U_v\), which has infinite order in characteristic \(\ell\). Among combinations of finitely many local separating elements, avoid the finitely many proper kernels of the character ratios modulo torsion. This gives a local element whose positive powers all separate. Reducibility on the cyclotomic kernel would give a quadratic permutation twist by Clifford theory, since the abelian quotient acts transitively on the constituents and their number divides four. For odd \(\ell\) the twist factors through the residual cyclotomic quotient. The fitted residual tests exclude it.

Let \(H\) be the CM subgroup and \(H'=H\cap\ker\varepsilon_\ell\). We check the noncommutativity hypothesis of Lemma~\ref{B6.49g} on \(H'\). If an open normal subgroup \(N\le H'\) had commuting image, its core under \(H\) would also have commuting image and would contain \(u^{[H':N]}\), where \(u\in H'\) is the separating element above. Indeed each \(H\)-conjugate of \(u\) lies in \(H'\) and its indicated power lies in \(N\). The core acts diagonally on the four eigenlines of this power. The lines are its distinct character spaces. Thus \(H\) permutes them transitively by irreducibility. Let \(P\) be the kernel of the permutation. Lemma~\ref{B6.49i}(2) with \(m=4\) shows that every ratio of the prescribed residual characters is trivial on \(P\). Such a ratio factors through a subgroup of \(S_4\). Its image in the multiplicative residue field is finite cyclic, and a preimage of a generator acts as a permutation of order at most four. This contradicts the fitted ratio whose first four powers are nontrivial. Apply Lemma~\ref{B6.49g} to the absolutely irreducible restriction to \(H'\), in the eigenbasis of \(u\). In a preserving coset it supplies \(M\) with nonzero diagonal and adjugate diagonal, since \(\operatorname{adj}M=(\det M)M^{-1}\). In an exchanging coset it supplies both off-diagonal blocks entrywise nonzero.

Take a power of the separating element into the open subgroup and kill torsion in the group generated by its eigenvalues \(a_i\). Map this finitely generated torsion-free group to \(\Z\) by a homomorphism with distinct eigenvalue weights \(w_i\). In the preserving case order them \(w_1>w_2>w_3>w_4\). The first, third and fourth elementary coefficients of \(M\operatorname{diag}(X^{w_i})\) have respective leading exponents \(w_1,w_1+w_2+w_3,\sum_iw_i\) with nonzero coefficients. The second coefficient has exponent at most \(w_1+w_2\). The Newton polygon at infinity gives unique extreme roots of exponents \(w_1,w_4\). A repeated middle pair would have exponent \((w_2+w_3)/2\), and in odd characteristic the second coefficient would have exponent \(w_1+(w_2+w_3)/2\). This lies strictly between the largest and next possible pair-weights \(w_1+w_2,w_1+w_3\), a contradiction.

In the exchanging case a block of the square is \(B\operatorname{diag}(X^{w'_i})C\operatorname{diag}(X^{w_i})\), with \(B,C\) the off-diagonal blocks. Order the weights within each block. Its trace has unique leading exponent \(w_1+w'_1\) with nonzero coefficient. The determinant has exponent \(w_1+w_2+w'_1+w'_2<2(w_1+w'_1)\) and nonzero coefficient. The square block has two distinct nonzero roots, so their opposite square roots give four distinct roots. In both cases the weighted discriminant is nonzero. Group its multivariate terms by the equal products of the \(a_i\). At least one group has nonzero sum of coefficients, since equal products have equal weights. Vandermonde in the distinct products gives a nonzero evaluation at a common integer power of the \(a_i\). This proves the required separability.
\leanref{UpperTheory.\allowbreak{}V.\allowbreak{}Ws.\allowbreak{}P1.\allowbreak{}irrL}{Development.\allowbreak{}AutomorphicArithmetic.\allowbreak{}AutomorphicVarieties.\allowbreak{}WeakShimuraTheory.\allowbreak{}ProjectiveLineFlags.\allowbreak{}TraceOfCp}{17--75}
\leanref{UpperTheory.\allowbreak{}V.\allowbreak{}Ws.\allowbreak{}h4}{Development.\allowbreak{}AutomorphicArithmetic.\allowbreak{}AutomorphicVarieties.\allowbreak{}WeakShimuraTheory.\allowbreak{}CompactIntegralFields.\allowbreak{}H4}{19--118}
\leanref{UpperTheory.\allowbreak{}V.\allowbreak{}AlgebraTools.\allowbreak{}regCosetFour}{Development.\allowbreak{}AutomorphicArithmetic.\allowbreak{}AutomorphicVarieties.\allowbreak{}AlgebraTools.\allowbreak{}MonomialRepresentations.\allowbreak{}RegCosetFour}{337--395}

Now \ref{B6.48aa}(a) applies to \(f_i\). For \(\mathfrak q\subseteq\ker(yj)\), the prime \(\phi_i^{-1}\mathfrak q\) lies below \(\ker f_i\), and \(f_i(J_i)=0\). Thus \(J_i\subseteq\phi_i^{-1}\mathfrak q\). Extending and summing gives \(J\subseteq\mathfrak q\), as required in \ref{B6.48b}.
\leanref{UpperTheory.\allowbreak{}V.\allowbreak{}Ws.\allowbreak{}modular\_\allowbreak{}deformation\_\allowbreak{}quotient\_\allowbreak{}regular\_\allowbreak{}at\_\allowbreak{}point\_\allowbreak{}alt}{Development.\allowbreak{}AutomorphicArithmetic.\allowbreak{}AutomorphicVarieties.\allowbreak{}ModularityArgument.\allowbreak{}ThresholdArguments.\allowbreak{}DenseThresholdCases.\allowbreak{}ModularDeformationQuotientRegularAtPoint}{25--59}

At the characteristic-zero seed, the algebraic twist for both links gives integral weights with gaps at least two. The attachment, ordinary flag and prescribed inertia give a point on each Hecke quotient. A first-order determinant lift with fixed weight gives a matrix lift by full span. The flag lifts by separation, and \eqref{eq:p6c5-local-finite} gives a finite adjoint class. The weight gaps exclude \(\psi_i=\psi_j\varepsilon_\ell\) on inertia. Local comparison relates the ordered characters to the toral support by their cyclotomic weights \(m_i\). At a Weil element of cyclotomic value one and positive Weil degree \(d_0\), the support moduli are \(|\iota(\psi_i)|q_v^{-m_i d_0}\). Reciprocity gives such an element for the flag. A finite-order ratio \(\psi_i/\psi_j\) on it would give equal character moduli, and \eqref{eq:p6c5-toral-support} would give \(|m_i-m_j|<1\). Thus the gaps give infinite-order separation. The residual tests give dense separability and the seed is irreducible on the cyclotomic kernel. By \ref{B6.53} the class is a coboundary, so its derivatives on all coefficients and characters vanish. The local determinant ring has embedding dimension at most the number of weight parameters. The Hecke quotient is finite torsion-free over the regular weight ring and gives the opposite dimension bound. Therefore the local determinant ring is regular and the Hecke quotient has zero kernel there by equality of dimensions. Each prime below the seed prime contains \(J_i\).
\leanref{UpperTheory.\allowbreak{}V.\allowbreak{}Th.\allowbreak{}local\_\allowbreak{}zero\_\allowbreak{}of\_\allowbreak{}modular\_\allowbreak{}p}{Development.\allowbreak{}AutomorphicArithmetic.\allowbreak{}AutomorphicVarieties.\allowbreak{}ModularityArgument.\allowbreak{}ThresholdArguments.\allowbreak{}DenseThresholdCases.\allowbreak{}BirthPm}{28--65}
\leanref{UpperTheory.\allowbreak{}V.\allowbreak{}Th.\allowbreak{}eigenbasis\_\allowbreak{}rank\_\allowbreak{}four}{Development.\allowbreak{}AutomorphicArithmetic.\allowbreak{}AutomorphicVarieties.\allowbreak{}ModularityArgument.\allowbreak{}ThresholdArguments.\allowbreak{}DenseThresholdCases.\allowbreak{}EigenbasisRankFour}{122--173}
\leanref{UpperTheory.\allowbreak{}V.\allowbreak{}Z.\allowbreak{}Win.\allowbreak{}RR.\allowbreak{}free\_\allowbreak{}core\_\allowbreak{}alt}{Development.\allowbreak{}AutomorphicArithmetic.\allowbreak{}AutomorphicVarieties.\allowbreak{}ModularityArgument.\allowbreak{}ThresholdArguments.\allowbreak{}DenseThresholdCases.\allowbreak{}RRFree}{97--116}

Contraction along the links and summation put \(J\) in each coefficient prime below the seed kernel. A minimal prime of \(A\) below the integral seed kernel contains \(JA\) and lies below the residual seed kernel. The standing hypothesis \(\mathscr H\) supplies the four-dimensional frame cohomology of Definition~\ref{B6.48be}, so Lemma~\ref{B6.48b} applies at the seed and at each link endpoint; hence every minimal prime below the residual seed kernel contains \(JA\). In a link \(A\to k_\Lambda[X]\), a minimal prime below the polynomial kernel lies below both endpoint kernels. The conclusion at one endpoint puts \(JA\) in this prime, so \ref{B6.48b} applies at the other endpoint. Follow the links back to \(q_0\). The prime kernel of the valuation-domain lift lies below the residual kernel at \(q_0\) and contains one of its minimal primes. It therefore contains \(JA\).
\leanref{UpperTheory.\allowbreak{}V.\allowbreak{}Ws.\allowbreak{}conn}{Development.\allowbreak{}AutomorphicArithmetic.\allowbreak{}AutomorphicVarieties.\allowbreak{}WeakShimuraTheory.\allowbreak{}CompactIntegralFields.\allowbreak{}AlgHomResUnique}{56--68}

In characteristic \(\ell\), the flag characters on the prescribed inertia of the fitted tower make its inertia generators in \eqref{eq:p6c5-potential-finite} vanish. The characteristic polynomials on \(G_F\) therefore equal the products prescribed by the residual characters. Every group element has a positive power in \(G_F\), so all its eigenvalues are algebraic over \(\F_\ell\). The trace difference from the residual sum is algebraic over \(\F_\ell\) and lies in the maximal valuation ideal. Hence it is zero, and multiplication of the residual characters gives
\[
 \operatorname{tr}\bigl(r(x)[r(g),r(h)]\bigr)=0
 \qquad(x,g,h\in G_\Q).
\]
Here the bracket denotes the additive commutator. The test is linear in the left multiplier, so it holds for that multiplier anywhere in the span of the group matrices. Under full span, testing against each matrix unit gives every entry of the commutator equal to zero. The group matrices would then span a commutative matrix algebra, contrary to full span in rank four. In the block case write \(A_x,B_x\) for the two diagonal blocks of \(r(x)\), and let \(p_B\) be the characteristic polynomial of the second block at the marking \(g_1\) (using the marked element for this datum). The group algebra contains the multiplier \(r(x)p_B(r(g_1))\). Its second diagonal block vanishes by Cayley--Hamilton. Its first is \(A_xp_B(A_{g_1})\), with \(p_B(A_{g_1})\) invertible by coprimality and B\'ezout. The trace test therefore gives \(\operatorname{tr}(A_xp_B(A_{g_1})[A_g,A_h])=0\) for every \(x\). As \(A_x\) spans the full two-by-two algebra, the trace pairing gives \([A_g,A_h]=0\) for every pair. This again contradicts full matrix span.

For the application of \ref{B6.49a}, the framed universal lift has a homogeneous presentation in the \(\Z^2\) weights of \ref{B6.48b}, with coordinates given by entries minus their residual lifts. Consider its points in characteristic-\(\ell\) complete DVRs. The marked projectors put the diagonal matrix units in the span. The relation \(i\leftarrow j\) defined by the presence of a nonzero \((i,j)\) entry is transitive, by multiplying matrices with a diagonal projector inserted. In the support classification the alternatives are full span, absence of a pair of opposite off-diagonal entries, and both entries \(02,20\) present with the matrices block triangular for \(\{0,2\},\{1,3\}\). In the last case projection to the two diagonal blocks preserves multiplication, similitude and the local flag conditions. The entries \(02,20\) and the marked projectors give full first-block span. Lemma~\ref{B6.49f} gives a nonsplit extension with the same characteristic polynomials. Lemmas~\ref{B6.49c}--\ref{B6.49e} give its starting point for the chain. The trace test with the marked projector applies to its block quotient. Thus the full-span and full-block alternatives are excluded in characteristic \(\ell\).
\leanref{UpperTheory.\allowbreak{}V.\allowbreak{}Ws.\allowbreak{}triV}{Development.\allowbreak{}AutomorphicArithmetic.\allowbreak{}AutomorphicVarieties.\allowbreak{}WeakShimuraTheory.\allowbreak{}CompactIntegralFields.\allowbreak{}CobOfCoord}{297--338}
\leanref{UpperTheory.\allowbreak{}V.\allowbreak{}Ws.\allowbreak{}blockPart}{Development.\allowbreak{}AutomorphicArithmetic.\allowbreak{}AutomorphicVarieties.\allowbreak{}WeakShimuraTheory.\allowbreak{}CompactIntegralFields.\allowbreak{}BlockPart}{192--267}
\leanref{UpperTheory.\allowbreak{}V.\allowbreak{}Ws.\allowbreak{}bl}{Development.\allowbreak{}AutomorphicArithmetic.\allowbreak{}AutomorphicVarieties.\allowbreak{}WeakShimuraTheory.\allowbreak{}CompactIntegralFields.\allowbreak{}NotOver}{51--85}

In the remaining support, each diagonal entry is a character because every product of opposite off-diagonal entries vanishes. Let \(U=\ker\bar r_D\cap\ker\vartheta\). On inertia in \(U\) the representation is trivial away from \(\ell\) by the frame conditions. At \(\ell\) it is unipotent in the ordinary flag, so its \(\ell^2\)-th power is one in dimension four. Consequently the diagonal characters raised to \([G_\Q:U]\ell^2\) kill all inertia. Apply this at every finite coefficient depth. An everywhere unramified finite character of \(G_\Q\) is trivial, and separation in the coefficient topology gives the same conclusion for the powered characters. Their values in the valuation ring are thus algebraic over \(\F_\ell\) and equal their residual lifts. The zero-weight coordinates vanish. The transitive relation on distinct indices is acyclic, so a nonzero product of off-diagonal entries cannot have weight zero. Indeed the weights of such a product sum to zero only if the flow between the two pairs of coordinate weights balances both coordinates. For the symplectic support with no opposite entries, a linear functional on the two weight coordinates is strictly positive on every present off-diagonal weight. To construct the functional, choose integers \(h_i\) strictly increasing along the strict partial order of the surviving entries, taking \(h_i>h_j\) for an entry in position \((i,j)\). The symplectic involution reverses this order. Thus \(H_i=h_i-h_{3-i}\) still satisfies \(H_i>H_j\) on every such entry and \(H_{3-i}=-H_i\). Since \(H_0+H_3=H_1+H_2\), the values \(H_i-H_3\) are exactly the values of a linear functional on the four displayed basis weights. Its value on every surviving off-diagonal entry weight is strictly positive. These values, up to a common shift, are obtained from a linear form on \((1,1),(0,1),(1,0),(0,0)\). Thus every positive-degree monomial of weight zero vanishes. Lemma \ref{B6.49a} makes the universal weight-zero part finite over \(O\), and its image contains the weight-zero monomials of every allowed lift.
\leanref{UpperTheory.\allowbreak{}V.\allowbreak{}Ws.\allowbreak{}Po.\allowbreak{}univGr}{Development.\allowbreak{}AutomorphicArithmetic.\allowbreak{}AutomorphicVarieties.\allowbreak{}WeakShimuraTheory.\allowbreak{}PolynomialTori.\allowbreak{}UnivGr}{57--85}
\leanref{UpperTheory.\allowbreak{}V.\allowbreak{}Ws.\allowbreak{}diagConst}{Development.\allowbreak{}AutomorphicArithmetic.\allowbreak{}AutomorphicVarieties.\allowbreak{}WeakShimuraTheory.\allowbreak{}CompactIntegralFields.\allowbreak{}DiagConst}{201--230}

In characteristic zero the two evaluations at an allowed point are arithmetic, with weights \(k+1,k,1,0\) up to the common twist. For an irreducible point on a quadratic restriction, passage to \(L\) preserves irreducibility. A nontrivial self-twist trivial on \(G_L\) would force trace zero on a coset, contradicting the independence of the restricted residual characters with unit coefficients. Apply \ref{B6.48a} to the two evaluations for the twisted representation. The good polynomials give the same attachment and the two distinguished places give the full witness by that lemma. Undo the twist and use solvable descent from \(L\) to the imaginary quadratic field by Proposition~\ref{B6.06d}(4), polarizing each step via the similitude pairing. This gives the attachment and witness on each irreducible quadratic restriction with the fixed \(\iota\).
\leanref{UpperTheory.\allowbreak{}V.\allowbreak{}Ws.\allowbreak{}contradiction\_\allowbreak{}star\_\allowbreak{}promodular}{Development.\allowbreak{}AutomorphicArithmetic.\allowbreak{}AutomorphicVarieties.\allowbreak{}ModularityArgument.\allowbreak{}ThresholdArguments.\allowbreak{}DenseThresholdCases.\allowbreak{}ContradictionStarPromodular}{183--222}
\leanref{UpperTheory.\allowbreak{}Gal.\allowbreak{}projectedFormsSystemPart.\allowbreak{}promodular\_\allowbreak{}existence\_\allowbreak{}of\_\allowbreak{}local\_\allowbreak{}zero}{Development.\allowbreak{}AutomorphicArithmetic.\allowbreak{}AutomorphicVarieties.\allowbreak{}ModularityArgument.\allowbreak{}ThresholdArguments.\allowbreak{}DenseThresholdCases.\allowbreak{}PmTwo}{21--51}
\leanref{UpperTheory.\allowbreak{}V.\allowbreak{}Ws.\allowbreak{}aut\_\allowbreak{}n\_\allowbreak{}alt}{Development.\allowbreak{}AutomorphicArithmetic.\allowbreak{}AutomorphicVarieties.\allowbreak{}ModularityArgument.\allowbreak{}ThresholdArguments.\allowbreak{}DenseThresholdCases.\allowbreak{}MapRepComp}{201--243}
\leanref{UpperTheory.\allowbreak{}V.\allowbreak{}Z.\allowbreak{}Win.\allowbreak{}Dn.\allowbreak{}pmDown\_\allowbreak{}alt}{Development.\allowbreak{}AutomorphicArithmetic.\allowbreak{}AutomorphicVarieties.\allowbreak{}ModularityArgument.\allowbreak{}ThresholdArguments.\allowbreak{}DenseThresholdCases.\allowbreak{}DnPmDown}{32--87}

For a characteristic-zero block point with full first-block span apply the chain to the companion of Lemma~\ref{B6.49f} and transfer its characteristic polynomials back to the block point. The arithmetic Hecke comparison gives a semisimple parameter whose constituents occur in strings \(\rho_i\varepsilon_\ell^{-j}\), \(0\le j<m_i\), by \ref{B6.04a} and \ref{B6.04aa}. If \(m_i>1\), reduction of consecutive constituent polynomials gives two residual characters related by \(\bar\varepsilon_\ell\) by uniqueness of residual semisimplification. Fittedness excludes this. Hence the pure constituents at the common arithmetic weight give, on an open subgroup at good Frobenius of norm \(Q\),
\begin{equation}\tag{6.49.4}\label{eq:p6c5-root-ratio}
 |x|<Q|y|
\end{equation}
for every pair of roots \(x,y\). The scalar twist preserves ratios, and passing to a Frobenius power raises both the ratios and the norm to that power. The two block determinants on their respective open subgroups are \(\varepsilon_\ell^{k+2}\) and \(\varepsilon_\ell^k\). Take a good Frobenius power in the intersection of the three subgroups. Apply \eqref{eq:p6c5-root-ratio} to the roots of the first block against those of the second and multiply the two inequalities. The result is \(Q^{k+2}<Q^2Q^k\), which proves (6.49.1).
\leanref{UpperTheory.\allowbreak{}V.\allowbreak{}Ws.\allowbreak{}frobIn}{Development.\allowbreak{}AutomorphicArithmetic.\allowbreak{}AutomorphicVarieties.\allowbreak{}WeakShimuraTheory.\allowbreak{}CompactIntegralFields.\allowbreak{}FrobIn}{31--59}
\leanref{UpperTheory.\allowbreak{}V.\allowbreak{}Ws.\allowbreak{}nePure\_\allowbreak{}w}{Development.\allowbreak{}AutomorphicArithmetic.\allowbreak{}Foundations.\allowbreak{}GaloisRepresentations.\allowbreak{}AdmissibleResidues.\allowbreak{}PosOfMulPosOfPos}{231--287}
\leanref{UpperTheory.\allowbreak{}V.\allowbreak{}Cu.\allowbreak{}modularity\_\allowbreak{}rank\_\allowbreak{}four\_\allowbreak{}of\_\allowbreak{}promodularity\_\allowbreak{}case\_\allowbreak{}three}{Development.\allowbreak{}AutomorphicArithmetic.\allowbreak{}AutomorphicVarieties.\allowbreak{}ModularityArgument.\allowbreak{}ThresholdArguments.\allowbreak{}DenseThresholdCases.\allowbreak{}ModularityRankFourOfPromodularity}{20--38}
\leanref{UpperTheory.\allowbreak{}V.\allowbreak{}Cu.\allowbreak{}promodularity\_\allowbreak{}case\_\allowbreak{}three\_\allowbreak{}of\_\allowbreak{}chain\_\allowbreak{}reduction}{Development.\allowbreak{}AutomorphicArithmetic.\allowbreak{}AutomorphicVarieties.\allowbreak{}ModularityArgument.\allowbreak{}ThresholdArguments.\allowbreak{}DenseThresholdCases.\allowbreak{}PromodularityCaseThreeOfLemmas}{131--153}
\leanref{UpperTheory.\allowbreak{}V.\allowbreak{}Ws.\allowbreak{}fin0\_\allowbreak{}no\_\allowbreak{}end\_\allowbreak{}all\_\allowbreak{}irr\_\allowbreak{}of\_\allowbreak{}good\_\allowbreak{}con\_\allowbreak{}purity\_\allowbreak{}aut}{Development.\allowbreak{}AutomorphicArithmetic.\allowbreak{}AutomorphicVarieties.\allowbreak{}ModularityArgument.\allowbreak{}ThresholdArguments.\allowbreak{}DenseThresholdCases.\allowbreak{}VWinPm}{75--106}
\end{proof}

\subsection*{Proof of \ref{B6.50}}\label{prf:B6.50}\phantomsection
\begin{proof}
Use the rank-\(n\) rational count of \ref{B6.41a} with \(n=4\). Retain seventeen of the eighteen fixed imaginary quadratic fields at \(p\). Each restricted multiset has at most sixteen lifts. The weights \((k+1,k,1,0)\) contain the consecutive pairs \(k+1,k\) and \(1,0\), so the count uses the second alternative of Lemma~\ref{B6.41b}(b): \(\ell\) divides \(M\) and every \(r\in X\) is unlinked at \(\ell\). The minimum weight in the integral comparison of forward roots is zero. This gives a conjugation-stable label set \(S\) with distinct labels for distinct isomorphism classes and integral traces \(t_s\). Put \(x=p^{(k+1)/2}\), \(y=p^k\). The witness and similitude give \(|\beta|\le\sqrt p\,x\) for the four forward roots and \(\prod|\beta|=x^4\), so one has modulus at most \(x\). The count and these bounds give
\begin{equation}\tag{6.50.3}\label{eq:p6c5-four-count}
 \#S\le c_0[M(k+1)]^{A_0},\qquad |t_s|\le3\sqrt p\,x+x\le\tfrac72 y .
\end{equation}
For \(a\in(\Z/h)^\times\) put \(F_{s,a}=t_s-(p^{k+1}z_1^a+p^kz_2^a+pz_3^a+1)\). These are integral. Since \(p+1<y/2\), \eqref{eq:p6c5-four-count} bounds the terms apart from the leading one by \(9y/2+p+1<py\). Hence \(0<|F_{s,a}|\le p^{2(k+1)}\). Galois automorphisms permute the factors, so \(D=\prod_{s,a}F_{s,a}\) is a nonzero integer, with empty product one. The size bound and \eqref{eq:p6c5-four-count} yield (6.50.1). The forward trace is the sum of the inverses of the roots in the inverse-Frobenius comparison. Thus \(\Delta_r=\iota^{-1}(F_{s,1})\) with distinct labels for \(r\in X\). All other factors have nonnegative valuation, giving (6.50.2).
\leanref{UpperTheory.\allowbreak{}GalN.\allowbreak{}V.\allowbreak{}Win.\allowbreak{}N.\allowbreak{}number\_\allowbreak{}sp4\_\allowbreak{}of\_\allowbreak{}admissible\_\allowbreak{}to\_\allowbreak{}admissible\_\allowbreak{}n}{Development.\allowbreak{}AutomorphicArithmetic.\allowbreak{}AutomorphicVarieties.\allowbreak{}ModularityArgument.\allowbreak{}ThresholdArguments.\allowbreak{}DenseThresholdCases.\allowbreak{}VWinNumN}{94--155}
\end{proof}

\subsection*{Proof of \ref{B6.51fin}}\label{prf:B6.51fin}\phantomsection
\begin{proof}
The monoid of nonnegative exponent vectors of weight zero is finitely generated. Its nonzero minimal vectors for the coordinatewise order form a finite set by Dickson's lemma, and subtracting such a vector preserves weight zero. Substitution of the corresponding monomials presents \(R_0\) as a quotient of a power series ring over \(O\). These monomials and \(\varpi\) generate its maximal ideal.

Let \(b\) be one of the monomial generators. If \(b\notin\sqrt{\varpi R}\), choose a prime maximal among those containing \(\varpi\) and avoiding \(b\). Its quotient is a local domain in which \(b\) is a nonzero nonunit. In dimension at least two, prime avoidance gives an element of the maximal ideal outside the finitely many height-one primes containing \(b\). A minimal prime over that element has height one and avoids \(b\), contrary to maximality of the chosen prime. Thus the quotient has dimension one. Its normalisation gives a local map to a complete DVR with finite residue field in which \(b\) survives. Every nonconstant monomial of weight zero contains a zero-weight factor or factors of both strict signs. Hence the hypothesis kills \(b\), a contradiction.

Projection to weight zero gives \(\varpi R\cap R_0=\varpi R_0\). Thus each monomial generator is nilpotent modulo \(\varpi R_0\), and the finitely generated maximal ideal of \(R_0/\varpi R_0\) is nilpotent. This quotient is finite-dimensional over the residue field of \(O\). Lift a basis and approximate successively modulo \(\varpi^jR_0\). The coefficients converge in \(O\) and the remainders tend to zero in the maximal-ideal topology of \(R_0\), so the lifts generate \(R_0\) over \(O\).
\end{proof}

\subsection*{Proof of \ref{B6.51}}\label{prf:B6.51}\phantomsection
\begin{proof}
Map the absolute Galois marking of \ref{B6.52}(4) to \(\Gamma_{\Q,S}\). The frame sign and character pull back to \(\eta_4,\Psi\), so the image satisfies \eqref{eq:p6c5-gaussian-marking}. A point of the marked frame pulls back to an ordinary point of that endpoint. Over a residue-characteristic DVR its polynomial identity on the specified finite-index subgroup, together with the marking and normalisation, gives upper or lower triangularity on the sign kernel. The ratio of each diagonal character to its prescribed character is trivial on the conjugates of the framed inertias. The frame quotient makes this ratio everywhere unramified over \(\Q(i)\). Class field theory and the class number of \(\Q(i)\) make it trivial. Use the adapted presentation in which each coordinate is an entry of the universal matrix minus the corresponding base entry, with weight given by diagonal conjugation. The diagonal equality kills the zero-weight coordinates, and triangularity kills one strict sign. Lemma~\ref{B6.51fin} gives finiteness of \(R_0\).
\leanref{UpperTheory.\allowbreak{}V.\allowbreak{}UnitaryThreeInertia.\allowbreak{}frame\_\allowbreak{}class\_\allowbreak{}one}{Development.\allowbreak{}AutomorphicArithmetic.\allowbreak{}AutomorphicVarieties.\allowbreak{}UnitaryThreeDeformations.\allowbreak{}FrameClassOne}{166--268}
\leanref{UpperTheory.\allowbreak{}V.\allowbreak{}UnitaryThreeInertia.\allowbreak{}universal\_\allowbreak{}diagonal\_\allowbreak{}of\_\allowbreak{}triangular\_\allowbreak{}points}{Development.\allowbreak{}AutomorphicArithmetic.\allowbreak{}AutomorphicVarieties.\allowbreak{}UnitaryThreeDeformations.\allowbreak{}UniversalDiagonalOfTriangularPoints}{22--73}

In characteristic zero, finiteness of \(R_0\) gives a model over a finite coefficient DVR after diagonal conjugation. Work with this model. Let \(U\) be the normal core in \(G_\Q\) of an open subgroup of \(G_K\). Choose \(\sigma\in I_\ell\) with \(\varepsilon_\ell(\sigma)=1+\ell\), and put
\[
 s=\sigma^{[G_\Q:U]\varphi(N)}\in U\cap I_\ell .
\]
Then \(\vartheta(s)=1\) and the three inertia eigenvalues at \(s\) are distinct. The local square test from \eqref{eq:p6c5-gaussian-host} gives \(x\in D_\ell\cap G_K\) with \(\bar\Psi(x)^2\ne1\). If restriction to \(U\) is reducible, every \(U\)-invariant subspace is spanned by a subset of the \(s\)-eigenlines. The minimal such nonzero subspaces have disjoint sets of eigenlines. Their sum is \(G_K\)-invariant by normality, hence is the whole space. The action on them is transitive. Since the dimension is three, reducibility gives three lines. The local flag is the flag of the first one and the first two eigenlines in the ordinary order. Its preservation by \(x\) forces \(x\) to fix the three lines in this permutation action. Choose \(g\in G_K\) acting by a three-cycle. The trace of \(r(gx^j)\) is zero for every \(j\ge0\). Reducing at \(j=0,1,2\) gives a Vandermonde system with nodes \(\bar\Psi(x),1,\bar\Psi(x)^{-1}\) and coefficients \(\bar\Psi(g),1,\bar\Psi(g)^{-1}\). The nodes are distinct and the coefficients nonzero. This contradiction proves irreducibility on \(U\) and on the original open subgroup. Apply \ref{B6.52}(4) and transport the attachment through the diagonal conjugation.
\leanref{UpperTheory.\allowbreak{}V.\allowbreak{}UnitaryThreeInertia.\allowbreak{}matIrr\_\allowbreak{}open}{Development.\allowbreak{}AutomorphicArithmetic.\allowbreak{}AutomorphicVarieties.\allowbreak{}UnitaryThreeDeformations.\allowbreak{}IrrOpen}{340--382}
\end{proof}

\subsection*{Proof of \ref{B6.51framea}}\label{prf:B6.51framea}\phantomsection
\begin{proof}
All dimensions in this proof are over \(k'\). Write \(h_H^i(M)=\dim H^i(H,M)\), and use a place in the subscript for its absolute local group. Enlarge \(S\) to \(S'\) by adjoining the finite primes ramified in the field fixed by \(K\). Put \(G=G_{F,S'}\). There is a surjection \(G\to\Gamma\). Its kernel \(J\) is topologically normally generated by the images of \(J_q=I_q\cap K\) for \(q\in S'\setminus S\), with one prolongation per place. Indeed outside \(S'\) the whole inertia group lies in \(K\), and all prolongations of a given place are conjugate.
\leanref{UpperTheory.\allowbreak{}V.\allowbreak{}GC.\allowbreak{}levelS}{Development.\allowbreak{}AutomorphicArithmetic.\allowbreak{}AutomorphicVarieties.\allowbreak{}GaloisCochains.\allowbreak{}LevelS}{20--22}

Fix an added place \(q\), and put \(D_q=G_{F_q}\). The group \(D_q/J_q\) is an extension of \(\widehat{\Z}\) by \(I_q/J_q\), whose order divides \([G_F:K]\). Averaging over this finite group and using the \(\ell\)-cohomological dimension one of \(\widehat{\Z}\) shows
\[
 H^2(D_q/J_q,V)=0,\qquad
 h^1_{D_q/J_q}(V)=h^0_{D_q}(V).
\]
For the second identity, Frobenius acts on \(V^{I_q/J_q}\); its coinvariant and invariant spaces have the same dimension by rank-nullity. Thus the local inflation--restriction sequence gives
\[
 \dim\operatorname{Hom}_{\rm cont}(J_q,V)^{D_q}
 =h_q^1(V)-h_q^0(V)=h_q^2(V).
\]
Here \(J_q\) acts trivially on \(V\), and the last equality is the local Euler characteristic formula at \(q\nmid\ell\) for a finite \(\ell\)-primary module \cite[I, Thm.~2.8]{milne-2006}. Moreover every class inflated from \(\Gamma\) restricts to zero in \(H^2(D_q,V)\), since the local homomorphism factors through \(D_q/J_q\).
\leanref{UpperTheory.\allowbreak{}V.\allowbreak{}GC.\allowbreak{}tameEq}{Development.\allowbreak{}AutomorphicArithmetic.\allowbreak{}AutomorphicVarieties.\allowbreak{}GaloisCochains.\allowbreak{}TameEq}{117--188}

Restriction to the \(J_q\) injects \(\operatorname{Hom}_{\rm cont}(J,V)^G\) into the sum of the local equivariant homomorphism spaces just computed. A homomorphism in the kernel vanishes on all conjugates of the generators and hence, by continuity, on \(J\). Write \(t\) for the dimension of the kernel of inflation \(H^2(\Gamma,V)\to H^2(G,V)\). The five-term sequence for \(J\to G\to\Gamma\) gives
\[
 t+h_G^1(V)-h_\Gamma^1(V)
 =\dim\operatorname{Hom}_{\rm cont}(J,V)^G
 \le\sum_{q\in S'\setminus S}h_q^2(V).
\]
In particular \(H^1(\Gamma,V)\) and \(H^2(\Gamma,V)\) are finite, by finite-coefficient global finiteness for \(G\).
\leanref{UpperTheory.\allowbreak{}V.\allowbreak{}GC.\allowbreak{}inflKer}{Development.\allowbreak{}AutomorphicArithmetic.\allowbreak{}AutomorphicVarieties.\allowbreak{}GaloisCochains.\allowbreak{}InflKer}{146--206}

Let \(s_\Gamma\) denote the dimension of the kernel of localization in degree two from \(\Gamma\) to the full local groups at \(S\), with coefficient \(V\). Define \(s_G\) using \(G\) and \(S'\). By the vanishing at the added places, inflation maps the first kernel to the second, so \(s_\Gamma\le t+s_G\). The cyclotomic character factors through \(\Gamma\), since it is unramified outside \(S\). The dual-invariant hypothesis therefore holds over \(G\) as well. For the underlying \(\F_\ell\)-module, composition with the field trace \(k'\to\F_\ell\) identifies the \(k'\)-linear cyclotomic dual with the \(\F_\ell\)-linear Tate dual. Apply Theorem~\ref{B6.09} to \(G\), with \(r=1\). Its last term vanishes, giving
\[
 s_G=h_G^2(V)-\sum_{q\in S'}h_q^2(V).
\]
Combining the last two displays gives
\[
 s_\Gamma+\sum_{v\in S}h_v^2(V)
 \le h_\Gamma^1(V)-h_G^1(V)+h_G^2(V).
\]
The global Euler characteristic formula \cite[I, Thm.~5.1]{milne-2006} applies since \(S'\) contains all places over \(\ell\), the action factors through \(G\), and \(\ell\) is odd. Writing \(c_w\) for complex conjugation at a real place and \(d=[F:\Q]\), it gives
\[
 h_G^2(V)-h_G^1(V)
 =\sum_{w\mid\infty}\dim V^{c_w=1}-d\dim V-h_\Gamma^0(V).
\]
\leanref{UpperTheory.\allowbreak{}V.\allowbreak{}GC.\allowbreak{}H2\_\allowbreak{}galois\_\allowbreak{}quotient\_\allowbreak{}eq}{Development.\allowbreak{}AutomorphicArithmetic.\allowbreak{}AutomorphicVarieties.\allowbreak{}GaloisCochains.\allowbreak{}H2GaloisQuotientEq}{87--165}

Put \(Q_v=V/b_v\), and let \(s_b\) be the kernel dimension of \(H^2(\Gamma,V)\to\bigoplus_{v\in S}H^2(F_v,Q_v)\). On each local group the coefficient map \(H^2(F_v,V)\to H^2(F_v,Q_v)\) is onto by local Tate duality: the dual map is the inclusion of the invariant duals. The additional kernel after composition has dimension at most the sum of the local kernel dimensions. Thus
\[
 s_b+\sum_{v\in S}h_v^2(Q_v)
 \le s_\Gamma+\sum_{v\in S}h_v^2(V).
\]
The long exact sequence of the shifted cone, through the image in \(H^2(\Gamma,V)\), proves finiteness in the asserted degrees and gives
\[
 \dim H_b^2-\dim H_b^1+\dim H_b^0
 =s_b-h_\Gamma^1(V)+h_\Gamma^0(V)
       +\sum_{v\in S}\bigl(h_v^1(Q_v)-h_v^0(Q_v)\bigr).
\]
At \(v\mid\ell\) the local Euler formula gives
\(h_v^1(Q_v)-h_v^0(Q_v)=h_v^2(Q_v)+d_v\dim Q_v\).
Substitution and \(\sum_{v\in S}d_v=d\) leave exactly
\(\sum_{w\mid\infty}\dim V^{c_w=1}-\sum_{v\in S}d_v\dim b_v\) as an upper bound.
\leanref{UpperTheory.\allowbreak{}V.\allowbreak{}GC.\allowbreak{}H2\_\allowbreak{}galois\_\allowbreak{}quotient\_\allowbreak{}surjective}{Development.\allowbreak{}AutomorphicArithmetic.\allowbreak{}AutomorphicVarieties.\allowbreak{}GaloisCochains.\allowbreak{}H2GaloisQuotientSurjective}{50--78}
\leanref{UpperTheory.\allowbreak{}V.\allowbreak{}GC.\allowbreak{}coneCount}{Development.\allowbreak{}AutomorphicArithmetic.\allowbreak{}AutomorphicVarieties.\allowbreak{}GaloisCochains.\allowbreak{}ConeCount}{349--375}
\end{proof}

\subsection*{Proof of \ref{B6.51aa}}\label{prf:B6.51aa}\phantomsection
\begin{proof}
First introduce a matrix \(Y\) and a column \(v\) of independent variables over \(S\). In \(S[Y,v]\) let \(J\) be generated by the entries of \(Yv\) and by \(\det Y\). We determine its radical. Write \(e_1\) for the first coordinate vector and \(E=\operatorname{diag}(0,1,\ldots,1)\). With independent \(g\)-by-\(g\) matrices \(N,U\) and a scalar variable \(t\), make the polynomial substitution
\[
 Y=N E\operatorname{adj}(U),\qquad v=tUe_1.
\]
It kills \(J\), since \(Ee_1=0\), \(\operatorname{adj}(U)U=(\det U)1_g\), and \(\det E=0\). Its kernel is prime since its target is a polynomial ring over the domain \(S\).

Let \(\mathfrak q\) be any prime over \(J\), and pass to the fraction field of \(S[Y,v]/\mathfrak q\). If \(v\ne0\), take \(w=v\) and \(t=1\). If \(v=0\), choose a nonzero \(w\in\ker Y\) and take \(t=0\). Complete \(w\) to the columns of an invertible \(U\), with \(w\) first, and set \(N=(\det U)^{-1}YU\). Since the first column of \(YU\) vanishes, these choices give precisely the displayed substitution. Thus its kernel lies in \(\mathfrak q\). The kernel is consequently \(P=\sqrt J\), and \(v_1\notin P\) since its image is \(tU_{11}\).
\leanref{UpperTheory.\allowbreak{}V.\allowbreak{}CA.\allowbreak{}g\_\allowbreak{}detrad}{Development.\allowbreak{}AutomorphicArithmetic.\allowbreak{}AutomorphicVarieties.\allowbreak{}CommutativeAlgebraHeight.\allowbreak{}GDetrad}{111--163}

Grade the polynomial ring by total degree in \(Y,v\), with \(S\) in degree zero. The ideal \(J\) is homogeneous, so \(P\) is homogeneous. In \(T=S[[Y,v]]\), a series belongs to \(PT\) precisely when each homogeneous part belongs to \(P\). Indeed, choose a finite homogeneous generating set of \(P\), express each part in those generators with homogeneous coefficients, and collect the coefficients degree by degree. For two series outside \(PT\), the product of their first homogeneous parts surviving modulo \(P\) survives, since \(S[Y,v]/P\) is a domain. Thus \(PT\) is prime. A power of \(P\) is contained in \(J\) by Noetherianity, whence \(PT=\sqrt{JT}\), and \(v_1\notin PT\).

Choose a prime \(Q\subseteq PT\) minimal over the \(g\) entries of \(Yv\) in \(T\). The adjugate identity gives \((\det Y)v_1\in Q\). Since \(v_1\notin Q\), we have \(\det Y\in Q\), hence \(Q=PT\). Krull's height theorem gives \(\operatorname{ht}(PT)\le g\).
\leanref{UpperTheory.\allowbreak{}V.\allowbreak{}CA.\allowbreak{}g\_\allowbreak{}detps}{Development.\allowbreak{}AutomorphicArithmetic.\allowbreak{}AutomorphicVarieties.\allowbreak{}CommutativeAlgebraHeight.\allowbreak{}GDetps}{45--90}

There are \(c=g^2+g\) new variables. The complete regular local ring \(T\) has dimension \(1+a+g+c\), so the complete local domain \(T/PT\) has dimension at least \(1+a+c\). In that domain impose
\[
 Y_{ij}-Z_{ij},\qquad v_j-u_j,\qquad f_1,\ldots,f_b.
\]
These \(c+b\) elements lie in the maximal ideal. Evaluation at \(Y=Z,\ v=u\) is defined by completeness. Its kernel in \(T\) is generated by the displayed differences, as one sees by translating the formal variables. Consequently this section and \(\operatorname{Spec}A\) have the same underlying closed space. Krull's height theorem and catenarity give component dimension at least \(1+a+c-(c+b)=D+1\). Grothendieck's connectedness theorem for a section of a complete local domain gives connectedness after removal of any closed subset of dimension less than \(D\) \cite[Exp.~XIII, Thm.~2.1]{grothendieck-1968}. If the asserted graph of minimal primes were disconnected, partition its vertices into two nonempty sets with no edge between them. The two unions of components would meet in a finite union of closed sets of dimension less than \(D\). Removing that union would disconnect the section. This proves the chain assertion.
\end{proof}

\subsection*{Proof of \ref{B6.51ab}}\label{prf:B6.51ab}\phantomsection
\begin{proof}
Write \(k'\) for the residue field. On both the relative cotangent space and the space of minimal relations the parity projections are \((1\pm\sigma)/2\), where \(\sigma\) is the involution. They exist since \(\ell\) is odd. Lift eigenbases using these projections. Complete Nakayama gives a surjection \(S=O'[[x,u]]\twoheadrightarrow R\) with even variables \(x\) and odd variables \(u\), and generators of its kernel \(J\) of definite parity. Write \(t_+=a,t_-=g\) for the variable counts and \(r_+,r_-\) for the minimal relation counts. The ideal of the fixed locus is generated by the \(u_j\), since taking the quotient by them makes \(\sigma\) the identity, and \(2u_j=u_j-\sigma(u_j)\). In matrix coordinates the same ideal is generated by the entries between the middle line and the outer corner. The character coordinates are even.

For parity \(\delta\), let \(V_\delta\) be the corresponding subspace of the residual adjoint matrices and \(b_{\delta,v}\) its subspace preserving the ordinary flag. Write \(h^i_\delta=\dim_{k'}H_b^i(V_\delta)\), with the cochain complex of Definition~\ref{B6.51frame}. The tangent and relation counts satisfy
\[
 r_\delta\le h^2_\delta,\qquad
 h^1_\delta\le t_\delta,\qquad h^1_++1\le t_++h^0_+ .
\]
Here is the deformation calculation. For a square-zero extension with kernel killed by the maximal ideal, lift the group matrices and the local flag bases. The defect of multiplication is a two-cochain. Modulo the flag-preserving matrices the local defects are its local primitives. Associativity gives the cocycle identity in the shifted cone. Changing the lifts adds a boundary, and a primitive of this cocycle corrects the lifts to a representation with the required flags. Apply this obstruction to the extension \(S/\mathfrak m_SJ\to S/J\), pushing its kernel out by a linear functional of parity \(\delta\) on \(J/\mathfrak m_SJ\). A vanishing obstruction would give a section of the pushout by universality. Changing the variable lifts by that square-zero kernel leaves every relation value fixed, since all relative linear coefficients of a relation vanish in \(k'\) in a minimal presentation. The value of a relation in the pushout is its image under the chosen functional. Hence a functional with zero obstruction is zero. This gives the first inequality.

A first-order deformation gives a cocycle and the first-order local flag changes. Conversely a cone cocycle gives the matrices \( (1+\epsilon z(s))\bar r(s)\) with their corrected flags over \(k'[\epsilon]/(\epsilon^2)\), up to infinitesimal change of basis. Put the marking in its prescribed normal form and project to the indicated parity. The local flag changes modulo their stabilizers are determined by the cocycle: a flag change fixed modulo its stabilizer has zero class in the tangent space to the flag variety, since a local marking element has distinct eigenvalues on the graded lines. Thus framed tangents map onto \(H_b^1\) in each parity. In even parity conjugation by \(1+\epsilon\operatorname{diag}(1,0,-1)\) preserves the marking and gives a coboundary. Either this gives a nonzero tangent in the kernel, or its nonzero generator fixes the residual representation and its flags and belongs to \(H_b^0\). This proves the last two inequalities.
\leanref{UpperTheory.\allowbreak{}V.\allowbreak{}Fr.\allowbreak{}ep\_\allowbreak{}mazur}{Development.\allowbreak{}AutomorphicArithmetic.\allowbreak{}AutomorphicVarieties.\allowbreak{}FrobeniusSmoothVectors.\allowbreak{}EpMazur}{154--284}
\leanref{UpperTheory.\allowbreak{}V.\allowbreak{}Fr.\allowbreak{}ep\_\allowbreak{}tan}{Development.\allowbreak{}AutomorphicArithmetic.\allowbreak{}AutomorphicVarieties.\allowbreak{}FrobeniusSmoothVectors.\allowbreak{}LieTMem}{29--78}

In an eigenbasis of the involution adapted to a local flag, two graded lines have one parity and one has the other. Among the upper triangular matrix entries, four have even parity and two have odd parity. At a real place the fixed adjoint matrices identify, by the polarizing form, with alternating forms. Their even part has dimension one, from the exterior square of the corner, and their odd part has dimension two, from pairing the corner with the middle line. On the sign kernel the composition characters of the residual adjoint are powers of \(\bar\Psi\). The condition on \(\bar\varepsilon\) in Definition~\ref{B6.51frame} therefore excludes a cyclotomic quotient in either parity. Applying \eqref{eq:p6c5-frame-euler} to the two spaces gives
\[
 h^2_++h^0_++3d\le h^1_+,\qquad
 h^2_-+h^0_-\le h^1_- .
\]
It follows that \(r_++3d+1\le a\) and \(r_-\le g\). Append zero relations to bring the two lists to lengths \(a-(3d+1)\) and \(g\).

Finally, if \(g=0\), the block ideal vanishes in the completion. Faithful flatness over the local ring makes that ideal vanish in the local ring as well. It is finitely generated, so it vanishes on a neighborhood of the point. This proves the assertion about \(g\).
\leanref{UpperTheory.\allowbreak{}V.\allowbreak{}Fr.\allowbreak{}gpos}{Development.\allowbreak{}AutomorphicArithmetic.\allowbreak{}AutomorphicVarieties.\allowbreak{}FrobeniusSmoothVectors.\allowbreak{}AnnOfFlat}{48--72}
\end{proof}

\subsection*{Proof of \ref{B6.51a}}\label{prf:B6.51a}\phantomsection
\begin{proof}
Use \(\iota\) to regard the finite seed characters as Dirichlet characters. Apply Lemma~\ref{B6.49h}(1) first to \(\phi=\chi'^2\eta'\) and then to \(\phi=\chi'\eta'\). In each application take \(D\mid D'\) to be the primitive conductor of \(\phi\). The first character is odd and the second is nontrivial by \(\chi'\ne\eta'\), so \(D>1\). In both cases \(\ell w\) is prime to \(D\), \(\phi(w)=1\), and \(\theta\) has conductor \(w\) and order \(\ell\), so \(\ell\mid w-1\). Take \(m=j+(\ell-1)t\) with \(t\ge0\) large enough for both weight thresholds of Lemma~\ref{B6.49h}(1) and \(m\ge2\). Use \(h=2m+1\) in the first application and \(h=m+1\) in the second; the parities of \(\phi\) are respectively \(-1\) and \((-1)^{j+1}\). In each case Lemma~\ref{B6.49h}(1) supplies the normalized cuspidal newform of weight \(h\), level \(Dw\) and nebentype \(\phi\theta\), together with the continuous irreducible representation \(\rho:G_\Q\to\GL_2(C_\ell)\) of determinant \(D_h=\varepsilon_\ell^{h-1}\phi\theta\), residual polynomial \((X-1)(X-\bar\Phi_h)\) for \(\Phi_h=\varepsilon_\ell^{h-1}\phi\), inertia type \((1,\phi\theta)\) away from \(\ell\), and ordinary flag at \(\ell\) of inertia weights \((h-1,0)\) and decomposition-group residual labels \((\bar\Phi_h,1)\).

For the first form put \(\Xi=\varepsilon_\ell^m\chi'\) and \(\widetilde r_2=\rho\Xi^{-1}\), with \(r_2=\widetilde r_2|_{G_L}\). Here \(\Phi_h=\Xi^2\eta'\), and on \(G_L\) we have \(\eta'=\theta=1\) and \(\bar\Xi=\bar\Psi\). At the seed element \(s\) of \ref{B6.51frame}, \(\bar\Phi_h(s)=\bar\Psi(s)^2\ne1\). Lemma~\ref{B6.49h}(1) gives irreducibility of \(r_2\). Its residual labels are \((\bar\Psi,\bar\Psi^{-1})\) and its tame types are \((\chi',\chi'^{-1})\). Its weights are \((m,-m)\). The determinant of \(\widetilde r_2\) is \(\eta'\theta\). For \(J=\left(\begin{smallmatrix}0&1\\-1&0\end{smallmatrix}\right)\) and \(S_c=\widetilde r_2(c)\), one has \(S_c^2=1\), \(\det S_c=-1\), and \(JS_c\) is symmetric. On \(G_L\) the identity \(R^tJR=(\det R)J\) gives the conjugate pairing
\[
 r_2(g)^t JS_c\,r_2(cgc)=JS_c.
\]

For the rank-two automorphic normalisation, since \(\ell\nmid D_{K_0}\), Lemma~\ref{B6.49h}(2) supplies an algebraic Hecke character of \(K_0\) of type \((-1,0)\) with \(\ell\)-adic realization \(\kappa_0\) of conjugate product \(\varepsilon_\ell^{-1}\), unramified off \(\ell D_{K_0}\), of prime-to-\(\ell\) inertia exponent away from \(\ell\), and with rational-idelic transfer \(\eta'\varepsilon_\ell^{-1}\) from \(G_\Q\). Restrict \(\kappa_0\) to \(G_L\) to obtain \(\kappa_2\). Quadratic base change to \(K_0\), the scalar twists, and abelian base change to \(L\) give the attachment of \(r_2\kappa_2\) by \cite{arthur-clozel-1989}. The irreducible restriction to \(G_L\) gives cuspidality throughout. The choice of \(\kappa_2\) is fixed in both Hecke comparisons and for every rank-two point.

By Lemma~\ref{B6.49h}(1), the initial newform has spherical components outside \(Dw\) and unlinked principal-series support at \(p\mid Dw\), with one unramified character and one of conductor exponent \(v_p(Dw)\). After a common global norm twist to the unitary normalisation, the support exponents lie in \((-1/2,1/2)\) by the local bound used in \ref{B9.38a}. The local base-change comparisons pull the two characters back by the norm. A place of residue degree \(f\) raises their uniformiser moduli to the \(f\)-th power and replaces \(q\) by \(q^f\), so the exponents are preserved. An algebraic scalar twist adds one common exponent at every finite place. This gives \eqref{eq:p6c5-toral-support} for \(r_2\kappa_2\) at every finite place of \(L\) with one \(t\). On an open subgroup of ordinary inertia the ratio of any ordered pair of flag characters is a cyclotomic power of exponent \(0\) or \(\pm2m\). None of these exponents equals one, and the cyclotomic character has infinite image there. Thus no positive power of a ratio equals the corresponding power of \(\varepsilon_\ell\). Every place above \(\ell\) is potentially unlinked, which completes the rank-two witness.
\leanref{UpperTheory.\allowbreak{}V.\allowbreak{}Sd.\allowbreak{}twist\_\allowbreak{}of\_\allowbreak{}printed\_\allowbreak{}results}{Development.\allowbreak{}AutomorphicArithmetic.\allowbreak{}AutomorphicVarieties.\allowbreak{}ModularityArgument.\allowbreak{}ThresholdArguments.\allowbreak{}DenseThresholdCases.\allowbreak{}TTwist}{111--178}
\leanref{UpperTheory.\allowbreak{}V.\allowbreak{}Z.\allowbreak{}Bc.\allowbreak{}second\_\allowbreak{}initial\_\allowbreak{}of\_\allowbreak{}prints}{Development.\allowbreak{}AutomorphicArithmetic.\allowbreak{}AutomorphicVarieties.\allowbreak{}ModularityArgument.\allowbreak{}ThresholdArguments.\allowbreak{}DenseThresholdCases.\allowbreak{}FinOrderOfKills}{296--453}

For the second form \(\bar\Phi_h|_{G_L}=\bar\Psi|_{G_L}\), so \(\rho\) stays irreducible on every open subgroup by Lemma~\ref{B6.49h}(1) and the square test. Put \(r_3=(\operatorname{Sym}^2\rho\otimes\det(\rho)^{-1})|_{G_L}\). For ordinary diagonal characters \(a,b\) of \(\rho\), the new characters are \(a/b,1,b/a\). The resulting residual labels are \((\bar\Psi,1,\bar\Psi^{-1})\) and the weights are \((m,0,-m)\). The tame types are \((\chi',1,\chi'^{-1})\). The discriminant of binary quadratics supplies an invariant nondegenerate symmetric form. An invariant plane in the three-dimensional vector space cuts the conic \([x^2:2xy:y^2]\) in a divisor of degree two. Its support gives a finite invariant set of lines in \(\rho\), so an open subgroup fixes one of them. An invariant line in \(r_3\) gives an invariant plane by orthogonal complementation. Thus \(r_3\) stays irreducible on every open subgroup of \(G_L\).

Base-change the second form to \(L\) and unitarise it. A nontrivial self-twist is quadratic by determinants. The corresponding intertwiner is scalar on the quadratic kernel by open irreducibility, so the self-twist character is trivial. Gelbart--Jacquet gives the cuspidal adjoint lift \cite[Thm.~9.3]{gelbart-jacquet-1978}. The scalar polarising twist on the base change cancels in its adjoint. The arithmetic normalisation gives the attachment for \(r_3\varepsilon_\ell^{-1}\).

The rank-two local comparison just used applies to the second newform as well. Its base change and algebraic polarising twist give toral support with exponents in a half-unit interval about a common exponent. In the adjoint comparison a pair of local characters \(\xi_1,\xi_2\) gives the unitary-normalised support \((\xi_1/\xi_2,1,\xi_2/\xi_1)\). At a place over \(M\), use the comparison with that place in the split set. Its local nonlinking test excludes \((\xi_1/\xi_2)^j=|\cdot|^{\pm1}\) for \(j=1,2\). Away from \(M\) the comparison gives spherical support. The good Frobenius polynomials identify the attachments from these comparisons by strong multiplicity one. The unitary-normalised adjoint attachment therefore has unlinked toral support at every finite place. Apply the strict half-unit bound of \ref{B9.38a} and undo the global norm twist to obtain \eqref{eq:p6c5-toral-support}. On an open subgroup of ordinary inertia, the ratios are cyclotomic powers with exponents given by differences of \(m,0,-m\). Since \(m\ge2\), these differences are distinct from one. The cyclotomic character has infinite image there, so the same comparison holds after taking positive powers. This gives potential unlinking above \(\ell\) and completes the rank-three witness.
\leanref{UpperTheory.\allowbreak{}V.\allowbreak{}Sd.\allowbreak{}ad0Irr}{Development.\allowbreak{}AutomorphicArithmetic.\allowbreak{}AutomorphicVarieties.\allowbreak{}IdeleGenerators.\allowbreak{}HolomorphicIrreducibility.\allowbreak{}Ad0IrrMulVec}{213--303}
\leanref{UpperTheory.\allowbreak{}V.\allowbreak{}Z.\allowbreak{}Sym.\allowbreak{}W.\allowbreak{}seedInW\_\allowbreak{}alt}{Development.\allowbreak{}AutomorphicArithmetic.\allowbreak{}AutomorphicVarieties.\allowbreak{}ModularityArgument.\allowbreak{}ThresholdArguments.\allowbreak{}DenseThresholdCases.\allowbreak{}TSeedInW}{117--270}
\leanref{UpperTheory.\allowbreak{}V.\allowbreak{}Z.\allowbreak{}Sym.\allowbreak{}W.\allowbreak{}adjoint\_\allowbreak{}lift\_\allowbreak{}from\_\allowbreak{}eisenstein\_\allowbreak{}of\_\allowbreak{}propagation}{Development.\allowbreak{}AutomorphicArithmetic.\allowbreak{}AutomorphicVarieties.\allowbreak{}ModularityArgument.\allowbreak{}ThresholdArguments.\allowbreak{}DenseThresholdCases.\allowbreak{}AdjointLiftFromEisensteinSeedOfHyps}{184--281}
\leanref{UpperTheory.\allowbreak{}V.\allowbreak{}Z.\allowbreak{}Sym.\allowbreak{}W.\allowbreak{}adjoint\_\allowbreak{}automorphy\_\allowbreak{}w\_\allowbreak{}of\_\allowbreak{}pure\_\allowbreak{}weight\_\allowbreak{}adjoint\_\allowbreak{}lift}{Development.\allowbreak{}AutomorphicArithmetic.\allowbreak{}AutomorphicVarieties.\allowbreak{}ModularityArgument.\allowbreak{}ThresholdArguments.\allowbreak{}DenseThresholdCases.\allowbreak{}MemOfUnder}{101--114}

For either seed, the excluded inertia away from \(\ell\) has \(\bar\Psi=1\), hence \(\chi'=1\) there by its prime-to-\(\ell\) order. The tame types give identity inertia matrices on this subgroup and hence factorisation through the frame. The residual determinant law is the product of the prescribed characters. Each factor of the ordered ordinary product relations sends the corresponding flag step into the preceding one, in the order of action on a column. After the twist the pairing has multiplier \(\varepsilon_\ell^{1-n}\). This gives a determinant point in \ref{B6.48c}. Its separated weights \(mw_{n,i}\) and cuspidal witness put it on the Hecke quotient by \ref{B6.54a}.
\leanref{UpperTheory.\allowbreak{}V.\allowbreak{}Pr.\allowbreak{}seedPtOf\_\allowbreak{}alt}{Development.\allowbreak{}AutomorphicArithmetic.\allowbreak{}AutomorphicVarieties.\allowbreak{}ModularityArgument.\allowbreak{}ThresholdArguments.\allowbreak{}DenseThresholdCases.\allowbreak{}PrSeedPtOf}{255--258}

Fix a Hecke comparison in each rank. Write \(B_{0,n}\) for its universal ordinary determinant ring and \(J_{0,n}\) for the Hecke ideal. Write \(A_n\) for the framed matrix ring before the link and \(B_n\subseteq A_n\) for the coefficient ring of weight-zero data. Let \(M_n\) be the ideal on \(B_n\) generated from \(J_{0,n}\) by the link. Denote the upper and lower triangular loci by the vanishing of the lower and upper entries, respectively. A prime of type \(I\) is off both loci and, in rank three, off the block locus where the middle line splits from the outer two-dimensional corner. A prime of type \(E\) is on this block locus and off both triangular loci.

Choose a minimal Hecke component through the seed. Its dimension is \(1+nd\) by the finite free weight module of \ref{B6.48c}, since there are \(nd\) weight variables over the coefficient DVR. Over the algebraic closure of the component's function field its determinant has a matrix realisation. Full span at the seed gives a nonvanishing trace-pairing minor, hence full span at the generic point. Untwist and use the marked basis. The resulting type \(I\) prime contracts to a coefficient prime containing a prime \(Z\supseteq M_n\) of dimension
\begin{equation}\tag{6.51a.1}\label{eq:p6c5-linked-seed-dimension}
 \dim B_n/Z=1+nd .
\end{equation}
\leanref{UpperTheory.\allowbreak{}V.\allowbreak{}Pr.\allowbreak{}seedOf}{Development.\allowbreak{}AutomorphicArithmetic.\allowbreak{}AutomorphicVarieties.\allowbreak{}CurveDVRSpecialization.\allowbreak{}SeedOf}{501--547}

In \(B_n\) let \(\mathfrak l\) impose the coefficient uniformiser \(\varpi\) and triviality of the flag characters on \(\ker\bar\Psi\) in inertia. Complete \(A_n\) at a residual point off the upper triangular locus and denote the resulting local ring by \(R\). Suppose a quotient of \(R\) has dimension at least \(nd+1\) and its contraction contains \(I\subseteq B_n\) satisfying \(\mathfrak m_{B_n}\subseteq\sqrt{I+\mathfrak l}\). For a specified pair of flag characters, coincidence at \(v\) means equality on \(N_v\). The dimension bounds are
\begin{equation}\tag{6.51a.2}\label{eq:p6c5-linked-exclusions}
 \begin{aligned}
 \dim(\text{lower triangular locus})&\le(n-1)d+2+3C_0+5<nd,\\
 \dim(\text{special-fibre coincidence at }v
 \text{ over }I\text{ off both triangular loci})&\le nd+5-d_v<nd.
 \end{aligned}
\end{equation}
For the first bound, the diagonal parameters consist of the uniformiser, \(d+1+C_0\) parameters in the first direction and \(d\) inertia parameters in each of the other \(n-2\) free directions. The lower cocycle fibre contributes at most \(2C_0\) by \eqref{eq:p6c5-frame-cohomology-bounds}, and the marked normalisation uses at most five further coordinates. For the second bound, equal characters on \(N_v\) eliminate the corresponding inertia parameters of one direction. Retain the inertia value outside the specified generators in \(N_v\) and use the radical condition on \(B_n\). The fibre off the triangular loci adds at most one dimension. The inequalities \(d>3C_0+7\) and \(d_v>5\) give the strict bounds in \eqref{eq:p6c5-linked-exclusions}. Apply Lemma~\ref{B6.48bb} with \(b=nd\), with the ideal of the given quotient as its initial ideal, and with the coefficient uniformiser as \(s\). Take the lower triangular ideal for the first list. For the second list take the sum of the extended coefficient ideal \(I\), the uniformiser ideal and the coincidence ideal for each place and each pair of distinct flag characters. The upper triangular locus is empty in this completion. The two lines of \eqref{eq:p6c5-linked-exclusions} give the required strict bounds, respectively, and the assumed quotient dimension is at least \(b+1\). This gives a one-dimensional prime containing the initial quotient ideal, off both triangular loci, with separation on every \(N_v\).
\leanref{UpperTheory.\allowbreak{}V.\allowbreak{}Fr.\allowbreak{}l51}{Development.\allowbreak{}AutomorphicArithmetic.\allowbreak{}AutomorphicVarieties.\allowbreak{}FrobeniusSmoothVectors.\allowbreak{}CharpolyCoeffsInSubalgebra}{50--85}
\leanref{UpperTheory.\allowbreak{}V.\allowbreak{}Fr.\allowbreak{}l52}{Development.\allowbreak{}AutomorphicArithmetic.\allowbreak{}AutomorphicVarieties.\allowbreak{}FrobeniusSmoothVectors.\allowbreak{}TameInertiaCondition}{257--281}
\leanref{UpperTheory.\allowbreak{}V.\allowbreak{}Pr.\allowbreak{}exists\_\allowbreak{}nice\_\allowbreak{}prime\_\allowbreak{}above\_\allowbreak{}of\_\allowbreak{}large\_\allowbreak{}dim}{Development.\allowbreak{}AutomorphicArithmetic.\allowbreak{}AutomorphicVarieties.\allowbreak{}CurveDVRSpecialization.\allowbreak{}ExistsNicePrimeAboveOfLargeDim}{7--70}

Normalise the curve over a finite residue extension and apply the link. Write \(H\) for the CM subgroup of the level group and \(H'=H\cap\ker\varepsilon_\ell\). The differences on \(N_v\) give distinct ordinary characters and unlinking, since the cyclotomic character is one there. Their nontrivial ratios are principal units of the equal-characteristic DVR and have infinite order. One can choose a separating element on \(N_v\) by avoiding the finitely many kernels of the ratios in the group generated by pairwise separating elements. In prime rank, reducibility of an irreducible representation on \(H'\) gives either scalar multiplicity or a transitive permutation of lines. An element with distinct eigenvalues in \(H'\) excludes scalar multiplicity. In the permutation case the ordinary flag forces the local group to fix the lines distinguished by the separating element. Since the quotient \(H/H'\) is abelian, the stabilisers are equal and the residual character test on the local group rules out the permutation case. For a type \(E\) curve the opposite entries in the two-dimensional corner give irreducibility, and the same argument applies to that corner. It has no intertwiner with the middle character. The prime-rank restriction argument applies also to \(H'\cap V\) for every open normal \(V\le H\): take a positive power of the same separating element in \(V\). More explicitly, the translates of a simple submodule span, so a reducible restriction in prime rank is scalar or gives a transitive permutation of lines. The separating power excludes scalars and identifies the lines with its eigenlines. The ordinary decomposition group fixes each line because it preserves their complete ordered flag. On the permutation kernel the reduced line characters and prescribed residual characters all coincide by Lemma~\ref{B6.49i}(2). This contradicts the two distinct residual characters on the decomposition group, also for the outer corner since \(\bar\Psi^2\ne1\) at its test element. Thus every open subgroup of \(H'\) acts absolutely irreducibly on the indicated prime-rank representation. In particular its image is noncommutative.

Take an open coset in \(H'\). In type \(I\) apply Lemma~\ref{B6.49g} in rank two or three. In type \(E\) apply it to the outer two-dimensional block in the same coset and include the nonzero middle character. In each case obtain a full matrix \(M\) with all principal minors nonzero. In rank three its two-row principal minors are \(\det(M)(M^{-1})_{ii}\), indexed by the omitted row. In the block case the minors are products of principal minors of the outer block and, when present, the middle scalar. Take a power of the separating element in the subgroup defining the coset, and a further power to kill torsion in the group generated by its eigenvalues \(a_i\). Their ratios have infinite order, including those involving the middle line. Choose a homomorphism from this finitely generated free abelian group to \(\Z\) taking distinct values \(w_i\) on the \(a_i\), by avoiding the finitely many hyperplanes defined by their differences. Order the indices by decreasing \(w_i\). In the degree-\(j\) elementary coefficient of \(M\operatorname{diag}(X^{w_i})\), the unique leading term comes from the first \(j\) indices, and its coefficient is their nonzero principal minor. The Newton polygon at infinity has successive segments of length one and distinct slopes. Hence the characteristic polynomial is separable. In its multivariate discriminant group the monomials by equal products of the \(a_i\). At least one coefficient sum is nonzero because the substitution \(X^{w_i}\) has nonzero discriminant and equal products have equal weights. If all evaluations at \((a_i^m)_i\) vanished, the Vandermonde equations for \(m=0,\ldots,b-1\), where \(b\) is the number of distinct products, would kill every coefficient sum. One of these powers therefore gives a separable element in the specified coset. An open subset of \(H\) meeting \(H'\) contains such a coset in its intersection with \(H'\).
\leanref{UpperTheory.\allowbreak{}V.\allowbreak{}Pr.\allowbreak{}regCosetE}{Development.\allowbreak{}AutomorphicArithmetic.\allowbreak{}AutomorphicVarieties.\allowbreak{}CurveDVRSpecialization.\allowbreak{}RegCosetE}{174--191}
 At the distinguished decomposition group \(D_0\), Lemma~\ref{B6.49i}(1) gives the zero intertwiner condition. These are the curve hypotheses in \ref{B6.48aa}(a) for type \(I\) and in \ref{B6.48aa}(b) for type \(E\).
\leanref{UpperTheory.\allowbreak{}V.\allowbreak{}Pr.\allowbreak{}irrNormal}{Development.\allowbreak{}AutomorphicArithmetic.\allowbreak{}AutomorphicVarieties.\allowbreak{}CurveDVRSpecialization.\allowbreak{}IrrNormal}{62--180}
\leanref{UpperTheory.\allowbreak{}V.\allowbreak{}Pr.\allowbreak{}nice\_\allowbreak{}specialization\_\allowbreak{}twist\_\allowbreak{}satisfies\_\allowbreak{}third\_\allowbreak{}condition}{Development.\allowbreak{}AutomorphicArithmetic.\allowbreak{}AutomorphicVarieties.\allowbreak{}CurveDVRSpecialization.\allowbreak{}NiceSpecializationTwistSatisfiesThirdCondition}{164--214}

Use the whole completion in rank two and at a nonblock residual point in rank three. Put \(D=nd+1\). The ungraded tangent and obstruction calculation gives at least \(D\) more variables than relations. Indeed, the flag-preserving matrices have dimension \(n(n+1)/2\), and the fixed matrices at a real involution identify with alternating forms, of dimension \(n(n-1)/2\). The cyclotomic character on the sign kernel differs from every power of \(\bar\Psi\), so the residual adjoint has no cyclotomic quotient. Equation~\eqref{eq:p6c5-frame-euler} gives \(h^2+h^0+nd\le h^1\). The square-zero calculation in the proof of Lemma~\ref{B6.51ab}, without the parity projections, gives \(r\le h^2\) for the minimal relation count and \(h^1+1\le t+h^0\) for the variable count: the infinitesimal marked torus has nonzero generator \(\operatorname{diag}(w_{n,i})\). Hence \(t-r\ge D\). Grothendieck's connectedness theorem applied to these \(r\) equations in the complete regular local ring of dimension \(1+t\) gives the chain bound below. \cite[Exp.~XIII, Thm.~2.1]{grothendieck-1968}
\leanref{UpperTheory.\allowbreak{}V.\allowbreak{}Fr.\allowbreak{}pres}{Development.\allowbreak{}AutomorphicArithmetic.\allowbreak{}AutomorphicVarieties.\allowbreak{}FrobeniusSmoothVectors.\allowbreak{}HasObs}{26--125}

At a residual block point use Lemma~\ref{B6.51ab}. The polynomial family \((ty_0,t^2a_0+a)\) constructed in this proof specializes a nonblock point to each such point, so \(g>0\). Assign each monomial of an odd relation to the first odd variable dividing it and factor out that variable. This writes the \(g\) odd relations as the entries of \(Zu\) with \(Z\) square. A padded zero relation gives a zero row. The constant terms of the entries of \(Z\) lie in the maximal ideal of the coefficient DVR, since they are relative linear coefficients of relations in a minimal choice of variables. Cut by the image of \(\det Z\). The identity \(\operatorname{adj}(Z)Zu=(\det Z)u\) shows that every prime with a surviving odd coordinate is in the section. Apply Lemma~\ref{B6.51aa} with \(a-b=3d+1\). In the whole completions just considered and in this determinantal section, every component has dimension at least \(nd+2\), and any two minimal primes are joined by a chain with
\leanref{UpperTheory.\allowbreak{}V.\allowbreak{}Fr.\allowbreak{}oddm}{Development.\allowbreak{}AutomorphicArithmetic.\allowbreak{}AutomorphicVarieties.\allowbreak{}FrobeniusSmoothVectors.\allowbreak{}Oddm}{5--71}
\begin{equation}\tag{6.51a.3}\label{eq:p6c5-linked-adjacent}
 \dim R/(P+Q)\ge nd+1
\end{equation}
for consecutive primes \(P,Q\). In rank two use \(I=M_2\) in the curve construction, since the finite weight module gives the radical condition. The minimal primes are of type \(I\) by \eqref{eq:p6c5-linked-exclusions}. If the first contains the extended Hecke ideal, the curve above \(P+Q\) kills that ideal. Apply \ref{B6.48aa}(a) on the coefficient ring to deduce containment on the second branch.

A rank-three minimal prime of the section on the block locus is thick. In its function field the odd coordinates vanish and \(\det Z=0\) supplies a nonzero tangent vector in the kernel of \(Z\). Clear denominators. In the resulting first-order evaluation the odd relations vanish by the kernel equation, and the even relations have zero odd derivative by parity. The ordinary and inertia relations are among the imposed relations. Projecting to an off-diagonal block gives an extension of the two blocks with the determinant law of the block sum. A splitting would make the first-order off-diagonal entries a conjugation derivative. The marked normal form eliminates such a derivative, since the marking separates the outer residual characters. The surviving odd derivative therefore gives a nonsplit extension.
\leanref{UpperTheory.\allowbreak{}V.\allowbreak{}Fr.\allowbreak{}sharp3}{Development.\allowbreak{}AutomorphicArithmetic.\allowbreak{}AutomorphicVarieties.\allowbreak{}FrobeniusSmoothVectors.\allowbreak{}ChainedOfComap}{72--143}
\leanref{UpperTheory.\allowbreak{}V.\allowbreak{}Fr.\allowbreak{}thick}{Development.\allowbreak{}AutomorphicArithmetic.\allowbreak{}AutomorphicVarieties.\allowbreak{}FrobeniusSmoothVectors.\allowbreak{}ThickOfDer}{224--256}

On a type \(E\) prime the corner gives a rank-two point off both triangular loci. Once rank-two propagation holds, its coefficient ideal gives an ideal \(I_3\subseteq B_3\) contained in the coefficient contractions of type \(E\) primes, with \(\mathfrak m_{B_3}\subseteq\sqrt{I_3+\mathfrak l}\). Indeed after imposing the inertia relations the corner coefficients become residual by rank-two finiteness. The middle character becomes residual by the abelian exponent condition, and the corner coefficients and middle character generate the block's weight-zero coefficients. Use \eqref{eq:p6c5-linked-adjacent} and \eqref{eq:p6c5-linked-exclusions} at successive intersections. A type \(I\) curve through an intersection gives propagation by \ref{B6.48aa}(a) to the coefficient contraction of the second branch.

For a string of block components use \(I_3\) to take consecutive type \(E\) curves \(x_i:R\to k_i[[t]]\) as above. Let \(\beta:B_{0,3}\to R\) be the link map to the completion. At an intervening block prime \(E_i\), the same thick prime lies in \(\ker x_i\) and \(\ker x_{i+1}\). Transport its nonsplit extension through the link. It has the determinant law of the block sum and satisfies the ordered products and the inertia relations, with their polarised adjoints. It therefore gives a module for the specialised Cayley--Hamilton algebra at \(\beta^{-1}(E_i)\). The marked odd derivative remains nonzero in the transported extension. Thus the algebra is nonsemisimple and the block determinant makes the contraction nonstable. The same thick contraction lies below both \(\ker(x_i\beta)\) and \(\ker(x_{i+1}\beta)\). The endpoint type \(I\) primes contract to stable primes by full matrix span. Containment of the extended ideal at the first endpoint puts \(J_{0,3}\) in its contraction. Apply \ref{B6.48aa}(b) to the curves \(x_i\beta\) and extend the containment from the last contraction to \(R\). A string with one curve uses the same assertion with no intervening prime.
\leanref{UpperTheory.\allowbreak{}V.\allowbreak{}Fr.\allowbreak{}l53}{Development.\allowbreak{}AutomorphicArithmetic.\allowbreak{}AutomorphicVarieties.\allowbreak{}FrobeniusSmoothVectors.\allowbreak{}QuotientHomNonvanishingBound}{221--237}
\leanref{UpperTheory.\allowbreak{}V.\allowbreak{}Pr.\allowbreak{}thickOf}{Development.\allowbreak{}AutomorphicArithmetic.\allowbreak{}AutomorphicVarieties.\allowbreak{}CurveDVRSpecialization.\allowbreak{}ThickOf}{142--198}
\leanref{UpperTheory.\allowbreak{}V.\allowbreak{}Pr.\allowbreak{}rankThree\_\allowbreak{}deformationSetup\_\allowbreak{}splitting\_\allowbreak{}condition}{Development.\allowbreak{}AutomorphicArithmetic.\allowbreak{}AutomorphicVarieties.\allowbreak{}CurveDVRSpecialization.\allowbreak{}SecondThirdPrimeConditionTransfer}{27--97}

For global propagation put \(C=V(\mathfrak m_{B_n}A_n)\setminus V(\mathrm{tri}^+)\), where \(\mathrm{tri}^+\) is the ideal of the upper triangular locus. Every type \(I\) component meets \(C\). Otherwise a power of the ideal of lower entries lies in its prime plus \(\mathfrak m_{B_n}A_n\). In sufficiently negative weight each monomial contains that many lower factors. Graded Nakayama then puts powers of the lower entries in the component prime, a contradiction. The seed component specialises to a closed point of \(C\) with a completed type \(I\) branch of dimension \(\dim B_n/Z+1\). The extra dimension is the matrix fibre off the triangular loci. Thus \eqref{eq:p6c5-linked-seed-dimension} supplies the starting branch for the local propagation.

Put \(P=\bar\Psi\) and write \(g^*=cg^{-1}c^{-1}\) for \(g\in\Gamma_0\). Take a nonzero strict class for \(P^{-2}\eta\) and normalise its cocycle \(z\) by \(z(c)=0\). This is possible since the coefficient character at \(c\) is \(-1\) and the characteristic is odd. The restriction to \(\Gamma_0\) is nonzero. If it vanished, \(z(c)=0\) would give \(z=0\) on both cosets and the local square tests would give zero strict scalars. Set \(b(g)=P(g)z(g)\) on \(\Gamma_0\). If \(D_v\) is the image of the ordinary local map, the cocycle and local coboundary laws give
\[
 b(gh)=P(h)b(g)+P(g)^{-1}b(h),\qquad b(g^*)=b(g),
 \qquad b|_{D_v}=\gamma_v(P-P^{-1}),
\]
where \(\gamma_v\) is the negative of the strict scalar for \(z\) at \(v\). These are the lower-entry conditions for diagonal \(P,P^{-1}\) in rank two. The first gives matrix multiplication, the second gives the marked polarisation, and the third gives a lower unipotent conjugate of the diagonal locally.

At a prime of \(C\), the weight-zero specialisation gives a lower triangular point over its residue field. In rank two interpolate its lower entry \(a\) by \(ta+(1-t)b\). The displayed conditions are linear, so the interpolation defines a framed point over the polynomial ring in \(t\). In rank three the lower entries in positions \(10,21,20\) are \(y(g),y(g^*),a(g)\), with diagonal \(P,1,P^{-1}\). Use
\[
 (y,a)\longmapsto (ty,t^2a+(1-t)b).
\]
In the \(20\) product law, \(y(g^*)y(h)\) scales by \(t^2\), and the added term obeys the linear law of \(b\). The marking equation on \(y\) scales by \(t\). The local unipotent conjugation parameters in positions \(10,21\) scale by \(t\), and that in position \(20\) scales by \(t^2\) with the addition of \((1-t)\gamma_v\). Thus the local conditions persist. The generic point of each polynomial family lies in \(C\) and specialises to both endpoints. The endpoint at \(t=0\) is the fixed point with lower entry \(b\) in rank two or with pair \((0,b)\) in rank three. Hence \(C\) is connected in each rank.
\leanref{UpperTheory.\allowbreak{}V.\allowbreak{}Fr.\allowbreak{}preconnected\_\allowbreak{}rank\_\allowbreak{}two}{Development.\allowbreak{}AutomorphicArithmetic.\allowbreak{}AutomorphicVarieties.\allowbreak{}FrobeniusSmoothVectors.\allowbreak{}PreconnectedRankTwo}{206--264}

For density off the rank-three block locus, use a nonzero strict class for \(P^{-1}\eta\), with cocycle \(x\) normalised by \(x(c)=0\). Put \(u=Px\) on \(\Gamma_0\). Then \(u\ne0\), \(u(gh)=P(h)u(g)+u(h)\), and \(u(g^*)=P(g)^{-1}u(g)\). At \(v\) write \(u|_{D_v}=\alpha_v(P-1)\). Choose \(\delta\in k\) with
\[
 \delta(P(g_0)^2-1)=-u(g_0)(P(g_0)-1).
\]
Define
\[
 y_0=u+\delta(P-1),\qquad
 a_0=\delta^2-\frac{\delta^2}{2}(P+P^{-1})
 -\delta(1+P^{-1})u+\frac12P^{-1}u^2.
\]
Substitution gives \(y_0(gh)=P(h)y_0(g)+y_0(h)\) and
\[
 a_0(gh)=P(h)a_0(g)+y_0(g^*)y_0(h)+P(g)^{-1}a_0(h),
 \qquad a_0(g^*)=a_0(g).
\]
The equation for \(\delta\) gives \(y_0(g_0)=y_0(g_0^*)\). On \(D_v\), substituting \(u=\alpha_v(P-1)\) gives the lower unipotent conjugate of the diagonal with conjugating entries \(\alpha_v+\delta\) in position \(10\), \(\alpha_v-\delta\) in position \(21\), and \(\alpha_v^2/2-\delta\alpha_v-\delta^2/2\) in position \(20\). Thus \((y_0,a_0)\) defines a framed point. If \(y_0=0\), its values at \(g_0,g_0^*\) give \(2\delta(P(g_0)^2-1)=0\). Hence \(\delta=0\) and \(u=0\), a contradiction.

At a block point \((0,a)\), the entry \(a\) satisfies the same linear conditions as \(b\). The polynomial point \((ty_0,t^2a_0+a)\) has a nonzero \(10\) entry at its generic point and specialises at \(t=0\) to \((0,a)\). Thus the nonblock points are dense in \(C\). Their local presentations put them on type \(I\) components. The finite union of these components covers \(C\). If some lie in the Hecke locus and others fail containment, the two corresponding closed unions meet at a closed point of \(C\). Lift the two primes to the completion by flatness. Apply local propagation and contract the containment. This proves global containment in rank two and then in rank three.
\leanref{UpperTheory.\allowbreak{}V.\allowbreak{}Fr.\allowbreak{}preconnected\_\allowbreak{}rank\_\allowbreak{}three}{Development.\allowbreak{}AutomorphicArithmetic.\allowbreak{}AutomorphicVarieties.\allowbreak{}FrobeniusSmoothVectors.\allowbreak{}ExistsNonzeroLowCocycle}{144--194}
\leanref{UpperTheory.\allowbreak{}V.\allowbreak{}Pr.\allowbreak{}components\_\allowbreak{}good\_\allowbreak{}rank\_\allowbreak{}two}{Development.\allowbreak{}AutomorphicArithmetic.\allowbreak{}AutomorphicVarieties.\allowbreak{}CurveDVRSpecialization.\allowbreak{}ComponentsGoodRankTwo}{129--164}
\leanref{UpperTheory.\allowbreak{}V.\allowbreak{}Pr.\allowbreak{}components\_\allowbreak{}good\_\allowbreak{}rank\_\allowbreak{}three}{Development.\allowbreak{}AutomorphicArithmetic.\allowbreak{}AutomorphicVarieties.\allowbreak{}CurveDVRSpecialization.\allowbreak{}PropagationTheoremGL3}{213--230}
\leanref{UpperTheory.\allowbreak{}V.\allowbreak{}Pr.\allowbreak{}globProp}{Development.\allowbreak{}AutomorphicArithmetic.\allowbreak{}AutomorphicVarieties.\allowbreak{}CurveDVRSpecialization.\allowbreak{}GlobProp}{3--47}

An irreducible arithmetic evaluation is of type \(I\). The link identifies its determinant with that of \(r\kappa_n\), of weights \(a_\tau-ew_{n,i}\) at the embedding \(\tau\), where \(a_\tau\) is the weight contribution of the fixed twist. These weights increase strictly in \(i\). Apply \ref{B6.48a} with the same twist and two different exceptional places. Its comparison for a fixed attachment gives the witness at every required place.

Let \(k=O/\mathfrak m_O\). Since \(\mathfrak m_OA=0\), there is a specified embedding \(j:k\to A\). Denote the evaluation of the framed matrix ring by \(\xi:A_3\to A\). Its restriction to \(B_3\) kills \(\mathfrak l\), since the inertia characters are one on the kernel of \(\bar\Psi\). Off both triangular loci its prime kernel contains \(I_3\cap M_3\) by the two component containments. A prime containing this intersection contains at least one of the two ideals. The radical condition for that ideal and \(\mathfrak l\) puts \(\mathfrak m_{B_3}\) in the kernel. For each \(b\in B_3\), subtract an \(O\)-lift of its residue to obtain \(\xi(b)=j(\bar b)\). The diagonal entries have weight zero. Every cycle product \(r_{i_1i_2}(g)\cdots r_{i_si_1}(g)\) likewise has weight zero, by telescoping the basis weights; its residue vanishes if any two consecutive indices differ. All three two-cycles and both three-cycles therefore vanish. In the expansion of \(\det(X-r(g))\) only the identity permutation remains, giving the prescribed product.\par At a point on either triangular locus put \(d_i(g)=r_{ii}(g)\). Matrix multiplication makes each \(d_i\) a character. The character \(e_i=d_i/j(\bar\Psi^{1-i})\) takes values in \(1+\mathfrak m_A\). At each ordinary place take the inertia generators in \(\ker\bar\Psi\) specified in \ref{B6.51frame}. All three flag characters are one on them, so each generator has characteristic polynomial \((X-1)^3\). Its triangular diagonal entries are one because \(A\) is a domain. The same holds on the conjugate generators by the polarisation relation, which reverses the indices and inverts the diagonal values. Modulo \(\mathfrak m_A^a\) the principal units form a finite abelian \(\ell\)-group. The generator condition of \ref{B6.51frame} therefore gives triviality on inertia at both lifts, and its global exponent condition gives a single \(h>0\) such that \(e_i(g)^h=1\bmod\mathfrak m_A^a\) for every \(a\). Separation gives the equality in \(A\). To see that it forces \(e_i(g)=1\), write \(h=\ell^b h'\) with \(\ell\nmid h'\). In characteristic \(\ell\), \((e_i(g)^{h'}-1)^{\ell^b}=0\). Also the quotient of \(X^{h'}-1\) by \(X-1\) takes the unit value \(h'\) modulo \(\mathfrak m_A\) at \(e_i(g)\). The claimed diagonal characters and their characteristic polynomial follow.\par
\leanref{UpperTheory.\allowbreak{}V.\allowbreak{}Pr.\allowbreak{}finPtOf}{Development.\allowbreak{}AutomorphicArithmetic.\allowbreak{}AutomorphicVarieties.\allowbreak{}CurveDVRSpecialization.\allowbreak{}FinPtOf}{183--342}
\end{proof}

\subsection*{Proof of \ref{B6.52a}}\label{prf:B6.52a}\phantomsection
\begin{proof}
Untwist on \(G_L\). Twisting all characters of a flag by one character leaves their ratios unchanged. On the automorphic side the local support characters are twisted together. The absolute value of a Hecke character on the idèles is a power of the global norm, since the norm-one idèle class group is compact. Thus untwisting changes the centre \(t\) of the witness interval by one constant at all finite places. We obtain an attachment with witness for \(W|_{G_L}\).

Write \(I^0\) for the inertia subgroup in (1). Its cyclotomic image is open in the cyclotomic image of inertia, so it contains an element of infinite cyclotomic order. The two characters on \(I^0\) are distinct. The line on which \(I^0\) acts by \(\varepsilon_\ell^k\) is unique and is preserved by the whole decomposition group, since \(I^0\) is normal there and the character is conjugation invariant. Conjugating gives such a flag at every prolongation. For its two quotient characters \(\lambda_i\), an equality \(\lambda_i^a=(\lambda_j\varepsilon_\ell)^a\), \(a\ge1\), would give \(\varepsilon_\ell^{a(s_i-s_j-1)}=1\) on \(I^0\), with \(s_i,s_j\in\{k,-k\}\). The exponent is nonzero. The same test works after any finite extension because a power of the chosen inertia element remains in each open subgroup. In particular all places above \(\ell\) in \(L\) are potentially unlinked.

By \ref{B6.06f} take the rational cuspidal representation \(\pi\) attached to \(W\). Its base change through \(K\) to \(L\) is cuspidal at each step. Indeed a self-twist by a nontrivial step character would, by Frobenius comparison, give the same self-twist on the Galois representation. On the upper subgroup an intertwiner would be scalar by irreducibility there, which precludes a nontrivial self-twist. The base change and the given attachment over \(L\) have the same good polynomials and thus agree by strong multiplicity one.

At a place of \(L\) above \(\ell\) the witness puts the character moduli into one open interval with endpoint ratio \(q_v\). Consequently their ratio cannot be \(|\cdot|_v^{\pm1}\). Their local parameter is a sum of characters with zero monodromy by \ref{B9.47}. The toral clause of \ref{B9.06} applies: the polarization is the exterior-product polarization used in \ref{B6.06f}. The filtration jumps are \(\{k,-k\}\), since on an open inertia subgroup the successive characters are the indicated cyclotomic powers and restriction to a finite extension preserves labeled jumps. By \ref{B9.04} and archimedean base change the infinitesimal pair of \(\pi\) is \(\{\tfrac12+k,\tfrac12-k\}\). The cuspidal newvector correspondence on \(\GL_2/\Q\) therefore gives a newform \(f\) of weight \(2k+1\), with \(j=k\) in \ref{B6.06d}(3) \cite[Ch.~III]{jacquet-langlands-1970}. At a good prime \(p\), twisting \(W\) by \(\varepsilon_\ell^{-k}\) multiplies its inverse arithmetic Frobenius roots by \(p^k\). Its polynomial is consequently the Hecke polynomial of this \(f\), as asserted.

To check the witness, let \(K'/\Q\) be any imaginary quadratic field with \(W|_{G_{K'}}\) irreducible. The self-twist argument gives cuspidal quadratic base change of \(\pi\). Every regular algebraic cuspidal attachment to \(W|_{G_{K'}}\) equals that base change, by strong multiplicity one. We compare its cuspidal support with that of the known base change to \(L\).

At \(p\ne\ell\), (2) makes the inertia image finite and central in the local image. After Frobenius semisimplification, inertia and Frobenius diagonalize together. At \(p=\ell\), use the stable local flag above. The representation is potentially semistable: this holds over \(L\) by the toral comparison, hence over a finite extension of \(\Q_\ell\). The associated Weil representation, semisimplified, is the sum of the two rank-one parameters of the flag by exactness of potentially semistable comparison. In both cases the Weil constituents over \(\Q_p\) are characters. For the descended newform the semisimplified local comparison at every prime is \cite[Thms.~2.1--2.2]{saito-2009}; the field is \(\Q\), of odd degree, and the weight \(2k+1\) is at least three. Duals and norm twists commute with this comparison. Under local base change the Weil representation restricts, with the same normalization \cite[Chs.~I, III]{arthur-clozel-1989}. Thus the support over \(K'\) also consists of characters at every place under consideration.

Write \(\xi_{p,i}\) for the two rational support characters in the automorphic normalization of \(\pi\). For a place \(u\) of \(L\) over \(p\) the absolute values of the upper support characters on a uniformizer are \(|\xi_{p,i}(p)|^{f(u/p)}\). This follows from the norm of a uniformizer and the fact that a smooth character has absolute value one on units. The witness over \(L\) gives one \(t\) for its attachment, and \(q_u=p^{f(u/p)}\). Taking the positive \(f(u/p)\)-th root gives
\[
             p^{t-1/2}<|\xi_{p,i}(p)|<p^{t+1/2}.
\]
At \(v\mid p\) of \(K'\), raise these inequalities to \(f(v/p)\). This gives the required interval with the same \(t\) at every \(v\). The flag argument already checked potential unlinking at every prolongation above \(\ell\). These are both clauses of Definition~\ref{def:p6c5-witness}.
\end{proof}

\subsection*{Proof of \ref{B6.52}}\label{prf:B6.52}\phantomsection
\begin{proof}
Put \(k_{\rm res}=O/\mathfrak m_O\) and \(\omega=\bar\varepsilon_\ell\). Take \(M=\ell N|D_K|\) for the first count and \(M=\ell N\) for the Gaussian host. For the first count take \(H_b=G_\Q\). For the Gaussian host take \(H_b\) to be the kernel of the \(\ell\)-primary part of \(\vartheta\) with its factor at \(\ell\) removed. Its quotient is an abelian \(\ell\)-group unramified at \(\ell\). On inertia in \(H_b\) away from \(\ell\), the value of the prime-to-\(\ell\) part determines the value of \(\Psi\). Write \(K_0=\Q(\sqrt{-D_0})\) for the sign field, with \(D_0>0\) squarefree. The two input tests are
\begin{equation}\tag{6.52.2}\label{eq:p6c5-link-input-tests}
 \ker\omega\cap\ker\bar\Psi\not\subseteq\ker\eta,\qquad
 \exists s\in D_B\cap\ker\eta,\quad\bar\Psi(s)^2\ne1.
\end{equation}
For the first count these are the sign-kernel and local unit conditions of the good datum. For the Gaussian host the first uses \(c_2(\vartheta)\ge4\) and the second is the decomposition test from \eqref{eq:p6c5-gaussian-host}. In both cases parity gives \(\Psi(c)=1\), \(\eta(c)=-1\), and the sign is unramified at \(\ell\). The prime-to-\(\ell\) part of a local cyclotomic character on \(I_\ell\) factors through \(\F_\ell^\times\), so \(\bar\Psi|_{I_\ell}=\omega^j|_{I_\ell}\) for some \(j\ge0\).
\leanref{UpperTheory.\allowbreak{}V.\allowbreak{}Fr.\allowbreak{}gen\_\allowbreak{}gd}{Development.\allowbreak{}AutomorphicArithmetic.\allowbreak{}AutomorphicVarieties.\allowbreak{}FrobeniusSmoothVectors.\allowbreak{}SqrtDiscr}{152--285}
\leanref{UpperTheory.\allowbreak{}V.\allowbreak{}Fr.\allowbreak{}gen\_\allowbreak{}host}{Development.\allowbreak{}AutomorphicArithmetic.\allowbreak{}AutomorphicVarieties.\allowbreak{}FrobeniusSmoothVectors.\allowbreak{}GenHost}{220--256}

Choose \(a\ge1\) so that \(\eta,\bar\Psi\) factor through \(\Q(\mu_{M^a})\). Write \(M^a\ell=\ell^bD'\) with \(\ell\nmid D'\). The Teichmüller lift \(\chi'=[\bar\Psi\omega^{-j}]\) and the sign lift \(\eta'=\eta\) factor through \(\Q(\mu_{D'})\), since both are trivial on the omitted \(\ell\)-inertia factor. The order of \(\chi'\) divides \(|k_{\rm res}|-1\), and \(\chi'(c)=(-1)^j\). The first test in \eqref{eq:p6c5-link-input-tests} excludes \(\chi'=\eta'\). Unramifiedness of \(\eta\) at \(\ell\) gives \(\ell\nmid D_{K_0}\).

By Dirichlet's theorem \cite{dirichlet-1837}, take a prime \(p>4D_0+M\) with \(p\equiv-1\pmod{4M}\), and a prime \(w>M^{a+1}D'\) with \(w\equiv1\pmod{M^{a+1}D'}\). Put \(H_{\rm q}=G_{\Q(\sqrt{D_0p})}\). At every prime dividing \(M\) the corresponding quadratic character agrees locally with \(\eta\), since \(-p\) is a local square, including at \(2\) by \(2\mid M\). At \(p\) the quadratic character of \(H_{\rm q}\) is nontrivial on inertia, while \(\eta,\bar\Psi,\omega\) are unramified. The prime \(w\) splits all three characters and \(\chi'\), and \(\ell\mid w-1\). Let \(H_w\) be the inverse image of the subgroup of \(\ell\)-th powers of \((\Z/w)^\times\), and take \(\theta\) faithful on its cyclic quotient of order \(\ell\). Then \(\theta\) has conductor \(w\).
\leanref{UpperTheory.\allowbreak{}V.\allowbreak{}Fr.\allowbreak{}exists\_\allowbreak{}twistCharacter\_\allowbreak{}continuousCohomology}{Development.\allowbreak{}AutomorphicArithmetic.\allowbreak{}AutomorphicVarieties.\allowbreak{}FrobeniusSmoothVectors.\allowbreak{}EtaleSimplicialCochainSchemeSectionExtensionNatu}{99--141}
\leanref{UpperTheory.\allowbreak{}V.\allowbreak{}Fr.\allowbreak{}auxw\_\allowbreak{}ex}{Development.\allowbreak{}AutomorphicArithmetic.\allowbreak{}AutomorphicVarieties.\allowbreak{}FrobeniusSmoothVectors.\allowbreak{}AuxwEx}{26--62}
\leanref{UpperTheory.\allowbreak{}V.\allowbreak{}Fr.\allowbreak{}seed\_\allowbreak{}cl}{Development.\allowbreak{}AutomorphicArithmetic.\allowbreak{}AutomorphicVarieties.\allowbreak{}FrobeniusSmoothVectors.\allowbreak{}SeedCl}{23--132}

For \(m\ge0\) put
\[
 H_{{\rm cyc},m}=\ker((\varepsilon_\ell\bmod\ell^{m+1})^{\ell-1}),
 \qquad U_m=H_{\rm q}\cap H_w\cap H_b\cap H_{{\rm cyc},m},
 \qquad F_m=\overline\Q^{\,U_m},\qquad L_m=F_mK_0 .
\]
The three factors after \(H_{\rm q}\) have abelian \(\ell\)-power quotient and contain complex conjugation. The real quadratic subgroup contains it as well. Thus \(F_m\) is totally real abelian, \(L_m\) is CM, \(G_{L_m}=U_m\cap\ker\eta\), and \(\theta|_{G_{L_m}}=1\). The sign is trivial on the local subgroups of \(U_m\) over \(M\) by the choice of \(H_{\rm q}\). For each prolongation over \(\ell\), conjugate the element \(s\) in \eqref{eq:p6c5-link-input-tests} to that prolongation and raise it to an \(\ell\)-power sufficient to enter the three \(\ell\)-power-index subgroups. It already belongs to \(H_{\rm q}\) and retains its square test because \(k_{\rm res}^\times\) has order prime to \(\ell\). At \(B\) take this element as \(g_1\). It is in \(G_{L_m}\) and serves as the seed element in \ref{B6.51frame}. For each \(v\mid\ell\), inertia surjects onto the principal-unit quotient modulo \(\ell^{m+1}\). Since \(U_m\subseteq H_{{\rm cyc},m}\), the local degree \(d_v=[F_{m,v}:\Q_\ell]\) satisfies
\begin{equation}\tag{6.52.3}\label{eq:p6c5-linked-local-degree}
 d_v\ge\ell^m.
\end{equation}
The mod-\(\ell\) cyclotomic character is nontrivial on the local group over \(v\). Inertia is in \(H_{\rm q}\cap H_b\cap H_w\), and \(H_{{\rm cyc},m}\) removes the principal-unit quotient while retaining the full tame cyclotomic quotient.

For the character tests on \(G_{L_m}\), take \(t\in\ker\omega\cap\ker\bar\Psi\) with \(\eta(t)=-1\). The element \(ct\) has \(\eta=1\), \(\bar\Psi=1\), and \(\omega=-1\). If necessary multiply it by an inertia element at \(p\) to enter \(H_{\rm q}\), preserving these values. Raising to a suitable \(\ell\)-power puts it in \(G_{L_m}\) and preserves those values. This tests \(\omega\) against every power of \(\bar\Psi\), and gives the required nontrivial characters involving \(\omega\) in \ref{B6.51frame}.
\leanref{UpperTheory.\allowbreak{}V.\allowbreak{}Fr.\allowbreak{}ufacts}{Development.\allowbreak{}AutomorphicArithmetic.\allowbreak{}AutomorphicVarieties.\allowbreak{}FrobeniusSmoothVectors.\allowbreak{}Ufacts}{78--242}
\leanref{UpperTheory.\allowbreak{}V.\allowbreak{}Fr.\allowbreak{}field\_\allowbreak{}ex}{Development.\allowbreak{}AutomorphicArithmetic.\allowbreak{}AutomorphicVarieties.\allowbreak{}FrobeniusSmoothVectors.\allowbreak{}FieldEx}{24--156}

Fix \(i\in\{1,2\}\) and put \(\chi=\bar\Psi^{-i}\eta\), \(V=k_{\rm res}(\chi)\). The character is odd and \(\chi\ne\omega\), as the element \(t\) of the first test has \(\chi(t)=-1\), \(\omega(t)=1\). Let \(S_\chi\) consist of \(w,\ell\) and the ramified primes of \(\chi\). Use the full local \(H^1\) at \(w\), zero at the other places of \(S_\chi\), and the unramified condition elsewhere. For local subspaces \(\mathcal L_v\) the finite-module Greenberg--Wiles formula is \cite[Prop.~1.6]{wiles-1995}
\begin{equation}\tag{6.52.4}\label{eq:p6c5-link-selmer-ratio}
 \frac{|\mathrm{Sel}_{\mathcal L}(V)|}
 {|\mathrm{Sel}_{\mathcal L^\perp}(V^\vee(1))|}
 =
 \frac{|H^0(\Q,V)|}{|H^0(\Q,V^\vee(1))|\,|H^0(\R,V)|}
 \prod_{v\in S_\chi}\frac{|\mathcal L_v|}{|H^0(\Q_v,V)|},
\end{equation}
where \(V^\vee\) is the \(\F_\ell\)-linear dual and \(\mathcal L^\perp\) is defined by local Tate duality. The set of places includes \(\ell\) and the ramification of \(V\), and the odd-prime hypothesis applies. Oddness kills the real and global invariants, and \(\chi\ne\omega\) kills the global dual invariants. At a ramified place an inertia element acts with \(\chi\ne1\). At \(\ell\) the marking \(g_1\) gives such an element in the decomposition group for both \(i\). At \(w\) the two characters \(\chi,\omega\) are trivial. Local Tate duality and the local Euler characteristic in \ref{B6.09} give
\[
 |H^1(\Q_w,V)|=|H^0(\Q_w,V)|\,|H^2(\Q_w,V)|=|k_{\rm res}|^2.
\]
Thus \eqref{eq:p6c5-link-selmer-ratio} is \(|k_{\rm res}|>1\).

Choose a representative \(y_i:G_\Q\to k_{\rm res}\) of a nonzero Selmer class. Its cocycle law is \(y_i(gh)=y_i(g)+\chi(g)y_i(h)\). It is a coboundary at \(\ell\) for every prolongation. Away from \(w,\ell\) it vanishes on inertia in \(\ker\eta\cap\ker\bar\Psi\), by the coboundary condition at ramified primes and the unramified condition at the others. At \(w\) it is additive on inertia. Wild inertia is pro-\(w\), and tame inertia has maximal elementary abelian \(\ell\)-quotient cyclic of order \(\ell\). The character cutting out \(H_w\) is nonzero on that quotient, so \(y_i=0\) on \(I_w\cap H_w\) at every prolongation.

The restriction to \(U_m\) remains nonprincipal. Otherwise subtract a coboundary to obtain \(z|_{U_m}=0\). At \(u=g_1\) one has \(\chi(u)\ne1\), and normality puts \(gug^{-1}\) in \(U_m\) for all \(g\). Expanding \(z(gu)=z((gug^{-1})g)\) gives \((1-\chi(u))z(g)=0\) for all \(g\), contrary to the nonzero class.
\leanref{UpperTheory.\allowbreak{}V.\allowbreak{}Fr.\allowbreak{}st\_\allowbreak{}loc}{Development.\allowbreak{}AutomorphicArithmetic.\allowbreak{}AutomorphicVarieties.\allowbreak{}FrobeniusSmoothVectors.\allowbreak{}DualNontriv}{14--56}
\leanref{UpperTheory.\allowbreak{}V.\allowbreak{}Fr.\allowbreak{}st\_\allowbreak{}count}{Development.\allowbreak{}AutomorphicArithmetic.\allowbreak{}AutomorphicVarieties.\allowbreak{}FrobeniusSmoothVectors.\allowbreak{}StCount}{10--77}
\leanref{UpperTheory.\allowbreak{}V.\allowbreak{}Fr.\allowbreak{}st\_\allowbreak{}sel}{Development.\allowbreak{}AutomorphicArithmetic.\allowbreak{}AutomorphicVarieties.\allowbreak{}FrobeniusSmoothVectors.\allowbreak{}StSel}{31--129}
\leanref{UpperTheory.\allowbreak{}V.\allowbreak{}Fr.\allowbreak{}st\_\allowbreak{}tame}{Development.\allowbreak{}AutomorphicArithmetic.\allowbreak{}AutomorphicVarieties.\allowbreak{}FrobeniusSmoothVectors.\allowbreak{}StTame}{47--80}
\leanref{UpperTheory.\allowbreak{}V.\allowbreak{}Fr.\allowbreak{}strict}{Development.\allowbreak{}AutomorphicArithmetic.\allowbreak{}AutomorphicVarieties.\allowbreak{}FrobeniusSmoothVectors.\allowbreak{}Strict}{81--175}

The frame group is the quotient of \(U_m=G_{F_m}\) by the closed normal subgroup generated by \(I_v\cap\ker\eta\cap\ker\bar\Psi\) at the finite places away from \(\ell\), with inertia taken over \(F_m\). Its ordinary local maps come from the decomposition groups above \(\ell\), with one prolongation for each place. Where \(\chi=1\) the zero set of \(y_i\) is a subgroup. If \(y_i(h)=0\) and \(\chi(h)=1\), then the cocycle law gives \(y_i(ghg^{-1})=0\). Local constancy gives closedness, so \(y_i\) descends. Choose the scalar giving its local coboundary at each ordinary map. The descended pair is a strict class of \ref{B6.51frame} and is nonzero by its pullback to \(U_m\). The two coefficient characters use the same \(w,F_m,L_m,g_1\).
\leanref{UpperTheory.\allowbreak{}V.\allowbreak{}Fr.\allowbreak{}link\_\allowbreak{}ex}{Development.\allowbreak{}AutomorphicArithmetic.\allowbreak{}AutomorphicVarieties.\allowbreak{}FrobeniusSmoothVectors.\allowbreak{}LinkEx}{150--309}
\leanref{UpperTheory.\allowbreak{}V.\allowbreak{}Fr.\allowbreak{}frobenius\_\allowbreak{}galoisField\_\allowbreak{}four\_\allowbreak{}compatibility}{Development.\allowbreak{}AutomorphicArithmetic.\allowbreak{}AutomorphicVarieties.\allowbreak{}FrobeniusSmoothVectors.\allowbreak{}FrobeniusGaloisFieldFourCompatibility}{105--161}

Let \(E\) be the fixed field of \(H_{\rm q}\cap H_w\cap H_b\cap\ker\eta\cap\ker\bar\Psi\cap\ker\omega\), and put \(E_m=E(\mu_{\ell^{m+1}})\). The common residual kernel at layer \(m\) fixes \(E_m\), since raising a principal unit to \(\ell-1\) preserves its congruence depth. The field \(E\) is finite abelian and contains \(\mu_\ell\). By \ref{B6.17} the \(\ell\)-ranks of the class groups of all \(E_m\) are bounded by one \(C_0\).
\leanref{UpperTheory.\allowbreak{}V.\allowbreak{}Fr.\allowbreak{}b0}{Development.\allowbreak{}AutomorphicArithmetic.\allowbreak{}AutomorphicVarieties.\allowbreak{}FrobeniusSmoothVectors.\allowbreak{}RankLeOfCardLe}{137--188}

For each of the first two groups in \eqref{eq:p6c5-frame-cohomology-bounds}, restriction to \(G_{E_m}\) is an unramified character. The strict condition gives zero on the local groups above \(\ell\), and the frame quotient kills inertia on the common kernel elsewhere. Restriction is injective on the global classes by averaging over the prime-to-\(\ell\) residual quotient, and the local square test determines the strict scalars. Class field theory gives the first two bounds in \eqref{eq:p6c5-frame-cohomology-bounds}. For the trivial character use the same Selmer formula over \(F_m\), with full conditions above \(\ell\) and unramified conditions elsewhere. Every frame character into the additive residue field satisfies these conditions. At \(v\mid\ell\), the local \(\omega\) test gives zero dual invariants, so the local factor by Euler characteristic is \(|k_{\rm res}|^{d_v}\). The real factors contribute \(|k_{\rm res}|^{-d}\), and the global factor is \(|k_{\rm res}|\). The dual strict classes for \(\omega\) inject into unramified characters of \(G_{E_m}\) by the same restriction argument. The resulting bound is the third of \eqref{eq:p6c5-frame-cohomology-bounds}.
\leanref{UpperTheory.\allowbreak{}V.\allowbreak{}GC.\allowbreak{}h1}{Development.\allowbreak{}AutomorphicArithmetic.\allowbreak{}AutomorphicVarieties.\allowbreak{}GaloisCochains.\allowbreak{}CastToZModPowOne}{18--48}

The linked frame uses the cyclotomic kernel in each ordinary inertia for \(N_v\), on which \(\bar\Psi=1\). By local reciprocity take generators of the pro-\(\ell\) inertia image in the local abelianisation with all but one in the cyclotomic kernel. The local residual cyclotomic test excludes \(\ell\)-power roots of unity in the local field and the pro-\(\ell\) unit group has rank \(d_v\). The generators can be taken in \(\ker\bar\Psi\) by removing their prime-to-\(\ell\) images. Class field theory gives the common exponent on the sign kernel from the finite unramified class group over \(L_m\), since the remaining allowed tame inertia has order prime to \(\ell\). Finiteness for open-subgroup character spaces follows from finite-coefficient global cohomology with ramification at finitely many places. For \eqref{eq:p6c5-frame-euler}, apply Lemma~\ref{B6.51framea} with \(F=F_m\) and \(K=G_{F_m}\cap\ker\eta\cap\ker\bar\Psi\). This is the kernel of the pair of finite characters. Its index divides \(2|k_{\rm res}^{\times}|\), prime to the odd prime \(\ell\). The field \(F_m\) is totally real. The inertia relations defining the frame are exactly those of the lemma, and the ordinary local groups identify with the full decomposition groups at the places above \(\ell\). The frame cyclotomic character and the real involutions are the images of the corresponding Galois data. The lemma therefore gives the required bound for every coefficient field, representation and family of subspaces in Definition~\ref{B6.51frame}.
\leanref{UpperTheory.\allowbreak{}V.\allowbreak{}GC.\allowbreak{}ec}{Development.\allowbreak{}AutomorphicArithmetic.\allowbreak{}AutomorphicVarieties.\allowbreak{}GaloisCochains.\allowbreak{}ResH1SurjectiveOfIsomorphism}{84--180}

Take \(m=3C_0+7\). By \eqref{eq:p6c5-linked-local-degree} and the existence of a place over \(\ell\),
\begin{equation}\tag{6.52.1}\label{eq:p6c5-linked-layer}
 d=\sum_vd_v\ge\ell^m>3C_0+7,\qquad d_v>5 .
\end{equation}
\leanref{UpperTheory.\allowbreak{}V.\allowbreak{}Fr.\allowbreak{}frame\_\allowbreak{}ex}{Development.\allowbreak{}AutomorphicArithmetic.\allowbreak{}AutomorphicVarieties.\allowbreak{}FrobeniusSmoothVectors.\allowbreak{}FrameEx}{126--161}
\leanref{UpperTheory.\allowbreak{}V.\allowbreak{}Fr.\allowbreak{}ok\_\allowbreak{}basic}{Development.\allowbreak{}AutomorphicArithmetic.\allowbreak{}AutomorphicVarieties.\allowbreak{}FrobeniusSmoothVectors.\allowbreak{}OkBasic}{46--114}

The link pulls an ordinary point into the framed problem of \ref{B6.51a}. Away from \(\ell\), on inertia in \(U_m\) the exact condition \(\Psi=1\) agrees with \(\bar\Psi=1\) because \(U_m\subseteq H_b\) removes the \(\ell\)-primary part there. At \(\ell\) the ordinary flag transports to each prolongation, and in residue characteristic the prescribed characters kill \(\ker\bar\Psi\) in inertia. For (1) take \(U=G_{L_m}\) before choosing \(A\) or the point. It has index \([L_m:\Q]\) and lies in \(\ker\eta\). Restrict the point to the group used for the link. Its kernel contains the closed normal subgroup imposed away from \(\ell\): each generating inertia element has \(\eta=\bar\Psi=1\), the link gives \(\Psi=1\) there, and the ordinary point kills it. The kernel is closed by continuity. Thus the representation descends to the frame quotient, and on its sign kernel the descended matrix pulled back to \(G_{L_m}\) is the original matrix. The link sends the chosen conjugation and marking to those of the frame, so their matrix equations and the prescribed diagonal reduction are preserved. Transport the flag at \(B\) to each prolongation above \(\ell\). Its inertia characters are the images of \(\Psi^{1-i}\), for \(i=0,1,2\); conjugate duality reverses the order and inverts the characters, preserving this list. Since \(O\to A\) factors through \(j:O/\mathfrak m_O\to A\), they equal \(j(\bar\Psi^{1-i})\) on inertia. This checks the residue-characteristic hypothesis of \ref{B6.51a}. Given \(g\in U\), use its image in the sign kernel of the frame in that lemma. Pullback gives \(\det(X-r(g))=\prod_{i=0}^2(X-j(\bar\Psi(g)^{1-i}))\) in \(A[X]\).\par
\leanref{UpperTheory.\allowbreak{}V.\allowbreak{}Pr.\allowbreak{}exists\_\allowbreak{}nrep3}{Development.\allowbreak{}AutomorphicArithmetic.\allowbreak{}AutomorphicVarieties.\allowbreak{}CurveDVRSpecialization.\allowbreak{}ExistsNrep3}{271--444}
\leanref{UpperTheory.\allowbreak{}V.\allowbreak{}Pr.\allowbreak{}finEnd\_\allowbreak{}of\_\allowbreak{}finPt}{Development.\allowbreak{}AutomorphicArithmetic.\allowbreak{}AutomorphicVarieties.\allowbreak{}CurveDVRSpecialization.\allowbreak{}FinEndOfFinPt}{9--48}
 For (2) the restriction to \(G_{L_m}\) is irreducible by hypothesis. Apply the rank-three conclusion, then solvable descent with the pairing and local flags to obtain the attachment with witness over \(K\).

For (3) the Type II corner gives \(W|_{G_K}\) irreducible. If it became reducible over \(L_m\), the abelian quotient \(G_K/G_{L_m}\) would give a nontrivial self-twist \(\lambda\) of the restriction to \(G_K\). Indeed, after semisimplifying that restriction to the normal subgroup \(G_{L_m}\), reducibility makes its endomorphism space larger than the scalars. The abelian quotient acts on this space by conjugation, and a nontrivial eigenspace gives the self-twist by irreducibility. For \(g\) with \(\lambda(g)\ne1\), both \(g\) and \(gg_1\) have trace zero because \(g_1\in G_{L_m}\). The residual corner trace at either element \(h\) is \(\bar\Psi(h)+\bar\Psi(h)^{-1}\), so \(\bar\Psi(g)^2=\bar\Psi(gg_1)^2=-1\). This contradicts \(\bar\Psi(g_1)^2\ne1\). Apply the rank-two conclusion of \ref{B6.51a}. Use the algebraic twist and witness on \(L_m\) furnished by \ref{B6.51a}. Continuity on the open subgroup \(G_K\) follows from the corner matrices, hence \(W\) is continuous. For \(g\in G_K\), the polarized matrix identity identifies the corner determinant at \(cg^{-1}c\) with that at \(g\), whereas the determinant character of \(W\) takes inverse values there. Hence \(\det W(g)^2=1\). Its diagonal residual corner has determinant one, and \(\ell\) is odd, so \(\det W(g)=1\). The involution \(W(c)\) has trace zero and determinant \(-1\). Thus \(\det W=\eta\), a finite cyclotomic character. At \(p\ne\ell\), inertia in the common kernel of \(\eta,\Psi\) acts trivially on the corner. The quotient through which inertia acts is finite, since the cyclotomic character is unramified there, and decomposition acts trivially on this quotient since both characters are global. At \(B\mid\ell\), intersect inertia with the open normal subgroup \(\ker\eta\cap\ker\vartheta\). The first line of the ordinary flag has character \(\varepsilon_\ell^k\). Its projection to the middle coordinate vanishes, since the middle character is one and a cyclotomic inertia element of infinite order distinguishes the two. It is therefore a line of the corner, whose quotient character is \(\varepsilon_\ell^{-k}\) by the determinant. All hypotheses of \ref{B6.52a} now hold for this same \(W\). The same link construction with Gaussian sign gives (4).
\leanref{UpperTheory.\allowbreak{}V.\allowbreak{}Fr.\allowbreak{}host\_\allowbreak{}ex}{Development.\allowbreak{}AutomorphicArithmetic.\allowbreak{}AutomorphicVarieties.\allowbreak{}FrobeniusSmoothVectors.\allowbreak{}HostEx}{8--26}
\leanref{UpperTheory.\allowbreak{}V.\allowbreak{}Z.\allowbreak{}Win.\allowbreak{}Dn.\allowbreak{}autEnd3\_\allowbreak{}alt}{Development.\allowbreak{}AutomorphicArithmetic.\allowbreak{}AutomorphicVarieties.\allowbreak{}ModularityArgument.\allowbreak{}ThresholdArguments.\allowbreak{}DenseThresholdCases.\allowbreak{}AutEnd3Alt}{13--106}
\leanref{UpperTheory.\allowbreak{}V.\allowbreak{}Z.\allowbreak{}Win.\allowbreak{}Dn.\allowbreak{}hostEnd\_\allowbreak{}of\_\allowbreak{}loc3\_\allowbreak{}frame\_\allowbreak{}linkedEnd\_\allowbreak{}downCsdK}{Development.\allowbreak{}AutomorphicArithmetic.\allowbreak{}AutomorphicVarieties.\allowbreak{}ModularityArgument.\allowbreak{}ThresholdArguments.\allowbreak{}DenseThresholdCases.\allowbreak{}DnHost}{17--67}
\leanref{UpperTheory.\allowbreak{}V.\allowbreak{}Z.\allowbreak{}Win.\allowbreak{}Dn.\allowbreak{}AutEnd2Ao}{Development.\allowbreak{}AutomorphicArithmetic.\allowbreak{}AutomorphicVarieties.\allowbreak{}ModularityArgument.\allowbreak{}AutomorphyLifting.\allowbreak{}AutomorphyDescent.\allowbreak{}AutEnd2Ao}{18--26}
\end{proof}

\subsection*{Proof of \ref{B6.53}}\label{prf:B6.53}\phantomsection
\begin{proof}
Let \(p:G_{L^+}\twoheadrightarrow\Gamma\) be the level quotient. It respects the quadratic character of \(L/L^+\) and the full \(\ell\)-adic cyclotomic character. In particular
\[
 p(G_L)=\Gamma_0,\qquad
 p(G_{L(\mu_{\ell^\infty})})=\Gamma'.
\]
For the second equality, lift any element of \(\Gamma'\) to \(G_L\). Its cyclotomic value equals that of its image and is therefore one. Hypothesis~3 gives absolute irreducibility of \(\rho\circ p\) on \(G_{L(\mu_{\ell^\infty})}\), since its coefficient field is \(\overline{\Q}_\ell\).

By hypothesis~1 and Definition~\ref{def:b6-cuspidal-witness}, choose a regular algebraic conjugate self-dual cuspidal attachment \(\pi\) over \(L\). At the unramified places in the definition, \(\rho\circ p\) and \(r_{\pi,\iota}\) have the same characteristic polynomial: at geometric Frobenius the roots in the complex normalisation are the Satake roots multiplied by \(q_v^{(n-1)/2}\). Chebotarev density and Brauer--Nesbitt identify the two semisimple representations. Thus the irreducibility required in \cite[Theorem~6.1]{thorne-2022} holds for \(r_{\pi,\iota}\).

We identify the action on the coefficient module of that theorem. On \(G_L\) it is conjugation on the entire matrix algebra. A polarising form gives the action at \(c\) by \(X\mapsto-P^{-1}X^tP\), by differentiating inverse transpose. Any two intertwiners between \(\rho^c\) and \(\rho^\vee\varepsilon_\ell^{1-n}\) differ by a scalar by irreducibility and Schur's lemma. The scalar cancels in this formula. Hence the extended adjoint used by Thorne is \(R\circ p\).

The pullback \(g\mapsto z(p(g))\) gives a class in \(H^1_f(L^+,R\circ p)\) by the finite local conditions in the statement. Theorem~6.1 of \cite{thorne-2022} gives zero for this group: \(L\) is CM, \(n\ge2\), the attachment is polarizable and the cyclotomic restriction is irreducible. The published theorem allows all such ranks and does not require the ordinary flag conditions. There is consequently an \(X\in M_n(\overline{\Q}_\ell)\) with
\[
 z(p(g))=R(p(g))X-X\qquad(g\in G_{L^+}).
\]
Surjectivity of \(p\) gives \eqref{eq:p6c5-adjoint-coboundary}.
\end{proof}

\subsection*{Proof of \ref{B6.53a}}\label{prf:B6.53a}\phantomsection
\begin{proof}
Put \(\partial(\rho(g))=z(g)\rho(g)\). If \(\sum_g a_g\rho(g)=0\) is a finite relation, apply the zero trace condition at \(gh\) and use the cocycle law to obtain
\[
 \operatorname{tr}\!\left(\sum_g a_g\partial(\rho(g))\rho(h)\right)
 =-\operatorname{tr}\!\left(\sum_g a_g\rho(g)\partial(\rho(h))\right)=0 .
\]
Full span and the trace pairing make \(\partial\) well defined. It satisfies the Leibniz rule by the cocycle identity. For the matrix units \(E_{ij}\), put \(A=\sum_i\partial(E_{i1})E_{1i}\). Applying the rule to \(E_{jk}E_{i1}\) gives \(AE_{jk}-E_{jk}A=\partial(E_{jk})\). Hence \(z=d_R(-A)\).
\end{proof}

\subsection*{Proof of \ref{B6.54a}}\label{prf:B6.54a}\phantomsection
\begin{proof}
Let \(\pi\) be the cuspidal attachment. Outside the level the determinant factors through the inertia quotient. For an inertia element \(u\) there, trace translation gives \(\operatorname{tr}(\rho(g)(\rho(u)-1))=0\) for all \(g\). Full span gives \(\rho(u)=1\), and local compatibility makes \(\pi\) spherical there. Apply weak descent \ref{B9.36a} to obtain the definite-unitary representation. The base-change comparison of \ref{B6.04a} and uniqueness at split places identify its split local components with those of \(\pi\), including at the places in the level. The infinitesimal character and the Hodge--Tate weights give the dual algebraic weight of the ordinary module. Triviality on \(I_0\) supplies the spherical vector at the distinguished place. At \(v\mid\ell\), local compatibility for the unlinked flag gives the ordinary vectors with diamond eigencharacters in the prescribed order. At the remaining level places the local comparison gives the specified tame types.

These vectors give a nonzero eigenspace in the arithmetic weight fibre. Its determinant agrees with \(D_x\) at the unramified Frobenii and hence everywhere by density and continuity. Its ordered characters agree by the ordinary comparison. Thus the action factors through the point \(x\).
\end{proof}

\subsection*{Proof of \ref{B6.54}}\label{prf:B6.54}\phantomsection
\begin{proof}
Fix an ordinary cocycle \(z\). Write \(K=\overline{\Q}_\ell\). In the dual numbers \(K[\eta]/(\eta^2)\) put
\[
 \rho_z(g)=(1+\eta z(g))\rho(g).
\]
The cocycle equation is exactly the multiplication law for \(\rho_z\). The equation involving \(c\) gives the polarisation of its determinant. At an ordinary place, conjugation by a matrix \(1+\eta Y_v\) subtracts a local coboundary, with the sign determined by the choice of \(Y_v\). Choose \(Y_v\) so that the corrected cocycle is upper triangular with zero diagonal on inertia. If its diagonal entries are \(d_{vi}(g)\), the characters of the lifted flag are \((1+\eta d_{vi})\kappa_{v,i}\). Their restrictions to inertia have zero derivative. On the designated tame kernels both the lift and its specialisation have the same action.

Here are the integral choices needed to apply \ref{B6.55}. The representation and the cocycle are defined over a common finite extension of \(\Q_\ell\). Indeed their joint image in the semidirect product of the matrices with \(\GL_n\) is a compact group. It is a countable union of its closed subgroups defined over finite coefficient extensions. One of these has interior by Baire's theorem, hence finite index. Adjoin the coefficients of coset representatives, the finitely many flag bases and local correcting matrices, and the roots at the determinant tests. A stable lattice exists since the image of \(\rho\) is compact.

After finite coefficient extension one can take the pairing to be perfect on a stable lattice. To see this, let \(D\) be a stable lattice and \(D^\vee\) its dual for \(P\). The multiplier is a unit and \(D\) is stable under the conjugate representation as well, so \(D^\vee\) is stable. Diagonalise the form on \(D\) in orthogonal lines in the symmetric case and orthogonal planes in the alternating case, after clearing one common denominator. This is the elementary divisor algorithm for a symmetric or alternating form over a DVR with \(2\) invertible: a pairing of least valuation splits off a line if it occurs on the diagonal, and otherwise splits a plane; in the symmetric case a sum of its two basis vectors gives a diagonal entry of least valuation. Write \(\varpi^{d_i}\) for the elementary divisors, repeated within each alternating plane. Stability of both lattices gives
\[
 v(\rho(g)_{ij})\ge\max(0,d_j-d_i).
\]
Adjoin \(\varpi^{1/2}\) and replace the basis vectors by \(\varpi^{-d_i/2}\) times those vectors. The new entries have valuations at least \(\max(0,d_j-d_i)+(d_i-d_j)/2\ge0\), and the form is perfect. Intersect each ordinary flag with this lattice. Consecutive quotients are torsion-free, so an integral basis adapts to each flag. Full rational span gives a nonzero multiple of the entire integral matrix algebra in the integral span, by clearing the denominators of expressions for the matrix units.

Denote the resulting coefficient DVR by \(O'\). Multiply \(z\) and the local corrections by a common nonzero scalar \(a\) to make them integral. The residual determinant and ordered residual characters are the original ones extended to the residue field of \(O'\), since conjugation preserves the determinant and the specialised flag characters. The integral determinant and flag characters over \(O'[\eta]/(\eta^2)\) thus give, by universality, an arithmetic specialisation \(x'\) and an \(x'\)-derivation \(\delta\) on the determinant ring for this residual datum. For every \(g\),
\[
 \delta(\tau^{\rm univ}(g))=a\,\operatorname{tr}(z(g)\rho(g)).
\]
The inertia derivatives vanish. At \(D_0\) choose an element with nontrivial residual cyclotomic value. Its residual eigenvalues are all the same nonzero scalar. The adjoint action there has all eigenvalues one modulo the maximal ideal. Its cyclotomic multiple minus the identity is invertible, which proves \eqref{eq:p6c5-distinguished}. The nonzero arithmetic weight differences separate the local characters on an inertia element, and \eqref{eq:p6c5-unlinked} gives the required unlinking. Determinant regularity persists over \(O'\): a rational solution of the coboundary equations after scalar extension descends by linear algebra.

Use the modules of \ref{B6.48c} after this coefficient extension. The determinant and flag of \(x'\) have the original cuspidal attachment. Triviality on \(I_0\) and \eqref{eq:p6c5-unlinked} give \ref{B6.54a}, so \(x'\) kills the base Hecke ideal. In an arithmetic weight fibre the determinant coefficients act through the classical eigensystems of \ref{B6.48c}. Thus a power acting by zero in characteristic zero has zero eigenvalue in each eigensystem, and the operator itself acts by zero. Integral weight fibres are free, so the implication also holds integrally. These are the level and arithmetic control hypotheses of \ref{B6.55}. It follows that \(\delta(\tau^{\rm univ}(g))=0\) for every \(g\). Cancel \(a\) and apply \ref{B6.53a} with the full matrix span in \eqref{eq:p6c5-unlinked}.
\end{proof}

\subsection*{Proof of \ref{B6.55}}\label{prf:B6.55}\phantomsection
\begin{proof}
Use \ref{B6.56} at the specified \(s_j\). Every element of its subgroup \(H\) has residual characteristic polynomial \((X-1)^n\). The finite residual image of \(H\) is an \(\ell\)-group: each of its elements is unipotent and has \(\ell\)-power order. The congruence kernel on an integral lattice is pro-\(\ell\), so \(\rho_x|_H\) factors through the maximal pro-\(\ell\) quotient. This gives hypothesis~3 of \ref{B6.57}. The separated flags, distinguished-place condition and determinant regularity in the present theorem are the hypotheses of \ref{B6.55a}. Its last inequality gives hypothesis~4 of \ref{B6.57}.

It remains to construct the operator in hypothesis~2 of \ref{B6.57}. Write \(I=\ker(x|_\Lambda)\), \(K=\operatorname{Frac}O\), \(V=(M/IM)\otimes_O K\). Let \(A\) be the image of \(B[1/\ell]\) in \(\operatorname{End}_K(V)\). The character \(x\) occurs in the support of \(V\): localise \(M\) at \(\ker x\), use \(J\subseteq\ker x\) and \(\operatorname{Supp}M=V(J)\), and apply Nakayama to \(I\). Thus \(x\) defines a factor of the finite algebra \(A\) with residue field \(K\).

The subalgebra generated by the determinant coefficients in \(A\) is reduced. For a nilpotent element in it, clear denominators to obtain an integral polynomial \(t\) in these coefficients. Then \(t^k\) kills \(V\) for some \(k\), and kills \(M/IM\) because that module is free over \(O\). Arithmetic control at \(x\) gives \(tM\subseteq IM\). A finite reduced commutative algebra over \(K\) is a product of separable fields, so every \(K\)-valued derivation at \(x\) kills this subalgebra. By the second assertion of \ref{B6.57a}, the pullback of such a derivation to \(B\) also kills the ordered characters, since its target is a characteristic-zero field. The coefficients and characters generate the image algebra. Indeed they generate \(B\) adically, the action on the weight fibre is continuous, and their \(K\)-algebra in the finite-dimensional endomorphism space is a closed vector subspace. Hence the derivation is zero.

Let \(\mathfrak n\) be the maximal ideal of the local factor at \(x\). Its cotangent dual is the space of these derivations. Nakayama applied to \(\mathfrak n/\mathfrak n^2=0\) gives \(\mathfrak n=0\). The idempotent of the factor therefore kills the image of \(\ker x\) on \(V\) and takes value one at \(x\). Lift a nonzero \(O\)-multiple of this idempotent to \(t_0\in B\). Freeness of \(M/IM\) gives \(t_0(\ker x)M\subseteq IM\), and \(x(t_0)\ne0\).

The base and auxiliary modules give hypothesis~1 of \ref{B6.57} by the level assumptions. Each weight parameter is a fixed unit times a prescribed inertia-character value minus one, so its derivative vanishes. Lemma~\ref{B6.57} now gives \eqref{eq:p6c5-inertia-to-trace}.
\end{proof}

\subsection*{Proof of \ref{B6.56}}\label{prf:B6.56}\phantomsection
\begin{proof}
Residual conjugate duality permutes the kernels of the residual characters under \(c\). Thus the displayed subgroup \(H\) is open and normal in \(\Gamma\). Write \(P=H/N_\ell(H)\). The given finiteness of \(\operatorname{Hom}_{\rm cont}(H,\F_\ell)\) makes \(P\) a finitely generated pro-\(\ell\) group. For \(r\ge1\) intersect the kernels of its continuous maps to groups of order at most \(\ell^r\). There are finitely many such kernels: there are finitely many groups of bounded order, and images of a finite generating set determine a map. Their intersection is open and characteristic. Its inverse image \(W_r\) in \(H\) is open normal in \(\Gamma\), and the \(W_r\) are cofinal among the neighbourhoods of \(N_\ell(H)\) in \(H\).

Let \(p:G_{L^+}\to\Gamma\) be the level map. Lift \(s_j\) to \(g_j\in G_L\). The cyclotomic value of \(g_j\) is one because the cyclotomic character factors through \(p\). Use the finite Galois quotient cut out by
\[
 U_r=p^{-1}(W_r)\cap\ker(\varepsilon_\ell\bmod\ell^r).
\]
By Chebotarev \cite[\S1A]{thorner-2019}, choose, successively for the finitely many \(j\), a prime \(v_{r,j}\) outside \(S\), distinct from the previous choices, with Frobenius conjugate to \(g_j^{-1}\) in this quotient. Choose its prolongation to realise that element of the conjugacy class. Its arithmetic Frobenius \(\sigma_{r,j}\) satisfies
\begin{equation}\tag{6.56.3}\label{eq:p6c5-aux-chebotarev}
 \sigma_{r,j}g_j\in U_r .
\end{equation}
Since \(U_r\subset G_L\), the prime splits in \(L\). For the place \(w_{r,j}\) given by the prolongation, the cyclotomic value of arithmetic Frobenius is \(\mathrm Nw_{r,j}\), giving \eqref{eq:p6c5-aux-norm}. Taking inverses in the quotient by \(W_r\) shows that the difference for inverse arithmetic Frobenius lies in \(W_r\). This proves both assertions about \(d_{r,j}\).

Allow ramification at the new places to obtain \(h_r\), and put \(H_r=h_r^{-1}H\). Its kernel over the old level is normally generated as a closed subgroup by the new inertia groups, by the definition of the unramified quotient. Each new inertia group lies in \(H_r\). In the maximal pro-\(\ell\) quotient of \(H_r\) its wild inertia and its prime-to-\(\ell\) tame inertia vanish. Choose \(\tau_{r,j}\) lifting a generator of the pro-\(\ell\) tame inertia. Arithmetic Frobenius conjugates this generator to its \(\mathrm Nw_{r,j}\)-th power. With \(\phi_{r,j}\) the inverse arithmetic Frobenius this is \eqref{eq:p6c5-tame-relation}. We obtain \(h_r(\tau_{r,j})=1\) and the asserted normal generation modulo \(N_\ell(H_r)\).

Finally a continuous character \(H_r\to\F_\ell\) kills \(N_\ell(H_r)\) and is invariant under conjugation within \(H_r\). Its restriction to \(\ker h_r\) is determined by its values on the conjugates of the \(q_0\) elements \(\tau_{r,j}\) by representatives of \(\Gamma_r/H_r\). The kernel of this restriction consists of characters inflated from \(H\). Since \([\Gamma_r:H_r]=[\Gamma:H]\), the number of characters is at most
\(\#\operatorname{Hom}_{\rm cont}(H,\F_\ell)\ell^{q_0[\Gamma:H]}\), uniformly in \(r\).
\end{proof}

\subsection*{Proof of \ref{B6.55a}}\label{prf:B6.55a}\phantomsection
\begin{proof}
Write \(M_m=M_n(A_m)\). Extend an allowed cocycle to the full group by assigning zero at \(c\). The polarisation identity in \eqref{eq:p6c5-ordinary-cocycle} is the compatibility needed for this extension. A coboundary has trace derivative zero, since \(\operatorname{tr}(\rho(g)X-X\rho(g))=0\). It suffices to bound the finite Selmer group for \(M_m\). Under the trace pairing its Cartier dual has the cyclotomic-twisted adjoint action. For the additive dual one can compose the matrix trace with a perfect additive pairing on \(O/\mathfrak m_O^m\), obtained from the inverse different of \(O/\Z_\ell\).

Use finite Poitou--Tate duality \ref{B6.09} with the allowed local subgroups \(\mathcal L_v\) and their orthogonal complements. Taking orders in its exact sequence and using the local and global Euler characteristic formulas \cite[Theorems~7.3.1 and~8.7.4]{neukirch-2008} gives the ratio of the primal to the dual Selmer order. The factors are the global ratio of invariant orders, the factors \(|\mathcal L_v|/|H^0(L_v^+,M_m)|\) at finite places, and the reciprocals of the invariant orders at the real places. The modules are finite of odd prime-power order, the ramification set contains the primes above \(\ell\), and outside it we use the unramified condition.

We use the following elementary calculation for the invariant factors. For equations on a free \(O\)-module of rank \(a\), take finitely many rows generating the row module. If their common kernel has dimension \(h\) over \(\operatorname{Frac}O\), and the nonzero Smith entries have valuations \(b_i\), the reduced kernel has order
\(Q^{mh+\sum_i\min(m,b_i)}\). For global invariants \(h=0\): full matrix span leaves the scalars on the CM subgroup, and \(c\) acts by minus one on scalars. This bounds the numerator of the global ratio.

At \(v\mid\ell\) let \(T_v\) be the upper triangular matrices in the flag lattice. Its rank is \(n+b\). In its filtration by matrix entries the characters are the trivial characters and \(\kappa_{v,i}/\kappa_{v,j}\) with \(i<j\). None equals the cyclotomic character, by unlinking and the infinite cyclotomic image on the local group. Thus the rational twisted dual has no invariants, and its reduced invariants have bounded order by the calculation above. Local duality and the Euler characteristic formula bound the order of the cocycles of \(T_v/\mathfrak m_O^mT_v\) by \(Q^{m(n+b)(1+d_v)+C_v}\), with \(d_v=[L_v^+:\Q_\ell]\). Dividing by triangular coboundaries and by the full invariant order costs a factor at least \(Q^{m(n+b)}\), since triangular invariants inject into full invariants. Requiring zero diagonal on \(\mu_v\) can only decrease the numerator.

At a real place the fixed subspace of negative transpose for the polarising form has rank at least \(b\). It is a direct summand integrally because \(2\) is invertible. At a tame type place the prescribed inertia quotient has order prime to \(\ell\). The unramified quotient on its inertia invariants has first cohomology the cokernel of Frobenius minus one, of the same order as its kernel. At the distinguished place the full-condition ratio is the order of twisted dual invariants, bounded by \eqref{eq:p6c5-distinguished} and the Smith calculation. At an auxiliary place, for fixed \(m\) and large \(r\), the twisted Frobenius on \(M_m\) is conjugation by \(\rho_x(s_j)\) or its inverse. Indeed \(\rho_x|_H\) has pro-\(\ell\) image, and \ref{B6.56} gives agreement modulo \(\mathfrak m_O^m\) and cyclotomic value one modulo that ideal. The centralizer of a matrix with \(n\) distinct roots has rational dimension \(n\). Consequently the bounds on the local factors are
\begin{equation}\tag{6.55.4}\label{eq:p6c5-selmer-local}
 \frac{|\mathcal L_v|}{|H^0(L_v^+,M_m)|}\le
 \begin{cases}
 Q^{m(n+b)d_v+C_v}&v\mid\ell,\\
 Q^{-mb}&v\text{ real},\\
 1&v\text{ of tame type},\\
 Q^{C_0}&v\text{ distinguished},\\
 Q^{mn+C_j}&v=v_{r,j}.
 \end{cases}
\end{equation}
At a real place the numerator denotes the order one.

It remains to bound the dual Selmer order uniformly. Subtract a coboundary to make a dual cocycle zero at \(c\). At each new place the orthogonal condition to the full condition is zero, so the cocycle is a local coboundary. It vanishes on new inertia, whose action is trivial. On \(H_r=h_r^{-1}H\) the action and the additive cocycle together map into a finite \(\ell\)-group. They kill \(N_\ell(H_r)\). The kernel description of \ref{B6.56} now gives descent of the cocycle to \(\Gamma\). At a fixed \(m\) there are finitely many descended cocycles: their restrictions to \(H\) are determined on finitely many pro-\(\ell\) generators, and finitely many cosets remain. For all sufficiently large \(r\), \ref{B6.56} therefore identifies the cocycle and action values at \(h_r(\phi_{r,j})\) and \(s_j\) simultaneously for all of them. Being a coboundary at the former makes every trace test at the latter vanish, by cyclic invariance of the trace and \(\varepsilon_\ell(s_j)=1\).

For clarity, the bound on reductions satisfying these tests is uniform in \(m\). Record values on pro-\(\ell\) generators and coset representatives in a finite free \(O\)-module. The cocycle law on relations, the value zero at \(c\) and the trace tests are linear equations in these coordinates. To obtain equations valid at every depth, start with a free pro-\(\ell\) presentation of \(H/N_\ell(H)\). Values on its generators extend to the free group by the cocycle law at every depth: the action together with the values has finite \(\ell\)-group image. Evaluation on each relator is an \(O\)-linear row, obtained by inverse limit of the evaluations at finite depth. Impose also the coset multiplication and conjugation relations. The row span is finitely generated over \(O\). Over the fraction field its kernel is precisely the space of normalised coboundaries by determinant regularity. A rational cobounding matrix can be taken fixed by the twisted action of \(c\), since the cocycle is zero there. Apply Smith reduction first to the equations and then to the coboundary map on this invariant lattice. The first calculation bounds the extra reduced kernel modulo reductions of the integral kernel; the second bounds the index of the integral coboundaries in that kernel. Both bounds depend only on the two integral matrices. Hence the number of dual classes is bounded independently of \(m\).

Multiply \eqref{eq:p6c5-selmer-local} and the global bounds. The \(\ell\)-adic degrees sum to \(d\), there are \(d\) real places, and the \(q_0\) auxiliary factors contribute \(mnq_0+O(1)\) to the exponent. This gives \(m((n+b)d-bd+nq_0)+O(1)=m(D+nq_0)+O(1)\), equivalently \eqref{eq:p6c5-selmer-count}.
\end{proof}

\subsection*{Proof of \ref{B6.57a}}\label{prf:B6.57a}\phantomsection
\begin{proof}
Let \(\tau\) be the determinant trace. The Cayley--Hamilton algebra of the split multiplicity-free residual determinant has orthogonal idempotents \(e_i\) with scalar diagonal corners and trace one \cite[Theorem~2.22]{chenevier-2014}. The specialised representation factors through it by full span and the trace pairing, and the idempotents split the specialised lattice into rank-one summands. Work in that basis. Take a fixed \(c_0\ne0\) with \(c_0M_n(O)\) in the specialised group-algebra image. By surjectivity of the level map the same \(c_0\) works at every level. Lift \(c_0E_{i1}\) and \(c_0E_{1i}\) to \(U_i\) and \(V_i\) in the corresponding corners. Write \(U_iV_i=\beta_ie_i\), multiply \(U_i\) by \(\prod_{j\ne i}\beta_j\), and put \(\beta=\prod_j\beta_j\). Now \(\sum_i U_iV_i=\beta\). For \(S(a)_{ij}=\tau(V_i aU_j)\) the scalar corner identities give
\[
 S(a)S(a')=\beta S(aa'),\qquad
 \operatorname{tr}S(a)=\beta\tau(a),\qquad
 x_R(S(a))=x_R(\beta)\rho_x(a).
\]
We use \(\rho_x\) also for the pullback to the new level and its \(O\)-linear extension to the group algebra. Put \(b_0=x_R(\beta)=c_0^{2n}\). Differentiation gives the matrix product derivative
\[
 d(a)=b_0\bigl(\delta S(a)-\delta(\beta)\rho_x(a)\bigr),\qquad
 \operatorname{tr}d(a)=b_0^2\delta\tau(a).
\]
Thus \(d(g)\rho_x(g)^{-1}\) is an adjoint cocycle. Average it with its polarised conjugate under \(c\) to obtain \eqref{eq:p6c5-ordinary-cocycle}. The trace derivative is preserved because the determinant satisfies the polarisation identity. Use \(d\) for the resulting matrix derivative as well. For target \(A_m\) it is locally constant on the group because its entries come from continuous trace functions and the specialised matrices, and the Leibniz rule makes \(\delta\) kill a sufficiently high power of the maximal ideal of \(R\).

For a prescribed inertia relation \(a=g-1\), or an ordered relation \(a=\prod_j(g_j-\psi_{v,j}(g_j))\), the ordinary pair has \(\tau(ha)=0\) for every \(h\) in the CM group. Specialising first gives \(\rho_x(a)=0\) by full span. Differentiating gives \(\operatorname{tr}(\rho_x(h)d(a))=0\), so \(c_0d(a)=0\) by the integral span and matrix units. The equality of inertia kernels gives the required tame condition. The ordered relation holds to first order for the matrix lift with derivative \(c_0d\) and character derivatives \(c_0b_0^2\delta\psi_{v,i}\).
\leanref{UpperTheory.\allowbreak{}V.\allowbreak{}Th.\allowbreak{}Coc.\allowbreak{}gma}{Development.\allowbreak{}AutomorphicArithmetic.\allowbreak{}AutomorphicVarieties.\allowbreak{}TheoremsMuLinked.\allowbreak{}DualModuleCocycles.\allowbreak{}LiftSurj}{33--93}
\leanref{UpperTheory.\allowbreak{}V.\allowbreak{}Th.\allowbreak{}Coc.\allowbreak{}liftOfS}{Development.\allowbreak{}AutomorphicArithmetic.\allowbreak{}AutomorphicVarieties.\allowbreak{}TheoremsMuLinked.\allowbreak{}DualModuleCocycles.\allowbreak{}LiftOfS}{4--38}

For the local calculation use a flag basis and take \(s\in D_v\) with distinct values \(\lambda_i=\kappa_{v,i}(s)\). Diagonalise \(\rho_x(s)\) by an upper triangular change of basis over \(K\). Rescaling the eigenvectors by successive powers of a uniformiser makes the specialised matrices integral in this basis. Clear the fixed denominators of the change of basis and its inverse on the derivative. Multiply further by the product of the nonzero \(\lambda_i-\lambda_j\), \(i<j\), and subtract a commutator to make the derivative at \(s\) diagonal. For entry \((i,j)\) of this commutator one uses the off-diagonal derivative at \(s\) multiplied by the clearing factor divided by \(\lambda_i-\lambda_j\). All these operations on the derivative use coefficients in \(O\).

In the differentiated ordered relation set \(g_i=g\) and \(g_j=s\) for \(j\ne i\). A term differentiating an \(s\)-factor is upper triangular. Its \((i,i)\) entry vanishes because the \(g\)-factor has \((i,i)\) entry zero. For a lower entry \((a,b)\) choose \(i=b\). The coefficient of the derivative of the \(g\)-factor is
\(\prod_{j<i}(\lambda_a-\lambda_j)\prod_{j>i}(\lambda_b-\lambda_j)\), which is nonzero. For diagonal entry \((i,i)\) the corresponding coefficient is \(\prod_{j\ne i}(\lambda_i-\lambda_j)\). Multiplying by a common multiple of these products kills the lower entries and equates each diagonal entry of the matrix derivative to the correspondingly scaled \(\delta\psi_{v,i}(g)\). Transfer back to the integral flag basis with fixed clearing factors. Right multiplication by \(\rho_x(g)^{-1}\) preserves upper triangularity and divides the \(i\)-th diagonal entry by \(\kappa_{v,i}(g)\). On \(\mu_v\) the universal character is prescribed, so this derivative vanishes. The product of the fixed clearing factors over all \(v\), including \(c_0b_0^2\), gives \(c_1\).
\leanref{UpperTheory.\allowbreak{}V.\allowbreak{}Th.\allowbreak{}Coc.\allowbreak{}locAt}{Development.\allowbreak{}AutomorphicArithmetic.\allowbreak{}AutomorphicVarieties.\allowbreak{}TheoremsMuLinked.\allowbreak{}DualModuleCocycles.\allowbreak{}BarOne}{35--107}

For the character comparison at \(v\), write \(\kappa_i=\kappa_{v,i}\) and assume all trace derivatives vanish. Multiply the ordered relation by \(s^a\) on the left and differentiate its trace. The trace terms have zero derivative and the remaining coefficient terms give
\[
 \sum_k\lambda_k^a\prod_{j\ne k}(\kappa_k(g_j)-\kappa_j(g_j))
 \delta(\psi_{v,k}(g_k))=0 .
\]
Combine these equations with the coefficients of \(\prod_{k\ne i}(X-\lambda_k)\). Taking \(g_i=g\) and \(g_j=s\) for \(j\ne i\) gives
\[
 \prod_{j\ne i}(\lambda_i-\lambda_j)^2\delta(\psi_{v,i}(g))=0 .
\]
The product of these nonzero coefficients over \(v,i\) gives \(c_2\).
\end{proof}

\subsection*{Proof of \ref{B6.57}}\label{prf:B6.57}\phantomsection
\begin{proof}
Put \(D_*=D+nq_0\), \(Q=|\mathbf k|\), and suppose \(d_0=\delta(\tau^{\rm univ}(g_0))\ne0\). We construct \(D_*+1\) directions at depth \(m\). Their common loss of precision will be independent of \(m\).

At level \(r\) the coefficient character is \(x\beta\). By \ref{B6.56} and hypothesis~3, the specialised Frobenius polynomials tend coefficientwise to the test polynomials. Each limiting polynomial has distinct roots \(\alpha_{ji}\) in \(O\). Apply Hensel's lemma around each root, with the coefficient errors of valuation greater than twice the valuation of the derivative there. For all large \(r\) this gives roots \(\alpha'_{ji}\in O\) tending to \(\alpha_{ji}\). Let \(f_r\) be the value of the comparison operator \(f\) at these ordered roots. It tends to the nonzero value \(F_*\) obtained with roots \(\alpha_{ji}\) and all \(Q_j=1\). In particular \(v(f_r)=v(F_*)\) eventually.

Choose an \(x\)-eigenfunctional \(\lambda:M\to O\) and \(v_0\in M\) with \(\lambda(v_0)\ne0\). Such a functional comes from the nonzero fibre at \(\ker x\), by taking a linear functional over \(\operatorname{Frac}O\) and clearing denominators on a finite set of \(\Lambda\)-generators. Evaluate the ordered splitting algebra at \(\alpha'\). The map from \(M_w\) to the splitting-algebra module pulls \(\lambda\) back to a functional \(L:M_1\to O\). Denote its character on the commuting operator algebra by \(x_T\). Its diamond values are one and its root values are \(\alpha'_{ji}\). Lift to \(v\in M_1\) the image of \(1\otimes v_0\) under the reverse comparison. The composite relation gives \(L(v)=f_r\lambda(v_0)\).

The coefficient map \(\beta\) is onto: its compact image contains the adic generators of \(B\) by compatibility of the pair and surjectivity of the level map. Lift \(t_0\) of hypothesis~2 to \(t_1\). In the operator algebra put
\[
 P_0=\prod_{j,i}\prod_{a\ne i}(u_{ji}-\alpha'_{ja}),\qquad
 t=t_1P_0 f^{E+1}.
\]
Write \(S=\Lambda[\Delta]\) and \(I=\ker(x_T|_S)\). For every operator \(b\) in the algebra generated by coefficients, roots and diamonds we have
\begin{equation}\tag{6.57.5}\label{eq:p6c5-operator-ideal}
 x_T(b)=0\quad\Longrightarrow\quad tbM_1\subseteq IM_1 .
\end{equation}
Indeed pass to diamond coinvariants and apply the comparison to the splitting-algebra module. Modulo \(\ker x\) the polynomial of the \(j\)-th Frobenius is \(\prod_a(X-\alpha'_{ja})\). Substitution of \(u_{ji}\) shows that \((u_{ji}-\alpha'_{ji})P_0=0\) there. Thus multiplication by \(P_0\) replaces every root by its chosen value modulo the extended \(\ker x\). Every diamond already acts as one. It follows that \(P_0b\) has image in the extended \(\ker x\) times that module. Multiplication by \(t_0\) carries it into the weight ideal by hypothesis~2 and scalar extension. Comparing back introduces a factor \(f\). The kernel of \(M_1\to M_w\) lies in \(IM_1\), which proves \eqref{eq:p6c5-operator-ideal}.

Take \(r\ge m\) large enough for the preceding constructions. For a chosen weight or diamond coordinate take its coordinate derivation \(d_S:S\to A_m\) at \(x_T|_S\), with value one on that coordinate. A diamond relation has derivative \(\ell^{a_j}d_S(y_{ji})=0\) since \(a_j\ge r\). Give \(O\oplus\eta A_m\), \(\eta^2=0\), the \(S\)-action \(x_T+\eta d_S\). Projectivity of \(M_1\) over \(S\) lifts \(L\) to this module. Its second component \(H_1\) satisfies
\[
 H_1(sw)=x_T(s)H_1(w)+d_S(s)L(w),\qquad H_1(I^2M_1)=0.
\]
Set \(\partial(b)=H_1((b-x_T(b))t^2v)\). The product of two centred operators times \(t^2\) maps \(M_1\) into \(I^2M_1\) by two applications of \eqref{eq:p6c5-operator-ideal}. Expanding \(bb'-x_T(bb')\) now gives the Leibniz rule for \(\partial\). On \(S\) it equals \(x_T(t)^2 f_r\lambda(v_0)d_S\).

The trace in brackets in \eqref{eq:p6c5-operator} has value zero at \(x_T\), and multiplication by \(t\) annihilates it. Thus \(\partial\) of this trace is zero. In the part involving \(\tau_{r,j}-1\), differentiating a polynomial coefficient gives zero because \(\tau_{r,j}\) specialises to the identity. In the part involving \(y_{ji}\), the coefficient of \(\partial(y_{ji})\) is \(p_{ji}^{E+1}\), where \(p_{ji}=\prod_{a\ne i}(\alpha'_{ji}-\alpha'_{ja})\). Indeed the polynomial has value zero on the other eigenlines and value \(p_{ji}^{E+1}\) on line \(i\). We therefore obtain the same coordinate responses on the weights and the coefficient-ring tests
\begin{equation}\tag{6.57.2}\label{eq:p6c5-trace-tests}
 T_{ji}=\operatorname{tr}_{u_r}\left[
 \left(\prod_{a\ne i}(\phi_{r,j}-\alpha'_{ja})\right)^{E+1}
 (\tau_{r,j}-1)\right],
\end{equation}
with the diamond response multiplied by \(p_{ji}^{E+1}\). Both \(P_0(\alpha')\) and each \(p_{ji}\) divide \(f_r\). The weight response divides \(x(t_0)^2\lambda(v_0)f_r^{2E+5}\), and every diamond response after the extra factor divides \(x(t_0)^2\lambda(v_0)f_r^{3E+6}\). Since \(v(f_r)=v(F_*)\), the fixed nonzero scalar
\[
 c_3=x(t_0)^2\lambda(v_0)F_*^{3E+6}
\]
is divisible by each response in \(O\), before reduction at depth \(m\). Scale the coordinate derivations by these integral quotients and restrict to \(R\). The resulting \(D_*\) derivations \(e_a\) have response matrix \(c_3\,\mathrm{Id}\).

Adjoin \(e_*=\delta\beta\bmod\mathfrak m_O^m\). It kills every weight by assumption and every \(T_{ji}\) because \(h_r(\tau_{r,j})=1\). Denote the resulting index set by \(\mathcal I\). Use the constants \(c_1,c_2\ne0\) of \ref{B6.57a}. Compatibility of the auxiliary pair gives the prescribed specialised representation and characters, and hypothesis~1 identifies the tame kernels. Hence the map
\begin{equation}\tag{6.57.3}\label{eq:p6c5-count-map}
 \Phi:O^{\mathcal I}\longrightarrow\mathcal S^{\rm tr}_{r,m},\qquad
 \Phi(b)(g)=c_1\sum_{a\in\mathcal I}b_a e_a(\operatorname{tr}_{u_r}(g))
\end{equation}
takes values in the stated trace set. If \(b\) is in its kernel, \(c_1\sum b_a e_a\) kills traces and therefore all \(T_{ji}\), by trace linearity. After multiplication by \(c_2\) it also kills the weight characters. The response matrix gives \(c_1c_2c_3 b_a=0\) in \(A_m\) for every \(a\ne *\). Evaluate \(\Phi(b)\) at a lift of \(g_0\) and multiply by \(c_2c_3\). The remaining equation is \(c_1c_2c_3d_0b_*=0\). With \(c_*=c_1c_2c_3d_0\) this proves
\begin{equation}\tag{6.57.4}\label{eq:p6c5-count-kernel}
 b\in\ker\Phi\quad\Longrightarrow\quad c_*b_a\in\mathfrak m_O^m
 \quad(a\in\mathcal I).
\end{equation}

Let \(v_0'=v_O(c_*)\), take \(h=v_0'D_*+C+1\) and set \(m=h+v_0'\). Now choose \(r\) sufficiently large for this \(m\), including the bound \eqref{eq:p6c5-count-bound}. Representatives modulo \(\mathfrak m_O^h\) in each coordinate have distinct images by \eqref{eq:p6c5-count-kernel}. We obtain \(Q^{h(D_*+1)}\le Q^{mD_*+C}\), whereas \(h(D_*+1)-(mD_*+C)=1\). This contradiction proves the assertion.
\end{proof}

\subsection*{Proof of \ref{B6.58c}}\label{prf:B6.58c}\phantomsection
\begin{proof}
Each \(A_j\) is finite over \(A_0\). In monomials in a fixed finite set of homogeneous generators consider the exponent vectors of weight \(j\). They have only finitely many coordinatewise minimal elements, and division by a minimal monomial leaves weight zero. Choose finite generating lists of \(A_{-2}\) and \(A_2\). Every two-sided point has \(fg\ne0\) for some pair from these lists, since a surviving generator of weight \(1\) or \(-1\) has a surviving square. Rescale the point so that \(f=1\). Over \(K\) the slice \(f=1,\ g\ne0\) on \(A/I\) is of finite type by hypothesis~(2). The orbit tangent has \(f\)-coordinate \(-2\). Hypothesis~(4) therefore gives zero tangent space to the slice at every closed point. Each such local ring is a field by Nakayama. A finite type zero-dimensional \(K\)-algebra is Artinian, so each slice is finite and reduced. The union of the slices over the finite generating pairs gives the required finite orbit cover \(T\).

The ideal \(\mathfrak a\) is finitely generated since \(A\) is Noetherian, and \(\mathfrak a/\mathfrak a^2\) is finite over \(A/\mathfrak a=O\). The saturated weight-\(-1\) cotangent lattice \(C_{-1}\) is free of rank \(r\) by the linear term condition in hypothesis~(1). Smith reduction on the \(r+1\) columns represented by \(F_i\) gives a basis with representatives \(y_i\in A_{-1}\) and an invertible change among the \(F_i\) satisfying \eqref{eq:p6c5-cot-smith}. The \(F_i\) span rationally by hypothesis~(1), hence all \(b_i\) are finite and nonnegative, which proves \eqref{eq:p6c5-cot-length}. The change preserves \(I\).

The \(m\) images of the \(Z_i\) have independent linear terms spanning the weight-zero variables. Use these images as coordinates by formal inversion and set them equal to zero to obtain \(\phi\). If \(\phi(a)=0\), add to \(a\) the product of an \(F_i\) with nonzero linear term and a weight-one element with nonzero linear term. Both exist by hypothesis~(1) and \(r\ge1\). The product has nonzero image in the power series domain, and the new \(a\) agrees with the old one on \(A/I\).

Finally, for \((s,L)\) admissible on \(N\), both \(N/sN\) and \(0:_N s\) are annihilated by a power of \(\mathfrak m_{A_0}R\,R_{++}\) modulo \((L)\), so their Koszul homologies with respect to \(L\) have finite \(O\)-length in coordinatewise large multidegrees. Tensoring the Koszul complex of \(L\) with the exact sequences relating multiplication by \(s\) on \(N\), its kernel \(0:_N s\), and its cokernel \(N/sN\) gives the two-term complex on \(s\) and proves \eqref{eq:p6c5-koszul-recursion}.
\end{proof}

\subsection*{Proof of \ref{B6.58}}\label{prf:B6.58}\phantomsection
\begin{proof}
First consider the framed presentation. In an Artinian test algebra the congruence subgroup of the matrix group is a finite \(\ell\)-group. On the open kernel of the residual representation one may therefore use generators of its maximal pro-\(\ell\) quotient. There are finitely many by the hypothesis on \(\F_\ell\)-characters. Adjoin entries for lifts of these generators and for a set of coset representatives. Impose the group relations, the form equations and the local equations, and adjoin entries for the frame changes with their intertwining equations. Use the closed ideal of these relations in the power series ring over \(B\). At every finite depth the equations give precisely the allowed continuous representations with their frames: on a free pro-\(\ell\) group generator images in a finite \(\ell\)-group determine a continuous map, and imposing the relators makes it descend. Taking inverse limits therefore gives the universal complete ring \(U\).

Over \(\mathbf k[\eta]/(\eta^2)\), with \(B\) acting through \(\mathbf k\), write a lift as \((1+\eta z)\bar r\) and the \(v\)-frame as \(1-\eta b_v\). Multiplication requires \(dz=0\), the form equation places \(z\) in the Lie algebra of isometries, and the local equations are \(z\alpha_v=\delta b_v\), \(b_{v_0}=0\), \(z\beta_q|_{N_q}=0\). In weight \(j\) these are exactly the pairs defining \(H^1_{\mathrm{fr},j}\). The entries of the universal lift and the frame changes separate these tangent vectors. Select a basis of the dual in each weight from their entry evaluations. The corresponding entries give a minimal relative presentation \(P=B[[X_1,\ldots,X_e]]\twoheadrightarrow U\) by complete Nakayama. Conjugation by the diagonal weight torus acts on an \(ab\)-entry by the character \(w_a-w_b\). The equations are stable under this action, so their closed ideal is homogeneous: at each finite power series depth take the coefficients of the independent torus characters and then pass to the inverse limit.

Write \(J_0=\ker(P\to U)\). A weight-\(j\) functional on \(J_0/\mathfrak m_PJ_0\) gives a small extension of \(U\) by \(\mathbf k\) of that weight, by pushing out \(P/\mathfrak m_PJ_0\to U\). Choose lifts of the universal matrices across it. For an isometry lift with form defect \(D\) in the square-zero kernel, right multiplication by \(1-\tfrac12 P_f^{-1}D\), where \(P_f\) is the form matrix, removes the defect. This works for both types of form and uses odd residual characteristic. Perform the calculation on the isometry factor relative to \(r_0\). The multiplication defects give a continuous two-cocycle \(c\), by associativity. Comparing with the prescribed framed lifts at \(v\) and with local solutions at \(q\) gives one-cochains \(u_v,u_q\). Such local solutions exist by the small-extension hypothesis. Their differentials are the restrictions of \(c\). These cochains can be chosen continuously by lifting on finite congruence quotients with discrete kernel and using the fixed lifts on the other coordinates.

Changing the choices by a global cochain \(u\), local frame corrections \(b_v\), and local solution differences \(e_q\) changes the triple by the boundary in the statement. At \(q\) precisely the cocycles vanishing on \(N_q\) preserve the local solution. Thus we obtain a class in \(\operatorname{Ob}_j\). The construction is linear in the functional, since one may compute the defects before pushing out and apply the functional to their entries. If the class vanishes, the corrections give a global representation with all the required local conditions. Simultaneous conjugation restores the identity frame at \(v_0\). Universality gives a section over \(B\) of the small extension. Composing this section with \(P\to U\) and comparing with the quotient map from \(P\) to the extension gives two maps equal on \(B\), differing on the variables by elements of the small kernel. Their difference vanishes on \(\mathfrak m_BP+\mathfrak m_P^2\), which contains \(J_0\) by relative minimality. On \(J_0\) their difference is the original functional, up to sign. It is therefore zero. This injects the dual minimal relation space into \(\operatorname{Ob}_j\). Choose homogeneous lifts of minimal relation generators and pad the lists with zeros to obtain (6.58.1).

For (6.58.2), an integral \(x\)-derivation kills \(\mathfrak a^2\) by the product rule and kills \(S_x\) since \(O\) is torsion-free. Its reduction is determined in weight \(j\), and a reduction killing all \(F_i\) factors there through \(C_j(F)\). Thus the given additive injection gives an injection into \(\operatorname{Hom}_O(C_j(F),O/(c))\). If \(C_j(F)\) has finite length, a cyclic summand \(O/(\pi^u)\) contributes \(\min(u,v_O(c))\) to the Hom length. Taking cardinalities proves (6.58.2). The inequality also holds when its right hand side is infinite.

For (6.58.3), apply Lemma~\ref{B6.58c} to obtain a finite two-sided orbit cover, the finite length of \(C_{-1}(F)\) in \eqref{eq:p6c5-cot-length}, the chart \(\phi\), the replacement of \(a\) modulo \(I\) with \(\phi(a)\ne0\), the graded algebra \(R\), the module \(M=R/J\), and the section list \(\mathbf t\). Replacing \(a\) modulo \(I\) preserves \(\xi(a)\) at every point \(\xi\) of \(A/I\), so we use the modified \(a\) for any given finite two-sided orbit cover \(T\). Lemma~\ref{B6.58b} gives admissibility of \((\mathbf t,a)\) on \(M\), so the hypothesis of Lemma~\ref{B6.58a} holds on \(M/aM\). For a multidegree \(d\) large enough in all three coordinates the two lemmas give
\[
 2^{r+r'}\sum_i b_i
 \le\chi_{\mathbf t}(M/aM)_d
 \le 2^{r+1+r'}\sum_{\xi\in T}v(\xi(a)).
\]
Division by \(2^{r+r'}\) and \eqref{eq:p6c5-cot-length} prove (6.58.3).
\end{proof}

\subsection*{Proof of \ref{B6.58a}}\label{prf:B6.58a}\phantomsection
\begin{proof}
For a monomial in the chart after the weight-zero variables have been set to zero, let \(u\) be the sum of the absolute weights of its negative variables, counted with multiplicity, and \(v\) the corresponding sum for its positive variables. Its weight is \(v-u\). In an element of \(A_{-2}\) the least possible \(u\) is \(2\). In \(\phi(F_{r+1})\) it is at least \(2\): a term of weight \(-1\) with \(u=1\) would be linear, and \eqref{eq:p6c5-cot-smith} kills the linear term. Thus in the third coefficient module the least possible \(u\) is \(4\). Its terms at \(u=4\) have \(v=2\). In degree \(d=(d_0,d_1,d_2)\) the least possible pair is consequently
\[
 u(d)=2d_0+4d_2,\qquad v(d)=2d_1+2d_2 .
\]
Taking the terms with this pair in each multidegree is multiplicative, since the two bounds are additive. Write \(P\) for the image. This quotient kills \(J\) and \(a\): a chart term of \(a\) has weight zero and zero constant term, hence positive \(u\). We obtain \(M/aM\twoheadrightarrow P\).

In each degree \(P_d\) is a finite \(O\)-module since \(R_d\) is finite over \(A_0\) and the action of \(A_0\) on leading terms is through \(x\). It embeds into a space of polynomials over \(\operatorname{Frac}O\), so is free. We compute its rank. Over the fraction field the cotangent hypothesis lifts any individual variable by a homogeneous element with that linear term. For a variable of absolute weight one the sole terms of the same sign and least absolute count are linear. For absolute weight two the additional terms of least count are products of two variables of absolute weight one and the same sign. Subtract products of their lifts to isolate the chosen weight-two variable. Clearing denominators gives nonzero scalar multiples integrally. This supplies every monomial block of absolute weight two on either side. Every monomial of even absolute weight on one side is a product of such blocks: keep each weight-two variable and pair the weight-one variables. For multidegree \(d\), allocate \(d_0\) negative blocks and \(d_1\) positive blocks to the first two directions. Allocate \(2d_2\) negative blocks and \(d_2\) positive blocks to the third, using \(A_{-2}^2A_2\). Hence \(P_d\) spans the whole space of bidegree \((u(d),v(d))\). With formal variables \(X,Y\) this gives
\begin{equation}\tag{6.58.8}\label{eq:p6c5-intersection-hilbert}
 \operatorname{length}_O(P_d/\pi P_d)=[X^{u(d)}Y^{v(d)}]H,\qquad
 H=\frac1{(1-X)^r(1-X^2)^s(1-Y)^{r'+1}(1-Y^2)^{s'+1}} .
\end{equation}

On \(P\) the section \(F_i^2T_0\) acts as \(\pi^{2b_i} y_i^2T_0\) for \(i\le r\), since the congruence in \eqref{eq:p6c5-cot-smith} identifies the linear terms in the rational chart. Remove these scalar factors one at a time. For a section \(z\) of degree \(e\) and the list \(L\) of other sections, factor multiplication by \(\pi^h z\) into multiplication by \(z\) followed by multiplication by \(\pi^h\). The exact triangle of the three mapping cones, tensored with the Koszul complex of \(L\), gives
\[
 \chi_{(\pi^h z,L)}(P)_d
   =\chi_{(z,L)}(P)_d+h\,\chi_L(P/\pi P)_d .
\]
Here multiplication by \(\pi\) is injective, and \(P/\pi^hP\) has \(h\) successive quotients \(P/\pi P\). The replacement enlarges the section ideal and preserves admissibility. On \(P/\pi P\), the \(A_0\)-action factors through its residue field, so each complex term has finite length.

When one first-direction section is removed, the other multidegrees consist of \(r-1+s\) first, \(r'+s'\) second, and one third direction. The alternating sum of lengths of complex terms is therefore the coefficient of \(X^{u(d)}Y^{v(d)}\) in
\[
 H(1-X^2)^{r-1+s}(1-Y^2)^{r'+s'}(1-X^4Y^2).
\]
This expression equals
\[
 (1+X)^{r-1}(1+Y)^{r'}
 \left(\frac1{(1-X)(1-Y)}
       +\frac{Y^2(1+X+X^2+X^3)}{(1-Y)(1-Y^2)}\right).
\]
The second term has bounded \(X\)-degree. In the first term the coefficient for large \(u(d),v(d)\) is the product of the sums of the two numerator coefficient lists, namely \(2^{r-1+r'}\). All shifts used in the Koszul complex have nonnegative remaining multidegrees when \(d\) is large enough. Removing the scalar for \(i\) consequently contributes \(2b_i\,2^{r-1+r'}\).

By the nonnegativity assertion of \ref{B6.58b}, the characteristic on \(P\) after all removals is eventually nonnegative. The same holds for the original list on the kernel of \(M/aM\to P\), since admissibility passes to submodules. Additivity in that exact sequence gives \eqref{eq:p6c5-intersection-lower}.
\end{proof}

\subsection*{Proof of \ref{B6.58b}}\label{prf:B6.58b}\phantomsection
\begin{proof}
Put \(B=A_0\), \(L=I\cap B\), \(D=A_{-2}A_2\subset B\), and \(W_d=(M/\mathbf t M)_d\). We first locate the primes over \(J+(\mathbf t)\) which avoid \(R_{++}\). Let \(P\) be such a prime. Choose \(fT_0,gT_1,hT_2\notin P\) in the three generating modules. Since
\[
 (F_{r+1}^2T_0)(hT_2)=(hT_0)(F_{r+1}^2T_2),
\]
we have \(F_{r+1}^2T_0\in P\). If \(cG_i\) has weight zero, multiplication by \(fT_0\) expresses its product as \(cf\cdot G_iT_0\), so \(cG_i\in P\). Use the same calculation on \((cF_i)^2\) when \(cF_i\) has weight zero. Use \(gT_1\) on the positive equations. Finally \(Z_i\in J\), since the chart sets these variables to zero. Decomposing coefficients into homogeneous parts in an expression for an element of \(I\cap B\) proves \(L\subseteq P\cap B\).

We claim
\[
 D\subseteq P\quad\Longrightarrow\quad P\cap B=\mathfrak m_B .
\]
Suppose \(D\subseteq P\) and \(\pi\notin P\). Since \(B/L\) is finite over \(O\), an embedding of \(\operatorname{Frac}(B/(P\cap B))\) into \(K\) gives a point \(\sigma\) of \(B\). We construct points of both strict signs over it. A monomial of weight \(-2n\) contains at least \(n\) negative blocks of weight \(-2\), formed from weight-\(-2\) generators or pairs of weight-\(-1\) generators. The remaining factors can be taken to have weight zero after extracting total negative weight \(2n\). Thus \(A_{-2n}=A_{-2}^n\): equivalently one extracts \(n\) blocks and absorbs the weight-zero remainder into a block for \(n>0\). The positive version holds as well.

On the even-weight subring send a homogeneous \(u\in A_{-2n}\) to \(uT_0^n\bmod P\) in \(\operatorname{Frac}(R/P)\), use the residue map on \(B\), and send positive weights to zero. For mixed multiplication with \(v\in A_{2k}\), \(0<k\le n\), multiply the proposed image of \(vu\) by \((fT_0)^k\). The result is \((vf^k)(uT_0^n)\), which vanishes since \(vf^k\in D^k\). For \(k>n\) the product has positive weight. This proves multiplicativity. All the negative even equations and the squares of the negative odd equations vanish. The entire graded algebra is integral over the even-weight subring, since every homogeneous generator has its square in that subring. By lying over, extend the point to an algebraically closed overfield, choosing the overfield compatibly with \(\sigma\). The resulting point kills \(I\) and every positive-weight element, and takes a nonzero value on \(f\). These conditions define a nonempty finite type locus in the fibre over \(\sigma\), hence give a \(K\)-point. Interchanging the two signs gives a positive point with \(g\ne0\). There is also a fixed point over \(\sigma\): kill all nonzero weights. For a product of opposite weights summing to zero its square belongs to \(D\), by the even-weight identity, so this prescription is multiplicative. Hypothesis~(3) of Lemma~\ref{B6.58c} forces this fixed point to be \(x_K\).

At the negative point over \(x_K|_B\), every term of weight \(-1\) in \(\mathfrak a^2\) vanishes. It has either a positive-weight factor or a weight-zero factor in \(\ker x\). The same holds for the saturation since nonzero elements of \(O\) are invertible in \(K\). The \(F_i\) span the rational weight-\(-1\) cotangent and vanish at the point, so all of \(A_{-1}\) vanishes there. In weight \(-2\) the additional possibility in \(\mathfrak a^2\) is a product of two weight-\(-1\) elements. It too vanishes. The \(G_i\) now give vanishing of \(A_{-2}\), contrary to \(f\ne0\). Therefore \(\pi\in P\). Since the local ring \(B/L\) is finite over \(O\), its unique prime over \(\pi\) is its maximal ideal. This proves the claim.

The two prime containments imply, for some \(h\) and all coordinatewise large \(d\),
\begin{equation}\tag{6.58.9}\label{eq:p6c5-intersection-support}
 L^h W_d=0,\qquad \mathfrak m_B^h W_d\subseteq DW_d .
\end{equation}
Indeed apply the radical criterion first to \(J+(\mathbf t)\), then to \(J+(\mathbf t)+D\), multiplying by powers of \(R_{++}\) to include the primes containing \(R_{++}\). The ideals are finitely generated, and sufficiently large pieces of this algebra generated in unit multidegrees are obtained by multiplication by any fixed power of \(R_{++}\). The first inclusion makes \(W_d\) finite over \(O\): it is finite over \(B/L^h\), and the filtration by \(L^i/L^{i+1}\) shows that \(B/L^h\) is finite over \(O\).

For a prime \(P\) as above with \(\pi\notin P\), set \(p=P\cap B=\ker\sigma\) with \(\sigma:B\to K\) as before. We have \(D\not\subseteq p\). Choose \(f\in A_{-2}\), \(g\in A_2\) with \(\sigma(fg)\ne0\). The map \(B/L\to A/I\) splits as a module map by projection onto weight zero. The fibre over \(\sigma\) is therefore a nonzero finite type \(K\)-algebra and has a \(K\)-point. Every such point is two-sided, so hypothesis~(5) of Lemma~\ref{B6.58c} gives \(\sigma(a)\ne0\). We have proved that every prime over \(J+(\mathbf t,a)\) contains \(R_{++}\) or \(\mathfrak m_B\). The radical criterion gives admissibility.

We next bound the length at a horizontal prime \(p\) in \(\operatorname{Supp}_B W_d\), meaning one omitting \(\pi\). By \eqref{eq:p6c5-intersection-support} and Nakayama, \(D\not\subseteq p\). Use \(f,g\) as above. On \(fg\ne0\) the algebra \((A/I)\otimes_O K\) is reduced. Indeed adjoining an invertible square root of the unit \(f\) is faithfully flat. The torus action then identifies this cover with the product of the slice \(f=1,\ g\ne0\) and \(\mathbb G_m\): the chosen root specifies the rescaling, and reversing the rescaling recovers the point and the root. The slice is reduced by Lemma~\ref{B6.58c}. Descent and scalar restriction give reducedness over \(\operatorname{Frac}O\) on this open set. The finite weight-zero subalgebra there is reduced, so its localisation at \(p\) is a field. In the graded algebra localised at \(B\setminus p\), take a nonzero element \(h_0\) of weight \(w\). Use powers of the unit \(f\) to get a unit \(e\) of weight \(-2w\). Then \(h_0^2e\) is nonzero by reducedness and belongs to the weight-zero field. Every \(h\) of weight \(w\) equals
\[
 (hh_0 e)(h_0^2e)^{-1}h_0 .
\]
Thus each localised weight piece of \(A/I\) has \(B_p\)-length at most one.

Replace an odd generator \(u\) of \(I\) by \(u^2\). In the algebra modulo the other equations, the kernel of the map from the quotient by \(u^2\) to the quotient by \(u\) is an image of the quotient by \(u\), by multiplication by \(u\) with its weight shift. The maximum length of a weight piece therefore at most doubles at each replacement. There are \(r+1+r'\) replacements. To apply this to \(W_d\), invert \(fg\). The three coefficient modules of \(R\) then have bases \(f,g,f^2g\) over the weight-zero ring, since \(f,g\) are units and \(f^2g\) belongs to the third module. In multidegree \(d\) with all coordinates positive, divide by the product of these bases with exponents \(d_0,d_1,d_2\). Each section relation then imposes its coefficient equation multiplied by the inverse basis of the corresponding weight. These are the weight-zero relations from the generators of \(I\) with the odd ones squared. The \(Z_i\) already vanish modulo \(J\). The remaining relations of \(J\) give a further quotient. Hence
\begin{equation}\tag{6.58.10}\label{eq:p6c5-intersection-multiplicity}
 \operatorname{length}_{B_p}(W_d)_p\le 2^{r+1+r'} .
\end{equation}

Use a prime filtration of \(W_d\) over \(B\). A vertical factor has finite \(O\)-length. Multiplication by \(a\) on it has kernel and cokernel of equal length. A horizontal factor \(B/p\) is finite torsion-free over \(O\). Multiplication by \(a\) on it has cokernel length equal to the valuation of its determinant, namely \(\sum_{\sigma}v(\sigma(a))\), where \(\sigma\) runs over the embeddings of its fraction field in \(K\). The multiplicity of the factor is \(\operatorname{length}_{B_p}(W_d)_p\), as seen by localising the filtration. Each \(\sigma\) extends to a two-sided point by the fibre argument, and restriction to \(B\) is constant on a torus orbit. Thus the restrictions of \(T\) cover these embeddings. Their valuations are nonnegative since \(B/L\) is integral over \(O\). Additivity and \eqref{eq:p6c5-intersection-multiplicity} give the upper bound for \(\chi_{(a)}(W_d)\).

Finally consider an admissible list on a finite graded module. Its Koszul homology has finite-length large pieces by the radical criterion. Take a finite homogeneous prime filtration and use additivity. Factors supported on \(V(R_{++})\) vanish in large multidegrees. On a remaining domain factor, a nonzero first section acts injectively, so \eqref{eq:p6c5-koszul-recursion} reduces nonnegativity to that for a shorter list. If the first section acts as zero, the tail is already admissible. The characteristic is then the difference of its tail characteristics in degrees \(d\) and \(d-e\), with \(e\) the first degree. Choose a monomial of degree \(e\) outside the prime, by taking surviving generators in the unit directions with their indicated exponents. Its multiplication sequence identifies the same difference with the tail characteristic on its quotient. Induction ends with the length of a module for the empty list. This proves eventual nonnegativity.

Since \(\phi(a)\ne0\), multiplication by \(a\) is injective on \(M\), a subring of a polynomial ring over the chart domain. Computing first with \(a\) gives \(\chi_{(\mathbf t,a)}(M)_d=\chi_{\mathbf t}(M/aM)_d\). Computing instead with the sections of \(\mathbf t\) successively gives at most \(\chi_{(a)}(W_d)\) by \eqref{eq:p6c5-koszul-recursion}, since each subtracted characteristic is eventually nonnegative. This proves both inequalities.
\end{proof}

\subsection*{Proof of \ref{B6.59}}\label{prf:B6.59}\phantomsection
\begin{proof}
On \(J_A\) the restriction of \(z\) is a homomorphism to a group of exponent \(\ell\). Its image has order dividing the index in \eqref{eq:p6c5-unrev}, so it is zero. Thus \(z|_{G_F}\) is an unramified homomorphism for \(F=\Q(\mu_N)\). The nontriviality of \(\vartheta\) gives \(N>2\). For each \(t\in(\Z/N)^\times\), choose a lift \(\sigma_t\) of the automorphism acting by \(t\) on \(\mu_N\).

Choose a positive odd \(a\) prime to \(N\) with \(\operatorname{red}\vartheta(a)\ne\operatorname{red}a\). If \(\ell\nmid N\), take \(a=\ell\). If \(\ell\mid N\), both residual characters factor through \((\Z/N)^\times\), and a class where they differ has an odd representative. With \(y=[a^{-1}t]_N\), the integral smoothing coefficients satisfy
\[
 E_t=(a-1)/2-(ay-[t]_N)/N,\qquad
 \sum_t E_t\vartheta^{-1}(t)
 =(1-a\vartheta^{-1}(a))B_{1,\vartheta^{-1}}.
 \tag{6.59.6}
\]
The constant term drops out by the vanishing character sum, and reindexing \(t=ay\bmod N\) gives the second term of the product. The right side is a unit. For \(u\in G_F\) the cocycle law gives \(z(\sigma_t^{-1}u\sigma_t)=\bar\vartheta(\sigma_t)^{-1}z(u)\). The integral Stickelberger relation of \ref{B6.10a} \cite[Theorem~6.10]{washington-1997} makes the sum of these conjugate values weighted by \(E_t\) zero. Reducing \((6.59.6)\) therefore gives \(z|_{G_F}=0\).
\leanref{UpperTheory.\allowbreak{}V.\allowbreak{}Odd.\allowbreak{}homKill}{Development.\allowbreak{}AutomorphicArithmetic.\allowbreak{}AutomorphicVarieties.\allowbreak{}OddResidueBounds.\allowbreak{}HomKill}{9--50}
Prime-to-\(\ell\) roots reduce injectively, so choose \(t_0\) with \(\bar\vartheta(\sigma_{t_0})\ne1\). The cocycle factors through the abelian quotient \(G_\Q/G_F\). Expanding its equal values at \(g\sigma_{t_0}\) and \(\sigma_{t_0}g\) gives \eqref{eq:p6c5-bernoulli-coboundary} with \(b=z(\sigma_{t_0})/(\bar\vartheta(\sigma_{t_0})-1)\).

Apply \ref{B6.59a} to \(\phi,\xi\) and to their inverses. Both products have order prime to \(\ell\), and (6.59.3) implies \(n\ge4\). A Galois element acting by \(1+2^{n-1}\) on \(\mu_{2^n}\) and trivially on \(\mu_{\ell f}\) has product character value \(-1\) by primitivity of \(\xi\). Thus either reduced product is nontrivial and differs from \(\bar\varepsilon_\ell\).
\end{proof}

\subsection*{Proof of \ref{B6.59a}}\label{prf:B6.59a}\phantomsection
\begin{proof}
Write \(M=2^n\), \(m=n-e\) and \(s=f2^e\). Then \(\ell\) is odd, \(m\ge e\ge2\) and \(1<s\le2^m\). Extend characters by zero off the units and put
\[
 A_u=\sum_{\substack{0\le j<fM\\j\equiv u\ (M)}}j\phi(j),\qquad
 g(u)=\frac{A_u-A_0}{fM},\qquad c=M^{-1}\sum_{j\bmod f}\phi(j).
\]
Translate the summands of \(A_{u-ft}\) by \(ft\). Their original indices run through all residues modulo \(f\), and \(\phi\) is unchanged by the translation. Wrapping a new index \(j\) into the interval \([0,fM)\) subtracts \(fM\phi(j)\) and gives exactly the indices in the sum below. Hence
\begin{equation}\tag{6.59.7}\label{eq:p6c5-dyadic-shift}
 g(u)-g(u-ft)=tc-\sum_{\substack{0\le j<ft\\j\equiv u\ (M)}}\phi(j)
 \qquad(0\le t<M).
\end{equation}
Take \(ft\equiv u\bmod M\). Since \(g(0)=0\), the formula makes \(g(u)\) integral and its reduction fixed by the \(Q\)-power Frobenius. Here \(M\) is a unit and \(\phi(j)^Q=\phi(j)\). The character sum of \(\xi\) is zero since \(\xi\) is primitive and \(n\ge4\). Grouping the defining Bernoulli sum by residues modulo \(M\) gives \(B_{1,\phi\xi}=\sum_u\xi(u)g(u)\).

Use bars for reduction into \(\overline{\F}_\ell\) and suppose the sum vanishes there. Reindexing gives \(\sum_u\bar\xi(u)\bar g(uy)=0\) for each unit \(y\). Successive differences of squares give \(v_2(Q^{2^j}-1)=e+j\). Thus the powers of \(Q\bmod M\) are \(1+2^eh\) for \(0\le h<2^m\). Average the \(2^m\) successive Frobenius conjugates of the sum. Each \(\bar g(uy)\) is fixed, and the geometric sum on the character values gives
\[
 U=\{u\in(\Z/M)^\times:\bar\xi(u)^{2^e}=1\},\qquad
 \sum_{u\in U}\bar\xi(u)\bar g(uy)=0 .
\]
The element \(h_0=1+M/2\) generates the last reduction kernel, so \(\bar\xi(h_0)=-1\) by primitivity. It is the unique involution of the cyclic subgroup of units congruent to \(1\bmod4\), on which \(\bar\xi\) is therefore faithful. For \(u\in U\), this gives \(u^{2^e}=1\bmod M\), since \(u^{2^e}\) belongs to that subgroup. In the factorisation of \(u^{2^e}-1\), the factors \(u^{2^i}+1\) for \(1\le i<e\) each have \(2\)-adic valuation one, as does one of \(u-1,u+1\). Hence \(u\equiv\delta_u\bmod 2^m\) for a unique \(\delta_u\in\{1,-1\}\).

Subtract the identities at \(y=1\) and \(y=1-s\). Since \(M\mid 2^ms\), we have \(u(1-s)\equiv u-\delta_us\bmod M\). For \(\delta_u=1\), apply \eqref{eq:p6c5-dyadic-shift} with \(t=2^e\). The boundary term occurs at \(u=1\), the sole residue of this sign with \([u]_M<s\). For \(\delta_u=-1\), apply the same shift at \(u+s\) with the opposite sign. The boundary term occurs at \(u=-1\), where \([u+s]_M=s-1\), the sole residue in \([0,s)\) congruent to \(s-1\bmod 2^m\). Thus
\[
 0=2^e\bar c\sum_{u\in U}\delta_u\bar\xi(u)
 -\bar\phi(1)+\bar\xi(-1)\bar\phi(s-1).
\]
Multiplication by \(h_0\) preserves \(U\) and \(\delta_u\) and negates the character value, so the sum cancels. Since \(s-1\equiv-1\bmod f\) and the product character is odd, the remaining terms equal \(-2\), a contradiction.
\end{proof}

\section*{Proofs for Chapter~\ref{ch:B6-6}: The exponent 401 and the downstream window}

\subsection*{Proof of \ref{B6.61}}\label{prf:B6.61}\phantomsection
\begin{proof}
Put $u=x^{1/100}\ge2$. The bounds on $\varepsilon$ give
\[
 f\le4u^{216},\quad k\le2u^{470},\quad r\le2k\le4u^{470},
 \quad m\ge u^{49},\quad L^2\le u^4.
\]
Indeed, $2+8\varepsilon\le216/100$,
$9/2+10\varepsilon\le470/100$,
$1/2-\varepsilon/16\ge49/100$, and $2\varepsilon\le4/100$.
Hence
\[
 \frac{f^2(k-1)^3L^2}{m}
 \le128u^{2\cdot216+3\cdot470+4-49}=128u^{1797},
 \qquad
 f+\frac{f^2(k-1)^3L^2}{m}\le132u^{1797}.
\]
Using $\lceil v\rceil<v+1$ in the definition of $n_W$ gives
\[
 n_W\le5280u^{2267}+8\le5288u^{2267},
 \qquad T_O\le2u^{37700},
 \qquad T_On_W\le10576u^{39967}.
\]
The rank identity now yields
\[
 t\le128u^{686}+4+10576u^{39967}
 \le10708u^{39967}\le u^{40000}=x^{400},
\]
where the last inequality uses $10708\le2^{33}\le u^{33}$.
\end{proof}

\section*{Proofs for Chapter~\ref{ch:green}: Green's classification and its use}

\subsection*{Proof of \ref{B7.11}}\label{prf:B7.11}\phantomsection
\begin{proof}
Conjugation transports the invariant subspaces and their restrictions. Induct on $\dim V$ and average \ref{B7.10} over $h$. The terms for $W\subseteq A$ sum to $s(c_{g|_A})$. For $0<W\cap A<W\subsetneq V$, the $h$ for which $g+h$ stabilizes $W$ form an affine space if nonempty. Restriction gives the uniform distribution on a translate of $\operatorname{Hom}(W/(W\cap A),W\cap A)$, since maps in this latter space extend to $V/A\to A$. Induction gives zero.

For $W\cap A=0$, the image $B$ of $W$ in $V/A$ must be invariant under the induced $\bar g$. Each of its $v^{\dim A\dim B}$ lifts is stable with probability $v^{-\dim A\dim B}$, since stability prescribes $h|_B$. The restriction is conjugate to $\bar g|_B$, so these terms sum to $s(c_{\bar g})$ by \ref{B7.10}. The zero subspace contributes $s(1)=0$. By the block factorization and additivity, $s(c_{g|_A})+s(c_{\bar g})=s(c_g)$.
\end{proof}

\subsection*{Proof of \ref{B7.13}}\label{prf:B7.13}\phantomsection
\begin{proof}
Choose a finite extension of $F$ containing the eigenvalues of every element of $H$. Apply Green's Theorem~1 \cite{green-1955}, based on Brauer's character criterion \cite{brauer-1947}, to the restriction of $t$ and the integral symmetric polynomial $\sum_i X_i^h$. This gives the virtual character $R_h$.

In \ref{B7.10}, the sum over invariant subspaces of dimension $a$ is the parabolic induction of $\Psi_h$ in dimension $a$ tensored with the trivial character in dimension $n-a$. Induction on $n$ proves the claim for $\Psi_h$.
\end{proof}

\subsection*{Proof of \ref{B7.14}}\label{prf:B7.14}\phantomsection
\begin{proof}
Let $\mu$ and $b$ be the masses of $1$ and $\Psi_s$ divided in each degree by the order of the general linear group. Write $*$ for convolution over monic factorizations. Counting endomorphisms with an invariant subspace gives $s\mu=b*\mu$ by \ref{B7.10}. In dimension $a$ the number of subspaces is $|\GL_n(F)|/(|\GL_a(F)||\GL_{n-a}(F)|v^{a(n-a)})$. Its last denominator cancels the number of off-diagonal blocks.

Let $(T\mu)(P)=v^{\deg P}\mu(P)$. Counting an endomorphism with a vector by the cyclic subspace it generates gives $T\mu=\mathbf1*\mu$, since cyclic pairs with a given monic polynomial constitute one free orbit by the companion basis. The degree factor makes $T$ multiplicative for convolution. Additivity gives
\[
 T(s\mu)=s(\mathbf1*\mu)=s*\mu+\mathbf1*(s\mu).
\]
Substitution and cancellation of $\mu$ give $(Tb-b)*\mathbf1=s$. The identity $s=\Lambda_s*\mathbf1$, where $\Lambda_s(Q^a)=s(Q)$ for $a>0$ and is zero elsewhere, follows from unique factorization. Cancel $\mathbf1$ and divide by $v^{\deg P}-1$.
\end{proof}

\subsection*{Proof of \ref{B7.17}}\label{prf:B7.17}\phantomsection
\begin{proof}
By \ref{B7.15}, $I(\phi,R_k)=I(\phi,R_l)$ whenever $kv^i\equiv l\bmod M$, since Frobenius permutes the roots of each $P_\alpha$. Thus \ref{B7.16} shows that the primitive difference pairs to zero with every $R_l$. Fourier inversion in \ref{B7.15} gives zero mass on each invertible primary fibre.

On $c_g=Q^a$, a primitive class function is determined by its single value on the class $Q(g)=0$. For a function with this value zero, descend on $\dim A$ for $A=\ker Q(g)$. In the variations $g+h$ of \ref{B7.11}, $Q(g+h)$ vanishes on $A$ and agrees with $Q(g)$ modulo $A$. If its kernel equals $A$, successive preimages of $A$ give the same higher-power kernels. These variations are conjugate to $g$. Those with larger kernel have value zero by descent. The primitivity identity thus forces zero at $g$. The primitive function $\Psi_{\deg}$ has nonzero mass by \ref{B7.14}, so on this at most one-dimensional space the mass determines the function. This proves the expansion on the primary fibres, and \ref{B7.15} supplies vanishing on the others.
\end{proof}

\subsection*{Proof of \ref{B7.24}}\label{prf:B7.24}\phantomsection
\begin{proof}
For the empty list the group is trivial. Otherwise let \(P\) be the standard parabolic with the given ordered block sizes. Each cuspidal factor of its Levi has Whittaker multiplicity one for the restriction of the prescribed nondegenerate character. This is the finite-field Gelfand--Graev theorem, in the form recalled in \cite[Prop.~9.4]{zelevinsky-1981} and \cite[\S1, p.~76]{deriziotis-gotsis-1998}.

In the Mackey formula for restriction of the induced representation to \(U_n\), take the permutation which reverses the order of the blocks and preserves the order inside each block. The intersection of \(U_n\) with the conjugate parabolic consists of the upper unitriangular groups inside the permuted blocks. The character restricts to their product of Whittaker characters. Frobenius reciprocity therefore gives a contribution of one from this double coset. All contributions are dimensions. Decomposing \(I(C)\) into irreducibles now gives a constituent of positive Whittaker multiplicity.
\end{proof}

\subsection*{Proof of \ref{B7.25}}\label{prf:B7.25}\phantomsection
\begin{proof}
Write $C'=(\sigma,C'_0)$. By reciprocity and Mackey restriction, the pairing is a sum over distributions of blocks into the two factors of the Levi for $C'$. Cuspidality kills every split block. If several whole blocks go to the first factor, their induced product pairs to zero with $\sigma$ by cuspidality and reciprocity. For one block, orthogonality retains precisely the positions $i$ with $C_i=\sigma$. The sum becomes
\[
 \sum_{i:\,C_i=\sigma}\langle I(C\setminus i),I(C'_0)\rangle .
\]
The empty lists pair to one. Induction gives the two assertions.
\end{proof}

\subsection*{Proof of \ref{B7.37}}\label{prf:B7.37}\phantomsection
\begin{proof}
Choose a root $x_P$ of each $P\in\mathcal P_d$. By \ref{B7.34} and \ref{B7.35}, $\tau_P=\tau_{\theta_{x_P}}$ is cuspidal of degree $\Delta_d$. Every characteristic root $z$ of elliptic $y$ lies in $\iota(K^\times)$ by \ref{B7.32} and \ref{B7.35}. The regular formula and \ref{B7.33} give \eqref{eq:green-label} at $x_P$. For $x=x_P^{p^a}$, \ref{B7.36} gives
\[
 \sum_{i<d}\Theta_d(x,z^{p^i})
 =\sum_{i<d}\Theta_d(x_P,z^{p^{i+a}})
 =\sum_{i<d}\Theta_d(x_P,z^{p^i}).
\]
Equal characters give conjugate parameters, hence conjugate roots by \ref{B7.33} and \ref{B7.34}. Their minimal polynomials agree. Conversely, a cuspidal is $\tau_{\theta_x}$ with $\theta_x$ regular. The degree $e$ of $x$ divides $d$. If $e<d$, then $x^{p^e}=x$ contradicts regularity. Thus its minimal polynomial is in $\mathcal P_d$ and its chosen root is in the same orbit. Extend the family arbitrarily off $\mathcal P_d$.
\end{proof}

\subsection*{Proof of \ref{B7.47a}}\label{prf:B7.47a}\phantomsection
\begin{proof}
Write $W=\{s,\ldots,s+k-1\}$ and keep the character given by the two chains defining $\phi_J$. At stage $i$, $0\le i\le s$, insert $W$ after the first $i$ indices of the other chain. Allow the strictly upper entries in that order except the entries from $W$ to the next index of the other chain. Stage $s$ gives $H_J$. Stage zero, with $W$ moved to the front, gives the subgroup and character of $e'_J$.

For a step from $i$ to $i-1$ put $a=i-1$ and let $b$ be the successor of $a$ in the other chain. Thus $b=i$ for $i<s$ and $b=s+k$ for $i=s$. The two patterns differ by replacing $(a,w)$ by $(w,b)$ for all $w\in W$. The union is a group $H$ under matrix multiplication. The two subgroups $H_1,H_2$ each extend their intersection by $p^k$ choices, and their characters agree on the intersection. The commutator of the two added rows of root coordinates has $(a,b)$ entry $\sum_{w\in W}x_wy_w$. The pairing $\eta(\sum_w x_wy_w)$ is nondegenerate since a nonzero coordinate can be paired with a variable taking all values in $\F_p$.

For the Mackey self-pairing on the first subgroup, any $h\in H$ outside it has some $h_{wb}\ne0$. Conjugating $1+tE_{aw}$ by $h$ in the form $h^{-1}(1+tE_{aw})h$ stays in the first subgroup. The new entries can occur only at predecessors of $a$ and successors of $w$ in the union pattern, with endpoints included. Among the two chains' consecutive entries this changes only $(a,b)$, by $t h_{wb}$. Choosing $t$ so that the character changes makes that double coset contribute zero. For the second subgroup use $1+tE_{wb}$ and a nonzero $(h^{-1})_{aw}$. Its conjugate stays in the second pattern and changes the character sum by $t(h^{-1})_{aw}$. Thus each induced character on $H$ has norm one. Their cross-pairing has a contribution of one from the identity double coset, so the induced characters agree. Inducing to $G_n$ and using reciprocity proves equality of the two traces for characters of $G_n$. They span the class functions, so the successive equalities prove the lemma.
\end{proof}

\subsection*{Proof of \ref{B7.47b}}\label{prf:B7.47b}\phantomsection
\begin{proof}
By linearity take $\chi$ to be a character. For $1\le j\le n$ let $P_j$ be the subgroup of $G_j$ fixing the last row covector. Let $T$ on $P_{k+1}$ be the character obtained from $\chi|_{P_n}$ by successively taking the nontrivial last-column character spaces down to $P_{k+1}$. At the step from $P_{j+1}$ to $P_j$, use $v\mapsto\eta(v_j)$ on the column $v\in\F_p^j$. Thus $T$ is the tail average over upper unitriangular matrices whose left upper $(k+1)$-block is the identity. The weights use the adjacent entries from that block to the end. Grouping the average for $e'_J$ by its left upper $k$-block gives
\[
 \chi(e'_J)=\operatorname{wh}_\eta\bigl(h\mapsto T(\operatorname{diag}(h,1))\bigr).
\]
The allowed entries in each fibre are exactly the tail entries, and the character factors as the product of the upper $U_k$ character and the tail character.

In $P_{j+1}=G_j\ltimes\F_p^j$ the dual of the column group has two $G_j$-orbits. The zero orbit gives the column invariants as a $G_j$-representation. The nonzero orbit is induced from $P_j\ltimes\F_p^j$, with the chosen character on the column and its multiplicity representation on $P_j$. Iterate this decomposition on the latter representation. For $T$ the summands are indexed by the rank $i$, $0\le i\le k$, at which the zero character is taken. In $P_{k+1}$ the summand is induced from the upper block subgroup with diagonal blocks $G_i$ and $U_{k+1-i}$. It uses $D_i^\eta\chi$ on $G_i$ and $\xi_\eta$ on $U_{k+1-i}$. It is trivial on the upper off-diagonal block. The selected columns impose the simple-root character within the tail, and the zero step imposes invariance on the remaining entries to the right of the first $i$ coordinates. Together these averages are precisely the unipotent restriction and Whittaker average defining $D_i^\eta$.

The inducing subgroup contains the last-column group of $P_{k+1}$. Restriction to $G_k$ therefore gives a single Mackey term, induced from the corresponding subgroup with diagonal blocks $G_i,U_{k-i}$. Induction in stages identifies its character with the Harish--Chandra product of $D_i^\eta\chi$ and $\Gamma_{\eta,k-i}=\operatorname{Ind}_{U_{k-i}}^{G_{k-i}}\xi_\eta$. Whittaker multiplicities multiply on that product, by the opposite-cell argument of \ref{B7.48}. The coefficient is $p^{k-i}-p^{k-i-1}$ by \ref{B7.26} for $i<k$. For $i=k$ the second factor is the trivial character in rank zero, of multiplicity one. Apply \ref{B7.47a}.
\end{proof}

\subsection*{Proof of \ref{B7.48}}\label{prf:B7.48}\phantomsection
\begin{proof}
Mackey restriction distributes the inducing blocks between the two Levi factors. Cuspidality kills every distribution splitting a block, and each other term is a product induced separately in the two factors. For Whittaker multiplicity of a parabolic induction, the opposite cell contributes the product of the multiplicities as in \ref{B7.24}. Every other cell has a simple root subgroup in the intersection with the conjugated unipotent radical, where the inducing representation is trivial and $\xi_\eta$ is nontrivial. Each surviving distribution therefore contributes one.
\end{proof}

\subsection*{Proof of \ref{B7.50}}\label{prf:B7.50}\phantomsection
\begin{proof}
The host fields of $D$ give distinct primes $q,d>4tr+1$, where $t$ is the host rank, and $2^{e+2}\mid Q_S$ with $e\ge2$. The normalized types give $p_j\equiv1\bmod Q_Sqd$, hence $16,q,d\mid p_j-1$ and $p_j\ge5$.

Apply \ref{B7.43} with the Teichmüller lift $t_j:\overline{\F}_{p_j}\to\C$ and the chosen generators $u_{j,h}$ of $\mu_{p_j^h-1}(\overline{\F}_{p_j})$, $h\ge1$. Their compatibility reads $u_{j,ha}^{(p_j^{ha}-1)/(p_j^h-1)}=u_{j,h}$ for every $a\ge1$. Here $r=k+2s+1$ with $k\ge1$ and $s=\lfloor m/2\rfloor-1$. The arithmetic projector is $E_{\sigma_*f}$ on admissible labels and zero elsewhere. The substitution preserves degrees and squarefreeness. By \ref{B7.42} these projectors sum to the identity, and by \ref{B7.44} the weights are nonnegative with positive denominator. The sum of the numerators is that denominator.
\end{proof}

\subsection*{Proof of \ref{B7.50a}}\label{prf:B7.50a}\phantomsection
\begin{proof}
Let $I_j$ count monic irreducibles of degree $j$ with nonzero constant term. They divide $X^{u^j-1}-1$, so $jI_j\le u^j-1$. The mass of a type with $m_j$ parts of size $j$ is
\[
 \prod_{j=1}^N\binom{I_j}{m_j}(u^j-1)^{-m_j}
 \le\prod_{j=1}^N(j^{m_j}m_j!)^{-1}.
\]
The right side is the permutation proportion because a cycle listing counts each permutation $\prod_j j^{m_j}m_j!$ times. Sum over the desired types. Division by $h$ injects the squarefree multiples of $h$ into $\mathcal S_{N-\deg h}(K)$. The identity $w(f)=w(h)w(f/h)$ gives the divisor bound.

Each squarefree admissible $f$ has weight at least $u^{-N}$. Of the $u^N$ monic polynomials, $u^{N-1}$ have zero constant term. For $N\ge2$ at most
\[
 \sum_{1\le j\le N/2}I_j u^{N-2j}
 \le (u-1)u^{N-2}+\tfrac12\sum_{2\le j\le N/2}u^{N-j}
 \le u^{N-1}
\]
are divisible by the square of an irreducible with nonzero constant term. In degree one there are no such squares. Multiply the lower count by $u^{-N}$.
\end{proof}

\subsection*{Proof of \ref{B7.50b}}\label{prf:B7.50b}\phantomsection
\begin{proof}
By \ref{B7.50a} it suffices to bound the probability of an invariant $b$-set for a uniform permutation. Let $X$ count such sets and $Y$ count those with a fixed point both inside and outside. A specified $b$-set is invariant with probability $b!(N-b)!/N!$, so $\mathbb E X=1$. Conditional on invariance, the restrictions are independent uniform permutations. On $s\ge2$ letters, the first two inclusion-exclusion terms bound the probability of a fixed point below by $1-\binom{s}{2}/(s(s-1))=1/2$. On one letter the probability is one. Thus $\mathbb E Y\ge1/4$. If $Y>0$, exchanging the two fixed points gives a second invariant set, so $2\mathbf1_{X>0}+Y\le2X$. The inequality holds for $Y=0$ as well. Taking expectations gives $7/8$.
\end{proof}

\subsection*{Proof of \ref{B7.50c}}\label{prf:B7.50c}\phantomsection
\begin{proof}
For each triple $(A,T,h)$ in the table start at $a=A$. While $a<T$, put $c=a+\lfloor a/10\rfloor$ and $s=\lfloor c/2\rfloor$, compute
\[
 L(a,c,h)=
 \sum_{\substack{s\le j<s+2^{\lfloor\log_2 a\rfloor}\\
 c<2j,\ j+3\le a,\ h<j<(h+1)^2\\
 \gcd(j,h!)=1}}
 \left\lfloor\frac{10^9}{j}\right\rfloor,
\]
and replace $a$ by $c+1$. The iteration counts and minima are
\[
\begin{array}{r r r|r r}
 A&T&h&\text{iterations}&\min L\\ \hline
100&196&13&7&108738643\\
196&1024&31&18&86865767\\
1024&20000&141&32&61695338
\end{array}
\]
A composite $j<(h+1)^2$ has a prime divisor at most $h$, so every summand comes from a prime in $\mathcal J_N$ for all $a\le N\le c$. The intervals cover $100\le N<20000$, and each minimum gives $H_N\ge1/20\ge1/(5\log N)$.

For $N\ge20000$ put $n=\lfloor N/2\rfloor$, $y=\sqrt{2n}$, $L=\log(2n)$, $l=\log2$, and $k=\pi(2n)-\pi(n)$. With $C_n=\binom{2n}{n}$,
\[
 4^n<nC_n,\qquad C_n\le(2n)^{\pi(y)+k}\,4^{\lfloor2n/3\rfloor}.
\]
The first bound starts at $n=4$ and propagates since $(n+1)C_{n+1}/(nC_n)>4$. Each term $\lfloor2n/t^j\rfloor-2\lfloor n/t^j\rfloor$ in the valuation at prime $t$ is zero or one. Thus primes $t\le y$ contribute at most $2n$ each. Above $y$ the exponent is at most one and vanishes on $(2n/3,n]$. The product of primes up to $m$ is at most $4^m$ by induction. For $m=2s+1$, the primes in $(s+1,2s+1]$ divide $\binom{2s+1}{s}\le2^{2s}$, and the smaller ones contribute at most $4^{s+1}$. For even $m\ge4$ the product is unchanged from $m-1$. This bounds the primes up to $2n/3$, and the final $k$ primes contribute at most $(2n)^k$.

Counting classes modulo $6$ gives $\pi(y)\le y/3+3$. Taking logs yields $y^2l/3<(y/3+4+k)L$. Since $y\ge141$, $\log(y/128)\le y/128-1$ gives $L\le14l+y/64-2$. With $0.69<l<0.70$ from the series for $\log2$,
\[
 (5l/3-1)y^2-5(y/3+7)(14l+y/64-2)-1
 \ge 3y^2/25-14y-275>0.
\]
Thus $5(k-3)L\ge2n+1$. Each prime in $(n,2n]$ exceeds $N/2$, and deleting the possible values $N-2,N-1,N$ leaves at least $k-3$ primes in $\mathcal J_N$. Consequently $5|\mathcal J_N|\log N\ge N$. Each reciprocal is at least $1/N$.
\end{proof}

\subsection*{Proof of \ref{B7.50e}}\label{prf:B7.50e}\phantomsection
\begin{proof}
Write $n_j=u^j-1$ and $a_\ell(n)=\ell^{v_\ell(n)}$. The prime-to-$\ell$ projection in the group of order $n_j$ has fibres of size $a_\ell(n_j)\le u^{j/4}$ for $\ell\notin B(u,N)$ and $1\le j\le N$.

If two roots of the same degree-$j$ factor collide, their common prime-to-$\ell$ part has Frobenius orbit of length $e$ properly dividing $j$. Each root of the factor solves $X^{(u^e-1)a_\ell(n_j)}=1$. For each $e\le j/2$ there are at most $u^{3j/4}$ solutions. There are at most $j$ choices of $e$. Since each factor contributes $j$ roots, there are at most $u^{3j/4}$ factors. By \ref{B7.50a}, each costs at most $1/n_j$. Such a collision requires $j\ge2$ and $\ell\mid n_j$.

For distinct factors of degrees $j\ge i$, fix the degree-$i$ factor and one root $y_0$. By Frobenius the other factor has a root in the fibre above $(y_0)_{\ell'}$ in the group of order $n_j$. This gives at most $u^{j/4}$ choices of that factor. There are at most $n_i$ choices of the first factor and each pair costs at most $1/(n_i n_j)$, with zero mass if $i+j>N$. The contribution for fixed $i,j,\ell$ is at most $u^{j/4}/n_j$. Here $\ell\mid n_i n_j$.

Put $z=u^{-1/4}\le1/8$. The number of prime divisors of $n_j$ is at most $(3/2)j\log u$, and that of $n_i n_j$ for $i\le j$ is at most $3j\log u$. Since $1/n_j\le2u^{-j}$, the two collision bounds sum to at most
\[
 3\log u\sum_{j\ge2}jz^j+
 6\log u\sum_{j\ge1}j^2z^{3j}
 \le (16z^2+30z^3)\log u
 \le23\log u/\sqrt u.
\]
For the middle inequality use $j\le2^j$ and sum the geometric series.

For a collision with $c$, the quotient $x/c$ has prime-power order dividing $n_j$, or order one. In either case the minimal polynomial divides $X^{a_s(n_j)}-c^{a_s(n_j)}$ for a prime $s\mid n_j$, since $16\mid u-1$. If $b_j$ counts the degree-$j$ factors of these binomials, counting roots gives $jb_j\le\sum_{s\mid n_j,\ s\ {\rm prime}}a_s(n_j)$. The mass is at most $\sum_{j=1}^N b_j/n_j$ by \ref{B7.50a}.

Set $A=u-1$ and $t=\min(q,d)$. Let $U\ge16$ and $V\ge t$ be the $2$-part and the $t$-part, respectively. Let $Z\ge\max(16,t)$ be the part at $\max(q,d)$. Let $R$ be the product of the remaining parts. Their sum is at most $R$ by $a+b\le ab$ for integers $a,b\ge2$. Each of $U,V,Z,R$ is at most $A/(16t)$ by the product of the other three. Thus $b_1/A\le1/(4t)$.

For $j\ge2$ and prime $s\mid A$, one has $v_s(u^j-1)=v_s(A)+v_s(j)$. If $s\mid h$ and $s\nmid a$, binomial expansion gives $v_s((1+h)^a-1)=v_s(h)$. Raising to $s$ adds one to the valuation for odd $s$, and for $s=2$ when $4\mid h$, as holds here. Thus the product of $a_s(n_j)$ over $s\mid A$ divides $jA$, and over $s\nmid A$ it divides $n_j/A$. Bounding each sum by the corresponding product gives
\[
 \frac{b_j}{n_j}\le\frac A{n_j}+\frac1{jA}
 \le u^{1-j}+\frac1{jA}.
\]
Sum from $j=2$, using $\sum_{j\ge2}u^{1-j}=1/A$ and $\sum_{j=2}^N1/j\le\log N$, and add the degree-one term.
\end{proof}

\subsection*{Proof of \ref{B7.51}}\label{prf:B7.51}\phantomsection
\begin{proof}
Fix $j$ and put $p=p_j$. Since $16\mid Q_S$ and $p\equiv1\pmod{Q_Sqd}$, each of $16,q,d$ divides $p-1$. Since $r\ge100$ and $t\ge1$, the host bound $q,d>4tr+1$ gives $q,d>16$.

Let $P$ be the $U$-law at $\eta_j$ and $Q_v$ the law of the idempotent at $v$. These are probabilities. The lower $r$ coefficients recover each monic polynomial, and admissible roots are nonzero of finite order. Hence the arithmetic tests give the same prime-to-$\ell$ parts and excluded primes as in \ref{B7.50e}, with one test for each twist and one for its inverse. The divisor tests use the same recovered polynomials.

For the product $P^{\otimes\nu}$, use \ref{B7.50y} at each place. The union of failures of squarefreeness, root separation and the two twist tests has weight at most
\[
 \frac{2\nu}{p}+\frac{23\nu\log p}{\sqrt p}
 +2\nu m\left(\frac1{4\min(q,d)}+\frac{1+\log r}{p-1}\right).
\]
On squarefree admissible tuples, failure for fixed $b$ requires a degree-$b$ divisor at every place. It has probability at most $(7/8)^\nu$. Sum over $1\le b\le\lfloor r/2\rfloor$. Failure of the long-factor test requires every entry to miss $\mathcal J_r$, with probability at most $(1-1/(5\log r))^\nu\le e^{-\nu/(5\log r)}$ on those tuples.

Let $E$ be this union of failures on admissible tuples. By \ref{B7.50}, $Q_v\ge(1-p^{-1})P$ pointwise. Multiply on tuples of $E^c$ and use total mass one to get
\[
 \left(\bigotimes_v Q_v\right)(E)
 \le1-(1-p^{-1})^\nu(1-P^{\otimes\nu}(E))
 \le P^{\otimes\nu}(E)+\nu/p .
\]
This gives $W$ since the weights vanish off admissible tuples and $\lfloor r/2\rfloor\le r/2$.
\end{proof}

\subsection*{Proof of \ref{B7.52}}\label{prf:B7.52}\phantomsection
\begin{proof}
The full cylinder gives $D_{\rm all}\ge(1-\epsilon_C)D_*>0$. Let $S_j$ be the finite set of tuples of all-row indices bad at $j$. Then $0\le D_{\rm all}-D_{\rm rows}\le\sum_j N_j(S_j)$ by the index condition. The finite-cylinder identity gives $D_j(S_j)=D_*P_j(S_j)$.

By \ref{def:p7c1-compatible-rows}, all places above $p_j$ use $e_J(\eta_j)$ for the same nontrivial character $\eta_j$. The two distinguished primes of the host split completely in $E$ and $M_{\rm ray}$ and are congruent to $1$ modulo $Q_Sqd$, as required in \ref{B7.51}. For $x^\alpha\ge3/2$ the size bound gives $p_j\ge4096$. By \ref{B7.51} and relative counting,
\[
 0\le\delta_{x,i}\le
 \frac{1+\epsilon_C}{1-\epsilon_C}\sum_{j=0}^1 W(\nu,m,r,p_j,q,d).
\]
Eventually $\epsilon_C\le1/2$ uniformly. At points with indices, $C_s\ge2$. By the size bounds and the decrease of $(\log p)/\sqrt p$ for $p\ge e^2$, the first three terms of $W$ total $O(x^{3C_s+3}e^{-x^{\alpha/6}})$. The last two are $O(x^{C_s}\exp(-c'x^2/\log x))$ for some $c'>0$, since $\nu/\log r\ge x^2/(C_s\log x)$. Both are $o(e^{-x^{\alpha/7}})$ uniformly in $i$ for $0<\alpha\le1$.
\end{proof}

\subsection*{Proof of \ref{B7.53}}\label{prf:B7.53}\phantomsection
\begin{proof}
Reduce the $\varepsilon$-threshold to at most $1$ and keep the size threshold $\mathcal K$ for each eligible $\varepsilon$. For a family and height exponent $C_H$ meeting the size thresholds $\mathcal K$ of \ref{B8.39}, write $x_1$ for the family's initial cutoff. As point threshold take the maximum of $x_c$, the selection threshold, $x_1+1$, $4$, $(3/2)^{1/\alpha}$ and the numerical threshold of \ref{B7.52}. Beyond it, the selected normalized types give \ref{B7.50} for the same rows.

The sizes give $n\ge x^{2+8\varepsilon-\varepsilon/100}$ and $x^{9/2+10\varepsilon-\varepsilon/100}\le k\le x^{9/2+10\varepsilon+\varepsilon/100}$, with $m\le f=2n\le k$ and $k\le r\le2k$. As $\varepsilon<1/52$ in the size data, $x\ge4$ gives $x^2\le\nu\le x^6$, $m\le x^6$ and $100\le r\le x^6$. From $(kb)^8\le p_j$ and $\alpha\le\varepsilon$ we have $p_j\ge e^{8x^\alpha}$. The lower bounds on $\log q,\log d$ are at least $x^{\alpha/6}$. Apply \ref{B7.52} with $C_s=6$.
\end{proof}

\subsection*{Proof of \ref{B7.54}}\label{prf:B7.54}\phantomsection
\begin{proof}
Put $\mathcal A_G=\mathcal A\cap\mathcal E$. Minkowski's rational bound says that the order of each finite subgroup of $\GL_n(\Q)$ divides $M(n)$ of \ref{B8.02} \cite[Thm.~1(i)]{serre-2007}. For every sufficiently large member and every word $u$, the design gives $D_{\Gamma,u}\mid M(2fN(m))$, where $D_{\Gamma,u}$ is the stabilizer denominator of \ref{B4.03}. The primes in $M(n)$ are at most $n+1$, and $M(n)\mid(2n)!$ since $\lfloor n/(p^a(p-1))\rfloor\le\lfloor2n/p^{a+1}\rfloor$ for primes $p$ and $a\ge0$. Thus $\log M(n)\le2n\log(2n)$. At $n=2fN(m)$ with $fN(m)\ge2$, this gives $\log D_{\Gamma,u}\le12fN(m)\log(fN(m))$ uniformly in $u$. The prime bound supplies smoothness.

On $\mathcal W_1(\mathcal A_G)$, if $C_{Tu}\ge c_T(\varepsilon,C_{YL})$, the level-prime hypothesis of \ref{B8.02} gives $b^{m/2}\le T_u$ and $\log T_u\le C_{Tu}mL$ for every word of every sufficiently large member, with $L=\log b=x^\varepsilon$. On this working class the common hypotheses of \ref{B8.02} supply the class-number row of \ref{B8.54}. The aligned selection supplies the symmetric columns of item~5 of \ref{B8.09}. Intersect their thresholds by \ref{B8.39}. Lemma~\ref{B5.32b}, applied to the word windows of \ref{B5.11}, gives $1/3\le f_u\Gamma^2/D^\star\le3$ uniformly in the word. This establishes the rank condition used in choosing the design.
\leanref{Formalism.\allowbreak{}MainNotions.\allowbreak{}LinkB.\allowbreak{}rank\_\allowbreak{}formula\_\allowbreak{}of\_\allowbreak{}part\_\allowbreak{}four}{Formalism.\allowbreak{}SharpRestriction.\allowbreak{}GlobalSchurBounds.\allowbreak{}TopDensityNGeSixteen}{95--103}
\leanref{Formalism.\allowbreak{}MainNotions.\allowbreak{}proviso\_\allowbreak{}e}{Formalism.\allowbreak{}SharpRestriction.\allowbreak{}GlobalSchurBounds.\allowbreak{}MainSup}{196--245}
\leanref{Formalism.\allowbreak{}MainNotions.\allowbreak{}StandardLattice.\allowbreak{}HarishChandraReduction.\allowbreak{}chosenProviso\_\allowbreak{}4f}{Formalism.\allowbreak{}SharpRestriction.\allowbreak{}GlobalSchurBounds.\allowbreak{}TheoremThreeFamilyOfFormalDegreeInput}{64--79}

For every sufficiently large member the design gives a scale $Y\in[x^{1/J},x]$ for the six sign-character zero-free boxes and the real host-character boxes of \eqref{eq:p8c5-boxes}. It gives a scale in the same interval for the prime and conductor conditions \eqref{eq:p8c5-qd}. At $x$ itself the primitive representatives of the $q$- and $d$-sign characters and their product satisfy \eqref{eq:p5c4-qd-input} with a common $C_5>0$. With the uniform zero estimates and explicit formula of \ref{B8.45}--\ref{B8.46}, the closed form of \ref{B8.47} gives a common $c_6>0$ such that every primitive host character satisfies $|g_\chi(y)|\le y e^{-c_6x^{\varepsilon/8}}$ for $y\ge e^{x^{\varepsilon/3}}/3$. These two character estimates and the rank bound above complete the design conditions at $3,C_H$. The remaining conditions from \ref{B8.02} are the lower bound on the thick set, the top density and weighted bad-orbit bound with a common $C_b>0$, and the lower bound $3\Gamma^2x^{C_H}$ on the number of words. All hold for every member beyond a common point threshold.
\leanref{Formalism.\allowbreak{}MainNotions.\allowbreak{}designExists\_\allowbreak{}of\_\allowbreak{}choices}{Formalism.\allowbreak{}SharpRestriction.\allowbreak{}DesignExistsOfChoices}{54--68}
\leanref{GlobalTheory.\allowbreak{}Top.\allowbreak{}designExists\_\allowbreak{}of\_\allowbreak{}tuSize\_\allowbreak{}aux4n\_\allowbreak{}litThin}{Development.\allowbreak{}AutomorphicArithmetic.\allowbreak{}ThetaCorrespondence.\allowbreak{}ThetaPairLifts.\allowbreak{}ThirdPlaceSeries.\allowbreak{}LinkEx}{15--37}

On $\mathcal D_3(\mathcal A)$ the row conditions of \ref{B8.02} and Green give the closed defect bound of \ref{B7.53}. The top density gives $\delta_1n/2\le n_T\le2\delta_1n$. Applied to the design's bad-orbit bound, the Mertens estimate bounds the weighted count of bad top positions in \ref{B5.21} by $n_Tx^{-2+\varepsilon}$. The selected second layer gives group compatibility. On the same class, \ref{B5.20a} gives the closed local lattice and gain estimates, with a common constant and an error scale $b^{\delta(x)}$ for some $\delta(x)\to0$.
\leanref{GlobalTheory.\allowbreak{}Top.\allowbreak{}closed\_\allowbreak{}linkG\_\allowbreak{}of\_\allowbreak{}duty}{Development.\allowbreak{}AutomorphicArithmetic.\allowbreak{}ThetaCorrespondence.\allowbreak{}ThetaPairLifts.\allowbreak{}ThirdPlaceSeries.\allowbreak{}LinkG}{13--24}

For the rank identity at a twisted place, take the maximum of the selection threshold and $x_1+1$, where $x_1$ is the family's initial cutoff. At each eligible member beyond this threshold, take the normalized local compatibility data and their Hecke ties. Fix $1\le j\le m$ and a place $w$ with nonreference tuple. Write $K_{\rm coeff}$ for the row coefficient field. Apply the normalized local comparison of the selected member to $W=M^{\mathrm{tw}}_{w,j}$. It assigns to each of its indices $i\in I_{j,w}$ a maximal ideal of the integral Hecke algebra on the scalar extension of $W$. The comparison is injective on this fibre and covers every maximal ideal at a nonreference tuple. For the member construction the map and its inverse are those of Lemma~\ref{B8.07t}. For a selected realization in other coordinates, its comparison identifies the local coefficient rings and gives an injective semilinear map from $\operatorname{Frac}(O_i)\otimes_{O_i}M_i$ onto the whole summand for the assigned maximal ideal. Thus an $O_i$-basis of $M_i$ gives a basis of that summand after scalar extension. The summands form a direct sum by the joint primary decomposition in Lemma~\ref{B8.07t}. The dimension of $W$ is preserved by extension of its coefficient field. Summing gives
\[
 \sum_{i\in I_{j,w}}\operatorname{rank}_{O_i}M_i
 =\dim_{K_{\rm coeff}} M^{\mathrm{tw}}_{w,j},
\]
where $M^{\mathrm{tw}}_{w,j}$ is the twisted layer space.
\leanref{GlobalTheory.\allowbreak{}Top.\allowbreak{}closed\_\allowbreak{}linkEta\_\allowbreak{}of\_\allowbreak{}duty}{Development.\allowbreak{}AutomorphicArithmetic.\allowbreak{}ThetaCorrespondence.\allowbreak{}ThetaPairLifts.\allowbreak{}ThirdPlaceSeries.\allowbreak{}LinkEta}{13--23}
\leanref{Formalism.\allowbreak{}L2.\allowbreak{}locsCoverTw\_\allowbreak{}of\_\allowbreak{}tie}{Formalism.\allowbreak{}LocalToGlobal.\allowbreak{}NormalizedHeckeFactors.\allowbreak{}CoverThm}{120--122}

Under the level-prime bounds, the spectral hypotheses of \ref{B8.02}, the signed nongeneric and equal-word bounds of \ref{B8.10d}, and the counts of \ref{B8.52} give the following closed estimate for members satisfying $1/3\le f_u\Gamma^2/D^\star\le3$ for every word. On the same selected realization, the real-place and row ties identify the layer projectors along with the volume and degree normalization. The diagonal, off-diagonal, loop and distinct-word estimates of \ref{B8.10}--\ref{B8.10b} can then be combined with those counts. After a common point threshold, for each member they give errors $\eta_0,\eta_1,e_0$ for the distinct-word bound, total dimension and squared diagonal error, respectively, with
\[
 e_0\ge0,\qquad \eta_0,\eta_1,\sqrt{e_0}\le e^{-x^{\varepsilon/8}}.
\]
\leanref{GlobalTheory.\allowbreak{}Top.\allowbreak{}link\_\allowbreak{}alt\_\allowbreak{}alt\_\allowbreak{}bis}{Development.\allowbreak{}AutomorphicArithmetic.\allowbreak{}ThetaCorrespondence.\allowbreak{}ThetaPairLifts.\allowbreak{}ThirdPlaceSeries.\allowbreak{}LinkAltAltBis}{19--48}
\leanref{Formalism.\allowbreak{}L2.\allowbreak{}SG.\allowbreak{}ThmSOwn.\allowbreak{}theoremS\_\allowbreak{}transport\_\allowbreak{}link\_\allowbreak{}zero}{Formalism.\allowbreak{}LocalToGlobal.\allowbreak{}SiegelLatticeTheory.\allowbreak{}FlagStabilizerTensors.\allowbreak{}LinkSPos}{148--162}

For every sufficiently large selected member and each $0\le k<m$, the level projection on row vectors at the base tuple $\mu_0$ is surjective by \ref{B8.02}. Apply this to the selected primitive rows with their level-three conditions. The row and division identifications of \ref{B8.08} give $\dim_{\C} Z_j(\mu_0)$ for the image dimension and $D_0$ for the target dimension. Thus $\dim_{\C} Z_j(\mu_0)=D_0$ in every lift layer, with closure on the design class.
\leanref{GlobalTheory.\allowbreak{}Top.\allowbreak{}link\_\allowbreak{}alt\_\allowbreak{}alt\_\allowbreak{}alt}{Development.\allowbreak{}AutomorphicArithmetic.\allowbreak{}ThetaCorrespondence.\allowbreak{}ThetaPairLifts.\allowbreak{}ThirdPlaceSeries.\allowbreak{}Link}{20--54}

Under the bounds on $T_u$ and $D_{\Gamma,u}$, the selected row estimates of \ref{B8.38} and their transfer in \ref{B5.10} give a closed conditional estimate for the period maps, with division exponent $C_{\rm abs}=2$. Its hypotheses are the twisted-place rank identity above and an eventual bound $\eta_j(\mathcal L)\le D_0 C x^\varepsilon(B_0+N_c)$ with one $C$ for all members and lift layers. Its conclusion is the period-map containment and the height and slope estimates of \ref{B5.10}, with a common constant. Under the level-prime bounds and the design conditions at $3,C_H$, the real-place second-layer identifications and the volume and mass hypotheses of \ref{B8.02} give fullness of $N'_u$ and positive rank for the input lattice of every word, as in \ref{B8.37}.
\leanref{Formalism.\allowbreak{}L2.\allowbreak{}DIV.\allowbreak{}closed\_\allowbreak{}linkDIVt\_\allowbreak{}of\_\allowbreak{}closed\_\allowbreak{}bridge\_\allowbreak{}kernel}{Development.\allowbreak{}AutomorphicArithmetic.\allowbreak{}ThetaCorrespondence.\allowbreak{}ThetaPairLifts.\allowbreak{}ThirdPlaceSeries.\allowbreak{}DIVRecBridgeSel}{197--205}

By the fifth family condition of \ref{B8.02}, every family in $\mathcal A$ has a threshold above which the weighted sum $R_W(i)$ of \ref{B5.22b} is at most $x^{-9}$ for every member $i\in I_x$. For the constructed class, that duty follows from \ref{B8.02a}. By \ref{B5.22b}, every sufficiently large window point has a common weight pair with a realizing model, and each model for the pair satisfies $R_W\le x^{-9}$. The top density, the weighted bad-orbit bound and the local comparison of \ref{B5.20a} now give the loss bound of \ref{B5.21} for each such model. The common-weight bound \ref{B5.22} gives $e_L=0$. With the remaining design conditions at $3,C_H$, Theorem~\ref{B5.33} gives a closed implication from the gain, frame, defect, division and slope estimates to a common chain realisation at every sufficiently large window point.
\leanref{Formalism.\allowbreak{}L2.\allowbreak{}GlobalTheoryTheory.\allowbreak{}Loss7.\allowbreak{}loss\_\allowbreak{}requirements\_\allowbreak{}of\_\allowbreak{}small\_\allowbreak{}shares}{Formalism.\allowbreak{}LocalToGlobal.\allowbreak{}GlobalTheory.\allowbreak{}LossRequirementsOfSmallShares}{337--340}
\leanref{GlobalTheory.\allowbreak{}Top.\allowbreak{}link\_\allowbreak{}gl}{Development.\allowbreak{}AbelianVarietyTheory.\allowbreak{}Proofs.\allowbreak{}GlobalInterpolation.\allowbreak{}LinkGl}{14--31}

Combine the closed assertions by \ref{B8.39}. Take the minimum of their positive $\varepsilon$-thresholds and the design threshold. For each smaller positive $\varepsilon$, take the componentwise maximum $\mathcal K$ of the size thresholds before choosing a family. The design gives $F_m\in\mathcal W_1(\mathcal A_G)$ and $C_H\ge0$ meeting the thresholds $\mathcal K$ of \ref{B8.39}, the design conditions at $3,C_H$, and the bounds on $T_u$ and $D_{\Gamma,u}$. This family belongs to $\mathcal D_3(\mathcal A)$ and $\mathcal E$, and meets all the closed thresholds.

Since $F_m\in\mathcal E$, the cohomological lengths of \ref{B6.01} are eventually finite at every member, and their logarithmically weighted sum over every finite set of primes is at most $\sqrt b$. Apply \ref{B8.40a} to the local length bound of \ref{B8.40} with $\dim_{\C}Z_j(\mu_0)=D_0$. This gives \eqref{eq:B8-divinN} in every lift layer. The twisted-place rank identity supplies the other hypothesis of the conditional division estimate, so the period maps satisfy \ref{B5.10} with division exponent $2$. Increase the point threshold until $C_{Tu}\le x^{\varepsilon/100}$. On the slope scale $X''$ of \ref{B5.06}, the upper bound $\log T_u\le C_{Tu}mL$ and the family size inequalities now give $\log T_u\le2X''$ for every word. Likewise $\log D_{\Gamma,u}\le C fN(m)\log(fN(m))$ and the family size inequalities give $\log D_{\Gamma,u}\le10|C|X''$.  The division estimate supplies the zero-layer condition, and full rank gives finite generation. Together with the slope-rank condition, \ref{B5.06} and \ref{B5.12} give the lattice slope comparison.
\leanref{Formalism.\allowbreak{}MainNotions.\allowbreak{}muTu\_\allowbreak{}of\_\allowbreak{}tuSize}{Formalism.\allowbreak{}SharpRestriction.\allowbreak{}MuTuOfTuSize}{27--40}
\leanref{Formalism.\allowbreak{}MainNotions.\allowbreak{}J.\allowbreak{}muDGam\_\allowbreak{}of\_\allowbreak{}fact}{Formalism.\allowbreak{}SharpRestriction.\allowbreak{}AdmissibleParameterJoins.\allowbreak{}SumLogPrimeFactorsLe}{756--766}
\leanref{Formalism.\allowbreak{}MainNotions.\allowbreak{}J.\allowbreak{}linkMu\_\allowbreak{}of\_\allowbreak{}linkDIV}{Formalism.\allowbreak{}SharpRestriction.\allowbreak{}AdmissibleParameterJoins.\allowbreak{}SumLogPrimeFactorsLe}{775--783}

Theorem~\ref{B5.33} now gives $K,c_{HH},e_K,C_L,x_0$ for this family. At every window point $x\ge x_0$, one pair of integer weight functions works at every translation, and the resulting models realise the common chain with division exponent $2$ and loss exponent $0$. This is the connectedness assertion used in \ref{B1.18}.
\leanref{UpperTheory.\allowbreak{}Top.\allowbreak{}topDesWin}{Development.\allowbreak{}AutomorphicArithmetic.\allowbreak{}CharacterBoxes.\allowbreak{}DesWin}{11--31}
\end{proof}

\section*{Proofs for Chapter~\ref{ch:p8c1}: A common witness for the 46 conditions}

\subsection*{Proof of \ref{B8.01}}\label{prf:B8.01}\phantomsection
\begin{proof}
By \ref{B8.14} and the convergent Siegel--Weil input of \(P\), the local theta hypotheses \(F\) give the separation assertion of \ref{B8.15}. Apply \ref{B8.18} with this separation, the matching and integral comparison of \ref{B8.11}, and the toric comparison and functional equation of \ref{B1.14k} and \ref{B1.15ac}. Its Siegel-set and real Howe hypotheses belong to \(F\), and its cuspidal Green hypothesis belongs to \(P\). At each admissible stage its finite system of tests has common kernel zero and takes values in the embedded coefficient field on the theta images. This gives the rational tests required in \ref{B8.03}.

Apply Lemma~\ref{B8.08b} to \(P\). Both statements quantify over the member points \(p\in\mathcal P\) of Definition~\ref{def:p8c3-admissible}. For fixed eligible \(F_m,C_H\), take the cutoff of \ref{B8.08b}. At every index beyond it, that lemma gives the level-one inequality at the same \(p\) for every \(C\in\R\). The fitted-grid condition of \ref{B8.22} restricts the points at which its conclusion is used. Thus \ref{B8.22} gives the fitted functional bound. Apply \ref{B8.23} with this bound and the discrete-series, globalization, decay and ray-class inputs of \(P\), and denote its threshold by \(C_{F0}\). Independently, \ref{B8.21} gives a threshold \(C_{F1}\) for the grid extension, using the same local matching, the toric comparisons and the separated tests. Put
\[
 C_F=\max(C_{F0},C_{F1}),\qquad c_{\rm Dep}=C_F.
\]
Take \(A\) in (8.6) to be the grid extension condition of \ref{B8.20} at \(C_F\). The inequality \(C_F\ge C_{F1}\) gives its eventual existence.

Forgetting the grid and then the origin gives a list extension at the same member indices. The transfer in \ref{B8.08} therefore supplies the twelve open rows on the augmented extension. Forgetting the origin supplies the list required in \ref{B8.23}. With \(C_F\ge C_{F0}\), this gives the functional implication on that extension, with the family division estimate as its antecedent. Fix a datum admitted by \(\mathcal N_A\). Forget its extra grid condition, then its origin, then its signed character list. The datum left in \(\mathcal N_{\rm list}\) after forgetting the origin already consists of the real expansion data, the rational-column law and that list. The last projection therefore gives the expansion data with their law at the original base, frame, twist, norm choices and bundle indices. This gives the first quantified premise on extensions in \ref{B8.34}. Apply that theorem with the twelve open rows just obtained and the assumed covolume bound of Definition~\ref{B8.24}. Choose the weak local correspondences of Theorem~\ref{B9.46}. Definition~\ref{B1.15ac} includes the small-prime hypotheses of Definition~\ref{B8.41h} in \(P\). Thus Theorem~\ref{B8.41} supplies the linked bound at \(2,3\) for every dependence parameter, in particular for \(c_{\rm Dep}=C_F\) used here. These are the hypotheses of \ref{B8.03} for the same augmented extension, so that theorem gives the required common witness.
\leanref{UpperTheory.\allowbreak{}N.\allowbreak{}boundAtTwoThree\_\allowbreak{}linked\_\allowbreak{}of\_\allowbreak{}recT}{Development.\allowbreak{}AutomorphicArithmetic.\allowbreak{}ThetaCorrespondence.\allowbreak{}ThetaPairLifts.\allowbreak{}ThirdPlaceSeries.\allowbreak{}At23OfAll}{24--26}
\leanref{UpperTheory.\allowbreak{}Lit.\allowbreak{}every\_\allowbreak{}character\_\allowbreak{}of\_\allowbreak{}theta\_\allowbreak{}integrals\_\allowbreak{}local\_\allowbreak{}matching}{Development.\allowbreak{}AutomorphicArithmetic.\allowbreak{}ThetaCorrespondence.\allowbreak{}ThetaPairLifts.\allowbreak{}ThirdPlaceSeries.\allowbreak{}SiegelEisensteinSecondTermHypotheses}{69--73}
\end{proof}

\subsection*{Proof of \ref{B8.02d}}\label{prf:B8.02d}\phantomsection
\begin{proof}
Since \(C_\ell>2\) and \(m\) is even, \(m\ge4\) and \(r_0\ge1\). The congruence modulo \(g^2\) gives the pair congruence modulo \(g\). It also excludes \(T_u=2\) and yields coprimality with \(g\). The scalar bounds of \ref{B8.02c} give
\[
 C_\ell m^6L_b\le m^7L_b\le m^5fL_b\le k\le k\log k,\qquad
 (\log f)^2\le r_0\le m^2L_b,\qquad
 \log f\le\sqrt{r_0}\le r_0^2.
\]
The two comparisons are equalities for a repeated word and follow from \ref{B5.11b} in both orders for the indicated distinct-word row.
\end{proof}

\subsection*{Proof of \ref{B8.02ba}}\label{prf:B8.02ba}\phantomsection
\begin{proof}
For a quotient of order \(i\), record the product of the prime ideals occurring in a composition series, with multiplicities. This ideal has norm \(i\). Over a rational prime \(p\), the number of ideals of norm \(p^e\) is at most the coefficient of \(z^e\) in \((1-z)^{-2f}\). Indeed, each prime above \(p\) of residue degree \(a\) contributes \((1-z^a)^{-1}\), and there are at most \(2f\) such primes. The coefficient is at most \((2f)^e\), by counting ordered lists of length \(e\). Thus the number of ideals of norm \(i\) is at most \((2f)^{\Omega(i)}\le i^\ell\), where \(\Omega(i)\) counts rational prime factors with multiplicity.

Fix one such ideal and order its prime factors. Every corresponding sublattice is reached by successive quotients \(O/\mathfrak p\) in this order, by the primary decomposition and a composition series in each primary part. At a step of norm \(q=N\mathfrak p\), the current lattice is locally free of rank \(t\). Its sublattices with quotient \(O/\mathfrak p\) are the inverse images of the hyperplanes of \(\F_q^t\). There are \((q^t-1)/(q-1)\le q^t\). Multiplication gives at most \(i^t\) choices for the fixed ideal.

Enlarge \(P\) to a maximal integral lattice \(J\). Every integral enlargement lies in \(P^\#\), so a maximal one exists. In the fixed space \(W\) all maximal integral lattices form one genus. At a split place their two components are mutually dual. At an inert place with base residue cardinality \(q\), the maximal local model is unimodular when the determinant has even valuation. For odd valuation it has a single elementary divisor of valuation one, and its discriminant module has order \(q^2\). These are the unramified models of \cite[\S3, (3.1)--(3.3) and Tables~1--2]{gan-hanke-yu-2001}. The norm-integrality convention there agrees with pairing integrality here. Indeed, applying norm integrality to \(x+\alpha y\) for all local integers \(\alpha\) tests the pairing against the trace dual of the local integer ring, and the relative different is trivial.

Let \(b=d(J)\). Gan, Hanke and Yu's formula for a totally definite maximal lattice, applied also to the standard lattice of the same rank and over the same field, gives
\[
 \frac{\operatorname{mass}(J)}{M_t}
 =\prod_{v}\frac{q_v^t-(-1)^t}{q_v+1},
 \qquad
 b=\prod_v q_v^2.
\]
Both products run over the inert places of odd determinant valuation. The global factors agree because the signatures and ranks agree. Every displayed local factor is positive and at most \(q_v^t\) \cite[Propositions~4.4--4.5]{gan-hanke-yu-2001}. Hence \(\operatorname{mass}(J)\le M_t b^{t/2}\).

Count inclusions \(P\subset J\) up to rational isometry, weighting an inclusion by its inverse stabilizer order. Over \(P\) this gives at least \(1/|\operatorname{Aut}P|\), since there is at least one maximal enlargement. Over \(J\) each \(P\) has the same index \(i\), determined by \(n=b i^2\). This identity follows by taking indices in \(P\subset J\subset J^\#\subset P^\#\), since dualizing preserves the index. The first assertion bounds the mass of the inclusions by
\[
 M_t b^{t/2}i^{t+\ell}\le M_t n^{(t+\ell)/2}.
\]
\end{proof}

\subsection*{Proof of \ref{B8.02bb}}\label{prf:B8.02bb}\phantomsection
\begin{proof}
Write \(P=L\cap U\). Then \(L/P\) is torsion-free. Conversely a primitive sublattice \(P\) equals \(L\cap EP\), so its span recovers \(U\).

First count the rational isometry types of the embedded subspace. Enlarge \(P\) inside \(U\) to a maximal integral lattice \(J\). Let \(\mathfrak b\) be the product in \(O\) of the primes where \(J\) is not self-dual. By the local description in \ref{B8.02ba},
\[
 N\mathfrak b=d(J),\qquad a=d(J)[J:P]^2.
\]
At inert places the support of \(\mathfrak b\) therefore specifies the parity of the determinant valuation. Units are norms in an unramified quadratic local extension. At split places the Hermitian space of rank \(j\) is unique, and at every real place the signature is \((j,0)\). Landherr's classification now determines \(U\) up to isometry from \(\mathfrak b\) \cite[\S3, Theorem~3.1]{gross-2021}. Its orthogonal complement is determined as well: its local determinant class is that of \(V\) divided by that of \(U\), and its signatures are \((s,0)\). Taking the orthogonal sum of the two isometries identifies the embeddings in \(V\). By the ideal count in the proof of \ref{B8.02ba}, there are at most \(\sum_{i\le a}i^\ell\le a^{1+\ell}\) orbits of \(U\).

Fix such a subspace. Write \(W=U^\perp\), \(Q=L\cap W\), \(b=d(Q)\), and \(H=L/(P\oplus Q)\). Integrality embeds \(H\) in \(P^\#/P\oplus Q^\#/Q\). Each projection of \(H\) is injective, since its kernel is represented by a vector in \(L\cap U\) or \(L\cap W\). Taking indices gives
\[
 ab=D|H|^2,\qquad |H|\le a,\qquad b\le aD.
\]

The discriminant module of \(Q\) needs at most \(s\) generators. Locally this follows by diagonalizing the inclusion \(Q\subset Q^\#\) over a discrete valuation ring, and the local lists combine by the Chinese remainder theorem. The orthogonal projection \(p:V\to U\) satisfies \(p(L^\#)=P^\#\). Indeed, an integral functional on \(P\) extends to \(L\) because \(L/P\) is projective over \(O\). Hermitian duality identifies restriction of functionals with \(p\). Now \(p(L)/P\simeq H\) embeds in \(Q^\#/Q\), so it too needs at most \(s\) generators. The quotient \(p(L^\#)/p(L)\) is an image of the cyclic module \(L^\#/L\). Thus \(P^\#/P\) needs at most \(s+1\) generators.

For fixed \(P,Q\), the possible \(L\)'s are determined by their submodules \(H\) of the direct sum of the two discriminant modules. A finite \(O\)-module of order \(m\) generated by \(e\) elements has at most \(m^e\) submodules. Each submodule has at most \(e\) generators by the local elementary divisor theorem and the Chinese remainder theorem, and hence is the span of an ordered \(e\)-tuple. Here \(e=2s+1\). Put \(X=aD\). The number of gluings is at most \((ab)^{2s+1}\le(aX)^{2s+1}\).

To sum the weights, observe that for any action of a finite group \(A\) on a finite set \(S\),
\(\sum_{[z]\in S/A}|\operatorname{Stab}_A z|^{-1}=|S|/|A|\).
Applied over \(L\), this counts pairs \((L,U)\) with the weight on the left of the lemma. Applied with \(U\) fixed, it divides the number of gluings for \(P,Q\) by \(|\operatorname{Aut}P|\,|\operatorname{Aut}Q|\), since the rational stabilizer of \(U\) is \({\rm U}(U)\times{\rm U}(W)\). Thus all stabilizers of glued lattices are included. Apply \ref{B8.02ba} in \(U\) at order \(a\) and in \(W\) at each order \(1\le b\le X\). The resulting bound is
\[
 a^{1+\ell}\,
 M_j a^{(j+\ell)/2}\,
 X M_s X^{(s+\ell)/2}\,
 (aX)^{2s+1}.
\]
Retain \(a^{j/2}\) and use \(a\le X\) in the remaining factors. The exponent of \(X\) is
\((1+\ell)+\ell/2+1+(s+\ell)/2+4s+2=(9s+8+4\ell)/2\).
\end{proof}

\subsection*{Proof of \ref{B8.02bc}}\label{prf:B8.02bc}\phantomsection
\begin{proof}
We first bound \(M_t\) for \(t=1,2\). On a lattice \(J\) of the standard genus take the Euclidean metric whose squared length is \(\operatorname{Tr}_{E/\Q}h(z,z)\). If \(\mathfrak D_E\) is the absolute different, the trace dual is \(\mathfrak D_E^{-1}J\), since \(J\) is Hermitian self-dual. For any nonzero \(z\) in the trace dual, \(\mathfrak D_E^2 h(z,z)\) is an integral ideal. Its norm is at least one. The arithmetic-geometric mean inequality at the \(2f\) embeddings gives
\[
 \operatorname{Tr}_{E/\Q}h(z,z)\ge 2f/\rho^2.
\]
The trace lattice is full of Euclidean dimension \(2ft\). Banaszczyk's successive-minima inequality \(\lambda_{2ft}(J)\lambda_1(J^*)\le2ft\) therefore gives \(2ft\) independent vectors of squared length at most \(2ft^2\rho^2\) \cite[Theorem~2.1]{banaszczyk-1993}. Choose \(t\) among them spanning over \(E\). At each embedding their Gram entries have absolute value at most \(B=2ft^2\rho^2\), by positivity and Cauchy--Schwarz.

There are at most \((1+2B)^{2f}\) possibilities for a single integral entry. In the product of the \(f\) complex discs, distinct algebraic integers have sup distance at least one by the norm inequality. The open balls of radius \(1/2\) around them fit inside the ball of radius \(B+1/2\). For a fixed Gram matrix let \(P\) be the lattice freely generated by the chosen vectors. The determinant bound gives \(d(P)\le(tB)^{2ft}\). The original lattice lies between \(P\) and \(P^\#\), and \(P^\#/P\) is generated by at most \(t\) elements. The submodule count in \ref{B8.02bb} gives at most \((tB)^{2ft^2}\) possibilities. The Gram matrix and this intermediate lattice determine the isometry class. Consequently
\[
 M_t\le (1+2B)^{2ft^2}(tB)^{2ft^2}
 \le(20f\rho^2)^{16f}\le M_*,
\]
using \(t\le2\), \(\rho\ge2\) and \(c\ge2\).

For \(i\ge3\) put \(R_i=M_i/M_{i-1}\), \(A=\rho^{f/2}\), and \(B_0=16^f\). The Euler factor \(L_\eta(E,i)\) of \ref{B8.48} is between \(2^{-f}\) and \(2^f\): both it and its reciprocal are at most \(\zeta_F(i)\le\zeta(i)^f\le2^f\). Formula \eqref{eq:p8c5-mass-ratio}, with \(v_E\ge1\), now gives
\[
 R_i^{-1}\le A B_0^i,\qquad
 R_i/R_{i+t}\le B_0^t\quad(t\ge1).
\]
For the first inequality use \((i-1)!(16v_E/\pi)^i\ge2\). For the second use the monotonicity of the factorial and \(4\le(16v_E/\pi)^t\), which absorbs the quotient of the two Euler factors.

Write \(j=k-s\). If \(s,j\ge2\), then
\[
 \frac{M_sM_j}{M_k}
 =\frac{M_2}{R_{j+1}R_{j+2}}
      \prod_{i=3}^{s}\frac{R_i}{R_{j+i}}
 \le M_* A^2 B_0^{\,2j+3+j(s-2)}
 \le M_* A^2 B_0^{sk}.
\]
If either rank is one, the quotient is \(M_1/R_k\le M_* A B_0^k\). Since \(A\ge1\) and \(s\ge1\), this gives the same bound.
\end{proof}

\subsection*{Proof of \ref{B8.02bd}}\label{prf:B8.02bd}\phantomsection
\begin{proof}
Choose one \(E\)-prime from each of the \(f\) conjugate pairs over \(T\). For each \(L\) in the genus, take the inverse image in \(L^\#\) of the sum of the chosen components of \(L^\#/L\). Each component is a copy of \(\F_T\). The discriminant pairing pairs opposite components perfectly, so the chosen sum equals its annihilator. Its inverse image \(J\) is self-dual and is determined by \(L\) and the choice of primes.

Conversely, for fixed \(J\), at each chosen prime one obtains \(L\) as the inverse image of a hyperplane of \(J/\mathfrak pJ\); the component at the conjugate prime stays fixed. Away from \(T\) one takes \(L=J\). These prescriptions give distinct lattices. They belong to the given genus: at a split place the stabilizer acts as the full general linear group on the chosen component and transitively on its residual hyperplanes. Away from \(T\) the local lattices are self-dual. The genus mass of \(J\) is \(M_k\), by the maximal-lattice formula in the proof of \ref{B8.02ba} with empty local product. Counting the fibres with inverse stabilizer weights gives the asserted equality and, in particular,
\[
 m_0\ge M_k T^{f(k-1)}=M_kD^{(k-1)/2}.
\]

All members of the genus have cyclic discriminant module of order \(D\), since their local discriminant modules agree. Apply \ref{B8.02bb} with \(j=k-s\), divide by the last lower bound and use \ref{B8.02bc}. The remaining powers satisfy
\[
 D^{-(k-1)/2}a^{(k-s)/2}(aD)^{(9s+8+4\ell)/2}
 \le (a/D)^{k/2}(aD)^{5s+5+2c}.
\]
Indeed, the right side divided by the left side is \(a^{s+3}D^{(s+5)/2}\ge1\).
\end{proof}

\subsection*{Proof of \ref{B8.02bf}}\label{prf:B8.02bf}\phantomsection
\begin{proof}
Write \(C,C'\) for the diagonal lattices of the two graphs. On \(U\) the compatibility of the graphs gives \(C\cap U=C'\cap U=P\). Put \(y=[L\cap U:P]\). Then \(1\le y\le x\) and \(d(P)=ay^2\).

We prove the estimate for this discriminant. Let \(p\) be orthogonal projection onto \(U\), and put \(W=U^\perp\). For \(z\in C'\), \(w\in C\), the pairs \((\eta z,z),(w,w)\) lie in \(\Gamma\). Preservation of the form modulo \(g\) gives \(g^{-1}(1-\eta)C'\subset C^\#\). Also \(C'\subset L\subset C^\#\). The lattice
\[
 B=(C'\cap U)\oplus(1-p)C'
\]
has the same covolume as \(C'\), by projection onto \(W\) with kernel \(C'\cap U\). The map \(S\) equal to the identity on \(U\) and to \((1-\eta)/g\) on \(W\) is invertible. It sends \(B\) into \(C^\#\), and \(p(SB)=P\). Since \(p(C^\#)=P^\#\) by the restriction-of-functionals argument of \ref{B8.02bb}, projection gives
\[
 d(P)\le[C^\#:SB]
 =|N_{E/\Q}\det S|\,[C^\#:C'].
\]
At every complex embedding \(\eta|_W\) is unitary, so \(|\det(1-\eta|_W)|\le2^s\). Moreover \([C^\#:C']=[C^\#:L^\#][L^\#:L][L:C']=xDx'\). Hence \(ay^2\le A_s\).

There are at most \(x\) indices \(y\). At any such index, \ref{B8.02ba} gives at most \(y^{k-s+c}\) possibilities for \(P\), with one possibility in rank zero. Since \(y\le x\le X\),
\[
 y^{k-s+c}\le (A_s/a)^{k/2}X^c.
\]
For fixed \(P\), project \(C\) to \(V/U\). Its image has index \(x/y\) in the image of \(L\). There are at most \((x/y)^{s+\ell}\) such images by \ref{B8.02ba}. For a fixed image, lifts are splittings modulo \(P\), differing by a homomorphism to \((L\cap U)/P\). The image is projective of rank \(s\), so there are at most \(y^s\) homomorphisms, by local freeness at the primes dividing \(y\). This bounds the number of \(C\)'s by \(x^{2s+c}\), also for \(s=0\).

For a fixed \(C\), one has \(C\subset K\subset C^\#\). The quotient \(C^\#/C\) has order \(Dx^2\) and at most \(2r+1\) generators, by the filtration through \(L,L^\#\). Its successive quotients need at most \(r,1,r\) generators, the last by duality. The submodule count of \ref{B8.02bb} gives at most \((Dx^2)^{2r+1}\) possibilities for \(K\). Choose at most \(2r+1\) generators \(z\) of \(K/C\). For each generator choose \(w\) with \((z,w)\in\Gamma\). Pairing with \((v,v)\), \(v\in C\), shows that \((z-w)/g\in C^\#\). Its class modulo \(L\) determines \(w\bmod gL\). There are \([C^\#:L]=Dx\) such classes. The graph is fixed on \(C\), so these images determine it. Thus the lattice and graph together contribute at most \((Dx)^{3(2r+1)}\) for fixed \(C\). The same count holds on the primed side.

For fixed data the possible isometries, if any, inject by division by one of them into the automorphisms of \(K'\) fixing \(U\) pointwise. This finite group acts faithfully on \(V/U\), since it preserves \(W\). Its rational representation has dimension \(2sf\). For \(s>0\), Minkowski's finite-subgroup bound \cite[Theorem~1(i)]{serre-2007} gives an order dividing the Minkowski number of \ref{B8.02}, which divides \((4sf)!\). For \(s=0\) the group is trivial. Multiplying the counts gives
\[
 x(A_s/a)^{k/2}X^c\,
 X^{2s+c}Y^e(4sf)^{4sf}.
\]
Use \(x\le X\) to obtain the first bound.

If \(U=0\) and a triple exists, the discriminant inequality gives \(1\le(2/g)^{2kf}DX\). Count \(C,C'\) directly in \(L\), using \ref{B8.02ba}; the unrestricted count of triples is at most \(X^{k+c}Y^e(4kf)^{4kf}\). Now \(X^c\le Y^c\), and \(X^{k/2}((2/g)^{kf}D)^k\ge1\), as is seen by squaring and using \(D\ge1\). This gives the second bound.
\end{proof}

\subsection*{Proof of \ref{B8.10k}}\label{prf:B8.10k}\phantomsection
\begin{proof}
Fix the length shells, write \(u=2^{i+i'}\), and put \(M=R^rE\). Partition \((0,M]\) into the intervals \((M2^{-a-1},M2^{-a}]\). Use indices \(a,a'\) for \(Y,Y'\), with lower endpoints \(y,y'\). Write \(w=yy'\) and \(N=2^j/D\). On a nonempty piece,
\[
 w\le N,\qquad w\le R^{2r}E^2 2^{-a-a'},\qquad
 1\le NR^{2r},\qquad s_*\le u^r.
\]
Indeed \(SS'\le u(2H)^2\le R^2\). Cauchy--Schwarz on the piece gives \(\sqrt{KK'}R^ew^h u^{-q}\). In \(w^h=w^{h-1}w\), use the first bound for \(w^{h-1}\), the second for \(w\), and multiply the resulting upper bound by \(NR^{2r}\ge1\). The result is bounded by
\[
 \sqrt{KK'}(2H)^{e+4r}E^2 N^h
 s_*^{-(q-e-4r-1)/r}\,2^{-i-i'-a-a'} .
\]
The four geometric series each have sum two. Positivity permits summation over the countable partition.
\end{proof}

\subsection*{Proof of \ref{B8.02bg}}\label{prf:B8.02bg}\phantomsection
\begin{proof}
Every translate of \(w\) is smooth. Linearity gives
\[
 \Phi(g)=\|v\|cL(\sigma(g)w).
\]
Squaring and integrating gives \(I_B=\|v\|^2|c|^2/d\). Positivity of \(I_B\) gives \(c\ne0\), and division by \(I_B\) proves the formula for \(\rho\).

For the positive moment, use
\[
 l_a=l_0\prod_{i=0}^{a-1}R_i.
\]
Hölder's inequality with \(a+1\) equal exponents gives
\[
 \int l_a^{2q}\,d\rho
 \le
 \left(\int l_0^{2(a+1)q}\,d\rho\right)^{1/(a+1)}
 \prod_{i=0}^{a-1}
 \left(\int R_i^{2(a+1)q}\,d\rho\right)^{1/(a+1)}
 \le A^qR^{2aq}.
\]
For the inverse moment,
\[
 l_a^{-1}=l_r^{-1}\prod_{i=a}^{r-1}R_i
 \le\prod_{i=a}^{r-1}R_i.
\]
When \(r\ge1\), apply Hölder with \(r\) equal exponents, filling the unused positions with the constant function \(1\). Since \(\rho\) has mass one,
\[
 \int l_a^{-2q}\,d\rho
 \le\prod_{i=a}^{r-1}
       \left(\int R_i^{2rq}\,d\rho\right)^{1/r}
 \le R^{2(r-a)q}\le R^{2rq}.
\]
For \(r=0\), the inverse bound follows from \(l_0\ge1\). For \(q=0\), both assertions follow from \(\rho(S\backslash G)=1\).
\end{proof}

\subsection*{Proof of \ref{B8.02bh}}\label{prf:B8.02bh}\phantomsection
\begin{proof}
Each factor is a nonzero integral ideal. Products of integral ideals are contained in each factor, so \(\mathfrak c\subset(T)\), \(\mathfrak c\subset\mathfrak a_l\), and \(\mathfrak c\subset\overline{\mathfrak a_l}\) for \(l\in S\). Conjugation exchanges the two factors in every pair and fixes \((T)\). Thus \(\bar{\mathfrak c}=\mathfrak c\). The defining identity for \(\mathfrak a_l\) gives \(\mathfrak a_l\mathfrak d^{2q_l}=(p_l)\). Its conjugate gives the second divisibility, since \(\bar{\mathfrak d}=\mathfrak d\). For \(l=0\), use \(q_0=0\) and \(p_0=1\). Conjugation induces a bijection \(\mathcal O_E/\mathfrak a_l\simeq\mathcal O_E/\overline{\mathfrak a_l}\), so the two ideals have the same norm. Multiplicativity gives
\[
 N(\mathfrak c)=T^{2f}\prod_{l\in S}N(\mathfrak a_l)^2=(D_bN_T)^2.
\]
Every nonzero integral ideal has positive integral norm, which proves the remaining inequalities.
\end{proof}

\subsection*{Proof of \ref{B8.10y}}\label{prf:B8.10y}\phantomsection
\begin{proof}[Proof (outline)]
Write the fixed and moving frames as \((e_0,\ldots,e_{r-1},z)\) and \((e'_0,\ldots,e'_{r-1},z')\) in the Hermitian space with form \(H_E\). The first \(r\) vectors span totally isotropic spaces \(X,X'\), and the final vectors are perpendicular to their respective spaces and have norm \(-T_u\). Put \(e_r=z,\ e'_r=z'\). Let \(X_a,X'_a\) be the spans of the first \(a\) vectors, \(I=X\cap X'\), and \(j=r-\dim_E I\). Write \(p_l=\det(H_E(e_a,e'_b))_{0\le a,b<l}\), with \(p_0=1\). The condition \(G(t)=0\) is \(p_{r+1}\bar p_{r+1}=T_u^2p_r\bar p_r\). The normal form is
\begin{equation}\label{eq:p8c1-common-flag}
 X=X_j\oplus I,\qquad X'=X'_j\oplus I,\qquad
 p_j\ne0,\qquad z'-cz\in X+X',\quad c\bar c=1 .
\end{equation}

For \(I\ne0\), take \(x_*\in I\setminus\{0\}\) of levels \(a,a'\) in the two flags, so \(x_*\equiv\alpha e_a\bmod X_a\) and \(x_*\equiv\alpha'e'_{a'}\bmod X'_{a'}\) with \(\alpha,\alpha'\ne0\). For \(v\perp X_{a+1}+X'_{a'+1}\), the skew-Hermitian map
\[
 w\longmapsto H_E(w,v)x_*-H_E(w,x_*)v
\]
lies in both infinitesimal flag stabilizers. Relevance means their infinitesimal characters agree on the intersection. In these coordinates the character pairs the level-\(a\) coefficient with \(e_{a+1}\) in the first flag, and the level-\(a'\) coefficient with \(e'_{a'+1}\) in the second. With all scalar multiples of \(v\), the equality and the nondegenerate trace pairing of \(E/F\) give
\(\alpha e_{a+1}-\alpha'e'_{a'+1}\in X_{a+1}+X'_{a'+1}\).
Subtract vectors of \(X_{a+1}\) and \(X'_{a'+1}\) to obtain a common vector one level higher with the same leading coefficients, for as long as both levels are below \(r\). Reaching the final vector on one side first would place it in \(X+X'\). Every vector of \(X+X'\) perpendicular to \(X\) is isotropic, as is every vector of \(X+X'\) perpendicular to \(X'\). The negative final norms therefore force \(a=a'\). Iteration to the last step gives \(z'-cz\in X+X'\) with \(c=\alpha/\alpha'\), uniquely since \(z\notin X+X'\).

Starting at the least level in \(I\) gives a vector at each subsequent level through \(r-1\). Elimination of leading coefficients shows that these vectors form a basis of \(I\), so the least level is \(j\) and the direct sums in \eqref{eq:p8c1-common-flag} hold. If \(p_j=0\), there is \(0\ne y\in X_j\) perpendicular to \(X'_j\). It is perpendicular to \(I\), hence to \(X'\), and the final-vector relation gives \(y\perp z'\). The perpendicular of the moving frame modulo \(X'\) is positive definite, which would put \(y\) in \(X_j\cap I=0\). Thus \(X_j+X'_j\) is nondegenerate. Project \(z,z'\) onto its perpendicular by subtracting vectors of \(X_j,X'_j\), respectively. These projections retain norm \(-T_u\), and the second differs from \(c\) times the first by an element of \(I\). Both are perpendicular to \(I\), so their equal norms give \(c\bar c=1\).

If \(I=0\), then \(j=r\). Were \(p_r=0\), \(G(t)=0\) would give \(p_{r+1}=0\). Choose \(0\ne y\in X\) perpendicular to \(X'\), and a nonzero \(y'+\beta z'\in X'+Ez'\) perpendicular to \(X+Ez\). Positivity of the perpendicular modulo \(X\) gives \(\beta\ne0\), because an isotropic vector there belongs to \(X\) and \(X\cap X'=0\). Positivity modulo \(X'\) gives \(H_E(y,z')\ne0\). This contradicts \(H_E(y,y'+\beta z')=0\). With \(p_r\ne0\), subtract vectors of \(X,X'\) from \(z,z'\) to project them to \((X+X')^\perp\). They retain norm \(-T_u\) and their mutual pairing is the Schur complement \(p_{r+1}/p_r\), up to conjugation. Hence their Gram determinant is zero. The complement has one negative direction. Subtract the multiple \(c\) of the first projection that makes the difference perpendicular to it. The difference is isotropic in a positive definite space and thus vanishes. Equality of the norms gives \eqref{eq:p8c1-common-flag}.
\leanref{Formalism.\allowbreak{}HermitianLatticeCounting.\allowbreak{}P.\allowbreak{}existsNormalForm}{Formalism.\allowbreak{}FlagUnipotentLattices.\allowbreak{}MinkowskiFactorialBounds.\allowbreak{}ExistsHasLeadingCoeffOfMemSpanBelow}{1421--1458}

If \(j=0\), a common vector of each level gives \(X_a=X'_a\) and \(e'_a-ce_a\in X_a\). In \(X^\perp/X\) the two final vectors differ by the scalar \(c\). Their primitive lines in the row lattice make \(c\) a unit. Every conjugate of \(c\) has absolute value one, so two of its powers coincide by discreteness of the embedded ring of integers. Pairing the level congruence with the fixed final vector gives \(T_u(c-1)=0\) modulo each prime above \(g_0g_1\). Both primes are prime to \(T_u\). For an element of order \(d>1\) prime to the residue characteristic in the kernel of reduction, \((X^d-1)/(X-1)\) vanishes at that element but reduces to \(d\), a contradiction. The order of \(c\) is consequently a power of each of the distinct primes \(g_0,g_1\), and \(c=1\).

Let \(\Lambda\) be the row lattice. On \(X\) the map \(U_X:e_a\mapsto e'_a\) is unitriangular. Locally the two frames rescaled by their graded ideals give bases of \(X\cap\Lambda\). Equality in each graded line identifies the rescaling ideals, so \(U_X\) is integral at each prime and hence on \(X\cap\Lambda\). This sublattice is saturated, so the quotient is projective over \(\mathcal O_E\) and gives an integral retraction \(P_0\) onto \(X\cap\Lambda\). Unimodularity makes its Hermitian adjoint integral. Choose \(\theta\in\mathcal O_E\) with \(\theta+\bar\theta=1\) by trace surjectivity for the finite unramified extension. Then \(P_X=P_0-\theta P_0P_0^*\) is the identity on \(X\) and satisfies \(P_XP_X^*=0\), by \(P_0^*P_0=0\) and \(P_0^2=P_0\). For \(D_X=(U_X-1)P_X\) we have \(D_XD_X^*=D_X^*D_X=0\) and \(D_X^{r+1}=0\). Hence \((1+D_X)(1+D_X^*)^{-1}\) is an integral isometry. Both inverses are finite geometric sums through degree \(r\). The isometry sends \(e_a\) to \(e'_a\), and for \(z\in X^\perp\) it changes \(z\) by an element of \(X\). Since \(z'-z\in X\), these are the frame conditions for the identity cell at the tied level. Thus \(j\ge1\) for the terms in the proposition.
\leanref{Formalism.\allowbreak{}MainNotions.\allowbreak{}BridgeLemma.\allowbreak{}l\_\allowbreak{}iso}{Formalism.\allowbreak{}SharpRestriction.\allowbreak{}CuspSpanNormBounds.\allowbreak{}BruhatCellsKronecker.\allowbreak{}LisoSmulMem}{15--66}
\leanref{Formalism.\allowbreak{}MainNotions.\allowbreak{}BridgeLemma.\allowbreak{}id\_\allowbreak{}lift}{Formalism.\allowbreak{}SharpRestriction.\allowbreak{}CuspSpanNormBounds.\allowbreak{}BruhatCellsKronecker.\allowbreak{}IdLift}{249--251}
\leanref{Formalism.\allowbreak{}L2.\allowbreak{}SG.\allowbreak{}KappaOffDiagonal.\allowbreak{}one\_\allowbreak{}le\_\allowbreak{}index\_\allowbreak{}of\_\allowbreak{}relevant}{Formalism.\allowbreak{}LocalToGlobal.\allowbreak{}SiegelLatticeTheory.\allowbreak{}CokernelGLCounts.\allowbreak{}RankPointCounts.\allowbreak{}IsIdentityOfEqualFlags}{71--87}

For a term \(t\) let \(w_t\) be the geometric mean of the two inverse row-stabilizer orders and \(\omega_t\) the infimum of the inverse relative stabilizer volumes over adapted forms. We prove
\begin{equation}\label{eq:p8c1-type-factors}
 \|v_t\|\le w_t\sqrt{\dim\pi_u}\,\omega_t A(j,\delta),\qquad
 \sum_{(t,t')\in U}w_tw_{t'}\omega_t\omega_{t'}
 \le\operatorname{mass}_u M(j,\delta)M(j',\delta').
\end{equation}

The graded adapted bases put the reduced pairing in block antidiagonal form, with nonzero leading minors at the block ends. At a place \(\tau\) put \(P_\tau=d_{\tau j}\) and
\[
 s_\tau=\max(1,\max_{1\le l\le j}d_{\tau l}/(2H_1)^{2r}),\qquad
 e_\tau=\max(1,\max_{\substack{1\le l\le j\\\delta_l\ne0}}d_{\tau l}^{-1}),
\]
where \(d_{\tau l}=0\) for a zero minor. The normalized nonzero minor ideals are integral, so \(\prod_\tau d_{\tau l}=\sqrt{d_l}\ge1\) and \(\prod_\tau P_\tau=\sqrt{D_\delta}\). For each minor the product of the factors above one bounds the product of the reciprocal factors below one. Taking the product over the minors gives
\begin{equation}\label{eq:p8c1-reciprocal-minors}
 \prod_\tau e_\tau\le\prod_\tau\bigl((2H_1)^{2r}s_\tau\bigr)^{r+1}.
\end{equation}

Put \(h=k-1-j\) and \(q=(\lceil C_s\rceil+2)kr\). The thresholds give \(q\le km^3\) and \(h\ge k/2\). At a real place put \(n_0=k+2r+1\), \(T=T_u\), \(\beta=b^{\max(2,C_w)}\), and \(q_0=km^5\). Use the normalized Bessel family through the given representation, indexed by integral shifts \(\nu\) with \(\|\nu\|_1\le q_0\). Its representation at \(\nu\) has parameter \(\lambda+\nu\), its standard period is \(L_\nu\), and \(d_\nu>0\) is its formal degree for the specified group Haar measure. The family includes the numerical gap conditions, the quotient-measure identity, and the factorization of the Bessel integral as the rank-one operator defined by \(L_\nu\). The original representation is at \(\nu=0\). In the numerical family the local scale \(b_v\) satisfies
\[
 10^6n_0^7\le b_v\le\beta,\qquad n_0\le\beta,\qquad 6b_v\le\beta^2.
\]
The spectral separation and compact-weight parameters are \(n_0^{19}\) and \(n_0m^{10}\), respectively. The size hypotheses give the required numerical conditions on the entire shift ball, including the large-gap condition with separation \(n_0^{19}/2\). Moving the character transports the unipotent measure by \(n\mapsto t^{-1}nt\) and the quotient measure by \(x\mapsto t^{-1}x\); the group Haar measure, and hence \(d_\nu\), stay fixed. The normalized periods are chosen for these transported measures.
\leanref{BranchingSupport.\allowbreak{}SideFam}{Development.\allowbreak{}AbelianVarietyTheory.\allowbreak{}Definitions.\allowbreak{}Jets.\allowbreak{}SideFam}{36--41}
\leanref{GlobalTheory.\allowbreak{}SizeNum.\allowbreak{}num\_\allowbreak{}of\_\allowbreak{}size}{Development.\allowbreak{}AbelianVarietyTheory.\allowbreak{}Proofs.\allowbreak{}FlagMeasures.\allowbreak{}MomentPieriNumericsOfSizeHyps}{175--485}
\leanref{GlobalTheory.\allowbreak{}Amp.\allowbreak{}LocalSetupGap.\allowbreak{}exists\_\allowbreak{}setup\_\allowbreak{}both}{Development.\allowbreak{}AbelianVarietyTheory.\allowbreak{}Proofs.\allowbreak{}FlagMeasures.\allowbreak{}AmplitudeLocalDataMomentPieriHypotheses}{59--98}
\leanref{BranchingSupport.\allowbreak{}Setup}{Development.\allowbreak{}AbelianVarietyTheory.\allowbreak{}Definitions.\allowbreak{}Jets.\allowbreak{}SideFam}{42--46}
\leanref{BranchingSupport.\allowbreak{}IsBeL}{Development.\allowbreak{}AbelianVarietyTheory.\allowbreak{}Definitions.\allowbreak{}Jets.\allowbreak{}SumIoPrAlt}{135--140}
\leanref{GlobalTheory.\allowbreak{}TwoSidedMomentsGap.\allowbreak{}two\_\allowbreak{}sided\_\allowbreak{}moments\_\allowbreak{}gap}{Development.\allowbreak{}AbelianVarietyTheory.\allowbreak{}Proofs.\allowbreak{}FlagMeasures.\allowbreak{}MajorantLengthTwoSidedMomentsLargeGap}{206--341}

Let \(\mu\) be the quotient measure and \(\Phi(x)=B(\sigma(x)w)\), where \(w\) is a lowest-compact-type unit vector. All square norms in this paragraph use \(\mu\). A zero coefficient contributes zero. Suppose \(0<\|\Phi\|_2^2<\infty\). The period nonvanishing identity applied to this coefficient produces a \(K\)-finite vector on which \(L_0\) is nonzero. Bessel multiplicity one gives \(B=cL_0\) on smooth vectors. The coefficient-period norm identity gives the normalization. Choose the \(K\)-finite vector \(v\) just obtained, so \(L_0v\ne0\), and write \(\mathcal B_{y,v}\) for the Bessel coefficient-period operator. Its rank-one factorization gives
\[
 \mathcal B_{y,v}(L_0v)=\|L_0v\|^2L_0y.
\]
The coefficient-period norm identity, with the group Haar measure defining \(d_0\) and the quotient measure \(\mu\), therefore reads
\[
 \int\|\mathcal B_{\sigma(x)w,v}(L_0v)\|^2\,d\mu(x)
 =\frac{\|w\|^2}{d_0}
   \operatorname{Re}\langle L_0v,\mathcal B_{v,v}(L_0v)\rangle
 =\frac{\|w\|^2\|L_0v\|^4}{d_0}.
\]
Substitute the rank-one formula and cancel the positive number \(\|L_0v\|^4\). This proves that the Bessel realization \(w\mapsto[x\mapsto\sqrt{d_0}L_0(\sigma(x)w)]\) is isometric on smooth vectors, and in particular
\[
 d_0\int\|L_0(\sigma(x)w)\|^2\,d\mu(x)=1.
\]
\leanref{GlobalTheory.\allowbreak{}scaled\_\allowbreak{}bessel\_\allowbreak{}isometric}{Development.\allowbreak{}AbelianVarietyTheory.\allowbreak{}Proofs.\allowbreak{}JetHarmonics.\allowbreak{}GapCoefficientDerivativeBound}{27--51}
Lemma~\ref{B8.02bg} consequently gives \(c\ne0\) and the probability measure
\[
 \|\Phi\|_2^2=\frac{|c|^2}{d_0},\qquad
 d\rho(x)=\frac{\|\Phi(x)\|^2}{\|\Phi\|_2^2}\,d\mu(x)
 =d_0\|L_0(\sigma(x)w)\|^2\,d\mu(x).
\]
\leanref{GlobalTheory.\allowbreak{}T.\allowbreak{}exists\_\allowbreak{}kFinite\_\allowbreak{}besselPeriod\_\allowbreak{}ne\_\allowbreak{}zero\_\allowbreak{}gap}{Development.\allowbreak{}AbelianVarietyTheory.\allowbreak{}Proofs.\allowbreak{}FlagMeasures.\allowbreak{}ExistsKFiniteBesselPeriodNeZeroGapStatement}{12--14}
\leanref{GlobalTheory.\allowbreak{}GapLine.\allowbreak{}bessel\_\allowbreak{}line\_\allowbreak{}gap}{Development.\allowbreak{}AbelianVarietyTheory.\allowbreak{}Proofs.\allowbreak{}FlagMeasures.\allowbreak{}BesselMapsProportionalDiscreteSeriesLargeGap}{45--54}
\leanref{Gp91.\allowbreak{}GlobalTheory.\allowbreak{}T.\allowbreak{}bessel\_\allowbreak{}functional\_\allowbreak{}l2\_\allowbreak{}bound\_\allowbreak{}and\_\allowbreak{}isometry}{CompanionTheory.\allowbreak{}AutomorphicTheory.\allowbreak{}PieriRules.\allowbreak{}ThetaBesselIsometries.\allowbreak{}FamilySetupBesselGrowthStatement}{16--18}
\leanref{GlobalTheory.\allowbreak{}MomMass.\allowbreak{}mom\_\allowbreak{}mass}{Development.\allowbreak{}AbelianVarietyTheory.\allowbreak{}Proofs.\allowbreak{}JetHarmonics.\allowbreak{}FamilySetupBesselGrowth}{21--57}

The majorant \(l_a(x)\) is the Euclidean distance of \(x^{-1}e_a\) from the span of the preceding transformed frame vectors. It is positive almost everywhere. In normalized character coordinates the two moment estimates, with their separate ranges, are, for integers \(p\ge0\),
\begin{equation}\label{eq:p8c1-cayley-moments}
\begin{aligned}
 &A_0=\max(1,T)\beta^{20(r+1)},\qquad R_0=\beta^2,\\
 &\mathbb E_\rho l_a^{2p}\le A_0^p
 &&\bigl(0\le a\le r,\ (r+1)p\le q_0,\ (r+1)p\le2n_0^2\bigr),\\
 &\mathbb E_\rho(l_{a+1}/l_a)^{2p}\le R_0^{2p}
 &&\bigl(0\le a<r,\ p\le2n_0^2\bigr).
\end{aligned}
\end{equation}
The positive-moment estimate includes the initial length \(a=0\). At a general shift \(\nu\), its spectral condition is \(\|\nu\|_1+(r+1)p\le q_0\); the ratio estimate requires \(\|\nu\|_1\le q_0\).
\leanref{BranchingSupport.\allowbreak{}lenA}{Development.\allowbreak{}AbelianVarietyTheory.\allowbreak{}Definitions.\allowbreak{}Jets.\allowbreak{}SumIoPrAlt}{150--154}
\leanref{BranchingSupport.\allowbreak{}Hsq}{Development.\allowbreak{}AbelianVarietyTheory.\allowbreak{}Definitions.\allowbreak{}Jets.\allowbreak{}SideFam}{34--35}
\leanref{GlobalTheory.\allowbreak{}T.\allowbreak{}oneMajorantLength\_\allowbreak{}moment\_\allowbreak{}lowestKType}{Development.\allowbreak{}AbelianVarietyTheory.\allowbreak{}Proofs.\allowbreak{}FlagMeasures.\allowbreak{}OneMajorantLengthMomentLowestKTypeStatement}{12--14}

Here is the iteration for the ratio estimate. Fix a unitary Cayley matrix and its infinitesimal matrix units \(E_{uv}\), so \([E_{uv},E_{st}]=\delta_{vs}E_{ut}-\delta_{tu}E_{sv}\). For a highest-weight unit vector in the lowest compact type, set \(R_a=l_{a+1}/l_a\) and
\[
 M_p(v)=\int R_a^{2p}d_\nu\|L_\nu(\sigma_\nu(x)v)\|^2\,d\mu(x).
\]
The pointwise side inequality is obtained by differentiating Bessel equivariance along the simple-root one-parameter subgroup at \(x\), expanding the derived action in the Cayley matrix units, and applying Cauchy--Schwarz to the expansion. It gives
\[
 R_a(x)^2\|L_\nu(\sigma_\nu(x)v)\|^2
 \le\sum_{u,v'}\|L_\nu(\sigma_\nu(x)E_{uv'}v)\|^2.
\]
For a vector of level \(i\le4n_0^2\), the diagonal operators contribute at most \(n_0(3b_v)^2\|v\|^2\); the off-diagonal operators contribute at most \((b_v/16)^2\|v\|^2\). The numerical scale gives
\[
 n_0(3b_v)^2+(b_v/16)^2
 \le b_v^4\le\beta^4.
\]
The general linear commutation relations and unitarity put each \(E_{uv'}v\) at level at most \(i+2\). The weight changes by a difference of two coordinate weights, of \(\ell^1\)-norm at most two; the compact Casimir bound changes from the threshold \(\sqrt{c_K}+4n_0i\) to at most \(\sqrt{c_K}+4n_0(i+2)\), where \(c_K\) is the lowest compact-type Casimir eigenvalue. Thus the pointwise inequality, the level estimate, and the mass bound give
\[
 M_{p+1}(v)\le\sum_{u,v'}M_p(E_{uv'}v),\qquad
 \sum_{u,v'}\|E_{uv'}v\|^2\le\beta^4\|v\|^2,
 \qquad M_0(v)\le\|v\|^2.
\]
For an initial vector at level \(i\), these estimates iterate through \(i+2p\le4n_0^2\), because each matrix unit raises the level by at most two. Induction gives \(M_p(v)\le\beta^{4p}\|v\|^2\). Taking \(i=0\) gives the second estimate in \eqref{eq:p8c1-cayley-moments}. These estimates hold for every unitary Cayley choice. To pass to the full lowest compact type, let \(\varphi\) be a nonnegative measurable function of the lengths and consider
\[
 \mathcal Q_\varphi(v)
 =\int\varphi(l_0(x),\ldots,l_r(x))\,
          d_\nu\|L_\nu(\sigma_\nu(x)v)\|^2\,d\mu(x).
\]
Right translation by the Cayley maximal compact subgroup preserves the Euclidean lengths and the quotient measure. Hence this quadratic form is compact-invariant. The vectors on which it is finite form a linear compact-invariant subspace. The bound holds for every highest-weight unit vector. Each irreducible copy of the lowest compact type contains such a vector and is therefore contained in this finite-integrability subspace. Every vector of the algebraic lowest isotypic space lies in a finite-dimensional compact-stable subspace. On that subspace the quadratic form is continuous and is represented by a positive self-adjoint operator commuting with the compact group. Its largest eigenspace is compact-stable and contains a highest-weight unit vector. The bound on such vectors therefore bounds the quadratic form on every unit vector in the subspace. All these highest-weight vectors have level zero. Apply this argument to \(\varphi=R_a^{2p}\); it will also apply to \(\varphi=l_a^{2p}\).
\leanref{GlobalTheory.\allowbreak{}MomentPointwise.\allowbreak{}pointwise}{Development.\allowbreak{}AbelianVarietyTheory.\allowbreak{}Proofs.\allowbreak{}JetHarmonics.\allowbreak{}BesselCayleySideBound}{52--111}
\leanref{GlobalTheory.\allowbreak{}lev\_\allowbreak{}gl\_\allowbreak{}size}{Development.\allowbreak{}AbelianVarietyTheory.\allowbreak{}Proofs.\allowbreak{}JetHarmonics.\allowbreak{}UnitaryGlModuleLevelSumNormSqLe}{11--58}
\leanref{GlobalTheory.\allowbreak{}lev\_\allowbreak{}E\_\allowbreak{}add\_\allowbreak{}two}{Development.\allowbreak{}AbelianVarietyTheory.\allowbreak{}Proofs.\allowbreak{}JetHarmonics.\allowbreak{}LevEAddTwo}{134--163}
\leanref{GlobalTheory.\allowbreak{}RatioMoment.\allowbreak{}iterate\_\allowbreak{}bound}{Development.\allowbreak{}AbelianVarietyTheory.\allowbreak{}Proofs.\allowbreak{}JetHarmonics.\allowbreak{}FamilySetupCayleyBound}{16--65}
\leanref{GlobalTheory.\allowbreak{}MomOne.\allowbreak{}lintegral\_\allowbreak{}le\_\allowbreak{}of\_\allowbreak{}highest}{Development.\allowbreak{}AbelianVarietyTheory.\allowbreak{}Proofs.\allowbreak{}FlagMeasures.\allowbreak{}LintegralLeOfHighest}{168--253}
\leanref{GlobalTheory.\allowbreak{}LowTypeMassGap.\allowbreak{}weightedMass\_\allowbreak{}bound\_\allowbreak{}lowestType\_\allowbreak{}largeGap}{Development.\allowbreak{}AbelianVarietyTheory.\allowbreak{}Proofs.\allowbreak{}FlagMeasures.\allowbreak{}WeightedMassHighestToLowestKTypeLargeGap}{63--166}
\leanref{Gp91.\allowbreak{}GlobalTheory.\allowbreak{}T.\allowbreak{}family\_\allowbreak{}setup\_\allowbreak{}cayley\_\allowbreak{}bound}{CompanionTheory.\allowbreak{}AutomorphicTheory.\allowbreak{}PieriRules.\allowbreak{}ThetaBesselIsometries.\allowbreak{}FamilySetupCayleyBoundStatement}{16--18}
\leanref{GlobalTheory.\allowbreak{}RatioMoment.\allowbreak{}ratio\_\allowbreak{}moment}{Development.\allowbreak{}AbelianVarietyTheory.\allowbreak{}Proofs.\allowbreak{}JetHarmonics.\allowbreak{}FamilySetupCayleyBound}{66--96}
\leanref{GlobalTheory.\allowbreak{}T.\allowbreak{}weightedMass\_\allowbreak{}highestToLowestKType\_\allowbreak{}largeGap}{Development.\allowbreak{}AbelianVarietyTheory.\allowbreak{}Proofs.\allowbreak{}FlagMeasures.\allowbreak{}WeightedMassHighestToLowestKTypeLargeGapStatement}{15--17}

The initial-length recursion has the sharper constant
\[
 A_{\rm first}=n_0^2T\beta^4(6b_v)^{2(2r+2)}.
\]
For a smooth level-\(i\) vector \(v\) in the shifted representation it gives
\[
 \int l_0^{2p}d_\nu\|L_\nu(\sigma_\nu(x)v)\|^2\,d\mu(x)
 \le A_{\rm first}^{p}\|v\|^2,
 \quad
 \|\nu\|_1+p\le q_0,\quad i+2p\le4n_0^2.
\]
The induction starts with the mass bound. Let \(\varepsilon_s\) be the \(s\)-th coordinate vector of the shift lattice. In the step from \(p\) to \(p+1\), the Pieri moment step applies the order-\(p\) bound to level-\(i+2\) vectors in representations of parameter \(\lambda+\nu-\varepsilon_s\), and multiplies its constant by \(A_{\rm first}\). Both reserves decrease by exactly the required amount: \(\|\nu-\varepsilon_s\|_1\le\|\nu\|_1+1\) and \((i+2)+2p=i+2(p+1)\).

For a level-zero unit vector \(v\) and a period nonzero on smooth vectors, the measure \(d\rho_{\nu,v}=d_\nu\|L_\nu(\sigma_\nu(x)v)\|^2\,d\mu(x)\) has mass one. Lemma~\ref{B8.02bg}, using first-length and ratio moments through order \((r+1)p\), gives
\[
 \mathbb E_{\rho_{\nu,v}} l_a^{2p}
 \le(A_{\rm first}\beta^{4a})^p
 \le(A_{\rm first}\beta^{4r})^p.
\]
If the period vanishes on all smooth vectors, the corresponding weighted integrals are zero. The uniform bound thus transfers to every lowest-type unit vector by the preceding compact-invariant quadratic-form argument. Finally,
\[
 \begin{aligned}
 A_{\rm first}\beta^{4r}
 &\le \max(1,T)\beta^{\,2+4+4(2r+2)+4r}\\
 &=\max(1,T)\beta^{12r+14}
 \le\max(1,T)\beta^{20(r+1)}=A_0.
 \end{aligned}
\]
This proves the first estimate in \eqref{eq:p8c1-cayley-moments}, with its stated spectral and degree reserves.
\leanref{GlobalTheory.\allowbreak{}Moment.\allowbreak{}moment\_\allowbreak{}first\_\allowbreak{}length}{Development.\allowbreak{}AbelianVarietyTheory.\allowbreak{}Proofs.\allowbreak{}JetHarmonics.\allowbreak{}FamilySetupCayleyRecursion}{24--42}
\leanref{Gp91.\allowbreak{}GlobalTheory.\allowbreak{}T.\allowbreak{}family\_\allowbreak{}setup\_\allowbreak{}cayley\_\allowbreak{}recursion}{CompanionTheory.\allowbreak{}AutomorphicTheory.\allowbreak{}PieriRules.\allowbreak{}ThetaBesselIsometries.\allowbreak{}FamilySetupCayleyRecursionStatement}{16--18}
\leanref{GlobalTheory.\allowbreak{}MomOne.\allowbreak{}single\_\allowbreak{}moment\_\allowbreak{}le}{Development.\allowbreak{}AbelianVarietyTheory.\allowbreak{}Proofs.\allowbreak{}FlagMeasures.\allowbreak{}MassBoundLowTypesSum}{39--78}
\leanref{GlobalTheory.\allowbreak{}MomNum.\allowbreak{}mom\_\allowbreak{}num}{Development.\allowbreak{}AbelianVarietyTheory.\allowbreak{}Proofs.\allowbreak{}JetHarmonics.\allowbreak{}NumericalFamilyProductBoundConstantLe}{9--37}
\leanref{GlobalTheory.\allowbreak{}MomOneGap.\allowbreak{}oneLength\_\allowbreak{}moment\_\allowbreak{}lowestType\_\allowbreak{}largeGap}{Development.\allowbreak{}AbelianVarietyTheory.\allowbreak{}Proofs.\allowbreak{}FlagMeasures.\allowbreak{}OneMajorantLengthMomentLowestKType}{62--181}

For the positive moment at order \(q\), apply the first estimate directly. The spectral and degree reserves are
\[
 (r+1)q\le km^4\le q_0,\qquad
 rq\le(r+1)q\le km^4\le k^2\le2n_0^2,
\]
since \(r+1\le m\) and \(m^4\le f\le k\). For the inverse moment use \(l_r^2\ge T\ge1\) and
\[
 l_a^{-1}\le\prod_{i=a}^{r-1}(l_{i+1}/l_i).
\]
Hölder with \(r\) equal exponents, filling unused factors with \(1\), applies the ratio estimate at order \(rq\). Thus
\[
 \mathbb E_\rho l_a^{2q}\le A_0^q,\qquad
 \mathbb E_\rho l_a^{-2q}\le\beta^{4rq}
 \quad(0\le a\le r).
\]
For \(r=0\), the inverse assertion follows from \(l_0\ge1\).
\leanref{GlobalTheory.\allowbreak{}InvLenMomentGap.\allowbreak{}inv\_\allowbreak{}len\_\allowbreak{}moment\_\allowbreak{}gap}{Development.\allowbreak{}AbelianVarietyTheory.\allowbreak{}Proofs.\allowbreak{}FlagMeasures.\allowbreak{}InverseMajorantLengthMomentLargeGap}{24--77}

Return to the specified character by the diagonal torus scaling with parameter \(\xi=\pi/(v_Eg_0g_1)\). Pulling back the quotient measure by the same translation preserves the coefficient's square norm. Each length and its reciprocal change by factors at most
\[
 K_T=\max(\xi,\xi^{-1})^r\le b^{(2|C_g|+2)r}.
\]
Since \(T\le b^{|C_T|m}\), \(m=2r+2\), and \(b\ge2\),
\[
 K_T\sqrt{A_0}\le
 b^{(2|C_g|+2)r+|C_T|m/2+10\max(2,C_w)(r+1)}\le H_1,
 \qquad K_T\beta^{2r}\le H_1
\]
for \(C_H=2|C_g|+2+|C_T|+10\max(2,C_w)\). For \(S=\max_a\{l_a,l_a^{-1}\}\), Markov's inequality and the union bound give
\begin{equation}\label{eq:p8c1-type-tail}
 \int l_a^{\pm2q}\|\Phi\|^2\,d\mu\le H_1^{2q}\|\Phi\|_2^2,\qquad
 \int_{S>\lambda H_1}\|\Phi\|^2\,d\mu
 \le 2(2r+2)\lambda^{-2q}\|\Phi\|_2^2\quad(\lambda\ge1).
\end{equation}
\leanref{GlobalTheory.\allowbreak{}TwoSidedMomentsGap.\allowbreak{}moment\_\allowbreak{}of\_\allowbreak{}bessel\_\allowbreak{}gap}{Development.\allowbreak{}AbelianVarietyTheory.\allowbreak{}Proofs.\allowbreak{}FlagMeasures.\allowbreak{}MajorantLengthTwoSidedMomentsLargeGap}{57--127}
\leanref{GlobalTheory.\allowbreak{}TwoSidedMomentsGap.\allowbreak{}inv\_\allowbreak{}moment\_\allowbreak{}of\_\allowbreak{}bessel\_\allowbreak{}gap}{Development.\allowbreak{}AbelianVarietyTheory.\allowbreak{}Proofs.\allowbreak{}FlagMeasures.\allowbreak{}MajorantLengthTwoSidedMomentsLargeGap}{129--198}
\leanref{GlobalTheory.\allowbreak{}InvLenMomentGap.\allowbreak{}inv\_\allowbreak{}len\_\allowbreak{}moment\_\allowbreak{}gap}{Development.\allowbreak{}AbelianVarietyTheory.\allowbreak{}Proofs.\allowbreak{}FlagMeasures.\allowbreak{}InverseMajorantLengthMomentLargeGap}{24--77}
\leanref{GlobalTheory.\allowbreak{}TwoSidedMomentsGap.\allowbreak{}two\_\allowbreak{}sided\_\allowbreak{}moments\_\allowbreak{}gap}{Development.\allowbreak{}AbelianVarietyTheory.\allowbreak{}Proofs.\allowbreak{}FlagMeasures.\allowbreak{}MajorantLengthTwoSidedMomentsLargeGap}{206--341}
\leanref{GlobalTheory.\allowbreak{}Amp.\allowbreak{}TailTwo.\allowbreak{}tailBd\_\allowbreak{}of\_\allowbreak{}moments}{Development.\allowbreak{}AbelianVarietyTheory.\allowbreak{}Proofs.\allowbreak{}FlagMeasures.\allowbreak{}MajorantLengthTailBoundOfMoments}{58--135}

Use length shells with lower endpoints \(2^iH_1,2^{i'}H_1\), \(i,i'\ge0\), including lengths below \(H_1\) in the bottom intervals. Put \(R=2^{i+i'+1}H_1\) and \(M_0=R^{2r}e_\tau\). In the local section the \(j\) vector columns have length at most \(M_0\), the triangular variables have radius \(2rM_0^2\), and the off-block variables have radius \(M_0\). These bounds come from solving the graded pairing equations with reciprocal leading minors bounded by \(e_\tau\).

Let \(Y\) be the Gram volume of the columns projected onto the positive coordinate space, and put \(d=k-1\). Projection onto a \(d\)-dimensional subspace of that space can decrease the Gram volume. The determinant-one change of column coordinates in the chart preserves Lebesgue measure, and the two complementary complex coordinates per column lie in disks of radius \(M_0\). To bound the projected volume, let \(V(d,j;R_1,t_1)\) be the volume of \(j\) columns in \(\C^d\) of lengths at most \(R_1\) and Gram volume at most \(t_1\). For a nonzero first column \(x\), split the remaining columns into their components along \(x\) and their projections onto \(x^\perp\). The first components cost \((\pi R_1^2)^{j-1}\), and the projected \(j-1\) columns have Gram volume at most \(t_1/\|x\|\). If \(h=d-j\), induction gives a factor \(\|x\|^{-2h}\) in the fibre bound. Radial integration gives
\[
 \int_{\|x\|\le R_1}\|x\|^{-2h}\,dx
 =\frac{\pi^d R_1^{2j}}{j(d-1)!}.
\]
Using \(j(d-1)!\ge h!\) and \(2(j-1)^2+2(j-1)+2j=2j^2\), the induction yields
\[
 V(d,j;R_1,t_1)\le
 \frac{t_1^{2(d-j)} R_1^{2j^2}\pi^{jd}}{((d-j)!)^j}.
\]
The empty tuple has volume at most one. The first-column exceptional set has measure zero. Multiplication by \((\pi R_1^2)^{2j}\) accounts for the complementary disks. In particular the zero Gram set is null when \(h>0\).

Take \(t_1=2y,\ R_1=M_0\). If the numbers of complex triangular, real diagonal and complex off-block coordinates are \(n_1,n_2,n_3\), respectively, the section box contributes
\[
 [\pi(2rM_0^2)^2]^{n_1}(4rM_0^2)^{n_2}(\pi M_0^2)^{n_3}.
\]
Here \(n_1,n_3\le r(r-1)/2\), \(n_2\le r\), and \(j\le r\). Using \(h!\ge(h/e)^h\) and \(M_0\ge1\), the product of the box and Gram bounds for \(y<Y\le2y\) is at most
\begin{equation}\label{eq:p8c1-gram-section}
 (4r)^{4r^2}M_0^{8r^2}(\pi e/h)^{jh}(2y)^{2h}.
\end{equation}

In the normal form the product coordinates are \(\rho s\), with \(\rho\) in the common unipotent stabilizer and \(s\) in the section. The stabilizer coordinates are the complex within-block upper entries, the \(r-j\) complex columns of length \(k+1\), the complex triangular entries with row index at least \(j\), and the corresponding real diagonal entries. Use product Lebesgue measure in these coordinates, transported by the conjugation that removes the second unipotent factor of the adapted frame. The section uses the complementary coordinates of the box above. Multiplication has triangular Jacobian one in the product coordinates. If the stabilizer measure is multiplied by \(s_R>0\), the quotient measure is divided by \(s_R\), so the bound in \eqref{eq:p8c1-gram-section} is divided by \(s_R\). The arithmetic covolume comparison below uses \(s_R=1\). The reversed section has the corresponding bound with Gram volume \(Y'\). Left translation by the relative frame transformation carries its quotient measure to the first section's measure.
\leanref{GlobalTheory.\allowbreak{}Amp.\allowbreak{}GramVol.\allowbreak{}vol\_\allowbreak{}le}{Development.\allowbreak{}AbelianVarietyTheory.\allowbreak{}Proofs.\allowbreak{}FlagMeasures.\allowbreak{}MeasurePreservingMulVec}{100--163}
\leanref{GlobalTheory.\allowbreak{}Amp.\allowbreak{}GramVol.\allowbreak{}vol\_\allowbreak{}gram}{Development.\allowbreak{}AbelianVarietyTheory.\allowbreak{}Proofs.\allowbreak{}FlagMeasures.\allowbreak{}MeasurePreservingMulVec}{189--227}
\leanref{GlobalTheory.\allowbreak{}Amp.\allowbreak{}FibreVolume.\allowbreak{}fibre\_\allowbreak{}volume}{Development.\allowbreak{}AbelianVarietyTheory.\allowbreak{}Proofs.\allowbreak{}FlagMeasures.\allowbreak{}FibreVolumeNormalFormStabiliser}{79--175}

The pairing and the two length maxima \(S_1,S_2\) give, in these same coordinates,
\begin{equation}\label{eq:p8c1-gram-pairing}
 YY'P_\tau\le2^j,\qquad
 s_\tau\le\max\bigl(1,(S_1S_2/(2H_1)^2)^r\bigr).
\end{equation}
The first inequality is the determinant inequality for the perpendicular projections, with a factor of at most two per paired vector. The second follows by bounding each order-\(l\) leading minor by \((S_1S_2)^l\).

Use descending dyadic shells for \(Y,Y'\) from \(R^re_\tau\), with indices \(a,a'\ge0\). Cauchy--Schwarz combines the two section bounds \eqref{eq:p8c1-gram-section} with the tails \(2^{-2qi'},2^{-2qi}\) from \eqref{eq:p8c1-type-tail}. Apply \ref{B8.10k} with \(D=P_\tau\), \(E=e_\tau\), \(H=H_1\), \(s_*=s_\tau\) and \(e=16r^3\). Indeed \(M_0^{8r^2}=R^{16r^3}e_\tau^{8r^2}\) in \eqref{eq:p8c1-gram-section}. The upper bound for the two Gram volumes is \(R^r e_\tau\). The leading minor of order \(j\) is at most \((S_1S_2)^j\le(S_1S_2)^r\), and \eqref{eq:p8c1-gram-pairing} gives the other two product bounds. The threshold on \(q\) gives \(q\ge16r^3+4r+2\). The four shell sums cost 16. Apart from \(\sqrt{\dim\pi_{u,\tau}}(2v_E)^{-rk}2^{r(k+1)+r(r-1)}\), the local factor is at most
\begin{equation}\label{eq:p8c1-type-shell}
 16(4r+4)(4r)^{4r^2}(2H_1)^{16r^3+4r}e_\tau^{8r^2+2}
 (8\pi e/h)^{jh}P_\tau^{-h}s_\tau^{-(q-16r^3-4r-1)/r}.
\end{equation}
The dimension and \(v_E\) factors here are the one-sided squared coefficient normalizations in the two section estimates. In the product of \eqref{eq:p8c1-type-shell}, \eqref{eq:p8c1-reciprocal-minors} costs \((r+1)(8r^2+2)\) powers of \(S_\delta\). Since
\[
 (r+1)(8r^2+2)+(16r^3+4r+1)/r\le2k,
\]
the net power is at most \(-C_sk\). The remaining factors of \(r\) and \(H_1\) cost at most \((8rH_1)^{80r^4f}\). With the product Lebesgue normalization above, the integral stabilizer covolume and the inverse unipotent covolume have factors, apart from their index weight, bounded respectively by
\begin{equation}\label{eq:p8c1-chart-volumes}
 v_E^{k(r-j)f}(8v_E)^{2(r+1)^2f}
 \quad\text{and}\quad (8v_E)^{2(r+1)^2f}.
\end{equation}
The \(k(r-j)\) directions per place are the definite components of the columns in the isotropic intersection. The remaining column components and the triangular and real coordinates are included in the \(2(r+1)^2\) exponents. These covolumes multiply the local quotient estimates with the same Haar normalization. The factor \((2v_E)^{-rkf}\) absorbs the positive \(k\)-power of \(v_E\) in \eqref{eq:p8c1-chart-volumes}. With \(1\le v_E\le b/2\), \(h\ge k/2\), and \(C\ge100(C_H+10)\), the remaining powers fit in \((C/k)^{jhf}H(C)^{Cr^4}\). The product of \(P_\tau^{-h}\) is \(D_\delta^{-h/2}\). Taking the infimum over adapted forms proves the amplitude inequality in \eqref{eq:p8c1-type-factors}.
\leanref{GlobalTheory.\allowbreak{}Amp.\allowbreak{}DegConv.\allowbreak{}deg\_\allowbreak{}core}{Development.\allowbreak{}AbelianVarietyTheory.\allowbreak{}Proofs.\allowbreak{}FlagMeasures.\allowbreak{}AmplitudeDegenerateOperatorNormLeRealPlace}{165--258}
\leanref{GlobalTheory.\allowbreak{}TypeAsm.\allowbreak{}CompatibleChartOperations.\allowbreak{}prod\_\allowbreak{}degBd\_\allowbreak{}le\_\allowbreak{}explicit}{Development.\allowbreak{}AbelianVarietyTheory.\allowbreak{}Proofs.\allowbreak{}FlagMeasures.\allowbreak{}LocalBoundProductOverPlacesLabel}{21--88}
\leanref{GlobalTheory.\allowbreak{}TypeAsm.\allowbreak{}AmpTerms.\allowbreak{}amplitudeBound\_\allowbreak{}of\_\allowbreak{}localData\_\allowbreak{}of\_\allowbreak{}tailBound}{Development.\allowbreak{}AbelianVarietyTheory.\allowbreak{}Proofs.\allowbreak{}FlagMeasures.\allowbreak{}AmplitudeBoundAtTermsOfLocalData}{113--226}

For the second inequality in \eqref{eq:p8c1-type-factors} put \(N_T=T_u^f\), \(g=g_0g_1\) and \(c=1+\log_2(2f)\). A pair of terms gives two embedded lattice data on the definite rank-\(k\) quotient, with congruence graphs and a relative isometry. Their indices are \(1\le x\le D_\delta\) and \(1\le x'\le D_{\delta'}\). Put \(D_b=\prod_{1\le l\le j,\,\delta_l\ne0}d_l\).

\begin{proposition}[Weighted fibre over an embedded datum]\label{B8.02b-fibre}
For a fixed row, type \((j,\delta)\) and embedded datum of index \(x\), the sum of relative inverse stabilizer weights of the inequivalent moving frames is bounded by
\begin{equation}\label{eq:p8c1-frame-count}
 F_\delta(x)=(D_\delta/x)^{k/2}8^{r+1}f^2D_\delta^{3c+6}N_T^2
 (D_bD_\delta N_T)^{10000(r+1)^3}.
\end{equation}
\end{proposition}
\leanref{GlobalTheory.\allowbreak{}DegAsm.\allowbreak{}sum\_\allowbreak{}weights\_\allowbreak{}per\_\allowbreak{}datum}{Development.\allowbreak{}AbelianVarietyTheory.\allowbreak{}Proofs.\allowbreak{}FlagMeasures.\allowbreak{}FlagFrameVectorMemCommonSubspace}{204--271}
\leanref{GlobalTheory.\allowbreak{}TypeAsm.\allowbreak{}perRowC}{Development.\allowbreak{}AbelianVarietyTheory.\allowbreak{}Proofs.\allowbreak{}FlagMeasures.\allowbreak{}TermPairsRowLatticeCountedViaEmbeddings}{114--212}

Fix the signed block matrix and the norm-one scalar first. For the matrix there are at most \(8^j\) choices. The block lengths, the possible roots of the block determinants, and the column signs each cost at most \(2^j\). The scalar has at most \(8f^2D_\delta^{2(2+c)}\) choices by the denominator-ideal count. For each denominator the norm-one relation determines the numerator ideal, and generators with the same principal ideal differ by a norm-one unit, hence by a root of unity.

For one pair of values take a common normal form. Each adapted pair is obtained from it by a unipotent map \(\sigma_i\). Transport the row lattice by \(\sigma_i^{-1}\). The transported lattices retain the prescribed definite lattice and congruence graph. An isometry \(\rho\) in the common unipotent stabilizer between two transported lattices gives \(\sigma_{i'}\rho\sigma_i^{-1}\) on the original lattice carrying the graded frames. The frame classes therefore inject into the lattice orbits for the common normal form. The frame weight is bounded by the inverse stabilizer covolume of its lattice with that normal form.

There are at most \(N_T^2\) possibilities for the second inner lattice after fixing the definite trace. Indeed its extension in the direction of the last vector is determined by a class in the dual of the fixed trace lattice modulo that lattice, of cardinality at most \(T_u^{2f}\). With both inner lattices fixed, group the lattice orbits by the two indices \(i_1,i_2\) of the definite and terminal lattice inclusions. The discriminant relation gives
\[
 x i_1^2 i_2^2\le D_\delta .
\]
Each index is a positive integer at most \(D_\delta\), so there are at most \(D_\delta^2\) pairs. At fixed indices the orbit sum is at most
\[
 (i_1i_2)^k i_1^c(D_bD_\delta N_T)^{10000(r+1)^3}.
\]
We describe the representatives and their weights. Write \(V_0=X+X'+Ez+Ez'\), \(W_{\rm sp}=V_0^\perp\), and \(I=X\cap X'\). In the fixed normal form let \(W_\theta\) be the positive lift of \(W_{\rm sp}/I\), let \(\Pi\) be its perpendicular, and let \(e_a^*\), \(j\le a<r\), be the isotropic dual vectors. They satisfy
\[
 H_E(e_b,e_a^*)=\delta_{ba},\quad
 H_E(e'_b,e_a^*)=0\ (b<j),\quad
 H_E(z,e_a^*)=0,\quad H_E(e_a^*,e_b^*)=0.
\]
Thus \(W_{\rm sp}=W_\theta\oplus I\), \(\Pi=V_0+\operatorname{span}_E\{e_a^*:j\le a<r\}\), and \(\dim_E\Pi=2r+1\). Here \(\mathfrak a_a=\mathfrak d^{-(r-a)}\) for \(0\le a<r\). A framed lattice has leading-coefficient ideal \(\mathfrak a_a\) on the \(a\)-th isotropic graded line and \(\mathcal O_E\) on the final line. Equivalently, its leading-coefficient module at \(e_a\) is \(\{t\in E:te_a\in\Lambda+X_a\}\), and likewise for the moving frame. The positive lift is specifically
\[
 W_\theta=(V_0+\operatorname{span}_E\{e_a^*:j\le a<r\})^\perp.
\]
\leanref{GlobalTheory.\allowbreak{}DegCnt.\allowbreak{}NF.\allowbreak{}aa}{Development.\allowbreak{}AbelianVarietyTheory.\allowbreak{}Definitions.\allowbreak{}SemilocalFrames.\allowbreak{}NormalFormDualOrthogonalLast}{62--63}
\leanref{GlobalTheory.\allowbreak{}HermLat.\allowbreak{}IsFrame}{Development.\allowbreak{}AbelianVarietyTheory.\allowbreak{}Definitions.\allowbreak{}SemilocalFrames.\allowbreak{}FlagPairDefinitePartLattice}{131--133}
\leanref{GlobalTheory.\allowbreak{}DegCnt.\allowbreak{}NF.\allowbreak{}Wt}{Development.\allowbreak{}AbelianVarietyTheory.\allowbreak{}Definitions.\allowbreak{}SemilocalFrames.\allowbreak{}NormalFormDualOrthogonalLast}{66--67}
The reference perpendicular lattice is
\[
 P_0=\sum_{0\le a<r}\mathfrak a_ae_a
       +\sum_{0\le a<j}\mathfrak a_ae'_a+\mathcal O_Ez
       +\sum_{j\le a<r}\mathfrak d^{r-a}e_a^*.
\]
A term lattice is unimodular, has both prescribed framed structures, and is admissible at level \(g\ge1\): if \(x'\in\Lambda\) has leading coefficient \(t\) at level \(a\) of the moving flag, then some \(x\in\Lambda\) has the same coefficient at level \(a\) of the fixed flag and satisfies \(x-x'\in g\Lambda\). Its split condition is
\[
 \Theta(\Lambda):=(\Lambda\cap W_{\rm sp})+I
       =(\Lambda\cap W_\theta)+I.
\]
Let \(R\) be the common flag-unipotent stabilizer and let \(R_s\) be its subgroup fixing \(W_\theta\) pointwise. These groups preserve admissibility, the embedded datum, \(\Theta(\Lambda)\), and the terminal index \(i_2\).
\leanref{GlobalTheory.\allowbreak{}DegCnt.\allowbreak{}NF}{Development.\allowbreak{}AbelianVarietyTheory.\allowbreak{}Definitions.\allowbreak{}SemilocalFrames.\allowbreak{}NormalFormDualOrthogonalLast}{33--57}
\leanref{GlobalTheory.\allowbreak{}T.\allowbreak{}hermFramePairNormalForm\_\allowbreak{}geometry\_\allowbreak{}summary}{Development.\allowbreak{}AbelianVarietyTheory.\allowbreak{}Proofs.\allowbreak{}FlagMeasures.\allowbreak{}HermFramePairNormalFormGeometrySummaryStatement}{11--13}
\leanref{GlobalTheory.\allowbreak{}DegCnt.\allowbreak{}NF.\allowbreak{}P0}{Development.\allowbreak{}AbelianVarietyTheory.\allowbreak{}Definitions.\allowbreak{}SemilocalFrames.\allowbreak{}NormalFormDualOrthogonalLast}{72--76}
\leanref{GlobalTheory.\allowbreak{}DegCnt.\allowbreak{}NF.\allowbreak{}IsTerm}{Development.\allowbreak{}AbelianVarietyTheory.\allowbreak{}Definitions.\allowbreak{}SemilocalFrames.\allowbreak{}NormalFormDualOrthogonalLast}{101--103}
\leanref{GlobalTheory.\allowbreak{}HermLat.\allowbreak{}IsAdm}{Development.\allowbreak{}AbelianVarietyTheory.\allowbreak{}Definitions.\allowbreak{}SemilocalFrames.\allowbreak{}FlagPairDefinitePartLattice}{134--137}
\leanref{GlobalTheory.\allowbreak{}DegCnt.\allowbreak{}NF.\allowbreak{}Split}{Development.\allowbreak{}AbelianVarietyTheory.\allowbreak{}Definitions.\allowbreak{}SemilocalFrames.\allowbreak{}NormalFormDualOrthogonalLast}{119--121}
\leanref{GlobalTheory.\allowbreak{}DegCnt.\allowbreak{}SplitLift.\allowbreak{}latTh\_\allowbreak{}eq}{Development.\allowbreak{}AbelianVarietyTheory.\allowbreak{}Proofs.\allowbreak{}FlagMeasures.\allowbreak{}HermFramePairExistsUnipotentSplitLattice}{130--147}
\leanref{GlobalTheory.\allowbreak{}DegCnt.\allowbreak{}NF.\allowbreak{}Rs}{Development.\allowbreak{}AbelianVarietyTheory.\allowbreak{}Definitions.\allowbreak{}SemilocalFrames.\allowbreak{}NormalFormDualOrthogonalLast}{97--98}
\leanref{GlobalTheory.\allowbreak{}DegCnt.\allowbreak{}NF.\allowbreak{}bI}{Development.\allowbreak{}AbelianVarietyTheory.\allowbreak{}Definitions.\allowbreak{}SemilocalFrames.\allowbreak{}NormalFormDualOrthogonalLast}{108--109}
\leanref{GlobalTheory.\allowbreak{}T.\allowbreak{}commonUnipotent\_\allowbreak{}stabiliser\_\allowbreak{}preserves\_\allowbreak{}terms\_\allowbreak{}and\_\allowbreak{}datum}{Development.\allowbreak{}AbelianVarietyTheory.\allowbreak{}Proofs.\allowbreak{}FlagMeasures.\allowbreak{}CommonUnipotentStabiliserPreservesTermsAndDatumStatement}{11--13}

First obtain a split representative. The saturated inclusion \(\Lambda\cap I\subset\Lambda\cap W_{\rm sp}\) has a finitely generated torsion-free quotient over the Dedekind domain \(\mathcal O_E\). That quotient is projective. Choose a splitting and let \(W\) be the \(E\)-span of its image. Then
\[
 W\cap I=0,\qquad W_{\rm sp}=W+I,
 \qquad
 \forall l\in\Lambda\cap W_{\rm sp}\ \exists w\in\Lambda\cap W: l-w\in I.
\]
To construct the required common unipotent map, write \(W_\theta\) as the graph \(w-\phi(w)\) of a linear map \(\phi:W\to I\). Choose a basis \(x_b\) of \(I\). The form on \(W\) is anisotropic, hence nondegenerate, so there are \(v_b\in W\) with
\[
 \phi(w)=-\sum_bH_E(w,v_b)x_b.
\]
For \(x\in I\), \(v\in W_{\rm sp}\), and \(c+\bar c=-H_E(v,v)\), the Eichler map
\[
 E_{x,v,c}(y)
 =y+H_E(y,v)x-H_E(y,x)v+cH_E(y,x)x
\]
is an isometry, with inverse \(E_{x,-v,\bar c}\). It fixes \(V_0\) pointwise and belongs to both flag-unipotent groups. On \(W_{\rm sp}\) it sends \(y\) to \(y+H_E(y,v)x\). Choose such a \(c\) for each pair \((x_b,v_b)\) and compose these maps. Since \(I\) is orthogonal to \(W_{\rm sp}\), their restrictions add the displayed corrections. The resulting \(s\in R\) fixes \(V_0\) and sends \(w\) to \(w-\phi(w)\), so \(sW=W_\theta\).
\leanref{Formalism.\allowbreak{}HermitianLatticeCounting.\allowbreak{}P.\allowbreak{}liftTransitive\_\allowbreak{}of}{Formalism.\allowbreak{}FlagUnipotentLattices.\allowbreak{}MinkowskiFactorialBounds.\allowbreak{}ExistsHasLeadingCoeffOfMemSpanBelow}{1635--1710}
\leanref{Formalism.\allowbreak{}HermitianLatticeCounting.\allowbreak{}P.\allowbreak{}good\_\allowbreak{}family}{Formalism.\allowbreak{}FlagUnipotentLattices.\allowbreak{}MinkowskiFactorialBounds.\allowbreak{}ExistsHasLeadingCoeffOfMemSpanBelow}{1571--1596} For \(l=w+i\) as above, \(sl=sw+i\) with \(sw\in s\Lambda\cap W_\theta\); hence \(s\Lambda\) is split.
\leanref{Formalism.\allowbreak{}HermitianLatticeCounting.\allowbreak{}P.\allowbreak{}complementLiftTransitive}{Formalism.\allowbreak{}FlagUnipotentLattices.\allowbreak{}MinkowskiFactorialBounds.\allowbreak{}ExistsHasLeadingCoeffOfMemSpanBelow}{1711--1712}
\leanref{GlobalTheory.\allowbreak{}DegCnt.\allowbreak{}SplitLift.\allowbreak{}main}{Development.\allowbreak{}AbelianVarietyTheory.\allowbreak{}Proofs.\allowbreak{}FlagMeasures.\allowbreak{}HermFramePairExistsUnipotentSplitLattice}{148--187}

For \(l\le j\), choose leading coefficients \(\alpha_a,\beta_a\in\mathfrak d^{-(r-a)}\), \(0\le a<l\), and lift \(\alpha_ae_a\) and \(\beta_ae'_a\) to vectors of \(\Lambda\) modulo the respective preceding flag spaces. The determinant of their pairings is
\[
 \left(\prod_{a<l}\alpha_a\right)p_l
 \overline{\left(\prod_{a<l}\beta_a\right)}.
\]
It is integral because \(\Lambda\) is unimodular. Products of these coefficients generate \(\mathfrak d^{-2q_l}\); conjugation preserves \(\mathfrak d\). Multilinearity therefore gives \((p_l)\mathfrak d^{-2q_l}\subset\mathcal O_E\).
\leanref{GlobalTheory.\allowbreak{}DegCnt.\allowbreak{}Absorb.\allowbreak{}Dl\_\allowbreak{}mul\_\allowbreak{}mem\_\allowbreak{}one\_\allowbreak{}of\_\allowbreak{}mem\_\allowbreak{}ipow}{Development.\allowbreak{}AbelianVarietyTheory.\allowbreak{}Proofs.\allowbreak{}FlagMeasures.\allowbreak{}ExistsStarStableIdealAbsorbingMinors}{309--349} Lemma~\ref{B8.02bh} supplies the same ideal \(\mathfrak c\) for all terms in the normal form, with \(N(\mathfrak c)=(D_bN_T)^2\). In particular \(D_b,D_\delta\ge1\).
\leanref{GlobalTheory.\allowbreak{}DegCnt.\allowbreak{}Absorb.\allowbreak{}exists\_\allowbreak{}ideal}{Development.\allowbreak{}AbelianVarietyTheory.\allowbreak{}Proofs.\allowbreak{}FlagMeasures.\allowbreak{}ExistsStarStableIdealAbsorbingMinors}{355--403}

Put \(e_0=2r+2\). Write the two flag lattices in unitriangular coordinates \(U,U'\). In each antidiagonal block of the reduced pairing, invert the corresponding diagonal block of \(U\); act on the opposite isotropic block by the inverse conjugate transpose with the pairing coefficients included. This defines an isometry \(s_1\in R_s\) fixing \(I+Ez+\operatorname{span}_E\{e_a^*\}\). If its two flag matrices are \(M,S\), then \(MU\) has no off-diagonal entries within one block. The integral mixed pairings and the nonzero leading block determinants give
\[
 \mathfrak c(MU)_{ba}\mathfrak a_a\subset\mathfrak a_b,\qquad
 \mathfrak c^2(SU')_{ba}\mathfrak a_a\subset\mathfrak a_b.
\]
The first estimate uses one cleared leading determinant; the second uses the two adjacent block ends and conjugation-stability of \(\mathfrak c\). Both matrices are unitriangular. Their finite inverse series and at most \(r+1\) filtration steps give denominator \((\mathfrak c^2)^{r+1}=\mathfrak c^{e_0}\) in both directions.
\leanref{GlobalTheory.\allowbreak{}DegCnt.\allowbreak{}PosJ.\allowbreak{}exists\_\allowbreak{}elt}{Development.\allowbreak{}AbelianVarietyTheory.\allowbreak{}Proofs.\allowbreak{}FlagMeasures.\allowbreak{}NormalFormFlagLatticesNearUnderStabiliser}{620--722}
\leanref{GlobalTheory.\allowbreak{}DegCnt.\allowbreak{}PosJ.\allowbreak{}denom}{Development.\allowbreak{}AbelianVarietyTheory.\allowbreak{}Proofs.\allowbreak{}FlagMeasures.\allowbreak{}NormalFormFlagLatticesNearUnderStabiliser}{802--872}
\leanref{GlobalTheory.\allowbreak{}DegCnt.\allowbreak{}PosJ.\allowbreak{}main}{Development.\allowbreak{}AbelianVarietyTheory.\allowbreak{}Proofs.\allowbreak{}FlagMeasures.\allowbreak{}NormalFormFlagLatticesNearUnderStabiliser}{875--920}
Thus, for \(\Lambda_1=s_1\Lambda\) and \(Y=X_j,X'_j\),
\[
 \mathfrak c^{e_0}(P_0\cap Y)\subset\Lambda_1,
 \qquad \mathfrak c^{e_0}(\Lambda_1\cap Y)\subset P_0.
\]
For the next reduction, choose adapted vectors \(v_a=e_a+\xi_a\), with \(\xi_a\in I\cap X_a\), for the lattice filtration in \(I\). Successive shears, taken from the highest index downwards, send \(v_a\) to \(e_a\), fix \(z\), and preserve the corrections already made. An Eichler translation removes the remaining \(I\)-component of an adapted last vector while fixing \(X\). Both operations fix \(X_j+X'_j\) and \(W_\theta\).
\leanref{GlobalTheory.\allowbreak{}DegCnt.\allowbreak{}PosIso.\allowbreak{}exists\_\allowbreak{}unitri}{Development.\allowbreak{}AbelianVarietyTheory.\allowbreak{}Proofs.\allowbreak{}FlagMeasures.\allowbreak{}NormalFormIdealDivisibility}{452--489}
\leanref{GlobalTheory.\allowbreak{}DegCnt.\allowbreak{}PosIso.\allowbreak{}exists\_\allowbreak{}move\_\allowbreak{}alt}{Development.\allowbreak{}AbelianVarietyTheory.\allowbreak{}Proofs.\allowbreak{}FlagMeasures.\allowbreak{}NormalFormIdealDivisibility}{435--451}
\leanref{GlobalTheory.\allowbreak{}DegCnt.\allowbreak{}PosIso.\allowbreak{}normal\_\allowbreak{}position}{Development.\allowbreak{}AbelianVarietyTheory.\allowbreak{}Proofs.\allowbreak{}FlagMeasures.\allowbreak{}NormalFormIdealDivisibility}{559--637}
Induction over the leading coefficients, using the clearing inclusion for \(p_j\), gives an isometry \(s_2\in R_s\) such that, with \(\Lambda_2=s_2\Lambda_1\),
\[
 \mathfrak c^2(P_0\cap I)\subset\Lambda_2,
 \quad \Lambda_2\cap I\subset P_0,
 \quad \mathfrak c\mathcal O_Ez\subset\Lambda_2.
\]
Both lattices remain split. The second map preserves their intersections with \(X_j,X'_j\), so the first pair of inclusions also holds for \(\Lambda_2\). Since the radical of \(V_0\) is \(I\), expanding the displayed generators of \(P_0\) gives
\[
 P_0\cap V_0\subset(P_0\cap X_j)+(P_0\cap X'_j)+(P_0\cap I)+\mathcal O_Ez.
\]
Indeed an isotropic-dual component lying in \(V_0\) pairs trivially with \(I\), which forces all its coefficients to vanish. Since \(e_0\ge2\), the preceding inclusions imply \(\mathfrak c^{e_0}(P_0\cap V_0)\subset\Lambda_2\). The skew reduction, applied with this exponent and the upper bounds on \(\Lambda_2\cap X_j\), \(\Lambda_2\cap X'_j\), and \(\Lambda_2\cap I\), gives \(s_3\in R_s\) and
\[
 \mathfrak c^{6e_0+30}P_0\subset s_3\Lambda_2\cap\Pi,
 \qquad
 \mathfrak c^{6e_0+30}(s_3\Lambda_2\cap\Pi)\subset P_0.
\]
The lifting step in this last reduction uses
\[
 P_*=\sum_{a\ge j}\mathfrak d^{r-a}e_a^*,\qquad
 P_0=(P_0\cap V_0)+P_*.
\]
The cleared discriminant modules satisfy
\[
 \mathfrak c^3\Theta(\Lambda_2)^\#\subset\Theta(\Lambda_2),\qquad
 \mathfrak cP_0^\#\subset P_0,\qquad
 \mathfrak cP_0\subset P_0^\#.
\]
Here, for an \(\mathcal O_E\)-submodule \(M\), write \(M^\#=\{y\in\operatorname{span}_E M:H_E(y,M)\subset\mathcal O_E\}\); this definition includes the radical of the restricted form. Unimodularity and \(\Lambda_2\cap I\subset P_0\) give an \(E\)-linear lift \(\Psi:V\to V\) with \(\Psi(\mathfrak cP_*)\subset\Lambda_2\) and \(H_E(\Psi y-y,e_a)=0\) for \(y\in P_*\), \(a\ge j\). Let \(\pi\) be projection to \(\Pi\) along \(W_\theta\). The split condition and the first discriminant inclusion show that a projection to \(\Pi\) of a lattice vector differs by a \(W_\theta\)-vector with denominator cleared by \(\mathfrak c^3\). Consequently
\[
 \mathfrak c^4\pi\Psi(P_*)\subset\Lambda_2,\qquad
 H_E(\pi\Psi y-y,e_a)=0.
\]
The skew-adjustment estimate applied to this lift and to \(\mathfrak c^{e_0}(P_0\cap V_0)\subset\Lambda_2\) gives the displayed exponent \(6e_0+30\). Its trace correction uses \(\vartheta\in\mathcal O_E\) with \(\vartheta+\bar\vartheta=1\), available because the CM extension is unramified. The correction fixes \(V_0\) and \(W_\theta\), and its displacement lies in \(I\). In particular \(s_3\) lies in the subgroup fixing both spaces, and the lattice \(s_3\Lambda_2\) remains split.
\leanref{GlobalTheory.\allowbreak{}HermLat.\allowbreak{}mem\_\allowbreak{}dual\_\allowbreak{}iff}{Development.\allowbreak{}AbelianVarietyTheory.\allowbreak{}Definitions.\allowbreak{}SemilocalFrames.\allowbreak{}FlagPairDefinitePartLattice}{38--43}
\leanref{GlobalTheory.\allowbreak{}DegCnt.\allowbreak{}DualPos.\allowbreak{}lift\_\allowbreak{}bd}{Development.\allowbreak{}AbelianVarietyTheory.\allowbreak{}Proofs.\allowbreak{}FlagMeasures.\allowbreak{}HermFramePairSplitTermNearReferenceBySkew}{400--434}
\leanref{GlobalTheory.\allowbreak{}DegCnt.\allowbreak{}DualPos.\allowbreak{}near\_\allowbreak{}of\_\allowbreak{}skew}{Development.\allowbreak{}AbelianVarietyTheory.\allowbreak{}Proofs.\allowbreak{}FlagMeasures.\allowbreak{}HermFramePairSplitTermNearReferenceBySkew}{509--583} As \(6e_0+30=12r+42\le12r+50\), the same inclusions hold with \(e=12r+50\). Thus every orbit has a split representative whose perpendicular lattice \(P\) satisfies \(\mathfrak c^eP_0\subset P\subset\mathfrak c^{-e}P_0\).
\leanref{GlobalTheory.\allowbreak{}T.\allowbreak{}normalForm\_\allowbreak{}flag\_\allowbreak{}lattices\_\allowbreak{}near\_\allowbreak{}under\_\allowbreak{}stabiliser}{Development.\allowbreak{}AbelianVarietyTheory.\allowbreak{}Proofs.\allowbreak{}FlagMeasures.\allowbreak{}NormalFormFlagLatticesNearUnderStabiliserStatement}{11--13}
\leanref{GlobalTheory.\allowbreak{}T.\allowbreak{}normal\_\allowbreak{}form\_\allowbreak{}ideal\_\allowbreak{}divisibility}{Development.\allowbreak{}AbelianVarietyTheory.\allowbreak{}Proofs.\allowbreak{}FlagMeasures.\allowbreak{}NormalFormIdealDivisibilityStatement}{11--13}
\leanref{GlobalTheory.\allowbreak{}DegCnt.\allowbreak{}NormalPos.\allowbreak{}P0\_\allowbreak{}inf\_\allowbreak{}V0\_\allowbreak{}le\_\allowbreak{}sup\_\allowbreak{}of\_\allowbreak{}radical}{Development.\allowbreak{}AbelianVarietyTheory.\allowbreak{}Proofs.\allowbreak{}FlagMeasures.\allowbreak{}HermFramePairNormalFormExistsConjNearReference}{89--128}
\leanref{GlobalTheory.\allowbreak{}T.\allowbreak{}hermFramePair\_\allowbreak{}split\_\allowbreak{}term\_\allowbreak{}near\_\allowbreak{}reference\_\allowbreak{}by\_\allowbreak{}skew}{Development.\allowbreak{}AbelianVarietyTheory.\allowbreak{}Proofs.\allowbreak{}FlagMeasures.\allowbreak{}HermFramePairSplitTermNearReferenceBySkewStatement}{11--13}
\leanref{GlobalTheory.\allowbreak{}DegCnt.\allowbreak{}NF.\allowbreak{}Near}{Development.\allowbreak{}AbelianVarietyTheory.\allowbreak{}Definitions.\allowbreak{}SemilocalFrames.\allowbreak{}NormalFormDualOrthogonalLast}{122--124}
\leanref{GlobalTheory.\allowbreak{}DegCnt.\allowbreak{}NormalPos.\allowbreak{}exists\_\allowbreak{}near}{Development.\allowbreak{}AbelianVarietyTheory.\allowbreak{}Proofs.\allowbreak{}FlagMeasures.\allowbreak{}HermFramePairNormalFormExistsConjNearReference}{143--175} The definite sublattice of index \(i_1\) has at most \(i_1^{k+c}\) choices. For fixed intersections \(A\) with the definite space and \(P\) with its perpendicular, the number of unimodular lattices inducing these intersections is at most \(|A^\#/A|^{2r+1}\). The definite intersection satisfies \(|A^\#/A|\le(D_\delta N_T)^2\). Hence there are at most \((D_\delta N_T)^{2(2r+1)}\) gluings.
\leanref{GlobalTheory.\allowbreak{}T.\allowbreak{}encard\_\allowbreak{}unimodular\_\allowbreak{}hermGluings\_\allowbreak{}le}{Development.\allowbreak{}AbelianVarietyTheory.\allowbreak{}Proofs.\allowbreak{}FlagMeasures.\allowbreak{}EncardUnimodularHermGluingsLeStatement}{11--13}
\leanref{GlobalTheory.\allowbreak{}DegCnt.\allowbreak{}OmCount.\allowbreak{}idx\_\allowbreak{}dual\_\allowbreak{}le}{Development.\allowbreak{}AbelianVarietyTheory.\allowbreak{}Proofs.\allowbreak{}FlagMeasures.\allowbreak{}HermFramePairOmSumSameDatumLe}{188--234} The number of perpendicular lattices is at most \(N(\mathfrak c)^{2e(2r+1)^2}\), by counting submodules between \(\mathfrak c^eP_0\) and \(\mathfrak c^{-e}P_0\).

For the weight, let \(L=(X+X')^\perp\cap(\Lambda+X)\). The fixed normal form satisfies \(e'_a=ce_a\) for \(j\le a\le r\); in particular \(z'=cz\), so \(z\) is orthogonal to \(X'\). Its reference completion is exactly
\[
 \Lambda_{\rm ref}=P_0+\bigl(L\cap(W_\theta+Ez)\bigr).
\]
\leanref{GlobalTheory.\allowbreak{}DegCnt.\allowbreak{}NF.\allowbreak{}Ref}{Development.\allowbreak{}AbelianVarietyTheory.\allowbreak{}Definitions.\allowbreak{}SemilocalFrames.\allowbreak{}NormalFormDualOrthogonalLast}{77--78} For any lattice \(A\), write \(\Gamma(A)=\{s\in R:sA=A\}\). The relative inverse covolume is normalized by
\[
 \omega(\Lambda)=\operatorname{iv}(\Gamma(\Lambda_{\rm ref}),\Gamma(\Lambda)),
 \qquad
 \operatorname{iv}(A,B)=\frac{[A:A\cap B]}{[B:A\cap B]}.
\]
The two stabilizers are commensurable because both lattices are full. The same holds after intersection with a subgroup of \(R\): its stabilizer acts on a finite set of lattices between two common scalar bounds. Conjugation by \(s\in R\) preserves this weight. To see the index assertion, put \(d=\operatorname{iv}(\Gamma(\Lambda),\Gamma(s\Lambda))>0\). Multiplicativity of relative indices gives \(\operatorname{iv}(\Gamma(\Lambda),\Gamma(s^t\Lambda))=d^t\). Since \((s-1)^{2r+1}=0\), the binomial expansions of \(s^t\) and \(s^{-t}\) have a common denominator independent of \(t\in\N\). The lattices \(s^t\Lambda\) therefore lie between two fixed lattices, and two coincide. It follows that \(d=1\). The embedded datum, and hence its reference completion, is preserved by \(R\).
\leanref{GlobalTheory.\allowbreak{}DegCnt.\allowbreak{}NF.\allowbreak{}om}{Development.\allowbreak{}AbelianVarietyTheory.\allowbreak{}Definitions.\allowbreak{}SemilocalFrames.\allowbreak{}NormalFormDualOrthogonalLast}{110--111}
\leanref{GlobalTheory.\allowbreak{}HL.\allowbreak{}invVol}{Development.\allowbreak{}AbelianVarietyTheory.\allowbreak{}Definitions.\allowbreak{}Foundations.\allowbreak{}IsHermND}{86--88}
\leanref{GlobalTheory.\allowbreak{}DegCnt.\allowbreak{}Unimod.\allowbreak{}invVol\_\allowbreak{}smul}{Development.\allowbreak{}AbelianVarietyTheory.\allowbreak{}Proofs.\allowbreak{}FlagMeasures.\allowbreak{}RelIndexNeZero}{169--212}

Use
\[
 R_1=\{s\in R:s|_{V_0}=1\},\qquad
 Z=\{s\in R_1:s|_{W_\theta}=1\}.
\]
They are the normal subgroups of the common stabilizer's skew, middle, and top filtration.
\leanref{GlobalTheory.\allowbreak{}DegCnt.\allowbreak{}NF.\allowbreak{}R1}{Development.\allowbreak{}AbelianVarietyTheory.\allowbreak{}Definitions.\allowbreak{}SemilocalFrames.\allowbreak{}NormalFormDualOrthogonalLast}{93--94}
\leanref{GlobalTheory.\allowbreak{}DegCnt.\allowbreak{}NF.\allowbreak{}ZI}{Development.\allowbreak{}AbelianVarietyTheory.\allowbreak{}Definitions.\allowbreak{}SemilocalFrames.\allowbreak{}NormalFormDualOrthogonalLast}{95--96} Put \(A=\Gamma(\Lambda_{\rm ref})\), \(B=\Gamma(\Lambda)\). Normality and multiplication of finite relative indices give
\[
 \operatorname{iv}(A,B)=
 \operatorname{iv}(A\cap Z,B\cap Z)\,
 \operatorname{iv}((A\cap R_1)Z,(B\cap R_1)Z)\,
 \operatorname{iv}(AR_1,BR_1).
\]
For every split term with \(\mathfrak c^eP_0\subset\Lambda\cap\Pi\) and \(\mathfrak c^e(\Lambda\cap\Pi)\subset P_0\), the two outer factors are bounded by finite congruence actions. First \(A\cap Z\) preserves \(P_0=\Lambda_{\rm ref}\cap\Pi\). An element acting trivially on \(P_0/\mathfrak c^{3e+10}P_0\) preserves \(\Lambda\), by the deep-congruence estimate. Thus it belongs to \(B\cap Z\). For a projective lattice of rank \(d\), the ring of endomorphisms modulo \(\mathfrak c^t\) has cardinality \(N(\mathfrak c)^{td^2}\), as is seen after localization at the primes dividing \(\mathfrak c\). Therefore the first factor, which is at most \([A\cap Z:A\cap B\cap Z]\), is bounded by \(N(\mathfrak c)^{(3e+10)(2r+1)^2}\).

For the top factor, \(AR_1\) preserves \(P_0\cap V_0\). If \(s\in AR_1\) is the identity modulo \(\mathfrak c^{3e+11}\) on this lattice, the inclusion \(\mathfrak c\subset(T_u)\) permits a lift \(s'\in R_s\) agreeing with \(s\) on \(V_0\) and acting trivially modulo \(\mathfrak c^{3e+10}\) on \(P_0\). The deep-congruence estimate puts \(s'\) in \(B\); since \(s'^{-1}s\in R_1\), it follows that \(s\in BR_1\). The same finite-module count, with \(\operatorname{rank}(P_0\cap V_0)\le2r+1\), bounds the top factor. Hence both outer factors are at most
\[
 N(\mathfrak c)^{(3e+11)(2r+1)^2}.
\]
\leanref{GlobalTheory.\allowbreak{}DegCnt.\allowbreak{}WtSkew.\allowbreak{}invVol\_\allowbreak{}skew\_\allowbreak{}le}{Development.\allowbreak{}AbelianVarietyTheory.\allowbreak{}Proofs.\allowbreak{}FlagMeasures.\allowbreak{}HermFramePairNormalFormInvVolSkewLayerLe}{74--127}
\leanref{GlobalTheory.\allowbreak{}DegCnt.\allowbreak{}invVol\_\allowbreak{}top\_\allowbreak{}layer\_\allowbreak{}le}{Development.\allowbreak{}AbelianVarietyTheory.\allowbreak{}Proofs.\allowbreak{}FlagMeasures.\allowbreak{}HermFramePairNormalFormInvVolTopLayerLe}{70--143}
For the middle factor use the additive rows
\[
 \rho_a(s)=s(e_a^*)-e_a^*\pmod I,\qquad j\le a<r,\quad s\in R_1.
\]
Their common kernel is \(Z\). Put \(M=\Lambda\cap W_\theta\) and \(L_\theta=L\cap W_\theta\). The definite-index bound gives \([L_\theta:M]\le\eta i_1\), where \(\eta\) is the congruence-index factor. Filter the rows from the highest index downwards, and let
\[
 C_a=\{H_E(e_a^*,x):x\in\Lambda\cap I,\
              H_E(e_b^*,x)=0\text{ for }b>a\}.
\]
The flag-lattice filtration gives
\[
 \prod_{a=j}^{r-1}[\mathfrak a_a:C_a\cap\mathfrak a_a]\le i_2.
\]
Let \(Q_a\) be the possible \(a\)-th rows of \((B\cap R_1)Z\) with all higher rows zero, and put \(P_a=Q_a+\mathfrak a_aL_\theta\), with both groups viewed in \(V/I\) and the products interpreted as additive spans of scalar products. The integral Eichler transvection \(E_{x,-v,-\vartheta H_E(v,v)}\), for \(x\in\Lambda\cap I\) and \(v\in M\), preserves \(\Lambda\) and gives row \(H_E(e_a^*,x)v\). Thus \(Q_a\) contains \((C_a\cap\mathfrak a_a)M\). The rows of \((A\cap R_1)Z\) lie in \(P_a\), and
\[
 [P_a:Q_a]\le
 [\mathfrak a_a:C_a\cap\mathfrak a_a]^k[L_\theta:M].
\]
Successive passage to the row kernels bounds the relative inverse covolume by the product of these indices. Multiplying over the \(r-j\) rows proves that the middle factor is at most \(i_2^k(\eta i_1)^{r-j}\).
\leanref{GlobalTheory.\allowbreak{}HermLat.\allowbreak{}flagPairDefinitePartLattice}{Development.\allowbreak{}AbelianVarietyTheory.\allowbreak{}Definitions.\allowbreak{}SemilocalFrames.\allowbreak{}FlagPairDefinitePartLattice}{151--152}
\leanref{GlobalTheory.\allowbreak{}DegCnt.\allowbreak{}MidLayer.\allowbreak{}exists\_\allowbreak{}eichler}{Development.\allowbreak{}AbelianVarietyTheory.\allowbreak{}Proofs.\allowbreak{}FlagMeasures.\allowbreak{}HaarInvVolumeMiddleLayerEichlerLe}{780--854}
\leanref{GlobalTheory.\allowbreak{}DegCnt.\allowbreak{}MidLayer.\allowbreak{}invVol\_\allowbreak{}le\_\allowbreak{}prod}{Development.\allowbreak{}AbelianVarietyTheory.\allowbreak{}Proofs.\allowbreak{}FlagMeasures.\allowbreak{}HaarInvVolumeMiddleLayerEichlerLe}{19--149}
\leanref{GlobalTheory.\allowbreak{}DegCnt.\allowbreak{}MidLayer.\allowbreak{}invVol\_\allowbreak{}mid\_\allowbreak{}layer\_\allowbreak{}le}{Development.\allowbreak{}AbelianVarietyTheory.\allowbreak{}Proofs.\allowbreak{}FlagMeasures.\allowbreak{}HaarInvVolumeMiddleLayerEichlerLe}{1008--1184} The definite-index relation gives \(1\le\eta i_1\le D_\delta\), so this factor is at most \(i_2^kD_\delta^r\). These estimates use the same nonzero conjugation-stable ideal \(\mathfrak c\subset(T_u)\) and both minor-clearing inclusions of Lemma~\ref{B8.02bh}. Multiplying the factors proves
\[
 \omega(\Lambda)\le i_2^kD_\delta^r
       N(\mathfrak c)^{2(3e+11)(2r+1)^2}.
\]
\leanref{GlobalTheory.\allowbreak{}DegCnt.\allowbreak{}Weight.\allowbreak{}invVol\_\allowbreak{}three}{Development.\allowbreak{}AbelianVarietyTheory.\allowbreak{}Proofs.\allowbreak{}FlagMeasures.\allowbreak{}SplitTermWeightLeNearReferenceLattice}{20--30}
\leanref{GlobalTheory.\allowbreak{}T.\allowbreak{}haarInvVolume\_\allowbreak{}middleLayer\_\allowbreak{}eichler\_\allowbreak{}le}{Development.\allowbreak{}AbelianVarietyTheory.\allowbreak{}Proofs.\allowbreak{}FlagMeasures.\allowbreak{}HaarInvVolumeMiddleLayerEichlerLeStatement}{11--13}
\leanref{GlobalTheory.\allowbreak{}T.\allowbreak{}hermFramePairNormalForm\_\allowbreak{}invVol\_\allowbreak{}skew\_\allowbreak{}layer\_\allowbreak{}le}{Development.\allowbreak{}AbelianVarietyTheory.\allowbreak{}Proofs.\allowbreak{}FlagMeasures.\allowbreak{}HermFramePairNormalFormInvVolSkewLayerLeStatement}{11--13}
\leanref{GlobalTheory.\allowbreak{}T.\allowbreak{}hermFramePairNormalForm\_\allowbreak{}invVol\_\allowbreak{}top\_\allowbreak{}layer\_\allowbreak{}le}{Development.\allowbreak{}AbelianVarietyTheory.\allowbreak{}Proofs.\allowbreak{}FlagMeasures.\allowbreak{}HermFramePairNormalFormInvVolTopLayerLeStatement}{11--13}
\leanref{GlobalTheory.\allowbreak{}T.\allowbreak{}flagPairDefiniteThetaLattice\_\allowbreak{}le\_\allowbreak{}fixed\_\allowbreak{}index\_\allowbreak{}bound}{Development.\allowbreak{}AbelianVarietyTheory.\allowbreak{}Proofs.\allowbreak{}FlagMeasures.\allowbreak{}FlagPairDefiniteThetaLatticeLeFixedIndexBoundStatement}{11--13}
\leanref{GlobalTheory.\allowbreak{}DegCnt.\allowbreak{}Weight.\allowbreak{}om\_\allowbreak{}le}{Development.\allowbreak{}AbelianVarietyTheory.\allowbreak{}Proofs.\allowbreak{}FlagMeasures.\allowbreak{}SplitTermWeightLeNearReferenceLattice}{35--77}

Put \(e=12r+50\), \(d_0=2r+1\), \(n=N(\mathfrak c)\), and \(M=D_bD_\delta N_T\). Multiplication of the number of representatives by the bound for each relative inverse stabilizer covolume gives
\[
 i_1^{k+c}n^{2ed_0^2}(D_\delta N_T)^{2d_0}
 \bigl(i_2^kD_\delta^r n^{2(3e+11)d_0^2}\bigr).
\]
Since \(D_b,D_\delta,N_T\ge1\), one has \(n\le M^2\), \(D_\delta N_T\le M\), and \(D_\delta\le M\). Thus this product is at most
\[
 (i_1i_2)^ki_1^cM^{r+4(4e+11)d_0^2+2d_0}
 \le(i_1i_2)^ki_1^cM^{10000(r+1)^3}.
\]
Indeed the first exponent equals \(768r^3+4144r^2+3573r+846\).
\leanref{GlobalTheory.\allowbreak{}DegCnt.\allowbreak{}OmCount.\allowbreak{}numeric}{Development.\allowbreak{}AbelianVarietyTheory.\allowbreak{}Proofs.\allowbreak{}FlagMeasures.\allowbreak{}HermFramePairOmSumSameDatumLe}{122--148} Now \((i_1i_2)^k\le(D_\delta/x)^{k/2}\) and \(i_1^c\le D_\delta^c\). Summing over the index pairs and inner lattices gives the bound for one matrix and scalar. Multiplying by their counts, with \(8^j\le8^r\), gives \eqref{eq:p8c1-frame-count}.
\leanref{GlobalTheory.\allowbreak{}DegAsm.\allowbreak{}FixedForm.\allowbreak{}main}{Development.\allowbreak{}AbelianVarietyTheory.\allowbreak{}Proofs.\allowbreak{}FlagMeasures.\allowbreak{}SumWeightsInequivalentHermFrameTermsLe}{44--101}
\leanref{GlobalTheory.\allowbreak{}DegCnt.\allowbreak{}omSum\_\allowbreak{}sameDat\_\allowbreak{}le}{Development.\allowbreak{}AbelianVarietyTheory.\allowbreak{}Proofs.\allowbreak{}FlagMeasures.\allowbreak{}HermFrameWeightSumSameDatumLe}{48--117}
\leanref{GlobalTheory.\allowbreak{}DegCnt.\allowbreak{}OmCount.\allowbreak{}main}{Development.\allowbreak{}AbelianVarietyTheory.\allowbreak{}Proofs.\allowbreak{}FlagMeasures.\allowbreak{}HermFramePairOmSumSameDatumLe}{237--318}

Put \(D_0=N_T^2,\ X_*=xx',\ Y_*=X_*D_0^2\), and \(M_*=(20f(2v_E)^2)^{2f(6+c)}\).

\begin{proposition}[Mass of triples of embedded data]\label{B8.02b-mass}
On the genus of the definite rank-\(k\) level lattice, weight each class by the reciprocal order of its isometry group. Divide the weighted count of triples of two embedded data of indices \(x,x'\) and their relative isometry by the genus mass. The result is at most
\begin{equation}\label{eq:p8c1-genus-triples}
 X_*^{k/2}Y_*^{6r+10+4c}
 \left[1+M_*(2v_E)^f\sum_{s=1}^{k-1}
 \big((32/g)^{kf}(4kf)^{4f}Y_*^7\big)^s+
 \big((2/g)^{kf}(4kf)^{4f}D_0X_*\big)^k\right].
\end{equation}
\end{proposition}
\leanref{GlobalTheory.\allowbreak{}DegCnt.\allowbreak{}PairMean.\allowbreak{}main}{Development.\allowbreak{}AbelianVarietyTheory.\allowbreak{}Proofs.\allowbreak{}FlagMeasures.\allowbreak{}DiscNResNeZero}{312--429}

\begin{proof}
Take \(T=T_u\), \(D=D_0=T^{2f}\), \(X=X_*\), \(Y=Y_*\), and \(e=3(2r+1)\). By \ref{B8.02c}, \(T\) splits completely in \(E\), \(g\ge2\) and \(v_E\ge1\). The row space is totally positive definite of rank \(k\ge m+3\). To identify its integral lattice, pass first to the perpendicular of the isotropic flag modulo that flag. The induced lattice is self-dual: restriction of integral functionals to the saturated isotropic lattice is surjective by projectivity, and dualizing the resulting perpendicular quotient identifies it with its dual. The last frame vector spans a primitive copy of \(Oz\) in this quotient and satisfies \(h(z,z)=-T\). Its perpendicular is the row lattice \(L_0\). In a self-dual lattice the projections onto a primitive sublattice and its perpendicular have images their respective duals. The quotient by the direct sum thus projects isomorphically onto both discriminant modules. It follows that \(L_0^\#/L_0\simeq (Oz)^\#/Oz\simeq O/(T)\).
\leanref{GlobalTheory.\allowbreak{}DegCnt.\allowbreak{}PrimPerp.\allowbreak{}perp\_\allowbreak{}prim}{Development.\allowbreak{}AbelianVarietyTheory.\allowbreak{}Proofs.\allowbreak{}FlagMeasures.\allowbreak{}FlagPairDefinitePartLatticeDiscriminantEqQuotien}{36--142}
\leanref{GlobalTheory.\allowbreak{}TypeAsm.\allowbreak{}lat0C}{Development.\allowbreak{}AbelianVarietyTheory.\allowbreak{}Proofs.\allowbreak{}FlagMeasures.\allowbreak{}BaseLatticeConstructionHolds}{163--184}
\leanref{GlobalTheory.\allowbreak{}PairsRow.\allowbreak{}main}{Development.\allowbreak{}AbelianVarietyTheory.\allowbreak{}Proofs.\allowbreak{}FlagMeasures.\allowbreak{}GenusMeanEmbeddedPairsLeMassMulPairBound}{13--37}

Count the triples of \ref{B8.02be} by their fixed space \(U\). For \(0<s=\operatorname{codim}U<k\), fix \(a=d(L\cap U)\). By \ref{B8.02bf}, a nonempty fibre has \(1\le a\le A_s\le DX\). Multiply its bound in \ref{B8.02bf} by the genus mean of \(N_{k-s,a}\) from \ref{B8.02bd}. The factors with exponent \(k/2\) give
\[
 16^{skf}(a/D)^{k/2}(A_s/a)^{k/2}
 =X^{k/2}(32/g)^{skf}.
\]
Since \(aD\le Y\) and \(X\le Y\), the remaining discriminant and graph factors satisfy
\[
 (aD)^{5s+5+2c}X^{2s+2c+1}Y^{3(2r+1)}
 \le Y^{\,7s+6r+9+4c}.
\]
There are at most \(DX\le Y\) values of \(a\), and \((4sf)^{4sf}\le(4kf)^{4sf}\). Thus the contribution at \(s\), per unit genus mass, is at most
\[
 X^{k/2}Y^{6r+10+4c}M_*\rho^f
 \left((32/g)^{kf}(4kf)^{4f}Y^7\right)^s .
\]

For \(U=V\) there is one possible space, \(a=D\), and \(A_0/a=X\). The first bound of \ref{B8.02bf} is \(X^{k/2}X^{2c+1}Y^e\le X^{k/2}Y^{6r+10+4c}\). For \(U=0\) the second bound of \ref{B8.02bf} applies, with \(c+e\le6r+10+4c\). Adding these endpoints to the sum over \(1\le s<k\) gives \eqref{eq:p8c1-genus-triples}.
\end{proof}



In the product of \eqref{eq:p8c1-genus-triples} and \(F_\delta(x)F_{\delta'}(x')\), the indices cancel to \((D_\delta D_{\delta'})^{k/2}\). Put \(C_a=\max(1,2(C_H+1),C_T)\), \(U_1=H(C_a)S_\delta^2\) and \(U_2=H(C_a)S_{\delta'}^2\). We have \(D_\delta,N_T\le U_1,\ D_b\le U_1^r\), and the corresponding bounds at \(U_2\). The thresholds give \(c\le4r\) and \(42C_am^2\log b\le k\). Using \(f\le k\) and the lower bound on \(g\) in \ref{B8.02c} gives
\[
 (32/g)^{kf}(4kf)^{4f}H(C_a)^{42}\le1/2.
\]
As \(Y_*\le H(C_a)^6(S_\delta S_{\delta'})^2\), the \(s\)-th power in the bracket of \eqref{eq:p8c1-genus-triples} is at most \(2^{-s}(S_\delta S_{\delta'})^{14k}\). The zero-fixed-space factor obeys the same comparison at \(s=k\). The prefactor \(Y_*^{6r+10+4c}\) and \(2+2M_*(2v_E)^f\) then cost at most \((U_1U_2)^{96r+183}\). The remaining factors of \eqref{eq:p8c1-frame-count} cost \(U_i^{250000r^4}\) per side and the index sum costs \(U_1U_2\). Take \(C_h=300000C_a\), \(C_p=16\), and \(k\ge300000r^4\). For \(S_i=S_\delta,S_{\delta'}\), the combined \(U_i\) powers and \(S_i^{14k}\) fit in \(H(C_h)^{C_hr^4}S_i^{16k}\). Sum over row representatives to recover \(\operatorname{mass}_u\) and obtain the second inequality of \eqref{eq:p8c1-type-factors}. Multiplication proves the proposition.
\leanref{GlobalTheory.\allowbreak{}TypeAsm.\allowbreak{}pair\_\allowbreak{}count\_\allowbreak{}bound\_\allowbreak{}of\_\allowbreak{}siegel\_\allowbreak{}mass}{Development.\allowbreak{}AbelianVarietyTheory.\allowbreak{}Proofs.\allowbreak{}FlagMeasures.\allowbreak{}TermPairCountBoundOfSiegelMassRatios}{70--162}
\leanref{GlobalTheory.\allowbreak{}DegAsm.\allowbreak{}sum\_\allowbreak{}pairBd\_\allowbreak{}mul\_\allowbreak{}datBd\_\allowbreak{}le}{Development.\allowbreak{}AbelianVarietyTheory.\allowbreak{}Proofs.\allowbreak{}FlagMeasures.\allowbreak{}SumPairBdMulDatBdLe}{307--392}
The complete proof is contained in \Lean{GlobalTheory.\allowbreak{}rowsPos\_\allowbreak{}typeP}.
\end{proof}

\subsection*{Proof of \ref{B8.02yaa}}\label{prf:B8.02yaa}\phantomsection
\begin{proof}
Choose \(\xi\) with \(A(x_1)\ne0\). The bound on \(a(x)\) gives absolute convergence of \(B(x)\). Unitarity gives \(\|A(x)\|\le\|\xi\|a(x)\). The invariance of \(a\) therefore gives
\[
 \|A(x)\|^2
 \le\|\xi\|^2\int_{U\times N}|m(anx)|a(anx)\,da\,dn.
\]
Integrating and applying Tonelli's theorem and the quotient integral formula yields
\[
 \mathcal N
 \le c_S\|\xi\|^2\int_G|m(y)|a(y)\,dy
 \le c_SC\|\xi\|^2\int_G|m(y)|\|y\|_{\rm op}^{M}\,dy<\infty.
\]
Choose \(\epsilon>0\) and a relatively compact open neighborhood \(O\) of \(x_1\) on which \(\|A\|>\epsilon\). Put
\[
 q_O(x)=\int_{U\times N}{\bf1}_O(anx)\,da\,dn.
\]
This is a nonnegative measurable function on the quotient, and
\[
 \int q_O\,d\nu=c_S\operatorname{vol}_G(O)>0.
\]
Thus \(P=\{x:q_O(x)>0\}\) has positive \(\nu\)-measure. For \(x\in P\), some \(anx\) belongs to \(O\). Equivariance and unitarity imply \(\|A(x)\|=\|A(anx)\|>\epsilon\). Hence \(\mathcal N\ge\epsilon^2\nu(P)>0\). Multiplication by \(\mathcal N^{-1/2}\) proves the normalization and preserves equivariance and the convolution identity.
\end{proof}

\subsection*{Proof of \ref{B8.02yab}}\label{prf:B8.02yab}\phantomsection
\begin{proof}
Write \(\Phi(x)=\ell(\sigma(x)w)\). Fix \(x_0\in G\) and \(\eta\in V\). Left translation in the reproducing integral, followed by the quotient integral formula, gives
\[
 \begin{aligned}
 \langle\eta,W(x_0)\rangle
 &=d\int_G
       \overline{\langle\sigma(x_0)w,\sigma(y)w\rangle}
       \langle\eta,W(y)\rangle\,dy\\
 &=d\int_{S\backslash G}
       \bigl\langle
          \mathcal P(\sigma(x)w,\sigma(x_0)w)\eta,W(x)
       \bigr\rangle\,d\nu(x).
 \end{aligned}
\]
Absolute convergence of the reproducing integral justifies the change of variables and Fubini's theorem. In the inner integral use the equivariance of \(W\), the unitarity of \(\pi\), and conjugate-linearity in the first argument of the inner product. The rank-one period identity then makes the last integrand
\[
 \langle\eta,\Phi(x_0)\rangle\langle\Phi(x),W(x)\rangle.
\]
If \(\Phi\) is zero, this formula gives \(W=0\), and \(\kappa=0\) works. If \(\Phi(x_*)\ne0\), take \(\eta=\Phi(x_*)\) and \(x_0=x_*\). The first factor is \(\|\Phi(x_*)\|^2>0\), so the preceding absolute convergence proves that \(x\mapsto\langle\Phi(x),W(x)\rangle\) is integrable. Set
\[
 \kappa=d\int_{S\backslash G}\langle\Phi(x),W(x)\rangle\,d\nu(x).
\]
The formula for arbitrary \(x_0,\eta\) gives
\(\langle\eta,W(x_0)\rangle=\langle\eta,\kappa\Phi(x_0)\rangle\), hence \(W=\kappa\Phi\).

If \(W\) is nonzero, choose \(x_0\) with \(\Phi(x_0)\ne0\) and put \(\xi=\Phi(x_0)\). The period identity gives \(A(x)=\|\xi\|^2\Phi(x)\) and \(A(x_0)=\|\xi\|^2\xi\ne0\). Apply the stated norm identity to \(v=\sigma(x_0)w\) and divide by \(\|\xi\|^2\). This proves \(\int\|\Phi\|^2\,d\nu=1/d\). The formulas for the norm of \(A\) and for \(W^0\) follow.
\end{proof}

\subsection*{Proof of \ref{B8.02ya}}\label{prf:B8.02ya}\phantomsection
\begin{proof}[Proof (outline)]
By \ref{B8.53a} and the Cartan integration formula on \(G_v^0\simeq\U(p_v,q_v)\), the spectral gap \(b_{\tau_v,j+1}-b_{\tau_v,j}\ge16n^{19}\) at \(n\ge1801\) implies that \(y\mapsto|\langle\sigma_v(x_0)w_v,\sigma_v(y)w_v\rangle|\,\|y\|_{\rm op}^{N_v}\) belongs to \(L^1(G_v^0)\). Parametrizing \(N_{u,v}^0\) in the graded skew-Hermitian Lie-algebra coordinates of the isotropic flag of length \(r_S\) and orthogonal negative-norm vector \(z_{u,v}\), and integrating over the compact group \(U_{u,v}^0\simeq\U(k)\), yields locally uniform integrable domination in \(y\) and the nilpotent chart bound \(a_v(y;x_0)\le C_v(x_0)\|y\|_{\rm op}^{N_v}\). Dominated convergence and right invariance of the product Haar measure \(da\,dn\) on the unimodular group \(S_{u,v}^0=U_{u,v}^0\ltimes N_{u,v}^0\) give continuity and \((S_{u,v}^0,\pi_{u,v}\otimes\psi_{u,v}^0)\)-equivariance of \(B_v(\cdot,x_0)\).

The discrete-series Bessel nonvanishing theorem applies to the even rank \(k\ge2\), the negative-norm frame vector, the nonzero character parameters, the regular Harish--Chandra parameters with the stated interlacing and parity, and the large-gap numerical data of \(X\). For the given unit lowest-type vector and the chosen Haar measures it gives \(x_1,x_0\in G_v^0\) such that
\[
 \int_{U_{u,v}^0\times N_{u,v}^0}
 \langle\sigma_v(x_0)w_v,\sigma_v(anx_1)w_v\rangle
 \overline{\psi_{u,v}^0(n)}\,\pi_{u,v}(a)^{-1}\,da\,dn\ne0
 \quad\text{in }\operatorname{End}_{\C}(V_{u,v}).
\]
This integral is \(B_v(x_1,x_0)\). Choose \(\xi\in V_{u,v}\) on which this nonzero operator is nonzero. This proves \textup{(2)}. The quotient integral formula here has \(c_S=1\) in Lemma~\ref{B8.02yaa}. Right invariance gives \(a_v(anx;x_0)=a_v(x;x_0)\). Apply that lemma with \(G=G_v^0\), \(S=S_{u,v}^0\), \(M=N_v\), and \(C=C_v(x_0)\), using the weighted coefficient integral from the first paragraph and the nonzero value just obtained. It gives
\[
 0<\int_{S_{u,v}^0\backslash G_v^0}\|A_v\|^2\,d\nu_{u,v}^0
 \le C_v(x_0)\|\xi\|^2
 \int_{G_v^0}
 |\langle\sigma_v(x_0)w_v,\sigma_v(y)w_v\rangle|
 \|y\|_{\rm op}^{N_v}\,dg_v^0(y)<\infty.
\]
Put \(\phi_v(g)=d_v\overline{\langle w_v,\sigma_v(g)w_v\rangle}\). For each \(x\in G_v^0\), the first paragraph gives
\[
 \int_{G_v^0}|\phi_v(g)|a_v(xg;x_0)\,dg_v^0(g)
 \le C_v(x_0)\|x\|_{\rm op}^{N_v}
 \int_{G_v^0}|\phi_v(g)|\|g\|_{\rm op}^{N_v}\,dg_v^0(g)<\infty.
\]
Thus Fubini's theorem applies to right convolution of \(A_v\) with \(\phi_v\). Schur orthogonality and \(\|w_v\|=1\) give
\[
 d_v\int_{G_v^0}\overline{\langle w_v,\sigma_v(g)w_v\rangle}
       \sigma_v(g)w_v\,dg_v^0(g)=w_v.
\]
Applying \(\sigma_v(anx)\), taking the inner product with \(\sigma_v(x_0)w_v\), and integrating over \(a,n\) gives
\(\int_{G_v^0}\phi_v(g)A_v(xg)\,dg_v^0(g)=A_v(x)\).
The same reproducing identity holds for \(W_v^0=A_v/\|A_v\|_{L^2(\nu_{u,v}^0)}\).
\leanref{WlocSk.\allowbreak{}ANz.\allowbreak{}aNonzero\_\allowbreak{}of}{Formalism.\allowbreak{}LocalToGlobal.\allowbreak{}SiegelLatticeTheory.\allowbreak{}HaarTransversals.\allowbreak{}UnitaryHaarProducts.\allowbreak{}LocalWeightCharts.\allowbreak{}ANonzeroOf}{34--148}
\leanref{WlocSk.\allowbreak{}AQ.\allowbreak{}wave\_\allowbreak{}packet\_\allowbreak{}l2\_\allowbreak{}finite}{Formalism.\allowbreak{}LocalToGlobal.\allowbreak{}SiegelLatticeTheory.\allowbreak{}HaarTransversals.\allowbreak{}UnitaryHaarProducts.\allowbreak{}LocalWeightCharts.\allowbreak{}WavePacketL1L2OfHypotheses}{397--479}
\leanref{WlocSk.\allowbreak{}scWlocShape\_\allowbreak{}of}{Formalism.\allowbreak{}LocalToGlobal.\allowbreak{}SiegelLatticeTheory.\allowbreak{}HaarTransversals.\allowbreak{}UnitaryHaarProducts.\allowbreak{}LocalWeightCharts.\allowbreak{}L2SGWfEqOne}{87--125}

When \(W_u\) does not vanish identically, isolating the \(v\)-th tensor factor at a point where \(W_u\ne0\) and passing to \(G_v^0\) shows that the transported local factor \(W_v\) is continuous, \((S_{u,v}^0,\pi_{u,v}\otimes\psi_{u,v}^0)\)-equivariant, square-integrable on \(S_{u,v}^0\backslash G_v^0\), and fixed under right convolution by \(d_v\overline{\langle w_v,\sigma_v(\cdot)w_v\rangle}\). Double-coset and smooth globalization bounds under the gap \(16n^{19}\) attach to every such function on \(G_v^0\) a continuous \((S_{u,v}^0,\pi_{u,v}\otimes\psi_{u,v}^0)\)-equivariant linear functional \(\mathcal B:\sigma_v^\infty\to V_{u,v}\) on the Fr\'echet space \(\sigma_v^\infty\) of smooth vectors satisfying \(W_v(x)=\mathcal B(\sigma_v(x)w_v)\). Applying this realization to \(W_v\) and to \(W_v^0\) gives continuous Bessel intertwiners \(\mathcal B,\mathcal B_0:\sigma_v^\infty\to V_{u,v}\) with \(\mathcal B_0\ne0\) since \(\|W_v^0\|_{L^2(\nu_{u,v}^0)}=1\). By the archimedean Bessel multiplicity-one theorem of Jiang, Sun and Zhu \cite[Theorem~A]{jiang-sun-zhu-2010}, the space \(\operatorname{Hom}_{S_{u,v}^0}(\sigma_v^\infty,V_{u,v}\otimes\psi_{u,v}^0)\) has dimension at most one; the archive proves the uniqueness it needs itself, under a spectral gap of at least \(120n^3\), which the gap \(16n^{19}\) supplies.
\leanref{CompanionTheory.\allowbreak{}Lime.\allowbreak{}BM.\allowbreak{}UniqGap}{CompanionTheory.\allowbreak{}AutomorphicTheory.\allowbreak{}LocalBesselUniqueness.\allowbreak{}BesselModels.\allowbreak{}AtFormAdmOneOfSteps}{15--27}
\leanref{CompanionTheory.\allowbreak{}Bmu.\allowbreak{}besselMapUniqGap}{CompanionTheory.\allowbreak{}AutomorphicTheory.\allowbreak{}BesselOrbitBounds.\allowbreak{}BesselMapUniqueness.\allowbreak{}BesselMapUniqGapAlt}{33--35}
 Hence \(\mathcal B\) is a scalar multiple of \(\mathcal B_0\), and \(W_v\) is the same scalar multiple of \(W_v^0\).
\leanref{Formalism.\allowbreak{}MainNotions.\allowbreak{}BridgeCS3dBOW.\allowbreak{}besselOfW\_\allowbreak{}holds}{CompanionTheory.\allowbreak{}AutomorphicTheory.\allowbreak{}LocalBesselUniqueness.\allowbreak{}BesselModels.\allowbreak{}DoubleCosetBesselBoundOfGapConstantLe}{22--32}
\leanref{Formalism.\allowbreak{}L2.\allowbreak{}SG.\allowbreak{}Pf.\allowbreak{}WuDec.\allowbreak{}BM.\allowbreak{}besselUniq\_\allowbreak{}core}{CompanionTheory.\allowbreak{}AutomorphicTheory.\allowbreak{}LocalBesselUniqueness.\allowbreak{}BesselModels.\allowbreak{}LocalUniquenessReduction}{18--67}
The complete nonvanishing and local proportionality proofs are contained in \Lean{ReciprocityTheory.\allowbreak{}ChartRestriction.\allowbreak{}DS.\allowbreak{}Lap.\allowbreak{}TopClosed.\allowbreak{}besselNVG} and \Lean{CompanionTheory.\allowbreak{}Lime.\allowbreak{}BM.\allowbreak{}wuDecayOfGap\_\allowbreak{}holds'}.
\end{proof}

\subsection*{Proof of \ref{B8.02y}}\label{prf:B8.02y}\phantomsection
\begin{proof}
At each infinite place \(v\) of \(F\), let \((p_v,q_v)\) be the signature of \(G_v\simeq\U(p_v,q_v)\) (\(p_v+q_v=n\ge1801\)), put \(q=\min(p_v,q_v)\), and let \((\sigma_v,\mathcal H_v)\) be the local discrete series representation of formal degree \(d_v>0\) with maximal compact subgroup \(K_v\le G_v\) and unit lowest-\(K_v\)-type vector \(w_v\in\mathcal H_v\). If \(q=0\), then \(G_v\) is compact and the continuous function \(g\mapsto|\langle w_v,\sigma_v(g)w_v\rangle|\,\|g\|_{\rm op}^N\) is integrable on \(G_v\). Suppose \(q>0\). Write \(g=k_1\exp(Y)k_2\) in the Cartan decomposition of \(G_v\), with \(k_1,k_2\in K_v\), singular coordinates \(a_1\ge\cdots\ge a_q\ge0\), and \(d=\|Y\|_{\rm HS}\). Up to the fixed conjugation identifying \(G_v\) with \(\U(p_v,q_v)\) in diagonal signature coordinates, \(k_1,k_2\) are Euclidean unitary and \(\|\exp(Y)\|_{\rm op}=e^{a_1}\le e^d\), so \(\|g\|_{\rm op}^N\le C_{0,v}e^{Nd}\) for a constant \(C_{0,v}>0\) independent of \(g\). By \ref{B8.53a}, for the spectral gap \(h\ge16n^{19}\) we have
\[
 |\langle w_v,\sigma_v(g)w_v\rangle|\le
 \exp\!\bigl[-(h/8)\min(d^2,d/\sqrt{2q})\bigr].
\]
In the Cartan integral formula the radial density is a constant multiple of
\[
 \prod_{1\le i<j\le q}\sinh^2(a_i-a_j)\sinh^2(a_i+a_j)
 \prod_{i=1}^q\sinh^{2|p_v-q_v|}(a_i)\sinh(2a_i).
\]
For \(1\le i<j\le q\) we have \(a_i\ge a_j\ge0\) and \(2(a_i-a_j)+2(a_i+a_j)=4a_i\). Since \(|\sinh t|\le e^{|t|}\) and \(|p_v-q_v|=n-2q\), the \(i\)-th coordinate contributes at most \(\exp\bigl((4(q-i)+2(n-2q)+2)a_i\bigr)=\exp\bigl((2n+2-4i)a_i\bigr)\). Because \(q\le n/2\), each coefficient \(2n+2-4i\ge2\) is positive, and \(a_i\le d\) gives \(\sum_{i=1}^q(2n+2-4i)a_i\le 2q(n-q)d\le2nqd\). For \(d>1\), since \(q\ge1\) gives \(\sqrt{2q}>1\), we have \(\min(d^2,d/\sqrt{2q})=d/\sqrt{2q}\); moreover \(1\le q\le q_v\le\max_{v'}q_{v'}\le n\) and \(n\ge1801\) give
\[
 \frac{h}{8\sqrt{2q}}\ge\frac{16n^{19}}{8\sqrt{2n}}=\sqrt2\,n^{37/2}
 >4\sqrt2\,n^{5/2}+2n^2\ge N+2nq.
\]
The region \(d\le1\) is compact, and on \(d>1\) the integral of \(e^{-(h/(8\sqrt{2q})-N-2nq)d}\) over the \(q\)-dimensional Cartan subspace converges. Therefore
\begin{equation}\label{eq:p8c1-weighted-coefficient}
 \int_{G_v}|\langle w_v,\sigma_v(g)w_v\rangle|\,\|g\|_{\rm op}^N\,dg_v<\infty.
\end{equation}
For any \(x_0\in G_v\), substituting \(g=x_0^{-1}y\) and using left Haar invariance with \(\|y\|_{\rm op}\le\|x_0\|_{\rm op}\|g\|_{\rm op}\) gives \(\int_{G_v}|\langle\sigma_v(x_0)w_v,\sigma_v(y)w_v\rangle|\,\|y\|_{\rm op}^N\,dg_v(y)<\infty\), and the same holds in Cayley coordinates \(G_v^0\). Every \(g\in G_v\) preserves a nondegenerate Hermitian form, so \(1=|\det g|\le\|g\|_{\rm op}^n\) and \(\|g\|_{\rm op}^N\ge1\). Thus \(\phi_v(g)=d_v\overline{\langle w_v,\sigma_v(g)w_v\rangle}\) lies in \(L^1(G_v)\) for each \(v\), and Tonelli's theorem gives \(\phi=\prod_v\phi_v\in L^1(G_\infty)\).
\leanref{Formalism.\allowbreak{}L2.\allowbreak{}SG.\allowbreak{}SplitPrimeNormLemmas.\allowbreak{}CoefWeightedLoc}{Formalism.\allowbreak{}LocalToGlobal.\allowbreak{}SiegelLatticeTheory.\allowbreak{}HaarTransversals.\allowbreak{}UnitaryHaarProducts.\allowbreak{}SettingBlocks}{323--334}
\leanref{Formalism.\allowbreak{}L2.\allowbreak{}SG.\allowbreak{}Pf.\allowbreak{}WuDec.\allowbreak{}phInt\_\allowbreak{}of\_\allowbreak{}gap}{Formalism.\allowbreak{}LocalToGlobal.\allowbreak{}SiegelLatticeTheory.\allowbreak{}FlagStabilizerTensors.\allowbreak{}RootLatticeThresholds.\allowbreak{}WuDecayAll}{136--139}
\leanref{ReciprocityTheory.\allowbreak{}ChartRestriction.\allowbreak{}DS.\allowbreak{}coefWeightedN}{ReciprocityTheory.\allowbreak{}ReductiveGroups.\allowbreak{}DiscreteSeries.\allowbreak{}FormalDegree.\allowbreak{}CoefWeightedG}{14--63}

Fix a word \(u\). If \(W_u\) vanishes identically on \(G_\infty\), then \(\int\|W_u(x)\|L_u(x)^N\,d\nu_u(x)=0<\infty\), and \(a_u=0\) on the compact set \(\{1\}\) satisfies all assertions. Assume that \(W_u\) does not vanish identically on \(G_\infty\). In the spectral realization \(X\), we have \(S_u=\prod_v S_{u,v}\) inside \(G_\infty=\prod_v G_v\), and \(W_u(x)=\bigotimes_v W_{u,v}(x_v)\) for continuous local functions \(W_{u,v}:G_v\to V_{u,v}\) under a map \((y_v)_v\mapsto\bigotimes_v y_v\) whose inner products factor as \(\langle\bigotimes_v y_v,\bigotimes_v y'_v\rangle=\prod_v\langle y_v,y'_v\rangle\) and whose image spans \(V_{\pi_u}\). This inner-product identity makes \((y_v)_v\mapsto\bigotimes_v y_v\) multilinear and gives \(\|W_u(x)\|=\prod_v\|W_{u,v}(x_v)\|\). Choose local Haar measures on the unimodular stabilizers \(S_{u,v}\) and local quotient measures \(\nu_{u,v}\) on \(S_{u,v}\backslash G_v\) disintegrating \(dg_v\). By uniqueness of Haar measure on \(S_u\) and Weil's quotient integral formula on \(S_u\backslash G_\infty\simeq\prod_v(S_{u,v}\backslash G_v)\), we have \(\nu_u=C_u\prod_v\nu_{u,v}\) for a finite constant \(C_u\ge0\).

Fix \(v\) and pass to the Cayley coordinates \(G_v^0,S_{u,v}^0=U_{u,v}^0\ltimes N_{u,v}^0,\nu_{u,v}^0\) of Proposition~\ref{B8.02ya} via the fixed conjugation matrix \(M_v\in\GL_n(\C)\). By Proposition~\ref{B8.02ya}\textup{(3)}, the transported local factor \(W_v\) satisfies \(W_v=\kappa_v W_v^0\) for a scalar \(\kappa_v\in\C\), where \(W_v^0=A_v/\|A_v\|_{L^2(\nu_{u,v}^0)}\) and \(A_v(x)=B_v(x,x_0)\xi\) from \eqref{eq:p8c1-bessel-average} satisfies \(\|A_v(x)\|\le\|\xi\|\,a_v(x;x_0)\). Put \(L_{u,v}^0(x)=\inf_{s\in S_{u,v}^0}\|sx\|_{\rm op}\). Since \(L_{u,v}^0(x)=L_{u,v}^0(anx)\le\|anx\|_{\rm op}\) for \((a,n)\in U_{u,v}^0\times N_{u,v}^0\),
\[
 \|A_v(x)\|\,L_{u,v}^0(x)^N\le\|\xi\|\int_{U_{u,v}^0\times N_{u,v}^0}
 |\langle\sigma_v(x_0)w_v,\sigma_v(anx)w_v\rangle|\,\|anx\|_{\rm op}^N\,da\,dn.
\]
Integrating over \(S_{u,v}^0\backslash G_v^0\) and unfolding to \(G_v^0\) bounds \(\int_{S_{u,v}^0\backslash G_v^0}\|A_v(x)\|\,L_{u,v}^0(x)^N\,d\nu_{u,v}^0(x)\) by a finite multiple of \(\int_{G_v^0}|\langle\sigma_v(x_0)w_v,\sigma_v(y)w_v\rangle|\,\|y\|_{\rm op}^N\,dg_v^0(y)\), which is finite by \eqref{eq:p8c1-weighted-coefficient}.

Let \(K_v^0\le G_v^0\) be the maximal compact subgroup with probability Haar measure \(dk\), and let \(d_\tau\) be the dimension of the lowest \(K_v^0\)-type of \(\sigma_v\). Because the lowest \(K_v^0\)-type has multiplicity one in \(\sigma_v\), Schur orthogonality on \(K_v^0\) gives \(\int_{K_v^0}d_\tau\overline{\langle w_v,\sigma_v(k)w_v\rangle}\,\sigma_v(k)v'\,dk=\langle w_v,v'\rangle w_v\) for all \(v'\in\mathcal H_v\). By strong continuity of \(\sigma_v\), choose a compact neighborhood \(U_1\subset G_v^0\) of \(1\) on which \(\|\sigma_v(y)w_v-w_v\|\le1/2\). Writing \(\mu_{G_v^0}\) for the Haar measure \(dg_v^0\), we have \(0<\mu_{G_v^0}(U_1)<\infty\), and \(v_1=\mu_{G_v^0}(U_1)^{-1}\int_{U_1}\sigma_v(y)w_v\,dg_v^0(y)\) satisfies \(|\langle w_v,v_1\rangle-1|\le1/2\), so \(c_1=\langle w_v,v_1\rangle\ne0\). Set
\[
 \alpha_v^0(g)=c_1^{-1}\mu_{G_v^0}(U_1)^{-1}\int_{K_v^0}
 d_\tau\overline{\langle w_v,\sigma_v(k)w_v\rangle}\,{\bf1}_{U_1}(k^{-1}g)\,dk.
\]
Then \(\alpha_v^0\) is measurable and bounded by \(|c_1|^{-1}\mu_{G_v^0}(U_1)^{-1}d_\tau\), has support in the compact set \(\Omega_v^0=K_v^0U_1\) containing \(1\), and satisfies \(\int_{G_v^0}\alpha_v^0(z)\sigma_v(z)w_v\,dg_v^0(z)=w_v\) by Fubini's theorem. For every \(x\in G_v^0\), Proposition~\ref{B8.02ya}\textup{(1)} gives
\[
 \int_{G_v^0}|\alpha_v^0(z)|\,a_v(xz;x_0)\,dg_v^0(z)
 \le C_v(x_0)\|x\|_{\rm op}^{N_v}
 \int_{\Omega_v^0}|\alpha_v^0(z)|\,\|z\|_{\rm op}^{N_v}\,dg_v^0(z)<\infty.
\]
By Tonelli's and Fubini's theorems applied to \eqref{eq:p8c1-bessel-average}, \(z\mapsto\alpha_v^0(z)A_v(xz)\) is Bochner integrable on \(G_v^0\) with integral \(A_v(x)\), and the same holds for \(W_v=\kappa_v W_v^0\). Conjugating back by \(M_v\) changes operator norms by at most \(\|M_v\|_{\rm op}\|M_v^{-1}\|_{\rm op}\) and scales the Haar and quotient measures by positive constants. Writing \(L_{u,v}(y)=\inf_{s\in S_{u,v}}\|sy\|_{\rm op}\), we obtain \(\int_{S_{u,v}\backslash G_v}\|W_{u,v}(y)\|\,L_{u,v}(y)^N\,d\nu_{u,v}(y)<\infty\) together with a bounded measurable function \(\alpha_v\) supported on a compact set \(\Omega_v\subset G_v\) containing \(1\) such that \(z_v\mapsto\alpha_v(z_v)W_{u,v}(yz_v)\) is Bochner integrable on \(G_v\) with integral \(W_{u,v}(y)\) for every \(y\in G_v\).

Since \(L_u(x)=\inf_{s\in S_u}\prod_v\|s_vx_v\|_{\rm op}\le\prod_v L_{u,v}(x_v)\), Tonelli's theorem gives
\[
 \int_{S_u\backslash G_\infty}\|W_u(x)\|\,L_u(x)^N\,d\nu_u(x)
 \le C_u\prod_v\int_{S_{u,v}\backslash G_v}\|W_{u,v}(y)\|\,L_{u,v}(y)^N\,d\nu_{u,v}(y)<\infty.
\]
Put \(a_u(z)=\prod_v\alpha_v(z_v)\) for \(z=(z_v)_v\in G_\infty\). Then \(a_u\) is bounded and measurable, with support in the compact set \(\Omega=\prod_v\Omega_v\) containing \(1\). Because \(a_u(z)W_u(xz)=\bigotimes_v\bigl(\alpha_v(z_v)W_{u,v}(x_vz_v)\bigr)\) and the image of \(\bigotimes_v\) spans \(V_{\pi_u}\), Fubini's theorem on \(G_\infty=\prod_v G_v\) yields absolute Bochner integrability of \(z\mapsto a_u(z)W_u(xz)\) and \(\int_{G_\infty}a_u(z)W_u(xz)\,dz=\bigotimes_v W_{u,v}(x_v)=W_u(x)\) for every \(x\in G_\infty\).
\leanref{Formalism.\allowbreak{}L2.\allowbreak{}SG.\allowbreak{}Pf.\allowbreak{}WuDec.\allowbreak{}QuotAll}{Formalism.\allowbreak{}LocalToGlobal.\allowbreak{}SiegelLatticeTheory.\allowbreak{}FlagStabilizerTensors.\allowbreak{}RootLatticeThresholds.\allowbreak{}WuDecayAll}{431--432}
\leanref{Formalism.\allowbreak{}L2.\allowbreak{}SG.\allowbreak{}Pf.\allowbreak{}WuDec.\allowbreak{}wuWeighted\_\allowbreak{}of\_\allowbreak{}loc}{Formalism.\allowbreak{}LocalToGlobal.\allowbreak{}SiegelLatticeTheory.\allowbreak{}FlagStabilizerTensors.\allowbreak{}RootLatticeThresholds.\allowbreak{}WuDecayAll}{467--490}
\leanref{Formalism.\allowbreak{}L2.\allowbreak{}SG.\allowbreak{}Pf.\allowbreak{}WuDec.\allowbreak{}wuReprod\_\allowbreak{}of\_\allowbreak{}loc}{Formalism.\allowbreak{}LocalToGlobal.\allowbreak{}SiegelLatticeTheory.\allowbreak{}FlagStabilizerTensors.\allowbreak{}RootLatticeThresholds.\allowbreak{}WuDecayAll}{491--519}
\end{proof}

\subsection*{Proof of \ref{B8.02a}}\label{prf:B8.02a}\phantomsection
\begin{proof}
At the preliminary point on a host, \ref{B5.22a} supplies two integer functions \(\alpha,\beta\) in its boxes whose prime sum is at most \(x^{-9}\). In the construction of \ref{B8.05} the same functions are prescribed at every translation of that host. By \ref{B8.02ac} they agree there with the member's top weights extended by zero, and the boxes and the prime sum agree with those at the preliminary point. In particular the member's own weights satisfy
\[
 \sum_{p\ {\rm candidate}}\frac{\log p}{p}
 \left(4\Sigma_p+2\mathbf1_{\{\Sigma_p>x^{-\varepsilon}\}}\right)
 \le x^{-9},
\]
where the candidate primes and shares are those of \ref{B5.22a}. This is \(R_W(i)\) in \eqref{5.19.1}. The final family cutoff in \ref{B8.05} already includes the preliminary choice cutoff, so the inequality holds at every indexed \(x>x_1\).
\end{proof}

\subsection*{Proof of \ref{B8.02ac}}\label{prf:B8.02ac}\phantomsection
\begin{proof}
The candidate tests of \ref{B3.51} agree by the common prime class, \(X'\), order of the Frobenius element \(h_p\) and band parameter \(\rho_B\). At a candidate prime let \(d_p=\operatorname{ord}(h_p)\). The frame identification preserves the least representative \(t_0\) of each coset of \(\langle h_p\rangle\), its marked positions \(S_{t_0}=\{0\le j<d_p:t_0+jh_p\in TT\}\), and the arrays \((\alpha(t_0+jh_p),\beta(t_0+jh_p))_{0\le j<d_p}\). These give the same window counts and good representatives. For \(p<p_0=k^2x^{26}\) the share is \(\min(1,K_{\kappa,II}|S_{t_0}|r_G^2kx^{-26})\), with \(r_G\) the local rank parameter of \ref{B3.51}. For \(p\ge p_0\) the arrays identify the preceding even mark and its weights, so \ref{B3.52a} gives equal \(D_I,c_6,t_I\) and equal constraints on the dominant base weights. The shares are the same infima because the nonnegative fractions passing the rank-one test at \((\lceil t_I\rceil_+,D_I,8(D_I+1))\) coincide for all the specified local rings. Summing over the good representatives and multiplying by \(m\) gives equal \(\Sigma_p\). The common \(s,n_H,k,\mathcal L\) give the same box endpoints on \(TT\), and outside \(TT\) the box is \(\{0\}\). The top-label identification transports the prescribed weights, hence also their extensions by zero.
\end{proof}

\subsection*{Proof of \ref{B8.03a}}\label{prf:B8.03a}\phantomsection
\begin{proof}
Choose a compact open subgroup \(K'\) fixing \(e\) by left translation. Such a subgroup exists by local constancy and compactness of the support. Let \(p=\phi\otimes{\bf1}_{K'}/\mu_f(K')\), where \(\mu_f\) is finite-adelic Haar measure. The formal-degree orthogonality relation, applied at each place with the two middle vectors equal to \(w_v\), gives \(\phi*\phi=\phi\). Hence \(p*(\phi\otimes e)=\phi\otimes e\). If \(h=R_eR_\phi h_0\) with \(h_0\in L^2(Y)\cap L^\infty(Y)\), convolution gives a bounded continuous representative of \(h\): its change under right translation by \(u\) is at most \(\|h_0\|_\infty\) times the \(L^1\) change of the translated kernel. This representative is right \(K'\)-invariant and satisfies \(h=R_ph=R_\phi h\).

We first prove the vanishing of the averages needed on the cusps. On a real component, let \(N\) be the unipotent radical of the stabilizer of a nonzero isotropic subspace. For a product of real groups take the product radical. We claim that \(R_\phi F=0\) for every bounded continuous left \(N\)-invariant function \(F\) on \(G_\infty\).

Choose a real factor with nonzero isotropic subspace \(W\). In a Witt decomposition let \(a(t)\) act by \(e^t\) on \(W\), by \(e^{-t}\) on its dual, and by the identity on the middle space and the other factors. Choose \(0\ne v_W\in\bigwedge^{\dim W}W\). For a maximal-compact-invariant Euclidean norm put
\[
 H(x)=-\frac{\log(\|x^{-1}v_W\|/\|v_W\|)}{\dim W}.
\]
Then \(H(mx)=H(x)\) for \(m\in N\), \(H(a(t)x)=t+H(x)\), and \(|H(xy)-H(x)|\le c(y)\) for a continuous \(c\), by the operator norms of \(y\) and \(y^{-1}\) on the exterior power. For every identity neighborhood \(U\), all sufficiently small \(t\) have the following property, uniformly in \(x\): some \(m\in N\) satisfies \(x^{-1}ma(t)x\in U\). Indeed write the chosen component of \(x\) as \(n l k\) in the parabolic decomposition, with \(l\) in the Levi subgroup and \(k\) in the maximal compact. Taking \(m=na(t)n^{-1}a(-t)\) reduces the conjugate to \(k^{-1}a(t)k\).

On bounded continuous functions invariant on the left by both \(N\) and \(a(\R)\), define
\[
 Q_sF(x)=\int\phi(y)F(xy)e^{is(H(xy)-H(x))}\,dy .
\]
The height identities preserve both invariances. The phases telescope in a double integral, so \(Q_s^2=Q_s\). The operators are norm continuous in \(s\), since their differences are bounded by the integral of
\(|\phi(y)|\min(2,|s-s'|c(y))\). For \(s>0\) large take \(t=\pi/s\) and \(u=x^{-1}ma(t)x\) as above. Translation by \(u\) in the integration variable reverses the sign of the phase and preserves the value of \(F\). Thus
\[
 2|Q_sF(x)|\le\|\phi-\phi(u\,\cdot)\|_1\|F\|_\infty .
\]
Continuity of translations in \(L^1\) gives \(\|Q_s\|\to0\). A nonzero idempotent has norm at least one, so the zero set of this continuous family of idempotents is open and closed in \(\R\). Hence \(Q_s=0\) for all \(s\).

For general bounded continuous left \(N\)-invariant \(F\), the function \(F_r(z)=F(a(r-H(z))z)\) has both invariances. Applying \(Q_s=0\) to \(F_r\), and using Fourier inversion for \(\rho\in C_c^\infty(\R)\), gives
\[
 \int \phi(y)F_r(xy)\rho(r-H(xy))\,dy=0.
\]
Integrate in \(r\) and substitute \(t=r-H(xy)\). Fubini applies with absolute integral at most \(\|\phi\|_1\|F\|_\infty\|\rho\|_1\). It follows that
\(\int\rho(t)(R_\phi F)(a(t)x)\,dt=0\).
The integrand without \(\rho\) is continuous, so its value at \(t=0\) vanishes. This proves the claim.
\leanref{CompanionTheory.\allowbreak{}D.\allowbreak{}gk.\allowbreak{}RF.\allowbreak{}selTw}{CompanionTheory.\allowbreak{}HarishChandraModules.\allowbreak{}KFiniteAnalyticVectors.\allowbreak{}FundamentalDomainCusps.\allowbreak{}IntegrableMulBc}{73--205}
\leanref{CompanionTheory.\allowbreak{}D.\allowbreak{}gk.\allowbreak{}RF.\allowbreak{}selP}{CompanionTheory.\allowbreak{}HarishChandraModules.\allowbreak{}KFiniteAnalyticVectors.\allowbreak{}FundamentalDomainCusps.\allowbreak{}IntegrableMulBc}{207--311}

Fix an identity neighborhood \(U\). There are finitely many components at level \(K'\) by \cite[Theorem~5.1]{borel-1963}, applied to \(\operatorname{Res}_{F/\Q}G\); passing between commensurable compact open levels preserves finiteness. Each component is a quotient of \(G_\infty\) by an arithmetic subgroup \(\Gamma'\). We describe the compact part and the cusp averages in one component. The group is reductive with anisotropic rational center, so the Siegel sets of \cite[Theorem~13.1]{borel-1969} have no unbounded split central coordinate. In left-quotient coordinates write a Siegel representative as \(\omega a k\), up to one of finitely many rational translates, with \(\omega\) in a fixed compact part of the minimal parabolic. Outside a compact set some simple root \(\alpha(a)\) is arbitrarily large. For the corresponding maximal rational parabolic let \(N\) be its real unipotent radical. The intersection \(\Gamma'\cap N\) is arithmetic in a unipotent group and has compact quotient by \cite[Theorem~8.7]{borel-1969}. Take bounded representatives for this quotient. Conjugation by \(\omega^{-1}\) keeps these representatives in a bounded part of \(N\). Each root in its Lie algebra contains \(\alpha\) with positive coefficient, and all simple roots of \(a\) are bounded below in the Siegel set. Thus conjugation by \(a^{-1}\) contracts the representatives into any given neighborhood as \(\alpha(a)\) increases. Conjugation by \(k\) preserves this conclusion uniformly. The finitely many rational translates give the same assertion with the conjugated parabolics and arithmetic subgroups.

Consequently, outside a compact part, the probability measure on \((\Gamma'\cap N)\backslash N\) can be represented by a measure \(\nu_x\) on elements \(m\) with \(x^{-1}mx\in U\). For the bounded representative above, \(F(y)=\int h(my)\,d\nu_x(m)\) is continuous and left \(N\)-invariant by right invariance of quotient Haar measure. Fubini and \(h=R_\phi h\) give \(F=R_\phi F\). The claim gives \(F=0\). On lifting to the adelic quotient, right \(K'\)-invariance gives
\[
 \int h(mxz)\,d\nu_x(m)=0\qquad(z\in G_\infty\times K').
\]
This proves the zero average for the image of each bounded input.
\leanref{CompanionTheory.\allowbreak{}D.\allowbreak{}gk.\allowbreak{}RF.\allowbreak{}cuspidal\_\allowbreak{}at\_\allowbreak{}infinity\_\allowbreak{}of}{CompanionTheory.\allowbreak{}HarishChandraModules.\allowbreak{}KFiniteAnalyticVectors.\allowbreak{}FundamentalDomainCusps.\allowbreak{}CuspidalAtInfinityOf}{14--67}
\leanref{CompanionTheory.\allowbreak{}D.\allowbreak{}gk.\allowbreak{}RF.\allowbreak{}adelic\_\allowbreak{}cusp\_\allowbreak{}of\_\allowbreak{}infinite\_\allowbreak{}cusp}{CompanionTheory.\allowbreak{}HarishChandraModules.\allowbreak{}KFiniteAnalyticVectors.\allowbreak{}FundamentalDomainCusps.\allowbreak{}PBAdel}{21--101}

Choose a continuous compactly supported kernel \(\alpha\) in \(G_\infty\times K'\) with \(\|p-\alpha\|_1\le1/8\). For \(u\) in a sufficiently small identity neighborhood,
\[
 |\alpha(u^{-1}z)-\alpha(z)|\le\eta{\bf1}_B(z),
 \qquad \eta\mu(B)\le1/8.
\]
Here \(B\) is a fixed compact set containing the nearby translates of the support, and uniform continuity permits the indicated choice of \(\eta\). On the cusp part just obtained for this neighborhood, the average of \(R_\alpha h(mx)\) vanishes. Subtracting it from \(R_\alpha h(x)\) bounds the latter in absolute value by \(R_{\eta{\bf1}_B}|h|(x)\).

On the remaining compact part choose a compact set \(C\) of representatives. On \(C\), the functions \(R_\alpha h\) with \(\|h\|_2\le1\) are uniformly equicontinuous. Indeed a difference at nearby points is bounded by the sup norm of the difference of the translated kernels times \(\int_{CB}|h|\), after enlarging \(B\) if necessary. The latter integral is uniformly bounded by Cauchy--Schwarz because a compact set has bounded multiplicity over the quotient. Choose finitely many test points \(y_1,\ldots,y_d\) so that vanishing of all \(R_\alpha h(y_j)\) bounds the integral of \(|R_\alpha h|^2\) on the compact part by \(\|h\|_2^2/8\).

Write \(S=\|h\|_2^2\). Right translation is unitary on the quotient, hence \(\|R_qh\|_2\le\|q\|_1\|h\|_2\) by the integral triangle inequality. The squared norms of the cusp bound and of \(R_{p-\alpha}h\) are each at most \(S/64\). If the tests vanish, \(h=R_\alpha h+R_{p-\alpha}h\) therefore gives
\[
 S\le2(S/8+S/64+S/64)=5S/16.
\]
We have proved
\begin{equation}\label{eq:p8c1-finite-tests}
 (R_\alpha h(y_j)=0\text{ for }1\le j\le d)\quad\Longrightarrow\quad h=0\text{ in }L^2(Y).
\end{equation}
The image of the bounded inputs thus embeds linearly in \(\C^d\) and is closed in \(L^2(Y)\). Bounded inputs are dense, and \(R_eR_\phi\) is continuous, so this image contains the image of every \(L^2\) input.
\leanref{CompanionTheory.\allowbreak{}D.\allowbreak{}gk.\allowbreak{}RF.\allowbreak{}discrete\_\allowbreak{}series\_\allowbreak{}bound\_\allowbreak{}of\_\allowbreak{}inputs}{CompanionTheory.\allowbreak{}HarishChandraModules.\allowbreak{}KFiniteAnalyticVectors.\allowbreak{}FundamentalDomainCusps.\allowbreak{}FiniteDimensionalFromLinearForms}{227--305}
\leanref{CompanionTheory.\allowbreak{}D.\allowbreak{}gk.\allowbreak{}RF.\allowbreak{}rowsFinDim\_\allowbreak{}of}{CompanionTheory.\allowbreak{}HarishChandraModules.\allowbreak{}KFiniteAnalyticVectors.\allowbreak{}FundamentalDomainCusps.\allowbreak{}RowsFinDimOf}{25--86}
\end{proof}

\subsection*{Proof of \ref{B8.03}}\label{prf:B8.03}\phantomsection
\begin{proof}
Apply \ref{B8.07} with the assumed rational tests and enlargement property to obtain the member and bundle supplies on the augmented extension. The nonendoscopic and geometric Galois hypotheses give the third-layer data. The series assertions for the member construction and for \ref{B8.43} are the conclusions of Theorem~\ref{B8.43h}. Its remaining hypotheses are the norm input, depth-zero support, local finiteness, local Tate duality, the local Euler characteristic and the assumed linked bound at \(2,3\). Hence we have the linked implication
\[
 \mathcal A\longmapsto \mathrm{NL}(\mathcal D_3(\mathcal A)). \tag{8.7}
\]
The projection of the augmented extension preserves the member indices. Apply the transfer assertion of \ref{B8.08} to (8.7) and to its open rows.

In particular we now have the label realization on this extension. With the convergent Siegel--Weil formula of \ref{B8.13}, it gives the counting row by \ref{B8.52}. The finite dimension used there is \ref{B8.03a}. Apply \ref{B8.06} with the transferred open rows, (8.7), the assumed functional and covolume rows, the reference index input, and the level-one estimate of Lemma~\ref{B8.08b}. For this last application the primitive and divisor comparisons have been transferred at the same indices, and the estimate holds for every \(C\) beyond a common point cutoff. This supplies all the member rows on the extension on which the bundles were chosen.

Apply \ref{B8.04}. Its separate finite-dimension hypothesis follows from \ref{B8.03a} for every kernel specified in \ref{B8.02}. The analytic input is \ref{B8.44}, the reproduction and integrability estimates are \ref{B8.02y}, and the minor-label and signed bounds are \ref{B8.10y} and \ref{B8.10d}. The adjacent-rank mass formula is \ref{B8.48}. The nonsingular Gram formula of \ref{B8.13} gives the volume comparison of \ref{B8.49}, and \ref{B8.50} gives the genus mass. With the member and bundle supplies above, \ref{B8.04} gives one family class and selection satisfying \ref{B8.02}.
\end{proof}

\subsection*{Proof of \ref{B8.04}}\label{prf:B8.04}\phantomsection
\begin{proof}
The joint counting row of \ref{B8.06} first supplies an enlargement and then \(c_{30},c_{\rm Ic}\) with their sign and size conditions. Fix them before choosing a family or selection. Conjoin this row with the other \(25\) member rows and the \(T_u\)-row of \ref{B8.51}, which uses item~3 with (8.8) and point condition \(\top\). Intersection as in (8.11) gives a single enlargement for the conjunction \(H\). Apply \ref{B8.05} to obtain
\[
 \mathcal A,\qquad \sigma,\qquad H(\mathcal A,\sigma). \tag{8.9}
\]
Projection gives all these rows for the same \(\sigma\). Choose a join at every selected point by the join condition of \ref{B8.08}. The extra selection inequality follows from \(x>F_m.x_1\). The definitions of \(W_1,W_3\) identify the working classes, so the closed implications retain their thresholds. The design supply and construction duty of \ref{B8.05} give the design supply, the seven family conditions and stability of the coefficient lattices, with the fifth family condition as in \ref{B8.02a}. Insert the hypotheses of items~2 and~3 in their fields of \ref{B8.02}.
\leanref{GlobalTheory.\allowbreak{}Fc.\allowbreak{}AllRowsAlt}{Development.\allowbreak{}AutomorphicArithmetic.\allowbreak{}ThetaCorrespondence.\allowbreak{}ThetaPairLifts.\allowbreak{}ThirdPlaceSeries.\allowbreak{}AllRows}{12--39}
\leanref{GlobalTheory.\allowbreak{}Fc.\allowbreak{}R1.\allowbreak{}CountsIrreducibilityConstants'}{Development.\allowbreak{}AutomorphicArithmetic.\allowbreak{}ThetaCorrespondence.\allowbreak{}ThetaPairLifts.\allowbreak{}ThirdPlaceSeries.\allowbreak{}CountsIrreducibilityConstants}{113--119}
\end{proof}

\subsection*{Proof of \ref{B8.05a}}\label{prf:B8.05a}\phantomsection
\begin{proof}
The rank bounds \eqref{eq:host-rank-ranges} give \(m\to\infty\), \(m\le x\), and \(k\ge x^{9/2+10\varepsilon}\) for large \(x\). In particular \(m^4/k\to0\). Put \(v_E=\operatorname{rd}(E)/2\). By \ref{B2.46}, \(E\) contains \(\Q(\sqrt{-d})\) and lies in \(\Q(\zeta_N)\) with \(N=Q_Sqd\). The discriminant tower and the cyclotomic discriminant formula give
\(\sqrt d\le2v_E\le N\). Thus \(v_E\ge\pi\) and \(\log(2v_E)\le L\) eventually, by the host modulus bound. The inertia groups in \ref{B2.03} omit complex conjugation: at \(d\) inertia is generated by the \(d\)-sign, and away from \(d\) inertia has trivial \(d\)-sign. Hence \(E/E^+\) is unramified at finite places.

Let \(F_j=\prod_{a=0}^{j-1}a!\). Apply \ref{B8.48} at \(j=k+1,\ldots,r\). In its Euler product each \(\log L_\eta(E,j)\) is at least \(-f\sum_{n\ge2}n^{-2}\ge-f\), since there are at most \(f\) primes above a rational prime and \(\log(1+y)\le y\). The factors \((v_E/\pi)^j\) are at least one. Multiplication therefore gives
\[
 \log(M_r/M_k)\ge f\log(F_r/F_k)
                     -f(m-1)\bigl(1+\tfrac12\log(2v_E)\bigr).
\]

The number of free entries in the column group is \(k(k-1)/2\), and in the row group it is \(r(r-1)/2-k\). Dividing \(|\GL_j(\F_g)|=g^{j^2}\prod_{i=1}^j(1-g^{-i})\) by the two subgroup orders gives
\[
 \frac{I_r(g)}{I_c(g)}
   =g^{km+m(m-1)/2}\prod_{i=k+1}^r(1-g^{-i})
   \ \ge\ \tfrac12 g^{km}.
\]
Here \(g\ge3\) and the product is at least \(1-\sum_{i\ge1}g^{-i}\ge1/2\).
\leanref{GlobalTheory.\allowbreak{}EX.\allowbreak{}WinLow.\allowbreak{}idx\_\allowbreak{}ratio}{Development.\allowbreak{}AbelianVarietyTheory.\allowbreak{}Proofs.\allowbreak{}GlobalInterpolation.\allowbreak{}DimVRowPos}{124--169}

It remains to bound the two Vandermondes before division by \(F_k,F_r\). Put \(a=xe^L\), \(\delta=\lceil x^{c_D}\rceil\) and \(S_c=4k\delta\). The preliminary coordinates are formed by the formulas of \ref{B4.02b}: \(v_i=4\delta(k-2-i)\) and \(U_{i+1}=v_i-\tfrac12-\xi_i\delta\), with \(\xi_i=0\) or \(1\). These assignments use the box weights and bits. At a mark the endpoints of \(U\) are \(a_s+(k-1)/2\) and \(-b_s-(k-1)/2\); elsewhere they are \(v_0+S_c+1/2\) and \(-1/2\). The box bounds give \(a_s\le2a\), \(b_s\le2a/\exp(x^{\varepsilon/3})\) eventually. Both \(k\) and \(S_c\) are polynomially bounded. Thus every gap meeting an endpoint is at most \(4a\), uniformly. For \(0\le i<j\le k-3\) the interior gap \(U_{i+1}-U_{j+1}\) is at most \(v_i-v_{j+1}\). These pairs \((i,j+1)\) are distinct. All unused positive-row gaps are at least one. The negative row entries have successive gaps \(2\lceil x^{c_h}\rceil\), and \(b_{0,\rho}\ge1\) in each of the three base cases in \ref{B4.02b}. Thus all cross gaps are at least one as well. Writing \(\Delta\) for the numerator of \(W\), we obtain
\[
        \frac{\Delta(\lambda_\rho)}{\Delta(U_{v,\rho})}
                         \ge (4a)^{-(2k-3)} .
\]
This uses no final member or word prime.
\leanref{GlobalTheory.\allowbreak{}EX.\allowbreak{}WinLow.\allowbreak{}places\_\allowbreak{}ok}{Development.\allowbreak{}AbelianVarietyTheory.\allowbreak{}Proofs.\allowbreak{}GlobalInterpolation.\allowbreak{}WindowLower}{68--147}

The factorials in the mass bound cancel the Weyl denominators. Since \(g_i\ge(ke^L)^8\) by \ref{B2.51} and \(\log(4a)\le2L\) eventually, the preceding inequalities give
\[
 \frac1{fk}\log(D^\star/d_v^\star)
 \ge16mL-4L-\frac{m-1}{k}(1+L/2)-\frac{2\log2}{k}
 \ge15mL
\]
eventually. Finally \(2\log\Gamma/(fk)=(2C_\Gamma m^4/k)mL\le mL\) beyond a cutoff depending on \(C_\Gamma\). All cutoffs preceded the choice of box weights and words.
\leanref{GlobalTheory.\allowbreak{}EX.\allowbreak{}WinLow.\allowbreak{}rank\_\allowbreak{}consts}{Development.\allowbreak{}AbelianVarietyTheory.\allowbreak{}Proofs.\allowbreak{}GlobalInterpolation.\allowbreak{}WindowLower}{148--176}
\end{proof}

\subsection*{Proof of \ref{B8.05}}\label{prf:B8.05}\phantomsection
\begin{proof}
Conjoin point and bundle conditions by intersection. Take the minimum of the positive \(\varepsilon\)-cutoffs and the maximum of the lower thresholds in dependency order \(c_D,c_h,x\). Take the unions of the required algebraic elements and of the corresponding lists of required or avoided primes. Intersect the two strengthenings \(R\) from bundle supply and \eqref{eq:p8c1-row-large}. With \(Q^*=Q'\land R\) this gives
\[
 \mathsf B(Q^*,L_k),\qquad \operatorname{Row}(Q^*,L_k,H).
 \tag{8.11}\label{eq:p8c1-common-row}
\]
Write \(\kappa.\mathrm{Pos}\) for \(C_5,c_6,C_b>0\) in the parameter record \(\kappa\). For the bundle condition \(E\) of the row, let \(\mathcal Q_E(\kappa,D,O)\) be
\[
 \kappa.\mathrm{Pos}\ \Longrightarrow\ \forall c\in\R\
 \exists\beta\in\mathrm{Bd}(D,O,c)\quad
 \left[
 \begin{gathered}
 \beta\text{ comes from a linked }Q^*\text{-point with }\kappa\\
 {}\land E(\kappa,D,O,c,\beta)
 \end{gathered}
 \right].
 \tag{8.12}\label{eq:p8c1-bundle-condition}
\]
Beyond the larger cutoff in \eqref{eq:p8c1-common-row}, bundle supply gives \(\beta\) and the row gives \(E\) on that bundle. Take \(\mathcal A\) to consist of families in \(\mathcal C(\mathcal Q_E(\kappa))\) for some linked \(\kappa\). The link gives \(\kappa.\mathrm{Pos}\). Choose \(\sigma\) at \(c=F_m.C_{\mathrm{sz}}\) by \eqref{eq:p8c1-bundle-condition}. Since \(Q^*\) implies \(Q'\), this gives item~3, and the row gives item~4.

The construction class supplies the fixed parameters, windows, auxiliary-prime conditions, tame labels and pointwise weight test. Projection gives \(q_0<X'\) in clause~3 of \ref{B8.01k}. The two tame lists have the same entries \(1,\chi_j,\chi_j^{-1}\). The square condition therefore transfers at the same cutoff as ordinary and prime-to-\(\ell\) injectivity. For clause~5 use \(x_1+1\), so that the pointwise test applies at every \(x\ge x_1+1\) and gives the small-share inequality of \ref{B8.02a}. The defining point condition gives tuple use and the twist witnesses for every \(x>x_1\). Use \(x_1\) as the cutoff in clauses~6 and~7, which retain the strict inequality. The same point condition gives coefficient-lattice stability with cutoff \(x_1+1\).
\leanref{GlobalTheory.\allowbreak{}EX.\allowbreak{}DQ.\allowbreak{}familyClass\_\allowbreak{}sevenDuties\_\allowbreak{}and\_\allowbreak{}stableLattices}{Development.\allowbreak{}AbelianVarietyTheory.\allowbreak{}Definitions.\allowbreak{}GlobalTwisting.\allowbreak{}FamilyClassSevenDutiesAndStableLattices}{55--59}
\leanref{GlobalTheory.\allowbreak{}EX.\allowbreak{}DQ.\allowbreak{}duty\_\allowbreak{}seven\_\allowbreak{}library}{Development.\allowbreak{}AbelianVarietyTheory.\allowbreak{}Proofs.\allowbreak{}GlobalInterpolation.\allowbreak{}DutySevenLibrary}{15--23}
\leanref{GlobalTheory.\allowbreak{}Fc.\allowbreak{}copies\_\allowbreak{}frozen}{Development.\allowbreak{}AbelianVarietyTheory.\allowbreak{}Proofs.\allowbreak{}GlobalInterpolation.\allowbreak{}WorkingFamClassCopiesEqFrozen}{9--17}

For item~1 assume Minkowski division. Fix \(\lambda>1\) and a requirement \(G\) whose hosts exist on windows of exponent \(\lambda\) in the sense following Definition~\ref{B2.58}. For the target character \(\psi\), this supply is Theorem~\ref{B2.61} at \(\lambda=3/2\) under the cohomological window hypothesis of Definition~\ref{B6.03}. In Theorem~\ref{B2.61} (via Theorems~\ref{B2.59}--\ref{B2.60}, \ref{B2.48}, \ref{B2.46}, \ref{B2.57}, and \ref{B2.56}), the good pairs \((q,d)\) on \([Y,x]\) are supplied unconditionally by Theorem~\ref{B2.43}, the top set \(TT\) is chosen by Theorem~\ref{B2.54}, and the character bounds at \(C_5=4\) are supplied as follows:
\begin{itemize}
 \item the pointwise square-root bounds \(|g_\chi(x)|\le x^{1/2+\varepsilon}\) and \(|(L'/L)(1,\chi)|\le 4\log x\) for every primitive representative \(\chi\) at level \(8qd\) of \(\chi_q=(\,\cdot\,/q)\), \(\chi_d=(\,\cdot\,/d)\), and \(\chi_q\chi_d\) are supplied in Theorem~\ref{B2.56} by Lemma~\ref{B2.55} from \eqref{eq:aux-pair-good-bounds} of Definition~\ref{B2.42} (with \(\xi=1\) modulo \(M\));
 \item the bound \(|\Re(L'/L)(1,\chi_d)|\le x^\varepsilon\) follows in Theorem~\ref{B2.56} from the logarithmic-derivative bound of Lemma~\ref{B2.55} for \(\chi_d\);
 \item the reality, conductor, and zero-free box conditions in \(\mathcal R(8X_1(Y)X_2(Y)^j)\) for the six quadratic twists of \(\chi_q,\chi_d,\chi_q\chi_d\) are supplied in Theorem~\ref{B2.56} by Lemma~\ref{B2.55} from \eqref{eq:aux-pair-good-boxes} of Definition~\ref{B2.42}, while the real-character zero-free boxes at height \(H_{Q_{\mathrm{odd}}}(j)\) and the thin-position bound \eqref{eq:thin-character-bound} are supplied by Theorem~\ref{B2.52} from \eqref{eq:aux-pair-good-boxes};
 \item the bounds \(|g_\chi(x)|\le x^{1/2+\varepsilon}\) and \(|(L'/L)(1,\chi)|\le 4\log x\) for every \(\theta\in\mathcal B_H(\mathcal X)\) in the cyclic host requirement \(G\) of Definition~\ref{B2.58} are supplied in Theorems~\ref{B2.59}--\ref{B2.60} by Theorem~\ref{B2.47} from \eqref{eq:aux-pair-good-bounds} of Definition~\ref{B2.42}.
\end{itemize}
Take \(\varepsilon>0\) at most the minimum of \(1/53\), the positive cutoff for point construction from item~4 of \ref{B8.04} and \eqref{eq:p8c1-common-row}, and the cutoff of \ref{B5.22a}. Fix a requested threshold record \(\mathcal K\). Host supply gives \(C_{T_2},C_8^0,C_5,c_6,C_b,x_H\) and \(S=(J,\mathcal W)\), with \(C_8^0\ge1\), \(C_5,c_6,C_b>0\), and \(J\ge1\). The set \(\mathcal W\) contains the intervals of \eqref{eq:window-set-property}. At every \(x\ge x_H\) in \(\mathcal W\) there are ranks and a completed host of Theorem~\ref{B2.56} satisfying \(G\) at \(s=(x,\varepsilon,J)\).

On such a host \(H\), write \(N=Q_Sqd\), and denote its distinguished primes \(p_{\rm de},p'_{\rm de}\) by \(p,p'\). The host conductor ideal \(\mathfrak c_K\subseteq\mathcal O_K\) equals \(\mathcal O_K\). Let \(r_0\) be the least prime not dividing \(N(p-1)(p'-1)\), and take \(q_0(H)\) to be the least prime \(1\) modulo \(r_0N\). A prime factor of \(N(p-1)(p'-1)+1\) bounds \(r_0\) by \(2N(p-1)(p'-1)\). Since \(N\ge2\) and \(p,p'\le N^{C_8^0}\) by the completed-host bounds of \ref{B2.51} and \ref{B2.56}, \(r_0N\le N^{3+2C_8^0}\). With the exponent \(A_0\ge1\) of \ref{B2.49}, we obtain
\begin{equation}\label{eq:p8c1-q0-power}
 q_0(H)\le N^{C_8},\qquad C_8=(3+2C_8^0)A_0\ge C_8^0 .
\end{equation}
The choice of \(r_0\) excludes \(p,p'\), and the congruence modulo \(N\) gives splitting in the host field. Replace the level set by \(\{g_0,g_1,p,p',q_0(H)\}\). This leaves \(q,d,Q_S\), the frame, the arithmetic datum, \(TT\), and the auxiliary rank unchanged, so the completed-host character bounds (including the bounds for \(\chi_q,\chi_d,\chi_q\chi_d\) and the six quadratic twists from Lemma~\ref{B2.55} and the real-character boxes and thin-position bound from Theorem~\ref{B2.52}) and the requirement \(G\) transfer verbatim. Thus this choice gives completed hosts with the specified \(q_0(H)\) on the same \(S\), with \(J\le C_8\).
\leanref{UpperTheory.\allowbreak{}Top.\allowbreak{}hostQ}{Development.\allowbreak{}AutomorphicArithmetic.\allowbreak{}CharacterBoxes.\allowbreak{}SmallestPrimePolynomialBound}{64--102}

At the fixed \(\varepsilon,C_8,C_5,c_6,C_b\), let \(\mathcal K_O\) be the maximum of the thresholds for constructing preliminary points and for constructing points with the required member and bundle conditions, in both cases with prescribed box weights. Let \(\mathcal K_W=\mathcal K_W(\varepsilon)\) be the numerical threshold of \ref{B5.20y}, used in \ref{B5.22a}. These records are chosen before either family. Put \(t_*=c_T(\varepsilon,9)\). Form \(\mathcal K'\) by evaluating the middle and final functions of \(\mathcal K\) with the \(C_{Tu}\)-entry of the outer list replaced by \(t_*\). Take the maximum \(\mathcal K_*=\mathcal K'\vee\mathcal K_O\vee\mathcal K_W\) as in \ref{B8.39}. Write \(\mathbin{\Vert}\) for list concatenation. Use the construction constants \(C\) given by
\[
 \begin{aligned}
 P_0&=\max(1000,\mathcal K_{*,\mathrm{pre}}),\\
 \mathbf o&=(401,2,C_{T_2},11\cdot405,405,9/10,9,C_8,1,P_0,1,12,377,25,6,C_8),\\
 U&=\max(1,\mathcal K_{*,\mathrm{mid}}(\mathbf o)),\qquad
 \mathbf u=(U,\ldots,U)\in\R^{10},\\
 V&=\max(1,\mathcal K_{*,\mathrm{last}}(\mathbf o\mathbin{\Vert}\mathbf u)).
 \end{aligned}
\]
The preliminary list is \((P_0)\), \(\mathbf o\) is the outer list, the middle list is \(\mathbf u\) at \(CH=C_F=U\), and every final entry is \(V\). With \(K\) the orbit parameter record, \(C_\Gamma=C_{\rm sz}=K.\mathrm{La}=U\). Also \(C_{YL}=9\), and the construction value of \(C_{Tu}\) is \(25\). The three inequalities of \ref{B8.39} follow from the maxima. For the point construction and its numerical size data, include \(0\) among the lower bounds for \(c_D\) and take their maximum. At that \(c_D\), include \(402\) among the lower bounds for \(c_h\) and take their maximum. Fix these exponents and write
\[
 \kappa=(\varepsilon,C,c_D,c_h,C_5,c_6,C_b).
\]
The two scale cutoffs \(x_O,x_N\) for the point construction and the numerical data now use the same \(\kappa\).
\leanref{GlobalTheory.\allowbreak{}EX.\allowbreak{}stdConst}{Development.\allowbreak{}AbelianVarietyTheory.\allowbreak{}Proofs.\allowbreak{}GlobalInterpolation.\allowbreak{}ExistsStdLarge}{10--27}
\leanref{GlobalTheory.\allowbreak{}EX.\allowbreak{}exists\_\allowbreak{}std\_\allowbreak{}large}{Development.\allowbreak{}AbelianVarietyTheory.\allowbreak{}Proofs.\allowbreak{}GlobalInterpolation.\allowbreak{}ExistsStdLarge}{39--67}

At a completed host above \(x_O\), prescribe any pair of integer weight functions in the intervals \eqref{eq:p5c2-common-box}, extended by zero off the top set. Choose an injection of \(M=\lceil3\Gamma^2x^{C_F}\rceil\) indices into binary words of length \(k-2\) and a translation \(0\le u<n\), where \(n\) is the frame order. Here \(L=x^\varepsilon\) and \(b=e^L\). Put \(\Gamma=\exp(C_\Gamma m^5 fL)\). For a coded word \(v\) let \(d_v^\star\) be the mass of the standard column lattice times its pre-level indices and archimedean weight dimensions. Let \(D^\star\) be the row mass times the two finite indices and the product of the row Weyl factors in Lemma~\ref{B8.05a}. The completed host, the prescribed box weights and the coded word are the preliminary data of that lemma. Enlarge \(x_O\) by its uniform cutoff. It gives
\[
 S_v=\left(\frac{D^\star}{d_v^\star\Gamma^2}\right)^{1/(fk)},
 \qquad \log S_v\ge mL.
\]
At \(x^{1/J}\le Y\le x\), the pair conditions of Lemma~\ref{B2.55} and Theorem~\ref{B2.56} give \(d\equiv7\bmod16\) and \(\log N\le2x^{\varepsilon/5}(\log x)^2\). They also give \(d\ge100\) eventually. By \eqref{eq:host-rank-ranges}, eventually \(m\ge x^{1/3}\) and \(1\le fk\le x^{12}\). The host bound \(g_i\le2(kb)^8\) of \ref{B2.51} gives \(g_i\le b^9\) for large \(x\) since \(k\) is polynomial in \(x\). Thus \ref{B8.51a} with \(c_s=1,\ C_Y=9,\ V=fk\) gives a norm prime in \((S_v,(1+1/(fk))S_v]\). It is at least \(X'=\exp(x^{\varepsilon/3})\) for large \(x\) and satisfies the required norm congruence modulo \(N(g_0g_1)^2\).
\leanref{GlobalTheory.\allowbreak{}EX.\allowbreak{}WinLow.\allowbreak{}window\_\allowbreak{}lower}{Development.\allowbreak{}AbelianVarietyTheory.\allowbreak{}Proofs.\allowbreak{}GlobalInterpolation.\allowbreak{}WindowLower}{177--295}

Fix these levels. Item~4 of \ref{B8.04} gives an untwisted \(Q\)-member on a \(KP\)-base containing the prescribed elements. Export its data and use them in the second-member clause, which preserves the frame and prime set. The resulting member satisfies \(Q\) and \(Q'\) on a \(KP\)-base. Since \(Q\Rightarrow_{KP}R\), it satisfies \(Q^*=Q'\land R\). Use the volume ratio on its tied datum and \eqref{eq:p8c1-bundle-condition} on its linked member. The layer construction carries tuple use, compatible twists and coefficient-lattice stability to the point. It has scale \(s\), the prescribed weights at the corresponding top labels, the prescribed code bits and translation \(u\). Its orbit data and size parameters are those in \ref{B8.02ac}. The cutoff \(x_O\) precedes the choices of host, weights, code and translation. The same supply gives preliminary points by retaining the scale, weight, code and translation data.
\leanref{GlobalTheory.\allowbreak{}Fc.\allowbreak{}memberTop\_\allowbreak{}final\_\allowbreak{}contextParameters}{Development.\allowbreak{}AbelianVarietyTheory.\allowbreak{}Proofs.\allowbreak{}IwasawaDecomposition.\allowbreak{}MemberTopFinalContextParameters}{14--23}
\leanref{GlobalTheory.\allowbreak{}Fc.\allowbreak{}objCP\_\allowbreak{}and\_\allowbreak{}of\_\allowbreak{}recExP}{Development.\allowbreak{}AbelianVarietyTheory.\allowbreak{}Proofs.\allowbreak{}GlobalInterpolation.\allowbreak{}ExistsObjectsOfAdmissibleTwistDatum}{70--119}
\leanref{UpperTheory.\allowbreak{}Top.\allowbreak{}objC\_\allowbreak{}true}{Development.\allowbreak{}AutomorphicArithmetic.\allowbreak{}CharacterBoxes.\allowbreak{}ObjCTrue}{16--28}

Above \(x_N\) the numerical witnesses use the host arithmetic of \ref{B2.56} and the prescribed scale, sizes and ranks. Both lists of local parameters have the entries fixed by \(C\). Put \(\delta_1=x^{-\varepsilon/8}\), \(a=xb\), and \(b'=a/X'\) as in \ref{B3.51}. The second box ends at \(b'+2n\delta_1x^\varepsilon\mathcal L b'/k\), where \(\mathcal L\) is the orbit factor of \ref{B3.51}. In \eqref{eq:p3c5-window}, \(u_1\delta_1\le2\) and \(M_0=O_\varepsilon(\log x)\), so \(N_{\rm win}=O_\varepsilon(\log x)\). Hence \(\mathcal L=O_{\varepsilon,C}(x^\varepsilon(\log x)^2)\). In the two terms of \(\rho_B\) in \ref{B3.51}, substitute \(n\le2x^{2+8\varepsilon}\) and \(k\ge x^{9/2+10\varepsilon}\) from \eqref{eq:host-rank-ranges}. This gives
\begin{equation}\label{eq:p8c1-rho-point}
 \rho_B=O_{\varepsilon,C}\bigl(x^{-5/2}(\log x)^2+
                     x^{-9/2-9\varepsilon}(\log x)^3\bigr)\le x^{-2}
\end{equation}
after enlarging \(x_N\).
\leanref{GlobalTheory.\allowbreak{}EX.\allowbreak{}szRow\_\allowbreak{}numAt}{Development.\allowbreak{}AbelianVarietyTheory.\allowbreak{}Proofs.\allowbreak{}GlobalInterpolation.\allowbreak{}SzRowNumAt}{14--47}

Use \(\lceil a\rceil,\lceil b'\rceil\) as preliminary weights on \(TT\), extended by zero. Both box widths of \ref{B5.22a} eventually exceed one since \(k\le x^5\) by \eqref{eq:host-rank-ranges}, \(n,\delta_1x^\varepsilon,\mathcal L\ge1\) for large \(x\), and \(b,b'\) exceed every fixed power of \(x\). Let \(x_w\) be a common cutoff. By \eqref{eq:host-rank-ranges} and \(f=2n\),
\[
 m^5 f L\le128x^{9/2+9\varepsilon},\qquad k\ge x^{9/2+10\varepsilon}.
\]
Consequently \(\log(3\Gamma^2x^{C_F})=o(k)\) uniformly in the ranks, and beyond a cutoff \(x_c\) we have \(M\le2^{k-2}\). Use the first \(M\) binary strings of length \(k-2\) in lexicographic order. Put \(\xi=\max(2,x_H,x_O,x_N,x_w,x_c)\).

For every \(z\in\mathcal W\) with \(z\ge\xi\), fix one host \(H_z\) and construct preliminary points \(P_{z,u}\) at all \(0\le u<n_z\), where \(n_z\) is its frame order. Use the same preliminary weights and code for every translation. The common labels and weight arrays identify the remaining fibre data by \ref{B8.02ac}. The code identifies the words by their indices. The offsets in \ref{B3.52} depend on the scale, layer and membership in \(TT\), so they agree at matching labels.
\leanref{GlobalTheory.\allowbreak{}EX.\allowbreak{}exists\_\allowbreak{}weights\_\allowbreak{}in\_\allowbreak{}windows}{Development.\allowbreak{}AbelianVarietyTheory.\allowbreak{}Proofs.\allowbreak{}GlobalInterpolation.\allowbreak{}ExistsWeightsInWindows}{18--78}
\leanref{GlobalTheory.\allowbreak{}EX.\allowbreak{}exists\_\allowbreak{}injective\_\allowbreak{}words}{Development.\allowbreak{}AbelianVarietyTheory.\allowbreak{}Proofs.\allowbreak{}GlobalInterpolation.\allowbreak{}MPowFiveMulFMulLLe}{40--94}
\leanref{GlobalTheory.\allowbreak{}EX.\allowbreak{}TieSame.\allowbreak{}same\_\allowbreak{}of\_\allowbreak{}from}{Development.\allowbreak{}AbelianVarietyTheory.\allowbreak{}Proofs.\allowbreak{}GlobalInterpolation.\allowbreak{}SamePointOfSameScale}{87--131}
\leanref{GlobalTheory.\allowbreak{}EX.\allowbreak{}DC.\allowbreak{}kit\_\allowbreak{}of\_\allowbreak{}points}{Development.\allowbreak{}AbelianVarietyTheory.\allowbreak{}Proofs.\allowbreak{}GlobalInterpolation.\allowbreak{}KitOfPoints}{16--28}

Choose \(y\ge\xi\) in \(\mathcal W\) by \eqref{eq:window-set-property} and define
\begin{equation}\label{eq:p8c1-family-indices}
 z(x)=\begin{cases}
 x,&x>\xi,\ x\in\mathcal W,\\
 y,&\text{otherwise},
 \end{cases}
 \qquad
 I_x=\{u\in\Z:0\le u<n_{z(x)},\ x\le\xi\ \text{or}\ x\in\mathcal W\}.
\end{equation}
Set \(F^0.x_1=\xi,\ F^0.J=J,\ F^0.\mathrm{Win}=\mathcal W\). At \(u\in I_x\) use \(P_{z(x),u}\), with the common top labels and box endpoint read at index \(0\) and the constants read from \(C\). The index \(0\) exists at \(x=0\) and at each window point. At indexed \(x>\xi\) the point scale is \((x,\varepsilon,J)\). Numerical witnesses are available at every \(z(x)\ge\xi\). The threshold lists agree with those of \(C\) since the window-dependent entry is \(\max(C_8,J)=C_8\). Thus \(F^0\) meets the thresholds \(\mathcal K_W\) at \(C_F\) in the sense of Definition~\ref{B8.39}.
\leanref{GlobalTheory.\allowbreak{}EX.\allowbreak{}cutScale}{Development.\allowbreak{}AbelianVarietyTheory.\allowbreak{}Proofs.\allowbreak{}GlobalInterpolation.\allowbreak{}ExistsFamPOfKits}{15--16}
\leanref{GlobalTheory.\allowbreak{}EX.\allowbreak{}famOfKits}{Development.\allowbreak{}AbelianVarietyTheory.\allowbreak{}Proofs.\allowbreak{}GlobalInterpolation.\allowbreak{}ExistsFamPOfKits}{56--157}
\leanref{GlobalTheory.\allowbreak{}EX.\allowbreak{}exists\_\allowbreak{}famP\_\allowbreak{}of\_\allowbreak{}kits}{Development.\allowbreak{}AbelianVarietyTheory.\allowbreak{}Proofs.\allowbreak{}GlobalInterpolation.\allowbreak{}ExistsFamPOfKits}{159--172}

The preliminary weights were chosen by rounding the lower box endpoints. For the resulting family \(F^0\), the pre-list, outer list, middle list at \(C_F\), and last list are the lists on which the construction constants passed \(\mathcal K_W\). Thus \ref{B5.22a} applies with \(CH=C_F\) and gives a single point cutoff \(x_W=x_W(F^0,C_F)\) for all its indices. The pair obtained at each subsequent scale is a trial pair on the preliminary point. It will be prescribed as the actual weights in the final construction. At each \(x\in\mathcal W\) with \(x\ge X_H=\max(\xi+1,x_W)\), take the trial pair \((\alpha_x,\beta_x)\) furnished at \(P_{x,0}\). By \ref{B8.02ac} it belongs to the host boxes, vanishes off \(TT\), and satisfies
\begin{equation}\label{eq:p8c1-host-share}
 R_{H_x}(\alpha_x,\beta_x)=R_{P_{x,0}}(\alpha_x,\beta_x)\le x^{-9}.
\end{equation}
\leanref{UpperTheory.\allowbreak{}Top.\allowbreak{}hostWW}{Development.\allowbreak{}AutomorphicArithmetic.\allowbreak{}CharacterBoxes.\allowbreak{}HostWW}{27--87}

For the design family take \(\xi'\ge\max(2,X_H,x_O,x_N,x_c)\) large enough that \(q_0(H_x)<X'\) on the window set. By \eqref{eq:p8c1-q0-power} and the conductor bound of \ref{B2.56},
\[
 \log q_0(H_x)\le C_8\log N\le2C_8x^{\varepsilon/5}(\log x)^2
 <x^{\varepsilon/3}
\]
for large \(x\). At each such \(x\), construct \(P'_{x,u}\) for every \(0\le u<n_x\) with the single host pair \((\alpha_x,\beta_x)\) and one injective code at that scale. By \ref{B8.02ac}, the pair gives the extended weights of each point, and \eqref{eq:p8c1-host-share} gives the small-share bound at each translation. Thus \(t\mapsto\alpha_x(t)\) is fixed throughout the translation average in \ref{B2.20}.

Assemble these fibres by \eqref{eq:p8c1-family-indices} with cutoff \(\xi'\) and a reference \(y'\in\mathcal W\) strictly above it. The numerical witnesses again give the constant lists. At every indexed \(x>\xi'\) the points retain \(G\) and \eqref{eq:p8c1-bundle-condition}. Each constructed point \(P'_{x,u}\) has the same primes \(q,d\) and conductor \(Q_S\) as \(H_x\), so the completed-host pointwise bounds for \(\chi_q\), \(\chi_d\), and \(\chi_q\chi_d\) with constant \(C_5\) (supplied in Theorem~\ref{B2.56} by Lemma~\ref{B2.55} from \eqref{eq:aux-pair-good-bounds} of Definition~\ref{B2.42}) transfer verbatim to item~4 of Definition~\ref{B8.02} at every \(x\ge\xi'+1\) and every \(i\in I_x\). The six quadratic-twist nonvanishing boxes and the real-character zero-free boxes (supplied in Theorem~\ref{B2.56} by Lemma~\ref{B2.55} and Theorem~\ref{B2.52} from \eqref{eq:aux-pair-good-boxes} of Definition~\ref{B2.42}) transfer verbatim to the family box and local \(QD\) conditions. Likewise, the host thickness, bad-orbit, and \(G\) conditions (including the square-root bounds for every character in \(\mathcal B_H(\mathcal X)\) from Theorem~\ref{B2.47} when \(G\) is the cyclic host requirement of Definition~\ref{B2.58}) transfer through the equal sets, prime data, and Frobenius labels. The prime \(q_0\) retains its level membership and congruence. Since \(X'\le b\), its logarithm is at most \(L\), and the bounds \(g_i\ge(kb)^8\) and \(T_u\ge X'\) give the remaining exclusions. The tame-generator limit is
\[
 l_J=\frac{mbL}{k(m+t)\log(3(b+\mathrm{Sc}+\eta+t+1))},
\]
where \(t\) is the auxiliary rank and \(\mathrm{Sc},\eta\) are the point sizes of \ref{B8.02ac}. At fixed construction parameters \(k,m,t,\mathrm{Sc},\eta\) grow at most polynomially by \eqref{eq:host-rank-ranges} and \ref{B8.02ac}. Thus \(\log l_J=L-O_{\varepsilon,C,c_D,c_h}(\log x)\) and \(X'\le l_J\) after enlarging \(\xi'\). The word bijection gives \(M\) labels at the actual scale, with \(M\le3\Gamma^2x^{C_F}+1\). The full member construction gives the twisting and stability conditions. Minkowski division gives the word-lattice denominator divisibility by \(M(2fN(m))\) of the working class in \ref{B8.02}.

The family has \(C_{YL}=9\) and \(CH=C_F\ge1\). Set its \(C_{Tu}\) to \(t_*\). This replaces entry \(13\) of the zero-indexed outer list. The middle and final tests of \ref{B8.39} for \(\mathcal K\) now evaluate the functions used for \(\mathcal K'\), and the preliminary test is unchanged. Restoring that entry to \(25\) in the parameter comparison gives \(\mathbf o\), which satisfies the linkage relation to \(\kappa\). The family is therefore in \(\mathcal A\), is large for \((\mathcal K,C_F)\), realizes \(G\) on \(\mathcal W\), and has \(C_{Tu}=c_T(\varepsilon,C_{YL})\). These choices give the design supply of Definition~\ref{B8.02}.
\leanref{UpperTheory.\allowbreak{}Top.\allowbreak{}kitS}{Development.\allowbreak{}AutomorphicArithmetic.\allowbreak{}CharacterBoxes.\allowbreak{}KitS}{14--62}
\leanref{GlobalTheory.\allowbreak{}EX.\allowbreak{}DC.\allowbreak{}has\_\allowbreak{}withDuty\_\allowbreak{}wall\_\allowbreak{}and\_\allowbreak{}choices}{Development.\allowbreak{}AbelianVarietyTheory.\allowbreak{}Proofs.\allowbreak{}GlobalInterpolation.\allowbreak{}HasWithDutyWallAndChoices}{20--60}
\leanref{UpperTheory.\allowbreak{}Top.\allowbreak{}designS}{Development.\allowbreak{}AutomorphicArithmetic.\allowbreak{}CharacterBoxes.\allowbreak{}DesignS}{21--98}
\leanref{UpperTheory.\allowbreak{}Top.\allowbreak{}recS}{Development.\allowbreak{}AutomorphicArithmetic.\allowbreak{}CharacterBoxes.\allowbreak{}ObjCTrue}{30--36}
\end{proof}

\subsection*{Proof of \ref{B8.06}}\label{prf:B8.06}\phantomsection
\begin{proof}
Use \ref{B8.09}, \ref{B8.10}, \ref{B8.54} and \ref{B8.49} for the arithmetic, spectral, split-prime and volume rows. For the type-prime period bounds, take \(R=\top\) in \eqref{eq:p8c1-lg1}. From the good member take its built witness \(b\) and the genuine column component furnished by that witness. The row component of \(p\) is \(R_{\rm row}\), so the comparison supplied by \(b\) applies to these arithmetic row coordinates. By Lemma~\ref{lem:p8c1-row-forms}, the reference and twisted identifications carry the rational row forms to their prescribed period coordinates in every column. The designated generators have rational preimages, and the integral row lattice is the inverse image of their integral span. Bundle linkage transports the genuine columns and their period identities to every linked bundle point, where the type-prime estimates of \ref{B4.23} and \ref{B4.23a} give the row in \ref{B8.02}.
\leanref{GlobalTheory.\allowbreak{}Fc.\allowbreak{}row\_\allowbreak{}aAtG\_\allowbreak{}at\_\allowbreak{}built\_\allowbreak{}alt}{Development.\allowbreak{}AutomorphicArithmetic.\allowbreak{}ThetaCorrespondence.\allowbreak{}ThetaPairLifts.\allowbreak{}ThirdPlaceSeries.\allowbreak{}RowAAtGAtBuilt}{10--18}

For the complementary bounds, use \(\varepsilon_1=1/100\) and the zero threshold record in \ref{B8.39}. Fix an eligible family in \(W_3\) and \(0<\varepsilon\le\varepsilon_1<1/52\). In the notation of \ref{B5.12k}, \(B_0=fmbL\) and \(D_{\Gamma,u}\) is the stabilizer denominator. The level and stabilizer size bounds in the definition of \(W_3\) give cutoffs \(x_T,x_\Gamma\) and a family constant \(C_2\ge0\) such that at every indexed point above the respective cutoff,
\begin{equation}\label{eq:p8c1-complement-input-sizes}
 \log T_u\le C_{Tu}mL,\qquad
 \log D_{\Gamma,u}\le C_2 fN_m\log(fN_m)
\end{equation}
for all words. The point data exclude \(T_u=g_i\). For the norm identity in \ref{B5.12k}, sum the valuation of each \(\tau(\bar\pi_i)\) over the complete fibre \(w\mid g_i\), where \(\pi_i\bar\pi_i=g_i\). In the normalization of \ref{B4.03} and \ref{B4.52}, the weighted sum for each \(\tau\) is \(\log g_i\). Multiplication by \(A_{u,j,\tau}\) and summation in \(\tau\) give \(A_{uj}\log g_i\).
\leanref{GlobalTheory.\allowbreak{}TypeValNN.\allowbreak{}typePrime\_\allowbreak{}valuationSum\_\allowbreak{}eq\_\allowbreak{}logGap\_\allowbreak{}all}{Development.\allowbreak{}AbelianVarietyTheory.\allowbreak{}Proofs.\allowbreak{}DivisorProofs.\allowbreak{}TypePrimeValuationSumEqLogGapFamily}{13--25}

For \(x\ge2\), put \(W_j=\sum_\tau w_{j\tau}\) and let \(W_j^{\rm off}\) be its sum over embeddings outside the top positions. The layer degree calculation sums the contributions of the initial rows and the added letters. Let \(d_j^{\rm Fock}\) be the sum, over all embeddings, of the absolute values of the entries in both Fock weight lists at layer \(j\). Subtracting this degree and \(S_a+W_j^{\rm off}\) gives
\begin{equation}\label{eq:p8c1-complement-degree}
 |A_{uj}-d_j^{\rm Fock}-S_a-W_j^{\rm off}|
 \le20\left[f(kS_c+N_m^2+mt_{\rm ch}\eta)+|TT|b'/x\right].
\end{equation}
The size parameters give \(N_m^2\le\eta\) and \(mt_{\rm ch}\ge1\). Since \(0\le W_j^{\rm off}\le W_j\), adding back \(W_j^{\rm off}\) in \eqref{eq:p8c1-complement-degree} gives \(R_j\) with \(C_{\deg}=40\).
\leanref{GlobalTheory.\allowbreak{}FockDegreeFine.\allowbreak{}degree\_\allowbreak{}fine}{Development.\allowbreak{}AbelianVarietyTheory.\allowbreak{}Proofs.\allowbreak{}DivisorProofs.\allowbreak{}NuAbsPlaceFine}{135--159}
\leanref{GlobalTheory.\allowbreak{}AdmissibleDegree.\allowbreak{}admissibleDegreeCount\_\allowbreak{}along\_\allowbreak{}family}{Development.\allowbreak{}AbelianVarietyTheory.\allowbreak{}Proofs.\allowbreak{}DivisorProofs.\allowbreak{}AdmissibleDegreeCountAlongFamily}{63--70}

Choose the type-prime constant after \(C_{\deg}\). The point size bounds give \(W_j\le3fb\) and \(|TT|b'/x\le fb\). They also give \(kS_c,mt_{\rm ch}\eta,kN_m\le b\). Therefore \(R_j\le(6+4C_{\deg})fb\). They also give \(f\eta\le b\), and hence \(|H_g|\le |c'_g|(1+m^2)b\). Since \(m\le f\) and \(\log g_i\le(C_{YL}+1)L\), these terms contribute at most \([(6+4C_{\deg})+4|c'_g|](C_{YL}+1)B_0\) to the type-prime bound. The denominator term satisfies
\[
 v_{g_i}(D_{\Gamma,u})\log g_i\le\log D_{\Gamma,u}
 \le C_2 fN_m\log(fN_m).
\]
For this term the size bound \(kN_m\le b\) gives \(fN_m\le b\) since \(f\le k\), so \(fN_m\log(fN_m)\le B_0\). Thus one may take
\begin{equation}\label{eq:p8c1-complement-type}
 C_g=\max\{0,(6+4C_{\deg}+4|c'_g|)(C_{YL}+1)+2C_2\}.
\end{equation}
For the stabilizer clause the family numerical sizes, using the level size bounds, also give \(N_m L\le b\) beyond a family cutoff. With \(f\le N_m\le b\) this gives \(fN_m\log(fN_m)\le2B_0\). The second inequality of \eqref{eq:p8c1-complement-input-sizes} then gives \(2\log D_{\Gamma,u}\le4C_2B_0\).
\leanref{GlobalTheory.\allowbreak{}TypeSize.\allowbreak{}typePrime\_\allowbreak{}sizes\_\allowbreak{}le\_\allowbreak{}budget\_\allowbreak{}single}{Development.\allowbreak{}AbelianVarietyTheory.\allowbreak{}Proofs.\allowbreak{}DivisorProofs.\allowbreak{}TypePrimeSizesLeBudgetFamily}{80--135}
\leanref{Formalism.\allowbreak{}MainNotions.\allowbreak{}J.\allowbreak{}gammaDiscriminant\_\allowbreak{}budget\_\allowbreak{}of\_\allowbreak{}sizes}{Formalism.\allowbreak{}SharpRestriction.\allowbreak{}AdmissibleParameterJoins.\allowbreak{}DIVfromGLbud}{63--124}

On the complete level fibre over \(T_u\), the nonnegative normalized losses of \ref{B4.03} sum to \(E^\pi_u\log T_u\) because the local valuations of the level generator sum to \(\log T_u\) with the weights of \ref{B4.52}. This also bounds the sum on each finite subset of the fibre. With \(W_{u,\tau}\) the archimedean level weight at \(\tau\), the full allowance of \ref{B5.12k} is
\[
 E^\pi_u=2\sum_\tau W_{u,\tau}+fN_m+
 \frac{|TT|(t_{\rm ch}+1)(m+1)}4.
\]
For \(x\ge2\), each \(W_{u,\tau}\le5b\), \(fN_m\le fb\), and \(|TT|(t_{\rm ch}+1)(m+1)\le fb\). Hence
\begin{equation}\label{eq:p8c1-complement-level}
 E^\pi_u\le2(5fb)+fb+fb=12fb.
\end{equation}
By \eqref{eq:p8c1-complement-input-sizes} and \eqref{eq:p8c1-complement-level}, take \(C_T=\max(12C_{Tu},0)\). The discriminant loss of \ref{B4.03} obeys \(|D_j^*|\le |TT|b+fk b'\). This follows by bounding the discriminant base term and the largest column correction on the ramified fibre. Since \(|TT|\le f\), \(kb'\le b\) and \(\log d\le L\) in the point sizes for \(x\ge2\), we obtain \(2\log d\,|D_j^*|\le4B_0\). Take \(C_d=4\).
\leanref{Formalism.\allowbreak{}MainNotions.\allowbreak{}product\_\allowbreak{}first\_\allowbreak{}zeroth\_\allowbreak{}parts.\allowbreak{}ETpi}{Formalism.\allowbreak{}SharpRestriction.\allowbreak{}GlobalSchurBounds.\allowbreak{}ProductFirstZerothParts}{307--309}
\leanref{GlobalTheory.\allowbreak{}SeenHalf.\allowbreak{}seen\_\allowbreak{}half}{Development.\allowbreak{}AbelianVarietyTheory.\allowbreak{}Proofs.\allowbreak{}DivisorProofs.\allowbreak{}VAddNonnegOfIsIntegral}{112--118}
\leanref{GlobalTheory.\allowbreak{}LevelBudget.\allowbreak{}admissible\_\allowbreak{}off\_\allowbreak{}type\_\allowbreak{}level\_\allowbreak{}budget}{Development.\allowbreak{}AbelianVarietyTheory.\allowbreak{}Proofs.\allowbreak{}DivisorProofs.\allowbreak{}AdmissibleOffTypeFamilyBound}{38--56}
\leanref{GlobalTheory.\allowbreak{}RamifiedOffD.\allowbreak{}offDiagonalRamified\_\allowbreak{}bound\_\allowbreak{}along\_\allowbreak{}family}{Development.\allowbreak{}AbelianVarietyTheory.\allowbreak{}Proofs.\allowbreak{}DivisorProofs.\allowbreak{}OffDiagonalRamifiedBoundAlongFamily}{50--54}

For the last clause of \ref{B5.12k}, partition the finite set of eligible primes \(P\) into \(P_C\), those in residue class zero with \(h_p=0\), and its complement \(P_I\). On \(P_C\) the loss is \(C_{j,p}^*\) of Definition~\ref{B4.03}. In \ref{B4.03} its digit contribution after multiplication by \(\log p\) is \(3fN_m^2\log a\), and its other contribution is \(W_j\log p/(p^2-1)\), with \(a=xb\). The size condition \(xN_m^2\le w''\) gives \(X''N_m^2\le b\), where \(X''=a/w''\). With \(W_j\le3fb\), \(\log p\le\log a\) and \(\sum_{n\ge2}(n^2-1)^{-1}\le1\), the sum is at most \(6fb\log a\). Use \(\log a\le(1+1/\varepsilon)L\) and \(m\ge1\) to take \(18(1+1/\varepsilon)B_0\) as an upper bound after the factor \(2\).

On \(P_I\) the loss is \(I_{j,p}^*\) of the same definition. The orbit average, the extra term on residue-class-zero primes of odd order, and the digit and size errors in \ref{B4.03} give
\begin{equation}\label{eq:p8c1-complement-inert}
 \sum_{p\in P_I}|I_{j,p}^*|\log p
 \le126fmb\sum_{p\in P_I}\frac{\log p}{p}
     +9fbk\sum_{p\in P_I\cap\mathcal C_s}\frac{\log p}{p^3}
     +70f\eta L\,|P_I|,
\end{equation}
where \(\mathcal C_s\) consists of the residue-class-zero primes for which \(h_p\) has odd order. The odd-order condition gives \(p\equiv\pm1\pmod{2^{e+2}}\) on this set, where \(2^e\) is the frame's two-power parameter. Put \(A=2^e\). Each such prime is in one of the two progressions \(4Aj\pm1\) with \(j\ge1\). For \(p=4Aj\pm1\), we have \(p\ge3Aj\) and \(\log p\le2(p/A)\log A\), by \(\log(p/A)\le p/A-1\) and \(\log A\ge1/2\). Thus
\[
 \frac{\log p}{p^3}\le\frac{\log A}{A^3}\frac{2A^2}{p^2}
 \le\frac{\log A}{4A^3j^2}.
\]
Summing the two progressions with \(\sum_{j\ge1}j^{-2}\le2\) gives \(\log A/A^3\). The frame bound \(n\le oA\), the size lower bound \(n\ge x^{2+\varepsilon/100}\), and \(o\le x^{\varepsilon/100}\) give \(x^2\le A\). Hence \(k\le x^5\le A^{5/2}\), and \(\log A\le2\sqrt A\) gives \(k\sum_{p\in P_I\cap\mathcal C_s}\log p/p^3\le2\).
\leanref{UpperTheory.\allowbreak{}sum\_\allowbreak{}log\_\allowbreak{}div\_\allowbreak{}cube\_\allowbreak{}le\_\allowbreak{}of\_\allowbreak{}two\_\allowbreak{}pow\_\allowbreak{}dvd}{Development.\allowbreak{}AutomorphicArithmetic.\allowbreak{}OrbitLabellings.\allowbreak{}SumLogDivCubeLeOfTwoPowDvd}{48--123}

The point sizes give \(X''\eta\le b\) and \(X''\le b\), so \(\eta|P_I|\le b\). The archive uses \(\sum_{p\le b}\log p/p\le\log b+\log4\le3L\) for \(L\ge1\). Rosser and Schoenfeld give the stronger upper bound \(\log b\) for \(b>1\) \cite[(3.24)]{rosser-1962}. With \(m\ge2\) and \(L\ge1\), the right side of \eqref{eq:p8c1-complement-inert} is at most \((378+9+35)B_0\le424B_0\). The two prime sets contribute disjoint sums, so take \(C_I=18(1+1/\varepsilon)+848\), with the larger of their cutoffs.
\leanref{GlobalTheory.\allowbreak{}ConjLossPrimes.\allowbreak{}loss\_\allowbreak{}sum\_\allowbreak{}primes}{Development.\allowbreak{}AbelianVarietyTheory.\allowbreak{}Proofs.\allowbreak{}DivisorProofs.\allowbreak{}LossSumPrimes}{103--112}
\leanref{GlobalTheory.\allowbreak{}OtherLoss.\allowbreak{}sum\_\allowbreak{}cube\_\allowbreak{}le}{Development.\allowbreak{}AbelianVarietyTheory.\allowbreak{}Proofs.\allowbreak{}DivisorProofs.\allowbreak{}OtherLossSum}{438--446}
\leanref{GlobalTheory.\allowbreak{}OtherLoss.\allowbreak{}sum\_\allowbreak{}eatI}{Development.\allowbreak{}AbelianVarietyTheory.\allowbreak{}Proofs.\allowbreak{}DivisorProofs.\allowbreak{}OtherLossSum}{447--498}
\leanref{GlobalTheory.\allowbreak{}InertLossAdm.\allowbreak{}inertPrimes\_\allowbreak{}loss\_\allowbreak{}log\_\allowbreak{}sum\_\allowbreak{}admissible}{Development.\allowbreak{}AbelianVarietyTheory.\allowbreak{}Proofs.\allowbreak{}DivisorProofs.\allowbreak{}InertPrimesLossLogSumAdmissible}{12--48}

Take the maximum of the six cutoffs. The constants depend on the fixed family and work at every point above that cutoff with \(x>F_m.x_1\). The bounds in \ref{B5.12k} hold for every word and layer and for every finite set specified there. Thus the complementary bounds satisfy \eqref{eq:p8c1-lg1} with \(R=\top\).
\leanref{GlobalTheory.\allowbreak{}familySecondComplementInputs\_\allowbreak{}of\_\allowbreak{}tuSize\_\allowbreak{}of\_\allowbreak{}dGamSize}{Development.\allowbreak{}AbelianVarietyTheory.\allowbreak{}Proofs.\allowbreak{}DivisorProofs.\allowbreak{}AdmissibleComplementTwoAlongFamily}{18--41}
\leanref{GlobalTheory.\allowbreak{}admissibleConstants\_\allowbreak{}closed\_\allowbreak{}under\_\allowbreak{}sizeBounds}{Development.\allowbreak{}AbelianVarietyTheory.\allowbreak{}Proofs.\allowbreak{}DivisorProofs.\allowbreak{}AdmissibleConstantsClosedUnderSizeBounds}{10--19}

Use the pointwise rows of \ref{B8.08}. Primitivity gives surjectivity of the reference embedding and of each layer projection. The comparison of the two realizations identifies the selected twists at the places in the reference twisting set. Use \ref{B8.55} for the extremal vectors and one-map identities, and \ref{B8.56} with reciprocity and linkage for the local lifts and the kernel conditions. For the counting row use the label realization among the assumed open conditions of \ref{B8.08}. On its original enlargement, \ref{B8.52} applies to that realization by the assumed convergent formula of \ref{B8.13}. It supplies the two counting constants before the member rows are intersected. The assumed norm condition gives \(NN(W_3(\mathcal A))\). For the simultaneous row (8.4), impose the primitivity, divisor, finite-type goodness, geometry, multiplicity, \(\kappa\) and definite-side conditions of \ref{B8.08} together with the covolume and functional bounds. Apply \ref{B8.38} on their common strengthening with the compatible summands. The paired-index and functional bounds use the family division estimate \eqref{eq:B8-divinN}. The lower weight bound uses the prescribed leaf map at every layer, as in \ref{B8.38}.
\end{proof}

\subsection*{Proof of \ref{B8.07a}}\label{prf:B8.07a}\phantomsection
\begin{proof}
Write \(K\subset\C\) for the embedded coefficient field and \(K_0\subset K\) for its specified fixed subfield. First choose the places on which to retain the row data. At the reference tuple retain \(v\) if, at \(v\) or its conjugate, some layer has nonzero Iwasawa index, nonzero lattice index in the row pairing, or nonzero prescribed content. The prescribed content at layer \(j\) depends on \(v\) through its residue characteristic.

In the untwisted case the choice is made uniformly as follows. For the fixed preliminary datum and frame, group the initial choices by the generators at the two distinguished primes and at the word primes. Their defining norm conditions give a finite set of generator tuples. For each tuple admitting adelic data choose one such realization. Include generators of its coefficient field among the finite requirements on \(K_0\), and include the residue characteristics of its nonzero reference indices among the required primes. The index supports are finite because they come from the elementary divisors of the finitely many reference lattices. Take the unions over the generator tuples and include the primes supporting the prescribed content. These choices precede the choice of \(w\).
\leanref{GlobalTheory.\allowbreak{}Mem.\allowbreak{}R3.\allowbreak{}Rw.\allowbreak{}Resize.\allowbreak{}finite\_\allowbreak{}gens}{Development.\allowbreak{}AbelianVarietyTheory.\allowbreak{}Proofs.\allowbreak{}IwasawaDecomposition.\allowbreak{}FiniteNonzeroIndexPlacesReferenceTuple}{114--120}
\leanref{GlobalTheory.\allowbreak{}Mem.\allowbreak{}R3.\allowbreak{}Rw.\allowbreak{}Resize.\allowbreak{}seedW\_\allowbreak{}bound}{Development.\allowbreak{}AbelianVarietyTheory.\allowbreak{}Proofs.\allowbreak{}IwasawaDecomposition.\allowbreak{}FiniteNonzeroIndexPlacesReferenceTuple}{166--175}
\leanref{GlobalTheory.\allowbreak{}Mem.\allowbreak{}R3.\allowbreak{}Rw.\allowbreak{}Resize.\allowbreak{}seeds\_\allowbreak{}uniform\_\allowbreak{}exists}{Development.\allowbreak{}AbelianVarietyTheory.\allowbreak{}Proofs.\allowbreak{}IwasawaDecomposition.\allowbreak{}ExistsUniformNormalisedAdelicSeeds}{12--42}

Extend the chosen reference realization to the coefficient data of \(w\), with its frame and generator tuple identified. Under this extension the support of the first two indices pulls back from the reference field. Its image in \(\C\) lies in \(K_0\). Thus two places of \(K\) with the same restriction to \(K_0\) lie over the same reference place. They also have the same residue characteristic, so they agree on membership in the content support. This proves the fibre assertion. All retained places have residue characteristic in the required prime set. In the twisted case the eligibility of \(w\) supplies adelic data with this support condition and trivial twist off the selected set. Restriction to the selected places preserves the indicated indices. The multiplicity identity and the rational tests of Theorem~\ref{B8.18} give the row lifting data there. Denote the resulting set of places by \(W\).
\leanref{GlobalTheory.\allowbreak{}Mem.\allowbreak{}R3.\allowbreak{}Dat.\allowbreak{}mem\_\allowbreak{}auxiliaryPlaces\_\allowbreak{}iff\_\allowbreak{}three\_\allowbreak{}parts}{Development.\allowbreak{}AbelianVarietyTheory.\allowbreak{}Proofs.\allowbreak{}IwasawaDecomposition.\allowbreak{}AuxiliaryPlacesEqUnionThreeParts}{43--46}
\leanref{GlobalTheory.\allowbreak{}Mem.\allowbreak{}R3.\allowbreak{}Rw.\allowbreak{}UntwSat.\allowbreak{}dat\_\allowbreak{}saturated\_\allowbreak{}of\_\allowbreak{}untw}{Development.\allowbreak{}AbelianVarietyTheory.\allowbreak{}Proofs.\allowbreak{}IwasawaDecomposition.\allowbreak{}UntwistedAdelicDataSaturatedOverFixedField}{47--98}
\leanref{GlobalTheory.\allowbreak{}Mem.\allowbreak{}R3.\allowbreak{}Rw.\allowbreak{}exists\_\allowbreak{}pre2\_\allowbreak{}of\_\allowbreak{}fine\_\allowbreak{}of\_\allowbreak{}resChar\_\allowbreak{}mem}{Development.\allowbreak{}AbelianVarietyTheory.\allowbreak{}Proofs.\allowbreak{}IwasawaDecomposition.\allowbreak{}ConstructionStartingChoicesOfFineAdelicData}{195--241}

We now specify the integral state over these rows. In the row construction the parameters \(\eta_{j,v}\) range over nonnegative rationals. Take the function \((j,v)\mapsto0\). Let \(b_1,\ldots,b_a\) be the chosen complex row basis, let \(\iota:K\hookrightarrow\C\) be the coefficient embedding, and put \(e(x)=\sum_i\iota(x_i)b_i\) on \(K^a\). Write \(R_{\mu(v)}\) for the prescribed row set at the tuple selected by \(v\). In the common row coordinates take the module
\[
 \left\{x\in K^a:\ \forall v\ \exists d_v\in\mathcal O_K\setminus\mathfrak p_v,\
 d_vx\in
 \operatorname{span}_{\mathcal O_K}\{y\in K^a:e(y)\in R_{\mu(v)}\}\right\}.
\]
It is closed under addition since two admissible denominators at \(v\) can be multiplied. For each word take the inverse image of its level-zero lattice under the inclusion of its ramified subspace. Give this subspace the restricted norm. On the row space choose a frame for the given norm at the distinguished embedding and keep the prescribed frames at the other infinite places. Take all word heights equal to zero and the separate real error parameter equal to \(6x^{-C_F}\). For a row pairing whose values on these coordinates lie in \(\iota(K)\), pull the values back by \(\iota^{-1}\). Injectivity of \(\iota\) gives bilinearity and the required pairing identity. These are the division data for the chosen zero function.
\leanref{GlobalTheory.\allowbreak{}Mem.\allowbreak{}R3.\allowbreak{}Rw.\allowbreak{}prescribedRowLat}{Development.\allowbreak{}AbelianVarietyTheory.\allowbreak{}Proofs.\allowbreak{}IwasawaDecomposition.\allowbreak{}ExistsDivStageOfPk}{54--71}
\leanref{GlobalTheory.\allowbreak{}Mem.\allowbreak{}R3.\allowbreak{}Rw.\allowbreak{}exists\_\allowbreak{}divStage\_\allowbreak{}of\_\allowbreak{}pk}{Development.\allowbreak{}AbelianVarietyTheory.\allowbreak{}Proofs.\allowbreak{}IwasawaDecomposition.\allowbreak{}ExistsDivStageOfPk}{76--131}

For the Galois data retain the given length data and take the set indexing local contributions to be \(\varnothing\). At \(g_0,g_1\) use the actual places above the split primes, with the tame generators and series of the frame. Identify the resulting finite group with the original one by its place coordinates. Pull the two actions of the centre of the group algebra back along this identification, on the ambient rows and on their multiplicity spaces. The distinguished idempotents have equal coefficients under the identification, so their operators agree in both actions. This checks the conditions that still involve the group when the local index set is empty. The equation at every \((j,v)\) now reads
\[
                  \eta_{j,v}=\sum_{\alpha\in\varnothing}0=0 .
\]
The conditions on an individual local index have empty domain. In particular this construction keeps the tame group data at the two distinguished primes.
\leanref{GlobalTheory.\allowbreak{}MemN.\allowbreak{}emptyLayer}{Development.\allowbreak{}AbelianVarietyTheory.\allowbreak{}Proofs.\allowbreak{}IwasawaDecomposition.\allowbreak{}ExistsUnramifiedDataCharacter}{17--105}
\leanref{GlobalTheory.\allowbreak{}Mem.\allowbreak{}Top.\allowbreak{}trueLayer}{Development.\allowbreak{}AbelianVarietyTheory.\allowbreak{}Proofs.\allowbreak{}IwasawaDecomposition.\allowbreak{}TypeIdempotentCoeffAcrossLayers}{196--210}
\leanref{GlobalTheory.\allowbreak{}Mem.\allowbreak{}R3.\allowbreak{}Rw.\allowbreak{}Bare.\allowbreak{}exists\_\allowbreak{}nStage\_\allowbreak{}act}{Development.\allowbreak{}AbelianVarietyTheory.\allowbreak{}Proofs.\allowbreak{}IwasawaDecomposition.\allowbreak{}LocalizationFreeStageTrueGroupColumnChoiceStep}{72--144}
\leanref{GlobalTheory.\allowbreak{}Mem.\allowbreak{}R3.\allowbreak{}Rw.\allowbreak{}Bare.\allowbreak{}exists\_\allowbreak{}bare}{Development.\allowbreak{}AbelianVarietyTheory.\allowbreak{}Proofs.\allowbreak{}IwasawaDecomposition.\allowbreak{}LocalizationFreeStageTrueGroupColumnChoiceStep}{167--194}

Use the spectral stage for the same rows and the original Haar datum. Its hypotheses are positivity of the starting columns and the inequality between the row dimension and the ambient dimension, supplied by the column and row frames. Keep the original geometric stage. The division coordinates just constructed, the group identification and the two central actions now fit these stages simultaneously. Transport the row lifts along these identifications and read their maps in the transported bases. This adds the row realization to the later state with the chosen division and Galois components. This gives the initial state over \(w\) with empty local index set and zero length coordinates. The localization at the last finite level uses the twisted row spaces and their integral generators as in Lemma~\ref{B8.07t}.
\leanref{GlobalTheory.\allowbreak{}Mem.\allowbreak{}SProv.\allowbreak{}exists\_\allowbreak{}sStage\_\allowbreak{}true}{Development.\allowbreak{}AbelianVarietyTheory.\allowbreak{}Proofs.\allowbreak{}IwasawaDecomposition.\allowbreak{}ExistsSStageTrue}{17--26}
\leanref{GlobalTheory.\allowbreak{}Mem.\allowbreak{}R3.\allowbreak{}Rw.\allowbreak{}st\_\allowbreak{}of\_\allowbreak{}fits}{Development.\allowbreak{}AbelianVarietyTheory.\allowbreak{}Proofs.\allowbreak{}IwasawaDecomposition.\allowbreak{}FittingLaterStagesRowIsConstructionState}{22--28}
\end{proof}

\subsection*{Proof of \ref{B8.07b}}\label{prf:B8.07b}\phantomsection
\begin{proof}
We first choose weights \(w_j=\ell^{e_j}\) for the substitution \(Z_j=(1+T)^{w_j}-1\). For each \(s\) and each monomial \(Z^\alpha\), record the order of its coefficient in \(F_{M,s}\), capped at \(\mu_s+1\), assigning this last value to a zero coefficient. There are countably many entries, each with finitely many possible values. Successive extraction and a diagonal subsequence give an unbounded sequence of \(M\)'s on which each entry is eventually constant. Among the finitely many monomials of degree at most \(D\), at least one has limiting order at most \(\mu_s\). Let \(h_s\) be the least limiting order among all monomials. Among those with order \(h_s\), choose \(\alpha_s\) lexicographically least, comparing the last coordinate first.

Choose successive powers \(w_j\) so that
\[
 w_j>\sum_{i<j}(\alpha_s)_i w_i
 \quad\text{for every }s.
\]
Additional lower bounds on any \(w_j\) are permitted. Then \(\alpha_s\) uniquely minimizes the weight \(\sum_i\alpha_iw_i\) among the monomials of limiting order \(h_s\). Indeed, at the last differing coordinate any other such monomial has a larger exponent, and its gain there exceeds the entire earlier weight of \(\alpha_s\).

Choose the first \(R\) weights by this rule. Put \(m_*=b\sum_{i=1}^R w_i\), \(V=\max(\{0\}\cup\{v_y\})\), and \(L=(V+1)(m_*+1)\). Choose the last \(R\) weights sufficiently large that every coefficient of degrees \(1,\ldots,L\) in \((1+T)^{w_j}-1\) belongs to the \((V+1)\)-st power of the maximal ideal of \(O\). This is possible since
\[
 \binom{\ell^n}{d}=\frac{\ell^n}{d}\binom{\ell^n-1}{d-1}
 \qquad(1\le d\le\ell^n).
\]
For sufficiently large \(M\) the same divisibility holds for \((1+T)^{w_i\ell^M}-1\) at every first-half index \(i\).

Write \(P^*(T)=P(((1+T)^{w_j}-1)_j)\). Set \(W_s=\sum_j(\alpha_s)_jw_j\). Terms of total degree greater than \(W_s\) cannot contribute to \([T^{W_s}]F_{M,s}^*\). For large \(M\) on our subsequence, all coefficients through total degree \(W_s\) have their limiting capped orders. Divide this truncation by a uniformizer of \(O_0\) to the power \(h_s\) and reduce. In residue characteristic \(\ell\), the substitution becomes \(Z_j=T^{w_j}\). By uniqueness of the least weight, the coefficient of \(T^{W_s}\) in the reduction is the nonzero coefficient contributed by \(\alpha_s\). Thus \(F_{M,s}^*\) has a coefficient of order \(h_s\) at degree \(W_s\).

For the second family, take \(M\) large enough that \(\ell^M>m_*\). The first nonzero residue coefficient of \(A^*\) is at degree \(\sum_i n_iw_i\le m_*\). For a nonzero polynomial \(H=\sum a_jT^j\) over \(O\) put
\[
 \omega(H)=\min_j\bigl(\operatorname{ord}_O(a_j)+j/(m_*+1)\bigr).
\]
This quantity adds under multiplication. In each factor choose the least degree realizing the minimum. At the sum of these degrees the product of the two chosen coefficients is the unique contribution of minimum weight. Hence
\[
 \omega(A^*g^*)\le m_* /(m_*+1)+v_y<v_y+1 .
\]
It follows that \(A^*g^*\) has a coefficient of order at most \(v_y\) at a degree at most \(L\). Through degree \(L\), the coefficients of \(f^*-A^*g^*\) are divisible by the \((V+1)\)-st power of the maximal ideal, by the choice of the last weights and of \(M\). The same coefficient bound therefore holds for \(f^*\).

It remains to choose the order of the root of unity. Let \(W\) bound \(L\) and all \(W_s\). Consider any one of the resulting univariate polynomials \(H=\sum a_jT^j\) over its coefficient ring \(A\). It has a coefficient at some degree \(d\le W\) of order at most its prescribed bound \(v\). For \(t>0\), consider the minimum of the lines \(\operatorname{ord}_A(a_j)+jt\) with \(a_j\ne0\). As \(t\) decreases, the degrees of the successive minimizing lines increase. At most \(W+1\) of those degrees are at most \(W\). A minimizing degree greater than \(W\) must have intercept strictly less than \(\operatorname{ord}_A(a_d)\), by comparison with the line of degree \(d\). The successive intercepts strictly decrease. Thus at most \(W+v+1\) positive values of \(t\) can give a tie. This bound is independent of \(M\).

Choose \(k_0>\max(\{0\}\cup\{e_j\})\) so large that \(W/\varphi(\ell^{k_0})\le\delta\) and \(w_j/\varphi(\ell^{k_0})\le\delta\) for all \(j\). Take \(M\) on the subsequence large enough that the interval \([k_0,M]\) contains more integers than the sum of the tie bounds for all the polynomials. For a primitive \(\ell^k\)-th root, \(v_\ell(\zeta-1)=1/\varphi(\ell^k)\). Indeed, all primitive powers give differences from one of the same valuation, since their ratios to \(\zeta-1\) are units, and their product up to sign is \(\Phi_{\ell^k}(1)=\ell\). The numbers \(d_A/\varphi(\ell^k)\) are distinct as \(k\) runs through the interval. Choose \(k\) avoiding the ties in every polynomial at its corresponding value of \(t\). Each \(H(\zeta-1)\) now has a unique term of least valuation. It is nonzero and its valuation is at most
\[
 v/d_A+d/\varphi(\ell^k)\le v/d_A+\delta .
\]
The substitution gives \((1+(\zeta-1))^{w_j}-1=z_j\), proving both assertions for every primitive root of this order.
\end{proof}

\subsection*{Proof of \ref{B8.07c}}\label{prf:B8.07c}\phantomsection
\begin{proof}
Apply Lemma~\ref{B8.07b} to all the indicated polynomials with this \(\delta\). In the coefficient ring at level \(k\), take its designated primitive \(\ell^k\)-th root \(\zeta\) and set
\[
 \xi(g)=\prod_{j=1}^{2R}
            \bigl(\zeta^{\ell^{e_j}}\bigr)^{d_j(g)} .
\]
Each exponent is well defined modulo \(\ell^M\) since \(k\le M\). The character is continuous, anticyclotomic and trivial above \(\ell\) by the corresponding properties of the \(d_j\). At a place of the pair of \(P_j\), choose an inertia element on which \(d_j=1\). Every other coordinate is zero there. Thus the inertia image has order exactly \(\ell^{k-e_j}>1\), and all inertia images outside \(Q\) are trivial. The character factors through the finite abelian Artin quotient unramified outside \(Q\) with the indicated anticyclotomic and above-\(\ell\) conditions. At a place in the pair of \(P_j\), the chosen inertia element bounds the tame valuation by
\[
 v_\ell(\zeta^{\ell^{e_j}}-1)=\ell^{e_j}/\varphi(\ell^k)\le\delta .
\]
The same computation applies after raising to \(h\). Since \(|Q|\le4R\), the tame contribution is at most \(4Rt\delta\le e\).

Fix \(s,h\) and use \(u_j=(\zeta^{h\sigma_s})^{\ell^{e_j}}\) in the presentation. The resulting determinant is \(F_{M,s}((u_j-1)_j)\). The root \(\zeta^{h\sigma_s}\) is primitive of order \(\ell^k\), so Lemma~\ref{B8.07b} bounds the valuation of this nonzero determinant by \(\nu_a+\delta\). To compute the length, diagonalize the specialized matrix over \(A\) by elementary divisors. A diagonal entry of uniformizer order \(n\) has kernel \(\mathfrak m_A^{-n}/A\) on \(\operatorname{Frac}(A)/A\), of length \(n\). Adding the lengths and dividing by \(d_A\) gives the normalized determinant valuation. Thus every indicated Selmer group has finite length, and \(|\mathcal S|\delta\le e\). Taking \(\zeta^{h}\) or \(\zeta^{-h}\) in the second conclusion of Lemma~\ref{B8.07b} proves the assertion for the theta polynomials on the same \(Q\).
\end{proof}

\subsection*{Proof of \ref{B8.07e}}\label{prf:B8.07e}\phantomsection
\begin{proof}
At a split unramified prime of norm \(q\) outside the exceptional set, write the unshifted parameters as \(\alpha_j\). At the conjugate prime they are \(\alpha_j^{-1}\). Thus the eigenvalues of \(T\) at the first prime are \(q^{(1-t)/2}\alpha_j\), and those of its Tate dual at the conjugate prime are
\[
        q\big(q^{(1-t)/2}\alpha_j^{-1}\big)^{-1}
             =q^tq^{(1-t)/2}\alpha_j .
\]
These are the eigenvalues of \(T(t)\). The traces agree on \(G_E\). Indeed, modulo each power of the maximal ideal the two actions factor through a finite quotient. Apply Lemma~\ref{B8.42p} to that quotient, excluding the exceptional primes and those over rational primes ramified in the host. A degree-one prime so obtained splits over the maximal real subfield, since its local degree over \(\Q\) is one and it is unramified. Equality modulo every power gives equality in \(O_0\).

For the integral conclusion, let \(R_1,R_2\) denote the two representations being compared. The matrices in each image span \(M_r(O_0)\). For \(T(t)\) this follows by reduction and Nakayama, since twisting multiplies each matrix by a unit. Inversion in the group, transposition and scalar multiplication give the same span for the conjugate Tate dual. Extend both representations linearly to \(O_0[G_E]\). If \(R_1(x)=0\), then
\(\operatorname{tr}(R_2(x)R_2(y))=\operatorname{tr}(R_1(xy))=0\) for every \(y\). Testing against the matrix units gives \(R_2(x)=0\). Interchanging the two actions gives equality of their kernels. The resulting automorphism of \(M_r(O_0)\) is inner. To see this, the images of the diagonal matrix units decompose \(O_0^r\) into \(r\) isomorphic summands. Each is free of rank one over the local ring. Transport a generator of the first summand by the images of the units \(E_{i1}\). In the resulting basis the images of all the matrix units act as the original units. This gives the stated integral isomorphism for \(T\). Since \(\eta^c=\eta^{-1}\), tensoring it with \(\eta^{-1}\) gives the isomorphism for \(V\).

An isomorphism \(V\simeq V(t)\) would give \(\varepsilon_\ell^{tr}=1\) by determinants. Here \(tr>0\) by the rank inequalities. The cyclotomic image of a number field is infinite, since its intersections with the fields of \(\ell\)-power roots of unity have bounded degree and the degrees of those fields tend to infinity. The same determinant calculation in the residue field gives \(\overline{\varepsilon}_\ell^{\,tr}=1\) if the reductions are equivalent. This is precisely the condition that every \(\varepsilon_\ell(g)^{tr}-1\) be divisible by \(\ell\) in \(\Z_\ell\), which gives \(\nu_a>0\).
\end{proof}

\subsection*{Proof of \ref{B8.07f}}\label{prf:B8.07f}\phantomsection
\begin{proof}
Choose a finite set \(S\), disjoint from \(T_0\), containing the primes above \(\ell\) and prime ideal representatives generating the class group. The representatives may be chosen away from \(T_0\) by weak approximation applied to fractional ideals. Then the ring of \(S\)-integers has trivial class group. Adjoin \(\mu_n\) and the \(n\)-th roots of a finite set of generators of its unit group, and take \(H\) to be the group fixing the resulting splitting field. Finite generation follows from the unit theorem and the valuation map from \(S\)-units to \(\Z^S\). At every place in \(T_0\), \(n\) is a unit and the elements whose roots are being adjoined are units. Their roots lie in the maximal unramified extension by reduction and Hensel's lemma.

For \(P\) as in the statement, its residue field contains \(\mu_n\) and every \(S\)-unit reduces to an \(n\)-th power. Take a surjection from the cyclic multiplicative group of the residue field to \(\Z/n\). For an ideal prime to \(P\), principalize it over the \(S\)-integers and apply the surjection to the reduction of a generator. A change of generator multiplies by an \(S\)-unit. This defines an ideal character, trivial on principal generators congruent to one at \(P\) and on the primes in \(S\). By the Artin reciprocity assertion of Definition~\ref{B1.15ac}, it gives a Galois character unramified outside \(P\), trivial on decomposition groups above \(\ell\), and surjective on inertia at \(P\). The last assertion follows from the surjection on the residue units at \(P\). Subtract the conjugate character to obtain \(d_P\). Since \(P\ne cP\), the subtracted character is unramified at \(P\). At \(cP\) the roles reverse.

Put \(W=U\cap H\cap G_{E(\mu_{2n})}\). The element \(z\) belongs to \(W\) except possibly for the condition \(z\in U\): it belongs to \(H\) by construction and fixes \(\mu_{2n}\) since the prime below \(T_0\) is prime to \(2n\). Apply Lemma~\ref{B8.42p} in \(G_E/W\) to its class, successively excluding previously chosen pairs, the original exclusions and their conjugates, and all primes over ramified rational primes. This gives the first primes with the required splitting and norm. Inertia at each lies in \(W\). On this inertia \(d_{P_j}\) is surjective for its own index and all the other first characters vanish. Hence their product maps \(W\) onto \((\Z/n)^R\). Multiply \(h_i\) by an element of \(W\) so that its vector of first character values is the \(i\)-th coordinate vector. Apply Lemma~\ref{B8.42p} to the corrected element modulo \(W\cap\bigcap_j\ker d_{P_j}\), with the same finite exclusions enlarged at each choice. Its images modulo \(U\), modulo \(H\) and on \(\mu_{2n}\) have been preserved. This gives the second primes after all the first ones have been fixed.
\end{proof}

\subsection*{Proof of \ref{B8.07g}}\label{prf:B8.07g}\phantomsection
\begin{proof}
Write \(k_0\) for the residue field of \(O_0\). The ring \(\Lambda\) is complete local and finite free over \(O_0\). Indeed, modulo a uniformizer its augmentation ideal is generated by the nilpotent elements \([b_j]-1\), where \(b_j\) are the coordinate generators. Its residue field is \(k_0\), and \(\kappa\) reduces to one. Put \(W=T\otimes k_0\). The full residual matrix algebra and \(r\ge2\) give \(H^0(E,W)=H^0(E,W^\vee(1))=0\). The same invariant vanishing holds for any simple coefficient subquotient, and hence for every discrete coefficient module by a finite-length filtration and passage to unions.

Here are the local counts entering the proof. At \(v\mid\ell\), let \(d_v=\dim H^1_f(E_v,T\otimes\operatorname{Frac}(O_0))\). The space of rational finite cocycles has dimension \(r+d_v\), since rational coboundaries have dimension \(r\). Its integral-valued elements form a saturated lattice of this rank. Indeed, finitely many coordinate evaluations separate a basis of the rational space, so the integral-valued elements lie in a lattice. Continuity on the compact decomposition group lets one clear denominators in each basis vector. Saturation follows by testing the values. Reduction of this lattice modulo the uniformizer injects into the functions with values in \(W\): if all values are divisible by the uniformizer, divide the cocycle by it. Subtracting the dimension \(r-\dim H^0(E_v,W)\) of the residual coboundaries gives the order of the residual local condition \(L_v\):
\[
           |L_v|=|H^0(E_v,W)|\,|k_0|^{d_v}.
\]
The local dimension formula and the conjugate-embedding hypothesis give \(\sum_{v\mid\ell}d_v=r[E:\Q]/2\). Thus the product of the local ratios is \(|W|^{[E:\Q]/2}\).

The functor \(S\) is additive and left exact. Global injectivity in degree one for an inclusion of coefficients follows from invariant vanishing in the quotient. Above \(\ell\), choose finitely many coordinate evaluations which give full column rank on a basis of the lattice of finite cocycles modulo the uniformizer. A square minor is a unit. Evaluating at these coordinates recovers the coefficients of a finite cocycle over any \(B\). Hence finite cocycles pull back under inclusions and lift under surjections. Coboundaries belong to the finite condition. Outside the ramification set, inertia acts trivially and the unramified test is vanishing on inertia. At places with unrestricted condition there is no further pullback test. These observations give left exactness with the local conditions imposed. The functor commutes with increasing unions of finite submodules: a continuous cocycle on the compact group has finite image, and its finitely many local coefficient witnesses above \(\ell\) lie in a finite submodule.

For a simple extension \(0\to k_0\to B\to B''\to0\) we claim
\[
                 |S(B'')/\operatorname{im}S(B)|\le |S(k_0)|.
\]
Let \(\Sigma\) contain the primes above \(\ell\) and the finite places where ramification is allowed, including \(Q\). Work on the Galois group unramified outside \(\Sigma\). Write \(H_G^i\) for its cohomology with coefficients \(W\), and put
\[
 D_2=\ker\!\left(H_G^2\longrightarrow\prod_{v\mid\ell}H^2(E_v,W)\right),\qquad
 C_1=\operatorname{coker}\!\left(H_G^1\longrightarrow
                 \prod_{v\mid\ell}H^1(E_v,W)/L_v\right).
\]
The connecting class of any element of \(S(B'')\) belongs to \(D_2\), since its finite local representatives lift above \(\ell\). In a fibre of the connecting map, subtract one fixed class and choose a global lift of the difference. Compare it at the places above \(\ell\) with finite local lifts. Their differences come from local classes for \(W\). Choosing such differences gives an element of \(C_1\). Two choices of global classes in the fibre that give the same element of \(C_1\) differ by a Selmer lift: subtract a global \(W\)-class realizing the difference of the local corrections modulo \(L_v\). Consequently the number of cosets is at most \(|D_2|\,|C_1|\).

The map from \(H_G^2\) onto the displayed local product is surjective by Theorem~\ref{B6.09}, applied to the finite \(\ell\)-torsion module \(W\) over the totally complex field \(E\); its global Tate dual has no invariants. The local and global Euler characteristic formulas of Definition~\ref{B1.15ac} give
\[
 |H_G^1|=|H_G^2|\,|W|^{[E:\Q]/2},\qquad
 |H^1(E_v,W)|=|H^0(E_v,W)|\,|H^2(E_v,W)|\,|W|^{[E_v:\Q_\ell]} .
\]
Inserting these identities and the product of the local ratios in the definition of \(C_1\) gives \(|S(k_0)|=|D_2|\,|C_1|\). This proves the claim.

For completeness, let \(I\) be the continuous additive dual of \(\Lambda\) with values in \(\Q_\ell/\Z_\ell\). The pairing \((u,f)\mapsto f(u)\) identifies \(\operatorname{End}_\Lambda(I)\) with \(\Lambda\) by duality. Put \(M=S(I)^\vee\), using the same additive dual. Its quotient modulo the maximal ideal \(\mathfrak m\) is dual to \(S(k_0)\): the socle of \(I\) is one copy of \(k_0\), and its image under \(S\) is the simultaneous kernel of multiplication by generators of \(\mathfrak m\). Thus \(M\) has \(a=\dim_{k_0}S(k_0)\) minimal generators. To see that lifts generate, pass to finite quotients of the compact module \(M\) and apply Nakayama. The image of the resulting map \(\Lambda^a\to M\) is compact, hence closed.

For a finite \(B\), present \(B^\vee\) by finite free \(\Lambda\)-modules and dualize. This exhibits \(B\) as a kernel between finite powers of \(I\). Applying \(S\) and using left exactness identifies \(S(B)\) with \(\operatorname{Hom}_{\rm cont}(M,B)\), where homomorphisms are \(\Lambda\)-linear. The identification is evaluation: a map \(B\to I\) acts on \(S(B)\) by functoriality. It is therefore natural and extends to discrete torsion modules by the union property.

Write \(N=\ker(\Lambda^a\to M)\), so \(N\subset\mathfrak m\Lambda^a\). Embed \(k_0\) as the socle of \(I\). A map \(N\to k_0\) extends to a continuous map \(\Lambda^a\to I\) after this embedding. Indeed continuous \(\Lambda\)-maps to \(I\) correspond by evaluation at \(1\) to continuous additive characters. Such a character on \(N\) factors through a finite quotient induced from \(\Lambda^a\), on which it extends by divisibility of \(\Q_\ell/\Z_\ell\). Reduction modulo \(k_0\) now gives a map \(M\to I/k_0\). Its class modulo maps lifting to \(I\) is independent of the extension. This embeds \(\operatorname{Hom}_\Lambda(N,k_0)\) into \(S(I/k_0)/\operatorname{im}S(I)\). In the kernel, the original map \(N\to k_0\) extends to \(\Lambda^a\to k_0\) and hence vanishes on \(N\). The simple-extension inequality gives \(\dim N/\mathfrak mN\le a\). The ring is noetherian, so Nakayama gives at most \(a\) generators of \(N\). Use them as rows and append zero rows to make a square matrix. A tuple in \(B^a\) lies in its kernel precisely when the map \(\Lambda^a\to B\) that it defines factors through \(M\).
\end{proof}

\subsection*{Proof of \ref{B8.07h}}\label{prf:B8.07h}\phantomsection
\begin{proof}
Let \(V\) be the residual Selmer space with the conditions relaxed at \(Q\), obtained by augmentation of the matrix. Choose inertia elements \(\gamma_i\) at \(P'_i\) with \(d_{P'_i}(\gamma_i)=1\). At \(cP'_i\) use their conjugates. Evaluation at these elements and at the indicated Frobenius elements is well defined on classes, since coboundaries vanish on the two decomposition groups. The map
\[
 I:V\longrightarrow (k_0^r\oplus k_0^r)^R,\qquad
 I(x)_i=(x(\gamma_i),x(c\gamma_i c^{-1}))
\]
has kernel \(S_0\). At a second prime the maximal pro-\(\ell\) quotient of tame inertia is procyclic, and the specified element generates it. At the first primes there is no residual local cohomology.

By the sum of the local invariants in Theorem~\ref{B6.09}, there are \(b_i\in k_0^\times\) such that, for \(x\in V,\ y\in S'_0\),
\[
 \sum_i b_i\{x(\gamma_i)\cdot y(\phi_i)
          -x(c\gamma_i c^{-1})\cdot y(c\phi_i c^{-1})\}=0.
\]
Indeed the local conditions away from \(Q\) are orthogonal. At \(P'_i\) the dual class is unramified. The tame cup invariant therefore pairs the inertia value of \(x\) with the Frobenius value of \(y\). Its coefficient is nonzero by the uniformizer normalization in Theorem~\ref{B6.09}. Using the conjugate inertia and Frobenius elements at the other prime changes the coefficient by \(-1\), since \(c\) inverts the roots of unity. The formula over \(k_0\) follows from the prime-field formula by testing all scalar multiples and using the nondegenerate finite-field trace pairing.

Define \(J:V\to(k_0^r\oplus k_0^r)^R\) by evaluating on the two Frobenius elements and multiplying the \(i\)-th pair by \(a_i/b_i\). Then
\[
 W=\{v\in\ker I:Jv\in I(V)\}
\]
has dimension at most \(b\). For \(Iv'=Jv\), substitute \(v'\) for \(x\) in the reciprocity identity. It puts \(v\) in the left kernel from the hypothesis. Applying \(J\) from \(\ker I\) to the quotient by \(I(V)\) also gives \(\dim V\le 2rR+b\).

Work over \(K=k_0((Z))\), where \(Z\) is an indeterminate. Put \(\beta_i=Za_i/b_i\). In the polynomial variables for the second primes take the direction \(\beta_i\); in those for the first primes take the direction
\[
              1-\sum_j\beta_j\,\overline{d_{P'_j}(\phi_i)} .
\]
The bar means reduction from \(\Z/\ell^M\) to \(k_0\). Write these directions as \(v_j\). On substituting \(X_j=v_jD\) with \(D^2=0\), the universal inverse character becomes \(1+D\lambda\), with \(\lambda(\phi_i)=-1\) and \(\lambda(\gamma_i)=-\beta_i\). This uses the Kronecker delta in Lemma~\ref{B8.07f}. The values at the conjugate elements have the opposite signs.

A class lifting in this direction has residual value \(x\in V\otimes K\). Scalar extension here commutes with the residual Selmer condition: a continuous cocycle into a discrete vector space has finite image, and expanding its coordinates in a \(k_0\)-basis of \(K\) gives the old cocycle and local linear equations coordinate by coordinate. The tame relation \(\phi_i\gamma_i\phi_i^{-1}=\gamma_i^{q_i}\) in the pro-\(\ell\) inertia quotient, \(q_i=\mathrm NP'_i\), gives
\[
 \lambda(\phi_i)x(\gamma_i)=\lambda(\gamma_i)x(\phi_i).
\]
Indeed in a cocycle with action \(1+D\lambda\), the first-order term on the \(q_i\)-th power contributes \(\binom{q_i}{2}\lambda(\gamma_i)x(\gamma_i)\) in addition to \(q_i\) times the first-order value. Here \(q_i=1\) and \(\binom{q_i}{2}=0\) in \(k_0\), since \(q_i\equiv1\bmod 2\ell^M\). The same computation at the conjugate prime gives \(Ix=ZJx\).

The kernel of \(I-ZJ\) over \(K\) has dimension at most \(b\). Its intersection with \(V[[Z]]\) is a saturated submodule of a finite free \(k_0[[Z]]\)-module. Reduction of that intersection therefore has dimension equal to its rank. If \(v_0+Zv_1+\cdots\) belongs to it, then \(Iv_0=0\) and \(Iv_1=Jv_0\), so its reduction lies in \(W\).

Lift the matrix entries to polynomials and substitute \(X_j=v_jD\), now in \(K[[D]]\). Its constant matrix has nullity \(d=\dim V\). In elementary-divisor form there are \(d\) entries divisible by \(D\), counting zero entries. The number of entries divisible by \(D^2\), again including zero entries, equals the dimension of the subspace of its residual kernel lifting modulo \(D^2\). This follows by solving each diagonal equation modulo \(D^2\). That number is at most \(b\) by the preceding paragraph and naturality in Lemma~\ref{B8.07g}. For \(b=0\), the determinant has order \(d\le 2rR\). For \(b=1\), discard an entry of largest order. The remaining product has order at most \(\max(0,d-1)\le 2rR\). The ideal of cofactors is invariant under invertible row and column operations, so some cofactor has at most this order. A matrix of size one has cofactor \(1\).

Finally, in \(K[D]/(D^{4rR+1})\) the substitutions \([b_j]\mapsto1+v_jD\) respect the group relations because \((1+v_jD)^{\ell^M}=1\). Thus every polynomial representative of the detected entry has the same nonzero image in this quotient. If all coefficients through total degree \(4rR\) lay in \(\mathfrak m_0\), that image would be zero.
\end{proof}

\subsection*{Proof of \ref{B8.07i}}\label{prf:B8.07i}\phantomsection
\begin{proof}
On \(L\) the cocycles are additive, and \(x(gug^{-1})=U(g)x(u)\), with the corresponding identity for \(y\). The scalar actions of \(z\) on \(U,U'\) are \(d,d^{-1}\); those of \(\bar z\) are \(d^{-1},d\), by the displayed polarization. Apply the diagonal identity to \(uv\) and subtract its values at \(u,v\). Do this a second time with \(u\) replaced by \(zuz^{-1}\). The two cross terms then acquire the factors \(d,d^{-1}\), interchanged on the conjugate side. Since \(d-d^{-1}\) is a unit, solving the two equations gives
\[
                    x(u)\cdot y(v)=x(\bar v)\cdot y(\bar u).
\]
Replace \(u\) by \(gug^{-1}\). With column-vector notation put
\[
 B=x(u)y(v)^{\mathsf t},\qquad
 B'=J^{-1}y(\bar u)x(\bar v)^{\mathsf t}J,\qquad \theta=\varepsilon_\ell^t .
\]
The cross-term identity reads
\[
            \operatorname{tr}(U(g)B)=
            \operatorname{tr}\big(\theta(g)U(g)B'\big).
\]

Let \(f,f':A[G_E]\to M_r(A)\) be induced by \(U,\theta U\). Both are surjective by the span hypothesis. Put \(I=\{a\in A:aB=0\}\). For \(v_0\in\ker f'\), apply the trace identity to \(p v_0 q\) with \(f(p),f(q)\) arbitrary matrix units. It gives \(f(v_0)_{ij}B_{kl}=0\) for all indices. Hence \(f\) modulo \(I\) factors through \(f'\). This factor map kills \(I M_r(A)\) and gives an \(A/I\)-algebra endomorphism of \(M_r(A/I)\) carrying \(f'(g)\) to \(f(g)\) modulo \(I\). If \(I=A\) the assertion already holds. Otherwise \(A/I\) is local. The matrix-unit argument in Lemma~\ref{B8.07e} applies to a unital endomorphism as well: the images of the diagonal units give \(r\) isomorphic projective summands of total rank \(r\). Consequently this endomorphism preserves determinants. Since \(\det U(g)\) is a unit,
\[
                              \theta(g)^r-1\in I .
\]
Choose \(g\) attaining the minimum in Definition~\ref{B8.01k}(7). The minimum is finite by the cyclotomic argument in Lemma~\ref{B8.07e}. In \(\Z_\ell\) one has \(\theta(g)^r-1=\ell^{\nu_a}u_0\) with \(u_0\) a unit. Its image annihilates \(B\), proving the formula entry by entry.
\end{proof}

\subsection*{Proof of \ref{B8.07j}}\label{prf:B8.07j}\phantomsection
\begin{proof}
Set the variables in \(J\) to zero and reduce modulo \(\varpi\). Some homogeneous part of degree at most \(w\) is nonzero. Over the algebraic closure of the residue field choose nonzero values for the remaining variables on which that part is nonzero. Such a point exists since the product of the homogeneous polynomial and the remaining coordinate functions is a nonzero polynomial over an infinite field. All the coordinates lie in a finite extension \(k'\) of the residue field.

Take a monic lift of the minimal polynomial of a generator of \(k'\). Adjoining a root to \(O_0/(\varpi^{d+1})\) gives a finite free algebra \(C\) with reduction \(k'\). Its ideal \(\varpi C\) is nilpotent, so \(C\) is local. Lift the chosen coordinates to units of \(C\), and put \(A=C[u]/(u^{2L+2})\). Reduction of \(G((a_i u)_i)\) modulo \(\varpi\) is \(u^j\) times a unit for some \(j\le w\). Modulo \(G((a_i u)_i)\) this gives \(u^j\in\varpi A/(G((a_i u)_i))\). Raising to the power \(d+1\) gives the divisibility. Finally \(1,u,\ldots,u^{2L+1}\) form a free \(C\)-basis and \(\ell^{\nu_a}\) is a unit times \(\varpi^d\), nonzero in \(C\).
\end{proof}

\subsection*{Proof of \ref{B8.07k}}\label{prf:B8.07k}\phantomsection
\begin{proof}
Use the parameter of the finite-type row occurring at \(s\), its coefficient embedding at the selected place, and the Galois realization of this parameter. These are taken from the same column. In the Galois comparison there is a scalar character between the row realization and \(\rho_s\). Reducing their matrix identity gives the identity of residual characteristic polynomials with this same scalar. At a type prime different from \(\ell\) the shift character is unramified. The inertia value of the block character there is the prescribed tame value or its inverse, according to the embedding. At the conjugate place these values are interchanged, by Artin reciprocity and the triviality of the anticyclotomic character on the real idèles.

Here is the residual image argument. At a designated place above a type prime \(g\), the label has degree \(r\), the dimension of the row parameter. The label counts the tame characters on the wild-invariant part. Equality of the degrees exhausts the dimension, so wild inertia has characteristic polynomial \((X-1)^r\). In the residual realization its image is a finite \(g\)-group if \(g\ne\ell\). A unipotent element of such a group in characteristic \(\ell\) is the identity. The eigenvalues of the tame generator are consequently the readings of the prime-to-\(\ell\) parts of the roots of the label. Choose \(g\) for which \(\ell\) is outside its exceptional set, using the disjointness in Lemma~\ref{B7.50e}. The eigenvalues are distinct. Over an algebraic closure of \(k_0\), a globally invariant subspace selects a subset of the eigenlines at every designated place above \(g\). Frobenius carries the tame generator to its \(g\)-th power modulo wild inertia. The subset therefore gives a divisor over \(\F_g\) of the label at every such place, of the dimension of the subspace. Its complement gives the complementary degree. The good-row condition in Lemma~\ref{B7.51} excludes a common degree between \(1\) and \(r/2\). This proves absolute irreducibility. This is the factor-degree argument of Lemma~\ref{B8.42pa} applied to the row's own residual realization. Scalars preserve its invariant subspaces. By the density theorem for a simple module, the algebra generated by the image over \(\overline{k}_0\) is the full matrix algebra. The linear span is already an algebra since the image is a group. Descent of dimension gives the full span over \(k_0\).

At \emph{either} type prime different from \(\ell\), the other root exclusion of the good row says that the tame generator has no eigenvalue equal to the block-character value. Thus \(T\) has no residual inertia invariants. The same holds for \(T'\), since the cyclotomic character is trivial there on inertia and inverse transpose preserves the exclusion of \(1\).

At \(D_0\), \(\rho_s\) is unramified and the inertia character in \(T_s\) is the prescribed order-fifteen character or its inverse. Choose \(p=3\) if \(\ell\ne3\), and \(p=5\) otherwise. At a place in \(D_0\) take an element mapping to order \(p\) under that character. In the procyclic closure of this element retain its prime-to-\(\ell\,\mathrm N v\) component. The order-\(p\) image survives since \(p\) is prime to \(\ell\,\mathrm N v\). This gives \(z\). Its character value and its square have nontrivial reduction, since reduction is injective on roots of unity of order prime to \(\ell\). The cyclotomic character is unramified at \(v\). This proves the scalar assertion. The same construction at each place in \(D_0\) proves the invariant exclusion there. At all the exceptional places in question local Tate duality now gives \(H^2=0\) on the residual modules. The local Euler characteristic away from \(\ell\) gives \(H^1=0\). These are the finite-module formulas of Definition~\ref{B1.15ac}.
\leanref{GlobalTheory.\allowbreak{}WildFull.\allowbreak{}main}{Development.\allowbreak{}AbelianVarietyTheory.\allowbreak{}Proofs.\allowbreak{}InertiaEigenvalues.\allowbreak{}FullDegreeLabelForcesTrivialWildInertia}{14--21}
\leanref{GlobalTheory.\allowbreak{}TwSys.\allowbreak{}exists\_\allowbreak{}tameInertiaElementAtAuxiliaryPrime}{Development.\allowbreak{}AbelianVarietyTheory.\allowbreak{}Proofs.\allowbreak{}InertiaEigenvalues.\allowbreak{}ExistsTameInertiaElementAtAuxiliaryPrime}{85--143}
\leanref{GlobalTheory.\allowbreak{}TwSys.\allowbreak{}TameScalar.\allowbreak{}exists\_\allowbreak{}scalar\_\allowbreak{}tmat\_\allowbreak{}of\_\allowbreak{}tame\_\allowbreak{}inertia}{Development.\allowbreak{}AbelianVarietyTheory.\allowbreak{}Proofs.\allowbreak{}InertiaEigenvalues.\allowbreak{}TameInertiaElementActsByScalar}{29--61}

For the characteristic-zero local comparison, use the cuspidal row first with the inverse of the frame character and then with the layer character. The Galois shifts for \(\rho_s\) and \(T_s\) are respectively \((1-r-t)/2\) and \((1-t)/2\). The split Hecke polynomial identities identify these semisimple representations with their attachments. Indeed, the scalar twists preserve semisimplicity, and split Frobenius elements are dense by Lemma~\ref{B8.42p}. At a spherical place the local parameter is unramified. At a type place its full degree and separated labels give zero monodromy and the tame parameter, with the same properties after the prescribed tame character twist. Thus the local clauses of \ref{B9.06} and \ref{B9.07} apply also when the coefficient prime is a type prime. They give de Rham comparison above \(\ell\), and unramifiedness of the corresponding integral actions at spherical places away from \(\ell\). Passing from one prolongation of a place to another conjugates the inertia group.

We compute the filtration. At a CM embedding, take the row
\[
       y\in\{v_0,\ldots,v_{k-2},-b_0,\ldots,-b_{m-1}\}
\]
with multiplicities, in the notation of Lemma~\ref{B4.02a}. For layer \(j+1\), the two conjugate infinitesimal parameters are \(y+e_j\) and \(-(y+e_j)\). Put \(h=(t+3)/2\). The inequalities of Lemma~\ref{B4.03a} give \(v_i>-b_0\), \(e_j-b_i\ge h\) for \(i\le j\), and \(e_j-b_i\le-h\) for \(i>j\). In particular \(|y+e_j|\ge h\). If \(z_\tau\) is one of the two infinitesimal parameters, its filtration degree on \(T\) is
\[
                         a_\tau=-z_\tau+(t-1)/2 .
\]
These degrees are integers by regular algebraicity and the odd ranks. The attachment identifies their multiplicities with the dimensions of the graded pieces. Every degree is at most \(-2\) or at least \(t+1\), and the conjugate multiset is the image under \(a\mapsto t-1-a\). In each matched pair exactly one degree is negative. Since the multiset has cardinality \(r\), the two dimensions of \(D_{\rm dR}/\mathrm{Fil}^0\) sum to \(r\). A rational invariant would supply a degree-zero line. For the Tate dual the degrees are \(-1-a_\tau\), which also exclude zero. Thus both rational invariant spaces vanish above \(\ell\). The local Bloch--Kato formula in \ref{B8.41h} applies to these de Rham representations.
\leanref{GlobalTheory.\allowbreak{}TripleLocal.\allowbreak{}rowFacts}{Development.\allowbreak{}AbelianVarietyTheory.\allowbreak{}Proofs.\allowbreak{}InertiaEigenvalues.\allowbreak{}RankOneCuspidalUnitaryCharEquivApplySymm}{260--303}
\leanref{GlobalTheory.\allowbreak{}TripleGal.\allowbreak{}Jumps.\allowbreak{}exists\_\allowbreak{}jumpsAre\_\allowbreak{}goodJ\_\allowbreak{}n}{Development.\allowbreak{}AutomorphicArithmetic.\allowbreak{}Foundations.\allowbreak{}GaloisRepresentations.\allowbreak{}StringRepresentations.\allowbreak{}FrobDualOfTripleAttN}{85--134}
\leanref{GlobalTheory.\allowbreak{}TripleGal.\allowbreak{}negCount\_\allowbreak{}of\_\allowbreak{}goodJ}{Development.\allowbreak{}AbelianVarietyTheory.\allowbreak{}Proofs.\allowbreak{}InertiaEigenvalues.\allowbreak{}ConjugateNegativeJumpCountOfGoodHodgeJumps}{26--36}

At a remaining place of norm \(q\), every eigenvalue of \(T(\Frob)\) is \(q^{(1-t)/2}\alpha\). The unramified bound of Definition~\ref{B1.15ac} for the unitary cuspidal row implies \(|\alpha|\le q^{(r+3)/2}\). Since \(t\ge r+5\), the resulting absolute value is less than one. An invariant of \(T\) would require eigenvalue \(1\), and an invariant of \(T'\) would require eigenvalue \(q\) on \(T\). At the exceptional places away from \(\ell\) the residual inertia exclusion already implies the rational exclusion, by scaling a fixed vector to be primitive. This proves the local invariant assertions at all places. Lemma~\ref{B8.07e} applies to the conjugate self-dual row with the shift just computed and gives its integral polarization.

Finally an abelian label gives a positive integer \(u\) whose double is its assigned parameter. The range in the row data is \(1\le u\le k+m(t+b_r+1)+1\), with \(b_r\) as in Definition~\ref{B8.07d}. The empty-conductor coefficient sum is a nonzero common scalar times the reading of the toric sum at \(u\). Unfolding the toric sum of Definition~\ref{B1.14k} at empty auxiliary level gives \(L(u,\chi)\). The character is unitary and has angular exponents \(2(u+\kappa_\tau)\) on the CM type, with \(\kappa_\tau\ge0\). Lemma~\ref{B8.11h} gives nonvanishing, including \(u=1\). The member's rank and type-prime conditions give \(r\ge100\) and \(2r\le g_i\). Together with the polarization and filtration computation these prove the three index conditions.
\leanref{GlobalTheory.\allowbreak{}TripleGal.\allowbreak{}pure\_\allowbreak{}of\_\allowbreak{}tripleAtt\_\allowbreak{}n}{Development.\allowbreak{}AutomorphicArithmetic.\allowbreak{}Foundations.\allowbreak{}GaloisRepresentations.\allowbreak{}StringRepresentations.\allowbreak{}FrobDualOfTripleAttN}{256--282}
\leanref{GlobalTheory.\allowbreak{}LamTie.\allowbreak{}sum\_\allowbreak{}th\_\allowbreak{}empty\_\allowbreak{}ne\_\allowbreak{}zero}{Development.\allowbreak{}AbelianVarietyTheory.\allowbreak{}Proofs.\allowbreak{}InertiaEigenvalues.\allowbreak{}SumThEmptyNeZero}{10--24}
\end{proof}

\subsection*{Proof of \ref{B8.07l}}\label{prf:B8.07l}\phantomsection
\begin{proof}
Write \(W=T\otimes k_0\), \(W'=T'\otimes k_0\), and use a generator of the residual Tate line to write the pairing as a dot product. At empty auxiliary conductor the Selmer spaces \(S,S'\) for these two representations are finite and of equal dimension. To see which dual occurs, let \(\varpi\) be a uniformizer. At \(v\mid\ell\), embed \(W'\) into \(T'\otimes(\operatorname{Frac}(O_0)/O_0)\) by division by \(\varpi\). A class annihilating the reductions of all integral finite cocycles maps into the annihilator of the integral finite condition. Integral and divisible Bloch--Kato duality in \ref{B8.41h} identifies this annihilator with the image of the rational finite condition. Choose a rational finite cocycle giving the image and adjust it by a rational coboundary lifting the difference of representatives in the divisible module. Its values now agree modulo the lattice with those of the residual cocycle divided by \(\varpi\). Multiplying by \(\varpi\) gives an integral finite lift. The reverse inclusion follows by the same pairing and its compatibility with division by \(\varpi\). This proves residual orthogonality, using the finite-field trace pairing to pass between additive and \(k_0\)-valued tests. At exceptional places away from \(\ell\) both local \(H^1\)'s vanish by Lemma~\ref{B8.07k}. Elsewhere the unramified conditions are orthogonal. The two global \(H^0\)'s vanish by absolute irreducibility and \(r\ge2\). In the finite Selmer order formula of \ref{B6.09} and \ref{B1.15ac}, the product of the local ratios above \(\ell\) is \(|W|^{[E:\Q]/2}\) by Lemma~\ref{B8.07g}. This cancels the factor at the complex places. The remaining ratios are one. Hence \(|S|=|S'|\), which gives equality of dimensions over the finite field.
\leanref{GlobalTheory.\allowbreak{}ResDual.\allowbreak{}card\_\allowbreak{}eq}{Development.\allowbreak{}AbelianVarietyTheory.\allowbreak{}Proofs.\allowbreak{}InertiaEigenvalues.\allowbreak{}ResidualSelmerFinrankEqDual}{155--250}
\leanref{GlobalTheory.\allowbreak{}ResDual.\allowbreak{}Perp.\allowbreak{}mem\_\allowbreak{}bkOn\_\allowbreak{}of\_\allowbreak{}orthogonal}{Development.\allowbreak{}AbelianVarietyTheory.\allowbreak{}Proofs.\allowbreak{}InertiaEigenvalues.\allowbreak{}MemDualBlochKatoOfOrthogonalResidual}{57--99}

Write \(\bar g=cgc^{-1}\), where \(c\) denotes complex conjugation. Put \(I=(\ell^{\nu_a})\mathfrak m_0\). Inside the ray subgroup \(H\) of Lemma~\ref{B8.07f}, intersect the kernels of the \(T_s\bmod I\), the residual cyclotomic kernel and the subgroup fixing \(\mu_{2\ell^M}\). Intersect also with the conjugate of the resulting subgroup, obtaining \(L\). It is open, normal and stable under \(c\). It contains every commutator \([g,z]\). Indeed, \(H\) contains the designated inertia and is normal, the cyclotomic quotients are abelian, and both \(z\) and \(\bar z\) act by scalars in the indicated matrix quotients. On \(L\), \(W\) and \(W'\) are trivial. If a cocycle \(x\) of either vanishes on \(L\), and the action of \(z\) is the scalar \(d\), comparison of \(x(gz)\) with \(x([g,z]zg)\) gives
\[
                  (d-1)x(g)=(g-1)x(z).
\]
Here \(d-1\) is invertible by Lemma~\ref{B8.07k}, also for the dual scalar. Thus \(x\) is a coboundary. Independent classes restrict to linearly independent vector-valued functions on \(L\). On this subgroup we have \(x(gug^{-1})=g x(u)\).
\leanref{GlobalTheory.\allowbreak{}ResSel.\allowbreak{}SmallSub.\allowbreak{}exists\_\allowbreak{}small}{Development.\allowbreak{}AbelianVarietyTheory.\allowbreak{}Proofs.\allowbreak{}InertiaEigenvalues.\allowbreak{}ExistsOpenNormalSubgroupUnderInertia}{132--185}

For \(u\in L\) put \(B_u(x,y)=x(u)\cdot y(u)-x(\bar u)\cdot y(\bar u)\). First suppose \(W\) and \(g\mapsto W'(\bar g)\) are inequivalent over \(k_0\). Take cocycles \(x_i,y_j\) lifting bases of \(S,S'\). Suppose \(C=(C_{ij})\) annihilates every matrix \((B_u(x_i,y_j))\) under entrywise pairing. Denote \(\sum_{ij}C_{ij}X_i\cdot Y_j\) by \(D(X,Y)\). Polarize the equality at \(uv\) and subtract the equalities at \(u,v\). Replace \(u\) by \(gug^{-1}\), first with \(g\), then with \(gz\). In the second equality the terms containing \(x(u)\) and \(y(u)\) acquire respectively the scalars \(d_s,d_s^{-1}\), interchanged on the conjugate side. Subtraction and division by \(d_s-d_s^{-1}\) give
\[
 D(W(g)x(u),y(v))=D(x(\bar v),W'(\bar g)y(\bar u)).
\]
The two sides are linear functionals of the respective representation matrices. A nonzero common matrix coefficient of this form would identify the two representations. Indeed, both group-algebra maps onto \(M_r(k_0)\) are surjective. If an element of the kernel of the second has nonzero image under the first, the ideal generated by this image contains all matrix units. Applying the equality to left and right multiples forces the first functional to vanish on every unit matrix \(E_{ij}\). Otherwise the first map factors through the second and the resulting automorphism of \(M_r(k_0)\) is inner by the matrix-unit argument of Lemma~\ref{B8.07e}. Thus \(D(Ax(u),y(v))=0\) for every matrix \(A\). Taking \(A=E_{pq}\) gives
\[
                 \sum_{ij}C_{ij}x_i(u)_q y_j(v)_p=0 .
\]
Independence on \(L\), first for the \(y_j\) and then for the \(x_i\), gives \(C=0\). The matrices of the \(B_u\) span all matrices, in particular the identity in the two equal-dimensional bases.

In the equivalent case choose an invertible matrix \(J\) with \(W'(\bar g)J=JW(g)\). Then
\[
                   W(\bar g)^{\mathsf t}JW(g)=\bar\varepsilon_\ell(g)J ,
\]
where \(\bar\varepsilon_\ell\) denotes reduction of the cyclotomic character. Replacing \(g\) by \(\bar g\) and transposing gives this equation for \(J^{\mathsf t}\) too. Hence \(J^{-1}J^{\mathsf t}\) centralizes the full matrix algebra. Write \(J^{\mathsf t}=bJ\). Transposing twice and taking determinants give \(b^2=b^r=1\). Since \(r\) is odd, \(b=1\), in every residue characteristic.

The map \(x\mapsto (g\mapsto Jx(\bar g))\) takes cocycles and coboundaries for \(W\) to those for \(W'\). Transport \(S'\) back by its inverse to a space \(S_2\), and choose cocycles \(x_i\) lifting a basis of \(S+S_2\). Use \(y_i(g)=Jx_i(\bar g)\). For an annihilating matrix \(C\), put
\[
 P(u,v)=\sum_{ij}(C_{ij}-C_{ji})x_i(u)^{\mathsf t}Jx_j(\bar v).
\]
Symmetry of \(J\) gives \(P(u,u)=0\). Additivity on \(L\) gives \(P(u,v)+P(v,u)=0\). Conjugating \(v\) by \(z\) gives \(d_s^{-1}P(u,v)+d_s P(v,u)=0\), hence \(P(u,v)=0\). Conjugate \(u\) by arbitrary \(g\). Full matrix span again permits every matrix unit on \(x_i(u)\). Independence on \(L\) for \(x_i\), invertibility of \(J\), and independence for \(x_j\circ c\) then give \(C_{ij}=C_{ji}\). Every alternating matrix annihilates every symmetric matrix under entrywise pairing, also in characteristic two since its diagonal is zero. By finite-dimensional duality the \(B_u\) therefore span all alternating forms on \(S+S_2\).

There is such a form with left kernel of dimension at most one on \(S\times S_2\). When \(S=S_2\), pair consecutive basis vectors in alternating planes. Otherwise take \(v\in S\setminus S_2\), \(w\in S_2\setminus S\). Choose \(\alpha,\beta\) with \(\alpha(v)=\beta(w)=1\), \(\alpha|_{S_2}=\beta|_S=0\). Pair \(v,w\) by \(\alpha\wedge\beta\) and apply induction to \(S\cap\ker\alpha,\ S_2\cap\ker\beta\) in \(\ker\alpha\cap\ker\beta\). The added pair has zero cross terms with that subspace. This proves the assertion, including dimension zero. By Lemma~\ref{B8.07e}, the equivalent case requires \(\nu_a>0\).
\leanref{GlobalTheory.\allowbreak{}ResSel.\allowbreak{}SpanOddPf.\allowbreak{}span\_\allowbreak{}abstract}{Development.\allowbreak{}AbelianVarietyTheory.\allowbreak{}Proofs.\allowbreak{}InertiaEigenvalues.\allowbreak{}CocycleFormsSpanOddResidueChar}{96--180}
\leanref{GlobalTheory.\allowbreak{}ResSel.\allowbreak{}AltSpan.\allowbreak{}symm\_\allowbreak{}of\_\allowbreak{}forms}{Development.\allowbreak{}AbelianVarietyTheory.\allowbreak{}Proofs.\allowbreak{}InertiaEigenvalues.\allowbreak{}SpanTwoOfGaloisHypothesis}{110--176}
\leanref{GlobalTheory.\allowbreak{}AltCorank.\allowbreak{}exists\_\allowbreak{}alt\_\allowbreak{}corank\_\allowbreak{}le\_\allowbreak{}one}{Development.\allowbreak{}AbelianVarietyTheory.\allowbreak{}Proofs.\allowbreak{}InertiaEigenvalues.\allowbreak{}ExistsAlternatingFormLeftCorankLeOne}{108--114}

For each signature express the required form as a linear combination of a basis selected from the \(B_u\). In the first case at most \((\dim S)^2\) elements suffice, and in the second at most \((\dim(S+S_2))^2\). Pad each list to that size with identities. These sizes are fixed before \(L\) or \(M\). Concatenate the lists and call their total size \(R_0\). In \(L\) intersect the kernels of all the chosen basis cocycles and their conjugates. This subgroup \(U\) is open by continuity and normal by the conjugation formula. If \(g\phi g^{-1}\equiv u\bmod U\), then \(B_u=\bar\varepsilon_\ell(g)B_\phi\) on the associated bases. The scalar is a unit, so the span is preserved. Apply Lemma~\ref{B8.07f}, intersecting \(U\) with the additional prescribed open normal subgroup. At the second primes inertia and Frobenius belong to \(L\), hence so does the decomposition group. At each first prime residual Frobenius on \(W\) is \(d_s\) or \(d_s^{-1}\), and its norm is \(1\bmod\ell\). Inertia is trivial, and neither \(W\) nor \(W'\) has a fixed vector. The local Euler formula gives \(H^1=0\). This verifies the other two local hypotheses of Lemma~\ref{B8.07h}.
\end{proof}

\subsection*{Proof of \ref{B8.07m}}\label{prf:B8.07m}\phantomsection
\begin{proof}
Set \(Y_n=0\) in the new group ring and reduce modulo \(I\), obtaining the old group ring over \(O_0/I\), denoted by \(A\). For an old Frobenius reading at the new prime, let \(b\) be the vector of its old ray coordinates. Put
\[
            p=\prod_j(1+Z_j)^{b_j}-1 ,
\]
where \(Z\) lists the old variables and the exponents are their least nonnegative representatives. Its constant term is zero and its \(X_n\)-coefficient is one. Let \(\Phi,\Phi'\) be an old and a new presentation over \(A\), respectively, and let \(M_\Phi,M_{\Phi'}\) be the free modules on their column indices modulo their row relations. For an \(A\)-module \(B\), homomorphisms from these modules to \(B\) are the respective Selmer classes. Here \(A\) is finite and killed by \(\ell^{\nu+1}\), so all the coefficient modules in this description may carry the discrete topology.

At the new pair the old character is unramified, and \(T\) is trivial over \(A\) on both decomposition groups. Inertia values of a relaxed class are therefore well defined on classes and are determined by one pro-\(\ell\) tame generator at each place. They give \(2r\) elements of the representing module \(M_{\Phi'}\), by applying evaluation to its universal class. Killing these elements imposes precisely the old unramified condition, so the quotient is \(M_\Phi\). Each element is killed by \(p\). Indeed, if \(\sigma,\gamma\) are Frobenius and tame inertia, the tame relation on a cocycle gives \(\kappa(\sigma)x(\gamma)=(\mathrm N v)x(\gamma)=x(\gamma)\). At the conjugate place the old Frobenius character is inverted. The two differences from one generate the same ideal \((p)\).

For \(i=0,1\) this gives
\[
       p^{2r}\operatorname{Fitt}_i(M_\Phi)
                    \subseteq \operatorname{Fitt}_i(M_{\Phi'}).
\]
To see the inclusion, use generators for \(M_\Phi\), lifted to \(M_{\Phi'}\), and adjoin the \(2r\) inertia generators. An old relation lifts after adding a linear combination of the latter. Include the relations \(p\) times each inertia generator. Taking a minor with all \(2r\) of these last rows and columns multiplies the corresponding old minor by \(p^{2r}\). Fitting ideals can be computed in these generators since the modules are finitely presented over the commutative ring \(A\). In the original square presentations the two ideals are generated by the determinant and the cofactors respectively. An old presentation of positive size can always be used, by adjoining a generator with relation \(1\) if necessary.

For the determinant calculation quotient the polynomial ring by \(I\) and all monomials of degree \(\ge c_0+2r\). For the cofactor calculation quotient instead by \(\mathfrak m_0\), those monomials, and \(X_j\) for \(j>n\). Both quotients receive the indicated maps from \(A\): for \(0<k<\ell^{M-\nu-1}\),
\[
 v_\ell\binom{\ell^M}{k}\ \ge\ M-v_\ell(k)\ \ge\ \nu+1 .
\]
This follows by multiplying the binomial coefficient by \(k\). Hence the group relations \((1+Z_j)^{\ell^M}=1\) hold in both quotients.

If the new determinant failed its test, a representative with all the indicated coefficients in \(I\) would give zero in the first quotient after setting \(Y_n=0\). If the new positive-size matrix failed its cofactor test, one could choose a failing representative for each cofactor. All would give zero in the second quotient. In either case the inclusion would kill \(p^{2r}F\) for an old representative \(F\) passing the relevant test. After the indicated variables are set to zero, \(F\) has a nonzero homogeneous term of least degree \(<c_0\). The linear part of \(p\) is monic in \(X_n\), so multiplication by it is injective on the polynomial ring even over \(O_0/I\). The least nonzero degree in \(p^{2r}F\) is therefore \(<c_0+2r\), a contradiction. The new size-zero case of the cofactor test holds by definition.
\end{proof}

\subsection*{Proof of \ref{B8.07n}}\label{prf:B8.07n}\phantomsection
\begin{proof}
Choose a failing presentation and determinant representative. Its size is positive, since augmentation of a representative of the empty determinant is \(1\). Choose a representative of the cofactor specified by the old test. Apply Lemma~\ref{B8.07j} with the unused first variables set to zero. All other variables specialize to \(a_i u\) with \(a_i\) a unit. These substitutions define a group-ring map: the binomial divisibility in Lemma~\ref{B8.07m} handles the low powers, and the high powers vanish by nilpotence. Every term of the failing determinant vanishes, by its coefficient in degrees \(<2l+2\) and by its degree thereafter. This gives all the displayed assertions. Write \(C=(\operatorname{adj}\Psi)_{ij}\), \(\lambda=\ell^\nu u^l\) and \(b=\lambda u\). Since \(C\mid u^l\), the inequality \(\ell^\nu u^{2l+1}\ne0\) gives \(bC\ne0\).

Let \(Y=\operatorname{Hom}_{\Z}(A,\Z/\ell^{\nu+1})\), with \(A\) acting on the argument. Evaluation is a perfect balanced pairing since \(A\) is finite and killed by \(\ell^{\nu+1}\). Denote by \(S(Y)\) the old Selmer group, and by \(D(A)\) its group of dual classes. The dual condition is unramified away from the indicated ramification set, zero at the unrestricted places, and the cup-product annihilator above \(\ell\). For this pair, and for the induced perfect pair \(Y/bY,\ A[b]\), the finite order formula gives
\[
                 |S(Y)|=|D(A)|\prod_{v\in Q}|H^0(E_v,T'\kappa^{-1}\otimes A)|.
\]
Here \(\kappa\) is the specialized inverse ray character and the tensor notation includes the original \(O_0\)-action on \(T'\). For the second formula replace \(Y,A\) by \(Y/bY,A[b]\). Indeed, the dot product with evaluation is perfect with values in the Tate line. Global invariants vanish on either side by the scalar at \(z\). At \(v\mid\ell\) the free lattice of integral finite cocycles in Lemma~\ref{B8.07g} evaluates injectively after tensoring with either coefficient module, by its unit evaluation minor. Dividing by coboundaries gives local ratio \(|Y|^{d_v}\), or \(|Y/bY|^{d_v}\); their products cancel the complex-place factors. At a relaxed place away from \(\ell\) the Euler formula leaves the order of the dual invariants. At the fixed exceptional places these are zero by Lemma~\ref{B8.07k} and a composition series. These observations give the displayed formula by \ref{B6.09}.
\leanref{GlobalTheory.\allowbreak{}SelOrder.\allowbreak{}orderFormula}{Development.\allowbreak{}AbelianVarietyTheory.\allowbreak{}Proofs.\allowbreak{}InertiaEigenvalues.\allowbreak{}SelmerOrderFormulaDualSelmer}{150--186}

Multiplication by \(u\) kills the dual local invariants at all the unrestricted places. At the fixed exceptional places this follows from their vanishing. At the first primes the residual Frobenius scalar differs from one, so Frobenius minus one is invertible over the local ring, also after an \(\ell\)-power character twist. At an old second prime take inertia of ray coordinate \(1\). All other coordinates vanish on it, and the untwisted dual is unramified. The scalar on the dual is \(1+a_i u\), with \(a_i\) a unit. This proves the assertion, also at the conjugate prime by inversion. In particular the local products in the two order formulas agree.
\leanref{GlobalTheory.\allowbreak{}NewPrime.\allowbreak{}LocInv.\allowbreak{}local\_\allowbreak{}invariants\_\allowbreak{}killed}{Development.\allowbreak{}AbelianVarietyTheory.\allowbreak{}Proofs.\allowbreak{}InertiaEigenvalues.\allowbreak{}LocalInvariantsKilled}{36--86}

Suppose that \(\lambda\) kills every dual class. For a representative \(y\) write \(\lambda y=\partial m\), where \(\partial m(g)=gm-m\). The image of \(m\) modulo \(\lambda A\) is globally invariant and hence zero, by the scalar at \(z\). Write \(m=\lambda m'\). Then \(y-\partial m'\) takes values in \(A[\lambda]^r\subset A[b]^r\). At an unrestricted place let \(m_v\) be a local primitive for \(y\). The vector \(\lambda(m_v-m')\) is locally invariant, so \(u\lambda(m_v-m')=0\). The corrected cocycle has a primitive in \(A[b]^r\) there. Away from the exceptional set inertia acts trivially. Above \(\ell\) every finite test with coefficients in \(Y/bY\) lifts to one with coefficients in \(Y\), by lifting the coefficients of its integral finite cocycles. Compatibility of the pairings now proves orthogonality for the corrected cocycle. Thus the dual classes with coefficients in \(A[b]\) map onto \(D(A)\). The order formulas give
\[
                 |\ker(\Psi:Y^a\to Y^a)|
                 \le |\ker(\Psi:(Y/bY)^a\to(Y/bY)^a)|,
\]
where \(a\) is the size of \(\Psi\). For a square endomorphism of a finite module its kernel and cokernel have the same order. The map on cokernels here is surjective, so the inequality makes it an isomorphism. In particular \(bY^a\subset\Psi Y^a\). Row \(i\) of the adjugate kills \(\Psi Y^a\). Testing it on \(b\) times the \(j\)-th coordinate copy of \(Y\) and using perfectness gives \(bC=0\), a contradiction. This constructs the required \(y\).

The column defining \(x\) lies in \(\ker\Psi\) by the adjugate identity. On \(L_1\) both specialized actions are trivial. This subgroup is normal and stable under \(c\), and contains \([g,z]\) since the extra kernels are abelian ray kernels. The twist is trivial at \(z\). Full residual span gives full matrix span over \(A\) by Nakayama. The polarization of Lemma~\ref{B8.07e} extends to the specialized anticyclotomic character by tensoring its integral intertwiner. If \(B_h(x,y)=0\) throughout \(L_1\), Lemma~\ref{B8.07i} gives \(\ell^\nu J\, y(L_1)=0\), where \(J\subset A\) is the ideal generated by the coordinates of \(x(L_1)\). Modulo \(J\), \(x\) vanishes on \(L_1\). The commutator computation of Lemma~\ref{B8.07l} works over \(A/J\), since \(d_s-1\) is a unit. Thus \(x\bmod J\) is a coboundary. Naturality and injectivity of the presentation on \(A/J\) put the whole adjugate column in \(J\). In particular \(C\in J\), so \(u^l\in J\). It follows that \(\lambda y\) vanishes on \(L_1\); the same commutator computation for the dual makes it a coboundary. This contradicts the choice of \(y\).
\end{proof}

\subsection*{Proof of \ref{B8.07o}}\label{prf:B8.07o}\phantomsection
\begin{proof}
Write \(l=(d+1)w\). If the signature already passes the determinant test at \(2l+2\), take the next second prime by the coordinate prescription of Lemma~\ref{B8.07f} with target element \(1\), keeping all existing primes. The unused first inertia groups lie in \(L\) and allow the next coordinate vector to be prescribed. Lemma~\ref{B8.07m} proves the result in this case.

Otherwise use \(A,\Psi,x,y,h,L_1\) from Lemma~\ref{B8.07n}. In \(L_1\) take the common kernel \(U\) of \(x,y\) and of their conjugate restrictions. This is open and normal by the restriction and conjugation identities of Lemma~\ref{B8.07l}. Every unused first inertia group lies in \(U\). Indeed, the active ray coordinates vanish there and the two actions are unramified there. For a cocycle, the tame relation forces its inertia value into the kernel of Frobenius minus the residue norm. The reduction of this difference is a nonzero scalar since the norm is \(1\bmod\ell\) and the Frobenius scalar differs from one. This applies at both members of the pair and to both cocycles.

The active old coordinates of \(h\) vanish. Multiply \(h\) by elements of unused first inertia groups to give it first-coordinate vector \(e_n\). Each such inertia group surjects onto its own coordinate and vanishes on the others. The correction lies in \(U\). Apply Lemma~\ref{B8.42p} modulo \(U\) intersected with all first-coordinate kernels. Exclude the previous pairs and their conjugates, the exceptional sets of Lemma~\ref{B8.07f} and the given exclusions. The resulting second prime \(P\) has decomposition group in \(L_1\), since inertia lies there and Frobenius does also. It is split of degree and ramification index one, and \(q=\mathrm NP\equiv1\bmod 2\ell^M\). Choose a tame inertia generator \(\gamma\) with new coordinate \(1\). Adjust the Frobenius lift \(\sigma\) by inertia to give it new coordinate zero. Its old coordinates and its class up to conjugacy modulo \(U\) are unchanged. On conjugating, both terms in \(B_h\) are multiplied by the same cyclotomic unit. Consequently \(B_\sigma(x,y)\ne0\).
\leanref{GlobalTheory.\allowbreak{}NewPrimePick.\allowbreak{}exists\_\allowbreak{}newPlace}{Development.\allowbreak{}AbelianVarietyTheory.\allowbreak{}Proofs.\allowbreak{}InertiaEigenvalues.\allowbreak{}ExistsNewPlaceFrobeniusConjugateModOpenNormalSub}{12--127}

Put \(m_1=2r\) and \(B=A[X,Y]/(X^2,Y^{m_1})\), with augmentation \(\pi:B\to A\). Give the new coordinate value \(X\), the first coordinate numbered \(n\) value \(Y\), and the other old coordinates their previous values. These are values of the polynomial variables, so a group generator maps to one plus its value. The binomial calculation of Lemma~\ref{B8.07m} defines the specialization. Its inverse character reduces to the old one. Its values at \(\gamma,\sigma\) are \((1+X)^{-1},(1+Y)^{-1}\), and at their conjugates the inverses of these. Choose \emph{any} new presentation and denote its matrix over \(B\) by \(\Phi'\).

For a square presentation \(P_0\) over a finite coefficient ring, a universal cocycle is obtained by taking a cocycle representing the identity of its representing module \(M_{P_0}\). Lifting its values to the free generators gives a matrix of values \(Z(g)\); multiplying it by any column in the kernel over a coefficient module gives the corresponding cocycle up to coboundary. Use \(Z,Z'\) for the old and new presentations. Set \(M'=B^{a'}/\operatorname{row}(\Phi')\), with \(a'\) the size of \(\Phi'\). Denote the \(m_1\) rows of \(Z'\) at \(\gamma,\bar\gamma\) by \(t_\alpha\), and those at \(\sigma,\bar\sigma\) by \(v_\alpha\), in matching order. In \(M'\) these rows satisfy
\[
       Xv_\alpha=Y\eta_\alpha t_\alpha,\qquad
       \eta_\alpha=
       \begin{cases}(1+Y)^{-1}&\text{at }P,\\1&\text{at }cP.\end{cases}
\]
Indeed the cocycle formula on \(\sigma\gamma\sigma^{-1}=\gamma^q\) in the pro-\(\ell\) tame quotient gives the relation. In computing a value on \(\gamma^q\), the sum of the \(q\) scalar powers is \(q\) plus a multiple of \(\binom q2 X\), since \(X^2=0\). In \(B\), \(q=1\) and \(\binom q2=0\) by \(q\equiv1\bmod 2\ell^M\). Thus at \(P\) the equation is \((1-(1+X)^{-1})v_\alpha=(1-(1+Y)^{-1})t_\alpha\); at \(cP\) both scalar actions are inverted.
\leanref{GlobalTheory.\allowbreak{}NewPrimeRel.\allowbreak{}rel}{Development.\allowbreak{}AbelianVarietyTheory.\allowbreak{}Proofs.\allowbreak{}InertiaEigenvalues.\allowbreak{}UniversalCocycleRowsRelationNewPlace}{37--123}

In \(M'\otimes_B A\) the quotient by the rows \(\pi t_\alpha\) represents the old condition. Indeed, the old specialized action is trivial on the new decomposition groups, and vanishing at the two tame generators is exactly unramifiedness. Choose lifts \(E_b\) of the old free generators to \(M'\otimes A\). Modulo the inertia rows, the old relations on the \(E_b\) are \(\Psi\), and \(\pi v_\alpha\) has old coordinates \(V_{\alpha b}\), the corresponding row of \(Z(\sigma)\) or \(Z(\bar\sigma)\). Evaluation here is independent of the choice of cocycle, because the action at these Frobenius elements is trivial over \(A\).

There is one more relation in \(M'\otimes A\). Put \(k_\alpha=y(\sigma)_i\) at \(P\), and \(k_\alpha=-y(\bar\sigma)_i\) at \(cP\), where \(i\) is the corresponding representation coordinate. Then
\[
                             \sum_\alpha k_\alpha\,\pi t_\alpha=0 .
\]
To prove this test the row on every relaxed class with coefficients in the additive dual of \(A\). Its cup product with \(y\) has vanishing local invariant away from the new pair by the old dual conditions. At each new place \(y\) is unramified. As in Lemma~\ref{B8.07h}, the normalized invariant pairs the inertia value with the dual Frobenius value up to a unit, and using the conjugate generators reverses its sign on the Tate line. The local sum in \ref{B6.09} is zero. Testing also all scalar multiples of each class, the perfect evaluation pairing gives the asserted relation. Indeed, the additive dual of \(A\) separates points of every finite \(A\)-module, by extending additive characters from a cyclic subgroup.
\leanref{GlobalTheory.\allowbreak{}NewPrime.\allowbreak{}comparison\_\allowbreak{}cocycles}{Development.\allowbreak{}AbelianVarietyTheory.\allowbreak{}Proofs.\allowbreak{}InertiaEigenvalues.\allowbreak{}ComparisonCocycles}{24--118}
\leanref{GlobalTheory.\allowbreak{}NewRecip.\allowbreak{}reciprocity}{Development.\allowbreak{}AbelianVarietyTheory.\allowbreak{}Proofs.\allowbreak{}InertiaEigenvalues.\allowbreak{}ReciprocityRelaxedSelmerDualSelmerNewPlace}{190--316}

Lift the \(E_b\) to \(M'\) and adjoin \(\eta_\alpha t_\alpha\) as generators. They generate after augmentation, hence generate \(M'\) since \((X,Y)\) is nilpotent. A relation modulo augmentation lifts with its specified reduced coefficients: lift the coefficients arbitrarily, express the error in \((X,Y)M'\) in these generators with coefficients in \((X,Y)\), and subtract. In the same way express \(v_\alpha\) in the generators by a row \(W_\alpha\) whose first \(a\) coefficients reduce to \(V_{\alpha b}\), where \(a\) is the size of \(\Psi\).

Form \(a\) top relation rows \(D\). In reduction their first block is \(\Psi\) with row \(j\) replaced by zero, where \(j\) is the column index of the old adjugate in Lemma~\ref{B8.07n}. For the other rows choose the inertia coefficients that express the old relation in \(M'\otimes A\). In row \(j\) take the inertia coefficients \(k_\alpha\), using the reciprocity relation. The bottom \(m_1\) relation rows are \(Y(0,1_{m_1})-XW\). Call this square relation matrix \(D'\). Its determinant belongs to \((\det\Phi')\): the rows are relations on a generating list for \(M'\), and its zeroth Fitting ideal in the original presentation is \((\det\Phi')\).

Expand \(\det D'\) by multilinearity in the bottom rows. The term with no \(W\)-row contains \(Y^{m_1}\), and any term with two \(W\)-rows contains \(X^2\). Thus \(\det D'=XY^{m_1-1}f\). To compute \(\pi f\), put the identity rows \((0,1_{m_1})\) in the bottom block and in turn replace bottom row \(\alpha\) by \(-\pi W_\alpha\). Swap this row with top row \(j\). The sign of the swap cancels the minus sign on the new top row. The lower-left block is zero, and the lower-right block is the identity with row \(\alpha\) replaced by \(k\). Its determinant is \(k_\alpha\). Expansion in row \(j\) of the upper-left block therefore gives
\[
 \pi f=\sum_\alpha k_\alpha\sum_b V_{\alpha b}(\operatorname{adj}\Psi)_{bj}
      =B_\sigma(x,y)\ne0 .
\]
In the free \(A\)-basis \(X^iY^j\), \(0\le i<2,\ 0\le j<m_1\), multiplication by \(XY^{m_1-1}\) retains exactly the augmentation coefficient. Hence \(\det D'\ne0\), and \(\det\Phi'\ne0\).

Every specialized variable is a multiple of \(u\), or is \(X\) or \(Y\). A surviving monomial has degree at most \((2l+1)+1+(m_1-1)=2l+m_1+1\). If a polynomial representative of the new determinant had all coefficients in \(I\) below \(2l+2+m_1\), its value would vanish in \(B\). This proves the requested test for every representative of every new presentation, on the single new pair chosen above.
\end{proof}

\subsection*{Proof of \ref{B8.07p}}\label{prf:B8.07p}\phantomsection
\begin{proof}
Use \(R_0\) of Lemma~\ref{B8.07l}. If \(\nu_a=0\), take \(R=R_0\), \(D=4rR_0\) and \(M_0=D+2\). For \(M\ge M_0\) we have \(4rR_0+2<\ell^M\). The primes of Lemma~\ref{B8.07l} meet the hypotheses of Lemma~\ref{B8.07h} with left kernel zero. That lemma gives a unit coefficient through degree \(D\).

Now let \(\nu=\nu_a>0\), \(d=\operatorname{ord}_{O_0}(\ell^\nu)\), and let \(t_0\) be the number of signatures. Set
\[
 R=R_0+t_0,\qquad w_0=4rR_0,\qquad
 w_{j+1}=2(d+1)w_j+2+2r,\qquad D=w_{t_0},\qquad M_0=D+\nu+3 .
\]
For \(M\ge M_0\) we have \(w_0+2<\ell^M\), \(D\le\ell^{M-\nu-1}\) and \(M\ge\nu+2\). In Lemma~\ref{B8.07f} pad the residual list with \(t_0\) identities. Keep all \(R\) first primes and initially keep the first \(R_0\) second primes. For every signature the cofactor test holds at \(w_0+1\). Indeed, set the unused first variables to zero. At their conductor pairs the remaining character is unramified. The possible inertia value of a cocycle is killed by Frobenius minus the norm; this operator is invertible modulo the maximal ideal by Lemma~\ref{B8.07l}, hence is invertible over the remaining group ring. Thus unrestricted classes there are already unramified. The specialized matrix presents the problem on the first \(R_0\) matched pairs in each half. Lemma~\ref{B8.07h} applies with \(b=1\). Specializing an adjugate commutes with taking it, and a monomial surviving this substitution uses none of the unused first variables.

Enumerate the signatures. After \(j\) additions, require every signature to pass the cofactor test at \(w_j+1\) and the first \(j\) signatures to pass the determinant test at \(w_j+1\). For signature \(j+1\), Lemma~\ref{B8.07o} gives the next prime and a determinant coefficient in degree \(<w_{j+1}\). For the same extension, Lemma~\ref{B8.07m} carries all the old tests to \(w_j+1+2r\le w_{j+1}\). Its numerical hypothesis follows from \(w_{j+1}\le D\). Increasing the cutoff to \(w_{j+1}+1\) gives the induction assertion. At the last stage the two halves both have size \(R\), and every determinant has a coefficient outside \(I\) of degree at most \(D\). In a discrete valuation ring, \(I=(\varpi^{d+1})\), so this is the asserted coefficient bound.

For the chosen primes use the matrix of Lemma~\ref{B8.07g}, whose hypotheses hold by Lemma~\ref{B8.07k}. Substituting the character values and taking the divisible coefficient module at the designated level gives the kernel required in Lemma~\ref{B8.07c}, with its local conditions at this \(Q\). The polynomial-to-group-ring map is onto since its images include all \([b_j]-1\). Thus polynomial representatives exist at each \(M\).
\end{proof}

\subsection*{Proof of \ref{B8.07r}}\label{prf:B8.07r}\phantomsection
\begin{proof}
Write \(H\) for the definite unitary group, \(\Gamma=H(\Q)\) in restriction-of-scalars notation, \(J\subset H(\mathbb A_f)\) for the level of the twist, and \(\lambda\) for its level character. Let \(\rho\) denote the action of \(\Gamma\) on coefficient polynomials. A class form satisfies
\[
         F(\gamma xk)=\lambda(k)\rho(\gamma)F(x)
         \qquad(\gamma\in\Gamma,\ k\in J).
\]
All the values of \(\lambda\) are in \(O_t^\times\). Indeed the fixed twist and the chosen twisting character take their values in the designated coefficient rings. Their local unit characters are obtained by evaluating those same characters on idèles supported at the specified places. The finite type character is defined in the coefficient field and has finite order. At a type prime above \(g_i\) the residue field has order \(g_i\), so a residue-unit character has order dividing \(g_i-1\). The member field contains the roots of unity of order dividing \(g_i(g_i-1)\). Each of these character values and its inverse is integral, by the equation \(X^d=1\) for its order \(d\). Taking the product gives the assertion about \(\lambda\).
\leanref{GlobalTheory.\allowbreak{}TwInputs.\allowbreak{}levLam\_\allowbreak{}mem\_\allowbreak{}eOk}{Development.\allowbreak{}AbelianVarietyTheory.\allowbreak{}Proofs.\allowbreak{}InertiaEigenvalues.\allowbreak{}AdmT4OfRootsInField}{91--109}

Let \(h\) be the finite point of a generator, including its allowed translation at the auxiliary split places, and let \(P\) be its polynomial. Then \(P\) belongs to the \(O_t\)-coefficient lattice at \(h\). To see the effect of the translation, let \(B\) be the product of the norms of its supporting primes. The original lattice and its translate coincide away from \(B\). They are full lattices, so a power of \(B\) takes either one into the other. The same assertion holds for their coefficient polynomials. Indeed, in lattice coordinates their change of basis and its inverse have denominators supported on \(B\). Substituting in each minor of a generating product and expanding by minors gives coefficients in \(\mathcal O_E[1/B]\). The determinant character, including a negative determinant power, has the same property by use of the inverse matrix. Clearing these finitely many denominators proves the assertion for each generator and hence for its span. The auxiliary places avoid the places over \(\ell\), so \(B\) is a unit in \(O_t\).
\leanref{GlobalTheory.\allowbreak{}LatTranslate.\allowbreak{}lat\_\allowbreak{}translate}{Development.\allowbreak{}AbelianVarietyTheory.\allowbreak{}Proofs.\allowbreak{}InertiaEigenvalues.\allowbreak{}CoefLatLeTranslateSplitPlacesPrimeToL}{49--70}

Define a class form by
\[
 F_{h,P}(x)=
   \sum_{\substack{\gamma\in\Gamma\\ k_\gamma=h^{-1}\gamma^{-1}x\in J}}
                    \lambda(k_\gamma)\rho(\gamma)P .
\]
When nonempty, the indexing set is a torsor under
\(\Gamma_x=\Gamma\cap xJx^{-1}\), acting on the left. This stabilizer is finite since the unitary group is definite at infinity and \(J\) is compact. Reindexing the sum at \(\gamma_0xk_0\) by \(\gamma=\gamma_0\gamma'\) gives the class-form law. In each summand \(x=\gamma hk_\gamma\). The rational change of variables carries the coefficient lattice at \(h\) to the one at \(\gamma h\), and \(k_\gamma\) preserves the level lattice. Multiplication by \(\lambda(k_\gamma)\) preserves integrality. Thus \(F_{h,P}\in M_t\).
\leanref{GlobalTheory.\allowbreak{}TwTraceCx.\allowbreak{}traceForm\_\allowbreak{}mem\_\allowbreak{}mtLat}{Development.\allowbreak{}AbelianVarietyTheory.\allowbreak{}Proofs.\allowbreak{}InertiaEigenvalues.\allowbreak{}TwistedThetaGeneratorEqClassSumIntegralForm}{115--136}

Write \(\Theta_t(x,P)\) for the kernel with the twisted splitting and the finite vector of the twist. Its transformation law is
\[
     \Theta_t(\gamma hk,\rho(\gamma)P)
                  =\lambda(k)^{-1}\Theta_t(h,P).
\]
Here the finite vector is in the inverse level-character space. The finite part of \(k\) acts on it by \(\lambda(k)^{-1}\). The infinite part of \(\gamma\) acts on the coefficient polynomial by substitution, and left multiplication by the rational element \(\gamma\) reindexes the theta sum. These give the displayed identity with the same \(J\) and \(\lambda\) as in the class form.
\leanref{GlobalTheory.\allowbreak{}TwTraceCx.\allowbreak{}dualLaw\_\allowbreak{}kern}{Development.\allowbreak{}AbelianVarietyTheory.\allowbreak{}Proofs.\allowbreak{}InertiaEigenvalues.\allowbreak{}TwistedThetaGeneratorEqClassSumIntegralForm}{147--196}

Choose representatives \(x\) for \(\Gamma\backslash H(\mathbb A_f)/J\). In
\[
       A_t(F_{h,P})=\sum_{[x]}|\Gamma_x|^{-1}
                                      \Theta_t(x,F_{h,P}(x))
\]
the only contributing class is that of \(h\). For its representative there are exactly \(|\Gamma_x|\) terms, each equal to \(\Theta_t(h,P)\). Thus \(A_t(F_{h,P})=\Theta_t(h,P)\), including when \(\ell\) divides \(|\Gamma_x|\).

Finally the member specifies its generator using the reference splitting and the determinant twist on the target. For either unitary group \(U\), put \(d_{t,U}(g)=t(z)\), where \(z/\bar z=\det g\). The character is trivial on the fixed-field idèles, so this value is independent of \(z\). Changing both splitting characters by \(t\) multiplies the action at \((h,g)\) by \(d_{t,U_{\rm target}}(g)d_{t,H}(h)^{-1}\), the quotient of their determinant characters in the Weil action. The member's generator is therefore \(d_{t,H}(h)\Theta_t(h,P)\). Its preimage is \(d_{t,H}(h)F_{h,P}\), still in \(M_t\) since the twisting value lies in \(O_t^\times\). The identification of the member's class-form lattice with \(M_t\) respects the class-sum map. Injectivity of its coefficient realization gives the asserted containment in the member coordinates.
\leanref{GlobalTheory.\allowbreak{}TwTraceCx.\allowbreak{}gen\_\allowbreak{}eq\_\allowbreak{}smul}{Development.\allowbreak{}AbelianVarietyTheory.\allowbreak{}Proofs.\allowbreak{}InertiaEigenvalues.\allowbreak{}TwistedThetaGeneratorEqClassSumIntegralForm}{78--105}
\leanref{GlobalTheory.\allowbreak{}ClassTrace.\allowbreak{}classSum\_\allowbreak{}traceForm}{Development.\allowbreak{}AbelianVarietyTheory.\allowbreak{}Proofs.\allowbreak{}InertiaEigenvalues.\allowbreak{}FiniteStabOfReps}{155--183}
\end{proof}

\subsection*{Proof of \ref{B8.07s}}\label{prf:B8.07s}\phantomsection
\begin{proof}
If every \(x_\alpha=0\), take \(N=0\). Otherwise take one nonzero \(x_\alpha\) and write \(\lambda r(x_\alpha)=r(y)\) with \(y\in B\). In the given basis let \(s_j,t_j\in K\) be the coordinates of \(x_\alpha,y\). Complex linear independence gives \(t_j=\lambda s_j\) for every \(j\). Choosing \(s_j\ne0\) gives \(\lambda=t_j/s_j\in K^\times\). For every index \(\beta\), injectivity now gives \(\lambda x_\beta\in B\). Choose \(N\ge0\) with \(N\operatorname{ord}_O(\ell)\ge\operatorname{ord}_O(\lambda)\). Multiplication by \(\ell^N/\lambda\in O\) gives the result. The coefficient at an image ray class is the sum over its fibre, so it belongs to the same lattice after multiplication by \(\ell^N\).
\end{proof}

\subsection*{Proof of \ref{B8.07q}}\label{prf:B8.07q}\phantomsection
\begin{proof}
Keep the real row, series realization and attached Galois realization of the starting column used in Lemma~\ref{B8.07a}. The local correspondence in Theorem~\ref{B8.07} is used on this realization. Lemma~\ref{B8.07k} verifies the index conditions at its eligible indices. A selected coefficient place has eligible index by restriction of its place to the fixed coefficient subfield.

First fix the integral data that precede the choice of a character. At an abelian label let \(\Delta_Q\) be the finite Galois ray quotient, and let \(\pi_Q:\operatorname{Pic}(O_Q)\to\Delta_Q\) be the map induced by geometric reciprocity. At a prime \(P\) outside the base exclusions and the auxiliary conductor, \(\pi_Q([\varpi_P])\) is the image of inverse arithmetic Frobenius. Use the fixed coefficient embedding and write \(T(c)\) for the resulting reading of the toric class term at \(c\). For the fixed nonzero scalar \(u_0\) of the label, the ray coefficient at \(\delta\) is \(u_0\sum_{\pi_Q(c^{-1})=\delta}T(c)\). Indeed evaluation at a ray character \(\rho\) weights the term at \(c\) by \(\rho(\pi_Q(c^{-1}))=(\rho\circ\pi_Q)(c)^{-1}\), exactly the weight of that Picard twist in Lemma~\ref{B8.43i}.
\leanref{GlobalTheory.\allowbreak{}ToricDenom.\allowbreak{}IsRecMap}{Development.\allowbreak{}AbelianVarietyTheory.\allowbreak{}Proofs.\allowbreak{}FoundationProofs.\allowbreak{}PlacesBelowSplitCoprimeConductorOfAdmissible}{71--74}
\leanref{GlobalTheory.\allowbreak{}ToricDistribution.\allowbreak{}thEval\_\allowbreak{}eq}{Development.\allowbreak{}AbelianVarietyTheory.\allowbreak{}Proofs.\allowbreak{}FoundationProofs.\allowbreak{}ExistsIsHeckeLOfHasInf}{213--247}
 The host is abelian CM, unramified over its real subfield, and has the distinguished imaginary quadratic subfield and induced type of Definition~\ref{B1.14k}. The base conductor has the prescribed split prime factors, and the auxiliary pairs avoid it and \(\ell\). The parameter \(u\) is positive and its angular weights are those of the label. The class-term bound \eqref{eq:p1c1-toric-denominator}, proved in Lemma~\ref{B1.15a}, is uniform in these auxiliary pairs. Taking the twelfth root in its valuation inequality bounds the denominators of the class terms relative to the fixed period. Clear the remaining common period and scalar in the coefficient reading by a power of \(\ell\). Finite sums of the resulting integral elements are integral, and a discrete valuation ring is integrally closed in its fraction field. This gives an integer for each label. Take their maximum, using zero for an empty set, to obtain \(N_0(a)\). Then all \(\ell^{N_0(a)}\)-scaled ray coefficients are integral, at every admissible auxiliary conductor.
\leanref{GlobalTheory.\allowbreak{}ToricDenom.\allowbreak{}exists\_\allowbreak{}bound\_\allowbreak{}at}{Development.\allowbreak{}AbelianVarietyTheory.\allowbreak{}Proofs.\allowbreak{}FoundationProofs.\allowbreak{}BoundedDenominatorsOfBeta}{133--155}
\leanref{GlobalTheory.\allowbreak{}ToricDenom.\allowbreak{}exists\_\allowbreak{}gen\_\allowbreak{}bound}{Development.\allowbreak{}AbelianVarietyTheory.\allowbreak{}Proofs.\allowbreak{}FoundationProofs.\allowbreak{}BoundedDenominatorsOfBeta}{156--178}

For each layer and sign, project the allowed untwisted vectors onto the prescribed finite-type part. Equivalently one can multiply the projector by the order of the finite type group, using its integral central group-algebra element. Multiplication by this nonzero scalar identifies the two image lattices. Apply Lemma~\ref{B8.07r} with empty auxiliary conductor and the trivial character. Under the map \(x\mapsto1\otimes x\), each allowed untwisted generator belongs to the theta image of the definite class-form lattice at that level. The coefficient extension at this level is finite and injective; tensoring with it is faithfully flat over the base discrete valuation ring. Thus the projected integral span is finitely generated and torsion free. Reduce the generating set modulo the maximal ideal, select a basis of the reduction from it, and lift the corresponding allowed entries. Nakayama gives generation; equality of the number of lifts with the rank gives independence. Take these finite sets as \(S(a,j,\sigma)\), and take the empty set at unused indices or layers.

At a selected place these spans have full rank. In the real realization each vector in the prescribed good summand is read as evaluation of its model vector at the fixed frame, multiplied by the nonzero comparison scalar. The column corresponding to the place has the base tuple of this realization. Its good projected lifts and their right translates span the model values: apply the theta lift in the prescribed range, with the level projections, and take the class sum at base level. In this operation averaging a right translate over the base level is a finite sum on its smooth vector, and averaging fixes the target base-level vector. Thus evaluating the projected lifts gives the span of the allowed class sums. The range here is that of the constructed theta model, with the member's level characters and weights; cuspidal rapid decrease in \ref{B1.15b} permits its integrations. The separating rational readings of \ref{B8.18} apply to the same row and identify its coefficient form. In a rational basis full rank is witnessed by a nonzero minor, which remains nonzero in the prescribed coefficient embedding. This proves full rank at the selected place.
\leanref{GlobalTheory.\allowbreak{}RankOnto.\allowbreak{}sumOnto\_\allowbreak{}of\_\allowbreak{}goodRange}{Development.\allowbreak{}AbelianVarietyTheory.\allowbreak{}Proofs.\allowbreak{}InertiaEigenvalues.\allowbreak{}GoodFormsClassSumsOfProjectedLiftRange}{19--30}
\leanref{GlobalTheory.\allowbreak{}TieRank.\allowbreak{}rank\_\allowbreak{}of\_\allowbreak{}ties}{Development.\allowbreak{}AbelianVarietyTheory.\allowbreak{}Proofs.\allowbreak{}InertiaEigenvalues.\allowbreak{}FullRankThetaSeriesFormsOfRationalScalars}{82--90}
\leanref{GlobalTheory.\allowbreak{}TwistChoice.\allowbreak{}exists\_\allowbreak{}basisOf}{Development.\allowbreak{}AbelianVarietyTheory.\allowbreak{}Proofs.\allowbreak{}InertiaEigenvalues.\allowbreak{}ExistsEquivMapLattice}{188--211}

We prove the uniform assertion of Definition~\ref{B8.07d}. Fix an eligible index satisfying its conditions, suppose the projected spans have full rank, and give any allowed integral bases, \(e>0\), and a finite avoidance set. Let \(O\) denote the coefficient ring for the theta polynomials in Lemma~\ref{B8.07b}. We describe the polynomials on the ray coordinates of the member. Use in Lemma~\ref{B8.07f} the admissible single-pair coordinates for its Artin maps. The coordinate on a retained pair is the same after restriction from a larger conductor: both maps on \(G_E\) are the chosen single-pair character. Surjectivity of the Artin quotient then gives the same identity there. Proposition~\ref{B8.07p} works for this family of coordinates.

For an abelian label \(\lambda\), use as \(f\) a polynomial reading of its scaled ray sum at all pairs, and as \(g\) its reading at the first pairs alone. Explicitly, push its finitely many integral coefficients into the product of the cyclic ray groups and replace a group element with coordinates \(c_i\) by \(\prod_i(1+Z_i)^{c_i}\), using least nonnegative exponents. For a split pair of norm \(q\) on which the testing character is unramified, passage from the smaller to the larger order multiplies the sum by
\[
                  (1-aq^{-u})(1-a^{-1}q^{-u})+(q-1)q^{-2u},
\]
where \(a\) is the product of the label and testing characters evaluated at arithmetic Frobenius at the chosen prime \(P\). In the common coefficient field the label factor is \(\chi_f(\varpi_P)^{-1}\), and the testing factor is \(\rho(\pi_Q([\varpi_P])^{-1})\) at the smaller conductor. Thus both factors use the same Frobenius. Replacing the testing character by its inverse inverts the testing factor. Indeed, locally the smaller order is the product of the two valuation rings and the new order requires equal residues. Among units this leaves the diagonal residue units. If both valuations are positive the congruence is automatic, with unit index \(q-1\). Summing the two valuation series and dividing by the unrestricted Euler series gives the displayed factor. This computation in the convergence half-plane extends to \(u\) by the toric class-sum formula of Definition~\ref{B1.14k} and \ref{B8.11h}. The expression is invariant under inversion of \(a\).

Require first Frobenius to agree up to conjugacy with \(z\) on the residual label characters. Their values at \(z\) differ from one: on designated inertia their ratio to the prescribed order-fifteen value or its inverse has residue-characteristic-power order, whereas the chosen component of \(z\) has order prime to that residue characteristic in every finite quotient and retains its nontrivial order-\(p\) value. Since \(q\equiv1\bmod\ell\), the first-pair factors at the trivial testing character are units. Consequently
\[
              \operatorname{ord}_O g(0)
                =\operatorname{ord}_O\!\left(\ell^{N_0(a)}
                             \sum_\delta\theta_{\varnothing,\lambda,\delta}\right).
\]
Here \(\theta_{\varnothing,\lambda,\delta}\) denotes an empty-conductor ray coefficient. The sum is nonzero by Lemma~\ref{B8.07k}. Set the second variables to zero. The Frobenius reading at second prime \(i\) of a character on the first pairs is \(1+X_i\). Modulo the first cyclic relations use \((1+X_i)^{\ell^M-1}\) for its inverse. The product \(A_\lambda\) of the displayed local factors gives \(f=A_\lambda g\) in that quotient. One may check this at all tuples of \(\ell^M\)-th roots in the designated extension: a polynomial reduced to degree \(<\ell^M\) separately in each variable and vanishing at every tuple is zero, by successive univariate interpolation in the fraction field. In the residue quotient each factor is
\[
              (1-\alpha(1+X_i))(1-\alpha^{-1}(1+X_i)^{-1}),
\]
with \(\alpha\) the residual label value. It has a unit constant term if \(\alpha\ne1\); for \(\alpha=1\) it is \(-X_i^2/(1+X_i)\). This gives the second hypothesis of Lemma~\ref{B8.07b} with exponent bound two.
\leanref{GlobalTheory.\allowbreak{}ToricDistribution.\allowbreak{}gen\_\allowbreak{}relation}{Development.\allowbreak{}AbelianVarietyTheory.\allowbreak{}Proofs.\allowbreak{}FoundationProofs.\allowbreak{}ExistsIsHeckeLOfHasInf}{314--394}
\leanref{GlobalTheory.\allowbreak{}ToricFactor.\allowbreak{}propE}{Development.\allowbreak{}AbelianVarietyTheory.\allowbreak{}Proofs.\allowbreak{}InertiaEigenvalues.\allowbreak{}PropEOfAdmissible}{420--562}

For a projected basis of size \(b\) use the lattice obtained by intersecting the full integral type lattice with the good fraction-field space. It is a full lattice there, since denominators of each vector in the subspace can be cleared in the full lattice. The integral multiple of the projector preserves the full lattice, takes it into this intersection and commutes with the Hecke operators. Indeed, integral finite-group translates preserve the lattice, and the Hecke operators in question are at split places outside the type levels. They preserve the full lattice by the integral Hecke comparison of \ref{B8.43h}. Fix a basis of the intersection lattice.

There is one power of \(\ell\) clearing the coordinates of the projected universal series for the given \(b\) entries, for all ray levels disjoint from their levels. Here one uses the integral series assertion on the pinned frame in \ref{B8.43h}. For each small-field generator apply Proposition~\ref{B8.43o} to its cyclic additive lattice at the chosen cover place. To obtain a prescribed universal coefficient, multiply the readings by the indicators of the fibres of the determinant ray coordinate and sum. The fibres partition the components, so precisely one indicator contributes at each point. These indicators are integral component functions on the ray cover. Right translation by a base level element permutes the local labels in this sum: its reduction preserves the test lattice and its upper left block at each type prime is the identity. The inverse change of the determinant label carries the fibre indicator to the indicator of the translated test. Reindexing gives invariance under the base compact subgroup. In the fixed coefficient valuation ring put \(e=\operatorname{ord}_O(\ell)\). Lemma~\ref{B8.43p} now gives a combination of base integral sections after the further factor \(\ell^{e^2}\) at every selected cover place, including ramified ones. This exponent was fixed before the auxiliary conductor. An allowed harmonic coefficient function is a finite integral linear combination of these generators in the coefficient lattice. The series is linear in that function, so clearing its finitely many coefficient denominators works for every ray class and every cover. Write the member's reading as \(r=\nu^{-1}\operatorname{ev}\), where \(\operatorname{ev}\) is extremal evaluation and \(\nu\ne0\) is the fixed comparison scalar. If \(\ell^c u\,r(x)\) is the evaluation of an element of the integral section lattice, then \((\ell^c u/\nu)r(x)\) belongs to the reading of the member's lattice. The rational cusp frame stays linearly independent in its complex realization. Lemma~\ref{B8.07s} now gives a single clearing power for this datum, including after adding coefficients in each fibre of the ray-coordinate map. Taking the maximum over the finite basis gives the asserted common power. Exclude its test levels and the places over the host primes \(q,d\).
\leanref{GlobalTheory.\allowbreak{}XiIntComb.\allowbreak{}xi\_\allowbreak{}int\_\allowbreak{}comb}{Development.\allowbreak{}AbelianVarietyTheory.\allowbreak{}Proofs.\allowbreak{}InertiaEigenvalues.\allowbreak{}XiIntComb}{78--109}
\leanref{GlobalTheory.\allowbreak{}XiBound.\allowbreak{}main}{Development.\allowbreak{}AbelianVarietyTheory.\allowbreak{}Proofs.\allowbreak{}InertiaEigenvalues.\allowbreak{}UniversalSeriesCoeffBoundedInSupOfTwistSeriesSid}{40--95}
\leanref{GlobalTheory.\allowbreak{}TwInputs.\allowbreak{}univSup\_\allowbreak{}of\_\allowbreak{}xi}{Development.\allowbreak{}AbelianVarietyTheory.\allowbreak{}Proofs.\allowbreak{}InertiaEigenvalues.\allowbreak{}UnivSupOfXi}{13--45}

Form the \(b\)-column coordinate matrix of these scaled projected vectors over the ray group ring. Write \(m_\varnothing,m_1,m_Q\) for the matrices at empty conductor, first pairs, and all pairs, using polynomial representatives. The determinant of \(m_\varnothing\) is nonzero by the basis property. At a split place \(P\) of norm \(q\) write \(H_j(P)\) for the matrix of the order-\(j\) normalized Hecke operator on the intersection lattice. Removing a pair on which the testing character is unramified expresses the larger-conductor columns as the smaller-conductor columns acted on by
\[
                 C_P(y)=1+\sum_{j=1}^{r}
                    (-1)^j q^{j(j-1)/2}H_j(P)y^j ,
\]
where \(y\) is the signed Frobenius value at the chosen member of the pair. Write \(O_P\) for the local integers, \(\varpi\) for a uniformizer, and \(N\) for the number of columns of the local test matrix. For \(Z\in M_{r,N}(O_P)\), summing the determinant test over its label counts the \((N-r)\)-planes complementary to the rows of \(\bar Z\), with weight \(q^{-r(N-r)}\) each. At rank \(r\), fix one complement. All complements are the graphs of the linear maps from it to the row space, giving \(q^{r(N-r)}\) choices. At smaller rank the two subspaces cannot span \(\F_q^N\). Thus summing out the label gives the indicator of row rank \(r\).\par For a right coset representative \(g\) of the order-\(j\) spherical double coset, \(L=gO_P^r\) runs once through the intermediate lattices \(\varpi O_P^r\subset L\subset O_P^r\) of index \(q^j\). Indeed the inverse image of any codimension-\(j\) residue subspace has these elementary divisors, and two representatives give the same lattice exactly when they differ on the right by \(\GL_r(O_P)\). The translated integral test requires \(g^{-1}Z\) integral. This says that every column of \(\bar Z\) lies in \(L/\varpi O_P^r\). Taking annihilators identifies the admissible lattices with the \(j\)-planes in \(\ker(\bar Z^{\mathsf t})\). If this kernel has dimension \(h\), the alternating weighted count is \(\sum_j{h\brack j}_q(-1)^jq^{\binom j2}=\prod_{a=0}^{h-1}(1-q^a)\). The polynomial identity before evaluation at one follows by induction from \({h\brack j}_q={h-1\brack j}_q+q^{h-j}{h-1\brack j-1}_q\), obtained by intersecting with a fixed hyperplane. The product equals the full-rank indicator, including at \(h=0\). For a nonintegral matrix every lattice test vanishes since \(L\subset O_P^r\).\par Let \(x_P\) be the uniformizer idèle on the chosen side, and let \(\mu\) be the splitting character of the theta frame. On this double coset the normalized Hecke sum of right translates has prefactor \(\mu(x_P)^{-j}q^{-j\lfloor N/2\rfloor}\). The determinant has valuation \(j\), and unit factors act trivially on the unramified character. This is the inverse of the splitting-character and modulus factor in the Weil action on the integral test. The theta series is linear in that test, so the preceding indicator identity becomes the Hecke identity with coefficient \((-1)^jq^{\binom j2}\). To track its ray value, write \(c\) for the remaining additive ray coordinates and \(\delta\) for their reading on \(x_P\). The smaller-level coefficient on the translated term is evaluated at \(c+j\delta\). The testing character factors through the smaller level. Reindexing its sum therefore multiplies the term by \(\chi(-j\delta)\). With the Artin convention of the ray coordinates, \(\delta\) is minus the coordinate of arithmetic Frobenius, so \(\chi(-j\delta)=y^j\) on the chosen side. The other member of the pair gives the inverse reading. Finally the projection onto the good type commutes with all these translates because its type places are disjoint from \(P\). Applying the identity to each basis vector gives left multiplication of the smaller-level coordinate matrix by \(C_P(y)\).
\leanref{GlobalTheory.\allowbreak{}DistRel.\allowbreak{}xi\_\allowbreak{}rel}{Development.\allowbreak{}AbelianVarietyTheory.\allowbreak{}Proofs.\allowbreak{}InertiaEigenvalues.\allowbreak{}IsDetRootOfDetAt}{190--313}
\leanref{GlobalTheory.\allowbreak{}DistRel.\allowbreak{}thURel\_\allowbreak{}of\_\allowbreak{}tie}{Development.\allowbreak{}AbelianVarietyTheory.\allowbreak{}Proofs.\allowbreak{}InertiaEigenvalues.\allowbreak{}UniversalTwistedSeriesDistributionOfCompatibilit}{152--364}
\leanref{GlobalTheory.\allowbreak{}TwistDet.\allowbreak{}det\_\allowbreak{}factor}{Development.\allowbreak{}AbelianVarietyTheory.\allowbreak{}Proofs.\allowbreak{}InertiaEigenvalues.\allowbreak{}AevalOneSubOne}{142--209}

Impose also the first-prime congruences for all residual Hecke readings and their conjugates. These are open normal congruences: take the kernels of the finitely many residual Galois representations giving the reverse Hecke polynomials, and their conjugates. On the split spherical primes the reverse polynomial is that of \(T_s(t)\) at inverse arithmetic Frobenius; coefficient \(j\) is \((-1)^j q^{j(j-1)/2}\) times the order-\(j\) Hecke eigenvalue. This follows by taking elementary symmetric functions of the attached Satake parameters with the shift of Lemma~\ref{B8.07k}. For a good signature \(s\), Lemma~\ref{B8.07k} gives \(T_s(z)=d_s1_r\), with \(d_s^2-1\) a unit and with cyclotomic character one at \(z\). Hence the extra Tate twist in the Hecke reading has value one there. In the finite quotient just chosen, arithmetic Frobenius at a first prime is conjugate to the image of \(z\), with conjugation of the place allowed. Thus its inverse in the residual reading is a scalar \(c_s1_r\) with \(c_s\in\{\bar d_s,\bar d_s^{-1}\}\). Evaluating the reverse polynomial at one gives \(\det(1-c_s1_r)=(1-c_s)^r\ne0\) for both choices.\par Extend the coefficient discrete valuation ring to split the finitely many characteristic-zero Hecke eigensystems on the good type space. A common invariant flag exists by simultaneous triangularization. Its intersections with the extended type lattice are saturated, since the flag lies in the fraction-field space. Their successive quotients are free of rank one and reduce to a full residual flag. On each quotient the diagonal of \(C_P(1)\) is the reading just computed for the corresponding signature in the row. The reduction of its determinant is the product of these nonzero readings, counted with their multiplicities. A nonunit in the original ring remains a nonunit in the extension. Hence \(\det C_P(1)\) is a unit in the original ring as well.
\leanref{GlobalTheory.\allowbreak{}TripleGal.\allowbreak{}HkFrobProof.\allowbreak{}hkFrob\_\allowbreak{}of\_\allowbreak{}tripleAtt\_\allowbreak{}n}{Development.\allowbreak{}AutomorphicArithmetic.\allowbreak{}Foundations.\allowbreak{}GaloisRepresentations.\allowbreak{}StringRepresentations.\allowbreak{}FrobDualOfTripleAttN}{189--251}
\leanref{GlobalTheory.\allowbreak{}TwistDet.\allowbreak{}exists\_\allowbreak{}subgroup\_\allowbreak{}isUnit}{Development.\allowbreak{}AbelianVarietyTheory.\allowbreak{}Proofs.\allowbreak{}InertiaEigenvalues.\allowbreak{}ExistsSubgroupIsUnit}{75--186}

Use \(f=\det m_Q,\ g=\det m_1\). At the trivial character, successive removal of the first pairs gives \(\operatorname{ord}_O g(0)=\operatorname{ord}_O\det m_\varnothing\). Removing the second pairs with \(Y_i=0\) gives \(f=Ag\) modulo the first cyclic relations, where \(A\) is the product of the determinants of \(C_{P_i'}((1+X_i)^{\pm1})\). This follows by the same interpolation argument as for the abelian labels. Over the residue field \(\det C_{P_i'}(y)\) is a polynomial of degree at most \(rb\) with value \(1\) at \(y=0\). Its order at \(y=1\) is at most \(rb\). Substituting \(y=(1+X_i)^{\pm1}\) preserves this order since \(y-1\) is \(X_i\) times a unit in the local power series ring. This yields the factor required in Lemma~\ref{B8.07b}, modulo \(X_i^{\ell^M}\), with exponent at most \(rb\). For \(b=0\) the determinant is \(1\).
\leanref{GlobalTheory.\allowbreak{}TwistDet.\allowbreak{}propABCD\_\allowbreak{}of\_\allowbreak{}rel}{Development.\allowbreak{}AbelianVarietyTheory.\allowbreak{}Proofs.\allowbreak{}InertiaEigenvalues.\allowbreak{}AevalOneSubOne}{293--466}

Intersect the open normal subgroups for these finitely many tests and unite the finite exclusion sets with the requested one. Proposition~\ref{B8.07p} supplies the Selmer matrices on primes satisfying precisely these extra restrictions. The constants just obtained for the theta polynomials depend on the fixed bases and labels and are independent of \(M\). Apply Lemma~\ref{B8.07c} to all these polynomials, taking
\[
                  \delta=e/(1+|\mathcal S|+2m+4Rt).
\]
It gives one triple \(Q,k',\xi\) and the Selmer and tame inequalities for every required signed coprime power. For an abelian label, undoing \(\ell^{N_0(a)}\) gives lower valuation bound \(-N_0(a)\) and upper bound \(c_0(a,\lambda)+\delta\), since the empty-conductor reference is nonzero and was scaled by the same power.

For a layer and sign, the specialized determinant is nonzero. Its normalized valuation is at most the normalized valuation of \(\det m_\varnothing\) plus \(\delta\). The two clearing powers contribute the same \(b\)-fold scalar to the two determinants. Thus the twisted basis columns span a full lattice with signed index at most the untwisted index plus \(\delta\). Each column is a specialized allowed theta vector because \(Q\) avoids its level. By Lemma~\ref{B8.07r}, the full projected twisted span lies in the projection of the theta image of the definite class-form lattice for this twist. That image is finitely generated. Its submodule is finitely generated over the coefficient discrete valuation ring and is torsion free in the coefficient space. It lies in the prescribed good space by the type of its generators. It contains the full-rank sublattice supplied by the specialized columns, so it is a lattice there and has signed index at most that of the sublattice. Removing the common multiple of the projector cancels the same scalar on the two sides. Over a finite injective extension of the coefficient discrete valuation rings, determinant orders and \(\operatorname{ord}(\ell)\) multiply by the same ramification index. A change to any test lattice shifts both signed indices by the same determinant valuation. This proves the required inequality for every such extension and test lattice. There are \(2m\) layer-sign tests, so their errors sum to at most \(e\).
\leanref{GlobalTheory.\allowbreak{}ExportFromInputs.\allowbreak{}export\_\allowbreak{}at}{Development.\allowbreak{}AbelianVarietyTheory.\allowbreak{}Proofs.\allowbreak{}InertiaEigenvalues.\allowbreak{}IsLatInOfIndexLe}{371--451}

The same triple works for all coprime powers because the specialization estimates hold at every primitive root of the chosen order, including its inverse. The integers \(N_0(a)\) and the selected basis sets were fixed before the exclusions or powers. This proves Definition~\ref{B8.07d}.
\end{proof}

\subsection*{Proof of \ref{B8.07t}}\label{prf:B8.07t}\phantomsection
\begin{proof}
The input to localization is the rational realization of the multiplicity spaces with its chosen tuple and conductor at \(w\). The character at the last finite level is the product of the reference character and the chosen twist, and its conductor set is the one used to select the split operators. Thus the integral Hecke action at that level applies to the twisted space \(W\). More explicitly, for \(s\in S\) and an allowed generator \(H\), the double-coset action on the embedded vector is a linear combination of the embedded generators with coefficients integral at \(w\). The rational map into the space of forms is injective and intertwines \(H\). Pulling the equality back gives the same combination for \(Hs\). It follows that
\(\Lambda=\sum_{s\in S}O(1\otimes s)\) is stable under \(\mathbb T\). These generators span \(W\) over \(K_0\) by the layer lattice condition. Thus \(\Lambda\) is a full lattice and each \(H\) has integral eigenvalues.
\leanref{GlobalTheory.\allowbreak{}HeightCounting.\allowbreak{}Fin3.\allowbreak{}mtInput\_\allowbreak{}of\_\allowbreak{}levelC}{Development.\allowbreak{}AutomorphicArithmetic.\allowbreak{}ThetaCorrespondence.\allowbreak{}ThetaPairLifts.\allowbreak{}ThirdPlaceSeries.\allowbreak{}RowsAtStart}{31--33}
\leanref{GlobalTheory.\allowbreak{}HeckeLocN.\allowbreak{}heckeIntTw\_\allowbreak{}base}{Development.\allowbreak{}AbelianVarietyTheory.\allowbreak{}Proofs.\allowbreak{}IwasawaDecomposition.\allowbreak{}HeckeOperatorIntegralOnLatticeAtTwistedPlaces}{23--33}
\leanref{Formalism.\allowbreak{}L2.\allowbreak{}Fam.\allowbreak{}NPC.\allowbreak{}LatticeFinitenessFree.\allowbreak{}span\_\allowbreak{}latL}{Formalism.\allowbreak{}LocalToGlobal.\allowbreak{}NumberFieldFamilies.\allowbreak{}IsLocalizedOfSpan}{105--111}

Operators at different places commute. At one split spherical place transpose fixes each double coset by the elementary divisor decomposition and reverses convolution. Thus the spherical algebra is commutative. To choose a splitting field at \(w\), take finitely many of the operators spanning their \(K_0\)-linear span in \(\operatorname{End}_{K_0}(W)\). Adjoin the roots of their characteristic polynomials to the completion at \(w\). Commuting split operators admit a common triangularization, obtained by taking a common eigenvector and passing to the quotient. Each remaining operator is a linear combination of the chosen ones, hence also splits there. Taking the compositum with the finite local fields required for the coefficient and character embeddings gives the \(K'\) used above.
\leanref{GlobalTheory.\allowbreak{}HeckeLocN.\allowbreak{}ndSpec\_\allowbreak{}exists\_\allowbreak{}alt}{Development.\allowbreak{}AutomorphicArithmetic.\allowbreak{}ThetaCorrespondence.\allowbreak{}ThetaPairLifts.\allowbreak{}ThirdPlaceSeries.\allowbreak{}NdSpecExists}{23--47}

Decompose \(V\) successively into the primary spaces of the chosen commuting operators. Each decomposition preserves the preceding ones. The resulting nonzero joint spaces have characters \(\lambda:\mathbb T\to O\), since the diagonal values of all generators lie in \(O\). Write \(J_\lambda=\ker\lambda\) and \(\mathfrak m_\lambda=\ker(\mathbb T\xrightarrow{\lambda}O\to k)\), where \(k\) is the residue field of \(O\). Every maximal ideal is some \(\mathfrak m_\lambda\). Indeed, on a common triangular flag the product of the kernels of the successive diagonal characters annihilates \(V\). The action of \(\mathbb T\) on \(V\) is faithful, so this product is zero in \(\mathbb T\). A maximal ideal must contain one \(J_\lambda\). Since \(\lambda\) is the identity on scalars, \(\mathbb T/J_\lambda=O\), which has just one maximal ideal. This proves the assertion. For \(V=0\) the endomorphism algebra is zero and the index set is empty.

Collect the joint spaces with the same \(\mathfrak m_\lambda\). This gives \(V=\bigoplus_{\mathfrak m}V_{\mathfrak m}\). In the construction a local index consists of a layer, a place with nonreference tuple, and a maximal ideal of this integral algebra. The fibre at \((j,w)\) is taken by equality of the first two entries. The map to maximal ideals forgets those entries and has inverse \(\mathfrak m\mapsto(j,w,\mathfrak m)\). This also explains why a reference place has empty fibre.
\leanref{Formalism.\allowbreak{}L2.\allowbreak{}Fam.\allowbreak{}NPC.\allowbreak{}locsC\_\allowbreak{}mem}{Formalism.\allowbreak{}LocalToGlobal.\allowbreak{}NumberFieldFamilies.\allowbreak{}NPartCNd}{153--156}
\leanref{Formalism.\allowbreak{}L2.\allowbreak{}Fam.\allowbreak{}NPC.\allowbreak{}mxC\_\allowbreak{}inj}{Formalism.\allowbreak{}LocalToGlobal.\allowbreak{}NumberFieldFamilies.\allowbreak{}NPartCNdTie}{56--68}
\leanref{Formalism.\allowbreak{}L2.\allowbreak{}Fam.\allowbreak{}NPC.\allowbreak{}mxC\_\allowbreak{}all}{Formalism.\allowbreak{}LocalToGlobal.\allowbreak{}NumberFieldFamilies.\allowbreak{}NPartCNdTie}{69--72}

For each \(\mathfrak m\), the projector \(p_{\mathfrak m}\) preserves \(\Lambda\). To see this, choose finitely many Hecke generators distinguishing all the distinct residual characters. For one chosen operator, group the linear factors of its characteristic polynomial according to their roots modulo the maximal ideal of \(O\), retaining their multiplicities. Factors belonging to different groups have unit resultant. The Chinese remainder theorem gives polynomials over \(O\) acting as the primary projectors for these groups. Multiply the projectors for the residues prescribed by \(\mathfrak m\) over the chosen operators. On each joint space the product is the identity or zero according to its residual character. Hence the product is \(p_{\mathfrak m}\), and belongs to \(\mathbb T\).

The construction takes \(M_i=p_{\mathfrak m}\Lambda\). It is generated by the indicated projected vectors, is torsion free over the discrete valuation ring \(O\), and spans \(V_{\mathfrak m}\). It is therefore finite free. Localization of its inclusion gives the displayed isomorphism. The local comparison sends a row vector to its projection followed by the inverse of this isomorphism. In particular its image in \(V\) differs from the original row vector by an element of \(\bigoplus_{\mathfrak n\ne\mathfrak m}V_{\mathfrak n}\). The comparison intertwines every Hecke generator because each commutes with \(p_{\mathfrak m}\). Taking dimensions, and using \(\dim_{K'}V=\dim_{K_0}W\), gives the rank formula.
\leanref{Formalism.\allowbreak{}L2.\allowbreak{}Fam.\allowbreak{}NPC.\allowbreak{}MtSub}{Formalism.\allowbreak{}LocalToGlobal.\allowbreak{}NumberFieldFamilies.\allowbreak{}NPartCNd}{185--187}
\leanref{Formalism.\allowbreak{}L2.\allowbreak{}Fam.\allowbreak{}NPC.\allowbreak{}prMtC}{Formalism.\allowbreak{}LocalToGlobal.\allowbreak{}NumberFieldFamilies.\allowbreak{}NPartCNd}{196--200}
\leanref{Formalism.\allowbreak{}L2.\allowbreak{}Fam.\allowbreak{}NPC.\allowbreak{}PsiC\_\allowbreak{}pr}{Formalism.\allowbreak{}LocalToGlobal.\allowbreak{}NumberFieldFamilies.\allowbreak{}NPartCNdTie}{88--96}

For the length assertion, use the good spectral points of Definition~\ref{B8.40} on each of these modules. Write \(G_i\) for the sum of their eigenspaces in \(K'\otimes_O M_i\) and \(H_i\) for the sum for the remaining points. The local spectral realization is split semisimple on these spaces. Under the comparison just constructed, \(1\otimes M_i\) has image \(p_{\mathfrak m}\Lambda\), and the eigenspaces have the same Hecke characters on both sides. Thus the local quotient in Definition~\ref{B8.40} is carried isomorphically to
\[
          \frac{G_i\cap(p_{\mathfrak m}\Lambda+H_i)}
               {p_{\mathfrak m}\Lambda\cap G_i},
\]
where we use \(G_i,H_i\) also for their images in \(V_{\mathfrak m}\). In particular, if \(G=\bigoplus_iG_i\) and \(H=\bigoplus_iH_i\), then
\[
 \frac{G\cap(\Lambda+H)}{\Lambda\cap G}
   \ \simeq\
   \bigoplus_i
   \frac{G_i\cap(p_{\mathfrak m}\Lambda+H_i)}
        {p_{\mathfrak m}\Lambda\cap G_i}.
\]
Indeed \(\Lambda=\bigoplus_{\mathfrak m}p_{\mathfrak m}\Lambda\), so both intersections can be taken componentwise. The quotients have finite length because the upper and lower lattices each span the good subspace. Division of the length by \(\operatorname{length}_O(O/\ell O)\) gives exactly the sum of normalized lengths. This normalization also agrees before and after finite extension of the local coefficient field. An elementary divisor of order \(a\) acquires order \(ea\) under an extension of ramification index \(e\), as does the order of \(\ell\) with the same factor.

At the last finite level the construction first forms the local modules and then assigns the good spectral points on those very modules. For \(1\le j\le m\) it sets the row coordinate equal to the sum of their normalized quotient lengths. It puts the new local data and this coordinate into the same member over the column. The subsequent spectral comparison uses its maximal-ideal map and its maps from the localized modules. Thus the comparison maps used for the ranks also carry the lattices and Hecke subspaces used for the lengths.
\leanref{Formalism.\allowbreak{}BasicNotions.\allowbreak{}NData.\allowbreak{}etaLoc}{Formalism.\allowbreak{}SharpRestriction.\allowbreak{}LevelLetterCocycles.\allowbreak{}DataN}{219--221}
\leanref{GlobalTheory.\allowbreak{}HeightCounting.\allowbreak{}NLay0.\allowbreak{}fixZ}{Development.\allowbreak{}AbelianVarietyTheory.\allowbreak{}Definitions.\allowbreak{}AutomorphicMembership.\allowbreak{}NStageTrivialCharacterCoeffEqAlt}{62--65}
\leanref{GlobalTheory.\allowbreak{}HeightCounting.\allowbreak{}setN}{Development.\allowbreak{}AbelianVarietyTheory.\allowbreak{}Definitions.\allowbreak{}AutomorphicMembership.\allowbreak{}NStageTrivialCharacterCoeffEqAlt}{125--127}
\leanref{GlobalTheory.\allowbreak{}HeightCounting.\allowbreak{}Fin3.\allowbreak{}oL}{Development.\allowbreak{}AbelianVarietyTheory.\allowbreak{}Proofs.\allowbreak{}IwasawaDecomposition.\allowbreak{}TwistDataAgreeAcrossCoefficientsLocalized}{50--52}
\leanref{GlobalTheory.\allowbreak{}EtaMem.\allowbreak{}exists\_\allowbreak{}etaJ\_\allowbreak{}alt}{Development.\allowbreak{}AutomorphicArithmetic.\allowbreak{}ThetaCorrespondence.\allowbreak{}ThetaPairLifts.\allowbreak{}ThirdPlaceSeries.\allowbreak{}NonGaloisEtaJMem}{29--92}
\end{proof}

\subsection*{Proof of \ref{B8.07}}\label{prf:B8.07}\phantomsection
\begin{proof}
For the first assertion apply Theorem~\ref{B8.48} to each field and rank in the mass assertion. Its field hypothesis is the prescribed unramified CM condition, and the adjacent-rank formula is required for \(M\ge3\).

First work on the origin extension of (8.6). The base assertion chooses a preliminary datum and frame before the two finite coefficient requirements are specified. Use the coordinate preliminary datum of the host and a frame on it. The first requirement puts its algebraic elements in the fixed coefficient subfield, and the second puts its elements in the full coefficient field. The frame is transported from the preliminary datum in both choices. On subsequent admissible enlargements the preliminary datum and the prescribed field of frame values are identified, so the base condition persists.
\leanref{GlobalTheory.\allowbreak{}Mem.\allowbreak{}Top.\allowbreak{}Bl.\allowbreak{}KPxCond.\allowbreak{}baseCond\_\allowbreak{}KPx}{Development.\allowbreak{}AbelianVarietyTheory.\allowbreak{}Proofs.\allowbreak{}IwasawaDecomposition.\allowbreak{}ExistsPrePin}{30--37}

On an eligible untwisted base containing the required primes, choose a starting column with the given frame. The untwisted character supplies its support and conjugacy conditions. In \(P\), the twisted Siegel--Weil identity gives the comparison for the theta realization, and twisted local multiplicity gives equality of the row multiplicities. Use the separating functionals of Theorem~\ref{B8.18} on this same realization. The row values are rational at its specified embedding. These are the multiplicity and rationality inputs of Lemma~\ref{B8.07a}; apply it to this column.
\leanref{GlobalTheory.\allowbreak{}Mem.\allowbreak{}member\_\allowbreak{}hodgeHyp\_\allowbreak{}of\_\allowbreak{}localInputs\_\allowbreak{}rationalWeights}{Development.\allowbreak{}AutomorphicArithmetic.\allowbreak{}ThetaCorrespondence.\allowbreak{}ThetaPairLifts.\allowbreak{}LocalTypeDictionary.\allowbreak{}MemberHodgeHypOfLocalInputsRationalWeights}{12--15}

Starting from its state with empty local index set, put the local correspondence for \(h_L\) on the later stage. Use its multiplicity spaces, rational row realization and cuspidal realization together as the intermediate level. Extend successively by the three local levels, carrying the preceding data at each extension. At the last level retain the selected places and their characters. For each twisted tuple apply Lemma~\ref{B8.07t} to its own multiplicity space. The tuple is the reference character times the chosen twist, and its conductor is the one defining the allowed double cosets. The integral action on the layer generators thus applies at precisely those places. Install the local modules, their Hecke actions and their projection maps together. The row coordinate at each lift layer is the sum of the normalized lengths on this local family. The resulting member lies over the original column. The theta data for its rows and columns have the same indexing, so adjoining the ambient coordinates and the division data gives a constructed point. Apply Proposition~\ref{B8.07q} to this point, using the real row and the series of the column just constructed and its Galois realization for \(h_L\). This gives the twisting data of Definition~\ref{B8.07d} for the point's own theta vectors.
\leanref{GlobalTheory.\allowbreak{}Mem.\allowbreak{}Top.\allowbreak{}Bl.\allowbreak{}has3\_\allowbreak{}neo\_\allowbreak{}of\_\allowbreak{}localLanglands\_\allowbreak{}alt}{Development.\allowbreak{}AutomorphicArithmetic.\allowbreak{}ThetaCorrespondence.\allowbreak{}ThetaPairLifts.\allowbreak{}ThirdPlaceSeries.\allowbreak{}Has3NeoOfSteps}{19--35}
\leanref{GlobalTheory.\allowbreak{}Mem.\allowbreak{}Top.\allowbreak{}Bl.\allowbreak{}exists\_\allowbreak{}baseLattice\_\allowbreak{}of\_\allowbreak{}pair\_\allowbreak{}alt}{Development.\allowbreak{}AutomorphicArithmetic.\allowbreak{}ThetaCorrespondence.\allowbreak{}ThetaPairLifts.\allowbreak{}ThirdPlaceSeries.\allowbreak{}ExistsBaseLatticeOfPrime}{36--56}
\leanref{GlobalTheory.\allowbreak{}Mem.\allowbreak{}Top.\allowbreak{}Bl.\allowbreak{}memberE\_\allowbreak{}of\_\allowbreak{}steps2}{Development.\allowbreak{}AbelianVarietyTheory.\allowbreak{}Proofs.\allowbreak{}IwasawaDecomposition.\allowbreak{}ExistsMemberOfConstructionStepsSecondPass}{38--62}

For the enlargement assertion, start with the given untwisted member and its export data. For each residual characteristic unite the requested avoidance set with the avoidance set for transporting the column. In the local character choice exclude every host place over these primes, as well as the places already excluded by the host. Unpack Definition~\ref{B8.07d} for this member and keep its integers \(N_0(a)\) and its basis sets. If \(w\) is selected, its index is eligible by restriction to the fixed coefficient subfield. The last condition at \(w\) gives full rank and the integral basis property for these particular sets, and the three index conditions hold there. Apply the uniform assertion with those sets, the enlarged exclusions and \(e=\ell_a^{-2}\). This is where the triple \((Q_a,k_a,\xi_a)\) is chosen. In the simultaneous specialization of Lemma~\ref{B8.07c}, its Selmer test is evaluated on \((\zeta^{h\sigma_s\ell^{e_j}}-1)_j\) for the sign \(\sigma_s\) of the signature. The theta tests use the same exponents and the same \(Q_a\). To make the choices simultaneously, choose one selected representative for each pair of conjugate restrictions to the fixed coefficient subfield, and use its index throughout that pair. Enumerate the distinct indices so obtained; there are finitely many selected places. At each index add all previously chosen conductor sets to the exclusions. The new conductor is disjoint from their union, and each test with the original exclusions still holds, since the added exclusions only strengthen the disjointness requirement. In particular the global Selmer inequalities and the theta-lattice inequalities hold for the same \(\xi_a\), with the bases fixed at the start.
\leanref{GlobalTheory.\allowbreak{}Mem.\allowbreak{}Top.\allowbreak{}ExportAt}{Development.\allowbreak{}AbelianVarietyTheory.\allowbreak{}Definitions.\allowbreak{}AutomorphicMembership.\allowbreak{}AmbientBaseHeq}{36--41}
\leanref{GlobalTheory.\allowbreak{}HeightCounting.\allowbreak{}Up.\allowbreak{}nonempty\_\allowbreak{}seed\_\allowbreak{}of\_\allowbreak{}kept}{Development.\allowbreak{}AbelianVarietyTheory.\allowbreak{}Proofs.\allowbreak{}IwasawaDecomposition.\allowbreak{}TwistDataAgreeAcrossCoefficientsSymm}{133--211}
 On the selected places use the character, conductor set and exponent so obtained. Off the selected places take the trivial character, the empty conductor set and exponent zero. Lemma~\ref{B8.07a} ensures that over a fixed coefficient place the two prescriptions agree: the whole fibre is selected or the whole fibre is omitted. In the selected case the character choice is constant on this fibre.
\leanref{GlobalTheory.\allowbreak{}Mem.\allowbreak{}Top.\allowbreak{}TwChar.\allowbreak{}exists\_\allowbreak{}twDat\_\allowbreak{}stagesOn}{Development.\allowbreak{}AbelianVarietyTheory.\allowbreak{}Proofs.\allowbreak{}IwasawaDecomposition.\allowbreak{}NotMemSqOfRamificationIdx}{414--461}

The twist has now been specified, so its character set can be used as the argument of the requested finite coefficient set. Adjoin that set together with the algebraic elements required for transport and the roots needed for the chosen exponents. Transport the frame by the identity of preliminary data and pull the twist back to the larger coefficient field. The set of rational primes of the twist is unchanged. At each new place its character data are read at the place below it, so it satisfies the original avoidance condition. The new character set identifies with the one used in choosing the coefficient elements. The enlarged coefficient field therefore contains precisely the requested elements in the sense of item~4 of Theorem~\ref{B8.04}.
\leanref{GlobalTheory.\allowbreak{}HeightCounting.\allowbreak{}Up.\allowbreak{}exists\_\allowbreak{}twisted\_\allowbreak{}cohomology\_\allowbreak{}class\_\allowbreak{}lift\_\allowbreak{}alt}{Development.\allowbreak{}AutomorphicArithmetic.\allowbreak{}ThetaCorrespondence.\allowbreak{}ThetaPairLifts.\allowbreak{}ThirdPlaceSeries.\allowbreak{}ExistsTwistedCohomologyClassLift}{15--53}

Use the same column generators and place selections on this enlargement. Pull back the selected set by restriction of places. Transport the row realization and its character readings to the new column. Fix a selected place of the enlarged coefficient field, lying over \(v\), and write \(\ell,k,Q,\xi\) for the residue prime, exponent, conductor set and ray character chosen below it. The enlarged coefficient field contains all \(\ell^k\)-th roots of unity. The two readings of their cyclic group therefore differ by raising to one power \(b\) prime to \(\ell\), determined by the image of a primitive root. To identify the characters on the entire ray quotient, take any element of that quotient and lift it to the absolute Galois group of \(E\). Apply \ref{B8.42p} with the open normal kernel of the ray map and with the finite exclusion set consisting of \(Q\) and the places above \(\ell\). The resulting Frobenius represents the chosen element. Both readings give the inverse of the value of the same complex Hecke character at this prime, after their respective root identifications. Thus the enlarged member reads \(\xi^b\) on every element of the quotient. Its signed \(h\)-th power, for \(h\) prime to \(\ell\), reads \(\xi^{bh}\) or \(\xi^{-bh}\). Both exponents are prime to \(\ell\). The transported signature representations and local conditions agree at these readings, as do the theta generators and their projections. The bounds furnished before enlargement apply to these very powers and to the transported bases.
\leanref{GlobalTheory.\allowbreak{}HeightCounting.\allowbreak{}Rigid.\allowbreak{}pow\_\allowbreak{}of\_\allowbreak{}tied}{Development.\allowbreak{}AbelianVarietyTheory.\allowbreak{}Proofs.\allowbreak{}IwasawaDecomposition.\allowbreak{}RawPackEqPow}{23--47}
\leanref{GlobalTheory.\allowbreak{}HeightCounting.\allowbreak{}Up.\allowbreak{}GoodPrimeSelection.\allowbreak{}goodO\_\allowbreak{}of\_\allowbreak{}seed}{Development.\allowbreak{}AbelianVarietyTheory.\allowbreak{}Proofs.\allowbreak{}IwasawaDecomposition.\allowbreak{}GoodTwistOfSeedCharacterPowers}{157--228}
 The resulting member on the geometric extension projects to a constructed member on the origin extension, with both the goodness and twisting conditions. This proves the second member assertion with the frame and prime set held fixed throughout.
\leanref{GlobalTheory.\allowbreak{}Mem.\allowbreak{}Top.\allowbreak{}ReW.\allowbreak{}extend\_\allowbreak{}pinned}{Development.\allowbreak{}AbelianVarietyTheory.\allowbreak{}Proofs.\allowbreak{}IwasawaDecomposition.\allowbreak{}GensJAlgebraic}{184--256}
\leanref{GlobalTheory.\allowbreak{}HeightCounting.\allowbreak{}Up.\allowbreak{}secondNo\_\allowbreak{}of\_\allowbreak{}sMt\_\allowbreak{}sRat\_\allowbreak{}sCusp\_\allowbreak{}sFinl}{Development.\allowbreak{}AutomorphicArithmetic.\allowbreak{}ThetaCorrespondence.\allowbreak{}ThetaPairLifts.\allowbreak{}ThirdPlaceSeries.\allowbreak{}FiniteLevelAndSecond}{72--115}
\leanref{GlobalTheory.\allowbreak{}Mem.\allowbreak{}Top.\allowbreak{}Closing.\allowbreak{}goodE\_\allowbreak{}of\_\allowbreak{}builtE\_\allowbreak{}alt}{Development.\allowbreak{}AutomorphicArithmetic.\allowbreak{}ThetaCorrespondence.\allowbreak{}ThetaPairLifts.\allowbreak{}ThirdPlaceSeries.\allowbreak{}GoodEOfBuiltE}{16--31}

Write \(Q_o,Q'_o\) for the origin constructed and good member conditions. For a scale context \(K\), a point \(X\) on base \(B\), frame \(f\), twist \(T\) and auxiliary datum \(c\), augmentation gives
\[
 Q_A(K,X,f,T,c)\ \Longleftrightarrow\
 Q_o(K,X,f,T,c)\land A(K,B,f,T).                         \tag{8.15}
\]
The same identity holds for \(Q'_A,Q'_o\). For the first member take the common cutoffs for its construction and for \(A\). For export use the first conjunct. For the second member include the finite requirements for \(A\) on the enlarged base in the unions just used. The base condition persists there, so the eventual hypothesis on \(A\) applies there as well. Thus the same enlarged member satisfies both augmented conditions.
\leanref{GlobalTheory.\allowbreak{}Mem.\allowbreak{}Top.\allowbreak{}GroupRetwisting.\allowbreak{}pass\_\allowbreak{}step}{Development.\allowbreak{}AbelianVarietyTheory.\allowbreak{}Proofs.\allowbreak{}IwasawaDecomposition.\allowbreak{}PassStepCongr}{16--48}
\leanref{GlobalTheory.\allowbreak{}Mem.\allowbreak{}LiteralAutomorphy.\allowbreak{}memberE\_\allowbreak{}extensionWithAtProperty\_\allowbreak{}of\_\allowbreak{}memberE}{Development.\allowbreak{}AutomorphicArithmetic.\allowbreak{}ThetaCorrespondence.\allowbreak{}ThetaPairLifts.\allowbreak{}ThirdPlaceSeries.\allowbreak{}IsBuiltEPrmAIff}{35--53}

It remains to construct the linked bundle. In the definition preceding Theorem~\ref{B8.05} take the additional base condition to be the constant true condition. Work above the common cutoffs for the spectral data and the row multiplicities. Fix a good augmented point \(\beta\) and then an arbitrary \(c\in\R\). Unpacking its good-member witness gives a weighted point \(q\) with the same host datum and arithmetic object as \(\beta\), and local place matrices giving its kernel entries. The augmentation retains this witness and its projection to the list extension at the same point indices.

For \(q\) one may compute row multiplicities on the stage with the same group and prescribed tame series and with empty local index set. Copy the row realization by the group identification. The character labels, levels and central idempotent actions are preserved. Twisted multiplicity one on the copied realization gives equality of the two multiplicities. Transport the equality back along the same identification. This construction applies on the augmented extension since it uses the group and rows of \(q\).
\leanref{GlobalTheory.\allowbreak{}Mem.\allowbreak{}Top.\allowbreak{}Bl.\allowbreak{}witExE\_\allowbreak{}of\_\allowbreak{}witMv}{Development.\allowbreak{}AbelianVarietyTheory.\allowbreak{}Proofs.\allowbreak{}IwasawaDecomposition.\allowbreak{}ExistsIsEOfHeq}{33--37}
\leanref{GlobalTheory.\allowbreak{}Mem.\allowbreak{}Top.\allowbreak{}Bl.\allowbreak{}MultQClose.\allowbreak{}multiplicative\_\allowbreak{}closure\_\allowbreak{}prime\_\allowbreak{}mem\_\allowbreak{}alt}{Development.\allowbreak{}AutomorphicArithmetic.\allowbreak{}ThetaCorrespondence.\allowbreak{}ThetaPairLifts.\allowbreak{}ThirdPlaceSeries.\allowbreak{}MultiplicativeClosurePrimeMem}{18--39}

The good host has level set \(\{g_0,g_1,p_{\rm de},p'_{\rm de},q_0\}\), so its level agrees with the level of \(q\). Take the spectral comparison on the arithmetic datum of \(q\) at that level. Its local indices, integral modules and maps from the twisted layers belong to this datum. Adjoining the kernel to the spectral comparison preserves its projection to the same spectral comparison, so these local data are also those used by the recursion. Combine them with the kernel of \(q\) using the common column frames: the coordinate changes in the kernel and spectral description are the identity at the pinned entries. The lower row maps in these coordinates agree on the reference and twisted rows. The Artin and ray-field inputs give the host data for the same place selections. Together these choices give a recursion realization \(R\) with the prescribed spectral and kernel components.
\leanref{GlobalTheory.\allowbreak{}Mem.\allowbreak{}Top.\allowbreak{}Bl.\allowbreak{}exists\_\allowbreak{}kernel\_\allowbreak{}spectral\_\allowbreak{}data\_\allowbreak{}at\_\allowbreak{}witnesses}{Development.\allowbreak{}AutomorphicArithmetic.\allowbreak{}ThetaCorrespondence.\allowbreak{}ThetaPairLifts.\allowbreak{}ThirdPlaceSeries.\allowbreak{}TiedSKAtWitnesses}{23--70}
\leanref{GlobalTheory.\allowbreak{}Mem.\allowbreak{}Top.\allowbreak{}Bl.\allowbreak{}WPt.\allowbreak{}exists\_\allowbreak{}recPinE\_\allowbreak{}alt}{Development.\allowbreak{}AutomorphicArithmetic.\allowbreak{}ThetaCorrespondence.\allowbreak{}ThetaPairLifts.\allowbreak{}ThirdPlaceSeries.\allowbreak{}ExistsRecPinE}{21--31}

Use the saturated spectral spaces already attached to \(q\). They apply to \(R\) because its column datum is the column datum of \(q\). Their columns lie in the specified Hilbert space. The multiplicity equality above gives the nonzero row pieces at each required place. Match the square-integrable columns of the prescribed kernel to these pieces, using the row and column realizations of \(q\). The matching supplies a definite face, its spectral identification and the isometry on the indefinite output with its prescribed norm normalization. Use the same matched data for the arithmetic pair of row sides. Its lower maps are the ones just identified. On the indefinite side take the transported row coordinates and the same place embeddings and finite indices as on \(q\). The character idempotents and the identified row multiplicities give the labels on each nonzero eigenspace. Thus the arithmetic and indefinite realizations have the same face and spectral component.
\leanref{GlobalTheory.\allowbreak{}Mem.\allowbreak{}R3.\allowbreak{}SelectionPairing.\allowbreak{}exists\_\allowbreak{}selection\_\allowbreak{}pair\_\allowbreak{}local\_\allowbreak{}langlands\_\allowbreak{}alt}{Development.\allowbreak{}AutomorphicArithmetic.\allowbreak{}ThetaCorrespondence.\allowbreak{}ThetaPairLifts.\allowbreak{}ThirdPlaceSeries.\allowbreak{}ExistsSelectionPairLocalLanglands}{69--87}
\leanref{GlobalTheory.\allowbreak{}Mem.\allowbreak{}R3.\allowbreak{}SelRealIn.\allowbreak{}realIn\_\allowbreak{}exists\_\allowbreak{}alt}{Development.\allowbreak{}AutomorphicArithmetic.\allowbreak{}ThetaCorrespondence.\allowbreak{}ThetaPairLifts.\allowbreak{}ThirdPlaceSeries.\allowbreak{}RealInExists}{54--92}
\leanref{GlobalTheory.\allowbreak{}Mem.\allowbreak{}R3.\allowbreak{}RecInAt.\allowbreak{}recIn\_\allowbreak{}parts\_\allowbreak{}exists\_\allowbreak{}alt}{Development.\allowbreak{}AutomorphicArithmetic.\allowbreak{}ThetaCorrespondence.\allowbreak{}ThetaPairLifts.\allowbreak{}ThirdPlaceSeries.\allowbreak{}RecInPartsExists}{30--69}

Finally take the integers \(d_0,h_0\) from the numerical data of \(\beta\), so that \(\Delta=d_0\) and the size coordinate is \(mh_0\). At the scale prescribed by \(d_0,h_0,c\), form the bundle of Definition~\ref{B8.02} from \(R\), the saturated spectral realization, the arithmetic row sides with their lower maps, these integers, the square-integrable columns and the matched indefinite realization. The last realization uses exactly the face, spectral space and scale of the preceding components. This gives a bundle over \(q\) with recursion component \(R\). Transport it across the equality of the host and arithmetic object with those of \(\beta\). The point \(q\), that bundle and this equality constitute the required linkage. The cutoffs were chosen before \(c\).
\leanref{GlobalTheory.\allowbreak{}Mem.\allowbreak{}R3.\allowbreak{}SelAll.\allowbreak{}bwit\_\allowbreak{}exists\_\allowbreak{}alt}{Development.\allowbreak{}AutomorphicArithmetic.\allowbreak{}ThetaCorrespondence.\allowbreak{}ThetaPairLifts.\allowbreak{}ThirdPlaceSeries.\allowbreak{}BwitExists}{53--78}
\leanref{GlobalTheory.\allowbreak{}Mem.\allowbreak{}Top.\allowbreak{}Bl.\allowbreak{}exBL\_\allowbreak{}isLinked\_\allowbreak{}alt}{Development.\allowbreak{}AutomorphicArithmetic.\allowbreak{}ThetaCorrespondence.\allowbreak{}ThetaPairLifts.\allowbreak{}ThirdPlaceSeries.\allowbreak{}ExistsLinkedOfEq}{31--51}
\end{proof}

\subsection*{Proof of \ref{B8.08a}}\label{prf:B8.08a}\phantomsection
\begin{proof}
Fix the word. Write \(F\) for the maximal totally real subfield of the host \(E\), and \(\xi_t\) for the finite-order character by which tuple \(t\) twists the reference tuple \(t_0\). Thus \(\xi_{t_0}=1\). Put \(m=2n\). Denote the level-measure scalar and the column-kernel scalar at step \(j\) by \(I_{t,j}\) and \(K_{t,j}\). Write \(V_{t,j}\) for the vector scalar of the double step beginning at even \(j\), and \(P_t\) for the scalar of the terminal period. The scalar that commutes the column change through the last lift is a single \(s_u\ne0\) for this word. With these conventions the comparison uses
\[
 c_t=s_u P_t\prod_{\substack{0\le j<m\\j\ {\rm even}}}
           \frac{V_{t,j}K_{t,j}K_{t,j+1}}{I_{t,j}I_{t,j+1}} .
\]
Indeed, a double step contributes its two kernel scalars and its vector scalar on the column side, and its two level measures on the row side. The terminal period occurs once.

First consider the vector and level factors. For each finite translating element \(b_j\) of the row chain choose an idele \(x_j\) with \(x_j/\bar x_j=\det b_j\). At the endpoint use the translating element of the terminal group for \(b_m\). In the Weil action the source determinant occurs with the inverse splitting character and the target determinant with the splitting character. Twisting both characters by \(\xi_t\) therefore multiplies the step-\(j\) vector by
\[
                   z_{t,j}=\xi_t(x_j)^{-1}\xi_t(x_{j+1}).
\]
The real changes of frame have determinant representatives with finite part one. The finite-order character is trivial on the connected infinite idele group of \(E\), so these changes give the same formula for every real test vector. In particular \(z_{t,j}^{\operatorname{ord}(\xi_t)}=1\). The reference vector scalar is common to the tuples and nonzero: at the type primes it is the scalar of the fixed reference vacuum, and the twist operators are supported away from those primes. The two sets of local operators act on distinct tensor factors.
\leanref{GlobalTheory.\allowbreak{}TupRoot.\allowbreak{}weil\_\allowbreak{}tuple\_\allowbreak{}root\_\allowbreak{}alt}{Development.\allowbreak{}AutomorphicArithmetic.\allowbreak{}ThetaCorrespondence.\allowbreak{}ThetaPairLifts.\allowbreak{}ThirdPlaceSeries.\allowbreak{}WeilTupleRoot}{28--83}
\leanref{GlobalTheory.\allowbreak{}StepVacForm.\allowbreak{}main\_\allowbreak{}alt\_\allowbreak{}alt}{Development.\allowbreak{}AutomorphicArithmetic.\allowbreak{}ThetaCorrespondence.\allowbreak{}ThetaPairLifts.\allowbreak{}ThirdPlaceSeries.\allowbreak{}ImportsBis}{25--75}

For the torus factor, let \(a_i\) be the idele of the specified local torus element at marked index \(i\). Its contribution at a double step is
\[
       M_{t,j}=\prod_i r_{ji}\,\mu_{t,j+1}(a_i)^{e_{ji}},
       \qquad r_{ji}>0,\quad e_{ji}\in\Z ,
\]
where \(r_{ji}\) and \(e_{ji}\) are the real multiplier and the character exponent of the fixed torus test. The tuple character at this step is \(\mu_{t_0,j+1}\xi_t\). Thus \(M_{t,j}/M_{t_0,j}=\prod_i\xi_t(a_i)^{e_{ji}}\), also killed by \(\operatorname{ord}(\xi_t)\). Let \(J_{t,j}\) be the product of the local indices of the new step level in the reference level, in the common level frames. Subdivision into cosets gives \(I_{t,j}J_{t,j}=I_{t_0,j}\). The two local averages in \(V_{t,j}\) contribute \((J_{t,j}J_{t,j+1})^{-1}\). Hence their contribution after division by the level measures is exactly
\[
 \frac{(J_{t,j}J_{t,j+1})^{-1}}{I_{t,j}I_{t,j+1}}
           =\frac{1}{I_{t_0,j}I_{t_0,j+1}} .
\]
The remaining factor of \(V_{t,j}\), apart from the two vector scalars and \(M_{t,j}\), is the common local integral with the fixed type-prime tests. The intermediate-step multiplier in that integral is \(1\) on the last pair. We have thus computed the relative contribution of \(V_{t,j}/(I_{t,j}I_{t,j+1})\).
\leanref{GlobalTheory.\allowbreak{}SeeSaw.\allowbreak{}mono\_\allowbreak{}pow\_\allowbreak{}eq}{Development.\allowbreak{}AbelianVarietyTheory.\allowbreak{}Proofs.\allowbreak{}DecoratedLatticeProofs.\allowbreak{}ListProdMapNeZero}{28--40}
\leanref{GlobalTheory.\allowbreak{}Fam.\allowbreak{}lawsLinks\_\allowbreak{}of\_\allowbreak{}geo\_\allowbreak{}rows\_\allowbreak{}cols\_\allowbreak{}alt}{Development.\allowbreak{}AutomorphicArithmetic.\allowbreak{}ThetaCorrespondence.\allowbreak{}ThetaPairLifts.\allowbreak{}ThirdPlaceSeries.\allowbreak{}LawsLinksOfGeoRowsCols}{36--189}

We compute the normalization of the column kernels. At a twisting place, a chart at level \(j\) is a choice of local coset representative in the finite level average. For source and target charts \(\alpha,\gamma\), write the splitting multiplier of the local kernel as a source factor times a target factor. Before the last two steps these are signs. More explicitly, if \(N_j=k+j(t_{\rm ch}+1)\) is the column rank, put \(h_j=N_j+m-1-j\) for \(j<m\), where \(t_{\rm ch}\) is the auxiliary rank. The source sign is the sign of the source chart raised to \(h_{j+1}\), and the target sign is the sign of the target chart raised to \(h_j\). These exponents come from the determinant on the other space, including its \(m-1-j\) extra coordinates at level \(j\). At an intermediate chart the two signs cancel since
\[
       h_j+h_{j+2}=2(k+m-1+(j+1)t_{\rm ch}).
\]

At the join of the last two steps, write the reduction of the change of parahoric representative as a block triangular matrix with diagonal blocks \(A,D\). Let \(\delta\) be the tame character of the twisting place, whose residue field has order \(q\). The parahoric line contributes \(\delta(\det D)\). The determinant in the splitting contributes \(\delta(\det g)^{-1}\), since the unramified part is one on the unit \(\det g\). Reduction of the determinant gives
\[
       \delta(\det D)\delta(\det g)^{-1}
                        =\delta(\det A)^{-1}.
\]
This has order dividing \(q-1\). Moving the same representative from one side of the join to the other gives its inverse. Consequently at each place the local multiplier can be written \(\sigma_j(\alpha)\tau_j(\gamma)\) with both factors roots of unity and \(\tau_j(\gamma)\sigma_{j+1}(\gamma)=1\). The ordinary joins use the sign calculation above. The last join uses the block determinant calculation. Multiplying over the twisting places gives factors \(p_j(\alpha),q_j(\gamma)\) with the same cancellation.
\leanref{GlobalTheory.\allowbreak{}LevelLaw.\allowbreak{}last\_\allowbreak{}steps\_\allowbreak{}law}{Development.\allowbreak{}AbelianVarietyTheory.\allowbreak{}Proofs.\allowbreak{}DecoratedLatticeProofs.\allowbreak{}ParahoricChainLastTwoStepsRootOfUnity}{127--197}
\leanref{GlobalTheory.\allowbreak{}TwPlaces.\allowbreak{}exists\_\allowbreak{}roots\_\allowbreak{}of\_\allowbreak{}unity\_\allowbreak{}signs\_\allowbreak{}facAt}{Development.\allowbreak{}AutomorphicArithmetic.\allowbreak{}ThetaCorrespondence.\allowbreak{}ThetaPairLifts.\allowbreak{}ThirdPlaceSeries.\allowbreak{}RatDatumAddTwistPlaces}{176--227}

There is also a determinant character in the reading of a chart. Write \(x_{j,\alpha}/\bar x_{j,\alpha}=\det g_{j,\alpha}\) for its determinant representative and \(\chi_{j+1}\) for the step character of the tuple under consideration. Set
\[
 A_j(\alpha)=p_j(\alpha)\chi_{j+1}(x_{j,\alpha}),\qquad
 B_j(\gamma)=q_j(\gamma)\chi_{j+1}(x_{j+1,\gamma})^{-1}.
\]
These are roots of unity. Indeed \(d=\det g_{j,\alpha}\) has infinite part one and finite part an integral unit: the finite matrix and its inverse preserve the standard lattice. A continuous circle character on the compact group of finite integral units has finite image, since this group is profinite and the circle has a neighborhood of one containing no nontrivial subgroup. Some positive power of the step character is trivial on the ideles of \(F\), by its prescribed restriction and the finite order of the twist. Now \(x_{j,\alpha}^2=d(x_{j,\alpha}\bar x_{j,\alpha})\), with the second factor an idele of \(F\). Taking a common multiple of the two character orders proves the assertion.
\leanref{GlobalTheory.\allowbreak{}CharRoot.\allowbreak{}pow\_\allowbreak{}eq\_\allowbreak{}one\_\allowbreak{}of\_\allowbreak{}isDetRoot}{Development.\allowbreak{}AbelianVarietyTheory.\allowbreak{}Proofs.\allowbreak{}DecoratedLatticeProofs.\allowbreak{}TupleCharacterAtBasePointDetRootRootOfUnity}{41--56}

The value \(B_j(\gamma)A_{j+1}(\gamma)\) is a constant \(k_j\) on the intermediate charts. The local factors \(q_jp_{j+1}\) cancel. The remaining character is \(\chi_{j+2}/\chi_{j+1}\) on \(x_{j+1,\gamma}\). Its restriction to the ideles of \(F\) is trivial, so the choice of determinant representative does not affect it. Before the penultimate level the charts give the same standard finite matrix. At the penultimate level their additional determinants are local units at the twisting places. The two characters have the same tame restriction there, so their quotient is one on these units. This proves constancy, and \(k_j\) is a root of unity. The initial chart is unique. Let \(d_0\) be its value of \(A_0\), put \(d_j=d_0\prod_{i<j}k_i\), and rescale the chart vectors by
\[
             s_0(\alpha)=1,\qquad s_{j+1}(\gamma)=d_j B_j(\gamma).
\]
Induction gives \(s_j(\alpha)A_j(\alpha)=d_j\). The residual multiplier of each rescaled kernel entry is therefore
\[
        \frac{s_j(\alpha)}{s_{j+1}(\gamma)}
                              A_j(\alpha)B_j(\gamma)=1.
\]
In this normalization \(K_{t,j}\) is the product of the common step volume, the common normalization of the reference kernel and its initial column scalar. The terminal scalar is the common terminal volume times \(\sum_\beta s_m(\beta)\). The terminal level has one local label at each twisting place, so the product of the local label sets is a singleton, also for the empty twisting set. Its one summand is a root of unity.
\leanref{GlobalTheory.\allowbreak{}ChartQuot.\allowbreak{}chr\_\allowbreak{}quot\_\allowbreak{}chart\_\allowbreak{}free\_\allowbreak{}alt}{Development.\allowbreak{}AutomorphicArithmetic.\allowbreak{}ThetaCorrespondence.\allowbreak{}ThetaPairLifts.\allowbreak{}ThirdPlaceSeries.\allowbreak{}ChrQuotChartFree}{24--71}
\leanref{GlobalTheory.\allowbreak{}Fam.\allowbreak{}QuadraticExtensionData.\allowbreak{}dat\_\allowbreak{}exists}{Development.\allowbreak{}AutomorphicArithmetic.\allowbreak{}ThetaCorrespondence.\allowbreak{}ThetaPairLifts.\allowbreak{}ThirdPlaceSeries.\allowbreak{}DatExists}{22--90}
\leanref{GlobalTheory.\allowbreak{}Fam.\allowbreak{}colSide\_\allowbreak{}alt}{Development.\allowbreak{}AutomorphicArithmetic.\allowbreak{}ThetaCorrespondence.\allowbreak{}ThetaPairLifts.\allowbreak{}ThirdPlaceSeries.\allowbreak{}ColSide}{35--222}

In the quotient \(c_t/c_{t_0}\) the common nonzero factors cancel. The product of the orders of the finitely many remaining roots of unity gives an admissible \(N>0\), with \(r=1\) for these choices. Complex conjugation preserves the power identity.

Finally let \(p\) be the residue characteristic of the coefficient place \(w\). Off the selected twisting places of the coefficient field, \(\mu_w=t_0\) and \(Q_{t_0}\) is empty. At a selected \(w\), the character is the product of the auxiliary character for \(p\) and the chosen ray twist, which is unramified above \(p\). The auxiliary character uses the other exceptional prime when \(p\) equals its first exceptional prime. Its local restriction above \(p\) is therefore unramified as well. Every place in \(Q_{\mu_w}\) is ramified for the product. Hence none lies above \(p\).
\end{proof}

\subsection*{Proof of \ref{B8.08d}}\label{prf:B8.08d}\phantomsection
\begin{proof}
First work before the good projector. Let \(\sigma\) be the product at infinity of the prescribed discrete series, \(w\) its product of unit lowest-type vectors, and \(P\) the convolution by its normalized matrix coefficient. Use the finite translation carrying the reference level and its character to the spectral level and its character. In the linked realization this translation is obtained by carrying the two local level frames to their specified frames and taking the identity at the places where they already agree. Right translation is invertible and preserves the quotient measure.

Evaluation at \(w\) identifies the space of level intertwiners with \(H\). In detail, formal-degree orthogonality gives
\[
              PR(g)P=\langle w,\sigma(g)w\rangle P .
\]
For \(z\in H\), this identity shows that \(\sigma(g)w\mapsto R(g)z\) preserves the Gram matrix up to the scalar \(\|z\|^2\). It therefore extends from the cyclic span of \(w\) to a bounded intertwiner. Its image has the required level since the finite action commutes with \(G_\infty\). Evaluation recovers \(z\); cyclicity gives uniqueness. This construction takes values in the cuspidal space: \ref{B8.02y} gives integrability at the rank and weight gaps of the construction, and the constant-term calculation in \ref{B8.03a} annihilates each proper constant term on the projected space. The latter proof identifies the projected image of bounded inputs with the entire image, so its vanishing applies to every \(z\). On smooth vectors the intertwiner gives smooth automorphic functions by differentiation under right translation and local Sobolev embedding. Conjugation and the inverse finite translation give precisely the reference multiplicity maps. At infinity conjugation identifies the prescribed contragredient models by their discrete-series parameters in \(P\). Thus translated evaluation gives an isomorphism \(J\) from the full space of reference multiplicity maps onto \(\overline H\), the complex conjugate of \(H\) in \(L^2\).
\leanref{GlobalTheory.\allowbreak{}Mem.\allowbreak{}SRe.\allowbreak{}exists\_\allowbreak{}multR\_\allowbreak{}eval\_\allowbreak{}of\_\allowbreak{}star\_\allowbreak{}mem}{Development.\allowbreak{}AbelianVarietyTheory.\allowbreak{}Proofs.\allowbreak{}IwasawaDecomposition.\allowbreak{}ExistsRowLevelTranslatedEvaluationOfConjRow}{28--52}
\leanref{GlobalTheory.\allowbreak{}Mem.\allowbreak{}SRe.\allowbreak{}evalOn\_\allowbreak{}all}{Development.\allowbreak{}AbelianVarietyTheory.\allowbreak{}Proofs.\allowbreak{}IwasawaDecomposition.\allowbreak{}LowestVectorEvaluationIdentifiesLevelRowsConj}{19--87}

Write \(e_0\) for the good idempotent acting on the values of a reference multiplicity map \(S\). In the identified level frames, its coefficients are the complex conjugates of those of \(G\). Evaluation commutes with each finite translate, with the two finite actions identified by the change of frame. Summing over the finite group gives
\[
                     J(e_0 S)=\overline{G\,\overline{J(S)}} .
\]
If \(z\in GH\), take \(S=J^{-1}(\bar z)\). The identity and injectivity of \(J\) give \(e_0S=S\). Conversely \(e_0S=S\) gives \(\overline{J(S)}\in GH\). Thus the fixed subspace of \(e_0\) maps onto \(\overline{GH}\).
\leanref{GlobalTheory.\allowbreak{}Mem.\allowbreak{}SRe.\allowbreak{}evalOn\_\allowbreak{}good}{Development.\allowbreak{}AbelianVarietyTheory.\allowbreak{}Proofs.\allowbreak{}IwasawaDecomposition.\allowbreak{}EvaluationGoodRowsToConjugateGoodRows}{58--216}

We identify the arithmetic space in this calculation. In the row coordinates of the linked point, let \(j\) be the reference coordinate isomorphism from \(E_{\rm ar}\) to the reference row of the chain in tensor coordinates. The latter row is the image of the \(e_0\)-fixed multiplicity maps under linearization of a multilinear map on the archimedean factors. Denote this linearization by \(L\), and the realization of an arithmetic coordinate \(x\) as a multiplicity map by \(S_x\). The compatibility of the reference row coordinates in \textup{(8.6)} reads \(j(x)=L(S_x)\). For an \(e_0\)-fixed map \(S\), the vector \(x=j^{-1}(L(S))\) therefore satisfies \(S_x=S\), since linearization on a tensor product is injective. Conversely \(j(x)\) lies in that reference row, so \(S_x\) is fixed. This gives equality of ranges before taking dimensions. The two row norms are the same Petersson pairing multiplied by their respective positive measure normalizations. Dividing the composition by its positive norm factor yields the specified isometry onto \(\overline{GH}\). Conjugating gives the spectral embedding onto \(GH\).
\leanref{GlobalTheory.\allowbreak{}Mem.\allowbreak{}R3.\allowbreak{}RowsLvl.\allowbreak{}range\_\allowbreak{}eE\_\allowbreak{}and\_\allowbreak{}norm}{Development.\allowbreak{}AbelianVarietyTheory.\allowbreak{}Proofs.\allowbreak{}IwasawaDecomposition.\allowbreak{}RangeEEAndNorm}{68--119}
\leanref{GlobalTheory.\allowbreak{}Mem.\allowbreak{}SRe.\allowbreak{}exists\_\allowbreak{}rowJ\_\allowbreak{}of\_\allowbreak{}evalOn}{Development.\allowbreak{}AbelianVarietyTheory.\allowbreak{}Definitions.\allowbreak{}AutomorphicMembership.\allowbreak{}ConcreteLevelAdelicSubset}{88--153}

The finite-type assignment places each transported basis vector in its assigned type space. Apply the single-label idempotent to a basis expansion of any vector in the space for a good \(t\). Orthogonality of the idempotents in \ref{B7.44} deletes exactly the terms whose labels differ from \(t\). The retained vectors therefore span that space, and they are independent by the isometry.
\leanref{Formalism.\allowbreak{}L2.\allowbreak{}SG.\allowbreak{}AllRowsOf.\allowbreak{}rowsLab\_\allowbreak{}of}{Formalism.\allowbreak{}LocalToGlobal.\allowbreak{}SiegelLatticeTheory.\allowbreak{}FlagStabilizerTensors.\allowbreak{}AllRowsOf}{864--882}
\end{proof}

\subsection*{Proof of \ref{B8.08g}}\label{prf:B8.08g}\phantomsection
\begin{proof}
Write \(a,b,c,d\) for the lengths of \(A,B,C,D\). The centered lists are regular Harish--Chandra parameters for discrete series on both sides. We apply Atobe's unconditional discrete-series criterion. Its branching input is He's theorem for pairs of discrete series, as Atobe explains immediately after the statement \cite[Definition~1.6 and Theorem~1.7]{atobe-2020}. For the smaller source the centered entries are integers and none is zero. Thus the integer \(k_\lambda\) of that criterion is \(-1\), and its two signature indices are \(a+d,b+c\). The desired target has signature
\[
                    (a+d+t,b+c).
\]
In part~(1) of the criterion its parameters are \(h=(t-1)/2\) and \(l=0\), where the target signature in that part is \((a+d+2h+l+1,b+c+l)\). The two exceptional sets counted there are empty. Indeed any entry \(\alpha\) they count would have \(0\le\pm\alpha-1<h\), contrary to the bound on \(|\alpha|\). The lift to this signature is nonzero. Removing \(i>0\) planes keeps \(h\) fixed and gives \(l=-i\), violating the necessary inequality \(l\ge0\).

In the reverse direction the centered list contains \(I_t\) in one sign. Neither neighbor \((t+1)/2,-(t+1)/2\) occurs, so this is the maximal central string and \(k_\lambda=t\). The signature indices after deleting the string are \(a+b,c+d\). The smaller target has exactly this signature. The equal-increment clause of part~(2) gives nonvanishing at \((a+b+l,c+d+l)\) precisely for \(l\ge0\). Again every lower member in this Witt tower has \(l=-i<0\).

For the global assertion, a nonzero lift to a lower member would give a nonzero local theta pairing at the chosen place. Indeed, on pure tensors fix all the other local inputs and evaluate the lifted form. Varying the two inputs at the chosen place gives a pairing invariant under the source action, by the change of integration variable in the theta integral. It factors through the local theta quotient there. The vanishing just proved kills this pairing for every choice of the other inputs. Linearity and continuity give vanishing for the allowed tests. Theorem~\ref{B3.00e} identifies the constant term along the parabolic of an isotropic \(i\)-plane with the lower lift. Thus all proper constant terms vanish. On an anisotropic target there are no such parabolics.
\end{proof}

\subsection*{Proof of \ref{B8.08e}}\label{prf:B8.08e}\phantomsection
\begin{proof}
On a right half-plane multiply the Euler factors at each \(w\). The resulting function is
\[
 L^S(B\otimes\mu_{j+1}^{-1},s+1/2)
 \prod_{i=1}^{j}\prod_{l=0}^{t-1}
       L^S(\mu_i/\mu_{j+1},s+t/2-l).
\]
The first factor is holomorphic and nonzero for \(\Re s>(r+1)/2\). To see this, put \(\sigma=\Re s+1/2-r/2>1\). The absolute logarithmic series on any closed half-plane with this lower bound is at most
\[
   \frac{r[E:\Q]}{1-2^{-\sigma}}\sum_{n=2}^{\infty}n^{-\sigma}.
\]
This uses unitarity of the twist and the fact that at most \([E:\Q]\) primes lie above a rational prime. Exponentiating the locally uniformly convergent logarithm gives the claim.

At \(s=t/2\) the remaining arguments are exactly \(t,t-1,\ldots,1\). The ratio of two distinct indexed characters has nonzero angular exponent at the specified embedding. Lemma~\ref{B8.11h} applies to each ratio, including at the argument one and with the finite Euler factors in \(S\) removed. All these finitely many functions are holomorphic and nonzero at their indicated arguments. Since \(t/2>(r+1)/2\), their product with the bottom factor gives the required continuation and value. For \(j=0\) the product over \(i\) is empty.
\end{proof}

\subsection*{Proof of \ref{B8.08j}}\label{prf:B8.08j}\phantomsection
\begin{proof}
Fix the levels away from \(T\). The two type primes split completely, and \(g_0\ge(kb)^8\) gives \(g_0-1>2fk\). The stabilizer argument of Theorem~\ref{B8.54} already applies at this level. Indeed, a nontrivial finite stabilizer would contain an element of order \(g_0\), whereas its rational representation has dimension \(2fk<g_0-1\). Thus each class contributes a whole copy of its coefficient space. The product of the number of these classes and the coefficient dimension is precisely the mass, type indices and Weyl product defining \(d_u^\star\) in Lemma~\ref{B8.05a}.

Over each class the new level specifies one nonzero vector of \(\F_T^k\) at each of the \(f\) split places. The absence of a stabilizer gives \((T^k-1)^f\) distinct choices. Coordinatewise scalar multiplication acts freely on these choices. The diamond translations act on the corresponding coefficient fibres by isomorphisms. On each orbit choose a basis over one point and transport it by translation. Freeness makes this a well-defined basis over the whole orbit. Taking one copy of the orbit for each basis vector gives the asserted function model. The level norm at a place is the sum of its diamond translations, which identifies its kernel in this model.

The level window gives \(T_*\le T\le(1+1/(fk))T_*\). Since \(T\ge2\) and \(f\le k\), we have \(T^k\ge2^k\ge2f\). Consequently
\[
 (T^k-1)^f=T^{fk}(1-T^{-k})^f
 \ge T^{fk}(1-fT^{-k})\ge\tfrac12 T_*^{fk}.
\]
The upper bound is \((T^k-1)^f\le T^{fk}\le\exp(1)T_*^{fk}\). Multiply by \(d_u^\star\) and use the definition of \(T_*\).
\end{proof}

\subsection*{Proof of \ref{B8.08k}}\label{prf:B8.08k}\phantomsection
\begin{proof}
Throughout the proof \(C'm^5fL\le k\). In particular \(a_2\le k\), since \(m^2+2m^4\le3m^5 fL\). Also \(1+\log_2(2f)\le m\) by \(\log f\le\sqrt r\) and \(m=2r+2\ge40\). To bound the ceiling defining \(q\), put \(t=r+1=m/2\). We have \(t\ge1\), \(k\ge2t+3\). The contribution before the ceiling that contains \(c\) is at most
\[
 4ct\{k+(k+2t)2t+2t+(t-1)^2\}
 \le 36ckt^2 .
\]
The remaining contribution including the ceiling is at most
\(tk/2+12t^3+1\le12kt^2\). This gives \(q\le(9c+3)km^2\).

Put \(P=b^{Am^3f}\). Comparing the cutoffs in each of the \(f(r+1)\) factors of \(S_c\) and \(S_{C_h}\) gives \(h_Z\le P S_Z^4\). For a finite set of occurring labels we claim
\[
 \sum_Z S_Z^{-12(r+1)^2}\le (41^f P^3)^{(r+1)^2}.
\]
Indeed, in a single entry lattice the number of vectors with \(|y_\tau|\le R_\tau\), \(R_\tau\ge1\), is at most \(9^f\prod_\tau R_\tau^2\). The products of open disks of radius \(1/2\) about such vectors are disjoint by the product lower bound for differences. Compare their volumes with the product of disks of radii \(R_\tau+1/2\). Integrating this count against \(\prod_\tau 3R_\tau^{-4}\,dR_\tau\) bounds the sum of \(\prod_\tau\max(1,|y_\tau|)^{-3}\) by \(27^f\). Each normalized entry has product height at most \(P S_Z^4\). Multiply the resulting inequalities over the \((r+1)^2\) entries and enlarge the set of entry tuples to the product of its coordinate projections. The entry tuple determines \(Z\), so \(27^f\le41^f\) proves the claim.
\leanref{Formalism.\allowbreak{}MainNotions.\allowbreak{}BridgeLemma₂.\allowbreak{}rows\_\allowbreak{}sizeSum}{Formalism.\allowbreak{}SharpRestriction.\allowbreak{}CuspSpanNormBounds.\allowbreak{}CayleyDegreeNondegeneracy.\allowbreak{}RowsSiz}{129--179}

Fix \(s,i\) and two labels. Apply Lemma~\ref{B8.10l} with \(d=i\) in that lemma, using the stated count and the two amplitude inequalities. The factors \(b_i(Z)^kb_i(Z')^k\) cancel. Dividing by \(\sqrt{n_Zn_{Z'}}\) cancels \((\sqrt{G_ZG_{Z'}}/N)^k\) and leaves \((kv/(\exp(1)\pi))^{kf}\). The two square roots contribute one factor \(\Gamma F_i\). For the endpoint masses use
\[
 \Xi_s(h_Zh_{Z'})\le\Xi_s(P^2)(S_ZS_{Z'})^{4ck}.
\]
This also holds at interior ranks, where the mass is independent of its argument. After bounding \(h_Z,h_{Z'}\), the remaining exponent on each size is at most
\[
 4ck+4e_s+4cr^2+a_i/2-q/(r+1)\le-12(r+1)^2 .
\]
Apply the label sum once for each label. Per unit of \(d\), the resulting bound is
\[
 (2/g)^{(s-i)kf}
 \left(\frac{kv}{\exp(1)\pi}\right)^{kf}
 \Xi_s(P^2)(csf)^{6sf}\Gamma F_i
 P^{2e_s+2cr^2}(41^fP^3)^{2(r+1)^2}.
\]

Write \(H=b^{Cm^2f}\), and put \(\Xi_s=M_sM_{k-s}/M_k\) for \(s\le k-2\), \(\Xi_s=(Ck)^{kf}(H^{Cm}N)^{Ck}\) at the endpoints. Since \(2A\le C^2\), we have \(P^2\le H^{Cm}\) and \(\Xi_s(P^2)\le\Xi_s\). Also
\[
 2e_s+2cr^2+6(r+1)^2\le(6c+2)(s+m)m .
\]
Since \(b\ge2\), the factors \(2\cdot41^{2f(r+1)^2}\) cost at most \(b^{1+12f(r+1)^2}\). The choice \(C^2\ge4+A(6c+2)\) thus absorbs these factors and all the powers of \(P\) in \(H^{C(s+m)m^2}\). Define
\[
 F_s(C,\Gamma)=(kv)^{kf}\Xi_s C^{kf}H^{C(s+m)m^2}
                      (Csf)^{6sf}\Gamma .
\]
The two finite sums of corank \(s\), including the factor \(d\), are consequently bounded by
\[
 \begin{split}
 dR\,(2/g)^{(s-1)kf}F_s(C,\Gamma),\\
 d\,(2/g)^{(s-2)kf}F_s(C,\Gamma)
                  e^{Cm^6fL}(C/k)^{kf}.
 \end{split}
\]
In particular the degenerate part has one factor of \(d\), including at \(s=2\), where its power of \(2/g\) is zero.

We give the numerical summation. The inequalities \(8C_\Gamma\le C'\) and \(C'm^5 fL\le k\) give \(\Gamma\le e^{k/8}\). Hence \(C^{kf}\Gamma\le(\exp(1)C)^{kf}\). Set \(D=\exp(1)C\). Increasing \(C\) to \(D\) in the other positive factors and evaluating with \(\Gamma=1\) gives an upper bound for either row. We have \(k\ge16D^{32}m^5fL\), \(1\le v\le b\), \(1\le N\le b^{Dmf}\), and \(g\ge(kb)^8\).

For \(2\le s\le k-2\) we need
\[
 (kv)^{kf}\Xi_s\le b^{D^2m^2f}D^{skf}.
\]
To see this, put \(R_j=M_j/M_{j-1}\) and expand
\(\Xi_s=M_2\prod_{j=3}^sR_j/\prod_{j=k-s+1}^kR_j\) by Theorem~\ref{B8.48}. Pair each numerator index \(j\) with \(j+k-s-2\). All paired factorial quotients are at most one. The two remaining denominator factorials have product \((k-2)!(k-1)!\ge(k/\exp(1))^k\) for \(k\ge6\), by integration of \(\log x\). The net power of \(\pi/v\), before raising to \(f\), is \(2k-1+(s-2)(k-s-2)\), between \(k\) and \((s+1)k\). The Euler factors and their reciprocals are at most \(2^f\) at each index by the Euler product estimate in Lemma~\ref{B8.02bc}. Thus the factors apart from \(M_2(2v)^f\) are absorbed in \(D^{skf}(kv)^{-kf}\). The rank-two bound in the same lemma gives
\(M_2\le(20f(2v)^2)^{16f}\). Since \(2v\le b\) and \(\log f\le\sqrt r\), this gives \(M_2(2v)^f\le b^{D^2m^2f}\).

Put \(h=\log k\) and \(\ell=\log g\). For the interior ranks this mass bound gives
\[
 \frac{\log F_s(D,1)}{kf}
 \le(s+1)\log D+
       \frac{D^2[m^2+(s+m)m^4]L}{k}
       +\frac{6s}{k}\log(Dsf)
 \le \frac{sh}{8}.
\]
Indeed \(\log D\le h/32\), the middle term is at most \((s+2)/16\), and the last is at most \(sh/32\) since \(s,f\le k\) and \(k\ge 400\). For \(s=k-1,k\), the contribution of \((kv)^{kf}\Xi_sD^{kf}\) to the left side is at most
\[
 2h+2\log D+(1+D^3m^3+D^2m)L.
\]
Here we used the endpoint formula and \(N\le b^{Dmf}\). The last term is at most \(k/64\) by the rank threshold. Since \(s\ge k-1\), these three terms together are at most \(sh/16\). The remaining height and counting terms \(D^2(s+m)m^4L/k\) and \(6s\log(Dsf)/k\) are each at most \(sh/32\). This proves the same bound at both endpoints.

Now \(\ell\ge8h\) and \(Dm^6L/k\le1/16\), since \(m\le f\). For \(s\ge2\),
\[
 \begin{aligned}
 -(s-1)(\ell-\log2)+sh/8
     &\le-\ell/2-h/(kf),\\
 -(s-2)(\ell-\log2)+sh/8+Dm^6L/k+\log D-h
     &\le-h/2-h/(kf).
 \end{aligned}
\]
Both left sides decrease with \(s\), so it suffices to substitute \(s=2\). The two corank sums therefore cost at most \(k^{-1}g^{-kf/2}\) and \(k^{-1}e^{-B/2}\) per rank, apart from \(dR\) and \(d\), respectively. There are fewer than \(k\) ranks. Summing and taking the supremum over finite subsets proves the assertion.
\end{proof}

\subsection*{Proof of \ref{B8.08l}}\label{prf:B8.08l}\phantomsection
\begin{proof}
Absolute summability gives \(\sum_{t,t'}|\langle v_t,v_{t'}\rangle|\le(\sum_t\|v_t\|)^2<\infty\). Hence \(P(U)\) and \(|P|(U)\) are finite for every \(U\), and \(|P(U)|\le|P|(U)\).

Fix \(E_1\in\{E_{\rm gen},E_{\rm ng}\}\), and let \(\mathcal A(E_1)\) be the set of same-row ordered pairs in \(E_1\). Expand the squared norms as absolutely convergent double series. Then
\[
 \sum_i\|w_i(E_1)\|^2
 =\sum_i\ \sum_{\substack{t,t'\in E_1\\ \iota(\operatorname{row}(t))=\iota(\operatorname{row}(t'))=i}}\langle v_t,v_{t'}\rangle
 =\sum_{(t,t')\in\mathcal A(E_1)}\langle v_t,v_{t'}\rangle .
\]
The second equality uses (iii). The left side is real, so the right side equals its real part. Every term of \(E_1\) has nonzero Gram invariant. By (iv), \(\mathcal A(E_1)\) is therefore the disjoint union of its pairs of corank one and its pairs of corank at least two. The latter pairs form \(\mathcal B\cap(E_1\times E_1)\). Write \(\mathcal U(E_1)\) and \(c(E_1)\) for \(\mathcal U_{\rm gen},c_{\rm gen}\) or \(\mathcal U_{\rm ng},c_{\rm ng}\). By (ii),
\[
 \|\operatorname{part}(E_1)\|_{\rm HS}^2
 \le P_g\bigl(P(\mathcal B\cap(E_1\times E_1))+P(\mathcal U(E_1))\bigr)
 \le P_g\bigl(|P|(\mathcal B\cap(E_1\times E_1))+c(E_1)\bigr).
\]
The sets \(E_{\rm gen}\) and \(E_{\rm ng}\) are disjoint. So the two sets \(\mathcal B\cap(E_1\times E_1)\) are disjoint subsets of \(\mathcal B\). Since the summands of \(|P|\) are nonnegative, their two values of \(|P|\) add up to at most \(|P|(\mathcal B)\le c_{\rm col}\). The inequality \(\|x+y\|^2\le2\|x\|^2+2\|y\|^2\), applied to \(X=\operatorname{part}(E_{\rm gen})+\operatorname{part}(E_{\rm ng})\), proves the first assertion.

For the numerical part, use \(1+x\le e^x\), \(f\le k\), \(k\ge1\) and \(b\ge2\):
\[
 \left(1+\frac3{g_j}\right)^f\le\exp\frac{3f}{g_j}\le\exp\frac{3k}{k^8b^8}\le\exp\frac{3}{256}\qquad(j=0,1).
\]
Thus \(P_g\le e^{3/128}<6/5\). Put \(C_*=\max(C_{\rm gen},C_{\rm ng})\), so \(C_1=C_*+1\ge1\). Since \(Y\ge1\), we have \(2\le e^Y\), and hence
\[
 e^{C_{\rm gen}Y}+e^{C_{\rm ng}Y}\le2e^{C_*Y}\le e^{(C_*+1)Y}=e^{a}.
\]
Also \((g_0g_1)^{-kf/2}\le1\). Substituting these bounds and \(2P_g\le12/5\) into the first assertion gives the displayed estimate.
\end{proof}

\subsection*{Proof of \ref{B8.08n}}\label{prf:B8.08n}\phantomsection
\begin{proof}
For \(t\in E_{\rm gen}\) its class is \(c(t)=\{t'\in E_{\rm gen}:\operatorname{row}(t')=\operatorname{row}(t),\ s(t,t')=1\}\). If \(t\sim t'\), symmetry and transitivity show \(c(t)=c(t')\). Conversely, reflexivity gives \(t'\in c(t')\), so equality of the classes implies \(t\sim t'\). Hence \(\mathcal U_{\rm gen}=\coprod_{c\in\mathcal C}(c\times c)\). The Cauchy--Schwarz inequality for each inner product gives
\[
 \sum_{t,t'}|\langle v_t,v_{t'}\rangle|\le\Bigl(\sum_t\|v_t\|\Bigr)^2<\infty .
\]
We may therefore regroup the signed pair sum by classes. For each class, expanding the squared norm gives \(\sum_{t,t'\in c}\Re\langle v_t,v_{t'}\rangle=\|\sum_{t\in c}v_t\|^2\). The subdivisions by point and normal labels preserve absolute summability, and their subtotals sum to \(\sum_{t\in c}v_t\). This proves the first identity.

If \(J_2(\gamma_0)>0\), the second identity is the definition of \(a_2\). If \(J_2(\gamma_0)=0\), every term in the defining subtotal vanishes by hypothesis, so both sides are zero. The definition also gives \(a_2\ge0\). Every nonzero term belongs to its unique class and carries its assigned normal label; conversely, a nonzero term in such a fibre represents its point label. This proves the assertion about \(\mathcal P_*\).

For the finite amplitude estimate, group \(S_0\) by its assigned tuples. The factorization and the fibre bound give
\[
 \sum_{(c,\nu_0)\in S}a_2(c,\gamma_0,\nu_0)
 =\sum_{\mathbf e}\Bigl(\sum_{\substack{(c,\nu_0)\in S_0\\ \mathbf d(c,\nu_0)=\mathbf e}}\rho_c\Bigr)\prod_{p\in\mathcal V}|a_p(e_p)|^2
 \le V_0\prod_{p\in\mathcal V}\Bigl(\sum_{e_p\in D_p}w_p(e_p)|a_p(e_p)|^2\Bigr).
\]
Put \(q_p=8H_p^{c_0(r+1)^2+2m}\) for \(p\in\Omega\) and \(q_p=1\) for \(p\notin\Omega\). The local assumptions bound the last expression by \(V_0\prod_{p}H_p^{c_1(r+1)^2}q_p\prod_pn_p^{\rm gen}\). For \(p\in\Omega\) the inequality \(32\le H_p^5\) gives \(H_p^{c_1(r+1)^2}q_p\le H_p^{e_7}\), and for \(p\notin\Omega\) the same inequality follows from \(H_p\ge1\). Multiplying over \(p\) gives \(\prod_pH_p^{c_1(r+1)^2}q_p\le H_{\rm tot}^{e_7}\), which proves the bound since \(M_{\rm gen}\ge0\).
\end{proof}

\subsection*{Proof of \ref{B8.08oa}}\label{prf:B8.08oa}\phantomsection
\begin{proof}
At each \(v\in\mathcal V_{\mathrm{typ}}\), let \(U_v\le\GL_k(\F_{q_v})\) be the upper unitriangular subgroup. For every \(y\in\GL_k(\F_{q_v})\), the conjugate \(y\cdot 1\cdot y^{-1}=1\) lies in \(U_v\), where the Whittaker character takes the value \(1\). The local induced factor of Lemma~\ref{B8.10f} at \(1\) is therefore \(|\GL_k(\F_{q_v})|/|U_v|\). Taking the product over \(v\in\mathcal V_{\mathrm{typ}}\) gives \(\Gamma_1=|\mathcal I|/|\mathcal U|=|\mathcal I/\mathcal U|\ge1\).

Let \(n=\dim_E V\), fix the model's \(E\)-basis identifying \(V\) with \(E^n\), let \(A_f\) be the finite adele ring of \(F\), and embed \(V\) coordinatewise into \((E\otimes_F A_f)^n\) via \(E\hookrightarrow E\otimes_F A_f\). Let \(L_f\subset(E\otimes_F A_f)^n\) be the standard finite-adelic lattice, \(G_f\) the finite-adelic unitary group, \(K_{\max}\le G_f\) the stabilizer of \(L_f\), \(K_2\le K_{\max}\) the \(g\)-adic mirabolic level subgroup, and \(x_u\in K_{\max}\) the word shift. The flag \(F_u\) has isotropic basis \(e_1,\ldots,e_r\) spanning \(X_u\), dual isotropic vectors \(f_1,\ldots,f_r\) satisfying \(h_V(e_a,f_b)=\delta_{ab}\) and \(h_V(f_a,f_b)=0\), and final vector \(z_u\in X_u^\perp\) orthogonal to \(f_1,\ldots,f_r\), with \(x_u\) fixing \(e_1,\ldots,e_r,f_1,\ldots,f_r\). Let \(H_{\mathrm{fin}}(u)\le G_f\) be the subgroup fixing the embedded vectors \(e_1,\ldots,e_r,f_1,\ldots,f_r,z_u\), put
\[
 K_u=\{h\in H_{\mathrm{fin}}(u):x_u^{-1}hx_u\in K_2\},
\]
and let \(\Gamma_u\) be the rational subgroup of \(G(F)\) whose finite-adelic projection \(\iota_f:\Gamma_u\to G_f\) lands in \(H_{\mathrm{fin}}(u)\). Neatness of \(K_u\) on the tied model identifies \(C_u\) with the direct sum of copies of \(\pi_u\) indexed by the finite class set \(\mathrm{Cl}_u\) of representatives \(h_\rho\in H_{\mathrm{fin}}(u)\) for \(\iota_f(\Gamma_u)\backslash H_{\mathrm{fin}}(u)/K_u\), so \(d=|\mathrm{Cl}_u|n_\pi\). Partition \(\mathrm{Cl}_u=\coprod_{\Lambda\in\mathcal R_0(u)}\mathrm{Cl}_u(\Lambda)\) according to the assigned representative lattice \(\Lambda_\rho\in\mathcal R_0(u)\). For \(h\in G_f\), set
\[
 \mathcal L_u(h)=\{y\in V:(hx_u)^{-1}y\in L_f\},
\]
so that \(\mathcal L_u(h_\rho)=\Lambda_\rho\) for \(\rho\in\mathrm{Cl}_u\). If \(\kappa\in K_u\), then \((x_u^{-1}\kappa x_u)^{-1}\in K_2\subseteq K_{\max}\) preserves \(L_f\); writing \((h\kappa x_u)^{-1}=(x_u^{-1}\kappa x_u)^{-1}(hx_u)^{-1}\) gives \(\mathcal L_u(h\kappa)=\mathcal L_u(h)\). If \(\gamma\in\Gamma_u\), its rational matrix \(M\in\GL_n(E)\) preserves the Hermitian Gram matrix and has finite-adelic image \(\iota_f(\gamma)\). Because \(\iota_f(\gamma)\in H_{\mathrm{fin}}(u)\) fixes the embedded coordinates of \(e_a,f_a,z_u\) and \(E\to E\otimes_F A_f\) is injective, \(M\) defines an \(E\)-linear Hermitian isometry \(g_\gamma\in\operatorname{Stab}(F_u)\) fixing \(e_a,f_a,z_u\) and satisfying \(((\iota_f(\gamma)hx_u)^{-1})y=(hx_u)^{-1}(g_\gamma^{-1}y)\), whence \(\mathcal L_u(\iota_f(\gamma)h)=g_\gamma(\mathcal L_u(h))\). Whenever \(\mathcal L_u(h)=\Lambda\in\mathcal R_0(u)\) and \(h=\iota_f(\gamma)h_\rho\kappa\) with \(\gamma\in\Gamma_u\), \(\rho\in\mathrm{Cl}_u\), and \(\kappa\in K_u\), it follows that \(\Lambda=g_\gamma(\Lambda_\rho)\), and uniqueness of orbit representatives in \(\mathcal R_0(u)\) forces \(\Lambda_\rho=\Lambda\), so \(\rho\in\mathrm{Cl}_u(\Lambda)\).

Fix \(\Lambda\in\mathcal R_0(u)\) with \(\operatorname{Aut}_{F_u}(\Lambda)\) finite, and let \(A_\Lambda\subseteq\iota_f(\Gamma_u)\) be the set of \(\eta=\iota_f(\gamma)\) for which there exist \(h,h'\in H_{\mathrm{fin}}(u)\) with \(\mathcal L_u(h)=\mathcal L_u(h')=\Lambda\) and \(\kappa\in K_u\) satisfying \(h'=\eta h\kappa\). For \(\eta=\iota_f(\gamma)\in A_\Lambda\), the invariance formulas give \(\Lambda=\mathcal L_u(h')=g_\gamma(\mathcal L_u(h))=g_\gamma(\Lambda)\). Because \(E\to E\otimes_F A_f\) is injective, \(\eta=\iota_f(\gamma)\) uniquely determines \(M\) and hence \(g_\gamma\). Since \(g_\gamma\) is an isometry fixing \(e_1,\ldots,e_r\), it fixes \(X_u\) pointwise and preserves \(X_u^\perp\); passing to \(X_u^\perp/X_u\) gives a well-defined map \(\varphi(\eta)\in\operatorname{Aut}_{F_u}(\Lambda)\). Suppose \(\varphi(\eta_1)=\varphi(\eta_2)\), and write \(g_i=g_{\gamma_i}\). For \(x\in X_u^\perp\), the difference \(g_1(x)-g_2(x)\) lies in \(X_u\), so \(g_1(x)-g_2(x)=\sum_{a=1}^r c_a e_a\) with \(c_a\in E\). Because \(g_1,g_2\) are isometries fixing every \(f_b\) and \(h_V(f_b,e_a)=\delta_{ab}\),
\[
 c_b=h_V\bigl(f_b,g_1(x)-g_2(x)\bigr)=h_V(g_1(f_b),g_1(x))-h_V(g_2(f_b),g_2(x))=0
\]
for \(1\le b\le r\). Thus \(g_1=g_2\) on \(X_u^\perp\). Every \(v\in V\) decomposes as \(\bigl(v-\sum_{a=1}^r h_V(e_a,v)f_a\bigr)+\sum_{a=1}^r h_V(e_a,v)f_a\), whose first summand pairs to zero with every \(e_b\) and therefore lies in \(X_u^\perp\). Since \(g_1,g_2\) agree on \(X_u^\perp\) and fix each \(f_a\), we obtain \(g_1=g_2\) on \(V\); hence their rational matrices coincide, so their finite-adelic images satisfy \(\eta_1=\eta_2\). Hence \(|A_\Lambda|\le|\operatorname{Aut}_{F_u}(\Lambda)|\).

Set \(c_u(s)=x_u^{-1}sx_u\) for \(s\in H_{\mathrm{fin}}(u)\), and put \(H_{\mathrm{int}}(u)=\{s\in H_{\mathrm{fin}}(u):c_u(s)\in K_{\max}\}\). Every \(c_u(s)\) with \(s\in H_{\mathrm{int}}(u)\) preserves \(L_f\) and fixes the finite set \(W_u\) consisting of the embedded vectors \(e_1,\ldots,e_r,f_1,\ldots,f_r\) and \(x_u^{-1}z_u\). By the parabolic block structure modulo \(g\), the \(g\)-adic Levi reduction \(\bar\ell(s)=\operatorname{levi}(c_u(s))\in\mathcal I\) is a group homomorphism on \(H_{\mathrm{int}}(u)\) that maps \(K_u\) into \(\mathcal U\). To see that \(\bar\ell\) maps \(H_{\mathrm{int}}(u)\) onto \(\mathcal I\), let \(e\in A_f\) be the characteristic idempotent of \(\mathcal V_{\mathrm{typ}}\), let \(\mathcal O_{E,f}\) be the integral subring of \(E\otimes_F A_f\) with CM involution \(\sigma_A\), and let \(e_w\in\mathcal O_{E,f}\) be the model's split-place idempotent above \(g\), with \(e_w^2=e_w\) and \(e_w\sigma_A(e_w)=0\). Put \(\varepsilon=(1\otimes e)e_w\in\mathcal O_{E,f}\), so \(\varepsilon^2=\varepsilon\) and \(\varepsilon\sigma_A(\varepsilon)=0\). The model supplies vectors \(\beta_1,\ldots,\beta_k\) on \(e_w\) and \(\beta_1^*,\ldots,\beta_k^*\) on \(\sigma_A(e_w)\) in \(L_f\), orthogonal modulo \(g\) to \(W_u\) and satisfying \(h_A(\beta_i^*,\beta_j)\equiv\delta_{ij}e_w\pmod g\) for the adelic Hermitian form \(h_A\). Since \(e_a,f_a\) form hyperbolic pairs and \(h_A(x_u^{-1}z_u,x_u^{-1}z_u)=-T_u\) is a \(v\)-adic unit at every \(v\in\mathcal V_{\mathrm{typ}}\) by \((T_u,g)=1\), orthogonal projection away from \(W_u\) on \(\mathcal V_{\mathrm{typ}}\) replaces \(\varepsilon\beta_j\) and \(\sigma_A(\varepsilon)\beta_l^*\) by vectors \(b_j,c_l\in L_f\) on \(\varepsilon,\sigma_A(\varepsilon)\) that are strictly orthogonal to \(W_u\) and satisfy \(G_{lj}=h_A(c_l,b_j)\equiv\delta_{lj}\varepsilon\pmod g\). On \(\varepsilon\mathcal O_{E,f}\), elements of \(1+g\mathcal O_{E,f}\) have \(g\)-adic unit norms and are invertible, so \(\varepsilon\det G\) has an inverse \(\varepsilon y\) with \(y\in\mathcal O_{E,f}\). Setting \(M=(\varepsilon y)\operatorname{adj}(G)\) and \(b_i^*=\sum_{l=1}^k\sigma_A(M_{il})c_l\in L_f\) gives dual isotropic tuples \(b_1,\ldots,b_k\) on \(\varepsilon\) and \(b_1^*,\ldots,b_k^*\) on \(\sigma_A(\varepsilon)\), orthogonal to \(W_u\), with \(h_A(b_i^*,b_j)=\delta_{ij}\varepsilon\) and modulo-\(g\) pairings \(\delta_{ij}e_w\) against \(\beta_j^*,\beta_i\). Given \(A\in\mathcal I\), lift \(A\) and \(A^{-1}\) to \(M_k(\varepsilon\mathcal O_{E,f})\) and invert their product on \(\varepsilon\mathcal O_{E,f}\) to obtain \(C_0,C'_0\in M_k(\varepsilon\mathcal O_{E,f})\) with \(C'_0C_0=C_0C'_0=\varepsilon I_k\) and reduction \(A\) at every \(v\in\mathcal V_{\mathrm{typ}}\). Set \(C=C_0+(1-\varepsilon)I_k\) and \(C'=C'_0+(1-\varepsilon)I_k\), so \(C'C=CC'=I_k\). Writing \(P_{x,y}(N)v=\sum_{i,j=1}^k x_i N_{ij}h_A(y_j,v)\), the endomorphism
\[
 s'=I_n+P_{b,b^*}(C-I_k)+P_{b^*,b}\bigl(\sigma_A(C')^T-I_k\bigr)
\]
preserves \(h_A\) and \(L_f\), with inverse obtained by interchanging \(C\) and \(C'\); it also fixes \(W_u\) pointwise and has Levi reduction \(A\in\mathcal I\). Thus \(s_F=x_u s' x_u^{-1}\in H_{\mathrm{int}}(u)\) satisfies \(\bar\ell(s_F)=A\).

Choose lifts \(a_1,\ldots,a_{\Gamma_1}\in H_{\mathrm{int}}(u)\) of the \(\Gamma_1\) left cosets in \(\mathcal I/\mathcal U\). Fix \(\Lambda\in\mathcal R_0(u)\). If \(\operatorname{Aut}_{F_u}(\Lambda)\) is infinite, \(w(\Lambda)=0\le|\mathrm{Cl}_u(\Lambda)|/\Gamma_1\). Assume \(\operatorname{Aut}_{F_u}(\Lambda)\) is finite, and let \(h_b=h_{\rho_{\mathrm{base}}(\Lambda)}\) be the base representative of \(\Lambda\), with \(\mathcal L_u(h_b)=\Lambda\). Put \(t_q=h_b a_q\in H_{\mathrm{fin}}(u)\) for \(1\le q\le\Gamma_1\). Because \(c_u(a_q)\in K_{\max}\) preserves \(L_f\), we have \(\mathcal L_u(t_q)=\mathcal L_u(h_b)=\Lambda\). If \(t_p^{-1}t_q=a_p^{-1}a_q\in K_u\), then \(\bar\ell(a_p)^{-1}\bar\ell(a_q)=\bar\ell(a_p^{-1}a_q)\in\mathcal U\), so \(p=q\). Each double coset \(\iota_f(\Gamma_u)t_q K_u\) has representative \(h_\rho\) with \(\rho\in\mathrm{Cl}_u(\Lambda)\); hence at most \(|\mathrm{Cl}_u(\Lambda)|\) distinct double cosets occur as \(q\) varies. In a single double coset \(\iota_f(\Gamma_u)t_{q_0}K_u\), each \(t_q=\eta_q t_{q_0}\kappa_q\) (\(\eta_q\in\iota_f(\Gamma_u)\), \(\kappa_q\in K_u\)) satisfies \(\eta_q\in A_\Lambda\); if \(\eta_q=\eta_{q'}\), then \(t_q^{-1}t_{q'}=\kappa_q^{-1}\kappa_{q'}\in K_u\) and \(q=q'\). Each fibre therefore has cardinality at most \(|A_\Lambda|\le|\operatorname{Aut}_{F_u}(\Lambda)|\), whence \(\Gamma_1\le|\operatorname{Aut}_{F_u}(\Lambda)|\,|\mathrm{Cl}_u(\Lambda)|\) and \(\Gamma_1 w(\Lambda)\le|\mathrm{Cl}_u(\Lambda)|\). Summing over \(\Lambda\in\mathcal R_0(u)\) and multiplying by \(n_\pi\) gives
\[
 \Gamma_1\mu_u n_\pi
 =\Bigl(\sum_{\Lambda\in\mathcal R_0(u)}\Gamma_1 w(\Lambda)\Bigr)n_\pi
 \le\Bigl(\sum_{\Lambda\in\mathcal R_0(u)}|\mathrm{Cl}_u(\Lambda)|\Bigr)n_\pi
 =|\mathrm{Cl}_u|n_\pi=d.
\]
\end{proof}

\subsection*{Proof of \ref{B8.08ob}}\label{prf:B8.08ob}\phantomsection
\begin{proof}
Let \(G_\infty\) be the archimedean unitary group of the model with Haar measure \(\mu_\infty\), let \(V_{\pi_u}\) be the space of \(\pi_u\) with complex inner product \(\langle\cdot,\cdot\rangle\) conjugate-linear in its first variable, and let \(W_u:G_\infty\to V_{\pi_u}\) be the continuous Bessel--Whittaker function of \(u\). Let \(\beta_t\) be the rational Hermitian isometry representing \(t\) (carrying \(F_u\) to \(F'_t\)), let \(\gamma_{t,\infty}\in G_\infty\) be its archimedean component, and let \(\Gamma_{\mathrm{orb}}(t)\le G_\infty\) be the archimedean image of the rational common unipotent lattice stabilizer of \(t\). For \(x\in G_\infty\), let \(f_t(x)\) be the continuous linear operator \(v\mapsto\langle W_u(\gamma_{t,\infty}^{-1}x),v\rangle W_u(x)\) on \(V_{\pi_u}\). By Lemma~\ref{B8.08m} and the definition of the normalized orbital amplitude, the single nonzero row coordinate of \(v_t\) is a scalar multiple of the Hilbert--Schmidt image of the Bochner orbital integral
\[
 \mathcal O_t=\int_{\Omega_t}f_t(x)\,d\mu_\infty(x)
\]
over a chosen measurable fundamental domain \(\Omega_t\subseteq G_\infty\) for left multiplication by \(\Gamma_{\mathrm{orb}}(t)\), with \(\mathcal O_t=0\) whenever no such fundamental domain exists or \(f_t\) is not Bochner integrable on \(\Omega_t\). Assume therefore that \(\Omega_t\) exists and \(f_t\) is integrable on \(\Omega_t\).

As a finite product of real unitary groups, \(G_\infty\) is unimodular. Let \(N_\infty(u)\le G_\infty\) be the closed unipotent radical stabilizing the archimedean image of \(F_u\). On \(N_\infty(u)\), \(W_u\) transforms on the left by a continuous unitary character \(\psi_{\infty,u}:N_\infty(u)\to\C^\times\). Set
\[
 R_t=\{r\in N_\infty(u):\gamma_{t,\infty}^{-1}r\gamma_{t,\infty}\in N_\infty(u)\},
\]
a closed subgroup of \(G_\infty\). On \(R_t\), define the continuous character
\[
 \chi_t(r)=\psi_{\infty,u}(r)\overline{\psi_{\infty,u}(\gamma_{t,\infty}^{-1}r\gamma_{t,\infty})}.
\]
For \(r\in R_t\) and \(x\in G_\infty\), the equivariance \(W_u(rx)=\psi_{\infty,u}(r)W_u(x)\) and \(W_u(\gamma_{t,\infty}^{-1}rx)=\psi_{\infty,u}(\gamma_{t,\infty}^{-1}r\gamma_{t,\infty})W_u(\gamma_{t,\infty}^{-1}x)\) gives \(f_t(rx)=\chi_t(r)f_t(x)\). The rational unipotent elements defining \(\Gamma_{\mathrm{orb}}(t)\) preserve the flag filtrations of both \(F_u\) and \(F'_t\), so \(\Gamma_{\mathrm{orb}}(t)\) is a discrete subgroup of \(R_t\). With \(r=\dim_E X_u\) as in Lemma~\ref{B8.08oa}, filtering the skew-Hermitian intersection Lie algebra \(\mathfrak n(F_u)\cap\mathfrak n(F'_t)\) in \(2r+1\) graded steps and applying rational lattice approximation and Cayley transforms on each step produces a compact subset of \(R_t\) whose left \(\Gamma_{\mathrm{orb}}(t)\)-translates cover \(R_t\), so \(\Gamma_{\mathrm{orb}}(t)\backslash R_t\) is compact. Moreover, for rational elements in \(\Gamma_{\mathrm{orb}}(t)\), the product formula for the adelic additive character and level integrality of the finite-adelic unipotent conjugates give \(\chi_t|_{\Gamma_{\mathrm{orb}}(t)}=1\).

Suppose \(t\) is not relevant. Choose \(Y\in\mathfrak n(F_u)\cap\mathfrak n(F'_t)\) with \(\operatorname{Tr}_{E/F}(\ell_{F_u}(Y))\ne\operatorname{Tr}_{E/F}(\ell_{F'_t}(Y))\). Write \(\ell=\ell_{F_u}(Y)\), \(\ell'=\ell_{F'_t}(Y)\), and \(\delta=(\ell-\ell')+\overline{\ell-\ell'}\in E\). Since \(\delta=0\) would force \(\operatorname{Tr}_{E/F}(\ell-\ell')=0\), we have \(\delta\ne0\), and \(\bar\delta=\delta\). Let \(N=[F:\Q]\) be the number of infinite places of \(F\), let \(g=g_0g_1\), put \(c=g/(2N)\in\Q^\times\), and set \(s=c\delta^{-1}\in E^\times\), so \(\bar s=s\) and \(s\delta=c\). Because \(Y\) is nilpotent and skew-Hermitian and \(\bar s=s\), the truncated exponential \(n=\exp(sY)\) on \(V\) is a rational Hermitian isometry in the unipotent stabilizer of \(F_u\), while \(\beta_t^{-1}n\beta_t=\exp(s\,\beta_t^{-1}Y\beta_t)\) also lies in the unipotent stabilizer of \(F_u\) because \(Y\in\mathfrak n(F'_t)\). Let \(r_0\in G_\infty\) be the archimedean point of \(n\); then \(r_0\in R_t\). In the root pairings defining \(\psi_{\infty,u}(r_0)\) and \(\psi_{\infty,u}(\gamma_{t,\infty}^{-1}r_0\gamma_{t,\infty})\), all quadratic and higher powers of \(Y\) map the flag steps into lower filtration subspaces orthogonal to the dual isotropic vectors, so only the linear term \(sY\) contributes. At each infinite place \(v\) of \(F\), the local additive functional \(\Phi_v:E\to\R\) satisfies \(\Phi_v(q)=q\) for \(q\in\Q\), and the local phase exponent is \((2\pi i/g)\Phi_v(s(\ell+\bar\ell))\) for \(F_u\) and \((2\pi i/g)\Phi_v(s(\ell'+\bar\ell'))\) for \(F'_t\). Their difference at each \(v\) is
\[
 \frac{2\pi i}{g}\Phi_v\bigl(s(\ell+\bar\ell-\ell'-\bar\ell')\bigr)
 =\frac{2\pi i}{g}\Phi_v(s\delta)=\frac{2\pi i}{g}\cdot\frac{g}{2N}=\frac{\pi i}{N}.
\]
Summing over the \(N\) infinite places of \(F\) gives \(\chi_t(r_0)=\exp(\pi i)=-1\ne1\).

Unimodularity of \(G_\infty\) and the triviality \(\chi_t|_{\Gamma_{\mathrm{orb}}(t)}=1\) allow the integral \(\mathcal O_t\) over \(\Omega_t\) to be averaged on the left over the compact quotient \(\Gamma_{\mathrm{orb}}(t)\backslash R_t\) against its right-\(R_t\)-invariant Haar measure. Because \(\chi_t\) is continuous on \(R_t\) and \(\chi_t(r_0)\ne1\), its integral over \(\Gamma_{\mathrm{orb}}(t)\backslash R_t\) vanishes, so \(\mathcal O_t=0\) and \(v_t=0\). Specializing to \(t\in\mathcal T_0\setminus\mathcal T_{\mathrm{rel}}\) gives the displayed assertion.
\end{proof}

\subsection*{Proof of \ref{B8.08oc}}\label{prf:B8.08oc}\phantomsection
\begin{proof}
Use the notation \(H_{\mathrm{fin}}(u)\), \(K_{\max}\), \(K_2\), \(K_u\), \(x_u\), \(\mathrm{Cl}_u\), \(h_\rho\), \(\mathcal R_0(u)\), \(\Lambda_\rho\), \(\rho_{\mathrm{base}}(\Lambda)\), and \(w(\Lambda)\) from Lemma~\ref{B8.08oa} and its proof. Under the orthogonal decomposition \(C_u\simeq\bigoplus_{\rho\in\mathrm{Cl}_u}V_{\pi_u}\), each continuous linear operator \(A\) on \(C_u\) has matrix blocks \(B_{\rho,\rho'}(A)\in\operatorname{End}_\C(V_{\pi_u})\) defined by \(B_{\rho,\rho'}(A)v=(A\delta_{\rho'}v)_\rho\). Taking the coordinate-singleton orthonormal basis built from an orthonormal basis of \(V_{\pi_u}\) and interchanging finite sums gives the block Parseval identity
\[
 \langle A,B\rangle_{\mathrm{HS}}
 =\sum_{\rho,\rho'\in\mathrm{Cl}_u}\bigl\langle B_{\rho,\rho'}(A),B_{\rho,\rho'}(B)\bigr\rangle_{\mathrm{HS}}.
\]
The level character \(\chi_2\) on \(K_2\) factors as \(\chi_2(\kappa)=\chi_N(\kappa)\chi_L(\kappa)\), where \(\chi_N\) is the unipotent root-pairing character modulo \(g\) and \(\chi_L(\kappa)=\psi_{\mathcal I}(\operatorname{levi}(\kappa))\) is the Levi Whittaker phase. If \(\kappa,k_1,k_2\in K_{\max}\) satisfy \(k_1^{-1}f_{0,a}\equiv f_{0,a}\pmod g\) and \(k_2 v_{0,a}\equiv v_{0,a}\pmod g\) for the adelic dual isotropic and next-step flag vectors \(f_{0,a},v_{0,a}\in L_f\) (\(1\le a\le r\)), then expanding \(h_A(f_{0,a},k_1\kappa k_2 v_{0,a})=h_A(k_1^{-1}f_{0,a},\kappa(k_2 v_{0,a}))\) shows that the sum of root pairings differs from that of \(\kappa\) by \(g\) times an element of \(\mathcal O_{E,f}\); dividing by \(g\) and applying the finite-adelic additive character, which is trivial on \(\mathcal O_{E,f}\), gives \(\chi_N(k_1\kappa k_2)=\chi_N(\kappa)\). For \(\rho\in\mathrm{Cl}_u\), write \(\delta(x,y)=(x x_u)^{-1}(y x_u)\) for \(x,y\in H_{\mathrm{fin}}(u)\) and put \(s_u(\rho)=\operatorname{levi}(\delta(h_{\rho_{\mathrm{base}}(\Lambda_\rho)},h_\rho))\in\mathcal I\). For a term \(t\) with lattice pair \(q_t=(\Lambda_t,\Lambda'_t)\in\mathcal R_0(u)^2\) and normalized Hilbert--Schmidt amplitude \(a_t\), the transition factors between \(\rho,\rho'\) and the base representatives of \(\Lambda_t,\Lambda'_t\) satisfy these modulo-\(g\) congruences whenever \(\Lambda_\rho=\Lambda_t\) and \(\Lambda_{\rho'}=\Lambda'_t\). Thus the \(\chi_N\)-factor is independent of \(\rho,\rho'\) in those fibres, while the relative finite-adelic element \(k(\rho,\rho',t)\) lies in \(K_2\) if and only if \((\Lambda_\rho,\Lambda_{\rho'})=q_t\) and \(\operatorname{levi}(k(\rho,\rho',t))=s_u(\rho)^{-1}\iota_t s_u(\rho')\in\mathcal U\). Consequently,
\[
 B_{\rho,\rho'}(O_t)=
 \begin{cases}
 \psi_{\mathcal I}\bigl(s_u(\rho)^{-1}\iota_t s_u(\rho')\bigr)\,a_t,&(\Lambda_\rho,\Lambda_{\rho'})=q_t,\\
 0,&(\Lambda_\rho,\Lambda_{\rho'})\ne q_t,
 \end{cases}
\]
and the \(q_t\)-coordinate of \(v_t\) is \(w_{q_t}a_t\) with \(w_{(\Lambda,\Lambda')}=\sqrt{w(\Lambda)w(\Lambda')}\).

For \(\Lambda\in\mathcal R_0(u)\) with base representative \(h_b=h_{\rho_{\mathrm{base}}(\Lambda)}\), define
\[
 A_u(\Lambda)=\bigl\{\operatorname{levi}(\delta(h_b,\iota_f(\gamma)h_b)):
   \gamma\in\Gamma_u,\ \delta(h_b,\iota_f(\gamma)h_b)\in K_{\max}\bigr\}\subseteq\mathcal I.
\]
We claim that the multiplication map
\[
 \Phi_\Lambda:A_u(\Lambda)\times\mathrm{Cl}_u(\Lambda)\times\mathcal U\longrightarrow\mathcal I,\qquad
 (a,\rho,w)\longmapsto a\,s_u(\rho)\,w,
\]
is bijective and that \(w(\Lambda)=|A_u(\Lambda)|^{-1}\). Suppose \(a_1 s_u(\rho_1)w_1=a_2 s_u(\rho_2)w_2\), with \(a_m=\operatorname{levi}(\delta(h_b,\iota_f(\gamma_m)h_b))\) (\(m=1,2\)). Put \(p_m=\iota_f(\gamma_m)h_{\rho_m}\in H_{\mathrm{fin}}(u)\). Since \(\mathcal L_u(h_{\rho_m})=\Lambda=\mathcal L_u(h_b)\), we have \(\delta(h_b,h_{\rho_m})\in K_{\max}\), and hence \(\delta(h_b,p_m)=\delta(h_b,\iota_f(\gamma_m)h_b)\delta(h_b,h_{\rho_m})\in K_{\max}\). Multiplicativity of \(\operatorname{levi}\) gives \(\operatorname{levi}(\delta(h_b,p_m))=a_m s_u(\rho_m)\), so \(\delta(p_1,p_2)=\delta(h_b,p_1)^{-1}\delta(h_b,p_2)\in K_{\max}\) has Levi image \(w_1 w_2^{-1}\in\mathcal U\). An element of \(K_{\max}\) satisfying the word-level flag congruences whose Levi reduction lies in \(\mathcal U\) belongs to \(K_2\), so \(p_1^{-1}p_2\in K_u\). Thus \(h_{\rho_2}=\iota_f(\gamma_2^{-1}\gamma_1)h_{\rho_1}(p_1^{-1}p_2)\); uniqueness of double-coset representatives in \(\mathrm{Cl}_u\) gives \(\rho_1=\rho_2\), and neatness of \(K_u\) on \(h_{\rho_1}^{-1}\iota_f(\gamma_1^{-1}\gamma_2)h_{\rho_1}=p_1^{-1}p_2\in K_u\) gives \(\gamma_1=\gamma_2\), whence \(a_1=a_2\) and \(w_1=w_2\). Conversely, given \(s\in\mathcal I\), the split-frame lift in the proof of Lemma~\ref{B8.08oa} gives \(s_F\in H_{\mathrm{int}}(u)\) with \(c_u(s_F)\in K_{\max}\) and \(\operatorname{levi}(c_u(s_F))=s\). Put \(h=h_b s_F\in H_{\mathrm{fin}}(u)\), so \(\delta(h_b,h)=c_u(s_F)\in K_{\max}\), \(\delta(h,h_b)=c_u(s_F)^{-1}\in K_{\max}\), and \(\operatorname{levi}(\delta(h_b,h))=s\). Decompose \(h=\iota_f(\gamma)h_\rho\kappa\) (\(\gamma\in\Gamma_u\), \(\rho\in\mathrm{Cl}_u\), \(\kappa\in K_u\)). Since \(\delta(h_b,h)\in K_{\max}\), we have \(\mathcal L_u(h)=\mathcal L_u(h_b)=\Lambda\), which as in Lemma~\ref{B8.08oa} forces \(\Lambda_\rho=\Lambda\) and \(\rho\in\mathrm{Cl}_u(\Lambda)\). Then \(\delta(h_b,\iota_f(\gamma)h_b)=\delta(h_b,h)\,\delta(1,\kappa)^{-1}\delta(h_\rho,h_b)\in K_{\max}\) defines \(a\in A_u(\Lambda)\), while \(w=\operatorname{levi}(\delta(1,\kappa))\in\mathcal U\) satisfies \(a\,s_u(\rho)\,w=\operatorname{levi}(\delta(h_b,h))=s\). Thus \(\Phi_\Lambda\) is bijective and \(|A_u(\Lambda)|\ge1\). Set \(\Gamma_u(\Lambda)=\{\gamma\in\Gamma_u:\delta(h_b,\iota_f(\gamma)h_b)\in K_{\max}\}=\{\gamma\in\Gamma_u:g_\gamma(\Lambda)=\Lambda\}\). The injectivity argument above with \(\rho_1=\rho_2=\rho_{\mathrm{base}}(\Lambda)\) shows via neatness of \(K_u\) that \(\gamma\mapsto\operatorname{levi}(\delta(h_b,\iota_f(\gamma)h_b))\) is a bijection \(\Gamma_u(\Lambda)\xrightarrow{\sim}A_u(\Lambda)\). Moreover, the map \(\gamma\mapsto\varphi(\iota_f(\gamma))\) from \(\Gamma_u(\Lambda)\) to \(\operatorname{Aut}_{F_u}(\Lambda)\) is injective by the proof of Lemma~\ref{B8.08oa}, and it is surjective: the standard lattice \(L_f\) splits orthogonally across \((X_u\oplus Y_u)\oplus(X_u\oplus Y_u)^\perp\) (where \(Y_u=\operatorname{span}(f_1,\ldots,f_r)\)), and \(h_b\in H_{\mathrm{fin}}(u)\) fixes \(X_u\oplus Y_u\) while \(x_u\in K_{\max}\) preserves \(L_f\), so \(\Lambda=\mathcal L_u(h_b)\) splits orthogonally as \((\Lambda\cap(X_u\oplus Y_u))\oplus(\Lambda\cap(X_u\oplus Y_u)^\perp)\); hence any \(g\in\operatorname{Stab}(F_u)\) preserving \(\Lambda\) (which fixes \(e_1,\ldots,e_r\) and \(z_u\bmod X_u\)) has a block-diagonal lift \(\tilde g\in\Gamma_u(\Lambda)\) fixing \(e_a,f_a,z_u\) and inducing the same automorphism on \(X_u^\perp/X_u\). Therefore \(\operatorname{Aut}_{F_u}(\Lambda)\) is finite of cardinality \(|A_u(\Lambda)|\), and \(w(\Lambda)=|A_u(\Lambda)|^{-1}>0\).

Fix \(S\in\{E_{\mathrm{gen}},E_{\mathrm{ng}},\mathcal T_0\}\). Because \(\sum_t\|v_t\|<\infty\) and \(w_q>0\) for all \(q\in\mathcal R_0(u)^2\), the subseries
\[
 Z_q(i)=\sum_{\substack{t\in S\\q_t=q,\ \iota_t=i}}a_t\qquad\bigl(q\in\mathcal R_0(u)^2,\ i\in\mathcal I\bigr)
\]
converge absolutely in Hilbert--Schmidt norm. Put \(G_q(i,j)=\langle Z_q(i),Z_q(j)\rangle_{\mathrm{HS}}\). Both the nonzero-Gram condition and the generic condition are invariant under rational Hermitian isometries of \(V\), and the identity cell is preserved by left and right multiplication by elements of \(\Gamma_u\le\operatorname{Stab}(F_u)\), so each of \(E_{\mathrm{gen}},E_{\mathrm{ng}},\mathcal T_0\) is preserved by rational row-stabilizer transport on either side. By the \(\chi_N\)-invariance above and independence of the normalized orbital amplitudes from choices of representatives and fundamental domains, left action by \(b\in A_u(\Lambda)\) at \(q=(\Lambda,\Lambda')\) gives a bijection \(\varphi_L\) of \(\{t\in S:q_t=q\}\) onto itself and a unitary operator \(U_L\) on Hilbert--Schmidt space such that \(\iota_{\varphi_L(t)}=b\,\iota_t\) and \(a_{\varphi_L(t)}=U_L(a_t)\). Right action by \(b\in A_u(\Lambda')\) similarly gives a bijection \(\varphi_R\) of \(\{t\in S:q_t=q\}\) onto itself and a unitary operator \(U_R\) such that \(\iota_{\varphi_R(t)}=\iota_t b^{-1}\) and \(a_{\varphi_R(t)}=U_R(a_t)\). Reindexing \(Z_q\) along \(\varphi_L^{-1}\) and \(\varphi_R^{-1}\) gives \(Z_q(bi)=U_L(Z_q(i))\) and \(Z_q(ib^{-1})=U_R(Z_q(i))\), and unitarity of \(U_L,U_R\) yields
\[
 G_q(bi,bj)=G_q(i,j)\quad\bigl(b\in A_u(\Lambda)\bigr),\qquad
 G_q(ib^{-1},jb^{-1})=G_q(i,j)\quad\bigl(b\in A_u(\Lambda')\bigr).
\]

Since \(\sum_{t\in S}\|O_t\|_{\mathrm{HS}}<\infty\) and \(B_{\rho,\rho'}\) is continuous and linear,
\[
 B_{\rho,\rho'}\bigl(\operatorname{part}(S)\bigr)
 =\sum_{i\in\mathcal I}\psi_{\mathcal I}\bigl(s_u(\rho)^{-1}i\,s_u(\rho')\bigr)\,Z_{(\Lambda_\rho,\Lambda_{\rho'})}(i).
\]
For \(q=(\Lambda,\Lambda')\in\mathcal R_0(u)^2\) and \(x,y\in\mathcal I\), put
\[
 F_q(x,y)=\sum_{i,j\in\mathcal I}\overline{\psi_{\mathcal I}(x^{-1}iy)}\,\psi_{\mathcal I}(x^{-1}jy)\,G_q(i,j),
\]
so \(\|\operatorname{part}(S)\|_{\mathrm{HS}}^2=\Re\sum_{\rho,\rho'\in\mathrm{Cl}_u}F_{(\Lambda_\rho,\Lambda_{\rho'})}(s_u(\rho),s_u(\rho'))\). For \(b\in A_u(\Lambda)\) and \(w\in\mathcal U\), the relation \(\psi_{\mathcal I}(w^{-1}z)=\psi_{\mathcal I}(w^{-1})\psi_{\mathcal I}(z)\) with \(|\psi_{\mathcal I}(w^{-1})|=1\), followed by the bijection \(i\mapsto bi\) and \(G_q(bi,bj)=G_q(i,j)\), gives \(F_q(bxw,y)=F_q(x,y)\). Similarly, for \(b'\in A_u(\Lambda')\) and \(w'\in\mathcal U\), the relation \(\psi_{\mathcal I}(zw')=\psi_{\mathcal I}(z)\psi_{\mathcal I}(w')\) with \(|\psi_{\mathcal I}(w')|=1\) and the bijection \(i\mapsto i(b')^{-1}\) with \(G_q(i(b')^{-1},j(b')^{-1})=G_q(i,j)\) give \(F_q(x,b'yw')=F_q(x,y)\). By the bijections \(\Phi_\Lambda\) and \(\Phi_{\Lambda'}\), each fibre over \((\rho,\rho')\in\mathrm{Cl}_u(\Lambda)\times\mathrm{Cl}_u(\Lambda')\) has size \(|A_u(\Lambda)|\,|\mathcal U|\cdot|A_u(\Lambda')|\,|\mathcal U|\), whence
\[
 |A_u(\Lambda)|\,|A_u(\Lambda')|\,|\mathcal U|^2
 \sum_{\substack{\rho\in\mathrm{Cl}_u(\Lambda)\\\rho'\in\mathrm{Cl}_u(\Lambda')}}F_q\bigl(s_u(\rho),s_u(\rho')\bigr)
 =\sum_{x,y\in\mathcal I}F_q(x,y).
\]
For fixed \(i,j\in\mathcal I\), \(\psi_{\mathcal I}(x^{-1}jy)\) vanishes unless \(u_0=x^{-1}jy\in\mathcal U\), that is, \(x=jyu_0^{-1}\) with \(u_0\in\mathcal U\). Substituting \(x=jyu_0^{-1}\) gives \(x^{-1}iy=u_0 y^{-1}j^{-1}iy\) and
\[
 \sum_{x,y\in\mathcal I}\overline{\psi_{\mathcal I}(x^{-1}iy)}\,\psi_{\mathcal I}(x^{-1}jy)
 =\sum_{y\in\mathcal I}\sum_{u_0\in\mathcal U}\overline{\psi_{\mathcal I}(u_0)\psi_{\mathcal I}(y^{-1}j^{-1}iy)}\,\psi_{\mathcal I}(u_0)
 =|\mathcal U|\sum_{y\in\mathcal I}\overline{\psi_{\mathcal I}(y^{-1}j^{-1}iy)}
 =|\mathcal U|^2\mathcal G_{ij}.
\]
Dividing by \(|A_u(\Lambda)|\,|A_u(\Lambda')|\,|\mathcal U|^2\) and using \(w_q^2=w(\Lambda)w(\Lambda')=(|A_u(\Lambda)|\,|A_u(\Lambda')|)^{-1}\) gives
\[
 \sum_{\substack{\rho\in\mathrm{Cl}_u(\Lambda)\\\rho'\in\mathrm{Cl}_u(\Lambda')}}F_q\bigl(s_u(\rho),s_u(\rho')\bigr)
 =w_q^2\sum_{i,j\in\mathcal I}\mathcal G_{ij}G_q(i,j).
\]
Projecting the absolutely convergent sum \(w_i(S)\) to coordinate \(q\) gives \(w_i(S)_q=w_q Z_q(i)\), so \(\langle w_i(S),w_j(S)\rangle=\sum_{q\in\mathcal R_0(u)^2}w_q^2 G_q(i,j)\). Summing the previous identity over \(q\in\mathcal R_0(u)^2\) and interchanging the finite sums over \(q\) and \(i,j\) gives \(\sum_{\rho,\rho'\in\mathrm{Cl}_u}F_{(\Lambda_\rho,\Lambda_{\rho'})}(s_u(\rho),s_u(\rho'))=\sum_{i,j\in\mathcal I}\mathcal G_{ij}\langle w_i(S),w_j(S)\rangle\); taking real parts yields
\[
 \|\operatorname{part}(S)\|_{\mathrm{HS}}^2
 =\Re\sum_{i,j\in\mathcal I}\mathcal G_{ij}\bigl\langle w_i(S),w_j(S)\bigr\rangle.
\]
Taking \(S=\mathcal T_0\) gives the final assertion.
\end{proof}

\subsection*{Proof of \ref{B8.08o}}\label{prf:B8.08o}\phantomsection
\begin{proof}
Put
\[
 C_c=21+4|C_H|,\qquad K_*=1+C+C^2+C_h^2+C_c^2,
\]
and define
\[
 M_*=\max\{16,K_*^8\},\qquad
 \eta_*=\max\{16K_*,M_*\log M_*\},\qquad
 C_{\ell,\mathrm{pair}}=\max\{16,\eta_*\}.
\]
These choices precede the choice of the model and the word. We prove the assertion for this threshold.

The normal form established in the proof of Proposition~\ref{B8.02b} gives \(1\le j\le r\) and \(\delta_j\ne0\) for every type \((j,\delta)\) represented in \(\mathcal T_{\mathrm{rel}}\). If \(r=0\), this set is empty, and \eqref{eq:p8c1-zero-normalization} makes every summand zero. Hence assume \(r\ge1\).

We first sum the minor labels. Put
\[
 q_l=\sum_{a=1}^l(r+1-a),\qquad
 H_1=b^{C_Hm},\qquad
 \mathfrak d=\mathfrak D_{F/\Q}\mathcal O_E.
\]
The integrality of the normalized nonzero minor ideals used in the proof of Proposition~\ref{B8.02b} gives
\((\delta_l)\mathfrak d^{-2q_l}\subseteq\mathcal O_E\) for a nonzero minor with \(1\le l\le j\). Since \(q_l\ge0\) and \(\mathfrak d\) is integral, it follows that
\[
 (\delta_l)\subseteq\mathfrak d^{2q_l}\subseteq\mathcal O_E.
\]
Thus every coordinate \(\delta_l\) is an algebraic integer. The coordinates \(\delta_1,\ldots,\delta_j\) determine the label, since \(\delta_0=1\) and \(\delta_l=0\) for \(l>j\).

Let \(\tau\) range over the fixed CM type, which has \(f\) embeddings. For distinct \(x,y\in\mathcal O_E\),
\[
 \prod_\tau|\tau(x-y)|^2=|N_{E/\Q}(x-y)|\ge1.
\]
The products of open complex disks of radius \(1/2\) about the embedded algebraic integers are therefore disjoint. If \(R_\tau\ge1\), the disks about points satisfying \(|\tau(x)|\le R_\tau\) lie in the product of disks of radii \(R_\tau+1/2\). Comparing volumes for each finite subset gives
\begin{equation}\label{eq:p8c1-zero-label-box}
 \#\{x\in\mathcal O_E:|\tau(x)|\le R_\tau\text{ for all }\tau\}
 \le\prod_\tau(2R_\tau+1)^2
 \le9^f\prod_\tau R_\tau^2.
\end{equation}
The same finite-subset bound proves that the set on the left is finite.

Fix \(R\ge1\). Partition \(\mathcal O_E\) by the integers \(n_\tau\ge0\) satisfying
\[
 2^{n_\tau}\le\max\{1,|\tau(x)|/R\}<2^{n_\tau+1}.
\]
By \eqref{eq:p8c1-zero-label-box}, a shell contains at most
\(36^fR^{2f}2^{2\sum_\tau n_\tau}\) points. Its weight
\(\prod_\tau\max\{1,|\tau(x)|/R\}^{-3}\) is at most
\(2^{-3\sum_\tau n_\tau}\). Summing the resulting geometric series gives
\begin{equation}\label{eq:p8c1-zero-label-tail}
 \begin{aligned}
 \sum_{x\in\mathcal O_E}\prod_\tau
                 \max\{1,|\tau(x)|/R\}^{-3}
 &\le36^fR^{2f}\sum_{(n_\tau)\in\N^f}2^{-\sum_\tau n_\tau}\\
 &=72^fR^{2f}\le(64R)^{2f}.
 \end{aligned}
\end{equation}

For \(1\le l\le j\), set
\[
 R_l=\max\{1,\rho^{2q_l}(2H_1)^{2r}\}.
\]
Let \(s_\tau=\max\bigl(1,\max_{\delta_l\ne0,\,1\le l\le j}d_{\tau l}/(2H_1)^{2r}\bigr)\) be the factor of \(S_\delta\) at \(\tau\) in Proposition~\ref{B8.02b}. For each \(l\le j\),
\[
 \max\{1,|\tau(\delta_l)|/R_l\}\le s_\tau.
\]
Using \(j\le r\) and \(s_\tau\ge1\), we obtain
\[
 S_\delta^{-3r}\le
 \prod_{l=1}^j\prod_\tau\max\{1,|\tau(\delta_l)|/R_l\}^{-3}.
\]
Let \(T_j\) be any finite set of represented labels of rank \(j\). Enlarge its set of coordinate tuples to the product of its coordinate projections and apply \eqref{eq:p8c1-zero-label-tail} to each coordinate. All summands are nonnegative, so
\begin{equation}\label{eq:p8c1-zero-label-product}
 \sum_{\delta\in T_j}S_\delta^{-3r}
 \le\prod_{l=1}^j(64R_l)^{2f}
 \le2^j\prod_{l=1}^j(64R_l)^{2f}.
\end{equation}

Put
\[
 B_0=\rho^{2r^2}\max\{1,2H_1\}^{2r},\qquad
 c=17+4/C_\ell+4|C_H|.
\]
Since \(q_l\le r^2\), \(\rho\ge2\), and \(H_1>0\), we have \(R_l\le B_0\). Also
\(\max\{1,2H_1\}\le2b^{|C_H|m}\). Hence \(j\le r\), \(f\ge1\), \(L>1\), and \(\log\rho\le L/C_\ell\) give
\[
 \begin{aligned}
 \log\bigl(2^j(64B_0)^{2fj}\bigr)
 &\le frL\left(13+4r+\frac{4r^2}{C_\ell}+4|C_H|rm\right)\\
 &\le c\,m^2frL
 \le c^2m^2frL
 \le C_c^2m^2frL.
 \end{aligned}
\]
For the second inequality use \(r\le m\) and \(13+4r\le17m^2\); for the last two use \(1\le c\le C_c\), which follows from \(C_\ell\ge1\). Exponentiation in \eqref{eq:p8c1-zero-label-product} proves the label bound
\begin{equation}\label{eq:p8c1-zero-label-sum}
 \sum_{\delta\in T_j}S_\delta^{-3r}\le H(C_c)^{C_cr}.
\end{equation}

We now sum the fixed-type estimates. Since \(\delta_j\ne0\) and \(\operatorname{disc}E=(\operatorname{disc}F)^2\) (from Definition~\ref{B8.02c}), we have \(N(\mathfrak d)=\operatorname{rd}(E)^{2f}\), so the product formula for the normalized terminal minor gives
\[
 0<D_\delta=\prod_\tau d_{\tau j}^{2}
 \le(2H_1)^{4rf}S_\delta^2\le H(C)S_\delta^2.
\]
The last inequality follows from \(2H_1\le b^{C_Hm+1}\) and
\[
 4rf(C_Hm+1)\le2(C_H+1)m^2f\le Cm^2f.
\]
Define
\[
 Q_j=(C/k)^{j(k-1-j)f}H(C)^{Cr^4}H(C_h)^{C_hr^4}.
\]
The factors \(A(j,\delta)\) and \(M(j,\delta)\) of Proposition~\ref{B8.02b} satisfy the exact identity
\[
 (AM)(j,\delta)=Q_jD_\delta^{(j+1)/2}S_\delta^{-4k},
\]
because \(C_s=C_p+4\). Since \(k\ge2r+5\), \(j\le r\), and \(S_\delta\ge1\),
\[
 \begin{aligned}
 (AM)(j,\delta)
 &\le Q_jH(C)^{(j+1)/2}S_\delta^{j+1-4k}\\
 &\le Q_jH(C)^{(j+1)/2}S_\delta^{-3r}.
 \end{aligned}
\]
Consequently, for every finite set \(T_j\) as above,
\[
 \sum_{\delta\in T_j}(AM)(j,\delta)\le A_j,
 \qquad
 A_j=Q_jH(C)^{(j+1)/2}H(C_c)^{C_cr}.
\]

Let \(U\) be a finite subset of \(\mathcal Q_0\). By \eqref{eq:p8c1-zero-normalization}, its pairs with an irrelevant entry contribute zero. Partition the remaining pairs by the ranks and labels of both entries. For each \(j\), let \(T_j\) and \(T'_j\) be the labels of rank \(j\) occurring as first and second entries. The fixed-type estimate applied to every fibre, followed by enlargement to the product of these two finite sets of types, gives
\[
 \begin{aligned}
 \sum_{(t,t')\in U}\|v_t\|\|v_{t'}\|
 &\le n_\pi\mu_u
 \left(\sum_{j=1}^r\sum_{\delta\in T_j}(AM)(j,\delta)\right)
 \left(\sum_{j=1}^r\sum_{\delta'\in T'_j}(AM)(j,\delta')\right)\\
 &\le n_\pi\mu_u\left(\sum_{j=1}^r A_j\right)^2.
 \end{aligned}
\]
Taking the supremum over finite \(U\) proves summability and
\begin{equation}\label{eq:p8c1-zero-unnormalized-pairs}
 \sum_{(t,t')\in\mathcal Q_0}\|v_t\|\|v_{t'}\|
 \le n_\pi\mu_u\left(\sum_{j=1}^r A_j\right)^2.
\end{equation}

The logarithm of the height factor in \(A_j\) is
\[
 m^2fL\left((C^2+C_h^2)r^4+C\frac{j+1}{2}+C_c^2r\right)
 \le K_*m^6fL.
\]
Here \(r^4\le m^4\), \((j+1)/2\le m^4\), and \(r\le m^4\). As \(C\le K_*\) and \(j(k-1-j)f\ge0\), it follows that
\begin{equation}\label{eq:p8c1-zero-rank-majorant}
 A_j\le\exp(K_*m^6fL)(K_*/k)^{j(k-1-j)f}.
\end{equation}
Write
\[
 S_*=\sum_{j=1}^r\exp(K_*m^6fL)(K_*/k)^{j(k-1-j)f}.
\]
The threshold and the source sizes give
\[
 \eta_*m^6L\le C_\ell m^6L\le k\log k.
\]
Since \(m^6L\ge1\), we have \(k\log k\ge M_*\log M_*\). The function \(x\log x\) is increasing for \(x\ge1\), so \(k\ge M_*\). Therefore
\[
 k\ge16,\qquad \log K_*\le\tfrac18\log k,\qquad
 K_*/k\le1,\qquad k-2\ge\tfrac78k.
\]
For \(1\le j\le r\),
\[
 j(k-1-j)-(k-2)=(j-1)(k-j-2)\ge0,
\]
and hence
\[
 S_*\le r\exp(K_*m^6fL)(K_*/k)^{(k-2)f}.
\]
Set \(B=kf\log k\). Since \(r\le k\), \(f\ge1\), and \(\eta_*\ge16K_*\),
\[
 \begin{aligned}
 \log(2S_*^2)
 &\le\log2+2\log r+2K_*m^6fL
       -2(k-2)f(\log k-\log K_*)\\
 &\le3\log k+\tfrac18 B-\tfrac{49}{32}B\\
 &\le\left(\tfrac3{16}+\tfrac18-\tfrac{49}{32}\right)B
 =-\tfrac{39}{32}B\le-B.
 \end{aligned}
\]
Thus
\begin{equation}\label{eq:p8c1-zero-numerical-saving}
 2\left(\sum_{j=1}^r A_j\right)^2\le2S_*^2\le e^{-kf\log k}.
\end{equation}

Finally, \(g_i\ge(kb)^8\), \(b\ge2\), and \(f\le k\) imply \(g_i\ge12f\). Hence
\[
 0<P_g\le\exp(3f/g_0+3f/g_1)\le e^{1/2}\le2.
\]
Multiplying \eqref{eq:p8c1-zero-unnormalized-pairs} by \(P_g\Gamma_1\), and then using \eqref{eq:p8c1-zero-normalization} and \eqref{eq:p8c1-zero-numerical-saving}, gives
\[
 \begin{aligned}
 P_g\Gamma_1\sum_{\mathcal Q_0}\|v_t\|\|v_{t'}\|
 &\le P_g\Gamma_1 n_\pi\mu_u\left(\sum_{j=1}^r A_j\right)^2\\
 &\le2d\left(\sum_{j=1}^r A_j\right)^2
 \le d e^{-kf\log k}.
 \end{aligned}
\]
This proves \eqref{eq:p8c1-zero-pair-amplitudes}.
\end{proof}

\subsection*{Proof of \ref{B8.08p}}\label{prf:B8.08p}\phantomsection
\begin{proof}
The operator subseries is absolutely convergent by comparison with the full series. The vector subseries defining the \(w_i\) are also absolutely convergent. Moreover, Cauchy--Schwarz gives
\[
 \sum_{t,t'}|\langle v_t,v_{t'}\rangle|
 \le\left(\sum_t\|v_t\|\right)^2<\infty.
\]
Thus the inner-product series can be regrouped by finite-type indices and row labels.

By Lemma~\ref{B8.10f}, the finite-type Gram matrix \(\mathcal G\) satisfies \(\mathcal G\preceq P_g\Gamma_1 I\) in positive semidefinite order. Expanding the finite family \((w_i)_{i\in\mathcal I}\) in an orthonormal basis of its finite-dimensional span and applying this matrix inequality to each coordinate vector in \eqref{eq:p8c1-zero-expansion} (equivalently, applying the second assertion of Lemma~\ref{B8.10f} to the cell \(\mathcal T_0\)) gives
\[
 \|Z\|_{\mathrm{HS}}^2\le P_g\Gamma_1\sum_i\|w_i\|^2.
\]

Expand each squared norm into its absolutely convergent double series. At a fixed finite-type index, distinct complete row labels have distinct lattice-pair coordinates \((\Lambda_t,\Lambda'_t)\). Their orbital vectors are orthogonal by Lemma~\ref{B8.08m}. Consequently
\[
 \sum_i\|w_i\|^2
 =\sum_{(t,t')\in\mathcal Q_0}\Re\langle v_t,v_{t'}\rangle.
\]
Since \(P_g\Gamma_1\ge0\), another application of Cauchy--Schwarz gives
\[
 \begin{aligned}
 \|Z\|_{\mathrm{HS}}^2
 &\le P_g\Gamma_1\sum_{(t,t')\in\mathcal Q_0}
                             \Re\langle v_t,v_{t'}\rangle\\
 &\le P_g\Gamma_1\sum_{(t,t')\in\mathcal Q_0}
                             |\langle v_t,v_{t'}\rangle|\\
 &\le P_g\Gamma_1\sum_{(t,t')\in\mathcal Q_0}\|v_t\|\|v_{t'}\|
 \le d e^{-kf\log k}.
 \end{aligned}
\]
\end{proof}

\subsection*{Proof of \ref{B8.08q}}\label{prf:B8.08q}\phantomsection
\begin{proof}
Apply Proposition~\ref{B8.02b} to the fixed constants to obtain \(C_H,C_h,C_p\ge0\). First choose
\[
 C_s=C_p+4.
\]
Let \(C_{\min}\) be the amplitude threshold supplied for this value of \(C_s\), and choose
\[
 C=\max\{C_{\min},2(C_H+1)\}.
\]
At this value of \(C\), Proposition~\ref{B8.02b} supplies thresholds \(C'_{\mathrm{type}}\) and \(C_{\ell,\mathrm{type}}\). Proposition~\ref{B8.08o} supplies \(C_{\ell,\mathrm{pair}}\) for the same constants. Set
\begin{equation}\label{eq:p8c1-zero-thresholds}
 C'_{\mathrm{req}}=\max\{C'_{\mathrm{type}},1\},\qquad
 C_{\ell,\mathrm{req}}
 =\max\{C_{\ell,\mathrm{type}},C_{\ell,\mathrm{pair}},16\}.
\end{equation}

Fix parameters above these thresholds and a model as in the statement. Write \(r=(m-2)/2\), \(L=\log b\), and \(\rho=\operatorname{rd}(E)\). By \eqref{eq:p8c1-zero-thresholds} and Definition~\ref{B8.02c}, \(m\ge C_\ell\ge16\), so \(r=(m-2)/2\ge1\) and \(k\ge f\ge m^4\ge16^4>e\), whence \(\log k\ge1\). Since \(C'\ge1\), \(C_\ell\le m\), and \(m^2\le m^4\le f\), the source inequalities of Definition~\ref{B8.02c} give
\[
 C_\ell m^6L\le m^7L\le m^5fL\le k\le k\log k.
\]
Combining \(\log f\le\sqrt r\) from Definition~\ref{B8.02c} with \(r\ge1\) gives
\[
 \log f\le\sqrt r\le r^2.
\]
Multiplying \(\log\rho\le L/C_\ell\) by \(C_\ell>0\) gives \(C_\ell\log\rho\le L\). The adjacent Harish--Chandra gaps in Definition~\ref{B8.02c} satisfy \(\lambda_i-\lambda_{i+1}>C_\ell km^4+nm^{10}>C_\ell km^3\); summing adjacent gaps gives the required bound between every two distinct parameters. The auxiliary-prime bounds and the upper bound \(T_w\le b^{C_Tm}\) for every word \(w\) are source inequalities. The distinct level and auxiliary primes give \((T_w,g_0g_1)=1\) for every word \(w\). The model has the stated unramified CM property. Thus the parameter conditions of Proposition~\ref{B8.02b} hold.

By \eqref{eq:p8c1-zero-thresholds}, the fixed-type estimate applies at \(C_s=C_p+4\) and the chosen \(C\). Lemmas~\ref{B8.08oa} and~\ref{B8.08ob} give \eqref{eq:p8c1-zero-normalization}, so Proposition~\ref{B8.08o} gives the summable pair-amplitude family and \eqref{eq:p8c1-zero-pair-amplitudes}. Under the assumed absolute convergence, Proposition~\ref{B8.08oc} gives \eqref{eq:p8c1-zero-expansion}. The tied local data have nontrivial selected additive characters, the primes \(g_0,g_1\) split completely in \(E\), and \(g_i\ge(kb)^8\ge3\); hence the selected local hypotheses of Theorem~\ref{B8.10} used in Lemma~\ref{B8.10f} hold. Since \(\Gamma_1\ge0\) by Lemma~\ref{B8.08oa}, Lemma~\ref{B8.08p} gives both conclusions of \eqref{eq:p8c1-zero-gram-bound}. Every constant was chosen before the datum, the tied model, and the word, which proves the stated uniformity.
\end{proof}

\subsection*{Proof of \ref{B8.08h}}\label{prf:B8.08h}\phantomsection
\begin{proof}
The level window and \((1+1/(fk))^{fk}\le\exp(1)\) give
\[
   0<d\le d^\star T^{fk}\le \exp(1)D^\star\Gamma^{-2}
                                \le D^\star e^{-2a}.
\]
For the last inequality use \(2(C_\Gamma-C_1)Y\ge1\). Cauchy--Schwarz gives
\(\|X+Z\|^2\le\frac32\|X\|^2+3\|Z\|^2\) in Hilbert--Schmidt norm. Since \(B/2\ge5\), we have \(e^{-B}\le(2/5)e^{-B/2}\). Substituting the two equal-word bounds gives
\[
 \|K-I\|_{\rm HS}^2
 \le 4(1+e^a)\frac{d^2}{D^\star}
                     +(4e^{-B/2}+2e^{-B})d .
\]
Division by \(d\) and the bound \(d/D^\star\le e^{-2a}\) give the stated \(e\).

On a free orbit each character of \((\F_T^\times)^f\) occurs once. Summation in coordinate \(i\) acts by \(T-1\) when the \(i\)-th character is trivial and by zero otherwise. Exactly \((T-2)^f\) of the \((T-1)^f\) characters survive in the common kernel. Summing over orbits gives
\[
       \dim U=d\left(1-\frac1{T-1}\right)^f
                  \ge d\left(1-\frac f{T-1}\right).
\]
Since \(T>1\), its congruence gives \(T-1\ge(g_0g_1)^2\ge4f\). Thus \(\delta\le1/4\). Both \(a\) and \(B/2\) are at least \(5\), so
\(2e\le28e^{-5}<3/4\), as required.
\end{proof}

\subsection*{Proof of \ref{B8.08f}}\label{prf:B8.08f}\phantomsection
\begin{proof}
Extend an orthonormal basis of \(\ker P_u\) to one of \(C_u\). On each of the first \(\dim\ker P_u\) basis vectors \(K_u-1\) acts as \(-1\). Hence
\[
       \dim\ker P_u\le\|K_u-1\|_{\rm HS}^2
       \le e d_0<(1-\delta)d_0\le\dim U_u^{\rm ram}.
\]
Choose \(\phi\in U_u^{\rm ram}\setminus\ker P_u\). At the linked realization let \(C\) represent this column in the model. The isometry identifying the unrestricted targets takes \(P_u\phi\) to its model period \(p\), so \(p\ne0\).

Write \(J_z\) for evaluation of a full reference multiplicity map at the prescribed unit lowest-type tensor, followed by the prescribed finite translation, the central translation \(z\), and complex conjugation in the row-function convention. The full-level evaluation in the proof of Lemma~\ref{B8.08d} is onto \(H\). The central translation acts bijectively on \(H\): it commutes with the level projector and its inverse is translation by \(z^{-1}\). Thus \(T_0=J_z^{-1}p\) is a full reference multiplicity map. Its model pairing with \(C\) is
\[
                        \langle J_zT_0,p\rangle=\|p\|^2>0 .
\]
\end{proof}

\subsection*{Proof of \ref{B8.08i}}\label{prf:B8.08i}\phantomsection
\begin{proof}
Write \(I_t=\prod_{j<m}I_{t,j}\) and
\[
        W_{t,j}=V_{t,j}K_{t,j}K_{t,j+1}\quad(j\ \text{even}),\qquad
        H_t=P_t\prod_{j\ {\rm even},\,j<m}W_{t,j}.
\]
Thus \(c_t=s_uH_t/I_t\). We describe how the descending chain acts on a linear relation. Allow a finite list of top rows, each with its own tuple and column. At even height \(j<m\), denote by \(D_{i,j}(y,v)\) the Bessel period for entry \(i\), with column evaluated in frame \(y\) and row evaluated at smooth vector \(v\). The row is descended from the top to \(j\). The column has been carried up from zero by successive pairs of kernels, with the prescribed column operator before the first kernel of each pair. At height \(m\) write \(D_{i,m}\) for the scalar terminal period of this column and the given top row.

At an interior double step the following implication holds for arbitrary complex coefficients \(\beta_i\), where \(t_i\) is the tuple of entry \(i\):
\[
 \begin{split}
 &\sum_i\beta_i W_{t_i,j} D_{i,j+2}(y',v')=0
                      \quad\text{for all }y',v'\\
 &\hspace{12mm}\Longrightarrow\
 \sum_i\beta_i I_{t_i,j}I_{t_i,j+1}D_{i,j}(y,v)=0
                      \quad\text{for all }y,v.
 \end{split}
\]
Here is the calculation explaining the common tests in this implication. Evaluate the lower row on a vector obtained by two polynomial theta tests applied to \(v'\). The two defining row-lift integrals absorb \(I_{t_i,j}I_{t_i,j+1}\). Unfold the resulting Bessel period first along the intermediate Fourier--Jacobi subgroup and then along the next Bessel subgroup. At fixed real translates, the finite average of the resulting expression has the form
\[
          \sum_{y'} a(y')\,W_{t_i,j}
                            D_{i,j+2}(y',v''),
\]
where \(v''\) is the real translate of \(v'\). To see the independence of \(a\), first do one definite class sum. On the fixed finite-dimensional polynomial coefficient space, the compact average with the prescribed polynomial test is a linear combination of frame evaluations: the evaluations on invertible frames span the dual, since a polynomial vanishing there vanishes identically. Choose this combination before inserting the class form. In the class sum its coefficients are multiplied by the source inverse determinant power and the target determinant power of the chosen frames. Both powers are fixed by the weights. Applying this calculation twice and multiplying the two finite matrices of coefficients gives \(a\). The finite chart coefficients of entry \(i\) are inside its two column kernels. Their scalar normalizations are \(K_{t_i,j}K_{t_i,j+1}\); the change of the test vector contributes \(V_{t_i,j}\). Real translation acts on the argument of each intertwiner by the same representation operator, giving the same \(v''\).

Sum over \(i\) before taking the quotient integrals. For the finite integration take the intersection of the finitely many level stabilizers and the union of the finite support sets. Subdivision into cosets computes each original integral on this common level. The weighted integrand after finite averaging vanishes by the upper relation for each \(y',v''\). The quotient integrations preserve the relation. The interchanges use cusp decay on the row and its intermediate lift, with the polynomial growth of the fixed theta kernels as in the proof of Lemma~\ref{B8.15f}. The polynomial theta tests span densely in the smooth row models: at each real place their translates generate the prescribed irreducible Harish--Chandra module, and its compact-finite vectors are dense in its smooth globalization. The period against fixed column data is continuous in each smooth argument by the same decay estimate. This extends the relation to every \(v\).
\leanref{GlobalTheory.\allowbreak{}SeeSaw.\allowbreak{}kerX\_\allowbreak{}of\_\allowbreak{}clsPinD}{Development.\allowbreak{}AbelianVarietyTheory.\allowbreak{}Proofs.\allowbreak{}DecoratedLatticeProofs.\allowbreak{}SeeSawKernelLiftCombinationOfDefiniteClassData}{14--49}
\leanref{GlobalTheory.\allowbreak{}SeeSawStep.\allowbreak{}stepX\_\allowbreak{}of\_\allowbreak{}laws}{Development.\allowbreak{}AbelianVarietyTheory.\allowbreak{}Proofs.\allowbreak{}DecoratedLatticeProofs.\allowbreak{}SeeSawDoubleStepLinearRelationsPeriods}{82--149}

For the terminal pair the upper expression is
\(W_{t_i,m-2}P_{t_i}D_{i,m}\). Indeed the terminal compact average at \(y',v'\) is a scalar multiple of the fixed branching map. Its multiplicity is one for the interlacing terminal weights. The scalar depends on \(y',v'\) and the weights. The finite class sum multiplying this branching map is \(P_{t_i}D_{i,m}\). Thus the extra scalar multiplies the entire upper relation at each test. Iterating the interior implication, beginning with this terminal pair, gives
\[
 \sum_i\beta_i H_{t_i}D_{i,m}=0
       \quad\Longrightarrow\quad
 \sum_i\beta_i I_{t_i}D_{i,0}(y,v)=0\quad\text{for all }y,v .
\]
At zero the chart set is a singleton. Its finite operator has coefficient one and acts by the single prescribed right translation. Hence \(D_{i,0}\) is the bottom Bessel period of the descended row with that translation.
\leanref{GlobalTheory.\allowbreak{}SeeSaw.\allowbreak{}perX\_\allowbreak{}of\_\allowbreak{}clsPinD}{Development.\allowbreak{}AbelianVarietyTheory.\allowbreak{}Proofs.\allowbreak{}DecoratedLatticeProofs.\allowbreak{}SeeSawLastPeriodFactorizationOfClassData}{72--142}
\leanref{GlobalTheory.\allowbreak{}SeeSaw.\allowbreak{}lastX\_\allowbreak{}of\_\allowbreak{}laws}{Development.\allowbreak{}AbelianVarietyTheory.\allowbreak{}Proofs.\allowbreak{}DecoratedLatticeProofs.\allowbreak{}SeeSawLastTwoStepsTransportLinearRelations}{93--137}
\leanref{GlobalTheory.\allowbreak{}SeeSaw.\allowbreak{}mainX\_\allowbreak{}of\_\allowbreak{}stepX\_\allowbreak{}lastX}{Development.\allowbreak{}AbelianVarietyTheory.\allowbreak{}Proofs.\allowbreak{}DecoratedLatticeProofs.\allowbreak{}MainXOfStepXLastX}{87--114}
\leanref{GlobalTheory.\allowbreak{}MiTop.\allowbreak{}linkZero\_\allowbreak{}of\_\allowbreak{}chartZero\_\allowbreak{}alt}{Development.\allowbreak{}AutomorphicArithmetic.\allowbreak{}ThetaCorrespondence.\allowbreak{}ThetaPairLifts.\allowbreak{}ThirdPlaceSeries.\allowbreak{}LinkZeroOfChartZeroAlt}{22--36}

A relation valid for all \(y,v\) gives the same relation between translated Petersson pairings. In the bottom Bessel integral, change variables by the finite column translation. The column becomes the model column after the bottom isometry. The row is translated by the prescribed row element and by a level element. The latter acts on every full reference multiplicity map by the same nonzero level character value, which cancels from the relation. The column coefficient vectors read in the frames \(y\) span the coefficient space. By linearity the relation holds at every coefficient vector. In the unfolding of the pairing with the period series, at a real group element \(g\) insert its column coefficient vector and translate \(v\) by \(g\). This gives the relation pointwise in the remaining real integral. Integration against its common measure proves the assertion for the Petersson pairings on the full level space.
\leanref{GlobalTheory.\allowbreak{}MiTop.\allowbreak{}botRel\_\allowbreak{}of\_\allowbreak{}pin\_\allowbreak{}of\_\allowbreak{}mem\_\allowbreak{}center}{Development.\allowbreak{}AutomorphicArithmetic.\allowbreak{}ThetaCorrespondence.\allowbreak{}ThetaPairLifts.\allowbreak{}ThirdPlaceSeries.\allowbreak{}BotRelTwo}{32--136}

Choose one reference top row descending to \(T_0\). Write \(S^{-1}\) for the column change in descent and \(R\) for the bottom isometry. Since \(S^{-1}\) multiplies norms by a positive constant, it is injective and hence onto the finite-dimensional permitted space. Choose \(\phi_0\) so that \(RS^{-1}\phi_0\) represents the column giving \(b_0\). For this reference top row and the column started at \(S^{-1}\phi_0\), put \(p=H_{t_0}D_{0,m}/I_{t_0}\).

For the second entry take the top row of \(F\) at \(t\). Its descent is the bottom row of \(\Theta_t j_t e_tF\): at every arrow the row map is the defining row lift, and the final row projection fixes its image. Start its column at \(S^{-1}\phi\). Commuting the column change through the terminal lift gives \(D_{1,m}=s_u A_t(\phi,F)\). Apply the preceding implication with coefficients \(\lambda/I_{t_0}\) and \(\lambda'/I_t\). We have proved
\[
       \lambda p+\lambda' c_tA_t(\phi,F)=0
          \quad\Longrightarrow\quad
       \lambda b_0+\lambda' B_t(\phi,F)=0.
\]
If \(p=0\), take \((\lambda,\lambda')=(1,0)\) to get \(b_0=0\). Hence \(p\ne0\). Taking \((\lambda,\lambda')=(c_t A_t(\phi,F),-p)\) proves the identity, including when the second pairing vanishes.
\leanref{GlobalTheory.\allowbreak{}MiTop.\allowbreak{}linkDown\_\allowbreak{}of\_\allowbreak{}ok}{Development.\allowbreak{}AutomorphicArithmetic.\allowbreak{}ThetaCorrespondence.\allowbreak{}ThetaPairLifts.\allowbreak{}ThirdPlaceSeries.\allowbreak{}GeometricDatumOfFamily}{77--86}
\end{proof}

\subsection*{Proof of \ref{B8.08}}\label{prf:B8.08}\phantomsection
\begin{proof}
Fix \(h_L,c_{\rm Dep}\) and use the list extension in (8.6). Its projection to the law keeps the point indices. The primitive row data and the definite evaluation are supplied by its parameters. The Weil action on those rows and the discrete-series realization give the recursion identification. Combining it with the definite evaluation gives the joins of Definition~\ref{B8.02}.

Fix a linked bundle in the common range of the row comparison. The matching and doubling identities at the prescribed levels give the surjections used in the downward seesaw. Use the row choices with common level data for the different tuples. The bottom comparison uses a single central finite translate for all words. Right translation by this element preserves the good subspace, since the translate is central, and preserves its norm by invariance of the Petersson measure. Denote its action on the output by \(V\).
\leanref{GlobalTheory.\allowbreak{}Fam.\allowbreak{}miStar2\_\allowbreak{}of\_\allowbreak{}allTrue\_\allowbreak{}alt\_\allowbreak{}alt}{Development.\allowbreak{}AutomorphicArithmetic.\allowbreak{}ThetaCorrespondence.\allowbreak{}ThetaPairLifts.\allowbreak{}ThirdPlaceSeries.\allowbreak{}MainTheoremFromPublishedResultsAltAltAltAlt}{23--28}
\leanref{GlobalTheory.\allowbreak{}MiTop.\allowbreak{}famAllQ\_\allowbreak{}of\_\allowbreak{}allTrue\_\allowbreak{}alt}{Development.\allowbreak{}AutomorphicArithmetic.\allowbreak{}ThetaCorrespondence.\allowbreak{}ThetaPairLifts.\allowbreak{}ThirdPlaceSeries.\allowbreak{}FamAllQOfAllTrue}{38--54}
\leanref{GlobalTheory.\allowbreak{}MiTop.\allowbreak{}exists\_\allowbreak{}isShift\_\allowbreak{}of\_\allowbreak{}center\_\allowbreak{}alt}{Development.\allowbreak{}AutomorphicArithmetic.\allowbreak{}ThetaCorrespondence.\allowbreak{}ThetaPairLifts.\allowbreak{}ThirdPlaceSeries.\allowbreak{}PostOnMemEEatOfCenter}{49--94}

For a fixed word \(u\), use the column map from the seesaw comparison. It multiplies norms by a positive constant. Divide it by this constant, compose with the bottom column isometry and express the result in the prescribed column coordinates. Its matrix \(U_u\) preserves the column norm. Choose the reference row in the full space of level-zero multiplicity maps. The reference descent has this entire space as target. Indeed, its arrows follow the reference row spaces at levels \(0,\ldots,m-1\), followed by the terminal space. On each bottom Hecke eigenspace the adjoint lift is injective by the nonvanishing in the doubling calculation. Distinct systems of unramified parameters stay distinct after adjoining the fixed chain parameters, since those parameters can be cancelled in the multisets at each place. The adjoint images of the distinct eigenspaces are therefore independent. The bottom space is their finite-dimensional span, so the adjoint has zero kernel and the descent is onto. This applies before imposing the good-type idempotent.
\leanref{GlobalTheory.\allowbreak{}ChainOnto.\allowbreak{}ontoAll\_\allowbreak{}of\_\allowbreak{}chainUp}{Development.\allowbreak{}AbelianVarietyTheory.\allowbreak{}Proofs.\allowbreak{}DecoratedLatticeProofs.\allowbreak{}ThetaLiftChainSurjectiveOfAdjointNonvanishing}{14--44}
\leanref{GlobalTheory.\allowbreak{}BaseChain.\allowbreak{}linkOnto\_\allowbreak{}of\_\allowbreak{}ontoAll\_\allowbreak{}third\_\allowbreak{}place\_\allowbreak{}series}{Development.\allowbreak{}AutomorphicArithmetic.\allowbreak{}ThetaCorrespondence.\allowbreak{}ThetaPairLifts.\allowbreak{}ThirdPlaceSeries.\allowbreak{}LinkOntoOfOntoAll}{46--54} By Lemma~\ref{B8.08f}, the diagonal Siegel--Weil estimate and the dimension bound for the ramified subspace give a column with nonzero period, and surjectivity of translated bottom-row evaluation gives a row pairing nontrivially with it. Use this full reference row as \(T_0\) in Lemma~\ref{B8.08i}. Reference descent supplies its top preimage. For the second row take the prescribed tuple and rational coordinate \(F\). Its image belongs to the good subspace by the row identification of Lemma~\ref{B8.08d} and the theta-row comparison in (8.6). Thus every arithmetic input uses the same reference pair for this word. In the complex coordinates the resulting identity reads
\[
 c_{u,t}\,\iota(\mathrm{Pt}_{u,t}(\phi,F))\,b_u
   =p_u\left\langle s_{EE}(V\Theta_tj_te_tF),\,
                         P_u(U_u\iota(\phi))\right\rangle ,
 \qquad p_u b_u\ne0 .
\]
More explicitly, let \(r_u>0\) be the norm multiplier of the column change and let \(\kappa\ne0\) convert the model pairing into the spectral pairing in the fixed row normalization. In the notation of Lemma~\ref{B8.08i}, the column isometry is \(r_u^{-1}RS^{-1}\). Substitution of \(r_u^{-1}\phi\) in that lemma gives \(p_u=p\), \(b_u=\kappa r_u^{-1}b_0\) and \(c_{u,t}=c_t\) in the direct complex coordinates. In the conjugate coordinates conjugate all three scalars and conjugate the column matrix entrywise. The model pairing is linear in its column, so the same displayed equality holds. Both normalization scalars are nonzero.
\leanref{GlobalTheory.\allowbreak{}MiTop.\allowbreak{}AtReal.\allowbreak{}anchor\_\allowbreak{}word\_\allowbreak{}alt}{Development.\allowbreak{}AutomorphicArithmetic.\allowbreak{}ThetaCorrespondence.\allowbreak{}ThetaPairLifts.\allowbreak{}ThirdPlaceSeries.\allowbreak{}MiStar2AtRealisation}{48--79}
\leanref{GlobalTheory.\allowbreak{}MiTop.\allowbreak{}miStar2\_\allowbreak{}at\_\allowbreak{}realisation\_\allowbreak{}alt}{Development.\allowbreak{}AutomorphicArithmetic.\allowbreak{}ThetaCorrespondence.\allowbreak{}ThetaPairLifts.\allowbreak{}ThirdPlaceSeries.\allowbreak{}MiStar2AtRealisation}{84--196}

Set
\[
       \zeta_{u,t}=\frac{c_{u,t}}{c_{u,\mu_0}},\qquad
       a'_u=\frac{p_u}{b_uc_{u,\mu_0}} .
\]
Division gives \eqref{eq:p8c1-multiplicity}, with \(\zeta_{u,\mu_0}=1\). Fix a finite place \(w\) and \(z\in K\) such that \(\iota(z)=\zeta_{u,\mu_w}\). Lemma~\ref{B8.08a} gives \(N>0\) and \(r\in\Q^\times\) with \(z^N=r\), by injectivity of \(\iota\). Every prime dividing the reduced numerator or denominator of \(r\) belongs to \(S_{\mu_w}\). The residue characteristic of \(w\) belongs to neither support. Hence both numerator and denominator are units in \(\mathcal O_{K,w}\), and
\[
                         N v_w(z)=v_w(r)=0 .
\]
This proves the finite-place clause with the same scalars as the period identity.
\leanref{GlobalTheory.\allowbreak{}MiTop.\allowbreak{}miStar2\_\allowbreak{}of\_\allowbreak{}anchored}{Development.\allowbreak{}AbelianVarietyTheory.\allowbreak{}Proofs.\allowbreak{}DecoratedLatticeProofs.\allowbreak{}MiStar2OfAnchored}{15--30}
\leanref{GlobalTheory.\allowbreak{}Fam.\allowbreak{}vw\_\allowbreak{}eq\_\allowbreak{}zero\_\allowbreak{}of\_\allowbreak{}unitRatio}{Development.\allowbreak{}AbelianVarietyTheory.\allowbreak{}Proofs.\allowbreak{}DecoratedLatticeProofs.\allowbreak{}ValuationEqZeroOfRootOfSUnitRatio}{58--94}

For the divisor row use the rational row projections and lifts of the law on this realization. The comparison then has the form \(\zeta_{u,t}\iota(B)=\iota(B')\) for the two rational pairing values on a given input. If some \(B\ne0\), its ratio \(B'/B\) is the unique field element representing \(\zeta_{u,t}\), and the comparison holds on every input. If \(\zeta_{u,t}\notin\iota(K)\), the displayed equation forces \(B=B'=0\) on every input; choose the field factor \(1\) there. This gives the rational factors for the divisor row, with valuation zero at the specified places.
\leanref{Formalism.\allowbreak{}IndefR.\allowbreak{}Cx.\allowbreak{}CxData.\allowbreak{}zetaOK\_\allowbreak{}of\_\allowbreak{}miUnit2}{Formalism.\allowbreak{}LocalToGlobal.\allowbreak{}IndefiniteReal.\allowbreak{}IsometricPeriodsOfUnitHypotheses}{42--58}
\leanref{Formalism.\allowbreak{}IndefR.\allowbreak{}Cx.\allowbreak{}CxData.\allowbreak{}matrix\_\allowbreak{}identity\_\allowbreak{}of\_\allowbreak{}unit\_\allowbreak{}hypotheses}{Formalism.\allowbreak{}LocalToGlobal.\allowbreak{}IndefiniteReal.\allowbreak{}IsometricPeriodsOfUnitHypotheses}{102--119}

The real-group decomposition gives the isotypic row spaces and their multiplicities. The Harish--Chandra parameters in that decomposition give the asserted infinite-place equality. For the finite types use the projection to the law, which preserves the assigned types and tame generators. The Green character formula in Theorem~\ref{B7.37} and the Zelevinsky input in \(P\) give the finite-type conditions on this assignment. Apply local Bessel uniqueness with these same types. These local identifications put each basis vector in its assigned single-label space. Apply Lemma~\ref{B8.08d} to the linked reference realization to obtain the isometry onto the full good subspace. In particular the number of selected basis indices of each good label is its global multiplicity. The recursion and divisor rows with the ray-class-field input give the geometric row. Theorem~\ref{B8.37} supplies the definite-side row with its cutoff for each \(C\ge C_0\).
\leanref{GlobalTheory.\allowbreak{}T.\allowbreak{}linked\_\allowbreak{}bundles\_\allowbreak{}good\_\allowbreak{}of\_\allowbreak{}local\_\allowbreak{}langlands\_\allowbreak{}alt\_\allowbreak{}alt}{Development.\allowbreak{}AutomorphicArithmetic.\allowbreak{}ThetaCorrespondence.\allowbreak{}ThetaPairLifts.\allowbreak{}ThirdPlaceSeries.\allowbreak{}LinkedBundlesGoodOfLocalLanglands}{96--97}
\leanref{GlobalTheory.\allowbreak{}Fc.\allowbreak{}LabLinked.\allowbreak{}ties\_\allowbreak{}of\_\allowbreak{}nAtData\_\allowbreak{}real}{Development.\allowbreak{}AutomorphicArithmetic.\allowbreak{}ThetaCorrespondence.\allowbreak{}ThetaPairLifts.\allowbreak{}ThirdPlaceSeries.\allowbreak{}RowLabOfLayer}{44--56}

All pointwise identifications were made at the indices of the linked bundle. Taking their eventual forms therefore gives the indicated rows on \(W_3\). For the transfer assertion, project a witness of the stronger extension to a witness of the old extension at the given indices. The linked bundle and the family selection then satisfy the old premise with its original cutoffs, which gives the claimed conclusion.
\end{proof}

\subsection*{Proof of \ref{B8.08b}}\label{prf:B8.08b}\phantomsection
\begin{proof}
Take the middle threshold of Definition~\ref{B8.39} to be at least \(8\). The family size constant \(C_{\rm sz}\) is a middle entry, so \(C_{\rm sz}\ge8\). Fix \(F_m,C_H\) in this range. We first check the sizes in \ref{B5.07y} uniformly in the member and layer. Write \(n=|A|\), \(r=k+m-1\), \(n_V=r+t\), and use \(b',w'',\eta,S_c\) from Definition~\ref{def:p8c2-point-sizes}. The rank ranges \eqref{eq:host-rank-ranges} give \(k\ge5\) and \(m\le k\) eventually. Thus \(r\le2k\). The rank identity of \ref{B8.59} gives \(n_V\ge8fr+4\), hence \(n_V\ge8\) and \(t\ge r+5\).

For the widths, \(b'/b=x/\exp(x^{\varepsilon/3})\) tends to zero and \(w''=b'/(xn)\). The width and spacing conditions in Definition~\ref{def:p8c2-point-sizes} give \(\eta\le w''\) and \(S_c\le w''\). The top-weight inequalities there give \(b_s\le(1+x^{-2}/4)b'\) at a top position \(s\). By the three intervals of \ref{B3.47}, the layer shift is positive and is bounded above by \(b_s+b'/x+\eta\), \(b+\eta+1\), or \(w''+\eta\), according to the position. All three are at most \(2b\) beyond one family cutoff. The cloud entries lie between zero and \(S_c\), and the auxiliary entries differ from the shift by at most \(\eta\), by the weight-box conditions. This gives the two bounds by \(b\) in \ref{B5.07y}. The type-prime condition in Definition~\ref{def:p8c2-point-sizes} gives \(g_\nu\ge(kb)^8\) for \(\nu=0,1\). These conditions concern the member before the point realization or \(C\) is chosen.
\leanref{GlobalTheory.\allowbreak{}LowRow.\allowbreak{}weightSizes\_\allowbreak{}along}{Development.\allowbreak{}AbelianVarietyTheory.\allowbreak{}Proofs.\allowbreak{}DivisorProofs.\allowbreak{}ThetaLiftWeightSizesEventuallyAlongFamily}{16--60}

Fix a point in \(\mathcal P\) at such a member. Its primitive row has separated types and the level-one Hecke decomposition of \ref{B8.08}. Its normalized realization supplies the standard definite lattice, base phase one, the extremal vectors, and the class expansion with its level action as in \ref{B8.55}. The twist set of the reference tuple is empty. Each marked prime lies over one of \(g_0,g_1\). The host level conditions make these primes coprime to \(Q_Sqd\) and to the conductor ideal. Thus the prime restriction in \ref{B5.07y} holds. The analytic inputs for that lemma are supplied by \(P\): rapid decay, the unitary Siegel--Weil formulas, the local doubling formula and its unramified evaluation, finite length, and the compact-dual formal-degree formula.
\leanref{GlobalTheory.\allowbreak{}LowRow.\allowbreak{}primeToHost\_\allowbreak{}base}{Development.\allowbreak{}AbelianVarietyTheory.\allowbreak{}Proofs.\allowbreak{}DivisorProofs.\allowbreak{}BadPlacesPrimeToHostBaseTuple}{13--18}
\leanref{GlobalTheory.\allowbreak{}LowRowAdel.\allowbreak{}adelC\_\allowbreak{}of\_\allowbreak{}siegelWeil\_\allowbreak{}of\_\allowbreak{}doubling\_\allowbreak{}alt}{Development.\allowbreak{}AutomorphicArithmetic.\allowbreak{}ThetaCorrespondence.\allowbreak{}ThetaPairLifts.\allowbreak{}LocalTypeDictionary.\allowbreak{}FinitenessInputsRealData}{13--27}

Put \(H_j=\sum_\rho\log\operatorname{Hook}(a^+_{j\rho})\), where \(a^+_{j\rho}\) is the positive partition of \ref{B5.07y}, and put \(D_j=H_j-8fn_V^2L'\), with \(L'=\log(b\,\operatorname{rd}(E)g_0g_1n_V)\). The coefficient of \(C\) in the sixth entry is zero at this scale. The identification of the positive partitions in \ref{B5.07b} therefore gives, for every \(C\),
\[
 L_{6,j}=H_j-C_{\rm sz}fn_V^2L'\le D_j .
\]
Indeed \(L'\ge0\), since \(b\ge1\), \(\operatorname{rd}(E)\ge1\) and the remaining factors are positive integers.

At the distinguished place the primitive row prescribes the scalar of the projected lift times its adjoint on each occurring block. At every other place it prescribes the scalar in the conjugate pairing on rational layer coordinates. These are the two scalar identities in \ref{B5.07b}. The one-map comparison and the phase normalization at our point give the comparison of the rational lift in \ref{B5.07f}. Thus \ref{B5.07b} applies, including its convention for an absent block. It gives the asserted bound at every place, with the cutoff already chosen.
\leanref{GlobalTheory.\allowbreak{}LowRow.\allowbreak{}lamLow\_\allowbreak{}closing\_\allowbreak{}alt}{Development.\allowbreak{}AutomorphicArithmetic.\allowbreak{}ThetaCorrespondence.\allowbreak{}ThetaPairLifts.\allowbreak{}LocalTypeDictionary.\allowbreak{}FinitenessInputsRealDataThree}{14--23}
\end{proof}

\subsection*{Proof of \ref{B8.08c}}\label{prf:B8.08c}\phantomsection
\begin{proof}
Subtract \(\sum_j2e_wF_{0,j,w}+c_w\) from the first inequality and take the positive part. Since all the \(u_{jw}\) are nonnegative, the result \(R_w\) satisfies
\[
 0\le R_w\le\sum_j u_{jw}.
\]
The index error prescribed at the point is \(e'_{jw}=R_w/m\). Its sum over the layers is \(R_w\), which proves the first two reference inequalities. For each layer \(h\), nonnegativity of \(J_{0,h,w}\) gives \(\sum_w\ell_wu_{hw}\le aU\). Hence
\[
 \sum_w\ell_we'_{jw}
 \le \frac1m\sum_{h=1}^{m}\sum_w\ell_wu_{hw}\le aU.
\]
The hypothesis for layer \(j\) also gives \(\sum_w\ell_we_wJ_{0,j,w}\le aU\). Addition proves the third inequality for \(C=2a\), and \(U\ge0\) gives it for every larger \(C\). Thus the same \(x_0\) works after \(C\) is quantified. The construction uses the thresholds of the hypotheses without increasing them.
\end{proof}

\subsection*{Proof of \ref{B8.09}}\label{prf:B8.09}\phantomsection
\begin{proof}
Take \(R=\top\) in (8.13). Good members supply the symmetric column system with finite-order terminal and level-zero characters valued in \(\iota_D(K_D)\), the reference forms map with its label injectivity, stabilizer bijections and norm bound \eqref{eq:p8c1-point-norm}, the integral interlacing lowering maps on the reference and used twisted chains, and the Petersson class decomposition. Lemma~\ref{lem:p8c1-row-forms} supplies rational row lifts of \(\mathcal O_{K_D}\)-combinations of the row generators, whose class polynomials lie in the local Weyl duals at all allowed places. Lemma~\ref{B4.18a} gives integrality of the rational point maps in items~1 and~2 and represents each period by \(y=\sum_c b_c/h_c\in K_D\), where \(h_c\mid D_{\Gamma,u}\) and \(b_c\) is integral at every allowed \(w\). This gives \eqref{eq:p8c1-period-valuation} simultaneously at these places. The real comparison gives item~5. Apply \ref{B4.22a} to the tied level chain for the reference scaling and the localized twisted representation.
\end{proof}

\subsection*{Proof of \ref{B8.10}}\label{prf:B8.10}\phantomsection
\begin{proof}
On \(W_3\) the level primes of the twisting selection eventually satisfy \(b^{m/2}\le T_w\) and \(\log T_w\le C_{Tu}m\log b\), uniformly in the point and word. Apply \ref{B8.49} to a unit hyperplane in the strict range of \ref{B8.13} for one represented vector and ambient rank at least three. For the level-frame row put \(g=g_0g_1,\ a=\pi/v_E,\ \xi=a/g\), and \(d_E=\operatorname{rd}(E)\). The critical level is the positive number defined by
\[
 T_{\mathrm{crit},\tau}(w)^k
 =a^{-[(m-1)k+m(m-1)/2]}d_E^{-(m-1)/2}
  g^{km+m(m-1)/2}\frac{VV_\tau}{V(w_\tau)},
 \qquad VV_\tau=d_{\sigma,\tau}\prod_{j=0}^{k+m-2}j!,
\]
where \(V\) is the Vandermonde product and \(d_{\sigma,\tau}\) is the local formal-degree factor in the column normalization.
\leanref{Formalism.\allowbreak{}MainNotions.\allowbreak{}S.\allowbreak{}Tcrit}{Formalism.\allowbreak{}SharpRestriction.\allowbreak{}CuspSpanNormBounds.\allowbreak{}StandingNumericalData}{39--43}

Raise the last threshold of Definition~\ref{B8.39} to at least \(40\). Each of the six entries defining \(C_\ell\) belongs to the last family list, so \(C_\ell\ge40\). Use \(C_3=1\), and cutoffs with \(2C_3\le m,\ C_u\log b\le m-1,\ C_0e^8\le k\), where \(C_u=3\) is the weight bound parameter and \(C_0=3\) is the dimension bound parameter. The mass and column-dimension formulas give
\(\prod_\tau(T_{\mathrm{crit},\tau}(w)/T_w)^k\le e^{f2^{3-k}}D^\star/d_w^S\).
The local splitting exponent bound gives
\(T_{\mathrm{crit},\tau}(w)\ge g^m a^{-(m-1)}\pi^{m-1}k^{3r_S}\) at each place. Since \(a\le\pi\), and the source dimension bound gives \(D^\star/d_w^S\le C_0e^{k/4}\), these estimates yield
\begin{equation}\label{eq:p8c1-critical-level}
 g^2\le \xi^{2r_S}T_{\mathrm{crit},\tau}(w),
 \qquad
 \prod_\tau(T_{\mathrm{crit},\tau}(w)/T_w)^k
 \le e^{f2^{3-k}}D^\star/d_w^S\le e^8 C_0e^{k/4}.
\end{equation}
If \(T_w\xi^{2r_S}<1\), each factor of the product is at least \((g^2)^k\). The upper bound is strictly smaller by \(C_0e^8\le k<g^k\) and \(e^{k/4}\le g^k\). This proves \eqref{eq:p8c1-level-row}.
\leanref{Formalism.\allowbreak{}MainNotions.\allowbreak{}BridgeLemma.\allowbreak{}TFrame\_\allowbreak{}of\_\allowbreak{}sizes}{Formalism.\allowbreak{}SharpRestriction.\allowbreak{}CuspSpanNormBounds.\allowbreak{}BruhatCellsKronecker.\allowbreak{}CriticalLevelLowOfRefined}{408--427}
\leanref{GlobalTheory.\allowbreak{}bothRowsFromSiegelWeil\_\allowbreak{}alt}{Development.\allowbreak{}AutomorphicArithmetic.\allowbreak{}ThetaCorrespondence.\allowbreak{}ThetaPairLifts.\allowbreak{}ThirdPlaceSeries.\allowbreak{}BothRowsFromSiegelWeil}{11--15}

Let \(L_*=F_m.c_{\rm out}\mathbin{\Vert}F_m.c_{\rm mid}(CH)\), where \(\Vert\) denotes list concatenation. Evaluate the loop thresholds at \(C_{YL},C_{Tu},a_F(F_m.c_{\rm out})\) and the row-bound threshold computed from \(a_F,L_*\). Impose these together with the size, volume and raised selection cutoffs. Then \ref{B8.10a} gives \eqref{eq:p8c1-loop-rows}. For \eqref{eq:p8c1-distinct-row}, use \ref{B5.11b} with \(C_0=3\) and the constant function \(a_*=3e^{7/4}\). It bounds the level ratio in either order by \(a_*\) and the Vandermonde ratio by \(a_*^k\). Use the loop row with \(C\ge4\sqrt{a_*}\). Entry \(17\) of \(L_*\), counted from zero, is \(C_\Gamma>0\). The list determines \(k_0,C_n,C'_0,C_{\ell,0}\) for \ref{B8.10b}. Raise the last threshold to \(\max(C'_0,29,C_{\ell,0},40)\). The entry \(C_k\) and all six entries in the minimum belong to the last list. Thus \(C'\ge\max(C'_0,29)\) and \(C_\ell\ge\max(C_{\ell,0},40)\). The source and selection supply the counting and gap hypotheses, ties and convergence. Apply \ref{B8.10b} with the nongeneric estimate of \ref{B8.10d}. Take the maximum of the point cutoffs and \(R=\top\) in (8.13).
\end{proof}

\subsection*{Proof of \ref{B8.10c}}\label{prf:B8.10c}\phantomsection
\begin{proof}
Write \(r=r_S,\ g=g_0g_1\). At a represented point \(w\), let \(D_i(w)\) be the leading minors. Its kind \(\kappa(w)\) is \(D_{r+1}/D_r\) for \(D_r\ne0\) and a separate symbol for \(D_r=0\). Set
\[
 G(w)=D_{r+1}\overline{D_{r+1}}-TT'D_r\overline{D_r},\qquad
 n(w)=c_{\rm count}\alpha |N_{E/\Q}(G(w)/g^m)|^{(k-1)/2}.
\]
The count constant \(c_{\rm count}\) is specified below. At a finite place \(v\), partition integral matrices in the lower and upper unipotent orbit of \(w\) into integral double orbits \(O\). Let \(h_v(w,O)\) be the orbital measure with its character and terminal phase, and \(\theta(O)\) the integral presentation class after the last diagonal shift. Put \(B_v(w,\theta)=\sum_O h_v(w,O)\#\mathrm{Iso}(\theta,\theta(O))\). An isomorphism set is an automorphism torsor for equal classes and empty otherwise. The orbital sum for \(\theta\) is therefore \(B_v(w,\theta)/|\operatorname{Aut}(\theta)|\).

Index the coefficients by \(j=(\kappa,\boldsymbol\theta)\) with the standard class outside finitely many places. Put \(v_j(w)={\bf1}_{\kappa(w)=\kappa}\prod_v B_v(w,\theta_v)\) and \(a_j=\prod_v|\operatorname{Aut}(\theta_v)|^{-1}\). Expansion counts the local isomorphisms between the two presentations with conjugate phase at the second point. The local square sums converge and the reference factor is one away from the exceptional places. Thus on each finite point list
\begin{equation}\label{eq:p8c1-point-kernel}
 K_{\rm pt}(w_x,w_y)=\sum_j a_jv_j(w_x)\overline{v_j(w_y)}.
\end{equation}
Removing finitely many summands leaves a positive semidefinite matrix.
\leanref{Formalism.\allowbreak{}L2.\allowbreak{}SG.\allowbreak{}Pf.\allowbreak{}ktop\_\allowbreak{}sub\_\allowbreak{}posSemidef}{Formalism.\allowbreak{}LocalToGlobal.\allowbreak{}SiegelLatticeTheory.\allowbreak{}CokernelGLCounts.\allowbreak{}RankPointCounts.\allowbreak{}SupportProposition}{533--596}

Let \(\mathcal E_x\) be the generic terms with \(G\ne0\) at \(w_x\). For \(t\) write \(\rho(t)\) for its scalar-partner class, \(m_t\) for its row mass divided by total row mass, \(h_t(x)\) for its orbital coefficient and \(j(t)\) for its kernel index. With \(d_j\) the density divided by \(\alpha\), the count identities at positive normalizers are
\begin{equation}\label{eq:p8c1-partner-fibre}
 \sum_{t'\in\mathcal E_y:\,\rho(t')=\rho(t)}h_{t'}(y)=v_{j(t)}(w_y),\qquad
 \sum_{t\in\mathcal E_x:\,j(t)=j}m_th_t(x)=n(w_x)d_j a_jv_j(w_x).
\end{equation}
For the first identity, fix a tuple of admissible double orbits at \(w_y\). At its nonstandard places and the exceptional reference places choose local unitriangular parameters and congruence depths realizing it. Approximate the parameters entrywise by a rational lower unitriangular matrix \(L\) and an upper one \(U\), with denominators supported on those places. For a point \(w\) with block data \((t_i,s_i)\), let \(M(w)\) be the block diagonal matrix whose \(i\)-th block is \(t_i\) times the order-reversing permutation matrix of size \(s_i\). Then \(LM(w_y)US_g\), with \(S_g\) the standard involution in point coordinates, realizes the tuple. At all other places the matrices and their inverses are integral and preserve the reference orbit. Intersect the local partner lattices with the rational space to obtain the partner frame lattice. Localization recovers the lattices with their identifications, so the partner count is the restricted product of local isomorphism counts. An inadmissible local orbit contributes zero. Weight by the local phase measures and sum over orbit tuples.

For the second identity, the perpendicular space of the rank-\(m\) frame is definite of rank \(k-1\). Unimodular extensions correspond to complement lattices with their discriminant modules glued to that of the frame. Count the gluings with inverse stabilizer weights. Nonintegral frames have count zero. At an integral place with \(c_v\) nonunit elementary divisors, put \(R=k+m-1\) and \(\alpha_v(a)=\prod_{i=1}^a(1-\eta_v^i q_v^{-i})\), where \(q_v\) is the residue cardinality and \(\eta_v=1,-1\) in the split and inert cases. The local complement ratio is
\[
 \frac{\alpha_v(k-1)}{\alpha_v(k-1-c_v)}
 =\frac{\alpha_v(k-1)}{\alpha_v(R)}
   \frac{\alpha_v(R)}{\alpha_v(k-1-c_v)} .
\]
The two factors give the product of \(L(i,\eta^i)\) for \(k\le i\le R\) and the represented densities, respectively. The determinant factor is
\(|\operatorname{disc}F|^{-2r(r+1)(k-1)}|N_{E/\Q}(G/g^m)|^{(k-1)/2}\), where \(F\) is the maximal real subfield of \(E\). Take \(c_{\rm count}\) to be the rank-\((k-1)\) mass constant times the field factor and the \(L\)-values divided by total row mass. Dividing the local orbital sums by the automorphism orders gives the second identity.
\leanref{Formalism.\allowbreak{}L2.\allowbreak{}SG.\allowbreak{}Pf.\allowbreak{}formulaExplicit\_\allowbreak{}of}{Formalism.\allowbreak{}LocalToGlobal.\allowbreak{}SiegelLatticeTheory.\allowbreak{}CokernelGLCounts.\allowbreak{}RankPointCounts.\allowbreak{}DualFamilyCounting.\allowbreak{}Assemble}{158--248}
\leanref{Formalism.\allowbreak{}L2.\allowbreak{}SG.\allowbreak{}LsDatumB.\allowbreak{}partnerSumC\_\allowbreak{}of}{Formalism.\allowbreak{}LocalToGlobal.\allowbreak{}SiegelLatticeTheory.\allowbreak{}CokernelGLCounts.\allowbreak{}RankPointCounts.\allowbreak{}LsDatumBPf}{27--82}
\leanref{Formalism.\allowbreak{}L2.\allowbreak{}SG.\allowbreak{}LsDatumB.\allowbreak{}fiberCntC\_\allowbreak{}of}{Formalism.\allowbreak{}LocalToGlobal.\allowbreak{}SiegelLatticeTheory.\allowbreak{}CokernelGLCounts.\allowbreak{}RankPointCounts.\allowbreak{}LsDatumBPf}{83--173}

The point-coordinate row is \(C_\rho(x)=n(w_x)^{-1/2}\sum_{\rho(t)=\rho}\sqrt{m_t}h_t(x)\) for \(n(w_x)>0\), and zero otherwise. Scalar partners have equal row mass. By \eqref{eq:p8c1-partner-fibre}, on positive support the Gram entry is
\[
 \sqrt{n(w_x)/n(w_y)}\sum_j d_j a_jv_j(w_x)\overline{v_j(w_y)}.
\]
The sum uses the indices of the represented classes. On the support of \(v_j\), the shifted determinants are locally associated. If \(D_r=0\), this fixes the absolute norm of \(D_{r+1}\). If \(D_r\ne0\), the common kind fixes the shift factor and the associations fix the absolute norm of \(D_r\). In that case \(G=D_r\overline{D_r}(\kappa\bar\kappa-TT')\). Hence \(n\) is constant on each support.
\leanref{Formalism.\allowbreak{}L2.\allowbreak{}SG.\allowbreak{}Pf.\allowbreak{}supportVec\_\allowbreak{}of\_\allowbreak{}normProd}{Formalism.\allowbreak{}LocalToGlobal.\allowbreak{}SiegelLatticeTheory.\allowbreak{}CokernelGLCounts.\allowbreak{}RankPointCounts.\allowbreak{}SupportProposition}{484--504}
\leanref{Formalism.\allowbreak{}L2.\allowbreak{}SG.\allowbreak{}DataCut.\allowbreak{}rows\_\allowbreak{}pair}{Formalism.\allowbreak{}LocalToGlobal.\allowbreak{}SiegelLatticeTheory.\allowbreak{}CokernelGLCounts.\allowbreak{}RankPointCounts.\allowbreak{}DataCut}{674--753}

Each rank-one summand of \eqref{eq:p8c1-point-kernel} therefore enters the row Gram matrix with a multiplier in \([0,s_d]\), with zero multiplier on zero support. Subtract the row Gram matrix from \(s_dK_{\rm pt}\) term by term. The difference is positive semidefinite. Pull back to the flag points and take the diagonal congruence by the inverse positive flag normalizers, with zero for flags without points. Restricting the rows to \(S\) adds positive Gram matrices to the difference. Evaluate at \(z\).
\end{proof}

\subsection*{Proof of \ref{B8.10ca}}\label{prf:B8.10ca}\phantomsection
\begin{proof}
Put \(M=\max(u_l,u'_l)\) and \(m_l=\min(u_l,u'_l)\). The displacement bound gives \(2-M\le y,y+1\le-2-m_l\) for \(y\in Y\). For \(i<l\) both weights have \(i\)-th coordinate at least \(M+1\), and for \(i>l\) both have \(i\)-th coordinate at most \(m_l-1\). Thus the two products in \(q(y)\) have opposite signs from coordinate \(l\), and each denominator factor has absolute value at least one. In particular \(q(y)>0\).
\leanref{CompanionTheory.\allowbreak{}WaveGrid.\allowbreak{}nq\_\allowbreak{}onCut\_\allowbreak{}pair}{CompanionTheory.\allowbreak{}AutomorphicTheory.\allowbreak{}BesselWaveGrids.\allowbreak{}NqMovedExample}{85--106}

For \(x=-y\) pair \(\lambda_i>x\) with the upper neighbour and \(\lambda_i\le x\) with the lower one. Monotonicity of \(\lambda_i\) makes the two lists of neighbours disjoint. For the paired factor the half-unit inequalities give \(|y+1/2+u_j|\ge|y+\lambda_i|\), and likewise for \(u'\). For \(|y+\lambda_i+e|+c\), \(e=0,1\), the nearest paired factor on either side of \(x\) costs at most \(2+c\). The \(h\)-th further factor has distance at least \(h(1+c)\) and costs at most \(1+1/h\). These factors cost at most \(k\) on each side. There are at most \(2r+3\) unpaired numerator factors, including the exceptional index, each at most \(3b\). Using the unit lower bounds on the remaining denominators gives a cost of at most \((2+c)^2 k^2(3b)^{2r+3}\le b^{4r+10}\). Use \(u\) at \(e=0\) and \(u'\) at \(e=1\) to obtain
\begin{equation}\label{eq:p8c1-pstar}
 P_*(y)P_*(y+1)\le b^{8r+20}q(y)/(T_pT_q).
\end{equation}
\leanref{GroundTheory.\allowbreak{}BlockD.\allowbreak{}FourierDecayAlt.\allowbreak{}WaveGridLiftGrid.\allowbreak{}pst\_\allowbreak{}mul\_\allowbreak{}le}{CompanionTheory.\allowbreak{}AutomorphicTheory.\allowbreak{}BesselWaveGrids.\allowbreak{}InterlacingGrids.\allowbreak{}PstMul}{23--56}

There are at least \(3(d+1)\) unused nodes since \(d\le2r\) and \(|w|\le N-3\). Spacing five permits the positive and negative extensions. At an unused node each existing letter contributes a positive factor and costs at most \(4/3\) in \(q(y)/q_w(y)\). Thus the square root in \eqref{eq:p8c1-wave-grid} is bounded by \(\Lambda\) by \eqref{eq:p8c1-pstar}. The degree-\(d\) coefficient of the \(j\)-th Lagrange basis polynomial in row \(a\) is \(L_{aj}\), with \(\sum_j|L_{aj}|\le(d+1)5^{-d}\le1\). Factor the three row sums.
\leanref{CompanionTheory.\allowbreak{}WaveGrid.\allowbreak{}gridOK\_\allowbreak{}of\_\allowbreak{}free}{CompanionTheory.\allowbreak{}AutomorphicTheory.\allowbreak{}BesselWaveGrids.\allowbreak{}Bridge}{58--113}
\end{proof}

\subsection*{Proof of \ref{B8.10o}}\label{prf:B8.10o}\phantomsection
\begin{proof}
If \(I_0=0\), the zero function gives the expansion. Otherwise discard zero-energy children and put \(V_\mu=(I_0/I_\mu)^{1/2}\mathcal J_\mu\) at each surviving index. These functions have energy \(I_0\) and \(V_0=\mathcal J_0\). The coefficient from \(\mu\) to \(\mu+\mathbf1_I\) becomes \(\tau_{\mu I}(I_{\mu+\mathbf1_I}/I_\mu)^{1/2}\). Summing over the at most \(n^s\) children gives
\[
 \sum_I|\tau_{\mu I}|(I_{\mu+\mathbf1_I}/I_\mu)^{1/2}
 \le (2 n^{3/2}H)^s .
\]
Here \(\sqrt{\exp(1)}<2\). We have \(2n^{3/2}\le b^{5/2}\). Since \(m=2s\ge2\),
\[
 s\big((C_h+10)m+5/2\big)\le (C_h/2+6)m^2 .
\]
Expand successively along paths of length \(a\). At each step the sum of absolute coefficients is multiplied by at most this one-step bound. Combining paths ending at the same index can only decrease that sum.
\end{proof}

\subsection*{Proof of \ref{B8.10hba}}\label{prf:B8.10hba}\phantomsection
\begin{proof}[Proof (outline)]
Choose \(\beta\in\C\) and strictly decreasing integers \(L_0,\ldots,L_{k-1}\) such that
\[
 u_{j+1}-k/2+1/2=\beta+L_j\quad(0\le j<k).
\]
Then \(L_i-L_j=u_{i+1}-u_{j+1}\), so \(L_0-L_{k-1}\le2b\). Choose an orthonormal basis \(e_0,\ldots,e_{k-1}\) of the positive Hermitian space \(E_z\), and put \(e=e_{k-1}\). Under the subgroup fixing \(e\), the Capelli eigenspaces \(V_M\), with shifted parameters \(\beta+M_j\), span \(V\), where
\[
 M\in\Z^{k-1},\qquad L_{j+1}<M_j\le L_j\quad(0\le j<k-1).
\]
Define
\begin{equation}\label{eq:p8c1-contraction}
 w_+(L,M,h)=\frac{\prod_{j=0}^{k-2}(M_j-1-L_h)}{\prod_{i\ne h}(L_i-L_h)},\qquad
 w_-(L,M,h)=\frac{\prod_{j=0}^{k-2}(M_j-L_h)}{\prod_{i\ne h}(L_i-L_h)}.
\end{equation}
\leanref{GlobalTheory.\allowbreak{}WigSh.\allowbreak{}wP}{Development.\allowbreak{}AbelianVarietyTheory.\allowbreak{}Definitions.\allowbreak{}Jets.\allowbreak{}Imports}{9--11}
\leanref{GlobalTheory.\allowbreak{}WigSh.\allowbreak{}wM}{Development.\allowbreak{}AbelianVarietyTheory.\allowbreak{}Definitions.\allowbreak{}Jets.\allowbreak{}Imports}{12--14}

For a strictly decreasing list \(L\) whose shift \(L+\mathbf1_h\) is strictly decreasing, the one-step raising construction gives a degree-one contraction from the type shifted by \(+\mathbf1_{h+1}\) into the original type, with
\(A(e)A(e)^*v=w_+(L,M,h)v\) on \(V_M\). The lowering construction has the same conclusion with \(-\mathbf1_{h+1}\) and \(w_-\), when \(L-\mathbf1_h\) is strictly decreasing. Composing these maps gives the chains. More precisely, if \(L_h+d<L_j\) for every \(j<h\), there is a degree-\(d\) raising contraction with
\[
 A(e)A(e)^*\bigm|_{V_M}
 =\prod_{a=0}^{d-1}w_+(L+a\mathbf1_h,M,h)\,\operatorname{id}_{V_M}.
\]
If \(L_j<L_h-d\) for every \(j>h\), the lowering chain has eigenvalue
\[
 \prod_{a=0}^{d-1}w_-(L-a\mathbf1_h,M,h).
\]
For the induction, the degree-zero map is the identity. If \(A\) is a chain and \(B\) the next one-step map, their composition has the stated polynomial degree, equivariance and real homogeneity. Its operator norm is bounded by the product of the two operator norms. The operator \(A(e)^*\) maps into the eigenspace with the same Capelli parameter in the intermediate representation because it intertwines the subgroup fixing \(e\). Thus
\((AB)(AB)^*=A(BB^*)A^*\) multiplies the preceding scalar by the next factor. The strict gap conditions ensure that every intermediate type and every one-step map is defined.
\leanref{GlobalTheory.\allowbreak{}WigSh.\allowbreak{}chain\_\allowbreak{}raise}{Development.\allowbreak{}AbelianVarietyTheory.\allowbreak{}Proofs.\allowbreak{}JetHarmonics.\allowbreak{}ExistsWignerContractionChainRaising}{15--98}
\leanref{GlobalTheory.\allowbreak{}WigSh.\allowbreak{}exists\_\allowbreak{}wigMap\_\allowbreak{}chain\_\allowbreak{}lowering}{Development.\allowbreak{}AbelianVarietyTheory.\allowbreak{}Proofs.\allowbreak{}JetHarmonics.\allowbreak{}ExistsWignerContractionChainLowering}{17--104}

The scalar estimate used at each step is as follows. Suppose \(D\ge1\), \(L_0-L_h\le D\), and \(L_h-L_{k-1}\le D\). Then
\[
 w_+(L,M,h)\ge D^{-2}\quad\hbox{if every }M_j\ne L_h+1,
 \qquad
 w_-(L,M,h)\ge D^{-2}\quad\hbox{if every }M_j\ne L_h.
\]
To prove this, set \(d_i=\max(1,|L_i-L_h|)\), so \(d_h=1\), and put
\[
 Q=\prod_{j<h}d_{j+1}\prod_{j\ge h}d_j,
\]
where both products use \(0\le j<k-1\). Integrality and interlacing show that each absolute numerator factor is at least \(d_{j+1}\) when \(j<h\), and at least \(d_j\) when \(j\ge h\). The stated nonzero-factor conditions handle the adjacent factor. The denominator has absolute value \(Qd_0d_{k-1}\). Numerator and denominator have the same sign, since each has \(k-1-h\) negative factors. Consequently each quotient is at least \((d_0d_{k-1})^{-1}\ge D^{-2}\).
\leanref{GlobalTheory.\allowbreak{}WigSh.\allowbreak{}inv\_\allowbreak{}sq\_\allowbreak{}le\_\allowbreak{}wP}{Development.\allowbreak{}AbelianVarietyTheory.\allowbreak{}Proofs.\allowbreak{}JetHarmonics.\allowbreak{}InvSqLeWignerCoeffSqRaising}{35--96}
\leanref{GlobalTheory.\allowbreak{}WigSh.\allowbreak{}inv\_\allowbreak{}sq\_\allowbreak{}le\_\allowbreak{}wM}{Development.\allowbreak{}AbelianVarietyTheory.\allowbreak{}Proofs.\allowbreak{}JetHarmonics.\allowbreak{}InvSqLeWignerCoeffSqLowering}{60--103}

Put \(i=l-1\). If \(M_i\ge L_{i+1}+d+1\), use the chain raising coordinate \(i+1\); otherwise use the chain lowering coordinate \(i\). For every \(0\le a<d\), the interlacing inequalities persist and the required numerator factors are nonzero. This follows for the distinguished coordinate from \(L_i-L_{i+1}\ge2d+1\); the other coordinates lie strictly above \(L_i\) or at most \(L_{i+1}\). The endpoint distances of the shifted list are at most \(2b+a\le3b\). Hence one of the two chain eigenvalues is at least \((3b)^{-2d}\) on each \(V_M\).

Let \(X=A_+(e)A_+(e)^*+A_-(e)A_-(e)^*\). Both summands are positive and self-adjoint. On every nonzero \(V_M\), \(X\) acts by the sum of the two nonnegative chain eigenvalues, hence by a scalar at least \((3b)^{-2d}\). Since these spaces span \(V\), the spectral theorem gives
\[
 (3b)^{-2d}\|v\|^2
 \le\Re\langle v,Xv\rangle
 =\|A_+(e)^*v\|^2+\|A_-(e)^*v\|^2.
\]
Finally write \(\omega=tae\), with \(a\in U_z\), \(t\in\R\), and \(t^2=\Re h(\omega,\omega)\). Equivariance and unitarity transfer the unit-vector inequality to \(ae\); homogeneity multiplies it by \(t^{2d}\). This proves the asserted bound.
\leanref{GlobalTheory.\allowbreak{}WigSh.\allowbreak{}chain\_\allowbreak{}lower\_\allowbreak{}bound}{Development.\allowbreak{}AbelianVarietyTheory.\allowbreak{}Proofs.\allowbreak{}JetHarmonics.\allowbreak{}ChainLowerBound}{22--116}
\leanref{GlobalTheory.\allowbreak{}norm\_\allowbreak{}sq\_\allowbreak{}adjoint\_\allowbreak{}add\_\allowbreak{}ge\_\allowbreak{}of\_\allowbreak{}joint\_\allowbreak{}eigenspaces}{Development.\allowbreak{}AbelianVarietyTheory.\allowbreak{}Proofs.\allowbreak{}JetHarmonics.\allowbreak{}NormSqAdjointAddGeOfJointEigenspaces}{7--83}
\leanref{GlobalTheory.\allowbreak{}WigSh.\allowbreak{}lower\_\allowbreak{}bound\_\allowbreak{}spreads}{Development.\allowbreak{}AbelianVarietyTheory.\allowbreak{}Proofs.\allowbreak{}JetHarmonics.\allowbreak{}NormAdjointApplyOfNormPreserving}{46--59}
The complete proof is contained in \Lean{GlobalTheory.\allowbreak{}WigSh.\allowbreak{}exists\_\allowbreak{}shifted\_\allowbreak{}types\_\allowbreak{}lower\_\allowbreak{}bound}.
\end{proof}

\subsection*{Proof of \ref{B8.10hb}}\label{prf:B8.10hb}\phantomsection
\begin{proof}[Proof (outline)]
Set \(p_0=2(r+1)\), \(K=3\), \(N=3(r+1)\), and for \(0\le e\le C_am^3\) take the displacement budget
\[
 q_0=(r+1)e+40(r+1)^3+(r+1)(r+2),
\]
using \(c=q_0\) in \(P_*(z)\) of Definition~\ref{B8.10ha}. Write \(G_{\rm loc}\) for the Gram polynomial in standardized frame coordinates and put \(H=G_{\rm loc}^{p_0}F\). Multiplication by a matrix entry or its conjugate tensors the coefficient with the standard representation or its dual. The factor \(G_{\rm loc}^{p_0}\) adds at most \(mp_0\le c\) such steps and decomposes \(H=\sum_\nu H_\nu\) into summands of infinitesimal character \(\lambda+\nu\) with \(|\nu_j|\le c\). Distinct infinitesimal characters are orthogonal for \(\beta\), and on \(H_\nu\) each factor \(C(z)\) in \(C_w\) has modulus at most \(P_*(z)\). Hence
\begin{equation}\label{eq:p8c1-endpoint}
 \Re\beta(C_wH,H)\le n_c(w)\sum_\nu\beta(H_\nu,H_\nu)=n_c(w)\beta(H,H),
\end{equation}
and the same bound holds when the letters of \(w\) are distributed between the two arguments.
\leanref{Formalism.\allowbreak{}FourierDecayAlt.\allowbreak{}P.\allowbreak{}CutBP.\allowbreak{}limBd\_\allowbreak{}of\_\allowbreak{}eigStr}{Formalism.\allowbreak{}HermitianCurvature.\allowbreak{}HermitianCurvatureProofs.\allowbreak{}CutLimBdPf}{179--223}

Put \(H_\epsilon=\chi(G_{\rm loc}/\epsilon)H\) for a smooth cutoff \(\chi\) vanishing near zero and equal to one above the cutoff interval. In the diagonal chart with coordinates \(c_{\rm diag}\), unfolding the top Capelli-weighted form over letter indices \(e'\) and compact branches \(\rho\) with diagonal weight \(d_{e',\rho}(c_{\rm diag})^2\) and applying the Leibniz rule in \(G\partial_G\) (\(G=G_{\rm loc}(c_{\rm diag})>0\)) produces summands \((G\partial_G)^q\chi(G/\epsilon)J_\rho(L_{e',q}(G_{\rm loc}^{p_0}F);c_{\rm diag})\) (\(0\le q\le Q_{e'}\)) with bounded Euler factors. Using \(G^{2p_0-d}\le G^{2p_0-N}(1+G^{2N})\) for \(d\le N\) and adjoining \(G^ND\) to each differential coefficient operator \(D\) yields an \(\epsilon\)-independent majorant
\begin{equation}\label{eq:p8c1-cut-domination}
 M\sum_D|G_{\rm loc}(c_{\rm diag})|^{2p_0-N}E(DF;c_{\rm diag}).
\end{equation}
The opposite-sign spectral gap supplies rapid invariant-derivative decay of \(F\) that covers the differential loss of order \(O((r+1)n(N+1))\), the \(mp_0\) steps of \(G_{\rm loc}^{p_0}\) and the \(20(r+1)n^3\) section margin. Thus \(|G_{\rm loc}|^{2p_0-N}E(DF;c_{\rm diag})\le E(DF;c_{\rm diag})+E(G_{\rm loc}^{p_0}DF;c_{\rm diag})\) is integrable, and dominated convergence passes from \(H_\epsilon\) to \(H\).
\leanref{Formalism.\allowbreak{}FourierDecayAlt.\allowbreak{}P.\allowbreak{}CutBP.\allowbreak{}nfDomPtwS\_\allowbreak{}of\_\allowbreak{}cut}{Formalism.\allowbreak{}HermitianCurvature.\allowbreak{}HermitianCurvatureProofs.\allowbreak{}CutNFDomS}{15--143}
\leanref{Formalism.\allowbreak{}FourierDecayAlt.\allowbreak{}P.\allowbreak{}CutBP.\allowbreak{}dominated\_\allowbreak{}L2\_\allowbreak{}operator\_\allowbreak{}bound\_\allowbreak{}holds}{Formalism.\allowbreak{}HermitianCurvature.\allowbreak{}HermitianCurvatureProofs.\allowbreak{}DominatedL2OperatorBoundHolds}{228--298}
\leanref{Formalism.\allowbreak{}FourierDecayAlt.\allowbreak{}P.\allowbreak{}CutBP.\allowbreak{}CompactFormNormalFormDomination.\allowbreak{}nfTopDomAllJ\_\allowbreak{}of\_\allowbreak{}cut}{CompanionTheory.\allowbreak{}AutomorphicTheory.\allowbreak{}CompactTopologicalForms.\allowbreak{}FinalS}{19--50}

For compactly supported approximants and any integer \(K\ge1\), let \(\mathcal W_w(\omega)\) be the Fourier transform in an appended singleton coordinate of the level-\(j\) partial pairing at wave multiplier \(\omega\), and let \(\mathcal R_w(t,t')\) be the partial pairing with last blocks \((s,t),(s,t')\). Repeated integration by parts in \(K\) rows of grid nodes \(\mathbf y=(y_{a j_a})_{a=1}^K\) from Lemma~\ref{B8.10ca} produces polynomials of degrees \(2j\) and \(2(j-s)+s\) whose leading Lagrange identities are
\[
 \begin{aligned}
 \omega^K\mathcal W_w(\omega)
   &=\sum_{j_1,\ldots,j_K}\Bigl(\prod_{a=1}^K L_{a j_a}\Bigr)
                     \mathcal W^{+,-}_{w,\mathbf y}(\omega),\\
 [(-1)^{s-1}(t^{-1}-{t'}^{-1})]^K\mathcal R_w(t,t')
   &=\sum_{j_1,\ldots,j_K}\Bigl(\prod_{a=1}^K L_{a j_a}\Bigr)
                     \mathcal R^{+,-}_{w,\mathbf y}(t,t').
 \end{aligned}
\]
Both identities are proved in the archive without further hypotheses: the wave identity from the slice-polynomial theorem, and the lift identity from the normalized left Fubini tower and the smooth right average.
\leanref{Formalism.\allowbreak{}FourierDecayAlt.\allowbreak{}P.\allowbreak{}degree\_\allowbreak{}of\_\allowbreak{}plane\_\allowbreak{}wave}{Formalism.\allowbreak{}HermitianCurvature.\allowbreak{}HermitianCurvatureProofs.\allowbreak{}PBGWaveMain}{55--56}
\leanref{Formalism.\allowbreak{}FourierDecayAlt.\allowbreak{}P.\allowbreak{}degree\_\allowbreak{}of\_\allowbreak{}lift\_\allowbreak{}pairing}{Formalism.\allowbreak{}HermitianCurvature.\allowbreak{}HermitianCurvatureProofs.\allowbreak{}PBGDegLift}{4--6}
Take \(K=3\), pass to the cut coefficient by \eqref{eq:p8c1-cut-domination}, and apply Cauchy--Schwarz with Lemma~\ref{B8.10ca} when \(u\ne u'\) to bound either step multiplier \(d_1\) by
\begin{equation}\label{eq:p8c1-wave-decay}
 \min\{1,(\Lambda/|d_1|)^3\}
\end{equation}
times an average of normalized diagonal energies on the enlarged letter set. Fourier inversion on singletons (\(\int_0^\infty\min(1,(\Lambda/t)^3)t\,dt=3\Lambda^2/2\)) and block shortening via the boundary-limit, Herglotz and planar-sum identities \eqref{eq:p8c1-boundary-limit}, \eqref{eq:p8c1-herglotz} and \eqref{eq:p8c1-planar-sum}, which are proved in the proof of Lemma~\ref{B8.10i} without using the present proposition, cost \((C_{\rm step}\Lambda^2)^{r+1}\) (for an absolute constant \(C_{\rm step}>0\)) across the \(r+1\) positions and reach \eqref{eq:p8c1-endpoint} with at most \(N\) letters.
\leanref{Formalism.\allowbreak{}FourierDecayAlt.\allowbreak{}KB.\allowbreak{}DegWave}{Formalism.\allowbreak{}HermitianCurvature.\allowbreak{}RealInputsClosedAdmissibleSettings}{155--167}
\leanref{Formalism.\allowbreak{}FourierDecayAlt.\allowbreak{}KB.\allowbreak{}DegLift}{Formalism.\allowbreak{}HermitianCurvature.\allowbreak{}RealInputsClosedAdmissibleSettings}{169--178}
\leanref{Formalism.\allowbreak{}FourierDecayAlt.\allowbreak{}P.\allowbreak{}TBBx5.\allowbreak{}integral\_\allowbreak{}three\_\allowbreak{}at\_\allowbreak{}place\_\allowbreak{}of\_\allowbreak{}parts}{Formalism.\allowbreak{}HermitianCurvature.\allowbreak{}HermitianCurvatureProofs.\allowbreak{}OnePointBasisFundamentalDomainVolume}{765--914}
\leanref{Formalism.\allowbreak{}FourierDecayAlt.\allowbreak{}P.\allowbreak{}TBBx5.\allowbreak{}integral\_\allowbreak{}four\_\allowbreak{}at\_\allowbreak{}place\_\allowbreak{}of\_\allowbreak{}parts}{Formalism.\allowbreak{}HermitianCurvature.\allowbreak{}HermitianCurvatureProofs.\allowbreak{}OnePointBasisFundamentalDomainVolume}{521--662}

When \(u=u'\), put \(n_0=30(r+1)^2+6\). The consecutive gap \(u_l-u_{l+1}\ge2n_0+1\) implies \(n_0\le b\), since \(u_l-u_{l+1}\le2b\). Apply Lemma~\ref{B8.10hba} to the right compact type at \(z'\). The perpendicular flag vector \(\omega_\gamma\) at a represented positive-Gram point has squared norm \(G_{\rm loc}(\gamma)/T_q\), so the lemma gives
\begin{equation}\label{eq:p8c1-shift-norm}
 \|v_1\|^2\le(3b)^{2n_0}(T_q/G_{\rm loc}(\gamma))^{n_0}
 \bigl(\|A_+(\omega_\gamma)^*v_1\|^2+\|A_-(\omega_\gamma)^*v_1\|^2\bigr).
\end{equation}
The corresponding orbital contraction identities are
\[
 J(FA_\pm;\gamma)=J(F;\gamma)A_\pm(\omega_\gamma),\qquad
 \|A_\pm(\omega_\gamma)\|^2\le(G_{\rm loc}(\gamma)/T_q)^{n_0}.
\]
Here \(FA_\pm\) denotes contraction of the coefficient with the transported degree-\(n_0\) flag map. Apply \eqref{eq:p8c1-shift-norm} to \(J(F;\gamma)^*e_i\) and sum over an orthonormal target basis. This gives the sum of the two shifted Hilbert--Schmidt energies. The equal-weight comparison uses the determinant expansion for the original compact parameters \(u,u'\) through order \(2p_0+n_0\), together with the one-index interlacing variant of Lemma~\ref{B8.10ca} for the shifted compact types.
\leanref{Formalism.\allowbreak{}FourierDecayAlt.\allowbreak{}P.\allowbreak{}CutBP2.\allowbreak{}AmapsSh}{Formalism.\allowbreak{}HermitianCurvature.\allowbreak{}HermitianCurvatureProofs.\allowbreak{}NormSqLeShifted}{39--52}
\leanref{Formalism.\allowbreak{}FourierDecayAlt.\allowbreak{}P.\allowbreak{}CutBPEq.\allowbreak{}CompactFormNormalFormDomination.\allowbreak{}thmBpAt\_\allowbreak{}eqR\_\allowbreak{}top}{CompanionTheory.\allowbreak{}AutomorphicTheory.\allowbreak{}CompactTopologicalForms.\allowbreak{}NormalFormTopDomination}{50--77}

For the minor power \(|D_l(\gamma)|^{2e}\), put \(s=r+1\) and \(n_0=30s^2+6\), so that \(m=2s\), \(p_0=2s\), and \(q_0=se+40s^3+s(s+1)\). For a coefficient with the original compact parameters \(u,u'\), spectral displacement \(\|\nu_0\|_1\le se\), and derivative order \(a_1\le2p_0+n_0\), one has
\[
 2p_0+n_0=4s+30s^2+6\le40s^2,\qquad
 \|\nu_0\|_1+sa_1+s(s+1)\le q_0.
\]
Also \(e\le C_am^3\) gives \(q_0+1\le(C_a/2+3)m^4\). Choose the fixed threshold at least \(\max(10^6,100(C_a/2+3))\). Put \(\Delta_\lambda=C_gm^4\) and \(\Delta_\pm=C_gm^{9/2}(3sn^2+n^3)\). Both gap parameters are at least \(100(q_0+1)\), and the splitting scale is at least \(10^6n^7\). The quarter-gap estimates for the Bessel period construction follow from
\[
 \Delta_\pm/4\ge120n^3,\qquad
 \frac{\Delta_\pm}{32\sqrt m}
 =\frac{C_gm^4(3sn^2+n^3)}{32}\ge125m^2n^3,
\]
which bound the growth losses in that construction. The coordinate \(\lambda_k\) gives
\[
 b\ge|\lambda_k|\ge b_v\ge10^6n^7\ge\max(2,n).
\]
The interlacing, parity, splitting and gap assumptions therefore give the numerical family conditions for both \(u\) and \(u'\).
\leanref{CompanionTheory.\allowbreak{}WaveGrid.\allowbreak{}GapSizes.\allowbreak{}gapOK\_\allowbreak{}of}{CompanionTheory.\allowbreak{}AutomorphicTheory.\allowbreak{}BesselWaveGrids.\allowbreak{}GapSizes}{54--122}

Fix one such coefficient \(F\), at spectral displacement \(\nu_0\). Write \(R_0\) for its given representation, of parameter \(\lambda+\nu_0\). For every \(\nu\in\Z^n\) with \(\|\nu\|_1\le q_0\), choose a discrete-series representation \(R_\nu\) of parameter \(\lambda+\nu\), with \(R_{\nu_0}=R_0\). The numerical family conditions give regular integral parameters, the required Cayley signature and the retained opposite-sign gap at every such \(\nu\). The discrete-series existence input supplies the other members. At indices outside this ball choose \(R_0\). Thus the starting representation is preserved exactly.
\leanref{Gp91.\allowbreak{}WG.\allowbreak{}family\_\allowbreak{}through\_\allowbreak{}discreteSeries}{CompanionTheory.\allowbreak{}AutomorphicTheory.\allowbreak{}BesselWaveGrids.\allowbreak{}WaveBallEstimateWeightedDualAlt}{70--96}

On the index ball \(\|\nu\|_1\le q_0\), choose the period-normalized Bessel maps \(\ell_\nu,\ell'_\nu\) and positive formal degrees \(d_\nu\). The coefficient-period identity assigns a scalar \(\alpha\) to a coefficient of \(R_\nu\); its squared big-cell energy is
\[
 I(F)=c_1\frac{|\alpha|^2}{d_\nu},\qquad c_1>0,
\]
with \(c_1\) depending on the local Bessel datum and its measures, uniformly in \(\nu\). The orbital vector of that coefficient is \(\alpha\) times its period-normalized Bessel orbital vector. Hence \(I(F)=0\) implies that the orbital vector vanishes at every represented positive-Gram big-cell point, and the required expansion is zero in this case.

Suppose \(I(F)>0\), and let \(\alpha_0\ne0\) be its period scalar. Set
\[
 t=\frac{d_{\nu_0}}{|\alpha_0|}>0,\qquad F_{\nu_0}=tF.
\]
For each remaining index in this ball where both Bessel maps are nonzero on smooth vectors, choose \(F_\nu\) with period scalar \(\alpha_\nu\) of modulus \(d_\nu\). Set \(F_\nu=0\) when either map is zero or \(\|\nu\|_1>q_0\). At the starting index both maps are nonzero because \(\alpha_0\ne0\). Each coefficient is continuous, and the big-cell energies are measurable. For \(\|\nu\|_1\le q_0\), the formal-degree formula for the fixed Haar measure is
\[
 d_\nu=c_2\mathcal D(\lambda+\nu),\qquad
 \mathcal D(\mu)=\frac{\prod_{1\le i<j\le n}|\mu_i-\mu_j|}
                         {\prod_{1\le i<j\le n}(j-i)},\qquad c_2>0.
\]
Consequently, for \(\|\nu\|_1\le q_0\), either \(F_\nu=0\) and \(I_\nu=0\), or
\[
 I_\nu:=I(F_\nu)=c_1c_2\mathcal D(\lambda+\nu)>0.
\]
\leanref{Gp91.\allowbreak{}WG.\allowbreak{}coef\_\allowbreak{}norm\_\allowbreak{}finite\_\allowbreak{}and\_\allowbreak{}zeroW}{CompanionTheory.\allowbreak{}AutomorphicTheory.\allowbreak{}BesselWaveGrids.\allowbreak{}WaveBallEstimateWeightedDualAlt}{376--405}
\leanref{Gp91.\allowbreak{}WG.\allowbreak{}exists\_\allowbreak{}coef\_\allowbreak{}family}{CompanionTheory.\allowbreak{}AutomorphicTheory.\allowbreak{}BesselWaveGrids.\allowbreak{}WaveBallEstimateWeightedDualAlt}{97--138}
\leanref{Gp91.\allowbreak{}WG.\allowbreak{}suppliedData\_\allowbreak{}existsW}{CompanionTheory.\allowbreak{}AutomorphicTheory.\allowbreak{}BesselWaveGrids.\allowbreak{}WaveBallEstimateWeightedDualAlt}{351--375}
\leanref{Gp91.\allowbreak{}WG.\allowbreak{}sup\_\allowbreak{}norm}{CompanionTheory.\allowbreak{}AutomorphicTheory.\allowbreak{}BesselWaveGrids.\allowbreak{}WaveBallEstimateWeightedDualAlt}{215--259}
\leanref{Gp91.\allowbreak{}WG.\allowbreak{}sup\_\allowbreak{}meas}{CompanionTheory.\allowbreak{}AutomorphicTheory.\allowbreak{}BesselWaveGrids.\allowbreak{}WaveBallEstimateWeightedDualAlt}{260--267}

For \(\mu=\lambda+\nu\) with \(\|\nu\|_1\le q_0\), every difference \(\mu_i-\mu_j\), \(i<j\), is an integer at least two. Therefore
\[
 \mathcal D(\mu)>0,\qquad
 \mathcal D(\mu+\mathbf1_A)\le n^{|A|}\mathcal D(\mu),
\]
and cancellation of the factors not involving \(i\) gives
\[
 \frac{\mathcal D(\mu-\mathbf1_i)}{\mathcal D(\mu)}
 =\prod_{j\ne i}\frac{|\mu_i-1-\mu_j|}{|\mu_i-\mu_j|}.
\]
For the upper bound, add the coordinates in \(A\) one at a time. At each step the factors involving \(j<i\) decrease. The product of the factors involving \(j>i\) is at most
\(\prod_{a=1}^{n-i}(1+1/a)=n-i+1\le n\), because the intermediate coordinates remain strictly decreasing with integral differences. Multiplication proves the bound. Thus, whenever \(I_\nu>0\) and both indices lie in the ball,
\[
 I_{\nu+\mathbf1_A}\le n^{|A|}I_\nu
 \le(\exp(1)n)^s I_\nu\qquad(|A|=s),
\]
including the case \(F_{\nu+\mathbf1_A}=0\).
\leanref{GlobalTheory.\allowbreak{}FamilyDegreeRatio.\allowbreak{}degree\_\allowbreak{}ratios}{Development.\allowbreak{}AbelianVarietyTheory.\allowbreak{}Proofs.\allowbreak{}JetHarmonics.\allowbreak{}AbsCoordLeOfLOneNormLe}{60--108}

Write \(\mathcal J_\nu\) for the big-cell orbital vector of \(F_\nu\), and use the common point weight in its squared integral \(I_\nu\). At each index with \(\|\nu\|_1+s(s+1)\le q_0\), the exterior-power decomposition of \(\bigwedge^s\C^n\otimes R_\nu\) has components \(R_{\nu+\mathbf1_A}\), \(|A|=s\), and signs \(\epsilon_{\nu A}\in\{1,-1\}\). Let \(q_z,q_{z'}\) be the exterior products of the two Bessel frames in their determinant order. In a branch where both \(F_\nu\) and \(F_{\nu+\mathbf1_A}\) are nonzero, the projection \(\operatorname{pr}_A\) to the indicated component satisfies
\[
 \ell_{\nu+\mathbf1_A}(\operatorname{pr}_A(q_z\otimes v))
   =\kappa_{\nu A}\ell_\nu(v),\qquad
 \ell'_{\nu+\mathbf1_A}(\operatorname{pr}_A(q_{z'}\otimes v))
   =\kappa'_{\nu A}\ell'_\nu(v)
\]
on smooth vectors. When both coefficients are nonzero, put
\[
 \tau_{\nu A}
 =\epsilon_{\nu A}\kappa'_{\nu A}\overline{\kappa_{\nu A}}
       \frac{\alpha_\nu}{\alpha_{\nu+\mathbf1_A}};
\]
put \(\tau_{\nu A}=0\) if either coefficient is zero. The tensor identity and the two Bessel moment bounds give, at represented positive-Gram points,
\[
 D_s\mathcal J_\nu=\sum_{|A|=s}\tau_{\nu A}\mathcal J_{\nu+\mathbf1_A},
 \qquad |\tau_{\nu A}|\le H^s,
 \qquad H=\sqrt{H_\square(T_p)H_\square(T_q)},
 \qquad H_\square(T)=\max(1,T)b^{20s}.
\]
These identities include the zero-coefficient branches. For \(C_h^*=2|C_h|\), the norm assumptions imply
\[
 H_\square(T_p),H_\square(T_q)\le b^{(C_h^*+10)m},
 \qquad H\le b^{(C_h^*+10)m}.
\]
\leanref{Gp91.\allowbreak{}WG.\allowbreak{}exists\_\allowbreak{}tauL}{CompanionTheory.\allowbreak{}AutomorphicTheory.\allowbreak{}BesselWaveGrids.\allowbreak{}WaveBallEstimateWeightedDualAlt}{686--765}
\leanref{Gp91.\allowbreak{}WG.\allowbreak{}supplied\_\allowbreak{}pieriL}{CompanionTheory.\allowbreak{}AutomorphicTheory.\allowbreak{}BesselWaveGrids.\allowbreak{}WaveBallEstimateWeightedDualAlt}{801--849}
\leanref{Gp91.\allowbreak{}WG.\allowbreak{}tau\_\allowbreak{}boundL}{CompanionTheory.\allowbreak{}AutomorphicTheory.\allowbreak{}BesselWaveGrids.\allowbreak{}WaveBallEstimateWeightedDualAlt}{766--800}
\leanref{GlobalTheory.\allowbreak{}TopNum.\allowbreak{}top\_\allowbreak{}num}{Development.\allowbreak{}AbelianVarietyTheory.\allowbreak{}Proofs.\allowbreak{}JetHarmonics.\allowbreak{}TopNum}{25--36}

Every index reached in \(a'\le a_1\) steps has distance at most \(sa'\) from \(\nu_0\). The inequality \(\|\nu_0\|_1+sa_1+s(s+1)\le q_0\) therefore supplies the projection reserve at every stage. Apply Lemma~\ref{B8.10o} with origin \(\nu_0\), norm exponent \(C_h^*\), and the family just constructed. Undo the common positive scalar \(t\) at the starting coefficient. The resulting functions have energy at most \(I(F)\), and their absolute coefficient sum for \(D_s^{a'}\) is at most
\[
 b^{(C_h^*/2+6)m^2a'}=b^{(|C_h|+6)m^2a'}.
\]
\leanref{Gp91.\allowbreak{}WG.\allowbreak{}familyDerivativeExpansion\_\allowbreak{}of\_\allowbreak{}isDS\_\allowbreak{}of\_\allowbreak{}coefOf}{CompanionTheory.\allowbreak{}AutomorphicTheory.\allowbreak{}BesselWaveGrids.\allowbreak{}WaveBallEstimateWeightedDualAlt}{406--511}

Combining the \((2n+1)^{2(r+1)e}\) tensor summands (\(2n+1\le b^2\) from \(b\ge b_v\ge C_gn^7\)), the grid and shift factors, and the two factors bounded by \(\pi^{2m^2}\le b^{4m^5}\) comparing the powers of \(\pi\) in the orbital point-weight and big-cell section-measure normalizations yields \eqref{eq:p8c1-minor-moment}.
\leanref{Formalism.\allowbreak{}FourierDecayAlt.\allowbreak{}P.\allowbreak{}CutBP.\allowbreak{}BigCellDerivativeExpansion}{Formalism.\allowbreak{}HermitianCurvature.\allowbreak{}HermitianCurvatureProofs.\allowbreak{}CutBP}{130--138}
\leanref{Gp91.\allowbreak{}WG.\allowbreak{}bigCellDerivativeExpansion\_\allowbreak{}of\_\allowbreak{}isDS\_\allowbreak{}of\_\allowbreak{}coefOf}{CompanionTheory.\allowbreak{}AutomorphicTheory.\allowbreak{}BesselWaveGrids.\allowbreak{}WaveBallEstimateWeightedDualAlt}{512--525}
\leanref{Formalism.\allowbreak{}FourierDecayAlt.\allowbreak{}P.\allowbreak{}CutBP.\allowbreak{}CompactFormNormalFormDomination.\allowbreak{}thmBpAt\_\allowbreak{}top}{CompanionTheory.\allowbreak{}AutomorphicTheory.\allowbreak{}CompactTopologicalForms.\allowbreak{}NormalFormTopDomination}{14--35}
\leanref{Formalism.\allowbreak{}FourierDecayAlt.\allowbreak{}P.\allowbreak{}CutBPEq.\allowbreak{}CompactFormNormalFormDomination.\allowbreak{}thmBpAt\_\allowbreak{}eqR\_\allowbreak{}top}{CompanionTheory.\allowbreak{}AutomorphicTheory.\allowbreak{}CompactTopologicalForms.\allowbreak{}NormalFormTopDomination}{50--77}
\leanref{CompanionTheory.\allowbreak{}WaveGrid.\allowbreak{}WGlue.\allowbreak{}wLoc1N\_\allowbreak{}of}{CompanionTheory.\allowbreak{}AutomorphicTheory.\allowbreak{}BesselWaveGrids.\allowbreak{}WeightedUniformRotatedOfExpansion}{30--160}
The complete proof is contained in \Lean{CompanionTheory.\allowbreak{}WaveGrid.\allowbreak{}WGlue.\allowbreak{}wLoc1N\_\allowbreak{}of}.
\end{proof}

\subsection*{Proof of \ref{B8.10hc}}\label{prf:B8.10hc}\phantomsection
\begin{proof}[Proof (outline)]
The rank-one torus comparison holds at every threshold \(C_{gf}\ge100\); take \(C_{gf}=100\), so that \(C_g\ge C_{g,1}\ge1000\ge C_{gf}\) supplies its opposite-sign gap hypothesis.
\leanref{Formalism.\allowbreak{}MainStatements.\allowbreak{}RankOne.\allowbreak{}one\_\allowbreak{}bessel\_\allowbreak{}torus\_\allowbreak{}bound}{CompanionTheory.\allowbreak{}AutomorphicTheory.\allowbreak{}RankOneTori.\allowbreak{}FrameBesselBoundOfLargeConstants}{223--225}

For a local datum \(L\), write
\[
 \mathcal E_L(F;\gamma)
 =p_{k,r}(T_p,T_q;\gamma)\|J(F;\gamma)\|_{\rm HS}^2,
\]
with the supremum convention of Definition~\ref{B8.10ha}.
First compare two choices of unipotent Haar measure related by
\(dn=s\,dn_{\rm std}\), where \(0<s<\infty\).
The unipotent measure occurs twice in the product measure defining the Bessel operator, so that product measure changes by \(s^2\).
The section operator changes by \(s^2\), and its squared integral satisfies
\[
 I=s^4 I_{\rm std}.
\]
For every admissible orbital quotient measure, multiplication by \(s^{-2}\) gives an admissible quotient measure for the standardized datum. The corresponding orbital operator changes by \(s^2\). The point weight is nonnegative on represented points, and both orbital suprema vanish on unrepresented points. Taking suprema and multiplying by the point weight gives
\[
 \mathcal E_L(F;\gamma)\le s^4\mathcal E_{L_{\rm std}}(F;\gamma).
\]
These two relations transfer the standardized inequality to arbitrary finite positive Haar scale.
\leanref{CompanionTheory.\allowbreak{}WaveGrid.\allowbreak{}Ties.\allowbreak{}pairMeas\_\allowbreak{}smul}{CompanionTheory.\allowbreak{}AutomorphicTheory.\allowbreak{}BesselWaveGrids.\allowbreak{}NormalizedMeasureScaling}{97--107}
\leanref{CompanionTheory.\allowbreak{}WaveGrid.\allowbreak{}Ties.\allowbreak{}orbSq\_\allowbreak{}le}{CompanionTheory.\allowbreak{}AutomorphicTheory.\allowbreak{}BesselWaveGrids.\allowbreak{}NormalizedMeasureScaling}{108--127}
\leanref{CompanionTheory.\allowbreak{}WaveGrid.\allowbreak{}Ties.\allowbreak{}jNorm2\_\allowbreak{}eq}{CompanionTheory.\allowbreak{}AutomorphicTheory.\allowbreak{}BesselWaveGrids.\allowbreak{}NormalizedMeasureScaling}{128--142}

Set \(\kappa=\kappa_0^{-1}\) and \(\tau=\kappa_0^r\).
The torus comparison supplies a datum \(L_1\), a compact-finite coefficient \(F_1\), a section, and \(K_1>0\). Its spectral parameter and compact weights equal those of \(L\), its character parameter is \(-2i\), and its final norms are
\[
 T_{p,1}=\tau^2T_p,\qquad T_{q,1}=\tau^2T_q.
\]
For its squared section integral \(I_1\) and the transformed block point \(\gamma_1\), the comparison gives
\[
 I_1=K_1I,\qquad
 K_1\mathcal E_L(F;\gamma)
 \le F_*(\kappa_0)\mathcal E_{L_1}(F_1;\gamma_1),
\]
\[
 G_{{\rm loc},1}(\gamma_1)
 =\kappa^{2r(r+1)}\tau^{4(r+1)}G_{\rm loc}(\gamma).
\]
Here the same \(K_1\) includes the square of the nonzero scalar in the section and orbital transport, together with the section Jacobian and weight factors.
\leanref{Formalism.\allowbreak{}MainNotions.\allowbreak{}OnePlaceExpansion.\allowbreak{}frame\_\allowbreak{}bound\_\allowbreak{}of\_\allowbreak{}bessel\_\allowbreak{}torus}{Formalism.\allowbreak{}SharpRestriction.\allowbreak{}CuspSpanNormBounds.\allowbreak{}OnePlaceExpansion}{71--116}

The hypotheses give \(T_{p,1},T_{q,1}\ge1\), and
\[
 T_{p,1},T_{q,1}
 \le b^{|C_h|(m+2r)}\le b^{2|C_h|m}.
\]
The leading minors satisfy
\[
 |D_l(\gamma)|\le b^{|C_h|m^2}|D_l(\gamma_1)|
 \quad(1\le l\le r+1).
\]
\leanref{CompanionTheory.\allowbreak{}WaveGrid.\allowbreak{}WGlue.\allowbreak{}dpt\_\allowbreak{}le}{CompanionTheory.\allowbreak{}AutomorphicTheory.\allowbreak{}BesselWaveGrids.\allowbreak{}WeightedUniformRotatedOfExpansion}{188--217}
Moreover,
\[
 g_*(\kappa_0)
 =(\kappa_0^{-1})^{2r(r+1)}(\kappa_0^r)^{4(r+1)}
 =\kappa_0^{2r(r+1)}.
\]
If \(Q=\max(\kappa_0,\kappa_0^{-1})\), then
\[
 F_*(\kappa_0)=Q^{2m^2+2rm}
 \le b^{|C_h|(2m^2+2rm)}
 \le b^{4|C_h|m^2}.
\]
\leanref{Formalism.\allowbreak{}MainNotions.\allowbreak{}TheoremBConstants.\allowbreak{}GsC}{Formalism.\allowbreak{}SharpRestriction.\allowbreak{}CuspSpanNormBounds.\allowbreak{}TheoremBConstants}{133--134}
\leanref{Formalism.\allowbreak{}MainNotions.\allowbreak{}TheoremBConstants.\allowbreak{}num\_\allowbreak{}Fr}{Formalism.\allowbreak{}SharpRestriction.\allowbreak{}CuspSpanNormBounds.\allowbreak{}TheoremBConstants}{796--827}

Apply Proposition~\ref{B8.10hb} to \(L_1,F_1,\gamma_1\), using the Haar-scale transfer above. Multiply its bound by
\(b^{2|C_h|m^2e}F_*(\kappa_0)\), substitute the minor comparison and \(I_1=K_1I\), and cancel \(K_1>0\). This gives the asserted inequality with \(g_*(\kappa_0)=\kappa_0^{2r(r+1)}\).
The complete proof is contained in \Lean{CompanionTheory.\allowbreak{}WaveGrid.\allowbreak{}WGlue.\allowbreak{}w\_\allowbreak{}in\_\allowbreak{}frame}.
\end{proof}

\subsection*{Proof of \ref{B8.10h}}\label{prf:B8.10h}\phantomsection
\begin{proof}[Proof (outline)]
Apply Proposition~\ref{B8.10hb} (established from the cut Gram forms of the Capelli-weighted partial orbital pairings of Definition~\ref{B8.10ha} and the \(K\)-variable wave and lift interpolation identities) with \(C_a=2(1+A_w+10)\) to obtain \(C\ge0\) and \(C_{g,1}\), and set \(C_{g,0}=C_{g,1}\ge1000\), which serves for both Proposition~\ref{B8.10hb} and Lemma~\ref{B8.10hc}. Fix \(C_g\ge C_{g,0}\). Put \(a=10m^2\), \(p=A_w(r+1)^2\), and at each infinite place \(v\) put \(d_v(\gamma)=|D_{r+1}(\gamma)|_v^2\) and \(\gamma_v\) for the rotated local block point of \(\gamma\).

Fix a rational block point \(\gamma\) in the sum and a place \(v\). If \(\|J_v(\gamma)\|_{\rm HS}^2=0\), the \(v\)-th factor vanishes. Otherwise \(\gamma_v\) has a unitary representative in the local Bessel datum of character parameter \(-2i\xi\), so \(G_{\rm loc}(\gamma_v)=\bigl(\prod_{i=0}^{r-1}\operatorname{rd}(E)^{-(r-i)}\bigr)^4 x_{G,v}\ge0\), where \(x_{G,v}=|D_{r+1}(\gamma)|_v^2-TT'|D_r(\gamma)|_v^2=v(G(\gamma))\). Since \(G(\gamma)\ne0\), we have \(0<|G(\gamma)|_v=x_{G,v}\le d_v(\gamma)\) and \(p_v(\gamma)\ge0\).

For \(1\le l\le r+1\), set \(s_{l-1}=\bigl(\prod_{i=0}^{l-1}\operatorname{rd}(E)^{-(r-i)}\bigr)^{-4}\), so \(1+s_{l-1}|D_l(\gamma_v)|^2=1+|D_l(\gamma)|_v^2\). From the coordinate and splitting-scale bounds we have \(b\ge b_v\ge C_gn^7\ge2\), and then \(\xi=\pi/(v_Eg_0g_1)>0\) with \(\xi^{\pm1}\le b^{C_h}\) gives \(C_h\ge0\); the root-discriminant bound \(\operatorname{rd}(E)=2\pi/(\xi g_0g_1)\le b^{3+C_h}\) then gives \(s_{l-1}\le b^{(3+C_h)m^2}\). Expanding \(w_v(\gamma)^2(1+d_v(\gamma))^aW_v(\gamma)^p\) in the block-end minors assigns exponents \(N_l\in\N\) with \(\sum_{l=1}^{r+1}N_l\le(1+A_w+10)m^3\) such that
\[
 w_v(\gamma)^2(1+d_v(\gamma))^aW_v(\gamma)^p
 \le\Bigl(\prod_{l=1}^{r+1}\bigl(1+s_{l-1}|D_l(\gamma_v)|^2\bigr)^{N_l}\Bigr)^2.
\]
Choose \(l_*\) maximizing \(|D_l(\gamma_v)|^2\), and set \(e=0\) when \(|D_{l_*}(\gamma_v)|\le1\) and \(e=2\sum_l N_l\le C_am^3\) otherwise; this bounds the right-hand side by \(\bigl(2b^{(3+C_h)m^2}\bigr)^{2\sum_l N_l}|D_{l_*}(\gamma_v)|^{2e}\le b^{(4+C_h)C_am^5}|D_{l_*}(\gamma_v)|^{2e}\). Applying Lemma~\ref{B8.10hc} at \(\kappa_0=\xi\) and using \(\max\bigl(1,|g_*(\xi)G_{\rm loc}(\gamma_v)|^{-1}\bigr)^{40(r+1)^2}\le b^{10(3+3C_h)m^5}\max(1,|G(\gamma)|_v^{-1})^a\) give
\begin{equation}\label{eq:p8c1-one-place-weight}
 w_v(\gamma)^2(1+d_v(\gamma))^ap_v(\gamma)W_v(\gamma)^p\|J_v(\gamma)\|_{\rm HS}^2
 \le b^{C_*m^5}\max(1,|G(\gamma)|_v^{-1})^aI_v
\end{equation}
with \(C_*=(4+C_h)C_a+2C_hC_a+4C_h+C+10(3+3C_h)\). Since \((1+d_v(\gamma))^{-1}\max(1,|G(\gamma)|_v^{-1})\le|G(\gamma)|_v^{-1}\) and \(\prod_v|G(\gamma)|_v=|N_{F/\Q}(G(\gamma))|\ge1\) for nonzero integral \(G(\gamma)\), taking the product of \eqref{eq:p8c1-one-place-weight} over \(v\) gives \(w(\gamma)^2p(\gamma)W(\gamma)^p\prod_v\|J_v(\gamma)\|_{\rm HS}^2\le b^{C_*m^5f}\prod_v I_v\).

For fixed block lengths \((s_i)\) and block-end minors, the ratios \(D_j/D_{j-s_i}\) determine \(t_i^{s_i}\) up to at most \(\prod_i s_i\le e^{(r+1)/e}\) root choices, across \(2^r\) compositions of \(r+1\). In each nonzero algebraic-integer box \(|v(t)|\le R_v\) with \(U=\prod_v R_v\ge1\), dilating the \(v\)-th plane by \(U^{1/f}/R_v\) separates distinct points by at least \(1\) in supremum norm; disjoint polydisks of radius \(1/2\) fit inside radius \(U^{1/f}+1/2\) and bound the count by \(9^fU^2\). Integrating against \(\prod_v 4R_v(1+R_v^2)^{-3}\,dR_v\) gives at most \(16^f\) per block-end minor, and hence
\begin{equation}\label{eq:p8c1-reciprocal-weight}
 \sum_\gamma w(\gamma)^{-1}
 \le 2^r e^{(r+1)/e}16^{f(r+1)}
 \le64^{f(r+1)}\le b^{6m^5f}.
\end{equation}
Multiplying the product bound by \(w(\gamma)^{-1}\) and summing over \(\gamma\) with \(C_b=C_*+6\) proves the inequality.
\leanref{CompanionTheory.\allowbreak{}WaveGrid.\allowbreak{}WGlue.\allowbreak{}weight\_\allowbreak{}le}{CompanionTheory.\allowbreak{}AutomorphicTheory.\allowbreak{}BesselWaveGrids.\allowbreak{}WeightedUniformRotatedOfExpansion}{275--292}
\leanref{Formalism.\allowbreak{}MainNotions.\allowbreak{}TheoremBConstants.\allowbreak{}weights\_\allowbreak{}le}{Formalism.\allowbreak{}SharpRestriction.\allowbreak{}CuspSpanNormBounds.\allowbreak{}TheoremBConstants}{636--665}
\leanref{Formalism.\allowbreak{}MainNotions.\allowbreak{}TheoremBConstants.\allowbreak{}sum\_\allowbreak{}weighted\_\allowbreak{}point\_\allowbreak{}counts\_\allowbreak{}le\_\allowbreak{}cube}{Formalism.\allowbreak{}SharpRestriction.\allowbreak{}CuspSpanNormBounds.\allowbreak{}TheoremBConstants}{957--976}
\leanref{Formalism.\allowbreak{}MainNotions.\allowbreak{}TheoremBConstants.\allowbreak{}num\_\allowbreak{}G}{Formalism.\allowbreak{}SharpRestriction.\allowbreak{}CuspSpanNormBounds.\allowbreak{}TheoremBConstants}{854--905}
\leanref{Formalism.\allowbreak{}MainNotions.\allowbreak{}BesselShellMomentBounds.\allowbreak{}theorem\_\allowbreak{}B\_\allowbreak{}weighted\_\allowbreak{}bound\_\allowbreak{}one\_\allowbreak{}place}{Formalism.\allowbreak{}SharpRestriction.\allowbreak{}CuspSpanNormBounds.\allowbreak{}BesselShellMomentBounds}{992--996}
\leanref{Formalism.\allowbreak{}MainNotions.\allowbreak{}BesselShellMomentBounds.\allowbreak{}theorem\_\allowbreak{}B\_\allowbreak{}second\_\allowbreak{}bound\_\allowbreak{}64}{Formalism.\allowbreak{}SharpRestriction.\allowbreak{}CuspSpanNormBounds.\allowbreak{}BesselShellMomentBounds}{265--286}
The complete proof is contained in \Lean{CompanionTheory.\allowbreak{}WaveGrid.\allowbreak{}weighted\_\allowbreak{}bound\_\allowbreak{}rotated\_\allowbreak{}uniform}.
\end{proof}

\subsection*{Proof of \ref{B8.10j}}\label{prf:B8.10j}\phantomsection
\begin{proof}
Projection by the \(\rho\) pairings with the orthonormal vectors splits \(U\) as \(O_E^\rho\) and its orthogonal complement. Expand \(q^a{\bf1}_Z\) in the additive characters of \(O/\varpi^a O\) evaluated on \(h(u,u)\). The trivial character contributes \(\mathbb E\chi={\bf1}_{\chi=1}\).

Consider a character of conductor \(\varpi^i\), \(1\le i\le a\). In one orthonormal coordinate its normalized sum has the form
\[
 S=\mathbb E_x\,\psi(x\bar x)\chi_0(x),
 \qquad x\in O_E/\varpi^N O_E ,
\]
where \(\chi_0\) is the linear character in that coordinate. Expand \(|S|^2\) using the difference \(y\) of the two summation variables. The inner average in \(x\) is the character average for \(\psi(\operatorname{Tr}_{E/F}(x\bar y))\). The trace pairing on \(O_E\) is perfect over \(O\), also at residue characteristic two. The average vanishes unless \(y\in\varpi^iO_E\). The proportion of such \(y\) is \(q^{-2i}\), so \(|S|\le q^{-i}\). Taking the product over the \(\rho\) coordinates and bounding the average on the complement by one gives \(q^{-i\rho}\). There are \(q^i-q^{i-1}\) characters of conductor \(\varpi^i\). Their total contribution is at most
\[
 \sum_{i\ge1}q^{-i(\rho-1)}
 \le \frac87 q^{1-\rho},
\]
since \(q^{1-\rho}\le1/8\).

If \(h(u,W)\) is proper, all the \(\rho\) coordinates belong to one maximal ideal \(\pi O_E\). In the field case take \(\pi=\varpi\), with \(\pi\bar\pi=\varpi^2\). In the split case the possibilities are \((\varpi,1)\) and \((1,\varpi)\), each with norm \(\varpi\). Write \(j\in\{1,2\}\) for the valuation of this norm. Restriction to the specified ideal in all coordinates costs \(q^{-j\rho}\). Conditional on that restriction, write \(x=\pi z\) with \(z\) uniform. The norm character of conductor \(i\) then has conductor exponent \((i-j)_+\). The same square calculation bounds the conditional product by \(q^{-(i-j)_+\rho}\). Including the trivial character, the sum of these bounds with multiplicities is at most
\[
 q^j+\sum_{i>j}q^i q^{-(i-j)\rho}
 \le (8/7)q^j .
\]
After the restriction cost, each ideal contributes at most \((8/7)q^{j(1-\rho)}\) to \(q^a\) times the probability. There are at most two ideals.
\end{proof}

\subsection*{Proof of \ref{B8.10n}}\label{prf:B8.10n}\phantomsection
\begin{proof}
If a solution exists, associate to it the integral ideal \(I=(x)J^{-1}\). Then
\[
             I\bar I=(a)(J\bar J)^{-1}=:\mathfrak b .
\]
In particular \(\mathfrak b\) is integral and \(I\) divides \(\mathfrak b\). If \(\mathfrak p^h\Vert\mathfrak b\), there are at most \(h+1\le (N\mathfrak p)^h\) possibilities for the exponent in \(I\). Thus the number of possible \(I\)'s is at most \(N\mathfrak b=|N_{E/\Q}a|/(NJ)^2\).

Two solutions giving the same \(I\) have quotient an integral unit \(u\) with \(u\bar u=1\). Every conjugate of \(u\) has absolute value one. The characteristic polynomials over \(\Q\) of its positive powers have integral coefficients bounded by the binomial coefficients of order \(2f\). They therefore have only finitely many roots in total. Two powers coincide, so \(u\) is a root of unity. If \(w\) is the number of roots of unity of \(E\), then \(\varphi(w)\le2f\), since a finite subgroup of the multiplicative group of a field is cyclic. For \(p^e\Vert w\) the ratio \(\varphi(p^e)^2/p^e=p^{e-2}(p-1)^2\) is at least one except for \(p^e=2\), when it equals \(1/2\). Hence \(w\le2\varphi(w)^2\le8f^2\). This proves the first estimate, also for fractional \(J\).

Apply it to each coordinate with \(a=1\). Injectivity gives the upper bound \(64f^4(TT')^{2f}P^4\). Since \(f\ge1\), induction gives \(8f^2\le81^f\): it holds at \(f=1\), and \((f+1)^2/f^2\le4\). Thus \(64f^4\le9^{4f}\le(TT')^{4f}\). Together with \(P^4\le P^6\) this proves the second estimate.
\end{proof}

\subsection*{Proof of \ref{B8.10dba}}\label{prf:B8.10dba}\phantomsection
\begin{proof}
Continuity is preserved by multiplication by \(a_w\). Scaling the two measures gives
\[
 \int_{q_w^{-1}(K)}\Phi_w\,d\mu_w
 =c_w\int_K f_w\,d\nu_w^0
 =\int_K (c_w/e_w)f_w\,d(e_w\nu_w^0).
\]
For each finite set \(S\) of finite places,
\[
 \prod_{v\in S}g_v(T_v)
 =\left(\prod_{v\in S}a_v\right)
  \left(\prod_{v\in S}f_v(T_v)\right).
\]
Taking the limit proves convergence and identifies the product. Multiplication by \(g_\infty(T_\infty)=a_\infty f_\infty(T_\infty)\), followed by the assumed product-one identity, proves the last formula.
\end{proof}

\subsection*{Proof of \ref{B8.10db}}\label{prf:B8.10db}\phantomsection
\begin{proof}[Proof (outline)]
Proposition~\ref{B8.02ya} gives archimedean Bessel uniqueness and unit-scale energy convergence on the tied Hida model. For the density construction, let \(E/F\) be CM, let the Hermitian form be nonsingular of rank \(M_0\), and take a nonsingular Gram matrix of rank \(n_0\) with \(2n_0<M_0\). Fix one global \(F\)-basis of \(E\), one admissible global additive character, and the compatible standard local measures. The test function is a compatible pure tensor \(\Phi=\Phi_\infty\Phi_f\): \(\Phi_\infty\) is continuous, rapidly decreasing, smooth and compactly supported; every finite-place factor \(\Phi_v\) is locally constant and compactly supported, almost all are the indicator of the integral coordinate lattice, and \(\Phi_f\) is their product. The coordinate Haar measures on the tuple space and the Hermitian matrix space formed from these choices admit Gram-fibre measures \(m_{w,S}\). For a continuous compactly supported test function \(\Phi_w\), their fibre integrals
\[
 f_w(S)=\int\Phi_w(X)\,dm_{w,S}(X)
\]
are continuous densities on the nonsingular locus. The prescribed source and target measures, self-dual for this same character, are respectively \(c_w\) and \(e_w\) times these coordinate measures. Their positive ratios \(a_w=c_w/e_w\) have convergent finite-place product and satisfy \(a_\infty\prod_v a_v=1\). Thus Lemma~\ref{B8.10dba} gives the normalized densities \(g_w=a_wf_w\) and preserves the product of the local factors. The inequality \(2n_0<M_0\) gives \(n_0=0\) or \(n_0+2\le M_0\), which is the range used for the Euler-product formula for the coordinate fibre measures. Integrating the pure tensor in that formula and applying the product-one identity gives the normalized theta-mean formula with these densities.
\leanref{GlobalTheory.\allowbreak{}CompanionTheory.\allowbreak{}SW.\allowbreak{}TopDens.\allowbreak{}siegelDensities}{CompanionTheory.\allowbreak{}SpectralGeometry.\allowbreak{}LocalGlobalProofs.\allowbreak{}AdelicAnalysis.\allowbreak{}SiegelDensitiesOfHermitianForm}{47--107}

Applying the convergent unitary density formula to the totally positive definite rank-\((k-1)\) orthogonal complement \(W_t^\perp\) expresses \(n_{\rm gen}(\gamma)\) as \(A_{\rm re}G_k(\gamma)\,\mathrm{dens}(\gamma)\) up to local exceptional factors bounded by \(H(\gamma)^{10(r+1)^2}\), compares \(n_p\) with \(n_p^{\rm gen}\) at \(p\in\Omega\) with exponent \(C'_0(r+1)^2+2m\), and combines \(\mathrm{dens}(\gamma)\le e_1H(\gamma)\) with \(t(\gamma)^{-1}\le H(\gamma)^{6m}\) to give \eqref{eq:p8c1-generic-mass-comparison}.
The complete proof is contained in \Lean{GlobalTheory.\allowbreak{}St.\allowbreak{}stepChoiceDataC\_\allowbreak{}of\_\allowbreak{}siegelWeil\_\allowbreak{}inInfinity}.
\end{proof}

\subsection*{Proof of \ref{B8.10dca}}\label{prf:B8.10dca}\phantomsection
\begin{proof}
Both \(P(\gamma)\) and \(Q(\gamma)\) are positive, so \(\vartheta>0\). The factors \(\vartheta\) and \(\vartheta^{-1}\) cancel in the disintegration identity, which proves the quotient-measure assertion. Linearity of the orbital integral in its measure gives
\[
 \mathcal A_\tau(\nu_{\tau,\vartheta})
   =\vartheta^{-1}\mathcal A_\tau(\nu_{\tau,1}),\qquad
 \|\mathcal A_\tau(\nu_{\tau,\vartheta})\|_{\rm HS}^2
   =\vartheta^{-2}\|\mathcal A_\tau(\nu_{\tau,1})\|_{\rm HS}^2.
\]
Since \(\vartheta^{2f}=P(\gamma)/Q(\gamma)\), multiplying over \(\tau\) gives \eqref{eq:p8c1-selected-orbital-factorization}. Multiplication by \(a_2(c,\gamma,\nu)\) proves the last identity.
\end{proof}

\subsection*{Proof of \ref{B8.10dc}}\label{prf:B8.10dc}\phantomsection
\begin{proof}[Proof (outline)]
First suppose \(Q(\gamma)>0\). By Proposition~\ref{B8.02ya}, uniform restriction bounds for the local row functions, and point uniformity between the pinned nongeneric normal-data weight (or equal-word generic weight) and the rotated block-point coordinates of Lemma~\ref{B8.10h}, the local row-function orbital factors admit unit-root quotient measures with positive product \(P(\gamma)\), and the ratio \(\vartheta=(P(\gamma)/Q(\gamma))^{1/(2f)}\) lies in \([b^{-C_1m^2},b^{C_1m^2}]\). Here \(C_1\) is the constant in the local root-measure normalization window. Lemma~\ref{B8.10dca} gives the factorization with the selected measures \(\vartheta^{-1}\nu_{\tau,1}\). The local Bessel comparison gives
\[
 \|\mathcal A_\tau(\nu_{\tau,\vartheta})\|_{\rm HS}^2
 \le K_\tau b^{C_am^2}\|J_\tau(\gamma)\|_{\rm HS}^2.
\]
Multiplication and Lemma~\ref{B8.10dca} give
\[
 Q(\gamma)\le C_J\prod_\tau\|J_\tau(\gamma)\|_{\rm HS}^2.
\]
Since the same pinned measure defines \(a_2=\|V\|^2/Q\), multiplication by \(a_2\ge0\) gives the subtotal estimate on \(Q>0\). If \(Q=0\), the displayed analytic inequality follows from the nonnegativity of its right-hand side.
\leanref{Formalism.\allowbreak{}L2.\allowbreak{}SG.\allowbreak{}Pf.\allowbreak{}BDatR.\allowbreak{}jointAmplitude\_\allowbreak{}bound\_\allowbreak{}rotated}{Formalism.\allowbreak{}LocalToGlobal.\allowbreak{}SiegelLatticeTheory.\allowbreak{}FlagStabilizerTensors.\allowbreak{}RootLatticeThresholds.\allowbreak{}BDataOfR}{46--79}

For every compatible family of big-cell sections, the two section bounds are \((\prod_\tau K_\tau)\prod_\tau I_\tau\le\prod_\tau M_\tau\) and \(\prod_\tau M_\tau\le b^{C_nm^2f}I_{\rm re}\). The constants \(K_\tau\) therefore cancel the same factors in \(C_J\). More precisely, if \(C_b^0\) is the constant from Lemma~\ref{B8.10h} at scale \(b^{C_\beta}\), with \(C_\beta=\max(2,C_w)\) and \(C_w\) the standing coordinate-size constant, the remaining factor is bounded by
\[
 \exp\!\left((\max(0,C_b^0)C_\beta+|C_a|+|C_n|)m^5f\log b\right)I_{\rm re}.
\]
This follows by multiplying the factors \(b^{C_am^2f}\), \(b^{C_b^0 C_\beta m^5f}\), and \(b^{C_nm^2f}\), and using \(m^2\le m^5\). For a real \(A_w\ge0\), the integer exponent \(\lceil A_w\rceil(r+1)^2\) gives the required bound because \(W(\gamma)\ge1\).
The complete proof is contained in \Lean{Formalism.\allowbreak{}MainStatements.\allowbreak{}RowRest.\allowbreak{}analytic\_\allowbreak{}pinnedTheoremB\_\allowbreak{}uniform} and \Lean{Formalism.\allowbreak{}MainStatements.\allowbreak{}RowRest.\allowbreak{}bDat\_\allowbreak{}theoremB\_\allowbreak{}global}.
\end{proof}

\subsection*{Proof of \ref{B8.10d}}\label{prf:B8.10d}\phantomsection
\begin{proof}[Proof (outline)]
Put \(r=r_S\) and \(\kappa=k-2r-5\), with \(k\ge4r+7\) from the rank thresholds. Use the partner classes \(c\), block points \(\gamma\), normal labels \(\nu\), subtotals \(V(c,\gamma,\nu)\) and local defect data of Definition~\ref{B8.10da}. By Lemma~\ref{B8.08n} and absolute summability,
\begin{equation}\label{eq:p8c1-ng-class-squares}
 P(U_{\rm ng})=\Gamma_1\sum_c\left\|\sum_{\gamma,\nu}V(c,\gamma,\nu)\right\|^2,
\end{equation}
and the corresponding class-square identity holds for \(P(U_{\rm sc})\) over the equal-word generic scalar-partner classes.

At a finite place \(p\) with height \(H_p=H_p(\gamma)\), residue cardinality \(q_p\) and unsigned reference mass \(n_p\), the local defect sums of Definition~\ref{B8.10da} satisfy
\begin{equation}\label{eq:p8c1-defect}
 \sum_{d\in D_p} w_d|a_N(d)|^2\le H_p^{39(r+1)^2}q_p^{-\kappa}n_p,\qquad
 \sum_{d\in D_p} w_d|a_Z(d)|^2\le H_p^{39(r+1)^2}n_p.
\end{equation}
For the first bound set \(A_0=H_p^{9(r+1)^2+12(r+1)}\) and \(A_1=4H_p^{4(r+1)}\). At defect \(\beta=q_p^i>1\) the total mass is at most \(A_1\beta^{-(k-1)}n_p\), while image cardinality \(q_p^l\) at fixed data costs at most \(A_0q_p^{l(k-1)}\); summing the proper-image geometric progression gives at most \(2A_0(\beta/q_p)^{k-1}\) and saves \(2A_0A_1q_p^{-(k-1)}\).

In a full-image family at \(\beta=q_p^i>1\), fix the last occupied layer \(J\), its coset count \(\beta_J\), and \(d\), so the image mass is at most \(2A_0\beta^{k-1}\). Write \(\Lambda\) for a starting lattice in the family, \(X_J,X'_J\) for the two flag prefix spans, \(Y=X_J+X'_J\), \(\theta=\theta_p\), \(\varpi\) for a uniformizer of \(F_p\), \(a\ge0\) for the norm-congruence exponent, and \(h\) for the local Hermitian pairing linear in its first variable. For \(b\in X_J\), \(b'\in X'_J\) with \(b+b'\in\theta\Lambda\) of exact level \(J\) modulo the frame lattice, \(v\in U=\Lambda\cap Y^\perp\) with \(h(v,v)\in\varpi^aO\), and trace-zero \(t_0\in E_p\) such that \(c_0=h(v,v)/2+t_0\in\varpi^aO_E\), the map \(T_{b,v,c_0}(z)=z+h(z,v)b-h(z,b)(v+c_0b)\) is a Hermitian isometry with inverse \(T_{b,-v,\bar c_0}\). Giving probability measure to the full parameter domain and multiplying by \(q_p^a\beta/\beta_J\) on targets \(\Lambda'=T_{b,v,c_0}^{-1}(\Lambda)\) remaining in the family defines a symmetric stabilizer-invariant kernel \(m(x,y)\) satisfying
\[
 \left|1-\sum_y m(x,y)\right|\le a_p,\qquad
 \left|\sum_y m(x,y)\omega_d(y)\right|\le a_p,\qquad
 a_p=\tfrac{24}{7}q_p^{1-(k-2r-3)},
\]
for the unit-modulus orbital phase \(\omega_d\). Indeed \(W=\Lambda\cap(F_u+F'_t)^\perp\subseteq U\) is saturated of codimension at most \(2r+2\) in rank \(k+2r+1\), so its modulo-\(\varpi\) reduction has radical of dimension at most \(2r+2\) and retains a nonsingular subspace of dimension at least \(\rho=k-2r-3\ge4\), which lifts by Hensel's lemma and surjective unramified unit norms (or dual split bases) to \(\rho\) orthonormal vectors in \(W\). Moves satisfying \(h(v,W)=O_E\) remain in the family, the \((b,b')\)-constraint has relative measure \(\beta_J/\beta\), the \(t_0\)-coset has measure \(1\), and Lemma~\ref{B8.10j} supplies \((8/7)q_p^{1-\rho}\) for the unrestricted trivial and nontrivial character averages and \((16/7)q_p^{1-\rho}\) for the \(h(v,W)\ne O_E\) restriction. Summing against the relative inverse stabilizer volumes \(\mu(x)\) over stabilizer orbits gives \(\sum_{[x]}\mu(x)\sum_y m(x,y)\omega_d(y)=\sum_{[x]}\mu(x)\omega_d(x)\sum_y m(x,y)\), whence \(|\sum_{[x]}\mu(x)\omega_d(x)|\le2a_p\sum_{[x]}\mu(x)\le7q_p^{1-(k-2r-3)}\sum_{[x]}\mu(x)\). Squaring and summing over the at most \(H_p^{2(r+2)}\) families with \(H_p\ge q_p\ge2\) gives the first inequality of \eqref{eq:p8c1-defect}, while at \(\beta=1\) the singleton-image bound gives the second.
\leanref{GlobalTheory.\allowbreak{}HL.\allowbreak{}exists\_\allowbreak{}symm\_\allowbreak{}invariant\_\allowbreak{}kernel}{Development.\allowbreak{}AbelianVarietyTheory.\allowbreak{}Proofs.\allowbreak{}FlagMeasures.\allowbreak{}ExistsSymmInvariantKernelHermitianFrames}{17--31}
\leanref{GlobalTheory.\allowbreak{}HL.\allowbreak{}FirstType.\allowbreak{}first\_\allowbreak{}type\_\allowbreak{}bound}{Development.\allowbreak{}AbelianVarietyTheory.\allowbreak{}Proofs.\allowbreak{}FlagMeasures.\allowbreak{}HermitianFramesFirstTypeAmpSqBound}{44--112}
\leanref{GlobalTheory.\allowbreak{}HL.\allowbreak{}AmpCancel.\allowbreak{}amp\_\allowbreak{}le\_\allowbreak{}mass}{Development.\allowbreak{}AbelianVarietyTheory.\allowbreak{}Proofs.\allowbreak{}FlagMeasures.\allowbreak{}AmplitudeLeMassSecondTypeFamily}{31--105}
\leanref{GlobalTheory.\allowbreak{}HL.\allowbreak{}FamilyBound.\allowbreak{}family\_\allowbreak{}bound}{Development.\allowbreak{}AbelianVarietyTheory.\allowbreak{}Proofs.\allowbreak{}FlagMeasures.\allowbreak{}HermitianDoubleCosetSecondTypeFamilyAmpSqLe}{59--125}
\leanref{GlobalTheory.\allowbreak{}HL.\allowbreak{}NonzeroDefect.\allowbreak{}ampSq\_\allowbreak{}nonzero\_\allowbreak{}defect}{Development.\allowbreak{}AbelianVarietyTheory.\allowbreak{}Proofs.\allowbreak{}FlagMeasures.\allowbreak{}NonzeroDefectHeightPowerAbsorbsConstants}{49--81}
\leanref{GlobalTheory.\allowbreak{}HL.\allowbreak{}NoCancellation.\allowbreak{}ampSq\_\allowbreak{}le\_\allowbreak{}height\_\allowbreak{}mul\_\allowbreak{}mass}{Development.\allowbreak{}AbelianVarietyTheory.\allowbreak{}Proofs.\allowbreak{}FlagMeasures.\allowbreak{}AmpSqLeHeightMulMass}{112--160}

Expanding the nongeneric finite-place amplitude over nonempty subsets of \(\Omega\) (with \(a_N\) on the subset and \(a_Z\) outside), applying Cauchy--Schwarz at cost \(4^{|\Omega|}\), and inserting \eqref{eq:p8c1-defect} together with the local mass comparison \(n_p\le8H_p^{C'_0(r+1)^2+2m}n_p^{\rm gen}\) of Lemma~\ref{B8.10db} (with \(4\cdot8\le H_p^5\) for \(p\in\Omega\)) bounds the finite-place squared amplitude sum at \(\gamma\) by \(n_{\rm gen}(\gamma)\) times
\begin{equation}\label{eq:p8c1-ng-height}
 2^{-\kappa}H(\gamma)^{e_7},\qquad e_7=(39+C'_0)(r+1)^2+2m+5,
\end{equation}
while for equal-word generic scalar classes Lemma~\ref{B8.08n} gives \(H(\gamma)^{e_7}n_{\rm gen}(\gamma)\) with saving factor \(1\).
\leanref{Formalism.\allowbreak{}FourierDecayAlt.\allowbreak{}P.\allowbreak{}HermitianFormNormLemmas.\allowbreak{}B.\allowbreak{}t\_\allowbreak{}MeanAmpJoint}{Formalism.\allowbreak{}HermitianCurvature.\allowbreak{}HermitianCurvatureProofs.\allowbreak{}UnipotentBlockMeasure}{326--368}

At each \((c,\gamma)\), Lemma~\ref{B8.10n} applied to the two norm-one coordinates of \(\nu\) bounds the number of normal labels by \(P_c N_G(\gamma)^6\), where the class factor \(P_c\) is absorbed into the class mass weight with global factor \((T_uT_v)^{6f}d_0^{12r(r+1)f}\) (\(d_0=\operatorname{rd}(E)\)). Apply Cauchy--Schwarz to the inner sum over \((\gamma,\nu)\) in \eqref{eq:p8c1-ng-class-squares} with weights \(w(\gamma)P_c N_G(\gamma)^6\); by \eqref{eq:p8c1-reciprocal-weight} the reciprocal sum is at most \(64^{f(r+1)}\). Choose \(A_w=39+C'_0+38\), so \(A_w(r+1)^2\ge e_7+e_N+6\) and \(H(\gamma)^{e_7+e_N}N_G(\gamma)^6\le g^{fA_w(r+1)^2}W(\gamma)^{A_w(r+1)^2}\). Combining \eqref{eq:p8c1-ng-height}, Lemma~\ref{B8.10db} \eqref{eq:p8c1-generic-mass-comparison} and Proposition~\ref{B8.10dc} bounds \(P(U_{\rm ng})\) and \(P(U_{\rm sc})\) by \(2^{-\kappa}\Delta_{u,v}\) and \(\Delta_{u,u}\), respectively, times
\begin{equation}\label{eq:p8c1-ng-global-factor}
 64^{f(r+1)}c_{\rm st}e_1d_0^{12r(r+1)f}T_*^{6f}
 g^{fA_w(r+1)^2}b^{C_bm^5f}(d_0/2)^{f(r+1)}g^{f(r+1)(2r+3)},
\end{equation}
where \(T_*=T_uT_v\) and the stabilizer factor satisfies \(c_{\rm st}\le b^{C_{\rm st}m^3f}\). Since \(e_1\le2\), \(d_0\le b\), \(g\le b^{2C_g}\), \(T_*\le b^{2C_Tm}\) and \(b\ge16\), the product \eqref{eq:p8c1-ng-global-factor} is bounded by \(e^{C_4m^5f\log b}\) in the first case and \(e^{C_Bm^5f\log b}\) in the second.
\leanref{GlobalTheory.\allowbreak{}rowNGP\_\allowbreak{}of\_\allowbreak{}siegelWeil\_\allowbreak{}of\_\allowbreak{}stepChoice}{Development.\allowbreak{}AbelianVarietyTheory.\allowbreak{}Proofs.\allowbreak{}FlagMeasures.\allowbreak{}RowPositivityOfSiegelWeilConstants}{11--19}
\leanref{GlobalTheory.\allowbreak{}rowMainEqP\_\allowbreak{}of\_\allowbreak{}siegelWeil\_\allowbreak{}of\_\allowbreak{}stepChoice}{Development.\allowbreak{}AbelianVarietyTheory.\allowbreak{}Proofs.\allowbreak{}FlagMeasures.\allowbreak{}RowMainEqPOfSiegelWeilConstants}{14--23}
\leanref{Formalism.\allowbreak{}FourierDecayAlt.\allowbreak{}P.\allowbreak{}HermitianFormNormLemmas.\allowbreak{}F.\allowbreak{}t\_\allowbreak{}RowNGJoint}{Formalism.\allowbreak{}HermitianCurvature.\allowbreak{}HermitianCurvatureProofs.\allowbreak{}UnipotentBlockMeasure}{1032--1075}
\leanref{Formalism.\allowbreak{}L2.\allowbreak{}SG.\allowbreak{}Pf.\allowbreak{}MainEqCut.\allowbreak{}rowGenJoint\_\allowbreak{}of}{Formalism.\allowbreak{}LocalToGlobal.\allowbreak{}SiegelLatticeTheory.\allowbreak{}CokernelGLCounts.\allowbreak{}RankPointCounts.\allowbreak{}MainEqCut}{259--299}
The complete proof is contained in \Lean{GlobalTheory.\allowbreak{}rowNGP\_\allowbreak{}of\_\allowbreak{}siegelWeil} and \Lean{GlobalTheory.\allowbreak{}rowMainEqP\_\allowbreak{}of\_\allowbreak{}siegelWeil}.
\end{proof}

\subsection*{Proof of \ref{B8.10m}}\label{prf:B8.10m}\phantomsection
\begin{proof}
Put \(z_{iv}=\min(1,s_{iv})\). For \(s>0\) one has \(\min(1,s)/s\le1\) and \(\sqrt{1+s^2}\min(1,s)/s\le\sqrt2\), by considering \(s\le1\) and \(s\ge1\). Therefore
\[
 z_{1v}z_{2v}\frac{c_{1v}+c_{2v}}{s_{1v}s_{2v}}\le 2\sqrt2\le4.
\]
Multiplication over \(v\) gives \((\prod_v z_{1v})(\prod_v z_{2v})\le(4g/\sqrt{TT'})^f\). The smaller of the two products has square at most this bound. Multiply by \(T'^f\) and use \(T'z_{iv}^2=\min(T's_{iv}^2,T')\). For a single side, \(\min(1,s)^2\le s\), so the hypothesized product of \(1/s_v\) gives the conclusion even with \(4\) replaced by \(1\).

Write \(q=(C/k)^{2r(1-\eta)}\). From \(k^{-1}\le C/k\le k^{-2/3}\) and \(3r/2\le2r(1-\eta)\le2r\) we obtain \(k^{-2r}\le q\le k^{-r}\). In particular \(T'a_vg^{-m}q\ge k^{3r}k^{-2r}\ge1\). Consequently
\[
 \min(1,\nu_v a_vg^{-m}q)
 \le \min(T',\nu_v)\,a_vg^{-m}q .
\]
Take the product, use the first conclusion and raise to the power \(k\). The product bound on \(a_v^k\) gives the displayed estimate. Both terms belong to the ordered word pair \((u,u')\), so \(T=T_u,\ T'=T_{u'}\) in either application.
\end{proof}

\subsection*{Proof of \ref{B8.10ac}}\label{prf:B8.10ac}\phantomsection
\begin{proof}[Proof (outline)]
Group each row by the reduced Gram labels \(Z,Z'\) and the rank-\(s\) exceptional lattice. Counting the lattice fibre with inverse stabilizer weights splits orthogonally into ranks \(s\) and \(k-s\) with mass factor \(\Xi_s=M_sM_{k-s}/M_k\) for \(2\le s\le k-2\), and uses the index bound \eqref{eq:p8c1-endpoint-index} at \(s=k-1,k\); index and generator choices in the rank-\(s\) quotient contribute \((Csf)^{6sf}\), while modulo-\(g\) congruences contribute \((2/g)^{(s-1)kf}\) on \(\mathcal P_{{\rm nd},s}\) and \((2/g)^{(s-2)kf}\) on \(\mathcal P_{{\rm deg},s}\). The count and amplitude inputs can be combined quantitatively as follows. Choose the count constant \(c\ge\max(1,C_{Tu})\) and the amplitude constants \(C_a,C_2,C_h\), with \(C_h\ge0\), before the point threshold. Put
\[
 \sigma_Z=S_{C_h}(Z),\qquad
 e_s=c(s+m)m+c(1+\log_2(2f)),\qquad
 P_h=b^{(c^2+2\max(C_h-c,0))m^3f},\qquad
 \zeta=12(r+1)^2 .
\]
For \(x>0\), define
\[
 \widehat\Xi_s(x)=
 \begin{cases}
 M_sM_{k-s}/M_k,&2\le s\le k-2,\\
 (ck)^{kf}(xN)^{ck},&s=k-1,k.
 \end{cases}
\]
Index the nondegenerate and degenerate cases by \(i=1,2\), respectively, and set
\[
 \begin{gathered}
 b_1(Z)=1,\qquad b_2(Z)=T_u^{-f/2}K(Z),\\
 \Delta_i=(2/g)^{(s-i)f}\sqrt{G_ZG_{Z'}}\,b_i(Z)b_i(Z'),\\
 \mathcal B_i(Z,Z')
 =(\Delta_i/N)^k\widehat\Xi_s(h_Zh_{Z'})
          (h_Zh_{Z'})^{e_s}(csf)^{6sf},\\
 a_1=C_2m^2,\qquad a_2=C_2(m^2+2m^4),\\
 \eta_1=2R',\qquad
 \eta_2=e^{2C_2m^6fL_b}(|C_a|/k)^{kf}.
 \end{gathered}
\]
Thus the two determinant cutoffs are
\[
 \Delta_1=(2/g)^{(s-1)f}\sqrt{G_ZG_{Z'}},\qquad
 \Delta_2=(2/g)^{(s-2)f}\sqrt{G_ZG_{Z'}}\,T_u^{-f}K(Z)K(Z').
\]
The mean-count estimate is used in the range
\(1\le\Delta_i\le(h_Zh_{Z'})^3\).
For a nonempty fibre, choose one represented pair. Its determinant factor and its two index factors are each at least one, and their product is at most \(\Delta_i\); hence \(1\le\Delta_i\).
For its upper endpoint, \(g\ge2\), \(T_u\ge1\) and \(0<G_Z\le h_Z\) give
\(\Delta_1\le(h_Zh_{Z'})^{1/2}\).
In the degenerate case the additional input
\(K(Z)^2\le9^fh_Z^3\), together with
\(9^f\le h_Zh_{Z'}\), gives
\[
 \Delta_2\le9^f(h_Zh_{Z'})^2\le(h_Zh_{Z'})^3 .
\]
For a row family of mass \(M\), the finite-fibre input is the weighted count bound \(M\mathcal B_i(Z,Z')\). Let \(A_i(Z)\ge0\) be the orbital amplitude majorants after the row weight and the coefficient dimension have been factored out, and suppose
\[
 n_Z^{(2)}A_i(Z)^2b_i(Z)^{2k}
 \le h_Z^{2cr^2}\Gamma\eta_i
             \sigma_Z^{\,a_i-2q/(r+1)}.
\]
These inputs give a fibre contribution at most
\(M\dim(\pi_u)\mathcal B_i(Z,Z')A_i(Z)A_i(Z')\).
Lemma~\ref{B8.10l}, or direct cancellation of
\[
 \sqrt{n_Z^{(2)}n_{Z'}^{(2)}}
 =\left(\frac{\exp(1)\pi}{kv}\right)^{kf}
       \left(\frac{\sqrt{G_ZG_{Z'}}}{N}\right)^k,
\]
gives
\[
 \begin{split}
 \mathcal B_i(Z,Z')A_i(Z)A_i(Z')
 \le{}&(2/g)^{(s-i)kf}
       \left(\frac{kv}{\exp(1)\pi}\right)^{kf}
       \widehat\Xi_s(h_Zh_{Z'})(h_Zh_{Z'})^{e_s+cr^2}\\
 &{}\cdot(csf)^{6sf}\Gamma\eta_i
       (\sigma_Z\sigma_{Z'})^{\,a_i/2-q/(r+1)} .
 \end{split}
\]
In particular, raising the determinant cutoff to the \(k\)-th power gives exactly the exponents \((s-1)kf\) and \((s-2)kf\).

Each factor defining \(S_c\) is at most
\(b^{\max(C_h-c,0)m^2}\) times the corresponding factor defining \(S_{C_h}\).
There are \(f(r+1)\) factors, so \(m=2r+2\) gives
\[
 1\le h_Z\le P_h\sigma_Z^4,\qquad \sigma_Z\ge1.
\]
The definition of \(\widehat\Xi_s\) then gives
\[
 \widehat\Xi_s(h_Zh_{Z'})
 \le\widehat\Xi_s(P_h^2)(\sigma_Z\sigma_{Z'})^{4ck}.
\]
Consequently, whenever
\[
 4ck+4e_s+4cr^2+\tfrac12a_i+\zeta\le q/(r+1),
\]
the preceding fibre product is at most
\[
 \begin{split}
 \mathcal R_i(s)\sigma_Z^{-\zeta}\sigma_{Z'}^{-\zeta},
 \qquad
 \mathcal R_i(s)={}&(2/g)^{(s-i)kf}
       \left(\frac{kv}{\exp(1)\pi}\right)^{kf}
       \widehat\Xi_s(P_h^2)\\
 &{}\cdot P_h^{2e_s+2cr^2}(csf)^{6sf}\Gamma\eta_i .
 \end{split}
\]
Indeed the remaining exponent of \(\sigma_Z\sigma_{Z'}\) is
\(4ck+4e_s+4cr^2+a_i/2-q/(r+1)\), which is at most \(-\zeta\).
The choice
\[
 q=\left\lceil(r+1)\left(
       4ck+4e_k+4cr^2+
       \tfrac12|C_2|(m^2+2m^4)+\zeta\right)\right\rceil
\]
works for both cases and every \(s\le k\).
The numerical conditions used to bound this moment order are
\[
 r\ge1,\quad m=2r+2,\quad m+3\le k,\quad
 \log f\le\sqrt r,\quad
 |C_2|(m^2+2m^4)\le k .
\]
The last inequality follows from the point threshold
\(3|C_2|m^5fL_b\le k\), since \(m\ge2\), \(f\ge1\), and \(L_b\ge\log2>1/2\).
Put \(c_f=1+\log_2(2f)\). Since \(\log2>1/2\),
\[
 c_f=2+\frac{\log f}{\log2}
 \le2+2\sqrt r\le2+2r=m .
\]
Let \(Q_0\) be the expression inside the ceiling defining \(q\), and put \(\theta=r+1=m/2\).
Using \(c_f\le2\theta\), \(r^2\le\theta^2\),
\(|C_2|(m^2+2m^4)\le k\), \(\theta\ge1\), \(2\theta\le k\) and \(5\le k\), termwise estimation gives
\[
 q\le Q_0+1
 \le(24c+7)k\theta^2
 \le(36c+12)k\theta^2
 =(9c+3)km^2 .
\]
Thus this choice lies in the allowed moment range.

The label-sum input is
\[
 \sum_{Z\in\mathcal Z}\sigma_Z^{-\zeta}
 \le\Phi_0,\qquad
 \Phi_0=(41^fP_h^3)^{(r+1)^2},
\]
for each finite set \(\mathcal Z\) of represented reduced labels. Grouping a finite pair sum by its two labels and applying the fibre estimate therefore gives
\(M\dim(\pi_u)\mathcal R_i(s)\Phi_0^2\).
The numerical comparison with the constants in the statement uses
\[
 c\le C,\quad 1\le C,\quad
 4+(c^2+2\max(C_h-c,0))(6c+2)\le C^2,\quad
 2C_2\le C,\quad |C_a|\le C,
\]
and yields
\(\mathcal R_1(s)\Phi_0^2\le R'U_s\) and
\(\mathcal R_2(s)\Phi_0^2\le V_s\).

 Multiplying by \(\Gamma_1\dim(\pi_u)\operatorname{mass}_u\le f_u\) from Lemma~\ref{B8.08oa} gives \eqref{eq:p8c1-loop-nd-deg}.
The complete proof is contained in \Lean{GlobalTheory.\allowbreak{}RowsFive.\allowbreak{}rowsND\_\allowbreak{}of} and \Lean{GlobalTheory.\allowbreak{}RowsFive.\allowbreak{}rowsDeg\_\allowbreak{}of}.
\end{proof}

\subsection*{Proof of \ref{B8.10ad}}\label{prf:B8.10ad}\phantomsection
\begin{proof}[Proof (outline)]
For a represented frame \((F_u,F'_t)\) of label \(Z\), row-lattice extensions correspond to definite rank-\((k-1)\) complement lattices in \(W_t^\perp\) glued along the frame discriminant module as in Lemma~\ref{B8.10c}; summing inverse stabilizer weights over the genus bounds the weighted count of first frames of label \(Z\) by \(M h_Z^{cr^2}n_Z\). Fixing \(t\), every corank-one partner \(t'\) in the same row lies in \(W_t=\operatorname{span}_E(F_u,F'_t)\) and acts as the identity on \(W_t^\perp\). Integrality in the row lattice confines its switch coordinates in \((F_u,F'_t)\) to congruence classes and archimedean boxes bounded by \(c^f(h_Zh_{Z'})^{cr^2}\), with at most one term per partner datum by injectivity of the basis expansion and uniqueness of the glued lattice on \(W_t\oplus W_t^\perp\). Multiplying the first-frame bound by this partner count gives the first expression in \eqref{eq:p8c1-switch-count}, reversing the roles of \(t\) and \(t'\) gives the second, and taking their square-root product gives \(M c^f(h_Zh_{Z'})^{3cr^2/2}\sqrt{n_Zn_{Z'}}\).

On every nonscalar generic corank-one pair \((t,t')\in S_{\rm nsc}\), the local off-diagonal parameters satisfy the product hypothesis of Lemma~\ref{B8.10m}. Split \(U\) into the subset where the first term satisfies the small-side product inequality of Lemma~\ref{B8.10m} and the complementary subset where the second term satisfies the symmetric inequality. On each subset, Lemma~\ref{B8.10m} multiplies the small-side squared amplitude bound by \(2C_0\Gamma^2(4\sqrt{T_{u'}/T_u}\,g^{1-m}k^{-r})^{kf}\) (with \(4\sqrt{T_{u'}/T_u}\le C\) in both orders from \(\max(T_u/T_{u'},T_{u'}/T_u)\le a_F\) and \(C\ge4\sqrt{a_F}\)), while \(\sqrt{n_Zn_{Z'}}\) from \eqref{eq:p8c1-switch-count} cancels \(n_Z^{-1/2}n_{Z'}^{-1/2}\) from the two amplitude bounds. Taking the square root gives \((Cg^{-(2r+1)}k^{-r})^{kf/2}\). High-order length moments absorb \((h_Zh_{Z'})^{3cr^2/2}\) and leave \(S_{C_h}(Z)^{-12(r+1)^2}S_{C_h}(Z')^{-12(r+1)^2}\) (for the cutoff constant \(C_h\ge0\) of Lemma~\ref{B8.08k}); summing over \((Z,Z')\) costs at most \(C^fH^{Cm^3}\) with decay and dimension losses bounded by \(e^{Cm^6fL_b}\), and summing over row classes via \(\Gamma_1\dim(\pi_u)\operatorname{mass}_u\le f_u\) yields \eqref{eq:p8c1-loop-one}.
\leanref{Formalism.\allowbreak{}L2.\allowbreak{}SG.\allowbreak{}RowsU.\allowbreak{}rowsOneU}{Formalism.\allowbreak{}LocalToGlobal.\allowbreak{}SiegelLatticeTheory.\allowbreak{}FlagStabilizerTensors.\allowbreak{}RootLatticeThresholds.\allowbreak{}DegenerateConstant}{297--422}
\leanref{Formalism.\allowbreak{}L2.\allowbreak{}SG.\allowbreak{}RowsU.\allowbreak{}rowsFU}{Formalism.\allowbreak{}LocalToGlobal.\allowbreak{}SiegelLatticeTheory.\allowbreak{}FlagStabilizerTensors.\allowbreak{}RootLatticeThresholds.\allowbreak{}DegenerateConstant}{425--442}
The complete proof is contained in \Lean{Formalism.\allowbreak{}L2.\allowbreak{}SG.\allowbreak{}RowsU.\allowbreak{}rowsOneU}, which proves the row bound \Lean{Formalism.\allowbreak{}MainNotions.\allowbreak{}BridgeLemma.\allowbreak{}RowsOne} and enters the loop estimate through \Lean{Formalism.\allowbreak{}L2.\allowbreak{}SG.\allowbreak{}RowsU.\allowbreak{}rowsFU}.
\end{proof}

\subsection*{Proof of \ref{B8.10a}}\label{prf:B8.10a}\phantomsection
\begin{proof}[Proof (outline)]
Use the notation and constant \(C\) of Definition~\ref{B8.10ab}, Proposition~\ref{B8.10ac} and Proposition~\ref{B8.10ad}, and put \(\Delta=f_uR'=\Delta_{u,u'}\) by the Hida-model identification \(d_w^C=d_w^S\), \(D_C^\star=D^\star\). Partition \(S_{\rm big}\cup S_{\rm nsc}\) into \(\coprod_{s=2}^k(\mathcal P_{{\rm nd},s}\sqcup\mathcal P_{{\rm deg},s})\sqcup S_{\rm nsc}\). On any finite subset \(U\subset S_{\rm big}\cup S_{\rm nsc}\), apply \eqref{eq:p8c1-loop-nd-deg} of Proposition~\ref{B8.10ac} to \(U\cap\mathcal P_{{\rm nd},s}\) and \(U\cap\mathcal P_{{\rm deg},s}\) for each \(2\le s\le k\), and \eqref{eq:p8c1-loop-one} of Proposition~\ref{B8.10ad} to \(U\cap S_{\rm nsc}\). From the source dimension bound \(1\le C_0R'\Gamma^2\) of Definition~\ref{B8.02c} we obtain \(f_uV_s\le\max(1,C_0)\Delta\Gamma^2V_s\). Summing over \(2\le s\le k\) and taking the supremum over finite subsets \(U\) gives
\begin{equation}\label{eq:p8c1-loop-majorant}
 \operatorname{PairAbs}_{u,u'}(S_{\rm big}\cup S_{\rm nsc})
 \le\max(1,C_0)\Delta\left[\sum_{s=2}^k(U_s+\Gamma^2V_s)+J\right].
\end{equation}

Put \(\gamma=e^{\min(C,C_\Gamma)m^5fL_b}\), so \(1\le\gamma\le\Gamma\). Because \(U_s\), \(\Gamma^2V_s\) and \(J\) depend on \(\Gamma\) through \(\Gamma\), \(\Gamma^3\) and \(\Gamma\), replacing \(\Gamma\) by \(\gamma\) inside the bracket of \eqref{eq:p8c1-loop-majorant} multiplies the bound by at most \((\Gamma/\gamma)^3\). Set \(A=\max(1,C_0)(\Gamma/\gamma)^3\). The source thresholds of Definition~\ref{B8.02c} give \(\log(\Gamma/\gamma)\le C_\Gamma m^5fL_b\le k/8\) and \(6\log(3A)\le P\). Evaluating \(F_s,U_s,V_s,J\) at \(\gamma\), Lemma~\ref{B8.02bc} and Theorem~\ref{B8.48} give \((kv)^{kf}\Xi_s\le H^CC^{skf}\) for \(2\le s\le k-2\) as in the proof of Lemma~\ref{B8.08k}, while at \(s=k-1,k\) we substitute \eqref{eq:p8c1-endpoint-index}, \(v\le b\), \(N\le b^{Cmf}\) and \(s-2\ge k-3\). With \(4\le m\le f\le k\), \(g\ge(kb)^8\), \(C_\ell m^6L_b\le k\log k\) and \(k\ge16C^{32}m^5fL_b\), taking logarithms at \(\gamma\) gives
\begin{equation}\label{eq:p8c1-loop-decay}
 \log U_s\le-\tfrac12kf\log g-\log k,\qquad
 \log V_s\le-P/2-2\log\gamma-\log k,\qquad
 \log J\le-P,
\end{equation}
where at \(s=2\) the factor \((C/k)^{kf}\) supplies the negative main term in \(\log V_2\). Summing over \(2\le s\le k\) bounds \(\sum_{s=2}^k U_s\), \(\sum_{s=2}^k\gamma^2V_s\) and \(J\) at \(\gamma\) by \(e^{-4P}\), \(e^{-P/2}\) and \(e^{-P}\), respectively. Substituting into \eqref{eq:p8c1-loop-majorant} gives at most \(A(e^{-4P}+e^{-P/2}+e^{-P})\Delta\le3Ae^{-P/2}\Delta\le e^{-P/3}\Delta\).
The complete proof is contained in \Lean{GlobalTheory.\allowbreak{}LoopRowBoundsGap.\allowbreak{}rows\_\allowbreak{}along\_\allowbreak{}family}.
\end{proof}

\subsection*{Proof of \ref{B8.10aa}}\label{prf:B8.10aa}\phantomsection
\begin{proof}
For \(0\ne l\in J\), the integral ideal \((l)J^{-1}\) has norm at least one, so \(\prod_v|v(l)|\ge\sqrt{N(J)}=:D\). If some \(z_v=0\) the sum vanishes. Otherwise partition the points by \(n_v=\max(0,\lceil\log_2(|v(l)|/z_v)\rceil)\). Put \(h=\sum_v n_v,\ R_v=2^{n_v}z_v\). A nonempty box has \(U=\prod_vR_v/D\ge1\). Subdivide the coordinate intervals \([-R_v,R_v]\) in each complex plane into intervals of length at most \(R_v/(2U^{1/f})\). Each product box contains at most one point since a nonzero difference in it would have product of absolute values strictly below \(D\). There are at most \(25^f U^2\) boxes. Each summand is at most \(2^{5f-5h}\) and \(U\le 2^h(a/\Theta)^f\). Since \(2^h\ge(\Theta/a)^f\), the shell contributes at most \(800^f(a/\Theta)^{5f/2}2^{-h/2}\). Sum over \((n_v)\) using \(\sum_{(n_v)}2^{-h/2}=(1-2^{-1/2})^{-f}\).

The Fourier character on \(E\otimes\R\) is \(\exp(2\pi i\sum_v2\Re(v(l)c_v)/g)\), which gives the stated dual. Since \(g^2\mathfrak d^\nu\subset L_{b_0s}\), trace duality gives \(L_{b_0s}^*\subset g^{-1}\mathfrak d^{-\nu-1}\). The inverse-difference lattice is contained in \(g^{-1}\mathfrak d^{-\nu}\) by definition.
\end{proof}

\subsection*{Proof of \ref{B8.10g}}\label{prf:B8.10g}\phantomsection
\begin{proof}
Make the changes \(X=tV,\ B=tb,\ \rho=taJ_n(1+V)\), and denote the last matrix entry by \(\sigma\). The resulting matrix is
\[
 M=\begin{pmatrix}J_n(tI+X)&J_n B\\ \rho&\sigma\end{pmatrix}.
\]
Write \(x_1,\ldots,x_n\) for the first superdiagonal entries of \((X\ B)\), including \(x_n=B_n\). The character is \(\psi_\beta((\rho_1+\sum_i x_i)/t)\). Indeed \(\rho_1=ta_n\). The Jacobian from the three changes before introducing \(\sigma\) is \(|t|^{-n(n+3)}\) in the integral.

Put \(D'=(-1)^nD\). The matrix \(J_s M\) has first subdiagonal \(t\), zeros below that subdiagonal, diagonal \((\rho_1,x_1,\ldots,x_n)\), and determinant \(D'\). Use its strictly upper entries \(y\) and the \(x_i\) as independent variables. In replacing \(\rho_1\) by \(\sigma\), the derivative has modulus \(|t|^n/|C_t|\), where \(C_t\) is the cofactor of \(\rho_1\) in \(J_sM\). This follows by differentiating the determinant in these two entries. The cofactor of \(\sigma\), up to sign, is \(t^n\). Thus the remaining Jacobian, after multiplication by \(|t|^{n(n+1)}\), is \(|C_t|^{-2}\). At \(t=0\), \(C_0=\prod_i x_i\).

Every independent matrix entry is bounded on the support. Expansion of the determinant at \(t=0\) gives \(D'=\rho_1\prod_i x_i+O(|t|)\) there. Since \(|D'|\) has a positive lower bound, each \(|x_i|\) and \(|C_t|\) has a positive lower bound for small \(t\). We can therefore solve for \(\rho_1\) throughout the support. Put
\[
 p(x)=\sum_i x_i+D'/\prod_i x_i .
\]
After removal of \(\psi_\beta(p(x)/t)\), the amplitude is
\[
 |C_t|^{-2}\,
 \psi_\beta\!\left(\frac{\rho_1-D'/\prod_i x_i}{t}\right)\Phi(M).
\]
The determinant equation is polynomial in \(t\) and affine in \(\rho_1\). Division by \(C_t\) shows that the divided difference in this amplitude extends smoothly to \(t=0\). A fixed smooth cutoff on a compact neighborhood of the support gives uniform bounds for all its derivatives in \(x,y\).

The critical equations are \(x_i=D'/\prod_jx_j\) for every \(i\). Thus \(x=(z,\ldots,z)\) with \(z^s=D'\), and \(p(x)=sz\). The complex Hessian there is \(z^{-1}(I+\mathbf1\mathbf1^{\mathsf T})\), of determinant \(s z^{-n}\). The roots stay separated on the compact parameter set. Outside fixed small neighborhoods of them the gradient is bounded below, so \(n+1\) integrations by parts give \(O(|t|^{n+1})\) in the transformed integral. At a root apply stationary phase in \(2n\) real dimensions to \(\Re(\beta |t|p/t)\). Its Hessian has signature zero and absolute determinant \(|\beta|^{2n}|s z^{-n}|^2\). The amplitude is compactly supported with uniformly bounded derivatives, and the inverse Hessians are uniformly bounded. The leading multiplier in the measure \(d\mu_\beta^n\) is \(|t|^n|z|^n/s\), with error \(O(|t|^{n+1})\). The scalar bound is uniform on compositions of \(\Phi\) with linear functionals of norm at most one, which gives the Banach-space estimate. \cite[Theorem~7.7.5]{hormander-2003}

At the critical point write \(J_s M|_{t=0}=zI+Y_0\), with \(Y_0\) strictly upper triangular. Differentiate the determinant with \(x,y,D'\) fixed. The subdiagonal has derivative one, and the first diagonal entry has derivative \(\partial_t\rho_1\). The trace identity for the derivative, using the first superdiagonal of \((zI+Y_0)^{-1}\), gives \(\partial_t\rho_1=\ell(Y_0)/z\). There the cofactor contributes \(|z|^{-2n}\). Substitution \(Y_0=zY\) contributes \(|z|^{n(n+1)}\) and turns the residual integral into \(L_\Phi(z)\). The powers of \(|z|\) give \(n-2n+n(n+1)=n^2\). Finally \(|t|^{n(n+1)}/|t|^n=|t|^{n^2}\), and the transformed remainder becomes \(O(|t|)\).
\end{proof}

\subsection*{Proof of \ref{B8.10i}}\label{prf:B8.10i}\phantomsection
\begin{proof}
For a compact annulus \(K\), set \(K'=\{z:\omega^{-1}z^s\in K\}\). Expand \(\int_K\|A_{p,x}\|^2\) and change variables by \(D=\omega^{-1}z^s\). The pairs of roots are indexed by \(\zeta^s=1\). The diagonal contribution is exactly \(\int_{K'}\|b_p\|^2\). Each other contribution is the Fourier integral of an \(L^1\) density evaluated at frequency \(s(\zeta-1)x\). Cauchy--Schwarz gives integrability on \(K'\), and the Riemann--Lebesgue lemma makes the contribution tend to zero. The annular approximation and the bound on \(U_{p,x}\) now give \(\int_{K'}\|b_p\|^2\le M_p\). Exhaust \(\C^\times\) by these annuli.

The function \(a_p(D)=\sum_{z^s=\omega D}\|B_p(z)\|\) satisfies \(\|a_p\|_2^2\le s\|b_p\|_2^2\) by the same Jacobian and the finite Cauchy--Schwarz inequality. It dominates every \(A_{p,x}\). In the \(D\)-coordinates of the first fibre the second boundary vector is dominated by \(|\lambda|a_{p'}(\lambda D)\), where \(\lambda=(x/x')^{s-1}\to1\). These dominators converge to \(a_{p'}\) in \(L^2\). This follows first for continuous compactly supported functions by change of variables, then by density and the isometry of normalized dilation. Given \(\epsilon>0\), choose a compact annulus on whose complement the dominators have \(L^2\)-norm below \(\epsilon\) for all sufficiently large \(|x|\). On the annulus the error in replacing \(U_{p,x}\) by \(A_{p,x}\) tends to zero. On its complement the mixed error is at most \((\sqrt{M_p}+\|a_p\|_2)\epsilon\). Interchanging the prefixes proves the other mixed comparison.

For the boundary-boundary pairing choose the root \(\rho_x^s=\lambda\) tending to one. The second roots are \(\zeta\rho_x z\). The change of variables just used gives densities
\[
 \left\langle b_p(z),|\rho_x|b_{p'}(\zeta\rho_x z)\right\rangle
\]
and phases \(\chi(s(\zeta\rho_x x'-x)z)\). The densities converge in \(L^1\) by \(L^2\)-continuity of dilations and rotations. Since \(\rho_x=1-(s-1)\sigma/(s x')+O(|x|^{-2})\), the frequency tends to \(\sigma\) for \(\zeta=1\). For the other roots its modulus tends to infinity. Fourier continuity and decay give the displayed integral. The block positivity follows by taking all the \(T_xU_{p,x},T_xA_{p,x}\) as vectors in the same \(L^2(dc)\).
\end{proof}

\subsection*{Proof of \ref{B8.10p}}\label{prf:B8.10p}\phantomsection
\begin{proof}
The argument of \(\theta\) is at least one on the smaller ball and nonpositive outside the larger one. The equation \(g^*\operatorname{diag}(1_p,-1_q)g=\operatorname{diag}(1_p,-1_q)\) defines a closed subset of matrix space and forces \(g\) to be invertible. Its intersection with the larger closed ball is compact and avoids zero. For fixed left and right word lengths the differentiated coefficient depends continuously on \(g\) and multilinearly on the Lie algebra directions. Take the maximum on this compact set times the corresponding closed unit balls, and then the maximum over the finitely many pairs of word lengths. Outside the support every derivative vanishes. The factor \(\|g\|_F^A\) is bounded on the same compact set.
\end{proof}

\subsection*{Proof of \ref{B8.10e}}\label{prf:B8.10e}\phantomsection
\begin{proof}
Put \(r=r_S,\ g=g_0g_1,\ v_E=\operatorname{rd}(E)/2,\ \kappa=\pi/v_E\), and choose a differing interior position \(l\) of the two binary words. At each place let \(v_l\) be the corresponding interlacing parameter of \ref{B8.02c}, and write \(u'\) for the weight of word \(v\). Use \(N=20(r+1)=10m\) letters and a cut of width \(I=\mathcal D/2-5/2\), where \(\mathcal D\) is the choice spacing. Set
\[
 P(y)=\prod_j(y+\lambda_j),\qquad
 q(y)=-T_uT_v\prod_i(y+1/2+u_i)(y+1/2+u'_i).
\]
Write \(B\) for the supremum of \(g^m\kappa^{1-m}|P(y)P(y+1)/q(y)|^{1/2}\) over \(5/2\le y+v_l\le\mathcal D/2\), and \(B_1\) for the supremum on the part of this interval with \(y+v_l\le N^3\). The cut avoids all roots of \(q\). For a threshold constant \(C_A\) put \(d=C_A\kappa N/(gI)\). Define \(\Theta^{-f}\) as the maximum of the products over places of the following factors, with the indices common to all places:
\begin{equation}\label{eq:p8c1-thresholds}
\begin{array}{ll}
 \widehat X_1=C_A B_1, &(i=1,\ s=1),\\
 \widehat X_i=C_A B d^{2i-2}, &(2\le i\le r+1,\ s=1),\\
 \widehat\Lambda_{is}=C_A B d^{2(i-1)+s-1}, &(i\ge1,\ s\ge2,\ i+s\le r+2),\\
 \widehat\Lambda_i'=(C_A B\kappa d^{2(i-1)}/\sqrt{T_uT_v})^{1/2}, &(1\le i\le r+1).
\end{array}
\end{equation}

For use in the arithmetic coordinates define \(J_{is}=g^{-1}\mathfrak d^{-\nu-1}\), \(\nu=2r+3-2i-s\). Divide the first three factors of \eqref{eq:p8c1-thresholds} by \(g\operatorname{rd}(E)^{\nu+1}\) to obtain the scales \(z_{is,v}\), with \(i\) now the block start. With \(s=r+2-i\), divide the last factor by \(g\operatorname{rd}(E)^\nu\) to obtain the terminal scale \(z'_{i,v}\). The product of these divisors is the inverse square root of the norm of the corresponding ideal \(J_{is}\) or \(g^{-1}\mathfrak d^{-\nu}\).

Put \(t_v=\sqrt{T_{\mathrm{crit},v}(u)T_{\mathrm{crit},v}(u')/(T_uT_v)}\) using the critical levels of \ref{B8.10}. For the cut estimate put \(x=-y\) and let \(\delta_i,\delta_i'\) be the binary choices at \(i\), so that the applicable weight coordinates are \(v_i-1/2-\delta_i\mathcal D\) and \(v_i-1/2-\delta_i'\mathcal D\). Pair the interlacing spectral roots with the weight roots by writing the interior factors as
\[
 R_i(x)=
 \frac{\sqrt{|(v_i-x)(v_i+1-x)|}}
      {\sqrt{|(v_i-x-\delta_i\mathcal D)(v_i-x-\delta_i'\mathcal D)|}} .
\]
The top-place last weight uses \(B_*+x-1/2\) in the denominator. For \(i<l\), the denominators are bounded below by the corresponding distances \(v_i-x-\mathcal D\), and the product of these factors is at most \(3.3\sqrt{k}\) by the spacing of the \(v_i\). For \(i>l\) in the interior, the ratios are at most one. At \(i=l\), exactly one denominator is shifted. With \(h=v_l+y\), this gives
\(
 R_l(x)^2=(h+1)/(\mathcal D-h)\le(14/5)h/\mathcal D
\)
from \(5/2\le h\le\mathcal D/2\).

The remaining spectral factors satisfy
\(\prod_j\sqrt{(B_j+x)(B_j+x-1)}\le e^\beta\prod_j B_j\), where \(\beta\) can be taken as the sum of the upper bounds \(x/B_j\) on the cut. The size bound \(B_j\ge C_\ell mk\mathcal D\) makes \(\beta\) bounded. In the edge comparison the first weight factor \(1/(u_1+1/2-x)\) times the last ratio is at most \(8/a_{\rm eff}\), where \(a_{\rm eff}\) is the edge denominator in the critical-level comparison. Off the top it is \(u_1+1/2\), and on the top it is \((u_1+1/2)B_*/v_{k-1}\) with the top edge factor included. The critical-level lower bounds compare \(g^m\kappa^{1-m}\prod_j B_j/a_{\rm eff}\) with each \(T_{\mathrm{crit},v}\) up to a fixed factor. Multiply the two comparisons and take the square root, then divide by \(\sqrt{T_uT_v}\). Consequently on the cut
\[
 g^m\kappa^{1-m}|P(y)P(y+1)/q(y)|^{1/2}
 \le C_2 t_v\sqrt{k}\sqrt{(v_l+y)/\mathcal D}.
\]
Taking suprema gives \(B\le C_2\sqrt{k}\,t_v\) and \(B_1\le C_2\sqrt{kN^3/\mathcal D}\,t_v\).
\leanref{Formalism.\allowbreak{}FourierDecayAlt.\allowbreak{}P.\allowbreak{}Sz7.\allowbreak{}betaOnCut\_\allowbreak{}of\_\allowbreak{}tcritLower}{Formalism.\allowbreak{}HermitianCurvature.\allowbreak{}HermitianCurvatureProofs.\allowbreak{}BetaBoundOnBoundary}{163--254}

Write \(T_*=T_uT_v\) and \(D=C_A\pi N/I\). Take the rank cutoff so that \(0<d\le D\le1\). Let \(K\ge1\) absorb the fixed constants in the two cut bounds above, including \(C_A\) and \(\kappa\le\pi\). Substitution in \eqref{eq:p8c1-thresholds} bounds the four local entries, in their displayed order, by \[ Kt_v\sqrt{kN^3/\mathcal D},\qquad Kt_v\sqrt{k}\,D^2,\qquad Kt_v\sqrt{k}\,D,\qquad K\sqrt{t_v}(k/T_*)^{1/4}. \] Indeed \(2i-2\ge2\) for a later singleton, \(2(i-1)+s-1\ge1\) for a longer block, and the power of \(d\) inside the terminal square root is nonnegative. The level, dimension and mass ties give \(\prod_v t_v\le C_t^f\). Thus taking the geometric mean over places replaces \(t_v\) by at most \(C_t\), and \(\sqrt{t_v}\) by at most \(\sqrt{C_t}\). Since \(\Theta^{-1}\) is the maximum of these geometric means and \(D^2\le D\), taking reciprocals gives
\begin{equation}\label{eq:p8c1-flag-size}
 \Theta\gg\min\{\sqrt{\mathcal D/(kN^3)},\,I/(\sqrt{k}N),\,(T_uT_v/k)^{1/4}\}.
\end{equation}
Choose \(C_A\) above the constants for the decay and recursion before taking the rank cutoff. With \(\mathcal D\ge(kf\log b)^5,\ m\le f\le k,\ b\ge k^{3/2}\), and \(T_w\ge b^{m/2}\), the three terms grow at least as \(k^{1/2},k^{1/2},k^{1/8}\). Thus the flag-size exponents \((5,\tfrac12,3)\) allow \(\Theta\ge k^{1/16}\).
\leanref{Formalism.\allowbreak{}L2.\allowbreak{}SG.\allowbreak{}ThmSOwn.\allowbreak{}SizesRowS}{Formalism.\allowbreak{}LocalToGlobal.\allowbreak{}SiegelLatticeTheory.\allowbreak{}FlagStabilizerTensors.\allowbreak{}ThmSOwn}{1008--1013}
\leanref{Formalism.\allowbreak{}FourierDecayAlt.\allowbreak{}P.\allowbreak{}threshold\_\allowbreak{}products}{Formalism.\allowbreak{}HermitianCurvature.\allowbreak{}HermitianCurvatureProofs.\allowbreak{}CmRingIdempotentUnitRatioIdentity}{241--334}
\leanref{Formalism.\allowbreak{}FourierDecayAlt.\allowbreak{}P.\allowbreak{}theta\_\allowbreak{}lower\_\allowbreak{}bound}{Formalism.\allowbreak{}HermitianCurvature.\allowbreak{}HermitianCurvatureProofs.\allowbreak{}CmRingIdempotentUnitRatioIdentity}{335--393}

Fix a finite set \(S\) of scalar-partner rows. For each start \(b_0\) and length \(s\) choose a finite \(S_{b_0s}\subset E^\times\). The set \(X_i\) consists of lists \(p=((s_1,t_1),\ldots,(s_d,t_d))\) with \(\sum s_a=i\) and \(t_a\in S_{1+\sum_{h<a}s_h,s_a}\). Thus \(X_0\) consists of the empty list. The sets of allowed entries are ordered by inclusion. Every finite set of terminal flags is contained in some \(X_{r+1}\).

Use the cut Gram forms of the partial orbital pairings from \ref{B8.10h} in the arithmetic coordinates. A parameter \(w\) is a tuple of local letter lists, and the cut form at each place is divided by its product of positive letter weights. Let \(H_i(w;p,q)\) be the product of these normalized local Gram forms with the square roots of the two point weights included. For a finite prefix set \(X\) and a kernel \(B_0\) on it put \([B_0,H]_X=\Re\sum_{p,q\in X}B_0(p,q)H(p,q)\). The product of two positive semidefinite kernels is the Gram matrix of tensor products of their Gram vectors. Hence this expression is nonnegative and monotone in \(B_0\) for positive semidefinite order. The triangle inequality for the sum of tensor vectors bounds the square root of the full form by the sum of the square roots of the forms restricted to each last-block length.

The arithmetic kernels \(K_i\) are the flag kernels on prefixes, with \(K_0=1\) on the empty flag and \(K_{r+1}=K_{\rm top}\) before density scaling. Normalize the empty-flag energy so that its uncut value is \(I_N\), the normalized archimedean integral. At the fixed tower choose \(\epsilon>0\) small enough that every supported terminal Gram entry agrees with the Hilbert--Schmidt pairing of the terminal orbital vectors \(V_p\). The cut is one on the finitely many supporting Gram values for small enough \(\epsilon\). Apply \ref{B8.10c} to the Hilbert-space coordinates of these vectors on the entire \(X_{r+1}\). This bounds the row square sum by the energy with top kernel \(e_1K_{\rm top}\), where \(e_1=\max(1,s_d)\). Use this scaled kernel at the top and put \(\eta_i=1\) below the top and \(\eta_{r+1}=e_1\).

Fix a last length \(s\), put \(b_0=i-s+1\), and write \(pt\) for \(p\) with \((s,t)\) appended. The diagonal quantity at \(t\) for a prefix kernel \(B_0\) is \(D_t(w;B_0)=[B_0,H_i(w;pt,qt)]_{X_{i-s}}\). For distinct \(t,t'\) the arithmetic kernel is supported on \(l=t^{-1}-{t'}^{-1}\in L'_{b_0s}\). At \(i=r+1,s=1\) its support requires \(t=t'\).

For the cross-term bound write \(M_z(l)=\prod_v\min(1,(z_v/|v(l)|)^5)\) with the inverse-difference scales \(z_v\). Use the lift interpolation identity of \ref{B8.10h} with five evaluation variables. In a variable of degree \(d_0=2(b_0-1)+s\), \(d_0+1\) grid nodes at spacing \(a\) have Lagrange coefficient sum at most \(2^{d_0}/(a^{d_0}d_0!)\). Use rows in the middle cut, with spacing of order \(I/((d_0+1)N)\) after removing used letters. Below the terminal level the multiplier is the inverse difference, and the scale is \(z_{b_0,s+1,v}\). At the terminal level with \(s\ge2\) use the diagonal interpolation of a single vector. Put \(d=r+1-s\), \(\tau=(g/(2\pi))^\nu i^{s-1}t\), and \(K=5\). Each of the \(K\) rows interpolates a polynomial of degree \(d\), with leading multiplier \((-1)^{s-1}/\tau\). If \(V\) is a terminal component and \(V_n\) its positive-letter extensions on the grid, extraction of the leading coefficients gives \(V=(((-1)^{s-1}/\tau)^K)^{-1}\sum_n L_n V_n\), where \(L_n\) is the product of the row interpolation coefficients. Write \(P_n=\prod_a|P(y_{a,n_a})|\). If the component's squared weight is \(d_0\) before extension and \(d_n\) afterwards, the positive-letter weight comparison is \(d_0P_n\le C_t^Kd_n\), with \(C_t\) the uniform upper bound for \(|P(y)|\) divided by the positive one-letter weight ratio on this grid. Apply weighted Cauchy--Schwarz to the interpolation sum, also after taking any linear combination of prefixes. With \(L_*=\sum_n|L_n|\) this bounds its weighted square by \(|\tau|^{2K}L_*^2 C_t^K\) times the average of the extended normalized squares with probabilities \(|L_n|/L_*\). Distinct old components extend to distinct components, so summing these inequalities preserves the bound. Here \(L_*\le(28N/I)^{dK}\); the resulting diagonal multiplier is \((|t|z'_{b_0,v})^{2K}\) with the terminal scale of \eqref{eq:p8c1-thresholds} as upper bound. For \(l=t^{-1}-{t'}^{-1}\ne0\), at least one of \(|t|,|t'|\) is at most \(2/|l|\). Use the diagonal estimate on that side and take a square root in Cauchy--Schwarz. This gives the terminal scale \(2z'_{b_0,v}\) in the off-diagonal bound. These scales are measured in the inverse-difference ideal of \ref{B8.10aa} because \(\nu\) decreases by one at the next length.

At a place with \(|v(l)|\le z_v\) use the unextended Cauchy--Schwarz inequality. Else use the one with five letters, which pairs a positive extension and its negative companion. Tensor the inequalities over the places and bound the geometric mean of the two diagonal evaluations by their arithmetic mean. For a fixed \(w\) there are at most three choices per place in these word averages, the original list and the two signs of extension. The cross term is bounded by \(M_z(l)\) times a probability average of half the sum of the two extended diagonal energies. For fixed \(t\), \(t'\mapsto l\) is injective, and the same assertion holds with \(t,t'\) interchanged. By \eqref{eq:p8c1-lattice-sum} the off-diagonal cost is at most \(3^fe\) times a uniform bound on the diagonal sum over the enlarged word set, where
\begin{equation}\label{eq:p8c1-small-frequency-error}
 e=(C_1\Theta^{-5/2})^f .
\end{equation}
The factor two in the terminal scale is absorbed in \(C_1\).
\leanref{Formalism.\allowbreak{}FourierDecayAlt.\allowbreak{}P.\allowbreak{}schurRecursion\_\allowbreak{}true}{Formalism.\allowbreak{}HermitianCurvature.\allowbreak{}HermitianCurvatureProofs.\allowbreak{}BlockStartPoints}{171--327}

On the fixed tower let \(A_{v,t}\) be the local diagonal kernel on \(X_{i-s}\) at a finite place \(v\), \(B_v\) its preceding kernel, and \(C_{v,\bar t}\) the local coset majorant. The local orbital description gives \(C_{v,\bar t}\succeq A_{v,t}\succeq0\) and \(\sum_{\bar t}C_{v,\bar t}\preceq B_v\), where a finite subset of distinct local cosets may be used in the sum. The coset kernel uses the phase and the integral presentation summed over the full residue class of the last entry. Restricting to a particular entry gives its diagonal kernel. In the Gram description the omitted terms are positive. Summing over cosets partitions the presentations in the preceding kernel. With the product normalization \(\alpha_i\) of the arithmetic kernels, take
\[
 \Theta_t=\alpha_i\prod_v C_{v,\bar t_v}.
\]
Include the top density scaling in \(\alpha_i\) at \(i=r+1\). The positive semidefinite order survives products, by telescoping the difference as a sum of products with one positive semidefinite difference factor. This gives \(\Theta_t-K_i|_{t,t}\succeq0\). Congruence modulo \(L_{b_0s}\) is equivalent to the local congruences, so equivalent entries give the same majorant. Distinct global cosets give distinct tuples of local cosets. On a finite set of representatives, extend their tuples to the product of the sets of local components that occur, and sum. The added kernels are positive semidefinite. The extended sum is the product of the local sums and is bounded by \(\alpha_i\prod_v B_v=\eta_iK_{i-s}\). The reference factors at unexceptional places use the same normalization in all three comparisons. Thus applying \([\,\cdot\,,H]\) to the sum of the majorants gives the preceding energy with multiplier \(\eta_i\).
\leanref{Formalism.\allowbreak{}FourierDecayAlt.\allowbreak{}TB.\allowbreak{}Data.\allowbreak{}ThetaDomination}{Formalism.\allowbreak{}HermitianCurvature.\allowbreak{}StatementTB}{141--149}
\leanref{Formalism.\allowbreak{}FourierDecayAlt.\allowbreak{}P.\allowbreak{}DensitySumsAndDomination.\allowbreak{}domAfterOfLocal}{Formalism.\allowbreak{}HermitianCurvature.\allowbreak{}HermitianCurvatureProofs.\allowbreak{}DensitySumsAndDomination}{112--142}

When \(s=1\), periodize the diagonal pairing in the new coordinate \(c\in\C^f\). With the character of \ref{B8.10aa}, the Fourier coefficient at \(l\in L_{i1}^*\) is its Fourier integral divided by the covolume of the period lattice. Translation by a coset representative contributes a phase of modulus one. The zero coefficient is the preceding energy by unfolding the singleton integral, whose normalization cancels the covolume. To bound the other coefficients use the first two scales of \eqref{eq:p8c1-thresholds}. The singleton interpolation identity with five variables of degree \(2(i-1)\) uses the short cut \(y+v_l\le N^3\) at \(i=1\) and the middle cut at the later positions. At each place use the stronger of the unextended bound and the bound with five letters. This gives \(M_z(l)\) times an average of the preceding diagonal energies, with at most \(2^f\) subsets of places using extensions. By \eqref{eq:p8c1-lattice-sum} for \(J_{i1}\), the nonzero frequencies cost at most \(2^fe\) times the uniform preceding bound and their sum is absolutely convergent.

In lattice coordinates, integrate over a fundamental parallelepiped and use the Fourier series of the periodization \cite[Ch.~5, \S3]{stein-shakarchi-2003}. For the cut pairing, integrability gives the coefficients by unfolding. The absolutely convergent series agrees almost everywhere with the periodization. A finite sum of nonnegative terms at a represented coset point is bounded by the Fourier majorant at that point, by continuity of its terms and approximation through the almost-everywhere set. Including the inverse-difference terms, the singleton costs at most \(\eta_i(1+3^fe)(1+2^fe)\).
\leanref{Formalism.\allowbreak{}FourierDecayAlt.\allowbreak{}P.\allowbreak{}singleBlockStep}{Formalism.\allowbreak{}HermitianCurvature.\allowbreak{}HermitianCurvatureProofs.\allowbreak{}FormDomination}{115--196}

For \(s\ge2\), at place \(v\) normalize the Fourier transform by \(j_{b_0s}^{-1}\), where
\[
 D_g=\frac{2}{\operatorname{rd}(E)g^3},\qquad
 j_{b_0s}=\frac{\operatorname{rd}(E)^{2\nu}}{D_g}
         =\frac{g^3}{2}\operatorname{rd}(E)^{2\nu+1},
 \qquad \nu=2r+3-2b_0-s .
\]
At each completely split prime dividing \(g\), the condition \(t-(-1)^s\bar t\in g^2\mathfrak d^\nu\) imposes one equality modulo that prime on each conjugate pair of components of \(t/g\). It has total index \(g^f\) in \(g\mathfrak d^\nu\). The latter lattice has covolume \(g^{2f}\operatorname{rd}(E)^{(2\nu+1)f}/2^f\) in the product complex Lebesgue measure. Thus \(\prod_v j_{b_0s}=\operatorname{covol}(L_{b_0s})\).

Write \(H^{(v)}\) for the local factor of \(H\). For fixed earlier flags \(p,q\), let \(D_{pq}(c)=H^{(v)}_{b_0+s-1}(w;p(s,c),q(s,c))\), where \(p(s,c)\) denotes appending \((s,c)\). Its normalized transform is
\[
 \widehat D_{pq}(\sigma)=
 \frac{2}{g^3\operatorname{rd}(E)^{2\nu+1}}
 \int_\C D_{pq}(c)\exp(2\pi i\,2\Re(\sigma c)/g)\,dc .
\]
In the inverse coordinate \(x\ne0\), the shorter Gram kernel \(U\) pairs \(p(s-1,x^{-1})\) and \(q(s-1,{x'}^{-1})\). There are cross and boundary kernels \(Q\) and \(P_b\) on \((p,x)\) with \(\begin{pmatrix}U&Q\\ Q^*&P_b\end{pmatrix}\) positive semidefinite on every finite set of such arguments. The boundary diagonal is bounded, and
\begin{equation}\label{eq:p8c1-boundary-limit}
 \lim_{|x|\to\infty}P_b(p,x;q,x+\sigma)
 =\lim_{|x|\to\infty}Q(p,x;q,x+\sigma)
 =\lim_{|x|\to\infty}Q^*(p,x;q,x+\sigma)
 =\widehat D_{pq}(\sigma).
\end{equation}

To construct the kernels, use the fibre over a singleton after the shortened block in the partial orbital integral. Denote its vector by \(V_p(x,c)\), with the square roots of the section weights included and the square root of the singleton measure factor divided out. Integration in \(c\) gives \(U\) by the singleton identity. On writing \(c=x^{s-1}D\), the vector with its measure factor is \( |x|^{s-1} V_p(x,x^{s-1}D)\). Its leading expression is
\[
 \sum_{z^s=\omega D} B_p(z)\exp(2\pi i\,2\Re(szx)/g),
\]
where \(|\omega|=1\) is the rotation of the block. The vector \(B_p(z)\) is the length-\(s\) orbital vector at the corresponding arithmetic entry, divided by \(\sqrt{j_{b_0s}}\) and by the square root \(s|z|^{s-1}\) of the root-map Jacobian.

For each fixed truncation, the group coefficient is smooth and compactly supported by Lemma~\ref{B8.10p}, and the lower-level vectors are obtained by integrating the cut section on the residual slices. On a compact annulus in \(D\), write \(t_{\rm std}\) for the standard shortened fibre variable. In standard coordinates the shortened entry and the long entry have scale factors \(l_1=(g/(2\pi))^s i^{s-2}\) and \(l_2=(g/(2\pi))^{s-1}i^{s-1}\). They satisfy \(l_2^s=(-1)^{s-1}l_1^{s-1}\). Thus \ref{B8.10g} at \(t_{\rm std}=l_1/x\), \(c=x^{s-1}D\), has long entries \(l_2 z\), \(z^s=D\). Since \(l_2/l_1=2\pi i/g\), the standard parameter \(2i\) gives the conjugate of the arithmetic phase \(\exp(4\pi i\Re(szx)/g)\) in the conjugate Gram coordinate. The annular error is \(O_K(|t_{\rm std}|)\). The rescaling gives \(|t_{\rm std}|=c_{b_0s}|x|^{-1}\) with \(c_{b_0s}\) fixed. To transfer this bound to the orbital vectors, take the finite Gram matrix indexed by the shortened point with its last singleton and the \(s\) root points. The quadratic form at the coefficient vector with shortened coefficient \(|x|^{s-1}\) and the negative root coefficients is the squared norm of the difference. In the slice model this form is bounded by the square of the slice difference times the bounded form norm. Thus the difference in the orbital model is at most \(C_K|x|^{-1}\) on the annulus \(K\).

At \(b_0+s-1=r+1\), the cut Gram matrix is a sum of Gram matrices of the projected top orbital operators. Write the terms as \(c_i\langle L_i J_i(P),L_i J_i(Q)\rangle\) with \(c_i\ge0\), including the point weights in \(J_i\). The same coefficient vector gives the difference in each term. The top section expression, after division by the matching section factors, varies by at most \(C_{i,K}|x|^{-1}\) from its root-slice expression. This follows by taking the difference along the rescaled fibre variable in its compactly supported section integral. The total difference in norm is bounded by \(\sum_i\sqrt{c_i}\|L_i\|C_{i,K}|x|^{-1}\). Use cutoffs that satisfy the larger of the lower- and top-level thresholds. For the finite prefix list use the preceding length's uniform diagonal bound. By the singleton integral identity it bounds the squared norm of each shortened fibre, after the isometric change \(c=x^{s-1}D\). The length-\(s\) vectors are square integrable on compact annuli by continuity there. Apply \ref{B8.10i} with the local root coefficients just computed. Its boundary Gram kernel gives \(P_b\), and its original-boundary pairing gives \(Q\).
\leanref{Formalism.\allowbreak{}FourierDecayAlt.\allowbreak{}P.\allowbreak{}FormNormSumIdentities.\allowbreak{}germPlace\_\allowbreak{}of}{Formalism.\allowbreak{}HermitianCurvature.\allowbreak{}HermitianCurvatureProofs.\allowbreak{}SliceAtApply}{110--403}
\leanref{Formalism.\allowbreak{}FourierDecayAlt.\allowbreak{}P.\allowbreak{}pb\_\allowbreak{}germRateLowB\_\allowbreak{}of\_\allowbreak{}parts}{Formalism.\allowbreak{}HermitianCurvature.\allowbreak{}HermitianCurvatureProofs.\allowbreak{}PbGermRateLowBOfParts}{487--490}
\leanref{Formalism.\allowbreak{}FourierDecayAlt.\allowbreak{}P.\allowbreak{}germ\_\allowbreak{}rate\_\allowbreak{}top\_\allowbreak{}bound}{Formalism.\allowbreak{}HermitianCurvature.\allowbreak{}HermitianCurvatureProofs.\allowbreak{}GermRateTopBound}{45--92}
\leanref{Formalism.\allowbreak{}FourierDecayAlt.\allowbreak{}P.\allowbreak{}GermRateSplit.\allowbreak{}germRateB\_\allowbreak{}of}{Formalism.\allowbreak{}HermitianCurvature.\allowbreak{}HermitianCurvatureProofs.\allowbreak{}GermRateSplit}{65--72}
\leanref{Formalism.\allowbreak{}FourierDecayAlt.\allowbreak{}P.\allowbreak{}germ\_\allowbreak{}parts\_\allowbreak{}compact\_\allowbreak{}support}{Formalism.\allowbreak{}HermitianCurvature.\allowbreak{}HermitianCurvatureProofs.\allowbreak{}GermPartsCompactSupport}{289--291}
\leanref{Formalism.\allowbreak{}FourierDecayAlt.\allowbreak{}P.\allowbreak{}GermRealisation.\allowbreak{}FibreGerm.\allowbreak{}toGermReal}{Formalism.\allowbreak{}HermitianCurvature.\allowbreak{}HermitianCurvatureProofs.\allowbreak{}GermLocalOfFibres}{93--153}

In pairing the leading expressions at \(x\) and \(x+\sigma\), keep the original integration coordinate \(c\) fixed. A root \(z\) in the first expression corresponds to the roots \(\zeta\rho_x z\) in the second, with \(\zeta^s=1\) and \(\rho_x^s=(x/(x+\sigma))^{s-1}\) on the branch tending to one. The expansion
\(
 \rho_x=1-\frac{s-1}{s}\frac{\sigma}{x+\sigma}+O(|x|^{-2})
\)
gives \(s\{\rho_x(x+\sigma)-x\}\to\sigma\). Thus for \(\zeta=1\) the limiting phase is the character of \(z\sigma\). For \(\zeta\ne1\) the frequency \(s\{\zeta\rho_x(x+\sigma)-x\}\) tends to infinity in modulus. After the change of variables to \(z\), continuity of dilations in \(L^2\) gives convergence of the product densities in \(L^1\). Fourier decay eliminates the terms with \(\zeta\ne1\). In the surviving term the Jacobian \(s^2|z|^{2(s-1)}\) cancels the two root-vector denominators and gives the length-\(s\) density with the factor \(j_{b_0s}^{-1}\). Its integral is \(\widehat D_{pq}(\sigma)\). The mixed-pairing bounds of \ref{B8.10i} give the remaining limits in \eqref{eq:p8c1-boundary-limit} with the same factor \(j_{b_0s}^{-1}\).
\leanref{Formalism.\allowbreak{}FourierDecayAlt.\allowbreak{}P.\allowbreak{}PBK.\allowbreak{}tendsto\_\allowbreak{}shift}{Formalism.\allowbreak{}HermitianCurvature.\allowbreak{}HermitianCurvatureProofs.\allowbreak{}NormOfMemRootSetOne}{188--222}
\leanref{Formalism.\allowbreak{}FourierDecayAlt.\allowbreak{}P.\allowbreak{}PBK.\allowbreak{}tendsto\_\allowbreak{}pairD}{Formalism.\allowbreak{}HermitianCurvature.\allowbreak{}HermitianCurvatureProofs.\allowbreak{}TendstoSmulCocompact}{107--147}

For the lattice estimate, let \(L\subset\C^f\) be the period lattice under consideration, and use its ordinary dual \(L^\vee\) for the character \(\exp(2\pi i\sum_v2\Re(\sigma_v c_v))\). Thus \(gL^\vee\) is the dual in the character with denominator \(g\). Fix a positive semidefinite prefix majorant \(\Theta_t\). Contract the tensor products of the local kernels against this majorant. Let \(u(x),b(x)\) be the resulting shorter and boundary vectors, with zero vectors at points where an inverse is undefined. Denote by \(D(c)\) the corresponding contracted diagonal density, and by \(\widehat D_0(\tau)\) the contraction of the product of normalized local transforms. Then the ordinary Fourier transform satisfies
\(\widehat D(\sigma)=(\prod_v j_{b_0s})\,\widehat D_0(g\sigma)\).

The lift interpolation bound gives
\[
 |\langle u(x),u(x+g\sigma)\rangle|\le F(\sigma)\qquad(\sigma\in L^\vee),
\]
where \(F(\sigma)\) is \(M_z(g\sigma)\) times the preceding diagonal bound averaged over the enlarged words. Take a translation-invariant mean on \(gL^\vee\) that agrees with limits as every coordinate tends to infinity in modulus. Such a mean is obtained as a limit of averages over growing boxes in a lattice basis. A fixed translate changes a proportion of the box tending to zero, and the proportion with a bounded coordinate tends to zero by counting along a lattice direction with nonzero projection to that coordinate. With \(r(x)=u(x)-b(x)\), set
\[
 A_0(\tau)=\operatorname{Mean}_x\langle r(x),r(x+\tau)\rangle .
\]
Translation invariance and the Gram inequality for finite lists \(r(x+\tau_i)\) show that \(A_0\) is positive definite. Expanding the residual pairing and using the three limits of \eqref{eq:p8c1-boundary-limit} gives
\[
 \widehat D_0(g\sigma)+A_0(g\sigma)
       =\operatorname{Mean}_x\langle u(x),u(x+g\sigma)\rangle .
\]
Thus \(\widehat D(\sigma)+A(\sigma)\), with \(A(\sigma)=(\prod_v j_{b_0s}) A_0(g\sigma)\), has absolute value at most \((\prod_v j_{b_0s}) F(\sigma)\).

For a finite \(Y_0\subset c_0+L\), periodize the nonnegative density \(D\) and take its product Fej\'er averages in lattice coordinates. The Fej\'er sum of \(A\) at each point is nonnegative by positive definiteness, since its coefficients are the autocorrelation weights of a finite lattice box. The Fourier coefficients of the averaged periodization are \(\widehat D(\sigma)/\operatorname{covol}(L)\) times the Fej\'er weights. Bound the averages by adding the Fej\'er sum of \(A/\operatorname{covol}(L)\) and using the summable majorant on the combined coefficients. A finite sum of translates of \(D\) is bounded by the periodization. Lower semicontinuity and concentration of the Fej\'er kernels bound that finite sum at \(c_0\) by the lower limit of its averages. Hence
\begin{equation}\label{eq:p8c1-herglotz}
 \operatorname{covol}(L)\sum_{c\in Y_0}D(c)
 \le \Bigl(\prod_v j_{b_0s}\Bigr)\sum_{\sigma\in L^\vee}F(\sigma).
\end{equation}
\leanref{Formalism.\allowbreak{}FourierDecayAlt.\allowbreak{}P.\allowbreak{}PBC.\allowbreak{}herglotz\_\allowbreak{}germ}{Formalism.\allowbreak{}HermitianCurvature.\allowbreak{}HermitianCurvatureProofs.\allowbreak{}PBCStep}{97--161}
\leanref{Formalism.\allowbreak{}FourierDecayAlt.\allowbreak{}P.\allowbreak{}PBC.\allowbreak{}germ\_\allowbreak{}global}{Formalism.\allowbreak{}HermitianCurvature.\allowbreak{}HermitianCurvatureProofs.\allowbreak{}ImageProdLoc}{55--129}
\leanref{Formalism.\allowbreak{}FourierDecayAlt.\allowbreak{}P.\allowbreak{}PBC.\allowbreak{}germ\_\allowbreak{}step}{Formalism.\allowbreak{}HermitianCurvature.\allowbreak{}HermitianCurvatureProofs.\allowbreak{}GermStep}{16--157}

For the pointwise shortening bound use square lattices of spacing \(R\) in the complex coordinates. The normalized dual sum, with inverse-difference scales \(z_v\), satisfies
\begin{equation}\label{eq:p8c1-planar-sum}
 R^{-2f}\sum_{\sigma\ {\rm dual}} M_z(g\sigma)
 \le R^{-2f}\prod_v C\{1+(Rz_v/g)^2\}.
\end{equation}
Beyond the radius \(\max(1,Rz_v/g)\) in integer lattice coordinates, the dyadic shell count grows as \(2^{2j}\), and the summand decays as \(2^{-5j}\). The same proof works for degree three. As \(R\) increases, the bound tends to at most \(C^f\prod_v(z_v/g)^2\). Restoring the section factors gives the volume
\begin{equation}\label{eq:p8c1-block-volume}
 \mathrm{vol}_{b_0,a}=\frac{(\prod_v z_{b_0,a,v})^2}
 {N(g^{-1}\mathfrak d^{-(2r+3-2b_0-a)-1})}\le\Theta^{-2f}
\end{equation}
for the step at length \(a\). This follows from \eqref{eq:p8c1-thresholds} and the divisors defining \(z_{b_0,a,v}\). The singleton integral bound starts the induction. Successive pointwise estimates on the square lattices give the bound \(C^{f(s-1)}\prod_{a=1}^{s-1}\mathrm{vol}_{b_0,a}\) at length \(s-1\) times an average of earlier energies. At the last step use \(L=L_{b_0s}\) in \eqref{eq:p8c1-herglotz}. The dual sum includes the zero term and is bounded by \(C^f\) by \ref{B8.10aa} when the last volume is at most one, enlarging the constant at the cutoff. The actual covolume \(\operatorname{covol}(L_{b_0s})=\prod_v j_{b_0s}\) cancels the transform factors in \eqref{eq:p8c1-herglotz}. Thus the coset sum costs at most \(C_4^{fs}\prod_{a=1}^{s-1}\mathrm{vol}_{b_0,a}\) times the previous energy bound. The averages over shorter words have mass one. Including the inverse-difference terms and the arithmetic majorants, the cost for this length is at most
\[
 \eta_i (1+3^fe) C_4^{fs}\prod_{a=1}^{s-1}\mathrm{vol}_{b_0,a}.
\]
\leanref{Formalism.\allowbreak{}FourierDecayAlt.\allowbreak{}P.\allowbreak{}PBC.\allowbreak{}gsup\_\allowbreak{}succ}{Formalism.\allowbreak{}HermitianCurvature.\allowbreak{}HermitianCurvatureProofs.\allowbreak{}PBCFinal}{18--46}
\leanref{Formalism.\allowbreak{}FourierDecayAlt.\allowbreak{}P.\allowbreak{}PBC.\allowbreak{}sumFormConcrete\_\allowbreak{}of}{Formalism.\allowbreak{}HermitianCurvature.\allowbreak{}HermitianCurvatureProofs.\allowbreak{}PBCFinal}{107--218}

Take square roots, sum over last lengths and divide by \(\prod_{j\le i}\sqrt{\eta_j}\). Use \(R_i\) for the resulting bound with the maximum taken over words reachable with the remaining letter budget. A singleton uses at most five letters in either one estimate and at most ten in their composition. A block of length \(s\) uses at most \(5(s+1)\) letters including the cross-term estimate, so a chain of total length \(r+1\) fits within \(N\). The singleton coefficient after taking the square root is at most \(\sqrt{(1+3^fe)(1+2^fe)}\le1+3^fe\). By \eqref{eq:p8c1-small-frequency-error}, \(3^fe=(3C_1\Theta^{-5/2})^f\). For \(s\ge2\), equation~\eqref{eq:p8c1-block-volume} bounds the square-root coefficient by \(\sqrt{1+3^fe}\,C_4^{fs/2}\Theta^{-f(s-1)}\). In dividing by the product of the \(\sqrt{\eta_j}\), the factor \(\sqrt{\eta_i}\) cancels. The skipped levels have \(\eta_j\ge1\). Take \(C\ge\max(3C_1,2\max(1,C_4))\) and the cutoff so that \(\Theta\ge1\) and \(3^fe\le1\). Since \(s/2\le s-1\), the long-block coefficient is at most \(((C/\Theta)^f)^{s-1}\), and the singleton coefficient is at most \(1+(C/\Theta)^f\). Summing the contributions for the possible last lengths gives
\begin{equation}\label{eq:p8c1-finite-energy-recursion}
 R_i\le(1+h)R_{i-1}+\sum_{s=2}^i h^{s-1}R_{i-s},
 \qquad h=(C/\Theta)^f.
\end{equation}
For \(h\le1/2\) the running maximum grows by at most \(1+h+h/(1-h)\) per step, so its square grows by at most \(\exp(6(r+1)h)\). Increase \(C\) to absorb the linear bound for the exponential error, then choose the cutoff so that \((r+1)h\le1/6\) and \(C^{f+1}(r+1)\le k^{f/32}\). With \(\Theta\ge k^{1/16}\) and \(C_f\ge256,\ C_f\log b\le f\log k\), the error is at most \(k^{-f/32}\le b^{-8}\).

At a fixed tower let \(\epsilon\downarrow0\). The finitely many normalized level-zero energies and their maximum tend to \(I_N\). Enlarge the tower after taking this limit. For each \(\rho\in S\) the error in the row vector is bounded by \(\sum_{p\notin X_{r+1}}|c_{\rho p}|\|V_p\|\), which tends to zero by absolute convergence of the orbital expansion. The finite sum over \(S\) therefore converges and satisfies the bound below.
\begin{equation}\label{eq:p8c1-flag-energy}
 \delta_s=C(r+1)(C/\Theta)^f\le b^{-8},\qquad
 \sum_{\rho\in S}\left\|\sum_p c_{\rho p}V_p\right\|^2
 \le (1+\delta_s)e_1 I_N .
\end{equation}
\leanref{Formalism.\allowbreak{}FourierDecayAlt.\allowbreak{}P.\allowbreak{}recursion\_\allowbreak{}with\_\allowbreak{}densities}{Formalism.\allowbreak{}HermitianCurvature.\allowbreak{}HermitianCurvatureProofs.\allowbreak{}RecursionWithDensities}{1162--1163}

At a finite place of residue cardinality \(q\) the logarithm of a local represented density is at most \(2q^{-(k-m)}\), and the normalizer is at least \(\exp(-2q^{-(k-m)})\). There are at most \(f\) places above a rational prime. For \(j=k-m\ge2\), the integral comparison gives \(\sum_{n\ge2} n^{-j}\le3\cdot2^{-j}\). Thus
\begin{equation}\label{eq:p8c1-density-product}
 \log e_1\le12f\,2^{-(k-m)} .
\end{equation}
\leanref{Formalism.\allowbreak{}L2.\allowbreak{}SG.\allowbreak{}Pf.\allowbreak{}DensLoc.\allowbreak{}densLocBound\_\allowbreak{}at}{Formalism.\allowbreak{}LocalToGlobal.\allowbreak{}SiegelLatticeTheory.\allowbreak{}CokernelGLCounts.\allowbreak{}RankPointCounts.\allowbreak{}DensLoc}{134--203}

For the mass conversion let \(K_V\) be the orbital quadratic form, \(c,c'\) the individual unfolding constants, and \(\rho_0\) the normalized row mass. The contribution is \(\Gamma_1\rho_0K_V\). Unfolding by scalar-partner classes and using \eqref{eq:p8c1-flag-energy} gives \(cc'K_V\le n_0 A I_N\), with \(n_0\) the count normalizer and \(A=(1+\delta_s)e_1\). Write \(V_{ZZ}\) for the real coordinate volume, \(V_{\rm loc}\) for the unit-level volume with its real factor, \(C_\infty\) for the Jacobian, \(V_B\) for the stabilizer volume factor, \(M\) for the product mass and \(\mu_K\) for the initial level measure. The measure and row identities are
\begin{equation}\label{eq:p8c1-mass-conversion}
 I_N=\frac{V_{ZZ}I_{\rm orb}}{V_{\rm loc}C_\infty},\qquad
 I_{\rm orb}\le\deg^{-1},\qquad
 n_0\rho_0=\frac{V_B C_\infty M}{V_{ZZ}\mu_K}.
\end{equation}
Here \(I_{\rm orb}\) is the orbital integral in isometric column coordinates, and \(\deg\) the formal degree. The middle bound is the squared coefficient orthogonality with this degree. Ordinary places contribute \(\mu_K^2/(\mu_S\mu_{S'})\) to \(M/V_{\rm loc}\), where \(\mu_S,\mu_{S'}\) are the two stabilizer measures. At the level places the natural measures differ from the typed ones by the finite ratio \(|\GL_k(\F_q)|/|U_k(\F_q)|\) for the level and its stabilizers. With double primes indicating the typed measures this gives
\[
 \frac{M}{V_{\rm loc}\mu_K}=\frac{\mu_{K''}}{\Gamma_1\mu_{S''}\mu_{S''}'},\qquad
 c^{-1}=\frac{d_u^S\mu_{S''}}{V_{H,\infty}},\qquad
 (c')^{-1}=\frac{d_v^S\mu_{S''}'}{V_{H,\infty}'} .
\]
The two real column volumes satisfy \(V_{H,\infty}V_{H,\infty}'=V_B V_G\), where \(V_G\) is the real group volume factor. Substitution into \eqref{eq:p8c1-mass-conversion} gives
\[
 \Gamma_1\rho_0 K_V\le \frac{\Gamma_1 A V_B}{cc'\deg}\frac{M}{V_{\rm loc}\mu_K}
 =\frac{A d_u^S d_v^S\mu_{K''}}{\deg V_G}=A\Delta_{u,v},
 \qquad D^\star=\deg V_G/\mu_{K''}.
\]
\leanref{Formalism.\allowbreak{}FourierDecayAlt.\allowbreak{}P.\allowbreak{}GrowthPolynomialOperator.\allowbreak{}t\_\allowbreak{}MainRow}{Formalism.\allowbreak{}HermitianCurvature.\allowbreak{}HermitianCurvatureProofs.\allowbreak{}GrowthOfBasic}{1070--1184}
\end{proof}

\subsection*{Proof of \ref{B8.10f}}\label{prf:B8.10f}\phantomsection
\begin{proof}
At a place over \(g_j\) put \(q=g_j\), \(G=\GL_k(\F_q)\), and let \(U\) be its upper unitriangular subgroup. The primes split completely, so the residue fields are prime fields. Let \(\xi(u)=\psi(\sum_{i<k}u_{i,i+1})\), where the selected additive character \(\psi\) is nontrivial. Extend \(\xi\) by zero off \(U\). The local factor in the finite Gram kernel is the induced character
\[
 \gamma(x)=|U|^{-1}\sum_{y\in G}\xi(yxy^{-1}),
 \qquad \gamma(1)=|G|/|U|.
\]
On a free level-group orbit there is exactly one group element carrying one specified index to another. In orbit coordinates the matrix is therefore convolution by \(\gamma\), up to complex conjugation. Entries between different orbits vanish.

On an irreducible summand with character \(\theta\), central convolution by \(\overline{\gamma}/\gamma(1)\), using counting measure, has scalar
\[
 \frac{\sum_{x\in G}\overline{\gamma(x)}\theta(x)}
      {\gamma(1)\theta(1)}
 =\frac{|U|}{\theta(1)}
   \dim\operatorname{Hom}_{U}(\xi,\theta|_U).
\]
Complex conjugation of the convolution kernel permutes the irreducibles and preserves this list of real eigenvalues. To compute the multiplicity, factor each monic degree-\(k\) polynomial with nonzero constant term. By \ref{B7.23} its irreducible factors with repetitions specify a multiset of cuspidals. By \ref{B7.24} its induced product contains a generic constituent, and \ref{B7.25} gives different constituents for different multisets. By \ref{B7.29} this accounts for \(q^k-q^{k-1}\) positive multiplicities. Their squares, summed over all irreducibles, total exactly \(q^k-q^{k-1}\) by \ref{B7.26}. Every positive multiplicity is one.

For the dimension bound, take a representation containing the specified Whittaker vector. In rank \(a+1\) restrict to the abelian last-column subgroup \(\F_q^a\). The restriction of the Whittaker character to this subgroup is nonzero. The \(q^a-1\) nonzero characters form one \(\GL_a(\F_q)\)-orbit, and their eigenspaces have the same dimension. The stabilizer of the specified character contains the rank-\(a\) mirabolic subgroup. Its upper unitriangular subgroup acts on the original vector by the rank-\(a\) Whittaker character. Repeat the last-column decomposition inside this eigenspace. At rank one the dimension is at least one. Multiplication gives \(\theta(1)\ge\prod_{a=1}^{k-1}(q^a-1)\). This argument also applies in the dual realization of the Whittaker functional.
\leanref{PI.\allowbreak{}GGMirBoundW.\allowbreak{}mirabolicBound\_\allowbreak{}of}{Formalism.\allowbreak{}PeriodIntegrals.\allowbreak{}UnitaryGroupMassFormula.\allowbreak{}MirabolicBound.\allowbreak{}Development}{112--128}

Since \(|U|=q^{k(k-1)/2}\), the normalized local eigenvalues are bounded by
\[
 \prod_{a=1}^{k-1}(1-q^{-a})^{-1}
 \le \left(1-\sum_{a\ge1}q^{-a}\right)^{-1}
 =1+\frac1{q-2}\le1+\frac3q .
\]
Here \(q\ge3\) by the level hypotheses. Take the tensor product over the \(f\) places above each distinguished prime. All eigenvalues are nonnegative. Restoring the product of the identity values, \(\Gamma_1\), gives \(0\preceq G_{\rm fin}\preceq P_g\Gamma_1 I\).

For a finite term set the norm expansion contracts \(G_{\rm fin}\) against the vectors \(w_i(E_1)\). In an orthonormal basis of their span, apply the matrix inequality to the coefficient vector for each basis coordinate and sum. The result is \(P_g\Gamma_1\sum_i\|w_i(E_1)\|^2\). Absolute summability gives convergence of both the orbital vector sums and the operator sum over finite subsets of \(E_1\). Continuity of the squared norms gives the same bound on \(E_1\).
\end{proof}

\subsection*{Proof of \ref{B8.10b}}\label{prf:B8.10b}\phantomsection
\begin{proof}
Let \(G\) be the unitary group of the selected row, \(F\) the real base field, and \(\Gamma=G(F)\). Write \(S,S'\subset\Gamma\) for the rational stabilizers. The kernel on column vectors is the pairing of their Poincar\'e series divided by \(\sqrt{c_uc_v}\), where \(c_u,c_v\) are the unfolding constants.

For compact \(\Omega,C'\subset G(\mathbb A)\), the translate multiplicity into \(\Omega\) at \(c(g,1)\) is at most \(C\|g\|_\infty^N\) for \(c\in C'\), with \(N\) as in \ref{B8.02y}. The level charts bound the rational translation parameters by archimedean norm balls with finite components in a fixed compact set. The stabilizer fundamental domain and the level ties express the resulting absolute integral as a finite-level integral of the real section. Shifting the representative through its stabilizer gives the quotient weight of \ref{B8.02y}, which gives finiteness. The compact reproducer of that lemma bounds the absolute section by a constant times its local average. Folding this bound over the compact neighbourhood gives \(L^2\) bounds for both absolute Poincar\'e series. Their product is integrable by Cauchy--Schwarz.

Unfold the first series over \(S\backslash G(\mathbb A)\) and partition the second by double cosets in \(S\backslash\Gamma/S'\), indexed after inversion. The translates for \(a\) tile a fundamental domain for \(S\cap aS'a^{-1}\). The integral pairs the first section at \(x\) with the second at \(a^{-1}x\). Tonelli gives absolute unfolding. Partition the finite components by the row lattices and stabilizers. The finite-character matrix and real orbital integral give \(\operatorname{op}(t)\). The sharp coordinates span the column spaces, so the term pairings recover the kernel form and absolute convergence permits grouping by cells.
\leanref{Formalism.\allowbreak{}L2.\allowbreak{}SG.\allowbreak{}Pf.\allowbreak{}Slab.\allowbreak{}dcSingleE}{Formalism.\allowbreak{}LocalToGlobal.\allowbreak{}SiegelLatticeTheory.\allowbreak{}ClosedLocalOrbits.\allowbreak{}GaugeLatticeShifts.\allowbreak{}DoubleCosetSingleUnfolding}{31--76}
\leanref{Formalism.\allowbreak{}L2.\allowbreak{}SG.\allowbreak{}TermOp.\allowbreak{}KExpCut.\allowbreak{}kExpansion\_\allowbreak{}of\_\allowbreak{}steps}{Formalism.\allowbreak{}LocalToGlobal.\allowbreak{}SiegelLatticeTheory.\allowbreak{}ClosedLocalOrbits.\allowbreak{}GaugeLatticeShifts.\allowbreak{}KExpCut}{773--790}
\leanref{Formalism.\allowbreak{}L2.\allowbreak{}SG.\allowbreak{}Closed.\allowbreak{}KExpansionC\_\allowbreak{}holds}{Formalism.\allowbreak{}LocalToGlobal.\allowbreak{}SiegelLatticeTheory.\allowbreak{}ClosedVolumeStatements.\allowbreak{}UnipotentInLevelSubgroup}{63--67}

In the identity cell conjugate the stabilizers into common coordinates and lift to the compatible level. The local type pairing is zero for distinct words. Summing in the lifted coordinates gives that the sum over the identity cell is zero in the continuous Hilbert--Schmidt embedding. Other terms with zero Gram value have zero orbital integral if irrelevant. For a relevant zero Gram term, the orbital integral is covariant for the real common pointwise stabilizer of the first frame and the translated second frame. Let \(W\) be the span of these two frames at a real place \(v_0\). The zero Gram normal form lets us subtract vectors of the respective isotropic flag spaces from the two orthonormal definite frames, so that the resulting frames \((e_a)\) and \((e'_a)\) lie in \(W^\perp\). They remain orthonormal because the subtracted vectors are isotropic and perpendicular to the corresponding original definite frame.

Put \(K_{ab}=\langle e_a,e'_b\rangle\) in Hermitian coordinates. To see \(K^*K=I\), use the first frame's hyperbolic partners to subtract from each \(e'_b\) an isotropic first-flag vector and project to the span of the first definite frame. The coefficients of the projection are \(K_{ab}\). Replacing the first definite vectors by \(e_a\) changes the residual vector within the first isotropic flag. This residual vector lies in \(W^\perp\), hence in the intersection of the two isotropic flag spaces by the zero Gram normal form. Its pairing with \(e'_a\) vanishes. Taking inner products with \(e'_a\) gives \(K^*K=I\).

For each \(M\in U(k)\), extend its action on the positive span of \((e_a)\) by the identity on the orthogonal complement. This fixes \(W\) pointwise and so gives an element \(h\) of the common stabilizer at \(v_0\). Take identity components at the other real places. The two actions on definite type coordinates are \(M\) and \(K^*MK\), since \(h\) fixes the isotropic corrections and acts on the second corrected frame by its matrix coefficients against the first. Covariance now yields an intertwiner after taking matrix entries on the tensor factors away from \(v_0\). Conjugation by \(K\) identifies the second action with its original \(U(k)\)-type. The two highest weights at \(v_0\) differ, so every such matrix entry vanishes by Schur's lemma. The orbital integral is zero, and the full operator equals \(K_0\).
\leanref{Formalism.\allowbreak{}L2.\allowbreak{}SG.\allowbreak{}FixedDecomposition.\allowbreak{}OOt\_\allowbreak{}eq\_\allowbreak{}zero0}{Formalism.\allowbreak{}LocalToGlobal.\allowbreak{}SiegelLatticeTheory.\allowbreak{}CokernelGLCounts.\allowbreak{}RankPointCounts.\allowbreak{}FixedDecomposition}{744--753}
\leanref{PI.\allowbreak{}SchurPlaces.\allowbreak{}Geo.\allowbreak{}geomB\_\allowbreak{}of}{Formalism.\allowbreak{}LocalToGlobal.\allowbreak{}SiegelLatticeTheory.\allowbreak{}CokernelGLCounts.\allowbreak{}RankPointCounts.\allowbreak{}SubOneMulVecFix}{207--508}
\leanref{PI.\allowbreak{}SchurPlaces.\allowbreak{}schurNFC\_\allowbreak{}of}{Formalism.\allowbreak{}LocalToGlobal.\allowbreak{}SiegelLatticeTheory.\allowbreak{}CokernelGLCounts.\allowbreak{}RankPointCounts.\allowbreak{}SchurNFCOf}{20--30}

Apply \ref{B8.10f} in the generic and nongeneric cells. Within each finite-type fibre, different rows are orthogonal. In a same-row pair with first term of nonzero Gram value the \(2r+2\) frame vectors of \(F_0,F_t\), \(r=r_S\), are linearly independent. The pair space lies in their perpendicular in dimension \(k+2r+1\), so it has dimension at most \(k-1\). The corank is positive even for identical terms.
\leanref{PI.\allowbreak{}PairCorankPos.\allowbreak{}W.\allowbreak{}corankPosG}{Formalism.\allowbreak{}PeriodIntegrals.\allowbreak{}UnitaryGroupMassFormula.\allowbreak{}PairCorankPositivity.\allowbreak{}FinrankOrthAddLe}{122--131}
\leanref{Formalism.\allowbreak{}L2.\allowbreak{}SG.\allowbreak{}Pf.\allowbreak{}SatGram.\allowbreak{}corankPos\_\allowbreak{}tO}{Formalism.\allowbreak{}LocalToGlobal.\allowbreak{}SiegelLatticeTheory.\allowbreak{}ClosedLocalOrbits.\allowbreak{}GaugeLatticeShifts.\allowbreak{}SatGram}{431--443}

Put \(\ell=e^{-kf\log k/3}\) and \(n=2^{-(k-2r-5)}e^{C_nm^5f\log b}\), with \(C_n\ge C_4\) of \ref{B8.10d}. The generic scalar contribution is bounded by \(A\Delta_{u,v}\) by \ref{B8.10e}. Each loop contribution is bounded by \(\ell\Delta_{u,v}\) by \ref{B8.10a}, since its absolute sum is over a subset of the union in that lemma. In the nongeneric cell the higher-corank pairs are in \(S_{\rm big}\). Apply \ref{B8.10d} to the signed corank-one sum. The source hypotheses give the counting and pair standing conditions, with ties and convergence from the selection. The triangle inequality gives
\begin{equation}\label{eq:p8c1-distinct-radicands}
 \|K_0\|_{\rm HS}^2
 \le\bigl(\sqrt{P_g(A+2\ell)}+\sqrt{P_g(\ell+n)}\bigr)^2\Delta_{u,v}.
\end{equation}
\leanref{GlobalTheory.\allowbreak{}NGRow.\allowbreak{}rowNG\_\allowbreak{}of\_\allowbreak{}cut2P}{Development.\allowbreak{}AbelianVarietyTheory.\allowbreak{}Proofs.\allowbreak{}FlagMeasures.\allowbreak{}NonGenericCorankOnePairsBoundOfSplitFiniteRealRo}{20--37}

Set \(z=b^{-7}\). The thresholds give \(k\ge\max(32,4r+12,70\log b)\), \((kb)^8\le g_j,\ f\le k,\ 8C_nm^5f\log b\le k\). Then \(P_g\le e^{6z}\). By \eqref{eq:p8c1-density-product}, \(\log e_1\le12k\,2^{-k/2-4}\le e^{-k/5}\). Also \(\log n\le-(k/2+1)\log2+k/8\le-k/5\) and \(\ell\le e^{-k/5}\). Thus \(\log e_1,n,\ell\le z^2\) and \(A\le(1+z)e^{z^2}\). For \(b\ge16\) the radicands in \eqref{eq:p8c1-distinct-radicands} are at most \(1+11z\) and \(3z^2\). Their roots sum to at most \(1+(11/2+7/4)z\), whose square is at most \(1+16z\le1+b^{-6}\).
\end{proof}

\section*{Proofs for Chapter~\ref{ch:p8c2}: Theta correspondence, test vectors and evaluation}

\subsection*{Proof of \ref{B8.11h}}\label{prf:B8.11h}\phantomsection
\begin{proof}
For a finite place \(w\), put \(q_w=Nw\). Let \(c_w=\chi_w(\varpi_w)\) if \(\chi_w\) is unramified, and let \(c_w=0\) otherwise. Unitarity gives \(|c_w|\le1\). For real \(\sigma>1\),
\[
 \sum_w\sum_{h\ge1}\frac{|c_w|^h}{h q_w^{h\sigma}}
 \le \frac{[E:\Q]}{1-2^{-\sigma}}\sum_{n\ge2}n^{-\sigma}<\infty .
\]
Indeed there are at most \([E:\Q]\) places over each rational prime \(p\), and their norms are at least \(p\). Exponentiating the logarithmic Euler series gives \(L(s,\chi)\ne0\) for \(\Re s>1\).

The angular character at a place with \(m_v\ne0\) is nontrivial on the unit circle. Thus neither \(\chi\) nor \(\chi^{-1}\) is a power of the idelic norm. Tate's global functional equation and local factors give entire completed functions for both characters. Dividing out their archimedean factors gives entire finite \(L\)-functions, since \(1/\Gamma\) is entire. For the trivial character the same theorem gives a simple pole of \(\zeta_E\) at one. The pole of its completion at zero is cancelled in the finite function by the reciprocal gamma factor at a complex place. \cite[Th.~4.4.1 and \S4.5]{tate-1967}

Suppose \(L(1,\chi)=0\). Then
\[
 F(s)=\zeta_E(s)^2 L(s,\chi)\overline{L(\bar s,\chi)}
\]
is entire. For \(\Re s>1\) its Euler products give
\[
 F(s)=\exp\!\left(\sum_w\sum_{h\ge1}a_{wh}q_w^{-hs}\right),
 \qquad a_{wh}=\frac{2+2\Re(c_w^h)}{h}\ge0.
\]
Expand the exponent at \(s=2-t\). For \(|t|<1\) its series converges locally absolutely. Its Taylor coefficients in \(t\) are nonnegative, since
\(q_w^{-h(2-t)}=q_w^{-2h}\exp(th\log q_w)\). The Taylor coefficients of \(F(2-t)\) dominate those of the exponent term by term: in the exponential series the linear term gives the exponent and all other terms have nonnegative coefficients. The Taylor series of the entire function \(F(2-t)\) converges at \(t=2\). Summing the dominated coefficients there and interchanging nonnegative sums gives
\[
 \sum_{w,h}a_{wh}
 =\sum_{w,h}a_{wh}q_w^{-2h}
       \sum_{n\ge0}\frac{(2h\log q_w)^n}{n!}<\infty .
\]
Fix any \(w\). For \(z=c_w^h\) and \(u=\Re z\),
\[
 h(a_{wh}+a_{w,2h})
 =3+2u+\Re(z^2)
 \ge 2+2u+2u^2\ge\tfrac32 .
\]
The two series on the left after division by \(h\) are summable, whereas \(\sum_{h\ge1}3/(2h)\) diverges. This proves nonvanishing at one and applies also to \(\chi^{-1}\).
\leanref{GlobalTheory.\allowbreak{}landau\_\allowbreak{}exp\_\allowbreak{}summable}{Development.\allowbreak{}AbelianVarietyTheory.\allowbreak{}Proofs.\allowbreak{}AdelicLatticeProofs.\allowbreak{}LandauSummableOfExpDirichletSeriesEntire}{84--293}
\leanref{GlobalTheory.\allowbreak{}eulerProduct\_\allowbreak{}ne\_\allowbreak{}zero\_\allowbreak{}one\_\allowbreak{}of\_\allowbreak{}entire}{Development.\allowbreak{}AbelianVarietyTheory.\allowbreak{}Proofs.\allowbreak{}AdelicLatticeProofs.\allowbreak{}EulerProductNeZeroOneOfEntire}{95--128}
\leanref{GlobalTheory.\allowbreak{}CMVal.\allowbreak{}IsHeckeL.\allowbreak{}ne\_\allowbreak{}zero\_\allowbreak{}nat}{Development.\allowbreak{}AbelianVarietyTheory.\allowbreak{}Proofs.\allowbreak{}DivisorProofs.\allowbreak{}HeckeLNeZeroNat}{10--22}

Now let \(k\le0\) satisfy the stated inequalities. At \(1-k\ge1\) the finite value for \(\chi^{-1}\) is nonzero. Its archimedean arguments \(1-k+|m_v|/2\) are positive. The functional equation therefore gives a nonzero completed value for \(\chi\) at \(k\), since the conductor factor and the root number are nonzero. Each \(\Gamma_\C(k+|m_v|/2)\) is finite and nonzero by the strict inequality at that place. Division gives \(L(k,\chi)\ne0\), including \(k=0\).

Finally, omitting a finite set \(S\) multiplies the continued finite \(L\)-function by \(\prod_{w\in S}(1-c_wq_w^{-s})\). This identity follows first from the Euler product and then by continuation. At \(s=k\ge1\) every factor is nonzero because \(|c_w|q_w^{-k}<1\).
\end{proof}

\subsection*{Proof of \ref{B8.11i}}\label{prf:B8.11i}\phantomsection
\begin{proof}
Split the variables into the \(R\) positive columns and the \(S\) negative columns. The source action on their tensor product is the sum of the two actions, which we denote by \(D_+\) and \(D_-\). We use the following polynomials \(x,y\) in the respective sets of columns. A determinant of size zero means \(1\).
\[
\begin{array}{c|l|l}
 &x&y\\ \hline
 T&
 \det(\text{positive rows, first }P\text{ positive columns})^n&
 \det(\text{negative rows, first }Q\text{ negative columns})^m\\
 T'&
 \det(\text{negative rows, last }Q'\text{ positive columns})^n&
 \det(\text{positive rows, last }P'\text{ negative columns})^m
\end{array}
\]
The identities \(R=P+Q'\), \(S=Q+P'\) ensure that all four minors exist.

We show that \(xy\) survives the source relations. Use \(h=c\) in the first row of the table and \(h=-c\) in the second. Let \(J\) be the diagonal source signature matrix. On the negative-column polynomials introduce the action
\[
 \widetilde D_-(Z)=h\operatorname{tr}Z-D_-(JZ^tJ).
\]
For either row, \(x\) under \(D_+\) and \(y\) under \(\widetilde D_-\) have the same character on the block diagonal source algebra. Indeed, on the positive and negative row blocks the vacuum weights for \(D_+\) are \(R/2,-R/2\), and those for \(D_-\) are \(-S/2,S/2\). In the first row the minors add \((n,0)\) and \((0,m)\) to these pairs. In the second they add \((0,-n)\) and \((-m,0)\). In both cases the sums are \((h,h)\). Off-diagonal units within either block act by zero on the determinants. The lower left source units kill \(x\) under \(D_+\) and \(y\) under \(\widetilde D_-\). In each case the relevant operator on the indicated half of the variables is a sum of second derivatives, one in each row sign, and the minor uses one row sign.

Give the polynomial spaces their positive Hermitian Fischer forms, with orthogonal monomials of squared norm \(\alpha!\) in multi-degree \(\alpha\). For both source actions the adjoint of \(E_{ij}\) is \(J_{ii}J_{jj}E_{ji}\), by the formulas for \(A_l,B_l\). Consequently the map \(Ux\mapsto Uy\), for words \(U\) in source matrices with the respective actions, defines an intertwiner of the cyclic spaces. To see that it respects every relation, normalize the two starting vectors to have norm one and compute the Gram entry of \(Ux,Vx\) by commuting the word \(U^*V\) into upper right, block diagonal, lower left order. Terms with a lower left factor on the right vanish. Terms with an upper right factor on the left have zero pairing with \(x\), by adjunction. The remaining terms evaluate in the common block diagonal character. This gives identical Gram matrices for the two cyclic spaces. Rescaling the intertwiner sends the original \(x\) to \(y\).

The cyclic space of \(x\) is graded by polynomial degree. On the positive columns each off-block source matrix unit is homogeneous of degree \(2\) or \(-2\), and block diagonal units preserve degree. Orthogonal projection onto the cyclic space therefore exists on polynomials, degree by degree. Its orthogonal complement is invariant by the adjoint identity. Composing this projection with the intertwiner gives a map \(L\) on all positive-column polynomials. For the bilinear Fischer form \([\, ,\,]\), with the same diagonal monomial entries \(\alpha!\), the functional
\[
 u\otimes v\longmapsto [Lu,v]
\]
has source character \(h\operatorname{tr}\). Indeed the transpose of \(D_-(Z)\) for this bilinear form is \(D_-(JZ^tJ)\). The functional takes \(x\otimes y\) to \([y,y]>0\), since \(y\) is nonzero with real coefficients. This proves nonvanishing in both quotients.

We compute the generating compact type in the untwisted polynomial models. For the pair \((\U(b,e),\U(R,S))\) the source vacuum shift is \(((R-S)/2,-(R-S)/2)\) on the two blocks. The target shift is \(((b-e)/2,-(b-e)/2)\). Subtracting the source shift from the constant weight \(h\) gives \((n^{[P]};m^{[Q]})\) on the first side and \(((-m)^{[P']};(-n)^{[Q']})\) on the second. The joint-harmonic correspondence gives the following target weights after subtraction of their respective vacuum shifts:
\[
 \big(n^{[P]},0^{[Q']};m^{[Q]},0^{[P']}\big),\qquad
 \big(0^{[P]},(-n)^{[Q']};0^{[Q]},(-m)^{[P']}\big).
\]
These are the two rectangular cases of \cite[Lemma~2.5(1)]{paul-1999}. Its length conditions here are \(P,Q'\le R\) and \(Q,P'\le S\). The displayed minors give the joint highest vectors: the upper compact units annihilate them, and every quadratic contraction differentiates a pair of variables of which at least one is absent. Their products have degrees \(Pn+Qm\) and \(Q'n+P'm\), respectively, the minimum degrees in \cite[Lemma~2.5(2)]{paul-1999}. For each pair the source character has a single compact type. The generation statement of Proposition~\ref{B1.15b}, or \cite[\S4, equation~(4.5)]{howe-1989}, therefore gives generation by the image of this joint type. We have proved that its highest vector has nonzero image.

Finally, put \(d'=(P'-Q')/2\). Since \(d-d'=(R-S)/2\), one has \(c+d'=d+n\) and \(c-d'=m-d\). Adding the vacuum shifts and the primed trace twist now gives the asserted equality at every coordinate.
\end{proof}

\subsection*{Proof of \ref{B8.11j}}\label{prf:B8.11j}\phantomsection
\begin{proof}
If \(S=0\), then \(\mathfrak p=0\). Otherwise fix one side. Let \(d_0\) be its minimum source-character degree, computed in Lemma~\ref{B8.11i}. All polynomials of degree less than \(d_0\) map to zero: their finite-dimensional source compact modules contain no copy of the source character, and the quotient map transforms by that character. Elements of \(L\) have representatives of degree \(d_0\). An off-block target unit acts as the sum of a quadratic multiplication and a quadratic differentiation. Thus \(\mathfrak pL\) is represented in degree \(d_0+2\).

In that degree take the orthogonal complement, for the Fischer form, of the images of the degree-raising parts of the off-block source units. Those images map to zero. Indeed, on a polynomial of degree \(d_0\) the full off-block source unit acts by zero in the quotient and its lowering part has degree \(d_0-2\) or vanishes. The complement is the joint kernel of the adjoint contractions. Explicitly these contract a positive and a negative source row over the \(R\) columns, or over the \(S\) columns separately. The complement is target compact stable, and its image contains \(\mathfrak pL\). Hence each irreducible type of \(\mathfrak p L\) has a representative of that type in this kernel, by complete reducibility in finite degree.

For such a type of weight \(\mu\), subtract the target vacuum shift and, on the primed side, the trace twist. Write \(\widehat\mu\) for the resulting pair of strings. The harmonic condition on the \(R\) columns and on the \(S\) columns gives
\[
 \sum_i|\widehat\mu_i|=d_0+2.
\]
More explicitly, in one block with \(t\) columns and with \(b_+,b_-\) copies of the defining and dual column representations, the kernel of the row contractions is the joint-harmonic space for the auxiliary pair with compact member \(\U(t)\) and other member \(\U(b_+,b_-)\), in this polarization. A shifted column weight there has at most \(b_+\) positive and \(b_-\) negative entries. Its homogeneous degree is the sum of their absolute values. To see the degree directly in the joint-harmonic formula, form highest minors of sizes given by the column lengths of the positive partition, and dual highest minors for the negative partition. They use disjoint column indices, at the beginning and end of the ordered column basis, so every contraction vanishes. Their degree is the asserted sum. The compact-type correspondence and multiplicity one exhaust the joint harmonics of that column type. This applies separately in both column blocks, whose degrees add. The bounds on the numbers of entries and the multiplicity assertion are precisely \cite[Theorem~2.2(1) and Lemma~2.5(1)]{paul-1999}, applied with one member compact. For a source of signature \((b,e)\), the two pairs of bounds are \((b,e)\) and \((e,b)\).

Now \(\mathfrak p L\) is an image of \(\mathfrak p\otimes L\). The two summands of \(\mathfrak p\) are a defining representation of one block tensored with the dual defining representation of the other. Thus \(\mu=\lambda+e_i-e_j\) is block dominant, with \(i,j\) in opposite blocks. This one-box rule also follows by multiplying the Weyl character numerator in each block by \(\sum x_i\) or \(\sum x_j^{-1}\): the alternating determinants with two equal exponents vanish, and the others give the indicated dominant changes.

For \(\lambda\) the sum of absolute shifted entries is \(d_0\). To increase it by two, each of the two changed entries must increase its absolute value by one. On the unprimed side all allowed positive entries are already present in each block. Thus the addition is to a positive entry and the subtraction is from a zero entry. On the primed side all allowed negative entries are present. The subtraction is from a negative entry and the addition is to a zero entry. In the common weight \(\lambda\), both rules add to the initial constant string of one block and subtract from the final constant string of the other. Dominance forces \(i\) to be the first position of its block and \(j\) the last of its block. The two possible values of \(\lambda_i-\lambda_j\) are \(n+2d\) and \(m-2d\). Both are positive, since \(P>Q\) by \eqref{eq:p8c2-triple-range}, and \(m>R\ge P\).

On a block of size \(t\), \(C_K\) has eigenvalue \(\sum_{k=1}^t\lambda_k(\lambda_k+t+1-2k)\). This follows by commuting the positive matrix units past the negative ones on a highest vector. Adding at the first position of one block and subtracting at the last of the other changes the sum by \(D+2(\lambda_i-\lambda_j)>D\).
\end{proof}

\subsection*{Proof of \ref{B8.11k}}\label{prf:B8.11k}\phantomsection
\begin{proof}
Identify the two generating compact types by their highest weight and let \(L\subset T\oplus T'\) be the graph of this identification. Write \(M=U(\mathfrak{gl}_D)L\) with the componentwise action. Both projections are surjective by Lemma~\ref{B8.11i}. We shall prove that the \(b_0\)-eigenspace of \(C_K\) in \(M\) is exactly \(L\).

Filter \(M\) by the spaces \(M_k\) spanned by words of at most \(k\) off-block matrices applied to \(L\). They are finite-dimensional compact modules since \([\mathfrak k,\mathfrak p]\subset\mathfrak p\). Their union is \(M\). For \(k\ge1\) multiplication gives a compact-equivariant surjection
\[
 \mathfrak p^{\otimes k}\otimes L\longrightarrow M_k/M_{k-1}.
\]
It is symmetric in the \(k\) slots: exchanging two adjacent factors introduces their commutator in \(\mathfrak k\), which can be moved to the right and gives terms of smaller filtration degree. Choose compact-invariant positive Hermitian forms on the factors and identify the quotient with the orthogonal complement \(V_k\) of the kernel.

Let \(B\) be the kernel of the action \(\mathfrak p\otimes L\to\mathfrak pL\) inside \(M\). A tensor with its last slot and its \(L\)-factor in \(B\) maps to zero, and symmetry gives the same conclusion in every slot. Thus \(V_k\) belongs to \(B^\perp\) in each such pair of slots, with the other slots as tensor factors. The diagonal Casimir on that pair, denoted \(C_{jL}\) for slot \(j\), has eigenvalues greater than \(b_0+D\) on \(B^\perp\). Indeed \(B^\perp\simeq\mathfrak pL\), and the latter embeds in the sum of the two first images from Lemma~\ref{B8.11j}.

The Casimir on \(\mathfrak p\) equals \(D\): the trace-form Casimir on the defining or dual defining representation of a block of size \(t\) is \(t\). Expanding the diagonal Casimir on the whole tensor gives
\[
 C_{\rm total}-b_0
   =C_{\mathfrak p^{\otimes k}}
       +\sum_{j=1}^k\big(C_{jL}-D-b_0\big).
\]
The first term has nonnegative quadratic form, since \(E_{ab}^*=E_{ba}\) for the compact action and the Casimir is the sum of \(E_{ab}E_{ab}^*\). Every summand in parentheses has strictly positive quadratic form on a nonzero vector of \(V_k\), by the preceding slot calculation. Hence \(M_k/M_{k-1}\) has no \(b_0\)-eigenvector. Since \(M_0=L\), this proves the eigenspace assertion.

A proper submodule \(U\) of \(M\) intersects \(L\) in zero, since \(L\) is compact irreducible and generates \(M\). For \(u\in U\cap M_k\), projection to the \(b_0\)-eigenspace is a polynomial in \(C_K|_{M_k}\), by finite-dimensional compact complete reducibility. The projection therefore belongs to \(U\cap L\) and vanishes. The sum \(J\) of all proper submodules still has zero projection to \(L\), as each of its elements is a finite sum. It follows that \(J\) is proper and \(M/J\) is simple. The two coordinate kernels are proper submodules of \(M\), so both projections induce surjections \(T\to M/J\), \(T'\to M/J\). Their equivariance follows from the componentwise action used to define \(M\).
\end{proof}

\subsection*{Proof of \ref{B8.11l}}\label{prf:B8.11l}\phantomsection
\begin{proof}
Put \(d=a+a'\) and take \((n,\chi)=(a,\nu^{-1})\) or \((a',\nu)\). On each side average a real polynomial of nonzero quotient image from Lemma~\ref{B8.11i} over the compact source. Use the vacua off \(T\). At \(v\in T\) take the indicator of \(U=Y_0+\varpi_v^e\operatorname{Mat}_{n,d}(\mathcal O_v)\), with \(Y_0=(I_n\ 0)\). If \({}^tgU\cap U\ne\varnothing\), comparison of the first \(n\) columns gives
\[
 {}^tg=(1+\varpi_v^e A)(1+\varpi_v^e B)^{-1}
\]
with \(A,B\) integral. Such a matrix stabilizes \(U\). Choose \(e\) large enough that the determinant character is trivial on this subgroup. The coefficient is zero off the stabilizer by disjoint supports and equals its value at one on the stabilizer. On the domain integral off \(T\) the support is compact and includes an identity neighborhood, and its doubling height equals one. Write \(H\ne0\) for its measure times the coefficient at one.

Double on the source against the constant function. Multiplying the unramified local doubling integrals outside \(T\) for the trivial representation \cite[\S4.2.1]{eischen-harris-li-skinner-2020} yields the partial Euler product \(L_E^T(\chi,z+n/2-i)\) omitting places above \(T\). Since \(\chi=\nu^{\mp1}\) is ramified at both split places of \(E\) above every \(v\in T\), every local Euler factor of the primitive Hecke \(L\)-function \(L_E(\chi,s)\) above \(T\) equals \(1\), so \(L_E^T(\chi,s)=L_E(\chi,s)\) identically. This gives
\begin{equation}\label{eq:p8c2-bottom-doubling}
 Z(z)D_n^T(z)=H\prod_{i=0}^{n-1}L_E(\chi,z+n/2-i),\qquad
 D_n^T(z)=\prod_{v\notin T}\prod_{j=1}^n
 (1-\eta_v(\varpi_v)^{n-j}(Nv)^{-(2z+j)})^{-1}.
\end{equation}
The normalizer converges without zeros for \(\Re z>0\). Continue to \(z=(d-n)/2\) by Theorem~\ref{B8.12}. Write \(m_w\) for the angular exponent of \(\chi\) at an infinite place \(w\), so \(|m_w|=m_{\alpha,w}-m_{\beta,w}>d\). The arguments \(k=d/2-i\) are integers with \(2k\ge-d\), so \(k+|m_w|/2>0\) at every infinite place. Because \eqref{eq:p8c2-bottom-doubling} contains the primitive \(L\)-function \(L_E(\chi,k)\) with no nontrivial Euler factor omitted, Lemma~\ref{lem:p8c2-hecke-nonvanishing} gives a nonzero value at every such integer \(k\), including \(k\le0\). The definite Siegel--Weil value \eqref{eq:p8c2-definite-sw} and the doubled product identify this value with a nonzero multiple of the squared norm of the theta integral. This gives a witness on each side. For the primed witness use Lemma~\ref{B8.11f} at \(v_1\) to take a full-rank representative. Absolute convergence for the definite source and source invariance preserve its integral.
\end{proof}

\subsection*{Proof of \ref{B8.11m}}\label{prf:B8.11m}\phantomsection
\begin{proof}
Enlarge \(E_{\rm exc}\) to include all finite places where \(P\) is not a single pure-tensor factor in the split model. Fix \(v\notin E_{\rm exc}\) above a rational prime split completely in \(E\), fix all local coordinates of \(x\) and \(P\) away from \(v\), and fix one map \(U_j\). Realize \(\Pi_v\) on the source opposite to \(S\) by \eqref{eq:p8c2-split-matching} and Lemma~\ref{B8.11f}; for \(S=V'\) use the unprimed block to read \(\Pi_v\), and for \(S=V\) use the primed block. Since the entire image of \(U_j\) lies in \(R_v\), pairing \(U_j\) on the opposite local Schwartz space against \(\overline{\theta_{\omega_S}}\) in the \(v\)-slot defines a distribution in the two local Schwartz arguments. With \(d=a+a'\), in the split models the two local arguments of the pairing occupy complementary row blocks of a square \(d\)-by-\(d\) matrix. Complex conjugation in the reverse lift gives the right covariance \(\zeta_v(\det h)|\det h|_v^{-d/2}\) under simultaneous right translation by \(h\in\GL_d(F_v)\), where \(\zeta_v\) is \(\nu_v\) or \(\nu_v^{-1}\) according to the chosen side and split embedding. On invertible matrices the distribution is twisted Haar measure and has the same left covariance. For a fixed left translate the defect of this covariance is supported on singular matrices and has the original right covariance. On a singular right orbit a stabilizer can scale a kernel coordinate by \(\varpi_v\). A covariant distribution there pulls back locally to twisted Haar measure on the group. The stabilizer twist must equal the positive quotient modulus. Here \(\zeta_v(\varpi_v)|\varpi_v|_v^{-d/2}\notin\R_{>0}\) by Lemma~\ref{B8.11f}. Apply this on the row-space charts of each stratum, in descending rank, to kill the defect. On the reverse source block the left character cancels the Weil multiplier. Summing over \(j\) and over finite sums of pure tensors for \(P\) proves invariance of \(T_Pu\) under \(\U(S)(F_v)\) for every completely split \(v\notin E_{\rm exc}\).

By \eqref{eq:p8c2-triple-range}, \(\dim_E S\ge2\). Strong approximation for the special unitary group at one good split place reduces constancy to density in the norm-one determinant torus \cite{platonov-1969}. A character of a proper quotient by the rational points and good local groups pulls back under \(z\mapsto z/\bar z\) to a nontrivial Hecke character of \(E\), by Hilbert 90. It is trivial at almost all absolute degree-one primes by Galoisness. The quotient of its \(L\)-function by \(\zeta_E\) is regular and nonzero at one, by absolute convergence of the degree-at-least-two factors near one. This contradicts holomorphy of its \(L\)-function there \cite{tate-1967}. The argument allows exclusion of any fixed finite set.
\end{proof}

\subsection*{Proof of \ref{B8.11}}\label{prf:B8.11}\phantomsection
\begin{proof}[Proof of Theorem~\ref{B8.11}]
At a ramified split place the operators in \eqref{eq:p8c2-split-matching} take values in the same normalized parabolic induction. For an upper block triangular \(p\) with diagonal blocks \(A,D_1\) of sizes \(a,a'\), substitute \(gA^{-t}\) on the first side and \(gD_1\) on the second. Both left multipliers are
\[
 \alpha_v(\det A)\beta_v(\det D_1)|\det A|_v^{a'/2}|\det D_1|_v^{-a/2}.
\]
An indicator of a small ball about each standard row block gives at \(h=1\) the positive volume of a principal congruence subgroup on which \(w=1\). Thus both maps are nonzero. Bernstein's irreducibility theorem \cite{bernstein-1984} gives surjectivity onto the common induction. At an unramified split place use the spherical model and identify the integral-matrix vacua.
\leanref{PositivityTheory.\allowbreak{}ThK.\allowbreak{}SplitAsm.\allowbreak{}finite\_\allowbreak{}ramified\_\allowbreak{}of\_\allowbreak{}bernstein}{Development.\allowbreak{}WeilBound.\allowbreak{}CoreStatements.\allowbreak{}CompactAveragingSplitting.\allowbreak{}EigenvectorLifting.\allowbreak{}SplitTop}{14--16}
\leanref{PositivityTheory.\allowbreak{}ThK.\allowbreak{}SplitAsm.\allowbreak{}finite\_\allowbreak{}split\_\allowbreak{}holds}{Development.\allowbreak{}WeilBound.\allowbreak{}CoreStatements.\allowbreak{}CompactAveragingSplitting.\allowbreak{}EigenvectorLifting.\allowbreak{}SplitTop}{18--21}

At an inert place take functionals \(l,l'\) killing the two source difference spaces, invariant under the target maximal compact \(K_v\) and normalized to take value one on the vacua \(s,s'\). They are obtained by averaging the zero evaluations over \(K_v\). For the Iwasawa parabolic \(P_v\) stabilizing a maximal isotropic space, the unaveraged normalized vacuum coefficients are right \(K_v\)-invariant and left \(P_v\)-equivariant. Their characters are \(e(p)=\beta(z)R^{a/2}\) and \(e'(p)=\alpha(z)R^{a'/2}\), where \(R=|N_{E_v/F_v}(\det t)|_v\) for the Levi block \(t\), and \(z\) is the corresponding local determinant idele. The unramified splitting characters give \(\alpha(z)\beta(z)=1\). Write \(J=P_v\cap K_v\), \(n_+(p)=[J:J\cap pJp^{-1}]\) and \(n_-(p)=n_+(p^{-1})\). The parabolic modulus gives \(n_-(p)=R^{D/2}n_+(p)\). Thus
\[
 n_+(p)e(p)=n_-(p)e'(p^{-1})
\]
by \(a/2+a'/2=D/2\). In the finite average on \(K_vhK_v/K_v\), group the Iwasawa representatives into \(J\)-orbits on \(P_v/J\). Inversion exchanges the \(J\)-double cosets of \(p\) and \(p^{-1}\), and the Cartan decomposition gives \(K_vh^{-1}K_v=K_vhK_v\). The displayed identity pairs the contributions, giving \(l(hs)=l'(hs')\). By Lemma~\ref{B8.14} and Corollary~\ref{B8.14a} both coefficient images are the span of the translates of this spherical function. For a nonzero coefficient, translate it to have nonzero value at one and average its test over \(K_v\). The resulting coefficient is a nonzero multiple of the vacuum coefficient. Hence the common image is irreducible.
\leanref{PositivityTheory.\allowbreak{}Sph.\allowbreak{}sphFun\_\allowbreak{}of}{Development.\allowbreak{}WeilBound.\allowbreak{}ThetaLifts.\allowbreak{}SphericalLattices.\allowbreak{}FunAsm}{113--191}
\leanref{PositivityTheory.\allowbreak{}Sph.\allowbreak{}finAt\_\allowbreak{}sph}{Development.\allowbreak{}WeilBound.\allowbreak{}ThetaLifts.\allowbreak{}SphericalLattices.\allowbreak{}FinInert}{120--217}

At each real place use the signature frames of the three spaces to identify the polynomial vectors with the two Fock models of Lemma~\ref{B8.11i}. The source splitting adds \(-c\operatorname{tr}\) on the first side and \(c\operatorname{tr}\) on the second, so the maps to \(T,T'\) kill the actual source Lie actions. Lemma~\ref{B8.11k} supplies surjections from both modules onto the same nonzero simple target module. Add \((m_{\beta,w}/2)\operatorname{tr}\) to its action. The unprimed action then has the prescribed target splitting, and the primed action has total added trace coefficient \(c+m_{\beta,w}/2=m_{\alpha,w}/2\). The intertwining identities hold on every matrix unit. The compact groups are connected, and the compact orbits of polynomial vectors span finite-dimensional spaces, so exponentiation gives target compact equivariance and source compact invariance. Tensor the surjections over the real places. Each local module is a quotient of a countable-dimensional polynomial space and is simple. Lemma~\ref{B1.15ba} therefore gives simplicity of their tensor product. At completely split good places Lemma~\ref{B8.11f} gives the kernel and full-rank surjectivity. The split matrix action preserves the rank of the selected primed factor.
\leanref{PositivityTheory.\allowbreak{}KRealA.\allowbreak{}kReal\_\allowbreak{}of\_\allowbreak{}lie}{Development.\allowbreak{}WeilBound.\allowbreak{}CoreStatements.\allowbreak{}GaussianSeminormBounds.\allowbreak{}RealSeminormFiltration.\allowbreak{}KRealAsm}{292--385}
\leanref{PositivityTheory.\allowbreak{}Sph.\allowbreak{}theoremK\_\allowbreak{}sph\_\allowbreak{}of}{Development.\allowbreak{}WeilBound.\allowbreak{}ThetaLifts.\allowbreak{}SphericalLattices.\allowbreak{}TheoremKSphOf}{10--16}

Lemma~\ref{B8.11e} gives absolute convergence on the unprimed side by \(a+a'+2p<2a\). It gives \eqref{eq:p8c2-full-rank-decay}, primed moderate growth and Schwartz continuity on the allowed spaces.

Induct on \(p+q\). At \(p=q=0\), Lemma~\ref{B8.11l} gives a polynomial witness on the unprimed side with \(I(\Phi,\cdot)\ne0\) and a polynomial witness \(\Phi'\in\mathcal S'_{t_1}\) on the primed side with \(I'(\Phi',\cdot)\ne0\). Surjectivity of the matching maps supplies polynomial partners for both witnesses, and Lemma~\ref{B8.11f} at \(v_1\) allows the primed partner of the unprimed witness to be chosen in \(\mathcal S'_{t_1}\).
\leanref{GlobalTheory.\allowbreak{}ThK.\allowbreak{}Dbl.\allowbreak{}bump}{Development.\allowbreak{}WeilBound.\allowbreak{}ThetaLifts.\allowbreak{}LatticeBasics.\allowbreak{}NearUnits}{113--185}
\leanref{GlobalTheory.\allowbreak{}ThK.\allowbreak{}Dbl.\allowbreak{}nonZeroBottom\_\allowbreak{}of\_\allowbreak{}doubling\_\allowbreak{}integral}{Development.\allowbreak{}WeilBound.\allowbreak{}ThetaLifts.\allowbreak{}LatticeBasics.\allowbreak{}BotAsm}{117--147}
\leanref{GlobalTheory.\allowbreak{}ThK.\allowbreak{}Dbl.\allowbreak{}pair\_\allowbreak{}ne\_\allowbreak{}zero}{Development.\allowbreak{}WeilBound.\allowbreak{}ThetaLifts.\allowbreak{}LatticeBasics.\allowbreak{}BotAsm}{58--116}
\leanref{GlobalTheory.\allowbreak{}ThK.\allowbreak{}Dbl.\allowbreak{}badFactor\_\allowbreak{}of}{Development.\allowbreak{}WeilBound.\allowbreak{}ThetaLifts.\allowbreak{}LatticeBasics.\allowbreak{}BadFactor}{96--192}

By Lemma~\ref{B8.11g}, the two integrals \(I_y,J_y\) of a polynomial match \(y\) are equivariant maps from the common restricted tensor product. The local spaces have at most countable complex dimension, from polynomial monomials and indicators of a countable basis of balls, so their simple quotients have at most countable complex dimension as well. Adjoining the nonzero vacua, Lemma~\ref{B1.15ba} gives simplicity of the restricted tensor product. At \(p=q=0\) the connected compact real target integrates its Lie action on compact-finite vectors.
\leanref{PositivityTheory.\allowbreak{}KReal.\allowbreak{}tensorIrr}{Development.\allowbreak{}WeilBound.\allowbreak{}CoreStatements.\allowbreak{}GaussianSeminormBounds.\allowbreak{}RealSeminormFiltration.\allowbreak{}Tensor}{131--166}
\leanref{GlobalTheory.\allowbreak{}ThK.\allowbreak{}Dbl.\allowbreak{}pairsIrr\_\allowbreak{}of\_\allowbreak{}arch}{Development.\allowbreak{}WeilBound.\allowbreak{}ThetaLifts.\allowbreak{}LatticeBasics.\allowbreak{}PairsIrrOfArch}{10--12}

The primed map \(J\) is a nonzero embedding. Any invariant Petersson form from the two maps is a scalar multiple of the positive form \(\langle J_y,J_z\rangle\), taken linear in the first argument. On the compact rational quotient, a compact type with fixed open finite level lies in a finite direct sum of its matrix-coefficient spaces. Represent the form on one nonzero such isotypic space by an operator relative to the positive form. Subtract an eigenvalue times the positive form. Compact invariance makes distinct types orthogonal, so the resulting radical contains an eigenvector. Simplicity gives the whole space as radical. First apply this to \(\langle I_y,J_z\rangle\) to choose \(c\), then to \(D_y=I_y-cJ_y\). This gives
\begin{equation}\label{eq:p8c2-hermitian-line}
 \langle D_y,J_z\rangle=0,\qquad
 \langle D_y,D_z\rangle=\lambda\langle J_y,J_z\rangle.
\end{equation}
\leanref{GlobalTheory.\allowbreak{}ThK.\allowbreak{}Dbl.\allowbreak{}hermLine\_\allowbreak{}of}{Development.\allowbreak{}WeilBound.\allowbreak{}ThetaLifts.\allowbreak{}LatticeBasics.\allowbreak{}RankPolData}{107--175}

For \(u=I_y\) or \(J_y\) and a polynomial primed test \(P\), the reverse lift \(T_Pu\) of Lemma~\ref{B8.11d} is constant on \(\U(V')(\mathbb A_F)\). Indeed \([\U(W)]\) is compact at \(p=q=0\), so every continuous function on \([\U(W)]\) is integrable against \(\overline{\theta_{\omega'}(x,\cdot;P)}\). Fixing the factors of the image of \(y\) away from each completely split good place \(v\notin T\cup\{v_1\}\) writes \(u=\sum_j U_j(\xi_j)\) via Lemma~\ref{B8.11g}, and Lemma~\ref{B8.11m} gives constancy.
\leanref{GlobalTheory.\allowbreak{}ThK.\allowbreak{}Dbl.\allowbreak{}LD.\allowbreak{}eq\_\allowbreak{}case}{Development.\allowbreak{}WeilBound.\allowbreak{}ThetaLifts.\allowbreak{}LatticeBasics.\allowbreak{}LocDistE}{106--144}
\leanref{GlobalTheory.\allowbreak{}ThK.\allowbreak{}Dbl.\allowbreak{}TC.\allowbreak{}inv\_\allowbreak{}loc}{Development.\allowbreak{}WeilBound.\allowbreak{}ThetaLifts.\allowbreak{}LatticeBasics.\allowbreak{}ThConstInv}{86--153}
\leanref{GlobalTheory.\allowbreak{}ThK.\allowbreak{}Dbl.\allowbreak{}thConst\_\allowbreak{}of'}{Development.\allowbreak{}WeilBound.\allowbreak{}ThetaLifts.\allowbreak{}LatticeBasics.\allowbreak{}ThConstFin}{13--16}
\leanref{GlobalTheory.\allowbreak{}ThK.\allowbreak{}Dbl.\allowbreak{}normOne\_\allowbreak{}mem}{Development.\allowbreak{}WeilBound.\allowbreak{}ThetaLifts.\allowbreak{}LatticeBasics.\allowbreak{}ConstIdele}{134--182}

Choose \(J_y\ne0\) and let \(P\) be its test. With \(v'=\operatorname{vol}[\U(V')]\), Fubini gives \(T_PD_y=0\) and \(T_PJ_y=\|J_y\|^2/v'\). Lemma~\ref{B8.11d} gives \(0=\lambda\|J_y\|^4/v'\). Hence \(\lambda=0\), and \eqref{eq:p8c2-hermitian-line} gives \(D_z=0\) for every \(z\). The unprimed witness gives \(c\ne0\).
\leanref{GlobalTheory.\allowbreak{}ThK.\allowbreak{}Dbl.\allowbreak{}polZero\_\allowbreak{}of}{Development.\allowbreak{}WeilBound.\allowbreak{}ThetaLifts.\allowbreak{}LatticeBasics.\allowbreak{}PolZero}{44--110}

For positive index sum take the adapted planes of Lemma~\ref{B8.11a}. Use \(S,C\) of Lemma~\ref{B8.11b} and put \(c=c_1/(CS)\) for \(q>0\) and \(c=CS c_1\) for \(q=0<p\). The primitive Gauss sums and the Euler product to the right of one are nonzero. Apply the Levi identities to compact translates of a polynomial match with \(\Phi'\in\mathcal S'_{t_1}\). Iwasawa decomposition and the two character laws give \(B=cB'\) for \(q>0\) and \(A=cA'\) for \(q=0<p\). The difference \(f=I(\Phi)-cI'(\Phi')\) is in \(\Sigma\) of Lemma~\ref{B8.11c} and has a line term of the remaining character type.

If \(q=0<p\), then \(B'=0\). The zero evaluation for \(B\) has target dimension \(d_1=a+a'-2\) and smaller target index \(p-1<p\). Since \(|m_{\beta,w}-m_{\alpha,w}|>d_1\) and \(d_1+2p<2a\), Lemma~\ref{lem:p8c2-real-vanishing} gives \(B=0\).

For \(q>0\) induct on \(p\) for differences with the remaining \(\alpha\)-type line term. At \(p=0\), definiteness gives \(A=0\). Apply the full-rank zero-evaluation rule in Lemma~\ref{B8.11b} to the primed test. It preserves the upper integral and makes the lower zero evaluation full rank. Lemma~\ref{lem:p8c2-real-vanishing} now kills \(A'\), since \(q-1<q\) and \(|m_{\alpha,w}-m_{\beta,w}|>d_1\).

At \(p>0\), move the line by Witt transitivity to the plane for removal on the \(V\)-side. At line coordinate one the line term of each compact translate is a lower difference with indices \((p-1,q)\), by the cone and zero evaluations. Its line term \(c_2\) has two lower Levi types. Let \(u\in\mathbb A_F^\times\) equal \(2\) at infinity and \(1\) at finite places. Let \(a_u\) scale the two isotropic lines by \(u,u^{-1}\) and their duals by the inverse conjugates. The two types have values \(1,\lambda\) at \(a_u\), where
\begin{equation}\label{eq:p8c2-two-step-contraction}
 \lambda=\alpha(u)\beta(u^{-1})|u|_E^{(a'-(a-2))/2},\qquad
 |\lambda|=|u|_E^{(a'-(a-2))/2}<1.
\end{equation}
The first type is the product of the two \(\alpha\) types. In the second the lower unprimed source has dimension \(a-2\). Interchanging the unipotent integrations realizes \(c_2\) as the upper-unipotent average of the two-plane constant term restricted to \(\mathrm{SL}_2(\mathbb A_E)\) on the two lines. The same character calculation for variable positive scale bounds this average on the positive Siegel ray, uniformly on compact right factors. The two coefficients are continuous, since they are recovered from \(c_2\) and its translate by \(a_u\).

Apply \eqref{eq:p8c2-smoothing} at each fixed target right translate to obtain a bounded smoothing of the two-plane term. Its smoothed unipotent average is bounded. The smoothed difference \(c_2(a_u g)-c_2(g)\) scales by \(\lambda\) under \(a_u\). Evaluating at \(a_u^{-k}g\) bounds its value at \(g\) by its supremum times \(|\lambda|^k\). Let \(k\to\infty\) and shrink the smoothing supports. The difference is zero, so the second lower type vanishes. The induction kills the lower difference, and Iwasawa decomposition kills the upper line term. Lemma~\ref{B8.11c} gives \(f=0\). Transport the comparison through the change of frames.

For \(q>0\) use the nonzero \(B\) witness of Lemma~\ref{B8.11b} and a primed partner of full rank at \(v_1\). The Levi identity gives \(B'\ne0\). For \(q=0<p\) the lemma gives an \(A'\) witness. Replace its primed \(v_1\)-factor by a full-rank representative. The source is definite, so absolute convergence and coinvariance preserve \(A'\). The Levi identity gives \(A\ne0\). These witnesses give uniqueness of \(c\), and continuity extends the comparison to general matches. Cancel the nonzero line terms in the Levi identities to obtain \eqref{eq:p8c2-beta-step} and \eqref{eq:p8c2-alpha-step}.
\end{proof}

\subsection*{Proof of \ref{B8.11e}}\label{prf:B8.11e}\phantomsection
\begin{proof}
Write \(n=a\) or \(a'\), \(d=a+a'\), \(e=[E:\Q]\), and \(r,l\) for the rational Witt indices of source and target. In rational Witt bases the positive chamber has source weights \(s_i,-s_i,0\), with \(s_0\ge\cdots\ge s_{r-1}\ge0\) and \(n-2r\) zero weights. Use \(t_j\) on the target. Bounded shifts account for the fixed lower bounds in the Siegel set. Put \(\rho_n(s)=\sum_{i<r}(n-2i-1)s_i\), and define \(\rho_d(t)\) likewise. The chamber densities are
\begin{equation}\label{eq:p8c2-theta-densities}
 e^{-e\rho_n(s)}\,ds,\qquad e^{-e\rho_d(t)}\,dt .
\end{equation}
The roots for a pair \(i<j\) contribute \(2s_i\) with multiplicity \(e\), the middle block contributes \(n-2r\) copies of \(s_i\) with that multiplicity, and \(2s_i\) occurs with multiplicity \(e/2\). The rank-zero chamber has mass one. The Siegel-set estimate of Proposition~\ref{B1.15b} uses a majorant uniform on the compact right factors. Its height comparisons give
\begin{equation}\label{eq:p8c2-theta-heights}
 H_{\rm quot}(g)\le C e^{A\sum_i s_i},\qquad
 e^{\sum_j t_j}\le C\operatorname{ht}(h)^B
\end{equation}
for positive \(C,A,B\). The first uses the reduced diagonal and compact bound. For the second the diagonal controls each \(e^{t_j}\) and the reduced height is bounded by a power of the original height.
\leanref{GlobalTheory.\allowbreak{}CompanionTheory.\allowbreak{}SiegelRed.\allowbreak{}siegelCore}{Development.\allowbreak{}AutomorphicSpectrum.\allowbreak{}Proofs.\allowbreak{}GlobalProofs.\allowbreak{}SiegelSetUnitary}{17--28}

Set \(S=\sum_i s_i,\ T_1=\sum_jt_j\), pad the sequences with zeros and put
\[
 P=dS+nT_1-2\sum_{i<r,j<l}\min(s_i,t_j),\qquad
 \Delta=\max_{0\le k\le r}\sum_{i<k}(s_i-t_i).
\]
For fixed finite factors and real Schwartz factor \(f\), the bounds are
\begin{equation}\label{eq:p8c2-theta-square-estimates}
 |\theta(g,h)|^2\le p_0(f)^2e^{eP},\qquad
 |\theta(g,h)|^2\le p_m(f)^2e^{eP-m\Delta}\quad(m\in\N).
\end{equation}
The \(p_m\) are continuous Schwartz seminorms, and the second estimate uses full-rank support at the chosen split place. It applies to either side in compatible models.

In Schr\"odinger coordinates put \(D(u,v)\Phi(x)=\psi(v\cdot(x+u/2))\Phi(x+u)\). Group the product of two rational theta sums, conjugate in the second factor, by the difference \(u\). Poisson summation in the other argument gives the coefficients \(\int D(u,v)\Phi_1(x)\overline{\Phi_2(x)}\,dx\) at rational \((u,v)\), up to unit phases and a measure factor. The rational lattice is its own annihilator. In the chosen polarization \((u,v)\) gives a matrix \(Z\in\operatorname{Mat}_{n,d}(E)\), which Weil covariance replaces by \(g^{-1}Zh^{-t}\). At finite places the coefficient has compact support, since translation must bring the supports together and modulation must be trivial on a common period lattice. At infinity, integrating by parts in \(x\) against the modulation character gives decay in \(v\). Schwartz decay gives the corresponding bound in \(u\). Together they bound the coefficients by any inverse power of \(1+\|g^{-1}Zh^{-t}\|_\infty\), times Schwartz seminorms. Compact right factors give bounded enlargements of the support and equivalent norms.

For the full-rank estimate partition the selected finite support into small cosets \(Y_0+L_0\) for which every element of \(2Y_0+L_0\) has rank \(n\). For each piece pair its test with its reflection \(x\mapsto-x\). Reflection commutes with the Weil action up to a unit scalar and permutes the summation lattice. The absolute product therefore equals the squared absolute theta sum of the piece. A contributing translation coordinate lies in \(\pm(2Y_0+L_0)\). In the barred model this coordinate is the conjugate split component of \(Z\) up to invertible form matrices. All contributing \(Z\) have rank \(n\). Sum the bounds for the absolute theta values of the pieces and then for the pure tensors.
\leanref{PositivityTheory.\allowbreak{}ThKEst.\allowbreak{}sqAll\_\allowbreak{}of'}{Development.\allowbreak{}WeilBound.\allowbreak{}ThetaLifts.\allowbreak{}FrameEstimates.\allowbreak{}SqAllC}{103--112}
\leanref{PositivityTheory.\allowbreak{}ThKEst.\allowbreak{}soCut\_\allowbreak{}of}{Development.\allowbreak{}WeilBound.\allowbreak{}ThetaLifts.\allowbreak{}FrameEstimates.\allowbreak{}HermitianFormDetNeZero}{76--155}
\leanref{PositivityTheory.\allowbreak{}ThKEst.\allowbreak{}sqSo\_\allowbreak{}of'}{Development.\allowbreak{}WeilBound.\allowbreak{}ThetaLifts.\allowbreak{}FrameEstimates.\allowbreak{}ThetaSumNeg}{136--168}

In the rational Witt bases the entries of \(Z\) lie in a fixed fractional ideal \(\mathfrak a\). Counting disjoint balls about its points and summing dyadic shells at scale \(\max(1,e^{b_1})\) gives, for \(M>e\),
\begin{equation}\label{eq:p8c2-theta-lattice-count}
 \sum_{z\in\mathfrak a}(1+e^{-b_1}\|z\|_\infty)^{-M}
 \ll_{\mathfrak a,M}e^{e\max(b_1,0)} .
\end{equation}
Take sufficient Schwartz order to bound a coefficient by the product of these entry weights. The four signed sums for a pair \(s_i,t_j\) contribute \(2\max(s_i,t_j)\) to the exponent sum, and the zero weights contribute \((d-2l)S+(n-2r)T_1\). The total is \(P\).

For a full-rank \(Z\), choose \(k\) attaining \(\Delta\). The first \(k\) source rows expanded by \(\exp s_i\) have a nonzero \(k\)-minor. Bounded denominators and the field norm bound its unscaled determinant below at some embedding. The scaling factor is at least \(\exp(\sum_{i<k}(s_i-t_i))\), since any \(k\) target weights have sum at most \(\sum_{i<k}t_i\). Including \(k=0\) this gives
\begin{equation}\label{eq:p8c2-theta-minor}
 e^\Delta\ll(1+\|g^{-1}Zh^{-t}\|_\infty)^{n+1}.
\end{equation}
Use \(m(n+1)\) additional Schwartz powers for this factor. This proves \eqref{eq:p8c2-theta-square-estimates}.
\leanref{PositivityTheory.\allowbreak{}ThKEst.\allowbreak{}minor\_\allowbreak{}of}{Development.\allowbreak{}WeilBound.\allowbreak{}ThetaLifts.\allowbreak{}FrameEstimates.\allowbreak{}ExistsColsDetNeZero}{77--167}

Since \(S\le T_1+\Delta\) and \(P\le dS+nT_1\), use \(m\ge ed+2AN\) in the square root of the full-rank estimate to obtain \eqref{eq:p8c2-full-rank-decay}. With \(m\ge ed+2\) the same estimate without the height multiplier gives
\begin{equation}\label{eq:p8c2-theta-moderate-density}
 |\theta(g,h)|\le p_m(f)\exp\bigl(-S+(en+m)T_1/2\bigr).
\end{equation}
Integration using \eqref{eq:p8c2-theta-densities} and the target height bound gives moderate growth.
\leanref{PositivityTheory.\allowbreak{}ThKEst.\allowbreak{}soSemiU\_\allowbreak{}of}{Development.\allowbreak{}WeilBound.\allowbreak{}ThetaLifts.\allowbreak{}FrameEstimates.\allowbreak{}SoSemi}{22--103}
\leanref{PositivityTheory.\allowbreak{}ThKEst.\allowbreak{}soDecay\_\allowbreak{}of}{Development.\allowbreak{}WeilBound.\allowbreak{}ThetaLifts.\allowbreak{}FrameEstimates.\allowbreak{}ThickFrames}{156--175}

For \(n=a\), use \(P\le dS+nT_1\) in the first bound. In the successive differences of chamber coordinates the coefficients after multiplication by the source density, and after source integration, squaring and multiplication by the target density, respectively, are
\begin{equation}\label{eq:p8c2-theta-negative-sums}
 ek(d/2-n+k)<0\quad(1\le k\le r),\qquad
 ek(n-d+k)<0\quad(1\le k\le l).
\end{equation}
These are the partial sums of the original exponents. The signs follow from \(r\le p,\ d+2p<2a,\ l\le p+q<a'\). The source calculation works for any compatible pair with \(d+2r<2n\).

For \(n=a'\), put \(Y=p+q\) and \(\epsilon=1/(4(r+1))\). Here \(r,l\le Y,\ 2Y+2\le n,\ 2n\le d\). Then
\begin{equation}\label{eq:p8c2-theta-primed-exponent}
 P/2-(d+r+1)\Delta\le
 \rho_n(s)-\epsilon S+\rho_d(t)/2-T_1/4 .
\end{equation}
Write \(H=\sum_{i,j<r}\min(s_i,s_j)=\sum_{i<r}(2i+1)s_i\) and \(C_1=\sum_{i<r,j<l}\min(s_i,t_j)\). For \(j<r\), if \(k\) of the \(s_i\) exceed \(t_j\), then
\[
 \sum_i(s_i-t_j)_+
 \le \Delta+\sum_{i<k}(t_i-t_j)\le\Delta+\sum_{i<j}t_i.
\]
Thus \(H-C_1\le r\Delta+\sum_{j<r}(r-1-j)t_j\). Put \(b_2=d/2-n+\epsilon\ge0\). From \(S\le\sum_{j<r}t_j+\Delta\) we get
\[
 b_2S+H-C_1\le(b_2+r)\Delta+\sum_{j<r}(b_2+r-1-j)t_j.
\]
The coefficient at \(j<r\) is at most \((d-2j-1-n)/2-1/4\), which is nonnegative for \(r\le j<l\). Using \(\rho_n=nS-H\) gives the claim. Take \(m\ge2e(d+r+1)\) in \eqref{eq:p8c2-theta-square-estimates}. The source integrand has factor \(e^{-e\epsilon S}\). After source integration and squaring the target density leaves \(e^{-eT_1/2}\). Both are integrable.

On a compact target set the source estimates are uniform. Dominated convergence and the seminorm bounds give the stated continuity. A definite source has \(r=0\) and its absolute-convergence estimate applies to all its tests.
\leanref{PositivityTheory.\allowbreak{}ThKEst.\allowbreak{}Sq.\allowbreak{}exponent\_\allowbreak{}bound\_\allowbreak{}orthogonal\_\allowbreak{}case}{Development.\allowbreak{}WeilBound.\allowbreak{}ThetaLifts.\allowbreak{}FrameEstimates.\allowbreak{}InitialPairs.\allowbreak{}SquareIntegrableIntegralHolds}{194--258}
\leanref{PositivityTheory.\allowbreak{}ThKEst.\allowbreak{}Sq.\allowbreak{}square\_\allowbreak{}integrable\_\allowbreak{}of\_\allowbreak{}lattice\_\allowbreak{}torus}{Development.\allowbreak{}WeilBound.\allowbreak{}ThetaLifts.\allowbreak{}FrameEstimates.\allowbreak{}InitialPairs.\allowbreak{}SquareIntegrableOfLatticeTorus}{134--152}
\leanref{PositivityTheory.\allowbreak{}ThKEst.\allowbreak{}Sq.\allowbreak{}square\_\allowbreak{}integrable\_\allowbreak{}orthogonal\_\allowbreak{}case}{Development.\allowbreak{}WeilBound.\allowbreak{}ThetaLifts.\allowbreak{}FrameEstimates.\allowbreak{}InitialPairs.\allowbreak{}SquareIntegrableOfLatticeTorus}{161--185}
\end{proof}

\subsection*{Proof of \ref{B8.11f}}\label{prf:B8.11f}\phantomsection
\begin{proof}
Suppose \(\phi\) is supported where the \(J\)-rows are independent and choose a row \(i\notin J\). Let \(T_b\) add \(\sum_{j\in J}b_jY_j\) to row \(i\), and let \(D_x\) scale that row by \(x\). Choose \(|x|\le1\) with \(\zeta(x)\ne1\). Invert the chosen element if needed to meet the size bound. Smoothness gives a lattice \(L_0\subset k^J\) of periods for the row additions. Choose a larger \(L\) containing the coefficients of every dependent \(i\)-row on the compact support. Boundedness follows from the nonzero \(J\)-minors on compact open charts. Put \(N_1=[L:L_0], N_2=[x^{-1}L:L], z=\zeta(x)|x|^{d/2}\) and
\[
 F=\sum_{b\in L/L_0}\phi\circ T_b,\qquad
 \Psi=\frac{z}{N_1(zN_2-1)}
 \left(\sum_{b\in x^{-1}L/L}F\circ T_b-F\circ D_x\right).
\]
The denominator is nonzero by unitarity and \(\zeta(x)\ne1\). In the coinvariants row additions act trivially and \([zF\circ D_x]=[F]\), so \([\Psi]=[\phi]\). On dependent \(i\)-rows, \(F\) is constant on \(L\) and zero elsewhere in the coefficient variable. Both terms in parentheses give the same constant on \(x^{-1}L\) and vanish outside. They also vanish on dependent \(J\)-rows. Start with \(J=\varnothing\) and adjoin the rows to obtain \(\mathcal S=\mathcal S^\circ+A(\mathcal S)\).
\leanref{PositivityTheory.\allowbreak{}KCo.\allowbreak{}step}{Development.\allowbreak{}WeilBound.\allowbreak{}CoreStatements.\allowbreak{}LocallyConstantCoinvariants.\allowbreak{}PfSurj}{236--316}
\leanref{PositivityTheory.\allowbreak{}KCo.\allowbreak{}surjectivity\_\allowbreak{}statement}{Development.\allowbreak{}WeilBound.\allowbreak{}CoreStatements.\allowbreak{}LocallyConstantCoinvariants.\allowbreak{}PfSurj}{317--331}

The full-rank row-space map has fibres \(\GL_n(k)\) over the Grassmannian. A nonzero minor gives a local section by making its columns the identity. On compact open chart pieces the tests are sums of products of base indicators and fibre functions. Weighted integration with \(\zeta(\det r)|\det r|^{d/2}\,dr\) has kernel the source differences. Indeed the weight converts them to ordinary translation differences, and an integral-zero fibre function is a sum of such differences by expressing it on equal-measure cosets and translating each indicator to a fixed coset. A fibre function of weighted integral one supplies every local section.

The section on the target has the form
\[
 \lambda(\det h)|\det h|^{n/2}
 \int_{\GL_n(k)}\phi(rX_0h)\zeta(\det r)|\det r|^{d/2}\,dr,\qquad X_0=(I_n\ 0).
\]
For the row stabilizer with diagonal blocks \(A,B\), substitution gives the multiplier
\[
 \lambda(\det A\det B)\zeta(\det A)^{-1}
 |\det A|^{-(d-n)/2}|\det B|^{n/2}.
\]
This is normalized induction of two unitary determinant characters. Bernstein's theorem of \ref{B8.43h} gives irreducibility \cite[\S8.2]{bernstein-1984}. The composite to \(\Pi_v\) is onto by the row argument. Since \(\Pi_v\ne0\), the composite and both arrows are isomorphisms.
\leanref{GlobalTheory.\allowbreak{}CompanionTheory.\allowbreak{}SoCoinvNarrow.\allowbreak{}irr\_\allowbreak{}ind}{Development.\allowbreak{}AutomorphicSpectrum.\allowbreak{}Proofs.\allowbreak{}BlockDecomposition.\allowbreak{}SoCoinvOfNarrow}{102--146}
\leanref{GlobalTheory.\allowbreak{}CompanionTheory.\allowbreak{}SoCoinvNarrow.\allowbreak{}irrF}{Development.\allowbreak{}AutomorphicSpectrum.\allowbreak{}Proofs.\allowbreak{}BlockDecomposition.\allowbreak{}SoCoinvOfNarrow}{155--190}

At a good completely split place the source character is \(\nu_v\) or its inverse. Let \(\pi\) generate a positive power of the selected prime, raised to a further power to kill the ramified unit characters. If \(\nu_v(\varpi_v)=1\), the Hecke law on \(\pi\), squared to remove the angular square roots, gives
\[
 \prod_w\left(\frac{\tau_w(\pi)}{\tau_w(\bar\pi)}\right)^{m_{\alpha,w}-m_{\beta,w}}=1.
\]
The primes of the numerator and denominator embeddings are all distinct by Galoisness and complete splitting. At a numerator prime the valuation is positive, a contradiction. The multiplication characters in any barred model separate the matrix points. A functional with the character of evaluation at \(Y\) kills tests vanishing at \(Y\), by partition into pieces where a multiplication character differs from its value at \(Y\). Thus changes of the allowed model preserve support in the corresponding matrix coordinates.
\leanref{PositivityTheory.\allowbreak{}KCo.\allowbreak{}nu\_\allowbreak{}statement}{Development.\allowbreak{}WeilBound.\allowbreak{}CoreStatements.\allowbreak{}LocallyConstantCoinvariants.\allowbreak{}ContinuousFinId}{198--268}
\leanref{PositivityTheory.\allowbreak{}KCo.\allowbreak{}full\_\allowbreak{}of\_\allowbreak{}soLoc}{Development.\allowbreak{}WeilBound.\allowbreak{}CoreStatements.\allowbreak{}LocallyConstantCoinvariants.\allowbreak{}FullOfSoLoc}{66--75}
\end{proof}

\subsection*{Proof of \ref{B8.11g}}\label{prf:B8.11g}\phantomsection
\begin{proof}
Fix every factor except one and write \(u(x)\) for the theta integral with that factor equal to \(x\). Keep the primed \(v_1\)-factor in the full-rank subspace. Lemma~\ref{B8.11e} gives the positive semidefinite form
\[
 B(x,y)=\int_{[\U(W)]}\overline{u(x)(h)}u(y)(h)\,dh
\]
with radical \(\ker u\), since a continuous function of squared norm zero vanishes everywhere. Source changes of variables give \(u(gx)=u(x)\), and right translation gives target invariance of \(B\). At split places with nontrivial source determinant character the required kernel assertion follows from Lemma~\ref{B8.11f}. At \(v_1\) on the primed side the differences use full-rank tests, and both terms of each difference are in the domain of convergence.

At the other finite places write \(\mathcal S\) for the test space, \(A\) for the source differences, \(G\) for the target and \(s\) for the vacuum. Then
\begin{equation}\label{eq:p8c2-theta-spherical-averages}
 \mathcal S=A+\operatorname{span}_\C Gs,\qquad
 \sum_{i=1}^b w_i k_i y-cs\in A,\quad w_i k_i s=s .
\end{equation}
The second assertion holds for each \(y\in\mathcal S\) with suitable \(b\ge1,\ k_i\in G,\ w_i\in\C^\times,\ c\in\C\). Corollary~\ref{B8.14a} gives this at inert places with \(w_i=1\).

At split places with trivial source determinant character, average over the integral source in the barred model. A left integral invariant is a linear combination of indicators that all rows belong to a prescribed lattice. Work on a finite quotient obtained from a common support and period lattice. Left integral orbits are indexed by the submodules generated by the rows. Row operations first select a basis modulo the maximal-ideal multiple of the generated submodule, then clear the redundant generators. Any two minimal generating lists differ by a matrix invertible modulo the maximal ideal. Containment indicators therefore span the orbit indicators by subtraction in increasing submodule order. Each containment indicator is a translate of the vacuum up to its determinant multiplier. The transposed column argument applies to right integral invariants. Average on the right with weights inverse to the right determinant character, using finitely many terms by smoothness. The weights cancel the character on the vacuum and the column argument gives a sum of source translates of it. Passing to coinvariants gives \eqref{eq:p8c2-theta-spherical-averages}.
\leanref{PositivityTheory.\allowbreak{}Sph.\allowbreak{}latSt}{Development.\allowbreak{}WeilBound.\allowbreak{}ThetaLifts.\allowbreak{}SphericalLattices.\allowbreak{}BelowIffRows}{120--240}
\leanref{PositivityTheory.\allowbreak{}Sph.\allowbreak{}SplitPl.\allowbreak{}sphAt\_\allowbreak{}of}{Development.\allowbreak{}WeilBound.\allowbreak{}ThetaLifts.\allowbreak{}SphericalLattices.\allowbreak{}SphSplitOf}{49--174}
\leanref{PositivityTheory.\allowbreak{}Sph.\allowbreak{}SplitPl.\allowbreak{}pair\_\allowbreak{}case2}{Development.\allowbreak{}WeilBound.\allowbreak{}ThetaLifts.\allowbreak{}SphericalSplitting.\allowbreak{}SphSplitGeo}{47--83}

The proper target-stable kernel \(N\) contains \(A\), and the spanning assertion gives \(s\notin N\). For \(y\in N\) the averaging assertion gives \(cs\in N\), hence \(c=0\). Pair with \(s\) to obtain \(0=\sum_i|w_i|^2B(s,y)\). The sum is positive. Apply this to target translates of \(y\), then use spanning and \(u(A)=0\). It follows that \(u(N)=0\).
\leanref{PositivityTheory.\allowbreak{}Sph.\allowbreak{}killFin3}{Development.\allowbreak{}WeilBound.\allowbreak{}ThetaLifts.\allowbreak{}SphericalLattices.\allowbreak{}ThetaKillsCuspFormsSpherical}{19--153}

At infinity compact source differences vanish under \(u\). Schwartz continuity allows differentiation, so source Lie derivatives also vanish under \(u\). Thus \(u,B\) descend to the real source coinvariants \(D\) of Proposition~\ref{B1.15b}. The proper invariant kernel \(N\) of \(D\to\Pi_\infty\) has zero projection onto the single generating compact type \(D[\tau]\) of that proposition \cite[equation~(4.5)]{howe-1989}. The compact isotypic projection is \(B\)-self-adjoint by compact invariance, so \(B(D[\tau],N)=0\). Differentiate right invariance. Lemma~\ref{B8.11e} and Cauchy--Schwarz give \(L^1\) bounds on the differentiated products, uniform on bounded intervals of the right flow. Fubini gives
\begin{equation}\label{eq:p8c2-theta-lie-adjoint}
 B(D_Xx,y)+B(x,D_Xy)=0
\end{equation}
for the real Lie action \(D_X\). The complexified adjoint is \(Z\mapsto JZ^*J\) for the signature matrix \(J\). Move the generating enveloping-algebra words across the pairing. Invariance of \(N\) gives \(B(D,N)=0\).
\leanref{PositivityTheory.\allowbreak{}ThKF.\allowbreak{}L2.\allowbreak{}skew}{Development.\allowbreak{}WeilBound.\allowbreak{}ThetaLifts.\allowbreak{}CartanFactorisation.\allowbreak{}IntegrableOnRight}{106--126}
\leanref{PositivityTheory.\allowbreak{}ThKF.\allowbreak{}maxInfP}{Development.\allowbreak{}WeilBound.\allowbreak{}ThetaLifts.\allowbreak{}CartanFactorisation.\allowbreak{}RinfZero}{148--164}
\leanref{PositivityTheory.\allowbreak{}ThKF.\allowbreak{}killInf'}{Development.\allowbreak{}WeilBound.\allowbreak{}ThetaLifts.\allowbreak{}CartanFactorisation.\allowbreak{}KillInf}{18--151}

Choose a finite set containing \(T,v_1\) and all exceptional factors. Nonzero vacuum images make the transitions of quotient tensors injective. A primed tensor sum in \(\mathcal S'_{t_1}\) can be written with each \(v_1\)-slot full rank by applying a linear projection onto that subspace in this slot. The projection fixes the test, since it fixes the generators of \(\mathcal S'_{t_1}\). Equality of test functions detects equality of tensors by separate point evaluations on the finite-dimensional spans of their factors. Choose a section \(s_j\) of each slot map \(p_j\). At the selected slot use the full-rank surjection. Replace \(x_j\) successively by \(s_jp_jx_j\). Every difference has one factor in a kernel and belongs to the allowed space, so its integral vanishes. For a zero quotient tensor the final sum is zero, which proves descent.
\leanref{PositivityTheory.\allowbreak{}ThKF.\allowbreak{}so\_\allowbreak{}proj}{Development.\allowbreak{}WeilBound.\allowbreak{}ThetaLifts.\allowbreak{}CartanFactorisation.\allowbreak{}Factor}{100--116}
\leanref{PositivityTheory.\allowbreak{}ThKF.\allowbreak{}vanish}{Development.\allowbreak{}WeilBound.\allowbreak{}ThetaLifts.\allowbreak{}CartanFactorisation.\allowbreak{}Factor}{36--61}
\end{proof}

\subsection*{Proof of \ref{B8.11d}}\label{prf:B8.11d}\phantomsection
\begin{proof}
Put \(d=a+a'\), \(z_0=(a'-d)/2\). Let \(e:G\times G\to\U(W\oplus W^-)\) be the doubling embedding, with \(W^-\) carrying the negative form. Use the flat section \(f_z\) of \(P\boxtimes\bar P\) at \(z_0\). It is \(K\)-finite and \(\alpha|_{\mathbb A_F^\times}=\eta^{a'}\). The pairing
\[
 K_z(u)=\int_{[G]^2}E(z,e(h_1,h_2);f_z)
 \alpha(\det h_2)^{-1}u(h_2)\overline{u(h_1)}\,dh_1\,dh_2
\]
continues meromorphically by Theorem~\ref{B8.12} and compactness. In the convergence half-plane the rational maximal isotropic spaces are graphs of isometries, by anisotropy of \(W\). Unfolding and putting \(h_2=h_1g\) gives
\[
 K_z(u)=\int_{G(\mathbb A_F)}
 f_z(e(1,g))\alpha(\det g)^{-1}\langle R(g)u,u\rangle\,dg .
\]
Right translates preserve polynomial matches by compactness at infinity and rank stability at \(v_1\). Thus the assumed proportionality compares the integrands at \(gy,y\). Continue to \(z_0\).
\leanref{GlobalTheory.\allowbreak{}ThK.\allowbreak{}Dbl.\allowbreak{}kpair\_\allowbreak{}cont}{Development.\allowbreak{}WeilBound.\allowbreak{}ThetaLifts.\allowbreak{}LatticeBasics.\allowbreak{}ContinuousFlat}{49--96}

At \(z_0\) the doubled product and the second case of \eqref{eq:p8c2-definite-sw} give
\[
 E(z_0,e(h_1,h_2);f_{z_0})\alpha(\det h_2)^{-1}
 =\frac{2}{\operatorname{vol}[H']}
 \int_{[H']}\theta_{\omega'}(x,h_1;P)\overline{\theta_{\omega'}(x,h_2;P)}\,dx .
\]
Fubini identifies \(K_{z_0}(u)\) with \(2\|T_Pu\|^2/\operatorname{vol}[H']\).
\end{proof}

\subsection*{Proof of \ref{B8.11a}}\label{prf:B8.11a}\phantomsection
\begin{proof}
A space losing a plane has dimension \(n\ge3\) and mixed real signs. For isotropy multiply a skew form by a trace-zero scalar. The rational quadratic form \(\sum_i d_iN_{E/F}(z_i)\) has dimension \(2n\ge6\), hence is isotropic at finite places. It is isotropic at infinity by the signatures. Hasse--Minkowski gives a rational isotropic line \(Ex\) \cite[Ch.~6]{scharlau-1985}.

Write \(h\) for the original form, linear in the first argument, and \(\Lambda\) for its self-dual lattice. Duality uses \(\mathcal O_E\) on a source and the absolute codifferent on the skew target. Scale \(x\) so that \(J=h(\Lambda,x)\) has valuation zero at the prescribed split primes and their conjugates. From \(JJ^{-1}=\mathcal O_E\), take \(y_0\in J^{-1}\Lambda\) with \(h(x,y_0)=1\). Pairing \(h(z,y_0)x\) against \(w\in\Lambda\) gives \(h(z,h(w,x)y_0)\) up to the Hermitian sign. Duality gives \(h(z,y_0)x\in\Lambda\). Choose \(\theta\in\mathcal O_E\) of trace one by unramifiedness and set \(y=y_0-\theta h(y_0,y_0)x\). The trace equation gives \(h(y,y)=0,\ h(x,y)=1\) for either sign. The projection just established gives \(Jy\subseteq\Lambda\), and duality gives \(h(z,y)x\in\Lambda\). Hence the plane projections preserve \(\Lambda\) and the complement is self-dual. The line ideals are units at the prescribed split places, so the plane lines and complement give the required frames. Choose Cartan involutions preserving the plane and complement and exchanging the two lines.

Each removal subtracts one from each sign. The left sides of the size inequalities in \eqref{eq:p8c2-triple-range} fall by two and the right sides by at most two. The character restrictions persist by \(\eta^2=1\).
\end{proof}

\subsection*{Proof of \ref{B8.11b}}\label{prf:B8.11b}\phantomsection
\begin{proof}
Let \(e_0\) set the removed target variable to zero in the line polarization. On the side losing a source plane \(H=(x,y)\), with target plane \(H'=(x',y')\), let \(e\) be evaluation in the mixed polarization for the two planes. Keep the tensor coordinates of the two complements and evaluate the other coordinates at zero after the partial Fourier transform. With \(K\) the chosen source compact group and \(\rho\) the Heisenberg translation, put
\[
 \mathcal T\phi=\int_K\int_{\mathbb A_E^\times}
 e\bigl(\rho(u(t))\omega(k,1)\phi\bigr)\nu(t)|t|_E^z\,d^\times t\,dk,\qquad
 u_\alpha(t)=txy'^{\,t},\quad u_\beta(t)=\bar t\,yx'^{\,t}.
\]
Use the Weil action of the side losing a plane.

Central character orthogonality in the line unipotent retains the isotropic source vectors. The zero orbit gives \(e_0\); unfold the nonzero orbit at its vector stabilizer by Witt transitivity to obtain \(\mathcal T\). On the primed side let \(\Sigma_A^{\rm all},\Sigma_B^{\rm all}\) denote the partial Fourier sums before deletion of the zero vector and \(\Theta_A(0),\Theta_B(0)\) their zero terms. Fourier inversion gives \(\Sigma_A^{\rm all}=\kappa\Sigma_B^{\rm all}\) for the ratio \(\kappa=C_A^{-1}C_B\) of Fourier constants. Reindex the second nonzero sum by \(t\mapsto\bar t^{-1}\). Full-rank support gives \(\Theta_B(0)=0\); therefore \(\Sigma_A^\times=\kappa\Sigma_B^\times-\Theta_A(0)\), with the subtraction in the \(A'\) type. A definite source contributes the zero orbit. Evaluating at one and at a norm different from one separates the characters. With the lower integrals evaluated at \(h_1\), the formulas are
\begin{equation}\label{eq:p8c2-levi-evaluations}
 \begin{array}{c|cc}
 &\text{unprimed}&\text{primed}\\ \hline
 \beta&B=C_0I_1(e_0\Phi)&B'=C_1I'_1(\mathcal T\Phi')\\
 \alpha&A=C_1I_1(\mathcal T\Phi)&A'=C_0I'_1(e_0\Phi'),\quad B'=0 .
 \end{array}
\end{equation}
The measure and evaluation constants are nonzero, with \(\kappa\) included in \(C_1\) for \(\beta\). The zero-orbit formula applies also in the presence of the other type and when removal is from the opposite source. On the side keeping its source it is the lower integral of \(e_0\) times the measure scalar.

For zero evaluation at \(v_1\), the barred fibre at \(Y_1\) consists of \(Y=Y_1\kappa+y\rho\), with \(y\) a column and \(\rho\) a fixed row complementary to the lower coordinates. Thus full rank upstairs implies rank at least \(n-1\) for \(Y_1\), where \(n\) is the common source rank for this evaluation. Cover the support of \(e_0\phi\) by compact open charts on which the rows other than some \(i\) are independent. Apply the row average of Lemma~\ref{B8.11f} after partitioning by these charts. The resulting lower test is full rank modulo a sum of lower source differences of chart tests.

Each such chart test \(\Psi\) has a full-rank upper preimage. In fibre coordinates multiply \(\Psi(Y_1)\) by a compactly supported fibre test of evaluation one at zero, translated by \(Ne_i\). The evaluation is \(\Psi\) by the translation law of zero evaluation. On the compact chart support the coefficients expressing a dependent row \(i\) in terms of the other rows are bounded by the nonvanishing minors. The fibre coordinates before translation are also bounded. An upper row relation would have nonzero coefficient on row \(i\), since the other lower rows are independent. Normalize it to one. The relation on the fibre column would express \(N\) as a bounded combination of the original fibre coordinates. Taking \(|N|_v\) larger gives full rank. Since \(e_0\) intertwines the source action, a lower difference then lifts to a difference of two full-rank upper tests. Subtract the sum of the lifted differences. This gives a full-rank upper \(\phi'\) with \(e_0\phi'\) full rank and with \(\phi-\phi'\) in the restricted source differences. The upper integral is preserved by absolute convergence and source invariance.
\leanref{PositivityTheory.\allowbreak{}ThKEnd.\allowbreak{}Comp.\allowbreak{}SoLower}{Development.\allowbreak{}WeilBound.\allowbreak{}ThetaLifts.\allowbreak{}EigenvalueEndgame.\allowbreak{}EigenKilling.\allowbreak{}ConeIntegralFactorisation}{29--35}
\leanref{PositivityTheory.\allowbreak{}ThKEnd.\allowbreak{}Comp.\allowbreak{}SoL.\allowbreak{}corank\_\allowbreak{}one}{Development.\allowbreak{}WeilBound.\allowbreak{}ThetaLifts.\allowbreak{}EigenvalueEndgame.\allowbreak{}EigenKilling.\allowbreak{}EigenKillSet}{375--403}
\leanref{PositivityTheory.\allowbreak{}ThKEnd.\allowbreak{}Comp.\allowbreak{}SoL.\allowbreak{}lift}{Development.\allowbreak{}WeilBound.\allowbreak{}ThetaLifts.\allowbreak{}EigenvalueEndgame.\allowbreak{}EigenKilling.\allowbreak{}EigenKillSet}{426--514}

On pure tensors \(\mathcal T\) factors locally. The vacuum valuation shells give the unramified Euler factors. Conjugation inverts \(\nu\), and the codifferent has value one. The initial shells for the line ideals give the powers \((1,1)\) for \(u_\alpha\) and \((-1,-1)\) for \(u_\beta\). At every \(v\in T\), the character \(\nu_v\) is ramified, so its conductor exponent satisfies \(c_v\ge1\). Taking additive volume one on \(\mathcal O_v\) and splitting \(\mathcal O_v^\times\) into cosets of \(1+\varpi_v^{c_v}\mathcal O_v\) gives
\[
 \int_{\mathcal O_v^\times}\nu_v(u)^{-1}\psi_v(u/\varpi_v^{c_v})\,du
 =(Nv)^{-c_v}\sum_{u\bmod\varpi_v^{c_v}}^\times
 \nu_v(u)^{-1}\psi_v(u/\varpi_v^{c_v}).
\]
Multiplication by \(\nu_v(\varpi_v)^{c_v}(Nv)^{c_vz}\) gives \(\epsilon_v(\psi,\nu_v,1-z)\). For the standard matched tests of \eqref{eq:p8c2-split-matching}, zero evaluation on \(\varpi_v^{c_v}\mathcal O_v\) in \(a\) coordinates for \(\beta\), or \(a'\) for \(\alpha\), gives the corresponding power of \((Nv)^{-c_v}\) times the lower standard vector. The cone calculation has the same volume and the remaining Gauss factor with \(s=1\). The volume cancels by nonvanishing of the lower standard image.
\leanref{PositivityTheory.\allowbreak{}ThKG.\allowbreak{}Std.\allowbreak{}stdB}{Development.\allowbreak{}WeilBound.\allowbreak{}ThetaLifts.\allowbreak{}LocalTheory.\allowbreak{}StandardValuations.\allowbreak{}ExistsVol}{162--217}
\leanref{PositivityTheory.\allowbreak{}ThKG.\allowbreak{}Std.\allowbreak{}stdA\_\allowbreak{}gen}{Development.\allowbreak{}WeilBound.\allowbreak{}ThetaLifts.\allowbreak{}LocalTheory.\allowbreak{}StandardValuations.\allowbreak{}ExistsVol}{219--283}

To pass from these tests to arbitrary matches, take the two evaluations as maps from the upper local quotient to the lower one for the local Levi, with the prescribed Levi character. Their descent follows by fixing the other evaluated factors with nonzero lower images. The lower integral is injective on the lower matching quotient by the inductive comparison, the nonzero witness and irreducibility. The upper kernel kills the line coefficient by descent of the upper integral, so its evaluated image in the lower quotient is zero. At \(v_1\) take a full-rank upper representative first. The allowed representatives give the same lower image by the kernel relation and compatibility with the source action. The cone evaluation there has lower full-rank support because the removed rows are accounted for by the plane blocks. At places in \(T\), multiplicity one for restriction makes the Levi maps proportional, with the scalar determined by the computed standard images. At places outside \(T\) the unramified restriction uniqueness gives proportionality, and the common vacuum fixes its scalar. This identifies the entire local maps, including on non-vacuum inputs outside \(T\). The real comparison is the scalar fixed by the real maps.

Tensor these identities for arbitrary pure tensors and extend linearly. The evaluated lower matches are
\[
 (C_2 S e_0\Phi,\mathcal T\Phi')\quad(\beta),\qquad
 (\mathcal T\Phi,C_2 S e_0\Phi')\quad(\alpha)
\]
for \(C_2\ne0\). The first primed test is allowed by the rank calculation. In the second pair one may take a full-rank preimage of the lower primed image at \(v_1\) by Lemma~\ref{B8.11f}; the lower source is definite. Lower comparison and \eqref{eq:p8c2-levi-evaluations} give the claimed identities with \(C=C_2C_1/C_0\). Zero evaluation surjects on lower polynomial tests by adjoining tests of evaluation one in the removed coordinates. Lift a nonzero lower unprimed test for \(\beta\), or a nonzero lower primed test for \(\alpha\) with a polynomial partner. Evaluating at a Levi lift of a nonzero value gives the witnesses.
\end{proof}

\subsection*{Proof of \ref{B8.11c}}\label{prf:B8.11c}\phantomsection
\begin{proof}
Witt transitivity gives the line case. Induct on \(j=\dim_E U\). Use a dual frame to embed the Levi \(\GL_j\) by \(m\). Each proper standard block average of \(g\mapsto C_Uf(m(g)h)\) factors through a constant term for a smaller initial subspace. If the restriction is nonzero, a translate has nonzero Whittaker coefficient. Expand down the last columns of the upper unipotent. The zero frequencies are the block averages, while a rational change of basis moves a nonzero frequency to the last coordinate at each step. Fourier uniqueness on the compact additive quotients gives the assertion. Denote the nondegenerate Whittaker average by \(W_f(h)\).

Choose a split finite place outside the exceptional places of the tests, frame and \(h\). In each local matrix model the distribution defining the coefficient has source determinant covariance and a character on the right shears. For columns \(Y_i\), the shear \(Y\mapsto Y+Y_1\rho\), \(\rho_1=0\), has character \(\xi(\rho_2)\) with \(\xi\ne1\). If \(Y_1,Y_2\) are independent, choose \(\lambda Y_1=0,\lambda Y_2=1\) and \(\xi(t)\ne1\). The shear \(\rho=t\lambda Y\) followed by the determinant-one left operation \(1-tY_1\lambda\) fixes \(Y\) with nontrivial character. For \(j\ge3\) and \(Y_2=sY_1\ne0\), or \(Y_2=0,\ Y_1\ne0\), add \(tY_2-tsY_1\) to column three with nontrivial character on the \(Y_2\) shear. For \(Y_1=0\) the first shear already fixes \(Y\) with nontrivial character.

The stabilizers kill the distribution on their strata. On minor charts the stabilizers vary continuously and their characters are locally constant. For commuting square-zero parts \(A,B\), the lattice \(L_0+AL_0+BL_0+ABL_0\) is preserved by both operations. On a sufficiently small coset about a point fixed by the operations, the indicator is fixed with the same nontrivial character. Partition compact supports into such cosets. For strata closed in the remaining support, read the distribution on restrictions by extending locally constant tests to the ambient open set. This first reduces to the dependent locus and then proves vanishing for \(j\ge3\).

For \(j=2\) write \(Q=Nv\) and let \(d\) be the conductor exponent of \(\xi\), trivial on \(\varpi^d\mathcal O_v\) and nontrivial on \(\varpi^{d-1}\mathcal O_v\). Put \(r=\nu_v(\varpi)Q^{(a-a')/2}\). Let \(D(P,M)\) be the unprimed distribution on the box with first two columns in \(\varpi^P\mathcal O_v,\varpi^M\mathcal O_v\) and all other entries integral. Set \(\Delta(P,M)=D(P,M)-D(P+1,M)\). Below \(M-P=d\) the box has a shear of nontrivial character. On the shell for the first column, partition by its first unit pivot after dividing by \(\varpi^P\). Scaling that row by \(\varpi\) gives the determinant multiplier times \(Q^{n-1}/Q^{a+a'-1}\) for source rank \(n\). On dependent pairs the entry of column two in this row determines column two. The trivial-character shears exchange its \(Q\) residue classes at \(M-P\ge d\). This gives
\begin{equation}\label{eq:p8c2-whittaker-shells}
 \begin{aligned}
 D(P,M)&=0 &&(M-P<d),\\
 \Delta(P,M)&=r\Delta(P+1,M)&&(M-P>d),\\
 \Delta(P,M)&=Q\Delta(P,M+1)&&(M-P\ge d).
 \end{aligned}
\end{equation}
On the boundary \(\Delta(M-d,M)=D(M-d,M)\). Applying the third line of \eqref{eq:p8c2-whittaker-shells} and then the second gives \(D(M-d,M)=(Qr)\,D(M+1-d,M+1)\). Thus \(D(M-d,M)=(Qr)^{-M}D(-d,0)\) for \(M\ge0\). Summing shells with \(P=-i,M=i\) gives vacuum values proportional to
\begin{equation}\label{eq:p8c2-whittaker-ray}
 Y_i(r)=(Qr)^{-i}\sum_{k=0}^{2i-d}r^k\qquad(i\ge0),
\end{equation}
with empty sum zero. These boxes are translates by \(b_i=\operatorname{diag}(\varpi^i,\varpi^{-i})\). The primed factor replaces \(r\) by \(r^{-1}\), and reversing the geometric sum gives \(Y_i(r^{-1})=r^dY_i(r)\). Thus \(W_f(m(b_i)h')=c(h')Y_i(r)\) on an identity neighborhood of the right factor at \(h\), by smoothness there. Since \(|r|>Q\), \(|Y_i(r)|\to\infty\).

Let \(\phi\) be continuous automorphic on \(\mathrm{SL}_2(\mathbb A_E)\) with upper-unipotent average \(\phi_U\). Write \(a_k=\operatorname{diag}(2^{k/2},2^{-k/2})\) at infinity and one elsewhere. Suppose \(\phi(a_k z)=O_\Omega(2^{kM})\) for compact right factors \(z\in\Omega\), \(k\ge0\), and integer \(M\ge1\). For an integral-one nonnegative smooth compact bump \(\eta\) write \(\phi*\eta(g)=\int\eta(v)\phi(gv)\,dv\). Then
\begin{equation}\label{eq:p8c2-smoothing}
 |(\phi*\eta)(a_k z)-(\phi_U*\eta)(a_k z)|
 \le C_{\eta,\Omega}2^{k(M-1)} .
\end{equation}
Use a bounded fundamental set of the additive quotient with finite coordinates in a sufficiently small open subgroup, by weak approximation. In the second integral move the unipotent \(n\) across \(a_kz\) and change \(v\). The difference of bumps is \(\eta(v)-\eta(z^{-1}a_k^{-1}n^{-1}a_kzv)\). The finite periods are unchanged and the real displacement is \(O_\Omega(2^{-k})\). The mean-value estimate on the enlarged compact support proves the bound. Right convolution commutes with the average. If the average is bounded on the ray, iteration lowers the growth exponent to zero with arbitrarily small total support, and reduction gives a bounded convolution.
\leanref{PositivityTheory.\allowbreak{}ThKEnd.\allowbreak{}Two.\allowbreak{}Dat.\allowbreak{}conv}{Development.\allowbreak{}WeilBound.\allowbreak{}ThetaLifts.\allowbreak{}EigenvalueEndgame.\allowbreak{}SignatureSplitting.\allowbreak{}ConvCont}{87--111}

Apply this to \(\phi(g)=C_Uf(m(g)h)\), of moderate growth by Lemma~\ref{B8.11e} and compactness of the unipotent quotient. Its upper-unipotent average is zero by the line case. For \(2i_0\ge d\), small convolutions preserve
\[
 Y_{i_0}(r)W_f(m(b_i)h')=Y_i(r)W_f(m(b_{i_0})h')
\]
at \(h\), since it holds on the neighborhood used in the convolution. The smoothed coefficient is bounded by boundedness of the smoothed \(\phi\) and compactness of the unipotent quotient. Thus the smoothed coefficient at \(m(b_{i_0})h\) vanishes. Shrinking the supports and using the relation at \(i=0\) gives \(W_f(h)=0\).

Suppose \(p+q>0\) and pair the resulting cusp form with the defining theta integrals. The reverse lift against the conjugate theta kernel is continuous and rationally invariant. At each completely split \(v\notin T\cup\{v_1\}\), varying the \(v\)-factors in the defining finite sum and using Lemma~\ref{B8.11g} gives an equivariant map \(T_v:\Pi_v^b\to Q_v\) whose image contains \(f\), where \(Q_v\) is the space of continuous automorphic functions on the target. Let \(R_v\subset Q_v\) be the stable space of rapidly decreasing functions. By semisimplicity of \(\Pi_v^b\) take an equivariant projection onto \(T_v^{-1}R_v\). Compose with \(T_v\) and restrict to the summands. This writes \(f=\sum_j U_j(\xi_j)\) with \(U_j(\Pi_v)\subseteq R_v\). Since every element of \(R_v\) is integrable against the moderate-growth conjugate Schwartz kernel, Lemma~\ref{B8.11m} applies to these \(U_j\) and shows that the reverse lift is constant.
\leanref{PositivityTheory.\allowbreak{}ThKEnd.\allowbreak{}Orth.\allowbreak{}locDecomp\_\allowbreak{}of}{Development.\allowbreak{}WeilBound.\allowbreak{}ThetaLifts.\allowbreak{}EigenvalueEndgame.\allowbreak{}OrthogonalDecomposition.\allowbreak{}LocDecompOf}{151--185}
\leanref{PositivityTheory.\allowbreak{}ThKEnd.\allowbreak{}Orth.\allowbreak{}sigThConst\_\allowbreak{}of}{Development.\allowbreak{}WeilBound.\allowbreak{}ThetaLifts.\allowbreak{}EigenvalueEndgame.\allowbreak{}OrthogonalDecomposition.\allowbreak{}SigThConstOf}{22--41}

For an indefinite source take a full-rank representative at \(v_1\) using Lemmas~\ref{B8.11f} and \ref{B8.11g}. Cuspidal decay \cite{harish-chandra-1968} and the full-rank estimate in Lemma~\ref{B8.11e} bound the reverse lift times every power of source quotient height. The rational hyperbolic plane gives a positive-rank rational Siegel ray \(a(s)\). Reduction gives \(e^{\sum_i s_i}\le C\operatorname{ht}(\gamma a(s))^k\) for every rational \(\gamma\). Taking the infimum shows that quotient height is unbounded along the ray, so the constant is zero. Fubini gives orthogonality to the indefinite-source integral.
\leanref{PositivityTheory.\allowbreak{}ThKEnd.\allowbreak{}Orth.\allowbreak{}qht\_\allowbreak{}unbdd}{Development.\allowbreak{}WeilBound.\allowbreak{}ThetaLifts.\allowbreak{}EigenvalueEndgame.\allowbreak{}OrthogonalDecomposition.\allowbreak{}QhtUnbdd}{25--49}

For a definite source \(H\), suppose \(\langle I_H(\Phi_2),f\rangle\ne0\) for a polynomial test. Fubini gives \(T_{\Phi_2}f(1)\ne0\). For any polynomial \(\Phi_1\) pair the doubled Eisenstein kernel of \(\Phi_1\boxtimes\overline{\Phi_2}\) in its second variable with \(f\), using the determinant twist of doubling. In the convergence half-plane, a proper constant term in the first variable unfolds over the rational Siegel cosets. Each contributing orbit has a second-variable stabilizer containing a proper unipotent radical, whose average on \(f\) vanishes. Rapid decay and Theorem~\ref{B8.12} continue the pairing and constant term to the Siegel--Weil point. There the paired value is
\[
 \frac{1+\mathbf 1_{\dim W\ge\dim H}}{\operatorname{vol}[H]}
 T_{\Phi_2}f(1)I_H(\Phi_1).
\]
Thus every \(I_H(\Phi_1)\) would be cuspidal. Start with the double-definite witness of Lemma~\ref{B8.11l} on \(H\) and increase the other signature index while keeping \(H\). Lift the nonzero test at each step by zero evaluation. Equation \eqref{eq:p8c2-levi-evaluations} then gives a nonzero line coefficient at a Levi lift at the prescribed positive index sum. This contradicts cuspidality of every \(I_H(\Phi_1)\). Both defining integrals are orthogonal to \(f\); hence \(\|f\|^2=0\), and continuity gives \(f=0\).
\end{proof}

\subsection*{Proof of \ref{B8.14}}\label{prf:B8.14}\phantomsection
\begin{proof}
For \(\Lambda=\operatorname{Mat}_{n'\times n}(\mathcal O)\), write \(V_{\bar x}=\C\rho(x)s\), with \(\bar x=x+\Lambda\). An integral translation by \(u\) acts on this line by \(\psi(H(u,x)+\sigma H(u,x))\). The conductor and unimodularity distinguish the characters. Thus \(M\) is the direct sum of these lines, with \(\C s\) the trivial-character line.

The group \(K_T\) permutes the lines. Put \(S=\operatorname{span}_\C\{\omega(g,1)s:g\in\U(J_S)\}\subseteq M^{K_T}\). On an orbit supporting an invariant, the stabilizer acts trivially on the corresponding line. For \(u=hx-x\in\Lambda\) the condition is \(\psi(-(H(x,u)+\sigma H(x,u))/2)=1\). Use the Cayley transform, absorbing \(J_T\) into the skew test matrix. For \(D_x=xJ_S^{-1}x^*\) this gives the phase \(\psi((z+\sigma z)/2)=1\) with \(z=\delta\operatorname{tr}(CD_x)\) whenever \(C^*=-C,\ C\in\varpi\operatorname{Mat}_{n'}(\mathcal O)\) and \(Cx\) is integral.

Use Smith coordinates with the forms transformed correspondingly. Subtract integral entries so the active diagonal entries of \(x\) are \(\varpi^{-a_i}\), \(a_i>0\). This leaves the phase test unchanged because \(C,Cx\) are integral and \(C^*=-C\). Put \(A_0=J_S^{-1}\) in these coordinates. Vary the skew entries of \(C\) in \(\varpi^{\max(a_i,a_j)}\mathcal O\) in the active block. The conductor and trace pairing give
\begin{equation}\label{eq:p8c2-smith-gram}
 (A_0)_{ij}\in\varpi^{\min(a_i,a_j)}\mathcal O.
\end{equation}
Here one varies a diagonal anti-invariant entry or an off-diagonal pair with its negative conjugate. Choose an integral \(U_0\) with \((A_0)_{ij}=\varpi^{a_i}(U_0)_{ij}+\varpi^{a_j}\sigma(U_0)_{ji}\) on the active block, assigning off-diagonal entries to the smaller exponent and dividing the diagonal by two. Set other entries to zero. For \(l=U_0A_0^{-1}\), the linear terms in \((x-l)A_0(x-l)^*\) cancel \(xA_0x^*\) and the remaining \(lA_0l^*\) is integral. Replace \(x\) by \(x-l\).
\leanref{GlobalTheory.\allowbreak{}Hk3.\allowbreak{}HoweDuality.\allowbreak{}exists\_\allowbreak{}int\_\allowbreak{}gram\_\allowbreak{}of\_\allowbreak{}phases}{Development.\allowbreak{}AbelianVarietyTheory.\allowbreak{}Proofs.\allowbreak{}SatakeParameters.\allowbreak{}ExistsIntegralGramOfPhaseConditions}{90--128}

Put \(U=\mathcal O^n,\ V=\mathcal O^{n'}\) and \(x^\dagger=J_S^{-1}x^*J_T\). Let \(N=U\cap x^{-1}V,\ A=N+x^\dagger V\). Integrality of \(xx^\dagger\) gives \(A\cap U=N\) and \(A\subseteq x^{-1}V\). From \(N^\vee=U+x^\dagger V\) we obtain \(A^\vee=(U+x^\dagger V)\cap x^{-1}V=A\). Take representatives \(Y\) modulo \(\Lambda\) of the maps integral on \(A\), including \(x\). This is a finite set. Choose columns \(M_i\in U\) generating the cyclic summands of \(U/N\) of orders \(\varpi^{e_i}\), \(e_i>0\), put \(M_0=(M_i)\), and set \(\Phi(y)=(yM_0)^*J_T(yM_0)\).

Let \(B\) be the skew Hermitian matrices \(G\) with entries in \(\varpi^{\max(e_i,e_j)}\mathcal O\) modulo those with entries in \(\varpi^{e_i+e_j}\mathcal O\). The character \(\chi_G(y)=\psi(-\delta\operatorname{tr}(G\Phi(y)))\) descends to \(B\) by the integrality of \(\varpi^{e_i}yM_i\). With \(E_1=M_0GM_0^*J_S,\ k=1+E_1\), we have \(E_1U\subset A,\ E_1A\subset U\) by the divisibilities, \(\varpi^{e_i}M_i\in A\) and self-duality. The matrix \(E_1\) is skew adjoint and divisible by \(\varpi\). Hence \(k^\dagger k=1-E_1^2\) preserves \(A\) with unit determinant, and \(kA\) is self-dual. The maps \(yk^{-1}\) represent the maps integral on \(kA\). Their shifts from \(y\) are integral. In the translation phase expand \(k^{-1}-1=-E_1+E_1^2k^{-1}\). The first term gives \(-\delta\operatorname{tr}(G\Phi(y))\); the second pairs integrally since \(yE_1,k^{-1}\) are integral and \(E_1^\dagger=-E_1\). Thus
\[
 w_G=\sum_{y\in Y}\chi_G(y)\rho(y)s
\]
is the vacuum sum fixed by translations integral on \(kA\). Those translations permute the summands with phase one by integrality of their pairings. Every unimodular Hermitian lattice here has an orthonormal frame, by splitting unit-norm vectors detected over the residue field and normalizing by the norm on units. The residue norm and successive trace corrections give its surjectivity. Consequently \(kA=gU\) for a source isometry \(g\), and covariance gives \(\omega(g^{-1},1)w_G\in\C s\). In particular \(w_G\in S\).

Average \(\chi_G(x)^{-1}w_G\) over \(B\). Character orthogonality retains the terms with the same Gram phase as \(x\). Testing diagonal and off-diagonal skew entries gives
\begin{equation}\label{eq:p8c2-gram-congruence}
 \varpi^{\max(e_i,e_j)}(\Phi(y)-\Phi(x))_{ij}\in\mathcal O.
\end{equation}
For \(m,m'\in U\) and \(c\in\mathcal O\) with \(cm,cm'\in N\), express \(m=M_0\alpha+u,\ m'=M_0\alpha'+u'\) with \(u,u'\in N\). The factor \(c\sigma(\alpha_i)\alpha_j'\) is divisible by \(\varpi^{\max(e_i,e_j)}\); the terms involving \(u,u'\) are integral after multiplication by \(c\). Hence \(c(\langle ym,ym'\rangle_{J_T}-\langle xm,xm'\rangle_{J_T})\in\mathcal O\).

Put \(N_y=U\cap y^{-1}V,\ \mu(\bar y)=[U:N_y]\). As \(N\subseteq N_y\), \(\mu(\bar y)\le\mu(\bar x)\). At equality \(N_y=N\). The columns \(X_i=\varpi^{e_i}xM_i,\ Y_i=\varpi^{e_i}yM_i\) are each independent modulo \(\varpi\), by independence of the order-\(\varpi\) generators of \(U/N\). Their Gram entries agree modulo \(\varpi^{\min(e_i,e_j)}\). Witt extension over the residue field gives an ambient isometry of the residual columns \cite{scharlau-1985}. Lift it by successive Gram corrections \(Z+Z^\dagger=H\) for Hermitian \(H\), solved by \(Z=H/2\).

If an integral isometry \(h\) matches the columns modulo \(\varpi^{\min(t,e_i)}\), the corrections \((Y_i-hX_i)/\varpi^t\) for \(e_i>t\) prescribe a skew pairing modulo \(\varpi\) by the Gram congruence. Extend from the independent columns to a basis and lift by antisymmetrization to a skew adjoint \(Z\). Left multiplication by \((1+\varpi^tZ/2)(1-\varpi^tZ/2)^{-1}\) makes the next correction and preserves the earlier ones. Iteration gives \(hX_i\equiv Y_i\bmod\varpi^{e_i}V\), or \(y-hx\in\Lambda\). With no positive exponents take \(h=1\).
\leanref{GlobalTheory.\allowbreak{}Hk3.\allowbreak{}GramSep.\allowbreak{}gram\_\allowbreak{}class\_\allowbreak{}family}{Development.\allowbreak{}AbelianVarietyTheory.\allowbreak{}Proofs.\allowbreak{}SatakeParameters.\allowbreak{}ExistsGramClassRepsSelfDualLatticeGroup}{26--257}
\leanref{GlobalTheory.\allowbreak{}Hk3.\allowbreak{}HoweG.\allowbreak{}exists\_\allowbreak{}intUnitary\_\allowbreak{}of\_\allowbreak{}gramClass}{Development.\allowbreak{}AbelianVarietyTheory.\allowbreak{}Proofs.\allowbreak{}SatakeParameters.\allowbreak{}ExistsIntUnitaryOfGramClass}{17--108}

The character average is in \(S\), hence \(K_T\)-invariant. Its part at smaller \(\mu\) belongs to \(S\) by induction. At equality the surviving terms lie on the orbit of \(\bar x\) and have a nonzero coefficient on its line. That coefficient determines an invariant on the orbit. Subtracting the smaller part gives its generator in \(S\). Hence \(M^{K_T}\subseteq S\). The adjoint model exchanges the spaces.
\end{proof}

\subsection*{Proof of \ref{B8.14a}}\label{prf:B8.14a}\phantomsection
\begin{proof}
The unramified quadratic extension of odd residue characteristic supplies a common uniformizer and an anti-invariant unit. The models have unimodular frames, unramified splitting characters and conductor \(\mathcal O_v\). The vacuum is fixed by both integral groups and generates through the Heisenberg translations. Apply Lemma~\ref{B8.14} to the two compact averages, each a finite mean by smoothness. The source average preserves the class of \(x\) and lies in the target span. The target average lies in the source span and is thus a scalar vacuum class modulo \(A\). Multiply by its number of terms.
\end{proof}

\subsection*{Proof of \ref{B8.15d}}\label{prf:B8.15d}\phantomsection
\begin{proof}
Put \(b(u)=(f,H(u)v_0)\). The operators \(H\) are unitary and
\(H(v)H(u)=\psi_0([v,u])H(u)H(v)\). Applying \(H(v)\) to both arguments of the pairing therefore gives \(b(u)=0\) off \(\Gamma^\circ\). The same calculation with a common finite open level of \(f,v_0\) confines \(u_f\) to a compact set.

For the real variables, differentiate the commutator identity at \(v=0\) in a real direction \(e\). Multiplication of \(b(u)\) by the real linear coordinate obtained from \([e,u]\) gives, up to a fixed scalar, a sum of two matrix coefficients with a derivative on \(f\) or \(v_0\). These coefficients are bounded by the products of the respective \(L^2\) norms. Take a real basis and iterate. Nondegeneracy of the form gives decay of every polynomial order in \(u_\infty\), uniformly in \(u_f\). Thus \(b\in L^1\).

Fix \(u\). The function \(f\,\overline{H(u)v_0}\) is periodic on \(X(F)\backslash X_{\mathbb A}\). Its Fourier characters are indexed by \(X(F)\) through \(\psi_0([\ ,\ ])\). The quasiperiodicity formula identifies its Fourier expansion, after multiplication by \(H(u)v_0(x)\), as
\[
 \sum_{\gamma\in X(F)}
 b(u+\gamma)H(u+\gamma)v_0(x)
 = f(x)\,|v_0(x+u)|^2 .
\]
Indeed \(H(u+\gamma)v_0(x)/H(u)v_0(x)\), interpreted as the multiplier before evaluation, is the character \(\psi_0([x,\gamma])\) times a scalar of modulus one independent of \(x\). That scalar cancels against its conjugate in \(b(u+\gamma)\). The finite support restriction and real decay just proved give absolute summability: the contributing rational points lie in a translate of a fixed fractional lattice in the real coordinates. Integrate the identity over a fundamental domain in \(u\). Unfolding the sum gives
\[
 \int_{X_{\mathbb A}} b(u)H(u)v_0(x)\,du=(v_0,v_0)f(x).
\]
The right side is a theta test. Apply the superposition property of \(L\) and the support restriction on \(b\). Finally \((v_0,v_0)>0\), since the continuous periodic function \(|v_0|^2\) is positive on a nonempty open set.
\end{proof}

\subsection*{Proof of \ref{B8.15e}}\label{prf:B8.15e}\phantomsection
\begin{proof}
First consider finite translations. At the selected factors, the coprimalities make the oscillator ideal a unit. There the integral translation lattice is \(e\widehat M\), and the translations fixing the coset indicator include \(e\widehat{\mathfrak m M}\). At the conjugate factors a translation acts by the character of the selected variable under the alternating pairing. For an integral translation that character is trivial on the integral coset. At the remaining factors the vector is the lattice vacuum. These formulas give the two invariances in the definition of \(D(\Gamma_i)\). In combining the factors the Heisenberg product has multiplier \(\psi_0([u,v]/2)\). For doubled integral shifts this multiplier is one by integrality, and for shifts in a single split half the form on that half is zero. Real creation operators commute with all these finite translations. Each \(v_\lambda\) is nonzero: its coset contains the projection of a representative in \(M\), and its other factors are nonzero vacua.

For a finite \(u\in\Gamma_i^\circ\), commuting with the integral translations on the conjugate half forces \(eu\in e\widehat M\) by lattice duality. Reduction modulo \(\mathfrak m\) gives an element \(a\in M/\mathfrak mM\), since completion induces an isomorphism on this finite quotient. Translation sends the indicator at \(\lambda\) to that at \(\lambda-a\). On the conjugate half its multiplier is constant on this coset, since it commutes with the translations by \(e\widehat{\mathfrak mM}\). Away from these pairs the commutant acts on the vacuum by a scalar. Indeed it annihilates the full integral lattice there: away from 2 divide a doubled shift by 2, and above 2 use the two single-half conditions, since the grid is split there. Self-duality puts the translating vector in the vacuum lattice. The central multiplier on recombining the three parts is a scalar as well. Thus \(H(u)\bar v_0\) belongs to the span of the \(\bar v_\lambda\), where the subscript \(0\) denotes the zero coset.
\leanref{GlobalTheory.\allowbreak{}CosShift.\allowbreak{}epart}{Development.\allowbreak{}AbelianVarietyTheory.\allowbreak{}Proofs.\allowbreak{}TwistedOperators.\allowbreak{}CommutantDepthLatticePermutesCosetVectors}{106--124}
\leanref{GlobalTheory.\allowbreak{}CosPerm.\allowbreak{}cosVec\_\allowbreak{}perm}{Development.\allowbreak{}AbelianVarietyTheory.\allowbreak{}Proofs.\allowbreak{}TwistedOperators.\allowbreak{}HeisenbergCommutantPermutesFirstRowCosetVectors}{49--95}

For the real calculation work in the model for \((-S,T)\), with the linear functional \(L'(f)=\overline{L(\bar f)}\). It has the superposition property of \(L\) by conjugation of the coefficient under the integral. Fix \(v=v_\lambda\). In the orthonormal Fock frames write \(z=(z_j)\) for the complex coordinates of a real translation, with coordinate vector \(e_j\) and complex structure \(J\). Put \(c_j=2\pi|\tau(\xi_0)|\) at the place \(\tau\) of \(j\), where \(\xi_0\) is the nonzero imaginary scalar in the trace alternating form. Put \(q(z)=\sum_j c_j|z_j|^2\). For each \(x\),
\[
 R_x(z)=e^{q(z)/2}H(z)v(x)
\]
is entire. To see this, differentiate the Heisenberg product law. The derivative of the central phase in the \(\bar z_j\) direction cancels the derivative of the Gaussian factor, leaving the action of \((dH(e_j)+i\,dH(Je_j))/2\) on \(v\), which is zero by the vacuum equation. The complementary derivative at zero gives
\[
 \partial^\alpha R_x(0)
   =\left(\prod_j c_j^{\alpha_j/2}\right)
        ((a^+)^\alpha v)(x),
 \qquad
 a_j^+=\frac{dH(e_j)-i\,dH(Je_j)}{2\sqrt{c_j}} .
\]
For repeated derivatives, moving a creation operator past \(H(z)\) adds a linear function of \(\bar z\). Holomorphic differentiation followed by evaluation at zero kills all such terms.

Set \(F(z)=L'(H(z)v)\). This is bounded. Gaussian integration of the Taylor series of \(R_x\) against \(\bar z^\alpha e^{-q(z)}\) gives
\[
 \int \bar z^\alpha e^{-q(z)/2}H(z)v(x)\,dz
 =\left(\prod_j\frac{\pi}{c_j}\,c_j^{-\alpha_j}\right)
       \partial^\alpha R_x(0).
\]
Here \(dz\) is product Lebesgue measure. Angular integration extracts the indicated coefficient and radial integration gives
\(\int_{\C}|z|^{2a}e^{-c|z|^2}\,dz=\pi a!c^{-a-1}\).
The bound \(|R_x(z)|\le\|v\|_\infty e^{q(z)/2}\) justifies the integration. If \(L\) kills all prescribed tests, applying \(L'\) shows that every moment
\(\int\bar z^\alpha e^{-q(z)/2}F(z)\,dz\) vanishes.

For any \(z_0\), multiply these moments by the coefficients of
\(\exp(\sum_j c_jz_{0j}\bar z_j)\) and sum. Absolute convergence follows from the integrable majorant
\(\|F\|_\infty\exp(-q(z)/2+\sum_j c_j|z_{0j}z_j|)\).
The same angular and radial integration, with \(R_x\) in place of the monomial, gives the reproducing identity
\[
 H(z_0)v(x)
 =e^{-q(z_0)/2}\prod_j(c_j/\pi)
   \int e^{\sum_jc_jz_{0j}\bar z_j-q(z)/2}H(z)v(x)\,dz .
\]
Its left side is a theta test, so applying \(L'\) gives \(F(z_0)=0\). This proves vanishing on every real translate of every conjugate coset vacuum.
\leanref{GlobalTheory.\allowbreak{}FockTransl.\allowbreak{}Jet.\allowbreak{}jet\_\allowbreak{}eq\_\allowbreak{}fock}{Development.\allowbreak{}AbelianVarietyTheory.\allowbreak{}Proofs.\allowbreak{}TwistedOperators.\allowbreak{}FockGaussianJetsAtOriginEqFockVectors}{42--88}
\leanref{GlobalTheory.\allowbreak{}FockTransl.\allowbreak{}tame\_\allowbreak{}kills\_\allowbreak{}real\_\allowbreak{}translates}{Development.\allowbreak{}AbelianVarietyTheory.\allowbreak{}Proofs.\allowbreak{}TwistedOperators.\allowbreak{}TameFunctionalKillsFockVacuumRealTranslates}{28--120}

Finally, if \(u\in\Gamma_i^\circ\), then \(u_f\in\Gamma_i^\circ\), since every level condition is finite. Real and finite directions are orthogonal, so \(H(u)=H(u_\infty)H(u_f)\). The finite calculation followed by the real calculation gives \(L(H(u)\bar v_0)=0\). Lemma~\ref{B8.15d} with this nonzero vacuum gives the assertion.
\end{proof}

\subsection*{Proof of \ref{B8.15f}}\label{prf:B8.15f}\phantomsection
\begin{proof}
Conjugate the downward identities and the bottom period. Use the conjugate Weil action, with both splitting characters inverted. If \(u_i\) is a partial row lift, the next lift with test \((\lambda,P)\), after conjugation, is the lift of \(\bar u_i\) with \(E_i(\lambda,P)\). A change of the test at \(i\) leaves \(u_i\) unchanged. Write \(\Theta_f\) for a lift with test \(f\), in the direction indicated by its argument.

We record why the period functionals in this argument admit the integrations of Lemma~\ref{B8.15d}. For a fixed theta test \(b\), with \(\omega\) its Weil action, the bound needed is
\[
 \sup_{w,x}|(\omega(g,h)H(w)b)(x)|
       \le C_b\bigl(H_S(g)H_T(h)\bigr)^{A_b},
\]
where \(H_S,H_T\) are quotient heights, the infima of the matrix heights over rational left translates. Intertwining moves \(H(w)\) to a translation of \(\omega(g,h)b\) with unit multiplier. For the resulting supremum use a fixed finite level of \(b\). On each finite-level component the Sobolev bound on the compact additive quotient uses finitely many real Heisenberg derivatives. Moving these derivatives through \(\omega(g,h)\) acts on their directions by the matrix action \(x\mapsto h^{-1}xg\). Each word therefore costs a polynomial in the matrix heights, and unitarity bounds the \(L^2\) norms of the fixed differentiated tests. Reduction to the finite set of level components gives the matrix-height bound. Rational left translation preserves the supremum, so representatives of height at most twice the infimum give the displayed bound.
\leanref{GlobalTheory.\allowbreak{}WeilWordGrowth.\allowbreak{}om\_\allowbreak{}pointwise}{Development.\allowbreak{}AbelianVarietyTheory.\allowbreak{}Proofs.\allowbreak{}TwistedOperators.\allowbreak{}HeisenbergWordsWeilOperatorSupNormPolynomialGrow}{65--130}
\leanref{GlobalTheory.\allowbreak{}KerTame.\allowbreak{}om\_\allowbreak{}heis\_\allowbreak{}bound}{Development.\allowbreak{}AbelianVarietyTheory.\allowbreak{}Proofs.\allowbreak{}TwistedOperators.\allowbreak{}WeilKernelFunctionalSuperpositionIsTame}{30--36}

In a Bessel integral the unipotent character has modulus one. In a Fourier--Jacobi integral its additional line theta factor satisfies the same bound for its fixed line test. The unipotent quotient is compact. In a fixed Witt decomposition the height of the embedded smaller group times a representative in this compact quotient is bounded by a power of its own quotient height: use a rational translate in the smaller group and the block embedding in the larger group. Rapid decrease of the two automorphic factors now gives an integrable majorant independent of \(w\). For a superposition with coefficient \(a(w)\) multiply this majorant by \(|a(w)|\). Weil covariance and Fubini give the superposition identity, and dominated convergence gives continuity of the orbit function. This argument applies also after contraction against a fixed line test and integration on the core.
\leanref{GlobalTheory.\allowbreak{}LiftTame.\allowbreak{}lift\_\allowbreak{}periods\_\allowbreak{}tame}{Development.\allowbreak{}AbelianVarietyTheory.\allowbreak{}Proofs.\allowbreak{}TwistedOperators.\allowbreak{}ThetaLiftPeriodsIsTameThetaFunctional}{95--154}

At the bottom write \(V_0=Y_0\perp L_0\), with \(L_0\) the distinguished line. The two orthogonal projections preserve the bottom depth lattice. In the grid frames they are integral summand projections and commute with the selected split idempotents. After the rank-zero Witt identification the two projected lattices are the first row depth lattice and the line depth lattice. Thus their direct sum is the core depth lattice. Choose nonzero vacua \(b,c\) on the two summands. On the row one may take the zero coset vacuum. On the line use the lattice vacuum, restricted to the induced depth conditions. The product \(b\otimes\bar c\) is a nonzero core depth test. Its translates by the core annihilator are products of translates at the respective depths, since the alternating form is the orthogonal sum. The sign change on the line is accounted for by conjugation.

As a functional of the core test, the conjugate bottom period takes a nonzero value at the conjugate downward test defining \(F_0\). By Lemma~\ref{B8.15d} it takes a nonzero value on one of the translated products. On a product \(b'\otimes\bar c'\), the kernel at the stabilizer is
\[
 (\omega_{\rm row}(g,h)b')(0)\,
       \overline{(\omega_{\rm line}(g,1)c')(0)}.
\]
The character on the orthogonal sum is the product of the two splitting characters with the line character inverted. Interchange the \(g\) and \(h\) integrations. The result is the Fourier--Jacobi period at height \(1\) of \(\bar F_1\) and \(\Theta_{b'}\bar u\) with line test \(c'\). The flag has length zero at this height. Fix \(c'\) and apply Lemma~\ref{B8.15e} to the functional of \(b'\). This gives a prescribed first test with nonzero period.
\leanref{GlobalTheory.\allowbreak{}LastSeeSaw.\allowbreak{}lastLat\_\allowbreak{}std}{Development.\allowbreak{}AbelianVarietyTheory.\allowbreak{}Proofs.\allowbreak{}TwistedOperators.\allowbreak{}LastSeesawLatticesNonzeroLineVacua}{125--146}
\leanref{GlobalTheory.\allowbreak{}LastSeesaw.\allowbreak{}upLast}{Development.\allowbreak{}AbelianVarietyTheory.\allowbreak{}Proofs.\allowbreak{}TwistedOperators.\allowbreak{}LastSeesawPeriodNonvanishing}{195--278}
\leanref{GlobalTheory.\allowbreak{}AscBottom.\allowbreak{}rung0}{Development.\allowbreak{}AbelianVarietyTheory.\allowbreak{}Proofs.\allowbreak{}TwistedOperators.\allowbreak{}NonvanishingThetaTowerBottomRung}{27--103}

Here are the two exchanges used above height \(1\). At height \(2r+1\) remove \(r\) hyperbolic planes from the row space. Its middle space is \(V_{2r+1}\). At height \(2r+2\) remove \(r\) hyperbolic planes. Its middle space is \(V_{2r+2}\perp L\), with \(L\) the distinguished line. Use the ordered isotropic columns \(e_1,\ldots,e_r\) of these Witt frames. Let \(\zeta_j=(-\xi_0)^j\zeta_0\) in the line frame. A Bessel period integrates over the flag unipotent and the unitary group fixing the distinguished line vector, against the character on the successive simple entries and the last entry paired with that vector. A Fourier--Jacobi period integrates over the flag unipotent and the unitary group of its middle space, against the flag character and the conjugate line theta kernel. All unipotent integrations here use their rational quotients.

For the exchange from height \(i\) to \(i+1\), put \(V=Y_{i+1}\), \(W=Y_i\), and denote their Witt middle spaces by \(V^\natural,W^\natural\). In the two mixed Schrödinger coordinates write a full test as \(b(X,Y,w)\), where \(X\) is the tuple of columns in \(W\) corresponding to the first isotropic block of \(V\), \(Y\) is the dual coordinate for the isotropic block of \(W\), and \(w\in\operatorname{Hom}_E(V^\natural,W^\natural)_{\mathbb A}\). With \(U_a\) the upper unitriangular group of size \(a=\#X\), its generic character is
\(\chi_U(u)=\psi_0(\operatorname{Tr}_{E/F}\sum_{j<a}u_{j,j+1})\), with indices starting at one. The restriction of a test needed after unfolding is
\[
 Rb(w)=\int_{U_a(\mathbb A_E)}
        \overline{\chi_U(u)}\,b(xu,y,w)\,du .
\]
For \(i=2r+1\), \(x=(e_1,\ldots,e_r)\) and \(y\) has last column \(-\zeta_{r+1}\) and all other columns zero. For \(r=0\) both tuples are empty. For \(i=2r+2\), \(x=(e_1,\ldots,e_r,\zeta_{r+1})\) and \(y=0\). Each \(x\) has independent columns. The stabilizer \(R_x=\{h:hx=xu,\ u\in U_a\}\) is the Jacobi group in the first case and the Bessel group in the second. This follows by writing a flag-preserving isometry in Witt blocks. In the second case its middle block must also fix \(\zeta_{r+1}\).

We describe the finite restrictions since they determine the level of the resulting test. Use the pullback of \(\Gamma_i\) to the first translation block as the period lattice \(A_1\) in \(X\), and its character dual as the opposite block. For the second pair of blocks use the pullbacks of the integral oscillator lattice with target lattice multiplied by \(\mathfrak m/\bar{\mathfrak m}\). They are dual under the alternating pairing. Denote the resulting period lattice in \(Y\) by \(B\), and let \(\Gamma_*\) be the pullback of \(\Gamma_i\) under extension by zero on the two hyperbolic parts. Use tests with finite support in \(x+A_1\) and \(y+B\) in these coordinates, invariant under the respective period lattices and with middle test in \(D(\Gamma_*)\). Equivalently the opposite translations have the characters prescribed by the centers \(x,y\). In the standard grid frames, decomposition into the two hyperbolic parts and the middle is integral. On the first dual block and the second dual block the characters of a row depth shift are trivial, since their pairings with \(x,y\) are integral for \(\psi_0\). Consequently these tests belong to \(D(\Gamma_i)\). They are stable under real Heisenberg translation and under extension by zero of a translation in \(\Gamma_*^\circ\), and
\[
 R(H(\widetilde w)b)=H(w)Rb\qquad(w\in\Gamma_*^\circ)
\]
for this extension \(\widetilde w\). These assertions also follow directly in mixed coordinates: a translation in the middle acts on the last argument and commutes with evaluation of the first two. Denote this class of full tests by \(\mathcal S\).
\leanref{GlobalTheory.\allowbreak{}PcBF.\allowbreak{}pieces\_\allowbreak{}bf}{Development.\allowbreak{}AbelianVarietyTheory.\allowbreak{}Proofs.\allowbreak{}TwistedOperators.\allowbreak{}BesselToFourierJacobiMixedPiecesAdmissible}{50--67}
\leanref{GlobalTheory.\allowbreak{}AscFin.\allowbreak{}piecesFB\_\allowbreak{}is\_\allowbreak{}and\_\allowbreak{}splits}{Development.\allowbreak{}AbelianVarietyTheory.\allowbreak{}Proofs.\allowbreak{}TwistedOperators.\allowbreak{}FourierJacobiStepPiecesAdmissibleSplits}{19--26}
\leanref{GlobalTheory.\allowbreak{}AscFin.\allowbreak{}GoodStd.\allowbreak{}loc\_\allowbreak{}cls}{Development.\allowbreak{}AbelianVarietyTheory.\allowbreak{}Proofs.\allowbreak{}TwistedOperators.\allowbreak{}StandardAuxiliaryFramesGoodRowGrids}{94--117}

There is a \(b_*\in\mathcal S\) with \(Rb_*\ne0\). Take the two finite coset indicators above, smooth compactly supported real functions equal to one at \(x,y\), and a nonzero middle lattice vacuum. The finite \(u\) for which \(xu-x\in A_1\) have \(\chi_U(u_f)=1\): the successive entries in a standard isotropic frame lie in the conductor of the trace character. At infinity use a nonnegative first function. Since \(x\) has a left inverse, its support can be chosen so small around \(x\) that \(\Re\chi_U(u_\infty)>0\) wherever it contributes. The integral defining \(Rb_*\) then has a scalar factor of positive real part.

Write \(K_i\) for the compact open row level at height \(i\). The restriction \(R(\omega(1,h)b)\), for \(b\in\mathcal S\), is supported on \(R_x\U(W)(F_\infty)K_i\). We give the finite-place calculation. Nonvanishing requires \(h^{-1}xu-x\in A_1\) locally for some unitriangular \(u\). Translation of \(u\) by a strictly upper triangular entry preserving this coset must have trivial simple-root character. The second dual block must preserve the character specified by \(y\). The same conditions hold after replacing \(h\) by \(hk\), \(k\in K_i\), since the row level fixes the test. At a marked split place this last invariance follows from the action on the selected coset modulo depth and triviality of the splitting characters on the principal determinant units. Off the marked places the self-dual lattice vacuum is fixed by the integral group.

At a marked place write the split action on the two components as a matrix and its inverse transpose in dual frames. After scaling the specified columns and dual rows to coordinate vectors and covectors, the condition on \(h^{-1}xu\) prescribes columns \(\xi_p\equiv e_{a_p}\) and covectors \(\eta_q\equiv e_{b_q}^*\) modulo the level ideal. Their pairings are \(\eta_q\xi_p=\delta_{b_q,a_p}\): the Gram matrix of \(x\) is zero except possibly at its last diagonal entry, which upper unitriangular changes preserve. Complete the prescribed columns by the unused standard columns to a matrix \(C\equiv1\). Replace rows \(b_q\) of the identity by \(\eta_q C\) to form \(D\equiv1\). The pairing identity gives \(D e_{a_p}=e_{a_p}\). Thus \(CD^{-1}\) realizes the prescribed columns and inverse row action and belongs to the principal congruence group. Removing this element leaves an element of \(R_x\).

At a good place use the Witt Iwasawa decomposition for the self-dual row lattice and remove the integral right factor. Write \(t_1,\ldots,t_r\) for the diagonal entries on the isotropic flag. The column condition gives \(t_j^{-1}\) integral. On the root \(j,j+1\) the translation condition forces \(t_j\) integral: \(t_j\) multiplies the ideal on which the primitive simple-root character is trivial into that same ideal. For the isotropic tuple the final entry uses the second character, paired with the primitive last column of \(y\). For the tuple with an appended line vector the final simple root is available as well. All the \(t_j\) are units. In the appended-line case the column condition also makes the transported line lattice integral in the middle lattice. It is a self-dual summand there, since the transported vector has its prescribed norm and the initial line is a self-dual summand. An integral isometry carries it, with its generator in a local line frame, to the original summand: split off the line and identify its unimodular orthogonal complement by the local classification. At a dyadic place the grid is split and this identification is ordinary completion to dual bases. The remaining middle isometry fixes the distinguished vector. In the isotropic-tuple case the entire middle group already belongs to \(R_x\). This proves the support assertion on taking the restricted product.
\leanref{GlobalTheory.\allowbreak{}AscFinNe.\allowbreak{}flat\_\allowbreak{}ne\_\allowbreak{}zero}{Development.\allowbreak{}AbelianVarietyTheory.\allowbreak{}Proofs.\allowbreak{}TwistedOperators.\allowbreak{}ExistsUnfoldingFlatMapNeZeroFullRank}{42--74}
\leanref{GlobalTheory.\allowbreak{}AscFin.\allowbreak{}OneOrbitAt.\allowbreak{}core\_\allowbreak{}coord}{Development.\allowbreak{}AbelianVarietyTheory.\allowbreak{}Proofs.\allowbreak{}TwistedOperators.\allowbreak{}SplitPlaceTuplesOneOrbitModUniformizerPower}{103--212}
\leanref{GlobalTheory.\allowbreak{}OneCosetGood.\allowbreak{}oneCoset\_\allowbreak{}good}{Development.\allowbreak{}AbelianVarietyTheory.\allowbreak{}Proofs.\allowbreak{}TwistedOperators.\allowbreak{}GoodPlaceOneCosetOfBesselLevel}{119--141}
\leanref{GlobalTheory.\allowbreak{}AscFin.\allowbreak{}OneCosetFJ.\allowbreak{}oneCoset\_\allowbreak{}good}{Development.\allowbreak{}AbelianVarietyTheory.\allowbreak{}Proofs.\allowbreak{}TwistedOperators.\allowbreak{}FourierJacobiLevelOneCosetAtGoodPlace}{241--254}

Suppose the outgoing period vanishes on \(\mathcal S\). Unfold its flag integral in the first mixed coordinate. Integration over the rational additive quotients selects the tuples with the prescribed Gram matrix and successive flag character. The surviving full-rank tuples form the orbit of \(x\) under rational isometries and upper unitriangular changes: an isometry between their column spans extends by Witt extension. A dependent tuple gives zero by integration over a stabilizing root group with nontrivial character. The integral is therefore over \(R_x(F)\backslash\U(W)(\mathbb A_F)\), with the integration in \(u\) appearing in \(R\). The second mixed coordinate is evaluated at \(y\) by additive character orthogonality.

For precision about cancellation in this integral, translate the test by a real vector \(c\) in the first opposite isotropic block. This keeps the test in \(\mathcal S\) and multiplies the unfolded integrand by the character of \(c\) paired with \((h^{-1}xu)_\infty\). All Fourier transforms of the pushforward of the integrand to this real matrix space thus vanish. The integrand is absolutely integrable by the height bound above and rapid decrease. Uniqueness of Fourier transform for finite measures gives zero integral on every strip \(R_x C\U(W)(\mathbb A_{F,f})\), \(C\) compact in the real group. Indeed membership is determined by whether the real matrix lies in \(C^{-1}xU_a(F_\infty\otimes_F E)\). This is a measurable set, since the full-rank matrix \(x\) has a left inverse and the unitriangular group is closed.

Average first along \(R_x(F)\backslash R_x(\mathbb A_F)\). The resulting mean is continuous, left \(R_x(\mathbb A_F)\)-invariant and right \(K_i\)-invariant. If its value at the identity were \(z\ne0\), choose a compact real neighborhood \(C\) of the identity on which the real part of the mean divided by \(z\) exceeds \(1/2\). By the support assertion, every nonzero value on the strip is represented by a point of \(C\): the stabilizer splits into its real and finite components and the right finite factor can be taken in \(K_i\). The strip integral divided by \(z\) would have positive real part, since \(K_i\) is open. Hence the mean at the identity vanishes. This mean is the incoming functional on \(Rb\). For the isotropic tuple it is the Fourier--Jacobi integral with the entire middle test \(Rb\). For the appended tuple it is the Bessel integral after contracting \(Rb\) with the fixed line test in the line variable.
\leanref{GlobalTheory.\allowbreak{}StripVanish.\allowbreak{}integral\_\allowbreak{}strip\_\allowbreak{}eq\_\allowbreak{}zero}{Development.\allowbreak{}AbelianVarietyTheory.\allowbreak{}Proofs.\allowbreak{}TwistedOperators.\allowbreak{}IntegralStripEqZeroOfRealModulationVanishing}{70--86}
\leanref{GlobalTheory.\allowbreak{}Asc.\allowbreak{}vanBF}{Development.\allowbreak{}AbelianVarietyTheory.\allowbreak{}Proofs.\allowbreak{}TwistedOperators.\allowbreak{}VanishingDescendsBesselToFourierJacobi}{127--224}
\leanref{GlobalTheory.\allowbreak{}Asc.\allowbreak{}VanJac.\allowbreak{}vanFB}{Development.\allowbreak{}AbelianVarietyTheory.\allowbreak{}Proofs.\allowbreak{}TwistedOperators.\allowbreak{}VanishingDescendsFourierJacobiToBessel}{237--325}

In particular the incoming functional vanishes on \(H(w)Rb_*\) for \(w\in\Gamma_*^\circ\), by translating \(b_*\) before restriction. Lemma~\ref{B8.15d} gives vanishing on every middle test in \(D(\Gamma_*)\). Evaluate this conclusion on the following products.

At \(i=2r+1\), let \(b_1\) be the conjugate downward core test and \(c\) the incoming line test. The link of lines sends \(\zeta_r\) to \(-\zeta_{r+1}\) by a similitude of multiplier \(\xi_0\bar\xi_0\). Composition with this map followed by the normalized Hermitian adjoint gives an invertible map \(k\) from the old line Hom space to the new one, reversing the alternating form. Set \(b_2=\overline{c\circ k^{-1}}\). Then \(\overline{b_2\circ k}=c\), with the original level. The middle projections send the middle depth shifts into the core depth shifts and, through the link, into the line depth shifts. Hence \(b_1\otimes b_2\in D(\Gamma_*)\). The Weil action on this product is the product of the core and line actions. Its line argument is \(k\) of the incoming argument and its central phase has opposite sign. Evaluating the middle functional therefore gives exactly the incoming Fourier--Jacobi period, with \(\bar F_i=\Theta_{b_1}\bar F_{i+1}\).

At \(i=2r+2\), choose a nonzero vacuum \(c\) at the new line depth, and use the middle test \(b_1\otimes c\). More explicitly, if \(\mathfrak a=(\mathfrak c_K\mathcal O_E)\mathfrak b^{-1}\) and \(\mathfrak s=\mathfrak a/\bar{\mathfrak m}\), the line lattice for flag length \(j=r+1\) is formed with the line module
\((\mathfrak s/\bar{\mathfrak s})^j(\mathfrak m/\bar{\mathfrak m})\mathcal O_E\zeta_0\).
This is the line component in the successive Witt frames. Projection to that component and to the core carries a middle depth shift into the respective lattices. At each split half the projections preserve that half. This gives \(b_1\otimes c\in D(\Gamma_*)\). Contracting the line factor against \(c\) gives \((c,c)\) by unitarity of the line action. The remaining integral is the incoming Bessel period with \(\bar F_i=\Theta_{b_1}\bar F_{i+1}\). Its multiplier \((c,c)\) is positive.
\leanref{GlobalTheory.\allowbreak{}LinkInvTheta.\allowbreak{}link\_\allowbreak{}inv\_\allowbreak{}theta}{Development.\allowbreak{}AbelianVarietyTheory.\allowbreak{}Proofs.\allowbreak{}TwistedOperators.\allowbreak{}LinkedLineVectorsInvertibleOnThetaFunctions}{69--131}
\leanref{GlobalTheory.\allowbreak{}LinkProd.\allowbreak{}perAll\_\allowbreak{}box\_\allowbreak{}eq\_\allowbreak{}per}{Development.\allowbreak{}AbelianVarietyTheory.\allowbreak{}Proofs.\allowbreak{}TwistedOperators.\allowbreak{}PeriodAllDataProductsEqFourierJacobiPeriod}{51--87}
\leanref{GlobalTheory.\allowbreak{}HandProd.\allowbreak{}handP\_\allowbreak{}boxT}{Development.\allowbreak{}AbelianVarietyTheory.\allowbreak{}Proofs.\allowbreak{}TwistedOperators.\allowbreak{}BesselHandedFunctionProductEqCoreLiftMulPairing}{16--38}

In each case a nonzero incoming period contradicts vanishing of the outgoing period on \(\mathcal S\). The latter therefore has a nonzero value in \(D(\Gamma_i)\). Lemma~\ref{B8.15e} gives a pair \((\lambda_i,P_i)\) with nonzero value. The partial row is right \(K_i\)-invariant at every step where the strip argument uses it: the preceding prescribed test is fixed by the output level, which can be moved onto the kernel by right translation. Updating the two entries at \(i\) extends the existing partial row lift.
\leanref{GlobalTheory.\allowbreak{}UpBessel.\allowbreak{}upBF}{Development.\allowbreak{}AbelianVarietyTheory.\allowbreak{}Proofs.\allowbreak{}TwistedOperators.\allowbreak{}FourierJacobiPeriodGivesBesselPeriodOfLift}{18--44}
\leanref{GlobalTheory.\allowbreak{}UpFromBessel.\allowbreak{}upFB}{Development.\allowbreak{}AbelianVarietyTheory.\allowbreak{}Proofs.\allowbreak{}TwistedOperators.\allowbreak{}BesselPeriodGivesFourierJacobiPeriodOfLift}{64--83}
\leanref{GlobalTheory.\allowbreak{}AscRungB.\allowbreak{}rungB}{Development.\allowbreak{}AbelianVarietyTheory.\allowbreak{}Proofs.\allowbreak{}TwistedOperators.\allowbreak{}FourierJacobiToBesselNonvanishingTower}{39--109}
\leanref{GlobalTheory.\allowbreak{}Asc.\allowbreak{}RungJac.\allowbreak{}rungJ}{Development.\allowbreak{}AbelianVarietyTheory.\allowbreak{}Proofs.\allowbreak{}TwistedOperators.\allowbreak{}BesselToFourierJacobiNonvanishingTower}{57--133}

Thus height \(1\) gives height \(2\), and a Bessel period at height \(2r+2<m\) gives a Fourier--Jacobi period at \(2r+3\) and then a Bessel period at \(2r+4\). The grid has even \(m\ge2\). At height \(m\) the frame character is the conjugate of the prescribed top character. Indeed the successive isotropic columns have scale factors \((-\xi_0)^j\), \(1\le j\le(m-2)/2\), and the marked line has the next scale factor. Each simple entry and the final line entry therefore acquire the same ratio \(-\xi_0\). Conjugating the integral gives the stated top period up to the positive ratio of the two Haar normalizations.
\leanref{GlobalTheory.\allowbreak{}AscFin.\allowbreak{}EndsPf.\allowbreak{}Chr.\allowbreak{}psiBes\_\allowbreak{}eq\_\allowbreak{}star\_\allowbreak{}besChar}{Development.\allowbreak{}AbelianVarietyTheory.\allowbreak{}Proofs.\allowbreak{}TwistedOperators.\allowbreak{}StandardAuxiliaryFramesFitGridEnds}{129--157}
\leanref{GlobalTheory.\allowbreak{}AscFin.\allowbreak{}EndsPf.\allowbreak{}per\_\allowbreak{}conj}{Development.\allowbreak{}AbelianVarietyTheory.\allowbreak{}Proofs.\allowbreak{}TwistedOperators.\allowbreak{}StandardAuxiliaryFramesFitGridEnds}{213--238}
\end{proof}

\subsection*{Proof of \ref{B8.15}}\label{prf:B8.15}\phantomsection
\begin{proof}
Fix \(w\ne0\) and a nonzero row \(t\). Use the threshold of Lemma~\ref{B8.15b}. The level primes \(p_{\rm de},p'_{\rm de}\) split away from \(Q_Sq\). The size hypotheses give \(44^4\le m^4\le f\le k\) and \(r=k+m-1>1800\). The grid separation with weight-gap constant \(10\) gives consecutive center differences greater than twice the cube of the grid rank bound, as required by Lemma~\ref{B8.15c}. The given realization supplies the bottom discrete-series model and its compatible Cartan frame.

Put \(u=t(w)\). Lemma~\ref{B8.15b} gives a core value with nonzero period against \(u\). By Lemma~\ref{B8.15c} it is a finite linear combination of bottom values of downward column chains. Each bottom value is bounded by rapid decrease, and \(u\) is integrable by cuspidal rapid decrease in Proposition~\ref{B1.15b}. Integrate the finite combination. One chain has nonzero bottom period. Fix that chain, including its terminal column \(\Psi\) and coefficient covector \(\xi\).
\leanref{GlobalTheory.\allowbreak{}RowSizes.\allowbreak{}sizes\_\allowbreak{}of\_\allowbreak{}sizeRows}{Development.\allowbreak{}AbelianVarietyTheory.\allowbreak{}Proofs.\allowbreak{}TwistedOperators.\allowbreak{}SizesOfSizeRows}{14--35}
\leanref{GlobalTheory.\allowbreak{}SepRows.\allowbreak{}WeakSepPf.\allowbreak{}exists\_\allowbreak{}reach\_\allowbreak{}of\_\allowbreak{}span}{Development.\allowbreak{}AbelianVarietyTheory.\allowbreak{}Proofs.\allowbreak{}TwistedOperators.\allowbreak{}RowSeparationByRealFunctionalsOfChainSteps}{40--68}

For the partial row lifts, at each arrow the real theta intertwiner takes the previous smooth vector and the polynomial of that arrow to a vector of the next discrete series. The finite coset test takes the previous multiplicity map to a cuspidal multiplicity map there. Indeed the signed source and target lists are the successive grid lists with the shift of the splitting character at \(m-i\); the two splitting exponents are that center with the half-unit correction prescribed by the ranks. The lower real lifts vanish by the real parameter rule, so the proper constant terms vanish by the tower property of Theorem~\ref{B3.00e}. The zero lift is included. At a split depth place each coset vacuum has the same real vacuum equations, since its construction changes the finite factor. Iterating gives convergence and expresses each partial lift as a smooth vector evaluated under a cuspidal multiplicity map. Cuspidal rapid decrease in Proposition~\ref{B1.15b} thus applies at every partial height. Lemma~\ref{B8.15f} gives the polynomials \(\rho=(P_i)_i\) and the cosets \(\Phi=(\lambda_i)_i\) for the fixed terminal column. In particular
\begin{equation}\label{eq:p8c2-weak-separation}
 \forall w\ne0\quad\forall t\ne0\quad
 \exists\rho,\xi,\Psi,\Phi\quad
 \mathcal B_{\rho,\xi,\Psi,\Phi}(t(w))\ne0 .
\end{equation}
\leanref{GlobalTheory.\allowbreak{}RowArrow.\allowbreak{}separatingFamily\_\allowbreak{}extends\_\allowbreak{}by\_\allowbreak{}twisted\_\allowbreak{}row}{Development.\allowbreak{}AutomorphicArithmetic.\allowbreak{}ThetaCorrespondence.\allowbreak{}ThetaPairLifts.\allowbreak{}LocalTypeDictionary.\allowbreak{}SeparatingFamilyExtendsByTwistedRow}{16--77}
\leanref{GlobalTheory.\allowbreak{}SepRows.\allowbreak{}RowRapidPf.\allowbreak{}rowRapid}{Development.\allowbreak{}AbelianVarietyTheory.\allowbreak{}Proofs.\allowbreak{}TwistedOperators.\allowbreak{}LevelTowerRapidDecay}{20--34}

It remains to fix \((\rho,\xi)\) before choosing \(t\). Write \(v_\rho(w)\) for the terminal vector obtained by composing the real intertwiners. It is independent of the finite cosets. At each real place the terminal row representation is a smooth discrete series, and the coefficient representation of the definite middle group is finite-dimensional and irreducible. The two are smooth admissible representations of moderate growth with infinitesimal character. The isotropic flag, negative distinguished vector and its positive definite orthogonal complement give the unitary Bessel subgroup. Its unitary character is nontrivial on each of the indicated simple-root coordinates. The terminal signed parameter is dual to the translated grid list. Translation and duality preserve the absolute differences on opposite-sign pairs, so the terminal gap in the grid gives
\begin{equation}\label{eq:p8c2-bessel-gap}
 \operatorname{gap}_{\rm opp}\ge120(\operatorname{rank}Y_m)^3 .
\end{equation}
Fixing the other real arguments of the period gives a continuous equivariant bilinear form. The space of these forms has dimension at most one by \cite[Theorem~A]{jiang-sun-zhu-2010}.
\leanref{UpperTheory.\allowbreak{}Lit.\allowbreak{}Bur.\allowbreak{}Gap.\allowbreak{}fit\_\allowbreak{}gap}{Development.\allowbreak{}AutomorphicArithmetic.\allowbreak{}ThetaCorrespondence.\allowbreak{}BesselModels.\allowbreak{}GapArith}{293--310}
\leanref{UpperTheory.\allowbreak{}Lit.\allowbreak{}Bur.\allowbreak{}burGap\_\allowbreak{}model}{Development.\allowbreak{}AutomorphicArithmetic.\allowbreak{}ThetaCorrespondence.\allowbreak{}BesselModels.\allowbreak{}SepGap}{34--60}

If the row space is zero, any real pair works. Otherwise \eqref{eq:p8c2-weak-separation} gives a nonzero global period, so each local space of forms is nonzero. Choose a generator at each real place and arguments on which it has value one. Evaluate successively at those arguments. Multilinearity on pure tensors and linearity on their span give
\begin{equation}\label{eq:p8c2-common-real-factor}
 \mathcal B_{\rho,\xi,\Psi,\Phi}(t(w))
    =A_w(\rho,\xi)c_{t,\Psi,\Phi},
 \qquad A_w(\rho,\xi)=\ell(\xi,v_\rho(w)),
\end{equation}
where \(\ell\) is the tensor product of the chosen generators. Apply \eqref{eq:p8c2-weak-separation} to one nonzero row to choose \((\rho,\xi)\) with \(A_w(\rho,\xi)\ne0\). For each nonzero \(t\), the same equation gives \((\Psi,\Phi)\) with \(c_{t,\Psi,\Phi}\ne0\). Their product detects \(t\) at the chosen real pair.
\end{proof}

\subsection*{Proof of \ref{B8.15g}}\label{prf:B8.15g}\phantomsection
\begin{proof}
The weights of the dual defining representation are \(-e_s\), each with multiplicity one. Every \(\lambda-e_s\) is regular and in the chamber of \(\lambda\), by the spacing assumption. The discrete-series coherent character identity therefore gives the sum of the characters of precisely these discrete series \cite[\S\S7--8, in particular (8.1)--(8.4)]{schmid-vilonen-1998}. Here the Cartan subgroup is the diagonal compact torus, and the shifts are integral characters of that torus. The decreasing lists \(\lambda-e_s\) are distinct, so the summands have distinct infinitesimal characters. Taking the primary spaces for the center gives a direct sum on compact-finite vectors, each primary space having length one. Smooth globalization gives the same decomposition on smooth vectors by \ref{B8.38a}. Write \(R_s\) for projection followed by an isomorphism to the smooth unitary model of \(\pi_s\).

Let \(\epsilon_l\) be the signature of coordinate \(l\). The tensor action preserves the nondegenerate Hermitian form
\[
 [y,x]=\sum_l\epsilon_l\langle y_l,x_l\rangle ,
\]
with the inner product linear in the second variable. Put \(D_{ab}=D(E_{ab})\) for the matrix units, and let \(\mathcal E\) be the operator matrix with entries \(D_{ab}\). The trace-form Casimir on the tensor is the sum of the two individual Casimirs minus \(2\mathcal E\), since the dual defining action of \(E_{ab}\) is \(-E_{ba}\). The Casimir eigenvalues on \(\pi\) and \(\pi_s\) are \(\|\lambda\|^2-\|\rho\|^2\) and \(\|\lambda-e_s\|^2-\|\rho\|^2\), where \(\rho\) is the half-sum of the positive roots. The eigenvalue on the dual defining representation is \(n\). Thus \(\mathcal E\) acts on summand \(s\) by \(\lambda_s+(n-1)/2\). It is symmetric for \([\, ,\,]\), since \(D_{ab}^*=\epsilon_a\epsilon_b D_{ba}\) on smooth vectors. Its distinct real eigenspaces are orthogonal for this form. Schur's lemma on the irreducible compact-finite module, followed by continuity on smooth vectors, now gives
\[
 [y,p_s x]=c_s\langle R_s y,R_s x\rangle .
\]
Nondegeneracy gives \(c_s\ne0\).

We compute \(p_s\). For ordered row and column lists \(I,J\) of equal length \(m\), form the minor
\[
 C_{I,J}(u)=
 \sum_{\tau\in S_m}\operatorname{sgn}(\tau)
 \prod_{k=1}^{m}
 \bigl(D_{I_{\tau(k)},J_k}
       +(u-k+1)\delta_{I_{\tau(k)},J_k}\bigr),
\]
with factors in increasing order of \(k\) from left to right. Let \(C(u)\) be the full minor, and \(A(u)_{cl}=(-1)^{c+l}C_{\widehat l,\widehat c}(u)\), where a hat means deletion from the full list. The quantum comatrix identity is
\[
 A(u)\bigl((u-n+1)I+\mathcal E\bigr)=C(u)I .
\]
This is the evaluation \(T_{ab}(u)=\delta_{ab}+D_{ab}/u\) of the comatrix identity, after clearing the scalar denominators \cite[Propositions~2.2 and 2.27, equation~(2.30)]{molev-2003}. The evaluation applies because the \(D_{ab}\) satisfy the matrix-unit commutation relations.

Put \(m_j=\lambda_j-(n-1)/2\). The scalar of \(C(u)\) is \(\prod_j(u+m_j)\). Indeed the quantum determinant is central. On a highest-weight vector of weight \(\lambda-\rho\) its diagonal product is \(\prod_j(u-j+1+\lambda_j-\rho_j)\), and in every other product the rightmost displaced column supplies a raising operator. This computes its Harish--Chandra polynomial and hence its scalar at the infinitesimal character of \(\pi\). On summand \(s\), the comatrix identity gives
\[
 (u+m_s)A(u)=\prod_j(u+m_j)\,I .
\]
Cancel the polynomial factor before substituting \(u=-m_s\). On all other summands \(A(-m_s)=0\), and consequently
\[
 p_s=\frac{A(-m_s)}{\prod_{j\ne s}(\lambda_j-\lambda_s)}.
\]
The \(zz\)-entry of \(A(u)\) is the full Capelli determinant on the complementary coordinates. A simultaneous permutation of its rows and columns preserves the minor, by alternation in each list. The map \(P\) commutes with the complementary derived matrix units, by differentiation of its \(H\)-equivariance. The preceding central-character calculation in rank \(n-1\) therefore evaluates the numerator on \(\sigma\) as
\(\prod_a(-m_s+\mu_a-(n-2)/2)\). This gives the stated \(\gamma_s\). Each denominator factor has absolute value at least one and each numerator factor has absolute value at most \(2B+1/2\le X\).

To bound the other factor, let \(K=\U(p)\times\U(q)\), and let \(V_k\) be the span of words of at most \(k\) derived matrix units applied to the lowest \(K\)-type. It is a finite-dimensional \(K\)-module. Write \(\rho_c,\rho_n\) for the compact and noncompact half-sums. The lowest highest weight is \(\lambda+\rho_n-\rho_c\) by \ref{B1.15b}. In passing from \(V_k\) to \(V_{k+1}\) the possible new types occur in the tensor product with the matrix representation under conjugation. On \(K\) its blocks are a defining representation tensored with a dual defining representation. The one-box Pieri rule on each unitary factor adds one coordinate and subtracts one, with the dominant intermediate weights retained. Thus every highest weight \(\nu\) in \(V_k\) satisfies
\[
 \|\nu+\rho_c\|\le\|\lambda+\rho_n\|+\sqrt2\,k .
\]
Let \(\Omega_K=\sum_{\epsilon_a=\epsilon_b}D_{ab}D_{ba}\) and \(\Omega_G=\sum_{a,b}D_{ab}D_{ba}\). The adjoint formula gives, for \(x\in V_k\),
\[
 \sum_{a,b}\|D_{ab}x\|^2
 =\langle x,(2\Omega_K-\Omega_G)x\rangle .
\]
The compact Casimir on type \(\nu\) is \(\|\nu+\rho_c\|^2-\|\rho_c\|^2\). Each coordinate of \(\rho_n\) and of \(\rho\) has modulus at most \(n/2\). The opposite-sign gap gives \(B\ge500n^3\). For \(k\le S\) the last display is consequently bounded by
\[
 \bigl(2(\sqrt n(B+n/2)+\sqrt2\,S)^2+n(n/2)^2\bigr)\|x\|^2
 \le 64n(B+S)^2\|x\|^2 .
\]
Cauchy--Schwarz gives the corresponding bound on \(D(A)x\) for \(\|A\|_{\rm HS}\le1\). Since \(D(A)V_k\subseteq V_{k+1}\), iteration proves the word estimate.

At \(u=-m_s\) every scalar in a minor of length \(n-1\) has modulus at most \(4B\). In each permutation product choose the factors contributing a matrix unit. The word estimate with \(S=n-1\), summed over these choices, gives
\[
 \|(p_s(e_l v))_l\|
 \le (n-1)!\bigl(4B+8\sqrt n(B+n-1)\bigr)^{n-1}
 \le X^{4n}.
\]
For the last inequality use \((n-1)!\le X^{n-1}\) and \(4B+8\sqrt n(B+n-1)\le X^2\). Taking \(x=y=e_l v\) in the signed form identity proves the asserted bound on \(|c_s|\|R_s(e_l v)\|^2\). Finally use the formal-degree product of \ref{B1.15b}. The Haar normalization cancels in the ratio. Only the \(n-1\) differences involving \(s\) change, and each old absolute difference is at most twice the new one, since the new one is at least one. This gives \(d_\pi/d_{\pi_s}\le2^{n-1}\le X^n\).
\end{proof}

\subsection*{Proof of \ref{B8.15h}}\label{prf:B8.15h}\phantomsection
\begin{proof}
For either unitary group, of matrix size \(a\le n\), the set \(\|g\|_{\rm op}\le R\), \(R\ge1\), is covered by at most
\[
 (1+2\sqrt n R^2/\delta)^{2n^2}
\]
left translates of the Hilbert--Schmidt \(\delta\)-ball at the identity. Indeed, take a maximal \(\delta/R\)-separated subset in the ambient matrix space. The disjoint balls of half that radius about the selected points lie in the ball of radius \(\sqrt n R+\delta/(2R)\). The covering follows from \(\|g^{-1}\|_{\rm op}=\|g\|_{\rm op}\), since inversion is adjoint followed by conjugation by the signature matrix.

Put \(k=\lfloor\log\|h\|_{\rm op}\rfloor\). In Cartan coordinates the length \(d\) of Lemma~\ref{B8.53a} satisfies \(d\ge\sqrt2\,k\), since the extreme singular values are reciprocal. The row signature has smaller index \(q\), and its opposite-sign gap \(h_0\ge1000n^3\) exceeds \(120(n-1)q^{3/2}\). For \(k\ge1\), the local compression bound of that lemma gives
\[
 |\langle w,\sigma(h)w\rangle|\le \exp(-125n^{5/2}k).
\]
For \(k=0\) use unitarity. Apply the covering bound with \(R=e^{k+1}\), \(\delta=1\). After division by \(\operatorname{vol}(B_H(1))\), the integral of the absolute coefficient is at most
\[
 (1+2\sqrt n e^2)^{2n^2}
 \sum_{k\ge0}\exp\bigl((4n^2-125n^{5/2})k\bigr)
 \le 2(1+2\sqrt n e^2)^{2n^2}\le X^{67n^2}.
\]
The squared coefficient is bounded by the absolute coefficient and integrates to \(1/d_\sigma\) by \ref{B1.15b}.

A coset \(Hg\) is determined by \(g^{-1}z=(u,v)\) in the positive and negative coordinate spaces. Here \(\|u\|^2-\|v\|^2=1\). Compact rotations put \(u,v\) on coordinate axes. The hyperbolic rotation with \(\cosh t=\|u\|\), \(\sinh t=\|v\|\) then gives a representative of operator norm \(e^t\le2\|g^{-1}z\|\). In \(D\) this is at most \(3X^{4n}\), so its Hilbert--Schmidt norm is at most \(X^{10n}\). Left multiplication by \(B_H(1)\) keeps its operator norm at most \(6X^{4n}\). Quotient integration bounds \(\operatorname{vol}(D)\operatorname{vol}(B_H(1))\) by the volume of this last operator-norm ball in \(G\).

Take \(\delta=(64n(B+1))^{-1}\). If \(\|g-1\|_{\rm HS}<\delta\), the logarithm series gives \(g=\exp A\) with \(\|A\|_{\rm HS}\le2\delta\) and \(A\) in the real Lie algebra of \(G\). For a lowest-type unit vector \(v_0\) the length-one estimate in Lemma~\ref{B8.15g}, integrated on \(\exp(tA)\), gives \(\|\pi(g)v_0-v_0\|\le16\delta\sqrt n(B+1)\le1/2\). Schur orthogonality now gives \(d_\pi\operatorname{vol}(B_G(\delta))\le4\). In the covering bound for \(R=6X^{4n}\) we have \(1+2\sqrt n R^2/\delta\le X^{12n}\), using \(n\ge4\) and \(B\ge500n^3\). This proves the quotient bound.
\end{proof}

\subsection*{Proof of \ref{B8.15i}}\label{prf:B8.15i}\phantomsection
\begin{proof}
Put \(\xi=iw\), \(P=i^*\), and \(\phi(h)=\langle w,\sigma(h)w\rangle\). By Lemma~\ref{B8.15h}, \(f=d_\sigma\overline{\phi}\) belongs to \(L^1(H)\). In any unitary representation \(\tau\) of \(H\) set \(Q_\tau=\int_H f(h)\tau(h)\,dh\). Schur orthogonality gives \(f^*=f\) and
\[
 f*\delta_h*f=\phi(h)f ,
\]
where \(\delta_h\) is unit point mass. Indeed the convolution at \(x\), divided by \(d_\sigma^2\), is the integral of
\(\overline{\langle w,\sigma(y)w\rangle}\langle\sigma(x)w,\sigma(y)\sigma(h)w\rangle\).
Thus \(Q_\tau\) is an orthogonal projection and \(Q_\tau\tau(h)Q_\tau=\phi(h)Q_\tau\).

A unit vector in the image has coefficient \(\phi\). The Gram matrices of its \(H\)-translates are therefore the Gram matrices of the translates of \(w\), so its cyclic subrepresentation is an isometric copy of \(\sigma\). He's discrete branching theorem gives multiplicity one for the stated interlacing \cite[Theorem~0.1]{he-2017}. On that copy, Schur orthogonality gives the projection onto \(\C iw\). Hence \(Q_{\pi|_H}u=\xi\langle\xi,u\rangle\). Choose a \(G\)-smooth \(\eta\) with \(a=\langle\xi,\eta\rangle\ne0\), using density.

Write \(L(g)\) for the tensor action. It satisfies \((L(g)e_l v)_z=(g^{-1})_{lz}\pi(g)v\). In the \(s\)-th unitary summand of Lemma~\ref{B8.15g} put
\[
 b_s=Q_{\pi_s|_H}R_s(e_z\eta).
\]
Since \(H\) fixes the positive coordinate \(z\), insertion in that coordinate intertwines the \(H\)-actions. For every smooth tuple \(x\), integrate the signed form identity of Lemma~\ref{B8.15g} with \(y=e_z\eta\) and with \(x\) replaced by \(L(h)x\), against \(f(h)\,dh\). Each side is an integrable scalar coefficient. Self-adjointness of the two projections gives
\[
 \overline a\,\langle\xi,(p_s x)_z\rangle
     =c_s\langle b_s,R_s x\rangle .
\]
At \(x=e_z\eta\), the left side equals \(\gamma_s|a|^2\), by the corner identity of Lemma~\ref{B8.15g} with \(P=i^*\). The right side is \(c_s\|b_s\|^2\). Consequently, at \(x=L(g)e_l v_0\), Schur orthogonality in \(\pi_s\) gives
\[
 \int_G|\langle\xi,(p_s L(g)e_l v_0)_z\rangle|^2\,dg
 =\frac{|\gamma_s|\,|c_s|\,\|R_s(e_l v_0)\|^2}{d_{\pi_s}}
 \le\frac{X^{6n}}{d_\pi}.
\]
The factors in this inequality have exponents \(n,4n,n\), respectively, by Lemma~\ref{B8.15g}. The choice of \(\eta\) has canceled.

Sum the projections to recover \((g^{-1})_{lz}\langle\xi,\pi(g)v_0\rangle\). Cauchy--Schwarz for the sum of \(n\) projections costs \(n\); summing the integral over \(s\) and \(l\) costs \(n^2\). Since \(n\le X\) and \(6n+3\le8n\), this proves both inequalities.
\end{proof}

\subsection*{Proof of \ref{B8.15a}}\label{prf:B8.15a}\phantomsection
\begin{proof}
At \(\tau\) let \(\pi\) be the core representation and \(\sigma\) the conjugate row representation. In the combined decreasing list of the two interlaced grid rows with the inserted positive entry, adjacent signs agree exactly when the entries come from different lists. The rule persists after lowering one core entry by one. The within-list spacings are at least two and opposite-sign gaps at least \(1000n^3\). The row has both signs and at least as many positive signs as negative. Eventually \(r\ge3\). Both lists have absolute entries at most \(B=b^2\). Put \(X_0=8nB\); the sizes give
\begin{equation}\label{eq:p8c2-branch-weight}
 \log X_0\le6\log b.
\end{equation}

For \(H=\U(Y_\tau)\), He supplies the isometric discrete branching embedding \(i:\sigma\to\pi|_H\) for this interlacing \cite[Theorem~0.1]{he-2017}. Choose lowest-type unit vectors \(v_0,w\) in \(\pi,\sigma\). For formal degrees with the chosen Haar measures, the density \(|\langle iw,\pi(g)v_0\rangle|^2\) has mass \(1/\deg\pi\). The inserted positive coordinate is \(z_0\), and deleting it gives the row signature. The spacing, opposite-sign gaps and bounds on the two lists above give the hypotheses of Lemma~\ref{B8.15g} with \(X=X_0\). Their interlacing gives the remaining hypothesis of Lemma~\ref{B8.15i}. Its column moment is at most \(X_0^{8n}/\deg\pi\).
\leanref{PositivityTheory.\allowbreak{}UnitaryDecomposition.\allowbreak{}momentHeA\_\allowbreak{}of}{Development.\allowbreak{}WeilBound.\allowbreak{}RepresentationTheoryGL.\allowbreak{}HolomorphicDiscreteSeries.\allowbreak{}SchurEigenProjections.\allowbreak{}HolomorphicMoments.\allowbreak{}MomentHigherEstimateHolds}{149--179}

Set \(\mathcal D=\{g:\|g^{-1}z_0\|^2\le2X_0^{8n}\}\). Since \(H\) fixes \(z_0\), \((hg)^{-1}z_0=g^{-1}z_0\). Thus the column bound cuts out a subset of the left quotient, whose preimage retains at least half the mass. Integrating on each left \(H\)-fibre gives \(\|i^*\pi(g)v_0\|^2/\deg\sigma\). With compatible measures the bounds are
\begin{equation}\label{eq:p8c2-branch-volumes}
 \begin{aligned}
 (\deg\sigma)^{-1}&\le X_0^{67n^2}\operatorname{vol}(\operatorname{ball}_H(1)),\\
 (\deg\pi)\operatorname{vol}(H\backslash\mathcal D)
 &\le4X_0^{24n^3}/\operatorname{vol}(\operatorname{ball}_H(1)).
 \end{aligned}
\end{equation}
Both volume inequalities follow from Lemma~\ref{B8.15h}. The integral of \(\|i^*\pi(g)v_0\|^2\) over \(H\backslash\mathcal D\) is at least \((\deg\sigma)/(2\deg\pi)\) by the fibre calculation. At some point its value is at least half its average. Substituting the two volume bounds cancels \(\operatorname{vol}(B_H(1))\). The norm of the projection is constant on the left coset by equivariance and unitarity. Taking the representative given by Lemma~\ref{B8.15h} gives \(g_\tau\), \(v=\pi(g_\tau)v_0\), with
\begin{equation}\label{eq:p8c2-branch-projection}
 \|g_\tau\|_{\rm HS}\le X_0^{10n},\qquad
 1\le16X_0^{67n^2+24n^3}\|i^*v\|^2 .
\end{equation}
\leanref{UpperTheory.\allowbreak{}Lit.\allowbreak{}Branch.\allowbreak{}projLarge\_\allowbreak{}of}{Development.\allowbreak{}AutomorphicArithmetic.\allowbreak{}ThetaCorrespondence.\allowbreak{}FockModels.\allowbreak{}RealBranch}{140--195}
\leanref{PositivityTheory.\allowbreak{}RgHe.\allowbreak{}A.\allowbreak{}Pf.\allowbreak{}bigProjP\_\allowbreak{}of}{Development.\allowbreak{}WeilBound.\allowbreak{}RepresentationTheoryGL.\allowbreak{}HolomorphicDiscreteSeries.\allowbreak{}SchurEigenProjections.\allowbreak{}SchurFiniteDimensionality.\allowbreak{}IsFdOfSchurK}{53--275}

With inner product linear in the second argument set \(C(g)=\langle v,\pi(g)v\rangle\) and \(A=\int_H\overline{C(jh)}\sigma(h)\,dh\). The inversion law makes \(A\) self-adjoint, and Plancherel orthogonality on the discrete-series coefficient space makes it Hilbert--Schmidt. The embedded \(\sigma\) gives the rank-one contribution and the remaining coefficient is of positive type. Thus
\begin{equation}\label{eq:p8c2-branch-eigenlower}
 \langle a,Aa\rangle\ge|\langle i^*v,a\rangle|^2/\deg\sigma .
\end{equation}
The compact spectral theorem gives \(w'\ne0\) with eigenvalue \(c_\tau\ge\|i^*v\|^2/(2\deg\sigma)>0\). The weak coefficient integral
\[
 Tb=c_\tau^{-1}\int_H\langle\pi(jh)v,b\rangle\,\sigma(h)w'\,dh
\]
defines a bounded map by Schur orthogonality and intertwines by change of variables. It sends \(v\) to \(w'\), which is therefore smooth. Normalize and transport back the conjugate of \(w'\) to the row model, with its Haar factor, and take \(k_\tau(g)=C(g)\) in the transported frame.
\leanref{GlobalTheory.\allowbreak{}CompanionTheory.\allowbreak{}RgHeA.\allowbreak{}eigA}{Development.\allowbreak{}AutomorphicSpectrum.\allowbreak{}Proofs.\allowbreak{}BlockDecomposition.\allowbreak{}ExistsEigOfLow}{72--93}
\leanref{UpperTheory.\allowbreak{}Lit.\allowbreak{}Branch.\allowbreak{}act\_\allowbreak{}of\_\allowbreak{}eigen}{Development.\allowbreak{}AutomorphicArithmetic.\allowbreak{}ThetaCorrespondence.\allowbreak{}RealBands.\allowbreak{}BranchPack}{60--105}

For an absolute \(C_1\) the two real estimates are
\begin{equation}\label{eq:p8c2-branch-real}
 \begin{aligned}
 \int_H|k_\tau(gj(h))|\,dh
 &\le c_\tau X_0^{C_1n^4}\|gz_0\|^{-(2n+5)},\\
 \int_H|k_\tau(j(\exp(sX)h))-k_\tau(j(h))|\,dh
 &\le c_\tau X_0^{C_1n^4}|s|\|X\|_{\rm HS}.
 \end{aligned}
\end{equation}
For the first, integrate the lowest-type decay of Lemma~\ref{B8.53a} on the stabilizer in Cartan coordinates. The extra column gives power \(2r+7=2n+5\). For the second, integrate the first derivative along the translation. The lowest-type first-order Lie estimate supplies a factor bounded by a power of \(X_0\) times \(\|X\|_{\rm HS}\). The remaining negative norm power is integrable by the opposite-sign gap, with integral bounded by a power of \(X_0\) times the radius-\(1/4\) stabilizer ball volume. The formal-degree bound for that ball and \eqref{eq:p8c2-branch-projection}--\eqref{eq:p8c2-branch-eigenlower} bound its volume divided by \(c_\tau\) by a power of \(X_0\). Include the conjugation factors from \(\|g_\tau\|_{\rm HS}\) in \(C_1n^4\).
\leanref{UpperTheory.\allowbreak{}Lit.\allowbreak{}Branch.\allowbreak{}realInputE}{Development.\allowbreak{}AutomorphicArithmetic.\allowbreak{}ThetaCorrespondence.\allowbreak{}RealBands.\allowbreak{}BranchInput}{41--139}
\leanref{UpperTheory.\allowbreak{}Lit.\allowbreak{}realInputE\_\allowbreak{}holds}{Development.\allowbreak{}AutomorphicArithmetic.\allowbreak{}ThetaCorrespondence.\allowbreak{}FockModels.\allowbreak{}RealBranch}{231--236}

Tensor the row vectors. Haar changes multiply the coefficient bounds and eigenvalues by the same factors. Equation \eqref{eq:p8c2-branch-weight} gives the total loss. The pullback of the finite core level fixes the row vector, and its indicator contributes its positive volume to \(c\). The support congruence gives \(z-z_0\in\mathfrak m\Lambda_{j(y)}\). This proves \eqref{eq:p8c2-row-reproduce} and \eqref{eq:p8c2-row-coefficient}. Translate \eqref{eq:p8c2-row-reproduce} by \(\exp(sX)\), subtract, and use the second integral bound of \eqref{eq:p8c2-branch-real}.
\end{proof}

\subsection*{Proof of \ref{B8.15b}}\label{prf:B8.15b}\phantomsection
\begin{proof}
Use \(u=t(w_*)\) of Lemma~\ref{B8.15a} and \(M=\sup|u|\), finite by cuspidal decay \cite{harish-chandra-1968} and positive by cyclicity. Choose \(x_*\) with \(|u(x_*)|>M/2\). Put \(N=\operatorname N\mathfrak m,\ g=g_0g_1,\ Y_*=r^2fg\operatorname{rd}(E)\), and \(s=(2e^{L_{\rm loss}}Y_*)^{-m}/2\). Then
\begin{equation}\label{eq:p8c2-depth-pays}
 0<s\le\tfrac12,\qquad e^{L_{\rm loss}}Y_*(2s)^{1/m}=\tfrac12,\qquad
 R:=s^f\sqrt N\ge2e^{L_{\rm loss}+(6n+9)f}>1 .
\end{equation}
Indeed the first auxiliary place gives \(\log N>C_Ff^2n^6\log b\). Use \(m\le f\le k,\ n=k+m\) and \eqref{eq:p8c2-size-rows-a}--\eqref{eq:p8c2-size-rows-c} to get \(\log Y_*\le n^3+21\log b\). For \(n\ge C+1,\ L\ge1\) the needed logarithmic bound is
\begin{equation}\label{eq:p8c2-depth-log}
 \log2+L_{\rm loss}+(6n+9)f
 +f\{\log2+m(\log2+L_{\rm loss}+\log Y_*)\}
 \le(2C+42)f^2n^5\log b\le\tfrac12\log N .
\end{equation}
\leanref{GlobalTheory.\allowbreak{}SepPer.\allowbreak{}depth\_\allowbreak{}pays}{Development.\allowbreak{}AbelianVarietyTheory.\allowbreak{}Proofs.\allowbreak{}TwistedOperators.\allowbreak{}GridDepthPaysEveryLoss}{148--277}

Let \(\Lambda_{Y,x_*}\) have finite completion \((x_*)_f\widehat\Lambda_Y\). For every integral ideal \(\mathfrak a\ne0\) and nonzero \(v\in\mathfrak a\Lambda_{Y,x_*}\), one has
\begin{equation}\label{eq:p8c2-deep-point}
 \prod_\tau\|(x_*)_\tau^{-1}v_\tau\|\ge s^f\sqrt{\operatorname N\mathfrak a}.
\end{equation}
Otherwise the least lattice scale \(\lambda<s\). The short-vector flag between \(\lambda\) and \(1/2\) has at most \(m\) active layers by the smaller signature index. The product of consecutive ratios is \(2\lambda\), so one adjacent lower-to-upper ratio is at most \((2\lambda)^{1/m}\). Its saturated lower lattice spans a nonzero isotropic rational \(W\).

For \(u_0\) in the adelic unipotent radical of \(W\), use unipotent strong approximation to remove its finite part up to the conjugated principal congruence subgroup of level \(g\). In the resulting real logarithm first reduce the image in the first layer by its lattice at level \(2g\). In the two-step exponential product the commutator with the rational first-layer correction is central. Reduce that central logarithm by the centre lattice at level \(g\). The first covering radius uses the adjacent flag ratio, and the centre radius uses the square of the lower scale. These two reductions give
\[
 u_0x_*=\gamma x_*(u_\infty,k_f),\qquad
 u_{\infty,\tau}=\exp A_\tau\exp Z_\tau
\]
with rational \(\gamma\), principal congruence \(k_f\) of level \(g\), and
\begin{equation}\label{eq:p8c2-short-line}
 \begin{aligned}
 \sqrt{\sum_\tau\|A_\tau\|_{\rm HS}^2}+\sqrt{\sum_\tau\|Z_\tau\|_{\rm HS}^2}
 &\le\tfrac14r^2\sqrt f\,g\operatorname{rd}(E)(2\lambda)^{1/m},\\
 \sum_\tau(\|A_\tau\|_{\rm HS}+\|Z_\tau\|_{\rm HS})
 &\le r^2fg\operatorname{rd}(E)(2s)^{1/m}.
 \end{aligned}
\end{equation}
The second inequality is Cauchy--Schwarz. Subtract the zero cuspidal average over the radical from \(u(x_*)\). Rational and level invariance reduce its integrand to the successive real translates. Equation \eqref{eq:p8c2-row-variation} with \eqref{eq:p8c2-depth-pays} bounds the difference by \(M/2\), a contradiction. Adjoining \(z_0\) preserves the product lower bound for \(\mathfrak a\Lambda_{j(x_*)}\). If the extra coordinate is zero, apply \eqref{eq:p8c2-deep-point}. A nonzero extra coordinate has product norm at least \(\sqrt{\operatorname N\mathfrak a}\).
\leanref{GlobalTheory.\allowbreak{}DeepPoint.\allowbreak{}deep\_\allowbreak{}point}{Development.\allowbreak{}AbelianVarietyTheory.\allowbreak{}Proofs.\allowbreak{}TwistedOperators.\allowbreak{}AutomorphicPointExistsIsotropicReductionOfSmallS}{141--196}
\leanref{GlobalTheory.\allowbreak{}SmallFund.\allowbreak{}small\_\allowbreak{}fundamental\_\allowbreak{}set}{Development.\allowbreak{}AbelianVarietyTheory.\allowbreak{}Proofs.\allowbreak{}FlagMeasures.\allowbreak{}SmallFundamentalSetRealUnipotent}{135--207}
\leanref{GlobalTheory.\allowbreak{}ShortLineCusp.\allowbreak{}norm\_\allowbreak{}le\_\allowbreak{}of\_\allowbreak{}not\_\allowbreak{}long}{Development.\allowbreak{}AbelianVarietyTheory.\allowbreak{}Proofs.\allowbreak{}TwistedOperators.\allowbreak{}NormCuspFormLeOfShortLine}{67--141}

Form the core Poincar\'e series from \(k\) translated by \(j(x_*)^{-1}\). In the unfolded period the rational cosets are indexed by \(z=\gamma z_0\). The identity term is \(cu(x_*)\). Other contributing terms have \(v=z-z_0\ne0\) in \(\mathfrak m\Lambda_{j(x_*)}\). Use \(N_\tau\) of Lemma~\ref{B8.15a} at \(x_*\). Then \(\prod_\tau N_\tau(v)\ge R\) and \(1+N_\tau(v)\le 3N_\tau(z)\), since \(z\) has Hermitian norm one. The tail of absolute integrals is at most
\begin{equation}\label{eq:p8c2-period-tail}
 M|c|e^{L_{\rm loss}}\sum_v\prod_\tau\left(\frac3{1+N_\tau(v)}\right)^{2n+5}
 \le M|c|e^{L_{\rm loss}+(6n+9)f}/R .
\end{equation}
Extract \(R^{-1}\) from one power of the denominator product. The products of radius-\(1/2\) Euclidean balls centered at the \(v\)'s have disjoint interiors by the lower bound on differences. Comparing the remaining weight inside a ball costs \((3/2)^{2n+4}\) per place. In dimension \(2n\), radial integration using \(t^{2n-1}(1+t)^{-2n-4}\le(1+t)^{-5}\) costs at most \(n4^n/2\) in units of ball volume. Thus the local cost is at most \(3^{2n+5}(3/2)^{2n+4}n4^n/2\le e^{6n+9}\). By \eqref{eq:p8c2-depth-pays}, the sum of absolute values of the nonidentity terms divided by \(M|c|\) is at most \(e^{L_{\rm loss}+(6n+9)f}/R\le1/2\). The identity term divided by the same positive number has modulus \(|u(x_*)|/M>1/2\). The unfolded period therefore converges absolutely and is nonzero.
\leanref{GlobalTheory.\allowbreak{}ProdNormSum.\allowbreak{}lattice\_\allowbreak{}sum\_\allowbreak{}le}{Development.\allowbreak{}AbelianVarietyTheory.\allowbreak{}Proofs.\allowbreak{}TwistedOperators.\allowbreak{}LatticeSumProdNormLeOfProdNormLowerBound}{225--279}
\leanref{GlobalTheory.\allowbreak{}IdDom.\allowbreak{}identity\_\allowbreak{}dominates}{Development.\allowbreak{}AbelianVarietyTheory.\allowbreak{}Proofs.\allowbreak{}TwistedOperators.\allowbreak{}PoincareSeriesIdentityTermDominates}{40--242}

The mixed coefficient satisfies
\begin{equation}\label{eq:p8c2-core-mixed-decay}
 |\langle v_0,\pi(g)a\rangle|
 \ll_a(1+\|g\|_{\rm HS})^{-(2n^2+1)}
\end{equation}
for \(K\)-finite \(a\). For \(a=v_0\) this follows from Lemma~\ref{B8.53a} and the gap. Every translate of \(v_0\) in the second slot satisfies the same decay with a translate-dependent constant, by the matrix norm inequality. Their span is dense by irreducibility. For \(u\) in that span average the positive rank-one operator of \(\pi(k)u\) over \(K\). Its image satisfies the mixed bound since the translate constants are bounded on \(K\). The sum \(Y\) of the isotypic spaces meeting a given \(K\)-finite \(a\) is finite dimensional by admissibility. The averages commute with \(K\) and preserve \(Y\). A vector in \(Y\) orthogonal to the sum of their images on \(Y\) satisfies \(\int_K|\langle\pi(k)u,z\rangle|^2\,dk=0\) for all such \(u\), hence \(z=0\) by continuity and density. Their images therefore span \(Y\).
\leanref{GlobalTheory.\allowbreak{}SepPer.\allowbreak{}KDecay}{Development.\allowbreak{}AbelianVarietyTheory.\allowbreak{}Definitions.\allowbreak{}TwistedInnerProducts.\allowbreak{}DepthPays}{58--62}
\leanref{GlobalTheory.\allowbreak{}KFinDecay.\allowbreak{}decayAt\_\allowbreak{}of\_\allowbreak{}kFin}{Development.\allowbreak{}AbelianVarietyTheory.\allowbreak{}Proofs.\allowbreak{}TwistedOperators.\allowbreak{}CoefficientDecayKFiniteSecondVector}{204--321}

The mixed bound gives convergence on \(K\)-finite tensors. Unfolding the proper constant terms gives the zero proper unipotent integrals of the discrete-series coefficients. The resulting cuspidal intertwiner at core level gives a smooth multiplicity map \(y\) by globalization. In \eqref{eq:p8c2-period-tail} the core vectors are \(\pi(g_\tau)v_0\) in both slots. Move the first translations to the base of the series and use the translated second vectors as the input \(a'\) of \(y\). Thus \(\operatorname{Per}_0(y(a'),t(w_*))\ne0\). For prescribed \(w\ne0\), cyclicity and continuity give a real translate of \(t(w)\) pairing nontrivially. Move it to the core side by right invariance. The period is continuous in the smooth core vector by boundedness of the restricted core function and rapid decay of the row form. Approximation by compact-type projections in the smooth topologies gives the asserted \(a\).
\leanref{GlobalTheory.\allowbreak{}PeriodNePf.\allowbreak{}exists\_\allowbreak{}kf}{Development.\allowbreak{}AbelianVarietyTheory.\allowbreak{}Proofs.\allowbreak{}TwistedOperators.\allowbreak{}PeriodNonvanishingOfArchimedeanInputs}{65--111}
\end{proof}

\subsection*{Proof of \ref{B8.15j}}\label{prf:B8.15j}\phantomsection
\begin{proof}
In the Schr\"odinger model the vacuum is the characteristic function of the \(a\)-fold product of the target lattice. The two compact stabilizers fix it. The target convolution \(\omega(1,h')\phi_0\) is fixed by both stabilizers. By Lemma~\ref{B8.14} it is a finite linear combination of source translates of \(\phi_0\). The extension is unramified of odd residue characteristic, both lattices are unimodular, and the integral translations fix the vacuum and its Heisenberg translates span the lattice model. These are the hypotheses of that lemma. Average the combination over the source stabilizer. Since \(\phi_0\) is also fixed on the right, the averaged translates are convolutions by normalized double-coset indicators. This gives \(h\).

For the transform use Liu's spherical-module theorem \cite[Definition~2.8 and Theorem~1.1]{liu-2022}. In his notation \(r=a\), the target Witt index is \(s=b\), and \(\epsilon(V)=+\) because the target lattice is self-dual. His Weil representation uses the conductor and trivial splitting characters specified here. Part~(2) says that the vacuum generates the spherical module when \(Q\) is odd. Thus \(h\otimes1-1\otimes h'\), which kills the vacuum, belongs to the annihilator of that module. Part~(1) identifies the annihilator with the kernel of the two Satake substitutions of Definition~2.8. The source substitution inverts the \(a\) variables, which fixes every spherical polynomial by Weyl invariance. The target substitution inserts \(Q^{2j-1}\) for \(1\le j\le b-a\) and retains the \(a\) variables. Applying these substitutions to the difference gives the displayed identity.
\end{proof}

\subsection*{Proof of \ref{B8.15c}}\label{prf:B8.15c}\phantomsection
\begin{proof}
Keep the chain notation above. The core \(V_i\) has rank \(n_i\) and rational Witt index less than \(d\). The characters \(\widetilde\chi_i\) are unramified off \(S\), trivial on \(\mathbb A_F^\times\), and have infinite exponents the negatives of the distinct grid centers. Reframing by \((\overline{\mathfrak A}\mathfrak A^{-1})^i\), for the core lattice-change ideal \(\mathfrak A\), preserves spherical stabilizers off \(S\) and the principal congruence levels \(J_i\) on \(S\).
\leanref{GlobalTheory.\allowbreak{}CoreOnto.\allowbreak{}comOK}{Development.\allowbreak{}AbelianVarietyTheory.\allowbreak{}Proofs.\allowbreak{}TwistedOperators.\allowbreak{}TestGridCommonDataPropertiesHold}{31--82}

Let \(C_i\) be the finite-dimensional cuspidal multiplicity space of the prescribed smooth discrete series at level \(J_i\) and trivial level character. Finiteness follows at fixed infinitesimal character and lowest type \cite{harish-chandra-1968}. Petersson products factor into its multiplicity pairing and the invariant real pairing. The spherical Hecke algebra off \(S\) is commutative and closed under adjoints, so \(C_0\) is spanned by simultaneous eigenvectors. Label their characters by the conjugate multisets \(B(w)\).

Choose finite tests before choosing an eigenvector. At depth places use the coset vacuum for a split lattice projection surjective modulo depth in the downward direction. Elsewhere take the standard vacuum. Both levels fix the test. For a source element integral off \(S\), equality of the Heisenberg characters on a nonzero self-coefficient forces preservation of the coset. Surjectivity of the lattice projection forces the source congruence. Thus on this domain the self-coefficient is supported on \(\U(V_i)(F_\infty)J_i\). Exchange of the adjoint lattice models makes each polynomial with this finite vacuum an allowed downward depth test.

Write \(A_i:C_{i+1}\to C_i\) for the downward multiplicity map and \(I_i(P)\) for the real intertwiner using polynomial input \(P\). The downward integral of \(y'(v')\) is \((A_i y')(I_i(P)v')\). For a nonzero eigenvector \(u\) reached from the bottom, \(A_i^*u\ne0\), with parameters given by spherical transfer. The induction gives
\begin{equation}\label{eq:p8c2-chain-satake}
 B_i(w)=B(w)\sqcup\bigsqcup_{j<i}
 \{\widetilde\chi_j(w)(Nw)^{(d-1)/2-l}:0\le l<d\}.
\end{equation}
The doubled numerator for the conjugate value is therefore
\begin{equation}\label{eq:p8c2-chain-numerator}
 H_i(s)=L^S(B\otimes\widetilde\chi_i^{-1},s+1/2)
 \prod_{j<i}\prod_{l=0}^{d-1}
 L^S(\widetilde\chi_j\widetilde\chi_i^{-1},s+d/2-l).
\end{equation}
Here \(L^S(B\otimes\widetilde\chi_i^{-1},z)\) has roots \(\beta\widetilde\chi_i(w)^{-1}\) and variable \((Nw)^{-z}\) off \(S\). By \eqref{eq:p8c2-bottom-satake} this first factor converges absolutely without zeros for \(\Re s>n/2\), including \(s_0=d/2\). The other arguments at \(s_0\) are \(d-l\ge1\) and their ratio characters have nontrivial angular type. Lemma~\ref{lem:p8c2-hecke-nonvanishing} gives nonvanishing. The normalizer \(D_i^S(s)=\prod_{l=0}^{n_i-1}L_F^S(2s+n_i-l,\eta^{n_{i+1}+l})\) has arguments greater than one there.
\leanref{GlobalTheory.\allowbreak{}EvenOnto.\allowbreak{}chainL\_\allowbreak{}and\_\allowbreak{}inj}{Development.\allowbreak{}AbelianVarietyTheory.\allowbreak{}Proofs.\allowbreak{}TwistedOperators.\allowbreak{}EvenThetaChainLValuesNeZeroAndParamsInjective}{95--174}

In the real source, after duality and removal of the common shift, write the signed lists as \((AB,CD)\), with \(A,C\) positive and \(B,D\) negative. The real theta rule sends this parameter to
\[
 (A[(d-1)/2,\ldots,-(d-1)/2]D,\ CB).
\]
The grid gives old entries of modulus at least \((d+3)/2\) and
\begin{equation}\label{eq:p8c2-chain-real-gap}
 \operatorname{gap}_{\rm opp}>n_{i+1}^3\quad\text{in both models},\qquad
 \operatorname{gap}_{\rm opp}\ge120n_i\min(p_i,q_i)^{3/2}
 \quad\text{in the source of signature }(p_i,q_i).
\end{equation}
The joint highest Fock line maps nontrivially to the two lowest types. Take a polynomial \(P_\tau\) on this line and a unit highest vector \(v_\tau\) of the source lowest type. Schur orthogonality integrates the Weil coefficient against the conjugate highest-vector coefficient to the Schur projector \(E_{v_\tau}\) divided by the source formal degree. The local integral is \(\|E_{v_\tau}P_\tau\|^2/\deg(\pi_\tau)>0\). Convergence follows from the gap.
\leanref{GlobalTheory.\allowbreak{}EvenOnto.\allowbreak{}upPair\_\allowbreak{}arch}{Development.\allowbreak{}AbelianVarietyTheory.\allowbreak{}Proofs.\allowbreak{}TwistedOperators.\allowbreak{}RealThetaCorrespondenceUpwardPairBothDirections}{60--76}
\leanref{GlobalTheory.\allowbreak{}ArchRec.\allowbreak{}lift\_\allowbreak{}down}{Development.\allowbreak{}AbelianVarietyTheory.\allowbreak{}Proofs.\allowbreak{}TwistedOperators.\allowbreak{}DiscreteSeriesHarishChandraModuleAdmissible}{290--314}

Tensor these polynomials with the finite test and take \(F=u(\bigotimes_\tau v_\tau)\). On the real group times \(J_i\) its Petersson coefficient is \(\|u\|^2\) times the real coefficients. The semilocal integral \(Z_{S,\infty}\) at \(s_0\) on the domain integral off \(S\) is positive by the support calculation and the real projectors. For its continuation, real Gaussian integration in Cartan coordinates bounds the Weil self-coefficient \(c_P(g)\) in terms of the doubled height \(0<h(g)\le1\):
\begin{equation}\label{eq:p8c2-chain-zeta-bound}
 |c_P(g)|\le C h(g)^{n_{i+1}/2},\qquad
 |c_P(g)|h(g)^{\Re s-s_0}\le C\quad(\Re s\ge0).
\end{equation}
The second uses \(n_{i+1}\ge d\). The integrable highest-vector discrete-series coefficient, by \eqref{eq:p8c2-chain-real-gap}, supplies the majorant.
\leanref{GlobalTheory.\allowbreak{}EvenOntoRealZeta.\allowbreak{}realZeta}{Development.\allowbreak{}AbelianVarietyTheory.\allowbreak{}Proofs.\allowbreak{}TwistedOperators.\allowbreak{}EvenThetaStepRealZetaNonvanishing}{40--134}
\leanref{GlobalTheory.\allowbreak{}EvenOnto.\allowbreak{}ZetaNeProof.\allowbreak{}zetaNe\_\allowbreak{}of\_\allowbreak{}real}{Development.\allowbreak{}AbelianVarietyTheory.\allowbreak{}Proofs.\allowbreak{}TwistedOperators.\allowbreak{}EvenThetaStepSemilocalZetaNeZeroOfReal}{63--116}

By the split and inert unramified local doubling input in \(\mathscr P\) \cite[\S4.2.1]{eischen-harris-li-skinner-2020}\cite[p.~334, Remark~3]{lapid-rallis-2005}, the normalized local doubling integrals off \(S\) multiply to \(H_i(s)/D_i^S(s)\) in the convergence half-plane. At \(s_0=d/2\), for every \(i\ge1\) we have \(2n_i-n_{i+1}=(r+1)+(i-1)(r+5)>0\), so \(n_{i+1}<2n_i\) and the first formula of Theorem~\ref{B8.13} applies since the target Witt index is less than \(d\). At \(i=0\), we have \(n_1=2n_0+4>2n_0\), and the absolutely convergent unitary Siegel--Weil formula for \(b>2a\) in \(\mathscr P\) \cite{weil-1965} applies with covolume one. Write \((\Theta_P^{\rm up}F)(h)=\int_{[\U(V_i)]}\overline{\theta_i(g,h;P)}F(g)\,dg\). Factoring the doubled pairing on \(\overline F\) gives
\begin{equation}\label{eq:p8c2-chain-norm}
 \|\Theta_P^{\rm up}F\|^2
 =\operatorname{vol}[\U(V_{i+1})]\frac{H_i(s_0)}{D_i^S(s_0)}Z_{S,\infty}\ne0 .
\end{equation}
\leanref{GlobalTheory.\allowbreak{}Core.\allowbreak{}Rk.\allowbreak{}IP.\allowbreak{}siegelWeilAll\_\allowbreak{}of\_\allowbreak{}hypotheses}{Development.\allowbreak{}AbelianVarietyTheory.\allowbreak{}Proofs.\allowbreak{}DecoratedLatticeProofs.\allowbreak{}SiegelWeilFockDoubledPairOfWeilAndIchino}{49--85}
\leanref{GlobalTheory.\allowbreak{}EvenOnto.\allowbreak{}upNe\_\allowbreak{}of\_\allowbreak{}heckeForm}{Development.\allowbreak{}AbelianVarietyTheory.\allowbreak{}Proofs.\allowbreak{}TwistedOperators.\allowbreak{}UpwardThetaLiftNeZeroOfLValueNeZero}{33--140}

The real parameter rule kills the lower Witt-tower lifts on Fock tests because the old entries avoid \(\pm(d+1)/2\). Density and continuity extend the vanishing to the allowed tests. The proper constant terms vanish because they expand through those lower lifts. Moving \(J_{i+1}\) onto \(P\) gives the new level. The real type puts the upward lift in the space of values of \(C_{i+1}\). By Fubini its pairing against \(y'(v')\) is the downward pairing \(\langle A_i y',u\rangle\langle I_i(P)v',v\rangle\), in the argument order of the downward map. Nondegeneracy and the adjoint relation identify the lift as a value of \(A_i^*u\). Equation \eqref{eq:p8c2-chain-norm} gives \(A_i^*u\ne0\), and the spherical transfer gives \eqref{eq:p8c2-chain-satake} at \(i+1\).
\leanref{GlobalTheory.\allowbreak{}EvenOnto.\allowbreak{}LiftUp.\allowbreak{}liftUp\_\allowbreak{}stage\_\allowbreak{}span}{Development.\allowbreak{}AbelianVarietyTheory.\allowbreak{}Proofs.\allowbreak{}TwistedOperators.\allowbreak{}LiftUpMultiplicitySpacesSpanKFiniteVectors}{43--91}
\leanref{GlobalTheory.\allowbreak{}EvenOnto.\allowbreak{}stepUp}{Development.\allowbreak{}AbelianVarietyTheory.\allowbreak{}Proofs.\allowbreak{}TwistedOperators.\allowbreak{}EvenThetaStepAdjointNonzeroTransferredParams}{38--57}

Let \(A=A_0\cdots A_{m-1}\). A nonzero vector orthogonal to its image decomposes into simultaneous nonzero Hecke components \(u_B\). Each \(A^*u_B\ne0\), and adding the same segments preserves distinct multisets. Thus these images are linearly independent, contradicting \(A^*u=0\). The argument works with arbitrary multiplicities.
\leanref{GlobalTheory.\allowbreak{}EvenOnto.\allowbreak{}ontoAll\_\allowbreak{}of}{Development.\allowbreak{}AbelianVarietyTheory.\allowbreak{}Proofs.\allowbreak{}TwistedOperators.\allowbreak{}EvenRankThetaChainLiftSurjectiveOfAdjointNonzero}{97--123}

Take a top preimage of \(y\). The span of \(I_i(P)v'\) for \(K\)-finite \(v'\) contains a nonzero joint-type vector and is Lie-stable by differentiating the Fock polynomial. Irreducibility gives the full \(K\)-finite tensor space. Express \(a\) as a finite sum of these values and repeat up the chain. The downward integrals give the allowed tests at every rung. At the top the perfect pairing with the coefficient representation gives each value as \(\phi_{\xi,\Psi}\) with a column of the prescribed level.
\leanref{GlobalTheory.\allowbreak{}ItOnto.\allowbreak{}kfinSpan\_\allowbreak{}le\_\allowbreak{}span\_\allowbreak{}it}{Development.\allowbreak{}AbelianVarietyTheory.\allowbreak{}Proofs.\allowbreak{}DecoratedLatticeProofs.\allowbreak{}KFiniteSpanLeSpanRealIntertwinerValues}{44--101}
\leanref{GlobalTheory.\allowbreak{}ChainLifts.\allowbreak{}chainEx}{Development.\allowbreak{}AbelianVarietyTheory.\allowbreak{}Proofs.\allowbreak{}TwistedOperators.\allowbreak{}CoreChainCompositesExistLiftCombinations}{36--154}
\leanref{GlobalTheory.\allowbreak{}CoreOnto.\allowbreak{}topCols}{Development.\allowbreak{}AbelianVarietyTheory.\allowbreak{}Proofs.\allowbreak{}TwistedOperators.\allowbreak{}DefiniteTopSpaceColumnValues}{65--72}
\end{proof}

\subsection*{Proof of \ref{B8.16d}}\label{prf:B8.16d}\phantomsection
\begin{proof}
Since \(d=u+v\) is even, \(v+2(2(p+q)+1)<u\), and \(2(p+q)+1<v\), the dimensions \((a,a')=(u,v)\) satisfy \(a'<a\), \(a+a'\equiv0\pmod2\), \(2(p+q)+1<(a-a')/2\), and \(2(p+q)+1<a'\). This gives \eqref{eq:p8c2-triple-range}. Fix a bijection \(e:\{1,\ldots,d\}\simeq\{1,\ldots,u+v\}\). Keeping \(J_V=J_W\) preserves the signature \((u-p,p)\), while setting \(J_{V'}=-J_{W_1}\) reverses \((v-q,q)\) to \((q,v-q)\) and matches the sign \(-\) of \(\omega'\). For \(\bar h^t J_X h=J_X\), taking the transpose gives \(h^t J_X^t \bar h=J_X^t\); multiplication by \(\xi_0\) with \(\bar\xi_0=-\xi_0\) makes \(J_{W^{\rm skew}}=\xi_0(J_X^t)^e\) skew-Hermitian and makes \(\kappa_X(h)=(\bar h)^e\) an isomorphism of algebraic groups \(\U(X)\xrightarrow{\sim}\U(W^{\rm skew})\) over \(F\), whence \(\operatorname{ht}(\kappa_X(h))\le C'H_X(h)^{N'}\) for suitable \(C'>0\) and \(N'\in\N\). At each infinite place \(w\) of \(F\) with chosen embedding \(\tau\in\theta\), Hermitian symmetry of \(J_X\) gives \(\bar\tau(J_X^t)=\tau(J_X)\), while \(\bar\xi_0=-\xi_0\) and \(\operatorname{Im}\tau(\xi_0)>0\) give \(i\bar\tau(\xi_0)=\operatorname{Im}\tau(\xi_0)>0\). Hence \(i\bar\tau(J_{W^{\rm skew}})\) is a positive scalar multiple of \((\tau(J_X))^e\) and has signature \((d-p-q,p+q)\). Under \(\bar\tau\) the angular exponents \(2a_\tau,2b_\tau\in\Z\) of \(\alpha,\beta\) become \(m_{\alpha,w}=-2a_\tau\) and \(m_{\beta,w}=-2b_\tau\), so \(2(a_\tau-b_\tau)<-d\) becomes \(m_{\alpha,w}-m_{\beta,w}=2(b_\tau-a_\tau)>d\).

By Definition~\ref{B8.16c}, \(\psi\) is continuous on \(\mathbb A_F/F\) with infinite part \(y\mapsto\exp(-2\pi i\operatorname{Tr}_{F_\infty/\R}y)\). Let \(\psi_{\Q,\rm std}\) be the standard character of \(\mathbb A_\Q/\Q\) with archimedean part \(t\mapsto e^{-2\pi i t}\). Then \(\psi\cdot(\psi_{\Q,\rm std}\circ\operatorname{Tr}_{F/\Q})^{-1}\) is continuous and trivial on \(F+F_\infty\), which is dense in \(\mathbb A_F\) by strong approximation for the additive group of \(F\). Thus \(\psi=\psi_{\Q,\rm std}\circ\operatorname{Tr}_{F/\Q}\), and at each finite place \(v_0\) of \(F\) its conductor is \(\mathfrak D_{F/\Q,v_0}^{-1}\), equal to \(\mathcal O_{v_0}\) whenever \(v_0\nmid\mathfrak D_{F/\Q}\). Let \(P=T\) (with \(S_0=\varnothing\) and \(Q=T\)) be the finite set of finite places of \(F\) where \(\alpha\) or \(\beta\) ramifies. Every \(v_0\in T\) splits in \(E/F\) and satisfies \(v_0\nmid\mathfrak D_{F/\Q}\), so \(\psi_{v_0}\) has conductor \(\mathcal O_{v_0}\) on \(P\). Since \(d=u+v\) is even, \((\alpha/\beta)|_{\mathbb A_F^\times}=\eta^{v-u}=1\), whence \((\alpha/\beta)(\bar t)=(\alpha/\beta)(t)^{-1}\) for \(t\in\mathbb A_E^\times\); hence ramification of \(\alpha/\beta\) at one split component above \(v_0\in T\) implies ramification at both split places of \(E\) above \(v_0\). In the arithmetic point, \(E/\Q\) is Galois with \(E/F\) unramified at finite places, and every inert place \(v_0\) of \(E/F\) is unramified over \(\Q\) of odd residue characteristic with \(Nv_0>3\) (since \(2\), the ramified primes of \(F/\Q\), and odd primes up to the host split cutoff at least \(11\) split in \(E/F\)), so \(\psi_{v_0}\) also has conductor \(\mathcal O_{v_0}\) at inert places. Moreover, the Cartan involutions and self-dual lattices of \((W,W_1,X)\) transfer to \((V,V',W^{\rm skew})\) because \(J_V=J_W\), \(J_{V'}=-J_{W_1}\), and at every inert place \(v_0\) both \(\xi_0\) and the absolute codifferent are units. Choosing split embeddings on \(P\) and a place \(v_1\notin T\) above a rational prime split completely in \(E\) supplies good base data preceding Theorem~\ref{B8.11} with the same additive character \(\psi\).

In the frames \((A,B)\) and \((A_1,B)\), put \(c_A(m)=P_e^{-1}(\bar m)^t J_W\) and \(c_B(m')=-P_e^{-1}(\bar m')^t J_{W_1}\), where \(P_e\) is the permutation matrix \((P_e)_{ij}=\delta_{i,e(j)}\) (so \(P_e^{-1}(\bar m)^t\) reindexes only the row index of \((\bar m)^t\) by \(e^{-1}\)). Since \(\bar g^t J_W=J_W g^{-1}\), \(\bar g_1^t J_{W_1}=J_{W_1}g_1^{-1}\), and \(P_e^{-1}\overline{\kappa_X(h)}=hP_e^{-1}\), these are \(\mathbb A_F\)-linear topological isomorphisms preserving the rational subspaces, intertwining \(m\mapsto\kappa_W(g)m\kappa_X(h)^t\) with \(\phi\mapsto h\phi g^{-1}\) (and similarly for \(c_B\)), and carrying the symplectic trace pairings of \(V\otimes_E W^{\rm skew}\) and \(V'\otimes_E W^{\rm skew}\) to the oscillator alternating pairings of \((W,X)\) and \((-W_1,X)\). For polynomial Schwartz vectors \(\Phi_D,\Phi'_D\), the sums defining \(\mathcal W_A(\Phi_D)\) and \(\mathcal W_B(\Phi'_D)\) converge absolutely, satisfy rational quasiperiodicity because rational Darboux shifts permute the summation lattice with trivial \(\psi\)-phase, and intertwine \(\rho_\psi,\rho'_\psi\) with Heisenberg translations along \(c_A,c_B\). Fourier inversion on the compact Pontryagin torus \((F\backslash\mathbb A_F)^{u\times(u+v)}\) (and its \(v\times(u+v)\) counterpart) in the complementary Darboux coordinate recovers \(\Phi_D,\Phi'_D\) and proves injectivity. Because \(c_A,c_B\) conjugate \(E\)-coordinates, they carry the complexified polarization at \(\bar\theta\) to the Fock polarization at \(\theta\), map finite Schwartz--Bruhat factors and archimedean Gaussian polynomials onto the permitted finite theta functions and real Fock polynomials, and, since \(\det\kappa_X(h)=\overline{\det h}=(\det h)^{-1}\) matches the opposite target determinant convention of the oscillator model while \(\kappa_W,\kappa_{W_1}\) preserve source determinants and \(c_A,c_B\) conjugate \(E\)-coordinates, the splitting pairs \((\beta,\beta/\alpha),(\alpha,\alpha/\beta)\) agree on both sides so that \((\mathcal W_A,\mathcal W_B)\) intertwine \((\tilde\omega(\kappa_W(g),\kappa_X(h)),\tilde\omega'(\kappa_{W_1}(g_1),\kappa_X(h)))\) with \((\omega(g,h),\omega'(g_1,h))\). Evaluation at \(0\) gives the two pointwise kernel identities. Under Proposition~\ref{B1.15b}, Lemma~\ref{B8.11e} ensures absolute convergence for unprimed polynomial vectors and for primed polynomial vectors in \(\mathcal S'_{t_1}\), and measurable change of variables along \(\kappa_W,\kappa_{W_1}\) gives the two integral identities.
\end{proof}

\subsection*{Proof of \ref{B8.16e}}\label{prf:B8.16e}\phantomsection
\begin{proof}
By Lemma~\ref{B8.16d}, \((\mathcal W_A,\mathcal W_B)\) map the full polynomial Schwartz spaces of \((\tilde\omega,\tilde\omega')\) bijectively onto the algebraic Fock spaces of \((\omega,\omega')\). Hence \(\tilde p=p\circ\mathcal W_A\) and \(\tilde p'=p'\circ\mathcal W_B\) are target-equivariant surjections onto \(\Pi\) (with target action transported along \(\kappa_X^{-1}\)) vanishing on the finite-adelic source differences and archimedean source Lie actions, just as \((q_{\rm ss},q'_{\rm ss})\) surject onto \(\Pi_{\rm ss}\). Fix either source side (say \(V\); the proof for \(V'\) is identical). Every element of \(\U(V)(\mathbb A_{F,f})\) stabilizes the local vacuum at almost all finite places, so a global finite source difference on a pure tensor expands as a finite sum of single-place source differences. At infinity, the source maximal compact group is connected and acts on finite-degree Fock orbits by exponentiating the source Lie action, so vanishing on the archimedean source Lie action implies compact source invariance. Thus the global source-coinvariant space coincides with \(M=M_\infty\otimes_\C\bigotimes'_{v_0}M_{v_0}\), where \(M_\infty\) is the archimedean source-coinvariant module and \(M_{v_0}\) is the local source-coinvariant module at each finite place \(v_0\) of \(F\) (with restricted product relative to the unramified vacua, which survive outside \(T\) since the local matching maps preceding Theorem~\ref{B8.11} send the vacua off \(T\) to nonzero vectors).

Let \(Q\in\{\Pi,\Pi_{\rm ss}\}\) be either nonzero jointly irreducible quotient of \(M\) with projection \(\pi_Q:M\twoheadrightarrow Q\). At infinity, the connected target maximal compact group integrates the compact Lie action on finite-degree polynomial orbits, so target Lie-invariant subspaces of quotients of \(M_\infty\) are compact-invariant. By real Howe duality in Proposition~\ref{B1.15b}, \(M_\infty\) has finite length and the sum of any two proper target-invariant submodules is proper; hence the sum \(N_\infty\) of all proper target-invariant submodules of \(M_\infty\) is proper and \(M_\infty/N_\infty\) is simple. Fixing finite factors yields a nonzero map \(M_\infty\to Q\) whose image has finite length and therefore contains a nonzero simple archimedean submodule. The sum of all simple archimedean submodules of \(Q\) is nonzero and invariant under the commuting action of \(\U(W^{\rm skew})(\mathbb A_{F,f})\), so it equals \(Q\) by joint irreducibility. Thus \(Q\) is archimedeanly semisimple, and projection onto simple summands shows that every fixed-finite-factor map \(M_\infty\to Q\) vanishes on \(N_\infty\).

At each finite place \(v_0\) of \(F\), put \(G_{v_0}=\U(W^{\rm skew})(F_{v_0})\). At split places with nontrivial source determinant character (including every \(v_0\in T\)), Lemma~\ref{B8.11f} shows that \(M_{v_0}\) has length one, so \(N_{v_0}=0\) is its unique maximal proper submodule and \(Q\) is spanned by simple \(G_{v_0}\)-images. At all remaining finite places (inert places and split places with trivial source determinant character; at split places \(M_{v_0}\) also has finite length by Bernstein--Zelevinsky finiteness in \(\mathscr P\), Definition~\ref{B1.15ac}), Corollary~\ref{B8.14a} and \eqref{eq:p8c2-theta-spherical-averages} in Lemma~\ref{B8.11g} give \(M_{v_0}=\operatorname{span}_\C G_{v_0}\bar s_{v_0}\). Let \(K_{v_0}\subset G_{v_0}\) be the integral maximal compact subgroup, and define on any smooth \(G_{v_0}\)-module (in particular on \(M_{v_0}\) and its smooth quotient \(Q\)) the Haar averaging operator \(e_{v_0}(y)=\int_{K_{v_0}}w_{v_0}(k)\,ky\,dk\) against the inverse right determinant character \(w_{v_0}\) on \(K_{v_0}\) (which is trivial at inert places and at split places off \(T\), and evaluates on each smooth vector \(y\) as a finite mean over cosets of an open subgroup of \(K_{v_0}\) fixing \(y\)). By Corollary~\ref{B8.14a} and the proof of \eqref{eq:p8c2-theta-spherical-averages} in Lemma~\ref{B8.11g}, \(e_{v_0}(M_{v_0})\subseteq\C\bar s_{v_0}\) and \(e_{v_0}(\bar s_{v_0})=\bar s_{v_0}\). Here \(N_{v_0}=\{y\in M_{v_0}:\forall g\in G_{v_0},\ e_{v_0}(gy)=0\}\) is the unique maximal proper \(G_{v_0}\)-submodule of \(M_{v_0}\), with simple quotient \(M_{v_0}/N_{v_0}\). In \(Q\), \(e_{v_0}\) commutes with \(\pi_Q\) and with all actions commuting with \(G_{v_0}\), so \(Q_0=\{q\in Q:\forall g\in G_{v_0},\ e_{v_0}(gq)=0\}\) is invariant under \(G_{v_0}\) and under all commuting actions at other finite and archimedean places; because \(M_{v_0}=\operatorname{span}_\C G_{v_0}\bar s_{v_0}\) and \(\pi_Q\) is \(G_{v_0}\)-equivariant and surjective onto \(Q\ne0\), some image \(\pi_Q(\bar s_{v_0}\otimes m^{v_0})\) is nonzero and fixed by \(e_{v_0}\), so \(Q_0\) is proper. Joint irreducibility forces \(Q_0=0\), so every fixed-other-places map \(M_{v_0}\to Q\) vanishes on \(N_{v_0}\) and factors through the simple module \(M_{v_0}/N_{v_0}\).

Consequently, both \(\ker\tilde p\) and \(\ker q_{\rm ss}\) contain \(\ker\bigl(M\to(M_\infty/N_\infty)\otimes_\C\bigotimes'_{v_0}(M_{v_0}/N_{v_0})\bigr)\). Because \(M_\infty/N_\infty\) and each \(M_{v_0}/N_{v_0}\) are simple of at most countable complex dimension, their restricted tensor product is simple by the countable-dimensional scalar-commutant and tensor-product argument in Theorem~\ref{B8.11}. Since \(\Pi\ne0\) and \(\Pi_{\rm ss}\ne0\), we obtain \(\ker\tilde p=\ker q_{\rm ss}\) and likewise \(\ker\tilde p'=\ker q'_{\rm ss}\). Descent gives target-equivariant isomorphisms \(\iota_A,\iota_B:\Pi\xrightarrow{\sim}\Pi_{\rm ss}\) with \(q_{\rm ss}=\iota_A\circ\tilde p\) and \(q'_{\rm ss}=\iota_B\circ\tilde p'\). Since \(\Pi\) has at most countable complex dimension, the countable-dimensional Schur argument gives \(\iota_B^{-1}\circ\iota_A=s\,\mathrm{id}_\Pi\) for a unique \(s\in\C^\times\). Whenever \(p(\Phi)=p'(\Phi_0)\), we have \(\tilde p(\Phi_D)=\tilde p'(\Phi_{0D})\) and hence \(q_{\rm ss}(\Phi_D)=\iota_A(\tilde p(\Phi_D))=\iota_B(s\,\tilde p'(\Phi_{0D}))=q'_{\rm ss}(s\Phi_{0D})\), so \((\Phi_D,s\Phi_{0D})\) is a polynomial match. Finally, Lemma~\ref{B8.11f} at the completely split good place \(v_1\notin T\) gives \(p'_{v_1}(\operatorname{span} U_{t_1})=\Pi_{v_1}\), whence \(q'_{\rm ss}\) maps the polynomial subspace of \(\mathcal S'_{t_1}\) onto \(\Pi_{\rm ss}\ne0\). Choosing a polynomial vector \(\Phi_{0D}\in\mathcal S'_{t_1}\) with \(q'_{\rm ss}(\Phi_{0D})\ne0\) and using surjectivity of \(p\) to pick \(\Phi=\mathcal W_A(\Phi_D)\) with \(p(\Phi)=p'(\mathcal W_B(\Phi_{0D}))\) shows that \(s\) is uniquely determined on algebraic pairs with \(\Phi_{0D}\in\mathcal S'_{t_1}\).
\end{proof}

\subsection*{Proof of \ref{B8.16}}\label{prf:B8.16}\phantomsection
\begin{proof}
Put \(n=\operatorname{rank}M\), \(d_X=\operatorname{rank}X\), and \(f_F=[F:\Q]\), and choose a maximal rational Witt frame of \(\U(M)\) of Witt index \(r\). At the infinite place where \(\operatorname{sig}(M)=(n-p,p)\), the two dual \(r\)-dimensional isotropic blocks span a subspace of signature \((r,r)\), so \(r\le p\) and \(d_X<2(n-p)\le2(n-r)\). By Proposition~\ref{B1.15b}, choose a closed Siegel set \(\mathfrak S\) of points \(p_0 a k\) with \(p_0\) in a compact subset \(\Omega\) of the norm-one minimal parabolic subgroup, \(k\) in a compact subgroup \(K\subset\U(M)(\mathbb A_F)\), and archimedean torus coordinates \(t_i\ge c_0 t_{i+1}\), \(t_{r-1}\ge c_0>0\), such that \(\mathfrak S\) meets every left \(\U(M)(F)\)-orbit, all conjugates \(a^{-1}p_0 a\) lie in a fixed compact set, and Haar integration over \(\mathfrak S\) is bounded up to a constant factor by cone integration against \(\prod_{i<r}t_i^{-1-2f_F(n-2i-1)}\prod_{i<r}dt_i\). Because \(g\mapsto\theta_\omega(g,h;\Phi)\) is continuous and left \(\U(M)(F)\)-invariant, its absolute integral over the fundamental domain \(\mathcal F\) is bounded by its integral over \(\mathfrak S\). On \(g=p_0 a k=a(a^{-1}p_0 a k)\), moving the bounded factor \(a^{-1}p_0 a k\) and \(h\) onto \(\Phi\) costs finitely many archimedean Schwartz--Sobolev seminorms growing polynomially in \(H_X(h)\), while the lattice count over contracted tuple coordinates as in Lemma~\ref{B8.11e} contributes \(t_i^{f_F d_X}\) per isotropic line. Hence there exist \(C>0\) and \(N\in\N\) such that
\begin{equation}\label{eq:p8c2-square-kernel}
 |\theta_\omega(p_0 a k,h;\Phi)|\le C H_X(h)^N\prod_{i<r}t_i^{f_Fd_X}.
\end{equation}
\leanref{GlobalTheory.\allowbreak{}T.\allowbreak{}hermitianWittIndex\_\allowbreak{}le\_\allowbreak{}negative\_\allowbreak{}signature}{Development.\allowbreak{}AbelianVarietyTheory.\allowbreak{}Proofs.\allowbreak{}TwistedOperators.\allowbreak{}HermitianWittIndexLeNegativeSignatureStatement}{17--19}
\leanref{GlobalTheory.\allowbreak{}T.\allowbreak{}thetaKernel\_\allowbreak{}siegelSet\_\allowbreak{}le\_\allowbreak{}torusPower\_\allowbreak{}height}{Development.\allowbreak{}AbelianVarietyTheory.\allowbreak{}Proofs.\allowbreak{}TwistedOperators.\allowbreak{}ThetaKernelSiegelSetLeTorusPowerHeightStatement}{17--19}
\leanref{GlobalTheory.\allowbreak{}T.\allowbreak{}lintegral\_\allowbreak{}fundamentalDomain\_\allowbreak{}le\_\allowbreak{}of\_\allowbreak{}meets\_\allowbreak{}all\_\allowbreak{}orbits}{Development.\allowbreak{}AbelianVarietyTheory.\allowbreak{}Proofs.\allowbreak{}TwistedOperators.\allowbreak{}LintegralFundamentalDomainLeOfMeetsOrbitsStatement}{7--9}
Multiplying \eqref{eq:p8c2-square-kernel} by the cone density gives exponents \(A_i-1\) in \(t_i\) with \(A_i=f_F(d_X-2n+4i+2)\), whose partial sums satisfy
\begin{equation}\label{eq:p8c2-cone-exponents}
 \sum_{i=0}^j A_i=f_F(j+1)\bigl(d_X-2n+2(j+1)\bigr)<0\qquad(0\le j<r)
\end{equation}
because \(j+1\le r\) and \(d_X<2(n-r)\). Successive integration from \(t_0\) to \(t_{r-1}\) via \(\int_b^\infty t^{B-1}\,dt=b^B/(-B)\) at each negative partial sum \(B\) (with cone factor \(1\) when \(r=0\)) proves \eqref{eq:p8c2-square-moderate}.
\leanref{GlobalTheory.\allowbreak{}T.\allowbreak{}lintegral\_\allowbreak{}siegelCone\_\allowbreak{}inv\_\allowbreak{}modulus\_\allowbreak{}lt\_\allowbreak{}top\_\allowbreak{}weil}{Development.\allowbreak{}AbelianVarietyTheory.\allowbreak{}Proofs.\allowbreak{}TwistedOperators.\allowbreak{}LintegralSiegelConeInvModulusLtTopWeilStatement}{8--10}

Now fix the standard additive character \(\psi\) of sign \(-1\), the frames \((A,A_1,B)\), and rational numbers \((a_\tau,b_\tau)\). For any twisted-square triple \((\alpha,\beta)\), existence of oscillator Weil representations in \(\mathscr P\) (Definition~\ref{B1.15ac}) supplies \(\omega,\omega'\) of signs \(+,-\) and splitting pairs \((\beta,\beta/\alpha),(\alpha,\alpha/\beta)\). Lemma~\ref{B8.16d} transfers \((W,W_1,X,\alpha,\beta,\omega,\omega')\) to a triple of Theorem~\ref{B8.11}, and pulling the resulting jointly irreducible seesaw quotient \(\Pi_{\rm ss}\) and matching maps \((q_{\rm ss},q'_{\rm ss})\) back along \((\mathcal W_A,\mathcal W_B)\) (extended linearly from the algebraic Fock subspaces) gives matching data \((\omega,\omega',\Pi,p,p')\). To construct \((P_{\rm real},q,q')\) prior to any match, set \(P_{\rm real}=0\) and \(q=0,q'=0\) if no matching data for \(\psi\) exist with exponents \((a_\tau,b_\tau)\). Otherwise choose one witness match \((\omega_0,\omega'_0,\Pi_0,p_0,p'_0)\) for \(\psi\) with exponents \((a_\tau,b_\tau)\). The archimedean source and target Lie actions on \(\mathcal P(W,X)\) and \(\mathcal P(W_1,X)\) depend solely on the Gram matrices, frames, and \((a_\tau,b_\tau)\). Let \(N_A\subset\mathcal P(W,X)\) and \(N_B\subset\mathcal P(W_1,X)\) be the sums of all proper target-Lie-invariant submodules containing the archimedean source Lie relations. Since \(\Pi_0\ne0\) is spanned on each side by images of Fock vectors, some fixed finite factor on each side produces a proper kernel containing the source relations; by real Howe duality in Proposition~\ref{B1.15b} (bounded chain length and properness of pairwise sums), \(N_A\) and \(N_B\) are proper, with simple countable-dimensional quotients \(P_A=\mathcal P(W,X)/N_A\) and \(P_B=\mathcal P(W_1,X)/N_B\). For any match \((\omega,\omega',\Pi,p,p')\) with exponents \((a_\tau,b_\tau)\), the archimedean finite-length and socle-semisimplicity argument of Lemma~\ref{B8.16e} shows that \(\Pi\) is semisimple for the archimedean target Lie action and that for every permitted finite factor \(f\) (resp.\ \(f'\)) the map \(P\mapsto p(\operatorname{Fock}(f,P))\) (resp.\ \(P'\mapsto p'(\operatorname{Fock}'(f',P'))\)) vanishes on \(N_A\) (resp.\ \(N_B\)). For \(\Pi_0\), images from each side span \(\Pi_0\), so every simple archimedean summand of \(\Pi_0\) is isomorphic to both \(P_A\) and \(P_B\). Fix a target-Lie isomorphism \(\iota_{\rm real}:P_B\xrightarrow{\sim}P_A\), put \(P_{\rm real}=P_A\), let \(q:\mathcal P(W,X)\twoheadrightarrow P_{\rm real}\) be the canonical projection, and put \(q'=\iota_{\rm real}\circ(\mathcal P(W_1,X)\twoheadrightarrow P_B)\).

On either side, the permitted finite theta functions form a complex subspace \(S\) (resp.\ \(S'\)) on which \(\operatorname{Fock}\) (resp.\ \(\operatorname{Fock}'\)) is \(\C\)-bilinear. For any match \((\omega,\omega',\Pi,p,p')\) with exponents \((a_\tau,b_\tau)\), every \(f\in S\) and \(f'\in S'\) induce unique target-Lie intertwiners \(\varphi_f,\varphi'_{f'}\in P_{\rm fin}=\operatorname{Hom}_{\rm Lie}(P_{\rm real},\Pi)\) satisfying \(p(\operatorname{Fock}(f,P))=\varphi_f(q(P))\) and \(p'(\operatorname{Fock}'(f',P'))=\varphi'_{f'}(q'(P'))\). Since a match exists for \((a_\tau,b_\tau)\), we are in the witness branch and \(P_{\rm real}\ne0\). Simplicity and countable complex dimension of \(P_{\rm real}\) give \(\operatorname{End}_{\rm Lie}(P_{\rm real})=\C\,\mathrm{id}\). Indeed, a nonscalar endomorphism \(T\) would have \(T-z\,\mathrm{id}\) bijective for all \(z\in\C\), and \(\{(T-z\,\mathrm{id})^{-1}v_0:z\in\C\}\) for \(v_0\ne0\) would be uncountably linearly independent. Because \(\Pi=\sum_{f\in S,\,\varphi_f\ne0}\varphi_f(P_{\rm real})\) is a sum of simple submodules, Zorn's lemma expresses \(\Pi=\bigoplus_{i\in I}\varphi_i(P_{\rm real})\) for a subfamily \(\{\varphi_i\}_{i\in I}\) with \(\varphi_i=\varphi_{f_i}\ne0\) and \(f_i\in S\). Any \(\psi_0\in P_{\rm fin}\) sends a nonzero \(x_0\in P_{\rm real}\) into a finite partial sum \(\bigoplus_{m=1}^k\varphi_{i_m}(P_{\rm real})\); simplicity of \(P_{\rm real}\) confines \(\psi_0(P_{\rm real})\) to those \(k\) summands, and \(\operatorname{End}_{\rm Lie}(P_{\rm real})=\C\,\mathrm{id}\) gives \(\psi_0=\sum_{m=1}^k c_m\varphi_{i_m}\). Thus \(f\mapsto\varphi_f\) maps \(S\) onto \(P_{\rm fin}\), and the identical argument on \(S'\) shows that \(f'\mapsto\varphi'_{f'}\) maps \(S'\) onto \(P_{\rm fin}\). Moreover \(\{\varphi_i\}_{i\in I}\) is a \(\C\)-basis of \(P_{\rm fin}\), so evaluation \(e:P_{\rm real}\otimes_\C P_{\rm fin}\to\Pi\), \(\sum_{m=1}^k y_m\otimes\varphi_{i_m}\mapsto\sum_{m=1}^k\varphi_{i_m}(y_m)\), is a linear isomorphism. Linear extension of \(f\mapsto\varphi_f\) and \(f'\mapsto\varphi'_{f'}\) across vector-space complements of \(S\) and \(S'\) gives \(r,r'\) satisfying the displayed tensor formulas.
\leanref{GlobalTheory.\allowbreak{}T.\allowbreak{}twist\_\allowbreak{}triple\_\allowbreak{}matches}{Development.\allowbreak{}AutomorphicArithmetic.\allowbreak{}ThetaCorrespondence.\allowbreak{}ThetaPairLifts.\allowbreak{}LocalTypeDictionary.\allowbreak{}TripleExistsMatching}{43--47}
\leanref{GlobalTheory.\allowbreak{}T.\allowbreak{}twisted\_\allowbreak{}quotient\_\allowbreak{}real}{Development.\allowbreak{}AutomorphicArithmetic.\allowbreak{}ThetaCorrespondence.\allowbreak{}ThetaPairLifts.\allowbreak{}LocalTypeDictionary.\allowbreak{}TwistedQuotientReal}{180--193}
\leanref{GlobalTheory.\allowbreak{}dixmier\_\allowbreak{}commutant\_\allowbreak{}scalar}{Development.\allowbreak{}AbelianVarietyTheory.\allowbreak{}Proofs.\allowbreak{}TwistedOperators.\allowbreak{}DixmierCommutantScalar}{43--110}
\leanref{GlobalTheory.\allowbreak{}FinFockLin.\allowbreak{}finF\_\allowbreak{}lin}{Development.\allowbreak{}AbelianVarietyTheory.\allowbreak{}Proofs.\allowbreak{}TwistedOperators.\allowbreak{}FinFLin}{22--51}
\leanref{GlobalTheory.\allowbreak{}SquareLM.\allowbreak{}productForm}{Development.\allowbreak{}AbelianVarietyTheory.\allowbreak{}Proofs.\allowbreak{}TwistedOperators.\allowbreak{}LsumMapDomainSome}{286--330}

Finally, fix an additive character \(\psi\) with a twisted-square triple \((W,W_1,X,\alpha,\beta)\), frames \((A,A_1,B)\), matching data \((\omega,\omega',\Pi,p,p')\), and Haar measures and fundamental domains \((\nu,\mathcal F),(\nu_1,\mathcal F_1)\). Apply Lemma~\ref{B8.16d} to obtain the seesaw transfer \((e,V,V',W^{\rm skew},\tilde\omega,\tilde\omega',\kappa_W,\kappa_{W_1},\kappa_X)\) and the bijections \((\mathcal W_A,\mathcal W_B)\) from the full polynomial Schwartz spaces onto the algebraic Fock spaces. Theorem~\ref{B8.11} provides matching maps \((q_{\rm ss},q'_{\rm ss})\) onto \(\Pi_{\rm ss}\), the decay and convergence bounds on \(\mathcal S'_{t_1}\), the surjectivity \(p'_{v_1}(U_{t_1})=\Pi_{v_1}\) and kernel identity \(\ker p'_{v_1}=A(\mathcal S_{v_1})\) at the completely split good place \(v_1\notin T\), and the unique nonzero scalar \(c_1\in\C^\times\) in \eqref{eq:p8c2-matched-integrals}. Lemma~\ref{B8.16e} supplies the seesaw comparison scalar \(s\in\C^\times\); put \(c=c_1 s\in\C^\times\).

Given any primed algebraic Fock vector \(\Phi'\), let \(\Phi'_D=\mathcal W_B^{-1}(\Phi')\) be its primed polynomial Schwartz preimage and write \(\Phi'_D=\sum_j\phi_{j,v_1}\otimes\Psi_j^{v_1}\). Since \(p'_{v_1}(U_{t_1})=\Pi_{v_1}\) and \(\ker p'_{v_1}=A(\mathcal S_{v_1})\), each \(\phi_{j,v_1}\) admits \(\phi_{0j,v_1}\in\operatorname{span}_\C U_{t_1}\) with the same image in \(\Pi_{v_1}\). For \(\Phi_{0D}=\sum_j\phi_{0j,v_1}\otimes\Psi_j^{v_1}\in\mathcal S'_{t_1}\), the difference \(\Phi'_D-\Phi_{0D}\) is a finite linear combination of local differences \(\tilde\omega'(g',1)\Psi_D-\Psi_D\) with \(g'\in\U(V')(F_{v_1})\) and \(\Psi_D\) in the primed polynomial Schwartz space. Put \(\Phi_0=\mathcal W_B(\Phi_{0D})\) in the primed algebraic Fock space. For each generator, \(u_f=\kappa_{W_1}^{-1}(g')\) lies in \(\U(W_1)(\mathbb A_{F,f})\), \(\kappa_X(1)=1\), and \(\Psi=\mathcal W_B(\Psi_D)\) is a primed algebraic Fock vector satisfying \(\mathcal W_B(\tilde\omega'(g',1)\Psi_D-\Psi_D)=\omega'(u_f,1)\Psi-\Psi\in\mathcal R_{\rm fin}\), whence \(\Phi'-\Phi_0\in\mathcal R_{\rm fin}\).

Since \(\Phi_{0D}\in\mathcal S'_{t_1}\), \eqref{eq:p8c2-full-rank-decay} at \(N=0\) and Lemma~\ref{B8.11e} give integrability on \(\kappa_{W_1}(\mathcal F_1)\) and a uniform bound \(|\theta_{\tilde\omega'}(g',h';\Phi_{0D})|\le C_0\operatorname{ht}(h')^{k_0}\). By the kernel identity and height bound \(\operatorname{ht}(\kappa_X(h))\le C'H_X(h)^{N'}\) of Lemma~\ref{B8.16d}, we have \(|\theta_{\omega'}(g_1,h;\Phi_0)|\le C_0(C')^{k_0}H_X(h)^{N'k_0}\) for all \(g_1\in\U(W_1)(\mathbb A_F)\) and \(h\in\U(X)(\mathbb A_F)\). By Proposition~\ref{B1.15b}, \(\U(W_1)(\mathbb A_F)\) admits a closed Siegel set of finite \(\nu_1\)-measure meeting every left \(\U(W_1)(F)\)-orbit, whence \(\nu_1(\mathcal F_1)<\infty\). Integration over \(\mathcal F_1\) gives \(\int_{\mathcal F_1}|\theta_{\omega'}(g_1,h;\Phi_0)|\,d\nu_1(g_1)\le\nu_1(\mathcal F_1)C_0(C')^{k_0}H_X(h)^{N'k_0}\), which proves \eqref{eq:p8c2-square-moderate} for \(\Phi_0\).
\leanref{GlobalTheory.\allowbreak{}T.\allowbreak{}theta\_\allowbreak{}integral\_\allowbreak{}moderate\_\allowbreak{}growth\_\allowbreak{}of\_\allowbreak{}bounded\_\allowbreak{}kernel}{Development.\allowbreak{}AbelianVarietyTheory.\allowbreak{}Proofs.\allowbreak{}TwistedOperators.\allowbreak{}ThetaIntegralModerateGrowthOfBoundedKernelStatement}{16--29}

For any unprimed algebraic Fock vector \(\Phi\) with \(p(\Phi)=p'(\Phi_0)\), put \(\Phi_D=\mathcal W_A^{-1}(\Phi)\). By Lemma~\ref{B8.16e}, \((\Phi_D,s\Phi_{0D})\) is a polynomial match for \((q_{\rm ss},q'_{\rm ss})\) with \(s\Phi_{0D}\in\mathcal S'_{t_1}\). By Lemma~\ref{B8.16d}, \eqref{eq:p8c2-matched-integrals}, and linearity of \(I'\), for every \(h\in\U(X)(\mathbb A_F)\) we obtain
\[
 \theta\operatorname{Int}_\omega(\Phi)(h)
 =I(\Phi_D,\kappa_X(h))
 =c_1 I'(s\Phi_{0D},\kappa_X(h))
 =c_1 s\,I'(\Phi_{0D},\kappa_X(h))
 =c\,\theta\operatorname{Int}_{\omega'}(\Phi_0)(h).
\]
Since \(h\) is arbitrary, \(\theta\operatorname{Int}_\omega(\Phi)=c\,\theta\operatorname{Int}_{\omega'}(\Phi_0)\).
\leanref{GlobalTheory.\allowbreak{}SquareAgree.\allowbreak{}I\_\allowbreak{}smul}{Development.\allowbreak{}AbelianVarietyTheory.\allowbreak{}Proofs.\allowbreak{}TwistedOperators.\allowbreak{}ThetaIntegralsOfMinguez}{26--32}
\end{proof}

\subsection*{Proof of \ref{B8.17a}}\label{prf:B8.17a}\phantomsection
\begin{proof}
We first reduce the set of tests to finitely many generators over a number field. On either side choose a common open finite level fixing the finite test vectors for all \(t\). It exists by smoothness and finiteness of \(T\). Finite translation at this level commutes with insertion of the polynomial at infinity. We can therefore shrink to a compact open subgroup \(J\) under which all the integrals are right invariant, for every \(k\).

The standard target has Gram matrix \(1\), so its real groups are compact. Its unitary group is linear algebraic over the number field \(F\). Borel's finite class-number theorem gives a finite set \(H\) of finite adelic representatives at level \(J\). Indeed one can intersect \(J\) with the standard integral level; the intersection has finite index in the latter. Apply the theorem at the integral level and subdivide its cosets. Thus, with finite adeles embedded with trivial archimedean component, every \(h\) in the statement has a decomposition
\[
 h=\gamma r h_0\kappa,\qquad
 \gamma\in\U(W)(F),\quad r\in\U(W)(F_\infty),\quad
 h_0\in H,\quad \kappa\in J .
\]
Here \(r=\gamma_\infty^{-1}\). This last restriction on \(r\) determines the coefficient field in the next step. \cite[Thm.~5.1]{borel-1963}
\leanref{GlobalTheory.\allowbreak{}BottomLevel.\allowbreak{}rightTranslate\_\allowbreak{}preserves\_\allowbreak{}level}{Development.\allowbreak{}AutomorphicArithmetic.\allowbreak{}ThetaCorrespondence.\allowbreak{}ThetaPairLifts.\allowbreak{}LocalTypeDictionary.\allowbreak{}BottomFormRightTranslateLevel}{21--45}
\leanref{GlobalTheory.\allowbreak{}BotClassFin.\allowbreak{}bottomClass\_\allowbreak{}exists\_\allowbreak{}finite\_\allowbreak{}reps}{Development.\allowbreak{}AutomorphicArithmetic.\allowbreak{}ThetaCorrespondence.\allowbreak{}ThetaPairLifts.\allowbreak{}LocalTypeDictionary.\allowbreak{}BottomClassExistsFiniteReps}{23--38}

Remove \(\gamma\) by rational left invariance. The Fock action of \(r\) substitutes its place matrices in the polynomial and multiplies by the character scalar of that action. In the standard Gram matrices and frames the entries of \(r\) and their conjugates lie in the compositum of the embedded conjugates of the host CM field \(E\), since \(\gamma\) is rational. The determinant character on each compact factor has integral exponent. The substitution and its scalar therefore have coefficients in a fixed number field \(M_0\); one may use the normal closure of \(E\) in \(\C\). Expanding the substituted polynomial and removing \(\kappa\) expresses \(I_t(h,k)\) as an \(M_0\)-linear combination of the \(I_t(h_0,k')\). The positive-sign model gives the same expansion for \(I'_t\) with its own polynomial variables.
\leanref{GlobalTheory.\allowbreak{}RealRatAct.\allowbreak{}rationalAction\_\allowbreak{}real\_\allowbreak{}translate}{Development.\allowbreak{}AutomorphicArithmetic.\allowbreak{}ThetaCorrespondence.\allowbreak{}ThetaPairLifts.\allowbreak{}LocalTypeDictionary.\allowbreak{}RealTranslatePreservesRationalAction}{61--133}
\leanref{GlobalTheory.\allowbreak{}BotSpanSnd.\allowbreak{}real\_\allowbreak{}right}{Development.\allowbreak{}AbelianVarietyTheory.\allowbreak{}Proofs.\allowbreak{}TwistedOperators.\allowbreak{}BottomTriangleDualAlgebraicityFiniteData}{414--486}

There are also only finitely many exponents \(k'\) with nonzero integral. At each infinite place \(\tau\), a central source rotation acts on a monomial by a fixed character times the rotation to the signed total degree \(d_\tau(k')\) in its variables at \(\tau\). Invariance of the integral forces this scalar to be one whenever the integral is nonzero. Comparing two surviving monomials and taking a rotation of infinite order gives equality of their degrees at \(\tau\). For each \(t\), if there is a surviving monomial, take its multidegree. Finitely many monomials in the given variables have that multidegree. Taking the union over \(T\) gives a finite set \(D_0\) on each side. Thus each family is spanned over \(M_0\) by the integrals indexed by \(T\times H\times D_0\).
\leanref{GlobalTheory.\allowbreak{}BotSpan.\allowbreak{}classSums\_\allowbreak{}span\_\allowbreak{}bottom}{Development.\allowbreak{}AutomorphicArithmetic.\allowbreak{}ThetaCorrespondence.\allowbreak{}ThetaPairLifts.\allowbreak{}LocalTypeDictionary.\allowbreak{}BottomSpannedByClassSums}{26--68}
\leanref{GlobalTheory.\allowbreak{}BotSpanSnd.\allowbreak{}twistedClassSums\_\allowbreak{}span\_\allowbreak{}bottom}{Development.\allowbreak{}AutomorphicArithmetic.\allowbreak{}ThetaCorrespondence.\allowbreak{}ThetaPairLifts.\allowbreak{}LocalTypeDictionary.\allowbreak{}BottomSpannedByTwistedClassSums}{32--72}

We next construct a period on each side. On the first source, take an open finite level fixing the test data. Decomposition of the definite source quotient into finite classes gives representatives \(z_l\), finite stabilizers of orders \(w_l>0\), a positive common measure factor \(v\), and probability Haar measure \(dg\) on the compact real group. If \(K_t(P;h,x)\) denotes the theta kernel evaluated with polynomial \(P\), finite target point \(h\) and source point \(x\), then
\[
 I_t(h,k)=v\sum_l\frac1{w_l}
           \int_{\U(A)(F_\infty)}K_t(X^k;h,z_lg)\,dg .
\]
The factor \(1/w_l\) divides out the stabilizer of that class. The same decomposition on \(A'\) gives \(v',w'_l,z'_l\) for the primed kernel. The sources are definite and the finite data are smooth at the chosen levels, so both formulas apply to all the displayed monomials.
\leanref{GlobalTheory.\allowbreak{}T.\allowbreak{}bottom\_\allowbreak{}class\_\allowbreak{}sum\_\allowbreak{}integral\_\allowbreak{}neg}{Development.\allowbreak{}AutomorphicArithmetic.\allowbreak{}ThetaCorrespondence.\allowbreak{}ThetaPairLifts.\allowbreak{}LocalTypeDictionary.\allowbreak{}BottomClassSumIntegralNeg}{121--128}
\leanref{GlobalTheory.\allowbreak{}T.\allowbreak{}twist\_\allowbreak{}integral\_\allowbreak{}eq\_\allowbreak{}bottom\_\allowbreak{}classSum}{Development.\allowbreak{}AutomorphicArithmetic.\allowbreak{}ThetaCorrespondence.\allowbreak{}ThetaPairLifts.\allowbreak{}LocalTypeDictionary.\allowbreak{}TwistIntegralEqBottomClassSum}{116--123}

On the first source the compact action on polynomials is a common character times polynomial substitution. Its central characters select a single multidegree \(d=(d_\tau)\) for the real average. If no degree survives, all averages vanish and we may take period \(1\). In the surviving degree the averaging matrix has algebraic entries. To see this, let \(P_d\) be the Haar average of the representation matrix. On the dense subgroup of compact unitary matrices with algebraic entries, the character and the polynomial action have algebraic matrices. For any automorphism \(\sigma\) of \(\C\) over \(\overline{\Q}\), the identities of left and right invariance for \(P_d\) thus hold for \(\sigma(P_d)\) on that subgroup. Continuity in the group variable extends both identities to the compact group. Averaging gives
\[
                P_d\,\sigma(P_d)=\sigma(P_d),
                \qquad \sigma(P_d)P_d=\sigma(P_d).
\]
Applying \(\sigma\) to the first identity for \(\sigma^{-1}\) also gives \(\sigma(P_d)P_d=P_d\). Thus \(P_d\) is fixed and has algebraic entries.
\leanref{GlobalTheory.\allowbreak{}RealAvg.\allowbreak{}real\_\allowbreak{}average\_\allowbreak{}kernel\_\allowbreak{}isAlgebraic}{Development.\allowbreak{}AutomorphicArithmetic.\allowbreak{}ThetaCorrespondence.\allowbreak{}ThetaPairLifts.\allowbreak{}LocalTypeDictionary.\allowbreak{}TwistedRealAverageKernelIsAlgebraic}{22--112}
\leanref{GlobalTheory.\allowbreak{}HaarAlg.\allowbreak{}isAlgebraic\_\allowbreak{}integral}{Development.\allowbreak{}AbelianVarietyTheory.\allowbreak{}Proofs.\allowbreak{}TwistedOperators.\allowbreak{}IsAlgebraicIntegral}{64--100}

In the finite test expand the level projector in Heisenberg translates of the standard vacuum. Its finite Fourier coefficients are algebraic. Moving a finite pair \((z_l,h)\) past these translates again gives a finite combination with algebraic coefficients. Evaluation at zero now asks for derivatives of degree \(d_\tau\) at each \(\tau\). In the standard self-dual lattices choose the CM vacuum with the prescribed nonvanishing two-torsion normalization. Comparing the Fock vacuum with this one contributes a fixed nonzero scalar \(a\). Write \(\omega_\tau\ne0\) for the CM period in the derivative coordinate at \(\tau\). A monomial of multidegree \(d\) brings the factor \(\prod_\tau\omega_\tau^{d_\tau}\). In the negative model the comparison is conjugated. The remaining value is a Taylor coefficient of a rational lattice translate of the normalized vacuum, or an algebraic linear combination of such coefficients. The CM Taylor formula applies to the standard self-dual lattices: the product of the source and target ranks is even, the translating points have an integral multiple in the lattice, and the translation phases at such points are roots of unity. It gives algebraic coefficients also after translation and finite algebraic linear combination. Hence one may take
\[
                  \Omega=v\,\overline{a\prod_\tau
                                           \omega_\tau^{d_\tau}} .
\]
It works for every term of every class sum. In particular \(h,t,k\) do not enter its choice.
\leanref{GlobalTheory.\allowbreak{}ThetaAlgebraic.\allowbreak{}taylor\_\allowbreak{}algebraic}{Development.\allowbreak{}AbelianVarietyTheory.\allowbreak{}Proofs.\allowbreak{}FoundationProofs.\allowbreak{}ThetaTaylorCoeffsAlgebraicAtComplexMultiplicatio}{594--635}
\leanref{GlobalTheory.\allowbreak{}BotTransJet.\allowbreak{}isAlgebraic\_\allowbreak{}pureTensorO\_\allowbreak{}div\_\allowbreak{}period}{Development.\allowbreak{}AutomorphicArithmetic.\allowbreak{}ThetaCorrespondence.\allowbreak{}ThetaPairLifts.\allowbreak{}LocalTypeDictionary.\allowbreak{}JetsDivPeriodIsAlgebraic}{20--86}
\leanref{GlobalTheory.\allowbreak{}BotAlgI.\allowbreak{}classSum\_\allowbreak{}kernel\_\allowbreak{}isAlgebraic}{Development.\allowbreak{}AutomorphicArithmetic.\allowbreak{}ThetaCorrespondence.\allowbreak{}ThetaPairLifts.\allowbreak{}LocalTypeDictionary.\allowbreak{}ClassSumKernelIsAlgebraic}{19--62}
\leanref{GlobalTheory.\allowbreak{}BotAlgI.\allowbreak{}classSums\_\allowbreak{}isAlgebraic}{Development.\allowbreak{}AutomorphicArithmetic.\allowbreak{}ThetaCorrespondence.\allowbreak{}ThetaPairLifts.\allowbreak{}LocalTypeDictionary.\allowbreak{}ClassSumKernelIsAlgebraic}{63--88}

On the primed source the real average is the projection for the determinant character. It writes the averaged monomial as an algebraic linear combination of products of maximal minors, with a fixed number \(n_\tau\) of minors at each infinite place. Every term thus has total degree \(a'\sum_\tau n_\tau\), where \(a'=\dim_E A'\). In the standard model the period per derivative coordinate on this side has a common value \(p\ne0\). If \(a_0\ne0\) compares the two vacua, the minor calculation gives the same factor \(a_0p^{a'\sum_\tau n_\tau}\) for all these products. Indeed, expanding a minor gives integer coefficients, and moving the finite Heisenberg translates through the differential polynomial leaves algebraic coefficients and rational lattice translates. Applying the same CM Taylor formula to these translates gives algebraic values after division by that factor. Inserting the class weights gives the period
\[
                          \Omega'=v'a_0p^{a'\sum_\tau n_\tau}.
\]
\leanref{GlobalTheory.\allowbreak{}RealProj.\allowbreak{}realProjection\_\allowbreak{}average\_\allowbreak{}monomial}{Development.\allowbreak{}AutomorphicArithmetic.\allowbreak{}ThetaCorrespondence.\allowbreak{}ThetaPairLifts.\allowbreak{}LocalTypeDictionary.\allowbreak{}TwistRealProjectionAverageMonomial}{15--38}
\leanref{GlobalTheory.\allowbreak{}BottomFieldDict.\allowbreak{}vacuum\_\allowbreak{}translate\_\allowbreak{}minor\_\allowbreak{}jet\_\allowbreak{}alt}{Development.\allowbreak{}AutomorphicArithmetic.\allowbreak{}ThetaCorrespondence.\allowbreak{}ThetaPairLifts.\allowbreak{}LocalTypeDictionary.\allowbreak{}BottomFieldMinorJetDescription}{16--100}
\leanref{GlobalTheory.\allowbreak{}BottomFieldDict.\allowbreak{}kernel\_\allowbreak{}minor\_\allowbreak{}polynomial\_\allowbreak{}eq\_\allowbreak{}jet}{Development.\allowbreak{}AutomorphicArithmetic.\allowbreak{}ThetaCorrespondence.\allowbreak{}ThetaPairLifts.\allowbreak{}LocalTypeDictionary.\allowbreak{}BottomFieldMinorJetDescription}{101--132}
\leanref{GlobalTheory.\allowbreak{}BotFormulaI'.\allowbreak{}bottomClassSum\_\allowbreak{}formula}{Development.\allowbreak{}AutomorphicArithmetic.\allowbreak{}ThetaCorrespondence.\allowbreak{}ThetaPairLifts.\allowbreak{}LocalTypeDictionary.\allowbreak{}BottomClassSumEqFormula}{17--53}

For each side adjoin to \(M_0\) the values at its finite set \(T\times H\times D_0\), divided by the period just constructed for that side. They are algebraic, so each extension is finite. Divide the spanning identities by the same period; their coefficients still lie in \(M_0\). All normalized integrals therefore lie in the compositum of these two finite extensions.
\leanref{GlobalTheory.\allowbreak{}RowEx.\allowbreak{}BotField.\allowbreak{}field\_\allowbreak{}of\_\allowbreak{}gens}{Development.\allowbreak{}AbelianVarietyTheory.\allowbreak{}Proofs.\allowbreak{}TwistedOperators.\allowbreak{}NormalFormBottom}{26--48}
\end{proof}

\subsection*{Proof of \ref{B8.17}}\label{prf:B8.17}\phantomsection
\begin{proof}
Take the maximum of the thresholds for grid existence, completion of the tests and the coefficient fields. For fixed \(\varepsilon,C_F\) above the required bounds, take the maximum of their scale cutoffs. Fix a scale beyond it, a preliminary datum with frame, and the finite set \(\Xi\). The information needed from the level of a realization is the set
\[
                         \kappa=\{(\xi_t,Q_t):t\in T\}.
\]
The level condition gives \(Q_t\subseteq\operatorname{Ram}(\xi_t)
\subseteq Q_t\cup\overline{Q_t}\). For each character of \(\Xi\) occurring in a realization its ramification set is therefore finite, and there are finitely many possibilities for \(Q_t\). Taking subsets of the finite set of possible pairs gives finitely many \(\kappa\).
\leanref{GlobalTheory.\allowbreak{}RowEx.\allowbreak{}keyFin\_\allowbreak{}holds}{Development.\allowbreak{}AbelianVarietyTheory.\allowbreak{}Proofs.\allowbreak{}TwistedOperators.\allowbreak{}KeyFinHolds}{22--45}

For each realized \(\kappa\) choose a representative. The grid construction gives its grid and Weil system. Theorems~\ref{B8.11} and~\ref{B8.16} give matching and the proportionality of the theta integrals for the ordinary and twisted tests. With separation from Theorem~\ref{B8.15} these complete the test choice at \((1,C_F,b(s))\). The representative satisfies the Hida level conditions, in particular finiteness of \(T\) and the prescribed character and ramification conditions.

For each \(0\le i,j<m\) choose the standard positive definite source spaces of the two bottom ranks and standard target of their sum rank. Use the identity Gram matrices, standard self-dual lattices and standard frames. The two level characters have the restrictions prescribed by the ranks, so give the two Weil actions. Use their vacua with the grid additive character, insert the local idempotents and the prescribed conductor powers, and take the quotient measures for these two sources. This gives a standard bottom triple for this pair. Lemma~\ref{B8.17a} gives its two periods and a number field. Take the compositum over the finitely many pairs. Each pair keeps its own triple and its own two periods.
\leanref{GlobalTheory.\allowbreak{}RowTest.\allowbreak{}exists\_\allowbreak{}bottomTri}{Development.\allowbreak{}AbelianVarietyTheory.\allowbreak{}Proofs.\allowbreak{}TwistedOperators.\allowbreak{}BottomExistsIntegrable}{104--135}
\leanref{GlobalTheory.\allowbreak{}RowEx.\allowbreak{}bottomField\_\allowbreak{}isNumberField}{Development.\allowbreak{}AutomorphicArithmetic.\allowbreak{}ThetaCorrespondence.\allowbreak{}ThetaPairLifts.\allowbreak{}LocalTypeDictionary.\allowbreak{}BottomFieldIsNumberField}{30--47}

Include the remaining coefficient fields of the representative. The normalized local values of the algebraic-type characters lie in number fields by their integral infinity types and finite conductors; the finite-order characters take values in cyclotomic fields. The finite sets of test places and depths specify finitely many orders of roots of unity and square roots of norms. For the phases, the finite support conditions give an annihilating integer on each supported term, so one cyclotomic field suffices after taking their product. For the jets, the coset expansions at the specified depth and the class representatives give finitely many entries. Each is a normalized definite-kernel coefficient on the prescribed lattice with a supported test vector. Its CM Taylor expansion has algebraic coefficients by the calculation used in Lemma~\ref{B8.17a}.
\leanref{GlobalTheory.\allowbreak{}RowExNF.\allowbreak{}charsIn\_\allowbreak{}nf}{Development.\allowbreak{}AbelianVarietyTheory.\allowbreak{}Proofs.\allowbreak{}TwistedOperators.\allowbreak{}NormalFormCharacters}{184--253}
\leanref{GlobalTheory.\allowbreak{}RowEx.\allowbreak{}nfRoot}{Development.\allowbreak{}AbelianVarietyTheory.\allowbreak{}Proofs.\allowbreak{}TwistedOperators.\allowbreak{}IsIntegralOfPowEqOne}{44--73}
\leanref{GlobalTheory.\allowbreak{}RowEx.\allowbreak{}nfPhase\_\allowbreak{}of\_\allowbreak{}sumOKK}{Development.\allowbreak{}AbelianVarietyTheory.\allowbreak{}Proofs.\allowbreak{}TwistedOperators.\allowbreak{}NormalFormPhase}{24--29}
\leanref{GlobalTheory.\allowbreak{}NFJetProof.\allowbreak{}jetCl\_\allowbreak{}of\_\allowbreak{}choice}{Development.\allowbreak{}AbelianVarietyTheory.\allowbreak{}Proofs.\allowbreak{}TwistedOperators.\allowbreak{}NormalFormJets}{40--79}

For the \(L\)-value ratios use the twisted transport assertion of Definition~\ref{B1.14k}, proved in Lemma~\ref{B1.15a}. The host provides the CM type induced from its distinguished imaginary quadratic subfield. The grid gives the prescribed base conductor, auxiliary split primes coprime to it, and angular weights in the allowed range for \(1\le s\le N_*+1\). The finite twist is trivial on principal ideles after its infinite component is removed and on the units of the indicated order, hence descends to its ideal class group. The untwisted toric sum reads the base \(L\)-value, and weighting its classes by this character reads the twisted value. Twisted transport gives algebraicity of their ratio when the denominator is nonzero. There are finitely many pairs, level indices and integers \(s\), and uniqueness of the associated \(L\)-function fixes each value. These ratios generate a finite extension. Let \(N_\kappa\subset\C\) be the compositum of all the fields constructed for this representative.
\leanref{GlobalTheory.\allowbreak{}RowEx.\allowbreak{}RatioField.\allowbreak{}inNF\_\allowbreak{}ratios}{Development.\allowbreak{}AbelianVarietyTheory.\allowbreak{}Proofs.\allowbreak{}TwistedOperators.\allowbreak{}NormalFormToricRatio}{149--170}
\leanref{GlobalTheory.\allowbreak{}RowEx.\allowbreak{}RatioField.\allowbreak{}twistRatio\_\allowbreak{}mem\_\allowbreak{}numberField}{Development.\allowbreak{}AutomorphicArithmetic.\allowbreak{}ThetaCorrespondence.\allowbreak{}ThetaPairLifts.\allowbreak{}LocalTypeDictionary.\allowbreak{}TwistRatioMemNumberField}{17--32}

Take \(S\) to be the union of the images in \(\C\) of rational bases of the \(N_\kappa\). An unrealized key contributes the empty set. Suppose now that \(\iota_X(K_X)\) contains \(S\), and take any realization in the statement. Match each of its indices to an index of the representative by the pair \((\xi_t,Q_t)\). The trivial character matches the base indices because the level character labeling is injective. With the common preliminary datum and frame, this matching transports the grid, the Weil actions and the finite tests. Complete the transported tests on the new realization. On each bottom pair the transport retains the spaces, measures and finite test functions and reindexes the two kernels by the matching of tuples. Thus the two periods chosen on the representative work on the new realization. The same reindexing preserves the local character values, the supports for the phases and jets, and the \(L\)-value ratios. Finally \(\iota_X(K_X)\) contains every rational linear combination of the chosen basis of \(N_\kappa\). It contains all the required values on this realization.
\leanref{GlobalTheory.\allowbreak{}RowEx.\allowbreak{}trBot\_\allowbreak{}holds}{Development.\allowbreak{}AbelianVarietyTheory.\allowbreak{}Proofs.\allowbreak{}TwistedOperators.\allowbreak{}TransferBottom}{40--157}
\leanref{GlobalTheory.\allowbreak{}RowEx.\allowbreak{}exKey\_\allowbreak{}of}{Development.\allowbreak{}AbelianVarietyTheory.\allowbreak{}Proofs.\allowbreak{}TwistedOperators.\allowbreak{}ExKeyOf}{23--90}
\leanref{GlobalTheory.\allowbreak{}RowEx.\allowbreak{}exField\_\allowbreak{}of\_\allowbreak{}keyFin\_\allowbreak{}exKey}{Development.\allowbreak{}AbelianVarietyTheory.\allowbreak{}Proofs.\allowbreak{}TwistedOperators.\allowbreak{}ExFieldOfKeyFinExKey}{22--37}
\end{proof}

\subsection*{Proof of \ref{B8.18}}\label{prf:B8.18}\phantomsection
\begin{proof}
At the joint threshold use the bare realization \(R\) of the stage group. Let \(H_E\) be the Hilbert-model row space of \(R\). Write \(e:H_E\to\mathcal E_0\) for its realization by bottom multiplicity maps, so that \(e(u)(w)\) is an automorphic function for \(u\in H_E\) and \(w\in\mathcal V\). A complete choice gives the real tests \(\rho,\xi\), a nonzero smooth tensor \(w_{\rm sm}\), and pairs \((\Psi_i,\Phi_i)\) of columns and grid data. Form
\[
 \ell_i(u)=\mathcal B_{\rho,\xi,\Psi_i,\Phi_i}\bigl(e(u)(w_{\rm sm})\bigr).
\]
The complete choice has the following vanishing property from \eqref{eq:p8c2-common-real-factor}. If all its tests vanish on \(u\) for the chosen \((\rho,\xi)\), they vanish on \(u\) after replacing \((\rho,\xi)\) by any real pair and keeping the finite data fixed. Theorem~\ref{B8.15} gives a nonzero pairing for \(u\ne0\), and the chosen columns span at each datum. Thus the common kernel is zero.

Let \(\iota_H:K\hookrightarrow\C\) be the Hilbert embedding. Project \(x\in FR_t\) by the idempotent for the good summand to \(x'\). Equivariance gives the same row lift. Write \(v_i(t,x')\) for the sum of phase, jet and lattice terms of the row pairing on this summand. For each layer pair compare the two models by Theorems~\ref{B8.11} and \ref{B8.16}, with scalar \(c_{i'j't}\ne0\). Common real and measure normalizations give a \(\kappa\) independent of \(t\), and finite character comparison contributes a root of unity \(z_t\in K\). Write \(j_t^H\) for realization of the rational source in the Hilbert model at \(\iota_H\), and \(\Theta_t^H\) for the lift from that model to \(H_E\). The expansion is
\begin{equation}\label{eq:p8c2-test-formula}
 \iota_H(z_t)c_t\,\ell_i\bigl(\Theta_t^H(j_t^H x)\bigr)=\kappa v_i(t,x'),
 \qquad c_t=\prod_{i',j'}c_{i'j't}\ne0 .
\end{equation}
The field condition and the local sum formulas give \(v_i(t,x')\in\iota_H(K)\). In the comparison with the base index \(\mu_0\), the chosen periods and fixed normalizations cancel. The finite character factors have their prescribed normalized values, and the \(L\)-ratios at the nonzero base factors are in the field. Hence \(c_t/c_{\mu_0}\in\iota_H(K)\). For \(\kappa\ne0\) take \(L_i=(c_{\mu_0}/\kappa)\ell_i\). Its theta value is \(v_i(t,x')/(\iota_H(z_t)(c_t/c_{\mu_0}))\). For \(\kappa=0\) take \(L_i=\ell_i\), with theta value zero. Both preserve separation. Let \(c_s\) be the identity or complex conjugation according to the stage sign, so that \(v_0=c_s\circ\iota_H\). Write \(u^{(s)}\) for the image of \(u\) in the corresponding scalar twist of the row space. Set \(\psi_i(u^{(s)})=c_s(L_i(u))\) and transfer along the common group to \(EE\).
\end{proof}

\subsection*{Proof of \ref{B8.21}}\label{prf:B8.21}\phantomsection
\begin{proof}
Use the prescribed \(G_0\) and its Weil system in the construction of Theorem~\ref{B8.17} to obtain \(N_1\). Lemma~\ref{lem:p8c2-square-field} gives \(N_2\). Level agreement transports both conditions. At fixed preliminary datum and gap there are finitely many \(w(v_0)\mid v_0\mid q_0\). The infinite types are fixed and conductors are bounded by a fixed ideal with exponents at most two. Ray class finiteness gives finite grid keys. Lemma~\ref{B8.20} gives the finite field requirement for prescribed grids at threshold \(C_{F1}\). For \(C_F\ge C_{F1}\) take \(\varepsilon_1=1/100\), zero auxiliary thresholds and the scale threshold in that finite field requirement. Use empty first three field requests. At the fourth use \(S(P,f,\Xi)\) when \(\Xi\) is finite and the empty request otherwise. Each completed stage has a finite twisting set, so prescribed grids extend there.
\end{proof}

\subsection*{Proof of \ref{B8.21a}}\label{prf:B8.21a}\phantomsection
\begin{proof}
Put \(B'=B_0/6\). Lemma~\ref{lem:p8c2-square-field} gives nonzero primitive Hecke values, algebraic \(\beta(v)\) and \(\kappa_\tau\ge1\). Apply the first assertion of Definition~\ref{B1.14k}, proved in Lemma~\ref{B1.15a}. The fitted host is abelian CM of degree \(4o\), unramified over its real subfield at finite places, with even discriminant. Here \(f=2o\ge2,\ d>4,\ d\equiv3\pmod4\) is prime. All embeddings in the CM type take the selected \(\sqrt{-d}\) to \(+i\sqrt d\), and their product on roots of unity is one. The grid conductor is proper and nonzero, with exponents at most two at odd completely split primes. Splitting of \(g_0,g_1,q_0\) follows from the congruences, and that of \(p_{\rm de},p'_{\rm de}\) from the fitted sizes. These are the field and conductor conditions of Lemma~\ref{B1.14c}. Each of the three characters is unitary, trivial on the real subfield and has the prescribed exponents and conductor. Fix the same rational partial zeta values \(z(C)=\zeta_F(1-2s,C)\) for all three applications. In each take \(Q=\varnothing\) and additional Picard character \(1\).
\leanref{GlobalTheory.\allowbreak{}SqTop.\allowbreak{}fieldDat}{Development.\allowbreak{}AbelianVarietyTheory.\allowbreak{}Proofs.\allowbreak{}DivisorProofs.\allowbreak{}NatCastMemBelow}{154--162}
\leanref{GlobalTheory.\allowbreak{}FunObj.\allowbreak{}SqNum.\allowbreak{}IsTor}{Development.\allowbreak{}AbelianVarietyTheory.\allowbreak{}Definitions.\allowbreak{}Divisors.\allowbreak{}FinitePrimesUnder}{270--272}

Put \(k_0=2sf,\ k_1=\sum_\tau\kappa_\tau\). Use \(a_l,c_l,e_l,b_l\) from the first assertion of Definition~\ref{B1.14k}. These are common to the three characters because the weights, conductor and partial zeta data are common. In the power-integrality assertion use \(j=1,\ r=12\lceil b_l\rceil\). Thus \(l^{12n_l}\beta(v)\) is integral above \(l\), where
\begin{equation}\label{eq:p8c2-toric-denominators}
 n_l=e_l+Ma_l+\lceil t_l\rceil,\qquad
 t_l=\begin{cases}
 0,&c_l=0,\\
 k_0c_l+k_1(2c_l+1+\mathbf1_{l=2}+\tfrac12\mathbf1_{l=d}),&c_l>0 .
 \end{cases}
\end{equation}
Let \(S_{\rm grid}=\{g_0,g_1,p_{\rm de},p'_{\rm de},q_0\}\) and \(N_m=N(m)\). The fitted sizes give
\begin{equation}\label{eq:p8c2-toric-sizes}
 \begin{gathered}
 \log l\le9L+1\quad(l\in S_{\rm grid}),\qquad
 \log\operatorname{rd} E\le\log\operatorname{cond} E\le L/100,\\
 s\le8N_m^2,\quad m\ge44,\quad mfN_m^2d\le b,\quad
 16\le k_0\le b,\quad \kappa_\tau\le5b,\quad L\ge1 .
 \end{gathered}
\end{equation}
Set \(K_1=k_0+5k_1+2M+1\le49fb\). At the grid primes \(a_l\le2\) and \(c_l\log l\le18L+2\). Equation \eqref{eq:p8c2-toric-denominators}, including the ceiling, costs at most \(K_1(18L+2)\) per grid prime beyond \(e_l\log l\). Elsewhere use \(c_l=v_l(|D_E|)/[E:\Q]\) and the root-discriminant and conductor sums to bound the remaining cost by \((k_0+5k_1+1)L/100\). The class zeta calculation of Lemma~\ref{B1.15a} gives \(\sum_l e_l\log l\le k_0\log(k_0+1)\) for \(k_0\ge16\). Thus on every finite set of rational primes
\begin{equation}\label{eq:p8c2-toric-den-sum}
 \sum_l n_l\log l\le b(L+1)+8K_1(18L+2)\le8000fb(L+\log x)\le B'
\end{equation}
provided \(C_1/6\ge8000\).

For each transported character put \(Z=\zeta(3/2)^{[E:\Q]}\), \(A=|D_E|N\mathfrak f_\chi\), with conductor \(\mathfrak f_\chi\). The Euler bound is \(|L(s,\chi)|\le Z\) for \(s\ge2\). At one use the Hecke completion in the strip from \(-1/2\) to \(3/2\), with right boundary bounded by the gamma integral and left boundary by the Hecke--Tate functional equation \cite{tate-1967}. Nontrivial angular type gives an entire completion bounded on closed strips by its theta Mellin integral. The strip bound and \(\Gamma(t+1/2)/\Gamma(t)\le\sqrt t\) give
\begin{equation}\label{eq:p8c2-toric-L-bound}
 |L(1,\chi)|\le ZA\prod_\tau(1+s+\kappa_\tau)\qquad(s=1).
\end{equation}
Let \(Q_{\rm lev}=\prod_{v\in S_0}(Nv)^{a_v}\), where \(a_{v_0}=1\) at the base place and \(a_v=2\) at every other grid prime. The conductor exponents give \(N\mathfrak f_\chi\le Q_{\rm lev}^2\). We have \(\log|D_E|\le2fL\), \(\log Q_{\rm lev}\le2f(21L+2)\), and \(1+s+\kappa_\tau\le7b\). Thus the Euler and strip bounds give
\begin{equation}\label{eq:p8c2-toric-log-L}
 \log|L(s,\chi)|\le101fL.
\end{equation}
\leanref{GlobalTheory.\allowbreak{}HeckeStrip.\allowbreak{}completedLBoundedInStrips}{Development.\allowbreak{}AbelianVarietyTheory.\allowbreak{}Proofs.\allowbreak{}DivisorProofs.\allowbreak{}CompletedHeckeLBoundedInVerticalStrips}{17--54}

Apply ordinary transport of Definition~\ref{B1.14k} and Lemma~\ref{B1.15a}. For an embedding fixing \(K_d\), use the resulting ideal \(\mathfrak a\), period lattice \(\Lambda_{\mathfrak a}\) and character \(\chi'\) with permuted weights. Then
\begin{equation}\label{eq:p8c2-cm-transport}
 \log|\sigma\beta(v)|=12\log C_0+12\log|L(s,\chi')|
 -M\log(N\mathfrak a^6|\operatorname{disc}(\Lambda_{\mathfrak a})|).
\end{equation}
Complex conjugation preserves log absolute values and pairs the embeddings over the two embeddings of \(K_d\). The class action of Definition~\ref{B1.14k} is equidistributed. For \(h(K_d)\) the class number, the logarithmic Chowla--Selberg formula for discriminant \(-d<-4\) gives \cite{selberg-chowla-1967}
\begin{equation}\label{eq:p8c2-cm-class-mean}
 \frac1{h(K_d)}\sum_{[\mathfrak a]\in\operatorname{Cl}(K_d)}
 \log(N\mathfrak a^6|\operatorname{disc}(\Lambda_{\mathfrak a})|)
 =12\log(2\pi)-6\log d+6\gamma-6D_1,\quad
 D_1=\frac{L'(1,\chi_{-d})}{L(1,\chi_{-d})}.
\end{equation}
Here \(\chi_{-d}\) is the quadratic character of \(K_d\). The mean in \eqref{eq:p8c2-cm-transport} divided by twelve is at most \(101fL+\log C_0+M(\log d+|D_1|)/2\). Use \(\log|D_F|\le fL,\ \log d,|D_1|\le L,\ M\le12fb,\ 2s+\kappa_\tau\le7b\). Use factorials for the gamma factors to bound the three terms of \(\log C_0\) by \(2fbL,\ 49fbL,\ 40fbL\). The final \(M\)-term costs at most \(12fbL\). The mean is therefore at most \(204fbL\). Hence
\begin{equation}\label{eq:p8c2-toric-arch-mean}
 \log|P_{\Q,v}(0)|/(12\deg P_{\Q,v})\le B',
 \qquad P_{\Q,v}=\min_{\Q}(\beta(v))\ \text{monic}.
\end{equation}
\leanref{GlobalTheory.\allowbreak{}GalMean.\allowbreak{}galmean}{Development.\allowbreak{}AbelianVarietyTheory.\allowbreak{}Proofs.\allowbreak{}DivisorProofs.\allowbreak{}MeanLogNormConjugatesOverIdealClasses}{187--257}
\leanref{GlobalTheory.\allowbreak{}T.\allowbreak{}classMean\_\allowbreak{}log\_\allowbreak{}norm\_\allowbreak{}disc\_\allowbreak{}chowlaSelberg}{Development.\allowbreak{}AbelianVarietyTheory.\allowbreak{}Proofs.\allowbreak{}DivisorProofs.\allowbreak{}ClassMeanLogNormDiscChowlaSelbergStatement}{9--11}

The norm to \(K\) gives \(y_w(v)\ge-n_lv_w(l)\) for \(w\mid l\). Using \(\sum_{w\mid l}\ell_wv_w(l)=\log l\), \eqref{eq:p8c2-toric-den-sum} bounds the weighted negative part by \(B'\). The product formula and transitivity of norms identify the weighted signed sum with the left side of \eqref{eq:p8c2-toric-arch-mean}. Therefore \(\sum_w\ell_w|y_w(v)|\le3B'\), with finite support. The identity \(\beta(u_c)=z_c^{12}\beta(u)\) gives \(y_w(u_c)=v_w(z_c)+y_w(u)\). Use \(|n+\operatorname{trunc}_0t|\le|n+t|+|t|\) to sum the two \(3B'\) bounds. For \(z_T=1\) use \(|\operatorname{trunc}_0t|\le|t|\); outside range each term is zero.
\end{proof}

\subsection*{Proof of \ref{B8.21fa}}\label{prf:B8.21fa}\phantomsection
\begin{proof}
After \(j\) top steps the loss is at most \(j\delta_T(w)\). After \(j\) row steps it is at most \(\delta_F(w)+j\delta_R(w)\). The terminal contraction therefore has loss at most
\[
 m\delta_T(w)+\delta_F(w)+m\delta_R(w)+\delta_P(w).
\]
Multiplication by the four scalar factors adds their four losses. The nonarchimedean triangle inequality preserves this common bound under each of the finite sums. It follows that every elementary test pairing, and hence the pairing with every element of \(L_{\rm form}\), has valuation at least \(-\delta(w)\) when nonzero.

Set \(r=\lfloor\delta(w)\rfloor\). Since valuations are integral, local duality for the full lattice gives \(y\in\mathfrak p_w^{-r}M_w\). Choose \(\pi\in K\) with \(v_w(\pi)=1\). Then \(\pi^ry\in M_w\), so some \(u\in\mathcal O_K\setminus\mathfrak p_w\) satisfies \(u\pi^ry\in M\). Taking \(s_w=u\pi^r\) gives \(v_w(s_w)=r\le\delta(w)\).

The support of \(\delta\) lies in the union of the eight finite supports. Adding the weighted bounds gives
\[
 \sum_w\ell_w\delta(w)
 \le4m\mathcal B+2m\mathcal B+2\mathcal B
 =(6m+2)\mathcal B\le8m\mathcal B.
\]
\end{proof}

\subsection*{Proof of \ref{B8.21f}}\label{prf:B8.21f}\phantomsection
\begin{proof}[Proof (outline)]
With \(\sigma(G)\) from the statement, put
\[
 B_1(C,b)=CfmN_*^2b\bigl(\log b+\sigma(G)\bigr).
\]
The depth \(d_{\rm pl}(q)\) is the prescribed depth of Definition~\ref{B8.21c}. For a prime \(\ell\), put
\[
 \Delta_G(\ell)=
 \max_{\substack{\mathfrak P\in W_0\\\mathfrak P\mid\ell}}
 d_{\rm pl}(\mathfrak P\cap\mathcal O_F),
\]
with value zero for an empty set. The local expansion theorem gives a finite supported term set and finite class sets. For every supported term \(\beta\) it gives a positive integer \(k_{t,\Phi,\beta}\) and a finite expansion into pure coset data at that modulus. At a place \(w\mid\ell\) satisfying the condition in the statement,
\[
 v_\ell(k_{t,\Phi,\beta})\le\Delta_G(\ell).
\]
The coefficient of each pure datum is the product of the corresponding local datum values. Thus its class pairing has the expansion in Lemma~\ref{B8.21fa}.
\leanref{GlobalTheory.\allowbreak{}SumSize.\allowbreak{}DepthAt}{Development.\allowbreak{}AbelianVarietyTheory.\allowbreak{}Statements.\allowbreak{}PlaceShiftStatements.\allowbreak{}InfClassAt}{97--101}
\leanref{GlobalTheory.\allowbreak{}SumSize.\allowbreak{}depL}{Development.\allowbreak{}AbelianVarietyTheory.\allowbreak{}Statements.\allowbreak{}PlaceShiftStatements.\allowbreak{}InfClassAt}{57--58}
\leanref{GlobalTheory.\allowbreak{}RowTest.\allowbreak{}IsCosetExpansion}{Development.\allowbreak{}AbelianVarietyTheory.\allowbreak{}Definitions.\allowbreak{}TwistedInnerProducts.\allowbreak{}MemHalfCosetSelf}{122--124}
\leanref{DenominatorLossAdditivity.\allowbreak{}GlobalTheory.\allowbreak{}CompanionTheory.\allowbreak{}FinDen.\allowbreak{}core\_\allowbreak{}loss}{Development.\allowbreak{}AbelianVarietyTheory.\allowbreak{}Proofs.\allowbreak{}DivisorProofs.\allowbreak{}FiniteSumDenominatorBoundOfFactorClaims}{280--339}

Under \(\mathscr P\), choose the common grid threshold for the finite support, depth and local factor estimates. At a fixed \(C_F\) above that threshold, let \(c_1,\ldots,c_{10}\) be their ten estimate thresholds and take
\[
 C_0=16\max(c_1,\ldots,c_{10}).
\]
For \(C_1\ge C_0\), apply each estimate at \(C_1/16\) and take the maximum of the resulting cutoffs in \(x\). Put \(\mathcal B=B_1(C_1/16,b)\). The scalar, coefficient, finite datum and stabilizer estimates give \(\delta_S,\delta_C,\delta_D,\delta_G\) with weighted bound \(m\mathcal B\) each. The top and row kernel estimates give \(\delta_T,\delta_R\) with bound \(\mathcal B\) for each step. The initial form and terminal contraction give \(\delta_F,\delta_P\) with bound \(\mathcal B\), and the initial test-vector value has loss zero. These bounds imply \(\mathcal B\ge0\). These assertions are uniform in the level index and the supported term and class indices, and give exactly the local assertions of Lemma~\ref{B8.21fa}.
\leanref{GlobalTheory.\allowbreak{}SumSize.\allowbreak{}FinScalAt}{Development.\allowbreak{}AbelianVarietyTheory.\allowbreak{}Statements.\allowbreak{}PlaceShiftStatements.\allowbreak{}InfClassAt}{102--103}
\leanref{GlobalTheory.\allowbreak{}SumSize.\allowbreak{}FinCoefAt}{Development.\allowbreak{}AbelianVarietyTheory.\allowbreak{}Statements.\allowbreak{}PlaceShiftStatements.\allowbreak{}InfClassAt}{104--106}
\leanref{GlobalTheory.\allowbreak{}SumSize.\allowbreak{}FinDatAt}{Development.\allowbreak{}AbelianVarietyTheory.\allowbreak{}Statements.\allowbreak{}PlaceShiftStatements.\allowbreak{}InfClassAt}{107--109}
\leanref{GlobalTheory.\allowbreak{}SumSize.\allowbreak{}FinGamAt}{Development.\allowbreak{}AbelianVarietyTheory.\allowbreak{}Statements.\allowbreak{}PlaceShiftStatements.\allowbreak{}InfClassAt}{110--112}
\leanref{GlobalTheory.\allowbreak{}SumSize.\allowbreak{}FinTopAt}{Development.\allowbreak{}AbelianVarietyTheory.\allowbreak{}Statements.\allowbreak{}PlaceShiftStatements.\allowbreak{}InfClassAt}{113--117}
\leanref{GlobalTheory.\allowbreak{}SumSize.\allowbreak{}FinRowAt}{Development.\allowbreak{}AbelianVarietyTheory.\allowbreak{}Statements.\allowbreak{}PlaceShiftStatements.\allowbreak{}InfClassAt}{118--123}
\leanref{GlobalTheory.\allowbreak{}SumSize.\allowbreak{}FinPairAt}{Development.\allowbreak{}AbelianVarietyTheory.\allowbreak{}Statements.\allowbreak{}PlaceShiftStatements.\allowbreak{}InfClassAt}{124--127}
\leanref{GlobalTheory.\allowbreak{}SumSize.\allowbreak{}FinPsiAt}{Development.\allowbreak{}AbelianVarietyTheory.\allowbreak{}Statements.\allowbreak{}PlaceShiftStatements.\allowbreak{}InfClassAt}{128--130}
\leanref{GlobalTheory.\allowbreak{}SumSize.\allowbreak{}FinFormAt}{Development.\allowbreak{}AbelianVarietyTheory.\allowbreak{}Statements.\allowbreak{}PlaceShiftStatements.\allowbreak{}InfClassAt}{131--134}
\leanref{DenominatorLossAdditivity.\allowbreak{}GlobalTheory.\allowbreak{}CompanionTheory.\allowbreak{}FinDen.\allowbreak{}fin\_\allowbreak{}of\_\allowbreak{}factors}{Development.\allowbreak{}AbelianVarietyTheory.\allowbreak{}Proofs.\allowbreak{}DivisorProofs.\allowbreak{}FiniteSumDenominatorBoundOfFactorClaims}{535--575}

For the tied test lattice, apply Lemma~\ref{B8.21fa} to the evaluation vector \(y\). Let \(\delta_{\rm sum}\) be the sum of the eight losses in that lemma and let \(T_{\rm sum}\) be its support. Since \(B_1\) is linear in its first argument,
\[
 \begin{split}
 2\sum_w\ell_w\delta_{\rm sum}(w)
 &\le16mB_1(C_1/16,b)=mB_1(C_1,b)\\
 &=C_1fbm^2N_*^2\bigl(L+\sigma(G)\bigr).
 \end{split}
\]
The same lemma gives the required scalar \(s_0\). Outside \(T_{\rm sum}\), its valuation is at most zero, and \(s_0y\in M\) implies \(y\in M_w\).
\leanref{GlobalTheory.\allowbreak{}SumSize.\allowbreak{}bOne}{Development.\allowbreak{}AbelianVarietyTheory.\allowbreak{}Statements.\allowbreak{}PlaceShiftStatements.\allowbreak{}InfClassAt}{48--49}
\leanref{DenominatorLossAdditivity.\allowbreak{}GlobalTheory.\allowbreak{}CompanionTheory.\allowbreak{}FinDen.\allowbreak{}fin\_\allowbreak{}at}{Development.\allowbreak{}AbelianVarietyTheory.\allowbreak{}Proofs.\allowbreak{}DivisorProofs.\allowbreak{}FiniteSumDenominatorBoundOfFactorClaims}{507--532}
The complete proof is given in \Lean{GlobalTheory.\allowbreak{}SumSize.\allowbreak{}Join.\allowbreak{}sum\_\allowbreak{}size\_\allowbreak{}finite\_\allowbreak{}away\_\allowbreak{}archimedean\_\allowbreak{}combined}.
\end{proof}

\subsection*{Proof of \ref{B8.21g}}\label{prf:B8.21g}\phantomsection
\begin{proof}[Proof (outline)]
Choose the grid constant above the finite-sum and place-avoidance thresholds. In the closed family construction, choose \(C_0\) above both corresponding estimate thresholds and above \(\max(C_F,0)\). At each fitted point beyond the common cutoff, the row-size and series conditions give the Hida row setting. The place-avoidance assertion says that, for each finite place \(w\) of \(K\) and every \(\mathfrak P\in Q_{\mu(w)}\), either \(\mathfrak P\cap\mathcal O_F\in S_0\), or the rational residue prime of \(\mathfrak P\) differs from that of \(w\). Thus Lemma~\ref{B8.21f} applies at \(t=\mu(w)\).
\leanref{GlobalTheory.\allowbreak{}SumSize.\allowbreak{}Away}{Development.\allowbreak{}AbelianVarietyTheory.\allowbreak{}Statements.\allowbreak{}PlaceShiftStatements.\allowbreak{}InfClassAt}{52--54}
\leanref{GlobalTheory.\allowbreak{}SumSize.\allowbreak{}AwayAt}{Development.\allowbreak{}AbelianVarietyTheory.\allowbreak{}Statements.\allowbreak{}PlaceShiftStatements.\allowbreak{}InfClassAt}{139--140}
\leanref{GlobalTheory.\allowbreak{}Adopt.\allowbreak{}twisted\_\allowbreak{}local\_\allowbreak{}factor\_\allowbreak{}away\_\allowbreak{}from\_\allowbreak{}level}{Development.\allowbreak{}AutomorphicArithmetic.\allowbreak{}ThetaCorrespondence.\allowbreak{}ThetaPairLifts.\allowbreak{}LocalTypeDictionary.\allowbreak{}TwistedLocalFactorAwayFromLevel}{51--52}
\leanref{LeafFixedPointPairing.\allowbreak{}GlobalTheory.\allowbreak{}CompanionTheory.\allowbreak{}LeafFixedPointPlaceNorm.\allowbreak{}fp\_\allowbreak{}tie\_\allowbreak{}alt}{Development.\allowbreak{}AutomorphicArithmetic.\allowbreak{}ThetaCorrespondence.\allowbreak{}ThetaPairLifts.\allowbreak{}ThirdPlaceSeries.\allowbreak{}LeafFixedPointOfFinKAwayK}{23--46}

Let \(\delta_{\rm sum}\) be its allowance, put \(T_F=\operatorname{supp}(\delta_{\rm sum})\), and set \(d_F(w)=\lfloor\delta_{\rm sum}(w)\rfloor\). At a member point the source lattice in Lemma~\ref{B8.21f} is \(M_t\), by Definition~\ref{B8.21d}(1). The evaluation tie therefore gives, for each \(x\in M_{\mu(w)}\), a nonzero \(s\) with
\[
 v_w(s)\le\delta_{\rm sum}(w),\qquad
 sV_{\mu(w)}^0x\in M.
\]
Since \(v_w(s)\in\Z\), it is at most \(d_F(w)\). Localization gives
\[
 V_{\mu(w)}^0(M_{\mu(w),w})\subseteq\mathfrak p_w^{-d_F(w)}M_w.
\]
The function \(d_F\) vanishes off \(T_F\).
\leanref{LeafFixedPointPairing.\allowbreak{}GlobalTheory.\allowbreak{}CompanionTheory.\allowbreak{}LeafFixedPointPlaceNorm.\allowbreak{}local\_\allowbreak{}clause}{Development.\allowbreak{}AbelianVarietyTheory.\allowbreak{}Proofs.\allowbreak{}DivisorProofs.\allowbreak{}FiniteAndAwayFormula}{71--94}
\leanref{LeafFixedPointPairing.\allowbreak{}GlobalTheory.\allowbreak{}CompanionTheory.\allowbreak{}LeafFixedPointPlaceNorm.\allowbreak{}leaf\_\allowbreak{}at}{Development.\allowbreak{}AbelianVarietyTheory.\allowbreak{}Proofs.\allowbreak{}DivisorProofs.\allowbreak{}FiniteAndAwayFormula}{125--137}

At the fitted grid, its depth sum is \(\sigma(G)=\Sigma_0(C_F,b)\). The base depth is monotone in its constant, so \(C_F\le C_1\) gives \(\sigma(G)\le\Sigma_0(C_1,b)\). Since \(C_1\ge0\), \(L\ge0\), and \(\Sigma_0\ge0\),
\[
 \begin{split}
 2\sum_{w\in T_F}\ell_wd_F(w)
 &\le2\sum_{w\in T_F}\ell_w\delta_{\rm sum}(w)\\
 &\le C_1fbm^2N_*^2\bigl(L+\sigma(G)\bigr)\\
 &\le C_1fbm^2N_*^2(L+\Sigma_0)\\
 &\le C_1fbm^2(N_*+N_*^2)(L+\Sigma_0)=B_F.
 \end{split}
\]
\leanref{LeafFixedPointPairing.\allowbreak{}GlobalTheory.\allowbreak{}CompanionTheory.\allowbreak{}LeafFixedPointPlaceNorm.\allowbreak{}sig\_\allowbreak{}le}{Development.\allowbreak{}AbelianVarietyTheory.\allowbreak{}Proofs.\allowbreak{}DivisorProofs.\allowbreak{}FiniteAndAwayFormula}{63--70}
\leanref{GlobalTheory.\allowbreak{}FinitePlacesProbe.\allowbreak{}sigma\_\allowbreak{}mono\_\allowbreak{}in\_\allowbreak{}constant}{Development.\allowbreak{}AbelianVarietyTheory.\allowbreak{}Counterexamples.\allowbreak{}RotationExamples.\allowbreak{}SigmaMonoInConstant}{31--42}
\leanref{LeafFixedPointPairing.\allowbreak{}GlobalTheory.\allowbreak{}CompanionTheory.\allowbreak{}LeafFixedPointPlaceNorm.\allowbreak{}sum\_\allowbreak{}clause}{Development.\allowbreak{}AbelianVarietyTheory.\allowbreak{}Proofs.\allowbreak{}DivisorProofs.\allowbreak{}FiniteAndAwayFormula}{95--124}

Finally, the inverse of a root of unity is integral over \(\mathcal O_K\). Multiplication by \(\zeta_t^{-1}\) preserves every localized \(\mathcal O_K\)-lattice and therefore preserves the inclusion and its unchanged allowance.
\leanref{GlobalTheory.\allowbreak{}finitePlaceBound\_\allowbreak{}inv\_\allowbreak{}rootOfUnity\_\allowbreak{}smul}{Development.\allowbreak{}AbelianVarietyTheory.\allowbreak{}Proofs.\allowbreak{}DivisorProofs.\allowbreak{}MatLatSmulLeOfIsIntegral}{24--33}
The complete proof is given in \Lean{LeafFixedPointPairing.\allowbreak{}GlobalTheory.\allowbreak{}CompanionTheory.\allowbreak{}LeafFixedPointPlaceNorm.\allowbreak{}fp\_\allowbreak{}tie\_\allowbreak{}alt}.
\end{proof}

\subsection*{Proof of \ref{B8.21h}}\label{prf:B8.21h}\phantomsection
\begin{proof}[Proof (outline)]
Choose \(C_F\) above the threshold for the combined finite and archimedean layer-sum estimates. Choose \(C_0\) above the corresponding archimedean constant and above \(\max(C_F,0)\), and fix \(C_1\ge C_0\). Pass above both the Hida row-setting cutoff and the archimedean-estimate cutoff.

The source-frame estimate used at a fitted point has the following quantitative form. Let \((M,X')\) be the tied test lattice, let \(X_v\) be invertible with
\[
 (X_v^{-1})^\dagger X_v^{-1}=H_V(\mu_0)_v,
\]
and let \(V\) be an evaluation matrix satisfying
\[
 \iota((Vx)_\alpha)=
 v_{\rm row}(a,\mu_0,\mathrm{form}_{F_{\rm fc}}(\mu_0,x))_\alpha
 \quad\text{for every source vector }x\text{ and coordinate }\alpha.
\]
The archimedean layer-sum estimate and its grid-budget comparison give
\[
 \begin{split}
 h_{X,X'}(V)
 &\le m\bigl(C_1fN_*b\log b+C_1mN_*^2\Sigma_0(C_1,b)\bigr)\\
 &=C_1fmN_*bL+C_1m^2N_*^2\Sigma_0(C_1,b)=H_{I,0}.
 \end{split}
\]
The estimate applies to every compatible source frame and every matrix with the displayed evaluation identity. Its two terms retain the factors \(m\) and \(m^2\) before passage to the quotient metric.
\leanref{GlobalTheory.\allowbreak{}SumSize.\allowbreak{}InfAt}{Development.\allowbreak{}AbelianVarietyTheory.\allowbreak{}Statements.\allowbreak{}PlaceShiftStatements.\allowbreak{}InfClassAt}{176--180}
\leanref{GlobalTheory.\allowbreak{}SumSize.\allowbreak{}Join.\allowbreak{}sum\_\allowbreak{}size\_\allowbreak{}finite\_\allowbreak{}away\_\allowbreak{}archimedean\_\allowbreak{}combined}{Development.\allowbreak{}AutomorphicArithmetic.\allowbreak{}ThetaCorrespondence.\allowbreak{}ThetaPairLifts.\allowbreak{}LocalTypeDictionary.\allowbreak{}SumSizeFiniteAwayArchimedeanCombined}{26--44}
\leanref{GlobalTheory.\allowbreak{}QuotientNormCoisometry.\allowbreak{}bound\_\allowbreak{}le}{Development.\allowbreak{}AbelianVarietyTheory.\allowbreak{}Proofs.\allowbreak{}DivisorProofs.\allowbreak{}InfinitePartIdentity}{95--107}

Take \(X_v=H_V(\mu_0)_v^{-1/2}\) and \(V=V_{\mu_0}^0\). Positive definiteness gives the required frame identity, and the fitted evaluation tie gives the displayed identity for \(V\). The map \(P_{\mu_0}\) is surjective and its target carries the direct quotient metric. Lemma~\ref{B8.21ha}, applied to \(B=V_b\), therefore gives
\[
 h(V_b)=h_{X,X'}(V_bP_{\mu_0})
       =h_{X,X'}(V_{\mu_0}^0)\le H_{I,0}.
\]
\leanref{Formalism.\allowbreak{}L2.\allowbreak{}DIV.\allowbreak{}frameOfSpec}{Formalism.\allowbreak{}LocalToGlobal.\allowbreak{}Divisibility.\allowbreak{}FramesLemmas}{77--90}
\leanref{GlobalTheory.\allowbreak{}QuotientNormCoisometry.\allowbreak{}htMap\_\allowbreak{}quot}{Development.\allowbreak{}AbelianVarietyTheory.\allowbreak{}Proofs.\allowbreak{}DivisorProofs.\allowbreak{}InfinitePartIdentity}{64--84}
The complete proof is given in \Lean{GlobalTheory.\allowbreak{}QuotientNormCoisometry.\allowbreak{}leafHeightMapBound\_\allowbreak{}tie\_\allowbreak{}of\_\allowbreak{}infK}.
\end{proof}

\subsection*{Proof of \ref{B8.21i}}\label{prf:B8.21i}\phantomsection
\begin{proof}[Proof (outline)]
The family row-size inequalities and the series condition of the tied complete choice \(a\) of \(d\) give the Hida row-setting condition at \(p\in\mathcal P\). By the test-lattice tie of Definition~\ref{B8.21d}(3), \((M,(X'_v)_{v\mid\infty})\) is the direct-sum Hermitian test \(\mathcal O_K\)-lattice in \(K^n\) indexed along \(e_a:\{1,\dots,n\}\simeq\mathcal I_{\rm test}(G)\times\mathcal I_{\rm coef}(G)\) on the initial top node \(X_0\) of rank \(n_p=k+m+m(k+m+4)\). For \(C_1\ge1\), Definition~\ref{B8.21e} gives \(R_p(C_1)\ge1\) and bounds the coordinate spread of the dominant weight \(\lambda_{P,G}\) across \(\tau\in\mathcal T\) (\(|\mathcal T|=f\)) by \((m+1)(\gamma_{\rm gap}(C_1)+1)+L_c+n_p=R_p(C_1)\), whence nondegeneracy of the constituent frames makes each \(\det X'_v\) a unit, \(M\) finitely generated and nonzero, and \(\mu_{\max}(M)\le4fn_pR_p\log(n_pR_p)=E_M\).
\leanref{GlobalTheory.\allowbreak{}BigTop.\allowbreak{}Fun.\allowbreak{}LeafMMTie.\allowbreak{}leafMM\_\allowbreak{}of\_\allowbreak{}test}{Development.\allowbreak{}AbelianVarietyTheory.\allowbreak{}Proofs.\allowbreak{}DivisorProofs.\allowbreak{}MultiplicityAtFittedPoints}{94--139}

The complete proof is given in \Lean{GlobalTheory.\allowbreak{}BigTop.\allowbreak{}Fun.\allowbreak{}LeafMMTie.\allowbreak{}mm\_\allowbreak{}tie\_\allowbreak{}alt}.
\end{proof}

\subsection*{Proof of \ref{B8.21ja}}\label{prf:B8.21ja}\phantomsection
\begin{proof}
Every index in the two tower products lies in \([1,h]\): their largest index is \(e+q=e_0+p\le h\), and the smallest index in the product of length \(p\) is \(e+q-p+1=e_0+1\). The right-hand side of the tower identity is nonzero, so \(R\ne0\). Put \(x_s=v(\rho_s)\), \(A_1=\sum_{s=1}^h a_s\), and \(N_1=\sum_{s=1}^h n_s\). Taking valuations and using the bound for \(\gamma\) gives
\[
 \begin{split}
 v(R)\le{}&v(A)+v(A')+6d+(p+q)B\\
 &+\sum_{s=1-e_0}^{h_0}x_s
  +\sum_{k=0}^{p-1}x_{e+q-k}
  -\sum_{k=0}^{q-1}x_{e+k+1}.
 \end{split}
\]
Reflection and the positive upper bound give
\[
 \sum_{s=1-e_0}^{h_0}x_s
 =\sum_{s=1}^{e_0}x_s+\sum_{s=1}^{h_0}x_s
 \le2A_1+(e_0+h_0)\eta.
\]
Each tower product has distinct indices in \([1,h]\), so
\[
 \sum_{k=0}^{p-1}x_{e+q-k}\le A_1+p\eta,
 \qquad
 -\sum_{k=0}^{q-1}x_{e+k+1}\le A_1+N_1.
\]
Thus the three displayed sums contribute at most
\[
 4A_1+N_1+(e_0+h_0+p)\eta
 \le4A_1+5N_1+2h\eta,
\]
since \(N_1,\eta\ge0\) and \(e_0+h_0+p\le2h\). Finally use \(v(A)\le|v(A)|\) and \(v(A')\le|v(A')|\).
\end{proof}

\subsection*{Proof of \ref{B8.21j}}\label{prf:B8.21j}\phantomsection
\begin{proof}[Proof (outline)]
The host parameters of every family member satisfy
\(\bigl|\Re(L'(1,\chi_d)/L(1,\chi_d))\bigr|\le x^\varepsilon\).
The family logarithmic-derivative theorem places this bound in the closed quantifiers of Definition~\ref{B8.39}. Together with the row-size and maximal-real-subfield conditions, it gives the square-recursive size datum. Take maxima of the finitely many cutoffs for the Hecke nonvanishing, toric denominator, bottom-pair, reflection and tower estimates below. This gives a common \(C_{F0}\), and, after \(C_F\ge C_{F0}\) is fixed, a common \(C_0\) and an eventual threshold. The estimates hold simultaneously for every fitted datum and every \(i,j,w\).
\leanref{GlobalTheory.\allowbreak{}logDeriv\_\allowbreak{}LFunction\_\allowbreak{}at\_\allowbreak{}one\_\allowbreak{}le\_\allowbreak{}along\_\allowbreak{}family\_\allowbreak{}alt}{Development.\allowbreak{}AutomorphicArithmetic.\allowbreak{}ThetaCorrespondence.\allowbreak{}ThetaPairLifts.\allowbreak{}ThirdPlaceSeries.\allowbreak{}LogDerivLFunctionAtOneLeAlongFamily}{12--17}
\leanref{GlobalTheory.\allowbreak{}SqTop.\allowbreak{}GBig}{Development.\allowbreak{}AbelianVarietyTheory.\allowbreak{}Statements.\allowbreak{}PlaceShiftStatements.\allowbreak{}NormSmall}{34--37}
\leanref{GlobalTheory.\allowbreak{}SqTop.\allowbreak{}Model.\allowbreak{}gBig\_\allowbreak{}and}{Development.\allowbreak{}AbelianVarietyTheory.\allowbreak{}Statements.\allowbreak{}PlaceShiftStatements.\allowbreak{}AuxiliaryPrimeNeTypePrime}{52--62}
\leanref{GlobalTheory.\allowbreak{}SqTop.\allowbreak{}square\_\allowbreak{}top\_\allowbreak{}term\_\allowbreak{}joined\_\allowbreak{}estimates}{Development.\allowbreak{}AutomorphicArithmetic.\allowbreak{}ThetaCorrespondence.\allowbreak{}ThetaPairLifts.\allowbreak{}LocalTypeDictionary.\allowbreak{}SquareTopTermJoinedEstimates}{13--35}

The CM theta input package is available for every arithmetic datum, so the bottom-pair and coordinate norm estimates apply to this member.
\leanref{GlobalTheory.\allowbreak{}NormSum.\allowbreak{}inp\_\allowbreak{}all}{Development.\allowbreak{}AbelianVarietyTheory.\allowbreak{}Proofs.\allowbreak{}FoundationProofs.\allowbreak{}CmThetaNormBoundHypothesesSatisfied}{29--52}

Fix such a member point and fitted datum, and fix \(i,j,w\). Suppress \(i,j\) in the notation. Put
\[
 \begin{gathered}
 h=h_{ij},\qquad y=m-1-j,\qquad
 p_*=\lfloor y/2\rfloor,\qquad q_*=y-p_*,\\
 e=\frac{r_{i+1,j+1}-r_{ij}}2,\qquad
 e_0=e+q_*-p_*,\qquad h_0=h-y.
 \end{gathered}
\]
The rank formulas and the parity assumptions give
\[
 e_0+p_*=e+q_*,\quad h=h_0+p_*+q_*,\quad
 1\le e_0<h_0,\quad e+q_*\le h\le N_*.
\]
The two bottom ranks are \(h_0+e_0\) and \(h_0-e_0\). For \(1\le s\le h\), define
\[
 \begin{split}
 a_s&=|v_w(z_T(s))+g_T(s,w)|
      +|v_w(z_{T,c(w)}(s))+g_T(s,w)|,\\
 n_s&=\sum_{l\in S_d}N_0(s,l)v_w(l),\qquad
 \eta_w=\sum_{l\in S_e}\frac{v_w(l)}{l^2},\qquad
 C_w=\sum_{q\in S_0}c_qv_w(N\mathfrak q).
 \end{split}
\]
Here \(N\mathfrak q\) is the absolute norm denoted by \(q\) in the statement. The fitted denominator and integrality data give \(n_s,\eta_w,d_\ell(w),C_w\ge0\). If \(\mu(w)=\mu_0\), the fitted normalization gives \(r_{ij,\mu(w)}=1\), and the claimed valuation inequality follows from these nonnegative terms.
\leanref{GlobalTheory.\allowbreak{}SqTop.\allowbreak{}Model.\allowbreak{}tower\_\allowbreak{}numbers}{Development.\allowbreak{}AbelianVarietyTheory.\allowbreak{}Statements.\allowbreak{}PlaceShiftStatements.\allowbreak{}AuxiliaryPrimeNeTypePrime}{63--83}

Suppose now that \(\tau=\mu(w)\ne\mu_0\). Let \(Q_{\rm tw}=\prod_{\mathfrak P\in Q_\tau}N\mathfrak P\), where \(Q_\tau\) is the twist's finite ramification-place set. The fitted twisting condition gives \(w\nmid Q_{\rm tw}\). Let \(l\) be the residue prime of \(w\). Since \(S_d,S_e\) consist of rational primes, and \(l\in S_e\),
\[
 n_s=N_0(s,l)v_w(l),\qquad \eta_w=\frac{v_w(l)}{l^2}.
\]
For the Hecke characters \(\nu_t=\mu_{0,j}\xi_t/\chi_{m-i}\), choose their Hecke functions \(\Lambda_t\). The functional equation and the square-recursive size conditions give \(\Lambda_\tau(s)\Lambda_{\mu_0}(s)\ne0\) for every integer \(1-e_0\le s\le h\). Write \(\Lambda_0=\Lambda_{\mu_0}\). The fitted field condition and the toric-value identities give, for \(1\le s\le h\), nonzero \(z_\psi(s),z_T(s),z_{T,c(w)}(s)\in K\) such that
\[
 \iota(z_\psi(s))\Lambda_0(s)=\Lambda_\tau(s),\qquad z_T(s)=1,
 \qquad
 \iota(z_{T,c(w)}(s))\Lambda_0(s)=\Lambda_d(s),
\]
where \(\Lambda_d\) is the Hecke function of \(\nu_{\mu_0}\delta_{c(w)}\). The selected prime twist \(\delta_{c(w)}\) is the twist attached to \(l\). The toric denominator estimate and the prime-to-\(l\) twisting estimate, applied to \(z_\psi(s)/z_{T,c(w)}(s)\), give
\[
 -n_s\le v_w(z_\psi(s))+g_T(s,w),\qquad
 v_w(z_\psi(s))\le v_w(z_{T,c(w)}(s))+\eta_w.
\]
Reflection gives nonzero \(E_s\in K\) for \(1-e_0\le s\le0\), with \(v_w(E_s)=0\), such that
\[
 \Lambda_\tau(s)\Lambda_0(1-s)
 =\iota(E_s)\Lambda_\tau(1-s)\Lambda_0(s).
\]
Define
\[
 \rho_s=
 \begin{cases}
 z_\psi(s)/z_T(s),&1\le s\le h,\\
 E_s\,z_\psi(1-s)/z_T(1-s),&1-e_0\le s\le0.
 \end{cases}
\]
Then \(\rho_s\ne0\), \(\iota(\rho_s)\Lambda_0(s)=\Lambda_\tau(s)\), and
\[
 \begin{gathered}
 v_w(\rho_s)\le a_s+\eta_w,\qquad
 -v_w(\rho_s)\le |v_w(z_T(s))+g_T(s,w)|+n_s\le a_s+n_s
       \quad(1\le s\le h),\\
 v_w(\rho_s)=v_w(\rho_{1-s})\quad(1-e_0\le s\le0).
 \end{gathered}
\]
The two positive-range inequalities follow by writing
\(v_w(\rho_s)=v_w(z_\psi(s))-v_w(z_T(s))\) and adding and subtracting \(g_T(s,w)\).
\leanref{GlobalTheory.\allowbreak{}SqTop.\allowbreak{}RhoLow}{Development.\allowbreak{}AbelianVarietyTheory.\allowbreak{}Statements.\allowbreak{}PlaceShiftStatements.\allowbreak{}NormSmall}{153--158}
\leanref{GlobalTheory.\allowbreak{}SqTop.\allowbreak{}Refl}{Development.\allowbreak{}AbelianVarietyTheory.\allowbreak{}Statements.\allowbreak{}PlaceShiftStatements.\allowbreak{}NormSmall}{148--152}
\leanref{GlobalTheory.\allowbreak{}Square.\allowbreak{}Ratio.\allowbreak{}exists\_\allowbreak{}ratios}{Development.\allowbreak{}AbelianVarietyTheory.\allowbreak{}Proofs.\allowbreak{}FoundationProofs.\allowbreak{}ExistsRatios}{33--84}
\leanref{GlobalTheory.\allowbreak{}SqTop.\allowbreak{}OneSqProof.\allowbreak{}wt\_\allowbreak{}at}{Development.\allowbreak{}AbelianVarietyTheory.\allowbreak{}Proofs.\allowbreak{}DivisorProofs.\allowbreak{}OneSqStOfClaims}{77--94}

Let \(P_0,P'_0\) be the fitted bottom pair at the base tuple, and let \(P_\tau,P'_\tau\) be a bottom pair at tuple \(\tau\) and level \(Q_{\rm lev}Q_{\rm tw}\), using the same period and two-torsion point. Put \(A=z_A(i,j)\) and \(A'=z'_A(i,j)\). Thus \(\iota(A)=\operatorname{norm}(P_0)\) and \(\iota(A')=\operatorname{norm}(P'_0)\). The larger-bottom norm formula supplies a nonzero complex number \(H\), independent of the tuple and its level, for which
\[
 \operatorname{norm}(P_u)=H\prod_{s=1-e_0}^{h_0}\Lambda_u(s)
       \qquad(u=\mu_0,\tau).
\]
Consequently
\[
 z=A\prod_{s=1-e_0}^{h_0}\rho_s\ne0,
 \qquad \iota(z)=\operatorname{norm}(P_\tau),\qquad A\ne0.
\]
The nonzero-coordinate property of \(P'_0\) makes its CM theta norm nonzero, so \(A'\ne0\). Let \(b(u)\) be the proportionality constant of a bottom theta triple at tuple \(u\), with the standard additive character and the fitted field of scalars. The comparison of bottom constants, applied to the two bottom pairs and the fitted coordinate-integrality allowance, gives \(b(\mu_0)\ne0\) and \(\gamma\in K^\times\) such that
\[
 \iota(\gamma)b(\mu_0)=b(\tau),\qquad
 v_w(\gamma)\le v_w(z)+v_w(A')+6d_\ell(w).
\]
The level comparison uses \(w\nmid Q_{\rm tw}\), so the \(w\)-adic allowance at \(Q_{\rm lev}Q_{\rm tw}\) is the allowance at \(Q_{\rm lev}\).
\leanref{GlobalTheory.\allowbreak{}SqTop.\allowbreak{}NormBig}{Development.\allowbreak{}AbelianVarietyTheory.\allowbreak{}Statements.\allowbreak{}PlaceShiftStatements.\allowbreak{}NormSmall}{114--117}
\leanref{GlobalTheory.\allowbreak{}T.\allowbreak{}definiteThetaPair\_\allowbreak{}norm\_\allowbreak{}ne\_\allowbreak{}zero\_\allowbreak{}of\_\allowbreak{}coord\_\allowbreak{}ne\_\allowbreak{}zero}{Development.\allowbreak{}AbelianVarietyTheory.\allowbreak{}Proofs.\allowbreak{}DivisorProofs.\allowbreak{}DefiniteThetaPairNormNeZeroOfCoordNeZeroStatement}{13--15}
\leanref{GlobalTheory.\allowbreak{}SqTop.\allowbreak{}BotConst}{Development.\allowbreak{}AbelianVarietyTheory.\allowbreak{}Statements.\allowbreak{}PlaceShiftStatements.\allowbreak{}NormSmall}{139--147}
\leanref{GlobalTheory.\allowbreak{}SqTop.\allowbreak{}botConst\_\allowbreak{}of\_\allowbreak{}inputs\_\allowbreak{}alt}{Development.\allowbreak{}AutomorphicArithmetic.\allowbreak{}ThetaCorrespondence.\allowbreak{}ThetaPairLifts.\allowbreak{}LocalTypeDictionary.\allowbreak{}BottomConstantOfTheta}{22--61}

Write \(c(u)\) for the square's top proportionality constants, so that \(c(\mu_0)\ne0\) and \(\iota(r_{ij,\tau})=c(\tau)/c(\mu_0)\). The seesaw tower has \(q_*\) small-rank steps and \(p_*\) large-rank steps. Its normalized step factors \(u_0,\ldots,u_{y-1}\in\C\) and its comparison factor \(\zeta\in K\) satisfy
\[
 \frac{c(\tau)}{c(\mu_0)}
 \prod_{k=0}^{q_*-1}\frac{\Lambda_\tau(e+k+1)}{\Lambda_0(e+k+1)}
 =\iota(\zeta)\left(\prod_{k=0}^{y-1}u_k\right)
 \frac{b(\tau)}{b(\mu_0)}
 \prod_{k=0}^{p_*-1}\frac{\Lambda_\tau(e+q_*-k)}{\Lambda_0(e+q_*-k)},
\]
with \(\zeta\ne0\) and \(v_w(\zeta)=0\). In this identity each step has the form
\[
 u_k=\iota\left(\epsilon_k
          \prod_{q\in S_0}\frac{a_{k,q}}{a'_{k,q}}\right),
 \qquad
 a_{k,q},a'_{k,q}\in\mathcal O_K,\qquad
 a_{k,q}\mid(N\mathfrak q)^{c_q},\quad
 a'_{k,q}\mid(N\mathfrak q)^{c_q},
\]
where \(\epsilon_k\ne0\) and \(v_w(\epsilon_k)=0\). The tower is obtained by applying the small-rank and large-rank norm steps to the top triple, comparing the descended triple with the chosen bottom triple, and dividing the resulting identity at \(\tau\) by the identity at \(\mu_0\); all common scalar factors cancel. The top and bottom comparison factors are units away from \(Q_{\rm tw}\), as are the factors \(\epsilon_k\).
\leanref{GlobalTheory.\allowbreak{}SqTop.\allowbreak{}Tower}{Development.\allowbreak{}AbelianVarietyTheory.\allowbreak{}Statements.\allowbreak{}PlaceShiftStatements.\allowbreak{}NormSmall}{128--138}
\leanref{GlobalTheory.\allowbreak{}SqTower.\allowbreak{}tower\_\allowbreak{}alt\_\allowbreak{}bis}{Development.\allowbreak{}AutomorphicArithmetic.\allowbreak{}ThetaCorrespondence.\allowbreak{}ThetaPairLifts.\allowbreak{}LocalTypeDictionary.\allowbreak{}TowerOfThetaMatching}{32--104}
\leanref{GlobalTheory.\allowbreak{}Square.\allowbreak{}StepRatio}{Development.\allowbreak{}AbelianVarietyTheory.\allowbreak{}Definitions.\allowbreak{}Foundations.\allowbreak{}TupleConstantsNeZero}{113--117}

Set \(t_k=\epsilon_k\prod_{q\in S_0}a_{k,q}/a'_{k,q}\). Every divisor in this formula is nonzero. Its valuation lies between \(0\) and \(c_qv_w(N\mathfrak q)\), so the triangle inequality gives
\[
 t_k\ne0,\qquad u_k=\iota(t_k),\qquad |v_w(t_k)|\le2C_w.
\]
Substitute these identities, the \(\rho_s\), and \(\gamma\) in the complex tower identity and use injectivity of \(\iota\). This gives the identity in \(K\)
\[
 r_{ij,\tau}\prod_{k=0}^{q_*-1}\rho_{e+k+1}
 =\zeta\left(\prod_{k=0}^{y-1}t_k\right)\gamma
             \prod_{k=0}^{p_*-1}\rho_{e+q_*-k}.
\]
Apply Lemma~\ref{B8.21ja} with \((p,q)=(p_*,q_*)\), \(B=2C_w\), \(d=d_\ell(w)\), and \(\eta=\eta_w\). It gives \(r_{ij,\tau}\ne0\) and the required bound. The four positive-range blocks in that lemma give the coefficient \(4\); its nonnegative denominator sum is enlarged to the prescribed coefficient \(5\); the lengths of the upper-bounded blocks give the error allowance \(2h\eta_w\); the bottom comparison contributes \(6d_\ell(w)\); and the \(y\) normalized steps contribute \(2yC_w\).
\leanref{GlobalTheory.\allowbreak{}Square.\allowbreak{}stepRatio\_\allowbreak{}vw\_\allowbreak{}bound}{Development.\allowbreak{}AbelianVarietyTheory.\allowbreak{}Proofs.\allowbreak{}FoundationProofs.\allowbreak{}StepRatioAbsVwProd}{60--85}
\leanref{GlobalTheory.\allowbreak{}SqTop.\allowbreak{}OneSqNoSmall.\allowbreak{}oneSqSt\_\allowbreak{}of\_\allowbreak{}claims'}{Development.\allowbreak{}AbelianVarietyTheory.\allowbreak{}Proofs.\allowbreak{}DivisorProofs.\allowbreak{}SquareRatioValuationWithoutNormFormula}{24--241}

Finally, the fitted supports consist of primes, and the nonnegative ceiling defining \(N_0\), together with the denominator-sum estimate, gives \(N_0(s,l)=0\) off \(S_d\) and \(\sum_{l\in S_d}N_0(s,l)\log l\le B_0\). Definition~\ref{B8.21e} gives \(c_q\le e_q\), while \(0\le y\) and \(y+1\le m\) follow from \(0\le j<m\). These statements and the rank calculation above give all the remaining assertions.
\leanref{GlobalTheory.\allowbreak{}BigTop.\allowbreak{}Fun.\allowbreak{}den\_\allowbreak{}tie\_\allowbreak{}alt}{Development.\allowbreak{}AutomorphicArithmetic.\allowbreak{}ThetaCorrespondence.\allowbreak{}ThetaPairLifts.\allowbreak{}ThirdPlaceSeries.\allowbreak{}DenTie}{14--26}
The complete proof is given in \Lean{GlobalTheory.\allowbreak{}SqTop.\allowbreak{}OneSqNoSmall.\allowbreak{}oneSqSt\_\allowbreak{}of\_\allowbreak{}claims'} and its family specialization \Lean{GlobalTheory.\allowbreak{}BigTop.\allowbreak{}Fun.\allowbreak{}oneSq\_\allowbreak{}of\_\allowbreak{}log\_\allowbreak{}derivative\_\allowbreak{}bound\_\allowbreak{}thA}.
\end{proof}

\subsection*{Proof of \ref{B8.21ka}}\label{prf:B8.21ka}\phantomsection
\begin{proof}
The definitions give nonnegativity and finite support. Put \(r=2fN\ge2\). Since \((4r)!\le(4r)^{4r}\),
\[
 \begin{split}
 \log M
 &=2r\log8+\log((4r)!)\\
 &\le2r\log8+4r\log(4r)
 =4r\log r+14r\log2
 \le18r\log r.
 \end{split}
\]
Prime factorization gives
\[
 \sum_{p\in S'}v_p(M)\log p=\log M,
 \qquad
 \sum_{p\in S'}v_p(Q)\log p=\log Q.
\]
Also \(\sum_{p\in S'}a_p\log p\le L\) and
\(\sum_{p\in S'}\mathbf1_{p=d}\log p/4\le\log d/4\le L/2\).
Substitution gives the first displayed bound. For the second, use
\(6fNBL\le36fNBL\), \(N\ge1\), \(fB\ge1\), and \(\log Q\ge0\) to obtain
\[
 4fNB\log Q+4N^2\log Q
 \le8fN^2B\log Q\le72fN^2B\log Q.
\]
\end{proof}

\subsection*{Proof of \ref{B8.21k}}\label{prf:B8.21k}\phantomsection
\begin{proof}[Proof (outline)]
The genus-one CM construction supplies a finite set \(S\) and nonnegative derivation-denominator exponents \(a_p\), zero outside \(S\), with
\[
 \sum_{p\in S}a_p\log p\le\log\operatorname{rd}(E)\le L,
 \qquad \log d\le2\log\operatorname{rd}(E)\le2L,
 \qquad f\ge1.
\]
Define \(M,S'\) and \(e_p\) as in Lemma~\ref{B8.21ka}, and set \(\delta_{\rm vp}(p)=e_p\). That lemma gives
\[
 \sum_{p\in S'}\delta_{\rm vp}(p)\log p\le I(18;f,N,B,L,Q).
\]
\leanref{GlobalTheory.\allowbreak{}NormSum.\allowbreak{}letters\_\allowbreak{}and\_\allowbreak{}sizes}{Development.\allowbreak{}AbelianVarietyTheory.\allowbreak{}Proofs.\allowbreak{}FoundationProofs.\allowbreak{}DerivationDenominatorsAndHostFieldSizes}{229--243}

For every prime \(p\), every \(p\)-compatible level \(Q'\), and every \((N,B)\)-small pair \(P\) at that level, the CM jet coordinate estimate supplies a representative integral over \(p\) in every target class and gives
\[
 \begin{gathered}
 1\le |\Gamma_D(J_P,h)|,\qquad
 v_p(|\Gamma_D(J_P,h)|)\le e_p,\\
 c_P(h,T)\in\overline{\Q},\qquad
 |\lambda(c_P(h,T))|_p\le p^{e_p}
 \end{gathered}
\]
for each integral representative \(h\), each minor tuple \(T\), and each \(\lambda:\overline{\Q}\to\C_p\). The stabilizer inequalities hold for every \(h\). The exponent in this estimate is precisely
\[
 v_p\bigl(8^{4fN}(8fN)!\bigr)
 +4fNB\left(a_p+\frac{\mathbf1_{p=d}}4+v_p(Q)\right)
 +4N^2v_p(Q).
\]
\leanref{GlobalTheory.\allowbreak{}NormSum.\allowbreak{}CoordBound.\allowbreak{}coordInt}{Development.\allowbreak{}AbelianVarietyTheory.\allowbreak{}Proofs.\allowbreak{}FoundationProofs.\allowbreak{}CoordIntOfJets}{176--178}

At level \(Q\), choose representatives integral over \(p\) and expand
\[
 \mathrm{norm}(P)=
 \sum_{q,T}|\Gamma_D(J_P,h_q)|^{-1}
             c_P(h_q,T)\overline{c_P(h_q,T)}.
\]
The expansion also shows that the complex value of \(\mathrm{norm}(P)\) belongs to \(\overline{\Q}\). The same coordinate bound applies to the conjugate coordinate through \(\lambda\circ\overline{\phantom z}\). Each summand therefore has \(p\)-adic absolute value at most \(p^{3e_p}\), and the ultrametric inequality gives
\[
 |\lambda(\mathrm{norm}(P))|_p\le p^{3e_p}.
\]
A nonzero coordinate gives \(\mathrm{norm}(P)>0\), so every element \(z\in K\) with \(\sigma(z)=\mathrm{norm}(P)\) is nonzero. For \(w\mid p\), transfer the last inequality to the embedding that induces \(w\); it becomes
\[
 v_w(z)\ge-3e_pv_w(p)=-3d_\ell(w),
 \qquad d_\ell(w)=\sum_{p\in S'}e_pv_w(p).
\]
The identity \(\sum_{w\mid p}\ell_wv_w(p)=\log p\) and nonnegativity give, for every finite set \(U\) of places,
\[
 \sum_{w\in U}\ell_wd_\ell(w)
 \le\sum_{p\in S'}e_p\log p
 \le I(18;f,N,B,L,Q).
\]
\leanref{GlobalTheory.\allowbreak{}NormSum.\allowbreak{}RepInv.\allowbreak{}rep\_\allowbreak{}indep}{Development.\allowbreak{}AbelianVarietyTheory.\allowbreak{}Proofs.\allowbreak{}FoundationProofs.\allowbreak{}DefiniteThetaPairTermsNormsIndependentRepresenta}{241--266}
\leanref{GlobalTheory.\allowbreak{}T.\allowbreak{}definiteThetaPairNorm\_\allowbreak{}padicValuation\_\allowbreak{}bound}{Development.\allowbreak{}AbelianVarietyTheory.\allowbreak{}Proofs.\allowbreak{}FoundationProofs.\allowbreak{}DefiniteThetaPairNormPadicValuationBoundStatement}{8--10}
\leanref{GlobalTheory.\allowbreak{}PadicPlaceBound.\allowbreak{}neg\_\allowbreak{}mul\_\allowbreak{}vw\_\allowbreak{}le\_\allowbreak{}vw}{Development.\allowbreak{}AbelianVarietyTheory.\allowbreak{}Proofs.\allowbreak{}FoundationProofs.\allowbreak{}NormLeOneOfValuationLeOne}{122--157}

Apply this construction with \(N=N_*\), \(B=B_w(c_{\rm gap})\), \(L=\log b\), and \(Q=Q_{\rm lev}\). Since the coefficient of \(C\) in \(I(C;f,N,B,L,Q)\) is nonnegative, the same allowance works for every \(C\ge18\).
\leanref{GlobalTheory.\allowbreak{}FunObj.\allowbreak{}SqNum.\allowbreak{}exists\_\allowbreak{}isIntExp}{Development.\allowbreak{}AbelianVarietyTheory.\allowbreak{}Definitions.\allowbreak{}Divisors.\allowbreak{}FinitePrimesUnder}{303--311}
\leanref{GlobalTheory.\allowbreak{}BigTop.\allowbreak{}Fun.\allowbreak{}TieModel.\allowbreak{}isIntExp\_\allowbreak{}mono}{Development.\allowbreak{}AbelianVarietyTheory.\allowbreak{}Definitions.\allowbreak{}Divisors.\allowbreak{}MulVecEqOfBasis}{38--51}
The complete proof is given in \Lean{GlobalTheory.\allowbreak{}NormSum.\allowbreak{}DenomBound.\allowbreak{}intSt\_\allowbreak{}of\_\allowbreak{}forall\_\allowbreak{}inp}.
\end{proof}

\subsection*{Proof of \ref{B8.21l}}\label{prf:B8.21l}\phantomsection
\begin{proof}[Proof (outline)]
At a square-recursive datum over a grid \(G\) at \((1,C_F,b)\), the grid level \(Q_{\rm lev}\) is admissible for \(D\) and every standard bottom pair \((P_{ij},P'_{ij})\) on \(G\) at \((\mu_0,i,j,\mathfrak m_E)\) over \((\mathcal O_{K_d},\delta,t_0,Q_{\rm lev})\) is \((N_*,B_w)\)-small; under \(\mathscr P\), the Hecke--Tate functional equation, and Theorem~\ref{B8.11}, nonvanishing of the Hecke \(L\)-values, real Howe duality, and bottom theta matching produce at least one nonzero coordinate on both \(P_{ij}\) and \(P'_{ij}\) for every \(0\le i,j<m\).
\leanref{GlobalTheory.\allowbreak{}SqTop.\allowbreak{}BotSmallProof.\allowbreak{}botSmall\_\allowbreak{}holds}{Development.\allowbreak{}AbelianVarietyTheory.\allowbreak{}Proofs.\allowbreak{}DivisorProofs.\allowbreak{}BottomRankSumLeTwoMulTotalRank}{128--136}
\leanref{GlobalTheory.\allowbreak{}SqTop.\allowbreak{}lNe\_\allowbreak{}of\_\allowbreak{}functionalEquation}{Development.\allowbreak{}AbelianVarietyTheory.\allowbreak{}Proofs.\allowbreak{}DivisorProofs.\allowbreak{}SquareBottomHalfRankDiffLeHalfRank}{171--181}
\leanref{GlobalTheory.\allowbreak{}T.\allowbreak{}bottomTheta\_\allowbreak{}nonvanishing}{Development.\allowbreak{}AutomorphicArithmetic.\allowbreak{}ThetaCorrespondence.\allowbreak{}ThetaPairLifts.\allowbreak{}LocalTypeDictionary.\allowbreak{}BottomThetaNeZeroOfInnerProductFormula}{81--82}

In a fitted evaluation datum \(d\) at \((p,C_F)\), the square-data tie of Definition~\ref{B8.21d}(3) supplies \((\delta,t_0)\) and standard bottom pairs \((P_{ij},P'_{ij})\) with \(\iota(z_A(i,j))=\mathrm{norm}(P_{ij})\) and \(\iota(z_A'(i,j))=\mathrm{norm}(P'_{ij})\), while \(d_\ell\) is a grid integrality allowance on \(G\) at constant \(C_F\) in the sense of Definition~\ref{B8.21ca}(4). Applying the norm valuation bound of \(d_\ell\) at \(N=N_*\), \(B=B_w\), \(L=\log b\ge1\), and \(Q=Q_{\rm lev}\) (\(\log Q=c_S\)) gives \(d_\ell(w)\ge0\), \(z_A(i,j),z_A'(i,j)\ne0\), \(v_w(z)\ge-3d_\ell(w)\), and \(\sum_{w\in U}\ell_w d_\ell(w)\le I(C_F;f,N_*,B_w,L,Q_{\rm lev})\le D_0\) whenever \(C_1\ge\max(C_F,0)\).
\leanref{GlobalTheory.\allowbreak{}BigTop.\allowbreak{}Fun.\allowbreak{}IntTie.\allowbreak{}integralNormBound\_\allowbreak{}mono}{Development.\allowbreak{}AbelianVarietyTheory.\allowbreak{}Proofs.\allowbreak{}DivisorProofs.\allowbreak{}IntegralityFittedIdentity}{35--38}

The complete proof is given in \Lean{GlobalTheory.\allowbreak{}BigTop.\allowbreak{}Fun.\allowbreak{}int\_\allowbreak{}tie\_\allowbreak{}alt}.
\end{proof}

\subsection*{Proof of \ref{B8.21ma}}\label{prf:B8.21ma}\phantomsection
\begin{proof}[Proof (outline)]
For a period lattice \(a(\Z+\tau\Z)\), write \(y=\Im\tau\) and let \(\Delta(\tau)\) be the normalized modular discriminant. The modular invariant function \(y^6|\Delta(\tau)|\) is bounded on a fundamental domain: it is continuous on each compact truncation and tends to zero at the cusp. Choose an upper bound \(M_\Delta\ge0\). The discriminant-one identity gives
\[
 \operatorname{area}(a(\Z+\tau\Z))^6
 =|a|^{12}y^6=(2\pi)^{12}y^6|\Delta(\tau)|.
\]
Thus \(C_A=\max\{1,(2\pi)^{12}M_\Delta\}\) has the required property.

Put \(R=\operatorname{rd}(E)\), \(\ell_d=\Re(L'(1,\chi_d)/L(1,\chi_d))\), and \(n_M=[M:\Q]\). The induced genus-one type is admissible for every arithmetic datum. The host-field size estimate gives \(R\ge1\) and \(\log d\le2\log R\). Conjugation of the curve datum realizes each \(\Lambda_\rho\) as a curve lattice for the same arithmetic datum. In particular, it has discriminant one, is stable under \(\mathcal O_{\Q(\sqrt{-d})}\), and has integral \(j\)-invariant. Hence \(0<A_\rho\le C_A\). Since \(t_\rho\) is unkilled and the lattice is stable under the CM order, \(t_\rho\notin\Lambda_\rho\).
\leanref{GlobalTheory.\allowbreak{}InducedType.\allowbreak{}induced\_\allowbreak{}type}{Development.\allowbreak{}AbelianVarietyTheory.\allowbreak{}Proofs.\allowbreak{}FoundationProofs.\allowbreak{}InducedTypeThetaGenusOnePermutes}{74--76}
\leanref{GlobalTheory.\allowbreak{}T.\allowbreak{}derivationDenominators\_\allowbreak{}and\_\allowbreak{}hostField\_\allowbreak{}sizes}{Development.\allowbreak{}AbelianVarietyTheory.\allowbreak{}Proofs.\allowbreak{}FoundationProofs.\allowbreak{}DerivationDenominatorsAndHostFieldSizesStatement}{8--10}
\leanref{GlobalTheory.\allowbreak{}ArchFree.\allowbreak{}means.\allowbreak{}exists\_\allowbreak{}crv\_\allowbreak{}conj}{Development.\allowbreak{}AbelianVarietyTheory.\allowbreak{}Proofs.\allowbreak{}FoundationProofs.\allowbreak{}ConjugateLatticeMeansPosLogAreaThetaBounded}{16--54}
\leanref{GlobalTheory.\allowbreak{}T.\allowbreak{}idealPeriodLattice\_\allowbreak{}of\_\allowbreak{}curveDatum}{Development.\allowbreak{}AbelianVarietyTheory.\allowbreak{}Proofs.\allowbreak{}FoundationProofs.\allowbreak{}IdealPeriodLatticeOfCurveDatumStatement}{8--10}
\leanref{GlobalTheory.\allowbreak{}T.\allowbreak{}admJ\_\allowbreak{}and\_\allowbreak{}torsion\_\allowbreak{}data\_\allowbreak{}algebraic}{Development.\allowbreak{}AbelianVarietyTheory.\allowbreak{}Proofs.\allowbreak{}FoundationProofs.\allowbreak{}IntegralJAndAlgebraicDataStatement}{8--10}
\leanref{GlobalTheory.\allowbreak{}PositiveMeans.\allowbreak{}not\_\allowbreak{}mem\_\allowbreak{}lattice\_\allowbreak{}of\_\allowbreak{}isCMTorsionBasePoint}{Development.\allowbreak{}AbelianVarietyTheory.\allowbreak{}Proofs.\allowbreak{}FoundationProofs.\allowbreak{}ExistsBoundMeansPosLogAreaNormalisedTheta}{205--232}

The period identity, with \(\gamma\) denoting Euler's constant, is
\[
 \frac1{n_M}\sum_\rho\log(A_\rho/\pi)
 =\gamma+\log(2\pi)-\frac12\log d-\ell_d.
\]
For each \(\rho\),
\[
 \log^+(A_\rho^{-1})
 =\log^+A_\rho-\log A_\rho
 \le\log^+C_A-\log(A_\rho/\pi)-\log\pi.
\]
Averaging, using \(\ell_d\le L\), \(\frac12\log d\le L\), \(L\ge1\), and \(\gamma,\log\pi,\log(2\pi)>0\), gives
\[
 \frac1{n_M}\sum_\rho\log^+(A_\rho^{-1})
 \le\log^+C_A+2L\le(\log^+C_A+2)L.
\]
\leanref{GlobalTheory.\allowbreak{}T.\allowbreak{}mean\_\allowbreak{}log\_\allowbreak{}period\_\allowbreak{}conjugates\_\allowbreak{}eq\_\allowbreak{}chowlaSelberg}{Development.\allowbreak{}AbelianVarietyTheory.\allowbreak{}Proofs.\allowbreak{}FoundationProofs.\allowbreak{}MeanLogPeriodConjugatesEqChowlaSelbergStatement}{8--10}

The genus-one value estimate bounds the squared value at zero of every unitary translate of the odd theta function by \(3(2R+1)\) times its Gaussian weighted mean square. For a discriminant-one lattice with integral \(j\)-invariant and a two-torsion point outside the lattice, the theta normalization therefore gives
\[
 \log^+\nu_\rho\le
 2\log^+\bigl(3(2R+1)\bigr)
 +\frac32\bigl(\log^+(2\pi^2)+\log^+(A_\rho^{-1})\bigr).
\]
Here the factor \(3/2\) combines the identity for the odd theta norm with the reciprocal bound for its translated value at the two-torsion point. Since \(3(2R+1)\le9R\) and \(\log9\le8\),
\[
 \log^+\bigl(3(2R+1)\bigr)\le9L.
\]
The preceding mean bound now gives
\[
 \frac1{n_M}\sum_\rho\log^+\nu_\rho
 \le18L+\frac32\log^+(2\pi^2)
       +\frac32(\log^+C_A+2)L
 \le C_{\rm gal}L.
\]
The same constant exceeds \(\log^+C_A+2\). All choices are uniform in \(D,M\), and the chosen torsion points.
\leanref{GlobalTheory.\allowbreak{}T.\allowbreak{}sigmaStar\_\allowbreak{}value\_\allowbreak{}le\_\allowbreak{}meanSquare\_\allowbreak{}genusOne}{Development.\allowbreak{}AbelianVarietyTheory.\allowbreak{}Proofs.\allowbreak{}FoundationProofs.\allowbreak{}SigmaStarValueLeMeanSquareGenusOneStatement}{8--10}
\leanref{GlobalTheory.\allowbreak{}PositiveMeans.\allowbreak{}posLog\_\allowbreak{}nrm\_\allowbreak{}le}{Development.\allowbreak{}AbelianVarietyTheory.\allowbreak{}Proofs.\allowbreak{}FoundationProofs.\allowbreak{}ExistsBoundMeansPosLogAreaNormalisedTheta}{102--128}
The complete proofs are \Lean{GlobalTheory.\allowbreak{}PositiveMeans.\allowbreak{}means\_\allowbreak{}le} and its unconditional version \Lean{GlobalTheory.\allowbreak{}ArchFree.\allowbreak{}means.\allowbreak{}means\_\allowbreak{}le}.
\end{proof}

\subsection*{Proof of \ref{B8.21m}}\label{prf:B8.21m}\phantomsection
\begin{proof}[Proof (outline)]
Write \(\log^+ u=\max(0,\log u)\) for \(u>0\), with \(\log^+0=0\). For a CM lattice \(\Lambda\), put \(A_\Lambda=\operatorname{area}(\Lambda)\), and fix the additive fundamental domain \(\mathcal F_\Lambda\) used in the theta normalization. For an unkilled two-torsion point \(t\), define
\[
 \nu(\Lambda,t)=\frac{1}{\operatorname{vol}(\mathcal F_\Lambda)}
 \int_{\mathcal F_\Lambda}
 |\theta_0(\Lambda,t)(z)|^2
 \exp\!\left(-\frac{\pi|z|^2}{A_\Lambda}\right)\,d\operatorname{vol}(z).
\]
Thus \(\nu\) is the Gaussian weighted mean square of the normalized theta function.
\leanref{Formalism.\allowbreak{}MainNotions.\allowbreak{}StandardLatticesGL3.\allowbreak{}nrm1Sq}{Formalism.\allowbreak{}SharpRestriction.\allowbreak{}GlobalSchurBounds.\allowbreak{}StandardLatticesGL3}{53--57}

Take \(C_A\) and \(C_{\rm gal}\) from Lemma~\ref{B8.21ma}. Its two mean bounds apply to every number field containing \(g_2(\Lambda),g_3(\Lambda),\wp_\Lambda(t_0)\) and every permitted family of conjugate torsion points. Set
\[
 C_{\rm arch}=1100+4C_{\rm gal},\qquad
 C_{\rm gal}=18+\frac32\bigl(\log^+(2\pi^2)+\log^+C_A+2\bigr).
\]

Fix one \((N,B)\)-small pair \(P\) and \(z\in K\) with \(\sigma(z)=\operatorname{norm}(P)\). Write \(a=a_P\) for its smaller rank, retain \(d_P\) for its larger rank, and put \(s=\sum_\tau n_{P,\tau}\) and \(R=\operatorname{rd}(E)\). Then
\[
 a\le d_P\le2N,\qquad s\le2fB,\qquad f,N,B,L,R\ge1.
\]
Let \(c_J=|\mathrm{Cls}_D(d_P,J_P)|\), \(c_T=|\mathrm{Tup}_D(a,d_P,n_P)|\), and \(c_u=|\mathrm{Cls}_D(u,G_{\rm int}(u))|\) for \(u=a,d_P\), where \(G_{\rm int}(u)\) is the standard integral compact subgroup. The class and minor counts give
\[
 \begin{split}
 c_J&\le c_{d_P}Q^{d_P^2+4},\qquad c_T\le2^{d_Ps},\\
 c_u&\le\bigl(12f^2u^4R^2\bigr)^{6fu^2}
       \quad(u=a,d_P).
 \end{split}
\]
\leanref{GlobalTheory.\allowbreak{}NormSum.\allowbreak{}cls\_\allowbreak{}finite\_\allowbreak{}card\_\allowbreak{}le}{Development.\allowbreak{}AbelianVarietyTheory.\allowbreak{}Proofs.\allowbreak{}FoundationProofs.\allowbreak{}UnitaryDoubleCosetClassesFiniteCardLe}{37--96}
\leanref{GlobalTheory.\allowbreak{}MinorBound.\allowbreak{}minor\_\allowbreak{}bound}{Development.\allowbreak{}AbelianVarietyTheory.\allowbreak{}Proofs.\allowbreak{}FoundationProofs.\allowbreak{}SumCoeffMonomial}{96--139}

Choose a finite Galois extension \(M/\Q\) in \(\C\) containing \(\sigma(K)\), the finitely many coordinates \(c_P(h_q,T)\), and \(g_2(\Lambda),g_3(\Lambda),\wp_\Lambda(t_0)\). Let \(c\) denote complex conjugation on \(M\), put \(w_{q,T}=|\Gamma_D(J_P,h_q)|^{-1}\), and define
\[
 S_\rho=\sum_{q,T}w_{q,T}|\rho(c_P(h_q,T))|^2.
\]
The stabilizer orders are positive integers, so \(0<w_{q,T}\le1\). The identity
\[
 \sigma(z)=\sum_{q,T}w_{q,T}c_P(h_q,T)c(c_P(h_q,T))
\]
and weighted Cauchy--Schwarz imply
\[
 |\rho(\sigma(z))|\le(S_\rho S_{\rho\circ c})^{1/2}.
\]
The map \(\rho\mapsto\rho\circ c\) permutes the embeddings of \(M\). Each embedding of \(K\) has \([M:\sigma(K)]\) extensions to \(M\), with \(K\) identified with \(\sigma(K)\). Since \(K\) is totally complex, averaging gives
\[
 \operatorname{mean}_{v\mid\infty}\log^+|z|_v
 \le\frac1{[M:\Q]}\sum_\rho\log^+S_\rho.
\]
\leanref{GlobalTheory.\allowbreak{}ArchMean.\allowbreak{}meanInf\_\allowbreak{}posLog\_\allowbreak{}le}{Development.\allowbreak{}AbelianVarietyTheory.\allowbreak{}Proofs.\allowbreak{}FoundationProofs.\allowbreak{}MeanInfPosLogLe}{75--132}

Write \(A_\rho=A_{\Lambda_\rho}\) and \(\nu_\rho=\nu(\Lambda_\rho,t_\rho)\). Extend \(\rho:M\hookrightarrow\C\) to an embedding \(\overline{\Q}\hookrightarrow\C\). Since \(g_2(\Lambda),g_3(\Lambda)\in M\), the conjugate lattice obtained from the extension is \(\Lambda_\rho\). The conjugate coordinate estimate therefore gives
\[
 |\rho(c_P(h_q,T))|^2\le
 c_a^2Q^{ad_P}(a!)^{2s}\,
 3(2fad_PR+1)^{fad_P}\nu_\rho^{fad_P}
 \left(\frac{2e\pi B}{A_\rho}\right)^{as}.
\]
\leanref{GlobalTheory.\allowbreak{}T.\allowbreak{}thetaPair\_\allowbreak{}coordinates\_\allowbreak{}conjugate\_\allowbreak{}bound}{Development.\allowbreak{}AbelianVarietyTheory.\allowbreak{}Proofs.\allowbreak{}FoundationProofs.\allowbreak{}ThetaPairCoordinatesConjugateBoundStatement}{8--10}

The logarithmic bounds for the factors are as follows. Put \(\ell_B=\log(2NB)\), \(\ell_f=\log(2fNB)\), and \(\ell_Q=\log Q\). The class counts, the factorial estimate \(a!\le a^a\) for \(a\ge1\) together with \(0!=1\), and \(a,d_P\le2N\), \(s\le2fB\), \(R\le e^L\) give
\[
 \begin{aligned}
 c_J&\le\exp\bigl(24fN^2(12+4\ell_f+2L)+8N^2\ell_Q\bigr),\\
 c_T&\le\exp(8fNB),\\
 c_a^2&\le\exp\bigl(48fN^2(12+4\ell_f+2L)\bigr),\\
 Q^{ad_P}&\le\exp(4N^2\ell_Q),\qquad
 (a!)^{2s}\le\exp(8fNB\ell_B),\\
 (2fad_PR+1)^{fad_P}&\le\exp\bigl(4fN^2(3+2\ell_f+L)\bigr),\\
 \nu_\rho^{fad_P}&\le\exp\bigl(4fN^2\log^+\nu_\rho\bigr),\\
 \left(\frac{2e\pi B}{A_\rho}\right)^{as}
 &\le\exp\bigl(4fNB(7+\ell_B+\log^+A_\rho^{-1})\bigr).
 \end{aligned}
\]
Using \(3\le e^3\) and \(w_{q,T}\le1\), their product yields
\[
 \log^+S_\rho\le E_{\rm raw}
       +4fN^2\log^+\nu_\rho+4fNB\log^+A_\rho^{-1},
\]
where
\[
 E_{\rm raw}=3+fNB(36+12\ell_B)
       +fN^2(876+296\ell_f+148L)+12N^2\ell_Q.
\]
Define
\[
 \Xi=2fNB(L+\ell_B)+fN^2(L+\ell_f)+N^2\ell_Q.
\]
Since \(fN^2\ge1\), \(L\ge1\), and \(\ell_B,\ell_f,\ell_Q\ge0\),
\[
 3+876fN^2+148fN^2L\le1027fN^2L,
 \qquad E_{\rm raw}\le1100\Xi.
\]
Thus
\[
 \log^+S_\rho\le1100\Xi+
        4fN^2\log^+\nu_\rho+4fNB\log^+A_\rho^{-1}.
\]
\leanref{GlobalTheory.\allowbreak{}ArchNorm.\allowbreak{}sizes}{Development.\allowbreak{}AbelianVarietyTheory.\allowbreak{}Proofs.\allowbreak{}FoundationProofs.\allowbreak{}LatticeUnitNormSqNonneg}{104--204}

Averaging the last inequality and applying the two defining bounds for \(C_{\rm gal}\) gives
\[
 \operatorname{mean}_{v\mid\infty}\log^+|z|_v
 \le1100\Xi+4C_{\rm gal}L(fN^2+fNB)
 \le(1100+4C_{\rm gal})\Xi.
\]
This proves (1), for each pair in either family.

For (2), apply the logarithmic-derivative theorem for the eligible family. Under \(\mathscr P\), in the closed sense of Definition~\ref{B8.39}, it gives
\[
 \left|\Re\frac{L'(1,\chi_d)}{L(1,\chi_d)}\right|
 \le x^\varepsilon=\log b
\]
eventually at every member point. The row-setting conditions give \(\log\operatorname{rd}(E)\le\log b\) and \(\log b\ge1\). The fitted bottom pairs are \((N_*,B_w)\)-small and satisfy \(B_w\ge1\). Apply (1) with \(N=N_*\), \(B=B_w\), \(L=\log b\), and \(Q=Q_{\rm lev}\). Since \(\log Q_{\rm lev}=c_S\), choosing \(C_0\ge C_{\rm arch}\) gives precisely \(A_0\) for every \(C_1\ge C_0\).
\leanref{GlobalTheory.\allowbreak{}T.\allowbreak{}abs\_\allowbreak{}logDerivativeAtOne\_\allowbreak{}le\_\allowbreak{}x\_\allowbreak{}rpow\_\allowbreak{}eps}{Development.\allowbreak{}AutomorphicArithmetic.\allowbreak{}ThetaCorrespondence.\allowbreak{}ThetaPairLifts.\allowbreak{}ThirdPlaceSeries.\allowbreak{}LogDerivLFunctionAtOneLeAlongFamily}{23--24}
\leanref{GlobalTheory.\allowbreak{}BigTop.\allowbreak{}Fun.\allowbreak{}arch\_\allowbreak{}tie\_\allowbreak{}of\_\allowbreak{}logDeriv\_\allowbreak{}bound}{Development.\allowbreak{}AutomorphicArithmetic.\allowbreak{}ThetaCorrespondence.\allowbreak{}ThetaPairLifts.\allowbreak{}ThirdPlaceSeries.\allowbreak{}ArchTie}{14--47}

The complete proof is given in \Lean{GlobalTheory.\allowbreak{}ArchNorm.\allowbreak{}norm\_\allowbreak{}meanInf\_\allowbreak{}posLog\_\allowbreak{}le} for (1) and in \Lean{GlobalTheory.\allowbreak{}BigTop.\allowbreak{}Fun.\allowbreak{}arch\_\allowbreak{}tie\_\allowbreak{}of\_\allowbreak{}logDeriv\_\allowbreak{}bound} for (2).
\end{proof}

\subsection*{Proof of \ref{B8.21n}}\label{prf:B8.21n}\phantomsection
\begin{proof}[Proof (outline)]
Under \(\mathscr P\), the Hecke--Tate functional equation, and Theorem~\ref{B8.11}, the local sum-formula estimates give finite support, rationality of the evaluated row sums \(v_{\rm row}(a,t,x')_i\in\iota(K)\) under the evaluation-field condition, and nonvanishing of the seesaw constants \(c_{ij,t}\); combining these with the seesaw step-ratio tower and bottom triple integral formula places every ratio \(c_{ij,t}/c_{ij,\mu_0}\) in \(\iota(K)\) at square-recursive points.
\leanref{GlobalTheory.\allowbreak{}RowClaim.\allowbreak{}JoinF.\allowbreak{}joined\_\allowbreak{}local\_\allowbreak{}sum\_\allowbreak{}estimates}{Development.\allowbreak{}AutomorphicArithmetic.\allowbreak{}ThetaCorrespondence.\allowbreak{}ThetaPairLifts.\allowbreak{}LocalTypeDictionary.\allowbreak{}JoinedLocalSumEstimates}{12--16}
\leanref{GlobalTheory.\allowbreak{}BigTop.\allowbreak{}Fun.\allowbreak{}sqRatio\_\allowbreak{}mem\_\allowbreak{}of}{Development.\allowbreak{}AbelianVarietyTheory.\allowbreak{}Proofs.\allowbreak{}DivisorProofs.\allowbreak{}SquareRatioInEmbeddingRange}{23--88}
\leanref{GlobalTheory.\allowbreak{}BigTop.\allowbreak{}Fun.\allowbreak{}square\_\allowbreak{}ratio\_\allowbreak{}at\_\allowbreak{}point\_\allowbreak{}of\_\allowbreak{}leading\_\allowbreak{}term\_\allowbreak{}alt}{Development.\allowbreak{}AutomorphicArithmetic.\allowbreak{}ThetaCorrespondence.\allowbreak{}ThetaPairLifts.\allowbreak{}ThirdPlaceSeries.\allowbreak{}SquareRatioAtPointOfLeadingTerm}{19--45}

Above the joint family and grid cutoffs (with \(C_F\ge C_{\rm int}\)), the point size datum is square-recursive with \(L=\log b\ge1\), \(\log\operatorname{rd}(E)\le L\), and admissible grid level \(Q_{\rm lev}\). Using the integral CM period lattice \(\mathcal O_{K_d}\), the ordinary transport and partial-zeta denominator inputs give nonzero toric values, minimal polynomials over \(K\) for \(\beta(v)\), and finitely supported prime-count functions \(N_0(i,j,s,\cdot)\); a common discriminant twelfth-root period \(\delta_{\rm per}\) of \(\mathcal O_{K_d}\) and unkilled two-torsion point \(t_0\) for \(\Lambda(\mathcal O_{K_d},\delta_{\rm per})\) admit standard bottom pairs \((P_{ij},P'_{ij})\) for all \(0\le i,j<m\), so the square-field condition of \(a\) puts the twisted-to-untwisted toric ratios and standard bottom pair norms in \(\iota(K)\), and Lemma~\ref{B8.21k} supplies a grid integrality allowance \(d_\ell\) on \(G\) at \(C_F\). Assembling these witnesses with the direct-sum test lattice \((M,(X'_v)_{v\mid\infty})\) and the matrices \(V_t^0\) defined on standard basis vectors by \(\iota((V_t^0 e_\alpha)_i)=v_{\rm row}(a,t,e_{V,0}(p,t,e_\alpha))_i\) satisfies all tie relations of Definition~\ref{B8.21d}(3).
\leanref{GlobalTheory.\allowbreak{}ToricExist.\allowbreak{}torEx}{Development.\allowbreak{}AbelianVarietyTheory.\allowbreak{}Proofs.\allowbreak{}DivisorProofs.\allowbreak{}ExistsRatRepr}{128--157}
\leanref{GlobalTheory.\allowbreak{}DenExist.\allowbreak{}denEx}{Development.\allowbreak{}AbelianVarietyTheory.\allowbreak{}Proofs.\allowbreak{}DivisorProofs.\allowbreak{}BetaBoundVp}{71--78}
\leanref{GlobalTheory.\allowbreak{}FunObj.\allowbreak{}SqNum.\allowbreak{}exists\_\allowbreak{}isIntExp}{Development.\allowbreak{}AbelianVarietyTheory.\allowbreak{}Definitions.\allowbreak{}Divisors.\allowbreak{}FinitePrimesUnder}{303--311}
\leanref{GlobalTheory.\allowbreak{}BigTop.\allowbreak{}Fun.\allowbreak{}TieModel.\allowbreak{}exists\_\allowbreak{}tie}{Development.\allowbreak{}AbelianVarietyTheory.\allowbreak{}Definitions.\allowbreak{}Divisors.\allowbreak{}MulVecEqOfBasis}{86--112}

The complete proof is given in \Lean{GlobalTheory.\allowbreak{}BigTop.\allowbreak{}Fun.\allowbreak{}ex\_\allowbreak{}of\_\allowbreak{}claims\_\allowbreak{}alt}.
\end{proof}

\subsection*{Proof of \ref{B8.21o}}\label{prf:B8.21o}\phantomsection
\begin{proof}[Proof (outline)]
At every \(p\in\mathcal P\), the primitive link conditions furnish the level-zero row data, giving positive row rank, injectivity of \(\iota_{\mu_0}\), surjectivity of \(P_{\mu_0}\) and \(\iota_{\mu_0}P_{\mu_0}\), and the local lattice identifications \(RR_w=\iota_{\mu(w)}P_{\mu(w)}M_{\mu(w),w}\) by lattice gluing.
\leanref{GlobalTheory.\allowbreak{}BigTop.\allowbreak{}Fun.\allowbreak{}bundle\_\allowbreak{}of\_\allowbreak{}rows\_\allowbreak{}alt}{Development.\allowbreak{}AutomorphicArithmetic.\allowbreak{}ThetaCorrespondence.\allowbreak{}ThetaPairLifts.\allowbreak{}ThirdPlaceSeries.\allowbreak{}BundleOfRows}{13--17}

Under \(\mathscr P\) and Theorem~\ref{B8.11}, the local sum estimates yield convergence along the row lifts, nonvanishing seesaw constants \(c_{ij,t}\ne0\) (giving \(r_{\mu_0}=1\) and \(r_{\mu(w)}\ne0\) via the square-data tie of Definition~\ref{B8.21d}(3)), and the expansion formula \eqref{eq:p8c2-test-formula} relating the evaluated row sums of Definition~\ref{B8.21d}(3) to the row test functionals with a common archimedean factor \(\kappa\) and finite-character roots of unity \(\zeta_t\in K\) satisfying \(\zeta_{\mu_0}=1\); under \(\mathscr P\), twisted theta lifts are cuspidal, so Bessel-period convergence and row separation in Theorem~\ref{B8.15} imply that the chosen functionals have zero common kernel on the row space.
\leanref{GlobalTheory.\allowbreak{}RowClaim.\allowbreak{}JoinF.\allowbreak{}joined\_\allowbreak{}local\_\allowbreak{}sum\_\allowbreak{}estimates}{Development.\allowbreak{}AutomorphicArithmetic.\allowbreak{}ThetaCorrespondence.\allowbreak{}ThetaPairLifts.\allowbreak{}LocalTypeDictionary.\allowbreak{}JoinedLocalSumEstimates}{12--16}
\leanref{GlobalTheory.\allowbreak{}RowTest.\allowbreak{}injK\_\allowbreak{}of\_\allowbreak{}separates}{Development.\allowbreak{}AbelianVarietyTheory.\allowbreak{}Proofs.\allowbreak{}TwistedOperators.\allowbreak{}InjKOfSeparates}{91--104}

Since \(RR\) has positive rank, \(\iota_{\mu_0}P_{\mu_0}\) is onto \(RR\otimes_{\mathcal O_K}K\ne0\), \(c_{\mu_0}\ne0\), \(\zeta_{\mu_0}=1\), and the chosen functionals have zero common kernel, the left-hand side of \eqref{eq:p8c2-test-formula} at \(t=\mu_0\) is not identically zero, forcing \(\kappa\ne0\). Dividing by \(\kappa\) and factoring the base map through the surjective map \(\iota_{\mu_0}P_{\mu_0}\) into \(K^n=M\otimes_{\mathcal O_K}K\) (since \(M\) has full rank \(n\) by Definition~\ref{B8.21d}(3)) defines an injective \(\Psi\) satisfying the comparison identity across \(t=\mu_0\) and \(t=\mu(w)\), and injectivity of \(\Psi\) and \(\iota_{\mu_0}\) forces \(\ker V_{\mu_0}^0=\ker P_{\mu_0}\).
\leanref{GlobalTheory.\allowbreak{}BigTop.\allowbreak{}Fun.\allowbreak{}exists\_\allowbreak{}unit\_\allowbreak{}root\_\allowbreak{}ze\_\allowbreak{}leafC\_\allowbreak{}of\_\allowbreak{}square\_\allowbreak{}nonzero}{Development.\allowbreak{}AutomorphicArithmetic.\allowbreak{}ThetaCorrespondence.\allowbreak{}ThetaPairLifts.\allowbreak{}ThirdPlaceSeries.\allowbreak{}OneMapOfRows}{21--127}
\leanref{GlobalTheory.\allowbreak{}BigTop.\allowbreak{}Fun.\allowbreak{}ker\_\allowbreak{}mulVecLin\_\allowbreak{}eq\_\allowbreak{}U\_\allowbreak{}of\_\allowbreak{}formulaAt}{Development.\allowbreak{}AutomorphicArithmetic.\allowbreak{}ThetaCorrespondence.\allowbreak{}ThetaPairLifts.\allowbreak{}ThirdPlaceSeries.\allowbreak{}KerEqOfFormula}{20--86}
\leanref{GlobalTheory.\allowbreak{}BigTop.\allowbreak{}Fun.\allowbreak{}leafKernel\_\allowbreak{}of\_\allowbreak{}linK\_\allowbreak{}injK\_\allowbreak{}constK\_\allowbreak{}formula2K}{Development.\allowbreak{}AutomorphicArithmetic.\allowbreak{}ThetaCorrespondence.\allowbreak{}ThetaPairLifts.\allowbreak{}ThirdPlaceSeries.\allowbreak{}KerEqOfFormula}{87--101}

The complete proof is given in \Lean{GlobalTheory.\allowbreak{}BigTop.\allowbreak{}Fun.\allowbreak{}oneMap\_\allowbreak{}of\_\allowbreak{}bundle\_\allowbreak{}setAt\_\allowbreak{}kernel\_\allowbreak{}conditions}.
\end{proof}

\subsection*{Proof of \ref{B8.21p}}\label{prf:B8.21p}\phantomsection
\begin{proof}[Proof (outline)]
For \(x>1\) and \(\varepsilon>0\), \(L=x^\varepsilon\ge1\), and \(\log x\le2x^\varepsilon\) for large \(x\) gives \(L+\log x\le3L\). The window inequalities of the family class give \(n_{\rm half}=f/2\ge x^{2+8\varepsilon}\) and \(m\le2\sqrt{x\delta_1}\) with \(\delta_1=x^{-\varepsilon/8}\le1\), whence \(m^4\le16x^2\le2x^{2+8\varepsilon}\le f\) once \(x\ge8^{1/\varepsilon}\), together with \(2\le m\), \(m+3\le k\), \(f\le k\), \(x^{377}\le t\le x^T\) for \(T=\lceil|C_t|\rceil+1\), and \(0\le\mathrm{cl}(\tau,i)\le S_c\), so \(L_c\le S_c\le b\) by Definition~\ref{B8.21e} and \eqref{eq:p8c2-size-rows-c}.
\leanref{GlobalTheory.\allowbreak{}BigTop.\allowbreak{}Fun.\allowbreak{}existsIntegralMultiple\_\allowbreak{}holds\_\allowbreak{}alt}{Development.\allowbreak{}AutomorphicArithmetic.\allowbreak{}ThetaCorrespondence.\allowbreak{}ThetaPairLifts.\allowbreak{}ThirdPlaceSeries.\allowbreak{}ExistsIntegralMultipleHolds}{12--26}
\leanref{GlobalTheory.\allowbreak{}BigTop.\allowbreak{}Fun.\allowbreak{}Close.\allowbreak{}desHLeaf\_\allowbreak{}alt}{Development.\allowbreak{}AutomorphicArithmetic.\allowbreak{}ThetaCorrespondence.\allowbreak{}ThetaPairLifts.\allowbreak{}ThirdPlaceSeries.\allowbreak{}DesHLeaf}{13--55}
\leanref{GlobalTheory.\allowbreak{}BigTop.\allowbreak{}Fun.\allowbreak{}Close.\allowbreak{}cloudIcc}{Development.\allowbreak{}AbelianVarietyTheory.\allowbreak{}Proofs.\allowbreak{}DivisorProofs.\allowbreak{}ZetaDenominatorsFitFunction}{65--81}
\leanref{GlobalTheory.\allowbreak{}BigTop.\allowbreak{}Fun.\allowbreak{}Lc\_\allowbreak{}le\_\allowbreak{}of\_\allowbreak{}cloud\_\allowbreak{}le}{Development.\allowbreak{}AbelianVarietyTheory.\allowbreak{}Proofs.\allowbreak{}DivisorProofs.\allowbreak{}RamifiedSumsFifthEstimate}{30--45}

Since \(b=\exp(x^\varepsilon)\) dominates every fixed power of \(x\) and \(L=x^\varepsilon\to\infty\) at large \(x\), one has \(mfN_*^2\le b\) and \(2\log(2C_1)\le L\); likewise \(R_*=k+m-1=O(x^{9/2+10\varepsilon})\) and \(N_*\ge mt\ge2x^{377}\) give \(80R_*^4\le N_*\) for small \(\varepsilon\) and large \(x\). From \(m\ge2\) and \(k\ge f\ge m^4\ge16\), one has \(R_*=k+m-1\ge17\), so Definition~\ref{B8.21e} gives \(1\le n_p=(R_*+1)+m(R_*+5)\le2mR_*\) (as \(2mR_*-n_p=(m-1)(R_*-5)-6\ge6\)). Writing \(u=\sqrt b\ge100\) and using \((R_*+1)^4\le16R_*^4\le N_*\), \((R_*+1)^7\le N_*^3\), \(\sqrt{L_c}\le u\), and \(\sqrt m\le R_*+1\) bounds \(\gamma_{\rm gap}(C_1)+1\le2C_1N_*^3u\), whence \(1\le R_p(C_1)=(m+1)(\gamma_{\rm gap}(C_1)+1)+L_c+n_p\le C_1mN_*^3b\) and \(n_pR_p(C_1)\le2C_1u^7\), giving \(\log(n_pR_p)\le\log(2C_1)+\tfrac72L\le4L\).
\leanref{GlobalTheory.\allowbreak{}SizeLeafCheck.\allowbreak{}tree\_\allowbreak{}leaf\_\allowbreak{}parameter\_\allowbreak{}arithmetic\_\allowbreak{}of\_\allowbreak{}bounds}{Development.\allowbreak{}AbelianVarietyTheory.\allowbreak{}Counterexamples.\allowbreak{}DegenerateExamples.\allowbreak{}TreeLeafParameterArithmeticOfBounds}{66--202}
\leanref{GlobalTheory.\allowbreak{}BigTop.\allowbreak{}Fun.\allowbreak{}rankGapBound\_\allowbreak{}alt}{Development.\allowbreak{}AutomorphicArithmetic.\allowbreak{}ThetaCorrespondence.\allowbreak{}ThetaPairLifts.\allowbreak{}ThirdPlaceSeries.\allowbreak{}RamifiedSumsFifthBound}{16--127}

Summing \(\log N\mathfrak q\) over \(q\in S_0\) using \(\log g_i\le9L+1\), \(\log q_0\le L\), and \(\log p_{\rm de},\log p'_{\rm de}\le L\) gives \(\sum_{q\in S_0}\log N\mathfrak q\le f(21L+2)\), so \(c_S=\sum_{q\in S_0}c_q\log N\mathfrak q\le 2\sum_{q\in S_0}\log N\mathfrak q\le 42fL+4f\le45fL\) once \(L\ge4/3\); for \(C_1\ge42\), \(b\ge1\), \(f\ge1\), and \(k+m\ge2\), Definition~\ref{B8.21e} gives \(P_{\rm pf}\ge C_1f^2(k+m)^6L\ge45fL\ge0\) and \(e_{v_0}(C_1,b)\log N\mathfrak q_{v_0}\le P_{\rm pf}+\log N\mathfrak q_{v_0}\), whence \(\Sigma_0\le P_{\rm pf}+c_S\le P_{\rm pf}+45fL\le2P_{\rm pf}\).
\leanref{GlobalTheory.\allowbreak{}BigTop.\allowbreak{}Fun.\allowbreak{}RamSum.\allowbreak{}le\_\allowbreak{}PPF}{Development.\allowbreak{}AbelianVarietyTheory.\allowbreak{}Proofs.\allowbreak{}DivisorProofs.\allowbreak{}RamifiedSumsSeventhEstimate}{104--124}
\leanref{GlobalTheory.\allowbreak{}BigTop.\allowbreak{}Fun.\allowbreak{}RamSum.\allowbreak{}leaves\_\allowbreak{}of\_\allowbreak{}sum}{Development.\allowbreak{}AbelianVarietyTheory.\allowbreak{}Proofs.\allowbreak{}DivisorProofs.\allowbreak{}RamifiedSumsSeventhEstimate}{39--69}
\leanref{GlobalTheory.\allowbreak{}BigTop.\allowbreak{}Fun.\allowbreak{}ramified\_\allowbreak{}sums\_\allowbreak{}seventh\_\allowbreak{}bound\_\allowbreak{}alt}{Development.\allowbreak{}AutomorphicArithmetic.\allowbreak{}ThetaCorrespondence.\allowbreak{}ThetaPairLifts.\allowbreak{}ThirdPlaceSeries.\allowbreak{}RamifiedSumsSeventhBound}{18--38}

The row size conditions bound the gap-\(1\) bottom weight by \(1\le B_w\le10b\le C_1b\), and for \(L\ge\log10\) one has \(\log B_w\le\log10+L\le2L\le C_1L\). Finally, setting \(Q=x^{T+9}\) bounds \(m,f,k,t,L\le Q\), \(N_*\le10Q^2\), \(R_*\le2Q\), \(R_*N_*L\le20Q^4\), and \(P_{\rm pf}\le1280C_1Q^9\); for \(C_H\ge16(T+9)+1\) and \(x\ge13000000000C_1^4\), the polynomial prefactor is at most \(13000000000C_1^4Q^{16}\le x\cdot x^{16(T+9)}\le x^{C_H}\).
\leanref{GlobalTheory.\allowbreak{}BigTop.\allowbreak{}Fun.\allowbreak{}maxWt\_\allowbreak{}le}{Development.\allowbreak{}AbelianVarietyTheory.\allowbreak{}Proofs.\allowbreak{}DivisorProofs.\allowbreak{}WeightOfSizeBound}{92--129}
\leanref{GlobalTheory.\allowbreak{}BigTop.\allowbreak{}Fun.\allowbreak{}wt\_\allowbreak{}tie\_\allowbreak{}alt}{Development.\allowbreak{}AutomorphicArithmetic.\allowbreak{}ThetaCorrespondence.\allowbreak{}ThetaPairLifts.\allowbreak{}ThirdPlaceSeries.\allowbreak{}WtTie}{20--41}
\leanref{GlobalTheory.\allowbreak{}BigTop.\allowbreak{}Fun.\allowbreak{}desPol\_\allowbreak{}alt}{Development.\allowbreak{}AutomorphicArithmetic.\allowbreak{}ThetaCorrespondence.\allowbreak{}ThetaPairLifts.\allowbreak{}ThirdPlaceSeries.\allowbreak{}DesPol}{17--58}

The complete proof combines \Lean{GlobalTheory.\allowbreak{}BigTop.\allowbreak{}Fun.\allowbreak{}existsIntegralMultiple\_\allowbreak{}holds\_\allowbreak{}alt}, \Lean{GlobalTheory.\allowbreak{}BigTop.\allowbreak{}Fun.\allowbreak{}Close.\allowbreak{}desHLeaf\_\allowbreak{}alt}, \Lean{GlobalTheory.\allowbreak{}BigTop.\allowbreak{}Fun.\allowbreak{}rankGapBound\_\allowbreak{}alt} (via \Lean{GlobalTheory.\allowbreak{}SizeLeafCheck.\allowbreak{}tree\_\allowbreak{}leaf\_\allowbreak{}parameter\_\allowbreak{}arithmetic\_\allowbreak{}of\_\allowbreak{}bounds}), \Lean{GlobalTheory.\allowbreak{}BigTop.\allowbreak{}Fun.\allowbreak{}ramified\_\allowbreak{}sums\_\allowbreak{}seventh\_\allowbreak{}bound\_\allowbreak{}alt}, \Lean{GlobalTheory.\allowbreak{}BigTop.\allowbreak{}Fun.\allowbreak{}wt\_\allowbreak{}tie\_\allowbreak{}alt}, and \Lean{GlobalTheory.\allowbreak{}BigTop.\allowbreak{}Fun.\allowbreak{}desPol\_\allowbreak{}alt}.
\end{proof}

\subsection*{Proof of \ref{B8.22}}\label{prf:B8.22}\phantomsection
\begin{proof}
The base kernel equality and surjectivity in Proposition~\ref{prop:p8c2-fitted-evaluation} make \(\Psi\) injective. Set \(a_w=v_w(r_{\mu(w)})+d_F(w)\), supported on the twist places and \(T_F\). Equations \eqref{eq:p8c2-psi-comparison} and \eqref{eq:p8c2-evaluation-estimates} give \(\Psi(RR_w)\subseteq\mathfrak p_w^{-a_w}M_w\). Gluing gives
\begin{equation}\label{eq:p8c2-denominator-inclusion}
 \Psi(RR)\subseteq\left(\prod_w\mathfrak p_w^{-a_w}\right)M .
\end{equation}

At infinity divide the squared source norm by \(\kappa_v\). Put \(D_0'=3C_1fN_*^2L\). The two levels have the same finite-place index selector, infinite-place normalizations and block eigenvalues. In particular \(\lambda^{(0)}_{j,\mu_0,p,v}=\lambda^{(1)}_{j,\mu_0,p,v}\). By Theorem~\ref{B5.07}, \(L_{6,j}=\sum_\rho\log\operatorname{Hook}(a^+_{j\rho})-C_{\rm sz}fn_V^2L_p\), where \(L_p=\log(b\,\operatorname{rd}(E)g_0g_1n_V)\). Each hook product is at least one. Use \(L_p\le3cL,\ n_V\le N_*,\ C_1\ge c\max(0,C_{\rm sz})\) for the family constants. Hence \(L_{6,j}\ge-D_0'\). Summing the assumed eigenvalue bound over the layers \(j\in\mathcal J\) of a nonzero base block gives
\begin{equation}\label{eq:p8c2-base-block}
 \sum_{j\in\mathcal J}\log\lambda^{(0)}_{j,\mu_0,p,v}
 \ge-|\mathcal J|D_0'\ge-mD_0'.
\end{equation}
If \(g(v)\) bounds the block sums below, the squared quotient-norm factor is at least \(\exp g(v)\). Minimizing over lifts gives \(\|\Psi\|_v\le\|\Psi\iota_{\mu_0}\|_v e^{-g(v)/2}\). Taking mean log squared norms with \(g(v)=-mD_0'\) gives \(h(\Psi)\le H_I=H_{I,0}+mD_0'\).
\leanref{Formalism.\allowbreak{}L2.\allowbreak{}Indef.\allowbreak{}Tree.\allowbreak{}ip\_\allowbreak{}first}{Formalism.\allowbreak{}LocalToGlobal.\allowbreak{}IndefiniteForms.\allowbreak{}IndefiniteTree.\allowbreak{}FunctionalLowerBoundOfConditionsNumeric}{826--909}

For each bottom number the weighted negative part is at most \(3D_0\), and the signed sum is at most \(A_0\) by the product formula and \eqref{eq:p8c2-bottom-budget}. Thus the absolute sum is at most \(A_0+6D_0\). For the base and both higher twists Lemma~\ref{B8.21a} gives \(B_0\) with common shift \(g_T\). Set \(B=\max(A_0+6D_0,B_0)\).

Sum \eqref{eq:p8c2-one-pair}. The two bottom terms cost \(2B\). Sum the selected twist over the subset where it is used to bound the higher terms at each \(s\) by \(4(B+2B)\). Together these fit \(2(2+8N_*)B\) per pair. The prime-count bound with coefficient five fits \(10N_*B_0\). For the degree term,
\begin{equation}\label{eq:p8c2-degree-sum}
 \sum_{p\in S_e}\frac1{p^2}\sum_{w\mid p}\ell_wv_w(p)
 =\sum_{p\in S_e}\frac{\log p}{p^2}
 \le\sum_{j=2}^4\frac{\log j}{j^2}+\frac{1+\log4}{4}<1 .
\end{equation}
For \(j\ge5\) the integral of \((\log t)/t^2\) from \(j-1\) to \(j\) bounds the summand and the integrals telescope. The degree term costs at most \(2N_*\). The last two terms cost at most \(6D_0\) and \(4m\Sigma_0\). Hence
\begin{equation}\label{eq:p8c2-grid-budget}
 \sum_w\ell_w v_w(r_{\mu(w)})\le
 G_*:=m^2\bigl(2(2+8N_*)B+10N_*B_0+2N_*+6D_0+4m\Sigma_0\bigr).
\end{equation}
\leanref{Formalism.\allowbreak{}PairingDivisionLocalBalance.\allowbreak{}wt}{Formalism.\allowbreak{}PairingDivisors.\allowbreak{}StatementH}{16--17}
\leanref{Formalism.\allowbreak{}ValuationGridEstimates.\allowbreak{}sum\_\allowbreak{}log\_\allowbreak{}div\_\allowbreak{}sq\_\allowbreak{}telescoping}{Formalism.\allowbreak{}TowerValuationBounds.\allowbreak{}TowerValuationProofs.\allowbreak{}SumAdditiveCharacterEqZero}{504--511}

The frame change adds \(\langle\log\kappa\rangle_\infty\) to the source slope. The ideal in \eqref{eq:p8c2-denominator-inclusion} adds \(2\sum_w\ell_w a_w\) to the target slope. For each rank \(r\) sublattice, the determinant norm of the injection at \(v\) is at most \(\|\Psi\|_v^r\). Its source slope is bounded by its image slope plus \(h(\Psi)\). Using the height, lattice and finite-place budgets gives
\begin{equation}\label{eq:p8c2-functional-budget}
 \mu_{\max,RR}+\langle\log\kappa\rangle_\infty
 \le E_M+H_I+2G_*+B_F .
\end{equation}

The prime-power modulus budget is
\[
 P_{\rm pf}=C_1f^2(k+m)^6
 \log\bigl(b(k+m)g_0g_1p_{\rm de}p'_{\rm de}|D_E|\bigr).
\]
For \(C_1\ge42\), Lemma~\ref{B8.21p} gives the fitted size bounds
\begin{equation}\label{eq:p8c2-functional-sizes}
 \begin{gathered}
 1\le f\le k,\quad 1\le m\le N_*,\quad R_*\le N_*,\quad mfN_*^2\le b,\quad
 L\ge1,\quad L+\log x\le3L,\\
 \Sigma_0\le2P_{\rm pf},\quad c_S\le45fL,\quad
 1\le B_w\le C_1b,\quad \log B_w\le C_1L,\\
 n_p\le2mR_*,\quad R_p\le C_1mN_*^3b,\quad \log(n_pR_p)\le4L .
 \end{gathered}
\end{equation}
In particular \(l_1\le1+L,\ l_2,l_3\le1+(C_1+1)L\). Set \(U=C_1^3fm^2R_*N_*^3bL\) and \(V=C_1^3fm^2N_*^2P_{\rm pf}b\). Substitution gives
\begin{equation}\label{eq:p8c2-budget-table}
 \begin{array}{ll}
 \text{Quantity}&\text{Upper bound}\\
 D_0&186C_1^2f^2N_*^2bL\\
 A_0&57C_1^3fN_*^2bL\\
 B_0&3C_1fbL\\
 E_M&32U\\
 H_I&4U+2V\\
 B_F&2U+4V\\
 G_*&(23520+30+2+1116)U+8V .
 \end{array}
\end{equation}
For the last row use \eqref{eq:p8c2-grid-budget}, \(B\le1176C_1^3f^2N_*^2bL,\ 2(2+8N_*)\le20N_*\) and \(f\le R_*\).

Let \(C_t\) be the family's growth exponent for the auxiliary rank \(t\). Set \(T=\lceil|C_t|\rceil+1\) and \(Q=x^{T+9}\). The family sizes give \(m\le x,\ f,k\le x^5,\ t\le x^T\) and \(L\le x\), hence all five quantities are at most \(Q\). The logarithms of \(b\) and the four auxiliary primes are at most \(x\) each. Also \(\log(k+m)\le2x^5\) and \(\log|D_E|\le6x^8\) by the root-discriminant bound. The logarithm in \(P_{\rm pf}\) is at most \(20Q\). Therefore \(N_*\le10Q^2,\ R_*\le2Q\). It follows that
\[
 R_*N_*L\le20Q^4,\qquad
 P_{\rm pf}\le C_1Q^2(2Q)^6(20Q)=1280C_1Q^9 .
\]
The prefactor \(100000 C_1^3fm^2N_*^2\) is at most \(10^7 C_1^3Q^7\). Since \(Q,C_1\ge1\), the product with the sum in parentheses is at most \(13000000000C_1^4Q^{16}\). The family threshold gives \(C_H\ge16(T+9)+1\) and the point threshold gives \(x\ge13000000000C_1^4\). With the preceding table this yields
\begin{equation}\label{eq:p8c2-functional-polynomial}
 E_M+H_I+2G_*+B_F
 \le100000C_1^3fm^2N_*^2(R_*N_*L+P_{\rm pf})b
 \le x^{C_H}b .
\end{equation}
Choose \(C_1\ge\max(42,c\max(0,C_{\rm sz}))\) above the joint estimate threshold and take the fitted-grid and size thresholds at that constant. Equations \eqref{eq:p8c2-functional-budget} and \eqref{eq:p8c2-functional-polynomial} give \eqref{eq:p8c2-fun} for every \(C\).
\end{proof}

\subsection*{Proof of \ref{B8.23}}\label{prf:B8.23}\phantomsection
\begin{proof}
Use the threshold of Theorem~\ref{B8.22}. The extension with list, ray-class fields and real-place compatibility transport the divisor row to the geometry row. Intersect with the goodness, primitive and divisor rows to obtain \(R,E\). The geometry row supplies a compatible selection \(\tau'\) of comparison data on the chosen members, with point summands, fitted grids and ties to the chosen bundles. Discrete-series realization and globalization give the eventual point-summand conditions from the real realization, quotient and division data. The class assumptions give the seven family conditions of Definition~\ref{B8.01k}. Apply Theorem~\ref{B8.22} at those fitted points using the assumed division bound. The ties in \(\tau'\) transfer the functional condition to the original bundles with \(C_0=0\).
\end{proof}

\section*{Proofs for Chapter~\ref{ch:B8-covolume}: Covolume bounds}

\subsection*{Proof of \ref{B8.24a}}\label{prf:B8.24a}\phantomsection
\begin{proof}
Under the flat-degree law, the difference of the conjugate Lie eigenlattice degrees at a CM point of good reduction contributes \((P-Q)\log\Omega/2\) by \ref{B5.07x}. The local different terms in the frame comparison contribute \(Q(\log d-\log\operatorname{rd}(E))\). The first Chern value is the negative of their sum.

Put \(\beta=P-Q\) and let \(\rho(\tau)\) be the row index matching \(\tau\). The host and anchor comparisons give
\begin{equation}\label{eq:p8c3-anchor}
 \begin{gathered}
 L_x\ge10,\qquad \log\log x\ge1,\qquad E_p=fmbL_x\log\log x,\\
 0\le\log\operatorname{rd}(E)\le2L_x,\qquad
 |\log d-\log\operatorname{rd}(E)|\le3L_x,\qquad |\log\Omega|\le3L_x,\\
 |a_\tau|\le4b,\qquad (P+Q)|a_\tau-w_{j,\rho(\tau)}|\le20b .
 \end{gathered}
\end{equation}
For the last comparison, the discrepancy before multiplication by \(r\) is bounded by the scale width plus \(n_V+m+t+2\). After multiplication the scale contributes at most \(16b\) and the other terms at most \(4rn_V\le4b\), by \eqref{eq:p8c3-host-size} at the family threshold. Summing over the \(f\) embeddings gives
\begin{equation}\label{eq:p8c3-anchor-errors}
 \left|M_{p,i}-\beta\log\Omega\sum_\tau a_\tau\right|
 \le60fbL_x\le100E_p,\qquad
 QL_x\sum_\tau |a_\tau|\le4fmbL_x\le100E_p .
\end{equation}

The integral twist adds the anchor at each embedding and the scalar exponent \(m_0\), where \(m_0\) is half the even column count. This gives \(G_s(-c_0)\) at the level of section exponents and preserves rank. If \(\mathcal S_{\rm bad}\) is the set of bad primes of the covering level and \(\ell_{\rm rd}=\log\operatorname{rd}(E)\), the degree comparison divided by the positive common rank gives
\begin{equation}\label{eq:p8c3-flat-twist}
 \begin{split}
 \mu_Y(G_s(-c_0))-\mu_Y(u_M)
 \ \ge\ &
 2\sum_\tau a_\tau c_{1,\tau}-2Q\ell_{\rm rd}\sum_\tau|a_\tau|\\
 &-2m_0fr\bigl(\tfrac12\log(2\pi)+\ell_{\rm rd}\bigr)
 -2\sum_{\ell\in\mathcal S_{\rm bad}}\log\ell .
 \end{split}
\end{equation}
By \eqref{eq:p8c3-flat-c1} and \eqref{eq:p8c3-anchor-errors}, the first term is at least \(-M_{p,i}-700E_p\). The second costs at most \(400E_p\) by \eqref{eq:p8c3-anchor}. The column and family size bounds, together with the deep goodness bound on the bad primes, give the following estimates using \(2m_0\le n_V\)
\begin{equation}\label{eq:p8c3-flat-costs}
 2m_0fr\bigl(\tfrac12\log(2\pi)+\ell_{\rm rd}\bigr)
 \le3frn_VL_x\le3b\le3E_p,\qquad
 2\sum_{\ell\in\mathcal S_{\rm bad}}\log\ell\le80L_x\le80E_p.
\end{equation}
Substitution in \eqref{eq:p8c3-flat-twist} gives the loss \(700+400+3+80=1183\).
\end{proof}

\subsection*{Proof of \ref{B8.24b}}\label{prf:B8.24b}\phantomsection
\begin{proof}
Work first over \(F'\). The coefficient field contains the images of the embeddings of \(E\), so the idempotents for the \(E\)-action split de Rham homology over \(F'\). Let \(H_\tau\) be its \(\tau\)-summand. The Hodge sequence for homology, the polarisation, and the fact that the Rosati involution acts on \(E\) by conjugation give an exact sequence
\[
                 0\longrightarrow V_{\bar\tau}^{\vee}
                 \longrightarrow H_\tau\longrightarrow V_\tau
                 \longrightarrow0 .
\]
In particular \(\det H_\tau={\cal L}_\tau\). In bases \(v,w\) of the two Lie eigenspaces a section has coefficient \(c(v,w)\) relative to \(\bigwedge v\otimes(\bigwedge w)^{-1}\). Changing the bases by \(D_1,D_2\) multiplies that coefficient by \(\det D_2/\det D_1\). This is the flat row character.

Fix a complex embedding of \(F'\). The monodromy of \(H_\tau\) is the \(\tau\)-realisation of the unitary lattice action in the uniformisation. Its determinant is one. Indeed, the stabilizer preserves the period lattice and acts trivially on its principal level of modulus \(M\ge3\). Localising the lattice at each prime of \(E\) gives
\(\det\gamma=1+Ma\) with \(a\in\mathcal O_E\). Unitarity gives
\((\det\gamma)\overline{\det\gamma}=1\), so every conjugate of \(a\) has absolute value at most \(2/M<1\). Its integral norm forces \(a=0\). Thus the exterior product of a flat basis gives a single-valued horizontal generator of \(\det H_\tau\) on the complex component.
\leanref{GlobalTheory.\allowbreak{}LawsDeep.\allowbreak{}laws\_\allowbreak{}deep}{Development.\allowbreak{}AbelianVarietyTheory.\allowbreak{}Proofs.\allowbreak{}DivisorProofs.\allowbreak{}ComplexUniformizationLawsOfGoodDeepCover}{41--52}
\leanref{GlobalTheory.\allowbreak{}UnitaryCongruenceDetAny.\allowbreak{}det\_\allowbreak{}eq\_\allowbreak{}one}{Development.\allowbreak{}AbelianVarietyTheory.\allowbreak{}Proofs.\allowbreak{}DivisorProofs.\allowbreak{}UnitaryCongruenceDetEqOneOfThreeLe}{11--91}

Here one can algebraize the horizontal generator by the regular-singular comparison theorem. The base is smooth of finite type in characteristic zero. The abelian scheme is smooth and proper, so it supplies its own relative smooth compactification with empty relative boundary. Its Gauss--Manin connection is regular \cite[II, Th.~7.9]{deligne-1970}. Duals, direct summands and exterior powers preserve regularity. Full faithfulness in \cite[II, Th.~5.9]{deligne-1970}, applied to the trivial connection and \(\det H_\tau\), identifies the single-valued horizontal sections with the algebraic horizontal sections. This space has dimension one because \(Z\) is geometrically connected. To descend to \(F'\), compute horizontal sections on a finite affine cover as the kernel of the connection and the differences of restrictions on overlaps. Extension of the coefficient field commutes with this kernel by flatness. A basis over \(F'\) therefore becomes a nonzero multiple of the horizontal generator over \(\C\). It generates at every point by faithful flatness. Its reciprocal is a regular section of the dual line.

At an infinite place let \(J\) be the involution which is \(+1\) on the positive Hodge subspace and \(-1\) on the negative one for the polarising Hermitian form \(H\). The positive Hermitian metric is \(H(\,\cdot\,,J\,\cdot\,)\), up to the fixed period normalisation at that place. Its determinant in a flat basis is \(|\det H|\), since \(|\det J|=1\). The determinant metric in the Hodge sequence is the product of the subspace and quotient metrics, as one sees in an orthogonal splitting. Hence in the Lie frames its squared norm is
\[
       |c(v,w)|^2\,\frac{\det\operatorname{Gram}(v)}
                              {\det\operatorname{Gram}(w)} .
\]
It is constant on \(Z\). The same calculation applies after every embedding of \(F'\). Extending the section and its reciprocal by zero to the other components gives two sections whose product is the idempotent of \(Z\).
\end{proof}

\subsection*{Proof of \ref{B8.24c}}\label{prf:B8.24c}\phantomsection
\begin{proof}
Completion preserves the valuation of a determinant. The maximal unramified order splits after tensoring with \(R\), since all its embeddings take values in \(R\) and their distinct residual embeddings are separated by units. In the selected factor the remaining extension is totally ramified and its integer ring is generated by \(\pi\). The resulting algebra is
\[
              B=R[t]\big/\prod_{\sigma\in S}(t-a_\sigma).
\]
All the following modules are taken in this factor. Write \(T=t-a_\tau\). Evaluation at \(a_\tau\) has kernel \(TB\). Since \(H\) is free over \(B\), it identifies \(H/TH\) integrally with \(H(\tau)\).

Put \(F=\ker(H\to L)\), and choose an \(R\)-basis of \(L[\tau]=\ker(T|L)\). Lift it to \(H\) and project the lifts to \(H(\tau)\). Adjoin to these lifts a basis of \(\operatorname{pr}_\tau F\). This last lattice identifies with \(L[\bar\tau]^\vee\). Indeed, projection of a functional to the \(\tau\)-summand amounts to restriction to the \(\bar\tau\)-summand of \(L_K\). Every integral functional on \(L[\bar\tau]\) extends to \(L\), since \(L[\bar\tau]\) is saturated in the free \(R\)-module \(L\). Thus the determinant of the assembled vectors is the determinant of the stated pair of Lie frames. Changing a lift adds a vector in \(\operatorname{pr}_\tau F\) and preserves this determinant.

Denote their span by \(A\). Taking the quotient first by \(F\), then by the lifted basis, gives
\[
              H(\tau)/A \simeq L/(TL+L[\tau])
                    =\operatorname{coker}\bigl(T:L/L[\tau]\to L/L[\tau]\bigr).
\]
The last module on which \(T\) acts is free. On its generic fibre \(T\) has eigenvalues \(a_\sigma-a_\tau\), with multiplicity \(h_\sigma\), for \(\sigma\ne\tau\) in \(S\). Diagonal reduction over \(R\) identifies the length of the cokernel with the valuation of that determinant. Dividing \(b\) by the determinant of the assembled frame gives the asserted sign.
\end{proof}

\subsection*{Proof of \ref{B8.24dd}}\label{prf:B8.24dd}\phantomsection
\begin{proof}
Every \(\sigma:E\hookrightarrow\Omega\) factors through \(F'\), and since \(O\) is integrally closed in \(F'\), \(\sigma(\mathcal O_E)\) lies in the image of \(O\). Let \(R\) be any commutative \(\F_\ell\)-algebra and \(B\) a \(d\)-dimensional formal biextension germ over \(R\). Write \(\mathbf x^{(\ell)}=(x_1^\ell,\ldots,x_d^\ell)^{\mathsf t}\) and \(M^{(\ell)}=(M_{ab}^\ell)\) for a matrix \(M\) over \(R\). Pullback of the basis of invariant differentials of \(\mu_B\) along \([\ell]_B\) multiplies them by \(\ell=0\), so all first partial derivatives of \([\ell]_B(\mathbf x)\) vanish; since every integer prime to \(\ell\) is invertible in \(R\), every nonzero monomial of \([\ell]_B(\mathbf x)\) has all exponents divisible by \(\ell\), giving \([\ell]_B(\mathbf x)\equiv V_1\mathbf x^{(\ell)}\pmod{(\mathbf x)^{2\ell}}\) for a unique \(V_1\in\operatorname{Mat}_d(R)\). Modulo \((\mathbf y)^2\), \(\varphi_1(\mathbf x',\mathbf x'',\mathbf y)\equiv1+c(\mathbf x',\mathbf x'')^{\mathsf t}\mathbf y\) defines a symmetric formal \(2\)-cocycle \(c\) and hence a \(2d\)-dimensional commutative formal group law on \(R[[\mathbf x]]\times\widehat{\mathbb G}_a^d\) whose multiplication by \(\ell\) on \((\mathbf x,0)\) has second component \(\sum_{k=1}^{\ell-1}c([k]_B\mathbf x,\mathbf x)\); vanishing of its invariant differentials under \([\ell]\) shows that \(\Psi_\ell(\mathbf x,\mathbf y)=\prod_{k=1}^{\ell-1}\varphi_1([k]_B\mathbf x,\mathbf x,\mathbf y)\) satisfies
\[
 \Psi_\ell(\mathbf x,\mathbf y)\equiv1+(\mathbf x^{(\ell)})^{\mathsf t}W_1\mathbf y\pmod{(\mathbf x)^{2\ell}+(\mathbf y)^2}
\]
for a unique \(W_1\in\operatorname{Mat}_d(R)\). By the same argument on the second factor, \(\widetilde\Psi_\ell(\mathbf x,\mathbf y)=\prod_{k=1}^{\ell-1}\varphi_2(\mathbf x,[k]_B\mathbf y,\mathbf y)\) satisfies
\[
 \widetilde\Psi_\ell(\mathbf x,\mathbf y)\equiv1+\mathbf x^{\mathsf t}\widetilde W_1\mathbf y^{(\ell)}\pmod{(\mathbf x)^2+(\mathbf y)^{2\ell}}
\]
for a unique \(\widetilde W_1\in\operatorname{Mat}_d(R)\). Put
\[
 \mathcal V_B=\begin{pmatrix}V_1\\ W_1^{\mathsf t}\end{pmatrix}:(R^d)^{(\ell)}\longrightarrow R^d\times R^d,\qquad
 \mathcal F_B=\begin{pmatrix}-\widetilde W_1^{\mathsf t} & V_1^{\mathsf t}\end{pmatrix}:R^d\times R^d\longrightarrow R^d.
\]
Telescoping the four-variable biextension interchange law first over \(\mathbf y'=[k]_B\mathbf y,\ \mathbf y''=\mathbf y\) (\(1\le k<\ell\)) and then over \(\mathbf x'=[k]_B\mathbf x,\ \mathbf x''=\mathbf x\) (\(1\le k<\ell\)) yields \(\Psi_\ell(\mathbf x,\mathbf y)^\ell\,\widetilde\Psi_\ell([\ell]_B\mathbf x,\mathbf y)=\widetilde\Psi_\ell(\mathbf x,\mathbf y)^\ell\,\Psi_\ell(\mathbf x,[\ell]_B\mathbf y)\). Since \(\varphi_1(\mathbf x',\mathbf x'',\mathbf y)-1\in(\mathbf x')(\mathbf x'')(\mathbf y)\), for every prime \(\ell\ge2\) one has \((\Psi_\ell(\mathbf x,\mathbf y)-1)^\ell\in(\mathbf x)^{2\ell}(\mathbf y)^\ell\) and \((\widetilde\Psi_\ell(\mathbf x,\mathbf y)-1)^\ell\in(\mathbf x)^\ell(\mathbf y)^{2\ell}\); comparing \((\mathbf x^{(\ell)},\mathbf y^{(\ell)})\)-coefficients modulo \((\mathbf x)^{2\ell}+(\mathbf y)^{2\ell}\) gives \(V_1^{\mathsf t}\widetilde W_1=W_1 V_1\), hence \(\mathcal F_B\circ\mathcal V_B=0\). Similarly, for any \(g(\mathbf x,\mathbf y)\in 1+(\mathbf x)(\mathbf y)R[[\mathbf x,\mathbf y]]\) and any prime \(\ell\ge2\), \(g(\mathbf x,\mathbf y)^\ell-1\in(\mathbf x)^\ell(\mathbf y)^\ell\) vanishes both modulo \((\mathbf y)^2\) and modulo \((\mathbf x)^2\). Thus telescoping the left Rosati cocycle modulo \((\mathbf x)^{2\ell}+(\mathbf y)^2\), the right Rosati cocycle modulo \((\mathbf x)^2+(\mathbf y)^{2\ell}\), the symmetry cocycle \(\widetilde\Psi_\ell(\mathbf x,\mathbf y)\,\mathrm{sy}(\mathbf x,[\ell]_B\mathbf y)=\mathrm{sy}(\mathbf x,\mathbf y)^\ell\,\Psi_\ell(\mathbf y,\mathbf x)\) modulo \((\mathbf x)^2+(\mathbf y)^{2\ell}\), and the \(\varphi_1,\varphi_2\) transition cocycles of any formal biextension isomorphism \((\varphi,\psi,\lambda):B\xrightarrow{\sim}B'\) gives, for all \(a\in\mathcal O_E\):
\[
 \begin{gathered}
 \rho_H(B)(a)\circ\mathcal V_B=\mathcal V_B\circ\rho_B(a)^{(\ell)},\qquad
 \mathcal F_B\circ\rho_H(B)(a)=(\rho_B(\bar a)^{(\ell)})^{\mathsf t}\circ\mathcal F_B,\\
 \mathcal F_B=\mathcal V_B^{\mathsf t}\begin{pmatrix}\operatorname{bil}(\mathrm{sy})_{1,1}^{\mathsf t} & I_d\\ -I_d & 0\end{pmatrix},\qquad
 \operatorname{transH}(\varphi,\psi,\lambda)\circ\mathcal V_{B'}=\mathcal V_B\circ D\varphi(0)^{(\ell)},\\
 \mathcal F_B\circ\operatorname{transH}(\varphi,\psi,\lambda)=(D\psi(0)^{(\ell)})^{\mathsf t}\circ\mathcal F_{B'}.
 \end{gathered}
\]

On the regular proper integral curve \(\widetilde\Gamma\) over \(k'\), let \(\kappa=k'(\widetilde\Gamma)\) be its function field and \(\mathcal O_i=\mathcal O_{\widetilde\Gamma,i}\) (\(i\in\widetilde\Gamma\) closed) its discrete valuation rings with residue fields \(k_i\). Base-changing the stalk germs and specialisation isomorphisms of \(S\) from Lemma~\ref{B8.24dc}(2) along \(\operatorname{Spec}\kappa\to X\) and \(\operatorname{Spec}\mathcal O_i\to X\) yields formal biextension germs \(B_i\) over \(\mathcal O_i\) with reductions \(B_{k_i}\) over \(k_i\) and compatible transitions to \(B_\kappa\). Fix a closed point \(i\) and an algebraic closure \(\bar k_i\) of \(k_i\), with geometric point \(\bar x:\operatorname{Spec}\bar k_i\to X\) and extension \(0\to T\to G_{\bar x}\xrightarrow{q_{\bar x}}B_{\rm ab}\to0\) as in Lemma~\ref{B8.24dc}(d). In formal parameters \((\mathbf x_T,\mathbf x_B)\) adapted to \(T\simeq\mathbb G_{m,\bar k_i}^{d_T}\) and \(B_{\rm ab}\) with trivialisation pulled back from \(\Lambda(M)\) on \(\widehat B_{\rm ab}\times\widehat B_{\rm ab}\), one has \([\ell]_T(\mathbf x_T)=\mathbf x_T^{(\ell)}\) so \(V_{1,T}=I_{d_T}\), while \(W_1\) vanishes on \(T\)-variables; hence \(\operatorname{rk}_{\bar k_i}\mathcal V_{B_{\bar x}}=d_T+\operatorname{rk}_{\bar k_i}\mathcal V_{B_{\rm ab}}\). Under the identification \(\operatorname{Lie}(B_{\rm ab})\times\operatorname{Lie}(B_{\rm ab})^\vee\simeq\operatorname{Lie}(E(B_{\rm ab}))\) of Lemma~\ref{B8.24dc}, evaluating \([\ell]_{E(B_{\rm ab})}\) on \((\mathbf x_B=\zeta w,1)\) over \(\bar k_i[\zeta]/(\zeta^{2\ell})\) (with \(\varepsilon=\zeta^\ell\)) identifies \(\mathcal V_{B_{\rm ab}}\) with the map \(\operatorname{Lie}(B_{\rm ab})^{(\ell)}\to\operatorname{Lie}(E(B_{\rm ab}))\) induced by Verschiebung on \(\mathbb D(B_{\rm ab}[\ell^\infty])(W(\bar k_i))^\vee/\ell\) \cite[Prop.~2.5.2]{berthelot-breen-messing-1982}, which is injective of rank \(d_B=\dim B_{\rm ab}\). Hence \(\operatorname{rk}_{k_i}\mathcal V_{B_{k_i}}=d_T+d_B=d\), and the matrix relation between \(\mathcal F_{B_{k_i}}\) and \(\mathcal V_{B_{k_i}}^{\mathsf t}\) gives \(\operatorname{rk}_{k_i}\mathcal F_{B_{k_i}}=d\). By Nakayama's lemma over each discrete valuation ring \(\mathcal O_i\) and compatibility with \(\operatorname{transH}\), the maps \(\mathcal V_{B_i}\) and \(\mathcal F_{B_i}\) glue to an \(\mathcal O_E\)-equivariant exact sequence of vector bundles on \(\widetilde\Gamma\),
\[
 0\longrightarrow F_{\widetilde\Gamma}^*\mathcal L_{\widetilde\Gamma}
 \xrightarrow{\;\mathcal V\;}\mathcal H_{\widetilde\Gamma}
 \xrightarrow{\;\mathcal F\;}F_{\widetilde\Gamma}^*\bigl(\mathcal L_{\widetilde\Gamma}^\vee(\mathrm{conj})\bigr)
 \longrightarrow0,
\]
where \(F_{\widetilde\Gamma}:\widetilde\Gamma\to\widetilde\Gamma\) is the absolute \(\ell\)-power Frobenius morphism and \(\mathcal L_{\widetilde\Gamma}=\nu^*(\mathcal L|_\Gamma)\), \(\mathcal H_{\widetilde\Gamma}=\nu^*(\mathcal H|_\Gamma)\).

For every \(a\in\mathcal O_E\), the characteristic polynomial \(P_a(T)\in\Z[T]\) of multiplication by \(a\) on \(\mathcal O_E\) splits in \(\Omega[T]\) as \(\prod_{\sigma':E\hookrightarrow\Omega}(T-\sigma'(a))\); since \(\sigma'(a)\in O\) and \(O\hookrightarrow\Omega\) is injective, \(P_a(T)\) splits in \(O[T]\) and hence over the subfield \(k_1\subseteq k_0\subseteq k'\) generated over \(\F_\ell\) by the reductions modulo \(\varpi\) of \(\{\sigma'(b):\sigma':E\hookrightarrow\Omega,\ b\in\mathcal O_E\}\). Because \(\mathcal O_E\) is finite over \(\Z\), \(k_1\) is a finite field. Thus \(\mathcal O_E\otimes_{\Z}k_1=\prod_{\nu=1}^s A_\nu\) decomposes into local Artin \(k_1\)-algebras \(A_\nu=\epsilon_\nu(\mathcal O_E\otimes_{\Z}k_1)\) with residue fields \(A_\nu/\mathfrak m_\nu=k_1\), and \(\mathrm{id}_{\mathcal O_E}\otimes F_{k_1}\) permutes the primitive idempotents by \(\epsilon_\nu\mapsto\epsilon_{\pi(\nu)}\). For any \(\mathcal O_E\)-bundle \(\mathcal E\) on \(\widetilde\Gamma\) we have \(\epsilon_{\pi(\nu)}F_{\widetilde\Gamma}^*\mathcal E=F_{\widetilde\Gamma}^*(\epsilon_\nu\mathcal E)\), while \(\deg_{\widetilde\Gamma}(F_{\widetilde\Gamma}^*\mathcal E)=\ell\deg_{\widetilde\Gamma}(\mathcal E)\) because \(F_{\widetilde\Gamma}^*\) raises transition determinants to the \(\ell\)-th power at every valuation of \(\kappa\). Comparing the conjugate sequence above with the pulled-back Hodge sequence \(0\to\mathcal L_{\widetilde\Gamma}^\vee(\mathrm{conj})\to\mathcal H_{\widetilde\Gamma}\to\mathcal L_{\widetilde\Gamma}\to0\) of Lemma~\ref{B8.24dc}(3) gives
\[
 \deg_{\widetilde\Gamma}(\epsilon_{\pi(\nu)}\mathcal H_{\widetilde\Gamma})
 =\ell\deg_{\widetilde\Gamma}(\epsilon_\nu\mathcal L_{\widetilde\Gamma})
  +\ell\deg_{\widetilde\Gamma}\bigl(\epsilon_\nu\mathcal L_{\widetilde\Gamma}^\vee(\mathrm{conj})\bigr)
 =\ell\deg_{\widetilde\Gamma}(\epsilon_\nu\mathcal H_{\widetilde\Gamma}).
\]
Iterating around the finite \(\pi\)-cycle of \(\nu\) yields \(\deg_{\widetilde\Gamma}(\epsilon_\nu\mathcal H_{\widetilde\Gamma})=0\) for every \(\nu\). Let \(A_{\nu_0}\) be the factor through which the residual homomorphism \(\bar\sigma:\mathcal O_E\to O\to k_1\) factors, and choose a \(k_1\)-composition series \(A_{\nu_0}=I_0\supset I_1\supset\cdots\supset I_{e_0}=0\) (\(e_0=\dim_{k_1}A_{\nu_0}\ge1\)), so each quotient \(I_{t-1}/I_t\) is isomorphic to \(k_1\) with \(\mathcal O_E\) acting via \(\bar\sigma\). By Lemma~\ref{B8.24dc}(3), \(\epsilon_{\nu_0}\mathcal H_{\widetilde\Gamma}\) is locally free of rank \(r\) over \(A_{\nu_0}\otimes_{k_1}\mathcal O_{\widetilde\Gamma}\), hence flat over \(A_{\nu_0}\); tensoring the composition series over \(A_{\nu_0}\) with \(\epsilon_{\nu_0}\mathcal H_{\widetilde\Gamma}\) filters it by \(e_0\) subbundles with successive quotients isomorphic to \(\nu^*(\mathcal H\otimes_{\mathcal O_E,\sigma}\mathcal O_X)\). Hence \(e_0\deg_{\widetilde\Gamma}(\nu^*\mathscr D_\sigma)=\deg_{\widetilde\Gamma}(\epsilon_{\nu_0}\mathcal H_{\widetilde\Gamma})=0\) and \(\deg_{\widetilde\Gamma}(\nu^*\mathscr D_\sigma)=0\).
\end{proof}

\subsection*{Proof of \ref{B8.24d}}\label{prf:B8.24d}\phantomsection
\begin{proof}
Let \(O=\mathcal O_{F',w}\subset F'\) be the valuation ring at \(w\), with uniformiser \(\varpi\mid\ell\) and residue field \(k_0=O/\varpi O\), and let \((P,Q)\) be the signature of \(Z\) with \(r=P+Q\). If \(r=0\), both Lie eigenspaces are zero, \(Q=0\), and \(s\in(F')^\times\) is a nonzero constant with \(\rho_w=v_w(s)\), vanishing outside finitely many places. Assume now that \(r\ge1\), and put \(g=\tfrac12[E:\Q]r\). By the definition of a good cover, the covering level is \(M'=g_0g_1 15^e\) where \(g_0,g_1>3\) are the two distinct odd base-level primes, and \(F'\) contains all complex embeddings of \(E\). Choose \(n\in\{g_0,g_1\}\) with \(n\ne\ell\); then \(n\ge5\) is odd and invertible in \(O\), and \(\zeta_n\in F'\) (for a general good cover without \(\zeta_n\in F'\), replacing \(F'\) by \(F'(\zeta_n)\) at a place \(w'\mid w\) preserves integrality of \(Z\) because \(F'\) is relatively algebraically closed in \(\Omega\) and yields \(\rho_w=\rho_{w'}/e(w'/w)\in\Q\)). Since \(O\) is integrally closed in \(F'\), every embedding \(\sigma:E\hookrightarrow\Omega\) into the function field \(\Omega\) of \(Z\) factors through \(F'\) and maps \(\mathcal O_E\) into \(O\), and \(\zeta_n\in O\).

Base-changing the integral Siegel toroidal compactification of Lemma~\ref{B8.25a} at odd full level \(n\) \cite[Chs.~IV--V]{faltings-chai-1990} along \(\Z[\zeta_n,1/n]\to O\) yields a projective smooth \(O\)-scheme \(\mathcal S\) containing \(\mathcal A_{g,n,O}\) as a fibrewise dense open subscheme equal to the abelian locus of the universal semiabelian scheme \(G_{\mathcal S}\to\mathcal S\). Because \(\mathcal S\) is regular and \(G_{\mathcal S}^{\times_{\mathcal S}m}\to\mathcal S\) (\(m\ge1\)) is smooth with geometrically integral fibres, every prime divisor of \(G_{\mathcal S}^{\times_{\mathcal S}m}\) outside \(\mathcal A_{g,n,O}\) is the preimage of a boundary divisor of \(\mathcal S\); extending a Weil divisor on the regular scheme \(G_{\mathcal S}\times_{\mathcal S}G_{\mathcal S}\) and twisting along \((e_{\mathcal S},e_{\mathcal S})\) therefore extends the universal symmetric Poincar\'e biextension on \(\mathcal A_{g,n,O}\) uniquely to a symmetric bi-rigidified commutative biextension \(\mathcal P_{\mathcal S}\) on \(G_{\mathcal S}\times_{\mathcal S}G_{\mathcal S}\). At any geometric point \(\bar s:\operatorname{Spec}\bar k\to\mathcal S\) with Chevalley extension \(0\to T\to G_{\mathcal S,\bar s}\xrightarrow{q_{\bar s}}B\to0\), choose a complete discrete valuation ring \(V_0\) with residue field \(\bar k\) and characteristic-zero fraction field \(K_0\) mapping into \(\mathcal A_{g,n,F'}\). By Grothendieck's duality and biextension extension over \(V_0\) \cite[\S7]{grothendieck-1972}, \cite[Ch.~III, \S7]{faltings-chai-1990}, the generic principal polarisation extends to an isomorphism \(G_{V_0}\xrightarrow{\sim}G_{V_0}^t\) inducing a principal polarisation \(\lambda_B:B\xrightarrow{\sim}B^\vee\) and extending the generic Poincar\'e biextension to a bi-rigidified biextension on \(G_{V_0}\times_{V_0}G_{V_0}\) whose closed-fibre restriction is \((q_{\bar s}\times q_{\bar s})^*(\mathrm{id}_B\times\lambda_B)^*\mathcal P\); uniqueness of rigidified extension on the regular scheme \(G_{V_0}\times_{V_0}G_{V_0}\) with irreducible closed fibre identifies \(\mathcal P_{\mathcal S}|_{G_{\bar s}\times G_{\bar s}}\) with \((q_{\bar s}\times q_{\bar s})^*(\mathrm{id}_B\times\lambda_B)^*\mathcal P\).

On the moduli open \(U\subset Z\) over \(F'\), the fixed Weil pairing and \(\zeta_n\in F'\) turn the level-\(n\) half-level data into a full symplectic trivialisation of \(A_U[n]\). Forgetting the \(\mathcal O_E\)-action and prime-to-\(n\) level gives a quasi-finite \(F'\)-morphism \(f_U:U\to\mathcal A_{g,n,F'}\), which is proper (and hence finite) by the valuative criterion: over any characteristic-zero discrete valuation ring \(R_0\) over \(F'\) on which \(A_{K_0}\) extends to a principally polarised abelian scheme \(A_{R_0}\), the \(\mathcal O_E\)-action extends uniquely by Lemma~\ref{B8.25a} \cite[Prop.~3.3.1.5]{lan-2013} and the \(K_0\)-rational half-level sections of the finite \'etale subgroup scheme of \(A_{R_0}[M']\) extend uniquely across \(R_0\) by properness of \(A_{R_0}\to\operatorname{Spec}R_0\), remaining disjoint on the closed fibre. Let \(\overline W\subset\mathcal S\) be the schematic closure of \(f_U(U)\), and let \(X\to\overline W\) be the normalisation of \(\overline W\) in \(\Omega\) \cite[Ths.~7.14 and 11.2]{lan-2016}. Then \(\pi_X:X\to\operatorname{Spec}O\) is integral, normal, flat and projective, and finiteness of \(f_U\) identifies the abelian locus of \(G=G_{\mathcal S}\times_{\mathcal S}X\) in \(X_{F'}\) with the dense open \(U\subset X_{F'}\). Pull \(\mathcal P_{\mathcal S}\) back to \(\mathcal P_X\) on \(G\times_X G\). Lemma~\ref{B8.25a} \cite[Prop.~3.3.1.5]{lan-2013} extends the generic \(\mathcal O_E\)-action on \(A_U\) uniquely to \(G\). For each \(a\in\mathcal O_E\), the generic Rosati isomorphism \(\operatorname{ad}_{a,U}\) is a rational section of the invertible sheaf \(\mathcal M=(a\times\mathrm{id}_G)^*\mathcal P_X^{-1}\otimes(\mathrm{id}_G\times\bar a)^*\mathcal P_X\) equal to \(1\) along \((e,e)\). At every codimension-one point \(\eta_D\) of the normal scheme \(G\times_X G\) outside \(G_U\times_U G_U\), \(\eta_D\) is the generic point of the smooth fibre over a height-one prime \(\mathfrak p\) of \(X\); since \(\mathcal O_{X,\mathfrak p}\) is a discrete valuation ring, \(G_{\mathcal O_{X,\mathfrak p}}\times_{\mathcal O_{X,\mathfrak p}}G_{\mathcal O_{X,\mathfrak p}}\) is regular, and evaluating along \((e,e)\) forces \(\operatorname{ord}_{\eta_D}(\operatorname{ad}_{a,U})=0\). Hence \(\operatorname{ad}_{a,U}\) extends to a bi-rigidified isomorphism \(\operatorname{ad}_a:(a\times\mathrm{id}_G)^*\mathcal P_X\xrightarrow{\sim}(\mathrm{id}_G\times\bar a)^*\mathcal P_X\) across \(G\times_X G\). Since the quadratic host \(\Q(\sqrt{-d})\subseteq E\) has odd discriminant \(-d\equiv1\pmod4\), the element \(w_0=(1+\sqrt{-d})/2\) lies in \(\mathcal O_E\) and satisfies \(w_0+\bar w_0=1\); define \(L_{\rm PP}=(\mathrm{id}_G,w_0)^*\mathcal P_X\) on \(G\), rigidified along \(e\). The biextension partial-group laws give \(([-1]\times[-1])^*\mathcal P_X\simeq\mathcal P_X\), making \(L_{\rm PP}\) symmetric, and the partial-group laws, symmetry and \(\operatorname{ad}_{w_0}\) identify its Mumford biextension as
\[
 \Lambda(L_{\rm PP})
 \simeq(\mathrm{id}_G\times w_0)^*\mathcal P_X\otimes(w_0\times\mathrm{id}_G)^*\mathcal P_X
 \simeq\bigl(\mathrm{id}_G\times(w_0+\bar w_0)\bigr)^*\mathcal P_X
 =\mathcal P_X,
\]
making \(L_{\rm PP}\) cubical with \(L_{\rm PP}\simeq(\mathrm{id}_G,w_0)^*\Lambda(L_{\rm PP})\). On any geometric special-fibre point \(\bar x:\operatorname{Spec}\bar k\to X\) with \(0\to T\to G_{\bar x}\xrightarrow{q_{\bar x}}B\to0\), \(M=(\mathrm{id}_B,w_{0,B})^*(\mathrm{id}_B\times\lambda_B)^*\mathcal P\) on \(B\) satisfies \(L_{\rm PP}|_{G_{\bar x}}\simeq q_{\bar x}^*M\) and \(\Lambda(M)\simeq(\mathrm{id}_B\times\lambda_B)^*\mathcal P\), inducing \(\lambda_B\).

Lemma~\ref{B8.24dc} therefore furnishes the coherent \(\mathcal O_E\)-module \(\mathcal H\) on \(X\), locally free of rank \(r\) over \(\mathcal O_E\otimes_{\Z}\mathcal O_X\), in the locally split \(\mathcal O_X\)-exact sequence \(0\to\mathcal L^\vee(\mathrm{conj})\to\mathcal H\to\mathcal L\to0\), where \(\mathcal L=\operatorname{Lie}(G/X)\). Form the invertible sheaf
\[
 \mathscr D_\tau=\det\bigl(\mathcal H\otimes_{\mathcal O_E,\tau}\mathcal O_X\bigr).
\]
Over \(X_{F'}\), splitting \(\mathcal O_E\otimes_{\Z}\mathcal O_{X_{F'}}\) identifies \(\mathscr D_\tau|_{X_{F'}}\) canonically with \(\det(\mathcal L_{X_{F'},\tau})\otimes\det(\mathcal L_{X_{F'},\bar\tau})^{-1}\), extending the line \(\mathcal L_\tau\) of Lemma~\ref{B8.24b} across \(X_{F'}\). Fix a complex embedding \(F'\hookrightarrow\C\) and choose an embedded characteristic-zero resolution \(\widetilde X_\C\to X_\C\) which is an isomorphism over \(U_\C\) and has simple normal crossings boundary \(\widetilde X_\C\setminus U_\C\) \cite{hironaka-1964}. On \(\mathcal S_\C\), the logarithmic canonical extension \(\mathscr H_{\rm can}\) of de Rham homology \cite[Ch.~VI, Th.~1.1]{faltings-chai-1990} has unipotent monodromy around the toroidal boundary (as the odd level \(n\ge5\) is neat) and sits in an exact sequence \(0\to\operatorname{Lie}(G_{\mathcal S,\C}/\mathcal S_\C)^\vee\to\mathscr H_{\rm can}\to\operatorname{Lie}(G_{\mathcal S,\C}/\mathcal S_\C)\to0\); because \(\widetilde X_\C\setminus U_\C\) is a simple normal crossings divisor and the monodromy is unipotent, the pullback of \(\mathscr H_{\rm can}\) to \(\widetilde X_\C\) coincides with Deligne's canonical extension of \(H_1^{\rm dR}(A_{U_\C}/U_\C)\), so the flat \(\mathcal O_E\)-action on \(U_\C\) extends uniquely to it compatibly with the pulled-back Hodge sequence, and the determinant of its \(\tau\)-summand is also canonically isomorphic to \(\det(\mathcal L_{\widetilde X_\C,\tau})\otimes\det(\mathcal L_{\widetilde X_\C,\bar\tau})^{-1}\). Around any boundary point of \(\widetilde X_\C\), the local canonical-extension frame is obtained from a multivalued flat basis by applying the exponential of a linear combination of the commuting nilpotent monodromy logarithms \(N_i\); since the flat \(\mathcal O_E\)-action commutes with \(N_i\), each \(N_i\) preserves the \(\tau\)-summand and has trace zero there, so \(\det\exp(\sum_i z_i N_i|_{\tau})=1\). Thus the horizontal generator \(s\) of Lemma~\ref{B8.24b} extends to a nowhere-vanishing holomorphic section on \(\widetilde X_\C\), algebraizes by projective GAGA \cite{serre-1956} (Lemma~\ref{B8.25a}), descends along \(\widetilde X_\C\to X_\C\) by normality of \(X_\C\), and by faithful flatness generates \(\mathscr D_\tau|_{X_{F'}}\) everywhere on \(X_{F'}\). Hence \(\operatorname{div}_X(s)\) is supported on the special fibre \(X_{k_0}\).
\leanref{GlobalTheory.\allowbreak{}CompanionTheory.\allowbreak{}Ar.\allowbreak{}Flat.\allowbreak{}A0.\allowbreak{}flat\_\allowbreak{}zero\_\allowbreak{}of\_\allowbreak{}resolution\_\allowbreak{}gaga}{Development.\allowbreak{}AutomorphicSpectrum.\allowbreak{}Proofs.\allowbreak{}GeometryProofs.\allowbreak{}FlatZeroFromResolutionAndGaga}{39--57}
\leanref{GlobalTheory.\allowbreak{}CompanionTheory.\allowbreak{}Ar.\allowbreak{}Flat.\allowbreak{}GenAgree.\allowbreak{}generic\_\allowbreak{}agree}{Development.\allowbreak{}AutomorphicSpectrum.\allowbreak{}Proofs.\allowbreak{}GeometryProofs.\allowbreak{}HodgeLatticeGenericAgree}{225--230}

Write \(\operatorname{div}_X(s)=\sum_i n_i X_i\) and \(\operatorname{div}_X(\varpi)=\sum_i m_i X_i\) (\(m_i\ge1\)) on the irreducible components \(X_i\) of \(X_{k_0}\), and put \(a/M=\min_i(n_i/m_i)\) with integers \(M>0\) and \(a\). By normality of \(X\), \(s^M/\varpi^a\) is a global regular section of \(\mathscr D_\tau^{\otimes M}\) whose effective Cartier divisor omits at least one component \(X_i\). On the projective \(k_0\)-variety \(X_{i,\rm red}\), the restricted effective Cartier divisor has degree zero on the normalisation of every complete curve after extending \(k_0\) to \(\bar k_0\), by Lemma~\ref{B8.24dd}. If that restriction were nonempty, choose an irreducible component \(Y\) of \((X_{i,\rm red})_{\bar k_0}\) meeting the divisor at a closed point \(p\): when \(\dim X_i=0\) any nonzero section on a point is invertible; when \(\dim X_i=1\) take \(\Gamma=Y\); and when \(\dim X_i\ge2\) intersect \(Y\) with \(\dim X_i-1\) hyperplanes through \(p\) meeting the divisor support in dimension \(0\) (by prime avoidance over the infinite field \(\bar k_0\)) and take a reduced irreducible component \(\Gamma\) through \(p\) (which has dimension \(1\) by Krull's principal ideal theorem). In each positive-dimensional case \(\Gamma\subset X_{\bar k_0}\) is a complete integral curve meeting the divisor properly at \(p\) with positive degree, a contradiction. Every omitted component is therefore disjoint from the divisor support. By Stein factorisation \cite[III$_1$, Th.~4.3.1]{grothendieck-1961}, \(\pi_{X,*}\mathcal O_X=O\) because \(F'\) is relatively algebraically closed in \(\Omega\) and \(O\) is normal, so \(X_{k_0}\) is connected; hence \(\operatorname{div}_X(s^M/\varpi^a)=0\) and \(s^M/\varpi^a\) generates \(\mathscr D_\tau^{\otimes M}\) everywhere on \(X\).

Let \(y\) be any integral frame on \(Z\) over a test trait \(\operatorname{Spec}V\) of ramification index \(e=v_y(\varpi)\) over \(w\), with fraction field \(K\) and semiabelian extension \(G_V\to\operatorname{Spec}V\). By the valuative criterion of properness, \(\operatorname{Spec}K\xrightarrow{y}U\hookrightarrow X\) extends uniquely to an \(O\)-morphism \(\tilde y:\operatorname{Spec}V\to X\), and Lemma~\ref{B8.25a} \cite[Prop.~3.3.1.5]{lan-2013} applied to \(\mathrm{id}_{A_K}\) identifies \(\tilde y^*G\xrightarrow{\sim}G_V\) compatibly with \(\mathcal O_E\) and hence \(\tilde y^*\mathcal L\xrightarrow{\sim}\operatorname{Lie}(G_V/V)\). Injectivity of \(\operatorname{Lie}(G_V/V)\hookrightarrow\operatorname{Lie}(A_K/K)\) identifies the integral \(\tau\)- and \(\bar\tau\)-eigenframes over \(V\) with bases of \(\operatorname{Lie}(G_V/V)\cap(\operatorname{Lie}(A_K/K))_\tau\) and \(\operatorname{Lie}(G_V/V)\cap(\operatorname{Lie}(A_K/K))_{\bar\tau}\). Pulling \(0\to\mathcal L^\vee(\mathrm{conj})\to\mathcal H\to\mathcal L\to0\) back along \(\tilde y\) gives an integral Hodge sequence over \(V\) in which \(s(y)\) has valuation \(ea/M\) relative to a \(V\)-generator of \(\tilde y^*\mathscr D_\tau=\det\bigl((\tilde y^*\mathcal H)(\tau)\bigr)\). Applying Lemma~\ref{B8.24c} over \(V\) yields \(v_y(s(y))=e\rho_w\) with
\[
 \rho_w=\frac aM-\sum_{\sigma\in S\setminus\{\tau\}}h_\sigma\,v_w(a_\sigma-a_\tau)\in\Q,
\]
since each \(a_\sigma\) lies in the integer ring of \(F'_w\) so \(v_y(a_\sigma-a_\tau)=e\,v_w(a_\sigma-a_\tau)\), and \(h_\sigma\in\{P,Q\}\) depends only on the generic signature. (Comparing with any other \(F'\)-rational horizontal generator \(s'=cs\) shifts \(\rho_w\) by \(+v_w(c)\) for \(c\in(F')^\times\).)
\leanref{GlobalTheory.\allowbreak{}CompanionTheory.\allowbreak{}Ar.\allowbreak{}Flat.\allowbreak{}CPoint.\allowbreak{}value\_\allowbreak{}at\_\allowbreak{}point}{Development.\allowbreak{}AutomorphicSpectrum.\allowbreak{}Proofs.\allowbreak{}GeometryProofs.\allowbreak{}EigenLinePAdicFrame}{72--339}

Choose a closed point \(y_1\) of \(U\) over a finite extension \(K_1/F'\). Spreading out the abelian variety \(A_{y_1}\), its \(\mathcal O_E\)-action, bases of its two Lie eigenspaces (with their inverse transition matrices) and the nonzero value \(s(y_1)\in K_1^\times\) after inverting \(\operatorname{disc}E\) and finitely many denominators produces an integral frame over \(\mathcal O_{K_1,w_1}\) with \(v_{w_1}(s(y_1))=0\) at all but finitely many places \(w_1\) of \(K_1\), giving \(\rho_w=0\) outside finitely many places \(w\) of \(F'\).

Finally, the row-metric defect is \(Q\) times the difference of the orders of the extended different \(\tau(\mathfrak D_{E/\Q})\) and of \(d\). Since \(E/\Q\) is Galois, taking the ideal norm identifies \(v_w(\tau(\mathfrak D_{E/\Q}))/v_w(\ell)\) at \(w\mid\ell\) with \(v_\ell(|\operatorname{disc}E|)/[E:\Q]\), while the order of \(d\) is \(v_w(\ell){\bf1}_{\ell=d}\). Subtracting this defect gives the stated coefficient, whose sum against \(\log Nw/[F':\Q]\) contributes \(Q(\log d-\log\operatorname{rd}(E))\).
\leanref{GlobalTheory.\allowbreak{}FlatLine.\allowbreak{}flDefect\_\allowbreak{}vw\_\allowbreak{}eq}{Development.\allowbreak{}AbelianVarietyTheory.\allowbreak{}Definitions.\allowbreak{}Foundations.\allowbreak{}ExistsFlatSecRankZero}{188--210}
\leanref{GlobalTheory.\allowbreak{}FlatLine.\allowbreak{}sum\_\allowbreak{}flDefect\_\allowbreak{}vw}{Development.\allowbreak{}AbelianVarietyTheory.\allowbreak{}Definitions.\allowbreak{}Foundations.\allowbreak{}ExistsFlatSecRankZero}{219--224}
\end{proof}

\subsection*{Proof of \ref{B8.25}}\label{prf:B8.25}\phantomsection
\begin{proof}
Let \(\mathcal H_X\) consist of the prescribed-fan model, Lie seminegativity, the separation conclusion under \eqref{eq:p8c3-separation}, the cuspidal Green parametrisation of \ref{B7.31} and the Fubini--Study formula for Chow height. Let \(\mathcal H_H\) consist of the Siegel model, the prescribed-fan boundary charts, their compatibility with the Siegel charts and the arithmetic degree formula. The section comparison uses the vanishing, extension, algebraisation and nef intersection statements of \(\mathcal H_P\), together with its Chow form and the hyperplane-section formula.

For a component \(Z\), write \(X\) for the projective image model and \(A\) for its Hermitian line. Choose an integral degree-one section \(w\) of nonzero generic value \(g_w\) on \(Z\), and let \(\psi\) be the embedding of the component function field in that of \(X\). A small section \(s\) of degree \(8a\), with generic value \(g_s\), gives the rational representative \(\psi(g_s)/\psi(g_w)^{8a}\) relative to the image of \(w^{8a}\). The regulator comparison gives regularity in each finite-place chart. The pointwise comparison for \(A\) puts the section in the unit ball of the \(8a\)-th power of \(A\) with metric shifted by the small-section threshold. Equal images give equal rational representatives. Cancel the nonzero denominator and use injectivity of \(\psi\) and of generic evaluation on the component space. This is the counting injection used with the Chow bound on self-intersection in \ref{B8.28}.
\leanref{PositivityTheory.\allowbreak{}Xj7N.\allowbreak{}Inj.\allowbreak{}small\_\allowbreak{}sections\_\allowbreak{}injective\_\allowbreak{}of\_\allowbreak{}self\_\allowbreak{}intersection\_\allowbreak{}bound}{Development.\allowbreak{}AutomorphicArithmetic.\allowbreak{}ThetaCorrespondence.\allowbreak{}ThetaPairLifts.\allowbreak{}HodgeBundleWeights.\allowbreak{}SmallSectionsInjectiveOfSelfIntersectionBound}{20--74}

Fix an eligible family and take its common threshold for the weight and cover conditions. Lemmas~\ref{B8.25b} and \ref{B8.25c} give the quotient weights and good covers. On every good cover at a nonterminal layer, the weight-gap inequalities give weakly decreasing rows and a wall gap at least one for \(u_M\) and for the rows at shifts \(-c_0,0,c_0\). For all four weights, the Petersson norms admit Gram frames and the integral cusp sections form a finitely generated lattice spanning the cusp space over the coefficient field. The same conclusions hold on each component for \(u_M\). The components form a finite nonempty set. The toroidal and nef comparisons give a nonzero cusp section on a component in the anchored weight. The flat twist carries it to a nonzero section in shift \(-c_0\), so the section lattice there has positive rank. This proves deep goodness as defined in \ref{def:p8c3-section-laws}.
\leanref{GlobalTheory.\allowbreak{}CovTop.\allowbreak{}DeepGoodAll.\allowbreak{}deepGood\_\allowbreak{}alt\_\allowbreak{}alt}{Development.\allowbreak{}AutomorphicArithmetic.\allowbreak{}ThetaCorrespondence.\allowbreak{}ThetaPairLifts.\allowbreak{}ThirdPlaceSeries.\allowbreak{}DeepGoodReductionStatement}{23--44}
\leanref{GlobalTheory.\allowbreak{}CovTop.\allowbreak{}LowPos.\allowbreak{}lowered\_\allowbreak{}rank\_\allowbreak{}pos\_\allowbreak{}alt\_\allowbreak{}alt}{Development.\allowbreak{}AutomorphicArithmetic.\allowbreak{}ThetaCorrespondence.\allowbreak{}ThetaPairLifts.\allowbreak{}ThirdPlaceSeries.\allowbreak{}RankPosNegativeWeightBundle}{26--48}

To prove the flat-degree law of \ref{def:p8c3-section-laws}, fix a component and a CM embedding. Use the rational generator on this component constructed in Lemma~\ref{B8.24b}, together with its reciprocal. Its norm at each infinite place is constant in the frame. Lemma~\ref{B8.24d} computes the common coefficient-place order \(\rho_w\) in integral Lie frames, proves finite support and gives the corrected order after subtraction of the different defect. The CM frame formula computes the degree from the archimedean norms and the \(\rho_w\) by the product formula. The different correction adds \(Q(\log d-\log\operatorname{rd}(E))\) to that degree. The corrected frame has degree \(-c_{1,\tau}\) of \eqref{eq:p8c3-flat-c1}. The frame formula gives this degree for every flat frame, and the correction gives a flat frame on each component at each CM embedding.
\leanref{GlobalTheory.\allowbreak{}T.\allowbreak{}flatDeg\_\allowbreak{}of\_\allowbreak{}scholzeWeinstein\_\allowbreak{}of\_\allowbreak{}prints}{Development.\allowbreak{}AutomorphicArithmetic.\allowbreak{}ThetaCorrespondence.\allowbreak{}ThetaPairLifts.\allowbreak{}ThirdPlaceSeries.\allowbreak{}FlatDegreeStatement}{22--23}
\leanref{GlobalTheory.\allowbreak{}FlatCore.\allowbreak{}flatDegree\_\allowbreak{}of\_\allowbreak{}steps}{Development.\allowbreak{}AbelianVarietyTheory.\allowbreak{}Proofs.\allowbreak{}DivisorProofs.\allowbreak{}FlatDegreeOfSteps}{11--27}

Apply \ref{B8.25e} to the degrees and ranks, and \ref{B8.26} and \ref{B8.28} to the heights. Proposition~\ref{B8.29} descends the cover slope with the flat comparison of \ref{B8.24a} and treats the terminal layer with the definite defect. The Gram frame and norm comparison of \ref{B8.25d} give \eqref{eq:p8c3-gram} and \eqref{eq:p8c3-core} in the realised lattices. Combine \eqref{eq:p8c3-sizes}, the bad-rank bound \ref{B8.30} and the projector inequality \ref{B8.31}. Lemma~\ref{B8.32} gives the layer estimate and \ref{B8.33} gives a common family threshold and constant.
\end{proof}

\subsection*{Proof of \ref{B8.25h}}\label{prf:B8.25h}\phantomsection
\begin{proof}
The degree-zero part of the localisation of the homogeneous coordinate ring at the coordinates selected by \(c\) is \(R_c\). The hypothesis on quotients gives \(\operatorname{Frac}(R_c)=\Omega\), by dividing the numerator and denominator by the same chart monomial. Thus \(X\) is integral. It is projective by the Segre embedding and flat over \(\Z\) since its affine rings are torsion-free. Each \(R_c\) is of finite type over \(\Z\), hence excellent. Its integral closure is finite, and integral closure commutes with localisation. Thus the inverse images of the indicated charts under \(\nu\) are \(\operatorname{Spec}\widetilde R_c\). In the rational trivialisation given by \(y\), the local frame of \(L^r\) is \(m_c^r\). The two formulas for global sections now follow by gluing inside the common function field. They apply also when some \(a_k\) are zero or negative.

In a rational presentation \(s=g(y)/h_1(y)\) of degree difference \(ra\), the unit-ball condition on the locus \(h_1\ne0\) is
\[
                 |g(z)|\le |h_1(z)|h(z)^r .
\]
Such presentations exist by clearing chart denominators. On either model the condition extends by continuity from the locus of definition of the quotient, which is dense in every complex component. Conversely, a bound for the norm of the section gives the displayed inequality for every presentation by multiplication by the denominator. A proper surjective pullback preserves a sup norm.

Here is a conductor independent of \(r\). On each chart clear the denominators of finitely many \(R_c\)-module generators of \(\widetilde R_c\). This gives \(0\ne b_c\in R_c\) with \(b_c\widetilde R_c\subset R_c\). Homogenise \(b_c\) to an integral form \(F_c\), using the selected coordinate in each block. Let \(F=\prod_c F_c\), and let \(H\) be the restriction to \(X\) of the line of its multidegree. Its coefficient in chart \(c\) is \(b_c\) times an element of \(R_c\). Every factor is nonzero at the generic point. Hence multiplication by \(F\) gives an injection
\[
 H^0(X^\nu,\nu^*L^r)\longrightarrow H^0(X,L^r\otimes H)
\]
for every \(r\). Choose a smooth metric on \(H\) and rescale it by a constant so that \(\|F\|_{\sup}\le1\). We obtain
\[
 \widehat h^0(X,r\overline L)\le
 \widehat h^0(X^\nu,r\nu^*\overline L)\le
 \widehat h^0(X,r\overline L+\overline H).
\]
Apply Moriwaki's convergence theorem to the subgroup of the Hermitian Picard group generated by \(\overline L,\overline H\). The two outer sequences of classes, divided by \(r\), tend to \(\overline L\) in its real span. The base is a projective arithmetic variety and the metrics are continuous. The two outer counts therefore have the same leading term, including when the volume is zero \cite[Th.~5.1]{moriwaki-2009b}. This proves the asserted little-oh comparison. The conductor here is a section of a fixed line bundle and works also along a horizontal nonnormal locus.

Choose a projective resolution of \(X^\nu_{\Q}\) by Hironaka's theorem in characteristic zero \cite{hironaka-1964}. Embed the resolution projectively over \(X^\nu_{\Q}\), take its integral closure as a closed subscheme in the sense of schematic closure in \(X^\nu\times\mathbb P^n_{\Z}\), and normalise that closure. Denote the resulting projective birational morphism to \(X^\nu\) by \(f\). Its generic fibre is the chosen resolution. Since \(X^\nu\) is normal, \(f_*\mathcal O=\mathcal O_{X^\nu}\), by Stein factorisation and finite birationality. The projection formula identifies the section spaces after pullback, with their sup norms.

On a complex chart of the resolution lift the map to the cone by setting \(z_{k c_k}=1\); denote this holomorphic lift by \(u\). For every \(F_0\in P_0\), \(F_0\circ u=0\), so \(du(v)\) is tangent to \(P_0\). The Levi form of the pulled-back local weight at \(v\) is
\(\mathcal L_{\log h}(u;du(v))\ge0\). Changing the lift adds a sum of logarithms of absolute values of holomorphic units to the weight. This leaves its Levi form unchanged.
\end{proof}

\subsection*{Proof of \ref{B8.25i}}\label{prf:B8.25i}\phantomsection
\begin{proof}
Shrink to an affine neighbourhood of \(x\) and choose \(0\ne t\in A\) whose principal open in \(\operatorname{Spec}A\) maps into \(U\). The residue field is finite. Excellence gives that \(R\) is a normal local domain.

We first record how to test a nonzero element \(w\) of a Noetherian local domain \(B\) of characteristic zero and finite residue field of characteristic \(p\). There is a one-dimensional local domain quotient in which \(pw\) survives. Indeed, while the dimension exceeds one, choose an element of the maximal ideal outside all height-one primes dividing \(pw\) and outside the zero prime. There are finitely many primes to avoid by the principal ideal theorem. Quotient by a prime minimal over the chosen element. It has height one, avoids \(pw\), and the dimension drops. In the resulting one-dimensional domain \(C\), the radical of \(pC\) is the maximal ideal. Thus \(C/pC\) is finite and the completion of \(C\) is module-finite over \(\Z_p\), by lifting generators modulo \(p\) and successive approximation. The completion is faithfully flat over \(C\). Choose a minimal prime of it avoiding \(pw\), which is possible since \(pw\) is a nonzero divisor. The quotient is a finite domain over \(\Z_p\). Its fraction field embeds in \(\overline{\Q}_p\), and its maximal ideal maps into the maximal ideal of the valuation ring. This gives a local test of \(B\) on which \(w\) survives. Taking \(B=R,w=t\) gives the first assertion.

Suppose \(s\notin aR\). In \(R[a/s]\), reduction modulo the maximal ideal of \(R\) and evaluation of \(a/s\) at zero define a homomorphism to the residue field. Indeed, a polynomial relation for \(a/s\) with unit constant term would give a monic equation for \(s/a\) by reversal and division by that unit. Normality would then give \(s/a\in R\). Localise \(R[a/s]\) at the kernel of this homomorphism. The resulting local domain \(B\) has the same residue field, and \(a=s u\) with \(u\) in its maximal ideal. Apply the preceding construction to \(w=st\). The resulting test has \(|v(a)|<|v(s)|\) and its generic point lies in \(U\), a contradiction.

Finally write the rational function as \(s/a\) with \(s,a\in A\), \(a\ne0\). Apply the divisibility assertion with the test open further restricted by \(a\ne0\). This gives \(s\in aR\). Faithful flatness gives \(aR\cap A=aA\).
\end{proof}

\subsection*{Proof of \ref{B8.25j}}\label{prf:B8.25j}\phantomsection
\begin{proof}
Write \(\Gamma=(\Z/rN\Z)^{2g}\) for the full group of level labels, with \(g\) the dimension of the principally polarised abelian variety. For a free isotropic subgroup \(T\subset\Gamma\) and a character \(\chi:T\to\mu_N\), let \(P_\chi=|T|^{-1}\sum_{t\in T}\chi(t)^{-1}\rho(t)\), where \(\rho\) is the theta action in the chosen basis of sections. On \(T\) the chosen lifts multiply as the group does. Write \(v_j(t)\) for the evaluation of basis section \(j\) at label \(t\) in its torsion frame. The entries of the Fourier tuple are
\[
 F_{i,\eta}=\sum_{t\in T}\eta(t)\sum_j(P_\chi)_{ji}v_j(t),
       \qquad \eta:T\longrightarrow\mu_{(rN)^2}.
\]
Here \(\eta\) ranges over all functions. The tuple is finite. To obtain the model in the statement, take the projective image given by all products of one entry from each nonzero tuple, including the theta-null tuple, and normalise in the function field of the component. The theta map is generically finite on this component. On a chart where a product entry generates, vary one factor at a time. Every ratio to the selected entry of that factor is regular. This construction uses only the level evaluations. Choose a dense open \(U\) where their nonzero numerators and denominators are units.

Fix a local test \(v\) of Lemma~\ref{B8.25i} with generic point in \(U\). The symmetric degeneration datum gives
\[
                        |v(s_\alpha)|=\exp(-q_v(\alpha)),
\]
where \(q_v\) is the positive definite quadratic form obtained by taking minus the logarithm of the absolute value of the cubical period. The polarisation identity follows from the biextension law and symmetry under inversion, as in Lemma~\ref{B8.25k} below. Thus \(q_v\) attains its minimum on each residue class. Write \(b_v\) for a minimiser of \(c\) on this test.

We describe the subgroup and character used for \(c\). The full level gives the quotient map from the torsion labels to \(X_*/rNX_*\). This map on the completed chart is read by pairing with the split toric \(rN\)-torsion. Indeed, evaluation against each character of \(X_*/rNX_*\) is the Weil pairing with the corresponding toric label. These pairings take values in a constant group of roots of unity since the level is invertible and split. They commute with base change, and the characters of the finite free module separate points. Thus the quotient map is the same at every test under the completed-chart identification of period lattices. Take \(T\) to be the perpendicular of its kernel for the level pairing. In symplectic coordinates the toric labels form a free isotropic summand, dual to the period quotient. Choose a full level label above a lift of \(c\) modulo \(rN\). Pair it with \(T\) and reduce the exponent modulo \(N\) to get \(\chi\). Adding an element of the kernel or changing the lift by \(N\) leaves \(\chi\) unchanged. This choice precedes the choice of \(b_v\).

Let \(\delta_v\) be the greatest norm of a theta null at \(v\). We give the finite Fourier calculation in the normalisation of the period sections:
\[
          \max_{i,\eta}|F_{i,\eta}(v)|
                   =\delta_v\exp(-q_v(b_v)/N).                 \tag{1}
\]
In an integral frame write the principal theta section as a scalar times a primitive section. The norm of this scalar is \(\delta_v\). To see the equality, the terms of the primitive theta expansion at a level point have norm at most one: the exponents after translation are the values of the nonnegative quadratic form on the translated lattice, with the torsion-frame factor included. In the integral torus direction the reduction is the theta section on the abelian part. It has a nonzero value at some prime-to-\(p\) level point. Indeed, for an invertible level divisor \(n\ge2\), the translates of the \(n\)-th power of a nonzero principal theta section by \(n\)-torsion span its level-\(n\) section space. The character projectors for one maximal isotropic subgroup project to lines and the complementary subgroup permutes these lines transitively. The line bundle of level \(n\) is generated by its sections, so the translates cannot all vanish at the origin. Lift a nonvanishing point by the prime-to-\(p\) torsion of the semiabelian scheme. The integral torsion frame then gives a unit value.

Divide the level basis by the same scalar. At the origin the integral Fourier expansions of level \(N\) are those whose coefficient at \(\beta\), in the cubical frame, has norm at most \(\exp(-q_v(\beta)/N)\). Indeed, choosing one representative in each class modulo \(N\) identifies an expansion with its coefficients on the abelian part at those representatives. Translation by \(Nz\) multiplies the coefficient by the period factor of log valuation
\[
                     \frac{q_v(\beta+Nz)-q_v(\beta)}{N}.
\]
This is the Fourier translation law \(D_2\psi=\tau\) for the cubical trivialisation \(\psi\) and its biextension trivialisation \(\tau\) \cite[Def.~4.2.1.13 and \S4.3]{lan-2013}. One may extend the valued field to realise the finitely many fractional weights of the representatives.

For clarity, the integral lattice in this calculation is generated by the normalised theta translates. Multiplication by \(m\) pulls the primitive principal theta section back to a primitive level-\(N\) vector \(s_0\): its exponent \(\beta\) becomes \(m\beta\), and \(q_v(m\beta)/N=q_v(\beta)\). Translation in a torsion frame preserves the weighted coefficient lattice by the same translation law. This lattice has an integral basis, obtained class by class from the section lattices of the ample coefficient lines on the abelian part after clearing the indicated weights. Take an integral coordinate functional \(f\) with \(f(s_0)\) a unit. If \(U_x\) are the theta lifts of \(x\in A[N]\) on the generic section space of dimension \(D=N^g\), the perfect commutator pairing gives
\(\operatorname{tr}(U_y^{-1}U_x)=D\,\mathbf 1_{x=y}\). For \(x\ne y\), conjugation by a lift having nontrivial pairing with \(x-y\) proves the vanishing. There are \(D^2\) operators, so they form a basis of the endomorphism algebra. Expanding the rank-one operator \(w\mapsto f(w)s\) in that basis and evaluating at \(s_0\) gives
\[
                   s=\sum_{x\in A[N]}
                    \frac{f(U_x^{-1}s)}{D f(s_0)}\,U_xs_0 .
\]
All coefficients are integral for integral \(s\), since the residue characteristic is prime to \(N\). The level basis spans these translates by its construction. This also identifies the scale of its coefficients with the scale of the theta nulls.

The projector \(P_\chi\) retains precisely the terms with exponent in \(c\). Evaluating at toric torsion multiplies their coefficients by units. The ultrametric inequality therefore gives the upper bound in (1). For attainment, prescribe at \(b_v\) a coefficient with norm \(\exp(-q_v(b_v)/N)\) at the origin and extend it by the translation law, setting the other classes to zero. To make this prescription, first choose a section of the coefficient line on the abelian part that is a unit at the given special-fibre point. Geometrically that line is a degree-zero twist of the \(N\)-th power of an ample principal line. It is a translate of that power and is globally generated since \(N\ge4\). Sections lift over the valuation ring by base change for the ample line on the abelian scheme, whose positive-degree cohomology vanishes on the fibres. Scale the section by an element of the prescribed valuation. The translation formula gives the required integral expansion with equality at \(b_v\).

In its values on \(T\), average against the inverse of the toric character of \(b_v\). Orthogonality removes exponents outside \(b_v+rNX_*\). For \(z\ne0\) one has
\[
 q_v(b_v+rNz)-q_v(b_v)
 =r\{q_v(b_v+Nz)-q_v(b_v)\}+r(r-1)N^2q_v(z)>0 .
\]
Thus the term at \(b_v\) survives with strictly greatest norm. The series converges because the form is positive definite, and \(|T|\) is a unit. At least one torsion evaluation of the prescribed expansion attains the upper bound. Since the expansion is an integral linear combination of the normalised basis sections and belongs to the \(\chi\)-space, some projected basis evaluation attains it. Finally, if the sum of those evaluations over \(t\), for that basis index, has smaller norm, change the phase at one attaining \(t\) from \(1\) to \(\zeta\). The difference has the required norm since \(1-\zeta\) is a unit. This proves (1), also in the presence of ties among shortest representatives.

Choose at \(x\) a generating theta null \(\theta_0\). All ratios to it are regular at \(x\), hence \(|\theta_0(v)|=\delta_v\) on every test. Equation (1) shows that the Fourier tuple for \(c\) is nonzero. Choose its generating entry \(F_0\) at \(x\) and put \(e_c=F_0/\theta_0\). All ratios \(F_{i,\eta}/F_0\) belong to \(A\). Consequently (1) gives, for every test,
\[
                |e_c(v)|^N=\max_{\alpha\bmod N=c}|v(s_\alpha)|\le1 .
\]
By Lemma~\ref{B8.25i}, \(a=e_c^N\) belongs to \(A\), and \(a\mid s_\alpha\) in \(R\) for every representative. Choose one test and an attaining representative \(\mu_c\) for it. Write \(s_{\mu_c}=a u\). On this test \(a\) is nonzero and \(|v(u)|=1\). The test homomorphism is local, so \(u\) is a unit of \(R\). This proves both assertions in the completed ring.
\end{proof}

\subsection*{Proof of \ref{B8.25k}}\label{prf:B8.25k}\phantomsection
\begin{proof}
Put \(A=\mathcal O_{X,x}\). The specialisation gives a map \(A\to V\). To compare it with \(R=\widehat A\), form \(V\otimes_A R\). This algebra is faithfully flat over \(V\), and its special fibre is reduced because the formal fibres of the excellent ring \(A\) are geometrically regular. Localise at a minimal prime of its special fibre. The quotient by \(\varpi\) is a field and \(\varpi\) is a nonzero divisor. Completing for \((\varpi)\) gives a complete discrete valuation ring \(V'\): each successive quotient \((\varpi^j)/(\varpi^{j+1})\) is one-dimensional over the residue field, so every nonzero element of the completion is a power of \(\varpi\) times a unit. We have compatible maps \(R\to V'\) and \(V\to V'\) with ramification index one.

Use split tori to read the period data, passing to an étale splitting extension as necessary. In a two-step degeneration the torus at the more general centre extends as a subtorus of the torus at the closed centre. Restriction of characters gives the indicated integral surjection. In the description by formal thickenings, a proper image on characters would leave a nonzero finite quotient of its cokernel. The dual finite diagonalizable group would map trivially on every thickening under the extended torus map. Since that group is finite and the semiabelian scheme is separated, equality on all thickenings implies equality of the two maps. On the fibre where the torus is embedded, the induced map is nontrivial, a contradiction. Endomorphisms commute with restriction. The polarisation changes their action on the dual torus by conjugation, giving the asserted \(\mathcal O_E\)-linearity. This is the character restriction in the two-step construction of \cite[Props.~3.3.1.5, 3.3.1.7 and \S4.5.6]{lan-2013}.

Write \(Q(y)\) for the order of the cubical period section on the trait in an integral frame. For a split polarised degeneration the Poincaré biextension trivialisation is the second difference of that section \cite[Def.~4.2.1.13(6)--(7), Th.~4.2.1.14]{lan-2013}. Hence its order is
\[
 b_V(y,z)=Q(y+z)-Q(y)-Q(z).
\]
Rigidification gives \(Q(0)=0\), and the symmetric integral line gives \(Q(-y)=Q(y)\). Substitution of \(z=-y\) gives \(b_V(y,y)=2Q(y)\). This accounts for the square in the order formula.

For the comparison under restriction, one can read the order of each biextension parameter modulo every \(\ell^j\), with \(\ell\) a prime invertible at both centres. After strict henselisation a unit has an \(\ell^j\)-th root. Tame inertia acting on an \(\ell^j\)-th root of \(u\varpi^a\) thus reads \(a\bmod\ell^j\) from its action on a root of \(\varpi\). The two uniformisations give the same polarised generic torsion, with toric characters related by restriction. Their Weil pairings give the same inertia character for the corresponding biextension parameters. The two integral orders agree modulo every \(\ell^j\), hence agree. Applying the preceding diagonal identity on both sides gives
\[
          2\operatorname{ord}_{V'}(s_\alpha)
                        =b_V(\pi_V\alpha,\pi_V\alpha).
\]
The local divisibility from Lemma~\ref{B8.25j} persists under \(R\to V'\). The associate of \(s_{\mu_c}\) is the image of the actual stalk element \(e_c^N\). Its order is therefore the order of \(e_c^N\) in \(V\), since \(\varpi\) remains a uniformiser. This proves the equality and inequality.

If the centre is \(x\), the map of the completed stalk to the extended trait ring is local. For \(\alpha\ne0\), positivity of the degeneration puts \(s_\alpha\) in the maximal ideal of \(R\). If \(\pi_V\alpha=0\), the diagonal identity makes its image a unit. Thus the surjection has zero kernel.
\end{proof}

\subsection*{Proof of \ref{B8.25m}}\label{prf:B8.25m}\phantomsection
\begin{proof}
First fix a cone with integral generators \(\gamma_1,\ldots,\gamma_m\) whose sum \(G\) is positive definite. In one residue class \(c\), associate to \(v\in c\) the tuple
\[
       \bigl(\gamma_1(v,v),\ldots,\gamma_m(v,v)\bigr)\in\N^m .
\]
There are finitely many coordinatewise minimal tuples. Indeed, every infinite sequence in \(\N\) has an infinite nondecreasing subsequence. Successive extraction in the \(m\) coordinates rules out an infinite antichain. Every tuple dominates a minimal one, by descent in the sum of its coordinates. In a strictly positive linear combination of the \(\gamma_i\), a minimiser must realise a minimal tuple: a strict decrease in any coordinate with no increase in the others lowers the value of the combination. The possible values of \(G(v,v)\) on these tuples are bounded. Since \(G\) is positive definite, its bounded sublevel sets on \(X\) are finite. Taking the product over the finite set \(X/NX\) gives finitely many complete minimum systems on this cone.

Every positive definite form is in the relative interior of a cone, and the sum of that cone's generators is positive definite. By finiteness in the definition of the polarised fan, a lattice automorphism carries its generators into a fixed finite set. There are finitely many resulting cones, and on each there are finitely many transported systems by the first paragraph. For a fixed transported system \((v_c)_c\), an inverse change of coordinates must carry every \(v_c\) into the finite set \(\{\mu_c\}_c\). The \(v_c\) span \(X\otimes\Q\): in lattice coordinates the representatives of the standard basis classes form a matrix with determinant \(1\bmod N\). Hence the images of the \(v_c\) determine the change of coordinates. This gives finitely many cones before transport.

Let \(B\) be any of the semidefinite forms in the statement. The minimum inequalities are linear in the form, so \(b_0+jB\) satisfies them for every positive integer \(j\). These forms are positive definite. One cone in the finite list contains them for unbounded \(j\). Since that cone is closed and is invariant under positive scaling, it contains \(\lim_{j\to\infty}(b_0+jB)/j=B\) along this subsequence.
\end{proof}

\subsection*{Proof of \ref{B8.25n}}\label{prf:B8.25n}\phantomsection
\begin{proof}
The inclusion \(J\mathcal O_Y\subseteq\mathcal O_Y(-D)\) follows from the definition. Fix a closed \(x\). Properness makes \(z\) closed. The residue field of \(A\) is finite, so \(A/\mathfrak m^j\) is finite for every \(j\). Diagonal extraction gives a subsequence \(F_{n_i}\), with \(n_i\to\infty\), converging in \(\widehat A\) to some \(F\).

For a closed \(x'\) over \(z\), the local homomorphism \(A\to\mathcal O_{Y,x'}\) induces \(\widehat A\to\widehat{\mathcal O}_{Y,x'}\). On any test at \(x'\) one has \(\rho_y<1\), and an element of \(\mathfrak m^j\widehat A\) has absolute value at most \(\rho_y^j\). Consequently
\[
 |F(y)|\le |s_{x'}(y)|,\qquad |F(y_0)|=|s_x(y_0)|>0 .
\]
Lemma~\ref{B8.25i}, applied at \(x'\) with the given dense open and \(s_{x'}\ne0\), shows that \(s_{x'}\) divides \(F\) in \(\widehat{\mathcal O}_{Y,x'}\). At \(x\) the quotient is a unit: its value on the local test \(y_0\) has absolute value one.

It remains to put the formal germ \(F\) in \(J_z\widehat A\). Base change to \(\operatorname{Spec}\widehat A\), and denote the pullback of \(D\) by \(D'\). At the points of the closed fibre the completed local rings of the base-changed \(Y\) are the original completed local rings. Indeed, modulo the \(j\)-th power of the point's maximal ideal the base action factors through \(A/\mathfrak m^j=\widehat A/\mathfrak m^j\widehat A\). The divisibilities just proved and faithful flatness of local completion show that the section \(F|_{D'}\) vanishes at the stalk of every closed point of the closed fibre. Its support is closed and proper over \(\widehat A\). A nonempty support would meet the closed fibre in a nonempty closed set and hence contain a closed point there. Thus \(F|_{D'}=0\).

Flat base change identifies the kernel of \(\widehat A\to H^0(D',\mathcal O_{D'})\) with \(J_z\widehat A\). For degree zero this follows by computing sections on a finite affine cover and its pairwise intersections and using flatness to preserve the equalizer. It follows that \(J\widehat{\mathcal O}_{Y,x}\) contains \(F\), which generates \(s_x\widehat{\mathcal O}_{Y,x}\). Faithful flatness gives equality of the two ideals at \(x\). The quotient of the two coherent ideals has zero stalk at every closed point of \(Y\), so it is zero.
\end{proof}

\subsection*{Proof of \ref{B8.25o}}\label{prf:B8.25o}\phantomsection
\begin{proof}
Write \(X=\operatorname{Spec}A\), \(X_1=\operatorname{Spec}B\) locally. For each finite tuple \(b_1,\ldots,b_m\) in \(I(B)\), take the coefficients of
\[
       \prod_{\sigma:L\hookrightarrow\overline K}
                  \left(\sum_{i=1}^m\sigma(b_i)T_i\right),
\]
where the embeddings fix \(K\). The coefficients lie in \(K\) by Galois invariance and are integral over \(A\), so they lie in \(A\). They generate \(J\). This prescription commutes with localisation: a tuple with localised entries can be multiplied by a common denominator \(s\in A\), and its polynomial is then multiplied by \(s^h\). The ideals therefore glue.

In a normal closure of \(L/K\) form the \(A\)-algebra \(B'\) generated by the images of \(B\) under all \(h\) embeddings. It is finite and integral over \(A\). Given \(A\to V\), choose a prime of \(B'\) over the kernel of \(A\to\operatorname{Frac}V\). The residue field at that prime is finite over the residue field at its contraction and embeds over that field in an algebraic closure of \(\operatorname{Frac}V\). Its compositum with \(\operatorname{Frac}V\) is finite over \(\operatorname{Frac}V\). The valuation ring \(V'\) in this extension is a complete discrete valuation ring. The images of \(B'\) lie in \(V'\) by integrality. Composition with the \(h\) maps \(B\to B'\) gives the required lifts.

Over a discrete valuation ring the ideal of coefficients of a product of polynomials is the product of their coefficient ideals. For nonzero factors, divide each by the common uniformiser power of its coefficients; the resulting polynomials have nonzero reductions, and their product still has nonzero reduction. A zero factor gives the same identity. Applying this to the displayed product gives one inclusion in the assertion for every tuple. Taking a tuple of generators of \(I(B)\) gives the reverse inclusion for all \(h\) lifts at once.
\end{proof}

\subsection*{Proof of \ref{B8.25a}}\label{prf:B8.25a}\phantomsection
\begin{proof}
For item~1 of Definition~\ref{B8.25l}, take the prescribed-fan model below and retain its smoothness, its simple normal crossings boundary and its logarithmic Kodaira--Spencer isomorphism on the designated interior. The volume assertions on an integral arithmetic variety give the lower bound of \ref{B8.27} for a semipositive fibrewise nef Hermitian line and equality with self-intersection for an arithmetically nef line. For such a nef line \(L\) of self-intersection \(a\), a fixed Hermitian twist \(N\) and generic dimension \(D\), the unit ball in \(H^0(L^j\otimes N)\) is eventually finite and its logarithmic count is
\begin{equation}\label{eq:p8c3-nef-count}
 \frac{a j^{D+1}}{(D+1)!}+o(j^{D+1}).
\end{equation}
The intersection laws cover change of model and metric, constant shifts, multilinearity, projection and isometry. Cycle heights obey projection. For a good metric the Chern-form integral computes the intersection of the extended classes. On a full flag bundle, insertion of the fibre point class and pushforward give the base intersection.

The ray-class assertion gives an abelian ray field for each nonzero modulus, unramified off the modulus. Its Frobenius elements generate the Galois group and the common kernel of its characters on coprime ideals is the ray-principal subgroup. Lemma~\ref{B8.25c} uses this for the covering level. Let \(K/\Q_p\) be finite and let \(A_K\) be the underlying abelian variety of the polarised datum. The base \(\operatorname{Spec}\mathcal O_K\) is regular, connected, Noetherian and one-dimensional. Grothendieck's semistable reduction theorem gives a finite Galois extension \(K'/K\) over whose integer ring the identity component of the Néron model is semiabelian with generic fibre \(A_{K'}\) \cite[Exp.~IX, Th.~3.6]{grothendieck-1972}. For the extension of homomorphisms, let \(G,H\) be semiabelian schemes over the same normal Noetherian scheme \(S\). If \(U\subset S\) is dense and open, restriction gives a bijection \(\operatorname{Hom}_S(G,H)\to\operatorname{Hom}_U(G_U,H_U)\) by \cite[Prop.~3.3.1.5]{lan-2013}. In particular the endomorphisms defining the coefficient action extend, and their addition and composition laws extend by uniqueness.
\leanref{GlobalTheory.\allowbreak{}SS.\allowbreak{}semistableReductionSt}{Development.\allowbreak{}AbelianVarietyTheory.\allowbreak{}Proofs.\allowbreak{}DegenerationProofs.\allowbreak{}SemistableReductionPpavPadicField}{19--21}
\leanref{GlobalTheory.\allowbreak{}CompanionTheory.\allowbreak{}SAH.\allowbreak{}semiAbelianHomExtension}{Development.\allowbreak{}AutomorphicSpectrum.\allowbreak{}Proofs.\allowbreak{}GlobalProofs.\allowbreak{}SemiAbelianHomExtension}{13--20}
 For a nondegenerate Hermitian space of rank at least two over a CM extension, rational special unitary points together with the omitted local factors are dense in the adelic group if the product of those local factors is noncompact. Nonsingular Hermitian forms of equal rank over a CM extension are isometric when their signatures agree and the ratio of their determinants is a norm. For \(x\delta\ne0\) in a number field, \(x=a^2-\delta b^2\) has a solution if it has one at every real and finite place.

The Green assertion of \ref{B7.31} gives distinct generic irreducible characters for its squarefree parameters. For \(\GL_N(\F_g)\), each such degree times the order of its torus is \(\prod_{i=1}^N(g^i-1)\). Over a field of characteristic \(p\), the induced module for a dominant \(\GL_N\)-weight \(\lambda\) is nonzero and simple when \(\lambda_i-\lambda_j+j-i\le p\) for all \(i<j\). Over an algebraically closed field a principal polarisation has a one-dimensional section space, and the geometric \(n\)-torsion of a \(g\)-dimensional abelian variety has order \(n^{2g}\) for invertible \(n\).

In item~5 of Definition~\ref{B8.25l}, the de Rham module over the affine base is finite projective of constant rank \(2g\), and the Hodge summand has constant rank \(g\). Over a nilpotent divided-power thickening with \(p\) nilpotent they give the crystal and the lifting of objects and homomorphisms with compatible Hodge filtrations. Let \(O=\mathcal O_C\) be the given \(p\)-adically complete valuation ring, with \(C\) algebraically closed of characteristic zero and \(p\) a nonunit. Put \(O^\flat=\varprojlim_{x\mapsto x^p}O/p\), \(A_{\inf}=W(O^\flat)\), and write \(\varphi\) for Witt Frobenius. From the compatible primitive roots of unity form \(\epsilon\in O^\flat\), \(\mu=[\epsilon]-1\), \(\xi=\mu/\varphi^{-1}(\mu)\), and \(\widetilde\xi=\varphi(\xi)\). The ring \(O\) is integral perfectoid: it is \(p\)-torsion-free and complete, every element of \(O\) has a \(p\)-th root in \(O\), and \(p^{1/p}\in O\). Fontaine's map \(A_{\inf}\to O\) has kernel \((\xi)\). Use the contravariant prismatic Dieudonn\'e functor on this perfectoid base \cite[Prop.~4.12, Prop.~4.48 and Th.~4.74]{anschutz-lebras-2023}. For a \(p\)-divisible group \(G\) of height \(h\) its module \(M_G\) is finite projective of rank \(h\), hence free since \(A_{\inf}\) is local. In the prism convention with ideal \((\widetilde\xi)\), its Frobenius \(F_G\) satisfies \(\widetilde\xi M_G\subseteq F_G(M_G)\subseteq M_G\). Here the image of the semilinear map is the image of its linearisation since \(\varphi\) is an automorphism. The crystalline comparison of \cite[Th.~17.5.2]{scholze-weinstein-2020}, dualised as in \cite[Prop.~4.48]{anschutz-lebras-2023}, identifies base change along Fontaine's map with the dual of the Lie algebra of the universal vector extension of \(G\). Dualising its Lie sequence gives the functorial exact sequence \(0\to\omega_G\to M_G/\xi M_G\to\operatorname{Lie}(G^\vee)\to0\). The constant group \(\Q_p/\Z_p\) has the unit module with Frobenius \(\varphi\). For \(\mu_{p^\infty}\) take the dual of the covariant Tate twist \(\mu^{-1}A_{\inf}(1)\) of \cite[Ex.~4.24]{bhatt-morrow-scholze-2018}. In the dual basis Frobenius is \((\varphi(\mu)/\mu)\varphi=\widetilde\xi\varphi\). Full faithfulness identifies each Frobenius-compatible map \(M_{\mu_{p^\infty}}\to M_{\Q_p/\Z_p}\) with a unique homomorphism \(\Q_p/\Z_p\to\mu_{p^\infty}\).
\leanref{UpperTheory.\allowbreak{}Lit.\allowbreak{}prismaticDieudonneHodge}{Development.\allowbreak{}AutomorphicArithmetic.\allowbreak{}ThetaCorrespondence.\allowbreak{}PrismaticDieudonne.\allowbreak{}PrismaticDieudonneHodge}{19--26}

The nef numerical assertions give nonnegative intersections of nef products and characterise bigness by positive top degree. They also give the leading section dimension. A geometrically integral projective image of dimension \(D\) in characteristic zero has a Chow form unique up to scalar, of degree equal to the image degree in each of its \(D+1\) blocks. The separation argument uses Nadel vanishing and the construction of an isolated non-klt locus. Principalisation makes the base ideal invertible. Let \(R\) be Noetherian and complete for an ideal \(I\), and put \(R_j=R/I^{j+1}\). For \(X\) proper and \(Z\) separated of finite type over \(R\), compatible \(R_j\)-morphisms \(X_{R_j}\to Z_{R_j}\) determine a unique \(R\)-morphism \(X\to Z\) by \cite[III$_1$, Th.~5.4.1]{grothendieck-1961}. For a system of proper \(R_j\)-schemes \(X_j\) with \(X_{j+1}\times_{R_{j+1}}R_j\simeq X_j\), choose compatible invertible sheaves \(L_j\). If \(L_0\) is relatively ample, the induced proper formal scheme has an invertible sheaf ample on its closed fibre. The algebraisation theorem produces a projective \(R\)-scheme with a relatively ample line inducing the given system \cite[III$_1$, Th.~5.4.5]{grothendieck-1961}.
\leanref{PositivityTheory.\allowbreak{}FormalThickening.\allowbreak{}formalAlgebraisation}{Development.\allowbreak{}WeilBound.\allowbreak{}CoreStatements.\allowbreak{}TemperedRepresentationTheory.\allowbreak{}FormalAlgebraisation.\allowbreak{}Development}{17--23}

For item~19 of Definition~\ref{B8.25l}, if \(X\) is smooth integral projective over an algebraically closed field of characteristic zero and \(N\) is nef and big, then \(H^a(X,K_X\otimes N)=0\) for \(a>0\). This is the line-bundle case of \cite{kawamata-1982}. The comparison in item~25 is degree-zero GAGA for a line bundle on a smooth projective complex variety \cite{serre-1956}.

For item~21 of Definition~\ref{B8.25l}, let \(X\) be a compact complex \(n\)-fold, \(n\ge1\), admitting a positive line. Take sections \(s_1,\ldots,s_n\) of a Hermitian line \(L\) with \(\sum_j|s_j|_L^2\le1\) and transverse finite common zero set \(Z\). Write \(\phi_L,\phi_E\) for the metric weights of \(L\) and of the extension line \(E\), and put \(\psi=n\log\sum_j|s_j|_L^2\). Suppose \(\phi_E+\alpha\psi\) is semipositive off \(Z\) for \(1\le\alpha\le1+\delta\), where \(\delta>0\). Values \(f(P)\) in local frames of \(K_X\otimes E\) extend to \(F\in H^0(X,K_X\otimes E)\) with
\begin{equation}\label{eq:p8c3-extension}
 2^n\int_X\frac{|F|^2e^{-\psi}}{1+\delta^2\psi^2}\,d\mu
 \le \frac{34}{\delta}\frac{2^{n+1}\pi^n}{(n-1)!}
 \sum_{P\in Z}\left|\frac{f(P)}{\operatorname{Jac}s(P)}\right|^2
 e^{-\phi_E(P)+n\phi_L(P)} .
\end{equation}
Here \(d\mu\) is the chart measure and the Jacobian is that of the local coefficients of the \(s_j\). This is the transverse finite-zero-set case of \cite[Th.~2.8]{demailly-2016} in the indicated residue normalisation.

The proper-intersection formula used with \(\mathcal H_P\) is a consequence of the intersection laws. For a finite proper cut of a projective integral \(d\)-fold by \(d\) sections of a line, the local rings of the \(d\)-fold at the cut are assumed Cohen--Macaulay. The local equations form a system of parameters and hence a regular sequence. The local top intersection multiplicity is the length of its quotient. Push the resulting zero-cycle to the ground field to obtain the sum of these lengths weighted by residue degrees.
\leanref{GlobalTheory.\allowbreak{}CovTop.\allowbreak{}PublishedHodgeBundleTheorems}{Development.\allowbreak{}AutomorphicArithmetic.\allowbreak{}ThetaCorrespondence.\allowbreak{}ThetaPairLifts.\allowbreak{}HodgeBundleWeights.\allowbreak{}ToKnownTheoremsOf}{7--35}
\leanref{GlobalTheory.\allowbreak{}CovTop.\allowbreak{}PublishedHodgeBundleTheorems.\allowbreak{}q12}{Development.\allowbreak{}AutomorphicArithmetic.\allowbreak{}ThetaCorrespondence.\allowbreak{}ThetaPairLifts.\allowbreak{}HodgeBundleWeights.\allowbreak{}ToKnownTheoremsOf}{38--38}
\leanref{GlobalTheory.\allowbreak{}Published.\allowbreak{}L2ExtensionLogCanonical}{Development.\allowbreak{}AbelianVarietyTheory.\allowbreak{}PublishedTheorems.\allowbreak{}Divisors₂.\allowbreak{}L2ExtensionLogCanonical}{27--54}

For a reduced separated scheme of finite type over a characteristic-zero field, resolution is an isomorphism over the smooth locus. For principalisation take a smooth separated quasi-compact scheme over such a field with simple normal crossings boundary. An ideal whose support misses the generic points of the irreducible components becomes effective Cartier with simple normal crossings support after blow-ups compatible with the boundary. The modification is an isomorphism off the support of the ideal. Item~27 of Definition~\ref{B8.25l} uses this for a nonzero ideal on a smooth integral projective variety. These are the characteristic-zero resolution and principalisation results of \cite{hironaka-1964}. Invariant Poincar\'e--Lelong gives the divisor current of a Hermitian section from the curvature and the \(dd^c\) of its log norm. Together with resolution it gives the intersection and restriction laws used here.

For the three arithmetic Hilbert--Samuel assertions, take the projective arithmetic variety \(X\) of Lemma~\ref{B8.25h} with generic dimension \(D\). Write \(I(A)=\widehat{\deg}(\widehat c_1(A)^{D+1})\) for its self-intersection on a smooth conjugation-invariant Hermitian line \(A\). The arithmetic Kleiman assertion is
\begin{equation}\label{eq:p8c3-kleiman}
 I(A)\ge0
\end{equation}
when \(A\) is semipositive, nef on fibres and of nonnegative degree on every horizontal arithmetic curve. Indeed these hypotheses make the Hermitian line arithmetically nef. Its top self-intersection equals its volume by \cite[Cor.~5.5]{moriwaki-2009}, and every unit ball contains zero. For the two upper bounds assume \(D\ge1\), use the homogeneous ideal \(P\) of Lemma~\ref{B8.25h}, and fix an ideal \(P_0\subseteq P\). In both bounds \(A\) is conjugation-invariant and its curvature is nonnegative at complex points of \(P\) along the cone-tangent directions of \(P_0\). The asserted inequality is
\begin{equation}\label{eq:p8c3-hs-upper}
 \limsup_{j\to\infty}\frac{(D+1)!\log\#\{s\in L_j:\|s\|_{\sup}\le1\}}{j^{D+1}}
 \le I(A).
\end{equation}
For the first bound take \(L_j=S_j\) of Lemma~\ref{B8.25h}. The coefficient field has characteristic zero, the generic degree is positive, and the integral sections of sup norm at most one specified in the hypothesis generate the line: their coefficients generate the unit ideal in each \(R_c\). For the second take \(L_j=\widetilde S_j\). The multidegree is nonnegative and each coordinate monomial of that multidegree has point norm at most one. These monomials generate, since on each chart one of them is the chart frame. In either case a generated line has nonnegative degree on vertical curves. On the normalisation of any horizontal arithmetic curve choose a generating section nonzero at its generic point. Its finite zero divisor is effective and each contribution \(-\log\|s\|\) at infinity is nonnegative. Thus the line has nonnegative arithmetic degree there. Lemma~\ref{B8.25h} gives smooth semipositive curvature on the generically smooth modification. Pullback preserves generation and the sup bounds. Moriwaki's nef Hilbert--Samuel formula, with a fixed Hermitian twist if required, now gives \eqref{eq:p8c3-hs-upper} and \eqref{eq:p8c3-nef-count} with denominator \(j^{D+1}/(D+1)!\) \cite[Cor.~5.5]{moriwaki-2009}. Zhang's relatively semipositive formula for a generically ample line uses the same intersection normalisation for the Euler characteristic of the sup-norm lattice \cite[Th.~1.4]{zhang-1995}.
\leanref{PositivityTheory.\allowbreak{}MbHsUse.\allowbreak{}HSUse}{Development.\allowbreak{}WeilBound.\allowbreak{}CoreStatements.\allowbreak{}AutomorphicKernels.\allowbreak{}MoriwakiBigness.\allowbreak{}HilbertSamuelApplications.\allowbreak{}SmallGeneratedCountUpperOfUse}{12--20}
\leanref{GroundTheory.\allowbreak{}Heights₂.\allowbreak{}Heights.\allowbreak{}HilbertSamuel.\allowbreak{}G7.\allowbreak{}Statement}{GroundTheory.\allowbreak{}HeightsAndTheta.\allowbreak{}Heights.\allowbreak{}HilbertSamuel.\allowbreak{}StatementsArith}{437--444}
\leanref{ReciprocityTheory.\allowbreak{}ChartRestriction.\allowbreak{}G.\allowbreak{}Geo.\allowbreak{}ArithmeticHilbertSamuelUpperBound}{ReciprocityTheory.\allowbreak{}ReductiveGroups.\allowbreak{}ArithmeticGeometry.\allowbreak{}BirationalGeometry.\allowbreak{}ArithmeticHilbertSamuelLowerBound}{135--145}
\leanref{PositivityTheory.\allowbreak{}Pf1.\allowbreak{}HilbertSamuelMoriwaki.\allowbreak{}hsUse\_\allowbreak{}outright}{Development.\allowbreak{}WeilBound.\allowbreak{}CoreStatements.\allowbreak{}AutomorphicKernels.\allowbreak{}MoriwakiBigness.\allowbreak{}MoriwakiBigTheorem.\allowbreak{}Outright}{13--16}
\leanref{PositivityTheory.\allowbreak{}Pf1.\allowbreak{}HilbertSamuelMoriwaki.\allowbreak{}moriwakiBig\_\allowbreak{}outright}{Development.\allowbreak{}WeilBound.\allowbreak{}CoreStatements.\allowbreak{}AutomorphicKernels.\allowbreak{}MoriwakiBigness.\allowbreak{}MoriwakiBigTheorem.\allowbreak{}Outright}{9--12}

For Fargues exactness, let \(C\) be an algebraically closed valued field of characteristic zero with \(v(p)<1\) and \(p\)-adically complete integer ring \(O_C\). Normalise the dual Tate pairing using a compatible system \(\zeta_0=1,\ \zeta_{a+1}^p=\zeta_a,\ \zeta_1\ne1\). For a \(p\)-divisible group \(G\) over \(O_C\), the linear Hodge--Tate map on \(T_pG\otimes O_C\) has kernel the annihilator of the kernel of the dual logarithmic differential on \(T_pG^\vee\otimes O_C\). Over \(C\) this is the duality of the Hodge--Tate subspaces in \cite[Prop.~5.2.8]{scholze-weinstein-2013}. The targets of the integral differentials are torsion-free modules of invariant differentials, so their kernels are the intersections of the kernels over \(C\) with the respective Tate lattices.
\leanref{UpperTheory.\allowbreak{}LiteratureTheorems.\allowbreak{}fargues\_\allowbreak{}exact\_\allowbreak{}of\_\allowbreak{}literature}{Development.\allowbreak{}AutomorphicArithmetic.\allowbreak{}ThetaCorrespondence.\allowbreak{}ThetaPairLifts.\allowbreak{}ThirdPlaceSeries.\allowbreak{}EveryCharacterOfLiteratureOfVWin}{549--550}

In Zarhin's retraction, \(A\) is an abelian scheme with a relatively ample line \(L\) over the integer ring of an algebraically closed valued field. Put \(P=(A\times A^\vee)^4\). The polarisation \(\ell:A\to A^\vee\) of \(L\) has a quasi-inverse \(\mu:A^\vee\to A\) with \(\mu\ell=[m]\) and \(\ell\mu=[m]\), \(m\ge1\). A quaternion matrix \(U\) from four squares summing to \(4m^2-1\) satisfies \(UU^{\mathsf t}=U^{\mathsf t}U=(4m^2-1)I_4\). On \(A^4\times(A^\vee)^4\) take the pullbacks of \(L\) by the four projections and their negatives, and by the four projections composed with \(\mu\) and their negatives. Tensor them with the pullbacks of the Poincar\'e line on each pair \((A_i,A^\vee_j)\) to the powers \(U_{ij}\). The theorem of the cube computes the associated map to the dual, with the factors in the order \((A^\vee)^4\times A^4\), as
\begin{equation}\label{eq:p8c3-zarhin}
 \begin{pmatrix} 2\ell I_4 & U\\ U^{\mathsf t} & 2m\mu I_4\end{pmatrix}.
\end{equation}
The candidate inverse has diagonal entries \(2m\mu I_4,2\ell I_4\) and off-diagonal entries \(-U,-U^{\mathsf t}\). Multiplication gives the identity by the two quasi-inverse identities and the quaternion norm identity. For ampleness on each geometric fibre, pull the line back by the isogeny \(1\times\ell^4:A^4\times A^4\to P\). Its matrix relative to the product polarisation is \(\left(\begin{smallmatrix}2I_4&U\\U^{\mathsf t}&2m^2I_4\end{smallmatrix}\right)\). Its Schur complement is \(\tfrac12 I_4\). Rational congruence diagonalises it with positive diagonal, so after clearing denominators an isogeny pulls the line back to the class of a product of positive powers of \(L\). Ampleness descends under the isogenies.

Translate the polarising line on the geometric generic fibre to a symmetric representative of the same polarisation. Extend its Picard class over the valuation ring and choose a line bundle \(M\) representing that class. Injectivity of restriction to the generic fibre for the Picard groups of \(P\) and \(P\times P\) gives symmetry and identifies the polarisation of \(M\) with \eqref{eq:p8c3-zarhin}. The base is local, so \(M\) can be rigidified at the zero section. On the closed fibre \(M\) induces the same polarisation as the ample line from the matrix construction and is therefore ample there. Fibrewise ampleness gives relative ampleness, and the principal polarisation gives trivial kernel. Finally take \(i:A\to P\) to be inclusion in the first \(A\)-factor and \(r:P\to A\) its projection. These maps are defined over the valuation ring and \(ri=1_A\).
\leanref{PositivityTheory.\allowbreak{}HtPt.\allowbreak{}Narrow.\allowbreak{}zarhin}{Development.\allowbreak{}WeilBound.\allowbreak{}CoreStatements.\allowbreak{}HodgeTatePoints.\allowbreak{}ZarTop}{6--8}
\leanref{PositivityTheory.\allowbreak{}PrwZar.\allowbreak{}zarhinStatement\_\allowbreak{}of\_\allowbreak{}dual\_\allowbreak{}cube\_\allowbreak{}qiR}{Development.\allowbreak{}WeilBound.\allowbreak{}CoreStatements.\allowbreak{}AbelianHeightsDuality.\allowbreak{}ZarhinTrick.\allowbreak{}ZarhinStatementOfDualCubeQiR}{27--66}
\leanref{GlobalTheory.\allowbreak{}CompanionTheory.\allowbreak{}Hypotheses₂.\allowbreak{}Zarhin.\allowbreak{}SymB.\allowbreak{}symmetricBundle}{Development.\allowbreak{}AutomorphicSpectrum.\allowbreak{}Proofs.\allowbreak{}AlgebraProofs.\allowbreak{}RelAmpleOfCls}{258--308}

For nef Riemann--Roch let \(y\) be birational coordinates of generic dimension \(D\) and \(a\) a generically nef multidegree. The leading term is
\begin{equation}\label{eq:p8c3-nef-rr}
 \lim_{t\to\infty}\frac{D!}{t^D}
 \dim_{\Q}\langle\widetilde S_t\rangle_{\Q}
 =\deg_{\rm gen}(a).
\end{equation}
Here \(\widetilde S_t\) is the lattice of sections on the normalisation defined in \ref{B8.25h}. Its rational span contains the span for the unshifted model and embeds, after multiplication by a fixed nonzero rational representative, in a span with fixed line-bundle twist. Both bounding dimensions have the leading coefficient in \eqref{eq:p8c3-nef-rr} by the nef dimension formula.
\leanref{PositivityTheory.\allowbreak{}Tls2.\allowbreak{}Nef.\allowbreak{}nefRR}{Development.\allowbreak{}WeilBound.\allowbreak{}CoreStatements.\allowbreak{}TemperedRepresentationTheory.\allowbreak{}NefDivisors.\allowbreak{}NefRR}{14--53}

For each Siegel genus \(g\ge1\) choose a smooth polarised fan without self-intersections. This choice works at every odd level \(n\ge3\). Over \(\Z[\zeta_n,1/n]\) the compactification is projective and smooth of relative dimension \(g(g+1)/2\). The open is fibrewise dense and is the abelian locus of the semiabelian extension. Let \(\omega\) be the determinant Hodge line. The polarised charts give an effective Cartier divisor \(D_\phi\), a generated power \(\omega^N\), \(N\ge1\), and an ample twist \(\omega^{Nk_0}(-D_\phi)\). Tensor this twist with the generated line \((\omega^N)^{k-k_0}\) to obtain ampleness for every \(k\ge k_0\) with the same divisor. This uses the integral Siegel construction of \cite[Chs.~IV--V]{faltings-chai-1990} at odd full level.
\leanref{GlobalTheory.\allowbreak{}siegelToroidalIntegral\_\allowbreak{}holds}{Development.\allowbreak{}AbelianVarietyTheory.\allowbreak{}Proofs.\allowbreak{}DivisorProofs.\allowbreak{}SiegelToroidalIntegralHolds}{14--16}
\leanref{GlobalTheory.\allowbreak{}PolDeg.\allowbreak{}siegelToroidalIntegral\_\allowbreak{}of\_\allowbreak{}polCharts}{Development.\allowbreak{}AbelianVarietyTheory.\allowbreak{}Proofs.\allowbreak{}DivisorProofs.\allowbreak{}SiegelIntegralToroidalOfPolarisedCharts}{16--29}

For the unitary assertion fix a CM field \(E\), host \(H\) and moduli datum of signature \((p,q)\), \(p+q\ge1\). A prescribed polarising function \(\phi\) is defined on the semidefinite balanced forms on finite projective \(\mathcal O_E\)-modules of rank at most \(\min(p,q)\). A symmetric form \(B\) is balanced when \(B(ax,y)=B(x,\bar a y)\). Pullback of a form along a surjection of lattices preserves the value of the function. On a finite projective lattice one may compute the value after pullback along a free presentation. Two such presentations have a common comparison by taking their direct sum and lifting the two surjections using projectivity. The function is homogeneous and superadditive. In particular the functions on cusp lattices agree on faces where the monodromy form factors through a quotient. For every integral positive semidefinite balanced form \(B\), the value \(\phi(B)\) is an integer, positive if \(B\ne0\). The fan covers the semidefinite cone with the rational boundary allowed by the lattice. Its cones are closed under faces and have finitely many types up to lattice automorphism. Their primitive generators extend to integral bases. Each cone satisfies the strict support condition. There is an integral linear form \(\ell\) and an integer \(N>0\) with \(\ell\ge N\phi\) on integral positive semidefinite balanced forms and equality on the cone generators. Equality for an integral form \(B\) forces some positive integer multiple of \(B\) to be a nonnegative integral combination of those generators. The no-self-intersections condition keeps distinct rays of a cone inequivalent under lattice automorphisms.
\leanref{GlobalTheory.\allowbreak{}PolFan.\allowbreak{}IsCone}{Development.\allowbreak{}AbelianVarietyTheory.\allowbreak{}Definitions.\allowbreak{}AutomorphicMembership.\allowbreak{}SumMulEqZeroOfSumToReal}{292--302}
\leanref{GlobalTheory.\allowbreak{}PolFan.\allowbreak{}IsPolFanOn}{Development.\allowbreak{}AbelianVarietyTheory.\allowbreak{}Definitions.\allowbreak{}AutomorphicMembership.\allowbreak{}SumMulEqZeroOfSumToReal}{405--422}

For each level of parameter \(M'\), open moduli scheme and coefficient field \(F\) of the datum, there is a model \(Y\) over \(F\), smooth of the moduli dimension with simple normal crossings boundary. There are \(d,k\ge1\) and an effective Cartier divisor \(D_\phi\) supported off the abelian locus. At every boundary point an \'etale completed chart \(k'[[t_1,\ldots,t_s,z_1,\ldots,z_e]]\) has
\begin{equation}\label{eq:p8c3-fan-chart}
 I_{\partial,\mathrm{red}}=(t_1\cdots t_s),\qquad
 \mathcal O(-D_\phi)=\left(\prod_{i=1}^s t_i^{a_i}\right),\qquad
 a_i=d\phi(\gamma_i),\quad \gamma_i=M'\delta_i .
\end{equation}
Here \(\delta_i\) are the primitive balanced forms of the chart cone, and \(\gamma_i\) are the monodromy forms on the level lattice. If the determinant Hodge module is a line \(\omega\), then \(\omega^k(-D_\phi)\) is ample. The logarithmic Kodaira--Spencer map is an isomorphism on the designated interior and agrees there with the map of the moduli open. The PEL toroidal construction corresponding to the characteristic-zero model is \cite[Th.~6.4.1.1]{lan-2013}.
\leanref{GlobalTheory.\allowbreak{}Published.\allowbreak{}ToroidalForFanKS}{Development.\allowbreak{}AbelianVarietyTheory.\allowbreak{}PublishedTheorems.\allowbreak{}Divisors₂.\allowbreak{}ToroidalCompactificationBis}{45--54}
\leanref{GlobalTheory.\allowbreak{}MumfordChart.\allowbreak{}HasChartOf}{Development.\allowbreak{}AbelianVarietyTheory.\allowbreak{}PublishedTheorems.\allowbreak{}AutomorphicForms₂.\allowbreak{}ToroidalChartsForFan}{17--30}

Start over \(\mathcal O_F[1/(3M'd_H\operatorname{disc}E)]\), where \(d_H\) is the host prime. Put \(g=[E:\Q](p+q)/2\). Choose an auxiliary theta level \(h=81\cdot3^j>81^g g!\). The theta map at this level has finite geometric fibres and target \(\mathbb P^n\), \(n+1=h^{4g}\). Normalisation in the theta map gives a normal model finite over projective space, with the moduli open dense on the generic fibre. The pullback of \(\mathcal O(8(n+1)!)\) is ample. Denote it by \(A\). The theta transformation law identifies its pullback on a model with Hodge line with \(\omega^N\), where \(N=4((n+1)!)^2\).
\leanref{GlobalTheory.\allowbreak{}umin\_\allowbreak{}theta\_\allowbreak{}compactification\_\allowbreak{}exists}{Development.\allowbreak{}AbelianVarietyTheory.\allowbreak{}Proofs.\allowbreak{}DivisorProofs.\allowbreak{}UnitaryIntegralCompactificationExistsViaTheta}{78--143}

On a component meeting the open, take the finite theta cover and the normal projective model of Lemma~\ref{B8.25j}. Fix a point and a closed specialisation \(x\) of it. Apply the lemma at \(x\), with \(N=N_0\) the theta modulus, to obtain \(\mu_c,e_c\). To identify the lattice, take a trait centred at \(x\) with generic point in the theta open. Such a trait is obtained by the one-dimensional quotient construction of Lemma~\ref{B8.25i} and completion, followed by finite extension for the polarised degeneration with level. Lemma~\ref{B8.25k} identifies its lattice \(X_0\) with the completed period lattice. Its monodromy form \(b_0\) is positive definite by positivity of the periods and is balanced by compatibility of the polarisation with conjugation on \(E\). The character lattice is torsion-free and finite over the Dedekind ring \(\mathcal O_E\), hence projective. Its rank is at most each of the two signature multiplicities: at either conjugate embedding its torus contributes that many dimensions to the corresponding Lie space. Thus its rank is at most \(\min(p,q)\).
\leanref{GlobalTheory.\allowbreak{}UTraitCen.\allowbreak{}exists\_\allowbreak{}uTrait\_\allowbreak{}cen}{Development.\allowbreak{}AbelianVarietyTheory.\allowbreak{}Proofs.\allowbreak{}DivisorProofs.\allowbreak{}ExistsUnitaryTraitCentredClosedPoint}{18--48}
\leanref{GlobalTheory.\allowbreak{}UPoinc.\allowbreak{}UTrait.\allowbreak{}finrank\_\allowbreak{}le\_\allowbreak{}min}{Development.\allowbreak{}AbelianVarietyTheory.\allowbreak{}Proofs.\allowbreak{}DivisorProofs.\allowbreak{}CharacterGroupFinrankLeMinSignatureUnitaryTrait}{17--51}

For each degeneration trait in Lemma~\ref{B8.25k}, pull its monodromy form back to \(X_0\) and call the result \(B\). The lemma gives \(B(\mu_c,\mu_c)\le B(\alpha,\alpha)\) for every \(\alpha\) in the class \(c\), including when \(B=b_0\). With the order normalised to be one on a uniformiser it gives
\begin{equation}\label{eq:p8c3-theta-orders}
 \operatorname{ord}(e_c^{m_1})=m_2 B(\mu_c,\mu_c),
 \qquad m_1=2N_0,\quad m_2=1 .
\end{equation}
A trait whose centre specialises to the originally fixed point also specialises to \(x\), by transitivity. The same choices therefore apply at that point.
\leanref{GlobalTheory.\allowbreak{}T.\allowbreak{}unitaryComparison\_\allowbreak{}lattice\_\allowbreak{}periods\_\allowbreak{}at\_\allowbreak{}point}{Development.\allowbreak{}AbelianVarietyTheory.\allowbreak{}Proofs.\allowbreak{}DivisorProofs.\allowbreak{}UnitaryComparisonLatticePeriodsAtPointStatement}{13--14}
\leanref{GlobalTheory.\allowbreak{}UPoinc.\allowbreak{}ClosedPt.\allowbreak{}exists\_\allowbreak{}lattice\_\allowbreak{}data}{Development.\allowbreak{}AbelianVarietyTheory.\allowbreak{}Proofs.\allowbreak{}DivisorProofs.\allowbreak{}ExistsBalancedLatticeDataClosedPoint}{61--146}
\leanref{GlobalTheory.\allowbreak{}unitaryStepA\_\allowbreak{}at\_\allowbreak{}point}{Development.\allowbreak{}AbelianVarietyTheory.\allowbreak{}Proofs.\allowbreak{}DivisorProofs.\allowbreak{}UnitaryComparisonLatticePeriodsAtPoint}{15--29}

Apply Lemma~\ref{B8.25m} on \(X_0\) with modulus \(N_0\ge3\). The trait through \(x\) gives the positive definite witness \(b_0\), and Lemma~\ref{B8.25k} gives its minimum inequalities for the chosen representatives. Thus a finite list of cones in \(X_0\) covers all the balanced semidefinite forms under consideration.

Clear denominators in the finitely many supporting linear forms. Their minimum on this region is a positive integer multiple of \(\phi\). Each linear form on balanced forms is a rational combination of the evaluations \(B\mapsto B(\mu_c,\mu_c)\). To see this, reduce modulo an odd prime dividing \(N_0\). The representatives exhaust the residue vectors there. Evaluation at the coordinate vectors and their pairwise sums spans the dual of the space of symmetric forms in odd characteristic, hence the evaluations span over \(\Q\) before restriction to balanced forms. The finite orbit set of minimiser systems and the finitely many isomorphism types of bounded-rank projective \(\mathcal O_E\)-lattices allow a common multiplier \(\iota_0>0\) and a finite list of coefficient sets. Transport in the finite residue modules permits the same list in lattice coordinates. Thus at the point one can choose a finite set \(\mathcal L\) of integral vectors \(\nu=(\nu_c)_c\) such that for every \(B\) in its minimiser region,
\begin{equation}\label{eq:p8c3-fan-minimum}
 \iota_0\phi(B)=\min_{\nu\in\mathcal L}\sum_c\nu_c B(\mu_c,\mu_c).
\end{equation}
\leanref{GlobalTheory.\allowbreak{}UPoinc.\allowbreak{}CellCones.\allowbreak{}exists\_\allowbreak{}finite\_\allowbreak{}cones}{Development.\allowbreak{}AbelianVarietyTheory.\allowbreak{}Proofs.\allowbreak{}DivisorProofs.\allowbreak{}ExistsFiniteConesPolarizedFanLevelStructure}{99--215}
\leanref{GlobalTheory.\allowbreak{}UPoinc.\allowbreak{}FinTypes.\allowbreak{}single}{Development.\allowbreak{}AbelianVarietyTheory.\allowbreak{}Proofs.\allowbreak{}DivisorProofs.\allowbreak{}LeastIntCombinationMinimiserCellUnitary}{21--58}
\leanref{GlobalTheory.\allowbreak{}UPoinc.\allowbreak{}finite\_\allowbreak{}types\_\allowbreak{}unitary}{Development.\allowbreak{}AbelianVarietyTheory.\allowbreak{}Proofs.\allowbreak{}DivisorProofs.\allowbreak{}LeastIntCombinationMinimiserCellUnitary}{126--186}

Take the tuple with entries \(\prod_c e_c^{m_1\nu_c}\), \(\nu\in\mathcal L\). Formula \eqref{eq:p8c3-theta-orders} gives the order of each entry by addition. The quotients are invertible at the generic point of each trait being tested, so negative powers also have these orders. By \eqref{eq:p8c3-fan-minimum} the least order is \(m_2\iota_0\phi(B)\). Invariance under the lattice surjection identifies \(\phi(B)\) with \(\phi(b)\) for the original monodromy form \(b\). The tuples come from a finite collection because the theta quotients and the coefficient sets range over finite sets.
\leanref{GlobalTheory.\allowbreak{}UPoinc.\allowbreak{}exists\_\allowbreak{}uLocTuples\_\allowbreak{}theta}{Development.\allowbreak{}AbelianVarietyTheory.\allowbreak{}Proofs.\allowbreak{}DivisorProofs.\allowbreak{}ExistsFiniteThetaTuplesLeastOrdersUnitary}{17--186}

Take a common normal projective model on which all the tuple ideals are principal. Their local Cartier equations glue to a boundary ideal with the prescribed trait orders. Descend a power of this ideal along the proper modification to an ideal \(J_1\) on the finite theta model, with orders \(\iota\phi(b)\) for a common integer \(\iota\ge1\).

Apply Lemma~\ref{B8.25o} to \(J_1\) on the finite theta model over this component. The two schemes are integral and normal, and their function fields have characteristic zero, so the generic extension is separable. Write \(h\) for its degree. On an extended trait the norm ideal becomes the product of the \(h\) branch ideals. After a further common extension choose semiabelian degeneration data for the moduli point. Every branch has that same underlying point; hence the order formula for \(J_1\) applies with the same monodromy form on all branches. If the total ramification index is \(e\), this form is \(eb\), since each order in the period pairing is multiplied by \(e\). Writing \(J\) for the norm ideal gives \(e\operatorname{ord}_V J=h\iota\phi(eb)=eh\iota\phi(b)\). Division by \(e\) gives the multiplier \(h\iota\) on the original trait.

On components disjoint from the moduli open take the unit ideal. By normality and noetherianity, the finitely many components are disjoint. Each is both open and closed. Raise the component ideals to powers giving a common multiple \(d\) of their order multipliers. The resulting ideal \(J\) has pullback
\begin{equation}\label{eq:p8c3-trait-ideal}
 J\mathcal O_{\rm trait}=(\varpi^{d\phi(b)})
\end{equation}
on each degeneration trait with parameter \(\varpi\).
\leanref{GlobalTheory.\allowbreak{}UPoinc.\allowbreak{}exists\_\allowbreak{}zhModel\_\allowbreak{}effCartier\_\allowbreak{}orders}{Development.\allowbreak{}AbelianVarietyTheory.\allowbreak{}Proofs.\allowbreak{}DivisorProofs.\allowbreak{}UnitaryExistsNormalModelEffCartierAlongTraits}{15--25}
\leanref{GlobalTheory.\allowbreak{}UPoinc.\allowbreak{}exists\_\allowbreak{}uSetup\_\allowbreak{}ideal\_\allowbreak{}over\_\allowbreak{}component}{Development.\allowbreak{}AbelianVarietyTheory.\allowbreak{}Proofs.\allowbreak{}DivisorProofs.\allowbreak{}ExistsUnitaryThetaSetupIdealOverComponent}{17--25}
\leanref{GlobalTheory.\allowbreak{}UPoinc.\allowbreak{}UNormDown.\allowbreak{}normDescent}{Development.\allowbreak{}AbelianVarietyTheory.\allowbreak{}Proofs.\allowbreak{}DivisorProofs.\allowbreak{}UnitaryNormDescentIdealToComponent}{26--76}
\leanref{GlobalTheory.\allowbreak{}UPoincDown.\allowbreak{}idealDescent}{Development.\allowbreak{}AbelianVarietyTheory.\allowbreak{}Proofs.\allowbreak{}DivisorProofs.\allowbreak{}IdealDescentToComponentCompactification}{45--72}
\leanref{GlobalTheory.\allowbreak{}UPoinc.\allowbreak{}uIdealSt\_\allowbreak{}holds}{Development.\allowbreak{}AbelianVarietyTheory.\allowbreak{}Proofs.\allowbreak{}DivisorProofs.\allowbreak{}UnitaryCompactificationExistsBoundaryIdeal}{19--39}

On the generic fibre \(Z\) take the normalisation \(S\) of the blow-up \(S_0=\operatorname{Proj}_Z\bigoplus_{j\ge0}J^j\) and let \(\pi:S\to Z\) be the projection. Over the unit-ideal locus this map is an isomorphism. The tautological line \(J\mathcal O_{S_0}=\mathcal O_{S_0}(1)\) is relatively ample. Its pullback by the finite normalisation is \(J\mathcal O_S=\mathcal O_S(-D_\phi)\) and is relatively ample over \(Z\). To see the absolute ampleness explicitly, choose \(v\) so large that \(J\otimes A^v\) is generated by global sections on the projective scheme \(Z\). A finite generating set gives a surjection from a polynomial algebra onto \(\bigoplus_{j\ge0}J^j\otimes A^{vj}\). Relative Proj embeds \(S_0\) as a closed subscheme of \(Z\times\mathbb P^m\) for some \(m\); twisting the graded pieces by \(A^{vj}\) leaves the underlying Proj unchanged. The restriction of the second factor's \(\mathcal O(1)\) is \(J\mathcal O_{S_0}\otimes\pi_0^*A^v\), where \(\pi_0:S_0\to Z\). For \(u>v\), the external product \(A^{u-v}\boxtimes\mathcal O(1)\) is ample. Its restriction to \(S_0\) and pullback under the finite normalisation give the ample line \(\pi^*A^u(-D_\phi)\). Write \(D_S\) for the Cartier divisor of \(J\mathcal O_S\). We compare the two models on one component of the moduli open. Let \(Y\) be the toroidal compactification for the prescribed fan. The lattice with its conjugate-adjoint action and principal polarisation is a PEL datum in characteristic zero. On a geometric component the Weil pairing is fixed. The specified half-level fixes one summand of the \(M\)-torsion, and the complementary summand is its perfect dual since \(\mathfrak b\bar{\mathfrak b}=(M)\) and the two ideals are coprime. An isometry fixing the first summand acts identically on the second. Thus the component group has principal congruence level \(M\ge3\), with \(M\) odd. It is neat: at any odd prime \(p\mid M\) the eigenvalues of a matrix congruent to the identity modulo \(p\) belong to \(1+p\overline{\Z}_p\). Their products and inverses lie there too, and the only root of unity in this group is one. This check may be made after adjoining the level roots of unity. Isomorphisms extending the identity on the moduli open descend uniquely. The strict supporting forms give a projective cone decomposition, smooth by the basis condition and compatible at the cusps by invariance under lattice quotients. Apply \cite[Ths.~6.4.1.1 and 7.3.3.4]{lan-2013}, in the characteristic-zero formulation of \cite[Construction~3.1 and Th.~3.10]{lan-2017}. The Siegel construction is \cite[Ch.~V, Th.~5.8]{faltings-chai-1990}. For some positive multiple \(d_0\) of the polarising function this gives on \(Y\) a divisor \(D_Y\) with orders \(d_0\phi(\gamma_i)\). Its negative is relatively ample over the minimal compactification, on which a positive Hodge power is ample. Consequently \(\omega^{Nud_0}(-dD_Y)\) is ample for all sufficiently large divisible \(u\).\par Let \(G\) be the normalisation of the closure of the common open in \(Y\times_F S\), with projections \(a:G\to Y\), \(b:G\to S\). Both are proper birational morphisms between normal varieties. On \(G\) use the pullback of the semiabelian family on \(Y\). Its abelian locus is the common open, since the graph is already closed over that open. The theta-line identification applied on this model to \(\pi b:G\to Z\) gives \(b^*\pi^*A\simeq a^*\omega^N\). We claim that \(d_0 b^*D_S=d a^*D_Y\). Test at the completed local ring of each prime divisor of \(G\), a complete discrete valuation ring with generic point on the moduli open. Over the interior both orders are zero. At a boundary centre use the toroidal chart of \(Y\); one may extend the valued field to split the cusp data. If the orders of its boundary parameters are \(v_i\), the pulled-back family has monodromy form \(B=\sum_i v_i\gamma_i\): this follows by taking orders of the monomials of the chart on each period. The form lies in the chart cone, where \(\phi\) is linear. Hence the order of \(a^*D_Y\) is \(d_0\sum_i v_i\phi(\gamma_i)=d_0\phi(B)\). Equation~\eqref{eq:p8c3-trait-ideal}, applied to the same trait through \(G\to S\to Z\) and the same semiabelian family, gives order \(d\phi(B)\) for \(b^*D_S\). A ramified extension multiplies both orders by its ramification index. Normality of \(G\) now proves the claimed equality of Cartier divisors.\par Choose \(u\) sufficiently large and divisible for both ample twists. The pullbacks to \(G\) of \(\omega^{Nud_0}(-dD_Y)\) and \((\pi^*A^u(-D_S))^{d_0}\) are isomorphic. For a proper birational morphism onto a normal variety the direct image of the structure sheaf is the structure sheaf. The projection formula thus identifies both graded section rings with the section ring of the common pullback, compatibly with multiplication. Taking \(\operatorname{Proj}\) recovers \(Y\) and \(S\) from their respective ample lines. The resulting isomorphism extends the original identification of the open, by construction on its graph. In particular \(Y\to S\) is finite. On this model \(\pi^*A\simeq\omega^N\), so \(k=Nu\); the divisor of \(J\) has exponents \(d\phi(\gamma_i)\) in \eqref{eq:p8c3-fan-chart}. This also gives the valuation normalisation for a ramified trait or level pullback. In a framed chart the logarithmic Kodaira--Spencer map sends a log derivation to its balanced symmetric matrix. Its source and target are finite projective of the same fibre rank. Injectivity on the fibres of the rigid charts gives the isomorphism by the local criterion for an isomorphism of vector bundles.
\leanref{GlobalTheory.\allowbreak{}UPoincTop.\allowbreak{}uPoinStX\_\allowbreak{}sep}{Development.\allowbreak{}AbelianVarietyTheory.\allowbreak{}Proofs.\allowbreak{}DivisorProofs.\allowbreak{}UnitaryToroidalSeparationOfBase}{23--88}
\leanref{GlobalTheory.\allowbreak{}BlowAmpleNoeth.\allowbreak{}ample\_\allowbreak{}twist\_\allowbreak{}of\_\allowbreak{}blowup}{Development.\allowbreak{}AbelianVarietyTheory.\allowbreak{}Proofs.\allowbreak{}DivisorProofs.\allowbreak{}AmpleTwistNormalisedBlowupNoetherianDomain}{21--67}
\leanref{GlobalTheory.\allowbreak{}UChartJoin.\allowbreak{}rigidForFan\_\allowbreak{}of\_\allowbreak{}uStB}{Development.\allowbreak{}AbelianVarietyTheory.\allowbreak{}Proofs.\allowbreak{}DivisorProofs.\allowbreak{}RigidBoundaryChartsForFanOfUnitarySeparation}{21--87}
\leanref{GlobalTheory.\allowbreak{}LogKSAssembly.\allowbreak{}hasLogKS}{Development.\allowbreak{}AbelianVarietyTheory.\allowbreak{}Proofs.\allowbreak{}DivisorProofs.\allowbreak{}LogKodairaSpencerIsoSmoothSncToroidalModel}{57--69}
\leanref{GlobalTheory.\allowbreak{}RigChart.\allowbreak{}toroidalForFanKS\_\allowbreak{}of\_\allowbreak{}rigidForFan}{Development.\allowbreak{}AbelianVarietyTheory.\allowbreak{}Proofs.\allowbreak{}DivisorProofs.\allowbreak{}ToroidalFanKodairaSpencerOfRigidBoundaryCharts}{16--20}

Let \(C\) be a smooth projective geometrically connected curve over a field of characteristic zero, and let \(G\to C\) be semiabelian with identity section \(e\). Bost's seminegativity consequence of \cite[Th.~3.2]{bost-2006} applies: the group scheme is smooth and its generic fibre is an extension of an abelian variety by a torus. Such a fibre has no nontrivial connected unipotent subgroup. Thus every subbundle of \(\operatorname{Lie}(G/C)\) has degree at most zero. For any locally free subsheaf \(W\) take its saturation \(W'\). On the smooth curve \(W'\) is a subbundle and \(\deg W=\deg W'-\deg(W'/W)\le0\). This also gives nefness of the Hodge bundle \(e^*\Omega^1_{G/C}\). On pulling back to a smooth projective curve over an algebraically closed extension of the ground field, the group scheme remains semiabelian and formation of its Lie bundle commutes with pullback. Dualising any quotient line of the pulled-back Hodge bundle gives a line subbundle of the pulled-back Lie bundle. Its degree is nonpositive, as required.
\leanref{GlobalTheory.\allowbreak{}Nef.\allowbreak{}linesLE\_\allowbreak{}zero\_\allowbreak{}of\_\allowbreak{}sub\_\allowbreak{}lie}{Development.\allowbreak{}AbelianVarietyTheory.\allowbreak{}Proofs.\allowbreak{}DivisorProofs.\allowbreak{}LieSubsheafLinesDegreeNonpos}{18--69}
\leanref{GlobalTheory.\allowbreak{}Nef.\allowbreak{}semiNeg\_\allowbreak{}lie\_\allowbreak{}of\_\allowbreak{}isAlgClosed}{Development.\allowbreak{}AbelianVarietyTheory.\allowbreak{}Proofs.\allowbreak{}DivisorProofs.\allowbreak{}LieSheafSemiAbelianSemiNegOfIsAlgClosed}{27--34}
\leanref{GlobalTheory.\allowbreak{}Nef.\allowbreak{}semiAbelianLieSeminegative}{Development.\allowbreak{}AbelianVarietyTheory.\allowbreak{}Proofs.\allowbreak{}DivisorProofs.\allowbreak{}SemiAbelianLieSeminegativeBis}{22--29}

For point separation let \((X,\Delta)\) be an integral normal projective klt pair over \(\C\), of dimension \(n\), with \(\Delta\) effective. Let \(G=m(K_X+\Delta)\) be Cartier with \(m\ge1\). Take a line bundle \(M\) and a nef and big Cartier divisor \(N_0\) with \(M\equiv G/m+N_0/r\), \(r\ge1\) integral. For distinct closed points \(x_1,x_2\), assume there are \(c(e)>0\), \(1\le e\le n\), with
\begin{equation}\label{eq:p8c3-separation}
 \frac{(N_0^e\cdot Z)}{r^e}>c(e)^e
 \quad\text{for every irreducible \(e\)-fold \(Z\) through either point},
 \qquad
 \sum_{e=1}^n \frac{2^{1/e}e}{c(e)}\le1 .
\end{equation}
Set \(\rho=1/r\). By integrality, \((N_0^e\cdot Z)\ge D_e=\lfloor r^ec(e)^e\rfloor+1\). For each \(e\) choose a rational \(h_e\) strictly between \(2^{1/e}e/D_e^{1/e}\) and \(2^{1/e}e/(rc(e))\). Then
\begin{equation}\label{eq:p8c3-separation-budget}
 \sum_e h_e<\rho,\qquad h_e^e(N_0^e\cdot Z)>2e^e .
\end{equation}
Write \(aN_0=A+E_1\) with \(A\) ample and \(E_1\) effective Cartier, by bigness. On a log resolution of the pair and \(E_1\), all log discrepancies stay positive after adding \(E_1/m'\) for sufficiently large \(m'\). By \eqref{eq:p8c3-separation-budget}, choose such a multiple \(m'\) of \(r\) with \(q=m'/r>a\) and \(\sum_e h_e<(q-a)/m'\). The ample divisor \(A'=A+(q-a)N_0\) gives
\begin{equation}\label{eq:p8c3-separation-ample}
 M\equiv K_X+\Delta+E_1/m'+A'/m',\qquad
 (A'^e\cdot Z)\ge(q-a)^e(N_0^e\cdot Z).
\end{equation}
The mixed terms of the intersection are nonnegative by nefness. Apply the ample two-point theorem \cite[Th.~1.2]{fujino-2010} to the log canonical pair \((X,\Delta+E_1/m')\) and ample part \(A'/m'\), with thresholds \(c'(e)=2^{1/e}e(q-a)/(m'h_e)\). Equations \eqref{eq:p8c3-separation-budget} and \eqref{eq:p8c3-separation-ample} give \(((A'/m')^e\cdot Z)>c'(e)^e\). The weighted sum of the reciprocals is \(m'\sum_e h_e/(q-a)<1\). The theorem gives a section of \(M\) nonzero at exactly one of the two points. The reduction to the ample case gives point separation for the nef and big adjoint part in \eqref{eq:p8c3-separation}.
\leanref{GlobalTheory.\allowbreak{}pointSeparationPair\_\allowbreak{}of\_\allowbreak{}vanishing}{Development.\allowbreak{}AbelianVarietyTheory.\allowbreak{}Proofs.\allowbreak{}DivisorProofs.\allowbreak{}KollarTwoPointSeparationOfNadelVanishing}{15--39}
\leanref{NadelVanishingSeparation.\allowbreak{}NSI.\allowbreak{}pointSeparationPair}{Development.\allowbreak{}WeilBound.\allowbreak{}CoreStatements.\allowbreak{}AutomorphicKernels.\allowbreak{}NadelVanishing.\allowbreak{}TopPointSep}{6--8}
\leanref{GlobalTheory.\allowbreak{}Published.\allowbreak{}PointSeparationPair}{Development.\allowbreak{}AbelianVarietyTheory.\allowbreak{}PublishedTheorems.\allowbreak{}Divisors₂.\allowbreak{}PointSeparationPair}{15--30}

The arithmetic degree assertion comprises preservation of self-intersection under birational pullback between generically smooth arithmetic varieties and equality of the complex top Chern integral of a smooth weight with the divisor's top intersection. The first follows by projection for cycle heights, the second by the complex Chern degree formula \cite{bost-gillet-soule-1994}. Grothendieck deformation uses the de Rham crystal over the nilpotent divided-power thickening and the lifting criterion for homomorphisms preserving the Hodge filtration.
\leanref{UpperTheory.\allowbreak{}LiteratureTheorems.\allowbreak{}ar\_\allowbreak{}degree\_\allowbreak{}of\_\allowbreak{}literature}{Development.\allowbreak{}AutomorphicArithmetic.\allowbreak{}ThetaCorrespondence.\allowbreak{}ThetaPairLifts.\allowbreak{}ThirdPlaceSeries.\allowbreak{}EveryCharacterOfLiteratureOfVWin}{251--252}
\leanref{GlobalTheory.\allowbreak{}Published.\allowbreak{}grothendieckDeformation\_\allowbreak{}of\_\allowbreak{}lifting}{Development.\allowbreak{}AbelianVarietyTheory.\allowbreak{}Proofs.\allowbreak{}DivisorProofs.\allowbreak{}GrothendieckDeformationOfLifting}{9--13}

Equations \eqref{eq:p8c3-kleiman} and \eqref{eq:p8c3-hs-upper} give the arithmetic Kleiman inequality and the two upper bounds for small-section counts required in \ref{B8.25}. The lower bound for arithmetic volume is Theorem~\ref{B8.27}. In the order of Theorem~\ref{B8.25}, resolution, principalisation and invariant Poincar\'e--Lelong supply the first three premises. The fourth is the conjunction of the three arithmetic inequalities just identified, and the fifth is Theorem~\ref{B8.27}. Hodge--Tate exactness supplies the sixth. The principally polarised retraction constructed using \eqref{eq:p8c3-zarhin} supplies the seventh, and \eqref{eq:p8c3-nef-rr} supplies the eighth. The integral Siegel compactification and the unitary compactification for the prescribed fan supply the ninth and tenth, respectively. Lie seminegativity, two-point separation and the arithmetic degree formula are the eleventh through thirteenth. The fourteenth is the conjunction in Definition~\ref{B8.25l}. The fifteenth uses de Rham comparison and lifting of abelian schemes and their homomorphisms across nilpotent divided-power thickenings.
\end{proof}

\subsection*{Proof of \ref{B8.25b}}\label{prf:B8.25b}\phantomsection
\begin{proof}
Put \(L=L_x\) and use \eqref{eq:p8c3-anchor}. Increase the threshold so that \(x\ge6\) and \((x^A)^2\le b\). The host and weight estimates give
\begin{equation}\label{eq:p8c3-host-size}
 \begin{gathered}
 f\ge8,\qquad (frL)^{12}\le P_*\le x^A,\qquad
 frn_VL\le x^A,\qquad 20fr\le b,\\
 \kappa\le4frb,\qquad \sum_\rho w_{j,\rho}\le\kappa+fr,\qquad b\le\kappa+r .
 \end{gathered}
\end{equation}
For the product estimate use the host bounds \(fr\le x^{20}\), \(n_V\le2t\le x^{1+C_t}\), and \(L\le x\). Use the absolute row sums and allow a correction of at most \(r\) at each embedding. The layer weight has size at least \(b\). Then \(b\le2\kappa\) and \(\sum_\rho w_{j,\rho}\le3\kappa\). By \eqref{eq:p8c3-anchor}, \(|M_{p,i}|\le9rL\kappa\le P_*\kappa\). Also \(9rL\kappa\le36fr^2Lb\le x^Bb\) by \eqref{eq:p8c3-host-size}. The bounds on \(E_p\) follow from its expression in \eqref{eq:p8c3-anchor} and \(fmL\le frn_VL\).
\leanref{GlobalTheory.\allowbreak{}CovTop.\allowbreak{}Sizes.\allowbreak{}hostN\_\allowbreak{}alt}{Development.\allowbreak{}AutomorphicArithmetic.\allowbreak{}ThetaCorrespondence.\allowbreak{}ThetaPairLifts.\allowbreak{}ThirdPlaceSeries.\allowbreak{}HostN}{13--38}
\leanref{GlobalTheory.\allowbreak{}CovTop.\allowbreak{}WtShape.\allowbreak{}wtNum\_\allowbreak{}at}{Development.\allowbreak{}AbelianVarietyTheory.\allowbreak{}Proofs.\allowbreak{}DivisorProofs.\allowbreak{}WeightBoundHolds}{45--186}
\leanref{GlobalTheory.\allowbreak{}CovTop.\allowbreak{}SizeRows.\allowbreak{}sum\_\allowbreak{}e\_\allowbreak{}le}{Development.\allowbreak{}AbelianVarietyTheory.\allowbreak{}Proofs.\allowbreak{}DivisorProofs.\allowbreak{}SizeBoundStatement}{196--222}

At a nonterminal layer the flag dimension is \(d_*=f\binom r2\). Use \(T=T_*=\lceil P_*^2d_*^2 k/m\rceil_{\N}\) as the weight-division length. The family parameters give the gap parameter \(h_0\in\Z\), and width parameters \(S,\eta\ge0\) with \(P_*(S+\eta)\le100w''\), where \(w''\) is the scale bound of the window. The spacing bounds are
\begin{equation}\label{eq:p8c3-spacing}
 T\ge4d_*+4,\quad T\ge P_*^2d_*^2,\quad
 h_0\ge4T^2,\quad h_0\ge P_*Tr f^2m^2kd_*,\quad P_*\ge6400r^2 .
\end{equation}
Sum the floored adjacent differences in \eqref{eq:p8c3-quotient} from the zero anchor at \(k-2\) to obtain \(u_1\). The weight-gap conditions give strict decrease in each block and an unshifted wall gap at least \(h_0/(2T)\). Subtracting \(3r\eta_H\) from \(u_1\) decreases each cross-wall difference by \(3r\) and leaves it positive. Its relative loss is at most \(6rT/h_0\). Within-block differences stay fixed. With \(V\) the product of differences over pairs of columns and embeddings, let \(\rho=1-V(u')/V(u_1)\). There are at most \(fr^2\) cross-wall factors, so the bound on the loss of a product gives
\begin{equation}\label{eq:p8c3-rho}
 0\le\rho<1,\qquad \rho h_0\le6fr^3T .
\end{equation}
Together \eqref{eq:p8c3-spacing} and \eqref{eq:p8c3-rho} give \(d_*\rho\le6/6400\), \(d_*^2/T\le1/100\) and \(T\ge3d_*+5\). Thus the parameter \(z=7(d_*^2/T+d_*\rho)\) of \eqref{eq:p8c3-z} is at most \(7(1/100+6/6400)<1/2\).
\end{proof}

\subsection*{Proof of \ref{B8.25c}}\label{prf:B8.25c}\phantomsection
\begin{proof}
In the quadratic host \(K=\Q(\sqrt{-d})\) of discriminant \(-d\), the splitting tests at \(3,5\) give \(s^2\equiv-d\bmod15\) with \((2s,15)=1\). Thus \(J=(15,s+\sqrt{-d})\) satisfies \(J+\bar J=(1)\). The generators of the product \(J\bar J\) are divisible by \(15\), and generate \(15^2\) and \(30s\), so \(J\bar J=(15)\). In particular \(J^e\) and \(\bar J^e\) are complementary halves of \((15^e)\), with generators in \(\Z[w_0]\) for \(w_0=(1+\sqrt{-d})/2\). They define the cover at \(3,5\). Retain the base level \(g_0g_1\) away from \(3,5\). The auxiliary level parameter is one, so the full covering level is \(M'\).

Let \(N=M'\) and \(R\) be the ray class field of \(K\) modulo \(N\mathcal O_K\). It contains the needed roots of unity and Weber values. Conjugation preserves the modulus. The ray class sequence bounds \([R:K]\) by \(h_K\#(\mathcal O_K/N\mathcal O_K)^\times\), where \(h_K\) is the class number. The quadratic class-number bound \(h_K\le d\) and the conductor-discriminant formula give
\begin{equation}\label{eq:p8c3-ray}
 [R:K]\le h_KN^2\le dN^2,\qquad
 \log\operatorname{rd}(R)\le\tfrac12\log d+\log N .
\end{equation}
Indeed every character conductor divides \(N\mathcal O_K\) of norm \(N^2\), and the discriminant formula divides its logarithmic contribution by \(2[R:K]\). Take \(F'=FR\) for the original coefficient field \(F\). The Weber values give relative algebraic closure of \(F'\) in the component function fields. Fix a lattice embedding for uniformisation. At each good embedding of \(F'\) take an orthogonal frame compatible with that same lattice embedding. These frames give the unit minors and models required for goodness.
\leanref{GlobalTheory.\allowbreak{}CovTop.\allowbreak{}Inhab.\allowbreak{}exists\_\allowbreak{}half}{Development.\allowbreak{}AbelianVarietyTheory.\allowbreak{}Definitions.\allowbreak{}Divisors.\allowbreak{}NormNegCmConstantEqAtAllEmbeddings}{324--385}
\leanref{GlobalTheory.\allowbreak{}CoverHalf.\allowbreak{}exists\_\allowbreak{}coverLevel}{Development.\allowbreak{}AbelianVarietyTheory.\allowbreak{}Proofs.\allowbreak{}FoundationProofs.\allowbreak{}ExistsHalfOfBasePins}{144--173}
\leanref{GlobalTheory.\allowbreak{}CovTop.\allowbreak{}SizedCover.\allowbreak{}exists\_\allowbreak{}ray}{Development.\allowbreak{}AbelianVarietyTheory.\allowbreak{}Proofs.\allowbreak{}DivisorProofs.\allowbreak{}ExistsGoodCoverOverSizedField}{28--108}
\leanref{GlobalTheory.\allowbreak{}CovTop.\allowbreak{}SizedCover.\allowbreak{}exists\_\allowbreak{}good\_\allowbreak{}sized}{Development.\allowbreak{}AbelianVarietyTheory.\allowbreak{}Proofs.\allowbreak{}DivisorProofs.\allowbreak{}ExistsGoodCoverOverSizedField}{111--161}

For a number field \(F_0\) write \(\operatorname{sz}(F_0)=\log[F_0:\Q]+\log\operatorname{rd}(F_0)\). The tied field and host estimates give
\begin{equation}\label{eq:p8c3-original-field}
 \operatorname{sz}(F)\le217(frL_x)^2,\qquad
 \log g_a\le9L_x+1\quad(a=0,1),\qquad \log d\le L_x .
\end{equation}
The degree of a compositum is at most the product of degrees and its root discriminant is at most the product of root discriminants, by the relative different bound under base change. Thus \eqref{eq:p8c3-ray} bounds \(\operatorname{sz}(F')-\operatorname{sz}(F)\) by \(20(\log N+\log d+1)\). The definition of \(e\) gives \(e\log15\le9L_x\), so by \eqref{eq:p8c3-original-field} the increase is at most \(20(28L_x+3)\le174(frL_x)^2\), since \(f\ge8\) and \(L_x\ge10\).
\leanref{GlobalTheory.\allowbreak{}CovTop.\allowbreak{}RecTied.\allowbreak{}core\_\allowbreak{}size\_\allowbreak{}le}{Development.\allowbreak{}AbelianVarietyTheory.\allowbreak{}Definitions.\allowbreak{}Divisors.\allowbreak{}CovolXOfCovol}{192--194}
\leanref{GlobalTheory.\allowbreak{}CovTop.\allowbreak{}DeepExist.\allowbreak{}size\_\allowbreak{}le}{Development.\allowbreak{}AbelianVarietyTheory.\allowbreak{}Proofs.\allowbreak{}DivisorProofs.\allowbreak{}DeepExistenceStatement}{29--57}
\end{proof}

\subsection*{Proof of \ref{B8.25d}}\label{prf:B8.25d}\phantomsection
\begin{proof}
Primitive uniformisation and the core lattice regularity give the core Gram frame and full cusp lattice. Pull back the Hilbert inner product by the row realisation. The ambient realised lattice is the integral span of the semilinear image of the core lattice. The primitive row laws give the good lattice of dimension \(d_W(i)=D_{\rm norm}>0\).
\leanref{GlobalTheory.\allowbreak{}CovTop.\allowbreak{}GramPt.\allowbreak{}hasGram\_\allowbreak{}of\_\allowbreak{}unif}{Development.\allowbreak{}AbelianVarietyTheory.\allowbreak{}Proofs.\allowbreak{}DivisorProofs.\allowbreak{}RowNormHasGram}{25--37}
\leanref{GlobalTheory.\allowbreak{}CovTop.\allowbreak{}latAll\_\allowbreak{}fg\_\allowbreak{}at\_\allowbreak{}alt}{Development.\allowbreak{}AutomorphicArithmetic.\allowbreak{}ThetaCorrespondence.\allowbreak{}ThetaPairLifts.\allowbreak{}ThirdPlaceSeries.\allowbreak{}LatAllFgAt}{15--30}
\leanref{GlobalTheory.\allowbreak{}CovTop.\allowbreak{}good\_\allowbreak{}lattice\_\allowbreak{}ne\_\allowbreak{}bot\_\allowbreak{}at\_\allowbreak{}point}{Development.\allowbreak{}AutomorphicArithmetic.\allowbreak{}ThetaCorrespondence.\allowbreak{}ThetaPairLifts.\allowbreak{}ThirdPlaceSeries.\allowbreak{}LatGoodNeBotAt}{12--24}

The squared-norm ratio is \(c^2=V\prod_\tau d_\tau\), where \(V\) is the row fundamental-domain volume and \(d_\tau\) is the dimension of the lowest compact type. Put \(L'=2L_x\). Tamagawa measure comparison bounds \(V\) by \(n_\Gamma\operatorname{vol}(U(r))^f\operatorname{mass}_I\), where \(n_\Gamma\) is the relative index of the row level in the full level and \(\operatorname{mass}_I\) is the remaining adelic mass. Indeed, the stabiliser with archimedean projection in a box of positive finite volume lies in the full-level set over that box. Its measure is at most \(n_\Gamma\) times the box volume in the row normalisation. Cancel the box volume in the Tamagawa comparison. The bounds for the index, compact volume factor and mass have logarithms at most \(60fr^2L'\), \(3fr^2L'\) and \(54fr^2L'\), respectively. This gives the volume bound in \eqref{eq:p8c3-meas-upper}.
\leanref{MainStatements.\allowbreak{}Tam.\allowbreak{}Plug.\allowbreak{}tamagawa\_\allowbreak{}pairs\_\allowbreak{}compatible}{Development.\allowbreak{}AutomorphicArithmetic.\allowbreak{}ThetaCorrespondence.\allowbreak{}ThetaPairLifts.\allowbreak{}HodgeBundleWeights.\allowbreak{}PlugPairs}{12--12}
\leanref{MainStatements.\allowbreak{}Tam.\allowbreak{}A.\allowbreak{}up\_\allowbreak{}core}{Development.\allowbreak{}WeilBound.\allowbreak{}RepresentationTheoryGL.\allowbreak{}TubeMeasures.\allowbreak{}DisjointTubes.\allowbreak{}ExponentialThreeBound.\allowbreak{}LeExpThree}{91--179}

For each \(\tau\), the Weyl dimension is the product over pairs \(a<b\) within the two compact blocks. If \(\kappa_a,\kappa_b\) are the corresponding entries of the lowest compact weight, each factor has absolute value at most \(1+|\kappa_a-\kappa_b|\), since \(b-a\ge1\). The row size estimate bounds the difference by \(e^{2L'}\). Since \(e^{L'}\ge2\), each factor contributes at most \(\log(1+e^{2L'})\le3L'\) to the logarithm. The two blocks supply at most \(r^2/2\) factors, so \(\log d_\tau\le3r^2L'/2\). Sum over the \(f\) embeddings to obtain the second bound:
\begin{equation}\label{eq:p8c3-meas-upper}
 \log V\le fr^2\log(r+3)+119fr^2L',\qquad
 \sum_\tau\log d_\tau\le\tfrac32fr^2L'.
\end{equation}
\leanref{GlobalTheory.\allowbreak{}OneMeasDim.\allowbreak{}log\_\allowbreak{}dimAt\_\allowbreak{}le}{Development.\allowbreak{}AbelianVarietyTheory.\allowbreak{}Proofs.\allowbreak{}DivisorProofs.\allowbreak{}LogLowestKTypeWeylDimensionLeOfSizeBounds}{77--143}
At the family threshold \(\log(r+3)\le L'\), this gives \(\log c^2\le122fr^2L'\). For the lower bound integrate an archimedean box over the level subgroup inside the stabiliser. Cancelling the box volume gives \(V\ge\operatorname{vol}(U(r))^f\operatorname{mass}_I\). The compact volume and mass bounds in the given measure normalisation yield an absolute \(C\ge0\) with
\begin{equation}\label{eq:p8c3-meas-lower}
 \log V\ge-C\{fr^2(\log(fr+3)+1)+\log|\operatorname{disc}E|\},
 \qquad \sum_\tau\log d_\tau\ge-fr^2\log(r+1).
\end{equation}
By \eqref{eq:p8c3-host-size} and the threshold in \ref{B8.25b}, \(r^2\le b\) and eventually \(\log(fr+3)\le2L_x\). Equation \eqref{eq:p8c3-anchor} gives \(\log|\operatorname{disc}E|\le4fL_x\). Substitute in \eqref{eq:p8c3-meas-upper} and \eqref{eq:p8c3-meas-lower} and use the formula for \(E_p\) in \eqref{eq:p8c3-anchor}. This gives \(|\log c|\le C_{\rm meas}E_p\), for example with \(C_{\rm meas}=(244+6C+3)/2\).
\leanref{GlobalTheory.\allowbreak{}OneMeasureRatio.\allowbreak{}log\_\allowbreak{}covolume\_\allowbreak{}ratio\_\allowbreak{}le}{Development.\allowbreak{}AutomorphicArithmetic.\allowbreak{}ThetaCorrespondence.\allowbreak{}ThetaPairLifts.\allowbreak{}HodgeBundleWeights.\allowbreak{}LogMeasureRatioBounded}{16--37}
\leanref{MainStatements.\allowbreak{}Tam.\allowbreak{}low}{Development.\allowbreak{}WeilBound.\allowbreak{}RepresentationTheoryGL.\allowbreak{}TubeMeasures.\allowbreak{}DisjointTubes.\allowbreak{}BoxExistence.\allowbreak{}ExistsBox}{64--127}
\leanref{GlobalTheory.\allowbreak{}CovTop.\allowbreak{}MeasUp.\allowbreak{}ratAt\_\allowbreak{}pos\_\allowbreak{}le}{Development.\allowbreak{}AbelianVarietyTheory.\allowbreak{}Proofs.\allowbreak{}DivisorProofs.\allowbreak{}LogRatioUpperBound}{166--216}
\leanref{GlobalTheory.\allowbreak{}CovTop.\allowbreak{}MeasLow.\allowbreak{}nef\_\allowbreak{}share\_\allowbreak{}lower\_\allowbreak{}bound\_\allowbreak{}alt}{Development.\allowbreak{}AutomorphicArithmetic.\allowbreak{}ThetaCorrespondence.\allowbreak{}ThetaPairLifts.\allowbreak{}ThirdPlaceSeries.\allowbreak{}CovolumeLowerBoundOfCover}{21--54}

In the Petersson estimate let \(s\) be the sum of the entries of the two algebraic highest-weight partitions over all embeddings, \(t=frn_V\) the column weight and \(q=Q\) the negative signature index. In the Petersson comparison, write \(\lambda_0\) for the mean logarithmic loss and \(s_0\) for the log supremum constant. Use the weight comparison \(s+t+1\le10frb\) and \eqref{eq:p8c3-host-size}, so \(\log(s+t+1)\le4L_x\) at the family threshold. The mean and supremum bounds are
\begin{equation}\label{eq:p8c3-petersson}
 \begin{split}
 \lambda_0&\le\bigl((1+|\log\Omega|)/2+2L_x+
 (2q+1)(6f+48)2L_x\bigr)(s+t),\\
 s_0&\le500f^2r^2L_x .
 \end{split}
\end{equation}
For the core cusp lattice \(\Lambda_{\rm core}\) over \(F\), the maximal slope bound is
\begin{equation}\label{eq:p8c3-core-max}
 \mu_{\max}(\Lambda_{\rm core})
 \le2\lambda_0+2s_0+\log\operatorname{rd}(F)\le x^{2A+4}b .
\end{equation}
In the last inequality, \(q\le r\) and \(f\ge8\) bound \((2q+1)(6f+48)\) by \(36fr\). With \eqref{eq:p8c3-anchor} and \eqref{eq:p8c3-host-size}, the first line of \eqref{eq:p8c3-petersson} is at most \(760x^{2A}b\). The second line and \eqref{eq:p8c3-original-field} bound the remaining terms in \eqref{eq:p8c3-core-max} by \(1217x^{2A}b\). Use \(x\ge L_x\ge10\).

Choose a sublattice of \(\Lambda_{\rm all}\) attaining its maximal slope and extend it to the normal closure of the realisation field over \(F\). With the extended realised norms, its normalised slope is preserved. Rescale those norms at each embedding by \(c^{-1}\), with \(c\) taken at the corresponding embedding of the realisation field. The rescaled norms on the extended core space are the scalar extensions of the Petersson norms. For a rank-\(a\) sublattice the rescaling changes the normalised logarithmic squared covolume by at most \(2aC_{\rm meas}E_p\). Its slope therefore changes by at most \(2C_{\rm meas}E_p\). The extended sublattice lies in the integral span of the extended core lattice. For the rescaled norms, Galois descent of the maximal slope bounds its slope by \(\mu_{\max}(\Lambda_{\rm core})\). Applying this to the chosen maximal-slope sublattice and using \eqref{eq:p8c3-size-main} and \eqref{eq:p8c3-core-max} gives
\begin{equation}\label{eq:p8c3-realised-max}
 \mu_{\max}(\Lambda_{\rm all})\le x^{2A+4}b+2C_{\rm meas}x^Bb .
\end{equation}
Increase the family exponent and threshold to absorb the sum.
\leanref{UpperTheory.\allowbreak{}Lit.\allowbreak{}Tam.\allowbreak{}slopeMax\_\allowbreak{}sq\_\allowbreak{}latW\_\allowbreak{}le\_\allowbreak{}pow\_\allowbreak{}mul}{Development.\allowbreak{}AutomorphicArithmetic.\allowbreak{}ThetaCorrespondence.\allowbreak{}FockModels.\allowbreak{}TamagawaPeterssonLambdaSqNonneg}{596--671}
\leanref{GlobalTheory.\allowbreak{}CovTop.\allowbreak{}LamCore.\allowbreak{}collect}{Development.\allowbreak{}AbelianVarietyTheory.\allowbreak{}Proofs.\allowbreak{}DivisorProofs.\allowbreak{}SlopeMaxSquaredBound}{49--71}
\leanref{GlobalTheory.\allowbreak{}SlopeMaxExt.\allowbreak{}slopeMax\_\allowbreak{}sq\_\allowbreak{}span\_\allowbreak{}le}{Development.\allowbreak{}AbelianVarietyTheory.\allowbreak{}Proofs.\allowbreak{}DivisorProofs.\allowbreak{}MaxSlopeSqFieldExtensionLeComparableNorms}{142--205}
\leanref{GlobalTheory.\allowbreak{}CovTop.\allowbreak{}LamAll.\allowbreak{}slopeMax\_\allowbreak{}all\_\allowbreak{}le}{Development.\allowbreak{}AbelianVarietyTheory.\allowbreak{}Proofs.\allowbreak{}DivisorProofs.\allowbreak{}LambdaBoundFromMeasure}{42--66}
\end{proof}

\subsection*{Proof of \ref{B8.26p}}\label{prf:B8.26p}\phantomsection
\begin{proof}
These two inequalities adapt Zhang's successive-minima intersection bounds \cite[Th.~5.2]{zhang-1995} to birational multihomogeneous models with cone metrics. For (1), subtract \(a\) from the logarithmic metric weight of \(A\). The curvature form and the geometric fibre degrees are unchanged, while on each horizontal curve \(W\) the arithmetic degree decreases by \(a\deg_{\C}W\) and remains nonnegative. By multilinearity and the constant-shift recursion, the top self-intersection decreases by \((D_A+1)a d_A\). Passing to the projective birational modification of Lemma~\ref{B8.25h} via characteristic-zero resolution produces a projective flat integral model with smooth generic fibre on which the pulled-back shifted metric is smooth and semipositive while fibrewise nefness, nonnegativity on horizontal curves and top self-intersection are preserved; arithmetic Kleiman \eqref{eq:p8c3-kleiman} gives nonnegativity of the shifted self-intersection, proving \eqref{eq:p8c3-constant-shift}.

For (2), when \(D_A=0\) the generic fibre \(X_{\Q}\) is a single horizontal orbit, so Zariski density puts \(X\) in the collection and \(I=h_A(X)d_A\le a d_A\). When \(D_A\ge1\), generation by coordinate monomials of nonnegative multidegree makes \(A\) nef on all geometric fibres and gives \(d_A\ge0\). Suppose \(I>(D_A+1)a d_A\), and choose \(a'>a\) and \(\varepsilon>0\) with \(I>(D_A+1)(a'+\varepsilon)d_A\). Subtracting \(a'+\varepsilon\) from the logarithmic metric weight leaves positive top self-intersection while preserving smoothness, conjugation invariance, fibrewise nefness and nonnegativity of the Levi form along \(P_0\subseteq P\). Applying Theorem~\ref{B8.27} to the shifted line produces an integer \(r\ge1\) with at least two sections in \(S_r\) of shifted sup norm at most \(1\), hence a nonzero global section \(s\in S_r=H^0(X,A^r)\) satisfying \(\|s\|_{\sup}\le e^{-r(a'+\varepsilon)}<e^{-ra'}\) in the original metric. Since the curves \(W\) with \(h_A(W)\le a\) are Zariski dense in \(X_{\Q}\), one such curve is not contained in the zero locus of \(s\). Restricting \(s\) to the normalisation of \(W\), effective finite intersection multiplicities and the archimedean bound \(-\log\|s\|>ra'\) give \(r\,II(W;A)>ra'\deg_{\C}W\), contradicting \(h_A(W)\le a<a'\).
\end{proof}

\subsection*{Proof of \ref{B8.26q}}\label{prf:B8.26q}\phantomsection
\begin{proof}
For (1), resolution and principalisation give embedded characteristic-zero resolutions with coordinate projections, so archimedean potential integrals converge and are independent of the chart. Given a horizontal prime ideal \(P\) of arithmetic dimension \(n+2\), smooth Hermitian lines \(L_1,\ldots,L_n,A,B\) with \(B\) of multidegree zero, and a multihomogeneous section \(s\notin P\) of \(A\), the set \(\mathcal H\) of horizontal prime components \(W\) of \(\operatorname{div}_P(s)\) is finite, and each \(W\in\mathcal H\) has arithmetic dimension \(n+1\). Choose a rational prime \(p\) such that \(-p\) has no square root in the function field of \(P\) or of any \(W\in\mathcal H\). Adjoining a \(\mathbb P^1\)-block cut out by the imaginary quadratic relation \(X_0^2+pX_1^2=0\) produces a horizontal prime ideal \(P_p\) of arithmetic dimension \(n+2\) whose two complex sheets \(X_0/X_1=\pm i\sqrt{p}\) are interchanged by complex conjugation, so each smooth Hermitian line glues with its complex conjugate across the two sheets to a smooth conjugation-invariant Hermitian line. On \(P_p\) and on the quadratic twist \(W_p\) of each \(W\in\mathcal H\), the archimedean potential integrals double, and \(W\mapsto W_p\) is a multiplicity-preserving bijection from \(\mathcal H\) onto the horizontal divisor support of \(s\) on \(P_p\). Applying conjugation-invariant Poincar\'e--Lelong on \(P_p\) and dividing by \(2\) gives unrestricted Poincar\'e--Lelong, from which divisor induction constructs \(II\).

For (2), cf.\ \cite[Prop.~3.3.1.5]{lan-2013}, smoothness of \(G\to S\) makes \(G\) normal and \(G_U\to G\) scheme-theoretically dominant. Let \(\Gamma\subset G\times_S G'\) be the schematic closure of the graph of \(f\), with projections \(p:\Gamma\to G\) and \(q:\Gamma\to G'\). Any point of \(G\setminus G_U\) is the closed-point image of a discrete valuation trait in \(G\) with generic point in \(G_U\), and any point of \(\Gamma\setminus G_U\) is the closed-point image of a discrete valuation trait in \(\Gamma\) with generic point in \(G_U\). Unique extension of semiabelian homomorphisms over discrete valuation rings \cite[Ch.~I, Prop.~2.7]{faltings-chai-1990} lifts traits from \(G\) to \(\Gamma\), so \(p\) is surjective, and for any integer \(n\ge1\) and geometric point \(x\) of \(\Gamma\) with \(p(x)^n=1\), applying it to a trait in \(\Gamma\) through \(x\) when \(p(x)\notin G_U\) (and using the homomorphism \(f\) directly when \(p(x)\in G_U\)) shows that \(q(x)^n=1\). Since the \(n\)-torsion on each geometric fibre of \(G'\) is finite, every point of \(\Gamma\) lying over a torsion point of a fibre of \(G\) lies in the quasi-finite locus of \(p\), on which \(p\) is an open immersion because \(G\) is normal and \(G_U\hookrightarrow\Gamma\) is dominant. Given \(w_1,w_2\in\Gamma\) with \(p(w_1)=p(w_2)=y\), if \(y\in G_U\) then \(w_1=w_2\) because \(p^{-1}(G_U)=\Gamma_f\simeq G_U\); otherwise \(w_1,w_2\in\Gamma\setminus G_U\), so realise both points via trait extensions and amalgamate their residue fields over \(\kappa(y)\) into an algebraically closed field so that both closed points induce the same geometric point \(\bar y\) in a common geometric fibre over \(S\). For every torsion point on that geometric fibre, its two graph lifts factor through the quasi-finite open immersion and agree after composition with \(p\), hence coincide; since torsion is schematically dense on semiabelian fibres, the two pulled-back homomorphisms coincide, so their values at \(\bar y\) agree and \(w_1=w_2\). Thus \(p\) is bijective and locally quasi-finite, hence an isomorphism, and \(F=q\circ p^{-1}\) is the unique homomorphism extending \(f\).

For (3), the kernel duality is the Fargues exactness premise of Theorem~\ref{B8.25}. By Fargues's valuation theorem for the canonical filtration and the Hodge--Tate map \cite[Th.~4 and Cor.~8]{fargues-2011}, every \(\kappa\in\mathcal O_C\) with \(v(\kappa)^{p-1}\le v(p)\) annihilates \(\operatorname{coker}(\alpha_G)\); when \(p=2\), \(v(2)<1\) gives \(v(4)=v(2)^2\le v(2)\), so \(v(\kappa)\le v(4)\) implies \(v(\kappa)^{2-1}\le v(2)\).

For (4), a smooth morphism \(X\to\operatorname{Spec}\C\) partitions \(X\) into pairwise disjoint clopen subschemes \(W_D\) smooth of constant relative dimension \(D\). When \(X\) is proper, each connected component \(Y\) of \(W_D\) is smooth and proper over \(\C\); passing to a smooth projective birational model \(\pi:\widetilde Y\to Y\) via Chow's lemma \cite[Th.~5.6.1]{grothendieck-dieudonne-1961} and characteristic-zero resolution \cite{hironaka-1964}, compatible holomorphic local sections algebraize on \(\widetilde Y\) by projective degree-zero GAGA \cite{serre-1956}, descend along \(\pi_*\mathcal O_{\widetilde Y}=\mathcal O_Y\) by normality of \(Y\), and glue uniquely across \(\bigsqcup_D W_D\). For an \(R\)-Hopf order \(\mathcal A\subset C^N\) spanned over \(R\) by the columns \(b_1,\ldots,b_{\#N}\) of a matrix \(B\in\GL_{\#N}(C)\), its \(R\)-dual lattice \(\mathcal A^*=\{u\in C^N:\sum_{a\in N}u(a)f(a)\in R\ \forall f\in\mathcal A\}\) is spanned by the columns \(u_i\) of \((B^{-1})^{\mathsf t}\) and satisfies \(\mathcal A^{**}=\mathcal A\). Under the Fourier transform \(\widehat u(h)=\sum_{a\in N}\langle h,a\rangle u(a)\), \(\mathcal A^\vee=\widehat{\mathcal A^*}\) is spanned over \(R\) by the columns of \(M=T(B^{-1})^{\mathsf t}\), where \(T=(\langle h,a\rangle)_{h\in H,a\in N}\) is the character table. Point mass at \(0\) maps to \(1\), evaluation at \(0\) gives the counit, convolution in \(\mathcal A^*\) (closed by comultiplication of \(\mathcal A\)) maps to pointwise multiplication, negation gives the antipode, and for \(u\in\mathcal A^*\) the dual-basis expansion \(\langle h_2,a\rangle=\sum_i\widehat u_i(h_2)b_i(a)\) together with \(b_i u\in\mathcal A^*\) yields \(\widehat u(h_1+h_2)=\sum_i\widehat{b_i u}(h_1)\widehat u_i(h_2)\), giving comultiplication on \(\mathcal A^\vee\). Fourier duality gives \((\mathcal A^\vee)^\vee=\mathcal A\) because \(T^{\mathsf t}(M^{-1})^{\mathsf t}=T^{\mathsf t}(T^{\mathsf t})^{-1}B=B\). Since \(\det(M)=\det(T)\det(B)^{-1}\) and \(\det(T)^2=\pm(\#N)^{\#N}\), the determinant valuations \(c=v(\det B)\) and \(c'=v(\det M)\) satisfy \((cc')^2=v(\#N)^{\#N}\). For a subgroup \(K\subseteq N\), the image of \(\mathcal A\) under restriction \(C^N\to C^K\) is an \(R\)-Hopf order in \(C^K\) whose \(R\)-dual is the preimage of \(\mathcal A^*\) under extension by zero \(C^K\hookrightarrow C^N\), while \(\{\bar f\in C^{N/K}:\bar f\circ\pi\in\mathcal A\}\) along \(\pi:N\to N/K\) equals the \(R\)-dual of coset pushforward \(\pi_*(\mathcal A^*)\) and forms a quotient \(R\)-Hopf order.

For (5), relative ampleness of \(L\) makes \(A\) projective locally over \(S_0\), so every finite \(G\)-orbit lies in an affine open subscheme of \(A\), and passage to the quotient by the free finite locally free translation equivalence relation \(G\times_{S_0}A\hookrightarrow A\times_{S_0}A\) produces a finite, flat, surjective \(S_0\)-morphism \(q:A\to A'\) with \(G\times_{S_0}A\simeq A\times_{A'}A\) \cite[Th.~4.1]{gabriel-1970}, \cite[\S12, Th.~1]{mumford-1970}. The commutative group operations on \(A\) descend along \(q\), making \(q\) a group-scheme homomorphism whose pullback along \(1_{A'}\) identifies \(\ker(q)\) functorially with \(G\). Properness and geometric connectedness of \(A'\to S_0\) follow from finite surjectivity of \(q\), flatness of \(A'\to S_0\) from faithful flatness of \(q\), and smoothness on each geometric fibre \(A'_{\bar s}\) by faithfully flat descent of regularity from \(A_{\bar s}\). Because \(q_*\mathcal O_A\) is finite locally free and rank-one projective modules over semilocal rings are free, the norm \(L'=\operatorname{Nm}_q(L)\) is an invertible sheaf on \(A'\). On each geometric fibre \(A_{\bar s}\) where \(G_{\bar s}\) has rank \(r_G\ge1\), \(q_{\bar s}^*L'_{\bar s}\) is algebraically equivalent to \(L_{\bar s}^{\otimes r_G}\) by the theorem of the square, hence ample, so \(L'_{\bar s}\) is ample on \(A'_{\bar s}\) by finite surjective descent \cite[\S7]{mumford-1970}; since \(A'\to S_0\) is proper and of finite presentation, fibrewise ampleness then gives relative ampleness of \(L'\) on \(A'\to S_0\) by Noetherian limit reduction from \cite[III$_1$, Th.~4.7.1]{grothendieck-1961}. When \(R_0=R\) is the valuation ring of \((C,v)\), applying Zarhin's retraction \eqref{eq:p8c3-zarhin} to \((A',L')\) embeds \(A'\) as a direct-summand group-scheme retract of a principally polarised abelian scheme over \(\operatorname{Spec}R\).
\end{proof}

\subsection*{Proof of \ref{B8.26u}}\label{prf:B8.26u}\phantomsection
\begin{proof}[Proof (outline)]
The steps are the following.
\begin{enumerate}
\item By Lemmas~\ref{B8.26p} and \ref{B8.26q}, the first eight premises of \ref{B8.25} give the arithmetic intersection system \(II\), the arithmetic Hilbert--Samuel theorem and the bigness criterion on the model, resolution of singularities in characteristic zero, extension of homomorphisms of semiabelian schemes over normal bases, the Hodge--Tate exact sequence of Fargues, algebraisation of sections by GAGA, finite locally free quotients of polarised abelian schemes, and Zarhin's polarised retraction.
\item Let \(F_c\) be the coefficient field and let \(Z_0\) be an algebraic component of the cover. On its moduli open let \(\mathcal A\) be the universal abelian scheme. For \(\tau\) in the CM type put \(V_\tau=(\operatorname{Lie}\mathcal A)_\tau\), \(V_{\bar\tau}=(\operatorname{Lie}\mathcal A)_{\bar\tau}\), \(p=p_{\rm sig}\), \(q=q_{\rm sig}\), and
\[
 \mathcal L_\tau=\det V_\tau\otimes(\det V_{\bar\tau})^{-1}.
\]
By Lemma~\ref{B8.24b}, choose an \(F_c\)-rational section \(s_\tau\) generating \(\mathcal L_\tau\) on this open. For a field extension \(\iota:F_c\hookrightarrow F\), an \(F\)-point \(y\) on the open, and eigenframes \(v=(v_1,\ldots,v_p)\), \(w=(w_1,\ldots,w_q)\) at \(y\), define \(\varphi_F(y;v,w)\in F^\times\) by
\[
 s_\tau(y)=\varphi_F(y;v,w)
 (v_1\wedge\cdots\wedge v_p)\otimes
 (w_1\wedge\cdots\wedge w_q)^{-1}.
\]
If \(v'_i=\sum_j A_{ij}v_j\) and \(w'_i=\sum_j B_{ij}w_j\), with \(A\in\GL_p(F)\) and \(B\in\GL_q(F)\), then the determinant frame changes by \(\det A/\det B\), and hence
\[
 \varphi_F(y;v',w')\det A
 =\varphi_F(y;v,w)\det B.
\]
For a field embedding \(\eta:F\hookrightarrow F'\), formation of exterior products and evaluation of \(s_\tau\) commute with base change, so the coefficient becomes \(\eta(\varphi_F(y;v,w))\). At the generic point, with function field \(\Omega\), use the same symbols for the generic fibres of the eigenbundles and write \(s_\tau=a\otimes b^{-1}\) with nonzero \(a\in\bigwedge^p_\Omega V_\tau\) and \(b\in\bigwedge^q_\Omega V_{\bar\tau}\). For bases \((b_{1,i})_{i=1}^p\) of \(V_\tau\) and \((b_{2,i})_{i=1}^q\) of \(V_{\bar\tau}\), write \(a=\alpha\,b_{1,1}\wedge\cdots\wedge b_{1,p}\) and \(b=\kappa\,b_{2,1}\wedge\cdots\wedge b_{2,q}\), where \(\alpha,\kappa\in\Omega^\times\). The frame coefficient is \(\alpha/\kappa\); at a point where these frames and this quotient specialise, its specialisation is \(\varphi_F(y;v,w)\).
\leanref{GlobalTheory.\allowbreak{}CompanionTheory.\allowbreak{}Ar.\allowbreak{}Fl.\allowbreak{}IsLineFn}{Development.\allowbreak{}AutomorphicSpectrum.\allowbreak{}Definitions.\allowbreak{}GlobalGeometry.\allowbreak{}IsConstMul}{30--38}
\leanref{GlobalTheory.\allowbreak{}CompanionTheory.\allowbreak{}Ar.\allowbreak{}Flat.\allowbreak{}Represents}{Development.\allowbreak{}AutomorphicSpectrum.\allowbreak{}Definitions.\allowbreak{}GlobalGeometry.\allowbreak{}IsFlatAnIffIsConstMul}{121--130}
For the local relative-generator condition, let \(V\) be the generic Lie space, let \(R_x\subset\Omega\) be a local coefficient ring, let \(A_x,B_x\subset V\) be the eigenlattices of ranks \(p,q\), and let \(a\in\bigwedge_\Omega^p V\) and \(b\in\bigwedge_\Omega^q V\) be nonzero. A scalar \(h\in\Omega\) is a relative coefficient at \(x\) if there are \(\alpha_x,\beta_x\in\Omega^\times\) such that
\[
 \bigwedge^p A_x=R_x(\alpha_x a),\qquad
 \bigwedge^q B_x=R_x(\beta_x b),\qquad
 h\alpha_x=\beta_x.
\]
For a set \(T\) of such points, the relative-generator condition requires this equality with \(h=1\) at every \(x\in T\). Thus at each point one common scalar generates both exterior lattices, and its cancellation makes \(a\otimes b^{-1}\) a generator of their determinant quotient.
\leanref{GlobalTheory.\allowbreak{}CompanionTheory.\allowbreak{}Ar.\allowbreak{}Flat.\allowbreak{}IsCoeffAt}{Development.\allowbreak{}AutomorphicSpectrum.\allowbreak{}Definitions.\allowbreak{}GlobalGeometry.\allowbreak{}IsLineGenOn}{13--16}
\leanref{GlobalTheory.\allowbreak{}CompanionTheory.\allowbreak{}Ar.\allowbreak{}Flat.\allowbreak{}IsLineGenOn}{Development.\allowbreak{}AutomorphicSpectrum.\allowbreak{}Definitions.\allowbreak{}GlobalGeometry.\allowbreak{}IsLineGenOn}{17--19}

For every embedding \(\sigma:F_c\hookrightarrow\C\), there is a real number \(t_\sigma\) with the following property. For either \(\sigma'=\sigma\) or \(\sigma'=\overline{\sigma}\), every complex point \(y\) of the component open over \(\sigma'\), and every pair of eigenframes \(v,w\) there, let \(\Gamma,\Delta\) be their Hodge Gram matrices. Then
\[
 \log|\varphi_\C(y;v,w)|
 +\frac12\bigl(\log|\det\Gamma|-\log|\det\Delta|\bigr)
 =t_\sigma.
\]
The determinant change-of-frame identity makes this expression independent of \(v,w\). Consequently the invariant point norm is the constant \(e^{t_\sigma}\) on each of these two sets of frames, and both supremum norms are finite. The constant-norm assertion follows from Lemma~\ref{B8.24b}; conjugation of the frame coefficient and the Hodge Gram matrices preserves the displayed expression and gives the same constant at the conjugate embedding.
\leanref{GlobalTheory.\allowbreak{}CompanionTheory.\allowbreak{}Ar.\allowbreak{}Fl.\allowbreak{}HasInfNorm}{Development.\allowbreak{}AutomorphicSpectrum.\allowbreak{}Definitions.\allowbreak{}GlobalGeometry.\allowbreak{}IsConstMul}{59--64}
\leanref{GlobalTheory.\allowbreak{}CompanionTheory.\allowbreak{}Ar.\allowbreak{}Flat.\allowbreak{}FlatLineGeneratorComponent}{Development.\allowbreak{}AutomorphicSpectrum.\allowbreak{}Statements.\allowbreak{}GeometryStatements.\allowbreak{}CardComapMulFinsumEqFinsumFibres}{18--20}
\leanref{GlobalTheory.\allowbreak{}T.\allowbreak{}inf\_\allowbreak{}norm\_\allowbreak{}conjugate\_\allowbreak{}invariant}{Development.\allowbreak{}AutomorphicSpectrum.\allowbreak{}Proofs.\allowbreak{}GeometryProofs.\allowbreak{}InfNormConjugateInvariantStatement}{15--17}

Choose the generator with the relative lattice property on the whole generic fibre of the split model in Step~4. It has the following uniformised determinant description. Write \(U\) for the component open, \(\iota_0:F_c\hookrightarrow\C\) for the distinguished inclusion, and \(K_\ell\subset F_c\) for the field of definition of the level datum. Let \(\mathrm{ty}:\{1,\ldots,g\}\to\operatorname{Hom}(E,\C)\) have multiplicity \(p\) on the distinguished CM type \(\Psi\) and \(q\) on \(\overline\Psi\), and choose any family of complex uniformisations over \(\iota_0|_{K_\ell}\) of this type. Fix a member \(i\), with period lattice \(\Lambda_i\), its \(E\)-valued polarisation pairing \(S_i:\Lambda_i\times\Lambda_i\to E\), and exponential map \(\exp_{i,P}:\C^g\to\mathcal A_y(\C)\) for each admissible period map \(P:\Lambda_i\to\C^g\). Assume that every point over \(\iota_0\) represented by this member belongs to \(U\). Put \(I=\{1,\ldots,p\}\sqcup\{1,\ldots,q\}\). There are an \(\mathcal O_E\)-linearly independent tuple \((e_m)_{m\in I}\) in \(\Lambda_i\) and a scalar \(\rho_i\in\C^\times\), fixed throughout this member, with the following property.

For an embedding \(\eta:E\hookrightarrow\C\), define the coordinate projection
\[
 (\operatorname{pr}_\eta z)_k=
 \begin{cases}z_k&\mathrm{ty}(k)=\eta,\\0&\mathrm{ty}(k)\ne\eta.\end{cases}
\]
For every admissible \(P\), every represented point \(y\) over \(\iota_0\), every pair of eigenframes \(v,w\) there, every lift \(V_a,W_j\in\C^g\), and matrices \(C\in\operatorname{Mat}_{p\times I}(\C)\), \(F,B\in\operatorname{Mat}_{q\times I}(\C)\) satisfying
\[
 \begin{gathered}
 d\exp_{i,P}(0)V_a=v_a,\qquad d\exp_{i,P}(0)W_j=w_j,\\
 \operatorname{pr}_\tau\!\left(\sum_{m\in I}C_{am}P(e_m)\right)=V_a,
 \qquad
 \operatorname{pr}_{\bar\tau}P(e_n)=\sum_{j=1}^q B_{jn}W_j,\\
 \sum_{m\in I}F_{jm}\tau\bigl(S_i(e_m,e_n)\bigr)=B_{jn},
 \end{gathered}
\]
one has
\[
 \varphi_\C(y;v,w)\det\begin{pmatrix}C\\F\end{pmatrix}=\rho_i.
\]
\leanref{GlobalTheory.\allowbreak{}CompanionTheory.\allowbreak{}Ar.\allowbreak{}Fl.\allowbreak{}sesq}{Development.\allowbreak{}AutomorphicSpectrum.\allowbreak{}Definitions.\allowbreak{}GlobalGeometry.\allowbreak{}IsHodgeGram}{43--45}
\leanref{GlobalTheory.\allowbreak{}CompanionTheory.\allowbreak{}Ar.\allowbreak{}Fl.\allowbreak{}tyPart}{Development.\allowbreak{}AutomorphicSpectrum.\allowbreak{}Definitions.\allowbreak{}GlobalGeometry.\allowbreak{}IsHodgeGram}{46--48}
\leanref{GlobalTheory.\allowbreak{}CompanionTheory.\allowbreak{}Ar.\allowbreak{}Fl.\allowbreak{}IsConstMul}{Development.\allowbreak{}AutomorphicSpectrum.\allowbreak{}Definitions.\allowbreak{}GlobalGeometry.\allowbreak{}IsConstMul}{39--51}
\leanref{GlobalTheory.\allowbreak{}CompanionTheory.\allowbreak{}Ar.\allowbreak{}Fl.\allowbreak{}IsFlatGen}{Development.\allowbreak{}AutomorphicSpectrum.\allowbreak{}Definitions.\allowbreak{}GlobalGeometry.\allowbreak{}IsConstMul}{52--58}

Here is the analytic construction used for this identity. More generally, let \(\sigma:F_c\hookrightarrow\C\), let \(\tau'=\sigma\circ\tau_{F_c}\), where \(\tau_{F_c}:E\hookrightarrow F_c\) is the lift of \(\tau\), and suppose that \(\tau'\) lies in the distinguished CM type. Use uniformisations over \(\sigma|_{K_\ell}\), replace \(\tau,\bar\tau\) in the projections and polarisation matrix by \(\tau',\bar\tau'\), and take algebraic eigenframes over \(\sigma\). For each linearly independent tuple \(e\), the displayed linear equations determine \(V,W,F,B\) and a nonzero scalar \(d(P,y,v,w)\) uniquely: a matrix \(C\) satisfying them exists, and every such \(C\) has \(\det\binom CF=d(P,y,v,w)\). There is a number \(a_i>0\), depending on the member and its tuple, such that a frame function \(\Phi\) satisfying \(\Phi d=\rho\) obeys
\[
 |\Phi(y;v,w)|^2|\det\Gamma|
 =|\rho|^2a_i|\det\Delta|.
\]
Choose \(\rho=a_i^{-1/2}\). The functions \(\rho/d\) assemble on all the uniformised members to a nowhere-zero frame function \(\Phi\), with the determinant change-of-frame rule and
\[
 |\Phi(y;v,w)|^2|\det\Gamma|=|\det\Delta|.
\]
Its holomorphy has the following precise meaning. If \(R\) is a commutative ring, \(x\) is an \(R\)-point of the moduli space over \(F_c\) with algebraic tangent tuples \(v_R,w_R\), and \(\psi_b:R\to\C\) is a family on an open set \(B\subset\C^m\) such that every \(b\mapsto\psi_b(r)\) is holomorphic, then \(b\mapsto\Phi(x_b;v_b,w_b)\) is holomorphic whenever base change gives eigenframes over \(\sigma\) throughout \(B\).

For the split model over \(\mathcal O_{F_c}\) constructed in Step~4, resolution and GAGA produce a relative determinant generator whose frame function equals \(\kappa_0\Phi\), with \(\kappa_0\ne0\), on a nonempty algebraic open. Since \(F_c\) is algebraically closed in the component function field, two generators satisfying the model's relative lattice condition throughout its generic fibre differ by a scalar in \(F_c^\times\). The frame comparison extends the equality over \(U\). Thus \(\varphi_\C=\kappa\Phi\) there for one \(\kappa\in\C^\times\). At \(\sigma=\iota_0\), substitution gives the required determinant identity with \(\rho_i=\kappa a_i^{-1/2}\). Taking logarithms of the norm identity gives \(t_\sigma=\log|\kappa|\) for embeddings with \(\tau'\) in the distinguished type; conjugation gives the paired norm identity above for the other embeddings.
\leanref{GlobalTheory.\allowbreak{}CompanionTheory.\allowbreak{}Ar.\allowbreak{}Flat.\allowbreak{}IsFlatAn}{Development.\allowbreak{}AutomorphicSpectrum.\allowbreak{}Definitions.\allowbreak{}GlobalGeometry.\allowbreak{}IsFlatAnIffIsConstMul}{19--32}
\leanref{GlobalTheory.\allowbreak{}CompanionTheory.\allowbreak{}Ar.\allowbreak{}Flat.\allowbreak{}IsUnitAn}{Development.\allowbreak{}AutomorphicSpectrum.\allowbreak{}Definitions.\allowbreak{}GlobalGeometry.\allowbreak{}IsFlatAnIffIsConstMul}{33--53}
\leanref{GlobalTheory.\allowbreak{}T.\allowbreak{}uniformised\_\allowbreak{}eigen\_\allowbreak{}frame\_\allowbreak{}one}{Development.\allowbreak{}AutomorphicSpectrum.\allowbreak{}Proofs.\allowbreak{}GeometryProofs.\allowbreak{}UniformisedEigenFrameOneStatement}{13--15}
\leanref{GlobalTheory.\allowbreak{}T.\allowbreak{}uniformised\_\allowbreak{}eigen\_\allowbreak{}frame\_\allowbreak{}two}{Development.\allowbreak{}AutomorphicSpectrum.\allowbreak{}Proofs.\allowbreak{}GeometryProofs.\allowbreak{}UniformisedEigenFrameTwoStatement}{13--15}
\leanref{GlobalTheory.\allowbreak{}CompanionTheory.\allowbreak{}Ar.\allowbreak{}Flat.\allowbreak{}anUnit}{Development.\allowbreak{}AutomorphicSpectrum.\allowbreak{}Proofs.\allowbreak{}GeometryProofs.\allowbreak{}ExistsUnitAnalyticFunction}{182--227}
\leanref{GlobalTheory.\allowbreak{}CompanionTheory.\allowbreak{}Ar.\allowbreak{}Flat.\allowbreak{}flat\_\allowbreak{}eq\_\allowbreak{}const\_\allowbreak{}mul\_\allowbreak{}unitAn}{Development.\allowbreak{}AutomorphicSpectrum.\allowbreak{}Proofs.\allowbreak{}GeometryProofs.\allowbreak{}FrameFunctionAnalyticNorm}{279--304}
\leanref{GlobalTheory.\allowbreak{}CompanionTheory.\allowbreak{}Ar.\allowbreak{}Flat.\allowbreak{}FlatFirstTheorem.\allowbreak{}exists\_\allowbreak{}infNorm\_\allowbreak{}of\_\allowbreak{}rt}{Development.\allowbreak{}AutomorphicSpectrum.\allowbreak{}Proofs.\allowbreak{}GeometryProofs.\allowbreak{}FlatFirstFromResolutionAndGaga}{77--102}
\item Let \(\mathcal Y\) be a global model with split components of a unitary datum of signature \((p,q)\), with \(p,q\ge1\), coefficient field \(F_c\), host CM type \(\Psi\), and level modulus \(M'\). Put \(n=p+q\) and \(d=|\Psi|\binom n2\). Take \(\operatorname{Comp}\) to be the finite nonempty set of component generic points of \(\mathcal Y\), and set
\[
 \mathcal B=\{\ell\text{ prime}:\ell\mid2|\operatorname{disc}E|M'\},
 \qquad
 S_N=\{\ell\text{ prime}:\ell\le N\}\cup\mathcal B,
 \qquad
 \nu_{\tau,i}=-i\quad(0\le i<n).
\]
Thus \(2\in\mathcal B\), every ramified prime belongs to every \(S_N\), and \(\nu\) is rational and strictly decreasing in each row.

For each \(Z\in\operatorname{Comp}\), the reference construction supplies a root model \(R_Z\), a compact complex space \(X_Z^{\rm an}\) of dimension \(d\), a field \(\Omega_Z^0\) and a morphism \(\xi_Z:\operatorname{Spec}\Omega_Z^0\to\mathcal Y\) with image \(Z\). Its fixed data include coordinates \(y_{Z,0},\ldots,y_{Z,m_Z}\) whose ratios generate \(\Omega_Z^0\), positive integers \(k_{A,Z},q'_Z\), a nonzero \(s_{E,Z}\in\Omega_Z^0\), and exponents
\[
 n_{\tau,i}=k_{A,Z}q'_Z\nu_{\tau,i}.
\]
The reference flag-cover degree is one, its auxiliary line is ample on the generic fibre, and its embedded analytic realisation maps bijectively to \(X_Z^{\rm an}\). At every \(x\in X_Z^{\rm an}\), the projective point map has a local nonzero holomorphic homogeneous lift \(\lambda\) such that, for each \(a\) with \(\lambda_a(x)\ne0\), the differential of \((\lambda_i/\lambda_a)_i\) at \(x\) is injective. For every \(\tau\) and every real regularisation parameter, both the terminal fundamental potential and its negative are plurisubharmonic.

For each \(N,Z\), the finite-prime construction applied to the reference data and \(S_N\), followed by passage to birational multihomogeneous coordinates, gives \(M(N,Z)\) with smooth generic fibre, a normal Hodge model \(\mathfrak M_{N,Z}\) over \(\mathcal Y\), and a field embedding \(g_{N,Z}:\Omega_Z^0\hookrightarrow\Omega_{N,Z}\) compatible with the generic morphisms and flag bases. If \(d_{N,Z}\) is the flag-cover degree, then
\[
 [\Omega_{N,Z}:\Omega_Z^0]=d_{N,Z},
 \qquad \delta_{N,Z}=[F_c:\Q]d_{N,Z}.
\]
For every \(\ell\in S_N\), the model parameter is \(4\) and the generic semiabelian scheme has full \(\ell^4\)-torsion. Patching the de Rham eigenlattices away from \(S_N\) with the Hodge--Tate eigenlattices at \(S_N\) gives locally free lattices of the prescribed eigenranks.

The flags supply root lattices \(\Lambda_{\tau,i}\), their multidegrees, rational factors, singular Hodge metrics and smooth conjugation-invariant root metrics \(x_{N,Z}(t,\tau,i)\). The auxiliary line has degree one in the coordinate block obtained from \((g_{N,Z}(y_{Z,i}))_i\) and degree zero in the other blocks. Write \(t_{0,Z}\) for the regularisation threshold of the fixed analytic system, \(H_{\tau,i}(t)\) for the regularised fundamental line, and \(h_L\) for the metric coefficient of a Hermitian line \(L\). At every multihomogeneous complex point \(z\),
\[
 \log h_{H_{\tau,i}(t)}(z)
 =-\sum_{k=0}^{i}\log h_{x_{N,Z}(t,\tau,k)}(z)
 \qquad(0<t\le t_{0,Z}).
\]
At analytic-realisation lifts whose image belongs to the open part of \(X_Z^{\rm an}\), the singular identity has the reference fundamental log metric plus its singular potential on the left and the negative sum of the singular root log metrics on the right. The analytic realisation has degree \(d_{N,Z}\) and carries the auxiliary metric from the fixed data. A represented point lies in the abelian interior exactly when its image belongs to that open part.

The section \(s_{E,Z}\) records effectivity. At a point over characteristic zero, on the projective coordinate chart indexed by \(a\), if \(u\) generates \(\bigotimes_{\tau,i}\Lambda_{\tau,i}^{-n_{\tau,i}}\), then
\[
 \frac{g_{N,Z}(s_{E,Z})\,g_{N,Z}(y_{Z,a})^{q'_Z}}{u}
\]
is regular and is a nonunit outside the abelian locus. Assemble the models with split index \(p\), component set \(\operatorname{Comp}\), bad-prime set \(\mathcal B\) and weight \(\nu\). The generic-point maps identify their components with those of \(\mathcal Y\); the fixed analytic systems and reference models give the origin data.
\leanref{GlobalTheory.\allowbreak{}CompanionTheory.\allowbreak{}Ar.\allowbreak{}TopTower.\allowbreak{}comp\_\allowbreak{}data}{Development.\allowbreak{}AutomorphicSpectrum.\allowbreak{}Proofs.\allowbreak{}GeometryProofs.\allowbreak{}ExistsGoodRootTower}{20--58}
\leanref{GlobalTheory.\allowbreak{}CompanionTheory.\allowbreak{}Ar.\allowbreak{}TopTower.\allowbreak{}level\_\allowbreak{}data}{Development.\allowbreak{}AutomorphicSpectrum.\allowbreak{}Proofs.\allowbreak{}GeometryProofs.\allowbreak{}ExistsGoodRootTower}{59--76}
\leanref{GlobalTheory.\allowbreak{}CompanionTheory.\allowbreak{}Ar.\allowbreak{}TopTower.\allowbreak{}exists\_\allowbreak{}good\_\allowbreak{}tower}{Development.\allowbreak{}AutomorphicSpectrum.\allowbreak{}Proofs.\allowbreak{}GeometryProofs.\allowbreak{}ExistsGoodRootTower}{77--145}
\leanref{GlobalTheory.\allowbreak{}CompanionTheory.\allowbreak{}Ar.\allowbreak{}FixedData}{Development.\allowbreak{}AutomorphicSpectrum.\allowbreak{}Definitions.\allowbreak{}GlobalGeometry.\allowbreak{}TowerOrigin}{18--55}
\leanref{GlobalTheory.\allowbreak{}CompanionTheory.\allowbreak{}Ar.\allowbreak{}RootOrigin}{Development.\allowbreak{}AutomorphicSpectrum.\allowbreak{}Definitions.\allowbreak{}GlobalGeometry.\allowbreak{}TowerOrigin}{58--104}
\leanref{GlobalTheory.\allowbreak{}CompanionTheory.\allowbreak{}Ar.\allowbreak{}TowerOrigin}{Development.\allowbreak{}AutomorphicSpectrum.\allowbreak{}Definitions.\allowbreak{}GlobalGeometry.\allowbreak{}TowerOrigin}{105--120}
\leanref{GlobalTheory.\allowbreak{}CompanionTheory.\allowbreak{}Ar.\allowbreak{}HodgeCpt}{Development.\allowbreak{}AutomorphicSpectrum.\allowbreak{}Definitions.\allowbreak{}GlobalGeometry.\allowbreak{}IsHodgeLine}{33--75}

Clause~2 of Definition~\ref{B8.26r} is the flag-height assertion of Lemma~\ref{B8.26s}(4), whose proof uses the arithmetic Kleiman premise. The tensor product of the metrised quotient lines in one row is the trivial Hermitian determinant line, so \(h_{\rm fl}(p_{\tau,1}Q)=0\). Modulo these determinant classes, a dominant weight is a nonnegative combination of the Pl\"ucker lines; arithmetic Kleiman gives nonnegative top height. On a wall the corresponding Pl\"ucker coefficient vanishes, the line descends to the partial flag variety omitting that step, and its top height is zero.
\item Let \(X_0\), \(F_c\) and \(\ell\) denote the moduli scheme, coefficient field and level of the deep cover. Choose a family \((X_{\ell'})_{\ell'}\) of moduli schemes for the same unitary datum whose member at \(\ell\) is \(X_0\). The level \(g_0g_1\) is good at every prime, multiplication of its modulus preserves this property, and the auxiliary level factor of the deep cover equals one. The global model construction gives a normal projective flat model \(Y_{\rm gl}\) of \(X_0\) over \(\mathcal O_{F_c}\), with its semiabelian extension, which descends to the moduli model and dominates the compatible local model at each valuation ring over a rational prime. The coefficient field is algebraically closed in the component function fields, so \(Y_{\rm gl}\) has split components. Apply Step~3 to \(Y_{\rm gl}\) with \((p,q)=(p_{\rm sig},q_{\rm sig})\). The comparison of component generic points gives a bijection \(\kappa:\operatorname{Comp}(X_0/F_c)\xrightarrow{\sim}\operatorname{Comp}(\mathscr T)\). Reindexing the tower by \(\kappa\) gives the asserted component set.
\leanref{Formalism.\allowbreak{}L2.\allowbreak{}Pj.\allowbreak{}goodAt\_\allowbreak{}of\_\allowbreak{}two\_\allowbreak{}primes}{Formalism.\allowbreak{}LocalToGlobal.\allowbreak{}ProjectiveLattices.\allowbreak{}ModelAt}{124--154}
\leanref{GlobalTheory.\allowbreak{}LawsDeep.\allowbreak{}goodAt\_\allowbreak{}of\_\allowbreak{}mul}{Development.\allowbreak{}AbelianVarietyTheory.\allowbreak{}Proofs.\allowbreak{}DivisorProofs.\allowbreak{}ComplexUniformizationLawsOfGoodDeepCover}{12--20}
\leanref{GlobalTheory.\allowbreak{}CompanionTheory.\allowbreak{}Ar.\allowbreak{}Rec.\allowbreak{}exists\_\allowbreak{}isGlobal}{Development.\allowbreak{}AutomorphicSpectrum.\allowbreak{}Proofs.\allowbreak{}GeometryProofs.\allowbreak{}ExistsIsGlobal}{33--39}
\leanref{GlobalTheory.\allowbreak{}T.\allowbreak{}split\_\allowbreak{}components\_\allowbreak{}of\_\allowbreak{}alg\_\allowbreak{}closed}{Development.\allowbreak{}AutomorphicSpectrum.\allowbreak{}Proofs.\allowbreak{}GeometryProofs.\allowbreak{}SplitComponentsOfAlgClosedStatement}{15--17}
\leanref{GlobalTheory.\allowbreak{}T.\allowbreak{}tower\_\allowbreak{}components\_\allowbreak{}equiv\_\allowbreak{}of\_\allowbreak{}given}{Development.\allowbreak{}AutomorphicSpectrum.\allowbreak{}Proofs.\allowbreak{}GeometryProofs.\allowbreak{}TowerComponentsEquivOfGivenStatement}{30--32}

For the remaining clauses put \(r=p_{\rm sig}\), \(s=q_{\rm sig}\), \(f=|\Psi|\), and \(d=f\binom n2\). The deep-cover assumptions give \(s=m-j\), \(n=k+m-1\), \(m\ge2\), \(m^2\le k\), and \(1\le j<m\). Hence
\[
 s\ge1,\qquad 3s+1\le n,\qquad s\le r,\qquad n\ge4.
\]
For the host order parameter \(n_{\rm ord}\ge1\), one has \(f=2n_{\rm ord}\). The host discriminant estimate and the bound of fifteen prime divisors of the level leave at most eighteen bad primes above \(n_{\rm ord}\). Thus
\[
 \#\mathcal B_{\mathscr T}\le n_{\rm ord}+19
 \le128n_{\rm ord}=64f\le4fn^2.
\]
\leanref{Formalism.\allowbreak{}BasicNotions.\allowbreak{}Host.\allowbreak{}card\_\allowbreak{}bad\_\allowbreak{}gt\_\allowbreak{}le}{Development.\allowbreak{}AutomorphicArithmetic.\allowbreak{}PrimeDiscriminants.\allowbreak{}BadPrimes}{74--79}

\item For clause~1, the analytic systems underlying the root models are uniformised with split index \(r\). Their degree polynomials and that of the complete flag variety are positive multiples of
\[
 \Delta(v)=\prod_{\tau\in\Psi}\prod_{0\le a<b<n}
                   (v_{\tau,a}-v_{\tau,b}).
\]
Evaluation at a fixed regular dominant weight fixes the positive proportionality constant. Geometric comparison with \(M(N,Z)\), including the normalisation by \(\delta_{N,Z}\) in Step~3, therefore gives \(V_Z>0\), independent of \(N\), such that
\[
 \deg_{M(N,Z)}(L_v^d)=V_ZD(L_v^d)
\]
for every real weight \(v\). Polarisation gives \(\deg_{M(N,Z)}(F)=V_ZD(F)\) for every homogeneous polynomial of degree \(d\); both functionals vanish in the other degrees. The comparison with the fixed analytic system also makes the geometric functional including the auxiliary class independent of the cutoff. The component comparison for a global split model with all local modification parameters equal to \(4\) gives \(V_Z=V_{Z'}\) whenever \(s\ge1\) and \(3s+1\le n\). The range proved above supplies these inequalities.
\leanref{GlobalTheory.\allowbreak{}T.\allowbreak{}good\_\allowbreak{}tower\_\allowbreak{}real\_\allowbreak{}agreement}{Development.\allowbreak{}AutomorphicSpectrum.\allowbreak{}Proofs.\allowbreak{}DivisibilityProofs.\allowbreak{}GoodTowerRealAgreementStatement}{16--18}
\leanref{GlobalTheory.\allowbreak{}T.\allowbreak{}proportional'}{Development.\allowbreak{}AutomorphicSpectrum.\allowbreak{}Proofs.\allowbreak{}MetricProofs.\allowbreak{}ProportionalStatement}{9--11}
\leanref{GlobalTheory.\allowbreak{}CompanionTheory.\allowbreak{}Ar.\allowbreak{}PB.\allowbreak{}proportional}{Development.\allowbreak{}AutomorphicSpectrum.\allowbreak{}Proofs.\allowbreak{}MetricProofs.\allowbreak{}Proportional}{78--94}
\leanref{GlobalTheory.\allowbreak{}T.\allowbreak{}real\_\allowbreak{}tower\_\allowbreak{}proportional}{Development.\allowbreak{}AutomorphicSpectrum.\allowbreak{}Proofs.\allowbreak{}DivisibilityProofs.\allowbreak{}RealTowerProportionalStatement}{7--9}
\leanref{GlobalTheory.\allowbreak{}CompanionTheory.\allowbreak{}Ar.\allowbreak{}Top.\allowbreak{}rootTower\_\allowbreak{}proportional\_\allowbreak{}of\_\allowbreak{}unif}{Development.\allowbreak{}AutomorphicSpectrum.\allowbreak{}Proofs.\allowbreak{}DivisibilityProofs.\allowbreak{}RealTowerProportional}{153--208}
\leanref{GlobalTheory.\allowbreak{}T.\allowbreak{}good\_\allowbreak{}tower\_\allowbreak{}intersection\_\allowbreak{}extra}{Development.\allowbreak{}AutomorphicSpectrum.\allowbreak{}Proofs.\allowbreak{}DivisibilityProofs.\allowbreak{}GoodTowerIntersectionExtraStatement}{16--18}

\item For clause~3, write
\[
 A_M(v)=P_M(L_v^{d+1}),\qquad
 H(v)=h_{\rm fl}(L_v^{d+1}),\qquad
 \delta(v)=D(L_v^d).
\]
Choose \(c_0>0\) and, for each \(\tau_0\), a compact patch with nonempty relative interior in the wall \(v_{\tau_0,r-1}=v_{\tau_0,r}\), on which \(\|v\|_\infty\le\tfrac12\) and every other adjacent gap is at least \(c_0\). The curve-height bounds and successive-minima inequalities of Lemma~\ref{B8.26p}, using Zhang's bounds \cite[Th.~5.2]{zhang-1995}, give a uniform wall estimate as follows. For rational \(v\) put
\[
 \theta=\frac{5f}{M},\qquad
 v'_{\tau,i}=v_{\tau,i}
       +2\theta\,\mathbf1_{\{\tau=\tau_0,\ i<r\}},
 \qquad \varepsilon=\theta^2.
\]
For sufficiently large \(M\), the perturbed weight \(v'+\varepsilon\nu\) has the required positive gaps, and
\[
 \|v'+\varepsilon\nu-v\|_\infty
       \le2\theta+\|\nu\|_\infty\theta^2,
 \qquad \log(M/\varepsilon)\le3\log M.
\]
Its geometric degree is at most a fixed multiple of \(V_Z\theta\). The component height estimate, regularisation error and Lipschitz bound for the arithmetic functional then give
\[
 |P_{M,Z}(L_v^{d+1})|\le C_Z\theta\log M.
\]
Sum over components and use continuity in \(v\). On the chosen wall patches this gives
\[
 |A_M(v)|\le C_{\rm wall}\frac{\log M}{M}.
\]

The full-flag intersection comparison and the alternating-polynomial perturbation estimate give real coefficients \(a_{M,\tau},b_{M,\tau}\) such that every coefficient of
\[
 A_M(v)-VH(v)
 -V(d+1)\delta(v)
   \sum_{\tau\in\Psi}
       \left(a_{M,\tau}\sum_{i<r}v_{\tau,i}
            +b_{M,\tau}\sum_{i\ge r}v_{\tau,i}\right)
\]
has absolute value at most \(C_1\log M/M\), with \(C_1\) fixed by the tower. The flag-height polynomial vanishes on these walls by clause~2; the perturbation estimate bounds this coefficient error by a fixed multiple of the wall error.

Twisting by a host automorphism gives an isomorphism of moduli problems and permutes the components and CM labels. The comparison of the corresponding root models says that the difference of their normalised intersections of degree \(d+1\) tends to zero with the smoothing parameter. Their local parameters agree, since they are all \(4\). Passing to the limiting functionals and summing over components therefore permits averaging the correction over the host-label action. The flag-height functional is invariant under this relabelling, and the averaged correction has two coefficients \(\alpha_M,\beta_M\), constant over \(\tau\).

For powers of \(L_v\), direct differentiation gives
\[
 \operatorname{Sh}_r(\tau,L_v^{d+1})
   =-(d+1)\left(\sum_{i<r}v_{\tau,i}\right)\delta(v),
 \qquad
 \operatorname{Tot}(\tau,L_v^{d+1})
   =-(d+1)\left(\sum_i v_{\tau,i}\right)\delta(v).
\]
Polarising the averaged coefficient estimate gives \(C\ge0\), \(N_1\ge2\) and, for every \(M\ge N_1\) and every monomial \(m\) of degree \(d+1\),
\[
 \left|P_M(m)-V\left(h_{\rm fl}(m)
       -(\alpha_M-\beta_M)\sum_\tau\operatorname{Sh}_r(\tau,m)
       -\beta_M\sum_\tau\operatorname{Tot}(\tau,m)\right)\right|
       \le C\frac{\log M}{M}.
\]
Extend the coefficient sequences arbitrarily below \(N_1\). This is clause~3.
\leanref{GlobalTheory.\allowbreak{}CompanionTheory.\allowbreak{}Ar.\allowbreak{}Top.\allowbreak{}thmW\_\allowbreak{}tower}{Development.\allowbreak{}AutomorphicSpectrum.\allowbreak{}Proofs.\allowbreak{}DivisibilityProofs.\allowbreak{}TowerWeilTheorem}{120--204}
\leanref{GroundTheory.\allowbreak{}Heights₂.\allowbreak{}Heights.\allowbreak{}AntisymPerturbation.\allowbreak{}Statement}{GroundTheory.\allowbreak{}HeightsAndTheta.\allowbreak{}Heights.\allowbreak{}Antisymmetric.\allowbreak{}Statements}{86--98}
\leanref{GlobalTheory.\allowbreak{}CompanionTheory.\allowbreak{}Ar.\allowbreak{}Cmp.\allowbreak{}CmpPair}{Development.\allowbreak{}AutomorphicSpectrum.\allowbreak{}Statements.\allowbreak{}GeometryStatements.\allowbreak{}CmpPair}{12--23}
\leanref{GlobalTheory.\allowbreak{}CompanionTheory.\allowbreak{}Ar.\allowbreak{}TwistAut}{Development.\allowbreak{}AutomorphicSpectrum.\allowbreak{}Statements.\allowbreak{}GeometryStatements.\allowbreak{}PadicIntegersHenselianFiniteResidue}{11--15}
\leanref{GlobalTheory.\allowbreak{}T.\allowbreak{}good\_\allowbreak{}tower\_\allowbreak{}propA\_\allowbreak{}of\_\allowbreak{}global\_\allowbreak{}extra}{Development.\allowbreak{}AutomorphicSpectrum.\allowbreak{}Proofs.\allowbreak{}DivisibilityProofs.\allowbreak{}GoodTowerGlobalExtraStatement}{18--20}
\leanref{GlobalTheory.\allowbreak{}CompanionTheory.\allowbreak{}Ar.\allowbreak{}Top.\allowbreak{}chain\_\allowbreak{}tower}{Development.\allowbreak{}AutomorphicSpectrum.\allowbreak{}Proofs.\allowbreak{}DivisibilityProofs.\allowbreak{}TowerWeilInputsCombined}{22--54}
\leanref{GlobalTheory.\allowbreak{}T.\allowbreak{}tower\_\allowbreak{}weil\_\allowbreak{}inputs\_\allowbreak{}combined}{Development.\allowbreak{}AutomorphicSpectrum.\allowbreak{}Proofs.\allowbreak{}DivisibilityProofs.\allowbreak{}TowerWeilInputsCombinedStatement}{7--9}
\leanref{GlobalTheory.\allowbreak{}CompanionTheory.\allowbreak{}Ar.\allowbreak{}Top.\allowbreak{}exists\_\allowbreak{}finite\_\allowbreak{}identity}{Development.\allowbreak{}AutomorphicSpectrum.\allowbreak{}Proofs.\allowbreak{}DivisibilityProofs.\allowbreak{}TowerIdentityAtCutoff}{97--150}

\item For clause~4, let \(\lambda_\epsilon:\mathcal O_E\to F_c\) be the embedding attached to a signed label \(\epsilon\in\Psi\sqcup\Psi\), with the right copy obtained by conjugation. Choose permutations \(\varphi_\ell\) of \(\Psi\sqcup\Psi\), for \(\ell\in\N\), such that for every prime \(\ell\), every prime ideal \(\mathfrak n\subset\mathcal O_{F_c}\) containing \(\ell\), every signed label \(\epsilon\), and every \(a\in\mathcal O_E\),
\[
 \lambda_{\varphi_\ell(\epsilon)}(a)
       \equiv\lambda_\epsilon(a)^\ell\pmod{\mathfrak n}.
\]
The congruence uses the integral representatives in \(\mathcal O_{F_c}\). The abelian Galois hypothesis gives these permutations. The distinguished type \(\Psi\) is induced by the host root: it consists of the embeddings sending the host root to its fixed complex value. A Frobenius element sends the host root to itself or to its negative, so it carries \(\Psi\) to \(\Psi\) or to \(\overline\Psi\). This makes the side of
\[
 \operatorname{Fr}_{\ell,h}(\tau)
       =\varphi_\ell^{\,h+1}(\operatorname{inl}\tau)
 \qquad(\ell,h\in\N,\ \tau\in\Psi)
\]
independent of \(\tau\). In particular \(h=0\) denotes the first arithmetic Frobenius iterate.

Fix \(2\le N\le M\) and \(Z\), and put
\[
 S_{N,M}=\{\ell:\ell\text{ prime},\ N<\ell\le M,
                              \ell\notin\mathcal B_{\mathscr T}\}.
\]
The two cutoff models admit a common root model with two collections of root lines. Their convergent limiting functionals are \(P_{N,Z}\) and \(P_{M,Z}\), and the common model has the geometric functional of \(M(N,Z)\). The collections differ only over \(S_{N,M}\): the cutoff sets are nested, contain the discriminant primes and \(2\), and their Hodge--Tate parameters agree on their intersection because they all equal \(4\).

There are comparison polynomials \(C_{\ell,k,\tau}\) with the following local property. At a discrete valuation ring over \(\ell\in S_{N,M}\), for \(k\ge1\), the degree-\(k\) Newton power-sum comparison of the two \(\tau\)-flags, intersected with any list of \(d+1-k\) further multidegrees, is the intersection of that list and \(C_{\ell,k,\tau}\) with the divisor of \(\ell\). On each special-fibre component and for every \(\tau\), an extension of the residue field supplies homogeneous coordinates cutting out the component and a basis of the trivial rank-\(n\) bundle that gives a flag frame of the prescribed root multidegrees. The Hodge--Tate input is Lemma~\ref{B8.26q}(3), including the Fargues annihilation bound \cite[Th.~4 and Cor.~8]{fargues-2011}.

For \(1\le k\le n\) and \(Q\) homogeneous of degree \(d+1-k\), the comparison polynomial has numerical evaluation
\[
 D(C_{\ell,k,\tau}Q)
 =k\left(\sum_{h\ge0}\ell^{-k(h+1)}
               D(P_{\ell,\tau,h+1,k-1}Q)
       -\frac{D(p_{\tau,k-1}^{<r}Q)}{\ell-1}\right).
\]
The local height-change identity holds at every positive regularisation parameter. Passing to the two limiting functionals gives
\[
 P_{N,Z}(p_{\tau,k}Q)-P_{M,Z}(p_{\tau,k}Q)
   =\sum_{\ell\in S_{N,M}}\log\ell\,
                       G_{\rm com}(C_{\ell,k,\tau}Q),
\]
where \(G_{\rm com}\) is the common geometric functional. By clause~1, \(G_{\rm com}=V_ZD\). Substitution proves clause~4 with its specified right-block sign.
\leanref{GlobalTheory.\allowbreak{}CompanionTheory.\allowbreak{}Ar.\allowbreak{}IsFrobLabels}{Development.\allowbreak{}AutomorphicSpectrum.\allowbreak{}Definitions.\allowbreak{}GlobalGeometry.\allowbreak{}IsLatFlagPair}{32--35}
\leanref{GlobalTheory.\allowbreak{}CompanionTheory.\allowbreak{}Ar.\allowbreak{}Num.\allowbreak{}FlagChernComparison}{Development.\allowbreak{}AutomorphicSpectrum.\allowbreak{}Definitions.\allowbreak{}DivisibilityTowers.\allowbreak{}BundleChernComparison}{120--125}
\leanref{GlobalTheory.\allowbreak{}CompanionTheory.\allowbreak{}Ar.\allowbreak{}Num.\allowbreak{}TrivFlagAt}{Development.\allowbreak{}AutomorphicSpectrum.\allowbreak{}Definitions.\allowbreak{}DivisibilityTowers.\allowbreak{}BundleChernComparison}{92--97}
\leanref{GlobalTheory.\allowbreak{}T.\allowbreak{}good\_\allowbreak{}tower\_\allowbreak{}global\_\allowbreak{}conditions}{Development.\allowbreak{}AutomorphicSpectrum.\allowbreak{}Proofs.\allowbreak{}DivisibilityProofs.\allowbreak{}GoodTowerGlobalConditionsStatement}{17--19}
\leanref{GlobalTheory.\allowbreak{}CompanionTheory.\allowbreak{}Ar.\allowbreak{}Top.\allowbreak{}Link.\allowbreak{}goodOrigin\_\allowbreak{}linked}{Development.\allowbreak{}AutomorphicSpectrum.\allowbreak{}Proofs.\allowbreak{}DivisibilityProofs.\allowbreak{}GoodTowerGlobalConditions}{197--231}
\leanref{GlobalTheory.\allowbreak{}T.\allowbreak{}linked\_\allowbreak{}tower\_\allowbreak{}power\_\allowbreak{}sum}{Development.\allowbreak{}AutomorphicSpectrum.\allowbreak{}Proofs.\allowbreak{}DivisibilityProofs.\allowbreak{}LinkedTowerPowerSumStatement}{8--10}
\leanref{GlobalTheory.\allowbreak{}CompanionTheory.\allowbreak{}Ar.\allowbreak{}Num.\allowbreak{}power\_\allowbreak{}sum\_\allowbreak{}term\_\allowbreak{}linked\_\allowbreak{}tower}{Development.\allowbreak{}AutomorphicSpectrum.\allowbreak{}Proofs.\allowbreak{}DivisibilityProofs.\allowbreak{}LinkedTowerPowerSum}{12--37}

\item For clause~5, the flat-degree theorem and the Hodge--Tate discrepancy bound of Lemma~\ref{B8.26q}(3) give, for each \(Z,\tau\), a flat Hermitian line of algebraic class \(\sum_iX_{\tau,i}\) on \(M(2,Z)\). Write \(h_{\rm flat}\) for its metric. The root metrics satisfy
\[
 \prod_{i=0}^{n-1}h_{x_{\tau,i}(t)}
             =h_{\rm flat}\exp(-v_t),
\]
where \(v_t\) is smooth, invariant under blockwise coordinate scaling, and has integral tending to zero against every degree-\(d\) monomial in the root curvature forms as \(t\downarrow0\). Denote its flat height by \(H_{Z,\tau}\). The flat pushforward formula and clause~1 give
\[
 P_{2,Z}(p_{\tau,1}Q)=H_{Z,\tau}V_ZD(Q)
 \qquad(\deg Q=d).
\]
Put \(u=\log\operatorname{rd}(E)\). The flat height has a decomposition
\[
 H_{Z,\tau}=h_{\rm rec}+\sum_{\ell\in\mathcal B_{\mathscr T}}d_\ell,
 \qquad |h_{\rm rec}|\le\frac n2(3u+|L_E|),
\]
with \(v_\ell\ge0\) such that
\[
 |d_\ell|\le n\left(\frac1{\ell-1}+v_\ell\right)\log\ell,
 \qquad
 \sum_{\ell\in\mathcal B_{\mathscr T}}v_\ell\log\ell\le u.
\]
For a good prime \(\ell\), let \(a_{\ell,\tau,h}\) equal \(r\) on the left Frobenius copy and \(s\) on the right, and put
\[
 e_{\ell,\tau}
   =\sum_{h\ge0}\ell^{-(h+1)}a_{\ell,\tau,h}
       -\frac r{\ell-1}.
\]
This series converges absolutely, \(|e_{\ell,\tau}|\le n/(\ell-1)\), and
\(\delta_\ell(\tau,1,Q)=e_{\ell,\tau}\log\ell\,D(Q)\).
For \(N\ge2\), define
\[
 h_{N,Z,\tau}=H_{Z,\tau}
   -\sum_{\substack{2<\ell\le N\ {\rm prime}\\
                     \ell\notin\mathcal B_{\mathscr T}}}
                  e_{\ell,\tau}\log\ell.
\]
Clause~4 with initial cutoff \(2\) and \(k=1\) gives
\(P_{N,Z}(p_{\tau,1}Q)=h_{N,Z,\tau}V_ZD(Q)\).

Extend the local discrepancies to the good primes at most \(N\) by \(d_\ell=-e_{\ell,\tau}\log\ell\) and \(v_\ell=0\). For \(N\ge\max(128,\#\mathcal B_{\mathscr T})\), equation~\eqref{eq:p8c3-prime-sum} and monotonicity of \(\log x/(x-1)\) give
\[
 \sum_{\substack{\ell\ {\rm prime}\\
                  \ell\le N\ {\rm or}\ \ell\in\mathcal B_{\mathscr T}}}
                   \frac{\log\ell}{\ell-1}
       \le\log N+2.2+\log(N+1)\le4\log N.
\]
There are at most \(N\) bad primes above \(N\), each contributing at most \(\log(N+1)/N\). Since \(\operatorname{rd}(E)\ge1\),
\[
 |h_{N,Z,\tau}|
 \le n\left(\frac52u+\frac12|L_E|+4\log N\right)
 \le4n\bigl(\log(N\operatorname{rd}(E))+|L_E|\bigr).
\]
This proves clause~5 with constant \(4\).
\leanref{GlobalTheory.\allowbreak{}CompanionTheory.\allowbreak{}Ar.\allowbreak{}Fl.\allowbreak{}FlatDegree}{Development.\allowbreak{}AutomorphicSpectrum.\allowbreak{}Statements.\allowbreak{}GeometryStatements.\allowbreak{}FlatDegree}{12--20}
\leanref{GlobalTheory.\allowbreak{}CompanionTheory.\allowbreak{}Ar.\allowbreak{}Fl.\allowbreak{}flatDegree}{Development.\allowbreak{}AutomorphicSpectrum.\allowbreak{}Proofs.\allowbreak{}GeometryProofs.\allowbreak{}FlatDegreeHolds}{60--119}
\leanref{GlobalTheory.\allowbreak{}T.\allowbreak{}good\_\allowbreak{}tower\_\allowbreak{}real\_\allowbreak{}extra}{Development.\allowbreak{}AutomorphicSpectrum.\allowbreak{}Proofs.\allowbreak{}DivisibilityProofs.\allowbreak{}GoodTowerRealExtraStatement}{16--18}
\leanref{GlobalTheory.\allowbreak{}CompanionTheory.\allowbreak{}Ar.\allowbreak{}Flat.\allowbreak{}RootFlat}{Development.\allowbreak{}AutomorphicSpectrum.\allowbreak{}Definitions.\allowbreak{}GlobalGeometry.\allowbreak{}NumberField}{12--29}
\leanref{GlobalTheory.\allowbreak{}CompanionTheory.\allowbreak{}Ar.\allowbreak{}Flat.\allowbreak{}NearRecord}{Development.\allowbreak{}AutomorphicSpectrum.\allowbreak{}Definitions.\allowbreak{}GlobalGeometry.\allowbreak{}NumberField}{34--38}
\leanref{GlobalTheory.\allowbreak{}T.\allowbreak{}PropFTop\_\allowbreak{}holds}{Development.\allowbreak{}AutomorphicSpectrum.\allowbreak{}Proofs.\allowbreak{}GeometryProofs.\allowbreak{}PropFTopHoldsStatement}{7--9}
\leanref{GlobalTheory.\allowbreak{}CompanionTheory.\allowbreak{}Ar.\allowbreak{}Flat.\allowbreak{}propFTop}{Development.\allowbreak{}AutomorphicSpectrum.\allowbreak{}Proofs.\allowbreak{}GeometryProofs.\allowbreak{}PropFTopHolds}{25--98}

\item For the base construction in clause~6, assume \(s\le r\) and fix \(N\ge\max(128,4fn^2,\#\mathcal B_{\mathscr T})\) and \(Z\). Let
\[
 S_N^{\rm good}=\{\ell:\ell\text{ prime},\ \ell\le N,
                                   \ell\notin\mathcal B_{\mathscr T}\}.
\]
The generic point of the cutoff-\(N\) Hodge model agrees with that of the fixed component model. At this point the base construction chooses the left labels for the rank-\(r\) block and the right labels for the rank-\(s\) block. It gives a horizontal base model of arithmetic dimension \(e+1\), where \(e=frs\), with smooth generic fibre, normal-crossing boundary charts, and an auxiliary positive Hermitian line \(A\) whose normalised height on every horizontal arithmetic curve is at least \(-K_A\), for a real constant \(K_A\).

The generic rank-\(n\) bundle at \(\tau\) is the sum of the two labelled Lie eigenspaces, with distinguished rank-\(r\) subspace the left eigenspace. The tensor product over \(\tau\) of these quotient determinants is the line of class \(\mu\). At every \(\ell\in S_N^{\rm good}\), the full \(\ell^R\)-torsion with \(R\ge4\) gives Hodge--Tate frames. The determinant of the change from an integral small-block frame to a Hodge--Tate frame is, up to a local unit, the rational section of the prescribed vertical divisor. These give the finite-place determinant corrections on the base model.

Choose the nonnegative cusp weight \(U\), a complex model, a tied regularisation of the original Hodge metric, and a tied bounded regularisation of its corrected logarithmic weight \(\phi_\mu+U/2\). Let \(w_{t'}\) be the latter regularised weight, \(h_{\rm ref}\) its smooth reference metric, and \(v_\ell\) the vertical determinant factors. For \(1\le a\le b\), the cleared-denominator line has generic class \(b\mu+aA\) and metric
\[
 \left(h_{\rm ref}e^{w_{t'}}
       \prod_{\ell\in S_N^{\rm good}}v_\ell\right)^b h_A^a
\]
on the model. Denote it by \(\overline M_{N,a,b,t'}\), and write \(h_{\rm reg}(W)=II(W;\overline M_{N,a,b,t'})/(b\deg_\C W)\) for its normalised height on a horizontal curve \(W\). For all sufficiently small \(t'>0\), these lines are smooth, conjugation-invariant, semipositive along the defining ideal, nef on fibres and ample on the generic fibre. The uniform height and cusp estimates give constants \(C_h,C_u\) such that on both interior and boundary horizontal curves
\[
 h_{\rm reg}\ge-C_hfs\lambda_N-\frac abK_A,
 \qquad
 \operatorname{mean}U\le C_ufs\lambda_N.
\]
The mean estimate uses \(3s+1\le n\). The constants \(C_h,C_u\) are those of the uniform base-height and cusp-mean estimates, and are independent of the component and cutoff.

Let \(g>0\) be the generic degree of the unperturbed determinant line and \(A_N\) its arithmetic intersection value from the tied original-metric regularisation. The comparison of generic points and flag fibre intersections gives
\[
 \frac{A_N}{(e+1)g}
   =\frac{P_{N,Z}(F_*)}{(e+1)V_ZD(\mu^eJ)},
\]
with \(J,F_*\) as in Definition~\ref{B8.26r}. Write \(I_{a,b,t'}\), \(D_{a,b,t'}\) for the arithmetic top intersection and generic top degree of the cleared-denominator line. The Stokes comparison states that, for every \(\eta\in\R\), if for every \(1\le a\le b\) there is \(t_0>0\) with
\[
 (e+1)\left(\eta-\frac abK_A\right)
        \frac{D_{a,b,t'}}{b^e}
       \le\frac{I_{a,b,t'}}{b^{e+1}}
 \qquad(0<t'\le t_0),
\]
then
\[
 \eta\le\frac{P_{N,Z}(F_*)}{(e+1)V_ZD(\mu^eJ)}
                    +\frac12\operatorname{mean}U.
\]
The inequalities \(N\ge4fn^2\), \(f\ge1\), \(n\ge2\) give \(N\ge4f+2\), as needed for positivity and Stokes. The lower arithmetic intersection inequality applied to the curve-height bound permits \(\eta=-C_hfs\lambda_N\). It yields the base slope lower bound with constant \(C_h+C_u/2\).
\leanref{GlobalTheory.\allowbreak{}T.\allowbreak{}base\_\allowbreak{}setting\_\allowbreak{}of\_\allowbreak{}good\_\allowbreak{}tower}{Development.\allowbreak{}AutomorphicSpectrum.\allowbreak{}Proofs.\allowbreak{}GeometryProofs.\allowbreak{}BaseSettingOfGoodTowerStatement}{24--26}
\leanref{GlobalTheory.\allowbreak{}CompanionTheory.\allowbreak{}Ar.\allowbreak{}Cusp.\allowbreak{}TopBase.\allowbreak{}base\_\allowbreak{}setting\_\allowbreak{}of\_\allowbreak{}good\_\allowbreak{}origin}{Development.\allowbreak{}AutomorphicSpectrum.\allowbreak{}Proofs.\allowbreak{}GeometryProofs.\allowbreak{}BaseSettingOfGoodTower}{45--107}
\leanref{GlobalTheory.\allowbreak{}CompanionTheory.\allowbreak{}Ar.\allowbreak{}CptOrigin}{Development.\allowbreak{}AutomorphicSpectrum.\allowbreak{}Definitions.\allowbreak{}GlobalGeometry.\allowbreak{}IsPeriodAtRestrict}{20--55}
\leanref{GroundTheory.\allowbreak{}Heights₂.\allowbreak{}Heights.\allowbreak{}LowerBoundFree.\allowbreak{}HodgeData}{GroundTheory.\allowbreak{}HeightsAndTheta.\allowbreak{}Heights.\allowbreak{}HodgeLowerBound.\allowbreak{}FrameDefinitions.\allowbreak{}IsFrameOverLine}{60--75}
\leanref{GroundTheory.\allowbreak{}Heights₂.\allowbreak{}Heights.\allowbreak{}LowerBoundFree.\allowbreak{}MUc}{GroundTheory.\allowbreak{}HeightsAndTheta.\allowbreak{}Heights.\allowbreak{}HodgeLowerBound.\allowbreak{}FrameDefinitions.\allowbreak{}CpRegPhRegOfNonneg}{77--82}
\leanref{GroundTheory.\allowbreak{}Heights₂.\allowbreak{}Heights.\allowbreak{}LowerBoundFree.\allowbreak{}htMUc}{GroundTheory.\allowbreak{}HeightsAndTheta.\allowbreak{}Heights.\allowbreak{}HodgeLowerBound.\allowbreak{}FrameDefinitions.\allowbreak{}CpRegPhRegOfNonneg}{89--91}
\leanref{GroundTheory.\allowbreak{}Heights₂.\allowbreak{}Heights.\allowbreak{}LowerBoundFree.\allowbreak{}CuspBundlePos}{GroundTheory.\allowbreak{}HeightsAndTheta.\allowbreak{}Heights.\allowbreak{}HodgeLowerBound.\allowbreak{}FlagIndices}{62--67}
\leanref{GroundTheory.\allowbreak{}Heights₂.\allowbreak{}Heights.\allowbreak{}LowerBoundFree.\allowbreak{}HtInterior}{GroundTheory.\allowbreak{}HeightsAndTheta.\allowbreak{}Heights.\allowbreak{}HodgeLowerBound.\allowbreak{}FlagIndices}{68--72}
\leanref{GroundTheory.\allowbreak{}Heights₂.\allowbreak{}Heights.\allowbreak{}LowerBoundFree.\allowbreak{}HtBoundary}{GroundTheory.\allowbreak{}HeightsAndTheta.\allowbreak{}Heights.\allowbreak{}HodgeLowerBound.\allowbreak{}FlagIndices}{73--77}
\leanref{GroundTheory.\allowbreak{}Heights₂.\allowbreak{}Heights.\allowbreak{}LowerBoundFree.\allowbreak{}CuspMean}{GroundTheory.\allowbreak{}HeightsAndTheta.\allowbreak{}Heights.\allowbreak{}HodgeLowerBound.\allowbreak{}FlagIndices}{78--81}
\leanref{GlobalTheory.\allowbreak{}CompanionTheory.\allowbreak{}Ar.\allowbreak{}Top.\allowbreak{}StokesAt}{Development.\allowbreak{}AutomorphicSpectrum.\allowbreak{}Definitions.\allowbreak{}GlobalGeometry.\allowbreak{}HasBaseDegreeSettingStrict}{9--15}
\leanref{GlobalTheory.\allowbreak{}CompanionTheory.\allowbreak{}Ar.\allowbreak{}Top.\allowbreak{}ppc\_\allowbreak{}lower\_\allowbreak{}of\_\allowbreak{}setting}{Development.\allowbreak{}AutomorphicSpectrum.\allowbreak{}Proofs.\allowbreak{}DivisibilityProofs.\allowbreak{}SlopeLowerBaseSettings}{76--109}

The assembly uses the same tower throughout, with base slope constant \(\max(C_h+C_u/2,0)\), flat-height constant \(4\), and final tower constant \(2\max(C_h+C_u/2,0)+3\cdot4+5\). It supplies the proportionality, equal component volumes, averaged intersection estimate and mean-slope bound used below.
\leanref{GlobalTheory.\allowbreak{}CompanionTheory.\allowbreak{}Ar.\allowbreak{}Cusp.\allowbreak{}volumeTheoremX\_\allowbreak{}of\_\allowbreak{}inputs}{Development.\allowbreak{}AutomorphicSpectrum.\allowbreak{}Proofs.\allowbreak{}DivisibilityProofs.\allowbreak{}VolumeTheoremOfInputs}{72--122}
\end{enumerate}
\leanref{GlobalTheory.\allowbreak{}T.\allowbreak{}given\_\allowbreak{}of\_\allowbreak{}all\_\allowbreak{}inputs}{Development.\allowbreak{}AutomorphicSpectrum.\allowbreak{}Proofs.\allowbreak{}GeometryProofs.\allowbreak{}GivenOfAllInputsStatement}{27--29}
\leanref{GlobalTheory.\allowbreak{}T.\allowbreak{}flat\_\allowbreak{}zero\_\allowbreak{}from\_\allowbreak{}resolution\_\allowbreak{}and\_\allowbreak{}gaga}{Development.\allowbreak{}AutomorphicSpectrum.\allowbreak{}Proofs.\allowbreak{}GeometryProofs.\allowbreak{}FlatZeroFromResolutionAndGagaStatement}{8--10}
\leanref{GlobalTheory.\allowbreak{}T.\allowbreak{}flat\_\allowbreak{}first\_\allowbreak{}from\_\allowbreak{}resolution\_\allowbreak{}and\_\allowbreak{}gaga}{Development.\allowbreak{}AutomorphicSpectrum.\allowbreak{}Proofs.\allowbreak{}GeometryProofs.\allowbreak{}FlatFirstFromResolutionAndGagaStatement}{8--10}
\leanref{GlobalTheory.\allowbreak{}CompanionTheory.\allowbreak{}Ar.\allowbreak{}Flat.\allowbreak{}FlatLineGenerator}{Development.\allowbreak{}AutomorphicSpectrum.\allowbreak{}Statements.\allowbreak{}GeometryStatements.\allowbreak{}FlatLineGenerator}{18--20}
\leanref{GlobalTheory.\allowbreak{}CompanionTheory.\allowbreak{}Ar.\allowbreak{}Flat.\allowbreak{}FlatOnModel}{Development.\allowbreak{}AutomorphicSpectrum.\allowbreak{}Definitions.\allowbreak{}GlobalGeometry.\allowbreak{}IsFlatAnIffIsConstMul}{164--167}
The complete proof is contained in \Lean{GlobalTheory.\allowbreak{}CompanionTheory.\allowbreak{}Ar.\allowbreak{}Link.\allowbreak{}good\_\allowbreak{}tower\_\allowbreak{}and\_\allowbreak{}volume\_\allowbreak{}final}.
\end{proof}

\subsection*{Proof of \ref{B8.26}}\label{prf:B8.26}\phantomsection
\begin{proof}
By Lemmas~\ref{B8.26p} and \ref{B8.26q}, the first eight premises of Theorem~\ref{B8.25} furnish the arithmetic intersection system \(II\), the horizontal curve height bounds, unique extension of semiabelian homomorphisms over normal bases, Fargues--Hodge--Tate exactness with sharp cokernel annihilation, clopen GAGA section algebraisation, Hopf-order duality, finite locally free quotients of polarised abelian schemes over affine bases and Zarhin's polarised retraction. Together with Theorem~\ref{B8.27} and nef Riemann--Roch \eqref{eq:p8c3-nef-rr}, these are the inputs of Proposition~\ref{B8.26u}, which supplies the good root tower over \(E\) in the sense of Definition~\ref{B8.26r}.
\leanref{GlobalTheory.\allowbreak{}CompanionTheory.\allowbreak{}Ar.\allowbreak{}Giv.\allowbreak{}hsLinks\_\allowbreak{}of\_\allowbreak{}use}{Development.\allowbreak{}WeilBound.\allowbreak{}CoreStatements.\allowbreak{}AutomorphicKernels.\allowbreak{}MoriwakiBigness.\allowbreak{}HilbertSamuelApplications.\allowbreak{}HsLinksOfUse}{9--20}
\leanref{GlobalTheory.\allowbreak{}CompanionTheory.\allowbreak{}Ar.\allowbreak{}Link.\allowbreak{}given\_\allowbreak{}of\_\allowbreak{}use}{Development.\allowbreak{}WeilBound.\allowbreak{}CoreStatements.\allowbreak{}AutomorphicKernels.\allowbreak{}MoriwakiBigness.\allowbreak{}HilbertSamuelApplications.\allowbreak{}GivenOfUse}{11--18}
\leanref{GlobalTheory.\allowbreak{}CompanionTheory.\allowbreak{}Ar.\allowbreak{}Link.\allowbreak{}good\_\allowbreak{}tower\_\allowbreak{}volume\_\allowbreak{}theorem\_\allowbreak{}bis}{Development.\allowbreak{}WeilBound.\allowbreak{}CoreStatements.\allowbreak{}AutomorphicKernels.\allowbreak{}MoriwakiBigness.\allowbreak{}HilbertSamuelApplications.\allowbreak{}GoodTowerVolumeTheoremBis}{9--20}

At each rational prime \(\ell\) use the prime-to-\(\ell\) part of the base level, which is at least three because the base level \(g_0g_1\) is a product of two distinct odd primes exceeding three, together with the auxiliary prime different from \(\ell\). The resulting localisation moduli generate the unit ideal of \(\Z\). Choose local ample PEL model systems over this open cover, pass to a common dominating projective model and normalise in the smooth generic cover over \(F\). By Lemma~\ref{B8.26q}(2), the local projections composed with the semiabelian comparisons agree with the prescribed prime-to-\(\ell\) level maps and their level sections across the normal flat projective global model.
\leanref{GlobalTheory.\allowbreak{}Models.\allowbreak{}exists\_\allowbreak{}localSystems\_\allowbreak{}ample}{Development.\allowbreak{}AbelianVarietyTheory.\allowbreak{}Proofs.\allowbreak{}ModuliProofs.\allowbreak{}ExistsLocalModelSystemsAmpleFrames}{30--48}
\leanref{GlobalTheory.\allowbreak{}GlobalModel.\allowbreak{}exists\_\allowbreak{}global\_\allowbreak{}model}{Development.\allowbreak{}AbelianVarietyTheory.\allowbreak{}Proofs.\allowbreak{}ModuliProofs.\allowbreak{}ExistsGlobalModelOfLocalModels}{155--232}
\leanref{GlobalTheory.\allowbreak{}CompanionTheory.\allowbreak{}Ar.\allowbreak{}Rec.\allowbreak{}exists\_\allowbreak{}isGlobal}{Development.\allowbreak{}AutomorphicSpectrum.\allowbreak{}Proofs.\allowbreak{}GeometryProofs.\allowbreak{}ExistsIsGlobal}{33--39}

The good cover datum identifies the algebraic components of \(Y\) with the component set \(\operatorname{Comp}\) of the associated good root tower \(\mathscr T\) of split index \(p_{\rm sig}=k+j-1\) and rank \(n=p_{\rm sig}+q_{\rm sig}\) in Definition~\ref{B8.26r}. From \(m\ge2\), \(m^2\le k\) and \(1\le j<m\) one has \(q_{\rm sig}=m-j\ge1\), \(2m\le m^2\le k\), \(3q_{\rm sig}+1\le n\) and \(n\ge4\), so \(n^2\ge16\), while \(|\Psi|=2n_{\rm ord}\ge2\) for the host order parameter \(n_{\rm ord}\ge1\). Every prime in the bad set \(\mathcal B_{\mathscr T}\) lies in the set \(S\) of prime factors of \(2|\operatorname{disc}E|M'\). Splitting \(S\) at \(n_{\rm ord}\), the host discriminant bound and \(\#\{\ell\mid M'\}\le15\) leave at most \(15+3=18\) primes of \(S\) above \(n_{\rm ord}\), whereas at most \(n_{\rm ord}+1\) primes lie in \([0,n_{\rm ord}]\). Since \(19+n_{\rm ord}\le128n_{\rm ord}=4|\Psi|\cdot16\le4|\Psi|n^2\), the cutoff \(B=B(v,g)\ge4|\Psi|n^2\) satisfies
\begin{equation}\label{eq:p8c3-deep-cutoff}
 B\ge128,\qquad
 \#\mathcal B_{\mathscr T}\le4|\Psi|n^2\le N=\lfloor B\rfloor .
\end{equation}
Let \(a_Z(v;N)\) be the component model slopes of Definition~\ref{B8.26r}.
\leanref{GlobalTheory.\allowbreak{}CompanionTheory.\allowbreak{}Ar.\allowbreak{}Plug.\allowbreak{}vs\_\allowbreak{}tower}{Development.\allowbreak{}AutomorphicSpectrum.\allowbreak{}Proofs.\allowbreak{}GeometryProofs.\allowbreak{}FullIntersectionGoodLevel}{18--51}
\leanref{Formalism.\allowbreak{}BasicNotions.\allowbreak{}Host.\allowbreak{}card\_\allowbreak{}bad\_\allowbreak{}gt\_\allowbreak{}le}{Development.\allowbreak{}AutomorphicArithmetic.\allowbreak{}PrimeDiscriminants.\allowbreak{}BadPrimes}{74--78}

Let \(d_0\) be the generic flag-cover degree on \(Z\), and fix \(\eta>0\). Set \(t'=t-2c_2(B)\|v\|_1-\eta\), and choose a cuspidal generic trace basis \(f_1,\ldots,f_{d_0}\) of common degree \(k_0\) divisible by \(8\). For any positive integer \(M_j\), nondegeneracy of the field trace pairing makes \(s\mapsto(M_j\operatorname{Tr}(sf_a))_{a=1}^{d_0}\) injective. In degree \(j+k_0\) divisible by \(8\), passing from the tower root model at cutoff \(N\) (which uses Hodge--Tate frames at good primes \(\ell\le N\) and integral model valuations elsewhere) to the patched valuation family of \(\nu_Z(v;B)\) (which uses Hodge--Tate valuations above \(S_B\) and record valuations outside \(S_B\)), together with clearing the \(S_B\)-frame denominators and the fixed denominators of \(f_1,\ldots,f_{d_0}\) by a positive integer \(M_j\) with \(\log M_j\le2(j+k_0)c_2(B)\|v\|_1+O(1)\) and bounding the archimedean factors of \(f_1,\ldots,f_{d_0}\), increases logarithmic archimedean norms by at most \(2(j+k_0)c_2(B)\|v\|_1+O(1)\), where the local Hodge--Tate comparison at primes in \(S_B\) uses Lemma~\ref{B8.26q}(3). For all large \(j\) divisible by \(8\), each cleared coordinate \(M_j\operatorname{Tr}(sf_a)\) of a model section \(s\) of degree \(j\) and logarithmic sup norm at most \(jt'\) therefore lies among the \(n_{j+k_0}(t)\) small cusp sections on \(Z\). Hence the model section count is at most \(n_{j+k_0}(t)^{d_0}\); taking logarithms cancels the extra factor \(d_0\) in the model degree normalisation, and the fixed shift \(k_0\) leaves the asymptotic leading ratio unchanged. Nef Riemann--Roch \eqref{eq:p8c3-nef-rr}, degree pushforward and Theorem~\ref{B8.27} bound the model volume below by \(\deg_Z(v)(a_Z(v;N)+t')\). Letting \(\eta\downarrow0\) and taking the infimum over \(t\in\R\) gives
\begin{equation}
 a_Z(v;N)-2c_2(B)\|v\|_1\le\nu_Z(v;B).
 \tag{8.26d}\label{eq:p8c3-model-mixed}
\end{equation}
\leanref{GlobalTheory.\allowbreak{}CompanionTheory.\allowbreak{}Ar.\allowbreak{}Plug.\allowbreak{}SlopeVolM.\allowbreak{}slope\_\allowbreak{}sub\_\allowbreak{}cost\_\allowbreak{}le\_\allowbreak{}compVolSlope}{Development.\allowbreak{}AutomorphicSpectrum.\allowbreak{}Proofs.\allowbreak{}GeometryProofs.\allowbreak{}UniformAllNormHeightBound}{300--372}

For \(N\in\N\), \(N\ge1\), the prime sum satisfies
\begin{equation}
 \sum_{\ell\le N}\frac{\log\ell}{\ell-1}\le\log N+2.2.
 \tag{8.26f}\label{eq:p8c3-prime-sum}
\end{equation}
At \(N=1\) the sum is empty. For \(N\ge2\), \(N/\ell\le v_\ell(N!)+1\) gives \(N\sum_{\ell\le N}(\log\ell)/\ell\le\log N!+\vartheta(N)\), where \(\vartheta(t)=\sum_{\ell\le t}\log\ell\le t\log4\) from \(\binom{2m}{m}\le4^m\) and \(\log N!\le(N+1)\log N-N+1\) by integration. In \(\sum_{\ell\le N}(\log\ell)/(\ell(\ell-1))\), bound \(\log\ell/\log2\) at \(\ell=2,3,5,7,11,13,17,19,23,29\) by \(1,2,3,3,4,4,5,5,5,5\), and bound the tail over integers \(a\ge31\) via integration by parts and \(1/(t(t+2))\le t^{-2}\) by \(\int_{29}^\infty\log(t+2)t^{-2}\,dt\le(\log31+1)/29\). Combining with \((\log N+1)/N\le(\log2)/N+1/2\le(\log2+1)/2\) for \(N\ge2\) gives \eqref{eq:p8c3-prime-sum}.
\leanref{ReciprocityTheory.\allowbreak{}ChartRestriction.\allowbreak{}G.\allowbreak{}H1.\allowbreak{}primeLogSum\_\allowbreak{}nat}{ReciprocityTheory.\allowbreak{}ReductiveGroups.\allowbreak{}ArithmeticGeometry.\allowbreak{}FirstCohomology.\allowbreak{}PrimeLogSum}{270--353}

By \eqref{eq:p8c3-deep-cutoff}, \(S_B\setminus\{\ell\le N\}\) contains at most \(\#S\le4|\Psi|n^2\le N\) primes, each strictly exceeding \(N\) and hence contributing at most \(\log(N+1)/N\) to \(\sum_{\ell\in S_B}(\log\ell)/(\ell-1)\). The discriminant valuations in \eqref{eq:p8c3-mixed-places} sum to \(40\log\operatorname{rd}(E)\), and \(\ell=2\) adds \(\log2\). Thus
\begin{equation}
 0\le c_2(B)\le60(\log B+\log\operatorname{rd}(E)).
 \tag{8.26e}\label{eq:p8c3-cost2}
\end{equation}
Set \(s_Z=a_Z(v;N)-2c_2(B)\|v\|_1\). Since \(B<N+1\), Lemma~\ref{B8.26a}, Lemma~\ref{B8.26b} and Proposition~\ref{B8.26c} together with \eqref{eq:p8c3-model-mixed} and \eqref{eq:p8c3-cost2} give \eqref{eq:p8c3-deep} with the absolute constant \(C_{\rm deep}=C_T+120=2C_h+141\ge0\), where \(C_h\ge0\) is the absolute constant of Lemma~\ref{B8.26a}, independent of the datum, chain, layer and cover.
\end{proof}

\subsection*{Proof of \ref{B8.26s}}\label{prf:B8.26s}\phantomsection
\begin{proof}
In (1), alternating \(L_v^d\) across \(S_n\) at each embedding produces the Vandermonde product divided by \(\prod_{a<b}(b-a)\) with multinomial prefactor \(d!\); for \(a<b\) each ratio \((v_{\tau,a}-v_{\tau,b})/(b-a)\) is positive under regular dominance and nonnegative under dominance.

In (2), the quotient of \(\R[X_{\tau,i}]\) by the symmetric power sums \(\{p_{\tau,k}:1\le k\le n\}\) is the cohomology ring of the product of full flags of dimension \(d\), so its degree-\((d+1)\) component is zero; taking homogeneous parts yields \(F=\sum_{\tau,k}p_{\tau,k}Q_{\tau,k}\). Full symmetric polynomials of positive degree commute with alternation and yield zero under \(D\). For \(\sigma\ne\tau\), \(\partial_{\sigma,i}p_{\tau,k}=0\) and \(D(p_{\tau,k}\partial_{\sigma,i}Q)=0\). For \(\sigma=\tau\), \(\sum_{i<r}\partial_{\tau,i}(p_{\tau,k}Q)=kp_{\tau,k-1}^{<r}Q+p_{\tau,k}\sum_{i<r}\partial_{\tau,i}Q\), whose second term vanishes under \(D\), while \(\sum_{i=0}^{n-1}\partial_{\tau,i}(p_{\tau,k}Q)=kp_{\tau,k-1}Q+p_{\tau,k}\sum_i\partial_{\tau,i}Q\) vanishes under \(D\) for \(k\ge2\) and equals \(nD(Q)\) for \(k=1\). Taking \(Q_1=(-1)^d\prod_{\sigma,i}X_{\sigma,i}^{n-1-i}\) gives \(D(Q_1)=1\), and reducing the exponent of \(X_{\tau_0,r-1}\) by one gives \(Q_2\) with \(X_{\tau_0,r-1}Q_2=Q_1\) and repeated exponents \(n-r-1\) (at indices \(r-1\) and \(r\)) in \(X_{\tau_0,i}Q_2\) for \(i<r-1\), whence \(D(p_{\tau_0,1}^{<r}Q_2)=1\).

In (3), when \(k=1\), \(P_{\ell,\tau,h+1,0}=\operatorname{rk}_{\rm Lie}(\operatorname{Fr}_{\ell,h}(\tau))=\operatorname{rk}_{\rm Lie}(\operatorname{Fr}_{\ell,h}(\tau_0))\) by clause~4 of Definition~\ref{B8.26r}. Hence \(\log\ell\sum_{h\ge0}\ell^{-(h+1)}D(P_{\ell,\tau,h+1,0}Q)=n\kappa_\ell D(Q)=\kappa_\ell T(p_{\tau,1}Q)\) and \(-\frac{\log\ell}{\ell-1}D(p_{\tau,0}^{<r}Q)=-\frac{\log\ell}{\ell-1}rD(Q)=-\frac{\log\ell}{\ell-1}S(p_{\tau,1}Q)\) by (2). When \(k\ge2\), \(T(p_{\tau,k}Q)=0\) and \(S(p_{\tau,k}Q)=kD(p_{\tau,k-1}^{<r}Q)\), giving the same formula with \(\rho_\ell(\tau,k,Q)\). Since \(P_{\ell,\tau,h,k-1}\) ranges over the finite set of \(2f\) signed block power sums, \(|D(P_{\ell,\tau,h,k-1}Q)|\) is bounded uniformly in \((\ell,h)\), and \(\sum_{h\ge1}\ell^{-hk}=\ell^{-k}/(1-\ell^{-k})\le2\ell^{-2}\) for \(k\ge2\) and \(\ell\ge2\).

In (4), orthogonal splitting of the standard Hermitian filtration on \(\mathcal O^n\) identifies \(\bigotimes_{i=0}^{n-1}\overline Q_{\tau,i}\simeq(\overline{\det\mathcal O^n})^\vee\) isometrically with the trivial Hermitian line of constant norm \(1\), whose standard basis generator has empty divisor and zero log norm; the rational-section divisor recursion therefore gives \(h_{\rm fl}(p_{\tau,1}Q)=0\) for every homogeneous \(Q\) of degree \(d\). Summation by parts rewrites \(L_v\) modulo \(\operatorname{span}\{p_{\tau,1}:\tau\in\Psi\}\) as \(\sum_{\tau\in\Psi}\sum_{j=1}^{n-1}(v_{\tau,j-1}-v_{\tau,j})\sum_{i=j}^{n-1}X_{\tau,i}\), where \(\sum_{i=j}^{n-1}X_{\tau,i}\) represents the pullback of \(\overline{\mathcal O_{\operatorname{Gr}(n-j,n)}(1)}\). Since \(\overline{\mathcal O_{\operatorname{Gr}(n-j,n)}(1)}\) is smooth, conjugation-invariant, semipositive, fibrewise nef and generated by integral Pl\"ucker coordinates of sup norm at most \(1\), clearing denominators for rational dominant \(v\), applying arithmetic Kleiman \eqref{eq:p8c3-kleiman} and extending by continuity to all real dominant \(v\) yields \(h_{\rm fl}(L_v^{d+1})\ge0\). When \(v_{\tau_0,a}=v_{\tau_0,a+1}\) for some \(\tau_0\in\Psi\) and \(0\le a<n-1\), the summand \(j=a+1\) at \(\tau_0\) vanishes, so for rational \(v\) on that wall (and by continuity for all real \(v\) there) the class pulls back from the smooth projective partial flag variety over \(\Z\) omitting \(E_{\tau_0,a+1}\), which has arithmetic dimension \(d\) and vanishing degree-\((d+1)\) arithmetic Chow group \cite[Th.~3.2.3]{bost-gillet-soule-1994}, giving \(h_{\rm fl}(L_v^{d+1})=0\).
\end{proof}

\subsection*{Proof of \ref{B8.26a}}\label{prf:B8.26a}\phantomsection
\begin{proof}
Work on the base Grassmannian model of generic dimension \(e=frs\) on a component \(Z\in\operatorname{Comp}\) provided by clause~6 of Definition~\ref{B8.26r}. Since \(3s+1\le n=r+s\), we have \(s<r\), so \(\min(r,s)=s\) and the smaller Hodge block at each embedding \(\tau\in\Psi\) has rank \(s\), with determinant class \(\mu=\sum_{\tau,i\ge r}X_{\tau,i}\). Regularise the singular Hodge metric on \(\mu\) by adding \(\tfrac12 U\), where the cusp weight \(U\ge0\) is defined at a uniformised frame of positive-majorant first minimum \(\lambda_1\) as the boundary logarithmic degeneration weight satisfying \eqref{eq:p8c3-cusp-point} below, and perturb by \((a/b)\) times the auxiliary ample line for integers \(1\le a\le b\). At finite places \(\ell\le N\) or \(\ell\in\mathcal B_{\mathscr T}\), equip the bundle with the adapted determinant metric given at good primes \(\ell\le N\) by the canonical-subgroup and Hodge--Tate frames of Lemma~\ref{B8.26q}(3) \cite{fargues-2011}. At a closed horizontal point choose independent period vectors spanning a finite-index \(\mathcal O_E\)-submodule, with \(t\le s\) projected vectors in the rank-\(s\) eigenspaces, and complete their projections to bases there. Index and determinant comparison bound the labelled fibre degree above by
\begin{equation}\label{eq:p8c3-label-degree}
 fsX-ft\log\operatorname{rd}(E)+x+\tfrac12ft\log(4\pi),\qquad
 X=2\sum_{\ell\le N\ {\rm or}\ \ell\in\mathcal B_{\mathscr T}}
 \frac{\log\ell}{\ell-1}+2\log\operatorname{rd}(E).
\end{equation}
Here \(x\) is the sum over \(\overline\Psi\) of the Arakelov degrees of the saturated rank-\(s\) Lie eigenlattices; the image lattices have the same rank and degree at most that of their saturations. Write \(y\) for the conjugate sum and \(h\) for the Faltings height (minus the total integral Lie degree) of the abelian fibre of dimension \(fn\). Determinant inclusion, conjugate degree comparison across \(\mathcal O_E\) and Bost's intrinsic Arakelov slope inequality for subbundles of Lie algebras of abelian schemes \cite[Th.~I]{bost-1996} give an absolute constant \(c\ge0\) with
\begin{equation}\label{eq:p8c3-eigen-degree}
 x+y\le-h,\quad |x-y|\le cfn\lambda_N,\quad
 x\le fsc\log(2+fn+\max(h,0)).
\end{equation}
If \(x>0\), adding the first two bounds gives \(2x-cfn\lambda_N\le-h\), so \(h\le cfn\lambda_N\). Since \(fn\le N\le e^{\lambda_N}\), in both cases \(x\le C_afs\lambda_N\) for an absolute constant \(C_a\ge0\). Equation~\eqref{eq:p8c3-prime-sum}, \(\#\mathcal B_{\mathscr T}\le N\) and \(t\le s\) bound \eqref{eq:p8c3-label-degree} by \(C_hfs\lambda_N\) with the absolute constant \(C_h\ge0\). Adding \(\tfrac12 U\) dominates the boundary logarithmic singularity of the adapted determinant metric, so on both interior and boundary horizontal curves the normalised height satisfies
\begin{equation}\label{eq:p8c3-point-lower}
 h_{\rm reg}\ge-C_hfs\lambda_N-(a/b)K_A,
\end{equation}
where \(-K_A\) is a lower bound for the auxiliary ample height on the component model.
\leanref{LevelTowerAlt.\allowbreak{}height\_\allowbreak{}comparison\_\allowbreak{}global\_\allowbreak{}model}{Development.\allowbreak{}WeilBound.\allowbreak{}CoreStatements.\allowbreak{}PlurisubharmonicCones.\allowbreak{}HeightComparisonGlobalModel}{80--245}
\leanref{GlobalTheory.\allowbreak{}CompanionTheory.\allowbreak{}Ar.\allowbreak{}Cusp.\allowbreak{}TopHt.\allowbreak{}htU\_\allowbreak{}of\_\allowbreak{}given}{Development.\allowbreak{}AutomorphicSpectrum.\allowbreak{}Proofs.\allowbreak{}GeometryProofs.\allowbreak{}CuspHeightUpper}{120--223}

To bound \(\operatorname{mean}U\), apply the Euclidean transference bound in the proof of Lemma~\ref{B8.25f} to \(\mathcal O_E\) of degree \(2f\) to obtain \(2f\) rationally independent algebraic integers \(\omega_1,\ldots,\omega_{2f}\in\mathcal O_E\) of Euclidean length at most \(v_0=2f\operatorname{rd}(E)\), so every conjugate satisfies \(|\sigma\omega_i|\le v_0\). For \(t>0\), let \(N_t\) count the nonzero isotropic vectors \(v\) in the arithmetic lattice \(\Lambda\) of positive-majorant length below \(t\) that satisfy \([\Lambda\cap Ev:\mathcal O_Ev]\le\lfloor(v_0\sqrt{2f})^{2f}\rfloor\). For any two vectors \(v,w\in\Lambda\) of length below \(t\le v_0^{-1/2}\), the \(\Z\)-valued trace pairings with \(\omega_1w,\ldots,\omega_{2f}w\) have modulus at most \(v_0\|v\|\|w\|<1\) and therefore vanish; since \(\omega_1,\ldots,\omega_{2f}\) span \(E\) over \(\Q\), all vectors in \(\Lambda\) of length below \(t\) are isotropic and mutually orthogonal for the Hermitian form, spanning a single isotropic \(E\)-subspace \(W_0\) of \(E\)-dimension at most \(s\) and real dimension \(d_0\le2fs\). For any shortest nonzero vector \(v\in\Lambda\) with \(\|v\|=\lambda_1<t\le v_0^{-1/2}\), the full-rank sublattice \(\bigoplus_{i=1}^{2f}\Z\omega_iv\subseteq\mathcal O_Ev\) has covolume at most \((v_0\lambda_1)^{2f}\) by Hadamard's inequality, while packing the half-cube \((-\lambda_1/(2\sqrt{2f}),\lambda_1/(2\sqrt{2f}))^{2f}\) in orthonormal coordinates on \((Ev)\otimes_\Q\R\) gives \(\operatorname{covol}(\Lambda\cap Ev)\ge(\lambda_1/\sqrt{2f})^{2f}\); dividing bounds \([\Lambda\cap Ev:\mathcal O_Ev]\le\lfloor(v_0\sqrt{2f})^{2f}\rfloor\), so \(v\) belongs to the counted set and \(N_t\ge1\). Since this set is invariant under the arithmetic unitary group, unfolding \(N_t\) over a fundamental set \(\mathcal F\) of finite positive Haar measure expresses \(\int_{\mathcal F}N_t\) as a sum over \(\Gamma\)-orbits of local orbital integrals on the isotropic cone, where scaling the length cutoff by \(\theta>0\) scales the invariant isotropic-cone measure of homogeneity degree \(2f(n-1)\) by \(\theta^{2f(n-1)}\). Normalise Haar measure on \(\mathcal F\) to a probability measure and put \(I_t=\mathbb E N_t<\infty\).
\leanref{GlobalTheory.\allowbreak{}CompanionTheory.\allowbreak{}Ar.\allowbreak{}Cusp.\allowbreak{}group\_\allowbreak{}mean\_\allowbreak{}logPlus\_\allowbreak{}le}{Development.\allowbreak{}AutomorphicSpectrum.\allowbreak{}Proofs.\allowbreak{}MetricProofs.\allowbreak{}FirstMinimumIntegralBound}{220--307}

Inside the \(d_0\)-dimensional Euclidean space \(W_0\), open balls of radius \(\lambda_1/2\) around distinct lattice points of length below \(t\) are disjoint and lie in the ball of radius \(\tfrac32 t\), giving \(1\le N_t\le(3t/\lambda_1)^{d_0}\le(3t/\lambda_1)^{2fs}\) on \(\lambda_1<t\le v_0^{-1/2}\), whereas \(N_t=0\) on \(\lambda_1\ge t\). At \(t=v_0^{-1/2}\), split \(I_t\) at \(\lambda_1=t/12\). The region \(\lambda_1\ge t/12\) contributes at most \(36^{2fs}\). On \(\lambda_1<t/12\), put \(a_*=2fs\) and \(k_*=2f(n-1)\); homogeneous scaling gives \(\Pr(\lambda_1<\theta t)\le\mathbb E N_{\theta t}=\theta^{k_*}I_t\) for \(0<\theta\le1\). Layer-cake integration of \((3t/\lambda_1)^{a_*}\) over \(u=t/\lambda_1>12\) bounds the remaining contribution by
\[
 3^{a_*}\left(12^{a_*-k_*}+a_*\int_{12}^\infty u^{a_*-k_*-1}\,du\right)I_t
 =3^{2fs}12^{-2f(n-1-s)}\left(1+\frac{s}{n-1-s}\right)I_t
 \le\tfrac12 I_t ,
\]
since \(n-1-s=r-1\ge2s\ge2\). Subtracting \(\tfrac12 I_t\) gives \(I_t\le2\cdot36^{2fs}\), so
\begin{equation}\label{eq:p8c3-cusp-tail}
 \Pr(Y_1>u)\le\min(1,2\cdot36^{2fs}e^{-2f(n-1)u}),
 \qquad Y_1=\log^+(v_0^{-1/2}/\lambda_1),\quad u\ge0 .
\end{equation}
Integrating \eqref{eq:p8c3-cusp-tail} over \(u\ge0\) with \(n-1\ge3s\) gives \(\mathbb E Y_1\le3\). Using the pointwise cusp-weight bound
\begin{equation}\label{eq:p8c3-cusp-point}
 U\le fs\bigl(2\log^+(1/\lambda_1)+\log(2fs\operatorname{rd}(E))\bigr)
\end{equation}
and \(\log^+(1/\lambda_1)\le Y_1+\tfrac12\log v_0\), averaging over fundamental sets and sheets yields
\begin{equation}\label{eq:p8c3-cusp-mean}
 \operatorname{mean}U\le fs\bigl(6+\log v_0+\log(2fs\operatorname{rd}(E))\bigr)\le4fs\lambda_N ,
\end{equation}
since \(2fs<fn\le\frac14 N\) and \(N\ge128\) give \(6+\log(4f^2s\operatorname{rd}(E)^2)\le4\log(N\operatorname{rd}(E))\le4\lambda_N\).
\leanref{GlobalTheory.\allowbreak{}CompanionTheory.\allowbreak{}Ar.\allowbreak{}Cusp.\allowbreak{}fund\_\allowbreak{}set\_\allowbreak{}mean\_\allowbreak{}cuspW}{Development.\allowbreak{}AutomorphicSpectrum.\allowbreak{}Proofs.\allowbreak{}MetricProofs.\allowbreak{}CuspWeightIntegrableOnFundamentalSet}{231--329}

Now fix integers \(1\le a\le b\). Since \(N\ge4fn^2\) and \(n\ge4\), we have \(4f+2\le N\). By clause~6 of Definition~\ref{B8.26r}, there exist positive thresholds \(t_1,t_2,t_3,t_4>0\) such that for every \(0<t'\le\min(t_1,t_2,t_3,t_4)\) the cleared-denominator smoothed bundle \(\overline M_{N,a,b,t'}\) on the base model of \(Z\) is smooth, conjugation-invariant, cone-semipositive along the defining ideal of the base model and fibrewise nef, and satisfies \(h(\overline M_{N,a,b,t'};W)\ge b e_0\) on every horizontal arithmetic curve \(W\), where \(e_0=-C_hfs\lambda_N-(a/b)K_A\) by \eqref{eq:p8c3-point-lower}. Let \(D_{t'}=b^e\deg_{\rm gen}(\mu+(a/b)A)\) and \(I_{t'}\) be its generic degree and top arithmetic self-intersection. Lemma~\ref{B8.26p}(1) in dimension \(e\) gives
\begin{equation}\label{eq:p8c3-smooth-lower}
 I_{t'}\ge(e+1)D_{t'}\,be_0 .
\end{equation}
Divide \eqref{eq:p8c3-smooth-lower} by \((e+1)bD_{t'}>0\) and pass to the Stokes limit \(t'\downarrow0\), then let \(a/b\to0\). By clause~6 of Definition~\ref{B8.26r} and \eqref{eq:p8c3-cusp-mean} (which supplies the cusp-mean constant \(C_u=4\), so \(\tfrac12 C_u=2\)),
\[
 \frac{P_{N,Z}(F_*)}{(e+1)V_ZD(\mu^eJ)}\ge-\bigl(C_h+\tfrac12 C_u\bigr)fs\lambda_N=-(C_h+2)fs\lambda_N .
\]
Since \(H_*=(e+1)D(\mu^eJ)>0\) and \(V_Z>0\), multiplying by \(V_ZH_*\) and summing over \(Z\in\operatorname{Comp}\) yields \(P_N(F_*)\ge-VH_*(C_h+2)fs\lambda_N\).
\leanref{GlobalTheory.\allowbreak{}CompanionTheory.\allowbreak{}Ar.\allowbreak{}Top.\allowbreak{}slX\_\allowbreak{}lower\_\allowbreak{}of\_\allowbreak{}stokesAt}{Development.\allowbreak{}AutomorphicSpectrum.\allowbreak{}Proofs.\allowbreak{}DivisibilityProofs.\allowbreak{}SlopeLowerBaseSettings}{16--75}
\leanref{GlobalTheory.\allowbreak{}CompanionTheory.\allowbreak{}Ar.\allowbreak{}Top.\allowbreak{}ppc\_\allowbreak{}lower\_\allowbreak{}of\_\allowbreak{}setting}{Development.\allowbreak{}AutomorphicSpectrum.\allowbreak{}Proofs.\allowbreak{}DivisibilityProofs.\allowbreak{}SlopeLowerBaseSettings}{76--109}
\end{proof}

\subsection*{Proof of \ref{B8.26d}}\label{prf:B8.26d}\phantomsection
\begin{proof}
Write \(t=\min(r,s)\), \(w=\max(r,s)\), and let \(R=(w^t)\) be the partition with \(t\) rows of length \(w\). For a partition \(\lambda\), write \(f^\lambda\) for its number of standard tableaux. A partition \(\nu\) of \(k\) with at most \(t\) rows gives the shape
\[
 R+\nu=(w+\nu_1,\ldots,w+\nu_t)
\]
after zero rows have been appended to \(\nu\). In a standard tableau of this shape, record the set of entries in the last \(k\) boxes, meaning the boxes strictly to the right of column \(w\). Standardise the entries in those boxes and in the rectangle separately, preserving their relative orders. The set and the two standard tableaux recover every entry of the original tableau. Thus
\[
 f^{R+\nu}\le \binom{a+k}{k}f^R f^\nu.
\]
For the hook \(\nu=(k-j,1^j)\), the corner contains \(1\), and the \(j\) entries below it determine the whole tableau. Consequently
\[
 \sum_{\substack{\nu=(k-j,1^j)\\0\le j<\min(t,k)}}f^{R+\nu}
 \le \binom{a+k}{k} f^R
      \sum_{j=0}^{\min(t,k)-1}\binom{k-1}{j}
 \le 2^{k-1}\binom{a+k}{k} f^R .
\]
\leanref{GroundTheory.\allowbreak{}Heights₂.\allowbreak{}Heights.\allowbreak{}RestTerms.\allowbreak{}LemmaC.\allowbreak{}proof}{GroundTheory.\allowbreak{}HeightsAndTheta.\allowbreak{}Heights.\allowbreak{}RemainderTerms.\allowbreak{}PLemmaC}{19--28}
\leanref{GroundTheory.\allowbreak{}Heights₂.\allowbreak{}Heights.\allowbreak{}RestTerms.\allowbreak{}HookTableauCount.\allowbreak{}proof}{GroundTheory.\allowbreak{}HeightsAndTheta.\allowbreak{}Heights.\allowbreak{}RemainderTerms.\allowbreak{}PHookTableauCount}{140--145}
\leanref{GroundTheory.\allowbreak{}Heights₂.\allowbreak{}Heights.\allowbreak{}RestTerms.\allowbreak{}HookSumBound.\allowbreak{}proof}{GroundTheory.\allowbreak{}HeightsAndTheta.\allowbreak{}Heights.\allowbreak{}RemainderTerms.\allowbreak{}PHookSumBound}{83--129}

To evaluate the insertion, alternate in the \(n\) roots at each embedding and divide by the corresponding Vandermonde. On polynomials of top flag degree this gives \(D\), up to a common normalising factor. Applied to \(PF_*\), it gives a symmetric polynomial of degree \(k\). Take the coefficient of the single power sum \(p_{\tau k}\) in the power-sum coordinates at all embeddings. These coordinates are algebraically independent through index \(n\) over \(\Q\). Full power sums commute with alternation. In the decomposition of \(F_*\), a summand divisible by \(p_{\eta j}\) can therefore contribute to that coefficient only for \((\eta,j)=(\tau,k)\). The remaining factor gives the top-degree integral \(D(PQ_{\tau k})\), with the same normalisation.

Here is the block calculation of the coefficient. In one embedding let \(Y\) denote the \(t\) variables of the smaller block. On multiplying \(s_\lambda(Y)\) by the two within-block Vandermondes, the two lists of exponents in the alternants are
\[
 \lambda_1+t-1,\ldots,\lambda_t
 \quad\hbox{and}\quad w-1,\ldots,0 .
\]
If \(\lambda_t<w\), an exponent occurs twice. Otherwise division of the full alternant by the full Vandermonde gives \(s_{\lambda-R}\) in all \(n\) variables, up to a common factor from the two block alternations. Reordering the blocks or changing the signs of all the variables in the chosen block contributes a sign in each homogeneous degree. All common factors cancel on division by \(H_*\).

In the expansion of a product of power sums \(p_\rho(Y)\), the coefficient of \(s_\lambda(Y)\) is the symmetric-group character \(\chi^\lambda(\rho)\). In particular the coefficient for \(p_1(Y)^a\) and \(\lambda=R\) is \(f^R\). In the inverse expansion the coefficient of the single power sum \(p_k\) in \(s_\nu\) is \(\chi^\nu((k))/k\). The character of a \(k\)-cycle is \((-1)^j\) for \(\nu=(k-j,1^j)\) and zero otherwise. These are the Frobenius formula and its single-cycle case for partitions of \(k\); here \(k\le n\), so the coefficient can be read in \(n\) variables. For every cycle type \(\rho\) of size \(|\lambda|\) one has \(|\chi^\lambda(\rho)|\le f^\lambda\), since a finite-order operator on the representation has eigenvalues of absolute value one. See \cite[Ch.~I, \S7, in particular (7.8)]{macdonald-1995}.
\leanref{GroundTheory.\allowbreak{}Heights₂.\allowbreak{}Heights.\allowbreak{}RestTerms.\allowbreak{}SchurSignCoefficient.\allowbreak{}proof}{GroundTheory.\allowbreak{}HeightsAndTheta.\allowbreak{}Heights.\allowbreak{}RemainderTerms.\allowbreak{}PSchurSignCoefficient}{42--59}
\leanref{GroundTheory.\allowbreak{}Heights₂.\allowbreak{}Heights.\allowbreak{}RestTerms.\allowbreak{}CharacterBound.\allowbreak{}proof}{GroundTheory.\allowbreak{}HeightsAndTheta.\allowbreak{}Heights.\allowbreak{}RemainderTerms.\allowbreak{}PCharacterBound}{15--24}

In \(D(\mu^eJ)\) every embedding must supply \(a\) linear factors. Its multinomial and tableau factors are \(e!/(a!)^f\) and \((f^R)^f\). Suppose first that \(\sigma=\tau\). The input embedding must supply \(a+1\) linear factors from \(\mu^{e+1}\), in addition to the insertion of degree \(k-1\). Every other embedding supplies \(a\) linear factors. After division by the multinomial in \(H_*\), the ratio is
\[
 \frac{(e+1)!}{(a+1)!(a!)^{f-1}}
 \frac{(a!)^f}{(e+1)e!}=\frac1{a+1}.
\]
The remaining character at the input has shape \(R+\nu\) and cycle type \((k-1,1^{a+1})\). The factor \(k\) in the assertion cancels the denominator \(k\) in the single-cycle coefficient. The absolute value is therefore bounded by \(1/(a+1)\) times the hook sum above divided by \(f^R\).

If \(\sigma\ne\tau\), the input supplies \(a+k\) linear factors and the output supplies \(a-k+1\), besides its insertion of degree \(k-1\). Here \(k-1\le r+s-1\le rs=a\). The input coefficient before extraction is \(f^{R+\nu}\); the output character has shape \(R\) and absolute value at most \(f^R\). The multinomial ratio is
\[
 \frac{(a!)^2}{(a+k)!(a-k+1)!}.
\]
Multiplying this by \(2^{k-1}\binom{a+k}{k}\) gives the second bound.
\leanref{GroundTheory.\allowbreak{}Heights₂.\allowbreak{}Heights.\allowbreak{}RestTerms.\allowbreak{}SameFactorFormula.\allowbreak{}proof}{GroundTheory.\allowbreak{}HeightsAndTheta.\allowbreak{}Heights.\allowbreak{}RemainderTerms.\allowbreak{}PSameFactorFormula}{158--254}
\leanref{GroundTheory.\allowbreak{}Heights₂.\allowbreak{}Heights.\allowbreak{}RestTerms.\allowbreak{}OtherFactorFormula.\allowbreak{}proof}{GroundTheory.\allowbreak{}HeightsAndTheta.\allowbreak{}Heights.\allowbreak{}RemainderTerms.\allowbreak{}POtherFactorFormula}{41--193}
\end{proof}

\subsection*{Proof of \ref{B8.26b}}\label{prf:B8.26b}\phantomsection
\begin{proof}
Let \((\alpha_M),(\beta_M)\), \(C\ge0\) and \(N_1\ge2\) be given by clause~3 of Definition~\ref{B8.26r}. Fix \(c_0>0\) and a wall embedding \(\tau_0\in\Psi\). On the compact patch of weights \(\|v\|_\infty\le1/2\) where \(v_{\tau_0,r-1}=v_{\tau_0,r}\) and all other wall gaps and within-block adjacent gaps are at least \(c_0\), Lemma~\ref{B8.26s}(1) and~(4) give \(\delta_{\rm fl}(v)=0\) (hence \(S(L_v^{d+1})=T(L_v^{d+1})=0\)) and \(h_{\rm fl}(L_v^{d+1})=0\), so clause~3 gives
\begin{equation}\label{eq:p8c3-wall}
 |(P_M-Vh_{\rm fl})(L_v^{d+1})|=O(\log M/M)
\end{equation}
uniformly on each wall patch. (On the moduli tower of Proposition~\ref{B8.26}, \eqref{eq:p8c3-wall} is proved first: for rational \(v\) in this patch and large cutoff \(M\), shifting its first block at \(\tau_0\) by \(2\theta\) with \(\theta=5f/M\) and subtracting \(z\nu\) for the rational tower weight \(\nu\) and \(z=\theta^2\) gives geometric degree \(O(\theta)\), normalised intersection \(O(\theta\log M)\) at smoothing \(\exp(-1/z^2)\) with smoothing error \(O(z)\), and return cost \(O((2\theta+z\|\nu\|_\infty+z)\log M)=O(\log M/M)\) since the cutoff-\(M\) intersection coefficients are \(O(\log M)\), whence rational density and continuity give \eqref{eq:p8c3-wall}; within-block alternation and translated wall factors then identify the common kernel of restriction to these wall patches with the Vandermonde times block-symmetric linear forms, so polarisation of \(L_v^{d+1}\) and averaging over \(\Psi\) yield \((\alpha_M),(\beta_M)\).) Setting
\[
 G_M(F)=P_M(F)-V\bigl(h_{\rm fl}(F)-(\alpha_M-\beta_M)S(F)-\beta_MT(F)\bigr),
\]
linearity and the monomial bound of clause~3 give \(|G_M(F)|\le\bigl(\sum_m|c_m(F)|\bigr)C(\log M)/M\to0\) as \(M\to\infty\) for every homogeneous \(F\) of degree \(d+1\).
\leanref{GlobalTheory.\allowbreak{}CompanionTheory.\allowbreak{}Ar.\allowbreak{}Top.\allowbreak{}thmW\_\allowbreak{}tower}{Development.\allowbreak{}AutomorphicSpectrum.\allowbreak{}Proofs.\allowbreak{}DivisibilityProofs.\allowbreak{}TowerWeilTheorem}{120--204}
\leanref{GlobalTheory.\allowbreak{}CompanionTheory.\allowbreak{}Ar.\allowbreak{}Top.\allowbreak{}tendsto\_\allowbreak{}const\_\allowbreak{}mul\_\allowbreak{}log\_\allowbreak{}div}{Development.\allowbreak{}AutomorphicSpectrum.\allowbreak{}Proofs.\allowbreak{}DivisibilityProofs.\allowbreak{}TowerIdentityAtCutoff}{73--80}

By Lemma~\ref{B8.26s}(2), every homogeneous \(F\) of degree \(d+1\) decomposes as \(F=\sum_{\tau,k}p_{\tau,k}Q_{\tau,k}\) with \(1\le k\le n\) and \(Q_{\tau,k}\) homogeneous of degree \(d+1-k\). For \(M\ge N\ge2\), summing clause~4 of Definition~\ref{B8.26r} over \(Z\in\operatorname{Comp}\) and substituting Lemma~\ref{B8.26s}(3) gives
\[
 P_N(p_{\tau,k}Q)=V\Bigl(h_{\rm fl}(p_{\tau,k}Q)-a_MS(p_{\tau,k}Q)-b_MT(p_{\tau,k}Q)
 +\sum_{\substack{N<\ell\le M\ {\rm prime}\\\ell\notin\mathcal B_{\mathscr T}}}\rho_\ell(\tau,k,Q)\Bigr)+G_M(p_{\tau,k}Q),
\]
where, with \(\kappa_\ell\) as in Lemma~\ref{B8.26s}(3),
\[
 a_M=\alpha_M-\beta_M+\sum_{\substack{N<\ell\le M\ {\rm prime}\\\ell\notin\mathcal B_{\mathscr T}}}\frac{\log\ell}{\ell-1},\qquad
 b_M=\beta_M-\sum_{\substack{N<\ell\le M\ {\rm prime}\\\ell\notin\mathcal B_{\mathscr T}}}\kappa_\ell .
\]
For \(2\le k\le n\), Lemma~\ref{B8.26s}(3) shows that as \(M\to\infty\) the finite prime sums converge absolutely to
\begin{equation}\label{eq:p8c3-cutoff-tail}
 R_N(p_{\tau,k}Q)=\sum_{\substack{\ell>N\ {\rm prime}\\\ell\notin\mathcal B_{\mathscr T}}}
 k\log\ell\sum_{h\ge1}\ell^{-hk}D(P_{\ell,\tau,h,k-1}Q),
\end{equation}
while for \(k=1\) one has \(R_N(p_{\tau,1}Q)=0\). Fix \(\tau_0\in\Psi\) and take the two test classes \(p_{\tau_0,1}Q_1\) and \(p_{\tau_0,2}Q_2\) of Lemma~\ref{B8.26s}(2), on which \((S,T)\) equal \((r,n)\) and \((2,0)\), respectively. Since \(V>0\), evaluating the finite-\(M\) identity on \(p_{\tau_0,2}Q_2\) (where \(T=0\) and \(S=2\ne0\)) and letting \(M\to\infty\) proves \(a_M\to A\) for a finite limit \(A\in\R\); evaluating next on \(p_{\tau_0,1}Q_1\) (where \(T=n\ne0\)) proves \(b_M\to b\) for a finite limit \(b\in\R\). Passing to the limit \(M\to\infty\) on each summand \(p_{\tau,k}Q_{\tau,k}\) and summing over \((\tau,k)\) establishes \eqref{eq:p8c3-cutoff} with \(R_N(F)=\sum_{\tau,k}R_N(p_{\tau,k}Q_{\tau,k})\).
\leanref{GlobalTheory.\allowbreak{}CompanionTheory.\allowbreak{}Ar.\allowbreak{}Top.\allowbreak{}exists\_\allowbreak{}finite\_\allowbreak{}identity}{Development.\allowbreak{}AutomorphicSpectrum.\allowbreak{}Proofs.\allowbreak{}DivisibilityProofs.\allowbreak{}TowerIdentityAtCutoff}{97--150}
\leanref{GlobalTheory.\allowbreak{}CompanionTheory.\allowbreak{}Ar.\allowbreak{}Top.\allowbreak{}tendsto\_\allowbreak{}coeffs\_\allowbreak{}of\_\allowbreak{}tests}{Development.\allowbreak{}AutomorphicSpectrum.\allowbreak{}Proofs.\allowbreak{}DivisibilityProofs.\allowbreak{}TowerIdentityAtCutoff}{37--60}
\leanref{GlobalTheory.\allowbreak{}CompanionTheory.\allowbreak{}Ar.\allowbreak{}Top.\allowbreak{}eq\_\allowbreak{}of\_\allowbreak{}tendsto\_\allowbreak{}coeffs}{Development.\allowbreak{}AutomorphicSpectrum.\allowbreak{}Proofs.\allowbreak{}DivisibilityProofs.\allowbreak{}TowerIdentityAtCutoff}{61--72}
\leanref{GlobalTheory.\allowbreak{}CompanionTheory.\allowbreak{}Ar.\allowbreak{}Top.\allowbreak{}identOnDecomposedForm}{Development.\allowbreak{}AutomorphicSpectrum.\allowbreak{}Proofs.\allowbreak{}DivisibilityProofs.\allowbreak{}TowerHeightIdentityExpansion}{14--16}

On the degree-one test \(p_{\tau_0,1}Q_1\), Lemma~\ref{B8.26s}(4) gives \(h_{\rm fl}(p_{\tau_0,1}Q_1)=0\), while \(R_N(p_{\tau_0,1}Q_1)=0\), \((S,T)=(r,n)\) and \(D(Q_1)=1\). By clause~5 of Definition~\ref{B8.26r}, since \(N\ge\max(128,\#\mathcal B_{\mathscr T})\) and \(D(Q_1)=1\), one has \(P_{N,Z}(p_{\tau_0,1}Q_1)=h_{N,Z,\tau_0}V_Z\) with \(|h_{N,Z,\tau_0}|\le4n\lambda_N\) (on the moduli tower of Proposition~\ref{B8.26}, \(h_{N,Z,\tau_0}=h_{\rm rec}+\sum_\ell d_\ell\) decomposes into the global flat record term and the local canonical-subgroup discrepancies at primes \(\ell\le N\) or \(\ell\in\mathcal B_{\mathscr T}\) \cite{fargues-2011}, which for \(u=\log\operatorname{rd}(E)\) satisfy
\begin{equation}\label{eq:p8c3-degree-one}
 |h_{\rm rec}|\le n(3u+|L_E|)/2,\quad
 |d_\ell|\le n((\ell-1)^{-1}+v_\ell)\log\ell,\quad
 \sum_\ell v_\ell\log\ell\le u ,
\end{equation}
so \eqref{eq:p8c3-prime-sum} together with \(\#\mathcal B_{\mathscr T}\le N\) and \(N\ge128\) yields \(|h_{N,Z,\tau_0}|\le4n\lambda_N\)). Substituting \(p_{\tau_0,1}Q_1\) into \eqref{eq:p8c3-cutoff} gives \(rA+nb=-V^{-1}\sum_{Z\in\operatorname{Comp}}h_{N,Z,\tau_0}V_Z\), whence \(|rA+nb|\le4n\lambda_N\) by the triangle inequality since \(V_Z>0\) and \(\sum_Z V_Z=V\).
\leanref{GlobalTheory.\allowbreak{}CompanionTheory.\allowbreak{}Ar.\allowbreak{}Top.\allowbreak{}flat\_\allowbreak{}85\_\allowbreak{}bound}{Development.\allowbreak{}AutomorphicSpectrum.\allowbreak{}Proofs.\allowbreak{}DivisibilityProofs.\allowbreak{}CutoffCoefficientsBound}{37--81}

Now evaluate \eqref{eq:p8c3-cutoff} on the Hodge test polynomial \(F_*=\mu^{e+1}J\) of Lemma~\ref{B8.26a}. Since \(3s+1\le n\), we have \(s<r\) and \(n\le2r\). Uniform shifts within the first \(r\) variables or across all \(n\) variables leave every within-block difference in \(J\) invariant, while \(\sum_{i<r}\partial_{\tau,i}\mu=0\) and \(\sum_{i=0}^{n-1}\partial_{\tau,i}\mu=s\); thus \(S(F_*)=0\) and \(T(F_*)=fs(e+1)D(\mu^eJ)=fsH_*\). Under the projection \(\pi:\prod_{\tau\in\Psi}\operatorname{Fl}(\Z^n)\to\prod_{\tau\in\Psi}\operatorname{Gr}(s,n)_{\Z}\), the class \(\mu=\sum_{\tau,i\ge r}X_{\tau,i}\) in Definition~\ref{B8.26r} pulls back the external tensor product \(\overline M\) of the Pl\"ucker Fubini--Study lines \(\overline{\mathcal O_{\operatorname{Gr}(s,n)}(1)}\). For any horizontal integral subvariety \(Y\subseteq\prod_{\tau\in\Psi}\operatorname{Gr}(s,n)_{\Z}\) of generic dimension \(m\in\{0,\dots,e\}\), write \(II(\pi^{-1}(Y);(\pi^*\overline M)^{m+1};J)\) for the multilinear evaluation of \(II\) on \(\pi^{-1}(Y)\) (of arithmetic dimension \(m+1+d-e\)) obtained by substituting the Hermitian root lines \((\overline Q_{\tau,i})\) into the homogeneous degree-\((d-e)\) polynomial \(J\) in the last \(d-e\) slots. By induction on \(m\in\{0,\dots,e\}\), choosing a Pl\"ucker coordinate section \(s_Y\) of \(\overline M\) not vanishing identically on \(Y\) and applying the rational-section divisor recursion simultaneously to \(\pi^*s_Y\) on \(\pi^{-1}(Y)\) and to \(s_Y\) on \(Y\) shows that \(II(\pi^{-1}(Y);(\pi^*\overline M)^{m+1};J)=c_{\rm fib}\,II(Y;\overline M^{m+1})\), where \(c_{\rm fib}=D(\mu^eJ)/\deg_{\rm gen}(\overline M)\) is the constant fibrewise degree of \(J\) on the smooth \((d-e)\)-dimensional fibres of \(\pi\): on every vertical component of \(\operatorname{div}_Y(s_Y)\) the factor \(c_{\rm fib}\) separates by the geometric projection formula, on every horizontal component (absent when \(m=0\)) it separates by induction, and in the archimedean integral on a characteristic-zero resolution \(\widetilde Y\to Y_{\C}\) the factors \(\log\|\pi^*s_Y\|\) and \(\pi^*c_1(\overline M)^m\) are pulled back from \(\widetilde Y(\C)\) while Fubini's theorem integrates the \((d-e,d-e)\)-form \(J(c_1(\overline Q_{\tau,i}))\) fibrewise to \(c_{\rm fib}\). Taking \(Y=\prod_{\tau\in\Psi}\operatorname{Gr}(s,n)_{\Z}\) (\(m=e\)) identifies \(h_{\rm fl}(F_*)\) with \(c_{\rm fib}\,II\bigl(\prod_{\tau\in\Psi}\operatorname{Gr}(s,n)_{\Z};\overline M^{e+1}\bigr)\). On \(\operatorname{Gr}(s,n)_{\C}\), affine chart matrices \(Z\in\C^{s\times r}\) whose entries are roots of unity define a Zariski-dense collection of horizontal orbits whose Pl\"ucker coordinates are algebraic integers including \(1\) and satisfy \(\det(I_s+ZZ^*)\le(1+r)^s\le(n+1)^s\) by Hadamard's inequality, giving normalised height at most \(\tfrac12 fs\log(n+1)\) on \(\prod_{\tau\in\Psi}\operatorname{Gr}(s,n)_{\Z}\). Lemma~\ref{B8.26p}(2) therefore yields
\begin{equation}\label{eq:p8c3-flag-test}
 h_{\rm fl}(F_*)\le \tfrac12 H_*fs\log(n+1).
\end{equation}
\leanref{GlobalTheory.\allowbreak{}CompanionTheory.\allowbreak{}Ar.\allowbreak{}Top.\allowbreak{}testPoly\_\allowbreak{}facts}{Development.\allowbreak{}AutomorphicSpectrum.\allowbreak{}Proofs.\allowbreak{}DivisibilityProofs.\allowbreak{}TestPolynomialShiftTotal}{13--15}

To bound \(R_N(F_*)\), put \(a=rs\) and \(X=2(a+n)\le\frac32 n^2\) (using \(a\le n^2/4\) and \(n\ge4\)). In each summand of \eqref{eq:p8c3-cutoff-tail} for a decomposition \(F_*=\sum_{\tau,k}p_{\tau,k}Q_{\tau,k}\) and \(2\le k\le n\), the full power sum \(p_{\sigma,k-1}\) integrates to zero against \(Q_{\tau,k}\) under \(D\), so \(P_{\ell,\tau,h,k-1}\) may be replaced up to sign by the degree-\((k-1)\) power sum on the smaller block of size \(s\). Lemma~\ref{B8.26d} therefore gives
\[
 \frac{k|D(P_{\ell,\tau,h,k-1}Q_{\tau,k})|}{H_*}
 \le
 \begin{cases}
 \displaystyle \frac{2^{k-1}}{a+1}\binom{a+k}{k},&\sigma=\tau,\\[4pt]
 \displaystyle \frac{2^{k-1}a(a-1)\cdots(a-k+2)}{k!},&\sigma\ne\tau ,
 \end{cases}
\]
where \(\sigma\in\Psi\) is the underlying embedding of \(\operatorname{Fr}_{\ell,h-1}(\tau)\). Both right-hand expressions are at most \(X^{k-1}/k!\), since \(k!\binom{a+k}{k}/(a+1)=\prod_{j=2}^k(a+j)\le(a+n)^{k-1}\) and \(a(a-1)\cdots(a-k+2)\le a^{k-1}\). For any prime \(\ell>N\ge4fn^2\), \(\sum_{h\ge1}\ell^{-hk}\le2/\ell^k\) and \(X/\ell\le\frac32 n^2/(4n^2)=3/8\), so
\[
 \sum_{k=2}^n\frac{2X^{k-1}}{\ell^k k!}
 \le\frac{2X}{\ell^2}\sum_{j\ge0}\frac{(X/\ell)^j}{(j+2)!}
 \le\frac{X}{\ell^2}e^{3/8}
 \le\frac{1.5X}{\ell^2}.
\]
Summing over the \(f\) embeddings \(\tau\in\Psi\), the contribution of \(\ell\) has modulus at most \(1.5fXH_*(\log\ell)/\ell^2\). Partial summation with \(\vartheta(t)=\sum_{\ell\le t}\log\ell\le t\log4\) gives
\begin{equation}\label{eq:p8c3-prime-tail}
 \sum_{\ell>N}\frac{\log\ell}{\ell^2}
 \le2\int_N^\infty\frac{\vartheta(t)}{t^3}\,dt
 \le\frac{2\log4}{N}.
\end{equation}
Since \(1.5fX\cdot(2\log4)/N\le\frac94 fn^2(2\log4)/(4fn^2)=\frac98\log4<1.6\), every finite partial sum over good primes \(\ell>N\), and hence the convergent series \(R_N(F_*)\), satisfies \(|R_N(F_*)|\le1.6H_*\).
\leanref{GlobalTheory.\allowbreak{}CompanionTheory.\allowbreak{}Ar.\allowbreak{}Top.\allowbreak{}testRest\_\allowbreak{}bound}{Development.\allowbreak{}AutomorphicSpectrum.\allowbreak{}Proofs.\allowbreak{}DivisibilityProofs.\allowbreak{}TestPolynomialRemainderBound}{417--443}

Combining Lemma~\ref{B8.26a} (\(P_N(F_*)/V\ge-C_0fs\lambda_NH_*\) with the absolute constant \(C_0=C_h+2\)) with \eqref{eq:p8c3-cutoff}, \(S(F_*)=0\), \(T(F_*)=fsH_*\), \eqref{eq:p8c3-flag-test} and \(|R_N(F_*)|\le1.6H_*\) gives
\[
 -C_0fs\lambda_NH_*\le\tfrac12 H_*fs\log(n+1)-b\,fsH_*+1.6H_* .
\]
Dividing by \(fsH_*>0\) and using \(1/(fs)\le1\) yields \(b\le C_0\lambda_N+\tfrac12\log(n+1)+1.6\). From \(rA+nb\ge-4n\lambda_N\) and the upper bound on \(b\), we get \(rA\ge-n\bigl((C_0+4)\lambda_N+\tfrac12\log(n+1)+1.6\bigr)\); since \(n\le2r\) and the bracket is nonnegative, dividing by \(r>0\) gives \(A\ge-2\bigl((C_0+4)\lambda_N+\tfrac12\log(n+1)+1.6\bigr)\). Finally \(n+1\le4fn^2\le N\), \(\operatorname{rd}(E)\ge1\) and \(N\ge128=2^7\) give \(\log(n+1)\le\log N\le\lambda_N\) and \(1.6<7\log2=\log128\le\lambda_N\), whence \(A\ge-\bigl(2(C_0+4)+3\bigr)\lambda_N=-K\lambda_N\).
\leanref{GlobalTheory.\allowbreak{}CompanionTheory.\allowbreak{}Ar.\allowbreak{}Top.\allowbreak{}lower\_\allowbreak{}bound\_\allowbreak{}key\_\allowbreak{}inequality}{Development.\allowbreak{}AutomorphicSpectrum.\allowbreak{}Proofs.\allowbreak{}DivisibilityProofs.\allowbreak{}TestHeightBound}{17--27}
\leanref{GlobalTheory.\allowbreak{}CompanionTheory.\allowbreak{}Ar.\allowbreak{}Top.\allowbreak{}lower\_\allowbreak{}bound\_\allowbreak{}finite\_\allowbreak{}arith}{Development.\allowbreak{}AutomorphicSpectrum.\allowbreak{}Proofs.\allowbreak{}DivisibilityProofs.\allowbreak{}TestHeightBound}{28--31}
\end{proof}

\subsection*{Proof of \ref{B8.26c}}\label{prf:B8.26c}\phantomsection
\begin{proof}
Since \(v\) has gap \(g>0\), each row \(v_\tau\) is strictly decreasing, so \(v\) is regular dominant and in particular dominant. Lemma~\ref{B8.26s}(1) therefore gives \(\delta_{\rm fl}(v)>0\), so \(H_v=(d+1)\delta_{\rm fl}(v)>0\), and Lemma~\ref{B8.26s}(4) gives \(h_{\rm fl}(L_v^{d+1})\ge0\). Because \(B(v,g)\ge4fn^2\) and the integer \(N\) satisfies \(N+1>B(v,g)\), we also have \(N\ge4fn^2\), so Lemma~\ref{B8.26b} applies at cutoff \(N\) to \(F=L_v^{d+1}\).

Write \(s_\tau=\sum_{i<r}v_{\tau,i}\), \(t_\tau=\sum_{i=0}^{n-1}v_{\tau,i}\) and \(a_\tau=\sum_{i=0}^{n-1}|v_{\tau,i}|\). Differentiating \(L_v^{d+1}\) gives \(\operatorname{Sh}_r(\tau,F)=-s_\tau H_v\) and \(\operatorname{Tot}(\tau,F)=-t_\tau H_v\), so \((-AS(F)-bT(F))/H_v=\sum_{\tau\in\Psi}(As_\tau+bt_\tau)\). Rewrite each summand as
\[
 As_\tau+bt_\tau=A\bigl(s_\tau-rt_\tau/n\bigr)+(rA+nb)t_\tau/n .
\]
Since \(v_\tau\) is nonincreasing, the mean of its first \(r\) entries is at least the mean of all \(n\) entries, while \(ns_\tau-rt_\tau=s\sum_{i<r}v_{\tau,i}-r\sum_{i\ge r}v_{\tau,i}\le\max(r,s)a_\tau\le na_\tau\); hence \(0\le s_\tau-rt_\tau/n\le a_\tau\), and \(|t_\tau|\le a_\tau\). Using \(A\ge-K\lambda_N\) and \(|rA+nb|\le4n\lambda_N\) from Lemma~\ref{B8.26b} gives \((-AS(F)-bT(F))/H_v\ge-(K+4)\lambda_N\|v\|_1\).

To bound \(R_N(F)\) in a decomposition \(F=\sum_{\tau,k}p_{\tau,k}Q_{\tau,k}\), recall that for \(k\ge2\) the complete power sum \(p_{\sigma,k-1}\) vanishes against \(Q_{\tau,k}\) under \(D\), so each \(P_{\ell,\tau,h,k-1}\) in \eqref{eq:p8c3-cutoff-tail} may be replaced up to sign by the degree-\((k-1)\) power sum on the smaller block of size \(s\). Fix \(\tau\in\Psi\) and translate row \(v_\tau\) into \([0,U_\tau]\), where \(U_\tau=v_{\tau,0}-v_{\tau,n-1}\) is the row spread; this changes \(L_v\) by a multiple of \(p_{\tau,1}\) and leaves \(Q_{\tau,k}\) (\(k\ge2\)) and \(H_v\) unchanged. For an index \((\sigma,i)\) in the smaller block (\(i\ge r\), where adjacent gaps are at least \(g\)), perturb the \((\sigma,i)\)-entry by \(+t\), using the translated row when \(\sigma=\tau\). Weyl alternation expresses \(kD(X_{\sigma,i}^{k-1}Q_{\tau,k})\), up to sign, as \((d+1)!(k-1)!\) times the \(t^{k-1}\)-coefficient of the corresponding sum of Vandermonde-type determinants in which one column exponent at \(\tau\) is raised by \(k\) with column-factorial denominators.

Reducing the raised monomial modulo \(\prod_{j=0}^{n-1}(X-v_{\tau,j}(t))\) factors out the Vandermonde. In column \(c\), starting at exponent \(n-1-c\), elementary and complete symmetric expansions majorise the coefficient of \(X^{n-1-c}\) coefficientwise in \(t\ge0\) by \(U_\tau^k\binom{k-1}{c}\binom{n-1-c+k}{k}e^{kt/U_\tau}\). Multiplying by \((n-1-c)!/(n-1-c+k)!\) replaces the second binomial by \(1/k!\), and summing over \(c\) gives \(2^{k-1}U_\tau^ke^{kt/U_\tau}/k!\). Dividing the perturbed Vandermonde at \(\sigma\) by its \(t=0\) value leaves factors \(1\pm t/(v_{\sigma,i}-v_{\sigma,j})\) (\(j\ne i\)), whose product is majorised coefficientwise by \(\prod_{j\ne i}(1+t/g)\le e^{nt/g}\) because \(i\ge r\) gives \(|v_{\sigma,i}-v_{\sigma,j}|\ge g\). Extracting \([t^{k-1}]\) yields
\begin{equation}\label{eq:p8c3-insertion}
 k|D(X_{\sigma,i}^{k-1}Q_{\tau,k})|
 \le H_v\frac{2^{k-1}}{k!}U_\tau(k+nU_\tau/g)^{k-1}.
\end{equation}
\leanref{RootFieldNumerics.\allowbreak{}M5.\allowbreak{}lam\_\allowbreak{}bound}{Development.\allowbreak{}WeilBound.\allowbreak{}CoreStatements.\allowbreak{}FlagImmersions.\allowbreak{}RestDegreePowerBound}{135--184}
\leanref{RootFieldNumerics.\allowbreak{}M5.\allowbreak{}lam\_\allowbreak{}shift}{Development.\allowbreak{}WeilBound.\allowbreak{}CoreStatements.\allowbreak{}FlagImmersions.\allowbreak{}RestDegreePowerBound}{199--210}
\leanref{RootFieldNumerics.\allowbreak{}M5.\allowbreak{}one\_\allowbreak{}bound}{Development.\allowbreak{}WeilBound.\allowbreak{}CoreStatements.\allowbreak{}FlagImmersions.\allowbreak{}RestDegreePowerBound}{212--248}
\leanref{RootFieldNumerics.\allowbreak{}M5.\allowbreak{}small\_\allowbreak{}block}{Development.\allowbreak{}WeilBound.\allowbreak{}CoreStatements.\allowbreak{}FlagImmersions.\allowbreak{}RestDegreePowerBound}{250--281}

Put \(D_*=n(1+2\|v\|_\infty/g)\) and \(B_*=B(v,g)=4fn^2(1+\|v\|_\infty/g)\). Since \(k\le n\) and \(U_\tau\le2\|v\|_\infty\), summing \eqref{eq:p8c3-insertion} over the \(s\) indices of the smaller block bounds \(k|D(P_{\ell,\tau,h,k-1}Q_{\tau,k})|\) by \(H_vsU_\tau(2D_*)^{k-1}/k!\). Every prime \(\ell>N\) satisfies \(\ell\ge N+1>B_*\), so \(\sum_{h\ge1}\ell^{-hk}\le2\ell^{-k}\le2B_*^{-(k-2)}\ell^{-2}\), and applying partial summation as in \eqref{eq:p8c3-prime-tail} from lower limit \(B_*\) bounds \(\sum_{\ell>N}2(\log\ell)/\ell^2\le\sum_{\ell>B_*}2(\log\ell)/\ell^2\) by \(4\log4/B_*\le5.56/B_*\). Since \(2D_*/B_*\le2n(1+2\|v\|_\infty/g)/(4fn^2(1+\|v\|_\infty/g))\le1/2\) and \(5.56\sum_{k\ge2}2^{-(k-1)}/k!<2\), summing over \(k\ge2\) and \(\tau\in\Psi\) gives
\begin{equation}\label{eq:p8c3-rest-bound}
 |R_N(F)|\le2H_vs\operatorname{spr}(v).
\end{equation}

By clause~1 of Definition~\ref{B8.26r}, there is a common volume \(V_0>0\) with \(V_Z=V_0\) for every \(Z\in\operatorname{Comp}\) and \(V=(\#\operatorname{Comp})V_0>0\). Hence
\[
 \frac1{\#\operatorname{Comp}}\sum_{Z\in\operatorname{Comp}}a_Z(v;N)
 =\frac1{\#\operatorname{Comp}}\sum_{Z\in\operatorname{Comp}}\frac{P_{N,Z}(F)}{V_0H_v}
 =\frac{P_N(F)}{VH_v}.
\]
Applying \eqref{eq:p8c3-cutoff} with \(h_{\rm fl}(F)\ge0\), \((-AS(F)-bT(F))/H_v\ge-(K+4)\lambda_N\|v\|_1\), \eqref{eq:p8c3-rest-bound} and the absolute constant \(C_T=K+4+2=2C_0+17=2C_h+21\ge2\) yields \eqref{eq:p8c3-model-mean}.
\end{proof}

\subsection*{Proof of \ref{B8.27}}\label{prf:B8.27}\phantomsection
\begin{proof}
Let \(f:Y\to X\) be the modification of Lemma~\ref{B8.25h}, including the normalisation. Every geometric fibre curve contracted by \(f\) has pullback degree zero. On a fibre curve mapping onto a curve the degree is the degree on its image multiplied by the degree of the map. Thus \(f^*A\) is nef on every geometric fibre. Its metric is smooth and semipositive by the Levi-form calculation in the lemma. Moriwaki's lower volume inequality applies on the projective flat integral model \(Y\) and gives \(\widehat{\deg}(\widehat c_1(f^*A)^{d+1})\le\widehat{\operatorname{vol}}(f^*A)\) \cite[Cor.~6.4]{moriwaki-2009}.

The arithmetic intersection on the left is \(II(P;A^{d+1})\). Indeed the divisor recursion for the intersection system counts each vertical component with its multiplicity and its degree times \(\log p\). On the complex fibre the recursion subtracts the integral of the log norm against the remaining Chern forms. This is the arithmetic divisor recursion in the unsquared norm convention. It determines the number by induction on dimension: for a possibly negative multidegree, tensor with a sufficiently positive coordinate line, use multilinearity, and use a coordinate monomial nonzero at the generic point in each resulting nonnegative degree. The arithmetic projection formula then gives the same number on \(Y\), since \(f_*[Y]=[X]\) \cite[\S3]{bost-gillet-soule-1994}. By Lemma~\ref{B8.25h} the volume on the right is the limit of the ordinary chart count in \eqref{eq:p8c3-moriwaki}. The normalisation uses the absolute arithmetic degree over \(\Z\) and the arithmetic dimension \(d+1\).
\end{proof}

\subsection*{Proof of \ref{B8.25p}}\label{prf:B8.25p}\phantomsection
\begin{proof}
Write \(H_a=\sum_{i=1}^a1/i\), with \(H_0=0\). For a complex homogeneous polynomial \(g\) of degree \(k\) in \(N\) variables, let \(m(g)\) be its spherical logarithmic mean. If \(u\) is a unit vector, then
\[
                 \log|g(u)|\le m(g)+kH_{N-1}/2 .
\]
Indeed, choose \(u\) as the first coordinate vector. Conditional on the first coordinate \(t\) of a uniform sphere point, the remaining coordinates are uniform on the sphere of radius \(\sqrt{1-|t|^2}\). Submean for the logarithm of the absolute value of a holomorphic function gives conditional mean at least \(\log|g(tu)|\). For \(N>1\), the density of \(|t|^2\) is \((N-1)(1-s)^{N-2}\,ds\) on \([0,1]\). Integration gives
\[
 \mathbb E[-\log |t|]
   =\frac12\int_0^1\frac{1-(1-s)^{N-1}}s\,ds
   =\frac{H_{N-1}}2 .
\]
For \(N=1\), \(|t|=1\). Submean may first be applied to regularised logarithms and then by decreasing convergence. Choosing a point of nonzero value gives finiteness of the mean. Iterating over the blocks gives the same inequality with \(k\) replaced by \(e(G)\).

The substituted polynomial is nonzero because the \(w_i\) span. Choose a nonzero coefficient \(a\). Its absolute value at a finite place is at most the local coefficient maximum. At an embedding, Cauchy's coefficient formula bounds \(|\sigma(a)|\) by the supremum on the unit polydisc of the substitution variables. Each argument of \(G\) there has norm at most \(n\Lambda\). Homogeneity and the sphere inequality bound the logarithm of that supremum by
\[
              m_\sigma(G)+e(G)\bigl(\log(n\Lambda)+H_{N-1}/2\bigr).
\]
Sum over the places and use the product formula for \(a\). Since \(N\le n\), \(H_{N-1}\le1+\log n\). This proves the lower bound. At a finite place the coefficient maximum is multiplicative. To see this, scale each polynomial to have integral coefficients and at least one unit coefficient, then multiply the nonzero reductions. At an embedding logarithms add almost everywhere and have finite integrals. This proves the product formula for \(h_0\).

Let \(e\) be the exterior product of the chosen basis and write \(\bigwedge^vM=\mathfrak a e\) for a fractional ideal \(\mathfrak a\). Every coefficient of the substituted \(De\) is a sum of exterior products of vectors of \(M\), hence belongs to \(\mathfrak a e\). At every archimedean argument on the product of spheres, Hadamard's inequality gives
\[
                |\sigma(D)|\,\|e\|_\sigma
                \le\prod_{j=1}^v\|P_j\|_\sigma
                \le\prod_{j=1}^vL_j .
\]
Integrate and sum. The finite contribution from \(\mathfrak a\) and the archimedean contribution from \(-\log\|e\|_\sigma\) sum to \(\widehat{\deg}M\). A change of basis multiplies the scalar determinant by an element of \(K^\times\), whose total logarithm is zero.
\end{proof}

\subsection*{Proof of \ref{B8.25e}}\label{prf:B8.25e}\phantomsection
\begin{proof}
Top Chern proportionality gives \(\deg_Z(v)=v_{0,Z}V(v)\), with \(v_{0,Z}>0\) proportional to the component volume. Volume equality gives the common \(\Delta'_Z\). The weight-gap criterion and \eqref{eq:p8c3-separation} give a birational map for the image. Comparison of moving degrees by generic multiplicities gives \(\deg_Z(u')\le\Delta_Z\le\deg_Z(u_1)\), with \(\deg_Z(u')/\deg_Z(u_1)=1-\rho\).

Put \(T=T_*\) and \(d=d_{\rm fl}\). The flat twist preserves rank. The nef section count for the generated Hodge model gives \(\binom{T-1}{d}\Delta_Z\le r_Z\). For the other bound set \(n_0=T+d+3\) and \(J_0=\lfloor n_0/(1-\rho)\rfloor\). Then \(n_0\deg_Z(u_1)<(J_0+1)\Delta_Z\). The adjacent gaps give the dominance comparisons of the weight with \(n_0u_1\) for jet interpolation. Evaluation on \(J_0\)-jets at the generic fibre of degree \(\Delta_Z\) is injective and bounds \(r_Z\) by \(\binom{J_0+d}{d}\Delta_Z\). The smallest factor of \(\binom{T-1}{d}/\binom{J_0+d}{d}\) is \((T-d)/(J_0+1)\). Since \(J_0(1-\rho)\le T+d+3\), \(d\ge2\) and \(J_0+1\ge T\),
\[
 J_0+d-T+1\le J_0\rho+2d+4\le(J_0+1)(\rho+4d/T).
\]
Each factor thus has relative loss at most \(\rho+4d/T\). The product is at least \(1-d(4d/T+\rho)\ge1-z\).
\leanref{GlobalTheory.\allowbreak{}CovTop.\allowbreak{}ShareRowsNef.\allowbreak{}nef\_\allowbreak{}share\_\allowbreak{}lower\_\allowbreak{}bound}{Development.\allowbreak{}AutomorphicArithmetic.\allowbreak{}ThetaCorrespondence.\allowbreak{}ThetaPairLifts.\allowbreak{}HodgeBundleWeights.\allowbreak{}ShareStatementOfKnownTheorems}{32--56}
\leanref{GlobalTheory.\allowbreak{}CovTop.\allowbreak{}ShareRowsNef.\allowbreak{}nef\_\allowbreak{}share\_\allowbreak{}upper\_\allowbreak{}bound}{Development.\allowbreak{}AutomorphicArithmetic.\allowbreak{}ThetaCorrespondence.\allowbreak{}ThetaPairLifts.\allowbreak{}HodgeBundleWeights.\allowbreak{}ShareStatementOfKnownTheorems}{55--91}

For the section degree put \(h=h_Z\), \(\delta=\Delta_Z\), \(v=r_Z\), \(B=\binom{T-1}{d}\delta\) and \(\zeta=1-B/v\). By \eqref{eq:p8c3-shares}, \(0\le\zeta\le z\). Use the component cusp lattices with their Petersson metrics. Let \(n_1\) be the size of the integral generating frame for \(u_1\), \(\Lambda\) an upper bound for its column norms, and \(\beta\) the Bergman upper comparison for multiplication by those sections. Thus multiplication by a linear combination of the frame columns has norm at most \(\sqrt\beta\) times the Petersson norm of the combination. Take \(n_1,\Lambda,\beta\ge1\). Put \(P=P_*\), \(q=\kappa/T\) and
\begin{equation}\label{eq:p8c3-hs-parameters}
 \begin{gathered}
 \Lambda_d=\Lambda_T=e^{P\kappa},\quad J=dP(q+1),\quad
 g=(d+1)(\log n_1+1)/2+1,\\
 H=\log\Lambda+26\log\{n_1d\delta(2+|h|+\log\Lambda)\},\\
 \log n_1,\ \log\beta,\ \log\delta,\ \log b\le P,\qquad
 \log\Lambda\le P(q+1),\qquad TP\le\sqrt b,\qquad zP\kappa\le b .
 \end{gathered}
\end{equation}
These frame and parameter comparisons are for the generated Hodge model, with a second component frame in weight \(u_2\). Denote the canonical flag weight by \(w_{\rm can}\). Use the weights \(w_j=w_0+ju_1\) for \(0\le j\le T\), with \(w_0=w_{\rm can}+u_2+du_1\). The terminal weight gives the untwisted lattice of degree \(D_0\) and rank \(v\). At degrees \(0,d,T\) integral spanning families have radii \(e^J,\Lambda_d,\Lambda_T\), respectively.
\leanref{GlobalTheory.\allowbreak{}HSShortRun.\allowbreak{}hilbertSamuel\_\allowbreak{}estimate\_\allowbreak{}at\_\allowbreak{}level}{Development.\allowbreak{}AutomorphicArithmetic.\allowbreak{}ThetaCorrespondence.\allowbreak{}ThetaPairLifts.\allowbreak{}HodgeBundleWeights.\allowbreak{}ExistsDegreeParametersBound}{33--76}
\leanref{GlobalTheory.\allowbreak{}HSShortRun.\allowbreak{}hsClaim\_\allowbreak{}short\_\allowbreak{}alt\_\allowbreak{}alt}{Development.\allowbreak{}AutomorphicArithmetic.\allowbreak{}ThetaCorrespondence.\allowbreak{}ThetaPairLifts.\allowbreak{}ThirdPlaceSeries.\allowbreak{}ExistsDegreeParametersBound}{34--70}

Take a transverse cut of the birational image by \(d\) linear sections \(\ell_1,\ldots,\ell_d\) in the span of the frame for \(u_1\). Work over a splitting field of the cut points \(y_1,\ldots,y_\delta\). The degree-zero restriction is onto their fibres by the nef and big adjoint restriction comparison. Choose \(\delta\) degree-zero sections \(s_1,\ldots,s_\delta\) from the short spanning family whose restrictions form a basis. If \(I\) is the cut ideal, the symbol at order \(j\) has kernel the sections of order at least \(j+1\). It sends a product with a cutting monomial of degree \(j\) to its normal monomial times the restriction of the coefficient to the cut.

For the deepest step use the kernel of restriction at each degree \(i>d\): a section \(s\) of degree \(i\) with zero restriction to the cut can be written \(s=\sum_{j=1}^d\ell_j s_j'\) with \(s_j'\) of degree \(i-1\). This follows from the Koszul resolution of the cut and adjoint vanishing for the lower weights in that resolution. Iterate this equality. If a section of degree \(T\) vanishes to order \(a+1\) and is written as a sum of cutting monomials of degree \(a\) with coefficients of degree \(T-a\), its order-\(a\) symbol is zero. The normal monomials are independent at each cut point, so all coefficients restrict to zero. The kernel equality divides each coefficient once more. Starting at \(a=0\) proves that order at least \(T-d\) is spanned by products of the degree-\(d\) sections with cutting monomials of degree \(T-d\). These restriction and symbol maps commute with base change to the splitting field.
\leanref{GlobalTheory.\allowbreak{}HSShort.\allowbreak{}flatAt\_\allowbreak{}of\_\allowbreak{}simpleAt}{Development.\allowbreak{}AbelianVarietyTheory.\allowbreak{}Proofs.\allowbreak{}DivisorProofs.\allowbreak{}HilbertSamuelFiltrationAtGoodSections}{27--75}
\leanref{GlobalTheory.\allowbreak{}HSCore.\allowbreak{}filMono\_\allowbreak{}of\_\allowbreak{}kerRes}{Development.\allowbreak{}AbelianVarietyTheory.\allowbreak{}Proofs.\allowbreak{}DivisorProofs.\allowbreak{}FilMonoOfKerRes}{24--72}

Choose a cut where the deepest step has maximum rank \(n_d\). For \(\alpha\in\N^d\) with \(|\alpha|<T-d\), use the columns
\[
 s_i\,\ell_1^{\alpha_1}\cdots\ell_d^{\alpha_d}\ell^{T-|\alpha|},
\]
with \(\ell\) a variable last linear section. Their symbols at degree \(j=|\alpha|\) give a block diagonal matrix with \(\binom{j+d-1}{d-1}\) copies of the restriction matrix multiplied on the left by \(\operatorname{diag}(\ell(y_a)^{T-j})\). For \(\ell\) nonzero at each cut point, these \(B\) columns are independent modulo the deepest step. Choose \(n_d\) independent products from the short degree-\(d\) family and the cutting monomials of degree \(T-d\). Append these columns and complete to a basis from the degree-\(T\) short family. The splitting-field comparison preserves independence of these choices over the coefficient field. Fix all the chosen sections and monomial types and allow the cuts to vary. The deepest rank and the relevant nonzero minors persist on an open set.
\leanref{GlobalTheory.\allowbreak{}HSShort.\allowbreak{}ChooseCols.\allowbreak{}exists\_\allowbreak{}cols}{Development.\allowbreak{}AbelianVarietyTheory.\allowbreak{}Proofs.\allowbreak{}DivisorProofs.\allowbreak{}ExistsShortColumnsDetNeZeroHilbertSamuel}{99--287}

Let \(D_{\rm col}\) be the determinant of these columns in a fixed terminal basis, as a polynomial in the \(d+1\) linear sections. Its specialisation is nonzero for generic \(\ell\), by the symbol matrices and the basis completion. Each symbol block supplies \(\prod_a\ell(y_a)^{T-j}\) to its determinant. The product over \(j\) is the \(u\)-th power of \(\prod_a\ell(y_a)\), where
\begin{equation}\label{eq:p8c3-chow-power}
 u=\sum_{j=0}^{T-d-1}\binom{j+d-1}{d-1}(T-j).
\end{equation}
For cuts in the indicated open set the Chow form \(F_{\rm Ch}\) specialises to a scalar times \(\prod_a\ell(y_a)\). Thus \(D_{\rm col}=F_{\rm Ch}^{u}G_{\rm col}\) for a nonzero multihomogeneous polynomial \(G_{\rm col}\). Divisibility holds over the function field of the first \(d\) linear sections, and the Chow form is primitive and irreducible in the full polynomial ring. The cut formula and the height normalisation are those of arithmetic Chow forms \cite{bost-gillet-soule-1994}.

The determinant height is \(u\,h(F_{\rm Ch})+h(G_{\rm col})\), by multiplicativity of the local coefficient norm at finite places and integration of the logarithm at infinite places. The binomial sum in \eqref{eq:p8c3-chow-power} gives \((d+1)u\delta=B(T+d(d+1))\). The first \(B\) columns have total degree \(T\) each, the next \(n_d\) have degree \(T-d\), and the terminal columns have degree zero. Hence the cofactor's total degree is
\[
 e_G=BT+n_d(T-d)-(d+1)u\delta
 =n_d(T-d)-Bd(d+1),\qquad 0\le e_G\le T(v-B).
\]
Nonnegativity follows from the multihomogeneous factorisation, and the upper bound uses \(n_d\le v-B\). Put
\[
 h'=\frac{h(F_{\rm Ch})}{(d+1)\delta},\qquad
 c=\log(n_1\Lambda)+\frac{\log n_1+1}{2}.
\]
Apply Lemma~\ref{B8.25p} to the nonzero multiforms \(F_{\rm Ch}\) and \(G_{\rm col}\) in the generating frame. Their total degrees are \((d+1)\delta\) and \(e_G\). Its lower bound gives \(h'\ge-c\) and \(h(G_{\rm col})\ge-ce_G\). The Chow height comparison gives \(h'\ge h-(\log n_1+1)/2\). The column coefficient norms in the frame metric are at most \(\beta^{T/2}e^J,\ \beta^{(T-d)/2}\Lambda_d,\ \Lambda_T\) in the three groups. After substituting the integral frame, each coefficient of a column is a sum of products of integral sections, hence belongs to the terminal lattice. On the product of unit spheres each multiplication by a cutting section costs at most \(\sqrt\beta\). The numbers of columns with the three displayed bounds are \(B,n_d,v-B-n_d\), respectively. Lemma~\ref{B8.25p} thus gives \(D_0\ge u h(F_{\rm Ch})+h(G_{\rm col})-B(J+T\log\beta/2)-n_d(\log\Lambda_d+(T-d)\log\beta/2)-(v-B-n_d)\log\Lambda_T\). Using the expression for \((d+1)u\delta\) and \(h'\ge-c\) gives
\[
 \begin{aligned}
 D_0\ge{}&B\{Th'-d(d+1)c-J-\tfrac T2\log\beta\}\\
 &-(v-B)\{Tc+\tfrac T2\log\beta+\log\Lambda_d+\log\Lambda_T\}.
 \end{aligned}
\]
In \eqref{eq:p8c3-hs-parameters}, \(H\ge c\) because the argument of the logarithm is at least \(2n_1\). Thus \(d(d+1)c\le26d^2(H+\log T)\), and the cost \(T(\log n_1+1)/2\) of replacing \(h'\) by \(h\) is at most \((T+d^2)(g+2)\). After division by \(v\) we may weaken the bound to
\begin{equation}\label{eq:p8c3-det-degree}
 \begin{aligned}
 D_0/v\ \ge{}&(1-\zeta)\{Th-3(d+1)|h|-J-26d^2(H+\log T)\\
 &\quad -(T+d+1)\log\beta-(T+d^2)(g+2)\}\\
 &-\zeta\{T(\log2+H+\tfrac12\log\beta)+\log\Lambda_d+\log\Lambda_T\}.
 \end{aligned}
\end{equation}
\leanref{GlobalTheory.\allowbreak{}HSShort.\allowbreak{}monomial\_\allowbreak{}counts}{Development.\allowbreak{}AbelianVarietyTheory.\allowbreak{}Proofs.\allowbreak{}DivisorProofs.\allowbreak{}MonomialCountsDegreeAndWeightedSum}{37--42}
\leanref{GlobalTheory.\allowbreak{}HSShortEnd.\allowbreak{}claimS\_\allowbreak{}of\_\allowbreak{}numbers}{Development.\allowbreak{}AbelianVarietyTheory.\allowbreak{}Proofs.\allowbreak{}DivisorProofs.\allowbreak{}ClaimSOfNumbers}{19--103}
\leanref{GlobalTheory.\allowbreak{}HSShort.\allowbreak{}SlopeTop.\allowbreak{}slopeBound}{Development.\allowbreak{}AbelianVarietyTheory.\allowbreak{}Proofs.\allowbreak{}DivisorProofs.\allowbreak{}SlopeLowerBoundTopLatticeOffHypersurface}{106--256}

Use \(H\le26\log(2+|h|)+H_0\), where \(H_0=27P(q+1)+65P\). This follows from \(2+|h|+\log\Lambda\le(2+|h|)(1+\log\Lambda)\), \(\log(1+\log\Lambda)\le\log\Lambda\), and \(\log d\le\log T\le P/2\). Multiplying \eqref{eq:p8c3-det-degree} by \(v/B\) puts the second brace at coefficient \(\zeta/(1-\zeta)\le2z\). By \eqref{eq:p8c3-z}, \eqref{eq:p8c3-quotient-size} and \eqref{eq:p8c3-hs-parameters},
\[
 d^2Pq\le b/7,\qquad d^2P\le b/14,\qquad TdP\le(TP)^2\le b .
\]
Substitution gives
\begin{equation}\label{eq:p8c3-det-costs}
 \begin{gathered}
 J\le9b,\qquad 26d^2(H_0+\log T)\le300b,\qquad
 (T+d+1)\log\beta\le2b,\\
 (T+d^2)(g+2)\le10b,\qquad
 2z\{T(\log2+H_0+\tfrac12\log\beta)+\log\Lambda_d+\log\Lambda_T\}\le200b .
 \end{gathered}
\end{equation}
For the last bound use \(zTPq=zP\kappa\le b\), \(2z\le1\) and \(TP\le\sqrt b\). Their sum is at most \(4000b\), and the coefficient of \(\log(2+|h|)\) is bounded by \(A=26(26d^2+2zT)\). By the flat degree comparison the twist lowers the degree by at most \(vC_p(frL_x)^2\) with \(C_p\ge0\) uniform. As \(v/B\le2\), take \(R=4000b+2C_p(frL_x)^2\). These choices are common to the components. Since \(1\le d\le T\) and \(2z\le1\), we have \(1\le A\le702T^2\). By \eqref{eq:p8c3-hs-parameters}, \(2\log T\le\log b\le P\), \(P\ge1\), and \(T^2P^2\le b\). Thus \(\log A\le\log702+2\log T\le8P\), and \(A\log A\le5616T^2P\le6000b\). Also \((frL_x)^2\le(frL_x)^{12}\le P\le TP\le\sqrt b\le b\), and the size bound of \ref{B8.25b} gives \(b\le E_p\). This proves \eqref{eq:p8c3-hs} with \(C_{\rm HS}=10000+2C_p\).

The component spaces sum to the cusp space and their ranks add. The integral direct sum is contained in the whole lattice, so its degree is at most the whole degree. Divide by rank and multiply by two in the squared metric to obtain \eqref{eq:p8c3-degree}.
\end{proof}

\subsection*{Proof of \ref{B8.25f}}\label{prf:B8.25f}\phantomsection
\begin{proof}
Equip \(K\otimes_\Q\R\) with the Euclidean norm squared \(\sum_\sigma|\sigma x|^2\), with conjugates counted separately. The Euclidean dual of \(\mathcal O_K\) in this norm is the image of the inverse different with conjugation in the complex factors, by the trace pairing. For a nonzero element \(y\) of the inverse different, the integral ideal \((y)\mathfrak D_K\) has norm at least one. Hence
\[
 \max_\sigma|\sigma y|\ge |N_{K/\Q}(y)|^{1/a}\ge\operatorname{rd}(K)^{-1}.
\]
The Euclidean transference bound \(\lambda_a(\Lambda)\lambda_1(\Lambda^\vee)\le a\) \cite{banaszczyk-1993} gives \(a\) rationally independent integers \(\omega_i\) of Euclidean length at most \(a\operatorname{rd}(K)\). In particular
\begin{equation}\label{eq:p8c3-short-integers}
 \max_\sigma|\sigma\omega_i|\le a^3\operatorname{rd}(K).
\end{equation}
For \(v=0\) the empty family works. Otherwise take a \(K\)-basis \(s_1,\ldots,s_v\) from \(S\). The \(av\) products \(e_{ij}=\omega_i s_j\) belong to \(L\) and form a rational basis. For each vector \(z=\sum c_{ij}e_{ij}\) of a \(\Z\)-basis of \(L\), put \(z'=z-\sum\lfloor c_{ij}\rfloor e_{ij}\). The remainders and the products generate every original basis vector. There are \(av\) remainders, since \(L\) has full rational span and is torsion-free. By \eqref{eq:p8c3-short-integers} each product has seminorm at most \(a^3\operatorname{rd}(K)H\). The triangle inequality with \(0\le c_{ij}-\lfloor c_{ij}\rfloor<1\) gives \(N_\sigma(z')\le a^4v\operatorname{rd}(K)H\).
\end{proof}

\subsection*{Proof of \ref{B8.25q}}\label{prf:B8.25q}\phantomsection
\begin{proof}
The transference calculation in Lemma~\ref{B8.25f}, applied to \(\mathcal O_E\), gives a rational spanning set \(\omega_i\in\mathcal O_E\) with \(\max_\sigma|\sigma(\omega_i)|\le2fR_E\le\Omega\). Thus each \(\omega_i\) acts on \(V\) with operator norm at most \(\Omega\).

Start with \(W=0\). At each stage take \(I=L\cap W\), maintain isotropy and \(\mathcal O_E\)-stability, and maintain the second displayed implication. Suppose \(x\in L\setminus I\) is at distance less than \(\theta\) from \(W\). Write \(x=v+h\), where \(v\in W\) and \(\|h\|<\theta\). Since \(B\) vanishes on \(W\), \(\|u(x)\|\le\|h\|<\theta\le(\Omega\theta)^{-1}\). Hence \(B(x,W)=0\). Adjointness gives \(B(ax,W)=0\) for every \(a\in\mathcal O_E\). In particular
\[
          |B(x,\omega_i x)|=|B(h,\omega_i h)|
                      \le\Omega\|h\|^2<1 .
\]
The left side is the absolute value of an integer. The rational spanning property and adjointness now give \(B(ax,bx)=0\) for all \(a,b\in\mathcal O_E\). Enlarge \(W\) to \(W'=W+\operatorname{span}_{\R}(\mathcal O_E x)\). This is an isotropic rational subspace of strictly larger dimension.

For \(y\in L\), restrict \(B(y,\cdot)\) to \(W'\). If its restriction to \(W\) is nonzero, its norm on \(W'\) is at least \((\Omega\theta)^{-1}\) by induction. If it vanishes on \(W\) and is nonzero on \(W'\), then \(B(y,\omega_i x)\) is a nonzero integer for some \(i\). The vector \(\omega_i h\) lies in \(W'\), and \(B(y,\omega_i h)=B(y,\omega_i x)\), since \(\omega_i v\in W\). The norm of the restriction is at least \(1/\|\omega_i h\|>(\Omega\theta)^{-1}\). This proves the induction step. The dimension bound on \(W\) terminates the construction. When \(\lambda\ge\theta\) the initial space already satisfies both conclusions. When \(\lambda<\theta\), the distance inequality puts every shortest vector in the final \(I\).

Let \(\rho=\dim_E(I\otimes\Q)\). Isotropy gives \(W\cap iW=0\), since \(B(w,iw)\) has absolute value \(\|w\|^2\). The complexification of \(W\) therefore embeds in \(V\). Each embedding of \(E\) occurs there with multiplicity \(\rho\), which proves both rank assertions.

Fix a shortest vector \(\xi\in I\). The integral functional \(a\mapsto B(x_j,a\xi)\) on \(\mathcal O_E\) has the form \(\operatorname{Tr}_{E/\Q}(a y_j)\) for a unique \(y_j\) in the inverse different. At least one \(y_j\) is nonzero. Otherwise \(B(ax_j,\xi)=B(x_j,\bar a\xi)=0\) for every \(a,j\), contrary to the spanning hypothesis. For each pair \(\{\sigma,\bar\sigma\}\) put \(s_\sigma=\|\xi_\sigma\|^2+\|\xi_{\bar\sigma}\|^2\). On \(E\otimes_\Q\R\), the map \(a\mapsto a\xi\) has squared norm \(\sum_{\{\sigma,\bar\sigma\}}s_\sigma|a_\sigma|^2\). The real linear functional \(a\mapsto\operatorname{Tr}(ay_j)\) has squared dual norm
\[
            4\sum_{\{\sigma,\bar\sigma\},\,s_\sigma>0}
                    |\sigma(y_j)|^2/s_\sigma
                 \le\|u(x_j)\|^2 .
\]
Indeed its value is \(\langle u(x_j),a\xi\rangle_\R\), and its coordinate at a pair is \(2\Re(a_\sigma\sigma(y_j))\). The functional vanishes on every coordinate with \(s_\sigma=0\). Since \(s_\sigma\le\lambda^2\), this gives \(2\sum_\sigma|\sigma(y_j)|^2\le\lambda^2\|u(x_j)\|^2\). For any nonzero \(y_j\), the integral ideal \(y_j\mathfrak D_E\) has norm at least one. The arithmetic-geometric mean inequality therefore gives
\[
       \sum_\sigma|\sigma(y_j)|^2
            \ge 2f\,|N_{E/\Q}y_j|^{1/f}
            \ge 2f R_E^{-2}.
\]
This proves the last assertion.
\end{proof}

\subsection*{Proof of \ref{B8.25r}}\label{prf:B8.25r}\phantomsection
\begin{proof}
At one embedding write \(J=\operatorname{diag}(\epsilon_i)\), with \(\epsilon_i=1\) for \(i<p\) and \(-1\) otherwise. Put \(L_j=\log|\det(n^*Jn)_{[0,j)}|\), including \(L_0=0\). A weight \(\nu\) has potential \(\sum_i\nu_i(L_{i+1}-L_i)\). The kernel potential \(R_x\) has weight \(\mathbf1_{i\ge p}\), since \(L_r=0\). The log density has weight \(2i-(r-1)\). For the latter formula, change from an orthogonal frame by an upper triangular matrix with diagonal entries \(b_i\). The tangent coordinates are the entries below the diagonal, on which the diagonal acts by \(b_a/b_c\), \(a>c\). The squared absolute Jacobian has exponent \(2i-(r-1)\) on \(|b_i|^2\), and \(\log |b_i|^2=L_{i+1}-L_i\).

For a strictly lower triangular tangent matrix \(H\), the Levi form of \(L_j\) at the origin is \(\sum_{c<j\le a}\epsilon_a\epsilon_c|H_{ac}|^2\). Indeed, in the first \(j\) columns of \(1+tH\) the upper \(j\)-rowed square block has determinant one. Factor it out of the Gram determinant. The lower rectangular block is \(t(H_{ac})_{a\ge j,c<j}+O(t^2)\), so the coefficient of \(|t|^2\) in the log determinant is the indicated signed sum. Telescoping in \(j\) gives the coefficient \(\epsilon_a\epsilon_c(\nu_c-\nu_a)\) on \(|H_{ac}|^2\). For \(\nu=\lambda\) this is at least \(T\) within a block by adding adjacent gaps. Across blocks it is \(\lambda_a-\lambda_c\ge\lambda_{r-1}-\lambda_0\ge T\). For the kernel weight the coefficient is zero or one. For the density weight it is at most \(2r\). Thus the lower bound at an orthogonal origin is \(T-b_*-2r\).

At a general point reorthogonalise the frame, writing \(xn=x'B\) with \(B\) upper triangular. In the overlap of the two holomorphic flag charts the metric potentials differ by logarithms of absolute squares of nonvanishing holomorphic diagonal factors. The log densities differ by the log squared absolute Jacobian. Their Levi forms are therefore transported by the chart differential. This proves nonnegativity throughout the chart.

Here is the norm of that differential on the specified ball. Set \(\delta=\|z\|\). Entrywise \(n^*Jn-J\) has absolute value at most \(2\delta+r\delta^2\le3\delta\). Since \(B^*JB=n^*Jn\), triangular elimination gives
\[
          \big||B_{aa}|^2-1\big|\le6\delta,\qquad
          |B_{ac}|\le6\delta\quad(a<c).
\]
Indeed, in row \(a\) the already eliminated terms contribute at most \(36r\delta^2\le0.36\delta\) to each entry. The diagonal equation first bounds the squared diagonal error by \(3.36\delta\). The off-diagonal equation then bounds the next entry by \(3.36\delta/\sqrt{1-6\delta}\le6\delta\). Write \(B=D_0+U_0\) with \(D_0\) diagonal. Its diagonal moduli lie in \([0.969,1.03]\), and \(\|U_0\|_{\rm F}\le6r\delta\le0.06\), where \(\|\cdot\|_{\rm F}\) is Frobenius norm.

Let \(H'\) be the tangent in the recentered chart and put \(Y=n^{-1}H\). Differentiating the change of frame gives \(BY=H'B+B_1\) with \(B_1\) upper triangular. Project to strictly lower entries. Since \(Y\) is strictly lower triangular, the norm inequality is
\[
        0.969\|Y\|_{\rm F}
          \le 1.09\|H'\|_{\rm F}+0.06\|Y\|_{\rm F}.
\]
Hence \(\|Y\|_{\rm F}\le1.2\|H'\|_{\rm F}\) and \(\|H\|_{\rm F}\le1.01\|Y\|_{\rm F}\). Square and sum over the embeddings. This gives \(\|H\|^2\le2\|H'\|^2\), as required.

To check the domain for the \(\bar\partial\) estimate, write the first \(p\) columns of \(n\) as \(\binom{n_1}{W_1}\), and put \(Z=W_1n_1^{-1}\). The matrix \(n_1\) is lower unitriangular. The condition on the first \(p\) columns is \(I-Z^*Z>0\). It implies all the other flag signature conditions, since the perpendicular complement of this positive \(p\)-plane is negative definite. In particular
\[
             R_x=-\sum_\tau\log\det(I-Z_\tau^*Z_\tau).
\]
Its Levi form is nonnegative by the kernel-weight computation at an orthogonal origin and change of chart. The function \(R_x+\log(1+\|z\|^2)\) is a plurisubharmonic exhaustion. A sublevel bounds \(\|z\|\), and it bounds the product of the determinants away from zero. Each factor is at most one, so a limit in this bounded set stays in the domain. This proves pseudoconvexity.

Let \(t=\|z\|^2/\rho^2\), \(F_0=\chi(t)\), and on the ball set
\[
                  \psi_j=D\chi(t)\log((t+1/(j+1))/2).
\]
Transport \(\psi_j\) to the translates of the ball and extend by zero. Injectivity makes the transported functions agree and the supports are locally finite by proper discontinuity. On \(t\le1/4\) the logarithm has nonnegative Levi form. On \(1/4\le t\le1\) its absolute value is at most \(\log8\), \(\partial\bar\partial t[v,\bar v]=\|v\|^2/\rho^2\), and \(|\partial t[v]|^2\le t\|v\|^2/\rho^2\). The two cross terms from differentiating the product cost at most \(2A\|v\|^2/\rho^2\), and the term with the differentiated cutoff costs at most \(2A\log8\,\|v\|^2/\rho^2\). Thus the negative part costs at most \(D(13A+20)\|v\|^2/\rho^2\). The gap hypothesis makes \(U+\psi_j\) plurisubharmonic everywhere. On the original cutoff annulus its Levi form dominates \(|\bar\partial F_0[v]|^2\). These comparisons persist under transport; outside the translates the added weight is zero.

Apply the weighted \(\bar\partial\) estimate of \cite[Th.~2.1]{blocki-2014} on this pseudoconvex chart to the smooth plurisubharmonic weight \(U+\psi_j\), the closed form \(\bar\partial F_0\), and the indicator of the original annulus as the bound for its squared dual norm. On that annulus \(-\psi_j\le3D\). The bounds on \(|B_{ii}|^2\) above give \(|w|\le\|\lambda\|_1\) and \(|R_x|\le f\) there, the latter by summing \(\log|B_{ii}|^2\) for \(i<p\). Consequently the solutions \(u_j\) satisfy
\[
       \int |u_j|^2 e^{-U-\psi_j}\,dz
       \le e^{3D+\|\lambda\|_1+b_*f}\mu(B(0,\rho)).
\]
For each fixed \(k\) the tail is bounded in \(L^2(e^{-U-\psi_k})\). Diagonal weak extraction gives a distributional limit \(u\) with \(\bar\partial u=\bar\partial F_0\). Weak lower semicontinuity followed by monotone convergence gives the same bound for the limiting weight. Near each orbit centre its singular factor is bounded below by a positive multiple of the inverse \(2D\)-th power of the distance. The solution is holomorphic there, and a nonzero constant value would give the divergent radial integral \(\int_0^\rho dt/t\). Hence it vanishes there. Now \(F=F_0-u\) has the asserted values. Use \(\psi_j\le0\) and \(|F|^2\le2|F_0|^2+2|u|^2\) for the integral bound. Since \(D\le fr^2\) and \(b_*f\le5(fr)^8\), its constant is uniform.
\end{proof}

\subsection*{Proof of \ref{B8.25s}}\label{prf:B8.25s}\phantomsection
\begin{proof}
The kernel depends on the positive planes. For graph matrices \(Z,Y\) at one embedding, the factor of \(K_y(z)^{1/b_*}\) is
\[
      \frac{\det(I-Z^*Z)\det(I-Y^*Y)}
           {|\det(I-Y^*Z)|^2}.
\]
Indeed the two Gram matrices of the graph bases are \(I-Z^*Z,I-Y^*Y\), and their mutual Gram matrix is \(I-Y^*Z\). This also gives symmetry and invariance. For \(z\in B_x\) in coordinates at \(x\), the formula \(Z=W_1n_1^{-1}\) in Lemma~\ref{B8.25r} gives \(\|Z\|_{\rm op}\le2/(100r^2)<1/2\). Thus the numerator factor \(\det(I-Z^*Z)\) is at least \((3/4)^p\), and \(|\det(I-Y^*Z)|\le(3/2)^p\). Multiplication over the embeddings gives \(K_y(z)\ge3^{-fpb_*}K_y(x)\).

For completeness, compute the mass against this kernel. Over the matrix ball at one embedding the invariant base density, equal to one at zero, is \(\det(I-Z^*Z)^{-r}\) times Lebesgue measure. The differential at zero of translation to \(Z\) acts on a rectangular tangent matrix \(H\), up to unitary factors, by
\((I-ZZ^*)^{1/2}H(I-Z^*Z)^{1/2}\). Its squared absolute Jacobian is \(\det(I-Z^*Z)^r\), giving the density. The fibre consists of a full flag in each of the two definite blocks. A compact flag of rank \(s\), with density one at the orthogonal origin in lower triangular coordinates, has mass \(\prod_{j=1}^{s-1}\pi^j/j!\). To see this, project to its first line. The invariant density on \(\mathbb P^j\) in the affine chart is \((1+\|z\|^2)^{-j-1}\), whose radial integral is \(\pi^j/j!\). The remaining flag has rank one less. Base and fibre measures multiply because their invariant densities agree at the origin, where the three blocks of lower triangular coordinates are disjoint. Since \(b_*\ge r\), the remaining base integral is bounded by the Euclidean volume of the product of \(p(r-p)\) unit discs. This gives \(\int K_y\,d\mu\le\pi^{fr^2}\).

On \(B_x\) the triangular bounds in Lemma~\ref{B8.25r} give \(|\log d|\le fr\), by the density-weight formula there. In dimension \(D=fr(r-1)/2\) it follows that
\[
           \mu(B_x)\ge e^{-fr}\frac{\pi^D}{D!(100r^2)^{2D}}
                     \ge e^{-C_1(fr)^8}
\]
for an absolute \(C_1\), using \(D!\le D^D\). Integrate the kernel comparison over the disjoint sets \(\gamma^{-1}B_x\). Symmetry and invariance give
\[
       3^{-fpb_*}\mu(B_x)\sum_\gamma K_x(\gamma y)
                            \le\int K_y\,d\mu .
\]
This proves the first bound.

For a finite sum of translates, weighted Cauchy--Schwarz bounds its squared norm at \(y\) by
\[
          \left(\sum_\gamma K_x(\gamma y)\right)
          \sum_\gamma\frac{\|F(\gamma y)\|^2}{K_x(\gamma y)}.
\]
Integration over a fundamental domain unfolds the second sum. The same inequality for tails gives convergence in \(L^2\) there. On a compact subset of the flag domain the translates have bounded overlap by proper discontinuity, so the tails converge in local \(L^2\) there as well. The holomorphic submean inequality in local frames gives locally uniform convergence. The estimate passes to the resulting section.
\end{proof}

\subsection*{Proof of \ref{B8.25t}}\label{prf:B8.25t}\phantomsection
\begin{proof}
Fix a rational isotropic line \(E\xi\), with \(\xi\ne0\) in the period lattice. Its intersection with that lattice is saturated and isotropic for the alternating form. Choose a symplectic basis \(e_j,e'_j\) in which this intersection lies in the first half. Such a basis is obtained by splitting off a primitive vector and a dual vector of pairing one, then repeating in their symplectic orthogonal complement. A change of symplectic basis changes the level theta function into a finite linear combination of theta functions with characteristics. This follows on translations from the theta translation formula and on interchange of the two halves from Poisson summation on the level cosets. Harmonicity gives the same harmonic differential after the change of period coordinates.

Write \(h\) for the fixed Hermitian form. In the expansion with one such characteristic let \(u_b\) be the integer summation vector plus the real characteristic in column \(b\), and put
\[
        L_{\tau}(u_b)=\sum_j u_{bj}\,\tau\bigl(h(e'_j,\xi)\bigr).
\]
Along the real centre of the cusp the frequency in coordinate \(\tau\) is a positive constant times \(\sum_{b=1}^N|L_\tau(u_b)|^2\). Indeed the change of the period matrix along that coordinate is a real multiple of the symmetric matrix
\(\bigl(\Re(\tau h(e'_j,\xi)\overline{\tau h(e'_k,\xi)})\bigr)_{jk}\).
Its quadratic form on a real column \(u_b\) is \(|L_\tau(u_b)|^2\). Thus zero frequency forces \(L_\tau(u_b)=0\) in every column.

We give the differential calculation at the embedding with \(a_{\tau,P}>0\). In frame coordinates let \(w=(w_1,w_2)\) be the coordinates of \(\xi\), and let \((\ell_1,\ell_2)\) be its Hermitian dual covector, with the blocks ordered so that the \(P\) differential directions come from the second block. Isotropy gives \(\ell_1w_1+\ell_2w_2=0\). Both \(w_2\) and \(\ell_1\) are nonzero, since a nonzero isotropic vector cannot lie in a definite summand. We may shear the frame by an upper block unipotent so that the complementary coordinates \(w_1,\ell_2\) vanish. Explicitly, take \(d_0w_2=1\), \(\ell_1 e_0=1\), and use the off-diagonal block
\[
 K=w_1d_0-e_0\ell_2-(\ell_1w_1)e_0d_0 .
\]
Then \(Kw_2=w_1\) and \(\ell_1K=-\ell_2\). This shear preserves the flag and its weight character.

In this frame let \(v\) be the period direction of \(\xi\) in the \(P\)-row block at \(\tau\). For clarity, write the variable part of a single theta summand as a product over its columns of
\[
       \exp\!\left(\tfrac{\pi}{2}z_b^t Bz_b
                         +2\pi i\,u_b^t H_0z_b\right).
\]
Here \(z_b\) is the period variable, \(H_0\) is the inverse first period matrix, and \(B\) is the symmetric matrix of the Gaussian prefactor in these coordinates. The period identities in the adapted basis give \(Bv=0\) and \(u_b^tH_0v\) proportional to \(\overline{L_\tau(u_b)}\). To see the two identities, the first period matrix sends the coordinates of each \(\alpha\xi\), \(\alpha\in E\), in the first half of the symplectic basis to its period vector. Applying the complementary period matrix instead gives zero after the shear. Separating the embeddings in these identities gives \(Bv=0\). Pairing \(\alpha\xi\) with the dual half \(e'_j\) reads its first-half coordinates by the alternating form. Separating the \(\tau\)-term gives the entries of \(H_0v\) as the same nonzero scalar times \(\overline{\tau h(e'_j,\xi)}\). The separation uses linear independence of the distinct embeddings of \(E\).
\leanref{GlobalTheory.\allowbreak{}ThFr.\allowbreak{}IsAdapted.\allowbreak{}hMat\_\allowbreak{}mulVec\_\allowbreak{}part}{Development.\allowbreak{}AbelianVarietyTheory.\allowbreak{}Proofs.\allowbreak{}CuspFormEvaluations.\allowbreak{}IsHarmonicAndMemMinorSpan}{300--326}
\leanref{GlobalTheory.\allowbreak{}ThFr.\allowbreak{}IsAdapted.\allowbreak{}splitMat\_\allowbreak{}mulVec\_\allowbreak{}part}{Development.\allowbreak{}AbelianVarietyTheory.\allowbreak{}Proofs.\allowbreak{}CuspFormEvaluations.\allowbreak{}IsHarmonicAndMemMinorSpan}{437--460}

At zero frequency each column factor is invariant under \(z_b\mapsto z_b+tv\). Consequently the same nonzero linear combination of the \(P\) row derivatives annihilates the summand in every column. Make this combination the last row by an invertible row operation. A full \(P\)-row determinant of the derivatives now annihilates the summand, by expansion along that row and commutation of constant-coefficient derivatives. The differential contains such a determinant to the positive power \(a_{\tau,P}\). Its zero-frequency coefficient therefore vanishes term by term.

At a fixed frame the differentiated terms are bounded by a polynomial times a positive-definite Gaussian in the summation vectors. Absolute convergence permits taking the mean along the real centre term by term, for instance by averaging over expanding boxes, where each nonzero exponential has mean tending to zero. Thus the zero Fourier--Jacobi coefficient vanishes for each rational isotropic line. At a higher-dimensional rational boundary component its constant term includes the average in any such line in that component. This proves cuspidality.
\end{proof}

\subsection*{Proof of \ref{B8.25gb}}\label{prf:B8.25gb}\phantomsection
\begin{proof}
For (1), put \(S_0=R_0^2=90(fr+D)\ge90(2+1)=270\), so \(R_0>0\). From \(1\le m\le m^4\le f\) and \(f^2\le R_E\) we obtain \(f\ge1\), \(f,fm,m^2,f^{3/2}\le R_E\) and \(d_*=2fm\le2R_E\). The hypothesis \(S_0R_E^{28}Dr\le q\) gives \(b_0(q,R_0)=q/(S_0R_E^{28})\ge Dr\ge1\), \(1\le D\le Dr\le q\) and \(t_0(q,R_0)=qb_0(q,R_0)\ge qDr\). Next, since \(R_E\ge16\),
\[
 d_*^2\Omega_*fS_0R_E^{21}\le(2R_E)^2(2R_E^2)(R_ES_0R_E^{21})=8S_0R_E^{26}<R_E^2S_0R_E^{26}=S_0R_E^{28}\le q,
\]
so \(d_*^2\Omega_*\theta(q,R_0)^2<1\). Also \(M_1=90m^2d_*^2R_E\le360R_E^4\le\frac{360}{16}R_E^5\le\frac12 S_0R_E^{28}\le q/2\). Using \(\log(q/16)\le q/16-1\), \(\log16\le15\) and \(\log D\le\log q\) gives \(\log(qD)\le q/8+28\), so \(4D(7+\log(qD))+d_*+1\le D(q/2+140+2R_E+1)\le qD\le t_0(q,R_0)\) because \(q\ge270R_E^{28}\ge4320R_E\). Finally \(\log2\le1\) and \(\log r\le r-1<r\le fr\) give \((D+1)(r\log r+3r+3)\le2Dr(fr+6)\), whence
\[
 fr\log2+2d_*+3+(D+1)(r\log r+3r+3)\le(3fr+4R_E+15)Dr\le qDr\le t_0(q,R_0).
\]
Thus \((E,r,D,q,R_0,m)\) is uniformly admissible.

For the bound on \(P_0=\mathcal P(E,r,q,m,5,1)\), every term is nonnegative since \(frq\ge1\), \(d_*\ge2\), \(A\ge625\) and \(K_*\ge16\). From \(e>2.71\) we have \(e^2>7.3\) and \(e^{10}>7.3^5>10000\), so \(\log2\le1\), \(\log5,\log6,\log7\le2\) and \(\log10000\le10\). Hence \(\log(1+6\sqrt{frq})\le\log(7\sqrt{frq})\le2+\frac12(\log f+\log r+\log q)\le2+37L\), and \(\log m\le\log r\le L\) gives \(p_2=\log2+d_*\log(6fm)\le1+2fr(2+3L)\), so \(p_3\le2r+fr^2(8+43L)\le fr^2(10+43L)\). With \(R=5\) (\(2R^4=1250\), \(2fr^2\log5\le4fr^2\)), \(A=625fmr\le625fr^2\) and \(\log(AK_*)=\log(10000rf^3m^3R_E^7)\le10+24L\), we obtain \(P_0\le23750fr^2+68750fr^2L\). Since \(f\ge1\) and \(L\ge2\) imply \(4fr^2\le(frL)^2\) and \(2fr^2L\le(frL)^2\), this gives \(P_0\le200000(frL)^2\).

For (2), Gram--Schmidt orthogonalisation in the signature flag produces an upper-triangular Borel matrix \(u_{\rm B}\) such that \(x=yu_{\rm B}\) lies in the domain on \(k\) with all distinguished Gram minors of modulus \(1\), preserving \(\lambda_1\) and hence \(\Phi\). Because \(f_0\) has section weight \((0,r)\) and its square \(f_0^2\) has section weight \((0,2r)\) (half-weight exponent \(2r\), central parameter \(r\)) with zero tuple weight, the Lie scale factor of \((0,2r)\) is \(1\) and the holomorphic Borel transformation law gives \(\|f_0^2(y)\|=\|f_0^2(x)\|=|f_0(x)|^2\). At the unit-minor frame \(x\), the analytic uniformisation supplies the period map \(\operatorname{Pm}:\Lambda_k\to\C^{fr}\), the Appell--Humbert Hermitian form \(H_m\), a holomorphic trivialisation of the pulled-back principal line with multiplier \(\chi_{\rm lin}(\operatorname{Pm}(u))=(-1)^{E_F(u,w_0u)}\) (where \(E_F\) is the Riemann form and \(2w_0=1+\sqrt{-d_H}\) in \(\mathcal O_E\) for the host prime \(d_H\)), and an entire principal theta function \(\theta_0:\C^{fr}\to\C\) of unit \(L^2\) norm against \(e^{-\pi\Re H_m(z,z)}(\operatorname{vol}\mathcal F_0)^{-1}\,dz\) on a fundamental domain \(\mathcal F_0\) for \(\operatorname{Pm}(\Lambda_k)\). Diagonalising \(H_m\) on the \(\mathcal O_E\)-eigenspaces yields a \(\C\)-linear isometry \(\psi:(\C^{fr},H_m)\xrightarrow{\sim}(V,H)\) onto a polarised CM period torus \(\mathcal T=(V,L_0,H,\varepsilon)\) of rank \(r\) with \(L_0=\psi(\operatorname{Pm}(\Lambda_k))\), \(d_\tau\) taking the values \(p_{\rm sig}\) and \(q_{\rm sig}\) on the CM type and its conjugate (so \(\min_\tau d_\tau=\min(p_{\rm sig},q_{\rm sig})\)), \(\lambda_1(\mathcal T)=\lambda_1(x)=\lambda_1(y)\) and \(\phi_{\theta(q,R_0)}(\mathcal T)=\Phi\), on which \(\theta_{\mathcal T}=\theta_0\circ\psi^{-1}\) is a unit-\(L^2\) principal theta section with point-norm square \(|\theta_0(z)|^2\exp(-\pi\Re H_m(z,z))\) at \(\psi(z)\).

Let \(\operatorname{ev}\) be the geometric full level-\(R^4\) structure induced by the level-\(M\) datum and theta rule \(f_0\) at the complex point of \(x\), and choose period representatives \(a:\Lambda/R^4\Lambda\to\Lambda_k\) with \(\operatorname{ev}(u)=\exp(R^{-4}\operatorname{Pm}(a(u)))\). Because \(R^2\) is odd and \(\operatorname{ev}(u)^{(R^2)^2}=1\), each \(\operatorname{ev}(u)\) admits a unique normalised \(R^2\)-null-value \(w(u)\), and single-term evaluation of the theta rule gives \(f_0(x)=\prod_{j=1}^r w(x_j)\). Put \(z_j=R^2R^{-4}\operatorname{Pm}(a(x_j))\in\C^{fr}\) and \(v_j=\psi(z_j)\in V\). Then \(R^2v_j=\psi(\operatorname{Pm}(a(x_j)))\in L_0\), and the \(R^2\)-null-value comparison gives \(|w(x_j)|=|\theta_0(z_j)|\exp(-\frac\pi2\Re H_m(z_j,z_j))\). Because \((x_1,\ldots,x_r)\) spans \(\Lambda/R^4\Lambda\) over \(\mathcal O_E\), every \(u\in\Lambda_k\) has the form \(\sum_{j=1}^r c_j a(x_j)+R^4v\) with \(c_j\in\mathcal O_E\) and \(v\in\Lambda_k\); applying \(\psi\circ\operatorname{Pm}\) expresses every vector of \(L_0\) as \(\sum_{j=1}^r c_j(R^2v_j)+R\cdot\psi(\operatorname{Pm}(R^3v))\), so \(R^2v_1,\ldots,R^2v_r\) span \(L_0/RL_0\) over \(\mathcal O_E\).

Assume now that \((E,r,D,q,R_0,m)\) is uniformly admissible and \(\min(p_{\rm sig},q_{\rm sig})\le m\), so \((q,R_0,m)\) is admissible on \(\mathcal T\). When \(\Phi=0\), we have \(\lambda_1(\mathcal T)\ge\theta(q,R_0)\) with \(b_0(q,R_0)\ge1\), and the pointwise Gaussian theta bound gives \(|w(x_j)|\le e^{p_1}\), whence \(|f_0(x)|^{2R^4}\le e^{2R^4rp_1}\le\exp(\mathcal P(E,r,q,m,R,1))\). When \(\Phi>0\), we have \(\lambda_1(\mathcal T)<\theta(q,R_0)\) and \(\Omega_*\theta(q,R_0)^2\le d_*^2\Omega_*\theta(q,R_0)^2<1\). Lemma~\ref{B8.25q} applied to \((V,L_0,H)\) produces a saturated \(\mathcal O_E\)-stable isotropic sublattice \(I\subset L_0\) of \(\Z\)-rank \(j_I\le d_*\) and real symplectic projection \(u:V\to W=I\otimes_{\Z}\R\). Expanding \(\theta_{\mathcal T}\) in Fourier modes along the real torus of \(I\) expresses the modes as translates of a quotient principal theta section \(\theta_I\) with Gaussian coset sums. If \(c_I\) is the Euclidean covolume of \(I\) in \(W\), unit \(L^2\) norm of \(\theta_{\mathcal T}\) gives \(c_I2^{-j_I/2}\|\theta_I\|_2^2\le1\), the quotient supremum bound gives \(\|\theta_I\|_{\sup}\le e^{p_1}\|\theta_I\|_2\), and the isotropic covolume bound gives \((c_I2^{-j_I/2})^{-1/2}\le e^{d_*}(K_*\Phi)^{d_*/4}\). Choosing \(y_j\in L_0\) to minimise \(D_j=\|u(v_j-y_j)\|^2\) bounds the Gaussian coset sum at \(v_j\) by \(2(3d_*)^{d_*}e^{-\pi D_j/2}\). Meanwhile \(R^2(v_j-y_j)\equiv R^2v_j\pmod{RL_0}\) since \(R^2y_j\in RL_0\), so the \(\mathcal O_E\)-multiples of \(R^2(v_1-y_1),\ldots,R^2(v_r-y_r)\) span \(L_0/RL_0\) and hence span \(V\) over \(\R\); nondegeneracy of \(\Im H\) on \(V\times W\) shows that their images under \(u\) span \(W\) over \(\R\), and Lemma~\ref{B8.25q} gives \(\Phi\le\frac12\sum_{j=1}^r\|u(R^2(v_j-y_j))\|^2=\frac12 R^4\sum_{j=1}^r D_j\). Combining these factors yields
\begin{equation}\label{eq:p8c3-theta-factor}
 |w(x_j)|\le
 e^{p_1}\,2(3d_*)^{d_*}\,e^{d_*}(K_*\Phi)^{d_*/4}e^{-\pi D_j/2},
 \qquad
 \Phi\le\frac{R^4}{2}\sum_{j=1}^r D_j .
\end{equation}
Multiplying over \(1\le j\le r\), inserting \(R^{2fr^2}\ge1\) and raising to \(2R^4\) gives
\[
 |f_0(x)|^{2R^4}\le\exp\bigl(2R^4(p_3+2fr^2\log R)+A\log(K_*\Phi)-2\pi\Phi\bigr).
\]
Using \(\log(\Phi/A)\le\Phi/A-1<\Phi/A\) gives \(A\log(K_*\Phi)\le A\log(AK_*)+\Phi\), and \(2\pi-1\ge1\) then gives \(|f_0(x)|^{2R^4}\le\exp(\mathcal P(E,r,q,m,R,1)-\Phi)\). Taking \(R^4\)-th roots completes the proof.
\end{proof}

\subsection*{Proof of \ref{B8.25gc}}\label{prf:B8.25gc}\phantomsection
\begin{proof}
For (1), since \(y_0\) lies in the signature domain on \(k_0\), every distinguished Gram minor at \(y_0\) is nonzero, so \(P_{k_0}(y_0)\ne0\) gives \(A=\|P(y)\|_\lambda>0\). The ratio hypothesis bounds \(\sup_z\|P(z)\|_\lambda\) over domain frames \(z\) on \(U\) by \(e^{-1}A\). Because \(\mu_Z\) has total mass at most \(1\) and \(\pi_\sigma\) lands in the signature domains, \(\|P\|_{2,Z}\le\sup_z\|P(z)\|_\lambda\le e^{-1}A\). Hence \(e^2\|P\|_{2,Z}^2\le A^2\), and since \(A>0\) this gives \(\|P(y)\|_\lambda^2/\|P\|_{2,Z}^2\ge e^2\) in \([0,\infty]\).

For (2), choose an orthogonal frame \(x_{0,k}\) and a measurable fundamental set \(\mathcal F_k\) on each analytic component \(k\) of \(U\), apply the hypothesis to \((x_{0,k_0},\mathcal F_{k_0})\) to obtain \(P_0\), set \(P'_{k_0}\) equal to \(P_0\) on the signature domain of \(k_0\) and \(0\) outside it, and set \(P'_k=0\) for \(k\ne k_0\). Then \(P'\) is a holomorphic cuspidal family of weight \(\lambda\) with \(\|P'(y)\|_\lambda=\|P_0(y_0)\|_{\lambda_{k_0}}>0\). Unfolding the adelic quotient measure \(\mu_{Q,1}\) into the fundamental sets \(\mathcal F_k\) yields a finite constant \(c_{\rm vol}>0\) with \(\mu_F(\mathcal F_k)<\infty\) such that
\[
 \int_{Q_W}\|P'(\pi_\sigma(q))\|_\lambda^2\,d\mu_{Q,1}(q)
 =c_{\rm vol}\int_{\mathcal F_{k_0}}\|P_0(x_{0,k_0}n(z))\|_{\lambda_{k_0}}^2\,d\mu_F(z)
\]
and \(\mu_{Q,1}(\{q\in Q_W:\operatorname{comp}(\pi_\sigma(q))=k_0\})=c_{\rm vol}\mu_F(\mathcal F_{k_0})\). Because \(y=(k_0,y_0)\) lies over \(Z\), every domain frame on the connected analytic component \(k_0\) lies over \(Z\), so \(\{q\in Q_W:\operatorname{comp}(\pi_\sigma(q))=k_0\}\subseteq\mathcal C_Z\) and \(c_{\rm vol}\mu_F(\mathcal F_{k_0})\le\mu_{Q,1}(\mathcal C_Z)\). If \(\mu_F(\mathcal F_{k_0})=0\), then \(\|P'\|_{2,Z}=0\) and \(\|P'(y)\|_\lambda^2/\|P'\|_{2,Z}^2=+\infty\ge\beta\). If \(\mu_F(\mathcal F_{k_0})>0\), then \(\mu_{Q,1}(\mathcal C_Z)\ge c_{\rm vol}\mu_F(\mathcal F_{k_0})>0\), so
\[
 \|P'\|_{2,Z}^2
 \le\mu_{Q,1}(\mathcal C_Z)^{-1}\int_{Q_W}\|P'(\pi_\sigma(q))\|_\lambda^2\,d\mu_{Q,1}(q)
 \le\frac{\int_{\mathcal F_{k_0}}\|P_0(x_{0,k_0}n(z))\|_{\lambda_{k_0}}^2\,d\mu_F(z)}{\mu_F(\mathcal F_{k_0})}.
\]
Multiplying by \(\beta\) and using the assumed integral bound gives \(\beta\|P'\|_{2,Z}^2\le\|P'(y)\|_\lambda^2\), whence \(G(y)\ge\beta\).
\end{proof}

\subsection*{Proof of \ref{B8.25gca}}\label{prf:B8.25gca}\phantomsection
\begin{proof}[Proof (outline)]
\begin{enumerate}
\item Take the supremum over all domain frames on all analytic components at the fixed embedding:
\[
 S=\sup_z\|s(z)\|_{(a,N)}.
\]
Cuspidality gives \(S<\infty\), and every point norm is at most \(S\). The coefficient bound gives \(\rho\ge0\). Since \(f,r\ge1\), \(r\le N\) and \(frN\le2\kappa\), one has \(\kappa>0\). If \(S=0\), the conclusion follows.

\item Suppose \(S>0\), and put \(A=e^L\). The parameter assumptions give
\[
 A\ge2,\qquad M\le A^{30},\qquad
 R_E\le A^3,\qquad fr\le A,\qquad \kappa\le A^{20}.
\]
The host comparison gives \(\sqrt{d_H}\le R_E\), where \(d_H\) is the positive prime parameter of the host. The sum of the weight spreads is at most
\[
 D+frN\le3\kappa.
\]
Using this spread bound and the trivial stabiliser character, the maximum-location construction at parameter \(3\kappa\) supplies an orthonormal domain frame \(x\) and a positive real number \(\varepsilon\) such that
\[
 S\le2|s(x)|,\qquad
 \lambda_1(x)^2\ge\frac{\varepsilon}{R_E^2},
 \qquad \varepsilon\ge A^{-57}.
\]
Here \(\lambda_1(x)\) is the first minimum in the positive majorant. The last inequality uses the bounds on \(M,\sqrt{d_H},fr,\kappa\) displayed above. Since
\[
 A^{-80}R_E^2\le A^{-80}A^6\le A^{-57},
\]
it follows that
\[
 \lambda_1(x)^2\ge A^{-80}=e^{-80L}.
\]

\item At this frame, put \(d=frN\) and \(\lambda=A^{-40}\). The differentiated theta estimate with polarisation scale \(2\xi_{\R}\), minimum parameter \(\lambda\), and coefficient sum at most \(\rho\) gives
\[
 (2\pi\xi_{\R})^{-D/2}S
 \le2\sqrt2\,\rho\,(eD)^{D/2}
 \left(1+\frac{8f^{3/2}N^2R_E\sqrt{2d}}{\sqrt\pi\,\lambda}\right)^d.
\]
The degree factor is interpreted as one when \(D=0\). Since \(frN\ge1\) and \(frN\le2\kappa\), one has \(\kappa\ge1/2\). Using \(f\le A\), \(N\le2\kappa\le2A^{20}\), \(R_E\le A^3\), and \(\sqrt{2d}\le2A^{10}\), the parenthesised factor is at most
\[
 1+64A^{95}\le A^{102}.
\]
If \(D\ge1\), then \(\log D\le\log\kappa\le20L\), so \((eD)^{D/2}\le e^{11\kappa L}\); the same upper bound holds when \(D=0\). Also \(A^{102d}\le e^{204\kappa L}\) and \(2\sqrt2\le e^{4\kappa L}\). Consequently
\[
 (2\pi\xi_{\R})^{-D/2}S
 \le\rho e^{219\kappa L}
 \le\rho\exp\!\bigl(9000(frL)^2\kappa\bigr),
\]
using \(fr,L\ge1\). The pointwise conclusion follows from \(\|s(z)\|_{(a,N)}\le S\).
\end{enumerate}
The complete proof is contained in \Lean{UpperTheory.\allowbreak{}Dv.\allowbreak{}Thin.\allowbreak{}theta\_\allowbreak{}sup\_\allowbreak{}le}.
\end{proof}

\subsection*{Proof of \ref{B8.25gcb}}\label{prf:B8.25gcb}\phantomsection
\begin{proof}[Proof (outline)]
\begin{enumerate}
\item Replace \(y\) by an orthogonal frame \(x\) on the same analytic component. The first minimum and the invariant norms of the sections \(s_u\) are unchanged. The level-geometry laws give the analytic uniformisation, its realisation by a test family, and the identity \(|\tau(\xi_{0,k})|=\xi_{\R}\) at every archimedean embedding. Choose a period uniformisation and a unit theta trivialisation at \(x\). Its Hermitian matrix is \(2\xi_{\R}\) times the identity. The period-coordinate isometry, with the embedding relabelling determined by \(\sigma\), identifies this with a polarised \(\mathcal O_E\)-torus \(T\) of rank \(r\), preserves the first minimum and truncated phase, and carries \(P\) to a homogeneous block polynomial \(P_T\) of degree \(D\) with the same disjoint column supports. The normalised principal theta section \(s\) on \(T\) has squared \(L^2\)-norm one and is nonzero. Multiplication by \(\alpha\) scales its Hermitian form by \(q\). Uniform admissibility, applied to \(T\) and \(q_{\rm sig}\le m\), supplies the admissible tuple of Definition~\ref{B8.25ga} with these \(q,R_0,m\) and phase \(\Phi_D(y)\); in particular \(D\ge1\).

\item Write \(L_T\) for the period lattice and \(\operatorname{mk}:\Lambda\to L_T\) for the marking. Choose representatives \(x_i(u)\in2\overline{\mathfrak b}^{\,2}\Lambda\) congruent to the characteristics \(u_i\) modulo \(\mathfrak b\Lambda\). The marking modulo \(q\) and the equality \(\mathfrak b=(\bar\alpha)\) give a finite set \(K\) of representatives of
\[
 (q^{-1}L_T)^N/(\alpha^{-1}L_T)^N,
\]
each of whose elements is of the form
\[
 b_i(u)=-\bar\alpha q^{-2}\operatorname{mk}(x_i(u)).
\]
Put \(s_0(z_1,\ldots,z_N)=\prod_{i=1}^Ns(\alpha z_i)\), and denote by \(U_q(b)s_0\) its theta translate with the Gaussian multiplier at level \(q\). The Gaussian torus estimate for disjoint column blocks, applied to the admissible tuple and \(\|s\|_2^2=1\), gives the precise intermediate inequality
\[
 J_y\le\sum_{b\in K}
 \left|\left.P_T(\partial)\bigl(U_q(b)s_0\bigr)\right|_{0}\right|^2.
\]
Here its constant is exactly \(4^{-N}K_6^{-N}2^{-(r+2)D}
\exp\{-2\pi NC_\chi\Phi_D(y)/q\}\). Every factor is positive since \(d_*>0\), so \(J_y>0\).

\item Compare theta-rule evaluation in the test family with differentiation in period coordinates. Homogeneity produces an \(\alpha\)-scaling scalar of squared modulus \(q^D\), while each period-coordinate direction contributes \(\sqrt{2\xi_{\R}}\). The resulting squared-value identity is
\[
 |s_u(x)|^2=
 \left(\frac{2\xi_{\R}}q\right)^D
 \left|\left.P_T(\partial)\bigl(U_q(b(u))s_0\bigr)\right|_{0}\right|^2.
\]
Every \(b\in K\) has a characteristic preimage. Choose one preimage for each \(b\); they are distinct, and all remaining squared values are nonnegative. Summing the identity therefore gives
\[
 \left(\frac{2\xi_{\R}}q\right)^D J_y
 \le\sum_u |s_u(x)|^2.
\]
The chosen orthogonal frame makes \(|s_u(x)|\) equal to the invariant point norm at \(y\). Multiplication by \(q^D(2\pi\xi_{\R})^{-D}\) proves the asserted sum bound.

\item Put \(S=q^D\sum_u v_u(y)^2\). Its lower bound is positive. Expanding the logarithm of the displayed constant gives
\[
 \begin{split}
 -\log\sqrt{\pi^{-D}J_y}
 &=\frac D2\log\pi+
 \frac N2\left(\log24+\frac{d_*}{2}\log(2d_*)\right)
 +\frac{(r+2)D}{2}\log2+c_0\Phi_D(y)\\
 &\le A_0D+c_0\Phi_D(y)
 \le AD+c\phi(y).
 \end{split}
\]
The first inequality uses \(D\ge1\) and nonnegativity of the parenthesised logarithmic term; the second uses \(0\le\Phi_D(y)\le\phi(y)\) and \(c\ge c_0\ge0\). Monotonicity of the logarithm now gives \(\sqrt S\ge e^{-AD-c\phi(y)}\). There are exactly \(q^{frN}\) characteristic tuples, so a tuple attaining the largest \(v_u(y)\) satisfies
\[
 v_u(y)\ge q^{-(D+frN)/2}\sqrt S
 \ge q^{-(D+frN)}e^{-AD-c\phi(y)},
\]
where the last inequality uses \(q\ge1\). This proves the required positive lower bound.
\end{enumerate}
The complete proof is contained in \Lean{UpperTheory.\allowbreak{}Dv.\allowbreak{}Thin.\allowbreak{}theta\_\allowbreak{}lower}.
\end{proof}

\subsection*{Proof of \ref{B8.25g}}\label{prf:B8.25g}\phantomsection
\begin{proof}
Write \(Y\) for the good cover, \(P=P_*\) and \(L=L_x\). By \eqref{eq:p8c3-anchor}, \eqref{eq:p8c3-host-size}, \eqref{eq:p8c3-original-field} and Lemmas~\ref{B8.25b} and \ref{B8.25c}, beyond a common family threshold the host, scale and cover parameters satisfy
\begin{equation}\label{eq:p8c3-berg-sizes}
 \begin{gathered}
 L\ge10,\quad b\ge1000,\quad 1\le m\le r,\quad
 \log R_E,\log f\le2L,\quad \log r,\log(fr)\le L,\\
 b^{16}\le Q_0\le b^{70},\quad Q_0b\le M,\quad
 m^4\le f,\quad f^2\le R_E,\quad 16\le R_E,\quad R_E^7\le b,\\
 (frL)^{12}\le P,\quad (frL)^{24}\le T,\qquad
 PTd_{\rm fl}^{\,2}L\le b.
 \end{gathered}
\end{equation}
Here the base primes \(g_0,g_1\) lie between \(b^8\) and \(b^{35}\); for the host conductor \(Q_Sqd\) of \ref{B2.46}, the inclusion \(E\subset\Q(\mu_{Q_Sqd})\) gives \(R_E\le Q_Sqd\), and the host size condition \((Q_Sqd)^{36}t^{150}\le\sqrt b\) of Definition~\ref{B8.01k} (with chain rank parameter \(t\ge1\)) gives \(R_E^7\le R_E^{28}\le b\). In Lemma~\ref{B8.25c}, \(M=Q_0\cdot15^e\) with \(e=\lfloor3L/\log3\rfloor\); since \(0<\log3\le2\) and \(L\ge10\), we have \(e>3L/\log3-1\ge14\), so \(e\ge15\ge4\), \(5^4\mid15^e\mid M\), and \(e\log3>3L-2\ge L\), whence \(15^e\ge3^e>e^L=b\) and \(Q_0b\le M\). Moreover, for every integer \(1\le D\le b^2\), \eqref{eq:p8c3-berg-sizes} gives \(R_E^{28}=(R_E^7)^4\le b^4\), \(r\le fr\le P\le b\) (since \(T,d_{\rm fl},L\ge1\)), \(90(fr+D)\le180b^2\) and \(b\ge1000\ge180\), so
\[
 90R_E^{28}(fr+D)Dr\le b^4(180b^2)b^2b=180b^9\le b^{10}\le b^{16}\le Q_0.
\]
Summing the floored differences in \eqref{eq:p8c3-quotient} from the anchor \(u_{1,\tau,k-2}=0\) with the family weight-gap bounds (cf.\ Lemma~\ref{B8.25b}) gives adjacent gaps \(u_{1,\tau,s}-u_{1,\tau,s+1}\ge T\) for \(0\le s<r-1\). Since \(p_{\rm sig}=k+j-1\ge k=(k-2)+2\), every index \(s\ge p_{\rm sig}\) satisfies \(u_{1,\tau,s}\le-2T<-3r\), so \(u'=u_1-3r\eta_H\) keeps \(u'_{\tau,s}=u_{1,\tau,s}+3r<0\) on \(s\ge p_{\rm sig}\) and agrees with \(u_1\) on \(s<p_{\rm sig}\), giving \(\|u'\|_1\le\|u_1\|_1\). With the width parameters \(S,\eta\ge0\) and window scale bound \(w''\le8b\) of Lemma~\ref{B8.25b}, the \(k-2=r-m-1\) indices \(s<k-2\) satisfy \(0\le u_{1,\tau,s}\le(S+2r)/T\) and the \(m\) indices \(s>k-2\) satisfy \(|u_{1,\tau,s}|\le(8b+c_0+4r)/T\) (with \(2(k-2)r+4mr=2(r+m-1)r\le4r^2\) since \(m\le r\)). Using \(P(S+\eta)\le100w''\le800b\), \(r\le P\), \(c_0\le22f^2r^4\), \(26f^2r^5\le26P\le PTd_{\rm fl}^{\,2}L\le b\) and \eqref{eq:p8c3-anchor} gives
\begin{equation}\label{eq:p8c3-row-sum}
 T\|u'\|_1\le T\|u_1\|_1
 \le f\{rS+8mb+r(c_0+4r)\}\le834fmb\le834E_p .
\end{equation}
\leanref{GlobalTheory.\allowbreak{}CovTop.\allowbreak{}Sizes.\allowbreak{}hostNum\_\allowbreak{}of\_\allowbreak{}sizes}{Development.\allowbreak{}AbelianVarietyTheory.\allowbreak{}Proofs.\allowbreak{}DivisorProofs.\allowbreak{}HostBoundHolds}{37--109}
\leanref{GlobalTheory.\allowbreak{}CovTop.\allowbreak{}WtShape.\allowbreak{}wtNum\_\allowbreak{}at}{Development.\allowbreak{}AbelianVarietyTheory.\allowbreak{}Proofs.\allowbreak{}DivisorProofs.\allowbreak{}WeightBoundHolds}{45--186}
\leanref{GlobalTheory.\allowbreak{}ThetaSizeBounds.\allowbreak{}level\_\allowbreak{}ge}{Development.\allowbreak{}AbelianVarietyTheory.\allowbreak{}Proofs.\allowbreak{}DivisorProofs.\allowbreak{}EtaBoundLengthUnit}{139--162}
\leanref{GlobalTheory.\allowbreak{}ThetaSizeBounds.\allowbreak{}size\_\allowbreak{}inequality\_\allowbreak{}real}{Development.\allowbreak{}AbelianVarietyTheory.\allowbreak{}Proofs.\allowbreak{}DivisorProofs.\allowbreak{}EtaBoundLengthUnit}{72--121}
\leanref{GlobalTheory.\allowbreak{}CovHt.\allowbreak{}DividedNum.\allowbreak{}lenT\_\allowbreak{}bounds\_\allowbreak{}of\_\allowbreak{}covers\_\allowbreak{}before\_\allowbreak{}last}{Development.\allowbreak{}AutomorphicArithmetic.\allowbreak{}ThetaCorrespondence.\allowbreak{}ThetaPairLifts.\allowbreak{}ThirdPlaceSeries.\allowbreak{}DividedWeight}{18--24}

Let \(C_{\rm seed},C_S\ge0\) be the absolute constants of Lemmas~\ref{B8.25r} (for a fixed cutoff) and \ref{B8.25s}, set \(C_{\rm peak}=C_{\rm seed}+C_S+1\), and choose the common family threshold before fixing the point and layer so that \((frL)^{24}\) exceeds both \(C_{\rm peak}(fr)^8\) and the gap threshold of Lemma~\ref{B8.25r} at \(\rho_*=1/(200r^2)\), \(b_*=4fr^2+1\) and \(d_{\C}=fr(r-1)/2\). Put \(P_0=\mathcal P(E,r,Q_0,m,5,1)\) in Definition~\ref{B8.25ga} and \(B'_0=C_{\rm peak}(\|u_1\|_1+(fr)^8)\ge0\). By \eqref{eq:p8c3-berg-sizes} and \eqref{eq:p8c3-host-size} (which with \(r=p_{\rm sig}+q_{\rm sig}\ge2\) give \(fr\ge8\cdot2\ge2\) and \(\log Q_0\le70L\)) and Lemma~\ref{B8.25gb}(1), \(R_0(D)=\sqrt{90(fr+D)}>0\) makes \((E,r,D,Q_0,R_0(D),m)\) uniformly admissible for every integer \(1\le D\le b^2\), and \(0\le P_0\le200000(frL)^2\). Since \((frL)^2,(fr)^8\le(frL)^{12}\le P\), \(d_{\rm fl}\ge1\), \(L\ge10\) and \(b\le E_p\) by \eqref{eq:p8c3-size-main}, \eqref{eq:p8c3-berg-sizes} and \eqref{eq:p8c3-row-sum} give \(TP_0\le200000\,TP\le200000b\le200000E_p\) and \(TB'_0\le C_{\rm peak}(834E_p+E_p)=835C_{\rm peak}E_p\), proving \eqref{eq:p8c3-berg-budget}. At a nonterminal layer \(q_{\rm sig}=m-j\le m\), so \(\min(p_{\rm sig},q_{\rm sig})\le m\); also \(5>1\) is odd and \(5^4\mid M\). Applying Lemma~\ref{B8.25gb}(2) on \(Y\) with \(D=1\), \(q=Q_0\), \(R=5\) (\(R^4=625\)) and \(R_0=R_0(1)=\sqrt{90(fr+1)}>0\) gives \(a=\|f_0^2(y)\|\le\exp((P_0-\Phi)/625)\le\exp(P_0/625)\) everywhere since \(\Phi\ge0\). If \(M\lambda_1^2<b\), then \(Q_0b\lambda_1^2\le M\lambda_1^2<b\) by \eqref{eq:p8c3-berg-sizes}, so \(Q_0\lambda_1^2<1\le fR_0^2R_E^{21}\), \(\lambda_1<\theta(Q_0,R_0)\), \(\Phi=\phi\), and Lemma~\ref{B8.25gb}(2) gives \(a\le\exp((P_0-\phi)/625)\).
\leanref{GlobalTheory.\allowbreak{}ThetaSizeBounds.\allowbreak{}sqrt\_\allowbreak{}lt\_\allowbreak{}th\_\allowbreak{}real}{Development.\allowbreak{}AbelianVarietyTheory.\allowbreak{}Proofs.\allowbreak{}DivisorProofs.\allowbreak{}EtaBoundLengthUnit}{122--138}

Suppose now that \(M\lambda_1^2<b\), and transport \(y\) along the covering map from \(Y\) to a domain frame \(\bar y\) at base level \(Q_0\). The transport preserves \(\lambda_1\), so \(Q_0\lambda_1^2<1\). Let \(\Psi\) be the CM type of size \(f\), let \(h_\sigma\) be the Hodge multiplicity at \(\sigma:E\to\C\) (\(p_{\rm sig}\) on \(\Psi\) and \(q_{\rm sig}\) on \(\overline\Psi\)), index pairs \((\sigma,i)\) by \(1\le i\le h_\sigma\), write \(v=fr\) for their number, and define the column count
\[
 N=8\left(\left\lfloor\frac18\sum_{\sigma,i}2i\right\rfloor+v+1\right).
\]
Then \(8\mid N\), \(N\ge\sum_{\sigma,i}2i\ge2fr\ge2r\), \(2N\ge fr(r+1)\), and \(N\le2fr^2+8fr+8\le20fr^2\), whence \(frN\le20f^2r^3\le P\le b\) by \eqref{eq:p8c3-berg-sizes}. With zero-based column indices on \(u_1\), put \(s_\tau=N/2+1-u_{1,\tau,p_{\rm sig}-1}\) for \(\tau\in\Psi\). Set \(a_{\tau,i}=(u_{1,\tau,i-1}+s_\tau)-N/2=u_{1,\tau,i-1}-u_{1,\tau,p_{\rm sig}-1}+1\) for \(1\le i\le p_{\rm sig}\) with \(a_{\tau,p_{\rm sig}+1}=0\), and at \(\bar\tau\) reverse and negate the shifted second block and subtract \(N/2\) to set \(a_{\bar\tau,i}=-(u_{1,\tau,r-i}+s_\tau)-N/2=u_{1,\tau,p_{\rm sig}-1}-u_{1,\tau,r-i}-N-1\) for \(1\le i\le q_{\rm sig}\) with \(a_{\bar\tau,q_{\rm sig}+1}=0\). Because \(u_{1,\tau}\) has adjacent gaps at least \(T\ge(frL)^{24}>N+1\), each row \(a_\sigma\) is nonincreasing and nonnegative with \(a_{\tau,p_{\rm sig}}=1>0\) on \(\Psi\). Put \(c_{\sigma,i}=a_{\sigma,i}-a_{\sigma,i+1}\ge0\) and \(D=\sum_{\sigma,i}ic_{\sigma,i}=\sum_{\sigma,i}a_{\sigma,i}\). Then \(D\ge f\ge1\), and \eqref{eq:p8c3-row-sum} with \eqref{eq:p8c3-berg-sizes} gives \(\|u_1\|_1<b\) and \(D\le(r+1)\|u_1\|_1+fr\le b^2\), so \((E,r,D,Q_0,R_0(D),m)\) is uniformly admissible. Since \(\sum_{\sigma,i}i\le N\), we may assign pairwise disjoint sets of \(i\) columns in \(\{1,\ldots,N\}\) to the pairs \((\sigma,i)\) and multiply the leading \(i\times i\) minors of \(\sigma\)-variables in those columns raised to the powers \(c_{\sigma,i}\). Specialising the assigned blocks to identity matrices gives value \(1\), while every Laplacian contraction differentiates \(\tau\)- and \(\bar\tau\)-variables in the same column and therefore annihilates the product; this defines a nonzero integral harmonic differential of theta weight \((a,N)\) and degree \(D\), whose \(\ell^1\) coefficient sum is at most \(\prod_{\sigma,i}(i!)^{c_{\sigma,i}}\le r^D\le e^{Db}\).

Let \(\alpha\in\mathcal O_E\) be the split half-level generator of the base level, with \(\alpha\bar\alpha=Q_0\) and half-level ideal \((\bar\alpha)\). Index the \(Q_0^{frN}\) half-level characteristic tuples by \(u\), and let \(s_u\) be the associated theta derivative of weight \((a,N)\). The polynomial is harmonic, \(N\) is even, and \(a_{\tau,p_{\rm sig}}=1>0\) on \(\Psi\), so Lemma~\ref{B8.25t} makes every \(s_u\) cuspidal. Put \(v_u(\bar z)=(2\pi\xi_\R)^{-D/2}\|s_u(\bar z)\|_{(a,N)}\), \(R_0=R_0(D)\), and let \(\Phi_D\) be the truncated phase at \(\bar y\) with threshold \(\theta(Q_0,R_0)\). The disjoint column blocks constructed above and \(q_{\rm sig}\le m\) verify the hypotheses of Lemma~\ref{B8.25gcb}. With its constants \(d_*=2fm\), \(K_6=6(2d_*)^{d_*/2}\) and \(C_\chi=8100m^4d_*^4R_E^7\), that lemma gives
\begin{equation}\label{eq:p8c3-derivative-test}
 J=4^{-N}\{6(2d_*)^{d_*/2}\}^{-N}2^{-(r+2)D}
 \exp\{-2\pi N(8100m^4d_*^4R_E^7)\Phi_D/Q_0\}>0
\end{equation}
and
\begin{equation}\label{eq:p8c3-derivative-sum}
 Q_0^D\sum_u v_u(\bar y)^2\ge\pi^{-D}J .
\end{equation}
Here \(N,r\le b\), \(d_*=2fm\le2fr\le2b\), and \(b\ge1000\), so the constant \(A_0\) of Lemma~\ref{B8.25gcb} satisfies
\[
 A_0\le\frac12\{\log\pi+b(\log24+b\log(4b))+(b+2)\log2\}
 \le3b^3.
\]
For the last inequality one may use \(\log\pi\le2\), \(\log24\le4\), \(\log2\le1\) and \(\log(4b)\le2b\). Moreover, \eqref{eq:p8c3-berg-sizes} gives \(m^4d_*^4=16f^4m^8\le(frL)^{12}\le b\), \(R_E^7\le b\) and \(Q_0\ge b^{16}\), whence
\[
 \frac{\pi NC_\chi}{Q_0}\le\frac{8100\pi b^3}{b^{16}}<\frac1{2500}.
\]
Apply the selected-characteristic conclusion of Lemma~\ref{B8.25gcb} with \(A=3b^3\) and \(c=1/2500\). It gives a tuple \(u\) with \(s_u(\bar y)\ne0\) satisfying the first bound in
\begin{equation}\label{eq:p8c3-derivative-ratio}
 \begin{split}
 v_u(\bar y)&\ge\exp\{-3b^3D-\phi/2500-(D+frN)\log Q_0\},\\
 v_u(\bar z)&\le\exp\{Db+9000(3fr\log b)^2b^2\},
 \end{split}
\end{equation}
The second bound follows from Lemma~\ref{B8.25gca} with \(L_{\rm up}=3\log b\), \(\kappa=b^2\) and \(\rho=e^{Db}\). Indeed, \(L_{\rm up}\ge1\), \(r\le N\), \(\log Q_0\le70\log b\le30L_{\rm up}\), \(\log R_E,\log(fr)\le\log b\le L_{\rm up}\), \(D\le\kappa\), \(frN\le b\le2\kappa\), and \(\log\kappa=2\log b\le20L_{\rm up}\). The integral harmonic polynomial has coefficient sum at most \(r^D\le e^{Db}\), and every \(s_u\) is cuspidal. From \(Q_0\lambda_1^2<1\) and \eqref{eq:p8c3-berg-sizes} we have \(\phi=2f/(R_E^2\lambda_1^2)>2fQ_0/R_E^2\ge2b^{15}\ge b^{14}\); since \(D\le b^2\), \(frN\le b\), \(\log Q_0\le70\log b\le70b\), \(fr\log b=frL\le b\) and \(b\ge1000\), the non-phase exponents in \eqref{eq:p8c3-derivative-ratio} sum to at most \(3b^5+70b(b^2+b)+b^3+81000b^6\le82000b^6\le\phi/2500\). Cancelling \((2\pi\xi_\R)^{-D/2}\) gives \(e^{-\phi/1250}\|s_u(\bar z)\|_{(a,N)}\le\|s_u(\bar y)\|_{(a,N)}\) for every base domain frame \(\bar z\). Pulling \(s_u\) back to \(Y\) preserves cuspidality, nonvanishing at \(y\) and the invariant point-norm ratio; multiplying on \(Z\) by the determinant section of central parameter \(-N/2\) (half-weight exponent \(-N\)), the flat sections of Lemma~\ref{B8.24b} of row exponents \((-s_\tau)_{\tau\in\Psi}\) and the component idempotent of \(Z\) produces a section \(g\) of weight \((u_1,0)\) supported on \(Z\) whose invariant point norm on \(Z\) is a constant \(K>0\) times that of the pullback. In each boundary Fourier--Jacobi coordinate supplied by \(\mathcal H_P\) and \(\mathcal H_H\), the determinant and flat multipliers are holomorphic of constant invariant norm and hence of at most polynomial growth under isotropic boundary shears; Cauchy's circle integral eliminates negative Fourier exponents in each multiplier, while the pulled-back cusp section has strictly positive Fourier exponents, so \(g\) is cuspidal. Since the adelic frame map of \(Y\) lands in the signature domains and \(g\) vanishes off \(Z\), applying Lemma~\ref{B8.25gc}(1) to the holomorphic cuspidal family of \(g\) with \(e=e^{-\phi/1250}\) yields \(G(y)\ge e^{-\phi/625}\).
\leanref{GlobalTheory.\allowbreak{}RefSum.\allowbreak{}GOf}{Development.\allowbreak{}AbelianVarietyTheory.\allowbreak{}Definitions.\allowbreak{}Divisors.\allowbreak{}EightDvdNCol}{74--75}
\leanref{GlobalTheory.\allowbreak{}T.\allowbreak{}thin\_\allowbreak{}divisor\_\allowbreak{}pivot\_\allowbreak{}bound\_\allowbreak{}alt}{Development.\allowbreak{}AutomorphicArithmetic.\allowbreak{}ThetaCorrespondence.\allowbreak{}ThetaPairLifts.\allowbreak{}ThirdPlaceSeries.\allowbreak{}CardMulSuccLeTwoMulNcolCopy}{149--175}
\leanref{GlobalTheory.\allowbreak{}T.\allowbreak{}cusp\_\allowbreak{}form\_\allowbreak{}bounded\_\allowbreak{}on\_\allowbreak{}frames\_\allowbreak{}alt\_\allowbreak{}alt}{Development.\allowbreak{}AutomorphicArithmetic.\allowbreak{}ThetaCorrespondence.\allowbreak{}ThetaPairLifts.\allowbreak{}ThirdPlaceSeries.\allowbreak{}CuspFormBoundedOnFrames}{36--36}
\leanref{GlobalTheory.\allowbreak{}ThinBerg.\allowbreak{}bergman\_\allowbreak{}kernel\_\allowbreak{}bound\_\allowbreak{}thin\_\allowbreak{}part\_\allowbreak{}alt}{Development.\allowbreak{}AutomorphicArithmetic.\allowbreak{}ThetaCorrespondence.\allowbreak{}ThetaPairLifts.\allowbreak{}ThirdPlaceSeries.\allowbreak{}BergmanKernelThinPartEstimate}{18--52}

Suppose finally that \(M\lambda_1^2\ge b\). Write \(y=(k_0,y_0)\) and let \(\lambda\) be the model weight of \((u_1,0)\). Reorthogonalising \(y_0=x_1u'_1\) on \(k_0\) with \(x_1\) orthogonal and \(u'_1\) upper triangular preserves \(\lambda_1^2\) and satisfies \(P_0(y_0)\ne0\) and \(\|P_0(y_0)\|_{\lambda_{k_0}}=\|P_0(x_1)\|_{\lambda_{k_0}}\) whenever \(P_0\) is a cusp form of weight \(\lambda_{k_0}\) on \(k_0\) with \(P_0(x_1)\ne0\). At a nonterminal layer both \(p_{\rm sig}\) and \(q_{\rm sig}\) are at least \(1\), and the component analytic uniformisations, adelic frame parametrisation and compatible orthogonal frames of the good cover \(Y\) verify the level-geometry laws of Lemma~\ref{B8.25gc}(2). Because \(g_0,g_1\ge b^8>3\) are distinct odd primes, the prime-to-\(\ell\) part of \(Q_0=g_0g_1\) is at least \(3\) at every prime \(\ell\), so \(Q_0\) and \(M=Q_0\cdot15^e\) are good at every prime, and \(M'=M\) since the auxiliary level parameter is \(1\). From \eqref{eq:p8c3-berg-sizes} and \(M\lambda_1^2\ge b\) we obtain \(M'\lambda_1(x_1)^2\ge b\ge100(fr)^4\). For any \(\gamma\in\Gamma_{k_0}\) carrying a point of the unitriangular chart ball \(\|z\|\le1/(100r^2)\) at \(x_1\) to a right Borel translate of another point of that ball, both points lie in the domain on \(k_0\); the good-level orbit decomposition writes \(\gamma\) up to right Borel action as a unitary lattice stabiliser \(\gamma'\equiv1\pmod{M'}\) on the trace-self-dual lattice, the congruence displacement bound forces \(\gamma'=1\), and freeness of \(\Gamma_{k_0}\) modulo the Borel subgroup gives \(\gamma=1\). Thus the radius-\(1/(100r^2)\) chart ball at \(x_1\) injects into the arithmetic quotient modulo the right Borel action.

Because the component Gram matrix is the negative transpose of the layer Gram matrix, the component domain on \(k_0\) has positive block size \(p_0=q_{\rm sig}\) (\(0<p_0<r\)) and CM type of size \(f\). Passing from \((u_1,0)\) to the model weight \(\lambda_{k_0}\) places \((u_{1,\tau,p_{\rm sig}+a})_{0\le a<q_{\rm sig}}\) in \(0\le a<p_0\) and \((u_{1,\tau,a-q_{\rm sig}})_{q_{\rm sig}\le a<r}\) in \(p_0\le a<r\), preserving \(\|\lambda_{k_0}\|_1=\|u_1\|_1\) and transferring the adjacent gaps \(\ge T\) of \(u_1\) into within-block adjacent gaps \(\lambda_{k_0,\tau,a}-\lambda_{k_0,\tau,a+1}\ge T\) for \(a+1\ne p_0\) and endpoint gaps \(\lambda_{k_0,\tau,r-1}-\lambda_{k_0,\tau,0}=u_{1,\tau,p_{\rm sig}-1}-u_{1,\tau,p_{\rm sig}}\ge T\). With \(\rho_*=1/(200r^2)\), \(b_*=4fr^2+1\) and chart dimension \(d_{\C}=fr(r-1)/2\le fr^2\), the threshold choice above ensures that \(T\ge(frL)^{24}\) exceeds both \(C_{\rm peak}(fr)^8\) and the gap threshold of Lemma~\ref{B8.25r}. Lemma~\ref{B8.25r} constructs a holomorphic seed \(F\) on the chart at \(x_1\) equal to \(1\) at \(0\), vanishing at all other \(\Gamma_{k_0}\)-translates of \(x_1\) in the chart, and satisfying
\begin{equation}\label{eq:p8c3-peak-seed}
 \int_{\rm chart}|F(z)|^2 e^{-w(z)}K_{x_1}(z)^{-1}\,d\mu_F(z)
 \le \exp\{C_{\rm seed}(\|u_1\|_1+(fr)^8)\}\mu_F(U_0),
\end{equation}
where \(w\) is the metric potential of \(\lambda_{k_0}\), \(K_{x_1}=\exp(-b_*R_{x_1})\), and \(U_0\) is the intersection of the ball \(\|z\|<1/(100r^2)\) with the chart domain (containing \(B(0,\rho_*)\)). Extending \(F\) across the lower-dimensional chart complement by local \(L^2\) removability of singularities preserves the \(\Gamma_{k_0}\)-invariant singular-weight integral from the proof of Lemma~\ref{B8.25r}, whose \(|z|^{-2d_{\C}}\) singularity at every orbit point forces the extended holomorphic section \(F\) to vanish at all nontrivial \(\Gamma_{k_0}\)-translates of \(x_1\) on the flag domain (including any in the chart complement); summing its \(\Gamma_{k_0}\)-translates with the weight character forms the Poincar\'e series \(P_F\) on \(k_0\). By the holomorphic submean inequality on fixed recentered chart balls (where the weight potential and \(K_{x_1}\) vary by bounded factors as in the proofs of Lemmas~\ref{B8.25r} and~\ref{B8.25s}), \eqref{eq:p8c3-peak-seed} gives \(\sup_w\|F(w)\|_{\lambda_{k_0}}^2 K_{x_1}(w)^{-1}\le M_F<\infty\), so \(\|F(w)\|_{\lambda_{k_0}}\le M_F^{1/2}K_{x_1}(w)^{1/2}\); since \(b_*/2=2fr^2+1/2\ge r\), the kernel-sum bound in the proof of Lemma~\ref{B8.25s} with exponent \(b_*/2\) and ball injectivity give \(\sup_y\sum_{\gamma\in\Gamma_{k_0}}K_{x_1}(\gamma y)^{1/2}<\infty\). Hence \(P_F\) is uniformly bounded in the invariant norm and therefore cuspidal on \(k_0\) by the within-block monotonicity and strict endpoint gap \(\lambda_{k_0,\tau,0}<\lambda_{k_0,\tau,r-1}\); since \(x_1\) is orthogonal and \(F\) vanishes at all nontrivial orbit points of \(x_1\), we have \(P_F(x_1)=1\) and \(\|P_F(x_1)\|_{\lambda_{k_0}}=1\). For any orthogonal frame \(x_0\) and measurable fundamental set \(\mathcal F\) on \(k_0\), Lemma~\ref{B8.25s} and \eqref{eq:p8c3-peak-seed} bound \(\int_{\mathcal F}\|P_F(x_0n(z))\|_{\lambda_{k_0}}^2\,d\mu_F(z)\) by \(\exp\{C_S(fr)^8+C_{\rm seed}(\|u_1\|_1+(fr)^8)\}\mu_F(U_0)\le e^{B'_0}\mu_F(U_0)\), and ball injectivity unfolds \(U_0\) into \(\mathcal F\) to give \(\mu_F(U_0)\le\mu_F(\mathcal F)\). Orthogonal transfer from \(x_1\) to \(y_0\) therefore verifies the hypothesis of Lemma~\ref{B8.25gc}(2) with \(\beta=e^{-B'_0}\), giving \(G(y)\ge e^{-B'_0}\).
\leanref{GlobalTheory.\allowbreak{}LawsDeep.\allowbreak{}levLaws}{Development.\allowbreak{}AbelianVarietyTheory.\allowbreak{}Proofs.\allowbreak{}DivisorProofs.\allowbreak{}ComplexUniformizationLawsOfGoodDeepCover}{21--40}
\leanref{Formalism.\allowbreak{}L2.\allowbreak{}Pj.\allowbreak{}goodAt\_\allowbreak{}of\_\allowbreak{}two\_\allowbreak{}primes}{Formalism.\allowbreak{}LocalToGlobal.\allowbreak{}ProjectiveLattices.\allowbreak{}ModelAt}{124--154}
\leanref{GlobalTheory.\allowbreak{}LawsDeep.\allowbreak{}goodAt\_\allowbreak{}of\_\allowbreak{}mul}{Development.\allowbreak{}AbelianVarietyTheory.\allowbreak{}Proofs.\allowbreak{}DivisorProofs.\allowbreak{}ComplexUniformizationLawsOfGoodDeepCover}{12--20}
\leanref{GlobalTheory.\allowbreak{}ThickBerg.\allowbreak{}orth\_\allowbreak{}transfer}{Development.\allowbreak{}AbelianVarietyTheory.\allowbreak{}Proofs.\allowbreak{}DivisorProofs.\allowbreak{}FrameExponentialVolumeSmall}{74--89}
\leanref{GlobalTheory.\allowbreak{}ThickBerg.\allowbreak{}ball\_\allowbreak{}injects}{Development.\allowbreak{}AbelianVarietyTheory.\allowbreak{}Proofs.\allowbreak{}DivisorProofs.\allowbreak{}FrameExponentialVolumeSmall}{40--62}
\leanref{GlobalTheory.\allowbreak{}ThickBerg.\allowbreak{}modelWt\_\allowbreak{}shape}{Development.\allowbreak{}AbelianVarietyTheory.\allowbreak{}Proofs.\allowbreak{}DivisorProofs.\allowbreak{}FrameExponentialVolumeSmall}{63--73}
\leanref{GlobalTheory.\allowbreak{}PeakThick.\allowbreak{}peak}{Development.\allowbreak{}AbelianVarietyTheory.\allowbreak{}Proofs.\allowbreak{}DivisorProofs.\allowbreak{}ExistsCuspFormPeakAtInjectiveChartFrame}{27--122}
\end{proof}

\subsection*{Proof of \ref{B8.28a}}\label{prf:B8.28a}\phantomsection
\begin{proof}
Put \(m=25\). In the semiabelian uniformisation, a level label lifts to a point of the extension of the abelian part by the torus whose \(m^2\)-multiple is a period \(y\in X\). The class of \(y\) modulo \(m^2X\) is independent of the lift. This map on the full geometric level is onto: a period admits an \(m^2\)-division point in the extension, since multiplication by \(m^2\) is surjective on the torus and on the abelian part over the algebraic closure. The map commutes with addition and the \(\mathcal O_E\)-action. Thus for the tuple in the statement its period labels \(y_b\) span \(X/m^2X\).
\leanref{GlobalTheory.\allowbreak{}EtaVal.\allowbreak{}exists\_\allowbreak{}qkPt\_\allowbreak{}span}{Development.\allowbreak{}AbelianVarietyTheory.\allowbreak{}Proofs.\allowbreak{}TorusSplitting.\allowbreak{}ExistsHeisenbergLiftsSpanLatticeClasses}{19--73}

We compute the order of one theta constant. The integral frame supplies an integral section of the principal theta line. Its coefficient at \(0\) in the Fourier expansion is integral on the abelian part: it is the weight-zero coefficient of the integral section on the formal torus torsor. Extraction of a character coefficient there is projection to a graded summand over \(V\). At level one all indices differ from \(0\) by periods. Translation by \(\chi\in X\) identifies the coefficient at \(\chi\) with a translate of the coefficient at \(0\) times the cubical period of order \(Q(\chi)\). Translation on the abelian part preserves integral sections; its points over the valued field extend to the valuation ring by properness. Thus the coefficient at \(\chi\), in an integral frame, has order at least \(Q(\chi)\).
\leanref{GlobalTheory.\allowbreak{}integralFamily\_\allowbreak{}levelOne}{Development.\allowbreak{}AbelianVarietyTheory.\allowbreak{}Proofs.\allowbreak{}RingsOfIntegers.\allowbreak{}IntegralFamilyLevelOne}{16--45}
\leanref{GlobalTheory.\allowbreak{}EtaVal.\allowbreak{}exists\_\allowbreak{}expansion\_\allowbreak{}nonneg}{Development.\allowbreak{}AbelianVarietyTheory.\allowbreak{}Proofs.\allowbreak{}TorusSplitting.\allowbreak{}ExistsExpansionNonneg}{30--58}

In evaluating the constant with label \(y_b\), the translate is the \(m\)-multiple of the chosen division point. On character \(\chi\) its toric contribution has order \(B(y_b,\chi)/m\): taking its \(m\)-th power gives the period pairing with \(y_b\). The symmetric torsion frame contributes \(Q(y_b)/m^2\). Indeed its order \(t\) satisfies \(2m^2t=B(y_b,y_b)\) by the symmetric cubical trivialisation, which identifies the square of the frame with the diagonal biextension frame. The rigidified multiplication isomorphism transports this frame to the evaluation at the \(m\)-multiple. The order of the term of index \(\chi\) is therefore at least
\[
          Q(\chi)+\frac{B(y_b,\chi)}m+\frac{Q(y_b)}{m^2}
                    =\frac{Q(m\chi+y_b)}{m^2}.
\]
The theta constant is the convergent sum of these terms. The valuation inequality for a sum gives the same lower bound with the minimum over \(\chi\). This computation uses multiplication on the characteristic-zero generic fibre and takes valuations of powers; it applies also in residue characteristic five.
\leanref{GlobalTheory.\allowbreak{}UT.\allowbreak{}TxVal.\allowbreak{}main}{Development.\allowbreak{}AbelianVarietyTheory.\allowbreak{}Proofs.\allowbreak{}UniformizationProofs.\allowbreak{}ValuationsSymmetricFrameAndTerms}{175--270}
\leanref{GlobalTheory.\allowbreak{}EtaVal.\allowbreak{}le\_\allowbreak{}val\_\allowbreak{}thetaConst}{Development.\allowbreak{}AbelianVarietyTheory.\allowbreak{}Proofs.\allowbreak{}TorusSplitting.\allowbreak{}LeValuationOddLevelThetaConstantIntegralExpansio}{25--37}

For each \(b\), choose \(z_b\in y_b+mX\) minimising \(Q\). A minimum exists since the values are nonnegative integers. The quotient of \(X\) by the \(\mathcal O_E\)-span of the \(z_b\) is a finitely generated abelian group on which multiplication by \(m\) is surjective. It has free rank zero. Hence the \(z_b\) span \(X\otimes\Q\) over \(E\).

Put \(w=(1+\sqrt{-d})/2\in\mathcal O_E\), so \(w+\bar w=1\). The functional
\[
                   \ell(C)=\sum_b C(wz_b,z_b)
                          =\tfrac12\sum_b C(z_b,z_b)
\]
is integral on integral forms by its first expression. For a positive semidefinite \(C\), equality \(\ell(C)=0\) places every \(z_b\) in the radical. That radical is \(\mathcal O_E\)-stable by the balance identity, so it is the whole rational span. Thus \(\ell\) satisfies the positivity condition defining \(\phi\), including when \(X=0\). It follows that
\[
             \phi\le\ell(B)=\sum_b Q(z_b)
                         \le m^2\sum_b v(\theta_b)=625\,v(f_0),
\]
where \(\theta_b\) are the constants in the product. If one factor vanishes the assertion follows directly.
\leanref{GlobalTheory.\allowbreak{}EtaVal.\allowbreak{}sum\_\allowbreak{}cosetMin\_\allowbreak{}ge\_\allowbreak{}phiC\_\allowbreak{}div\_\allowbreak{}sq}{Development.\allowbreak{}AbelianVarietyTheory.\allowbreak{}Proofs.\allowbreak{}TorusSplitting.\allowbreak{}SumCosetMinGePhiCDivSq}{15--51}
\end{proof}

\subsection*{Proof of \ref{B8.28b}}\label{prf:B8.28b}\phantomsection
\begin{proof}
We use three elementary prime sums. For \(n\ge2\),
\[
 \sum_{\ell\le n}\log\ell\ll n,\qquad
 \sum_{\ell\le n}\frac{\log\ell}{\ell-1}\ll\log(2n),\qquad
 \sum_{\ell\le n}\frac1{\ell-1}\ll1+\log\log(3n),
\]
where the sums run over primes and the constants are absolute. For the first, the primes in \((j,2j]\) divide \(\binom{2j}{j}\le4^j\); sum over dyadic intervals up to the next power of two. For the second, use \(\lfloor n/\ell\rfloor\ge n/(2\ell)\) in the prime factorisation of \(n!\). For the third take \(s=1+1/\log(3n)\). Then \(1/\ell\le e/\ell^s\) for \(\ell\le n\), and
\(\sum_\ell \ell^{-s}\le\log\zeta(s)\le\log(1+1/(s-1))\) by the Euler product and the integral test.

The size inequalities \eqref{eq:p8c3-spacing} and \eqref{eq:p8c3-berg-sizes} give \(T\le b\), \(Tr^3\le b\), and \(kr\le b\). The width bound in \ref{B8.25b} gives \(kS\le P_*(S+\eta)\le800b\). Consequently
\[
 TU\le9b,\qquad T(A+V)\ll fmb,\qquad TJ\le17b,\qquad
 \log(2J)\ll L .
\]
Also \(k\) grows at most polynomially in \(x\) by \ref{B8.01k}, so \(1+\log\log(3k)\ll\log\log x\) eventually. In the estimates below a constant may depend on the family.

The contribution at \(\ell\mid Q_S\) is at most
\[
 (A+V)\log Q_S+\tfrac12 A\log\operatorname{rd}(E),
\]
since \(\sum_\ell\delta_\ell\log\ell=\log\operatorname{rd}(E)\). Both logarithms are \(O(L)\) by the host size conditions. At \(q,d\) use \(q,d>4tr+1\) from the host construction and \(\log(qd)\le L\) from \((Q_Sqd)^{36}t^{150}\le\sqrt b\). This part costs \(O((A+V)L)\).

The ordinary fractions with denominator \(\ell-1\) cost \(O((A+V)\log(2J))\) by the second prime sum. The extra numerator \(Ah_\ell/2\) occurs only for \(\ell<k/m\). Its contribution is \(O(A\log(2J)(1+\log\log(3k)))\) by the third prime sum. In that same range \((h_\ell+1)\log\ell\le2\log(2J)\). There are at most \(k/m\) such primes, so the terms with coefficient \(f(m+1)r^2\) cost \(O(fr^3\log(2J))\). The unit terms cost \(O(J)\) by the first prime sum. Multiplication by \(T\) bounds their total by a constant times \(fmbL\log\log x=E_p\).

Finally \(\log(g_0g_1)=O(L)\) by \ref{B8.01k}. The sum of \(\log\ell\) on \(\mathcal B\) is at most \(\log(2dD_EM')=O(fL)\), by \eqref{eq:p8c3-short-conditions} and \(M'=M\) on a good cover. The different part costs at most \((N/2)fr\log\operatorname{rd}(E)\). The column count in \ref{B8.25g} gives \(N\le20fr^2\). Thus these last costs are \(O((frL)^3)\) before multiplication by \(T\), and \(T(frL)^3\le TP_*\le b\) by \eqref{eq:p8c3-berg-sizes}. Enlarging the family constant gives the assertion.
\end{proof}

\subsection*{Proof of \ref{B8.28c}}\label{prf:B8.28c}\phantomsection
\begin{proof}
The character lattice of the toric part has the same rational rank as the period lattice. Its Lie algebra in the Raynaud extension injects into the Lie algebra of the uniformised variety. Each embedding of \(E\) occurs in the toric Lie algebra with multiplicity \(s/n_E\), since the rational cocharacter space is an \(E\)-vector space. An embedding of Lie multiplicity \(Q\) gives the rank bound. For \(s=0\) take the empty basis.

Give \(X_{\R}\) the Euclidean form \(B_X\). Fix an integral functional \(l\) strictly positive on the nonzero positive semidefinite balanced forms, as in Lemma~\ref{B8.28a}. It is a sum of evaluations \(C\mapsto\sum_j C(x_j,y_j)\) with \(x_j,y_j\in X\). Indeed, the integral balanced symmetric forms constitute a saturated submodule of the integral matrices in any basis: symmetry and balance are integral linear equations. An integer functional on this submodule extends to the full matrix lattice.

Let \(u\ne0\) belong to the dual lattice. For \(z\in X\), define \(t_z\in E\) by \(\operatorname{Tr}(a t_z)=u(az)\) for \(a\in E\), extending \(u\) rationally. Thus \(t_z\) lies in the inverse different and \(t_{az}=a t_z\). The form
\[
 C_u(x,y)=\sum_{\sigma:E\to\C}
                  \Re\bigl(\sigma(t_x)\overline{\sigma(t_y)}\bigr)
\]
is positive semidefinite and balanced. Set
\[
 t=\sum_j(t_{x_j}\bar t_{y_j}+\bar t_{x_j}t_{y_j}).
\]
It lies in \(\mathfrak D_{E/\Q}^{-2}\). Each \(\sigma(t)>0\), by applying \(l\) to the nonzero positive semidefinite form
\(\Re(\sigma(t_x)\overline{\sigma(t_y)})\). Hence
\[
 2l(C_u)=\operatorname{Tr}(t)
       \ge n_E |N_{E/\Q}(t)|^{1/n_E}
       \ge n_E/\operatorname{rd}(E)^2 .
\]
For the last inequality use the inclusion \((t)\subset\mathfrak D_{E/\Q}^{-2}\).

The real summands belonging to the different complex places are mutually orthogonal for \(B_X\), by balance. On such a summand \(u=2\Re(\sigma(t_{\,\cdot\,}))\), and multiplication by \(i\) is an isometry. The real and imaginary parts are therefore orthogonal covectors of the same norm. It follows, on summing over the complex places, that \(C_u\le \|u\|^2 B_X/2\), where the norm on covectors is the dual norm. Applying \(l\) proves
\(\|u\|^2\ge n_E/(\operatorname{rd}(E)^2l(B_X))\). The Euclidean transference inequality \(\lambda_s(X)\lambda_1(X^\vee)\le s\) applies to this full lattice and its Euclidean dual \cite[Th.~2.1]{banaszczyk-1993}. Take the infimum over \(l\) to obtain
\(\lambda_s(X)^2\le 2\kappa s^2\phi\).

To pass from independent short vectors to a basis, take successively a shortest \(u_i\in X\) outside the span of its predecessors. Let \(W_i=\operatorname{span}_{\R}(u_1,\ldots,u_i)\). The projection of \(X\cap W_i\) perpendicular to \(W_{i-1}\) is cyclic. Choose a shortest lift \(b_i\) of a generator. Its perpendicular component has length at most \(\|u_i\|\). In the preceding Gram--Schmidt directions one can subtract earlier basis vectors, in reverse order, to make each coordinate at most half the length of the corresponding perpendicular component. It follows that
\(\|b_i\|^2\le(1+(i-1)/4)\|u_i\|^2\le s\|u_i\|^2\).
These lifts form a basis by induction along the saturated flag \(X\cap W_i\). If \(z\in X\cap W_{i-1}\), comparison with \(b_i\pm z\) gives \(2|B_X(b_i,z)|\le\|z\|^2\). Otherwise \(\|u_i\|\le\|z\|\) and Cauchy--Schwarz gives \(|B_X(b_i,z)|\le\sqrt s\,\|z\|^2\). Use \(2Q_X(z)=\|z\|^2\) in both inequalities.
\end{proof}

\subsection*{Proof of \ref{B8.28d}}\label{prf:B8.28d}\phantomsection
\begin{proof}
Use the basis of Lemma~\ref{B8.28c} in each copy. Both bounds for a basis vector persist on the product, since \(Q_X\) is nonnegative on every factor. Take \(\chi=\sum_i n_i b_i\), \(n_i\ge0\), \(\sum_i n_i\le D\). The triangle inequality gives \(Q_X(\chi)\le D^2\Phi\). If \(z\ne0\), then
\[
 M Q_X(z)+B_X(\chi,z)
 \ge (M-2D\sqrt{s})Q_X(z)>0 .
\]
Thus \(\chi\) is the unique minimum in its class modulo \(MX\).

Specify the coefficient at \(\chi\) by an integral section on the abelian part. Its line is the \(M\)-th power of the principal polarising line \(L\) tensored with the degree-zero line of the character \(\chi\). Its sections give arbitrary \(D\)-jets at the origin. Here are the details on a geometric fibre. The cube \(L^3\) is very ample and \(L^j\) is generated for \(j\ge2\) \cite[\S17]{mumford-1970}. Choose a section of \(L^3\) nonzero at the point and sections whose ratios to it give local parameters. Degree-\(D\) products of these sections, including copies of the nonvanishing one, span the truncated local ring after trivialisation. The residual power \(M-3D\ge4\), with its degree-zero twist, has a section nonzero at the point: the twisted power is a translate, by surjectivity of its polarisation isogeny. Multiplication by its jet is invertible. Sections lift from the special fibre, since a positive power with a degree-zero twist has vanishing higher cohomology on an abelian variety \cite[\S16]{mumford-1970}. Base change and Nakayama's lemma then give the integral jet assertion.

Extend the coefficient at \(\chi\) by the period translation law and put zero on the other classes. The increment of valuation at \(\chi+Mz\) is
\[
       M Q_X(z)+B_X(\chi,z)
           =\{Q_X(\chi+Mz)-Q_X(\chi)\}/M .
\]
Translation on the abelian part preserves its integral lattice. These coefficients give a theta section by the Fourier expansion of the semiabelian uniformisation: they satisfy the descent law, and their orders tend quadratically to infinity on the class. On the special fibre only the term at \(\chi\) survives. In a local trivialisation of the split torus torsor its germ is the character of \(\chi\) times the prescribed abelian germ, up to the common line frame and a unit from the base, which can be absorbed in that germ \cite[Ch.~III, \S6]{faltings-chai-1990}.

Write \(z_i\) for the torus characters of the basis. Together with parameters on the abelian part, \(z_i-1\) are parameters at the identity. Expanding \(\prod_i(z_i-1)^{n_i}\) gives characters with nonnegative exponent sum at most \(\sum_i n_i\). The special-fibre germs just constructed therefore span the whole ring of \(D\)-jets, including the mixed toric and abelian monomials. This ring is finite free over the trait, so the integral jets before reduction span by Nakayama's lemma. Finally multiply the entire linear combination by one scalar of valuation \(\lfloor D^2\Phi\rfloor/M\), adjoining a root of the trait uniformiser if necessary. For each participating class this is at least \(Q_X(\chi)/M\), since \(Q_X\) is integer-valued. The increment formula puts the resulting section in the asserted lattice.
\end{proof}

\subsection*{Proof of \ref{B8.28e}}\label{prf:B8.28e}\phantomsection
\begin{proof}
Use the differential of Lemma~\ref{B8.25g}. For each pair \((\tau,i)\) it has a separate block of \(i\) columns with factor \(\det(\partial_{aj})^{c_{\tau i}}\), where \(c_{\tau i}=a_{\tau,i}-a_{\tau,i+1}\). Its total degree is \(D=\sum_{\tau,i}i c_{\tau i}\). There are finitely many level tuples. The associated rational theta derivatives are cusp sections by Lemma~\ref{B8.25t}. Pull them back and multiply on each component by the fixed weight-changing sections of Lemma~\ref{B8.25g}. Include in \(\Sigma\) the primes of the level and discriminant, the primes dividing any \(c_{\tau i}!\), and the primes needed to spread out this finite list of sections and the weight-changing isomorphisms.

For clarity, the last operation concerns fixed algebraic sections. On the compactified generic flag model the cusp sections extend across the boundary. The determinant and flat weight changes and their inverses extend as isomorphisms of their respective lines there: in the determinant case a power identifies the two trivialisations, and in the flat case this follows from the canonical extension in Lemma~\ref{B8.24d}. After inverting one integer all these sections and identities extend on a finite affine cover, by clearing the denominators in their coefficients and transition laws. At the remaining good places the saturated eigenspaces give the ordinary flag model, since the embedding projectors are integral. Thus the weight changes have unit coefficients in integral frames there. One may first multiply all the rational cusp sections by a common denominator to put them in the global cusp lattice. At the finitely many excluded places this clearing can also be seen by taking the closure of the generic flags in the product of the Grassmannians of the integral Lie bundle. Every saturated trait flag gives a point of this proper model. A section regular on the whole generic fibre of an invertible sheaf extends after multiplication by a power of the base uniformiser, by the same finite affine cover argument. The determinant of the change from the projected Hodge lattice to the integral eigenframes has bounded order by Lemma~\ref{B8.24c}. Include the prime divisors of the common denominator in \(\Sigma\).

Fix \(\ell\notin\Sigma\). Forget the extra level, retaining \(M=g_0g_1\). The same semiabelian degeneration and frame give the base evaluation and the same period form. After a finite extension of ramification index \(b_1\) take a split chart. Orders, \(v(\ell)\), and \(\phi\) all multiply by \(b_1\); for \(\phi\) this follows from its definition as an infimum of linear functionals of the period form.

In the \(j\)-th column of the block \((\tau,i)\) choose an integral linear coordinate whose differential on the first \(i\) framed vectors is the \(j\)-th coordinate functional. Such coordinates exist because the frame is a basis of a saturated summand. Take the product, over the columns in the block, of the corresponding coordinates to the power \(c_{\tau i}\), and then the product over blocks. This jet is homogeneous of degree \(D\). At the origin only the principal symbol of the theta differential acts on it: every term of order \(<D\) has value zero. In each determinant power the diagonal monomial has coefficient one. The value is consequently
\[
                             \prod_{\tau,i}(c_{\tau i}!)^i ,
\]
a unit. This test uses the individual minor exponents even when several of them are nonzero.

To apply Lemma~\ref{B8.28d}, put \(s_*=2fm\). The inequalities of \eqref{eq:p8c3-berg-sizes} and Lemma~\ref{B8.25g} give, beyond a common threshold,
\[
 s\le s_*,\quad s_*^3\le b,\quad \kappa\le b,\quad
 D\le b^2,\quad g_0,g_1\ge b^8 .
\]
Indeed \(s_*\le2fr\), \((frL_x)^{24}\le T\le b\), and \(\operatorname{rd}(E)^7\le b\). These bounds give \(4D+4\le M\), \(2D\sqrt{s}<M\), and \(2500D^2\kappa s^3/M\le1\). Lemma~\ref{B8.28d} provides a section in the Fourier lattice whose differentiated value has order at most \(b_1\phi/2500\).

At this prime \(M\) is a unit. We explain why a tuple has differentiated value of at most the same order. Over a splitting field the two complementary halves of the level act through the theta group. Translation preserves the Fourier lattice: translating the index and changing its cubical frame changes the coefficient order and \(Q_X(z)/M\) by the same amount. The character projector of one half is the sum of its translations weighted by inverse character values, divided by the order of the half. This denominator is a unit. The projected spaces are lines: the invariant line descends through the half-level isogeny to the principal theta line, of section dimension one by Lemma~\ref{B8.25l}, and the commutator pairing translates it to each character. The pullback of an integral primitive principal theta section is primitive for the Fourier lattice. Pullback multiplies the period form by \(M\) on its index image, and a primitive coefficient on the abelian part stays primitive under the surjective isogeny on that part. Its translated copies are therefore generators of the integral character lines. This argument applies on the product, with its product principal polarisation. Hence the section just obtained is an integral linear combination of the translated generators. Differentiating the combination and using the valuation inequality selects a tuple. In passing from its pulled-back differential to the theta rule, each framed derivative contributes its embedding of a half-level factor; these factors divide \(M\) and are units. The cubical translations use finite-order normalisations, also units in this comparison.

The weight change on the selected component and the common denominator are units at \(\ell\). We have thus found a global integral section with order at most \(b_1\phi/2500\) upstairs. Dividing by \(b_1v(\ell)\) proves the claim on the original trait.
\end{proof}

\subsection*{Proof of \ref{B8.28f}}\label{prf:B8.28f}\phantomsection
\begin{proof}
The balance identity gives \(Q_X(\alpha z)=gQ_X(z)\). Express \(\alpha^{-1}y\) in the basis \(b_i\). Round each coordinate independently to one of its two neighbouring integers with probabilities giving the original coordinate as mean. The errors have mean zero and variance at most \(1/4\). The expected quadratic value of the error after multiplication by \(\alpha\) is at most \(g\sum_i Q_X(b_i)/4\), since all cross terms have mean zero. Some representative therefore satisfies the first bound. Choose \(z_0\) to minimise \(Q_X\) in the class; the values are nonnegative integers. Comparing with \(z_0\pm\alpha z\) gives
\(\lvert B_X(z_0,\alpha z)\rvert\le Q_X(\alpha z)\). Expansion of \(Q_X(z_0+g'\alpha z)\) proves the second inequality. It is strictly positive on the right for \(z\ne0\). The ultrametric inequality for the convergent expansion now gives the final assertion.
\end{proof}

\subsection*{Proof of \ref{B8.28h}}\label{prf:B8.28h}\phantomsection
\begin{proof}
Write \(\widetilde F(t)\) for the composite series in the definition. Write \(\det(\partial/\partial t_{ab})^k=\sum_\alpha p_\alpha\partial^\alpha\). The coefficients \(p_\alpha\) are integers, and every \(\alpha\) has total degree \(ik\). Then
\[
  \det(\partial/\partial t_{ab})^k\,\widetilde F(0)=\sum_\alpha p_\alpha\,\alpha!\,[t^\alpha]\widetilde F,\qquad \alpha!=\textstyle\prod_{a,b}\alpha_{ab}!.
\]

Let \(G\) be the formal group law in the chosen coordinates. For an integral Lie vector \(u\) the invariant derivation is
\[
 D_uF(X)=\sum_j\Bigl(\sum_r u_r\,\frac{\partial G_j}{\partial Y_r}(X,0)\Bigr)\frac{\partial F}{\partial X_j}.
\]
It maps integral series to integral series. The map \(s\mapsto\exp(su)\) is a homomorphism from the additive group, so \(\exp(x+su)=\exp(x)+_G\exp(su)\). Hence \(\partial/\partial t_{ab}\) acting on \(\widetilde F\) corresponds to \(D_{u_{\tau,a}}\) acting on the \(b\)-th copy. These derivations commute. Consequently \(\alpha!\,[t^\alpha]\widetilde F=\partial^\alpha\widetilde F(0)\) is the constant term of an integral series. This proves \(\nu_{\tau,i}(k;y)\ge0\).

For the upper bound, fix a field and integral coordinates. The vectors \(u_{\tau,1},\dots,u_{\tau,i}\) span a saturated summand. Hence their coordinate matrix has an \(i\times i\) minor \(m\) that is a unit. Let \(x_{r_1},\dots,x_{r_i}\) be the corresponding coordinates and put \(z_a=\sum_{s}(m^{-1})_{as}x_{r_s}\). These functions are integral and satisfy
\[
 z_a\bigl(\exp(\textstyle\sum_{a'}t_{a'}u_{\tau,a'})\bigr)=t_a+O(t^2).
\]
Take \(F=\prod_{a,b}z_a(\text{copy }b)^{\alpha_{ab}}\). Then \(\widetilde F=t^\alpha\) plus terms of degree \(>ik\). The operator is homogeneous of degree \(ik\), so the value is \(p_\alpha\,\alpha!\).

Take the diagonal exponent \(\alpha_{ab}=k\delta_{ab}\). A term \(\prod_{j=1}^k\prod_a\partial_{a,\pi_j(a)}\) of \(\det^k\), with permutations \(\pi_1,\dots,\pi_k\), is the diagonal monomial only when every \(\pi_j\) is the identity. Hence \(p_\alpha=1\), and \(\alpha!=(k!)^i\). Therefore \(\nu_{\tau,i}(k;y)\le i\,v_\ell(k!)\). By Legendre's formula, \(v_\ell(k!)=\sum_{j\ge1}\lfloor k/\ell^j\rfloor\le k/(\ell-1)\), and \(v_\ell(k!)=0\) for \(k<\ell\).
\end{proof}

\subsection*{Proof of \ref{B8.28i}}\label{prf:B8.28i}\phantomsection
\begin{proof}
By Lemma~\ref{B8.28h}, \(\nu_c(y)\ge0\). By hypothesis, \(\nu_c(y)\ge\sum_fA_f\). Hence \(\nu_c(y)\ge\max(0,\sum_fA_f)\ge A_\lambda\).

At an algebraic frame,
\[
 \nu_c(y)\le\sum_fB_f=\sum_fA_f+\sum_f(B_f-A_f)\le\max\Bigl(0,\sum_fA_f\Bigr)+\sum_f(B_f-A_f),
\]
and \(\max(0,\sum_fA_f)<A_\lambda+1\).

For the last two assertions take \(A_\lambda=0\). Lemma~\ref{B8.28h} gives \(0\le\nu_f(c_f;y)\le i\,c_f/(\ell-1)\) for \(f=(\tau,i)\), with equality to zero when \(c_f<\ell\). Summing over the blocks gives \(0\le\nu_c(y)\le D/(\ell-1)\). It also gives \(\nu_c(y)=0\) when \(\ell>\max_fc_f\).
\end{proof}

\subsection*{Proof of \ref{B8.28j}}\label{prf:B8.28j}\phantomsection
\begin{proof}
Put \(m'(\lambda)=k\) and \(m'(w')=m(w')\) for \(w'\ne\lambda\). Let \(S\) be the support of \(m'\) together with \(\lambda\). For \(w'\in S\) put \(n_{w'}=-\lfloor m'(w')\rfloor\).

We first find \(\sigma\in F_c^\times\) with \(v_{w'}(\sigma)=n_{w'}\) for \(w'\in S\) and \(v_{w'}(\sigma)\ge0\) for \(w'\notin S\). Let \(L\) be the product of the rational primes below \(S\). Choose \(j\ge0\) with \(j\,v_{w'}(L)+n_{w'}\ge0\) for all \(w'\in S\). By the Chinese remainder theorem there is \(\beta\in\mathcal O_{F_c}\) such that
\begin{itemize}
\item \(\beta\equiv\pi_{w'}^{\,jv_{w'}(L)+n_{w'}}\) modulo \(\mathfrak p_{w'}^{\,jv_{w'}(L)+n_{w'}+1}\) for \(w'\in S\), where \(\pi_{w'}\) is a uniformiser, and
\item \(\beta\equiv0\) modulo \(\mathfrak p_{w'}^{\,jv_{w'}(L)}\) for the primes \(w'\notin S\) dividing \(L\).
\end{itemize}
Then \(\beta\ne0\), and \(\sigma=\beta/L^j\) has valuation \(n_{w'}\) at \(w'\in S\). Its valuation is at least \(0\) at the primes outside \(S\) dividing \(L\). At the remaining primes its valuation equals \(v_{w'}(\beta)\ge0\).

Now let \(y\) be an algebraic integral frame over \(w'\) with ramification index \(e\). Then
\[
 v(\sigma f(y))/e=v_{w'}(\sigma)+v(f(y))/e.
\]
For \(w'\in S\) this is at least \(-\lfloor m'(w')\rfloor+m'(w')\ge0\). For \(w'\notin S\) both terms are nonnegative, since \(m(w')=0\) there. Hence \(\sigma f\) is integral, and \(v_\lambda(\sigma)=-k\).
\end{proof}

\subsection*{Proof of \ref{B8.28k}}\label{prf:B8.28k}\phantomsection
\begin{proof}
(a) Let \(\mathfrak w\) be the prime of \(\mathcal O_E\) on which \(\tau\), followed by reduction modulo \(\lambda\), vanishes. The extension \(E/E^+\) is unramified at every finite place, so each prime \(\mathfrak q\mid\ell\) of \(\mathcal O_E\) is either inert over \(E^+\), with \(\bar{\mathfrak q}=\mathfrak q\), or split, with \(\bar{\mathfrak q}\ne\mathfrak q\). For each such prime and each \(n\ge1\) we choose classes \(u,t\) modulo \(\mathfrak q^n\) and modulo \(\bar{\mathfrak q}^n\) such that \(u\equiv1\) and \(u\bar u+t\bar t\equiv0\). The Chinese remainder theorem combines these classes.

Suppose first \(\bar{\mathfrak q}\ne\mathfrak q\). Take \(u\equiv1\), take \(t\equiv-1\) modulo \(\mathfrak q^n\), and take \(t\equiv1\) modulo \(\bar{\mathfrak q}^n\). Then \(\bar t\equiv1\) modulo \(\mathfrak q^n\) and \(\bar t\equiv-1\) modulo \(\bar{\mathfrak q}^n\). Hence \(1+t\bar t\) vanishes modulo both ideals.

Suppose now \(\bar{\mathfrak q}=\mathfrak q\). Conjugation induces the nontrivial involution of the residue field \(k_{\mathfrak q}\), which is quadratic over the residue field of \(E^+\). The norm of finite fields is onto, so \(-1=t_1\bar t_1\) modulo \(\mathfrak q\) for some \(t_1\). Assume \(t_n\bar t_n=-1+\varepsilon\) with \(\varepsilon\in\mathfrak q^n\). Then \(\bar\varepsilon=\varepsilon\). Since \(\mathfrak q\) is unramified over \(E^+\), we may write \(\varepsilon=\pi^n\varepsilon_0\), where \(\pi\) is a uniformiser of \(E^+\) at \(\mathfrak q\cap E^+\) and \(\varepsilon_0\) is invariant under conjugation. The residue trace is onto, so there is \(t_0\) with \(t_0+\bar t_0\equiv\varepsilon_0\) modulo \(\mathfrak q\). Put \(t_{n+1}=t_n(1+\pi^nt_0)\). Then
\[
 t_{n+1}\bar t_{n+1}\equiv(-1+\varepsilon)\bigl(1+\pi^n(t_0+\bar t_0)\bigr)\equiv-1\pmod{\mathfrak q^{n+1}} .
\]

At \(\mathfrak w\) these choices satisfy \(t_{n+1}\equiv t_n\) modulo \(\mathfrak w^n\). Define \(\alpha_\tau\) as the limit of \(\tau(t_n)\) in \(V_\lambda\). In the split case \(\alpha_\tau=-1\). Given \(n_0\), take \(n\ge n_0v_{\mathfrak q}(\ell)\) at every \(\mathfrak q\mid\ell\), and combine. The ideal \(\mathfrak w^n\) maps into \(\ell^{n_0}V_\lambda\), which gives the two congruences in \(V_\lambda\).

(b) By the Cauchy--Binet formula, each determinant in \(P_{\rm col}\) is a \(V_\lambda\)-linear combination of leading minors in subsets of the columns of its block. Hence \(P_{\rm col}\) lies in the span of products of leading minors. The row \((\tau,i')\) occurs exactly in the determinants of the blocks \((\tau,i)\) with \(i\ge i'\), with degree one in each. Its total degree is therefore \(\sum_{i\ge i'}c_{\tau,i}=a_{\tau,i'}\).

For harmonicity, fix a column \(b\). At most one block uses \(b\), and that factor involves the variables of a single embedding. Hence no factor depends on both \(X_{\tau,\cdot,b}\) and \(X_{\bar\tau,\cdot,b}\). By the Leibniz rule, \(\partial^2/\partial X_{\tau,i,b}\partial X_{\bar\tau,i',b}\) annihilates the product.

The coefficients of \(P_{\rm col}\) lie in \(V_\lambda\). For the approximants, apply (a) with \(n_0\). In row \(j\) put \(x_{j,b_j}=u\), \(x_{j,b_{i+j}}=t\), and zero elsewhere. Then \(\sum_bx_{j,b}\bar x_{j,b}=u\bar u+t\bar t\). For \(j\ne j'\) the supports are disjoint and the sum vanishes.

(c) Let \(W_{\Q}\) be the space of rational polynomials of row degrees \(a\) that lie in the span of products of leading minors and are harmonic. Put \(W_{\Z}=W_{\Q}\cap\Z[X]\). The submodule \(W_{\Z}\) is saturated in the finite free module of integral polynomials of these row degrees, hence a direct summand. Both defining conditions of \(W_{\Q}\) are given over \(\Q\), and \(V_\lambda\) is torsion free. Therefore
\[
 \operatorname{harm}_{V_\lambda}(a,N)=(W_{\Q}\otimes\operatorname{Frac}V_\lambda)\cap V_\lambda[X]=W_{\Z}\otimes V_\lambda .
\]
Let \(P_1',\dots,P_K'\) be a \(\Z\)-basis of \(W_{\Z}\subset\operatorname{harm}(a,N)\), and write \(P_{\rm col}=\sum_i\gamma_iP_i'\) with \(\gamma_i\in V_\lambda\). Since \(\mathcal O_{F_c}\) is dense in \(V_\lambda\), there are \(\xi_i\in\mathcal O_{F_c}\) with \(\xi_i\equiv\gamma_i\) modulo \(\ell^{n_0}V_\lambda\). Put \(P_1=\sum_i\xi_iP_i'\). Then
\[
 P_1-P_{\rm col}=\sum_i(\xi_i-\gamma_i)P_i'\in\ell^{n_0}\operatorname{harm}_{V_\lambda}(a,N).
\]
Applying the map \(V_\lambda\to V\) of a trait over \(\lambda\) transports this identity.
\end{proof}

\subsection*{Proof of \ref{B8.28l}}\label{prf:B8.28l}\phantomsection
\begin{proof}
Let \(\lambda'\mid\ell\) be a place of \(F'\), \(V\) a trait over \(\lambda'\), and \(y'\) an algebraic integral frame on \(Y\) over \(V\) with phase \(\phi\). Forget the additional level structure of the cover and restrict the coefficients from \(F'\) to \(F_c\). This gives an algebraic integral frame \(y\) on the base moduli over the same \(V\). It has the same semiabelian scheme, hence the same period form and the same phase. Through \(F_c\subset F'\), the ring \(V\) is a trait over the place \(\lambda\) of \(F_c\) below \(\lambda'\).

Let \(f\) be an integral cusp section of weight \(w\) on the base. Its pullback \(g\) to \(Y\) is a cusp section of weight \(w\). It satisfies \(g(z')=f(z)\) for every frame \(z'\) on \(Y\) with image \(z\). Every algebraic integral frame on \(Y\) has an algebraic integral image on the base, so \(g\) is integral. Hence the infimum over the integral cusp sections on \(Y\) at \(y'\) is at most \(v(f(y))\), for every such \(f\). Both reference depths use the valuation of \(V\) and the same phase, which gives the inequality.
\end{proof}

\subsection*{Proof of \ref{B8.28m}}\label{prf:B8.28m}\phantomsection
\begin{proof}
(1) Pass to a finite extension of \(V\) over which the frame has a split uniformisation and the points of order dividing \(M^2\) are rational. The image frame is again algebraic and integral over \(\lambda\). Values of sections and valuations normalised by \(v(\ell)\) do not change. Since \(\ell\nmid M\), the factor \(M^{-D}\) and the normalising constant of each \(s_{x_b}\) are units at \(\lambda\).

The germ of \(s_{x_b}\) at the origin, in the integral frame of the \(M^2\)-th power, is integral. To see this, expand \(s_{x_b}\) in the characters \(\chi\) of the period lattice. The abelian coefficients are integral. They are multiplied by constants of order at least \(Q_X(M\chi+y_b)/M^2\ge0\), where \(y_b\) is the period label of the translating point. This is the order computation in the proof of Lemma~\ref{B8.28a} with \(m=M\). When the period lattice is zero, it reduces to the statement there that translation on the abelian part preserves integral sections. The characters have integral germs on the formal group of the Raynaud extension, which is the formal group of the frame.

Hence \(s_{x_1}\boxtimes\dots\boxtimes s_{x_N}\) has an integral germ on \(N\) copies. By definition of \(\delta_{M^2}\), the normalised order of \(P(\nabla)\) applied to this product at the origin is at least \(\delta_{M^2}(P;y)\ge A\). Multiplying by the unit \(M^{-D}\) proves (1).

(2) Let \(F\) be an integral power series on \(\widehat G^N\) and put
\[
 x_0=\bigl[P(D)\bigl(F\,e^{M\sum_bh_L(x_b)}\bigr)\bigr](0).
\]
We may assume \(x_0\ne0\). Every monomial of \(P\) has degree \(D\), so \(x_0\) depends only on the jet of order \(D\) of \(F\) at the origin. Apply Lemma~\ref{B8.28d} on \(N\) copies, with the integral frame of the theta line at the identity. After a scalar extension there are a scalar \(\gamma\) with \(v(\gamma)\le D^2\Phi/M\) and a level-\(M\) theta section \(s\) on the product, in the Fourier lattice, whose jet of order \(D\) in that frame is \(\gamma\) times the jet of \(F\). Hence \(P(\nabla)s\) has value \(\gamma x_0\) at the origin.

The level \(M\) is a unit in \(V\). Over a splitting field the two complementary halves of the level act through the theta group. Translation preserves the Fourier lattice, because translating the index and changing its cubical frame change the coefficient order and \(Q_X(z)/M\) by the same amount. The character projector of one half is the average of its translations, weighted by inverse character values. Its denominator is the order of the half, which is a unit.

Each projected space is a line. It is generated by a translate of the pullback of an integral primitive principal theta section. These translates generate the integral character lines, by the argument in the proof of Lemma~\ref{B8.28e}. Hence \(s\) is a linear combination, with integral coefficients, of translated generators.

Take a translated generator with characteristic tuple \(x\). The value of \(P(\nabla)\) on it at the origin is \(\theta_{x,P}(y)\) times an integral scalar. This scalar is a unit times the product of the images of the half-level factor \(\alpha\), \(\alpha\bar\alpha=M\), under the embeddings, raised to the row degrees. The ultrametric inequality therefore gives a tuple \(x\) with
\[
 v(\theta_{x,P}(y))\le v(\gamma x_0)\le v(x_0)+D^2\Phi/M .
\]
Normalised valuations do not change under the scalar extension. Dividing by \(v(\ell)\) and taking the infimum over \(F\) proves (2).
\end{proof}

\subsection*{Proof of \ref{B8.28n}}\label{prf:B8.28n}\phantomsection
\begin{proof}
Fix an algebraic integral frame \(y\) over \(\lambda\) on a trait \(V\), with phase \(\phi\). After a finite extension \(V'\) of ramification index \(r'\), the semiabelian degeneration of the frame has a split uniformisation \cite[Ch.~III]{faltings-chai-1990}. Let \(y'\) be the image frame. As in the proof of Lemma~\ref{B8.28e}, orders of sections, \(v(\ell)\) and the phase multiply by \(r'\). The frame \(y'\) is algebraic and integral over \(\lambda\).

Let \(s\) and \(\Phi=\kappa s^3r'\phi\) be as in Lemma~\ref{B8.28c} on \(V'\). At a nonterminal layer \(q_{\rm sig}\le m\), so \(s\le2fq_{\rm sig}\le2fm\). Hence \(2D\sqrt s<M\), and
\[
 \frac{D^2\Phi}{M\,r'v(\ell)}\le\frac{\kappa(2fm)^3D^2\phi}{M\,v(\ell)}\le\frac{c\,\phi}{v(\ell)} .
\]
The level \(M=g_0g_1\) is prime to \(\ell\). Lemma~\ref{B8.28m}(2) on \(V'\) and the upper bound in (i) at \(y'\) give
\begin{equation}\label{eq:p8c3-ref-away-up}
 \inf_x\frac{v(\theta_{x,P_1}(y))}{v(\ell)}\le A+\varrho+\frac{c\,\phi}{v(\ell)} .
\end{equation}
Here we use that the value at \(y'\) is the image of the value at \(y\), and that valuations normalised by the valuation of \(\ell\) agree.

Fix a tuple \(x\) and put \(\theta=\theta_{x,P_1}\). Every element of \(\operatorname{harm}(a,N)\) is a harmonic linear combination of products of leading minors. Each product contains minors of size \(P\) at an embedding of the CM type with \(a_{\tau,P}>0\). The proof of Lemma~\ref{B8.25t} applies to such a combination. Harmonicity is used there for the change of symplectic basis, and the zero-frequency term vanishes for each product separately. Since \(N\) is even, \(\theta\) is a cusp section of weight \((a,N)\).

By Lemma~\ref{B8.28m}(1) and the lower bound in (i), \(v(\theta(y''))\ge A\,v(\ell)\) at every algebraic integral frame \(y''\) over \(\lambda\). Equivalently, \(v(\theta(y''))/e''\ge A\,v_\lambda(\ell)\) for the ramification index \(e''\). Lemma~\ref{B8.28j}, with \(k=A\,v_\lambda(\ell)\) and the function \(m_x\) from (ii), gives \(\sigma\in F_c^\times\) with \(v_\lambda(\sigma)=-A\,v_\lambda(\ell)\) such that \(\sigma\theta\) is integral. On \(V\) we have \(v(\sigma)=e\,v_\lambda(\sigma)=-A\,v(\ell)\). Hence
\[
 \inf_f\frac{v(f(y))}{v(\ell)}\le\frac{v(\sigma\theta(y))}{v(\ell)}=-A+\frac{v(\theta(y))}{v(\ell)} .
\]
Taking the infimum over \(x\) and using \eqref{eq:p8c3-ref-away-up} proves the inequality. The last assertion follows from Definition~\ref{B8.28g}(e).
\end{proof}

\subsection*{Proof of \ref{B8.28ua}}\label{prf:B8.28ua}\phantomsection
\begin{proof}
The defining congruences are preserved by addition, subtraction and multiplication, and contain \(1\), so \(S\) is a subring. Fix \(u\in S\cap V^\times\) and \(n\ge0\). Choose \(a\in R\) such that \(\varphi(a)-u\in\varphi(q)^{n+1}V\). Since \(\varphi(q)\) lies in the maximal ideal \(\mathfrak m_V\), the element \(\varphi(a)\) is a unit. The ideal \(\mathfrak p=\varphi^{-1}(\mathfrak m_V)\) is a nonzero prime ideal of \(R\), because it contains \(q\) and \(\varphi(q)\ne0\). Hence \(\mathfrak p\) is maximal. There are \(b\in R\) and \(i\in\mathfrak p\) with \(ab+i=1\).

Choose an integer \(s\ge1\) such that \(\varphi(i)^s\in\varphi(q)^nV\), and put \(b_n=b\sum_{j=0}^{s-1}i^j\). The geometric-sum identity gives
\[
 \varphi(a)\varphi(b_n)=1-\varphi(i)^s
 \equiv1\pmod{\varphi(q)^nV}.
\]
Since \(\varphi(a)\equiv u\pmod{\varphi(q)^nV}\), multiplication by \(u^{-1}\in V\) yields
\(\varphi(b_n)\equiv u^{-1}\pmod{\varphi(q)^nV}\).
Thus \(u^{-1}\in S\). The entries of \(\operatorname{adj}(M)\) and \(\det M\) lie in \(S\), and
\[
 M^{-1}=(\det M)^{-1}\operatorname{adj}(M)
\]
proves the matrix assertion.
\end{proof}

\subsection*{Proof of \ref{B8.28ub}}\label{prf:B8.28ub}\phantomsection
\begin{proof}
The constant term of \(U\) is \(1\). Its multiplicative inverse has coefficients in \(\mathcal O_C\), by recursive comparison of coefficients in \(UU^{-1}=1\). Substitution by \(\rho\) preserves integrality, so \(G\) is integral. The two curve identities give
\[
 \begin{aligned}
 G(\mathbf c(t))W(\mathbf c(t))^k
 &=F(\mathbf e(t))U(\mathbf e(t))^{-k}
   W(\psi(\mathbf e(t)))^k\\
 &=F(\mathbf e(t)).
 \end{aligned}
\]
A polynomial differential operator at the origin is a linear functional on a finite jet, which proves the last assertion.
\end{proof}

\subsection*{Proof of \ref{B8.28uc}}\label{prf:B8.28uc}\phantomsection
\begin{proof}
Put \(\mathfrak p=\varphi^{-1}(\mathfrak m_V)\). It is nonzero because it contains \(q\ne0\), and hence it is maximal. Write \((q)=\mathfrak p^rJ\), with \(J+\mathfrak p=R\).

First suppose that \(b\in\mathfrak p^a\), \(\varphi(b)\ne0\), and \(\varphi(c)=\varphi(b)z\). Then \(c\in\mathfrak p^a\). Indeed, for \(c\ne0\), write \((c)=\mathfrak p^iJ'\) with \(J'+\mathfrak p=R\). If \(i<a\), choose \(s\in J'\setminus\mathfrak p\). Since
\[
 (b)J'\subseteq\mathfrak p^aJ'\subseteq(c)\mathfrak p,
\]
there is \(m\in\mathfrak p\) with \(bs=cm\). Cancelling \(\varphi(b)\) gives
\(\varphi(s)=z\varphi(m)\in\mathfrak m_V\), a contradiction.

For each \(N\), let \(k_N=v_V(\varphi(q)^N)\), where \(v_V\) gives a uniformiser value \(1\). The ideals \(\mathfrak p^{k_N}\) and \(J^N\) are comaximal. Choose \(r_N\in\mathfrak p^{k_N}\) and \(\varepsilon_N\in J^N\) with \(r_N+\varepsilon_N=1\). Since
\(\varphi(r_N)\in\mathfrak m_V^{k_N}=\varphi(q)^NV\), the first congruence follows. If \(M\le N\) and \(\varphi(a)\in\varphi(q)^MV\), the preceding argument with \(b=q^M\) gives \(a\in\mathfrak p^{rM}\). Therefore
\[
 \varepsilon_Na\in J^M\mathfrak p^{rM}=(q^M).
\]

For consecutive indices, use
\[
 \varepsilon_{N+1}x_{N+1}-\varepsilon_Nx_N
 =\varepsilon_N(x_{N+1}-x_N)
 +\bigl(\varepsilon_{N+1}(1-\varepsilon_N)
        -\varepsilon_N(1-\varepsilon_{N+1})\bigr)x_{N+1}.
\]
Each term on the right belongs to \(q^NR\), by the divisibility assertion just proved. Summing these differences proves the assertion for \(m\le n\). Multiplying the congruences for \(\varepsilon_N\) and \(x_N\) proves the final assertion.
\end{proof}

\subsection*{Proof of \ref{B8.28ud}}\label{prf:B8.28ud}\phantomsection
\begin{proof}
The embedding of \(\mathcal O_E\) into \(\operatorname{Frac}V\) has image in \(V\), by integral closedness. Apply Lemma~\ref{B8.28uc} with \(R=\mathcal O_E\) and \(q=\ell\). Since \(a\mapsto d[a]\) is an additive map to integral matrices and \(d[\ell^Na]=\ell^Nd[a]\), the matrices \(d[\varepsilon_{\tau,N}x_N]\) form an entrywise Cauchy sequence. Completeness gives an integral limit \(A\) with the stated error bound. On a \(\tau\)-eigenvector these matrices act by \(\tau(\varepsilon_{\tau,N}x_N)\). These scalars converge to \(t\), proving the eigenvector identity.

For the second assertion, the estimate at \(N=0\) makes every entry of \(A_b\) integral. Logarithms turn the formal group law into addition and the CM action into its differential:
\[
 \exp_{\widehat G}(w+w')
   =\exp_{\widehat G}(w)+_{\widehat G}\exp_{\widehat G}(w'),
 \qquad
 \exp_{\widehat G}(d[a]w)=[a]\exp_{\widehat G}(w).
\]
Thus \(\exp_{\widehat G}(\sum_bd[a_{N,b}]w_b)\) is integral for every \(N\).
\leanref{GlobalTheory.\allowbreak{}IsoTriv.\allowbreak{}log\_\allowbreak{}sum\_\allowbreak{}additive\_\allowbreak{}preserves\_\allowbreak{}integrality}{Development.\allowbreak{}AbelianVarietyTheory.\allowbreak{}Proofs.\allowbreak{}DivisorProofs.\allowbreak{}LogarithmSumIntegralAdditive}{34--82}
\leanref{GlobalTheory.\allowbreak{}IsoTriv.\allowbreak{}actSt}{Development.\allowbreak{}AbelianVarietyTheory.\allowbreak{}Proofs.\allowbreak{}DivisorProofs.\allowbreak{}ActionLinearLogCoordinatesIntegral}{15--67}

Fix an output coefficient. Because the \(w_b\) have zero constant term, this coefficient of \(\exp_{\widehat G}(\sum_bM_bw_b)\) is a polynomial \(Q\) in the entries of the finitely many matrices \(M_b\). Choose \(K\ge0\) such that every coefficient of \(Q\) has valuation at least \(-K\). On integral matrix arrays, the difference of two monomials has valuation at least the smallest valuation of an entrywise difference. Therefore the difference of the chosen coefficients for \(M_b=A_b\) and \(M_b=d[a_{N,b}]\) has valuation at least \(N-K\). Taking \(N\ge K\) makes it integral. The approximating coefficient is integral, so the limiting coefficient is integral. This proves the assertion coefficient by coefficient.
\leanref{GlobalTheory.\allowbreak{}IsoTriv.\allowbreak{}contSt\_\allowbreak{}holds}{Development.\allowbreak{}AbelianVarietyTheory.\allowbreak{}Proofs.\allowbreak{}DivisorProofs.\allowbreak{}SeriesCoeffUniformlyContinuousInCoeffsProof}{98--110}
\end{proof}

\subsection*{Proof of \ref{B8.28ue}}\label{prf:B8.28ue}\phantomsection
\begin{proof}
Let \(\mathcal U\) be the additive subgroup of zero-constant series \(s\in C[[t]]\) for which \(\exp(s)\) is an integral unit. It is closed under addition and negation, because \(\exp(s+s')=\exp(s)\exp(s')\) and an integral series with constant term \(1\) has integral inverse. Put
\[
 H(w,z)=h\bigl(\exp_{\widehat G}w,\exp_{\widehat G}z\bigr).
\]
For integral formal points, the biextension cocycle identities and adjointness give, modulo \(\mathcal U\),
\[
 H(w+w',z)=H(w,z)+H(w',z),\qquad
 H(z,w+w')=H(z,w)+H(z,w'),
\]
and
\[
 H(d[a]w,z)=H(w,d[\bar a]z).
\]
Each difference is the logarithm of the corresponding integral biextension cocycle evaluated on integral formal points; its exponential has constant term \(1\).
\leanref{GlobalTheory.\allowbreak{}IsoTriv.\allowbreak{}GaugeBiadd.\allowbreak{}biadd\_\allowbreak{}of\_\allowbreak{}add\_\allowbreak{}of\_\allowbreak{}act}{Development.\allowbreak{}AbelianVarietyTheory.\allowbreak{}Proofs.\allowbreak{}DivisorProofs.\allowbreak{}CanonGaugeBiadditiveOfLogSumActionLinear}{94--151}

Choose generators \(e_1,\ldots,e_r\) of the finite \(\Z\)-module \(\mathcal O_E\). Isotropy permits expressions
\[
 \left(\sum_j\overline{a_{N,b,j}}a_{N,b',j}\right)w_0
   =\ell^N\sum_{i=1}^r n_{N,b,b',i}e_i,
 \qquad n_{N,b,b',i}\in\Z.
\]
Write \(y_{N,j}=\sum_bd[a_{N,b,j}]w_b\) and
\[
 T_N=\ell^N\sum_{b,b',i}n_{N,b,b',i}
        H(w_b,d[e_i]w_{b'}).
\]
Expanding the preceding identities gives
\[
 \sum_jH(y_{N,j},d[w_0]y_{N,j})-T_N\in\mathcal U.
\]
Hence the frame evaluated at \((\exp_{\widehat G}y_{N,j})_j\) is
\(U_N\exp(T_N)\), where \(U_N\) is an integral unit.
\leanref{GlobalTheory.\allowbreak{}IsoTriv.\allowbreak{}HermRead.\allowbreak{}finSt\_\allowbreak{}of}{Development.\allowbreak{}AbelianVarietyTheory.\allowbreak{}Proofs.\allowbreak{}DivisorProofs.\allowbreak{}GaugeHermitianCombinationOfSteps}{52--111}
\leanref{GlobalTheory.\allowbreak{}IsoTriv.\allowbreak{}frameSt\_\allowbreak{}of}{Development.\allowbreak{}AbelianVarietyTheory.\allowbreak{}Proofs.\allowbreak{}DivisorProofs.\allowbreak{}FrameSeriesExpSumOfActionLinearLogUnit}{56--100}

The approximation at \(N=0\) makes every entry of each \(A_{b,j}\) integral. Fix a coefficient indexed by \(\alpha\). For every \(\gamma\le\alpha\), the coefficient of \(t^\gamma\) in \(\exp(T_N)\) is a polynomial in the finitely many scalars \(\ell^Nn_{N,b,b',i}\). These scalars have valuation at least \(N\), so the finite-jet estimate in Lemma~\ref{B8.28ud}, compared with all scalars equal to zero, makes these coefficients integral for sufficiently large \(N\). The coefficient of \(t^\alpha\) in \(U_N\exp(T_N)\) is then integral by the finite product formula. The same finite-jet estimate, applied to the frame after the formal exponential and to the matrices \(A_{b,j}\) and \(d[a_{N,b,j}]\), makes the difference of their \(\alpha\)-coefficients integral for sufficiently large \(N\). Thus the limiting frame's \(\alpha\)-coefficient is integral. Its constant term is \(1\) because \(h_L(0)=0\).
\end{proof}

\subsection*{Proof of \ref{B8.28u}}\label{prf:B8.28u}\phantomsection
\begin{proof}[Proof (outline)]
The algebraically closed complete valued extension \((C,v_\ell)\) of \(\operatorname{Frac}V\) exists over any trait above \(\lambda\mid\ell\). The integral formal coordinates of \(\widehat G\) over \(V\) descend the formal biextension data (group law, \(\mathcal O_E\)-action, cocycles, symmetry, and adjointness) to \(V\), with canonical gauge \(h\) on \(\widehat G\times\widehat G\) satisfying \(h(0,0)=0\) and \(h_L(x)=h(x,[w_0]x)\) having zero constant term. The differential \(d[a]\) of the action of \(a\in\mathcal O_E\) acts on each Lie vector \(u_{\tau,i'}\) by multiplication by \(\tau(a)\).
\leanref{GlobalTheory.\allowbreak{}FactorDepth.\allowbreak{}nonempty\_\allowbreak{}valExt\_\allowbreak{}of\_\allowbreak{}overLa}{Development.\allowbreak{}AbelianVarietyTheory.\allowbreak{}Definitions.\allowbreak{}Divisors.\allowbreak{}NonemptyValExtOfOverLa}{320--322}
\leanref{GlobalTheory.\allowbreak{}integralChart\_\allowbreak{}formalBiext}{Development.\allowbreak{}AbelianVarietyTheory.\allowbreak{}Proofs.\allowbreak{}JetHarmonics.\allowbreak{}IntegralChartFormalBiext}{338--348}
\leanref{GlobalTheory.\allowbreak{}PowGerm.\allowbreak{}eigDirSt}{Development.\allowbreak{}AbelianVarietyTheory.\allowbreak{}Proofs.\allowbreak{}JetHarmonics.\allowbreak{}IntegralChartDirectionsAreEigenvectors}{619--628}

For (i), let \(k\ge0\) and let \(G\) be an integral power series over \(\mathcal O_C\) on \(\widehat G^N\). Given any matrix \(x=(x_{\alpha,b})\) over \(\mathcal O_E\), the biextension cocycle laws produce a formal power series \(U\) with coefficients in \(V\) and constant term \(1\) on \(\widehat G^{\mathcal C}\) satisfying
\[
 U(y)=\exp\Biggl(\sum_{b=1}^N h_L\Bigl(\sum_{\alpha\in\mathcal C}{}_{\widehat G}\,[x_{\alpha,b}]y_\alpha\Bigr)-\sum_{\alpha,\alpha'\in\mathcal C}h\Bigl(y_\alpha,\Bigl[w_0\sum_{b=1}^N\bar x_{\alpha,b}x_{\alpha',b}\Bigr]y_{\alpha'}\Bigr)\Biggr),
\]
where \(\sum_{\widehat G}\) denotes formal-group addition.
\leanref{GlobalTheory.\allowbreak{}exists\_\allowbreak{}integral\_\allowbreak{}unit\_\allowbreak{}matrix\_\allowbreak{}gauge}{Development.\allowbreak{}AbelianVarietyTheory.\allowbreak{}Proofs.\allowbreak{}DivisorProofs.\allowbreak{}ExistsIntegralUnitMatrixGauge}{309--320}
Define the power series on \(\widehat G^{\mathcal C}\) by
\[
 F(y)=G\Bigl(\Bigl(\sum_{\alpha\in\mathcal C}{}_{\widehat G}\,[x_{\alpha,b}]y_\alpha\Bigr)_{b=1}^N\Bigr)\,U(y)^k.
\]
Since the formal group law and \(\mathcal O_E\)-action have coefficients in \(V\) and zero constant term, substitution preserves integrality over \(\mathcal O_C\), so \(F\) is integral over \(\mathcal O_C\). Let \(\exp\) be the formal exponential of \(\widehat G\) over \(C\). Because the \(\mathcal O_E\)-action is linear in formal logarithm coordinates and \(u_{\tau,a}\) is a \(\tau\)-eigenvector, substituting \(y_\alpha=\exp(\sum_{a=1}^i t_{\alpha,a}u_{\tau,a})\) sends \(\sum_{\alpha\in\mathcal C}{}_{\widehat G}\,[x_{\alpha,b}]y_\alpha\) to \(\exp(\sum_{\alpha,a}\tau(x_{\alpha,b})t_{\alpha,a}u_{\tau,a})\) and sends \([w_0\sum_b\bar x_{\alpha,b}x_{\alpha',b}]y_{\alpha'}\) to \(\exp(\sum_a\tau'(w_0\sum_b\bar x_{\alpha,b}x_{\alpha',b})t_{\alpha',a}u_{\tau',a})\). On the other hand, writing invariant derivations in exponential coordinates identifies \(\mathcal N_k(G)\) with the derivative at \(0\), by the block-determinant polynomial of \(\mathcal D_c\), of the same universal power series evaluated at linear substitution coefficients \((B_{\tau,i})_{j,b}\) and cross-coefficients \(0\).
\leanref{GlobalTheory.\allowbreak{}ExpCoord.\allowbreak{}constantCoeff\_\allowbreak{}dpolyN\_\allowbreak{}eq\_\allowbreak{}derivAt0}{Development.\allowbreak{}AbelianVarietyTheory.\allowbreak{}Proofs.\allowbreak{}DivisorProofs.\allowbreak{}ConstantCoeffDpolyNEqDerivAt0}{226--308}
Since this derivative at \(0\) depends only on a finite jet and is polynomial in the substitution coefficients, it is Lipschitz in \(v_\ell\) on integral coefficient arrays with a constant \(K\ge0\) depending on \(G\), \(k\), \(P_{\rm col}\), \(\widehat G\), and \(h_L\). Comparing the two coefficient arrays yields \(v_\ell(\mathcal N_k(G)-\mathcal D_c(F))\ge n'-K\).
\leanref{GlobalTheory.\allowbreak{}FGDepth.\allowbreak{}derivAt0\_\allowbreak{}subst\_\allowbreak{}linear\_\allowbreak{}lipschitz}{Development.\allowbreak{}AbelianVarietyTheory.\allowbreak{}Proofs.\allowbreak{}DivisorProofs.\allowbreak{}DerivAt0SubstLinearLipschitz}{77--85}
\leanref{GlobalTheory.\allowbreak{}ColApprox.\allowbreak{}exists\_\allowbreak{}germ\_\allowbreak{}close}{Development.\allowbreak{}AbelianVarietyTheory.\allowbreak{}Proofs.\allowbreak{}DivisorProofs.\allowbreak{}SubstTwoPointsPowerSeriesZeroOfVanishFirstAxis}{331--429}

For (ii), for each block \(f=(\tau,i)\) and every \(n_0\ge0\), choose the isotropic approximants \(x_{n_0,\alpha,b}\) of Lemma~\ref{B8.28k}(b), where \(\alpha=(f,j)\in\mathcal C\). Lemma~\ref{B8.28ud} supplies matrices \(A_{\alpha,b}\in M_d(\mathcal O_C)\) such that
\[
 v_\ell\bigl((A_{\alpha,b}-d[\varepsilon_{\tau,n_0}x_{n_0,\alpha,b}])_{rs}\bigr)\ge n_0,
 \qquad
 A_{\alpha,b}u_{\tau,a}=(B_{\tau,i})_{j,b}u_{\tau,a}.
\]
Multiplication of every row in the block \(f\) by \(\varepsilon_{\tau,n_0}\) preserves its Hermitian Gram divisibility by \(\ell^{n_0}\); different blocks have disjoint supports, so their mixed Gram sums vanish. Define
\[
 \psi_b(Y)=\exp_{\widehat G}\left(
       \sum_{\alpha\in\mathcal C}A_{\alpha,b}\log_{\widehat G}Y_\alpha\right).
\]
The second assertion of Lemma~\ref{B8.28ud}, applied to \(w_\alpha=\log_{\widehat G}Y_\alpha\), makes \(\psi\) integral, with zero constant term. Its eigenvector identity gives \(\psi(\mathbf e(t))=\mathbf c(t)\), where
\[
 \mathbf e_\alpha(t)=\exp_{\widehat G}\left(\sum_a t_{\alpha,a}u_{\tau,a}\right),
 \qquad
 \mathbf c_b(t)=\exp_{\widehat G}\left(
       \sum_{\alpha=(\tau,i,j),a}(B_{\tau,i})_{j,b}t_{\alpha,a}u_{\tau,a}\right).
\]
Lemma~\ref{B8.28ue} shows that \(U_0=W\circ\psi\) is an integral unit, where \(W=\exp(\sum_bh_L(x_b))\).
\leanref{GlobalTheory.\allowbreak{}IsoTriv.\allowbreak{}Incl.\allowbreak{}inclSt\_\allowbreak{}of\_\allowbreak{}lim}{Development.\allowbreak{}AbelianVarietyTheory.\allowbreak{}Proofs.\allowbreak{}DivisorProofs.\allowbreak{}ColumnInclusionIntegralOfLimitSteps}{77--99}
\leanref{GlobalTheory.\allowbreak{}T.\allowbreak{}limitIsotropicFrameIntegral\_\allowbreak{}of\_\allowbreak{}steps}{Development.\allowbreak{}AbelianVarietyTheory.\allowbreak{}Proofs.\allowbreak{}DivisorProofs.\allowbreak{}LimitIsotropicFrameIntegralOfStepsStatement}{6--8}
\leanref{GlobalTheory.\allowbreak{}T.\allowbreak{}columnInclusionIntegral\_\allowbreak{}of\_\allowbreak{}limitSteps}{Development.\allowbreak{}AbelianVarietyTheory.\allowbreak{}Proofs.\allowbreak{}DivisorProofs.\allowbreak{}ColumnInclusionIntegralOfLimitStepsStatement}{6--8}
For each block \(f=(\tau,i)\), choose its assigned leading columns \(r_f(1),\ldots,r_f(i)\) and put \(M_f=( (B_{\tau,i})_{j,r_f(j')} )_{j,j'}\). This is a unit minor, and its entries are approximable through \(\tau(\mathcal O_E)\). By Lemma~\ref{B8.28ua}, the entries of \(M_f^{-1}\) are approximable. Apply Lemma~\ref{B8.28ud} to these entries and obtain matrices \(A'_{f,j',j}\) acting on the \(\tau\)-eigenvectors by \((M_f^{-1})_{j',j}\). Define
\[
 \rho_{f,j}(X)=\exp_{\widehat G}\left(
      \sum_{j'=1}^iA'_{f,j',j}\log_{\widehat G}X_{r_f(j')}\right).
\]
These series are integral and have zero constant term by Lemma~\ref{B8.28ud}. On the column curve, all other blocks vanish in the selected columns, and the identity \(M_fM_f^{-1}=1\) gives \(\rho(\mathbf c(t))=\mathbf e(t)\).
\leanref{GlobalTheory.\allowbreak{}IsoTriv.\allowbreak{}retrSt\_\allowbreak{}of}{Development.\allowbreak{}AbelianVarietyTheory.\allowbreak{}Proofs.\allowbreak{}DivisorProofs.\allowbreak{}ColumnInclusionRetractionOfLimitSteps}{47--117}
\leanref{GlobalTheory.\allowbreak{}IsoTriv.\allowbreak{}limitEntriesMatrixInverseClosed\_\allowbreak{}proof}{Development.\allowbreak{}AbelianVarietyTheory.\allowbreak{}Proofs.\allowbreak{}DivisorProofs.\allowbreak{}LimitEntriesMatrixInverseClosedProof}{75--99}
\leanref{GlobalTheory.\allowbreak{}T.\allowbreak{}columnInclusionRetraction\_\allowbreak{}of\_\allowbreak{}limitSteps}{Development.\allowbreak{}AbelianVarietyTheory.\allowbreak{}Proofs.\allowbreak{}DivisorProofs.\allowbreak{}ColumnInclusionRetractionOfLimitStepsStatement}{6--8}
Let \(W=e^{\sum_{b=1}^N h_L(x_b)}\), so that \(U_0=\psi^*W\), and let \(U'=U_0^{-1}\). Apply Lemma~\ref{B8.28ub} with the block exponential curve as \(\mathbf e\), the column exponential curve as \(\mathbf c\), and the maps \(\psi,\rho\) above. It gives an integral series \(G=\rho^*(F(U')^k)\) whose product with \(W^k\), restricted to the column curve, equals \(F\) restricted to the block exponential curve. Applying the block determinant differential and the exponential-coordinate identity gives \(\mathcal N_k(G)=\mathcal D_c(F)\).
\leanref{GlobalTheory.\allowbreak{}T.\allowbreak{}logIntegralUnitsSubgroupExpHom\_\allowbreak{}proof}{Development.\allowbreak{}AbelianVarietyTheory.\allowbreak{}Proofs.\allowbreak{}DivisorProofs.\allowbreak{}LogIntegralUnitsSubgroupExpHomProofStatement}{6--8}
\leanref{GlobalTheory.\allowbreak{}T.\allowbreak{}nablaOperatorFrameSeries\_\allowbreak{}of\_\allowbreak{}steps}{Development.\allowbreak{}AbelianVarietyTheory.\allowbreak{}Proofs.\allowbreak{}DivisorProofs.\allowbreak{}NablaOperatorFrameSeriesOfStepsStatement}{6--8}
\leanref{GlobalTheory.\allowbreak{}ColumnGerm.\allowbreak{}exists\_\allowbreak{}nabX\_\allowbreak{}eq}{Development.\allowbreak{}AbelianVarietyTheory.\allowbreak{}Proofs.\allowbreak{}DivisorProofs.\allowbreak{}ExistsIntegralGermOnCopiesSameValue}{47--87}

The complete proof in the archive is contained in \Lean{GlobalTheory.\allowbreak{}ColApprox.\allowbreak{}exists\_\allowbreak{}germ\_\allowbreak{}cols} and \Lean{GlobalTheory.\allowbreak{}ColumnGerm.\allowbreak{}exists\_\allowbreak{}germ\_\allowbreak{}on\_\allowbreak{}copies}.
\end{proof}

\subsection*{Proof of \ref{B8.28o}}\label{prf:B8.28o}\phantomsection
\begin{proof}
Fix \(k\ge0\), and choose the field \((C,v_\ell)\) and integral formal \(V\)-coordinates of \(\widehat G\) provided by Proposition~\ref{B8.28u}. Any integral formal coordinates of \(\widehat G\) over an algebraically closed complete valued extension \((C',v_\ell')\) of \(\operatorname{Frac}V\) differ from the \(V\)-coordinates by an invertible integral power-series substitution with zero constant term on each copy of \(\widehat G\). Such a substitution permutes the integral power series over \(\mathcal O_{C'}\) while leaving the invariant derivations and \(h_L\) unchanged. In the \(V\)-coordinates we have \(h_L(0)=0\) and \(P_{\rm col}(D)\) has order \(D\). For any integral series \(G=\sum_\mu c_\mu X^\mu\) over \(\mathcal O_{C'}\), the value \(\mathcal N_k(G)\) equals the finite sum \(\sum_{|\mu|\le D}c_\mu\mathcal N_k(X^\mu)\). Each monomial \(X^\mu\) has coefficients in \(V\), so \(\mathcal N_k(X^\mu)\in\operatorname{Frac}V\) has the same normalised valuation in \(C'\) and in \(C\), while \(v_\ell'(c_\mu)\ge0\). By the ultrametric inequality, \(\delta_k(P_{\rm col};y)\) equals the infimum of \(v_\ell(\mathcal N_k(G))\) over power series \(G\) with coefficients in \(V\) in the \(V\)-coordinates, and also equals the infimum over integral power series \(G\) over \(\mathcal O_C\) in the \(V\)-coordinates.

Given an integral power series \(G\) over \(\mathcal O_C\) on \(\widehat G^N\) and \(m_0\ge0\), take \(K\) from Proposition~\ref{B8.28u}(i) and choose an integer \(n'\ge m_0+K\). For each block \((\tau,i)\) with \(c_{\tau,i}\ne0\), Lemma~\ref{B8.28k}(b) with \(n_0=n'\) supplies \(x_{(\tau,i,j),b}\in\mathcal O_E\) supported on the \(2i\) columns assigned to that block, with \(\tau(x_{(\tau,i,j),b})\equiv(B_{\tau,i})_{j,b}\) modulo \(\ell^{n'}V_\lambda\) and \(\sum_{b=1}^N x_{(\tau,i,j),b}\bar x_{(\tau,i,j'),b}\in\ell^{n'}\mathcal O_E\). For \(\alpha,\alpha'\in\mathcal C\) in distinct blocks the column supports are disjoint, so \(\sum_{b=1}^N\bar x_{\alpha,b}x_{\alpha',b}=0\). Within a single block, conjugation places \(\sum_{b=1}^N\bar x_{\alpha,b}x_{\alpha',b}\) in \(\ell^{n'}\mathcal O_E\), and multiplication by \(w_0\in\mathcal O_E\) preserves \(\ell^{n'}\mathcal O_E\). Passing to \(V\) and \(C\) gives both valuation hypotheses of Proposition~\ref{B8.28u}(i). Hence there is an integral series \(F\) over \(\mathcal O_C\) on \(\widehat G^{\mathcal C}\) with \(v_\ell(\mathcal N_k(G)-\mathcal D_c(F))\ge n'-K\ge m_0\). Whenever \(\mathcal N_k(G)\ne0\), choosing \(m_0>v_\ell(\mathcal N_k(G))\) yields \(v_\ell(\mathcal D_c(F))=v_\ell(\mathcal N_k(G))\) by the ultrametric equality. Combining this with Proposition~\ref{B8.28u}(ii) identifies \(\inf_G v_\ell(\mathcal N_k(G))\) with \(\inf_F v_\ell(\mathcal D_c(F))\).

By the coordinate and monomial argument of the first paragraph on \(i\) copies of \(\widehat G\), each block depth \(\nu_{\tau,i}(c_{\tau,i};y)\) equals the corresponding infimum in the \(V\)-coordinates over \(C\). It is finite and nonnegative by Lemma~\ref{B8.28h}. Testing \(\mathcal D_c\) on external products \(F=\prod_{c_{\tau,i}\ne0}F_{\tau,i}\) of integral series on the individual blocks gives \(\inf_F v_\ell(\mathcal D_c(F))\le\sum_{c_{\tau,i}\ne0}\nu_{\tau,i}(c_{\tau,i};y)\). Conversely, expanding any integral series \(F\) on \(\widehat G^{\mathcal C}\) in the coordinates of each block makes \(\mathcal D_c(F)\) a finite \(\mathcal O_C\)-linear combination of product evaluations, so the ultrametric inequality gives the reverse bound. Since \(\nu_{\tau,i}(0;y)=0\) by Lemma~\ref{B8.28h}, adding the blocks with \(c_{\tau,i}=0\) gives \(\nu_c(y)\).
\end{proof}

\subsection*{Proof of \ref{B8.28va}}\label{prf:B8.28va}\phantomsection
\begin{proof}
Each summand of the second-order expression is zero. Its coefficient vanishes for a nonconjugate pair, and the assumed harmonic identity makes the inner sum zero for a conjugate pair. If harmonicity is initially stated for one member of each conjugate pair, the other identities follow by interchanging the two partial derivatives.

For nonnegative integers \(A,B\), one has \(2\min(A,B)\le A+B\). Hence
\[
 2e(\mu)
 \le\sum_{b,q}(n^+_{q,b}(\mu)+n^-_{q,b}(\mu))
 =\sum_{\ell,b}\mu_{\ell,b}
 =\sum_\ell a_\ell=D.
\]
Since \(e(\mu)\) is an integer, this proves the bound on the number of opposite pairs.

The ultrametric inequality gives closure of \(J_\eta\) under addition. Multiplication by an element of \(V\) preserves its defining bound, and \(\eta\ge0\) gives \(1\in J_\eta\). A product of \(m\) elements of \(J_\eta\) has normalized valuation at least \(-m\eta\). Every finite \(V\)-linear combination of such products has the same lower bound. Finally, \(\lfloor D/2\rfloor\le D/2\) and \(\eta\ge0\) give the last assertion.
\end{proof}

\subsection*{Proof of \ref{B8.28vb}}\label{prf:B8.28vb}\phantomsection
\begin{proof}
Write \(a_\ell=\nabla_{\xi^0_\ell}\) and
\(b_\ell=\nabla_{(0,\psi_\ell)}\). Within either sign block, the opposite-pair assumptions allow the \(b\)'s to pass the \(a\)'s. Moving a plus-block \(b_\ell\) past a minus-block \(a_{\ell'}\) contributes
\[
 m\beta((0,\psi_\ell),\xi^0_{\ell'}).
\]
At the right end every remaining \(b\) kills \(\iota g\).
Consequently, if \(\Delta_0\) is the single-column second-order operator with precisely these contraction coefficients, then
\[
 O_{\xi^0+(0,\psi)}(Q)\iota g
 =\sum_{r=0}^{\deg Q}
       O_{\xi^0}\left(\frac{\Delta_0^rQ}{r!}\right)\iota g.
\]
This identity follows by induction from
\(ba^t=a^tb+t[b,a]a^{t-1}\), since each displayed commutator is scalar.

For (a), take \(\xi^0_\ell=(u_\ell,0)\) and
\(\psi_\ell=\phi_\ell\). The horizontal evaluation identity turns each term into an invariant differential value. Operators on different columns commute. The multinomial formula combines the product of their finite exponential expansions into \(\sum_r\Delta^rP/r!\). Since \(\Delta P=0\), the term with \(r=0\) gives the asserted equality.

For (b), every contraction coefficient belongs to \(J\), since it equals
\(-m\sum_i\psi_{\ell,i}u_{\ell',i}\).
The factorial in the expansion preserves integral polynomial coefficients. More explicitly, for contraction multiplicities \(t_{ij}\) from plus indices \(i\) to minus indices \(j\), let
\[
 r_i=\sum_jt_{ij},\qquad c_j=\sum_it_{ij},\qquad
 r=\sum_{ij}t_{ij}.
\]
On a monomial with plus exponents \(a_i\) and minus exponents \(b_j\), the coefficient accompanying
\(\prod_{ij}\kappa_{ij}^{t_{ij}}\) in \(\Delta_0^r/r!\) is
\[
 \frac{\prod_i(a_i)_{r_i}\prod_j(b_j)_{c_j}}
      {\prod_{ij}t_{ij}!}
 =\prod_i\left(\binom{a_i}{r_i}
              \frac{r_i!}{\prod_jt_{ij}!}\right)
       \prod_j(b_j)_{c_j}\in\Z,
\]
where \((a)_r=a(a-1)\cdots(a-r+1)\), and inadmissible contractions give zero. Hence the coefficients of the \(r\)-th term lie in \(J^r\). The support condition on \(\kappa\) bounds the number of contractions in each class by the smaller of its two sign totals. The stated bound therefore makes the total number over all columns at most \(e\).

The conjugation hypothesis says that every \(a_\ell\) preserves
\((\iota E_0)V[[x,z]]\). Every ordered horizontal word applied to
\(\iota(E_0f_b)\) consequently has constant coefficient in \(V\).
A term with \(r\) contractions over all columns lies in
\(J^r\subseteq J^e\), using \(1\in J\). Finite sums and multiplication by the coefficients of \(P\) preserve \(J^e\), proving (b).
\end{proof}

\subsection*{Proof of \ref{B8.28vc}}\label{prf:B8.28vc}\phantomsection
\begin{proof}[Proof (outline)]
The formal-group construction gives vector fields \(\mathcal V_\xi\) on \(K[[x,z]]\), linear in \(\xi=(u,\phi)\), and the connection has the form
\[
 \nabla_\xi=\mathcal V_\xi-m\lambda_u,
 \qquad \lambda_u=\sum_i(z_i-\sigma_i(x))u_i,
\]
where \(\lambda_u\) acts by multiplication. The vector fields restrict to \(D_u\) on series in \(x\). Commutativity of the invariant derivations gives
\[
 [\mathcal V_\xi,\mathcal V_\eta]=0,
 \qquad
 \mathcal V_\xi(\lambda_{u'})-\mathcal V_\eta(\lambda_u)
       =-\beta(\xi,\eta),\quad \eta=(u',\phi').
\]
Expanding the commutator of
\(\mathcal V_\xi-m\lambda_u\) and
\(\mathcal V_\eta-m\lambda_{u'}\) proves the commutator identity.
Taking \(u=0\) proves vertical annihilation.

The gauge's zero-axis identity gives a substitution by zero-constant series
\(\mathcal E:K[[x,z]]\to K[[x]]\) that fixes series in \(x\) and satisfies
\[
 \mathcal E\nabla_{(u,0)}=D_u\mathcal E.
\]
Iterating this identity and taking constant coefficients proves horizontal evaluation.
\leanref{ReciprocityTheory.\allowbreak{}ChartRestriction.\allowbreak{}G.\allowbreak{}TB.\allowbreak{}WickVanishing.\allowbreak{}Val.\allowbreak{}E\_\allowbreak{}liftX}{ReciprocityTheory.\allowbreak{}ReductiveGroups.\allowbreak{}ArithmeticGeometry.\allowbreak{}TaylorDiscriminants.\allowbreak{}WickVanishing.\allowbreak{}HarmonicValueIntegral}{62--84}
\leanref{ReciprocityTheory.\allowbreak{}ChartRestriction.\allowbreak{}G.\allowbreak{}TB.\allowbreak{}WickVanishing.\allowbreak{}Val.\allowbreak{}E\_\allowbreak{}nablaPoly}{ReciprocityTheory.\allowbreak{}ReductiveGroups.\allowbreak{}ArithmeticGeometry.\allowbreak{}TaylorDiscriminants.\allowbreak{}WickVanishing.\allowbreak{}HarmonicValueIntegral}{133--141}

Every element of \(K\) becomes integral after multiplication by a positive power of the residue characteristic. For \(u,\phi\in V^d\), differentiating the canonical gauge's integral cocycle and adjointness identities, with \(w_0+\bar w_0=1\), gives
\[
 (\iota E_0)^{-1}\nabla_{(u,\phi-Su)}((\iota E_0)G)
       \in V[[x,z]]
 \qquad(G\in V[[x,z]]).
\]
This is the required integral conjugation identity.

The complete connection-identity proofs in the archive are contained in \Lean{ReciprocityTheory.\allowbreak{}ChartRestriction.\allowbreak{}G.\allowbreak{}TB.\allowbreak{}WickVanishing.\allowbreak{}Nab.\allowbreak{}connection\_\allowbreak{}second\_\allowbreak{}identity}, \Lean{ReciprocityTheory.\allowbreak{}ChartRestriction.\allowbreak{}G.\allowbreak{}TB.\allowbreak{}WickVanishing.\allowbreak{}Nab.\allowbreak{}connection\_\allowbreak{}kills\_\allowbreak{}series}, \Lean{ReciprocityTheory.\allowbreak{}ChartRestriction.\allowbreak{}G.\allowbreak{}TB.\allowbreak{}WickVanishing.\allowbreak{}Val.\allowbreak{}harmonic\_\allowbreak{}value\_\allowbreak{}integral}, and \Lean{ReciprocityTheory.\allowbreak{}ChartRestriction.\allowbreak{}G.\allowbreak{}TB.\allowbreak{}WickVanishing.\allowbreak{}Nab.\allowbreak{}connection\_\allowbreak{}first\_\allowbreak{}identity}.
\end{proof}

\subsection*{Proof of \ref{B8.28v}}\label{prf:B8.28v}\phantomsection
\begin{proof}[Proof (outline)]
In both cases, expanding \(F\) in monomials and applying the ultrametric inequality reduces the valuation bound to external products \(F=f_1\boxtimes\dots\boxtimes f_N\) of single-factor power series \(f_b\) with coefficients in \(V\). The integral chart supplies a descended formal biextension over \(V\) whose canonical gauge \(h\) has constant term \(0\) and whose \(\mathcal O_E\)-action has the frame directions \(u_{\tau,i}\) as \(\tau\)-eigenvectors.
\leanref{GlobalTheory.\allowbreak{}PlaceInt.\allowbreak{}integralBiext\_\allowbreak{}gauge\_\allowbreak{}eig\_\allowbreak{}bound}{Development.\allowbreak{}AbelianVarietyTheory.\allowbreak{}Proofs.\allowbreak{}RingsOfIntegers.\allowbreak{}IntegralChartFormalBiextensionGaugeBound}{50--109}

For (1), suppose \(\operatorname{disc}E\in V^\times\). If \(V\) has residue characteristic \(0\), then \(\Q\subset V\); a canonical gauge \(h_0\) exists directly over \(V\), and harmonicity makes the evaluation independent of the choice of canonical gauge, so \(z\in V\).
\leanref{GlobalTheory.\allowbreak{}ChartUnram.\allowbreak{}value\_\allowbreak{}mem\_\allowbreak{}rat}{Development.\allowbreak{}AbelianVarietyTheory.\allowbreak{}Proofs.\allowbreak{}RingsOfIntegers.\allowbreak{}HarmonicValueIntegralOfUnitDiscriminant}{68--97}
If \(V\) has positive residue characteristic, then \(\operatorname{Frac}V\) is obtained from \(V\) by inverting positive integers. Because \(|\operatorname{disc}E|\in V^\times\), the trace-dual basis of \(\mathcal O_E\) produces integral idempotent data over \(V\) for each embedding \(\tau\). Applying these idempotents to the formal-biextension \(\mathcal O_E\)-action lifts the linear biextension vector attached to each frame direction \(\ell=(\tau,i')\) to a vector \(\varphi_\ell\) over \(V\) whose difference from that vector pairs to \(0\) under the canonical sesquilinear form against every non-conjugate eigenspace. Apply Lemmas~\ref{B8.28vc} and \ref{B8.28vb} with \(J=V\), \(\psi_\ell=0\), \(\xi^0_\ell=(u_\ell,\varphi_\ell-Su_\ell)\), and \(E_0=\exp(kh_L)\). Take the class map to identify each conjugate pair of embeddings and choose opposite signs on its two members. The eigenlift pairing identity gives the opposite-pair vanishing hypotheses of Lemma~\ref{B8.28vb}(a), and together with harmonicity it gives its second-order vanishing hypothesis. The integral conjugation identity, opposite-pair vanishing, and the count bound of Lemma~\ref{B8.28va} give part (b), so \(z\in V\) and \(v(z)\ge0\).
\leanref{ReciprocityTheory.\allowbreak{}ChartRestriction.\allowbreak{}G.\allowbreak{}PullbackSections.\allowbreak{}Pf.\allowbreak{}EigA.\allowbreak{}idem\_\allowbreak{}data}{ReciprocityTheory.\allowbreak{}ReductiveGroups.\allowbreak{}ArithmeticGeometry.\allowbreak{}PullbackSections.\allowbreak{}HolomorphicFamilies.\allowbreak{}EigenLiftsOfDirectionsBiext}{56--74}
\leanref{ReciprocityTheory.\allowbreak{}ChartRestriction.\allowbreak{}G.\allowbreak{}PullbackSections.\allowbreak{}Pf.\allowbreak{}EigA.\allowbreak{}eigLift\_\allowbreak{}core}{ReciprocityTheory.\allowbreak{}ReductiveGroups.\allowbreak{}ArithmeticGeometry.\allowbreak{}PullbackSections.\allowbreak{}HolomorphicFamilies.\allowbreak{}BiextCoeffSecondLinear}{169--239}
\leanref{GlobalTheory.\allowbreak{}ChartUnram.\allowbreak{}value\_\allowbreak{}mem\_\allowbreak{}idem}{Development.\allowbreak{}AbelianVarietyTheory.\allowbreak{}Proofs.\allowbreak{}RingsOfIntegers.\allowbreak{}HarmonicValueIntegralOfUnitDiscriminant}{98--161}

For (2), let \(p\mid Q_Sqd\) be a prime that is not a unit in \(V\), and write \(D_E=p^{a_0}m_0\) with \(p\nmid m_0\), so that \(\delta_p=a_0/[E:\Q]\). By the host construction, \(E/\Q\) is Galois, every prime dividing \(\operatorname{disc}E\) divides \(Q_Sqd\), and complex conjugation moves every prime ideal \(\mathfrak p\mid p\) of \(\mathcal O_E\).
\leanref{GlobalTheory.\allowbreak{}PlaceInt.\allowbreak{}host\_\allowbreak{}galois\_\allowbreak{}discr\_\allowbreak{}moves}{Development.\allowbreak{}AbelianVarietyTheory.\allowbreak{}Proofs.\allowbreak{}RingsOfIntegers.\allowbreak{}HostFieldGaloisDegreeDiscrConjMovesPrimes}{26--55}
Consequently, pulling back the maximal ideal of \(V\) along the embeddings \(\mathcal O_E\to V\) partitions the embeddings into conjugate-opposite pairs of prime ideals \((\mathfrak p,\bar{\mathfrak p})\) above \(p\), separated by integral place idempotents. Combining these place idempotents with generators of the different ideal yields for each direction \(\ell\) a decomposition of the lifted biextension vector into an integral part \(\varphi_\ell\) over \(V\) plus a polar part \(\psi_\ell\) over \(\operatorname{Frac}V\) satisfying the opposite-pair vanishing laws and the entry bound \(v(\psi_{\ell,i})/e\ge-(v(p)/e)\,\delta_p\).
\leanref{GlobalTheory.\allowbreak{}PlaceInt.\allowbreak{}exists\_\allowbreak{}ramified\_\allowbreak{}eigenLifts}{Development.\allowbreak{}AbelianVarietyTheory.\allowbreak{}Proofs.\allowbreak{}RingsOfIntegers.\allowbreak{}ExistsRamifiedEigenLiftsOfChartDirections}{177--301}
\leanref{GlobalTheory.\allowbreak{}T.\allowbreak{}exists\_\allowbreak{}ramified\_\allowbreak{}eigen\_\allowbreak{}lifts\_\allowbreak{}of\_\allowbreak{}chart\_\allowbreak{}directions}{Development.\allowbreak{}AbelianVarietyTheory.\allowbreak{}Proofs.\allowbreak{}RingsOfIntegers.\allowbreak{}ExistsRamifiedEigenLiftsOfChartDirectionsStatement}{13--15}
Let \(K=\operatorname{Frac}V\), and let \(S\) and the connection operators be as in Lemma~\ref{B8.28vc}, with \(m=k\). Put
\[
 J=\{x\in K:v(x)/e\ge-(v(p)/e)\delta_p\},\quad
 \xi^0_\ell=(u_\ell,\varphi_\ell-Su_\ell),\quad
 \eta_\ell=\varphi_\ell+\psi_\ell-Su_\ell,
 \qquad E_0=\exp(kh_L).
\]
Then \(J\) is a \(V\)-submodule containing \(1\), \(\xi^0_\ell\in H_V\), and \(\psi_\ell\in J^d\). Let \(L\) be the finite set of frame directions \(\ell=(\tau,i)\), let \(\Pi\) be their finite set of prime-pair classes, and write \(c:L\to\Pi\) and \(s:L\to\{+,-\}\) for the class and sign maps supplied by the eigenlifts. The contraction coefficient in Lemma~\ref{B8.28vb}(a) is
\[
 \kappa_{\ell\ell'}=
 \begin{cases}
 k\beta((0,\eta_\ell),(u_{\ell'},0)),&c(\ell)=c(\ell'),\ s(\ell)=+,\ s(\ell')=-,\\
 0,&\text{otherwise}.
 \end{cases}
\]
The eigenlift pairing identity makes it vanish unless the types of \(\ell\) and \(\ell'\) are conjugate. Lemma~\ref{B8.28va} therefore gives \(\sum_{b,\ell,\ell'}\kappa_{\ell\ell'}\partial_{X_{\ell,b}}\partial_{X_{\ell',b}}R=0\). The eigenlift pairing identity \(\sum_i\eta_{\ell',i}u_{\ell,i}=0\) holds whenever the type of \(\ell'\) is not the conjugate of the type of \(\ell\); since \(c(\bar\tau)=c(\tau)\) and \(s(\bar\tau)\ne s(\tau)\), it gives the opposite-pair vanishing hypotheses of Lemma~\ref{B8.28vb}(a). Lemma~\ref{B8.28vb}(a), applied to \(g_b=E_0f_b\), identifies \(z\) with \(\mathcal E_{\xi^0+(0,\psi)}(R;E_0f_1,\ldots,E_0f_N)\).

The three opposite-pair vanishing conditions are precisely those supplied by the eigenlifts. Lemma~\ref{B8.28vc} gives integral conjugation for \(\xi^0_\ell\), and Lemma~\ref{B8.28va} bounds the sum of the columnwise opposite-pair counts by \(\lfloor D/2\rfloor\). Lemma~\ref{B8.28vb}(b) now places \(z\) in \(J^{\lfloor D/2\rfloor}\). Its valuation satisfies
\[
 v(z)/e\ge-\lfloor D/2\rfloor\,(v(p)/e)\delta_p
          \ge-(v(p)/e)\delta_p D/2
\]
by the last assertion of Lemma~\ref{B8.28va}.
\leanref{ReciprocityTheory.\allowbreak{}ChartRestriction.\allowbreak{}G.\allowbreak{}PullbackSections.\allowbreak{}Pf.\allowbreak{}Jv}{ReciprocityTheory.\allowbreak{}ReductiveGroups.\allowbreak{}ArithmeticGeometry.\allowbreak{}PullbackSections.\allowbreak{}HolomorphicFamilies.\allowbreak{}ChartValueRamifiedOfLifts}{35--46}
\leanref{ReciprocityTheory.\allowbreak{}ChartRestriction.\allowbreak{}G.\allowbreak{}PullbackSections.\allowbreak{}Pf.\allowbreak{}sum\_\allowbreak{}pairMin\_\allowbreak{}le}{ReciprocityTheory.\allowbreak{}ReductiveGroups.\allowbreak{}ArithmeticGeometry.\allowbreak{}PullbackSections.\allowbreak{}HolomorphicFamilies.\allowbreak{}ChartValueRamifiedOfLifts}{110--127}
\leanref{ReciprocityTheory.\allowbreak{}ChartRestriction.\allowbreak{}G.\allowbreak{}PullbackSections.\allowbreak{}Pf.\allowbreak{}chartValJ}{ReciprocityTheory.\allowbreak{}ReductiveGroups.\allowbreak{}ArithmeticGeometry.\allowbreak{}PullbackSections.\allowbreak{}HolomorphicFamilies.\allowbreak{}ChartValJ}{31--159}
\leanref{ReciprocityTheory.\allowbreak{}ChartRestriction.\allowbreak{}G.\allowbreak{}PullbackSections.\allowbreak{}Pf.\allowbreak{}mem\_\allowbreak{}Jv\_\allowbreak{}pow}{ReciprocityTheory.\allowbreak{}ReductiveGroups.\allowbreak{}ArithmeticGeometry.\allowbreak{}PullbackSections.\allowbreak{}HolomorphicFamilies.\allowbreak{}ChartValueRamifiedOfLifts}{56--81}

The complete proof in the archive is contained in \Lean{GlobalTheory.\allowbreak{}PlaceInt.\allowbreak{}chBound\_\allowbreak{}zero\_\allowbreak{}of\_\allowbreak{}isUnit\_\allowbreak{}discr} and \Lean{GlobalTheory.\allowbreak{}PlaceInt.\allowbreak{}ChartRam.\allowbreak{}chBound\_\allowbreak{}ramified}.
\end{proof}

\subsection*{Proof of \ref{B8.28p}}\label{prf:B8.28p}\phantomsection
\begin{proof}
(1) By \ref{B2.46}, \(E\subset\Q(\mu_{Q_Sqd})\), so every prime divisor of \(\operatorname{disc}E\) divides \(Q_Sqd\), and each summand \((v(p)/e)\,\delta_p\,D/2\) is nonnegative. If every prime \(p\mid Q_Sqd\) is a unit in \(V\), then \(D_E\) is a product of units in \(V\), so \(\operatorname{disc}E\in V^\times\) and Proposition~\ref{B8.28v}(1) gives \(v(z)/e\ge0\), which dominates the negated sum. Otherwise some prime \(p\mid Q_Sqd\) is not a unit in \(V\); Proposition~\ref{B8.28v}(2) gives \(v(z)/e\ge-(v(p)/e)\,\delta_p\,D/2\), and subtracting the remaining nonnegative summands preserves the inequality.

(2) Put \(\rho_0=\sum_{p\mid Q_Sqd}\delta_p\,D/2\) and choose a nonnegative integer \(n_0>B+\rho_0\). Lemma~\ref{B8.28k}(c) gives \(P_1\in\operatorname{harm}(a,N)\) such that, on every trait \(V\) over \(\lambda\), \(P_1-P_{\rm col}=\ell^{n_0}R\) with \(R\in\operatorname{harm}_V(a,N)\). Fix a trait \(V\) over \(\lambda\), an algebraic integral frame \(y\) on \(V\), and \(k\ge0\). Since \(\lambda\mid\ell\), we have \(v(\ell)>0\) and \(\ell\) lies in the maximal ideal of \(V\). Any prime \(p\ne\ell\) is coprime to \(\ell\) in \(\Z\), hence also in \(V\), so \(p\in V^\times\) and \(v(p)=0\le v(\ell)\); for \(p=\ell\) equality holds. Applying (1) with \(e=v(\ell)\) in the integral formal \(V\)-coordinates of Proposition~\ref{B8.28u} therefore gives \(v(z)/v(\ell)\ge-\rho_0\) for every power series \(F\) over \(V\) on \(\widehat G^N\), where \(z=[R(D)(F\,e^{k\sum_{b=1}^N h_L(x_b)})](0)\).

As in the first paragraph of the proof of Lemma~\ref{B8.28o}, both \(\delta_k(P_1;y)\) and \(\delta_k(P_{\rm col};y)\) are computed as infima over power series \(F\) with coefficients in \(V\). For every such \(F\), linearity in the differential gives
\[
 \bigl[P_1(D)\bigl(F\,e^{k\sum_{b=1}^N h_L(x_b)}\bigr)\bigr](0)
 =\bigl[P_{\rm col}(D)\bigl(F\,e^{k\sum_{b=1}^N h_L(x_b)}\bigr)\bigr](0)+\ell^{n_0}z.
\]
Here \(v(\ell^{n_0}z)/v(\ell)=n_0+v(z)/v(\ell)\ge n_0-\rho_0>B\). By Lemma~\ref{B8.28o}, \(\delta_k(P_{\rm col};y)=\nu_c(y)\in[A,B]\). Since the perturbation has normalised valuation strictly greater than \(B\ge\delta_k(P_{\rm col};y)\), the ultrametric inequality on the infimum over \(F\) gives \(\delta_k(P_1;y)=\delta_k(P_{\rm col};y)\in[A,B]\). Taking \(A=A_\lambda\) and \(B=A_\lambda+\varrho\) from Definition~\ref{B8.28g}(b) yields the final statement.
\end{proof}

\subsection*{Proof of \ref{B8.28wa}}\label{prf:B8.28wa}\phantomsection
\begin{proof}[Proof (outline)]
Taking the product with the conjugate in \(m=\beta\alpha n'\) gives
\(m^2=n'^2AB\). In particular \(n'^2\), \(A\) and \(n'^2A\) are odd. Define
\[
 S_0=\mathcal P_{n',1}s_0,\qquad
 S_{1,b}=\mathcal T_{[\alpha\beta]z_{1,b},n'^2}S_0,
\]
\[
 S_{2,b}=\mathcal P_{\alpha,n'^2}S_{1,b},\qquad
 T_b=\mathcal T_{[\beta]z_{2,b},n'^2A}S_{2,b}.
\]
The displayed torsion conditions make both translations defined.
Normalised pullbacks compose, and for the torsion points and levels here they satisfy
\[
 \mathcal P_{\gamma,k}\mathcal T_{[\gamma]z,k}
 =\mathcal T_{z,k\gamma\bar\gamma}\mathcal P_{\gamma,k}.
\]
Apply this identity first with \(\gamma=\alpha\beta\), then with \(\gamma=\beta\). Uniqueness of the normalised pullbacks and symmetric translates gives
\[
 \mathcal P_{\beta,n'^2A}T_b=F_b.
\]

The four operations in the construction of \(T_b\) preserve the integral theta lattice of the split uniformisation. More precisely, a theta family of level \(k\) is in this lattice when its coefficient at every period index \(\chi\), measured in integral frames on the abelian part, has valuation at least \(Q_X(\chi)/k\). The principal family has this property. Normalised pullback preserves it, as does a symmetric finite-order torsion translation. The pullbacks used above start at level \(1\) or at the odd square \(n'^2\), which is \(1\) or at least \(4\); these are the level ranges in the integral pullback statement. Induction through the four operations gives the coefficient inequalities for \(T_b\). The theta expansion criterion on the integral model makes \(T_b\) regular. Its compatible finite-order germs are integral and have a unique inverse limit \(f_b\in V[[X]]\); multiplication of a section by \(r\in V\) multiplies its germ by \(r\).

The value of a harmonic differential on these sections is its value on their product germ in the integral biextension chart. This follows by identifying the formal group law and the \(\mathcal O_E\)-action in the generic and integral charts. The normalised pullback chain rule gives
\[
 \mathcal H_{m^2}(P;F_1,\ldots,F_n)
 =\nu_a(\beta)\mathcal H_{n'^2A}(P;T_1,\ldots,T_n).
\]
Indeed, \([\beta]\) acts on a direction of type \(\tau\) by \(\tau(\beta)\), and row homogeneity gives
\[
 \prod_\tau\tau(\beta)^{\sum_i a_{\tau,i}}=\nu_a(\beta).
\]
The rigidification fixes the scalar in the line-bundle pullback.
The identity for a polynomial with coefficients in the coefficient field follows by clearing its finitely many coefficient denominators and cancelling that nonzero scalar.

For \(P\ne0\), its total degree is \(D\). Multiplicativity of \(\nu_a\) gives
\[
 m^D=\nu_a(m)=\nu_a(\beta)\nu_a(\alpha)n'^D,
\]
which proves the asserted cancellation. For \(P=0\), both values vanish. Since \(n'\) is a unit, its inverse power has valuation zero. The assumed chart bound applies to the integral series
\(\prod_b f_b(X_b)\) and proves the conclusion.

The complete proof in the archive is contained in \Lean{GlobalTheory.\allowbreak{}PlaceInt.\allowbreak{}TermBound.\allowbreak{}term\_\allowbreak{}bound}.
\end{proof}

\subsection*{Proof of \ref{B8.28w}}\label{prf:B8.28w}\phantomsection
\begin{proof}
The base level is odd and squarefree, and
\(\mathfrak b\bar{\mathfrak b}=(M)\). If \(\ell_0\nmid M\), take
\[
 n'=M,\qquad A=B=1,\qquad\alpha=\beta=1.
\]
If \(\ell_0=g_i\), with \(\{i,j\}=\{0,1\}\), take
\[
 n'=g_j,\qquad A=B=g_i,\qquad
 \alpha=\bar\pi_i,\qquad\beta=\pi_i.
\]
In both cases \(n'\in(V')^\times\), \(M=\beta\alpha n'\),
\(\alpha\bar\alpha=A\), and \(\beta\bar\beta=B\).
\leanref{GlobalTheory.\allowbreak{}TorsFac.\allowbreak{}factorisation}{Development.\allowbreak{}AbelianVarietyTheory.\allowbreak{}Proofs.\allowbreak{}RingsOfIntegers.\allowbreak{}DvdOfCastDvd}{91--172}

Rational \(M^2\)-torsion and the theta normalisation give, for every characteristic \(x_b\), a point \(t_b\) killed by \(\mathfrak b^2\) such that
\[
 s_{x_b}=\mathcal T_{t_b,M^2}\mathcal T_{0,M^2}
                  \mathcal P_{M,1}s_0.
\]
Since \((M^2)\subset\mathfrak b^2\), one has \([M^2]t_b=0\).
Away from the level, the remaining torsion condition is this same equality. At \(\ell_0=g_i\), one has
\[
 n'\beta=g_j\pi_i=(\pi_i\pi_j)\bar\pi_j\in\mathfrak b,
 \qquad M\in\mathfrak b.
\]
Consequently \(n'^2A\beta=Mn'\beta\in\mathfrak b^2\), so
\([n'^2A\beta]t_b=0\). The other translation point is zero, hence it satisfies \([n'^2\alpha\beta]0=0\).
\leanref{GlobalTheory.\allowbreak{}RatData.\allowbreak{}exists\_\allowbreak{}thetaVal\_\allowbreak{}data}{Development.\allowbreak{}AbelianVarietyTheory.\allowbreak{}Proofs.\allowbreak{}RingsOfIntegers.\allowbreak{}ExistsThetaValData}{90--131}
\leanref{GlobalTheory.\allowbreak{}PlaceInt.\allowbreak{}uu\_\allowbreak{}through\_\allowbreak{}moves}{Development.\allowbreak{}AbelianVarietyTheory.\allowbreak{}Proofs.\allowbreak{}RingsOfIntegers.\allowbreak{}SymmetricTranslatesMultiplicationPullbackDescrip}{85--97}

Apply Lemma~\ref{B8.28wa} with \(m=M\), \(n=N\), \(z_{1,b}=0\), \(z_{2,b}=t_b\), and
\[
 \rho=\sum_{p\mid Q_Sqd}\frac{v'(p)}{e'}\delta_p\frac D2.
\]
Its chart hypothesis is Lemma~\ref{B8.28p}(1), on the integral chart of Proposition~\ref{B8.28u}. By the normalisation in Lemma~\ref{B8.28m},
\[
 \theta_{x,P}(y')
 =M^{-D}\mathcal H_{M^2}(P;s_{x_1},\ldots,s_{x_N}).
\]
The conclusion of Lemma~\ref{B8.28wa} is precisely the asserted inequality.
\end{proof}

\subsection*{Proof of \ref{B8.28q}}\label{prf:B8.28q}\phantomsection
\begin{proof}
(1) Since \((\bar\pi_0\bar\pi_1)(\pi_0\pi_1)=g_0g_1\ne0\), every factor \(\tau(\bar\pi_0\bar\pi_1)\) is nonzero, so \(\nu\in\mathcal O_{F_c}\) is nonzero. Each prime \(p\mid Q_Sqd\) and \(\nu\) have nonzero valuation at finitely many places of \(F_c\). Outside that finite union both \(\rho(w)\) and \(v_w(\nu)\) vanish.

(2) Fix \(P\in\operatorname{harm}(a,N)\), a tuple \(x\), a finite place \(w\), and an algebraic integral frame \(y\) on a trait \(V\) of ramification index \(e\) over \(w\). Passing to a finite extension \(V'\) of \(V\) of relative ramification index \(r'\ge1\) that splits the torus of the semiabelian degeneration \cite[Ch.~III]{faltings-chai-1990} and adjoins the coordinates of the \(M^2\)-torsion, the image algebraic integral frame \(y'\) has a split uniformisation and rational \(M^2\)-torsion; \(V'\) has ramification index \(e'=r'e\) over \(w\) and residue characteristic \(\ell_0\), the rational prime below \(w\). Take \(\alpha\) as in Proposition~\ref{B8.28w}. Since \(g_0\) and \(g_1\) are distinct primes, there is \(\alpha'\in\mathcal O_E\setminus\{0\}\) with \(\alpha\alpha'=\bar\pi_0\bar\pi_1\): for \(\ell_0\nmid M\) take \(\alpha'=\bar\pi_0\bar\pi_1\), while for \(\ell_0=g_0\) or \(\ell_0=g_1\) take \(\alpha'=\bar\pi_1\) or \(\alpha'=\bar\pi_0\). Thus \(\nu=\nu_a(\alpha)\nu_a(\alpha')\) with \(\nu_a(\alpha),\nu_a(\alpha')\in\mathcal O_{F_c}\setminus\{0\}\), so \(v_w(\nu_a(\alpha'))\ge0\) and \(v_w(\nu_a(\alpha))\le v_w(\nu)\).

On \(V'\) we have \(v'(p)/e'=v_w(p)\) for every prime \(p\) and \(v'(\nu_a(\alpha))/e'=v_w(\nu_a(\alpha))\). Moreover \(\theta_{x,P}(y')\) is the image of \(\theta_{x,P}(y)\), so \(v'(\theta_{x,P}(y'))/e'=v(\theta_{x,P}(y))/e\). Proposition~\ref{B8.28w} on \(V'\) therefore gives
\[
 v_w(\nu_a(\alpha))+\frac{v(\theta_{x,P}(y))}{e}\ge-\rho(w),
\]
whence \(v(\theta_{x,P}(y))/e\ge-\rho(w)-v_w(\nu_a(\alpha))\ge-(\rho(w)+v_w(\nu))=m(w)\).

(3) By assumption \(\gcd(g_0g_1,Q_Sqd)=1\). If \(w\mid g_0\) or \(w\mid g_1\), every prime \(p\mid Q_Sqd\) is a unit at \(w\), so \(v_w(p)=0\) and \(\rho(w)=0\), which leaves \(m(w)=-v_w(\nu)\).
\end{proof}

\subsection*{Proof of \ref{B8.28xe}}\label{prf:B8.28xe}\phantomsection
\begin{proof}[Proof (outline)]
The exponent hypotheses give \(U,V\ge0\). Since the set of blocks is finite, its exponents have an integer upper bound. Fix a place \(\lambda\mid\ell\). Lemma~\ref{B8.28h} gives
\(\nu_\beta(c_\beta;y)\ge0\) for every integral frame, and hence
\(\nu_c(y)\ge0\). We take \(N_\lambda=0\).

For (a), fix an algebraic integral frame \(y\), choose an integral formal chart, and pass to an algebraically closed valued extension with \(v_\ell(\ell)=1\). The split-place hypothesis supplies the oriented primitive factors \(R_s,\bar R_s\) for the embeddings of this chart. The intermediate estimate on these factors is a family of nonnegative real numbers \(t_\tau\), indexed by all embeddings of \(E\), with
\[
 \sum_\tau t_\tau
 \le Q\sum_s n_s(1+\delta_s/2)
 =fQ(1+\delta/2),
\]
and
\[
 \nu_{\tau,i}(n;y)\le nt_\tau
 \qquad(\tau\in\Psi,\ 1\le i\le P,\ n\ge0).
\]
These numbers may depend on the chart and frame. Their total bound depends on the conjugate Lie rank \(Q\) and the oriented factors' rank-weighted discriminant depths. The same \(t_\tau\) bounds every flag step in the \(\tau\)-row; its coefficient is independent of the step size \(i\).

Multiply the row estimate by the exponent bounds and sum:
\[
 \begin{aligned}
 \sum_{\beta\ \mathrm{big}}\nu_\beta(c_\beta;y)
 &\le\sum_{\tau\in\Psi}t_\tau
           \sum_{i=k-1}^{P}c_{\tau,i}\\
 &\le U\sum_\tau t_\tau
 \le fQU(1+\delta/2)
 \le fmU(1+\delta/2).
 \end{aligned}
\]
Here \(Q=m-j\le m\), and \(1+\delta/2\ge0\).
For the small blocks, Lemma~\ref{B8.28h} and \(\ell\ge2\) give
\[
 \sum_{\beta=(\tau,i)\ \mathrm{small}}\nu_\beta(c_\beta;y)
 \le\frac{\sum_{\beta\ \mathrm{small}}i c_\beta}{\ell-1}
 \le V.
\]
Thus \(0\le\nu_c(y)\le fmU(1+\delta/2)+V\) with the required lower and upper frame quantifiers. The lower endpoint is already an integer, so the displayed width needs no rounding term.

For (b), Lemma~\ref{B8.28xf} gives
\(\nu_\beta(c_\beta;y)\le c_\beta Q/(\ell-1)\) on big blocks at algebraic integral frames. The universal estimate of Lemma~\ref{B8.28h} gives
\(\nu_{\tau,i}(c_{\tau,i};y)\le i c_{\tau,i}/(\ell-1)\) on the small blocks. Since \(\#\Psi=f\),
\[
 \begin{aligned}
 \nu_c(y)
 &\le\frac{Q\sum_{\tau\in\Psi}\sum_{i=k-1}^{P}c_{\tau,i}
                   +\sum_{\beta\ \mathrm{small}}i c_\beta}{\ell-1}\\
 &\le\frac{fQU+V}{\ell-1}
 \le\frac{fmU+V}{\ell-1}.
 \end{aligned}
\]
Together with the lower bound at every integral frame, this proves (b).

The complete proofs are contained in \Lean{GlobalTheory.\allowbreak{}RefSum.\allowbreak{}oscLe\_\allowbreak{}two} and \Lean{GlobalTheory.\allowbreak{}OscSplit.\allowbreak{}oscLe\_\allowbreak{}split}.
\end{proof}

\subsection*{Proof of \ref{B8.28xfa}}\label{prf:B8.28xfa}\phantomsection
\begin{proof}[Proof (outline)]
Reduce the integral embeddings \(\mathcal O_E\to O\) modulo
\(\mathfrak m_O\). The discriminant hypothesis makes these reductions distinct. They exhaust the homomorphisms to the residue field, and Frobenius permutes them. If the kernel of the reduction of \(\tau\) has residue field of cardinality \(p^h\), its Frobenius orbit has length \(h\). Lift this orbit through the bijection of embeddings and their reductions. On the quadratic subfield \(K\), Frobenius fixes the two embeddings in the split case and interchanges them in the inert case. This proves the assertions about \(\Psi\).

For the character assertion, take the uniformisation chart supplied by algebraicity of \(y\). If its torus rank is \(t\) and its abelian quotient is \(B\), its formal maps
\[
 \widehat{\mathbf G}_m^t\xrightarrow{J}G\xrightarrow{q}\widehat B
\]
have a split exact tangent sequence over \(O\). Both maps respect the
\(\mathcal O_E\)-action. For every \(m\), the coordinate algebra of the
\(p^m\)-torsion is finite free over \(O\), and its \(O\)-valued points separate its functions. The polarized Hodge--Tate datum of \(B\) consists of its Tate lattice with its alternating polarization pairing and the Hodge--Tate map to the cotangent lattice, together with their integral torsion data. Its Hodge--Tate columns are differentials of integral formal characters. Combining this datum with the tangent sequence, the torsion property, and the integral action on the torus character lattice gives finitely many integral characters
\(\eta_1,\ldots,\eta_N\) of \(G\).

Let \(M_\tau\subset O^{\mu_\tau}\) be generated by the columns
\[
 \bigl(d\eta_j(u_{\tau,a})\bigr)_{1\le a\le\mu_\tau}.
\]
The unramified eigenspace length calculation gives a positive determinant length \(x_\tau\) satisfying
\[
 |p|^{R_\tau}\le x_\tau.
\]
Equivalently, the minimum valuation of the maximal minors of this finite matrix is at most \(R_\tau\). The inputs to that calculation are the split exact tangent sequence, the polarized Hodge--Tate datum, \(p\ne2\), the unit discriminant, and the Frobenius congruences above. Its sequence of multiplicities is
\(\mu_{\overline{\tau_{m+1}}}\), the conjugate Lie multiplicities of the whole formal group.

Choose a maximal minor with valuation at most \(R_\tau\) and expand it along the first \(i\) rows. Every complementary minor is integral. The ultrametric inequality gives an \(i\)-column minor on those rows with valuation at most \(R_\tau\). Use the corresponding characters as
\(\chi_1,\ldots,\chi_i\). In multiplicative notation the resulting determinant \(D\) satisfies
\[
 |D|\ge x_\tau\ge|p|^{R_\tau}.
\]
Taking logarithms and dividing by \(\log|p|<0\) gives \(v(D)\le R_\tau\).

The complete orbit construction, character construction, and passage to a flag minor are contained in \Lean{GlobalTheory.\allowbreak{}FrobOrbit.\allowbreak{}frobOnType\_\allowbreak{}and\_\allowbreak{}orbit}, \Lean{GlobalTheory.\allowbreak{}ChF.\allowbreak{}unramified\_\allowbreak{}companion\_\allowbreak{}chart\_\allowbreak{}characters}, \Lean{GlobalTheory.\allowbreak{}CharMinor.\allowbreak{}colength\_\allowbreak{}bound}, and \Lean{GlobalTheory.\allowbreak{}RefSum.\allowbreak{}hasChars\_\allowbreak{}frame\_\allowbreak{}of\_\allowbreak{}hasAF}.
\end{proof}

\subsection*{Proof of \ref{B8.28xf}}\label{prf:B8.28xf}\phantomsection
\begin{proof}
Fix an algebraic integral frame \(y\), a chart \(G\), and a valued extension as in Lemma~\ref{B8.28xfa}. For \(i\) copies of \(G\), let \(D_{u_a}^{(b)}\) be the invariant derivation in the
\(u_a=u_{\tau,a}\) direction on copy \(b\). Apply the determinant operator
\[
 \mathcal D=\det\bigl(D_{u_a}^{(b)}\bigr)_{a,b=1}^i
\]
to the integral series
\(F=\prod_{b=1}^i\chi_b(X^{(b)})\).
A formal character satisfies \(D_u\chi=d\chi(u)\chi\); hence
\[
 \mathcal D^nF(0)=\det\bigl(d\chi_b(u_a)\bigr)^n.
\]
Taking the infimum that defines the chart depth, and then the infimum over charts, gives
\[
 \nu_{\tau,i}(n;y)\le nR_\tau.
\]

If \(p=\ell\) splits in \(K\), every \(\tau_s\) lies in \(\Psi\), since
\(\tau\in\Psi\). Thus
\[
 R_\tau=Q\sum_{m\ge0}\ell^{-m-1}=\frac{Q}{\ell-1}.
\]
If \(\ell\) is inert in \(K\), membership in \(\Psi\) alternates. The conjugate multiplicity in the summand is \(P\) for even \(m\) and \(Q\) for odd \(m\), so
\[
 R_\tau
 =\sum_{a\ge0}\left(P\ell^{-2a-1}+Q\ell^{-2a-2}\right)
 =\frac{P\ell+Q}{\ell^2-1}
 =\frac{Q}{\ell-1}+\frac{(P-Q)\ell}{\ell^2-1}.
\]

The character calculation first controls the determinant colength on the full eigenspace, and then passes to each flag minor. Thus its loss in the split case is \(Q/(\ell-1)\) for every block size \(i\); taking an \(i\)-dimensional flag step adds no factor \(i\). This is where the conjugate multiplicity \(Q\) enters in place of the block size in Lemma~\ref{B8.28h}. That lemma supplies nonnegativity for every integral frame and depth zero when \(n=0\).
\end{proof}

\subsection*{Proof of \ref{B8.28xga}}\label{prf:B8.28xga}\phantomsection
\begin{proof}[Proof (outline)]
The unramified completed order \(\mathcal O_E\otimes\Z_p\) has a factor belonging to the prime specified by the orbit of \(\tau\). Its residue field has cardinality \(p^h\). The formal factor associated with this component carries the Teichm\"uller roots of unity of order \(p^h-1\). Their eigencharacters on its Lie algebra are
\(\zeta\mapsto\zeta^{p^s}\), with multiplicities \(\mu_{\tau_s}\).
Choose integral coordinates indexed by these characters. Their logarithms give the displayed formula for \([\zeta]\).

The integral inclusion and projection of this formal factor give the Lie map \(\Pi\). On the \(\tau_s\)-eigenframe, the projected vectors form a saturated frame of type \(s\), meaning that their entries are integral and one maximal minor is a unit. The formal inclusion and projection identify the two determinant depths for every exponent. This construction is made on the formal group of the integral frame, with its complete \(\mathcal O_E\)-action; the multiplicities and the depth identity refer to that formal group, including its torus part.

The complete typed-factor construction is contained in \Lean{GlobalTheory.\allowbreak{}T.\allowbreak{}formalGroup\_\allowbreak{}rootsOfUnity\_\allowbreak{}typed\_\allowbreak{}lieRanks}.
\end{proof}

\subsection*{Proof of \ref{B8.28xgb}}\label{prf:B8.28xgb}\phantomsection
\begin{proof}[Proof (outline)]
Apply Lemma~\ref{B8.28xga} to \(\tau\). Since Frobenius alternates the two CM types, \(h\) is even and
\[
 \tau_s\in\Psi\quad\Longleftrightarrow\quad
 \widetilde s\text{ is even}.
\]
The projected frame \(u'\) is saturated and supported in type zero, and every odd type has exactly \(Q\) coordinates.

The formal-group calculation used here has the following inputs: an algebraically closed field of characteristic zero with \(v(\ell)=1\), an integral formal group with the displayed
\((\ell^h-1)\)-st-root action, even \(h\), at most \(Q\) coordinates in every odd type, and a saturated frame of size \(i\) supported in an even type.
For \(\ell\ge3\) and \(n\ge1\), it gives a real correction \(c\) such that
\[
 c\le Q\left(\frac32+\frac12\log_\ell n\right),
\]
and
\[
 \operatorname{depth}_{G'}(\det D_{u'}^{\,n})
 \ge(i-c)\frac{n\ell}{\ell^2-1}
       -i^2(\log_\ell n+1).
\]
The coefficient of \(i-c\) is positive. Substitute the upper bound for \(c\) and use the depth identity in Lemma~\ref{B8.28xga}. Every input in this calculation concerns the typed formal factor of the given integral frame, so the conclusion applies to each chart and each trait.

The complete lower-depth calculation is contained in \Lean{GlobalTheory.\allowbreak{}T.\allowbreak{}detDepth\_\allowbreak{}lower\_\allowbreak{}bound\_\allowbreak{}saturated\_\allowbreak{}even\_\allowbreak{}type}, and its application here is contained in \Lean{GlobalTheory.\allowbreak{}RefSum.\allowbreak{}facBetween\_\allowbreak{}fine}.
\end{proof}

\subsection*{Proof of \ref{B8.28xg}}\label{prf:B8.28xg}\phantomsection
\begin{proof}
Set
\[
 a=\left(i-Q\left(\frac32+\frac L2\right)\right)
           \frac{n\ell}{\ell^2-1}-i^2(L+1),
 \qquad
 b=n\left(\frac Q{\ell-1}
              +\frac{(P-Q)\ell}{\ell^2-1}\right).
\]
These numbers are fixed before choosing a trait or a frame.
Lemma~\ref{B8.28xgb}, applied to every chart and valued extension, gives
\(a\le\nu_{\tau,i}(n;y)\) for every integral frame.
Lemma~\ref{B8.28xf} gives \(\nu_{\tau,i}(n;y)\le b\) for every algebraic integral frame.

The typed-factor construction in the lower bound includes the torus part of each formal group. The character construction in the upper bound starts from the split exact tangent sequence with that torus. Both bounds use the prescribed Lie multiplicities \(P,Q\) of the complete formal group. Thus this single interval serves all traits, including those with nonzero torus rank.

Since \(L\ge0\) and \(i\le P\),
\[
 b-a=\frac{nQ}{\ell-1}
   +\frac{n\ell}{\ell^2-1}
          \left(P-i+\frac Q2+\frac{QL}{2}\right)
   +i^2(L+1)\ge0.
\]
Put \(s=\ell/(\ell^2-1)\) and \(t=1/(\ell-1)\); then \(0<s\le t\). Direct subtraction gives
\[
 W-(b-a)
   =nQ\left(\frac t2+
             \left(\frac32+\frac L2\right)(t-s)\right)\ge0.
\]
This proves \(a\le b\) and \(b-a\le W\).
The first term of \(W\) contains the real number \(P-i-Q\), which is negative when \(i>P-Q\). It remains negative in this calculation; the displayed identity proves the bound for the complete sum without replacing that term by zero.
\end{proof}

\subsection*{Proof of \ref{B8.28xha}}\label{prf:B8.28xha}\phantomsection
\begin{proof}
The diagonal monomial \(\prod_a Z_{aa}^{n}\) has coefficient one, so the support is finite and nonempty. If \(m\) is a nonzero integer, the normalization \(v(p)=1\) gives \(v(m)=v_p(m)\): integers have nonnegative valuation, and B\'ezout's identity gives valuation zero for an integer prime to \(p\). Apply this to the nonzero integers \(c_\alpha\alpha!\) and take the minimum over the fixed support.

For the depth bound, invert a unit maximal minor of the frame to obtain integral linear forms \(\lambda_a\) with
\(\lambda_a(u_b)=\delta_{ab}\). Choose a supported multi-index \(\alpha\) attaining \(\kappa_{p,i,n}\), and use
\[
 \prod_{a,b}\lambda_a(X^{(b)})^{\alpha_{ab}},
\]
of degree \(in\), as an integral test function on \(i\) copies of \(G\).
At the origin, the order-\(in\) determinant operator on this homogeneous test function is its constant-coefficient principal part; terms with fewer derivatives vanish. Its value is \(c_\alpha\alpha!\).
The infimum defining depth is therefore at most \(\kappa_{p,i,n}\).
\end{proof}

\subsection*{Proof of \ref{B8.28xhb}}\label{prf:B8.28xhb}\phantomsection
\begin{proof}[Proof (outline)]
Choose integral curves \(\gamma_1,\ldots,\gamma_d\) through the identity with
\[
 \log_G\gamma_b(T)=\sum_{m\ge0}L_{b,m}T^{p^m}.
\]
The adapted typical-curve construction gives an injection
\(r:\{1,\ldots,i\}\hookrightarrow\{1,\ldots,d\}\) and an integral matrix
\(g\) of unit determinant such that the tangent matrix
\(T_{be}=L_{b,0,e}\) is integral with unit determinant,
\[
 L_{r(a),0}=\sum_{a'}g_{aa'}u_{a'},
 \qquad
 L_{r(a),1}=0\quad(a>j).
\]
Replacing \(u\) by these selected tangents preserves determinant depth, since its determinant operator is multiplied by the unit \(\det g\). The bound \(j\) is the dimension bound on the next type in the root-of-unity action.

The typical-curve structure constants \(s_{q,e,b}\) are integral and satisfy
\[
 pL_{b,m+1}=\sum_{q=0}^m\sum_e
          s_{m-q,e,b}^{p^q}L_{e,q}.
\]
Induction on \(m\) gives \(v(L_{b,m,e})\ge-m\).
If \(L_{b,1}=0\), the equation for \(m=0\) and invertibility of the tangent matrix give \(s_{0,e,b}=0\) for all \(e\). In the remaining equations the term \(q=m\) vanishes, so
\[
 v(L_{b,m,e})\ge1-m\qquad(m\ge1)
\]
for these curves. Call such a curve good.

Let
\[
 \Phi(X)=\gamma_1(X_1)+_G\cdots+_G\gamma_d(X_d).
\]
It and its inverse are integral, because its tangent matrix is invertible over \(O\). In these coordinates the formal map with logarithm
\(\sum_a L_{r(a),0}Y_a\) has coordinates \(z_b(Y)\), with zero constant term. Multiplying the logarithm equation by \(T^{-1}\) yields
\[
 z_b+\sum_{a=1}^d\sum_{m\ge1}c_{a,m,b}z_a^{p^m}
     =\sum_{q=1}^i\delta_{r(q),b}Y_q,
 \qquad
 c_{a,m,b}=\sum_e L_{a,m,e}(T^{-1})_{eb}.
\]
Thus \(v(c_{a,m,b})\ge-m\), with the stronger bound \(1-m\) for good \(a\).

Assign the radii
\[
 R_b=
 \begin{cases}
 2/[p(p-1)],&b\text{ good},\\
 2/(p-1),&b\text{ otherwise},
 \end{cases}
 \qquad
 \rho_q=R_{r(q)}.
\]
A series \(f=\sum_\alpha f_\alpha Y^\alpha\) has weighted order at least \(R\) when
\[
 v(f_\alpha)+\sum_q\rho_q\alpha_q\ge R
\]
for every \(\alpha\). The right side of the logarithm equation has weighted order at least \(R_b\). For \(m\ge1\),
\[
 -m+p^m\frac2{p-1}\ge\frac2{p-1},
 \qquad
 1-m+p^m\frac2{p(p-1)}\ge\frac2{p-1}.
\]
Induction on total degree in the logarithm equation proves that \(z_b\) has weighted order at least \(R_b\): every nonlinear coefficient of degree \(N\) uses coefficients of \(z_a\) of degree below \(N\). In particular, each \(z_b\) has weighted order at least zero.

An integral series on \(i\) copies of \(G\), pulled back through these coordinates, has its coefficient of multi-index
\(\alpha=(\alpha_{bq})\) of valuation at least
\[
 -\sum_{b,q}\rho_q\alpha_{bq}.
\]
Every monomial in \(H_{i,n}\) has degree \(n\) in each frame index \(q\), so for its supported indices this lower bound is
\(-n\sum_q\rho_q\). Applying \(H_{i,n}(\partial)\) at the origin and using Lemma~\ref{B8.28xha} gives
\[
 \operatorname{depth}_{G}(\det D_u^{\,n})
    \ge\kappa_{p,i,n}-n\sum_q\rho_q.
\]
At most \(j\) of the selected curves are bad, whence
\[
 n\sum_q\rho_q\le
       \frac{2in}{p(p-1)}+\frac{2jn}{p-1}.
\]
The claimed inequality follows.

The complete adapted-curve and depth calculation is contained in \Lean{GlobalTheory.\allowbreak{}FGDepth.\allowbreak{}InertBound.\allowbreak{}content\_\allowbreak{}le\_\allowbreak{}detDepth\_\allowbreak{}add}.
\end{proof}

\subsection*{Proof of \ref{B8.28xh}}\label{prf:B8.28xh}\phantomsection
\begin{proof}
For \(n=0\), take \(a=b=0\) by Lemma~\ref{B8.28h}.
For any \(n\ge0\), let \(\kappa=\kappa_{\ell,i,n}\) be the number in Lemma~\ref{B8.28xha} and put
\[
 W_0=\frac{2in}{\ell(\ell-1)}+\frac{2Qn}{\ell-1},
 \qquad a=\kappa-W_0,\qquad b=\kappa.
\]
Fix an arbitrary integral frame, chart and valued extension over \(\lambda\). In the typed factor from Lemma~\ref{B8.28xga}, the block frame has type zero. Since \(\tau\in\Psi\) and \(\ell\) is inert, the next Frobenius type lies outside \(\Psi\), so its Lie multiplicity is \(Q\).
Apply Lemma~\ref{B8.28xhb} with \(j=Q\) and transfer its lower bound by the depth identity. It gives \(a\le\nu_{\tau,i}(n;y)\) after taking the infimum over all charts and valued extensions.

For the upper bound, choose an integral chart and a valued extension over \(\lambda\). The saturated-frame bound in Lemma~\ref{B8.28xha} gives
\[
 \nu_{\tau,i}(n;y)\le
 \operatorname{depth}_{G}(\det D_u^{\,n})
 \le\operatorname{cont}_v(H_{i,n})=\kappa=b.
\]
Thus the same constants work for every trait. Since \(W_0\ge0\), they satisfy \(a\le b\) and \(b-a=W_0\).
The next-type rank \(Q\) bounds the number of selected curves with radius \(2/(\ell-1)\); all remaining curves use radius
\(2/[\ell(\ell-1)]\). This gives the additional divisor \(\ell\) on the term involving \(i\).

Under the three size conditions in the statement, \(\ell\ge2\) gives
\[
 2i+2Q\ell
 \le2(k+m)+2m\ell
 \le4m\ell+2m
 \le5m\ell.
\]
Multiply by \(n/[\ell(\ell-1)]\) to obtain
\(b-a\le5mn/(\ell-1)\).
\end{proof}

\subsection*{Proof of \ref{B8.28x}}\label{prf:B8.28x}\phantomsection
\begin{proof}
Write \(\Psi\) for the CM type and \(2f=[E:\Q]\). Then
\[
 P=k+j-1,\qquad Q=m-j,\qquad r=P+Q=k+m-1,
 \qquad A=fmU,\qquad h_\ell=\frac{\log J}{\log\ell}.
\]
There are \(f\) embeddings in \(\Psi\), so the row bounds give
\[
 \sum_{\beta\ \mathrm{big}}c_\beta\le fU.
\]
Also \(Q\le m\). Since \(\ell\le J\) and \(\ell\ge2\), we have \(J\ge2\) and \(h_\ell\ge0\).

(1) Fix \(C,O,\tau_O\) as in Lemma~\ref{B8.28xe}(a). Let \(\eta\in\mathcal O_E\) satisfy \(\eta^2=-d\), and choose \(\tau_0\in\Psi\). The host datum makes \(\ell\) split in \(K\): for odd \(\ell\mid Q_S\) the host requires \((-d/\ell)=1\), and for \(\ell=2\) the congruence \(d\equiv7\pmod 8\) gives \(-d\equiv1\pmod 8\). This gives a root \(\alpha\in\Z_\ell^\times\) of \(X^2+d\). Choose its sign so that its image in \(O\) is \(\tau_O(\tau_0)(\eta)\). Put \(R=\Z_\ell\otimes_\Z\mathcal O_E\). The element
\[
 e=\frac{\alpha+\eta}{2\alpha}\in R
\]
is integral. For odd \(\ell\), this follows because \(2\alpha\) is a unit. For \(\ell=2\), use \(w_0=(1+\eta)/2\in\mathcal O_E\) and write
\[
 e=\frac{(\alpha-1)/2+w_0}{\alpha}.
\]
The root \(\alpha\) is odd, so this expression belongs to \(R\). Direct calculation gives
\[
 e^2=e,\qquad \bar e=1-e,\qquad \eta e=\alpha e.
\]
On each primitive factor, \(e\) is either zero or one. Orient each conjugate pair so that \(e=1\) on \(R_s\). An embedding through \(R_s\) sends \(\eta\) to \(\alpha\); an embedding through \(\bar R_s\) sends it to \(-\alpha\). The compatibility condition imposed on \(\tau_O\) identifies these two classes with \(\tau_O(\Psi)\) and \(\tau_O(\Psi^c)\).

There are \(2f\) homomorphisms \(\mathcal O_E\to O\): they are the restrictions of the \(2f\) embeddings \(E\to C\). Hence the injection \(\tau_O\) is a bijection, and every embedding has the required orientation. Conjugation identifies the trace pairings on paired factors. Reorienting the factors therefore preserves their ranks and discriminant valuations. With
\[
 \delta_s=
 \frac{v_\ell(\operatorname{disc}(R_s/\Z_\ell))}{n_s},
\]
the hypotheses give
\[
 \sum_s n_s=f,\qquad
 \sum_s n_s\delta_s=f\delta_\ell.
\]
Since \(\delta_\ell\ge0\), Lemma~\ref{B8.28xe}(a) gives the required integer lower anchor and the width
\[
 fmU(1+\delta_\ell/2)+V=A(1+\delta_\ell/2)+V.
\]
\leanref{UpperTheory.\allowbreak{}splitPlaces\_\allowbreak{}discriminantDepth\_\allowbreak{}of\_\allowbreak{}splitAt}{Development.\allowbreak{}AutomorphicArithmetic.\allowbreak{}CMTraceForms.\allowbreak{}SplitPlacesDiscriminantDepthOfSplitAt}{59--150}
\leanref{GlobalTheory.\allowbreak{}CharDepth.\allowbreak{}twoPlaces\_\allowbreak{}host}{Development.\allowbreak{}AbelianVarietyTheory.\allowbreak{}Proofs.\allowbreak{}JetHarmonics.\allowbreak{}HostFieldTwoConjugatePlacesOverTwo}{16--27}

(2) A big block \(\beta=(\tau,i)\) has \(\tau\in\Psi\) and \(1\le i\le P\). The split case of Lemma~\ref{B8.28xf}, with exponent \(c_\beta\), gives
\[
 \nu_\beta(c_\beta;y)\le\frac{Qc_\beta}{\ell-1}
\]
at every algebraic integral frame. Lemma~\ref{B8.28h} bounds each small block by \(ic_\beta/(\ell-1)\), and gives nonnegativity at every integral frame. Consequently
\[
 0\le\nu_c(y)\le\frac{fQU+V}{\ell-1}
 \le\frac{A+V}{\ell-1}
\]
at algebraic integral frames, with lower bound zero at every integral frame, as in Lemma~\ref{B8.28xe}(b).

(3) Fix \(\lambda\mid\ell\). For every big block with \(c_\beta\ge1\), take the uniform endpoints from Lemma~\ref{B8.28xg}. Set
\[
 L_\beta=\frac{\log c_\beta}{\log\ell},\qquad
 t_\beta=r-2Q-i.
\]
The hypotheses give \(0\le L_\beta\le h_\ell\), \(i\le r\), and
\[
 t_\beta=k+m-1-2Q-i\le m.
\]
If \(t_\beta\le0\), then \(t_\beta\ell/(\ell^2-1)\le0\). If \(t_\beta>0\), then
\[
 \frac{t_\beta\ell}{\ell^2-1}\le\frac{m}{\ell-1}.
\]
In either case the width bound of Lemma~\ref{B8.28xg} yields
\[
 B_\beta-A_\beta
 \le\frac{mc_\beta(4+h_\ell/2)}{\ell-1}
       +r^2(h_\ell+1).
\]
For each \(\tau\in\Psi\), the possible big sizes are \(k-1,\ldots,P\), whose number is \(j+1\le m+1\). Thus there are at most \(f(m+1)\) big blocks. Summing the preceding inequality and using \(\sum_{\mathrm{big}}c_\beta\le fU\) gives
\[
 \sum_{\substack{\beta\ \mathrm{big}\\c_\beta\ge1}}
       (B_\beta-A_\beta)
 \le\frac{A(4+h_\ell/2)}{\ell-1}
       +f(m+1)r^2(h_\ell+1).
\]
\leanref{GlobalTheory.\allowbreak{}RefSum.\allowbreak{}OscInert.\allowbreak{}sum\_\allowbreak{}bound\_\allowbreak{}le}{Development.\allowbreak{}AbelianVarietyTheory.\allowbreak{}Proofs.\allowbreak{}DivisorProofs.\allowbreak{}DepthOscillationLeAtSmallOddInertPrime}{112--175}

(4) For a big block of size \(i\),
\[
 i+1\le P+1=k+j\le k+m,\qquad Q\le m,\qquad k\le m\ell.
\]
Lemma~\ref{B8.28xh} supplies uniform endpoints whose width is at most
\[
 W_\beta=\frac{2ic_\beta}{\ell(\ell-1)}
             +\frac{2Qc_\beta}{\ell-1}.
\]
The displayed size bounds and \(\ell\ge2\) give
\[
 2i+2Q\ell\le2(k+m)+2m\ell
             \le4m\ell+2m\le5m\ell.
\]
Multiplying by \(c_\beta/(\ell(\ell-1))\ge0\) proves
\[
 W_\beta\le\frac{5mc_\beta}{\ell-1}.
\]
For \(c_\beta=0\), Lemma~\ref{B8.28h} gives endpoints zero.

The connection with theta differentials uses Propositions~\ref{B8.28u}, \ref{B8.28v} and \ref{B8.28w}. Given a theta weight \((a,N)\) with gaps \(c\) and
\(N\ge\sum_{c_{\tau,i}\ne0}2i\), choose \(P_{\rm col}\) as in Lemma~\ref{B8.28k}(b). Proposition~\ref{B8.28u} and Lemma~\ref{B8.28o} identify \(\nu_c(y)\) with \(\delta_t(P_{\rm col};y)\) for every integer \(t\ge0\). In (1) and (2), use the total-depth intervals already obtained. In (3) and (4), add the small-block bounds and apply Lemma~\ref{B8.28i}. Proposition~\ref{B8.28v} supplies the error estimate used in Lemma~\ref{B8.28p}(2) to obtain a global harmonic differential with the resulting uniform interval. For the split uniformisations and half-level factors of Proposition~\ref{B8.28w}, that proposition transfers the chart lower bound to translated theta derivatives. For a trait above \(\ell\), dividing an inequality in \(v/e\) by \(v(\ell)/e>0\) expresses it in the normalization \(v_\ell(\ell)=1\).
\end{proof}

\subsection*{Proof of \ref{B8.28r}}\label{prf:B8.28r}\phantomsection
\begin{proof}
(1) Recall the construction of \(a\) in Lemma~\ref{B8.25g}. For \(\tau\) in the CM type, \(s_\tau=N/2+1-u_{1,\tau,P-1}\) gives \(a_{\tau,i}=u_{1,\tau,i-1}-u_{1,\tau,P-1}+1\) for \(1\le i\le P\) and \(a_{\tau,P+1}=0\). Since \(|u_{1,\tau,s}|\le B_0\), we have \(a_{\tau,i}\le 2B_0+1=J\). Telescoping from \(k-1\) to \(P\) and using \(u_{1,\tau,k-2}=0\) gives
\[
 \sum_{i=k-1}^P c_{\tau,i}=a_{\tau,k-1}=u_{1,\tau,k-2}-u_{1,\tau,P-1}+1=-u_{1,\tau,P-1}+1\le B_0+1=U.
\]
For \(1\le i\le k-2\) we have \(c_{\tau,i}=u_{1,\tau,i-1}-u_{1,\tau,i}\), so summation by parts and \(u_{1,\tau,k-2}=0\) give
\[
 \sum_{i=1}^{k-2}i\,c_{\tau,i}=\sum_{0\le s<k-2}u_{1,\tau,s}-(k-2)u_{1,\tau,k-2}\le\frac{k(S+2r)}{T}.
\]
At the conjugate embedding \(\bar\tau\), reversing and negating the second block and subtracting \(N/2\) gives \(a_{\bar\tau,i}=u_{1,\tau,P-1}-u_{1,\tau,P+Q-i}-N-1\) for \(1\le i\le Q\) and \(a_{\bar\tau,Q+1}=0\). Monotonicity of \(u_{1,\tau}\) and \(u_{1,\tau,P-1}-u_{1,\tau,P}\ge N+1\) make \(a_{\bar\tau,i}\) nonincreasing and nonnegative, with \(a_{\bar\tau,i}\le 2B_0\le J\) and \(\sum_{i=1}^Q i\,c_{\bar\tau,i}=\sum_{i=1}^Q a_{\bar\tau,i}\le 2QB_0=2q_{\rm sig}B_0\). Since \(c_{\sigma,i}\le a_{\sigma,i}\le J\) everywhere, summing the small blocks over the \(f\) embeddings \(\tau\) in the CM type and their \(f\) conjugates \(\bar\tau\) gives \(\sum_{\rm small}i\,c_{\sigma,i}\le V\). On every big block \(i\le P\le r\), so
\[
 D=\sum_{\rm big}i\,c_{\tau,i}+\sum_{\rm small}i\,c_{\sigma,i}\le r\sum_{\tau\text{ in CM type}}\sum_{i=k-1}^P c_{\tau,i}+V\le frU+V.
\]

(2) By hypothesis \(2\mid Q_S\), and by \ref{B2.46} we have \(E\subset\Q(\mu_{Q_Sqd})\), so every prime divisor of \(\operatorname{disc}E\) divides \(Q_Sqd\) and primes outside \(Q_S\) and \(\{q,d\}\) are odd and unramified in \(E\). Fix a prime \(\ell\mid Q_S\); by hypothesis \(\ell\) splits \(K=\Q(\sqrt{-d})\) over \(\Z_\ell\). Choose \(\alpha\in\Z_\ell\) with \(\alpha^2=-d\). Since \(\ell\) splits in \(\mathcal O_K\) of discriminant \(-d\), we have \(\alpha\in\Z_\ell^\times\), and when \(\ell=2\) we have \(\alpha\equiv1\pmod2\) and \(w_0=(1+\sqrt{-d})/2\in\mathcal O_K\), so \(e=(\alpha+\sqrt{-d})/(2\alpha)=((\alpha-1)/2+w_0)/\alpha\) lies in \(\Z_\ell\otimes_{\Z}\mathcal O_K\subset\Z_\ell\otimes_{\Z}\mathcal O_E\) and satisfies \(e^2=e\) and \(e+\bar e=1\). Decomposing \(e=\sum_{s=1}^u g_s\) into primitive orthogonal idempotents pairs each factor \(R_s=g_s(\Z_\ell\otimes_{\Z}\mathcal O_E)\) with its conjugate \(\bar R_s=\bar g_s(\Z_\ell\otimes_{\Z}\mathcal O_E)\) (with \(\bar g_s\ne g_s\) since \(e\bar e=0\)). Because conjugation is an isomorphism \(R_s\simeq\bar R_s\) of \(\Z_\ell\)-orders and \(\Z_\ell\otimes_{\Z}\mathcal O_E\cong\prod_{s=1}^u(R_s\times\bar R_s)\), we obtain \(2\sum_{s=1}^u n_s=[E:\Q]=2f\) and \(2\sum_{s=1}^u v_\ell(\operatorname{disc}(R_s/\Z_\ell))=v_\ell(D_E)=2f\delta_\ell\).

Now fix a prime \(\ell\). If \(\ell>J\), then \(\ell>\max_\beta c_\beta\) by (1), and Lemma~\ref{B8.28i} gives oscillation at most \(0=e_\ell\). Suppose \(\ell\le J\), so \(J\ge\ell\ge2\) and \(h_\ell=\log J/\log\ell\ge0\). If \(\ell\mid Q_S\), Proposition~\ref{B8.28x}(1) gives oscillation at most \(A(1+\delta_\ell/2)+V=e_\ell\). If \(\ell\nmid Q_S\) and \(\ell\in\{q,d\}\), Lemma~\ref{B8.28i} and \(D\le frU+V\) give oscillation at most \(D/(\ell-1)\le(frU+V)/(\ell-1)=e_\ell\).

Assume finally that \(\ell\le J\) with \(\ell\nmid Q_S\) and \(\ell\notin\{q,d\}\). Then \(\ell\nmid 2\operatorname{disc}E\) and \(\ell\ne d\), so \(\ell\) either splits or is inert in \(K\). On every small block \(\beta=(\sigma,i)\), Lemma~\ref{B8.28h} gives \(0\le\nu_\beta(c_\beta;y)\le i\,c_\beta/(\ell-1)\), and (1) bounds the sum of these upper bounds over small blocks by \(V/(\ell-1)\). If \(\ell\) splits in \(K\), Proposition~\ref{B8.28x}(2) and \(Q=m-j\le m\), summed with (1) and \(A_\lambda=0\), bound \(\nu_c(y)\) at every algebraic integral frame over \(\lambda\mid\ell\) by
\[
 \frac{Q\sum_{\rm big}c_\beta+\sum_{\rm small}i\,c_\beta}{\ell-1}\le\frac{fmU+V}{\ell-1}=\frac{A+V}{\ell-1},
\]
which is at most \(e_\ell\) whether \(m\ell<k\) or \(m\ell\ge k\) since \(A\ge0\) and \(h_\ell\ge0\). If \(\ell\) is inert in \(K\) and \(m\ell<k\), apply Lemma~\ref{B8.28i} with \((A_\beta,B_\beta)\) from Proposition~\ref{B8.28x}(3) on big blocks with \(c_\beta\ge1\), \((0,0)\) on big blocks with \(c_\beta=0\) by Lemma~\ref{B8.28h}, and \((0,i\,c_\beta/(\ell-1))\) on small blocks; this gives oscillation at most \(\{A(4+h_\ell/2)+V\}/(\ell-1)+f(m+1)r^2(h_\ell+1)+1=e_\ell\). If \(\ell\) is inert in \(K\) and \(m\ell\ge k\), apply Lemma~\ref{B8.28i} with Proposition~\ref{B8.28x}(4) on big blocks and \((0,i\,c_\beta/(\ell-1))\) on small blocks. By (1), \(\sum_{\rm big}5mc_\beta\le 5fmU\le 6A\), so \(c\) has oscillation at most \((6A+V)/(\ell-1)+1=e_\ell\).
\end{proof}

\subsection*{Proof of \ref{B8.28ya}}\label{prf:B8.28ya}\phantomsection
\begin{proof}
Every factor is a determinant of distinct variables and has integer coefficients. Each determinant is a nonzero homogeneous polynomial of degree equal to its size. Their product is nonzero in the polynomial ring over \(\Z\).

The determinants containing row \((\tau,i)\) are exactly those of sizes \(j\ge i\). Thus its row degree is
\[
 \sum_{j=i}^{\mu_\tau}c_{\tau,j}
 =\sum_{j=i}^{\mu_\tau}(a_{\tau,j}-a_{\tau,j+1})
 =a_{\tau,i}.
\]
Summing over the rows gives
\[
 \sum_\tau\sum_{j=1}^{\mu_\tau}j\,c_{\tau,j}
 =\sum_\tau\sum_{i=1}^{\mu_\tau}a_{\tau,i}=D.
\]
The polynomial is a product of the leading minors required in Definition~\ref{B8.28g}(c). Adding a multiple of row \(i'\) to row \(i\), with \(i'<i\), preserves each such minor: a smaller minor omits row \(i\), and a larger one contains both rows.

Conjugate embeddings use disjoint column sets. For every \(\tau,i,i',b\), at least one of the variables \(X_{\tau,i,b}\), \(X_{\bar\tau,i',b}\) is absent from \(P_S\). Hence
\[
 \frac{\partial^2P_S}
 {\partial X_{\tau,i,b}\,\partial X_{\bar\tau,i',b}}=0.
\]
Summing over \(b\) proves all the required harmonic identities. Together with the row degrees and integral coefficients, this proves membership in \(\operatorname{harm}(a,N)\).

Determinants, products and nonnegative integer powers commute with ring homomorphisms. This proves the assertion about change of coefficients.
\end{proof}

\subsection*{Proof of \ref{B8.28yb}}\label{prf:B8.28yb}\phantomsection
\begin{proof}[Proof (outline)]
Complete splitting supplies integral maps
\(\sigma_\tau:\mathcal O_E\to\mathcal O_F\), integral eigenvectors
\(e_{\tau,i}\) giving the chosen Lie directions, an integral left inverse to this system of vectors, and elements
\(\epsilon_\tau\in\mathcal O_E\) satisfying
\[
 d[a]e_{\tau,i}=\sigma_\tau(a)e_{\tau,i},
 \qquad
 \sigma_{\tau'}(\epsilon_\tau)
       \equiv\delta_{\tau,\tau'}\pmod{g\mathcal O_F}.
\]
For every \(a\in\mathcal O_E\), a positive valuation of
\(\sigma_\tau(a)\) implies \(\sigma_\tau(a)\in g\mathcal O_F\), while a nonpositive valuation makes it a unit.
Coprimality of \(\pi,\bar\pi\) and \(v_F(g)>0\) imply that exactly one of
\(\sigma_\tau(\bar\pi)\), \(\sigma_\tau(\pi)\) has positive valuation.
Thus \(S\) contains one embedding from each conjugate pair.
\leanref{GlobalTheory.\allowbreak{}T.\allowbreak{}exists\_\allowbreak{}integral\_\allowbreak{}types\_\allowbreak{}and\_\allowbreak{}directions'}{Development.\allowbreak{}AbelianVarietyTheory.\allowbreak{}Proofs.\allowbreak{}JetHarmonics.\allowbreak{}ExistsIntegralTypesAndDirectionsStatement}{14--16}

The chart character theorem supplies finitely many multiplicative formal series \(u_j\in\mathcal O_F[[X]]\), meaning
\[
 u_j(0)=1,\qquad
 u_j(\widehat G(X,Y))=u_j(X)u_j(Y),
\]
such that
\[
 c\,\mathcal O_F^d\subset
 \operatorname{span}_{\mathcal O_F}\{d\log u_j\}
 \quad\text{whenever }(g-1)v_F(c)\ge v_F(g).
\]
Here \(d\) is the chart dimension and \(d\log u_j\) is the differential at the identity. This theorem applies to the formal group of the split semiabelian extension, including its torus. It follows by combining the equivariant extension of the torus and the formal group of the abelian part with the polarised Hodge--Tate character datum of that abelian part.
\leanref{GlobalTheory.\allowbreak{}ChF.\allowbreak{}hodgeTate\_\allowbreak{}inclusion\_\allowbreak{}chart}{Development.\allowbreak{}AbelianVarietyTheory.\allowbreak{}Proofs.\allowbreak{}UnramifiedDepth.\allowbreak{}HodgeTateInclusionChartCharacters}{49--93}

Choose \(c\in\mathcal O_F\) with
\((g-1)v_F(c)=v_F(g)\); a \((g-1)\)-st root of \(g\) has this property. Pairing the inclusion above with the integral left inverses gives the same inclusion in each type of Lie directions. Pulling the characters back by \([\epsilon_\tau]\) separates those types modulo \(g\).
The simultaneous column selection for these characters yields integral directional-value arrays \(Y_{\tau,i,b}\), with columns assigned as in \(P_S\), such that their leading minors \(Y_{\tau,j}\) satisfy
\[
 v_F(\det Y_{\tau,j})\le jv_F(c)
 \qquad(1\le j\le\mu_\tau,\ \tau\in S).
\]
The construction uses an integral left inverse in each type, the congruences for \(\epsilon_\tau\), injectivity of the column assignment, and
\(\sigma_\tau(\bar\pi)\in g\mathcal O_F\) for \(\tau\in S\).
Multiplicative-series comparison converts its columns to the characters used by finite-order translates at \(\pi\)-torsion. The unit types use the other column set. The determinant contribution to the value is bounded by
\[
 \sum_{\tau\in S,j}j\,c_{\tau,j}v_F(c)
       \le Dv_F(c)=\frac{v(g)D}{g-1}.
\]
\leanref{GlobalTheory.\allowbreak{}ChartCharacter.\allowbreak{}characters\_\allowbreak{}exist\_\allowbreak{}of\_\allowbreak{}chartBound}{Development.\allowbreak{}AbelianVarietyTheory.\allowbreak{}Proofs.\allowbreak{}JetHarmonics.\allowbreak{}ExistsCharactersBoundedMinorsOfHodgeTateChartBou}{295--343}

Apply the point construction to \(\bar\pi\), level \(g'\), the integral directions and type separators just obtained, the first \(n_1\) columns, the exponents \(c_{\tau,j}\), and the short period basis.
Its output consists, for each column \(b\), of a family
\(s_b\in\mathscr T_{g'}^+\), its normalised \([\bar\pi]\)-pullback, a finite-order translate at a point killed by \([\pi]\), and a germ \(f_b\).

The period representative used in each of the first \(n_1\) columns has
\[
 Q_X(z_0)/(gg')\le s\Phi/(4g')
\]
by Lemma~\ref{B8.28f}. Its separation inequality makes the selected coefficient the term of least valuation in the corresponding supported theta expansion. Their total contribution is at most
\(n_1s\Phi/(4g')\).
For the integral jets, Lemma~\ref{B8.28d} gives the contribution
\(D^2\Phi/g'\). The hypotheses \(4D+4\le g'\) and
\(2D\sqrt s<g'\) are exactly its level and short-basis conditions.
Together with the character contribution and \(D<g-1\), the point construction gives a finite harmonic value of valuation
\[
 \rho\le \frac{v(g)D}{g-1}
              +\frac{n_1s\Phi}{4g'}+\frac{D^2\Phi}{g'}.
\]
\leanref{GlobalTheory.\allowbreak{}T.\allowbreak{}depth\_\allowbreak{}bound\_\allowbreak{}from\_\allowbreak{}characters\_\allowbreak{}pointDatum}{Development.\allowbreak{}AbelianVarietyTheory.\allowbreak{}Proofs.\allowbreak{}JetHarmonics.\allowbreak{}DepthBoundFromCharactersPointDatumStatement}{13--15}

Form the external products of these families and translations.
The comparison between the one-column chart and the chosen integral product frame gives one integral series \(u\) with \(u(0)=1\) and a product germ
\[
 h(X_1,\ldots,X_N)=\prod_{b=1}^N f_b(X_b)u(X_b).
\]
For a harmonic polynomial the value in this product frame equals the column value in the canonical gauge. The two assigned column sets are disjoint, so the opposite-type mixed derivatives of \(P_S\) vanish by Lemma~\ref{B8.28ya}. Hence the bound for \(\rho\) passes to
\(\mathcal H_M(P_S;h)\). Finally \(D^2\Phi/g'\ge0\), so replacing its coefficient \(1\) by \(2\) gives the stated bound.
\leanref{GlobalTheory.\allowbreak{}T.\allowbreak{}integralChart\_\allowbreak{}columns\_\allowbreak{}to\_\allowbreak{}tuple\_\allowbreak{}harmonicValues}{Development.\allowbreak{}AbelianVarietyTheory.\allowbreak{}Proofs.\allowbreak{}JetHarmonics.\allowbreak{}IntegralChartColumnsToTupleHarmonicValuesStatement}{14--16}

The complete proof in the archive is contained in \Lean{GlobalTheory.\allowbreak{}PowGerm.\allowbreak{}lemmaBT\_\allowbreak{}of\_\allowbreak{}characters}.
\end{proof}

\subsection*{Proof of \ref{B8.28yc}}\label{prf:B8.28yc}\phantomsection
\begin{proof}[Proof (outline)]
Use the product degeneration datum, with period form \(Q_N\), the product polarisation and the componentwise \(\mathcal O_E\)-action.
The integral generation statement for this datum takes as inputs
\(s\in\mathscr T_{g'}^+\), its normalised \([\bar\pi]\)-pullback and its finite-order translate at \(\pi\)-torsion. It supplies a principal theta frame \(\vartheta\), its normalised
\([\bar\pi\bar\pi']\)-pullback, and finitely many translates at
\(\pi\pi'\)-torsion. In compatible product frames its conclusion is the exact identity
\[
 h=\sum_i a_i h_i,\qquad v_F(a_i)\ge0.
\]
The frame and germ statements apply on every finite thickening of the identity. Compatibility under truncation gives the displayed identity of power series. The coefficients are in the valuation ring; this is the integrality assertion in the generation statement.
\leanref{GlobalTheory.\allowbreak{}PowGerm.\allowbreak{}pointStatements\_\allowbreak{}on\_\allowbreak{}copies}{Development.\allowbreak{}AbelianVarietyTheory.\allowbreak{}Proofs.\allowbreak{}JetHarmonics.\allowbreak{}PointStatementsOnCopiesHold}{24--34}

Fix one of these translated families. Splitting its product frame into the one-factor frames writes its germ as a scalar of valuation zero times the external product of the corresponding one-factor germs. The comparison of each one-factor finite-order theta action with the symmetric theta-group action introduces a root of unity. Their product has valuation zero. The product therefore corresponds, up to one scalar \(u_i\) of valuation zero, to the product of the translated pullback sections associated with a characteristic tuple \(x_i\).

Apply the normalised pullback chain rule to the row-homogeneous polynomial \(P\). The factor on the row of type \(\tau\) is \(\tau(\alpha)\), hence the product factor is
\[
 \prod_\tau\tau(\alpha)^{\sum_j a_{\tau,j}}=\nu_a(\alpha).
\]
The identification of harmonic values of tied sections with the product-germ value then gives
\[
 \mathcal H_M(P;h_i)=u_i\nu_a(\alpha)\theta_{x_i,P}(y).
\]
The characteristic tuple and \(u_i\) are obtained from the sections and frames, so they serve all the harmonic polynomials in the assertion.
\leanref{GlobalTheory.\allowbreak{}ColumnSplit.\allowbreak{}main}{Development.\allowbreak{}AbelianVarietyTheory.\allowbreak{}Proofs.\allowbreak{}JetHarmonics.\allowbreak{}TranslateThetaFrameGermEqUnitMulBoxProduct}{51--113}
\leanref{GlobalTheory.\allowbreak{}OneColTranslate.\allowbreak{}oneColumn\_\allowbreak{}translate}{Development.\allowbreak{}AbelianVarietyTheory.\allowbreak{}Proofs.\allowbreak{}JetHarmonics.\allowbreak{}OneColumnTranslateSectionRootOfUnityMul}{45--104}
\leanref{GlobalTheory.\allowbreak{}PtTie.\allowbreak{}harmonicValue\_\allowbreak{}translate}{Development.\allowbreak{}AbelianVarietyTheory.\allowbreak{}Proofs.\allowbreak{}JetHarmonics.\allowbreak{}HarmonicValueTranslateEqUnitMulNormaliserThetaVa}{62--131}

Linearity in the germ gives
\[
 \mathcal H_M(P;h)=\sum_i a_i\mathcal H_M(P;h_i).
\]
If every \(\mathcal H_M(P;h_i)\) had valuation greater than \(\rho\), then every summand would have valuation greater than \(\rho\), since
\(v_F(a_i)\ge0\). The ultrametric inequality would contradict the definition of \(\rho\). Thus one index has value at most \(\rho\).
At that index \(v_F(u_i)=0\), and \(v_F\) restricts to \(v\), proving the final inequality.

The complete proofs in the archive are contained in \Lean{GlobalTheory.\allowbreak{}PowGerm.\allowbreak{}pointStatements\_\allowbreak{}on\_\allowbreak{}copies} and \Lean{GlobalTheory.\allowbreak{}PtTie.\allowbreak{}harmonicValue\_\allowbreak{}translate}.
\end{proof}

\subsection*{Proof of \ref{B8.28y}}\label{prf:B8.28y}\phantomsection
\begin{proof}
Order the half-level generators so that
\(\pi\bar\pi=g\), \(\pi'\bar\pi'=g'\) and
\(\mathfrak b=(\pi\pi')\). Coprimality of
\(\mathfrak b\) and \(\bar{\mathfrak b}\) implies both
\[
 (\pi)+(\bar\pi)=\mathcal O_E,\qquad
 (\pi')+(\bar\pi')=\mathcal O_E.
\]
Choose an algebraically closed valued extension \(F\) of
\(\operatorname{Frac}V'\), with rational-valued valuation extending \(v'\), a \(g\)-adically complete valuation ring, and a PEL Tate isometry compatible with the Weil pairings. Equip the \(N\)-fold product with the product degeneration datum and the componentwise action.
\leanref{GlobalTheory.\allowbreak{}T.\allowbreak{}exists\_\allowbreak{}algClosed\_\allowbreak{}valuedField\_\allowbreak{}over\_\allowbreak{}completeDVR}{Development.\allowbreak{}AbelianVarietyTheory.\allowbreak{}Proofs.\allowbreak{}JetHarmonics.\allowbreak{}ExistsAlgClosedValuedFieldOverCompleteDVRStatement}{15--17}
\leanref{GlobalTheory.\allowbreak{}T.\allowbreak{}nonempty\_\allowbreak{}tateIsometry\_\allowbreak{}integralFramedPoint\_\allowbreak{}chart}{Development.\allowbreak{}AbelianVarietyTheory.\allowbreak{}Proofs.\allowbreak{}JetHarmonics.\allowbreak{}NonemptyTateIsometryIntegralFramedPointChartStatement}{15--19}
\leanref{GlobalTheory.\allowbreak{}T.\allowbreak{}exists\_\allowbreak{}productData\_\allowbreak{}copies\_\allowbreak{}trivialisation\_\allowbreak{}valuation}{Development.\allowbreak{}AbelianVarietyTheory.\allowbreak{}Proofs.\allowbreak{}JetHarmonics.\allowbreak{}ExistsProductDataCopiesTrivialisationValuationStatement}{15--19}

Put \(n_1=r\) and \(n_2=N-r\). The hypothesis \(N\ge2r\) and
\(\mu_\tau\le r\) give \(\mu_\tau\le n_1,n_2\). Let
\[
 S=\{\tau:v_F(\tau(\bar\pi))>0\}.
\]
Since \(v_F(g)>0\), coprimality and
\(\tau(\pi)\tau(\bar\pi)=g\) imply that \(S\) contains exactly one member of each conjugate pair. Lemma~\ref{B8.28ya} supplies
\(P_1=P_S\in\operatorname{harm}(a,N)\), with degree \(D\) and the two required sets of columns.

The complete splitting of \(g\), the short basis of parameter \(\Phi'\), and the three numerical hypotheses allow Lemma~\ref{B8.28yb} to be applied. It gives the compatible germ \(h\) of a finite-order translate of a \([\bar\pi]\)-pullback of a family in \(\mathscr T_{g'}^+\), with
\[
 \rho=v_F(\mathcal H_M(P_1;h))
 \le \frac{v'(g)D}{g-1}
              +\frac{rs\Phi'}{4g'}+\frac{2D^2\Phi'}{g'}.
\]
Lemma~\ref{B8.28yc}, applied with
\(\alpha=\bar\pi\bar\pi'\), supplies a characteristic tuple \(x\) such that
\[
 v'(\theta_{x,P_1}(y'))+v'(\nu_a(\bar\pi\bar\pi'))\le\rho.
\]
Since \(\nu_a(\bar\pi\bar\pi')=\nu\), these inequalities prove the proposition.
\end{proof}

\subsection*{Proof of \ref{B8.28s}}\label{prf:B8.28s}\phantomsection
\begin{proof}
Take \((g,g')=(g_0,g_1)\) or \((g_1,g_0)\); interchanging \(\pi_0\) and \(\pi_1\) leaves \(\bar\pi_0\bar\pi_1\) and \(\nu=\nu_a(\bar\pi_0\bar\pi_1)\) unchanged. Fix a place \(\lambda\mid g\) of \(F_c\), a trait \(V\) of ramification index \(e\) over \(\lambda\), and an algebraic integral frame \(y\) on \(V\) with phase \(\phi\ge0\). Put \(e_0=v(g)=e\,v_\lambda(g)>0\). Pass to a finite extension \(V'\) of relative ramification index \(r'\ge1\) over \(V\) on which the image algebraic integral frame \(y'\) has a split uniformisation \cite[Ch.~III]{faltings-chai-1990}, with phase \(r'\phi\) and \(v'(g)=r'e_0\). Let \(s\le 2fq_{\rm sig}\le 2fm\) and \(\Phi'=\kappa s^3r'\phi\) be as in Lemma~\ref{B8.28c}. Since \(4D+4\le g\) implies \(D<g-1\), while \(4D+4\le g'\) and \(2D\sqrt s\le 2D\sqrt{2fm}<g'\), Proposition~\ref{B8.28y} on \(V'\) supplies \(P_1\in\operatorname{harm}(a,N)\) and a tuple \(x\) such that, after dividing by \(v'(g)=r'e_0\),
\[
 \frac{v(\theta_{x,P_1}(y))}{v(g)}+\frac{v_\lambda(\nu)}{v_\lambda(g)}
 \le\frac{D}{g-1}+\frac{\kappa s^3(rs/4+2D^2)\phi}{g'\,v(g)}.
\]
Here \(D/(g-1)<1\) because \(4D+4\le g\), and \(s\le 2fm\) together with the hypothesis on \(c\) gives \(\kappa s^3(rs/4+2D^2)\phi/(g'v(g))\le c\,\phi/v(g)\).

Put \(\theta=\theta_{x,P_1}\). As in the proof of Proposition~\ref{B8.28n}, Lemma~\ref{B8.25t} shows that \(\theta\) is a cusp section of weight \((a,N)\) over \(F_c\). Since \(\gcd(g_0g_1,Q_Sqd)=1\), Lemma~\ref{B8.28q}(1)--(3) gives the finitely supported lower bound \(m\) on the finite places of \(F_c\) with \(m(\lambda)=-v_\lambda(\nu)\in\Z\). Lemma~\ref{B8.28j} with \(k=-v_\lambda(\nu)\) therefore gives \(\sigma\in F_c^\times\) with \(v_\lambda(\sigma)=v_\lambda(\nu)\) such that \(f=\sigma\theta\) is an integral cusp section of weight \((a,N)\). On \(V\), \(v(\sigma)/v(g)=v_\lambda(\nu)/v_\lambda(g)\), so
\[
 \frac{v(f(y))}{v(g)}=\frac{v_\lambda(\nu)}{v_\lambda(g)}+\frac{v(\theta(y))}{v(g)}
 \le 1+\frac{c\,\phi}{v(g)}.
\]
Subtracting \(c\,\phi/v(g)\) and taking the supremum in Definition~\ref{B8.28g}(e) gives \(d^c_g(a,N)\le1\).
\end{proof}

\subsection*{Proof of \ref{B8.28z}}\label{prf:B8.28z}\phantomsection
\begin{proof}[Proof (outline)]
By the construction of \((a,N)\) from \(u_1\) in Lemma~\ref{B8.25g}, passing from weight \((a,N)\) to weight \(u_1\) on \(Y\) shifts the scalar exponent by \(-N/2\) and translates each CM row \(\tau\) by \(-s_\tau=u_{1,\tau,P-1}-N/2-1\). On each component of \(Y\), multiplying an integral cusp section of weight \((a,N)\) by the determinant section of parameter \(-N/2\) and the flat section of row \((-s_\tau)_\tau\) (supplied by the flat-degree law of Definition~\ref{def:p8c3-section-laws} and Lemmas~\ref{B8.24b} and~\ref{B8.24d}) produces a cusp section of weight \(u_1\), as in the proof of Lemma~\ref{B8.25g}.

Fix a prime \(\ell\). If \(\ell\notin\mathcal B\), then \(\ell\nmid 2dD_EM'\), so \(\ell\) is unramified in \(E\) and prime to the covering level \(M'\). At each place \(\lambda'\mid\ell\), the smooth local PEL model of Lemma~\ref{B8.24d} has reduced special fibre and \(S=\{\tau\}\), so on each component the determinant multiplier has valuation zero and the flat multiplier has a constant integer valuation (normalised by the trait ramification index) across all integral frames over \(\lambda'\); applying Lemma~\ref{B8.28j} at \(\lambda'\) with that integer cancels it without rounding error while clearing denominators away from \(\lambda'\), giving \(d^{1/2500}_\ell(u_1)\le d^{1/2500}_\ell(a,N)\).

If \(\ell\in\mathcal B\), comparing the projected Hodge lattice with the integral eigenframes via Lemma~\ref{B8.24c} and clearing denominators at the finite places away from each place \(\lambda'\mid\ell\) via Lemma~\ref{B8.28j} introduces a valuation shift at \(\ell\) bounded by \(1+\frac N2fr\,\delta_\ell\), where the term \(\frac N2fr\,\delta_\ell\) is the local different contribution from the determinant twist of exponent \(N/2\) on the rank-\(fr\) Hodge bundle and the additive \(1\) accounts for integer rounding in clearing denominators. Adding this shift to a prime-sum bound for \((a,N)\) increases the log-weighted sum by at most \(\frac N2fr\log\operatorname{rd}(E)+\sum_{\ell\in\mathcal B}\log\ell\).
\leanref{GlobalTheory.\allowbreak{}RefSum.\allowbreak{}ToWeight.\allowbreak{}sum\_\allowbreak{}different\_\allowbreak{}le}{Development.\allowbreak{}AbelianVarietyTheory.\allowbreak{}Proofs.\allowbreak{}DivisorProofs.\allowbreak{}ReferenceWeightOfTheta}{32--45}
\leanref{GlobalTheory.\allowbreak{}RefSum.\allowbreak{}ToWeight.\allowbreak{}primeSum\_\allowbreak{}add}{Development.\allowbreak{}AbelianVarietyTheory.\allowbreak{}Proofs.\allowbreak{}DivisorProofs.\allowbreak{}ReferenceWeightOfTheta}{51--71}

The complete proof in the archive is contained in \Lean{GlobalTheory.\allowbreak{}T.\allowbreak{}refW\_\allowbreak{}of\_\allowbreak{}prints\_\allowbreak{}size\_\allowbreak{}flatDeg\_\allowbreak{}refTheta}.
\end{proof}

\subsection*{Proof of \ref{B8.28t}}\label{prf:B8.28t}\phantomsection
\begin{proof}
(1) In Lemma~\ref{B8.25g}, the differential weight \((a,N)\) has \(a_{\tau,P}=1>0\) on the CM type, \(N\ge\sum_{\tau,i}2i\ge\sum_{c_{\tau,i}\ne0}2i\), \(N\ge2r\), and \(1\le D\le b^2\). By \eqref{eq:p8c3-berg-sizes}, \(b\ge1000\), \(r\le b\), \((2fm)^3\le(frL)^{12}\le P_*\le b\), \(\kappa=R_E^2/(4f)\le R_E^7\le b\), and \(g_0,g_1\ge b^8\). Hence for both \(g\in\{g_0,g_1\}\),
\[
 4D+4\le 4b^2+4\le g,\qquad 2D\sqrt{2fm}\le 2b^2\sqrt b<g,
\]
and
\[
 \kappa(2fm)^3\Bigl(\frac{r(2fm)}{4}+2D^2\Bigr)\le b\cdot b\cdot\Bigl(\frac{b^2}{4}+2b^4\Bigr)\le 3b^6\le\frac{g}{2500}.
\]
The same bounds hold with \(g_0g_1\ge b^{16}\) in place of \(g\), and \((Q_Sqd)^{36}t^{150}\le\sqrt b\) from Lemma~\ref{B8.28b} (with \(t\ge1\)) gives \(Q_Sqd\le\sqrt b<b^8\le g_0,g_1\), hence \(\gcd(g_0g_1,Q_Sqd)=1\).

At a nonterminal layer \(1\le j<m\), and \(p_{\rm sig}=k+j-1\ge k=(k-2)+2\) in Lemma~\ref{B8.25g} gives \(k\ge2\). By Lemmas~\ref{B8.25b} and~\ref{B8.25g} together with \eqref{eq:p8c3-spacing}, \eqref{eq:p8c3-quotient}, and \eqref{eq:p8c3-berg-sizes}, each row of \(u_1\) is decreasing with \(u_{1,\tau,k-2}=0\), \(0\le u_{1,\tau,s}\le(S+2r)/T\le 8b/T\) on \(0\le s<k-2\) (so \(\sum_{0\le s<k-2}u_{1,\tau,s}\le k(S+2r)/T\)), \(|u_{1,\tau,s}|\le(8b+c_0+4r)/(T+d_{\rm fl}+1)\le 8b/T\) on \(s>k-2\) (since \(c_0+4r\le 26P_*\) and \(26P_*T\le P_*Td_{\rm fl}^2L\le b<8(d_{\rm fl}+1)b\)), whence \(|u_{1,\tau,s}|\le\lfloor 8b/T\rfloor=B_0\) by integrality, and wall gap \(u_{1,\tau,P-1}-u_{1,\tau,P}\ge h_0/(2T)\ge 20fr^2+1\ge N+1\); thus the row hypotheses of Lemma~\ref{B8.28r}(1) hold. By the host construction (\ref{B2.46}), \(2\mid Q_S\) and every prime \(\ell\mid Q_S\) splits \(K=\Q(\sqrt{-d})\) over \(\Z_\ell\), so the host hypotheses of Lemma~\ref{B8.28r}(2) hold.

For \(\ell\in\{g_0,g_1\}\), Proposition~\ref{B8.28s} with \(c=1/2500\) gives \(d^{1/2500}_\ell(a,N)\le 1\le e_\ell+1\) on the base moduli since \(e_\ell\ge0\). For \(\ell\notin\{g_0,g_1\}\), Lemma~\ref{B8.28r} gives that the gap tuple of \((a,N)\) has oscillation at most \(e_\ell\) at \(\ell\); applying Lemma~\ref{B8.28p}(2) at each place \(\lambda\mid\ell\) establishes hypothesis~(i) of Proposition~\ref{B8.28n} with \(\varrho=e_\ell\) at every place above \(\ell\), and Lemma~\ref{B8.28q}(1)--(2) establishes hypothesis~(ii), so Proposition~\ref{B8.28n} gives \(d^{1/2500}_\ell(a,N)\le e_\ell\) on the base moduli. Lemma~\ref{B8.28l} then carries the inequality for every prime \(\ell\) from the base moduli to \(Y\).

(2) Combining (1) on \(Y\) with Proposition~\ref{B8.28z} gives
\[
 d^{1/2500}_\ell(u_1)\le e_\ell+{\bf1}_{\ell\in\{g_0,g_1\}}+{\bf1}_{\ell\in\mathcal B}\bigl(1+\tfrac N2fr\,\delta_\ell\bigr)=a(\ell)
\]
for every prime \(\ell\). Since \(e_\ell\) and \(\delta_\ell\) are nonnegative and \(e_\ell=0\) for \(\ell>J\), each \(a(\ell)\) is nonnegative and vanishes outside \(\mathcal S_{\rm ref}\). Splitting \(a(\ell)\log\ell\) across the support of \(e_\ell\) and \(\{g_0,g_1\}\cup\mathcal B\), with \(\sum_{\ell\in\mathcal B}\delta_\ell\log\ell\le\log\operatorname{rd}(E)\), reduces the weighted prime sum to Lemma~\ref{B8.28b}, which gives \(T\sum_{\ell\in\mathcal S_{\rm ref}}a(\ell)\log\ell\le C_RE_p\).
\end{proof}

\subsection*{Proof of \ref{B8.28}}\label{prf:B8.28}\phantomsection
\begin{proof}
Put \(n=k+m-1\), \(T=T_*\), \(P=P_*\), \(L=L_x\) and \(Z_0=fnL\). Both signature indices are positive and \(f>2\). The size, level and flat bounds in \ref{B8.25b}, \ref{B8.25c} and \ref{B8.24a} give
\begin{equation}\label{eq:p8c3-short-conditions}
 \begin{gathered}
 u_1\ \text{decreasing with wall gap}\ \ge8fn+10,\qquad
 \log\|u_1\|_1\le10L,\quad \log M\le30L,\quad \log(fn)\le L,\\
 \log[E:\Q]+\log\operatorname{rd}(E)\le3L,\qquad
 \log[F':\Q]+\log\operatorname{rd}(F')\le391Z_0^2,\qquad
 |c_1|\le9nL ,
 \end{gathered}
\end{equation}
where \(c_1\) is the flat degree and \(F'\) is the cover field. Set \(m_0=fn+1\). The wall gap is at least \(2(4m_0)+2\), so use \(8m_0\) theta columns. Translate row \(\tau\) by \(n_\tau=1+4m_0-u_{1,\tau,p_{\rm sig}-1}\). Subtract \(4m_0\) from the translated first block and from the reversed negative second block to obtain the theta exponents at the CM embedding and its conjugate. Let \(a\) denote this exponent array. For the exponent array \(a\) and these \(8m_0\) columns, the theta spanning sections have central parameter \(4m_0\) and row entries \(u_{1,\tau k}+n_\tau\). To clear finite denominators, multiply the raw theta family by the common positive integer \(N_a\) of the integrality bound and, for each lattice in the family, a nonzero short scalar in its reflex ideal. These scalars preserve the span over \(F'\). Their log norms are included in the theta spanning constant \(C_G\).
\leanref{GlobalTheory.\allowbreak{}CovGmin.\allowbreak{}minimal\_\allowbreak{}level\_\allowbreak{}cover\_\allowbreak{}is\_\allowbreak{}good}{Development.\allowbreak{}AutomorphicArithmetic.\allowbreak{}ThetaCorrespondence.\allowbreak{}ThetaPairLifts.\allowbreak{}HodgeBundleWeights.\allowbreak{}ExistsGoodCoverAtMinimalLevel}{20--57}

Multiply the resulting integral sections by the determinant section of parameter \(-4m_0\) and the flat section of row \((-n_\tau)_\tau\) on \(Z\). The parameters cancel and the resulting row is \(u_1\). The two factors are invertible on the component for this change of weight, and the integral choice of twist gives nonnegative lower orders. Thus the products belong to the integral lattice and span over \(F'\).

Write \(S=\|u_1\|_1\), \(D_F=[F':\Q]\) and \(v\) for the component cusp rank. Let \(S'=\sum_{\tau,k}|u_{1,\tau k}+n_\tau|\) be the absolute sum of the translated row entries. The three log norm costs are bounded by
\begin{equation}\label{eq:p8c3-short-twist-costs}
 \begin{gathered}
 C_G Z_0^2 S',\qquad 2m_0fn\log(2\pi),\\
 (|c_1|+q_{\rm sig}\log\operatorname{rd}(E))\sum_\tau|n_\tau|
 +\sum_{\ell\ {\rm bad}}\log\ell+\tfrac12\log\operatorname{rd}(F')+1 .
 \end{gathered}
\end{equation}
The bad-prime log sum is at most \(C_{\rm bad}Z_0^5\) by the good-cover bound. The formula for \(n_\tau\) gives \(\sum_\tau|n_\tau|\le S+f(4m_0+1)\). Summing the translated absolute entries gives
\[
 S'\le S+n\sum_\tau|n_\tau|
 \le(1+n)S+nf(4m_0+1).
\]
With \eqref{eq:p8c3-short-conditions}, the sum of the costs is at most \(C_2(Z_0^3S+Z_0^5)\), where one may take \(C_2=11C_G+C_{\rm bad}+500\). Apply \ref{B8.25f} in \(F'\) with the component Petersson norms. The lattice is finitely generated by the component lattice comparison, and the Gram norm gives subadditivity. There are at most \(2D_Fv\) generators, at extra log radius at most \(4\log D_F+\log v+\log\operatorname{rd}(F')\) when \(v>0\). The field bound in \eqref{eq:p8c3-short-conditions} and the rank bound \(\log v\le C_b Z_0^3\) from Tamagawa comparison give
\[
 c_s=C_g(Z_0^3S+Z_0^5),\qquad C_g=2000+C_b+C_2 .
\]
The metric shifted by this log radius has short integral generators of norm at most one. With the conjugation-compatible Petersson metric these give \(h_Z+c_s\ge0\) by the height comparison on the image.
\leanref{GlobalTheory.\allowbreak{}spanShortSize}{Development.\allowbreak{}AbelianVarietyTheory.\allowbreak{}Proofs.\allowbreak{}DivisorProofs.\allowbreak{}ThetaWeightAndFlatExponentSizeBound}{32--67}
\leanref{GlobalTheory.\allowbreak{}ShortSpanComp.\allowbreak{}short\_\allowbreak{}span\_\allowbreak{}comp\_\allowbreak{}mem}{Development.\allowbreak{}AutomorphicArithmetic.\allowbreak{}ThetaCorrespondence.\allowbreak{}ThetaPairLifts.\allowbreak{}HodgeBundleWeights.\allowbreak{}ShortSpanCompositionClosed}{24--154}
\leanref{GlobalTheory.\allowbreak{}ShortGenComp.\allowbreak{}short\_\allowbreak{}generator\_\allowbreak{}comparison}{Development.\allowbreak{}AutomorphicArithmetic.\allowbreak{}ThetaCorrespondence.\allowbreak{}ThetaPairLifts.\allowbreak{}HodgeBundleWeights.\allowbreak{}ShortGeneratorNormComparison}{23--102}
\leanref{GlobalTheory.\allowbreak{}CovHt.\allowbreak{}ShortPf.\allowbreak{}weight\_\allowbreak{}rows}{Development.\allowbreak{}AbelianVarietyTheory.\allowbreak{}Proofs.\allowbreak{}DivisorProofs.\allowbreak{}ShortOfInputs}{112--202}
\leanref{GlobalTheory.\allowbreak{}CovHt.\allowbreak{}ShortPf.\allowbreak{}cover\_\allowbreak{}height\_\allowbreak{}short\_\allowbreak{}section\_\allowbreak{}alt}{Development.\allowbreak{}AutomorphicArithmetic.\allowbreak{}ThetaCorrespondence.\allowbreak{}ThetaPairLifts.\allowbreak{}ThirdPlaceSeries.\allowbreak{}ExistsShortSectionCoverHeight}{21--47}

Let \(d=d_{\rm fl}\) and \(J=6d+10+T(z+2\rho)\). The divided weight and size estimates give \(T\le b\), \(TS\le8fnb\), \(d^2P\le T\) and \(d\rho P\le6n^2\). Expanding \(z\) gives \(PJ\le23Pd^2+9PTd\rho\le77n^2T\). Since \(Z_0^3n^3,Z_0^5n^2\le P\),
\begin{equation}\label{eq:p8c3-short-budget}
 (Z_0^3 S+Z_0^5)J\le77(8fb+fb)=693fb\le693E_p .
\end{equation}

By \ref{B8.25c}, \(5^4\mid M\) when \(e\ge4\), and the complementary halves give the linear level. Choose a rule \(f_0\) from spanning classes with nonzero generic evaluation on \(Z\). With \(a,G,\phi,P_0,B'_0\) as in \ref{B8.25g}, put \(\beta=P_0+B'_0\). In the thin case \(3(P_0-\phi)/625\le\beta-\phi/625\), and in the thick case \(3P_0/625\le\beta-B'_0\). Hence \ref{B8.25g} gives
\begin{equation}\label{eq:p8c3-theta-mult}
 a^3\le e^\beta G,\qquad a\le e^\beta
\end{equation}
at every frame on \(Z\). Choose the determinant twist \(K\) of power \(3n\) and the constant flat section \(S_Z\) of row \(-12n\) at each embedding, supported on \(Z\). Their point norms are constant in the frame on \(Z\), and their generic evaluations are invertible there. The section
\[
 \theta_Z=(f_0^2)^{12}KS_Z
\]
has weight \(24n\eta_H=8(u_1-u')\), nonzero generic evaluation, and squared point norm at most \(e^{8\beta}G^8|KS_Z|^2\), since \((a^3)^8=a^{24}\). Bounded multiplication in this weight sends cusp sections to cusp sections.

At a finite place \(\lambda\), let \(e_0\) be the ramification index of the test trait and \(\phi_{\rm fin}\ge0\) its degeneration parameter. The chosen tuple spans the full level module, and the trait and its integral frame realise the phase \(\phi_{\rm fin}\). Apply Lemma~\ref{B8.28a} with the valuation taking value one on a trait uniformiser. Division by \(e_0\) gives \(\operatorname{ord}(f_0)\ge\phi_{\rm fin}/(625e_0)\) in the coefficient-place normalisation. The determinant twist has nonnegative lower order, and the flat twist has lower order \(-8e_\lambda\) for its finite allowance \(e_\lambda\). Orders add, and \(24/625\ge8/2500\), so
\begin{equation}\label{eq:p8c3-carrying-order}
 \operatorname{ord}(\theta_Z)\ge 8\phi_{\rm fin}/(2500e_0)-8e_\lambda .
\end{equation}
The cost of the finite allowances and the archimedean bound before correcting the reference exponent is bounded by
\[
 \beta/2+(3n)(2fn\log(2\pi))/8-\deg(S_Z)/8 .
\]
The flat degree and size bounds of \ref{B8.24a} and \ref{B8.25b} bound the last two terms together by \(10P\). Thus \(c_{\rm raw}=\beta/2+10P\) is a common bound before correction.

Proposition~\ref{B8.28t}(2) gives a finite set \(\mathcal S_{\rm ref}\) of rational primes and nonnegative \(a(q)\), zero off that set, such that
\[
 d_q(u_1)\le a(q),\qquad
 \sum_{q\in\mathcal S_{\rm ref}}a(q)\log q\le s_{\rm ref},\qquad
 Ts_{\rm ref}\le C_RE_p .
\]
For the prime weights displayed in Lemma~\ref{B8.28b}, one may take \(s_{\rm ref}=\sum_q a(q)\log q\) in the two prime-sum inequalities. Here \(d_q(u_1)\) is the supremum over the coefficient places above \(q\), complete discrete valuation traits above those places, and algebraic integral frames on them with phase \(\phi\), of the local section depth in weight \(u_1\) minus \(\phi/(2500v(q))\). The depth uses evaluation valuations divided by the trait valuation \(v(q)\). Take \(e'_\lambda=e_\lambda+\sum_{q\in\mathcal S_{\rm ref}}v_\lambda(q)a(q)\). At each \(\lambda\mid q\), \(e'_\lambda\ge e_\lambda+v_\lambda(q)a(q)\). The identity \(\sum_{\lambda\mid q}v_\lambda(q)\log N\lambda=[F':\Q]\log q\) gives an added cost at most \(s_{\rm ref}\) in the normalised degree.

For a record valuation \(v\) over \(\lambda\mid q\), extend its integral frame to a trait with phase \(\phi_{\rm fin}\) and ramification index \(e_0\), preserving the normalised valuation. If a cusp section on \(Z\) has finite \(v\)-value, its evaluation on the lifted frame is nonzero and the generic point is on \(Z\). Let \(d_v\) be the infimum of the valuations of the integral cusp sections of target weight \(u_1\) on \(Z\). The reference bound and restriction to \(Z\) give
\[
 d_v\le B_v=\phi_{\rm fin}/(2500e_0)+v_\lambda(q)a(q).
\]
Thus for every coefficient place and every prescribed valuation there, if some cusp section on \(Z\) of any weight has finite valuation, an integral cusp section on \(Z\) of weight \(u_1\) has finite valuation. Applying \eqref{eq:p8c3-carrying-order} on the same trait gives \(v(\theta_Z)\ge8(B_v-e'_\lambda)\), the order comparison used in multiplication. The corrected cost bound \(c_{\rm mult}=c_{\rm raw}+s_{\rm ref}\) satisfies
\begin{equation}\label{eq:p8c3-mult-budget}
 Tc_{\rm mult}\le(C_\beta+10+C_R)E_p,\qquad C_\beta=200000+835C_{\rm peak} ,
\end{equation}
by \eqref{eq:p8c3-berg-budget} and \eqref{eq:p8c3-berg-sizes}.
\leanref{GlobalTheory.\allowbreak{}T.\allowbreak{}eta\_\allowbreak{}bound\_\allowbreak{}length\_\allowbreak{}unit\_\allowbreak{}alt}{Development.\allowbreak{}AutomorphicArithmetic.\allowbreak{}ThetaCorrespondence.\allowbreak{}ThetaPairLifts.\allowbreak{}ThirdPlaceSeries.\allowbreak{}RunParamsStatement}{11--12}
\leanref{GlobalTheory.\allowbreak{}T.\allowbreak{}bergman\_\allowbreak{}kernel\_\allowbreak{}thin\_\allowbreak{}part\_\allowbreak{}estimate\_\allowbreak{}alt}{Development.\allowbreak{}AutomorphicArithmetic.\allowbreak{}ThetaCorrespondence.\allowbreak{}ThetaPairLifts.\allowbreak{}ThirdPlaceSeries.\allowbreak{}BergmanKernelThinPartEstimate}{58--59}
\leanref{GlobalTheory.\allowbreak{}T.\allowbreak{}exp\_\allowbreak{}neg\_\allowbreak{}B\_\allowbreak{}le\_\allowbreak{}bergOn\_\allowbreak{}of\_\allowbreak{}la1sq\_\allowbreak{}bound}{Development.\allowbreak{}AutomorphicArithmetic.\allowbreak{}ThetaCorrespondence.\allowbreak{}ThetaPairLifts.\allowbreak{}ThirdPlaceSeries.\allowbreak{}BergmanKernelThickPartEstimate}{87--89}
\leanref{GlobalTheory.\allowbreak{}CovHt.\allowbreak{}KeySection.\allowbreak{}theta\_\allowbreak{}function\_\allowbreak{}le\_\allowbreak{}bergman\_\allowbreak{}kernel\_\allowbreak{}alt}{Development.\allowbreak{}AutomorphicArithmetic.\allowbreak{}ThetaCorrespondence.\allowbreak{}ThetaPairLifts.\allowbreak{}ThirdPlaceSeries.\allowbreak{}ExistsSectionBoundedByBergmanKernel}{24--73}
\leanref{GlobalTheory.\allowbreak{}CovHt.\allowbreak{}KeySection.\allowbreak{}exists\_\allowbreak{}section}{Development.\allowbreak{}AbelianVarietyTheory.\allowbreak{}Proofs.\allowbreak{}DivisorProofs.\allowbreak{}ExistsCarryingSectionOrders}{126--267}
\leanref{GlobalTheory.\allowbreak{}CovHt.\allowbreak{}KeySection.\allowbreak{}cover\_\allowbreak{}height\_\allowbreak{}section\_\allowbreak{}bound\_\allowbreak{}alt}{Development.\allowbreak{}AutomorphicArithmetic.\allowbreak{}ThetaCorrespondence.\allowbreak{}ThetaPairLifts.\allowbreak{}ThirdPlaceSeries.\allowbreak{}ExistsSectionBoundedByBergmanKernel}{75--116}
\leanref{GlobalTheory.\allowbreak{}CovMult.\allowbreak{}multC\_\allowbreak{}of\_\allowbreak{}refW\_\allowbreak{}alt}{Development.\allowbreak{}AutomorphicArithmetic.\allowbreak{}ThetaCorrespondence.\allowbreak{}ThetaPairLifts.\allowbreak{}ThirdPlaceSeries.\allowbreak{}MultCOfRefW}{21--29}

For \(u'\) let \(g=h_0/(4T)\), where \(h_0\) is from \eqref{eq:p8c3-spacing}. Starting from any smaller-block index, take the adjacent step towards the comparison index. The adjacent-difference estimate with the wall correction gives a gap at least \(g\) on that step. The rest of the row is monotone, so this is a gap in the sense of \eqref{eq:p8c3-error}. The absolute entries are at most \(8b/T\). This bound and \(fn^2\le P\le h_0\) give \(1\le B(u',g)\le132b\). In a row of node scale \(w_\tau\) the spread is at most \(3w_\tau/(2T)\) and the \(m\) node entries have absolute value at least \(w_\tau/(4T)\). Since \(q_{\rm sig}\le m\),
\begin{equation}\label{eq:p8c3-row-budget}
 q_{\rm sig}\operatorname{spr}(u')\le6\|u'\|_1,\qquad \log B+\ell(D)+6\le100L .
\end{equation}
For the second inequality use the period identity and \(|\log\Omega|\le3L\), \(0\le\log\operatorname{rd}(E)\le2L\), which give \(|L_E|\le7L\).

At most four primes divide \(M'\). By \eqref{eq:p8c3-prime-sum} the primes up to \(B\) cost at most \(\log B+3\) in the reciprocal-prime sum. There are at most \(2+4\log\operatorname{rd}(E)+\#\{\ell:\ell\mid M'\}\) further primes, each costing at most one. Including the different gives
\begin{equation}\label{eq:p8c3-local-budget}
 \sum_{\ell\in S_B}h_E(\ell)\log\ell\le\log B+20+44\log\operatorname{rd}(E)\le100L .
\end{equation}
Put \(H_{\rm loc}=\|u'\|_1\sum_{\ell\in S_B}h_E(\ell)\log\ell\). Equations \eqref{eq:p8c3-row-sum}, \eqref{eq:p8c3-row-budget} and \eqref{eq:p8c3-local-budget} give
\begin{equation}\label{eq:p8c3-cost-budget}
 TH_{\rm loc},\ T\mathcal E_{\ell(D)}(u',g)\le83400E_p .
\end{equation}
\leanref{GlobalTheory.\allowbreak{}MeanTop.\allowbreak{}spread\_\allowbreak{}le\_\allowbreak{}at}{Development.\allowbreak{}AbelianVarietyTheory.\allowbreak{}Proofs.\allowbreak{}DivisorProofs.\allowbreak{}SpreadLeAt}{10--38}
\leanref{GlobalTheory.\allowbreak{}CovHt.\allowbreak{}BudgetPf.\allowbreak{}sizes}{Development.\allowbreak{}AbelianVarietyTheory.\allowbreak{}Proofs.\allowbreak{}DivisorProofs.\allowbreak{}BudgetOfSizeAndWeight}{151--245}
\leanref{GlobalTheory.\allowbreak{}CovHt.\allowbreak{}BudgetPf.\allowbreak{}budgetAt}{Development.\allowbreak{}AbelianVarietyTheory.\allowbreak{}Proofs.\allowbreak{}DivisorProofs.\allowbreak{}BudgetOfSizeAndWeight}{246--314}

The family has \(m^2\le k\). Apply \ref{B8.26} to \(u'\) and set
\begin{equation}
 c_c=c_{\rm mult}+H_{\rm loc},\qquad X=C_{\rm deep}\mathcal E_{\ell(D)}(u',g).
 \tag{8.28e}\label{eq:p8c3-height-costs}
\end{equation}
The resulting \(s_Z\) have mean at least \(-X\). For a section \(f_1\) of weight \(u\) and central parameter zero, let \(v\) be a record frame valuation over \(\lambda\mid\ell\). As \(E/\Q\) is Galois, lift the frame over a completed algebraic closure and use the Hodge--Tate frame valuation \(w\) over the same prime. The triangular change has diagonal entries \(t_s\) with \(|\log|t_s|/\log|\ell||\le h_E(\ell)\) and scalar entry of absolute value one. The diagonal character has absolute exponent sum at most \(\|u\|_1\). Hence
\begin{equation}\label{eq:p8c3-frame-comparison}
 w(f_1)\le v(f_1)+v_\lambda(\ell)h_E(\ell)\|u\|_1 .
\end{equation}
\leanref{GlobalTheory.\allowbreak{}CovHt.\allowbreak{}FrameCmp.\allowbreak{}exists\_\allowbreak{}frBor}{Development.\allowbreak{}AbelianVarietyTheory.\allowbreak{}Proofs.\allowbreak{}DivisorProofs.\allowbreak{}SumAbsExOfLe}{278--345}
\leanref{GlobalTheory.\allowbreak{}CovHt.\allowbreak{}FrameCmp.\allowbreak{}val\_\allowbreak{}le}{Development.\allowbreak{}AbelianVarietyTheory.\allowbreak{}Proofs.\allowbreak{}DivisorProofs.\allowbreak{}SumAbsExOfLe}{346--373}
\leanref{GlobalTheory.\allowbreak{}CovHt.\allowbreak{}FrameCmp.\allowbreak{}htCmp}{Development.\allowbreak{}AbelianVarietyTheory.\allowbreak{}Proofs.\allowbreak{}DivisorProofs.\allowbreak{}SumAbsExOfLe}{378--388}

In degree \(8r'\) balance the corrected finite exponents and the infinite weight of \(\theta_Z^{r'}\) by a scalar. Its infinite log bound has an additive constant depending on the fixed component and divisor, uniform in \(r'\). Multiply the scalar by \(\prod_{\ell\in S_B}\ell^{\lceil8r'h_E(\ell)\|u'\|_1\rceil}\) to clear the extra denominators in \eqref{eq:p8c3-frame-comparison}. The additional log cost is at most \(8r'H_{\rm loc}+\sum_{\ell\in S_B}\log\ell\). Denote the resulting scalar by \(\alpha_{8r'}\). For any \(\eta>0\), the fixed terms divided by \(8r'\) are at most \(\eta\) for large \(r'\). Multiplication by \(\alpha_{8r'}\theta_Z^{r'}\) then maps mixed small sections of \(u'\) at \(c_s\) to small sections of \(u_1\) at \(t=c_s+c_c+\eta\). It is injective by generic evaluation.

The base-point condition and the Gram norm give the image model of degree \(\Delta_Z\). Its shifted Fubini--Study line is nef by the integral generator bound at \(c_s\). Apply the counting injection of \ref{B8.25} and the nef volume equality of \ref{B8.25a}. The self-intersection enters divided by \([F':\Q](d_{\rm fl}+1)\), since the arithmetic-volume normalisation uses \((d_{\rm fl}+1)!\) and the count uses \(d_{\rm fl}!\). The Chow-height comparison bounds this quotient by \(\Delta_Z(h_Z+t)\). By \(s_Z\le\nu_Z(u';B)\), the lower count on \(u'\) gives
\begin{equation}
 \max(0,\Delta'_Z(s_Z+c_s))\le\Delta_Z(h_Z+c_s+c_c+\eta).
 \tag{8.28f}\label{eq:p8c3-height-volume}
\end{equation}
\leanref{GlobalTheory.\allowbreak{}CovMult.\allowbreak{}exists\_\allowbreak{}bound}{Development.\allowbreak{}AbelianVarietyTheory.\allowbreak{}Proofs.\allowbreak{}DivisorProofs.\allowbreak{}CarryingSectionIntegralFrame}{145--171}
\leanref{GlobalTheory.\allowbreak{}CovMult.\allowbreak{}mult\_\allowbreak{}of}{Development.\allowbreak{}AbelianVarietyTheory.\allowbreak{}Proofs.\allowbreak{}DivisorProofs.\allowbreak{}CarryingSectionIntegralFrame}{172--194}
\leanref{SelfIntersectionVolumeEstimates.\allowbreak{}ImageNefG.\allowbreak{}exists\_\allowbreak{}nef\_\allowbreak{}model\_\allowbreak{}of\_\allowbreak{}image}{Development.\allowbreak{}AutomorphicArithmetic.\allowbreak{}ThetaCorrespondence.\allowbreak{}ThetaPairLifts.\allowbreak{}HodgeBundleWeights.\allowbreak{}ExistsNefModelOfImage}{13--41}
\leanref{SelfIntersectionVolumeEstimates.\allowbreak{}VolUpG.\allowbreak{}volume\_\allowbreak{}line\_\allowbreak{}bundle\_\allowbreak{}le}{Development.\allowbreak{}AutomorphicArithmetic.\allowbreak{}ThetaCorrespondence.\allowbreak{}ThetaPairLifts.\allowbreak{}HodgeBundleWeights.\allowbreak{}VolumeLineBundleLe}{11--22}

Let \(\eta\downarrow0\) and use \eqref{eq:p8c3-shift-degree}. Equations \eqref{eq:p8c3-short-budget}, \eqref{eq:p8c3-mult-budget} and \eqref{eq:p8c3-cost-budget} give \(C_{\rm ht}=693C_g+C_\beta+10+C_R+83400(1+C_{\rm deep})\). The maximum of the input thresholds works for all the covers and weights.
\end{proof}

\subsection*{Proof of \ref{B8.29}}\label{prf:B8.29}\phantomsection
\begin{proof}
Put \(T=T_*\), \(d=d_{\rm fl}\), \(w_Z=\mathcal B\Delta_Z\), \(a=T-3d-4\ge0\), and \(L_A=A\log\max(A,1)\). Use \(\log y\le y-1\) at \(y=(2+t)/\max(A,1)\) to get
\begin{equation}\label{eq:p8c3-log-t}
 A\log(2+t)\le t+L_A+2\qquad(t\ge0).
\end{equation}
Since \(|h_Z|\le h_Z+2c_s\), \eqref{eq:p8c3-hs} gives
\begin{equation}\label{eq:p8c3-dZ}
 d_Z\ge w_Z\bigl(a(h_Z+c_s)-N_0\bigr),\qquad
 N_0=ac_s+2c_s(3d+4)+L_A+2+R .
\end{equation}
By \ref{B8.28} the nonnegative \(h_Z+c_s\) have mean at least \(m_0=\max(0,(1-\rho)c_s-X-c_c)\). By \eqref{eq:p8c3-shift-degree}, \(w_Z\ge\mathcal B\Delta_*\) and \((1-\rho)\sum_Zw_Z\le\#\operatorname{Comp}(Y)\mathcal B\Delta_*\). Thus
\begin{equation}\label{eq:p8c3-degree-weighted}
 \sum_Zd_Z\ge\left(\sum_Zw_Z\right)\bigl(a(1-\rho)m_0-N_0\bigr).
\end{equation}
Use \((1-\rho)m_0\ge(1-2\rho)c_s-X-c_c\) and \(a\le T\) to bound the bracket below by \(-H\), where
\begin{equation}\label{eq:p8c3-degree-budget}
 H=T(X+c_c)+c_s(6d+10+T(z+2\rho))+L_A+R+4 .
\end{equation}
Equations \eqref{eq:p8c3-shares}, \eqref{eq:p8c3-hs} and \eqref{eq:p8c3-height-budget} give \(0<\sum_Zw_Z\le\sum_Zr_Z\) and \(H\le(3C_{\rm ht}+C_{\rm HS}+4)E_p\). Apply \eqref{eq:p8c3-degree}.

For the nonterminal descent choose a good cover. Let \(A_-\) be its nonzero lattice at shift \(-c_0\) and \(B\) its deep type lattice at shift zero. By \ref{B8.24a} and \eqref{eq:p8c3-cover-slope},
\begin{equation}\label{eq:p8c3-shift-slope}
 \mu(A_-)\ge-M_{p,i}-(C_F+C_U)E_p .
\end{equation}
Put \(\mathcal K=P_*\kappa\), so by \ref{B8.25b}
\begin{equation}\label{eq:p8c3-small-b}
 0<b\le E_p,\qquad b\le\mathcal K,\qquad |M_{p,i}|\le\mathcal K .
\end{equation}
Let \(S\) be the horizontal subgroup of the level automorphisms and put \(s=|S|\). Orbit separation gives a cusp section \(u_0\) over the cover field of weight \(c_0\eta_H\) and scalar exponent zero. Its translation determinant has nonzero generic value on every component. Expand \(u_0\) in the short integral generators of the component lattices. Fix all coefficients except that of a generator supported on \(W\). The translations by \(S\) send \(W\) to at most \(s\) components, so the determinant can vary on at most \(s\) components. On each of them the translation matrix at the generic point has entries affine in the coefficient. Its determinant is a polynomial of degree at most \(s\) and is nonzero at the current coefficient. Thus at most \(s^2\) integers are excluded. Choose a coefficient in \(\{0,\ldots,s^2\}\) outside this set and repeat for every generator. The resulting section \(u\) is integral, and its translation determinant is a nonzero divisor by generic evaluation on the components. If each component has at most \(N_{\rm gen}\) generators of Petersson norm at most \(\Lambda\), and the point norm is bounded by \(\beta\) times the Petersson norm on that component, then every point norm of \(u\) is at most \(N_{\rm gen}s^2\beta\Lambda\). The generator families have disjoint component supports, so at most \(N_{\rm gen}\) summands contribute at any frame.
\leanref{GlobalTheory.\allowbreak{}T.\allowbreak{}orbits\_\allowbreak{}separated\_\allowbreak{}by\_\allowbreak{}regular\_\allowbreak{}cusp}{Development.\allowbreak{}AutomorphicArithmetic.\allowbreak{}ThetaCorrespondence.\allowbreak{}ThetaPairLifts.\allowbreak{}HodgeBundleWeights.\allowbreak{}ExistsCuspRegularLevelDeterminant}{205--206}
\leanref{GlobalTheory.\allowbreak{}ShortDet.\allowbreak{}exists\_\allowbreak{}nat\_\allowbreak{}step}{Development.\allowbreak{}AbelianVarietyTheory.\allowbreak{}Proofs.\allowbreak{}DivisorProofs.\allowbreak{}CardLeOfDetEqZero}{178--221}
\leanref{GlobalTheory.\allowbreak{}ShortDet.\allowbreak{}exists\_\allowbreak{}short}{Development.\allowbreak{}AbelianVarietyTheory.\allowbreak{}Proofs.\allowbreak{}DivisorProofs.\allowbreak{}CardLeOfDetEqZero}{279--337}

Multiplication by the translates of \(u\) and projection with the horizontal character give integral maps \(\phi_h:A_-\to B\). The determinant gives joint injectivity. Let \(B_0\) bound the multiplication norm of a translate and put \(Z_0=frL_x\). The short-section estimates give \(B_0\le\exp(CZ_0^9)\) with a uniform \(C\ge0\), and the group bound gives \(3\log s\le135Z_0^2\). At the family threshold \(Z_0\ge135+2C\) and \(Z_0^{12}\le b\). Thus \(s^3B_0^2\le e^{2b}\). Cauchy--Schwarz bounds the squared joint norm by \(s^3B_0^2\), and the translate multiplications and the cover rank comparison give the rank bound in
\begin{equation}\label{eq:p8c3-joint}
 \sum_h\|\phi_h f_1\|_\sigma^2\le e^{2b}\|f_1\|_\sigma^2,\qquad
 s\operatorname{rk}B\le r_Y(G_s(c_0))
 \le(1+b/\mathcal K)\operatorname{rk}A_- .
\end{equation}
\leanref{GlobalTheory.\allowbreak{}T.\allowbreak{}exists\_\allowbreak{}short\_\allowbreak{}integral\_\allowbreak{}vector\_\allowbreak{}of\_\allowbreak{}determinant\_\allowbreak{}bound}{Development.\allowbreak{}AutomorphicArithmetic.\allowbreak{}ThetaCorrespondence.\allowbreak{}ThetaPairLifts.\allowbreak{}HodgeBundleWeights.\allowbreak{}ExistsShortIntegralVectorOfDeterminantBound}{126--127}
So the joint map has height at most \(2b\). The short spanning vectors in \(B\) have norms at most \(e^{\mathcal K}\). Taking a basis among their images in any quotient of \(B^s\) bounds its determinant norm by the product of these norms and gives \(\mu_{\min}(B^s)\ge-2\mathcal K\). Put \(r=\operatorname{rk}A_-\), \(N=s\operatorname{rk}B\). The degree inequality for the image and quotient gives \(r(\mu(A_-)-2b)\le N\mu(B)+2(N-r)\mathcal K\). With \(r\le N\le(1+b/\mathcal K)r\), for all \(t\le\mu(A_-)-2b\),
\begin{equation}\label{eq:p8c3-descent}
 \mu(B)\ge t-(b/\mathcal K)(|t|+2\mathcal K).
\end{equation}
Take \(t=-M_{p,i}-(C_F+C_U)E_p-2b\). Equation \eqref{eq:p8c3-small-b} gives \(\mu(B)\ge-M_{p,i}-2(C_F+C_U)E_p-7b\). The ring comparison gives \(\mu_{\rm core}\ge\mu(B)-2b\).
\leanref{GlobalTheory.\allowbreak{}CovTop.\allowbreak{}RingAt}{Development.\allowbreak{}AbelianVarietyTheory.\allowbreak{}Definitions.\allowbreak{}Divisors.\allowbreak{}CovolXOfCovol}{116--117}
\leanref{GlobalTheory.\allowbreak{}CovTop.\allowbreak{}DescentMaps.\allowbreak{}covering\_\allowbreak{}descent\_\allowbreak{}maps\_\allowbreak{}compatible\_\allowbreak{}alt}{Development.\allowbreak{}AutomorphicArithmetic.\allowbreak{}ThetaCorrespondence.\allowbreak{}ThetaPairLifts.\allowbreak{}ThirdPlaceSeries.\allowbreak{}CoveringDescentMapsCompatible}{27--118}

At the terminal layer \(q_{\rm sig}=0\). For the definite-side bound \(D_{\rm def}\), the definite Gram comparison gives \(\mu_{\rm core}\ge2D_{\rm def}-C_0E_p\) with a uniform \(C_0\). Let \(a_\tau\) be the anchor of \ref{B8.24a}, \(\lambda_{\tau a}\) the unshifted entries, \(n_{\rm col}\) the column count and \(V^{\rm Weyl}_\tau\) the dual Weyl module of the row partition. The definite-side estimate is
\begin{equation}\label{eq:p8c3-terminal}
 \begin{split}
 2D_{\rm def}\ge{}&-p_{\rm sig}\log\Omega\sum_\tau a_\tau
 -|\log\Omega|\sum_{\tau,a}|\lambda_{\tau a}-a_\tau|\\
 &-\sum_\tau\log\dim V^{\rm Weyl}_\tau
 -n_{\rm col}fp_{\rm sig}|\log\operatorname{rd}(E)| .
 \end{split}
\end{equation}
The scale width is \(s_{\rm sc}=4k\Delta\), where \(\Delta\) is the spacing parameter of the point, and \(ks_{\rm sc}\le b\). The first \(k-1\) entries in each row differ from the anchor by at most \(s_{\rm sc}+2r\). Each of the remaining \(m\) entries differs by at most \(8b\). Since \(r^2L_x\le frn_VL_x\le b\), \(m\ge2\) and \(L_x\ge1\), this gives
\[
 \sum_a|\lambda_{\tau a}-a_\tau|
 \le(k-1)(s_{\rm sc}+2r)+8mb\le10mb .
\]
For the row partition \(\nu_\tau\) take \(w_\tau=a_\tau-n_{\rm col}/2\). Then \(\sum_a|\nu_{\tau a}-w_\tau|=\sum_a|\lambda_{\tau a}-a_\tau|\). Each factor in the Weyl dimension product is at most \(1+\sum_a|\nu_{\tau a}-w_\tau|\), since its denominator is at least one and the difference of two entries is at most the sum of their absolute deviations from \(w_\tau\). Thus \(\dim V^{\rm Weyl}_\tau\le(1+10mb)^{\binom r2}\). By \(11m\le2fr\), \(\log(fr)\le L_x\) and \(\log b=L_x\), its logarithm is at most \(r^2(1+2L_x)\le3b\). Also \(n_{\rm col}\le n_V\) and \eqref{eq:p8c3-anchor} give \(n_{\rm col}fp_{\rm sig}|\log\operatorname{rd}(E)|\le2frn_VL_x\le2b\). The three errors in \eqref{eq:p8c3-terminal} therefore total at most \(30fmbL_x+3fb+2b\le35E_p\). The anchor bound in \ref{B8.24a} gives \(|M_{p,i}-p_{\rm sig}\log\Omega\sum_\tau a_\tau|\le C_AE_p\). Take \(C_{\rm term}=C_A+35+C_0\). Combine the two cases using \(b\le E_p\).
\end{proof}

\subsection*{Proof of \ref{B8.30}}\label{prf:B8.30}\phantomsection
\begin{proof}
For large \(x\), \(r\ge100\) and \(g_a\ge4096\) for \(a=0,1\). Apply \ref{B7.50} and \ref{B7.51} with the nontrivial additive characters and bad labels at both primes. The local shares \(u_a\) satisfy \(u_a\le x^{21}e^{-x^{e_z/6}}\le\delta/20\) eventually, hence
\begin{equation}\label{eq:p8c3-share-bad}
 u_0+u_1\le\delta .
\end{equation}
For a nonterminal cover let \(\pi:\Gamma\to G_0\) be the level quotient onto the product of the \(\GL_r(\F_{g_a})\) at the selected places. It has a section commuting with \(\ker\pi\). Let \(P\subset G_0\) be the product of upper unitriangular groups with its Whittaker character and \(H_a\subset G_0\) be unrestricted at \(g_a\) and upper unitriangular elsewhere. Let \(\widetilde P,\widetilde H_a\) be the inverse images. In the cover space \(V\) write \(V_\lambda\) for the pulled-back character space of \(\widetilde P\). Its dimension is the type rank. The bad rank is the dimension cut out in it by the complement of the good idempotent.

For \(\widetilde P\) and each \(\widetilde H_a\), use the orbit separation in the proof of \ref{B8.29} to choose a section \(u\) of shift \(c_0\). The matrix \((u((hg)^{-1}\cdot))_{h,g}\) has determinant a nonzero divisor in the section ring. Multiplication by the translates gives a joint injection equivariant for the regular permutation on the indices, with the common character twist absorbed in that action. For \(\widetilde H_a\), multiply from shifts \(0\) to \(c_0\) and \(c_0\) to \(2c_0\). This embeds the space in \(n_a\) regular copies of total dimension at most \(r_Y(G_s(2c_0))\). For \(\widetilde P\), multiply from \(-c_0\) to zero and apply the character projections. Joint injectivity gives the last bound below, and the endpoint rank comparison gives the middle one.
\begin{equation}\label{eq:p8c3-reg-copy}
 n_a|\widetilde H_a|\le r_Y(G_s(2c_0))
 \le\tfrac32 r_Y(G_s(-c_0))
 \le\tfrac32|\widetilde P|\dim V_\lambda .
\end{equation}
In \(G_0\) the bad projector at \(g_a\) times the character projector is supported on \(H_a\) and has identity coefficient \(u_a/|P|\). Lift it by the section and average over \(\ker\pi\). The identity coefficient on \(\widetilde H_a\) is \(u_a/|\widetilde P|\). The trace on the regular copies and \eqref{eq:p8c3-reg-copy} bound the rank on \(V_\lambda\) by \(\tfrac32u_a\dim V_\lambda\). The projectors commute, so the two ranks bound the loss from the common good part by their sum. At the terminal layer the freely permuted deformation-class summands give regular copies of total dimension \(|\widetilde P|\dim V_\lambda\). The same trace formula bounds the loss there by \((u_0+u_1)\dim V_\lambda\). Both bounds imply \eqref{eq:p8c3-bad-rank}.
\end{proof}

\subsection*{Proof of \ref{B8.31}}\label{prf:B8.31}\phantomsection
\begin{proof}
Write \(L_g=L\cap\operatorname{im}e\) and \(L_b=L\cap\ker e\), with ranks \(d_g,d_b\). The inclusion \(gL\subseteq L_g\oplus L_b\) gives
\begin{equation}\label{eq:p8c3-cut-ranks}
 \operatorname{rk}L=d_g+d_b .
\end{equation}
The map \(1-e\) embeds \(L/(L_g\oplus L_b)\) in \(g^{-1}L_b/L_b\), of order \(g^{[K:\Q]d_b}\). Orthogonality and the squared metric then give
\begin{equation}\label{eq:p8c3-cut-degree}
 (d_g+d_b)\mu(L)\le d_g\mu(L_g)+d_b(\mu_{\max}(M)+2\log g).
\end{equation}
The rank bound yields \(d_g>0\) and
\begin{equation}\label{eq:p8c3-cut-rank-ratio}
 d_b\le\frac{3\delta}{1-3\delta}d_g\le4\delta d_g .
\end{equation}
Thus \(d_g(\mu(L)-\mu(L_g))\le3d_b\mathcal Q\le12\delta d_g\mathcal Q\).
\end{proof}

\subsection*{Proof of \ref{B8.34a}}\label{prf:B8.34a}\phantomsection
\begin{proof}
Put \(n=g_0g_1p_{\rm de}p'_{\rm de}\). Let \(R/E\) be the ray class field of \(n\mathcal O_E\), and write \(h=[R:E]\), \(h_E=\#\operatorname{Cl}(E)\). Define
\[
 c=\prod_{i=0}^1|\GL_r(\F_{g_i})|^f,\qquad
 a=\operatorname{lcm}(c,3^4 5^4).
\]
Let \(A/\Q(j)\) be the elliptic curve of the field datum. Its \(j\)-invariant is that of \(\C/\mathcal O_K\), and it has good reduction away from \(d\). Put \(M_0=K(j,A[12])\), adjoining the coordinates of all the indicated torsion points. Let \(V_h\) be generated over \(E(\mu_h)\) by all roots \(z^h=y\) with \(y\in E^\times\) and \((y)\) an \(h\)-th power of a fractional ideal of \(E\). We first show that
\[
                         F\subset M_0 R\,\Q(\mu_a)V_h .
\]

The coefficient list contains the embedded host, the ray and elliptic-curve data, roots of unity, the prescribed character values, and the roots coming from an ideal-class basis at each signature. The elliptic-curve data include the twelfth roots of the Weierstrass discriminant. They lie in \(M_0\) by \cite[Th.~5.1 and Rem.~5.3]{fukuda-yoshikawa-2017}: the ground field has characteristic zero and all the \(3\)- and \(4\)-torsion is rational over \(M_0\). The level roots have orders dividing \(a\). In particular \(g_i(g_i-1)\) divides \(|\GL_r(\F_{g_i})|\), by the factor \(g_i^r-g_i^{r-1}\) in the order formula. This covers the values of the characters of \(\F_{g_i}^{\times}\) as well.

For the character values, write \(j_0\) for the integer specifying the norm twist \(N\mathfrak p^{-j_0/2}\), and \(n_\tau\) for the angular exponents, one at each complex place. The initial layer characters have even \(n_\tau\) and \(j_0=0\), the unramified character with half twist has \(n_\tau=j_0=1\), and the two finite-order characters have \(n_\tau=j_0=0\). They are tame at the primes of \(n\) and unramified elsewhere. At a prime \(\mathfrak p\) where the character is unramified, choose \(k_{\mathfrak p}\) generating \(\mathfrak p^h\) and congruent to \(1\) at the other primes of the modulus, to their prescribed powers. This is possible because the ray class group for the modulus with the \(\mathfrak p\)-part removed is a quotient of the group of order \(h\). The principal idèle relation gives the \(h\)-th power of the twisted value as
\[
 \prod_{\tau}
   \tau(k_{\mathfrak p})^{-n_\tau}
   \bigl(\tau(k_{\mathfrak p})\bar\tau(k_{\mathfrak p})\bigr)^{(n_\tau-j_0)/2}.
\]
Indeed the inverse angular character supplies the powers \(-n_\tau\), and
\((N\mathfrak p)^h=\prod_\tau|\tau(k_{\mathfrak p})|^2\).
All displayed exponents are integers. Since \(E/\Q\) is Galois, every factor is in the embedded \(E\), and its fractional ideal is an \(h\)-th power. The twisted value therefore lies in \(V_h\).

For a class of order \(k_i\) choose an ideal \(\mathfrak a_i\) and \(\gamma_i\) with \(\mathfrak a_i^{k_i}=(\gamma_i)\). The generators used at a signature are \(k_i\)-th roots of products of conjugates of \(\gamma_i\) to integer powers. Their radicands have \(k_i\)-th-power ideals. Since \(k_i\mid h_E\mid h\), each such root is an \(h\)-th root of an element with \(h\)-th-power ideal. This proves the inclusion.
\leanref{GlobalTheory.\allowbreak{}Mem.\allowbreak{}FieldsEx.\allowbreak{}listField\_\allowbreak{}le\_\allowbreak{}Om}{Development.\allowbreak{}AbelianVarietyTheory.\allowbreak{}Proofs.\allowbreak{}IwasawaDecomposition.\allowbreak{}ClassNumberDvdHstar}{370--372}
\leanref{GlobalTheory.\allowbreak{}Mem.\allowbreak{}FieldsEx.\allowbreak{}valTw\_\allowbreak{}mem\_\allowbreak{}kummerField}{Development.\allowbreak{}AbelianVarietyTheory.\allowbreak{}Proofs.\allowbreak{}IwasawaDecomposition.\allowbreak{}AlgebraicHeckeCharacterValueMemKummerField}{25--55}

Write \(u=\log\operatorname{rd}(E)\), \(H=\log h_E\), \(s=\log h\), and \(v=\log(2f)\). We give separate bounds for the log degree and log root discriminant of the four fields, in that order:
\[
\begin{array}{c|cc}
 &\log[\ \cdot\ :\Q]&\log\operatorname{rd}(\ \cdot\ )\\ \hline
 M_0&\log d+20&3\log d+20\\
 R&v+s&u+\log n+s\\
 \Q(\mu_a)&\log a&\log a\\
 V_h&v+(\operatorname{rank}\mathcal O_E^\times+2)s+H&
       u+(\operatorname{rank}\mathcal O_E^\times+2)s+H .
\end{array}
\]

For the first row put \(H_K=K(j)\). By \ref{B6.05f}(2) it is unramified over \(K\) and \([H_K:K]=h_K\), since the period lattice has the maximal order. We have \(h_K\le d\) by the quadratic ideal-class bound in \ref{B8.25c}. The action on \(A[12]\) embeds \(\Gal(M_0/H_K)\) in \(\GL_2(\Z/12)\), so its order \(b_0\) is at most \(12^4\). At primes away from \(6d\) the curve has good reduction and multiplication by \(12\) is finite étale on its abelian model. Hence \(M_0/H_K\) is unramified there.

We use the following local different calculation for the first two rows. For a Galois extension of number fields of degree \(b_1\), unramified outside the primes dividing an integer \(n_1>0\), the relative different divides \((n_1b_1)\). At a ramified prime remove the maximal unramified subextension. A uniformiser of the totally ramified part has Eisenstein polynomial of degree \(e\) and generates the local integers. In its derivative the nonzero terms have distinct valuations modulo \(e\), and the leading term has valuation \(e-1+v_{\rm top}(e)\). This bounds the different exponent by \(e-1+v_{\rm top}(b_1)\), since \(e\mid b_1\). Also \(v_{\rm top}(n_1)\ge e\) there. Taking norms proves that the log root discriminant increases by at most \(\log(n_1b_1)\). For \(M_0\) this gives \(\tfrac12\log d+\log(6db_0)\), and its degree is at most \(2db_0\). Both fit the first row. For \(R\), the ray class exact sequence gives
\(h\le h_E|(\mathcal O_E/n\mathcal O_E)^\times|\le h_E n^{2f}\).
Applying the different calculation with \(n_1=n\), \(b_1=h\) gives the second row. For the cyclotomic row, the degree is at most \(a\), and differentiating \(X^a-1\) at a primitive root shows that the different divides \((a)\).
\leanref{GlobalTheory.\allowbreak{}Mem.\allowbreak{}Fld.\allowbreak{}TorsionSize.\allowbreak{}sz\_\allowbreak{}discField}{Development.\allowbreak{}AbelianVarietyTheory.\allowbreak{}Proofs.\allowbreak{}IwasawaDecomposition.\allowbreak{}FieldDatumSizeLogDiscriminant}{183--252}
\leanref{GlobalTheory.\allowbreak{}Mem.\allowbreak{}Fld.\allowbreak{}RaySize.\allowbreak{}degOver\_\allowbreak{}le\_\allowbreak{}and\_\allowbreak{}sz}{Development.\allowbreak{}AbelianVarietyTheory.\allowbreak{}Proofs.\allowbreak{}IwasawaDecomposition.\allowbreak{}RayClassFieldDegreeLeAndRootDiscrBound}{19--45}

For \(V_h\), adjoin \(\mu_h\), one \(h\)-th root of each fundamental unit and of a generator of the torsion units, and one \(k_i\)-th root of each \(\gamma_i\) for a cyclic decomposition of the ideal class group. To see that these suffice, write \((y)=\mathfrak b^h\) and
\(\mathfrak b=(t)\prod_i\mathfrak a_i^{a_i}\).
Then \(k_i\mid ha_i\) for every \(i\), and \(y\) differs from
\(t^h\prod_i\gamma_i^{ha_i/k_i}\) by a unit. This constructs one root of \(y\), and \(\mu_h\) supplies the others. Adjoining a \(k\)-th root of an element whose ideal is a \(k\)-th power increases each of the two logarithms by at most \(\log k\). For the different bound, locally divide the root by a generator of the corresponding fractional ideal. The resulting element \(z\) is a unit and satisfies \(z^k=u_0\) for a unit \(u_0\). Its minimal polynomial divides \(X^k-u_0\) integrally. The derivative \(kz^{k-1}\) belongs to the different, by the derivative formula for the generated order and inclusion in the field different. Thus the different divides \((k)\). These arguments also apply over every intermediate field in the tower. Since \(\prod_i k_i=h_E\), they give the last row.
\leanref{GlobalTheory.\allowbreak{}FieldSize.\allowbreak{}sz\_\allowbreak{}kummerField}{Development.\allowbreak{}AbelianVarietyTheory.\allowbreak{}Proofs.\allowbreak{}IwasawaDecomposition.\allowbreak{}DegreeRootDiscrBoundKummerFieldOfPowerIdealRoots}{124--219}

It remains to bound the entries of the table. Every ideal class of \(E\) contains an integral ideal of norm at most \(X=\sqrt{|D_E|}\), by the Minkowski bound. Explicitly, apply the convex-body lattice bound to a product of \(f\) complex discs in an ideal \(I\), whose covolume is \(2^{-f}\sqrt{|D_E|}\,NI\); the resulting principal multiple of \(I^{-1}\) has norm at most \(X\). The convex-body bound follows by placing half the body in a fundamental parallelepiped modulo the lattice and taking the difference of two overlapping points. The number of ideals of norm at most \(X\) is at most \(X(1+\log X)^{2f}\): at each rational prime the Euler coefficients are bounded by those of \(\zeta(s)^{2f}\), and one can count the resulting tuples by summing \(X/(n_1\cdots n_{2f-1})\). Hence
\[
 H\le fu+2f\log(1+fu)\le10fL .
\]
Here \(1+fu\le3fL\) and \(\log(3L)\le3L\). The prime hypotheses give \(\log n\le20L+2\), whence \(0\le s\le H+2f\log n\le51fL\). The quadratic subfield gives \(\tfrac12\log d\le u\), by discriminant transitivity. Also \(v\le1+L\) and \(\operatorname{rank}\mathcal O_E^\times\le2f\). Finally \(|\GL_r(\F_g)|\le g^{r^2}\), and therefore
\[
             \log a\le fr^2(18L+2)+24\le19fr^2L .
\]
The last inequality uses \((L-2)fr^2\ge24\).

Log degrees and log root discriminants are subadditive under compositum, by the degree formula and divisibility of the relative different after base change as in \ref{B8.25c}. In each column of the table the sum excluding the cyclotomic row is at most
\[
       102f^2L+163fL+42L\le198f^2L .
\]
Indeed the combined coefficient of \(s\) is at most \(2f+3\), the remaining class-number term is \(H\), and the terms independent of \(s,H\) are bounded by \(42L\) in either column for \(L\ge4\). Use \(f\ge8\) in the last inequality. Passing to \(F\) can only decrease degree and root discriminant.
\leanref{GlobalTheory.\allowbreak{}FieldSize.\allowbreak{}om\_\allowbreak{}sizes\_\allowbreak{}le\_\allowbreak{}poR}{Development.\allowbreak{}AbelianVarietyTheory.\allowbreak{}Proofs.\allowbreak{}IwasawaDecomposition.\allowbreak{}OmegaFieldDegreeRootDiscrBoundedByRingPolynomial}{24--50}
\end{proof}

\subsection*{Proof of \ref{B8.34}}\label{prf:B8.34}\phantomsection
\begin{proof}
Use the ray-class-field and height-counting statements on the open divisor row of \ref{B8.08}. Intersect the class-sum row with the primitive, divisor, goodness and geometric rows to get the common witness. For a selection satisfying it, use the same component for the primitive-row comparison and class-sum identity. Use the primitive and divisor comparisons on the component chosen by the geometric row. The class-sum identity supplies the divisor comparison needed for the rational row maps. Use the embedding of the coefficient ring at the selected member. Its field identification gives precisely the field of the coefficient list in Lemma~\ref{B8.34a}. To apply that lemma put \(L=L_x=x^\varepsilon\). The host rank has \(f\ge8\) by the frame order condition. The rank bounds of \ref{B8.01k} give \(r\ge62\) eventually, and the polynomial bound on \(f\) in \eqref{eq:p8c3-host-size} gives \(\log f\le L\) eventually. Increase the threshold also to have \(L\ge4\). Clause~1 of \ref{B8.01k} gives \(g_i\le b^9\) and \(p_{\rm de},p'_{\rm de}<X'=\exp(x^{\varepsilon/3})\). By \ref{B2.46}, \(E\subset\Q(\mu_{Q_Sqd})\), so \(\operatorname{rd}(E)\le Q_Sqd\) by the cyclotomic different calculation in \ref{B8.34a}. The host size condition \((Q_Sqd)^{36}t^{150}\le\sqrt b\), with \(t\ge1\), gives \(\log\operatorname{rd}(E)\le L\). Lemma~\ref{B8.34a} now gives both bounds at this index. The series identification and these two bounds give the two arithmetic requirements in Definition~\ref{def:p8c3-admissible} at the selected index. The representation-theoretic inputs give the eventual row conditions at that same point. The seven family conditions of Definition~\ref{B8.01k} put the selected family in \(\mathcal F_{\rm cov}\). Take \(C_0\) from \ref{B8.24}. For \(x_0\), take the maximum of the covolume threshold and the thresholds for the point conditions and ties. Apply the layer inequality of \ref{B8.24} with \((C,C_H)\) to the reference indefinite realisation of the selected summand. Its right side is \(M_{p,i}+CE_p\).
\end{proof}

\section*{Proofs for Chapter~\ref{ch:B8-definite-local}: The definite side and local errors}

\subsection*{Proof of \ref{B8.35a}}\label{prf:B8.35a}\phantomsection
\begin{proof}
Use the parameter choice of the selected point in \ref{B8.05}; the family uses its size datum at every indexed \(x>x_1\). The host of \ref{B2.56} satisfies \(|D_1|\le L\) and \(r_d\le (Q_S/2)\sqrt{qd}\le Q_Sqd\). By \ref{B2.55} the logarithm of the latter conductor is at most \(2x^{\varepsilon/5}(\log x)^2\), which is at most \(L\) eventually. The inequality \(\log(4x)\le L\) also holds eventually for fixed \(\varepsilon>0\). The two primes of \ref{B2.51} lie in \([Y,2Y]\) with \(Y=(kb)^8\). Since \(k\) is polynomial in \(x\), \(\log Y=8L+8\log k\le9L=C_{YL}L\) eventually.

For the node-size threshold in \ref{B8.05} one may take \(c_h\ge1300+c_D\), \(c_D\ge0\). The rank bounds of \ref{B8.59} give \(t\le x^{401}\), \(N\le x^{402}\) eventually. Thus \(4N^3+2t+6\le x^{1300}\le h_0=\eta/m\) eventually. Both \(kS_c\) and \(m(t+1)\eta\) are polynomially bounded by \ref{B4.02b}, whereas \(\log w''=L-x^{\varepsilon/3}-\log n\). This proves the two \(w''\) inequalities. The second box of \ref{B3.50} gives \(b_s\le(1+R)b'\) with \(4R\le x^{-2}\). The spacing is the ceiling in \ref{B4.02b}, and the thick upper bound is the upper endpoint in \ref{B3.47}. All cutoffs depend on the fixed construction constants and \(\varepsilon\).
\end{proof}

\subsection*{Proof of \ref{B8.36}}\label{prf:B8.36}\phantomsection
\begin{proof}
Transport the spectral spaces, sides and face along the division ties. The column at \(w_U(u)\) is then the column of \(u\). Let \(q_0,\ldots,q_{n-1}\) be the theta steps and \(q_{\rm cpl}:\mathcal F_n\to\mathcal S\) the complement. The initial ramified inclusion \(f_0\) is an isometry, and point evaluation gives a map \(g:\mathcal S\to(K^{n_0})^\vee\) with
\begin{equation}\label{eq:B8-period-kernel}
 P_{w_U(u),\mu_0}(\phi,F)=g(q_{\rm cpl}q_{n-1}\cdots q_0 f_0\phi)(F).
\end{equation}
Take matrices of the arrows in bases framed by their Hermitian forms. Include \(f_0\) in the first arrow, which is the complement when \(n=0\). Take \(P_{\rm end}\) to be the matrix of \(g\).
\leanref{GlobalTheory.\allowbreak{}DefRowsChain.\allowbreak{}chainOf}{Development.\allowbreak{}AbelianVarietyTheory.\allowbreak{}Proofs.\allowbreak{}DivisorProofs.\allowbreak{}KerChainMatricesWithNormBounds}{172--192}
\leanref{GlobalTheory.\allowbreak{}ColFirstIso.\allowbreak{}colsIn\_\allowbreak{}isometry}{Development.\allowbreak{}AbelianVarietyTheory.\allowbreak{}Proofs.\allowbreak{}DivisorProofs.\allowbreak{}ColsInIsometry}{100--123}

Apply \ref{B5.07c} to the interlacing terminal tuple and the common determinant shift. The terminal parameter at \(\tau\) is \(\tau(\bar\varpi)\) with \(\varpi\bar\varpi=T_u\). Take \(A\) to be the log squared terminal factor and use the witnesses
\[
 c_\tau^*=(N-1)D_{1,\tau}\log(2a),\qquad
 \xi_\tau^*=\tfrac12N(N-1)\log(2a),\qquad e_{\rm ev}=C_{\rm ev}=0.
\]
The level-group average gives reciprocal stabilizer weights and fixes the reference vectors with divided Gram form \(H_{\mu_0}\). The bound of \ref{B5.07c} in these frames is
\begin{equation}\label{eq:B8-terminal-period}
 |v(g(\psi)(F))|^2\le e^A\|\psi\|_{\mathcal S,v}^2\|Y_v^{-1}v(F)\|_2^2.
\end{equation}
Density extends this to the complex coordinate spaces.
\leanref{GlobalTheory.\allowbreak{}ColRows.\allowbreak{}LastSideProof.\allowbreak{}lastSide\_\allowbreak{}of\_\allowbreak{}logT\_\allowbreak{}le}{Development.\allowbreak{}AbelianVarietyTheory.\allowbreak{}Proofs.\allowbreak{}DivisorProofs.\allowbreak{}LastMapHalfRecNumsOfLogLevelPrimeLe}{14--193}

Let \(C_K\) be the kernel matrix and coupling factor constant, and \(C_c,C_i,C_y,C_t\) the class-number, intermediate-index, coupling-coset and normalized theta norm constants, respectively. Set \(E_*=C_c+C_i+C_y\) and \(U_*=fm^2N^2\log(g_0g_1T)\). At one embedding of a common coefficient field, apply Cauchy--Schwarz over the input and output class sums, then over the input and output indices. For class counts \(h_i,h_{i+1}\) and index counts \(n_i,n_{i+1}\), the factors are \(\sqrt{h_i h_{i+1}}\) and \(\sqrt{n_i n_{i+1}}\). The coupled end uses the coupling coset count in the second factor. Thus
\begin{equation}\label{eq:B8-block-cost}
 M_{\ell,v}\le C_KU_*+\tfrac12\log(h_i h_{i+1})
 +\tfrac12\log(n_i n_{i+1})\le(C_K+2E_*)U_*.
\end{equation}
The same bound holds after averaging the embeddings above \(v\). The place weights have uniform mean over all coefficient embeddings, so \ref{B5.07x} gives \(\bar\omega_\tau=\log\Omega\). The normalized theta logs divide out the product \(g_\ell^0\) of the adjacent ranks. Their mean bound and the kernel size bounds of Definition~\ref{def:B8-column-comparisons}, applied to the volume, the degrees and the two primes, give
\begin{equation}\label{eq:B8-kernel-sizes-special}
 \bar\vartheta\le2\max(C_t,0)L,\quad
 \log(8+32p_\ell V)\le8L,\quad
 \log(2n_{\ell,\tau}+2)\le5L,\quad
 0\le g_\ast=\log(g_0g_1)\le(2\max(C_Y,0)+2)L.
\end{equation}
Take \(h_\ell=h'_\ell=0\). The squared factors at an embedding multiply, with log bounds given by \eqref{eq:B8-height-size}. If layer \(\ell\) has index \(j\) in the kernel calculation, assign it to step \(\iota(\ell)=\min(j,n)\). This collects the ordinary factors separately and puts the remaining factors into the complement. It gives
\begin{equation}\label{eq:B8-column-step}
 \|X_{\alpha+1,v}^{-1}v(K_\alpha)X_{\alpha,v}\|^2
 \le\exp\left(\sum_{\iota(\ell)=\alpha}s_{\ell,v}\right).
\end{equation}
\leanref{GlobalTheory.\allowbreak{}KernelHalfS.\allowbreak{}kerSide\_\allowbreak{}chainS}{Development.\allowbreak{}AbelianVarietyTheory.\allowbreak{}Proofs.\allowbreak{}DivisorProofs.\allowbreak{}KernelHalfChainNumbersAnyLevelPrime}{219--331}
\leanref{GlobalTheory.\allowbreak{}ColRows.\allowbreak{}FactorBounds}{Development.\allowbreak{}AbelianVarietyTheory.\allowbreak{}Definitions.\allowbreak{}Divisors.\allowbreak{}FactorBounds}{40--43}

With \(\Gamma=C_K+2E_*+2\max(C_t,0)+20+2\max(C_Y,0)\), the constants in \eqref{eq:B8-table-constants} are covered by
\begin{equation}\label{eq:B8-column-threshold}
 C_0=\max\{2(16+C_T),6\Gamma+1+2\Gamma^2,4\Gamma(C_T+\Gamma)\}.
\end{equation}
\leanref{GlobalTheory.\allowbreak{}ColRowsTopS.\allowbreak{}rowsBody\_\allowbreak{}colFace}{Development.\allowbreak{}AbelianVarietyTheory.\allowbreak{}Proofs.\allowbreak{}DivisorProofs.\allowbreak{}ColumnsMemberRowsOfLogLevelLe}{25--58}
\leanref{GlobalTheory.\allowbreak{}ColRows.\allowbreak{}exists\_\allowbreak{}htNums\_\allowbreak{}of\_\allowbreak{}kerSide\_\allowbreak{}lastSide}{Development.\allowbreak{}AbelianVarietyTheory.\allowbreak{}Proofs.\allowbreak{}DivisorProofs.\allowbreak{}ExistsHtNumsOfKerSideLastSide}{11--27}
Take one coefficient \(X_{\alpha,v}^{-1}z\) at each stage, with weight and phase one, the transition in \eqref{eq:B8-column-step}, and \(t_{\max}(\alpha,v)=\exp(\sum_{\iota(\ell)=\alpha}s_{\ell,v})\). Its logarithm gives \eqref{eq:B8-layer-bound}.
\leanref{GlobalTheory.\allowbreak{}DefRowsChain.\allowbreak{}rowsBody\_\allowbreak{}of\_\allowbreak{}framed}{Development.\allowbreak{}AbelianVarietyTheory.\allowbreak{}Proofs.\allowbreak{}DivisorProofs.\allowbreak{}KerChainMatricesWithNormBounds}{227--269}
\leanref{GlobalTheory.\allowbreak{}DefRowsCB.\allowbreak{}exists\_\allowbreak{}cbData}{Development.\allowbreak{}AbelianVarietyTheory.\allowbreak{}Proofs.\allowbreak{}DivisorProofs.\allowbreak{}ExistsCbData}{13--56}

For \(0\ne z\in I_u\), the lattice \(N'_u\) contains \(zN_u^0\), and the two span the same \(K\)-space. Noetherianity gives finite generation, and transport by \(e_u\) gives the required full lattice of positive dimension.
\leanref{GlobalTheory.\allowbreak{}ColRows.\allowbreak{}nuLat\_\allowbreak{}comap\_\allowbreak{}isLattice}{Development.\allowbreak{}AbelianVarietyTheory.\allowbreak{}Proofs.\allowbreak{}DivisorProofs.\allowbreak{}IsLatticeOfEqSpanFracSmul}{52--60}
\end{proof}

\subsection*{Proof of \ref{B8.37}}\label{prf:B8.37}\phantomsection
\begin{proof}
Linkage supplies point evaluation and the size implication of \ref{B8.36}. Eventually \(m\ge4,\ 1\le f\le N(m),\ 3N(m)\le b,\ 24x\le b\), and the word weights satisfy \(a_{\max}<6xb\). The top window \eqref{eq:p3c5-top-window} gives $a_{\max}\le2xb$. For the selected weights, the common entry bounds are $-(2b+b')\le\lambda_{j,\tau,h}\le2xb+(b')^2/64$. On a marked type the negative of the last initial weight is between zero and $2b'$. In the maximum defining $W_{u,\tau}$, each retained marked entry has absolute value at most $2b'$ and the subtracted shift is at most $b'$, so $W_{u,\tau}\le3b'$. The size bounds also give $N_m\le b'\le b$. Since $T_u\ge b^{m/2}$, Lemma~\ref{B5.05b} supplies the level conditions in \ref{B5.05a}.
\leanref{Formalism.\allowbreak{}MainNotions.\allowbreak{}GlobalTheoryTheory.\allowbreak{}Alc.\allowbreak{}entry\_\allowbreak{}bd}{Formalism.\allowbreak{}SharpRestriction.\allowbreak{}GlobalSchurBounds.\allowbreak{}AlcoveOf}{65--122}
\leanref{Formalism.\allowbreak{}MainNotions.\allowbreak{}J.\allowbreak{}Wt\_\allowbreak{}top\_\allowbreak{}le}{Formalism.\allowbreak{}SharpRestriction.\allowbreak{}AdmissibleParameterJoins.\allowbreak{}DIVfromGLo}{112--138} With \(H=T_u-1\) and \(d_M=\dim_\C M^0_{u,\C}\), it gives
\begin{equation}\label{eq:B8-ram-rank}
 d_u=d_M(1-1/H)^f,\qquad
 \mu_{\rm ram}=\mu_0-\frac{f\log H}{H-1},\qquad \mu_{\max,\rm ram}\le\mu_{\max,0}.
\end{equation}
\leanref{Formalism.\allowbreak{}L2.\allowbreak{}GlobalTheoryTheory.\allowbreak{}Places.\allowbreak{}settingF\_\allowbreak{}of}{Formalism.\allowbreak{}LocalToGlobal.\allowbreak{}GlobalTheory.\allowbreak{}PlacesP}{759--777}
\leanref{Formalism.\allowbreak{}MainNotions.\allowbreak{}J.\allowbreak{}levelPrime\_\allowbreak{}setting\_\allowbreak{}of\_\allowbreak{}sizes}{Formalism.\allowbreak{}SharpRestriction.\allowbreak{}AdmissibleParameterJoins.\allowbreak{}DIVfromGLo}{249--293}
The local inclusions in \ref{B5.05a} put \(N_u^0\) in the ramified lattice. At a noninert place off \(T_u\) the localizations agree. Clear denominators for a finite set of generators of the ramified lattice at that place to obtain a nonzero multiple contained in \(N_u^0\). Multiplication by \(I_u\) preserves fullness. The lower rank bound in the selected-family design gives \(d_M>0\), and \(T_u>2\) gives \(d_u>0\).
\leanref{Formalism.\allowbreak{}MainNotions.\allowbreak{}Lnk.\allowbreak{}nuCommens\_\allowbreak{}of}{Formalism.\allowbreak{}SharpRestriction.\allowbreak{}GlobalSchurBounds.\allowbreak{}LinkPk}{147--162}
\leanref{Formalism.\allowbreak{}MainNotions.\allowbreak{}Wire.\allowbreak{}family\_\allowbreak{}theorems\_\allowbreak{}two\_\allowbreak{}and\_\allowbreak{}five}{Formalism.\allowbreak{}SharpRestriction.\allowbreak{}GlobalSchurBounds.\allowbreak{}LogXLowerFamilyOfTheorems}{77--81}
\leanref{Formalism.\allowbreak{}MainNotions.\allowbreak{}muRank\_\allowbreak{}items}{Formalism.\allowbreak{}SharpRestriction.\allowbreak{}GlobalSchurBounds.\allowbreak{}LinkMin}{739--745}

      Lemma~\ref{B8.35a} supplies the first condition of \eqref{eq:B8-five-size} on this member, with \(C_Y=C_{YL}\). Eventually \(X'\ge8,\ x\ge9,\ x^2\le b\), hence \(4a<b^2\le T_u\). The upper level bound is \ref{B8.51}. Use the middle threshold \eqref{eq:B8-column-threshold} with \(C_Y,C_T\) from the zero-indexed outer entries \(6,13\). Take the nonnegative part of entry \(13\). The maximum of the fullness, rank and size cutoffs works in \ref{B8.36}.
\leanref{GlobalTheory.\allowbreak{}CoreHt.\allowbreak{}famRowsSizes}{Development.\allowbreak{}AbelianVarietyTheory.\allowbreak{}Proofs.\allowbreak{}DivisorProofs.\allowbreak{}FamRowsSizes}{13--68}
\end{proof}

\subsection*{Proof of \ref{B8.38}}\label{prf:B8.38}\phantomsection
\begin{proof}
The unitary input uses the uniformly moderate-growth case of Miller--Schmid's cuspidal decay theorem \cite[Theorem~A]{miller-2012}. If \(H_{\rm quot}\) is the infimum of the height over rational left translates, then \(\sup_g|\phi(g)|H_{\rm quot}(g)^N<\infty\) for each \(N\in\N\). For a nonzero multilinear intertwiner on finite vectors, fix a tuple with nonzero image. Fixing all but one entry gives reference embeddings with the lower Sobolev bounds of \ref{B8.38a}. Extend each slice by that proposition. When one finite argument tends to a smooth argument, the slices in the other argument converge pointwise by the preceding extension. Uniform boundedness gives continuity of the limit slice. Iterate over the positions. Equivariance passes by density. The reference embeddings have continuous inverses on closed images, so continuous seminorms there are dominated by finite sums of the target seminorms. In smooth coordinates these give the required bounds by finitely many derived words. The zero map extends by zero.
\leanref{GlobalTheory.\allowbreak{}FactorExtension.\allowbreak{}extend\_\allowbreak{}positions}{Development.\allowbreak{}AbelianVarietyTheory.\allowbreak{}Proofs.\allowbreak{}TwistedOperators.\allowbreak{}MultilinearGKMapExtendsToSmoothVectors}{69--259}

Transport the primitive and divisor data and use the seven member duties of \ref{B8.02}. The resulting family belongs to the covolume class \(\mathcal F_{\rm cov}\) of Definition~\ref{B8.24}. This class imposes the page and window conditions, small first modulus, tame distinctness, the chosen wall, boxes, the \(q,d\) rows, \(C_{YL}\le9\), smooth stabilizer orders, twist goodness, the selected-family design with constants \(3,C_H'\) for some \(C_H'\), and level and stabilizer size. It requires eventually at every member and word
\[
 |D\Gamma_u|\mid M(2fN(m)),\qquad
 M(n)=\prod_{\substack{p\le n+1\\p\ {\rm prime}}}
 p^{\sum_{a\ge0}\lfloor n/(p^a(p-1))\rfloor}.
\]
At a member point \(p\), write \(\mathbf N_p(C,C_H)\) for its numerical datum and \(S_{p,0}\) for its reference realization. Under the family division estimate \eqref{eq:B8-divinN}, the third row is the following closed assertion on \(\mathcal F_{\rm cov}\).
\begin{equation}\label{eq:B8-surj-quantifiers}
\begin{gathered}
 \exists C_0\ge0,\ \exists x_0\in\R,\ \forall x>x_1,\ x\ge x_0,\
 \forall i\in I_x,\ \forall p\text{ at }(F_m,x,i),\\
 p\in\mathcal P\cap\mathcal R
 \ \Longrightarrow\ 
 \forall C\ge C_0,\quad \text{\eqref{eq:B8-paired-index}--\eqref{eq:B8-paired-error} hold on }S_{p,0}\text{ at }\mathbf N_p(C,C_H).
\end{gathered}
\end{equation}
Here \(\mathcal P\) and \(\mathcal R\) are the sets of member points in Definition~\ref{def:B8-reference-points}.  Linkage supplies membership in \(\mathcal R\) and transports the three bounds of Proposition~\ref{B8.38b}.
\leanref{GlobalTheory.\allowbreak{}BigOn.\allowbreak{}SurjectiveOntoX}{Development.\allowbreak{}AutomorphicArithmetic.\allowbreak{}ThetaCorrespondence.\allowbreak{}ThetaPairLifts.\allowbreak{}ThirdPlaceSeries.\allowbreak{}ExportSt}{47--48}
\leanref{GlobalTheory.\allowbreak{}BigTop.\allowbreak{}FamHyp}{Development.\allowbreak{}AbelianVarietyTheory.\allowbreak{}Definitions.\allowbreak{}Divisors.\allowbreak{}TupleReadBaseTupleEqFamily}{503--507}
\leanref{Formalism.\allowbreak{}HermitianLatticeCounting.\allowbreak{}S6.\allowbreak{}minkowskiM}{Formalism.\allowbreak{}FlagUnipotentLattices.\allowbreak{}MemOfOrthogonalAnisotropic}{498--500}
\leanref{GlobalTheory.\allowbreak{}BigOn.\allowbreak{}More}{Development.\allowbreak{}AbelianVarietyTheory.\allowbreak{}Statements.\allowbreak{}PlaceShiftStatements.\allowbreak{}SingleIndexPlacesAreTwisted}{82--82}
\leanref{GlobalTheory.\allowbreak{}BigTop.\allowbreak{}EventuallyAtPoints}{Development.\allowbreak{}AutomorphicArithmetic.\allowbreak{}ThetaCorrespondence.\allowbreak{}ThetaPairLifts.\allowbreak{}ThirdPlaceSeries.\allowbreak{}Realised}{117--118}

The fourth row is closed on \(\mathcal F_{\rm cov}\). Eventually, at every \(p\in\mathcal P\) and for every \(C\in\R\), the logarithm of the selected weight dominates the layer lower bound \(\ell_{6,j}(p;C,C_H)\) in the numerical datum
\[
 \ell_{6,j}(p;C,C_H)\le\log\lambda_{j,\mu_0,b_0,v}
 \qquad\text{for all }j,b_0,v .
\]
The eight rows in order are the primitive isometry and passage, covolume, surjectivity, lower weight, functional, definite, multiplicity and \(\kappa\) rows. Rows \(3,5\) use the division estimate \eqref{eq:B8-divinN}, row \(4\) uses the leaf hypothesis in transport, and row \(6\) is \ref{B8.37}. Apply Definition~\ref{B8.39} to their conjunction and restrict to \(W_3(\mathcal A)\). Take the maximum of the first six \(C\)-thresholds, with the first nonnegative, then the maximum of the eight \(x\)-cutoffs and the selection cutoff.
\leanref{GlobalTheory.\allowbreak{}BigTop.\allowbreak{}LambdaLowerBound₂}{Development.\allowbreak{}AutomorphicArithmetic.\allowbreak{}ThetaCorrespondence.\allowbreak{}ThetaPairLifts.\allowbreak{}ThirdPlaceSeries.\allowbreak{}Realised}{126--126}
\leanref{Formalism.\allowbreak{}L2.\allowbreak{}Indef.\allowbreak{}IndefData.\allowbreak{}L6}{Formalism.\allowbreak{}LocalToGlobal.\allowbreak{}IndefiniteForms.\allowbreak{}IndefSide}{325--326}
\leanref{GlobalTheory.\allowbreak{}Fc.\allowbreak{}S0.\allowbreak{}ontoX\_\allowbreak{}and\_\allowbreak{}lamLow\_\allowbreak{}of\_\allowbreak{}unitaryGroupHypotheses'}{Development.\allowbreak{}AutomorphicArithmetic.\allowbreak{}ThetaCorrespondence.\allowbreak{}ThetaPairLifts.\allowbreak{}ThirdPlaceSeries.\allowbreak{}SurjectiveAndLowWeightOfUnramifiedExponents}{16--29}
\end{proof}

\subsection*{Proof of \ref{B8.38b}}\label{prf:B8.38b}\phantomsection
\begin{proof}
Fix an auxiliary \(C_*\ge1\) above the input thresholds and take the maximum of the associated cutoffs. The following estimates are at the resulting point record.

\emph{Remainders for the two runs.} In layer \(j\) let \(n_j\) be the ambient cusp-form dimension, \(B\) the ambient lattice and \(\Theta\) the initial theta lattice of the run. At a base place \(\mathfrak p\) round the prescribed depth times the theta power up to \(E_{\mathfrak p j}\). The two-run input bounds the index on every subspace by \(E_{\mathfrak p j}\dim V+r_{\mathfrak p j}\). On extension take \(E_{vj}=e(v/\mathfrak p)E_{\mathfrak p j}\) and \(r_{a,vj}=e(v/\mathfrak p)r_{\mathfrak p j}\) for run \(a\). Then for \(a=1,2\)
\begin{equation}\label{eq:B8-theta-remainders}
 [B_V:\Theta_V]_v\le E_{vj}\dim V+r_{a,vj},\qquad
 r_{a,vj}\ge0,\qquad
 \sum_{v\in S}\ell_v r_{a,vj}\le n_j fmb\,\delta_2
 \quad(0\le\delta_2\le1)
\end{equation}
for every finite place set \(S\), by \ref{B4.52}. Use run \(2\) at its distinguished characteristic and run \(1\) elsewhere. The remainders \(r_{vj}\) of this choice cost at most \(2n_jfmb\).
\leanref{GlobalTheory.\allowbreak{}BigOn.\allowbreak{}hasIndexBound\_\allowbreak{}of\_\allowbreak{}runs}{Development.\allowbreak{}AbelianVarietyTheory.\allowbreak{}Proofs.\allowbreak{}DivisorProofs.\allowbreak{}HasIndexBoundOfRuns}{129--191}
\leanref{GlobalTheory.\allowbreak{}TraceClassRank.\allowbreak{}on\_\allowbreak{}core\_\allowbreak{}atPlaces}{Development.\allowbreak{}AbelianVarietyTheory.\allowbreak{}Proofs.\allowbreak{}FoundationProofs.\allowbreak{}TraceLatticeLeThetaLatticeAtPlace}{233--243}

Apply \ref{B4.54a} with mirabolic \(P\), unipotent \(U\), retained projector \(e\), and \(s=[P:U]\). The product character meets the character condition in a factor in which a coset lies outside \(U\). The normalized stable theta lattice is contained in \(\sigma_{vj}\Theta_K\), and the integral character sum puts \(|U|(\Theta_K)_W\) into \(\Theta_0\) in the notation of that lemma. The twisted projection adds a normalized error \(\xi_{vj}\ge0\) with \(\sum_j\xi_{vj}\le p_v^{-2}\). This follows by restricting the two-character goodness budget to the chosen character and the injective layer labels. To identify the lattice in the index, let \(\phi_j\) be the coordinate embedding into the ambient cusp-form space, let \(Z\) be its image, and let \(B\) be the ambient integral lattice. The coordinate identification sends \(M_j\) to \(B\cap Z\). At \(v=w,\bar w\) it sends \(A_{j,t_w}\) to \(\theta'\pi e_X M_t\), where \(M_t\) is the class lattice, \(e_X\) the source projector, \(\pi\) the map to the chosen representatives, and \(\theta'\) the ambient theta map. For \(x\in e_XM_t\), \(x-\pi x\in\ker\theta'\). Since \(\theta'e_X=e\theta'\), the image is exactly \(e\theta'M_t\). Its localization contains the comparison lattice \(X_3\) of \ref{B4.54a}. Thus the index bounded there is the index of the same image that occurs in the determinant formula of \ref{B5.08a}. Write \(u_v=v(|U|)/e_v\) and \(d_{vj}=(E_{vj}-v(\sigma_{vj}))/e_v\). Using \eqref{eq:B8-theta-remainders} in \ref{B4.54a} gives, for \(v=w,\bar w\),
\begin{equation}\label{eq:B8-single-index}
 [M_j:A_{j,t_w}]_v\le e_v\left((d_{vj}+u_v)D_R+
 \frac{r_{vj}/e_v+n_j u_v}{s}+\xi_{vj}\right).
\end{equation}
Here \(t_{\bar w}=t_w\) by \ref{B4.36a}. If \(d_{0j}\) is the retained dimension before rank loss, the rank and share rows give \((1-y)n_j\le |U|d_{0j}\), \((1-y)|U|\le s\), and \((1-\delta_j)d_{0j}\le D_R\) for \(y=1/(4(fk)^3)\), \(0\le\delta_j\le1/4\). In particular \(n_j/s\le3D_R\). At the theta primes \(u_v=fr(r-1)/2\) with \(r=k+m-1\), and elsewhere \(u_v=0\).
\leanref{GlobalTheory.\allowbreak{}TwistIndex.\allowbreak{}exists\_\allowbreak{}leafTW1\_\allowbreak{}of\_\allowbreak{}twGoodAt}{Development.\allowbreak{}AbelianVarietyTheory.\allowbreak{}Proofs.\allowbreak{}FoundationProofs.\allowbreak{}LeafTW1OfAllConjugatesGood}{439--465}
\leanref{GlobalTheory.\allowbreak{}OnAtPlacesBase.\allowbreak{}on\_\allowbreak{}base\_\allowbreak{}atPlaces}{Development.\allowbreak{}AbelianVarietyTheory.\allowbreak{}Proofs.\allowbreak{}DivisorProofs.\allowbreak{}TwistedPlacesBoundChainTwoRunChain}{43--63}

\emph{Depths on ordinary and exceptional primes.} Let \(w\) be twisted and let \(p=p_w\) lie outside the small primes and \(q,d,g_0,g_1\). For \(p\le Y\), both \(d_{wj}\) and \(d_{\bar wj}\) are at most \(F_{0,j,w}+a_j(p)+1\). For \(p>Y\), both depths vanish and \(F_{0,j,w}\ge0\). The nonnegative errors \(\vartheta_p\) and \(\alpha_p\) in \(a_j(p)\) satisfy, for every finite set \(S\) of primes,
\begin{equation}\label{eq:B8-depth-sums}
 \sum_{p\in S}\vartheta_p\log p\le C_*/k,\qquad
 T\sum_{p\in S}\alpha_p\log p\le C_*mb .
\end{equation}
Here \(T\) is the theta power, \(Y\) the depth cutoff, and, with \(W_j\) the layer weight and \(S_e=S_c+\eta_{\rm sz}\),
\[
 A_1=C_*(mW_j\log k+fkS_e),\quad A_2=C_*kW_j,\quad
 a_j(p)=A_1/p+A_2\vartheta_p+T\alpha_p+\mathbf1_{p>X''}A_2/p^2 .
\]
Set \(s_B=W_j+fS_e,\ s_C=C_*m(2+\log_2 Q_S)s_B\), and \(\beta_p=\mathbf1_{p=d}fkb'\). The depth comparisons at \(w,\bar w\) bound their combined excess over \(2F_{0,j,w}\) by \(H_j(p)\). The bounds include the two unipotent terms in \eqref{eq:B8-single-index}.
\begin{equation}\label{eq:B8-paired-depth}
\begin{array}{c|l|l}
 p&\text{depth and base term}&H_j(p)\\ \hline
 g_0,g_1 &d_{vj}\le 2T-\nu_{vj},\
 \nu_{wj}+\nu_{\bar wj}=-2F_{0,j,w}&4T+4fr^2\\
 \mathrm{small}&d_{vj}\le s_C+T+1,\ F_{0,j,w}=0&2s_C+2T+2\\
 q,d&d_{vj}\le s_B+T+1,\ -2\beta_p\le F_{0,j,w}\le0&2(s_B+T+1)+4\beta_p\\
 \text{other }p\le Y&d_{vj}\le F_{0,j,w}+a_j(p)+1&2a_j(p)+2\\
 \text{other }p>Y&d_{vj}=0,\ F_{0,j,w}\ge0&0 .
\end{array}
\end{equation}
In the theta row let \(n_j^{\rm lev}\) denote the layer's level normalization factor, and put \(\nu_{vj}=v(n_j^{\rm lev})/e_v\). The level normalization gives \(v(\sigma_{vj})/e_v=\nu_{vj}+s_{vj}\) with \(s_{vj}\ge0\). The two valuations of the level factor add to the sum of the positive and negative Fock degrees, so \(\nu_{wj}+\nu_{\bar wj}=-2F_{0,j,w}\) by \ref{B4.02}. For the depth before subtracting \(v(\sigma_{vj})\), the unrounded normalized bound is \(T+TD_L/(p-1)\). Here \(TD_L\le D_\kappa<(p-2)/2\). Rounding on the lower field adds at most the reciprocal of the ramification index there, which is at most \(1/(p-1)\) at a type prime. Hence \(E_{vj}/e_v\le T+1\le2T\). Subtraction gives \(d_{vj}\le2T-\nu_{vj}\). In particular the sum of the two depths is at most \(4T+2F_{0,j,w}\). The two unipotent contributions are at most \(8u_w\le4fr^2\). The small-prime comparison uses the split colength and Hodge--Tate inclusion. The ordinary comparison uses \(x\ge3,\ kb'\le b,\ C_{nW}\ge10,\ x^{e_z/5}\ge3\).
\leanref{GlobalTheory.\allowbreak{}RunDepth.\allowbreak{}leafDepth\_\allowbreak{}all\_\allowbreak{}run}{Development.\allowbreak{}AbelianVarietyTheory.\allowbreak{}Proofs.\allowbreak{}JetHarmonics.\allowbreak{}LeafDepthAllRun}{13--21}
\leanref{GlobalTheory.\allowbreak{}BigOn.\allowbreak{}leafDepth\_\allowbreak{}all\_\allowbreak{}record}{Development.\allowbreak{}AbelianVarietyTheory.\allowbreak{}Proofs.\allowbreak{}DivisorProofs.\allowbreak{}LeafDepthAllRecord}{17--26}
\leanref{GlobalTheory.\allowbreak{}BigOn.\allowbreak{}geoSt\_\allowbreak{}of\_\allowbreak{}pageSt\_\allowbreak{}alt}{Development.\allowbreak{}AutomorphicArithmetic.\allowbreak{}ThetaCorrespondence.\allowbreak{}ThetaPairLifts.\allowbreak{}ThirdPlaceSeries.\allowbreak{}GeoStOfPageSt}{21--46}

Put \(b_{jw}=H_j(p_w)+(r_{wj}/e_w+r_{\bar wj}/e_{\bar w})/(sD_R)\). The twist errors summed over both places and all layers cost at most \(2e_w/(D_Rp_w^2)\le c_w\). Thus
\begin{equation}\label{eq:B8-paired-aux-index}
 D_R^{-1}\sum_j i_{jw}\le\sum_j(2e_wF_{0,j,w}+e_w b_{jw})+c_w .
\end{equation}
The remainders cost at most \(4n_jfmb/(sD_R)\le12fmb\), since \(e_w,\ell_w\) are preserved by conjugation. Write \(Z_0=fmb\). The page bounds \(W_j\le2fb,\ S_e\le b,\ T^2\le b,\ kb'\le b,\ \log(qd)\le m\) give \(14Z_0\) at \(q,d\). At the small primes \((2+\log_2 Q_S)\sum_{\rm small}\log p\le11L\) and \(fm\ge352\) give \((66C_*+4)Z_0L\). At \(g_0,g_1\), \((frL)^{12}\le b\), \(f\ge8,\ r\ge9,\ L\ge10\), and \(\log(g_0g_1)\le18L\) give \(4b\), using \(36TL,36fr^2L\le b\).
\leanref{Formalism.\allowbreak{}L2.\allowbreak{}Indef.\allowbreak{}Tree.\allowbreak{}OnPage.\allowbreak{}leafBad\_\allowbreak{}of\_\allowbreak{}page}{Formalism.\allowbreak{}LocalToGlobal.\allowbreak{}IndefiniteForms.\allowbreak{}IndefiniteTree.\allowbreak{}ONPage}{226--277}

The prime estimates \(\sum_{p\le Y}\log(p)/p<\log Y\) and \(\sum_{p\le t}\log p<1.01624t\) are \cite[(3.24), (3.32)]{rosser-schoenfeld-1962}. Partial summation gives \(\sum_{p>X''}\log(p)/p^2<2.04/X''\). By \eqref{eq:B8-depth-sums} the ordinary cost is at most
\begin{equation}\label{eq:B8-paired-ord}
 2(A_1\log Y+A_2 C_*/k+C_*mb+2.04 A_2/X''+Y\log4).
\end{equation}
Use \(Y\le mb,\ \log Y\le2L,\ \log k\le L,\ X''\ge k,\ kS_e L\le mb\). The first term costs at most \(8C_*Z_0L^2+4C_*Z_0\), and the remaining terms at most \((4C_*^2+10.16C_*+2.78)Z_0\).

\emph{Exponents of the cuts.} The exponent input compares \(J_{0,j,w}\) to the model exponent. Let \(S_j\) be the distinguished set of primes and \(T_{2,j}\) the thick positions in the weight data of layer \(j\). Its model exponent \(J_j(p)\) equals \(J_A(p)\) when \(p\le l_J\), when \(p\) is excluded, or when \(p\in S_j\). Otherwise it is the coarse weight bound. The small primes are below the cutoff, and the extra exclusions are \(q,d\) and the two recursion primes. The recursion gives \(|S_j|\le5\) and \(|T_{2,j}|\le f\). The exceptional intervals cost \(30|T_{2,j}|mbL^2\) and the tail costs \(2mb\). Hence
\begin{equation}\label{eq:B8-paired-exponent}
 0\le J_{0,j,w}\le J_j(p_w),\qquad
 \sum_{p\in S}J_j(p)\log p\le30|T_{2,j}|mbL^2+2mb
\end{equation}
for every finite prime set \(S\).
\leanref{GlobalTheory.\allowbreak{}ExpTie.\allowbreak{}Asm.\allowbreak{}hasExpBound\_\allowbreak{}at}{Development.\allowbreak{}AbelianVarietyTheory.\allowbreak{}Proofs.\allowbreak{}DivisorProofs.\allowbreak{}ExpansionOfNumericsModel}{50--77}
\leanref{UpperTheory.\allowbreak{}ExcSet.\allowbreak{}sum\_\allowbreak{}exponent\_\allowbreak{}log\_\allowbreak{}le}{Development.\allowbreak{}AutomorphicArithmetic.\allowbreak{}PrimeDiscriminants.\allowbreak{}ExcSet}{50--73}
\leanref{GlobalTheory.\allowbreak{}sum\_\allowbreak{}exponent\_\allowbreak{}log\_\allowbreak{}le\_\allowbreak{}of\_\allowbreak{}cut}{Development.\allowbreak{}AbelianVarietyTheory.\allowbreak{}Proofs.\allowbreak{}DivisorProofs.\allowbreak{}SumExponentLogLeOfCut}{12--30}

By \ref{B4.52}, the exponent and remainder costs are at most \(32Z_0L^2\) and \(12Z_0\). The preceding prime estimates then bound \(\sum_{w\in\mathrm{Tw}}\ell_we_w(b_{jw}+J_{0,j,w})\) by \(E=(4C_*^2+25C_*+C_b+44)L(B_0+N_c)\), where \(C_b=14+(66C_*+4)+4\).
\leanref{GlobalTheory.\allowbreak{}BigOn.\allowbreak{}pageSt\_\allowbreak{}alt}{Development.\allowbreak{}AutomorphicArithmetic.\allowbreak{}ThetaCorrespondence.\allowbreak{}ThetaPairLifts.\allowbreak{}ThirdPlaceSeries.\allowbreak{}PageSt}{21--53}
By \eqref{eq:B8-paired-aux-index}, \(R_w\le\sum_j e_w b_{jw}\). The positive-part formula for \(R_w\) gives \eqref{eq:B8-paired-index} and \eqref{eq:B8-paired-positive}. Averaging the bounds for \(b_{jw}\) gives \(\sum_w\ell_w e'_{jw}\le E\). At each layer the exponent cost is at most \(E\). Thus \eqref{eq:B8-paired-error} holds with \(C\ge 2(4C_*^2+25C_*+C_b+44)\) on the interval fixed for \(C_*\).
\end{proof}

\subsection*{Proof of \ref{B8.38a}}\label{prf:B8.38a}\phantomsection
\begin{proof}
Schwartz integration from \eqref{eq:B8-moderate-growth} and the Casselman--Wallach globalization theorem give the smooth extension \cite[\S2.5, Theorem~10.6 and Proposition~11.2]{bernstein-2014}. The unitary-source extension must also satisfy the explicit lower Sobolev bounds.
\leanref{PositivityTheory.\allowbreak{}CW.\allowbreak{}lift\_\allowbreak{}of\_\allowbreak{}steps}{Development.\allowbreak{}WeilBound.\allowbreak{}CoreStatements.\allowbreak{}PublishedTheorems.\allowbreak{}CasselmanWallach.\allowbreak{}CasselmanWallachAdmissible}{38--59}

For a defining seminorm \(p\) nonzero on \(T(V)\), let \(Q(z)=\sup_g h(g)^{-N}p(\pi(g)z)\) with the exponent of \(p\). Then \(Q\) is continuous, \(K\)-invariant and \(Q(\pi(g)z)\le h(g)^NQ(z)\). Convolve the dense orbit span in \(H\) with compactly supported smooth functions and use an approximate identity. The span of each nonzero orbit in \(V\) is then dense in \(H^\infty\). Thus \(P=Q\circ T\) is a norm on \(V\).
\leanref{PositivityTheory.\allowbreak{}CW.\allowbreak{}e1a}{Development.\allowbreak{}WeilBound.\allowbreak{}CoreStatements.\allowbreak{}PublishedTheorems.\allowbreak{}CasselmanWallach.\allowbreak{}ExistsFunctional}{27--80}
\leanref{PositivityTheory.\allowbreak{}CW.\allowbreak{}e1b}{Development.\allowbreak{}WeilBound.\allowbreak{}CoreStatements.\allowbreak{}PublishedTheorems.\allowbreak{}CasselmanWallach.\allowbreak{}ExponentialKSet}{38--68}

For each \(a\ne0\) in a compact Casimir eigenspace choose \(\lambda\) by Hahn--Banach with \(|\lambda|\le P\) and \(\lambda(a)=P(a)\). The \(K\)-average \(B(x,y)\) of \(\overline{\lambda(\sigma(k)x)}\lambda(\sigma(k)y)\) satisfies \(|B(x,y)|\le P(x)P(y)\) and \(B(a,a)>0\). The compact Casimir is symmetric for \(B\). The functionals \(B(a,-)\) span the dual of each finite-dimensional eigenspace and vanish on the other eigenspaces. Decompose \(\eta\in V\) and extend by density to obtain \(|\langle\eta,w\rangle|\le c_\eta Q(Tw)\) on \(H^\infty\).
\leanref{PositivityTheory.\allowbreak{}CW.\allowbreak{}e2c}{Development.\allowbreak{}WeilBound.\allowbreak{}CoreStatements.\allowbreak{}PublishedTheorems.\allowbreak{}CasselmanWallach.\allowbreak{}SeminormApproximation}{46--119}
\leanref{PositivityTheory.\allowbreak{}CW.\allowbreak{}c5}{Development.\allowbreak{}WeilBound.\allowbreak{}CoreStatements.\allowbreak{}PublishedTheorems.\allowbreak{}CasselmanWallach.\allowbreak{}Imports}{11--52}

For a finite generating set \(I\subset V\) the map \((f_e)\mapsto\sum_{e\in I}\sigma(f_e)e\) gives the Schwartz quotient topology \cite[\S6.1]{bernstein-2014}. The quotient seminorm from \(\sum_e\int|f_e|h^N\) is bounded by \(M\|Az\|\), \(A=(1+\Omega_K)^m\), by the Casimir Sobolev estimate \cite[Proposition~3.9]{bernstein-2014}. Consequently if \(|\langle w,\sigma(g)e\rangle|\le B_0 h(g)^N\) for every \(g,e\), then \(|\langle w,z\rangle|\le M B_0\|Az\|\). By unitarity and \(h(g^{-1})=h(g)\), apply this with \(w=Ay,\ z=A^{-1}y\), \(y\in V\), and \(B_0=Q(TAy)\sum_e c_e\). Division by \(\|y\|\) gives the bound for the identity word since \(TAy=A_E Ty\) for the derived polynomial \(A_E\).
\leanref{PositivityTheory.\allowbreak{}CWB.\allowbreak{}maxPr}{Development.\allowbreak{}WeilBound.\allowbreak{}CoreStatements.\allowbreak{}PublishedTheorems.\allowbreak{}CasselmanWallach.\allowbreak{}EtaleCohomologyNegativeVanishingDegreeZero}{91--92}
\leanref{PositivityTheory.\allowbreak{}CW.\allowbreak{}e3c}{Development.\allowbreak{}WeilBound.\allowbreak{}CoreStatements.\allowbreak{}PublishedTheorems.\allowbreak{}CasselmanWallach.\allowbreak{}EigenspacesOrthogonal}{23--115}

Apply this bound to \(d\sigma(D)y\) and extend by density. If \(Tw_i\) converges, then \(w_i\) is Cauchy in the complete smooth topology. Its limit maps to the given limit. The zero map has closed image as well.
\end{proof}

\subsection*{Proof of \ref{B8.40b}}\label{prf:B8.40b}\phantomsection
\begin{proof}
Put \(w=\iota^{-1}v\). The two coordinates give the same prime ideal, hence the same logarithmic weight and ramification index \(e_w\). In the tied quotients the normalized local length of Definition~\ref{B8.40} is the ordinary length divided by \(e_w\), the length of the residue characteristic quotient of the local DVR. Summing over the local indices gives \(a_{jv}=e_w\eta_{a(j)}(w)\). The selector has nonreference type at each twisted place by Lemma~\ref{B4.36a}. Thus reindexing gives the subsum over \(\iota^{-1}(\mathrm{Tw})\) of the sum defining \(\eta_{a(j)}(\mathcal L)\). Every remaining term is nonnegative.
\end{proof}

\subsection*{Proof of \ref{B8.40y}}\label{prf:B8.40y}\phantomsection
\begin{proof}
Apply \ref{B6.01} in the member's local data and take the supremum of the nonnegative finite sums. Eventually \((\log x)^2\le L\) and \(\sqrt b\le L(B_0+N_c)\), so \eqref{eq:B8-linkN} gives \eqref{eq:B8-divinN}. Apply Lemma~\ref{B8.40b} with the reference-row type of \ref{B4.36a}. The division tie identifies its local quotients, and the twisted-support condition puts their arithmetic places in \(S_{\mathcal L}\). Thus the twisted sum is bounded by \(\eta_{a(j)}(\mathcal L)\) in \eqref{eq:B8-divinN}. For \ref{B7.54}, use \ref{B8.43} and the \(\sqrt b\) budget of \ref{B2.58} supplied by \ref{B2.59}. The primitive, lattice and division ties identify the image of the layer-\(j-1\) projection at \(\mu_0\) with \(Z_j(\mu_0)\), of dimension \(d_{j,0}\). By the surjectivity assertion of Definition~\ref{B8.02}, this image is the full reference target, so \(d_{j,0}=D_0\). The bound enters row \(3\) of \ref{B8.38} and the cut term in \(E_{2,j}\) of \ref{B5.08x}.
\leanref{GlobalTheory.\allowbreak{}Top.\allowbreak{}smallLevelRankEq\_\allowbreak{}closed\_\allowbreak{}of\_\allowbreak{}levelProjection\_\allowbreak{}surjective}{Development.\allowbreak{}AbelianVarietyTheory.\allowbreak{}Proofs.\allowbreak{}GlobalInterpolation.\allowbreak{}SmallLevelRankEqClosedOfLevelProjectionSurjective}{13--48}
\leanref{UpperTheory.\allowbreak{}Top.\allowbreak{}hypV\_\allowbreak{}of\_\allowbreak{}cls}{Development.\allowbreak{}AutomorphicArithmetic.\allowbreak{}CharacterBoxes.\allowbreak{}ZeroFreeAllOfTopG}{16--17}
\end{proof}

\subsection*{Proof of \ref{B8.41}}\label{prf:B8.41}\phantomsection
\begin{proof}
Put \(s_\ell=\max(1,v_\ell(N))\) and \(\lambda_\ell=s_\ell+\log_\ell(2H_B)+1\). At \(2\), \(4\mid Q_S\) and \(s_0=s_2\ge2\). At \(3\) outside the small set the local fields are unramified, \(s_0=s_3=1\), and \(({-d}/3)=1\) by the split bound.
\leanref{UpperTheory.\allowbreak{}N.\allowbreak{}onAt\_\allowbreak{}eq\_\allowbreak{}two\_\allowbreak{}bounds\_\allowbreak{}of\_\allowbreak{}setup}{Development.\allowbreak{}AutomorphicArithmetic.\allowbreak{}ThetaCorrespondence.\allowbreak{}ThetaPairLifts.\allowbreak{}ThirdPlaceSeries.\allowbreak{}NRowTwo}{11--43}
\leanref{UpperTheory.\allowbreak{}N.\allowbreak{}rowThreeAt\_\allowbreak{}of\_\allowbreak{}setup\_\allowbreak{}alt}{Development.\allowbreak{}AutomorphicArithmetic.\allowbreak{}ThetaCorrespondence.\allowbreak{}ThetaPairLifts.\allowbreak{}ThirdPlaceSeries.\allowbreak{}NRowThree}{12--50}
\leanref{GlobalTheory.\allowbreak{}HubLoc.\allowbreak{}ramIdx\_\allowbreak{}eq\_\allowbreak{}one}{Development.\allowbreak{}AbelianVarietyTheory.\allowbreak{}Definitions.\allowbreak{}NonGaloisDescent.\allowbreak{}FinrankMaxUnramifiedSubfieldPadicFinExtPos}{94--112}

Since \(N^{36}t^{150}\le\sqrt b\), \(N\le b\). The exponents of \ref{B8.41a} used in \ref{B8.41b} have \(D_h\le3\lambda_\ell\) and \(\lambda_\ell\le6H_Bs_\ell\). With \(c=1000\) the bounds are
\begin{equation}\label{eq:B8-small-exponents}
 e'\le cH_Bs_\ell,\quad c_0\le cfn_V\lambda_\ell,\quad h_W(u)\le rc_0,\quad
 c_B\le ctf n_V\lambda_\ell,\quad h_0\le c\lambda_\ell,\quad\delta\le t\lambda_\ell .
\end{equation}
For \(c_B\), substitute \(a_0=2fn_V\) to bound \(40(t+1)(a_0+1)D_h\) by \(720tf n_V\lambda_\ell\). On \((-e)\) the inertia difference has valuation \(s_0+v_\ell(e)\). This gives \(h_0\le s_0+\lfloor\log_\ell t\rfloor\) and \(\delta\le t(s_0+1)\).
\leanref{UpperTheory.\allowbreak{}N.\allowbreak{}h0\_\allowbreak{}le}{Development.\allowbreak{}AutomorphicArithmetic.\allowbreak{}ResidueCharacters.\allowbreak{}LenUnitInvariantsLe}{55--71}

Each degree sum and cardinality of places is at most \(2f\). In \eqref{eq:B8-split-quotient} the \(e'\)-terms together cost at most \(6cfmH_Bs_\ell\). The terms \(d_um\delta,\ 2\delta,\ h_W+rc_0\) and the three terms in the final line cost at most \(2,4,4,6,12,2\), respectively, in units of \(cf^2t^2n_V\lambda_\ell\) by \(m\le f,\ r\le t\). The tame term is at most \(\ell^{-2}\) by \eqref{eq:B8-killing}. Thus the allowance in \eqref{eq:B8-Bd} is at most \(AfmH_Bs_\ell+Bf^2t^2n_V\lambda_\ell\) for \(A=6000,\ B=30000\). The same calculation applies to the odd small split primes with \(s_0=s_\ell\).
\leanref{UpperTheory.\allowbreak{}N.\allowbreak{}Num.\allowbreak{}boundSplit\_\allowbreak{}le}{Development.\allowbreak{}AutomorphicArithmetic.\allowbreak{}ResidueCharacters.\allowbreak{}NumSplit}{170--209}

For a family exponent \(C_{\rm par}\ge1\) bound \(f,n,t\) by \(x^{C_{\rm par}}\) and \(n_V\) by \(2x^{C_{\rm par}}\). The bounds \(\ell^{v_\ell(N)}\le Q_S\le n^2\) in this range give \(s_\ell\le3C_{\rm par}\log x\) eventually. Since \(H_B\le2b\), the two allowances are bounded by
\begin{equation}\label{eq:B8-cost-growth}
 18AC_{\rm par}fmb\log x,\qquad
 2B(3C_{\rm par}+9)x^{5C_{\rm par}+1}L .
\end{equation}
Enlarge the constant and take the common family cutoff. The closed row uses cutoff \(1\) for \(\varepsilon\) and zero thresholds.
\leanref{UpperTheory.\allowbreak{}N.\allowbreak{}Num.\allowbreak{}cost\_\allowbreak{}le'}{Development.\allowbreak{}AutomorphicArithmetic.\allowbreak{}ResidueCharacters.\allowbreak{}NumCost}{55--129}
\leanref{UpperTheory.\allowbreak{}N.\allowbreak{}boundAtTwoThree\_\allowbreak{}of\_\allowbreak{}local\_\allowbreak{}bounds\_\allowbreak{}at\_\allowbreak{}two\_\allowbreak{}three}{Development.\allowbreak{}AutomorphicArithmetic.\allowbreak{}ThetaCorrespondence.\allowbreak{}ThetaPairLifts.\allowbreak{}ThirdPlaceSeries.\allowbreak{}FamilyBoundAtTwoThreeOfLocalBounds}{37--54}
\end{proof}

\subsection*{Proof of \ref{B8.41ca}}\label{prf:B8.41ca}\phantomsection
\begin{proof}
Write \(S=W(k')[[Z]]\). Give \(S\) Witt Frobenius on the coefficients and put \(\varphi(Z)=(1+Z)^\ell-1\). The map \(1+Z\mapsto\zeta_{\ell^s}\), for a primitive root \(\zeta_{\ell^s}\), identifies \(S/(Q)\) with \(\mathcal O_{K'}\). Choose compatible primitive roots \((\zeta_{\ell^n})_{n\ge1}\), and let \(\epsilon=(1,\zeta_\ell,\zeta_{\ell^2},\ldots)\) in the tilt of the completed algebraic closure. The map from the cyclotomic prism \((S,(Q))\) to the Fontaine prism sends \(1+Z\) to the Teichmüller lift \([\epsilon^{1/\ell^s}]\). It sends \(Q\) to \(([\epsilon]-1)/([\epsilon^{1/\ell}]-1)\), a generator of the kernel of Fontaine's map \(\theta\).

Apply \ref{B8.41h} to the underlying \(\Z_\ell\)-lattice. Let \(\mathcal E\) be its crystal and \(M=\mathcal E(S)\). Full faithfulness lifts the \(O\)-action and every equivariant map of lattices. Since \(S\) is local, its finite projective module \(M\) is free. The action of
\[
 O\otimes_{\Z_\ell}W(k')=\prod_{i:W(k')\hookrightarrow O}O
\]
splits \(M\) into \(A\)-modules \(N_i(T)\). Multiplication by \(Z\) is injective on each summand. Its quotient by \(Z\) is a direct summand of \(M/ZM\), hence is finite and \(\ell\)-torsion free. It is consequently free over the discrete valuation ring \(O\). Lift an \(O\)-basis of this quotient. The lifts generate over \(A\) by Nakayama's lemma. In a relation among the lifts all coefficients are divisible by \(Z\). Cancelling \(Z\) and repeating puts every coefficient in \(\bigcap_n Z^nA=0\). Thus the lifts form an \(A\)-basis.

The linearized Frobenius of \(M\) is invertible after inverting \(Q\), by the definition of an \(F\)-crystal. Modulo \(\ell\), \(Q=Z^{\ell^{s-1}(\ell-1)}\), so inverting \(Q\) and completing \(\ell\)-adically gives the same coefficient ring as inverting \(Z\) and completing. The étale evaluation on this cyclotomic prism is the cyclotomic \((\varphi,\Gamma)\)-module, by \cite[Theorem~5.2 and Corollary~5.8]{wu-2021}. Here \(k'\) is perfect, \(K'/W(k')[1/\ell]\) is totally ramified, and the tower is \(K'(\mu_{\ell^\infty})\). Over the Witt vectors of the tilted algebraic closure the crystal comparison identifies the resulting module with the scalar extension of \(T\). Taking the \(i\)-component gives the asserted rank. All these identifications commute with lattice maps.

First specialize at \(Z=0\). Constant term gives a morphism of prisms
\[
 (S,(Q))\longrightarrow (W(k'),(\ell))
\]
over \(\mathcal O_{K'}\), since \(Q(0)=\ell\) and the induced map on the quotients is residue reduction. The action of \(\gamma\) on \(S\) fixes \(W(k')\) and sends \(1+Z\) to \((1+Z)^a\). Composing this action with constant term leaves the morphism unchanged. Crystal base change therefore identifies \(M/ZM\) with \(\mathcal E(W(k'))\) and gives the identity action there.

At \(\zeta=1\), write the action in a basis over \(L[[t_1]]\) as \(v(t_1)\mapsto G(t_1)v(a t_1)\). We have \(G(0)=I\). There is a unique matrix \(H(t_1)\) with \(H(0)=I\) and \(G(t_1)H(a t_1)=H(t_1)\). In degree \(j>0\) the equation determines \((1-a^j)H_j\) from coefficients of smaller degree. Our depth hypothesis gives \(a^j\ne1\) for every \(j>0\). The columns of \(H\) give the required fixed basis.

Now let \(\zeta\) have order \(\ell^b\), with \(1\le b<s\). It is fixed by \(\gamma\). Frobenius transports the expansion at \(\zeta^\ell\) in the predecessor component to the expansion at \(\zeta\), by \(t_{\zeta^\ell}\mapsto\ell t_\zeta\). This is a change of parameter over \(L\). The constant term of \(Q\) in the latter expansion is \(Q(\zeta-1)=\ell\), so the linearized Frobenius is an isomorphism there. It commutes with \(\gamma\). Induction on \(b\) transports the fixed basis to every lower root. In this basis \(\gamma\) multiplies the coefficient of degree \(j\) by \(a^j\).

For \(y\in N_i(T)\), the constant term of \((\gamma-1)y\) vanishes at each lower root. The change between the fixed basis and any \(A\)-basis is invertible over \(L[[t_\zeta]]\). Thus each coordinate in the \(A\)-basis vanishes at \(Z=\zeta-1\). Weierstrass division by the monic distinguished polynomial \(X\) gives a remainder of degree \(<\ell^{s-1}\). The remainder vanishes at all \(\ell^{s-1}\) distinct roots of \(X\), so it is zero. This proves the first inclusion over \(O\). Finally \(\gamma(X)=(1+X)^a-1\equiv aX\bmod X^2\). Applying this identity to \(X^jy\) gives the second inclusion. Applying an equivariant map to both inclusions gives them on its image.
\end{proof}

\subsection*{Proof of \ref{B8.41cb}}\label{prf:B8.41cb}\phantomsection
\begin{proof}
Let \(C^\flat\) be the tilt of the completed algebraic closure of \(\Q_\ell\), with valuation ring \(\mathcal O^\flat\). Use the rings
\[
 R_0=O\otimes_{\Z_\ell}W(\mathcal O^\flat),\qquad
 R_1=O\otimes_{\Z_\ell}W(C^\flat).
\]
Frobenius acts on the Witt factor. The étale comparison of Lemma~\ref{B8.41ca} identifies \(R_1\otimes_B\mathcal D\) with \(R_1\otimes_O T\), with Frobenius fixing the \(T\)-factor. The scalar maps \(A_*\to R_0\), \(B\to R_1\) are compatible with Frobenius.

We first prove an integrality property. Suppose a finite matrix \(V\) over \(R_1\) satisfies \(\varphi(V)=H V H'+Y\), where \(H,H',Y\) have entries in \(R_0\). Then \(V\) has entries in \(R_0\). First take \(O=\Z_\ell\) and reduce the equation modulo \(\ell\). If the largest absolute value \(b\) of an entry of \(\bar V\) exceeds \(1\), a largest entry gives \(b^\ell\le b\), a contradiction. Lift the entries of \(\bar V\) by their integral Teichmüller representatives \(V_0\), and write \(V=V_0+\ell V_1\). The equation for \(V_1\) has the same coefficients \(H,H'\) and integral inhomogeneous term
\[
 (HV_0H'+Y-\varphi(V_0))/\ell .
\]
Indeed the numerator reduces to zero already over \(\mathcal O^\flat\), and reduction of its Witt ring has kernel \(\ell W(\mathcal O^\flat)\), since \(\mathcal O^\flat\) is perfect. Iteration and \(\ell\)-adic completeness prove integrality. For general \(O\), choose a \(\Z_\ell\)-basis and replace each entry by the matrix of multiplication on the \(O\)-factor. This commutes with Frobenius and multiplication, so the preceding proof applies to the enlarged matrices. Evaluating a multiplication matrix on \(1\in O\) recovers the original entry over \(R_0\).
\leanref{GlobalTheory.\allowbreak{}FrobInt.\allowbreak{}integral\_\allowbreak{}of\_\allowbreak{}frobenius}{Development.\allowbreak{}AbelianVarietyTheory.\allowbreak{}Proofs.\allowbreak{}InertiaEigenvalues.\allowbreak{}FrobeniusIntegralityWittTiltAinf}{239--260}
\leanref{GlobalTheory.\allowbreak{}FrobIntegralCoeff.\allowbreak{}frobenius\_\allowbreak{}integrality\_\allowbreak{}coeff}{Development.\allowbreak{}AbelianVarietyTheory.\allowbreak{}Proofs.\allowbreak{}InertiaEigenvalues.\allowbreak{}FrobeniusIntegralityTensorAinfCoefficients}{126--132}

We also need the scalar intersection \(B\cap R_0=A_*\) inside \(R_1\). Let \(\varpi\) be a uniformizer of \(O\), with residue field \(k\). Modulo \(\varpi\) the source inclusion is \(\prod_i k[[Z]]\subset\prod_i k((Z))\). On the period-ring side it is \(k\otimes_{\F_\ell}\mathcal O^\flat\subset k\otimes_{\F_\ell} C^\flat\), a product of valuation-ring inclusions indexed by embeddings of \(k\). Each component \(i\) is seen by a residue embedding extending the corresponding unramified embedding. In this factor \(Z\) maps to \(z=\epsilon^{1/\ell^s}-1\), of finite strictly positive valuation, with \(\epsilon\) as in Lemma~\ref{B8.41ca}. For a nonzero Laurent series with first exponent \(j\), the valuation of its image is \(jv(z)\): its first coefficient is a unit and subsequent exponents have strictly larger valuations. Thus the residue intersection consists exactly of the power series, and the residue map on the Laurent series is injective. Given an element of \(B\) with image in \(R_0\), lift its residue from \(A_*\), subtract and divide by \(\varpi\). The new element again has image in \(R_0\), by injectivity of \(R_0/\varpi\to R_1/\varpi\). Repetition gives a \(\varpi\)-adically convergent expansion in \(A_*\). The rings are separated, and the same residue argument gives injectivity of \(B\to R_1\).
\leanref{GlobalTheory.\allowbreak{}NL.\allowbreak{}IoXQ.\allowbreak{}ioE\_\allowbreak{}inj\_\allowbreak{}inter}{Development.\allowbreak{}AbelianVarietyTheory.\allowbreak{}Proofs.\allowbreak{}InertiaEigenvalues.\allowbreak{}WachInclusionXQFrobenius}{160--229}

Choose an \(O\)-basis of \(T\). For \(M=L,L'\) choose a basis of \(M\) over \(A_*\), assembling bases of the components, which all have rank \(\operatorname{rank}_O T\). Let \(U_M\) be the matrix whose rows are its coordinates in the chosen basis of \(T\) after extension to \(R_1\). It is invertible over \(R_1\). The Frobenius hypothesis, in both directions on the finitely many basis vectors, supplies \(h_M\ge0\) and matrices \(P_M,H_M\) over \(A_*\) such that
\[
 Q^{h_M}\varphi(U_M)=P_M U_M,\qquad
 Q^{h_M}U_M=H_M\varphi(U_M).
\]
Writing \(x\) for the image of \(X\), we have \(\varphi(x)=xQ\). Hence
\[
 \varphi(x^{h_M}U_M)=P_M(x^{h_M}U_M),\qquad
 \varphi(x^{h_M}U_M^{-1})=(x^{h_M}U_M^{-1})H_M .
\]
The integrality property applies to both matrices. Put \(a=h_L+h_{L'}\). The products \(x^a U_L U_{L'}^{-1}\) and \(x^a U_{L'}U_L^{-1}\) have entries in \(R_0\). These are the images of the two scaled transition matrices over \(B\), because both bases are bases of \(\mathcal D\). The scalar intersection puts their entries in \(A_*\). This proves both inclusions.
\leanref{GlobalTheory.\allowbreak{}Sandwich.\allowbreak{}upper}{Development.\allowbreak{}AbelianVarietyTheory.\allowbreak{}Proofs.\allowbreak{}InertiaEigenvalues.\allowbreak{}WachLatticeSandwichAinf}{39--69}
\leanref{GlobalTheory.\allowbreak{}Sandwich.\allowbreak{}lower}{Development.\allowbreak{}AbelianVarietyTheory.\allowbreak{}Proofs.\allowbreak{}InertiaEigenvalues.\allowbreak{}WachLatticeSandwichAinf}{71--127}

For the last assertion, write either transition matrix on component \(i\) as \(C_i=X^{-a}B_i\), with \(B_i\) integral. At \(Z=0\) use \(t=\log(1+Z)\) over \(\operatorname{Frac}(O)\). If \(C_i\) has a pole there with first term \(t^{-k}C_0\), \(k>0\), the equivariance equation for \(\gamma\) gives \(u^{-k}C_0=C_0\). Indeed \(\gamma(t)=ut\), and both lattice matrices of \(\gamma\) have constant term the identity. The depth assumption gives \(u^k\ne1\), since \(1+\ell^s\Z_\ell\) has injective logarithm for \(s\ge1\) at odd \(\ell\) and for \(s\ge2\) at \(2\). Thus all transitions are regular at \(Z=0\).

At \(Z=\zeta-1\) with \(\zeta^{\ell^{s-1}}=1\) one has \(Q(\zeta-1)=\ell\). Both Frobenius matrices and their inverses are regular in the characteristic-zero expansion there. Frobenius carries the expansion at \(\zeta^\ell\) in the predecessor component to the expansion at \(\zeta\), by \(t_{\zeta^\ell}\mapsto\ell t_\zeta\). Its transition equation therefore propagates regularity from \(\zeta^\ell\) to \(\zeta\). Induction gives regularity at every root of \(X\). Each entry of \(B_i\) vanishes to order at least \(a\) at all these simple roots. The remainder on Weierstrass division by \(X^a\) over \(O\) has the same vanishings and degree less than \(a\deg X\), hence is zero. For \(a=0\) the transition was already integral. This applies in both directions and proves equality.
\leanref{GlobalTheory.\allowbreak{}NL.\allowbreak{}WachMat.\allowbreak{}X\_\allowbreak{}pow\_\allowbreak{}dvd\_\allowbreak{}entry}{Development.\allowbreak{}AbelianVarietyTheory.\allowbreak{}Proofs.\allowbreak{}InertiaEigenvalues.\allowbreak{}MapConstantCoeffEqOne}{103--182}
\end{proof}

\subsection*{Proof of \ref{B8.41c}}\label{prf:B8.41c}\phantomsection
\begin{proof}
Apply Lemma~\ref{B8.41ca} with \(a=u\) to the source and target. For any of the Wach systems in \ref{B8.41}, the same lower-root calculation applies: its action is trivial modulo \(Z\), its linearized Frobenius is invertible off \(Q\), and Frobenius commutes with \(\gamma\). The map \(f\) is equivariant, so on either lattice and on its image we have
\begin{equation}\label{eq:B8-wach-graded}
 (\gamma-u^j)X^jN\subset X^{j+1}N\qquad(j\ge0).
\end{equation}
\leanref{GlobalTheory.\allowbreak{}PrismS.\allowbreak{}exists\_\allowbreak{}prism}{Development.\allowbreak{}AbelianVarietyTheory.\allowbreak{}Proofs.\allowbreak{}InertiaEigenvalues.\allowbreak{}ExistsWachPrismOverIntegers}{331--370}
\leanref{GlobalTheory.\allowbreak{}WachLattice.\allowbreak{}build\_\allowbreak{}prism}{Development.\allowbreak{}AbelianVarietyTheory.\allowbreak{}Proofs.\allowbreak{}InertiaEigenvalues.\allowbreak{}WachLatticesExistAllLayersOfBhattScholze}{26--169}
\leanref{GlobalTheory.\allowbreak{}NL.\allowbreak{}WachLatticeTrivialModUniformizer.\allowbreak{}wachLattice\_\allowbreak{}galois\_\allowbreak{}trivial\_\allowbreak{}mod\_\allowbreak{}uniformizer}{Development.\allowbreak{}AbelianVarietyTheory.\allowbreak{}Proofs.\allowbreak{}InertiaEigenvalues.\allowbreak{}WachLatticeGaloisTrivialModUniformizer}{12--18}
\leanref{GlobalTheory.\allowbreak{}NL.\allowbreak{}GamModX.\allowbreak{}sub\_\allowbreak{}mem\_\allowbreak{}X\_\allowbreak{}of\_\allowbreak{}val}{Development.\allowbreak{}AbelianVarietyTheory.\allowbreak{}Proofs.\allowbreak{}InertiaEigenvalues.\allowbreak{}ExistsSmulOfDvdRepr}{70--96}

Put \(P=fN_1\) and \(S=\{y\in N_2:\ell^cy,Z^cy\in P\text{ for some }c\ge0\}\). Work componentwise and let \(\varpi\) be a uniformizer of \(O\). The module \(S\) is finite since \(O[[Z]]\) is noetherian. If \(Zx=\varpi y\) in \(S\), freeness of \(N_2\) gives \(x=\varpi z,\ y=Zz\) in \(N_2\). A power of \(\ell\) sends \(z\) into \(P\) by the condition on \(x\) and \(\varpi\mid\ell\). A power of \(Z\) does so by the condition on \(y\). Thus \(z\in S\), and \(S/\varpi S\) is torsion free over the discrete valuation ring \((O/\varpi)[[Z]]\). Lift a basis of this free module. The lifts generate by Nakayama's lemma, and reducing a relation repeatedly modulo \(\varpi\) proves independence. Étale base change is flat, being localization followed by adic completion of a noetherian ring. By exactness of the étale module functor and surjectivity of \(f\), the base change of \(P\) fills that of \(N_2\). The same holds for the intermediate module \(S\).

The finite module \(S/P\) is supported at the closed point: powers of \(\ell\) and \(Z\) kill each of a finite set of generators. Its Frobenius pullback is also supported there, since \(\varphi(Z)\equiv Z^\ell\pmod\varpi\). Off \(Q\) both modules vanish, since \(Q\in(\varpi,Z)\). Frobenius pullback is flat: on each component the scalar ring is free over its Frobenius image on \((1+Z)^j\), \(0\le j<\ell\), by monic division in \(1+Z\) with \((1+Z)^\ell=1+\varphi(Z)\). The linearized Frobenius on \(P[1/Q]\) is consequently an isomorphism, by taking images in the equivariant map between the two Wach lattices, whose Frobenius maps are isomorphisms there. The equalities with \(S\) and its pullback off \(Q\) give the same property for \(S\). For \(y\in S\) choose \(c\) with \(\ell^cy,Z^cy\in P\). Since \(\gamma(Z)/Z\) is a unit and \(P\) is stable, \(Z^c\gamma y\in P\). Write \(z=(\gamma-1)y/Z\in N_2\), using \(Z\mid X\) and \eqref{eq:B8-wach-graded}. The same inclusion on \(P\) gives \(\ell^cz\in P\), and \(Z^{c+1}z=Z^c(\gamma y-y)\in P\). Thus \(z\in S\). Both free lattices now satisfy Lemma~\ref{B8.41cb}, which gives \(S=N_2\). By the definition of \(S\), a common power of \(Z\) sends a finite basis of every component of \(N_2\) into \(P\).
\leanref{GlobalTheory.\allowbreak{}NL.\allowbreak{}SatFree.\allowbreak{}satM\_\allowbreak{}free\_\allowbreak{}loc}{Development.\allowbreak{}AbelianVarietyTheory.\allowbreak{}Proofs.\allowbreak{}InertiaEigenvalues.\allowbreak{}WachSaturationImageFreeAndSpans}{86--166}
\leanref{GlobalTheory.\allowbreak{}T.\allowbreak{}wachModule\_\allowbreak{}uniqueness\_\allowbreak{}characterisation}{Development.\allowbreak{}AbelianVarietyTheory.\allowbreak{}Proofs.\allowbreak{}InertiaEigenvalues.\allowbreak{}WachModuleUniquenessCharacterisationStatement}{7--8}
\leanref{GlobalTheory.\allowbreak{}NL.\allowbreak{}FrobeniusMatrixUniqueness.\allowbreak{}main}{Development.\allowbreak{}AbelianVarietyTheory.\allowbreak{}Proofs.\allowbreak{}InertiaEigenvalues.\allowbreak{}WachModuleUniquenessCharacterisation}{247--269}
\leanref{GlobalTheory.\allowbreak{}NL.\allowbreak{}CokerFin.\allowbreak{}coker\_\allowbreak{}killed}{Development.\allowbreak{}AbelianVarietyTheory.\allowbreak{}Proofs.\allowbreak{}InertiaEigenvalues.\allowbreak{}WachNCokernelKilledByPower}{19--38}

At a primitive \(\ell^s\)-th root the completed Wach lattice has degree-\(k\) coefficients in \(\mathrm{Fil}^{-k}D_{\rm dR}\). Frobenius transports the predecessor's ordinary power series to this root and sends \(t_{\zeta^\ell}\) to \(\ell t_\zeta\). The source bounds and strictness for the quotient give
\begin{equation}\label{eq:B8-wach-heights}
 Q^{M_i}\varphi N_{j,i-1}\subset N_{ji},\qquad Q^{h_i}N_{ji}\subset A\varphi N_{j,i-1}
 \quad(j=1,2).
\end{equation}
The required orders follow from the filtration, and divisibility in a destination basis follows by monic division at the simple roots of \(Q\).
\leanref{GlobalTheory.\allowbreak{}NL.\allowbreak{}Heights.\allowbreak{}heights\_\allowbreak{}from\_\allowbreak{}jumps}{Development.\allowbreak{}AbelianVarietyTheory.\allowbreak{}Proofs.\allowbreak{}InertiaEigenvalues.\allowbreak{}WachLatticeHeightsFromHodgeJumps}{382--474}

Let \(\mathfrak a_i=\{g:gN_{2i}\subset P_i\}\). The scalar ring is free over \(\varphi(A)\) on \(t_a=(1+Z)^a,\ 0\le a<\ell\). For \(g\in\mathfrak a_i\) write \(Q^Hg=\sum_a t_a\varphi(g_a)\). For \(y\in N_{2,i-1}\) lift \(gQ^{M_i}\varphi y\) to the source and apply the second inclusion of \eqref{eq:B8-wach-heights}. Multiplication by \(Q^{H-M_i-h_i}\) puts \(\sum_a t_a\varphi(g_a y)\) in \(A\varphi(P_{i-1})\). Uniqueness of the scalar expansion in a Frobenius basis gives \(g_a y\in P_{i-1}\). Thus \(Q^H\mathfrak a_i\subset A\varphi(\mathfrak a_{i-1})\).
\leanref{GlobalTheory.\allowbreak{}NL.\allowbreak{}CokerKill.\allowbreak{}key\_\allowbreak{}inclusion}{Development.\allowbreak{}AbelianVarietyTheory.\allowbreak{}Proofs.\allowbreak{}InertiaEigenvalues.\allowbreak{}CokernelCrystallineSurjectionKilledByXPow}{233--319}
If \(X^H\notin\mathfrak a_i\), let \(J_k=(\mathfrak a_{i+k}:X^H)\). Then \(J_0\subset(\varpi,Z)\) and every Frobenius coordinate of an element of \(J_{k+1}\) belongs to \(J_k\), since \(Q^HX^H=\varphi(X^H)\). Reduction modulo \(\varpi\) gives \(J_k\subset(\varpi,Z^{k+1})\), contrary to \(Z^c\in J_c\).
\leanref{GlobalTheory.\allowbreak{}Cyc.\allowbreak{}XX\_\allowbreak{}pow\_\allowbreak{}mem\_\allowbreak{}of\_\allowbreak{}QQ\_\allowbreak{}pow\_\allowbreak{}mul\_\allowbreak{}mem\_\allowbreak{}map\_\allowbreak{}phi}{Development.\allowbreak{}AbelianVarietyTheory.\allowbreak{}Proofs.\allowbreak{}InertiaEigenvalues.\allowbreak{}XXPowMemOfQQPowMulMemMapPhi}{80--117}

For \(H>0\) form \(P_0(U)=\prod_{j=0}^{H-1}(U-u^j)\), \(P_1(U)=\prod_{j=H}^{2H-1}(U-u^j)\). Both \(P_0(\gamma)X^Hv\) and \(P_1(\gamma)X^Hv\) belong to \(X^HP_i\) for \(v\in N_{2i}\), by \eqref{eq:B8-wach-graded} on \(P_i\) and \(N_{2i}\), respectively. Binomial expansion, using \(s\ge2\) at \(2\), gives
\begin{equation}\label{eq:B8-cyc-val}
 v_\ell(u^k-1)=s+v_\ell(k)\quad(k\ne0).
\end{equation}
The \(i\)-th denominator for interpolating \(1/P_1\) at \(u^i\), \(0\le i<H\), has valuation \((2H-1)s+v_\ell(i!(2H-1-i)!)\le w\). Thus an integral polynomial \(B\) satisfies \(P_0\mid \ell^w-BP_1\). Evaluate at \(\gamma\) on \(X^Hv\), then cancel \(X^H\). At \(H=0\) the first inclusion gives equality.
\leanref{Formalism.\allowbreak{}LIFT.\allowbreak{}gammaGradedStep\_\allowbreak{}holds}{Formalism.\allowbreak{}LagrangeIdealLifts.\allowbreak{}CocyclePolynomialDivision.\allowbreak{}Graded}{5--26}
\leanref{Formalism.\allowbreak{}LIFT.\allowbreak{}wBTail\_\allowbreak{}of}{Formalism.\allowbreak{}LagrangeIdealLifts.\allowbreak{}CocyclePolynomialDivision.\allowbreak{}Graded}{115--158}
\end{proof}

\subsection*{Proof of \ref{B8.41aa}}\label{prf:B8.41aa}\phantomsection
\begin{proof}
Write \(\chi_\ell\) for the cyclotomic character. The field \(K_0\) is unramified and the extensions \(\Q_\ell(\mu_{\ell^j})/\Q_\ell\) are totally ramified. Hence the image of \(G_{K'}\) under \(\chi_\ell\) is \(1+\ell^s\Z_\ell\). For an element \(g\) with \(v_\ell(\chi_\ell(g)-1)=s\), equation~\eqref{eq:B8-cyc-val} gives
\[
 v_\ell(\chi_\ell(g)^{e+1}-1)=b.
\]
It follows that \(\ell^b\) annihilates the invariants of \((\Q_\ell/\Z_\ell)(e+1)\).

Apply Theorem~\ref{B8.41p} over \(K'\) to \((\Z/\ell^j)(-e)\) for each \(j\ge1\). Its dual under the local pairing is \((\Z/\ell^j)(e+1)\). Taking the inverse limit on \(H^2\) and the direct limit on \(H^0\) identifies \(H^2(K',\Z_\ell(-e))\) with the Pontryagin dual of the preceding invariant group. The finite \(H^1\)-groups in Theorem~\ref{B8.41p} satisfy the Mittag--Leffler condition, so the inverse limit computes integral \(H^2\). In particular \(\ell^b\) kills the obstruction to lifting \([\bar z]\) through \(\Q_\ell(-e)\). Choose a rational cocycle lifting \(\ell^b[\bar z]\). The difference between its reduction and \(\ell^b\bar z\) is the boundary of a scalar in \((\Q_\ell/\Z_\ell)(-e)\). Subtract the boundary of any rational lift of that scalar.
\end{proof}

\subsection*{Proof of \ref{B8.41ab}}\label{prf:B8.41ab}\phantomsection
\begin{proof}
Solve \(A_i\gamma(U_i)=U_i\) with \(U_i(0)=1\). In degree \(j>0\) the coefficient to invert is \(1-u^j\). All coefficients multiplying earlier entries are integral, since \(\binom{u}{h}\in\Z_\ell\). Writing \(j=\ell^a j_0\) with \(\ell\nmid j_0\), the binomial expansion gives \(v_\ell(u^{j_0}-1)=s\). Raising a unit of this depth to the \(\ell\)-th power increases the depth by one. For \(\ell=2\) this uses \(s\ge2\). Thus the coefficient of \(Z^j\) in \(U_i\) has norm at most \(\ell^{sj+v_\ell(j!)}\), by induction. Since \(v_\ell(j!)\le j/(\ell-1)\), all the matrices converge on a common disc of positive radius. Their bottom row is \((0,1)\) by uniqueness. The negative of the upper part of the last column is \(w_i\).

The commutation identity makes \(B_i\varphi(U_{i-1})\) another matrix of \(\gamma\)-invariant columns. A column is uniquely determined by its constant coefficient, by the same recursion. Therefore
\[
                   U_i=B_i\varphi(U_{i-1})B_i(0)^{-1}.
\]
Set \(M=m+|J|\), \(\tau=t_0/X\), and \(D=\operatorname{diag}(\tau^M 1,1)\). The series \(\tau\) is analytic on the open unit disc, since it equals \(\ell^{1-s}\log(1+X)/X\). For each \(j\ge0\), put \(X_j=\varphi^j(X)\). From \(X_{j+1}=X\prod_{h=0}^{j}\varphi^h(Q)\) we get
\[
 \frac{\tau}{\varphi^j(Q)}
   =\frac{\ell^{-(s+j)}\log(1+X_{j+1})}{X_{j+1}}
      \prod_{h<j}\varphi^h(Q).
\]
The quotient of the logarithm by its argument has a removable singularity at zero. It follows that \(\tau^M/\Lambda\) is analytic on the open unit disc, also when \(J\) contains repetitions. Since \(\varphi(\tau)=\ell\tau/Q\), the matrix \(DB_i\varphi(D)^{-1}\) has upper blocks \(Q^{|J|}P_i/\ell^M\) and \(\tau^M y_i/\Lambda\). Both are analytic there.

Now \(D U_i=(D B_i\varphi(D)^{-1})\varphi(D U_{i-1})B_i(0)^{-1}\). For any \(0<R<1\), substitution by \(\varphi\) maps the closed disc of radius \(R\) into the disc of radius \(R\max(\ell^{-1},R)\), by the binomial formula. Iteration reaches the common initial disc because \(\max(\ell^{-1},R)<1\). The displayed matrix identity then gives convergence of all \(D U_i\) on the disc of radius \(R\), working back through the finitely many iterations. In particular \(\tau^M w_i\) converges there. Multiplication by \(X^M\) gives the assertion for \(t_0^M w_i\).
\end{proof}

\subsection*{Proof of \ref{B8.41a}}\label{prf:B8.41a}\phantomsection
\begin{proof}
The prime-to-\(\ell\) roots generate an unramified field and \(Q=\Phi_{\ell^s}(1+Z)\) is Eisenstein of degree \((\ell-1)\ell^{s-1}\). A primitive root of order \(\operatorname{lcm}(\ell,N)\) generates \(K'\). Hence
\begin{equation}\label{eq:B8-cyc-degree}
 e(K'/\Q_\ell)=(\ell-1)\ell^{s-1},\qquad v_\ell([K':K])\le j(N)+1.
\end{equation}
Restriction and corestriction preserve period coboundaries, and their composite multiplies by \([K':K]\). Jumps pull back under scalar extension.
\leanref{UpperTheory.\allowbreak{}N.\allowbreak{}Cyc.\allowbreak{}val\_\allowbreak{}index\_\allowbreak{}le}{Development.\allowbreak{}AutomorphicArithmetic.\allowbreak{}ResidueCharacters.\allowbreak{}FiniteRoots}{117--125}

Work over \(K'\), with \(H=2h\) and \(w=2Hs+v_\ell((2H)!)\). For \(T_0=\ker(T_1\to T_2)\) write \(N_j=\mathcal N(T_j)\). On \(D(T)\) use
\[
 \partial b=((\varphi-1)b,(\gamma-1)b),\qquad d(y,z)=(\gamma-1)y-(\varphi-1)z.
\]
Herr's theorem \cite[Theorem~2.1]{herr-1998} computes cohomology modulo powers of \(\ell\); the finite degree-zero groups give the Mittag--Leffler condition for passage to the lattice. Write \(A_h(M)=Q^{-h}M\). The comparison on \(A_h(M)\times XM\) gives
\begin{equation}\label{eq:B8-finite-pairs}
 d(y,z)=0\ \Rightarrow\ [y,z]\in H^1_f(K',T),\qquad
 [y,z]=0\ \Rightarrow\ (y,z)=\partial b_0\ \text{with }b_0\in M
\end{equation}
for closed pairs in this domain. Solve \((\varphi-1)b=y\) in the Witt tensor. The Frobenius denominator bound puts a product of Frobenius translates of \(X\) times \(b\) in the period comparison domain. Those factors are units in \(B_{\rm cris}\), so division gives a period vector \(b'\) with the Galois differences corresponding to those of \(b\). In component bases of the Wach lattice, augment the matrices of \(\gamma\) and \(\varphi\) by their respective cocycle columns \(z\) and \(y\), and add the bottom row \((0,1)\). The equality \(d(y,z)=0\) is exactly the upper-right block of the commutation identity in Lemma~\ref{B8.41ab}. Here \(z\in XM\) has zero constant term. By Lemma~\ref{B8.41c} a power \(Q^m\) clears the Frobenius matrix, which is invertible after inverting \(Q\), and \(Q(0)=\ell\ne0\) in the coefficient field. For the column \(y\) take \(J\) to consist of \(h\) copies of zero. Lemma~\ref{B8.41ab} gives a zero-constant primitive \(w\) with \(\log(1+Z)^{m+h}w\) convergent on the open unit disc. To evaluate it, put \(z_s=[\underline\epsilon^{1/\ell^s}]-1\) in \(A_{\rm cris}\), where \(\underline\epsilon\) is the chosen compatible system of cyclotomic roots. Modulo \(\ell\), \(z_s^{\ell^s}\) equals \([\underline\epsilon]-1\). The latter belongs to the divided-power ideal, so its \(\ell\)-th power belongs to \(\ell A_{\rm cris}\). Consequently \(z_s^{\ell^{s+1}}\in\ell A_{\rm cris}\). For an open-disc series \(\sum a_nZ^n\), choose \(R\) with \(\ell^{-1/\ell^{s+1}}<R<1\). Then \(|a_n|R^n\to0\), and the evaluated terms tend to zero in the complete coefficient tensor of \(A_{\rm cris}[1/\ell]\), with a common denominator. Evaluation respects multiplication and the cyclotomic substitutions by convergence. The logarithm evaluates to \(\ell^{-s}t_{\rm cris}\), where \(t_{\rm cris}=\log[\underline\epsilon]\) is invertible in \(B_{\rm cris}\). Dividing by its indicated power gives the period coordinates of \(w'\in B_{\rm cris}\otimes T[1/\ell]\). For \(g\in G_{K'}\) the corresponding cyclotomic value solves \((\gamma-1)z_g=(g-1)z\) in the same zero-constant domain. Thus \((g-1)w'\) is its period image. The cocycle of the pair has period primitive \(w'-b'\). For a zero class take a Herr primitive \(b_0\in D(T)\). In the étale module of the crystalline lattice \(T\oplus O\), let \(e\) generate the trivial summand and compare the two lattices \(M\oplus A_*e\) and \(M\oplus A_*(e+b_0)\), with \(A_*\) as in Lemma~\ref{B8.41cb}. Both are componentwise free and span the same étale module. In the second basis the extra Frobenius column is \(y\in Q^{-h}M\) and the bottom diagonal entry is one, so linearized Frobenius is invertible after inverting \(Q\) on both lattices. The extra cyclotomic column is \(z\in XM\subset ZM\), so \(\gamma\) preserves both lattices and acts trivially modulo \(Z\) on each. Lemma~\ref{B8.41cb} identifies the two lattices. Subtracting \(e\) from \(e+b_0\) now gives \(b_0\in M\). The kernel comparison gives \(N_0=N_1\cap D(T_0)\).
\leanref{Formalism.\allowbreak{}MainNotions.\allowbreak{}LiftingTheory.\allowbreak{}coreWX}{Formalism.\allowbreak{}SharpRestriction.\allowbreak{}LevelLetterCocycles.\allowbreak{}LiftingTheory}{100--103}
\leanref{Formalism.\allowbreak{}MainNotions.\allowbreak{}LiftingTheory.\allowbreak{}coreWB\_\allowbreak{}of}{Formalism.\allowbreak{}SharpRestriction.\allowbreak{}LevelLetterCocycles.\allowbreak{}FrobeniusPowUniformizerFactorisation}{3449--3457}
\leanref{GlobalTheory.\allowbreak{}NL.\allowbreak{}IntComplex.\allowbreak{}main}{Development.\allowbreak{}AbelianVarietyTheory.\allowbreak{}Proofs.\allowbreak{}InertiaEigenvalues.\allowbreak{}WachIntegralComplexCocycleClassFinite}{40--58}
\leanref{GlobalTheory.\allowbreak{}NL.\allowbreak{}PoleCrys.\allowbreak{}main}{Development.\allowbreak{}AbelianVarietyTheory.\allowbreak{}Proofs.\allowbreak{}InertiaEigenvalues.\allowbreak{}HerrCocyclePolesFrobeniusQCrystalline}{539--562}
\leanref{GlobalTheory.\allowbreak{}NL.\allowbreak{}CocycleInj.\allowbreak{}main}{Development.\allowbreak{}AbelianVarietyTheory.\allowbreak{}Proofs.\allowbreak{}InertiaEigenvalues.\allowbreak{}MemWachOfFrobeniusGamma}{12--16}
\leanref{GlobalTheory.\allowbreak{}NL.\allowbreak{}Exact.\allowbreak{}exactness\_\allowbreak{}of\_\allowbreak{}N}{Development.\allowbreak{}AbelianVarietyTheory.\allowbreak{}Proofs.\allowbreak{}InertiaEigenvalues.\allowbreak{}WachModuleSaturatedExactnessCrystalline}{110--135}

Represent a finite target class \(c\) by its crystalline extension of the unit lattice. Lift \(\ell^w\) times the unit in the Wach lattice of the extension by \ref{B8.41c}. Its two differences give a closed pair \((y,z)\in A_h(N_2)\times XN_2\) representing \(\ell^wc\). Lift \(\ell^w Q^hy\) and \(\ell^w z/X\) to \(N_1\) by the same lemma, then divide the first by \(Q^h\) and multiply the second by \(X\) to get a pair \(v\). Since \(\gamma(Q)/Q\equiv1\bmod X\) and \(\varphi(X)=XQ\), its defect \(\omega=dv\) lies in \(XA_h(N_0)\). Some \(\ell^k\omega\) is \(d(Y,Xn)\) in the kernel. The Bloch--Kato presentation using \(D_{\rm cris}\oplus D_{\rm dR}/\mathrm{Fil}^0\) and strictness give right exactness of rational finite cohomology. Clear denominators of the lifted cocycle and kill the torsion difference downstairs. Represent a further multiple by the extension construction. Its difference from the corresponding multiple of the image of \(v\) is the boundary of an element of \(N_2\) by \eqref{eq:B8-finite-pairs}. Lift a further \(\ell^w\) times that element and subtract its boundary. The difference upstairs is now a kernel pair.
\leanref{GlobalTheory.\allowbreak{}NL.\allowbreak{}RightExactKs.\allowbreak{}h1fKs\_\allowbreak{}onto\_\allowbreak{}up\_\allowbreak{}to\_\allowbreak{}pow}{Development.\allowbreak{}AbelianVarietyTheory.\allowbreak{}Proofs.\allowbreak{}InertiaEigenvalues.\allowbreak{}H1fKsOntoUpToPow}{14--24}
\leanref{GlobalTheory.\allowbreak{}NL.\allowbreak{}DefectTorsion.\allowbreak{}defect\_\allowbreak{}of\_\allowbreak{}lift}{Development.\allowbreak{}AbelianVarietyTheory.\allowbreak{}Proofs.\allowbreak{}InertiaEigenvalues.\allowbreak{}DefectTorsionSecondCohomologyKernel}{286--334}

For \(M=N_0\) and the union \(J\) of source jump sets let
\[
 J^*=((J+[1,h])\setminus\{0\})\cup[1,h],\quad
 P(U)=\prod_{j\in J^*}(U-u^{-j}),\quad p_\gamma=P(\gamma),\quad
 c_\gamma=\prod_{j=1}^h(u^j-1).
\]
At a root of \(Q\), a degree-\(k\) coefficient of \(Q^hy\), for \(y\in A_h(M)\), belongs to \(\mathrm{Fil}^{-k}\), and integrality of \(y\) requires \(\mathrm{Fil}^{h-k}\). A failure uses a jump \(a\in J\) with \(h-k-a\in[1,h]\). On the first such degree of \(y\) the scalar is \(u^{k-h}\) by expansion of \(y\). The jump filtration therefore gives \((\gamma-1)p_\gamma A_h(M)\subset M\), by applying the factors in degree order and then using root divisibility.
\leanref{GlobalTheory.\allowbreak{}NL.\allowbreak{}KillLevels.\allowbreak{}annihilate}{Development.\allowbreak{}AbelianVarietyTheory.\allowbreak{}Proofs.\allowbreak{}InertiaEigenvalues.\allowbreak{}ProdGammaSubJumpUnitsAnnihilatesQuotient}{285--355}
\leanref{Formalism.\allowbreak{}MainNotions.\allowbreak{}LiftingTheory.\allowbreak{}liftb\_\allowbreak{}semi}{Formalism.\allowbreak{}SharpRestriction.\allowbreak{}LevelLetterCocycles.\allowbreak{}BaseChangePreservesUnitaryOf}{404--473}

Successive use of \eqref{eq:B8-wach-graded} in degrees \(1,\ldots,h\) gives
\begin{equation}\label{eq:B8-graded-reduction}
 c_\gamma Xn=(\gamma-1)Xx+X^{h+1}n',\qquad x,n'\in M .
\end{equation}
In \(M\)-coordinates the Frobenius term \(Q^{h+1}\varphi\) is contracting by \eqref{eq:B8-wach-heights}, so \(1-\varphi\) is onto on \(X^{h+1}M\). Apply \(p_\gamma\) to \(c_\gamma\ell^k\omega=(\gamma-1)(c_\gamma Y-(\varphi-1)Xx)-(\varphi-1)X^{h+1}n'\). The result is integral. In \(A/(Q^h)\), \(\ell\) acts injectively and \(X\) becomes invertible after inverting \(\ell\). Cancellation therefore gives \(p_\gamma\omega\in XM\). Apply \eqref{eq:B8-graded-reduction} again and invert \(1-\varphi\) to write \(c_\gamma p_\gamma\omega=d(Xx_0,X^{h+1}z_0)\). With \(S(U)=(P(U)-P(1))/(U-1)\) one obtains
\[
 c_\gamma P(1)\omega=d(Xx_0-c_\gamma S(\gamma)\omega,\ X^{h+1}z_0).
\]
Subtract this kernel pair from \(c_\gamma P(1)v\) and apply \eqref{eq:B8-finite-pairs}. The image is \(c_\gamma P(1)\ell^{2w}c\). Its valuation
\begin{equation}\label{eq:B8-wach-exponent}
 2w+\sum_{j\in J^*}(s+v_\ell(j))+hs+v_\ell(h!)
 \le32(2hs+1)
\end{equation}
can be seen term by term. The two Wach lifts contribute \(8hs+2v_\ell((4h)!)\). The nonzero shifts contribute at most \(3hs+v_\ell(h!)+v_\ell((2h)!)\), by counting the negative and positive shifts separately. The last product contributes \(hs+v_\ell(h!)\). Since \(v_\ell(k!)=\sum_{j\ge1}\lfloor k/\ell^j\rfloor\le k/(\ell-1)\), the total is at most \(12hs+12h/(\ell-1)\le24hs\). Restrict a target class to \(K'\), take its scaled finite lift there, and corestrict that lift. Its image is the same multiple of \([K':K]\) times the original class. Equation~\eqref{eq:B8-cyc-degree} bounds the additional valuation by \(j(N)+1\le D_h\). This proves the displayed numerical allowance also for \(h=0\).
\leanref{Formalism.\allowbreak{}LIFT.\allowbreak{}CocyclePolynomialSubmodules.\allowbreak{}main}{Formalism.\allowbreak{}LagrangeIdealLifts.\allowbreak{}CocyclePolynomialDivision.\allowbreak{}UClaim}{131--191}
\leanref{Formalism.\allowbreak{}MainNotions.\allowbreak{}LiftingTheory.\allowbreak{}lqb\_\allowbreak{}lifts}{Formalism.\allowbreak{}SharpRestriction.\allowbreak{}LevelLetterCocycles.\allowbreak{}BaseChangePreservesUnitaryOf}{4254--4352}
\leanref{Formalism.\allowbreak{}MainNotions.\allowbreak{}LiftingTheory.\allowbreak{}coreLiftB\_\allowbreak{}of'}{Formalism.\allowbreak{}SharpRestriction.\allowbreak{}LevelLetterCocycles.\allowbreak{}BaseChangePreservesUnitaryOf}{4357--4430}
\leanref{UpperTheory.\allowbreak{}N.\allowbreak{}liftK\_\allowbreak{}of}{Development.\allowbreak{}AutomorphicArithmetic.\allowbreak{}ResidueCharacters.\allowbreak{}LiftKOf}{12--41}

For the weight calculation take cyclotomic depth \(m'=2s+n_0+1\), where \(n_0\) is least with \(2h\le\ell^{n_0}(\ell-1)\). Then \(m'\le2s+j(2h)+2\). Separating weight \(i\) costs \(R_i=\sum_{j\in J,\ j\ne i}(m'+v_\ell(j-i))\) for the jump set of the lattice in question. The Hodge--Tate coordinate ball contains the \(\ell\)-multiple of its integral frame lattice. Scaling by \(\ell\) puts the differences in the ball, and the weight factors give a primitive after a further \(\ell^{1+\max(1,R_i)}\). The residual fixed vector belongs to weight \(i\). With the degree cost to reach this depth, the exponent to treat \(k\le h+1\) successive weights \(0,\ldots,k-1\) is bounded by
\begin{equation}\label{eq:B8-frame-cost}
 c(k)\le 2D_h+\sum_{i=0}^{k-1}(2+\max(1,R_i))
 \le 2D_h+3k(a+1)D_h .
\end{equation}
Indeed, for \(0\le i<k\le h+1\) and a jump \(j\ne i\), one has \(1\le|j-i|\le2h\). Thus \(R_i\le a(m'+j(2h))\le2aD_h\). When \(h=0\) the sum is empty. Each summand in \eqref{eq:B8-frame-cost} is at most \(3(a+1)D_h\). For the absent weight zero take \(k=1\). The bound \(2D_h+3(a+1)D_h\le I_{h,a}\) also kills the invariants over \(K\), since they form a subgroup of the invariants over \(K'\).
\leanref{GlobalTheory.\allowbreak{}FrameInv.\allowbreak{}pow\_\allowbreak{}smul\_\allowbreak{}mem\_\allowbreak{}of\_\allowbreak{}invariant\_\allowbreak{}mod}{Development.\allowbreak{}AbelianVarietyTheory.\allowbreak{}Proofs.\allowbreak{}InertiaEigenvalues.\allowbreak{}PrimePowSmulMemOfTwistedInvariantModLattice}{254--320}
\leanref{UpperTheory.\allowbreak{}N.\allowbreak{}invariants\_\allowbreak{}killed\_\allowbreak{}over\_\allowbreak{}cyclotomic\_\allowbreak{}subfield\_\allowbreak{}of\_\allowbreak{}admissible}{Development.\allowbreak{}AutomorphicArithmetic.\allowbreak{}ResidueCharacters.\allowbreak{}InvariantsKilledOverCyclotomic}{14--41}

For the character block let \(f(g)=gy_0-y_0\) in \(Y_0\cap\ker p\). Restrict to the cyclotomic subgroup and tensor \(f\) with \(1\) in the integral period quotient truncated through weight \(e\). Crystallinity of the extension makes this class torsion in the truncated quotient, and \eqref{eq:B8-frame-cost} with \(k=e+1\) makes it a coboundary there. In integral bases lift the coordinates of \(\phi\) to the coefficient field. The differences between the lifted map and its Galois translates are integral on the kernel lattice. Applying the lifts to the truncated primitive gives, for each scalar coordinate, the twisted difference of a period of filtration zero, modulo the integral period image and \(\mathrm{Fil}^{e+1}\). The scalar representative is a twisted cocycle modulo the integers. Apply Lemma~\ref{B8.41aa} over \(K'=K_0(\mu_{\ell^s})\) to each scalar coordinate. The resulting rational cocycle reduces pointwise to the prescribed coordinate multiplied by \(\ell^{s+v_\ell(e+1)}\), so it has the same scaled congruence in the period quotient. Apply the truncated primitive to take a representative with zero dual exponential at cost \((e+2)s+v_\ell(e)+\lfloor e/(\ell-1)\rfloor+1\). The zero dual exponential gives its \(B_{\rm cris}\) primitive. Thus the pushout is finite after the sum of these two costs and \(c(e+1)\), and corestriction adds at most \(j(N)+1\). Here \(v_\ell(0)=0\) in the scalar cost. Put \(D=D_h\). The first cost is at most \(2D+3t(a+1)D\) by \eqref{eq:B8-frame-cost}. The two scalar costs together are at most \((t+2)D+3t+1\), using \(s\le D\), \(v_\ell(e)\le e\), \(v_\ell(e+1)\le e+1\), and \(e<t\). The final degree costs at most \(D\). Their sum is at most \(40(t+1)(a+1)D=T_{h,a,t}\). For \(t=0\) there is no integer \(e\) in the stated range.
\leanref{GlobalTheory.\allowbreak{}Cyc.\allowbreak{}ClassBound.\allowbreak{}class\_\allowbreak{}bound}{Development.\allowbreak{}AbelianVarietyTheory.\allowbreak{}Proofs.\allowbreak{}InertiaEigenvalues.\allowbreak{}CrystallineExtensionClassBound}{359--387}
\leanref{UpperTheory.\allowbreak{}N.\allowbreak{}TBLoc.\allowbreak{}tb\_\allowbreak{}part1}{Development.\allowbreak{}AutomorphicArithmetic.\allowbreak{}ResidueCharacters.\allowbreak{}LenLeTorsion}{176--204}
\leanref{UpperTheory.\allowbreak{}N.\allowbreak{}LocNum.\allowbreak{}cLL\_\allowbreak{}le}{Development.\allowbreak{}AutomorphicArithmetic.\allowbreak{}ResidueCharacters.\allowbreak{}LocNum}{48--69}
\leanref{UpperTheory.\allowbreak{}N.\allowbreak{}LocNum.\allowbreak{}cTB\_\allowbreak{}le}{Development.\allowbreak{}AutomorphicArithmetic.\allowbreak{}ResidueCharacters.\allowbreak{}LocNum}{80--117}
\leanref{UpperTheory.\allowbreak{}N.\allowbreak{}tbK\_\allowbreak{}of}{Development.\allowbreak{}AutomorphicArithmetic.\allowbreak{}ResidueCharacters.\allowbreak{}TbKOf}{19--77}

By \eqref{eq:B8-BK-dim} the rational finite condition is zero for \(e>0\) and unramified for \(e=0\). In the latter case Frobenius evaluation lifts every divisible unramified class. The quotient in \eqref{eq:B8-tbk} injects into the \(\ell^{c_B}\)-torsion of \(H^1/H^1_f\). Duality identifies its dual with the integral finite group on \(O(e+1)\) modulo \(\ell^{c_B}\). Its free rank is at most \([K:\Q_\ell]\). Its torsion comes from the invariants in a rank-one divisible module and contributes at most \(c_B\) after reduction.
\leanref{UpperTheory.\allowbreak{}N.\allowbreak{}TBLoc.\allowbreak{}tbe\_\allowbreak{}clauses}{Development.\allowbreak{}AutomorphicArithmetic.\allowbreak{}ResidueCharacters.\allowbreak{}LenLeTorsion}{83--112}
\end{proof}

\subsection*{Proof of \ref{B8.41b}}\label{prf:B8.41b}\phantomsection
\begin{proof}
The common embedding sign in the jump conditions determines \(A,A'\) before \(\pi\) is chosen. Conjugate places have opposite types. The jump dictionary gives crystalline \(W\) on columns and crystalline \(W^\vee(1)\) on rows, with zero absent. The completions embed in \(K'\) of \ref{B8.41a}. That lemma gives the column invariant bound. Its exponents here are upper bounds for the cyclotomic estimates used in this split calculation. The canonical hub entries in \ref{B8.41h} are fixed by their separate formulas.
\leanref{UpperTheory.\allowbreak{}N.\allowbreak{}Two.\allowbreak{}decl\_\allowbreak{}split}{Development.\allowbreak{}AutomorphicArithmetic.\allowbreak{}ResidueCharacters.\allowbreak{}TwoDecl}{29--68}
\leanref{UpperTheory.\allowbreak{}N.\allowbreak{}decl\_\allowbreak{}split}{Development.\allowbreak{}AutomorphicArithmetic.\allowbreak{}ResidueCharacters.\allowbreak{}PlaceTypeConj}{69--108}

For the column module on \(g\) generators over the integral matrix algebra \(\mathcal A\), let \(\mathrm{Rel}\subset\mathcal A^g\) be the module of row vectors annihilating the generators. Take finitely many square arrays with every row in \(\mathrm{Rel}\), including arrays chosen for each block below. Form their linear combination with independent variables. At a point off \(Z(i)\) there is a nonzero solution of the relation equations, so the determinant of this universal matrix vanishes as a polynomial there. Its coefficients therefore belong to the Hecke ideal supported on \(Z(i)\). At the target point the matrix action splits by blocks and the determinant is a product of their determinants. In a DVR the least coefficient valuation is additive on products, by dividing out each least power and reducing to nonzero polynomials over the residue field. In each block the matrix span gives an array whose evaluated rows are an elementary-divisor basis of relations. Specializing to it bounds that block's least coefficient valuation by the length of its solution module in the corresponding power of \(K_i/O_i\). This bounds the length of the single evaluated congruence ideal by the sum of the block solution lengths.
\leanref{Formalism.\allowbreak{}RepresentationBoundsCohomology.\allowbreak{}ringDataObjects.\allowbreak{}displayFree}{Formalism.\allowbreak{}WachLatticeTwists.\allowbreak{}RamifiedBlockRings.\allowbreak{}RamifiedBlockProofs.\allowbreak{}CompositeOfChainHolds}{60--80}
\leanref{GlobalTheory.\allowbreak{}DetGen.\allowbreak{}length\_\allowbreak{}quot\_\allowbreak{}le\_\allowbreak{}sum\_\allowbreak{}length}{Development.\allowbreak{}AbelianVarietyTheory.\allowbreak{}Proofs.\allowbreak{}InertiaEigenvalues.\allowbreak{}DeterminantLemmaDiscreteValuationRingQuotientLen}{259--271}

The normalized full representation has blocks \(W\) and \((-e)\) for \(0\le e<t\). The separating projector has denominator, up to a unit, \(\Delta=\prod_{e=1}^{t-1}(1-u^{-e})\) for \(u\) of depth \(s_0(i)\); \eqref{eq:B8-cyc-val} gives \(v_\ell(\Delta)=\delta\). Evaluation on Galois differences gives cocycles. The quotient by the intersection of local kernels injects into the sum of their images. The classes satisfying the local conditions contribute \(s_\pi\) on \(W\) and the length \(s_{\ell,e}(E)\) of Definition~\ref{B8.40} on \((-e)\) for \(e>0\). On the trivial character they contribute \(v_\ell(h_E)\) by reciprocity. Inertia exclusion gives the local conditions at \(S_i\); at \(Q_i\) the character images are killed by \(\zeta_{i,v}-1\).

The column finite rank \(j_u\) of \(W\) is at most \(d_um\) by \eqref{eq:B8-BK-dim} and the number of contributing jumps. The free coordinates cost \(j_u(e'+\delta)\), and the coboundary kernel costs \(h_W(u)\). At a row the reflected jumps \(j\mapsto-1-j\) for the Tate dual give free cost \(d_wme'\), with invariant cost at most \(rc_0\). Pairing the rows with columns bounds these by \(\sum_{u\in A}(d_um e'+rc_0)\).
\leanref{Formalism.\allowbreak{}MainNotions.\allowbreak{}EmbeddingTypes.\allowbreak{}length\_\allowbreak{}bound\_\allowbreak{}of\_\allowbreak{}admitted}{Formalism.\allowbreak{}SharpRestriction.\allowbreak{}LevelLetterCocycles.\allowbreak{}JumpAndLengthOfAdmitted}{961--993}
\leanref{Formalism.\allowbreak{}MainNotions.\allowbreak{}EmbeddingTypes.\allowbreak{}sum\_\allowbreak{}lenW\_\allowbreak{}split}{Formalism.\allowbreak{}SharpRestriction.\allowbreak{}LevelLetterCocycles.\allowbreak{}JumpAndLengthOfAdmitted}{688--740}

At a column the trivial block lift and its projectors cost \(e'+2\delta\); the positive blocks have kernel cost at most \(h_0\) each. For each row the character-block lemma makes the cocycle finite after multiplication by \(\ell^{c_B}\). By \eqref{eq:B8-tbk}, the nonfinite quotient costs \((d_w+1)c_B\). One additional scalar direction costs \(c_B\). A dimension-\(d\) block has projector cost at most \(4d(|A|+1)\delta\), since clearing the projector on evaluation costs at most \(|A|d\delta\) and representatives killed by \(\Delta^2\) cost at most \(2d\delta\). Sum over the blocks, with total dimension \(n_V\), to obtain \eqref{eq:B8-split-quotient}.
\leanref{Formalism.\allowbreak{}NG.\allowbreak{}collectAlt}{Formalism.\allowbreak{}CharacterColumnKernels.\allowbreak{}IdempotentNormalForms.\allowbreak{}CharacterColumnKernelsOfBlockChiCOfBlockChiOne}{33--99}
\end{proof}

\subsection*{Proof of \ref{B8.42pa}}\label{prf:B8.42pa}\phantomsection
\begin{proof}
Choose a good point with its good list. Its squarefree degree-\(r\) label \(f_v\) at \(v\mid g_k\) factors as \(\prod_j f_{v,j}^{d_j}\) by the parameter parts of ranks \(a_j\) and multiplicities \(d_j\). Since \(\deg f_{v,j}\le a_j\) and \(\sum_jd_ja_j=r\), all degrees agree. Squarefreeness gives \(d_j=1\). Two or more parts would give a common divisor degree \(1\le a_j\le r/2\) at every \(v\mid g_k\), contrary to \ref{B7.51}.
\leanref{Formalism.\allowbreak{}MainNotions.\allowbreak{}AdmissibleThreePrime.\allowbreak{}cusp\_\allowbreak{}of}{Formalism.\allowbreak{}SharpRestriction.\allowbreak{}LevelLetterCocycles.\allowbreak{}ResidualIrreducibleOfTheoremThreePrime}{29--99}

Choose \(g=g_k\) with \(\ell\notin B(g_k,r)\), so \(g\ne\ell\). Since \(\rho\) realizes the list, its reduced eigenvalues at the tame generator \(t_v\) come from the prime-to-\(\ell\) components \(z_{\ell'}\) of roots \(z\) of \(f_v\), via the fixed torsion embedding. They are distinct by \ref{B7.51}. A stable subspace is a sum of eigenlines. The relation \(s_vt_vs_v^{-1}=t_v^g\) in this realization makes the corresponding roots stable under \(z\mapsto z^g\). The chosen roots and their complement give divisors of \(f_v\) over \(\F_g\). A proper nonzero subspace of dimension \(d\) gives degree \(\min(d,r-d)\) at every \(v\mid g\), which the good labels exclude.
\leanref{Formalism.\allowbreak{}CyclotomicNorms.\allowbreak{}CycleTypeFiniteOrder.\allowbreak{}t1}{Formalism.\allowbreak{}CyclotomicNorms.\allowbreak{}NonsquarefreeBounds.\allowbreak{}IrreducibilityOfRestriction}{427--637}

At either \(g_k\ne\ell\), a prescribed character eigenvalue would identify \(z_{\ell'}\) with the prime-to-\(\ell\) component of the corresponding character value or its inverse. Both are excluded by the root test. Twisting by \(\bar\xi\) preserves stable subspaces and leaves these inertia matrices unchanged.
\leanref{Formalism.\allowbreak{}MainNotions.\allowbreak{}AdmissibleThreePrime.\allowbreak{}not\_\allowbreak{}root\_\allowbreak{}of}{Formalism.\allowbreak{}SharpRestriction.\allowbreak{}LevelLetterCocycles.\allowbreak{}ResidualIrreducibleOfTheoremThreePrime}{129--172}

Since \(\rho\otimes\xi\) realizes \(x\), cancel the common monic character factor in its reduced Frobenius polynomials and those of the residual realization. Split density extends equality to their common finite quotient. Brauer--Nesbitt identifies the semisimplifications over \(\overline{\F}_\ell\) \cite{brauer-1941}. The irreducible side has the full dimension \(r\), so the other side is irreducible. Polynomial equality transfers the inertia exclusion.
\leanref{Formalism.\allowbreak{}MainNotions.\allowbreak{}N.\allowbreak{}residual\_\allowbreak{}irreducible\_\allowbreak{}of\_\allowbreak{}theorem}{Formalism.\allowbreak{}SharpRestriction.\allowbreak{}LevelLetterCocycles.\allowbreak{}JointConditionsZpowers.\allowbreak{}ResidualIrreducibleOfTheorem}{101--199}
\end{proof}

\subsection*{Proof of \ref{B8.42}}\label{prf:B8.42}\phantomsection
\begin{proof}
Take \(X_i\) to consist of points having a type vector \(F\) with \(e_0F\ne0\), where \(e_0\) is the good paired-character idempotent. The eigenspan tie gives the required inclusions. For a central idempotent \(e\) of the product type group \(\Gamma\), put \(M_0=|\Gamma|\). Then \((h^{-1}e)^{M_0+1}=h^{-1}e\). Its regular trace is \(M_0[h]e\) and its eigenvalues are zero or \(M_0\)-th roots of unity. Thus \(M_0e\) has integral cyclotomic coefficients constant on conjugacy classes. Evaluate at the chosen primitive root, and sum the resulting integral lifts over good tuples to obtain \(\alpha_i\).
\leanref{GlobalTheory.\allowbreak{}CentralIdem.\allowbreak{}centralIdemCoeff}{Development.\allowbreak{}AbelianVarietyTheory.\allowbreak{}Proofs.\allowbreak{}InertiaEigenvalues.\allowbreak{}TraceMemAdjoin}{49--62}
\leanref{UpperTheory.\allowbreak{}GalC.\allowbreak{}exists\_\allowbreak{}unit\_\allowbreak{}one\_\allowbreak{}add\_\allowbreak{}generator}{Development.\allowbreak{}AutomorphicArithmetic.\allowbreak{}Foundations.\allowbreak{}GaloisRepresentations.\allowbreak{}HeckeGaloisCriteria.\allowbreak{}ExistsUnitOneAddGenerator}{17--56}
For a lift with \(e_0F\ne0\), some summand of the good-tuple projector acts nontrivially. At a type prime \(g_k\) take its finite-type label \(f\). Goodness gives a squarefree polynomial with nonzero constant term and excludes the prescribed character value and its inverse from its roots. The assumed depth-zero type identity gives zero monodromy and the label \(f(X)(X-u)^t\), with \(u\in\F_{g_k}^\times\) the character value supplied by the identity. At the conjugate place the root labels are inverted. The tame characters attached to these roots are trivial on wild inertia. On a tame generator the label map is injective and multiplicative. Every repeated label therefore comes from \(u\), or from \(u^{-1}\) on the conjugate side, and its character value has order dividing \(g_k-1\). Complete splitting gives \(q_w=g_k\). These are the inertia conditions of the cut. The first provision identifies the other lists at the same spectral point up to permutation.
\leanref{UpperTheory.\allowbreak{}GalS.\allowbreak{}nc\_\allowbreak{}of\_\allowbreak{}label}{Development.\allowbreak{}AutomorphicArithmetic.\allowbreak{}Foundations.\allowbreak{}GaloisRepresentations.\allowbreak{}TypeConditions.\allowbreak{}CardRootsBar}{92--144}
\leanref{UpperTheory.\allowbreak{}GalC.\allowbreak{}nc\_\allowbreak{}of\_\allowbreak{}type}{Development.\allowbreak{}AutomorphicArithmetic.\allowbreak{}Foundations.\allowbreak{}GaloisRepresentations.\allowbreak{}HeckeGaloisCriteria.\allowbreak{}ExistsUnitOneAddGenerator}{143--152}

At such a place these tests pass to each cuspidal summand in the parameter list. Monodromy is a direct sum, the inertia polynomials multiply, and unramified twists preserve both tests. A repeated root of a factor is a repeated root of the product. For a cuspidal automorphic summand the normalized local component is generic, as in \ref{B9.07}, and \(E_w\simeq\Q_{g_k}\) by complete splitting. Lemma~\ref{B9.55a} gives depth-zero cuspidal support. Consequently the value of the weak correspondence used for this comparison is independent of its choice by \ref{B9.55}.
\leanref{UpperTheory.\allowbreak{}GalS.\allowbreak{}nc\_\allowbreak{}sum}{Development.\allowbreak{}AutomorphicArithmetic.\allowbreak{}Foundations.\allowbreak{}GaloisRepresentations.\allowbreak{}TypeConditions.\allowbreak{}AutomorphicGaloisAdmissible}{91--93}
\leanref{UpperTheory.\allowbreak{}GalS.\allowbreak{}nc\_\allowbreak{}tw}{Development.\allowbreak{}AutomorphicArithmetic.\allowbreak{}Foundations.\allowbreak{}GaloisRepresentations.\allowbreak{}TypeConditions.\allowbreak{}AutomorphicGaloisAdmissible}{49--50}
\leanref{UpperTheory.\allowbreak{}GalC.\allowbreak{}exists\_\allowbreak{}good\_\allowbreak{}alt}{Development.\allowbreak{}AutomorphicArithmetic.\allowbreak{}ThetaCorrespondence.\allowbreak{}ThetaPairLifts.\allowbreak{}ThirdPlaceSeries.\allowbreak{}ExistsGood}{15--54}

At \(w\mid\ell_i\) use the cut clause at type primes and unramified parameters elsewhere. Local compatibility and the dictionaries give admissible jumps, and the good spectrum gives residual existence. At \(S_i\) the type-prime exclusion is \ref{B8.42pa}. At the auxiliary prime cancel the common character factor from \((U-1)^r(U-\bar\chi_{i,1}(\gamma_{i,w}))^t\) for the specified inertia element. The rank-\(r\) polynomial is \((U-1)^r\) and \(\bar\chi_{i,1}(\gamma_{i,w})\ne1\).
\leanref{UpperTheory.\allowbreak{}GalS.\allowbreak{}eigPd\_\allowbreak{}holds\_\allowbreak{}alt}{Development.\allowbreak{}AutomorphicArithmetic.\allowbreak{}ThetaCorrespondence.\allowbreak{}ThetaPairLifts.\allowbreak{}ThirdPlaceSeries.\allowbreak{}EigPdHolds}{17--65}
Every residual cyclotomic twist of the normalized rank-\(r\) block and of its Tate dual has zero invariants at \(S_i\) by the exclusion. Local duality and \eqref{eq:B8-euler} give \(H^2=H^1=0\). This gives the fifth provision. Admissibility and residual existence give the sixth; together with the \(Z\)-conclusion these give the second. The first and the four ties restrict to the cut. Density, coefficient largeness and the jump and irreducibility conditions give its coefficient duties.
\leanref{UpperTheory.\allowbreak{}GalS.\allowbreak{}prvG5\_\allowbreak{}of\_\allowbreak{}admissible\_\allowbreak{}cutRec}{Development.\allowbreak{}AutomorphicArithmetic.\allowbreak{}ThetaCorrespondence.\allowbreak{}ThetaPairLifts.\allowbreak{}ThirdPlaceSeries.\allowbreak{}AdmissibleJOfRestricted}{28--44}
\leanref{GlobalTheory.\allowbreak{}StringRepresentationDecomposition.\allowbreak{}HOneTwo}{Development.\allowbreak{}AbelianVarietyTheory.\allowbreak{}Definitions.\allowbreak{}GaloisGeometry.\allowbreak{}HeckeCharEqOfLAdicRealisation}{129--136}
\leanref{Formalism.\allowbreak{}MainNotions.\allowbreak{}SumRepresentations.\allowbreak{}HStarZero\_\allowbreak{}iff}{Formalism.\allowbreak{}SharpRestriction.\allowbreak{}LevelLetterCocycles.\allowbreak{}ExistsInertiaCharpolyFactorOfTameType}{1479--1485}
\leanref{UpperTheory.\allowbreak{}GalC.\allowbreak{}facts\_\allowbreak{}of\_\allowbreak{}galois\_\allowbreak{}preconditions\_\allowbreak{}of\_\allowbreak{}cutRec}{Development.\allowbreak{}AutomorphicArithmetic.\allowbreak{}ThetaCorrespondence.\allowbreak{}ThetaPairLifts.\allowbreak{}ThirdPlaceSeries.\allowbreak{}CFacts}{45--57}
\leanref{UpperTheory.\allowbreak{}GalC.\allowbreak{}compatibilityTies\_\allowbreak{}eigenRestriction}{Development.\allowbreak{}AutomorphicArithmetic.\allowbreak{}Foundations.\allowbreak{}GaloisRepresentations.\allowbreak{}HeckeGaloisCriteria.\allowbreak{}CompatibilityTiesEigenRestriction}{86--88}

Choose the type threshold \(C_T\) and paired-character constant \(8C_T\) before choosing the cuts. Use the good cut at \(x>x_1\) and the full cut elsewhere. The transported family is twist-good.
\leanref{GlobalTheory.\allowbreak{}ExtraFromFacts.\allowbreak{}main\_\allowbreak{}alt\_\allowbreak{}alt}{Development.\allowbreak{}AutomorphicArithmetic.\allowbreak{}ThetaCorrespondence.\allowbreak{}ThetaPairLifts.\allowbreak{}ThirdPlaceSeries.\allowbreak{}ImportsExtraAlt}{18--27}
Use \eqref{eq:B8-nu-sum} and prime factorization for \(h_E,M_0\) to sum \eqref{eq:B8-cost-prime} for this family. The square term satisfies
\begin{equation}\label{eq:B8-prime-square-sum}
 \sum_{\ell\ {\rm prime}}\frac{2\log\ell}{\ell^2}\le
 \sum_{n\ge2}\frac{2\log n}{n^2}\le6 .
\end{equation}
Thus \eqref{eq:B8-cost-sum} gives \eqref{eq:B8-linkN} on the cut, under its finiteness and summability hypotheses.
\leanref{GlobalTheory.\allowbreak{}NLinkClose.\allowbreak{}main\_\allowbreak{}alt\_\allowbreak{}alt}{Development.\allowbreak{}AutomorphicArithmetic.\allowbreak{}ThetaCorrespondence.\allowbreak{}ThetaPairLifts.\allowbreak{}ThirdPlaceSeries.\allowbreak{}ImportsFurtherTail}{22--109}
\leanref{GlobalTheory.\allowbreak{}LinkFam.\allowbreak{}main}{Development.\allowbreak{}AbelianVarietyTheory.\allowbreak{}Proofs.\allowbreak{}InertiaEigenvalues.\allowbreak{}LinkNAlongFamilyOfLocalBounds}{11--17}

Project \(U_i/L_i\) to its cut counterpart. For a kernel representative \(u\), lift its cut image to an integral lattice vector \(m_0\). Then \(\alpha_i(u-m_0)=0\), hence \(M_0u=\alpha_i m_0\in L_i\). The preimage of this kernel is free of rank at most the dimension of the good sum. Indices without good points contribute zero. Consequently
\begin{equation}\label{eq:B8-eta-cut}
 \eta_j(w)\le\eta_j^c(w)+v_{\ell(w)}(M_0)
 \sum_{i\in I_{j,w}}\sum_{\pi\in Z_{\rm good}(i)}\mathrm{mult}(i,\pi).
\end{equation}
By \eqref{eq:B8-mult-rank}, the rank row and \ref{B4.52}, this gives
\begin{equation}\label{eq:B8-cut-sum}
 \eta_j(\mathcal L)\le\eta_j^c(\mathcal L)+d_{j,0}\log M_0 .
\end{equation}
\leanref{UpperTheory.\allowbreak{}GalC.\allowbreak{}Alg.\allowbreak{}length\_\allowbreak{}le}{Development.\allowbreak{}AutomorphicArithmetic.\allowbreak{}Foundations.\allowbreak{}GaloisRepresentations.\allowbreak{}HeckeGaloisCriteria.\allowbreak{}LengthLeOfSpan}{186--292}
\leanref{UpperTheory.\allowbreak{}GalC.\allowbreak{}etaJ\_\allowbreak{}le}{Development.\allowbreak{}AutomorphicArithmetic.\allowbreak{}Foundations.\allowbreak{}GaloisRepresentations.\allowbreak{}HeckeGaloisCriteria.\allowbreak{}IsHeckeGaloisIff}{73--115}
\leanref{UpperTheory.\allowbreak{}GalC.\allowbreak{}sum\_\allowbreak{}val\_\allowbreak{}le}{Development.\allowbreak{}AutomorphicArithmetic.\allowbreak{}Foundations.\allowbreak{}GaloisRepresentations.\allowbreak{}HeckeGaloisCriteria.\allowbreak{}EvPowLeExp}{16--45}
\leanref{UpperTheory.\allowbreak{}GalC.\allowbreak{}etaSum\_\allowbreak{}le\_\allowbreak{}add}{Development.\allowbreak{}AutomorphicArithmetic.\allowbreak{}Foundations.\allowbreak{}GaloisRepresentations.\allowbreak{}HeckeGaloisCriteria.\allowbreak{}EvPowLeExp}{48--57}
The cut preserves \(N_c,d_{j,0},v_\ell(D)\). Eventually \(\log M_0\le C_1x^{C_1}L\le C_1bL^2\), since the member scale gives \(\varepsilon>0\). Absorb this in the family constant.
\leanref{UpperTheory.\allowbreak{}GalC.\allowbreak{}linkN\_\allowbreak{}of\_\allowbreak{}gd}{Development.\allowbreak{}AutomorphicArithmetic.\allowbreak{}Foundations.\allowbreak{}GaloisRepresentations.\allowbreak{}HeckeGaloisCriteria.\allowbreak{}LinkNOfGd}{38--93}
\end{proof}

\subsection*{Proof of \ref{B8.42a}}\label{prf:B8.42a}\phantomsection
\begin{proof}
For a good point take the parameter list and its good triple \(s\). Write \(f:O_0\to O_i\) for the coefficient map. Cancel the common character factors in the tied Hecke polynomials, and use \ref{B8.42p} in the finite quotients modulo every power of the maximal ideal. Separation gives trace equality of \(\rho_\pi\) and \(f_*\rho_s\) on \(G_E\). Absolute residual irreducibility by \ref{B8.42pa} gives the full matrix span over the residue field by matrix-algebra density \cite[Chapter~II]{jacobson-1956}, and over \(O_i\) by Nakayama.

Let \(\phi,\phi'\) be their group-algebra maps and \(s_1\) an \(O_i\)-linear section of \(\phi\). Put \(\theta=\phi's_1\). The map \(T(A)_{ab}=\operatorname{tr}(A\theta(E_{ba}))\) satisfies \(T\phi'=\phi\) by trace equality applied to \(x s_1(E_{ba})\), hence \(T\theta=1\). Finite equal ranks give invertibility, and \(\phi'=\theta\phi\) makes \(\theta\) an algebra isomorphism. Take a unit diagonal entry \(u=\theta(E_{11})_{yy}\) using its trace \(1\). The matrices \(U_{xa}=\theta(E_{a1})_{xy}\) and \(V_{bx}=\theta(E_{1b})_{yx}\) satisfy \(VU=uI\) and \(\theta(A)U=UA\). The character tie now gives the integral identification
\[
 \chi_{i,1}=f_*(\chi_0)\psi_L,\qquad
 \rho_\pi\otimes\chi_{i,1}^{-1}\simeq
 f_*(\rho_s\otimes\chi_0^{-1})\otimes\psi_L^{-1}.
\]
\leanref{GlobalTheory.\allowbreak{}TraceConjugacy.\allowbreak{}exists\_\allowbreak{}conj\_\allowbreak{}of\_\allowbreak{}surjective}{Development.\allowbreak{}AbelianVarietyTheory.\allowbreak{}Proofs.\allowbreak{}FoundationProofs.\allowbreak{}ExistsConjOfTraceEqResiduallyAbsIrreducibleDVR}{10--79}
\leanref{GlobalTheory.\allowbreak{}TraceConjugacy.\allowbreak{}exists\_\allowbreak{}conj\_\allowbreak{}of\_\allowbreak{}trace\_\allowbreak{}eq\_\allowbreak{}local}{Development.\allowbreak{}AbelianVarietyTheory.\allowbreak{}Proofs.\allowbreak{}FoundationProofs.\allowbreak{}ExistsConjOfTraceEqResiduallyAbsIrreducibleDVR}{93--134}
\leanref{GlobalTheory.\allowbreak{}PtCore.\allowbreak{}twistedRho\_\allowbreak{}conj\_\allowbreak{}latTpsi}{Development.\allowbreak{}AbelianVarietyTheory.\allowbreak{}Proofs.\allowbreak{}InertiaEigenvalues.\allowbreak{}TwistedRhoConjLatTpsi}{33--44}
The auxiliary twist \(\psi_L^{-1}\) is trivial on decomposition groups above \(\ell\) and preserves the required unramifiedness. The signature budget in the twist-good condition gives finite length at most \(\nu_\ell+a_s\) with \(a_s\ge0,\ \sum_s a_s\le\ell^{-2}\). The conductor-root tie gives the tame bound.
\leanref{GlobalTheory.\allowbreak{}PtSel.\allowbreak{}sel\_\allowbreak{}at}{Development.\allowbreak{}AbelianVarietyTheory.\allowbreak{}Proofs.\allowbreak{}InertiaEigenvalues.\allowbreak{}SameCharEqOne}{27--64}
\leanref{GlobalTheory.\allowbreak{}JointSix.\allowbreak{}main}{Development.\allowbreak{}AbelianVarietyTheory.\allowbreak{}Proofs.\allowbreak{}InertiaEigenvalues.\allowbreak{}PsiAtPsiAt}{23--38}
At odd \(\ell\nmid Q_Sqd\), the cyclotomic residue image is full and \(tr\) is odd. At \(q,d\) the image contains the squares and \((\ell-1)/2\nmid tr\). These give \(\nu_\ell=0\). For \(\ell\mid Q_S\), a unit of depth \(v_\ell(Q_S)\) trivial on the other conductor factors gives \(\nu_\ell\le v_\ell(Q_S)+v_\ell(tr)\) by \eqref{eq:B8-cyc-val}; summation proves \eqref{eq:B8-nu-sum}.
\leanref{UpperTheory.\allowbreak{}Mu.\allowbreak{}sum\_\allowbreak{}nu\_\allowbreak{}le}{Development.\allowbreak{}AutomorphicArithmetic.\allowbreak{}CyclotomicBounds.\allowbreak{}RootsOfUnity.\allowbreak{}NuRecord}{53--62}
The \(Z\)-identity bounds the contribution at \(i\) by the sum of dimensions of its good isotypic spaces. Localization uniqueness makes their type sets disjoint in a fibre. Their direct sum lies in \(M_j\), whose theta image is \(Z_j(\mu_w)\) and whose theta map is injective by the nonzero scalars. An empty fibre contributes zero.
\leanref{GlobalTheory.\allowbreak{}MultiplicityBound.\allowbreak{}sum\_\allowbreak{}mult\_\allowbreak{}le}{Development.\allowbreak{}AbelianVarietyTheory.\allowbreak{}Proofs.\allowbreak{}InertiaEigenvalues.\allowbreak{}GoodPointMultiplicityLeRankProof}{46--64}
\leanref{Formalism.\allowbreak{}MainNotions.\allowbreak{}J.\allowbreak{}theta\_\allowbreak{}injective\_\allowbreak{}of\_\allowbreak{}scalars}{Formalism.\allowbreak{}SharpRestriction.\allowbreak{}AdmissibleParameterJoins.\allowbreak{}NandDIVb}{38--43}
\end{proof}

\subsection*{Proof of \ref{B8.42ba}}\label{prf:B8.42ba}\phantomsection
\begin{proof}
Let \(m_i\) be the least \(m\) with \(Z^mN_i\subset P_i\). It exists and is at most \(H_*\). We first calculate the cost of lowering this power of \(Z\). Suppose \(bZ^{j+1}N_i\subset P_i\) for \(b\in\Z_\ell\). For \(x\in N_i\), apply \((\gamma-1)/Z\) to \(bZ^{j+1}x\). The result belongs to \(P_i\). Using \(\gamma x\equiv x\bmod ZN_i\) and \(\gamma Z/Z\equiv u\bmod Z\), the result equals
\[
 b(u^{j+1}-1)Z^jx+bZ^{j+1}x',\qquad x'\in N_i.
\]
The second term belongs to \(P_i\). Descending from \(m_i\) to \(j\) shows that \(Z^j\prod_{j<k\le m_i}(u^k-1)\) annihilates \(N_i/P_i\).

For \(q_0\nmid k\) the factor \(u^k-1\) is a unit. For \(q_0\mid k\) it has valuation \(1+v_\ell(k/q_0)\), by binomial expansion at \(u^{q_0}\). Thus the valuation of the product is the sum, over \(v\ge0\), of the numbers of multiples of \(q_0\ell^v\) in \((j,m_i]\). Each count is at most \((m_i-j)/(q_0\ell^v)+1\), and at most \(C_h\) of the counts are nonzero. Consequently
\[
 \ell^a Z^j N_i\subset P_i
 \quad\text{if}\quad
 0\le j\le m_i,\qquad
 a\ge (m_i-j)\ell/q_0^2+C_h.
\]
For \(j=0\) the integer valuation of the product itself is at most \(m_i\ell/q_0^2+C_h\).

Put \(t=\lceil m_i/q_0\rceil\) and \(L=C_h+2+t\). We claim that \(Z^tQ^LN_i\subset P_i\). Write \(Q=Z^{q_0}+\ell w\) with \(w\in A\), using the binomial coefficients of \((1+Z)^\ell\). A term of \(Z^tQ^L\) is divisible by \(\ell^{L-v}Z^j\), where \(j=t+q_0v\) and \(0\le v\le L\). If \(j\ge m_i\) it annihilates the quotient. For \(j<m_i\), use \(m_i\le q_0t\) to obtain
\[
 (m_i-j)\ell/q_0^2
 \le t(1-q_0^{-2})-v(1+q_0^{-1})\le t+2-v.
\]
The preceding valuation bound applies with \(a=L-v\). This proves the claim.

For \(x\in N_{i-1}\), first apply the claim to \(Q^{M_i}\varphi x\in N_i\). Multiply by \(Z^{H_i+C_h+2}\), and then use \(Q^{h_i}P_i\subset A\varphi P_{i-1}\). Since \(\varphi(Z)=ZQ\), this gives
\[
 \varphi(Z^{H_i+C_h+2+t}x)\in A\varphi P_{i-1}.
\]
The ring \(A\) is free over \(\varphi(A)\) on \(1,Z,\ldots,Z^{\ell-1}\), by monic division using \((1+Z)^\ell-1\). Linearized Frobenius is injective, so comparison of the coordinate on \(1\) puts \(Z^{H_i+C_h+2+t}x\) in \(P_{i-1}\). We have proved
\[
 m_{i-1}\le H_i+C_h+2+\lceil m_i/q_0\rceil
 \le H_i+C_h+3+m_i/q_0 .
\]
Iterate \(J\) times starting at \(c\). The first \(J\) numerators are at most \(2g+C_h+3\), and the remaining term is \(m_{c+J}/q_0^J\le2h/q_0^J\). Sum the geometric progression and apply the valuation estimate at \(j=0\).
\end{proof}

\subsection*{Proof of \ref{B8.42bb}}\label{prf:B8.42bb}\phantomsection
\begin{proof}
Write \(y=\Lambda x\). At position \(i\) choose an \(O\)-basis adapted to \(V_i\), which is a direct summand by saturation. In each complementary coordinate divide \(x\) by the distinguished polynomial defining \(D_i\). The remainder has integral coefficients. If it is constant after inverting \(\varpi\), its positive-degree coefficients already vanish and its constant coefficient belongs to \(O\). When \(b_i=0\) the residue condition is empty. We thus have \(v_i\in V\) with
\[
 x-v_i\in(V_i)_A+D_iV_A .
\]

Let \(I=D_0A+D_1A\). The polynomials defining \(D_0,D_1\) are relatively prime over the fraction field. Clearing the denominators in their polynomial B\'ezout identity shows that \(I\cap O=\varpi^\alpha O\) for some \(\alpha\ge0\). Choose \(f_\alpha\in D_1A\) with \(f_\alpha-\varpi^\alpha\in D_0A\).

Subtract the two residue conditions. They give
\(v_0-v_1\in(V_0+V_1)_A+IV_A\). We claim that its constant vector lies in \(V_0+V_1+\varpi^\alpha V\). Take Smith coordinates for \(V_0+V_1\subset V\). In one coordinate the assertion reads
\[
 c\in O,\quad c=ag+h,\quad a\in O,\quad g\in A,\quad h\in I
 \quad\Longrightarrow\quad c\in aO+\varpi^\alpha O.
\]
For \(a=0\) this is the formula for \(I\cap O\). If \(a\mid c\) it holds directly. Otherwise \(v_\varpi(a)>v_\varpi(c)\), and \(1-(a/c)g\) is a unit of \(A\). Since its product with \(c\) is \(h\), we have \(c\in I\cap O\). This proves the claim in every coordinate.

Choose \(a_i\in V_i\) and \(v'\in V\) with
\(v_0-v_1=a_0'+a_1'+\varpi^\alpha v'\), where \(a_i'=a_i\). Put \(v_*=v_1+a_1'\). Modulo \(D_1\) the vector \(x-v_*-f_\alpha v'\) belongs to \((V_1)_A\). Modulo \(D_0\) it belongs to \((V_0)_A\), since \(v_*+\varpi^\alpha v'=v_0-a_0'\). Hence
\[
 y=\Lambda v_*+\Lambda f_\alpha v'+\nu,\qquad \nu\in{\cal N}(V).
\]
Choose \(w_*,w'\in W\) lifting \(\ell^{2a_0}v_*,\ell^{2a_0}v'\), respectively. The vector \(y_1=\Lambda(w_*+f_\alpha w')\) has constant residues \(w_*\) and \(w_*+\varpi^\alpha w'\) at the two positions. It belongs to \({\cal I}(W)\) and gives the asserted equality. The two constant vectors were lifted with the same power of \(\ell\).
\end{proof}

\subsection*{Proof of \ref{B8.42bc}}\label{prf:B8.42bc}\phantomsection
\begin{proof}
Counting multiples of \(\ell^j\) gives \(v_\ell(k!)=\sum_{j\ge1}\lfloor k/\ell^j\rfloor\le k/2\). Thus \(R\le6x+3/2\) and \(D\le6x\). Also
\[
 x\le \tfrac32 eh/\ell+\tfrac12 eS+1,\qquad
 \log_\ell h'\le\log_\ell h+S+1 .
\]
The first inequality uses \(e/(\ell-1)\le1\) in the constant term. For the second, \(h+S+1\le\max(1,h)\ell^{S+1}\).

Write \(j'=\lfloor\log_\ell h'\rfloor\). If \(h'<\ell\), then \(j'=0\). Otherwise \(h+S+1\ge\ell\), whence \(3h/\ell+S\ge2\). To see this for \(S=0,1\), substitute \(h\ge\ell-S-1\). For \(S\ge2\) it follows from \(h\ge0\). The logarithmic inequality therefore gives, in both cases,
\[
 j'\le\log_\ell h+\tfrac32 S+\tfrac32 h/\ell .
\]
Bounding the maximum in \(Y\) by the sum of its two entries now gives
\[
 \begin{split}
 Y&\le48x+\tfrac92+\eta+\log_\ell h'
                  +e\bigl(\tfrac32S+j'\bigr)\\
  &\le\tfrac{147}{2}eh/\ell+(27e+1)S
       +(e+1)\log_\ell h+\tfrac{107}{2}+\eta
   \le80B .
 \end{split}
\]
The negative and positive elements of \(J\) contribute at most \(L(h)\) and \(L(2h)\), respectively. Hence \(E'\le2L(4h)+2L(h)+L(2h)\le18h\).

The quotient \(C/C_0\) is a quotient of at most \(2(ne+1)r\) copies of \(O/\ell^YO\). The submodule \(C_0\) is a quotient of at most \(ne(e+1)m\) copies of \(O/\ell^{E'}O\). Add their normalized lengths, use \(Y\le80B\) in the first term and \(18h\le160(h+1)\) in the second.
\end{proof}

\subsection*{Proof of \ref{B8.42bd}}\label{prf:B8.42bd}\phantomsection
\begin{proof}
In the Artin description of the host in Theorem~\ref{B2.46}, the factor at \(\ell\) acts through its quadratic residue character. The other factors of the conductor are prime to \(\ell\) and give unramified local extensions. Thus the ramification index is at most two. The quadratic Gauss sum in the proof of Theorem~\ref{B2.46} belongs to \(E\), has square \(-\ell\), and belongs to \(\Q_\ell(\mu_\ell)\) under the chosen local embedding. Its valuation is \(1/2\) for \(v_\ell(\ell)=1\). Hence \(K_0(\sqrt{-\ell})\) has degree two over \(K_0\) and exhausts \(K\). The polynomial \(\Phi_\ell(1+Z)\) is Eisenstein over \(K_0\), giving the asserted index. In particular restriction to the level-one cyclotomic field followed by corestriction has multiplier prime to \(\ell\).

Use the pointwise local compatibility on the cut in Theorem~\ref{B8.42}. Above the coefficient prime, its clause away from the two type primes gives a crystalline full realization with the level-zero Hodge weights. It applies to every \(x\) here since \(q,d\) are distinct from \(g_0,g_1\). At a good point the semisimple decomposition of the full realization contains \(\rho_i(\pi)\) and the character lines \(\chi_{i,c}\), including \(c=1\).

For clarity, if \(0\to U\to V\to A\to0\) is exact with \(V\) crystalline, left exactness of crystalline periods gives
\[
 \dim D_{\rm cris}(V)
 \le \dim D_{\rm cris}(U)+\dim D_{\rm cris}(A)
 \le \dim U+\dim A .
\]
Here dimensions of periods are over \(K_0\) and representation dimensions are over \(\Q_\ell\), after restriction of coefficient scalars. The end terms agree. Both constituents are therefore crystalline. Applying this along a composition series gives crystallinity of the character and rank-\(r\) semisimple constituents above. Dualizing or tensoring crystalline comparison isomorphisms gives the comparison isomorphism for the dual or tensor product. This proves the claims about \(V_x,W_\pi\) and their Tate duals. Apply the same dimension argument to the specified subrepresentation of the direct sum. The de Rham filtration on a subrepresentation is the intersection filtration. Its graded pieces inject into those of the sum. Dualizing negates the jumps and the Tate twist subtracts one in the convention of Definition~\ref{B9.04}.
\end{proof}

\subsection*{Proof of \ref{B8.42bn}}\label{prf:B8.42bn}\phantomsection
\begin{proof}
Since \(t\le\eta\), we have \(0\le c_t\le\eta\) and \(t-1\le c_t+\eta\). The bounds on \(F_l\) therefore give
\[
 c_t-\eta\le m_l\le0,\qquad t-1\le M_l\le c_t+\eta,\qquad \eta'\le2\eta+1 .
\]
Hence \(\eta'+4\sigma+6\le2\eta+8B+7<\ell\), which is (P4). Every \(z\in F_l\cup F_c\) satisfies \(|z-c_t|\ge2B+t/2+3\); for \(z\in F_c\) this follows from \(|z-c_t|\ge\delta(z)\ge W_0\). An integer \(y'\in[0,t-1]\) has \(|y'-c_t|\le c_t\), so \(|z-y'|\ge2B+7/2>\sigma+2\). With \(x=z+w\) and \(y=y'+w\) this gives (P5), even for every \(x\in U\). The map \(z\mapsto z+w\) is injective on the set \(F_l\cup F_c\), which gives the multiplicities.

Let \(\varsigma=1\). An element of \(F_c\) satisfies \(z\ge c_t+W_0>0\ge m_l\). Hence every \(z\in F_l\cup F_c\) lies in \([m_l,M_a]\), and \(z+w\in[1,h-1]\). This is (P1). Next \(w=1-m_l\ge1\), the inequality \(w+t\le h\) reads \(t-1\le M_a\), and \(w+t-1\le\eta'\) reads \(t-1\le M_l\). These give (P2) and (P3). Every element of \(F_l\cup F_c\cup\{t-1\}\) is at most \(c_t+W_0+S+\eta+1\), and \(-m_l\le\eta-c_t\). Hence \(h\le W_0+S+2\eta+3\). For (P6) let \(x=z+w\notin[1,\eta']\). Then \(z\notin[m_l,M_l]\), so \(z\in F_c\) and \(\bar x=z+1-m_l\). Suppose \(\bar x\equiv m\pmod\ell\) with \(m\) in the range of (P6). Then \(\ell\) divides \(N=z+1-m_l-m\). The bounds \(m\le M_l-m_l+2\sigma+4\), \(M_l\le c_t+\eta\) and \(z\ge c_t+W_0\) give the lower bound below. The bounds \(m\ge-2\sigma-2\), \(-m_l\le\eta-c_t\) and \(z\le c_t+W_0+S+\eta+1\) give the upper bound:
\[
 W_0-\eta-2\sigma-3\le N\le W_0+S+2\eta+2\sigma+4 .
\]
Since \(\sigma\le\sigma_P\), the integer \(N\) lies in \eqref{eq:B8-pos-window}, a contradiction.

Let \(\varsigma=-1\). An element of \(F_c\) satisfies \(z\le c_t-W_0<c_t\le M_l\). Hence every \(z\in F_l\cup F_c\) lies in \([m_a,M_l]\), and \(z+w\in[1,h-1]\), which is (P1). Next \(w=1-m_a\ge1\), the inequality \(w+t\le h\) reads \(t-1\le M_l\), and \(h-\eta'=w+m_l\le w\). These give (P2) and (P3). Every element of \(F_l\cup F_c\cup\{0\}\) is at least \(c_t-W_0-S-\eta-1\), and \(M_l\le c_t+\eta\); hence \(h\le W_0+S+2\eta+3\). For (P6) let \(x=z+w\) lie outside \([h-\eta',h-1]=[w+m_l,w+M_l]\). Then \(z\in F_c\) and \(\bar x=h-x=M_l+1-z\). If \(\bar x\equiv m\pmod\ell\) with \(m\) in the range of (P6), then \(\ell\) divides \(N=M_l+1-z-m\), and
\[
 W_0-\eta-2\sigma-3\le m_l-z-2\sigma-3\le N\le M_l+1-z+2\sigma+2\le W_0+S+2\eta+2\sigma+4 ,
\]
by \(m\le M_l-m_l+2\sigma+4\), \(m_l\ge c_t-\eta\), \(c_t-W_0-S-\eta-1\le z\le c_t-W_0\) and \(M_l\le c_t+\eta\). Again \(N\) lies in \eqref{eq:B8-pos-window}, a contradiction.
\end{proof}

\subsection*{Proof of \ref{B8.42bod}}\label{prf:B8.42bod}\phantomsection
\begin{proof}
The substituted series defining \(\delta_z\) has zero constant coefficient and nonzero linear coefficient \(\bar z\). The binomial identities give \(\delta_z\delta_w=\delta_{zw}\), so \(\delta_{z^{-1}}\) is its inverse. These identities follow coefficientwise by approximating \(z,w\in\Z_\ell\) by nonnegative integers. Thus the displayed substitutions define an action on \(A_\pi\) and its fraction field.

If \(y\in B_\pi^\Delta\), choose \(r\ge0\) so that \(\iota(X)^ry\in A_\pi\). This product is invariant, hence belongs to \(\iota(A)\). Dividing by \(\iota(X)^r\) proves \(\iota(B)=B_\pi^\Delta\). The identity intertwining \(\gamma\) and \(\gamma_\pi\) extends from \(A\) to \(B\) by taking fractions. Finally write \(b=X^m a\), where \(m=v_X(b)\) and \(a\in A^\times\). Since \(\varphi\) preserves \(A\), it carries \(a\) and \(a^{-1}\) to mutually inverse elements of \(A\). Therefore
\[
 \varphi(b)=X^{\ell m}u^m\varphi(a),\qquad
 v_X(\varphi(b))=\ell m.
\]
Multiplication by \(X^a\) is invertible on \(E\), so a shifted lattice is again a lattice; the kernel intersection defining \(K_a\) is a \(k\)-subspace.
\end{proof}

\subsection*{Proof of \ref{B8.42boe}}\label{prf:B8.42boe}\phantomsection
\begin{proof}[Proof (outline)]
First, grouping Laurent coefficients by their exponents modulo \(\ell\) gives the basis \(1,\pi,\ldots,\pi^{\ell-1}\) over \(\varphi_\pi(B_\pi)\). The change to \(1,(1+\pi),\ldots,(1+\pi)^{\ell-1}\) is the triangular matrix \((\binom jm)_{0\le m,j<\ell}\), whose determinant is \(1\). This proves existence and uniqueness of the displayed expansion. Its zeroth coordinate defines \(\psi_\pi\), and
\[
 \psi_\pi(\varphi_\pi(a)b)=a\psi_\pi(b),\qquad
 \psi_\pi\varphi_\pi=1.
\]
The substitutions \(\delta_z\) permute the nonzero residue classes of the exponent \(j\) modulo \(\ell\); the substitution \(\gamma_\pi\) preserves those classes. Applying the expansion to these substitutions gives
\[
 \psi_\pi\delta_z=\delta_z\psi_\pi,\qquad
 \psi_\pi\gamma_\pi=\gamma_\pi\psi_\pi.
\]
Thus \(\psi_\pi\) preserves \(\iota(B)\) and descends to \(\psi\). This proves the three operator identities for \(\psi\).
\leanref{GlobalTheory.\allowbreak{}PsiLemmaCompare.\allowbreak{}assemble\_\allowbreak{}bijective}{Development.\allowbreak{}AbelianVarietyTheory.\allowbreak{}Proofs.\allowbreak{}InertiaEigenvalues.\allowbreak{}CyclotomicExistsIff}{158--184}
\leanref{GlobalTheory.\allowbreak{}PsiShiftProof.\allowbreak{}psi\_\allowbreak{}del}{Development.\allowbreak{}AbelianVarietyTheory.\allowbreak{}Proofs.\allowbreak{}InertiaEigenvalues.\allowbreak{}MemOfAssembleMem}{333--348}
\leanref{GlobalTheory.\allowbreak{}PsiShiftProof.\allowbreak{}exists\_\allowbreak{}psi}{Development.\allowbreak{}AbelianVarietyTheory.\allowbreak{}Proofs.\allowbreak{}InertiaEigenvalues.\allowbreak{}ExistsPullBasis}{31--40}

For every integer \(q\) and \(0\le r<\ell\), one has \(\psi_\pi(\pi^{\ell q+r})=(-1)^r\pi^q\). Hence \(v_\pi(\psi_\pi(y))\ge\lfloor v_\pi(y)/\ell\rfloor\). If \(2-\ell\le a\) and \(x\in X^aA\), then \(v_\pi(\iota x)\ge(\ell-1)a\). The last bound gives \(v_\pi(\iota\psi(x))>-(\ell-1)\). Every nonzero element of \(\iota(B)\) has valuation divisible by \(\ell-1\), so \(\psi(x)\in A\). Moreover, the constant term of \(\psi_\pi(\iota X)\) is the nonzero leading coefficient of \(\iota X\); therefore \(\psi(X)\in A^\times\). For \(b\in A\),
\[
 \psi\bigl(\varphi(b/\psi(X))X\bigr)=b.
\]
The argument of \(\psi\) lies in \(XA\subset X^aA\) when \(a\le1\). Thus \(\psi(X^aA)=A\) in the stated range.
\leanref{GlobalTheory.\allowbreak{}PsiShiftProof.\allowbreak{}psiRing}{Development.\allowbreak{}AbelianVarietyTheory.\allowbreak{}Proofs.\allowbreak{}InertiaEigenvalues.\allowbreak{}ExistsPullBasis}{78--98}

Choose a basis \(b_1,\ldots,b_n\) of the lattice \(M_0\). Since \(Bf(E)=E\), the vectors \(c_i=f(b_i)\) form a \(B\)-basis, and \(P=\bigoplus_iAc_i\). Define
\[
 \Psi\left(\sum_i x_ic_i\right)=\sum_i\psi(x_i)b_i.
\]
Semilinearity of \(f\) proves \(\Psi(bf(x))=\psi(b)x\). This identity determines \(\Psi\), since \(f(E)\) spans \(E\) over \(B\); hence the construction is independent of the basis and of the lattice \(M_0\). The relations \(fg_0=gf\) and \(\psi\gamma=\gamma\psi\) give \(\Psi g=g_0\Psi\). Applying \(\psi(X^aA)=A\) to each coordinate proves \(\Psi(X^aP)=M_0\).
\leanref{GlobalTheory.\allowbreak{}PsiShiftProof.\allowbreak{}exists\_\allowbreak{}lat\_\allowbreak{}basis}{Development.\allowbreak{}AbelianVarietyTheory.\allowbreak{}Proofs.\allowbreak{}InertiaEigenvalues.\allowbreak{}ExistsPullBasis}{206--217}
\leanref{GlobalTheory.\allowbreak{}PsiShiftProof.\allowbreak{}exists\_\allowbreak{}pull\_\allowbreak{}basis}{Development.\allowbreak{}AbelianVarietyTheory.\allowbreak{}Proofs.\allowbreak{}InertiaEigenvalues.\allowbreak{}ExistsPullBasis}{164--177}
\leanref{GlobalTheory.\allowbreak{}PsiShiftProof.\allowbreak{}exists\_\allowbreak{}Psi}{Development.\allowbreak{}AbelianVarietyTheory.\allowbreak{}Proofs.\allowbreak{}InertiaEigenvalues.\allowbreak{}ExistsPullBasis}{179--196}
\leanref{GlobalTheory.\allowbreak{}PsiShiftProof.\allowbreak{}pull\_\allowbreak{}span}{Development.\allowbreak{}AbelianVarietyTheory.\allowbreak{}Proofs.\allowbreak{}InertiaEigenvalues.\allowbreak{}ExistsPullBasis}{219--236}

Suppose \((g_0-1)M_0\subset XM_0\), and define \(G=(G_{ij})\) by \(g_0(b_j)=\sum_iG_{ij}b_i\). Then \(G-I_n\in XM_n(A)\). If \(x=\sum_jx_jc_j\), the \(i\)-th \(c\)-coordinate of \((g-1)x\) is
\begin{equation}\label{eq:B8-psi-coordinate-operator}
 T_G(x)_i=\sum_j\varphi(G_{ij})\gamma(x_j)-x_i.
\end{equation}
Also \(\Psi(x)=0\) exactly when every \(\psi(x_i)=0\).

Here is the scalar matrix estimate used in \eqref{eq:B8-psi-coordinate-operator}. For \(H-I_n\in\pi^2M_n(A_\pi)\), \(1\le j<\ell\), and \(s\in\Z\), the operator
\[
 y\longmapsto (1+\pi)^jH\gamma_\pi(y)-y
\]
is a bijection from \(\pi^sA_\pi^n\) to \(\pi^{s+1}A_\pi^n\). Indeed, its first coefficient is multiplication by \(j\), which is nonzero in \(k\); successive coefficients are determined recursively. The identities \(\gamma_\pi(\pi)-\pi\in\pi^\ell A_\pi\) and \(H-I_n\in\pi^2M_n(A_\pi)\) ensure that every later equation involves the next unknown with the same nonzero coefficient \(j\).
\leanref{GlobalTheory.\allowbreak{}PsiShiftProof.\allowbreak{}twist\_\allowbreak{}bij}{Development.\allowbreak{}AbelianVarietyTheory.\allowbreak{}Proofs.\allowbreak{}InertiaEigenvalues.\allowbreak{}MemOfAssembleMem}{457--475}

In the binomial expansion, \(\psi_\pi(y)=0\) means that its zeroth coordinate vanishes. The remaining coordinates transform by the preceding bijections. Furthermore,
\[
 \sum_{j=0}^{\ell-1}(1+\pi)^j\varphi_\pi(y_j)
 \in\pi^{\ell s}A_\pi
 \quad\Longleftrightarrow\quad
 y_j\in\pi^sA_\pi\quad\text{for every }j.
\]
Consequently \(y\mapsto\varphi_\pi(H)\gamma_\pi(y)-y\) maps
\(\pi^{\ell s}A_\pi^n\cap\ker\psi_\pi\) onto
\(\pi^{\ell(s+1)}A_\pi^n\cap\ker\psi_\pi\), and it is injective on the full kernel of \(\psi_\pi\).
\leanref{GlobalTheory.\allowbreak{}PsiShiftProof.\allowbreak{}op\_\allowbreak{}assemble}{Development.\allowbreak{}AbelianVarietyTheory.\allowbreak{}Proofs.\allowbreak{}InertiaEigenvalues.\allowbreak{}MemOfAssembleMem}{476--484}
\leanref{GlobalTheory.\allowbreak{}PsiShiftProof.\allowbreak{}core\_\allowbreak{}maps}{Development.\allowbreak{}AbelianVarietyTheory.\allowbreak{}Proofs.\allowbreak{}InertiaEigenvalues.\allowbreak{}MemOfAssembleMem}{496--506}
\leanref{GlobalTheory.\allowbreak{}PsiShiftProof.\allowbreak{}core\_\allowbreak{}inj}{Development.\allowbreak{}AbelianVarietyTheory.\allowbreak{}Proofs.\allowbreak{}InertiaEigenvalues.\allowbreak{}MemOfAssembleMem}{507--523}
\leanref{GlobalTheory.\allowbreak{}PsiShiftProof.\allowbreak{}core\_\allowbreak{}surj}{Development.\allowbreak{}AbelianVarietyTheory.\allowbreak{}Proofs.\allowbreak{}InertiaEigenvalues.\allowbreak{}MemOfAssembleMem}{524--541}

Apply this with \(H=\iota(G)\). The bound \(\ell\ge5\) gives \(H-I_n\in\pi^{\ell-1}M_n(A_\pi)\subset\pi^2M_n(A_\pi)\). For \(2-\ell\le a\le1\) and \(\psi(x)=0\),
\[
 x\in X^aA\quad\Longleftrightarrow\quad
 \iota(x)\in\pi^{\ell a}A_\pi.
\]
For \(a\le0\), this follows by comparing multiples of \(\ell-1\) with the two thresholds \((\ell-1)a\) and \(\ell a\). For \(a=1\), the coefficient of \(\pi^{\ell-1}\) vanishes because \(\psi_\pi(\iota x)=0\). The matrix operator commutes with \(\Delta\); its unique preimage of a \(\Delta\)-fixed vector is therefore \(\Delta\)-fixed. Descending through \(\iota\) yields
\[
 T_G\bigl((X^aA\cap\ker\psi)^n\bigr)
       =(X^{a+1}A\cap\ker\psi)^n
 \qquad(2-\ell\le a\le0).
\]
The coordinate formula and \(\Psi g=g_0\Psi\) now give \eqref{eq:B8-psi-shift-exact}. The complete proof is contained in \Lean{GlobalTheory.\allowbreak{}PsiShiftProof.\allowbreak{}lemmaPsi}.
\end{proof}

\subsection*{Proof of \ref{B8.42bof}}\label{prf:B8.42bof}\phantomsection
\begin{proof}[Proof (outline)]
Exactness makes \(D(j)_i\) injective and \(D(p)_i\) surjective. Identify \(D(U_1)_i\) with its image under \(D(j)_i\). Its intersection with \(N_i^b\) is finitely generated, since \(A\) is noetherian, and spans \(D(U_1)_i\), since multiplying any vector by a sufficiently large power of \(X\) puts it in \(N_i^b\). Its image under \(D(p)_i\) is therefore finitely generated and spans \(D(U)_i\). This proves the lattice assertion in \eqref{eq:B8-residue-model}. The residue comparison image defining \(\overline N_i(T)\) is likewise finitely generated over \(A\) and spans \(D(\overline T)_i\) over \(B\).

The cocycle identity is exactly the identity needed for the displayed action on \(E_z\) to be a group action. Exactness supplies lifts \(e_i\), and their two differences lie in \(D(U)_i\) because their images in \(D(k)_i\) vanish. Commutation of \(\varphi_i\) and \(\gamma_i\) gives the Herr relation. The same commutation shows that every denominator pair satisfies that relation, so the displayed quotient is defined. Replacing \(e_i\) by \(e_i+x_i\) adds precisely the denominator pair indexed by \((x_i)_i\), proving independence of the lifts. Finally, the reduction map \(\Z\to k\) defines the reduced multiset and its multiplicities without choices. The complete residue-lattice comparison used here is proved in \Lean{GlobalTheory.\allowbreak{}Present.\allowbreak{}RedLat.\allowbreak{}red\_\allowbreak{}isLat}.
\end{proof}

\subsection*{Proof of \ref{B8.42bobb}}\label{prf:B8.42bobb}\phantomsection
\begin{proof}
Put \(m=h_i-w_i-a\). By (P3),
\(0\le m\le h_i-w_i\le\eta'_i\). The assumed divisibility says
\(h_i-u\equiv m\pmod\ell\). If \(u\in U_i\) lies outside the low
window, this contradicts (P6), since
\[
 [0,\eta'_i]\subset[-2\sigma_i-2,\eta'_i+2\sigma_i+3].
\]
For a jump in the low window, or in the character block, (P1)--(P3)
give \(1\le h_i-u\le\eta'_i\). Hence
\(\lvert(h_i-u)-m\rvert\le\eta'_i<\ell\), where the last inequality
follows from (P4). An integer divisible by \(\ell\) and of absolute
value less than \(\ell\) is zero. Thus \(h_i-u=m\), which is the
required equality.
\end{proof}

\subsection*{Proof of \ref{B8.42bobc}}\label{prf:B8.42bobc}\phantomsection
\begin{proof}
The quotient map \(T_0\to T_0/\varpi T_0\), followed by the twist,
identifies the coefficient representation with \(\overline W\).
Choose a cocycle representing each class. The exponent-one hypothesis
makes its extension admissible, so the assumed model construction
supplies its model lattice. An invertible equivariant change of
coordinates preserves cocycles, coboundaries and linear independence.
This gives the first assertion.

The passage from a cocycle to its Herr class is linear. Suppose a
linear combination has zero Herr class. In its extension, subtract
suitable elements of \(D(\overline W)_i\) from the chosen lifts of
\(\mathbf1_i\). The resulting lifts \(e_i\) satisfy
\[
 \varphi_i(e_{i-1})=e_i,\qquad \gamma_i(e_i)=e_i.
\]
The fixed-vector property gives an equivariant map \(k\to E_z\)
whose realization takes \(1\) to \(e_i\). By \(k\)-linearity of the realization, its composite with \(E_z\to k\) is multiplication by a scalar; this scalar is \(1\), since the image of each \(e_i\) is \(\mathbf1_i\). Thus the extension splits and the original cocycle
is a coboundary. Linear independence of the original classes forces
all coefficients in the relation to be zero.
\end{proof}

\subsection*{Proof of \ref{B8.42bobe}}\label{prf:B8.42bobe}\phantomsection
\begin{proof}
Choose a Frobenius element \(F\) and the continuous unramified
homomorphism \(\nu:G_K\to\Z_\ell\) normalized by \(\nu(F)=1\).
Let \(\ell^aX=0\), with \(a\ge1\), and put \(q=\ell^a\). The
unramified character and its divided difference
\[
 \mu(g)=(1+q)^{\nu(g)},\qquad
 w(g)=\frac{\mu(g)-1}{q}
\]
satisfy \(w(g)\equiv\nu(g)\pmod{qO}\). Indeed,
\(\binom zj\in\Z_\ell\) for \(z\in\Z_\ell\), by continuity from
nonnegative integers. Thus every term
\(\binom{\nu(g)}j q^{j-1}\), for \(j\ge2\), belongs to
\(q\Z_\ell\). Every unramified homomorphism \(z:G_K\to X\) is therefore
\(z(g)=w(g)z(F)\).

In the sum of the characters \(\mu\) and \(1\), choose eigenvectors
\(e_\mu,e_1\), and form the stable lattice with basis
\(v=q e_\mu\), \(u=e_\mu+e_1\). Its action is
\[
 gv=\mu(g)v,\qquad gu=u+w(g)v.
\]
Modulo \(q\), this is the extension with cocycle \(w\bmod q\).
Twist by \(\chi_1\). The two rational characters are
\(\chi_1\mu\) and \(\chi_1\). Each is a summand of a full
representation: replace the \(\chi_1\)-summand in
\(V\oplus\bigoplus_cL'(\chi_c)\) by \(L'(\chi_1\mu)\) for the
first, and retain the original sum for the second. The jumps agree,
and \(\mu\equiv1\pmod\varpi\) preserves the residual characteristic
polynomials. Extend the displayed lattice by stable lattices in the
other summands. Its quotient supported on these two characters is
admissible. Push out its lower copy of \(O/qO\) by
\(1\mapsto z(F)\in X\). Closure under finite sums and quotients
makes the resulting extension admissible. As \(qO\) is a power of
\(\varpi O\), it is an extension at an exponent annihilating \(X\).
This proves the inclusion.

For the converse, write an admissible class as a continuous
homomorphism \(f:G_K\to X\), since the twist of \(\chi_1\) by
\(\chi_1^{-1}\) is trivial. Subtract the unramified homomorphism
\(g\mapsto\nu(g)f(F)\), and call the result \(h\). Its class remains
admissible and \(h(F)=0\). Suppose \(h\ne0\). Let \(Y\subset X\)
be the \(O\)-submodule spanned by its values. Nakayama's lemma gives
a nonzero \(O\)-linear map \(\lambda:Y\to k\). Since the values
of \(h\) span \(Y\), the homomorphism \(\lambda h\) is nonzero.
Pullback along \(Y\hookrightarrow X\), followed by functoriality
under \(\lambda\), makes its class admissible over \(k\).
The two admissible homomorphisms \(\bar\nu\) and \(\lambda h\)
are linearly independent: evaluation at \(F\) kills the first
coefficient of a relation, and evaluation at an element where
\(\lambda h\ne0\) kills the second. This contradicts the residue
dimension bound. Hence \(h=0\), so \(f\) is unramified. Finally,
evaluation at \(F\) identifies
\(\operatorname{Hom}_{\rm cont}(G_K/I_K,k)\) with \(k\), proving the
claimed dimension.
\end{proof}

\subsection*{Proof of \ref{B8.42bobf}}\label{prf:B8.42bobf}\phantomsection
\begin{proof}
Write \(S_i=\sum_{j\ge1}h_{i+j}\ell^{-j}\). Periodicity bounds the summands by a convergent geometric series, and
\(\ell S_{i-1}=h_i+S_i\). Thus the integer \(h_i+\sigma_i\) is strictly less than \(\ell(\sigma_{i-1}+1)\), which gives the carry inequality.

Let \(y\) lie in the displayed band. If \(\varsigma_i=1\), then
\(-\sigma_i-1\le y\le\eta'_i+\sigma_i+2\).
If \(\varsigma_i=-1\), the same inequalities hold for \(h_i-y\).
For a jump outside the low window, (P6) therefore excludes
\(u\equiv y\pmod\ell\). For a jump in the low window, the bounds just obtained give
\(\lvert u-y\rvert\le\eta'_i+\sigma_i+1<\ell\), by (P4).
The congruence would force \(u=y\). There is an integer
\(z\in[w_i,w_i+t-1]\) with \(\lvert y-z\rvert\le\sigma_i+2\), contradicting (P5). Taking \(y=w_i+a\) proves the vanishing of \(\rho_i(\bar a)\).

At an upward index, \(w_i\le\eta'_i<\ell\). In the first block sum, a jump contributes precisely when \(u\equiv j\pmod\ell\) for some \(0\le j<w_i\). Condition (P6) excludes every cloud jump. For the remaining jumps, \(1\le u\le\eta'_i<\ell\), so the congruence means \(u=j\). Each jump below \(w_i\) is counted once, with its multiplicity.

At a downward index, \(h_i-w_i\le\eta'_i<\ell\). In the second block sum a jump contributes when
\(u\equiv w_i+a\pmod\ell\), with \(0\le a\le h_i-w_i\).
For a cloud jump this gives
\(h_i-u\equiv h_i-w_i-a\in[0,\eta'_i]\), contrary to (P6).
For a low-window jump, both \(u\) and \(w_i+a\) lie in
\([h_i-\eta'_i,h_i]\), so the congruence is equality. The sum therefore counts exactly the jumps at least \(w_i\); adding \(d_i\) gives \(|U_i|\).
\end{proof}

\subsection*{Proof of \ref{B8.42bobg}}\label{prf:B8.42bobg}\phantomsection
\begin{proof}[Proof (outline)]
The reduced Wach realization supplies the model structure of the base pairs. The largest-model theorem, applied to the model \(N^{\rm std}\), gives a model \(N^{\max}\) with
\(N_i\subset N_i^{\max}\) for every model \(N\) and every \(i\).

For an admissible cocycle, apply the extension-model comparison to its extension by the trivial line. It gives a model \(N\) of \(\overline W\) and a representative \((b,g)\) satisfying the Herr relation, \(g_i\in XN_i\), the upward inclusion \(b_i\in X^{-w_i}N_i\), and the downward inclusion \(b_i\in X^{-(h_i-w_i)}A\varphi_i(N_{i-1})\). Each inclusion persists under \(N\subset N^{\max}\); Frobenius spans also preserve inclusions. This puts all representatives in the stated common lattice.

The height and \(\gamma\) conditions pass from a base pair to its inverse image under an injection and then to its image under a surjection. At a downward index, choose a base presentation of \(N^{\max}\). Its polynomial relation has \(f_i(0)\ne0\). The maps in the presentation commute with \(\gamma_i\) and are \(B\)-linear, so they commute with \(\nabla_i\) and with \(f_i(\nabla_i)\). Pulling back and pushing forward the relation proves the stated polynomial inclusion on \(N^{\max}\). Finally, a relation among the Herr classes would give a split extension for the corresponding linear combination, by the fixed-vector argument of Lemma~\ref{B8.42bobc}. The original residue classes are independent, so all coefficients in the relation vanish. The complete common-model reduction is contained in \Lean{GlobalTheory.\allowbreak{}Models.\allowbreak{}AdmBound.\allowbreak{}main}, using \Lean{GlobalTheory.\allowbreak{}Models.\allowbreak{}AdmBound.\allowbreak{}inZZ\_\allowbreak{}mono} and \Lean{GlobalTheory.\allowbreak{}Models.\allowbreak{}AdmBound.\allowbreak{}indep\_\allowbreak{}pairs}.
\end{proof}

\subsection*{Proof of \ref{B8.42bobh}}\label{prf:B8.42bobh}\phantomsection
\begin{proof}[Proof (outline)]
Suppose \(p\ne0\). Its one-dimensional target makes \(p\) surjective. Pull the standard lattice \(N^{\rm std}\) back under \(D(j)\), then push it forward under \(D(p)\). This gives the model
\[
 N_i^q=D(p)_i\bigl(D(j)_i^{-1}(N_i^{\rm std})\bigr)
       \subset D(k(\omega^{-e}))_i.
\]
Both \(N^q\) and the standard character lattice \(C^{(e)}\) satisfy the height bounds with shifts \(w_i,h_i-w_i\) and the condition \((\gamma_i-1)N_i\subset XN_i\).

Apply the carry comparison for these two lattices at \(i=0\), followed by the multiplicity inequality for the displayed subquotient of \(N^{\rm std}\). The character lattice has its single residue weight at \(\bar e\), while the standard target has reduced multiplicities \(r_0(U_0,w_0;\alpha)\). The comparison therefore supplies an integer \(a\) such that
\[
 |a|\le\sigma_0,\qquad
 r_0(U_0,w_0;\bar e+\bar a)>0.
\]
By the definition of this multiplicity there is \(u\in U_0\) with
\(u\equiv w_0+e+a\pmod\ell\). Since \(-1\le e<t\), the integer \(w_0+e+a\) belongs to the excluded band of Lemma~\ref{B8.42bobf}. This contradiction proves \(p=0\). The complete proof is contained in \Lean{GlobalTheory.\allowbreak{}Models.\allowbreak{}ModTh.\allowbreak{}hom\_\allowbreak{}chr\_\allowbreak{}eq\_\allowbreak{}zero}.
\end{proof}

\subsection*{Proof of \ref{B8.42bobi}}\label{prf:B8.42bobi}\phantomsection
\begin{proof}[Proof (outline)]
First let \(e=c-1\), so \(1\le e<t\), and let \(z\) be an admissible cocycle with coefficient line \(k(\omega^{-e})\). The extension-model comparison gives a Herr pair \((b,g)\) on a model of this line. The character-model equality identifies that model with \(C^{(e)}\). Put
\[
 J_i=\begin{cases}w_i,&\varsigma_i=1,\\
 h_i-w_i-e,&\varsigma_i=-1.
 \end{cases}
\]
These are nonnegative integers. The pair satisfies
\[
 g_i\in XC_i^{(e)},\qquad b_i\in X^{-J_i}C_i^{(e)},\qquad
 (\gamma_i-1)b_i=\varphi_i(g_{i-1})-g_i.
\]
Indeed, at a downward index use
\(A\varphi_i(C_{i-1}^{(e)})=X^eC_i^{(e)}\).
For every \(0\le j\le J_i\), (P3) and (P4) give
\[
 0<e+j\le\eta'_i<\ell.
\]

Choose an \(A\)-basis vector \(\epsilon_i\) of \(C_i^{(e)}\), and write
\(\varphi_i(\epsilon_{i-1})=X^ea_i\epsilon_i\),
\(\gamma_i(\epsilon_i)=c_i\epsilon_i\), with \(a_i\in A^\times\) and \(c_i\in1+XA\). Commutation gives
\[
 \varphi(c_{i-1})a_iX^e=\gamma(a_i)\gamma(X)^e c_i.
\]
The scalar identities of Definition~\ref{B8.42bod} imply
\(c_i\equiv1+eX\pmod{X^2A}\).
If the lowest term of \(b_i\) were \(dX^{-j}\epsilon_i\), with \(1\le j\le J_i\) and \(d\ne0\), the lowest term of \((\gamma_i-1)b_i\) would be
\(d(e+j)X^{1-j}\epsilon_i\).
Its coefficient is nonzero by the preceding inequality. This contradicts
\(\varphi_i(g_{i-1})-g_i\in XC_i^{(e)}\). Hence \(b_i\in C_i^{(e)}\).

On \(C=\prod_iC_i^{(e)}\), put
\((\Phi x)_i=\varphi_i(x_{i-1})\) and \((\Gamma x)_i=\gamma_i(x_i)\).
Frobenius sends \(X^jC\) into \(X^{\ell j+e}C\) for \(j\ge0\).
Consequently
\[
 m=-\sum_{n\ge0}\Phi^n b\in C
\]
converges and satisfies \((\Phi-1)m=b\). The element
\(r=(\Gamma-1)m-g\) lies in \(XC\). The Herr relation gives
\((\Phi-1)r=0\), and strict Frobenius contraction on \(XC\) gives \(r=0\). Thus \((b,g)\) is a Herr boundary. The fixed-vector property makes \(z\) a coboundary.

For the trivial coefficient line, the constant-pair reduction of a modelled extension gives a Herr representative \((b,g)\), vectors \(m_i\in D(k)_i\), and scalars \(c_i\in k\) satisfying
\[
 b_i=\varphi_i(m_{i-1})-m_i+c_i\mathbf1_i,
 \qquad g_i=\gamma_i(m_i)-m_i.
\]
Thus the class has a representative \((c_i\mathbf1_i,0)_i\).
If \(\sum_i c_i=0\), choose scalars \(d_i\) with
\(d_{i-1}-d_i=c_i\). Explicitly, for representatives \(0,\ldots,d'-1\), take \(d_0=0\) and
\(d_j=-\sum_{a=1}^jc_a\). The equation at \(i=0\) follows from the zero sum. Since both operators fix the distinguished units, this constant pair is the Herr boundary of \((d_i\mathbf1_i)_i\).

Given two admissible classes, choose such constant representatives and write their sums as \(s_0,s_1\). If \(s_0=0\), the first class is zero. Otherwise \(s_0\) times the second class minus \(s_1\) times the first has a zero-sum constant representative and is zero. Hence the residue dimension is at most one. The hypotheses of Lemma~\ref{B8.42bobe} are the crystalline sum and twist rules, the closure of admissible modules, the admissibility of \(O(\chi_1)\otimes X\) and the one-exponent criterion; these are part~(a) of Lemma~\ref{B8.42bob}, whose proof does not use the present lemma. The unramified construction of Lemma~\ref{B8.42bobe} then supplies a nonzero admissible class, so the dimension is one. The complete proofs are contained in \Lean{GlobalTheory.\allowbreak{}CharSplit.\allowbreak{}admissibleCocycle\_\allowbreak{}negCharPower\_\allowbreak{}mem\_\allowbreak{}coboundaries} and \Lean{GlobalTheory.\allowbreak{}TopSteps.\allowbreak{}TrivLine.\allowbreak{}main}.
\end{proof}

\subsection*{Proof of \ref{B8.42bob}}\label{prf:B8.42bob}\phantomsection
\begin{proof}[Proof (outline)]
\emph{Part (a).} Since \(V\) is crystalline with jump multiset \(U_i\) at \(\tau_i\), we have \(|U_i|=r\) for every \(i\in\Z/d'\Z\).
\leanref{GlobalTheory.\allowbreak{}T.\allowbreak{}card\_\allowbreak{}hodgeTate\_\allowbreak{}jumps\_\allowbreak{}crystalline\_\allowbreak{}eq\_\allowbreak{}dim}{Development.\allowbreak{}AbelianVarietyTheory.\allowbreak{}Proofs.\allowbreak{}InertiaEigenvalues.\allowbreak{}CardHodgeTateJumpsCrystallineEqDimStatement}{21--29}
\leanref{GlobalTheory.\allowbreak{}T.\allowbreak{}exists\_\allowbreak{}numerical\_\allowbreak{}galoisData\_\allowbreak{}localSituation}{Development.\allowbreak{}AbelianVarietyTheory.\allowbreak{}Proofs.\allowbreak{}InertiaEigenvalues.\allowbreak{}ExistsNumericalGaloisDataLocalSituationStatement}{23--36}
For an injective equivariant map \(f:N\hookrightarrow N'\) into an admissible module presented by \(p:L_1\twoheadrightarrow N'\) with \(\ker p=L_2\), the preimage \(L_0=\{x\in L_1:p(x)\in f(N)\}\) is a stable \(O\)-lattice presenting \(N\simeq L_0/L_2\); for a surjective equivariant map \(f:N\twoheadrightarrow N'\) from \(N\simeq L_1/L_2\), the kernel \(L_0=\ker(f\circ p)\) presents \(N'\simeq L_1/L_0\); direct sums of presentations present finite direct sums. Crystalline representations are closed under direct sums, crystalline character twists, and unramified twists, with jump multiplicities adding under direct sums and shifting under character twists. Thus \(V\oplus\bigoplus_{c=1}^t L'(\chi_c)\) is a full representation, so \(T_0\otimes_O\mathcal M\) and \(O(\chi_c)\otimes_O\mathcal M\) are admissible. Because \(W\) is free over \(O\), subquotient and product closure make the set of classes with \(\mathcal E_{z,a}\) admissible an \(O\)-submodule independent of \(a\ge1\) (whenever \(\varpi^a\mathcal M=0\)), functorial in \(\mathcal M\), and preserved under pullback along injections \(i:\mathcal M\hookrightarrow\mathcal M'\); hence a cocycle \(z\) is admissible in Definition~\ref{B8.42bo} if and only if its class lies in \(\operatorname{Adm}(W,\mathcal M)\).

\emph{Part (b).} Reducing the level-one Wach comparison over \(K(\mu_\ell)\) modulo \(\varpi\) yields an exact cyclotomic \((\varphi,\Gamma)\)-module functor from finite-dimensional \(k[G_K]\)-representations to étale modules over \(\prod_{i\in\Z/d'\Z}k((X))\), together with the reduced Wach lattice of every crystalline \(O\)-lattice.
\leanref{GlobalTheory.\allowbreak{}T.\allowbreak{}wachLattice\_\allowbreak{}usedProperties\_\allowbreak{}of\_\allowbreak{}admitted}{Development.\allowbreak{}AbelianVarietyTheory.\allowbreak{}Proofs.\allowbreak{}InertiaEigenvalues.\allowbreak{}WachLatticeUsedPropertiesOfAdmittedStatement}{7--9}
Use the cyclotomic scalar data and canonical modulo-\(\varpi\) residue realization \(D\) of Definition~\ref{B8.42bof} on the class \(\mathcal C\) of all finite-dimensional continuous \(k[G_K]\)-representations, taking \(\mathscr B\) to be the collection of pairs \((\overline T,\overline N(T))\) for stable \(O\)-lattices \(T\) in finite direct sums of \(\chi_1^{-1}\)-twisted full representations. For any \(U\in\mathcal C\) of the form \(\chi_1^{-1}M\) with \(M\) an admissible \(k[G_K]\)-module (in particular \(U=\overline W:=W_0/\varpi W_0\), \(U=k(\bar\chi_\ell^{1-c})\) for \(2\le c\le t\), \(U=k\) for \(c=1\), and every extension \(E_z\) representing a class in \(\operatorname{Adm}(-,k)\)), Definition~\ref{B8.42boa} and (a) supply a model lattice \(N=(N_i)_i\) of \(D_i=D(U)_i\) from \(\mathscr B\). For a base pair arising from a stable lattice in a sum of \(\mu\) full representations, put
\(J_i=U_i+\{w_i,\ldots,w_i+t-1\}\). Twisting by \(\chi_1^{-1}\) gives the jump multiset \(\mu(J_i-w_i)\), contained in \([-w_i,h_i-w_i]\). The integral Wach height relations are
\[
 Q^{w_i}\varphi^*N(T)_i\subset N(T)_i,
 \qquad Q^{h_i-w_i}N(T)_i\subset\varphi^*N(T)_i.
\]
Here \(Q=\Phi_\ell(1+Z)\), and \(\varphi^*N(T)_i\) is the span of Frobenius from the preceding component. A basis in the \(\Delta\)-invariants generates \(N(T)\). Since the reduction of \(Q\) is a unit times \(X\), these relations descend to
\[
 X^{h_i-w_i}N_i^b\subset P_i^b\subset X^{-w_i}N_i^b,
 \qquad P_i^b=A\varphi_i(N_{i-1}^b).
\]
\leanref{GlobalTheory.\allowbreak{}WachQN.\allowbreak{}qN\_\allowbreak{}of\_\allowbreak{}admitted}{Development.\allowbreak{}AbelianVarietyTheory.\allowbreak{}Proofs.\allowbreak{}InertiaEigenvalues.\allowbreak{}WachLatticeUsedPropertiesOfAdmitted}{83--93}
\leanref{GlobalTheory.\allowbreak{}BasePairs.\allowbreak{}base\_\allowbreak{}abc}{Development.\allowbreak{}AbelianVarietyTheory.\allowbreak{}Proofs.\allowbreak{}InertiaEigenvalues.\allowbreak{}ExistsQuotOfAdmissible}{243--261}

At a downward index put \(s=h_i-w_i+1\) and \(\nabla_i=X^{-1}(\gamma_i-1)\). Lemma~\ref{B8.42bobb} excludes every spurious zero residue among the differences \(a-(u-w_i)\), for \(0\le a<s\) and \(u\in J_i\). Applying the integral Sen polynomial relation at this cutoff, and separating its possible zero factor, gives a polynomial \(f_i\in k[Y]\) with
\[
 f_i(0)\ne0,
 \qquad\nabla_i f_i(\nabla_i)N_i^b\subset X^{h_i-w_i+1}N_i^b.
\]
\leanref{GlobalTheory.\allowbreak{}BasePairs.\allowbreak{}base\_\allowbreak{}e}{Development.\allowbreak{}AbelianVarietyTheory.\allowbreak{}Proofs.\allowbreak{}InertiaEigenvalues.\allowbreak{}ExistsQuotOfAdmissible}{281--294}
\leanref{GlobalTheory.\allowbreak{}Models.\allowbreak{}EPoly}{Development.\allowbreak{}AbelianVarietyTheory.\allowbreak{}Definitions.\allowbreak{}NonGaloisDescent.\allowbreak{}TwistedModelAdmissibleCocycles}{107--109}

By Lemma~\ref{B8.41ca} at depth \(1\), \((\gamma-1)N(T)\subset ZN(T)\). The equality
\(N(T)^\Delta\cap ZN(T)=\pi_0N(T)^\Delta\), with \(\pi_0\) reducing to \(X\), gives \((\gamma_i-1)N_i^b\subset XN_i^b\). Pulling this inclusion back under \(D(j)_i\) and pushing it forward under \(D(p)_i\) proves \((\gamma_i-1)N_i\subset XN_i\) for every model in \eqref{eq:B8-residue-model}. For \(-1\le e<t\), the rank-one character \(k(\bar\chi_\ell^{-e})\) has a lattice \(C_i^{(e)}\) with
\[
 A\varphi_i(C_{i-1}^{(e)})=X^eC_i^{(e)},
 \qquad (\gamma_i-1)C_i^{(e)}\subset XC_i^{(e)}.
\]
For \(0\le e<t\) these are model lattices, and \(C_i^{(0)}=A\mathbf1_i\). These statements follow by reducing the rank-one Wach equality. Saturation and compatibility with sums then permit the subquotient construction \eqref{eq:B8-residue-model} for every admissible extension. Lemma~\ref{B8.42bobc} transports independent admissible classes to independent modelled Herr classes.
\leanref{GlobalTheory.\allowbreak{}BasePairs.\allowbreak{}exists\_\allowbreak{}modTh\_\allowbreak{}of\_\allowbreak{}pres}{Development.\allowbreak{}AbelianVarietyTheory.\allowbreak{}Proofs.\allowbreak{}InertiaEigenvalues.\allowbreak{}ExistsModThOfPres}{28--71}
Put \(P_i=A\varphi_i(N_{i-1})\). Choosing \(B\)-linear coordinate isomorphisms \(D_j\xrightarrow{\sim}B^{\dim_k U}\) and applying Lemma~\ref{B8.42boe} on \(E=B^{\dim_k U}\) gives \(k\)-linear maps \(\Psi_i:D_i\to D_{i-1}\) with
\[
 \Psi_i(b\varphi_i(x))=\psi(b)x\quad(b\in B,\ x\in D_{i-1}),
 \qquad \Psi_i\gamma_i=\gamma_{i-1}\Psi_i,
 \qquad \Psi_i(X^aP_i)=N_{i-1}\quad(2-\ell\le a\le1),
\]
and the exact shift identity
\[
 (\gamma_i-1)(X^aP_i\cap\ker\Psi_i)
       =X^{a+1}P_i\cap\ker\Psi_i
       \qquad(2-\ell\le a\le0).
\]
Lemma~\ref{B8.42bobg} puts the Herr representatives in a common largest model, with the stated Frobenius, \(\gamma\) and downward polynomial conditions. Lemma~\ref{B8.42bobf} gives the carry recurrence, the excluded residue band and the standard block counts. The modelled-cocycle dimension estimate applies with the full model structure and (P4)--(P6): every independent family of admissible residue classes with coefficient module \(\overline W\) has cardinality at most
\[
 \sum_i d_i=\sum_i\sum_{u<w_i}\operatorname{mult}_{U_i}(u)=d_-.
\]
The fixed-vector argument transfers this estimate to \(\operatorname{Adm}(T_0,k)\). The character bounds, including the value one for \(\chi_1\), follow from Lemma~\ref{B8.42bobi}.
\leanref{GlobalTheory.\allowbreak{}Models.\allowbreak{}AdmBound.\allowbreak{}main}{Development.\allowbreak{}AbelianVarietyTheory.\allowbreak{}Proofs.\allowbreak{}InertiaEigenvalues.\allowbreak{}AdmissibleCocyclesCardLeSumDims}{73--82}
\leanref{GlobalTheory.\allowbreak{}BasePairs.\allowbreak{}real\_\allowbreak{}of\_\allowbreak{}presented}{Development.\allowbreak{}AbelianVarietyTheory.\allowbreak{}Proofs.\allowbreak{}InertiaEigenvalues.\allowbreak{}ExistsModThOfPres}{72--82}
\leanref{GlobalTheory.\allowbreak{}Models.\allowbreak{}ModTh}{Development.\allowbreak{}AbelianVarietyTheory.\allowbreak{}Definitions.\allowbreak{}NonGaloisDescent.\allowbreak{}TwistedModelAdmissibleCocycles}{117--159}
\leanref{GlobalTheory.\allowbreak{}Models.\allowbreak{}IsModel}{Development.\allowbreak{}AbelianVarietyTheory.\allowbreak{}Definitions.\allowbreak{}NonGaloisDescent.\allowbreak{}TwistedModelAdmissibleCocycles}{99--103}
\leanref{GlobalTheory.\allowbreak{}PsiShiftProof.\allowbreak{}lemmaPsi}{Development.\allowbreak{}AbelianVarietyTheory.\allowbreak{}Proofs.\allowbreak{}InertiaEigenvalues.\allowbreak{}ExistsPullBasis}{239--295}

\emph{Part (c).} Lemma~\ref{B8.42bobh} gives the character obstruction for every \(-1\le e<t\). Apply its cases \(e=0\) and \(e=-1\) as in Lemma~\ref{B8.42bobd}. They give the two equivariant-map exclusions and
\(H^0(K,\overline W)=H^2(K,\overline W)=0\).
\leanref{GlobalTheory.\allowbreak{}T.\allowbreak{}hom\_\allowbreak{}to\_\allowbreak{}characterLine\_\allowbreak{}eq\_\allowbreak{}zero}{Development.\allowbreak{}AbelianVarietyTheory.\allowbreak{}Proofs.\allowbreak{}InertiaEigenvalues.\allowbreak{}HomToCharacterLineEqZeroStatement}{8--10} By \eqref{eq:B8-euler} and Nakayama's lemma, \(H^1(K,W_0)\) is finite free over \(O\) of rank \(d'r\) and \(H^1(K,W_0)\otimes_O\mathcal M\xrightarrow{\sim}H^1(K,W_0\otimes_O\mathcal M)\) for every finite \(O\)-module \(\mathcal M\). Since \(H^0(K,\chi_1^{-1}V)=0\) and \(\chi_1\) has jump \(w_i\) at \(\tau_i\), \eqref{eq:B8-BK-dim} and saturation of integral finite cohomology make \(H^1_f(K,W_0)\) an \(O\)-direct summand of \(H^1(K,W_0)\) of rank \(d_-\). Crystalline lattice extensions \(0\to T_0\to E\to O(\chi_1)\to0\) realize all classes in \(\operatorname{im}(H^1_f(K,W_0)\otimes_O\mathcal M)\) inside \(\operatorname{Adm}(T_0,\mathcal M)\). At \(\mathcal M=k\), this image has \(k\)-dimension \(d_-\), so the bound \(\dim_k\operatorname{Adm}(T_0,k)\le d_-\) of (b) is an equality; devissage on the length of \(\mathcal M\) using (a) and freeness of \(H^1(K,W_0)\) gives the formula for \(\operatorname{Adm}(T_0,\mathcal M)\). For \(2\le c\le t\), \(\operatorname{Adm}(O(\chi_c),k)=\{0\}\) from (b) and injection-pullback from (a) give \(\operatorname{Adm}(O(\chi_c),\mathcal M)=\{0\}\) by induction on length. For \(c=1\), the crystalline rules and admissibility properties from (a) give the hypotheses of Lemma~\ref{B8.42bobe}. Its two-character lattice construction puts every unramified class in \(\operatorname{Adm}(O(\chi_1),\mathcal M)\). The residue bound of at most one then gives
\(\operatorname{Adm}(O(\chi_1),\mathcal M)=H^1_{\rm ur}(K,\mathcal M)\)
for every finite \(\mathcal M\). Evaluation at a Frobenius element identifies \(H^1_{\rm ur}(K,k)\) with \(k\), so the residue dimension in (b) is exactly one.

The complete proof is contained in \Lean{GlobalTheory.\allowbreak{}TopBridge.\allowbreak{}stmtAt\_\allowbreak{}of\_\allowbreak{}real}, \Lean{GlobalTheory.\allowbreak{}TopSteps.\allowbreak{}CharVanish.\allowbreak{}stmt\_\allowbreak{}of\_\allowbreak{}real}, and \Lean{GlobalTheory.\allowbreak{}TopSteps.\allowbreak{}admissibleUntwistedCharUnramified\_\allowbreak{}of\_\allowbreak{}residueRealisation}.
\end{proof}

\subsection*{Proof of \ref{B8.42bo1}}\label{prf:B8.42bo1}\phantomsection
\begin{proof}
Fix data \((K,L',O,\varpi,\tau_i,t,h,w,\eta',\varsigma,U,V,T_0,\chi_1,\dots,\chi_t)\) of Definition~\ref{B8.42bo} and a finite \(O\)-module \(\mathcal M\). For any stable \(O\)-lattice \(L\), continuous cocycles \(z:G_K\to L\otimes_O\mathcal M\) for the action \(g\mapsto\chi_1(g)^{-1}g\) coincide with continuous cocycles of the twisted coefficient module \((\chi_1^{-1}L)\otimes_O\mathcal M\).
\leanref{GlobalTheory.\allowbreak{}T.\allowbreak{}twistedAction\_\allowbreak{}eq\_\allowbreak{}twistedCoefficients\_\allowbreak{}cocycle\_\allowbreak{}iff}{Development.\allowbreak{}AbelianVarietyTheory.\allowbreak{}Proofs.\allowbreak{}InertiaEigenvalues.\allowbreak{}TwistedActionEqTwistedCoefficientsCocycleIffStatement}{21--28}

Put \(W_0=\chi_1^{-1}T_0\). By Lemma~\ref{B8.42bob}(a) and (c), a cocycle \(z:G_K\to T_0\otimes_O\mathcal M\) is admissible if and only if its class in \(H^1(K,W_0\otimes_O\mathcal M)\) lies in \(\operatorname{Adm}(T_0,\mathcal M)=\operatorname{im}\bigl(H^1_f(K,W_0)\otimes_O\mathcal M\to H^1(K,W_0\otimes_O\mathcal M)\bigr)\). Membership of \([z]\) in this image holds if and only if \(z=\sum_j\xi_j\otimes m_j+z_0\) for elements \(m_j\in\mathcal M\), a coboundary \(z_0\), and integral cocycles \(\xi_j:G_K\to T_0\) of \(\chi_1^{-1}V\) whose classes are finite in the sense of Definition~\ref{B8.41h}. This proves (T1).
\leanref{GlobalTheory.\allowbreak{}T.\allowbreak{}period\_\allowbreak{}data\_\allowbreak{}dual\_\allowbreak{}exact}{Development.\allowbreak{}AbelianVarietyTheory.\allowbreak{}Proofs.\allowbreak{}InertiaEigenvalues.\allowbreak{}PeriodDataDualExactStatement}{21--35}

Now let \(1\le c\le t\) and let \(L_c\) be a stable \(O\)-lattice in \(L'(\chi_c)\). Because \(O\) is a discrete valuation ring and \(L_c\) is a free rank-one \(O\)-module with \(G_K\) acting via \(\chi_c:G_K\to O^\times\), there is an \(O[G_K]\)-equivariant \(O\)-linear isomorphism \(e:O(\chi_c)\xrightarrow{\sim}L_c\). For any continuous cocycle \(z:G_K\to L_c\otimes_O\mathcal M\) for \(g\mapsto\chi_1(g)^{-1}g\), set \(z_0=(e^{-1}\otimes\operatorname{id}_{\mathcal M})\circ z:G_K\to O(\chi_c)\otimes_O\mathcal M\). Then \(z_0\) is a continuous cocycle, and for every integer \(a\ge1\) with \(\varpi^a\mathcal M=0\) the map \((e\otimes\operatorname{id}_{\mathcal M})\oplus\operatorname{id}_{O/\varpi^a}\) is an \(O[G_K]\)-equivariant isomorphism \(\mathcal E_{z_0,a}\xrightarrow{\sim}\mathcal E_{z,a}\). Thus \(z\) is admissible if and only if \(z_0\) is admissible, hence by Lemma~\ref{B8.42bob}(a) if and only if \([z_0]\in\operatorname{Adm}(O(\chi_c),\mathcal M)\). Moreover, \(e\otimes\operatorname{id}_{\mathcal M}\) preserves coboundaries on \(G_K\) and coboundaries on \(I_K\).
\leanref{GlobalTheory.\allowbreak{}T.\allowbreak{}exists\_\allowbreak{}equiv\_\allowbreak{}stableLattice\_\allowbreak{}integralCharacter}{Development.\allowbreak{}AbelianVarietyTheory.\allowbreak{}Proofs.\allowbreak{}InertiaEigenvalues.\allowbreak{}ExistsEquivStableLatticeIntegralCharacterStatement}{21--25}
\leanref{GlobalTheory.\allowbreak{}T.\allowbreak{}equivariantIso\_\allowbreak{}transport\_\allowbreak{}cocycles\_\allowbreak{}extensions}{Development.\allowbreak{}AbelianVarietyTheory.\allowbreak{}Proofs.\allowbreak{}InertiaEigenvalues.\allowbreak{}EquivariantIsoTransportCocyclesExtensionsStatement}{21--34}
\leanref{GlobalTheory.\allowbreak{}TopDict.\allowbreak{}admissibleExtension\_\allowbreak{}iff\_\allowbreak{}equivariantCoefIso}{Development.\allowbreak{}AbelianVarietyTheory.\allowbreak{}Proofs.\allowbreak{}InertiaEigenvalues.\allowbreak{}ExactLocalConditionsTopCore}{47--64}

If \(2\le c\le t\), Lemma~\ref{B8.42bob}(c) gives \(\operatorname{Adm}(O(\chi_c),\mathcal M)=\{0\}\). Hence \(z\) is admissible if and only if \([z_0]=0\), that is, \(z_0\) is a coboundary on \(G_K\), which holds if and only if \(z\) is a coboundary on \(G_K\). This proves (T2).
\leanref{GlobalTheory.\allowbreak{}TopDict.\allowbreak{}WachMod\_\allowbreak{}mem\_\allowbreak{}B1\_\allowbreak{}iff}{Development.\allowbreak{}AbelianVarietyTheory.\allowbreak{}Proofs.\allowbreak{}InertiaEigenvalues.\allowbreak{}ExactLocalConditionsTopCore}{65--72}

If \(c=1\), then \(\chi_1^{-1}\chi_1=1\), and Lemma~\ref{B8.42bob}(c) gives \(\operatorname{Adm}(O(\chi_1),\mathcal M)=H^1_{\rm ur}(K,\mathcal M)\). A cocycle \(z_0:G_K\to O(\chi_1)\otimes_O\mathcal M\) represents a class in \(H^1_{\rm ur}(K,\mathcal M)\) if and only if its restriction to \(I_K\) is a coboundary: if \(z'\) is cohomologous to \(z_0\) via \(z'(g)-z_0(g)=gb'-b'\) for all \(g\in G_K\) and \(z'(g)=gb-b\) for all \(g\in I_K\), then \(z_0(g)=g(b-b')-(b-b')\) for all \(g\in I_K\); conversely, if \(z_0(g)=gb-b\) for all \(g\in I_K\), then \([z_0|_{I_K}]=0\) in \(H^1(I_K,\mathcal M)\), so \([z_0]\in H^1_{\rm ur}(K,\mathcal M)\). Transporting along \(e\otimes\operatorname{id}_{\mathcal M}\) shows that \(z\) is admissible if and only if \(z|_{I_K}\) is a coboundary on \(I_K\). This proves (T3).
\leanref{GlobalTheory.\allowbreak{}IndefiniteRowSideFullHecke.\allowbreak{}AdmissibleCharUnramified.\allowbreak{}homModInertia\_\allowbreak{}eq\_\allowbreak{}unr}{Development.\allowbreak{}AbelianVarietyTheory.\allowbreak{}BlockTheory.\allowbreak{}NonGaloisCusps.\allowbreak{}AdmissibleCharUnramified}{37--47}
\leanref{GlobalTheory.\allowbreak{}TopDict.\allowbreak{}mk\_\allowbreak{}mem\_\allowbreak{}unr\_\allowbreak{}iff}{Development.\allowbreak{}AbelianVarietyTheory.\allowbreak{}Proofs.\allowbreak{}InertiaEigenvalues.\allowbreak{}ExactLocalConditionsTopCore}{32--46}
\leanref{GlobalTheory.\allowbreak{}TopDict.\allowbreak{}topCore\_\allowbreak{}of\_\allowbreak{}typings}{Development.\allowbreak{}AbelianVarietyTheory.\allowbreak{}Proofs.\allowbreak{}InertiaEigenvalues.\allowbreak{}ExactLocalConditionsTopCore}{75--172}
\leanref{GlobalTheory.\allowbreak{}TopBridge.\allowbreak{}twistedLatticeCohomology\_\allowbreak{}theoremTOP}{Development.\allowbreak{}AbelianVarietyTheory.\allowbreak{}Proofs.\allowbreak{}InertiaEigenvalues.\allowbreak{}TwistedLatticeCohomologyTheoremTOP}{67--83}
\end{proof}

\subsection*{Proof of \ref{B8.42bp}}\label{prf:B8.42bp}\phantomsection
\begin{proof}
\emph{Jump sets and size facts.} For a layer \(j\) and a weight type \(\tau\), the de Rham jump conditions of Theorem~\ref{B8.42} describe the jumps at an embedding of sign \(1\) as a set of integers formed by the low jumps, indexed by the \(m\) row parameters of \(\tau\), and the cloud jumps, indexed by its \(k-1\) column parameters. At an embedding of sign \(-1\) these sets are transported by \(u\mapsto t-1-u\). Let \(J_i\) be the resulting set for the type and sign at \(i\), with low part \(F_l(i)\) and cloud part \(F_c(i)\), and put \(c_t=(t-1)/2\). The size rows of the member give the following facts for every \(x>x_1\), every layer, type and sign.
(F1) Every \(z\in F_l(i)\) satisfies \(2b/Z+t/2+3\le|z-c_t|\le\eta\), and there is a base size \(W_0(i)\) with \(W_0(i)\le\varsigma_i(z-c_t)\le W_0(i)+S_c+\eta+1\) for every \(z\in F_c(i)\).
(F2) Every base size \(W_0\) satisfies \(2b/Z+t/2+3\le W_0\le b\).
(F3) \(1\le t\le\eta\), \(b>0\), \(1<Z\), \(2\eta+t+8b/Z+7<Z\) and \((b+S_c+2\eta+3)Z\le2b(Z-1)\).
Beyond a family cutoff one has also:
(F4) \(q\), \(d\) and the primes of the small set are less than \(Z\).
\leanref{GlobalTheory.\allowbreak{}T.\allowbreak{}lowJumps\_\allowbreak{}and\_\allowbreak{}cloud\_\allowbreak{}in\_\allowbreak{}windows\_\allowbreak{}along\_\allowbreak{}family}{Development.\allowbreak{}AbelianVarietyTheory.\allowbreak{}Proofs.\allowbreak{}InertiaEigenvalues.\allowbreak{}LowJumpsAndCloudInWindowsAlongFamilyStatement}{12--23}
\leanref{GlobalTheory.\allowbreak{}T.\allowbreak{}baseSize\_\allowbreak{}bounds\_\allowbreak{}along\_\allowbreak{}family}{Development.\allowbreak{}AbelianVarietyTheory.\allowbreak{}Proofs.\allowbreak{}InertiaEigenvalues.\allowbreak{}ParametersBaseSizeBoundsAlongFamilyStatement}{12--15}
\leanref{GlobalTheory.\allowbreak{}T.\allowbreak{}specialPrimeThreshold\_\allowbreak{}numerics\_\allowbreak{}along\_\allowbreak{}family}{Development.\allowbreak{}AbelianVarietyTheory.\allowbreak{}Proofs.\allowbreak{}InertiaEigenvalues.\allowbreak{}SpecialPrimeThresholdNumericsAlongFamilyStatement}{12--19}
\leanref{GlobalTheory.\allowbreak{}T.\allowbreak{}conductorPrimes\_\allowbreak{}lt\_\allowbreak{}specialPrimeThreshold\_\allowbreak{}eventually}{Development.\allowbreak{}AbelianVarietyTheory.\allowbreak{}Proofs.\allowbreak{}InertiaEigenvalues.\allowbreak{}ConductorPrimesLtSpecialPrimeThresholdEventuallyStatement}{12--16}

\emph{Part (a).} By (F3), \(Z>2\eta+t+7>8\). Hence every prime \(\ell>Z\) satisfies \(\ell\ge11\) and \(t+1<\ell\). By (F4), \(\ell\notin\{q,d\}\) and \(\ell\) lies outside the small set. The completions of \(E\) at such primes are unramified, by the description of the host in Theorem~\ref{B2.46}.
\leanref{GlobalTheory.\allowbreak{}UseLoc.\allowbreak{}exists\_\allowbreak{}caseA\_\allowbreak{}enl}{Development.\allowbreak{}AbelianVarietyTheory.\allowbreak{}Proofs.\allowbreak{}InertiaEigenvalues.\allowbreak{}FiniteDimensionalKsOne}{33--53}

Suppose \(\ell\) is not exceptional, and fix \(j\), \(d'\), the signs and the types. Fix \(i\) and put \(\eta=\eta_{\rm sz}\), \(S=S_c\), \(B=b/Z\), \(W_0=W_0(i)\) and \(\sigma_P=a_*\). The interval \eqref{eq:B8-pos-window} is then \eqref{eq:B8-exception-intervals} for the base size \(W_0(i)\), so \(\ell\) divides no integer in it. By (F1)--(F3), the other hypotheses of Lemma~\ref{B8.42bn} hold for \(F_l=F_l(i)\), \(F_c=F_c(i)\), \(\varsigma=\varsigma_i\) and \(\sigma=0\); in particular \(2\eta+8B+7<Z<\ell\). The construction of that lemma defines \((h_i,w_i,\eta'_i,U_i)\), which do not depend on \(\sigma\), and gives \(h_i\le W_0(i)+S_c+2\eta+3\). By (F2), \(h_c\le H_0\) for every \(c\), where \(H_0=b+S_c+2\eta+3\). Hence
\[
 \sigma_h(i)\le\sum_{j\ge1}H_0\ell^{-j}=\frac{H_0}{\ell-1}\le\frac{H_0}{Z-1}\le\frac{2b}Z ,
\]
by \(\ell>Z>1\) and the last inequality in (F3). Since \(\sigma_h(i)\) is an integer, \(\sigma_h(i)\le a_*\).
\leanref{Formalism.\allowbreak{}IrreducibleCountingLemmas.\allowbreak{}TOP.\allowbreak{}carryFloorLeGeom}{Formalism.\allowbreak{}PadicSubsetSums.\allowbreak{}StableGradedCounts.\allowbreak{}TOPSgLe}{19--41}
\leanref{Formalism.\allowbreak{}IrreducibleCountingLemmas.\allowbreak{}TOP.\allowbreak{}geomRatioBound}{Formalism.\allowbreak{}PadicSubsetSums.\allowbreak{}StableGradedCounts.\allowbreak{}TOPSgRange}{5--18}

Thus \(\sigma=\sigma_h(i)\) satisfies \(\sigma\le\min(2B,a_*)\), and Lemma~\ref{B8.42bn} with this \(\sigma\) gives (P1)--(P6) at \(i\) and the multiplicities.

\emph{Part (b), Hodge--Tate data.} Put \(K=E_w\), of degree \(d'=d_u\) over \(\Q_\ell\), unramified by (a). Number its embeddings \(\tau_\nu\), \(\nu\in\Z/d'\Z\), in Frobenius order. Since \(\ell\notin\{g_0,g_1\}\), the local compatibility on the cut makes every \(R_x\) crystalline at \(w\). As in the proof of Lemma~\ref{B8.42bd}, the semisimple decomposition of \(R_\pi\) makes \(W_\pi\) and the characters \(\chi_{i,c}\chi_{i,1}^{-1}\) crystalline at \(w\). The de Rham jump conditions give the jumps of \(W_\pi\) at \(\tau_\nu\) as the set \(J_\nu\) of (a), with multiplicity one, for the layer of the index and the sign at \(\tau_\nu\) given by the embedding type. They give the jumps of \(V_x\) at \(\tau_\nu\) as \(J_\nu\) and \(0,\dots,t-1\), with multiplicities added. By the weight dictionary, \(\chi_{i,c}\chi_{i,1}^{-1}=\chi_\ell^{1-c}\) on \(G_K\), with the single jump \(c-1\). Part (a) gives position data with \(\operatorname{mult}_{U_\nu}(b+w_\nu)=\mathbf 1_{J_\nu}(b)\).

By Lubin--Tate theory over the unramified field \(K\), there is a continuous crystalline character \(\theta:G_K\to O^\times\) with the single jump \(w_\nu\) at \(\tau_\nu\) for every \(\nu\). Products of Galois conjugates of the Lubin--Tate character of \(K\) realize every vector of jumps. Jumps add under tensor product with a crystalline character. Put \(V=W_\pi\otimes\theta\), with lattice \(T_0=W_\pi\otimes_{O_i}O\otimes\theta\), and \(\chi'_c=\chi_{i,c}\chi_{i,1}^{-1}\theta\). Then \(m_V(\tau_\nu,a)=\operatorname{mult}_{U_\nu}(a)\), \(\chi'_c=\chi'_1\chi_\ell^{1-c}\), and \(\chi'_c\) has the single jump \(w_\nu+c-1\). Each \(V_x\otimes\theta\) is crystalline with \(m_{V_x\otimes\theta}(\tau_\nu,a)=\mathbf 1_{J_\nu}(a-w_\nu)+\mathbf 1_{[0,t-1]}(a-w_\nu)=\tilde m_\nu(a)\).
\leanref{GlobalTheory.\allowbreak{}UseLocInst.\allowbreak{}exists\_\allowbreak{}setC}{Development.\allowbreak{}AbelianVarietyTheory.\allowbreak{}Proofs.\allowbreak{}InertiaEigenvalues.\allowbreak{}ExistsExactLocalConditionDataAtPlace}{22--111}
\leanref{GlobalTheory.\allowbreak{}T.\allowbreak{}crystalline\_\allowbreak{}characters\_\allowbreak{}given\_\allowbreak{}jumps\_\allowbreak{}unramified\_\allowbreak{}twists}{Development.\allowbreak{}AbelianVarietyTheory.\allowbreak{}Proofs.\allowbreak{}InertiaEigenvalues.\allowbreak{}CrystallineCharactersGivenJumpsUnramifiedTwistsStatement}{8--10}
\leanref{GlobalTheory.\allowbreak{}jTil\_\allowbreak{}shift\_\allowbreak{}block}{Development.\allowbreak{}AbelianVarietyTheory.\allowbreak{}Proofs.\allowbreak{}InertiaEigenvalues.\allowbreak{}JTilShiftBlock}{7--61}

\emph{The congruence.} Let \(x\in X_i\) and \(g\in G_E\). The points of \(X_i\) share the residual Hecke eigensystem of the index. At \(\pi\) the reduced full Hecke polynomial at a split place \(v\) is the residual Frobenius polynomial of \(\rho_i(\pi)\) times the character factors, by Lemma~\ref{B8.42pa}. Hence \(\det(U-\mathrm{Frob}_v\mid V_x)\) and \(\det(U-\mathrm{Frob}_v\mid W_\pi)\prod_c(U-\chi_{i,c}\chi_{i,1}^{-1}(\mathrm{Frob}_v))\) agree modulo \(\varpi\) at these places. Both sides, reduced modulo \(\varpi\), are constant on the cosets of an open normal subgroup. Lemma~\ref{B8.42p} therefore extends the congruence to every \(g\). Twisting by \(\theta\), each \(V_x\otimes\theta\) is a full representation for \((V,T_0,\chi'_1,\dots,\chi'_t)\) and the position data of (a).
\leanref{GlobalTheory.\allowbreak{}UseLoc.\allowbreak{}InstPoly.\allowbreak{}charpoly\_\allowbreak{}congr}{Development.\allowbreak{}AbelianVarietyTheory.\allowbreak{}Proofs.\allowbreak{}InertiaEigenvalues.\allowbreak{}ScaleRootsProdIcc}{222--259}

\emph{Lifting the cut cocycle.} Extend scalars from \(O_i\) to \(O\) throughout. Let \(L\) be the lattice \(M\otimes_{O_i}O\otimes\theta\) of \(M\otimes\theta\). On it the action \(\chi'_1(g)^{-1}g\) is the action \(\mu\) of \(M\), since \(\chi'_1=\theta\). The module \(\phi(\mathcal Y)\) is finitely generated and torsion, hence finite and killed by a power of \(\varpi\). Since \(G_E\) acts continuously on the lattice \(\mathcal Y+O_iv_0\), the cocycle \(\kappa\) is locally constant. Choose \(a\ge1\) with \(\varpi^a\bar\kappa=0\), put \(\mathcal M=\varpi^{-a}O/O\) and \(m_0=\varpi^{-a}\bmod O\). The \(\varpi^a\)-torsion of \(M\otimes L'/O\) is identified with \(L\otimes\mathcal M\) by \(u\otimes m_0\mapsto\varpi^{-a}u\). Hence \(\bar\kappa\) is the image of a unique cocycle \(z:G_K\to L\otimes\mathcal M\) for \(\mu\). Replace \(\mathcal Y\) by its \(O\)-submodule generated by the vectors \(gv_0-v_0\), \(g\in G_K\). It is finitely generated and stable, since \(g'(gv_0-v_0)=(g'g)v_0-v_0-(g'v_0-v_0)\). On it \(\phi\) takes values in the \(\varpi^a\)-torsion, so it defines \(\phi':\mathcal Y\to L\otimes\mathcal M\) with \(z(g)=\phi'(gv_0-v_0)\).
\leanref{GlobalTheory.\allowbreak{}T.\allowbreak{}cutCocycles\_\allowbreak{}locallyConstant\_\allowbreak{}killed\_\allowbreak{}by\_\allowbreak{}uniformiserPow}{Development.\allowbreak{}AbelianVarietyTheory.\allowbreak{}Proofs.\allowbreak{}InertiaEigenvalues.\allowbreak{}CutCocyclesLocallyConstantKilledByUniformiserPow}{12--29}
\leanref{GlobalTheory.\allowbreak{}T.\allowbreak{}exists\_\allowbreak{}integralTensorTorsion\_\allowbreak{}cocycle\_\allowbreak{}of\_\allowbreak{}torsion\_\allowbreak{}localCocycle}{Development.\allowbreak{}AbelianVarietyTheory.\allowbreak{}Proofs.\allowbreak{}InertiaEigenvalues.\allowbreak{}ExistsIntegralTensorTorsionCocycleOfTorsionLocalStatement}{10--18}

\emph{Admissibility.} The semisimple decomposition of \(R_\pi\) gives an injective equivariant map \(\iota:M\otimes K_i\to V_\pi\) and a nonzero vector \(e\in V_\pi\) fixed by \(G_E\), namely a generator of the summand \(c=1\). Let \(A=\mathcal Y\oplus L\oplus O\) with the action
\[
 g(y,u,s)=\bigl(\theta(g)(gy+s(gv_0-v_0)),\ gu,\ \theta(g)s\bigr),
\]
where \(gu\) is the action on \(L\). Put \(\mathcal V=\bigoplus_j(V_{x_j}\otimes\theta)\oplus(V_\pi\otimes\theta)^2\), a finite direct sum of full representations, and
\[
 \Psi(y,u,s)=\bigl((y_j+sv_{0,j})_j,\ \iota(u),\ se\bigr)\in\mathcal V .
\]
A direct computation gives \(\Psi(g\alpha)=g\,\Psi(\alpha)\), and \(\Psi\) is injective: \(se=0\) forces \(s=0\), \(\iota\) is injective, and then \(y=0\). For \(a''\ge a\) let \(q:A\to\mathcal E_{z,a''}\) be \(q(y,u,s)=(\phi'(y)+u\otimes m_0,\ s\bmod\varpi^{a''})\). It is surjective, since \(\mathcal M\) is generated by \(m_0\). It is equivariant, since
\[
 q(g(y,u,s))=\bigl(g(\phi'(y)+u\otimes m_0)+s\theta(g)z(g),\ \theta(g)s\bigr).
\]
Choose a stable \(O\)-lattice \(\Lambda\subset\mathcal V\) containing \(\Psi(A)\). By the Artin--Rees lemma there is \(n_0\) with \(\varpi^{n_1}\Lambda\cap \Psi(A)\subset\varpi^{n_1-n_0}\Psi(A)\) for \(n_1\ge n_0\). Take \(n_1=n_0+a''\) and put
\[
 L_1=\Psi(A)+\varpi^{n_1}\Lambda,\qquad L_2=\Psi(\ker q)+\varpi^{n_1}\Lambda .
\]
Both are finitely generated, stable and torsion free, and they span \(\mathcal V\), so they are stable lattices. Since \(\varpi^{a''}\) kills \(\mathcal E_{z,a''}\), we have \(\varpi^{n_1}\Lambda\cap \Psi(A)\subset\varpi^{a''}\Psi(A)\subset \Psi(\ker q)\). Therefore \(L_1/L_2\simeq \Psi(A)/\Psi(\ker q)\simeq\mathcal E_{z,a''}\) as \(O[G_K]\)-modules, for every \(a''\ge a\). Thus \(z\) is admissible for the data of the previous paragraphs.
\leanref{GlobalTheory.\allowbreak{}T.\allowbreak{}cocycleExtension\_\allowbreak{}twistedWeilCoefficients\_\allowbreak{}admissible}{Development.\allowbreak{}AbelianVarietyTheory.\allowbreak{}Proofs.\allowbreak{}InertiaEigenvalues.\allowbreak{}CocycleExtensionTwistedWeilCoefficientsAdmissibl}{60--86}
\leanref{GlobalTheory.\allowbreak{}T.\allowbreak{}cocycleExtension\_\allowbreak{}conductorTwistedCharacter\_\allowbreak{}admissible}{Development.\allowbreak{}AbelianVarietyTheory.\allowbreak{}Proofs.\allowbreak{}InertiaEigenvalues.\allowbreak{}CocycleExtensionConductorTwistedCharacterAdmissi}{53--87}
\leanref{UpperTheory.\allowbreak{}KN.\allowbreak{}exists\_\allowbreak{}admissible\_\allowbreak{}extension\_\allowbreak{}latticeTensor}{Development.\allowbreak{}AutomorphicArithmetic.\allowbreak{}CyclotomicBounds.\allowbreak{}DeRhamLattices.\allowbreak{}ExistsAdmissibleExtensionLatticeTensor}{25--125}

\emph{Applying (T).} If \(M=W_\pi\), (T1) writes \(z=\sum_j\xi_j\otimes m_j+z_0\) with \(\xi_j\) cocycles of \(\theta^{-1}V=W_\pi\) whose classes are finite. Write \(m_j=\alpha_jm_0\). The image of \(\xi_j\otimes m_j\) in \(M\otimes L'/O\) is the image of the rational cocycle \(\alpha_j\varpi^{-a}\xi_j\), whose class is finite. Coboundaries map to coboundaries. Hence \(\bar\kappa\) is finite. If \(2\le c\le t\), (T2) makes \(z\), and hence \(\bar\kappa\), a coboundary on \(G_K\). If \(c=1\), (T3) makes the restriction of \(\bar\kappa\) to inertia a coboundary.
\leanref{GlobalTheory.\allowbreak{}UseLoc.\allowbreak{}Cases.\allowbreak{}partW}{Development.\allowbreak{}AbelianVarietyTheory.\allowbreak{}Proofs.\allowbreak{}InertiaEigenvalues.\allowbreak{}UnramifiedEnlargementCutCocyclesBlochKato}{25--57}
\leanref{GlobalTheory.\allowbreak{}UseLoc.\allowbreak{}Cases.\allowbreak{}partC}{Development.\allowbreak{}AbelianVarietyTheory.\allowbreak{}Proofs.\allowbreak{}InertiaEigenvalues.\allowbreak{}UnramifiedEnlargementCutCocyclesBlochKato}{58--91}
\leanref{GlobalTheory.\allowbreak{}UseLoc.\allowbreak{}Cases.\allowbreak{}partU}{Development.\allowbreak{}AbelianVarietyTheory.\allowbreak{}Proofs.\allowbreak{}InertiaEigenvalues.\allowbreak{}UnramifiedEnlargementCutCocyclesBlochKato}{93--125}
\leanref{GlobalTheory.\allowbreak{}T.\allowbreak{}isBlochKato\_\allowbreak{}read\_\allowbreak{}of\_\allowbreak{}tensor\_\allowbreak{}selmer\_\allowbreak{}add\_\allowbreak{}coboundary}{Development.\allowbreak{}AbelianVarietyTheory.\allowbreak{}Proofs.\allowbreak{}InertiaEigenvalues.\allowbreak{}IsBlochKatoReadOfTensorSelmerAddCoboundaryStatement}{10--16}
\end{proof}

\subsection*{Proof of \ref{B8.42bqa}}\label{prf:B8.42bqa}\phantomsection
\begin{proof}[Proof (outline)]
Since \(K_0\subset K\subset K_0(\mu_\ell)\) and \([K_0(\mu_\ell):K_0]=\ell-1\), tower multiplicativity gives \(e\mid\ell-1\). Forming the crystalline kernel \(0\to\ker F\to T_1\to T_2\to0\) and using additivity of nonnegative jump multiplicities shows that every jump of \(T_2\) is a jump of \(T_1\); taking \(\eta_1=\min(\eta,h)\), \(h'=h+S+1\) and \(\beta'_\tau=\min(\beta_\tau,h+1)\) gives (a): \(\eta_1<\beta'_\tau\) since \(\beta_\tau>\eta\ge\eta_1\) and \(h+1>\eta_1\); \(\beta'_\tau+S\le h'\); all jumps lie in \([-h,h]\), and if \(\beta_\tau>h+1\) the shifted window contains no jump, so the jumps at \(\tau\) lie in \([-\eta,\eta]\cap[-h,h]=[-\eta_1,\eta_1]\); and \([-\eta_1,\eta_1]\subset[-\eta,\eta]\) keeps the bound \(m\). Part (b) follows from tame descent of the cyclotomic prism over \(K'/K\), where \(|\Delta|=s\) is prime to \(\ell\) and \(E_c\) are the conjugate Eisenstein polynomials of a uniformizer of \(K\) over \(K_0\).

For (c), the Frobenius pole bound and the identity \(\varphi(\pi')=\pi'E_\pi^su\), with \(u\in A_*^\times\), give \(\varphi(U_M(T))\subset E_\pi^{sM-h'}U_M(T)\). Since \(sM\ge4h'\), Frobenius is contracting on \(U_M(T)\) for the \((\ell,Z)\)-adic topology. The convergent geometric series gives \((\varphi-1)^{-1}=-\sum_{j\ge0}\varphi^j\) there. For a numerator vector \(w\), set \(z=(\varphi-1)^{-1}(\gamma-1)w\). The period-primitive argument of \eqref{eq:B8-finite-pairs} (extended across \(\Delta\) via \(|\Delta|=s\in O^\times\)) shows that \(\operatorname{cl}_{K,T}(w,z)\in H_f^1(K,T)\), represented by a finite cocycle. If \(w\) changes by \(u_0\in\pi'^MN_c(T)^\Delta\), put \(v=(\varphi-1)^{-1}u_0\). The pair changes by \(((\varphi-1)v,(\gamma-1)v)\), so its class is unchanged. This proves that \(\operatorname{cl}_c^H\) is well-defined; linearity and naturality follow from the same formula.
\leanref{GlobalTheory.\allowbreak{}LatExact.\allowbreak{}W\_\allowbreak{}exact}{Development.\allowbreak{}AbelianVarietyTheory.\allowbreak{}Proofs.\allowbreak{}InertiaEigenvalues.\allowbreak{}EtaleFunctorExactOnLattices}{12--16}
\leanref{Formalism.\allowbreak{}MainNotions.\allowbreak{}N.\allowbreak{}Ln34.\allowbreak{}frobPiNCondition}{Formalism.\allowbreak{}SharpRestriction.\allowbreak{}LevelLetterCocycles.\allowbreak{}DvrLatticeRank.\allowbreak{}FrobPiNPf}{345--355}
Realize any \(\kappa\in Z_f(T)\) by a crystalline extension \(0\to T\to T'\to O\to0\), whose jumps still lie in \([-h',h']\) since \(O\) has single jump \(0\), and choose a \(\Delta\)-invariant generator of the unit Wach module.
\leanref{Formalism.\allowbreak{}MainNotions.\allowbreak{}N.\allowbreak{}Ln34.\allowbreak{}ExtPf.\allowbreak{}ext\_\allowbreak{}pf}{Formalism.\allowbreak{}SharpRestriction.\allowbreak{}LevelLetterCocycles.\allowbreak{}DvrLatticeRank.\allowbreak{}ExtPfMultVEqZero}{58--75}
\leanref{Formalism.\allowbreak{}MainNotions.\allowbreak{}N.\allowbreak{}Ln34.\allowbreak{}OneGenPf.\allowbreak{}onegen\_\allowbreak{}pf}{Formalism.\allowbreak{}SharpRestriction.\allowbreak{}LevelLetterCocycles.\allowbreak{}DvrLatticeRank.\allowbreak{}OneGenPf}{10--14}
Averaging the cokernel annihilator over \(\Delta\) (using \(s\in O^\times\)) gives \(\pi'^{a'}\)-annihilation on \(\Delta\)-invariant components with \(a'=\lceil 2eh'/(\ell-1)\rceil\), and Lagrange interpolation over \(\{\chi_\ell(\gamma)^{sj}:0\le j<2a'\}\) clears Vandermonde denominators dividing \(\ell^{\rho_A}\) to lift \(\Delta\)-invariant target vectors after multiplication by \(\ell^{\rho_A}\).
\leanref{Formalism.\allowbreak{}MainNotions.\allowbreak{}N.\allowbreak{}Ln34.\allowbreak{}w\_\allowbreak{}pf}{Formalism.\allowbreak{}SharpRestriction.\allowbreak{}LevelLetterCocycles.\allowbreak{}DvrLatticeRank.\allowbreak{}WLocal}{312--385}
Combining this \(\ell^{\rho_A}\)-lift with saturated Frobenius-span decomposition, the pole bound \(E(\pi')^{h'}\), and \(\gamma\)-graded shift annihilation by \(\ell^{\sigma_A}\) up to \(\pi'^M\) decomposes \(\ell^{\rho_A+\sigma_A}\kappa\) into \(\sum_c\mathcal J_c(T)+B(T)\).
\leanref{Formalism.\allowbreak{}MainNotions.\allowbreak{}N.\allowbreak{}Ln34.\allowbreak{}kerlam2\_\allowbreak{}pf}{Formalism.\allowbreak{}SharpRestriction.\allowbreak{}LevelLetterCocycles.\allowbreak{}DvrLatticeRank.\allowbreak{}FrobeniusSpanDecomposition}{146--162}
\leanref{Formalism.\allowbreak{}MainNotions.\allowbreak{}N.\allowbreak{}Ln34.\allowbreak{}ninlam\_\allowbreak{}pf}{Formalism.\allowbreak{}SharpRestriction.\allowbreak{}LevelLetterCocycles.\allowbreak{}DvrLatticeRank.\allowbreak{}NinLamPf}{51--61}
\leanref{Formalism.\allowbreak{}MainNotions.\allowbreak{}N.\allowbreak{}Ln34.\allowbreak{}gradNCondition}{Formalism.\allowbreak{}SharpRestriction.\allowbreak{}LevelLetterCocycles.\allowbreak{}DvrLatticeRank.\allowbreak{}GradNPf}{9--28}
\leanref{Formalism.\allowbreak{}MainNotions.\allowbreak{}N.\allowbreak{}Ln34.\allowbreak{}gamPiCondition}{Formalism.\allowbreak{}SharpRestriction.\allowbreak{}LevelLetterCocycles.\allowbreak{}DvrLatticeRank.\allowbreak{}GamPiPf}{80--97}

The complete proof is contained in \Lean{Formalism.\allowbreak{}MainNotions.\allowbreak{}N.\allowbreak{}Ln34.\allowbreak{}local\_\allowbreak{}theorem\_\allowbreak{}of\_\allowbreak{}leaf\_\allowbreak{}lemmas} and \Lean{GlobalTheory.\allowbreak{}theoremA\_\allowbreak{}closing}.
\end{proof}

\subsection*{Proof of \ref{B8.42bqd}}\label{prf:B8.42bqd}\phantomsection
\begin{proof}
Since \(v_\ell(\zeta-1)>0\), the substitution defining \(\operatorname{ev}_{c,\zeta}\) converges coefficientwise. The identity \(Q(Z)=((1+Z)^\ell-1)/Z\) gives
\[
 \operatorname{ev}_{c,\zeta}(Q)(t)
   =\frac{\exp(\ell t)-1}{\zeta\exp(t)-1},
 \qquad
 [t]\operatorname{ev}_{c,\zeta}(Q)=\frac{\ell}{\zeta-1}\ne0.
\]
The comparison expansions have support bounded below. Hence the formula for \(v_{c,\zeta,T}(z)\) is defined whenever a regularizing exponent exists, and choosing the least exponent specifies its value uniquely.

The map defining \(R_c(S_0)\) is \(O\)-bilinear before passage to the tensor product. The other coordinate modules are images or inverse images under \(O\)-linear maps, and the displayed intersections are submodules of the respective numerators. Thus all the quotients are defined. The map \(f_P\) preserves Wach lattices and commutes with \(\Delta\), Frobenius, and scalar multiplication. It therefore descends to \(f_0\), and tensoring with the identity gives \(f_L\).
\end{proof}

\subsection*{Proof of \ref{B8.42bqe}}\label{prf:B8.42bqe}\phantomsection
\begin{proof}[Proof (outline)]
At an evaluation pair, write \(\operatorname{ev}_{c,\zeta}(E_\pi)=tU(t)\) with \(U(0)\ne0\). The defining filtration inequalities for numerator coefficients are preserved by multiplication by \(A^\Delta\), by the substitution \(t\mapsto\chi_\ell(\gamma)t\), and by division by a nonzero scalar in \(L'\). They give the inclusions and saturation in (a). Division by the distinguished polynomial defining \(E_\pi\) gives finitely many \(O\)-generators for \(\mathcal Q_c(T)\). Saturation of \(\pi'\mathcal M_c(T)\) in \(\Lambda_c^L(T)\) gives torsion freeness, and finite torsion-free modules over \(O\) are free.

For (b), put \(u=\chi_\ell(\gamma)\) and \(g(t)=f(ut)-f(t)\). The numerator condition gives
\[
 [t^j]\bigl(t^{h'}U(t)^{h'}g(t)\bigr)\in\mathscr F_\tau^{h'-j}(T).
\]
Thus \([t^j](U(t)^{h'}g(t))\in\mathscr F_\tau^{-j}(T)\). Induction from the lower bound of the Laurent series, using \(U(0)\ne0\) and the decreasing filtration, gives \([t^i]g(t)\in\mathscr F_\tau^{-i}(T)\). This coefficient is \((u^i-1)[t^i]f(t)\), and \(u^i-1\ne0\) for \(i\ne0\).

The expansion identity in Definition~\ref{B8.42bqb}(b), with the jumps contained in \([-h',h']\), identifies the model coefficients with Wach vectors. Multiplication by \(\pi'^j\) gives the two implications in (c); the remainders lie in \(J_c^L(T)\) and \(J_c^P(T)\), which are contained in the respective \(M\)-th powers. Commutation with \(\gamma\) then gives the invariant conditions. Quotienting by the first powers defines \(\alpha_{1,T}\) and \(\beta_{1,T}\); their compositions and naturality follow from those of the frame.

The complete proof is contained in \Lean{Formalism.\allowbreak{}MainNotions.\allowbreak{}N.\allowbreak{}Ln34.\allowbreak{}modell\_\allowbreak{}pf}, \Lean{Formalism.\allowbreak{}MainNotions.\allowbreak{}N.\allowbreak{}Ln34.\allowbreak{}modelflink\_\allowbreak{}pf}, and \Lean{Formalism.\allowbreak{}MainNotions.\allowbreak{}N.\allowbreak{}Ln34.\allowbreak{}modelqlink\_\allowbreak{}pf}.
\end{proof}

\subsection*{Proof of \ref{B8.42bqca}}\label{prf:B8.42bqca}\phantomsection
\begin{proof}
Write the \(c\)-component of \(a\) as \(F(\pi')\), where \(F(X)=\sum_{j\ge0}a_jX^j\in O[[X]]\). Since \(r^{2k}\in\ell^kO\), the following sums converge in \(O\):
\[
 F(r)=\sum_{j\ge0}a_jr^j,\qquad
 G_i=\sum_{j>i}a_jr^{j-i-1}.
\]
The identity \(F(X)=F(r)+(X-r)\sum_{i\ge0}G_iX^i\) follows coefficient by coefficient. Repeat this division \(d\) times and substitute \(X=\pi'\). This constructs the stated coefficients and remainder.

At a root \(r\), the identity
\[
 \Phi_\ell(\zeta\exp(t))
   =\frac{\exp(\ell t)-1}{\zeta\exp(t)-1}
\]
shows that the evaluated cyclotomic polynomial has a simple zero at \(t=0\). Its unit factor in \(E_c(\pi')\) has nonzero constant coefficient. Factoring \(E_c(X)\) by \(X-r\) therefore proves that the evaluation of \(\pi'-r\) has the form \(tU(t)\) with \(U(0)\ne0\). Comparison of the first nonzero coefficient proves the last assertion.
\end{proof}

\subsection*{Proof of \ref{B8.42bqcb}}\label{prf:B8.42bqcb}\phantomsection
\begin{proof}[Proof (outline)]
Both evaluation values are roots of the degree-two polynomial \(E_c\). Divide successively by \(X-r_U\) and \(X-r_D\) to obtain the displayed factorization. Put \(Q=\Phi_\ell(1+Z)\). Write \(\pi'=wZ^s\) and \(Q=w'E_\pi\), with \(w,w'\) units. Let \(w'_0\in O^\times\) be the constant coefficient of \(w'\). Choose \(t_U,t_D\in O\) with \(r_U^2=\ell t_U\), \(r_D^2=\ell t_D\). The constant-coefficient identity is \(\ell=w'_0a_2r_Ur_D\). Squaring and cancelling \(\ell^2\) gives
\[
 a_2\bigl(a_2(w'_0)^2t_Ut_D\bigr)=1,
\]
so \(a_2\) is a unit.

Choose the root of unity \(z\in O\) associated with \(r_U\), so \(z^\ell=1\), \(z\ne1\), and \(r_U=\pi'(z-1)\). Let \(u_0,u_1\in O\) be the constant and linear coefficients of the expansion of \(w\) at this point, and let \(u'_0\in O\) be the constant coefficient for \(w'\). Comparing linear coefficients in \(Q=w'E_\pi\) gives
\[
 \ell=(z-1)^s\delta_0\Omega,\qquad
 \Omega=-u'_0a_2\bigl(u_1(z-1)+u_0sz\bigr)\in O.
\]
The cyclotomic identity gives \((z-1)^{2s}=\ell t_0\) for some \(t_0\in O\). Squaring the preceding equality and cancelling the nonzero element \(\ell\) yields
\(\ell=\delta_0^2t_0\Omega^2\), which proves the divisibility and valuation bound.

For the last assertion, the numerical hypothesis gives \(y\ge v\) and \(2(y-v)\ge d\). The right side of the displayed formula therefore belongs to \(O\). Substitution of \(o_0=x\ell^vc'\delta_0^d\) and \(\delta_0^2r_0=\ell\) proves the equality. The complete proofs are contained in \Lean{Formalism.\allowbreak{}MainNotions.\allowbreak{}N.\allowbreak{}Ln34.\allowbreak{}LevelExpansions.\allowbreak{}fact2} and \Lean{Formalism.\allowbreak{}MainNotions.\allowbreak{}N.\allowbreak{}Ln34.\allowbreak{}LevelExpansions.\allowbreak{}caseC}.
\end{proof}

\subsection*{Proof of \ref{B8.42bqc}}\label{prf:B8.42bqc}\phantomsection
\begin{proof}[Proof (outline)]
The backward Frobenius construction on a basis of \(L_c(T)\) gives maps to the quotients by \(J_c^L(T)\) and \(J_c^P(T)\). Its two compositions are multiplication by \(\ell^{2\sigma_A}\). The identity \(\varphi(\pi')^M=\pi'^ME_\pi^{4h'}u\), with \(u\in A^\Delta\), places the remainders in these denominator modules. The construction commutes with scalar multiplication, \(\gamma\), and maps of lattices, and its expansion through degree \(k+4h'-1\) is the formula in Definition~\ref{B8.42bqb}(b). Thus it supplies the frame in (a).
\leanref{Formalism.\allowbreak{}MainNotions.\allowbreak{}N.\allowbreak{}Ln34.\allowbreak{}frame\_\allowbreak{}pf}{Formalism.\allowbreak{}SharpRestriction.\allowbreak{}LevelLetterCocycles.\allowbreak{}DvrLatticeRank.\allowbreak{}FramePf}{64--69}

On \(\pi'^jN_c(T)^\Delta/\pi'^{j+1}N_c(T)^\Delta\), \(\gamma-1\) acts as \(\chi_\ell(\gamma)^{sj}-1\). The graded modules are torsion-free over \(O\). Consequently an invariant element of \(\pi'N_c(T)^\Delta/\pi'^MN_c(T)^\Delta\) is zero. The cokernel of \(\gamma-1\) on this module is killed by
\[
 \prod_{j=1}^{M-1}(\chi_\ell(\gamma)^{sj}-1),
 \qquad
 v_\ell\!\left(\prod_{j=1}^{M-1}(\chi_\ell(\gamma)^{sj}-1)\right)
 =\sum_{j=1}^{M-1}(1+v_\ell(j))=\sigma_A.
\]
The empty product applies when \(M=1\). Given a representative \(w\) of a class in \(\mathcal P_c^{(1)}(T)\), solve for \(k\in\pi'N_c(T)^\Delta\) so that
\[
 (\gamma-1)k\equiv\ell^{\sigma_A}(\gamma-1)w
                  \pmod{\pi'^MN_c(T)^\Delta}.
\]
Then \(\ell^{\sigma_A}w-k\) is an invariant lift. The vanishing of the invariant kernel proves uniqueness after quotienting, and hence linearity and naturality. This proves (a).
\leanref{Formalism.\allowbreak{}MainNotions.\allowbreak{}N.\allowbreak{}Ln34.\allowbreak{}invLiftCondition}{Formalism.\allowbreak{}SharpRestriction.\allowbreak{}LevelLetterCocycles.\allowbreak{}DvrLatticeRank.\allowbreak{}InvLiftPf}{14--84}

For (b), the unprojected values \(v_{c,\zeta,T_2}(z)\) are independent of the evaluation pair above \(c\). They define an \(O\)-linear injection on \(L_c(T_2)\), which remains injective after extending scalars to \(L'\). Projection to the \(\tau\)-component identifies their image with \(D_{\rm dR}(T_2[1/\ell])_\tau\). Put \(W_c=L'\otimes_OL_c(T_2)\), and let \(U_\tau^a\subset W_c\) be the inverse image of \(\mathscr F_\tau^a(T_2)\) under evaluation. The jump conditions give
\[
 \dim_{L'}(W_c/U_\tau^{\beta'_\tau})\le m
   \quad(\varsigma_\tau=1),\qquad
 \dim_{L'}U_\tau^{-\eta_1}\le m
   \quad(\varsigma_\tau=-1).
\]
There are at most \(e\) embeddings above \(c\). Decompose the residue lattice using the intersection of the upward subspaces and the sum of the downward central subspaces. This gives the generators \(\bar z_j\) and a residual submodule \(L_r\) of rank at most \(em\). Since \(L_c(T_2)\) is finite free over the discrete valuation ring \(O\), so is \(L_r\).

The coefficient construction selects a finite subset \(\Omega_c^0\subset\Omega_c\). Put \(d_c^0=\prod_{i\in\Omega_c^0}(1-\chi_\ell(\gamma)^i)\). Apply \(\prod_{i\in\Omega_c^0}(\gamma_L-\chi_\ell(\gamma)^i)\) to an invariant coordinate representative. On Laurent coefficients this removes the selected exponents; on \(\mathcal H_c(T_2)\) it is multiplication by \(d_c^0\). In the residue directions \(\bar z_j\), the resulting coefficients have upward order at least \(-\beta'_\tau\), and at a downward pair they are a constant plus a series of order at least \(\beta'_\tau+S\). Their \(\gamma-1\) differences have numerators divisible by \(\pi'\). These are exactly the conditions in Definition~\ref{B8.42bqb}(d), and the remaining coordinate lies in \(\mathcal I_c^L(T_2;L_r)\).

Since \(\Omega_c^0\subset\Omega_c\), one has \(d_c=d_c^0t_c\), where
\[
 t_c=\prod_{i\in\Omega_c\setminus\Omega_c^0}(1-\chi_\ell(\gamma)^i)\in O.
\]
Apply this construction to \(t_cw\). The invariant module is an \(O\)-module, and the submodule \(L_r\) and the vectors \(z_j\) were fixed before choosing \(w\). It therefore gives admissible coefficients and the equality in (b) with \(d_c\).
\leanref{Formalism.\allowbreak{}MainNotions.\allowbreak{}N.\allowbreak{}Ln34.\allowbreak{}u1link\_\allowbreak{}pf}{Formalism.\allowbreak{}SharpRestriction.\allowbreak{}LevelLetterCocycles.\allowbreak{}DvrLatticeRank.\allowbreak{}U1linkPf}{16--68}
\leanref{Formalism.\allowbreak{}MainNotions.\allowbreak{}N.\allowbreak{}Ln34.\allowbreak{}FirstPropositionFrame}{Formalism.\allowbreak{}SharpRestriction.\allowbreak{}LevelLetterCocycles.\allowbreak{}DownSlotProposition}{216--232}
\leanref{Formalism.\allowbreak{}MainNotions.\allowbreak{}N.\allowbreak{}Ln34.\allowbreak{}LevelExpansions.\allowbreak{}valH\_\allowbreak{}of}{Formalism.\allowbreak{}SharpRestriction.\allowbreak{}LevelLetterCocycles.\allowbreak{}DvrLatticeRank.\allowbreak{}PiNSOfPi}{179--236}
\leanref{Formalism.\allowbreak{}MainNotions.\allowbreak{}N.\allowbreak{}Ln34.\allowbreak{}kapPiPf\_\allowbreak{}main}{Formalism.\allowbreak{}SharpRestriction.\allowbreak{}LevelLetterCocycles.\allowbreak{}DvrLatticeRank.\allowbreak{}KapPiPf}{217--284}

For (c), inclusion induces an injection
\[
 \mathcal R_c(T_2;S_0)\longrightarrow\mathcal Q_c(T_2).
\]
Lemma~\ref{B8.42bqe}(a) therefore gives finite generation and torsion freeness. The coefficient-ring filtration used for this injection has scalar \(\gamma\)-action on each graded quotient. There are at most \(e+1\) graded quotients with scalar \(1\), and each contributes at most \(\operatorname{rank}_OS_0\) to the rank of the invariant quotient. Thus its rank is at most \((e+1)\operatorname{rank}_OS_0\). The injection into \(\mathcal H_c(T_2)\) follows from the same inclusion and the same intersection with \(\pi'\mathcal M_c(T_2)\).
\leanref{Formalism.\allowbreak{}MainNotions.\allowbreak{}N.\allowbreak{}Ln34.\allowbreak{}RankBounds.\allowbreak{}u2rank\_\allowbreak{}local}{Formalism.\allowbreak{}SharpRestriction.\allowbreak{}LevelLetterCocycles.\allowbreak{}DvrLatticeRank.\allowbreak{}RankBounds}{295--326}
\leanref{Formalism.\allowbreak{}MainNotions.\allowbreak{}N.\allowbreak{}Ln34.\allowbreak{}RankLink.\allowbreak{}rank\_\allowbreak{}link\_\allowbreak{}of\_\allowbreak{}model}{Formalism.\allowbreak{}SharpRestriction.\allowbreak{}LevelLetterCocycles.\allowbreak{}DvrLatticeRank.\allowbreak{}RankLinkProof}{61--80}

For (d), every \(i\in\Omega_c\) is nonzero and \(|i|\le h'\). Counting multiples of \(\ell^a\) in an interval of \(S\) consecutive integers gives
\[
 \sum_{i\in\Omega_c}(1+v_\ell(|i|))
 \le e\bigl(S+v_\ell(S!)+\lfloor\log_\ell h'\rfloor\bigr)=\kappa_A.
\]
Since \(v_\ell(1-\chi_\ell(\gamma)^i)=1+v_\ell(|i|)\), the scalar \(d_c\) is a unit times \(\ell^v\) for an integer \(v\le\kappa_A\). Multiplication by the inverse unit and by \(\ell^{\kappa_A-v}\) gives \(a_c\).

For an elementary slot tensor \(f\otimes\bar z\), write \(E_\pi^{h'}(\gamma f-f)=\pi'g\). At an upward pair, the expansion of \(g\) has order at least \(h'-\beta'_\tau\), and \(v_{c,\zeta,T_2}(z)\in\mathscr F_\tau^{\beta'_\tau}(T_2)\). At a downward pair the constant term in the Laurent expansion of \(f\) cancels under \(\gamma-1\), so the expansion of \(g\) has order at least \(h'+\beta'_\tau+S\). Here \(v_{c,\zeta,T_2}(z)\in\mathscr F_\tau^{-\beta'_\tau-S}(T_2)\) by the lower jump bound. The factor \(\operatorname{ev}_{c,\zeta}(\pi')\) has nonzero constant coefficient. These inequalities put \(E_\pi^{-h'}g\otimes\bar z\) in \(\mathcal M_c(T_2)\), so the slot tensor belongs to \(\mathcal I_c^L(T_2;L_c(T_2))\).

The Wach lifting estimate gives \(z_1\in N_{c-1}(T_1)^\Delta\) with \(F_Pz_1=\ell^{\rho_A}z\). When all signs are downward, the preceding argument puts \(f\otimes\bar z_1\) in the invariant coordinate module, and naturality of the frame gives the required lift. When all signs are upward, \(f\otimes\bar z\in\mathcal M_c(T_2)\); its image under \(\alpha_{T_2}\) has a Wach representative. Lift that representative by the same Wach estimate at component \(c\). The action of \(\gamma\) is trivial on the Wach lattice modulo \(\pi'\), so the resulting vector represents a class in \(\mathcal P_c^{(1)}(T_1)\). Multiplying either lift by \(\ell^{y_1-\rho_A}\) gives the asserted exponent.

We first describe the passage from a source residue vector to the required slot representative. Suppose \(a\ge0\), \(z_*\in N_{c-1}(T_1)^\Delta\), \(F_Pz_*-\ell^az\in\pi'N_{c-1}(T_2)^\Delta\), and \(f\otimes\bar z_*\) belongs to the invariant coordinate module. Choose \(y_*\in\Lambda_c^P(T_1)\) satisfying
\[
 q_P(y_*)=\alpha_{T_1}(f\otimes\bar z_*).
\]
The invariant clause of Lemma~\ref{B8.42bqe}(c) gives \((\gamma-1)y_*\in\pi'N_c(T_1)^\Delta\). Naturality of the frame gives \(q_P(F_Py_*-\ell^ay)=0\). The height bound \(E_\pi^{h'}P_c^\varphi(T_2)\subset N_c(T_2)^\Delta\) implies
\[
 J_c^P(T_2)=\pi'^ME_\pi^{3h'}P_c^\varphi(T_2)
       \subset\pi'^MN_c(T_2)^\Delta\subset\pi'N_c(T_2)^\Delta.
\]
Hence \(F_Py_*-\ell^ay\in\pi'N_c(T_2)^\Delta\), which is the required slot lift.
\leanref{Formalism.\allowbreak{}MainNotions.\allowbreak{}N.\allowbreak{}Ln34.\allowbreak{}IsFrame}{Formalism.\allowbreak{}SharpRestriction.\allowbreak{}LevelLetterCocycles.\allowbreak{}DownSlotProposition}{122--152}
\leanref{Formalism.\allowbreak{}MainNotions.\allowbreak{}N.\allowbreak{}Ln34.\allowbreak{}ModelF}{Formalism.\allowbreak{}SharpRestriction.\allowbreak{}LevelLetterCocycles.\allowbreak{}DownSlotProposition}{198--212}

For mixed signs with \(e=2\), let \(r_{\zeta,\tau}=[t^0]\operatorname{ev}_{c,\zeta}(\pi')\). The relation \(\pi'=wZ^s\), with \(w\) a unit and \(2s=\ell-1\), gives \(r_{\zeta,\tau}=w_c(\zeta-1)(\zeta-1)^s\in O\). The cyclotomic identity \((\zeta-1)^{\ell-1}\in\ell O\) therefore gives \(\ell\mid r_{\zeta,\tau}^2\). First suppose that every evaluation has the same value as some downward evaluation. Apply Lemma~\ref{B8.42bqca} to the numerator of \(E_\pi^{h'}(\gamma f-f)/\pi'\). Its downward vanishing order transfers to every evaluation with that value. The lower jump bound for \(T_1\) then puts \(f\otimes\bar z_1\) in the invariant coordinate module, for the previously chosen lift \(F_Pz_1=\ell^{\rho_A}z\). This gives a lift at exponent \(\rho_A\).

Otherwise choose an upward evaluation \((\zeta_U,\tau_U)\) whose value differs from every downward value, and choose a downward evaluation \((\zeta_D,\tau_D)\). Put
\[
 r_U=r_{\zeta_U,\tau_U},\quad r_D=r_{\zeta_D,\tau_D},\quad
 b=\beta'_{\tau_U},\quad B=\beta'_{\tau_D}+S,\quad
 \delta_0=r_D-r_U.
\]
The normalized jump inequalities give \(1\le B\le h'\) and \(b\le h'\). There are two embeddings above \(c\). Every upward embedding is \(\tau_U\), every downward embedding is \(\tau_D\), and every evaluation value is \(r_U\) or \(r_D\). Lemma~\ref{B8.42bqcb} gives
\(E_c(X)=a_2(X-r_U)(X-r_D)\), with \(a_2\in O^\times\) and \(\delta_0^2\mid\ell\).

Suppose first that \(B+2\eta_1<b\). Set
\[
 a_0=\epsilon_c(\pi'-r_U)^{h'-b}(\pi'-r_D)^{h'+B}.
\]
The upward and downward orders put \((a_0/E_\pi^{h'})\otimes\bar z\) in \(\mathcal M_c(T_2)\). Apply \(\alpha_{T_2}\), choose its Wach representative, lift that representative through \(F_P\) after multiplication by \(\ell^{\rho_A}\), and apply \(\beta_{T_1}\). Write the resulting model representative as
\[
 \nu=\sum_j(E_\pi^{-h'}a_j)\otimes\bar z_j,
 \qquad a_j\in A^\Delta,\quad z_j\in N_{c-1}(T_1)^\Delta.
\]
For an evaluation pair, write \(v_j=v_{c,\zeta,T_1}(z_j)\) and \(v'_j=v_{c,\zeta,T_2}(F_Pz_j)\). The model filtration conditions give
\(\sum_j[t^i]\operatorname{ev}(a_j)v_j\in\mathscr F_\tau^{h'-i}(T_1)\). Naturality and the composition identity for the frame give, for \(0\le i\le h'\),
\[
 \sum_j[t^i]\operatorname{ev}(a_j)v'_j
   =\ell^{\rho_A+2\sigma_A}[t^i]\operatorname{ev}(a_0)
       v_{c,\zeta,T_2}(z).
\]
The denominator terms have order at least \(4h'\), so they contribute no coefficient in this range.
\leanref{Formalism.\allowbreak{}MainNotions.\allowbreak{}N.\allowbreak{}Ln34.\allowbreak{}LevelExpansions.\allowbreak{}compare}{Formalism.\allowbreak{}SharpRestriction.\allowbreak{}LevelLetterCocycles.\allowbreak{}DvrLatticeRank.\allowbreak{}TmulMkLMemInvLOfValAtMemFilC}{302--382}

The coefficient extraction in the next two divisions is triangular. At a root put \(Y=\operatorname{ev}(\pi'-r)=tU(t)\), where \(U(0)\ne0\). If a vector series has expansion \(\sum_j P_jY^j\), then
\[
 [t^k]\sum_jP_jY^j=U(0)^kP_k+\sum_{j<k}u_{kj}P_j
\]
for scalars \(u_{kj}\in L'\). Induction on \(k\), and the inclusions \(\mathscr F^{h'-j}\subset\mathscr F^{h'-k}\) for \(j<k\), recover the filtration condition on \(P_k\). The same induction recovers coefficient identities. Division by a power series with nonzero constant coefficient uses this recurrence as well.

Divide each \(a_j\) at \(r_D\), using Lemma~\ref{B8.42bqca}, through degree \(h'-\eta_1-1\):
\[
 a_j=\sum_{k<h'-\eta_1}p_{jk}(\pi'-r_D)^k
       +(\pi'-r_D)^{h'-\eta_1}a'_j
\]
on component \(c\). For each such \(k\), the downward filtration condition puts \(\sum_jp_{jk}v_j\) in \(\mathscr F_{\tau_D}^{\eta_1+1}(T_1)\). Its projection to the \(\tau_D\)-component is zero. Projection to this component is injective on the evaluation image, so the entire evaluation vector is zero. The kernel of evaluation on \(N_{c-1}(T_1)^\Delta\) is \(\pi'N_{c-1}(T_1)^\Delta\), hence \(\sum_jp_{jk}\bar z_j=0\). Naturality gives zero in the target residue lattice as well. The unprojected evaluation vector is independent of the evaluation pair above \(c\), so both zero identities also hold at the upward evaluation. Thus these polynomial terms disappear from both coefficient identities. At the upward evaluation, \(\pi'-r_D\) has nonzero constant coefficient, so division by its \((h'-\eta_1)\)-st power preserves the filtration conditions and the identities through degree \(h'\).

Put \(d_*=h'-b+B+\eta_1\), and divide each \(a'_j\) at \(r_U\) with remainder exponent \(d_*+1\):
\[
 a'_j=\sum_{k=0}^{d_*}q_{jk}(\pi'-r_U)^k
          +(\pi'-r_U)^{d_*+1}a''_j.
\]
Let \(o_j=q_{j,d_*}\in O\), and define
\[
 z'_1=\sum_jo_jz_j.
\]
After the first division, the target polynomial is
\[
 \ell^{\rho_A+2\sigma_A}
       X^{h'-b}(X+r_U-r_D)^{B+\eta_1}.
\]
Its coefficient of degree \(d_*\) is \(\ell^{\rho_A+2\sigma_A}\). Therefore
\[
 F_Pz'_1-\ell^{\rho_A+2\sigma_A}z
       \in\pi'N_{c-1}(T_2)^\Delta.
\]
The source coefficient belongs to
\(\mathscr F_{\tau_U}^{h'-d_*}(T_1)\subset
\mathscr F_{\tau_U}^{\eta_1+1}(T_1)\). There are no jumps between \(\eta_1+1\) and \(b\), so it belongs to \(\mathscr F_{\tau_U}^{b}(T_1)\). Independence of the unprojected evaluation vector gives this filtration condition at every evaluation with embedding \(\tau_U\), including those whose value of \(\pi'\) is \(r_D\). Apply the source-residue construction above with \(z_*=z'_1\) and \(a=\rho_A+2\sigma_A\). This constructs the invariant source slot at that exponent.
\leanref{Formalism.\allowbreak{}MainNotions.\allowbreak{}N.\allowbreak{}Ln34.\allowbreak{}LevelExpansions.\allowbreak{}ValH}{Formalism.\allowbreak{}SharpRestriction.\allowbreak{}LevelLetterCocycles.\allowbreak{}DvrLatticeRank.\allowbreak{}PiNSOfPi}{103--121}
\leanref{Formalism.\allowbreak{}MainNotions.\allowbreak{}N.\allowbreak{}Ln34.\allowbreak{}LevelExpansions.\allowbreak{}valH\_\allowbreak{}of}{Formalism.\allowbreak{}SharpRestriction.\allowbreak{}LevelLetterCocycles.\allowbreak{}DvrLatticeRank.\allowbreak{}PiNSOfPi}{179--236}
\leanref{Formalism.\allowbreak{}MainNotions.\allowbreak{}N.\allowbreak{}Ln34.\allowbreak{}LevelExpansions.\allowbreak{}caseB}{Formalism.\allowbreak{}SharpRestriction.\allowbreak{}LevelLetterCocycles.\allowbreak{}DvrLatticeRank.\allowbreak{}PolyExpand}{352--412}

Now suppose \(b\le B+2\eta_1\). Write \(f=a/E_\pi^{h'}\) on component \(c\), and let \(b_0\in L'\) be its downward constant. Division at \(r_U\) gives \(a=(\pi'-r_U)^{h'-b}A\) with \(A\) integral. At \(r_D\), the defining congruence for \(f\) gives
\[
 A\equiv b_0a_2^{h'}(\pi'-r_U)^b(\pi'-r_D)^{h'}
       \pmod{(\pi'-r_D)^{h'+B}}.
\]
Consequently, comparison of the coefficient of degree \(h'+j\), for \(0\le j<B\), gives
\[
 b_0a_2^{h'}\binom bj\delta_0^{b-j}\in O.
\]
If \(b<B\), take \(j=b\), so \(b_0\in O\) and set \(y_0=0\). If \(B\le b\), take \(j=B-1\). The factorial formula for the valuation of the binomial coefficient gives
\(v_\ell\binom b{B-1}\le\lfloor\log_\ell h'\rfloor\), and hence
\[
 b-B+1+2v_\ell\binom b{B-1}
   \le2\eta_1+1+2\log_\ell h'\le2y_c(\eta_1).
\]
Lemma~\ref{B8.42bqcb}, followed by division by the unit \(a_2^{h'}\), constructs \(o=\ell^{y_0}b_0\in O\) with \(y_0=y_c(\eta_1)\).

Define the corrected numerator
\[
 a_N=\ell^{y_0}a-o\epsilon_cE_\pi^{h'}.
\]
Its upward order is at least \(h'-b\), and its downward order is at least \(h'+B\). The transfer assertion in Lemma~\ref{B8.42bqca} supplies these conditions at every evaluation of the respective root. Thus
\[
 (\ell^{y_0}f-o\epsilon_c)\otimes\bar z\in\mathcal M_c(T_2).
\]
Lift this model term using its Wach representative and lift the constant slot \(o\epsilon_c\otimes\bar z\) using \(z_1\). Both require the factor \(\ell^{\rho_A}\). Their sum gives a lift at exponent \(\rho_A+y_0\). The two cases are bounded by \(\rho_A+\max(2\sigma_A,y_c(\eta_1))=y_1\).
\leanref{Formalism.\allowbreak{}MainNotions.\allowbreak{}N.\allowbreak{}Ln34.\allowbreak{}LevelExpansions.\allowbreak{}mix\_\allowbreak{}local}{Formalism.\allowbreak{}SharpRestriction.\allowbreak{}LevelLetterCocycles.\allowbreak{}DvrLatticeRank.\allowbreak{}MixLocal}{33--337}

For (e), let \(x_1,\ldots,x_d\in\mathcal I_c^L(T_2;L_r)\) represent an \(O\)-basis of \(\mathcal I_{\rm rest}\), so \(d\le e(e+1)m\). Put \(v_i=\alpha_{1,T_2}([x_i])\), choose a representative \(Y_i\) of \(\operatorname{lift}_c(v_i)\), and complete it to a Herr pair \((Y_i,Z_i)\) with \(Z_i\in U_M(T_2)\). Choose a finite cocycle \(\kappa_i\) representing its class and put \(\Sigma_c=\{\kappa_1,\ldots,\kappa_d\}\).

Let \(\kappa\) be a generator of \(\mathcal J_c(T_2)\), represented by a Herr pair \((w,z)\) as in Lemma~\ref{B8.42bqa}(c), and write \(\bar w\) for the reduction of \([w]\) to \(\mathcal P_c^{(1)}(T_2)\). Parts (b) and (c) give \(X\in\mathcal C_{\rm lift}(T_2,c)\) and \(b_i\in O\) with
\[
 \beta_{1,T_2}(d_c\bar w)=X+\sum_i b_i[x_i].
\]
Apply \(\alpha_{1,T_2}\), and then multiply by \(a_c\ell^{y_1}\). Part (d) supplies \(u_1\in\mathcal P_c^{(1)}(T_1)\) such that
\[
 \ell^{2\sigma_A+\kappa_A+y_1}\bar w
   =F_*u_1+\sum_i a_c\ell^{y_1}b_i v_i.
\]
Apply \(\operatorname{lift}_c\), and choose a representative \(W_1\) of \(\operatorname{lift}_c(u_1)\). By (a), with \(A=\sigma_A+2\sigma_A+\kappa_A+y_1\),
\[
 w'=\ell^Aw-F_PW_1-\sum_i a_c\ell^{y_1}b_iY_i\in U_M(T_2).
\]
Complete \(W_1\) to a Herr pair \((W_1,Z_1)\), where \(Z_1\in U_M(T_1)\), and choose a finite cocycle \(\kappa_1'\) representing it. The corresponding second-coordinate difference
\[
 z'=\ell^Az-F_PZ_1-\sum_i a_c\ell^{y_1}b_iZ_i
\]
also belongs to \(U_M(T_2)\). Set \(v=(\varphi-1)^{-1}w'\). Since \((w',z')\) is a Herr pair and \(\varphi-1\) is injective on \(U_M(T_2)\), one has \(z'=(\gamma-1)v\). This pair has zero cohomology class. Naturality of the Herr class map gives
\[
 \ell^A\kappa-F(\kappa_1')-\sum_i a_c\ell^{y_1}b_i\kappa_i\in B(T_2).
\]
Extend this relation \(O\)-linearly over the defining span of \(\mathcal J_c(T_2)\). This proves (e).

The complete proof is contained in \Lean{Formalism.\allowbreak{}MainNotions.\allowbreak{}N.\allowbreak{}Ln34.\allowbreak{}u1link\_\allowbreak{}pf} and \Lean{Formalism.\allowbreak{}MainNotions.\allowbreak{}N.\allowbreak{}Ln34.\allowbreak{}univlink\_\allowbreak{}pf}.
\end{proof}

\subsection*{Proof of \ref{B8.42bq1}}\label{prf:B8.42bq1}\phantomsection
\begin{proof}
Let \(\ell\) be odd and \(K_0\subset K\subset K_0(\mu_\ell)\) as in Definition~\ref{B8.42bq}; for \(\ell\in\{q,d\}\) and any place \(w\mid\ell\) of \(E\), Lemma~\ref{B8.42bd} shows that \(\ell\) is odd and \(K=E_w\) is of this form with \(e=[K:K_0]=2\) and \(s=(\ell-1)/2\).

Fix admissible data \((T_1,T_2,F,r,S,\eta,m,h,\varsigma_\tau,\beta_\tau)\) at \(K\). By Lemma~\ref{B8.42bqa}(a), \(es=\ell-1\) and there exist normalized integers \(0\le\eta_1\le\eta\), \(1\le h'\le h+S+1\), and \(\beta'_\tau\) (\(\tau:K\to L'\)) satisfying \(\eta_1<\beta'_\tau\) and \(\beta'_\tau+S\le h'\), such that the jumps of \(T_1\) and \(T_2\) lie in the corresponding shifted and central intervals inside \([-h',h']\) and \(T_2\) has at most \(m\) jumps in \([-\eta_1,\eta_1]\) at each \(\tau\). Choose \((\pi',E)\) as in Lemma~\ref{B8.42bqa}(b), and define \(M,\rho_A,\sigma_A\) and the submodules \(\mathcal J_c(T_2)\subset Z_f(T_2)\) (\(c\in\Z/n_K\Z\)) as in Lemma~\ref{B8.42bqa}(c), together with \(\kappa_A,y_c(\eta_1),y_1\) as in Proposition~\ref{B8.42bqc}.

For each \(c\in\Z/n_K\Z\), Proposition~\ref{B8.42bqc}(e) gives a subset \(\Sigma_c\subset Z_f(T_2)\) of cardinality at most \(e(e+1)m\) such that \(\ell^{\sigma_A+2\sigma_A+\kappa_A+y_1}\mathcal J_c(T_2)\subset I_F+O\Sigma_c\). Put \(\Sigma_0=\bigcup_{c\in\Z/n_K\Z}\Sigma_c\subset Z_f(T_2)\); then
\[
 |\Sigma_0|\le\sum_{c\in\Z/n_K\Z}|\Sigma_c|\le n_Ke(e+1)m .
\]
Set \(a=\ell^{\rho_A+\sigma_A}\) and \(b=\ell^{\sigma_A+2\sigma_A+\kappa_A+y_1}\). For any \(\kappa\in Z_f(T_2)\), Lemma~\ref{B8.42bqa}(c) writes \(a\kappa=\sum_{c\in\Z/n_K\Z}x_c+\beta_0\) with \(x_c\in\mathcal J_c(T_2)\) and \(\beta_0\in B(T_2)\subset I_F\). Multiplying by \(b\) gives \(bx_c\in I_F+O\Sigma_c\subset O\Sigma_0+I_F\) for every \(c\) and \(b\beta_0\in I_F\), whence \(ba\kappa\in O\Sigma_0+I_F\).
\leanref{Formalism.\allowbreak{}MainNotions.\allowbreak{}N.\allowbreak{}Ln34.\allowbreak{}annihilation\_\allowbreak{}of\_\allowbreak{}scaled\_\allowbreak{}spans}{Formalism.\allowbreak{}SharpRestriction.\allowbreak{}LevelLetterCocycles.\allowbreak{}MixedPositionHypothesis}{81--102}

Comparing exponents with the formula of Lemma~\ref{B8.42bc} gives
\[
 (\rho_A+\sigma_A)+\bigl(\sigma_A+2\sigma_A+\kappa_A+\rho_A+\max(2\sigma_A,y_c(\eta_1))\bigr)
 =2\rho_A+4\sigma_A+\kappa_A+\max\bigl(2\sigma_A,y_c(\eta_1)\bigr)
 =Y(\ell,e,s,h',S,\eta_1).
\]
Since \(\eta_1\le\eta\), we have \(y_c(\eta_1)\le y_c(\eta)\) and therefore \(Y(\ell,e,s,h',S,\eta_1)\le Y(\ell,e,s,h',S,\eta)\). Multiplying \(ba\kappa\in O\Sigma_0+I_F\) by \(\ell^{Y(\ell,e,s,h',S,\eta)-Y(\ell,e,s,h',S,\eta_1)}\) yields \(\ell^{Y(\ell,e,s,h',S,\eta)}Z_f(T_2)\subset O\Sigma_0+I_F\). This establishes statement (A) at \(K\).
\leanref{Formalism.\allowbreak{}MainNotions.\allowbreak{}N.\allowbreak{}Ln34.\allowbreak{}YA\_\allowbreak{}split}{Formalism.\allowbreak{}SharpRestriction.\allowbreak{}LevelLetterCocycles.\allowbreak{}MixedPositionHypothesis}{70--75}
\leanref{Formalism.\allowbreak{}MainNotions.\allowbreak{}N.\allowbreak{}Ln34.\allowbreak{}YA\_\allowbreak{}mono}{Formalism.\allowbreak{}SharpRestriction.\allowbreak{}LevelLetterCocycles.\allowbreak{}MixedPositionHypothesis}{76--80}
\leanref{Formalism.\allowbreak{}MainNotions.\allowbreak{}N.\allowbreak{}Ln34.\allowbreak{}killR\_\allowbreak{}of\_\allowbreak{}steps}{Formalism.\allowbreak{}SharpRestriction.\allowbreak{}LevelLetterCocycles.\allowbreak{}MixedPositionHypothesis}{105--171}
\end{proof}

\subsection*{Proof of \ref{B8.42br}}\label{prf:B8.42br}\phantomsection
\begin{proof}
\emph{Generators.} Put \(V_2=T_2[1/\ell]\). The quotient \(Z_f(T_2)/B(T_2)\) is the integral finite cohomology \(H^1_f(K,T_2)\), a submodule of \(H^1(K,T_2)\). The latter is finitely generated over \(O\), by Theorem~\ref{B8.41p} for the finite modules \(T_2/\varpi^nT_2\), the Mittag--Leffler property of their cohomology and Nakayama's lemma, where \(\varpi\) is a uniformizer of \(O\). Its torsion submodule is the image of \(H^0(K,V_2/T_2)\) under the connecting map, which kills the image of \(H^0(K,V_2)\). The maximal divisible submodule \(D\) of \(H^0(K,V_2/T_2)\) is a direct summand with finite complement \(Q'\). The image of \(H^0(K,V_2)\) lies in \(D\), and \(D\) modulo this image is finite and divisible, hence zero. Thus the torsion is isomorphic to \(Q'\). Since \(Q'[\varpi]\subset(V_2/T_2)[\varpi]\simeq(O/\varpi)^r\), the torsion needs at most \(r\) generators. By \eqref{eq:B8-BK-dim}, the rank of \(H^1_f(K,T_2)\) is \(\dim_{L'}H^0(K,V_2)+\dim_{L'}D_{\rm dR}(V_2)/\mathrm{Fil}^0\le r+n_Ker\). Hence \(Z_f(T_2)\) is generated modulo \(B(T_2)\) by at most \((n_Ke+1)r+r\le2(n_Ke+1)r\) elements.
\leanref{Formalism.\allowbreak{}MainNotions.\allowbreak{}N.\allowbreak{}Ln34.\allowbreak{}genA}{Formalism.\allowbreak{}SharpRestriction.\allowbreak{}LevelLetterCocycles.\allowbreak{}LocalFormOfThreeInputs}{13--20}

\emph{Lifting exponent.} Put \(K'=K_0(\mu_\ell)\). It contains \(K\) and equals \(\Q_\ell(\mu_{\ell N})\) with \(N=\ell^{n_K}-1\) prime to \(\ell\). Hence its depth in Lemma~\ref{B8.41c} is one, and \([K':K]=s\) is prime to \(\ell\). Apply the proof of the first assertion of Lemma~\ref{B8.41a} over \(K'\) to \(F\), whose source jumps lie in \([-h,h]\). It scales a finite target class \(c\) to \(c_\gamma P(1)\ell^{2w}c\), with \(w=4h+v_\ell((4h)!)\), and lifts the result to a finite source class. The valuation of the scalar is the left side of \eqref{eq:B8-wach-exponent} at depth one:
\[
 E'=2L(4h)+\sum_{j\in J^*}(1+v_\ell(j))+L(h),\qquad L(k)=k+v_\ell(k!).
\]
Here \(J^*=((J+[1,h])\setminus\{0\})\cup[1,h]\) for the union \(J\subset[-h,h]\) of source jump sets, so \(J^*\subset[-h,2h]\setminus\{0\}\). Restriction to \(K'\) followed by corestriction multiplies classes by \(s\), a unit of \(O\), and corestriction preserves finite classes. Hence \(\ell^{E'}Z_f(T_2)\subset I_F\).
\leanref{Formalism.\allowbreak{}MainNotions.\allowbreak{}N.\allowbreak{}Ln34.\allowbreak{}liftA}{Formalism.\allowbreak{}SharpRestriction.\allowbreak{}LevelLetterCocycles.\allowbreak{}LocalFormOfThreeInputs}{21--27}
\leanref{Formalism.\allowbreak{}MainNotions.\allowbreak{}ElementDenominator.\allowbreak{}CongruenceModules.\allowbreak{}expLift}{Formalism.\allowbreak{}SharpRestriction.\allowbreak{}LevelLetterCocycles.\allowbreak{}LocalTheoremOfGoodExtensions}{153--156}
\leanref{Formalism.\allowbreak{}WeaklyAdmissibleFiltrations.\allowbreak{}LatHom.\allowbreak{}LiftF}{Formalism.\allowbreak{}WeaklyAdmissibleFiltrations.\allowbreak{}Statement}{265--267}

\emph{Conclusion.} Let \(C=Z_f(T_2)/I_F\) and let \(C_0\) be the image of \(O\Sigma_0\), with \(h'\) and \(\Sigma_0\) from (A). Since \(B(T_2)\subset I_F\), the module \(C\) is generated by at most \(2(n_Ke+1)r\) elements. The submodule \(C_0\) is generated by at most \(n_Ke(e+1)m\) elements. Statement (A) gives \(\ell^YC\subset C_0\), and the lifting exponent gives \(\ell^{E'}C_0=0\). The integers satisfy \(es=\ell-1\), \(1\le h'\le h+S+1\) and \(J^*\subset[-h,2h]\setminus\{0\}\). Lemma~\ref{B8.42bc} with \(n=n_K\) gives the bound.
\leanref{Formalism.\allowbreak{}WeaklyAdmissibleFiltrations.\allowbreak{}LatHom.\allowbreak{}col}{Formalism.\allowbreak{}WeaklyAdmissibleFiltrations.\allowbreak{}Statement}{262--264}
\end{proof}

\subsection*{Proof of \ref{B8.42bs1a}}\label{prf:B8.42bs1a}\phantomsection
\begin{proof}
The Hecke algebra separates the points of \(X_i\), so its eigenspaces have direct sum. Commutation with the Hecke algebra makes every member of \(\mathcal A\) preserve each \(V_x\). If \(u=m+v\in U_{\mathcal S}\), with \(m\in M_i\) and \(v\in V_{\mathcal S^c}\), then
\(\alpha u=\alpha m\in M_i\cap V_{\mathcal S}=L_{\mathcal S}\)
for every \(\alpha\in\mathcal A\).

Order \(\mathcal S=\{\pi_1,\ldots,\pi_q\}\), set
\(F_j=\bigoplus_{a\le j}V_{\pi_a}\), and put
\(U_j=U_{\mathcal S}\cap F_j\), \(L_j=L_{\mathcal S}\cap F_j\).
Let \(p_j:F_j\to V_{\pi_j}\) be the projection with kernel \(F_{j-1}\). The preceding inclusion and preservation of \(F_j\) give
\[
 a_{\pi_j}p_j(U_j)\subset p_j(L_j)\subset p_j(U_j).
\]
The module \(p_j(U_j)\) is a free \(O_i\)-lattice of rank
\(\dim_{K_i}V_{\pi_j}\). Consequently
\[
 \lambda_i\bigl(p_j(U_j)/p_j(L_j)\bigr)
 \le\dim_{K_i}V_{\pi_j}\,\lambda_i(O_i/a_{\pi_j}O_i).
\]
The sequence
\[
 0\longrightarrow U_{j-1}/L_{j-1}
 \longrightarrow U_j/L_j
 \xrightarrow{p_j}p_j(U_j)/p_j(L_j)
 \longrightarrow0
\]
is exact. Add its length identities for \(1\le j\le q\). The empty set gives the zero quotient.
\end{proof}

\subsection*{Proof of \ref{B8.42bs1b}}\label{prf:B8.42bs1b}\phantomsection
\begin{proof}
For each \(x\in Y\setminus Z(i)\), choose \(f_x\in T_i\) with
\(f_x(\pi)\ne f_x(x)\). The product of \(f_x-f_x(x)\), with empty product \(1\), vanishes on \(Y\setminus Z(i)\) and is nonzero at \(\pi\). Thus
\(\mathfrak c_{Y,\pi}\ne0\). Since \(O_i\) is a discrete valuation ring, write
\(\mathfrak c_{Y,\pi}=cO_i\) with \(c\ne0\); the definition of the ideal supplies \(g\).

Both \(g\) and \(\beta\) commute with the Hecke algebra, so \(\alpha\) does also. It acts by \(cb\) on \(V_\pi\). For \(x\notin Y\), it vanishes because \(\beta\) vanishes. For \(x\in Y\setminus Z_{\rm good}(i)\), the assumption on \(Y\cap Z(i)\) gives \(x\notin Z(i)\), so \(g(x)=0\). Since \(\beta\) preserves \(V_x\), \(\alpha\) vanishes there too. Finally,
\[
 \lambda_i(O_i/cbO_i)
 =\lambda_i(O_i/cO_i)+\lambda_i(O_i/bO_i)
 \le v_{Y,\pi}+v_{g_k}(M_0).
\]
\end{proof}

\subsection*{Proof of \ref{B8.42bs1c}}\label{prf:B8.42bs1c}\phantomsection
\begin{proof}
Put \(A=H_f^1(K,L_{U'})\). Local Tate duality and the exact-annihilator assertion in Definition~\ref{B8.41h} give an injection
\[
 H^1(K,U/L_U)/H_f^1(K,U/L_U)
 \longrightarrow\operatorname{Hom}_O(A,L/O).
\]
Coboundaries lie in \(Z_f^1\), so \((S+Z_f^1)/Z_f^1\) injects into the \(\ell^e\)-torsion of the target. This torsion is
\(\operatorname{Hom}_O(A/\ell^eA,L/O)\).

The module \(A\) is finitely generated. Let \(\varpi\) be a uniformizer of \(O\). The exact sequence for multiplication by \(\varpi\) on \(L_{U'}\) gives a surjection
\[
 H^0(K,L_{U'}/\varpi L_{U'})
 \longrightarrow H^1(K,L_{U'})[\varpi].
\]
The source has residue-field dimension at most \(r_U\). Hence the torsion submodule of \(A\) has at most \(r_U\) generators. If \(d=\operatorname{rank}_O A\), the structure theorem over \(O\) gives
\[
 \operatorname{length}_O(A/\ell^eA)
 \le e(d+r_U)\operatorname{length}_O(O/\ell O).
\]
The dual of a finite-length \(O\)-module in \(L/O\) has the same length. Divide by \(\operatorname{length}_O(O/\ell O)\) to obtain the assertion.
\end{proof}

\subsection*{Proof of \ref{B8.42bs1d}}\label{prf:B8.42bs1d}\phantomsection
\begin{proof}[Proof (outline)]
The characters in \(\Sigma\) factor through a common finite quotient of \(I\). Character averaging on that quotient shows that their distinct eigenspaces have direct sum. Thus the displayed decomposition is equivalent to the sum condition on the inertia eigenspaces. The local cohomological inputs of \(\mathscr H_2\) give the lifting hypotheses for the primewise Wach system. The integral prismatic input supplies crystalline lattices with coefficients, and the Fontaine input supplies the comparison for their rational representations.

Put \(a=(\ell-1)!\) and \(e=\ell^a-1\). Then \(\ell\nmid a\), \(e>0\), and \(e\) is prime to \(\ell\). Choose a finite coefficient extension containing the unramified extension of \(\Q_\ell\) of degree \(a\). Its residue field contains \(\F_{\ell^a}\), and its units contain the cyclic group of Teichmüller lifts of \(\F_{\ell^a}^\times\), of order \(e\). Finite-cocycle lifting descends along this coefficient extension.

The numerical hypothesis gives \(40(HN^2+N^3)<\ell\) and \(N<\ell\). The characters of \(\Sigma\) have order prime to \(\ell\) and are trivial on \(I\cap H'\). Choose a compatible \(e\)-th root of \(\ell\). The tame-descent construction replaces the inertia characters by characters of the resulting tame model, preserves the maximum orbit length of their ratios, and gives good crystalline lattice structures for both \(T_1\) and \(T_2\) after the coefficient extension. The jump set remains \(J\). The good-lattice lifting estimate with these data gives the displayed finite-cocycle inclusion with coefficient \(40\), and the orbit estimate gives \(4r_{\rm orb}\le N^2\). Multiplication by the remaining power of \(\ell\), using \(C_T\ge40\), gives the asserted exponent. Descend along the coefficient extension. The complete proof is contained in \Lean{GlobalTheory.\allowbreak{}InertiaEigenvalues.\allowbreak{}theoremT\_\allowbreak{}all\_\allowbreak{}records}, using \Lean{GlobalTheory.\allowbreak{}TheoremTClosing.\allowbreak{}main}.
\end{proof}

\subsection*{Proof of \ref{B8.42bs1e}}\label{prf:B8.42bs1e}\phantomsection
\begin{proof}[Proof (outline)]
Fontaine comparison carries the inclusions of \(W_j\) into finite direct sums to inclusions of their crystalline period modules over \(K'\). Consequently \(W_j\) is crystalline over \(K'\), its period module is the sum of its \(\Sigma\)-eigenspaces, and its Hodge--Tate jumps belong to \(J\). Apply Lemma~\ref{B8.42bs1d} to the lattices \(L_1,L_2\). Its orbit estimate gives
\[
 r_{\rm orb}\le r',\qquad
 C_T(Nr_{\rm orb}+\eta'(N+r_{\rm orb}))\le e_{\rm fam}.
\]
The sum \(f_*Z_f^1(L_1)+B^1(L_2)\) is an \(O\)-submodule. Multiplication by the difference of the two exponents therefore gives the claimed inclusion. If the family is empty, its finite direct sums and their subrepresentations are zero; equivalently the character set may be replaced by \(\{1\}\). The complete proof is contained in \Lean{GlobalTheory.\allowbreak{}StepLiftPf.\allowbreak{}liftFam\_\allowbreak{}of\_\allowbreak{}eq}.
\end{proof}

\subsection*{Proof of \ref{B8.42bs1f}}\label{prf:B8.42bs1f}\phantomsection
\begin{proof}[Proof (outline)]
A cut cocycle is obtained from finitely many \(x_1,\ldots,x_b\in X_s\), a vector \(E\in\bigoplus_jV_{x_j}\), a finitely generated stable \(O_i\)-submodule \(Y\) containing \((g-1)E\) for every global Galois element \(g\), and an equivariant map \(Y\to U/L_U\), followed by restriction to \(G_{E_w}\). Let \(X\) be the local coefficient-field span of \(Y\). The Fontaine comparison for the common potentially crystalline family puts each of these cocycles in the span of finite pushouts from lattices in \(X\).

In the direct case, \(X\) and \(U\) are subrepresentations of finite sums of the given family. They satisfy the direct alternative of the local pushout hypothesis. The common crystalline extension makes \(U\) de Rham. Choose its Tate dual \(U'=U^\vee(1)\) and the dual lattice \(L_{U'}\). The Bloch--Kato dimension formula and the injection of \(H^0(E_w,U^{\prime})\) into \(\operatorname{Fil}^0D_{\rm dR}(U^{\prime})\), together with \([E_w:\Q_\ell]=1\), give
\[
 \operatorname{rank}_{O_i}H_f^1(E_w,L_{U'})\le k,
 \qquad \operatorname{rank}_{O_i}L_U=k.
\]
The local pushout bound therefore gives \(e(k+k)=2ke\) for this fixed choice of \((x_j,E,Y)\). For two choices, take the concatenation of the spectral lists, the pair of their vectors, and the direct sum of their modules. Projection to either summand recovers its pushouts. Thus these pushout sets form a directed family. Every finite collection of cut cocycles belongs to one such set, and the same length bound holds for their full span.

For the dual case, fix a finite collection of cut cocycles. The common-source construction puts its span in the finite-pushout span of a direct summand \(X\) of a finite sum of the \(V_x\). The representations \(X,U\) are de Rham, and their Tate duals are subrepresentations of finite sums of the dual family. The assumed lifting property therefore gives the dual alternative of the local pushout hypothesis. The same rank estimate and the same pushout bound give \(2ke\). Every finitely generated submodule of \((S_{w,U}+Z_f^1)/Z_f^1\) satisfies this bound; taking their supremum proves the assertion. The complete proofs are contained in \Lean{GlobalTheory.\allowbreak{}Pass.\allowbreak{}LocDirect.\allowbreak{}main} and \Lean{GlobalTheory.\allowbreak{}Pass.\allowbreak{}LocDual.\allowbreak{}main}.
\end{proof}

\subsection*{Proof of \ref{B8.42bs1}}\label{prf:B8.42bs1}\phantomsection
\begin{proof}[Proof (outline)]
\emph{Local and residual packages.}
At the member, the prismatic \(F\)-crystal input of \(\mathscr H_2\) supplies a primewise Wach system with its Fontaine functor and comparison. Fix an integer \(C_T\ge40\) and a real \(C'\ge8C_T\). The four local inputs used below are derived consequences of \(\mathscr H_2\). The local lifting input uses prismatic coefficients, Fontaine comparison and local cohomological hypotheses obtained from Bloch--Kato duality, local Tate duality and the Euler characteristic formula. The pushout input uses the rational and integral finite-condition orthogonality, the top-degree Tate pairing, local finiteness and de Rham duality. The character-block and full local Bloch--Kato inputs are derived from the same local cohomological inputs. First use Lemma~\ref{B8.42bs1d} with parameters \(N,\eta^{\prime},H\). It gives the local lifting input for the specified extension and period datum, with the prime-to-\(\ell\) condition on the inertia characters stated explicitly. The numerical hypotheses give
\[
 160H\le C_THN^2<C_T(HN^2+N^3)<\ell,
 \qquad H<\ell/160,
\]
since \(C_T\ge40\) and \(N\ge2\). For fixed \(N,\eta^{\prime},r_{\rm orb}\), the upper jump \(H\) enters the range condition and leaves the exponent in Lemma~\ref{B8.42bs1d} unchanged. Second, by Bloch--Kato duality over \(K/\Q_\ell\), if \(\ell^e\) lifts finite cohomology across lattice surjections in subrepresentations of finite sums of de Rham \(X\oplus U\) (or of \(X^\vee(1)\oplus U^\vee(1)\)), then for dual lattices \(L_U\subset U\) and \(L_{U'}\subset U^\vee(1)\), the \(O\)-span \(S\) of pushouts of integral finite cocycles of lattices in \(X\) along equivariant maps to \(U/L_U\) satisfies \(\lambda_O((S+Z^1_f)/Z^1_f)\le e(\operatorname{rank}_O H^1_f(K,L_{U'})+\operatorname{rank}_O L_U)\). Lemma~\ref{B8.42bs1c} gives the length deduction from the absorption condition \(\ell^eS\subset Z_f^1\), including the contribution from the torsion of integral finite cohomology. Third, Lemma~\ref{B8.41a} gives the character-block inclusions on divisible cyclotomic characters at level one. Fourth, local Tate duality and the Euler characteristic formula give the local Bloch--Kato finiteness and duality conditions at every place.
\leanref{GlobalTheory.\allowbreak{}T.\allowbreak{}crystallinePeriodBound\_\allowbreak{}all\_\allowbreak{}constant\_\allowbreak{}ge\_\allowbreak{}forty}{Development.\allowbreak{}AbelianVarietyTheory.\allowbreak{}Proofs.\allowbreak{}InertiaEigenvalues.\allowbreak{}CrystallinePeriodBoundAllConstantGeFortyStatement}{11--13}
\leanref{GlobalTheory.\allowbreak{}T.\allowbreak{}local\_\allowbreak{}absorption\_\allowbreak{}length\_\allowbreak{}bound\_\allowbreak{}of\_\allowbreak{}blochKato}{Development.\allowbreak{}AbelianVarietyTheory.\allowbreak{}Proofs.\allowbreak{}InertiaEigenvalues.\allowbreak{}LocalAbsorptionLengthBoundOfBlochKatoStatement}{8--10}
\leanref{GlobalTheory.\allowbreak{}T.\allowbreak{}theorem\_\allowbreak{}tb\_\allowbreak{}of\_\allowbreak{}local\_\allowbreak{}blochKato\_\allowbreak{}duality\_\allowbreak{}euler}{Development.\allowbreak{}AbelianVarietyTheory.\allowbreak{}Proofs.\allowbreak{}InertiaEigenvalues.\allowbreak{}TheoremTbOfLocalBlochKatoDualityEulerStatement}{8--10}
\leanref{GlobalTheory.\allowbreak{}T.\allowbreak{}blochKatoLocal\_\allowbreak{}all\_\allowbreak{}of\_\allowbreak{}tateDuality\_\allowbreak{}eulerChar}{Development.\allowbreak{}AbelianVarietyTheory.\allowbreak{}Proofs.\allowbreak{}InertiaEigenvalues.\allowbreak{}BlochKatoLocalAllOfTateDualityEulerCharStatement}{8--10}

The member rank bounds give \(r\ge2\). At every index \(i\) of nonempty good spectrum, residual irreducibility from the norm-choice conditions (Lemma~\ref{B8.42pa}) and matrix-algebra density show that the image of every residual realization spans the full \(r\times r\) matrix algebra over the residue field of \(O_i\). At such an index \(i\) with \(\ell_i=g_k\), strong multiplicity one, the first provision, residual irreducibility and the wild idempotent choice supply, for each \(\pi\in Z(i)\) and realizing families \((R_x),(\rho_x)\) with semisimple decomposition at \(\pi\), block-separating projector data of denominator valuation at most \(\delta_X\), together with the associated quotient-homomorphism and coboundary bounds. The first provision also gives membership of the characteristic-polynomial coefficients \(x\mapsto\operatorname{coeff}_m\det(1-UR_x(g))\) in the integral Hecke algebra \(T_i\) on \(X_i\), automorphic parameter compatibility, residual eigensystem constancy on \(X_i\), and unramifiedness of \(R_x\) outside \(\{v\mid\ell_i\}\cup S_i\cup Q_i\).
\leanref{GlobalTheory.\allowbreak{}Pass.\allowbreak{}residualRep\_\allowbreak{}spans\_\allowbreak{}matrixAlgebra\_\allowbreak{}of\_\allowbreak{}goodPoint}{Development.\allowbreak{}AbelianVarietyTheory.\allowbreak{}Proofs.\allowbreak{}InertiaEigenvalues.\allowbreak{}ResidualRepSpansMatrixAlgebraOfGoodPoint}{15--25}
\leanref{GlobalTheory.\allowbreak{}typePrimeCongruences\_\allowbreak{}forall}{Development.\allowbreak{}AutomorphicArithmetic.\allowbreak{}ThetaCorrespondence.\allowbreak{}ThetaPairLifts.\allowbreak{}ThirdPlaceSeries.\allowbreak{}LemmaOpenOfIntegrality}{16--17}
\leanref{GlobalTheory.\allowbreak{}ClosingLocalInputs.\allowbreak{}galoisFamilyConditions\_\allowbreak{}of\_\allowbreak{}sumRepresentations\_\allowbreak{}alt}{Development.\allowbreak{}AutomorphicArithmetic.\allowbreak{}ThetaCorrespondence.\allowbreak{}ThetaPairLifts.\allowbreak{}ThirdPlaceSeries.\allowbreak{}AdmG14OfAux}{13--14}

\emph{Per-point type-locus cuts.}
Global Artin reciprocity supplies the local-to-global congruence step, together with the integrality of \(M_0e\) for central idempotents \(e\) of the finite type group \(\Gamma\simeq\GL_r(\F_{g_0})^f\times\GL_r(\F_{g_1})^f\) of order \(M_0=|\GL_r(\F_{g_0})|^f|\GL_r(\F_{g_1})|^f\) and the local reduction steps below.
\leanref{GlobalTheory.\allowbreak{}T.\allowbreak{}no\_\allowbreak{}cut\_\allowbreak{}passes\_\allowbreak{}of\_\allowbreak{}artin\_\allowbreak{}reciprocity}{Development.\allowbreak{}AutomorphicArithmetic.\allowbreak{}ThetaCorrespondence.\allowbreak{}ThetaPairLifts.\allowbreak{}ThirdPlaceSeries.\allowbreak{}StepsNoCutOfArtinStatement}{11--12}
Using the local series compatibility, base change, local-global compatibility, the parameter-Hecke tie, strong multiplicity one and the eigenspan hypothesis, for each index \(i\) with \(\ell_i=g_k\) and each \(\pi\in Z_{\rm good}(i)\) there exist a spectral subset \(X_\pi\subset X_i\) with \(\pi\in X_\pi\) and \(X_\pi\cap Z(i)\subset Z_{\rm good}(i)\), a Hecke-commuting endomorphism \(\beta_\pi\) of the integral Hecke module, and a nonzero scalar \(b_\pi\in O_i\) (namely the image of \(M_0\)) with the following properties. The subset \(X_\pi\) is a type locus at \(g_k\), meaning that there is one assignment \((s_v)_v\) of squarefree monic degree-\(r\) depth-zero labels over \(\F_{g_k}\) at the designated places \(v\) above \(g_k\), each avoiding both the prescribed character value and its inverse in \(\F_{g_k}^\times\), such that every \(x\in X_\pi\) has an automorphic lift carrying \(s_v\) at every \(v\). Moreover \(\beta_\pi\) acts by \(b_\pi\) on the \(\pi\)-eigenspace and by \(0\) on every eigenspace outside \(X_\pi\), with \(\lambda_i(O_i/b_\pi O_i)\le v_{g_k}(M_0)\).
\leanref{GlobalTheory.\allowbreak{}T.\allowbreak{}admBC\_\allowbreak{}of\_\allowbreak{}serLoc}{Development.\allowbreak{}AutomorphicArithmetic.\allowbreak{}ThetaCorrespondence.\allowbreak{}ThetaPairLifts.\allowbreak{}ThirdPlaceSeries.\allowbreak{}BcAllStatement}{11--12}

\emph{Congruence bound on a type locus.}
Fix \(i\) with \(\ell=\ell_i=g_k\), representations \((R_x)\) and \((\rho_x)\), parameters \((\eta_0,s,S,\mathrm{sg},U_w)\) satisfying (i)--(iii) of Definition~\ref{B8.42bs}, a point \(\pi\in Z_{\rm good}(i)\subset Z(i)\), and a type locus \(X_s\subset X_i\) at \(g_k\) containing \(\pi\). Write \(T_i\subset O_i^{X_i}\) for the \(O_i\)-subalgebra of integral Hecke operators on the finite spectrum \(X_i\), and \(T_i|_{X_s}=\{f|_{X_s}:f\in T_i\}\subset O_i^{X_s}\) for its restriction to \(X_s\). Let \(\mathfrak c_{i,X_s,\pi}=\{f(\pi):f\in T_i,\ f|_{X_s\setminus Z(i)}=0\}=\{\bar f(\pi):\bar f\in T_i|_{X_s},\ \bar f|_{X_s\setminus Z(i)}=0\}\subset O_i\) be the evaluation at \(\pi\) of the ideal of operators in \(T_i\) supported on \((X_i\cap Z(i))\cup(X_i\setminus X_s)\) (equivalently, the congruence ideal at \(\pi\) for \(T_i|_{X_s}\) with cut locus \(X_s\cap Z(i)\)), and set \(v_{cc}(i,X_s,\pi)=\lambda_i(O_i/\mathfrak c_{i,X_s,\pi})\). Since \(g_k\) splits completely in \(E\), there are \(2f\) places \(w\mid\ell\), each with \(E_w\simeq\Q_\ell\), and \(\ell=g_k\ne2\).
\leanref{Formalism.\allowbreak{}L2.\allowbreak{}Pj.\allowbreak{}g\_\allowbreak{}ne\_\allowbreak{}two}{Formalism.\allowbreak{}LocalToGlobal.\allowbreak{}ProjectiveLattices.\allowbreak{}ExistsParameterConditions}{23--28}

Fix a place \(w\mid\ell\) and set \(V_x=R_x\chi_{i,1}^{-1}\) for \(x\in X_s\). Combining the depth-zero type identity with local-global compatibility gives zero monodromy and identifies the Weil--Deligne inertia polynomials on \(X_s\). Because all \(x\in X_s\) share the same depth-zero labels \((s_v)_v\), the inertia polynomial tie and paired-character tie show that \((V_x)_{x\in X_s}\) is potentially crystalline over a common finite Galois extension of \(E_w\simeq\Q_\ell\) with inertia acting through a single Frobenius-stable character set \(\Sigma\) containing \(1\) of size at most \(N=n_V+1\). Since this extension is Galois over \(\Q_\ell\), its maximal unramified subfield is Galois and unramified, so the image of its standard embedding in \(B_{\rm cris}\) is stable under \(G_{\Q_\ell}\); this is the period-stability hypothesis of Lemma~\ref{B8.42bs1d}. The characters in \(\Sigma\) are those of the depth-zero type, so they satisfy the inertia conditions of that lemma: they are trivial on \(I\cap G_{K'}\), they are trivial on wild inertia and therefore have order prime to \(\ell\), and the period module is the sum of their eigenspaces; the same holds for \((V_x^\vee(1))_{x\in X_s}\) with \(\Sigma^{-1}\). Put
\[
 N=n_V+1,\qquad r'=\lfloor N^2/4\rfloor,\qquad
 e_{\max}=C_T\bigl(Nr'+(\eta_0+t+1)(N+r')\bigr)\in\N.
\]
Since \(r\ge2\) and \(t\ge1\), we have \(N\ge4\) and \(r'\ge4\). By hypothesis (i) and \(S\ge0\), \(1\le\eta_0+t+1<s+S\) and \(C_T((s+S)N^2+N^3)<\ell\). Let \(M\) be either a stable \(O_i[G_E]\)-lattice \(W_\pi\) in \(\rho_\pi\chi_{i,1}^{-1}\) (rank \(k_M=r\)) or \(O_i(\chi_{i,c}\chi_{i,1}^{-1})\) for \(1\le c\le t\) (rank \(k_M=1\)); since \(\pi\in Z(i)\), the semisimple decomposition of \(R_\pi\) embeds \(M\otimes_{O_i}K_i\) equivariantly into \(V_\pi\). Let \(S_{w,M}\subset Z^1(E_w,M\otimes_{O_i}K_i/O_i)\) be the \(O_i\)-span of cut cocycles of \(M\) at \(w\) formed from \(X_s\) (defined as in Lemma~\ref{B8.42bp}(b) with \(x_1,\dots,x_{a'}\in X_s\)). Write \(Z^1_f\) for the submodule of cocycles whose classes lie in \(H^1_f\), and \(B^1\) for coboundaries. Write \(Z^1_{\rm ur}\) for cocycles vanishing on \(I_{E_w}\).
\begin{itemize}
\item If \(\mathrm{sg}(w)=1\), hypothesis (ii) and \(t\ge1\) put all jumps at \(w\) of every \(V_x\) (\(x\in X_s\)) into \(U_w\cup\{0\}\subset[-(\eta_0+t+1),\,s+S]\), of size at most \(n_V+1\le N+1\). Lemma~\ref{B8.42bs1e}, with \(\eta^{\prime}=\eta_0+t+1\) and \(H=s+S\), gives the common lifting exponent \(e_{\max}\) for subrepresentations of finite sums of \((V_x)_{x\in X_s}\). The direct branch of Lemma~\ref{B8.42bs1f} gives \(\lambda_i((S_{w,M}+Z^1_f)/Z^1_f)\le2k_Me_{\max}\).
\item If \(\mathrm{sg}(w)=-1\), dual jump reflection \(b\mapsto-1-b\) and \(t\ge1\) put all jumps at \(w\) of every \(V_x^\vee(1)\) (\(x\in X_s\)) into \(\{0\}\cup\{-1-b:b\in U_w\}\subset[-(\eta_0+t+1),\,s+S]\), of size at most \(n_V+1\le N+1\). Apply Lemma~\ref{B8.42bs1e} to \((V_x^\vee(1))_{x\in X_s}\) with the same parameters. The dual branch of Lemma~\ref{B8.42bs1f} gives \(\lambda_i((S_{w,M}+Z^1_f)/Z^1_f)\le2k_Me_{\max}\).
\end{itemize}
For the character blocks \(O_i(\chi_{i,c}\chi_{i,1}^{-1})\), the character-block theorem gives \(Z^1_f\subset B^1\) for \(2\le c\le t\) and \(Z^1_f\subset Z^1_{\rm ur}\) for \(c=1\).
\leanref{GlobalTheory.\allowbreak{}Closing.\allowbreak{}closing\_\allowbreak{}at\_\allowbreak{}record}{Development.\allowbreak{}AbelianVarietyTheory.\allowbreak{}Proofs.\allowbreak{}InertiaEigenvalues.\allowbreak{}ClosingAtRecord}{11--20}

Apply the global congruence-module estimate on \(X_s\) with the restricted Hecke algebra \(T_i|_{X_s}\) at \(\pi\in X_s\cap Z(i)\) (noting \(Z_{\rm good}(i)\ne\varnothing\)), using membership of the restricted characteristic-polynomial coefficients of \(R_x(g)\) (\(x\in X_s\)) in \(T_i|_{X_s}\), absence of nonzero \(G_E\)-invariants in \(V_x\) for \(x\in X_s\setminus Z(i)\) (the sixth provision, ensuring the relation determinants vanish on \(X_s\setminus Z(i)\) and hence evaluate into \(\mathfrak c_{i,X_s,\pi}\)), the full residual matrix-algebra span of \(\bar\rho_\pi\), the block-separating projector at \(\pi\) of denominator valuation at most \(\delta_X\), unramifiedness outside \(\{v\mid\ell\}\cup S_i\cup Q_i\), inertia annihilation of the character-block cocycles at \(S_i\) and quadratic inertia bounds at \(Q_i\) from the fifth provision and tame generator conditions, global Artin reciprocity on the unramified \(c=1\) block, and finiteness of \(s_\pi\) and \(s_{\ell,e}(E)\). Summing the local bounds over the \(2f\) places \(w\mid\ell\) gives
\[
 v_{cc}(i,X_s,\pi)\le
 2f(2re_{\max}+2te_{\max})+s_\pi+\sum_{0<e<t}s_{\ell,e}(E)+v_\ell(h_E)
 +4n_V\delta_X+t\sum_{v\in Q_i}\lambda_i(O_i/(\zeta_{i,v}-1)),
\]
where \(2f(2r+2t)e_{\max}=4fn_VC_T(Nr'+(\eta_0+t+1)(N+r'))\).
\leanref{GlobalTheory.\allowbreak{}TameDesign.\allowbreak{}main}{Development.\allowbreak{}AbelianVarietyTheory.\allowbreak{}Proofs.\allowbreak{}InertiaEigenvalues.\allowbreak{}TameDescentConditionsEveryLocalisation}{137--145}
\leanref{GlobalTheory.\allowbreak{}StepGlobalPf.\allowbreak{}main\_\allowbreak{}alt\_\allowbreak{}alt}{Development.\allowbreak{}AutomorphicArithmetic.\allowbreak{}ThetaCorrespondence.\allowbreak{}ThetaPairLifts.\allowbreak{}ThirdPlaceSeries.\allowbreak{}ImportsMoreMore}{15--22}
\leanref{GlobalTheory.\allowbreak{}CutCore.\allowbreak{}main}{Development.\allowbreak{}AbelianVarietyTheory.\allowbreak{}Proofs.\allowbreak{}InertiaEigenvalues.\allowbreak{}LengthQuotEqZeroIff}{175--244}
\leanref{GlobalTheory.\allowbreak{}Pass.\allowbreak{}cutBound\_\allowbreak{}of\_\allowbreak{}steps\_\allowbreak{}alt}{Development.\allowbreak{}AutomorphicArithmetic.\allowbreak{}ThetaCorrespondence.\allowbreak{}ThetaPairLifts.\allowbreak{}ThirdPlaceSeries.\allowbreak{}CutBoundOfSteps}{15--78}
\leanref{GlobalTheory.\allowbreak{}conclusion\_\allowbreak{}of\_\allowbreak{}typeLocus\_\allowbreak{}alt}{Development.\allowbreak{}AutomorphicArithmetic.\allowbreak{}ThetaCorrespondence.\allowbreak{}ThetaPairLifts.\allowbreak{}ThirdPlaceSeries.\allowbreak{}ConclusionOfAdmissibles}{16--28}

\emph{Passage to \(\lambda_i(U_i/L_i)\).}
If \(Z_{\rm good}(i)=\varnothing\), then \(V_Z=0\) in Definition~\ref{B8.40}, so \(U_i=L_i=0\) and both sides of (P) vanish. Assume henceforth that \(Z_{\rm good}(i)\ne\varnothing\). For each \(\pi\in Z_{\rm good}(i)\), take the type-locus cut \((X_\pi,\beta_\pi,b_\pi)\) constructed above. Lemma~\ref{B8.42bs1b}, with \(Y=X_\pi\), gives \(g_\pi\in T_i\), \(c_\pi=g_\pi(\pi)\ne0\), and a Hecke-commuting endomorphism \(\alpha_\pi=g_\pi\beta_\pi\) which vanishes outside \(Z_{\rm good}(i)\) and acts by \(a_\pi=c_\pi b_\pi\ne0\) on the \(\pi\)-eigenspace. It satisfies
\[
 \lambda_i(O_i/a_\pi O_i)
 \le v_{cc}(i,X_\pi,\pi)+v_\ell(M_0).
\]
\leanref{GlobalTheory.\allowbreak{}CutFam.\allowbreak{}perPoint}{Development.\allowbreak{}AbelianVarietyTheory.\allowbreak{}Proofs.\allowbreak{}InertiaEigenvalues.\allowbreak{}ExistsOfMemCcAt}{114--154}
\leanref{GlobalTheory.\allowbreak{}CutFam.\allowbreak{}fam\_\allowbreak{}of\_\allowbreak{}cut}{Development.\allowbreak{}AbelianVarietyTheory.\allowbreak{}Proofs.\allowbreak{}InertiaEigenvalues.\allowbreak{}ExistsOfMemCcAt}{155--177}

Lemma~\ref{B8.42bs1a}, with \(\mathcal S=Z_{\rm good}(i)\) and the endomorphisms \(\alpha_\pi\) above, gives
\[
 \lambda_i(U_i/L_i)\le\sum_{\pi\in Z_{\rm good}(i)}\mathrm{mult}(i,\pi)\,\lambda_i(O_i/a_\pi O_i)
 \le\sum_{\pi\in Z_{\rm good}(i)}\mathrm{mult}(i,\pi)\bigl(v_{cc}(i,X_\pi,\pi)+v_\ell(M_0)\bigr).
\]
\leanref{GlobalTheory.\allowbreak{}EtaFromCut.\allowbreak{}lenUnit\_\allowbreak{}le\_\allowbreak{}of\_\allowbreak{}family}{Development.\allowbreak{}AbelianVarietyTheory.\allowbreak{}Proofs.\allowbreak{}InertiaEigenvalues.\allowbreak{}LenUnitLeOfFamily}{64--94}
\leanref{GlobalTheory.\allowbreak{}EtaFromCut.\allowbreak{}etaBound\_\allowbreak{}of\_\allowbreak{}fam}{Development.\allowbreak{}AbelianVarietyTheory.\allowbreak{}Proofs.\allowbreak{}InertiaEigenvalues.\allowbreak{}LenUnitLeOfFamily}{97--106}
\leanref{GlobalTheory.\allowbreak{}EtaFromCut.\allowbreak{}etaBound\_\allowbreak{}of\_\allowbreak{}cut}{Development.\allowbreak{}AbelianVarietyTheory.\allowbreak{}Proofs.\allowbreak{}InertiaEigenvalues.\allowbreak{}LenUnitLeOfFamily}{107--109}
Finally, \(N=n_V+1\le2n_V\) and \(r'\le N^2/4\le n_V^2\) give \(Nr'+(\eta_0+t+1)(N+r')\le2n_V^3+3n_V^2(n_V+\eta_0+t)\le5n_V^2(n_V+\eta_0+t)\), which multiplied by \(4fn_VC_T\) is at most \(4fn_VC'n_V^2(n_V+\eta_0+t)\) since \(C'\ge8C_T\ge5C_T\). From \(\log g_{1-k}\le2\log g_k\) we get \(g_{1-k}\le g_k^2\), so \(v_{g_k}(g_{1-k}^a-1)\le2a\) for \(1\le a\le r\) and \(v_\ell(M_0)=fr(r-1)/2+f\sum_{a=1}^r v_{g_k}(g_{1-k}^a-1)\le fr(r-1)+fr(r+1)=2fr^2\). Hence each summand \(v_{cc}(i,X_\pi,\pi)+v_\ell(M_0)\) is at most \(R_i(\pi)\), proving \(\lambda_i(U_i/L_i)\le\sum_{\pi\in Z_{\rm good}(i)}\mathrm{mult}(i,\pi)R_i(\pi)\).
\leanref{GlobalTheory.\allowbreak{}typePrimeBound\_\allowbreak{}of\_\allowbreak{}etaBound\_\allowbreak{}alt}{Development.\allowbreak{}AutomorphicArithmetic.\allowbreak{}ThetaCorrespondence.\allowbreak{}ThetaPairLifts.\allowbreak{}ThirdPlaceSeries.\allowbreak{}TypePrimeBoundOfEtaBound}{13--30}
\leanref{GlobalTheory.\allowbreak{}admT\_\allowbreak{}of\_\allowbreak{}adm5\_\allowbreak{}eigSpan\_\allowbreak{}tieInePoly}{Development.\allowbreak{}AutomorphicArithmetic.\allowbreak{}ThetaCorrespondence.\allowbreak{}ThetaPairLifts.\allowbreak{}ThirdPlaceSeries.\allowbreak{}ConclusionOfAdmissibles}{29--41}

The complete proof is contained in \Lean{GlobalTheory.\allowbreak{}ClosingLocalInputs.\allowbreak{}main\_\allowbreak{}alt\_\allowbreak{}alt}.
\end{proof}

\subsection*{Proof of \ref{B8.42bt}}\label{prf:B8.42bt}\phantomsection
\begin{proof}
If \(Z_{\rm good}(i)\) is empty, then \(V_Z=0\) in Definition~\ref{B8.40}, so \(U_i=L_i=0\) and both sides vanish. Otherwise the type-prime clause of the parameter split bound in the setup supplies the hub families \((R_x)\), \((\rho_x)\) with their semisimple decompositions on \(Z(i)\). It also supplies \(s,S,\mathrm{sg},U_w\) with \(\eta_X+t+1<s\), \(C_T((s+S)(n_V+1)^2+(n_V+1)^3)<g_k\), opposite signs at conjugate places, and (ii) of Definition~\ref{B8.42bs} with \(\eta_0=\eta_X\). The size conditions on the type primes give \(g_{1-k}\le2g_k\), hence \(\log g_{1-k}\le\log2+\log g_k\le2\log g_k\). The setup assumes finite Selmer lengths, and the finiteness hypothesis of Definition~\ref{B8.40} gives finite \(s_{\ell,e}(E)\). Statement (P) now gives the bound.
\leanref{GlobalTheory.\allowbreak{}T.\allowbreak{}split\_\allowbreak{}bound\_\allowbreak{}of\_\allowbreak{}params\_\allowbreak{}alt}{Development.\allowbreak{}AutomorphicArithmetic.\allowbreak{}ThetaCorrespondence.\allowbreak{}ThetaPairLifts.\allowbreak{}ThirdPlaceSeries.\allowbreak{}ImportsMoreExtraStatement}{13--14}
\leanref{GlobalTheory.\allowbreak{}NLinkTyp.\allowbreak{}main\_\allowbreak{}alt\_\allowbreak{}alt}{Development.\allowbreak{}AutomorphicArithmetic.\allowbreak{}ThetaCorrespondence.\allowbreak{}ThetaPairLifts.\allowbreak{}ThirdPlaceSeries.\allowbreak{}EtaLocLeRhs}{60--90}
\end{proof}

\subsection*{Proof of \ref{B8.42bua}}\label{prf:B8.42bua}\phantomsection
\begin{proof}
Set \(Q=\sum_{i\in S}Ob_i\subset M\). The family \((b_i)_{i\in S}\) is \(O\)-linearly independent and spans \(Q\), so it forms an \(O\)-basis of \(Q\). The submodule \(P'=P\), viewed inside \(Q\), is \(O\)-saturated in \(Q\): if \(a\in O\setminus\{0\}\), \(x\in Q\), and \(ax\in P'\), then \(x\in P\) by saturation of \(P\) in \(M\). Over the discrete valuation ring \(O\), a saturated submodule of a finite free module is a direct summand with an adapted basis, so \(Q\) admits an \(O\)-basis \((c_j)_{j\in S}\) and a subset \(S''\subset S\) with \(P=\sum_{j\in S''}Oc_j\).

Define \(b'_i=c_i\in Q\subset M\) for \(i\in S\) and \(b'_i=b_i\) for \(i\notin S\). Every \(y\in Q\) has zero \(b_i\)-coordinate for all \(i\notin S\). Hence if \(\sum_{i\in I}g_ib'_i=0\) with \(g_i\in O\), reading the \(b_{i_0}\)-coordinate for each \(i_0\notin S\) gives \(g_{i_0}=0\), and then \(\sum_{j\in S}g_jc_j=0\) in \(Q\) gives \(g_j=0\) for all \(j\in S\). Moreover, for \(i\in S\) the vector \(b_i\in Q\) lies in \(\sum_{j\in S}Oc_j=\sum_{j\in S}Ob'_j\), while for \(i\notin S\) one has \(b_i=b'_i\). Thus \((b'_i)_{i\in I}\) spans \(M\) and is an \(O\)-basis of \(M\), satisfying \(b'_i=b_i\) for \(i\notin S\) and \(\sum_{i\in S}Ob'_i=\sum_{j\in S}Oc_j=Q\). Finally, \(P=\sum_{j\in S''}Ob'_j\); since \((b'_i)_{i\in I}\) is a basis of \(M\), a basis vector \(b'_i\) belongs to \(\sum_{j\in S''}Ob'_j\) if and only if \(i\in S''\), so \(P=\sum_{i\in I,\,b'_i\in P}Ob'_i\).
\end{proof}

\subsection*{Proof of \ref{B8.42bud}}\label{prf:B8.42bud}\phantomsection
\begin{proof}[Proof (outline)]
The crystalline-extension and integral Herr comparisons, applied to a basis of \(P\), give \(O\)-linear maps \(y,z,\mu_i,\nu_i\) and an integer \(N\ge0\). Put \(y_i=e_i y\) and \(z_i=e_i z\). The pair \((y(f),z(f))\) represents \([g(f)]\), and \(\mu_i(f),\nu_i(f)\in\mathcal D_i^\Delta(T_2)\). Put
\[
 r_i=\ell^Ny_i+\varphi\mu_{i-1}-\mu_i,
 \qquad (M_i,h_i)=
 \begin{cases}(\beta_i+S,\delta),&i\in J_-,\\
 (\delta,\beta_i+S),&i\notin J_-.
 \end{cases}
\]
The extension construction gives, for every coefficient vector,
\[
 \ell^Nz_i+(\gamma-1)\mu_i\in\mathcal U^1(T_2),\quad
 Q_0^{M_i}r_i\in N(T_2),\quad
 \ell^{s_i}Q_0^{h_i}r_i\in\mathcal R_{i-1}^{(1)}(T_2)
\]
for some \(s_i\ge0\), and
\(\ell^N\nu_i-Q_0^a\mu_i\in N(T_2)\).
The integer \(N\) belongs to this auxiliary presentation and will disappear from the final exponent.
\leanref{Formalism.\allowbreak{}MainNotions.\allowbreak{}RamifiedCoefficientChains.\allowbreak{}Sh.\allowbreak{}ExtData}{Formalism.\allowbreak{}SharpRestriction.\allowbreak{}LevelLetterCocycles.\allowbreak{}ExpRules}{106--117}

\emph{The first correction.}
On the localized invariant component put \(D=\pi_0^{-1}(\gamma-1)\), and define \(c(t)\in O\), for \(t\in\Z\), by
\[
 \ell c(t)=1-\chi_\ell(\gamma)^{-t}.
\]
Then \(c(0)=0\), and \(v_\ell(c(t))=v_\ell(t)\) for \(t\ne0\). For \(c\in J_-\), let \(S_c\) be the set of nonzero integers \(u+j\), where \(u\) is a Hodge--Tate jump of \(T_2\) over the \(c\)-component and \(1\le j\le a\). Exactness of the jumps and the two prescribed intervals give
\[
 S_c\subset
 [1-\beta_c-S,a-\beta_c]\cup[1-\delta,a+\delta],
 \qquad a+\delta<\beta_c.
\]
Set \(u_1=\max_{c\in J_-}\sum_{t\in S_c}v_\ell(t)\). Counting multiples of \(\ell^j\) in these intervals gives
\[
 \sum_{t\in S_c}v_\ell(t)
 \le\frac{2a+2\delta+S}{q}
       +\log_\ell(2H)+1=B_1.
\]
Choose \(b_c\in O\) so that
\[
 R_c(X)=b_c\prod_{t\in S_c}(X-c(t)),\qquad R_c(0)=\ell^{u_1},
 \qquad R_{c,1}(X)=\frac{R_c(X)-R_c(0)}X.
\]
Such \(b_c\) is integral because the valuation of the product of the constant terms is at most \(u_1\). The letter-operator condition gives
\[
 D R_c(D)x\in N(T_2)
 \quad\text{if }x\in\mathcal D_c^\Delta(T_2),\quad Q_0^a x\in N(T_2).
\]
Define the correction, before restricting \(P\), by
\[
 m_c=R_c(D)(Q_0^{-a}\nu_c)
       +R_{c,1}(D)(\pi_0^{-1}z_c),\qquad
 u_c=\ell^Nm_c-\ell^{u_1}\mu_c.
\]
For \(u_{c,0}=\ell^NQ_0^{-a}\nu_c-\mu_c\) and
\(v_c=\pi_0^{-1}(\ell^Nz_c+(\gamma-1)\mu_c)\), one has
\[
 u_c=R_c(D)u_{c,0}+R_{c,1}(D)v_c.
\]
Here \(Q_0^au_{c,0}\) and \(v_c\) are integral. The killing identity and preservation of the Wach lattice by \(D\) imply
\[
 Q_0^au_c\in N(T_2),\quad
 (\gamma-1)u_c\in\mathcal U^1(T_2),\quad
 \zeta_c:=\ell^{u_1}z_c+(\gamma-1)m_c\in\mathcal U^1(T_2).
\]
For the last assertion, use
\(\ell^N\pi_0^{-1}\zeta_c=\ell^{u_1}v_c+Du_c\)
and saturation of the invariant Wach lattice.

Let \(\mathcal F_c^j\) be the image of the \(c\)-component of the crystalline filtration \(\operatorname{Fil}^j\) in the de Rham period space. At a primitive \(\ell\)-th root, expand \(u_c\) as a Laurent series. Its constant coefficient belongs to \(\mathcal F_c^{-a}\), and \(u_c\in N(T_2)\) if and only if that coefficient belongs to \(\mathcal F_c^0\). The quotient \(\mathcal F_c^{-a}/\mathcal F_c^0\) has dimension
\(\sum_{-a\le j<0}\operatorname{mult}_{\tau_c}(j)\le m\).
Intersect the kernels of these constant-coefficient maps successively. This defines \(P_1\subset P\), with one generator step for each \(c\in J_-\), and \(u_c(P_1)\subset N(T_2)\).
\leanref{Formalism.\allowbreak{}MainNotions.\allowbreak{}RamifiedCoefficientChains.\allowbreak{}S9.\allowbreak{}ExtensionLemmaAlt.\allowbreak{}first\_\allowbreak{}lift\_\allowbreak{}of\_\allowbreak{}hypotheses}{Formalism.\allowbreak{}SharpRestriction.\allowbreak{}LevelLetterCocycles.\allowbreak{}FirstLiftOfHypotheses}{19--190}
\leanref{Formalism.\allowbreak{}MainNotions.\allowbreak{}RamifiedCoefficientChains.\allowbreak{}S9.\allowbreak{}ExtensionLemmaAlt.\allowbreak{}step1Pass\_\allowbreak{}pass}{Formalism.\allowbreak{}SharpRestriction.\allowbreak{}LevelLetterCocycles.\allowbreak{}MultVNonneg}{99--173}

\emph{The second correction.}
Put \(s=\widehat\sigma\). The level-one uniformizer truncation with parameters \((M,\delta,s)\), applied to \(\zeta_c\), gives maps \(t_c:P_1\to\mathcal R_c^{(1)}(T_2)\), linear after choices on an \(O\)-basis, such that
\[
 \xi_c:=\ell^s\varphi\zeta_c-(\gamma-1)t_c
       \in\pi_0^M Q_0^\delta\mathcal R_c^{(1)}(T_2)
       \subset\mathcal U^M(T_2).
\]
The maps preserve components, with \(t_c,\xi_c\) in component \(c+1\). The truncation is linear modulo \(\pi_0^M Q_0^\delta\mathcal R_c^{(1)}\), which is the precision used here. Define
\[
 \widehat m_c=\ell^s\varphi m_c+\ell^{u_1+s}y_{c+1}-t_c,
 \qquad
 \widehat u_c=\ell^N\widehat m_c-\ell^{u_1+s}\mu_{c+1}.
\]
The identity
\[
 \widehat u_c=\ell^s\varphi u_c+\ell^{u_1+s}r_{c+1}-\ell^Nt_c
\]
and the Frobenius denominator bound \(\delta\) on component \(c+1\) give
\(Q_0^\delta\widehat u_c\in N(T_2)\).
Also \((\gamma-1)\widehat u_c\in N(T_2)\), by the formulas for \(\xi_c\) and the original extension coordinates. The constant coefficient of its Laurent expansion lies in
\(\mathcal F_{c+1}^{-\delta}\), and its image modulo \(\mathcal F_{c+1}^0\) vanishes exactly when \(\widehat u_c\) is integral. This quotient has dimension at most \(m\). Its successive kernels define \(P_2\subset P_1\), with a further \(|J_-|\) generator steps.

On \(P_2\) set
\[
 w_c=\varphi\widehat m_{c-2}
       +\ell^{u_1+s}y_c-\ell^sm_c.
\]
Its two useful identities are
\[
 \ell^Nw_c=\varphi\widehat u_{c-2}-\ell^su_c+\ell^{u_1+s}r_c,
 \qquad
 (\gamma-1)w_c=\varphi\xi_{c-2}-\ell^s\zeta_c.
\]
The first, the height bounds, and saturation give
\(Q_0^\delta w_c\in\mathcal R_{c-1}^{(1)}(T_2)\) and
\(Q_0^{\beta_c+S}w_c\in N(T_2)\). The latter identity supplies the finite-pole condition with \(k=1\) and \(a_0=\beta_c+S\). The second displayed identity gives \((\gamma-1)w_c\in\mathcal U^1(T_2)\), since Frobenius preserves \(\mathcal U^M\) and \(M\ge1\). Thus \(w_c\in\mathcal V_c(T_2)\).

Put \(Y_c=\varphi w_c+t_c\). Direct expansion gives
\[
 Y_c=\ell^{u_1+s}(y_{c+1}+\varphi y_c)
            +\varphi^2\widehat m_{c-2}-\widehat m_c,
 \qquad(\gamma-1)Y_c=\varphi^2\xi_{c-2}-\xi_c.
\]
The bijection \(\varphi-1\) on \(\mathcal U^M(T_2)\) completes \(Y_c\) to a Herr cocycle. Applying the same truncation to \(w_c\) gives a Herr cocycle \(X_c=(z_c^{\rm H},v_c^{\rm H})\) whose first coordinate differs from \(\ell^sY_c\) by an element of \(\mathcal U^M\), and whose second coordinate belongs to \(\mathcal U^M\). Thus its class is \(\ell^s\) times the class just constructed. Summing over \(c\in J_-\), the shift \(c\mapsto c-2\) permutes \(J_-\). With
\[
 m_E=\sum_{c\in J_-}
       (\ell^{u_1+s}y_c+\widehat m_c+\varphi\widehat m_c),
\]
one has \(\sum_cY_c=\ell^{u_1+s}y+(\varphi-1)m_E\), and the corresponding second-coordinate identity holds modulo \(\mathcal U^M\). Herr boundaries and the bijection \(\varphi-1\) therefore give the asserted sum of classes at exponent \(u_1+2s\).

For the lifting assertion, truncate \(u\) in the source with the same parameters. The truncation preserves finite Frobenius poles, and the finite-class criterion for Herr cocycles gives a source cocycle \(Y=(z_u^{\rm H},v_u^{\rm H})\) with finite class and \(v_u^{\rm H}\in\mathcal U^M(T_1)\). Naturality, linearity modulo \(\pi_0^M Q_0^\delta\mathcal R_c^{(1)}\), and vanishing of a truncation of an element of \(\mathcal U^1\) give
\[
 \ell^{e_*}z_c^{\rm H}-F_Dz_u^{\rm H}
       \in\pi_0^M Q_0^\delta\mathcal R_c^{(1)}
       \subset\mathcal U^M(T_2).
\]
Together with the second-coordinate memberships, this gives the required congruence for \(Y\). There are \(2|J_-|=n\) generator steps, and
\(u_1+2s\le B_1+2\sigma+2=B_{\rm red}\).
The complete proof is contained in \Lean{Formalism.\allowbreak{}MainNotions.\allowbreak{}RamifiedCoefficientChains.\allowbreak{}S9.\allowbreak{}ExtensionLemmaAlt.\allowbreak{}cocyclesToWachVectorsBound\_\allowbreak{}of\_\allowbreak{}extTool}, using \Lean{Formalism.\allowbreak{}MainNotions.\allowbreak{}RamifiedCoefficientChains.\allowbreak{}S9.\allowbreak{}LocalLemmasFromLiterature.\allowbreak{}local\_\allowbreak{}lemmas\_\allowbreak{}of\_\allowbreak{}literature} and \Lean{Formalism.\allowbreak{}MainNotions.\allowbreak{}N.\allowbreak{}Pv.\allowbreak{}LocalLengthBound.\allowbreak{}rhoLift\_\allowbreak{}of}.
\end{proof}

\subsection*{Proof of \ref{B8.42buea}}\label{prf:B8.42buea}\phantomsection
\begin{proof}[Proof (outline)]
Put \(Q_1=\varphi(Q_0)\) and \(h_2=\beta_{i-1}+S\). Apply the two-Frobenius uniformizer truncation to \(y\), with parameters \((2,1,2H,s_2)\). It gives \(\lambda_y\in\mathcal D_i^\Delta(T_2)\) such that
\[
 Q_1^{h_2}\lambda_y\in\mathcal R_{i-1}^{(1)}(T_2),\qquad
 Q_1^{h_2}(\gamma-1)\lambda_y
       =\pi_0Q_0^{a+\beta_i}v
       \quad(v\in\mathcal N_i(T_2)).
\]
Set \(x=Q_0^{-\beta_i}Q_1^{h_2}\lambda_y\) in the localized module. The height bound gives \(Q_0^{\beta_i+\delta}x\in N(T_2)\). We next prove that \(x\) itself is integral.

At a primitive root, write the expansion of \(Q_0\) as \(tU(t)\), with \(U(0)\ne0\), and let \(d_x(k)\) be the Laurent coefficients of \(x\). The truncated horizontal comparison implies that the coefficients of \(U(t)^{\beta_i+\delta}d_x(t)\) belong to
\[
 V_k=\begin{cases}
 0,&k<-\beta_i,\\
 L'\cdot\ell^{s_2}\iota(d),&-\beta_i\le k<2H-\beta_i,\\
 \text{the ambient period space},&k\ge2H-\beta_i,
 \end{cases}
\]
where \(\iota:D_{\rm cris}\hookrightarrow D_{\rm dR}\) is the comparison inclusion. The spaces \(V_k\) increase with \(k\). Solving the coefficient equations successively, starting at the lower bound of the Laurent support, shows that \(d_x(k)\in V_k\): the coefficient to be solved has nonzero scalar multiplier \(U(0)^{\beta_i+\delta}\), and the earlier coefficients lie in \(V_k\). In the middle range, \(d\in\operatorname{Fil}^{\beta_i}\) gives \(d_x(k)\in\operatorname{Fil}^{-k}\). In the final range, \(2H-\beta_i\ge\delta\), and the lower jump bound \(-\delta\) gives the same inclusion. The expansion criterion for Wach membership therefore gives \(x\in\mathcal N_i(T_2)\).

At the expansion with root \(1\), both \(Q_0\) and \(Q_1\) have constant coefficient \(\ell\). The negative coefficients vanish. Comparing constant coefficients in \(Q_0^{\beta_i}x=Q_1^{h_2}\lambda_y\) gives
\[
 \operatorname{cmp}_{T_2}(x)=\ell^{s_2+h_2-\beta_i}d.
\]

To lift \(x\), put \(D_j=\pi_0^{-1}(\gamma-1)\) on \(\mathcal D(T_j)^\Delta\), define \(c(t)\) as in the preceding proof, and set
\[
 I=[\beta_i-a+1,\beta_i]\cap\Z,\quad
 J=([\beta_i+1,\beta_i+S+a]\cup[1-\delta,\delta+a])\cap\Z,
 \quad P_U(X)=\prod_{u\in U}(X-c(u)).
\]
The thickness inequality gives \(2a+\delta\le\beta_i\), so \(I\cap J=\varnothing\). The displayed equation for \((\gamma-1)\lambda_y\) gives
\(P_I(D_2)x=Q_0^aw\) with \(w\in\mathcal N_i(T_2)\). The Wach cokernel bound lifts \(Q_0^aw\) and \(Q_0^ax\) to invariant vectors \(y_I,y_0\in\mathcal N_i(T_1)\); average any chosen lifts over \(\Delta\), using that \(|\Delta|\) is a unit. The letter-operator condition gives
\[
 y_J=P_J(D_1)(Q_0^{-a}y_0)\in\mathcal N_i(T_1),
 \quad F_Dy_I=P_I(D_2)x,
 \quad F_Dy_J=P_J(D_2)x.
\]
For \(j\in I\), the valuation identity
\(v_\ell(c(j)-c(k))=v_\ell(j-k)\) gives
\[
 v_\ell\bigl(P_J(c(j))P_I'(c(j))\bigr)
 =\sum_{k\in J}v_\ell(j-k)+\sum_{k\in I\setminus\{j\}}v_\ell(j-k)
 \le\frac{3a+S+2\delta-1}{q}+\log_\ell H.
\]
The left side is an integer, hence is at most \(E_i\). Consequently
\[
 B(X)=\sum_{j\in I}
 \frac{\ell^{E_i}}{P_J(c(j))P_I'(c(j))}
       \prod_{k\in I\setminus\{j\}}(X-c(k))
\]
has coefficients in \(O\). It satisfies \(B(c(j))P_J(c(j))=\ell^{E_i}\) for all \(j\in I\). Division by the monic polynomial \(P_I\) gives \(A_0\in O[X]\) with
\(\ell^{E_i}=A_0P_I+BP_J\).
Thus
\(z=A_0(D_1)y_I+B(D_1)y_J\)
is integral, invariant, and satisfies \(F_Dz=\ell^{E_i}x\). This proves the comparison formula. Here \(H\ge1\) and \(a\ge3\), so \(E_i\ge0\). The estimate for \(E_i\) follows from \(a<A+1\), \(q\ge10\), and \(\lambda=1+\log_\ell(2H)\). The complete proof is contained in \Lean{Formalism.\allowbreak{}MainNotions.\allowbreak{}RamifiedCoefficientChains.\allowbreak{}Asm.\allowbreak{}lemmaGs\_\allowbreak{}rec}, using \Lean{Formalism.\allowbreak{}MainNotions.\allowbreak{}RamifiedCoefficientChains.\allowbreak{}S9.\allowbreak{}lemGA\_\allowbreak{}of\_\allowbreak{}pieces} and \Lean{Formalism.\allowbreak{}MainNotions.\allowbreak{}RamifiedCoefficientChains.\allowbreak{}S9.\allowbreak{}lemGB\_\allowbreak{}of}.
\end{proof}

\subsection*{Proof of \ref{B8.42bue}}\label{prf:B8.42bue}\phantomsection
\begin{proof}[Proof (outline)]
Fix \(c\in J_-\). Write \(B=\beta_c+S\), \(h_0=\beta_{c-2}+S\), and
\(\mathcal F^j=\iota(D_{{\rm cris},c}(T_2[1/\ell]))\cap\operatorname{Fil}^jD_{\rm dR}\).
Choose a primitive \(\ell\)-th root \(\zeta\), and expand \(w(f)\) at \(\zeta\) after clearing its pole by \(Q_0^{\delta+B}\). Division by the resulting power of \(tU(t)\), where \(U(0)\ne0\), gives Laurent coefficients \(D_f(k)\) which depend \(O\)-linearly on \(f\). The relation \((\gamma-1)w(f)\in N(T_2)\), together with the action \(\gamma:t\mapsto\chi_\ell(\gamma)t\), gives
\[
 D_f(k)\in\mathcal F^{-k}\quad(k\ne0).
\]

The constant-coefficient reduction gives an integer \(a_c\ge0\) and \(x_f\in\mathcal N_{c-2}(T_2)\) with
\[
 a_c\le\sigma_2+\beta_{c-1}+S-\frac{\beta_c}{\ell}+\lambda+3,
 \qquad
 \ell^{a_c}D_f(0)-\iota\varphi^2\operatorname{cmp}_{T_2}(x_f)
       \in\mathcal F^{-\delta}.
\]
\leanref{Formalism.\allowbreak{}MainNotions.\allowbreak{}N.\allowbreak{}Pv.\allowbreak{}LocalLengthBound.\allowbreak{}H2Bound}{Formalism.\allowbreak{}SharpRestriction.\allowbreak{}LevelLetterCocycles.\allowbreak{}QuotientImageDivisibility.\allowbreak{}H2Bound}{20--36}
\leanref{Formalism.\allowbreak{}MainNotions.\allowbreak{}RamifiedCoefficientChains.\allowbreak{}Asm.\allowbreak{}h2\_\allowbreak{}bound\_\allowbreak{}of\_\allowbreak{}admitted}{Formalism.\allowbreak{}SharpRestriction.\allowbreak{}LevelLetterCocycles.\allowbreak{}FrobeniusLiftLemmasOfAdmitted}{141--175}
Lemma~\ref{B8.42buea}, followed by one level-one truncation, gives an integer \(k_c\) satisfying
\[
 k_c\le\frac{\beta_c+S+1}{q}+E_G+\sigma_2
                    +\beta_{c-2}+S-\beta_{c-1}+3.
\]
More precisely, for every
\(d\in\varphi^2\mathcal C_{c-3}(T_2)\cap\operatorname{Fil}^{\beta_{c-1}}_{\rm cris}\)
it gives \(v_d\in\mathcal V_c(T_1)\) such that \(Q_0^BF_Dv_d\in N(T_2)\). At every primitive root, the Laurent expansion of \(F_Dv_d\) has coefficients
\[
 b_d(k)=0\quad(k<B,\ k\ne0),\qquad
 b_d(0)=\ell^{k_c}\iota\varphi d,
 \qquad b_d(k)\in\iota D_{{\rm cris},c}(T_2[1/\ell]).
\]
Indeed, lift \(d\) by Lemma~\ref{B8.42buea} at \(c-1\), then truncate its source lift with exponent
\(\lceil(B+1)/q+1\rceil\). Naturality of this truncation leaves an error in \(\pi_0Q_0^B\mathcal R_{c-1}^{(1)}\). After multiplication by \(Q_0^B\), that error starts in degree \(2B\). Division by \((tU)^B\) proves the displayed coefficient assertions.
\leanref{Formalism.\allowbreak{}MainNotions.\allowbreak{}N.\allowbreak{}Pv.\allowbreak{}LocalLengthBound.\allowbreak{}contentBoundLevel\_\allowbreak{}of}{Formalism.\allowbreak{}SharpRestriction.\allowbreak{}LevelLetterCocycles.\allowbreak{}QuotientImageDivisibility.\allowbreak{}SmulSeriesDeltaZero}{273--518}

Put \(e_{\rm vec}=\lfloor B_{\rm vec}\rfloor\). The identity
\[
 \frac{\beta_c+S+1}{q}-\frac{\beta_c}{\ell}
       =\frac{\beta_c}{\ell q}+\frac{S+1}{q}
\]
and \(\beta_c\le h\) show that
\(e_{\rm vec}\ge a_c+k_c-h_0\).
Multiply \(x_f\) by \(\ell^{e_{\rm vec}-a_c-k_c+h_0}\). The resulting relation is
\[
 \ell^{e_{\rm vec}}D_f(0)
       -\ell^{k_c-h_0}\iota\varphi^2\operatorname{cmp}_{T_2}(x_f)
       \in\mathcal F^{-\delta}.
\]
Define
\[
 K_3=\{x\in\mathcal N_{c-2}(T_2):
       \iota\varphi^2\operatorname{cmp}_{T_2}(x)\in\mathcal F^{-\delta}\}.
\]
The class \(x_f+K_3\) is uniquely determined by this relation and gives an \(O\)-linear map. Apply the determinant quotient estimate to
\[
 \mathcal P=\operatorname{span}_{\mathcal A^\Delta}\{\varphi z:z\in\mathcal N_{c-3}(T_2)\},\qquad
 X=\pi_0\mathcal N_{c-2}(T_2)+Q_0^{h_0}\mathcal P.
\]
Its at most \(m\) auxiliary vectors and its quotient of normalized length at most \(C_{\rm len}\), pulled back as in Lemma~\ref{B8.42bufa}, give one generator step and one length step. Apply that lemma with \(X+K_3\); adding \(K_3\) can only decrease the quotient length. On the resulting submodule, \(x_f\) may be represented modulo \(K_3\) by an element of \(X\). Specialization then gives
\(\operatorname{cmp}(x_f)=\ell^{h_0}\varphi\operatorname{cmp}(t_f)\)
modulo the same ambiguity, with \(t_f\in\mathcal N_{c-3}(T_2)\).

Impose the further condition
\(\varphi^2\operatorname{cmp}(t_f)\in\operatorname{Fil}^{\beta_{c-1}}_{\rm cris}\).
Put
\[
 A_3=\{t\in\mathcal N_{c-3}(T_2):\iota\varphi^3\operatorname{cmp}(t)\in\mathcal F^{-\delta}\},\qquad
 C_3=\{t\in\mathcal N_{c-3}(T_2):\varphi^2\operatorname{cmp}(t)\in\operatorname{Fil}^{\beta_{c-1}}_{\rm cris}\}.
\]
The class \(t_f+A_3\) is uniquely determined and linear. Take its preimage of \((C_3+A_3)/A_3\). The quotient \(\mathcal N_{c-3}(T_2)/C_3\) maps into
\(D_{{\rm cris},c-1}/\operatorname{Fil}^{\beta_{c-1}}_{\rm cris}\), whose dimension is at most \(m\). Specialization factors through the \(r\) constant coordinates of an invariant Wach basis, so this quotient is finitely generated over \(O\), and hence generated by at most \(m\) elements. This gives a second generator step. Denote the resulting coefficient module by \(U_3\), and choose the vectors
\(d_f=\varphi^2\operatorname{cmp}(t_f)\)
on an \(O\)-basis of \(U_3\), then extend linearly. They satisfy
\[
 \ell^{e_{\rm vec}}D_f(0)-\ell^{k_c}\iota\varphi d_f
       \in\mathcal F^{-\delta}.
\]

Choose the corresponding source vectors \(v_{d_f}\) on this basis and extend linearly to a map \(y_c:U_3\to\mathcal V_c(T_1)\). Define \(U'\) as the kernel of
\[
 f\longmapsto
 \ell^{e_{\rm vec}}D_f(0)-\ell^{k_c}\iota\varphi d_f
       \pmod{\mathcal F^0}.
\]
Its image lies in \(\mathcal F^{-\delta}/\mathcal F^0\), of dimension at most \(m\), so this is the third generator step. On \(U'\), every coefficient of
\(\ell^{e_{\rm vec}}w-F_Dy_c\) lies in the required filtration: the degree-zero assertion follows from the last kernel, nonzero degrees below \(B\) follow from \(b_{d_f}(k)=0\), and degrees at least \(B\) follow from the lower jump bound \(-B\). The expansion criterion, transported between primitive roots by \(\Delta\), gives the asserted integral residual. The complete proof is contained in \Lean{Formalism.\allowbreak{}MainNotions.\allowbreak{}N.\allowbreak{}Pv.\allowbreak{}LocalLengthBound.\allowbreak{}length\_\allowbreak{}bound\_\allowbreak{}of\_\allowbreak{}hypotheses}, using \Lean{Formalism.\allowbreak{}MainNotions.\allowbreak{}N.\allowbreak{}Pv.\allowbreak{}LocalLengthBound.\allowbreak{}pass3} and \Lean{Formalism.\allowbreak{}MainNotions.\allowbreak{}N.\allowbreak{}Pv.\allowbreak{}LocalLengthBound.\allowbreak{}pass4}.
\end{proof}

\subsection*{Proof of \ref{B8.42bufb}}\label{prf:B8.42bufb}\phantomsection
\begin{proof}
The height hypothesis puts \(Q_0^B\varphi z\) in the \(c\)-component of the invariant Wach lattice. The action of \(\gamma\) on this lattice is trivial modulo \(\pi_0\), so \(x\in\mathcal V_c(T_1)\). The same height bound for \(T_2\), together with naturality of Frobenius, shows that \(F_Dx\in\mathcal N_c(T_2)\). Weighted specialization gives
\[
 \operatorname{cmp}_{T_2}(F_Dx)=\ell^{t+B+k}\varphi^3d_3.
\]
The first three displayed identities give
\[
 \ell^{e+s_1+s_2}\operatorname{cmp}_{T_2}(r)
   =\ell^{e+B+h_0-b_0}\varphi^3d_3
   =\ell^{t+B+k}\varphi^3d_3.
\]
Thus the difference lies in the kernel of specialization on
\(\mathcal N_c(T_2)\), which is \(\pi_0\mathcal N_c(T_2)\). This is contained in \(\mathcal U^1(T_2)\).
\end{proof}

\subsection*{Proof of \ref{B8.42buf}}\label{prf:B8.42buf}\phantomsection
\begin{proof}[Proof (outline)]
For an index \(i\), let
\[
 \mathcal P_i=\operatorname{span}_{A^\Delta}
       \{\varphi z:z\in\mathcal N_{i-1}(T_2)\},\qquad
 X_i(k)=\pi_0\mathcal N_i(T_2)+Q_0^k\mathcal P_i,\qquad
 K_i=\ker(\operatorname{cmp}_{T_2}|_{\mathcal N_i(T_2)}).
\]
The height bound gives \(X_i(\beta_i+S)\subset\mathcal N_i(T_2)\), and specialization gives \(K_i\subset X_i(\beta_i+S)\). The local quotient estimate at \(i=c,c-2\) supplies a set \(S_i\subset\mathcal N_i(T_2)\), with \(|S_i|\le m\), for which
\[
 \lambda_O\bigl(\mathcal N_i(T_2)/(X_i(\beta_i+S)+OS_i)\bigr)
       \le C_{\rm len}.
\]
\leanref{Formalism.\allowbreak{}MainNotions.\allowbreak{}RamifiedCoefficientChains.\allowbreak{}S9.\allowbreak{}LemmaDN}{Formalism.\allowbreak{}SharpRestriction.\allowbreak{}LevelLetterCocycles.\allowbreak{}LatticeBoundsOfThreeParts}{492--503}
\leanref{Formalism.\allowbreak{}MainNotions.\allowbreak{}RamifiedCoefficientChains.\allowbreak{}S9.\allowbreak{}bigCellDetWeightedEnergyIntegral.\allowbreak{}dual\_\allowbreak{}lemma\_\allowbreak{}of\_\allowbreak{}admitted}{Formalism.\allowbreak{}SharpRestriction.\allowbreak{}LevelLetterCocycles.\allowbreak{}UniformizerClauseOfAdmitted}{50--119}

Define \(U_5\subset U\) as the inverse image of \(X_c(\beta_c+S)/K_c\) under \(f\mapsto r(f)\bmod K_c\). Lemma~\ref{B8.42bufa} constructs this restriction with one length step and one generator step. Write \(r(f)\), modulo \(\pi_0\mathcal N_c(T_2)\), as a sum of terms
\(Q_0^{\beta_c+S}a\varphi z\), with \(a\in A^\Delta\) and
\(z\in\mathcal N_{c-1}(T_2)\). Replace each \(a\) by its constant coefficient when applying specialization. This constructs
\(d_1(f)\in\mathcal C_{c-1}(T_2)\) such that
\[
 \operatorname{cmp}_{T_2}(r(f))
       =\ell^{\beta_c+S}\varphi d_1(f).
\]
The upper height inclusion at \(c-1\), followed by the same specialization, constructs
\(d_2(f)\in\mathcal C_{c-2}(T_2)\) with
\[
 \ell^{\beta_{c-1}+S}d_1(f)=\varphi d_2(f).
\]
Both maps are \(O\)-linear because Frobenius is injective. If
\(y_2(f)\in\mathcal N_{c-2}(T_2)\) represents \(d_2(f)\), its class modulo \(K_{c-2}\) is uniquely determined. Define \(U_6\subset U_5\) as the inverse image of \(X_{c-2}(\beta_{c-2}+S)/K_{c-2}\) under this linear map. Lemma~\ref{B8.42bufa} gives a second length step and a second generator step. On \(U_6\), the same constant-coefficient construction gives an \(O\)-linear map
\(d_3(f)\in\mathcal C_{c-3}(T_2)\) with
\[
 d_2(f)=\ell^{\beta_{c-2}+S}\varphi d_3(f).
\]
\leanref{Formalism.\allowbreak{}MainNotions.\allowbreak{}N.\allowbreak{}Pv.\allowbreak{}LocalLengthBound.\allowbreak{}cmp\_\allowbreak{}of\_\allowbreak{}mem\_\allowbreak{}X}{Formalism.\allowbreak{}SharpRestriction.\allowbreak{}LevelLetterCocycles.\allowbreak{}QuotientImageDivisibility.\allowbreak{}CmpOfMemX}{116--172}
\leanref{Formalism.\allowbreak{}MainNotions.\allowbreak{}N.\allowbreak{}Pv.\allowbreak{}LocalLengthBound.\allowbreak{}cmp\_\allowbreak{}up}{Formalism.\allowbreak{}SharpRestriction.\allowbreak{}LevelLetterCocycles.\allowbreak{}QuotientImageDivisibility.\allowbreak{}CmpOfMemX}{173--226}

Define
\[
 U''=\ker\left(U_6\longrightarrow
   D_{\rm cris}(T_2[1/\ell])/\operatorname{Fil}^{\beta_{c-1}}_{\rm cris},
       \quad f\longmapsto\varphi^2d_3(f)\right).
\]
Its image lies in the \((c-1)\)-component modulo the indicated filtration, whose \(L'\)-dimension is at most \(m\). Since \(U_6\) is finitely generated over the discrete valuation ring \(O\), this image is generated by at most \(m\) elements. Passing to \(U''\) gives the third generator step. Each of the two length steps has normalized length at most \(C_{\rm len}\).

Let \(E_{c-1}\) be the integer in the crystalline interpolation estimate of Lemma~\ref{B8.42buea}, and set
\[
 k=\lceil E_{c-1}+\sigma_2\rceil+\beta_{c-2}+S-\beta_{c-1},
 \qquad
 t=\lceil E_G+\sigma_2\rceil-\lceil E_{c-1}+\sigma_2\rceil\ge0.
\]
For \(f\in U''\), choose \(z_f\in\mathcal N_{c-1}(T_1)\) whose image under \(F_D\) is integral and satisfies
\[
 \operatorname{cmp}_{T_2}(F_Dz_f)=\ell^k\varphi^2d_3(f).
\]
Define the correction explicitly by
\[
 x_f=\ell^tQ_0^{\beta_c+S}\varphi z_f.
\]
The identity
\[
 e_{\rm int}+\beta_{c-2}+S=t+k+\beta_{c-1}+S
\]
allows Lemma~\ref{B8.42bufb} to be applied with \(s_1=s_2=0\). It gives
\(x_f\in\mathcal V_c(T_1)\) and
\(\ell^{e_{\rm int}}r(f)-F_Dx_f\in\mathcal U^1(T_2)\).
Finally,
\(e_{\rm int}<S+E_G+\sigma_2+1\le B_{\rm int}\), since \(\lambda\ge1\).
\leanref{Formalism.\allowbreak{}MainNotions.\allowbreak{}N.\allowbreak{}Pv.\allowbreak{}LocalLengthBound.\allowbreak{}frobenius\_\allowbreak{}sq\_\allowbreak{}bound\_\allowbreak{}at}{Formalism.\allowbreak{}SharpRestriction.\allowbreak{}LevelLetterCocycles.\allowbreak{}QuotientImageDivisibility.\allowbreak{}FrobeniusSqComparisonLetter}{54--245}

For \(w=w_c|_V\), the two vector corrections give
\[
 \ell^{e_{\rm vec}+e_{\rm int}}w_c(f)-F_Du_f
       =\ell^{e_{\rm int}}r_c(f)-F_Dx_f\in\mathcal U^1(T_2),
\]
and since \(u_f\in\mathcal V_c(T_1)\), the lifting property of Lemma~\ref{B8.42bud} supplies \(Y_c(f)\in\mathcal H(T_1)\) with finite class and places both coordinates of
\(\ell^{e_{\rm vec}+e_{\rm int}}X_c(f)-F_{\mathcal H}Y_c(f)\)
in \(\mathcal U^M(T_2)\). On this module \(\varphi-1\) is bijective. If \((z,v)\) is a Herr cocycle with both coordinates in \(\mathcal U^M(T_2)\), choose \(t\in\mathcal U^M(T_2)\) with \((\varphi-1)t=z\). Then
\[
 (\varphi-1)((\gamma-1)t-v)=0.
\]
The module \(\mathcal U^M(T_2)\) is \(\gamma\)-stable, so injectivity of \(\varphi-1\) gives \((\gamma-1)t=v\). Hence this Herr cocycle is a boundary. Naturality now gives
\[
 \ell^{e_{\rm vec}+e_{\rm int}}[\kappa_c(f)]
     =F_*\operatorname{cl}_{T_1}(Y_c(f)),
\]
which is equivalent to the asserted membership in \(I_F\). The complete constructions are contained in \Lean{Formalism.\allowbreak{}MainNotions.\allowbreak{}N.\allowbreak{}Pv.\allowbreak{}LocalLengthBound.\allowbreak{}frobenius\_\allowbreak{}sq\_\allowbreak{}bound\_\allowbreak{}of\_\allowbreak{}local} and \Lean{Formalism.\allowbreak{}MainNotions.\allowbreak{}N.\allowbreak{}Pv.\allowbreak{}LocalLengthBound.\allowbreak{}rhoLift\_\allowbreak{}of}.
\end{proof}

\subsection*{Proof of \ref{B8.42bu1a}}\label{prf:B8.42bu1a}\phantomsection
\begin{proof}
Choose Smith normal form for \(P\subset M\). It gives a basis
\((c_i)_{i\in I}\) of \(M\), a basis \((p_j)_{j\in J_0}\) of \(P\), an injection \(\alpha:J_0\to I\), and scalars \(a_j\in O\) with
\[
 p_j=a_jc_{\alpha(j)}.
\]
Each \(a_j\) is nonzero, since \(p_j\ne0\). Saturation gives
\(c_{\alpha(j)}\in P\). Conversely, expansion in the basis \((p_j)\) expresses every element of \(P\) as an \(O\)-linear combination of these
\(c_{\alpha(j)}\). Take \(J=\alpha(J_0)\). Both bases are bases of the same \(O\)-module, so their transition matrices are inverse matrices over \(O\). Their determinants are units. This gives the valuation assertions.
\end{proof}

\subsection*{Proof of \ref{B8.42bu1b}}\label{prf:B8.42bu1b}\phantomsection
\begin{proof}
Regard \(P\) as a submodule of \(Q\). It is saturated there: if \(a\ne0\) and \(ax\in P\) with \(x\in Q\), saturation in \(M\) gives \(x\in P\).
Apply Lemma~\ref{B8.42bu1a} to the basis \((b_i)_{i\in S}\) of \(Q\).
Obtain a basis \((c_i)_{i\in S}\) and \(J\subset S\) with
\(P=\bigoplus_{i\in J}Oc_i\). Define
\[
 b'_i=\begin{cases}c_i,&i\in S,\\ b_i,&i\notin S.\end{cases}
\]
The direct sum decomposition
\(M=Q\oplus\bigoplus_{i\notin S}Ob_i\) proves that \((b'_i)\) is a basis and proves the first two assertions. Every basis vector indexed by \(J\) lies in \(P\); every basis vector that lies in \(P\) belongs to its span. Hence replacing \(J\) by \(\{i:b'_i\in P\}\) leaves the span unchanged. In this direct sum decomposition the transition matrix has blocks equal to the transition matrix on \(Q\) and the identity on its complement. Lemma~\ref{B8.42bu1a} gives the asserted integrality and determinant valuation.
\end{proof}

\subsection*{Proof of \ref{B8.42bu1c}}\label{prf:B8.42bu1c}\phantomsection
\begin{proof}
The vectors \(w_k=\iota\varphi\operatorname{cmp}_{T_2}(e_k)\) form a basis of the image of the \(i\)-component of \(D_{\rm cris}\). To see independence, clear denominators in a putative linear relation. A constant \(O\)-linear combination of the \(e_k\) whose specialization is zero lies in \(\pi_0\mathcal N_{i-1}(T_2)\). Uniqueness of its power-series coordinates makes all its constant coefficients zero. Frobenius is injective. The height inclusions give surjectivity onto the next crystalline component, so these vectors span that component as well.

Let \(\Psi:(L')^r\to\operatorname{span}_{L'}\{w_k\}\) send the standard basis to the \(w_k\), and put
\[
 G^u=\Psi^{-1}(\operatorname{Fil}^uD_{\rm dR}),\qquad
 F^u=O^r\cap G^u.
\]
Each \(F^u\) is saturated in \(O^r\), because \(G^u\) is an \(L'\)-subspace. Apply Lemma~\ref{B8.42bu1b} successively to the descending chain
\(F^{-\delta}\supset\cdots\supset F^{\delta+1}=0\).
This gives an \(O\)-basis \(b_1,\ldots,b_r\) simultaneously adapted to the chain. Let \(S_u=\{j:b_j\in F^u\}\). Then
\[
 |S_u|=\dim_{L'}G^u,\qquad S_{u+1}\subset S_u,
 \qquad S_{\delta+1}=\varnothing.
\]
Order the indices so that \(S_{-\delta}\) comes first, put \(m'=|S_{-\delta}|\), and for these indices define
\(u_j=\max\{u\in[-\delta,\delta]:j\in S_u\}\).
The central-jump hypothesis gives \(m'\le m\). Use the coordinate rows \(b_j\) as the matrix \(B\). Then \(\Psi(b_j)\in\operatorname{Fil}^{u_j}\), proving the asserted filtration membership of \(e'_j\).

The number of indices with \(u_j=u\) is
\[
 |S_u|-|S_{u+1}|=\dim G^u-\dim G^{u+1}
       =\operatorname{mult}_{\tau_i}(u).
\]
Multiplying by \(h_i+u\) and summing proves the displayed bound, with equality for this construction. The inverse coordinate matrix belongs to \(\mathrm M_r(O)\) because the \(b_j\) form an \(O\)-basis. Thus \(B\in\GL_r(O)\), and extending this constant basis change to \(O[[\pi_0]]\) preserves the Wach basis and has valuation cost zero.
\end{proof}

\subsection*{Proof of \ref{B8.42bu1}}\label{prf:B8.42bu1}\phantomsection
\begin{proof}[Proof (outline)]
For all data of Definition~\ref{B8.42bu}, put \(\Delta=\Gal(K(\mu_\ell)/K)\), \(\Gamma=\Gal(K(\mu_{\ell^\infty})/K(\mu_\ell))\), and let \(\pi_0\in O[[Z]]^\Delta\) be the \(\Delta\)-invariant Wach parameter over \(K\) with \(O[[Z]]^\Delta=O[[\pi_0]]\) and \(Q\) a unit times \(\pi_0+\ell\) (so \(\pi_0(\zeta_\ell-1)=-\ell\)). From \(\mathscr H_2\), assemble the foundational packages in Steps 1--5 and apply the local uniformizer-lifting and Frobenius-chain steps in Steps 6--9.

\emph{Step 1: Local admissibility and rank bounds.}
From \(\mathscr H_2\) (Bloch--Kato local duality and dimension, local Tate duality, and the local Euler characteristic formula), deduce the local admissibility package with constants \((C_R,C_A,C_B,C_K)=(10,23,32,32)\). This package provides the Wach-module structural clauses, the logarithmic size and restriction-dimension bounds, the level-one lifting and crystalline-lift bounds, the Bloch--Kato dimension formula and crystalline extension criterion, the Herr description and naturality of \(H^1\), the divided difference \((\gamma-1)/\pi_0\) for \(\gamma\in\Gamma\) on \(\Delta\)-invariant Wach components, and the generator bound for \(H^1_f\).
\leanref{GlobalTheory.\allowbreak{}T.\allowbreak{}localLifting\_\allowbreak{}theoremC\_\allowbreak{}hypotheses\_\allowbreak{}admitted}{Development.\allowbreak{}AbelianVarietyTheory.\allowbreak{}Proofs.\allowbreak{}InertiaEigenvalues.\allowbreak{}LocalLiftingTheoremCHypothesesAdmittedStatement}{8--10}

\emph{Step 2: Exactness of Hodge--Tate jumps.}
In any short exact sequence \(0\to T_0\to T_1\to T_2\to0\) of crystalline \(O\)-lattices over \(K\), the jumps of \(T_1\) at each embedding \(\tau:K\to L'\) form the union of the jumps of \(T_0\) and \(T_2\).
\leanref{GlobalTheory.\allowbreak{}jumps\_\allowbreak{}exact\_\allowbreak{}adm}{Development.\allowbreak{}AbelianVarietyTheory.\allowbreak{}Proofs.\allowbreak{}InertiaEigenvalues.\allowbreak{}JumpsExactAdm}{9--14}

\emph{Step 3: Wach-lattice properties and the letter operator condition.}
Establish the structural clauses for level-one Wach lattices of crystalline representations: componentwise freeness, triviality of \(\gamma\) modulo \(Z\), comparison with \(D_{\rm cris}\), low-level expansions, height bounds from Hodge jumps, saturated exactness, Frobenius-pole bounds on Herr cocycles, tame descent, uniqueness characterization, cokernel annihilation by \(Z^{H_*}\) (as in Lemma~\ref{B8.41c}), and the Nygaard extension comparison.
\leanref{GlobalTheory.\allowbreak{}T.\allowbreak{}wachLattice\_\allowbreak{}laterClauses\_\allowbreak{}hold}{Development.\allowbreak{}AbelianVarietyTheory.\allowbreak{}Proofs.\allowbreak{}InertiaEigenvalues.\allowbreak{}WachLatticeLaterClausesHoldStatement}{8--10}
From these clauses, deduce the letter operator condition over unramified \(K\): for any crystalline lattice \(T\), component \(i\in\Z/d'\Z\), integer \(y\ge1\), element \(\gamma\in\Gamma\), operator \(D\) on the \(i\)-th \(\Delta\)-invariant Wach component lattice \(N_i(T)^\Delta\) (and its localization) satisfying \(\pi_0 D=\gamma-1\), scalars \(c_x\in O\) with \(-\ell c_x=\chi_\ell(\gamma)^{-x}-1\), and finite set \(X\subset\Z\) containing \(a+j\) for every jump \(a\) of \(T\) at \(\tau_i\) and \(1\le j\le y\), if \(Q^yz\in N_i(T)^\Delta\) then \(\prod_{x\in X}(D-c_x)(z)\in N_i(T)^\Delta\).
\leanref{GlobalTheory.\allowbreak{}T.\allowbreak{}letterConditionGlobal\_\allowbreak{}of\_\allowbreak{}N1used}{Development.\allowbreak{}AbelianVarietyTheory.\allowbreak{}Proofs.\allowbreak{}InertiaEigenvalues.\allowbreak{}EigenvalueProductMapsLetterLatticeOfConclusionsAStatement}{8--10}

\emph{Step 4: Herr-complex naturality and extension classes.}
From Fontaine's decompletion and the Herr comparison at level \(1\), establish naturality of the Herr class map and the extension-class formula up to a unit: given a crystalline extension \(0\to T\to T'\to O\to0\) realizing a cocycle \(c\), a \((\varphi,\Gamma,\Delta)\)-invariant \(O\)-generator \(e\) of the unit étale module, and \(a\in O\), lifting \(a\cdot e\) to a \(\Delta\)-invariant element \(x\) of the étale module of \(T'\) produces a Herr cocycle \(((\varphi-1)x,(\gamma-1)x)\) in the étale module of \(T\) whose class equals \(ua[c]\) for a unit \(u\in O^\times\).
\leanref{GlobalTheory.\allowbreak{}AdmHerrLocal.\allowbreak{}main}{Development.\allowbreak{}AbelianVarietyTheory.\allowbreak{}Proofs.\allowbreak{}InertiaEigenvalues.\allowbreak{}ImportsBisRest}{9--14}

\emph{Step 5: Fontaine's inequality and the \((\varphi,\Gamma)\)-module core.}
Establish Fontaine's period-dimension inequality from the crystalline period comparison.
\leanref{GlobalTheory.\allowbreak{}FontaineIneq.\allowbreak{}one\_\allowbreak{}place\_\allowbreak{}admissible\_\allowbreak{}of\_\allowbreak{}gap\_\allowbreak{}constant\_\allowbreak{}le}{Development.\allowbreak{}AbelianVarietyTheory.\allowbreak{}Proofs.\allowbreak{}InertiaEigenvalues.\allowbreak{}OnePlaceAdmissibleOfGapConstantLe}{10--13}
Over unramified \(K\) at level \(1\), assemble the \((\varphi,\Gamma)\)-module core: full faithfulness and exactness of the étale \((\varphi,\Gamma)\)-module functor on crystalline lattices, freeness of the coefficient ring over its Frobenius image on \(\{(1+Z)^j:0\le j<\ell\}\), and existence of the \(\Gamma\)-equivariant left-inverse \(\psi\)-operator on both the ring and the étale modules.
\leanref{GlobalTheory.\allowbreak{}T.\allowbreak{}fontaine\_\allowbreak{}fullFaithful\_\allowbreak{}exact\_\allowbreak{}psi}{Development.\allowbreak{}AbelianVarietyTheory.\allowbreak{}Proofs.\allowbreak{}InertiaEigenvalues.\allowbreak{}FontaineFullFaithfulExactPsiStatement}{8--10}

\emph{Step 6: Invariant components and coefficient modules.}
Use the modules and constants of Definition~\ref{B8.42buc}. Since \(|\Delta|=q\) is a unit in \(O\), the averaging operator
\[
 e_\Delta=\frac1q\sum_{\varepsilon\in\Delta}\varepsilon
\]
is an idempotent with image the \(\Delta\)-invariants. It commutes with \(F_D\) and Frobenius. At unramified level one, \(\mathcal A^\Delta=\prod_iO[[\pi_0]]\), each \(\mathcal N_i(T_j)\) is free over \(O[[\pi_0]]\), and the cyclotomic polynomial is a unit times \(Q_0=\pi_0+\ell\). The chosen jump windows give \(M=\lfloor h/q\rfloor+1\), \(a=\lceil A\rceil\), and \(\widehat\sigma=\lceil\sigma\rceil\). These are the integer parameters in the following three constructions. Lemma~\ref{B8.42bu1c} gives the basis adapted to the central jumps by a change in \(\GL_r(O)\); this change has determinant valuation zero. The recursive source lifts used in both Frobenius corrections are those of Lemma~\ref{B8.42buea}.

\emph{Step 7: Reduction of finite cocycles.}
By Lemma~\ref{B8.42bud}, for every finite free \(O\)-module \(P\) and map \(g:P\to Z_f(T_2)\), there is a coefficient submodule \(P_2\) reached by at most \(n\) generator steps. There is an integer
\[
 s_{\rm red}\le\frac{2a+2\delta+S}{q}+\lambda+2\sigma+2
\]
such that \(\ell^{s_{\rm red}}g(f)\) is, modulo \(B(T_2)\), a signed sum of the cocycles \(\kappa_c(f)\), \(c\in J_-\). Their chosen coordinate maps satisfy
\[
 w_c(P_2)\subset\mathcal V_c(T_2),\qquad
 Q_0^\delta w_c(P_2)\subset\mathcal R_{c-1}^{(1)}(T_2).
\]
The corresponding Herr cocycles have second coordinate in \(\mathcal U^M(T_2)\).

\emph{Step 8: The two Frobenius corrections.}
For each \(c\in J_-\), apply Lemma~\ref{B8.42bue} to the restriction of \(w_c\) to the current coefficient module. It gives at most three generator steps and one length step, and a source map \(y_c\) whose residual
\[
 r_c=\ell^{e_{\rm vec}}w_c-F_Dy_c
\]
lies in \(\mathcal N_c(T_2)\), where
\[
 e_{\rm vec}\le\frac{h}{\ell q}+\frac{S+1}{q}+S+E_G+2\sigma_2+6+\lambda.
\]
Apply Lemma~\ref{B8.42buf} to this integral residual. This gives at most three further generator steps and two further length steps, at exponent
\[
 e_{\rm int}=S+\lceil E_G+\sigma_2\rceil
       \le S+E_G+\sigma_2+\lambda.
\]
By the second part of that lemma,
\(\ell^{e_{\rm vec}+e_{\rm int}}\kappa_c(f)\in I_F\)
on the resulting coefficient module. Every length step in both constructions has normalized length at most \(C_{\rm len}\).

\emph{Step 9: Assembly of the chain and its exponent.}
Let \(V_2=T_2[1/\ell]\), let \(\varpi\) be a uniformizer of \(O\), let \(k_O=O/\varpi O\) be the residue field, and set
\[
 a_0=\dim_{L'}H^0(K,V_2),\qquad
 t_0=\dim_{L'}\bigl(D_{\rm dR}(V_2)/\operatorname{Fil}^0\bigr),\qquad
 b_0=\dim_{k_O}H^0(K,T_2/\varpi T_2).
\]
The local finiteness and Bloch--Kato dimension statements in Step 1 give
\(\operatorname{rank}_OH_f^1(K,T_2)=a_0+t_0\).
The cohomology sequence for multiplication by \(\varpi\) gives
\(\dim_{k_O}H^1(K,T_2)[\varpi]=b_0-a_0\).
Every torsion class lies in \(H_f^1(K,T_2)\), so this module has at most
\(t_0+b_0\le nr+r\) generators. Choose finite cocycles representing these generators and obtain a finite free presentation
\(\jmath:F_0\to Z_f(T_2)\), surjective modulo \(B(T_2)\), with
\(\operatorname{rank}_OF_0\le(n+1)r\). Apply Step 7 with \(P=F_0\), and perform the two corrections in Step 8 successively for the \(n/2\) indices in \(J_-\), restricting all previously chosen maps at each stage. The final submodule \(F_k\) satisfies
\[
 \ell^e\jmath(F_k)\subset I_F,
 \qquad e=s_{\rm red}+e_{\rm vec}+e_{\rm int}.
\]
Indeed, the same exponent \(e_{\rm vec}+e_{\rm int}\) lifts each term in the signed sum of Step 7. The chain has at most
\[
 n+(3+3)\frac n2=4n\quad\text{generator steps},\qquad
 (1+2)\frac n2=\frac{3n}{2}\quad\text{length steps}.
\]
The quotient by \(F_{i+1}+\jmath^{-1}(I_F)\), intersected with \(F_i\), is a quotient of \(F_i/F_{i+1}\), so each length-step bound implies the corresponding bound in (R).

To match the real witnesses in (R), put
\[
 a_1=a,\qquad
 s_1=\frac{2a+2\delta+S}{q}+2\lambda,\qquad
 E_1=E_G+2\lambda=\frac{3a+S+2\delta}{q}+3\lambda,\qquad
 e_2=e_{\rm vec}.
\]
Then \(a_1\le A+1\). Since \(q\ge10\), \(S\ge0\), and \(\lambda\ge1\),
\[
 S-\frac{S+1}{q}+2\lambda-1\ge0,
\]
so the bound for \(e_{\rm vec}\) gives
\[
 e_2\le\frac{h}{\ell q}+E_1+2\sigma_2+2S+5+\lambda.
\]
Also \(2\lambda-1\ge0\), and therefore
\[
 s_{\rm red}+e_{\rm int}
    \le s_1+2\sigma+S+E_1+\sigma_2+1.
\]
Adding \(e_2\) proves the required estimate for \(e\), and the constructed chain has all the properties of (R).
\leanref{Formalism.\allowbreak{}MainNotions.\allowbreak{}RamifiedCoefficientChains.\allowbreak{}Asm.\allowbreak{}theorem\_\allowbreak{}of\_\allowbreak{}local\_\allowbreak{}lemmas}{Formalism.\allowbreak{}SharpRestriction.\allowbreak{}LevelLetterCocycles.\allowbreak{}FrobeniusLiftLemmasOfAdmitted}{272--291}
\leanref{Formalism.\allowbreak{}MainNotions.\allowbreak{}RamifiedCoefficientChains.\allowbreak{}S9.\allowbreak{}lattice\_\allowbreak{}bound\_\allowbreak{}three\_\allowbreak{}six\_\allowbreak{}of}{Formalism.\allowbreak{}SharpRestriction.\allowbreak{}LevelLetterCocycles.\allowbreak{}LatticeBoundsOfThreeParts}{414--462}
\leanref{Formalism.\allowbreak{}MainNotions.\allowbreak{}RamifiedCoefficientChains.\allowbreak{}stmt\_\allowbreak{}of\_\allowbreak{}stepsAdm}{Formalism.\allowbreak{}SharpRestriction.\allowbreak{}LevelLetterCocycles.\allowbreak{}LatticeFiniteCohomologyGeneratorBound}{327--332}

The complete proof is contained in \Lean{GlobalTheory.\allowbreak{}ramifiedCoefficientChains\_\allowbreak{}prv\_\allowbreak{}twenty\_\allowbreak{}five} via \Lean{Formalism.\allowbreak{}MainNotions.\allowbreak{}RamifiedCoefficientChains.\allowbreak{}Asm.\allowbreak{}theorem\_\allowbreak{}of\_\allowbreak{}literature\_\allowbreak{}hypotheses}.
\end{proof}

\subsection*{Proof of \ref{B8.42bv}}\label{prf:B8.42bv}\phantomsection
\begin{proof}[Proof (outline)]
Let \(\mathscr C\subset\Z/d_u\Z\) be the cut set, and let \(v(c)\) be the backward interval length at \(c\in\mathscr C\). The cut configuration has
\[
 \mathscr C\ne\varnothing,\qquad
 J\le v(c)\le d_u,\qquad c-v(c)\in\mathscr C,
 \qquad c-s\notin\mathscr C\quad(1\le s<v(c)).
\]
Every \(c+k\), \(1\le k\le J\), is nonthick. Each backward interval
\(c,c-1,\ldots,c-v(c)+1\) contains at most two thick positions, and
\(v_{\max}:=\max_{c\in\mathscr C}v(c)\le3J+1\).
The cuts lie after the last thick position of each cluster and in the intervening nonthick runs. Put \(H_j(a)=M_j(a)+h_j(a)\), where the jumps of \(T_j\) at \(a\) lie in \([-M_j(a),h_j(a)]\). Strictness of the crystalline filtration under \(F\) gives the target enclosure from the source enclosure. Thus
\[
 H_j(a)\le2h,\qquad H_j(c+k)\le2g\quad
 (j=1,2,\ c\in\mathscr C,\ 1\le k\le J).
\]
\leanref{Formalism.\allowbreak{}MainNotions.\allowbreak{}LocalTheoryWindows.\allowbreak{}CutsOK}{Formalism.\allowbreak{}SharpRestriction.\allowbreak{}LevelLetterCocycles.\allowbreak{}LocalTheoryWindows}{49--55}
\leanref{Formalism.\allowbreak{}MainNotions.\allowbreak{}WachCoordinates.\allowbreak{}ThickLe2}{Formalism.\allowbreak{}SharpRestriction.\allowbreak{}LevelLetterCocycles.\allowbreak{}WachCoordinates}{21--23}
Put \(c_h=\lceil1+\log_\ell(2h+1)\rceil\). Applying Lemma~\ref{B8.42ba} at cyclotomic level one (and averaging over \(\Gal(K(\mu_\ell)/K)\) of unit order \(q_\ell\)) to \(F\) and to the unit projection of the crystalline extension realizing a target class bounds the cut remainder exponents by
\begin{equation}\label{eq:B8-cluster-remainder}
 \rho_c\le\frac{\ell}{q_\ell^2}\left((2g+c_h+3)\frac{q_\ell}{q_\ell-1}
 +\frac{2h}{q_\ell^J}\frac{q_\ell}{q_\ell-1}\right)+c_h .
\end{equation}
Let \(\rho_*\) be the maximum of \(\rho_c\) over cuts and both lattices, and set
\[
 M_*=\lfloor(h+2g_*)/q_\ell\rfloor+1,\qquad
 \sigma_*=\left\lceil(3g_*+M_*/\ell^{J-1})\ell/q_\ell^2+1\right\rceil,\qquad
 a_0=2\sigma_*+\rho_*+2g_*v_{\max}.
\]
Across a cut interval of length \(v\), the cyclotomic shift-killing polynomial for the shift set \(U\subset[-2h,2h]\) (lying in three intervals of length at most \(3g_*\)) has valuation at \(1\) equal to
\begin{equation}\label{eq:B8-killer-valuation}
 v+\sum_{i\in U\setminus\{0\}}(v+1+v_\ell(i)),
\end{equation}
bounded uniformly by \(\kappa_*\) where
\begin{equation}\label{eq:B8-cluster-kappa}
 \kappa_*\le v_{\max}+20g_*(1+v_{\max})+5\{4g_*/q_\ell+\log_\ell(2h)\}.
\end{equation}

Fix \(\kappa\in Z_f(T_2)\). Construct Frobenius-horizontal frames across each cut interval. Propagating endpoint coordinates to the preceding cut costs two factors \(\ell^{\sigma_*}\) and reduces \(\ell^{\rho_*+2\sigma_*}\kappa\) modulo coboundaries to cut vectors \(w_c\). The source and target are equipped with the same cut configuration and the enclosure heights just specified. The opposite-sign conditions used in the interval lifting have the following precise form. For \(0\le s_0<s_1<v(c)\), write \(b_0=s_{c-s_0}\) and \(b_1=s_{c-s_1}\) when these positions are thick. Then
\[
 \begin{aligned}
 &(\varsigma_{c-s_1},\varsigma_{c-s_0})=(-1,1)
       &&\Longrightarrow b_0<\ell b_1,\\
 &(\varsigma_{c-s_1},\varsigma_{c-s_0})=(1,-1)
       &&\Longrightarrow 2b_0\le\ell b_1\quad\text{and}\quad2a_0<b_1.
 \end{aligned}
\]
The cluster-size inequality with \(C_2\ge C_0\) supplies these comparisons. Apply the target shift killer \(P\), clear the window denominators by \(\ell^{\sigma_*}\), and lift the at most two downward positions at exponent \(2a_0\) by Lemma~\ref{B8.42bb}. Lift the upward positions in their filtered submodules using the displayed comparisons. The remaining interval-divisibility vector costs \(\ell^{\rho_*}\), after which the source shift killer \(P'\) gives a preimage in \(I_F\). The resulting integer exponent is
\[
 e_{\rm cl}=2\rho_*+4\sigma_*+2a_0+2\kappa_*,
 \qquad a_0=2\sigma_*+\rho_*+2g_*v_{\max}.
\]
\leanref{Formalism.\allowbreak{}MainNotions.\allowbreak{}WachCoordinates.\allowbreak{}WordOK}{Formalism.\allowbreak{}SharpRestriction.\allowbreak{}LevelLetterCocycles.\allowbreak{}WachCoordinates}{24--30}
\leanref{Formalism.\allowbreak{}MainNotions.\allowbreak{}WachCoordinates.\allowbreak{}initialExponentBound}{Formalism.\allowbreak{}SharpRestriction.\allowbreak{}LevelLetterCocycles.\allowbreak{}WachCoordinates}{65--68}
\leanref{Formalism.\allowbreak{}MainNotions.\allowbreak{}WachCoordinates.\allowbreak{}ePar}{Formalism.\allowbreak{}SharpRestriction.\allowbreak{}LevelLetterCocycles.\allowbreak{}WachCoordinates}{69--73}
\leanref{Formalism.\allowbreak{}MainNotions.\allowbreak{}WachCoordinates.\allowbreak{}ThmELiftAt}{Formalism.\allowbreak{}SharpRestriction.\allowbreak{}LevelLetterCocycles.\allowbreak{}WachCoordinates}{123--140}
\leanref{Formalism.\allowbreak{}MainNotions.\allowbreak{}WachCoordinates.\allowbreak{}ThmEBound}{Formalism.\allowbreak{}SharpRestriction.\allowbreak{}LevelLetterCocycles.\allowbreak{}WachCoordinates}{74--84}
Substituting \eqref{eq:B8-cluster-remainder}, \(\sigma_*\), and \eqref{eq:B8-cluster-kappa} gives
\begin{equation}\label{eq:B8-cluster-exp}
 2\rho_*+4\sigma_*+2a_0+2\kappa_*\le800\{h/q_\ell^{J+1}+g_*v_{\max}+\log_\ell h\}.
\end{equation}
Thus \(e_{\rm cl}=2\rho_*+4\sigma_*+2a_0+2\kappa_*\) satisfies \(\ell^{e_{\rm cl}}Z_f(T_2)\subset I_F\) and, since \(v_{\max}\le3J+1\), the asserted bound. The quotient \(Q=Z_f(T_2)/I_F\) is generated by at most
\(d=(d_u+1)r\) elements, since
\(Z_f(T_2)/B(T_2)\simeq H_f^1(K,T_2)\) has that generator bound and
\(B(T_2)\subset I_F\). The annihilation
\(\ell^{e_{\rm cl}}Q=0\) makes a chosen surjection \(O^d\to Q\)
factor through \((O/\ell^{e_{\rm cl}}O)^d\). Consequently
\[
 \lambda_O(Q)\le d e_{\rm cl}.
\]
Put \(R=h/q_\ell^{J+1}+(g+1)(J+1)+\log_\ell h\). Every summand of \(R\) is nonnegative, including the convention at \(h=0\). Since
\(g_*v_{\max}\le g_*(3J+1)\le3(g+1)(J+1)\), the exponent bound gives
\(\lambda_O(Q)\le2400(d_u+1)rR\). The threshold can be taken to be
\(C_0=\max(C',3200)\), where \(C'\ge1\) is the constant for the cut configuration. Thus every \(C_2\ge C_0\) gives \eqref{eq:B8-cluster-col}. Enlarging \(C_2\) strengthens the size hypothesis and increases the bound by a nonnegative multiple of \(R\).
\leanref{Formalism.\allowbreak{}MainNotions.\allowbreak{}WachCoordinates.\allowbreak{}corollary\_\allowbreak{}of\_\allowbreak{}Theorem\_\allowbreak{}E}{Formalism.\allowbreak{}SharpRestriction.\allowbreak{}LevelLetterCocycles.\allowbreak{}BoundedNumeratorPowerSummability}{722--730}
\leanref{Formalism.\allowbreak{}MainNotions.\allowbreak{}N.\allowbreak{}Ln03.\allowbreak{}corollary\_\allowbreak{}local\_\allowbreak{}exists\_\allowbreak{}constant}{Formalism.\allowbreak{}SharpRestriction.\allowbreak{}LevelLetterCocycles.\allowbreak{}LocalTheoremOfInputs}{112--133}


The complete proof is contained in \Lean{GlobalTheory.\allowbreak{}T.\allowbreak{}theoremE\_\allowbreak{}corollaryA\_\allowbreak{}closing\_\allowbreak{}all\_\allowbreak{}data}.
\end{proof}

\subsection*{Proof of \ref{B8.42bw}}\label{prf:B8.42bw}\phantomsection
\begin{proof}[Proof (outline)]
For (a), level-one unramified comparison combines two Wach lifts of cost \(\lfloor H'\ell/q_\ell^2+3+\log_\ell(H'+1)\rfloor\) with the shifted-weight kernel correction on \(\mathcal J\subset[-H',2H']\setminus\{0\}\) of valuation at most \(v_\ell(H'!)+v_\ell((2H')!)+\lfloor H'/q_\ell\rfloor+1+v_\ell((\lfloor H'/q_\ell\rfloor+1)!)\), which gives \eqref{eq:B8-unrestricted-exp}.

For (b), if a thick position exists and \(b_*/4<40U_0\), the unrestricted level-one lift with \(M=\max_a M_a\le h\) and \(H'\le h+\delta\) counts shifted valuations on the two signed ranges to lift all finite classes at an exponent
\begin{equation}\label{eq:B8-short-large-gap-exp}
 e_0\le\tfrac{71}{10}(h/\ell+1)+2\log_\ell(2h)\le160U_0,
\end{equation}
by \(b_*<160U_0\), \(h\le11b_*/10\), and \(\ell\ge11\). Otherwise \(40U_0\le b_*/4\) whenever a thick position occurs. Set \(M'_0=\lfloor2h/q_\ell\rfloor+1\) and \(\sigma_0=\lceil(2(w_*+S)+M'_0)\ell/q_\ell^2+1\rceil\); both \(\sigma_0\) and the cut remainder exponent \(r_0\) are at most \(3U_0\). Cut \(\Z/d_u\Z\) at each \(c\) with \(c+1\) nonthick; this partitions \(\Z/d_u\Z\) into intervals of lengths \(v\in\{1,2,3\}\) summing to \(d_u\). Across a cut interval of length \(v\in\{1,2,3\}\), take a Frobenius-horizontal basis at the preceding cut and split off the central jump directions in \([-\delta,\delta]\), which contribute at most \(4mv\) central classes at that cut (hence at most \(4md_u\) in total, each annihilated by \(\ell^{e'}\) from (a)). On the complementary coordinates, arc-series coordinate lifting lifts each remaining prescribed horizontal coordinate at an exponent \(e_1\) satisfying
\begin{equation}\label{eq:B8-short-coordinate-exp}
 \begin{array}{c|ccc}v&1&2&3\\\hline
 e_1&\le6U_0+1&\le15U_0+1+\delta I&\le34U_0+1+\delta I .
 \end{array}
\end{equation}
as follows: for \(v=1\) by direct filtered or constant-plus-remainder lifting; for \(v=2\) by subtracting the downward constant before transporting across the opposite-sign window, where the near-window comparison contributes \(\delta I\) when \eqref{eq:B8-local-band} holds; and for \(v=3\) (where two thick letters precede a nonthick letter) by comparing the middle or initial shift against the transported window threshold. For the union \(J_c\) of the at most three half-open shift intervals of width \(S\) on the cut interval, the cyclotomic shift-killing polynomial has valuation \(\kappa_c\) at \(1\) satisfying
\begin{equation}\label{eq:B8-short-kappa}
 \kappa_c\le13S+3\log_\ell(2h)+3 .
\end{equation}
Assembling the coordinate lifts at exponent \(\lfloor40U_0+\delta I\rfloor\) with \(r_0\), \(\max_c\kappa_c\), and \(3\sigma_0\) gives reduced exponent at most \(55U_0+13S+\delta I\), which together with \eqref{eq:B8-short-large-gap-exp} and the \((d_u+1)r\) generator bound yields \eqref{eq:B8-short-col}.

The complete proof is contained in \Lean{GlobalTheory.\allowbreak{}T.\allowbreak{}star\_\allowbreak{}ad\_\allowbreak{}and\_\allowbreak{}bb\_\allowbreak{}le\_\allowbreak{}of\_\allowbreak{}joint\_\allowbreak{}cond} and \Lean{GlobalTheory.\allowbreak{}T.\allowbreak{}mediumInert\_\allowbreak{}smallInert\_\allowbreak{}lengthBound\_\allowbreak{}tabX}.
\end{proof}

\subsection*{Proof of \ref{B8.42bxa}}\label{prf:B8.42bxa}\phantomsection
\begin{proof}
The parameter conditions give \(f,n,t,n_V,L\ge1\), \(m,r\le t\), and \(0\le\delta_{\rm pr}\le2t\). Since \(d_u\ge1\) and the sum of the local degrees over all places above \(\ell\) is \(2f\),
\[
 |\mathcal C|\le2f,\qquad
 \sum_{u\in\mathcal C}d_u\le2f,\qquad
 \sum_{w\in\mathcal R}(d_w+2)\le6f.
\]
The three summands of \(D_0\) are consequently bounded by
\[
 6B_c,\qquad 8n_Vt(2f+1)\le24B_c,\qquad
 2cf t^2\le2B_c,
\]
respectively. Their sum is at most \(32B_c\).

For the additional column terms, the degree contribution is at most
\(2fm\delta_{\rm pr}\le4B_c\), the constant projector contribution is at most \(2\delta_{\rm pr}|\mathcal C|\le8B_c\), and the two invariant contributions together are at most
\[
 |\mathcal C|\,2crn n_VL\le4cfrn n_VL\le4B_c.
\]
Thus \(D_{\rm sp}-D_0\le16B_c\). Finally \(48\le60\) and \(L\ge1\), so both bounds are at most \(60cfn t^2n_VL^2\). Taking \(c=120\) proves the last assertion.
\end{proof}

\subsection*{Proof of \ref{B8.42bx}}\label{prf:B8.42bx}\phantomsection
\begin{proof}[Proof (outline)]
\emph{The common local word.}
Take the maximum of the cutoffs for the long-orbit and remaining inert-word estimates. Their inputs are the local de Rham jump comparison, Fontaine comparison at the places above \(\ell\), the common embedding-sign condition, and the member scale inequalities. These inputs hold on the chosen twist-good cut family. The inert-prime cutoff gives \(\ell\ge11\), and the exclusion of the ramified host primes makes \(E_u/\Q_\ell\) unramified.

For every good point, place and finite splitting coefficient extension, the common comparison gives one cyclic word. Its letter at \(a\) consists of \(s_a\), its sign, and the two enclosure heights \(M_a,h_a\). The same word works for every stable source lattice in a subrepresentation of a finite sum of twisted dual full representations and for every stable Tate-dual target lattice. The source and target are crystalline with jumps in the displayed signed windows. The comparison also bounds the sum of the central multiplicities over the embeddings above each component by \(m\). Since the field is unramified, precisely one embedding lies over each component, so this is the required bound at \(\tau_a\). The word-size bounds give \eqref{eq:B8-word-box}; at a nonthick letter they give \(|s_a-y|\le T_0\) with \(w''(n)\le y\le w_B\).
\leanref{GlobalTheory.\allowbreak{}HTW.\allowbreak{}RecS.\allowbreak{}WordBase}{Development.\allowbreak{}AbelianVarietyTheory.\allowbreak{}Definitions.\allowbreak{}GaloisGeometry.\allowbreak{}ThickLetterWordLatticeSurjectionHypotheses}{30--44}
\leanref{GlobalTheory.\allowbreak{}Survey.\allowbreak{}ClusteredLetterwiseBoxes.\allowbreak{}base\_\allowbreak{}of}{Development.\allowbreak{}AbelianVarietyTheory.\allowbreak{}Counterexamples.\allowbreak{}HeckeCharacterExamples.\allowbreak{}ClusteredLetterwiseBoxes}{100--115}

In the long-orbit case, apply the scale estimate with constant \(\max(100,2C_2)\). Its conclusions are \(J+1\le d_u\), the inequality \eqref{eq:B8-long-word-condition}, and the marked-cluster condition with separation \(J\). The progression-window inequality gives \(2J+2\le L_w\). Passing from marked clusters to letters gives the stated singletons or pairs separated by \(J\) nonthick positions.
\leanref{GlobalTheory.\allowbreak{}HTW.\allowbreak{}RecS.\allowbreak{}jointB\_\allowbreak{}of\_\allowbreak{}large}{Development.\allowbreak{}AbelianVarietyTheory.\allowbreak{}Proofs.\allowbreak{}SatakeParameters.\allowbreak{}InertLongOrbitWordConditionsOfLargeScale}{20--41}

For the remaining case use the scale estimate with constant \(100\). Inert Frobenius, and therefore its inverse, has sign \(-1\), so successive signs alternate. The word estimate gives either the nonthick alternative directly, with \eqref{eq:B8-short-word-condition} and no three consecutive thick letters, or an even alternating word with
\[
 s_a\ge\frac{M+S}{2},\qquad s_a\ge2A+\delta+1
 \quad\text{for all }a,\qquad M=\max_a s_a,
\]
together with the size conditions \(1\le b_*\), \(1\le w_*\), \(20w_*\le b_*\), \(20\delta\le b_*\), \(20S\le b_*\) and \(s_a>\delta\) for all \(a\). In the latter alternative, if a thick letter exists then \(M\ge b_*\). A nonthick letter would satisfy both \(s_a\ge b_*/2\) and \(s_a\le w_*\le b_*/20\), a contradiction since \(b_*\ge1\). Thus every letter is thick. If no thick letter exists, every letter is nonthick, and the short-word alternative holds. This proves the stated dichotomy. The crystalline comparison is compatible with each lattice surjection, so the corresponding local lifting inputs apply to the chosen \(F\).
\leanref{GlobalTheory.\allowbreak{}HTW.\allowbreak{}RecS.\allowbreak{}jointCD\_\allowbreak{}of\_\allowbreak{}large}{Development.\allowbreak{}AbelianVarietyTheory.\allowbreak{}Proofs.\allowbreak{}SatakeParameters.\allowbreak{}InertPrimeWordConditionsOfLargeScale}{21--48}
\leanref{GlobalTheory.\allowbreak{}Survey.\allowbreak{}ClusteredLetterwiseBoxes.\allowbreak{}frozenCD\_\allowbreak{}of}{Development.\allowbreak{}AbelianVarietyTheory.\allowbreak{}Counterexamples.\allowbreak{}HeckeCharacterExamples.\allowbreak{}ClusteredLetterwiseBoxes}{123--160}

\emph{The arithmetic passage.}
Let \(\mathcal R,\mathcal C,D_0,D_{\rm sp}\) have the meaning of Lemma~\ref{B8.42bxa}. The character-block and invariant estimates give
\[
 c_B\le120tn n_VL,\quad h_0\le120t,\quad
 c_0(u)\le120n n_VL,\quad h_W(u)\le120rn n_VL.
\]
The last two estimates are used on the split columns; at an inert index the column set is empty. Let \(W\) be the normalized rank-\(r\) lattice in \(\rho_i(\pi)\chi_{i,1}^{-1}\). For the last inequality, the invariant-lattice estimate kills
\(H^0(E_u,W\otimes(K_i/O_i))\) by \(\ell^{c_0(u)}\). This subgroup lies in the \(\ell^{c_0(u)}\)-torsion of a divisible module of corank \(r\), so its normalized length is at most \(rc_0(u)\). The block-character invariant estimate gives the stronger bound \(h_0\le2t\), which implies the displayed bound. The uniform estimates for \(c_B,c_0\) use the canonical parameter choice and the size and jump conditions of the cut.
\leanref{GlobalTheory.\allowbreak{}NSizes.\allowbreak{}main\_\allowbreak{}alt\_\allowbreak{}alt}{Development.\allowbreak{}AutomorphicArithmetic.\allowbreak{}ThetaCorrespondence.\allowbreak{}ThetaPairLifts.\allowbreak{}ThirdPlaceSeries.\allowbreak{}ImportsFurtherRest}{13--44}

Suppose first that \(\ell\) splits. In the split congruence inequality, the contribution of a column place \(u\) is
\[
 e_A+d_um(2e_A+\delta_{\rm pr})+2\delta_{\rm pr}
     +h_W(u)+rc_0(u)
 =L_{i,u}+d_um\delta_{\rm pr}+2\delta_{\rm pr}
     +h_W(u)+rc_0(u).
\]
The row and column sets cover the places above \(\ell\), and the allowance is zero outside the columns. Hence the sum of the first terms is exactly \(\sum_{u\mid\ell}L_{i,u}\); all remaining row and column terms are \(D_{\rm sp}\). The split congruence comparison and global reciprocity therefore give
\[
 v_{cc}(i,\pi)\le s_\pi+
 \sum_{1\le e<t}s_{\ell,e}(E)+v_\ell(h_E)
 +\sum_{u\mid\ell}L_{i,u}+D_{\rm sp}
 +t\sum_{v\in Q_i}\lambda_i(O_i/(\zeta_{i,v}-1)).
\]

If \(\ell\) is inert, the column set is empty. The row set is then the whole set of places above \(\ell\), and the metric-integrality and determinant comparison has local contribution \(\sum_{u\mid\ell}L_{i,u}\) and remaining contribution \(D_0\). It gives the same inequality with \(D_{\rm sp}\) replaced by \(D_0\). The global character terms are finite by the assumed finiteness of the odd \(\Sha\)-groups. In both cases \eqref{eq:B8-killing} bounds the last sum by \(\ell^{-2}\), and Lemma~\ref{B8.42bxa} bounds the remaining contribution by \(7200fn t^2n_VL^2\). Definition~\eqref{eq:B8-Bd} gives the required \(\mathcal B_{i,\pi}\) assertion.
\leanref{GlobalTheory.\allowbreak{}NLink.\allowbreak{}starAD\_\allowbreak{}of\_\allowbreak{}global}{Development.\allowbreak{}AbelianVarietyTheory.\allowbreak{}Proofs.\allowbreak{}InertiaEigenvalues.\allowbreak{}StarInequalityRowsADOfGlobalTheorems}{8--110}

The complete proof is contained in \Lean{GlobalTheory.\allowbreak{}HTW.\allowbreak{}RecS.\allowbreak{}fam\_\allowbreak{}jointBCD\_\allowbreak{}closed\_\allowbreak{}alt} and \Lean{GlobalTheory.\allowbreak{}NLink.\allowbreak{}rowsAD\_\allowbreak{}in\_\allowbreak{}setting\_\allowbreak{}alt}, using \Lean{GlobalTheory.\allowbreak{}NLink.\allowbreak{}starAD\_\allowbreak{}of\_\allowbreak{}global} and \Lean{GlobalTheory.\allowbreak{}NLink.\allowbreak{}BBSize.\allowbreak{}main}.
\end{proof}

\subsection*{Proof of \ref{B8.42b}}\label{prf:B8.42b}\phantomsection
\begin{proof}
Choose \(C_2=\max(1,C_0)\), where \(C_0\) is the threshold of Lemma~\ref{B8.42bv}. Take the canonical hub realizing the good points and the twist of Lemma~\ref{B8.42a}, whose killing condition is \eqref{eq:B8-killing}. Work above the common family cutoff of Theorem~\ref{B8.41}, Lemma~\ref{B8.42bp}(a), and Lemma~\ref{B8.42bx}, as well as the parameter, exceptional-prime, and small-prime family cutoffs.
\leanref{GlobalTheory.\allowbreak{}T.\allowbreak{}theoremE\_\allowbreak{}corollaryA\_\allowbreak{}closing\_\allowbreak{}all\_\allowbreak{}data}{Development.\allowbreak{}AbelianVarietyTheory.\allowbreak{}Proofs.\allowbreak{}InertiaEigenvalues.\allowbreak{}TheoremECorollaryAClosingAllDataStatement}{8--10}
\leanref{GlobalTheory.\allowbreak{}T.\allowbreak{}exists\_\allowbreak{}hubData\_\allowbreak{}admJ\_\allowbreak{}setup}{Development.\allowbreak{}AutomorphicArithmetic.\allowbreak{}ThetaCorrespondence.\allowbreak{}ThetaPairLifts.\allowbreak{}ThirdPlaceSeries.\allowbreak{}ExistsSetupLargeStatement}{12--13}
\leanref{GlobalTheory.\allowbreak{}T.\allowbreak{}excludedPrimes\_\allowbreak{}few\_\allowbreak{}and\_\allowbreak{}small\_\allowbreak{}along\_\allowbreak{}family}{Development.\allowbreak{}AbelianVarietyTheory.\allowbreak{}Proofs.\allowbreak{}InertiaEigenvalues.\allowbreak{}ParametersExcludedPrimesFewAndSmallAlongFamilyStatement}{13--15}
\leanref{GlobalTheory.\allowbreak{}T.\allowbreak{}specialPrimes\_\allowbreak{}small\_\allowbreak{}along\_\allowbreak{}family}{Development.\allowbreak{}AbelianVarietyTheory.\allowbreak{}Proofs.\allowbreak{}InertiaEigenvalues.\allowbreak{}SpecialPrimesSmallAlongFamilyStatement}{13--15}
\leanref{GlobalTheory.\allowbreak{}T.\allowbreak{}admissible\_\allowbreak{}parameters\_\allowbreak{}family\_\allowbreak{}const}{Development.\allowbreak{}AbelianVarietyTheory.\allowbreak{}Proofs.\allowbreak{}InertiaEigenvalues.\allowbreak{}AdmissibleParametersFamilyConstStatement}{13--15}

Use the size width \(S_c\) and error width \(\eta=\eta_{\rm sz}\) of the member. Put \(n=f/2\) and \(n_V=r+t\). Take \(Z\), \(a_*\), the base sizes, the intervals \eqref{eq:B8-exception-intervals}, and the exceptional primes of Lemma~\ref{B8.42bp}, and define
\[
 Y=\frac{kL}{m\delta_2},
\]
where \(\delta_2=|T_2|/n>0\) is the density of thick host positions. For a sufficiently large family constant \(C>0\) chosen below (in particular \(C\ge\max(3C_2+1,22365,7200)\)), define \(A_\ell\) by the first applicable row of
\begin{equation}\label{eq:B8-prime-table}
\begin{array}{c|l}
 \ell&A_\ell\\ \hline
 g_0,g_1&Cx^C\eta\\
 \ell\mid Q_S\text{ or }\ell=3&Cfmb\log x+Cx^CL^2\\
 q,d&2Cfmb\\
 \ell>Z\text{ outside the exceptional set}&0\\
 \text{otherwise}&Cfn t^2n_VL^2+D_\ell .
\end{array}
\end{equation}
For the final row set \(P_\ell=C(1+H_B/\ell+\log_\ell H_B)\), \(J=\max(1,\lceil\log_{\ell-1}Y\rceil)\), and
\[
 \eta_3=\eta+t+2,\quad
 w_B=\max\{w''(n)+\eta,\max_T(b_T+b'/x+\eta)\},\quad
 w_P=w_B+(t+1)/2,\quad G=w_P+S_c+\eta_3.
\]
Let \(I_\ell=1\) if some \(w''(n)\le y\le w_B\) satisfies
\begin{equation}\label{eq:B8-global-band}
 |\ell y-b|\le C(\ell\eta_3+\ell S_c+b/\ell+y+\ell L),
\end{equation}
and \(I_\ell=0\) otherwise. Put \(F_\ell=(w_P+\eta_3)/\ell+S_c+L+\eta_3 I_\ell\). Let \(o_\ell\) be the order of \(\ell\) modulo \(2^{e+2}\) up to sign for the host tower exponent \(e\), and \(L_w\) the host progression-window length. Split and inert refer to \(\Q(\sqrt{-d})\). The place part is
\begin{equation}\label{eq:B8-prime-table-depth}
 D_\ell=\begin{cases}
 (2fm+f)P_\ell&\text{split},\\
 4fCr\{H_B/(\ell-1)^{J+1}+(G+1)(J+1)+L\}&\text{inert},\ \ell<Y,\ o_\ell>4L_w,\\
 4fCr(H_B/\ell^2+F_\ell)+2fCmP_\ell&\text{other inert}.
 \end{cases}
\end{equation}
\leanref{Formalism.\allowbreak{}MainNotions.\allowbreak{}TwistedPlaces.\allowbreak{}lenB}{Formalism.\allowbreak{}SharpRestriction.\allowbreak{}LevelLetterCocycles.\allowbreak{}MemPrimesOverOfMem}{249--254}
\leanref{Formalism.\allowbreak{}MainNotions.\allowbreak{}TwistedPlaces.\allowbreak{}lamB}{Formalism.\allowbreak{}SharpRestriction.\allowbreak{}LevelLetterCocycles.\allowbreak{}MemPrimesOverOfMem}{259--264}

Off the type primes, compose Hecke elements supported on \(Z(i)\) with \(\alpha_i\). Their support becomes good and their good evaluations are multiplied by \(M_0\). An evaluated ideal of valuation \(a\) costs at most \(ha\) on a lattice of rank \(h\) by elementary divisors. The eigenspace filtration yields
\begin{equation}\label{eq:B8-eta-passage}
 \lambda_i(U_i/L_i)\le\sum_{\pi\in Z_{\rm good}(i)}\mathrm{mult}(i,\pi)
 (v_{cc}(i,\pi)+v_\ell(M_0)).
\end{equation}
Combine \eqref{eq:B8-eta-passage} with \eqref{eq:B8-mult-rank}, the rank row (used as in the proof of Theorem~\ref{B8.42} to pass from \(\dim Z_j(\mu_w)\) to \(d_{j,0}\)), and the logarithmic place-weight sum of \ref{B4.52}. At \(\ell=2\) and at \(\ell=3\nmid Q_S\), Theorem~\ref{B8.41} gives \eqref{eq:B8-Bd23}; for odd \(\ell\mid Q_S\), the same split congruence bound \eqref{eq:B8-cost-growth} with \(s_0=s_\ell\) and \(\ell^{v_\ell(Q_S)}\le Q_S\le n^2\) gives the second row of \eqref{eq:B8-prime-table} once \(C\) exceeds the small-prime constants.
\leanref{GlobalTheory.\allowbreak{}EtaPassage.\allowbreak{}main}{Development.\allowbreak{}AbelianVarietyTheory.\allowbreak{}Proofs.\allowbreak{}InertiaEigenvalues.\allowbreak{}ImportsTerMore}{11--21}
\leanref{GlobalTheory.\allowbreak{}T.\allowbreak{}goodPointMultiplicityLeRank\_\allowbreak{}proof}{Development.\allowbreak{}AbelianVarietyTheory.\allowbreak{}Proofs.\allowbreak{}InertiaEigenvalues.\allowbreak{}GoodPointMultiplicityLeRankProofStatement}{13--15}
\leanref{UpperTheory.\allowbreak{}N.\allowbreak{}starSF\_\allowbreak{}of\_\allowbreak{}prints}{Development.\allowbreak{}AutomorphicArithmetic.\allowbreak{}ThetaCorrespondence.\allowbreak{}ThetaPairLifts.\allowbreak{}ThirdPlaceSeries.\allowbreak{}NStarS}{11--13}
\leanref{UpperTheory.\allowbreak{}N.\allowbreak{}rowC}{Development.\allowbreak{}AutomorphicArithmetic.\allowbreak{}ResidueCharacters.\allowbreak{}RowAtOfSetup}{43--46}

\paragraph{Large primes outside the exceptional set.}
For \(\ell>Z\) outside the exceptional set and outside \(\{g_0,g_1\}\), Lemma~\ref{B8.42bp}(a) shows that \(\ell\ge11\), \(t+1<\ell\), \(\ell\notin\{q,d\}\), \(\ell\nmid Q_S\), and every completion \(E_w\) (\(w\mid\ell\)) is unramified over \(\Q_\ell\), and constructs position data at every embedding. Over a finite extension \(L'/K_i\) containing all embeddings of \(E_w(\mu_\ell)\) with integers \(O\), Lemma~\ref{B8.42bp}(b) (using Proposition~\ref{B8.42bo1}) shows that the image of every cut cocycle at \(w\) is finite in the normalized rank-\(r\) block, a coboundary on \(G_{E_w}\) in each character block \(2\le c\le t\), and a coboundary on \(I_{E_w}\) (hence zero on \(I_{E_w}\), as the action is trivial) in the \(c=1\) block. Finite classes descend from \(L'\) to \(K_i\) by period comparison; applying a \(K_i\)-linear retraction preserving integer rings descends coboundaries on \(2\le c\le t\); and injectivity on divisible coefficients descends zero inertia restriction on \(c=1\). Since \(t+1<\ell\) and \(E_w/\Q_\ell\) is unramified, a cyclotomic inertia generator separates the \(t\) scalar characters by units, giving unit projector denominators in the determinant argument of Lemma~\ref{B8.41b}. Inertia exclusion or the unit tame root difference kills restrictions at \(S_i\), while at \(Q_i\) each character is killed by \(\zeta_{i,v}-1\). With the global Selmer, Sha, and class-number bounds and \eqref{eq:B8-killing}, this yields
\begin{equation}\label{eq:B8-zero-cost}
 \mathcal B_{i,\pi}(\rho;0)\qquad(\pi\in Z_{\rm good}(i)).
\end{equation}
\leanref{GlobalTheory.\allowbreak{}T.\allowbreak{}eventually\_\allowbreak{}congBound\_\allowbreak{}of\_\allowbreak{}facts}{Development.\allowbreak{}AutomorphicArithmetic.\allowbreak{}ThetaCorrespondence.\allowbreak{}ThetaPairLifts.\allowbreak{}ThirdPlaceSeries.\allowbreak{}AdmissibleParameterEventuallyStatement}{11--12}
\leanref{GlobalTheory.\allowbreak{}T.\allowbreak{}localBoundZeroAtLargePrimes\_\allowbreak{}of\_\allowbreak{}congruenceBound}{Development.\allowbreak{}AbelianVarietyTheory.\allowbreak{}Proofs.\allowbreak{}InertiaEigenvalues.\allowbreak{}LocalBoundZeroAtLargePrimesOfCongruenceBoundStatement}{13--15}

\paragraph{The primes \(q,d\).}
At \(\ell\in\{q,d\}\), Lemma~\ref{B8.42bd} supplies crystalline realizations over \(K=E_w\subset K_0(\mu_\ell)\) with \(e=2\) and \(s=(\ell-1)/2\). Proposition~\ref{B8.42bq1} and Lemma~\ref{B8.42br} give the ramified local lifting bound with constant \(160\). Applying it to column coordinates and to the row Tate dual, with their jump bounds, gives a cost \(R_W\le160\{3fr(4b/\ell+2S_c+2\log_\ell(2b)+1+\eta+t)+6fm(2b+1)\}+f\). Character-block finiteness from Lemma~\ref{B8.41a}, local duality and the Euler formula give the blocked character cost at most \(240f^2t^2(r+t)L\), and the projector contribution is at most \(4(r+t)t\). Since \(m\le r\), \(\ell\ge3\), \(L>1\) and \(\log_\ell(2b)\le2L\), the brace in \(R_W\) is at most \(13f\{r(b/\ell+S_c+t+\eta+L)+mb\}\), and the \(f\) term and the projector contribution fit in \(20f^2t^2(r+t)L\). With \eqref{eq:B8-killing} (where \(\nu_\ell=0\) by Lemma~\ref{B8.42a}), this bounds the allowance in \eqref{eq:B8-Bd} by
\begin{equation}\label{eq:B8-qd-poly}
 2080f\{r(b/\ell+S_c+t+\eta+L)+mb\}+260f^2t^2(r+t)L ,
\end{equation}
Because the prime lower bound \(\ell>x^{2C_t+5}\), where \(C_t\) bounds the polynomial rank growth in the auxiliary-prime row, absorbs \(frb/\ell\) and \(S_c+\eta\le x^{C'}\exp(x^{\varepsilon/4})\) makes the remaining polynomial terms eventually smaller than \(fmb\), this allowance is at most \(2Cfmb\) for \(C\) above the hub constant.
\leanref{GlobalTheory.\allowbreak{}T.\allowbreak{}embedding\_\allowbreak{}jumps\_\allowbreak{}aux2\_\allowbreak{}of\_\allowbreak{}setup}{Development.\allowbreak{}AutomorphicArithmetic.\allowbreak{}ThetaCorrespondence.\allowbreak{}ThetaPairLifts.\allowbreak{}ThirdPlaceSeries.\allowbreak{}PolynomialBoundSetup}{11--12}
\leanref{GlobalTheory.\allowbreak{}NLink.\allowbreak{}starG\_\allowbreak{}of\_\allowbreak{}thmB}{Development.\allowbreak{}AbelianVarietyTheory.\allowbreak{}Proofs.\allowbreak{}InertiaEigenvalues.\allowbreak{}PolyValLeSumMulPow}{35--115}

\paragraph{The theta primes.}
Let \(C_T\ge40\) be the type threshold fixed in the proof of Theorem~\ref{B8.42}; it depends on \(\mathscr H_2\) alone, so \(40\le C_T\le b\) eventually. By Proposition~\ref{B8.42bs1}, statement (P) of Definition~\ref{B8.42bs} holds at \(\ell=g_k\) (\(k\in\{0,1\}\)), and Lemma~\ref{B8.42bt} bounds \(\lambda_i(U_i/L_i)\) by the multiplicity-weighted sum of
\begin{equation}\label{eq:B8-theta-cost}
 4fn_V C'n_V^2(n_V+\eta_X+t)+2fr^2+4n_V\delta_X
 +s_\pi+\sum_{1\le e<t}s_{\ell,e}(E)+v_\ell(h_E)+
 t\sum_{v\in Q_i}\lambda_i(O_i/(\zeta_{i,v}-1)).
\end{equation}
Here \(\eta_X=\lfloor\eta\rfloor\le\eta\) and \(s_0(i)=1\) gives \(\delta_X\le2t\). Since \(t\le\eta\) and the ranks grow polynomially in \(x\), the first three terms fit in \(Cx^C\eta\), and \eqref{eq:B8-killing} with \(\nu_{g_k}=0\) gives the first row of \eqref{eq:B8-prime-table} via \eqref{eq:B8-mult-rank}.
\leanref{GlobalTheory.\allowbreak{}T.\allowbreak{}star\_\allowbreak{}bound\_\allowbreak{}of\_\allowbreak{}paired\_\allowbreak{}character\_\allowbreak{}condition}{Development.\allowbreak{}AutomorphicArithmetic.\allowbreak{}ThetaCorrespondence.\allowbreak{}ThetaPairLifts.\allowbreak{}ThirdPlaceSeries.\allowbreak{}EtaLocLeRhsStatement}{12--13}
\leanref{Formalism.\allowbreak{}MainNotions.\allowbreak{}TwistedPlaces.\allowbreak{}etaAt\_\allowbreak{}le\_\allowbreak{}g}{Formalism.\allowbreak{}SharpRestriction.\allowbreak{}LevelLetterCocycles.\allowbreak{}MemPrimesOverOfMem}{108--119}

\paragraph{The final row at split primes.}
Let \(\ell\ge5\) with \(\ell\nmid Q_Sqdg_0g_1\) split in \(\Q(\sqrt{-d})\), and put \(h_N=\lfloor H_B\rfloor\) and \(e_A=\lceil23\{2+2h_N/\ell+\log_\ell(2h_N+1)\}\rceil\). Since \(H_B\ge1\) and \(\ell\ge5\), we have \(\log_\ell(2h_N+1)\le1+\log_\ell H_B\), so \(e_A\le23\{3+2H_B/\ell+\log_\ell H_B\}+1\le70(1+H_B/\ell+\log_\ell H_B)\le P_\ell\) whenever \(C\ge70\). Complex conjugation pairs the column set \(A\) of Lemma~\ref{B8.42bx}(b) with its complement of equal local degrees \(d_u\ge1\), where \(\sum_{u\mid\ell}d_u=2f\). Hence the place sum in Lemma~\ref{B8.42bx}(b) satisfies
\begin{equation}\label{eq:B8-split-place-cost}
 \sum_{u\mid\ell}L_{i,u}=\tfrac12\sum_{u\mid\ell}(2d_um+1)e_A\le\tfrac12\sum_{u\mid\ell}(2d_um+1)P_\ell\le(2fm+f)P_\ell.
\end{equation}
\leanref{GlobalTheory.\allowbreak{}NLink.\allowbreak{}RowA.\allowbreak{}main}{Development.\allowbreak{}AbelianVarietyTheory.\allowbreak{}Proofs.\allowbreak{}InertiaEigenvalues.\allowbreak{}TwoEALe}{49--79}

\paragraph{The final row at inert primes.}
Now let \(\ell\ge5\) with \(\ell\nmid Q_Sqdg_0g_1\) be inert in \(\Q(\sqrt{-d})\). Fix a place \(u\mid\ell\) of \(E\) and a good point \(\pi\in Z_{\rm good}(i)\), and take the unramified field \(K=E_u\) of degree \(d_u\) (with \(\ell\ge11\)) and the crystalline surjection \(F:T_1\to T_2\) of Lemma~\ref{B8.42bx}(a), writing \(q(F)=\lambda_O(Z_f(T_2)/I_F)\). Let \(a_0=\dim_{L'}H^0(K,T_2[1/\ell])\); by \eqref{eq:B8-BK-dim}, \(\operatorname{rank}_O H^1_f(K,T_2)\le d_ur+a_0\), while saturation of the invariant lattice embeds the torsion of \(H^1_f(K,T_2)\) into a divisible \(O\)-module of corank \(r-a_0\), so \(Z_f(T_2)/B(T_2)\) is generated over \(O\) by at most \((d_u+1)r\) elements. Consequently, whenever \(\ell^e\) lifts every class in \(Z_f(T_2)\) modulo \(I_F\) plus the \(O\)-span of \(z\) elements of \(Z_f(T_2)\) and \(\ell^{e'}\) lifts every class in \(Z_f(T_2)\) into \(I_F\), we have
\begin{equation}\label{eq:B8-col-count}
 q(F)\le(d_u+1)re+ze'.
\end{equation}

\emph{Case 1: \(\ell<Y\) and \(o_\ell>4L_w\).}
By Lemma~\ref{B8.42bx}(a), \eqref{eq:B8-word-box} and \eqref{eq:B8-long-word-condition} hold and the thick positions form singletons or pairs separated by at least \(J\) nonthick positions. Applying Lemma~\ref{B8.42bv} gives \eqref{eq:B8-cluster-col}. By \eqref{eq:B8-word-box}, \(h\le H_B\) and \(g=w_*+S+\delta\le w_P+S_c+\eta_3=G\), while \(\log_\ell H_B\le L\) (from \(H_B\le2b=2e^L\) and \(\ell\ge11\)). Moreover \(\log_\ell h\le h/q_\ell^{J+1}+J+1\): this is immediate if \(h<\ell^{J+1}\), and if \(h\ge\ell^{J+1}\) then \(\log_\ell h=(J+1)+\log_\ell(h/\ell^{J+1})\le J+1+h/\ell^{J+1}\le J+1+h/q_\ell^{J+1}\). Thus \eqref{eq:B8-cluster-col} gives \(q(F)\le 3C_2B_u\), where
\begin{equation}\label{eq:B8-long-allowance}
 B_u=(d_u+1)r\{H_B/q_\ell^{J+1}+(G+1)(J+1)+L\}.
\end{equation}
Choose a rational allowance \(L_{i,u}\in[3C_2B_u,(3C_2+1)B_u]\), uniform over \(\pi\in Z_{\rm good}(i)\).
\leanref{GlobalTheory.\allowbreak{}NLink.\allowbreak{}localBoundTableFromData\_\allowbreak{}rowB\_\allowbreak{}rowEntries}{Development.\allowbreak{}AbelianVarietyTheory.\allowbreak{}Proofs.\allowbreak{}InertiaEigenvalues.\allowbreak{}LocalBoundTableFromDataRowBRowEntries}{90--114}

\emph{Case 2: Other inert primes, word with a nonthick position.}
Here Lemma~\ref{B8.42bx}(a) gives alternating signs and \eqref{eq:B8-short-word-condition}, with at least one nonthick position and no three consecutive thick positions. Lemma~\ref{B8.42bw}(a) and (b) then give \eqref{eq:B8-unrestricted-exp} with \(H'\le2H_B\) (by \eqref{eq:B8-word-box}) and the colength bound \eqref{eq:B8-short-col}.

\emph{Case 3: Other inert primes, every position thick.}
By Lemma~\ref{B8.42bx}(a), \(d_u\) is even, the signs alternate, and setting \(h=\max_a s_a+S\), \(H=h+\delta\), \(\lambda=1+\log_\ell(2H)\), and \(A=H(2\ell-1)/q_\ell^2+3+\lambda\) gives \(h/2\le s_a\) and \(2A+\delta+1\le s_a\) for all \(a\in\Z/d_u\Z\). Put \(\sigma=(\lfloor h/q_\ell\rfloor+1+\delta)/q_\ell+1\) and \(\sigma_2=2H/q_\ell^2+1\). By Proposition~\ref{B8.42bu1}, statement (R) of Definition~\ref{B8.42bu} applies with \(\beta_a=s_a\), central width \(\delta\), \(d'=d_u\), and \(p_0\in\{0,1\}\) chosen so that \(\varsigma_a=-1\) exactly when \(a+p_0\) is even, supplying \(a_1\le A+1\), \(s_1,E_1,e_2\), and \(e\le s_1+2\sigma+e_2+S+E_1+\sigma_2+1\). Substituting the bounds of Definition~\ref{B8.42bu} for \(s_1,E_1,e_2\) gives
\begin{equation}\label{eq:B8-all-thick-exp}
 e\le32H/\ell^2+4S+10\delta/q_\ell+25\lambda ,
\end{equation}
since for \(\ell\ge11\) (\(q_\ell\ge10\), \(\lambda\ge1\), \(\lfloor h/q_\ell\rfloor\le H/q_\ell\)) the coefficient of \(H\) times \(\ell^2\) is at most \(8(21/10)(11/10)^2+8(11/10)^2+11/10<32\), the coefficient of \(S\) is at most \(3+3/q_\ell\le4\), the coefficient of \(\delta\) is at most \(10/q_\ell\), and the remaining constants and logarithms fit in \(25\lambda\). Statement (R) also provides the free module \(F_0\) of rank at most \((d_u+1)r\), the surjection \(\jmath:F_0\to Z_f(T_2)\) modulo \(B(T_2)\), the chain \(F_0\supset\cdots\supset F_k\) with \(\ell^e\jmath(F_k)\subset I_F\), and \(\mathcal G\subset\{0,\dots,k-1\}\) with \(|\mathcal G|\le4d_u\) and \(k-|\mathcal G|\le3d_u/2\). Setting \(I'=\jmath^{-1}(I_F)\) identifies \(F_0/I'\simeq Z_f(T_2)/I_F\): the bottom piece \(F_k/(F_k\cap I')\) has normalized length at most \((d_u+1)re\); each of the at most \(4d_u\) pieces \(i\in\mathcal G\) is generated modulo the next by at most \(m\) elements annihilated in \(Z_f(T_2)/I_F\) by \(\ell^{e'}\) from Lemma~\ref{B8.42bw}(a) (with \(H'\le2H_B\)); and each of the at most \(3d_u/2\) remaining pieces has normalized length at most \(rS+2m((H+1)/q_\ell+1)\). Additivity of length gives
\begin{equation}\label{eq:B8-all-thick-col}
 q(F)\le(d_u+1)re+4d_ume'+\tfrac32 d_u\{rS+2m((H+1)/q_\ell+1)\}.
\end{equation}

\emph{Uniform inert allowances and the band condition.}
Let \(b_u\in\{0,1\}\) be the indicator that Case 2 occurs with \(I=1\) at some \(\pi\in Z_{\rm good}(i)\) at \(u\). Put \(E_\ell=1+H_B/\ell+\log_\ell H_B\) and
\[
 A_u=(d_u+1)rH_B/\ell^2,\qquad M_u=d_umE_\ell,\qquad
 F_u=(d_u+1)r\{(w_P+\eta_3)/\ell+S_c+L+\eta_3 b_u\}.
\]
From \(1\le H_B\le2e^L\) and \(\ell\ge11\), \eqref{eq:B8-unrestricted-exp} with \(H'\le2H_B\) gives \(e'\le69E_\ell\). In Case 2, \(\log_\ell(2h)\le3+L\) and \eqref{eq:B8-word-box} bound \eqref{eq:B8-short-col} by \(320A_u+1920F_u+22080M_u\). In Case 3, \(H\le H_B+\eta_3\), \(\log_\ell(2H)\le7+L\), \(32\eta_3/\ell^2\le3\eta_3/\ell\), and \(10\delta/q_\ell\le11\eta_3/\ell\) give \(e\le32H_B/\ell^2+225F_u/((d_u+1)r)\) in \eqref{eq:B8-all-thick-exp} and \((H+1)/q_\ell+1\le4E_\ell\) in \eqref{eq:B8-all-thick-col}, so \(q(F)\le32A_u+227F_u+288M_u\). Thus in both Case 2 and Case 3,
\begin{equation}\label{eq:B8-short-allowance}
 q(F)\le368A_u+2082F_u+22365M_u .
\end{equation}
Choose a rational allowance \(L_{i,u}\in[368A_u+2082F_u+22365M_u,\ 368A_u+2083F_u+22365M_u]\) (and \(L_{i,u}=0\) if \(Z_{\rm good}(i)=\varnothing\)). If \(b_u=1\), take the witnessing nonthick position \(a\) in \eqref{eq:B8-local-band} and its \(y\in[w''(n),w_B]\) from Lemma~\ref{B8.42bx}(a) with \(|s_a-y|\le T_0\): since \(|b_*-b|\le T_0\), \((\ell+1)T_0\le\ell\eta_3\), and \(\log_\ell(2h)\le L\), the brace \(D_{\rm band}=\ell\eta_3+\ell S_c+b/\ell+y+\ell L\) of \eqref{eq:B8-global-band} bounds the brace of \eqref{eq:B8-local-band} by \(D_{\rm band}+\ell\eta_3\le2D_{\rm band}\), whence \(|\ell y-b|\le|\ell s_a-b_*|+\ell T_0+T_0\le641D_{\rm band}\) and therefore \(b_u\le I_\ell\) whenever \(C\ge641\).
\leanref{Formalism.\allowbreak{}MainNotions.\allowbreak{}N.\allowbreak{}KillingConditions.\allowbreak{}RowK.\allowbreak{}Le}{Formalism.\allowbreak{}SharpRestriction.\allowbreak{}LevelLetterCocycles.\allowbreak{}JointConditionsZpowers.\allowbreak{}CommutingEndomorphismCondition}{180--183}
\leanref{GlobalTheory.\allowbreak{}T.\allowbreak{}mediumInert\_\allowbreak{}smallInert\_\allowbreak{}lengthBound\_\allowbreak{}tabX}{Development.\allowbreak{}AbelianVarietyTheory.\allowbreak{}Proofs.\allowbreak{}InertiaEigenvalues.\allowbreak{}LocalBoundTableRowsCDAdmissibleStatement}{7--9}
\leanref{GlobalTheory.\allowbreak{}T.\allowbreak{}localBoundTableFromData\_\allowbreak{}bandCondition}{Development.\allowbreak{}AbelianVarietyTheory.\allowbreak{}Proofs.\allowbreak{}InertiaEigenvalues.\allowbreak{}LocalBoundTableFromDataBandConditionStatement}{7--9}

Require \(C\ge\max(3C_2+1,22365,7200)\). Using \(\sum_{u\mid\ell}d_u=2f\), \(\sum_{u\mid\ell}(d_u+1)\le4f\), and \(b_u\le I_\ell\), summing \(L_{i,u}\) over \(u\mid\ell\) in \eqref{eq:B8-split-place-cost}, \eqref{eq:B8-long-allowance}, and \eqref{eq:B8-short-allowance} gives \(\sum_{u\mid\ell}L_{i,u}\le D_\ell\) in all three branches of \eqref{eq:B8-prime-table-depth}. Since \(7200\le C\le Cn\), Lemma~\ref{B8.42bx}(b) yields
\begin{equation}\label{eq:B8-middle-cost}
 \mathcal B_{i,\pi}(\rho;Cfn t^2n_VL^2+D_\ell)\qquad(\pi\in Z_{\rm good}(i)).
\end{equation}
\leanref{GlobalTheory.\allowbreak{}T.\allowbreak{}sizeConditions\_\allowbreak{}mono}{Development.\allowbreak{}AbelianVarietyTheory.\allowbreak{}Proofs.\allowbreak{}InertiaEigenvalues.\allowbreak{}ParametersConditionsMonotoneInConstantStatement}{13--15}
\leanref{Formalism.\allowbreak{}MainNotions.\allowbreak{}TwistedPlaces.\allowbreak{}loc\_\allowbreak{}le\_\allowbreak{}lamB}{Formalism.\allowbreak{}SharpRestriction.\allowbreak{}LevelLetterCocycles.\allowbreak{}MemPrimesOverOfMem}{384--500}
Combining \eqref{eq:B8-eta-passage}, \eqref{eq:B8-Bd23} (and its odd-\(Q_S\) counterpart), \eqref{eq:B8-zero-cost}, \eqref{eq:B8-qd-poly}, \eqref{eq:B8-theta-cost}, and \eqref{eq:B8-middle-cost} with \eqref{eq:B8-killing} and \eqref{eq:B8-mult-rank} proves \eqref{eq:B8-cost-prime}.
\leanref{Formalism.\allowbreak{}MainNotions.\allowbreak{}TwistedPlaces.\allowbreak{}etaAt\_\allowbreak{}le}{Formalism.\allowbreak{}SharpRestriction.\allowbreak{}LevelLetterCocycles.\allowbreak{}MemPrimesOverOfMem}{57--107}
\leanref{Formalism.\allowbreak{}MainNotions.\allowbreak{}TwistedPlaces.\allowbreak{}etaSum\_\allowbreak{}le}{Formalism.\allowbreak{}SharpRestriction.\allowbreak{}LevelLetterCocycles.\allowbreak{}MemPrimesOverOfMem}{558--648}

\paragraph{Summation of the table.}
Put \(X_0=fmbL^2\) and \(U=x^{C+2}\exp(x^{\varepsilon/4})\). Since \(L=x^\varepsilon\) and \(b=e^L\) with \(\varepsilon>0\), eventually
\begin{equation}\label{eq:B8-growth-sum}
 4\cdot10^6(C+1)^2U^{18}\le b,\qquad
 (C+3)^2(1+\log x)^2\le L,\qquad x^CL\le b .
\end{equation}
\leanref{Formalism.\allowbreak{}MainNotions.\allowbreak{}TwistedPlaces.\allowbreak{}big\_\allowbreak{}le\_\allowbreak{}b}{Formalism.\allowbreak{}SharpRestriction.\allowbreak{}LevelLetterCocycles.\allowbreak{}ExportBoundOfAdmissible}{501--541}
Polynomial growth and the width and type-prime size bounds give \(\sum_{\ell\in\{g_0,g_1\}}A_\ell\log\ell\le2(C+1)^2U^{18}\). The prime divisors of \(3Q_S\) have total logarithm at most \(\log3+2\log n\), so their contribution is at most \(8CX_0\); at \(q,d\), \(\log q,\log d\le L\) gives at most \(4CX_0\).
\leanref{Formalism.\allowbreak{}MainNotions.\allowbreak{}TwistedPlaces.\allowbreak{}sum\_\allowbreak{}g}{Formalism.\allowbreak{}SharpRestriction.\allowbreak{}LevelLetterCocycles.\allowbreak{}ExportBoundOfAdmissible}{651--706}
\leanref{Formalism.\allowbreak{}MainNotions.\allowbreak{}TwistedPlaces.\allowbreak{}sum\_\allowbreak{}23}{Formalism.\allowbreak{}SharpRestriction.\allowbreak{}LevelLetterCocycles.\allowbreak{}ExportBoundOfAdmissible}{733--753}
\leanref{Formalism.\allowbreak{}MainNotions.\allowbreak{}TwistedPlaces.\allowbreak{}sum\_\allowbreak{}qd}{Formalism.\allowbreak{}SharpRestriction.\allowbreak{}LevelLetterCocycles.\allowbreak{}ExportBoundOfAdmissible}{798--824}

The intervals \eqref{eq:B8-exception-intervals} contain at most \((|TT|+2)(S_c+3\eta+t+8b/Z+9)\le388U^8\) integers in \([1,3b]\). Because \(Z^2>3b\), each such integer has at most one prime divisor \(\ell>Z\). Using \(b/Z\le10U^7\) in the final-row formulas gives \(A_\ell\le1000(C+1)^2U^9\) and \(\log\ell\le\log(3b)\le3U\) at every exceptional \(\ell>Z\), for a total contribution of at most \(1164000(C+1)^2U^{18}\).
\leanref{Formalism.\allowbreak{}MainNotions.\allowbreak{}TwistedPlaces.\allowbreak{}sum\_\allowbreak{}exc}{Formalism.\allowbreak{}SharpRestriction.\allowbreak{}LevelLetterCocycles.\allowbreak{}ExportBoundOfAdmissible}{611--650}

For primes \(5\le\ell\le Z\), the definition of \(Z\) and \(\sum_{\ell\le Z}\log\ell\le Z\log Z\le ZL\) bound the non-place sum \(\sum_{5\le\ell\le Z}Cfnt^2n_VL^2\log\ell\) by \(CX_0\). Integral comparison gives
\begin{equation}\label{eq:B8-summation-comparisons}
 \sum_{2\le a\le Z}\frac{\log a}{a}\le(1+\log Z)^2,\quad
 \sum_{5\le a<Y}\frac{\log a}{a-1}\le2(1+\log Y)^2,\quad
 \sum_{a\ge Y}\frac{\log a}{a^2}\le\frac{4(\log Y+2)}Y .
\end{equation}
From the first inequality in \eqref{eq:B8-summation-comparisons}, \(H_B\le2b\), and \(ZL\le b\),
\begin{equation}\label{eq:B8-P-sum}
 \sum_{\substack{5\le\ell\le Z\\\ell\ {\rm prime}}} P_\ell\log\ell\le11CbL^2.
\end{equation}
In the long-order inert row (\(\ell<Y\), \(o_\ell>4L_w\)), use \((\ell-1)^J\ge Y\), \(r/Y\le2m/L\), \(rb'Y\le mb\), \(G+1\le11b'\), and \(J+1\le3+\log Y\) with \((3+\log Y)^2\le L\); the second sum in \eqref{eq:B8-summation-comparisons} and \(\sum_{a<Y}\log a\le Y\log Y\) bound its total by \(80CX_0\). For inert \(\ell\ge Y\), the third sum in \eqref{eq:B8-summation-comparisons} bounds the \(4fCr(H_B/\ell^2)\log\ell\) terms by \(64CX_0\).

For inert \(\ell<Y\) with \(o_\ell\le4L_w\), the \(2\)-adic valuation of \(\ell^{o_\ell}\mp1\) puts \(\ell\equiv\pm1\pmod M\), where \(M=2^{\max(0,e+2-\lfloor\log_2(4L_w)\rfloor)}\). With \(R=n/L\), the host gives \(R+1\le M\), and summing along both progressions yields
\begin{equation}\label{eq:B8-short-order-sum}
 \sum_{\substack{\ell<Y\ {\rm inert}\\o_\ell\le4L_w}}\frac{\log\ell}{\ell^2}
 \le\frac{2\log R}{R^2}+\frac{2(\log R+2)}{RM}.
\end{equation}
Since \(r\le2mR^2\), \(n\le oML_w\), \((\log R+2)L_w\le L\), and \(\log R+o\le L^2\) for the host order parameter \(o\), multiplying \eqref{eq:B8-short-order-sum} by \(r\) gives at most \(4mL^2\), whose contribution to the table sum is at most \(32CX_0\).

For \(F_\ell\), use \(r(w_P+\eta_3)\le9mb\) and \(r(S_c+L)Z\le2mb\). When \(I_\ell=1\) and \(\ell>2C\), \eqref{eq:B8-global-band} gives \((\ell-C)y\le(1+C/\ell)b+C\ell(\eta_3+S_c+L)\) with \(C(\eta_3+S_c+L)\le w''(n)/8\le y/8\), whence \(w''(n)\le y\le4b/\ell\); combined with \(r\eta_3\le4mw''(n)\) and \(r\eta_3\le4mb\), this yields
\begin{equation}\label{eq:B8-flat-band-sum}
 r\eta_3 I_\ell\le16mb/\ell+4mb\,\mathbf1_{\ell\le2C}.
\end{equation}
Summing via \eqref{eq:B8-summation-comparisons} and \(\sum_{a\le2C}\log a\le(2C+1)^2\) bounds the \(F_\ell\) contribution by \(4C[102+4(2C+1)^2]X_0\). Adding all contributions over \(5\le\ell\le Z\) gives coefficient at most
\[
 C+11C(3+2C)+80C+64C+32C+4C[102+4(2C+1)^2]\le1000(C+1)^3
\]
times \(X_0\).
\leanref{Formalism.\allowbreak{}MainNotions.\allowbreak{}TwistedPlaces.\allowbreak{}sum\_\allowbreak{}small}{Formalism.\allowbreak{}SharpRestriction.\allowbreak{}LevelLetterCocycles.\allowbreak{}ExportBoundOfAdmissible}{18--80}
Finally, \(Q_S\le n^2\), \(\log h_E\le Cx^C\), and \(\log M_0\le Cx^CL\) give \(6+\log(Q_Str)+\log h_E+\log M_0\le20(C+1)^2U^{18}\). By \eqref{eq:B8-growth-sum} and \(b\le X_0\), the sum \((20+2+1164000)(C+1)^2U^{18}\) is at most \(b\le X_0\), and \((1000(C+1)^3+8C+4C+1)X_0\le2000(C+1)^3X_0\). This proves \eqref{eq:B8-cost-sum} for every finite set \(S\) of primes.
\leanref{Formalism.\allowbreak{}MainNotions.\allowbreak{}TwistedPlaces.\allowbreak{}glob\_\allowbreak{}le}{Formalism.\allowbreak{}SharpRestriction.\allowbreak{}LevelLetterCocycles.\allowbreak{}ExportBoundOfAdmissible}{825--871}
\leanref{Formalism.\allowbreak{}MainNotions.\allowbreak{}TwistedPlaces.\allowbreak{}point\_\allowbreak{}bound}{Formalism.\allowbreak{}SharpRestriction.\allowbreak{}LevelLetterCocycles.\allowbreak{}ExportBoundOfAdmissible}{893--997}
\leanref{GlobalTheory.\allowbreak{}NLinkClose.\allowbreak{}main\_\allowbreak{}alt\_\allowbreak{}alt}{Development.\allowbreak{}AutomorphicArithmetic.\allowbreak{}ThetaCorrespondence.\allowbreak{}ThetaPairLifts.\allowbreak{}ThirdPlaceSeries.\allowbreak{}ImportsFurtherTail}{22--109}
\end{proof}

\subsection*{Proof of \ref{B8.43i}}\label{prf:B8.43i}\phantomsection
\begin{proof}
Let \(T_\psi=\sum_q\psi(q)^{-1}T(q)\). Apply Definition~\ref{B1.14k}(3), proved in Lemma~\ref{B1.15a}, first to \((Q,\psi)\) and then to \((\varnothing,1)\). The second application has the same field, CM type, base conductor and weights as the first. The empty set satisfies the condition on auxiliary primes. Thus there is one nonzero period \(\Omega_0\) for which both normalized sums are algebraic. The function \(x_{\Omega_0}\) multiplies each sum by the same nonzero scalar. Since \(T_0\ne0\), their quotient is \(R_\psi\), which belongs to the chosen algebraic closure of \(K\).

Write \(g^{\rm ab}=g|_{K^{\rm ab}}\). The hypothesis on \(t\) is
\[
                 \operatorname{Art}_K((1,t))^{-1}=g^{\rm ab},
\]
where \(1\) is the archimedean component. This is the hypothesis on the finite idèle in Definition~\ref{B1.14k}(3). Extend \(g\) from the embedded algebraic closure to \(\C\) by fixing a transcendence basis and then extending to its algebraic closure. The extension fixes each CM embedding pointwise on its image. It also fixes the prescribed values \(\chi_f(y)|y|_{E,f}^{s}\). Hence the conjugate weights and base character in the transport formula are the original ones. The conjugate Picard character is \(g\psi\), given by \((g\psi)(q)=g(\psi(q))\).

Apply \eqref{eq:p1c1-twisted-transport} twice with this \(g\) and \(t\). To see the character factor, write \(q'=q[t]^{-1}\) for the transported class. Then
\[
             g(\psi(q)^{-1})
               =(g\psi)(q')^{-1}(g\psi)([t])^{-1}.
\]
Translation by \([t]^{-1}\) permutes the classes. Put
\(A=|t|_{K,f}^{-K_\kappa}l_\chi^{-1}\), where \(l_\chi\) is evaluated on the image of \(t\) in the finite idèles of \(E\). The two transport identities are
\[
 \begin{aligned}
  g x_{\Omega_0}(T_\psi)
    &=x_{\Omega_1}\bigl(A(g\psi)([t])^{-1}T_{g\psi}\bigr),\\
  g x_{\Omega_0}(T_0)&=x_{\Omega_1}(A T_0).
 \end{aligned}
\]
The target period is common because it is chosen before the conductor and character. The norm of an idèle is positive and a character takes values in the multiplicative group, so \(A\ne0\). Division gives
\[
                   gR_\psi=(g\psi)([t])^{-1}R_{g\psi}.
\]
Both sides belong to the embedded algebraic closure, where the extension of \(g\) agrees with the given automorphism.
\end{proof}

\subsection*{Proof of \ref{B8.43j}}\label{prf:B8.43j}\phantomsection
\begin{proof}
Choose integral generators \(g_1,\ldots,g_K\) of the underlying tensor lattice before completion. One can take the products of integral additive generators of the coefficient ideal and the two input lattices. They generate the completed lattice modulo \(M\) for every \(M\), since reduction of a completed integral lattice modulo \(M\) is reduction of the original lattice. This follows prime by prime from the Chinese remainder theorem. Projecting gives generators \(eg_j\) of \(L_+/ML_+\).

For \(c\in(\Z/M\Z)^K\), use representatives between \(0\) and \(M-1\) and set \(y_c=M^{-1}\sum_j c_jeg_j\). For \(t\in L_-\),
\[
  \sum_c\psi(\langle t,y_c\rangle)=
  \prod_{j=1}^K\sum_{b=0}^{M-1}
        \psi(\langle t,eg_j/M\rangle)^b
  =M^K{\bf1}_{ML_-}(t).
\]
Indeed each inner character value has \(M\)-th power one. If all of them are one, \(t/M\) annihilates \(L_+\). The conjugate halves pair perfectly by self-duality, so \(t/M\in L_-\). Conversely every \(t\in ML_-\) gives the trivial character.

Take representatives \(v\) of \(L_-/ML_-\) and put
\[
                   b_c=\sum_v\Phi(v)\psi(-\langle v,y_c\rangle).
\]
Each \(b_c\) belongs to \(R\). By the displayed orthogonality identity,
\(\sum_c b_c\psi(\langle t,y_c\rangle)=M^K\Phi(t)\) on \(L_-\). In the half-lattice model the translate of the vacuum by \(y_c\) has coefficient function \({\bf1}_{L_-}(t)\psi(\langle t,y_c\rangle)\). Equality of coefficient functions gives equality of the vectors. Finally \(M^K/\ell^{K\operatorname{ord}_\ell(M_0)}\) is a unit at every place over \(\ell\).
\end{proof}

\subsection*{Proof of \ref{B8.43k}}\label{prf:B8.43k}\phantomsection
\begin{proof}
Choose additive generators of \(L_0\). Their images involve finitely many polynomial coefficients. If there is an embedded number field as in the statement, each such coefficient \(z\) is algebraic. Choose a positive integer \(a_z\) such that \(a_zz\) is an algebraic integer, and take \(c_1=\max_z\operatorname{ord}_\ell a_z\), with maximum zero for an empty list. The factor of \(a_z\) prime to \(\ell\) is invertible at every tested place. Thus \(\ell^{c_1}z\) is integral there. Additivity gives the assertion on \(L_0\).

For the second assertion \(I\) divides \(|\mathcal O_E^n/N\mathcal O_E^n|=N^{[E:\Q]n}\). The group \(A\) is a quotient of \(\mathcal L/M\mathcal L\), so \(|A|\) divides \(M^s\). Taking valuations in the definition of \(H\) and using \(\operatorname{ord}_\ell N=\operatorname{ord}_\ell M=\operatorname{ord}_\ell M_0\) gives the formula.
\end{proof}

\subsection*{Proof of \ref{B8.43l}}\label{prf:B8.43l}\phantomsection
\begin{proof}
The restriction to \(L^\times\) of the valuation of \(O\) has image \(a_0\Z\) for a positive integer \(a_0\), since it is nonzero on \(\ell\). Dividing by \(a_0\) gives \(\operatorname{ord}_w\), with valuation ring \(O\cap L\). In particular \(e_w=\operatorname{ord}_w(\ell)\le e\).

By the Chinese remainder theorem choose \(z\in\mathcal O_L\) congruent to \(1\) at \(w\) and to \(0\) modulo \((w')^b\) at each of the other primes above \(v\). Then \(\operatorname{Tr}_{L/K}(czy)\in\mathcal O_{K,v}\) for every \(c\in\mathcal O_L\). Localize at \(v\) and complete. The finite algebra of integers becomes \(\prod_{w'\mid v}\mathcal O_{L_{w'}}\), and its trace is the sum of the local traces. Testing the factor supported at \(w\) shows that \(zy\) belongs to the inverse different of \(L_w/K_v\), the trace dual of \(\mathcal O_{L_w}\). Since \(z\) is a unit there, \(y\) belongs to the same inverse different.

For completeness, the exponent of the absolute different of \(L_w/\Q_\ell\) is less than \(e_w^2\). Let \(U\) be its maximal unramified subextension. Its residue field is the residue field of \(L_w\): lift a separable generator of that finite field by Hensel's lemma. Thus \([L_w:U]=e_w\). Choose \(s\in\mathcal O_U\) whose residue has nonzero trace to \(\F_\ell\), which is possible by separability. The trace of \(s\) to \(\Z_\ell\) from \(U\) is a unit, as one sees by reducing the multiplication matrix in an integral basis. If the different exponent were at least \(e_w^2\), the absolute inverse different would contain \(\ell^{-e_w}\). Its pairing with \(s\) would be integral. That pairing is
\[
          \ell^{-e_w}e_w\operatorname{Tr}_{U/\Q_\ell}(s),
\]
whose valuation is negative since \(\operatorname{ord}_\ell(e_w)<e_w\). This proves the bound. The relative different divides the absolute different by transitivity. Since \(\operatorname{ord}_w(\ell^{e^2})=e_we^2\), multiplication by \(\ell^{e^2}\) clears the relative inverse different. In the unramified case the trace pairing on the integers is perfect: its reduction is the nondegenerate residue-field trace pairing. The relative different is then the unit ideal.
\end{proof}

\subsection*{Proof of \ref{B8.43m}}\label{prf:B8.43m}\phantomsection
\begin{proof}
Pass to a finite extension over which the \(N^2\)-torsion is rational. It suffices to prove the result there, since extension multiplies the valuation on the original field by a positive integer. Denote the abelian scheme by \(B\), its rigidified principal theta line by \(L\), and the theta generator in the frame by \(s\). The other entries of the frame are integral bases of the prescribed eigenspaces of \(\operatorname{Lie} B\).

Here is the sum to be estimated. Its first index is \(q=(q_b)\in A^n\). Choose the designated \(N\)-division point of the marked point for \(q_b\), in the primary torsion belonging to the marking ideal. The second index runs through the tuples \(z=(z_b)\in B[N^2]^n\) cut out by \(\Lambda\). Explicitly, for every \(E\)-matrix \(H\) with \(H\Lambda\subseteq\mathcal O_E^n\) and \(NH\) integral, the tuple satisfies \((NH)z=0\). For column \(b\), pull \(s\) back by \([N]\), identify \([N]^*L\) with \(L^{N^2}\) by the rigidification, and make the symmetric theta-group translations first by \(z_b\), then by the chosen division point. Denote the resulting section by \(F_{qzb}\). Both points are killed by \(N^2\), since the marked point is killed by \(M\mid N\). Write \(P=\sum_\alpha P_\alpha X^\alpha\) in the framed tangent variables, and let \(\alpha_b\) be the exponents in column \(b\). If \(\nabla^{\alpha_b}F(0)\) denotes the corresponding covariant derivative in the theta frame at the identity, the short value is
\[
 v=\frac{1}{I^r|A|^n N^D}
       \sum_{q\in A^n} w(q)\sum_z\sum_\alpha P_\alpha
                    \prod_{b=1}^n\nabla^{\alpha_b}F_{qzb}(0).
\]
For \(P=0\) take \(v=0\). For \(P\ne0\), homogeneity gives precisely the exponent \(D\) in this expression.
\leanref{Formalism.\allowbreak{}L2.\allowbreak{}Pj.\allowbreak{}ShortValAt}{Formalism.\allowbreak{}LocalToGlobal.\allowbreak{}ProjectiveLattices.\allowbreak{}ThShort}{244--261}

Each \(F_{qzb}\) has an integral germ at the identity. Indeed, the torsion points extend by properness. The rigidified multiplication isomorphism extends by the cubical law. To see integrality of a symmetric translation, compare the translated line with the original one on the abelian scheme. Their generic isomorphism can differ from an integral isomorphism only by a multiple of the special fibre. The special fibre is geometrically integral, and translation preserves its order. A power of the symmetric lift is the identity up to a root of unity. Adding the order of the isomorphism under iteration therefore gives zero, so its order is zero. This argument applies in residue characteristic dividing \(N\). Restricting the two translated sections to the formal completion gives the asserted integral germs.
\leanref{GlobalTheory.\allowbreak{}ReachRegular.\allowbreak{}regular\_\allowbreak{}of\_\allowbreak{}reach}{Development.\allowbreak{}AbelianVarietyTheory.\allowbreak{}Proofs.\allowbreak{}RingsOfIntegers.\allowbreak{}ReachedSectionRegularOnIntegralModel}{71--100}
\leanref{GlobalTheory.\allowbreak{}GoodThetaBound.\allowbreak{}reach\_\allowbreak{}of\_\allowbreak{}isUUOver}{Development.\allowbreak{}AbelianVarietyTheory.\allowbreak{}Proofs.\allowbreak{}InertiaEigenvalues.\allowbreak{}ShortThetaValueBoundedGoodReduction}{53--60}

We give the differential calculation since its denominator must be independent of \(N\). Use integral formal parameters for \(B\) and a trivialisation of its formal biextension. Let \(F(X,Y)\) be the formal addition law. In characteristic zero a logarithmic gauge is a series \(g(X,Y)\), zero on the axes, whose two additivity defects are the logarithms of the units giving biextension multiplication. The symmetry and adjoint-action identities are the ones for the given polarisation. Put \(\beta_{ij}=\partial_{X_i}\partial_{Y_j}g(0,0)\), and write \(A(a)\) for the tangent matrix of \(a\in\mathcal O_E\). Taking the mixed coefficient in the symmetry and adjoint identities gives, entrywise,
\[
       \beta^t-\beta\in V,\qquad A(a)^t\beta-\beta A(\bar a)\in V .
\]
The correction terms are integral because the logarithmic derivative of a formal unit \(u\) is \(u^{-1}du\), and the units in these identities are one on the axes.

Take an integral basis \(x_i\) of \(\mathcal O_E\) and multiply its trace-dual basis by \(\Delta\) to obtain integral \(y_i\). In the separability tensor one has \(\sum_i x_i y_i=\Delta\) and \(\sum_i ax_i\otimes y_i=\sum_i x_i\otimes ay_i\). Put
\[
 B_0=\sum_i A(x_i)^t\beta A(\bar y_i),
 \qquad B_1=(B_0+B_0^t)/(2\Delta).
\]
The tensor identity gives \(A(a)^tB_1=B_1A(\bar a)\). Moreover
\[
 \Delta\beta-B_0
   =\sum_i A(x_i)^t\bigl(A(y_i)^t\beta-\beta A(\bar y_i)\bigr).
\]
Consequently \(2\Delta(\beta-B_1)\) is integral. Subtract from \(g\) the bilinear form with matrix \(B_1\) in formal logarithm coordinates. It is symmetric, additive in both variables, and balanced for conjugation. Hence the new gauge \(g'\) satisfies the same gauge identities. In harmonic evaluation this change contributes contractions between conjugate tangent variables. Each contraction annihilates \(P\) by harmonicity. Equivalently, on expanding the exponential of the added quadratic form, every term of positive degree pairs with an iterate of a contraction on \(P\).
\leanref{GlobalTheory.\allowbreak{}GaugeBoundedDenom.\allowbreak{}exists\_\allowbreak{}symm\_\allowbreak{}balanced}{Development.\allowbreak{}AbelianVarietyTheory.\allowbreak{}Proofs.\allowbreak{}InertiaEigenvalues.\allowbreak{}ExistsCanonGaugeNabValBoundedDenom}{20--70}
\leanref{GlobalTheory.\allowbreak{}GoodThetaBound.\allowbreak{}hv\_\allowbreak{}bound}{Development.\allowbreak{}AbelianVarietyTheory.\allowbreak{}Proofs.\allowbreak{}InertiaEigenvalues.\allowbreak{}ShortThetaValueBoundedGoodReduction}{98--157}

Let \(w_0\in\mathcal O_E\) be the polarisation element with \(w_0+\bar w_0=1\), put \(q(X)=g'(X,[w_0]X)\), and denote the invariant derivative in an integral tangent direction \(u\) by \(\partial_u\). We have \(q(0)=\partial_uq(0)=0\). Set \(\delta=2\Delta\). Then \(\delta\,\partial_u\partial_v q\) is an integral series. To compute it modulo integral series, expand \(q(F(X,Y))\) by additivity in the two arguments of \(g'\). The four terms are \(q(X),q(Y),g'(X,[w_0]Y),g'(Y,[w_0]X)\), up to the logarithms of integral units. Differentiate in \(Y\) at zero and then invariantly in \(X\). Derivatives of the logarithmic defects are integral. Applying additivity once more to each mixed term leaves, modulo integral series, the constant
\[
       (\beta-B_1)(u,A(w_0)v)+(\beta-B_1)(v,A(w_0)u).
\]
Multiplication by \(\delta\) clears it. This computation uses integral invariant vector fields, obtained by differentiating the integral addition law.

The subring generated by the integral series and the \(\delta\partial_v q\) has integral constant terms and is stable under \(\delta\partial_u\). For a section of \(L^k\) with integral germ \(f\), its derivatives in this gauge are the derivatives at zero of \(f\exp(kq)\). After removing \(\exp(kq)\), one more scaled derivative takes \(h\) to
\[
                  \delta\partial_u h+k h\,\delta\partial_u q .
\]
Induction gives integrality at zero after \(|\alpha_b|\) such factors. Multiplying over the columns costs \(\delta^{\sum_b|\alpha_b|}=\delta^D\). The coefficients \(w(q)P_\alpha\) are integral by hypothesis. Substitution in the formula for \(v\) proves the assertion.
\leanref{GlobalTheory.\allowbreak{}GaugeBoundedDenom.\allowbreak{}pow\_\allowbreak{}mul\_\allowbreak{}const\_\allowbreak{}dmonOf\_\allowbreak{}mem}{Development.\allowbreak{}AbelianVarietyTheory.\allowbreak{}Proofs.\allowbreak{}InertiaEigenvalues.\allowbreak{}ExistsCanonGaugeNabValBoundedDenom}{76--150}
\end{proof}

\subsection*{Proof of \ref{B8.43n}}\label{prf:B8.43n}\phantomsection
\begin{proof}
Replace \(f\) by \(af\). Write \(V\) for the test ring, \(\pi\) for its uniformiser, and \(K_V\) for its fraction field. Use the proper semiabelian model at \(\lambda\). The prime-to-\(\ell\) marking has level at least three, as required for this model. The valuative criterion extends the underlying moduli point. If the extension lies in the abelian locus, the hypothesis applies. In definite signature the torus part vanishes: the CM action on the Lie algebra at the generic point has just one of the two eigenvalues of \(w_0\), hence acts there by a scalar. The scalar is integral, and the same equality holds on the integral Lie module. A nonzero torus character lattice is torsion free over the quadratic integer ring generated by \(w_0\). Locally at the residue characteristic it is projective over that ring, and its differential representation contains the regular quadratic algebra after reduction, on which \(w_0\) is not scalar. This excludes a torus part.

In the remaining case pass to a finite valued extension to split the semiabelian degeneration, rationalise the finite marking and make the orders of its quadratic periods divisible by twice the full level. These operations are possible because the character lattice is finitely generated, the required division points are algebraic, and ramification multiplies all period orders by its index. Continue to denote the new ring by \(V\). Let \(X\) be the torus character lattice, \(Q(\alpha)\) the order of the rigidified period at \(\alpha\), and \(B(\alpha,\beta)=Q(\alpha+\beta)-Q(\alpha)-Q(\beta)\). This is the symmetric conjugate-balanced monodromy pairing. In particular \(B(\alpha,w_0\alpha)=Q(\alpha)\). In the degeneration data remove the uniformiser powers from the periods and reinsert them with a variable \(t\) in place of \(\pi\). Thus the factors for the quadratic and bilinear periods are \(t^{Q(\alpha)}\) and \(t^{B(\alpha,\beta)}\), respectively, in the corresponding rigidifications. Additivity of \(B\) and the displayed polarisation identity are precisely the multiplication identities for the new periods. Divisibility by twice the full level gives integral exponents for the roots in the marking. This is the one-parameter toroidal chart of the split degeneration. It specialises at \(t=\pi\) to the given point. Where \(t\) is a unit its model lies in the abelian locus. The generic moduli map and its frame are defined over \(V[[t]][1/(\pi t)]\). For the frame take a generator on the abelian part and constant bases in the tangent eigenspaces of the extension. One can arrange that their specialisation and the given tangent bases differ by a triangular matrix: intersect each subspace of the given flag with the free integral tangent module and choose a basis successively in the resulting saturated submodules. The two integral generators differ by a unit. The diagonal entries of a triangular change between two integral bases are units as well.
\leanref{GlobalTheory.\allowbreak{}ArcOfPoint.\allowbreak{}exists\_\allowbreak{}arcIn\_\allowbreak{}of\_\allowbreak{}unif}{Development.\allowbreak{}AbelianVarietyTheory.\allowbreak{}Proofs.\allowbreak{}InertiaEigenvalues.\allowbreak{}ExistsArcDataOfBadReductionFiniteExtension}{67--236}
\leanref{GlobalTheory.\allowbreak{}GoodTransfer.\allowbreak{}inV\_\allowbreak{}of\_\allowbreak{}arc}{Development.\allowbreak{}AbelianVarietyTheory.\allowbreak{}Proofs.\allowbreak{}InertiaEigenvalues.\allowbreak{}CuspFormBoundGoodPointsExtendsAllPoints}{349--455}

Evaluation in this frame gives \(h(t)/(\pi t)^d\), with \(h\in V[[t]]\) and \(d\ge0\), by regularity of the section on the indicated localisation. We show separately that \(t^d\mid h\) and that \(\pi^d\mid h/t^d\).

Embed \(K_V\) into \(\C\) over the designated coefficient embedding. Such an abstract embedding exists since \(K_V\) is a finite extension of a local number field and its transcendence degree over the coefficient field is at most that of \(\C\). Apply it coefficientwise to the formal arc. In the complex uniformisation put \(t=\exp(2\pi i z)\), with the period of \(z\) rescaled if the marking requires it. The Laurent expansion of the evaluated section is its expansion in this parameter: the torus characters become the corresponding exponential monomials, the marked roots have the chosen exponents, and the constant tangent frame is their frame under the exponential map. The theta series of the period lattice converges normally for sufficiently large \(\Im z\), by positivity of \(Q\) on its nonzero lattice. Thus this identification uses the holomorphic arc attached to the degeneration; it does not require the abstract embedding to preserve a valued-field limit.

Here are the bounds on this arc. The imaginary part \(Y(z)\) of the normalised period matrix is affine in \(\Im z\), with positive semidefinite slope. Indeed the period is affine in \(z\), and positivity for every \(\Re z\) and every sufficiently large \(\Im z\), tested on each real vector, kills the \(\Re z\) coefficient and makes the other slope nonnegative. Fix \(z_0\) in this range. Then \(Y(z)\ge Y(z_0)>0\) for \(\Im z\ge\Im z_0\). The squared norm of a constant tangent vector in period coordinates is given by \(Y(z)^{-1}\) on its real and imaginary parts, so it is bounded by the value for \(Y(z_0)^{-1}\). In an orthonormal flag frame this bounds each diagonal tangent multiplier, since its squared absolute value times the fixed positive norm of the basis vector is a summand in that squared norm. For the theta generator the squared mean norm is
\[
             \frac{|\kappa|^2}{\sqrt{\,2^{g}\det Y(z)\,}},
\]
where \(g=\dim B\) and \(\kappa\) is the constant in the theta trivialisation of the arc. Integrating the squared theta series over a real period parallelepiped kills the terms with unequal Fourier indices. The remaining translates of the second half of the parallelepiped tile \(\R^g\), giving the Gaussian integral displayed above. This bounds the generator multiplier as well. The cusp section has bounded norm on unit frames. Its value in the arc frame is obtained by multiplying by the indicated diagonal factors to the nonnegative exponents of its dominant weight. It is therefore bounded as \(t\to0\). Cauchy's coefficient formula on circles tending to zero kills each negative Laurent coefficient. Injectivity of the coefficient embedding gives \(h=t^dh_1\) in \(V[[t]]\).
\leanref{GlobalTheory.\allowbreak{}ArcLieBound.\allowbreak{}norm\_\allowbreak{}diag\_\allowbreak{}le}{Development.\allowbreak{}AbelianVarietyTheory.\allowbreak{}Proofs.\allowbreak{}InertiaEigenvalues.\allowbreak{}NormDiagTangentUnitFrameHolomorphicArcLe}{293--308}
\leanref{GlobalTheory.\allowbreak{}ArcGenBound.\allowbreak{}gen\_\allowbreak{}ratio\_\allowbreak{}bounded}{Development.\allowbreak{}AbelianVarietyTheory.\allowbreak{}Proofs.\allowbreak{}InertiaEigenvalues.\allowbreak{}ArcGeneratorRatioBoundedUnitFrame}{74--100}
\leanref{GlobalTheory.\allowbreak{}GoodTransfer.\allowbreak{}coeff\_\allowbreak{}neg\_\allowbreak{}eq\_\allowbreak{}zero}{Development.\allowbreak{}AbelianVarietyTheory.\allowbreak{}Proofs.\allowbreak{}InertiaEigenvalues.\allowbreak{}CuspFormBoundGoodPointsExtendsAllPoints}{198--229}

For the second divisibility take the completion \(W\) of \(V[[t]]_{(\pi)}\). It is a discrete valuation ring with uniformiser \(\pi\), and \(t\) is a unit. An element of \(V[[t]]\) has order there equal to the minimum of the orders of its coefficients. The frame over \(W\) lies in the abelian locus. Its residue field can be infinite, so we explain why the good-point hypothesis applies. Use an affine chart of the framed abelian locus on which the integral frame coordinates are integral. Such charts are obtained by adjoining coordinates of the Lie flag and the generator of the invertible direct image of the theta line. At a valued point one takes a unit flag minor and a generating entry of the line. Refining on the abelian locus ensures that all integral points of the chart used for the test have good reduction.

On this chart, a failure of integrality persists at a test ring with finite residue field. To see this, take the closure of the graph of the evaluated rational function in the chart times \(\mathbb P^1\), over \(V\). A valued point with negative order specialises to infinity. In the closure of its centre take a closed point in the same chart at infinity. Its residue field is finite. Work on the horizontal irreducible component through which the valued point factors. In the local domain at the closed point choose a nonzero element \(b\) cutting out the complement of the required generic open, including the denominator of the chosen affine chart at infinity. Complete a nonzero multiple of \(\pi b\) to a system of parameters. A minimal prime over the other parameters avoids \(\pi b\): otherwise the parameter ideal would have height at most the number of those other parameters. The resulting one-dimensional quotient has a horizontal branch. Normalising and completing that branch gives a complete discrete valuation ring with finite residue field, generic point in the required open and closed point at infinity. The evaluated function has negative order there, contrary to the hypothesis. This proves the good-point test for the valued point of the chart over \(W\).
\leanref{GlobalTheory.\allowbreak{}GoodTransfer.\allowbreak{}inV\_\allowbreak{}of\_\allowbreak{}extIn}{Development.\allowbreak{}AbelianVarietyTheory.\allowbreak{}Proofs.\allowbreak{}InertiaEigenvalues.\allowbreak{}CuspFormBoundGoodPointsExtendsAllPoints}{120--150}
\leanref{GlobalTheory.\allowbreak{}ChartIntegral.\allowbreak{}integral\_\allowbreak{}of\_\allowbreak{}integral\_\allowbreak{}at\_\allowbreak{}algebraic\_\allowbreak{}points}{Development.\allowbreak{}AbelianVarietyTheory.\allowbreak{}Proofs.\allowbreak{}InertiaEigenvalues.\allowbreak{}IntegralOfIntegralAtAlgebraicPointsSmooth}{37--131}

It follows that \(h_1/\pi^d\in W\). Every coefficient of \(h_1\) is divisible by \(\pi^d\). The original value at \(t=\pi\) is now the evaluation of an element of \(V[[t]]\), hence belongs to \(V\). The triangular unit change does not alter integrality. Finally divide the order by the ramification index of the initial extension.
\leanref{GlobalTheory.\allowbreak{}GaussPoint.\allowbreak{}exists\_\allowbreak{}gaussPoint}{Development.\allowbreak{}AbelianVarietyTheory.\allowbreak{}Proofs.\allowbreak{}InertiaEigenvalues.\allowbreak{}ExistsGaussPoint}{64--108}
\end{proof}

\subsection*{Proof of \ref{B8.43o}}\label{prf:B8.43o}\phantomsection
\begin{proof}
Write \(M_0=g_0g_1\) and \(M=M_0d'\) for the base and cover levels. The uniformisation identifies the source with its standard self-dual lattice. The target rank \(r\) is odd, the source rank \(n\) is even, and the differential weight is dominant with positive final row degree in every nonempty block in the CM type. These are the signature and weight conditions on the standard layer. Write \(D\) for the total row degree.

We first fix the scalar. Use the coefficient basis of products of minors over the base field. On the pinned lowest-weight tensor the polynomial map in the Fock model is \(\kappa J\) in the standard letters, where \(J\) sends this basis to its prescribed harmonic polynomials. Indeed the normalised extremal vector spans its weight line. The basis change from the prescribed source frame to the standard one is unitary. The two isometric identifications of that line differ by a scalar of modulus one, giving \(|\kappa|=1\) on the whole coefficient map by linearity. The matrix of \(J\) on the minor basis is over the base field. A nonsingular coefficient minor expressing linear independence of the basis shows that \(Jv\) is over any embedded field containing the base field and the coefficients of \(v\).

Let \(\phi\) be the comparison of the vacuum with the geometric generator in the base frame, including the action of that unitary basis change on the vacuum line. The latter action preserves the norm; the base pin has unit modulus, so \(|\phi|=1\). To compare the component constants, evaluate the vacuum identity on a rational period vector where the analytic theta generator is nonzero. Such a vector exists since rational combinations of periods are dense and the generator is a nonzero holomorphic function. Take the same vector in each source column, using the inverse source basis matrix to prescribe its columns. The product of theta values is nonzero. The positive Gaussian and its reciprocal cancel in the two readings, so the ratio of the component constants equals \(\phi\) on every component met by the adelic frame.

The same \(\phi\) works on a cover. The vacuum identity pulls back with the same rigidified generator and the same periods. After division by \(\phi\) its component constants agree on pullbacks of base sections. To see that this determines the cover constants, one may use weight zero in the tangent variables and the prescribed even theta degree. On a geometric fibre \(B\) in any component, with principal theta line \(L\) and generator \(s_0\), let \(\pi:B\to B'=B/A\) be the quotient by the marked maximal isotropic subgroup \(A\subset B[M_0]\) and \(\hat\pi:B'\to B\) the complementary isogeny with \(\hat\pi\circ\pi=[M_0]\). Then \(\hat\pi^*L\) is algebraically equivalent to the \(M_0\)-th power of the descended principal polarisation on \(B'\), hence base-point free since \(M_0=g_0g_1\ge2\); the \(B'[M_0]\)-translates of \(\hat\pi^*s_0\) span \(H^0(B',\hat\pi^*L)\) by irreducibility of the theta-group action, so some \(a\in\hat\pi(B'[M_0])=A\) avoids the zero divisor of \(s_0\). Evaluating the \(n\)-th power of the translated theta generator at such a marked point \(a\in A\), as in Lemma~\ref{B8.43m} with \(P=1\), \(N=M_0\), and \(\Lambda=\mathcal O_E^n\), defines a base-level section of this weight over a finite Galois extension of the base coefficient field; taking Galois traces against a dual basis gives base-level sections over that field, and at least one coordinate has nonzero evaluation since extending scalars recovers the original section. Equivalently these are the level theta evaluations in the pinned section space. A real change of orthonormal frame moves the point of nonzero evaluation without changing its finite component. Apply both pin identities to the pullback there. The common nonzero evaluation cancels. This uses the pin for every dominant weight of the fixed theta degree. We take \(u=(\kappa\phi)^{-1}\).
\leanref{GlobalTheory.\allowbreak{}SerIntLetters.\allowbreak{}io\_\allowbreak{}eq\_\allowbreak{}phase\_\allowbreak{}smul\_\allowbreak{}harmLetters}{Development.\allowbreak{}AbelianVarietyTheory.\allowbreak{}Proofs.\allowbreak{}InertiaEigenvalues.\allowbreak{}HarmonicMapExtremalVectorEqPhaseSmulLetters}{22--81}
\leanref{GlobalTheory.\allowbreak{}LevelSeries.\allowbreak{}PinBase.\allowbreak{}baseCc\_\allowbreak{}of\_\allowbreak{}framePin}{Development.\allowbreak{}AbelianVarietyTheory.\allowbreak{}Proofs.\allowbreak{}InertiaEigenvalues.\allowbreak{}CoreConstantsEqVacuumIdentityUpToPhase}{168--255}
\leanref{GlobalTheory.\allowbreak{}LevelSeries.\allowbreak{}cc\_\allowbreak{}unique\_\allowbreak{}of\_\allowbreak{}pinned}{Development.\allowbreak{}AbelianVarietyTheory.\allowbreak{}Proofs.\allowbreak{}InertiaEigenvalues.\allowbreak{}CoverConstantsUniqueOnDictionaryComponents}{22--95}
\leanref{GlobalTheory.\allowbreak{}LevelSeries.\allowbreak{}covCc\_\allowbreak{}of\_\allowbreak{}pinned}{Development.\allowbreak{}AbelianVarietyTheory.\allowbreak{}Proofs.\allowbreak{}InertiaEigenvalues.\allowbreak{}VacuumIdentityConstantsCoverOfDominantPins}{19--33}

Fix \(\ell,s,L_0\). Apply the first assertion of Lemma~\ref{B8.43k} to \(J\) and \(L_0\); denote its exponent by \(c_1\). We may use \(\ell^{c_1}Jv\) as an integral polynomial at each chosen place. Denote its valuation ring in the designated embedding by \(O\). The coefficient field contains the \(M\)-th roots of unity, by the marking, and they are units. Lemma~\ref{B8.43j} gives the cost \(K\operatorname{ord}_\ell M_0\) for writing the finite test as an integral combination of vacuum translates.

The half-lattice in this expansion can differ from the half specified by the marking ideal. We record the cost of changing it. At a set \(T\) of wrongly oriented split primes dividing \(m\), use the central source element which scales the two halves by reciprocal powers of \(m\). Its determinant idèle on the selected factors is \(m^{-n}\). The scalar for rewriting the translates in the other half is
\[
 \chi(x)^{-1}\left(\prod_{P\in T}
               (\mathrm NP)^{\,\operatorname{ord}_P(m)nr}\right)^{-1/2},
 \qquad x_P=m^{-n}\ (P\in T),\quad x_P=1\ (P\notin T),
\]
where \(\chi\) is the reference splitting character on these idèles. Indeed, in each of the \(nr\) lattice coordinates one expands the characteristic function of a coset by its finite additive characters. The index in the display counts these characters. Unitary Fourier normalisation gives the inverse square root of the index, and the determinant action gives the first factor. Translating the characteristic function multiplies the Fourier coefficients by \(M\)-th roots of unity.

Make the change first on \(T\) over \(M_0\), then over \(d'\). The two supports are disjoint. For \(m=M_0\) there are finitely many choices of \(T\), hence finitely many scalars and central translates \(z_T\). These scalars belong to every tested coefficient field: choose a principal idèle with the specified local congruences at \(M_0\), use the character product formula, and combine the half norms at the remaining unramified places with the square root of the displayed index. The normalised finite values belong to the field by hypothesis, the remaining infinity factors are products of its CM embeddings, and \(r\) is odd. Clear these finitely many algebraic numbers as in Lemma~\ref{B8.43k}, at cost \(c_f\). For \(m=d'\) the character is unramified on \(T\). Put \(r=2s_0+1\) and \(a_P=n\operatorname{ord}_P(d')\). The factor at \(P\) is exactly
\[
       \bigl(\chi(P)\sqrt{\mathrm NP}\bigr)^{a_P}
                   (\mathrm NP)^{-(s_0+1)a_P}.
\]
It is integral at \(\ell\). For the unit assertion about the character, raise \(P\) to a principal power with generator \(\alpha\). Since \(P\nmid\ell\), all conjugates of \(\alpha\) are units above \(\ell\). The character product formula expresses a positive power of \(\chi(P)\) as a product of roots of unity from local units and powers of the normalised archimedean conjugates of \(\alpha\). Those conjugates and their absolute values are units at every extension of the valuation, since the squares of the absolute values are products of conjugates. The same holds for \(\sqrt{\mathrm NP}\). The normalised value lies in the tested field by hypothesis. The auxiliary central translate \(z'\) equals one above \(\ell\).
\leanref{GlobalTheory.\allowbreak{}LevelSeries.\allowbreak{}exists\_\allowbreak{}flip\_\allowbreak{}theta}{Development.\allowbreak{}AbelianVarietyTheory.\allowbreak{}Proofs.\allowbreak{}InertiaEigenvalues.\allowbreak{}ExistsCentralFlipThetaTranslateHalfLevel}{38--112}
\leanref{GlobalTheory.\allowbreak{}LevelSeries.\allowbreak{}flipScal\_\allowbreak{}mem\_\allowbreak{}valRing\_\allowbreak{}of\_\allowbreak{}unit}{Development.\allowbreak{}AbelianVarietyTheory.\allowbreak{}Proofs.\allowbreak{}InertiaEigenvalues.\allowbreak{}FlipScalarMemCoefficientRingAuxiliaryLevel}{124--228}

We describe the finite-point change because bounding a matrix chosen separately at every level would give the wrong quantifiers. Use the open subgroup fixing the standard lattice and all the translates occurring at level \(M\). In the double class of \(sz_Tz'\) choose a representative \(h\) supported at split primes away from \(2\ell M\). One construction is to approximate the prescribed finite matrix by a rational unitary matrix at the primes of the level and of its denominators. Weak approximation here follows by using the Cayley coordinates on the skew-Hermitian matrices, translating the open chart if necessary. To remove unwanted nonsplit denominators, write the rational unitary matrix as a product of unitary reflections in nonisotropic vectors. Such a factor acts by \(u\), \(u\bar u=1\), on its line and by one on the orthogonal complement. One obtains the factorisation by matching one vector of an orthogonal basis at a time: a reflection in the difference of two vectors of equal nonzero norm matches them, with an intermediate scalar multiple if their difference is isotropic. Write \(u=z/\bar z\). For each factor approximate \(z\) by weak approximation so as to retain the required congruences, and make \(u-1\) divisible by the denominator of the line projection at the other nonsplit primes in that denominator. At all remaining nonsplit primes a norm-one number is a unit and the projection is integral. Congruences may be imposed to greater precision by the finite denominators of the other factors in the product. Thus the remaining nonintegral factors away from the approximation set have split support. Put the entire remaining matrix in \(h\) at those split primes and put its integral, level-congruent part in the open subgroup. This gives
\[
                    sz_Tz'=\gamma\,h\,a_\infty k
\]
with \(\gamma\) rational, \(a_\infty\) the inverse infinite component of \(\gamma\), and \(k\) in the chosen stabiliser. At the primes over \(\ell\), the factors \(h,z'\) equal one and \(k,k^{-1}\) are integral. Thus one exponent \(d_0\) clears both \(\gamma,\gamma^{-1}\) there: take the maximum of the clearing exponents for \(sz_T,(sz_T)^{-1}\) over the finite list of \(T\). Rational left translation reindexes the theta sum. The action of \(a_\infty\) substitutes entries of \(\gamma\) and \(\gamma^{-1}\) in its degree-\(D\) polynomial and multiplies the vacuum by \(\prod_\tau\det(a_{\infty,\tau})^{k_\tau}\), where \(k_\tau\) is its fixed determinant weight. Consequently
\[
                 c_a=d_0\bigl(D+n\sum_\tau |k_\tau|\bigr)
\]
clears the transformed polynomial. This follows by expanding each determinant, using the inverse matrix for a negative power. The transformed polynomial is harmonic because the conjugate sets of variables undergo mutually contragredient substitutions.
\leanref{GlobalTheory.\allowbreak{}SplitRepsOpen.\allowbreak{}exists\_\allowbreak{}rep\_\allowbreak{}splitOffBad}{Development.\allowbreak{}AbelianVarietyTheory.\allowbreak{}Proofs.\allowbreak{}IwasawaDecomposition.\allowbreak{}TendstoSmulZeroOfCompact}{86--91}
\leanref{GlobalTheory.\allowbreak{}MoveDenom.\allowbreak{}denomAt\_\allowbreak{}of\_\allowbreak{}move}{Development.\allowbreak{}AbelianVarietyTheory.\allowbreak{}Proofs.\allowbreak{}InertiaEigenvalues.\allowbreak{}RationalPartDenominatorBoundedByPrimePower}{29--73}
\leanref{GlobalTheory.\allowbreak{}LevelSeries.\allowbreak{}ArchAct.\allowbreak{}real\_\allowbreak{}act\_\allowbreak{}harm}{Development.\allowbreak{}AbelianVarietyTheory.\allowbreak{}Proofs.\allowbreak{}InertiaEigenvalues.\allowbreak{}RealRationalPointActsHarmonicSubstitution}{50--98}

Take for \(\Lambda\) the conjugate of the lattice at \(h\). It is self-dual by unitarity and standard above \(2\ell M\). Clearing the finitely many elementary divisors of this lattice and its inverse away from \(2\ell M\) gives an odd \(u_0\) prime to \(\ell M\) for which \(N=Mu_0\) satisfies Lemma~\ref{B8.43m}. Every translate in the prescribed half modulo the vacuum lattice is represented by a marked point. To write one such translate as the mean over the \(|A|^n\) marking tuples, give its chosen representative coefficient \(|A|^n\) and give the others coefficient zero. The difference between the original representative and the marked one is a lattice vector in the same isotropic half; it fixes the vacuum and contributes trivial commutator. By addition we obtain an integral weight function for the mean.

Expand the products of the finitely many weights and polynomial coefficients in a basis of their span over the coefficient valuation ring. This span is finite and torsion free over a discrete valuation ring, hence free. Its basis elements belong to \(O\), and the coordinate polynomials satisfy the same harmonic and row-weight equations because these equations are linear over the coefficient field. For each coordinate use the geometric sum in Lemma~\ref{B8.43m}, with the weight indicating its marking tuple. It defines a section: after adjoining the finite torsion data the formula is natural under base change, and changes of the torsion enumeration permute the sum. Changing the division point within the designated primary part permutes the lattice translates with their theta-group actions. Thus the two pullbacks of the sum agree on the fibre product of the marking cover. Faithfully flat descent gives its value over each base test algebra. A triangular change of tangent basis acts by the differential weight and scaling the theta generator acts by its \(n\)-th power. These are the two defining laws of a section. At each good complex embedding the differential evaluation is a sum of theta derivatives with finite characteristics. In a nondefinite component its zero Fourier coefficient vanishes by the calculation of Lemma~\ref{B8.25t}: the number of columns is even, and positive last row degree supplies a full minor at an embedding in the CM type. The calculation there works for each characteristic individually and hence after a change of marking level. A definite component has no nonzero rational isotropic line. This proves cuspidality.
\leanref{GlobalTheory.\allowbreak{}ThetaSumDescent.\allowbreak{}exists\_\allowbreak{}cusp\_\allowbreak{}thetaSum}{Development.\allowbreak{}AbelianVarietyTheory.\allowbreak{}Proofs.\allowbreak{}InertiaEigenvalues.\allowbreak{}ExistsCuspCoefficient}{15--34}
\leanref{GlobalTheory.\allowbreak{}LevelSeries.\allowbreak{}SerId.\allowbreak{}exists\_\allowbreak{}forms\_\allowbreak{}of\_\allowbreak{}transSum}{Development.\allowbreak{}AbelianVarietyTheory.\allowbreak{}Proofs.\allowbreak{}InertiaEigenvalues.\allowbreak{}ExistsFormsThetaShortSeriesIdentity}{231--283}

For clarity, the scalar in the reading of this sum can be computed before taking valuations. Put \(J_0=[\Lambda:\Lambda\cap\mathcal O_E^n]\), and let \(y\) be the determinant idèle of \(h\) on its selected split factors. On the geometric side the unweighted count of lattice tuples is \([\Lambda:N\mathcal O_E^n]^r=I^rJ_0^r\); the short-section mean divides by \(I^r|A|^n\). On the Weil side the self-dual lattice change has factor \(\chi(y)^{-1}\sqrt{J_0^r}\), with the same marking mean. Applying a homogeneous harmonic differential to the Gaussian gives the common framed Lie factor and the sign \((-1)^D\). The ratio of the two readings, after removal of the pinned scalars, is therefore \((-1)^D\chi(y)^{-1}/\sqrt{J_0^r}\). To see it is in \(O\), take elementary divisors at each pair of conjugate split places. If the exponents on one side are \(b_i\), those on the other side are \(-b_i\). The exponent in \(J_0\) over this pair is \(\sum_i|b_i|\), congruent to \(\sum_i b_i\) modulo two. Thus \(J_0=(\mathrm N\mathfrak a)t_0^2\), where the determinant ideal is represented by \(\mathfrak a/\bar{\mathfrak a}\) with relatively prime numerator and denominator. All factors are prime to \(\ell\). Pairing the square roots at these primes with the normalised character values gives products of the units proved above and rational powers with integral exponents of primes different from \(\ell\).
\leanref{GlobalTheory.\allowbreak{}LevelSeries.\allowbreak{}SerId.\allowbreak{}thetaSum\_\allowbreak{}value\_\allowbreak{}emb}{Development.\allowbreak{}AbelianVarietyTheory.\allowbreak{}Proofs.\allowbreak{}InertiaEigenvalues.\allowbreak{}ExistsFormsThetaShortSeriesIdentity}{35--170}
\leanref{GlobalTheory.\allowbreak{}LevelSeries.\allowbreak{}clsScal\_\allowbreak{}mem\_\allowbreak{}valRing}{Development.\allowbreak{}AbelianVarietyTheory.\allowbreak{}Proofs.\allowbreak{}InertiaEigenvalues.\allowbreak{}SplitSupportedClassScalarMemValuationRing}{268--393}

By Lemma~\ref{B8.43m} the geometric sum multiplied by \(I^r|A|^n(2\Delta N)^D\) is integral at good frames. Lemma~\ref{B8.43n} extends this to all integral test frames. Its prime-to-\(\ell\) level condition holds since at least one of the distinct odd primes \(g_0,g_1\) remains after removing the \(\ell\)-part of \(M\). Apply the second assertion of Lemma~\ref{B8.43k} with this \(N,\Lambda\). Its exponent \(c_3\) suffices at \(\ell\), since the complementary factor of the denominator is a unit there. We have obtained the single exponent
\[
            c=c_1+K\operatorname{ord}_\ell(M_0)+c_f+c_a+c_3 .
\]
Every term was chosen before the cover and the place on it.

The same finite expansion over the embedded coefficient field gives a rational section; here the valuation ring in the coefficient expansion can be replaced by the whole field. At a set \(S\) of places use this rational section and the local expansions just obtained. They have the same section at every place: the readings agree in all adelic frames, orthonormal frames meet every flag up to triangular change, and evaluation on geometric points detects a section on a reduced test algebra. Its values are therefore integral simultaneously at \(S\). This argument includes \(S=\varnothing\). For a permitted determinant twist the multiplier on a component is a root of unity in the cover field. Indeed it factors through the finite norm-one ray quotient by rational elements and level congruences. The full ray field in the hypotheses defines the marked components over the cover field, so multiplication component by component preserves the rational section and every tested integral value.

For the rationality assertion the good cover has odd squarefree split auxiliary level, the indicated pins, the full ray field and the values of the two reference characters in their respective weights \(0\) and \(1/2\). Their quotient gives the normalised values used in the construction. Choose a prime larger than \(d'\), use the finite expansion over the whole coefficient field at this prime, and divide back by its nonzero clearing power in that field. The determinant character of an allowed finite-order splitting twist \(\delta\) on a unitary group \(U\) is \(d_{\delta,U}(g)=\delta(z)\) for \(z/\bar z=\det g\). This is independent of \(z\) since \(\delta\) is trivial on fixed-field idèles. Changing both splittings multiplies the series by \(d_{\delta,V}(s)^{-1}d_{\delta,W}(g)\). The first factor is in the cover field by the value condition on \(\delta\); the second is the component multiplier just considered. This proves the assertion for both rational twist laws.
\leanref{GlobalTheory.\allowbreak{}SerStd.\allowbreak{}untwRatS}{Development.\allowbreak{}AbelianVarietyTheory.\allowbreak{}Proofs.\allowbreak{}IwasawaDecomposition.\allowbreak{}UntwistedThetaRationalOnGoodCovers}{61--171}
\leanref{GlobalTheory.\allowbreak{}SerStdTwist.\allowbreak{}seriesRatS\_\allowbreak{}of\_\allowbreak{}untwRatS}{Development.\allowbreak{}AbelianVarietyTheory.\allowbreak{}Proofs.\allowbreak{}IwasawaDecomposition.\allowbreak{}ThetaSeriesRationalOnGoodCoversOfUntwisted}{18--32}
\leanref{GlobalTheory.\allowbreak{}TwistInt.\allowbreak{}mul\_\allowbreak{}twist\_\allowbreak{}mem\_\allowbreak{}intForms}{Development.\allowbreak{}AbelianVarietyTheory.\allowbreak{}Proofs.\allowbreak{}InertiaEigenvalues.\allowbreak{}TwistDeterminantMulIntegralFormsRayClassCover}{26--83}
\end{proof}

\subsection*{Proof of \ref{B8.43p}}\label{prf:B8.43p}\phantomsection
\begin{proof}
After coefficient extension to \(L\), forgetting the additional marking is finite étale on the generic fibre. Two lifts of a framed geometric point differ by the change of marking from the base compact subgroup. The determinant factor in the base reading is invariant under that subgroup, and the pin identifies readings of pullbacks. Thus invariance makes the two pullbacks of the section equal on the fibre product, first on geometric frames and then on reduced test algebras. It descends over \(L\) at base level. The cusp condition descends too: the Fourier expansion at a lifted isotropic line is the original expansion with its period group restricted to a finite-index subgroup, so the constant term of a pullback is the pullback of the constant term. Scalar extension of the space of sections commutes with this descent, by the equaliser for sections on an affine cover and flatness over \(K\). The same conclusion for the cusp subspace follows by imposing the vanishing constant terms at all conjugate embeddings. We therefore have an \(L\)-linear combination of base cusp sections.
\leanref{GlobalTheory.\allowbreak{}CoverDescRat.\allowbreak{}cover\_\allowbreak{}desc\_\allowbreak{}rat}{Development.\allowbreak{}AbelianVarietyTheory.\allowbreak{}Proofs.\allowbreak{}InertiaEigenvalues.\allowbreak{}CoverFormRationalDescentToBase}{55--65}

Fix \(w\in S\) and put \(v=w|_K\). Choose an \(\mathcal O_{K,v}\)-basis \(f_1,\ldots,f_m\) of the base integral cusp lattice. This is the lattice defined by evaluation on all algebraic integral frames. To recall the lattice argument, finite dimensionality and separating evaluations bound it inside a finite free module: choose finitely many evaluations of full rank over an extension and invert a nonzero minor. Each cusp section has an integral multiple, by clearing denominators on finitely many integral flag charts and extending across the boundary as in Lemma~\ref{B8.43n}. Thus the lattice spans the cusp space, and a finitely generated torsion-free module over the local discrete valuation ring is free. Write \(F=\sum_j y_j f_j\), \(y_j\in L\), using pullbacks in this equality.

At each \(w'\mid v\) different from \(w\), choose \(b\) large enough that \((w')^b y_j\) is integral there for all \(j\). For \(a\in\mathcal O_L\) in all these powers we claim that
\[
                      \sum_j\operatorname{Tr}_{L/K}(a y_j) f_j
\]
is in the base integral lattice. Test a base integral frame \(x\) over a complete discrete valuation ring over \(v\) with finite residue field. In a finite normal extension of its fraction field split all the embeddings of \(L\) over \(K\). For an embedding inducing \(w\), lift \(x\) after a further finite extension to the marking cover. Such a lift exists because the generic marking cover is surjective and finite, so its fibre has a point over a finite extension. Pull back the tangent basis and the generator through the isomorphism of the underlying polarised objects. These are integral: the semiabelian model, its Lie lattice and its rigidified generator are carried along the same base change, and forgetting the marking recovers the original frame. Evaluation of the cover section at this frame is the image of \(\sum_j y_j f_j(x)\), hence integral. Multiplication by \(a\) preserves integrality. For an embedding inducing \(w'\ne w\), each individual product \(a y_j f_j(x)\) is integral by the choice of \(b\). Summing over the embeddings gives the asserted integral value of the trace section. Positive ramification indices give integrality back on the original test ring.
\leanref{GlobalTheory.\allowbreak{}CoverPoint.\allowbreak{}exists\_\allowbreak{}ext\_\allowbreak{}lift\_\allowbreak{}to\_\allowbreak{}cover}{Development.\allowbreak{}AbelianVarietyTheory.\allowbreak{}Proofs.\allowbreak{}RingsOfIntegers.\allowbreak{}ExistsFiniteExtensionIntegralFramedPointLiftCove}{79--124}
\leanref{GlobalTheory.\allowbreak{}CoverLiftVal.\allowbreak{}le\_\allowbreak{}one\_\allowbreak{}at\_\allowbreak{}lift}{Development.\allowbreak{}AbelianVarietyTheory.\allowbreak{}Proofs.\allowbreak{}InertiaEigenvalues.\allowbreak{}BaseSectionsValuesIntegralAtCoverLift}{31--74}
\leanref{GlobalTheory.\allowbreak{}TraceIdem.\allowbreak{}trace\_\allowbreak{}comb\_\allowbreak{}integral}{Development.\allowbreak{}AbelianVarietyTheory.\allowbreak{}Proofs.\allowbreak{}InertiaEigenvalues.\allowbreak{}TraceCombinationCoverFormIntegralOverPrime}{37--56}

Coordinates in the chosen integral basis give \(\operatorname{Tr}_{L/K}(a y_j)\in\mathcal O_{K,v}\) for every \(j\). These are exactly the tests in Lemma~\ref{B8.43l}. In the unramified case it gives \(y_j\in\mathcal O_{L,w}\), taking \(O=\mathcal O_{L,w}\) in that lemma. In the fixed-ring case it gives \(\ell^{e^2}y_j\in\mathcal O_{L,w}\).

Finally \(S\) is finite. In the decomposition for \(w\) multiply each base form by a base integer which is a unit at \(v\) and sufficiently divisible at the other places below \(S\), and divide its coefficient by that integer. The resulting base forms are integral below all of \(S\). Denote the finite set of new coefficients in this decomposition by \(Y_w\). Choose a common exponent \(b'\) clearing the negative orders of the elements of all the \(Y_w\) at all the other places. The Chinese remainder theorem gives \(d_w\in\mathcal O_L\) equal to one modulo \(w^{b'+1}\) and to zero modulo \((w')^{b'+1}\) for \(w'\in S\setminus\{w\}\). Their sum is a unit at every place in \(S\). Multiply the decomposition for \(w\) by \(d_w/\sum_{w'}d_{w'}\) and add. The multipliers sum to one and all the resulting coefficients are integral at \(S\). For empty \(S\) the rational decomposition suffices.
\leanref{GlobalTheory.\allowbreak{}DescGlue.\allowbreak{}glue}{Development.\allowbreak{}AbelianVarietyTheory.\allowbreak{}Proofs.\allowbreak{}InertiaEigenvalues.\allowbreak{}GlueIntegralDecompositionsOverRationalPrime}{41--103}
\end{proof}

\subsection*{Proof of \ref{B8.43q}}\label{prf:B8.43q}\phantomsection
\begin{proof}
Every kernel is a subgroup of \(B[p^2]\), so their number is at most \(2^{(p^2)^{2g}}\). The Weil pairing is perfect and the stipulated isotropic order is the square root of \(|B[p^2]|\). The polarisation multiplied by \(p^2\) therefore descends to a principal polarisation on \(B/C\): its kernel on the quotient is \(C^\perp/C=0\). The endomorphisms descend by stability. On the marked torsion use \(p^{-1}\Psi\), where the inverse of \(p\) is taken modulo the marking level. Its Weil multiplier is one there. This gives the unique marked quotient by the moduli property. In characteristic zero the quotient is étale. Its differential is consequently an isomorphism, commutes with \(\mathcal O_E\), and gives the stated basis in each eigenspace.

We include the line-bundle calculation. Write \(\mathcal P_L\) for the biextension \(m^*L\otimes\mathrm{pr}_1^*L^{-1}\otimes\mathrm{pr}_2^*L^{-1}\otimes L(0)\), where \(m\) is addition. The principal line in this moduli problem has class \((1,w_0)^*\mathcal P_L\). The Weil-pairing equality gives \((\Psi,\Psi)^*\mathcal P_{L'}=\mathcal P_L^{p^2}\): these rigidified biextensions represent the polarisation homomorphisms, and the homomorphisms agree on the prime-to-\(p\) torsion by the pairing and hence on the abelian variety. Restrict along \((1,w_0)\) and use \(\Psi w_0=w_0\Psi\). This gives \(\Psi^*L'\simeq L^{p^2}\). Normalisation at the identity makes the isomorphism unique since the invertible functions on a proper geometrically connected abelian variety are constant. The same normalisation gives \([p]^*L\simeq L^{p^2}\).
\leanref{GlobalTheory.\allowbreak{}ThetaTransport.\allowbreak{}exists\_\allowbreak{}isoGen}{Development.\allowbreak{}AbelianVarietyTheory.\allowbreak{}Proofs.\allowbreak{}TwistedOperators.\allowbreak{}HeckeIsogenyThetaGeneratorTransport}{178--242}

Let \(v\) denote the ratio of the two isomorphisms from the multiplication and quotient pullbacks to \(L^{p^2}\), in local frames. For \(a\in C\cap B[p]\), the ratio \(v(x+a)/v(x)\) is a constant \(u_a\), since both pullbacks are unchanged and the ratio is an invertible function on \(B\). The same holds for \(w_0a\). The primitive line identity used above, written on a point \(x\), involves
\[
             L((1+w_0)x)\,L(0)\,L(x)^{-2}\,L(w_0x)^{-1}.
\]
Use its trivialisation on the two targets and take the quotient. The resulting invertible function on \(B\) is constant. Comparing it at \(x+a\) and \(x\) gives \(u_a u_{w_0a}=u_a^2 u_{w_0a}\), hence \(u_a=1\). Thus each summand depends only on the class of \(c\). Translation by any element of \(C\) permutes the summands. The sum descends uniquely, since \(\Psi\) is faithfully flat and its fibre product with itself is the disjoint union of the translations by \(C\).
\leanref{GlobalTheory.\allowbreak{}ThetaTransport.\allowbreak{}quot\_\allowbreak{}transl}{Development.\allowbreak{}AbelianVarietyTheory.\allowbreak{}Proofs.\allowbreak{}TwistedOperators.\allowbreak{}HeckeIsogenyThetaGeneratorTransport}{346--417}
\leanref{GlobalTheory.\allowbreak{}ThetaTransport.\allowbreak{}exists\_\allowbreak{}trace}{Development.\allowbreak{}AbelianVarietyTheory.\allowbreak{}Proofs.\allowbreak{}TwistedOperators.\allowbreak{}HeckeIsogenyThetaGeneratorTransport}{619--659}

To prove nonvanishing, translate the pullback sum by \(b\in B[p]\). Relative to the unshifted frame, the summand indexed by \(c\) has multiplier proportional to a character of \(b\). Indeed the quotient of the two translated trivialisations is a trivialisation on \([p]^*(t_{pc}^*L\otimes L^{-1})\), up to taking its dual. Translation by \(b\) changes it by an invertible constant, and composition gives the character law. If the characters for \(c,c'\) agree, the quotient of these trivialisations descends through \([p]\). It trivialises \(t_{p(c-c')}^*L\otimes L^{-1}\) up to translation. The principal polarisation then gives \(p(c-c')=0\). The characters indexed by distinct classes are therefore distinct. If the sum vanished identically, translating by all \(b\in B[p]\) and using character orthogonality would force every summand to vanish. Multiplication by \(p\) and translation are surjective and \(s\ne0\). This is a contradiction. Since \(H^0(B',L')\) has dimension one, the descended section is a generator.
\leanref{GlobalTheory.\allowbreak{}ThetaTransport.\allowbreak{}transl\_\allowbreak{}quot}{Development.\allowbreak{}AbelianVarietyTheory.\allowbreak{}Proofs.\allowbreak{}TwistedOperators.\allowbreak{}HeckeIsogenyThetaGeneratorTransport}{496--554}
\leanref{GlobalTheory.\allowbreak{}ThetaTransport.\allowbreak{}eq\_\allowbreak{}zero\_\allowbreak{}of\_\allowbreak{}trace}{Development.\allowbreak{}AbelianVarietyTheory.\allowbreak{}Proofs.\allowbreak{}TwistedOperators.\allowbreak{}HeckeIsogenyThetaGeneratorTransport}{576--617}
\end{proof}

\subsection*{Proof of \ref{B8.43r}}\label{prf:B8.43r}\phantomsection
\begin{proof}
Enlarge the field to define the kernel, its framed quotient, the \(p^2\)-torsion of the source and the \(p^3\)-torsion of the target. These are finite extensions since the marking, morphisms, and frames are of finite presentation and the points used are algebraic. The isogeny and the complementary isogeny \(\Psi'\), with \(\Psi'\Psi=[p^2]\) and \(\Psi\Psi'=[p^2]\), extend on the semiabelian models after semistable reduction. On the identity components this follows by extending the homomorphisms to the smooth group models; their restrictions to the generic fibre determine their compositions by separatedness. Their tangent maps thus have composites \(p^2\). As \(p\in V^\times\), \(p^{-1}d\Psi\) identifies the two integral Lie lattices and all their eigenspaces. Use this basis on the target. For good reduction the quotient extends directly as the quotient of the abelian scheme by the finite étale subgroup of its prime-to-residue-characteristic torsion.

Write \(s'=a s_0\) relative to an integral theta generator on the target. We prove both valuation inequalities. In a split semiabelian uniformisation write \(Q\) for the quadratic period order of the character lattice. For a section of \(L^k\) the integral Fourier bound used here requires its coefficient of character \(x\) to belong to the abelian-part integral section lattice with order at least \(Q(x)/k\). The framed generator has this bound at \(k=1\): its constant coefficient is integral and the period equation gives the bound at \(x\). A multiplication pullback by \([m]\) moves character \(x\) to \(mx\) and degree \(k\) to \(m^2 k\), preserving \(Q(x)/k\). A symmetric torsion translation shifts the character by a period label \(y\); in degree \(k\) its linear toric multiplier and its rigidifying multiplier supply
\[
             \frac{B(x,y)}k+\frac{Q(y)}k
                    =\frac{Q(x+y)-Q(x)}k .
\]
Indeed the \(k\)-th power of the torus evaluation is its period pairing, and symmetry determines the order of the rigidifying factor by its square, with \(B(y,y)=2Q(y)\). Ambiguities in the finite-order lift are roots of unity. Translation on the abelian part preserves its integral lattice by properness. Thus each summand of the pulled-back trace extends regularly. In the formal torus chart this criterion for regularity follows by writing the section in the character grading: the period equation gives all coefficients from the finitely many classes modulo \(kX\), and the bound makes every coefficient integral. The positive quadratic growth gives convergence in the completed chart. Descent from a finite extension follows by testing the coefficients in the original fraction field.
\leanref{GlobalTheory.\allowbreak{}ReachRegular.\allowbreak{}regular\_\allowbreak{}of\_\allowbreak{}reach}{Development.\allowbreak{}AbelianVarietyTheory.\allowbreak{}Proofs.\allowbreak{}RingsOfIntegers.\allowbreak{}ReachedSectionRegularOnIntegralModel}{71--100}
\leanref{GlobalTheory.\allowbreak{}PlaceInt.\allowbreak{}nuGe\_\allowbreak{}zero\_\allowbreak{}of\_\allowbreak{}moves}{Development.\allowbreak{}AbelianVarietyTheory.\allowbreak{}Proofs.\allowbreak{}RingsOfIntegers.\allowbreak{}SymmetricTorsionTranslateSectionNuNonneg}{34--57}

The regular pullback of \(a s_0\) implies \(v(a)\ge0\). In detail, the reduction of the primitive section \(s_0\) is nonzero somewhere on the special fibre. The extended isogeny is surjective there since its composite with the complementary isogeny is multiplication by \(p^2\), which is surjective on the semiabelian fibre. At a point above the nonvanishing locus, the rigidified pullback comparison is a unit. This last assertion can also be seen by comparing local generators: a generic isomorphism of the two lines has a vertical divisor on the smooth model with connected special fibre, and its order along that fibre is zero by the rigidification at the identity. Thus a pole of \(a\) would give a pole of the regular pullback.

For the opposite inequality put \(\Psi_2=[p]\Psi'\) and \(D_2=\ker\Psi_2\). It is a subgroup of \(B'[p^3]\), so \(|D_2|=p^m\) for some \(m\). Translate the multiplication pullback \([p^2]^*s_0\) by all elements of \(D_2\), using the descent lifts for \(\Psi_2\), and sum. Substitution of the forward trace of Lemma~\ref{B8.43q} into this sum gives the pullback of \(b s\) for a scalar satisfying
\[
                             ab=p^m\zeta,\qquad \zeta^p=1 .
\]
To see the cancellation, for each forward representative \(c\) the inner sum over \(d\in D_2\) has character given by the theta-group commutator of \(\Psi'd\) and \(c\) for \(L^{p^2}\). The image of \(D_2\) under \(\Psi'\) is \(B[p]\), by surjectivity of \(\Psi'\). The annihilator of \(B[p]\) in \(B[p^2]\) under the perfect \(p^2\)-pairing is \(B[p]\). Thus every representative with \(pc\ne0\) gives zero. On the identity class the sum has \(|D_2|\) identical terms, up to the ratio of the two lifts on a \(p\)-torsion representative. Its \(p\)-th power is one, by iteration of translation. This proves the displayed identity. The reverse sum is regular by the same coefficient calculation, now for \([p^2]\). The nonvanishing test for its primitive target gives \(v(b)\ge0\). Since \(v(p)=v(\zeta)=0\), we obtain \(v(a)=0\). Splitting the two models in different finite extensions gives the same conclusion because each ramification index is positive.
\leanref{GlobalTheory.\allowbreak{}TraceOrth.\allowbreak{}thetaTrace\_\allowbreak{}orthogonality}{Development.\allowbreak{}AbelianVarietyTheory.\allowbreak{}Proofs.\allowbreak{}TwistedOperators.\allowbreak{}ThetaTraceOrthogonality}{538--623}
\leanref{GlobalTheory.\allowbreak{}HkIntegral.\allowbreak{}theta\_\allowbreak{}unit}{Development.\allowbreak{}AbelianVarietyTheory.\allowbreak{}Proofs.\allowbreak{}TwistedOperators.\allowbreak{}IsUnitNatCastUp}{242--270}

Finally the frame character of any section is a product of powers of the diagonal units. Its integral value in the integral target frame remains integral in the transported one. It lies in the valuation ring of a finite extension of \(\operatorname{Frac}V\), whose elements are integral over \(V\). This proves the value assertion on both classes of test points.
\end{proof}

\subsection*{Proof of \ref{B8.43s}}\label{prf:B8.43s}\phantomsection
\begin{proof}
Over a reduced coefficient test algebra \(R\), listing the geometric kernels defines a finite étale correspondence over the source. Indeed the kernels are subgroups of the finite étale \(p^2\)-torsion. After a finite étale trivialisation of that torsion their defining conditions in Lemma~\ref{B8.43q} select a constant finite set. Let \(B_R\) be the resulting finite étale algebra. On its spectrum the universal subgroup has a quotient abelian scheme with the transported marking. Étale locally list its elements and locally choose a generator \(s_0\) on the target. If the list has length \(m\), summing over the full list repeats the trace over the quotient by \(C\cap B[p]\) exactly \(m/q_P^i\) times, since \(|C|=|B[p]|\). This positive integer is invertible over the coefficient field. The quotient of the full sum by the pullback of \(s_0\) is a regular scalar times this integer. Indeed, on the locus where that pullback is nonzero one may divide regular functions; this smooth locus meets each geometric fibre since \(s_0\) generates the theta line and the isogeny is surjective. On pairs of points above the same base point the quotients agree by Lemma~\ref{B8.43q}. They agree as functions on the smooth fibre product by reducedness. Faithfully flat descent gives the regular scalar on the base. Dividing by \(m/q_P^i\) gives the scalar \(a_0\) of the transported generator.
\leanref{GlobalTheory.\allowbreak{}HeckeScal.\allowbreak{}exists\_\allowbreak{}scal\_\allowbreak{}of\_\allowbreak{}family}{Development.\allowbreak{}AbelianVarietyTheory.\allowbreak{}Proofs.\allowbreak{}TwistedOperators.\allowbreak{}ExistsScalarHeckeTraceThetaGeneratorFamily}{224--277}

Use \(p^{-1}d\Psi\) on the Lie frame. The local expression for a term of the Hecke sum is \(a_0^n f(B',p^{-1}d\Psi\,\mathbf u,s_0)\). A change of \(s_0\) cancels between the two factors, and relisting the kernel leaves \(a_0\) fixed. Hence the term descends to an element \(u_f\in B_R\). Its trace to \(R\) specialises to the sum of its values at all the geometric kernels, since over an algebraic closure a finite étale algebra is the product of one copy of the field for each point. This trace commutes with base change. A triangular change of the original frame induces the same Lie-basis change at each target and scales the transported generator by the same scalar, by linearity of its defining trace. The trace consequently has exactly the section weight. Multiplying by \(bq_P^{-in/2}\) gives the asserted section. Equality on geometric frames implies uniqueness since the test rings are reduced. At \(\lambda\nmid p\), \(q_P\) is a unit. Lemma~\ref{B8.43r} gives integrality of every term over the test ring. The sum with its integral multiplier belongs to the original fraction field and is integral there, so belongs to the test ring.
\leanref{GlobalTheory.\allowbreak{}HeckeEx.\allowbreak{}exists\_\allowbreak{}sum}{Development.\allowbreak{}AbelianVarietyTheory.\allowbreak{}Proofs.\allowbreak{}TwistedOperators.\allowbreak{}HeckeTransformOfPointLaws}{194--223}

Fix a complex embedding in the cusp test. At unit frames normalise the theta generators to have squared mean norm one on their complex tori. If \(t_C\) is the theta scalar of the transported frame relative to the target unit frame, then \(|t_C|\le q_P^i\). Indeed let \(H\) be the linear map of universal covers inducing \(\Psi\), and \(u,u'\) the pointwise Gaussian norms of the two unit generators. The pullback identity of Lemma~\ref{B8.43q} and the relation \(H^*H_{B'}=p^2 H_B\) between their polarisation metrics give
\[
              |t_C|u'(Hz)\le\sum_{j=1}^{q_P^i}u(pz+v_j)
\]
for suitable translates \(v_j\). The rigidified line comparison is an isometry for these metrics: its curvature vanishes on the quotient of the two lines and its norm at the identity is one. Square the inequality and use Cauchy--Schwarz with \(q_P^i\) summands. Each translated square on the right has mean one, since multiplication by \(p\) and translation preserve normalised Haar measure. The square on the left has mean \(|t_C|^2\), since the quotient isogeny also pushes Haar probability to Haar probability. This gives the bound.
\leanref{GlobalTheory.\allowbreak{}PeriodicMean.\allowbreak{}mean\_\allowbreak{}factor\_\allowbreak{}le}{Development.\allowbreak{}AbelianVarietyTheory.\allowbreak{}Proofs.\allowbreak{}TwistedOperators.\allowbreak{}PeriodicThetaNormMeanFactorLe}{9--75}

Write \(\Lambda_j\) for the period lattices of the finitely many components. The map on periods of \(\Psi\) is a similitude of multiplier \(p^2\) and its image contains \(p^2\Lambda_{j'}\), by the complementary isogeny. Dividing its matrix by \(p\) gives a rational isometry \(A\) with \(pA\Lambda_j\subseteq\Lambda_{j'}\) and \(pA^{-1}\Lambda_{j'}\subseteq\Lambda_j\). Indeed the equality of the alternating trace forms, tested also after multiplying the first vector by each element of \(E\), recovers the Hermitian form by nondegeneracy of the field trace. The two possible signs of the purely imaginary polarisation parameter are distinguished by positivity on the complex tangent space: the two Riemann forms there are positive and \(H\) is complex linear. The period relation carries the flag to the target flag. Between unit flag frames its triangular matrix has diagonal entries of modulus one, by orthonormality on each successive quotient. Its weight character therefore has modulus one.

The bound on \(t_C\), the bound on the number of kernels in Lemma~\ref{B8.43q}, and the common frame weight give, for each source frame \(x\),
\[
 |(T_{P,i,b}f)_j(x)|\le
 |bq_P^{-in/2}|\,2^{(p^2)^{2g}} q_P^{in}
           \sup_{j',A}|f_{j'}(Ax)|.
\]
It suffices to prove this at unit frames, since a change to a general frame contributes the same weight factor on both sides. Likewise the factor converting an absolute value into the invariant section norm is the same at \(x\) and \(Ax\), by the isometry of the Hermitian forms. Hence boundedness of the original cusp section in invariant norm gives boundedness of the transform.

There are finitely many \(A\)'s modulo the target arithmetic group. The lattice \(A^{-1}\Lambda_{j'}\) lies between \(p\Lambda_j\) and \(p^{-1}\Lambda_j\). This finite quotient has finitely many subgroups. Maps with the same inverse lattice differ on the left by a target lattice isometry, and the marking subgroup has finite index in those isometries. Fix representatives. A change by the arithmetic group multiplies the reading by a character of absolute value one.

For a nonzero isotropic lattice vector \(\xi\), let \(x(t)\) be the frame at imaginary parameter \(it\) in its central direction. The limit of a bounded holomorphic section on this path is its constant Fourier--Jacobi coefficient. Indeed its expansion has factors \(\exp(-4\pi t\sum_\tau\nu_\tau)\). Boundedness excludes a frequency with a negative coordinate, by taking the Fourier integral in the real directions after shifting that imaginary coordinate. The Fourier coefficients at an interior frame are absolutely summable, by shifting slightly within the domain. Dominated convergence leaves just \(\nu=0\). Under \(A\) this path is the path in the integral nonzero isotropic direction \(pA\xi\) at time \(t/p^2\), because the central displacement is quadratic in its vector. Every one of the finitely many terms bounding the transform tends to zero, by cuspidality of \(f\) on its target component. Thus the transform has zero constant term at \(\xi\). This holds on every component and at every designated embedding.
\leanref{GlobalTheory.\allowbreak{}HeckeEx.\allowbreak{}rdU\_\allowbreak{}mem\_\allowbreak{}SS}{Development.\allowbreak{}AbelianVarietyTheory.\allowbreak{}Proofs.\allowbreak{}TwistedOperators.\allowbreak{}HeckeTransformOfPointLaws}{334--398}
\leanref{GlobalTheory.\allowbreak{}HeckeEx.\allowbreak{}mem\_\allowbreak{}SS\_\allowbreak{}of\_\allowbreak{}dom}{Development.\allowbreak{}AbelianVarietyTheory.\allowbreak{}Proofs.\allowbreak{}TwistedOperators.\allowbreak{}MemCuspSectionsOfDominatedRationalIsometries}{83--159}
\end{proof}

\subsection*{Proof of \ref{B8.43t}}\label{prf:B8.43t}\phantomsection
\begin{proof}
First match kernels with representatives. In a marking of the period lattice \(\Lambda\) the quotient by \(C\subset B[p^2]\) replaces \(\Lambda\) by the inverse image of \(C\) in \(p^{-2}\Lambda\). Dividing the quotient differential by \(p\) amounts to using \(p\) times this enlarged lattice, with its original rational Hermitian normalisation. At \(P,\bar P\) the two factors are perfect duals. The inclusion of the \((p)P^{-1}\)-torsion fixes one end of the lattice interval, the inclusion in the \((p)\bar P\)-torsion fixes the other, and isotropy determines the second factor from the first. On the variable factor the interval has quotient a vector space of dimension \(r\) over \(\mathcal O_E/P\); the index condition selects the elementary divisor of order \(i\). In bases adapted to the subspace, the change on the contracting half is diagonal with \(i\) uniformisers and \(r-i\) ones, with inverse dual change on the other half. Away from this pair the normalised lattice is unchanged. Integral changes of adapted basis act transitively on these subspaces. The stabiliser of the reference lattice change in \(K\) is \(K\cap\alpha K\alpha^{-1}\), giving exactly \(K\alpha K/K\). The marking away from \(p\) is transported by \(p^{-1}\Psi\), as in Lemma~\ref{B8.43q}. Thus the frame at \(ha\) for the matched representative and the frame \(y_C\) agree up to a triangular change whose Lie diagonal is one. In period coordinates this is the equality of the divided map on the successive flagged quotients. Conversely, two such divided maps giving the same flagged frame have isometric differentials differing by a unipotent triangular matrix. A unipotent unitary matrix for a positive Hermitian metric is the identity. Their maps of complex tori and hence their kernels agree. This also gives the matching when a different set of right-coset representatives is used.
\leanref{GlobalTheory.\allowbreak{}HeckeTermIdent.\allowbreak{}heckeKernelCoset}{Development.\allowbreak{}AbelianVarietyTheory.\allowbreak{}Proofs.\allowbreak{}TwistedOperators.\allowbreak{}HeckeKernelComplexPointMatchesCosets}{73--117}

Put \(q=\mathrm NP\). Let \(\mathscr L\) be the tensor lattice of the source and target and \(\Delta_a\) representatives of \(a\mathscr L/(a\mathscr L\cap\mathscr L)\) in the expanding half. For the vacuum \(f_0\) and the Heisenberg translation \(\rho(x)\), the split lattice model gives
\[
       \omega(1,a)f_0=\mu_j(P)^i q^{-in/2}
                              \sum_{x\in\Delta_a}\rho(x)f_0 .
\]
Indeed there is one quotient of order \(q\) for each of the \(i\) changed directions in each of the \(n\) columns. The characteristic function on the enlarged half is the sum of the \(q^{in}\) coset characteristic functions. On its dual half the finite character sum kills precisely the points outside the contracted lattice. The unitary Fourier factor is \(1/\sqrt{q^{in}}\), and the splitting character on the determinant gives \(\mu_j(P)^i\). The two integral factors in \(a\in K\alpha K\) fix the vacuum and transport the representative set. This proves the identity for the actual representatives.
\leanref{GlobalTheory.\allowbreak{}HeckeVacuum.\allowbreak{}omega\_\allowbreak{}vacuum\_\allowbreak{}hecke}{Development.\allowbreak{}AbelianVarietyTheory.\allowbreak{}Proofs.\allowbreak{}TwistedOperators.\allowbreak{}SqrtDivMulEqSqrtMul}{359--401}

Let \(C_h\) be the product of the pinned component constant and the adelic multiplier at \(h\). Denote by \(t\) the scalar of the transported theta generator relative to the generator at \(ha\). We compute \(t^n\). In the unit frame at \(ha\), write \(\vartheta(z)\) for the theta generator and \(H(z,z)\) for the polarisation Hermitian form. Choose rational period vectors \(z_d\), one for each source column. Put
\[
      E_z=\exp\!\left(\frac{\pi}{2}\sum_d H(z_d,z_d)\right),
      \qquad V_z=\prod_d\vartheta(z_d).
\]
Let \(W_z\) denote the sum of readings of the translated vacua in the preceding display, before multiplication by its scalar. The evaluation point is the rational matrix with these period columns in the pinned orthonormal basis; the same rational target change is used in the pin and in \(W_z\). With \(\phi_0\) the phase in that pin, the vacuum identity at \(ha\) reads
\[
              E_z\,\mu_j(P)^i q^{-in/2} W_z
                        =\phi_0 C_{ha}V_z.
\]
The geometric trace in Lemma~\ref{B8.43q} gives
\[
                        W_z=t^n\phi_0 C_h E_z^{-1}V_z .
\]

We detail the identification of the second sum. In one column the pullback generator is summed at the \(q^i\) points of \(C/(C\cap B[p])\). Choose lifts \(v\) to the complex universal cover. If \(H_\Psi\) denotes the linear lift of the isogeny, \(H_\Psi v\) is a target period. In an Appell--Humbert frame the input theta argument is \(p(w+v)\), and the target trivialisation is evaluated at \(H_\Psi(w+v)\). Dividing by its value at \(H_\Psi w\) supplies the period multiplier. The identity \(H_\Psi^*H'=p^2H\) cancels the real exponent after inserting the Gaussian on each side. The remaining unit multiplier is the Heisenberg multiplier. Explicitly, write \([x,y]=\langle x,y\rangle/2\) for half the alternating form and \(\psi\) for the additive character. If \(g\) is the target change of frame, represent a finite translation \(x\) as \(g^{-1}L_f+u_x\), with \(L\) rational, \(L_f=L-L_\infty\), and \(\rho(u_x)f_0=f_0\). This representation exists by additive approximation modulo the open vacuum stabiliser. Evaluated at \(X_\infty\), the translated term is the unshifted term at \((X-L)_\infty\) times
\[
              \psi\bigl(-[L_\infty,(X-L)_\infty]-[L_f,gu_x]\bigr).
\]
Indeed \(\rho(u)f(w)=\psi([w,u])f(w+u)\), and rational translation by \(L\) contributes \(\psi(-[L,w])\). These two substitutions give the displayed factor. The same substitutions in the period multiplier give its infinite factor from translation by \(L_\infty\) and its finite factor from the vacuum stabiliser; the product character is trivial on rational pairings. Multiplying the one-column sums enumerates the tuples in \((C/(C\cap B[p]))^n\). Their finite half-lattice classes are distinct by the period-lattice description of \(C\), and both this set and \(\Delta_a\) have \(q^{in}\) elements. This gives the second identity with the indicated \(t^n\).
\leanref{GlobalTheory.\allowbreak{}HeckeGeoRead.\allowbreak{}theta\_\allowbreak{}trace\_\allowbreak{}eq\_\allowbreak{}phased\_\allowbreak{}sum}{Development.\allowbreak{}AbelianVarietyTheory.\allowbreak{}Proofs.\allowbreak{}TwistedOperators.\allowbreak{}HeckePointThetaTraceEqPhasedTranslateSum}{16--65}
\leanref{GlobalTheory.\allowbreak{}HeckeGeoRead.\allowbreak{}theta\_\allowbreak{}trace\_\allowbreak{}read}{Development.\allowbreak{}AbelianVarietyTheory.\allowbreak{}Proofs.\allowbreak{}TwistedOperators.\allowbreak{}ThetaTraceRead}{19--72}

Take all \(z_d\) equal to a rational period vector with \(\vartheta(z_d)\ne0\). Density and nonvanishing give such a vector as in Proposition~\ref{B8.43o}. The orthonormal source basis is invertible over \(E\), so a rational matrix with these prescribed columns exists. Substituting the second identity in the first and cancelling the nonzero \(\phi_0V_z\) gives
\[
                         t^n C_h q^{-in/2}=\mu_j(P)^{-i} C_{ha}.
\]
For any section of the given weight, the Lie diagonal of the triangular frame change is one, so its multiplier is exactly \(t^n\). Multiply the last identity by the common Lie normalisation and by the evaluation of \(f\) at \(ha\). It now identifies one summand in Proposition~\ref{B8.43s} with \(\iota(b)\mu_j(P)^{-i}R_f(ha)\). Summing over the bijection proves the formula.
\leanref{GlobalTheory.\allowbreak{}HeckeGenConst.\allowbreak{}generator\_\allowbreak{}constant}{Development.\allowbreak{}AbelianVarietyTheory.\allowbreak{}Proofs.\allowbreak{}TwistedOperators.\allowbreak{}HeckeTermGeneratorConstantViaFrame}{126--172}
\end{proof}

\subsection*{Proof of \ref{B8.43h}}\label{prf:B8.43h}\phantomsection
\begin{proof}
The base frame, coefficient map and extremal tensor are those of the standard layer. Its base level is the product of the two distinct odd type primes. On a tower cover the auxiliary factor is squarefree and prime to the excluded primes, and the field and pin assumptions are exactly those of Proposition~\ref{B8.43o}. That proposition gives the rational and integral readings with the stated order of choices. For descent the form is invariant under the base compact subgroup and each tested extension is unramified, so Lemma~\ref{B8.43p} applies without a clearing power. For the Hecke assertions the split primes satisfy \(P+\bar P=\mathcal O_E\) and are outside the full level. Apply Lemmas~\ref{B8.43q} and~\ref{B8.43r} and Proposition~\ref{B8.43s}. The member has even source rank and its pin uses the reference unramified character, as required in Proposition~\ref{B8.43t}. Finally the toric transport in \(P\) and the nonvanishing of \(T_0\) are the hypotheses of Lemma~\ref{B8.43i}; applying it to the pulled-back character gives the stated action.\end{proof}

\subsection*{Proof of \ref{B8.43}}\label{prf:B8.43}\phantomsection
\begin{proof}
Take the member point in its default extension with the archimedean realization and column pin. The third-layer and geometric fields give linked local and Galois data. The series fields associate a triple to each parameter list on a good finite-type row; the twist place meets the positive valuation condition and the all-good clause gives the type-idempotent vector with its prescribed Hecke eigenvalues. This supplies the good triples and the good labels for both character choices.
\leanref{GlobalTheory.\allowbreak{}HeightCounting.\allowbreak{}Up.\allowbreak{}TripFin.\allowbreak{}finiteTriples\_\allowbreak{}satisfy\_\allowbreak{}condition\_\allowbreak{}alt}{Development.\allowbreak{}AutomorphicArithmetic.\allowbreak{}ThetaCorrespondence.\allowbreak{}ThetaPairLifts.\allowbreak{}ThirdPlaceSeries.\allowbreak{}FiniteTriplesSatisfyCondition}{21--28}

Linkage gives the point conditions and first provision. The label factorization, tame identities and good triples supply the realized lists and their twists for \ref{B8.42pa}; admissibility gives split density. Thus every eligible residual realization has the required irreducibility and inertia exclusion. Reduced multiplicity and the nonzero theta scalars give the \(Z\)-identity in the internal isotypic decomposition. Read the selected dictionaries from the block data and apply depth-zero support at \(j-1\) for \(1\le j\le m\). The Hecke realization gives localization uniqueness and the extension supplies coefficient largeness. Norm-choice transport supplies the norm and same-twist ties. This establishes the Galois hypotheses of Theorem~\ref{B8.42} for the chosen norm.
\leanref{Formalism.\allowbreak{}MainNotions.\allowbreak{}AdmissibleThreePrime.\allowbreak{}residual\_\allowbreak{}irreducible\_\allowbreak{}of\_\allowbreak{}theorem\_\allowbreak{}three\_\allowbreak{}prime\_\allowbreak{}alt}{Development.\allowbreak{}AutomorphicArithmetic.\allowbreak{}ThetaCorrespondence.\allowbreak{}ThetaPairLifts.\allowbreak{}ThirdPlaceSeries.\allowbreak{}ResidualIrreducibleOfTheoremThreePrime}{12--14}
\leanref{UpperTheory.\allowbreak{}GalS.\allowbreak{}core0\_\allowbreak{}of\_\allowbreak{}local\_\allowbreak{}duality\_\allowbreak{}inputs}{Development.\allowbreak{}AutomorphicArithmetic.\allowbreak{}ThetaCorrespondence.\allowbreak{}ThetaPairLifts.\allowbreak{}ThirdPlaceSeries.\allowbreak{}PreAtMem}{54--85}
\leanref{UpperTheory.\allowbreak{}GalS.\allowbreak{}galois\_\allowbreak{}preconditions\_\allowbreak{}at\_\allowbreak{}mem}{Development.\allowbreak{}AutomorphicArithmetic.\allowbreak{}ThetaCorrespondence.\allowbreak{}ThetaPairLifts.\allowbreak{}ThirdPlaceSeries.\allowbreak{}PreAtMem}{86--103}

The jump and depth-zero data give the paired character, the depth-zero type and its eigenspan required in Theorem~\ref{B8.42}. The character pin gives its inertia polynomial. The local series identifies its label by the nonzero action of the corresponding central operator on a parameter vector. Transport these four compatibilities and the Galois hypotheses to the list extension with geometric closure. Retain the cutoff and thresholds of the linked \(2,3\)-condition. For each family in \(W_3(\mathcal A)\), apply \ref{B8.42} with \ref{B8.41} on the cut family.
\leanref{GlobalTheory.\allowbreak{}TypePrimesFinished.\allowbreak{}locComp\_\allowbreak{}of\_\allowbreak{}isRc\_\allowbreak{}classMapD\_\allowbreak{}cenLoc}{Development.\allowbreak{}AutomorphicArithmetic.\allowbreak{}ThetaCorrespondence.\allowbreak{}ThetaPairLifts.\allowbreak{}ThirdPlaceSeries.\allowbreak{}LocalTypeDeterminedAtAllPrimes}{22--29}
\leanref{UpperTheory.\allowbreak{}GalC.\allowbreak{}linkedAt\_\allowbreak{}linkN\_\allowbreak{}of\_\allowbreak{}boundAtTwoThree}{Development.\allowbreak{}AutomorphicArithmetic.\allowbreak{}ThetaCorrespondence.\allowbreak{}ThetaPairLifts.\allowbreak{}ThirdPlaceSeries.\allowbreak{}LinkNFam}{64--73}

In \ref{B8.03}, the local cohomology and comparison results assumed in \(P\) supply \(\mathscr H_2\). Its nonendoscopic and geometric Galois hypotheses give the third-layer data. Theorem~\ref{B8.43h} supplies the six series assertions for these data. In particular the member receives the section action constructed by geometric transport in Lemma~\ref{B8.43q} and Proposition~\ref{B8.43s}, with the same cusp subspace and the reading of Proposition~\ref{B8.43t}. Apply the closed bound at \(h_L,c_{\rm Dep}\) and strengthen the list extension by the resulting member property. Project to the selection of \ref{B8.04}. The restricted selection therefore satisfies the closed family local-error estimate on \(W_3(\mathcal A)\).
\leanref{GlobalTheory.\allowbreak{}MemberClose.\allowbreak{}exists\_\allowbreak{}hyp\_\allowbreak{}n\_\allowbreak{}of\_\allowbreak{}geometric\_\allowbreak{}galois\_\allowbreak{}compat}{Development.\allowbreak{}AutomorphicArithmetic.\allowbreak{}ThetaCorrespondence.\allowbreak{}ThetaPairLifts.\allowbreak{}ThirdPlaceSeries.\allowbreak{}ThetaLiftHypothesesOfNonvanishing}{24--46}
\leanref{GlobalTheory.\allowbreak{}Fc.\allowbreak{}S0.\allowbreak{}StageU.\allowbreak{}hypotheses\_\allowbreak{}nonempty\_\allowbreak{}of\_\allowbreak{}unramified\_\allowbreak{}exponents\_\allowbreak{}alt}{Development.\allowbreak{}AutomorphicArithmetic.\allowbreak{}ThetaCorrespondence.\allowbreak{}ThetaPairLifts.\allowbreak{}ThirdPlaceSeries.\allowbreak{}HypothesesNonemptyOfUnramifiedExponents}{41--61}
\leanref{GlobalTheory.\allowbreak{}Fc.\allowbreak{}AllAtMemberStage0.\allowbreak{}allRows\_\allowbreak{}of\_\allowbreak{}open0\_\allowbreak{}alt}{Development.\allowbreak{}AutomorphicArithmetic.\allowbreak{}ThetaCorrespondence.\allowbreak{}ThetaPairLifts.\allowbreak{}ThirdPlaceSeries.\allowbreak{}LinkROfIsLinked}{50--86}
\leanref{GlobalTheory.\allowbreak{}Fc.\allowbreak{}nonempty\_\allowbreak{}hyps\_\allowbreak{}of\_\allowbreak{}allRows\_\allowbreak{}alt}{Development.\allowbreak{}AutomorphicArithmetic.\allowbreak{}ThetaCorrespondence.\allowbreak{}ThetaPairLifts.\allowbreak{}ThirdPlaceSeries.\allowbreak{}NonemptyHypsOfAllRows}{19--101}
\end{proof}

\section*{Proofs for Chapter~\ref{ch:p8c5}: Primes in short intervals, masses and counts}

\subsection*{Proof of \ref{B8.47}}\label{prf:B8.47}\phantomsection
\begin{proof}
The conductor divides \(N_E\). If \(\chi\) is real and \(q,d\nmid N\), then \(N\mid8s_D\). Take \(x>x_1\); for every \(i\in I_x\), the family bounds \(o\le C_o\) and \(Q_S\le C_o|A|\), together with \(2^{e+2}\mid Q_S\) (with \(Q_S\ge1\) and odd part \(s_D\)) and \(|A|=2^e o\), give
\[
 2^{e+2}s_D\le Q_S\le C_o|A|=2^e C_o o\le 2^e C_o^2,
\]
and hence \(s_D\le C_o^2\). Nonvanishing at \(1\) gives a uniform positive distance from \(1\) to the real zeros of these finitely many nonprincipal characters. For the other real characters use \eqref{eq:p8c5-boxes} to exclude the exceptional zero of \ref{B8.45}, after shrinking \(c\) below \(\log(3)/400\). Together with the nonreal and zeta regions this gives a uniform \(c_0>0\) with
\[
 \Re\rho\le1-\frac{c_0}{\lambda_E+\log(|\Im\rho|+3)}.
 \tag{8.47.5}\label{eq:p8c5-region}
\]
Since \(1-\rho\) is a zero for \(\chi^{-1}\), which also belongs to \(D\), the functional equation gives
\[
 |\Im\rho|\le2\ \Longrightarrow\ \Re\rho\ge\kappa/\lambda_E,\qquad \kappa=c_0/3.
 \tag{8.47.6}\label{eq:p8c5-left}
\]
Group zeros by \(j=\lfloor|\Im\rho|\rfloor\). Each group has multiplicity at most \(2C_{11}(\lambda_E+\log(\lfloor T\rfloor+2))\) by \eqref{eq:p8c5-window}. Bound the reciprocal weight by \({\bf1}_{j\le2}\lambda_E/\kappa+2/(j+1)\) and sum to obtain, for \(T\ge1\),
\[
 \sum_{\rho\in Z_\chi(T)}\frac{m_\chi(\rho)}{|\rho|}
 \le C_2(\lambda_E+\log T)^2,\qquad C_2=6C_{11}(3/\kappa+4).
 \tag{8.47.7}\label{eq:p8c5-reciprocals}
\]
\leanref{Formalism.\allowbreak{}MainNotions.\allowbreak{}ShortIntervalPrimes.\allowbreak{}zeroSumE\_\allowbreak{}of}{Formalism.\allowbreak{}SharpRestriction.\allowbreak{}GlobalSchurBounds.\allowbreak{}FirstClaimLemmaTwoFourOfLemmaTwoThree}{2355--2377}

For \(v=\log u\ge1\), put \(T=e^{\sqrt v}\) and \(H=e^{-c_1v/(\lambda_E+\sqrt v)}\), \(c_1=\min(\kappa,1)\). Then \(1/T\le H\) and \(uH\ge1\). On \(Z_\chi(T)\) we have \(u^{\Re\rho}\le uH\). Equations \eqref{eq:p8c5-explicit} and \eqref{eq:p8c5-reciprocals} give
\[
 |\psi_0(\chi,u)|\le[C_7(5+C_2)+C_2]u(v+\lambda_E)^2H .
 \tag{8.47.8}\label{eq:p8c5-psi}
\]
\leanref{Formalism.\allowbreak{}MainNotions.\allowbreak{}ShortIntervalPrimes.\allowbreak{}psi\_\allowbreak{}thin\_\allowbreak{}core}{Formalism.\allowbreak{}SharpRestriction.\allowbreak{}GlobalSchurBounds.\allowbreak{}FirstClaimLemmaTwoFourOfLemmaTwoThree}{2028--2148}

For \(s>1\), integrating \(\psi_0(\chi,u)u^{-s-1}\) gives \(-L'(s,\chi)/(sL(s,\chi))-{\bf1}_{\chi=1}/(s-1)\). By \eqref{eq:p8c5-psi} we may pass to \(s=1\). The limits are \(-L'(1,\chi)/L(1,\chi)\) and, for zeta, \(-1-\gamma\). Partial summation in \eqref{eq:p8c5-centred-sum} gives
\[
 g_\chi(y)=-y\int_y^\infty\psi_0(\chi,u)\,\frac{du}{u^2}.
 \tag{8.47.9}\label{eq:p8c5-tail}
\]
\leanref{Formalism.\allowbreak{}MainNotions.\allowbreak{}ShortIntervalPrimes.\allowbreak{}l21a}{Formalism.\allowbreak{}SharpRestriction.\allowbreak{}GlobalSchurBounds.\allowbreak{}FirstClaimLemmaTwoFourOfLemmaTwoThree}{368--402}

For large \(x\), \eqref{eq:p8c5-qd} gives
\[
 \lambda_E\le8x^{\varepsilon/5}(\log x)^2.
 \tag{8.47.10}\label{eq:p8c5-conductor}
\]
In \eqref{eq:p8c5-tail} integrate the bound \eqref{eq:p8c5-psi} from \(V=x^{\varepsilon/3}-\log3\) in the variable \(v=\log u\). Up to \(\lambda_E^2\), use \(H\le e^{-c_1v/(2\lambda_E)}\). Above \(\max(V,\lambda_E^2)\), use \(H\le e^{-c_1\sqrt v/2}\). Integration in \(v\) and \(\sqrt v\) gives polynomial factors times \(e^{-c_1V/(2\lambda_E)}\) and \(e^{-c_1\sqrt V/2}\). Both \(V/\lambda_E\gg x^{2\varepsilon/15}/(\log x)^2\) and \(\sqrt V\gg x^{\varepsilon/6}\) dominate \(x^{\varepsilon/8}\) and the logarithms of the polynomial factors. This proves \eqref{eq:p8c5-thin}.
\leanref{Formalism.\allowbreak{}MainNotions.\allowbreak{}ShortIntervalPrimes.\allowbreak{}thin\_\allowbreak{}of\_\allowbreak{}zeroFree}{Formalism.\allowbreak{}SharpRestriction.\allowbreak{}GlobalSchurBounds.\allowbreak{}FirstClaimLemmaTwoFourOfLemmaTwoThree}{2379--2397}
\leanref{Formalism.\allowbreak{}MainNotions.\allowbreak{}ShortIntervalPrimes.\allowbreak{}thin\_\allowbreak{}of\_\allowbreak{}psiThin}{Formalism.\allowbreak{}SharpRestriction.\allowbreak{}GlobalSchurBounds.\allowbreak{}FirstClaimLemmaTwoFourOfLemmaTwoThree}{1949--1992}
\end{proof}

\subsection*{Proof of \ref{B8.49}}\label{prf:B8.49}\phantomsection
\begin{proof}
Diagonalize \(W\) in the standard signature form. Let \(L_1\) be the coordinate lattice and \(x\in L_1\) a coordinate vector of norm \(s\). The base lattice of \(W'\) is \(L_1\cap x^\perp\). Set \(W_1=(W,h,L_1)\). The two self-dual base lattices have conjugate local stabilizers at every finite place, and the real normalization depends on the signature. Thus \(\operatorname{vol}_D(W_1)=\operatorname{vol}_D(W)\). Apply the one-vector Gram formula of \ref{B8.13} to \(W_1\) with a real test of positive fibre integral \(a\). The mean is
\[
 a\prod_{v<\infty}\alpha_v
 \tag{8.49.2}\label{eq:p8c5-mean}
\]
for the finite densities \(\alpha_v\).
\leanref{GlobalTheory.\allowbreak{}UnitPerp.\allowbreak{}unitVectorPerpOver}{Development.\allowbreak{}AbelianVarietyTheory.\allowbreak{}Proofs.\allowbreak{}GlobalInterpolation.\allowbreak{}HasComplementOfBasis}{248--255}

The ball of squared radius \(t\) on the compact dual has self-dual measure \((2\pi)^Mt^M/M!\). Differentiating at \(1\) gives the fibre factor \(c_M=(2\pi)^M/(M-1)!\), also the factor in signature frames for norm \(s\). Unfolding along the stabilizer of \(x\) gives the same mean as
\[
 \frac{\operatorname{vol}_D(W')}{\operatorname{vol}_D(W_1)}c_M^{-f}a .
 \tag{8.49.3}\label{eq:p8c5-unfold}
\]
\leanref{GlobalTheory.\allowbreak{}RealPlaceMeasuresProof.\allowbreak{}realPlaceMeasures}{Development.\allowbreak{}AbelianVarietyTheory.\allowbreak{}Proofs.\allowbreak{}GlobalInterpolation.\allowbreak{}RealPlaceMeasures}{138--140}

A nontrivial additive character summed over one residue norm gives \(q_v\) in the split case by summing over one factor, and \(-q_v\) in the inert case since each nonzero value has \(q_v+1\) preimages. Additive orthogonality gives \(q_v^{2M-1}(1-\eta_v^M q_v^{-M})\) vectors of norm \(s\). Surjectivity of the norm differential at a vector of unit norm gives \(q_v^{2M-1}\) lifts per step. If \(i_v\) is the norm of the local different of \(F\), the domain and value measures contribute \(i_v^{-M}/i_v^{-1/2}\). Thus
\[
 \prod_{v<\infty}\alpha_v
 =\prod_{v<\infty}i_v^{1/2-M}(1-\eta_v^M q_v^{-M})
 =\frac{(2v_E)^{f(1/2-M)}}{L_\eta(E,M)} .
 \tag{8.49.4}\label{eq:p8c5-density-product}
\]
Cancel \(a\) between \eqref{eq:p8c5-mean} and \eqref{eq:p8c5-unfold} and substitute the volume of \(W\).
\leanref{GlobalTheory.\allowbreak{}LatticeDensity.\allowbreak{}localDensity\_\allowbreak{}of\_\allowbreak{}selfDualLatticeLemmas}{Development.\allowbreak{}AbelianVarietyTheory.\allowbreak{}Proofs.\allowbreak{}GlobalInterpolation.\allowbreak{}LocalDensityOfSelfDualLatticeLemmas}{208--295}
\leanref{GlobalTheory.\allowbreak{}VolRatioUnfold.\allowbreak{}ratio\_\allowbreak{}of\_\allowbreak{}setting}{Development.\allowbreak{}AbelianVarietyTheory.\allowbreak{}Proofs.\allowbreak{}GlobalInterpolation.\allowbreak{}VolumeRatioOfUnfolding}{295--391}
\leanref{GlobalTheory.\allowbreak{}VolRatioUnfold.\allowbreak{}volumeRatio\_\allowbreak{}of\_\allowbreak{}unfolding}{Development.\allowbreak{}AbelianVarietyTheory.\allowbreak{}Proofs.\allowbreak{}GlobalInterpolation.\allowbreak{}VolumeRatioOfUnfolding}{393--421}
\end{proof}

\subsection*{Proof of \ref{B8.50}}\label{prf:B8.50}\phantomsection
\begin{proof}
At inert places the determinant of a self-dual Gram matrix is a unit and hence a norm. At split places the forms are isometric, and at infinity their signatures agree. Landherr's theorem gives a rational isometry with \(h_M\) \cite[Ch.~10]{scharlau-1985}. The transported self-dual lattice belongs to the standard genus by the local classification \cite{jacobowitz-1962}, and the isometry preserves the automorphism orders.
\leanref{GlobalTheory.\allowbreak{}DefiniteMass.\allowbreak{}exists\_\allowbreak{}isometry}{Development.\allowbreak{}AbelianVarietyTheory.\allowbreak{}Proofs.\allowbreak{}GlobalInterpolation.\allowbreak{}UnramifiedOfIsCMUnram}{45--94}
\leanref{GlobalTheory.\allowbreak{}DefiniteMass.\allowbreak{}sameGenus\_\allowbreak{}stdLat}{Development.\allowbreak{}AbelianVarietyTheory.\allowbreak{}Proofs.\allowbreak{}GlobalInterpolation.\allowbreak{}UnramifiedOfIsCMUnram}{126--145}
\end{proof}

\subsection*{Proof of \ref{B8.50a}}\label{prf:B8.50a}\phantomsection
\begin{proof}
The gap exceeds \(32rpq+2\), and \(\lambda_i\in(r-1)/2+\Z\). Set \(\lambda_i(u)=\lambda_i+u(r-1-i)\), \(0\le i<r\), and let \(d_u\) be its formal degree. The one-coordinate increment formula gives the quotient \(d(\mu+e_i)/d(\mu)=\Delta(\mu+e_i)/\Delta(\mu)\) on this path. At each step the increments from the first index to the last preserve order and the required gap. Hence \(d_u=C_\lambda\Delta(\lambda(u))\).
\leanref{ReciprocityTheory.\allowbreak{}ChartRestriction.\allowbreak{}DS.\allowbreak{}fdc\_\allowbreak{}cd}{ReciprocityTheory.\allowbreak{}ReductiveGroups.\allowbreak{}DiscreteSeries.\allowbreak{}FormalDegree.\allowbreak{}FdCons}{17--53}
\leanref{ReciprocityTheory.\allowbreak{}ChartRestriction.\allowbreak{}DS.\allowbreak{}fdCons\_\allowbreak{}cd}{ReciprocityTheory.\allowbreak{}ReductiveGroups.\allowbreak{}DiscreteSeries.\allowbreak{}FormalDegree.\allowbreak{}FdCons}{54--87}

Write \(\tau_u\) for the lowest type. In polar coordinates \(k\exp Z(z)\) use one complex coordinate \(z_{ab}\) for each \(a<b\) of opposite signs, with \(Z(z)\) Hermitian. The compact measure is probability measure. The real density is \(c_{pq}J(z)\), where \(J(0)=1\) and \(c_{pq}=F_r/(\pi^{pq}F_pF_q)\). Let \(\Phi_u\) be the negative logarithm of the weight in \eqref{eq:p8c5-compression} and \(a_u(z)\) the coefficient at a unit highest vector of \(\tau_u\). Schur orthogonality, followed by the compact average of \eqref{eq:p8c5-compression}, gives
\[
 c_{pq}\int |a_u(z)|^2J(z)\,dz
 \ \le\ \frac{\dim\tau_u}{d_u}
 \ \le\ c_{pq}\int e^{-2\Phi_u(Z(z))}J(z)\,dz .
\]
For the lower bound integrate the squared compression applied to the highest vector. For the upper bound its Hilbert--Schmidt norm squared is at most \(\dim\tau_u\) times the compact average of \(e^{-2\Phi_u}\), by Cauchy--Schwarz and Schur orthogonality on the compact type. Conjugation by the compact group preserves \(J(z)\,dz\).

After \(z\mapsto u^{-1/2}z\), the second-order term of the log of the \(i\)-th minor is \(\frac{1}{2u}\sum_{a\le i<b}|z_{ab}|^2\), with the sum over opposite signs. Thus \(2\Phi_u\) tends to \(Q(z)=\sum_{a<b,\ s_a\ne s_b}(b-a)|z_{ab}|^2\). In the highest-vector coefficient the triangular factorization of \(\exp(u^{-1/2}Z(z))\) has zero first-order weight term. Its quadratic weight term is \(-Q(z)/2\), from the commutators of opposite matrix units, and the terms of order three and higher contribute \(O_z(u^{-1/2})\) in the tangent coefficient estimate. Hence \(a_u(u^{-1/2}z)\to e^{-Q(z)/2}\). For the upper integral, the minor bound \eqref{eq:p8c5-minor-decay} gives a lower bound for \(2\Phi_u\) by a positive rank-dependent constant times \(u\min(\|z\|^2,\|z\|)\). The polar Jacobian grows at most exponentially in \(\|z\|\), by its product of \(\sinh v/v\) factors. The rescaled upper integrand is therefore bounded by an integrable exponential for all sufficiently large \(u\). Dominated convergence above and Fatou below give the same limit for \(u^{pq}\dim\tau_u/d_u\),
\[
 c_{pq}\int_{\C^{pq}} e^{-Q(z)}\,dz
 =\frac{F_r}{F_pF_q}\prod_{a<b,\ s_a\ne s_b}(b-a)^{-1}.
\]
\leanref{ReciprocityTheory.\allowbreak{}ChartRestriction.\allowbreak{}DS.\allowbreak{}matrix\_\allowbreak{}coefficient\_\allowbreak{}pointwise\_\allowbreak{}limit}{ReciprocityTheory.\allowbreak{}ReductiveGroups.\allowbreak{}DiscreteSeries.\allowbreak{}FormalDegree.\allowbreak{}FdPointwise}{120--168}
\leanref{ReciprocityTheory.\allowbreak{}ChartRestriction.\allowbreak{}DS.\allowbreak{}formal\_\allowbreak{}degree\_\allowbreak{}limsup\_\allowbreak{}upper\_\allowbreak{}bound}{ReciprocityTheory.\allowbreak{}ReductiveGroups.\allowbreak{}DiscreteSeries.\allowbreak{}FormalDegree.\allowbreak{}FdUpper}{95--140}

The compact Weyl numerator has same-sign factors \(\lambda_a(u)-\lambda_b(u)+\rho_N(a)-\rho_N(b)\), where \(\rho_N\) is the half-sum of positive noncompact roots. Their ratios to the corresponding factors of \(\Delta\) tend to one. The opposite-sign factors of \(\Delta/u^{pq}\) contribute \(\prod(b-a)\) after the same-sign factors are cancelled. The Weyl limit is therefore \(1/(C_\lambda F_pF_q)\prod(b-a)^{-1}\). Comparison gives \(C_\lambda=1/F_r\). At \(u=0\) the shifts \(\rho_N(a)-\rho_N(b)\) for same signs are nonnegative, and the \(pq\) other differences are bounded by \(H\). Since \(F_r\ge F_pF_q\), \(d_{\Sigma,\rho}=V(\lambda)\le d_\tau H^{pq}\), which implies \eqref{eq:p8c5-degree}.
\end{proof}

\subsection*{Proof of \ref{B8.51b}}\label{prf:B8.51b}\phantomsection
\begin{proof}
Write \(M=P+Q\). If \(Q>0\), take the orthogonal complement \(W'\) of a negative unit vector supplied by \ref{B8.49}. Its signature is \((P,Q-1)\). If \(Q=0\) and \(P>2\), use the positive unit vector in that proposition; its complement has signature \((P-1,0)\). In either case the complement is admissible, and \ref{B8.48} identifies the quotient in the volume recursion as
\begin{equation}\label{eq:p8c5-reverse-ratio}
 \frac{\operatorname{vol}_D(W')}{\operatorname{vol}_D(W)}
   =\frac{M_{M-1}}{M_M}.
 \tag{8.51.4}
\end{equation}
All masses are positive. Thus the ratio of volume to mass is unchanged upon passing to \(W'\). After \(Q\) negative and \(P-2\) positive steps this ratio is the definite rank-two ratio, which is one. The starting row has constant signature with positive index at least two and the sum of its indices is \(r\).
\end{proof}

\subsection*{Proof of \ref{B8.51}}\label{prf:B8.51}\phantomsection
\begin{proof}
Set \(C_Y=\max(5,C_{YL}+1)\) and \(C=c_{T,\min}(C_{YL})\). Lemma~\ref{B8.51b} gives the starting volume in the product formula \eqref{eq:B5-count-product}. Use the family estimates of \ref{B5.11} above the cutoff of \ref{B5.11c}, with \(\Gamma_u\) common to the words. Its upper product bound is
\begin{equation}\label{eq:p8c5-xu-upper}
 \log X_u\le(15+3C_Y)mL .
 \tag{8.51.5}
\end{equation}
For the lower bound, the weights in \(\mathcal W\) satisfy \(w_{j\rho}\ge1\) and \(v_E\ge\pi\), and \(\ell_g\ge16L\). Hence the lower product bound of \ref{B5.11} gives \(\log X_u\ge16mL-11L\). Since \(T_u\ge X_u\Gamma_u^{-2/(fk)}\) and \(C\ge11\), this implies
\begin{equation}\label{eq:p8c5-torus-lower}
 \log T_u+\frac{2\log\Gamma_u}{fk}\ge16mL-C(L+1).
 \tag{8.51.7}
\end{equation}
The denominator bound in \ref{B5.11} gives \(2\log\Gamma_u/(fk)\le2\), so \(\log T_u\ge16mL-11L-2\ge mL/2\) for \(m,L\ge1\). At the upper endpoint, \(\Gamma_u\ge1\) and \(\log(1+1/(fk))\le1\). Thus \(\log T_u\le(15+3C_Y)mL+1\le C_{Tu}mL\). Since \(b=e^L\), these are \eqref{eq:p8c5-level-scale}.
\end{proof}

\subsection*{Proof of \ref{B8.51a}}\label{prf:B8.51a}\phantomsection
\begin{proof}
Put \(M=N(g_0g_1)^2\) and \(D=4dM^2\). The lower bound on \(S\) gives
\[
        \sqrt{\log S}\ge\sqrt{c_s}\,x^{1/6+\varepsilon/2}.
\]
Since \(d\mid N\), the modulus and width satisfy
\[
 \log D\le\log4+9x^{\varepsilon/5}(\log x)^2+8C_Yx^\varepsilon,
 \qquad \log V\le12\log x .
\]
After division by \(x^{1/6+\varepsilon/2}\) both right sides tend to zero, since \(\varepsilon<1/52\). Choose the cutoff so that both quantities are at most \(c\sqrt{\log S}\), with the absolute \(c\) of \ref{B2.50}. Thus \(h=1/V\) belongs to \([e^{-c\sqrt{\log S}},1]\). Also \(M\ge d\ge100\), \(S>1\), and \(d\) is prime and \(3\bmod4\). Apply \ref{B2.50} with modulus \(M\), starting point \(S\), and width \(h\). Its lower bound \(c'hS/(D^A\log S)\) is positive. For a prime counted there take integers \(A_1,B_1\) with
\[
  A_1^2+dB_1^2=4T,\qquad M\mid B_1,\qquad 2M\mid A_1-2-B_1.
\]
Then
\(\pi_T=(A_1-B_1)/2+B_1(1+\sqrt{-d})/2\) is integral, has norm \(T\), and is \(1\bmod M\mathcal O_K\).
\end{proof}

\subsection*{Proof of \ref{B8.52a}}\label{prf:B8.52a}\phantomsection
\begin{proof}
Here \(L^2(X)\) denotes the automorphic \(L^2\) space with the quotient and measure of \(X\). The rank and gap assumptions give integrability of the real coefficients by \ref{B8.02y}. The finite projectors are locally constant and supported on the lattice stabilizer. Lemma~\ref{B8.03a} therefore applies, including to \(e\) itself.

For the evaluation assertion let \(P=T_\phi\), and write \(R(g)\) for right translation at infinity. Formal-degree orthogonality gives the matrix-unit identity
\[
 P R(g)P=\langle w,\sigma(g)w\rangle P
\]
with the convention for the coefficient in \ref{B8.03a}. Indeed the convolution of the normalized conjugate coefficient with its \(g\)-translate is that scalar times the same coefficient, by integration at each real place. If \(z\in V_S\), both \(P\) and \(R_{e_S}\) fix \(z\). The identity gives
\(\langle R(g)z,R(h)z\rangle=\|z\|^2\langle\sigma(g)w,\sigma(h)w\rangle\).
Thus \(\sigma(g)w\mapsto R(g)z\) extends linearly and continuously from the cyclic span of \(w\) to an intertwiner. Its image is fixed by \(R_{e_S}\), since the finite and real actions commute. Conversely evaluation of a projected intertwiner lies in \(V_S\), by the same orthogonality relation in \(\sigma\). An intertwiner is determined by its value on the cyclic vector \(w\).
\leanref{Formalism.\allowbreak{}L2.\allowbreak{}SG.\allowbreak{}Closed.\allowbreak{}RowsBij\_\allowbreak{}holds}{Formalism.\allowbreak{}LocalToGlobal.\allowbreak{}SiegelLatticeTheory.\allowbreak{}ClosedVolumeStatements.\allowbreak{}LevelSubgroupProperties}{55--56}
\leanref{Formalism.\allowbreak{}L2.\allowbreak{}SG.\allowbreak{}Pf.\allowbreak{}AllRowsDcount.\allowbreak{}dcountEq\_\allowbreak{}holds}{Formalism.\allowbreak{}LocalToGlobal.\allowbreak{}SiegelLatticeTheory.\allowbreak{}FlagStabilizerTensors.\allowbreak{}RootLatticeThresholds.\allowbreak{}AllRowsDcount}{58--68}

The single-label projectors are mutually orthogonal idempotents by \ref{B7.44}, applied in each finite factor. Their sum over all labels acts as the identity on the image of \(e\). Their sum over \(S\) has image \(\bigoplus_{t\in S}V_{\{t\}}\). This proves the dimension formula. The arithmetic identification in \ref{B8.08} identifies the dimension for the full set with the prescribed total row dimension.

For the good labels, the isometry and row identification of \ref{B8.08} give a basis of the good subspace by transporting the selected basis, with complex conjugation in the row-function convention. Conjugation preserves linear independence, and surjectivity onto the good subspace gives spanning. The label identification in that same hypothesis puts each transported vector in its specified single-label space. Applying the projector for a good \(t\) to a basis expansion retains precisely the vectors labelled \(t\). Their number is therefore \(d_t\). Take \(d_t\) indices for each label, and for good \(t\) match these indices bijectively with the selected basis indices labelled \(t\). The resulting injection has image exactly the good indices.
\leanref{Formalism.\allowbreak{}L2.\allowbreak{}SG.\allowbreak{}RowDict.\allowbreak{}range\_\allowbreak{}sEE}{Formalism.\allowbreak{}LocalToGlobal.\allowbreak{}SiegelLatticeTheory.\allowbreak{}FlagStabilizerTensors.\allowbreak{}RowDict}{148--163}
\leanref{Formalism.\allowbreak{}L2.\allowbreak{}SG.\allowbreak{}AllRowsOf.\allowbreak{}rowsLab\_\allowbreak{}of}{Formalism.\allowbreak{}LocalToGlobal.\allowbreak{}SiegelLatticeTheory.\allowbreak{}FlagStabilizerTensors.\allowbreak{}AllRowsOf}{864--884}
\end{proof}

\subsection*{Proof of \ref{B8.52}}\label{prf:B8.52}\phantomsection
\begin{proof}
Put \(A_0(\ell)=2(\ell[6]+1)\), with zero-based indexing and default entry zero. The seventh entry of \(\ell_{\rm out}\) is \(C_{YL}\). In the order \((C_{\rm res},n_0,A,C_b,C,c,c_b)\), take
\[
 c_{30}(\ell)=(0,n_0(A_0(\ell)),A_0(\ell),0,C_0(A_0(\ell)),1/16,1/16),
 \qquad c_{\rm Ic}=1/16 ,
 \tag{8.52.3}\label{eq:p8c5-count-constants}
\]
using the cutoffs of \ref{B8.53}.
\leanref{GlobalTheory.\allowbreak{}CountConstants.\allowbreak{}count\_\allowbreak{}constants\_\allowbreak{}of\_\allowbreak{}siegelWeil}{Development.\allowbreak{}AbelianVarietyTheory.\allowbreak{}Proofs.\allowbreak{}FlagMeasures.\allowbreak{}CountConstantsOfSiegelWeil}{11--24}

For the numerical gap required by \ref{B8.02y}, use the level-size bound included in membership of \(\mathcal D_3\) in \ref{B8.02}, and enlarge the last threshold of \ref{B8.39} to include \(29\) and \(8\sum_{c\in\ell}|c|\) for the concatenated outer and middle lists. Write \(b_{\rho,j}\) for the \(j\)-th auxiliary weight in the Harish--Chandra parameters at \(\rho\), and suppress \(\rho\) when the place is fixed. For the family constants of Theorem~\ref{B8.10}, the enlarged threshold gives \(C'\ge29,\ C_\ell\ge16\) and
\[
 r\ge k\ge C'm^5fL,\qquad
 b_{\rho,j+1}-b_{\rho,j}
 \ge C_\ell r^{19}f^5L.
 \tag{8.52.4}\label{eq:p8c5-selected-gap}
\]
Since \(m\ge2,\ f\ge m^4,\ L>1\), this gives \(r\ge1801\) and auxiliary gaps at least \(16r^{19}\).
\leanref{Formalism.\allowbreak{}L2.\allowbreak{}SG.\allowbreak{}GapClosed.\allowbreak{}gapEv\_\allowbreak{}closed\_\allowbreak{}of}{Formalism.\allowbreak{}LocalToGlobal.\allowbreak{}SiegelLatticeTheory.\allowbreak{}FlagStabilizerTensors.\allowbreak{}RootLatticeThresholds.\allowbreak{}GapClosed}{10--15}
\leanref{Formalism.\allowbreak{}L2.\allowbreak{}SG.\allowbreak{}ThmSOwn.\allowbreak{}gapRowT}{Formalism.\allowbreak{}LocalToGlobal.\allowbreak{}SiegelLatticeTheory.\allowbreak{}FlagStabilizerTensors.\allowbreak{}LinkSGap}{17--57}

For the counting hypotheses choose a second threshold with \(C_\ell\ge\max(2,2C)\). Require also \(n_0\le k\) and \(\log L\le c_bk\). The size conditions give \(\log\mathrm{rd}\le L/C_\ell\). Thus \(C\log(b\,\mathrm{rd})\le C_\ell L\). At this threshold the adjacent auxiliary gaps are at least \(G=C_\ell r^{19}f^5L\). In the row order the initial block has one sign and nonnegative row parameters, and the subsequent block has parameters \(-b_j\) with alternating signs starting with the initial sign. The first auxiliary weight is nonnegative. Every pair of opposite signs inside the auxiliary block has distance at least \(G\) by summing consecutive gaps. Every opposite-sign pair involving the initial block uses \(j\ge2\), so its distance is at least \(b_j\ge b_1+G\). The size conditions also give \(|\lambda_{v,a}|\le b^2\), \(q_v^2\le r\), \(\log f\le\sqrt r\), and \(\log\log\mathrm{rd}\le r\). Since \(k\le rf\log r\), the bound on \(\log L\) gives the last inequality of \eqref{eq:p8c5-count-gap}. At the two distinguished primes the bounds give \(g_0g_1\le b^A,\ 4r<g_0,g_1\), with both primes completely split in \(E\).
\leanref{GlobalTheory.\allowbreak{}countClauses\_\allowbreak{}of\_\allowbreak{}sizeHyp}{Development.\allowbreak{}AbelianVarietyTheory.\allowbreak{}Proofs.\allowbreak{}FlagMeasures.\allowbreak{}CountParameterConditionsOfSizeHypotheses}{9--83}
\leanref{GlobalTheory.\allowbreak{}IrrRowConcrete.\allowbreak{}irr\_\allowbreak{}row\_\allowbreak{}closed}{Development.\allowbreak{}AbelianVarietyTheory.\allowbreak{}Proofs.\allowbreak{}FlagMeasures.\allowbreak{}IrreducibilityRowOfConcreteAutomorphicData}{18--41}

Intersect the numerical and realization size tests by \ref{B8.39}. For a family passing them, take the maximum of the gap, counting and realization point cutoffs and the family starting scale plus one. At each point take the realization supplied by the assumed label row. Equation \eqref{eq:p8c5-selected-gap} permits the use of \ref{B8.52a} on this realization. Use its original quotient, measures and finite projectors for the arithmetic row system. At every place over \(g_\nu\) use the same character \(a\mapsto\exp(2\pi i a/g_\nu)\) of \(\Z/g_\nu\Z\), which is nontrivial at \(1\). The realization identifies the archimedean parameters, the split primes, the total row dimension and the positive total main term with those of the point. In the finite quotient let \(P_j\) be the subgroup-character projector at \(g_j\), and let \(Z_S\) be the sum of the central projectors for the types in \(S\) at \(g_{k_0}\). The cylinder projector has factor \(Z_SP_{k_0}\) at \(g_{k_0}\) and factor \(P_{1-k_0}\) at the opposite prime. These are commuting orthogonal projectors by \ref{B7.44}. On the labelled space of Lemma~\ref{B8.52a} their product fixes precisely the basis vectors whose labels belong to the cylinder and kills the others. To lift the finite convolution idempotent, divide its value on each congruence coset by the Haar volume of the kernel and extend by zero outside the lattice stabilizer. Coset integration then gives the finite convolution product, so the lift is a self-adjoint idempotent bi-invariant under the congruence kernel. At the opposite prime its factor is the original character projector on the subgroup with unipotent reduction. This factorization also gives the counting condition for a cylinder: the opposite factor is a subgroup-character idempotent and the chosen factor is the product of a central idempotent and a subgroup-character idempotent. The size and opposite-sign gap inequalities above apply at \(n=r\), with \(g_\circ=g_{1-k_0}>4r\) split completely. Thus \ref{B8.53} applies to this very projector. Its counted dimension is the full cylinder dimension by \ref{B8.52a}, and its main term uses the same formal degree and quotient volume. For the lifted projector \(e\), self-adjointness and idempotence give \(e(1)=(e^**e)(1)=\int |e(h)|^2\,d\mu_f(h)\ge0\), including when the cylinder is empty. This gives \eqref{eq:p8c5-cylinder-count}. For the full type set, \(Z_S=1\). Each subgroup-character projector has positive identity coefficient, so their product has positive main term \(D_A^\star(e')\). With \(D_{\rm all}\) the full row dimension, the same estimate gives \( |D_{\rm all}-D_A^\star(e')|\le e^{-rf\log r/16}D_A^\star(e')\).
\leanref{Formalism.\allowbreak{}L2.\allowbreak{}SG.\allowbreak{}Pf.\allowbreak{}AllRowsClosed.\allowbreak{}isAdmIdem\_\allowbreak{}cyl}{Formalism.\allowbreak{}LocalToGlobal.\allowbreak{}SiegelLatticeTheory.\allowbreak{}FlagStabilizerTensors.\allowbreak{}RootLatticeThresholds.\allowbreak{}AllRowsClosed}{30--31}
\leanref{Formalism.\allowbreak{}L2.\allowbreak{}SG.\allowbreak{}Pf.\allowbreak{}AllRowsPCount.\allowbreak{}isCountIdem\_\allowbreak{}cyl}{Formalism.\allowbreak{}LocalToGlobal.\allowbreak{}SiegelLatticeTheory.\allowbreak{}FlagStabilizerTensors.\allowbreak{}RootLatticeThresholds.\allowbreak{}AllRowsPCount}{611--693}
\leanref{Formalism.\allowbreak{}MainNotions.\allowbreak{}N.\allowbreak{}CylinderCountFamily.\allowbreak{}countCylAt\_\allowbreak{}of}{Formalism.\allowbreak{}SharpRestriction.\allowbreak{}LevelLetterCocycles.\allowbreak{}JointConditionsZpowers.\allowbreak{}JCountP}{16--22}
\leanref{GlobalTheory.\allowbreak{}AllRowsAtRecord.\allowbreak{}allRows\_\allowbreak{}at\_\allowbreak{}record}{Development.\allowbreak{}AbelianVarietyTheory.\allowbreak{}Proofs.\allowbreak{}FlagMeasures.\allowbreak{}AllRowsAtRecord}{67--95}
\end{proof}

\subsection*{Proof of \ref{B8.53}}\label{prf:B8.53}\phantomsection
\begin{proof}
Evaluation at \(w\) identifies the counted space with the image of convolution by
\(\phi(y)=d_\infty\overline{\prod_v\langle w_v,\sigma_v(y_v)w_v\rangle}e(y_f)\), by Schur orthogonality. Unfolding gives
\[
 K(x,y)=\sum_{\gamma\in U(F)}\phi(x^{-1}\gamma y).
 \tag{8.53.5}\label{eq:p8c5-kernel}
\]
The diagonal identity term is \(a=d_\infty\Re e(1)\). In definite signature the traces \(\operatorname{tr}(\gamma^j)-n\) of a supported element, for \(1\le j\le n\), are integral and vanish modulo \(g_\circ\). Their conjugates have absolute value at most \(2n<g_\circ\), so the norm forces each trace difference to vanish. Newton's identities give unipotence, and compact unitarity gives \(\gamma=1\). The compact quotient then has trace \(aV\).

In indefinite signature take \(n_0\ge8\), so \(1\le q\le p\). By \ref{B8.53e} the image has a finite orthonormal basis of continuous representatives \(F_\nu\), with
\[
 K(x,x)=\sum_\nu|F_\nu(x)|^2,\qquad D_{\rm count}=\int_{\mathscr F}K(x,x)\,dx.
 \tag{8.53.6}\label{eq:p8c5-trace}
\]
In the diagonal sum \eqref{eq:p8c5-kernel}, let \(S_{\rm near}(x)\) be the sum of the absolute values of the nonidentity terms with \(d_v(x_v^{-1}\gamma x_v)\le\eta\) at every real place. Let \(S_{\rm far}(x)\) be the corresponding sum over the remaining terms with \(\gamma\ne1\). Put \(\alpha=e^{-H/4}\), \(\beta=e^{-N}\). We prove
\[
 \begin{split}
 \int_{\rm moderate}S_{\rm near}(x)\,dx&\le\alpha D^\star,\qquad
 \int_{\rm moderate}S_{\rm far}(x)\,dx\le\beta D^\star,\\
 \operatorname{vol}(\mathrm{deep})&\le\alpha V,\qquad
 \int_{\rm deep}K(x,x)\,dx\le\beta D_{\rm count}.
 \end{split}
 \tag{8.53.7}\label{eq:p8c5-regions}
\]
By \ref{B8.53a} and \(|e(y_f)|\le e(1)\) from positive type,
\[
 |\phi(y)|\le a\prod_v\exp[-(h/8)\min(d_v(y_v)^2,d_v(y_v)/\sqrt{2q})].
 \tag{8.53.8}\label{eq:p8c5-point-decay}
\]

For a near term the trace differences have conjugates bounded by \(n(e^{n\eta}+1)\le3n<g_\circ\). Thus \(y=x^{-1}\gamma x\) is unipotent, and \(\|y_v-1\|_{\rm HS}^2=\|y_v\|_{\rm HS}^2-n\le n(e^{2\eta}-1)\). Put \(\delta=3\sqrt{n\eta}<1/4\). Then \(\gamma-1\) maps \(G_{\le t}\) into \(G_{\le\delta t}\) by \eqref{eq:p8c5-size}. If the least size is at least \(1/2\), the separation in \ref{B8.53a} forces the element to be the identity.

Otherwise the isotropic part of the filtration below \(1/2\) has sizes \(t_1<\cdots<t_s<1/2\) and positive increments \(j_i\) with \(J=\sum_i j_i\le q\). Complete it by orthogonal complements. The middle quotient between the two isotropic parts has sizes between \(1/2\) and \(2\), by self-duality. Use these cutoffs without subdividing that quotient. Order strictly lowering blocks by their distance down the filtration. In a fibre over earlier blocks, the difference at the next block is an integral map on the corresponding grades. The unitary relation determines the dual block from that difference because the product terms involve earlier blocks. A block paired with itself has skew-adjoint difference. A nonzero difference has geometric mean of operator norms at least the least target size divided by the largest source size, by the determinant estimate on its image and source quotient.

The scales \(\lambda\) for the flag density and the corresponding real dimensions \(m_\lambda\) are
\[
 \begin{array}{c|c}
 \lambda&m_\lambda\\ \hline
 t_i/t_k,\ t_it_k\ (i<k)&2fj_ij_k\ \text{each}\\
 t_i&2fj_i(n-2J)\\
 t_i^2&fj_i^2 .
 \end{array}
 \tag{8.53.20}\label{eq:p8c5-layer-scales}
\]
The separation scale \(\lambda'\) is the ratio of target lower to source upper cutoff. It equals the tabulated value for the blocks between small grades or their duals. For maps involving the middle quotient one can use \(\lambda'=t_i/2\), so in all cases \(\lambda\le3\lambda'\). Differences from a fixed representative in the fibre have maximum norm at most \(2\delta\). Packing disjoint balls of radius \(\lambda'/2\) gives \((1+4\delta/\lambda')^{m_\lambda}\le(1+12\delta/\lambda)^{m_\lambda}\) possibilities. Multiply by the density of \ref{B8.53b}. Since \(\lambda\le1\), \(\sum m_\lambda\le3Jnf\) and \(\delta qnf\le1/100\),
\[
 \prod_\lambda (1+12\delta/\lambda)^{m_\lambda}\lambda^{m_\lambda}
 \le e^{36\delta qnf}\le2.
 \tag{8.53.12}\label{eq:p8c5-near-cancel}
\]
\leanref{GlobalTheory.\allowbreak{}Count.\allowbreak{}layerProd}{Development.\allowbreak{}AbelianVarietyTheory.\allowbreak{}Definitions.\allowbreak{}SemilocalFrames.\allowbreak{}LayerProd}{32--35}
\leanref{GlobalTheory.\allowbreak{}NearCount.\allowbreak{}near\_\allowbreak{}count}{Development.\allowbreak{}AbelianVarietyTheory.\allowbreak{}Proofs.\allowbreak{}FlagMeasures.\allowbreak{}PowLePowRatio}{479--582}

On the moderate region \(t_i\ge\tau_*\). Each logarithmic integral is at most \(Y=22fq^2\log X\). Thus one flag type contributes at most \(2e^{-JH/2}Y^sD^\star\). From the size hypotheses, \(\log X\le\log n+\sqrt n+e^n+(1+A)e^{H/16}\), so \(Y\le e^{H/8}\) for large \(n\) at fixed \(A\). Summing over fewer than \(2^q\) types gives \(2^{q+1}e^{-3H/8}D^\star\le\alpha D^\star\).
\leanref{GlobalTheory.\allowbreak{}NearType.\allowbreak{}integral\_\allowbreak{}count\_\allowbreak{}dens}{Development.\allowbreak{}AbelianVarietyTheory.\allowbreak{}Proofs.\allowbreak{}FlagMeasures.\allowbreak{}NearTermsModeratePointsOneTypeBound}{40--97}
\leanref{GlobalTheory.\allowbreak{}NearMod.\allowbreak{}nearModerate}{Development.\allowbreak{}AbelianVarietyTheory.\allowbreak{}Proofs.\allowbreak{}FlagMeasures.\allowbreak{}OneTypeArith}{57--157}

By \ref{B8.53a}, the number of supported rational elements with all \(d_v\le R_v\) is at most
\[
 A_x e^{2n^2\sum_vR_v},\qquad A_x=(3\sqrt n/\tau_1(x)^2)^{2N}.
 \tag{8.53.9}\label{eq:p8c5-ball}
\]
Set \(T=h\eta^2/16\), and take \(C_0\ge\max(0,8\cdot10^{14}(A+4))\). Then \(T\ge1\) and \(2n^2\eta\le\min(\log2,T/2)\). Use \([0,\eta]\) and \((2^{j-1}\eta,2^j\eta]\), \(j\ge1\), as shells at each place. On the positive shells \(\min(d^2,d/\sqrt{2q})\ge\eta d\). Combining \eqref{eq:p8c5-point-decay} with the ball count leaves weights at most \(e^{-T2^{j-1}}\), with sum at most \(2e^{-T}\). The zero-shell weight is at most \(2\). Designating one positive-shell place gives
\[
 S_{\rm far}(x)\le aA_x\,2f3^f e^{-T}.
 \tag{8.53.10}\label{eq:p8c5-far}
\]
On the moderate region it suffices that
\[
 2N\log(3\sqrt n/\tau_*^2)+\log(2f)+f\log3+N\le T.
 \tag{8.53.21}\label{eq:p8c5-far-exponent}
\]
The right side is at least \(50(A+4)n^3 f\log(b\,\mathrm{rd})\). Divide the left side by \(n^3f\), use \(q\le\sqrt n\) and \(\log X\le\log n+\sqrt n+(1+A)\log(b\,\mathrm{rd})\), and obtain \(44+o(1)+44(1+A)n^{-1/2}\log(b\,\mathrm{rd})\). 
\leanref{GlobalTheory.\allowbreak{}FarTerms.\allowbreak{}farBound\_\allowbreak{}of\_\allowbreak{}ptBound}{Development.\allowbreak{}AbelianVarietyTheory.\allowbreak{}Proofs.\allowbreak{}FlagMeasures.\allowbreak{}SummableOfSupport}{42--124}
\leanref{GlobalTheory.\allowbreak{}Deep.\allowbreak{}farModerate\_\allowbreak{}of\_\allowbreak{}farBound}{Development.\allowbreak{}AbelianVarietyTheory.\allowbreak{}Proofs.\allowbreak{}FlagMeasures.\allowbreak{}FarTermsModerateRegionOfFlagBound}{53--113}

The deep volume follows from \ref{B8.53b}, since \(\tau_*<1/2\). Its trace bound is \ref{B8.53e}.
\leanref{GlobalTheory.\allowbreak{}CountFinal.\allowbreak{}countingTheorem\_\allowbreak{}of\_\allowbreak{}volumeStep}{Development.\allowbreak{}AbelianVarietyTheory.\allowbreak{}Proofs.\allowbreak{}FlagMeasures.\allowbreak{}UnitaryAutomorphicCountingOfCompactDualVolumeRat}{37--49}

Integrate \(K(x,x)-a\) using \eqref{eq:p8c5-regions} to get
\[
 |D_{\rm count}-D^\star|\le(2\alpha+\beta)D^\star+\beta D_{\rm count}.
 \tag{8.53.11}\label{eq:p8c5-count-combination}
\]
For large \(n_0\), \(\alpha,\beta\le e^{-H/16}/8\) and \(1-\beta\ge1/2\). Substitute \(D_{\rm count}\le D^\star+|D_{\rm count}-D^\star|\) and divide by \(1-\beta\).
\leanref{GlobalTheory.\allowbreak{}CountSum.\allowbreak{}countingTheorem\_\allowbreak{}of\_\allowbreak{}finite\_\allowbreak{}of\_\allowbreak{}bounds}{Development.\allowbreak{}AbelianVarietyTheory.\allowbreak{}Proofs.\allowbreak{}FlagMeasures.\allowbreak{}DiscreteSeriesCountingOfFiniteAndBounds}{34--142}
\leanref{GlobalTheory.\allowbreak{}abs\_\allowbreak{}sub\_\allowbreak{}le\_\allowbreak{}exp\_\allowbreak{}neg\_\allowbreak{}sixteenth\_\allowbreak{}mul}{Development.\allowbreak{}AbelianVarietyTheory.\allowbreak{}Proofs.\allowbreak{}FlagMeasures.\allowbreak{}AbsSubLeExpNegSixteenthMul}{7--43}
\end{proof}

\subsection*{Proof of \ref{B8.53a-jet}}\label{prf:B8.53a-jet}\phantomsection
\begin{proof}
Induct on the number \(d\) of noncompact derivatives of the difference. Compact derivatives act by covariance. Interchanging two noncompact derivatives gives a compact commutator and reduces the derivative order. Hence the order-\(d\) map \(T:\mathfrak p_\C^{\otimes d}\otimes S\to S\) is symmetric and \(K\)-equivariant. The differentiated Schmid equation kills every insertion of \(V_-\).

If \(z\perp\ker T\), it is symmetric with every one-slot slice in \(V_-^\perp\). Equivariance gives \((\Omega_K-c_K(\Lambda_S))z\in\ker T\), where \(c_K(\mu)=\|\mu+\rho_K\|^2-\|\rho_K\|^2\) for the trace pairing. The Casimir on the noncompact tensor factors is nonnegative. For \(\alpha=e_a-e_b\) with \(a<b\) of opposite signs, \(\|\alpha\|^2=2\), and the compact Casimir of the matrix-unit representation is \(n\). The one-slot mixed operator on the constituent of weight \(\Lambda_S+\alpha\) thus has scalar
\[
 \frac{c_K(\Lambda_S+\alpha)-c_K(\Lambda_S)-n}{2}
 =\langle\Lambda_S+\rho_K,\alpha\rangle+1-\frac n2 .
\]
Since \(\Lambda_S+\rho_K=\lambda+\rho_N\) with \(|\rho_N(i)|\le n/2\), this is at least \(g-2n\). Summing the mixed terms yields
\[
 \Re\langle z,(\Omega_K-c_K(\Lambda_S))z\rangle\ge2d(g-2n)\|z\|^2 .
 \tag{8.53.22}\label{eq:p8c5-jet-casimir}
\]
The left side is zero, so \(T=0\). Analyticity gives equality.
\leanref{Formalism.\allowbreak{}REAL.\allowbreak{}nPos\_\allowbreak{}of}{Formalism.\allowbreak{}RealPlaceAnalysis.\allowbreak{}RealPlaceTaylor.\allowbreak{}FirstSlotSliceSlotMixedCasimirZero}{517--547}
\leanref{GlobalTheory.\allowbreak{}Jet.\allowbreak{}jet\_\allowbreak{}map\_\allowbreak{}eq\_\allowbreak{}zero}{Development.\allowbreak{}AbelianVarietyTheory.\allowbreak{}Proofs.\allowbreak{}FlagMeasures.\allowbreak{}SymmetricKMapEqZeroOfVminusSlices}{114--226}
\end{proof}

\subsection*{Proof of \ref{B8.53a}}\label{prf:B8.53a}\phantomsection
\begin{proof}
Write \(y=k\exp Z\), with \(k\) compact and \(Z\) Hermitian and off the signature blocks, so \(d=\|Z\|_{\rm HS}\). For a highest-weight unit vector \(u\) of \(S\), the compression on \(\exp Z\) is the compact probability average of the rank-one projections onto \(lu\), multiplied by \(\dim S\) and weighted by
\[
 \prod_{i=1}^{n-1}\Delta_i(e^{l^{-1}Zl})^{-(\lambda_i-\lambda_{i+1}+1)},
 \tag{8.53.13}\label{eq:p8c5-compression}
\]
where \(\Delta_i\) is the leading principal minor. Near \(1\) use the corresponding minor character in triangular coordinates with the lowest-type determinant character. The diagonal weights give lowest-type covariance after compact averaging. At \(1\) the average is the identity by Schur orthogonality. For a direction \(X\), differentiation of the conjugated, adjointed minor character gives coefficients \((l^{-1}Xl)_{cd}\ell_{cd}\), with \(\ell_{cd}=0\) for \(d<c\) by triangular multiplication. For \(\sum E_{ab}\otimes G_{ab}\in V_-\), the corresponding coefficient \(\sum(l^{-1}E_{ab}l)_{cd}\langle lu,G_{ab}\rangle\) vanishes for \(c\le d\) by the lowering weights. The differentiated average therefore satisfies the Schmid equation, as does the compression. Apply \ref{B8.53a-jet} and continue along \(\exp(tZ)\), where all leading minors are positive.
\leanref{GlobalTheory.\allowbreak{}Jet.\allowbreak{}SchmidIntegral.\allowbreak{}eab\_\allowbreak{}integral\_\allowbreak{}eq\_\allowbreak{}zero}{Development.\allowbreak{}AbelianVarietyTheory.\allowbreak{}Proofs.\allowbreak{}FlagMeasures.\allowbreak{}HighestWeightEventuallyZero}{223--265}

Let \(Y=l^{-1}Zl\). By Cauchy--Binet, \(\Delta_i(e^{tY})=\mathbb E e^{tU_i}\), with \(U_i\) a sum of \(i\) eigenvalues weighted by squared eigenbasis minors. Differentiation gives \(\mathbb E U_i=0\) and \(\mathbb E U_i^2=\sum_{a>i\ge b}|Y_{ab}|^2\). Since \(\operatorname{rank}Y\le2q\), \(|U_i|\le\sqrt{2q}\,d\). For \(d>0\) take \(t=\min(1,1/(\sqrt{2q}\,d))\). The quadratic bound for \(e^z\) on \([-1,1]\) and concavity of \(u^t\) give
\[
 \log\Delta_i(e^Y)\ge t^{-1}\log(1+t^2\mathbb EU_i^2/3)
 \ge t\mathbb EU_i^2/4.
 \tag{8.53.23}\label{eq:p8c5-minor-decay}
\]
Summing the weights of \eqref{eq:p8c5-compression} over \(b\le i<a\) gives at least \(h\) for each nonzero \(Y_{ab}\) with \(a>b\). Their squared moduli sum to \(d^2/2\). The unweighted projection average times \(\dim S\) is the identity, so the norm is at most \(e^{-htd^2/8}\). At \(d=0\) use unitarity.

For the arithmetic ball a nonzero difference induces a full-rank map from its source quotient to its image. The least target size is at least \(\tau_1(x)\) and the largest source size at most \(\tau_1(x)^{-1}\) by duality. Taking determinants gives the separation. Rescale at \(v\) by \(e^{\bar R-R_v}\), \(\bar R=f^{-1}\sum R_v\). Their product is one, and the points are \(\tau_1(x)^2\)-separated in maximum norm inside radius \(\sqrt n\,e^{\bar R}\). Packing in real dimension \(2N\) gives \((1+2\sqrt n\,e^{\bar R}/\tau_1(x)^2)^{2N}\le A_xe^{2n^2\sum R_v}\).
\end{proof}

\subsection*{Proof of \ref{B8.53b}}\label{prf:B8.53b}\phantomsection
\begin{proof}
Let \(\mathcal V(I)=P_0V/\mu_f(K_f)\). In rank \(M\ge3\), choose a coordinate unit vector and its integral complement \(I'\) as in \ref{B8.49}. Unfold the Gram mean and cancel the real fibre \cite[\S52, Theorem~5, (41)]{weil-1965}. The finite product is \eqref{eq:p8c5-density-product}. In density-one measure the compact gauge ratio is \((2\pi)^M/(M-1)!\), from the unit-vector quotient with the Hilbert--Schmidt tangent metric. Consequently
\[
 \mathcal V(I)=\mathcal V(I')\,r^{f(M-1/2)}L_\eta(E,M).
 \tag{8.53.26}\label{eq:p8c5-gauge-step}
\]
\leanref{GlobalTheory.\allowbreak{}GaussMean.\allowbreak{}Step.\allowbreak{}volumeStep\_\allowbreak{}of\_\allowbreak{}siegelMean}{Development.\allowbreak{}AbelianVarietyTheory.\allowbreak{}Proofs.\allowbreak{}FlagMeasures.\allowbreak{}EnnrealVolumeRatioStepCancel}{77--119}

Descend to signature \((2,0)\). Put \(Y=r^f,\ L_0=1+\log Y\). In rank two Minkowski's theorem gives a nonzero vector whose squared Hermitian length is at most \(r\) at each real place, by using the product of complex two-dimensional balls of squared radius \(r\). The saturated line and its orthogonal line have common discriminant ideal \(\mathfrak n=\mathfrak a\mathcal O_E\) generated by real elements. The norm of \(\mathfrak a\) over \(F\) is at most the product of the squared lengths, hence at most \(Y\). Thus \(\mathrm N_E\mathfrak n=(\mathrm N_F\mathfrak a)^2\le Y^2\). Counting ideals in \(F\) gives at most \(YL_0^f\) choices for \(\mathfrak a\). In a field of degree \(a\), the ideal-count coefficients are bounded by those of \(\zeta(s)^a\), whose sum up to \(Y\) is at most \(Y(1+\log Y)^{a-1}\) by summing over the first \(a-1\) factors. Minkowski's ideal-class bound in \(E\) then gives at most \(YL_0^{2f}\) ideal classes \cite[Ch.~I, Theorem~6.3]{neukirch-1999}. Each line class has in addition at most \(2^f\) unit squareclasses. The unimodular gluing is the graph of a map of cyclic discriminant modules, determined by the image of a generator with at most \(\mathrm N_E\mathfrak n\) choices. Thus the class count is at most
\[
 YL_0^f(2^fYL_0^{2f})^2Y^2=4^fY^5L_0^{5f}.
 \tag{8.53.27}\label{eq:p8c5-rank-two-bound}
\]
\leanref{GlobalTheory.\allowbreak{}RankTwoCount.\allowbreak{}rankTwoClasses}{Development.\allowbreak{}AbelianVarietyTheory.\allowbreak{}Proofs.\allowbreak{}FlagMeasures.\allowbreak{}RankTwoDefiniteUnimodularHermitianClassBound}{225--254}

The rank-two compact gauge volume is \((8\pi^3)^f\). By the Euler factors, \(\prod_{M=3}^n L_\eta(E,M)\le(\prod_{l\ge3}\zeta(l))^f\). Summing to \(50\) and using the remainder bound \(50^{1-l}/(l-1)\) gives \(\prod_{l=3}^8\zeta(l)\le1.391\). For \(j\ge0\), the term at \(2\) and the sum and integral starting at \(3\) give
\[
 \zeta(9+j)\le 1+2^{-j}(2^{-9}+3^{-9}+3^{-8}/8)
 \le1+0.00235\,2^{-j}.
\]
The product is at most \(1.391e^{0.0047}<1.398<1.4\). With \(L_0\le Z^2\) and \(8\pi^3\cdot4\cdot1.4\le Z^{11}\), \eqref{eq:p8c5-gauge-step} gives \(\mathcal V(I)\le Z^{f(26+(n^2-4)/2)}\le Z^N\). Reduction in the lattice endomorphism ring gives \([K_f:K(g)]\le g^{2N}\).
\leanref{GlobalTheory.\allowbreak{}ZetaProd.\allowbreak{}zeta\_\allowbreak{}tail\_\allowbreak{}le}{Development.\allowbreak{}AbelianVarietyTheory.\allowbreak{}Proofs.\allowbreak{}FlagMeasures.\allowbreak{}ZetaProdSphereVolumeBounds}{81--99}

For the flag measure, bound the indicator of the filtration event by the sum over rational isotropic flags of the given type and sizes, and unfold along the stabilizer. Its unipotent blocks have the scales of \eqref{eq:p8c5-layer-scales}. The determinant directions give \(2fj_i\,dt_i/t_i\). The determinant-norm-one Levi block of rank \(j\) has volume
\[
 V_{\mathrm{GL},j}=(\operatorname{res}_1\zeta_E)(r/2)^{fj^2}
 \prod_{l=2}^{j}\zeta_E(l).
\]
For rank one the ideal classes and the unit fundamental domain give \((r/2)^f\operatorname{res}_1\zeta_E\). The Siegel unfolding of a primitive first column from rank \(j-1\) to \(j\) multiplies the volume by \((r/2)^{f(2j-1)}\zeta_E(j)\). The remaining unitary factor cancels by \eqref{eq:p8c5-gauge-step}.

Put \(S_l=\prod_{u=1}^l2\pi^u/(u-1)!\) and \(S=\sum_i j_i^2\). Set \(d_U=4\sum_{i<k}j_ij_k+2J(n-2J)+S\). At each real place, the compact quotient in Iwasawa coordinates has factor
\[
 2^{-d_U/2+nJ-J^2-J/2-S}\,
 \frac{S_pS_q}{S_{p-J}S_{q-J}\prod_i S_{j_i}} .
\]
The compact gauge volume in rank \(l\) is \(2^{(l^2-l)/2}S_l\). The compact \(\mathrm{GL}_j\) block contributes \(2^{j^2-j/2}S_j\) to the denominator, and the unipotent Jacobian contributes \(2^{-d_U/2}\). The arithmetic unipotent covolume contributes \(r^{fd_U/2}\). Multiply by \(2^{fj_i^2}V_{\mathrm{GL},j_i}\) for each Levi block to convert its measure to the flag gauge. Divide by the rank factors of \eqref{eq:p8c5-gauge-step} for \(M=n-2J+1,\ldots,n\). Their powers of \(r\) sum to \(2fnJ-2fJ^2\). Using \(d_U=2nJ-2J^2-S\) gives
\[
 \kappa=(r/2)^{-fd_U/2}2^{-f(nJ-J^2+J/2)}
 \left(\frac{S_pS_q}{S_{p-J}S_{q-J}\prod_i S_{j_i}}\right)^f
 \frac{(\operatorname{res}_1\zeta_E)^a\prod_i\prod_{l=2}^{j_i}\zeta_E(l)}
 {\prod_{M=n-2J+1}^n L_\eta(E,M)} .
 \tag{8.53.14}\label{eq:p8c5-flag-density}
\]
\leanref{GlobalTheory.\allowbreak{}FlagMeas.\allowbreak{}genLinVolume}{Development.\allowbreak{}AbelianVarietyTheory.\allowbreak{}Proofs.\allowbreak{}FlagMeasures.\allowbreak{}GenLinTamagawaVolume}{11--30}
\leanref{GlobalTheory.\allowbreak{}FiveFactors.\allowbreak{}prod\_\allowbreak{}five\_\allowbreak{}factors}{Development.\allowbreak{}AbelianVarietyTheory.\allowbreak{}Proofs.\allowbreak{}FlagMeasures.\allowbreak{}ProdFiveFactorsHermitianVolume}{80--127}
\leanref{GlobalTheory.\allowbreak{}FlagMeas.\allowbreak{}Bound.\allowbreak{}flagMeasureBound}{Development.\allowbreak{}AbelianVarietyTheory.\allowbreak{}Proofs.\allowbreak{}FlagMeasures.\allowbreak{}IsotropicFlagProbabilityBoundOfVolumeRatio}{64--103}

For the residue bound, let \(h_E,R_E,w_E\) be the class number, regulator and number of roots of unity. In an inverse ideal class choose an integral ideal \(J_0\). For any \(\sum_v x_v=0\), Minkowski applied to the discs \(|z_v|^2\le e^{x_v}(Y\mathrm N J_0)^{1/f}\) in \(J_0\) gives a nonzero \(a_0\), hence \(I_0=(a_0)J_0^{-1}\) of norm at most \(Y\). For fixed \(I_0\) the possible \(a_0\) form one unit orbit. The corresponding regions in the \(x_v\) are unit translates of a simplex of volume \(\log(Y/\mathrm N I_0)^{f-1}/(f-1)!\). They cover a regulator domain. Sum over ideals and classes and use \((\log u)^{f-1}/(f-1)!\le u^s\) for \(u\ge1,\ s>1\). This gives
\[
 h_ER_E\le Y^s\zeta_E(s)\le Y^s(s/(s-1))^{2f}.
 \tag{8.53.28}\label{eq:p8c5-regulator}
\]
Minkowski gives \(r\ge3/2\). In the class-number formula take \(s=1+1/\log r\) to get
\[
 \operatorname{res}_1\zeta_E=\frac{(2\pi)^fh_ER_E}{w_Er^f}
 \le [2\pi e(1+\log r)^2]^f\le(16\log r)^{2f}.
 \tag{8.53.29}\label{eq:p8c5-residue}
\]
\leanref{GlobalTheory.\allowbreak{}dedekindZeta\_\allowbreak{}residue\_\allowbreak{}le\_\allowbreak{}log\_\allowbreak{}rootDiscr\_\allowbreak{}pow}{Development.\allowbreak{}AbelianVarietyTheory.\allowbreak{}Proofs.\allowbreak{}FlagMeasures.\allowbreak{}DedekindZetaResidueLeSixteenLogRootDiscrPow}{48--75}

Put \(s_0=\sqrt n,\ l=\log n,\ c_0=\log(2\pi)\). In the compact factor the sphere factors \(2\pi^u/(u-1)!\) have consecutive ratio \(\pi/u\), so the maximum is at \(u=4\) and its logarithm is below \(3.49\). Use that bound for \(S_q/S_{q-J}\). For \(S_p/S_{p-J}\), use \((u-1)!\ge((u-1)/e)^{u-1}\) and \(u-1\ge p-J\ge n-2s_0\ge n/2\). The reciprocal factors \(1/S_{j_i}\) cost at most \(s_0l/4\) per unit of \(J\), from \(\log(u!)\le u\log u\) and \(j_i\le s_0\). The six log contributions in \eqref{eq:p8c5-flag-density}, divided by \(fJ\), are therefore bounded by
\[
 n/3,\quad0,\quad c_0-(n-2s_0)(l-c_0-1)+3.49+s_0l/4,\quad
 2(3+n),\quad2,\quad1 .
 \tag{8.53.30}\label{eq:p8c5-flag-coefficient}
\]
Here the first uses \(d_U\le2Jn\), \(r/2\ge3/4\). The last three follow from \eqref{eq:p8c5-residue}, \(\log r\le e^n\), \(\zeta_E(l)\le\zeta(l)^{2f}\) and \(L_\eta(E,M)^{-1}\le\zeta(M)^f\). For \(l\ge15\) their sum is at most \(-nl/2\). This gives \(\kappa\le e^{-JH/2}\). In \eqref{eq:p8c5-flag-measure} the middle-block factors include \(t_i^{2fj_i(n-2J)}\), and the other scales are at most one. The density with \(\kappa=1\) integrates to at most one by comparison with \(\prod_i 2fj_i t_i^{2fj_i-1}\,dt_i\) on the unit cube. A least size below \(1/2\) gives a nonempty flag. Summing over \(2^{J-1}\) compositions of \(J\) gives the claimed probability.
\end{proof}

\subsection*{Proof of \ref{B8.53d}}\label{prf:B8.53d}\phantomsection
\begin{proof}
At a single place the sum of squared derivatives in an orthonormal Lie basis is the quadratic form of \(2\Omega_K-\Omega_G\). A compact highest weight through order \(j\) satisfies \(\|\mu+\rho_K\|\le\|\lambda\|+\|\rho_N\|+\sqrt2 j\), since each derivative adds an adjoint weight. With \(a_0=\|\rho_N\|+\sqrt2j\), the Casimir expression is bounded by \(2(\|\lambda\|+a_0)^2-\|\lambda\|^2+n^3/12\le(\|\lambda\|+2a_0+0.3n^{3/2})^2\). Use \(\|\rho_N\|\le0.58n^{3/2}\) and \(\|\lambda\|\le\sqrt n\,b^2\) to bound this by \(B_j^2\). Summing over places gives
\[
 \|d\sigma(Z)u\|\le\sqrt f\,B_j\|Z\|_{\rm HS}\|u\|\qquad(u\in V_j).
 \tag{8.53.15}\label{eq:p8c5-differentiation}
\]

Use the ordered exponential cube of side \(s=1/B\) in the \(N\) real Lie coordinates and thicken by \(K(g)\). From \(2b^2\ge h\) and the gap we get \(12n^3s,8n^5s\le1\). Translated back to the identity, each differential column is its basis vector conjugated by the subsequent factors. The product estimate for the conjugations and the exponential coordinate map bounds the chart perturbation by \(n(e^{3n^2s}-1)\le1/2\). The quotient of two image points has operator norm at most \(e^{2n^2s}\le2\). The differential columns differ from the original columns by at most \(e^{2n^2s}-1\) in Hilbert--Schmidt norm. The resulting perturbation in dimension \(n^2\) has norm at most \(n(e^{2n^2s}-1)\), so the inverse determinant is at most \(\exp(2n^3(e^{2n^2s}-1))\le e\) at each place. These three inequalities follow from \(e^u-1\le2u\) for \(0\le u\le1\) and the two bounds on \(s\). Thus the chart integral costs at most \(e^fP_0/k_R\). Quotient overlaps are bounded by \((1+4\sqrt n/\tau_1(y)^2)^{2N}\) by the packing argument of \ref{B8.53a}.
\leanref{GlobalTheory.\allowbreak{}ProductChartAll.\allowbreak{}lintegral\_\allowbreak{}chart\_\allowbreak{}all\_\allowbreak{}le}{Development.\allowbreak{}AbelianVarietyTheory.\allowbreak{}Proofs.\allowbreak{}FlagMeasures.\allowbreak{}LintegralProductExpChartAllRealPlacesLe}{124--154}

In one coordinate the fundamental theorem of calculus and Cauchy--Schwarz give
\[
 |v(0)|^2\le(2/s)\int_0^s|v(t)|^2\,dt+2s\int_0^s|v'(t)|^2\,dt .
 \tag{8.53.32}\label{eq:p8c5-one-coordinate}
\]
Iterate in all coordinates using the one-place derivative bound \(B\) through order \(k+N\). Integration in the quotient and the isometries give, for \(u\in V_k\),
\[
 \|F(u,y)\|^2\le(6B)^Ne^f
 \left(1+\frac{4\sqrt n}{\tau_1(y)^2}\right)^{2N}
 \frac{DP_0}{k_R}\|u\|^2 .
 \tag{8.53.16}\label{eq:p8c5-cube}
\]
\leanref{GlobalTheory.\allowbreak{}SobolevBound.\allowbreak{}sobolev\_\allowbreak{}bound}{Development.\allowbreak{}AbelianVarietyTheory.\allowbreak{}Proofs.\allowbreak{}FlagMeasures.\allowbreak{}MeasurePreservingReindex}{239--285}

Append \(1/2\) to the small flag sizes starting at \(\tau\). At most \(q\) successive ratios have product \(2\tau\), so there are consecutive sizes \(t,t'\) with \(t/t'\le s_0=(2\tau)^{1/q}\), and \(t^2\le s_0/4\). Let \(R\) be the corresponding saturated isotropic sublattice, of rank \(j\le q\). Its largest size is at most \(t\), and the quotient above it has least size at least \(t'\).

For a Hom lattice from source of rank \(a\) and least size \(t'\) to target of rank \(b'\) and largest size \(t\), the trace-dual is the lattice of reverse maps taking the target into the source multiplied by the inverse different. On the nonzero image of a reverse map the determinant inequality gives geometric mean operator norm at least \(t'/(rt)\). The different contributes \(r\) because its norm over \(E\) is \(r^{2f}\). The Hilbert--Schmidt length across the real places is therefore at least \(\lambda=\sqrt f\,t'/(rt)\). The Euclidean transference inequality for covering radius and first dual minimum gives radius at most \(fab'/(2\lambda)\) in this normalization. The real dimension is \(2fab'\), and the trace pairs the real coordinates with a factor \(2\) compared with the Euclidean pairing. For the first layer \(\operatorname{Hom}_E(R^\perp/R,R)\) take \(a=n-2j,\ b'=j\). Lifting a block adds its adjoint and costs \(\sqrt2\) in norm. Since \((n-2j)j\le n^2/8\), the radius modulo the centre is at most
\[
 \frac{n^2\sqrt f\,r}{8\sqrt2}\frac{t}{t'} .
 \tag{8.53.17}\label{eq:p8c5-first-cover}
\]
\leanref{GlobalTheory.\allowbreak{}HomCover.\allowbreak{}exists\_\allowbreak{}hom\_\allowbreak{}near}{Development.\allowbreak{}AbelianVarietyTheory.\allowbreak{}Proofs.\allowbreak{}FlagMeasures.\allowbreak{}LatticePairVolumeBounds}{138--192}
\leanref{GlobalTheory.\allowbreak{}Count.\allowbreak{}FirstLayer.\allowbreak{}main}{Development.\allowbreak{}AbelianVarietyTheory.\allowbreak{}Proofs.\allowbreak{}FlagMeasures.\allowbreak{}NilradicalFirstLayerCoveringRadiusLe}{107--185}

The centre consists of skew-adjoint maps from \(\Lambda_x/R^\perp\) to \(R\). The source has least size at least \(1/t\). Approximate half of a central element by an integral \(\Phi\) and use \(\Phi-\Phi^*\). With \(a=b'=j\) the error after skew-symmetrization is at most \(j^2\sqrt f\,rt^2\), hence at most
\[
 \frac{n^2\sqrt f\,rs_0}{16}.
 \tag{8.53.18}\label{eq:p8c5-centre-cover}
\]
\leanref{GlobalTheory.\allowbreak{}Count.\allowbreak{}CentralCover.\allowbreak{}main}{Development.\allowbreak{}AbelianVarietyTheory.\allowbreak{}Proofs.\allowbreak{}FlagMeasures.\allowbreak{}NilradicalCenterIntegralCoveringRadius}{201--310}

Reduce first by \(2g\) times the integral first layer. Include the half-commutator in the central remainder and reduce by \(g\) times the integral centre. The linear term of the exponential is divisible by \(g\). The factor \(2g\) also makes its quadratic term divisible by \(g\). Modulo rational translation and \(K(g)\), the real orbit is covered by paths \(x\exp Z_1\exp Z_2\) with sum of exponent norms at most \(n^2\sqrt f\,grs_0(1/(4\sqrt2)+1/16)\le\rho\). For the adelic orbit, first use finite-place approximation in the additive first layer to move its finite part to the congruence subgroup. Then approximate in the additive centre to remove the central remainder. This reduces to the real cover.
\leanref{GlobalTheory.\allowbreak{}SmallFund.\allowbreak{}small\_\allowbreak{}fundamental\_\allowbreak{}set}{Development.\allowbreak{}AbelianVarietyTheory.\allowbreak{}Proofs.\allowbreak{}FlagMeasures.\allowbreak{}SmallFundamentalSetRealUnipotent}{135--218}

Along the paths, \(\tau_1(y)\ge\tau e^{-\rho}\ge\tau/2\). Thus \eqref{eq:p8c5-cube} starts an iteration on \(V_k\) with squared constant \((6B)^Ne^f(17\sqrt n/\tau^2)^{2N}DP_0/k_R\). If the bound on \(V_{j+1}\) is \(M'\|u\|\), integration of \eqref{eq:p8c5-differentiation} along the paths bounds \(F(u,y)-F(u,x)\) on \(V_j\) by \(\rho\sqrt f B_jM'\|u\|\). The constant term is zero, so averaging gives the same bound for \(F(u,x)\). The bound on the paths is then \(2\rho\sqrt f B_jM'\|u\|\). Iterate \(k\) times and square at \(w\).
\end{proof}

\subsection*{Proof of \ref{B8.53e}}\label{prf:B8.53e}\phantomsection
\begin{proof}
On each finite-dimensional subspace of counted intertwiners take the inner product \(T^*S=\langle T,S\rangle I\). An orthonormal basis consists of isometries with mutually orthogonal, \(K(g)\)-fixed images.

At each real place the compact-type projection writes any vector of the lowest type as a compact matrix-coefficient average of translates of \(w_v\). Convolve with a smooth bump supported near \(1\) and project to the type. The resulting operator on the finite-dimensional type can be chosen arbitrarily close to the identity and hence invertible there. This gives a smooth compactly supported group average for each lowest-type vector. Take products of these averages over the real places. A change of variables transfers the Lie derivative to the averaging function and gives smooth averages for the vectors of \(V_j\). Convolution with these functions and the normalized finite level indicator gives continuous representatives and the asserted derivative identity. For \(w\) itself, a fixed smooth average and Cauchy--Schwarz give
\[
 |f_i(x)|\le C\,\tau_1(x)^{-2N}\|T_iw\|,\qquad f_i=F_i(w).
\]
Cover the compact support of the average by finitely many chart translates of the kind in \ref{B8.53d}. The overlap count is \(O(\tau_1(x)^{-4N})\) by \ref{B8.53a}. Taking a square root gives the bound. The constants depend on the count data and the chosen averages.
\leanref{GlobalTheory.\allowbreak{}SmoothRepsType.\allowbreak{}repsExist}{Development.\allowbreak{}AbelianVarietyTheory.\allowbreak{}Proofs.\allowbreak{}FlagMeasures.\allowbreak{}ExistsContinuousRowRepresentativesLowestKType}{87--137}
\leanref{GlobalTheory.\allowbreak{}StepAverage.\allowbreak{}exists\_\allowbreak{}smooth\_\allowbreak{}average\_\allowbreak{}of\_\allowbreak{}isKTypeOf}{Development.\allowbreak{}AbelianVarietyTheory.\allowbreak{}Proofs.\allowbreak{}FlagMeasures.\allowbreak{}ExistsSmoothAverageLowestKType}{48--84}

Schur orthogonality gives \(f_i=R_\phi f_i\). The growth bound and \eqref{eq:p8c5-point-decay} give an absolutely convergent continuous convolution. In the radical constant term, real translation in the supported finite level changes \(\tau_1^{-1}\) by at most the geometric mean of the real operator norms and their inverses. The exponential decay in \eqref{eq:p8c5-point-decay} at large radius dominates this growth and the radical volume growth under \eqref{eq:p8c5-count-gap}. Thus the constant-term integral and convolution can be interchanged. The real integral of the translated discrete-series coefficient along a proper radical vanishes \cite{wallach-1992}. The constant term of \(f_i\) is zero. Its translates generate \(T_i(H_\sigma)\), and the cuspidal subspace is closed and invariant.
\leanref{GlobalTheory.\allowbreak{}RowCuspKType.\allowbreak{}row\_\allowbreak{}isCuspFn}{Development.\allowbreak{}AbelianVarietyTheory.\allowbreak{}Proofs.\allowbreak{}FlagMeasures.\allowbreak{}CountedVectorsAreCuspFormsSingleLowestKType}{85--114}

For a basis of size \(D\), apply \ref{B8.53d}, \ref{B8.53c} and \eqref{eq:p8c5-absolute-volume} with \(Z=X\). On the deep region they give
\[
 \sum_{i=1}^D|F_i(w)(x)|^2\le e^{-N}D/V.
 \tag{8.53.34}\label{eq:p8c5-deep-point}
\]
On the moderate region \eqref{eq:p8c5-cube} for a single normalized intertwiner bounds the evaluation norm uniformly, since \(\tau_1\ge\tau_*\). By Bessel the summed squared evaluations are then at most some constant \(C_1^2\) independent of \(D\). Their total integral is \(D\), so \((1-e^{-N})D\le VC_1^2\). Hence the full space is finite dimensional.

Its finite basis gives a continuous projector kernel. On compact sets the finite components range over finitely many cosets of the open level, and the real radii increase by bounded amounts. The ball and shell estimates thus give locally uniform convergence of the orbit sum \eqref{eq:p8c5-kernel}. On compactly supported tests unfolding identifies it with the projector. The continuous kernels therefore agree, and integrating \eqref{eq:p8c5-deep-point} with \(D=D_{\rm count}\) gives \eqref{eq:p8c5-deep-integral}.
\leanref{GlobalTheory.\allowbreak{}KernelExpansion.\allowbreak{}kernelB}{Development.\allowbreak{}AbelianVarietyTheory.\allowbreak{}Proofs.\allowbreak{}FlagMeasures.\allowbreak{}IntegrablePh}{166--224}
\end{proof}

\subsection*{Proof of \ref{B8.54}}\label{prf:B8.54}\phantomsection
\begin{proof}
The splitting follows from \ref{B8.51a}. Let \(K\) stabilize the self-dual overlattice of the definite first space, and \(U\subset K\) be the column level. Its unitriangular factors give (8.54.1), and fixing a nonzero vector at each of the \(f\) places over \(T_u\) gives
\[
 [K:U]=I_{01}(T_u^k-1)^f.
 \tag{8.54.4}
\]
At the prime in \eqref{eq:p8c5-stabilizer-prime}, a stabilizer for the action of a genus automorphism group on its level cosets embeds in the pro-\(g_\nu\) preimage of a unipotent group. If nontrivial, it contains an element of order \(g_\nu\), which in rational dimension \(2fk\) would require the cyclotomic factor of degree \(g_\nu-1\). Thus the action is free. Summing the level index divided by each genus automorphism order gives \([K:U]M_k\) by \ref{B8.50}.
\leanref{GlobalTheory.\allowbreak{}ColClassNumber.\allowbreak{}mass\_\allowbreak{}selfDualLattice\_\allowbreak{}firstSpace\_\allowbreak{}eq}{Development.\allowbreak{}AbelianVarietyTheory.\allowbreak{}Proofs.\allowbreak{}FoundationProofs.\allowbreak{}MassSelfDualLatticeFirstSpaceEq}{93--100}
\leanref{GlobalTheory.\allowbreak{}ColClassNumber.\allowbreak{}classNumber}{Development.\allowbreak{}AbelianVarietyTheory.\allowbreak{}Proofs.\allowbreak{}FoundationProofs.\allowbreak{}MassSelfDualLatticeFirstSpaceEq}{102--130}
\leanref{GlobalTheory.\allowbreak{}ColClassNumber.\allowbreak{}classNumber\_\allowbreak{}closed}{Development.\allowbreak{}AbelianVarietyTheory.\allowbreak{}Proofs.\allowbreak{}GlobalInterpolation.\allowbreak{}ClassNumberClosed}{10--13}
\end{proof}

\subsection*{Proof of \ref{B8.55b}}\label{prf:B8.55b}\phantomsection
\begin{proof}[Proof (outline)]
Section base change and uniformization imply that at every complex realization of an algebraic test point, evaluation of any section of the given weight is one common scalar times its value in a complex frame. Suppose \(\sum_i z_iR(s_i)=0\), with \(z_i\in\C\). Evaluate this relation at each such test point and apply any \(K\)-linear functional \(\lambda:\C\to K\) to the complex coefficients. Descent from complex points of a reduced, finitely generated \(K\)-algebra gives the identity of algebraic section rules
\(\sum_i\lambda(z_i)s_i=0\).
Independence over \(K\) gives \(\lambda(z_i)=0\) for every \(\lambda\). Linear functionals separate vectors, so every \(z_i=0\). Thus the full frame readouts are linearly independent over \(\C\).

Now suppose \(\sum_i z_iE(s_i)=0\). The function \(\sum_i z_iR(s_i)\) is a holomorphic cusp section with the prescribed arithmetic-group and Borel equivariance. Every frame in every component is carried to a comparison frame by an arithmetic-group element and an invertible Borel change of frame. The equivariance multipliers are nonzero, so vanishing on the comparison frames implies vanishing throughout every frame domain. Outside those domains the readouts are zero by definition. The independence just proved gives \(z_i=0\) for every \(i\). The complete proof is \Lean{GlobalTheory.\allowbreak{}ClassSumSpan.\allowbreak{}atDict\_\allowbreak{}indep}.
\end{proof}

\subsection*{Proof of \ref{B8.56f}}\label{prf:B8.56f}\phantomsection
\begin{proof}
The endomorphism ring contains the maximal order \(\Z[\omega]\). The CM invariant \(J\) is integral. Global reciprocity identifies the ideal-class quotient with the Galois group of the Hilbert class field, and the class invariants of the maximal order belong to this field \cite[§11, Theorem~11.1]{cox-2013}. In particular \(J\in K^{\rm ur}\) at \(p\nmid d\). In terms of Frobenius, the argument used here is to take the finite quotient recording both the ideal class and the permutation of the class invariants. At an inertia element the class coordinate is trivial. By Lemma~\ref{B8.42p} one can represent its conjugacy class by degree-one Frobenius avoiding the discriminant and the differences of distinct invariants. The CM ideal-isogeny congruence at this prime acts by a principal ideal, hence fixes the reduced invariants. Their reductions are distinct, so the permutation is trivial.

Normalize \(v_p(p)=1\). We claim
\[
  3\mid v_p(J)\quad(J\ne0),\qquad
  2\mid v_p(J-1728)\quad(J\ne1728).
  \tag{8.56.7}\label{eq:p8c5-cm-divisibilities}
\]
For \(J\ne0,1728\) start with \(g_2=g_3=c=27J/(J-1728)\). Scale by \(\mu^4,\mu^6\) to discriminant one. If \(v_p(J)>0\), then \(12v_p(\mu)=-2v_p(J)\), and the reduction has \(g_2=0,\ g_3\ne0\). Choose a cyclic CM kernel of odd order \(q\equiv2\bmod3\), prime to \(p\). For example, for \(\Q(\sqrt{-d})\ne\Q(\sqrt{-3})\), choose \(q\) split in the former field and inert in the latter, avoiding \(2,3,p,d\), by Lemma~\ref{B8.42p}. Take the kernel of a prime ideal over \(q\). In the excluded field the maximal-order invariant is \(0\). The chosen kernel is defined over \(K^{\rm ur}\) on the model with invariants \(c,c\): this field contains the local CM field and the endomorphisms are determined by their action on the invariant differential. Its nonzero points form \(t=(q-1)/2\) opposite pairs, with \(3\nmid t\). Their abscissas on the scaled model are integral units. Indeed prime-to-\(p\) torsion reduces injectively, and a point of abscissa zero on the reduction has order three by the tangent law. The product of the unscaled abscissas is Galois invariant. Its valuation \(-2t v_p(\mu)=t v_p(J)/3\) is an integer.

If \(v_p(J-1728)>0\), then \(4v_p(\mu)=v_p(J-1728)\), and the reduction has \(g_3=0,\ g_2\ne0\). Take the kernel of \(\sqrt{-d}\). Its order \(d\) is odd and prime to \(p\), and \(t=(d-1)/2\) is odd. On this reduction a point of abscissa zero has order two. The same product gives \(-t v_p(J-1728)/2\in\Z\). This proves \eqref{eq:p8c5-cm-divisibilities}, including the exceptional invariants since \(1728\) is a unit. Hensel lifting of the unit parts gives \(G_2^3=J/1728,\ 27G_3^2=J/1728-1\) in \(W\), using \(2,3\in W^\times\) and the algebraically closed residue field. These invariants have the required \(j\) and discriminant.
\leanref{ReciprocityTheory.\allowbreak{}ChartRestriction.\allowbreak{}L1.\allowbreak{}valDiv3}{ReciprocityTheory.\allowbreak{}ReductiveGroups.\allowbreak{}UnramifiedIntegrals.\allowbreak{}NormWeierstrassCoeffSmallJ}{98--186}
\leanref{ReciprocityTheory.\allowbreak{}ChartRestriction.\allowbreak{}L1.\allowbreak{}valDiv2}{ReciprocityTheory.\allowbreak{}ReductiveGroups.\allowbreak{}UnramifiedIntegrals.\allowbreak{}Norm1728EqOneOfFiveLe}{45--149}

Now suppose \(p\) inert. Let \(\sigma\) be Witt Frobenius and \(\tau f=\sigma^2(f)(T^{p^2})\). The CM endomorphism \(f^{-1}(\omega f)\) is integral on the good model. One can obtain this directly by extending the endomorphism to the good elliptic scheme and completing at zero. Its tangent and its conjugate differ by a unit. Define \(c_j\), in increasing order, by cancelling degree \(p^{2j+2}\) in
\((p-\sum_{i=0}^jc_i\tau^{i+1})f\). Each \(\tau^{j+1}f\) has leading coefficient one. We show that all \(c_j\in W\) and that all remaining coefficients are in \(pW\).

Use induction in the degree \(n\), retaining the terms with \(p^{2i+2}<n\), and let \(a_n\) be the current coefficient. The invariant differential gives \(k[T^k]f\in W\). Consequently substituting two integral zero-constant series congruent modulo \(p\) in any Frobenius twist of \(f\) gives results congruent modulo \(p\): in the binomial difference the term with \(r\ge1\) factors of \(p\) is divisible by \(p^r/r\). Apply this to the formal group law and the CM endomorphism after raising to \(p^{2i+2}\). Thus the additive defect and the \(\omega\)-linearity defect of every retained twist lie in \(pW\). Since the lower coefficients of the remainder lie in \(pW\), its additive defect gives \(a_n\binom nk\in pW\) for \(0<k<n\). Outside powers of \(p\) a middle binomial coefficient is a unit. At a power of \(p\) one has valuation one, so \(a_n\in W\). At an odd power \(p^{2i+1}\) the second defect gives \((\omega^n-\omega)a_n\in pW\), hence \(a_n\in pW\). At the even powers take \(c_i=a_n\). The constant and linear terms satisfy the assertion directly.

If \(c_0\in pW\), all coefficients of \(pf\) through degree \(p^2\) belong to \(pW\). From \(pf=f([p](T))\) the same follows successively for \([p](T)\), since its lower terms are then multiples of \(p\) and \(k[T^k]f\in W\). The special-fibre multiplication-by-\(p\) map would have local degree greater than its degree \(p^2\). Thus \(c_0\in W^\times\).

Multiply the type \(p-\sum c_i\tau^{i+1}\) on the left by \(\sum t_i\tau^i\), \(t_0=1\). To cancel the coefficients of \(\tau^{j+1}\) for \(j\ge1\), iterate the map on sequences given successively by
\[
 t'_0=1,\qquad
 t'_j=\frac{p\,t_{j+1}-\sum_{i<j}t'_i\,\sigma^{2i}(c_{j-i})}
                {\sigma^{2j}(c_0)} .
\]
The denominators are units. Induction on \(j\) shows that agreement of two inputs modulo \(p^k\) gives agreement of their outputs modulo \(p^{k+1}\). Completeness gives a fixed point, reducing the type to \(p-a\tau\), \(a=c_0-pt_1\in W^\times\). Solve \(\epsilon=a\sigma^2(\epsilon)\) in \(W^\times\): solve first in the nonzero residue classes, then lift by solving an additive equation \(z-z^{p^2}=b\) in \(\overline{\F}_p\) at each precision. Both \(f\) and \(\epsilon\ell\) now solve the same functional equation with integral inhomogeneous terms, since \((1-a\tau/p)f\) is integral and \((1-a\tau/p)(\epsilon\ell)=\epsilon T\). The functional-equation lemma for \(W\), its Frobenius square and the ideal \((p)\) gives \(f^{-1}(\epsilon\ell)\in W[[T]]\) \cite[Lemma~2.3]{hazewinkel-1979}. Its tangent is \(\epsilon\), so its inverse is integral as well.
\leanref{GlobalTheory.\allowbreak{}Honda.\allowbreak{}exists\_\allowbreak{}even\_\allowbreak{}type}{Development.\allowbreak{}AbelianVarietyTheory.\allowbreak{}Proofs.\allowbreak{}DecoratedLatticeProofs.\allowbreak{}ExistsHondaTypeEvenFrobeniusPowers}{259--395}
\leanref{GlobalTheory.\allowbreak{}Honda.\allowbreak{}OneTerm.\allowbreak{}exists\_\allowbreak{}one\_\allowbreak{}term}{Development.\allowbreak{}AbelianVarietyTheory.\allowbreak{}Proofs.\allowbreak{}DecoratedLatticeProofs.\allowbreak{}HondaTypeReductionToOneTerm}{109--127}
\leanref{GlobalTheory.\allowbreak{}TwV.\allowbreak{}Lt.\allowbreak{}exists\_\allowbreak{}genusOne\_\allowbreak{}formal\_\allowbreak{}iso\_\allowbreak{}of\_\allowbreak{}unramified\_\allowbreak{}good\_\allowbreak{}reduction}{Development.\allowbreak{}AbelianVarietyTheory.\allowbreak{}Proofs.\allowbreak{}SatakeParameters.\allowbreak{}ExistsGenusOneFormalIsoOfUnramifiedGoodReduction}{17--24}
\end{proof}

\subsection*{Proof of \ref{B8.56}}\label{prf:B8.56}\phantomsection
\begin{proof}
Use global reciprocity in the form of Proposition~\ref{B1.15b} and Definition~\ref{B1.15ac}. The idèle map onto the abelianized Galois group has kernel the closure of the principal idèles times the connected archimedean subgroup. Local units map onto the inertia image, and a uniformizer acts by inverse arithmetic Frobenius modulo inertia. Passage to the ideal-class quotient gives the reciprocity input of Lemma~\ref{B8.56f}. Thus it supplies the good CM model, and that lemma supplies the inert formal parameter. For the extension parameters use Theorem~\ref{B9.46} at each specified finite place and prolongation, with the correspondence of Definition~\ref{B9.47}. Its rank-one normalization is Theorem~\ref{B9.50}. In this proposition the correspondence enters through the norm choices of the extension in \((8.6)\). Given that extension, compatibility of the bundle identifies its underlying point and its reciprocity datum with those of a constructed point. Substituting this identification identifies the period lattice, determinant classes and selected column kernel with their standard arithmetic versions.
\leanref{CompanionTheory.\allowbreak{}FamInh.\allowbreak{}globalArtin}{CompanionTheory.\allowbreak{}AutomorphicTheory.\allowbreak{}PlaceSeparation.\allowbreak{}IsOpenQuotientMapArt0}{233--234}
\leanref{UpperTheory.\allowbreak{}Zero.\allowbreak{}LLC.\allowbreak{}recTExists}{Development.\allowbreak{}AutomorphicArithmetic.\allowbreak{}Foundations.\allowbreak{}LocalLanglands.\allowbreak{}LocalCentralUnits.\allowbreak{}RecTExists}{20--66}
\leanref{GlobalTheory.\allowbreak{}Svd.\allowbreak{}LinkStd.\allowbreak{}stdTie\_\allowbreak{}of\_\allowbreak{}eq\_\allowbreak{}alt}{Development.\allowbreak{}AutomorphicArithmetic.\allowbreak{}ThetaCorrespondence.\allowbreak{}ThetaPairLifts.\allowbreak{}ThirdPlaceSeries.\allowbreak{}LinkRStdTie}{16--19}

For the first conclusion fix a family and then a member. Mark the column representatives as in the proof of Proposition~\ref{B8.56c}. The only reading changed by this operation is the choice of representative of the finite class datum. The equalities of its two endpoint lattices and of the resulting theta functions with the specified column readings are part of the construction in that proof. For a used twist and a permitted \((u,w)\) outside its exceptional primes take the levels and the chain obtained there. For the other indices take the original levels and chain. At the reference twist the level-zero identification agrees with the original one. Once \(x\ge71/\varepsilon\), \(2r_G\le k\): otherwise the lower rank bound gives \(x\le k<2r_G\), while \(r_G\le8(4+8\varepsilon)/\varepsilon+2\). This gives \(x\varepsilon\le64+132\varepsilon<71\) since \(\varepsilon<1/52\). The stabilizer at zero is trivial. Indeed forgetting the additional twist frames identifies its rational component with an element of the reference initial stabilizer, which is trivial by the initial level condition. With that component fixed, the frame and the empty initial index determine the remaining component. Proposition~\ref{B8.56c} therefore supplies both admissibility and comparison on the chosen chain.
\leanref{GlobalTheory.\allowbreak{}Svd.\allowbreak{}jointsAt\_\allowbreak{}all\_\allowbreak{}std}{Development.\allowbreak{}AbelianVarietyTheory.\allowbreak{}Proofs.\allowbreak{}SatakeParameters.\allowbreak{}JointsAtAllStd}{11--29}
\leanref{GlobalTheory.\allowbreak{}TwV.\allowbreak{}modelsOK\_\allowbreak{}of\_\allowbreak{}joints}{Development.\allowbreak{}AbelianVarietyTheory.\allowbreak{}Proofs.\allowbreak{}SatakeParameters.\allowbreak{}ThetaLocalModelsOfSplitUnramifiedJoints}{11--29}
\leanref{GlobalTheory.\allowbreak{}TwU.\allowbreak{}two\_\allowbreak{}rG\_\allowbreak{}le\_\allowbreak{}k\_\allowbreak{}large}{Development.\allowbreak{}AbelianVarietyTheory.\allowbreak{}Proofs.\allowbreak{}SatakeParameters.\allowbreak{}TwoRGLeKOfSize}{38--46}

We record the size conditions used on this chain. Write \(w_z\) for the weight at stage \(j+1\), \(s_z\) for the diagram size and \(r_{z,0},r_{z,1}\) for its first two rows. Write \(c_z\) for its first row off the top and its second row on the top. Let \(\Delta\) be the spacing of the two interior word-coordinate choices, as in \ref{B8.02c}. With \(B_j(z)\) the filling bound in \eqref{eq:p4c2-shape}, the conditions are
\[
\begin{gathered}
 w_z>0,\qquad 0\le s_z-(k-1)w_z\le4w_z\quad(z\text{ off the top}),\\
 r_{z,0}\ge2w_z,\qquad
 0\le s_z-r_{z,0}-(k-3)w_z\le4w_z\quad(z\text{ on the top}),\\
 s_z\le a^2,\quad 2\log_5(a^2)+2\le k-1,\quad 1<a,\quad X''\le a,\\
 c_z\le b+\eta_{\rm Sz}+1+12k\Delta,\quad c_z\le B_j(z)+c_G,\qquad
 \bar\alpha_t\le r_{(0,t),0}-r_{(0,t),1}\le2a .
\end{gathered}
\tag{8.56.3}\label{eq:p8c5-entry-shape}
\]
Here the lower gap bound is on the marked first-sheet positions. One may choose a threshold depending on \(\varepsilon\) before the family: take \(x\ge4\), \(x^{\varepsilon/5}\ge3\), and \(x^6\le\exp(x^{\varepsilon/3})\) throughout the remaining half-line. The scale bounds then give \(kN_m\le k\Delta\), \(4k^2\Delta\le w''\), \(m(t+1)\eta_{\rm Sz}\le w''\) and \(1\le w''\le w_z\). Sum the long rows from the weight windows and the appended rows, each in \([0,\eta_{\rm Sz}]\). Off the top this gives the interval about \((k-1)w_z\) in \eqref{eq:p8c5-entry-shape}. On the top remove the first row before summing and use its prescribed length \(a_t+w_z-\kappa_{j+1}\), with \(\kappa_{j+1}=(k+(2j+1)(t+1))/2\). This gives the interval about \((k-3)w_z\). The top weight window and \(x\ge4\) give \(r_{z,0}\ge2w_z\); subtracting the second-row lower bound gives \(r_{z,0}-r_{z,1}\le2a\). The two cap bounds use that same second row on the top and the first row elsewhere. The filling lower gap is the one in Definition~\ref{B3.47}. These estimates use the diagram of the member and therefore apply to every entry realization at the point. The split-prime threshold is at least three by \(x^{\varepsilon/5}\ge3\).
\leanref{GlobalTheory.\allowbreak{}ShapeUniform.\allowbreak{}shape\_\allowbreak{}uniform}{Development.\allowbreak{}AbelianVarietyTheory.\allowbreak{}Proofs.\allowbreak{}SatakeParameters.\allowbreak{}ThreeLeSplitAt}{389--431}

For an unrepresented position of the filling the orbit lies on the second sheet. At \(p\le X''\) take the lower exponent equal to the orbit base term less its error, as defined in Theorem~\ref{B4.16}. The orbit has even length at least two, \(p\ge5\), and \(h_P\ne0\). The position identification respects the cyclic order and the row lengths. Apply the refined cyclic estimate of Proposition~\ref{B4.01b} with \(C_6=11,C=1\) to the integral cofactors. Its weight hypotheses are \eqref{eq:p8c5-entry-shape}; the fixed constants satisfy \(11\le12\) and \(3\le\min(1+d_{\ell1},3)\). On a second-sheet orbit there is no first-sheet marked gain. Summing one term per complementary orbit gives the required lower bound as in Proposition~\ref{B4.11a}(2). For \(p>X''\) take zero by the one-place bound in Proposition~\ref{B8.56c}. Its denominator exponents sum to \((m-1)v_p(D_{\Gamma,u})\). Together with the comparisons this gives \eqref{eq:p8c5-twist}.
\leanref{GlobalTheory.\allowbreak{}TwU.\allowbreak{}LowF.\allowbreak{}lowAt\_\allowbreak{}F}{Development.\allowbreak{}AbelianVarietyTheory.\allowbreak{}Proofs.\allowbreak{}SatakeParameters.\allowbreak{}LocalModelBoundPlaceOverSecondEmbedding}{66--95}
\leanref{GlobalTheory.\allowbreak{}Proof.\allowbreak{}UnitaryWeightModels.\allowbreak{}rowTwU\_\allowbreak{}of\_\allowbreak{}models}{Development.\allowbreak{}AbelianVarietyTheory.\allowbreak{}Proofs.\allowbreak{}SatakeParameters.\allowbreak{}UnitaryWeightModels}{15--31}

For the second conclusion take the entry of the selected bundle. At a level-free place the determinant of a class basis measures exactly the local index of that lattice in the reference lattice. Evaluating the type of the step character on this determinant cancels the corresponding basis determinant powers in the class scalar; its finite-order factor has norm one. This gives the unit assertion for every tuple. For a used tuple the matrix block between two classes is defined over the coefficient field in the two rational row readings. Applying that block to the source reading gives the required rational target reading. The sum of the blocks divided by the source stabilizer orders is the lift on the embedded complex spans, at the reference tuple as well as at the used twists.
\leanref{GlobalTheory.\allowbreak{}KerJ.\allowbreak{}unUnit\_\allowbreak{}of\_\allowbreak{}mtType}{Development.\allowbreak{}AbelianVarietyTheory.\allowbreak{}Proofs.\allowbreak{}SatakeParameters.\allowbreak{}KernelEntryUnramifiedUnit}{12--17}
\leanref{GlobalTheory.\allowbreak{}PdvClose.\allowbreak{}at\_\allowbreak{}of\_\allowbreak{}true\_\allowbreak{}columns}{Development.\allowbreak{}AbelianVarietyTheory.\allowbreak{}Proofs.\allowbreak{}SatakeParameters.\allowbreak{}AtOfTrueColumns}{30--47}

For the inert even-order and signal branches, use Lemmas~\ref{B8.56a} and \ref{B8.56b} on the normalized translates, using the same curve for all pairs at \(w\). In trace-dual coordinates average the paired phase over the residue-field torus as in Lemma~\ref{B8.56d}. Its order \(p^{n_P}-1\) is a unit, so replacing the phase by its average changes the cofactor by an integral exponential. This gives the all-place formula with sign \(-1\) and the common coordinate unit of the curve. The trace construction works with every choice of orbit representatives. At \(h_P=0\) the orbit is a singleton. The Hermitian form in the conjugate coordinate is obtained by substituting the same variable into the two paired coordinates with a unit multiplier. Polarization gives half the paired phase, as in the one-position expansion of Proposition~\ref{B4.11a}. For nonzero odd-order \(h_P\), use the trace form of the odd unramified orbit extension. A unit on its fixed side is a norm of a unit on the quadratic side, by the residue norm and Hensel lifting. Diagonalizing the unimodular form identifies it with the orthogonal trace model. The integral unitary substitution in Lemma~\ref{B8.56a} then gives the all-real orbit product. These last two arguments use the inert residue class and their respective orders of \(h_P\), independently of the all-place branch. Lemma~\ref{B8.56d}, with \eqref{eq:p8c5-entry-shape} and \(d_{\ell1}\ge5,d_{\ell2}\ge4\), gives the asserted minor bound for every series with the specified first-sheet expansion.
\leanref{GlobalTheory.\allowbreak{}Ggeo.\allowbreak{}frobeniusConjugation\_\allowbreak{}conditions\_\allowbreak{}adaptedKernelDatum\_\allowbreak{}alt}{Development.\allowbreak{}AbelianVarietyTheory.\allowbreak{}Proofs.\allowbreak{}SatakeParameters.\allowbreak{}FrobeniusConjugationConditionsAdaptedKernelDatum}{34--62}
\leanref{UpperTheory.\allowbreak{}lattice\_\allowbreak{}pairs\_\allowbreak{}odd\_\allowbreak{}degree\_\allowbreak{}isometric}{Development.\allowbreak{}AutomorphicArithmetic.\allowbreak{}OrbitLabellings.\allowbreak{}LatticePairsOddDegreeIsometric}{69--87}

For the jet assertions, write \(P\) for the corrected curve phase in the formal parameter, \(l\) for the good-model logarithm and \(S=(P\circ l^{-1})''(0)\). At odd primes the two addition defects and the adjoint defect give the following derivative calculation. Derivatives of order at least three come from the defect in three additive variables: differentiate once in each of the first two variables and the remaining times in the third. The exponential of the defect is integral in the formal parameter, and logarithms of integral one-units have integral derivatives at zero. Passing to logarithmic coordinates preserves integral derivatives since \(l'\) is an integral unit with tangent normalization one. Reflection controls the first derivative since \(2\) is invertible. Differentiating the adjoint defect once in each variable gives \((\alpha-\bar\alpha)S\) integral for a CM multiplier \(\alpha\) for which the endomorphism and the norm law hold. At \(p\ge5,\ p\nmid d\) take a prime-to-\(p\) multiplier with \(\alpha-\bar\alpha\) a unit, for instance an integral translate of \((1+\sqrt{-d})/2\) avoiding the two possible zero residues of its norm. At \(p=3\), the split-prime threshold gives \(d\equiv2\bmod3\). Use \(\alpha=\sqrt{-d}\). Its norm and \(\alpha-\bar\alpha\) are units, and the Weierstrass correction \((\alpha^2-1)b_2/12=-((d+1)/3)b_2/4\) is integral. In both cases the isogeny norm law follows from the sigma product over the nonzero kernel pairs, whose abscissas are integral. This proves integrality of all curve jets at these primes.

At \(p=d\) use a unit CM multiplier with \((\alpha-\bar\alpha)^2=-d\). The same calculation gives \(|S|_p\le p^{1/2}\) and integral jets after removing \(SZ^2/2\). In the two dual frames the square form polarizes to
\(2\sum_tu_t\sum_{a,b}X_{(0,t),a,b}X_{(1,t),a,b}\), with \(|u_t|_p\le1\) since the frame maps preserve the respective localized lattices. The scalar \(S\) is a jet of the one curve, and the \(u_t\) come from the fixed dual frames. Hence they serve every translate in this entry. The product rule and the chain rule preserve integral jets. Factoring the common quadratic exponential out of the finite integral sum proves \eqref{eq:p8c5-ramified}.
\leanref{Formalism.\allowbreak{}L2.\allowbreak{}Th1.\allowbreak{}G.\allowbreak{}ramified\_\allowbreak{}shape\_\allowbreak{}holds}{Formalism.\allowbreak{}LocalToGlobal.\allowbreak{}TorsionIntegrality.\allowbreak{}RamifiedCurveKernels.\allowbreak{}PairedSquareJets.\allowbreak{}InverseEndoIntegralHolds}{117--159}
\leanref{GlobalTheory.\allowbreak{}ThLoc.\allowbreak{}jets\_\allowbreak{}phW\_\allowbreak{}three}{Development.\allowbreak{}AbelianVarietyTheory.\allowbreak{}Proofs.\allowbreak{}SatakeParameters.\allowbreak{}IntSExpMulTranslate}{75--134}
\leanref{GlobalTheory.\allowbreak{}Proof.\allowbreak{}shape\_\allowbreak{}normalisedThetaCombinations\_\allowbreak{}dualBases\_\allowbreak{}at\_\allowbreak{}d.\allowbreak{}shapeComb}{Development.\allowbreak{}AbelianVarietyTheory.\allowbreak{}Proofs.\allowbreak{}SatakeParameters.\allowbreak{}ShapeNormalisedThetaCombinationsDualBasesAtD}{197--232}

At \(p=2\) the host has \(d\equiv7\bmod8\). Use the ordinary integral model with \(a_1\) a unit. In this case use the even theta at the selected half-period. The two CM scalars of \(\omega=(1+\sqrt{-d})/2\) have odd trace and even norm. Order them as \(a,b\in\Z_2\) with \(a\) a unit and \(b\) even. Neither CM multiple of the selected half-period is zero. On the formal subgroup there is a unique nonzero point of order two, and multiplication by the unit scalar fixes it. The half-period therefore reduces away from the identity: otherwise its other CM multiple would be zero since the norm is even. The same argument applied to the inverse image of the half-period under the unit multiplier puts the abscissas of that fibre in the integral ring. Pair the points in this fibre by opposite signs. The isogeny product for the shifted even section gives the phase norm law \(Q([a]t)\sim abQ(t)\); its normalized factors have unit constant term and integral coefficients since the fibre abscissas are integral. For the polarized phase, \(\Z_2\)-linearity follows from integer multiplication by continuity and binomial powering. Its norm law then makes \(a(a-b)\) times the polarized phase equivalent to zero. This scalar is a \(2\)-adic unit, so the polarized phase vanishes in the quotient. Now \(Q([a]t)\sim aQ(t)\) by additivity and continuity. Subtract the norm law and divide by the unit \(a(1-b)\). The normalized even section is an integral one-unit in the formal parameter. Its odd-order translates have integral coefficients by the torsion translation product. Substituting \(l^{-1}\) gives integral jets since \(l'\) and \(1/l'\) are integral. Apply the four-copy argument of Lemma~\ref{B8.56a} to the jets: for four separate variable sets the derivative of multi-order \((\nu,\nu,\nu,\nu)\) is the fourth power of the derivative of order \(\nu\). This argument also transfers the odd-prime jet calculations to the matrix section. Local equality of lattice pairs allows division of each direction by a prime-to-\(p\) comparison integer. The finite sum and the integer minor operators therefore preserve exponent zero.
\leanref{Formalism.\allowbreak{}L2.\allowbreak{}Th1.\allowbreak{}B.\allowbreak{}An2.\allowbreak{}Th0TwoAnalytic\_\allowbreak{}of}{Formalism.\allowbreak{}LocalToGlobal.\allowbreak{}TorsionIntegrality.\allowbreak{}FrobeniusFibreCurves.\allowbreak{}LogarithmicSubstitution.\allowbreak{}Th0TwoAnalyticOf}{25--108}
\leanref{Formalism.\allowbreak{}L2.\allowbreak{}Th1.\allowbreak{}Top.\allowbreak{}integral\_\allowbreak{}operator\_\allowbreak{}two\_\allowbreak{}orthogonal}{Formalism.\allowbreak{}LocalToGlobal.\allowbreak{}TorsionIntegrality.\allowbreak{}UniformGaussModels.\allowbreak{}OrthogonalGaussModels.\allowbreak{}IntegralOperatorTwoOrthogonal}{54--55}
\leanref{GlobalTheory.\allowbreak{}ThLoc.\allowbreak{}jetsIntegral\_\allowbreak{}away\_\allowbreak{}disc}{Development.\allowbreak{}AbelianVarietyTheory.\allowbreak{}Proofs.\allowbreak{}SatakeParameters.\allowbreak{}IntegralJetsNormalisedThetaTranslatesAwayDisc}{57--76}

Finally impose the two family conclusions together. For the local systems take the maximum of the shape and rank thresholds. For the selected entries impose the uniform shape threshold in the choice of eligible points, before evaluating the selection. Thus the entry conclusion holds at every point beyond the selection threshold. Restricting the resulting family assertion to \(\mathcal D_3\) preserves both conclusions with the quantifiers of Definition~\ref{B8.39}.
\leanref{GlobalTheory.\allowbreak{}Svd.\allowbreak{}servH\_\allowbreak{}row\_\allowbreak{}of\_\allowbreak{}artin}{Development.\allowbreak{}AbelianVarietyTheory.\allowbreak{}Proofs.\allowbreak{}SatakeParameters.\allowbreak{}BundleBothPartsOfStandardKernelOfGlobalArtinReci}{15--24}
\leanref{GlobalTheory.\allowbreak{}row\_\allowbreak{}hServ\_\allowbreak{}of\_\allowbreak{}link\_\allowbreak{}alt}{Development.\allowbreak{}AutomorphicArithmetic.\allowbreak{}ThetaCorrespondence.\allowbreak{}ThetaPairLifts.\allowbreak{}ThirdPlaceSeries.\allowbreak{}RowHServOfLink}{17--23}
\leanref{GlobalTheory.\allowbreak{}Fc.\allowbreak{}R1.\allowbreak{}row\_\allowbreak{}hServ\_\allowbreak{}of\_\allowbreak{}kerServ\_\allowbreak{}alt}{Development.\allowbreak{}AutomorphicArithmetic.\allowbreak{}ThetaCorrespondence.\allowbreak{}ThetaPairLifts.\allowbreak{}ThirdPlaceSeries.\allowbreak{}RowHServOfKerServ}{13--22}
\end{proof}

\subsection*{Proof of \ref{B8.56e}}\label{prf:B8.56e}\phantomsection
\begin{proof}
In the short model \(y^2=x^3-g'_2x/4-g'_3/4\) with \(t=-x/y\), \(D(t)=t^2x(t)\) is the even integral solution of \(D^3-D^2-(g'_2/4)t^4D-(g'_3/4)t^6=0\) with constant term one. The derivative in \(D\) there is one. The curve law \(\oplus\) is integral with inverse \(-t\). For
\(R(t,u)=P_E(t\oplus u)+P_E(t\oplus(-u))-2P_E(t)-2P_E(u)\),
the sigma addition formula after homothety gives
\[
 \exp(-R(t,u))=\frac{t^2D(u)-u^2D(t)}{(t\oplus u)(t\oplus(-u))}.
\]
The quadratic exponent cancels from the sigma quotient \cite[Eq.~23.10.3]{reinhardt-walker-2010}. The numerator is \(t^2-u^2\) times an integral one-unit by evenness. The two denominator factors divided by \(t+u,t-u\) are integral one-units by the inverse and linear terms of the group law. Thus \(R\sim0\).
\leanref{Formalism.\allowbreak{}L2.\allowbreak{}Th1.\allowbreak{}B.\allowbreak{}ParO.\allowbreak{}par\_\allowbreak{}eq}{Formalism.\allowbreak{}LocalToGlobal.\allowbreak{}TorsionIntegrality.\allowbreak{}FrobeniusFibreCurves.\allowbreak{}LogarithmicSubstitution.\allowbreak{}InVarAddInVarNeZero}{200--302}

Use formal curve addition in the following identity. With \(b_0=(y+z)/2,\ c_0=(y-z)/2,\ a_0=x+b_0\), expansion gives
\[
 B_E(x+y,z)-B_E(x,z)-B_E(y,z)
 =R(a_0,b_0)-R(a_0,c_0)+R(b_0,c_0)-R(b_0,b_0).
\]
Halving is integral and each substituted point has zero constant term, so the right side is equivalent to zero. Symmetry gives additivity in the other slot. Setting \(u=t\) in \(R\) gives \(P_E(m_2(t))\sim4P_E(t)\).
\leanref{Formalism.\allowbreak{}L2.\allowbreak{}Th1.\allowbreak{}B.\allowbreak{}cube\_\allowbreak{}identity}{Formalism.\allowbreak{}LocalToGlobal.\allowbreak{}TorsionIntegrality.\allowbreak{}FrobeniusFibreCurves.\allowbreak{}LogarithmicSubstitution.\allowbreak{}ConstantCoeffEqZeroOfParallelogram}{115--170}

For \(0\ne\alpha\in\mathcal O_K\) put \(n=\alpha\bar\alpha\) and \(T_\alpha=(\alpha^{-1}\Lambda/\Lambda)\setminus\{0\}\). The isogeny identity is
\[
 \Sigma_\Lambda(\alpha z)^2
 =\alpha^2\Sigma_\Lambda(z)^{2n}
 \prod_{q\in T_\alpha}(\wp_\Lambda(z)-\wp_\Lambda(q)).
\]
The Legendre relation gives \(\eta_\Lambda(\omega)=s_2\omega+\pi\bar\omega/A_\Lambda\), \(A_\Lambda=\Im(\bar\omega_1\omega_2)\). The squared quasi-period multipliers agree since \(\alpha\bar\alpha=n\). Each pair \(q,-q\) gives a double zero, and a single point of order two gives a double zero by ramification of \(\wp_\Lambda\). At zero both sides have leading term \(\alpha^2z^2\), since \(|T_\alpha|=n-1\). Their quotient is elliptic and holomorphic, hence one.
\leanref{Formalism.\allowbreak{}L2.\allowbreak{}Th1.\allowbreak{}F.\allowbreak{}isogenySigma\_\allowbreak{}of}{Formalism.\allowbreak{}LocalToGlobal.\allowbreak{}TorsionIntegrality.\allowbreak{}EisensteinTranslationProducts.\allowbreak{}AnalyticIdentityTheorems.\allowbreak{}IsogenySigma}{394--408}

Comparing \(z^4\) gives \((\alpha^2-n)s_2=\sum_{q\in T_\alpha}\wp_\Lambda(q)\). Take \(\alpha=\sqrt{-d}\); the multiplier is \(-2d\) and the kernel is killed by \(d\). The discriminant-one invariants are algebraic by the algebraicity of the CM \(j\)-invariant. The scaled abscissas \(x_q=e^2\lambda(\wp_\Lambda(q))\) belong to \(W\), by finite étaleness of the prime-to-\(p\) torsion and nonidentity reduction. Thus \(e^2\lambda(s_2)\in K_0\). The normalized Laurent solution for \(\wp\) belongs to \(K_0((z))\). Integrate \(-(\wp-z^{-2})\) twice with zero constants to recover \(\log(\sigma(z)/z)\). This determines the transferred \(\Sigma\) and \(P_E\) over \(K_0\).
\leanref{Formalism.\allowbreak{}L2.\allowbreak{}KerInert.\allowbreak{}Sigma.\allowbreak{}s2\_\allowbreak{}in\_\allowbreak{}unramified\_\allowbreak{}closure}{Formalism.\allowbreak{}LocalToGlobal.\allowbreak{}InertKernelTheory.\allowbreak{}InertCoprimeLemmas.\allowbreak{}PadicNormTwoHundredSixteenEqOneOfFiveLe}{181--226}

For \(p\nmid n\), matching principal parts gives
\(\alpha^2\wp(\alpha z)=\wp(z)+\sum_{q\in T_\alpha}(\wp(z+q)-\wp(q))\) in model coordinates. The translated abscissas are regular over \(W\) on the formal disc, since the torsion reduces away from the origin. Using \(dx/(2y)=dl\), we obtain \(X_\alpha=x(m_\alpha(t))\in t^{-2}W[[t]]\) and \(Y_\alpha=(2\alpha l'(t))^{-1}dX_\alpha/dt\in t^{-3}W[[t]]\). The leading terms \(\alpha^{-2}t^{-2},-\alpha^{-3}t^{-3}\) show that \(m_\alpha=-X_\alpha/Y_\alpha\) is integral and \(m_\alpha/(\alpha t)\) is a one-unit. The sigma identity then gives
\[
 \left(\frac{m_\alpha(t)}{\alpha t}\right)^2
 e^{-2(P_E(m_\alpha(t))-nP_E(t))}
 =\prod_{q\in T_\alpha}(D(t)-x_qt^2).
\]
The right side is an integral one-unit, so its normalized square root proves the norm law for \(\alpha\).
\leanref{Formalism.\allowbreak{}L2.\allowbreak{}Th1.\allowbreak{}C.\allowbreak{}phNormUnit\_\allowbreak{}of}{Formalism.\allowbreak{}LocalToGlobal.\allowbreak{}TorsionIntegrality.\allowbreak{}SecondCoefficientNorms.\allowbreak{}BinomialSeriesCoefficients.\allowbreak{}USideB}{226--295}
\leanref{Formalism.\allowbreak{}L2.\allowbreak{}KerInert.\allowbreak{}PN.\allowbreak{}phNormOKInert\_\allowbreak{}of}{Formalism.\allowbreak{}LocalToGlobal.\allowbreak{}InertKernelTheory.\allowbreak{}InertCoprimeLemmas.\allowbreak{}PhNorm}{371--541}

An integral parameter change \(q(t)\) of unit tangent \(c\) has \(l_{\rm new}(q(t))=cl_{\rm old}(t)\) and \(\Sigma_{\rm new}(cz)=c\Sigma_{\rm old}(z)\). The phase difference is \(\log(q(t)/(ct))\), so the laws transfer to the specified good model. The changes to short form use the units \(2,3\).

The completion of \(\mathcal O_K\) is \(W^{\sigma^2}\), with conjugation \(\sigma\). Approximate each unit \(u\) by elements of \(\mathcal O_K\) of norm prime to \(p\). The coefficients of \(m_a\) are polynomial in \(a\), and those of the exponentiated norm defect are continuous in \(a,\bar a\). Thus integrality and the norm law hold for \(u\). For \(k\in\Z_p\), biadditivity and continuity from integer multiplication give \(B_E(t,m_k(s))\sim kB_E(t,s)\), with binomial powering on integral one-units. Polarize the unit norm law at \(t,m_{u^{-1}}(s)\) to obtain \(B_E(m_u(t),s)\sim B_E(t,m_{\bar u}(s))\). Each \(a\in W^{\sigma^2}\) is the sum of two units, by choosing a nonzero residue different from that of \(a\). Formal addition gives the integral \(m_a\), and biadditivity gives the adjoint law. Doubling gives \(B_E(t,t)\sim2P_E(t)\), hence
\[
 2P_E(m_a(t))\sim B_E(m_a(t),m_a(t))
 \sim B_E(t,m_{a\bar a}(t))\sim2a\bar aP_E(t).
\]
Take the normalized square root.

The equality \(l(v(T))=\epsilon\ell(T)\) gives the signed identity
\[
 P_E(v(T))-P(T)=\log(e_0v(T)/T).
\]
The exponential is an integral one-unit. Since \(v\) intertwines addition and scalars, \(b^\ell=B_E(v(X),v(Y))\sim B_P\) and the laws transfer. The exponentials of both additivity defects and of the adjoint defect are integral after logarithmic substitution. Their coefficients lie in \(K_0\cap\mathcal O_{\C_p}=W\). Since the series vanishes on both axes, it gives the pairing of \ref{B3.05}.
\leanref{Formalism.\allowbreak{}L2.\allowbreak{}KerInert.\allowbreak{}BeCob.\allowbreak{}phase\_\allowbreak{}compatibility}{Formalism.\allowbreak{}LocalToGlobal.\allowbreak{}InertKernelTheory.\allowbreak{}InertCoprimeLemmas.\allowbreak{}ExistsWittSeriesInverseFormalGroupIso}{327--432}
\leanref{Formalism.\allowbreak{}L2.\allowbreak{}KerInert.\allowbreak{}be\_\allowbreak{}props\_\allowbreak{}at}{Formalism.\allowbreak{}LocalToGlobal.\allowbreak{}InertKernelTheory.\allowbreak{}InertCoprimeLemmas.\allowbreak{}BeProp}{96--185}
\end{proof}

\subsection*{Proof of \ref{B8.56d}}\label{prf:B8.56d}\phantomsection
\begin{proof}
Write \(d_0=2h\). Let \(b\) satisfy \ref{B3.05} and put \(b_0=b^\ell\). Take the \(\sigma^{d_0}\)-fixed trace-dual linear forms \(L_i^X,L_i^Y\) of \ref{B3.08}. The phase \(\beta\) is the average of \(\sum_i b(L_i^X,L_i^Y)\) under \(Z_s\mapsto a^{(-1)^sp^s}Z_s\) for \(a^{p^{d_0}-1}=1\). Let \(\mathcal L_s(Z)=\sum_{k\ge0}Z_{s-2k}^{p^{2k}}/p^k\), with cyclic indices. The expansion, with \(a,b'\) indexing matrix entries, is
\[
 \Theta(\xi\mathcal L(Z))=U(Z)
 \exp\left(-\sum_{a,b'}\beta(\mathcal L(Z_{\bullet,a,b'}))\right),
 \qquad U\in\mathcal O_{\C_p}[[Z]] .
 \tag{8.56.4}\label{eq:p8c5-orbit-series}
\]

Take \(z_s\) to be the inverse trace-dual transform of the logarithms of paired formal variables \(T_i,T_i'\). Put \(q_0=p^2\), \((\Psi z)_s=\sigma^2(z_{s-2})(T^{q_0})\), where the substitution raises all variables to \(q_0\). The inverse matrix cycles under \(\sigma^2\), so \((1-p^{-1}\Psi)z\) is integral. Solve \((1-p^{-1}\Psi)\mathcal L(t)=(1-p^{-1}\Psi)z\) by degree. Its linear-in-\(t\) term is \(t_s\). Each higher term is
\[
 p^{-k}\bigl(t_{s-2k}^{q_0^k}-(\sigma^2(t_{s-2k})(T^{q_0}))^{q_0^{k-1}}\bigr)
 \qquad(k\ge1).
\]
It uses lower coefficients. The two bases before taking the \(q_0^{k-1}\)-st power agree modulo \(p\), so their powers agree modulo \(p^k\). This gives integral \(t_s\) of zero constant term. Injectivity of \(1-p^{-1}\Psi\) by lowest degree yields \(\mathcal L(t)=z\).

Let \(S=\sum b_0(T_i,T_i')\) over the basis indices and entries, and \(B\) be the averaged phase at \(z\). For a torus parameter \(a\), the even matrix in trace coordinates has entries \(\sum_{j=0}^{h-1}\sigma^{2j}(a\omega_i^*\omega_k)\). The odd matrix is its inverse adjoint for \(\sum_i\sigma(A_i)B_i\) by trace duality. By \ref{B3.05} the transformed pairing thus reduces to the identity Gram matrix modulo integral exponentials. Multiplying over the torus gives integral \(\exp((p^{d_0}-1)(S-B))\). Its normalized root is integral since \(p^{d_0}-1\) is a unit. Thus \(U_0=U(t)\exp(S-B)\) is integral. Substituting \(\ell^{-1}(L_i^X(Z)),\ell^{-1}(L_i^Y(Z))\) into \(U_0e^{-S}\) recovers \(\Theta(\xi Z)\).
\leanref{GlobalTheory.\allowbreak{}KerStd.\allowbreak{}integral\_\allowbreak{}cofactor\_\allowbreak{}converse}{Development.\allowbreak{}AbelianVarietyTheory.\allowbreak{}Proofs.\allowbreak{}SatakeParameters.\allowbreak{}ExistsIntegralOrbitSeries}{343--372}

By \ref{B3.05b}, \(b\) and \(\ell^{-1}\) have integral derivatives at zero. By the product and chain rules, the displayed operations preserve integrality of the derivatives. This applies to exponentiation of a zero-constant phase. The integer minors give exponent zero by \ref{B3.11}.
\leanref{GlobalTheory.\allowbreak{}KerZero.\allowbreak{}kvalGe\_\allowbreak{}zero}{Development.\allowbreak{}AbelianVarietyTheory.\allowbreak{}Proofs.\allowbreak{}SatakeParameters.\allowbreak{}StandardOrbitSeriesValuationBoundZero}{363--380}

For \(\pi^{p^2-1}=p\), scale by \(\pi\) and use \ref{B3.08}. The cofactor stays integral since \(\ell^{-1}(\pi T)/\pi\) is integral. The exponential in the odd-power polynomial \(H\) has the \(p^{-k}H^k\) expansion of \ref{B3.07} since \(v_p(\pi^{(p+1)k}/k!)=s_p(k)/(p-1)\ge0\), with \(s_p(k)\) the base-\(p\) digit sum. Apply \ref{B3.26} with ranks \(N_j,N_{j+1}\) and orbit length \(d_0\). This gives the exponent \(g_A\) of \ref{B3.13}. Homogeneity removes the norm-one scale \(\xi\).
\leanref{GlobalTheory.\allowbreak{}KerGaussA.\allowbreak{}kvalGe\_\allowbreak{}gainA}{Development.\allowbreak{}AbelianVarietyTheory.\allowbreak{}Proofs.\allowbreak{}SatakeParameters.\allowbreak{}KvalGeOrbitGainWithoutOperator}{106--160}

Let \(\nu_s\) be the layer diagram and \(T\) the set of first-sheet top positions in the orbit. Put \(\alpha_s=\nu_{s,1}-\nu_{s,2}\) on \(T\), zero elsewhere. After removing height-one columns at the marked positions, the local caps are \(\nu_{s,2}\) on \(T\) and \(\nu_{s,1}\) off \(T\). By \eqref{eq:p8c5-entry-shape}, \(\bar\alpha\) of \ref{B3.47} is at most \(\alpha_s\) on \(T\), and the local caps are bounded by the global caps. A positive odd distance \(e<d_0\) between distinct top positions is their first positive cyclic distance and is excluded in \ref{B3.48}. Thus each retained global positive part is bounded by a retained local term of \ref{B3.13}; convergence follows from the geometric decay \(p^{-e}\). Put \(z=\log_p(2a)\ge1\) and \(M=N_m\). The bounds \(\alpha_s\le2a,\ |\nu_s|-\alpha_s\le a^2\) give digit lengths at most \(2z,3z\) by \ref{B3.03}. In \ref{B3.13} their total error is at most \(5d_0M^2z\), plus \(4|T|z^2\) for the squared marked lengths. Since the error in \ref{B3.47} dominates their sum, \(g_0(t,j+1)\le g_A\). Combine this with exponent zero.
\leanref{GlobalTheory.\allowbreak{}KerStd.\allowbreak{}WordFree.\allowbreak{}wordFreeNumber\_\allowbreak{}le\_\allowbreak{}orbitBound}{Development.\allowbreak{}AbelianVarietyTheory.\allowbreak{}Proofs.\allowbreak{}SatakeParameters.\allowbreak{}WordFreeNumberLeOrbitBound}{168--248}
\end{proof}

\subsection*{Proof of \ref{B8.56a}}\label{prf:B8.56a}\phantomsection
\begin{proof}
The good-model assertion is Lemma~\ref{B8.56f}. Extend \(\lambda\) to an isomorphism \(\iota:\C\to\C_p\) by extending transcendence bases. At an inert prime the ideal \((p)\) is prime in \(\mathcal O_K\). Write \(\mathfrak c=(p)^h\mathfrak c'\) with \(p\nmid\mathrm N\mathfrak c'\). The replacement \((\delta,\mathfrak c)\mapsto(p^h\delta,\mathfrak c')\) preserves the actual period lattice, its discriminant and its torsion classes. It permits the use of Lemma~\ref{B8.56e}.

First use a short good model, and normalize the transferred odd sigma section to have tangent one in its differential coordinate. For a torsion point \(t\) write \(G_t(T)\) for its translated sigma section at \(l_W(T)\), divided by the correction \(A(T)\) and normalized by the same differential scale. Thus at a period \(G_0(T)\) is \(T\) times a scalar of norm one. Suppose \(t\) has exact order \(n>1\), \(p\nmid n\), and put \(c_t=G_t(0)\ne0\). For \(1\le j<n\) the sigma addition identity gives
\[
 G_{(j+1)t}G_{(j-1)t}
 = c_t^2\bigl(x(t)-x(T+jt)\bigr)G_{jt}^2 ,
\]
using compatible translation representatives. Here \(x(T+jt)\) is the abscissa on the translated formal disc. The quadratic correction cancels from the sigma quotient \cite[Eq.~23.10.3]{reinhardt-walker-2010}. Multiply these identities with exponents \(n-j\). The intermediate terms cancel and \(nt\) is a period, giving, for a scalar \(|u|_p=1\),
\[
 G_t^n\prod_{j=1}^{n-1}
 \bigl(c_t^2(x(t)-x(T+jt))\bigr)^{n-j}=uT^n .
 \tag{8.56.6}\label{eq:p8c5-sigma-product}
\]
Every abscissa difference here is \(T^a\) times an integral series with unit constant term. If \(jt\ne\pm t\), the abscissas remain distinct after reduction because reduction is injective on the subgroup of order \(n\). If \(jt=\pm t\) and \(2t\ne0\), the first derivative is a unit since the reduced ordinate is nonzero. If \(2t=0\), the reduced abscissa is a simple root of the cubic and the abscissa difference has a double zero with unit leading coefficient. The leading coefficients in \eqref{eq:p8c5-sigma-product} give \(|c_t|_p^{\,n+2\sum_{j=1}^{n-1}(n-j)}=|c_t|_p^{n^2}=1\). After cancellation of \(T^n\), \((G_t/c_t)^n\) is integral. Its root with constant term one is integral by coefficient recursion, since \(n\) is a unit. This proves the assertion for the odd sigma section. For the normalized even section at a nonzero half-period \(h\), the denominator is the value at zero of the \(h\)-translate, a unit by the case \(n=2\). A further translate by \(s\), with \(Ns=0\) and \(p\nmid N\), has numerator a root of unity times the translate at \(s+h\); \(2N(s+h)=0\). This includes \(s+h=0\).
\leanref{Formalism.\allowbreak{}L2.\allowbreak{}Th1.\allowbreak{}E.\allowbreak{}joint\_\allowbreak{}integrality}{Formalism.\allowbreak{}LocalToGlobal.\allowbreak{}TorsionIntegrality.\allowbreak{}TorsionTranslationPowers.\allowbreak{}NonsingularJetNorms.\allowbreak{}Joints}{128--129}

An integral unitary change \((a_{ij})\) of orthonormal basis acts by formal substitution \(T_i\mapsto\bigoplus_j[a_{ij}]T_j\). The ratio of the products of corrections is an integral one-unit. Indeed the diagonal terms in the phase calculation of Lemma~\ref{B8.56e} have multipliers \(\sum_i\bar a_{ij}a_{ij}=1\), and the mixed pairing terms have second-slot multipliers \(\sum_i\bar a_{ij}a_{ik}=0\). A change of good parameter with tangent \(u_0\) changes the correction by the quotient of that parameter by \(u_0T\), an integral one-unit. Thus the calculation applies in the prescribed model. For \(p\ge5\) passage to a short equation uses only division by \(2,3\).

For the matrix section, start with one matrix entry in the trace space over \(K\). The four-copy normalization in Lemma~\ref{B4.22c} identifies its fourth tensor power with the normalized average of a product of elliptic sections, pulled back by an isometry. Choose that isometry to be integral on the two lattices at \(p\). Here the local trace lattice is unimodular. Over the unramified quadratic integer ring a unimodular Hermitian lattice has an orthonormal basis: reduction supplies a vector of nonzero norm, the unit norm map is onto by the residue norm and Hensel lifting, and its orthogonal complement is unimodular. Repeat in the complement. A rational unitary change can realize this local lattice identification. Indeed the stabilizer of a lattice is open, and on the dense open Cayley chart approximation reduces to approximation of the entries of a skew-Hermitian matrix. The same argument works at the finitely many primes above \(p\). Thus the relative index of the standard lattice and the moved lattice is prime to \(p\). The normalization scalar has norm one: in the fourth-power identity of Lemma~\ref{B4.22c} the factors of the initial scalar are supported at split rational primes; the change-of-isometry scalar has eighth power the fourth power of a norm-one determinant in \(K\). That determinant is a unit at an inert prime, since its valuation equals the valuation of its conjugate.
\leanref{Formalism.\allowbreak{}L2.\allowbreak{}KerInert.\allowbreak{}Kap.\allowbreak{}kappa\_\allowbreak{}K1\_\allowbreak{}hypothesis}{Formalism.\allowbreak{}LocalToGlobal.\allowbreak{}InertKernelTheory.\allowbreak{}InertCoprimeLemmas.\allowbreak{}KappaSelfDualFourthPowerLattice}{317--367}

Translate the four-copy identity in all four sets of variables. Its average has unit index denominator, and its shifts have prime-to-\(p\) order since their orders divide that index. After the orthonormal change the right side has integral corrected coefficients by the elliptic calculation. Its coefficient at \((\nu,\nu,\nu,\nu)\) on the left is the fourth power of the one-entry corrected coefficient at \(\nu\). This proves integrality for that entry without a choice of fourth root. For a standard pair of even ranks, decompose the matrix into \(2\)-by-\(2\) blocks. The section on each block and its correction are products over the four entries; translation projects to the four corresponding shifts.

For general endpoints, choose rational unitary changes making them standard at the primes above \(p\) as above. The finite lattice moves of Lemma~\ref{B3.35d} to this pair have indices prime to \(p\). Their determinant-ideal values are units by the inert type-norm calculation in Lemma~\ref{B3.35a-arithmetic}. The resulting sums have integral multipliers and prime-to-\(p\) shifts. Pullback carries the basis and torsion shift by the same isometry. If two endpoints agree after localization at \(p\), an integer prime to \(p\) carries either lattice into the other in both directions. The product of the two endpoint integers does this for their Hom lattices, so a shift of prime-to-\(p\) order for the reference pair still has prime-to-\(p\) order for the translated pair. Apply the preceding calculation term by term. Multiplication by the common correction and extraction of a coefficient commute with the finite sum, proving \eqref{eq:p8c5-orthogonal-taylor}.
\leanref{Formalism.\allowbreak{}L2.\allowbreak{}KerInert.\allowbreak{}weak\_\allowbreak{}approximation\_\allowbreak{}places}{Formalism.\allowbreak{}LocalToGlobal.\allowbreak{}InertKernelTheory.\allowbreak{}InertCoprimeLemmas.\allowbreak{}IntWAPlaces}{465--573}
\leanref{Formalism.\allowbreak{}L2.\allowbreak{}KerInert.\allowbreak{}intInert\_\allowbreak{}of}{Formalism.\allowbreak{}LocalToGlobal.\allowbreak{}InertKernelTheory.\allowbreak{}InertCoprimeLemmas.\allowbreak{}IntAsm}{79--92}
\end{proof}

\subsection*{Proof of \ref{B8.56b}}\label{prf:B8.56b}\phantomsection
\begin{proof}
Work in the additive group of phases modulo zero-constant series with integral exponential. Substitution of integral zero-constant series preserves this equivalence. Write \(Q(t)\) for the class of \(P(t)\), \(t\oplus u=\ell^{-1}(\ell(t)+\ell(u))\), and \(B(t,u)=Q(t\oplus u)-Q(t)-Q(u)\). For \(a,b'\in W^{\sigma^2}\) the hypotheses give
\[
 Q([a]t)=a\sigma(a)Q(t),\qquad
 B([a]t,[b']u)=B(t,[\sigma(a)b']u).
\]
The first scalar belongs to \(\Z_p\); it acts on integral one-units by binomial powering. In the second identity the scalar inside the bracket acts by formal multiplication on the argument.

Order the columns by putting all \(f_e\) before all \(f'_e\). Denote the resulting columns by \(a_j\) and the corresponding formal parameters by \(t_j\). Put \(z_i=\bigoplus_j[a_{ij}]t_j\). The \(z_i\) are integral, have zero constant term, and satisfy \(\ell(z_i)=\sum_j a_{ij}\ell(t_j)\). Expanding one formal sum at a time and using biadditivity gives
\[
 \sum_iQ(z_i)=
 \sum_j\langle a_j,a_j\rangle Q(t_j)
   +\sum_{j<k}B\bigl(t_j,[\langle a_j,a_k\rangle]t_k\bigr).
\]
All diagonal and same-family terms vanish. In the cross terms the multiplier is \(\delta_{eh}\), so the right side is represented by \(\sum_e b(\ell(t_e),\ell(t'_e))=S^\ell\). Therefore
\[
 \Theta(L(\ell(t)))\exp(S^\ell)
       =H(z)\exp\left(S^\ell-\sum_iP(z_i)\right)
       \in\mathcal O_{\C_p}[[t]] .
\]
The first factor is integral by substitution, the second by the computed phase equivalence. Finally the adjoint law with scalar \(-1\) gives \(b(-X,Y)\simeq b(X,-Y)\). Substitute \(-Y\) for \(Y\) and sum over \(e\). This proves the required evenness, with the same pairing throughout.
\end{proof}

\subsection*{Proof of \ref{B8.56c}}\label{prf:B8.56c}\phantomsection
\begin{proof}
Take the curve of Lemma~\ref{B8.56f} for the period lattice of the selected column. The permission condition of Theorem~\ref{B3.45} gives \(p\ge5\) and inertness in the imaginary quadratic field. The curve ideal at this site has norm prime to \(p\). The local comparison specifies the place of the global coefficient field. Extend its embedding to \(\lambda:\overline\Q\to\C_p\), then to a field isomorphism \(\iota:\C\to\C_p\) by transcendence bases. Thus the two readings of a global coefficient agree. Take the finite tests and step characters of Lemma~\ref{B3.35e}. All of these choices precede the layer index.

Use the level and class of the decorated source column for each word and twist. Choose representatives in the finite adelic class sets and orthogonal coordinates for their Hermitian forms. At the reference tuple use the identification on the initial level. At another tuple use its own level group and the same determinant-class map. A determinant lift gives the same class as its norm-one idèle, since two lifts differ in the kernel of \(x\mapsto x/\bar x\) in the determinant quotient. The character values required at these classes are in the global coefficient field by Lemma~\ref{B3.35e}. At a permitted place outside the twist primes, retain the representative off \(p\) and use the split chart of \eqref{eq:p3c3-chart-lattice} at each position above \(p\). The source and target positions, their embeddings, and their successor order are those of the prescribed filling. The initial class is the original initial class under this construction. The group action on each tuple has finite stabilizer: it is the intersection of a lattice stabilizer at finite places and a compact definite unitary group at infinity. Its orbit classes are the finite adelic classes just used. This gives the charts simultaneously for the successive lifts. Replacing marked class representatives by the representatives in these charts preserves the class kernel, because their lattice and theta readings agree under the rational changes of Lemma~\ref{B3.35d}.
\leanref{GlobalTheory.\allowbreak{}Pdv.\allowbreak{}stdKer\_\allowbreak{}stdRd\_\allowbreak{}mark}{Development.\allowbreak{}AbelianVarietyTheory.\allowbreak{}Proofs.\allowbreak{}SatakeParameters.\allowbreak{}StdKerStdRdMark}{17--46}
\leanref{GlobalTheory.\allowbreak{}Port.\allowbreak{}jIn\_\allowbreak{}nonempty\_\allowbreak{}true}{Development.\allowbreak{}AbelianVarietyTheory.\allowbreak{}Proofs.\allowbreak{}SatakeParameters.\allowbreak{}TrueColumnExistsChainChartsKernelDataForms}{23--102}
\leanref{GlobalTheory.\allowbreak{}TwV.\allowbreak{}joint\_\allowbreak{}at\_\allowbreak{}nonempty}{Development.\allowbreak{}AbelianVarietyTheory.\allowbreak{}Proofs.\allowbreak{}SatakeParameters.\allowbreak{}JointDataAtPAndAdelicDatumNonempty}{16--22}

At a pair of lattice tuples take the normalized Taylor series \eqref{eq:p3c3-chart-series}, with the chart substitution and restriction there. Its scalar has norm one. Indeed the principal step-character computation in Lemma~\ref{B3.35a-arithmetic}, applied to the quotient of the two determinant idèles, gives valuation \(\sum_v\Delta_v(j+1)(s_v-s'_v)\). The explicit power of \(p\) has the opposite valuation, and the determinant-ideal values are units. The test primes and the test-level denominators avoid \(p\), so the finite test gives a combination with integral coefficients and prime-to-\(p\) torsion shifts. Its two kernel lattices are self-dual and of even rank by the chart construction. Lemma~\ref{B8.56a} applies. Transfer the good parameter by \(f(v(T))=\epsilon\ell(T)\); this gives the common differential scale \(e\epsilon\), where \(e\) is the scale in \eqref{eq:p8c5-orthogonal-taylor}.
\leanref{GlobalTheory.\allowbreak{}TwV.\allowbreak{}Xc.\allowbreak{}scal\_\allowbreak{}of\_\allowbreak{}pos\_\allowbreak{}pair}{Development.\allowbreak{}AbelianVarietyTheory.\allowbreak{}Proofs.\allowbreak{}SatakeParameters.\allowbreak{}ThetaSeriesScalarBoundedOfPosPair}{22--87}

Choose the pairing of Lemma~\ref{B8.56e} on this one curve. In any framed matrix block, take matrix units on the first factor of the split algebra and dual units on its conjugate factor. Multiply by a trace basis over the unramified quadratic ring on the first factor and its trace dual on the second. The resulting coordinates in an orthonormal basis are integral. The Hermitian form pairs two vectors on the same factor to zero. On opposite factors its value is the product of the delta for the place, the two matrix-index deltas and the trace-basis delta. These are exactly the Gram identities of Lemma~\ref{B8.56b}. That lemma gives the framed generator in Definition~\ref{B3.05} with \(\Phi=\sum_e b(X_e,Y_e)\). Setting the variables at the other places to zero gives the one-place condition, with the same \(b\). Formula \eqref{eq:p3c3-chart-action} identifies the fixed-coordinate series with its transport from representatives, including at a stabilizer. The odd-prime exponential criterion is Lemma~\ref{B3.06}.

For the one-place exponent zero in Definition~\ref{B3.05a}, the derivatives of \(b\) and of \(\ell^{-1}\) are integral by Lemma~\ref{B3.05b}. An integral series after logarithmic substitution has the form \(V(\ell^{-1})\) with \(V\) integral. Products and compositions of series with integral derivatives at zero have integral derivatives there, by the product and chain rules. The same holds for exponentiating a zero-constant such series. The trace-basis coefficients in the one-place substitution are integral. Applying any of the integer minor operators of Definition~\ref{B3.11} therefore gives exponent zero for every cofactor in Definition~\ref{B3.12}. For the full product, absorb the phase difference into the integral cofactor, expand the cofactor in monomials and use the preceding argument on the factor from each place. For an operator of order \(D\) only monomials of degree at most \(D\) contribute. This proves the full zero-index entry bound as well.

We construct the common reading maps explicitly. Given a function on the complex classes of level \(j+1\), take its polynomial value in the transported target frame. For functions in the image of frame evaluation this polynomial is uniquely determined by evaluation on the frames. Use a linear projection onto that image to extend the inverse reading linearly to all class functions, and apply \(\iota\) to the coefficients. Outside lattice tuples take zero. Call the resulting semilinear map \(R_{j,c',L'}\). On a rational input it gives precisely the local comparison of the same input. Indeed both are obtained by writing the adelic point as \(\gamma y_C k\), acting on its class polynomial by \(\gamma\), and multiplying by the level character of \(k\); the coefficient embeddings agree. This is the reading construction of Lemma~\ref{B4.07c} at the target level. Define, at every layer,
\[
 S_j(L,L')=
 \left(\prod_v p^{\Delta_v(j+1)(s_v-s'_v)}\right)\Theta^F_{j,L,L'} .
\]
The length factor is nonzero. Dividing by it recovers the chart series. These choices give one family of readings and one family of series for the entire chain.
\leanref{Formalism.\allowbreak{}L2.\allowbreak{}KerInert.\allowbreak{}SecondKernelProjection.\allowbreak{}rdJ}{Formalism.\allowbreak{}LocalToGlobal.\allowbreak{}InertKernelTheory.\allowbreak{}ProjRSurj}{120--124}
\leanref{Formalism.\allowbreak{}L2.\allowbreak{}KerInert.\allowbreak{}SecondKernelProjection.\allowbreak{}series\_\allowbreak{}reading\_\allowbreak{}entries}{Formalism.\allowbreak{}LocalToGlobal.\allowbreak{}InertKernelTheory.\allowbreak{}InertCoprimeLemmas.\allowbreak{}SecondKernelProjection}{143--167}
\leanref{Formalism.\allowbreak{}L2.\allowbreak{}KerInert.\allowbreak{}SecondKernelSeries.\allowbreak{}srOf}{Formalism.\allowbreak{}LocalToGlobal.\allowbreak{}InertKernelTheory.\allowbreak{}TieSerPin}{11--14}

For clarity, consider one rational input \(x\), one source representative \(A\), and one target tuple \(B\). Express the source value of \(x\) by its weight polynomial. The contribution of its class to the complex lift at \(B\) is the contraction of this polynomial with the Taylor series of the finite test, with the source and target determinant characters. Transport this target polynomial to the chosen frame. A linear substitution in a Taylor series transforms its differential symbol by the transpose substitution. Applying this identity to each determinant in the symbol of Lemma~\ref{B3.10a} gives the contraction of the source reading with \(S_j(A,B)\). The scalar left over is constant at the layer. More explicitly, if \(k_s,k_t\) are the source and target determinant factors of the orthogonal frame, \(d_\tau\) the derivation generators, \(m_\tau\) the degrees of the symbol blocks, and \(\delta_j\in\{1,-1\}\) the sign of the row reversals, it is
\[
 c_j=\iota\!\left(
 \frac{k_t}{k_s\,\delta_j\prod_\tau d_\tau^{-m_\tau}}\right)
           (e\epsilon)^{-\sum_\tau m_\tau}.
\]
The frame factors here are the products of the determinants of the basis changes to the powers of the source and target determinant twists. The orthogonal basis used for these factors is the same in all components of a level. The rational unitary change contributes its type norm, a unit at this inert prime; the local integral basis changes have unit determinants. The derivation generators are units at the permitted prime, and the coordinate scale has norm one by discriminant normalization and the unit tangent in Lemma~\ref{B8.56f}. Thus \(|c_j|_p=1\), also for \(c_j^{-1}\). If a component set is empty use \(c_j=1\). In any two decompositions of the adelic point the target polynomials just obtained have the same image under \(\iota\), namely this contraction. Injectivity coefficientwise gives the same polynomial before applying \(\iota\). This proves the reading identity for every decomposition, including the initial lift.
\leanref{GlobalTheory.\allowbreak{}TwV.\allowbreak{}DelA.\allowbreak{}cKs}{Development.\allowbreak{}AbelianVarietyTheory.\allowbreak{}Definitions.\allowbreak{}GaloisGeometry.\allowbreak{}RealAt}{65--69}
\leanref{GlobalTheory.\allowbreak{}TwV.\allowbreak{}DelA.\allowbreak{}tieDel\_\allowbreak{}all\_\allowbreak{}steps}{Development.\allowbreak{}AbelianVarietyTheory.\allowbreak{}Proofs.\allowbreak{}SatakeParameters.\allowbreak{}KernelSiteFrameValuesAllStepsLtM}{23--94}

Sum the identity over source classes with reciprocal stabilizer weights. On rational generators this is the defining finite sum for the arithmetic lift. The rational class maps carry such generators to rational target vectors; hence their comparisons lie in the rational target span. Linearity gives both preservation of that span and \eqref{eq:p3c3-class-expansion} for all inputs. At \(j=m-1\) use the bijection between the rational classes and the local source representatives. Take the inner set in \eqref{eq:p3c3-terminal-expansion} to be a singleton and \(w_x=|\operatorname{Stab}(x)|\). The bijection identifies the stabilizer groups. At positive levels their orders divide \(D_{\Gamma,u}\); at zero they are one. Thus the choice \(d_G(0)=0\), \(d_G(j)=v_p(D_{\Gamma,u})\) for \(j>0\) works simultaneously. The same representatives work for every input and target.
\leanref{Formalism.\allowbreak{}L2.\allowbreak{}KerInert.\allowbreak{}SecondKernelCoupled.\allowbreak{}second\_\allowbreak{}kernel\_\allowbreak{}coupled}{Formalism.\allowbreak{}LocalToGlobal.\allowbreak{}InertKernelTheory.\allowbreak{}InertCoprimeLemmas.\allowbreak{}SecondKernelCoupled}{44--75}

For an invariant function on edges, the sum at a fixed endpoint divided by its stabilizer order contributes \(1/|\operatorname{Stab}(\text{edge})|\) per edge orbit. Counting from the other endpoint gives the identical weight. For the kernel use \eqref{eq:p3c3-chart-action} on the source and the transpose substitution on its input polynomial. The level character cancels, and the determinant and length factors are those already included in the contraction. This gives the required source invariance and the terminal neighbour identity after reindexing by the same class bijection. The indicator assertion is Lemma~\ref{B3.35a-arithmetic}. For the frame assertion, at each place lift to the target lattice a basis adapted to the given lower neighbour modulo \(p\). Multiplication of its first \(i_v\) vectors by \(p\) gives a basis of the neighbour. Lemma~\ref{B3.10a} applies to this single frame for every series, exponent and input vector. Lastly, moves at two distinct places commute by interchanging their two independent choices of neighbours. The edge-lattice calculation \eqref{eq:p3c3-edge-exponent} in Theorem~\ref{B3.34}, applied at level zero, gives integral preservation by each normalized input operator. This proves all conditions of Definition~\ref{B3.05a}. The initial comparison is the fixed one by construction, the class-sum calculation intertwines the adjacent lifts, and Lemma~\ref{B4.13a} transfers the terminal valuation bound with its original exponent.
\leanref{Formalism.\allowbreak{}L2.\allowbreak{}KerInert.\allowbreak{}SecondKernelCoupled.\allowbreak{}hecke\_\allowbreak{}coupled\_\allowbreak{}entry}{Formalism.\allowbreak{}LocalToGlobal.\allowbreak{}InertKernelTheory.\allowbreak{}InertCoprimeLemmas.\allowbreak{}SecondKernelCoupled}{104--136}
\leanref{GlobalTheory.\allowbreak{}TwV.\allowbreak{}admissible\_\allowbreak{}at\_\allowbreak{}each\_\allowbreak{}site}{Development.\allowbreak{}AbelianVarietyTheory.\allowbreak{}Proofs.\allowbreak{}SatakeParameters.\allowbreak{}JointSiteChainComparedCarriesAdmittedSystem}{29--76}
\end{proof}

\section*{Proofs for Chapter~\ref{ch:p8c6}: The chain length}

\subsection*{Proof of \ref{B8.59}}\label{prf:B8.59}\phantomsection
\begin{proof}
By the ceiling bound above, $a\le e(a)<a+2\le2a$ for $a\ge2$.
Here $2\le\alpha\le x$ and $\beta\ge x^4$.
Indeed, $1/2\le1-\varepsilon/8\le2$ and $x\ge16$ give
$4\le x^{1-\varepsilon/8}\le x^2$.
Choose $m=e(\alpha)$ and $k=e(\beta)$.
Since $2\le2+8\varepsilon\le3$, the order bounds give
$x^2\le n\le2x^3$. Hence
\[
 2\le m\le2x\le2x^2\le2n\le4x^3\le x^4\le k.
\]
Also $m+3\le2x+3\le x^4\le k$.

Define $r=k+m-1$ and $t=8fr+4+TW-r$.
Since $f\ge m\ge2$ and $r\ge m$, the difference is at least
$15r+4$, so it has its ordinary integer value.
Since $k$ and $m$ are even, $r$ is odd. Since $8\mid W$,
the integer $t$ is odd. Moreover,
\[
 r+5\le t,\qquad 3m<t,\qquad
 m^2\le m(t+1)\le k+m(t+1).
\]
The bounds and parity conditions make $(k,m,t)$ rank data on $\varphi$. The formula for $t$
gives the identity for $r+t$.  The choices of $m$ and $k$
give their respective bounds.
\end{proof}

\section*{Proofs for Chapter~\ref{ch:galois-polarized}: Galois representations attached to polarized cuspidal representations}

\subsection*{Proof of \ref{B9.07a}}\label{prf:B9.07a}\phantomsection
\begin{proof}
Write $b_v$ for the exponent of $\chi_v$ on the positive reals at each real place $v$ of $F$. Since $a\ge1$, $\pi$ has at least one infinitesimal weight at each embedding, and each weight lies in $(a-1)/2+\Z$, so the multisets $H_\tau$ of Definition~\ref{B9.04} consist of integers. At the pair of embeddings above $v$, polarization pairs infinitesimal weights by $z\mapsto b_v-z$, so $b_v=z+(b_v-z)$ is integral. The finite part of $\chi$ is trivial on a finite-index subgroup of the totally positive units of $F$. The vector $(b_v)_v$ annihilates their logarithms, which span the hyperplane of coordinate sum zero by Dirichlet's unit theorem \cite[Ch.~I, \S7]{neukirch-1999}. Thus all $b_v$ equal an integer $b$. Subtracting the paired weights from $(a-1)/2$ gives $H_{\tau c}=a-1-b-H_\tau$.
\end{proof}

\subsection*{Proof of \ref{B9.07}}\label{prf:B9.07}\phantomsection
\begin{proof}
Choose $r$ by \ref{B9.38}, and take $w_{\rm gt}=a-1-b$ and the integer multisets $H_\tau$ of \ref{B9.04} by Lemma~\ref{B9.07a}.

Fix a prolongation and a map $R$ of \ref{B9.47}. Set $\sigma^*=\pi_w\otimes|\det|_w^{(1-a)/2}$ and $y=R(\sigma^*)$. The dimension law gives $\dim y=a$, the dimension of the parameter away from $\ell$. The spherical assertion above $\ell$ follows from \ref{B9.38}.

Suppose $N_y=0$. Cuspidality gives genericity of $\pi_w$ and its twist $\sigma^*$ \cite{shalika-1974}. Under the commutator test, \ref{B9.55a} gives unlinked rank-one support. Proposition~\ref{B9.38c} gives the toral conclusions, including the jumps above $\ell$.

For the tame clause, assume $E_w\simeq\Q_p$ continuously. Every wild inertia element has finite-order image and all eigenvalues one, so the action of $P_w$ is trivial. By \ref{B9.55a}, $\sigma^*$ has unlinked depth-zero support and first principal congruence invariants. Its support terms of rank at least two are pairwise inertially inequivalent. The unramified normalizing twist preserves all these properties. Use \ref{B9.38d} to take the lift $\Pi$ to $M$, the place $u$ above $w$ with $M_u/E_w$ unramified, and the map $R'$ at $u$. Put $\Pi_u^*=\Pi_u\otimes|\det|_u^{(1-a)/2}$. Proposition~\ref{B9.38c} gives $x'=R'(\Pi_u^*)$ for the upper Galois parameter, with $N_{x'}=0$.

Above $\ell$ the upper comparison gives potential semistability downstairs, and the specified period bridge extends to $M_u$. In both cases the parameter $x'$ restricts the Weil action of the downstairs parameter $x$ and retains its monodromy operator \cite{fontaine-2004}. Thus $N_x=0$. For $t\in I_w=I_{M_u}$ the two comparisons give
\[
 \det(X-\rho_x(t))=\det(X-\rho_{x'}(t))
 =\det(X-\rho_{R'(\Pi_u^*)}(t))=\det(X-\rho_y(t)).
\]
At $t=1$ this gives the rank. Above $\ell$, scalar extension of the filtered de Rham module preserves the jumps at each induced embedding \cite{fontaine-2004}. The pullback of infinitesimal characters in \ref{B9.38d} identifies the upper jumps with $H_\tau$.
\end{proof}

\subsection*{Proof of \ref{B9.08}}\label{prf:B9.08}\phantomsection
\begin{proof}
Take the corresponding plan of \ref{B9.09} and the semisimplified eigenspace of \ref{B9.15}. Let $\kappa_\ell$ be the $\ell$-adic character of the plan's algebraic Hecke character $\kappa$. At a place $w$ where $\kappa$ is unramified, write $\kappa(w)$ for its value on a uniformizer. For $w\nmid\ell$, the value of $\kappa_\ell$ at inverse arithmetic Frobenius is $\kappa(w)$ after transport by $\iota$. For $w\mid\ell$, the unramified Weil parameter of the crystalline character has transported value $\kappa(w)$ at degree $-1$. The character has jump $-m_{\iota\circ\tau}$ at $\tau$, with $m_{\iota\circ\tau}$ the corresponding infinity exponent of $\kappa$ \cite[Ch.~III]{serre-1998}\cite{fontaine-2004}.

Tensor the eigenspace representation by $\kappa_\ell$. Proposition~\ref{B9.17} gives the annihilators away from $\ell$. At a Hodge place, \ref{B9.18a} and \ref{B9.19} give degree $j+m_{\iota\circ\tau}$ for $j\in J(\lambda)$. Tensoring gives jump $j$. At $v'$ the prescribed set supplies the lacunary subset. For the root condition above $\ell$, \ref{B9.19} gives the root multiset bound before twisting. Tensoring multiplies every root at degree $-1$ by $\kappa(w)$, so it gives the asserted divisibility for $A_a(w,B)$.
\end{proof}

\subsection*{Proof of \ref{B9.09a}}\label{prf:B9.09a}\phantomsection
\begin{proof}
Write \(f_U:H^1\to W^\vee\) for the given form. On a rational isometry its left transformation law is the dual weight action, which is the restriction of \(\xi\). Choose its level \(K\subset H\) with \(\nu(K)\subset\widehat{\Z}^{\times}\). At \(w_0\) take the kernel of the type character for this level, and retain the action of the full type subgroup by right translations. At the spherical places use the integral general linear group. Thus \(f_U\) is right invariant under \(K\cap H^1\).

Define a function \(f_0\) on \(H\) by
\[
 f_0(\gamma u k)=\xi(\gamma)f_U(u)
 \quad(\gamma\in\Gamma,\ u\in H^1,\ k\in K),\qquad
 f_0(h)=0\quad(h\notin\Gamma H^1K).
\]
If \(\gamma u k=\gamma' u'k'\), then
\(\nu(\gamma'^{-1}\gamma)=\nu(k'k^{-1})\). The left side is a positive rational number and the right side is a unit at every finite prime. Hence both equal one. Now
\[
             u'=(\gamma'^{-1}\gamma)u(kk'^{-1}),
\]
where the first factor is a rational isometry and the last lies in \(K\cap H^1\). The two transformation laws of \(f_U\) give the same value of \(f_0\). In particular \(f_0|_{H^1}=f_U\).
\leanref{UpperTheory.\allowbreak{}Gal.\allowbreak{}CertificateRingHomMap.\allowbreak{}pKerFormU\_\allowbreak{}of}{Development.\allowbreak{}AutomorphicArithmetic.\allowbreak{}Foundations.\allowbreak{}GaloisRepresentations.\allowbreak{}ProjectiveAutomorphicHecke.\allowbreak{}PKerFormU}{93--132}
\leanref{UpperTheory.\allowbreak{}Gal.\allowbreak{}CertificateRingHomMap.\allowbreak{}pExtZero\_\allowbreak{}holds}{Development.\allowbreak{}AutomorphicArithmetic.\allowbreak{}Foundations.\allowbreak{}GaloisRepresentations.\allowbreak{}ProjectiveAutomorphicHecke.\allowbreak{}PExtZero}{65--96}

The type translations at \(w_0\) have multiplier one and preserve the kernel of the type character. At a split spherical place the Hecke sum has right-coset representatives in the local general linear group embedded in \(H^1\). The integral subgroup permutes these cosets, so the sum preserves right \(K\)-invariance. On \(H^1\) each operator therefore has its original eigenvalue on \(f_0\). The eigenvalue identity extends to \(\Gamma H^1K\) by left equivariance and right invariance. Outside this set both sides vanish, since right multiplication by an element of \(H^1\) preserves each multiplier. This also proves that right translation by a central unitary element acts on \(f_0\) through the finite central character of \(U\).
\leanref{UpperTheory.\allowbreak{}Gal.\allowbreak{}CertificateRingHomMap.\allowbreak{}pFormGU\_\allowbreak{}of}{Development.\allowbreak{}AutomorphicArithmetic.\allowbreak{}Foundations.\allowbreak{}GaloisRepresentations.\allowbreak{}ProjectiveAutomorphicHecke.\allowbreak{}PFormGU}{35--153}

Consider the span of the functions \(h\mapsto f_0(h\zeta(y))\). Put \(O=\zeta^{-1}(K)\), an open subgroup of the finite idèles. Its elements act trivially on this span. For \(t\in E^\times\), the element \(\zeta(t)\) is rational, and \(N_{E/\Q}t>0\) since \(E\) is CM. Its action is scalar. Indeed the dual weight at \(v\) takes the scalar matrix \(z1_a\) to \(z^{-\sum_i\lambda_{v,i}}\). Consequently the scalar is
\[
 (N_{E/\Q}t)^b
 \prod_v\tau_v(\overline{N_{E/I}t})^{-\sum_i\lambda_{v,i}}
       =\prod_{\tau:E\hookrightarrow\C}\tau(t)^{m_\tau}.
\]
The equality follows by expanding each norm over the embeddings extending its indicated embedding of \(I\).

The quotient \(\mathbf A_{E,f}^\times/(E^\times O)\) is finite. The valuation map reduces the assertion to finiteness of the ideal class group of the number field \(E\) and finiteness of the index of \(O\cap\widehat{\mathcal O}_E^\times\) in the compact group \(\widehat{\mathcal O}_E^\times\). The first is \cite[Ch.~I, Thm.~6.3]{neukirch-1999}, and the second follows from openness. The central translates thus span a nonzero finite-dimensional space. They commute with all the multiplier-one operators above. Their actions on this space are commuting automorphisms, so they have a common eigenvector \(\Phi\). Indeed, restricting to an eigenspace of any nonscalar member reduces the dimension and preserves all the other actions. Write \(\psi(y)\) for the eigenvalue of translation by \(\zeta(y)\). The group law makes \(\psi\) a character. It is trivial on \(O\), and on the elements with \(\zeta(y)\in H^1\) it agrees with the central character of \(U\) evaluated at \(\zeta(y)\).
\leanref{UpperTheory.\allowbreak{}Gal.\allowbreak{}CertificateRingHomMap.\allowbreak{}pExtEig2\_\allowbreak{}of}{Development.\allowbreak{}AutomorphicArithmetic.\allowbreak{}Foundations.\allowbreak{}GaloisRepresentations.\allowbreak{}ProjectiveAutomorphicHecke.\allowbreak{}PExtEig2Of}{39--149}
\leanref{UpperTheory.\allowbreak{}Gal.\allowbreak{}CertificateRingHomMap.\allowbreak{}pCentralEig2\_\allowbreak{}of}{Development.\allowbreak{}AutomorphicArithmetic.\allowbreak{}Foundations.\allowbreak{}GaloisRepresentations.\allowbreak{}ProjectiveAutomorphicHecke.\allowbreak{}PCentralEig2Of}{52--90}

For an idèle \(y=(y_f,y_\infty)\) define
\[
             \kappa(y)=\psi(y_f)^{-1}
                      \prod_{\tau:E\hookrightarrow\C}
                                  y_{\infty,\tau}^{m_\tau},
 \qquad y_{\infty,\bar\tau}=\overline{y_{\infty,\tau}}.
\]
The rational scalar calculation gives \(\kappa(t)=1\) for every \(t\in E^\times\). This defines the asserted Hecke character and fixes the inverse in its finite part.
\leanref{UpperTheory.\allowbreak{}Gal.\allowbreak{}CertificateRingHomMap.\allowbreak{}zetaEx\_\allowbreak{}holds}{Development.\allowbreak{}AutomorphicArithmetic.\allowbreak{}Foundations.\allowbreak{}GaloisRepresentations.\allowbreak{}ProjectiveAutomorphicHecke.\allowbreak{}ZetaMain}{56--101}
\leanref{UpperTheory.\allowbreak{}Gal.\allowbreak{}CertificateRingHomMap.\allowbreak{}pKapOf\_\allowbreak{}of}{Development.\allowbreak{}AutomorphicArithmetic.\allowbreak{}Foundations.\allowbreak{}GaloisRepresentations.\allowbreak{}ProjectiveAutomorphicHecke.\allowbreak{}PKapOf}{37--139}
\end{proof}

\subsection*{Proof of \ref{B9.09ab}}\label{prf:B9.09ab}\phantomsection
\begin{proof}
For every rational prime \(q\) with \(-d\in(\Q_q^\times)^2\), the algebra \(I\otimes_\Q\Q_q\simeq\Q_q[T]/(T^2+d)\) splits, so \(q\) splits in \(I=\Q(\omega)\) and does not divide the discriminant of \(\mathcal O_I=\Z[\omega]\); hence the minimal polynomial of \(\omega\) has distinct roots modulo \(q\). By Chebotarev's density theorem for \(E/F\), choose a finite place \(\mathfrak p_1\) of \(F\) inert in \(E/F\) whose residue characteristic \(q_1\) is odd, unramified in \(E\), and outside \(X_0\). Let \(\chi_{E/F}:\mathbf A_F^\times/F^\times\to\{\pm1\}\) be the quadratic character of \(E/F\). Define \(z\in\mathbf A_F^\times\) with \(z_{v_0}=1\), \(z_v=-1\) at the \([F:\Q]-1\) real places \(v\ne v_0\), \(z_{\mathfrak p_1}\) a uniformizer at \(\mathfrak p_1\), and \(z_{\mathfrak p}=1\) elsewhere. Because \([F:\Q]\) is even, \(\prod_v\chi_{E/F,v}(z_v)=(-1)^{[F:\Q]-1}(-1)=1\). By global norm reciprocity \cite[Ch.~VI, \S5]{neukirch-1999}, \(z=\mu\,N_{E/F}(w)\) for some \(\mu\in F^\times\) and \(w\in\mathbf A_E^\times\). Then \(\mu_{v_0}>0\), \(\mu_v<0\) for real \(v\ne v_0\), and \(\mu\) is a local norm from \(E\) at every finite place \(\mathfrak p\ne\mathfrak p_1\) of \(F\). Set \(h_{\rm std}=\mu\,1_a\). Since \(\mu\in F^\times\) is fixed by conjugation, \({}^t\bar g(\mu\,1_a)g=s(\mu\,1_a)\) is equivalent to \({}^t\bar g g=s\,1_a\), so \(\mathrm{GU}^+(h_{\rm std})(\Q)=\Gamma\).

By weak approximation at the real places of \(F\), choose \(\alpha\in F^\times\) with \(\alpha_{v_0}<0\) and \(\alpha_v>0\) for every real \(v\ne v_0\), and set \(h_0=\operatorname{diag}(\mu,\alpha\mu,\dots,\alpha\mu)\) with one entry \(\mu\) and \(a-1\) entries \(\alpha\mu\). Then \(h_0\) has signature \((1,a-1)\) at \(v_0\) and \((0,a)\) at real \(v\ne v_0\). Because \(a-1\) is even, \(\det(h_0)/\det(h_{\rm std})=\alpha^{a-1}=N_{E/F}(\alpha^{(a-1)/2})\) is a local norm at every finite place of \(F\). By the local Hermitian classification in Definition~\ref{B1.15ac}, \((E^a\otimes_F F_{\mathfrak p},h_0)\) and \((E^a\otimes_F F_{\mathfrak p},h_{\rm std})\) are isometric at every finite place \(\mathfrak p\) of \(F\), and at every \(\mathfrak p\ne\mathfrak p_1\) both are isometric to \(1_a\).

At each \(q\in X_0\), we have \(q\ne q_1\) and \(\mathcal O_E\otimes_\Z\Z_q=\mathcal O_{E,q}^+\times\mathcal O_{E,q}^-\) according to the two roots of \(\omega\) modulo \(q\), interchanged by conjugation. Taking \(L_q^+=(\mathcal O_{E,q}^+)^a\) and \(L_q^-\) to be its \(\Z_q\)-dual under \((u,v)\mapsto\operatorname{Tr}_{E/\Q}(x\,h_0(u,v))\) gives an \((\mathcal O_E\otimes\Z_q)\)-lattice \(L_q=L_q^+\oplus L_q^-\) on which the alternating pairing is \(\Z_q\)-valued and unimodular. Modify \(\mathcal O_E^a\) at the primes \(q\in X_0\) to match \(L_q\), and multiply \(h_0\) and \(\mu\) by the square of a positive integer prime to \(\prod_{q\in X_0}q\) clearing denominators of the pairing away from \(X_0\) to obtain \(h_{\rm PEL}\), \(\Lambda\), and \(e=[\Lambda^\vee:\Lambda]\) prime to \(\prod_{q\in X_0}q\). At every \(q\nmid e\), \(e_\Lambda\) is perfect over \(\Z_q\). By Chebotarev's density theorem for \(I/\Q\), infinitely many primes \(p_0\) split in \(I\), hence satisfy \(-d\in(\Q_{p_0}^\times)^2\); choose \(p_0\notin X_1\) with \(p_0\nmid q_1 e\). Away from \(p_0\), choosing compatible local \(\mathcal O_E\)-lattice bases at \(q\ne q_1\) and matching the common local Hermitian space of \(h_{\rm PEL}\) and \(h_{\rm std}\) at \(\mathfrak p_1\) identifies the away-from-\(p_0\) PEL similitude group with \(H^{p_0}\) and places \(h_{\rm std}\) as the standard definite summand.
\end{proof}

\subsection*{Proof of \ref{B9.09ac}}\label{prf:B9.09ac}\phantomsection
\begin{proof}
By Chebotarev's density theorem for \(I/\Q\), choose a rational prime \(r_0\notin\{p_0,\ell\}\) splitting in \(I\) and avoiding the finite set of residue characteristics of \(H\cup C\), bad level places of \(\Phi\), and prime divisors of \(e\). Choose a place \(W\mid r_0\) of \(E\) on the selected side of \(I\); then \(r_0\) splits in \(I\), so \(W\ne\bar W\) and \(E_W\simeq F_{W|_F}\). Let \(I_W\subset\GL_a(\mathcal O_{E_W})\) be the upper-triangular Iwahori subgroup. Because \(r_0\mid N\), the principal congruence subgroup modulo \(N\) at \(W\) is contained in \(1+r_0 M_a(\mathcal O_{E_W})\subset I_W\). Define \(e_{\rm typ}\) in the away-from-\(p_0\) Hecke algebra of level \(K(N)\) as the tensor product of the elliptic type projector at \(w_0\) (transported by the geometric coordinate map), the normalized indicator \(\operatorname{vol}(I_W)^{-1}\mathbf1_{I_W}\) (embedded on \(E_W\times E_{\bar W}\) via \(g\mapsto(g,{}^t\bar g^{-1})\) of multiplier \(1\)), and the normalized indicator of \(K(N)\) away from \(\{w_0,W\}\). Since \(\Phi\) has the prescribed elliptic type at \(w_0\) and is \(\GL_a(\mathcal O_{E_W})\)-fixed at \(W\), \(e_{\rm typ}\) fixes the transported form of \(\Phi\).

Choose a uniformizer \(\varpi_W\) of \(E_W\), set \(t_W=\operatorname{diag}(1,\varpi_W^{-1},\dots,\varpi_W^{1-a})\in\GL_a(E_W)\), and define \(u\) to be \(\operatorname{vol}(I_W)^{-1}\mathbf1_{I_W t_W I_W}\) at \(W\) tensored with the identity of the level away from \(W\). By Bernstein's presentation of the Iwahori--Hecke algebra \(\mathcal H(\GL_a(E_W),I_W)\) \cite[\S3]{borel-1976}, the operator of a dominant translation is invertible in \(\mathcal H(\GL_a(E_W),I_W)\), providing \(u'\) supported at \(W\) such that \(u'ue_{\rm typ}=e_{\rm typ}\), and length additivity on the dominant chamber gives \((I_W t_W I_W)^n=I_W t_W^n I_W\) up to positive scalar for every \(n\ge1\). Deck elements of sites at residue characteristics \(q\ne r_0\) are congruent to \(1\) modulo \(N/q^{\operatorname{ord}_q N}\), hence modulo \(r_0\), so their \(W\)-components lie in \(I_W\); spherical operators at places of residue characteristic \(q\ne r_0\) have unit component at \(W\). Thus \(u\) commutes with \(e_{\rm typ}\) and with all these operators. Let \(n\ge1\) and \(x=i_1t_W^ni_2\) with \(i_1,i_2\in I_W\). Then \(x\) is conjugate to \(t_W^nJ\) with \(J=i_2i_1\in I_W\). Since \(J\) is upper triangular modulo \(\varpi_W\) with unit diagonal, every principal minor \(\det J[S,S]\) is a unit. The coefficient of \(X^{a-k}\) in the characteristic polynomial of \(t_W^nJ\) is \((-1)^k\sum_{|S|=k}\varpi_W^{-n\sum_{i\in S}(i-1)}\det J[S,S]\). Its term for \(S=\{a-k+1,\dots,a\}\) has the unique smallest valuation, so the coefficient has valuation \(-nk(2a-k-1)/2\) for \(1\le k\le a\). Hence the Newton polygon over \(E_W\) has \(a\) distinct slopes \(0,-n,\dots,-(a-1)n\) of multiplicity \(1\), and the characteristic polynomial has \(a\) distinct roots. Multiplying at \(W\) by elements of \(I_W\) preserves \(I_W t_W^n I_W\), so every element in the away-from-\(p_0\) support of \(b(T)u^ne_{\rm typ}\) has separable characteristic polynomial of degree \(a\).
\end{proof}

\subsection*{Proof of \ref{B9.09}}\label{prf:B9.09}\phantomsection
\begin{proof}
Let $H,C$ be the disjoint sets of prescribed Hodge and root places. Apply Lemma~\ref{B9.09ab} with $X_0=\{\ell\}\cup\{\operatorname{char}k(w):w\in H\cup C\}$ to obtain the ambient PEL datum $(E^a,\mathcal O_E,\Lambda,e_\Lambda,\omega)$, the integer $e$, and the standard definite form $h_{\rm std}=\mu\,1_a$. Choose local model charts on the same side of $\Q(x)$ for the places in $H\cup C$ of each residue characteristic, using at $\ell$ the side determined by $w_0$. Choose a weight vector of $U$ in the elliptic type at $w_0$ and spherical at the prescribed places. Apply Lemma~\ref{B9.09a} to obtain the similitude form $\Phi$ and its character $\kappa$. The two partitions at the embeddings over $v$ are $(b_v^a)$ and $\lambda_v+(b_v^a)$, both with nonnegative entries by the choice of $b_v$. Choose $g_0\in H$ with $\Phi(g_0)\ne0$. By Lemma~\ref{B9.09ab}(2), choose the model prime $p_0$ outside $X_0\cup\{q:q\mid e\}$, the residue characteristics of the bad level places of $\Phi$, and the finitely many primes where $g_0$ is not integral and unimodular, obtaining the compact open subgroup $K_{p_0}\subset G_{p_0}$ and the away-from-$p_0$ frame. By Lemma~\ref{B9.09ac}, choose the auxiliary split prime $r_0\notin\{p_0,\ell\}$ and place $W\mid r_0$ outside $X_0$, the bad level primes of $\Phi$, and the prime factors of $e$, and take a common principal level $N\ge3$ prime to $p_0$ and divisible by $r_0$ whose congruence subgroup $K(N)$ fixes $\Phi$ away from $w_0$ and admits all prescribed charts.

The deck correspondences that fix the vector have eigenvalue one. Choose spherical correspondences at charts compatible with $H$ and outside $H$, all of residue characteristic different from $r_0$. In the coordinates of Lemma~\ref{B9.09a}, the operator of index $i$ uses the multiplier-one spherical cosets together with a central translation $z_w^i$. Put $q_w=\#k(w)$ and choose a uniformizer $\varpi_w$, viewed as a finite idèle supported at $w$. On the side of $I=\Q(x)$ conjugate to $w$, the scalar of $z_w$ is $q_w$ at the places over the same rational prime, and on the remaining factors it is one. The scalar of $\zeta(\varpi_w)$ on the conjugate side has the same valuation, by the local norm formula, and it is one off that side. Their quotient therefore has integral unit components and is one at the exceptional level places by the choice of chart, so it fixes $\Phi$. The same argument with a local unit shows that $\kappa$ is unramified at $w$. Thus $z_w^i$ contributes $\kappa(w)^{-i}$, and the eigenvalue upstairs is the original spherical eigenvalue times $\kappa(w)^{-i}$ and the level index.
\leanref{UpperTheory.\allowbreak{}Gal.\allowbreak{}CertificateRingHomMap.\allowbreak{}pKapOf\_\allowbreak{}of}{Development.\allowbreak{}AutomorphicArithmetic.\allowbreak{}Foundations.\allowbreak{}GaloisRepresentations.\allowbreak{}ProjectiveAutomorphicHecke.\allowbreak{}PKapOf}{37--140}

To put this vector in the definite trace, use the geometric coordinates $(g,s)\mapsto({}^tg^{-1},s^{-1})$. Pullback gives left weight $\xi^\sharp(g,s)=\xi({}^tg^{-1},s^{-1})$. Transport the test by the same coordinate map. On the standard definite summand $(E^a,h_{\rm std})$ of Lemma~\ref{B9.09ab}(3), whose rational positive-multiplier similitude group is $\Gamma$, write $\phi(h)$ for the transported form read through the away-from-$p_0$ frame on $h\in H^{p_0}$. If $\delta\in\Gamma$ satisfies $\delta_{p_0}\in K_{p_0}$, then $\delta$ is integral at $p_0$ and $\phi(\delta^{p_0}h)=\xi^\sharp(\delta)\phi(h)$ by rational equivariance and right invariance under the away-from-$p_0$ frame discrepancy. On $H^{p_0}\times G_{p_0}$ set
\[
 \widetilde\phi(h,h_0)=\xi^\sharp(\delta)\phi((\delta^{p_0})^{-1}h)
 \quad\text{if }h_0\in\delta_{p_0}K_{p_0},
\]
and set it to zero if there is no such $\delta$. Two choices of $\delta$ differ by an element of $\Gamma$ with component in $K_{p_0}$, so the displayed equivariance makes $\widetilde\phi$ well defined. It is a rationally equivariant form of product level $K(N)\times K_{p_0}$. Its slice at $h_0=1$ is nonzero by the choice of $p_0$ at $g_0$ where $\Phi(g_0)\ne0$. On every nonzero slice the formula transports the type and Hecke equations, so the chosen eigensystem occurs in the standard definite summand of the trace, with the indicator $\mathbf1_{K_{p_0}}$ as the test factor at $p_0$. Let $e=e_{\rm typ}$ and $u,u'$ be the type projector and Iwahori operators at $W\mid r_0$ provided by Lemma~\ref{B9.09ac}.
\leanref{UpperTheory.\allowbreak{}Gal.\allowbreak{}Spc.\allowbreak{}eigEx}{Development.\allowbreak{}AutomorphicArithmetic.\allowbreak{}Foundations.\allowbreak{}GaloisRepresentations.\allowbreak{}ScalarCommutation.\allowbreak{}FrameTh}{52--92}
\leanref{UpperTheory.\allowbreak{}Gal.\allowbreak{}W2.\allowbreak{}Seed.\allowbreak{}pHasEigOf}{Development.\allowbreak{}AutomorphicArithmetic.\allowbreak{}Foundations.\allowbreak{}GaloisRepresentations.\allowbreak{}SphericalCoordinates.\allowbreak{}AlgebraEquivSetup.\allowbreak{}PHasEigOf}{45--87}
\leanref{UpperTheory.\allowbreak{}Gal.\allowbreak{}Spc.\allowbreak{}specOfF}{Development.\allowbreak{}AutomorphicArithmetic.\allowbreak{}Foundations.\allowbreak{}GaloisRepresentations.\allowbreak{}ScalarCommutation.\allowbreak{}SpecOfF}{179--228}
 The operators $T_ge$ commute on the direct sum over the definite form classes at this weight and product level $K(N)\times K_{p_0}$. Their finite-dimensional commutative algebra gives the joint generalized eigenspace projector $P_\theta$ as a word polynomial, and $\widetilde\phi$ lies in $\operatorname{im}(P_\theta e)$. By Lemma~\ref{B9.09ac}(1), $u$ commutes with $e$ and the selected Hecke operators $T_g$ and satisfies $u'ue=e$, so its restriction to the nonzero space $\operatorname{im}(P_\theta e)$ is invertible. Newton's identities give an integer $n\ge1$ with $\operatorname{tr}(u^nP_\theta e)\ne0$. Set $f=u^ne\otimes\mathbf1_{K_{p_0}}$. For any word polynomial $s$ with $s(\theta)\ne0$, the operator $s(T)-s(\theta)$ on $\operatorname{im}(P_\theta e)$ is nilpotent and commutes with $u$, so its product with $u^n$ is nilpotent and
\[
 \operatorname{tr}(P_\theta s(T)u^ne)
 =s(\theta)\operatorname{tr}(u^nP_\theta e)\ne0.
\]
Take the word polynomial for $P_\theta$ as $a_0$. By Lemma~\ref{B9.09ac}(2), every element in the away-from-$p_0$ support of $a_0(T)s(T)u^ne$ has separable characteristic polynomial of degree $a$. At the split type place $w_0$ the slice of $e$ is a cuspidal test function: its constant term along every proper parabolic vanishes. Convolution with the selected operators preserves this property, since integrating a translate along a unipotent radical gives the corresponding translate of a constant term along a conjugate parabolic. A regular element with reducible characteristic polynomial is conjugate into a proper block Levi, and its orbital integral factors through the constant term on that Levi because on the off-diagonal blocks the change of unipotent variables has invertible linear factors given by differences of eigenvalues in distinct blocks; hence its orbital integral vanishes. This removes the marked geometric terms with reducible polynomial at $w_0$. On the definite side, Lemma~\ref{B9.09ac}(2) forces every contributing rational similitude to have a separable characteristic polynomial of full degree $a$, and the same constant-term calculation at $w_0$ removes those with reducible polynomial.
\leanref{UpperTheory.\allowbreak{}Gal.\allowbreak{}Spc.\allowbreak{}vanPairs}{Development.\allowbreak{}AutomorphicArithmetic.\allowbreak{}Foundations.\allowbreak{}GaloisRepresentations.\allowbreak{}ScalarCommutation.\allowbreak{}VanPairs}{306--322}
\leanref{UpperTheory.\allowbreak{}Gal.\allowbreak{}Spc.\allowbreak{}vanForms}{Development.\allowbreak{}AutomorphicArithmetic.\allowbreak{}Foundations.\allowbreak{}GaloisRepresentations.\allowbreak{}ScalarCommutation.\allowbreak{}VanForms}{127--192}
\leanref{UpperTheory.\allowbreak{}Gal.\allowbreak{}Int.\allowbreak{}Geo.\allowbreak{}Pl.\allowbreak{}master}{Development.\allowbreak{}AutomorphicArithmetic.\allowbreak{}Foundations.\allowbreak{}GaloisRepresentations.\allowbreak{}IntegralStructures.\allowbreak{}GeometricFlagBases.\allowbreak{}PrimeLevelLocalization.\allowbreak{}SameKTrans}{86--206}
\leanref{UpperTheory.\allowbreak{}Gal.\allowbreak{}Spc.\allowbreak{}specTopG}{Development.\allowbreak{}AutomorphicArithmetic.\allowbreak{}Foundations.\allowbreak{}GaloisRepresentations.\allowbreak{}ScalarCommutation.\allowbreak{}SpecTopG}{9--22}

At a site in $H$, Proposition~\ref{B9.12a} gives the common smooth projective lower model and the integral correspondences, since the residue characteristic of $H\cup C$ is distinct from $r_0$. The deck maps act above its generic fibre, and the weights are those of \ref{B9.18a}. At a site in $C$ division by the level indices gives the normalized spherical Hecke coefficients, with roots $\kappa(w)^{-1}A_a(w,B)$. At usable places off the prime divisors of $Np_0 e$ and the bad places of $\Phi$, Lemma~\ref{B9.09ab}(1) gives goodness of the PEL datum and the same calculation in the chart compatible with $H$ applies. Take $(H,C)=(\{v'\},\varnothing)$ for the initial plan, $(\{v'\},\{w\})$ for an additional root identity, and $(\{v',w\},\varnothing)$ for two Hodge places.
\end{proof}

\subsection*{Proof of \ref{B9.10a}}\label{prf:B9.10a}\phantomsection
\begin{proof}
First rigidify the moduli problem. Choose a positive integer \(e\), prime to \(p\), killing the cokernel of \(\Lambda\to\Lambda^\vee\), and an odd prime \(M>\max(e,p,3)\). Add a symplectic similitude frame of level \(M\). Forgetting the endomorphisms now maps the moduli problem to the fine Siegel moduli scheme of polarized abelian schemes with bounded polarization kernel and auxiliary frame. One may use the separate polarization types here. Their elementary divisors divide \(e\), and there are finitely many since the dimension is fixed. The Siegel construction with invertible auxiliary level is given in \cite[Ch.~I, \S4]{faltings-chai-1990}.

Here is the relative construction over that scheme. Choose an integral basis of the order. Its images are homomorphisms of the universal abelian scheme. In the scheme of graphs of homomorphisms impose the multiplication table and the adjunction identities for this basis. The Rosati norm of each basis element is fixed by the trace of \(c^*c\) on the rational datum. Thus the graphs have bounded degree for the product polarization. Equivalently, their degrees are bounded by the elementary symmetric functions of the positive operator \(1+c^*c\) with bounded trace. This restricts the construction to finitely many Hilbert polynomials. The homomorphism condition is rigid, and homomorphisms between abelian schemes over a discrete valuation ring extend uniquely from the fraction field. These bounded choices are therefore finite over the Siegel base. The level-\(N\) sections are choices in the finite étale group \(A[N]\). The order and pairing identities cut out their prescribed locus. The prime-to-\(p\) lattice type is constant under geometric specialization, by transport of the prime-to-\(p\) Tate module with its form and action. On the Lie bundle impose the factorization through the sum of the indicated eigenspaces. The maps \(t_i\) are pairwise comaximal, so on this quotient the projections are the Chinese remainder idempotents. Impose the indicated ranks on their summands. This gives the required family with auxiliary frame. In particular the Lie condition is imposed on the bundle, hence on points over every \(R\)-algebra.
\leanref{GlobalTheory.\allowbreak{}UPel.\allowbreak{}HT.\allowbreak{}relRep\_\allowbreak{}htPull}{Development.\allowbreak{}AbelianVarietyTheory.\allowbreak{}Proofs.\allowbreak{}DivisorProofs.\allowbreak{}HarrisTaylorStructuresRelativelyRepresentable}{18--66}

Changes of the auxiliary frame act freely. An automorphism preserving the polarization has finite order, and an automorphism acting trivially on a full level \(N\ge3\) is the identity by the congruence lemma \cite[Rem.~1.4.1.9]{lan-2013}. The auxiliary frames of any given object form a finite étale torsor, allowing the primitive root in the pairing to vary. The quotient of the quasi-projective framed scheme by this finite group exists. The abelian scheme and its homomorphisms descend along the torsor. The bundle defining the polarization descends with its rigidification on the axes, and the level sections descend. Conversely, adding frames étale locally to any object gives classifying maps that agree in the quotient. Descent gives one classifying map and its uniqueness. This proves representability and quasi-projectivity.

For smoothness consider a square-zero lifting problem \(C\twoheadrightarrow C_0\) over an Artinian local ring \((C,\mathfrak m_C,k)\) over \(R\). In residue characteristic \(p\), the defining quadratic polynomial of \(\omega\) has distinct roots modulo \(p\), so Hensel's lemma splits \(\mathcal O_{K_0}\otimes_{\Z}\Z_p=\mathcal O_{K_0}^+\times\mathcal O_{K_0}^-\) into the selected and conjugate \(\omega\)-summands, interchanged by the involution. Accordingly the crystal value \(\mathbf D(A_0)(C)\) splits into two \(\omega\)-summands in perfect duality under the prime-to-\(p\) polarization. Using \(\mathcal O\otimes\Z_p\simeq M_m(\mathcal O_{K_0}\otimes\Z_p)\), pass by Morita equivalence on the selected side to one matrix column \(M=\mathbf D(A_0)(C)_1^+\). Because \(A_0/C_0\) is an abelian scheme and \(p\) is nilpotent in \(C\), \(\mathbf D(A_0)(C)\) is finite free over \(C\), and its direct summand \(M\) is finite projective, hence free, over \(C\). We show that \(M\) is free of rank \(n\) over \(\mathcal O_{K_0}^+\otimes_{\Z_p}C\), without assuming \(K_0/\Q_p\) unramified. Let \(\bar k\) be an algebraic closure of \(k\), and let \(D=\mathbf D(A_{\bar k})(W(\bar k))_1^+\) be the selected column of the contravariant Dieudonné module over \(W(\bar k)\). Decompose \(\mathcal O_{K_0}^+=\prod_{\mathfrak w}\mathcal O_{\mathfrak w}\) over the selected primes \(\mathfrak w\mid p\) of \(K_0\); their conjugates \(\bar{\mathfrak w}\) lie on the opposite \(\omega\)-side and are disjoint from them. For each selected \(\mathfrak w\), let \(L_{\mathfrak w}/\Q_p\) be the maximal unramified subextension of \(K_{0,\mathfrak w}\), of degree \(f_{\mathfrak w}\), and let \(\varpi_{\mathfrak w}\) be a uniformizer of \(\mathcal O_{\mathfrak w}\). Then
\[
 \mathcal O_{\mathfrak w}\otimes_{\Z_p}W(\bar k)
 \simeq\prod_{\tau:\mathcal O_{L_{\mathfrak w}}\to W(\bar k)} R_{\mathfrak w,\tau},
 \qquad R_{\mathfrak w,\tau}=\mathcal O_{\mathfrak w}\otimes_{\mathcal O_{L_{\mathfrak w}},\tau}W(\bar k),
\]
where each \(R_{\mathfrak w,\tau}\) is a complete discrete valuation ring with uniformizer \(\varpi_{\mathfrak w}\otimes 1\) because \(\mathcal O_{\mathfrak w}/\mathcal O_{L_{\mathfrak w}}\) is totally ramified and Eisenstein. Since \(D\) is finite free over \(W(\bar k)\), its summand \(D_{\mathfrak w,\tau}\) is torsion-free and finitely generated over \(R_{\mathfrak w,\tau}\), hence free of some rank \(r_{\mathfrak w,\tau}\). The \(\sigma\)-semilinear Frobenius isomorphism on \(D_{\mathfrak w}[1/p]\) permutes the embeddings \(\tau\) transitively and commutes with \(K_{0,\mathfrak w}\), so \(r_{\mathfrak w,\tau}=r_{\mathfrak w}\) is independent of \(\tau\), and \(D_{\mathfrak w}\) is free of rank \(r_{\mathfrak w}\) over \(\mathcal O_{\mathfrak w}\otimes_{\Z_p}W(\bar k)\). By weak approximation in \(K_0\), choose \(c\in\mathcal O_{K_0}\setminus\{0\}\) with \(v_{\mathfrak w}(c)=1\) and \(v_{\mathfrak v}(c)=0\) at every prime \(\mathfrak v\mid p\) of \(K_0\) different from \(\mathfrak w\) (including \(\bar{\mathfrak w}\)). For any prime \(\ell'\ne p\), the prime-to-\(p\) level identifies \(V_{\ell'}(A_{\bar k})\) with \(\Lambda\otimes_{\Z}\Q_{\ell'}\), a free \((K_0\otimes_\Q\Q_{\ell'})\)-module of rank \(mn\), giving \(\deg(c\mid A_{\bar k})=N_{K_0/\Q}(c)^{mn}\) and \(v_p(\deg(c\mid A_{\bar k}))=mnf_{\mathfrak w}\). On the Dieudonné module \(\mathbf D(A_{\bar k})(W(\bar k))\simeq(D)^m\oplus(\mathbf D_1^-)^m\), polarization duality shows that \(c\) acts invertibly on \(\mathbf D_1^-\) (since \(\bar c\) is a unit at all selected primes) and on \(D_{\mathfrak w'}\) for \(\mathfrak w'\ne\mathfrak w\), so \(v_p(\deg(c\mid A_{\bar k}))=\operatorname{length}_{W(\bar k)}(\mathbf D(A_{\bar k})(W(\bar k))/c\,\mathbf D(A_{\bar k})(W(\bar k)))=mr_{\mathfrak w}f_{\mathfrak w}\). Hence \(r_{\mathfrak w}=n\) for every \(\mathfrak w\), and \(M\otimes_C\bar k\simeq D/pD\) is free of rank \(n\) over \(\mathcal O_{K_0}^+\otimes_{\Z_p}\bar k\). Over the finite commutative Artinian \(k\)-algebra \(B_k=\mathcal O_{K_0}^+\otimes_{\Z_p}k=\prod_r B_{k,r}\), each local factor \(B_{k,r}\) has faithfully flat base change \(B_{k,r}\otimes_k\bar k\), so \((M\otimes_C k)_r\) is flat, hence free of rank \(n\), over \(B_{k,r}\); thus \(M\otimes_C k\) is free of rank \(n\) over \(B_k\). Lifting a \(B_k\)-basis gives a map \(\varphi:(\mathcal O_{K_0}^+\otimes_{\Z_p}C)^n\to M\), surjective by Nakayama's lemma over \(C\). Since \(M\) is free over \(C\), \(\operatorname{Tor}_1^C(M,k)=0\), so \((\ker\varphi)\otimes_C k=0\) and \(\ker\varphi=0\) by Nakayama's lemma; therefore \(M\) is free of rank \(n\) over \(\mathcal O_{K_0}^+\otimes_{\Z_p}C\).

On this free module \(M\), the Hodge quotient of type \(t_i\) over \(C_0\) is a line quotient after tensoring along \(t_i:\mathcal O_{K_0}^+\to C\) (which factors through \(\mathcal O_{K_0}^+\) since \(t_i(\omega)\) selects the chosen side). Lift this quotient over \(C\) by lifting its coordinates, one of which is a unit. The choices for different \(i\) are independent by comaximality of the maps \(t_i\). Taking the kernels, restoring the \(m\) columns, and taking the annihilator on the conjugate side gives an order-stable isotropic lift of the Hodge subbundle. In residue characteristic zero the same calculation uses the embedding summands of \(K_0\otimes_\Q C\). By the Grothendieck--Messing theorem for polarized abelian schemes \cite[Ch.~V, \S1]{messing-1972}, the abelian scheme, polarization, and \(\mathcal O\)-action lift uniquely compatibly with this isotropic order-stable subbundle. Prime-to-\(p\) level lifts uniquely in residue characteristic \(p\), and in characteristic zero the entire finite level lifts uniquely. On each of the \(b\) types \(t_i\), the deformation space of a line quotient of a free \(C_0\)-module of rank \(n\) has tangent space \(\operatorname{Hom}_{C_0}(\ker,\text{quot})\) of rank \(n-1\), proving smoothness of relative dimension \(b(n-1)\).
\leanref{GlobalTheory.\allowbreak{}UPel.\allowbreak{}HT.\allowbreak{}exists\_\allowbreak{}sqZero\_\allowbreak{}lift}{Development.\allowbreak{}AbelianVarietyTheory.\allowbreak{}Proofs.\allowbreak{}DivisorProofs.\allowbreak{}HarrisTaylorExistsLiftSquareZeroArtinian}{19--56}

For properness apply the valuative criterion, allowing an extension of the test valuation ring. After finite extension the abelian variety has semiabelian reduction with split torus \(T\), and the order acts on its character lattice \(X^*(T)\). We give the obstruction to \(T\ne0\) in the first case, since it applies to the commutative order. Some embedding of \(K_0\) on the selected side is absent from the generic Lie algebra. Over a splitting extension the idempotent of that embedding annihilates the Lie algebra. In an integral basis of \(\mathcal O_{K_0}\), scale a nonzero such relation so its coefficients belong to the valuation ring and at least one is a unit. The relation holds on the integral Lie module by torsion-freeness, hence on the Lie algebra of the special torus. Its reduction is a nonzero element of \(\mathcal O_{K_0}\otimes k\), where \(k\) is the residue field. If \(X^*(T)\ne0\), it is a finite torsion-free \(\mathcal O_{K_0}\)-module of positive rank. It is locally free of that rank at every prime. Thus \(\mathcal O_{K_0}\otimes k\) acts faithfully on \(X^*(T)\otimes k\), and dually on \(\operatorname{Lie}T\). This contradicts the relation.
\leanref{GlobalTheory.\allowbreak{}PotGoodLie.\allowbreak{}finrank\_\allowbreak{}zero}{Development.\allowbreak{}AbelianVarietyTheory.\allowbreak{}Proofs.\allowbreak{}DivisorProofs.\allowbreak{}PotentialGoodReductionOfLieMultZero}{108--156}

In the division-algebra case any nonzero rational character lattice of \(T\) has dimension at least \(m^2[K_0:\Q]\), the dimension of the division algebra over \(\Q\). But
\[
 \dim T\le\dim A=mn[K_0:\Q]/2<m^2[K_0:\Q].
\]
In both cases the semiabelian model is abelian. The order action and the polarization extend by rigidity, and the prime-to-residue-characteristic level extends by finite étaleness. Here the test trait has residue characteristic \(p\) or zero, since it lies over \(R\). The Lie conditions persist by the closed equations and locally constant ranks in the construction. This gives the valuative extension. A proper quasi-projective morphism is projective.
\leanref{GlobalTheory.\allowbreak{}UPel.\allowbreak{}HT.\allowbreak{}isProper\_\allowbreak{}model\_\allowbreak{}of\_\allowbreak{}compact}{Development.\allowbreak{}AbelianVarietyTheory.\allowbreak{}Proofs.\allowbreak{}DivisorProofs.\allowbreak{}HarrisTaylorIntegralModelIsProperOfCompact}{27--59}
\end{proof}

\subsection*{Proof of \ref{B9.11}}\label{prf:B9.11}\phantomsection
\begin{proof}
Apply \ref{B9.10a} to the given algebra. All nonzero integers are units, and there are no residue maps of characteristic $p$. To compare the two functors, restrict a full Tate frame to its prime-to-$p$ factors. In the reverse direction the polarization kernel is prime to $p$. At a characteristic-zero geometric point the $p$-primary lattice is free after passage to one matrix column at each prime of the centre above $p$. The two sides of $\omega$ are in perfect duality. Choose bases on one side and the dual bases on the other, and restore the matrix columns. This supplies compatible frames at all powers of $p$. Combining them with the given prime-to-$p$ frames by the Chinese remainder theorem preserves the sections of level $N$. Conversely a full frame gives the prescribed alternating lattice at $p$, so the polarization kernel has prime-to-$p$ exponent. The conditions on objects coincide and the isomorphisms preserve the same sections and polarization. The universal property now gives the result.
\leanref{GlobalTheory.\allowbreak{}UPel.\allowbreak{}HT.\allowbreak{}GenCmp.\allowbreak{}geomFull}{Development.\allowbreak{}AbelianVarietyTheory.\allowbreak{}Proofs.\allowbreak{}DivisorProofs.\allowbreak{}GenModelEqGenericFibre}{275--378}
\leanref{GlobalTheory.\allowbreak{}UPel.\allowbreak{}HT.\allowbreak{}GenCmp.\allowbreak{}cmp\_\allowbreak{}of}{Development.\allowbreak{}AbelianVarietyTheory.\allowbreak{}Proofs.\allowbreak{}DivisorProofs.\allowbreak{}GenModelEqGenericFibre}{171--201}
\end{proof}

\subsection*{Proof of \ref{B9.12}}\label{prf:B9.12}\phantomsection
\begin{proof}
Apply \ref{B9.10a} over the given field. The structural maps agree with those of the generic datum. The comparison in \ref{B9.11} carries the universal abelian scheme to the generic one. Equivalently, the resulting map from the generic moduli scheme into the scalar extension of the integral moduli scheme is an open and closed immersion. Smoothness restricts along the open immersion and projectivity along the closed immersion.
\leanref{UpperTheory.\allowbreak{}Gal.\allowbreak{}G3.\allowbreak{}pGenSm}{Development.\allowbreak{}AutomorphicArithmetic.\allowbreak{}Foundations.\allowbreak{}GaloisRepresentations.\allowbreak{}PAdicWittLattices.\allowbreak{}PGenSm}{15--32}

For the universal abelian scheme \(A\to X\), its \(k\)-fold fibre product is smooth over \(X\). The exterior product of the relatively ample diagonal bundles on its factors is relatively ample. Thus the fibre product is projective over \(X\), and hence over the ground field. For \(k=0\) it is \(X\). If the level is \(K\) and the Hecke element is \(g\), the source projection forgets level from \(K\cap gKg^{-1}\) to \(K\), with the conjugation convention for the chosen leg. At a normal level contained in this intersection the lifts are the cosets of the intersection in \(K\). They form a finite étale scheme of constant degree \([K:K\cap gKg^{-1}]\). The source projection on the fibre power is its base change. The target projection uses the isogeny on each abelian factor.
\end{proof}

\subsection*{Proof of \ref{B9.12a}}\label{prf:B9.12a}\phantomsection
\begin{proof}
The chart in \ref{B9.09} supplies a simple root \(r\) modulo \(p\) of the polynomial of \(\omega\). Put \(\mathfrak u=(p,\omega-r)\). Then \(\mathfrak u\subseteq\mathfrak w\). The alternating lattice is perfect at \(p\). If \(p_0\) denotes the auxiliary model prime, the lower level has the form \(N_0\mathfrak a\), where \(N_0\ge3\) is prime to \(pp_0\), \(\mathfrak a+\mathfrak w=\mathcal O_{K_0}\), and \(\mathfrak u^e\subseteq\mathfrak a\) for some \(e\). These are the local chart and level conditions used in choosing the common level. There is one structural map, the composite \(\mathcal O_{K_0}\simeq\mathcal O_E\to R\), and the matrix size in \ref{B9.10a} is one. Reduction of this map has kernel \(\mathfrak w\). All integers prime to \(p\) are units in \(R\). Finally
\[
 2<[K_0:\Q]=2[F:\Q]
\]
because \([F:\Q]\) is positive and even. Thus the first projectivity condition of \ref{B9.10a} applies.

On the chosen side, the local factors at primes different from \(\mathfrak w\) have zero Lie algebra, so their torsion is étale. The \(\mathfrak a\)-level is a frame on such torsion. Its frame scheme is finite étale, as can be seen after trivializing the torsion, when the permissible frames form a finite constant set. Adjoining this level preserves smoothness and projectivity. Taking the fibre powers gives the first assertion by the exterior product argument in \ref{B9.12}.

Over \(E\), trivialize \(A[\mathfrak w]\) étale locally. For every \(j\)-plane choose the corresponding subgroup on the selected side. The polarization determines the complementary subgroup on the conjugate side by taking the annihilator for the pairing on \(A[q^j]\). Dividing by the resulting kernel gives a polarized object, with its transported level and an isogeny whose composite with the reverse isogeny is multiplication by \(q^j\). The quotient by a specified kernel is unique. The choices therefore descend from the constant Grassmannian. Reversing the isogeny gives finite étaleness for the other projection as well.
\leanref{UpperTheory.\allowbreak{}Cb.\allowbreak{}exists\_\allowbreak{}heckeCorrespondence\_\allowbreak{}on\_\allowbreak{}integral\_\allowbreak{}shimuraModel}{Development.\allowbreak{}AutomorphicArithmetic.\allowbreak{}Foundations.\allowbreak{}EtaleClopenDescent.\allowbreak{}ExistsHeckeCorrespondenceOnIntegralShimuraModel}{243--246}

To extend the correspondence, normalize \(X\) in this generic finite étale cover, component by component. The normalization is finite since \(R\) is a localization of a number ring and \(X\) is of finite type over \(R\). Its local rings have no \(R\)-torsion, so it is flat over \(R\).

We give the finiteness argument that extends the second leg. Take the closure of the generic image in \(X\times_R X\). At a geometric point of a special fibre, a limiting isogeny of multiplier \(c=q^j\) has kernel in \(A[c]\). On the chosen side decompose \(A[c]\) by the local idempotents of \(\mathcal O_{K_0}/c\). On each factor the tangent dimension is at most one, by the Lie condition and the single structural map. At the identity its finite local coordinate ring therefore has principal maximal ideal. Indeed a lift of a generator modulo the square generates the ideal by Nakayama. Every nonzero ideal is a power of the maximal ideal, by taking an element of least order. The finitely many other geometric points give translates of the same assertion. There are consequently finitely many choices for the selected part of the kernel even at a nonordinary point. The pairing determines the opposite part. The quotient determines its target and the transported level. Applying this argument to the reverse isogeny gives finitely many sources for a given target. In reversing, multiplication by \(c\) on level \(N_0\) is an automorphism since \((c,N_0)=1\).
\leanref{GlobalTheory.\allowbreak{}UPel.\allowbreak{}HT.\allowbreak{}ES.\allowbreak{}KerFin.\allowbreak{}finite\_\allowbreak{}kernels}{Development.\allowbreak{}AbelianVarietyTheory.\allowbreak{}Proofs.\allowbreak{}DivisorProofs.\allowbreak{}FiniteKernelsPIsogeniesHarrisTaylorCharP}{482--595}

For a specialization along a discrete valuation ring the generic kernel extends by schematic closure in the finite torsion scheme. The closure is finite and torsion-free over the trait, hence flat. The quotient is an abelian scheme, the polarization descends, and separatedness identifies the target with its specified limit. Thus the preceding finite lists contain the geometric fibres of the closed graph under either projection. The projections are proper and quasi-finite, hence finite. Normalizing the graph in the original generic cover now gives the same normalization as for the first projection, by transitivity of integral dependence. This supplies \(t_j\) on \(Y_j\). The homomorphism of the two pulled-back abelian schemes extends uniquely over the normal base. Its identities with the order action, the reverse isogeny, and the polarization extend from the dense generic fibre. Auxiliary étale level is transported through this construction. On a fibre power the first map is obtained by base change along the source, and the second applies the isogeny to each coordinate. This proves the assertions about finiteness and the degree.

At a Hodge place the level in the chart was chosen simultaneously for the selected operators. Their level changes use the prime-to-\(p\) torsion and the étale factors on the chosen side away from \(\mathfrak w\). The first projection amounts to choosing these additional levels or subgroups in finite étale groups. It is finite étale over \(R\). Two such choices give the same lower correspondence precisely when they differ by the subgroup preserving the lower level. Thus the map from an upper correspondence to the pullback of its lower correspondence has degree the corresponding subgroup index. This is the multiplicity used in \ref{B9.19}.
\end{proof}

\subsection*{Proof of \ref{B9.15}}\label{prf:B9.15}\phantomsection
\begin{proof}
The smooth projective geometric fibre is a finite disjoint union of integral components because it is regular and Noetherian. On each component \ref{B9.26} bounds the degrees and gives finite dimension. The Galois action is continuous by the inverse-limit construction from finite-coefficient cohomology \cite{milne-1980}. The correspondences are defined over $E$, so pullback and trace commute with this action. Take the common eigenspace and the direct sum of the actions on the quotients of a composition flag.
\end{proof}

\subsection*{Proof of \ref{B9.16}}\label{prf:B9.16}\phantomsection
\begin{proof}
Write \(q=q_w\). At a geometric special point \(x\), the \(\mathfrak w\)-part of the divisible group has \(\mathcal O_{E_w}\)-height \(a\). Its Lie algebra is one-dimensional and the action on it is through the structural map. Here ordinarity means
\[
                 \#A_x[\mathfrak w](\bar k)=q^{a-1}.
\]
Equivalently, the connected formal \(\mathcal O_{E_w}\)-module has height one.

We explain why this applies above the generic point of each special-fibre component. Work first at a closed geometric point of the lower model of \ref{B9.12a}. Its selected divisible module has connected height \(h\), with \(1\le h\le a\), and étale height \(a-h\). The other selected factors are étale and their conjugates are determined by the polarization. The Serre--Tate deformation theorem therefore identifies deformations of the PEL point with deformations of this one-dimensional divisible \(\mathcal O_{E_w}\)-module. Indeed the theorem is functorial for homomorphisms, the action on the Lie line is structural, and the additional étale level lifts uniquely \cite[Thm.~1.2.1]{katz-1981}. For a uniformizer \(\varpi\) of \(E_w\), the completed deformation ring over the completed maximal unramified base has \(h-1\) parameters for the formal module and \(a-h\) parameters for the extension of the étale part \cite[\S4, Props.~4.2 and 4.5]{drinfeld-1974}. If \(h>1\), the coefficient of degree \(q\) in \([\varpi](X)\) on the special fibre is a deformation parameter up to a unit, in the coordinate of the proof of Proposition~4.2. Where it is nonzero the connected height is one. If \(h=1\) that coefficient is already a unit. Thus every special-fibre component has ordinary generic point. Étale auxiliary level and the smooth abelian fibre powers preserve this assertion for the underlying base point.
\leanref{UpperTheory.\allowbreak{}Gal.\allowbreak{}Ord.\allowbreak{}ordGen}{Development.\allowbreak{}AutomorphicArithmetic.\allowbreak{}Foundations.\allowbreak{}GaloisRepresentations.\allowbreak{}RadicalStalks.\allowbreak{}OrdGen}{68--174}

Pull the integral correspondences of \ref{B9.12a} back to the valuation ring in the statement. Every geometric generic branch has a unique specialization, by finiteness of the first projection. Identify the generic \(\mathfrak w\)-torsion with \(H=k(w)^a\). Reduction maps it onto the geometric points of the special \(\mathfrak w\)-torsion. Surjectivity follows from finite flatness of the torsion group over the valuation ring with algebraically closed fraction field. The kernel is a line \(C\), by the displayed cardinality.

For \(D\subset H\) let \(\bar D\) be its image in \(H/C\). If \(D\cap C=0\), the closure of \(D\) has \(|D|\) distinct special points and its reduction is the étale subgroup specified by \(\bar D\). Write \(b(\bar D)\) for the resulting target, including its fibre coordinates. If \(C\subset D\), the selected kernel additionally contains the connected subgroup of order \(q\). This is the kernel of relative \(q\)-Frobenius on the height-one formal module. The quotient on the selected side is consequently the relative Frobenius of the quotient for \(\bar D\). On the conjugate side the annihilator rule gives the complementary kernel for the multiplier \(q^{\dim D}\). More explicitly, composing the quotient for \(\bar D\) with relative \(q\)-Frobenius multiplies its polarization multiplier by \(q\) and gives exactly these two kernels. The level and all fibre coordinates are transported by the composite, so the target is \(F_*b(\bar D)\).
\leanref{UpperTheory.\allowbreak{}Gal.\allowbreak{}crystallineComparison.\allowbreak{}Mod.\allowbreak{}caseE\_\allowbreak{}of}{Development.\allowbreak{}AutomorphicArithmetic.\allowbreak{}Foundations.\allowbreak{}GaloisRepresentations.\allowbreak{}AnnihilatorModels.\allowbreak{}MCaseE}{13--60}
\leanref{UpperTheory.\allowbreak{}Gal.\allowbreak{}crystallineComparison.\allowbreak{}Mod.\allowbreak{}caseW\_\allowbreak{}of}{Development.\allowbreak{}AutomorphicArithmetic.\allowbreak{}Foundations.\allowbreak{}GaloisRepresentations.\allowbreak{}AnnihilatorModels.\allowbreak{}MCaseW}{13--86}

Let \(A_j\) be the multiset of \(b(U)\) for the \(j\)-planes \(U\subset H/C\). There is one lift containing \(C\) of every \(j\)-plane of \(H/C\), of dimension \(j+1\). The transverse lifts of a fixed \((j+1)\)-plane form a torsor under \(\operatorname{Hom}_{k(w)}(U,C)\), of cardinality \(q^{j+1}\). This proves
\[
 T_{j+1}=F_*A_j+q^{j+1}A_{j+1},\qquad A_0=T_0,\qquad A_a=0.
\]
Substitute into the asserted sum. The two terms involving \(A_j\) have opposite signs and the same coefficient because \(\binom j2+j=\binom{j+1}2\). The boundary values give the cancellations at both ends. The count uses each generic branch once, including branches with the same reduction, so it gives the asserted integral identity of multisets.
\end{proof}

\subsection*{Proof of \ref{B9.17}}\label{prf:B9.17}\phantomsection
\begin{proof}
Use the lower model of \ref{B9.12a}. The deck elements occur among the correspondences with eigenvalue one by \ref{B9.09}. On the image of the generic finite étale level cover, pullback after trace is the sum of the deck actions. Dividing the trace by the deck-group order thus writes each eigenclass as a pullback from the lower generic fibre. If the image is a union of components, extend the descended class by zero on its complement. At the geometric generic points of its special fibre, use the ordinary branches of \ref{B9.16}. Their weighted equality gives the equality of correspondence cycles because the generic branches of each finite correspondence give the multiplicities of its components. The trace formula for branches in \ref{B9.14} gives the induced operator relation. By \ref{B9.27} it holds on the generic fibre. Inertia acts trivially through this specialization. Residue Frobenius acts on the cohomology as inverse arithmetic Frobenius in the geometric pullback convention.

Write $\pi$ for the upper-to-lower level map and $T_j^{\rm up},T_j^{\rm low}$ for the corresponding $j$-th Hecke operators. Let $m_j$ be the finite étale degree from the upper correspondence to the pullback of the lower correspondence. Transitivity of trace gives
\[
 T_j^{\rm up}\pi^*=m_j\pi^*T_j^{\rm low}.
\]
Write the upper eigenvalue as $m_j\vartheta_j$, so the lower one is $\vartheta_j$. The branch relation gives the annihilator $\sum_{j=0}^a(-1)^jq_w^{\binom j2}\vartheta_j X^{a-j}$. Both this relation and trivial inertia pass to the semisimplification. Divide by $\vartheta_0\ne0$ to get the root polynomial of \ref{B9.09}. After transport by $\iota$, the unramified twist multiplies the roots by $\kappa(w)$ to give $P_{w,B}$.
\end{proof}

\subsection*{Proof of \ref{B9.18b}}\label{prf:B9.18b}\phantomsection
\begin{proof}
Fix \(P\). Write \(p_0\) for the model prime. Include in the finite test its factor at \(p_0\), the characteristic function of the prescribed compact subgroup there. In complex uniformization the first homology of a PEL point gives its polarized rational space and finite lattice data. The Hodge filtration is prescribed at the definite embeddings. At the distinguished embedding it amounts to a positive line, whose orthogonal complement specifies the conjugate part of the filtration. Conversely the resulting Hodge structure and lattice give the polarized complex torus by the Riemann bilinear relations. The two legs of a Hecke correspondence change the finite lattice data according to the two level maps. A returning coset therefore gives a rational similitude \((\gamma,c_0)\) preserving the positive line. Here \(c_0\in\Q_{>0}\), and
\[
                         h(\gamma v,\gamma v')=c_0 h(v,v')
\]
for the Hermitian form \(h\). The eigenvalue \(\mu\) on the line is its mark.
\leanref{UpperTheory.\allowbreak{}Gal.\allowbreak{}W2.\allowbreak{}fixedPoints\_\allowbreak{}integral\_\allowbreak{}base\_\allowbreak{}of\_\allowbreak{}uniformisation}{Development.\allowbreak{}AutomorphicArithmetic.\allowbreak{}Foundations.\allowbreak{}GaloisRepresentations.\allowbreak{}SphericalCoordinates.\allowbreak{}FixedPointsIntegralBaseOfUniformisation}{70--190}

The multiplier condition can be read before summing over conjugacy classes. On a double coset \(K t K\) of the full finite test, a returning argument \(g^{-1}(\gamma,c_0)g\), written in the finite trivialization, must have multiplier in \(\nu(t)\nu(K)\). Its multiplier is the diagonal image of \(c_0\). The test at \(p_0\) imposes its own compact-level condition on the localization of this same pair. Since the level multipliers are finite integral units, at most one positive rational \(c_0\) can occur for a given multiplier coset. These assertions apply to each double coset in a linear combination. Unfolding by rational conjugacy divides each finite-level orbit by its stabilizer order. This gives the orbital weight of \((\gamma,c_0)\). The marked reducible terms vanish by \ref{B9.09}.

Suppose \(P\) is irreducible. Identify the CM coefficient field of the PEL datum with \(E\) as in \ref{B9.12a}. Then \(B=E[\gamma]\simeq E[T]/(P)\) is a field of degree \(a\) over \(E\), and the underlying \(E\)-space is a line over \(B\). The adjoint involution on \(B\) sends \(\gamma\) to \(c_0\gamma^{-1}\). At a definite embedding all eigenvalues have squared absolute value \(c_0\). At the distinguished embedding the same holds on the positive line and on its definite orthogonal complement. Thus at every embedding of \(B\) the adjoint acts as complex conjugation. Write \(B^+\) for its fixed field. In particular \(c_0=|\mu|^2\) is fixed by the chosen polynomial and mark.

Fix a definite pair of polynomial \(P\), and denote its form by \(h\). In the polarization convention for the definite pairs, \(h\) is positive definite at the distinguished real place \(v_0\) of \(F\) and negative definite at every other real place. Over \(v_0\) the \(a\) embeddings of \(B\) extending the chosen embedding of \(E\) give the \(a\) eigenlines of \(\gamma\). They give distinct real places of \(B^+\), since conjugating an embedding of \(B\) changes its restriction to \(E\). Give the multiplier \(d\) sign \(+1\) at the place belonging to \(\mu\), sign \(-1\) at the other \(a-1\) places above \(v_0\), and sign \(+1\) at every place over \(v\ne v_0\). Thus the products of the signs of \(h\) and \(d\) leave exactly the marked line positive at \(v_0\) and leave the form negative definite at the other places. There exists \(d\in(B^+)^\times\) with these signs which is a norm from \(B\) at every finite place. Indeed the idèle with the prescribed signs and all finite components one has Artin symbol \((-1)^{a-1}=1\) in the quadratic extension \(B/B^+\). By global reciprocity it is a principal idèle times an idèlic norm. Its principal factor has the required signs, since norms at real places are positive, and its finite components are local norms \cite[Ch.~VI, \S5]{neukirch-1999}.

Replace \(h\) by \(h_d(v,v')=h(dv,v')\). Since \(d\) is self-adjoint this is Hermitian. It has the marked signature, and
\[
             h_d(\gamma v,\gamma v')=c_0 h_d(v,v').
\]
The new pair has exactly the same \(\gamma\) and \(c_0\). Use one \(d\) in the abstract field \(E[T]/(P)\) for this choice of \(\mu\). Its image in any representative is a polynomial in \(\gamma\), so the construction respects conjugacy. Replacing \(d\) by \(d^{-1}\) gives the inverse on marked classes.
\leanref{GlobalTheory.\allowbreak{}LKDef.\allowbreak{}move\_\allowbreak{}definite}{Development.\allowbreak{}AbelianVarietyTheory.\allowbreak{}Proofs.\allowbreak{}DivisorProofs.\allowbreak{}ExistsPolynomialMoveDefiniteSimilitudePairMarked}{750--777}

At each finite place write \(d=u\bar u\) in the local algebra of \(B\). The local choices can be taken from the idèlic norm above. Multiplication by \(u\) identifies \(h_d\) with \(h\) and commutes with \(\gamma\). Transporting the finite trivialization by this isometry preserves the argument of the test away from \(p_0\). At \(p_0\) the separate argument is the localization of \((\gamma,c_0)\) on both sides. The rational centralizer on either side is
\[
       \{(z,s)\in B^\times\times\Q_{>0}:z\bar z=s\}.
\]
Indeed an \(E\)-linear map commuting with \(\gamma\) is multiplication by an element of \(B\), and its similitude condition is \(z\bar z=s\) for both forms. The multiplier \(s\) varies in this centralizer. The element used in the test is the particular member \((\gamma,c_0)\). Thus the local isometries identify the orbits and their stabilizers, including all multiplier conditions.
\leanref{GlobalTheory.\allowbreak{}LKDefTerm.\allowbreak{}cent\_\allowbreak{}eq}{Development.\allowbreak{}AbelianVarietyTheory.\allowbreak{}Proofs.\allowbreak{}DivisorProofs.\allowbreak{}KottwitzTermEqDefinitePairOfLocalIsometry}{90--110}
\leanref{GlobalTheory.\allowbreak{}LKDefTerm.\allowbreak{}term\_\allowbreak{}eq\_\allowbreak{}of\_\allowbreak{}local\_\allowbreak{}isom}{Development.\allowbreak{}AbelianVarietyTheory.\allowbreak{}Proofs.\allowbreak{}DivisorProofs.\allowbreak{}KottwitzTermEqDefinitePairOfLocalIsometry}{185--199}

The definite trace in Proposition~\ref{B9.09} sums over the rational classes of definite forms on \(E^a\) of signature \((a,0)\) at \(v_0\) and \((0,a)\) at the \([F:\Q]-1\) other real places, with the away-from-\(p_0\) finite invariants matching the ambient PEL form \(h_{\rm PEL}\) of Lemma~\ref{B9.09ab}(1); both signatures have real discriminant sign product \((-1)^{[F:\Q]-1}=-1\). Equivalence allows rational change of basis and a positive rational multiplier. Since \(d\in(B^+)^\times\) is a local norm at every finite place, replacing \(h_{\rm PEL}\) on a marked pair by \(h_{d^{-1}}\) preserves all finite local invariants while changing the signature at \(v_0\) from \((1,a-1)\) to \((a,0)\), supplying a definite class with its transported finite trivialization. Lemma~\ref{B9.09ab}(3) provides the standard definite summand \(h_{\rm std}=\mu\,1_a\) in this sum, on which Proposition~\ref{B9.09} places the chosen eigenvector. For each mark the inverse constructions identify the orbital weights in the full sum term by term.
\leanref{GlobalTheory.\allowbreak{}LKDef.\allowbreak{}dTrace}{Development.\allowbreak{}AbelianVarietyTheory.\allowbreak{}Statements.\allowbreak{}PlaceShiftStatements.\allowbreak{}TraceDefinitePairsEqAutomorphicTrace}{80--83} Summing over the distinct roots of \(P\) gives each definite orbital weight the factor \(\sum_{P(\mu)=0}c(P,\mu)\), as asserted.
\leanref{GlobalTheory.\allowbreak{}LKDef.\allowbreak{}cls\_\allowbreak{}equiv\_\allowbreak{}dcls\_\allowbreak{}irr}{Development.\allowbreak{}AbelianVarietyTheory.\allowbreak{}Proofs.\allowbreak{}DivisorProofs.\allowbreak{}MarkedFieldPairsEquivDefiniteFieldPairsOdd}{81--187}

For its application to the trace, one can identify the coefficient on this very pair. In the geometric coordinates the weight in \ref{B9.09} is
\[
          \xi^\sharp(\gamma,c_0)=\xi({}^t\gamma^{-1},c_0^{-1}),
\]
using the formula of \ref{B9.09a} on the ambient matrix and scalar group. Let \(t_v=(t_{v,1},\ldots,t_{v,a})\) be the roots of \(\tau_v P\). Then
\[
 \operatorname{tr}\xi^\sharp(\gamma,c_0)
     =c_0^{-b}\prod_v s_{\lambda_v}(t_v),
\]
where \(s_{\lambda_v}\) is the character of \(W_v\) on diagonal matrices. Indeed taking the dual of the inverse transpose gives the trace of \(W_v\) on a matrix with eigenvalues \(t_v\). The reciprocal polynomial has root lists \(t_v^{-1}\) and \(t_v/c_0\) at the two conjugate embeddings. For the rectangle \((b_v^a)=(b_v,\ldots,b_v)\), the two character factors on these lists multiply to
\[
 s_{(b_v^a)}(t_v^{-1})\,
 s_{\lambda_v+(b_v^a)}(t_v/c_0)
   =c_0^{-\sum_i(\lambda_{v,i}+b_v)}s_{\lambda_v}(t_v).
\]
This follows by factoring the determinant character of power \(b_v\) out of the second factor and using homogeneity. Multiplication over \(v\) gives precisely the trace above. Finally, expanding the diagonal entries of the definite Hecke operator gives this coefficient trace on every returning coset, with the stabilizer weight already computed. Its sum over the definite classes is \(D(F)\), since \ref{B9.09} gives vanishing off the regular irreducible locus. This identifies the definite trace used in the marked-root calculation of \ref{B9.18a}.
\leanref{UpperTheory.\allowbreak{}Gal.\allowbreak{}CertificateRingHomMap.\allowbreak{}pWtCore\_\allowbreak{}of}{Development.\allowbreak{}AutomorphicArithmetic.\allowbreak{}Foundations.\allowbreak{}GaloisRepresentations.\allowbreak{}ProjectiveAutomorphicHecke.\allowbreak{}PWtCore}{45--128}
\leanref{GlobalTheory.\allowbreak{}LKDef.\allowbreak{}trace\_\allowbreak{}definite\_\allowbreak{}forms\_\allowbreak{}identity}{Development.\allowbreak{}AbelianVarietyTheory.\allowbreak{}Proofs.\allowbreak{}DivisorProofs.\allowbreak{}LieModelBaseChangeMap}{14--23}
\end{proof}

\subsection*{Proof of \ref{B9.18a}}\label{prf:B9.18a}\phantomsection
\begin{proof}
Use the test $F=a_0sf$ of \ref{B9.09} at the transported eigensystem. On fibre factor $u$ compose the correspondence with multiplication by $\alpha_u\in2\mathcal O\setminus\{0\}$, where $\mathcal O$ is the order acting on the universal abelian scheme. At a fixed point of the base, let $A_{\rm fib}$ denote the resulting endomorphism on the fibre tangent space. Its eigenvalues are twice algebraic integers since the endomorphism with $\alpha_u/2$ preserves the lattice. Thus the difference from the identity is an isogeny on the fibre. Its kernel has size $|\det(1-A_{\rm fib})|^2$ by the determinant on the underlying integral lattice. In \ref{B9.18} this cancels the fibre denominator and leaves $\overline{\det(1-A_{\rm fib})}\det(1-yA_{\rm fib})$ in the generating function at $T=-y$. The multiplicity of the finer levels cancels the relative Hecke index.

For a regular orbit polynomial $P$ of degree $n$, let $\widetilde P$ be its transport to the field of the PEL datum and put
\[
 P^\vee(X)=\widetilde P(0)^{-1}X^n\widetilde P(X^{-1}).
\]
The mark $\mu$ of $P$ gives the root $\mu^{-1}$ of the reciprocal polynomial on the chosen side of the imaginary quadratic subfield and $\bar\mu^{-1}$ on the other side. Let $t$ be the embedding indexed by $i_0$ and $\bar t$ its conjugate. At each embedding $\eta$, write $r_\eta$ for the roots of $\eta(P^\vee)$. At $t$ use the list $x=(x_0,\ldots,x_{n-1})=r_t$, with mark $x_i$. The base tangent eigenvalues in these coordinates are $x_i/x_j$ for $j\ne i$. The base factor in \ref{B9.18} is therefore $\prod_{j\ne i}(x_j-yx_i)/(x_j-x_i)$. The fibre factors have the form $\prod_j(1-\eta(\alpha_u)r_{\eta,j})$. Scale the marked root by $y$ at $t$ and all roots except the conjugate mark at $\bar t$. At the other embeddings, scale all or none of the roots according to the signature.

Let $\lambda_\eta$ be the conjugated partition at $\eta$ and $s_{\lambda_\eta}$ its Schur polynomial. Each Schur polynomial is a polynomial in elementary symmetric functions. These functions are the signed coefficients of the single-variable determinant polynomial $\prod_j(1-zx_j)$. Interpolate in the multiplication parameters to obtain products of these coefficients. Use a separate fibre factor for each determinant polynomial in a product. The joint embedding coordinates of $2\mathcal O$ are Zariski dense, since an integral basis gives a full lattice and a polynomial vanishing on all integer tuples vanishes identically. Thus the interpolation can use nonzero $\alpha_u$ from $2\mathcal O$. The finitely many conjugated weights and signatures admit a common number $k$ of factors. Denote by $C$ the resulting linear combination of multiplication pullbacks.

For a partition $\lambda$ at $t$, write $x^{(i)}$ for $x$ with $x_i$ replaced by $yx_i$. The identity for the sum over marks is
\[
 \sum_{i=0}^{n-1}s_\lambda(x^{(i)})
 \prod_{j\ne i}\frac{x_j-yx_i}{x_j-x_i}
 =s_\lambda(x)\sum_{i=0}^{n-1}y^{\lambda_i+n-1-i}.
\]
The product is the quotient of the scaled and original Vandermonde determinants. Multiply by the original Vandermonde and expand the Schur alternant with column exponents $\lambda_j+n-1-j$. In the term for mark $i$, the row $i$ contributes the power of $y$ of its assigned column. Summing over $i$ gives the sum over all column exponents.

At $\bar t$, all the roots except the mark are scaled. In the conjugate root variables the base factors are $(x_i-yx_j)/(x_i-x_j)$ for $j\ne i$, since the roots at $t$ have a common squared absolute value by the polarization of the orbit. The scaled Vandermonde quotient includes an additional $y^{\binom{n-1}2}$ from pairs with both roots scaled. In each determinant term the scaled rows contribute the sum of their assigned exponents. After division by the additional Vandermonde power, omitting the column indexed by $i$ gives exponent $\sum_{j\ne i}\lambda_{\bar t,j}+i$.

Put $b_0=\sum_{\eta\ne t,\bar t}\varepsilon_\eta|\lambda_\eta|$, where $\varepsilon_\eta=1$ on the conjugate side of the datum and $0$ on the chosen side, and $|\lambda_\eta|$ is the sum of its parts. A constant partition $m$ contributes the $m$-th power of the determinant. When $\lambda_{\bar t}$ is constant $m$, the exponent vector is $h_i=(n-1)m+b_0+\lambda_{t,i}+n-1-i$. When $\lambda_t$ is constant $m$, it is $h_i=m+b_0+\sum_{j\ne i}\lambda_{\bar t,j}+i$. These are the two weight transports of the plan. In either case the marked sum after applying $C$ is $\prod_\eta s_{\lambda_\eta}(r_\eta)\sum_i y^{h_i}$.

Apply \ref{B9.18b} with $c(P,\mu)$ equal to the interpolated fixed-point contribution just calculated. Its sum over the roots is $\prod_\eta s_{\lambda_\eta}(r_\eta)\sum_i y^{h_i}$. The product of Schur polynomials is the trace of the definite algebraic weight on this class. Summing it with the definite orbital weights gives $D(F)$ by the trace on the finite set of definite level classes: expand the diagonal entries of the Hecke matrix, summing over returning cosets and dividing by their stabilizer orders. Thus \ref{B9.18} gives
\[
 \sum_{j,b}(-y)^j(-1)^b\operatorname{tr}
 (T_F C\mid H^b(Z,\Omega_Z^j))
 =D(F)\sum_i y^{h_i}.
\]
In degree $h_i$, take the finite-dimensional span of all Hecke words applied to the images of $T_f$ in the groups $H^b(Z,\Omega_Z^{h_i})$. The operators commute there by \ref{B9.09}, so the span is the direct sum of joint primary spaces. Choose $s$ with $s(\theta)=1$ and value zero at the other characters. Raise the vanishing factors to powers large enough to kill their primary spaces. The coefficient of $y^{h_i}$ on the right is a positive integer times $D(F)$ and is nonzero. Since $T_F$ is the composite $a_0sT_f$, the image of $sT_f$ is nonzero in some $b$. On its $\theta$-primary space the commuting operators $T_g-\sigma(\iota(\theta_g))$ are nilpotent. Their common kernel contains a nonzero eigenclass.
\end{proof}

\subsection*{Proof of \ref{B9.19}}\label{prf:B9.19}\phantomsection
\begin{proof}
Use the smooth projective lower fibre power over $\mathcal O_{E,(w)}$ of \ref{B9.12a}, and base change it to the completed valuation ring at $w$. Its geometric generic fibre is the scalar extension of the lower generic fibre used in \ref{B9.17}. The assumed surjectivity of the level cover and the deck eigenvalues one give descent of the invariant piece by normalized trace. Smooth proper comparison on the good model gives crystallinity \cite{bhatt-morrow-scholze-2018}. For the jump calculation use the embedding $\sigma\circ\tau$ and eigenvalues $\sigma(\iota(\theta_g))$ together as in \ref{B9.18a}. The same trace and pullback calculation descends the Hodge eigenclass. Its eigenvalue for the lower correspondence is $\sigma(\iota(\theta_g/m_g))$, where $m_g$ is the degree from the upper correspondence to the pullback of the lower one.

Let $V$ be the lower model's total $\ell$-adic cohomology with coefficients extended to $\bar{\Q}_\ell$. Finite dimensionality gives a finite set $D$ of correspondences that defines the same common kernel as the full family. The descended eigenspace is therefore the kernel of the Galois-equivariant map
\[
 V\longrightarrow\bigoplus_{g\in D}V,\qquad
 v\longmapsto\bigl((T_g^{\rm low}-\theta_g/m_g)v\bigr)_{g\in D}.
\]
Comparison is functorial for the correspondence operators. Exactness and strictness of filtered de Rham comparison identify the $h$-th graded piece of this kernel with the kernel of the graded map \cite{fontaine-2004}. After scalar extension through the transported embedding, the descended Hodge eigenclass is a nonzero vector in this graded kernel. Pullback injects the descended eigenspace into the upper eigenspace, so jump $h$ occurs in the upper piece. Exactness on its composition flag preserves the dimensions of the graded pieces upon semisimplification.

On crystalline cohomology the ordinary branch identity of \ref{B9.16} gives the relation with Hecke coefficients by comparing the reductions of the correspondences. Divide the Hecke eigenvalues by the correspondence degrees as in \ref{B9.17} and the resulting polynomial by its leading coefficient. Write $P_R(X)=\prod_{\alpha\in R}(X-\alpha)$ for this normalized polynomial, and let $\Phi$ be the period operator at inverse arithmetic Frobenius. The relation on the geometric comparison module gives nilpotence of $P_R(\Phi)$ on the period module of the piece via its composition flag. Hence every eigenvalue of $\Phi$ lies in $R$. Taking $m$ equal to the dimension of the piece gives the asserted polynomial divisibility, with coefficients transported by $\iota$.
\end{proof}

\subsection*{Proof of \ref{B9.22}}\label{prf:B9.22}\phantomsection
\begin{proof}
The morphism $\operatorname{Spec}B\to\operatorname{Spec}T$ is smooth. At a geometric point of $\operatorname{Spec}B$ let $A$ be its strict henselian local ring and let $\mathfrak p$ be the image prime in $T$. The local ring $T_{\mathfrak p}$ is strictly henselian, since every monic polynomial over it splits. The algebraically closed field $K$ defines a geometric generic point of this strict localization of the base. Local acyclicity of smooth morphisms for locally constant torsion sheaves of invertible order gives
\[
 R\Gamma(\operatorname{Spec}A,\Z/b(i))
 \xrightarrow{\sim}
 R\Gamma(\operatorname{Spec}(K\otimes_T A),\Z/b(i)).
\]
This is the stalk comparison for the map from $\Z/b(i)$ to its derived pushforward from the generic fibre \cite[Prop.~59.93.2 and Lem.~59.93.3]{stacks-2026}. Taking global cohomology gives the assertion.
\leanref{GlobalTheory.\allowbreak{}LocAcycDegOne.\allowbreak{}tateTwistCohomology\_\allowbreak{}genericFibre\_\allowbreak{}henselLine\_\allowbreak{}subsingleton}{Development.\allowbreak{}AbelianVarietyTheory.\allowbreak{}Proofs.\allowbreak{}DivisorProofs.\allowbreak{}TateTwistCohomologyGenericFibreHenselLineSubsing}{36--45}
\end{proof}

\subsection*{Proof of \ref{B9.25}}\label{prf:B9.25}\phantomsection
\begin{proof}
For the trivial bundle choose bounded-below injective complex models for the shifted Tate twists and the projective-space cohomology. Represent the sum of pullbacks followed by cup products with $c_1(L)^j$ by a map of these complexes. By the scheme projective bundle formula, the map is a quasi-isomorphism on evaluation at each representable on which $b$ is invertible \cite{milne-1980}. The mapping cone is thus locally exact on the support of any discrete presheaf satisfying the invertibility condition. The injective-complex comparison gives zero hypercohomology for the cone on that presheaf. In the classifying-space model, a bundle map out of a discrete presheaf factors through the unique vertex of the general linear classifying space and gives the trivial bundle.

Apply the discrete formula in each simplicial degree. Each simplex has vertices, so the invertibility hypothesis applies in every degree. Choose the cocycle representatives of the Chern powers on the space classifying the tautological line, and pull them back to make the maps compatible with the simplicial structure. Comparison by simplicial degree gives the quasi-isomorphism on the bounded-below total complexes.
\end{proof}

\subsection*{Proof of \ref{B9.28}}\label{prf:B9.28}\phantomsection
\begin{proof}
Take $Z_0=\varnothing$. Translate the least jump to zero. For two disjoint nonempty subsets, the largest jump in their union exceeds $a$ times every smaller jump. Its subset has the larger mean. Thus the stronger gaps imply lacunarity, which is preserved by reflection. Inverse transpose relates the two split coordinates. Changing to a conjugate auxiliary place therefore preserves sphericity and the type condition. It reflects the jumps up to translation. Choose coordinates that meet the restrictions to $\Q(\sqrt{-d})$ in \ref{B9.08}. If a target place above $\ell$ is specified, use its side.

Let $\omega$ be the Hecke character obtained from the centre of $U'$ in split coordinates. The $\ell$-adic realization $\eta$ of $\omega|\cdot|^{-\binom a2}$ has transported value
\begin{equation}\label{eq:p9c1-seed-det}
 \iota(\eta(s^{-1}))=q_w^{\binom a2}\prod_{\beta\in B}\beta
 =\prod_{\alpha\in A_a(w,B)}\alpha
\end{equation}
at a usable spherical $w\nmid\ell$ and an arithmetic Frobenius lift $s$. This uses reciprocity with inverse arithmetic Frobenius, as in \ref{B9.50}. At a spherical place above $\ell$ the same character is crystalline, with the right-hand side of \eqref{eq:p9c1-seed-det} as its transported Weil value at degree $-1$.

Apply \ref{B9.29} to a geometric witness of \ref{B9.08} with determinant character $\eta$. The cofinite Frobenius tests give a dense set by \ref{B8.42p}. One can take places of absolute residue degree one, unramified over $\Q$, outside the rational primes below the finite exclusions and the primes dividing $2d\ell$. Their completions are $\Q_p$ and contain the image of $\sqrt{-d}$, so they are usable. On the separable locus of $r$, a monic annihilator of degree $a$ is its characteristic polynomial. Take $V$ to be the inverse of this locus and $X$ the exceptional set of the witness. For any other witness of \ref{B9.08}, rank extraction with the same $\eta$ gives the same traces on the intersection of the two open separable loci and the common cofinite Frobenius tests. This intersection is dense by the same degree-one test. Continuity and semisimplicity identify the extracted representations. At a specified place away from $\ell$, apply \ref{B9.30} to its annihilator and \eqref{eq:p9c1-seed-det}. At a specified place above $\ell$, apply the jump and root witnesses as in the proof of \ref{B9.30}.
\end{proof}

\subsection*{Proof of \ref{B9.29}}\label{prf:B9.29}\phantomsection
\begin{proof}
Put $H=\overline{\rho(G_E)}^{\,\mathrm{Zar}}$ in the general linear group and let $\mathfrak h$ be its Lie algebra. Dependence of $1,M,\ldots,M^a$ is closed, so every $M\in H$ has minimal polynomial of degree at most $a$. The same bound holds on $\mathfrak h$. Use a formal arc $M(t)$ through the identity with tangent $D$, apply dependence to $(M(t)-1)/t$, and specialize at zero. Such an arc exists by smoothness in characteristic zero. Sen's theorem puts the Hodge--Tate operator, with eigenvalues the actual jumps in the chosen convention, in $\mathfrak h\otimes\mathbf C_\ell$ \cite{sen-1980}. If $e_1,\ldots,e_a$ are the selected jumps, the equations $\det(e_i I-D)=0$ on $\mathfrak h$ therefore have a common point $D$ over $\bar{\Q}_\ell$. The degree bound makes $D$ semisimple with exactly these eigenvalues. Write $p_i$ for its spectral projectors.

A member $D'$ of $\mathfrak h$ commuting with $D$ has scalar blocks on the $p_i$-spaces. Indeed $D+tD'$ has pairwise disjoint block spectra for generic $t$, and the sum of the block minimal polynomial degrees is at most $a$. Each nonzero adjoint eigenvalue $e_i-e_j$ determines its ordered pair, since equality of two different differences would give disjoint subsets with equal means. The spectral projections of $\operatorname{ad}D$ thus put each off-diagonal block of $\mathfrak h$ in $\mathfrak h$. Products of composable off-diagonal blocks through three distinct indices are brackets, and an opposite product is the diagonal block of a bracket, hence scalar. The algebra $A$ generated by $\mathfrak h$ and the $p_i$ consequently has $p_iAp_i=\bar{\Q}_\ell p_i$. Conjugation by $H$ preserves $A$, since the $p_i$ are polynomials in $D$. The radical of $A$ annihilates a nonzero subspace on every irreducible summand of $\rho$, by nilpotence. This subspace is $G_E$-stable, so the radical acts by zero on every summand and is zero.

Fix $h\in H$. The primitive idempotents $h^{-1}p_jh$ can be matched to the $p_i$ in each matrix factor of $A$. If $p_iAhp_i\ne0$, then $h^{-1}p_ih$ and $p_i$ are in the same factor, and the matching may prescribe this pair. A unit implementing the matching can be chosen of determinant one in each factor by rescaling the matched basis vectors. It is then a product $u$ of transvections. Each transvection belongs to $H^\circ$, since its square-zero off-diagonal block belongs to $\mathfrak h$. With the matching oriented so that $hu$ permutes the projectors, let $T$ be the identity component of the diagonal subgroup of $H$ in the projectors. Its Lie algebra contains $D$.

For $t=\sum_i t_i p_i\in T$, the eigenvalues of $hut$ on a cycle $C$ of length $b$ are the $b$-th roots of $\lambda\prod_{i\in C}t_i$, for the eigenvalues $\lambda$ of the return map of $hu$ on one block of $C$. A collision with a cycle $C'$ of length $b'$ imposes
\[
 \left(\lambda\prod_{i\in C}t_i\right)^{b'}
 =\left(\lambda'\prod_{i\in C'}t_i\right)^b .
\]
An identity on $T$ would give $b'\sum_{i\in C}e_i=b\sum_{i\in C'}e_i$ by differentiation along $D$, contrary to lacunarity. Choose $t$ outside these finitely many closed sets. Each cycle contributes at least its length to the number of distinct eigenvalues. The degree bound gives exactly $a$ distinct eigenvalues and semisimplicity. In particular, if the matching fixes $i$, the block at $i$ is scalar and $p_i h=c_i p_i(ut)^{-1}\in p_i A$ for some $c_i$.

The span $B=\sum_{h\in H}Ah$ is an algebra normalized by $H$. The linear span of $H$ contains $\mathfrak h$ and the polynomials in $D$, so it equals $B$. For a matrix unit $e_{ij}\in p_iAp_j$ with $i\ne j$, apply the fixed-block conclusion to $(1+xe_{ij})h$ for $x\in\bar{\Q}_\ell$. Its $i$-th diagonal block is scalar, by the fixed-block conclusion if it is nonzero. Subtract two such blocks to get $e_{ij}hp_i\in\bar{\Q}_\ell p_i$. Taking $x=0$ gives the diagonal term. These matrix units span the corners of $A$, so $p_iBp_i=\bar{\Q}_\ell p_i$. The radical argument for $A$ applies to $B$. Its $a$ projectors are primitive and sum to one. One standard module of each matrix factor gives a representation of dimension $a$, realized inside the original module, which is a sum of copies of those standard modules. This gives the two embeddings. Semisimplicity follows since $\rho(G_E)$ spans $B$.

In a fixed inclusion $I$ and left inverse $L$, the representation is $r(g)=L\rho(g)I$ and is continuous. Every finite-index subgroup has Zariski closure containing $H^\circ$. The calculation of $hut$ gives a matrix with nonzero discriminant on the extracted module in the Zariski closure of every open coset. The separable locus is thus dense and open. On its intersection with the original dense set, $Q_g$ is the characteristic polynomial of $r(g)$. Comparing constant terms and using continuity gives $\det r=\eta$.
\end{proof}

\subsection*{Proof of \ref{B9.30}}\label{prf:B9.30}\phantomsection
\begin{proof}
The annihilator gives the support condition and the determinant gives the product condition. The last factor of \eqref{eq:p9c1-general-position} supplies uniqueness in \ref{B9.28}.

Above $\ell$, crystallinity passes to the subrepresentation of each witness. For the jump witness, $\rho\hookrightarrow r^{\oplus m}$ forces every prescribed jump to occur in $r$. The $a$ prescribed jumps are distinct and exhaust its rank. For the root witness, \ref{B9.19} gives the support condition on the crystalline module at a degree-$-1$ Weil element. The determinant comparison and \eqref{eq:p9c1-seed-det} give the product there. The same uniqueness argument identifies the polynomial, and unramifiedness of the crystalline Weil action gives the polynomial at every Weil element.
\end{proof}

\subsection*{Proof of \ref{B9.30a}}\label{prf:B9.30a}\phantomsection
\begin{proof}
Write $O=\mathcal O_{E_{w_0}}$. Smoothness gives a subgroup $J_0=1+\ell^b M_a(O)$ fixing $f$, with $b\ge2$. Put $J_n=1+\ell^{b+n}M_a(O)$, $H_n=J_0/J_n$, and $C_n=KJ_n$. Successive quotients are additive matrix groups in characteristic $\ell$, so $H_n$ is an $\ell$-group. The matrix logarithm and exponential on $J_0$ and $\ell^b M_a(O)$ are inverse and identify the powers of any one element with the integer multiples of its logarithm. Hence $J_0$ is torsion free.

Fix $w\in D$ and put
\[
 A(z)=\prod_{\pi\in\mathfrak S_a}\Phi_w(z_{\pi(1)},\ldots,z_{\pi(a)}),
 \qquad A^*(z)=\overline{A(\bar z^{-1})},
 \qquad p_k=(1-cAA^*)^k .
\]
Here inversion is coordinatewise and $c=(1+\max_{|z_i|=1}|A(z)|^2)^{-1}$. The symmetric Laurent polynomial $p_k$ equals one wherever $A=0$, including off the unit torus. On the torus, $0\le1-c|A|^2\le1$ and $A$ has zero set of Haar measure zero. For the induction on the number of variables, apply the one-variable root bound off the common zero set of the coefficients. Dominated convergence gives $\int|p_k|^2\,dz\to0$ for Haar probability measure $dz$.

For $\GL_a(E_w)$ with hyperspecial volume one, Macdonald's spherical Plancherel theorem gives the normalized Satake transform and the probability density proportional to
\[
 \prod_{i<j}\left|\frac{1-z_i/z_j}{1-q_w^{-1}z_i/z_j}\right|^2
\]
on the unit torus \cite[Ch.~V, \S5.1]{macdonald-1971}. This is its split type-$A$ case. If $h$ is the operator with transform $p_k$, write it as a sum over representatives $g$ of right hyperspecial cosets, with coefficients $h_g$. The density has a finite maximum $M_w$ since $q_w>1$. Plancherel gives
\begin{equation}\label{eq:p9c1-spherical-norm}
 \sum_g|h_g|^2\le M_w\int |p_k|^2\,dz .
\end{equation}
The operator is the identity on the $C_n$-fixed forms in every exceptional summand at $w$.

Let $h_0$ be the number of double cosets for $G(F)\backslash G(\mathbf A_F)/(G_\infty C_0)$ and choose representatives $s$. This quotient is finite by compactness of the definite adelic quotient and openness of the level. At level $C_n$ the representatives $sj$, for $j$ lifted from $H_n$, cover the double cosets. For fixed $h$ and sufficiently large $n$, the translates $sjg$ of each such representative belong to distinct double cosets as $g$ ranges over the support. A collision with $sjg'$ gives a rational point of the form
\[
 \gamma=s g' j c^{-1}j^{-1}g^{-1}x^{-1}s^{-1},
 \qquad c\in C_n,\quad x\in G_\infty ,
\]
by commutation of the distinct local factors. As $jC_nj^{-1}\subset C_0$, the possible $\gamma$ lie in a fixed compact set and are finite by discreteness. The $w_0$ component of $s^{-1}\gamma s$ belongs to $J_n$. Shrinking $J_n$ excludes each nonidentity point in the finite list, since a rational matrix equal to the identity at one completion is the identity. When $\gamma=1$, the two right hyperspecial cosets at $w$ coincide.

Choose actual double-coset representatives from the cover and use a basis of $V$ of dimension $d_V$ to identify the weight forms with a subspace of their coordinate evaluations. The archimedean matrices in these evaluations have entries bounded by a constant $M_V$ by compactness of $G_\infty$. Each row of the Hecke matrix uses at most one translate for any given destination, by separation. Thus its squared Hilbert--Schmidt norm is at most $h_0|H_n|d_V^2M_V^2\sum_g|h_g|^2$. The dimension of a space on which a matrix acts as the identity is at most its squared Hilbert--Schmidt norm, by extending an orthonormal basis of that space. By \eqref{eq:p9c1-spherical-norm}, choose $k$ at each $w$ so that the bound is at most $|H_n|/[2(1+|D|)]$, and then choose $n$ for all the supports. The exceptional span at level $C_n$ satisfies
\begin{equation}\label{eq:p9c1-bad-dimension}
 \dim(\text{span of summands bad somewhere in }D)\le |H_n|/2 .
\end{equation}
\leanref{UpperTheory.\allowbreak{}Gal.\allowbreak{}Gen.\allowbreak{}small\_\allowbreak{}pf}{Development.\allowbreak{}AutomorphicArithmetic.\allowbreak{}Foundations.\allowbreak{}GaloisRepresentations.\allowbreak{}SeparatingAlgebras.\allowbreak{}SmallPf}{5--70}
\leanref{UpperTheory.\allowbreak{}Gal.\allowbreak{}Gen.\allowbreak{}countCore\_\allowbreak{}pf}{Development.\allowbreak{}AutomorphicArithmetic.\allowbreak{}Foundations.\allowbreak{}GaloisRepresentations.\allowbreak{}SeparatingAlgebras.\allowbreak{}CountCore}{86--131}
\leanref{UpperTheory.\allowbreak{}Gal.\allowbreak{}Gen.\allowbreak{}countGen\_\allowbreak{}pf}{Development.\allowbreak{}AutomorphicArithmetic.\allowbreak{}Foundations.\allowbreak{}GaloisRepresentations.\allowbreak{}SeparatingAlgebras.\allowbreak{}CountGen}{39--86}

Put $H=H_n$. On the double coset of $s$ at level $C_0$, define $b(\gamma s x c)$ as the image of the $w_0$ component of $c$ in $H$. A stabilizer at $s$ is finite by discreteness and compactness, so its projection to the torsion-free $J_0$ is trivial. Thus $b$ is well-defined and $b(gc)=b(g)\bar c$ for $c\in C_0$. For a function $v_0:H\to\C$, the form $v_0(b(g))f$ belongs to level $C_n$. Let $P$ be the space of $v_0$ for which it belongs to the good span. The automorphic decomposition at level $C_n$ and \eqref{eq:p9c1-bad-dimension} give $\dim P\ge |H|/2$. The good span and $f$ are preserved by the right action of $C_0$, so $P$ is stable under right translation by $H$.

Transport coefficients by $\iota^{-1}$ and use the supremum norm on functions on $H$. For an integral row $(u_h)_{h\in H}$ annihilating $P$, the matrix $B=(u_{hg^{-1}})_{h,g\in H}$ has rank at most $|H|-\dim P$, since its columns are translates of the annihilating row. It is a linear combination of regular permutation matrices with coefficients $u_h$. Over $\F_\ell$ the regular representation of $H$ has a full invariant flag with trivial successive quotients. Indeed a central element of order $\ell$ has a nonzero stable fixed space on every nonzero representation. Apply induction on the group order to that fixed space and repeat on the successive quotients of the representation. Lift the flag basis over $\Z_\ell$. In that basis an integral correction $\ell D_1$ makes $B$ upper triangular with constant diagonal $s_0=\sum_h u_h$. Every surviving term in the determinant expansion uses at least $\dim P$ correction columns. Hence
\begin{equation}\label{eq:p9c1-group-det}
 |s_0|_\ell^{|H|}=|\det(B+\ell D_1)|_\ell
 \le\ell^{-\dim P},\qquad |s_0|_\ell\le\ell^{-1/2}.
\end{equation}
This gives $v_0\in P$ at distance at most $\ell^{-1/2}$ from the constant function one. For this duality calculation, normalize a largest coordinate of a nonzero vector of $P$ to one. Use the vector to eliminate that coordinate from the target and the other generators. The coordinate projection has norm at most one. An integral row annihilating the reduced subspace lifts to an integral annihilating row by subtracting its value on the pivot vector in the pivot coordinate. Induction reduces to coordinate rows when the approximating subspace is zero, where \eqref{eq:p9c1-group-det} bounds the remaining coordinates of the target.
\leanref{UpperTheory.\allowbreak{}Gal.\allowbreak{}Gen.\allowbreak{}groupDet\_\allowbreak{}pf}{Development.\allowbreak{}AutomorphicArithmetic.\allowbreak{}Foundations.\allowbreak{}GaloisRepresentations.\allowbreak{}SeparatingAlgebras.\allowbreak{}GroupDet}{19--92}
\leanref{UpperTheory.\allowbreak{}Gal.\allowbreak{}Gen.\allowbreak{}alg\_\allowbreak{}pf}{Development.\allowbreak{}AutomorphicArithmetic.\allowbreak{}Foundations.\allowbreak{}GaloisRepresentations.\allowbreak{}SeparatingAlgebras.\allowbreak{}AlgPf}{4--52}
\leanref{UpperTheory.\allowbreak{}Gal.\allowbreak{}Gen.\allowbreak{}stepCore\_\allowbreak{}pf}{Development.\allowbreak{}AutomorphicArithmetic.\allowbreak{}Foundations.\allowbreak{}GaloisRepresentations.\allowbreak{}SeparatingAlgebras.\allowbreak{}StepCore}{9--56}

Subtract $v_0(b(g))f$. The remainder is a pointwise multiple of $f$ with multiplier of norm at most $\ell^{-1/2}$, and is smooth at $w_0$ and $K$-fixed. Repeat on the remainder with its own starting subgroup $J_0$. After $m$ steps the product of the multipliers has norm at most $\ell^{-m/2}$. Taking $m$ with $\ell^{-m/2}\le\epsilon$ proves \eqref{eq:p9c1-generality-relative}.
\end{proof}

\subsection*{Proof of \ref{B9.31}}\label{prf:B9.31}\phantomsection
\begin{proof}
Choose supporting compact opens away from $w_0$ using a restricted tensor of local fixed vectors. Take the hyperspecial subgroups at spherical split places and the unramified levels at nonsplit places off $T$. At $v'$ choose a supporting subgroup $K_v$ inside the hyperspecial subgroup. Let $K$ denote the level with that factor enlarged to the hyperspecial subgroup, and let $K'=K\cap\operatorname{pr}_{v'}^{-1}(K_v)$. Both levels have component $1$ at $w_0$. The representation $U$ has a $K'$-fixed vector.
\leanref{UpperTheory.\allowbreak{}Gal.\allowbreak{}Int.\allowbreak{}Cong.\allowbreak{}exists\_\allowbreak{}level}{Development.\allowbreak{}AutomorphicArithmetic.\allowbreak{}Foundations.\allowbreak{}GaloisRepresentations.\allowbreak{}IntegralStructures.\allowbreak{}CongruenceOddLevel.\allowbreak{}OpenCompactMap}{77--155}

Put $D=Z\cup W$. At $z\in D$ let $\Phi_z$ be \eqref{eq:p9c1-general-position}. It is nonzero since the two monomials in each factor have distinct exponent vectors. Add the spherical operator whose transform is $\prod_\pi\Phi_z(\pi x)$ in Satake variables $x$, for each $z\in D$. At $w\in W$ use the Laurent spherical algebra of \ref{B9.30a} to add the operator with transform $\sum_i(q_w^{(a-1)/2}x_i)^n$ for every $n\in\Z$. Write each operator as a finite sum of right coset translations with its coefficients. Use the same finite-translation description for the given $f_j$.

For a complex scalar $z$ put $|z|_*=|\iota^{-1}z|_\ell$. On an algebraic weight representation $V$ use the coefficient lattice $L_V$ in its polynomial realization with determinant twists, transported by $\iota$. Let $\kappa(g)$ be the tuple of matrices at places above $\ell$, transported to the complex weight factors using their inducing embeddings, and let $g_\infty$ be the corresponding archimedean tuple. For the weight action, also denoted by $V$, put
\[
 b_g=V(g_\infty^{-1}\kappa(g))L_V,\qquad
 \|f\|=\sup_{g,\ v\in b_g}|f(v,g)|_* .
\]
This is finite at a full supporting level. One can use finitely many double-coset representatives. Rational translation cancels in $g_\infty^{-1}\kappa(g)$, the archimedean factor cancels by equivariance, and the compact finite level changes $L_V$ within a bounded multiple of itself. The lattices span $V$, so their evaluations separate forms.

For a finite translation $h$ and its coefficient $t(h)$, $b_{gh}=V(g_\infty^{-1}\kappa(g))V(\kappa(h))L_V$. Thus the contraction criterion for each term is
\begin{equation}\label{eq:p9c1-norm-contraction}
 t(h)V(\kappa(h))^{-1}L_V\subset L_V,\qquad
 \|t(h)f(\,\cdot\,,\,\cdot\,h)\|\le\|f\|.
\end{equation}
The given $f_j$ satisfy it since $w_j\nmid\ell$. For each added operator multiply by a power of $\ell$ clearing the denominators in its finitely many matrices. The denominators come from the polynomial action of the inverses and their determinant powers, in the weight factors at which the translations act. All added operators have trivial component below $v'$. Thus the scalars still work when the weight changes at the archimedean factor of an embedding inducing $v'$. Denote these scalars by $c_{w,n}$ for the power sums and $c_z$ for the nonvanishing tests. They are chosen before $J$.
\leanref{UpperTheory.\allowbreak{}Gal.\allowbreak{}Gen.\allowbreak{}frameW\_\allowbreak{}pf}{Development.\allowbreak{}AutomorphicArithmetic.\allowbreak{}Foundations.\allowbreak{}GaloisRepresentations.\allowbreak{}SeparatingAlgebras.\allowbreak{}FrameW}{14--120}
\leanref{UpperTheory.\allowbreak{}Gal.\allowbreak{}Int.\allowbreak{}Cong.\allowbreak{}exists\_\allowbreak{}powOp}{Development.\allowbreak{}AutomorphicArithmetic.\allowbreak{}Foundations.\allowbreak{}GaloisRepresentations.\allowbreak{}IntegralStructures.\allowbreak{}CongruenceOddLevel.\allowbreak{}ExistsTestOp}{64--90}
\leanref{UpperTheory.\allowbreak{}Gal.\allowbreak{}Int.\allowbreak{}Cong.\allowbreak{}exists\_\allowbreak{}testOp}{Development.\allowbreak{}AutomorphicArithmetic.\allowbreak{}Foundations.\allowbreak{}GaloisRepresentations.\allowbreak{}IntegralStructures.\allowbreak{}CongruenceOddLevel.\allowbreak{}ExistsTestOp}{93--138}
\leanref{UpperTheory.\allowbreak{}Gal.\allowbreak{}Int.\allowbreak{}Cong.\allowbreak{}exists\_\allowbreak{}scale}{Development.\allowbreak{}AutomorphicArithmetic.\allowbreak{}Foundations.\allowbreak{}GaloisRepresentations.\allowbreak{}IntegralStructures.\allowbreak{}CongruenceOddLevel.\allowbreak{}ExistsTestOp}{141--173}

Use $I=J\sqcup(W\times\Z)\sqcup D$ for all coordinates and write $\theta^{\rm aug}$ for the tuple including the tests. Normalize a $K'$-fixed weight form of $U$ to one at a nonzero lattice evaluation, and write $R\ge1$ for its norm. Apply \ref{B9.30a} with $\epsilon/R$ at level $K'$, which is hyperspecial at $D$. Write its approximant as $\sum_i F_i$ in automorphic summands, with tuples $\eta_i$ and transported evaluations $b_i$ at the normalization point. By \eqref{eq:p9c1-generality-relative} and \eqref{eq:p9c1-norm-contraction}, applying any ordered monomial in the operators to the difference preserves the bound $\epsilon$. Integer linear combinations preserve it by the ultrametric inequality. Evaluation gives
\begin{equation}\label{eq:p9c1-evaluation-approx}
 \left|P(\theta^{\rm aug})-\sum_i b_iP(\eta_i)\right|_\ell\le\epsilon
 \qquad(P\in\Z[X_j:j\in I]).
\end{equation}
On each summand the value of a monomial depends on its exponents, since the action is scalar. The weight and type approximations use the same evaluation argument and operators.
\leanref{UpperTheory.\allowbreak{}Gal.\allowbreak{}Gen.\allowbreak{}toPolyX\_\allowbreak{}pf}{Development.\allowbreak{}AutomorphicArithmetic.\allowbreak{}Foundations.\allowbreak{}GaloisRepresentations.\allowbreak{}SeparatingAlgebras.\allowbreak{}ToPolyX}{11--90}

Fix a normalized intermediate eigenform $f$ at level $K'$ and an embedding $\tau$ inducing $v'$. Choose a torsion-free principal congruence subgroup $J_0$ fixing $f$ at $w_0$ and put $C=KJ_0$. In the decomposition $g=\gamma s x c$ at this level write $b(g)=c$, with one chosen $s$ per double coset. The finite stabilizer projects trivially to $J_0$. Its rational points are then the identity, since they are the identity at $w_0$. Thus $b$ is left rational invariant and right $G_\infty$-invariant, with $b(gc)=b(g)c$. For $c\in C$ put $f_c(v,g)=f(v,g b(g)^{-1}c)$. These are $C$-fixed forms of the original weight. There are finitely many distinct $f_c$, since $f$ is fixed by $C\cap\operatorname{pr}_{v'}^{-1}(K_v)$. Increase $R$ to a common bound for their norms and that of $f$.

Let $k(g)$ be the matrix at $v'$ transported by the embedding extending $\iota^{-1}\tau$ and then by $\iota$. Write $g_\tau$ for the archimedean matrix, and put $M(g)=g_\tau^{-1}k(g)$ and $n(g)=k(b(g))M(g)^{-1}$. Then $n$ is left rational invariant and right $C$-invariant, with $n(gx)=n(g)x_\tau$ for $x\in G_\infty$. Choose $0<e<1$ with $eR\le\epsilon$. Partition $C$ into finitely many cosets $c_q C_1$, where $C_1$ intersects the inverse image of $K_v$ with the subgroup $|k(c)-1|_*\le e$ in entrywise maximum norm. All matrices $k(c)$ and their inverses are integral. The images of the cosets under $k$ are disjoint compact sets.

There are polynomials $P_q$ in matrix entries with
\begin{equation}\label{eq:p9c1-weight-separator}
 \left|P_q(k(c))-{\bf1}_{c_q C_1}(c)\right|_*\le e
 \qquad(c\in C).
\end{equation}
To construct one, take $d_0>0$ below the distance between the image of its coset and the image of the complement. Cover the coset image by finitely many balls of radius $ed_0$ with centres in that image. For each centre use the coordinate differences from it. On the compact range of a coordinate difference, construct a polynomial bounded by one, within $e$ of one at $|x|_*\le ed_0$, and of norm at most $e$ at $|x|_*\ge d_0$. Cover the last set by finitely many balls $|x-c|_*\le |c|_*/2$. Start with $p=1$ and take the centres in increasing order of norm. At each centre multiply by $(1-x/c)^m$ with $m$ chosen so that $2^{-m}\sup_{|x|_*\le|c|_*}|p(x)|_*\le e$. Later factors have norm at most one on the earlier balls, and all are within $e$ of one on the innermost ball. Their product has the required bounds. Take the product of the coordinate tests at each centre, then one minus the product of the complements of these tests. If the coset image is all of $k(C)$ use the constant one.

For a dominant integral weight $\nu=(\nu_0,\ldots,\nu_{a-1})$, put
\[
 D_\nu(h)=\det(h)^{\nu_{a-1}}
 \prod_{j=1}^{a-1}\det(h_{[1,j],[1,j]})^{\nu_{j-1}-\nu_j}.
\]
For integral matrices $h$ with integral inverse, $|D_\nu(h)|_*\le1$. If $|h-1|_*\le e$, every leading minor is within $e$ of one, so $|D_\nu(h)-1|_*\le e$, including for the signed determinant power. Define
\[
 B(y_1,y_2;v,g)
 =\sum_q P_q(n(g)y_2)\,
 D_\nu\big(k(c_q)^{-1}n(g)y_1\big)f_{c_q}(v,g)
 \qquad(y_1,y_2\in\GL_a(\C)).
\]
At $y_1=y_2=M(g)$ the matrix $n(g)M(g)$ is $k(b(g))$. For the coset of $b(g)$ the minor factor is within $e$ of one and $f_{c_q}(v,g)=f(v,g)$, since $k(c_q^{-1}b(g))$ comes from $K_v$. By \eqref{eq:p9c1-weight-separator} the other terms have norm at most $eR$ on the lattice evaluations. Hence
\begin{equation}\label{eq:p9c1-weight-error}
 \sup_{g,\ v\in b_g}|B(M(g),M(g);v,g)-f(v,g)|_*\le eR .
\end{equation}
\leanref{UpperTheory.\allowbreak{}Gal.\allowbreak{}Wt.\allowbreak{}famOf\_\allowbreak{}apply}{Development.\allowbreak{}AutomorphicArithmetic.\allowbreak{}Foundations.\allowbreak{}GaloisRepresentations.\allowbreak{}WeightCharacters.\allowbreak{}FamOfApply}{80--83}
\leanref{UpperTheory.\allowbreak{}Gal.\allowbreak{}Wt.\allowbreak{}famEx\_\allowbreak{}of}{Development.\allowbreak{}AutomorphicArithmetic.\allowbreak{}Foundations.\allowbreak{}GaloisRepresentations.\allowbreak{}WeightCharacters.\allowbreak{}FamExOf}{13--80}

For fixed parameters $y_1,y_2,v$ these are automorphic functions at level $C$. The archimedean translates span a finite-dimensional space, since they range over translates of polynomials of bounded degree with fixed determinant power and of the original weight vectors. The functions at other archimedean places retain the original weight. Project the family to the finitely many automorphic summands at this level and these archimedean types before substituting $y_i=M(g)$. Each projection has a scalar Hecke tuple. The substitution preserves the scalar action since $M(gh)=M(g)$ for every translation used by the operators.

Put $s_k(\lambda)=\sum_{i<k}\lambda_i$. For a constant $d$ independent of $\nu$, every nonzero projected weight $\mu$ at $\tau$ satisfies
\begin{equation}\label{eq:p9c1-weight-shape}
 |s_k(\mu)-s_k(\nu)|\le d\quad(1\le k\le a),
 \qquad |\mu_i-\nu_i|\le2d\quad(0\le i<a).
\end{equation}
For the upper bound use the circle $h_k(z)=\operatorname{diag}(z I_k,I_{a-k})$, $|z|=1$, in the chosen archimedean coordinates. By expansion of minors of products, a translated $j$-row minor contributes at most $\min(j,k)$ to the exponent. Add the determinant contribution $k\nu_{a-1}$ to obtain the upper exponent $s_k(\nu)$. Take $d$ to bound the degrees of the $P_q$ plus the absolute circle exponents on $V$. Translates of a nonzero projection span its irreducible weight factor, whose largest circle exponent is $s_k(\mu)$. Thus $s_k(\mu)\le s_k(\nu)+d$. For the reverse bound some value of the projection with $y_1=1$ is nonzero. This follows by unitary times upper triangular decomposition of $y_1$ and the identity $D_\nu(hb)=D_\nu(h)\prod_i b_{ii}^{\nu_i}$ for upper triangular $b$, using $n(gx)=n(g)x_\tau$. Acting by the circle on this value pulls out $z^{s_k(\nu)}$ from the minor factor. All remaining exponents are at least $-d$, and projection commutes with the circle. Thus $s_k(\mu)\ge s_k(\nu)-d$. Subtracting consecutive sums proves the second bound.
\leanref{UpperTheory.\allowbreak{}Gal.\allowbreak{}Wt.\allowbreak{}famSum\_\allowbreak{}of}{Development.\allowbreak{}AutomorphicArithmetic.\allowbreak{}Foundations.\allowbreak{}GaloisRepresentations.\allowbreak{}WeightCharacters.\allowbreak{}FamSum}{13--61}

Take $N=(a+1)(4d+a)+1$ and $\nu_i=-N(a+1)^i$. By \eqref{eq:p9c1-weight-shape}, the jumps of each projection are $j_i=i+N(a+1)^i+e_i$, with $|e_i|\le2d$. For $i+1<a$,
\[
 (j_{i+1}-j_0)-a(j_i-j_0)
 \ge N((a+1)^i+a-1)+1-(a-1)i-4ad>0 .
\]
The rank-one gap condition is empty. The summands have $K$-fixed vectors, hence are spherical at $v'$. Apply the evaluation argument of \eqref{eq:p9c1-evaluation-approx} with \eqref{eq:p9c1-weight-error}. The difference is a function of the original weight because $M(gx)=x_\tau^{-1}M(g)$ cancels the archimedean change in $n(gx)$ after substitution. Its tested translates satisfy \eqref{eq:p9c1-norm-contraction}. On the projections the substitution preserves the eigenvalue calculation. This gives polynomial approximation by the tuples of the new weights.
\leanref{UpperTheory.\allowbreak{}Gal.\allowbreak{}Wt.\allowbreak{}top\_\allowbreak{}of}{Development.\allowbreak{}AutomorphicArithmetic.\allowbreak{}Foundations.\allowbreak{}GaloisRepresentations.\allowbreak{}WeightCharacters.\allowbreak{}TopOf}{16--105}

To impose the type, take a form of one of the new weights and choose $J_m=1+\ell^m M_a(O)$ in its stabilizer at $w_0$, where $O=\mathcal O_{E_{w_0}}$ and $m\ge2$ can be arbitrarily large. Choose an additive map $p:O\to\Z/\ell^m$ with $p(1)=1$ by extending $1$ to a $\Z_\ell$-basis of $O$. Take $B_0\in M_a(O)$ with irreducible characteristic polynomial on reduction over the residue field, for example a companion matrix. Define
\[
 \lambda(1+\ell^m X)=p(\operatorname{tr}(B_0 X)).
\]
The product cross term maps to zero since $p(\ell^m O)=0$. For $a\ge2$ and a nonzero proper subspace $L\subset E_{w_0}^a$, the lattice $L\cap O^a$ is a saturated direct summand. Its reduction is a proper nonzero subspace and fails to be stable under $\bar B_0$. Choose $v\in L\cap O^a$ and an integral row $t$ killing $L$ with $tB_0v$ a unit, and rescale to $tB_0v=1$. Then $1+\ell^m vt$ has $\lambda$-value one. Taking $L$ to be the first proper step of a parabolic flag makes this element a member of its unipotent radical. For a nontrivial character of $\lambda$, a vector transforming by that character has zero image in the unipotent coinvariants. Each translate has the same property by applying the calculation to the conjugated flag. The translates span the irreducible local representation, so its proper Jacquet modules vanish. The same calculation proves surjectivity of $\lambda$. In rank one take $B_0=1$ and use $p(1)=1$. The Jacquet condition is vacuous in that rank. Torsion freeness follows as in \ref{B9.30a}.
\leanref{GlobalTheory.\allowbreak{}CongW.\allowbreak{}TypeExists.\allowbreak{}exists\_\allowbreak{}type\_\allowbreak{}surj}{Development.\allowbreak{}AbelianVarietyTheory.\allowbreak{}Proofs.\allowbreak{}DivisorProofs.\allowbreak{}ExistsTorsionFreeSurjZmod}{168--355}
\leanref{UpperTheory.\allowbreak{}Gal.\allowbreak{}Int.\allowbreak{}Cong.\allowbreak{}exists\_\allowbreak{}ellType}{Development.\allowbreak{}AutomorphicArithmetic.\allowbreak{}Foundations.\allowbreak{}GaloisRepresentations.\allowbreak{}IntegralStructures.\allowbreak{}CongruenceOddLevel.\allowbreak{}ExistsEllType}{15--50}

At level $KJ_m$ define $b_\lambda(\gamma s x c)=\lambda(c_{w_0})$ using double-coset representatives. As in \ref{B9.30a}, stabilizers have trivial image. For each character $\chi$ of $\Z/\ell^m$, the form $\chi(b_\lambda(g))f$ transforms by $\chi\circ\lambda$ at $J_m$. The negative sum over $\chi\ne1$ has the elliptic type in each nonzero constituent, and its difference from $f$ is the scalar multiple $\ell^m{\bf1}_{b_\lambda(g)=0}f$ by the finite geometric sum. Its relative error is at most $\ell^{-m}$. It preserves the weight and $K$-level. Apply \eqref{eq:p9c1-evaluation-approx} with $m$ large enough to give the desired error.
\leanref{UpperTheory.\allowbreak{}Gal.\allowbreak{}Gen.\allowbreak{}tyStep\_\allowbreak{}pf}{Development.\allowbreak{}AutomorphicArithmetic.\allowbreak{}Foundations.\allowbreak{}GaloisRepresentations.\allowbreak{}SeparatingAlgebras.\allowbreak{}TyStep}{9--59}
\leanref{UpperTheory.\allowbreak{}Gal.\allowbreak{}Int.\allowbreak{}Cong.\allowbreak{}stepSt\_\allowbreak{}of\_\allowbreak{}ex}{Development.\allowbreak{}AutomorphicArithmetic.\allowbreak{}Foundations.\allowbreak{}GaloisRepresentations.\allowbreak{}IntegralStructures.\allowbreak{}CongruenceOddLevel.\allowbreak{}NotApartOfOverPair}{76--109}

Two polynomial approximations compose. If the outer error is $\epsilon/2$ with weights $b_i$, using inner error $\delta=\epsilon/[2(1+\sum_i|b_i|_\ell)]$ and multiplying the weights gives error at most
\begin{equation}\label{eq:p9c1-compose}
 \epsilon/2+\delta\sum_i|b_i|_\ell\le\epsilon .
\end{equation}
Compose the weight and type steps in this way. For the good intermediate tuples $\eta_i$ from \eqref{eq:p9c1-evaluation-approx}, put $Q=\prod_{z\in D}X_z$ in the test variables, with empty product one. Each $Q(\eta_i)$ is nonzero. At inner tolerance $\delta$ approximate $\eta_i$ by the weight and type steps to tolerance $\delta|Q(\eta_i)|_\ell$. Test $PQ$ and divide by $Q(\eta_i)$. This gives the same approximation at tolerance $\delta$ with weights $c_xQ(x)/Q(\eta_i)$ in place of $c_x$, so every failing tuple has weight zero. Apply \eqref{eq:p9c1-compose} to the first approximation and omit the test coordinates. The $K$-fixed vectors supply all the retained level conditions.

For rational $f_j$, choose one nonzero eigenform in $U$ at a full supporting level. Its complex conjugates under $\operatorname{Aut}(\C/\Q)$ on the finite adeles extend to forms at the same level and in the finite orbit of the archimedean weight. Indeed the corresponding algebraic weight models can be conjugated through the embeddings of a number field of definition. At fixed level and this finite set of weights the definite class set gives a finite-dimensional space. The rational translations commute with conjugation. Distinct conjugate eigenvalue tuples give linearly independent eigenvectors there, by applying an operator separating two distinct tuples and inducting on the number of vectors. The tuple therefore has finite orbit. Each coordinate is algebraic over $\Q$, and every finite subtuple generates a field of degree at most the orbit size. Choose a finite subtuple attaining the maximal degree. Its field contains every coordinate, and its image under $\iota^{-1}$ generates a finite extension of $\Q_\ell$.
\leanref{GlobalTheory.\allowbreak{}Cong.\allowbreak{}finU\_\allowbreak{}holds}{Development.\allowbreak{}AbelianVarietyTheory.\allowbreak{}Proofs.\allowbreak{}SatakeParameters.\allowbreak{}LinearIndependentSimulEigen}{49--81}
\end{proof}

\subsection*{Proof of \ref{B9.33}}\label{prf:B9.33}\phantomsection
\begin{proof}
Apply \ref{B9.32} to $X_i-X_j$. At each tolerance the finite trace sum is continuous, so samples in a sufficiently small neighborhood of any group element satisfy the Cauchy condition. The limit lies in the complete field $L$ and agrees with every sampled value. The same inequalities give continuity. Take limits in each finite tuple of samples to obtain \ref{B9.32} on $\Gamma$. The trace value at the identity follows by testing $X_1-a$ with the variable indexed by the identity. Cyclic trace identities and the alternating degree-$(a+1)$ trace identity are integer polynomials in the trace variables of group elements and their products. They vanish on the family and hence on the limit. Apply \ref{B9.34}.
\end{proof}

\subsection*{Proof of \ref{B9.35}}\label{prf:B9.35}\phantomsection
\begin{proof}
Let $S$ be the finite set of split places where $U$ has no spherical vector. It is conjugation-stable since inverse transpose preserves the integral subgroup. For each finite specified set $Y$, take $Z_0(Y)$ from \ref{B9.28}. Let $Z(Y)$ adjoin to it the split spherical members of $Y$ away from $\ell$. Let $W(Y)$ be the spherical members of $Y$ above $\ell$ separated from the pair. By the field hypotheses, every place above $\ell$ is split, and $Z_0(Y)$ consists of split spherical places separated from the pair. Thus \ref{B9.31} applies to $Z(Y),W(Y)$. Its level, weight, type and general-position conditions give the hypotheses of \ref{B9.28}. The fixed levels give sphericity off $S$ at split places separated from the pair, and sphericity at $v'$ also gives it at its conjugate.

Take indices $j=(w,A,s,k)$ with $w\nmid\ell$ split spherical for $U$, $A$ a prolongation, $s$ an arithmetic Frobenius lift at $A$, and $k\ge1$. Put $x_j=s^{-k}$ and $\theta_j=\iota^{-1}\big(\sum_{\alpha\in A_a(w,B)}\alpha^k\big)$. The spherical Hecke operator for the $i$-th elementary symmetric function of the algebraic roots is $q_w^{i(i-1)/2}$ times the double-coset generator with $i$ uniformizer entries. Indeed the normalized Satake transform of the generator is $q_w^{i(a-i)/2}$ times the elementary symmetric function of the unscaled roots. Newton's identities express the positive power sums as integer polynomials in these generators. They give rational $\ell$-integral operators since $q_w$ is a unit at $\ell$. The levels are hyperspecial by \ref{B9.31}, whose rationality clause for $Y=\varnothing$ puts all the $\theta_j$ in one finite extension $L/\Q_\ell$.
\leanref{GlobalTheory.\allowbreak{}FrobOpsIdx.\allowbreak{}frobOpsIdx}{Development.\allowbreak{}AbelianVarietyTheory.\allowbreak{}Proofs.\allowbreak{}SatakeParameters.\allowbreak{}FrobOpsIdx}{20--94}

Index the seeds by $i=(Y,U',\eta_i)$ with ordinary Hecke tuple $\eta_i$ of an approximant at that $Y$. Choose $\rho_i$ by \ref{B9.28} and put $D_i=\{j:\eta_i(j)=\operatorname{tr}\rho_i(x_j)\}$. For every finite $H$ of seed indices, $x(\bigcap_{i\in H}D_i)$ is dense. Take an open coset in $G_E$ for the arithmetic Frobenius lifts. Intersect it with the dense open sets in \ref{B9.28} for $i\in H$, and choose a smaller coset of an open normal subgroup inside the intersection. Apply \ref{B8.42p} over $E$ away from the ramified primes. Exclude the rational primes below $S$ and the finite exceptional sets, as well as those dividing $2d\ell$. The place has completion $\Q_p$, so it is usable. Taking $k=1$ in the inverse coset and using the characteristic polynomials in \ref{B9.28} gives the simultaneous trace equalities.

For $g\in G_E$, let $x_j$ approach $g$ with $j$ eventually in every specified finite intersection of the $D_i$. Density makes these requirements consistent. Given $\delta>0$, choose a polynomial approximation of $\theta$ at precision $\delta$ from the empty-$Y$ family. On the intersection of the $D_i$ for its finite support, testing $X_j$ bounds the distance from $\theta_j$ to the continuous trace sum evaluated at $x_j$ by $\delta$. Near $g$ that trace sum is within $\delta$ of its value at $g$. Thus $\theta_j$ is Cauchy under the stated requirements. Let $t(g)\in L$ be its limit.

For any other polynomial approximation, impose the $D_i$ of its support in the same limit, including when its specified set is different. For a polynomial in finitely many group variables, choose a sample for each variable in this common intersection, substitute the corresponding indices, and take limits. This gives \ref{B9.32} on $G_E$ with the original weights and tolerance of each approximation. In particular,
\begin{equation}\label{eq:p9c1-global-trace-bound}
 \left|t(g)-\sum_i b_i\operatorname{tr}\rho_i(g)\right|_\ell\le\delta
 \qquad(g\in G_E).
\end{equation}
Apply \ref{B9.33} with the identity sample map to obtain the continuous semisimple $r$ of trace $t$.

Fix a usable spherical $w\nmid\ell$. All seeds in the family $Y=\{w\}$ are unramified there by \ref{B9.28}, since $w\in Z(Y)$ and is separated from $S$. For $s\in I_w$, use their approximations in \eqref{eq:p9c1-global-trace-bound} at $g$ and $gs$ and let $\delta$ tend to zero. This gives $t(gs)=t(g)$. By the trace pairing of \ref{B9.34}, $r(s)-1$ is zero, since it pairs to zero with the spanning matrices $r(g)$. For each index $j$ at $w$, the specified-place assertion gives $j\in D_i$ for every seed in this family. Testing $X_j$ in the ordinary Hecke tuple and using \eqref{eq:p9c1-global-trace-bound} at $x_j$ identifies $\theta_j=t(x_j)$ in the limit. The first $a$ powers give $P_{w,B}$ by Newton's identities.

At spherical $w\mid\ell$ separated from the pair, use $Y=\{w\}$ and $W(Y)=\{w\}$. By \ref{B9.28} the seeds are crystalline with the jumps of $U$ at the induced embeddings. Fix the local data, bridge and $\sigma$, and put $k_j=-j\deg\sigma$ and $u_j=\iota^{-1}(c_{w,k_j})^{-1}$. The unramified Weil formula in \ref{B9.28} identifies the transported Weil trace at $\sigma^j$ with $u_j$ times the augmented seed coordinate $(w,k_j)$. The same multiplication gives $S_j$ from the target coordinate. For $0\le j\le N$, use \ref{B9.31} on all coordinates at tolerance $\epsilon/(1+\sum_{j=0}^N|u_j|_\ell)$. The ordinary coordinates give \eqref{eq:p9c1-global-trace-bound}, and testing $X_{(w,k_j)}$ and multiplying by $u_j$ gives error at most $\epsilon$ at every tested power with these same weights. Sum the weights of indices that give the same representation. The joint approximation satisfies \ref{B9.36}, so the required crystalline and de Rham conclusion follows.
\leanref{UpperTheory.\allowbreak{}Gal.\allowbreak{}Lim.\allowbreak{}joint}{Development.\allowbreak{}AutomorphicArithmetic.\allowbreak{}Foundations.\allowbreak{}GaloisRepresentations.\allowbreak{}TraceLimits.\allowbreak{}SlopeTestRemainingChoices}{50--90}
\end{proof}

\subsection*{Proof of \ref{B9.35aa}}\label{prf:B9.35aa}\phantomsection
\begin{proof}
Let \(\mathfrak M\) denote the value of the crystal of \(\Lambda\) on \((\mathfrak S,(E))\). Its value on the Fontaine prism is \(A_{\inf}\otimes_{\mathfrak S}\mathfrak M\). In the chosen basis the semilinear action on this module has matrix \(\varepsilon_g\), since the images of the basis columns under the action form \(\rho(g)g(T)=T\varepsilon_g\).

Take \(\epsilon=(1,\zeta_\ell,\zeta_{\ell^2},\ldots)\) with compatible primitive roots, and put \(\mu=[\epsilon]-1\). For weights in \([0,h]\) in the convention of Gee and Liu, their crystalline lattice theorem gives
\[
 (g-1)\mathfrak M\ \subset\
 u\phi^{-1}(\mu)(A_{\inf}\otimes_{\mathfrak S}\mathfrak M)
 \qquad(g\in G_K).
\]
This is the crystalline implication of \cite[Definition~F.9, equation~(F.10), and Theorem~F.11]{gee-liu-2022}. The theorem supplies the descents for every choice of uniformizer roots for the given crystalline lattice. We take the descent for \(\pi^\flat\). Its embedding \(u\mapsto[\pi^\flat]\) is the one above, with no Frobenius twist. The identification with \(\mathfrak M\) respects the étale comparison by \cite[Theorem~7.2 and Remark~7.11]{bhatt-scholze-2023}.

For arbitrary weights choose a Tate line \(D\) of sufficiently large nonnegative weight that \(D\) and \(\Lambda\otimes D\) both have weights in \([0,h]\) for some \(h\). Apply the displayed inclusion to both lattices. In a tensor product basis, the matrix for \(\Lambda\) is the matrix for \(\Lambda\otimes D\) times the inverse of the scalar matrix coefficient for \(D\), by tensor compatibility of the crystalline realization. Both factors are the identity modulo \(u\). A change of \(\mathfrak S\)-basis preserves this congruence since \(g\) fixes \(W(k)\) and \(g(u)\in uA_{\inf}\). This proves the assertion about \(\varepsilon_g\).

For the second assertion put \(R_j=\mathfrak S[1/\ell]/(E^j)\). We describe first the rational comparison used to obtain a matrix over \(R_j\). For the Frobenius pullback \(\phi^*\mathfrak M\), the comparison near \(E=0\) identifies its rational basis space with the scalar extension of its constant isocrystal, the Frobenius pullback of \((\mathfrak M/u\mathfrak M)[1/\ell]\). This is the isomorphism lifting the identity at \(u=0\) in the proof of \cite[Theorem~7.9]{bhatt-scholze-2023}. In detail, in Construction~6.5 of that paper the rational crystalline module on the disc \(|E|\le|\ell|\) is obtained by modifying a constant module along \(E=0\). Frobenius pullback moves the modification to \(\phi(E)=0\). At \(u=\pi\) the value of \(\phi(E)\) is nonzero: the roots of the coefficient-Frobenius conjugate of \(E\) have valuation \(v_\ell(\pi)\), whereas the argument \(\pi^\ell\) has valuation \(\ell v_\ell(\pi)\). Thus the comparison and its inverse are regular on every \(R_j\). Restriction of disc functions to \(R_j\) is defined since \(E/\ell\) is nilpotent there. Under evaluation in \(B_{\rm dR}^+/(E^j)\) the constant basis vectors are fixed by \(G_K\), by the compatibility of the crystalline comparison with this lattice \cite[Proposition~5.3 and Remark~5.4]{bhatt-scholze-2023}.

Consequently their coordinate matrix \(H\) in the Frobenius pullback basis belongs to \(\GL_d(R_j)\) and satisfies \(\phi(\varepsilon_g)g(H)=H\) after evaluation modulo \(E^j\). Its reduction \(H_0\) is invertible over \(K\). The residue map \(R_j\to K\) has a \(W(k)[1/\ell]\)-algebra section: lift \(\pi\) to the unique root of \(E\) congruent to \(u\) modulo the nilpotent ideal \((E)\), using \(E'(\pi)\ne0\). After evaluation the image of this section is fixed by \(G_K\), by the same uniqueness. Multiply \(H\) on the right by the section of \(H_0^{-1}\). It now reduces to \(I_d\) and satisfies the same invariance equation.

Choose \(k_j\) clearing the denominators of this finite matrix over \(R_j\) and lift the result to \(B\) over \(\mathfrak S\). Reduction in \(\mathfrak S/(E)=\mathcal O_K\) gives \(B\equiv\ell^{k_j}I_d\). Each entry of its invariance defect belongs to \(A_{\inf}\) and to \(E^jB_{\rm dR}^+\). Their intersection is \(E^jA_{\inf}\): for \(j=1\) apply \(\theta\), whose kernel on \(A_{\inf}\) is \((E)\), then divide by \(E\) and iterate. This proves both congruences.
\end{proof}

\subsection*{Proof of \ref{B9.35ab}}\label{prf:B9.35ab}\phantomsection
\begin{proof}
We first choose one \(g\in G_K\) such that \(\lambda=\theta(g(E)/E)\) has infinite order. The quotient \(g(E)/E\) is a unit because \(g\) preserves \(\ker\theta\). Use \(\epsilon,\mu\) from the proof of Lemma~\ref{B9.35aa}. We have
\(\theta(\mu/E)\ne0\): the element \(\mu/\phi^{-1}(\mu)\) also generates \(\ker\theta\), and \(\theta(\phi^{-1}(\mu))=\zeta_\ell-1\). Normalize \(v_\ell(\ell)=1\) and put
\[
 A_0=\theta(\mu/E)\,\pi E'(\pi)\ne0.
\]
Choose \(b\ge0\) with \(v_\ell(\ell^b A_0)>1/(\ell-1)\). There exists \(g\in G_K\) with \(g(u)=[\epsilon]^{\ell^b}u\). Indeed, at level \(n\) the requested image of \(\pi^{1/\ell^n}\) is \(\zeta_{\ell^n}^{\ell^b}\pi^{1/\ell^n}\). It is a conjugate over \(K\) by the Eisenstein polynomial \(Y^{\ell^n}-\pi\). The conditions on \(g\) are nested nonempty closed subsets of the compact group \(G_K\). For a common \(g\), Taylor expansion modulo \(E^2\) gives
\[
 g(E)\equiv E+\ell^b\mu u E'(u)\pmod{E^2},
 \qquad \lambda=1+\ell^b A_0 .
\]
The \(\ell\)-adic logarithm is injective on \(v_\ell(x-1)>1/(\ell-1)\). Thus \(\lambda^i-1\ne0\) for every \(i\ge1\).

For \(1\le i<J\) choose the integer \(a_i=1+\lfloor v_\ell(\lambda^i-1)\rfloor\). Set \(s_1=0\), \(s_{i+1}=s_i+a_i\), and \(c_J=s_J\). If \(k<c_J\), multiply \(B\) by \(\ell^{c_J-k}\). Suppose henceforth that \(k\ge c_J\). Write \(D_h(U)=P_hh(U)-U\) after evaluating entries in \(A_{\inf}\). Reduction of \(D_h(B)\) modulo \(E\) gives \(\ell^k(\theta(P_h)-I_d)=0\), so \(\theta(P_h)=I_d\).

We prove successively that there exist \(C_i\in M_d(\mathfrak S)\) with
\[
 B\equiv\ell^{k-s_i}C_i\pmod{E^i}\qquad(1\le i\le J).
\]
Take \(C_1=I_d\). For \(i<J\), put \(n=k-s_i\) and write \(B=\ell^nC_i+E^i\Delta\) over \(\mathfrak S\). Put \(w=g(E)/E\). Then
\[
 D_g(B)=\ell^nD_g(C_i)+
 E^i\bigl(w^iP_g g(\Delta)-\Delta\bigr).
\]
The quotient \(A_{\inf}/(E^i)\) has no \(\ell\)-torsion. For \(i=1\) it is \(\mathcal O_C\), and the successive quotients \((E^t)/(E^{t+1})\) are copies of \(\mathcal O_C\) because \(E\) is a nonzero divisor. The displayed identity therefore gives \(D_g(C_i)=E^i U\) for an integral matrix \(U\). Divide the identity by \(E^i\) and apply \(\theta\). Its left side vanishes since \(i<J\). The entries of \(\theta(\Delta)\) lie in \(\mathcal O_K\) and are fixed by \(g\). Hence
\[
 (\lambda^i-1)\theta(\Delta)=-\ell^n\theta(U).
\]
Since \(n\ge a_i\), every entry of \(\theta(\Delta)\) lies in
\(\ell^{n-a_i}\mathcal O_C\cap K=\ell^{n-a_i}\mathcal O_K\).
Using \(\mathfrak S/(E)=\mathcal O_K\), choose \(D\in M_d(\mathfrak S)\) with
\(\Delta\equiv\ell^{n-a_i}D\bmod E\). Then
\[
 C_{i+1}=\ell^{a_i}C_i+E^iD
\]
has the required property. This proves the induction, including \(J=1\) with \(c_1=0\).

At \(i=J\) we have \(B\equiv\ell^{k-c_J}C_J\bmod E^J\). Reducing modulo \(E\) and cancelling \(\ell^{k-c_J}\) in \(\mathcal O_K\) gives \(C_J\equiv\ell^{c_J}I_d\bmod E\). For every \(h\in G_K\), the ideal \((E^J)A_{\inf}\) is \(h\)-stable. Consequently \(\ell^{k-c_J}D_h(C_J)\equiv D_h(B)\equiv0\bmod E^J\). The same torsion-freeness gives \(D_h(C_J)\equiv0\bmod E^J\), as required.
\end{proof}

\subsection*{Proof of \ref{B9.35a}}\label{prf:B9.35a}\phantomsection
\begin{proof}
Choose finite coefficient fields for the continuous representations under consideration. The preimages of the matrix groups over the countably many finite local extensions are closed subgroups and cover $G_E$. Baire's theorem gives one open preimage, and finitely many coset representatives enlarge its coefficient field to one for the full representation. Enlarge these fields to contain the images of all embeddings of the local field in the chosen bridge. For $r$ take the lattice $\Lambda_0$ spanned over $\Z_\ell$ by $r(g)$, with the left multiplication action. It is finite free in its $\Q_\ell$-span by compactness and contains the identity. The trace pairing in \ref{B9.34} shows that finitely many functions $x\mapsto\operatorname{tr}(x r(g_j))$ separate points of this span. On the unit sphere for an integral basis of $\Lambda_0$ their maximum norm is bounded below by a positive number by compactness. Thus for some integers $c,K_0\ge0$ and any $m\ge1$,
\begin{equation}\label{eq:p9c1-trace-loss}
 \max_j|\operatorname{tr}(x r(g_j))|_\ell\le\ell^{-(m+K_0+c)}
 \ \Longrightarrow\ x\in\ell^{m+K_0}\Lambda_0,\qquad
 |\operatorname{tr}(x)|_\ell\le\ell^{K_0}\quad(x\in\Lambda_0),
\end{equation}
where the implication applies to $x$ in the span and the last bound to $x$ in the lattice.

At level $m$ choose an approximation at precision $\ell^{-(m+K_0+c)}$ as in \ref{B9.36}. Keep all earlier seeds, with weight zero when unused. Write them as $\rho_k$, with weights $b_k$, and let $\Lambda_{\rm tup}$ be the span of their tuples $(\rho_k(g))_k$ over $\Z_\ell$. Let $\Lambda_{\rm full}$ be the product of full stable matrix lattices over the coefficient fields and put $\Lambda_m=\Lambda_{\rm tup}+\ell^{b}\Lambda_{\rm full}$. One can take the matrix lattices from stable lattices for the seed actions, obtained by spanning the translates of an initial lattice. Choose $b$ so large that the intersection of the added part with $\Lambda_{\rm tup}$ lies in $\ell^{m+K_0}\Lambda_{\rm tup}$, by elementary divisors. Also choose it so that the weighted trace on the added part has norm at most $\ell^{-m}$, which is possible because the list of weights is finite.

For a tuple relation $\sum_g z_g(\rho_k(g))_k=0$ with $z_g\in\Z_\ell$, test its translates by $g_j$ against the uniform trace bound. The corresponding $\sum_g z_g r(g)$ satisfies \eqref{eq:p9c1-trace-loss}. Hence evaluation at $r$ on the tuple part and zero on the added part define a map modulo $\ell^{m+K_0}$. Lift its matrix to $\Z_\ell$ in bases of $\Lambda_m$ and $\Lambda_0$. $\Pi_m$ is surjective modulo $\ell^m$. It intertwines the actions modulo $\ell^m$ and preserves the identity column modulo $\ell^m$. Write $t_i$ for the trace of the $i$-th basis matrix of $\Lambda_0$ and $t'_{hk}$ for the trace in seed component $k$ of the $h$-th basis vector of $\Lambda_m$. Testing the tuple generators and using the bounds on the added part and on the lift in \eqref{eq:p9c1-trace-loss} gives
\begin{equation}\label{eq:p9c1-layer-traces}
 \left|\sum_i(\Pi_m)_{ih}t_i-\sum_k b_k t'_{hk}\right|_\ell\le\ell^{-m},
 \qquad
 \left|\vartheta-\sum_k b_k p_k\right|_\ell\le\ell^{-m}.
\end{equation}
In the second estimate, $\vartheta=S_1$ is the target at the fixed Weil element and $p_k$ is its trace on seed $k$ transported by $\iota^{-1}$. For $\sigma^j$ use $\vartheta=S_j$. Choose the lattice exponents $b$ nondecreasing so that forgetting the new seed components gives an integral equivariant transition $f_m:\Lambda_{m+1}\to\Lambda_m$. Evaluation on tuples and vanishing on the added parts give $\Pi_m f_m\equiv\Pi_{m+1}\pmod{\ell^m}$. Each rational layer is the full product of the matrix spaces with the left actions. Thus it consists of scalar restrictions of copies of the crystalline seeds. At the limit the span of $r(g)$ over a coefficient field containing its entries is a faithful left matrix module for $r$ and contains every irreducible constituent of $r$.

Let $K=E_w$ in the chosen local realization, with residue field $k$, and put $\mathfrak S=W(k)[[u]]$ and $\mathcal O_{\mathcal E}=\widehat{\mathfrak S[1/u]}$ with $\ell$-adic completion. Use $A_{\inf}=W(\mathcal O_{\mathbf C_\ell^\flat})$ and $\mathcal W=W(\mathbf C_\ell^\flat)$. Evaluate $u$ at the Teichmüller lift of compatible roots of a uniformizer and use $u$ also for this image. Take an Eisenstein polynomial $E(u)$ with a nonzero coefficient of degree prime to $\ell$. Translating a uniformizer by an element of $\ell W(k)$ can give a nonzero linear coefficient, since the derivative of its polynomial has finitely many roots. Write $\phi$ for Witt Frobenius and its compatible actions, with $\phi(u)=u^\ell$. Use the Breuil--Kisin period $\tau$ with $\phi(\tau)=E(u)\tau$ and $\phi(\tau)$ a unit times $[\varepsilon]-1$, for a compatible primitive system of $\ell$-power roots of unity. Then $\tau$ is a unit in $\mathcal W$, its image under $\theta:A_{\inf}\to\mathbf C_\ell$ is nonzero, and the image of $\phi(\tau)$ in $B_{\rm cris}$ is invertible.

Take $N=\max(\{0\}\cup\{|j|: j\in J_\eta,\ \eta\text{ induces }w\})$. Each rational layer is crystalline with jumps in $[-N,N]$, since scalar restriction collects the jumps over the coefficient embeddings. Use the integral crystalline realization on the stable lattices, in the form of Breuil--Kisin modules and their Galois comparison \cite{kisin-2006,bhatt-scholze-2023}. The prismatic realization of stable lattices is used here on the Breuil--Kisin prism. In a basis write $T_m$ for the étale period matrix, a unit over $\mathcal W$. Its Frobenius and inverse Frobenius numerators satisfy
\begin{equation}\label{eq:p9c1-bk-height}
 E^N\phi(T_m)=T_mA_m,\qquad A_m A'_m=E^{2N}I,\qquad A_m,A'_m\in M_{d_m}(\mathfrak S),
\end{equation}
where $d_m=\operatorname{rank}_{\Z_\ell}\Lambda_m$. The elementary divisors at $(E)$ are the labeled jumps up to the sign of the filtration convention. Both the linearized Frobenius and its inverse thus become integral there on multiplication by $E^N$. Their poles lie at $E$, and divisibility by powers of the prime $E$ at $(E)$ descends to $\mathfrak S$. This gives the common height in \eqref{eq:p9c1-bk-height}.
\leanref{GlobalTheory.\allowbreak{}BKLim.\allowbreak{}exists\_\allowbreak{}frame\_\allowbreak{}sst}{Development.\allowbreak{}AbelianVarietyTheory.\allowbreak{}Proofs.\allowbreak{}InertiaEigenvalues.\allowbreak{}ExistsBreuilKisinFrameSemistableBoundedJumps}{109--135}
\leanref{GlobalTheory.\allowbreak{}BKLim.\allowbreak{}Lower.\allowbreak{}exists\_\allowbreak{}frame\_\allowbreak{}of\_\allowbreak{}jumps\_\allowbreak{}le}{Development.\allowbreak{}AbelianVarietyTheory.\allowbreak{}Proofs.\allowbreak{}InertiaEigenvalues.\allowbreak{}ExistsBreuilKisinFrameHeightOfJumpsLe}{76--90}

Put $Y_m=\tau^NT_m$ and $Y'_m=\tau^NT_m^{-1}$. By \eqref{eq:p9c1-bk-height}, $\phi(Y_m)=Y_mA_m$ and $\phi(Y'_m)=A'_mY'_m$. For a row satisfying $\phi(x)=xA+b$ with integral $A,b$, the largest norm $B$ of its entries modulo $\ell$ satisfies $B^\ell\le\max(B,1)$. Thus its first Witt digits are integral. Subtract their Teichmüller lifts and divide by $\ell$. The new equation has the same linear part and an integral inhomogeneous term, since its numerator has zero reduction. Iterate and apply the transposed calculation for columns. This gives
\begin{equation}\label{eq:p9c1-bk-poles}
 Y_m,Y'_m\text{ have entries in }A_{\inf},\qquad
 Y_mY'_m=Y'_mY_m=\tau^{2N}I .
\end{equation}
\leanref{GlobalTheory.\allowbreak{}BKLim.\allowbreak{}pre\_\allowbreak{}bounded\_\allowbreak{}poles}{Development.\allowbreak{}AbelianVarietyTheory.\allowbreak{}Proofs.\allowbreak{}InertiaEigenvalues.\allowbreak{}BreuilKisinPreframeBoundedPoles}{13--80}

Use a common étale basis $T_*$ for the target lattice. The étale comparison for the group fixing the uniformizer roots gives $T_*R_m\equiv\Pi_m T_m\pmod{\ell^m}$, with $R_m$ over $\mathcal O_{\mathcal E}$ admitting a right inverse modulo $\ell^m$ there. One construction is to quotient the étale layer module by the kernel of $\Pi_m$ modulo $\ell^m$. Its extension to $\mathcal W/\ell^m$ is free of the target rank by surjectivity. The quotient is therefore free over $\mathcal O_{\mathcal E}/\ell^m$ by elementary divisors. The consecutive quotients agree at the lower precision by $\Pi_m f_m\equiv\Pi_{m+1}$. Adjusting bases successively by lifts of the change-of-basis matrices gives compatible period matrices with invertible limit $T_*$. This uses the étale module equivalence and base change for the uniformizer tower in the Breuil--Kisin realization \cite{kisin-2006}. By functoriality of the crystalline realization, $D_m=T_m^{-1} f_m T_{m+1}$ has entries in $\mathfrak S$. Hence $R_mD_m\equiv R_{m+1}\pmod{\ell^m}$.

Fix $n\ge1$. The images $I_{n,m}=R_m\mathfrak S^{d_m}\bmod\ell^n$ for $m\ge n$ are descending. Set $a_0=\tau^{2N}$. Since $\bar a_0\ne0$ in the tilt and $0<|\bar u|<1$, choose $b_0$ with $\bar u^{b_0}/\bar a_0$ integral. Lifting this quotient gives $u^{b_0}=h a_0+\ell z$ with $h,z$ integral. In its $n$-th power, the term without $a_0$ is divisible by $\ell^n$. Thus for $c=b_0n$ one has $u^c\equiv a_0 y\pmod{\ell^n}$ for some integral $y$.

For $v\in\mathfrak S^{d_n}$ choose a right inverse $S_m$ of $\Pi_m$ over $\Z_\ell$ modulo $\ell^m$. The column $v'=yY'_m S_m\Pi_nY_nv$ is integral by \eqref{eq:p9c1-bk-poles} and satisfies
\[
 T_*R_mv'\equiv\Pi_m T_m v'\equiv
 y\tau^{2N}\Pi_nT_nv\equiv T_*R_n u^c v\pmod{\ell^n}.
\]
To descend the coefficients of $v'$, choose a topological $\mathfrak S$-basis of $A_{\inf}$ by lifting a $k$-basis of $A_{\inf}/(\ell,u)$ that contains $1$. Successive expansions modulo $(\ell,u)^j$ exist and are unique. Both rings are complete for this ideal, and $(\ell,u)$ is a regular sequence in each, since $\bar u\ne0$ in the tilt valuation ring. The associated graded rings are polynomial rings in the two initial forms. The coordinate of $1$ gives an $\mathfrak S$-linear retraction. Clear a power of $u$ from $R_mv'\equiv R_nu^c v$ at precision $n$ to put the matrix and target over $\mathfrak S/\ell^n$. Integral Witt vectors are $\ell$-saturated in $\mathcal W$, so applying the retraction preserves the congruence. Cancel the cleared power of $u$ over $\mathcal O_{\mathcal E}$. This proves
\begin{equation}\label{eq:p9c1-lattice-stabilization}
 u^c I_{n,n}\subset I_{n,m}\subset I_{n,n}\qquad(m\ge n).
\end{equation}
The quotient by $u^c I_{n,n}$ is a quotient of $(\mathfrak S/(\ell^n,u^c))^{d_n}$ and is finite. Thus the images stabilize for each $n$.
\leanref{GlobalTheory.\allowbreak{}BKLim.\allowbreak{}flatPow\_\allowbreak{}dvd\_\allowbreak{}perPow\_\allowbreak{}mod}{Development.\allowbreak{}AbelianVarietyTheory.\allowbreak{}Proofs.\allowbreak{}InertiaEigenvalues.\allowbreak{}FlatPowDvdPerPowMod}{15--25}
\leanref{GlobalTheory.\allowbreak{}CoefDescent.\allowbreak{}descent}{Development.\allowbreak{}AbelianVarietyTheory.\allowbreak{}Proofs.\allowbreak{}InertiaEigenvalues.\allowbreak{}BreuilKisinDescentCoefficientsAinfModPow}{68--120}
\leanref{GlobalTheory.\allowbreak{}BKLim.\allowbreak{}mittagLeffler\_\allowbreak{}approx}{Development.\allowbreak{}AbelianVarietyTheory.\allowbreak{}Proofs.\allowbreak{}InertiaEigenvalues.\allowbreak{}LatticeApproximationMittagLeffler}{30--157}

The stable images at consecutive precisions map onto one another, by using the same coefficient column at a layer beyond both stabilization indices. At precision one clear a power of $u$ from a right inverse of a stabilized $R_m$ over $k((u))$. Lift the resulting columns of $u^h I$ successively through the stable images, where $h$ is the clearing exponent. These lifts give an invertible matrix over $\mathcal O_{\mathcal E}$ since its reduction is invertible there. Each lifted column $v$ satisfies $\tau^NT_*v\in A_{\inf}^{d_0}$, where $d_0$ is the rank of $\Lambda_0$. At every precision its image is a column combination of $\Pi_mY_m$ with integral coefficients.

Consider the saturated lattice
\[
 \mathcal L=A_{\inf}^{d_0}\cap\tau^N T_*\mathcal O_{\mathcal E}^{d_0}
 \subset\mathcal W^{d_0}.
\]
In the coordinates of $\tau^N T_*$ its reduction is a $k[[u]]$-submodule of $k((u))^{d_0}$ with bounded denominators. The reduced inverse matrix has bounded tilt norm and the input entries have norm at most one. The lifted columns give full rank. The reduction is free of rank $d_0$ over $k[[u]]$. Lift a basis to the columns of a matrix $Y$ in $\mathcal L$. Subtracting its multiples and dividing the remainder by $\ell$ can be repeated inside $\mathcal L$ by the $\ell$-saturation of $A_{\inf}$ and $\mathcal O_{\mathcal E}$ in $\mathcal W$. Completeness gives coefficients in $\mathfrak S$, so $Y$ is a basis and $T=\tau^{-N}Y$ differs from $T_*$ by an $\mathcal O_{\mathcal E}$-unit.

At precision $n$, take $m\ge n$ beyond stabilization. Lift the columns of $R_m$ from the stable images and express their integral period vectors in the basis $Y$. The resulting matrix $R$ over $\mathfrak S$ gives $TR\equiv\Pi_m T_m\pmod{\ell^n}$. Put $Q_m^*=Y'_m S_m$. Then $YRQ_m^*\equiv\tau^{2N}I\pmod{\ell^n}$. It follows that $Y'=Y^{-1}\tau^{2N}$ in $\mathcal W$ is integral by the digit criterion at every precision. The Frobenius matrices
\[
 A=E^NT^{-1}\phi(T),\qquad A'=E^N\phi(T)^{-1}T
\]
have entries in $\mathcal O_{\mathcal E}$ by étale descent. The matrices $\tau^{2N}A=Y'\phi(Y)$ and $\tau^{2N}A'=\phi(Y')Y$ are integral.

In the topological basis used in \eqref{eq:p9c1-lattice-stabilization}, the ideal $I$ generated by the coefficients of $\tau^{2N}$ contains $\ell^r$ for some $r\ge0$. Indeed, its saturation for $\ell$ is principal, say $(g)$. Lift a generator of the reduced ideal in $k[[u]]$ and apply saturation and Nakayama's lemma. The generator is nonzero and $\ell\nmid g$. Frobenius of the expansion and $\phi(\tau^{2N})=E^{2N}\tau^{2N}$ give $\phi(g)\mid E^{2N}I$. Since $\ell^j g\in I$ for some $j$ and $\ell\nmid\phi(g)$, one has $\phi(g)\mid E^{2N}g$. If $g$ is a nonunit, its Weierstrass polynomial has roots of norm less than one. A zero root is excluded by comparing orders at $u=0$, since $E(0)\ne0$. Choose a root of largest norm and take its coefficient-Frobenius conjugate. Every $\ell$-th root of this conjugate or of its field conjugates has strictly larger norm and is a root of $E$. Their nonempty Galois-stable set must be all roots of $E$ by irreducibility. Its monic polynomial is the minimal polynomial of the conjugate evaluated at $u^\ell$, contrary to the chosen exponent of $E$. Thus $g$ is a unit.

Choose a finite $\mathfrak S$-linear combination $p$ of basis coordinates with $p(\tau^{2N})=\ell^r$. For $x\in\mathcal O_{\mathcal E}$ with $\tau^{2N}x\in A_{\inf}$, clear powers of $u$ at each precision to get $\ell^r x=p(\tau^{2N}x)\in\mathfrak S$. Since $\mathfrak S[1/\ell]\cap\mathcal O_{\mathcal E}=\mathfrak S$ by successive reduction modulo $\ell$, this descends $A$ and $A'$ to $\mathfrak S$. Also $RQ_m^*Y\equiv\tau^{2N}I$ by cancelling $Y$ over $\mathcal W$. With $G=p(Q_m^*Y)$ the links satisfy
\begin{equation}\label{eq:p9c1-frame-link}
 \forall n\ge1\ \exists m\ge n,\ R,G\text{ over }\mathfrak S,\qquad
 TR\equiv\Pi_m T_m,\qquad RG\equiv\ell^r I\pmod{\ell^n}.
\end{equation}
The exponent $r$ depends only on $\tau,N$.
\leanref{GlobalTheory.\allowbreak{}BSP.\allowbreak{}pow\_\allowbreak{}l\_\allowbreak{}mem\_\allowbreak{}span\_\allowbreak{}coeff}{Development.\allowbreak{}AbelianVarietyTheory.\allowbreak{}Proofs.\allowbreak{}InertiaEigenvalues.\allowbreak{}ExistsPowMemSpanCoeffPeriodPow}{83--160}
\leanref{GlobalTheory.\allowbreak{}LemZ.\allowbreak{}isUnit\_\allowbreak{}of\_\allowbreak{}frob\_\allowbreak{}dvd}{Development.\allowbreak{}AbelianVarietyTheory.\allowbreak{}Proofs.\allowbreak{}InertiaEigenvalues.\allowbreak{}IsUnitOfFrobeniusDvdEisensteinPowMul}{51--151}
\leanref{GlobalTheory.\allowbreak{}BKLim.\allowbreak{}descend\_\allowbreak{}rightInverse}{Development.\allowbreak{}AbelianVarietyTheory.\allowbreak{}Proofs.\allowbreak{}InertiaEigenvalues.\allowbreak{}DescendRightInverse}{12--43}

Write $\rho_0,\rho_m$ for the left actions on the limit and layer lattices. For $g\in G_K$ put $\varepsilon_g=T^{-1}\rho_0(g)g(T)$ and $\varepsilon_{m,g}=T_m^{-1}\rho_m(g)g(T_m)$. The lattice $\Lambda_m$ is stable and finite free over $\Z_\ell$. Its rational representation is a finite direct sum of scalar restrictions of crystalline seeds. Thus Lemma~\ref{B9.35aa}, with its chosen Breuil--Kisin basis and $d=d_m$, gives $\varepsilon_{m,g}-I_{d_m}\in uM_{d_m}(A_{\inf})$ for every $g\in G_K$. The equivariance of $\Pi_m$ and \eqref{eq:p9c1-frame-link} give, over $\mathcal W$,
\begin{equation}\label{eq:p9c1-cocycle-link}
 \varepsilon_g g(R)\equiv R\varepsilon_{m,g}\pmod{\ell^n}.
\end{equation}
Multiply by $g(G)$ and use $g(G)-G\in u M(A_{\inf})$, since $g$ fixes $W(k)$ and $g(u)$ is divisible by $u$. It follows that $\ell^r(\varepsilon_g-I)$ is congruent to a multiple of $u$ with integral entries at every precision. Both $A_{\inf}$ and $uA_{\inf}$ are $\ell$-adically closed and $\ell$-saturated in $\mathcal W$. For $uA_{\inf}$ use the same properties for integral Witt digits after division by the unit $u$ of $\mathcal W$. Hence $\varepsilon_g\in I+u M_{d_0}(A_{\inf})$.

A rank-$d$ frame with cocycle $\varepsilon_g$ satisfies the de Rham congruences if for every integer $j\ge0$ there are an integer $k\ge0$ and a matrix $\beta\in M_d(\mathfrak S)$ such that
\[
 \beta\equiv\ell^k I_d\pmod E,\qquad
 \phi(\varepsilon_g)\,g(\beta)\equiv\beta\pmod{E^j}\quad(g\in G_K).
\]
The first congruence is in $\mathfrak S$. In the second, evaluate $\mathfrak S$ in $A_{\inf}$ before applying $g$ and take the congruence entrywise in $A_{\inf}$.
\leanref{GlobalTheory.\allowbreak{}BKLim.\allowbreak{}Frame.\allowbreak{}IsBeta}{Development.\allowbreak{}AbelianVarietyTheory.\allowbreak{}Definitions.\allowbreak{}NonGaloisDescent.\allowbreak{}TrivIsCr}{74--76}
\leanref{GlobalTheory.\allowbreak{}BKLim.\allowbreak{}Frame.\allowbreak{}IsDR}{Development.\allowbreak{}AbelianVarietyTheory.\allowbreak{}Definitions.\allowbreak{}NonGaloisDescent.\allowbreak{}TrivIsCr}{78--78}

We prove these congruences for $T$. Fix $M\ge0$ and take $c=c_{M+1}$ from Lemma~\ref{B9.35ab} for the fixed $K,\pi^\flat,E$. For each layer, Lemma~\ref{B9.35aa} gives a solution at order $M+1$ with some exponent. Apply Lemma~\ref{B9.35ab} to that solution with $d=d_m$ and $P_g=\phi(\varepsilon_{m,g})$. This gives matrices $\beta_m$ over $\mathfrak S$ with
\[
 \beta_m\equiv\ell^c I\pmod E,\qquad
 \phi(\varepsilon_{m,g})g(\beta_m)\equiv\beta_m\pmod{E^{M+1}}
 \quad(g\in G_K).
\]
The second congruence uses the evaluation in $A_{\inf}$. Also, for every $x\in\mathfrak S$ there is a lift $\gamma(x)\in\mathfrak S$ of $\ell^{k_0}x$ modulo $E$, invariant after evaluation modulo $E^{M+1}$ for all $g\in G_K$. The exponent $k_0$ depends on $M,K,E$ and is uniform in $x$. For example, after inverting $\ell$, lift the uniformizer in $\mathfrak S/(E^{M+1})$ to a root of its Eisenstein polynomial using the invertible derivative. This gives the section from $K$ whose evaluation is fixed modulo $E^{M+1}$ in $B_{\rm dR}^+$. Clear denominators for a $W(k)$-basis of $\mathcal O_K$ to choose $k_0$. Divisibility by $E^{M+1}$ of the evaluated integral differences descends to $A_{\inf}$.

For each precision $n$, take $m,R,G$ as in \eqref{eq:p9c1-frame-link}, apply $\gamma$ entrywise to $\phi(G)$, and put $B_n=\phi(R)\beta_m\gamma(\phi(G))$. Reduction modulo $(\ell^n,E)$ gives
\[
 B_n\equiv\ell^{c+k_0}\phi(RG)\equiv\ell^{c+k_0+r}I\pmod{(\ell^n,E)} .
\]
Frobenius commutes with the Galois action after evaluation. Apply it to \eqref{eq:p9c1-cocycle-link}, which also holds modulo $\ell^n$ in $A_{\inf}$ by saturation. Then modulo $(\ell^n,E^{M+1})$ in $A_{\inf}$,
\[
 \phi(\varepsilon_g)g(B_n)
 \equiv\phi(R)\phi(\varepsilon_{m,g})g(\beta_m)\gamma(\phi(G))
 \equiv B_n
 \qquad(g\in G_K).
\]
These two tests define nested nonempty closed subsets of the matrix space over $\mathfrak S/(E^{M+1})$, since at each $n$ they depend on the class modulo $(\ell^n,E^{M+1})$. The quotient by $E^{M+1}$ is finite free over $W(k)$ by Weierstrass division, hence compact. Choose a common class and lift it to $\mathfrak S$. In $A_{\inf}$,
\[
 \bigcap_{n\ge1}(\ell^n,E^j)=E^j A_{\inf}\qquad(j\ge1).
\]
For $j=1$ this follows from $\ell$-adic separatedness after applying $\theta$. For the induction, divide by $E$ after using $E\mid\ell^n x\Rightarrow E\mid x$, which follows since $A_{\inf}/E$ has no $\ell$-torsion. Also $E A_{\inf}\cap\mathfrak S=E\mathfrak S$. The common class therefore gives a matrix congruent to $\ell^{c+k_0+r}I$ modulo $E$ with the twisted invariance modulo $E^{M+1}$. Taking $M+1\ge j$ gives the required congruences at every precision $j$.
\leanref{GlobalTheory.\allowbreak{}BKLim.\allowbreak{}limit\_\allowbreak{}pre\_\allowbreak{}exists}{Development.\allowbreak{}AbelianVarietyTheory.\allowbreak{}Proofs.\allowbreak{}InertiaEigenvalues.\allowbreak{}ExistsBreuilKisinBasisLimitLattice}{35--63}
\leanref{GlobalTheory.\allowbreak{}BKLim.\allowbreak{}limit\_\allowbreak{}frame\_\allowbreak{}exists}{Development.\allowbreak{}AbelianVarietyTheory.\allowbreak{}Proofs.\allowbreak{}InertiaEigenvalues.\allowbreak{}ExistsDeRhamFrameLimitLatticeBreuilKisinTower}{16--24}
\leanref{GlobalTheory.\allowbreak{}BKLim.\allowbreak{}isDR\_\allowbreak{}of\_\allowbreak{}links}{Development.\allowbreak{}AbelianVarietyTheory.\allowbreak{}Proofs.\allowbreak{}InertiaEigenvalues.\allowbreak{}IsDROfLinks}{107--224}
\leanref{GlobalTheory.\allowbreak{}BKLim.\allowbreak{}beta\_\allowbreak{}bdd\_\allowbreak{}denominators}{Development.\allowbreak{}AbelianVarietyTheory.\allowbreak{}Proofs.\allowbreak{}InertiaEigenvalues.\allowbreak{}BreuilKisinBetaBoundedDenominators}{41--105}
\leanref{GlobalTheory.\allowbreak{}BSP.\allowbreak{}BK.\allowbreak{}exists\_\allowbreak{}inv\_\allowbreak{}lift}{Development.\allowbreak{}AbelianVarietyTheory.\allowbreak{}Proofs.\allowbreak{}InertiaEigenvalues.\allowbreak{}ExistsInvLift}{82--116}

 In $A_{\rm cris}$ one has $E(u)\equiv u^{e_0}\pmod\ell$, with $e_0=\deg E$. The divided powers give $\ell\mid E(u)^\ell$, since $E(u)$ is in the divided-power ideal. Thus $\ell\mid u^{e_0\ell}$. Choose $n_0\ge1$ with $e_0\ell((2N+1)n+1)\le\ell^n-1$ for all $n\ge n_0$. Then
\begin{equation}\label{eq:p9c1-divided-power}
 \ell^{(2N+1)n+1}\mid u^{\ell^n-1}\qquad(n\ge n_0).
\end{equation}
Let $f_A:\mathfrak S\to A_{\rm cris}$ be evaluation followed by crystalline Frobenius, and $f_0$ the same map after taking constant terms. Write $f_A(E)=\ell v$ and $f_0(E)=\ell v_0$. For the first quotient,
\[
 v=\frac{\phi(E)}{\ell}=\delta(E)+\frac{E^\ell}{\ell}.
\]
The Eisenstein prism condition makes $\delta(E)$ a unit. The divided-power identity
\[
 (E^\ell/\ell)^\ell=\frac{(\ell^2)!}{\ell^\ell}\gamma_{\ell^2}(E)
\]
puts this power in $\ell A_{\rm cris}$. Hence $v$ is a unit modulo $\ell$, and $\ell$-adic completeness gives its inverse in $A_{\rm cris}$. The quotient $v_0$ is the Frobenius of the unit $E(0)/\ell$. For each frame, with matrices $A,A'$ as above, set
\[
 \alpha=v^{-N}f_A(A),\quad \alpha'=v^{-N}f_A(A'),\qquad
 \beta=v_0^{-N}f_0(A),\quad\beta'=v_0^{-N}f_0(A').
\]
Thus $\alpha\alpha'=\beta\beta'=\ell^{2N}I$. The differences $f_A(A)-f_0(A)$ and $\ell(v-v_0)$ are divisible by $u^\ell$, since the corresponding differences before Frobenius are divisible by $u$. Expanding the difference of the unit powers gives $\ell(\alpha-\beta)\in u^\ell M(A_{\rm cris})$ and likewise for $\alpha',\beta'$. Starting with $P_0=P'_0=I$, iterate
\[
 P_{j+1}=\alpha\phi(P_j)\beta',\qquad
 P'_{j+1}=\beta\phi(P'_j)\alpha'.
\]
Induction gives $P_jP'_j=P'_jP_j=\ell^{4Nj}I$ and
\begin{equation}\label{eq:p9c1-p-recurrence}
 \ell(P_{j+1}-\ell^{2N}P_j)\in u^{\ell^{j+1}}M(A_{\rm cris}),
 \qquad
 \ell(P'_{j+1}-\ell^{2N}P'_j)\in u^{\ell^{j+1}}M(A_{\rm cris}).
\end{equation}
For the first recurrence start with $\ell(\alpha-\beta)\beta'$, then apply $\phi$ and the same two end multipliers at each step. The second starts with $\ell(\beta-\alpha)\alpha'$. By \eqref{eq:p9c1-divided-power}, $\ell^{-2Ni}P_{n_0+i}$ and its primed version are integral and converge $\ell$-adically. Multiply both limits by $\ell^{n_0}$ to obtain $Q,Q'$ with
\begin{equation}\label{eq:p9c1-constant-frame}
 QQ'=Q'Q=\ell^{2k_0}I,\qquad Q\beta=\alpha\phi(Q),\qquad
 Q,Q'\equiv\ell^{k_0}I\pmod u,\qquad k_0=(2N+1)n_0 .
\end{equation}
For the congruences modulo $u$, the scaled differences in \eqref{eq:p9c1-p-recurrence} are divisible by $u$ with quotients tending to zero by \eqref{eq:p9c1-divided-power}. The initial term $\ell^{n_0}P_{n_0}$ is congruent to $\ell^{k_0}I$ by the finite telescoping sum of \eqref{eq:p9c1-p-recurrence}, and likewise for the primed term. The Frobenius relation follows from the recurrence by passage to the limit and cancellation of $\ell^{2N}$.
\leanref{GlobalTheory.\allowbreak{}AcrisFrobE.\allowbreak{}main}{Development.\allowbreak{}AbelianVarietyTheory.\allowbreak{}Proofs.\allowbreak{}InertiaEigenvalues.\allowbreak{}AcrisTorsionFreeCompleteFrobeniusProperties}{94--143}

Write $\epsilon_g'=\phi(\varepsilon_g)$ in $A_{\rm cris}$. Compatibility of Galois and Frobenius on $T$ gives $\epsilon_g' g(\alpha)=\alpha\phi(\epsilon_g')$. The constants $\beta,\beta'$ are fixed by $G_K$. For $D_g=\epsilon_g' g(Q)-Q$, \eqref{eq:p9c1-constant-frame} therefore gives $D_g\beta=\alpha\phi(D_g)$ and $\ell^{2N}D_g=\alpha\phi(D_g)\beta'$. The cocycle condition and \eqref{eq:p9c1-constant-frame} give $D_g\in uM(A_{\rm cris})$. Iteration puts $\ell^{2Nj}D_g$ in $u^{\ell^j}M(A_{\rm cris})$. By \eqref{eq:p9c1-divided-power}, cancellation of $\ell^{2Nj}$ makes $D_g$ divisible by arbitrarily high powers of $\ell$. Separatedness gives $D_g=0$.
\leanref{UpperTheory.\allowbreak{}Gal.\allowbreak{}ClosedPhiFromClosedTrace.\allowbreak{}invQ\_\allowbreak{}holds}{Development.\allowbreak{}AutomorphicArithmetic.\allowbreak{}Foundations.\allowbreak{}GaloisRepresentations.\allowbreak{}ClosedPhiFunctoriality.\allowbreak{}InvQHolds}{16--122}

For each frame use $Y,Y'$ for its integral matrices in \eqref{eq:p9c1-bk-poles} and the limit construction, and put
\[
 t=\phi(\tau),\qquad Z=\phi(Y)Q,\qquad Z'=Q'\phi(Y'),\qquad
 b=t^{2N}\ell^{2k_0},\qquad U=t^{-N}Z .
\]
Then $ZZ'=Z'Z=bI$. Since $t$ and $\ell$ are invertible in $B_{\rm cris}$, $U$ is a basis there. The equation $D_g=0$ gives $\rho_0(g)g(U)=U$ in the limit and the corresponding equation for each layer. Hence the rational matrix lattice of the limit is crystalline. Its left action contains every irreducible constituent of $r$ after extending coefficients, so $r_w$ is crystalline by stability under scalar extension and subobjects, and hence de Rham.

By \eqref{eq:p9c1-frame-link} and \eqref{eq:p9c1-bk-height}, the Frobenius matrices satisfy
\[
 A\phi(R)\equiv R A_m,\qquad A'R\equiv\phi(R)A'_m\pmod{\ell^n}.
\]
For the second congruence multiply the first by the inverse numerators and cancel $E^{2N}$ in $\mathcal W$, where $E$ has nonzero reduction and is a unit. The congruences descend to $\mathfrak S$ by its $\ell$-saturation in $\mathcal W$. Put $H=f_A(R)$ and $H_0=f_0(R)$. The two congruences give $\alpha\phi(H)\equiv H\alpha_m$ and $\beta'H_0\equiv\phi(H_0)\beta'_m$. Consequently $D_j=H P_{m,j}-P_jH_0$ satisfies $D_{j+1}\equiv\alpha\phi(D_j)\beta'_m\pmod{\ell^n}$, with $D_0=H-H_0$ divisible by $u^\ell$. Thus $D_j$ is divisible by $u^{\ell^{j+1}}$ modulo $\ell^n$. Given $L$, choose $j$ large enough that the scaled finite iterates approximate $Q,Q_m$ modulo $\ell^L$ by \eqref{eq:p9c1-p-recurrence} and \eqref{eq:p9c1-divided-power}. Take $n$ at least $L$ plus the denominator exponent of this scaling and use \eqref{eq:p9c1-frame-link} at that precision. The bound on $D_j$ gives $H Q_m\equiv QH_0\pmod{\ell^L}$. Applying $\phi$ to $\Pi_mY_m\equiv YR$ then gives
\begin{equation}\label{eq:p9c1-frame-Z-link}
 \Pi_m Z_m\equiv ZH_0\pmod{\ell^L}.
\end{equation}
Here $m\ge L$ and $H_0$ has coefficients in $W(k)$.
\leanref{UpperTheory.\allowbreak{}Gal.\allowbreak{}ClosedPhiFromClosedTrace.\allowbreak{}linkZ\_\allowbreak{}holds}{Development.\allowbreak{}AutomorphicArithmetic.\allowbreak{}Foundations.\allowbreak{}GaloisRepresentations.\allowbreak{}ClosedPhiFunctoriality.\allowbreak{}LinkZHolds}{10--41}

Let $n_1,n_2$ be the negative and positive parts of $[k:\F_\ell]\deg\sigma$. On an unramified crystalline parameter the Weil operator is the action of $\phi^{n_1-n_2}$ on crystalline invariants, with the degree convention of \ref{B9.06}. For each frame define
\begin{equation}\label{eq:p9c1-denominator}
 V=\phi^{n_2}(b)\phi^{n_1}(Z)\phi^{n_2}(Z'),\qquad
 s=\phi^{n_2}(t)^{3N}\ell^{4k_0}\phi^{n_1}(t)^N .
\end{equation}
The identity $Z'Z=bI$ gives $V\phi^{n_2}(U)=s\phi^{n_1}(U)$. This represents the required operator on the invariant columns after base extension. The power $\phi^{n_2}$ is invertible on the invariant module and both powers fix its unramified coefficient field. Each factor of $s$ is invertible in $B_{\rm cris}$ and all layers use the same $s$.

Both powers of Frobenius fix $H_0$ and $\Pi_m$. Apply them to \eqref{eq:p9c1-frame-Z-link}. Multiply the congruence for $n_2$ on the left by $\phi^{n_2}(Z')$ and on the right by $\phi^{n_2}(Z'_m)$ to get
\[
 \phi^{n_2}(b)\phi^{n_2}(Z')\Pi_m
 \equiv \phi^{n_2}(b)H_0\phi^{n_2}(Z'_m)\pmod{\ell^L}.
\]
Use the congruence for $n_1$ on the right after left multiplication by $\phi^{n_1}(Z)$ to obtain $V\Pi_m\equiv\Pi_m V_m$. The identity column congruence holds by construction of $\Pi_m$. This proves \eqref{eq:p9c1-weil-intertwine}.
\end{proof}

\subsection*{Proof of \ref{B9.36}}\label{prf:B9.36}\phantomsection
\begin{proof}
Use the frames and comparison of \ref{B9.35a}. Fix an induced embedding $\eta$ and transport the coefficient traces to $\mathbf C_\ell$ by an embedding that identifies this copy of $E_w$ with the chosen local copy $K$. This transport preserves the absolute value. Continue to write $\theta$ for reduction from $B_{\rm dR}^+$ and $E$ for the evaluated Eisenstein polynomial, which generates $\ker(\theta|_{A_{\inf}})$.

Take the element $g\in G_K$ constructed in the proof of Lemma~\ref{B9.35ab} for this uniformizer and its compatible roots. Put $w_g=g(E)/E$ and $z=\theta(w_g)^{-1}$. Then $z$ has infinite order. The same $g,z$ apply to the limit and every layer.
\leanref{GlobalTheory.\allowbreak{}BKLim.\allowbreak{}IsZ}{Development.\allowbreak{}AbelianVarietyTheory.\allowbreak{}Definitions.\allowbreak{}NonGaloisDescent.\allowbreak{}TrivIsCr}{182--183}
\leanref{GlobalTheory.\allowbreak{}BKLim.\allowbreak{}exists\_\allowbreak{}isZ\_\allowbreak{}not\_\allowbreak{}root}{Development.\allowbreak{}AbelianVarietyTheory.\allowbreak{}Proofs.\allowbreak{}InertiaEigenvalues.\allowbreak{}ExistsGaloisElemCyclotomicPeriodNotRootOfUnity}{108--162}

For any of these frames, write $T,Y,\varepsilon_g,\rho$ for its period matrix, integral numerator, cocycle and lattice action. Choose an adapted de Rham basis with invariant matrix $V\operatorname{diag}(E^{p_i-q})$ in $B_{\rm dR}$, where $V\in\GL_d(B_{\rm dR}^+)$ and $p_i-q$ are the jumps of this lattice representation over $K$. Here $d$ is its rank and $p_i,q$ are nonnegative integers. Put $V_0=\tau^N Y^{-1}V$ over $B_{\rm dR}^+$. Its reduction is fixed entrywise by $G_K$. To see this, the de Rham comparison for the frame gives an invariant basis $(E\tau)^{-N}YAB$ with $B\in\GL_d(B_{\rm dR}^+)$ admitting $\mathfrak S$-approximations modulo each power of $E$ after clearing a power of $\ell$. For each precision take a matrix $\beta$ from the de Rham congruences in the proof of Lemma~\ref{B9.35a} and divide it by its normalizing factor $\ell^k$. The resulting matrices reduce to the identity modulo $E$. Their classes agree at every common precision, since the leading coefficient of a difference of order $j\ge1$ would lie in $K$ and be fixed by multiplication by $\theta(w_g)^j\ne1$. Their limit gives $B$, and the Frobenius cocycle relation gives invariance of $(E\tau)^{-N}YAB$. The two invariant bases thus differ by a matrix $D\in\GL_d(K)$, so
\[
 E^q A B D=E^N V_0\operatorname{diag}(E^{p_i}).
\]
For each column $i$, approximate the left side modulo $E^{N+p_i+1}$ by $\mathfrak S$ after multiplying by a power of $\ell$. For the entries of $D$, use the section of $K$ modulo powers of $E$ from the invariant lift construction in \ref{B9.35a}. The $\mathfrak S$ approximant is divisible by $E^{N+p_i}$ by Weierstrass division. Divide and reduce by $\theta$. This puts the entries of $\theta(V_0)$ in $K$.
\leanref{GlobalTheory.\allowbreak{}BKLim.\allowbreak{}InvBasis.\allowbreak{}exists\_\allowbreak{}invariant\_\allowbreak{}basis}{Development.\allowbreak{}AbelianVarietyTheory.\allowbreak{}Proofs.\allowbreak{}InertiaEigenvalues.\allowbreak{}ExistsInvariantBasisDeRhamBreuilKisin}{88--158}

Put $y=\theta(Y)$, $v=\theta(V)$, $v_0=\theta(V_0)$, $a_\tau=\theta(\tau)$ and $b_\tau=\theta(g(\tau))$. Use $g(y)$ and $g(v)$ for the entrywise action on these reductions. From $YV_0=\tau^N V$ and its Galois conjugate,
\[
 yv_0=a_\tau^N v,\qquad g(y)v_0=b_\tau^N g(v).
\]
The cocycle identity on $T=\tau^{-N}Y$ gives
\[
 a_\tau^N\rho(g)g(y)=b_\tau^N y\theta(\varepsilon_g).
\]
The invariance of $V\operatorname{diag}(E^{p_i-q})$ gives, after cancellation of the column powers of $E$ and reduction, $\rho(g)g(v)=v\operatorname{diag}(z^{p_i-q})$. Multiply the cocycle identity on the right by $v_0$ and substitute $g(y)v_0=b_\tau^N g(v)$. Cancel the nonzero factor $b_\tau^N$. The remaining equality, together with $yv_0=a_\tau^N v$, gives
\[
 U y=y\theta(\varepsilon_g),\qquad
 Uv=v\operatorname{diag}(z^{p_i-q}),\qquad
 U=y\theta(\varepsilon_g)y^{-1}.
\]
\leanref{GlobalTheory.\allowbreak{}BKLim.\allowbreak{}adaptedBasis\_\allowbreak{}theta\_\allowbreak{}and\_\allowbreak{}poles}{Development.\allowbreak{}AbelianVarietyTheory.\allowbreak{}Proofs.\allowbreak{}InertiaEigenvalues.\allowbreak{}AdaptedBasisThetaGaloisFixedPoles}{33--132}

Write $U_0,U_m$ for these operators on the limit and the layers, and set $d_*=|\theta(\tau^{2N})|_\ell^{-1}\ge1$. Use the entrywise maximum norm, which is submultiplicative over an ultrametric field. By \eqref{eq:p9c1-bk-poles} and the limit construction the norm of each reduced inverse of $Y$ is at most $d_*$. Each reduced $Y$ and cocycle has norm at most one. Put $y_0=\theta(Y)$ for the limit, $y_m=\theta(Y_m)$, and $R_\theta=\theta(R)$ for the link \eqref{eq:p9c1-frame-link}. Both
\[
 \Delta=\Pi_m y_m-y_0 R_\theta,\qquad
 \Delta'=\theta(\varepsilon_g)R_\theta-R_\theta\theta(\varepsilon_{m,g})
\]
have norm at most $\ell^{-n}$, by the link and \eqref{eq:p9c1-cocycle-link}, since $\theta(g(R))=\theta(R)$. The identity
\[
 (U_0\Pi_m-\Pi_m U_m)y_m
 =U_0\Delta+y_0\Delta'-\Delta\theta(\varepsilon_{m,g})
\]
and the inverse bound give, for the linked $m\ge n$,
\begin{equation}\label{eq:p9c1-hodge-bound}
 \|U_0\|,\ \|U_m\|\le d_*,
 \qquad \|U_0\Pi_m-\Pi_m U_m\|\le d_*^2\ell^{-n}.
\end{equation}
\leanref{GlobalTheory.\allowbreak{}HTBound.\allowbreak{}bounded\_\allowbreak{}compat}{Development.\allowbreak{}AbelianVarietyTheory.\allowbreak{}Proofs.\allowbreak{}InertiaEigenvalues.\allowbreak{}ThetaConjBoundedCompatibleBreuilKisinFrames}{117--228}

Telescoping gives $\|U_0^j\Pi_m-\Pi_m U_m^j\|\le d_*^{j+1}\ell^{-n}$ for $j\ge1$. The matrix readings into the seed and limit matrix spaces are equivariant. After coefficient extension, functoriality of the filtered comparison sends each Hodge--Tate summand to the summand of the same jump. The operator on that summand is multiplication by $z^h$ for jump $h$, so the readings intertwine the operators. Evaluation at $e_m$ and then taking the trace in seed component $k$ thus gives $\sum_{h\in J_\eta}z^{jh}$, and the corresponding evaluation at $e$ gives the power sum for the actual jumps of $r_w$ at $\eta$. Use $\Pi_m e_m\equiv e$, \eqref{eq:p9c1-hodge-bound}, and the first bound of \eqref{eq:p9c1-layer-traces} on the columns of $U_m^j e_m$. For fixed $j$, the error from the projected column is bounded by $d_*^{j+1}\ell^{-n}\max_i|t_i|_\ell$, and the weighted trace error on the layer column is at most $d_*^j\ell^{-m}$. At the identity the first inequality of \eqref{eq:p9c1-joint} gives $|a(1-\sum_k b_k)|_\ell\le\ell^{-m}$. Since $a>0$, the sum of the weights tends to one. Thus the power sums on $r$ are $\sum_{h\in J_\eta}z^{jh}$. Newton's identities and the injectivity of $h\mapsto z^h$ recover the multiset $J_\eta$ with multiplicities.
\leanref{GlobalTheory.\allowbreak{}CovClose.\allowbreak{}closed\_\allowbreak{}of\_\allowbreak{}cover}{Development.\allowbreak{}AbelianVarietyTheory.\allowbreak{}Proofs.\allowbreak{}InertiaEigenvalues.\allowbreak{}ClosedOfCover}{159--301}

Fix a power $\sigma^j$. Apply \ref{B9.35a} to it with target $\vartheta=S_j$ and write $p$ for its transported trace on $r_w$ and $p_k$ for that on seed $k$. Let $K^{\rm ur}\subset\bar{\Q}_\ell$ be the maximal unramified algebraic extension of $\Q_\ell$, with period embedding $u_{\rm ur}:K^{\rm ur}\to B_{\rm cris}$ from the compatible data. For a $K^{\rm ur}$-linear functional $\lambda:\bar{\Q}_\ell\to K^{\rm ur}$, the comparison \eqref{eq:p9c1-denominator} gives
\begin{equation}\label{eq:p9c1-scalar-period-trace}
 \sum_i(Ve)_i\,u_{\rm ur}(\lambda(t_i))=s\,u_{\rm ur}(\lambda(p)).
\end{equation}
After coefficient extension, the matrix reading intertwines the period operators. Apply the operator to the identity matrix column and take its ordinary matrix trace. This gives the Weil trace, since the reading space has the left multiplication action. Applying the coefficient functional $\lambda$ and the period embedding gives \eqref{eq:p9c1-scalar-period-trace}. The same calculation on the layer with $V_m e_m,t'_{hk},p_k$ holds in each component. Apply it there with $x\mapsto\lambda(b_k x)$ and sum over $k$.

For the layer at precision $L$ furnished by \eqref{eq:p9c1-weil-intertwine}, write $\Pi=\Pi_m$ and put
\[
 \delta_h=\sum_i\Pi_{ih}t_i-\sum_k b_k t'_{hk},\qquad
 \eta_0=\vartheta-\sum_k b_k p_k .
\]
Subtract the summed layer identities from \eqref{eq:p9c1-scalar-period-trace} to obtain
\begin{equation}\label{eq:p9c1-weil-trace-error}
 s\,u_{\rm ur}(\lambda(p-\vartheta))
 =\sum_i(Ve-\Pi V_m e_m)_i u_{\rm ur}(\lambda(t_i))
 +\sum_h(V_m e_m)_h u_{\rm ur}(\lambda(\delta_h))
 -s\,u_{\rm ur}(\lambda(\eta_0)).
\end{equation}
The first column on the right is divisible by $\ell^L$ by \eqref{eq:p9c1-weil-intertwine}. By \eqref{eq:p9c1-layer-traces}, $\delta_h/\ell^L$ and $\eta_0/\ell^L$ are integral since $m\ge L$.

If $p\ne\vartheta$, choose integers $c_1,c_2\ge0$ so that $\ell^{c_1}t_i$ and $x=\ell^{c_1}(p-\vartheta)$ are integral and $|x|_\ell\ge\ell^{-c_2}$. A finite extension of $K^{\rm ur}$ containing these numbers and the divided errors admits an integral basis over the unramified valuation ring. Take an integral basis of a finite local field containing the listed numbers over its maximal unramified subfield and extend scalars. Choose the coordinate of largest norm in the expansion of $x$, and extend this projection linearly to $\bar{\Q}_\ell$. The chosen $\lambda$ takes all the listed integral values to integral elements of $K^{\rm ur}$, and $q_L=\lambda(x)$ satisfies $|q_L|_\ell\ge|x|_\ell$, by the ultrametric inequality on the integral basis. Integral elements of $K^{\rm ur}$ map to $A_{\rm cris}$. Multiplying \eqref{eq:p9c1-weil-trace-error} by $\ell^{c_1}$ gives $s\,u_{\rm ur}(q_L)\in\ell^L A_{\rm cris}$. Since $\ell^{c_2}/q_L$ is integral, multiplication by its period image gives $\ell^{c_2}s\in\ell^L A_{\rm cris}$ for every $L$. Separatedness and $\ell$-torsion freeness give $s=0$, contrary to \eqref{eq:p9c1-denominator}. Thus $p=\vartheta$ for each $j$.
\end{proof}

\subsection*{Proof of \ref{B9.36c}}\label{prf:B9.36c}\phantomsection
\begin{proof}
Take a split place $w_*$ with both places of its pair outside $S$. Let $J$ be the group of finite unit idèles with $w_*$ and $\bar w_*$ coordinates one. Choose an open subgroup $V\subset\mathcal O_{E_{w_*}}^\times$ whose intersection with the global roots of unity is $\{1\}$. Let $D$ be the open subgroup of idèles with finite unit components for which the quotient of the $w_*$ component by the conjugate $\bar w_*$ component lies in $V$. On $D$ prescribe the given infinity character, trivial on finite components. On the product $E^\times\mathbf A_F^\times J$, prescribe $\eta_{E/F}(b)$ on $kbu$ with $k\in E^\times$, $b\in\mathbf A_F^\times$ and $u\in J$.

If $x=kbu\in D$, then $k/c(k)$ is a unit at every finite place, by taking the quotients $x/c(x)$ and $u/c(u)$. Its complex absolute values are all one, so it is a root of unity. Its component at $w_*$ belongs to $V$, hence $k=c(k)\in F^\times$. Now $kb$ has finite unit components. Its quadratic character value is the product of its real signs, which agrees with the infinity prescription since all $e_v$ are odd. Also $\eta_{E/F}(k)=1$. The two prescriptions agree on the intersection and define a unitary character on the product with $D$. Extend it to the idèle group by divisibility of the unit circle. Agreement on the open subgroup $D$ gives continuity. The extension is trivial on principal idèles and on $J$, hence unramified on $S$. At a nonsplit place the units lie in $J$ and a uniformizer from $F$ is a uniformizer of $E_w$ with quadratic value $-1$.
\end{proof}

\subsection*{Proof of \ref{B9.36b}}\label{prf:B9.36b}\phantomsection
\begin{proof}
Choose one representative of each bad split pair and one further good split place $w_0$. At each bad representative the indicator of an open subgroup of the standard maximal compact fixing a nonzero vector has positive trace. At $w_0$, the Iwahori Hecke operator of the decreasing translation $(0,-1,\ldots,1-n)$ is invertible on the nonzero Iwahori invariants. If the dimension is $d$, vanishing of its first $d$ power traces would give determinant zero by Newton's identities. Choose a positive power of nonzero trace. Translation lengths add in this chamber, so the power is a nonzero multiple of the indicator of the double coset with multiplied exponents. The Newton polygon gives distinct root valuations on this coset.
\leanref{UpperTheory.\allowbreak{}V.\allowbreak{}Z.\allowbreak{}Tf.\allowbreak{}GlobalTraceComparison.\allowbreak{}Asm.\allowbreak{}rChoice}{Development.\allowbreak{}AutomorphicArithmetic.\allowbreak{}AutomorphicVarieties.\allowbreak{}ModularityArgument.\allowbreak{}TraceFormula.\allowbreak{}GlobalTraceComparison.\allowbreak{}TraceAsymptotics.\allowbreak{}RChoice}{13--74}

Use the infinite weight of the twisted character of $\pi$, the split normalized constant term and the nonsplit hyperspecial matching of \ref{B6.04b}. For a split block $(a,b)$ the multiplier of the Levi element $(m_1,m_2)$ is $\mu_w((-1)^{ab}\det(m_1)^b\det(m_2)^a)$. This gives the twisted Satake union. At $w_0$, a nonzero constant-term integrand comes from a block triangular matrix in the regular support of $f_{w_0}$, so its product polynomial is separable.
\leanref{UpperTheory.\allowbreak{}ConvolutionTheory.\allowbreak{}G2.\allowbreak{}matchSplit}{Development.\allowbreak{}AutomorphicArithmetic.\allowbreak{}Foundations.\allowbreak{}Convolution.\allowbreak{}WeightRankUnits.\allowbreak{}MatchSplit}{12--21}
\leanref{PositivityTheory.\allowbreak{}Rg.\allowbreak{}Bc.\allowbreak{}ArchbcB.\allowbreak{}B.\allowbreak{}archimedean\_\allowbreak{}base\_\allowbreak{}change\_\allowbreak{}first}{Development.\allowbreak{}WeilBound.\allowbreak{}RepresentationTheoryGL.\allowbreak{}ArchimedeanBaseChange.\allowbreak{}ArchimedeanBaseChangeProof.\allowbreak{}IwahoriFactorisation.\allowbreak{}MultiplicityOneBase.\allowbreak{}ArchimedeanBaseChangeFirst}{29--61}

The definite space at the fixed weight and level is finite-dimensional. At a good split pair the Satake multisets are mutually inverse and the spherical fixed space is a line. The nonsplit integral levels select unramified summands. Factoring the trace gives the unitary sum. The fixed tests contribute $A_h$.

Fix the level and the archimedean compact types from the tests at $Z$ and at infinity. The tests at the varying places will be spherical at this same level. For a cuspidal block of rank $d>0$ with prescribed infinitesimal multisets $m_\tau$, normalization on the positive split centre leaves a twist $|\det|^{z_0}$ with $z_0=(d[E:\Q])^{-1}\sum_\tau\sum_{x\in m_\tau}x$. For a sequence approaching these multisets the twists converge to $z_0$, by the same formula. Every normalized exponent on its tail has modulus at most $2+|z_0|+\max_{\tau,x\in m_\tau}|x|$. Lemma~\ref{B6.04u} applies at the fixed open level with these compact types and this bound. Its finite list contains a constant subsequence of normalized cuspidal representations. Restoring the twist gives convergence of their spherical coordinates at each place outside the level. On the orthogonal complement of $\pi$ and its split-centre translates, omit the normalized class of $\pi$. For a noncuspidal eigenform choose a nonzero proper constant term of ranks $r,s$. There are finitely many choices of $r+s=d$ and of the submultisets at each embedding. In the unnormalized constant term the two lower infinitesimal multisets contribute with shifts $-s/2$ and $r/2$. Hence each partition fixes both lower limiting multisets. The lower levels are fixed as well: in the compact realization take the intersections of the Levi with the conjugates of the given level for representatives of the finite quotient of the maximal compact by that level. The restrictions of the finitely many archimedean compact types give finitely many lower types. Induction on rank gives a finite union of products of the two lower lists. Rank zero contributes the empty list. Taking the product for the two original blocks completes the construction before choosing $T$ or its polynomials.
\leanref{UpperTheory.\allowbreak{}V.\allowbreak{}Z.\allowbreak{}Tf.\allowbreak{}GaloisFromTraceFormula.\allowbreak{}cuspReadI\_\allowbreak{}of}{Development.\allowbreak{}AutomorphicArithmetic.\allowbreak{}AutomorphicVarieties.\allowbreak{}ModularityArgument.\allowbreak{}TraceFormula.\allowbreak{}GaloisFromTraceFormula.\allowbreak{}CuspReadIOf}{63--74}

In the commutative algebra of Hecke operators, adjoints and archimedean multipliers, the character-space fibre at the prescribed level and archimedean values is compact. Its evaluations come from the preceding lists, so the tested operator has modulus less than $\rho$ on the fibre. On the closed set where its modulus is at least $\rho$, compactness gives finitely many differences $d_j$ from the prescribed data with no common zero. After multiplication by the level projector, a sufficiently high power of $1-\sum_jd_j^*d_j/M$, with $M>\|\sum_jd_j^*d_j\|$, suppresses this closed set. The multiplier is one on the prescribed fibre and leaves the stable test unchanged. Call the resulting operator $U$. On the block space to which the competing list applies, its norm is less than $\rho$.
\leanref{UpperTheory.\allowbreak{}Zero.\allowbreak{}ConvolutionTheory.\allowbreak{}EB.\allowbreak{}rSpec\_\allowbreak{}of}{Development.\allowbreak{}AutomorphicArithmetic.\allowbreak{}Foundations.\allowbreak{}Convolution.\allowbreak{}EndomorphismSlices.\allowbreak{}REndoOf}{24--47}

The fixed end kernels give $T_1U^kT_2$. Each logarithmic support size is at most $r_0+kr_1$ for fixed $r_0,r_1\ge0$ by subadditivity. Regular support at $w_0$ and the very cuspidal archimedean kernels give the slice hypotheses of \ref{B6.04c}. Haar conversion from the elliptic orbital sum to the stable term is a fixed nonzero factor. Apply \ref{B6.04c} on the full space of the proper block to obtain \eqref{eq:p9c1-proper-block}.

At the distinguished end let $\operatorname{ldet}x$ be the adelic log determinant modulus, and let $Q_\pi$ be the closed span of $\psi(x)u(\operatorname{ldet}x)$ with $\psi$ in the normalized $\pi$ and $u$ continuous of compact support. Both $Q_\pi$ and its orthogonal complement are preserved by convolution and its adjoint. Use spectral separation on $Q_\pi^\perp$. On $Q_\pi$ the full twisted diagonal contributes half the twisted trace of the normalized cusp form. The involution reverses $\operatorname{ldet}$, so a central test $b$ of integral one contributes $\int_\R b(-2s-c')\,ds=1/2$ for any shift $c'$. The archimedean weight trace is nonzero for this conjugate self-dual $\pi$. The compact conjugation average preserves each split trace. With $c_0$ the fixed nonzero weight and volume factor and $\kappa$ the conversion from stable term to elliptic orbital sum, the contribution is $Pv(B)^k$ with $c=c_0/(2\kappa)$.

Put $R_k=(1+k)(1+r_0+r_1)$ and use \eqref{eq:p6c1-cut-trace} with $\tau=C'R_k$, where $C'$ is a sufficiently large fixed constant. Insert the orthogonal projection onto $Q_\pi^\perp$ in the middle of the row and column estimate \eqref{eq:p6c1-diagonal-count}. By \eqref{eq:p6c1-integrated-count}, the cut integral on this summand is bounded by $C(1+k)^N\|U^k|_{Q_\pi^\perp}\|$. Both projections commute with bounded functions of $\operatorname{ldet}$, by the description of $Q_\pi$. The untwisted kernel has log determinant difference bounded by $cR_k$, for a fixed $c>0$, and the projections preserve that support. The twist reverses the second log determinant. Hence the projected twisted diagonal is supported in $|\operatorname{ldet}x|\le cR_k$.

Write $\mathrm{red}(x)$ for the infimum of the logarithmic size of $\gamma x$ over rational translates, as in Lemma~\ref{B6.04c}. Let $a_c(s)$ be the positive split central element of log determinant $s$, with scalar $\exp(s/(n[E:\Q]))$ at each infinite place and finite components one. The determinant-one coordinate has reduced logarithmic size
\[
 h(x)=\mathrm{red}\bigl(a_c(-\operatorname{ldet}x)x\bigr).
\]
The logarithmic size of $a_c(s)$ is bounded by a constant times $|s|$. Since $a_c(s)$ commutes with rational translates, $\mathrm{red}(x)\le h(x)+c(1+|\operatorname{ldet}x|)$ after increasing $c$. The degree bounds in \ref{B6.04c} give, for a fixed $c_1>0$,
\begin{equation}\label{eq:p9c1-cut-reduction}
 c_1(1+\mathrm{red}(x))\le\tau
 \quad\Longrightarrow\quad x\in\operatorname{Cut}(\tau).
\end{equation}
Indeed, a representative within one of the infimum bounds the degree of each rational subspace from above and the total degree from below. On the determinant band, $h(x)<R_k$ would give $1+\mathrm{red}(x)\le(2+c+c^2)R_k$. Take $C'\ge c_1(2+c+c^2)$ as well as large enough for the cut identity. Then $h(x)\ge R_k$ on the part of the band outside this same cut.
\leanref{UpperTheory.\allowbreak{}V.\allowbreak{}Z.\allowbreak{}Tf.\allowbreak{}GlobalTraceComparison.\allowbreak{}Pi.\allowbreak{}r0}{Development.\allowbreak{}AutomorphicArithmetic.\allowbreak{}AutomorphicVarieties.\allowbreak{}ModularityArgument.\allowbreak{}TraceFormula.\allowbreak{}GlobalTraceComparison.\allowbreak{}AutomorphicRepresentations.\allowbreak{}EvalBasis}{18--21}

On $Q_\pi$ expand the left end kernel on the determinant-one quotient in a finite basis of cusp forms of the fixed level and type. Its coefficient kernels are evaluations at fixed points translated by $a_c(\operatorname{ldet}x)$. In the row and column count of \ref{B6.04c}, these translated rows have bounded norms, since the split centre commutes with conjugation by rational elements. Thus the absolute diagonal is bounded by a constant times $\|U^k\|n_\Omega(x)$ times the sum of the absolute values of the basis forms at the determinant-one coordinate. Rapid decay bounds each basis form there by $C_L e^{-(L+1)h(x)}$ \cite[Thm.~A]{miller-schmid-2012}. On the determinant band, $H=1+\mathrm{red}(x)\ll(1+cR_k)(1+h(x))$. Use the integrable weight $H^{-4n}n_\Omega$ from \ref{B6.04c} and the boundedness of $(1+h)^{4n}e^{-h}$ for $h\ge0$. The absolute integral of the cuspidal diagonal outside the cut is at most
\begin{equation}\label{eq:p9c1-cusp-tail}
 C_L(1+cR_k)^{4n}e^{-LR_k}\|U^k\|.
\end{equation}
Set $\eta=\rho/(1+\|U\|)$ and $L=1+\max(0,-\log\eta)$. Since $R_k\ge1+k$, the factor $e^{-R_k}$ absorbs the polynomial, and $e^{-(L-1)R_k}\le\eta^k$. Thus \eqref{eq:p9c1-cusp-tail} is at most $C_\eta\eta^k\|U^k\|\le C_\eta\rho^k$. Add the bound on the cut integral on $Q_\pi^\perp$, using $\|U|_{Q_\pi^\perp}\|<\rho$, and divide by the fixed Haar conversion factor. This gives \eqref{eq:p9c1-end-block}.
\leanref{UpperTheory.\allowbreak{}V.\allowbreak{}Z.\allowbreak{}Tf.\allowbreak{}GlobalTraceComparison.\allowbreak{}Pi.\allowbreak{}growthPiJ\_\allowbreak{}holds}{Development.\allowbreak{}AutomorphicArithmetic.\allowbreak{}AutomorphicVarieties.\allowbreak{}ModularityArgument.\allowbreak{}TraceFormula.\allowbreak{}GlobalTraceComparison.\allowbreak{}AutomorphicRepresentations.\allowbreak{}GrowthPiJHolds}{52--54}

For a proper block impose inversion invariance of each of $X_1(w)$ and $X_2(w)$ at nonsplit places. For the $(n,0)$ end impose this where $B(w)=B(w)^{-1}$. For each competitor failing the applicable condition choose one nonsplit witness. An elementary symmetric difference $d$ between its multiset and its inverse distinguishes the two, after clearing determinants. The polynomial $1-d/d(X)$ at the witness kills that entry and equals one on the imposed locus. Conjugate inverse transpose interchanges the terms of the difference, so the modification preserves the twisted orbital matching. Insert the finitely many modifications before varying the split tests. On the distinguished spherical line each inserted factor is one by $\pi^c\simeq\pi^\vee$. At nonsplit places the twisting values $\mu_w(\varpi_w)=-1$ preserve inversion invariance. The weight equality of \ref{B6.04b} gives regularity of the shifted infinitesimal characters. At the ends the rank-zero polynomial and character are one and the end unipotent radical is trivial. Interchanging the factors identifies the two stable end terms.
\leanref{UpperTheory.\allowbreak{}V.\allowbreak{}Z.\allowbreak{}Tf.\allowbreak{}GaloisFromTraceFormula.\allowbreak{}endoS}{Development.\allowbreak{}AutomorphicArithmetic.\allowbreak{}AutomorphicVarieties.\allowbreak{}ModularityArgument.\allowbreak{}TraceFormula.\allowbreak{}GaloisFromTraceFormula.\allowbreak{}EndoS}{10--44}
\leanref{UpperTheory.\allowbreak{}V.\allowbreak{}Z.\allowbreak{}Tf.\allowbreak{}GaloisFromTraceFormula.\allowbreak{}stableS}{Development.\allowbreak{}AutomorphicArithmetic.\allowbreak{}AutomorphicVarieties.\allowbreak{}ModularityArgument.\allowbreak{}TraceFormula.\allowbreak{}GaloisFromTraceFormula.\allowbreak{}StableS}{10--55}
\end{proof}

\subsection*{Proof of \ref{B9.36a}}\label{prf:B9.36a}\phantomsection
\begin{proof}
Take $\mu$ by \ref{B9.36c} with $S=\varnothing$ and use the lists of \ref{B9.36b}. Write $n$ for the rank of $\pi$. Suppose neither a unitary summand nor a competing cuspidal list agrees with $B$ at almost every split place. Each proper-block factor has rank smaller than $n$, so its twisted list omits $\pi$. For each competitor $X$ choose a split witness $w$ outside $Z\cup\bar Z$ with $X(w)\ne B(w)$. A separating elementary symmetric coordinate $H$ gives the test factor
\[
 \frac{H-H(X(w))}{H(B(w))-H(X(w))}.
\]
For a witness on the conjugate side of the section use the inverted multisets in the factor for $e_{2,\bar w}$. Multiplying the factors gives $v(B)=1$ and $v(X)=0$ for each competitor.
\leanref{UpperTheory.\allowbreak{}V.\allowbreak{}Z.\allowbreak{}Tf.\allowbreak{}GlobalTraceComparison.\allowbreak{}Asm.\allowbreak{}rSep}{Development.\allowbreak{}AutomorphicArithmetic.\allowbreak{}AutomorphicVarieties.\allowbreak{}ModularityArgument.\allowbreak{}TraceFormula.\allowbreak{}GlobalTraceComparison.\allowbreak{}TraceAsymptotics.\allowbreak{}ConstantTermParabolicAverage}{82--112}

For $k\ge1$ the unitary trace vanishes. Apply the signed identity \ref{B6.04b} and \eqref{eq:p9c1-proper-block} and \eqref{eq:p9c1-end-block} with $\rho=1/2$. Each proper-block term tends to zero. For $(0,n)$ the selected minor is empty at every real place. For $(n,0)$ it uses the indices $0,\ldots,n-1$. In the sign of Lemma~\ref{B6.04b} the exponent is respectively $0$ and $n(n-1)$, so both signs are $+1$. Each ordered end occurs once. By Lemma~\ref{B9.36b} their unsigned terms are equal. In its notation each has limiting contribution $P=c_0(2\kappa)^{-1}\prod_{w\in Z}\operatorname{tr}\pi_w(f_w)$, since $v(B)=1$. The factor $1/2$ comes from integrating the central test under $s\mapsto-2s-c'$. Thus the right side tends to $2P\ne0$, whereas the left side is zero for every $k\ge1$.
\leanref{UpperTheory.\allowbreak{}ConvolutionTheory.\allowbreak{}G2.\allowbreak{}uSide}{Development.\allowbreak{}AutomorphicArithmetic.\allowbreak{}Foundations.\allowbreak{}Convolution.\allowbreak{}WeightRankUnits.\allowbreak{}USide}{9--49}

Suppose a competing cuspidal list matches. Write a rank-$d$ factor as a normalized cusp form twisted by $|\det|^z$. The normalized exponents sum to zero over all embeddings. The shifted exponents are real, so $d[E:\Q]\Im z=0$. At conjugate embeddings the normalized exponents pair as $x,-\bar x$. Since the shifted exponents are real and both $x+z$ and $-x+z$ belong to $(n-1)/2+\Z$, one has $2z\in\Z$. At a matching split place the geometric mean of this factor's root moduli is $q_w^{-z}$ by unitarity of its normalized central character. Conjugate self-duality gives unitarity of $\pi$. The Jacquet--Shalika bound puts the root moduli of $B(w)$ in $(q_w^{-1/2},q_w^{1/2})$ \cite{jacquet-shalika-1981}. Thus $|z|<1/2$ and $z=0$.
\leanref{UpperTheory.\allowbreak{}V.\allowbreak{}Z.\allowbreak{}Tf.\allowbreak{}GaloisFromTraceFormula.\allowbreak{}partNorm}{Development.\allowbreak{}AutomorphicArithmetic.\allowbreak{}AutomorphicVarieties.\allowbreak{}ModularityArgument.\allowbreak{}TraceFormula.\allowbreak{}GaloisFromTraceFormula.\allowbreak{}PartNorm}{166--197}

Let $j$ be the number of factors and set $M=16(j+1)^2$. The twists $\pi\otimes\mu^{2i}$, $0\le i<M$, are distinct unitary cusp forms because their sums of infinitesimal exponents differ. At nonsplit places $\mu_w^2=1$. Omit the finitely many exceptional places of these systems and of the factors. For a unitary cusp form $\sigma$ use coordinates $p_{w,r}(\sigma)=\sum_{\alpha\in B_\sigma(w)}\alpha^r$. Give them inner product weights $q_w^{-rs}/[r(-\log(s-1))]$ for $r\ge1$ and real $s>1$ near one. The partial Rankin--Selberg product has a simple pole for equivalent factors and is regular and nonzero at one for inequivalent factors. Taking its logarithm gives
\begin{equation}\label{eq:p9c1-rankin-inner}
 \lim_{s\downarrow1}\langle p(\sigma),p(\sigma')\rangle_s
 =\begin{cases}1&\sigma\simeq\sigma',\\0&\sigma\not\simeq\sigma'.\end{cases}
\end{equation}
Here the second coordinate is conjugated in the inner product \cite{jacquet-shalika-1981}.
\leanref{UpperTheory.\allowbreak{}V.\allowbreak{}Z.\allowbreak{}Tf.\allowbreak{}SR.\allowbreak{}antiTw}{Development.\allowbreak{}AutomorphicArithmetic.\allowbreak{}AutomorphicVarieties.\allowbreak{}ModularityArgument.\allowbreak{}TraceFormula.\allowbreak{}SplitRigidity.\allowbreak{}AntiTw}{30--60}

By \eqref{eq:p9c1-rankin-inner} the average of the $M$ twists has limiting squared norm $1/M$. Its nonsplit coordinates coincide with those of $\pi$, so the nonsplit squared norm of $\pi$ has limit superior at most $1/M$. Pair $\pi$ with the difference of its coordinates and the sum of the $j$ factors' coordinates. If every factor differs from $\pi$, the limit is one by \eqref{eq:p9c1-rankin-inner}. The split coordinates of the difference vanish at the places retained. Cauchy--Schwarz on the nonsplit coordinates bounds the absolute pairing in the limit by $1/M+j/\sqrt M<1$. Thus a factor is $\pi$, contrary to its omission. The unitary match gives the assertion.
\leanref{UpperTheory.\allowbreak{}V.\allowbreak{}Z.\allowbreak{}Tf.\allowbreak{}SR.\allowbreak{}splitRigid}{Development.\allowbreak{}AutomorphicArithmetic.\allowbreak{}AutomorphicVarieties.\allowbreak{}ModularityArgument.\allowbreak{}TraceFormula.\allowbreak{}SplitRigidity.\allowbreak{}SplitRigid}{12--13}
\leanref{UpperTheory.\allowbreak{}Gal.\allowbreak{}Desc.\allowbreak{}weakAll\_\allowbreak{}of}{Development.\allowbreak{}AutomorphicArithmetic.\allowbreak{}Foundations.\allowbreak{}GaloisRepresentations.\allowbreak{}BaseChangeDescent.\allowbreak{}Strong}{24--33}
\end{proof}

\subsection*{Proof of \ref{B9.37}}\label{prf:B9.37}\phantomsection
\begin{proof}
The character \(\chi_{E/F}\) is trivial on norms. The polarization therefore gives \(\pi^c\simeq\pi^\vee\). The extension \(E/F\) is unramified at the split finite places and is unramified at the remaining finite places by hypothesis. Lemma~\ref{B9.36a} gives an automorphic representation \(U\) of the definite unitary group which is spherical at every nonsplit finite place and matches \(\pi\) at almost every split place.

We apply \cite[Remark~5.2 and Corollary~5.3]{labesse-2011} to \(U\). The ground field \(F\) is totally real, \(E/F\) is totally imaginary quadratic, and the positive integer \([F:\Q]\) is even, hence at least two. The definite group is anisotropic modulo its centre, so \(U\) belongs to the discrete spectrum. At every real place the group is compact. Each irreducible archimedean component is consequently in the discrete series. To check Labesse's condition \((*)\), take Haar probability measure on the product of the real groups and use the complex conjugate of the character of \(U_\infty\) as pseudo-coefficient. Character orthogonality gives trace one on \(U_\infty\) and zero on every other irreducible archimedean representation. Thus every nonzero trace appearing in \((*)\) has the required sign.

Corollary~5.3 gives an isobaric transfer whose factors are discrete general linear representations. It gives local base change at every split place and compatibility of infinitesimal characters at infinity. By the discrete-spectrum theorem of \cite[p.~606]{moeglin-waldspurger-1989}, applied over the number field \(E\) to each discrete factor with its unitary central character, a factor has the form \(\operatorname{Speh}(\sigma,d)\) for a unitary cuspidal \(\sigma\) and \(d\ge1\). Choose a split place of matching outside the ramification sets of all factors. Such places exist by Chebotarev. Write \(q_w\) for the residue cardinality. For every Satake root \(\alpha\) of \(\sigma_w\) the root multiset of the transfer contains
\[
 \alpha q_w^{(d-1)/2-j},\qquad 0\le j<d.
\]
The common multiset is that of \(\pi_w\). Conjugate self-duality makes \(\pi\) unitary: the real modulus exponent of its central character is preserved by \(c\) and negated by duality, so it is zero. The local component is spherical and generic by cuspidality. The Jacquet--Shalika bound for \(\pi\) and its contragredient puts every root modulus strictly between \(q_w^{-1/2}\) and \(q_w^{1/2}\) \cite[Part~I, Corollary~2.5]{jacquet-shalika-1981}. If \(d\ge2\), the roots at \(j=0,1\) have modulus ratio \(q_w\), contradicting these strict inequalities. Thus every \(d\) equals one.

The transfer is now a sum of unitary cuspidal factors of total rank \(n=\operatorname{rank}\pi\). Their Satake union agrees with that of \(\pi\) at almost every split place. Apply the split-place Rankin--Selberg argument of Lemma~\ref{B9.36a}, equation~\eqref{eq:p9c1-rankin-inner}, to this finite list. One factor is isomorphic to \(\pi\). It has rank \(n\), so it accounts for the entire transfer. The split local base change and the infinitesimal-character assertion of Corollary~5.3 therefore give the claimed matching, including at ramified split places. The representation \(U\) obtained above is spherical at all the nonsplit finite places, so \(T=\varnothing\) works for every \(Z\).
\end{proof}

\subsection*{Proof of \ref{B9.38}}\label{prf:B9.38}\phantomsection
\begin{proof}
For quadratic polarization, first work over a field satisfying \ref{B9.37}, the inert-place conditions of \ref{B9.41}, and $x^2=-d\in E$, $-d\in(\Q_\ell^\times)^2$, with three pairwise separated places above $\ell$. Let $Q$ be the usable places for $d$. Take $U$ by \ref{B9.37}. Degree-one Chebotarev gives dense inverse Frobenius from $Q$ after finite removals. In odd rank apply \ref{B9.35}. For a target above $\ell$ take two of the three auxiliary places separated from it and use their conjugates as needed to agree with the target on $\Q(x)$. The joint approximation in the proof of \ref{B9.35} gives \ref{B9.36} for the same representation, with the Satake power sums as targets. Newton's identities give the stated Weil polynomial. The cofinite Frobenius data and semisimplicity identify the representations for the different choices of auxiliary pair. In even rank use \ref{B9.45}.

For descent, use two preliminary real quadratic steps to split the prescribed places over $\ell$ and obtain at least three separated places there. At either step one can adjoin the square root of a positive rational prime congruent to one modulo $8$ and the odd rational primes below the places to be preserved. Choose it outside the primes ramified in the normal closure of the current field and outside the primes below its nonspherical places. The square condition gives splitting at the prescribed places and valuation one at the fresh prime gives ramification at a spherical place. Varying the fresh prime gives infinitely many choices. In the second step write $Y=Y_0(\delta)$ with $Y_0$ the preceding CM field and $\delta^2$ the chosen prime. The normal closure of $Y_0$ is unramified at that prime, so it is disjoint from $\Q(\delta)$ over $\Q$. The sign change of $\delta$ therefore extends to an automorphism $s$ of the algebraic closure fixing every conjugate of $Y_0$ pointwise.

Fix a finite set $T$ of places of $Y$ containing the target, all places above $\ell$, and all places without Satake data. Take $c$ to be a positive multiple of $6$ divisible by the norms of the places in $T$ and of the places whose inertia over $Y^+$ or $\Q$ acts nontrivially on $Y$. In a common normal closure containing all conjugates of $Y$ and $\mu_{8c}$, let $c_\infty$ be complex conjugation. Choose unramified rational primes $p_1<p_2$ with Frobenius representatives $c_\infty s$ and $s^{-1}$, respectively. Chebotarev permits both primes to be arbitrarily large. Each representative sends $\delta$ to $-\delta$ and its square fixes every conjugate of $Y$. For $c_\infty s$ this uses that $c_\infty$ preserves the normal closure of $Y_0$ and fixes $\delta$. These two properties hold for their conjugates as well. The decomposition group in the normal closure is generated by its Frobenius, so at either prime a decomposition automorphism fixing $\delta$ fixes $Y$. On $\mu_{8c}$ the product of the two Frobenius restrictions is inversion, which gives $p_1p_2\equiv-1\bmod 8c$.
\leanref{UpperTheory.\allowbreak{}Gal.\allowbreak{}UnramifiedFromSplittingCriterion.\allowbreak{}twoPrimes}{Development.\allowbreak{}AutomorphicArithmetic.\allowbreak{}Foundations.\allowbreak{}GaloisRepresentations.\allowbreak{}InertiaQuadraticRamification.\allowbreak{}TwoPrimes}{37--89}

Put $d=p_1p_2$ and adjoin $x=\sqrt{-d}$. The congruence makes $-d$ a square in $\Q_p$ for every $p\mid c$, including at $2$ by the condition modulo $8$. All places above $T$ split in the extension of $Y$. They also split over the upper real subfield because their decomposition groups fix $x$.

At a nonsplit place of the upper CM field a decomposition automorphism over the upper real subfield moves $x$ and $Y$. Indeed $Y$ together with that real subfield generates the upper CM field. The place lies away from $c$ by the square conditions. It lies away from $p_1p_2$ since the automorphism fixes the real element $\delta$, which at either of those primes would force it to fix $Y$. At the remaining places inertia fixes $Y$ by the choice of $c$ and fixes $x$ since $2d$ is a unit. Hence the upper CM extension is unramified over its real subfield at the nonsplit places, and the real subfield is unramified over $\Q$ there. The residue characteristic is odd and the residue norm in the real subfield exceeds three. Each such place lies over a nonsplit place of $Y$, since the decomposition automorphism over the upper real subfield acts nontrivially over $Y^+$. The lower component is spherical by the choice of $T$. At a place of $Y$ over $p_1$, $-d$ has valuation one. This supplies a fresh ramified spherical place for the quadratic lift. Including prior ramification witnesses in $T$ gives infinitely many distinct choices.
\leanref{UpperTheory.\allowbreak{}Gal.\allowbreak{}Int.\allowbreak{}Core.\allowbreak{}exists\_\allowbreak{}quadratic\_\allowbreak{}CM\_\allowbreak{}ext\_\allowbreak{}split\_\allowbreak{}at}{Development.\allowbreak{}AutomorphicArithmetic.\allowbreak{}Foundations.\allowbreak{}GaloisRepresentations.\allowbreak{}IntegralStructures.\allowbreak{}CrystallineDescent.\allowbreak{}ExistsQuadraticCMExtSplitAt}{12--156}

At each step, \ref{B9.37a} gives cuspidal lifts. On the compositum of any two auxiliary fields, Chebotarev and the residue-degree rule identify the traces of the two upper representations. The same argument gives invariance under conjugation over the ground field. Apply the index-two patching lemma to the union of the auxiliary families for the targets at this step. Each target splits in infinitely many members, and every finite extension contains finitely many quadratic subfields. The patching lemma gives a continuous semisimple representation whose restriction is isomorphic to the chosen upper representation for all but finitely many members \cite[{\S}4.3, Lem.~1]{sorensen-2010}. For each target choose a member splitting it outside this exceptional set. The completions agree there, so the local data descend. This gives weak attachment and the stated local data at every spherical target. Cofinite Frobenius polynomials and semisimplicity give uniqueness up to isomorphism.

For general polarization first use fields with an imaginary quadratic subfield and even real degree, with nonsplit finite places unramified and away from $\ell$. Require sphericality at the nonsplit places. Let $b$ be the parallel infinity exponent of $\chi$ on positive real idèles from Lemma~\ref{B9.07a}, and put $\psi=\chi^{-1}\eta_{E/F}^{a+|b|}$. At each real place $\psi_v(x)=x^{-b}$ for $x\in\R^\times$. At nonsplit places $\psi$ is trivial on units by sphericality and polarization, since the unramified norm is onto on units.

Choose an auxiliary split pair above a real-subfield place distinct from that below the specified split place $w_0$. On the subgroup $D$ of idèles with arbitrary infinite components and finite unit components, impose deep congruence conditions at the auxiliary pair and, as needed, at other split places away from $w_0$. Keep all units at $w_0$ and at nonsplit places. Choose the levels so that $\psi$ is trivial on the finite part of $D\cap\mathbf A_F^\times$ and so that the congruence group at the auxiliary pair meets the global roots of unity trivially. At the place below $w_0$ one may impose the conductor condition at its conjugate place. Prescribe $\prod_{v\mid\infty}z_v^{-b}$ on $D$, using one complex embedding above each real place, prescribe $\psi$ on $\mathbf A_F^\times$, and prescribe the trivial character on $E^\times$.

For a relation $kbu=1$ with $k\in E^\times$, $b\in\mathbf A_F^\times$ and $u\in D$, the ratio $k/c(k)=c(u)/u$ is a unit at every finite place and has absolute value one at every complex embedding. It is a root of unity in the auxiliary congruence group, hence equals one. Thus $k\in F^\times$ and $kb=u^{-1}$ belongs to $D\cap\mathbf A_F^\times$. The infinity prescription on this intersection agrees with $\psi$ by its real exponents and signs, and the finite restriction is trivial. The prescriptions therefore give one character on the product subgroup. Extend it to the idèles by divisibility of $\C^\times$. Agreement on the open subgroup $D$ gives continuity. The resulting algebraic Hecke character $\kappa$ is unramified at $w_0$ and at nonsplit places. It satisfies $\kappa\kappa^c=\psi\circ N_{E/F}=(\chi\circ N_{E/F})^{-1}$.
\leanref{GlobalTheory.\allowbreak{}HeckeUnrExt.\allowbreak{}exists\_\allowbreak{}ext\_\allowbreak{}unr}{Development.\allowbreak{}AbelianVarietyTheory.\allowbreak{}Proofs.\allowbreak{}SatakeParameters.\allowbreak{}ExistsAlgebraicHeckeCharSplit}{27--150}

The twist $\pi\otimes\kappa$ has quadratic polarization. Twist its representation back by $(\kappa^{-1})_\ell$. At a spherical target above $\ell$ this character has a crystalline realization by the algebraic Hecke character comparison. Its Weil eigenvalue at degree $-1$ is the local unramified value under reciprocity and its jump is the negative of its infinity exponent at the inducing embedding. Tensoring cancels the Satake scaling and subtracts the jump shift of the original twist. For arbitrary fields use quadratic base change to fields with the required properties and descend at splitting places by the same patching lemma.
\end{proof}

\subsection*{Proof of \ref{B9.38a}}\label{prf:B9.38a}\phantomsection
\begin{proof}
The real exponent of the central character divided by $a$ is $t=(a-1-w_{\rm gt})/2$, since the sum of infinitesimal weights at an embedding and its conjugate is $a(a-1-w_{\rm gt})$. Thus $\pi|\det|^{-t}$ belongs to the unitary cuspidal spectrum and $\sigma^*|\det|^{w_{\rm gt}/2}$ is unitarizable and generic. The generic unitary dual has unitary discrete-series blocks with real twists $e$ satisfying $|e|<1/2$ \cite{tadic-1986}. A block of length at least two contributes adjacent $\nu_K$-twists to the support. Unlinking forces length one, and the character hypothesis forces rank one. Removing the common twist gives \eqref{eq:p9c1-unlinked-exponents}.
\end{proof}

\subsection*{Proof of \ref{B9.38e}}\label{prf:B9.38e}\phantomsection
\begin{proof}
Use multiplicativity of local \(\gamma\)-factors for the generic constituent of normalized parabolic induction \cite[Thm.~3.1]{jacquet-piatetskii-shapiro-shalika-1983}. It applies with a character in the other variable. The inducing supercuspidals are generic, and the generic constituent with support \(S\) is the given representation. A supercuspidal of rank greater than one paired with a character has both standard \(L\)-factors equal to one, as does a ramified character \cite[\S8]{jacquet-piatetskii-shapiro-shalika-1983}\cite[\S2.5]{tate-1967}. Its \(\gamma\)-factor is its epsilon monomial. For an unramified character of uniformizer value \(u\), the quotient of the dual and direct \(L\)-factors in Tate's functional equation is
\[
 \frac{1-uT}{1-u^{-1}q^{-1}T^{-1}}
             =-uqT\,\frac{1-uT}{1-uqT}.
\]
Multiplication over the support proves the formula, including multiplicities. It allows cancellation between the numerator and denominator for adjacent characters in a segment.
\end{proof}

\subsection*{Proof of \ref{B9.38b}}\label{prf:B9.38b}\phantomsection
\begin{proof}
Put \(s_0=(1-a)/2\) and \(\eta_j=\xi_j|\cdot|_K^{-s_0}\). These are the support characters before normalization. Let
\[
 U=\bigcap_j\ker(\eta_j|_{\mathcal O_K^\times}),\qquad
 H=U\langle b_K(g)\rangle .
\]
The group \(U\) is open of finite index in \(\mathcal O_K^\times\). Since \(\operatorname{ord}_K b_K(g)=-\deg_K g\ne0\), \(H\) is open of finite index in \(K^\times\) and \(H\cap\mathcal O_K^\times=U\). Local reciprocity gives a finite abelian extension \(L/K\) in the chosen closure with \(W_L=b_K^{-1}(H)\). In particular \(g\in W_L\), and every \(\eta_j\circ N_{L/K}\) is unramified.

We realize this local extension in a tower preserving cuspidality. At a CM stage \(D\), write \(r_D\) for the restricted representation. Its trace is \(a\) at the identity, so there is an open normal subgroup on which the trace never vanishes. Let \(D_1/D\) be its fixed field. In a finite Galois closure over \(D^+\) containing \(D_1\), choose by Chebotarev auxiliary places, away from the target, with Frobenius representatives generating the subgroup fixing \(D\). These places split in \(D\). Requiring them to split in a real abelian extension \(F'/D^+\) forces \(DF'\cap D_1=D\). Indeed their chosen Frobenius elements restrict trivially to the intersection and generate on \(D_1\).

First prescribe the local extension \(K/F_v\) at the place of the real subfield below the target, the trivial extension at these auxiliary places, and splitting at all real places. In the compositum with \(E\) the completion at the chosen prolongation is still \(K\) and the place is now split over the real subfield. Repeat there, prescribing \(L/K\), with the auxiliary places for the new trace subgroup. At each of the two stages, let $L_0/K_0$ be the prescribed extension at the target. Choose generators of the character group of $K_0^\times/N_{L_0/K_0}(L_0^\times)$ and prescribe the trivial character at each auxiliary and real place for every generator. The prescribed set of places is finite. By Grunwald--Wang, each generator of order $m$ extends with these local components to an idèle class character of order dividing $m$, or dividing $2m$ in the exceptional case \cite[Thm.~2.2]{wang-2015}. Take the compositum of the fields cut out by these characters. At the target the intersection of their local kernels is $N_{L_0/K_0}(L_0^\times)$, so local reciprocity identifies the completion with $L_0$. The prescribed trivial real characters make the auxiliary fields totally real. The local extensions are normal, so the prescribed image in the chosen local closure is \(L\). The resulting CM compositum \(M\) is fixed by \(g\).

Factor the two abelian towers into cyclic steps of prime degree. At each step the trace subgroup from the start of that tower surjects onto the remaining Galois group after intersection with the current absolute Galois group. This follows from the disjointness just proved. A nontrivial self-twist by a character \(\omega\) of the step would give
\[
           \operatorname{tr}r_D(\gamma)
          =\omega(\gamma)\operatorname{tr}r_D(\gamma)
\]
on all of \(G_D\), by weak attachment at good places and Chebotarev. Choosing \(\gamma\) in the trace subgroup with \(\omega(\gamma)\ne1\) contradicts this identity. Hence prime-degree cyclic base change is cuspidal \cite[Ch.~III, Thms.~4.2 and 5.1]{arthur-clozel-1989}. Strong multiplicity one transports the polarization, with character pulled back by the norm. Restriction preserves semisimplicity at a finite-index normal subgroup. At good places of residue degree \(f\), both the Satake values and the inverse-Frobenius eigenvalues pass to their \(f\)-th powers. This proves weak attachment at every stage.

We compute the support at the target one step at a time. In rank one the lift is the norm pullback of the Hecke character, by equality at almost every place and multiplicity one. At a split step the local fields and their components identify by the split-place base-change identity \cite[Ch.~III, Thm.~5.1]{arthur-clozel-1989}. Consider a nonsplit step \(k'/k\) of degree \(b\), residue degree \(f\), with norm \(N\) and lower residue cardinality \(q\). Write \(S,S'\) for the lower and upper support, and \(\omega\) for the local character of the step, obtained from a faithful character of the global cyclic quotient. We prove for every smooth \(\lambda\) of \(k^\times\) that
\begin{equation}\label{eq:p9c1-local-euler-polynomials}
 D_{k'}(S',\lambda\circ N;T^f)
       =\prod_{i=0}^{b-1}D_k(S,\lambda\omega^i;T).
\end{equation}

Prescribe \(\lambda\) on the target units. At every other exceptional lower place prescribe a finite-order unit character avoiding the inverse unit restrictions of the finitely many rank-one support terms in all lower twists. Require its norm pullbacks to avoid the corresponding upper restrictions. These choices exist: the norm image on units is open, so characters of sufficiently large conductor remain outside each of the finitely many forbidden restrictions upon pullback. Include in the exceptional set the ramification of the extension and of the additive characters. Extend these finite unit prescriptions to a global finite-order character by \cite[Ch.~X, Thm.~5]{artin-tate-2009}, first extending each prescription trivially on a uniformizer. At any new ramified place of the twist the original components were spherical and the field extension unramified or split. The norm is then onto on units. Thus the support polynomials there are one on both sides. An unramified difference at the target with uniformizer value \(v\) replaces \(T\) by \(vT\) on both sides, which is an invertible substitution.

Omit these finitely many places from both Euler products. The product for the \(b\) lower twists and the product for the upper representation agree in a right half-plane: at a good place the cyclic character values give \(\prod_{\zeta^f=1}(1-\zeta z)=1-z^f\) for each place above it. The same calculation applies to the contragredients. Apply the global standard functional equation to these cuspidal twists \cite[Thm.~13.8]{godement-jacquet-1972}. In the quotient of the two functional equations the incomplete products cancel. For the omitted finite factors use Lemma~\ref{B9.38e}. If \(P_0,P_1\) are the two sides of \eqref{eq:p9c1-local-euler-polynomials} and \(H=P_0/P_1\), the contribution of the target, up to a Laurent monomial, is
\[
                            H(T)/H(qT).
\]
Indeed upstairs the substitution is \(T^f\) and the residue cardinality is \(q^f\), so the denominator there is evaluated at \((qT)^f\). All the remaining omitted finite factors contribute monomials in their respective residue cardinalities to the power \(-s\).

Divide the resulting meromorphic identity at \(s+1\) by the identity at \(s\). The quotient of every exponential factor is constant. At a complex place the archimedean factors are products and quotients of shifted complex gamma functions, so their quotients under an integral shift are rational functions of \(s\) up to constants, by \(\Gamma(z+1)=z\Gamma(z)\). All infinite places in the comparison are complex. We obtain, up to a nonzero constant and inversion, a rational function of \(s\) equal to
\[
                    \frac{H(q^{-1}T)H(qT)}{H(T)^2}.
\]
Since \(H(0)=1\), this expression tends to one with error \(O(q^{-s})\) as real \(s\) tends to \(+\infty\). A nonconstant rational function tending to a finite limit has a nonzero leading power of \(1/s\) in its difference from the limit. Hence the displayed expression is identically one. For any \(z\ne0\), the integers \(\operatorname{ord}_{q^jz}H\) now have zero second difference in \(j\). They have finite support, so they all vanish. The value at zero gives \(H=1\). This proves \eqref{eq:p9c1-local-euler-polynomials}.
\leanref{UpperTheory.\allowbreak{}V.\allowbreak{}Z.\allowbreak{}ImaginaryQuadraticEtaCharacter.\allowbreak{}matchBCX}{Development.\allowbreak{}AutomorphicArithmetic.\allowbreak{}AutomorphicVarieties.\allowbreak{}ModularityArgument.\allowbreak{}LeviSmoothLocalization.\allowbreak{}PrincipalSeriesMatching.\allowbreak{}MatchBCX}{302--395}

Suppose inductively that \(S\) consists of characters. For \(\mu\in S\), the condition that \((\mu\lambda)\circ N\) be unramified is equivalent to triviality of \(\mu\lambda\) on \(N(\mathcal O_{k'}^\times)\). By reciprocity this is exactly the condition that some \(\mu\lambda\omega^i\) be unramified. If one such \(i\) exists, there are \(f\) of them, differing by the unramified characters of the cyclic quotient. Their uniformizer values differ by all \(f\)-th roots of unity. Their product in the right side of \eqref{eq:p9c1-local-euler-polynomials} is therefore
\[
                  1-((\mu\lambda)\circ N)(\varpi_{k'})T^f.
\]
Here an unramified character takes on \(N\varpi_{k'}\) its lower uniformizer value to the power \(f\), and \(\omega\) is trivial on norms.

For each distinct unit restriction of a pulled-back \(\mu\), choose \(\lambda\) inverse on units to a lower representative. The identity identifies the uniformizer values, with multiplicity, of all upper characters in that class. The values on units and on a uniformizer determine the characters. Summing over these disjoint classes accounts for \(a\) upper rank-one terms. The sum of the ranks of all upper support terms is \(a\), so the list is exhaustive. This proves the induction without assuming anything about upper segment lengths.

At the top the raw support is \(\{\eta_j\circ N_{L/K}\}_j\) and the normalized support is \(\{\xi_j\circ N_{L/K}\}_j\). The first list is unramified. By \ref{B9.38a} the real exponents differ pairwise by less than one, and norm pullback preserves the exponents since \(|N_{L/K}z|_K=|z|_L\). The pulled-back characters are therefore unlinked. Their normalized induction is irreducible \cite[Thm.~4.2]{zelevinsky-1980} and has a spherical line. Its normalized Satake values are \((\xi_j\circ N_{L/K})(\varpi_L)\). For \(h\in W_L\), use \(b_K(h)=N_{L/K}b_L(h)\) and \(\operatorname{ord}_L b_L(h)=-\deg_L h\). This gives the stated polynomial for every \(h\), including \(g\).
\end{proof}

\subsection*{Proof of \ref{B9.38c}}\label{prf:B9.38c}\phantomsection
\begin{proof}
For $g$ of nonzero degree choose $M,u$ by \ref{B9.38b}, with $L=M_u$. Away from $\ell$, \ref{B9.38} makes the restriction unramified and gives its polynomial for $B$. Above $\ell$ it makes the restriction crystalline and gives its period polynomial on $W_L$. Thus $r_w$ is potentially semistable above $\ell$, and the chosen bridge extends to $L$. Write $x$ for the downstairs parameter. Weil--Deligne restriction identifies the monodromy operators and restricts the Weil action. For the period parameter this follows by scalar extension of the potentially semistable module with descent action. Hence $N_x=0$. Since $g\in W_L$, \ref{B9.38b} identifies the characteristic polynomials of $\rho_x(g)$ and $\rho_Y(g)$.

For $t$ of degree zero take $\phi$ of degree one. Put $A=\rho_x(\phi)$, $B_0=\rho_Y(\phi)$ and $d_j=\operatorname{tr}\rho_x(t\phi^j)-\operatorname{tr}\rho_Y(t\phi^j)$. Thus $d_j=0$ for $j\ge1$. Write $\det(X-A)\det(X-B_0)=\sum_jc_jX^j$. Apply Cayley--Hamilton on $A$ and $B_0$, multiply by $\rho_x(t)$ and $\rho_Y(t)$ respectively, and subtract the traces. This gives $\sum_j c_jd_j=0$. Invertibility gives $c_0\ne0$, so $d_0=0$. Applying the argument to powers of $t$ gives the inertia polynomials by Newton's identities. Brauer--Nesbitt identifies the semisimple Weil actions. Both monodromy operators are zero.

For the jumps use the field chosen for $\phi$. The filtered de Rham module descends after passage to a Galois closure. Each induced coefficient embedding downstairs extends upstairs after increasing the coefficient field. The infinitesimal character of $\Pi$ at the extended embedding is that of $\pi$ at its restriction. Thus \ref{B9.38} gives the desired jumps upstairs and faithful scalar extension preserves the graded dimensions.

Every segment of $\sigma^*$ has length one because a longer segment contains adjacent support characters. With the reciprocity map of \ref{B9.50}, the segment and rank-one laws of \ref{B9.47} give $R(\sigma^*)=Y$.
\end{proof}

\subsection*{Proof of \ref{B9.38d}}\label{prf:B9.38d}\phantomsection
\begin{proof}
The place of $F$ below $w$ splits in $E$ since $[E_w:\Q_p]=1$. Let $S$ be the finite set where $\pi$ has no Satake multiset. Prescribe an unramified character at the target of $F$ taking Frobenius to a primitive $f$-th root of unity and trivial characters at the real places. The prescribed characters have finite order and the target together with the real places is a finite set. The character form of Grunwald--Wang gives a finite-order character $c$ with exactly these restrictions \cite[Thm.~2.2]{wang-2015}. Its exponent may be twice the least common multiple of the local orders; the descent below starts with the exponent of $c$.

Start with an exponent $m$ of $c$. At a stage with field cut out by $c^m$, if $m>1$ choose a prime $b\mid m$ and put $m'=m/b$. Choose an even auxiliary character $\psi$ of order $m$, trivial at the target, whose inertia image has order $m$ at a fresh place outside $S$ split in $E$ and in the field cut out by $c$. Choose a rational prime that splits completely in the normal closures of those fields and in $\Q(\mu_{2m})$. Require it also to split completely in the splitting fields of $X^m-h$ for rational primes $h$ below the places of prescribed triviality. A residue character of order $m$ there is even and each prescribed $h$ is an $m$-th power. Its restriction to $F$ has full inertia order since the rational prime splits completely in $F$.

The characters satisfy $(c\psi)^m=c^m$. The extension cut out by $(c\psi)^{m'}$ over the field cut out by $c^m$ has degree dividing $b$, and the new inertia image has order $b$. Thus it is cyclic of degree $b$. Repeat with $c\psi,m'$. After compositing with $E$, the fresh ramification place is still spherical before its step because it is outside $S$ and split in the preceding fields. A nontrivial step character ramifies the support there, so it cannot be a self-twist. Cyclic base change stays cuspidal. Polarization and weak attachment follow as in \ref{B9.38b}, where infinitesimal characters also pull back. The final real character is unramified at the target with Frobenius of order $f$. Its field gives the specified completion. For $f=1$ take $M=E$.

Write $\pi_w$ as the irreducible induction of its unlinked cuspidal terms. Evaluate the inducing functions on the finite flag variety to identify first congruence invariants with finite parabolic induction of the invariants of the terms. Each term is depth zero. For a term of rank $d$ let $\vartheta$ be its Green character of $\F_{p^d}^\times$ up to $p$-power conjugacy, as in \ref{B7.31} and \ref{B9.51}. In rank one use its residue character. For a unit character occurring $h$ times among the rank-one terms, the corresponding determinant character of a block of size $h$ occurs in the lower finite induction. After twisting off this character, constants on the flag variety give the trivial representation. Use the finite character of the resulting lower label, with the higher-rank cuspidal constituents, in the trace test.

At a nonsplit step $k'/k$ at the target, inflate the conjugate of the lower label character to the maximal compact and extend by zero. The upper matching function is the inflation of its Shintani lift. Let $\gamma$ be an integral regular norm of unit determinant. In the order $\mathcal O_k[\gamma]$ after unramified scalar extension there is a unit $\delta$ with norm $\gamma$. Norm is onto on the finite residue-field factors. On successive nilpotent or congruence quotients its linear term is the surjective trace of the unramified scalar extension. This gives the lift in the complete order. Both $\delta$ and $\delta^{-1}$ preserve every upper lattice stable under $\gamma,\gamma^{-1}$. A lattice fixed by $\delta$ composed with Frobenius is stable under the norm $\gamma$, hence under $\delta$ and Frobenius separately. It descends to a $\gamma$-stable lattice, and each such descended lattice is fixed by the twisted operator. The stabilizers have the same orbital weights with hyperspecial volume one. Reduction gives the finite Shintani norm on each lattice, and both sums vanish off the integral locus.

The cyclic base-change character identity carries the nonzero lower trace to a nonzero twisted upper trace \cite[Ch.~I, {\S}6]{arthur-clozel-1989}. On upper first congruence invariants decompose the finite group representation into irreducible isotypic components. An orbit of distinct components contributes zero because a group element acts within the components and the Frobenius intertwiner permutes them. On a stable component Schur's lemma gives a multiple of its twisted irreducible character. Twisted character orthogonality for the Shintani norm selects a stable constituent pairing nontrivially with the transferred lower label. Finite parabolic induction expresses the label character using the Green orbit lift on each block. The same orthogonality gives nonzero pairing with the lifted label, which retains the higher-rank Green terms whose orbits have yet to split \cite[(1.1) and Addendum]{silberger-zink-2008}. Use this label at the next step. At a split step use the identical component.
\leanref{UpperTheory.\allowbreak{}Gal.\allowbreak{}InertialTypesAlt.\allowbreak{}lblOn}{Development.\allowbreak{}AutomorphicArithmetic.\allowbreak{}Foundations.\allowbreak{}GaloisRepresentations.\allowbreak{}RegularSemisimpleFibres.\allowbreak{}PLblOn}{36--65}

At total residue degree $f=a!$, every $d\le a$ divides $f$. Thus every original Green orbit splits into singleton labels
\begin{equation}\label{eq:p9c1-green-lift}
 \vartheta^{p^i}\circ N_{\F_{p^f}/\F_{p^d}},
 \qquad 0\le i<d .
\end{equation}
The rank-one determinant blocks lift by norm on the determinant. Norm surjectivity and \ref{B9.51} preserve the original orbit lengths under $p$-Frobenius. The higher-rank orbits remain distinct by inertial inequivalence. The resulting blocks have distinct unit characters. Finite Frobenius reciprocity gives a vector for the inverse image of the upper parabolic transforming by the product of these determinant characters. On the Iwahori its diagonal unit character is a depth-zero minimal type. The minimal-type theorem puts the upper representation in an induction from smooth characters extending those unit characters \cite[Thm.~6.11(2)]{moy-prasad-1996}.

Group the generic upper segments by their unit characters. Each segment has constant unit character, and Whittaker heredity gives genericity of each group \cite{zelevinsky-1980}. A double coset in finite parabolic induction that pairs with the determinant vector must match unit characters on every intersection of Levi blocks, by varying a single diagonal unit there. Since the block characters differ, whole blocks must match. Each group thus has its determinant vector. Twisting by the inverse character gives a spherical generic block. A segment of length greater than one would contribute a Steinberg factor with zero hyperspecial invariants. The induced hyperspecial invariants would then vanish. Hence the segments are singletons, and genericity excludes linked singletons \cite[Thm.~9.7]{zelevinsky-1980}. Twisting back preserves unlinking within each group, and the unit characters distinguish different groups.

For $t\in I_K=I_L$ put $z=\overline{b_L(t)}$. For each rank $d$ take an intermediate field $E'$ of the cyclic extension $M/E$ with completion the degree-$d$ unramified subextension of $L/K$. If the global group has order $m_0$ and generator $\gamma_0$, the decomposition group is generated by $\gamma_0^{m_0/f}$. The subgroup of index $(m_0/f)d$ intersects that decomposition group in its subgroup of index $d$, so its fixed field works. Reduction $\mathcal O_{E'}\twoheadrightarrow\F_{p^d}$ supplies the residue representatives. The uniformizer is $p$, and the reciprocity value of inertia has residue $N_{\F_{p^f}/\F_{p^d}}z$ there. With these data \ref{B9.49} gives the lower rank-$d$ contribution
\[
 \prod_{i=0}^{d-1}\bigl(X-\vartheta((N_{\F_{p^f}/\F_{p^d}}z)^{p^i})\bigr).
\]
For rank one use the character law of \ref{B9.47}. Upstairs \eqref{eq:p9c1-green-lift} gives the same factors. Multiply over the support with multiplicities. The normalizing twists are trivial on inertia. Choose $R'$ by \ref{B9.46} and apply its rank-one and segment laws.
\end{proof}

\subsection*{Proof of \ref{B9.39}}\label{prf:B9.39}\phantomsection
\begin{proof}
Write $e$ for the doubling embedding and $a$ for the idèlic modulus of the Levi determinant in its Iwasawa decomposition. The height is $h(g)=a(e(g,1))$. The identity
\[
 e(g,1)=e(g,g)e(1,g^{-1})
\]
gives $a(e(g,1))=a(e(1,g^{-1}))$, since $e(g,g)$ stabilizes the isotropic diagonal and acts on it by $g$, whose determinant has modulus one.
\leanref{UpperTheory.\allowbreak{}TP.\allowbreak{}hgt\_\allowbreak{}swap\_\allowbreak{}pk}{Development.\allowbreak{}AutomorphicArithmetic.\allowbreak{}Foundations.\allowbreak{}GaloisRepresentations.\allowbreak{}ThetaParseval.\allowbreak{}ThetaKernels.\allowbreak{}TPGHgtFst}{18--27}

For $e(1,u)$ with $u=g^{-1}$, the inverse compact factor applied to the diagonal has blocks $T$ and $T'=u^{-1}T$, with $h(g)=|\det T|$. Put $\lambda=N_{E/F}(\det T)$ in the adèles of $F$. Unitarity gives $N_{E/F}(\det T')=\lambda$. The adjugate identities
\[
 \lambda u=\overline{\det T'}\,T\operatorname{adj}(T'),\qquad
 \lambda u^{-1}=\overline{\det T}\,T'\operatorname{adj}(T)
\]
put both scaled matrices in a fixed compact adelic set. The determinant norm $\lambda$ also belongs to a fixed compact set.

Suppose $|\lambda|\ge2^{-j}$. Clear the finite denominators of the compact set of determinant norms by a fixed integer. The product of the inverse finite absolute values of the resulting integral idèle bounds an integer $m$ that clears the finite denominators of its inverse. The bounded archimedean components and $|\lambda|\ge2^{-j}$ give $1\le m\le2^{j+i_0}$ for a fixed $i_0$. Let $\varrho$ have component $1/2$ at each real place and $1$ at the finite places. With $l=m\varrho^{2(j+i_0)}$, both $l$ and $l/\lambda$ range over fixed compact sets, after increasing $i_0$. Indeed, $m$ clears the inverse finite denominators after the fixed initial scaling, and the inverse real components after that scaling are $O(2^j)$. Thus the real factor $\varrho^{2(j+i_0)}$ bounds the real components of $l/\lambda$ uniformly. It follows that $lu$ and $lu^{-1}$ lie in a fixed compact set.

For each $m$, cover each coordinate of $lu$ by translates of $l^2(V-V)$, where $V$ is a fixed compact neighbourhood of zero in the adèles of $F$. A maximal family of disjoint translates of $l^2V$ with centres in the fixed coordinate set gives such a cover. Their union lies in a fixed compact set because $l$ ranges over a compact set. The Haar volume of $l^2V$ is $|\varrho|^{4(j+i_0)}\operatorname{vol}(V)$ by the product formula for $m$. Hence the number of coordinate cells is bounded by $C_0C_1^j$ in the fixed matrix dimension. If $u,u'$ lie in the same cell, write $l(u'-u)=l^2D$ with $D$ in a fixed compact set. Then $u^{-1}u'=1+(lu^{-1})D$ lies in a fixed compact matrix set, and the same calculation bounds its inverse. Each occupied cell is therefore contained in a translate of a fixed compact subset of the unitary group. Haar invariance bounds its volume uniformly. Summing over $m$ still gives an exponential bound in $j$. In particular for $R=2$ and some $C,b>0$,
\begin{equation}\label{eq:p9c1-doubling-height}
 \operatorname{vol}\{g:R^{-j-1}<h(g)\le R^{-j}\}\le CR^{bj}
 \qquad(j\ge0).
\end{equation}
The compact factor also bounds $h$ above, and the region with $h>1$ has finite volume by the estimate at $j=0$ for $|\lambda|\ge1$. For $t>b$, the contribution from the layers in \eqref{eq:p9c1-doubling-height} is bounded by the sum of $CR^{(b-t)j}$. Thus $\int h^t<\infty$ for all $t\ge t_0$ with any fixed $t_0>b$. Choose an integer $N>2t_0$ and set $d_0=\max(N,2n)$.
\leanref{GlobalTheory.\allowbreak{}ThK.\allowbreak{}Dbl.\allowbreak{}hgtInt\_\allowbreak{}of\_\allowbreak{}k}{Development.\allowbreak{}AutomorphicArithmetic.\allowbreak{}Foundations.\allowbreak{}GaloisRepresentations.\allowbreak{}ThetaParseval.\allowbreak{}ThetaKernels.\allowbreak{}ThkBaseHgtInt}{17--74}
\leanref{GlobalTheory.\allowbreak{}ThK.\allowbreak{}Dbl.\allowbreak{}Hgt.\allowbreak{}count\_\allowbreak{}ad}{Development.\allowbreak{}WeilBound.\allowbreak{}ThetaLifts.\allowbreak{}LatticeBasics.\allowbreak{}CompactBound}{135--194}

In the polynomial model the source partition has entries $\lambda_{v,k}-(d+m_W(v))/2$. Append $d-n$ zero entries and add $(n+m_X(v))/2$ to each entry. This gives the displayed target weight. Apply \ref{B9.40} to the source archimedean summand. Projection of the nonzero theta image to an irreducible constituent gives $U'$. At each split embedding the global pairing uses the split Weil model. Fixing data in the other local factors of a nonzero pairing gives the stated pairing at the chosen place.
\end{proof}

\subsection*{Proof of \ref{B9.40}}\label{prf:B9.40}\phantomsection
\begin{proof}
Apply \ref{B8.12} to the doubled section of $\Phi\otimes\bar\Phi$ with definite target and splitting character of restriction $\eta_{E/F}^d$. At $s_0=(d-n)/2$ the Eisenstein value is the theta average with factor $1$ for $d>n$ and $2$ otherwise \cite[Thm.~7.1]{gan-qiu-takeda-2014}. With doubling embedding $\epsilon$ and corresponding Eisenstein series $E(s,\cdot)$, pair by
\[
 \int_{[G]^2}E(s,\epsilon(g_1,g_2))\,\chi_W(\det g_2)^{-1}
 \varphi(g_2)\overline{\varphi(g_1)}\,dg_1\,dg_2.
\]
The inverse character cancels the splitting character in the doubled action. Unfolding far to the right gives a nonzero constant times $\int h^{s-s_0}C_\Phi\overline{M_\varphi}$. The section has Siegel exponent $s_0+n/2=d/2$ and its restriction to the compact Iwasawa factor is bounded. With Cauchy--Schwarz this gives
\begin{equation}\label{eq:p9c1-theta-coeff}
 |C_\Phi(g)|\le C h(g)^{d/2},\qquad
 |M_\varphi(g)|\le M_\varphi(1).
\end{equation}
Height integrability gives holomorphy for $\Re(s-s_0)+d/2>t_0$. At $s_0$ Fubini on the compact quotients identifies the paired theta average with a multiple of the squared lift norm.

For $d\ge2n$, let $U_{\rm int}$ be the subgroup integral at all finite places. For finite $S$ containing the bad levels of $\varphi$ and the ramification of the splitting data, put
\[
 A_S=\{g:g_v\text{ is integral for every }v\in S\}\setminus U_{\rm int}.
\]
The sets decrease to the empty set as $S$ increases. Continuity from above for the finite measure $h^{d/2}\,dg$ gives $S$ with $\int_{A_S}h^{d/2}<\operatorname{vol}(U_{\rm int})$. Decompose Haar measure into the factors at the real places, at the finite places in $S$, and off $S$. In both integrals the first two factors range over the same compact real and integral subgroups. The height is invariant under right translation by these subgroups. Cancel their common positive volume and let $K$ be the integral subgroup in the off-$S$ adèles. On this factor the inequality becomes
\begin{equation}\label{eq:p9c1-theta-tail}
 \int_{K^c}h^{d/2}<\operatorname{vol}(K).
\end{equation}

Use the usual vacua off $S$. At $v\in S$, let $\Lambda$ be the self-dual lattice for the local Heisenberg pairing and let $\mathcal A$ be the integral matrix order acting on it. Take an integral full-rank isotropic matrix $x_1\in\Lambda$ with a right inverse $r$ such that $\Lambda r\subset\mathcal A$. Such a matrix comes from the integral hyperbolic pairs in the target lattice in the range $d\ge2n$, by taking one isotropic vector in each of $n$ pairs and using the dual vectors to read the $n$ coordinates. Let $s=\varpi^k$ for a local uniformizer $\varpi$ and put $x_0=s^{-1}x_1$. Let $H_v$ consist of local source isometries $g$ with $g-1,g^{-1}-1\in s\mathcal A$. For large $k$ this subgroup is contained in the stabilizers of the vacuum and of $\varphi$. Translate the vacuum by $x_0$ in the Heisenberg action and denote the resulting vector by $b_v$. If $g=1+sA\in H_v$, then $gx_0=x_0+Ax_1$. The translation by $Ax_1$ lies in $\Lambda$ and has vacuum character one, since the Gram matrix of $x_1$ is zero. Its Heisenberg phase with $x_0$ is also one. Thus $H_v$ fixes $b_v$.

The characters of $b_v$ and $\omega(g)b_v$ on $\Lambda\cap g\Lambda$ agree only if $g^{-1}x_0-x_0$ belongs to the dual of $\Lambda\cap g^{-1}\Lambda$. The quadratic vacuum phases agree on this intersection by invariance of the pairing. The dual is $\Lambda+g^{-1}\Lambda$, by self-duality of $\Lambda$ and preservation of the pairing by $g$. Hence equality gives $g^{-1}x_0-x_0=a+g^{-1}b$ for some $a,b\in\Lambda$. Multiply by $s$ and then on the right by $r$ to obtain
\[
 g^{-1}(x_1-sb)=x_1+sa,\qquad
 g^{-1}(I-sbr)=I+sar.
\]
For sufficiently large $k$ the matrices $I-sbr$ and $I+sar$ are units of $\mathcal A$ congruent to $I$ modulo $s\mathcal A$. It follows that both $g^{-1}$ and $g$ are congruent to $I$ there, so $g\in H_v$. For $g\notin H_v$, a lattice translation therefore acts on the two vectors through distinct unitary characters, and their inner product vanishes. For $g\in H_v$ the coefficient is $\|b_v\|^2$. The local integral supplies $\|b_v\|^2\operatorname{vol}(H_v)>0$. At infinity Schur orthogonality supplies a nonzero factor from the common vector.
\leanref{UpperTheory.\allowbreak{}TP.\allowbreak{}bumpAny\_\allowbreak{}k}{Development.\allowbreak{}AutomorphicArithmetic.\allowbreak{}Foundations.\allowbreak{}GaloisRepresentations.\allowbreak{}ThetaParseval.\allowbreak{}ThetaKernels.\allowbreak{}TPGBumpAny}{182--329}

Divide by these factors. On $K$ the normalized vacuum coefficient is one and $\varphi$ is fixed. The contribution is $M_\varphi(1)\operatorname{vol}(K)$. Off $K$ the normalized vacuum coefficient is bounded by $h^{d/2}$ by the lattice intersection volume in the doubled section. By \eqref{eq:p9c1-theta-coeff} the absolute value of the remaining contribution is at most $M_\varphi(1)\int_{K^c}h^{d/2}$. Inequality \eqref{eq:p9c1-theta-tail} proves nonvanishing.
\end{proof}

\subsection*{Proof of \ref{B9.40a}}\label{prf:B9.40a}\phantomsection
\begin{proof}
Let $\nu=\chi_W^{-1}$ and $b=\chi_X$ at the split place. Convert the conjugate-linear input of the pairing to a linear input by taking the unitary dual of the source. In the split Schwartz model the candidate for the smooth dual of the upper component is normalized induction from the dual of the modified source and the rank-$t$ character $b^{-1}\circ\det$. Here the modified source is the linear source composed with inverse transpose and twisted by $(\nu b)\circ\det$.

To obtain the intertwining map, put $d=n+t$ and filter the Schwartz functions by vanishing on matrices of rank below $j$. Suppose the map kills functions supported on the locus of rank greater than $j$, where $j<n$. In a small neighbourhood of the standard rank-$j$ matrix, the remaining average is indexed by the $j$-by-$(n-j)$ matrices $\rho$ over the residue field. The source changes of variables $g_\rho$ have determinant $\varpi^{-(n-j)}$. Their normalized Weil scalar $s_j$ has modulus $q^{d(n-j)/2}$, from the determinant modulus to the power $d/2$ and the unitary splitting character. On functions vanishing at ranks below $j$, the sum of the inverse-normalized translates agrees with the original function modulo functions supported at ranks greater than $j$. Thus each vector $v$ obtained by evaluating the pairing there satisfies
\[
 v=s_j^{-1}\sum_\rho \sigma(g_\rho)v,\qquad
 \|v\|\le q^{j(n-j)-d(n-j)/2}\|v\|,
\]
where $\sigma$ is the unitary source action in the linear model. The exponent is negative since $d>2n$ and $j<n$. Hence $v=0$. Translate these standard neighbourhoods by the two general linear groups and use a finite partition of the support to obtain the same vanishing for functions supported at ranks at least $j$. If the full-rank restriction of the pairing vanished, descent from $j=n-1$ to $j=0$ would make the whole pairing zero.
\leanref{PositivityTheory.\allowbreak{}Rg.\allowbreak{}Ming.\allowbreak{}core}{Development.\allowbreak{}WeilBound.\allowbreak{}RepresentationTheoryGL.\allowbreak{}PointAveraging.\allowbreak{}MinkowskiBalls.\allowbreak{}ValLeOneOfLc}{407--460}
\leanref{PositivityTheory.\allowbreak{}Rg.\allowbreak{}Ming.\allowbreak{}step}{Development.\allowbreak{}WeilBound.\allowbreak{}RepresentationTheoryGL.\allowbreak{}PointAveraging.\allowbreak{}MinkowskiBalls.\allowbreak{}MemBallMSelf}{134--194}

On the full-rank orbit the stabilizer has Levi factors of ranks $n$ and $t$. Its action on the full-rank evaluation gives the indicated determinant twists, so Frobenius reciprocity gives a nonzero map into the indicated normalized induction. The unitary factors give irreducible induction, which identifies it with the upper dual. Taking its contragredient removes the dual on the modified source and replaces $b^{-1}\circ\det$ by $b\circ\det$. The unitary dual and inverse transpose cancel on the source Satake multiset, so the source contributes $(xy^{-1})B$. The rank-$t$ determinant character contributes $xq^{(t-1)/2-j}$ for $0\le j<t$ by the half modulus of the upper triangular Borel. Normalized induction takes the union of the two multisets.
\leanref{UpperTheory.\allowbreak{}Lit.\allowbreak{}HoweDuality.\allowbreak{}lift\_\allowbreak{}of\_\allowbreak{}irr}{Development.\allowbreak{}AutomorphicArithmetic.\allowbreak{}ThetaCorrespondence.\allowbreak{}FockModels.\allowbreak{}MingZel}{31--44}
\end{proof}

\subsection*{Proof of \ref{B9.41}}\label{prf:B9.41}\phantomsection
\begin{proof}
Let $d_0$ be the bound of \ref{B9.39}, take $t=2(d_0+n)+1$, and let $b$ be a lower bound for the source weights at the chosen embeddings. By \ref{B9.36c} take $\chi$ with $S=\{w_0\}$ and infinity exponent $-1$ at each chosen embedding. Choose a positive odd $K$ with $-K\le2b-(n+t)$. For $u=0,1$ put $\chi_{X,u}=\chi^{2u}$ and $\chi_{W,u}=\chi^{K+2u}$. Their restrictions have the required parity and their exponents meet the inequality of \ref{B9.39}. Apply that proposition with $d=n+t>2n$.

Write $|\cdot|_{\mathbf A_E}$ for the idèlic norm. Define $\kappa=\chi^K|\cdot|_{\mathbf A_E}^{t/2}$ and
\[
 \psi_{u+1,j}=\chi^{2u}|\cdot|_{\mathbf A_E}^{-(n+2t-2)/2+j},
 \qquad 0\le j<t .
\]
The integer powers in the added characters and the parity of $K,t$ in $\kappa$ give algebraic infinity types. At a split place, unramifiedness of the added characters implies that of $\chi_{X,u}$, and unramifiedness of $\kappa$ then gives that of $\chi_{W,u}$. Put $x=\chi_{X,u}(w)$, $y=\chi_{W,u}(w)$ and $q=q_w$. Scale the multiset in \ref{B9.40a} by $q^{(n+t-1)/2}$. The source contribution is $(xy^{-1}q^{t/2})A_n(w,B)$ and the added roots are $xq^{(n+2t-2)/2-j}$. These are the stated character values. At conjugate embeddings the exponent difference for an added line is zero for $u=0$ and nonzero for $u=1$, since the norm has equal exponents and $\chi^2$ has opposite nonzero exponents. The families are disjoint. They and $\kappa$ are unramified at $w_0$ by the choice of $\chi$. To get the displayed jumps, subtract the weights of \ref{B9.39} from their indices. Use the dual weights at conjugate embeddings.
\end{proof}

\subsection*{Proof of \ref{B9.44}}\label{prf:B9.44}\phantomsection
\begin{proof}
Apply odd-rank compatibility to both lifts to obtain $r_i$. Algebraic character realization is injective, so the realizations of the line families remain disjoint. The cofinite Satake unions give \ref{B9.42}. For a Satake multiset at $w_0$, the first lift and its lines have the local spherical data there by \ref{B9.41}. Subtract by \ref{B9.43}, and twist by $\kappa_\ell$ to cancel $\kappa(w_0)^{-1}$ and the jump shift $m_{\iota\circ\tau}$. The same subtraction and twist at good places give weak compatibility.
\end{proof}

\section*{Proofs for Chapter~\ref{ch:weak-local-langlands}: A weak local Langlands correspondence}

\subsection*{Proof of \ref{B9.46}}\label{prf:B9.46}\phantomsection
\begin{proof}
Extend \(r\) of \ref{B9.53} to \(R\) by \ref{B9.54}. The real twist law identifies the Weil constituents of the block for \([u,m]\) with \(R(u\nu_K^j)\), \(0\le j<m\). For the cuspidal support multiset \(s_x\) and \(\sigma\in W_K\),
\[
 \det(X-\rho_{R(x)}(\sigma))
 =\prod_{v\in s_x}\det(X-\rho_{R(v)}(\sigma)).
 \tag{1}\label{eq:p9c2-support}
\]
The same formula holds for any map in \ref{B9.47}. Degrees give the dimension law, and \(\omega_x=\prod_{v\in s_x}\omega_v\), with \(\omega_x\) the central character of \(x\), gives both determinant laws.

Twisting the first terms gives the twist laws. If \(x\) has good cuspidal support in the sense of Definition~\ref{B9.48}, the dual first term \((u\nu_K^{m-1})^\vee\) has parameter \(r(u)^\vee(1-m)\). The centering shift and self-duality of \(\operatorname{Sp}_m\) give the dual block.

If \(R(x)=R(y)\) with \(y\) good, uniqueness of the indecomposable decomposition matches their blocks \cite[\S4.4]{rodier-1982}. The block \(B\) of \([u,m]\) has nilpotence index \(m\) and Weil action \(r(u)\) on \(B/NB\). By \ref{B9.53} the matched first terms coincide, so \(x=y\).

The compact picture identifies the \(K_1\)-invariants of parabolic induction with finite parabolic induction of the terms' invariants. Exactness gives depth zero for each cuspidal term when \(x^{K_1}\ne0\). For good \(x\), (1), \ref{B9.53} and semisimplicity give tameness. The cuspidal laws pass to \(R\). The type law follows because an inflated cuspidal type forces cuspidality \cite[Prop.~6.6]{moy-prasad-1996}.
\end{proof}

\subsection*{Proof of \ref{B9.51}}\label{prf:B9.51}\phantomsection
\begin{proof}
Put \(n=d(u_0)\). The depth-zero cuspidals are compact inductions of \(\widetilde\tau(zk)=\omega(z)\tau(\bar k)\) from \(K^\times K_0\), with \(\tau\) finite cuspidal and \(\omega\) extending its central character \cite[Props.~6.6 and 6.8]{moy-prasad-1996}. For a Cartan representative \(g=\operatorname{diag}(\varpi^{a_j})\), a gap \(a_i<a_{i+1}\) puts the upper unipotent radical of block sizes \(i,n-i\) in the reduction of \(K_0\cap gK_1g^{-1}\). Cuspidality forces invariant functions to vanish on such cosets. Evaluation at \(1\) gives \(u_0^{K_1}\simeq\tau\).

Take the Green label \(\theta\) of \(\tau\) by \ref{B7.31} and a generator \(\zeta\) of \(k_n^\times\). For a proper subfield of degree \(d\), we have \(d\mid n\) and \(d\le n/2\). Its multiplicative group has index
\[
 (p^n-1)/(p^d-1)\ge 1+2^{n-d}>n.
 \tag{4a}\label{eq:p9c2-index}
\]
Thus \(\zeta^r\) has degree \(n\) for \(1\le r\le n\). Put \(\alpha_j=\theta(\zeta^{p^j})\). Some \(\sum_j\alpha_j^r\) in this range is nonzero, since vanishing for all \(r\) would give \(\prod_j\alpha_j=0\) by Newton's identities. Let \(g\) be multiplication by this \(\zeta^r\). Both \(g\) and \(cg\) for \(c\in\F_p^\times\) are elliptic. By (3), \(\operatorname{tr}\tau(cg)=\theta(c)\operatorname{tr}\tau(g)\ne0\). Schur's lemma gives central character \(\theta|_{\F_p^\times}\).

Lift a monic irreducible polynomial of degree \(n\) over \(\F_p\) to \(\mathcal O_E[X]\) and adjoin a root for a datum \(\mathscr D\) of \ref{B9.49}. On common inertia \(R_{\mathscr D}(t)=\overline{a'(t)}\) maps onto \(k_n^\times\) with open kernel by reciprocity. Its restriction to the pro-\(p\) group \(P_K\) is trivial since \(p\nmid p^n-1\). Set \(N=0\) and use inertia lines with characters \(t\mapsto\theta(R_{\mathscr D}(t)^{p^j})\), \(0\le j<n\). The Weil relation lets degree-one Frobenius \(\phi\) cycle these distinct lines, so the action is irreducible. The cyclic matrix with wrap coefficient \(b=(-1)^{n-1}\omega_{u_0}(a(\phi))\) has the required determinant and \(\rho(\phi)^n=b\,1\). Let \(h\) be the exponent of the inertia image. For \(\deg\sigma=k\), write \(\sigma^n=\phi^{nk}t\) with \(t\in I_K\). Then \(\rho(\sigma)^{nh}=b^{kh}1\). The polynomial \(X^{nh}-b^{kh}\) has simple roots, so \(\rho(\sigma)\) is semisimple. Norm compatibility gives
\[
 \det\rho(t)=\theta(N_{k_n/\F_p}R_{\mathscr D}(t))=\omega_{u_0}(a(t)).
\]

Choose \(\beta\in\Omega\) with \(\beta^{p^n-1}=p\). In the residue field of \(A\), tame reciprocity gives \(R_{\mathscr D}(t)=\overline{t(\beta)/\beta}\) for every datum \(\mathscr D\) \cite{serre-1979}. Density of \(\Q\) fixes \(i\), so each datum has type \(u_0^{K_1}\). Formula (3) on multiplication by \(\zeta^r\) gives equal first \(n\) power sums for any two labels of that type. By Newton's identities their orbits coincide, so (4) follows.

In a tame irreducible with inertia orbit length \(\ell\), take a \(\phi^\ell\)-eigenline in one inertia eigenspace. Its \(\phi\)-translates span an invariant space of dimension \(\ell\), so the inertia characters have multiplicity one. Their orbit, read from the inertia polynomials, and the Frobenius determinant determine the cyclic action up to diagonal conjugacy. Conversely, surjectivity of \(R_{\mathscr D}\) and \(a\) recovers the Green orbit and \(\omega_{u_0}\) from inertia and the determinant. These determine the compact induction.

A twist by \(\lambda\) leaves the Cartan unipotent action unchanged. The identity coset contributes exactly when \(\lambda\) is trivial on \(\det K_1=1+\mathfrak p_K\). Then the label gains \(\bar\lambda\circ N_{k_n/\F_p}\), with \(\bar\lambda\) the reduction on units, and the central character gains \(\lambda^n\). Thus \(V_{u_0\otimes\lambda}=V_{u_0}\otimes(\lambda\circ a)\). Any two depth-zero bases of one class differ by such a twist. In rank one the determinant gives \(V_{u_0}=(u_0\circ a,0)\). These identities give the assignment \(r_0\) and its twist law.
\end{proof}

\subsection*{Proof of \ref{B9.52}}\label{prf:B9.52}\phantomsection
\begin{proof}
Let \(L/K\) be unramified of degree \(n\), with residue field \(k_n\). For a uniformizer \(\varpi\) of \(K\) and \(M=q^n-1\), put
\[
 b(u)=\overline{\frac{u^M-1}{M\varpi}},\qquad u\in\mathcal O_L^\times .
\]
Multiplication modulo \(\varpi^2\) gives additivity and \(b(1+\varpi x)=\bar x\). Thus \(\pi=b\circ a_L|_{I_K}\) maps onto \(k_n^+\) with open kernel, where \(a_L\) is reciprocity and \(I_K=I_L\). For degree-one Frobenius \(\phi\),
\[
 \pi(\phi t\phi^{-1})=f(\pi(t)),\qquad f(z)=z^q,
 \tag{5}
\]
since \(\phi\) fixes \(\varpi\) and induces \(f\) on \(k_n\). For \(\lambda\in k_n\) of degree \(n\) over \(\F_q\), the trace pairing gives orbit length \(n\) for \(\alpha(z)=\exp(2\pi i\operatorname{Tr}_{k_n/\F_p}(\lambda z)/p)\).

Let inertia act on \(e_j\) by \(\eta_j=\alpha\circ f^j\circ\pi\), \(0\le j<n\). Put \(\rho(\phi)e_j=e_{j-1}\) for \(j>0\) and \(\rho(\phi)e_0=(-1)^{n-1}e_{n-1}\). Equation (5) gives the Weil relation. This finite-order matrix of determinant \(1\) normalizes the finite inertia image, so the Weil image is finite. Transitivity on the distinct inertia lines gives irreducible \(J_n=(\rho,0)\) with unitary determinant.

The \(p\)-group \(k_n^+/\pi(P_K)\) is a quotient of tame inertia, hence trivial. Choose \(t\in P_K\) with \(\eta_0(t)\ne\eta_1(t)\). Thus \(\rho_{J_n}(t)\) is nonscalar. Every character twist of a tame parameter acts at \(t\) by a scalar.
\end{proof}

\subsection*{Proof of \ref{B9.53}}\label{prf:B9.53}\phantomsection
\begin{proof}
Smooth central characters are unitary on units, so \(e(u)\) exists. Twists preserve goodness. The central character of \(u\nu_K^s\) is \(\omega_u\nu_K^{ns}\), hence
\[
 e(u\nu_K^s)=e(u)+s.
 \tag{7}
\]
This gives the real twist law on the second branch, where \(n\ge2\). Dimension and irreducibility follow from \ref{B9.52}, and
\[
 |\det\rho_{J_n(e(u))}(\sigma)|=q^{ne(u)\deg\sigma}=|\omega_u(a(\sigma))|.
 \tag{8}
\]
On good classes \ref{B9.51} gives the twist and determinant laws. It also gives the type law. Invert the label and central character of a depth-zero base, or the character in rank one, to get duality. Depth-zero good cuspidals have tame parameters by \ref{B9.51} in higher rank and by reciprocity in rank one.

Suppose \(r(u)=r(v)\) with \(v\) good. Their ranks agree. In rank one use surjectivity of \(a\). For good \(u=u_0\otimes\chi\) and \(v=v_0\otimes\chi'\) of higher rank, \(V_{u_0}=V_{v_0}\otimes\xi\), where \(\xi=\lambda\circ a\) and \(\lambda=\chi'\chi^{-1}\). Tameness makes \(\xi\) trivial on \(P_K\). Frobenius conjugation on \(I_K/P_K\) is the \(q\)-power map, so the character \(\xi\) satisfies \(\xi(t)^{q-1}=1\) on \(I_K\). Surjectivity and the degree formula give \(a(I_K)=\mathcal O_K^\times\). The image \(\lambda(1+\mathfrak p_K)\) is a finite \(p\)-group by smoothness, and its exponent divides \(q-1\). Thus it is trivial. By \ref{B9.51}, \(V_{u_0}=V_{v_0\otimes\lambda}\), so \(u_0=v_0\otimes\lambda\) and \(u=v\).

If \(u\) is on the second branch, \(v=v_0\otimes\chi\) over \(K\simeq\Q_p\), and equality would give \(J_{d(u)}=V_{v_0}\otimes((\chi\circ a)\nu_W^{-e(u)})\), contrary to \ref{B9.52}.
\end{proof}

\subsection*{Proof of \ref{B9.55a}}\label{prf:B9.55a}\phantomsection
\begin{proof}
Genericity gives an unlinked multisegment \cite[Thm.~9.7]{zelevinsky-1980}. Zero monodromy in (2) forces singleton segments, so \(R(x)=\bigoplus_{u\in s_x}R(u)\). The commutator test and semisimplicity give trivial commutator action. Each \(R(u)\) is then abelian and irreducible, hence one-dimensional.

In the tame case fix a support term \(u\) of rank \(d\) and a Frobenius \(\phi\) of degree \(-1\). Put \(D=\det\rho_{R(u)}(\phi)\). The cyclic description in \ref{B9.51} injects tame irreducibles of rank \(d\) and determinant \(D\) at \(\phi\) into length-\(d\) Frobenius orbits of inertia characters. These factor through the quotient of \(I_K/P_K\) of order \(q^d-1\), identified with \(\F_{q^d}^\times\) by reciprocity. For each orbit take its Green type and compactly induce as in \ref{B9.51}, with \(\omega(a(\phi))=D\) on the uniformizer \(a(\phi)\). The \(K_1\)-invariants distinguish these good depth-zero cuspidals. Their images under \(R\) are distinct tame irreducibles of the prescribed rank and determinant by \ref{B9.47}. The resulting injection of the finite orbit set into itself is onto. The depth-zero preimage of \(R(u)\) equals \(u\) by good injectivity. Thus all terms have depth zero. The unlinked induction is irreducible, so the compact picture gives \(x^{K_1}\ne0\).

Inertially equivalent terms have equal inertia polynomials by the good twist law. Their eigenvalues at a lift of a topological generator are repeated for \(R(x)\), hence fixed by \(z\mapsto z^q\). Their inertia characters are Frobenius invariant, so the terms have rank one.
\end{proof}

% ---- appendix chapter p1c5 ----
\chapter{The window criterion in the archive}\label{ch:p1c5}

The archive proves Theorem~\ref{B1.31} by bounding high derivatives of the Laplace transform on the windows. This chapter gives that proof, statement by statement, with the Lean declarations; Chapter~\ref{ch:p1c4} gives a shorter proof of the same theorem. Lemmas~\ref{B1.54} and~\ref{B1.55} are stated in Chapter~\ref{ch:p1c4}.

Let \(g:\R\to\R\) be Borel measurable, with
\(|g(u)|\le C_0e^{2u}\) for \(u\ge0\), where \(C_0\ge0\).
For \(k\in\N\), \(a>0\), and \(u\ge0\), put
\[
 K_{k,a}(u)=\frac{u^ke^{-au}}{k!},\qquad
 B_k(z)=\int_0^\infty g(u)\frac{u^ke^{-zu}}{k!}\,du
       \quad(\Re z>2).
\]
When \((g,\Phi,Z)\) satisfies the pole data of \ref{B1.33}, set
\(A_k(z)=(-1)^k\Phi^{(k)}(z)/k!\).
For a character family \(F\), Definition~\ref{B1.26} gives the
arithmetic example \(g(u)=G_m(F,e^u)\). All integrals are Lebesgue
integrals.

\begin{theorem}[Derivative bound from integral windows]\label{B1.47}
Suppose that \((g,\Phi,Z)\) satisfies \ref{B1.33}.
Fix \(0<\theta<1\) and \(\delta>0\). Assume the logarithmic window
condition \(\mathrm{WH}_U(g;\theta,\delta)\) of Lemma~\ref{B1.29}.
Then there exist \(\sigma>2\) and \(C\in\R\) such that, for every
\(\eta>0\) and \(U_0\in\R\), one can find \(U\ge U_0\) and
\(k\in\N\), \(k\ge U\), for which
\[
 |A_k(\sigma+it)|+\Re A_k(\sigma)
       \le \frac{Ce^{\eta U}}{(\sigma-\theta)^k}
       \qquad\text{for every }t\in\R.
\]
\end{theorem}
\begin{proof}
Take \(\sigma\) and the positive witness \(D\) displayed in
Lemma~\ref{B1.52}. Apply the window condition with lower bound
\(\max(U_0,1)\), and choose \(k\ge U\) from that lemma.
By \ref{B1.49} and \ref{B1.53}, for every real \(t\),
\[
 |A_k(\sigma+it)|+\Re A_k(\sigma)
 \le \frac{2e^{\eta U}}{(\sigma-\theta)^k}
                         +\frac{2C_0D}{\sigma^k}.
\]
Since \(0<\sigma-\theta<\sigma\) and \(e^{\eta U}\ge1\), take
\(C=2+2C_0D\).
\end{proof}

\section{The differentiated Laplace kernel}

\begin{lemma}[Differentiation under the Laplace integral]\label{B1.48}
For \(g\) as above and \(j\in\N\), define
\[
 H_j(z)=\int_0^\infty g(u)(-u)^je^{-zu}\,du
                  \qquad(\Re z>2).
\]
These integrals converge absolutely. For every \(j\in\N\) and
\(s\in\C\) with \(\Re s>2\),
\[
 H_j'(s)=H_{j+1}(s),\qquad H_0^{(j)}(s)=H_j(s).
\]
\end{lemma}
Apply Widder's Laplace--Stieltjes differentiation formula
\cite[Chapter II, \S5]{widder-1941} to the locally absolutely
continuous function \(\alpha(u)=\int_0^u g(v)\,dv\) on \([0,\infty)\).
% [Widder, title page, contents and introduction](https://api.pageplace.de/preview/DT0400.9781400876457_A26058229/preview-9781400876457_A26058229.pdf).

\begin{lemma}[Normalized coefficients]\label{B1.49}
If \((g,\Phi,Z)\) satisfies \ref{B1.33}, then
\(A_k(z)=B_k(z)\) for every \(k\in\N\) and \(\Re z>2\).
In particular, for \(\sigma>2\) and \(t\in\R\),
\[
 \begin{aligned}
 A_k(\sigma+it)&=\int_0^\infty g(u)K_{k,\sigma}(u)e^{-itu}\,du,\\
 \Re A_k(\sigma)&=\int_0^\infty g(u)K_{k,\sigma}(u)\,du.
 \end{aligned}
\]
\end{lemma}
On \(\Re z>2\), the pole data identify \(\Phi\) with \(H_0\), so
\(A_k=B_k\) is Widder's differentiation formula
\cite[Chapter II, \S5]{widder-1941} in Lemma~\ref{B1.48} with the
normalization \((-1)^k/k!\).

\begin{lemma}[Mass and pointwise bound for the kernel]\label{B1.50}
For \(a>0\) and \(k\in\N\), the kernel \(K_{k,a}\) is integrable
on \((0,\infty)\), and
\[
 \int_0^\infty K_{k,a}(u)\,du=a^{-k-1},\qquad
 0\le K_{k,a}(u)\le a^{-k}\quad(u\ge0).
\]
\end{lemma}
The integral is the specialization \(\mu=1\), \(\nu=k+1\), \(z=a\)
of Euler's gamma integral
\cite[Section 5.9(i), equation (5.9.1)]{olver-2010}, divided by \(k!\).
The pointwise bound follows from \((au)^k/k!\le e^{au}\) after
multiplication by \(a^{-k}e^{-au}\).
% [DLMF, gamma integral (5.9.1)](https://dlmf.nist.gov/5.9.E1).

\section{Both tails at a common derivative order}

\begin{lemma}[Exponential bounds for the two tails]\label{B1.51}
For \(k\in\N\), \(0\le\ell<a\), and \(U,V\in\R\),
\[
 \int_{(0,\infty)\setminus[U,V]}K_{k,a}(u)\,du
 \le \frac{e^{-\ell V}}{(a-\ell)^{k+1}}
                +\frac{e^{\ell U}}{(a+\ell)^{k+1}}.
\]
Furthermore, for \(b,\sigma>0\) and \(x,c\in\R\),
\[
 x+k\frac{\sigma-b}{b}\le c
 \quad\Longrightarrow\quad
 \frac{e^x}{b^{k+1}}\le\frac{e^c}{b\sigma^k}.
\]
\end{lemma}
\begin{proof}
If \(u>0\) lies outside \([U,V]\), then \(u<U\) or \(u>V\). Hence
\[
 K_{k,a}(u)\le
 e^{\ell U}K_{k,a+\ell}(u)+e^{-\ell V}K_{k,a-\ell}(u).
\]
By nonnegativity, the integral of this majorant over the complement
is at most its integral over \((0,\infty)\).
Lemma~\ref{B1.50} evaluates the latter.

For the second assertion, \(1+y\le e^y\), with
\(y=(\sigma-b)/b\), gives
\[
 \frac{e^x}{b^{k+1}}
 =\frac{e^x}{b\sigma^k}\left(\frac{\sigma}{b}\right)^k
 \le\frac{\exp\!\left(x+k(\sigma-b)/b\right)}{b\sigma^k}
 \le\frac{e^c}{b\sigma^k}.
\]
\end{proof}

\begin{lemma}[Choosing the order from the window]\label{B1.52}
For every \(\delta>0\), there exist \(\sigma>2\) and \(D\in\R\)
such that, for every real \(U\ge1\), there is \(k\in\N\) with
\(k\ge U\) and
\[
 \int_{(0,\infty)\setminus[U,(1+\delta)U]}
             K_{k,\sigma-2}(u)\,du\le\frac{D}{\sigma^k}.
\]
One may take \(\ell=2+6/\delta\), \(M=\ell(1+\delta)\), and
\(\sigma=M(\ell+2)\), with
\[
 k=\left\lceil(1+\delta/2)\sigma U\right\rceil,\qquad
 D=\frac{e}{\sigma-2-\ell}+\frac1{\sigma-2+\ell}.
\]
\end{lemma}
\begin{proof}
Write \(p=1+\delta/2\) and \(b_\pm=\sigma-2\pm\ell\).
Since \(M>\ell>2\), we have
\(b_-=(M-1)(\ell+2)>0\). Thus \(\sigma>2\) and \(D>0\).
Direct computation gives
\[
 p(\ell+2)=M-1,\qquad p(\ell-2)=\ell+1,\qquad
 p\sigma=M(M-1).
\]
The ceiling satisfies \(p\sigma U\le k<p\sigma U+1\), so \(k\ge U\).
By \ref{B1.51}, with \(a=\sigma-2\), the tail integral is at most
\(T_++T_-\), where
\[
 T_+=\frac{e^{-MU}}{b_-^{k+1}},\qquad
 T_-=\frac{e^{\ell U}}{b_+^{k+1}}.
\]
For the upper tail, \((\sigma-b_-)/b_-=1/(M-1)\). The upper
ceiling bound gives
\[
 \begin{aligned}
 -MU+k\frac{\sigma-b_-}{b_-}
 &=-MU+\frac{k}{M-1}\\
 &<-MU+\frac{p\sigma U+1}{M-1}
   =\frac1{M-1}\le1.
 \end{aligned}
\]
The second bound of \ref{B1.51}, with \(x=-MU\) and \(c=1\),
gives \(T_+\le e/(b_-\sigma^k)\).

For the lower tail, \((\sigma-b_+)/b_+=-(\ell-2)/b_+\).
Since \(\sigma\ge\ell(\ell-2)\), we have
\(\ell b_+\le(\ell+1)\sigma\). The lower ceiling bound gives
\[
 \frac{k(\ell-2)}{b_+}
 \ge\frac{p\sigma U(\ell-2)}{b_+}
 =\frac{(\ell+1)\sigma U}{b_+}\ge\ell U.
\]
Thus \(\ell U+k(\sigma-b_+)/b_+\le0\).
Apply the same bound with \(x=\ell U\) and \(c=0\) to obtain
\(T_-\le1/(b_+\sigma^k)\). Adding proves the claim.
\end{proof}

\section{The one-sided estimate}

\begin{lemma}[Combining the window with its complement]\label{B1.53}
Let \(g\) satisfy the growth and measurability assumptions above.
Let \(\sigma>2\), \(\theta<\sigma\), \(U>0\), and \(V\in\R\),
and put \(I=[U,V]\). Suppose that \(R:\R\to\R\) satisfies
\(R(u)\ge\max(g(u),0)\) on \(I\), and that
\(R(u)e^{-\theta u}\) is integrable on \(I\).
Then, for every \(k\in\N\) and \(t\in\R\),
\[
 \begin{split}
 |B_k(\sigma+it)|+\Re B_k(\sigma)
 \le{}&\frac{2}{(\sigma-\theta)^k}
                 \int_I R(u)e^{-\theta u}\,du\\
 &+2C_0\int_{(0,\infty)\setminus I}K_{k,\sigma-2}(u)\,du.
 \end{split}
\]
\end{lemma}
\begin{proof}
Set \(w(u)=K_{k,\sigma}(u)\).
The bound \(|g(u)|w(u)\le C_0K_{k,\sigma-2}(u)\) for \(u>0\),
together with \ref{B1.50}, gives absolute convergence of the
integrals below. Since \(|e^{-itu}|=1\), the triangle inequality gives
\[
 \begin{aligned}
 |B_k(\sigma+it)|+\Re B_k(\sigma)
 &\le\int_0^\infty |g(u)|w(u)\,du+\int_0^\infty g(u)w(u)\,du\\
 &=\int_0^\infty (|g(u)|+g(u))w(u)\,du.
 \end{aligned}
\]
On \(I\), the envelope assumption and \ref{B1.50} give
\[
 \begin{aligned}
 |g(u)|+g(u)&=2\max(g(u),0)\le2R(u),\\
 w(u)&=K_{k,\sigma-\theta}(u)e^{-\theta u}
       \le(\sigma-\theta)^{-k}e^{-\theta u}.
 \end{aligned}
\]
Since \(R\ge0\) and \(w\ge0\) on \(I\), multiplying these
inequalities and integrating gives the first term.
On the complement,
\((|g(u)|+g(u))w(u)\le2|g(u)|w(u)\le2C_0K_{k,\sigma-2}(u)\).
Adding the integral over \((0,\infty)\setminus I\) proves the inequality.
\end{proof}


\section{From the derivative bounds to the exclusion of poles}

\begin{lemma}[Dichotomy at a real centre]\label{B1.56}
Suppose $Z$ satisfies the local condition of Lemma~\ref{B1.55}.
Fix $0<\theta<\sigma$ with $Z\subset\{w:\Re w<\sigma\}$.
If $\Re\rho_0>\theta$ for some $\rho_0\in Z$, at least one of the
following alternatives holds.
\begin{enumerate}
\item There are $D\le\sigma-\theta$ and $\rho_1\in Z$ such that
$Z\cap B(\sigma,D)=\varnothing$ and $\sigma-\Re\rho_1<D$.
\item There are $\beta\in\R$ and $r\in\R$ such that
$\beta\in Z$, $\beta>\theta$, $\sigma-\beta<r\le\sigma-\theta$,
and $Z\cap B(\sigma,r)=\{\beta\}$.
\end{enumerate}
\end{lemma}

\begin{proof}
Set $R=\sigma-\theta$ and $F=Z\cap\overline B(\sigma,R)$.
Lemma~\ref{B1.55} makes $F$ finite. If $F$ is empty, take
$D=R$ and $\rho_1=\rho_0$.

Otherwise choose $\rho_*\in F$ minimising distance to $\sigma$,
and put $D=|\sigma-\rho_*|>0$. Then $D\le R$ and
$Z\cap B(\sigma,D)=\varnothing$. If some $\rho_1\in Z$ has
$\sigma-\Re\rho_1<D$, the first alternative holds.
In the remaining case every $w\in Z$ satisfies
$D\le\sigma-\Re w$. For every minimiser $w\in F$,
\[
 D\le\sigma-\Re w\le|\sigma-w|=D,\qquad
 |\sigma-w|^2=(\sigma-\Re w)^2+(\Im w)^2.
\]
Hence every minimiser equals the real point $\beta=\sigma-D$.
Applying $D\le\sigma-\Re\rho_0$ gives $D<R$, so $\beta>\theta$.
Every point of $F\setminus\{\beta\}$ has distance strictly greater
than $D$. Choose $r$ with $D<r\le R$ no larger than any of these
finitely many distances. Points of $Z$ outside $F$ have distance
greater than $R$. Thus $Z\cap B(\sigma,r)=\{\beta\}$.
\end{proof}

\begin{lemma}[An isolated pole after moving the centre]\label{B1.57}
Suppose $Z$ satisfies the local condition of Lemma~\ref{B1.55},
$\Re w<\sigma$ for all $w\in Z$, and $D<\sigma$.
If $\rho_1\in Z$ satisfies $\sigma-\Re\rho_1<D$, then there exist
$t\in\R$, $\rho\in Z$, and $r\in\R$ such that, for $z=\sigma+it$,
\[
 |z-\rho|<r\le D,\qquad Z\cap B(z,r)=\{\rho\}.
\]
\end{lemma}

\begin{proof}
Put $a=\sigma-\Re\rho_1$, so $0<a<D$, and choose
\[
 \xi=\tfrac12\min\{D-a,\sigma-D\}>0,\qquad
 z_0=\sigma+i\Im\rho_1.
\]
Since $D+\xi<\sigma$, the candidate set
$F=Z\cap\overline B(z_0,D+\xi)$ is finite and contains
$\rho_1$. For distinct $v,w\in F$, equality of distances to
$\sigma+it$ is equivalent to
\[
 2t(\Im v-\Im w)
 =(\sigma-\Re v)^2+(\Im v)^2
  -(\sigma-\Re w)^2-(\Im w)^2.
\]
If $\Im v\ne\Im w$, the displayed equation determines a unique
height at which the distances tie.
If $\Im v=\Im w$, equal distances would imply
$(\sigma-\Re v)^2=(\sigma-\Re w)^2$. Both bases are positive,
so $v=w$. Thus each distinct pair excludes at most one height.
Choose $t\in(\Im\rho_1-\xi,\Im\rho_1+\xi)$ outside this finite
set of heights and put $z=\sigma+it$.
All distances from $z$ to points of $F$ are distinct, and
\[
 |z-\rho_1|\le |z-z_0|+|z_0-\rho_1|
 <\xi+a<D.
\]
Let $\rho$ be the unique nearest point in $F$ and set
\[
 r=\min\bigl(\{D\}\cup\{|z-w|:w\in F\setminus\{\rho\}\}\bigr).
\]
Since $|z-\rho|\le|z-\rho_1|<D$, the choice of $\rho$ gives
$|z-\rho|<r\le D$. If $w\in Z\cap B(z,r)$, then
$|w-z_0|<D+\xi$, so $w\in F$ and $w=\rho$.
\end{proof}

For any $\Phi:\C\to\C$, write $A_k(z)=(-1)^k\Phi^{(k)}(z)/k!$ for
$k\in\N$. In Lemmas~\ref{B1.58} and~\ref{B1.59}, $\Phi$ is arbitrary.

\begin{lemma}[Cauchy's coefficient bound]\label{B1.58}
If $\Phi$ is holomorphic on $B(z,r)$ and $0<r'<r$, then some
$M\in\R$ satisfies
\[
 \frac{|\Phi^{(k)}(z)|}{k!}\le\frac{M}{(r')^k}
 \qquad(k\in\N).
\]
\end{lemma}

Ahlfors gives Cauchy's estimate with the circle maximum.
Taking absolute values inside his integral formula gives $M=(2\pi)^{-1}\int_0^{2\pi}|\Phi(z+r'e^{it})|\,dt$
\cite[Chapter~IV, \S2.3, formulas~(24)--(25)]{ahlfors-1979}.

\begin{lemma}[Coefficients of an isolated simple pole]\label{B1.59}
Let $\rho\in B(z,r)$, $\rho\ne z$, and $c\in\C$.
Suppose $\Phi$ is holomorphic on $B(z,r)\setminus\{\rho\}$ and
agrees with $c/(s-\rho)+h(s)$ on a punctured neighbourhood of
$\rho$, where $h$ is holomorphic at $\rho$.
For every $0<r'<r$ there is $M\in\R$ such that
\[
 \left|A_k(z)-\frac{c}{(z-\rho)^{k+1}}\right|
 \le\frac{M}{(r')^k}\qquad(k\in\N).
\]
\end{lemma}

After subtraction of the simple principal part, the remainder extends
across $\rho$ with value $h(\rho)$. Lemma~\ref{B1.58} applies on
the whole disc \cite[Chapter~IV, \S\S2.3, 3.2]{ahlfors-1979}.
The normalisation follows from
\[
 \frac{d^k}{ds^k}\frac{c}{s-\rho}
   =\frac{(-1)^k k!\,c}{(s-\rho)^{k+1}}.
\]

\begin{lemma}[Incompatible exponential rates]\label{B1.60}
Let $c>0$ and $0<d<e$. For every $C\in\R$ there are
$\eta>0$ and $U_*\in\R$ such that, for every $U\ge U_*$
and every $k\in\N$ with $k\ge U$,
\[
 C\,\frac{\exp(\eta U)}{e^k}<\frac{c}{d^k}.
\]
Consequently, for $e\le E$ and fixed $M,C\in\R$, the assertion
\[
 \forall\eta>0\ \forall U_0\in\R\ \exists U\ge U_0\
 \exists k\in\N,\quad k\ge U,\quad
 \frac{c}{d^k}\le\frac{M}{e^k}
                    +C\frac{\exp(\eta U)}{E^k}
 \tag{*}
\]
is false.
\end{lemma}

\begin{proof}
Write $L=\log(e/d)>0$ and $\eta=L/2$.
After multiplication by $d^k$, the first left-hand side is at most
\[
 |C|\exp(\eta U-kL)\le |C|\exp(-\eta U),
\]
which is less than $c$ for every sufficiently large $U$.
The first assertion with coefficient $|M|+|C|$ gives
$\eta>0$ and $U_*$. Apply $(*)$ with this $\eta$ and
$U_0=\max\{U_*,0\}$. For the resulting $U,k$, the right-hand
side of $(*)$ is at most
$(|M|+|C|)\exp(\eta U)/e^k<c/d^k$, a contradiction.
\end{proof}

Assume the integral window hypothesis of~\ref{B1.29}, with
$0<\theta<1$ and $\delta>0$. By~\ref{B1.47}, there are fixed
$\sigma>2$ and $C\in\R$ such that, for every $\eta>0$ and
$U_0\in\R$, some $U\ge U_0$ and $k\in\N$, $k\ge U$, satisfy
\begin{equation}\label{eq:p1c5-window}
 |A_k(\sigma+it)|+\Re A_k(\sigma)
 \le C\frac{\exp(\eta U)}{(\sigma-\theta)^k}
 \qquad\text{for every }t\in\R.
\end{equation}
Lemma~\ref{B1.54} gives $\Re\rho<\sigma$ for every $\rho\in Z$.

Theorem~\ref{B1.31} follows if both alternatives of
Lemma~\ref{B1.56} are excluded for the fixed $\sigma,C$ in
\eqref{eq:p1c5-window}.

\begin{lemma}[The moving-centre case]\label{B1.61}
Under these hypotheses, the first alternative of
Lemma~\ref{B1.56} is impossible.
\end{lemma}

\begin{proof}
Let $D$ and $\rho_1$ be supplied by that alternative.
Since $D\le\sigma-\theta<\sigma$, Lemma~\ref{B1.57} supplies
$z=\sigma+it$, $\rho\in Z$, and $r\le D$ with
$d=|z-\rho|<r$ and $Z\cap B(z,r)=\{\rho\}$.
Here $d>0$, since $\Re\rho<\sigma$. Set
\[
 e=\frac{d+r}{2},\qquad E=\sigma-\theta,
 \qquad 0<d<e<r\le D\le E.
\]
The discs $B(z,r)$ and $B(\sigma,D)$ lie in $H$ because their centres
have real part $\sigma$ and their radii are at most $\sigma-\theta$.
By Definition~\ref{B1.33}, $\Phi$ is holomorphic on
$B(z,r)\setminus\{\rho\}$ and on $B(\sigma,D)$.
Lemmas~\ref{B1.59} and~\ref{B1.58} give constants
$M_1,M_2$, independent of $k$, with
\[
 \left|A_k(z)-\frac{c_\rho}{(z-\rho)^{k+1}}\right|
 \le\frac{M_1}{e^k},\qquad |A_k(\sigma)|\le\frac{M_2}{e^k}.
\]
The triangle inequality and $\Re A_k(\sigma)\ge-|A_k(\sigma)|$
therefore give
\[
 \frac{|c_\rho|/d}{d^k}
 \le |A_k(z)|+\Re A_k(\sigma)+\frac{M_1+M_2}{e^k}.
\]
For the $U,k$ supplied by~\eqref{eq:p1c5-window}, this is $(*)$
with $c=|c_\rho|/d>0$. Lemma~\ref{B1.60} gives the contradiction.
\end{proof}

\begin{lemma}[The real-pole case]\label{B1.62}
Under the same hypotheses, the second alternative of
Lemma~\ref{B1.56} is impossible. Hence every $\rho\in Z$ satisfies
$\Re\rho\le\theta$.
\end{lemma}

\begin{proof}
Let $\beta\in Z\cap\R$, $\beta>\theta$, and $r$ be supplied by the
second alternative. Put
\[
 d=\sigma-\beta>0,\qquad e=\frac{d+r}{2},\qquad
 E=\sigma-\theta.
\]
Then $d<e<r\le E$, and $B(\sigma,r)\subset H$ contains the single
pole $\beta$. Write $c_\beta>0$ for its real residue.
Lemma~\ref{B1.59} gives a fixed $M$ such that
\[
 \left|A_k(\sigma)-\frac{c_\beta}{d^{k+1}}\right|
 \le\frac{M}{e^k}.
\]
Since $c_\beta/d^{k+1}$ is real, the bound gives
$c_\beta/d^{k+1}-\Re A_k(\sigma)\le M/e^k$.
Adding $|A_k(\sigma)|$ to the right gives
\[
 \frac{c_\beta/d}{d^k}
 \le \Re A_k(\sigma)+|A_k(\sigma)|+\frac{M}{e^k}.
\]
Use~\eqref{eq:p1c5-window} at $t=0$.
For every $\eta>0$ and every lower bound on $U$, the resulting
inequality is $(*)$ with $c=c_\beta/d>0$.
Lemma~\ref{B1.60} gives a contradiction.

If any $\rho_0\in Z$ had $\Re\rho_0>\theta$,
Lemma~\ref{B1.56} would give one of the two excluded alternatives.
This proves Theorem~\ref{B1.31}. For a family $F$ satisfying the
window hypothesis of Theorem~\ref{B1.28}, Lemma~\ref{B1.29} gives
the logarithmic window condition for $g_F$. Apply
Theorem~\ref{B1.31} to the pole data of Theorem~\ref{B1.30}.
A zero with $\theta<\Re s<1$ would lie in $Z_F$ and contradict
the pole bound. Lemma~\ref{B1.32} excludes zeros with $\Re s\ge1$.
\end{proof}


\begin{proof}[Proof of Theorem~\ref{B1.31} along this route]
Suppose \(\rho_0\in Z\), \(\Re\rho_0>\theta\).
Take \(\sigma>2\) from \ref{B1.47}. All poles lie to its left
by \ref{B1.54}. In \ref{B1.56}, the first alternative contradicts
\ref{B1.61}, and the second contradicts \ref{B1.62} by positivity
of the residue at its real pole.
\end{proof}

% ---- lean index ----
\chapter*{Index of Lean declarations}\addcontentsline{toc}{chapter}{Index of Lean declarations}

For each statement of the main text, in the order of the book, the Lean declarations that prove it, with their files and line numbers in the archive.

\section*{Chapter~\ref{ch:p1c1}: Five statements and the critical line}

{\footnotesize\raggedright\noindent Chapter introduction.\quad \texttt{riemannZeta\_\allowbreak{}one} (\texttt{Mathlib.\allowbreak{}NumberTheory.\allowbreak{}Harmonic.\allowbreak{}ZetaAsymp}, lines 418--422)\\
\texttt{riemannZeta\_\allowbreak{}one\_\allowbreak{}ne\_\allowbreak{}zero} (\texttt{Mathlib.\allowbreak{}NumberTheory.\allowbreak{}Harmonic.\allowbreak{}ZetaAsymp}, lines 441--450)\par}

{\footnotesize\raggedright\noindent \ref{B1.01}.\quad \texttt{riemannHypothesis} (\texttt{Main}, lines 4--4)\par}

{\footnotesize\raggedright\noindent \ref{B1.02}.\quad \texttt{riemannZeta\_\allowbreak{}eq\_\allowbreak{}zero\_\allowbreak{}re} (\texttt{Main}, lines 6--6)\\
\texttt{RiemannZetaZerosOnLine} (\texttt{Statements}, lines 3--4)\par}

{\footnotesize\raggedright\noindent \ref{B1.03}.\quad \texttt{LFunction\_\allowbreak{}ne\_\allowbreak{}zero\_\allowbreak{}of\_\allowbreak{}half\_\allowbreak{}lt\_\allowbreak{}re} (\texttt{Main}, lines 8--8)\\
\texttt{DirichletLNonvanishingRightOfHalf} (\texttt{Statements}, lines 6--8)\par}

{\footnotesize\raggedright\noindent \ref{B1.04}.\quad \texttt{LFunction\_\allowbreak{}eq\_\allowbreak{}zero\_\allowbreak{}re} (\texttt{Main}, lines 10--10)\\
\texttt{DirichletLZerosOnLine} (\texttt{Statements}, lines 10--12)\par}

{\footnotesize\raggedright\noindent \ref{B1.05}.\quad \texttt{LFunction\_\allowbreak{}eq\_\allowbreak{}zero\_\allowbreak{}re\_\allowbreak{}of\_\allowbreak{}isPrimitive} (\texttt{Main}, lines 12--12)\\
\texttt{PrimitiveDirichletLZerosOnLine} (\texttt{Statements}, lines 14--16)\par}

{\footnotesize\raggedright\noindent \ref{B1.06}.\quad \texttt{Formalism.\allowbreak{}VW.\allowbreak{}rh\_\allowbreak{}of\_\allowbreak{}every} (\texttt{Formalism.\allowbreak{}GammaFactorNonvanishing.\allowbreak{}CritLine}, lines 29--32)\par}

{\footnotesize\raggedright\noindent \ref{B1.07}.\quad \texttt{Claim.\allowbreak{}strip} (\texttt{Development.\allowbreak{}AutomorphicArithmetic.\allowbreak{}CharacterBoxes.\allowbreak{}ClaimStrip}, lines 59--73)\par}

{\footnotesize\raggedright\noindent \ref{B1.08}.\quad \texttt{Formalism.\allowbreak{}VW.\allowbreak{}critLine} (\texttt{Formalism.\allowbreak{}GammaFactorNonvanishing.\allowbreak{}CritLine}, lines 11--27)\par}

{\footnotesize\raggedright\noindent \ref{B1.09}.\quad \texttt{allChi} (\texttt{Supplements.\allowbreak{}Kernels.\allowbreak{}MainTheorems}, lines 7--7)\par}

{\footnotesize\raggedright\noindent \ref{B1.10}.\quad \texttt{UpperTheory.\allowbreak{}EveryCharacter} (\texttt{Development.\allowbreak{}AutomorphicArithmetic.\allowbreak{}CharacterBoxes.\allowbreak{}EveryCharacter}, lines 13--14)\par}

{\footnotesize\raggedright\noindent \ref{B1.11}.\quad \texttt{Formalism.\allowbreak{}VW.\allowbreak{}rh\_\allowbreak{}of\_\allowbreak{}right} (\texttt{Formalism.\allowbreak{}GammaFactorNonvanishing.\allowbreak{}RhOfRight}, lines 5--62)\par}

{\footnotesize\raggedright\noindent \ref{B1.12}.\quad \texttt{DirichletCharacter.\allowbreak{}LFunction\_\allowbreak{}changeLevel} (\texttt{Mathlib.\allowbreak{}NumberTheory.\allowbreak{}LSeries.\allowbreak{}DirichletContinuation}, lines 150--162)\\
\texttt{Claim.\allowbreak{}factor\_\allowbreak{}ne\_\allowbreak{}zero} (\texttt{Development.\allowbreak{}AutomorphicArithmetic.\allowbreak{}CharacterBoxes.\allowbreak{}ClaimStrip}, lines 7--21)\\
\texttt{Claim.\allowbreak{}prim\_\allowbreak{}zero} (\texttt{Development.\allowbreak{}AutomorphicArithmetic.\allowbreak{}CharacterBoxes.\allowbreak{}ClaimStrip}, lines 23--33)\par}

{\footnotesize\raggedright\noindent \ref{B1.13}.\quad \texttt{Claim.\allowbreak{}prim\_\allowbreak{}reflect} (\texttt{Development.\allowbreak{}AutomorphicArithmetic.\allowbreak{}CharacterBoxes.\allowbreak{}ClaimStrip}, lines 42--57)\\
\texttt{Claim.\allowbreak{}gamma\_\allowbreak{}ne\_\allowbreak{}zero} (\texttt{Development.\allowbreak{}AutomorphicArithmetic.\allowbreak{}CharacterBoxes.\allowbreak{}ClaimStrip}, lines 35--40)\\
\texttt{Formalism.\allowbreak{}VW.\allowbreak{}gammaFactor\_\allowbreak{}ne\_\allowbreak{}zero} (\texttt{Formalism.\allowbreak{}GammaFactorNonvanishing.\allowbreak{}CritLine}, lines 6--9)\\
\texttt{DirichletCharacter.\allowbreak{}LFunction\_\allowbreak{}eq\_\allowbreak{}completed\_\allowbreak{}div\_\allowbreak{}gammaFactor} (\texttt{Mathlib.\allowbreak{}NumberTheory.\allowbreak{}LSeries.\allowbreak{}DirichletContinuation}, lines 259--264)\par}

{\footnotesize\raggedright\noindent \ref{B1.13}.\quad \texttt{DirichletCharacter.\allowbreak{}IsPrimitive.\allowbreak{}completedLFunction\_\allowbreak{}one\_\allowbreak{}sub} (\texttt{Mathlib.\allowbreak{}NumberTheory.\allowbreak{}LSeries.\allowbreak{}DirichletContinuation}, lines 284--313)\par}

{\footnotesize\raggedright\noindent \ref{B1.14}.\quad \texttt{UpperTheory.\allowbreak{}Top.\allowbreak{}imprimitive} (\texttt{Development.\allowbreak{}AutomorphicArithmetic.\allowbreak{}CharacterBoxes.\allowbreak{}Imprimitive}, lines 5--35)\par}

{\footnotesize\raggedright\noindent \ref{B1.14k}.\quad \texttt{GlobalTheory.\allowbreak{}RowTest.\allowbreak{}NumbersIn} (\texttt{Development.\allowbreak{}AbelianVarietyTheory.\allowbreak{}Definitions.\allowbreak{}TwistedInnerProducts.\allowbreak{}MemHalfCosetSelf}, lines 245--261)\par}

{\footnotesize\raggedright\noindent \ref{B1.14k}.\quad \texttt{GlobalTheory.\allowbreak{}Mem.\allowbreak{}LiteralAutomorphy.\allowbreak{}JoinOut} (\texttt{Development.\allowbreak{}AbelianVarietyTheory.\allowbreak{}Definitions.\allowbreak{}Divisors.\allowbreak{}HypothesisBundleConversions}, lines 7--9)\\
\texttt{GlobalTheory.\allowbreak{}IndefiniteRowSideFullHecke.\allowbreak{}ClassZetaValuationBounds.\allowbreak{}ToricDenominator} (\texttt{Development.\allowbreak{}AbelianVarietyTheory.\allowbreak{}BlockTheory.\allowbreak{}DivisorBase.\allowbreak{}VpBounds}, lines 77--82)\\
\texttt{GlobalTheory.\allowbreak{}IndefiniteRowSideFullHecke.\allowbreak{}LatticeNormTransport.\allowbreak{}Transport} (\texttt{Development.\allowbreak{}AbelianVarietyTheory.\allowbreak{}BlockTheory.\allowbreak{}DivisorBase.\allowbreak{}LatticeNormTransport}, lines 86--98)\\
\texttt{GlobalTheory.\allowbreak{}ToricTransport.\allowbreak{}TwistedTransport} (\texttt{Development.\allowbreak{}AbelianVarietyTheory.\allowbreak{}Proofs.\allowbreak{}DivisorProofs.\allowbreak{}TwistedTransport}, lines 21--40)\\
\texttt{GlobalTheory.\allowbreak{}IndefiniteRowSideFullHecke.\allowbreak{}ZetaValueIdeals.\allowbreak{}ZetaDenominators} (\texttt{Development.\allowbreak{}AbelianVarietyTheory.\allowbreak{}BlockTheory.\allowbreak{}DivisorBase.\allowbreak{}ZetaValueIdeals}, lines 45--46)\\
\texttt{GlobalTheory.\allowbreak{}SqTop.\allowbreak{}Refl} (\texttt{Development.\allowbreak{}AbelianVarietyTheory.\allowbreak{}Statements.\allowbreak{}PlaceShiftStatements.\allowbreak{}NormSmall}, lines 148--151)\par}

{\footnotesize\raggedright\noindent \ref{B1.14k}.\quad \texttt{Formalism.\allowbreak{}MainNotions.\allowbreak{}LinkN} (\texttt{Formalism.\allowbreak{}SharpRestriction.\allowbreak{}GlobalTheoremInputs}, lines 45--47)\\
\texttt{GlobalTheory.\allowbreak{}Fc.\allowbreak{}HypsAlt} (\texttt{Development.\allowbreak{}AutomorphicArithmetic.\allowbreak{}ThetaCorrespondence.\allowbreak{}ThetaPairLifts.\allowbreak{}ThirdPlaceSeries.\allowbreak{}Hyps}, lines 22--90)\par}

{\footnotesize\raggedright\noindent \ref{B1.14a}.\quad \texttt{GlobalTheory.\allowbreak{}QFam.\allowbreak{}sz\_\allowbreak{}value\_\allowbreak{}of\_\allowbreak{}family} (\texttt{Development.\allowbreak{}AbelianVarietyTheory.\allowbreak{}Proofs.\allowbreak{}DivisorProofs.\allowbreak{}SzValueOfFamily}, lines 23--43)\par}

{\footnotesize\raggedright\noindent \ref{B1.14b}.\quad \texttt{GlobalTheory.\allowbreak{}CMSize.\allowbreak{}main} (\texttt{Development.\allowbreak{}AbelianVarietyTheory.\allowbreak{}Proofs.\allowbreak{}DivisorProofs.\allowbreak{}DenominatorBoundWeierstrassCmTorsion}, lines 30--96)\par}

{\footnotesize\raggedright\noindent \ref{B1.14b}.\quad \texttt{GlobalTheory.\allowbreak{}IsogFormula.\allowbreak{}main} (\texttt{Development.\allowbreak{}AbelianVarietyTheory.\allowbreak{}Proofs.\allowbreak{}DivisorProofs.\allowbreak{}IsogenyFormulaWeierstrassPIntegral}, lines 16--73)\par}

{\footnotesize\raggedright\noindent \ref{B1.14c}.\quad \texttt{GlobalTheory.\allowbreak{}LatJet.\allowbreak{}JetInt} (\texttt{Development.\allowbreak{}AbelianVarietyTheory.\allowbreak{}Definitions.\allowbreak{}Divisors.\allowbreak{}JetAnalytic}, lines 239--246)\\
\texttt{GlobalTheory.\allowbreak{}JetIntElliptic.\allowbreak{}jetInt} (\texttt{Development.\allowbreak{}AbelianVarietyTheory.\allowbreak{}Proofs.\allowbreak{}DivisorProofs.\allowbreak{}CMLatticeDividedPowerJetIntegralityOfHyp}, lines 251--252)\par}

{\footnotesize\raggedright\noindent \ref{B1.14c}.\quad \texttt{GlobalTheory.\allowbreak{}JetIntElliptic.\allowbreak{}core} (\texttt{Development.\allowbreak{}AbelianVarietyTheory.\allowbreak{}Proofs.\allowbreak{}DivisorProofs.\allowbreak{}CMLatticeDividedPowerJetIntegralityOfHyp}, lines 142--249)\par}

{\footnotesize\raggedright\noindent \ref{B1.15a}.\quad \texttt{GlobalTheory.\allowbreak{}JoinFun.\allowbreak{}joinOut\_\allowbreak{}of\_\allowbreak{}five} (\texttt{Development.\allowbreak{}AbelianVarietyTheory.\allowbreak{}Proofs.\allowbreak{}DivisorProofs.\allowbreak{}JoinOutOfFive}, lines 15--17)\\
\texttt{UpperTheory.\allowbreak{}LiteratureTheorems.\allowbreak{}join\_\allowbreak{}out\_\allowbreak{}of\_\allowbreak{}literature} (\texttt{Development.\allowbreak{}AutomorphicArithmetic.\allowbreak{}ThetaCorrespondence.\allowbreak{}ThetaPairLifts.\allowbreak{}ThirdPlaceSeries.\allowbreak{}EveryCharacterOfLiteratureOfVWin}, lines 233--234)\par}

{\footnotesize\raggedright\noindent \ref{B1.15ac}.\quad \texttt{UpperTheory.\allowbreak{}GalS.\allowbreak{}thr} (\texttt{Development.\allowbreak{}AutomorphicArithmetic.\allowbreak{}Foundations.\allowbreak{}GaloisRepresentations.\allowbreak{}TypeConditions.\allowbreak{}Transported}, lines 11--13)\\
\texttt{UpperTheory.\allowbreak{}GalS.\allowbreak{}NcW} (\texttt{Development.\allowbreak{}AutomorphicArithmetic.\allowbreak{}Foundations.\allowbreak{}GaloisRepresentations.\allowbreak{}TypeConditions.\allowbreak{}WeightsConj}, lines 19--28)\par}

{\footnotesize\raggedright\noindent \ref{B1.15ac}.\quad \texttt{GlobalTheory.\allowbreak{}Fc.\allowbreak{}S0.\allowbreak{}UnitaryGroupThirdPlaceHypotheses} (\texttt{Development.\allowbreak{}AutomorphicArithmetic.\allowbreak{}ThetaCorrespondence.\allowbreak{}ThetaPairLifts.\allowbreak{}ThirdPlaceSeries.\allowbreak{}BernsteinUnitaryProductIrreducible}, lines 22--82)\\
\texttt{GlobalTheory.\allowbreak{}FunPr.\allowbreak{}ThetaLocalInputs} (\texttt{Development.\allowbreak{}AutomorphicArithmetic.\allowbreak{}ThetaCorrespondence.\allowbreak{}ThetaPairLifts.\allowbreak{}LocalTypeDictionary.\allowbreak{}BernsteinUnitaryProductIrreducible}, lines 6--59)\\
\texttt{GlobalTheory.\allowbreak{}Mem.\allowbreak{}LiteralAutomorphy.\allowbreak{}NonEndoscopicHypothesesAlt} (\texttt{Development.\allowbreak{}AutomorphicArithmetic.\allowbreak{}ThetaCorrespondence.\allowbreak{}ThetaPairLifts.\allowbreak{}ThirdPlaceSeries.\allowbreak{}BernsteinUnitaryProductIrreducible}, lines 12--17)\\
\texttt{GlobalTheory.\allowbreak{}NRows.\allowbreak{}GeometricGaloisRepCompatibility} (\texttt{Development.\allowbreak{}AutomorphicArithmetic.\allowbreak{}ThetaCorrespondence.\allowbreak{}ThetaPairLifts.\allowbreak{}LocalTypeDictionary.\allowbreak{}GeometricGaloisRepCompatibility}, lines 6--18)\\
\texttt{UpperTheory.\allowbreak{}LiteratureTheorems.\allowbreak{}published\_\allowbreak{}results\_\allowbreak{}local} (\texttt{Development.\allowbreak{}AutomorphicArithmetic.\allowbreak{}ThetaCorrespondence.\allowbreak{}ThetaPairLifts.\allowbreak{}ThirdPlaceSeries.\allowbreak{}EveryCharacterOfLiteratureOfVWin}, lines 658--665)\par}

{\footnotesize\raggedright\noindent \ref{B1.15ba}.\quad \texttt{PositivityTheory.\allowbreak{}KReal.\allowbreak{}tensorIrr} (\texttt{Development.\allowbreak{}WeilBound.\allowbreak{}CoreStatements.\allowbreak{}GaussianSeminormBounds.\allowbreak{}RealSeminormFiltration.\allowbreak{}Tensor}, lines 131--166)\\
\texttt{GlobalTheory.\allowbreak{}ThK.\allowbreak{}Dbl.\allowbreak{}pairsIrr\_\allowbreak{}of\_\allowbreak{}arch} (\texttt{Development.\allowbreak{}WeilBound.\allowbreak{}ThetaLifts.\allowbreak{}LatticeBasics.\allowbreak{}PairsIrrOfArch}, lines 10--12)\par}

{\footnotesize\raggedright\noindent \ref{B1.15}.\quad \texttt{UpperTheory.\allowbreak{}FinalTheorems.\allowbreak{}everyCharacter\_\allowbreak{}of\_\allowbreak{}row} (\texttt{Development.\allowbreak{}AutomorphicArithmetic.\allowbreak{}CharacterBoxes.\allowbreak{}MainConclusion}, lines 7--8)\\
\texttt{UpperTheory.\allowbreak{}FinalTheorems.\allowbreak{}everyCharacter} (\texttt{Development.\allowbreak{}AutomorphicArithmetic.\allowbreak{}CharacterBoxes.\allowbreak{}MainConclusion}, lines 10--11)\par}

\section*{Chapter~\ref{ch:arithmetic-windows}: How arithmetic windows enter the zero-free criterion}

{\footnotesize\raggedright\noindent \ref{def:p1c2-host}.\quad \texttt{GlobalTheory.\allowbreak{}EX.\allowbreak{}HostVal} (\texttt{Development.\allowbreak{}AbelianVarietyTheory.\allowbreak{}Definitions.\allowbreak{}GlobalTwisting.\allowbreak{}PairRowIffQdRows}, lines 86--92)\\
\texttt{GlobalTheory.\allowbreak{}EX.\allowbreak{}HostGood} (\texttt{Development.\allowbreak{}AbelianVarietyTheory.\allowbreak{}Definitions.\allowbreak{}GlobalTwisting.\allowbreak{}ThinScalesIffHostBound}, lines 101--110)\\
\texttt{GlobalTheory.\allowbreak{}EX.\allowbreak{}HostExW} (\texttt{Development.\allowbreak{}AbelianVarietyTheory.\allowbreak{}Definitions.\allowbreak{}GlobalTwisting.\allowbreak{}ThinScalesIffHostBound}, lines 120--124)\\
\texttt{UpperTheory.\allowbreak{}Top.\allowbreak{}FReq.\allowbreak{}req} (\texttt{Development.\allowbreak{}AutomorphicArithmetic.\allowbreak{}CharacterBoxes.\allowbreak{}Supply}, lines 16--17)\\
\texttt{UpperTheory.\allowbreak{}Top.\allowbreak{}GChi} (\texttt{Development.\allowbreak{}AutomorphicArithmetic.\allowbreak{}CharacterBoxes.\allowbreak{}Supply}, lines 42--44)\\
\texttt{UpperTheory.\allowbreak{}Top.\allowbreak{}VVle} (\texttt{Development.\allowbreak{}AutomorphicArithmetic.\allowbreak{}CharacterBoxes.\allowbreak{}ExistsPageN}, lines 12--18)\par}

{\footnotesize\raggedright\noindent \ref{B1.18}.\quad \texttt{UpperTheory.\allowbreak{}Top.\allowbreak{}zeroFreeAll\_\allowbreak{}of\_\allowbreak{}topG} (\texttt{Development.\allowbreak{}AutomorphicArithmetic.\allowbreak{}CharacterBoxes.\allowbreak{}ZeroFreeAllOfTopG}, lines 26--35)\\
\texttt{UpperTheory.\allowbreak{}Top.\allowbreak{}TopG} (\texttt{Development.\allowbreak{}AutomorphicArithmetic.\allowbreak{}CharacterBoxes.\allowbreak{}ZeroFreeAllOfTopG}, lines 12--15)\par}

{\footnotesize\raggedright\noindent \ref{B1.18}.\quad \texttt{UpperTheory.\allowbreak{}Comparison₂.\allowbreak{}level\_\allowbreak{}gm\_\allowbreak{}le\_\allowbreak{}on\_\allowbreak{}family\_\allowbreak{}scales} (\texttt{Development.\allowbreak{}AutomorphicArithmetic.\allowbreak{}PrimitiveCosets.\allowbreak{}LevelGmLeOnFamilyScales}, lines 28--31)\par}

{\footnotesize\raggedright\noindent \ref{B1.19}.\quad \texttt{Formalism.\allowbreak{}HC.\allowbreak{}zeroFreeOn\_\allowbreak{}of\_\allowbreak{}tiedWin\_\allowbreak{}cyc} (\texttt{Formalism.\allowbreak{}SharpRestriction.\allowbreak{}RamificationFrobenius.\allowbreak{}TotientGaloisClasses.\allowbreak{}TopWin}, lines 90--94)\\
\texttt{Formalism.\allowbreak{}HC.\allowbreak{}zeroFreeOn\_\allowbreak{}of\_\allowbreak{}tiedWin} (\texttt{Formalism.\allowbreak{}SharpRestriction.\allowbreak{}RamificationFrobenius.\allowbreak{}TotientGaloisClasses.\allowbreak{}TopWin}, lines 83--89)\par}

{\footnotesize\raggedright\noindent \ref{B1.19}.\quad \texttt{Formalism.\allowbreak{}HC.\allowbreak{}Level.\allowbreak{}mem\_\allowbreak{}cyc} (\texttt{Formalism.\allowbreak{}SharpRestriction.\allowbreak{}RamificationFrobenius.\allowbreak{}TotientGaloisClasses.\allowbreak{}\allowbreak ExplicitFormulaErrorOfHostConditions}, lines 123--125)\par}

{\footnotesize\raggedright\noindent \ref{B1.20}.\quad \texttt{Formalism.\allowbreak{}HC.\allowbreak{}level\_\allowbreak{}of\_\allowbreak{}tiedWin} (\texttt{Formalism.\allowbreak{}SharpRestriction.\allowbreak{}RamificationFrobenius.\allowbreak{}TotientGaloisClasses.\allowbreak{}TopWin}, lines 68--82)\par}

{\footnotesize\raggedright\noindent \ref{B1.21}.\quad \texttt{Formalism.\allowbreak{}HC.\allowbreak{}KSet} (\texttt{Formalism.\allowbreak{}SharpRestriction.\allowbreak{}RamificationFrobenius.\allowbreak{}TotientGaloisClasses.\allowbreak{}PrimePairSizeCondition}, lines 39--46)\\
\texttt{Formalism.\allowbreak{}HC.\allowbreak{}Level} (\texttt{Formalism.\allowbreak{}SharpRestriction.\allowbreak{}RamificationFrobenius.\allowbreak{}TotientGaloisClasses.\allowbreak{}\allowbreak ExplicitFormulaErrorOfHostConditions}, lines 8--10)\\
\texttt{Formalism.\allowbreak{}HC.\allowbreak{}Level.\allowbreak{}wt} (\texttt{Formalism.\allowbreak{}SharpRestriction.\allowbreak{}RamificationFrobenius.\allowbreak{}TotientGaloisClasses.\allowbreak{}\allowbreak ExplicitFormulaErrorOfHostConditions}, lines 36--37)\\
\texttt{Formalism.\allowbreak{}HC.\allowbreak{}Level.\allowbreak{}ofSub} (\texttt{Formalism.\allowbreak{}SharpRestriction.\allowbreak{}RamificationFrobenius.\allowbreak{}TotientGaloisClasses.\allowbreak{}\allowbreak ExplicitFormulaErrorOfHostConditions}, lines 105--115)\\
\texttt{Formalism.\allowbreak{}HC.\allowbreak{}Level.\allowbreak{}Gm} (\texttt{Formalism.\allowbreak{}SharpRestriction.\allowbreak{}RamificationFrobenius.\allowbreak{}TotientGaloisClasses.\allowbreak{}\allowbreak ExplicitFormulaErrorOfHostConditions}, lines 38--39)\\
\texttt{Formalism.\allowbreak{}HC.\allowbreak{}Level.\allowbreak{}cyc} (\texttt{Formalism.\allowbreak{}SharpRestriction.\allowbreak{}RamificationFrobenius.\allowbreak{}TotientGaloisClasses.\allowbreak{}\allowbreak ExplicitFormulaErrorOfHostConditions}, lines 121--122)\par}

{\footnotesize\raggedright\noindent \ref{B1.22}.\quad \texttt{UpperTheory.\allowbreak{}Comparison₂.\allowbreak{}formB} (\texttt{Development.\allowbreak{}AutomorphicArithmetic.\allowbreak{}PrimitiveCosets.\allowbreak{}Imports}, lines 8--20)\\
\texttt{Formalism.\allowbreak{}HC.\allowbreak{}Level.\allowbreak{}FormBAt} (\texttt{Formalism.\allowbreak{}SharpRestriction.\allowbreak{}RamificationFrobenius.\allowbreak{}TotientGaloisClasses.\allowbreak{}Window}, lines 30--32)\\
\texttt{Formalism.\allowbreak{}HC.\allowbreak{}Level.\allowbreak{}FormB} (\texttt{Formalism.\allowbreak{}SharpRestriction.\allowbreak{}RamificationFrobenius.\allowbreak{}TotientGaloisClasses.\allowbreak{}Window}, lines 33--34)\par}

{\footnotesize\raggedright\noindent \ref{B1.23}.\quad \texttt{UpperTheory.\allowbreak{}Comparison₂.\allowbreak{}level\_\allowbreak{}family} (\texttt{Development.\allowbreak{}AutomorphicArithmetic.\allowbreak{}PrimitiveCosets.\allowbreak{}LevelFamily}, lines 41--53)\par}

{\footnotesize\raggedright\noindent \ref{B1.24}.\quad \texttt{Deduction.\allowbreak{}zeroFree\_\allowbreak{}of\_\allowbreak{}rpow\_\allowbreak{}window} (\texttt{Development.\allowbreak{}AutomorphicArithmetic.\allowbreak{}DedekindLFunctions.\allowbreak{}LandauTheorems.\allowbreak{}ZeroFreeOfRpowWindow}, lines 11--15)\par}

{\footnotesize\raggedright\noindent \ref{B1.25}.\quad \texttt{Deduction.\allowbreak{}CharFamily} (\texttt{Development.\allowbreak{}AutomorphicArithmetic.\allowbreak{}DedekindLFunctions.\allowbreak{}CharFamily}, lines 24--33)\\
\texttt{Deduction.\allowbreak{}ZeroFree} (\texttt{Development.\allowbreak{}AutomorphicArithmetic.\allowbreak{}DedekindLFunctions.\allowbreak{}CharFamily}, lines 38--40)\par}

{\footnotesize\raggedright\noindent \ref{B1.26}.\quad \texttt{Deduction.\allowbreak{}ramp} (\texttt{Development.\allowbreak{}AutomorphicArithmetic.\allowbreak{}DedekindLFunctions.\allowbreak{}CharFamily}, lines 10--12)\\
\texttt{Deduction.\allowbreak{}Gchi} (\texttt{Development.\allowbreak{}AutomorphicArithmetic.\allowbreak{}DedekindLFunctions.\allowbreak{}CharFamily}, lines 15--20)\\
\texttt{Deduction.\allowbreak{}Gm} (\texttt{Development.\allowbreak{}AutomorphicArithmetic.\allowbreak{}DedekindLFunctions.\allowbreak{}CharFamily}, lines 36--37)\\
\texttt{Deduction.\allowbreak{}gm} (\texttt{Development.\allowbreak{}AutomorphicArithmetic.\allowbreak{}DedekindLFunctions.\allowbreak{}LandauZeroFreeCriterion}, lines 9--10)\par}

{\footnotesize\raggedright\noindent \ref{B1.27}.\quad \texttt{Deduction.\allowbreak{}WindowHyp} (\texttt{Development.\allowbreak{}AutomorphicArithmetic.\allowbreak{}DedekindLFunctions.\allowbreak{}CharFamily}, lines 41--46)\\
\texttt{Deduction.\allowbreak{}windowHyp\_\allowbreak{}of\_\allowbreak{}rpow} (\texttt{Development.\allowbreak{}AutomorphicArithmetic.\allowbreak{}DedekindLFunctions.\allowbreak{}LandauTheorems.\allowbreak{}WindowHypOfRpow}, lines 27--58)\par}

{\footnotesize\raggedright\noindent \ref{B1.28}.\quad \texttt{Deduction.\allowbreak{}lemmaL} (\texttt{Development.\allowbreak{}AutomorphicArithmetic.\allowbreak{}DedekindLFunctions.\allowbreak{}LandauTheorems.\allowbreak{}lemmaL}, lines 13--23)\par}

{\footnotesize\raggedright\noindent \ref{B1.29}.\quad \texttt{Deduction.\allowbreak{}WindowHypU} (\texttt{Development.\allowbreak{}AutomorphicArithmetic.\allowbreak{}DedekindLFunctions.\allowbreak{}LandauZeroFreeCriterion}, lines 25--30)\\
\texttt{Deduction.\allowbreak{}windowHypU\_\allowbreak{}of\_\allowbreak{}windowHyp} (\texttt{Development.\allowbreak{}AutomorphicArithmetic.\allowbreak{}DedekindLFunctions.\allowbreak{}LandauTheorems.\allowbreak{}WindowHypUOfWindowHyp}, lines 29--57)\par}

\section*{Chapter~\ref{ch:zeros-become-poles}: Zeros become poles}

{\footnotesize\raggedright\noindent \ref{B1.30}.\quad \texttt{Deduction.\allowbreak{}poleData} (\texttt{Development.\allowbreak{}AutomorphicArithmetic.\allowbreak{}DedekindLFunctions.\allowbreak{}LandauTheorems.\allowbreak{}poleData}, lines 248--401)\par}

{\footnotesize\raggedright\noindent \ref{B1.31}.\quad \texttt{Deduction.\allowbreak{}landau\_\allowbreak{}window} (\texttt{Development.\allowbreak{}AutomorphicArithmetic.\allowbreak{}DedekindLFunctions.\allowbreak{}LandauTheorems.\allowbreak{}landauWindow}, lines 83--167)\par}

{\footnotesize\raggedright\noindent \ref{B1.32}.\quad \texttt{Deduction.\allowbreak{}charFamily\_\allowbreak{}LFunction\_\allowbreak{}ne\_\allowbreak{}zero\_\allowbreak{}of\_\allowbreak{}one\_\allowbreak{}le\_\allowbreak{}re} (\texttt{Development.\allowbreak{}AutomorphicArithmetic.\allowbreak{}DedekindLFunctions.\allowbreak{}LandauTheorems.\allowbreak{}\allowbreak CharFamilyLFunctionNeZeroOfOneLeRe}, lines 9--21)\par}

{\footnotesize\raggedright\noindent \ref{B1.33}.\quad \texttt{Deduction.\allowbreak{}zeroSet} (\texttt{Development.\allowbreak{}AutomorphicArithmetic.\allowbreak{}DedekindLFunctions.\allowbreak{}LandauZeroFreeCriterion}, lines 11--14)\\
\texttt{Deduction.\allowbreak{}PoleData} (\texttt{Development.\allowbreak{}AutomorphicArithmetic.\allowbreak{}DedekindLFunctions.\allowbreak{}LandauZeroFreeCriterion}, lines 15--24)\par}

{\footnotesize\raggedright\noindent \ref{B1.34}.\quad \texttt{Deduction.\allowbreak{}poleData\_\allowbreak{}char} (\texttt{Development.\allowbreak{}AutomorphicArithmetic.\allowbreak{}DedekindLFunctions.\allowbreak{}LandauTheorems.\allowbreak{}poleData}, lines 137--240)\par}

{\footnotesize\raggedright\noindent \ref{B1.35}.\quad \texttt{Deduction.\allowbreak{}logDeriv\_\allowbreak{}local} (\texttt{Development.\allowbreak{}AutomorphicArithmetic.\allowbreak{}DedekindLFunctions.\allowbreak{}LandauTheorems.\allowbreak{}LogDerivLocal}, lines 9--40)\par}

{\footnotesize\raggedright\noindent \ref{B1.36}.\quad \texttt{Deduction.\allowbreak{}poleData\_\allowbreak{}abstract} (\texttt{Development.\allowbreak{}AutomorphicArithmetic.\allowbreak{}DedekindLFunctions.\allowbreak{}LandauTheorems.\allowbreak{}poleData}, lines 22--90)\\
\texttt{Deduction.\allowbreak{}poleData\_\allowbreak{}dslope} (\texttt{Development.\allowbreak{}AutomorphicArithmetic.\allowbreak{}DedekindLFunctions.\allowbreak{}LandauTheorems.\allowbreak{}poleData}, lines 17--21)\par}

{\footnotesize\raggedright\noindent \ref{B1.37}.\quad \texttt{Deduction.\allowbreak{}poleData\_\allowbreak{}change} (\texttt{Development.\allowbreak{}AutomorphicArithmetic.\allowbreak{}DedekindLFunctions.\allowbreak{}LandauTheorems.\allowbreak{}poleData}, lines 114--136)\\
\texttt{Deduction.\allowbreak{}poleData\_\allowbreak{}integrable} (\texttt{Development.\allowbreak{}AutomorphicArithmetic.\allowbreak{}DedekindLFunctions.\allowbreak{}LandauTheorems.\allowbreak{}poleData}, lines 91--113)\par}

{\footnotesize\raggedright\noindent \ref{B1.38}.\quad \texttt{Deduction.\allowbreak{}measurable\_\allowbreak{}Gchi} (\texttt{Development.\allowbreak{}AutomorphicArithmetic.\allowbreak{}DedekindLFunctions.\allowbreak{}LandauTheorems.\allowbreak{}measurableGchi}, lines 24--29)\\
\texttt{Deduction.\allowbreak{}Gchi\_\allowbreak{}growth\_\allowbreak{}ramp} (\texttt{Development.\allowbreak{}AutomorphicArithmetic.\allowbreak{}DedekindLFunctions.\allowbreak{}LandauTheorems.\allowbreak{}GchiGrowth}, lines 9--36)\\
\texttt{Deduction.\allowbreak{}Gchi\_\allowbreak{}growth} (\texttt{Development.\allowbreak{}AutomorphicArithmetic.\allowbreak{}DedekindLFunctions.\allowbreak{}LandauTheorems.\allowbreak{}GchiGrowth}, lines 37--67)\par}

{\footnotesize\raggedright\noindent \ref{B1.39}.\quad \texttt{Deduction.\allowbreak{}mellin\_\allowbreak{}ramp} (\texttt{Development.\allowbreak{}AutomorphicArithmetic.\allowbreak{}DedekindLFunctions.\allowbreak{}LandauTheorems.\allowbreak{}MellinRamp}, lines 10--52)\\
\texttt{Deduction.\allowbreak{}mellin\_\allowbreak{}Gchi\_\allowbreak{}of\_\allowbreak{}ne\_\allowbreak{}one\_\allowbreak{}norm} (\texttt{Development.\allowbreak{}AutomorphicArithmetic.\allowbreak{}DedekindLFunctions.\allowbreak{}LandauTheorems.\allowbreak{}mellinGchiOfNeOne}, lines 27--42)\par}

{\footnotesize\raggedright\noindent \ref{B1.40}.\quad \texttt{Deduction.\allowbreak{}mellin\_\allowbreak{}Gchi\_\allowbreak{}of\_\allowbreak{}ne\_\allowbreak{}one\_\allowbreak{}ramp} (\texttt{Development.\allowbreak{}AutomorphicArithmetic.\allowbreak{}DedekindLFunctions.\allowbreak{}LandauTheorems.\allowbreak{}mellinGchiOfNeOne}, lines 65--115)\\
\texttt{Deduction.\allowbreak{}mellin\_\allowbreak{}Gchi\_\allowbreak{}modOne\_\allowbreak{}sum} (\texttt{Development.\allowbreak{}AutomorphicArithmetic.\allowbreak{}DedekindLFunctions.\allowbreak{}LandauTheorems.\allowbreak{}mellinGchiModOne}, lines 10--26)\par}

{\footnotesize\raggedright\noindent \ref{B1.41}.\quad \texttt{Deduction.\allowbreak{}mellin\_\allowbreak{}Gchi\_\allowbreak{}of\_\allowbreak{}ne\_\allowbreak{}one} (\texttt{Development.\allowbreak{}AutomorphicArithmetic.\allowbreak{}DedekindLFunctions.\allowbreak{}LandauTheorems.\allowbreak{}mellinGchiOfNeOne}, lines 116--143)\par}

{\footnotesize\raggedright\noindent \ref{B1.42}.\quad \texttt{Deduction.\allowbreak{}mellin\_\allowbreak{}Gchi\_\allowbreak{}modOne\_\allowbreak{}log\_\allowbreak{}int} (\texttt{Development.\allowbreak{}AutomorphicArithmetic.\allowbreak{}DedekindLFunctions.\allowbreak{}LandauTheorems.\allowbreak{}mellinGchiModOne}, lines 57--84)\\
\texttt{Deduction.\allowbreak{}mellin\_\allowbreak{}Gchi\_\allowbreak{}modOne\_\allowbreak{}log} (\texttt{Development.\allowbreak{}AutomorphicArithmetic.\allowbreak{}DedekindLFunctions.\allowbreak{}LandauTheorems.\allowbreak{}mellinGchiModOne}, lines 85--114)\par}

{\footnotesize\raggedright\noindent \ref{B1.43}.\quad \texttt{Deduction.\allowbreak{}mellin\_\allowbreak{}Gchi\_\allowbreak{}modOne} (\texttt{Development.\allowbreak{}AutomorphicArithmetic.\allowbreak{}DedekindLFunctions.\allowbreak{}LandauTheorems.\allowbreak{}mellinGchiModOne}, lines 115--200)\par}

{\footnotesize\raggedright\noindent \ref{B1.44}.\quad \texttt{Deduction.\allowbreak{}LFunctionConjugation.\allowbreak{}hurwitz\_\allowbreak{}conj} (\texttt{Development.\allowbreak{}AutomorphicArithmetic.\allowbreak{}DedekindLFunctions.\allowbreak{}LandauTheorems.\allowbreak{}LFunctionConj}, lines 76--85)\\
\texttt{Deduction.\allowbreak{}LFunctionConjugation.\allowbreak{}Gammaℝ\_\allowbreak{}conj} (\texttt{Development.\allowbreak{}AutomorphicArithmetic.\allowbreak{}DedekindLFunctions.\allowbreak{}LandauTheorems.\allowbreak{}LFunctionConj}, lines 19--25)\\
\texttt{Deduction.\allowbreak{}LFunctionConjugation.\allowbreak{}differentiable\_\allowbreak{}reg} (\texttt{Development.\allowbreak{}AutomorphicArithmetic.\allowbreak{}DedekindLFunctions.\allowbreak{}LandauTheorems.\allowbreak{}LFunctionConj}, lines 38--49)\\
\texttt{HurwitzZeta.\allowbreak{}tendsto\_\allowbreak{}hurwitzZeta\_\allowbreak{}sub\_\allowbreak{}one\_\allowbreak{}div\_\allowbreak{}nhds\_\allowbreak{}one} (\texttt{Mathlib.\allowbreak{}NumberTheory.\allowbreak{}LSeries.\allowbreak{}HurwitzZeta}, lines 96--100)\par}

{\footnotesize\raggedright\noindent \ref{B1.45}.\quad \texttt{Deduction.\allowbreak{}LFunction\_\allowbreak{}conj} (\texttt{Development.\allowbreak{}AutomorphicArithmetic.\allowbreak{}DedekindLFunctions.\allowbreak{}LandauTheorems.\allowbreak{}LFunctionConj}, lines 88--97)\\
\texttt{Deduction.\allowbreak{}Gchi\_\allowbreak{}conj} (\texttt{Development.\allowbreak{}AutomorphicArithmetic.\allowbreak{}DedekindLFunctions.\allowbreak{}LandauTheorems.\allowbreak{}GchiConj}, lines 38--52)\\
\texttt{Deduction.\allowbreak{}poleData\_\allowbreak{}inv\_\allowbreak{}zero} (\texttt{Development.\allowbreak{}AutomorphicArithmetic.\allowbreak{}DedekindLFunctions.\allowbreak{}LandauTheorems.\allowbreak{}poleData}, lines 241--247)\par}

{\footnotesize\raggedright\noindent \ref{B1.46}.\quad \texttt{Deduction.\allowbreak{}poleData} (\texttt{Development.\allowbreak{}AutomorphicArithmetic.\allowbreak{}DedekindLFunctions.\allowbreak{}LandauTheorems.\allowbreak{}poleData}, lines 336--400)\par}

\section*{Chapter~\ref{ch:p1c4}: A Gaussian window excludes poles}

{\footnotesize\raggedright\noindent \ref{B1.54}.\quad \texttt{Deduction.\allowbreak{}landau\_\allowbreak{}pole\_\allowbreak{}re\_\allowbreak{}le\_\allowbreak{}bound} (\texttt{Development.\allowbreak{}AutomorphicArithmetic.\allowbreak{}DedekindLFunctions.\allowbreak{}LandauTheorems.\allowbreak{}LandauPoleReLe}, lines 9--28)\\
\texttt{Deduction.\allowbreak{}landau\_\allowbreak{}pole\_\allowbreak{}re\_\allowbreak{}le} (\texttt{Development.\allowbreak{}AutomorphicArithmetic.\allowbreak{}DedekindLFunctions.\allowbreak{}LandauTheorems.\allowbreak{}LandauPoleReLe}, lines 29--54)\par}

{\footnotesize\raggedright\noindent \ref{B1.55}.\quad \texttt{Deduction.\allowbreak{}landau\_\allowbreak{}finite} (\texttt{Development.\allowbreak{}AutomorphicArithmetic.\allowbreak{}DedekindLFunctions.\allowbreak{}LandauTheorems.\allowbreak{}LandauFinite}, lines 9--27)\par}

\section*{Chapter~\ref{ch:p2c1}: Counting primes class by class}

{\footnotesize\raggedright\noindent \ref{B2.01}.\quad \texttt{Formalism.\allowbreak{}HC.\allowbreak{}Frame} (\texttt{Formalism.\allowbreak{}SharpRestriction.\allowbreak{}RamificationFrobenius.\allowbreak{}TotientGaloisClasses.\allowbreak{}ZsmulGeneratorEqZeroOfProjection}, lines 15--24)\\
\texttt{Formalism.\allowbreak{}HC.\allowbreak{}HostRaw} (\texttt{Formalism.\allowbreak{}SharpRestriction.\allowbreak{}RamificationFrobenius.\allowbreak{}TotientGaloisClasses.\allowbreak{}ConjEMemHostFIff}, lines 12--47)\\
\texttt{Formalism.\allowbreak{}HC.\allowbreak{}HostDer} (\texttt{Formalism.\allowbreak{}SharpRestriction.\allowbreak{}RamificationFrobenius.\allowbreak{}TotientGaloisClasses.\allowbreak{}ConjEMemHostFIff}, lines 87--171)\\
\texttt{Formalism.\allowbreak{}HC.\allowbreak{}HostAx} (\texttt{Formalism.\allowbreak{}SharpRestriction.\allowbreak{}RamificationFrobenius.\allowbreak{}TotientGaloisClasses.\allowbreak{}ConjEMemHostFIff}, lines 174--176)\par}

{\footnotesize\raggedright\noindent \ref{B2.02}.\quad \texttt{UpperTheory.\allowbreak{}Comparison₂.\allowbreak{}toE\_\allowbreak{}bij} (\texttt{Development.\allowbreak{}AutomorphicArithmetic.\allowbreak{}PrimitiveCosets.\allowbreak{}ToEBij}, lines 37--85)\\
\texttt{UpperTheory.\allowbreak{}Comparison₂.\allowbreak{}chr\_\allowbreak{}bij} (\texttt{Development.\allowbreak{}AutomorphicArithmetic.\allowbreak{}PrimitiveCosets.\allowbreak{}ChrBij}, lines 21--91)\par}

{\footnotesize\raggedright\noindent \ref{B2.03}.\quad \texttt{UpperTheory.\allowbreak{}Comparison₂.\allowbreak{}hostFactsS} (\texttt{Development.\allowbreak{}AutomorphicArithmetic.\allowbreak{}PrimitiveCosets.\allowbreak{}hostFactsS}, lines 127--137)\\
\texttt{Deduction.\allowbreak{}CH.\allowbreak{}HostFactsS} (\texttt{Development.\allowbreak{}AutomorphicArithmetic.\allowbreak{}DedekindLFunctions.\allowbreak{}SplitInertia}, lines 14--23)\par}

{\footnotesize\raggedright\noindent \ref{B2.04}.\quad \texttt{Formalism.\allowbreak{}LFunctionChainParameters.\allowbreak{}characterSumData.\allowbreak{}primVal} (\texttt{Formalism.\allowbreak{}SizeWindowBrackets.\allowbreak{}FirstPrimeNeSecondPrime}, lines 65--67)\\
\texttt{Formalism.\allowbreak{}LFunctionChainParameters.\allowbreak{}characterSumData.\allowbreak{}GPsi} (\texttt{Formalism.\allowbreak{}SizeWindowBrackets.\allowbreak{}FirstPrimeNeSecondPrime}, lines 68--71)\\
\texttt{Formalism.\allowbreak{}LFunctionChainParameters.\allowbreak{}characterSumData.\allowbreak{}ASK} (\texttt{Formalism.\allowbreak{}SizeWindowBrackets.\allowbreak{}FirstPrimeNeSecondPrime}, lines 84--88)\\
\texttt{Formalism.\allowbreak{}LFunctionChainParameters.\allowbreak{}characterSumData.\allowbreak{}MAIN} (\texttt{Formalism.\allowbreak{}SizeWindowBrackets.\allowbreak{}FirstPrimeNeSecondPrime}, lines 89--92)\\
\texttt{Formalism.\allowbreak{}LFunctionChainParameters.\allowbreak{}characterSumData.\allowbreak{}De} (\texttt{Formalism.\allowbreak{}SizeWindowBrackets.\allowbreak{}FirstPrimeNeSecondPrime}, lines 93--97)\par}

{\footnotesize\raggedright\noindent \ref{B2.05}.\quad \texttt{Formalism.\allowbreak{}LFunctionChainParameters.\allowbreak{}characterSumData.\allowbreak{}orthogonality\_\allowbreak{}ramified} (\texttt{Formalism.\allowbreak{}SizeWindowBrackets.\allowbreak{}ChainSlopeWindows.\allowbreak{}StronglyIrreducible}, lines 311--317)\\
\texttt{Formalism.\allowbreak{}LFunctionChainParameters.\allowbreak{}characterSumData.\allowbreak{}OrthoRam} (\texttt{Formalism.\allowbreak{}SizeWindowBrackets.\allowbreak{}FirstPrimeNeSecondPrime}, lines 116--122)\par}

{\footnotesize\raggedright\noindent \ref{B2.06}.\quad \texttt{Deduction.\allowbreak{}CharacterBalance.\allowbreak{}askExp} (\texttt{Development.\allowbreak{}AutomorphicArithmetic.\allowbreak{}DedekindLFunctions.\allowbreak{}LandauTheorems.\allowbreak{}ChBalance}, lines 122--197)\par}

{\footnotesize\raggedright\noindent \ref{B2.07}.\quad \texttt{Deduction.\allowbreak{}ch\_\allowbreak{}balance} (\texttt{Development.\allowbreak{}AutomorphicArithmetic.\allowbreak{}DedekindLFunctions.\allowbreak{}LandauTheorems.\allowbreak{}ChBalance}, lines 200--261)\par}

{\footnotesize\raggedright\noindent \ref{B2.08}.\quad \texttt{Deduction.\allowbreak{}CharacterBalanceMoves.\allowbreak{}sum\_\allowbreak{}A} (\texttt{Development.\allowbreak{}AutomorphicArithmetic.\allowbreak{}DedekindLFunctions.\allowbreak{}LandauTheorems.\allowbreak{}ChCore}, lines 28--55)\\
\texttt{Deduction.\allowbreak{}CharacterBalanceMoves.\allowbreak{}sum\_\allowbreak{}Th} (\texttt{Development.\allowbreak{}AutomorphicArithmetic.\allowbreak{}DedekindLFunctions.\allowbreak{}LandauTheorems.\allowbreak{}ChCore}, lines 56--74)\par}

{\footnotesize\raggedright\noindent \ref{B2.09}.\quad \texttt{Deduction.\allowbreak{}CharacterBalanceMoves.\allowbreak{}th\_\allowbreak{}le} (\texttt{Development.\allowbreak{}AutomorphicArithmetic.\allowbreak{}DedekindLFunctions.\allowbreak{}LandauTheorems.\allowbreak{}ChCore}, lines 75--153)\par}

{\footnotesize\raggedright\noindent \ref{B2.10}.\quad \texttt{Deduction.\allowbreak{}CharacterBalanceMoves.\allowbreak{}core} (\texttt{Development.\allowbreak{}AutomorphicArithmetic.\allowbreak{}DedekindLFunctions.\allowbreak{}LandauTheorems.\allowbreak{}ChCore}, lines 283--339)\par}

\section*{Chapter~\ref{ch:p2c2}: Reading off the target sum}

{\footnotesize\raggedright\noindent \ref{B2.11}.\quad \texttt{Deduction.\allowbreak{}CharacterBalanceMoves.\allowbreak{}move} (\texttt{Development.\allowbreak{}AutomorphicArithmetic.\allowbreak{}DedekindLFunctions.\allowbreak{}LandauTheorems.\allowbreak{}ChCore}, lines 154--166)\\
\texttt{Deduction.\allowbreak{}CharacterBalanceMoves.\allowbreak{}triv} (\texttt{Development.\allowbreak{}AutomorphicArithmetic.\allowbreak{}DedekindLFunctions.\allowbreak{}LandauTheorems.\allowbreak{}ChCore}, lines 167--180)\\
\texttt{Formalism.\allowbreak{}LFunctionChainParameters.\allowbreak{}ThickSetMeanBounds.\allowbreak{}DerivInt} (\texttt{Formalism.\allowbreak{}SizeWindowBrackets.\allowbreak{}GrowthBoundGivesRiemannHypothesis}, lines 62--65)\par}

{\footnotesize\raggedright\noindent \ref{B2.12}.\quad \texttt{Deduction.\allowbreak{}CharacterBalanceMoves.\allowbreak{}one\_\allowbreak{}char} (\texttt{Development.\allowbreak{}AutomorphicArithmetic.\allowbreak{}DedekindLFunctions.\allowbreak{}LandauTheorems.\allowbreak{}ChCore}, lines 202--282)\\
\texttt{Deduction.\allowbreak{}CharacterBalanceMoves.\allowbreak{}sum\_\allowbreak{}vv} (\texttt{Development.\allowbreak{}AutomorphicArithmetic.\allowbreak{}DedekindLFunctions.\allowbreak{}LandauTheorems.\allowbreak{}ChCore}, lines 187--200)\par}

{\footnotesize\raggedright\noindent \ref{B2.13}.\quad \texttt{Deduction.\allowbreak{}ch\_\allowbreak{}top\_\allowbreak{}main\_\allowbreak{}eq} (\texttt{Development.\allowbreak{}AutomorphicArithmetic.\allowbreak{}DedekindLFunctions.\allowbreak{}LandauTheorems.\allowbreak{}ChTop}, lines 11--35)\\
\texttt{Deduction.\allowbreak{}ch\_\allowbreak{}top} (\texttt{Development.\allowbreak{}AutomorphicArithmetic.\allowbreak{}DedekindLFunctions.\allowbreak{}LandauTheorems.\allowbreak{}ChTop}, lines 37--72)\par}

{\footnotesize\raggedright\noindent \ref{lem:p2c2-mertens}.\quad \texttt{Formalism.\allowbreak{}LogFactorialMertensBounds.\allowbreak{}mertens\_\allowbreak{}lower} (\texttt{Formalism.\allowbreak{}ZetaNonvanishing.\allowbreak{}GzLipschitz}, lines 86--105)\\
\texttt{Formalism.\allowbreak{}LogFactorialMertensBounds.\allowbreak{}mertens\_\allowbreak{}upper} (\texttt{Formalism.\allowbreak{}ZetaNonvanishing.\allowbreak{}GzLipschitz}, lines 55--84)\par}

{\footnotesize\raggedright\noindent \ref{B2.14}.\quad \texttt{Deduction.\allowbreak{}ch\_\allowbreak{}core} (\texttt{Development.\allowbreak{}AutomorphicArithmetic.\allowbreak{}DedekindLFunctions.\allowbreak{}LandauTheorems.\allowbreak{}ChCore}, lines 343--442)\par}

{\footnotesize\raggedright\noindent \ref{B2.15}.\quad \texttt{UpperTheory.\allowbreak{}Comparison₂.\allowbreak{}prim\_\allowbreak{}read} (\texttt{Development.\allowbreak{}AutomorphicArithmetic.\allowbreak{}PrimitiveCosets.\allowbreak{}PrimRead}, lines 48--97)\\
\texttt{UpperTheory.\allowbreak{}Comparison₂.\allowbreak{}prim\_\allowbreak{}inj} (\texttt{Development.\allowbreak{}AutomorphicArithmetic.\allowbreak{}PrimitiveCosets.\allowbreak{}PrimInj}, lines 6--27)\\
\texttt{Formalism.\allowbreak{}LFunctionChainParameters.\allowbreak{}LHost.\allowbreak{}primitive\_\allowbreak{}value\_\allowbreak{}primitive} (\texttt{Formalism.\allowbreak{}SizeWindowBrackets.\allowbreak{}ChainSlopeWindows.\allowbreak{}GaloisMapComplexConjugation}, lines 233--245)\par}

{\footnotesize\raggedright\noindent \ref{B2.16}.\quad \texttt{UpperTheory.\allowbreak{}Comparison₂.\allowbreak{}card\_\allowbreak{}PsiA} (\texttt{Development.\allowbreak{}AutomorphicArithmetic.\allowbreak{}PrimitiveCosets.\allowbreak{}CardPsiA}, lines 60--63)\\
\texttt{card\_\allowbreak{}PsiA\_\allowbreak{}index} (\texttt{Development.\allowbreak{}AutomorphicArithmetic.\allowbreak{}PrimitiveCosets.\allowbreak{}CardPsiA}, lines 26--58)\\
\texttt{Formalism.\allowbreak{}HC.\allowbreak{}Frame.\allowbreak{}index\_\allowbreak{}A0\_\allowbreak{}eq\_\allowbreak{}o} (\texttt{Formalism.\allowbreak{}SharpRestriction.\allowbreak{}RamificationFrobenius.\allowbreak{}TotientGaloisClasses.\allowbreak{}FrameFacts}, lines 214--218)\par}

{\footnotesize\raggedright\noindent \ref{B2.17}.\quad \texttt{UpperTheory.\allowbreak{}Comparison₂.\allowbreak{}charsA} (\texttt{Development.\allowbreak{}AutomorphicArithmetic.\allowbreak{}PrimitiveCosets.\allowbreak{}charsA}, lines 80--125)\par}

{\footnotesize\raggedright\noindent \ref{B2.18}.\quad \texttt{UpperTheory.\allowbreak{}Comparison₂.\allowbreak{}gmA\_\allowbreak{}eq} (\texttt{Development.\allowbreak{}AutomorphicArithmetic.\allowbreak{}PrimitiveCosets.\allowbreak{}GmAEq}, lines 61--149)\par}

{\footnotesize\raggedright\noindent \ref{B2.19}.\quad \texttt{UpperTheory.\allowbreak{}Comparison₂.\allowbreak{}sigA\_\allowbreak{}le} (\texttt{Development.\allowbreak{}AutomorphicArithmetic.\allowbreak{}PrimitiveCosets.\allowbreak{}SigALe}, lines 38--60)\par}

{\footnotesize\raggedright\noindent \ref{B2.19}.\quad \texttt{UpperTheory.\allowbreak{}Top.\allowbreak{}boundedTw\_\allowbreak{}key} (\texttt{Development.\allowbreak{}AutomorphicArithmetic.\allowbreak{}CharacterBoxes.\allowbreak{}BoundedTw}, lines 5--43)\\
\texttt{UpperTheory.\allowbreak{}T3.\allowbreak{}supplyOfAdmOf} (\texttt{Development.\allowbreak{}AutomorphicArithmetic.\allowbreak{}ShortNorms.\allowbreak{}supplyOfAdmOf}, lines 13--136)\par}

{\footnotesize\raggedright\noindent \ref{def:p2c2-core-at-host}.\quad \texttt{UpperTheory.\allowbreak{}Comparison₂.\allowbreak{}CoreAtH} (\texttt{Development.\allowbreak{}AutomorphicArithmetic.\allowbreak{}PrimitiveCosets.\allowbreak{}CoreAtH}, lines 29--41)\par}

{\footnotesize\raggedright\noindent \ref{B2.20}.\quad \texttt{UpperTheory.\allowbreak{}Comparison₂.\allowbreak{}core\_\allowbreak{}gm} (\texttt{Development.\allowbreak{}AutomorphicArithmetic.\allowbreak{}PrimitiveCosets.\allowbreak{}CoreGm}, lines 12--74)\par}

\section*{Chapter~\ref{ch:p2c3}: Choosing the auxiliary primes \(q\) and \(d\)}

{\footnotesize\raggedright\noindent \ref{B2.21}.\quad \texttt{QH.\allowbreak{}S} (\texttt{Development.\allowbreak{}AutomorphicArithmetic.\allowbreak{}MeanSquareBounds.\allowbreak{}PrimeSetTwoBounds}, lines 13--15)\\
\texttt{QH.\allowbreak{}admitted} (\texttt{Development.\allowbreak{}AutomorphicArithmetic.\allowbreak{}MeanSquareBounds.\allowbreak{}PrimeSetTwoBounds}, lines 18--25)\\
\texttt{QH.\allowbreak{}MeanSq} (\texttt{Development.\allowbreak{}AutomorphicArithmetic.\allowbreak{}MeanSquareBounds.\allowbreak{}PrimeSetTwoBounds}, lines 26--28)\\
\texttt{QH.\allowbreak{}BoxAll} (\texttt{Development.\allowbreak{}AutomorphicArithmetic.\allowbreak{}MeanSquareBounds.\allowbreak{}PrimeSetTwoBounds}, lines 38--40)\\
\texttt{Formalism.\allowbreak{}HermitianLattices.\allowbreak{}ZeroFreeRegionGL2.\allowbreak{}PPset} (\texttt{Formalism.\allowbreak{}GL3MeanSquares.\allowbreak{}SecondEstimateOfIntervalHypothesis}, lines 50--52)\par}

{\footnotesize\raggedright\noindent \ref{B2.22}.\quad \texttt{Formalism.\allowbreak{}HermitianLattices.\allowbreak{}Ext.\allowbreak{}PrimesInProgressions.\allowbreak{}DeuringHeilbronnMB} (\texttt{Formalism.\allowbreak{}GL3MeanSquares.\allowbreak{}GenusMeanAsymptotics.\allowbreak{}JacobiCharacterSums.\allowbreak{}PNTDHBSt}, lines 11--18)\\
\texttt{Formalism.\allowbreak{}HermitianLattices.\allowbreak{}Ext.\allowbreak{}PrimesInProgressions.\allowbreak{}deuringHeilbronnMB\_\allowbreak{}proof'} (\texttt{Formalism.\allowbreak{}GL3MeanSquares.\allowbreak{}GenusMeanAsymptotics.\allowbreak{}JacobiCharacterSums.\allowbreak{}PNTFinal}, lines 20--21)\par}

{\footnotesize\raggedright\noindent \ref{B2.23}.\quad \texttt{Formalism.\allowbreak{}HermitianLattices.\allowbreak{}Ext.\allowbreak{}LogarithmicIntegralBounds.\allowbreak{}PNTAPRelW} (\texttt{Formalism.\allowbreak{}GL3MeanSquares.\allowbreak{}GenusMeanAsymptotics.\allowbreak{}JacobiCharacterSums.\allowbreak{}MeanSquareHypWIff}, lines 38--47)\\
\texttt{Formalism.\allowbreak{}HermitianLattices.\allowbreak{}Ext.\allowbreak{}PrimesInProgressions.\allowbreak{}pnt\_\allowbreak{}progressions\_\allowbreak{}relative\_\allowbreak{}weak\_\allowbreak{}holds} (\texttt{Formalism.\allowbreak{}GL3MeanSquares.\allowbreak{}GenusMeanAsymptotics.\allowbreak{}JacobiCharacterSums.\allowbreak{}PNTFinal}, lines 22--25)\par}

{\footnotesize\raggedright\noindent \ref{B2.23a}.\quad \texttt{Formalism.\allowbreak{}ZeroFreeDirichlet.\allowbreak{}zFRDirichlet} (\texttt{Formalism.\allowbreak{}LandauZeroFreeBounds.\allowbreak{}NoZeroComplex}, lines 32--34)\par}

{\footnotesize\raggedright\noindent \ref{B2.24}.\quad \texttt{Formalism.\allowbreak{}HermitianLattices.\allowbreak{}Ext.\allowbreak{}LogarithmicIntegralBounds.\allowbreak{}page\_\allowbreak{}across} (\texttt{Formalism.\allowbreak{}GL3MeanSquares.\allowbreak{}GenusMeanAsymptotics.\allowbreak{}JacobiCharacterSums.\allowbreak{}MSqPage}, lines 37--66)\\
\texttt{Formalism.\allowbreak{}ZeroFreeDirichlet.\allowbreak{}excBound} (\texttt{Formalism.\allowbreak{}LandauZeroFreeBounds.\allowbreak{}SiegelWalfiszAPOfSiegelLower}, lines 32--33)\par}

{\footnotesize\raggedright\noindent \ref{B2.24}.\quad \texttt{Formalism.\allowbreak{}ZeroFreeDirichlet.\allowbreak{}zFRDirichlet} (\texttt{Formalism.\allowbreak{}LandauZeroFreeBounds.\allowbreak{}NoZeroComplex}, lines 32--34)\par}

{\footnotesize\raggedright\noindent \ref{B2.25}.\quad \texttt{Formalism.\allowbreak{}HermitianLattices.\allowbreak{}Ext.\allowbreak{}LogarithmicIntegralBounds.\allowbreak{}count\_\allowbreak{}a} (\texttt{Formalism.\allowbreak{}GL3MeanSquares.\allowbreak{}GenusMeanAsymptotics.\allowbreak{}JacobiCharacterSums.\allowbreak{}MSqCount}, lines 346--398)\par}

{\footnotesize\raggedright\noindent \ref{B2.26}.\quad \texttt{QH.\allowbreak{}twistK} (\texttt{Development.\allowbreak{}AutomorphicArithmetic.\allowbreak{}MeanSquareBounds.\allowbreak{}twistK}, lines 148--354)\par}

{\footnotesize\raggedright\noindent \ref{B2.27}.\quad \texttt{QH.\allowbreak{}meanSq\_\allowbreak{}abs} (\texttt{Development.\allowbreak{}AutomorphicArithmetic.\allowbreak{}MeanSquareBounds.\allowbreak{}MeanSqAbs}, lines 8--59)\par}

{\footnotesize\raggedright\noindent \ref{B2.28}.\quad \texttt{QH.\allowbreak{}RR\_\allowbreak{}le} (\texttt{Development.\allowbreak{}AutomorphicArithmetic.\allowbreak{}MeanSquareBounds.\allowbreak{}RRLe}, lines 9--100)\\
\texttt{QH.\allowbreak{}sPsq\_\allowbreak{}le} (\texttt{Development.\allowbreak{}AutomorphicArithmetic.\allowbreak{}MeanSquareBounds.\allowbreak{}SPsqLe}, lines 161--172)\\
\texttt{QH.\allowbreak{}exists\_\allowbreak{}primePowerSum\_\allowbreak{}le\_\allowbreak{}mul\_\allowbreak{}log} (\texttt{Development.\allowbreak{}AutomorphicArithmetic.\allowbreak{}MeanSquareBounds.\allowbreak{}SumMinLeDiv}, lines 71--102)\\
\texttt{QH.\allowbreak{}sum\_\allowbreak{}cN\_\allowbreak{}le} (\texttt{Development.\allowbreak{}AutomorphicArithmetic.\allowbreak{}MeanSquareBounds.\allowbreak{}SumCNLe}, lines 8--28)\par}

{\footnotesize\raggedright\noindent \ref{B2.29}.\quad \texttt{QH.\allowbreak{}meanSqK} (\texttt{Development.\allowbreak{}AutomorphicArithmetic.\allowbreak{}MeanSquareBounds.\allowbreak{}meanSqK}, lines 51--222)\par}

{\footnotesize\raggedright\noindent \ref{B2.30}.\quad \texttt{QH.\allowbreak{}bad\_\allowbreak{}share} (\texttt{Development.\allowbreak{}AutomorphicArithmetic.\allowbreak{}MeanSquareBounds.\allowbreak{}BadShare}, lines 251--319)\par}

{\footnotesize\raggedright\noindent \ref{B2.31}.\quad \texttt{Formalism.\allowbreak{}HermitianLattices.\allowbreak{}Ext.\allowbreak{}PrimitiveCharacterCasts.\allowbreak{}ZeroDensityBox} (\texttt{Formalism.\allowbreak{}GL3MeanSquares.\allowbreak{}GenusMeanAsymptotics.\allowbreak{}JacobiCharacterSums.\allowbreak{}CountSt}, lines 7--12)\\
\texttt{Formalism.\allowbreak{}HermitianLattices.\allowbreak{}Ext.\allowbreak{}PrimitiveCharacterCasts.\allowbreak{}ZD.\allowbreak{}zero\_\allowbreak{}density\_\allowbreak{}box} (\texttt{Formalism.\allowbreak{}GL3MeanSquares.\allowbreak{}GenusMeanAsymptotics.\allowbreak{}JacobiCharacterSums.\allowbreak{}CountZAll}, lines 13--16)\par}

{\footnotesize\raggedright\noindent \ref{B2.32}.\quad \texttt{QH.\allowbreak{}count\_\allowbreak{}bad\_\allowbreak{}core} (\texttt{Development.\allowbreak{}AutomorphicArithmetic.\allowbreak{}MeanSquareBounds.\allowbreak{}CountBadCore}, lines 18--60)\\
\texttt{QH.\allowbreak{}count\_\allowbreak{}bad} (\texttt{Development.\allowbreak{}AutomorphicArithmetic.\allowbreak{}MeanSquareBounds.\allowbreak{}CountBad}, lines 9--17)\par}

{\footnotesize\raggedright\noindent \ref{B2.33}.\quad \texttt{QH.\allowbreak{}twist\_\allowbreak{}prim} (\texttt{Development.\allowbreak{}AutomorphicArithmetic.\allowbreak{}MeanSquareBounds.\allowbreak{}TwistPrim}, lines 9--19)\\
\texttt{QH.\allowbreak{}ind\_\allowbreak{}unique} (\texttt{Development.\allowbreak{}AutomorphicArithmetic.\allowbreak{}MeanSquareBounds.\allowbreak{}IndUnique}, lines 7--52)\par}

{\footnotesize\raggedright\noindent \ref{B2.34}.\quad \texttt{QH.\allowbreak{}tailK} (\texttt{Development.\allowbreak{}AutomorphicArithmetic.\allowbreak{}MeanSquareBounds.\allowbreak{}TailKExp}, lines 92--146)\par}

{\footnotesize\raggedright\noindent \ref{B2.35}.\quad \texttt{QH.\allowbreak{}kappaC} (\texttt{Development.\allowbreak{}AutomorphicArithmetic.\allowbreak{}MeanSquareBounds.\allowbreak{}kappaC}, lines 9--63)\\
\texttt{QH.\allowbreak{}kappa4} (\texttt{Development.\allowbreak{}AutomorphicArithmetic.\allowbreak{}MeanSquareBounds.\allowbreak{}Kappa4Loglog}, lines 126--154)\par}

{\footnotesize\raggedright\noindent \ref{B2.36}.\quad \texttt{QH.\allowbreak{}box\_\allowbreak{}of\_\allowbreak{}S} (\texttt{Development.\allowbreak{}AutomorphicArithmetic.\allowbreak{}MeanSquareBounds.\allowbreak{}BoxOfS}, lines 28--61)\par}

{\footnotesize\raggedright\noindent \ref{B2.37}.\quad \texttt{QH.\allowbreak{}stage} (\texttt{Development.\allowbreak{}AutomorphicArithmetic.\allowbreak{}MeanSquareBounds.\allowbreak{}StagePow}, lines 41--171)\par}

{\footnotesize\raggedright\noindent \ref{def:aux-twist-selection}.\quad \texttt{QH.\allowbreak{}GoodQ} (\texttt{Development.\allowbreak{}AutomorphicArithmetic.\allowbreak{}MeanSquareBounds.\allowbreak{}PrimeSetTwoBounds}, lines 48--50)\\
\texttt{QH.\allowbreak{}GoodD} (\texttt{Development.\allowbreak{}AutomorphicArithmetic.\allowbreak{}MeanSquareBounds.\allowbreak{}PrimeSetTwoBounds}, lines 51--54)\par}

{\footnotesize\raggedright\noindent \ref{B2.38}.\quad \texttt{QH.\allowbreak{}stage1} (\texttt{Development.\allowbreak{}AutomorphicArithmetic.\allowbreak{}MeanSquareBounds.\allowbreak{}Stage1Succ}, lines 36--96)\par}

{\footnotesize\raggedright\noindent \ref{B2.39}.\quad \texttt{QH.\allowbreak{}stage2} (\texttt{Development.\allowbreak{}AutomorphicArithmetic.\allowbreak{}MeanSquareBounds.\allowbreak{}CharOf}, lines 61--218)\par}

{\footnotesize\raggedright\noindent \ref{B2.40}.\quad \texttt{QH.\allowbreak{}pair\_\allowbreak{}share} (\texttt{Development.\allowbreak{}AutomorphicArithmetic.\allowbreak{}MeanSquareBounds.\allowbreak{}PairShare}, lines 26--53)\par}

{\footnotesize\raggedright\noindent \ref{B2.41}.\quad \texttt{QH.\allowbreak{}pair\_\allowbreak{}exists} (\texttt{Development.\allowbreak{}AutomorphicArithmetic.\allowbreak{}MeanSquareBounds.\allowbreak{}PairExists}, lines 9--50)\par}

{\footnotesize\raggedright\noindent \ref{B2.42}.\quad \texttt{UpperTheory.\allowbreak{}Top.\allowbreak{}WinPair} (\texttt{Development.\allowbreak{}AutomorphicArithmetic.\allowbreak{}CharacterBoxes.\allowbreak{}ExistsPageN}, lines 27--38)\\
\texttt{UpperTheory.\allowbreak{}Top.\allowbreak{}PairGood} (\texttt{Development.\allowbreak{}AutomorphicArithmetic.\allowbreak{}CharacterBoxes.\allowbreak{}PairGood}, lines 18--27)\par}

{\footnotesize\raggedright\noindent \ref{B2.43}.\quad \texttt{QH.\allowbreak{}pairEx} (\texttt{Development.\allowbreak{}AutomorphicArithmetic.\allowbreak{}MeanSquareBounds.\allowbreak{}pairEx}, lines 112--170)\\
\texttt{UpperTheory.\allowbreak{}Top.\allowbreak{}PairEx} (\texttt{Development.\allowbreak{}AutomorphicArithmetic.\allowbreak{}CharacterBoxes.\allowbreak{}PairGood}, lines 28--31)\par}

\section*{Chapter~\ref{ch:arithmetic-window-host}: Building the arithmetic host of a window}

{\footnotesize\raggedright\noindent \ref{B2.44}.\quad \texttt{UpperTheory.\allowbreak{}Top.\allowbreak{}ordAny} (\texttt{Development.\allowbreak{}AutomorphicArithmetic.\allowbreak{}CharacterBoxes.\allowbreak{}OrdAny}, lines 86--108)\\
\texttt{Formalism.\allowbreak{}HC.\allowbreak{}AnyGrp.\allowbreak{}o\_\allowbreak{}eq} (\texttt{Formalism.\allowbreak{}SharpRestriction.\allowbreak{}RamificationFrobenius.\allowbreak{}TotientGaloisClasses.\allowbreak{}InertiaGaloisSymmetry.\allowbreak{}AnyLevel}, lines 29--35)\\
\texttt{UpperTheory.\allowbreak{}Top.\allowbreak{}kerChiAny} (\texttt{Development.\allowbreak{}AutomorphicArithmetic.\allowbreak{}CharacterBoxes.\allowbreak{}KerChiAny}, lines 104--142)\par}

{\footnotesize\raggedright\noindent \ref{B2.45}.\quad \texttt{UpperTheory.\allowbreak{}Top.\allowbreak{}twistAny} (\texttt{Development.\allowbreak{}AutomorphicArithmetic.\allowbreak{}CharacterBoxes.\allowbreak{}TwistAny}, lines 115--169)\\
\texttt{UpperTheory.\allowbreak{}Top.\allowbreak{}TwistForm} (\texttt{Development.\allowbreak{}AutomorphicArithmetic.\allowbreak{}CharacterBoxes.\allowbreak{}TwistForm}, lines 14--18)\par}

{\footnotesize\raggedright\noindent \ref{B2.46}.\quad \texttt{UpperTheory.\allowbreak{}Top.\allowbreak{}fieldChi} (\texttt{Development.\allowbreak{}AutomorphicArithmetic.\allowbreak{}CharacterBoxes.\allowbreak{}FieldChi}, lines 20--60)\par}

{\footnotesize\raggedright\noindent \ref{B2.47}.\quad \texttt{UpperTheory.\allowbreak{}Top.\allowbreak{}boundedTw} (\texttt{Development.\allowbreak{}AutomorphicArithmetic.\allowbreak{}CharacterBoxes.\allowbreak{}BoundedTw}, lines 45--57)\par}

{\footnotesize\raggedright\noindent \ref{B2.48a}.\quad \texttt{UpperTheory.\allowbreak{}Top.\allowbreak{}winPairs} (\texttt{Development.\allowbreak{}AutomorphicArithmetic.\allowbreak{}CharacterBoxes.\allowbreak{}WinPairs}, lines 19--81)\par}

{\footnotesize\raggedright\noindent \ref{B2.48}.\quad \texttt{UpperTheory.\allowbreak{}Top.\allowbreak{}winPairs} (\texttt{Development.\allowbreak{}AutomorphicArithmetic.\allowbreak{}CharacterBoxes.\allowbreak{}WinPairs}, lines 5--82)\\
\texttt{UpperTheory.\allowbreak{}Top.\allowbreak{}WinDatum} (\texttt{Development.\allowbreak{}AutomorphicArithmetic.\allowbreak{}CharacterBoxes.\allowbreak{}TwistForm}, lines 19--27)\par}

{\footnotesize\raggedright\noindent \ref{B2.49}.\quad \texttt{UpperTheory.\allowbreak{}T3.\allowbreak{}smallOne} (\texttt{Development.\allowbreak{}AutomorphicArithmetic.\allowbreak{}ShortNorms.\allowbreak{}Primes}, lines 7--8)\par}

{\footnotesize\raggedright\noindent \ref{B2.49}.\quad \texttt{GlobalTheory.\allowbreak{}T.\allowbreak{}primesInProgressions\_\allowbreak{}relativeError} (\texttt{Development.\allowbreak{}AbelianVarietyTheory.\allowbreak{}Proofs.\allowbreak{}GlobalInterpolation.\allowbreak{}PrimesInProgressionsRelativeErrorStatement}, lines 8--10)\\
\texttt{GlobalTheory.\allowbreak{}Lin.\allowbreak{}PrimesRel} (\texttt{Development.\allowbreak{}AbelianVarietyTheory.\allowbreak{}Definitions.\allowbreak{}GlobalTwisting.\allowbreak{}ZerosDisc}, lines 35--44)\par}

{\footnotesize\raggedright\noindent \ref{B2.49}.\quad \texttt{GlobalTheory.\allowbreak{}Hk.\allowbreak{}LFam} (\texttt{Development.\allowbreak{}AbelianVarietyTheory.\allowbreak{}Definitions.\allowbreak{}GlobalTwisting.\allowbreak{}LamIdealOfNotPrimePow}, lines 12--41)\par}

{\footnotesize\raggedright\noindent \ref{B2.49a}.\quad \texttt{GlobalTheory.\allowbreak{}Hk.\allowbreak{}ShortInt.\allowbreak{}zeros} (\texttt{Development.\allowbreak{}AbelianVarietyTheory.\allowbreak{}Proofs.\allowbreak{}GlobalInterpolation.\allowbreak{}TrivialClassShortInterval}, lines 20--102)\par}

{\footnotesize\raggedright\noindent \ref{B2.49b}.\quad \texttt{GlobalTheory.\allowbreak{}Hk.\allowbreak{}PsiExc.\allowbreak{}psi\_\allowbreak{}exceptional} (\texttt{Development.\allowbreak{}AbelianVarietyTheory.\allowbreak{}Proofs.\allowbreak{}GlobalInterpolation.\allowbreak{}RayClassFamilyPsiWithExceptionalZero}, lines 438--503)\par}

{\footnotesize\raggedright\noindent \ref{B2.49c}.\quad \texttt{GlobalTheory.\allowbreak{}Hk.\allowbreak{}ShortInt.\allowbreak{}trivialClass\_\allowbreak{}shortInterval} (\texttt{Development.\allowbreak{}AbelianVarietyTheory.\allowbreak{}Proofs.\allowbreak{}GlobalInterpolation.\allowbreak{}TrivialClassShortInterval}, lines 218--282)\\
\texttt{GlobalTheory.\allowbreak{}T.\allowbreak{}exists\_\allowbreak{}lam\_\allowbreak{}zero\_\allowbreak{}sum\_\allowbreak{}shortInterval\_\allowbreak{}ge} (\texttt{Development.\allowbreak{}AbelianVarietyTheory.\allowbreak{}Proofs.\allowbreak{}GlobalInterpolation.\allowbreak{}TrivialClassShortIntervalStatement}, lines 8--9)\par}

{\footnotesize\raggedright\noindent \ref{B2.50}.\quad \texttt{UpperTheory.\allowbreak{}T3.\allowbreak{}shortNorm} (\texttt{Development.\allowbreak{}AutomorphicArithmetic.\allowbreak{}ShortNorms.\allowbreak{}Primes}, lines 5--6)\\
\texttt{GlobalTheory.\allowbreak{}T.\allowbreak{}normPrimesImagQuad\_\allowbreak{}short\_\allowbreak{}interval\_\allowbreak{}lower\_\allowbreak{}bound} (\texttt{Development.\allowbreak{}AbelianVarietyTheory.\allowbreak{}Proofs.\allowbreak{}GlobalInterpolation.\allowbreak{}NormPrimesImagQuadShortIntervalLowerBoundStatement}, lines 8--10)\\
\texttt{GlobalTheory.\allowbreak{}Mem.\allowbreak{}An.\allowbreak{}NormOne} (\texttt{Development.\allowbreak{}AbelianVarietyTheory.\allowbreak{}Definitions.\allowbreak{}Foundations.\allowbreak{}RayPrimes}, lines 40--42)\par}

{\footnotesize\raggedright\noindent \ref{B2.50}.\quad \texttt{GlobalTheory.\allowbreak{}T.\allowbreak{}exists\_\allowbreak{}lFam\_\allowbreak{}normOne\_\allowbreak{}of\_\allowbreak{}mod\_\allowbreak{}four\_\allowbreak{}eq\_\allowbreak{}three} (\texttt{Development.\allowbreak{}AbelianVarietyTheory.\allowbreak{}Proofs.\allowbreak{}GlobalInterpolation.\allowbreak{}ExistsLFamRayStatement}, lines 8--9)\\
\texttt{GlobalTheory.\allowbreak{}Hk.\allowbreak{}normPrimes\_\allowbreak{}short\_\allowbreak{}interval\_\allowbreak{}lower} (\texttt{Development.\allowbreak{}AbelianVarietyTheory.\allowbreak{}Proofs.\allowbreak{}GlobalInterpolation.\allowbreak{}NormPrimesImagQuadShortIntervalLowerBound}, lines 98--148)\par}

{\footnotesize\raggedright\noindent \ref{B2.51}.\quad \texttt{UpperTheory.\allowbreak{}T3.\allowbreak{}levelData} (\texttt{Development.\allowbreak{}AutomorphicArithmetic.\allowbreak{}ShortNorms.\allowbreak{}levelData}, lines 36--78)\\
\texttt{UpperTheory.\allowbreak{}T3.\allowbreak{}typePrimesN} (\texttt{Development.\allowbreak{}AutomorphicArithmetic.\allowbreak{}ShortNorms.\allowbreak{}typePrimesN}, lines 78--130)\par}

{\footnotesize\raggedright\noindent \ref{B2.52}.\quad \texttt{UpperTheory.\allowbreak{}Top.\allowbreak{}realBox} (\texttt{Development.\allowbreak{}AutomorphicArithmetic.\allowbreak{}CharacterBoxes.\allowbreak{}RealBox}, lines 99--113)\\
\texttt{QH.\allowbreak{}thinH} (\texttt{Development.\allowbreak{}AutomorphicArithmetic.\allowbreak{}MeanSquareBounds.\allowbreak{}GaussSumFunctionEq}, lines 55--108)\\
\texttt{Formalism.\allowbreak{}HC.\allowbreak{}HostRaw.\allowbreak{}ChebyshevCharacterBound} (\texttt{Formalism.\allowbreak{}SharpRestriction.\allowbreak{}RamificationFrobenius.\allowbreak{}TotientGaloisClasses.\allowbreak{}Analytic}, lines 145--148)\par}

{\footnotesize\raggedright\noindent \ref{B2.53}.\quad \texttt{UpperTheory.\allowbreak{}BernoulliSumsFrobeniusOrbits.\allowbreak{}thickFrame2} (\texttt{Development.\allowbreak{}AutomorphicArithmetic.\allowbreak{}ThickFrames.\allowbreak{}ThickFrame2Large}, lines 29--199)\\
\texttt{Formalism.\allowbreak{}FrobeniusOrbits.\allowbreak{}NoRel} (\texttt{Formalism.\allowbreak{}FrobeniusOrbits.\allowbreak{}HostThickExists}, lines 13--17)\par}

{\footnotesize\raggedright\noindent \ref{B2.54}.\quad \texttt{UpperTheory.\allowbreak{}BernoulliSumsFrobeniusOrbits.\allowbreak{}topsChoiceH} (\texttt{Development.\allowbreak{}AutomorphicArithmetic.\allowbreak{}ThickFrames.\allowbreak{}topsChoiceH}, lines 30--68)\par}

{\footnotesize\raggedright\noindent \ref{B2.55}.\quad \texttt{UpperTheory.\allowbreak{}Top.\allowbreak{}pairRows} (\texttt{Development.\allowbreak{}AutomorphicArithmetic.\allowbreak{}CharacterBoxes.\allowbreak{}PairRows}, lines 334--388)\par}

{\footnotesize\raggedright\noindent \ref{B2.56}.\quad \texttt{UpperTheory.\allowbreak{}Top.\allowbreak{}hostOfField} (\texttt{Development.\allowbreak{}AutomorphicArithmetic.\allowbreak{}CharacterBoxes.\allowbreak{}HostOfField}, lines 8--16)\\
\texttt{UpperTheory.\allowbreak{}T3.\allowbreak{}hostOfFieldOf} (\texttt{Development.\allowbreak{}AutomorphicArithmetic.\allowbreak{}ShortNorms.\allowbreak{}hostOfFieldOf}, lines 159--256)\par}

{\footnotesize\raggedright\noindent \ref{B2.57}.\quad \texttt{UpperTheory.\allowbreak{}Top.\allowbreak{}admissibleClassesTwistedChar} (\texttt{Development.\allowbreak{}AutomorphicArithmetic.\allowbreak{}CoprimeRanks.\allowbreak{}AdmChi}, lines 13--21)\par}

{\footnotesize\raggedright\noindent \ref{B2.58}.\quad \texttt{Formalism.\allowbreak{}HC.\allowbreak{}EAtH} (\texttt{Formalism.\allowbreak{}SharpRestriction.\allowbreak{}RamificationFrobenius.\allowbreak{}TotientGaloisClasses.\allowbreak{}ExplicitFormulaErrorOfHostConditions}, lines 144--152)\\
\texttt{UpperTheory.\allowbreak{}Top.\allowbreak{}GChi} (\texttt{Development.\allowbreak{}AutomorphicArithmetic.\allowbreak{}CharacterBoxes.\allowbreak{}Supply}, lines 42--44)\par}

{\footnotesize\raggedright\noindent \ref{B2.58}.\quad \texttt{GlobalTheory.\allowbreak{}EX.\allowbreak{}HostExW} (\texttt{Development.\allowbreak{}AbelianVarietyTheory.\allowbreak{}Definitions.\allowbreak{}GlobalTwisting.\allowbreak{}ThinScalesIffHostBound}, lines 120--124)\par}

{\footnotesize\raggedright\noindent \ref{B2.59}.\quad \texttt{UpperTheory.\allowbreak{}T3.\allowbreak{}supplyOfAdmOf} (\texttt{Development.\allowbreak{}AutomorphicArithmetic.\allowbreak{}ShortNorms.\allowbreak{}supplyOfAdmOf}, lines 13--136)\\
\texttt{UpperTheory.\allowbreak{}Top.\allowbreak{}supplyOfAdm} (\texttt{Development.\allowbreak{}AutomorphicArithmetic.\allowbreak{}CharacterBoxes.\allowbreak{}HostOfField}, lines 17--24)\par}

{\footnotesize\raggedright\noindent \ref{B2.61}.\quad \texttt{UpperTheory.\allowbreak{}Top.\allowbreak{}supplyChi} (\texttt{Development.\allowbreak{}AutomorphicArithmetic.\allowbreak{}CharacterBoxes.\allowbreak{}SupplyChi}, lines 8--17)\par}

\section*{Chapter~\ref{ch:p3c1}: The local kernel and the differential operator}

{\footnotesize\raggedright\noindent \ref{B3.01a}.\quad \texttt{Formalism.\allowbreak{}MainNotions.\allowbreak{}G.\allowbreak{}liftWt} (\texttt{Formalism.\allowbreak{}SharpRestriction.\allowbreak{}GrowthBounds.\allowbreak{}DataGChain}, lines 24--26)\\
\texttt{Formalism.\allowbreak{}MainNotions.\allowbreak{}G.\allowbreak{}deltaL} (\texttt{Formalism.\allowbreak{}SharpRestriction.\allowbreak{}GrowthBounds.\allowbreak{}DataG}, lines 11--15)\\
\texttt{Formalism.\allowbreak{}GradedWeylDivisibility.\allowbreak{}weyl} (\texttt{Formalism.\allowbreak{}GradedWeylDivisibility.\allowbreak{}StatementA}, lines 71--74)\\
\texttt{Formalism.\allowbreak{}GradedWeylDivisibility.\allowbreak{}rowGen} (\texttt{Formalism.\allowbreak{}GradedWeylDivisibility.\allowbreak{}StatementA}, lines 35--39)\\
\texttt{Formalism.\allowbreak{}MainNotions.\allowbreak{}G.\allowbreak{}Places.\allowbreak{}HtupW} (\texttt{Formalism.\allowbreak{}SharpRestriction.\allowbreak{}GrowthBounds.\allowbreak{}DataG}, lines 79--81)\\
\texttt{Formalism.\allowbreak{}L2.\allowbreak{}PairingIndexBounds.\allowbreak{}Ent.\allowbreak{}delta\_\allowbreak{}zero\_\allowbreak{}eq\_\allowbreak{}twisted\_\allowbreak{}delta} (\texttt{Formalism.\allowbreak{}LocalToGlobal.\allowbreak{}PairingIndexBounds.\allowbreak{}RowActionCoefficients.\allowbreak{}AdmissibleInterpolation.\allowbreak{}Bridge}, lines 45--63)\par}

{\footnotesize\raggedright\noindent \ref{B3.01}.\quad \texttt{Formalism.\allowbreak{}MainNotions.\allowbreak{}FormalGroupSystem.\allowbreak{}Std} (\texttt{Formalism.\allowbreak{}SharpRestriction.\allowbreak{}GrowthBounds.\allowbreak{}FormalGroupSystem}, lines 89--102)\\
\texttt{Formalism.\allowbreak{}MainNotions.\allowbreak{}FormalGroupSystem.\allowbreak{}Sys} (\texttt{Formalism.\allowbreak{}SharpRestriction.\allowbreak{}GrowthBounds.\allowbreak{}FormalGroupSystem}, lines 264--278)\\
\texttt{Formalism.\allowbreak{}MainNotions.\allowbreak{}FormalGroupSystem.\allowbreak{}Sys.\allowbreak{}vsc} (\texttt{Formalism.\allowbreak{}SharpRestriction.\allowbreak{}GrowthBounds.\allowbreak{}FormalGroupSystem}, lines 309--310)\\
\texttt{Formalism.\allowbreak{}MainNotions.\allowbreak{}FormalGroupSystem.\allowbreak{}Sys.\allowbreak{}vGe} (\texttt{Formalism.\allowbreak{}SharpRestriction.\allowbreak{}GrowthBounds.\allowbreak{}FormalGroupSystem}, lines 311--312)\\
\texttt{Formalism.\allowbreak{}MainNotions.\allowbreak{}FormalGroupSystem.\allowbreak{}Pend} (\texttt{Formalism.\allowbreak{}SharpRestriction.\allowbreak{}GrowthBounds.\allowbreak{}FormalGroupSystem}, lines 447--494)\par}

{\footnotesize\raggedright\noindent \ref{B3.02}.\quad \texttt{Formalism.\allowbreak{}MainNotions.\allowbreak{}FormalGroupSystem.\allowbreak{}NN} (\texttt{Formalism.\allowbreak{}SharpRestriction.\allowbreak{}GrowthBounds.\allowbreak{}FormalGroupSystem}, lines 26--27)\\
\texttt{Formalism.\allowbreak{}MainNotions.\allowbreak{}FormalGroupSystem.\allowbreak{}nu} (\texttt{Formalism.\allowbreak{}SharpRestriction.\allowbreak{}GrowthBounds.\allowbreak{}FormalGroupSystem}, lines 28--32)\\
\texttt{Formalism.\allowbreak{}MainNotions.\allowbreak{}FormalGroupSystem.\allowbreak{}al} (\texttt{Formalism.\allowbreak{}SharpRestriction.\allowbreak{}GrowthBounds.\allowbreak{}FormalGroupSystem}, lines 33--35)\\
\texttt{Formalism.\allowbreak{}MainNotions.\allowbreak{}FormalGroupSystem.\allowbreak{}cap} (\texttt{Formalism.\allowbreak{}SharpRestriction.\allowbreak{}GrowthBounds.\allowbreak{}FormalGroupSystem}, lines 36--39)\\
\texttt{Formalism.\allowbreak{}MainNotions.\allowbreak{}FormalGroupSystem.\allowbreak{}lsz} (\texttt{Formalism.\allowbreak{}SharpRestriction.\allowbreak{}GrowthBounds.\allowbreak{}FormalGroupSystem}, lines 40--43)\par}

{\footnotesize\raggedright\noindent \ref{B3.03}.\quad \texttt{Formalism.\allowbreak{}MainNotions.\allowbreak{}FormalGroupSystem.\allowbreak{}lgp} (\texttt{Formalism.\allowbreak{}SharpRestriction.\allowbreak{}GrowthBounds.\allowbreak{}FormalGroupSystem}, lines 44--45)\\
\texttt{Formalism.\allowbreak{}MainNotions.\allowbreak{}FormalGroupSystem.\allowbreak{}de} (\texttt{Formalism.\allowbreak{}SharpRestriction.\allowbreak{}GrowthBounds.\allowbreak{}FormalGroupSystem}, lines 46--50)\\
\texttt{Formalism.\allowbreak{}InvariantFactorPaths.\allowbreak{}G7.\allowbreak{}deLoss} (\texttt{Formalism.\allowbreak{}ThetaLogCoordinates.\allowbreak{}UniformThetaSubstitution.\allowbreak{}StatementAB}, lines 100--102)\par}

{\footnotesize\raggedright\noindent \ref{B3.04}.\quad \texttt{Formalism.\allowbreak{}MainNotions.\allowbreak{}FormalGroupSystem.\allowbreak{}logF} (\texttt{Formalism.\allowbreak{}SharpRestriction.\allowbreak{}GrowthBounds.\allowbreak{}FormalGroupSystem}, lines 126--128)\\
\texttt{Formalism.\allowbreak{}MainNotions.\allowbreak{}FormalGroupSystem.\allowbreak{}inParam} (\texttt{Formalism.\allowbreak{}SharpRestriction.\allowbreak{}GrowthBounds.\allowbreak{}FormalGroupSystem}, lines 129--131)\\
\texttt{Formalism.\allowbreak{}MainNotions.\allowbreak{}FormalGroupSystem.\allowbreak{}inParamK} (\texttt{Formalism.\allowbreak{}SharpRestriction.\allowbreak{}GrowthBounds.\allowbreak{}FormalGroupSystem}, lines 132--135)\\
\texttt{Formalism.\allowbreak{}MainNotions.\allowbreak{}FormalGroupSystem.\allowbreak{}expS} (\texttt{Formalism.\allowbreak{}SharpRestriction.\allowbreak{}GrowthBounds.\allowbreak{}FormalGroupSystem}, lines 136--138)\\
\texttt{Formalism.\allowbreak{}MainNotions.\allowbreak{}FormalGroupSystem.\allowbreak{}Simeq} (\texttt{Formalism.\allowbreak{}SharpRestriction.\allowbreak{}GrowthBounds.\allowbreak{}FormalGroupSystem}, lines 139--143)\par}

{\footnotesize\raggedright\noindent \ref{B3.05}.\quad \texttt{Formalism.\allowbreak{}MainNotions.\allowbreak{}FormalGroupSystem.\allowbreak{}BeProp} (\texttt{Formalism.\allowbreak{}SharpRestriction.\allowbreak{}GrowthBounds.\allowbreak{}FormalGroupSystem}, lines 147--159)\\
\texttt{Formalism.\allowbreak{}MainNotions.\allowbreak{}FormalGroupSystem.\allowbreak{}ExtensionLocalCondition} (\texttt{Formalism.\allowbreak{}SharpRestriction.\allowbreak{}GrowthBounds.\allowbreak{}FormalGroupSystem}, lines 162--171)\\
\texttt{Formalism.\allowbreak{}MainNotions.\allowbreak{}FormalGroupSystem.\allowbreak{}ExtensionGeneratorCondition} (\texttt{Formalism.\allowbreak{}SharpRestriction.\allowbreak{}GrowthBounds.\allowbreak{}FormalGroupSystem}, lines 237--244)\\
\texttt{Formalism.\allowbreak{}MainNotions.\allowbreak{}FormalGroupSystem.\allowbreak{}adm2} (\texttt{Formalism.\allowbreak{}SharpRestriction.\allowbreak{}GrowthBounds.\allowbreak{}FormalGroupSystem}, lines 548--551)\\
\texttt{Formalism.\allowbreak{}MainNotions.\allowbreak{}FormalGroupSystem.\allowbreak{}adm2b} (\texttt{Formalism.\allowbreak{}SharpRestriction.\allowbreak{}GrowthBounds.\allowbreak{}FormalGroupSystem}, lines 552--555)\par}

{\footnotesize\raggedright\noindent \ref{B3.05b}.\quad \texttt{Formalism.\allowbreak{}L2.\allowbreak{}KerInert.\allowbreak{}Jet.\allowbreak{}beJet\_\allowbreak{}of\_\allowbreak{}prop} (\texttt{Formalism.\allowbreak{}LocalToGlobal.\allowbreak{}InertKernelTheory.\allowbreak{}InertCoprimeLemmas.\allowbreak{}HurwitzSubstituteTwoSeries}, lines 999--1002)\par}

{\footnotesize\raggedright\noindent \ref{def:p3c1-rational-weight}.\quad \texttt{Formalism.\allowbreak{}MainNotions.\allowbreak{}G.\allowbreak{}liftWt} (\texttt{Formalism.\allowbreak{}SharpRestriction.\allowbreak{}GrowthBounds.\allowbreak{}DataGChain}, lines 23--25)\\
\texttt{Formalism.\allowbreak{}GradedWeylDivisibility.\allowbreak{}dualWeyl} (\texttt{Formalism.\allowbreak{}GradedWeylDivisibility.\allowbreak{}StatementA}, lines 24--30)\\
\texttt{Formalism.\allowbreak{}MainNotions.\allowbreak{}G.\allowbreak{}Vsp} (\texttt{Formalism.\allowbreak{}SharpRestriction.\allowbreak{}GrowthBounds.\allowbreak{}UnramifiedQuadraticSetting}, lines 84--87)\\
\texttt{Formalism.\allowbreak{}MainNotions.\allowbreak{}G.\allowbreak{}Places.\allowbreak{}Vtup} (\texttt{Formalism.\allowbreak{}SharpRestriction.\allowbreak{}GrowthBounds.\allowbreak{}DataG}, lines 87--90)\par}

{\footnotesize\raggedright\noindent \ref{B3.05aa}.\quad \texttt{Formalism.\allowbreak{}L2.\allowbreak{}G.\allowbreak{}Pend.\allowbreak{}subD} (\texttt{Formalism.\allowbreak{}LocalToGlobal.\allowbreak{}GroupLattices.\allowbreak{}PendingLemmas.\allowbreak{}DetYtMatNeZero}, lines 63--65)\\
\texttt{Formalism.\allowbreak{}L2.\allowbreak{}G.\allowbreak{}Pend.\allowbreak{}numD} (\texttt{Formalism.\allowbreak{}LocalToGlobal.\allowbreak{}GroupLattices.\allowbreak{}PendingLemmas.\allowbreak{}DetYtMatNeZero}, lines 75--79)\\
\texttt{Formalism.\allowbreak{}L2.\allowbreak{}G.\allowbreak{}Pend.\allowbreak{}divD} (\texttt{Formalism.\allowbreak{}LocalToGlobal.\allowbreak{}GroupLattices.\allowbreak{}PendingLemmas.\allowbreak{}DetYtMatNeZero}, lines 87--92)\\
\texttt{Formalism.\allowbreak{}L2.\allowbreak{}G.\allowbreak{}Pend.\allowbreak{}polDel} (\texttt{Formalism.\allowbreak{}LocalToGlobal.\allowbreak{}GroupLattices.\allowbreak{}PendingLemmas.\allowbreak{}DetYtMatNeZero}, lines 100--102)\\
\texttt{Formalism.\allowbreak{}L2.\allowbreak{}G.\allowbreak{}Pend.\allowbreak{}delF} (\texttt{Formalism.\allowbreak{}LocalToGlobal.\allowbreak{}GroupLattices.\allowbreak{}PendingLemmas.\allowbreak{}DetYtMatNeZero}, lines 165--171)\par}

{\footnotesize\raggedright\noindent \ref{B3.05a}.\quad \texttt{Formalism.\allowbreak{}MainNotions.\allowbreak{}FormalGroupSystem.\allowbreak{}admitted} (\texttt{Formalism.\allowbreak{}SharpRestriction.\allowbreak{}GrowthBounds.\allowbreak{}FormalGroupSystem}, lines 771--773)\par}

{\footnotesize\raggedright\noindent \ref{B3.06}.\quad \texttt{Formalism.\allowbreak{}MainNotions.\allowbreak{}FormalGroupSystem.\allowbreak{}DworkLemma} (\texttt{Formalism.\allowbreak{}SharpRestriction.\allowbreak{}GrowthBounds.\allowbreak{}FormalGroupSystem}, lines 225--230)\\
\texttt{Formalism.\allowbreak{}InvariantFactorPaths.\allowbreak{}G7.\allowbreak{}dwork\_\allowbreak{}multivariable} (\texttt{Formalism.\allowbreak{}ThetaLogCoordinates.\allowbreak{}UniformThetaSubstitution.\allowbreak{}FrobeniusLiftNorms.\allowbreak{}LocalComponentGL}, lines 448--480)\par}

{\footnotesize\raggedright\noindent \ref{B3.06}.\quad \texttt{Formalism.\allowbreak{}MainNotions.\allowbreak{}FormalGroupSystem.\allowbreak{}Pf.\allowbreak{}sx\_\allowbreak{}dwork} (\texttt{Formalism.\allowbreak{}SharpRestriction.\allowbreak{}GrowthBounds.\allowbreak{}MultisetValuation.\allowbreak{}PositivityConditionHolds}, lines 107--113)\par}

{\footnotesize\raggedright\noindent \ref{B3.07}.\quad \texttt{Formalism.\allowbreak{}MainNotions.\allowbreak{}FormalGroupSystem.\allowbreak{}sSer} (\texttt{Formalism.\allowbreak{}SharpRestriction.\allowbreak{}GrowthBounds.\allowbreak{}FormalGroupSystem}, lines 200--204)\\
\texttt{Formalism.\allowbreak{}MainNotions.\allowbreak{}FormalGroupSystem.\allowbreak{}hhS} (\texttt{Formalism.\allowbreak{}SharpRestriction.\allowbreak{}GrowthBounds.\allowbreak{}FormalGroupSystem}, lines 205--209)\\
\texttt{Formalism.\allowbreak{}MainNotions.\allowbreak{}FormalGroupSystem.\allowbreak{}InClass} (\texttt{Formalism.\allowbreak{}SharpRestriction.\allowbreak{}GrowthBounds.\allowbreak{}FormalGroupSystem}, lines 210--215)\\
\texttt{Formalism.\allowbreak{}InvariantFactorPaths.\allowbreak{}G7.\allowbreak{}InClass} (\texttt{Formalism.\allowbreak{}ThetaLogCoordinates.\allowbreak{}UniformThetaSubstitution.\allowbreak{}OddLogRelation}, lines 76--80)\par}

{\footnotesize\raggedright\noindent \ref{B3.08}.\quad \texttt{Formalism.\allowbreak{}MainNotions.\allowbreak{}FormalGroupSystem.\allowbreak{}Pf.\allowbreak{}restBeta} (\texttt{Formalism.\allowbreak{}SharpRestriction.\allowbreak{}GrowthBounds.\allowbreak{}MultisetValuation.\allowbreak{}BiadditiveSecondVariable}, lines 137--154)\\
\texttt{Formalism.\allowbreak{}MainNotions.\allowbreak{}FormalGroupSystem.\allowbreak{}Pf.\allowbreak{}restGc} (\texttt{Formalism.\allowbreak{}SharpRestriction.\allowbreak{}GrowthBounds.\allowbreak{}MultisetValuation.\allowbreak{}CommonEigenvectorThird}, lines 21--23)\\
\texttt{Formalism.\allowbreak{}MainNotions.\allowbreak{}FormalGroupSystem.\allowbreak{}Pf.\allowbreak{}restSExp\_\allowbreak{}of} (\texttt{Formalism.\allowbreak{}SharpRestriction.\allowbreak{}GrowthBounds.\allowbreak{}MultisetValuation.\allowbreak{}DworkExponentialRestrictions}, lines 95--106)\\
\texttt{Formalism.\allowbreak{}MainNotions.\allowbreak{}FormalGroupSystem.\allowbreak{}Pf.\allowbreak{}restTwoSums\_\allowbreak{}of} (\texttt{Formalism.\allowbreak{}SharpRestriction.\allowbreak{}GrowthBounds.\allowbreak{}MultisetValuation.\allowbreak{}TripleSumComm}, lines 205--361)\\
\texttt{Formalism.\allowbreak{}MainNotions.\allowbreak{}FormalGroupSystem.\allowbreak{}Pf.\allowbreak{}piFact} (\texttt{Formalism.\allowbreak{}SharpRestriction.\allowbreak{}GrowthBounds.\allowbreak{}MultisetValuation.\allowbreak{}NatCastInvMemRp}, lines 52--71)\par}

{\footnotesize\raggedright\noindent \ref{B3.09}.\quad \texttt{Formalism.\allowbreak{}MainNotions.\allowbreak{}FormalGroupSystem.\allowbreak{}Pf.\allowbreak{}theta\_\allowbreak{}kernel\_\allowbreak{}expansion} (\texttt{Formalism.\allowbreak{}SharpRestriction.\allowbreak{}GrowthBounds.\allowbreak{}MultisetValuation.\allowbreak{}CommonEigenvectorThird}, lines 32--34)\\
\texttt{Formalism.\allowbreak{}MainNotions.\allowbreak{}FormalGroupSystem.\allowbreak{}Pf.\allowbreak{}expansions\_\allowbreak{}multiply} (\texttt{Formalism.\allowbreak{}SharpRestriction.\allowbreak{}GrowthBounds.\allowbreak{}MultisetValuation.\allowbreak{}CommonEigenvectorThird}, lines 35--37)\par}

{\footnotesize\raggedright\noindent \ref{B3.10}.\quad \texttt{Formalism.\allowbreak{}MainNotions.\allowbreak{}FormalGroupSystem.\allowbreak{}Sys.\allowbreak{}colHt} (\texttt{Formalism.\allowbreak{}SharpRestriction.\allowbreak{}GrowthBounds.\allowbreak{}FormalGroupSystem}, lines 398--400)\\
\texttt{Formalism.\allowbreak{}MainNotions.\allowbreak{}FormalGroupSystem.\allowbreak{}Sys.\allowbreak{}colFac} (\texttt{Formalism.\allowbreak{}SharpRestriction.\allowbreak{}GrowthBounds.\allowbreak{}FormalGroupSystem}, lines 405--411)\\
\texttt{Formalism.\allowbreak{}MainNotions.\allowbreak{}FormalGroupSystem.\allowbreak{}Sys.\allowbreak{}colExp} (\texttt{Formalism.\allowbreak{}SharpRestriction.\allowbreak{}GrowthBounds.\allowbreak{}FormalGroupSystem}, lines 412--416)\\
\texttt{Formalism.\allowbreak{}MainNotions.\allowbreak{}FormalGroupSystem.\allowbreak{}Sys.\allowbreak{}PI} (\texttt{Formalism.\allowbreak{}SharpRestriction.\allowbreak{}GrowthBounds.\allowbreak{}FormalGroupSystem}, lines 417--420)\\
\texttt{Formalism.\allowbreak{}MainNotions.\allowbreak{}FormalGroupSystem.\allowbreak{}Sys.\allowbreak{}PIv} (\texttt{Formalism.\allowbreak{}SharpRestriction.\allowbreak{}GrowthBounds.\allowbreak{}FormalGroupSystem}, lines 421--423)\par}

{\footnotesize\raggedright\noindent \ref{B3.10b}.\quad \texttt{Formalism.\allowbreak{}L2.\allowbreak{}PairingIndexBounds.\allowbreak{}Ent.\allowbreak{}MiddleFrameLemmas.\allowbreak{}blkX} (\texttt{Formalism.\allowbreak{}LocalToGlobal.\allowbreak{}PairingIndexBounds.\allowbreak{}RowActionCoefficients.\allowbreak{}AdmissibleInterpolation.\allowbreak{}ProdItePow}, lines 372--414)\par}

{\footnotesize\raggedright\noindent \ref{B3.10a}.\quad \texttt{Formalism.\allowbreak{}L2.\allowbreak{}PairingIndexBounds.\allowbreak{}Ent.\allowbreak{}admissible\_\allowbreak{}general\_\allowbreak{}lattice\_\allowbreak{}tuple} (\texttt{Formalism.\allowbreak{}LocalToGlobal.\allowbreak{}PairingIndexBounds.\allowbreak{}RowActionCoefficients.\allowbreak{}AdmissibleInterpolation.\allowbreak{}AdmissibleGeneralLatticeTuple}, lines 17--80)\\
\texttt{Formalism.\allowbreak{}L2.\allowbreak{}PairingIndexBounds.\allowbreak{}Ent.\allowbreak{}pending\_\allowbreak{}tuple\_\allowbreak{}admissible} (\texttt{Formalism.\allowbreak{}LocalToGlobal.\allowbreak{}PairingIndexBounds.\allowbreak{}RowActionCoefficients.\allowbreak{}LatticeShapeMiddleHolds}, lines 26--28)\par}

{\footnotesize\raggedright\noindent \ref{B3.11}.\quad \texttt{Formalism.\allowbreak{}MainNotions.\allowbreak{}FormalGroupSystem.\allowbreak{}applyD} (\texttt{Formalism.\allowbreak{}SharpRestriction.\allowbreak{}GrowthBounds.\allowbreak{}FormalGroupSystem}, lines 256--260)\par}

{\footnotesize\raggedright\noindent \ref{B3.12}.\quad \texttt{Formalism.\allowbreak{}MainNotions.\allowbreak{}FormalGroupSystem.\allowbreak{}Pend.\allowbreak{}PhAt} (\texttt{Formalism.\allowbreak{}SharpRestriction.\allowbreak{}GrowthBounds.\allowbreak{}FormalGroupSystem}, lines 506--509)\\
\texttt{Formalism.\allowbreak{}MainNotions.\allowbreak{}FormalGroupSystem.\allowbreak{}Pend.\allowbreak{}geC} (\texttt{Formalism.\allowbreak{}SharpRestriction.\allowbreak{}GrowthBounds.\allowbreak{}FormalGroupSystem}, lines 522--524)\\
\texttt{Formalism.\allowbreak{}MainNotions.\allowbreak{}FormalGroupSystem.\allowbreak{}Pend.\allowbreak{}LowAt} (\texttt{Formalism.\allowbreak{}SharpRestriction.\allowbreak{}GrowthBounds.\allowbreak{}FormalGroupSystem}, lines 529--534)\par}

\section*{Chapter~\ref{ch:b3-entry-estimates}: Entry estimates: factorials against denominators}

{\footnotesize\raggedright\noindent \ref{B3.13}.\quad \texttt{Formalism.\allowbreak{}MainNotions.\allowbreak{}FormalGroupSystem.\allowbreak{}dSet} (\texttt{Formalism.\allowbreak{}SharpRestriction.\allowbreak{}GrowthBounds.\allowbreak{}FormalGroupSystem}, lines 63--64)\\
\texttt{Formalism.\allowbreak{}MainNotions.\allowbreak{}FormalGroupSystem.\allowbreak{}gA} (\texttt{Formalism.\allowbreak{}SharpRestriction.\allowbreak{}GrowthBounds.\allowbreak{}FormalGroupSystem}, lines 66--70)\\
\texttt{Formalism.\allowbreak{}MainNotions.\allowbreak{}FormalGroupSystem.\allowbreak{}Sys.\allowbreak{}g0v} (\texttt{Formalism.\allowbreak{}SharpRestriction.\allowbreak{}GrowthBounds.\allowbreak{}FormalGroupSystem}, lines 325--325)\par}

{\footnotesize\raggedright\noindent \ref{B3.14}.\quad \texttt{Formalism.\allowbreak{}MainNotions.\allowbreak{}FormalGroupSystem.\allowbreak{}gB} (\texttt{Formalism.\allowbreak{}SharpRestriction.\allowbreak{}GrowthBounds.\allowbreak{}FormalGroupSystem}, lines 57--61)\\
\texttt{Formalism.\allowbreak{}MainNotions.\allowbreak{}FormalGroupSystem.\allowbreak{}Sys.\allowbreak{}gv} (\texttt{Formalism.\allowbreak{}SharpRestriction.\allowbreak{}GrowthBounds.\allowbreak{}FormalGroupSystem}, lines 323--323)\\
\texttt{Formalism.\allowbreak{}MainNotions.\allowbreak{}FormalGroupSystem.\allowbreak{}Sys.\allowbreak{}alEven} (\texttt{Formalism.\allowbreak{}SharpRestriction.\allowbreak{}GrowthBounds.\allowbreak{}FormalGroupSystem}, lines 355--355)\par}

{\footnotesize\raggedright\noindent \ref{B3.15}.\quad \texttt{Formalism.\allowbreak{}InvariantFactorPaths.\allowbreak{}G7.\allowbreak{}Top.\allowbreak{}e\_\allowbreak{}one\_\allowbreak{}identity} (\texttt{Formalism.\allowbreak{}ThetaLogCoordinates.\allowbreak{}UniformThetaSubstitution.\allowbreak{}KOneConsequence}, lines 25--25)\\
\texttt{Formalism.\allowbreak{}InvariantFactorPaths.\allowbreak{}G7.\allowbreak{}Top.\allowbreak{}e\_\allowbreak{}one\_\allowbreak{}integral} (\texttt{Formalism.\allowbreak{}ThetaLogCoordinates.\allowbreak{}UniformThetaSubstitution.\allowbreak{}KOneConsequence}, lines 26--26)\par}

{\footnotesize\raggedright\noindent \ref{B3.16}.\quad \texttt{Formalism.\allowbreak{}InvariantFactorPaths.\allowbreak{}G7.\allowbreak{}Top.\allowbreak{}minorHits} (\texttt{Formalism.\allowbreak{}ThetaLogCoordinates.\allowbreak{}UniformThetaSubstitution.\allowbreak{}KOneConsequence}, lines 28--28)\\
\texttt{Formalism.\allowbreak{}InvariantFactorPaths.\allowbreak{}G7.\allowbreak{}Top.\allowbreak{}hitsCapPos} (\texttt{Formalism.\allowbreak{}ThetaLogCoordinates.\allowbreak{}UniformThetaSubstitution.\allowbreak{}KOneConsequence}, lines 33--33)\\
\texttt{Formalism.\allowbreak{}InvariantFactorPaths.\allowbreak{}G7.\allowbreak{}Top.\allowbreak{}hitsColumnsPos} (\texttt{Formalism.\allowbreak{}ThetaLogCoordinates.\allowbreak{}UniformThetaSubstitution.\allowbreak{}KOneConsequence}, lines 34--34)\par}

{\footnotesize\raggedright\noindent \ref{B3.17}.\quad \texttt{Formalism.\allowbreak{}InvariantFactorPaths.\allowbreak{}G7.\allowbreak{}Top.\allowbreak{}hhDecomp} (\texttt{Formalism.\allowbreak{}ThetaLogCoordinates.\allowbreak{}UniformThetaSubstitution.\allowbreak{}KOneConsequence}, lines 38--38)\\
\texttt{Formalism.\allowbreak{}InvariantFactorPaths.\allowbreak{}G7.\allowbreak{}decomposition\_\allowbreak{}of\_\allowbreak{}hh} (\texttt{Formalism.\allowbreak{}ThetaLogCoordinates.\allowbreak{}UniformThetaSubstitution.\allowbreak{}FrobeniusLiftNorms.\allowbreak{}EDecomp}, lines 498--546)\par}

{\footnotesize\raggedright\noindent \ref{B3.18}.\quad \texttt{Formalism.\allowbreak{}InvariantFactorPaths.\allowbreak{}G7.\allowbreak{}Top.\allowbreak{}termBound} (\texttt{Formalism.\allowbreak{}ThetaLogCoordinates.\allowbreak{}UniformThetaSubstitution.\allowbreak{}KOneConsequence}, lines 37--37)\\
\texttt{Formalism.\allowbreak{}InvariantFactorPaths.\allowbreak{}G7.\allowbreak{}TermBound\_\allowbreak{}of} (\texttt{Formalism.\allowbreak{}ThetaLogCoordinates.\allowbreak{}UniformThetaSubstitution.\allowbreak{}FrobeniusLiftNorms.\allowbreak{}NormNatCastLeOfDvd}, lines 495--543)\par}

{\footnotesize\raggedright\noindent \ref{B3.19}.\quad \texttt{Formalism.\allowbreak{}InvariantFactorPaths.\allowbreak{}G7.\allowbreak{}factorialValuationLowerBound} (\texttt{Formalism.\allowbreak{}ThetaLogCoordinates.\allowbreak{}UniformThetaSubstitution.\allowbreak{}FrobeniusLiftNorms.\allowbreak{}ABArith}, lines 35--53)\par}

{\footnotesize\raggedright\noindent \ref{B3.20}.\quad \texttt{Formalism.\allowbreak{}InvariantFactorPaths.\allowbreak{}G7.\allowbreak{}gainFactorialInequality} (\texttt{Formalism.\allowbreak{}ThetaLogCoordinates.\allowbreak{}UniformThetaSubstitution.\allowbreak{}FrobeniusLiftNorms.\allowbreak{}ABArith}, lines 103--183)\par}

{\footnotesize\raggedright\noindent \ref{B3.21}.\quad \texttt{Formalism.\allowbreak{}InvariantFactorPaths.\allowbreak{}G7.\allowbreak{}EntryHyp} (\texttt{Formalism.\allowbreak{}ThetaLogCoordinates.\allowbreak{}UniformThetaSubstitution.\allowbreak{}StatementE}, lines 170--174)\\
\texttt{Formalism.\allowbreak{}InvariantFactorPaths.\allowbreak{}G7.\allowbreak{}Top.\allowbreak{}entryStar} (\texttt{Formalism.\allowbreak{}ThetaLogCoordinates.\allowbreak{}UniformThetaSubstitution.\allowbreak{}KOneConsequence}, lines 39--41)\par}

{\footnotesize\raggedright\noindent \ref{B3.22}.\quad \texttt{Formalism.\allowbreak{}InvariantFactorPaths.\allowbreak{}G7.\allowbreak{}cyclicShiftLeast} (\texttt{Formalism.\allowbreak{}ThetaLogCoordinates.\allowbreak{}UniformThetaSubstitution.\allowbreak{}FrobeniusLiftNorms.\allowbreak{}ABArith}, lines 185--191)\\
\texttt{Formalism.\allowbreak{}InvariantFactorPaths.\allowbreak{}G7.\allowbreak{}excessGainBound} (\texttt{Formalism.\allowbreak{}ThetaLogCoordinates.\allowbreak{}UniformThetaSubstitution.\allowbreak{}FrobeniusLiftNorms.\allowbreak{}ABArith}, lines 193--238)\par}

{\footnotesize\raggedright\noindent \ref{B3.23}.\quad \texttt{Formalism.\allowbreak{}InvariantFactorPaths.\allowbreak{}G7.\allowbreak{}Top.\allowbreak{}entry\_\allowbreak{}a\_\allowbreak{}one} (\texttt{Formalism.\allowbreak{}ThetaLogCoordinates.\allowbreak{}UniformThetaSubstitution.\allowbreak{}KOneConsequence}, lines 43--43)\\
\texttt{Formalism.\allowbreak{}InvariantFactorPaths.\allowbreak{}G7.\allowbreak{}entryExcessBound\_\allowbreak{}of} (\texttt{Formalism.\allowbreak{}ThetaLogCoordinates.\allowbreak{}UniformThetaSubstitution.\allowbreak{}FrobeniusLiftNorms.\allowbreak{}EntryAOneParity}, lines 177--312)\par}

{\footnotesize\raggedright\noindent \ref{B3.24}.\quad \texttt{Formalism.\allowbreak{}InvariantFactorPaths.\allowbreak{}G7.\allowbreak{}Top.\allowbreak{}entry\_\allowbreak{}a\_\allowbreak{}two} (\texttt{Formalism.\allowbreak{}ThetaLogCoordinates.\allowbreak{}UniformThetaSubstitution.\allowbreak{}KOneConsequence}, lines 44--44)\\
\texttt{Formalism.\allowbreak{}InvariantFactorPaths.\allowbreak{}G7.\allowbreak{}entryCappedBound\_\allowbreak{}of} (\texttt{Formalism.\allowbreak{}ThetaLogCoordinates.\allowbreak{}UniformThetaSubstitution.\allowbreak{}FrobeniusLiftNorms.\allowbreak{}EntryATwoAndB}, lines 220--266)\par}

{\footnotesize\raggedright\noindent \ref{B3.25}.\quad \texttt{Formalism.\allowbreak{}InvariantFactorPaths.\allowbreak{}G7.\allowbreak{}Top.\allowbreak{}bCore} (\texttt{Formalism.\allowbreak{}ThetaLogCoordinates.\allowbreak{}UniformThetaSubstitution.\allowbreak{}KOneConsequence}, lines 32--32)\\
\texttt{Formalism.\allowbreak{}InvariantFactorPaths.\allowbreak{}G7.\allowbreak{}Top.\allowbreak{}entryB} (\texttt{Formalism.\allowbreak{}ThetaLogCoordinates.\allowbreak{}UniformThetaSubstitution.\allowbreak{}KOneConsequence}, lines 45--45)\\
\texttt{Formalism.\allowbreak{}InvariantFactorPaths.\allowbreak{}G7.\allowbreak{}EntryB\_\allowbreak{}of} (\texttt{Formalism.\allowbreak{}ThetaLogCoordinates.\allowbreak{}UniformThetaSubstitution.\allowbreak{}FrobeniusLiftNorms.\allowbreak{}EntryATwoAndB}, lines 268--345)\par}

{\footnotesize\raggedright\noindent \ref{B3.26}.\quad \texttt{Formalism.\allowbreak{}InvariantFactorPaths.\allowbreak{}G7.\allowbreak{}kernel\_\allowbreak{}a\_\allowbreak{}one\_\allowbreak{}generic} (\texttt{Formalism.\allowbreak{}ThetaLogCoordinates.\allowbreak{}UniformThetaSubstitution.\allowbreak{}FrobeniusLiftNorms.\allowbreak{}EGaussK}, lines 213--215)\\
\texttt{Formalism.\allowbreak{}InvariantFactorPaths.\allowbreak{}G7.\allowbreak{}kernel\_\allowbreak{}a\_\allowbreak{}two\_\allowbreak{}generic} (\texttt{Formalism.\allowbreak{}ThetaLogCoordinates.\allowbreak{}UniformThetaSubstitution.\allowbreak{}FrobeniusLiftNorms.\allowbreak{}EGaussK}, lines 216--218)\\
\texttt{Formalism.\allowbreak{}InvariantFactorPaths.\allowbreak{}G7.\allowbreak{}kernel\_\allowbreak{}b\_\allowbreak{}generic} (\texttt{Formalism.\allowbreak{}ThetaLogCoordinates.\allowbreak{}UniformThetaSubstitution.\allowbreak{}FrobeniusLiftNorms.\allowbreak{}EGaussK}, lines 219--220)\par}

\section*{Chapter~\ref{ch:b3-one-lift}: One lift on a lattice and the Hecke commutation}

{\footnotesize\raggedright\noindent \ref{B3.27}.\quad \texttt{Formalism.\allowbreak{}MainNotions.\allowbreak{}FormalGroupSystem.\allowbreak{}Pf.\allowbreak{}first\_\allowbreak{}theorem\_\allowbreak{}place\_\allowbreak{}bound} (\texttt{Formalism.\allowbreak{}SharpRestriction.\allowbreak{}GrowthBounds.\allowbreak{}MultisetValuation.\allowbreak{}FirstTheoremPlaceBound}, lines 20--23)\\
\texttt{Formalism.\allowbreak{}MainNotions.\allowbreak{}FormalGroupSystem.\allowbreak{}Pf.\allowbreak{}eb\_\allowbreak{}jet} (\texttt{Formalism.\allowbreak{}SharpRestriction.\allowbreak{}GrowthBounds.\allowbreak{}MultisetValuation.\allowbreak{}EbBrJet}, lines 35--39)\\
\texttt{Formalism.\allowbreak{}MainNotions.\allowbreak{}FormalGroupSystem.\allowbreak{}Pf.\allowbreak{}ea\_\allowbreak{}brNum} (\texttt{Formalism.\allowbreak{}SharpRestriction.\allowbreak{}GrowthBounds.\allowbreak{}MultisetValuation.\allowbreak{}GainOfCapacityLoss}, lines 252--256)\par}

{\footnotesize\raggedright\noindent \ref{B3.28}.\quad \texttt{Formalism.\allowbreak{}MainNotions.\allowbreak{}FormalGroupSystem.\allowbreak{}Pf.\allowbreak{}second\_\allowbreak{}lemma\_\allowbreak{}place\_\allowbreak{}bound} (\texttt{Formalism.\allowbreak{}SharpRestriction.\allowbreak{}GrowthBounds.\allowbreak{}MultisetValuation.\allowbreak{}FirstTheoremPlaceBound}, lines 24--27)\\
\texttt{Formalism.\allowbreak{}MainNotions.\allowbreak{}FormalGroupSystem.\allowbreak{}Pf.\allowbreak{}hecke\_\allowbreak{}place\_\allowbreak{}bound} (\texttt{Formalism.\allowbreak{}SharpRestriction.\allowbreak{}GrowthBounds.\allowbreak{}MultisetValuation.\allowbreak{}FirstTheoremPlaceBound}, lines 28--31)\par}

{\footnotesize\raggedright\noindent \ref{B3.29}.\quad \texttt{Formalism.\allowbreak{}MainNotions.\allowbreak{}FormalGroupSystem.\allowbreak{}Pf.\allowbreak{}kernel\_\allowbreak{}product\_\allowbreak{}form} (\texttt{Formalism.\allowbreak{}SharpRestriction.\allowbreak{}GrowthBounds.\allowbreak{}MultisetValuation.\allowbreak{}RestEntOf}, lines 20--59)\\
\texttt{Formalism.\allowbreak{}MainNotions.\allowbreak{}FormalGroupSystem.\allowbreak{}Pf.\allowbreak{}restEnt\_\allowbreak{}of} (\texttt{Formalism.\allowbreak{}SharpRestriction.\allowbreak{}GrowthBounds.\allowbreak{}MultisetValuation.\allowbreak{}RestEntOf}, lines 62--64)\\
\texttt{Formalism.\allowbreak{}MainNotions.\allowbreak{}FormalGroupSystem.\allowbreak{}Pf.\allowbreak{}en\_\allowbreak{}places\_\allowbreak{}add} (\texttt{Formalism.\allowbreak{}SharpRestriction.\allowbreak{}GrowthBounds.\allowbreak{}MultisetValuation.\allowbreak{}KernelProductForm}, lines 51--110)\par}

{\footnotesize\raggedright\noindent \ref{B3.30}.\quad \texttt{Formalism.\allowbreak{}MainNotions.\allowbreak{}FormalGroupSystem.\allowbreak{}ExpansionEquality} (\texttt{Formalism.\allowbreak{}SharpRestriction.\allowbreak{}GrowthBounds.\allowbreak{}FormalGroupSystem}, lines 566--571)\\
\texttt{Formalism.\allowbreak{}MainNotions.\allowbreak{}FormalGroupSystem.\allowbreak{}ExpansionDecomposition} (\texttt{Formalism.\allowbreak{}SharpRestriction.\allowbreak{}GrowthBounds.\allowbreak{}FormalGroupSystem}, lines 572--586)\\
\texttt{Formalism.\allowbreak{}MainNotions.\allowbreak{}FormalGroupSystem.\allowbreak{}Pf.\allowbreak{}formB\_\allowbreak{}of} (\texttt{Formalism.\allowbreak{}SharpRestriction.\allowbreak{}GrowthBounds.\allowbreak{}MultisetValuation.\allowbreak{}KernelSumOfClassData}, lines 255--278)\par}

{\footnotesize\raggedright\noindent \ref{B3.31}.\quad \texttt{Formalism.\allowbreak{}MainNotions.\allowbreak{}FormalGroupSystem.\allowbreak{}Pf.\allowbreak{}ent\_\allowbreak{}strict} (\texttt{Formalism.\allowbreak{}SharpRestriction.\allowbreak{}GrowthBounds.\allowbreak{}MultisetValuation.\allowbreak{}RamifiedLatticeOfChainEstimates}, lines 46--64)\\
\texttt{Formalism.\allowbreak{}MainNotions.\allowbreak{}FormalGroupSystem.\allowbreak{}Pf.\allowbreak{}formB\_\allowbreak{}lim} (\texttt{Formalism.\allowbreak{}SharpRestriction.\allowbreak{}GrowthBounds.\allowbreak{}MultisetValuation.\allowbreak{}RamifiedLatticeOfChainEstimates}, lines 65--89)\\
\texttt{Formalism.\allowbreak{}MainNotions.\allowbreak{}FormalGroupSystem.\allowbreak{}Pf.\allowbreak{}form\_\allowbreak{}of\_\allowbreak{}place} (\texttt{Formalism.\allowbreak{}SharpRestriction.\allowbreak{}GrowthBounds.\allowbreak{}MultisetValuation.\allowbreak{}RamifiedLatticeOfChainEstimates}, lines 90--95)\par}

{\footnotesize\raggedright\noindent \ref{B3.32}.\quad \texttt{Formalism.\allowbreak{}MainNotions.\allowbreak{}FormalGroupSystem.\allowbreak{}Pf.\allowbreak{}theorem\_\allowbreak{}A\_\allowbreak{}entries\_\allowbreak{}of\_\allowbreak{}corollary} (\texttt{Formalism.\allowbreak{}SharpRestriction.\allowbreak{}GrowthBounds.\allowbreak{}MultisetValuation.\allowbreak{}RamifiedLatticeOfChainEstimates}, lines 104--115)\\
\texttt{Formalism.\allowbreak{}MainNotions.\allowbreak{}FormalGroupSystem.\allowbreak{}Pf.\allowbreak{}slopeEstimate\_\allowbreak{}of\_\allowbreak{}lemmaB\_\allowbreak{}slopeLine} (\texttt{Formalism.\allowbreak{}SharpRestriction.\allowbreak{}GrowthBounds.\allowbreak{}MultisetValuation.\allowbreak{}RamifiedLatticeOfChainEstimates}, lines 116--121)\\
\texttt{Formalism.\allowbreak{}MainNotions.\allowbreak{}FormalGroupSystem.\allowbreak{}Pf.\allowbreak{}form\_\allowbreak{}valuation\_\allowbreak{}of\_\allowbreak{}bound} (\texttt{Formalism.\allowbreak{}SharpRestriction.\allowbreak{}GrowthBounds.\allowbreak{}MultisetValuation.\allowbreak{}RamifiedLatticeOfChainEstimates}, lines 96--102)\par}

{\footnotesize\raggedright\noindent \ref{B3.33}.\quad \texttt{Formalism.\allowbreak{}MainNotions.\allowbreak{}FormalGroupSystem.\allowbreak{}Pf.\allowbreak{}hecke\_\allowbreak{}integral\_\allowbreak{}of\_\allowbreak{}forms} (\texttt{Formalism.\allowbreak{}SharpRestriction.\allowbreak{}GrowthBounds.\allowbreak{}MultisetValuation.\allowbreak{}RamifiedLatticeOfChainEstimates}, lines 122--141)\par}

{\footnotesize\raggedright\noindent \ref{B3.34}.\quad \texttt{Formalism.\allowbreak{}MainNotions.\allowbreak{}FormalGroupSystem.\allowbreak{}Pf.\allowbreak{}latticeChain\_\allowbreak{}succ\_\allowbreak{}of\_\allowbreak{}partial} (\texttt{Formalism.\allowbreak{}SharpRestriction.\allowbreak{}GrowthBounds.\allowbreak{}MultisetValuation.\allowbreak{}RamifiedLatticeOfChainEstimates}, lines 165--168)\\
\texttt{Formalism.\allowbreak{}MainNotions.\allowbreak{}FormalGroupSystem.\allowbreak{}Pf.\allowbreak{}restricted\_\allowbreak{}hecke\_\allowbreak{}integral} (\texttt{Formalism.\allowbreak{}SharpRestriction.\allowbreak{}GrowthBounds.\allowbreak{}MultisetValuation.\allowbreak{}RestrictedHeckeIntegralOfEdges}, lines 289--290)\par}

{\footnotesize\raggedright\noindent \ref{B3.35}.\quad \texttt{Formalism.\allowbreak{}MainNotions.\allowbreak{}G.\allowbreak{}Level.\allowbreak{}yv} (\texttt{Formalism.\allowbreak{}SharpRestriction.\allowbreak{}GrowthBounds.\allowbreak{}DataG}, lines 148--152)\\
\texttt{Formalism.\allowbreak{}MainNotions.\allowbreak{}FormalGroupSystem.\allowbreak{}Sys.\allowbreak{}Tt} (\texttt{Formalism.\allowbreak{}SharpRestriction.\allowbreak{}GrowthBounds.\allowbreak{}FormalGroupSystem}, lines 331--333)\\
\texttt{Formalism.\allowbreak{}MainNotions.\allowbreak{}FormalGroupSystem.\allowbreak{}Sys.\allowbreak{}cl} (\texttt{Formalism.\allowbreak{}SharpRestriction.\allowbreak{}GrowthBounds.\allowbreak{}FormalGroupSystem}, lines 344--348)\\
\texttt{Formalism.\allowbreak{}MainNotions.\allowbreak{}FormalGroupSystem.\allowbreak{}Sys.\allowbreak{}Ttu} (\texttt{Formalism.\allowbreak{}SharpRestriction.\allowbreak{}GrowthBounds.\allowbreak{}FormalGroupSystem}, lines 349--351)\\
\texttt{Formalism.\allowbreak{}MainNotions.\allowbreak{}FormalGroupSystem.\allowbreak{}Sys.\allowbreak{}Zop} (\texttt{Formalism.\allowbreak{}SharpRestriction.\allowbreak{}GrowthBounds.\allowbreak{}FormalGroupSystem}, lines 352--354)\par}

{\footnotesize\raggedright\noindent \ref{B3.35a}.\quad \texttt{Formalism.\allowbreak{}MainNotions.\allowbreak{}FormalGroupSystem.\allowbreak{}Pend} (\texttt{Formalism.\allowbreak{}SharpRestriction.\allowbreak{}GrowthBounds.\allowbreak{}FormalGroupSystem}, lines 447--494)\\
\texttt{Formalism.\allowbreak{}MainNotions.\allowbreak{}FormalGroupSystem.\allowbreak{}adm5} (\texttt{Formalism.\allowbreak{}SharpRestriction.\allowbreak{}GrowthBounds.\allowbreak{}FormalGroupSystem}, lines 590--596)\\
\texttt{GlobalTheory.\allowbreak{}XiGen.\allowbreak{}xiEx\_\allowbreak{}theta\_\allowbreak{}general} (\texttt{Development.\allowbreak{}AbelianVarietyTheory.\allowbreak{}Proofs.\allowbreak{}SatakeParameters.\allowbreak{}XiExThetaGeneral}, lines 86--304)\par}

{\footnotesize\raggedright\noindent \ref{B3.35d}.\quad \texttt{GlobalTheory.\allowbreak{}Mem.\allowbreak{}NatEqui.\allowbreak{}eqP} (\texttt{Development.\allowbreak{}AbelianVarietyTheory.\allowbreak{}Proofs.\allowbreak{}IwasawaDecomposition.\allowbreak{}NormalisedFamilySatisfiesEquivarianceLaws}, lines 59--117)\\
\texttt{GlobalTheory.\allowbreak{}TrueTheta.\allowbreak{}exists\_\allowbreak{}family\_\allowbreak{}of\_\allowbreak{}curve} (\texttt{Development.\allowbreak{}AbelianVarietyTheory.\allowbreak{}Proofs.\allowbreak{}IwasawaDecomposition.\allowbreak{}ExistsThetaVectorFamilyOfCurveDatum}, lines 15--108)\par}

{\footnotesize\raggedright\noindent \ref{B3.35f}.\quad \texttt{Formalism.\allowbreak{}L2.\allowbreak{}Col.\allowbreak{}Loc.\allowbreak{}E1} (\texttt{Formalism.\allowbreak{}LocalToGlobal.\allowbreak{}Columns.\allowbreak{}FiniteSumNonzeroLargeDepth}, lines 577--580)\\
\texttt{Formalism.\allowbreak{}L2.\allowbreak{}Col.\allowbreak{}XiAt} (\texttt{Formalism.\allowbreak{}LocalToGlobal.\allowbreak{}Columns.\allowbreak{}FiniteSumNonzeroLargeDepth}, lines 946--948)\\
\texttt{Formalism.\allowbreak{}L2.\allowbreak{}Col.\allowbreak{}XiEntry} (\texttt{Formalism.\allowbreak{}LocalToGlobal.\allowbreak{}Columns.\allowbreak{}FiniteSumNonzeroLargeDepth}, lines 1116--1119)\\
\texttt{Formalism.\allowbreak{}L2.\allowbreak{}Col.\allowbreak{}XiCpl1} (\texttt{Formalism.\allowbreak{}LocalToGlobal.\allowbreak{}Columns.\allowbreak{}FiniteSumNonzeroLargeDepth}, lines 1134--1138)\\
\texttt{Formalism.\allowbreak{}L2.\allowbreak{}Col.\allowbreak{}XiTopG} (\texttt{Formalism.\allowbreak{}LocalToGlobal.\allowbreak{}Columns.\allowbreak{}ChainInstY}, lines 50--56)\\
\texttt{GlobalTheory.\allowbreak{}TwV.\allowbreak{}xiAd} (\texttt{Development.\allowbreak{}AbelianVarietyTheory.\allowbreak{}Definitions.\allowbreak{}GaloisGeometry.\allowbreak{}UnderTie}, lines 48--53)\par}

{\footnotesize\raggedright\noindent \ref{B3.35e}.\quad \texttt{GlobalTheory.\allowbreak{}Mem.\allowbreak{}ChiDVal.\allowbreak{}chiD\_\allowbreak{}mem\_\allowbreak{}of} (\texttt{Development.\allowbreak{}AbelianVarietyTheory.\allowbreak{}Proofs.\allowbreak{}IwasawaDecomposition.\allowbreak{}ChiDValuesAndCcTop}, lines 101--112)\\
\texttt{GlobalTheory.\allowbreak{}Port.\allowbreak{}chRat\_\allowbreak{}ofChiD} (\texttt{Development.\allowbreak{}AbelianVarietyTheory.\allowbreak{}Definitions.\allowbreak{}GaloisGeometry.\allowbreak{}ExistsUnifPow}, lines 25--33)\\
\texttt{GlobalTheory.\allowbreak{}TwV.\allowbreak{}Xc.\allowbreak{}chiI\_\allowbreak{}pow\_\allowbreak{}rat} (\texttt{Development.\allowbreak{}AbelianVarietyTheory.\allowbreak{}Proofs.\allowbreak{}SatakeParameters.\allowbreak{}KernelSiteNormalisingCharPowValuesInCoeffField}, lines 107--128)\\
\texttt{GlobalTheory.\allowbreak{}TwV.\allowbreak{}adLawAt\_\allowbreak{}xiAd} (\texttt{Development.\allowbreak{}AbelianVarietyTheory.\allowbreak{}Proofs.\allowbreak{}SatakeParameters.\allowbreak{}AdelicChainDatumCharacterLawsJointKernel}, lines 62--116)\par}

{\footnotesize\raggedright\noindent \ref{B3.35g}.\quad \texttt{GlobalTheory.\allowbreak{}KerCommon.\allowbreak{}kerTh\_\allowbreak{}common\_\allowbreak{}op} (\texttt{Development.\allowbreak{}AbelianVarietyTheory.\allowbreak{}Proofs.\allowbreak{}SatakeParameters.\allowbreak{}KernelFunctionsCommonLinearOperatorOffSplitPrime}, lines 128--175)\\
\texttt{GlobalTheory.\allowbreak{}kerTh\_\allowbreak{}entire\_\allowbreak{}algCoef} (\texttt{Development.\allowbreak{}AbelianVarietyTheory.\allowbreak{}Proofs.\allowbreak{}SatakeParameters.\allowbreak{}KerThEntireAlgCoef}, lines 22--39)\par}

{\footnotesize\raggedright\noindent \ref{B3.35a-arithmetic}.\quad \texttt{GlobalTheory.\allowbreak{}TwV.\allowbreak{}site\_\allowbreak{}law\_\allowbreak{}xiIn} (\texttt{Development.\allowbreak{}AbelianVarietyTheory.\allowbreak{}Proofs.\allowbreak{}SatakeParameters.\allowbreak{}KernelSiteLawAndFixedCoordinateStabiliserLaw}, lines 20--49)\par}

{\footnotesize\raggedright\noindent \ref{B3.35b}.\quad \texttt{Formalism.\allowbreak{}MainNotions.\allowbreak{}FormalGroupSystem.\allowbreak{}Pf.\allowbreak{}kernel\_\allowbreak{}eq\_\allowbreak{}theta} (\texttt{Formalism.\allowbreak{}SharpRestriction.\allowbreak{}GrowthBounds.\allowbreak{}MultisetValuation.\allowbreak{}CoupledSumOfEntries}, lines 23--29)\\
\texttt{Formalism.\allowbreak{}MainNotions.\allowbreak{}FormalGroupSystem.\allowbreak{}Pf.\allowbreak{}integral\_\allowbreak{}span\_\allowbreak{}coordinates} (\texttt{Formalism.\allowbreak{}SharpRestriction.\allowbreak{}GrowthBounds.\allowbreak{}MultisetValuation.\allowbreak{}IntegralSpanCoordinates}, lines 62--67)\\
\texttt{Formalism.\allowbreak{}MainNotions.\allowbreak{}FormalGroupSystem.\allowbreak{}Pf.\allowbreak{}kernel\_\allowbreak{}nsmul} (\texttt{Formalism.\allowbreak{}SharpRestriction.\allowbreak{}GrowthBounds.\allowbreak{}MultisetValuation.\allowbreak{}IntegralSpanCoordinates}, lines 68--73)\par}

{\footnotesize\raggedright\noindent \ref{B3.35c}.\quad \texttt{Formalism.\allowbreak{}L2.\allowbreak{}G.\allowbreak{}Pend.\allowbreak{}up\_\allowbreak{}down\_\allowbreak{}adjoint} (\texttt{Formalism.\allowbreak{}LocalToGlobal.\allowbreak{}GroupLattices.\allowbreak{}PendingLemmas.\allowbreak{}OrbitSumStabilizer}, lines 257--299)\\
\texttt{Formalism.\allowbreak{}MainNotions.\allowbreak{}FormalGroupSystem.\allowbreak{}MinkowskiCounting} (\texttt{Formalism.\allowbreak{}SharpRestriction.\allowbreak{}GrowthBounds.\allowbreak{}FormalGroupSystem}, lines 597--602)\\
\texttt{Formalism.\allowbreak{}MainNotions.\allowbreak{}FormalGroupSystem.\allowbreak{}ExpansionDecomposition} (\texttt{Formalism.\allowbreak{}SharpRestriction.\allowbreak{}GrowthBounds.\allowbreak{}FormalGroupSystem}, lines 572--586)\\
\texttt{Formalism.\allowbreak{}MainNotions.\allowbreak{}FormalGroupSystem.\allowbreak{}Pf.\allowbreak{}hk\_\allowbreak{}adj} (\texttt{Formalism.\allowbreak{}SharpRestriction.\allowbreak{}GrowthBounds.\allowbreak{}MultisetValuation.\allowbreak{}CoupledSumOfEntries}, lines 65--110)\\
\texttt{Formalism.\allowbreak{}MainNotions.\allowbreak{}FormalGroupSystem.\allowbreak{}Pf.\allowbreak{}ksa\_\allowbreak{}of\_\allowbreak{}adm} (\texttt{Formalism.\allowbreak{}SharpRestriction.\allowbreak{}GrowthBounds.\allowbreak{}MultisetValuation.\allowbreak{}CoupledSumOfEntries}, lines 140--151)\\
\texttt{Formalism.\allowbreak{}L2.\allowbreak{}KerInert.\allowbreak{}SecondKernelCoupled.\allowbreak{}hecke\_\allowbreak{}coupled\_\allowbreak{}entry} (\texttt{Formalism.\allowbreak{}LocalToGlobal.\allowbreak{}InertKernelTheory.\allowbreak{}InertCoprimeLemmas.\allowbreak{}SecondKernelCoupled}, lines 104--137)\par}

{\footnotesize\raggedright\noindent \ref{B3.36}.\quad \texttt{Formalism.\allowbreak{}MainNotions.\allowbreak{}FormalGroupSystem.\allowbreak{}Pf.\allowbreak{}countLoc} (\texttt{Formalism.\allowbreak{}SharpRestriction.\allowbreak{}GrowthBounds.\allowbreak{}MultisetValuation.\allowbreak{}SmulPsiGUp}, lines 276--289)\\
\texttt{Formalism.\allowbreak{}MainNotions.\allowbreak{}FormalGroupSystem.\allowbreak{}Pf.\allowbreak{}heckeCondition\_\allowbreak{}of\_\allowbreak{}adm} (\texttt{Formalism.\allowbreak{}SharpRestriction.\allowbreak{}GrowthBounds.\allowbreak{}MultisetValuation.\allowbreak{}HeckeConditionOfAdm}, lines 75--147)\par}

{\footnotesize\raggedright\noindent \ref{B3.37}.\quad \texttt{Formalism.\allowbreak{}MainNotions.\allowbreak{}FormalGroupSystem.\allowbreak{}PrZ} (\texttt{Formalism.\allowbreak{}SharpRestriction.\allowbreak{}GrowthBounds.\allowbreak{}FormalGroupSystem}, lines 76--80)\\
\texttt{Formalism.\allowbreak{}MainNotions.\allowbreak{}FormalGroupSystem.\allowbreak{}Pr} (\texttt{Formalism.\allowbreak{}SharpRestriction.\allowbreak{}GrowthBounds.\allowbreak{}FormalGroupSystem}, lines 81--85)\\
\texttt{Formalism.\allowbreak{}MainNotions.\allowbreak{}FormalGroupSystem.\allowbreak{}Pf.\allowbreak{}PrZ\_\allowbreak{}cast} (\texttt{Formalism.\allowbreak{}SharpRestriction.\allowbreak{}GrowthBounds.\allowbreak{}MultisetValuation.\allowbreak{}LiftBoundOfIntegrality}, lines 91--101)\par}

{\footnotesize\raggedright\noindent \ref{B3.38}.\quad \texttt{Formalism.\allowbreak{}MainNotions.\allowbreak{}FormalGroupSystem.\allowbreak{}Conclusion4} (\texttt{Formalism.\allowbreak{}SharpRestriction.\allowbreak{}GrowthBounds.\allowbreak{}FormalGroupSystem}, lines 671--680)\\
\texttt{Formalism.\allowbreak{}MainNotions.\allowbreak{}FormalGroupSystem.\allowbreak{}Pf.\allowbreak{}Ci} (\texttt{Formalism.\allowbreak{}SharpRestriction.\allowbreak{}GrowthBounds.\allowbreak{}MultisetValuation.\allowbreak{}HeckeValueEstimateOfRecordValue}, lines 195--270)\\
\texttt{Formalism.\allowbreak{}MainNotions.\allowbreak{}FormalGroupSystem.\allowbreak{}Pf.\allowbreak{}Cii} (\texttt{Formalism.\allowbreak{}SharpRestriction.\allowbreak{}GrowthBounds.\allowbreak{}MultisetValuation.\allowbreak{}TheoremOfAlternateForm}, lines 538--549)\\
\texttt{Formalism.\allowbreak{}MainNotions.\allowbreak{}FormalGroupSystem.\allowbreak{}Pf.\allowbreak{}Ciii} (\texttt{Formalism.\allowbreak{}SharpRestriction.\allowbreak{}GrowthBounds.\allowbreak{}MultisetValuation.\allowbreak{}TheoremOfAlternateForm}, lines 608--617)\par}

\section*{Chapter~\ref{ch:chain-induction}: Induction along the chain}

{\footnotesize\raggedright\noindent \ref{B3.39}.\quad \texttt{Formalism.\allowbreak{}MainNotions.\allowbreak{}FormalGroupSystem.\allowbreak{}Pf.\allowbreak{}ZL\_\allowbreak{}chain} (\texttt{Formalism.\allowbreak{}SharpRestriction.\allowbreak{}GrowthBounds.\allowbreak{}MultisetValuation.\allowbreak{}YEPEComm}, lines 110--123)\\
\texttt{Formalism.\allowbreak{}MainNotions.\allowbreak{}FormalGroupSystem.\allowbreak{}Pf.\allowbreak{}ZL\_\allowbreak{}bound} (\texttt{Formalism.\allowbreak{}SharpRestriction.\allowbreak{}GrowthBounds.\allowbreak{}MultisetValuation.\allowbreak{}AofUpdate}, lines 172--217)\par}

{\footnotesize\raggedright\noindent \ref{B3.40}.\quad \texttt{Formalism.\allowbreak{}MainNotions.\allowbreak{}FormalGroupSystem.\allowbreak{}Sys.\allowbreak{}Iprod} (\texttt{Formalism.\allowbreak{}SharpRestriction.\allowbreak{}GrowthBounds.\allowbreak{}FormalGroupSystem}, lines 361--364)\par}

{\footnotesize\raggedright\noindent \ref{B3.41}.\quad \texttt{Formalism.\allowbreak{}MainNotions.\allowbreak{}FormalGroupSystem.\allowbreak{}Pf.\allowbreak{}bbJ} (\texttt{Formalism.\allowbreak{}SharpRestriction.\allowbreak{}GrowthBounds.\allowbreak{}MultisetValuation.\allowbreak{}YEPEComm}, lines 127--129)\\
\texttt{Formalism.\allowbreak{}MainNotions.\allowbreak{}FormalGroupSystem.\allowbreak{}Pf.\allowbreak{}main\_\allowbreak{}ind} (\texttt{Formalism.\allowbreak{}SharpRestriction.\allowbreak{}GrowthBounds.\allowbreak{}MultisetValuation.\allowbreak{}YEPEComm}, lines 131--185)\par}

{\footnotesize\raggedright\noindent \ref{B3.42}.\quad \texttt{Formalism.\allowbreak{}MainNotions.\allowbreak{}FormalGroupSystem.\allowbreak{}LowTermBound} (\texttt{Formalism.\allowbreak{}SharpRestriction.\allowbreak{}GrowthBounds.\allowbreak{}FormalGroupSystem}, lines 708--716)\\
\texttt{Formalism.\allowbreak{}MainNotions.\allowbreak{}FormalGroupSystem.\allowbreak{}Pf.\allowbreak{}lift\_\allowbreak{}bound\_\allowbreak{}of\_\allowbreak{}integrality} (\texttt{Formalism.\allowbreak{}SharpRestriction.\allowbreak{}GrowthBounds.\allowbreak{}MultisetValuation.\allowbreak{}LiftBoundOfIntegrality}, lines 103--174)\par}

{\footnotesize\raggedright\noindent \ref{B3.43}.\quad \texttt{Formalism.\allowbreak{}MainNotions.\allowbreak{}FormalGroupSystem.\allowbreak{}Sv} (\texttt{Formalism.\allowbreak{}SharpRestriction.\allowbreak{}GrowthBounds.\allowbreak{}FormalGroupSystem}, lines 86--87)\\
\texttt{Formalism.\allowbreak{}MainNotions.\allowbreak{}FormalGroupSystem.\allowbreak{}RelativePresentationLiterature} (\texttt{Formalism.\allowbreak{}SharpRestriction.\allowbreak{}GrowthBounds.\allowbreak{}FormalGroupSystem}, lines 703--706)\\
\texttt{Formalism.\allowbreak{}MainNotions.\allowbreak{}FormalGroupSystem.\allowbreak{}Pf.\allowbreak{}relative\_\allowbreak{}presentation\_\allowbreak{}of\_\allowbreak{}d7} (\texttt{Formalism.\allowbreak{}SharpRestriction.\allowbreak{}GrowthBounds.\allowbreak{}MultisetValuation.\allowbreak{}RelativePresentationOfD7}, lines 110--126)\par}

{\footnotesize\raggedright\noindent \ref{B3.44}.\quad \texttt{Formalism.\allowbreak{}MainNotions.\allowbreak{}FormalGroupSystem.\allowbreak{}Conclusion5} (\texttt{Formalism.\allowbreak{}SharpRestriction.\allowbreak{}GrowthBounds.\allowbreak{}FormalGroupSystem}, lines 697--701)\\
\texttt{Formalism.\allowbreak{}MainNotions.\allowbreak{}FormalGroupSystem.\allowbreak{}Pf.\allowbreak{}ramifiedLattice\_\allowbreak{}of\_\allowbreak{}chain\_\allowbreak{}estimates} (\texttt{Formalism.\allowbreak{}SharpRestriction.\allowbreak{}GrowthBounds.\allowbreak{}MultisetValuation.\allowbreak{}RamifiedLatticeOfChainEstimates}, lines 169--173)\par}

{\footnotesize\raggedright\noindent \ref{B3.45a}.\quad \texttt{Formalism.\allowbreak{}MainNotions.\allowbreak{}FormalGroupSystem.\allowbreak{}ErrorExponentStatement} (\texttt{Formalism.\allowbreak{}SharpRestriction.\allowbreak{}GrowthBounds.\allowbreak{}FormalGroupSystem}, lines 785--787)\par}

{\footnotesize\raggedright\noindent \ref{B3.45}.\quad \texttt{Formalism.\allowbreak{}MainNotions.\allowbreak{}FormalGroupSystem.\allowbreak{}admitted} (\texttt{Formalism.\allowbreak{}SharpRestriction.\allowbreak{}GrowthBounds.\allowbreak{}FormalGroupSystem}, lines 771--773)\\
\texttt{Formalism.\allowbreak{}MainNotions.\allowbreak{}FormalGroupSystem.\allowbreak{}Conclusion₂} (\texttt{Formalism.\allowbreak{}SharpRestriction.\allowbreak{}GrowthBounds.\allowbreak{}FormalGroupSystem}, lines 775--777)\\
\texttt{Formalism.\allowbreak{}MainNotions.\allowbreak{}FormalGroupSystem.\allowbreak{}Stmt} (\texttt{Formalism.\allowbreak{}SharpRestriction.\allowbreak{}GrowthBounds.\allowbreak{}FormalGroupSystem}, lines 779--780)\\
\texttt{Formalism.\allowbreak{}MainNotions.\allowbreak{}FormalGroupSystem.\allowbreak{}Pf.\allowbreak{}stmt\_\allowbreak{}proved} (\texttt{Formalism.\allowbreak{}SharpRestriction.\allowbreak{}GrowthBounds.\allowbreak{}MultisetValuation.\allowbreak{}PrvAndKernelCommutationUpToUnitsOfRests}, lines 107--108)\\
\texttt{Formalism.\allowbreak{}MainNotions.\allowbreak{}FormalGroupSystem.\allowbreak{}Pf.\allowbreak{}stmtAt\_\allowbreak{}proved} (\texttt{Formalism.\allowbreak{}SharpRestriction.\allowbreak{}GrowthBounds.\allowbreak{}MultisetValuation.\allowbreak{}PrvAndKernelCommutationUpToUnitsOfRests}, lines 110--111)\par}

\section*{Chapter~\ref{ch:global-frobenius-orbit-parameters}: Global Frobenius-orbit parameters}

{\footnotesize\raggedright\noindent \ref{B3.46}.\quad \texttt{Formalism.\allowbreak{}MainNotions.\allowbreak{}G.\allowbreak{}frobSet} (\texttt{Formalism.\allowbreak{}SharpRestriction.\allowbreak{}GrowthBounds.\allowbreak{}IsResidueClassZero}, lines 93--95)\\
\texttt{Formalism.\allowbreak{}MainNotions.\allowbreak{}G.\allowbreak{}hP} (\texttt{Formalism.\allowbreak{}SharpRestriction.\allowbreak{}GrowthBounds.\allowbreak{}IsResidueClassZero}, lines 97--98)\\
\texttt{Formalism.\allowbreak{}MainNotions.\allowbreak{}G.\allowbreak{}partner} (\texttt{Formalism.\allowbreak{}SharpRestriction.\allowbreak{}GrowthBounds.\allowbreak{}IsResidueClassZero}, lines 131--132)\\
\texttt{Formalism.\allowbreak{}MainNotions.\allowbreak{}G.\allowbreak{}stepP} (\texttt{Formalism.\allowbreak{}SharpRestriction.\allowbreak{}GrowthBounds.\allowbreak{}IsResidueClassZero}, lines 161--161)\\
\texttt{Formalism.\allowbreak{}MainNotions.\allowbreak{}G.\allowbreak{}nP} (\texttt{Formalism.\allowbreak{}SharpRestriction.\allowbreak{}GrowthBounds.\allowbreak{}IsResidueClassZero}, lines 163--164)\\
\texttt{Formalism.\allowbreak{}MainNotions.\allowbreak{}G.\allowbreak{}reps} (\texttt{Formalism.\allowbreak{}SharpRestriction.\allowbreak{}GrowthBounds.\allowbreak{}IsResidueClassZero}, lines 166--166)\\
\texttt{Formalism.\allowbreak{}MainNotions.\allowbreak{}G.\allowbreak{}topsOn} (\texttt{Formalism.\allowbreak{}SharpRestriction.\allowbreak{}GrowthBounds.\allowbreak{}IsResidueClassZero}, lines 168--169)\\
\texttt{Formalism.\allowbreak{}MainNotions.\allowbreak{}G.\allowbreak{}IsV} (\texttt{Formalism.\allowbreak{}SharpRestriction.\allowbreak{}GrowthBounds.\allowbreak{}IsResidueClassZero}, lines 179--180)\\
\texttt{Formalism.\allowbreak{}BasicNotions.\allowbreak{}primeClass} (\texttt{Formalism.\allowbreak{}LocalRingCardinality.\allowbreak{}DataHost}, lines 33--37)\\
\texttt{Formalism.\allowbreak{}HC.\allowbreak{}HostRaw.\allowbreak{}psAt} (\texttt{Formalism.\allowbreak{}SharpRestriction.\allowbreak{}RamificationFrobenius.\allowbreak{}TotientGaloisClasses.\allowbreak{}ConjEMemHostFIff}, lines 68--69)\par}

{\footnotesize\raggedright\noindent \ref{B3.47}.\quad \texttt{Formalism.\allowbreak{}MainNotions.\allowbreak{}G.\allowbreak{}Inst} (\texttt{Formalism.\allowbreak{}SharpRestriction.\allowbreak{}GrowthBounds.\allowbreak{}IsResidueClassZero}, lines 141--150)\\
\texttt{Formalism.\allowbreak{}MainNotions.\allowbreak{}G.\allowbreak{}fill} (\texttt{Formalism.\allowbreak{}SharpRestriction.\allowbreak{}GrowthBounds.\allowbreak{}IsResidueClassZero}, lines 154--155)\\
\texttt{Formalism.\allowbreak{}MainNotions.\allowbreak{}G.\allowbreak{}fillG} (\texttt{Formalism.\allowbreak{}SharpRestriction.\allowbreak{}GrowthBounds.\allowbreak{}IsResidueClassZero}, lines 159--159)\\
\texttt{Formalism.\allowbreak{}MainNotions.\allowbreak{}G.\allowbreak{}al} (\texttt{Formalism.\allowbreak{}SharpRestriction.\allowbreak{}GrowthBounds.\allowbreak{}IsResidueClassZero}, lines 230--230)\\
\texttt{Formalism.\allowbreak{}MainNotions.\allowbreak{}G.\allowbreak{}capU} (\texttt{Formalism.\allowbreak{}SharpRestriction.\allowbreak{}GrowthBounds.\allowbreak{}IsResidueClassZero}, lines 232--233)\\
\texttt{Formalism.\allowbreak{}MainNotions.\allowbreak{}G.\allowbreak{}dl} (\texttt{Formalism.\allowbreak{}SharpRestriction.\allowbreak{}GrowthBounds.\allowbreak{}IsResidueClassZero}, lines 241--243)\par}

{\footnotesize\raggedright\noindent \ref{B3.48}.\quad \texttt{Formalism.\allowbreak{}MainNotions.\allowbreak{}G.\allowbreak{}IsDist} (\texttt{Formalism.\allowbreak{}SharpRestriction.\allowbreak{}GrowthBounds.\allowbreak{}IsResidueClassZero}, lines 235--236)\\
\texttt{Formalism.\allowbreak{}MainNotions.\allowbreak{}G.\allowbreak{}Excl} (\texttt{Formalism.\allowbreak{}SharpRestriction.\allowbreak{}GrowthBounds.\allowbreak{}IsResidueClassZero}, lines 238--239)\\
\texttt{Formalism.\allowbreak{}MainNotions.\allowbreak{}G.\allowbreak{}g0} (\texttt{Formalism.\allowbreak{}SharpRestriction.\allowbreak{}GrowthBounds.\allowbreak{}IsResidueClassZero}, lines 247--250)\par}

{\footnotesize\raggedright\noindent \ref{B3.49}.\quad \texttt{Formalism.\allowbreak{}MainNotions.\allowbreak{}G.\allowbreak{}gU} (\texttt{Formalism.\allowbreak{}SharpRestriction.\allowbreak{}GrowthBounds.\allowbreak{}IsResidueClassZero}, lines 254--258)\par}

{\footnotesize\raggedright\noindent \ref{B3.49}.\quad \texttt{Formalism.\allowbreak{}MainNotions.\allowbreak{}FormalGroupSystem.\allowbreak{}Pf.\allowbreak{}nc\_\allowbreak{}summable\_\allowbreak{}min} (\texttt{Formalism.\allowbreak{}SharpRestriction.\allowbreak{}GrowthBounds.\allowbreak{}MultisetValuation.\allowbreak{}NcSummableMin}, lines 40--53)\par}

{\footnotesize\raggedright\noindent \ref{def:p3c5-orbit-model}.\quad \texttt{Formalism.\allowbreak{}BasicNotions.\allowbreak{}GatTies} (\texttt{Formalism.\allowbreak{}LocalRingCardinality.\allowbreak{}DataGatTies}, lines 7--29)\\
\texttt{Formalism.\allowbreak{}MainNotions.\allowbreak{}G.\allowbreak{}LatTied} (\texttt{Formalism.\allowbreak{}SharpRestriction.\allowbreak{}GrowthBounds.\allowbreak{}IsResidueClassZero}, lines 430--451)\par}

{\footnotesize\raggedright\noindent \ref{def:p3c5-size}.\quad \texttt{Formalism.\allowbreak{}MainNotions.\allowbreak{}SizesAt} (\texttt{Formalism.\allowbreak{}SharpRestriction.\allowbreak{}ElevenLe}, lines 134--138)\par}

{\footnotesize\raggedright\noindent \ref{B3.50}.\quad \texttt{Formalism.\allowbreak{}MainNotions.\allowbreak{}G.\allowbreak{}Pr} (\texttt{Formalism.\allowbreak{}SharpRestriction.\allowbreak{}GrowthBounds.\allowbreak{}IsResidueClassZero}, lines 262--267)\\
\texttt{Formalism.\allowbreak{}MainNotions.\allowbreak{}FormalGroupSystem.\allowbreak{}Pf.\allowbreak{}price\_\allowbreak{}bound\_\allowbreak{}large\_\allowbreak{}primes} (\texttt{Formalism.\allowbreak{}SharpRestriction.\allowbreak{}GrowthBounds.\allowbreak{}MultisetValuation.\allowbreak{}PriceBoundLargePrimes}, lines 117--180)\\
\texttt{Formalism.\allowbreak{}MainNotions.\allowbreak{}FormalGroupSystem.\allowbreak{}Pf.\allowbreak{}price\_\allowbreak{}bound\_\allowbreak{}small\_\allowbreak{}primes} (\texttt{Formalism.\allowbreak{}SharpRestriction.\allowbreak{}GrowthBounds.\allowbreak{}MultisetValuation.\allowbreak{}PriceLinearCombinationLe}, lines 89--142)\\
\texttt{Formalism.\allowbreak{}MainNotions.\allowbreak{}SizeRho} (\texttt{Formalism.\allowbreak{}SharpRestriction.\allowbreak{}ElevenLe}, lines 70--72)\\
\texttt{Formalism.\allowbreak{}MainNotions.\allowbreak{}FormalGroupSystem.\allowbreak{}Pf.\allowbreak{}ne\_\allowbreak{}bT\_\allowbreak{}le} (\texttt{Formalism.\allowbreak{}SharpRestriction.\allowbreak{}GrowthBounds.\allowbreak{}MultisetValuation.\allowbreak{}NeSumLe}, lines 87--105)\par}

{\footnotesize\raggedright\noindent \ref{B3.50}.\quad \texttt{Formalism.\allowbreak{}MainNotions.\allowbreak{}G.\allowbreak{}IsGoodOrbit} (\texttt{Formalism.\allowbreak{}SharpRestriction.\allowbreak{}GrowthBounds.\allowbreak{}UnramifiedQuadraticSetting}, lines 250--252)\\
\texttt{Formalism.\allowbreak{}G0.\allowbreak{}GoodOrbitIntervalTileCondition} (\texttt{Formalism.\allowbreak{}SubspaceCounting.\allowbreak{}GoodOrbitIntervalTileCondition}, lines 24--25)\par}

{\footnotesize\raggedright\noindent \ref{B3.51}.\quad \texttt{Formalism.\allowbreak{}MainNotions.\allowbreak{}G.\allowbreak{}kapv} (\texttt{Formalism.\allowbreak{}SharpRestriction.\allowbreak{}GrowthBounds.\allowbreak{}IsResidueClassZero}, lines 276--277)\\
\texttt{Formalism.\allowbreak{}MainNotions.\allowbreak{}G.\allowbreak{}Cand} (\texttt{Formalism.\allowbreak{}SharpRestriction.\allowbreak{}GrowthBounds.\allowbreak{}IsResidueClassZero}, lines 279--281)\\
\texttt{Formalism.\allowbreak{}MainNotions.\allowbreak{}G.\allowbreak{}VV} (\texttt{Formalism.\allowbreak{}SharpRestriction.\allowbreak{}GrowthBounds.\allowbreak{}IsResidueClassZero}, lines 292--293)\\
\texttt{Formalism.\allowbreak{}MainNotions.\allowbreak{}G.\allowbreak{}gainP} (\texttt{Formalism.\allowbreak{}SharpRestriction.\allowbreak{}GrowthBounds.\allowbreak{}IsResidueClassZero}, lines 355--358)\par}

{\footnotesize\raggedright\noindent \ref{B3.52}.\quad \texttt{Formalism.\allowbreak{}MainNotions.\allowbreak{}G.\allowbreak{}ParametersFill} (\texttt{Formalism.\allowbreak{}SharpRestriction.\allowbreak{}GrowthBounds.\allowbreak{}IsResidueClassZero}, lines 401--410)\\
\texttt{Formalism.\allowbreak{}MainNotions.\allowbreak{}G.\allowbreak{}TTP} (\texttt{Formalism.\allowbreak{}SharpRestriction.\allowbreak{}GrowthBounds.\allowbreak{}IsResidueClassZero}, lines 362--364)\\
\texttt{Formalism.\allowbreak{}MainNotions.\allowbreak{}G.\allowbreak{}Sv} (\texttt{Formalism.\allowbreak{}SharpRestriction.\allowbreak{}GrowthBounds.\allowbreak{}IsResidueClassZero}, lines 336--337)\\
\texttt{Formalism.\allowbreak{}MainNotions.\allowbreak{}G.\allowbreak{}tauP} (\texttt{Formalism.\allowbreak{}SharpRestriction.\allowbreak{}GrowthBounds.\allowbreak{}IsResidueClassZero}, lines 341--344)\\
\texttt{Formalism.\allowbreak{}MainNotions.\allowbreak{}G.\allowbreak{}ioP} (\texttt{Formalism.\allowbreak{}SharpRestriction.\allowbreak{}GrowthBounds.\allowbreak{}IsResidueClassZero}, lines 348--351)\\
\texttt{Formalism.\allowbreak{}MainNotions.\allowbreak{}G.\allowbreak{}tIIf} (\texttt{Formalism.\allowbreak{}SharpRestriction.\allowbreak{}GrowthBounds.\allowbreak{}IsResidueClassZero}, lines 197--200)\\
\texttt{Formalism.\allowbreak{}MainNotions.\allowbreak{}G.\allowbreak{}DIIf} (\texttt{Formalism.\allowbreak{}SharpRestriction.\allowbreak{}GrowthBounds.\allowbreak{}IsResidueClassZero}, lines 188--191)\\
\texttt{Formalism.\allowbreak{}MainNotions.\allowbreak{}G.\allowbreak{}tIf} (\texttt{Formalism.\allowbreak{}SharpRestriction.\allowbreak{}GrowthBounds.\allowbreak{}IsResidueClassZero}, lines 217--225)\\
\texttt{Formalism.\allowbreak{}MainNotions.\allowbreak{}G.\allowbreak{}DIf} (\texttt{Formalism.\allowbreak{}SharpRestriction.\allowbreak{}GrowthBounds.\allowbreak{}IsResidueClassZero}, lines 193--195)\par}

{\footnotesize\raggedright\noindent \ref{B3.52a}.\quad \texttt{Formalism.\allowbreak{}MainNotions.\allowbreak{}G.\allowbreak{}EpDef} (\texttt{Formalism.\allowbreak{}SharpRestriction.\allowbreak{}GrowthBounds.\allowbreak{}DefsInst}, lines 18--19)\\
\texttt{Formalism.\allowbreak{}MainNotions.\allowbreak{}G.\allowbreak{}GoodRankK} (\texttt{Formalism.\allowbreak{}SharpRestriction.\allowbreak{}GrowthBounds.\allowbreak{}UnramifiedQuadraticSetting}, lines 200--203)\\
\texttt{Formalism.\allowbreak{}MainNotions.\allowbreak{}G.\allowbreak{}LPprimeAt} (\texttt{Formalism.\allowbreak{}SharpRestriction.\allowbreak{}GrowthBounds.\allowbreak{}UnramifiedQuadraticSetting}, lines 205--214)\\
\texttt{Formalism.\allowbreak{}MainNotions.\allowbreak{}G.\allowbreak{}LPprime} (\texttt{Formalism.\allowbreak{}SharpRestriction.\allowbreak{}GrowthBounds.\allowbreak{}UnramifiedQuadraticSetting}, lines 216--220)\\
\texttt{Formalism.\allowbreak{}MainNotions.\allowbreak{}G.\allowbreak{}kapISet} (\texttt{Formalism.\allowbreak{}SharpRestriction.\allowbreak{}GrowthBounds.\allowbreak{}DefsInst}, lines 38--45)\\
\texttt{Formalism.\allowbreak{}MainNotions.\allowbreak{}G.\allowbreak{}sInf\_\allowbreak{}mem\_\allowbreak{}kapISet} (\texttt{Formalism.\allowbreak{}SharpRestriction.\allowbreak{}GrowthBounds.\allowbreak{}DefsInst}, lines 68--115)\\
\texttt{Formalism.\allowbreak{}MainNotions.\allowbreak{}G.\allowbreak{}KapIDef} (\texttt{Formalism.\allowbreak{}SharpRestriction.\allowbreak{}GrowthBounds.\allowbreak{}DefsInst}, lines 54--55)\\
\texttt{Formalism.\allowbreak{}MainNotions.\allowbreak{}G.\allowbreak{}exists\_\allowbreak{}inst\_\allowbreak{}defs} (\texttt{Formalism.\allowbreak{}SharpRestriction.\allowbreak{}GrowthBounds.\allowbreak{}DefsInst}, lines 120--124)\\
\texttt{Formalism.\allowbreak{}MainNotions.\allowbreak{}G.\allowbreak{}LatticeSystems.\allowbreak{}H0\_\allowbreak{}eq\_\allowbreak{}twisted\_\allowbreak{}nabla} (\texttt{Formalism.\allowbreak{}SharpRestriction.\allowbreak{}GrowthBounds.\allowbreak{}LatticeSystems}, lines 1403--1409)\\
\texttt{Formalism.\allowbreak{}MainNotions.\allowbreak{}G.\allowbreak{}LatticeSystems.\allowbreak{}la\_\allowbreak{}indep} (\texttt{Formalism.\allowbreak{}SharpRestriction.\allowbreak{}GrowthBounds.\allowbreak{}LatticeSystems}, lines 1504--1546)\par}

{\footnotesize\raggedright\noindent \ref{def:p3c5-parameters}.\quad \texttt{Formalism.\allowbreak{}BasicNotions.\allowbreak{}ParameterChoice'} (\texttt{Formalism.\allowbreak{}LocalRingCardinality.\allowbreak{}DataParams}, lines 62--291)\par}

{\footnotesize\raggedright\noindent \ref{B3.52b}.\quad \texttt{Formalism.\allowbreak{}L2.\allowbreak{}G.\allowbreak{}Model.\allowbreak{}sandwich} (\texttt{Formalism.\allowbreak{}LocalToGlobal.\allowbreak{}GroupLattices.\allowbreak{}FockModels.\allowbreak{}Sandwich}, lines 187--206)\\
\texttt{Formalism.\allowbreak{}MainNotions.\allowbreak{}G.\allowbreak{}Joints.\allowbreak{}weyl\_\allowbreak{}sandwich\_\allowbreak{}of\_\allowbreak{}factors} (\texttt{Formalism.\allowbreak{}SharpRestriction.\allowbreak{}GrowthBounds.\allowbreak{}Joints}, lines 1494--1534)\\
\texttt{Formalism.\allowbreak{}MainNotions.\allowbreak{}G.\allowbreak{}Joints.\allowbreak{}expBound\_\allowbreak{}of} (\texttt{Formalism.\allowbreak{}SharpRestriction.\allowbreak{}GrowthBounds.\allowbreak{}Joints}, lines 1848--1925)\par}

{\footnotesize\raggedright\noindent \ref{B3.53}.\quad \texttt{Formalism.\allowbreak{}MainNotions.\allowbreak{}FormalGroupSystem.\allowbreak{}SmallResidueCharComparison} (\texttt{Formalism.\allowbreak{}SharpRestriction.\allowbreak{}GrowthBounds.\allowbreak{}FormalGroupSystem}, lines 741--754)\\
\texttt{Formalism.\allowbreak{}MainNotions.\allowbreak{}FormalGroupSystem.\allowbreak{}Pf.\allowbreak{}numeric\_\allowbreak{}comparison\_\allowbreak{}of\_\allowbreak{}constants} (\texttt{Formalism.\allowbreak{}SharpRestriction.\allowbreak{}GrowthBounds.\allowbreak{}ProofLW}, lines 207--219)\par}

{\footnotesize\raggedright\noindent \ref{B3.53}.\quad \texttt{Formalism.\allowbreak{}MainNotions.\allowbreak{}FormalGroupSystem.\allowbreak{}Pf.\allowbreak{}nc\_\allowbreak{}loc} (\texttt{Formalism.\allowbreak{}SharpRestriction.\allowbreak{}GrowthBounds.\allowbreak{}MultisetValuation.\allowbreak{}NumericalClosedOfNeed}, lines 12--19)\\
\texttt{Formalism.\allowbreak{}MainNotions.\allowbreak{}FormalGroupSystem.\allowbreak{}Pf.\allowbreak{}lw\_\allowbreak{}size} (\texttt{Formalism.\allowbreak{}SharpRestriction.\allowbreak{}GrowthBounds.\allowbreak{}ProofLW}, lines 135--203)\\
\texttt{Formalism.\allowbreak{}MainNotions.\allowbreak{}FormalGroupSystem.\allowbreak{}Pf.\allowbreak{}lw\_\allowbreak{}de} (\texttt{Formalism.\allowbreak{}SharpRestriction.\allowbreak{}GrowthBounds.\allowbreak{}ProofLW}, lines 62--104)\\
\texttt{Formalism.\allowbreak{}MainNotions.\allowbreak{}FormalGroupSystem.\allowbreak{}Pf.\allowbreak{}nc\_\allowbreak{}gB\_\allowbreak{}ge} (\texttt{Formalism.\allowbreak{}SharpRestriction.\allowbreak{}GrowthBounds.\allowbreak{}MultisetValuation.\allowbreak{}NcSummableMin}, lines 120--146)\\
\texttt{Formalism.\allowbreak{}MainNotions.\allowbreak{}FormalGroupSystem.\allowbreak{}Pf.\allowbreak{}nc\_\allowbreak{}gA\_\allowbreak{}ge} (\texttt{Formalism.\allowbreak{}SharpRestriction.\allowbreak{}GrowthBounds.\allowbreak{}MultisetValuation.\allowbreak{}NcSummableMin}, lines 147--166)\par}

{\footnotesize\raggedright\noindent \ref{B3.54}.\quad \texttt{Formalism.\allowbreak{}MainNotions.\allowbreak{}G.\allowbreak{}TopChoiceBounds.\allowbreak{}hLarge\_\allowbreak{}proved} (\texttt{Formalism.\allowbreak{}SharpRestriction.\allowbreak{}GrowthBounds.\allowbreak{}TopChoiceBounds}, lines 4226--4373)\par}

{\footnotesize\raggedright\noindent \ref{B3.54k}.\quad \texttt{Formalism.\allowbreak{}MainNotions.\allowbreak{}WeilJacquetShalikaLaw.\allowbreak{}TwLoc} (\texttt{Formalism.\allowbreak{}SharpRestriction.\allowbreak{}GrowthBounds.\allowbreak{}WeilJacquetShalikaLaw}, lines 79--83)\\
\texttt{Formalism.\allowbreak{}MainNotions.\allowbreak{}WeilJacquetShalikaLaw.\allowbreak{}TwLoc.\allowbreak{}Yt} (\texttt{Formalism.\allowbreak{}SharpRestriction.\allowbreak{}GrowthBounds.\allowbreak{}WeilJacquetShalikaLaw}, lines 85--90)\\
\texttt{Formalism.\allowbreak{}MainNotions.\allowbreak{}WeilJacquetShalikaLaw.\allowbreak{}Ties9} (\texttt{Formalism.\allowbreak{}SharpRestriction.\allowbreak{}GrowthBounds.\allowbreak{}WeilJacquetShalikaLaw}, lines 92--95)\\
\texttt{Formalism.\allowbreak{}BasicNotions.\allowbreak{}GatTies} (\texttt{Formalism.\allowbreak{}LocalRingCardinality.\allowbreak{}DataGatTies}, lines 7--29)\\
\texttt{Formalism.\allowbreak{}MainNotions.\allowbreak{}G.\allowbreak{}LatTied} (\texttt{Formalism.\allowbreak{}SharpRestriction.\allowbreak{}GrowthBounds.\allowbreak{}IsResidueClassZero}, lines 430--451)\\
\texttt{Formalism.\allowbreak{}BasicNotions.\allowbreak{}LatCmp.\allowbreak{}VramDef} (\texttt{Formalism.\allowbreak{}LocalRingCardinality.\allowbreak{}DataLatDefs}, lines 37--40)\\
\texttt{Formalism.\allowbreak{}MainNotions.\allowbreak{}WeilJacquetShalikaLaw.\allowbreak{}CmpV} (\texttt{Formalism.\allowbreak{}SharpRestriction.\allowbreak{}GrowthBounds.\allowbreak{}WeilJacquetShalikaLaw}, lines 156--159)\\
\texttt{Formalism.\allowbreak{}MainNotions.\allowbreak{}Pk.\allowbreak{}reW} (\texttt{Formalism.\allowbreak{}SharpRestriction.\allowbreak{}SmoothDesignExistence.\allowbreak{}SumOverPrimesHypotheses}, lines 96--96)\par}

{\footnotesize\raggedright\noindent \ref{B3.55}.\quad \texttt{Formalism.\allowbreak{}MainNotions.\allowbreak{}G.\allowbreak{}Joints.\allowbreak{}SysUAt} (\texttt{Formalism.\allowbreak{}SharpRestriction.\allowbreak{}GrowthBounds.\allowbreak{}JointsTw}, lines 9--19)\\
\texttt{Formalism.\allowbreak{}MainNotions.\allowbreak{}FormalGroupSystem.\allowbreak{}Sys} (\texttt{Formalism.\allowbreak{}SharpRestriction.\allowbreak{}GrowthBounds.\allowbreak{}FormalGroupSystem}, lines 264--277)\\
\texttt{Formalism.\allowbreak{}MainNotions.\allowbreak{}FormalGroupSystem.\allowbreak{}OmBasis} (\texttt{Formalism.\allowbreak{}SharpRestriction.\allowbreak{}GrowthBounds.\allowbreak{}FormalGroupSystem}, lines 172--177)\\
\texttt{Formalism.\allowbreak{}MainNotions.\allowbreak{}G.\allowbreak{}yProd} (\texttt{Formalism.\allowbreak{}SharpRestriction.\allowbreak{}GrowthBounds.\allowbreak{}IsResidueClassZero}, lines 453--457)\\
\texttt{Formalism.\allowbreak{}MainNotions.\allowbreak{}G.\allowbreak{}Joints.\allowbreak{}UAt\_\allowbreak{}of} (\texttt{Formalism.\allowbreak{}SharpRestriction.\allowbreak{}GrowthBounds.\allowbreak{}JointsTw}, lines 23--215)\\
\texttt{Formalism.\allowbreak{}MainNotions.\allowbreak{}WeilJacquetShalikaLaw.\allowbreak{}UAt} (\texttt{Formalism.\allowbreak{}SharpRestriction.\allowbreak{}GrowthBounds.\allowbreak{}WeilJacquetShalikaLaw}, lines 161--170)\par}

{\footnotesize\raggedright\noindent \ref{B3.56}.\quad \texttt{Formalism.\allowbreak{}MainNotions.\allowbreak{}G.\allowbreak{}Joints.\allowbreak{}unit\_\allowbreak{}bound\_\allowbreak{}of\_\allowbreak{}sum\_\allowbreak{}over\_\allowbreak{}primes} (\texttt{Formalism.\allowbreak{}SharpRestriction.\allowbreak{}GrowthBounds.\allowbreak{}JointsTw}, lines 243--270)\\
\texttt{Formalism.\allowbreak{}MainNotions.\allowbreak{}Pk.\allowbreak{}SysTw} (\texttt{Formalism.\allowbreak{}SharpRestriction.\allowbreak{}SmoothDesignExistence.\allowbreak{}SumOverPrimesHypotheses}, lines 116--130)\\
\texttt{Formalism.\allowbreak{}MainNotions.\allowbreak{}Pk.\allowbreak{}sum\_\allowbreak{}over\_\allowbreak{}primes\_\allowbreak{}service} (\texttt{Formalism.\allowbreak{}SharpRestriction.\allowbreak{}SmoothDesignExistence.\allowbreak{}SumOverPrimesHypotheses}, lines 187--221)\\
\texttt{Formalism.\allowbreak{}MainNotions.\allowbreak{}WeilJacquetShalikaLaw.\allowbreak{}Serv} (\texttt{Formalism.\allowbreak{}SharpRestriction.\allowbreak{}GrowthBounds.\allowbreak{}WeilJacquetShalikaLaw}, lines 241--243)\\
\texttt{Formalism.\allowbreak{}MainNotions.\allowbreak{}WeilJacquetShalikaLaw.\allowbreak{}ServF} (\texttt{Formalism.\allowbreak{}SharpRestriction.\allowbreak{}GrowthBounds.\allowbreak{}WeilJacquetShalikaLaw}, lines 252--255)\\
\texttt{Formalism.\allowbreak{}MainNotions.\allowbreak{}WeilJacquetShalikaLaw.\allowbreak{}local\_\allowbreak{}theorem\_\allowbreak{}of\_\allowbreak{}statement} (\texttt{Formalism.\allowbreak{}SharpRestriction.\allowbreak{}GrowthBounds.\allowbreak{}WeilJacquetShalikaLaw}, lines 371--375)\\
\texttt{Formalism.\allowbreak{}MainNotions.\allowbreak{}WeilJacquetShalikaLaw.\allowbreak{}serv\_\allowbreak{}of} (\texttt{Formalism.\allowbreak{}SharpRestriction.\allowbreak{}GrowthBounds.\allowbreak{}WeilJacquetShalikaLaw}, lines 376--379)\\
\texttt{Formalism.\allowbreak{}MainNotions.\allowbreak{}Pk.\allowbreak{}serv\_\allowbreak{}re} (\texttt{Formalism.\allowbreak{}SharpRestriction.\allowbreak{}SmoothDesignExistence.\allowbreak{}SumOverPrimesHypotheses}, lines 98--107)\par}

\section*{Chapter~\ref{ch:layer-records}: What each layer records}

{\footnotesize\raggedright\noindent \ref{B4.01}.\quad \texttt{Formalism.\allowbreak{}BasicNotions.\allowbreak{}lpF} (\texttt{Formalism.\allowbreak{}LocalRingCardinality.\allowbreak{}DataParams₂}, lines 58--60)\\
\texttt{Formalism.\allowbreak{}BasicNotions.\allowbreak{}FOrb} (\texttt{Formalism.\allowbreak{}LocalRingCardinality.\allowbreak{}DataForced}, lines 99--100)\par}

{\footnotesize\raggedright\noindent \ref{B4.01d}.\quad \texttt{Formalism.\allowbreak{}IncomingBounds.\allowbreak{}FlowLoadBounds.\allowbreak{}exists\_\allowbreak{}opt\_\allowbreak{}flow} (\texttt{Formalism.\allowbreak{}IncomingBounds.\allowbreak{}FlowDegreeBounds.\allowbreak{}MaxFlowValueAttained}, lines 108--136)\par}

{\footnotesize\raggedright\noindent \ref{B4.01d}.\quad \texttt{Formalism.\allowbreak{}BasicNotions.\allowbreak{}FOrb\_\allowbreak{}step} (\texttt{Formalism.\allowbreak{}LocalRingCardinality.\allowbreak{}DataForced}, lines 102--108)\par}

{\footnotesize\raggedright\noindent \ref{B4.01a}.\quad \texttt{UpperTheory.\allowbreak{}exists\_\allowbreak{}maxFlow\_\allowbreak{}attained\_\allowbreak{}and\_\allowbreak{}isOptDual} (\texttt{Development.\allowbreak{}AutomorphicArithmetic.\allowbreak{}OrbitLabellings.\allowbreak{}ExistsMaxFlowAttainedAndIsOptDual}, lines 8--18)\\
\texttt{UpperTheory.\allowbreak{}sum\_\allowbreak{}demand\_\allowbreak{}filter\_\allowbreak{}two\_\allowbreak{}div\_\allowbreak{}le\_\allowbreak{}of\_\allowbreak{}isOptDual} (\texttt{Development.\allowbreak{}AutomorphicArithmetic.\allowbreak{}OrbitLabellings.\allowbreak{}SumDemandFilterTwoDivLeOfIsOptDual}, lines 9--57)\\
\texttt{Formalism.\allowbreak{}DworkPowerSeriesWeights.\allowbreak{}OddLevelGainBounds.\allowbreak{}proposition\_\allowbreak{}six} (\texttt{Formalism.\allowbreak{}GaussJetBounds.\allowbreak{}OneLetterGaussJets.\allowbreak{}OneLetterGaussProofs.\allowbreak{}CompanionTheory}, lines 38--147)\\
\texttt{Formalism.\allowbreak{}MainNotions.\allowbreak{}DIV.\allowbreak{}Bridge.\allowbreak{}DworkPowerSeriesWeights.\allowbreak{}local\_\allowbreak{}bound\_\allowbreak{}eleven\_\allowbreak{}comparison} (\texttt{Formalism.\allowbreak{}SharpRestriction.\allowbreak{}Divisibility.\allowbreak{}JetPolynomialLogs.\allowbreak{}LocalBoundElevenComparison}, lines 33--55)\par}

{\footnotesize\raggedright\noindent \ref{B4.01c}.\quad \texttt{Formalism.\allowbreak{}DworkPowerSeriesWeights.\allowbreak{}generalized\_\allowbreak{}log\_\allowbreak{}inverse} (\texttt{Formalism.\allowbreak{}GaussJetBounds.\allowbreak{}GaussJetProofs.\allowbreak{}LogarithmWeightComparison}, lines 227--251)\\
\texttt{Formalism.\allowbreak{}DworkPowerSeriesWeights.\allowbreak{}generalized\_\allowbreak{}log\_\allowbreak{}weight} (\texttt{Formalism.\allowbreak{}GaussJetBounds.\allowbreak{}GaussJetProofs.\allowbreak{}LogarithmWeightComparison}, lines 252--331)\\
\texttt{Formalism.\allowbreak{}DworkPowerSeriesWeights.\allowbreak{}lemmaS} (\texttt{Formalism.\allowbreak{}GaussJetBounds.\allowbreak{}OneLetterRadius}, lines 17--20)\par}

{\footnotesize\raggedright\noindent \ref{B4.01b}.\quad \texttt{Formalism.\allowbreak{}MainNotions.\allowbreak{}DIV.\allowbreak{}Bridge.\allowbreak{}DworkPowerSeriesWeights.\allowbreak{}split\_\allowbreak{}bound\_\allowbreak{}eleven\_\allowbreak{}one} (\texttt{Formalism.\allowbreak{}SharpRestriction.\allowbreak{}Divisibility.\allowbreak{}JetPolynomialLogs.\allowbreak{}LocalBoundElevenComparison}, lines 64--68)\par}

{\footnotesize\raggedright\noindent \ref{B4.02z}.\quad \texttt{Formalism.\allowbreak{}BasicNotions.\allowbreak{}DV.\allowbreak{}KapDef} (\texttt{Formalism.\allowbreak{}LocalRingCardinality.\allowbreak{}DataKap}, lines 13--35)\\
\texttt{Formalism.\allowbreak{}BasicNotions.\allowbreak{}DV.\allowbreak{}aPlusQ\_\allowbreak{}eq} (\texttt{Formalism.\allowbreak{}LocalRingCardinality.\allowbreak{}DataKap}, lines 60--65)\par}

{\footnotesize\raggedright\noindent \ref{B4.03a}.\quad \texttt{Formalism.\allowbreak{}BasicNotions.\allowbreak{}HCParam} (\texttt{Formalism.\allowbreak{}LocalRingCardinality.\allowbreak{}Params}, lines 50--69)\par}

{\footnotesize\raggedright\noindent \ref{B4.02a}.\quad \texttt{Formalism.\allowbreak{}BasicNotions.\allowbreak{}HCParam} (\texttt{Formalism.\allowbreak{}LocalRingCardinality.\allowbreak{}Params}, lines 50--69)\\
\texttt{Formalism.\allowbreak{}BasicNotions.\allowbreak{}DV.\allowbreak{}KapDef} (\texttt{Formalism.\allowbreak{}LocalRingCardinality.\allowbreak{}DataKap}, lines 31--35)\\
\texttt{Formalism.\allowbreak{}BasicNotions.\allowbreak{}DV.\allowbreak{}aPlusQ\_\allowbreak{}eq} (\texttt{Formalism.\allowbreak{}LocalRingCardinality.\allowbreak{}DataKap}, lines 60--65)\\
\texttt{GlobalTheory.\allowbreak{}Mem.\allowbreak{}R3.\allowbreak{}Rw.\allowbreak{}Ex.\allowbreak{}kapDef\_\allowbreak{}of\_\allowbreak{}ka} (\texttt{Development.\allowbreak{}AbelianVarietyTheory.\allowbreak{}Proofs.\allowbreak{}IwasawaDecomposition.\allowbreak{}ExistsRowsDatumPrescribedTuplesDefinedWeights}, lines 31--43)\\
\texttt{Formalism.\allowbreak{}BasicNotions.\allowbreak{}RowWeightLemmas.\allowbreak{}nuRow\_\allowbreak{}head} (\texttt{Formalism.\allowbreak{}LocalRingCardinality.\allowbreak{}CoeNeZeroOfNotMem}, lines 12--21)\\
\texttt{Formalism.\allowbreak{}BasicNotions.\allowbreak{}RowWeightLemmas.\allowbreak{}nuRow\_\allowbreak{}block} (\texttt{Formalism.\allowbreak{}LocalRingCardinality.\allowbreak{}CoeNeZeroOfNotMem}, lines 22--40)\\
\texttt{Formalism.\allowbreak{}BasicNotions.\allowbreak{}RowWeightLemmas.\allowbreak{}nuRow\_\allowbreak{}pos} (\texttt{Formalism.\allowbreak{}LocalRingCardinality.\allowbreak{}CoeNeZeroOfNotMem}, lines 180--183)\par}

{\footnotesize\raggedright\noindent \ref{B4.02b}.\quad \texttt{GlobalTheory.\allowbreak{}EX.\allowbreak{}HCVal} (\texttt{Development.\allowbreak{}AbelianVarietyTheory.\allowbreak{}Definitions.\allowbreak{}GlobalTwisting.\allowbreak{}PairRowIffQdRows}, lines 93--106)\\
\texttt{GlobalTheory.\allowbreak{}EX.\allowbreak{}exists\_\allowbreak{}hcParam\_\allowbreak{}of\_\allowbreak{}code} (\texttt{Development.\allowbreak{}AbelianVarietyTheory.\allowbreak{}Proofs.\allowbreak{}GlobalInterpolation.\allowbreak{}ExistsHcParamOfCode}, lines 12--76)\\
\texttt{GlobalTheory.\allowbreak{}EX.\allowbreak{}exists\_\allowbreak{}hc\_\allowbreak{}family} (\texttt{Development.\allowbreak{}AbelianVarietyTheory.\allowbreak{}Proofs.\allowbreak{}GlobalInterpolation.\allowbreak{}ExistsHcParamOfCode}, lines 128--242)\par}

{\footnotesize\raggedright\noindent \ref{B4.02}.\quad \texttt{Formalism.\allowbreak{}BasicNotions.\allowbreak{}SizeData₂.\allowbreak{}FDef} (\texttt{Formalism.\allowbreak{}LocalRingCardinality.\allowbreak{}DataForced}, lines 134--135)\par}

{\footnotesize\raggedright\noindent \ref{B4.02y}.\quad \texttt{Formalism.\allowbreak{}BasicNotions.\allowbreak{}sum\_\allowbreak{}orbits} (\texttt{Formalism.\allowbreak{}LocalRingCardinality.\allowbreak{}DataForced}, lines 170--189)\\
\texttt{Formalism.\allowbreak{}BasicNotions.\allowbreak{}primeClass} (\texttt{Formalism.\allowbreak{}LocalRingCardinality.\allowbreak{}DataHost}, lines 33--37)\par}

{\footnotesize\raggedright\noindent \ref{B4.03}.\quad \texttt{Formalism.\allowbreak{}BasicNotions.\allowbreak{}ErrX} (\texttt{Formalism.\allowbreak{}LocalRingCardinality.\allowbreak{}DataErr}, lines 20--35)\\
\texttt{Formalism.\allowbreak{}BasicNotions.\allowbreak{}ErrDefs} (\texttt{Formalism.\allowbreak{}LocalRingCardinality.\allowbreak{}DataErr}, lines 207--214)\\
\texttt{Formalism.\allowbreak{}BasicNotions.\allowbreak{}SizeData₂.\allowbreak{}EDef} (\texttt{Formalism.\allowbreak{}LocalRingCardinality.\allowbreak{}DataErr}, lines 171--172)\\
\texttt{Formalism.\allowbreak{}BasicNotions.\allowbreak{}DGamDef} (\texttt{Formalism.\allowbreak{}LocalRingCardinality.\allowbreak{}DataErr}, lines 74--75)\\
\texttt{Formalism.\allowbreak{}BasicNotions.\allowbreak{}Eval} (\texttt{Formalism.\allowbreak{}LocalRingCardinality.\allowbreak{}DataErr}, lines 159--169)\par}

{\footnotesize\raggedright\noindent \ref{B4.03b}.\quad \texttt{Formalism.\allowbreak{}BasicNotions.\allowbreak{}DGamDef.\allowbreak{}one} (\texttt{Formalism.\allowbreak{}LocalRingCardinality.\allowbreak{}DataErr}, lines 77--80)\par}

{\footnotesize\raggedright\noindent \ref{B4.04}.\quad \texttt{Formalism.\allowbreak{}BasicNotions.\allowbreak{}SizeData₂.\allowbreak{}gainD} (\texttt{Formalism.\allowbreak{}LocalRingCardinality.\allowbreak{}DataParams₂}, lines 212--213)\\
\texttt{Formalism.\allowbreak{}MainNotions.\allowbreak{}LocalQuadraticParts.\allowbreak{}DropDef} (\texttt{Formalism.\allowbreak{}SharpRestriction.\allowbreak{}Divisibility.\allowbreak{}LocalQuadraticParts}, lines 517--518)\\
\texttt{Formalism.\allowbreak{}MainNotions.\allowbreak{}NoGain} (\texttt{Formalism.\allowbreak{}SharpRestriction.\allowbreak{}ParConsts}, lines 24--29)\\
\texttt{Formalism.\allowbreak{}MainNotions.\allowbreak{}G.\allowbreak{}ParametersFill} (\texttt{Formalism.\allowbreak{}SharpRestriction.\allowbreak{}GrowthBounds.\allowbreak{}IsResidueClassZero}, lines 401--411)\par}

{\footnotesize\raggedright\noindent \ref{B4.04a}.\quad \texttt{Formalism.\allowbreak{}BasicNotions.\allowbreak{}SizeData₂} (\texttt{Formalism.\allowbreak{}LocalRingCardinality.\allowbreak{}DataParams₂}, lines 94--104)\\
\texttt{Formalism.\allowbreak{}MainNotions.\allowbreak{}LocalQuadraticParts.\allowbreak{}DropDef} (\texttt{Formalism.\allowbreak{}SharpRestriction.\allowbreak{}Divisibility.\allowbreak{}LocalQuadraticParts}, lines 517--518)\par}

{\footnotesize\raggedright\noindent \ref{B4.05}.\quad \texttt{Formalism.\allowbreak{}BasicNotions.\allowbreak{}SizeData₂.\allowbreak{}sg} (\texttt{Formalism.\allowbreak{}LocalRingCardinality.\allowbreak{}DataParams₂}, lines 215--217)\par}

{\footnotesize\raggedright\noindent \ref{B4.05a}.\quad \texttt{Formalism.\allowbreak{}BasicNotions.\allowbreak{}SizeData₂} (\texttt{Formalism.\allowbreak{}LocalRingCardinality.\allowbreak{}DataParams₂}, lines 94--102)\\
\texttt{Formalism.\allowbreak{}BasicNotions.\allowbreak{}SizeData₂.\allowbreak{}sg} (\texttt{Formalism.\allowbreak{}LocalRingCardinality.\allowbreak{}DataParams₂}, lines 215--217)\par}

\section*{Chapter~\ref{ch:local-to-column}: From local lifts to bounds on whole columns}

{\footnotesize\raggedright\noindent \ref{B4.06}.\quad \texttt{Formalism.\allowbreak{}MainNotions.\allowbreak{}WeilJacquetShalikaLaw.\allowbreak{}VGe} (\texttt{Formalism.\allowbreak{}SharpRestriction.\allowbreak{}GrowthBounds.\allowbreak{}WeilJacquetShalikaLaw}, lines 19--23)\\
\texttt{Formalism.\allowbreak{}MainNotions.\allowbreak{}WeilJacquetShalikaLaw.\allowbreak{}ColGe0} (\texttt{Formalism.\allowbreak{}SharpRestriction.\allowbreak{}GrowthBounds.\allowbreak{}WeilJacquetShalikaLaw}, lines 24--26)\\
\texttt{Formalism.\allowbreak{}MainNotions.\allowbreak{}WeilJacquetShalikaLaw.\allowbreak{}ColGeTw} (\texttt{Formalism.\allowbreak{}SharpRestriction.\allowbreak{}GrowthBounds.\allowbreak{}WeilJacquetShalikaLaw}, lines 28--31)\\
\texttt{Formalism.\allowbreak{}MainNotions.\allowbreak{}WeilJacquetShalikaLaw.\allowbreak{}ColGe} (\texttt{Formalism.\allowbreak{}SharpRestriction.\allowbreak{}GrowthBounds.\allowbreak{}WeilJacquetShalikaLaw}, lines 33--35)\par}

{\footnotesize\raggedright\noindent \ref{B4.07}.\quad \texttt{Formalism.\allowbreak{}MainNotions.\allowbreak{}G.\allowbreak{}NP} (\texttt{Formalism.\allowbreak{}SharpRestriction.\allowbreak{}GrowthBounds.\allowbreak{}IsResidueClassZero}, lines 464--473)\\
\texttt{Formalism.\allowbreak{}MainNotions.\allowbreak{}WeilJacquetShalikaLaw.\allowbreak{}InNP} (\texttt{Formalism.\allowbreak{}SharpRestriction.\allowbreak{}GrowthBounds.\allowbreak{}WeilJacquetShalikaLaw}, lines 60--64)\par}

{\footnotesize\raggedright\noindent \ref{B4.07c}.\quad \texttt{Formalism.\allowbreak{}MainNotions.\allowbreak{}G.\allowbreak{}Joints.\allowbreak{}Lit.\allowbreak{}NPDesc\_\allowbreak{}of} (\texttt{Formalism.\allowbreak{}SharpRestriction.\allowbreak{}GrowthBounds.\allowbreak{}JointsLit}, lines 256--347)\\
\texttt{Formalism.\allowbreak{}BasicNotions.\allowbreak{}LatCmp.\allowbreak{}M\_\allowbreak{}inf\_\allowbreak{}Vram} (\texttt{Formalism.\allowbreak{}LocalRingCardinality.\allowbreak{}Lemmas}, lines 968--983)\\
\texttt{GlobalTheory.\allowbreak{}EX.\allowbreak{}OB.\allowbreak{}NP\_\allowbreak{}full} (\texttt{Development.\allowbreak{}AbelianVarietyTheory.\allowbreak{}Proofs.\allowbreak{}GlobalInterpolation.\allowbreak{}MemSpanOfSemilinear}, lines 79--127)\par}

{\footnotesize\raggedright\noindent \ref{B4.07b}.\quad \texttt{Formalism.\allowbreak{}MainNotions.\allowbreak{}G.\allowbreak{}LatticeSystems.\allowbreak{}lb\_\allowbreak{}steps} (\texttt{Formalism.\allowbreak{}SharpRestriction.\allowbreak{}GrowthBounds.\allowbreak{}LatticeSystems}, lines 1790--1821)\\
\texttt{Formalism.\allowbreak{}MainNotions.\allowbreak{}G.\allowbreak{}LatticeSystems.\allowbreak{}ch\_\allowbreak{}locChain} (\texttt{Formalism.\allowbreak{}SharpRestriction.\allowbreak{}GrowthBounds.\allowbreak{}LatticeSystems}, lines 856--896)\\
\texttt{Formalism.\allowbreak{}MainNotions.\allowbreak{}G.\allowbreak{}LatticeSystems.\allowbreak{}jm\_\allowbreak{}Re} (\texttt{Formalism.\allowbreak{}SharpRestriction.\allowbreak{}GrowthBounds.\allowbreak{}LatticeSystems}, lines 3727--3745)\\
\texttt{Formalism.\allowbreak{}MainNotions.\allowbreak{}G.\allowbreak{}LatticeSystems.\allowbreak{}jStep1} (\texttt{Formalism.\allowbreak{}SharpRestriction.\allowbreak{}GrowthBounds.\allowbreak{}LatticeSystems}, lines 3339--3460)\par}

{\footnotesize\raggedright\noindent \ref{B4.07a}.\quad \texttt{Formalism.\allowbreak{}MainNotions.\allowbreak{}G.\allowbreak{}LPprime} (\texttt{Formalism.\allowbreak{}SharpRestriction.\allowbreak{}GrowthBounds.\allowbreak{}UnramifiedQuadraticSetting}, lines 216--221)\\
\texttt{Formalism.\allowbreak{}MainNotions.\allowbreak{}G.\allowbreak{}LatticeSystems.\allowbreak{}jSmith} (\texttt{Formalism.\allowbreak{}SharpRestriction.\allowbreak{}GrowthBounds.\allowbreak{}LatticeSystems}, lines 3746--3817)\\
\texttt{Formalism.\allowbreak{}MainNotions.\allowbreak{}G.\allowbreak{}LatticeSystems.\allowbreak{}jl\_\allowbreak{}good} (\texttt{Formalism.\allowbreak{}SharpRestriction.\allowbreak{}GrowthBounds.\allowbreak{}LatticeSystems}, lines 3875--3931)\\
\texttt{Formalism.\allowbreak{}MainNotions.\allowbreak{}G.\allowbreak{}LatticeSystems.\allowbreak{}jLocal} (\texttt{Formalism.\allowbreak{}SharpRestriction.\allowbreak{}GrowthBounds.\allowbreak{}LatticeSystems}, lines 3972--4063)\\
\texttt{Formalism.\allowbreak{}MainNotions.\allowbreak{}G.\allowbreak{}LatticeSystems.\allowbreak{}jRank} (\texttt{Formalism.\allowbreak{}SharpRestriction.\allowbreak{}GrowthBounds.\allowbreak{}LatticeSystems}, lines 4194--4237)\\
\texttt{Formalism.\allowbreak{}MainNotions.\allowbreak{}G.\allowbreak{}LatticeSystems.\allowbreak{}ja\_\allowbreak{}restJ} (\texttt{Formalism.\allowbreak{}SharpRestriction.\allowbreak{}GrowthBounds.\allowbreak{}LatticeSystems}, lines 4300--4389)\\
\texttt{Formalism.\allowbreak{}G0.\allowbreak{}WalkSharpP.\allowbreak{}main} (\texttt{Formalism.\allowbreak{}SubspaceCounting.\allowbreak{}SubspaceCountingProofs.\allowbreak{}WalkSharp}, lines 379--719)\\
\texttt{Formalism.\allowbreak{}MainNotions.\allowbreak{}G.\allowbreak{}LatticeSystems.\allowbreak{}walk\_\allowbreak{}allBut\_\allowbreak{}bound} (\texttt{Formalism.\allowbreak{}SharpRestriction.\allowbreak{}GrowthBounds.\allowbreak{}LatticeSystems}, lines 217--313)\\
\texttt{Formalism.\allowbreak{}MainNotions.\allowbreak{}FormalGroupSystem.\allowbreak{}Pf.\allowbreak{}stmt6G} (\texttt{Formalism.\allowbreak{}SharpRestriction.\allowbreak{}GrowthBounds.\allowbreak{}MultisetValuation.\allowbreak{}StabilizerCount}, lines 384--491)\par}

{\footnotesize\raggedright\noindent \ref{B4.08}.\quad \texttt{Formalism.\allowbreak{}MainNotions.\allowbreak{}WeilJacquetShalikaLaw.\allowbreak{}lowS} (\texttt{Formalism.\allowbreak{}SharpRestriction.\allowbreak{}GrowthBounds.\allowbreak{}WeilJacquetShalikaLaw}, lines 73--75)\\
\texttt{Formalism.\allowbreak{}MainNotions.\allowbreak{}WeilJacquetShalikaLaw.\allowbreak{}totS} (\texttt{Formalism.\allowbreak{}SharpRestriction.\allowbreak{}GrowthBounds.\allowbreak{}WeilJacquetShalikaLaw}, lines 76--78)\\
\texttt{Formalism.\allowbreak{}MainNotions.\allowbreak{}WeilJacquetShalikaLaw.\allowbreak{}gainS} (\texttt{Formalism.\allowbreak{}SharpRestriction.\allowbreak{}GrowthBounds.\allowbreak{}WeilJacquetShalikaLaw}, lines 49--51)\\
\texttt{Formalism.\allowbreak{}MainNotions.\allowbreak{}G.\allowbreak{}gainP} (\texttt{Formalism.\allowbreak{}SharpRestriction.\allowbreak{}GrowthBounds.\allowbreak{}IsResidueClassZero}, lines 355--358)\par}

{\footnotesize\raggedright\noindent \ref{B4.08}.\quad \texttt{Formalism.\allowbreak{}MainNotions.\allowbreak{}WeilJacquetShalikaLaw.\allowbreak{}TwLoc} (\texttt{Formalism.\allowbreak{}SharpRestriction.\allowbreak{}GrowthBounds.\allowbreak{}WeilJacquetShalikaLaw}, lines 79--83)\\
\texttt{Formalism.\allowbreak{}MainNotions.\allowbreak{}WeilJacquetShalikaLaw.\allowbreak{}TwLoc.\allowbreak{}Yt} (\texttt{Formalism.\allowbreak{}SharpRestriction.\allowbreak{}GrowthBounds.\allowbreak{}WeilJacquetShalikaLaw}, lines 85--90)\par}

{\footnotesize\raggedright\noindent \ref{B4.09}.\quad \texttt{Formalism.\allowbreak{}MainNotions.\allowbreak{}G.\allowbreak{}Joints.\allowbreak{}unit\_\allowbreak{}bound\_\allowbreak{}of\_\allowbreak{}sum\_\allowbreak{}over\_\allowbreak{}primes} (\texttt{Formalism.\allowbreak{}SharpRestriction.\allowbreak{}GrowthBounds.\allowbreak{}JointsTw}, lines 243--271)\\
\texttt{Formalism.\allowbreak{}MainNotions.\allowbreak{}WeilJacquetShalikaLaw.\allowbreak{}UAt} (\texttt{Formalism.\allowbreak{}SharpRestriction.\allowbreak{}GrowthBounds.\allowbreak{}WeilJacquetShalikaLaw}, lines 161--171)\par}

{\footnotesize\raggedright\noindent \ref{B4.10}.\quad \texttt{Formalism.\allowbreak{}MainNotions.\allowbreak{}WeilJacquetShalikaLaw.\allowbreak{}Proof.\allowbreak{}first\_\allowbreak{}in\_\allowbreak{}local\_\allowbreak{}models} (\texttt{Formalism.\allowbreak{}SharpRestriction.\allowbreak{}GrowthBounds.\allowbreak{}AdmissibleIntegers.\allowbreak{}StatementsOfRemainingBound}, lines 111--152)\par}

{\footnotesize\raggedright\noindent \ref{B4.10b}.\quad \texttt{GlobalTheory.\allowbreak{}KerLatS.\allowbreak{}kerEntHyp\_\allowbreak{}latS} (\texttt{Development.\allowbreak{}AbelianVarietyTheory.\allowbreak{}Proofs.\allowbreak{}IwasawaDecomposition.\allowbreak{}TrueChainsKernelDatumSourceSideLattice}, lines 36--65)\\
\texttt{GlobalTheory.\allowbreak{}Mem.\allowbreak{}Col2.\allowbreak{}TC.\allowbreak{}kerEnt\_\allowbreak{}latT} (\texttt{Development.\allowbreak{}AbelianVarietyTheory.\allowbreak{}Proofs.\allowbreak{}IwasawaDecomposition.\allowbreak{}KernelDatumLatticeTargetSideTrueChains}, lines 18--37)\\
\texttt{GlobalTheory.\allowbreak{}KerVal.\allowbreak{}onMin\_\allowbreak{}canInv\_\allowbreak{}mem} (\texttt{Development.\allowbreak{}AbelianVarietyTheory.\allowbreak{}Proofs.\allowbreak{}IwasawaDecomposition.\allowbreak{}TargetPairingSizesInCoefficientFieldValueGroup}, lines 89--159)\\
\texttt{GlobalTheory.\allowbreak{}KerVal.\allowbreak{}prTC\_\allowbreak{}val} (\texttt{Development.\allowbreak{}AbelianVarietyTheory.\allowbreak{}Proofs.\allowbreak{}IwasawaDecomposition.\allowbreak{}TargetPairingSizesInCoefficientFieldValueGroup}, lines 207--214)\par}

{\footnotesize\raggedright\noindent \ref{B4.10a}.\quad \texttt{Formalism.\allowbreak{}MainNotions.\allowbreak{}DIV.\allowbreak{}Lemmas.\allowbreak{}theta\_\allowbreak{}isometry\_\allowbreak{}facts\_\allowbreak{}hold} (\texttt{Formalism.\allowbreak{}SharpRestriction.\allowbreak{}Divisibility.\allowbreak{}ChainOperators.\allowbreak{}GenericFramesLetters}, lines 10--13)\\
\texttt{Formalism.\allowbreak{}MainNotions.\allowbreak{}OrbitPlaces.\allowbreak{}kform\_\allowbreak{}of\_\allowbreak{}ker} (\texttt{Formalism.\allowbreak{}SharpRestriction.\allowbreak{}Divisibility.\allowbreak{}OrbitPlaces}, lines 1621--1657)\par}

{\footnotesize\raggedright\noindent \ref{def:p4c2-orbit}.\quad \texttt{Formalism.\allowbreak{}MainNotions.\allowbreak{}OrbitPlaces.\allowbreak{}Ent} (\texttt{Formalism.\allowbreak{}SharpRestriction.\allowbreak{}Divisibility.\allowbreak{}OrbitPlaces}, lines 34--50)\\
\texttt{Formalism.\allowbreak{}MainNotions.\allowbreak{}OrbitPlaces.\allowbreak{}KForm} (\texttt{Formalism.\allowbreak{}SharpRestriction.\allowbreak{}Divisibility.\allowbreak{}OrbitPlaces}, lines 71--73)\\
\texttt{Formalism.\allowbreak{}MainNotions.\allowbreak{}OrbitPlaces.\allowbreak{}AdmK} (\texttt{Formalism.\allowbreak{}SharpRestriction.\allowbreak{}Divisibility.\allowbreak{}OrbitPlaces}, lines 475--477)\\
\texttt{Formalism.\allowbreak{}MainNotions.\allowbreak{}OrbitPlaces.\allowbreak{}admitted} (\texttt{Formalism.\allowbreak{}SharpRestriction.\allowbreak{}Divisibility.\allowbreak{}OrbitPlaces}, lines 479--481)\\
\texttt{Formalism.\allowbreak{}MainNotions.\allowbreak{}OrbitPlaces.\allowbreak{}Core} (\texttt{Formalism.\allowbreak{}SharpRestriction.\allowbreak{}Divisibility.\allowbreak{}OrbitPlaces}, lines 1323--1324)\par}

{\footnotesize\raggedright\noindent \ref{B4.11}.\quad \texttt{Formalism.\allowbreak{}MainNotions.\allowbreak{}OrbitPlaces.\allowbreak{}lift\_\allowbreak{}bounds\_\allowbreak{}of\_\allowbreak{}essential} (\texttt{Formalism.\allowbreak{}SharpRestriction.\allowbreak{}Divisibility.\allowbreak{}OrbitPlaces}, lines 1340--1344)\\
\texttt{Formalism.\allowbreak{}MainNotions.\allowbreak{}J.\allowbreak{}four\_\allowbreak{}orbit\_\allowbreak{}conditions\_\allowbreak{}of\_\allowbreak{}proof} (\texttt{Formalism.\allowbreak{}SharpRestriction.\allowbreak{}AdmissibleParameterJoins.\allowbreak{}GandSab}, lines 11--17)\par}

{\footnotesize\raggedright\noindent \ref{B4.11a}.\quad \texttt{PI.\allowbreak{}OrbBoundFW.\allowbreak{}prvF\_\allowbreak{}of\_\allowbreak{}orbitLetterBounds} (\texttt{Formalism.\allowbreak{}PeriodIntegrals.\allowbreak{}UnitaryGroupMassFormula.\allowbreak{}OrbitValuationBounds.\allowbreak{}PrvFOfOrbitLetterBounds}, lines 17--59)\\
\texttt{PI.\allowbreak{}OrbBoundFW.\allowbreak{}orbit\_\allowbreak{}place\_\allowbreak{}term\_\allowbreak{}le\_\allowbreak{}sum} (\texttt{Formalism.\allowbreak{}PeriodIntegrals.\allowbreak{}UnitaryGroupMassFormula.\allowbreak{}OrbitValuationBounds.\allowbreak{}SummableTruncatedGeometric}, lines 136--218)\\
\texttt{SumOrbCW.\allowbreak{}sumOrbC} (\texttt{Formalism.\allowbreak{}PeriodIntegrals.\allowbreak{}UnitaryGroupMassFormula.\allowbreak{}OrbitSumCount.\allowbreak{}SizeGe}, lines 36--78)\\
\texttt{UpperTheory.\allowbreak{}norm\_\allowbreak{}mfact\_\allowbreak{}mul\_\allowbreak{}coeff\_\allowbreak{}le\_\allowbreak{}of\_\allowbreak{}B0S} (\texttt{Development.\allowbreak{}AutomorphicArithmetic.\allowbreak{}OrbitLabellings.\allowbreak{}NormMfactMulCoeffLeOfB0S}, lines 69--137)\\
\texttt{UpperTheory.\allowbreak{}fInert\_\allowbreak{}sub\_\allowbreak{}eatI\_\allowbreak{}le\_\allowbreak{}orbit\_\allowbreak{}average\_\allowbreak{}sub\_\allowbreak{}errors} (\texttt{Development.\allowbreak{}AutomorphicArithmetic.\allowbreak{}OrbitLabellings.\allowbreak{}FInertSubEatILeOrbitAverageSubErrors}, lines 7--71)\\
\texttt{PI.\allowbreak{}OrbBoundF} (\texttt{Formalism.\allowbreak{}PeriodIntegrals.\allowbreak{}OrbitalBoundsConclusion.\allowbreak{}OrbBoundF}, lines 3--5)\\
\texttt{UpperTheory.\allowbreak{}kvalGe\_\allowbreak{}theta\_\allowbreak{}of\_\allowbreak{}shape} (\texttt{Development.\allowbreak{}AutomorphicArithmetic.\allowbreak{}OrbitLabellings.\allowbreak{}KvalGeThetaOfShape}, lines 155--205)\\
\texttt{jet\_\allowbreak{}bound\_\allowbreak{}ramified\_\allowbreak{}prime} (\texttt{Formalism.\allowbreak{}PeriodIntegrals.\allowbreak{}OrbitalBoundsConclusion.\allowbreak{}EntAtD}, lines 11--14)\\
\texttt{Formalism.\allowbreak{}MainNotions.\allowbreak{}DIV.\allowbreak{}Lemmas.\allowbreak{}OrbitEntryBounds.\allowbreak{}orbit\_\allowbreak{}entry\_\allowbreak{}bounds\_\allowbreak{}hold} (\texttt{Formalism.\allowbreak{}SharpRestriction.\allowbreak{}Divisibility.\allowbreak{}ChainOperators.\allowbreak{}OrbitEntryBounds}, lines 49--53)\\
\texttt{PI.\allowbreak{}OrbBoundFW.\allowbreak{}prvF\_\allowbreak{}of\_\allowbreak{}orbitLetterBounds} (\texttt{Formalism.\allowbreak{}PeriodIntegrals.\allowbreak{}UnitaryGroupMassFormula.\allowbreak{}OrbitValuationBounds.\allowbreak{}PrvFOfOrbitLetterBounds}, lines 17--59)\\
\texttt{jet\_\allowbreak{}bound\_\allowbreak{}ramified\_\allowbreak{}prime} (\texttt{Formalism.\allowbreak{}PeriodIntegrals.\allowbreak{}OrbitalBoundsConclusion.\allowbreak{}EntAtD}, lines 11--13)\par}

{\footnotesize\raggedright\noindent \ref{B4.12}.\quad \texttt{Formalism.\allowbreak{}MainNotions.\allowbreak{}WeilJacquetShalikaLaw.\allowbreak{}Proof.\allowbreak{}chain\_\allowbreak{}ge} (\texttt{Formalism.\allowbreak{}SharpRestriction.\allowbreak{}GrowthBounds.\allowbreak{}AdmissibleIntegers.\allowbreak{}RemainingPrimesOfAdmitted}, lines 144--162)\\
\texttt{Formalism.\allowbreak{}MainNotions.\allowbreak{}WeilJacquetShalikaLaw.\allowbreak{}Proof.\allowbreak{}remaining\_\allowbreak{}primes\_\allowbreak{}of\_\allowbreak{}admitted} (\texttt{Formalism.\allowbreak{}SharpRestriction.\allowbreak{}GrowthBounds.\allowbreak{}AdmissibleIntegers.\allowbreak{}RemainingPrimesOfAdmitted}, lines 193--209)\par}

{\footnotesize\raggedright\noindent \ref{B4.13}.\quad \texttt{Formalism.\allowbreak{}MainNotions.\allowbreak{}WeilJacquetShalikaLaw.\allowbreak{}Proof.\allowbreak{}first\_\allowbreak{}clause\_\allowbreak{}cases} (\texttt{Formalism.\allowbreak{}SharpRestriction.\allowbreak{}GrowthBounds.\allowbreak{}AdmissibleIntegers.\allowbreak{}StatementsOfRemainingBound}, lines 25--96)\par}

{\footnotesize\raggedright\noindent \ref{def:p4c2-charts}.\quad \texttt{Formalism.\allowbreak{}MainNotions.\allowbreak{}G.\allowbreak{}CmpJ.\allowbreak{}CmpJHyp} (\texttt{Formalism.\allowbreak{}SharpRestriction.\allowbreak{}GrowthBounds.\allowbreak{}ChartComparison}, lines 82--96)\\
\texttt{Formalism.\allowbreak{}MainNotions.\allowbreak{}G.\allowbreak{}CmpJ.\allowbreak{}StepIn.\allowbreak{}Core} (\texttt{Formalism.\allowbreak{}SharpRestriction.\allowbreak{}GrowthBounds.\allowbreak{}ChartComparison.\allowbreak{}StepIn}, lines 139--174)\\
\texttt{Formalism.\allowbreak{}L2.\allowbreak{}Forms.\allowbreak{}ValuationCompatibleChart} (\texttt{Formalism.\allowbreak{}LocalToGlobal.\allowbreak{}HermitianForms.\allowbreak{}TiedTwGen}, lines 72--88)\\
\texttt{Formalism.\allowbreak{}MainNotions.\allowbreak{}G.\allowbreak{}CmpJ.\allowbreak{}StepIn.\allowbreak{}TwUnit} (\texttt{Formalism.\allowbreak{}SharpRestriction.\allowbreak{}GrowthBounds.\allowbreak{}ChartComparison.\allowbreak{}IntAtJIffN}, lines 43--46)\\
\texttt{Formalism.\allowbreak{}MainNotions.\allowbreak{}G.\allowbreak{}CmpJ.\allowbreak{}StepIn.\allowbreak{}ValIntN} (\texttt{Formalism.\allowbreak{}SharpRestriction.\allowbreak{}GrowthBounds.\allowbreak{}ChartComparison.\allowbreak{}IntAtJIffN}, lines 89--90)\\
\texttt{Formalism.\allowbreak{}MainNotions.\allowbreak{}G.\allowbreak{}CmpJ.\allowbreak{}StepIn.\allowbreak{}IwF} (\texttt{Formalism.\allowbreak{}SharpRestriction.\allowbreak{}GrowthBounds.\allowbreak{}ChartComparison.\allowbreak{}CmpJMemSpan}, lines 54--56)\par}

{\footnotesize\raggedright\noindent \ref{B4.13a}.\quad \texttt{GlobalTheory.\allowbreak{}Proof.\allowbreak{}ColumnWeightComparison.\allowbreak{}cmpV\_\allowbreak{}of\_\allowbreak{}cmpLite} (\texttt{Development.\allowbreak{}AbelianVarietyTheory.\allowbreak{}Proofs.\allowbreak{}SatakeParameters.\allowbreak{}ColumnWeightComparison}, lines 37--50)\par}

{\footnotesize\raggedright\noindent \ref{B4.14}.\quad \texttt{Formalism.\allowbreak{}MainNotions.\allowbreak{}WeilJacquetShalikaLaw.\allowbreak{}Ties9} (\texttt{Formalism.\allowbreak{}SharpRestriction.\allowbreak{}GrowthBounds.\allowbreak{}WeilJacquetShalikaLaw}, lines 92--95)\\
\texttt{Formalism.\allowbreak{}MainNotions.\allowbreak{}WeilJacquetShalikaLaw.\allowbreak{}Proof.\allowbreak{}statements\_\allowbreak{}of\_\allowbreak{}remaining\_\allowbreak{}bound} (\texttt{Formalism.\allowbreak{}SharpRestriction.\allowbreak{}GrowthBounds.\allowbreak{}AdmissibleIntegers.\allowbreak{}StatementsOfRemainingBound}, lines 194--200)\\
\texttt{Formalism.\allowbreak{}MainNotions.\allowbreak{}WeilJacquetShalikaLaw.\allowbreak{}Proof.\allowbreak{}baseChange\_\allowbreak{}preserves\_\allowbreak{}unitary\_\allowbreak{}of} (\texttt{Formalism.\allowbreak{}SharpRestriction.\allowbreak{}GrowthBounds.\allowbreak{}AdmissibleIntegers.\allowbreak{}StatementsOfRemainingBound}, lines 162--193)\par}

{\footnotesize\raggedright\noindent \ref{B4.15}.\quad \texttt{Formalism.\allowbreak{}MainNotions.\allowbreak{}Pk.\allowbreak{}admissible\_\allowbreak{}size\_\allowbreak{}family} (\texttt{Formalism.\allowbreak{}SharpRestriction.\allowbreak{}SmoothDesignExistence.\allowbreak{}SumOverPrimesHypotheses}, lines 57--80)\\
\texttt{Formalism.\allowbreak{}MainNotions.\allowbreak{}Pk.\allowbreak{}szEatI\_\allowbreak{}fam} (\texttt{Formalism.\allowbreak{}SharpRestriction.\allowbreak{}SmoothDesignExistence.\allowbreak{}SumOverPrimesHypotheses}, lines 38--48)\\
\texttt{UpperTheory.\allowbreak{}G2.\allowbreak{}szEatI\_\allowbreak{}of\_\allowbreak{}stepA} (\texttt{Development.\allowbreak{}AutomorphicArithmetic.\allowbreak{}OrbitLabellings.\allowbreak{}SzEatIFree}, lines 11--68)\\
\texttt{PI.\allowbreak{}SzEatI.\allowbreak{}W.\allowbreak{}stepA} (\texttt{Formalism.\allowbreak{}PeriodIntegrals.\allowbreak{}UnitaryGroupMassFormula.\allowbreak{}SizeAbsorption.\allowbreak{}SizeDel}, lines 121--344)\\
\texttt{UpperTheory.\allowbreak{}sum\_\allowbreak{}log\_\allowbreak{}mul\_\allowbreak{}max\_\allowbreak{}zero\_\allowbreak{}le\_\allowbreak{}of\_\allowbreak{}odd\_\allowbreak{}addOrderOf} (\texttt{Development.\allowbreak{}AutomorphicArithmetic.\allowbreak{}OrbitLabellings.\allowbreak{}SumLogMulMaxZeroLeOfOddAddOrderOf}, lines 101--170)\par}

{\footnotesize\raggedright\noindent \ref{B4.16}.\quad \texttt{Formalism.\allowbreak{}MainNotions.\allowbreak{}Pk.\allowbreak{}sum\_\allowbreak{}over\_\allowbreak{}primes\_\allowbreak{}service} (\texttt{Formalism.\allowbreak{}SharpRestriction.\allowbreak{}SmoothDesignExistence.\allowbreak{}SumOverPrimesHypotheses}, lines 187--222)\\
\texttt{Formalism.\allowbreak{}L2.\allowbreak{}G.\allowbreak{}ServT.\allowbreak{}closed\_\allowbreak{}servF\_\allowbreak{}of\_\allowbreak{}servT\_\allowbreak{}of\_\allowbreak{}dGamSize} (\texttt{Development.\allowbreak{}AutomorphicArithmetic.\allowbreak{}ThetaCorrespondence.\allowbreak{}ThetaPairLifts.\allowbreak{}ThirdPlaceSeries.\allowbreak{}GServT}, lines 31--33)\par}

\section*{Chapter~\ref{ch:columns-to-periods}: From columns to periods}

{\footnotesize\raggedright\noindent \ref{B4.17}.\quad \texttt{Formalism.\allowbreak{}MainNotions.\allowbreak{}WeightDegreeSizes.\allowbreak{}pair\_\allowbreak{}core} (\texttt{Formalism.\allowbreak{}SharpRestriction.\allowbreak{}Divisibility.\allowbreak{}WeightDegreeSizes}, lines 783--801)\par}

{\footnotesize\raggedright\noindent \ref{B4.18}.\quad \texttt{Formalism.\allowbreak{}L2.\allowbreak{}G.\allowbreak{}Adm2InU.\allowbreak{}pair0U} (\texttt{Formalism.\allowbreak{}LocalToGlobal.\allowbreak{}GroupLattices.\allowbreak{}AdmissibleAtGeneratorPlaces}, lines 39--54)\\
\texttt{Formalism.\allowbreak{}L2.\allowbreak{}G.\allowbreak{}Adm2InU.\allowbreak{}pairTwU} (\texttt{Formalism.\allowbreak{}LocalToGlobal.\allowbreak{}GroupLattices.\allowbreak{}AdmissibleAtGeneratorPlaces}, lines 55--71)\par}

{\footnotesize\raggedright\noindent \ref{B4.18a}.\quad \texttt{GlobalTheory.\allowbreak{}RowIntQ.\allowbreak{}rowInt\_\allowbreak{}over} (\texttt{Development.\allowbreak{}AbelianVarietyTheory.\allowbreak{}Proofs.\allowbreak{}GlobalInterpolation.\allowbreak{}PeriodRowLatticesIntegralProvisionalRows}, lines 47--64)\\
\texttt{GlobalTheory.\allowbreak{}PairAtClass.\allowbreak{}pair\_\allowbreak{}integral} (\texttt{Development.\allowbreak{}AbelianVarietyTheory.\allowbreak{}Proofs.\allowbreak{}GlobalInterpolation.\allowbreak{}ClassPairingIntegralFormRowIsIntegral}, lines 49--70)\\
\texttt{GlobalTheory.\allowbreak{}PerRed.\allowbreak{}admissiblePeriodZero\_\allowbreak{}of\_\allowbreak{}classPeriodDecomposition} (\texttt{Development.\allowbreak{}AbelianVarietyTheory.\allowbreak{}Proofs.\allowbreak{}GlobalInterpolation.\allowbreak{}AdmissiblePeriodZeroOfClassPeriodDecompositionAl}, lines 196--240)\\
\texttt{GlobalTheory.\allowbreak{}PerRedU.\allowbreak{}perTwAt\_\allowbreak{}of\_\allowbreak{}perDecTw} (\texttt{Development.\allowbreak{}AbelianVarietyTheory.\allowbreak{}Proofs.\allowbreak{}GlobalInterpolation.\allowbreak{}AdmissiblePeriodTwistsReductionPlaceTuples}, lines 12--40)\\
\texttt{GlobalTheory.\allowbreak{}Mem.\allowbreak{}Col2.\allowbreak{}TC.\allowbreak{}formsMap\_\allowbreak{}at\_\allowbreak{}every\_\allowbreak{}trueColumnMember} (\texttt{Development.\allowbreak{}AbelianVarietyTheory.\allowbreak{}Proofs.\allowbreak{}GlobalInterpolation.\allowbreak{}FormsMapAtEveryTrueColumnMember}, lines 15--27)\\
\texttt{GlobalTheory.\allowbreak{}PerQ.\allowbreak{}twistedPointIntegral\_\allowbreak{}of\_\allowbreak{}TC\_\allowbreak{}ok} (\texttt{Development.\allowbreak{}AbelianVarietyTheory.\allowbreak{}Proofs.\allowbreak{}GlobalInterpolation.\allowbreak{}TwistedPointRowAdmissibleTrueColumns}, lines 12--16)\\
\texttt{GlobalTheory.\allowbreak{}RedPt0Pend.\allowbreak{}exists\_\allowbreak{}integralPendingObject\_\allowbreak{}symmetricFormsLevel} (\texttt{Development.\allowbreak{}AbelianVarietyTheory.\allowbreak{}Proofs.\allowbreak{}GlobalInterpolation.\allowbreak{}ExistsIntegralPendingObjectSymmetricFormsLevel}, lines 291--333)\par}

{\footnotesize\raggedright\noindent \ref{B4.19a}.\quad \texttt{GlobalTheory.\allowbreak{}Mem.\allowbreak{}DivModel.\allowbreak{}ramEmb} (\texttt{Development.\allowbreak{}AbelianVarietyTheory.\allowbreak{}Definitions.\allowbreak{}AutomorphicMembership.\allowbreak{}ExistsDivStage}, lines 315--317)\\
\texttt{GlobalTheory.\allowbreak{}Mem.\allowbreak{}DivModel.\allowbreak{}exists\_\allowbreak{}divStage} (\texttt{Development.\allowbreak{}AbelianVarietyTheory.\allowbreak{}Definitions.\allowbreak{}AutomorphicMembership.\allowbreak{}ExistsDivStage}, lines 334--381)\\
\texttt{GlobalTheory.\allowbreak{}PerQ.\allowbreak{}Model.\allowbreak{}reRows} (\texttt{Development.\allowbreak{}AbelianVarietyTheory.\allowbreak{}Definitions.\allowbreak{}GlobalTwisting.\allowbreak{}ExistsRowValNeZero}, lines 110--122)\par}

{\footnotesize\raggedright\noindent \ref{B4.19}.\quad \texttt{Formalism.\allowbreak{}L2.\allowbreak{}G.\allowbreak{}Adm2InU.\allowbreak{}period\_\allowbreak{}geTU} (\texttt{Formalism.\allowbreak{}LocalToGlobal.\allowbreak{}GroupLattices.\allowbreak{}AdmissibleAtGeneratorPlaces}, lines 72--91)\par}

{\footnotesize\raggedright\noindent \ref{B4.20}.\quad \texttt{Formalism.\allowbreak{}MainNotions.\allowbreak{}J.\allowbreak{}sizeData\_\allowbreak{}at\_\allowbreak{}level\_\allowbreak{}primes} (\texttt{Formalism.\allowbreak{}SharpRestriction.\allowbreak{}AdmissibleParameterJoins.\allowbreak{}DIVfromGLs}, lines 25--28)\\
\texttt{Formalism.\allowbreak{}MainNotions.\allowbreak{}WeightDegreeSizes.\allowbreak{}sg\_\allowbreak{}atT} (\texttt{Formalism.\allowbreak{}SharpRestriction.\allowbreak{}Divisibility.\allowbreak{}WeightDegreeSizes}, lines 887--902)\par}

{\footnotesize\raggedright\noindent \ref{B4.21b}.\quad \texttt{Formalism.\allowbreak{}BasicNotions.\allowbreak{}Ranks} (\texttt{Formalism.\allowbreak{}LocalRingCardinality.\allowbreak{}Params}, lines 6--23)\\
\texttt{Formalism.\allowbreak{}MainNotions.\allowbreak{}SizeCond} (\texttt{Formalism.\allowbreak{}SharpRestriction.\allowbreak{}ElevenLe}, lines 54--55)\\
\texttt{Formalism.\allowbreak{}MainNotions.\allowbreak{}eventually\_\allowbreak{}level\_\allowbreak{}prime\_\allowbreak{}coprime\_\allowbreak{}stabiliser\_\allowbreak{}of\_\allowbreak{}design} (\texttt{Formalism.\allowbreak{}SharpRestriction.\allowbreak{}GlobalSchurBounds.\allowbreak{}PlugDgt}, lines 10--101)\par}

{\footnotesize\raggedright\noindent \ref{def:p4c3-design}.\quad \texttt{Formalism.\allowbreak{}MainNotions.\allowbreak{}Des} (\texttt{Formalism.\allowbreak{}SharpRestriction.\allowbreak{}SmoothDesignExistence.\allowbreak{}DesignExistsDes}, lines 10--12)\\
\texttt{Formalism.\allowbreak{}MainNotions.\allowbreak{}ChosenE} (\texttt{Formalism.\allowbreak{}SharpRestriction.\allowbreak{}GlobalSchurBounds.\allowbreak{}Chosen}, lines 83--85)\\
\texttt{Formalism.\allowbreak{}MainNotions.\allowbreak{}Pk.\allowbreak{}DGamSmooth} (\texttt{Formalism.\allowbreak{}SharpRestriction.\allowbreak{}SmoothDesignExistence.\allowbreak{}DGamSmooth}, lines 23--26)\\
\texttt{Formalism.\allowbreak{}MainNotions.\allowbreak{}TuSize} (\texttt{Formalism.\allowbreak{}SharpRestriction.\allowbreak{}ParamsHeightRankConstant}, lines 180--183)\\
\texttt{Formalism.\allowbreak{}MainNotions.\allowbreak{}DGamSize} (\texttt{Formalism.\allowbreak{}SharpRestriction.\allowbreak{}ParamsHeightRankConstant}, lines 184--186)\\
\texttt{Formalism.\allowbreak{}BasicNotions.\allowbreak{}FactDGam} (\texttt{Formalism.\allowbreak{}LocalRingCardinality.\allowbreak{}RhoDetPrimeEqMuTildeMulChiPow}, lines 381--382)\\
\texttt{GlobalTheory.\allowbreak{}Cp.\allowbreak{}WithChN} (\texttt{Development.\allowbreak{}AbelianVarietyTheory.\allowbreak{}Definitions.\allowbreak{}GlobalTwisting.\allowbreak{}CountsIrreducibilityConstants}, lines 32--34)\\
\texttt{Formalism.\allowbreak{}MainNotions.\allowbreak{}TameDistinct} (\texttt{Formalism.\allowbreak{}SharpRestriction.\allowbreak{}LevelLetterCocycles.\allowbreak{}AlignN}, lines 14--20)\par}

{\footnotesize\raggedright\noindent \ref{B4.21}.\quad \texttt{Formalism.\allowbreak{}MainNotions.\allowbreak{}eventually\_\allowbreak{}level\_\allowbreak{}prime\_\allowbreak{}coprime\_\allowbreak{}stabiliser\_\allowbreak{}of\_\allowbreak{}design} (\texttt{Formalism.\allowbreak{}SharpRestriction.\allowbreak{}GlobalSchurBounds.\allowbreak{}PlugDgt}, lines 10--101)\\
\texttt{Formalism.\allowbreak{}MainNotions.\allowbreak{}closed\_\allowbreak{}eventually\_\allowbreak{}levelPrimeCoprimeStabiliser} (\texttt{Formalism.\allowbreak{}SharpRestriction.\allowbreak{}GlobalSchurBounds.\allowbreak{}PlugDgt}, lines 102--107)\par}

{\footnotesize\raggedright\noindent \ref{B4.21a}.\quad \texttt{GlobalTheory.\allowbreak{}EX.\allowbreak{}Led.\allowbreak{}exists\_\allowbreak{}lat} (\texttt{Development.\allowbreak{}AbelianVarietyTheory.\allowbreak{}Proofs.\allowbreak{}GlobalInterpolation.\allowbreak{}ExistsLedPart}, lines 119--186)\\
\texttt{Formalism.\allowbreak{}BasicNotions.\allowbreak{}LatCmp.\allowbreak{}NuServedDef} (\texttt{Formalism.\allowbreak{}LocalRingCardinality.\allowbreak{}DataLatDefs}, lines 42--46)\\
\texttt{Formalism.\allowbreak{}BasicNotions.\allowbreak{}NuAtTDef} (\texttt{Formalism.\allowbreak{}LocalRingCardinality.\allowbreak{}DataLatDefs}, lines 111--114)\\
\texttt{Formalism.\allowbreak{}MainNotions.\allowbreak{}WeightDegreeSizes.\allowbreak{}NuInM} (\texttt{Formalism.\allowbreak{}SharpRestriction.\allowbreak{}Divisibility.\allowbreak{}WeightDegreeSizes}, lines 852--854)\par}

{\footnotesize\raggedright\noindent \ref{B4.22}.\quad \texttt{Formalism.\allowbreak{}MainNotions.\allowbreak{}J.\allowbreak{}ColGe0\_\allowbreak{}atT} (\texttt{Formalism.\allowbreak{}SharpRestriction.\allowbreak{}AdmissibleParameterJoins.\allowbreak{}DIVfromGLi}, lines 34--68)\\
\texttt{Formalism.\allowbreak{}L2.\allowbreak{}G.\allowbreak{}Adm2InU.\allowbreak{}admissible\_\allowbreak{}at\_\allowbreak{}twist\_\allowbreak{}places} (\texttt{Formalism.\allowbreak{}LocalToGlobal.\allowbreak{}GroupLattices.\allowbreak{}AdmissibleAtGeneratorPlaces}, lines 92--112)\par}

{\footnotesize\raggedright\noindent \ref{B4.22d}.\quad \texttt{GlobalTheory.\allowbreak{}MomCrit.\allowbreak{}momInt\_\allowbreak{}iff\_\allowbreak{}norm\_\allowbreak{}mahler\_\allowbreak{}le} (\texttt{Development.\allowbreak{}AbelianVarietyTheory.\allowbreak{}Proofs.\allowbreak{}GlobalInterpolation.\allowbreak{}MahlerCoefficientNormLeOneIffJetsBoundedMeasureM}, lines 217--266)\\
\texttt{GlobalTheory.\allowbreak{}MomT.\allowbreak{}Riemann.\allowbreak{}riemann\_\allowbreak{}sums} (\texttt{Development.\allowbreak{}AbelianVarietyTheory.\allowbreak{}Proofs.\allowbreak{}GlobalInterpolation.\allowbreak{}IsLocallyConstantToZModPow}, lines 55--103)\par}

{\footnotesize\raggedright\noindent \ref{B4.22e}.\quad \texttt{GlobalTheory.\allowbreak{}CmThetaMoment.\allowbreak{}isomStd} (\texttt{Development.\allowbreak{}AbelianVarietyTheory.\allowbreak{}Proofs.\allowbreak{}GlobalInterpolation.\allowbreak{}ExistsStandardIsometryFourthPowerSplitPrime}, lines 255--272)\\
\texttt{Formalism.\allowbreak{}L2.\allowbreak{}Ker.\allowbreak{}Cd.\allowbreak{}PairingIndexBounds.\allowbreak{}Eq.\allowbreak{}J.\allowbreak{}first\_\allowbreak{}constant\_\allowbreak{}of\_\allowbreak{}pullback} (\texttt{Formalism.\allowbreak{}LocalToGlobal.\allowbreak{}KernelLattices.\allowbreak{}ColumnData.\allowbreak{}ColumnDataEquations.\allowbreak{}JFinal}, lines 23--25)\\
\texttt{Formalism.\allowbreak{}L2.\allowbreak{}Ker.\allowbreak{}Cd.\allowbreak{}PairingIndexBounds.\allowbreak{}Eq.\allowbreak{}J.\allowbreak{}second\_\allowbreak{}constant\_\allowbreak{}of\_\allowbreak{}pullback} (\texttt{Formalism.\allowbreak{}LocalToGlobal.\allowbreak{}KernelLattices.\allowbreak{}ColumnData.\allowbreak{}ColumnDataEquations.\allowbreak{}JFinal}, lines 26--28)\par}

{\footnotesize\raggedright\noindent \ref{B4.22c}.\quad \texttt{GlobalTheory.\allowbreak{}MomN.\allowbreak{}transPrime\_\allowbreak{}int} (\texttt{Development.\allowbreak{}AbelianVarietyTheory.\allowbreak{}Proofs.\allowbreak{}GlobalInterpolation.\allowbreak{}NormalizedVectorTranslatesPrimeOrderIntegral}, lines 32--83)\\
\texttt{GlobalTheory.\allowbreak{}MomCrit.\allowbreak{}momInt\_\allowbreak{}iff\_\allowbreak{}norm\_\allowbreak{}mahler\_\allowbreak{}le} (\texttt{Development.\allowbreak{}AbelianVarietyTheory.\allowbreak{}Proofs.\allowbreak{}GlobalInterpolation.\allowbreak{}MahlerCoefficientNormLeOneIffJetsBoundedMeasureM}, lines 217--266)\par}

{\footnotesize\raggedright\noindent \ref{B4.22f}.\quad \texttt{GlobalTheory.\allowbreak{}MomOne.\allowbreak{}Norm.\allowbreak{}exists\_\allowbreak{}normalised\_\allowbreak{}chart} (\texttt{Development.\allowbreak{}AbelianVarietyTheory.\allowbreak{}Proofs.\allowbreak{}JetHarmonics.\allowbreak{}ExistsNormalisedMultiplicativeChart}, lines 200--277)\\
\texttt{GlobalTheory.\allowbreak{}MomOne.\allowbreak{}ThShift.\allowbreak{}theta\_\allowbreak{}shift} (\texttt{Development.\allowbreak{}AbelianVarietyTheory.\allowbreak{}Proofs.\allowbreak{}JetHarmonics.\allowbreak{}CmLatticeWeylOperatorEqThetaRootOfUnityShift}, lines 158--281)\\
\texttt{GlobalTheory.\allowbreak{}MomOne.\allowbreak{}TranslateLaw.\allowbreak{}main} (\texttt{Development.\allowbreak{}AbelianVarietyTheory.\allowbreak{}Proofs.\allowbreak{}JetHarmonics.\allowbreak{}WeilOperatorThetaTranslateShiftPrimeToQ}, lines 61--144)\par}

{\footnotesize\raggedright\noindent \ref{B4.22g}.\quad \texttt{GlobalTheory.\allowbreak{}MomFin.\allowbreak{}moments\_\allowbreak{}fin} (\texttt{Development.\allowbreak{}AbelianVarietyTheory.\allowbreak{}Proofs.\allowbreak{}GlobalInterpolation.\allowbreak{}UnitPeriodTranslatesMomentsFinite}, lines 77--177)\\
\texttt{GlobalTheory.\allowbreak{}MomT.\allowbreak{}PtMoments.\allowbreak{}ptMom} (\texttt{Development.\allowbreak{}AbelianVarietyTheory.\allowbreak{}Proofs.\allowbreak{}GlobalInterpolation.\allowbreak{}ContinuousToZModComp}, lines 342--444)\par}

{\footnotesize\raggedright\noindent \ref{B4.22h}.\quad \texttt{GlobalTheory.\allowbreak{}MomT.\allowbreak{}IdxBound.\allowbreak{}relIndex\_\allowbreak{}per\_\allowbreak{}dvd} (\texttt{Development.\allowbreak{}AbelianVarietyTheory.\allowbreak{}Proofs.\allowbreak{}GlobalInterpolation.\allowbreak{}RelIndexPerDvd}, lines 91--118)\\
\texttt{GlobalTheory.\allowbreak{}MomT.\allowbreak{}CosetSum.\allowbreak{}moments\_\allowbreak{}coset\_\allowbreak{}sum} (\texttt{Development.\allowbreak{}AbelianVarietyTheory.\allowbreak{}Proofs.\allowbreak{}GlobalInterpolation.\allowbreak{}MomentsCosetSum}, lines 363--492)\par}

{\footnotesize\raggedright\noindent \ref{B4.22b}.\quad \texttt{GlobalTheory.\allowbreak{}KerT.\allowbreak{}PStd.\allowbreak{}pTerm\_\allowbreak{}std} (\texttt{Development.\allowbreak{}AbelianVarietyTheory.\allowbreak{}Proofs.\allowbreak{}GlobalInterpolation.\allowbreak{}ExistsMulEqOfUnif}, lines 340--424)\\
\texttt{GlobalTheory.\allowbreak{}BarStd.\allowbreak{}barStd} (\texttt{Development.\allowbreak{}AbelianVarietyTheory.\allowbreak{}Proofs.\allowbreak{}GlobalInterpolation.\allowbreak{}AdmissibleAdjointRationalPointPadicMatrix}, lines 388--458)\\
\texttt{GlobalTheory.\allowbreak{}LvT.\allowbreak{}Built.\allowbreak{}UpLaw.\allowbreak{}kAt\_\allowbreak{}uk} (\texttt{Development.\allowbreak{}AbelianVarietyTheory.\allowbreak{}Proofs.\allowbreak{}GlobalInterpolation.\allowbreak{}SameOfAbove}, lines 61--78)\par}

{\footnotesize\raggedright\noindent \ref{B4.22a}.\quad \texttt{GlobalTheory.\allowbreak{}LvT.\allowbreak{}Built.\allowbreak{}TrueLaws.\allowbreak{}laws\_\allowbreak{}true\_\allowbreak{}fam} (\texttt{Development.\allowbreak{}AbelianVarietyTheory.\allowbreak{}Proofs.\allowbreak{}GlobalInterpolation.\allowbreak{}FivePointLawsTrueChains}, lines 46--63)\\
\texttt{GlobalTheory.\allowbreak{}LvT.\allowbreak{}UKS.\allowbreak{}ukSharp} (\texttt{Development.\allowbreak{}AbelianVarietyTheory.\allowbreak{}Proofs.\allowbreak{}GlobalInterpolation.\allowbreak{}ExtInOfValIn}, lines 155--199)\\
\texttt{GlobalTheory.\allowbreak{}LvT.\allowbreak{}levelChainSharpColumnValuation\_\allowbreak{}of\_\allowbreak{}stability} (\texttt{Development.\allowbreak{}AbelianVarietyTheory.\allowbreak{}Proofs.\allowbreak{}GlobalInterpolation.\allowbreak{}LevelChainSharpColumnValuationOfStability}, lines 174--194)\\
\texttt{GlobalTheory.\allowbreak{}LvT.\allowbreak{}Built.\allowbreak{}TrueCols.\allowbreak{}row\_\allowbreak{}true} (\texttt{Development.\allowbreak{}AbelianVarietyTheory.\allowbreak{}Proofs.\allowbreak{}GlobalInterpolation.\allowbreak{}TrueColumnsRowSharpBounds}, lines 36--51)\\
\texttt{GlobalTheory.\allowbreak{}LvT.\allowbreak{}StabOrder.\allowbreak{}stabUnit\_\allowbreak{}of\_\allowbreak{}embed} (\texttt{Development.\allowbreak{}AbelianVarietyTheory.\allowbreak{}Proofs.\allowbreak{}GlobalInterpolation.\allowbreak{}NotDvdCardStabilizerForm}, lines 157--164)\par}

{\footnotesize\raggedright\noindent \ref{B4.23}.\quad \texttt{Formalism.\allowbreak{}MainNotions.\allowbreak{}WeightDegreeSizes.\allowbreak{}sg\_\allowbreak{}atG} (\texttt{Formalism.\allowbreak{}SharpRestriction.\allowbreak{}Divisibility.\allowbreak{}WeightDegreeSizes}, lines 1007--1017)\\
\texttt{Formalism.\allowbreak{}L2.\allowbreak{}G.\allowbreak{}Adm2InU.\allowbreak{}atG\_\allowbreak{}coreU} (\texttt{Formalism.\allowbreak{}LocalToGlobal.\allowbreak{}GroupLattices.\allowbreak{}AdmissibleAtGeneratorPlaces}, lines 113--133)\par}

{\footnotesize\raggedright\noindent \ref{B4.23a}.\quad \texttt{GlobalTheory.\allowbreak{}JetsFromTables.\allowbreak{}jetLossAll\_\allowbreak{}of\_\allowbreak{}jetSplitAmb} (\texttt{Development.\allowbreak{}AbelianVarietyTheory.\allowbreak{}Proofs.\allowbreak{}SatakeParameters.\allowbreak{}ThetaJetIntegralityTypePrimeOfSplitTableBound}, lines 221--293)\\
\texttt{GlobalTheory.\allowbreak{}ColPeriodTies.\allowbreak{}boundAt\_\allowbreak{}of\_\allowbreak{}chainTie} (\texttt{Development.\allowbreak{}AbelianVarietyTheory.\allowbreak{}Proofs.\allowbreak{}SatakeParameters.\allowbreak{}PeriodValuationBoundTypePrimePlaceOfChainAgreeme}, lines 19--95)\\
\texttt{GlobalTheory.\allowbreak{}ColPeriod.\allowbreak{}vG\_\allowbreak{}loss\_\allowbreak{}le} (\texttt{Development.\allowbreak{}AbelianVarietyTheory.\allowbreak{}Proofs.\allowbreak{}SatakeParameters.\allowbreak{}PeriodBoundOfJetLoss}, lines 87--142)\\
\texttt{Formalism.\allowbreak{}MainNotions.\allowbreak{}J.\allowbreak{}sumFE\_\allowbreak{}atG\_\allowbreak{}le} (\texttt{Formalism.\allowbreak{}SharpRestriction.\allowbreak{}AdmissibleParameterJoins.\allowbreak{}GandSx}, lines 12--85)\\
\texttt{GlobalTheory.\allowbreak{}AtG.\allowbreak{}row\_\allowbreak{}atG\_\allowbreak{}of\_\allowbreak{}jetSplitAmb} (\texttt{Development.\allowbreak{}AbelianVarietyTheory.\allowbreak{}Proofs.\allowbreak{}SatakeParameters.\allowbreak{}TypePrimePeriodRowOfConstructedPointChains}, lines 15--29)\\
\texttt{GlobalTheory.\allowbreak{}AtG.\allowbreak{}jetSplitAmb\_\allowbreak{}holds} (\texttt{Development.\allowbreak{}AbelianVarietyTheory.\allowbreak{}Proofs.\allowbreak{}SatakeParameters.\allowbreak{}JetTableBoundSplitPrime}, lines 7--9)\\
\texttt{GlobalTheory.\allowbreak{}JetsFromTables.\allowbreak{}datum\_\allowbreak{}bound} (\texttt{Development.\allowbreak{}AbelianVarietyTheory.\allowbreak{}Proofs.\allowbreak{}SatakeParameters.\allowbreak{}ThetaJetIntegralityTypePrimeOfSplitTableBound}, lines 108--165)\par}

{\footnotesize\raggedright\noindent \ref{B4.24}.\quad \texttt{Formalism.\allowbreak{}L2.\allowbreak{}G.\allowbreak{}Adm2InU.\allowbreak{}admissible\_\allowbreak{}at\_\allowbreak{}generator\_\allowbreak{}places} (\texttt{Formalism.\allowbreak{}LocalToGlobal.\allowbreak{}GroupLattices.\allowbreak{}AdmissibleAtGeneratorPlaces}, lines 134--170)\par}

{\footnotesize\raggedright\noindent \ref{B4.25}.\quad \texttt{Formalism.\allowbreak{}MainNotions.\allowbreak{}WeilJacquetShalikaLaw.\allowbreak{}gainD\_\allowbreak{}le\_\allowbreak{}gainS} (\texttt{Formalism.\allowbreak{}SharpRestriction.\allowbreak{}GrowthBounds.\allowbreak{}WeilJacquetShalikaLaw}, lines 552--576)\\
\texttt{Formalism.\allowbreak{}MainNotions.\allowbreak{}WeightDegreeSizes.\allowbreak{}gainLe\_\allowbreak{}of\_\allowbreak{}fills} (\texttt{Formalism.\allowbreak{}SharpRestriction.\allowbreak{}Divisibility.\allowbreak{}WeightDegreeSizes}, lines 931--934)\par}

{\footnotesize\raggedright\noindent \ref{B4.26}.\quad \texttt{Formalism.\allowbreak{}L2.\allowbreak{}G.\allowbreak{}Adm2InU.\allowbreak{}admissible\_\allowbreak{}off\_\allowbreak{}twist\_\allowbreak{}places} (\texttt{Formalism.\allowbreak{}LocalToGlobal.\allowbreak{}GroupLattices.\allowbreak{}AdmissibleAtGeneratorPlaces}, lines 171--197)\par}

{\footnotesize\raggedright\noindent \ref{B4.27}.\quad \texttt{Formalism.\allowbreak{}L2.\allowbreak{}G.\allowbreak{}Adm2InU.\allowbreak{}admissible\_\allowbreak{}away\_\allowbreak{}from\_\allowbreak{}twist} (\texttt{Formalism.\allowbreak{}LocalToGlobal.\allowbreak{}GroupLattices.\allowbreak{}AdmissibleAtGeneratorPlaces}, lines 198--213)\par}

{\footnotesize\raggedright\noindent \ref{B4.28}.\quad \texttt{Formalism.\allowbreak{}MainNotions.\allowbreak{}LocalQuadraticParts.\allowbreak{}UnroundedValuationBound} (\texttt{Formalism.\allowbreak{}SharpRestriction.\allowbreak{}Divisibility.\allowbreak{}LocalQuadraticParts}, lines 212--217)\\
\texttt{Formalism.\allowbreak{}MainNotions.\allowbreak{}LocalQuadraticParts.\allowbreak{}adm2} (\texttt{Formalism.\allowbreak{}SharpRestriction.\allowbreak{}Divisibility.\allowbreak{}LocalQuadraticParts}, lines 206--211)\\
\texttt{Formalism.\allowbreak{}MainNotions.\allowbreak{}LocalQuadraticParts.\allowbreak{}valuation\_\allowbreak{}bound\_\allowbreak{}of\_\allowbreak{}unrounded} (\texttt{Formalism.\allowbreak{}SharpRestriction.\allowbreak{}Divisibility.\allowbreak{}LocalQuadraticParts}, lines 218--222)\par}

{\footnotesize\raggedright\noindent \ref{B4.29}.\quad \texttt{Formalism.\allowbreak{}L2.\allowbreak{}G.\allowbreak{}Adm2InU.\allowbreak{}admissible\_\allowbreak{}level\_\allowbreak{}two} (\texttt{Formalism.\allowbreak{}LocalToGlobal.\allowbreak{}GroupLattices.\allowbreak{}AdmissibleAtGeneratorPlaces}, lines 214--246)\\
\texttt{Formalism.\allowbreak{}MainNotions.\allowbreak{}LocalQuadraticParts.\allowbreak{}adm12} (\texttt{Formalism.\allowbreak{}SharpRestriction.\allowbreak{}Divisibility.\allowbreak{}LocalQuadraticParts}, lines 323--329)\par}

{\footnotesize\raggedright\noindent \ref{B4.30}.\quad \texttt{Formalism.\allowbreak{}L2.\allowbreak{}G.\allowbreak{}Adm2InU.\allowbreak{}famU} (\texttt{Formalism.\allowbreak{}LocalToGlobal.\allowbreak{}GroupLattices.\allowbreak{}AdmissibleAtGeneratorPlaces}, lines 263--301)\\
\texttt{Formalism.\allowbreak{}L2.\allowbreak{}G.\allowbreak{}Adm2InU.\allowbreak{}admissible\_\allowbreak{}family\_\allowbreak{}of\_\allowbreak{}inclusion} (\texttt{Formalism.\allowbreak{}LocalToGlobal.\allowbreak{}GroupLattices.\allowbreak{}AdmissibleAtGeneratorPlaces}, lines 315--339)\par}

{\footnotesize\raggedright\noindent \ref{B4.31}.\quad \texttt{GlobalTheory.\allowbreak{}Top.\allowbreak{}closed\_\allowbreak{}family\_\allowbreak{}admissible\_\allowbreak{}twist\_\allowbreak{}link\_\allowbreak{}alt} (\texttt{Development.\allowbreak{}AutomorphicArithmetic.\allowbreak{}ThetaCorrespondence.\allowbreak{}ThetaPairLifts.\allowbreak{}ThirdPlaceSeries.\allowbreak{}ClosedFamilyAdmissibleTwistLink}, lines 12--31)\par}

\section*{Chapter~\ref{ch:hecke-row-gluing}: Gluing row lattices across Hecke characters}

{\footnotesize\raggedright\noindent \ref{B4.32}.\quad \texttt{Formalism.\allowbreak{}L2.\allowbreak{}DIV.\allowbreak{}finChar} (\texttt{Formalism.\allowbreak{}LocalToGlobal.\allowbreak{}Divisibility.\allowbreak{}Tuples}, lines 107--109)\\
\texttt{Formalism.\allowbreak{}L2.\allowbreak{}DIV.\allowbreak{}tupGrp} (\texttt{Formalism.\allowbreak{}LocalToGlobal.\allowbreak{}Divisibility.\allowbreak{}Tuples}, lines 112--115)\\
\texttt{Formalism.\allowbreak{}L2.\allowbreak{}DIV.\allowbreak{}Tup} (\texttt{Formalism.\allowbreak{}LocalToGlobal.\allowbreak{}Divisibility.\allowbreak{}Tuples}, lines 118--119)\par}

{\footnotesize\raggedright\noindent \ref{B4.33}.\quad \texttt{Formalism.\allowbreak{}L2.\allowbreak{}DIV.\allowbreak{}TupF} (\texttt{Formalism.\allowbreak{}LocalToGlobal.\allowbreak{}Divisibility.\allowbreak{}LiftLatInst}, lines 58--59)\\
\texttt{Formalism.\allowbreak{}L2.\allowbreak{}DIV.\allowbreak{}colIn3} (\texttt{Formalism.\allowbreak{}LocalToGlobal.\allowbreak{}Divisibility.\allowbreak{}HeckeCSDICommTrRep}, lines 130--141)\\
\texttt{Formalism.\allowbreak{}L2.\allowbreak{}DIV.\allowbreak{}DivisionKernelData} (\texttt{Development.\allowbreak{}AutomorphicArithmetic.\allowbreak{}ThetaCorrespondence.\allowbreak{}ThetaPairLifts.\allowbreak{}ThirdPlaceSeries.\allowbreak{}DIVRecBridgeSel}, lines 32--57)\\
\texttt{GlobalTheory.\allowbreak{}Mem.\allowbreak{}Col2.\allowbreak{}Reads} (\texttt{Development.\allowbreak{}AbelianVarietyTheory.\allowbreak{}Definitions.\allowbreak{}AutomorphicMembership.\allowbreak{}TrueColumnChoicesHeckeCharReferenceEqOne}, lines 48--49)\\
\texttt{GlobalTheory.\allowbreak{}Mem.\allowbreak{}Col2.\allowbreak{}ColIn.\allowbreak{}Ok.\allowbreak{}xi\_\allowbreak{}reads} (\texttt{Development.\allowbreak{}AbelianVarietyTheory.\allowbreak{}Definitions.\allowbreak{}AutomorphicMembership.\allowbreak{}TrueColumnChoicesHeckeCharReferenceEqOne}, lines 237--237)\\
\texttt{GlobalTheory.\allowbreak{}Mem.\allowbreak{}R3.\allowbreak{}tup\_\allowbreak{}ixOf} (\texttt{Development.\allowbreak{}AbelianVarietyTheory.\allowbreak{}Definitions.\allowbreak{}AutomorphicMembership.\allowbreak{}IsAbsIrrCharCompIff}, lines 105--112)\par}

{\footnotesize\raggedright\noindent \ref{B4.34}.\quad \texttt{Formalism.\allowbreak{}L2.\allowbreak{}Indef.\allowbreak{}DefSide.\allowbreak{}abelian₂} (\texttt{Formalism.\allowbreak{}LocalToGlobal.\allowbreak{}IndefiniteForms.\allowbreak{}IndefSide}, lines 68--69)\\
\texttt{Formalism.\allowbreak{}L2.\allowbreak{}Indef.\allowbreak{}IndefData.\allowbreak{}RR} (\texttt{Formalism.\allowbreak{}LocalToGlobal.\allowbreak{}IndefiniteForms.\allowbreak{}IndefSide}, lines 281--283)\\
\texttt{Formalism.\allowbreak{}DivisionValues.\allowbreak{}glueLat} (\texttt{Formalism.\allowbreak{}QuotientMetricDivisors.\allowbreak{}LatticeGlueUnique}, lines 11--14)\par}

{\footnotesize\raggedright\noindent \ref{B4.35}.\quad \texttt{Formalism.\allowbreak{}DivisionValues.\allowbreak{}lattice\_\allowbreak{}gluing} (\texttt{Formalism.\allowbreak{}QuotientMetricDivisors.\allowbreak{}HermitianTwistDivisors.\allowbreak{}SecondStepGluing}, lines 256--259)\\
\texttt{Formalism.\allowbreak{}DivisionValues.\allowbreak{}GluingInclusion.\allowbreak{}glue\_\allowbreak{}spec} (\texttt{Formalism.\allowbreak{}QuotientMetricDivisors.\allowbreak{}HermitianTwistDivisors.\allowbreak{}SecondStepGluing}, lines 74--183)\par}

{\footnotesize\raggedright\noindent \ref{B4.36}.\quad \texttt{Formalism.\allowbreak{}L2.\allowbreak{}Indef.\allowbreak{}IndefRows.\allowbreak{}RR\_\allowbreak{}spec} (\texttt{Formalism.\allowbreak{}LocalToGlobal.\allowbreak{}IndefiniteForms.\allowbreak{}IndefUse}, lines 219--230)\\
\texttt{Formalism.\allowbreak{}L2.\allowbreak{}Indef.\allowbreak{}IndefColsR.\allowbreak{}RR\_\allowbreak{}spec} (\texttt{Formalism.\allowbreak{}LocalToGlobal.\allowbreak{}IndefiniteForms.\allowbreak{}PairingConditionOfInverseMatrixIdentity}, lines 332--343)\par}

{\footnotesize\raggedright\noindent \ref{B4.36a}.\quad \texttt{GlobalTheory.\allowbreak{}Mem.\allowbreak{}Top.\allowbreak{}TwFin.\allowbreak{}tw\_\allowbreak{}fin} (\texttt{Development.\allowbreak{}AbelianVarietyTheory.\allowbreak{}Proofs.\allowbreak{}IwasawaDecomposition.\allowbreak{}ResCharConj}, lines 58--79)\\
\texttt{GlobalTheory.\allowbreak{}Mem.\allowbreak{}R3.\allowbreak{}tupL\_\allowbreak{}of\_\allowbreak{}twisted} (\texttt{Development.\allowbreak{}AbelianVarietyTheory.\allowbreak{}Proofs.\allowbreak{}IwasawaDecomposition.\allowbreak{}MemHeightOneSpectrumIffAdditiveValuationPos}, lines 82--189)\\
\texttt{Formalism.\allowbreak{}IndefR.\allowbreak{}Links.\allowbreak{}RealDat3.\allowbreak{}height\_\allowbreak{}change\_\allowbreak{}of\_\allowbreak{}D0} (\texttt{Formalism.\allowbreak{}LocalToGlobal.\allowbreak{}IndefiniteReal.\allowbreak{}BaseLiftsPackageOfPrimitiveLinks}, lines 105--109)\\
\texttt{Formalism.\allowbreak{}L2.\allowbreak{}Indef.\allowbreak{}IndefData.\allowbreak{}H0} (\texttt{Formalism.\allowbreak{}LocalToGlobal.\allowbreak{}IndefiniteForms.\allowbreak{}IndefSide}, lines 294--296)\\
\texttt{GlobalTheory.\allowbreak{}Mem.\allowbreak{}Col2.\allowbreak{}InEx.\allowbreak{}lg\_\allowbreak{}colIn} (\texttt{Development.\allowbreak{}AbelianVarietyTheory.\allowbreak{}Proofs.\allowbreak{}IwasawaDecomposition.\allowbreak{}ExistsColumnInputAssembled}, lines 328--373)\par}

{\footnotesize\raggedright\noindent \ref{B4.37}.\quad \texttt{Formalism.\allowbreak{}L2.\allowbreak{}G.\allowbreak{}ServJoint.\allowbreak{}CutUp} (\texttt{Formalism.\allowbreak{}LocalToGlobal.\allowbreak{}GroupLattices.\allowbreak{}JointServices}, lines 65--71)\\
\texttt{Formalism.\allowbreak{}L2.\allowbreak{}G.\allowbreak{}ServJoint.\allowbreak{}cutUp\_\allowbreak{}of\_\allowbreak{}stable} (\texttt{Formalism.\allowbreak{}LocalToGlobal.\allowbreak{}GroupLattices.\allowbreak{}JointServices.\allowbreak{}Cut}, lines 33--46)\par}

{\footnotesize\raggedright\noindent \ref{B4.38}.\quad \texttt{Formalism.\allowbreak{}L2.\allowbreak{}G.\allowbreak{}ServJointF.\allowbreak{}ColBrF.\allowbreak{}Tied} (\texttt{Formalism.\allowbreak{}LocalToGlobal.\allowbreak{}GroupLattices.\allowbreak{}ServJointF}, lines 30--40)\\
\texttt{Formalism.\allowbreak{}L2.\allowbreak{}G.\allowbreak{}ServJointF.\allowbreak{}perToVal2F} (\texttt{Formalism.\allowbreak{}LocalToGlobal.\allowbreak{}GroupLattices.\allowbreak{}ServJointF}, lines 42--83)\\
\texttt{Formalism.\allowbreak{}L2.\allowbreak{}G.\allowbreak{}ServJointF.\allowbreak{}val2AllAtRF} (\texttt{Formalism.\allowbreak{}LocalToGlobal.\allowbreak{}GroupLattices.\allowbreak{}ServJointF}, lines 84--103)\\
\texttt{Formalism.\allowbreak{}L2.\allowbreak{}G.\allowbreak{}ServJoint.\allowbreak{}vAdd\_\allowbreak{}eq\_\allowbreak{}vw} (\texttt{Formalism.\allowbreak{}LocalToGlobal.\allowbreak{}GroupLattices.\allowbreak{}JointServices.\allowbreak{}VAddEqVw}, lines 37--52)\\
\texttt{Formalism.\allowbreak{}L2.\allowbreak{}G.\allowbreak{}ServJoint.\allowbreak{}perCol\_\allowbreak{}lin} (\texttt{Formalism.\allowbreak{}LocalToGlobal.\allowbreak{}GroupLattices.\allowbreak{}JointServices.\allowbreak{}ValIntNonneg}, lines 12--23)\par}

\section*{Chapter~\ref{ch:pairing-after-division}: The pairing after division}

{\footnotesize\raggedright\noindent \ref{B4.39}.\quad \texttt{Formalism.\allowbreak{}BasicNotions.\allowbreak{}DV.\allowbreak{}Iu} (\texttt{Formalism.\allowbreak{}LocalRingCardinality.\allowbreak{}DataDIVDefs}, lines 17--20)\\
\texttt{Formalism.\allowbreak{}BasicNotions.\allowbreak{}DV.\allowbreak{}NprimeDef} (\texttt{Formalism.\allowbreak{}LocalRingCardinality.\allowbreak{}DataDIVDefs}, lines 22--24)\par}

{\footnotesize\raggedright\noindent \ref{B4.40}.\quad \texttt{Formalism.\allowbreak{}DivisionValues.\allowbreak{}localizationAt\_\allowbreak{}fractionalIdeal\_\allowbreak{}smul\_\allowbreak{}eq\_\allowbreak{}smul} (\texttt{Formalism.\allowbreak{}QuotientMetricDivisors.\allowbreak{}HermitianTwistDivisors.\allowbreak{}TwistedLatticeNaturality}, lines 547--572)\\
\texttt{Formalism.\allowbreak{}DivisionValues.\allowbreak{}shift\_\allowbreak{}existence} (\texttt{Formalism.\allowbreak{}QuotientMetricDivisors.\allowbreak{}HermitianTwistDivisors.\allowbreak{}TwistedLatticeNaturality}, lines 649--662)\par}

{\footnotesize\raggedright\noindent \ref{B4.40}.\quad \texttt{Formalism.\allowbreak{}DivisionValues.\allowbreak{}IsShiftAt} (\texttt{Formalism.\allowbreak{}QuotientMetricDivisors.\allowbreak{}DivIntoDualShift}, lines 18--21)\\
\texttt{Formalism.\allowbreak{}DivisionValues.\allowbreak{}IsShift} (\texttt{Formalism.\allowbreak{}QuotientMetricDivisors.\allowbreak{}DivIntoDualShift}, lines 23--24)\par}

{\footnotesize\raggedright\noindent \ref{B4.41}.\quad \texttt{Formalism.\allowbreak{}L2.\allowbreak{}DIV.\allowbreak{}sgRec} (\texttt{Formalism.\allowbreak{}LocalToGlobal.\allowbreak{}Divisibility.\allowbreak{}RecBridge}, lines 387--388)\\
\texttt{Formalism.\allowbreak{}L2.\allowbreak{}DIV.\allowbreak{}shiftZ} (\texttt{Formalism.\allowbreak{}LocalToGlobal.\allowbreak{}Divisibility.\allowbreak{}Inputs}, lines 381--381)\\
\texttt{Formalism.\allowbreak{}L2.\allowbreak{}DIV.\allowbreak{}shiftLat} (\texttt{Formalism.\allowbreak{}LocalToGlobal.\allowbreak{}Divisibility.\allowbreak{}Inputs}, lines 383--385)\\
\texttt{Formalism.\allowbreak{}L2.\allowbreak{}DIV.\allowbreak{}Cons.\allowbreak{}shift\_\allowbreak{}nuOf} (\texttt{Formalism.\allowbreak{}LocalToGlobal.\allowbreak{}Divisibility.\allowbreak{}DivisibilityConstruction.\allowbreak{}ShiftOfColTie}, lines 53--79)\par}

{\footnotesize\raggedright\noindent \ref{B4.42}.\quad \texttt{Formalism.\allowbreak{}DivisionValues.\allowbreak{}shift\_\allowbreak{}pairing\_\allowbreak{}formula} (\texttt{Formalism.\allowbreak{}QuotientMetricDivisors.\allowbreak{}HermitianTwistDivisors.\allowbreak{}TwistedLatticeNaturality}, lines 684--718)\par}

{\footnotesize\raggedright\noindent \ref{B4.43}.\quad \texttt{Formalism.\allowbreak{}DivisionValues.\allowbreak{}reduction\_\allowbreak{}first\_\allowbreak{}step} (\texttt{Formalism.\allowbreak{}QuotientMetricDivisors.\allowbreak{}HermitianTwistDivisors.\allowbreak{}VwEqZeroOfNotMem}, lines 34--60)\\
\texttt{Formalism.\allowbreak{}DivisionValues.\allowbreak{}divisor\_\allowbreak{}into\_\allowbreak{}dual\_\allowbreak{}lattice} (\texttt{Formalism.\allowbreak{}QuotientMetricDivisors.\allowbreak{}HermitianTwistDivisors.\allowbreak{}DivIntoDual}, lines 49--61)\par}

{\footnotesize\raggedright\noindent \ref{B4.44}.\quad \texttt{Formalism.\allowbreak{}DivisionValues.\allowbreak{}divisor\_\allowbreak{}into\_\allowbreak{}dual\_\allowbreak{}shift} (\texttt{Formalism.\allowbreak{}QuotientMetricDivisors.\allowbreak{}HermitianTwistDivisors.\allowbreak{}TwistedLatticeNaturality}, lines 720--730)\par}

{\footnotesize\raggedright\noindent \ref{B4.45}.\quad \texttt{Formalism.\allowbreak{}L2.\allowbreak{}Indef.\allowbreak{}IndefColsR.\allowbreak{}intoDual} (\texttt{Formalism.\allowbreak{}LocalToGlobal.\allowbreak{}IndefiniteForms.\allowbreak{}PairingConditionOfInverseMatrixIdentity}, lines 393--408)\\
\texttt{Formalism.\allowbreak{}L2.\allowbreak{}DIV.\allowbreak{}oneMap\_\allowbreak{}of\_\allowbreak{}coord} (\texttt{Formalism.\allowbreak{}LocalToGlobal.\allowbreak{}Divisibility.\allowbreak{}RecBridge}, lines 33--50)\par}

\section*{Chapter~\ref{ch:p4c6}: Slope change and the data for step 5}

{\footnotesize\raggedright\noindent \ref{B4.46}.\quad \texttt{Formalism.\allowbreak{}MainNotions.\allowbreak{}J.\allowbreak{}mu\_\allowbreak{}smul\_\allowbreak{}ideal} (\texttt{Formalism.\allowbreak{}SharpRestriction.\allowbreak{}AdmissibleParameterJoins.\allowbreak{}SumLogPrimeFactorsLe}, lines 30--72)\par}

{\footnotesize\raggedright\noindent \ref{B4.47}.\quad \texttt{Formalism.\allowbreak{}MainNotions.\allowbreak{}J.\allowbreak{}mu\_\allowbreak{}frac} (\texttt{Formalism.\allowbreak{}SharpRestriction.\allowbreak{}AdmissibleParameterJoins.\allowbreak{}SumLogPrimeFactorsLe}, lines 73--171)\par}

{\footnotesize\raggedright\noindent \ref{B4.48}.\quad \texttt{Formalism.\allowbreak{}MainNotions.\allowbreak{}J.\allowbreak{}Iu\_\allowbreak{}ne\_\allowbreak{}zero} (\texttt{Formalism.\allowbreak{}SharpRestriction.\allowbreak{}AdmissibleParameterJoins.\allowbreak{}SumLogPrimeFactorsLe}, lines 352--357)\\
\texttt{Formalism.\allowbreak{}MainNotions.\allowbreak{}J.\allowbreak{}log\_\allowbreak{}absNorm\_\allowbreak{}Iu} (\texttt{Formalism.\allowbreak{}SharpRestriction.\allowbreak{}AdmissibleParameterJoins.\allowbreak{}SumLogPrimeFactorsLe}, lines 358--379)\par}

{\footnotesize\raggedright\noindent \ref{B4.49}.\quad \texttt{Formalism.\allowbreak{}MainNotions.\allowbreak{}J.\allowbreak{}mu\_\allowbreak{}diff} (\texttt{Formalism.\allowbreak{}SharpRestriction.\allowbreak{}AdmissibleParameterJoins.\allowbreak{}SumLogPrimeFactorsLe}, lines 380--398)\\
\texttt{Formalism.\allowbreak{}BasicNotions.\allowbreak{}SizeData₂.\allowbreak{}NuH} (\texttt{Formalism.\allowbreak{}LocalRingCardinality.\allowbreak{}DataParams₂}, lines 236--237)\par}

{\footnotesize\raggedright\noindent \ref{B4.50}.\quad \texttt{Formalism.\allowbreak{}DivisionValues.\allowbreak{}shift\_\allowbreak{}index\_\allowbreak{}formula} (\texttt{Formalism.\allowbreak{}QuotientMetricDivisors.\allowbreak{}HermitianTwistDivisors.\allowbreak{}TwistedLatticeNaturality}, lines 90--98)\\
\texttt{Formalism.\allowbreak{}DivisionValues.\allowbreak{}mu\_\allowbreak{}log\_\allowbreak{}conductor} (\texttt{Formalism.\allowbreak{}QuotientMetricDivisors.\allowbreak{}HermitianTwistDivisors.\allowbreak{}TwistedLatticeNaturality}, lines 101--129)\\
\texttt{Formalism.\allowbreak{}DivisionValues.\allowbreak{}shift\_\allowbreak{}slope\_\allowbreak{}mu} (\texttt{Formalism.\allowbreak{}QuotientMetricDivisors.\allowbreak{}HermitianTwistDivisors.\allowbreak{}TwistedLatticeNaturality}, lines 145--183)\par}

{\footnotesize\raggedright\noindent \ref{B4.51}.\quad \texttt{Formalism.\allowbreak{}MainNotions.\allowbreak{}J.\allowbreak{}round\_\allowbreak{}le} (\texttt{Formalism.\allowbreak{}SharpRestriction.\allowbreak{}AdmissibleParameterJoins.\allowbreak{}SumLogPrimeFactorsLe}, lines 401--430)\par}

{\footnotesize\raggedright\noindent \ref{B4.52}.\quad \texttt{Formalism.\allowbreak{}MainNotions.\allowbreak{}J.\allowbreak{}sum\_\allowbreak{}e\_\allowbreak{}f} (\texttt{Formalism.\allowbreak{}SharpRestriction.\allowbreak{}AdmissibleParameterJoins.\allowbreak{}SumLogPrimeFactorsLe}, lines 202--217)\\
\texttt{Formalism.\allowbreak{}MainNotions.\allowbreak{}J.\allowbreak{}sum\_\allowbreak{}lw\_\allowbreak{}e} (\texttt{Formalism.\allowbreak{}SharpRestriction.\allowbreak{}AdmissibleParameterJoins.\allowbreak{}SumLogPrimeFactorsLe}, lines 218--236)\par}

{\footnotesize\raggedright\noindent \ref{B4.52a}.\quad \texttt{Formalism.\allowbreak{}MainNotions.\allowbreak{}WeightDegreeSizes.\allowbreak{}gain\_\allowbreak{}above} (\texttt{Formalism.\allowbreak{}SharpRestriction.\allowbreak{}Divisibility.\allowbreak{}WeightDegreeSizes}, lines 294--304)\\
\texttt{Formalism.\allowbreak{}BasicNotions.\allowbreak{}Eval\_\allowbreak{}above} (\texttt{Formalism.\allowbreak{}LocalRingCardinality.\allowbreak{}DataErr}, lines 174--188)\\
\texttt{Formalism.\allowbreak{}MainNotions.\allowbreak{}WeightDegreeSizes.\allowbreak{}zero\_\allowbreak{}gain\_\allowbreak{}above\_\allowbreak{}family} (\texttt{Formalism.\allowbreak{}SharpRestriction.\allowbreak{}Divisibility.\allowbreak{}WeightDegreeSizes}, lines 399--401)\\
\texttt{Formalism.\allowbreak{}L2.\allowbreak{}DIV.\allowbreak{}Wire.\allowbreak{}za\_\allowbreak{}of} (\texttt{Formalism.\allowbreak{}LocalToGlobal.\allowbreak{}Divisibility.\allowbreak{}RecBridgeWire}, lines 19--21)\par}

{\footnotesize\raggedright\noindent \ref{B4.53}.\quad \texttt{Formalism.\allowbreak{}MainNotions.\allowbreak{}J.\allowbreak{}sg\_\allowbreak{}eq\_\allowbreak{}zero} (\texttt{Formalism.\allowbreak{}SharpRestriction.\allowbreak{}AdmissibleParameterJoins.\allowbreak{}SumLogPrimeFactorsLe}, lines 431--447)\\
\texttt{Formalism.\allowbreak{}MainNotions.\allowbreak{}J.\allowbreak{}nprimeMu\_\allowbreak{}at} (\texttt{Formalism.\allowbreak{}SharpRestriction.\allowbreak{}AdmissibleParameterJoins.\allowbreak{}SumLogPrimeFactorsLe}, lines 448--493)\par}

{\footnotesize\raggedright\noindent \ref{def:p4c6-prime-sizes}.\quad \texttt{Formalism.\allowbreak{}MainNotions.\allowbreak{}SizeQD} (\texttt{Formalism.\allowbreak{}SharpRestriction.\allowbreak{}ElevenLe}, lines 56--63)\\
\texttt{Formalism.\allowbreak{}MainNotions.\allowbreak{}SizeG} (\texttt{Formalism.\allowbreak{}SharpRestriction.\allowbreak{}ElevenLe}, lines 123--125)\par}

{\footnotesize\raggedright\noindent \ref{B4.54}.\quad \texttt{Formalism.\allowbreak{}MainNotions.\allowbreak{}J.\allowbreak{}nprimeMuR\_\allowbreak{}at} (\texttt{Formalism.\allowbreak{}SharpRestriction.\allowbreak{}AdmissibleParameterJoins.\allowbreak{}SumLogPrimeFactorsLe}, lines 694--710)\par}

{\footnotesize\raggedright\noindent \ref{B4.54a}.\quad \texttt{Formalism.\allowbreak{}L2.\allowbreak{}Indef.\allowbreak{}Tree.\allowbreak{}local\_\allowbreak{}bound} (\texttt{Formalism.\allowbreak{}LocalToGlobal.\allowbreak{}IndefiniteForms.\allowbreak{}IndefiniteTree.\allowbreak{}ExampleSatisfiesConditionsAtFive}, lines 409--453)\\
\texttt{Formalism.\allowbreak{}EigenvalueLevelCeiling.\allowbreak{}lemmaGG\_\allowbreak{}of} (\texttt{Formalism.\allowbreak{}MirabolicMackeyTheory.\allowbreak{}MirabolicMackeyProofs.\allowbreak{}LemmaGG}, lines 218--366)\\
\texttt{Formalism.\allowbreak{}EigenvalueLevelCeiling.\allowbreak{}idxPsiGood\_\allowbreak{}of} (\texttt{Formalism.\allowbreak{}MirabolicMackeyTheory.\allowbreak{}MirabolicMackeyProofs.\allowbreak{}IdxPsiGood}, lines 18--122)\par}

{\footnotesize\raggedright\noindent \ref{def:p4c6-core}.\quad \texttt{Formalism.\allowbreak{}L2.\allowbreak{}Indef.\allowbreak{}IndefData} (\texttt{Formalism.\allowbreak{}LocalToGlobal.\allowbreak{}IndefiniteForms.\allowbreak{}IndefSide}, lines 241--269)\\
\texttt{Formalism.\allowbreak{}L2.\allowbreak{}Indef.\allowbreak{}IndefSide} (\texttt{Formalism.\allowbreak{}LocalToGlobal.\allowbreak{}IndefiniteForms.\allowbreak{}IndefSide}, lines 363--367)\\
\texttt{Formalism.\allowbreak{}L2.\allowbreak{}Indef.\allowbreak{}ColData.\allowbreak{}L1e} (\texttt{Formalism.\allowbreak{}LocalToGlobal.\allowbreak{}IndefiniteForms.\allowbreak{}IndefSide}, lines 388--394)\\
\texttt{Formalism.\allowbreak{}L2.\allowbreak{}Indef.\allowbreak{}IndefColsR} (\texttt{Formalism.\allowbreak{}LocalToGlobal.\allowbreak{}IndefiniteForms.\allowbreak{}PairingConditionOfInverseMatrixIdentity}, lines 218--237)\\
\texttt{Formalism.\allowbreak{}L2.\allowbreak{}Indef.\allowbreak{}IndefExtraR} (\texttt{Formalism.\allowbreak{}LocalToGlobal.\allowbreak{}IndefiniteForms.\allowbreak{}PairingConditionOfInverseMatrixIdentity}, lines 239--248)\\
\texttt{Formalism.\allowbreak{}L2.\allowbreak{}Indef.\allowbreak{}IndefCoreR} (\texttt{Formalism.\allowbreak{}LocalToGlobal.\allowbreak{}IndefiniteForms.\allowbreak{}PairingConditionOfInverseMatrixIdentity}, lines 250--256)\par}

{\footnotesize\raggedright\noindent \ref{def:p4c6-definite-blocks}.\quad \texttt{Formalism.\allowbreak{}L2.\allowbreak{}DIV.\allowbreak{}DefFam3.\allowbreak{}Ok3} (\texttt{Formalism.\allowbreak{}LocalToGlobal.\allowbreak{}Divisibility.\allowbreak{}InertRationalAdjoint}, lines 139--145)\\
\texttt{Formalism.\allowbreak{}L2.\allowbreak{}DIV.\allowbreak{}InertDat3.\allowbreak{}Ok3} (\texttt{Formalism.\allowbreak{}LocalToGlobal.\allowbreak{}Divisibility.\allowbreak{}InertRationalAdjoint}, lines 53--61)\\
\texttt{Formalism.\allowbreak{}L2.\allowbreak{}DIV.\allowbreak{}DefFam.\allowbreak{}Ok} (\texttt{Formalism.\allowbreak{}LocalToGlobal.\allowbreak{}Divisibility.\allowbreak{}RowLat}, lines 267--277)\\
\texttt{Formalism.\allowbreak{}L2.\allowbreak{}DIV.\allowbreak{}BlockDat.\allowbreak{}Ok} (\texttt{Formalism.\allowbreak{}LocalToGlobal.\allowbreak{}Divisibility.\allowbreak{}RowLat}, lines 176--198)\\
\texttt{Formalism.\allowbreak{}L2.\allowbreak{}DIV.\allowbreak{}CSD.\allowbreak{}Good} (\texttt{Formalism.\allowbreak{}LocalToGlobal.\allowbreak{}Divisibility.\allowbreak{}ClassSum}, lines 92--102)\\
\texttt{Formalism.\allowbreak{}L2.\allowbreak{}DIV.\allowbreak{}CSD.\allowbreak{}ActGood} (\texttt{Formalism.\allowbreak{}LocalToGlobal.\allowbreak{}Divisibility.\allowbreak{}GoodPartHermitianSpec}, lines 76--82)\\
\texttt{Formalism.\allowbreak{}L2.\allowbreak{}DIV.\allowbreak{}CSD.\allowbreak{}ActRat} (\texttt{Formalism.\allowbreak{}LocalToGlobal.\allowbreak{}Divisibility.\allowbreak{}GoodPartHermitianSpec}, lines 84--84)\\
\texttt{Formalism.\allowbreak{}L2.\allowbreak{}Col.\allowbreak{}BlockOk} (\texttt{Formalism.\allowbreak{}LocalToGlobal.\allowbreak{}Columns.\allowbreak{}HeckeSp}, lines 580--583)\par}

{\footnotesize\raggedright\noindent \ref{def:p4c6-division}.\quad \texttt{Formalism.\allowbreak{}L2.\allowbreak{}DIV.\allowbreak{}DivInputs0R} (\texttt{Formalism.\allowbreak{}LocalToGlobal.\allowbreak{}Divisibility.\allowbreak{}ConstantGeTwentyEight}, lines 26--103)\\
\texttt{Formalism.\allowbreak{}L2.\allowbreak{}DIV.\allowbreak{}DivInputsDef3.\allowbreak{}toDivInputs0R} (\texttt{Formalism.\allowbreak{}LocalToGlobal.\allowbreak{}Divisibility.\allowbreak{}TwistedPlaceEtaSumLeExtended}, lines 247--287)\\
\texttt{Formalism.\allowbreak{}DivisorInequality.\allowbreak{}places\_\allowbreak{}ell\_\allowbreak{}weighted\_\allowbreak{}sum\_\allowbreak{}le\_\allowbreak{}one} (\texttt{Formalism.\allowbreak{}DivisorInequality.\allowbreak{}DivisorInequalityProofs.\allowbreak{}BaseLevelBoundsWithCovolumeHolds}, lines 35--46)\par}

{\footnotesize\raggedright\noindent \ref{B4.55}.\quad \texttt{Formalism.\allowbreak{}L2.\allowbreak{}DIV.\allowbreak{}divClausesR\_\allowbreak{}of} (\texttt{Formalism.\allowbreak{}LocalToGlobal.\allowbreak{}Divisibility.\allowbreak{}ConstantGeTwentyEight}, lines 184--248)\\
\texttt{Formalism.\allowbreak{}L2.\allowbreak{}DIV.\allowbreak{}DivClausesR} (\texttt{Formalism.\allowbreak{}LocalToGlobal.\allowbreak{}Divisibility.\allowbreak{}ConstantGeTwentyEight}, lines 111--122)\\
\texttt{Formalism.\allowbreak{}L2.\allowbreak{}DIV.\allowbreak{}HtNums.\allowbreak{}le\_\allowbreak{}hu} (\texttt{Formalism.\allowbreak{}LocalToGlobal.\allowbreak{}Divisibility.\allowbreak{}DivisionHeightBoundStatement}, lines 113--123)\par}

{\footnotesize\raggedright\noindent \ref{B4.55b}.\quad \texttt{Formalism.\allowbreak{}L2.\allowbreak{}DIV.\allowbreak{}two\_\allowbreak{}sumTw\_\allowbreak{}eq\_\allowbreak{}Fo} (\texttt{Formalism.\allowbreak{}LocalToGlobal.\allowbreak{}Divisibility.\allowbreak{}SmallTies}, lines 49--80)\\
\texttt{Formalism.\allowbreak{}L2.\allowbreak{}DIV.\allowbreak{}two\_\allowbreak{}FoOfR\_\allowbreak{}eq\_\allowbreak{}Fo} (\texttt{Formalism.\allowbreak{}LocalToGlobal.\allowbreak{}Divisibility.\allowbreak{}SmallTies}, lines 86--89)\\
\texttt{Formalism.\allowbreak{}L2.\allowbreak{}DIV.\allowbreak{}fixedScale\_\allowbreak{}tied} (\texttt{Formalism.\allowbreak{}LocalToGlobal.\allowbreak{}Divisibility.\allowbreak{}EigenvalueSumLeClosedOf}, lines 350--353)\\
\texttt{Formalism.\allowbreak{}L2.\allowbreak{}DIV.\allowbreak{}ofRecord\_\allowbreak{}Fval} (\texttt{Formalism.\allowbreak{}LocalToGlobal.\allowbreak{}Divisibility.\allowbreak{}NumbersLemmas}, lines 780--782)\\
\texttt{Formalism.\allowbreak{}L2.\allowbreak{}DIV.\allowbreak{}Fj\_\allowbreak{}ofRecord\_\allowbreak{}eq} (\texttt{Formalism.\allowbreak{}LocalToGlobal.\allowbreak{}Divisibility.\allowbreak{}EigenvalueSumLeClosedOf}, lines 332--340)\par}

{\footnotesize\raggedright\noindent \ref{B4.55a}.\quad \texttt{Formalism.\allowbreak{}IndefR.\allowbreak{}LatticeTransfer.\allowbreak{}levelAmb\_\allowbreak{}transfer} (\texttt{Formalism.\allowbreak{}LocalToGlobal.\allowbreak{}IndefiniteReal.\allowbreak{}LevelAmbTransfer}, lines 13--55)\\
\texttt{Formalism.\allowbreak{}BranchingDeterminantCap.\allowbreak{}meanTransferBound} (\texttt{Formalism.\allowbreak{}BranchSandwichDivisors.\allowbreak{}BranchIsometryProofs.\allowbreak{}TransferMean}, lines 19--65)\\
\texttt{Formalism.\allowbreak{}IndefR.\allowbreak{}CxPass.\allowbreak{}In3.\allowbreak{}passage\_\allowbreak{}of\_\allowbreak{}steps} (\texttt{Formalism.\allowbreak{}LocalToGlobal.\allowbreak{}IndefiniteReal.\allowbreak{}CutPassage3SOfSteps}, lines 49--109)\par}

{\footnotesize\raggedright\noindent \ref{B4.56}.\quad \texttt{Formalism.\allowbreak{}L2.\allowbreak{}DIV.\allowbreak{}DivisionPackageFace.\allowbreak{}wordFacts} (\texttt{Formalism.\allowbreak{}LocalToGlobal.\allowbreak{}Divisibility.\allowbreak{}RecBridge}, lines 633--642)\\
\texttt{Formalism.\allowbreak{}L2.\allowbreak{}DIV.\allowbreak{}WordFacts} (\texttt{Formalism.\allowbreak{}LocalToGlobal.\allowbreak{}Divisibility.\allowbreak{}RecBridge}, lines 128--134)\par}

\section*{Chapter~\ref{ch:lattice-after-division}: The lattice after division}

{\footnotesize\raggedright\noindent \ref{B5.01}.\quad \texttt{Formalism.\allowbreak{}LFunctionChainParameters.\allowbreak{}characterSumData.\allowbreak{}ASK} (\texttt{Formalism.\allowbreak{}SizeWindowBrackets.\allowbreak{}FirstPrimeNeSecondPrime}, lines 84--88)\\
\texttt{Formalism.\allowbreak{}LFunctionChainParameters.\allowbreak{}characterSumData.\allowbreak{}askP} (\texttt{Formalism.\allowbreak{}SizeWindowBrackets.\allowbreak{}FirstPrimeNeSecondPrime}, lines 75--77)\\
\texttt{Formalism.\allowbreak{}LFunctionChainParameters.\allowbreak{}characterSumData.\allowbreak{}askD} (\texttt{Formalism.\allowbreak{}SizeWindowBrackets.\allowbreak{}FirstPrimeNeSecondPrime}, lines 79--82)\\
\texttt{Formalism.\allowbreak{}MainNotions.\allowbreak{}Lnk.\allowbreak{}Mod.\allowbreak{}data} (\texttt{Formalism.\allowbreak{}SharpRestriction.\allowbreak{}GlobalSchurBounds.\allowbreak{}QuadraticLogDerivCongr}, lines 464--505)\par}

{\footnotesize\raggedright\noindent \ref{B5.01a}.\quad \texttt{Formalism.\allowbreak{}LFunctionChainParameters.\allowbreak{}characterSumData.\allowbreak{}askP} (\texttt{Formalism.\allowbreak{}SizeWindowBrackets.\allowbreak{}FirstPrimeNeSecondPrime}, lines 75--77)\\
\texttt{Formalism.\allowbreak{}LFunctionChainParameters.\allowbreak{}characterSumData.\allowbreak{}askD} (\texttt{Formalism.\allowbreak{}SizeWindowBrackets.\allowbreak{}FirstPrimeNeSecondPrime}, lines 79--82)\par}

{\footnotesize\raggedright\noindent \ref{B5.02}.\quad \texttt{Formalism.\allowbreak{}BasicNotions.\allowbreak{}SizeData₂.\allowbreak{}sg} (\texttt{Formalism.\allowbreak{}LocalRingCardinality.\allowbreak{}DataParams₂}, lines 215--217)\\
\texttt{Formalism.\allowbreak{}BasicNotions.\allowbreak{}SizeData₂.\allowbreak{}sz} (\texttt{Formalism.\allowbreak{}LocalRingCardinality.\allowbreak{}DataParams₂}, lines 205--206)\\
\texttt{Formalism.\allowbreak{}BasicNotions.\allowbreak{}SizeData₂.\allowbreak{}gainD} (\texttt{Formalism.\allowbreak{}LocalRingCardinality.\allowbreak{}DataParams₂}, lines 210--213)\\
\texttt{Formalism.\allowbreak{}BasicNotions.\allowbreak{}SizeData₂.\allowbreak{}Fo} (\texttt{Formalism.\allowbreak{}LocalRingCardinality.\allowbreak{}DataParams₂}, lines 218--219)\\
\texttt{Formalism.\allowbreak{}BasicNotions.\allowbreak{}SizeData₂.\allowbreak{}Ga} (\texttt{Formalism.\allowbreak{}LocalRingCardinality.\allowbreak{}DataParams₂}, lines 220--221)\\
\texttt{Formalism.\allowbreak{}BasicNotions.\allowbreak{}SizeData₂.\allowbreak{}Ee} (\texttt{Formalism.\allowbreak{}LocalRingCardinality.\allowbreak{}DataParams₂}, lines 222--223)\\
\texttt{Formalism.\allowbreak{}MainNotions.\allowbreak{}Lnk.\allowbreak{}IOsum} (\texttt{Formalism.\allowbreak{}SharpRestriction.\allowbreak{}GlobalSchurBounds.\allowbreak{}QuadraticLogDerivCongr}, lines 163--164)\par}

{\footnotesize\raggedright\noindent \ref{B5.02a}.\quad \texttt{Formalism.\allowbreak{}BasicNotions.\allowbreak{}SizeData₂.\allowbreak{}Ga} (\texttt{Formalism.\allowbreak{}LocalRingCardinality.\allowbreak{}DataParams₂}, lines 220--221)\\
\texttt{Formalism.\allowbreak{}BasicNotions.\allowbreak{}SizeData₂.\allowbreak{}Fo} (\texttt{Formalism.\allowbreak{}LocalRingCardinality.\allowbreak{}DataParams₂}, lines 218--219)\par}

{\footnotesize\raggedright\noindent \ref{B5.03}.\quad \texttt{Formalism.\allowbreak{}MainNotions.\allowbreak{}LocalQuadraticParts.\allowbreak{}UnroundedValuationBound} (\texttt{Formalism.\allowbreak{}SharpRestriction.\allowbreak{}Divisibility.\allowbreak{}LocalQuadraticParts}, lines 212--217)\\
\texttt{Formalism.\allowbreak{}BasicNotions.\allowbreak{}DV.\allowbreak{}Iu} (\texttt{Formalism.\allowbreak{}LocalRingCardinality.\allowbreak{}DataDIVDefs}, lines 17--21)\\
\texttt{Formalism.\allowbreak{}BasicNotions.\allowbreak{}DV.\allowbreak{}NprimeDef} (\texttt{Formalism.\allowbreak{}LocalRingCardinality.\allowbreak{}DataDIVDefs}, lines 22--25)\par}

{\footnotesize\raggedright\noindent \ref{B5.04}.\quad \texttt{Formalism.\allowbreak{}L2.\allowbreak{}GlobalTheoryTheory.\allowbreak{}Places.\allowbreak{}mu\_\allowbreak{}index\_\allowbreak{}sum} (\texttt{Formalism.\allowbreak{}LocalToGlobal.\allowbreak{}GlobalTheory.\allowbreak{}PlacesP}, lines 61--103)\\
\texttt{Formalism.\allowbreak{}L2.\allowbreak{}GlobalTheoryTheory.\allowbreak{}Places.\allowbreak{}mu\_\allowbreak{}eq} (\texttt{Formalism.\allowbreak{}LocalToGlobal.\allowbreak{}GlobalTheory.\allowbreak{}PlacesP}, lines 129--170)\par}

{\footnotesize\raggedright\noindent \ref{B5.05}.\quad \texttt{Formalism.\allowbreak{}L2.\allowbreak{}GlobalTheoryTheory.\allowbreak{}Places.\allowbreak{}muMax\_\allowbreak{}le} (\texttt{Formalism.\allowbreak{}LocalToGlobal.\allowbreak{}GlobalTheory.\allowbreak{}PlacesP}, lines 171--238)\par}

{\footnotesize\raggedright\noindent \ref{B5.05b}.\quad \texttt{Formalism.\allowbreak{}MainNotions.\allowbreak{}product\_\allowbreak{}first\_\allowbreak{}zeroth\_\allowbreak{}parts.\allowbreak{}HidaSetting} (\texttt{Formalism.\allowbreak{}SharpRestriction.\allowbreak{}GlobalSchurBounds.\allowbreak{}ProductFirstZerothParts}, lines 131--137)\\
\texttt{Formalism.\allowbreak{}MainNotions.\allowbreak{}GlobalTheoryTheory.\allowbreak{}Alc.\allowbreak{}entry\_\allowbreak{}bd} (\texttt{Formalism.\allowbreak{}SharpRestriction.\allowbreak{}GlobalSchurBounds.\allowbreak{}AlcoveOf}, lines 65--122)\\
\texttt{Formalism.\allowbreak{}MainNotions.\allowbreak{}GlobalTheoryTheory.\allowbreak{}Alc.\allowbreak{}alcove\_\allowbreak{}of} (\texttt{Formalism.\allowbreak{}SharpRestriction.\allowbreak{}GlobalSchurBounds.\allowbreak{}AlcoveOf}, lines 123--155)\\
\texttt{Formalism.\allowbreak{}MainNotions.\allowbreak{}J.\allowbreak{}levelPrime\_\allowbreak{}setting\_\allowbreak{}of\_\allowbreak{}sizes} (\texttt{Formalism.\allowbreak{}SharpRestriction.\allowbreak{}AdmissibleParameterJoins.\allowbreak{}DIVfromGLo}, lines 249--293)\par}

{\footnotesize\raggedright\noindent \ref{B5.05a}.\quad \texttt{Formalism.\allowbreak{}L2.\allowbreak{}GlobalTheoryTheory.\allowbreak{}LocalLatticeBases.\allowbreak{}congruence\_\allowbreak{}no\_\allowbreak{}kernel\_\allowbreak{}minimal} (\texttt{Formalism.\allowbreak{}LocalToGlobal.\allowbreak{}GlobalTheory.\allowbreak{}GL6.\allowbreak{}CongruenceNoKernelOfHeckeRamified}, lines 75--159)\\
\texttt{Formalism.\allowbreak{}MainNotions.\allowbreak{}J.\allowbreak{}division\_\allowbreak{}rank\_\allowbreak{}slopes\_\allowbreak{}of\_\allowbreak{}lines} (\texttt{Formalism.\allowbreak{}SharpRestriction.\allowbreak{}GlobalSchurBounds.\allowbreak{}LinkMin}, lines 497--583)\\
\texttt{Formalism.\allowbreak{}MainNotions.\allowbreak{}J.\allowbreak{}lemma\_\allowbreak{}five\_\allowbreak{}one\_\allowbreak{}of\_\allowbreak{}theorem\_\allowbreak{}two} (\texttt{Formalism.\allowbreak{}SharpRestriction.\allowbreak{}GlobalSchurBounds.\allowbreak{}LinkMin}, lines 587--613)\par}

{\footnotesize\raggedright\noindent \ref{B5.06}.\quad \texttt{Formalism.\allowbreak{}MainNotions.\allowbreak{}Lnk.\allowbreak{}slopesN\_\allowbreak{}at\_\allowbreak{}of} (\texttt{Formalism.\allowbreak{}SharpRestriction.\allowbreak{}GlobalSchurBounds.\allowbreak{}TopDensityNGeSixteen}, lines 1529--1540)\\
\texttt{Formalism.\allowbreak{}LFunctionChainParameters.\allowbreak{}WeightedSlopeBounds.\allowbreak{}QuotientSlopeBoundsWith} (\texttt{Formalism.\allowbreak{}SizeWindowBrackets.\allowbreak{}QuotientSlopeBoundsWith}, lines 19--23)\\
\texttt{Formalism.\allowbreak{}L2.\allowbreak{}GlobalTheoryTheory.\allowbreak{}Places.\allowbreak{}places\_\allowbreak{}at\_\allowbreak{}of} (\texttt{Formalism.\allowbreak{}LocalToGlobal.\allowbreak{}GlobalTheory.\allowbreak{}PlacesP}, lines 297--436)\par}

{\footnotesize\raggedright\noindent \ref{B5.07}.\quad \texttt{Formalism.\allowbreak{}L2.\allowbreak{}DIV.\allowbreak{}height\_\allowbreak{}chain\_\allowbreak{}atR} (\texttt{Formalism.\allowbreak{}LocalToGlobal.\allowbreak{}Divisibility.\allowbreak{}ConstantGeTwentyEight}, lines 123--183)\\
\texttt{Formalism.\allowbreak{}BasicNotions.\allowbreak{}DV.\allowbreak{}hExpl6} (\texttt{Formalism.\allowbreak{}LocalRingCardinality.\allowbreak{}KernelPlaceMeanBound}, lines 25--28)\\
\texttt{Formalism.\allowbreak{}BasicNotions.\allowbreak{}DV.\allowbreak{}logHook} (\texttt{Formalism.\allowbreak{}LocalRingCardinality.\allowbreak{}KernelPlaceMeanBound}, lines 9--11)\\
\texttt{Formalism.\allowbreak{}BasicNotions.\allowbreak{}DV.\allowbreak{}Lprime} (\texttt{Formalism.\allowbreak{}LocalRingCardinality.\allowbreak{}KernelPlaceMeanBound}, lines 13--14)\\
\texttt{Formalism.\allowbreak{}BasicNotions.\allowbreak{}DV.\allowbreak{}L6lo} (\texttt{Formalism.\allowbreak{}LocalRingCardinality.\allowbreak{}KernelPlaceMeanBound}, lines 16--18)\\
\texttt{Formalism.\allowbreak{}BasicNotions.\allowbreak{}DV.\allowbreak{}kernelPlaceMeanBound} (\texttt{Formalism.\allowbreak{}LocalRingCardinality.\allowbreak{}KernelPlaceMeanBound}, lines 20--23)\par}

{\footnotesize\raggedright\noindent \ref{B5.07c}.\quad \texttt{GlobalTheory.\allowbreak{}BranchNorm.\allowbreak{}brNorm} (\texttt{Development.\allowbreak{}AbelianVarietyTheory.\allowbreak{}Proofs.\allowbreak{}DivisorProofs.\allowbreak{}LastMapNorm}, lines 336--394)\\
\texttt{GlobalTheory.\allowbreak{}LastPeriod.\allowbreak{}coefOpLe\_\allowbreak{}brAt} (\texttt{Development.\allowbreak{}AbelianVarietyTheory.\allowbreak{}Proofs.\allowbreak{}DivisorProofs.\allowbreak{}PeriodBoundLastLog}, lines 26--44)\\
\texttt{GlobalTheory.\allowbreak{}ColRowsTop.\allowbreak{}periodFun\_\allowbreak{}le} (\texttt{Development.\allowbreak{}AbelianVarietyTheory.\allowbreak{}Proofs.\allowbreak{}DivisorProofs.\allowbreak{}IpAtNonnegOfDatFramed}, lines 26--86)\par}

{\footnotesize\raggedright\noindent \ref{B5.07d}.\quad \texttt{GlobalTheory.\allowbreak{}LevelSupport.\allowbreak{}sum\_\allowbreak{}heisO\_\allowbreak{}along\_\allowbreak{}eq\_\allowbreak{}zero\_\allowbreak{}level} (\texttt{Development.\allowbreak{}AbelianVarietyTheory.\allowbreak{}Proofs.\allowbreak{}DivisorProofs.\allowbreak{}LevelFockVectorHeisenbergAverageIsotropicEqZero}, lines 274--319)\\
\texttt{GlobalTheory.\allowbreak{}DefLiftCusp.\allowbreak{}zcusp\_\allowbreak{}lift} (\texttt{Development.\allowbreak{}AbelianVarietyTheory.\allowbreak{}Proofs.\allowbreak{}IwasawaDecomposition.\allowbreak{}ThetaLiftDefiniteConstantTermEqZero}, lines 66--128)\\
\texttt{GlobalTheory.\allowbreak{}TransvCusp.\allowbreak{}isCuspFun\_\allowbreak{}of\_\allowbreak{}zCusp} (\texttt{Development.\allowbreak{}AbelianVarietyTheory.\allowbreak{}Proofs.\allowbreak{}DivisorProofs.\allowbreak{}IsCuspFunOfTransvectionConstantTermsZero}, lines 142--163)\par}

{\footnotesize\raggedright\noindent \ref{B5.07e}.\quad \texttt{GlobalTheory.\allowbreak{}VacCoef.\allowbreak{}vac\_\allowbreak{}coef} (\texttt{Development.\allowbreak{}AbelianVarietyTheory.\allowbreak{}Proofs.\allowbreak{}TwistedOperators.\allowbreak{}WeilMatrixCoefficientVacuaUnramifiedPlace}, lines 109--171)\\
\texttt{GlobalTheory.\allowbreak{}LocZeta.\allowbreak{}hasSum\_\allowbreak{}localZeta} (\texttt{Development.\allowbreak{}AbelianVarietyTheory.\allowbreak{}Proofs.\allowbreak{}TwistedOperators.\allowbreak{}HasSumLocalSphericalZetaUnramified}, lines 101--197)\par}

{\footnotesize\raggedright\noindent \ref{B5.07y}.\quad \texttt{GlobalTheory.\allowbreak{}LowRow.\allowbreak{}scalar\_\allowbreak{}of\_\allowbreak{}form\_\allowbreak{}key\_\allowbreak{}lowC} (\texttt{Development.\allowbreak{}AbelianVarietyTheory.\allowbreak{}Proofs.\allowbreak{}DivisorProofs.\allowbreak{}AdelCOfInnerProductFormula}, lines 214--243)\\
\texttt{GlobalTheory.\allowbreak{}LowForm.\allowbreak{}exists\_\allowbreak{}full\_\allowbreak{}lift\_\allowbreak{}six} (\texttt{Development.\allowbreak{}AbelianVarietyTheory.\allowbreak{}Proofs.\allowbreak{}DivisorProofs.\allowbreak{}RationalLiftLowerBoundOffBasePlace}, lines 152--177)\\
\texttt{GlobalTheory.\allowbreak{}LowForm.\allowbreak{}conjLift\_\allowbreak{}lower\_\allowbreak{}six} (\texttt{Development.\allowbreak{}AbelianVarietyTheory.\allowbreak{}Proofs.\allowbreak{}DivisorProofs.\allowbreak{}RationalLiftLowerBoundOffBasePlace}, lines 196--207)\par}

{\footnotesize\raggedright\noindent \ref{B5.07a}.\quad \texttt{Formalism.\allowbreak{}DivisorInequality.\allowbreak{}height\_\allowbreak{}bound\_\allowbreak{}at\_\allowbreak{}set} (\texttt{Formalism.\allowbreak{}DivisorInequality.\allowbreak{}DivisorInequalityProofs.\allowbreak{}EvenDegreeWeightBounds}, lines 154--288)\\
\texttt{Formalism.\allowbreak{}DivisorInequality.\allowbreak{}UniformCovolumeTail.\allowbreak{}stepDiSharp} (\texttt{Formalism.\allowbreak{}DivisorInequality.\allowbreak{}DivisorInequalityProofs.\allowbreak{}EvenDegreeWeightBounds}, lines 27--149)\\
\texttt{Formalism.\allowbreak{}L2.\allowbreak{}DIV.\allowbreak{}nK\_\allowbreak{}of\_\allowbreak{}cb} (\texttt{Formalism.\allowbreak{}LocalToGlobal.\allowbreak{}Divisibility.\allowbreak{}DivisionHeightBoundStatement}, lines 565--584)\par}

{\footnotesize\raggedright\noindent \ref{B5.07fa}.\quad \texttt{GlobalTheory.\allowbreak{}GalOne.\allowbreak{}ClassConj.\allowbreak{}class\_\allowbreak{}scalar\_\allowbreak{}conj\_\allowbreak{}ofFive} (\texttt{Development.\allowbreak{}AbelianVarietyTheory.\allowbreak{}Proofs.\allowbreak{}DivisorProofs.\allowbreak{}CoefMapNormEquivariantBasisOfFive}, lines 2202--2272)\par}

{\footnotesize\raggedright\noindent \ref{B5.07f}.\quad \texttt{GlobalTheory.\allowbreak{}ConjRatLiftFive.\allowbreak{}base\_\allowbreak{}lift5} (\texttt{Development.\allowbreak{}AbelianVarietyTheory.\allowbreak{}Proofs.\allowbreak{}DivisorProofs.\allowbreak{}ConjugatesOfRationalThetaLiftFromCharacterisatio}, lines 46--68)\\
\texttt{GlobalTheory.\allowbreak{}ConjRatLiftFive.\allowbreak{}conj\_\allowbreak{}lift5} (\texttt{Development.\allowbreak{}AbelianVarietyTheory.\allowbreak{}Proofs.\allowbreak{}DivisorProofs.\allowbreak{}ConjugatesOfRationalThetaLiftFromCharacterisatio}, lines 71--139)\par}

{\footnotesize\raggedright\noindent \ref{B5.07g}.\quad \texttt{Formalism.\allowbreak{}IndefR.\allowbreak{}Links.\allowbreak{}RealDat3.\allowbreak{}PartAt} (\texttt{Formalism.\allowbreak{}LocalToGlobal.\allowbreak{}IndefiniteReal.\allowbreak{}ZeroPartEigenvalueConvention}, lines 190--193)\\
\texttt{Formalism.\allowbreak{}IndefR.\allowbreak{}Links.\allowbreak{}RealDat3.\allowbreak{}scal\_\allowbreak{}off} (\texttt{Formalism.\allowbreak{}LocalToGlobal.\allowbreak{}IndefiniteReal.\allowbreak{}ZeroPartEigenvalueConvention}, lines 202--209)\\
\texttt{Formalism.\allowbreak{}IndefR.\allowbreak{}Links.\allowbreak{}RealDat3.\allowbreak{}scal\_\allowbreak{}abs} (\texttt{Formalism.\allowbreak{}LocalToGlobal.\allowbreak{}IndefiniteReal.\allowbreak{}ZeroPartEigenvalueConvention}, lines 212--214)\\
\texttt{GlobalTheory.\allowbreak{}ConjScalarOff.\allowbreak{}galD\_\allowbreak{}base} (\texttt{Development.\allowbreak{}AbelianVarietyTheory.\allowbreak{}Proofs.\allowbreak{}TwistedOperators.\allowbreak{}LiftScalarIdentityOffFixedPlaceBaseLevel}, lines 40--73)\par}

{\footnotesize\raggedright\noindent \ref{B5.07b}.\quad \texttt{GlobalTheory.\allowbreak{}LowRow.\allowbreak{}lamLow\_\allowbreak{}of\_\allowbreak{}weak} (\texttt{Development.\allowbreak{}AbelianVarietyTheory.\allowbreak{}Proofs.\allowbreak{}DivisorProofs.\allowbreak{}LiftEigenvalueLowerBoundOfWeakCharacterisation}, lines 20--41)\\
\texttt{GlobalTheory.\allowbreak{}LowRow.\allowbreak{}lamLow\_\allowbreak{}closing\_\allowbreak{}alt} (\texttt{Development.\allowbreak{}AutomorphicArithmetic.\allowbreak{}ThetaCorrespondence.\allowbreak{}ThetaPairLifts.\allowbreak{}LocalTypeDictionary.\allowbreak{}FinitenessInputsRealDataThree}, lines 14--23)\\
\texttt{Formalism.\allowbreak{}L2.\allowbreak{}Indef.\allowbreak{}IndefColsR.\allowbreak{}log\_\allowbreak{}eigenvalue\_\allowbreak{}sum\_\allowbreak{}le\_\allowbreak{}min} (\texttt{Formalism.\allowbreak{}LocalToGlobal.\allowbreak{}IndefiniteForms.\allowbreak{}PairingConditionOfInverseMatrixIdentity}, lines 328--331)\par}

{\footnotesize\raggedright\noindent \ref{B5.07x}.\quad \texttt{GlobalTheory.\allowbreak{}MeanPeriodClosed.\allowbreak{}meanPeriodAt} (\texttt{Development.\allowbreak{}AbelianVarietyTheory.\allowbreak{}Proofs.\allowbreak{}DivisorProofs.\allowbreak{}MeanLogPeriodChowlaSelberg}, lines 58--82)\\
\texttt{GlobalTheory.\allowbreak{}MeanPeriodGenerator.\allowbreak{}mean\_\allowbreak{}logMet\_\allowbreak{}eq} (\texttt{Development.\allowbreak{}AbelianVarietyTheory.\allowbreak{}Proofs.\allowbreak{}DivisorProofs.\allowbreak{}MeanLogMetricGeneratorEqMeanPeriod}, lines 45--57)\par}

{\footnotesize\raggedright\noindent \ref{B5.08d}.\quad \texttt{GlobalTheory.\allowbreak{}EX.\allowbreak{}exists\_\allowbreak{}hc\_\allowbreak{}family} (\texttt{Development.\allowbreak{}AbelianVarietyTheory.\allowbreak{}Proofs.\allowbreak{}GlobalInterpolation.\allowbreak{}ExistsHcParamOfCode}, lines 128--242)\par}

{\footnotesize\raggedright\noindent \ref{B5.08b}.\quad \texttt{Formalism.\allowbreak{}L2.\allowbreak{}DIV.\allowbreak{}NumDatum.\allowbreak{}NumFacts} (\texttt{Formalism.\allowbreak{}LocalToGlobal.\allowbreak{}Divisibility.\allowbreak{}NumbersLink}, lines 43--83)\\
\texttt{Formalism.\allowbreak{}L2.\allowbreak{}DIV.\allowbreak{}Arith4S.\allowbreak{}numFacts\_\allowbreak{}rows} (\texttt{Formalism.\allowbreak{}LocalToGlobal.\allowbreak{}Divisibility.\allowbreak{}EigenvalueSumLeClosedOf}, lines 1462--1522)\\
\texttt{Formalism.\allowbreak{}L2.\allowbreak{}DIV.\allowbreak{}Arith4S.\allowbreak{}numFacts\_\allowbreak{}rec} (\texttt{Formalism.\allowbreak{}LocalToGlobal.\allowbreak{}Divisibility.\allowbreak{}EigenvalueSumLeClosedOf}, lines 1969--1980)\\
\texttt{Formalism.\allowbreak{}L2.\allowbreak{}DIV.\allowbreak{}RecNums.\allowbreak{}ofNum} (\texttt{Formalism.\allowbreak{}LocalToGlobal.\allowbreak{}Divisibility.\allowbreak{}NumbersLink}, lines 541--650)\par}

{\footnotesize\raggedright\noindent \ref{B5.08c}.\quad \texttt{Formalism.\allowbreak{}L2.\allowbreak{}DIV.\allowbreak{}NumDatum.\allowbreak{}eEN} (\texttt{Formalism.\allowbreak{}LocalToGlobal.\allowbreak{}Divisibility.\allowbreak{}NumbersLink}, lines 35--37)\\
\texttt{Formalism.\allowbreak{}L2.\allowbreak{}DIV.\allowbreak{}NumDatum.\allowbreak{}dN} (\texttt{Formalism.\allowbreak{}LocalToGlobal.\allowbreak{}Divisibility.\allowbreak{}NumbersLink}, lines 38--40)\\
\texttt{Formalism.\allowbreak{}DivisionValues.\allowbreak{}young\_\allowbreak{}diagram\_\allowbreak{}hook} (\texttt{Formalism.\allowbreak{}QuotientMetricDivisors.\allowbreak{}HermitianTwistDivisors.\allowbreak{}YoungDiagramHook}, lines 386--423)\\
\texttt{Formalism.\allowbreak{}DivisionValues.\allowbreak{}young\_\allowbreak{}diagram\_\allowbreak{}factorial} (\texttt{Formalism.\allowbreak{}QuotientMetricDivisors.\allowbreak{}HermitianTwistDivisors.\allowbreak{}YoungDiagramHook}, lines 250--332)\\
\texttt{Formalism.\allowbreak{}DivisionValues.\allowbreak{}young\_\allowbreak{}diagram\_\allowbreak{}period} (\texttt{Formalism.\allowbreak{}QuotientMetricDivisors.\allowbreak{}HermitianTwistDivisors.\allowbreak{}YoungDiagramHook}, lines 350--376)\\
\texttt{Formalism.\allowbreak{}DivisionValues.\allowbreak{}young\_\allowbreak{}diagram\_\allowbreak{}step\_\allowbreak{}origin} (\texttt{Formalism.\allowbreak{}QuotientMetricDivisors.\allowbreak{}HermitianTwistDivisors.\allowbreak{}YoungDiagramTopExact}, lines 254--296)\\
\texttt{Formalism.\allowbreak{}L2.\allowbreak{}DIV.\allowbreak{}NumDatum.\allowbreak{}stepE\_\allowbreak{}num} (\texttt{Formalism.\allowbreak{}LocalToGlobal.\allowbreak{}Divisibility.\allowbreak{}NumbersLink}, lines 388--538)\par}

{\footnotesize\raggedright\noindent \ref{B5.08}.\quad \texttt{Formalism.\allowbreak{}L2.\allowbreak{}DIV.\allowbreak{}RecNums.\allowbreak{}divIneq} (\texttt{Formalism.\allowbreak{}LocalToGlobal.\allowbreak{}Divisibility.\allowbreak{}DivisionHeightBoundStatement}, lines 131--181)\\
\texttt{Formalism.\allowbreak{}DivisorInequality.\allowbreak{}slope\_\allowbreak{}bounds} (\texttt{Formalism.\allowbreak{}DivisorInequality.\allowbreak{}DivisorInequalityProofs.\allowbreak{}BaseLevelBoundsWithCovolumeHolds}, lines 47--70)\\
\texttt{Formalism.\allowbreak{}DivisionValues.\allowbreak{}lemmaRb\_\allowbreak{}of} (\texttt{Formalism.\allowbreak{}QuotientMetricDivisors.\allowbreak{}HermitianTwistDivisors.\allowbreak{}ReductionLemmaSecond}, lines 163--287)\\
\texttt{Formalism.\allowbreak{}L2.\allowbreak{}Indef.\allowbreak{}IndefSide.\allowbreak{}boundsMu} (\texttt{Formalism.\allowbreak{}LocalToGlobal.\allowbreak{}IndefiniteForms.\allowbreak{}IndefUse}, lines 271--335)\par}

{\footnotesize\raggedright\noindent \ref{B5.08a}.\quad \texttt{Formalism.\allowbreak{}DivisionValues.\allowbreak{}StepBAux.\allowbreak{}idx\_\allowbreak{}change\_\allowbreak{}img} (\texttt{Formalism.\allowbreak{}QuotientMetricDivisors.\allowbreak{}HermitianTwistDivisors.\allowbreak{}ReductionSecondStep}, lines 379--399)\\
\texttt{Formalism.\allowbreak{}DivisionValues.\allowbreak{}ConjugateWeightGram.\allowbreak{}local\_\allowbreak{}idW} (\texttt{Formalism.\allowbreak{}QuotientMetricDivisors.\allowbreak{}HermitianTwistDivisors.\allowbreak{}ConjugateWeightLatticeOfGeneral}, lines 174--318)\\
\texttt{Formalism.\allowbreak{}DivisionValues.\allowbreak{}stepBW\_\allowbreak{}of} (\texttt{Formalism.\allowbreak{}QuotientMetricDivisors.\allowbreak{}HermitianTwistDivisors.\allowbreak{}ConjugateWeightLatticeOfGeneral}, lines 324--499)\par}

{\footnotesize\raggedright\noindent \ref{B5.08x}.\quad \texttt{Formalism.\allowbreak{}DivisionValues.\allowbreak{}ConjugateWeightGram.\allowbreak{}local\_\allowbreak{}idW} (\texttt{Formalism.\allowbreak{}QuotientMetricDivisors.\allowbreak{}HermitianTwistDivisors.\allowbreak{}ConjugateWeightLatticeOfGeneral}, lines 174--318)\\
\texttt{Formalism.\allowbreak{}DivisionValues.\allowbreak{}second\_\allowbreak{}step\_\allowbreak{}refined\_\allowbreak{}weak} (\texttt{Formalism.\allowbreak{}QuotientMetricDivisors.\allowbreak{}HermitianTwistDivisors.\allowbreak{}ConjugateWeightLatticeOfGeneral}, lines 506--517)\\
\texttt{Formalism.\allowbreak{}L2.\allowbreak{}Indef.\allowbreak{}IndefSide.\allowbreak{}boundsMu} (\texttt{Formalism.\allowbreak{}LocalToGlobal.\allowbreak{}IndefiniteForms.\allowbreak{}IndefUse}, lines 271--335)\\
\texttt{Formalism.\allowbreak{}L2.\allowbreak{}DIV.\allowbreak{}DivInputsDef3.\allowbreak{}toDivInputs0R} (\texttt{Formalism.\allowbreak{}LocalToGlobal.\allowbreak{}Divisibility.\allowbreak{}TwistedPlaceEtaSumLeExtended}, lines 247--287)\par}

{\footnotesize\raggedright\noindent \ref{B5.09}.\quad \texttt{Formalism.\allowbreak{}L2.\allowbreak{}DIV.\allowbreak{}DivClausesR} (\texttt{Formalism.\allowbreak{}LocalToGlobal.\allowbreak{}Divisibility.\allowbreak{}ConstantGeTwentyEight}, lines 111--122)\\
\texttt{Formalism.\allowbreak{}L2.\allowbreak{}DIV.\allowbreak{}divClausesR\_\allowbreak{}of} (\texttt{Formalism.\allowbreak{}LocalToGlobal.\allowbreak{}Divisibility.\allowbreak{}ConstantGeTwentyEight}, lines 184--248)\par}

{\footnotesize\raggedright\noindent \ref{B5.09a}.\quad \texttt{GlobalTheory.\allowbreak{}ColumnMultSum.\allowbreak{}columnMult} (\texttt{Development.\allowbreak{}AbelianVarietyTheory.\allowbreak{}Proofs.\allowbreak{}DivisorProofs.\allowbreak{}ReordEqConj}, lines 35--78)\\
\texttt{Formalism.\allowbreak{}L2.\allowbreak{}DIV.\allowbreak{}Arith4S.\allowbreak{}q\_\allowbreak{}rec} (\texttt{Formalism.\allowbreak{}LocalToGlobal.\allowbreak{}Divisibility.\allowbreak{}EigenvalueSumLeClosedOf}, lines 1885--1954)\\
\texttt{Formalism.\allowbreak{}L2.\allowbreak{}DIV.\allowbreak{}hg\_\allowbreak{}of\_\allowbreak{}sizeG} (\texttt{Formalism.\allowbreak{}LocalToGlobal.\allowbreak{}Divisibility.\allowbreak{}EigenvalueSumLeClosedOf}, lines 295--325)\\
\texttt{Formalism.\allowbreak{}L2.\allowbreak{}DIV.\allowbreak{}delta\_\allowbreak{}mul\_\allowbreak{}Q\_\allowbreak{}le\_\allowbreak{}one} (\texttt{Formalism.\allowbreak{}LocalToGlobal.\allowbreak{}Divisibility.\allowbreak{}EigenvalueSumLeClosedOf}, lines 126--139)\par}

{\footnotesize\raggedright\noindent \ref{B5.10b}.\quad \texttt{Formalism.\allowbreak{}MainNotions.\allowbreak{}LinkDIV\_\allowbreak{}at} (\texttt{Formalism.\allowbreak{}SharpRestriction.\allowbreak{}GlobalTheoremInputs}, lines 77--90)\par}

{\footnotesize\raggedright\noindent \ref{B5.10}.\quad \texttt{Formalism.\allowbreak{}MainNotions.\allowbreak{}LinkDIV\_\allowbreak{}at} (\texttt{Formalism.\allowbreak{}SharpRestriction.\allowbreak{}GlobalTheoremInputs}, lines 77--90)\\
\texttt{Formalism.\allowbreak{}L2.\allowbreak{}DIV.\allowbreak{}wordFacts\_\allowbreak{}of\_\allowbreak{}face} (\texttt{Formalism.\allowbreak{}LocalToGlobal.\allowbreak{}Divisibility.\allowbreak{}RecBridge}, lines 204--226)\\
\texttt{Formalism.\allowbreak{}L2.\allowbreak{}DIV.\allowbreak{}linkDIV\_\allowbreak{}at\_\allowbreak{}words} (\texttt{Formalism.\allowbreak{}LocalToGlobal.\allowbreak{}Divisibility.\allowbreak{}RecBridge}, lines 327--340)\\
\texttt{Formalism.\allowbreak{}L2.\allowbreak{}DIV.\allowbreak{}DivisionPackageFace.\allowbreak{}wordFacts} (\texttt{Formalism.\allowbreak{}LocalToGlobal.\allowbreak{}Divisibility.\allowbreak{}RecBridge}, lines 633--641)\par}

{\footnotesize\raggedright\noindent \ref{B5.10a}.\quad \texttt{Formalism.\allowbreak{}L2.\allowbreak{}DIV.\allowbreak{}FunSharpTop.\allowbreak{}closed\_\allowbreak{}funSharp\_\allowbreak{}alt} (\texttt{Development.\allowbreak{}AutomorphicArithmetic.\allowbreak{}ThetaCorrespondence.\allowbreak{}ThetaPairLifts.\allowbreak{}ThirdPlaceSeries.\allowbreak{}HefunFunTop}, lines 88--116)\\
\texttt{Formalism.\allowbreak{}MainNotions.\allowbreak{}DIV.\allowbreak{}Proofs.\allowbreak{}height\_\allowbreak{}bound\_\allowbreak{}of\_\allowbreak{}component\_\allowbreak{}bounds} (\texttt{Formalism.\allowbreak{}SharpRestriction.\allowbreak{}Divisibility.\allowbreak{}SlopeExcess.\allowbreak{}HeightBoundOfComponentBounds}, lines 136--151)\par}

{\footnotesize\raggedright\noindent \ref{B5.11c}.\quad \texttt{Formalism.\allowbreak{}MainNotions.\allowbreak{}StandardLattice.\allowbreak{}DataB} (\texttt{Formalism.\allowbreak{}SharpRestriction.\allowbreak{}GlobalSchurBounds.\allowbreak{}StandardLattice}, lines 500--517)\\
\texttt{Formalism.\allowbreak{}MainNotions.\allowbreak{}StandardLattice.\allowbreak{}RowsA} (\texttt{Formalism.\allowbreak{}SharpRestriction.\allowbreak{}GlobalSchurBounds.\allowbreak{}TheoremThreeOfExplicitAndCompact}, lines 115--122)\\
\texttt{Formalism.\allowbreak{}MainNotions.\allowbreak{}StandardLattice.\allowbreak{}DataA} (\texttt{Formalism.\allowbreak{}SharpRestriction.\allowbreak{}GlobalSchurBounds.\allowbreak{}TheoremThreeOfExplicitAndCompact}, lines 123--125)\\
\texttt{Formalism.\allowbreak{}MainNotions.\allowbreak{}GlobalTheoryTheory.\allowbreak{}JDB.\allowbreak{}dataA\_\allowbreak{}of} (\texttt{Formalism.\allowbreak{}SharpRestriction.\allowbreak{}GlobalSchurBounds.\allowbreak{}GlobalRankLatticeProofs.\allowbreak{}FamilyPartFiveOfExplicitConstants}, lines 299--302)\par}

{\footnotesize\raggedright\noindent \ref{B5.11}.\quad \texttt{Formalism.\allowbreak{}MainNotions.\allowbreak{}GlobalTheoryTheory.\allowbreak{}count\_\allowbreak{}at\_\allowbreak{}of} (\texttt{Formalism.\allowbreak{}SharpRestriction.\allowbreak{}GlobalSchurBounds.\allowbreak{}JointsL}, lines 20--143)\\
\texttt{Formalism.\allowbreak{}MainNotions.\allowbreak{}StandardLattice.\allowbreak{}division\_\allowbreak{}level\_\allowbreak{}prime\_\allowbreak{}family} (\texttt{Formalism.\allowbreak{}SharpRestriction.\allowbreak{}GlobalSchurBounds.\allowbreak{}LinkMin}, lines 425--454)\\
\texttt{Formalism.\allowbreak{}MainNotions.\allowbreak{}GlobalTheoryTheory.\allowbreak{}countF\_\allowbreak{}fam} (\texttt{Formalism.\allowbreak{}SharpRestriction.\allowbreak{}GlobalSchurBounds.\allowbreak{}JointsL}, lines 1152--1155)\par}

{\footnotesize\raggedright\noindent \ref{B5.11b}.\quad \texttt{Formalism.\allowbreak{}L2.\allowbreak{}SG.\allowbreak{}Pf.\allowbreak{}SzHypFam.\allowbreak{}sizes3b\_\allowbreak{}of\_\allowbreak{}ranks} (\texttt{Formalism.\allowbreak{}LocalToGlobal.\allowbreak{}SiegelLatticeTheory.\allowbreak{}FlagStabilizerTensors.\allowbreak{}RootLatticeThresholds.\allowbreak{}SzHypFam}, lines 78--93)\\
\texttt{Formalism.\allowbreak{}MainNotions.\allowbreak{}BridgeLemma.\allowbreak{}level\_\allowbreak{}ratio\_\allowbreak{}bound} (\texttt{Formalism.\allowbreak{}SharpRestriction.\allowbreak{}CuspSpanNormBounds.\allowbreak{}BruhatCellsKronecker.\allowbreak{}DegreeRatioFormulaStatement}, lines 398--524)\par}

{\footnotesize\raggedright\noindent \ref{B5.11a}.\quad \texttt{Formalism.\allowbreak{}L2.\allowbreak{}GlobalTheoryTheory.\allowbreak{}TraceIntegrality.\allowbreak{}bc\_\allowbreak{}mu} (\texttt{Formalism.\allowbreak{}LocalToGlobal.\allowbreak{}GlobalTheory.\allowbreak{}GL4.\allowbreak{}ColumnLatticeExactSlope}, lines 98--177)\\
\texttt{Formalism.\allowbreak{}L2.\allowbreak{}GlobalTheoryTheory.\allowbreak{}TraceIntegrality.\allowbreak{}eLat\_\allowbreak{}mu} (\texttt{Formalism.\allowbreak{}LocalToGlobal.\allowbreak{}GlobalTheory.\allowbreak{}GL4.\allowbreak{}MemDualCrdLat}, lines 71--128)\\
\texttt{Formalism.\allowbreak{}L2.\allowbreak{}GlobalTheoryTheory.\allowbreak{}TraceIntegrality.\allowbreak{}lower\_\allowbreak{}slope\_\allowbreak{}inequality\_\allowbreak{}family} (\texttt{Formalism.\allowbreak{}LocalToGlobal.\allowbreak{}GlobalTheory.\allowbreak{}GL4.\allowbreak{}LowerSlopeInequalityFamily}, lines 14--15)\\
\texttt{Formalism.\allowbreak{}L2.\allowbreak{}GlobalTheoryTheory.\allowbreak{}TraceIntegrality.\allowbreak{}slope\_\allowbreak{}inequality\_\allowbreak{}constant\_\allowbreak{}four} (\texttt{Formalism.\allowbreak{}LocalToGlobal.\allowbreak{}GlobalTheory.\allowbreak{}GL4.\allowbreak{}PStmtFree}, lines 14--15)\par}

{\footnotesize\raggedright\noindent \ref{B5.12}.\quad \texttt{Formalism.\allowbreak{}L2.\allowbreak{}GlobalTheoryTheory.\allowbreak{}TraceIntegrality.\allowbreak{}q\_\allowbreak{}lattice\_\allowbreak{}slope} (\texttt{Formalism.\allowbreak{}LocalToGlobal.\allowbreak{}GlobalTheory.\allowbreak{}GL4.\allowbreak{}ColumnLatticeExactSlope}, lines 183--203)\\
\texttt{Formalism.\allowbreak{}MainNotions.\allowbreak{}flipTwoThree.\allowbreak{}slopesMF\_\allowbreak{}of\_\allowbreak{}count\_\allowbreak{}fam} (\texttt{Formalism.\allowbreak{}SharpRestriction.\allowbreak{}GlobalSchurBounds.\allowbreak{}GlobalRankLatticeProofs.\allowbreak{}SlopesMFOfCountFam}, lines 25--33)\\
\texttt{Formalism.\allowbreak{}MainNotions.\allowbreak{}GlobalTheoryTheory.\allowbreak{}JSL.\allowbreak{}slopesM\_\allowbreak{}at\_\allowbreak{}of} (\texttt{Formalism.\allowbreak{}SharpRestriction.\allowbreak{}GlobalSchurBounds.\allowbreak{}JointsL}, lines 1518--1589)\par}

{\footnotesize\raggedright\noindent \ref{B5.12k}.\quad \texttt{Formalism.\allowbreak{}MainNotions.\allowbreak{}WeightDegreeSizes.\allowbreak{}SecondComplementInputs} (\texttt{Formalism.\allowbreak{}SharpRestriction.\allowbreak{}Divisibility.\allowbreak{}WeightDegreeSizes}, lines 673--675)\\
\texttt{Formalism.\allowbreak{}MainNotions.\allowbreak{}WeightDegreeSizes.\allowbreak{}FamilySecondComplementInputs} (\texttt{Formalism.\allowbreak{}SharpRestriction.\allowbreak{}Divisibility.\allowbreak{}WeightDegreeSizes}, lines 683--685)\\
\texttt{Formalism.\allowbreak{}MainNotions.\allowbreak{}WeightDegreeSizes.\allowbreak{}errG} (\texttt{Formalism.\allowbreak{}SharpRestriction.\allowbreak{}Divisibility.\allowbreak{}WeightDegreeSizes}, lines 408--411)\\
\texttt{Formalism.\allowbreak{}MainNotions.\allowbreak{}product\_\allowbreak{}first\_\allowbreak{}zeroth\_\allowbreak{}parts.\allowbreak{}WeightBudgetBound} (\texttt{Formalism.\allowbreak{}SharpRestriction.\allowbreak{}GlobalSchurBounds.\allowbreak{}ProductFirstZerothParts}, lines 324--328)\par}

{\footnotesize\raggedright\noindent \ref{B5.13}.\quad \texttt{Formalism.\allowbreak{}MainNotions.\allowbreak{}LocalQuadraticParts.\allowbreak{}adm11} (\texttt{Formalism.\allowbreak{}SharpRestriction.\allowbreak{}Divisibility.\allowbreak{}LocalQuadraticParts}, lines 316--322)\\
\texttt{Formalism.\allowbreak{}MainNotions.\allowbreak{}WeightDegreeSizes.\allowbreak{}secondComplement\_\allowbreak{}of\_\allowbreak{}inputs} (\texttt{Formalism.\allowbreak{}SharpRestriction.\allowbreak{}Divisibility.\allowbreak{}WeightDegreeSizes}, lines 677--681)\\
\texttt{Formalism.\allowbreak{}MainNotions.\allowbreak{}Lnk.\allowbreak{}Mod.\allowbreak{}data\_\allowbreak{}EComp5} (\texttt{Formalism.\allowbreak{}SharpRestriction.\allowbreak{}GlobalSchurBounds.\allowbreak{}QuadraticLogDerivCongr}, lines 594--622)\par}

{\footnotesize\raggedright\noindent \ref{B5.14}.\quad \texttt{Formalism.\allowbreak{}MainNotions.\allowbreak{}DIV.\allowbreak{}HookProducts.\allowbreak{}RBCompl3.\allowbreak{}compl3\_\allowbreak{}diff} (\texttt{Formalism.\allowbreak{}SharpRestriction.\allowbreak{}Divisibility.\allowbreak{}HookProducts.\allowbreak{}Compl3Record}, lines 15--22)\\
\texttt{Formalism.\allowbreak{}MainNotions.\allowbreak{}DIV.\allowbreak{}HookProducts.\allowbreak{}RBCompl3.\allowbreak{}compl3\_\allowbreak{}record} (\texttt{Formalism.\allowbreak{}SharpRestriction.\allowbreak{}Divisibility.\allowbreak{}HookProducts.\allowbreak{}Compl3Record}, lines 217--281)\\
\texttt{Formalism.\allowbreak{}MainNotions.\allowbreak{}Lnk.\allowbreak{}Mod.\allowbreak{}data\_\allowbreak{}HComp} (\texttt{Formalism.\allowbreak{}SharpRestriction.\allowbreak{}GlobalSchurBounds.\allowbreak{}QuadraticLogDerivCongr}, lines 624--639)\par}

\section*{Chapter~\ref{ch:prime-sums-net-gain}: Prime sums and net gain}

{\footnotesize\raggedright\noindent \ref{B5.15}.\quad \texttt{Formalism.\allowbreak{}LFunctionChainParameters.\allowbreak{}DeficitSplitLemmas.\allowbreak{}splitEq\_\allowbreak{}valid} (\texttt{Formalism.\allowbreak{}SizeWindowBrackets.\allowbreak{}ChainSlopeWindows.\allowbreak{}AdmissibleSchurLemma}, lines 63--68)\\
\texttt{Formalism.\allowbreak{}L2.\allowbreak{}GlobalTheoryTheory.\allowbreak{}Loss.\allowbreak{}Ga\_\allowbreak{}at} (\texttt{Formalism.\allowbreak{}LocalToGlobal.\allowbreak{}GlobalTheory.\allowbreak{}PowFourLtCeilThresholdPrime}, lines 210--245)\\
\texttt{Formalism.\allowbreak{}L2.\allowbreak{}GlobalTheoryTheory.\allowbreak{}Loss.\allowbreak{}ASK\_\allowbreak{}at} (\texttt{Formalism.\allowbreak{}LocalToGlobal.\allowbreak{}GlobalTheory.\allowbreak{}PowFourLtCeilThresholdPrime}, lines 386--395)\par}

{\footnotesize\raggedright\noindent \ref{B5.16}.\quad \texttt{Formalism.\allowbreak{}LFunctionChainParameters.\allowbreak{}DeficitSplitLemmas.\allowbreak{}split\_\allowbreak{}classes\_\allowbreak{}count} (\texttt{Formalism.\allowbreak{}SizeWindowBrackets.\allowbreak{}ChainSlopeWindows.\allowbreak{}AdmissibleSchurLemma}, lines 110--139)\par}

{\footnotesize\raggedright\noindent \ref{B5.16}.\quad \texttt{Formalism.\allowbreak{}L2.\allowbreak{}GlobalTheoryTheory.\allowbreak{}Loss.\allowbreak{}clAt} (\texttt{Formalism.\allowbreak{}LocalToGlobal.\allowbreak{}GlobalTheory.\allowbreak{}SuppSat}, lines 43--47)\par}

{\footnotesize\raggedright\noindent \ref{B5.17}.\quad \texttt{Formalism.\allowbreak{}LFunctionChainParameters.\allowbreak{}DeficitSplitLemmas.\allowbreak{}loss\_\allowbreak{}bound} (\texttt{Formalism.\allowbreak{}SizeWindowBrackets.\allowbreak{}ChainSlopeWindows.\allowbreak{}AdmissibleSchurLemma}, lines 762--918)\\
\texttt{Formalism.\allowbreak{}LFunctionChainParameters.\allowbreak{}DeficitSplitLemmas.\allowbreak{}SecondKindLoss} (\texttt{Formalism.\allowbreak{}SizeWindowBrackets.\allowbreak{}SecondKindBucket}, lines 261--293)\par}

{\footnotesize\raggedright\noindent \ref{B5.17a}.\quad \texttt{Formalism.\allowbreak{}L2.\allowbreak{}GlobalTheoryTheory.\allowbreak{}Loss.\allowbreak{}pd\_\allowbreak{}fill} (\texttt{Formalism.\allowbreak{}LocalToGlobal.\allowbreak{}GlobalTheory.\allowbreak{}PowFourLtCeilThresholdPrime}, lines 1025--1037)\\
\texttt{Formalism.\allowbreak{}L2.\allowbreak{}GlobalTheoryTheory.\allowbreak{}Loss.\allowbreak{}pd\_\allowbreak{}wz} (\texttt{Formalism.\allowbreak{}LocalToGlobal.\allowbreak{}GlobalTheory.\allowbreak{}PowFourLtCeilThresholdPrime}, lines 1038--1045)\par}

{\footnotesize\raggedright\noindent \ref{B5.18}.\quad \texttt{Formalism.\allowbreak{}LFunctionChainParameters.\allowbreak{}PrimeWindowSizeBounds.\allowbreak{}first\_\allowbreak{}kind\_\allowbreak{}bound} (\texttt{Formalism.\allowbreak{}SizeWindowBrackets.\allowbreak{}ChainSlopeWindows.\allowbreak{}OneLeThresholdMulPowers}, lines 955--986)\\
\texttt{Formalism.\allowbreak{}LFunctionChainParameters.\allowbreak{}PrimeWindowSizeBounds.\allowbreak{}first\_\allowbreak{}le} (\texttt{Formalism.\allowbreak{}SizeWindowBrackets.\allowbreak{}ChainSlopeWindows.\allowbreak{}OneLeThresholdMulPowers}, lines 1036--1100)\par}

{\footnotesize\raggedright\noindent \ref{B5.19}.\quad \texttt{Formalism.\allowbreak{}LFunctionChainParameters.\allowbreak{}SmallOrderPrimeLosses.\allowbreak{}out\_\allowbreak{}discriminant\_\allowbreak{}bound} (\texttt{Formalism.\allowbreak{}SizeWindowBrackets.\allowbreak{}ChainSlopeWindows.\allowbreak{}DiscriminantSettingSatisfied}, lines 32--74)\\
\texttt{Formalism.\allowbreak{}LFunctionChainParameters.\allowbreak{}SmallOrderPrimeLosses.\allowbreak{}second\_\allowbreak{}discriminant\_\allowbreak{}bound} (\texttt{Formalism.\allowbreak{}SizeWindowBrackets.\allowbreak{}ChainSlopeWindows.\allowbreak{}DiscriminantSettingSatisfied}, lines 76--94)\\
\texttt{Formalism.\allowbreak{}LFunctionChainParameters.\allowbreak{}SmallOrderPrimeLosses.\allowbreak{}io\_\allowbreak{}discriminant\_\allowbreak{}bound} (\texttt{Formalism.\allowbreak{}SizeWindowBrackets.\allowbreak{}ChainSlopeWindows.\allowbreak{}DiscriminantSettingSatisfied}, lines 96--112)\par}

{\footnotesize\raggedright\noindent \ref{B5.20}.\quad \texttt{Formalism.\allowbreak{}LFunctionChainParameters.\allowbreak{}SmallOrderPrimeLosses.\allowbreak{}joint\_\allowbreak{}core} (\texttt{Formalism.\allowbreak{}SizeWindowBrackets.\allowbreak{}ChainSlopeWindows.\allowbreak{}DiscriminantSettingSatisfied}, lines 178--226)\\
\texttt{Formalism.\allowbreak{}LFunctionChainParameters.\allowbreak{}SmallOrderPrimeLosses.\allowbreak{}joint\_\allowbreak{}discriminant\_\allowbreak{}bound} (\texttt{Formalism.\allowbreak{}SizeWindowBrackets.\allowbreak{}ChainSlopeWindows.\allowbreak{}DiscriminantSettingSatisfied}, lines 228--232)\\
\texttt{Formalism.\allowbreak{}LFunctionChainParameters.\allowbreak{}SmallOrderPrimeLosses.\allowbreak{}joint\_\allowbreak{}discriminant\_\allowbreak{}eventually} (\texttt{Formalism.\allowbreak{}SizeWindowBrackets.\allowbreak{}ChainSlopeWindows.\allowbreak{}DiscriminantSettingSatisfied}, lines 234--241)\par}

{\footnotesize\raggedright\noindent \ref{B5.20k}.\quad \texttt{Formalism.\allowbreak{}MainNotions.\allowbreak{}flipOneThree.\allowbreak{}Pf.\allowbreak{}KP.\allowbreak{}blockSplit} (\texttt{Formalism.\allowbreak{}SharpRestriction.\allowbreak{}GrowthBounds.\allowbreak{}ColumnSpanDensity.\allowbreak{}BlockSplit}, lines 85--135)\par}

{\footnotesize\raggedright\noindent \ref{B5.20k}.\quad \texttt{Formalism.\allowbreak{}MainNotions.\allowbreak{}G.\allowbreak{}LPprimeAt} (\texttt{Formalism.\allowbreak{}SharpRestriction.\allowbreak{}GrowthBounds.\allowbreak{}UnramifiedQuadraticSetting}, lines 205--214)\\
\texttt{Formalism.\allowbreak{}MainNotions.\allowbreak{}G.\allowbreak{}LPprime} (\texttt{Formalism.\allowbreak{}SharpRestriction.\allowbreak{}GrowthBounds.\allowbreak{}UnramifiedQuadraticSetting}, lines 216--220)\par}

{\footnotesize\raggedright\noindent \ref{B5.20d}.\quad \texttt{Formalism.\allowbreak{}CountingSumsFrames.\allowbreak{}Top.\allowbreak{}localizerVanishes} (\texttt{Formalism.\allowbreak{}PolynomialDimension.\allowbreak{}DetMulEqSumAdjugate}, lines 55--103)\\
\texttt{Formalism.\allowbreak{}MainNotions.\allowbreak{}G.\allowbreak{}WittFrameLocaliserOrders.\allowbreak{}localiser\_\allowbreak{}vanishing\_\allowbreak{}of\_\allowbreak{}distinct\_\allowbreak{}planes} (\texttt{Formalism.\allowbreak{}SharpRestriction.\allowbreak{}GrowthBounds.\allowbreak{}WittFrameLocaliserOrders}, lines 570--685)\par}

{\footnotesize\raggedright\noindent \ref{B5.20e}.\quad \texttt{Formalism.\allowbreak{}MainNotions.\allowbreak{}Pf14.\allowbreak{}goodRank\_\allowbreak{}of\_\allowbreak{}prim} (\texttt{Formalism.\allowbreak{}SharpRestriction.\allowbreak{}GrowthBounds.\allowbreak{}ScaleEdges.\allowbreak{}GoodRankOfPrim}, lines 29--75)\\
\texttt{Formalism.\allowbreak{}MainNotions.\allowbreak{}Pf14.\allowbreak{}s7} (\texttt{Formalism.\allowbreak{}SharpRestriction.\allowbreak{}GrowthBounds.\allowbreak{}ScaleEdges.\allowbreak{}GoodRankOfPrim}, lines 102--170)\par}

{\footnotesize\raggedright\noindent \ref{B5.20f}.\quad \texttt{Formalism.\allowbreak{}GM.\allowbreak{}Top.\allowbreak{}windowFullAP} (\texttt{Formalism.\allowbreak{}MahlerMeasure.\allowbreak{}CoarseLevelSlackInequalities}, lines 10--10)\\
\texttt{Formalism.\allowbreak{}MainNotions.\allowbreak{}G.\allowbreak{}GramWindowAlt.\allowbreak{}gramWindow2\_\allowbreak{}proved} (\texttt{Formalism.\allowbreak{}SharpRestriction.\allowbreak{}GrowthBounds.\allowbreak{}FockGramWindow.\allowbreak{}GramWindow2Proved}, lines 138--236)\par}

{\footnotesize\raggedright\noindent \ref{B5.20g}.\quad \texttt{Formalism.\allowbreak{}MainNotions.\allowbreak{}flipOneThree.\allowbreak{}Pf.\allowbreak{}avg\_\allowbreak{}abs} (\texttt{Formalism.\allowbreak{}SharpRestriction.\allowbreak{}GrowthBounds.\allowbreak{}ColumnSpanDensity.\allowbreak{}DensityNumericalFacts}, lines 80--160)\\
\texttt{Formalism.\allowbreak{}MainNotions.\allowbreak{}flipOneThree.\allowbreak{}Pf.\allowbreak{}mean\_\allowbreak{}of} (\texttt{Formalism.\allowbreak{}SharpRestriction.\allowbreak{}GrowthBounds.\allowbreak{}ColumnSpanDensity.\allowbreak{}ColumnSpanNumericalFacts}, lines 218--300)\par}

{\footnotesize\raggedright\noindent \ref{B5.20l}.\quad \texttt{Formalism.\allowbreak{}MainNotions.\allowbreak{}G.\allowbreak{}PPath} (\texttt{Formalism.\allowbreak{}SharpRestriction.\allowbreak{}GrowthBounds.\allowbreak{}ObjPol}, lines 300--304)\\
\texttt{Formalism.\allowbreak{}MainNotions.\allowbreak{}G.\allowbreak{}IsSt} (\texttt{Formalism.\allowbreak{}SharpRestriction.\allowbreak{}GrowthBounds.\allowbreak{}ObjPol}, lines 460--463)\\
\texttt{Formalism.\allowbreak{}MainNotions.\allowbreak{}SelectionLemmas.\allowbreak{}Pf.\allowbreak{}Dat.\allowbreak{}Tset} (\texttt{Formalism.\allowbreak{}SharpRestriction.\allowbreak{}GrowthBounds.\allowbreak{}SegmentDegree.\allowbreak{}GoodPathWeight}, lines 41--43)\par}

{\footnotesize\raggedright\noindent \ref{B5.20m}.\quad \texttt{Formalism.\allowbreak{}InvariantFactorPaths.\allowbreak{}BranchingDimensions.\allowbreak{}pathWeightBound\_\allowbreak{}of} (\texttt{Formalism.\allowbreak{}ThetaLogCoordinates.\allowbreak{}ResidueClassEstimates.\allowbreak{}VandermondeRankEstimates.\allowbreak{}AtypicalWeightBoundOf}, lines 353--677)\\
\texttt{Formalism.\allowbreak{}MainNotions.\allowbreak{}SelectionLemmas.\allowbreak{}Pf.\allowbreak{}T.\allowbreak{}path\_\allowbreak{}count\_\allowbreak{}finite} (\texttt{Formalism.\allowbreak{}SharpRestriction.\allowbreak{}GrowthBounds.\allowbreak{}SegmentDegree.\allowbreak{}WindowSmall}, lines 87--178)\\
\texttt{Formalism.\allowbreak{}MainNotions.\allowbreak{}SelectionLemmas.\allowbreak{}Pf.\allowbreak{}T.\allowbreak{}lemT\_\allowbreak{}of} (\texttt{Formalism.\allowbreak{}SharpRestriction.\allowbreak{}GrowthBounds.\allowbreak{}SegmentDegree.\allowbreak{}AtypicalPathSmallWeight}, lines 112--203)\par}

{\footnotesize\raggedright\noindent \ref{B5.20n}.\quad \texttt{Formalism.\allowbreak{}MainNotions.\allowbreak{}SelectionLemmas.\allowbreak{}Pf.\allowbreak{}residue\_\allowbreak{}class\_\allowbreak{}count} (\texttt{Formalism.\allowbreak{}SharpRestriction.\allowbreak{}GrowthBounds.\allowbreak{}SegmentDegree.\allowbreak{}DminusLeBeS}, lines 10--42)\\
\texttt{Formalism.\allowbreak{}MainNotions.\allowbreak{}SelectionLemmas.\allowbreak{}Pf.\allowbreak{}restMon\_\allowbreak{}of} (\texttt{Formalism.\allowbreak{}SharpRestriction.\allowbreak{}GrowthBounds.\allowbreak{}SegmentDegree.\allowbreak{}TauDegreeOnPathBound}, lines 94--265)\par}

{\footnotesize\raggedright\noindent \ref{B5.20ha}.\quad \texttt{Formalism.\allowbreak{}MainNotions.\allowbreak{}G.\allowbreak{}GLWeightTheory.\allowbreak{}unit\_\allowbreak{}function\_\allowbreak{}mahler\_\allowbreak{}decay\_\allowbreak{}of} (\texttt{Formalism.\allowbreak{}SharpRestriction.\allowbreak{}GrowthBounds.\allowbreak{}RootWeights.\allowbreak{}MemBoundedDegreeExponentsIff}, lines 314--352)\\
\texttt{Formalism.\allowbreak{}MainNotions.\allowbreak{}Proposition5Numerics.\allowbreak{}digQ} (\texttt{Formalism.\allowbreak{}SharpRestriction.\allowbreak{}GrowthBounds.\allowbreak{}DiagonalSupport.\allowbreak{}BinomCombFOnFinset}, lines 22--105)\par}

{\footnotesize\raggedright\noindent \ref{B5.20he}.\quad \texttt{Formalism.\allowbreak{}MainNotions.\allowbreak{}CovariantDegrees.\allowbreak{}CanonicalTranslateJetIdentity} (\texttt{Formalism.\allowbreak{}SharpRestriction.\allowbreak{}GrowthBounds.\allowbreak{}StatementOfDegreeAndBound}, lines 417--436)\\
\texttt{Formalism.\allowbreak{}MainNotions.\allowbreak{}CovariantDegrees.\allowbreak{}CoefficientDependence} (\texttt{Formalism.\allowbreak{}SharpRestriction.\allowbreak{}GrowthBounds.\allowbreak{}StatementOfDegreeAndBound}, lines 407--416)\\
\texttt{Formalism.\allowbreak{}MainNotions.\allowbreak{}G.\allowbreak{}WittFrameLocaliserOrders.\allowbreak{}subst\_\allowbreak{}twisted\_\allowbreak{}row\_\allowbreak{}even} (\texttt{Formalism.\allowbreak{}SharpRestriction.\allowbreak{}GrowthBounds.\allowbreak{}WittFrameLocaliserOrders}, lines 483--507)\\
\texttt{Formalism.\allowbreak{}MainNotions.\allowbreak{}G.\allowbreak{}WittFrameLocaliserOrders.\allowbreak{}residue\_\allowbreak{}integral\_\allowbreak{}primal\_\allowbreak{}row} (\texttt{Formalism.\allowbreak{}SharpRestriction.\allowbreak{}GrowthBounds.\allowbreak{}WittFrameLocaliserOrders}, lines 508--532)\\
\texttt{Formalism.\allowbreak{}MainNotions.\allowbreak{}G.\allowbreak{}WittFrameLocaliserOrders.\allowbreak{}residue\_\allowbreak{}zero\_\allowbreak{}mem\_\allowbreak{}order\_\allowbreak{}one} (\texttt{Formalism.\allowbreak{}SharpRestriction.\allowbreak{}GrowthBounds.\allowbreak{}WittFrameLocaliserOrders}, lines 446--457)\par}

{\footnotesize\raggedright\noindent \ref{B5.20h}.\quad \texttt{Formalism.\allowbreak{}MainNotions.\allowbreak{}G.\allowbreak{}WittFrameLocaliserOrders.\allowbreak{}localiser\_\allowbreak{}vanishing\_\allowbreak{}of\_\allowbreak{}distinct\_\allowbreak{}planes} (\texttt{Formalism.\allowbreak{}SharpRestriction.\allowbreak{}GrowthBounds.\allowbreak{}WittFrameLocaliserOrders}, lines 570--685)\\
\texttt{Formalism.\allowbreak{}MainNotions.\allowbreak{}G.\allowbreak{}second\_\allowbreak{}theorem\_\allowbreak{}closed\_\allowbreak{}everywhere} (\texttt{Formalism.\allowbreak{}SharpRestriction.\allowbreak{}GrowthBounds.\allowbreak{}PlaneCountTheoremWithConstants}, lines 1121--1135)\\
\texttt{Formalism.\allowbreak{}MainNotions.\allowbreak{}G.\allowbreak{}second\_\allowbreak{}theorem\_\allowbreak{}absolute\_\allowbreak{}constant} (\texttt{Formalism.\allowbreak{}SharpRestriction.\allowbreak{}GrowthBounds.\allowbreak{}PlaneCountTheoremWithConstants}, lines 1137--1142)\par}

{\footnotesize\raggedright\noindent \ref{B5.20hb}.\quad \texttt{Formalism.\allowbreak{}MainNotions.\allowbreak{}PairMembershipBlocks.\allowbreak{}restMEM\_\allowbreak{}of\_\allowbreak{}last} (\texttt{Formalism.\allowbreak{}SharpRestriction.\allowbreak{}GrowthBounds.\allowbreak{}PlaneCoefficients}, lines 34--40)\\
\texttt{Formalism.\allowbreak{}MainNotions.\allowbreak{}Proposition5Numerics.\allowbreak{}pick\_\allowbreak{}core} (\texttt{Formalism.\allowbreak{}SharpRestriction.\allowbreak{}GrowthBounds.\allowbreak{}DiagonalSupport.\allowbreak{}PickCore}, lines 23--166)\\
\texttt{Formalism.\allowbreak{}MainNotions.\allowbreak{}PlaneGram.\allowbreak{}stepB\_\allowbreak{}of} (\texttt{Formalism.\allowbreak{}SharpRestriction.\allowbreak{}GrowthBounds.\allowbreak{}MonomialOrderSegments.\allowbreak{}CardLinkIdxMulLe}, lines 127--233)\par}

{\footnotesize\raggedright\noindent \ref{B5.20hc}.\quad \texttt{Formalism.\allowbreak{}MainNotions.\allowbreak{}PlaneGram.\allowbreak{}segStep\_\allowbreak{}of} (\texttt{Formalism.\allowbreak{}SharpRestriction.\allowbreak{}GrowthBounds.\allowbreak{}MonomialOrderSegments.\allowbreak{}RenamePullLCoefOut}, lines 150--197)\\
\texttt{Formalism.\allowbreak{}MainNotions.\allowbreak{}PlaneGram.\allowbreak{}stepC\_\allowbreak{}of} (\texttt{Formalism.\allowbreak{}SharpRestriction.\allowbreak{}GrowthBounds.\allowbreak{}MonomialOrderSegments.\allowbreak{}ExistsOfMemSupportLin}, lines 409--426)\par}

{\footnotesize\raggedright\noindent \ref{B5.20hd}.\quad \texttt{Formalism.\allowbreak{}MainNotions.\allowbreak{}PlaneGram.\allowbreak{}stepE\_\allowbreak{}of} (\texttt{Formalism.\allowbreak{}SharpRestriction.\allowbreak{}GrowthBounds.\allowbreak{}MonomialOrderSegments.\allowbreak{}SumOnLetterApplyInl}, lines 310--335)\\
\texttt{Formalism.\allowbreak{}MainNotions.\allowbreak{}PlaneGram.\allowbreak{}TwistData.\allowbreak{}value\_\allowbreak{}monomial} (\texttt{Formalism.\allowbreak{}SharpRestriction.\allowbreak{}GrowthBounds.\allowbreak{}MonomialOrderSegments.\allowbreak{}SubstitutedExponentApplyLabel}, lines 237--327)\\
\texttt{Formalism.\allowbreak{}MainNotions.\allowbreak{}PlaneGram.\allowbreak{}omega\_\allowbreak{}value} (\texttt{Formalism.\allowbreak{}SharpRestriction.\allowbreak{}GrowthBounds.\allowbreak{}MonomialOrderSegments.\allowbreak{}OmegaValue}, lines 12--106)\par}

{\footnotesize\raggedright\noindent \ref{B5.20o}.\quad \texttt{Formalism.\allowbreak{}GM.\allowbreak{}M.\allowbreak{}MMwZ} (\texttt{Formalism.\allowbreak{}MahlerMeasure.\allowbreak{}CsSaturated}, lines 50--52)\\
\texttt{Formalism.\allowbreak{}GM.\allowbreak{}M.\allowbreak{}FilE} (\texttt{Formalism.\allowbreak{}MahlerMeasure.\allowbreak{}CsSaturated}, lines 54--57)\\
\texttt{Formalism.\allowbreak{}GM.\allowbreak{}M.\allowbreak{}grLat} (\texttt{Formalism.\allowbreak{}MahlerMeasure.\allowbreak{}CsSaturated}, lines 59--62)\par}

{\footnotesize\raggedright\noindent \ref{B5.20t}.\quad \texttt{Formalism.\allowbreak{}InvariantFactorPaths.\allowbreak{}PW.\allowbreak{}Top.\allowbreak{}typeInt} (\texttt{Formalism.\allowbreak{}ThetaLogCoordinates.\allowbreak{}TraceIsometryWeights.\allowbreak{}FinalP}, lines 22--22)\\
\texttt{Formalism.\allowbreak{}InvariantFactorPaths.\allowbreak{}PW.\allowbreak{}Top.\allowbreak{}weylInt} (\texttt{Formalism.\allowbreak{}ThetaLogCoordinates.\allowbreak{}TraceIsometryWeights.\allowbreak{}DischargeHeightBound}, lines 8--8)\par}

{\footnotesize\raggedright\noindent \ref{B5.20u}.\quad \texttt{Formalism.\allowbreak{}InvariantFactorPaths.\allowbreak{}L4.\allowbreak{}restriction\_\allowbreak{}into\_\allowbreak{}holds} (\texttt{Formalism.\allowbreak{}ThetaLogCoordinates.\allowbreak{}FockResolutionOnto.\allowbreak{}FockMonomialProofs.\allowbreak{}ResInto}, lines 430--473)\\
\texttt{Formalism.\allowbreak{}InvariantFactorPaths.\allowbreak{}L4.\allowbreak{}restriction\_\allowbreak{}surjective} (\texttt{Formalism.\allowbreak{}ThetaLogCoordinates.\allowbreak{}FockResolutionOnto.\allowbreak{}FockMonomialProofs.\allowbreak{}DerivativeKeepsSpaceHolds}, lines 208--298)\\
\texttt{Formalism.\allowbreak{}InvariantFactorPaths.\allowbreak{}L4.\allowbreak{}restricted\_\allowbreak{}saturation\_\allowbreak{}holds} (\texttt{Formalism.\allowbreak{}ThetaLogCoordinates.\allowbreak{}FockResolutionOnto.\allowbreak{}FockMonomialProofs.\allowbreak{}RestrictedSaturationHolds}, lines 91--132)\par}

{\footnotesize\raggedright\noindent \ref{B5.20v}.\quad \texttt{Formalism.\allowbreak{}InvariantFactorPaths.\allowbreak{}L4.\allowbreak{}grSpan\_\allowbreak{}eq} (\texttt{Formalism.\allowbreak{}ThetaLogCoordinates.\allowbreak{}FockResolutionOnto.\allowbreak{}FockMonomialProofs.\allowbreak{}GramBoundOfRestrictionOfProj}, lines 200--238)\\
\texttt{Formalism.\allowbreak{}InvariantFactorPaths.\allowbreak{}L4.\allowbreak{}satAtL} (\texttt{Formalism.\allowbreak{}ThetaLogCoordinates.\allowbreak{}FockResolutionOnto.\allowbreak{}FockMonomialProofs.\allowbreak{}GramBoundUnconditionalOfSat}, lines 90--105)\\
\texttt{Formalism.\allowbreak{}InvariantFactorPaths.\allowbreak{}L4.\allowbreak{}graded\_\allowbreak{}bound\_\allowbreak{}pos\_\allowbreak{}of\_\allowbreak{}proj} (\texttt{Formalism.\allowbreak{}ThetaLogCoordinates.\allowbreak{}FockResolutionOnto.\allowbreak{}FockMonomialProofs.\allowbreak{}GramBoundPosOfProjSat}, lines 51--68)\par}

{\footnotesize\raggedright\noindent \ref{B5.20p}.\quad \texttt{Formalism.\allowbreak{}InvariantFactorPaths.\allowbreak{}L4.\allowbreak{}graded\_\allowbreak{}piece\_\allowbreak{}bound\_\allowbreak{}holds} (\texttt{Formalism.\allowbreak{}ThetaLogCoordinates.\allowbreak{}FockResolutionOnto.\allowbreak{}FockMonomialProofs.\allowbreak{}PartnerDeterminantHolds}, lines 114--198)\\
\texttt{Formalism.\allowbreak{}InvariantFactorPaths.\allowbreak{}L4.\allowbreak{}partner\_\allowbreak{}determinant\_\allowbreak{}holds} (\texttt{Formalism.\allowbreak{}ThetaLogCoordinates.\allowbreak{}FockResolutionOnto.\allowbreak{}FockMonomialProofs.\allowbreak{}PartnerDeterminantHolds}, lines 10--77)\par}

{\footnotesize\raggedright\noindent \ref{B5.20q}.\quad \texttt{Formalism.\allowbreak{}GM.\allowbreak{}M.\allowbreak{}tieGam} (\texttt{Formalism.\allowbreak{}MahlerMeasure.\allowbreak{}HasOrder}, lines 14--19)\\
\texttt{Formalism.\allowbreak{}GM.\allowbreak{}gamma\_\allowbreak{}degree\_\allowbreak{}bound} (\texttt{Formalism.\allowbreak{}MahlerMeasure.\allowbreak{}MahlerMeasureProofs.\allowbreak{}GamDeg}, lines 1153--1171)\\
\texttt{Formalism.\allowbreak{}GM.\allowbreak{}gamma\_\allowbreak{}leading\_\allowbreak{}term} (\texttt{Formalism.\allowbreak{}MahlerMeasure.\allowbreak{}MahlerMeasureProofs.\allowbreak{}GamLead}, lines 1852--1947)\\
\texttt{Formalism.\allowbreak{}GM.\allowbreak{}MeanGammaLemmas.\allowbreak{}lead\_\allowbreak{}part} (\texttt{Formalism.\allowbreak{}MahlerMeasure.\allowbreak{}MahlerMeasureProofs.\allowbreak{}MeanGam}, lines 400--539)\par}

{\footnotesize\raggedright\noindent \ref{B5.20r}.\quad \texttt{Formalism.\allowbreak{}GM.\allowbreak{}DimensionMonotonicity.\allowbreak{}highest\_\allowbreak{}weight\_\allowbreak{}dim\_\allowbreak{}monotone} (\texttt{Formalism.\allowbreak{}MahlerMeasure.\allowbreak{}MahlerMeasureProofs.\allowbreak{}WeightSpaceDimensionMonotone}, lines 2606--2856)\par}

{\footnotesize\raggedright\noindent \ref{B5.20s}.\quad \texttt{Formalism.\allowbreak{}MainNotions.\allowbreak{}flipOneThree.\allowbreak{}Pf.\allowbreak{}contraction\_\allowbreak{}of\_\allowbreak{}pairing\_\allowbreak{}lemmas} (\texttt{Formalism.\allowbreak{}SharpRestriction.\allowbreak{}GrowthBounds.\allowbreak{}HermitianTensor.\allowbreak{}ContractionOfPairingLemmas}, lines 9--173)\\
\texttt{Formalism.\allowbreak{}MainNotions.\allowbreak{}flipOneThree.\allowbreak{}Pf.\allowbreak{}cloudOne\_\allowbreak{}of} (\texttt{Formalism.\allowbreak{}SharpRestriction.\allowbreak{}GrowthBounds.\allowbreak{}HermitianTensor.\allowbreak{}CloudA}, lines 9--97)\\
\texttt{Formalism.\allowbreak{}MainNotions.\allowbreak{}flipOneThree.\allowbreak{}Pf.\allowbreak{}iso\_\allowbreak{}one} (\texttt{Formalism.\allowbreak{}SharpRestriction.\allowbreak{}GrowthBounds.\allowbreak{}HermitianTensor.\allowbreak{}CloudIso}, lines 33--59)\par}

{\footnotesize\raggedright\noindent \ref{B5.20w}.\quad \texttt{Formalism.\allowbreak{}CountingSumsFrames.\allowbreak{}TO.\allowbreak{}tie\_\allowbreak{}core} (\texttt{Formalism.\allowbreak{}PolynomialDimension.\allowbreak{}MultinomialDepth.\allowbreak{}TieOps}, lines 3104--3287)\\
\texttt{Formalism.\allowbreak{}InvariantFactorPaths.\allowbreak{}MC.\allowbreak{}partnerMX\_\allowbreak{}of\_\allowbreak{}parts} (\texttt{Formalism.\allowbreak{}ThetaLogCoordinates.\allowbreak{}PartnerMeanCoordinates.\allowbreak{}PartnerMeanProofs.\allowbreak{}GammaZeroIntegral}, lines 253--376)\\
\texttt{Formalism.\allowbreak{}InvariantFactorPaths.\allowbreak{}MC.\allowbreak{}pointMX\_\allowbreak{}of\_\allowbreak{}partner} (\texttt{Formalism.\allowbreak{}ThetaLogCoordinates.\allowbreak{}PartnerMeanCoordinates.\allowbreak{}PartnerMeanProofs.\allowbreak{}GammaZeroIntegral}, lines 378--398)\par}

{\footnotesize\raggedright\noindent \ref{B5.20x}.\quad \texttt{Formalism.\allowbreak{}MainNotions.\allowbreak{}ColumnSystems.\allowbreak{}Pf.\allowbreak{}loc4One} (\texttt{Formalism.\allowbreak{}SharpRestriction.\allowbreak{}GrowthBounds.\allowbreak{}CoarseMultiplicity.\allowbreak{}KeyOne}, lines 101--128)\\
\texttt{Formalism.\allowbreak{}MainNotions.\allowbreak{}ColumnSystems.\allowbreak{}Pf.\allowbreak{}loc4Res\_\allowbreak{}of} (\texttt{Formalism.\allowbreak{}SharpRestriction.\allowbreak{}GrowthBounds.\allowbreak{}CoarseMultiplicity.\allowbreak{}Loc4ResOfLine}, lines 33--52)\\
\texttt{Formalism.\allowbreak{}MainNotions.\allowbreak{}ColumnSystems.\allowbreak{}Pf.\allowbreak{}loc4Ord\_\allowbreak{}of} (\texttt{Formalism.\allowbreak{}SharpRestriction.\allowbreak{}GrowthBounds.\allowbreak{}CoarseMultiplicity.\allowbreak{}LiftSubstStart}, lines 43--50)\par}

{\footnotesize\raggedright\noindent \ref{B5.20ia}.\quad \texttt{Formalism.\allowbreak{}MainNotions.\allowbreak{}ColumnSystems.\allowbreak{}Pf.\allowbreak{}SegHyp.\allowbreak{}seg\_\allowbreak{}ord} (\texttt{Formalism.\allowbreak{}SharpRestriction.\allowbreak{}GrowthBounds.\allowbreak{}BinomialRatio.\allowbreak{}TopDegreeLeTwoB}, lines 436--514)\\
\texttt{Formalism.\allowbreak{}MainNotions.\allowbreak{}ColumnSystems.\allowbreak{}Pf.\allowbreak{}ord\_\allowbreak{}all} (\texttt{Formalism.\allowbreak{}SharpRestriction.\allowbreak{}GrowthBounds.\allowbreak{}BinomialRatio.\allowbreak{}TopDegreeLeTwoB}, lines 570--628)\par}

{\footnotesize\raggedright\noindent \ref{B5.20ib}.\quad \texttt{Formalism.\allowbreak{}InvariantFactorPaths.\allowbreak{}MC.\allowbreak{}gamma\_\allowbreak{}integral\_\allowbreak{}holds} (\texttt{Formalism.\allowbreak{}ThetaLogCoordinates.\allowbreak{}PartnerMeanCoordinates.\allowbreak{}PartnerMeanProofs.\allowbreak{}GamIntC}, lines 302--365)\\
\texttt{Formalism.\allowbreak{}GM.\allowbreak{}gamma\_\allowbreak{}degree\_\allowbreak{}bound} (\texttt{Formalism.\allowbreak{}MahlerMeasure.\allowbreak{}MahlerMeasureProofs.\allowbreak{}GamDeg}, lines 1153--1171)\par}

{\footnotesize\raggedright\noindent \ref{B5.20ic}.\quad \texttt{Formalism.\allowbreak{}MainNotions.\allowbreak{}ColumnSystems.\allowbreak{}Pf.\allowbreak{}cth\_\allowbreak{}SSt} (\texttt{Formalism.\allowbreak{}SharpRestriction.\allowbreak{}GrowthBounds.\allowbreak{}CoarseMultiplicity.\allowbreak{}OtherEFree}, lines 172--186)\\
\texttt{Formalism.\allowbreak{}MainNotions.\allowbreak{}ColumnSystems.\allowbreak{}Pf.\allowbreak{}op\_\allowbreak{}SSt} (\texttt{Formalism.\allowbreak{}SharpRestriction.\allowbreak{}GrowthBounds.\allowbreak{}CoarseMultiplicity.\allowbreak{}GTheta}, lines 27--58)\\
\texttt{Formalism.\allowbreak{}MainNotions.\allowbreak{}ColumnSystems.\allowbreak{}Pf.\allowbreak{}sat\_\allowbreak{}block} (\texttt{Formalism.\allowbreak{}SharpRestriction.\allowbreak{}GrowthBounds.\allowbreak{}CoarseMultiplicity.\allowbreak{}SatBlock}, lines 102--155)\par}

{\footnotesize\raggedright\noindent \ref{B5.20i}.\quad \texttt{GlobalTheory.\allowbreak{}EX.\allowbreak{}RangeOne.\allowbreak{}rangeOne\_\allowbreak{}abs} (\texttt{Development.\allowbreak{}AbelianVarietyTheory.\allowbreak{}Proofs.\allowbreak{}GlobalInterpolation.\allowbreak{}TensorOfPairLift}, lines 637--640)\\
\texttt{Formalism.\allowbreak{}MainNotions.\allowbreak{}flipOneThree.\allowbreak{}Conclusion2} (\texttt{Formalism.\allowbreak{}SharpRestriction.\allowbreak{}GrowthBounds.\allowbreak{}ProjectionLemmaBaseLevel}, lines 78--91)\par}

{\footnotesize\raggedright\noindent \ref{B5.20aa}.\quad \texttt{Formalism.\allowbreak{}MainNotions.\allowbreak{}AuxBlock8.\allowbreak{}Pf.\allowbreak{}St.\allowbreak{}starAt\_\allowbreak{}D} (\texttt{Formalism.\allowbreak{}SharpRestriction.\allowbreak{}GrowthBounds.\allowbreak{}LogarithmicSums.\allowbreak{}PartnerCountOfCentral}, lines 597--689)\\
\texttt{Formalism.\allowbreak{}MainNotions.\allowbreak{}AuxBlock8.\allowbreak{}Pf.\allowbreak{}St.\allowbreak{}starAt\_\allowbreak{}E} (\texttt{Formalism.\allowbreak{}SharpRestriction.\allowbreak{}GrowthBounds.\allowbreak{}LogarithmicSums.\allowbreak{}PartnerCountOfCentral}, lines 696--841)\\
\texttt{Formalism.\allowbreak{}MainNotions.\allowbreak{}AuxBlock8.\allowbreak{}Pf.\allowbreak{}St.\allowbreak{}star\_\allowbreak{}fam} (\texttt{Formalism.\allowbreak{}SharpRestriction.\allowbreak{}GrowthBounds.\allowbreak{}LogarithmicSums.\allowbreak{}StarFam}, lines 11--37)\par}

{\footnotesize\raggedright\noindent \ref{B5.20a}.\quad \texttt{GlobalTheory.\allowbreak{}Top.\allowbreak{}closed\_\allowbreak{}linkG\_\allowbreak{}of\_\allowbreak{}duty} (\texttt{Development.\allowbreak{}AutomorphicArithmetic.\allowbreak{}ThetaCorrespondence.\allowbreak{}ThetaPairLifts.\allowbreak{}ThirdPlaceSeries.\allowbreak{}LinkG}, lines 13--24)\\
\texttt{Formalism.\allowbreak{}MainNotions.\allowbreak{}LinkG\_\allowbreak{}at} (\texttt{Formalism.\allowbreak{}SharpRestriction.\allowbreak{}GlobalTheoremInputs}, lines 122--136)\par}

{\footnotesize\raggedright\noindent \ref{B5.20a}.\quad \texttt{Formalism.\allowbreak{}MainNotions.\allowbreak{}G.\allowbreak{}Sv} (\texttt{Formalism.\allowbreak{}SharpRestriction.\allowbreak{}GrowthBounds.\allowbreak{}IsResidueClassZero}, lines 336--337)\\
\texttt{Formalism.\allowbreak{}MainNotions.\allowbreak{}G.\allowbreak{}Joints.\allowbreak{}linear\_\allowbreak{}programming\_\allowbreak{}input\_\allowbreak{}of\_\allowbreak{}hypotheses} (\texttt{Formalism.\allowbreak{}SharpRestriction.\allowbreak{}GrowthBounds.\allowbreak{}Joints}, lines 2040--2128)\\
\texttt{Formalism.\allowbreak{}MainNotions.\allowbreak{}AuxBlock8.\allowbreak{}Pf.\allowbreak{}gain\_\allowbreak{}bound\_\allowbreak{}at\_\allowbreak{}of\_\allowbreak{}star} (\texttt{Formalism.\allowbreak{}SharpRestriction.\allowbreak{}GrowthBounds.\allowbreak{}LogarithmicSums.\allowbreak{}NotSingularOfCentral}, lines 619--729)\\
\texttt{Formalism.\allowbreak{}L2.\allowbreak{}GlobalTheoryTheory.\allowbreak{}Loss.\allowbreak{}loss\_\allowbreak{}bound\_\allowbreak{}at} (\texttt{Formalism.\allowbreak{}LocalToGlobal.\allowbreak{}GlobalTheory.\allowbreak{}PowFourLtCeilThresholdPrime}, lines 766--817)\par}

{\footnotesize\raggedright\noindent \ref{B5.20y}.\quad \texttt{Formalism.\allowbreak{}MainNotions.\allowbreak{}G.\allowbreak{}first\_\allowbreak{}theorem\_\allowbreak{}closed\_\allowbreak{}form} (\texttt{Formalism.\allowbreak{}SharpRestriction.\allowbreak{}GrowthBounds.\allowbreak{}PlaneCountTheoremWithConstants}, lines 1761--1781)\\
\texttt{Formalism.\allowbreak{}MainNotions.\allowbreak{}G.\allowbreak{}signs\_\allowbreak{}closed} (\texttt{Formalism.\allowbreak{}SharpRestriction.\allowbreak{}GrowthBounds.\allowbreak{}PlaneCountTheoremWithConstants}, lines 1783--1789)\\
\texttt{Formalism.\allowbreak{}MainNotions.\allowbreak{}G.\allowbreak{}TopChoiceBounds.\allowbreak{}family\_\allowbreak{}log\_\allowbreak{}bound\_\allowbreak{}closed} (\texttt{Formalism.\allowbreak{}SharpRestriction.\allowbreak{}GrowthBounds.\allowbreak{}TopChoiceBounds}, lines 4428--4432)\par}

{\footnotesize\raggedright\noindent \ref{B5.20j}.\quad \texttt{Formalism.\allowbreak{}MainNotions.\allowbreak{}G.\allowbreak{}Joints.\allowbreak{}hypIIAt\_\allowbreak{}ev} (\texttt{Formalism.\allowbreak{}SharpRestriction.\allowbreak{}GrowthBounds.\allowbreak{}SizesAtRecord}, lines 567--626)\\
\texttt{UpperTheory.\allowbreak{}even\_\allowbreak{}addOrderOf\_\allowbreak{}hP\_\allowbreak{}of\_\allowbreak{}pow\_\allowbreak{}le} (\texttt{Development.\allowbreak{}AutomorphicArithmetic.\allowbreak{}OrbitLabellings.\allowbreak{}EvenAddOrderOfHPOfPowLe}, lines 8--28)\par}

{\footnotesize\raggedright\noindent \ref{B5.20b}.\quad \texttt{Formalism.\allowbreak{}MainNotions.\allowbreak{}G.\allowbreak{}Joints.\allowbreak{}type\_\allowbreak{}two\_\allowbreak{}hypothesis\_\allowbreak{}of\_\allowbreak{}counts} (\texttt{Formalism.\allowbreak{}SharpRestriction.\allowbreak{}GrowthBounds.\allowbreak{}Joints}, lines 426--582)\par}

{\footnotesize\raggedright\noindent \ref{B5.20c}.\quad \texttt{GlobalTheory.\allowbreak{}EX.\allowbreak{}RangeOne.\allowbreak{}rangeOne\_\allowbreak{}abs} (\texttt{Development.\allowbreak{}AbelianVarietyTheory.\allowbreak{}Proofs.\allowbreak{}GlobalInterpolation.\allowbreak{}TensorOfPairLift}, lines 637--640)\\
\texttt{Formalism.\allowbreak{}MainNotions.\allowbreak{}G.\allowbreak{}Joints.\allowbreak{}type\_\allowbreak{}one\_\allowbreak{}hypothesis\_\allowbreak{}of\_\allowbreak{}segments} (\texttt{Formalism.\allowbreak{}SharpRestriction.\allowbreak{}GrowthBounds.\allowbreak{}Joints}, lines 1283--1325)\\
\texttt{Formalism.\allowbreak{}MainNotions.\allowbreak{}G.\allowbreak{}first\_\allowbreak{}theorem\_\allowbreak{}closed\_\allowbreak{}form} (\texttt{Formalism.\allowbreak{}SharpRestriction.\allowbreak{}GrowthBounds.\allowbreak{}PlaneCountTheoremWithConstants}, lines 1761--1781)\par}

{\footnotesize\raggedright\noindent \ref{B5.21}.\quad \texttt{Formalism.\allowbreak{}L2.\allowbreak{}GlobalTheoryTheory.\allowbreak{}Loss7.\allowbreak{}losses\_\allowbreak{}of\_\allowbreak{}small\_\allowbreak{}share\_\allowbreak{}weight} (\texttt{Formalism.\allowbreak{}LocalToGlobal.\allowbreak{}GlobalTheory.\allowbreak{}LossRequirementsOfSmallShares}, lines 86--161)\par}

{\footnotesize\raggedright\noindent \ref{B5.21a}.\quad \texttt{Formalism.\allowbreak{}MainNotions.\allowbreak{}Lnk.\allowbreak{}image\_\allowbreak{}plMap} (\texttt{Formalism.\allowbreak{}SharpRestriction.\allowbreak{}GlobalSchurBounds.\allowbreak{}QuadraticLogDerivCongr}, lines 351--360)\\
\texttt{Formalism.\allowbreak{}MainNotions.\allowbreak{}Lnk.\allowbreak{}Mod.\allowbreak{}w\_\allowbreak{}off} (\texttt{Formalism.\allowbreak{}SharpRestriction.\allowbreak{}GlobalSchurBounds.\allowbreak{}QuadraticLogDerivCongr}, lines 416--422)\\
\texttt{Formalism.\allowbreak{}MainNotions.\allowbreak{}Lnk.\allowbreak{}Mod.\allowbreak{}w\_\allowbreak{}plMap} (\texttt{Formalism.\allowbreak{}SharpRestriction.\allowbreak{}GlobalSchurBounds.\allowbreak{}QuadraticLogDerivCongr}, lines 423--426)\\
\texttt{Formalism.\allowbreak{}MainNotions.\allowbreak{}Lnk.\allowbreak{}numN\_\allowbreak{}at} (\texttt{Formalism.\allowbreak{}SharpRestriction.\allowbreak{}GlobalSchurBounds.\allowbreak{}QuadraticLogDerivCongr}, lines 1261--1313)\\
\texttt{Formalism.\allowbreak{}MainNotions.\allowbreak{}Lnk.\allowbreak{}nodes\_\allowbreak{}at} (\texttt{Formalism.\allowbreak{}SharpRestriction.\allowbreak{}GlobalSchurBounds.\allowbreak{}QuadraticLogDerivCongr}, lines 1198--1249)\\
\texttt{Formalism.\allowbreak{}MainNotions.\allowbreak{}Lnk.\allowbreak{}nodesF} (\texttt{Formalism.\allowbreak{}SharpRestriction.\allowbreak{}GlobalSchurBounds.\allowbreak{}QuadraticLogDerivCongr}, lines 1314--1322)\par}

{\footnotesize\raggedright\noindent \ref{B5.21a}.\quad \texttt{Formalism.\allowbreak{}MainNotions.\allowbreak{}FamCore.\allowbreak{}bTop\_\allowbreak{}eq} (\texttt{Formalism.\allowbreak{}SharpRestriction.\allowbreak{}FamCore}, lines 124--125)\par}

{\footnotesize\raggedright\noindent \ref{B5.22}.\quad \texttt{Formalism.\allowbreak{}MainNotions.\allowbreak{}Lnk.\allowbreak{}lossF\_\allowbreak{}of\_\allowbreak{}7} (\texttt{Formalism.\allowbreak{}SharpRestriction.\allowbreak{}GlobalSchurBounds.\allowbreak{}LossesSmallShares}, lines 308--313)\\
\texttt{Formalism.\allowbreak{}MainNotions.\allowbreak{}ChosenWS} (\texttt{Formalism.\allowbreak{}SharpRestriction.\allowbreak{}GlobalSchurBounds.\allowbreak{}WeightVectorSmallSharesSeven}, lines 51--59)\par}

{\footnotesize\raggedright\noindent \ref{B5.22}.\quad \texttt{Formalism.\allowbreak{}MainNotions.\allowbreak{}FamCore.\allowbreak{}coverS\_\allowbreak{}carry} (\texttt{Formalism.\allowbreak{}SharpRestriction.\allowbreak{}FamCore}, lines 187--205)\par}

{\footnotesize\raggedright\noindent \ref{B5.22a}.\quad \texttt{GlobalTheory.\allowbreak{}EX.\allowbreak{}G4.\allowbreak{}wClosed} (\texttt{Development.\allowbreak{}AbelianVarietyTheory.\allowbreak{}Proofs.\allowbreak{}GlobalInterpolation.\allowbreak{}WindowShareRequirementClosedFormProof}, lines 40--42)\par}

{\footnotesize\raggedright\noindent \ref{B5.22a}.\quad \texttt{Formalism.\allowbreak{}MainNotions.\allowbreak{}AuxBlock8.\allowbreak{}Pf.\allowbreak{}sum\_\allowbreak{}Sig\_\allowbreak{}le} (\texttt{Formalism.\allowbreak{}SharpRestriction.\allowbreak{}GrowthBounds.\allowbreak{}LogarithmicSums.\allowbreak{}SlopeEstimateOfLemmaBSlopeLine}, lines 329--359)\\
\texttt{Formalism.\allowbreak{}MainNotions.\allowbreak{}AuxBlock8.\allowbreak{}Pf.\allowbreak{}meanEx\_\allowbreak{}of} (\texttt{Formalism.\allowbreak{}SharpRestriction.\allowbreak{}GrowthBounds.\allowbreak{}LogarithmicSums.\allowbreak{}SlopeEstimateOfLemmaBSlopeLine}, lines 435--481)\\
\texttt{Formalism.\allowbreak{}MainNotions.\allowbreak{}AuxBlock8.\allowbreak{}Pf.\allowbreak{}weight\_\allowbreak{}condition\_\allowbreak{}of\_\allowbreak{}mean\_\allowbreak{}excess} (\texttt{Formalism.\allowbreak{}SharpRestriction.\allowbreak{}GrowthBounds.\allowbreak{}LogarithmicSums.\allowbreak{}SlopeEstimateOfLemmaBSlopeLine}, lines 483--501)\par}

{\footnotesize\raggedright\noindent \ref{B5.22b}.\quad \texttt{Formalism.\allowbreak{}MainNotions.\allowbreak{}chosenWS\_\allowbreak{}of\_\allowbreak{}wall} (\texttt{Formalism.\allowbreak{}SharpRestriction.\allowbreak{}GlobalSchurBounds.\allowbreak{}WeightVectorSmallSharesEight}, lines 22--31)\par}

\section*{Chapter~\ref{ch:p5c3}: A slope bound for families of maps}

{\footnotesize\raggedright\noindent \ref{B5.23}.\quad \texttt{Formalism.\allowbreak{}HN.\allowbreak{}slope\_\allowbreak{}inequality\_\allowbreak{}mean} (\texttt{Formalism.\allowbreak{}HarderNarasimhan.\allowbreak{}HarderNarasimhanProofs.\allowbreak{}SlopeIneqMean}, lines 249--252)\\
\texttt{Formalism.\allowbreak{}HN.\allowbreak{}SlopeIneqMean} (\texttt{Formalism.\allowbreak{}HarderNarasimhan.\allowbreak{}StatementPl}, lines 148--160)\par}

{\footnotesize\raggedright\noindent \ref{B5.24}.\quad \texttt{Formalism.\allowbreak{}HN.\allowbreak{}frameB\_\allowbreak{}count\_\allowbreak{}lt} (\texttt{Formalism.\allowbreak{}HarderNarasimhan.\allowbreak{}HarderNarasimhanProofs.\allowbreak{}FrameHermSharpB}, lines 62--76)\\
\texttt{Formalism.\allowbreak{}HN.\allowbreak{}frameB\_\allowbreak{}sum\_\allowbreak{}sl} (\texttt{Formalism.\allowbreak{}HarderNarasimhan.\allowbreak{}HarderNarasimhanProofs.\allowbreak{}FrameHermSharpB}, lines 115--125)\\
\texttt{Formalism.\allowbreak{}HN.\allowbreak{}frameB\_\allowbreak{}hnStep} (\texttt{Formalism.\allowbreak{}HarderNarasimhan.\allowbreak{}HarderNarasimhanProofs.\allowbreak{}FrameHermSharpB}, lines 77--79)\\
\texttt{Formalism.\allowbreak{}HN.\allowbreak{}polygon\_\allowbreak{}extreme\_\allowbreak{}slopes} (\texttt{Formalism.\allowbreak{}HarderNarasimhan.\allowbreak{}HarderNarasimhanProofs.\allowbreak{}HNExtreme}, lines 114--116)\par}

{\footnotesize\raggedright\noindent \ref{B5.25}.\quad \texttt{Formalism.\allowbreak{}HN.\allowbreak{}thmF\_\allowbreak{}share} (\texttt{Formalism.\allowbreak{}HarderNarasimhan.\allowbreak{}HarderNarasimhanProofs.\allowbreak{}TheoremFSharp}, lines 11--34)\par}

{\footnotesize\raggedright\noindent \ref{B5.26}.\quad \texttt{Formalism.\allowbreak{}HN.\allowbreak{}frame\_\allowbreak{}linear\_\allowbreak{}algebra\_\allowbreak{}first} (\texttt{Formalism.\allowbreak{}HarderNarasimhan.\allowbreak{}HarderNarasimhanProofs.\allowbreak{}FrameLinearAlgebraFirst}, lines 77--222)\par}

{\footnotesize\raggedright\noindent \ref{B5.27}.\quad \texttt{Formalism.\allowbreak{}HN.\allowbreak{}thmF\_\allowbreak{}place} (\texttt{Formalism.\allowbreak{}HarderNarasimhan.\allowbreak{}HarderNarasimhanProofs.\allowbreak{}TheoremFSharp}, lines 35--160)\par}

{\footnotesize\raggedright\noindent \ref{B5.28}.\quad \texttt{Formalism.\allowbreak{}CorankShareLayerCake.\allowbreak{}FR.\allowbreak{}layer\_\allowbreak{}cake\_\allowbreak{}formula} (\texttt{Formalism.\allowbreak{}CorankShare.\allowbreak{}CorankShareProofs.\allowbreak{}LayerCakeFormula}, lines 57--98)\par}

{\footnotesize\raggedright\noindent \ref{B5.29}.\quad \texttt{Formalism.\allowbreak{}HN.\allowbreak{}theoremFSharp\_\allowbreak{}of} (\texttt{Formalism.\allowbreak{}HarderNarasimhan.\allowbreak{}HarderNarasimhanProofs.\allowbreak{}TheoremFSharp}, lines 168--387)\\
\texttt{Formalism.\allowbreak{}HN.\allowbreak{}TheoremFSharp} (\texttt{Formalism.\allowbreak{}HarderNarasimhan.\allowbreak{}StatementFrameF}, lines 64--100)\\
\texttt{Formalism.\allowbreak{}HN.\allowbreak{}frame\_\allowbreak{}theorem\_\allowbreak{}sharp} (\texttt{Formalism.\allowbreak{}HarderNarasimhan.\allowbreak{}HarderNarasimhanProofs.\allowbreak{}TheoremFSharp}, lines 388--389)\\
\texttt{Formalism.\allowbreak{}HN.\allowbreak{}theoremF\_\allowbreak{}of\_\allowbreak{}sharp} (\texttt{Formalism.\allowbreak{}HarderNarasimhan.\allowbreak{}HarderNarasimhanProofs.\allowbreak{}TheoremFSharp}, lines 390--399)\\
\texttt{Formalism.\allowbreak{}HN.\allowbreak{}frame\_\allowbreak{}theorem\_\allowbreak{}log\_\allowbreak{}ideal} (\texttt{Formalism.\allowbreak{}HarderNarasimhan.\allowbreak{}HarderNarasimhanProofs.\allowbreak{}TheoremFSharp}, lines 400--401)\par}

{\footnotesize\raggedright\noindent \ref{B5.31}.\quad \texttt{Formalism.\allowbreak{}MainNotions.\allowbreak{}Lnk.\allowbreak{}frame\_\allowbreak{}at\_\allowbreak{}of} (\texttt{Formalism.\allowbreak{}SharpRestriction.\allowbreak{}GlobalSchurBounds.\allowbreak{}TopDensityNGeSixteen}, lines 538--601)\\
\texttt{Formalism.\allowbreak{}LFunctionChainParameters.\allowbreak{}MatrixLatticeEstimates.\allowbreak{}joint\_\allowbreak{}a\_\allowbreak{}ramified\_\allowbreak{}holds} (\texttt{Formalism.\allowbreak{}SizeWindowBrackets.\allowbreak{}ChainSlopeWindows.\allowbreak{}JointDModelAppliesUnitLattice}, lines 254--259)\\
\texttt{Formalism.\allowbreak{}LFunctionChainParameters.\allowbreak{}MatrixLatticeEstimates.\allowbreak{}ramification\_\allowbreak{}linear\_\allowbreak{}case\_\allowbreak{}holds} (\texttt{Formalism.\allowbreak{}SizeWindowBrackets.\allowbreak{}ChainSlopeWindows.\allowbreak{}JointDModelAppliesUnitLattice}, lines 164--186)\\
\texttt{Formalism.\allowbreak{}MainNotions.\allowbreak{}Lnk.\allowbreak{}FrameReal\_\allowbreak{}at} (\texttt{Formalism.\allowbreak{}SharpRestriction.\allowbreak{}GlobalSchurBounds.\allowbreak{}TopDensityNGeSixteen}, lines 519--532)\\
\texttt{Formalism.\allowbreak{}MainNotions.\allowbreak{}GamL} (\texttt{Formalism.\allowbreak{}SharpRestriction.\allowbreak{}GlobalSchurBounds.\allowbreak{}Chosen}, lines 14--15)\par}

{\footnotesize\raggedright\noindent \ref{B5.31a}.\quad \texttt{Formalism.\allowbreak{}L2.\allowbreak{}SG.\allowbreak{}RowsPos.\allowbreak{}kernel\_\allowbreak{}diagonal\_\allowbreak{}positive} (\texttt{Formalism.\allowbreak{}LocalToGlobal.\allowbreak{}SiegelLatticeTheory.\allowbreak{}FlagStabilizerTensors.\allowbreak{}RowsPos}, lines 202--209)\\
\texttt{Formalism.\allowbreak{}MainNotions.\allowbreak{}S.\allowbreak{}linkS\_\allowbreak{}concl} (\texttt{Formalism.\allowbreak{}SharpRestriction.\allowbreak{}CuspSpanNormBounds.\allowbreak{}StandingHypothesesHoldFamily}, lines 186--205)\\
\texttt{Formalism.\allowbreak{}MainNotions.\allowbreak{}Lnk.\allowbreak{}matK\_\allowbreak{}at\_\allowbreak{}of} (\texttt{Formalism.\allowbreak{}SharpRestriction.\allowbreak{}GlobalSchurBounds.\allowbreak{}TopDensityNGeSixteen}, lines 1097--1108)\\
\texttt{Formalism.\allowbreak{}MainNotions.\allowbreak{}product\_\allowbreak{}first\_\allowbreak{}zeroth\_\allowbreak{}parts.\allowbreak{}rank\_\allowbreak{}nat} (\texttt{Formalism.\allowbreak{}SharpRestriction.\allowbreak{}GlobalSchurBounds.\allowbreak{}FiniteModelOfAdmissibleTransports}, lines 782--841)\\
\texttt{Formalism.\allowbreak{}L2.\allowbreak{}SG.\allowbreak{}KappaOffDiagonal.\allowbreak{}cosetsC\_\allowbreak{}holds} (\texttt{Formalism.\allowbreak{}LocalToGlobal.\allowbreak{}SiegelLatticeTheory.\allowbreak{}CokernelGLCounts.\allowbreak{}RankPointCounts.\allowbreak{}KloostermanCosetSumLe}, lines 186--213)\\
\texttt{GlobalTheory.\allowbreak{}rowMainEqP\_\allowbreak{}of\_\allowbreak{}siegelWeil} (\texttt{Development.\allowbreak{}AbelianVarietyTheory.\allowbreak{}Proofs.\allowbreak{}FlagMeasures.\allowbreak{}RowMainEqPOfSiegelWeil}, lines 10--17)\par}

{\footnotesize\raggedright\noindent \ref{B5.32}.\quad \texttt{Formalism.\allowbreak{}MainNotions.\allowbreak{}Lnk.\allowbreak{}frameF\_\allowbreak{}of} (\texttt{Formalism.\allowbreak{}SharpRestriction.\allowbreak{}GlobalSchurBounds.\allowbreak{}TopDensityNGeSixteen}, lines 606--647)\\
\texttt{Formalism.\allowbreak{}MainNotions.\allowbreak{}Lnk.\allowbreak{}Mod.\allowbreak{}data\_\allowbreak{}HHBound} (\texttt{Formalism.\allowbreak{}SharpRestriction.\allowbreak{}GlobalSchurBounds.\allowbreak{}QuadraticLogDerivCongr}, lines 714--729)\\
\texttt{Formalism.\allowbreak{}MainNotions.\allowbreak{}Lnk.\allowbreak{}lat\_\allowbreak{}at} (\texttt{Formalism.\allowbreak{}SharpRestriction.\allowbreak{}GlobalSchurBounds.\allowbreak{}QuadraticLogDerivCongr}, lines 1323--1333)\\
\texttt{Formalism.\allowbreak{}MainNotions.\allowbreak{}Lnk.\allowbreak{}Mod.\allowbreak{}data} (\texttt{Formalism.\allowbreak{}SharpRestriction.\allowbreak{}GlobalSchurBounds.\allowbreak{}QuadraticLogDerivCongr}, lines 464--505)\par}

{\footnotesize\raggedright\noindent \ref{B5.32b}.\quad \texttt{Formalism.\allowbreak{}MainNotions.\allowbreak{}Ext.\allowbreak{}Lw.\allowbreak{}mu0Tie\_\allowbreak{}of\_\allowbreak{}sym} (\texttt{Formalism.\allowbreak{}SharpRestriction.\allowbreak{}GlobalSchurBounds.\allowbreak{}GlobalRankLatticeProofs.\allowbreak{}UnitTraceSumBounds.\allowbreak{}RankConditionsOfLayerTwo}, lines 284--289)\\
\texttt{Formalism.\allowbreak{}MainNotions.\allowbreak{}StandardLattice.\allowbreak{}HarishChandraReduction.\allowbreak{}chosenProviso\_\allowbreak{}min} (\texttt{Formalism.\allowbreak{}SharpRestriction.\allowbreak{}GlobalSchurBounds.\allowbreak{}TheoremThreeFamilyOfFormalDegreeInput}, lines 76--79)\par}

{\footnotesize\raggedright\noindent \ref{B5.32a}.\quad \texttt{UpperTheory.\allowbreak{}Top.\allowbreak{}designS} (\texttt{Development.\allowbreak{}AutomorphicArithmetic.\allowbreak{}CharacterBoxes.\allowbreak{}DesignS}, lines 21--99)\\
\texttt{Formalism.\allowbreak{}MainNotions.\allowbreak{}Lnk.\allowbreak{}matRealF\_\allowbreak{}of} (\texttt{Formalism.\allowbreak{}SharpRestriction.\allowbreak{}GlobalSchurBounds.\allowbreak{}TopDensityNGeSixteen}, lines 750--806)\par}

\section*{Chapter~\ref{ch:p5c4}: The pointwise bound and its eight error terms}

{\footnotesize\raggedright\noindent Chapter introduction.\quad \texttt{Formalism.\allowbreak{}RankGapBounds.\allowbreak{}WeightedMeanData.\allowbreak{}Data.\allowbreak{}s} (\texttt{Formalism.\allowbreak{}RankGapBounds.\allowbreak{}MeanBoundSecondForm}, lines 92--95)\\
\texttt{Formalism.\allowbreak{}MainNotions.\allowbreak{}Lnk.\allowbreak{}Mod.\allowbreak{}data} (\texttt{Formalism.\allowbreak{}SharpRestriction.\allowbreak{}GlobalSchurBounds.\allowbreak{}QuadraticLogDerivCongr}, lines 464--505)\par}

{\footnotesize\raggedright\noindent \ref{eq:p5c4-Lo}.\quad \texttt{Formalism.\allowbreak{}RankGapBounds.\allowbreak{}WeightedMeanData.\allowbreak{}Data.\allowbreak{}SDiv} (\texttt{Formalism.\allowbreak{}RankGapBounds.\allowbreak{}RankGapProofs.\allowbreak{}WeightedMeanData}, lines 129--133)\\
\texttt{Formalism.\allowbreak{}RankGapBounds.\allowbreak{}WeightedMeanData.\allowbreak{}Data.\allowbreak{}LoP} (\texttt{Formalism.\allowbreak{}RankGapBounds.\allowbreak{}MeanBoundSecondForm}, lines 105--106)\\
\texttt{Formalism.\allowbreak{}LFunctionChainParameters.\allowbreak{}WeightedSlopeBounds.\allowbreak{}Lo5} (\texttt{Formalism.\allowbreak{}SizeWindowBrackets.\allowbreak{}QuotientSlopeBoundsWith}, lines 15--17)\par}

{\footnotesize\raggedright\noindent \ref{hyp:p5c4-parameters}.\quad \texttt{Formalism.\allowbreak{}RankGapBounds.\allowbreak{}WeightedMeanData.\allowbreak{}Data.\allowbreak{}HidaSetting} (\texttt{Formalism.\allowbreak{}RankGapBounds.\allowbreak{}MeanBoundSecondForm}, lines 111--117)\\
\texttt{Formalism.\allowbreak{}LFunctionChainParameters.\allowbreak{}WeightedSlopeBounds.\allowbreak{}PointwiseBoundVariant} (\texttt{Formalism.\allowbreak{}SizeWindowBrackets.\allowbreak{}QuotientSlopeBoundsWith}, lines 72--75)\par}

{\footnotesize\raggedright\noindent \ref{B5.33}.\quad \texttt{Formalism.\allowbreak{}MainNotions.\allowbreak{}Lnk.\allowbreak{}linkCore} (\texttt{Formalism.\allowbreak{}SharpRestriction.\allowbreak{}GlobalSchurBounds.\allowbreak{}LossesSmallShares}, lines 38--297)\\
\texttt{Formalism.\allowbreak{}MainNotions.\allowbreak{}linkGLOBAL7\_\allowbreak{}items} (\texttt{Formalism.\allowbreak{}SharpRestriction.\allowbreak{}GlobalSchurBounds.\allowbreak{}LinkMin}, lines 693--722)\\
\texttt{Formalism.\allowbreak{}MainNotions.\allowbreak{}ChainTied} (\texttt{Formalism.\allowbreak{}SharpRestriction.\allowbreak{}GlobalSchurBounds.\allowbreak{}TieGlobal}, lines 97--115)\par}

{\footnotesize\raggedright\noindent \ref{B5.33a}.\quad \texttt{Formalism.\allowbreak{}MainNotions.\allowbreak{}LnkBr.\allowbreak{}prim\_\allowbreak{}unique} (\texttt{Formalism.\allowbreak{}SharpRestriction.\allowbreak{}GlobalSchurBounds.\allowbreak{}TopDensityNGeSixteen}, lines 133--158)\\
\texttt{Formalism.\allowbreak{}MainNotions.\allowbreak{}Lnk.\allowbreak{}char\_\allowbreak{}at} (\texttt{Formalism.\allowbreak{}SharpRestriction.\allowbreak{}GlobalSchurBounds.\allowbreak{}TopDensityNGeSixteen}, lines 476--505)\\
\texttt{Formalism.\allowbreak{}MainNotions.\allowbreak{}Lnk.\allowbreak{}charF} (\texttt{Formalism.\allowbreak{}SharpRestriction.\allowbreak{}GlobalSchurBounds.\allowbreak{}TopDensityNGeSixteen}, lines 506--508)\par}

{\footnotesize\raggedright\noindent \ref{B5.34}.\quad \texttt{Formalism.\allowbreak{}LFunctionChainParameters.\allowbreak{}WeightedSlopeBounds.\allowbreak{}Lam5} (\texttt{Formalism.\allowbreak{}SizeWindowBrackets.\allowbreak{}ChainSlopeWindows.\allowbreak{}LevelOfCongruence}, lines 17--19)\\
\texttt{Formalism.\allowbreak{}RankGapBounds.\allowbreak{}WeightedMeanData.\allowbreak{}Data.\allowbreak{}Main} (\texttt{Formalism.\allowbreak{}RankGapBounds.\allowbreak{}MeanBoundSecondForm}, lines 107--110)\par}

{\footnotesize\raggedright\noindent \ref{B5.35}.\quad \texttt{Formalism.\allowbreak{}LFunctionChainParameters.\allowbreak{}WeightedSlopeBounds.\allowbreak{}key\_\allowbreak{}identity5} (\texttt{Formalism.\allowbreak{}SizeWindowBrackets.\allowbreak{}ChainSlopeWindows.\allowbreak{}LevelOfCongruence}, lines 105--111)\par}

{\footnotesize\raggedright\noindent \ref{B5.35}.\quad \texttt{Formalism.\allowbreak{}RankGapBounds.\allowbreak{}WeightedMeanData.\allowbreak{}Data.\allowbreak{}key\_\allowbreak{}identity} (\texttt{Formalism.\allowbreak{}RankGapBounds.\allowbreak{}RankGapProofs.\allowbreak{}WeightedMeanData}, lines 170--217)\par}

{\footnotesize\raggedright\noindent \ref{B5.36}.\quad \texttt{Formalism.\allowbreak{}LFunctionChainParameters.\allowbreak{}WeightedSlopeBounds.\allowbreak{}s\_\allowbreak{}lower5} (\texttt{Formalism.\allowbreak{}SizeWindowBrackets.\allowbreak{}ChainSlopeWindows.\allowbreak{}LevelOfCongruence}, lines 50--66)\par}

{\footnotesize\raggedright\noindent \ref{B5.37}.\quad \texttt{Formalism.\allowbreak{}LFunctionChainParameters.\allowbreak{}WeightedSlopeBounds.\allowbreak{}Vvar\_\allowbreak{}bound5} (\texttt{Formalism.\allowbreak{}SizeWindowBrackets.\allowbreak{}ChainSlopeWindows.\allowbreak{}LevelOfCongruence}, lines 68--79)\par}

{\footnotesize\raggedright\noindent \ref{B5.38}.\quad \texttt{Formalism.\allowbreak{}LFunctionChainParameters.\allowbreak{}WeightedSlopeBounds.\allowbreak{}u\_\allowbreak{}function\_\allowbreak{}le\_\allowbreak{}mean} (\texttt{Formalism.\allowbreak{}SizeWindowBrackets.\allowbreak{}ChainSlopeWindows.\allowbreak{}LevelOfCongruence}, lines 81--103)\\
\texttt{Formalism.\allowbreak{}RankGapBounds.\allowbreak{}WeightedMeanData.\allowbreak{}Data.\allowbreak{}uF\_\allowbreak{}nonneg} (\texttt{Formalism.\allowbreak{}RankGapBounds.\allowbreak{}RankGapProofs.\allowbreak{}WeightedMeanData}, lines 159--169)\par}

{\footnotesize\raggedright\noindent \ref{B5.39}.\quad \texttt{Formalism.\allowbreak{}RankGapBounds.\allowbreak{}WeightedMeanData.\allowbreak{}phi\_\allowbreak{}bound} (\texttt{Formalism.\allowbreak{}RankGapBounds.\allowbreak{}RankGapProofs.\allowbreak{}WeightedMeanData}, lines 81--118)\par}

{\footnotesize\raggedright\noindent \ref{B5.40}.\quad \texttt{Formalism.\allowbreak{}LFunctionChainParameters.\allowbreak{}WeightedSlopeBounds.\allowbreak{}Ga\_\allowbreak{}bound5} (\texttt{Formalism.\allowbreak{}SizeWindowBrackets.\allowbreak{}ChainSlopeWindows.\allowbreak{}LevelOfCongruence}, lines 113--161)\par}

{\footnotesize\raggedright\noindent \ref{B5.41}.\quad \texttt{Formalism.\allowbreak{}LFunctionChainParameters.\allowbreak{}WeightedSlopeBounds.\allowbreak{}weighted\_\allowbreak{}mean\_\allowbreak{}lower\_\allowbreak{}bound\_\allowbreak{}le} (\texttt{Formalism.\allowbreak{}SizeWindowBrackets.\allowbreak{}ChainSlopeWindows.\allowbreak{}LevelOfCongruence}, lines 162--171)\\
\texttt{Formalism.\allowbreak{}LFunctionChainParameters.\allowbreak{}WeightedSlopeBounds.\allowbreak{}each\_\allowbreak{}unit\_\allowbreak{}bound\_\allowbreak{}statement} (\texttt{Formalism.\allowbreak{}SizeWindowBrackets.\allowbreak{}ChainSlopeWindows.\allowbreak{}LevelOfCongruence}, lines 188--191)\par}

{\footnotesize\raggedright\noindent \ref{B5.42}.\quad \texttt{Formalism.\allowbreak{}LFunctionChainParameters.\allowbreak{}WeightedSlopeBounds.\allowbreak{}ErrX5at} (\texttt{Formalism.\allowbreak{}SizeWindowBrackets.\allowbreak{}QuotientSlopeBoundsWith}, lines 45--55)\\
\texttt{Formalism.\allowbreak{}LFunctionChainParameters.\allowbreak{}WeightedSlopeBounds.\allowbreak{}UBar5} (\texttt{Formalism.\allowbreak{}SizeWindowBrackets.\allowbreak{}QuotientSlopeBoundsWith}, lines 41--43)\par}

{\footnotesize\raggedright\noindent \ref{B5.43}.\quad \texttt{Formalism.\allowbreak{}LFunctionChainParameters.\allowbreak{}WeightedSlopeBounds.\allowbreak{}each\_\allowbreak{}unit\_\allowbreak{}bound\_\allowbreak{}statement} (\texttt{Formalism.\allowbreak{}SizeWindowBrackets.\allowbreak{}ChainSlopeWindows.\allowbreak{}LevelOfCongruence}, lines 175--197)\\
\texttt{Formalism.\allowbreak{}LFunctionChainParameters.\allowbreak{}WeightedSlopeBounds.\allowbreak{}PointwiseBoundVariant} (\texttt{Formalism.\allowbreak{}SizeWindowBrackets.\allowbreak{}QuotientSlopeBoundsWith}, lines 72--75)\par}

\section*{Chapter~\ref{ch:p5c5}: Four terms and the square-root bound}

{\footnotesize\raggedright\noindent \ref{B5.44}.\quad \texttt{Formalism.\allowbreak{}LFunctionChainParameters.\allowbreak{}LErr.\allowbreak{}core\_\allowbreak{}ineq} (\texttt{Formalism.\allowbreak{}SizeWindowBrackets.\allowbreak{}ChainSlopeWindows.\allowbreak{}SubThreeMulLogTwoLe}, lines 13--208)\par}

{\footnotesize\raggedright\noindent \ref{B5.45}.\quad \texttt{Formalism.\allowbreak{}LFunctionChainParameters.\allowbreak{}LChain.\allowbreak{}errX5at\_\allowbreak{}le} (\texttt{Formalism.\allowbreak{}SizeWindowBrackets.\allowbreak{}ChainSlopeWindows.\allowbreak{}PChain}, lines 613--673)\par}

{\footnotesize\raggedright\noindent \ref{B5.45}.\quad \texttt{Formalism.\allowbreak{}MainNotions.\allowbreak{}Lnk.\allowbreak{}Mod.\allowbreak{}sum\_\allowbreak{}places} (\texttt{Formalism.\allowbreak{}SharpRestriction.\allowbreak{}GlobalSchurBounds.\allowbreak{}QuadraticLogDerivCongr}, lines 458--463)\\
\texttt{Formalism.\allowbreak{}MainNotions.\allowbreak{}Lnk.\allowbreak{}Mod.\allowbreak{}ident} (\texttt{Formalism.\allowbreak{}SharpRestriction.\allowbreak{}GlobalSchurBounds.\allowbreak{}QuadraticLogDerivCongr}, lines 792--798)\par}

{\footnotesize\raggedright\noindent \ref{B5.46}.\quad \texttt{Deduction.\allowbreak{}ch\_\allowbreak{}top\_\allowbreak{}main\_\allowbreak{}eq} (\texttt{Development.\allowbreak{}AutomorphicArithmetic.\allowbreak{}DedekindLFunctions.\allowbreak{}LandauTheorems.\allowbreak{}ChTop}, lines 11--36)\par}

{\footnotesize\raggedright\noindent \ref{B5.47}.\quad \texttt{Deduction.\allowbreak{}ch\_\allowbreak{}top} (\texttt{Development.\allowbreak{}AutomorphicArithmetic.\allowbreak{}DedekindLFunctions.\allowbreak{}LandauTheorems.\allowbreak{}ChTop}, lines 37--73)\par}

{\footnotesize\raggedright\noindent \ref{B5.48}.\quad \texttt{Deduction.\allowbreak{}ch\_\allowbreak{}balance} (\texttt{Development.\allowbreak{}AutomorphicArithmetic.\allowbreak{}DedekindLFunctions.\allowbreak{}LandauTheorems.\allowbreak{}ChBalance}, lines 200--261)\par}

{\footnotesize\raggedright\noindent \ref{B5.49}.\quad \texttt{Deduction.\allowbreak{}CharacterBalanceMoves.\allowbreak{}core} (\texttt{Development.\allowbreak{}AutomorphicArithmetic.\allowbreak{}DedekindLFunctions.\allowbreak{}LandauTheorems.\allowbreak{}ChCore}, lines 283--339)\\
\texttt{Deduction.\allowbreak{}CharacterBalanceMoves.\allowbreak{}th\_\allowbreak{}le} (\texttt{Development.\allowbreak{}AutomorphicArithmetic.\allowbreak{}DedekindLFunctions.\allowbreak{}LandauTheorems.\allowbreak{}ChCore}, lines 75--153)\par}

{\footnotesize\raggedright\noindent \ref{B5.50}.\quad \texttt{Formalism.\allowbreak{}MainNotions.\allowbreak{}Lnk.\allowbreak{}winA} (\texttt{Formalism.\allowbreak{}SharpRestriction.\allowbreak{}GlobalSchurBounds.\allowbreak{}QuadraticLogDerivCongr}, lines 1414--1416)\\
\texttt{Formalism.\allowbreak{}MainNotions.\allowbreak{}Lnk.\allowbreak{}winL} (\texttt{Formalism.\allowbreak{}SharpRestriction.\allowbreak{}GlobalSchurBounds.\allowbreak{}QuadraticLogDerivCongr}, lines 1417--1419)\par}

{\footnotesize\raggedright\noindent \ref{B5.51}.\quad \texttt{Formalism.\allowbreak{}LFunctionChainParameters.\allowbreak{}ThickSetMeanBounds.\allowbreak{}SZ.\allowbreak{}T1} (\texttt{Formalism.\allowbreak{}SizeWindowBrackets.\allowbreak{}ChainSlopeWindows.\allowbreak{}LipschitzWeightAndMinimum}, lines 31--40)\par}

{\footnotesize\raggedright\noindent \ref{B5.52}.\quad \texttt{Formalism.\allowbreak{}LFunctionChainParameters.\allowbreak{}ThickSetMeanBounds.\allowbreak{}SZ.\allowbreak{}T2} (\texttt{Formalism.\allowbreak{}SizeWindowBrackets.\allowbreak{}ChainSlopeWindows.\allowbreak{}LipschitzWeightAndMinimum}, lines 41--52)\par}

{\footnotesize\raggedright\noindent \ref{B5.53}.\quad \texttt{Formalism.\allowbreak{}LFunctionChainParameters.\allowbreak{}ThickSetMeanBounds.\allowbreak{}SZ.\allowbreak{}T3} (\texttt{Formalism.\allowbreak{}SizeWindowBrackets.\allowbreak{}ChainSlopeWindows.\allowbreak{}LipschitzWeightAndMinimum}, lines 53--76)\par}

{\footnotesize\raggedright\noindent \ref{B5.54}.\quad \texttt{Formalism.\allowbreak{}LFunctionChainParameters.\allowbreak{}ThickSetMeanBounds.\allowbreak{}SZ.\allowbreak{}T4} (\texttt{Formalism.\allowbreak{}SizeWindowBrackets.\allowbreak{}ChainSlopeWindows.\allowbreak{}LipschitzWeightAndMinimum}, lines 77--114)\par}

{\footnotesize\raggedright\noindent \ref{B5.55}.\quad \texttt{Formalism.\allowbreak{}LFunctionChainParameters.\allowbreak{}ThickSetMeanBounds.\allowbreak{}proposition\_\allowbreak{}size\_\allowbreak{}bound} (\texttt{Formalism.\allowbreak{}SizeWindowBrackets.\allowbreak{}ChainSlopeWindows.\allowbreak{}LipschitzWeightAndMinimum}, lines 507--593)\par}

{\footnotesize\raggedright\noindent \ref{B5.56}.\quad \texttt{Deduction.\allowbreak{}ch\_\allowbreak{}core} (\texttt{Development.\allowbreak{}AutomorphicArithmetic.\allowbreak{}DedekindLFunctions.\allowbreak{}LandauTheorems.\allowbreak{}ChCore}, lines 343--443)\\
\texttt{Deduction.\allowbreak{}CharacterBalanceMoves.\allowbreak{}one\_\allowbreak{}char} (\texttt{Development.\allowbreak{}AutomorphicArithmetic.\allowbreak{}DedekindLFunctions.\allowbreak{}LandauTheorems.\allowbreak{}ChCore}, lines 202--282)\par}

{\footnotesize\raggedright\noindent \ref{B5.56}.\quad \texttt{UpperTheory.\allowbreak{}Comparison₂.\allowbreak{}core\_\allowbreak{}gm} (\texttt{Development.\allowbreak{}AutomorphicArithmetic.\allowbreak{}PrimitiveCosets.\allowbreak{}CoreGm}, lines 12--74)\par}

\section*{Chapter~\ref{ch:window-inputs}: The window bound and its inputs}

{\footnotesize\raggedright\noindent \ref{B6.01}.\quad \texttt{Formalism.\allowbreak{}GaloisCohomology.\allowbreak{}Sha} (\texttt{Formalism.\allowbreak{}GaloisCohomology.\allowbreak{}Statement}, lines 75--76)\\
\texttt{Formalism.\allowbreak{}GaloisCohomology.\allowbreak{}shaLen} (\texttt{Formalism.\allowbreak{}GaloisCohomology.\allowbreak{}Statement}, lines 78)\\
\texttt{UpperTheory.\allowbreak{}Top.\allowbreak{}VVle} (\texttt{Development.\allowbreak{}AutomorphicArithmetic.\allowbreak{}CharacterBoxes.\allowbreak{}ExistsPageN}, lines 12--18)\par}

{\footnotesize\raggedright\noindent \ref{B6.02}.\quad \texttt{UpperTheory.\allowbreak{}Top.\allowbreak{}IsChiField} (\texttt{Development.\allowbreak{}AutomorphicArithmetic.\allowbreak{}CharacterBoxes.\allowbreak{}ExistsPageN}, lines 19--25)\\
\texttt{UpperTheory.\allowbreak{}Top.\allowbreak{}WinPair} (\texttt{Development.\allowbreak{}AutomorphicArithmetic.\allowbreak{}CharacterBoxes.\allowbreak{}ExistsPageN}, lines 27--37)\\
\texttt{UpperTheory.\allowbreak{}Top.\allowbreak{}PageN} (\texttt{Development.\allowbreak{}AutomorphicArithmetic.\allowbreak{}CharacterBoxes.\allowbreak{}ExistsPageN}, lines 39)\par}

{\footnotesize\raggedright\noindent \ref{B6.03}.\quad \texttt{UpperTheory.\allowbreak{}Top.\allowbreak{}VAt} (\texttt{Development.\allowbreak{}AutomorphicArithmetic.\allowbreak{}CharacterBoxes.\allowbreak{}ExistsPageN}, lines 41--42)\\
\texttt{UpperTheory.\allowbreak{}Top.\allowbreak{}VWinAt} (\texttt{Development.\allowbreak{}AutomorphicArithmetic.\allowbreak{}CharacterBoxes.\allowbreak{}ExistsPageN}, lines 44--47)\\
\texttt{UpperTheory.\allowbreak{}Top.\allowbreak{}VWin} (\texttt{Development.\allowbreak{}AutomorphicArithmetic.\allowbreak{}CharacterBoxes.\allowbreak{}ExistsPageN}, lines 49)\par}

{\footnotesize\raggedright\noindent \ref{B6.04}.\quad \texttt{UpperTheory.\allowbreak{}One.\allowbreak{}vwin\_\allowbreak{}of\_\allowbreak{}row} (\texttt{Development.\allowbreak{}AutomorphicArithmetic.\allowbreak{}AutomorphicVarieties.\allowbreak{}ModularityArgument.\allowbreak{}ThresholdArguments.\allowbreak{}DenseThresholdCases.\allowbreak{}VwinOfRow}, lines 4--5)\par}

{\footnotesize\raggedright\noindent \ref{B6.04}.\quad \texttt{UpperTheory.\allowbreak{}LiteratureTheorems.\allowbreak{}OpenImageCondition} (\texttt{Development.\allowbreak{}AutomorphicArithmetic.\allowbreak{}AutomorphicVarieties.\allowbreak{}ModularityArgument.\allowbreak{}ThresholdArguments.\allowbreak{}DenseThresholdCases.\allowbreak{}TemperedProductOfSquareIntegrable}, lines 207--208)\par}

{\footnotesize\raggedright\noindent \ref{B6.04b}.\quad \texttt{UpperTheory.\allowbreak{}ConvolutionTheory.\allowbreak{}G2.\allowbreak{}uSide} (\texttt{Development.\allowbreak{}AutomorphicArithmetic.\allowbreak{}Foundations.\allowbreak{}Convolution.\allowbreak{}WeightRankUnits.\allowbreak{}USide}, lines 9--48)\\
\texttt{UpperTheory.\allowbreak{}ConvolutionTheory.\allowbreak{}G2.\allowbreak{}g2s} (\texttt{Development.\allowbreak{}AutomorphicArithmetic.\allowbreak{}Foundations.\allowbreak{}Convolution.\allowbreak{}WeightRankUnits.\allowbreak{}NeZeroOfFinprod}, lines 40--72)\par}

{\footnotesize\raggedright\noindent \ref{B6.04b}.\quad \texttt{UpperTheory.\allowbreak{}ConvolutionTheory.\allowbreak{}G2.\allowbreak{}matchSplit} (\texttt{Development.\allowbreak{}AutomorphicArithmetic.\allowbreak{}Foundations.\allowbreak{}Convolution.\allowbreak{}WeightRankUnits.\allowbreak{}MatchSplit}, lines 12--21)\par}

{\footnotesize\raggedright\noindent \ref{B6.04b}.\quad \texttt{GlobalTheory.\allowbreak{}UnitLemma.\allowbreak{}InvAlg.\allowbreak{}kapSum} (\texttt{Development.\allowbreak{}AbelianVarietyTheory.\allowbreak{}Definitions.\allowbreak{}NonGaloisDescent.\allowbreak{}UnitLemmaGroup}, lines 145--147)\\
\texttt{GlobalTheory.\allowbreak{}UnitLemma.\allowbreak{}InvAlg.\allowbreak{}stSum} (\texttt{Development.\allowbreak{}AbelianVarietyTheory.\allowbreak{}Definitions.\allowbreak{}NonGaloisDescent.\allowbreak{}UnitLemmaGroup}, lines 143--144)\par}

{\footnotesize\raggedright\noindent \ref{B6.04e}.\quad \texttt{PositivityTheory.\allowbreak{}Rg.\allowbreak{}Bc.\allowbreak{}ArchbcB.\allowbreak{}B.\allowbreak{}archimedean\_\allowbreak{}base\_\allowbreak{}change\_\allowbreak{}first} (\texttt{Development.\allowbreak{}WeilBound.\allowbreak{}RepresentationTheoryGL.\allowbreak{}ArchimedeanBaseChange.\allowbreak{}ArchimedeanBaseChangeProof.\allowbreak{}IwahoriFactorisation.\allowbreak{}MultiplicityOneBase.\allowbreak{}ArchimedeanBaseChangeFirst}, lines 29--62)\par}

{\footnotesize\raggedright\noindent \ref{B6.04ea}.\quad \texttt{GlobalTheory.\allowbreak{}CountId.\allowbreak{}countIdentity\_\allowbreak{}of\_\allowbreak{}unitMatch} (\texttt{Development.\allowbreak{}AbelianVarietyTheory.\allowbreak{}Proofs.\allowbreak{}InertiaEigenvalues.\allowbreak{}RelativeFundamentalLemmaCountIdentityOfUnitMatch}, lines 328--425)\par}

{\footnotesize\raggedright\noindent \ref{B6.04eb}.\quad \texttt{GlobalTheory.\allowbreak{}StableLatCount.\allowbreak{}card\_\allowbreak{}stable\_\allowbreak{}lattices} (\texttt{Development.\allowbreak{}AbelianVarietyTheory.\allowbreak{}Proofs.\allowbreak{}InertiaEigenvalues.\allowbreak{}CardStableLattices}, lines 16--34)\\
\texttt{GlobalTheory.\allowbreak{}UnitLemma.\allowbreak{}FibSum.\allowbreak{}fib\_\allowbreak{}card} (\texttt{Development.\allowbreak{}AbelianVarietyTheory.\allowbreak{}Proofs.\allowbreak{}InertiaEigenvalues.\allowbreak{}SumStableLatticesProductEqQPowMul}, lines 15--53)\par}

{\footnotesize\raggedright\noindent \ref{B6.04ec}.\quad \texttt{GlobalTheory.\allowbreak{}UnitLemma.\allowbreak{}unitLemma\_\allowbreak{}lie\_\allowbreak{}and\_\allowbreak{}group} (\texttt{Development.\allowbreak{}AbelianVarietyTheory.\allowbreak{}Proofs.\allowbreak{}InertiaEigenvalues.\allowbreak{}UnitLemmaLieAlgebraUnitaryGroupOfJacquetRallis}, lines 14--18)\par}

{\footnotesize\raggedright\noindent \ref{B6.04ed}.\quad \texttt{GlobalTheory.\allowbreak{}UnitLemma.\allowbreak{}unitLemma\_\allowbreak{}identity\_\allowbreak{}unconditional} (\texttt{Development.\allowbreak{}AbelianVarietyTheory.\allowbreak{}Proofs.\allowbreak{}DivisorProofs.\allowbreak{}UnitLemmaUnramifiedProductUnitaryGroups}, lines 10--14)\par}

{\footnotesize\raggedright\noindent \ref{B6.04ed}.\quad \texttt{GlobalTheory.\allowbreak{}InertUnitOrb.\allowbreak{}kapOrb\_\allowbreak{}unit\_\allowbreak{}eq\_\allowbreak{}kapSum} (\texttt{Development.\allowbreak{}AbelianVarietyTheory.\allowbreak{}Proofs.\allowbreak{}DivisorProofs.\allowbreak{}InertPlaceKappaOrbitalUnitEqSignedLatticeSum}, lines 68--213)\par}

{\footnotesize\raggedright\noindent \ref{B6.04ed}.\quad \texttt{GlobalTheory.\allowbreak{}UnitLemma.\allowbreak{}UnitBase.\allowbreak{}exists\_\allowbreak{}unit\_\allowbreak{}base} (\texttt{Development.\allowbreak{}AbelianVarietyTheory.\allowbreak{}Proofs.\allowbreak{}InertiaEigenvalues.\allowbreak{}ExistsUnitBasePointAdjoinTwistedSigma}, lines 168--213)\par}

{\footnotesize\raggedright\noindent \ref{B6.04ed}.\quad \texttt{GlobalTheory.\allowbreak{}TwClassForms.\allowbreak{}twCl\_\allowbreak{}equiv\_\allowbreak{}cl} (\texttt{Development.\allowbreak{}AbelianVarietyTheory.\allowbreak{}Proofs.\allowbreak{}DivisorProofs.\allowbreak{}TwistedConjugacyClassesEquivUnitsModNormsNonspli}, lines 79--204)\\
\texttt{GlobalTheory.\allowbreak{}InertTwUnit.\allowbreak{}twisted\_\allowbreak{}unit\_\allowbreak{}eq\_\allowbreak{}stSum} (\texttt{Development.\allowbreak{}AbelianVarietyTheory.\allowbreak{}Proofs.\allowbreak{}DivisorProofs.\allowbreak{}InertPlaceStableTwistedOrbitalUnitEqLatticeSum}, lines 28--77)\par}

{\footnotesize\raggedright\noindent \ref{B6.04c}.\quad \texttt{PositivityTheory.\allowbreak{}Rg.\allowbreak{}BcB.\allowbreak{}E1.\allowbreak{}growthJ\_\allowbreak{}holds} (\texttt{Development.\allowbreak{}WeilBound.\allowbreak{}RepresentationTheoryGL.\allowbreak{}ArchimedeanBaseChange.\allowbreak{}ArchimedeanBaseChangeProof.\allowbreak{}ThetaCellCovers.\allowbreak{}GrowthOut}, lines 6)\par}

{\footnotesize\raggedright\noindent \ref{B6.04c}.\quad \texttt{PositivityTheory.\allowbreak{}Rg.\allowbreak{}BcB.\allowbreak{}E1.\allowbreak{}wtCount\_\allowbreak{}of} (\texttt{Development.\allowbreak{}WeilBound.\allowbreak{}RepresentationTheoryGL.\allowbreak{}ArchimedeanBaseChange.\allowbreak{}ArchimedeanBaseChangeProof.\allowbreak{}ThetaCellCovers.\allowbreak{}CutCountOf}, lines 56--94)\\
\texttt{PositivityTheory.\allowbreak{}Rg.\allowbreak{}BcB.\allowbreak{}E1.\allowbreak{}cutCount\_\allowbreak{}of} (\texttt{Development.\allowbreak{}WeilBound.\allowbreak{}RepresentationTheoryGL.\allowbreak{}ArchimedeanBaseChange.\allowbreak{}ArchimedeanBaseChangeProof.\allowbreak{}ThetaCellCovers.\allowbreak{}CutCountOf}, lines 96--115)\\
\texttt{PositivityTheory.\allowbreak{}Rg.\allowbreak{}BcB.\allowbreak{}E1.\allowbreak{}growthJ\_\allowbreak{}of} (\texttt{Development.\allowbreak{}WeilBound.\allowbreak{}RepresentationTheoryGL.\allowbreak{}ArchimedeanBaseChange.\allowbreak{}ArchimedeanBaseChangeProof.\allowbreak{}ThetaCellCovers.\allowbreak{}KernelDiagonalIntegrableOf}, lines 78--97)\par}

{\footnotesize\raggedright\noindent \ref{B6.04l}.\quad \texttt{PositivityTheory.\allowbreak{}RgBcC.\allowbreak{}lefschetz\_\allowbreak{}holds} (\texttt{Development.\allowbreak{}WeilBound.\allowbreak{}RepresentationTheoryGL.\allowbreak{}ArchimedeanBaseChange.\allowbreak{}ArchimedeanBaseChangeProof.\allowbreak{}ArchimedeanGelfandKernels.\allowbreak{}LeftEquivariantFunctions.\allowbreak{}LefschetzHolds}, lines 6)\par}

{\footnotesize\raggedright\noindent \ref{B6.04m}.\quad \texttt{PositivityTheory.\allowbreak{}Rg.\allowbreak{}BcB.\allowbreak{}G3.\allowbreak{}g3} (\texttt{Development.\allowbreak{}WeilBound.\allowbreak{}RepresentationTheoryGL.\allowbreak{}ArchimedeanBaseChange.\allowbreak{}ArchimedeanBaseChangeProof.\allowbreak{}ThetaNormCommutation.\allowbreak{}FactorFields.\allowbreak{}Count1}, lines 35--36)\\
\texttt{PositivityTheory.\allowbreak{}Rg.\allowbreak{}BcB.\allowbreak{}G8.\allowbreak{}global\_\allowbreak{}character\_\allowbreak{}holds} (\texttt{Development.\allowbreak{}WeilBound.\allowbreak{}RepresentationTheoryGL.\allowbreak{}ArchimedeanBaseChange.\allowbreak{}ArchimedeanBaseChangeProof.\allowbreak{}ThetaNormCommutation.\allowbreak{}LocalComparisonSplitting.\allowbreak{}GlobalCharacterHolds}, lines 6)\par}

{\footnotesize\raggedright\noindent \ref{B6.04da}.\quad \texttt{PositivityTheory.\allowbreak{}Rg.\allowbreak{}BcB.\allowbreak{}F.\allowbreak{}RunFrame.\allowbreak{}g} (\texttt{Development.\allowbreak{}WeilBound.\allowbreak{}RepresentationTheoryGL.\allowbreak{}ArchimedeanBaseChange.\allowbreak{}ArchimedeanBaseChangeProof.\allowbreak{}SatakeExponents.\allowbreak{}UniformizerValuations.\allowbreak{}FrameBoundSpaces.\allowbreak{}InstIPSFB}, lines 47--64)\\
\texttt{PositivityTheory.\allowbreak{}Rg.\allowbreak{}BcB.\allowbreak{}F.\allowbreak{}FrameOk.\allowbreak{}comm\_\allowbreak{}gen} (\texttt{Development.\allowbreak{}WeilBound.\allowbreak{}RepresentationTheoryGL.\allowbreak{}ArchimedeanBaseChange.\allowbreak{}ArchimedeanBaseChangeProof.\allowbreak{}SatakeExponents.\allowbreak{}UniformizerValuations.\allowbreak{}AdelicHeckeBlocks.\allowbreak{}PStruct}, lines 134--142)\\
\texttt{PositivityTheory.\allowbreak{}Rg.\allowbreak{}BcB.\allowbreak{}F.\allowbreak{}FrameOk.\allowbreak{}comm\_\allowbreak{}beta} (\texttt{Development.\allowbreak{}WeilBound.\allowbreak{}RepresentationTheoryGL.\allowbreak{}ArchimedeanBaseChange.\allowbreak{}ArchimedeanBaseChangeProof.\allowbreak{}SatakeExponents.\allowbreak{}UniformizerValuations.\allowbreak{}AdelicHeckeBlocks.\allowbreak{}PStruct}, lines 144--152)\par}

{\footnotesize\raggedright\noindent \ref{B6.04db}.\quad \texttt{PositivityTheory.\allowbreak{}Rg.\allowbreak{}BcB.\allowbreak{}F.\allowbreak{}Tail.\allowbreak{}sat\_\allowbreak{}of\_\allowbreak{}eig} (\texttt{Development.\allowbreak{}WeilBound.\allowbreak{}RepresentationTheoryGL.\allowbreak{}ArchimedeanBaseChange.\allowbreak{}ArchimedeanBaseChangeProof.\allowbreak{}SatakeExponents.\allowbreak{}UniformizerValuations.\allowbreak{}FrameReading.\allowbreak{}RExpReadOf}, lines 71--110)\\
\texttt{PositivityTheory.\allowbreak{}Rg.\allowbreak{}BcB.\allowbreak{}N3.\allowbreak{}Final.\allowbreak{}n3} (\texttt{Development.\allowbreak{}WeilBound.\allowbreak{}RepresentationTheoryGL.\allowbreak{}ArchimedeanBaseChange.\allowbreak{}ArchimedeanBaseChangeProof.\allowbreak{}SatakeExponents.\allowbreak{}SphericalEmbeddings.\allowbreak{}Imports}, lines 12--14)\par}

{\footnotesize\raggedright\noindent \ref{B6.04u}.\quad \texttt{PositivityTheory.\allowbreak{}Rg.\allowbreak{}BcB.\allowbreak{}SpectralAccumulation.\allowbreak{}FiniteDiscreteTypes.\allowbreak{}finDisc} (\texttt{Development.\allowbreak{}WeilBound.\allowbreak{}RepresentationTheoryGL.\allowbreak{}ArchimedeanBaseChange.\allowbreak{}ArchimedeanBaseChangeProof.\allowbreak{}InfinitesimalCharacterLimits.\allowbreak{}IsometrySymmetry.\allowbreak{}IsoVSymm}, lines 29--126)\par}

{\footnotesize\raggedright\noindent \ref{B6.04dc}.\quad \texttt{PositivityTheory.\allowbreak{}Rg.\allowbreak{}Bc.\allowbreak{}ArchbcB.\allowbreak{}B.\allowbreak{}M4.\allowbreak{}m4} (\texttt{Development.\allowbreak{}WeilBound.\allowbreak{}RepresentationTheoryGL.\allowbreak{}ArchimedeanBaseChange.\allowbreak{}ArchimedeanBaseChangeProof.\allowbreak{}IwahoriFactorisation.\allowbreak{}TwistedConjugationNorms.\allowbreak{}Development}, lines 1166--1177)\\
\texttt{GlobalTheory.\allowbreak{}CompanionTheory.\allowbreak{}MatchPow.\allowbreak{}pow\_\allowbreak{}claim} (\texttt{Development.\allowbreak{}AutomorphicSpectrum.\allowbreak{}Proofs.\allowbreak{}BlockDecomposition.\allowbreak{}MatchingTransferOf}, lines 456--493)\par}

{\footnotesize\raggedright\noindent \ref{B6.04d}.\quad \texttt{UpperTheory.\allowbreak{}BaseChangeIdentities.\allowbreak{}FrameWordIdentity.\allowbreak{}frameBound} (\texttt{Development.\allowbreak{}AutomorphicArithmetic.\allowbreak{}Foundations.\allowbreak{}BaseChangeIdentities.\allowbreak{}FrameMatching.\allowbreak{}HeckeWordBounds.\allowbreak{}FrameBound}, lines 42--43)\\
\texttt{PositivityTheory.\allowbreak{}Rg.\allowbreak{}BcB.\allowbreak{}F.\allowbreak{}B3.\allowbreak{}PTop.\allowbreak{}hW\_\allowbreak{}frame} (\texttt{Development.\allowbreak{}WeilBound.\allowbreak{}RepresentationTheoryGL.\allowbreak{}ArchimedeanBaseChange.\allowbreak{}ArchimedeanBaseChangeProof.\allowbreak{}SatakeExponents.\allowbreak{}UniformizerValuations.\allowbreak{}WordFrameBounds.\allowbreak{}FrameBoundOf}, lines 171--234)\par}

{\footnotesize\raggedright\noindent \ref{B6.04a}.\quad \texttt{PositivityTheory.\allowbreak{}Rg.\allowbreak{}BcN.\allowbreak{}IsBaseChangeN} (\texttt{Development.\allowbreak{}WeilBound.\allowbreak{}RepresentationTheoryGL.\allowbreak{}ArchimedeanBaseChange.\allowbreak{}ArchimedeanBaseChangeProof.\allowbreak{}BaseChangeIsobaricNOf}, lines 21--37)\\
\texttt{PositivityTheory.\allowbreak{}Rg.\allowbreak{}BcN.\allowbreak{}BaseChangeIsobaricN} (\texttt{Development.\allowbreak{}WeilBound.\allowbreak{}RepresentationTheoryGL.\allowbreak{}ArchimedeanBaseChange.\allowbreak{}ArchimedeanBaseChangeProof.\allowbreak{}BaseChangeIsobaricNOf}, lines 39--46)\\
\texttt{UpperTheory.\allowbreak{}BaseChangeIdentities.\allowbreak{}baseChangeN} (\texttt{Development.\allowbreak{}AutomorphicArithmetic.\allowbreak{}Foundations.\allowbreak{}BaseChangeIdentities.\allowbreak{}DefiniteNHolds}, lines 5--6)\par}

{\footnotesize\raggedright\noindent \ref{B6.04aa}.\quad \texttt{PositivityTheory.\allowbreak{}Rg.\allowbreak{}BcN.\allowbreak{}IsBaseChangeN} (\texttt{Development.\allowbreak{}WeilBound.\allowbreak{}RepresentationTheoryGL.\allowbreak{}ArchimedeanBaseChange.\allowbreak{}ArchimedeanBaseChangeProof.\allowbreak{}BaseChangeIsobaricNOf}, lines 21--37)\\
\texttt{GlobalTheory.\allowbreak{}UnitaryBC.\allowbreak{}infOfWeight} (\texttt{Development.\allowbreak{}AbelianVarietyTheory.\allowbreak{}PublishedTheorems.\allowbreak{}GaloisCohomology₂.\allowbreak{}BaseChangeDefiniteUnitary}, lines 142--143)\\
\texttt{PositivityTheory.\allowbreak{}Rg.\allowbreak{}Rs.\allowbreak{}AscEmb} (\texttt{Development.\allowbreak{}WeilBound.\allowbreak{}RepresentationTheoryGL.\allowbreak{}ResidualSpectrum.\allowbreak{}SliceTensors.\allowbreak{}IsSliceMap}, lines 53--55)\\
\texttt{UpperTheory.\allowbreak{}GalStar.\allowbreak{}Row0b} (\texttt{Development.\allowbreak{}AutomorphicArithmetic.\allowbreak{}Foundations.\allowbreak{}GaloisRepresentations.\allowbreak{}TypeConditions.\allowbreak{}WeightsConjBis}, lines 69--73)\\
\texttt{UpperTheory.\allowbreak{}Gal.\allowbreak{}Int.\allowbreak{}galoisLocalGlobalStrong\_\allowbreak{}recParameter} (\texttt{Development.\allowbreak{}AutomorphicArithmetic.\allowbreak{}Foundations.\allowbreak{}GaloisRepresentations.\allowbreak{}IntegralStructures.\allowbreak{}RowSrc}, lines 6--6)\\
\texttt{UpperTheory.\allowbreak{}GalStar.\allowbreak{}PropsB} (\texttt{Development.\allowbreak{}AutomorphicArithmetic.\allowbreak{}Foundations.\allowbreak{}GaloisRepresentations.\allowbreak{}TypeConditions.\allowbreak{}WeightsConjBis}, lines 54--68)\par}

{\footnotesize\raggedright\noindent \ref{B6.04g}.\quad \texttt{PositivityTheory.\allowbreak{}Rg.\allowbreak{}BcB.\allowbreak{}F.\allowbreak{}peak\_\allowbreak{}trace\_\allowbreak{}bound\_\allowbreak{}holds} (\texttt{Development.\allowbreak{}WeilBound.\allowbreak{}RepresentationTheoryGL.\allowbreak{}ArchimedeanBaseChange.\allowbreak{}ArchimedeanBaseChangeProof.\allowbreak{}SatakeExponents.\allowbreak{}UniformizerValuations.\allowbreak{}PeakTraceBoundHolds}, lines 368--412)\par}

{\footnotesize\raggedright\noindent \ref{B6.04fa}.\quad \texttt{PositivityTheory.\allowbreak{}Rg.\allowbreak{}BcB.\allowbreak{}F.\allowbreak{}Rd.\allowbreak{}B.\allowbreak{}realise\_\allowbreak{}commuting\_\allowbreak{}family} (\texttt{Development.\allowbreak{}WeilBound.\allowbreak{}RepresentationTheoryGL.\allowbreak{}ArchimedeanBaseChange.\allowbreak{}ArchimedeanBaseChangeProof.\allowbreak{}SatakeExponents.\allowbreak{}UniformizerValuations.\allowbreak{}FrameReading.\allowbreak{}LimitbAlg}, lines 52--115)\par}

{\footnotesize\raggedright\noindent \ref{B6.04f}.\quad \texttt{PositivityTheory.\allowbreak{}Rg.\allowbreak{}BcB.\allowbreak{}F.\allowbreak{}Rd.\allowbreak{}BlockOut} (\texttt{Development.\allowbreak{}WeilBound.\allowbreak{}RepresentationTheoryGL.\allowbreak{}ArchimedeanBaseChange.\allowbreak{}ArchimedeanBaseChangeProof.\allowbreak{}SatakeExponents.\allowbreak{}UniformizerValuations.\allowbreak{}FrameReading.\allowbreak{}ReadLimit}, lines 36--66)\par}

{\footnotesize\raggedright\noindent \ref{B6.04f}.\quad \texttt{PositivityTheory.\allowbreak{}Rg.\allowbreak{}BcB.\allowbreak{}SpectralAccumulation.\allowbreak{}eFinal\_\allowbreak{}of} (\texttt{Development.\allowbreak{}WeilBound.\allowbreak{}RepresentationTheoryGL.\allowbreak{}ArchimedeanBaseChange.\allowbreak{}ArchimedeanBaseChangeProof.\allowbreak{}InfinitesimalCharacterLimits.\allowbreak{}EFinalOf}, lines 11--14)\par}

{\footnotesize\raggedright\noindent \ref{B6.04f}.\quad \texttt{PositivityTheory.\allowbreak{}Rg.\allowbreak{}BcB.\allowbreak{}F.\allowbreak{}Fok.\allowbreak{}Mu3.\allowbreak{}muExists} (\texttt{Development.\allowbreak{}WeilBound.\allowbreak{}RepresentationTheoryGL.\allowbreak{}ArchimedeanBaseChange.\allowbreak{}ArchimedeanBaseChangeProof.\allowbreak{}SatakeExponents.\allowbreak{}UniformizerValuations.\allowbreak{}QuadraticCharacterNorms.\allowbreak{}QuadCharNorm}, lines 24--57)\par}

{\footnotesize\raggedright\noindent \ref{B6.04f}.\quad \texttt{PositivityTheory.\allowbreak{}Rg.\allowbreak{}BcB.\allowbreak{}F.\allowbreak{}Tail.\allowbreak{}rReadT} (\texttt{Development.\allowbreak{}WeilBound.\allowbreak{}RepresentationTheoryGL.\allowbreak{}ArchimedeanBaseChange.\allowbreak{}ArchimedeanBaseChangeProof.\allowbreak{}SatakeExponents.\allowbreak{}UniformizerValuations.\allowbreak{}FrameReading.\allowbreak{}PairDoubleCosetFiniteSupport}, lines 278--338)\\
\texttt{GlobalTheory.\allowbreak{}CompanionTheory.\allowbreak{}LinAlg.\allowbreak{}exists\_\allowbreak{}joint\_\allowbreak{}eigenvector\_\allowbreak{}combination\_\allowbreak{}ne\_\allowbreak{}zero} (\texttt{Development.\allowbreak{}AutomorphicSpectrum.\allowbreak{}Proofs.\allowbreak{}BlockDecomposition.\allowbreak{}ExistsJointEigenvectorNonzero}, lines 85--109)\par}

{\footnotesize\raggedright\noindent \ref{B6.04f}.\quad \texttt{PositivityTheory.\allowbreak{}Rg.\allowbreak{}BcB.\allowbreak{}F.\allowbreak{}PCandLh.\allowbreak{}inert\_\allowbreak{}d} (\texttt{Development.\allowbreak{}WeilBound.\allowbreak{}RepresentationTheoryGL.\allowbreak{}ArchimedeanBaseChange.\allowbreak{}ArchimedeanBaseChangeProof.\allowbreak{}SatakeExponents.\allowbreak{}UniformizerValuations.\allowbreak{}FrameReading.\allowbreak{}PCandLh}, lines 48--122)\\
\texttt{PositivityTheory.\allowbreak{}Rg.\allowbreak{}BcB.\allowbreak{}F.\allowbreak{}PCandLh.\allowbreak{}candLev\_\allowbreak{}of} (\texttt{Development.\allowbreak{}WeilBound.\allowbreak{}RepresentationTheoryGL.\allowbreak{}ArchimedeanBaseChange.\allowbreak{}ArchimedeanBaseChangeProof.\allowbreak{}SatakeExponents.\allowbreak{}UniformizerValuations.\allowbreak{}FrameReading.\allowbreak{}PCandLh}, lines 149--211)\\
\texttt{UpperTheory.\allowbreak{}BaseChangeIdentities.\allowbreak{}GelfandReading.\allowbreak{}frameRead} (\texttt{Development.\allowbreak{}AutomorphicArithmetic.\allowbreak{}Foundations.\allowbreak{}BaseChangeIdentities.\allowbreak{}FrameMatching.\allowbreak{}GelfandReading.\allowbreak{}FrameRead}, lines 16--17)\\
\texttt{PositivityTheory.\allowbreak{}Rg.\allowbreak{}BcB.\allowbreak{}F.\allowbreak{}runT\_\allowbreak{}of} (\texttt{Development.\allowbreak{}WeilBound.\allowbreak{}RepresentationTheoryGL.\allowbreak{}ArchimedeanBaseChange.\allowbreak{}ArchimedeanBaseChangeProof.\allowbreak{}SatakeExponents.\allowbreak{}UniformizerValuations.\allowbreak{}SetupEx}, lines 64--72)\\
\texttt{PositivityTheory.\allowbreak{}Rg.\allowbreak{}BcB.\allowbreak{}Plug2.\allowbreak{}bci\_\allowbreak{}of} (\texttt{Development.\allowbreak{}WeilBound.\allowbreak{}RepresentationTheoryGL.\allowbreak{}ArchimedeanBaseChange.\allowbreak{}ArchimedeanBaseChangeProof.\allowbreak{}SatakeExponents.\allowbreak{}BciOf}, lines 22--26)\par}

{\footnotesize\raggedright\noindent \ref{B6.04f}.\quad \texttt{PositivityTheory.\allowbreak{}Rg.\allowbreak{}BcB.\allowbreak{}F.\allowbreak{}F5.\allowbreak{}runFin\_\allowbreak{}of} (\texttt{Development.\allowbreak{}WeilBound.\allowbreak{}RepresentationTheoryGL.\allowbreak{}ArchimedeanBaseChange.\allowbreak{}ArchimedeanBaseChangeProof.\allowbreak{}SatakeExponents.\allowbreak{}FiniteDiscreteFamilies.\allowbreak{}RunFinOf}, lines 42--65)\par}

{\footnotesize\raggedright\noindent \ref{B6.04f}.\quad \texttt{PositivityTheory.\allowbreak{}Rg.\allowbreak{}BcB.\allowbreak{}One.\allowbreak{}fieldN\_\allowbreak{}of} (\texttt{Development.\allowbreak{}WeilBound.\allowbreak{}RepresentationTheoryGL.\allowbreak{}ArchimedeanBaseChange.\allowbreak{}ArchimedeanBaseChangeProof.\allowbreak{}FieldNOf}, lines 52--82)\\
\texttt{PositivityTheory.\allowbreak{}Rg.\allowbreak{}BcB.\allowbreak{}F.\allowbreak{}Rig.\allowbreak{}rCuspCsd} (\texttt{Development.\allowbreak{}WeilBound.\allowbreak{}RepresentationTheoryGL.\allowbreak{}ArchimedeanBaseChange.\allowbreak{}ArchimedeanBaseChangeProof.\allowbreak{}SatakeExponents.\allowbreak{}SatakeRigidity.\allowbreak{}LRCuspCsd}, lines 103--195)\par}

{\footnotesize\raggedright\noindent \ref{B6.05a}.\quad \texttt{UpperTheory.\allowbreak{}One.\allowbreak{}Join.\allowbreak{}local\_\allowbreak{}inputs\_\allowbreak{}from\_\allowbreak{}literature} (\texttt{Development.\allowbreak{}AutomorphicArithmetic.\allowbreak{}AutomorphicVarieties.\allowbreak{}ModularityArgument.\allowbreak{}ThresholdArguments.\allowbreak{}DenseThresholdCases.\allowbreak{}TopMainZj}, lines 18--37)\\
\texttt{GlobalTheory.\allowbreak{}CompanionTheory.\allowbreak{}Hypotheses₂.\allowbreak{}totallyDisjointProductIrreducible\_\allowbreak{}of\_\allowbreak{}polynomialRing} (\texttt{Development.\allowbreak{}AutomorphicSpectrum.\allowbreak{}Proofs.\allowbreak{}GlobalProofs.\allowbreak{}TotallyDisjointProductIrreducibleOfPolynomialRin}, lines 146--157)\par}

{\footnotesize\raggedright\noindent \ref{B6.05b}.\quad \texttt{GlobalTheory.\allowbreak{}CompanionTheory.\allowbreak{}MoyPrasadTop.\allowbreak{}moyPrasadJacquetIso} (\texttt{Development.\allowbreak{}AutomorphicSpectrum.\allowbreak{}Proofs.\allowbreak{}BlockDecomposition.\allowbreak{}MoyPrasadJacquetIso}, lines 16--32)\\
\texttt{GlobalTheory.\allowbreak{}DepthZeroPass.\allowbreak{}depthZeroInduction} (\texttt{Development.\allowbreak{}AbelianVarietyTheory.\allowbreak{}Proofs.\allowbreak{}InertiaEigenvalues.\allowbreak{}FixedNeBotOfIsSubquotient}, lines 64--82)\\
\texttt{DiscreteGroups.\allowbreak{}sglob.\allowbreak{}LConv.\allowbreak{}localConverseGL} (\texttt{Development.\allowbreak{}DiscreteGroups.\allowbreak{}GlobalIrreducibility.\allowbreak{}LConvFinal}, lines 18--21)\\
\texttt{GlobalTheory.\allowbreak{}BrauerZero.\allowbreak{}brauerDegreeZero} (\texttt{Development.\allowbreak{}AbelianVarietyTheory.\allowbreak{}Proofs.\allowbreak{}InertiaEigenvalues.\allowbreak{}BrauerDegreeZero}, lines 203--216)\\
\texttt{UpperTheory.\allowbreak{}LiteratureTheorems.\allowbreak{}chevalley\_\allowbreak{}units\_\allowbreak{}of\_\allowbreak{}open} (\texttt{Development.\allowbreak{}AutomorphicArithmetic.\allowbreak{}AutomorphicVarieties.\allowbreak{}ModularityArgument.\allowbreak{}ThresholdArguments.\allowbreak{}DenseThresholdCases.\allowbreak{}TemperedProductOfSquareIntegrable}, lines 253--254)\par}

{\footnotesize\raggedright\noindent \ref{B6.05c}.\quad \texttt{UpperTheory.\allowbreak{}One.\allowbreak{}Join.\allowbreak{}automorphic\_\allowbreak{}inputs\_\allowbreak{}from\_\allowbreak{}literature} (\texttt{Development.\allowbreak{}AutomorphicArithmetic.\allowbreak{}AutomorphicVarieties.\allowbreak{}ModularityArgument.\allowbreak{}ThresholdArguments.\allowbreak{}DenseThresholdCases.\allowbreak{}TopMainZj}, lines 38--54)\\
\texttt{UpperTheory.\allowbreak{}LiteratureTheorems.\allowbreak{}hecke\_\allowbreak{}ladic\_\allowbreak{}character\_\allowbreak{}above\_\allowbreak{}of\_\allowbreak{}open} (\texttt{Development.\allowbreak{}AutomorphicArithmetic.\allowbreak{}AutomorphicVarieties.\allowbreak{}ModularityArgument.\allowbreak{}ThresholdArguments.\allowbreak{}DenseThresholdCases.\allowbreak{}TemperedProductOfSquareIntegrable}, lines 295--296)\\
\texttt{UpperTheory.\allowbreak{}LiteratureTheorems.\allowbreak{}clozel\_\allowbreak{}purity\_\allowbreak{}of\_\allowbreak{}open} (\texttt{Development.\allowbreak{}AutomorphicArithmetic.\allowbreak{}AutomorphicVarieties.\allowbreak{}ModularityArgument.\allowbreak{}ThresholdArguments.\allowbreak{}DenseThresholdCases.\allowbreak{}TemperedProductOfSquareIntegrable}, lines 279--280)\par}

{\footnotesize\raggedright\noindent \ref{B6.05d}.\quad \texttt{UpperTheory.\allowbreak{}LiteratureTheorems.\allowbreak{}galois\_\allowbreak{}cohomology\_\allowbreak{}from\_\allowbreak{}literature} (\texttt{Development.\allowbreak{}AutomorphicArithmetic.\allowbreak{}AutomorphicVarieties.\allowbreak{}ModularityArgument.\allowbreak{}ThresholdArguments.\allowbreak{}DenseThresholdCases.\allowbreak{}TemperedProductOfSquareIntegrable}, lines 331--343)\\
\texttt{ReciprocityTheory.\allowbreak{}Gal.\allowbreak{}UnramifiedOrthogonalityAt} (\texttt{ReciprocityTheory.\allowbreak{}Hypotheses.\allowbreak{}GalCoh}, lines 36--53)\\
\texttt{GlobalTheory.\allowbreak{}unramifiedOrthogonality\_\allowbreak{}holds} (\texttt{Development.\allowbreak{}AbelianVarietyTheory.\allowbreak{}Proofs.\allowbreak{}InertiaEigenvalues.\allowbreak{}UnramifiedOrthogonalityHolds}, lines 7--9)\par}

{\footnotesize\raggedright\noindent \ref{B6.05e}.\quad \texttt{UpperTheory.\allowbreak{}One.\allowbreak{}Join.\allowbreak{}cm\_\allowbreak{}abelian\_\allowbreak{}inputs\_\allowbreak{}from\_\allowbreak{}literature} (\texttt{Development.\allowbreak{}AutomorphicArithmetic.\allowbreak{}AutomorphicVarieties.\allowbreak{}ModularityArgument.\allowbreak{}ThresholdArguments.\allowbreak{}DenseThresholdCases.\allowbreak{}TopMainZj}, lines 59--70)\\
\texttt{GlobalTheory.\allowbreak{}Published.\allowbreak{}grothendieckDeformation\_\allowbreak{}of\_\allowbreak{}lifting} (\texttt{Development.\allowbreak{}AbelianVarietyTheory.\allowbreak{}Proofs.\allowbreak{}DivisorProofs.\allowbreak{}GrothendieckDeformationOfLifting}, lines 9--13)\\
\texttt{GlobalTheory.\allowbreak{}AnPer.\allowbreak{}pTaylor} (\texttt{Development.\allowbreak{}AbelianVarietyTheory.\allowbreak{}Proofs.\allowbreak{}DivisorProofs.\allowbreak{}ConstantCycleExtensionTaylorPeriodOfThickeningCo}, lines 26--56)\par}

{\footnotesize\raggedright\noindent \ref{B6.05f}.\quad \texttt{UpperTheory.\allowbreak{}One.\allowbreak{}Join.\allowbreak{}cm\_\allowbreak{}abelian\_\allowbreak{}inputs\_\allowbreak{}from\_\allowbreak{}literature} (\texttt{Development.\allowbreak{}AutomorphicArithmetic.\allowbreak{}AutomorphicVarieties.\allowbreak{}ModularityArgument.\allowbreak{}ThresholdArguments.\allowbreak{}DenseThresholdCases.\allowbreak{}TopMainZj}, lines 59--70)\par}

{\footnotesize\raggedright\noindent \ref{def:p6c1-classical-inputs}.\quad \texttt{GlobalTheory.\allowbreak{}Mem.\allowbreak{}LiteralAutomorphy.\allowbreak{}LitLocAlt} (\texttt{Development.\allowbreak{}AutomorphicArithmetic.\allowbreak{}AutomorphicVarieties.\allowbreak{}ModularityArgument.\allowbreak{}ThresholdArguments.\allowbreak{}DenseThresholdCases.\allowbreak{}LitAut}, lines 12--30)\\
\texttt{GlobalTheory.\allowbreak{}Mem.\allowbreak{}LiteralAutomorphy.\allowbreak{}LiteratureAutomorphyHypothesis} (\texttt{Development.\allowbreak{}AutomorphicArithmetic.\allowbreak{}AutomorphicVarieties.\allowbreak{}ModularityArgument.\allowbreak{}ThresholdArguments.\allowbreak{}DenseThresholdCases.\allowbreak{}LitAut}, lines 31--46)\\
\texttt{GlobalTheory.\allowbreak{}Mem.\allowbreak{}LiteralAutomorphy.\allowbreak{}LitCoh} (\texttt{Development.\allowbreak{}AbelianVarietyTheory.\allowbreak{}Definitions.\allowbreak{}AutomorphicMembership.\allowbreak{}Brauer}, lines 105--117)\\
\texttt{GlobalTheory.\allowbreak{}Mem.\allowbreak{}LiteralAutomorphy.\allowbreak{}LiteratureCMHypothesis} (\texttt{Development.\allowbreak{}AutomorphicArithmetic.\allowbreak{}AutomorphicVarieties.\allowbreak{}ModularityArgument.\allowbreak{}ThresholdArguments.\allowbreak{}DenseThresholdCases.\allowbreak{}LitAut}, lines 47--58)\\
\texttt{GlobalTheory.\allowbreak{}Mem.\allowbreak{}LiteralAutomorphy.\allowbreak{}LiteratureRankNHypothesis} (\texttt{Development.\allowbreak{}AutomorphicArithmetic.\allowbreak{}AutomorphicVarieties.\allowbreak{}ModularityArgument.\allowbreak{}ThresholdArguments.\allowbreak{}DenseThresholdCases.\allowbreak{}LitAut}, lines 59--63)\par}

{\footnotesize\raggedright\noindent \ref{B6.05ca}.\quad \texttt{UpperTheory.\allowbreak{}One.\allowbreak{}Join.\allowbreak{}automorphic\_\allowbreak{}window\_\allowbreak{}of\_\allowbreak{}inputs} (\texttt{Development.\allowbreak{}AutomorphicArithmetic.\allowbreak{}AutomorphicVarieties.\allowbreak{}ModularityArgument.\allowbreak{}ThresholdArguments.\allowbreak{}DenseThresholdCases.\allowbreak{}TopMainZj}, lines 71--74)\\
\texttt{UpperTheory.\allowbreak{}Gal.\allowbreak{}Int.\allowbreak{}galoisLocalGlobalStrong\_\allowbreak{}recParameter} (\texttt{Development.\allowbreak{}AutomorphicArithmetic.\allowbreak{}Foundations.\allowbreak{}GaloisRepresentations.\allowbreak{}IntegralStructures.\allowbreak{}RowSrc}, lines 6--7)\par}

{\footnotesize\raggedright\noindent \ref{B6.05}.\quad \texttt{UpperTheory.\allowbreak{}One.\allowbreak{}Join.\allowbreak{}automorphic\_\allowbreak{}window\_\allowbreak{}of\_\allowbreak{}inputs} (\texttt{Development.\allowbreak{}AutomorphicArithmetic.\allowbreak{}AutomorphicVarieties.\allowbreak{}ModularityArgument.\allowbreak{}ThresholdArguments.\allowbreak{}DenseThresholdCases.\allowbreak{}TopMainZj}, lines 71--74)\\
\texttt{GlobalTheory.\allowbreak{}Mem.\allowbreak{}LiteralAutomorphy.\allowbreak{}LiteratureRankNHypothesis} (\texttt{Development.\allowbreak{}AutomorphicArithmetic.\allowbreak{}AutomorphicVarieties.\allowbreak{}ModularityArgument.\allowbreak{}ThresholdArguments.\allowbreak{}DenseThresholdCases.\allowbreak{}LitAut}, lines 59--63)\\
\texttt{GlobalTheory.\allowbreak{}Mem.\allowbreak{}LiteralAutomorphy.\allowbreak{}LitCoh} (\texttt{Development.\allowbreak{}AbelianVarietyTheory.\allowbreak{}Definitions.\allowbreak{}AutomorphicMembership.\allowbreak{}Brauer}, lines 105--117)\par}

{\footnotesize\raggedright\noindent \ref{B6.05}.\quad \texttt{UpperTheory.\allowbreak{}One.\allowbreak{}Join.\allowbreak{}cm\_\allowbreak{}abelian\_\allowbreak{}inputs\_\allowbreak{}from\_\allowbreak{}literature} (\texttt{Development.\allowbreak{}AutomorphicArithmetic.\allowbreak{}AutomorphicVarieties.\allowbreak{}ModularityArgument.\allowbreak{}ThresholdArguments.\allowbreak{}DenseThresholdCases.\allowbreak{}TopMainZj}, lines 59--70)\par}

{\footnotesize\raggedright\noindent \ref{B6.06a}.\quad \texttt{UpperTheory.\allowbreak{}V.\allowbreak{}TR.\allowbreak{}Cen.\allowbreak{}exists\_\allowbreak{}cap\_\allowbreak{}scalar} (\texttt{Development.\allowbreak{}AutomorphicArithmetic.\allowbreak{}AutomorphicVarieties.\allowbreak{}AutomorphicForms.\allowbreak{}CapCentral}, lines 153--158)\\
\texttt{UpperTheory.\allowbreak{}V.\allowbreak{}Af.\allowbreak{}cuspFormsFiniteReal} (\texttt{Development.\allowbreak{}AutomorphicArithmetic.\allowbreak{}AutomorphicVarieties.\allowbreak{}AutomorphicForms.\allowbreak{}TemperedLocalComponents.\allowbreak{}PCuspFin}, lines 31--83)\\
\texttt{UpperTheory.\allowbreak{}V.\allowbreak{}Af.\allowbreak{}admissibleReal} (\texttt{Development.\allowbreak{}AutomorphicArithmetic.\allowbreak{}AutomorphicVarieties.\allowbreak{}AutomorphicForms.\allowbreak{}TemperedLocalComponents.\allowbreak{}ExistsContRep}, lines 39--63)\par}

{\footnotesize\raggedright\noindent \ref{B6.06b}.\quad \texttt{UpperTheory.\allowbreak{}V.\allowbreak{}Af.\allowbreak{}clozelCohC\_\allowbreak{}alt} (\texttt{Development.\allowbreak{}AutomorphicArithmetic.\allowbreak{}AutomorphicVarieties.\allowbreak{}ModularityArgument.\allowbreak{}ThresholdArguments.\allowbreak{}DenseThresholdCases.\allowbreak{}TPCohC}, lines 10--30)\\
\texttt{UpperTheory.\allowbreak{}V.\allowbreak{}Z.\allowbreak{}Coh.\allowbreak{}clozelConjC} (\texttt{Development.\allowbreak{}AutomorphicArithmetic.\allowbreak{}AutomorphicVarieties.\allowbreak{}ModularityArgument.\allowbreak{}ClozelCohomology.\allowbreak{}ClozelConjC}, lines 1--32)\\
\texttt{UpperTheory.\allowbreak{}V.\allowbreak{}Lit.\allowbreak{}Geo.\allowbreak{}borelInjC} (\texttt{Development.\allowbreak{}AutomorphicArithmetic.\allowbreak{}AutomorphicVarieties.\allowbreak{}LiteratureResults.\allowbreak{}LieGeometryGrowth.\allowbreak{}BorelInjC}, lines 74--97)\par}

{\footnotesize\raggedright\noindent \ref{B6.06c}.\quad \texttt{GlobalTheory.\allowbreak{}T.\allowbreak{}galois\_\allowbreak{}rep\_\allowbreak{}patching\_\allowbreak{}infinite\_\allowbreak{}family} (\texttt{PositivityTheory.\allowbreak{}AutomorphicForms.\allowbreak{}PublishedInputs.\allowbreak{}RealGaloisRepresentations.\allowbreak{}IwahoriSphericalTransfer.\allowbreak{}GaloisRepPatchingInfiniteFamily}, lines 12--13)\par}

{\footnotesize\raggedright\noindent \ref{B6.06d}.\quad \texttt{UpperTheory.\allowbreak{}V.\allowbreak{}HV.\allowbreak{}DescentVanishingReductions.\allowbreak{}published\_\allowbreak{}hypotheses\_\allowbreak{}from\_\allowbreak{}patching\_\allowbreak{}alt} (\texttt{Development.\allowbreak{}AutomorphicArithmetic.\allowbreak{}AutomorphicVarieties.\allowbreak{}ModularityArgument.\allowbreak{}ThresholdArguments.\allowbreak{}DenseThresholdCases.\allowbreak{}PublishedHypothesesFromPatching}, lines 98--103)\\
\texttt{UpperTheory.\allowbreak{}V.\allowbreak{}Af.\allowbreak{}newformRepOf} (\texttt{Development.\allowbreak{}AutomorphicArithmetic.\allowbreak{}AutomorphicVarieties.\allowbreak{}AutomorphicForms.\allowbreak{}TemperedLocalComponents.\allowbreak{}NewformRepOf}, lines 114--139)\\
\texttt{UpperTheory.\allowbreak{}ConvolutionTheory.\allowbreak{}RankTwoDescent.\allowbreak{}descS} (\texttt{Development.\allowbreak{}AutomorphicArithmetic.\allowbreak{}Foundations.\allowbreak{}Convolution.\allowbreak{}SplitEigenfunctions.\allowbreak{}AbsNormConj}, lines 40--92)\par}

{\footnotesize\raggedright\noindent \ref{B6.06e}.\quad \texttt{UpperTheory.\allowbreak{}Gal.\allowbreak{}SV.\allowbreak{}galois\_\allowbreak{}rep\_\allowbreak{}exists\_\allowbreak{}totally\_\allowbreak{}real\_\allowbreak{}rank\_\allowbreak{}two\_\allowbreak{}of\_\allowbreak{}base\_\allowbreak{}change} (\texttt{Development.\allowbreak{}AutomorphicArithmetic.\allowbreak{}AutomorphicVarieties.\allowbreak{}ModularityArgument.\allowbreak{}ThresholdArguments.\allowbreak{}DenseThresholdCases.\allowbreak{}SVPatch}, lines 103--165)\par}

{\footnotesize\raggedright\noindent \ref{B6.06f}.\quad \texttt{UpperTheory.\allowbreak{}V.\allowbreak{}Z.\allowbreak{}Bc.\allowbreak{}isAutR\_\allowbreak{}of\_\allowbreak{}isAutC\_\allowbreak{}comp\_\allowbreak{}galEmb\_\allowbreak{}of\_\allowbreak{}odd\_\allowbreak{}quadratic\_\allowbreak{}char} (\texttt{Development.\allowbreak{}AutomorphicArithmetic.\allowbreak{}AutomorphicVarieties.\allowbreak{}ModularityArgument.\allowbreak{}ThresholdArguments.\allowbreak{}DenseThresholdCases.\allowbreak{}RankTwoDescentSplit}, lines 88--114)\\
\texttt{UpperTheory.\allowbreak{}V.\allowbreak{}Z.\allowbreak{}Bc.\allowbreak{}stepC\_\allowbreak{}alt} (\texttt{Development.\allowbreak{}AutomorphicArithmetic.\allowbreak{}AutomorphicVarieties.\allowbreak{}ModularityArgument.\allowbreak{}ThresholdArguments.\allowbreak{}DenseThresholdCases.\allowbreak{}CuspUnique}, lines 24--65)\\
\texttt{UpperTheory.\allowbreak{}V.\allowbreak{}Z.\allowbreak{}Bc.\allowbreak{}exDownC\_\allowbreak{}sv\_\allowbreak{}alt} (\texttt{Development.\allowbreak{}AutomorphicArithmetic.\allowbreak{}AutomorphicVarieties.\allowbreak{}ModularityArgument.\allowbreak{}ThresholdArguments.\allowbreak{}DenseThresholdCases.\allowbreak{}SVRSol}, lines 16--62)\\
\texttt{UpperTheory.\allowbreak{}V.\allowbreak{}Z.\allowbreak{}Bc.\allowbreak{}automorphic\_\allowbreak{}rank\_\allowbreak{}two\_\allowbreak{}of\_\allowbreak{}automorphic\_\allowbreak{}restrict\_\allowbreak{}of\_\allowbreak{}finite\_\allowbreak{}det} (\texttt{Development.\allowbreak{}AutomorphicArithmetic.\allowbreak{}AutomorphicVarieties.\allowbreak{}ModularityArgument.\allowbreak{}ThresholdArguments.\allowbreak{}DenseThresholdCases.\allowbreak{}DownQuadS}, lines 70--122)\par}

{\footnotesize\raggedright\noindent \ref{B6.06}.\quad \texttt{UpperTheory.\allowbreak{}V.\allowbreak{}HV.\allowbreak{}DescentVanishingReductions.\allowbreak{}published\_\allowbreak{}hypotheses\_\allowbreak{}from\_\allowbreak{}patching\_\allowbreak{}alt} (\texttt{Development.\allowbreak{}AutomorphicArithmetic.\allowbreak{}AutomorphicVarieties.\allowbreak{}ModularityArgument.\allowbreak{}ThresholdArguments.\allowbreak{}DenseThresholdCases.\allowbreak{}PublishedHypothesesFromPatching}, lines 98--103)\\
\texttt{UpperTheory.\allowbreak{}V.\allowbreak{}HV.\allowbreak{}DescentVanishingReductions.\allowbreak{}vwin\_\allowbreak{}of\_\allowbreak{}hypotheses\_\allowbreak{}H\_\allowbreak{}E\_\allowbreak{}alt} (\texttt{Development.\allowbreak{}AutomorphicArithmetic.\allowbreak{}AutomorphicVarieties.\allowbreak{}ModularityArgument.\allowbreak{}ThresholdArguments.\allowbreak{}DenseThresholdCases.\allowbreak{}WindowBoundOfLiteratureLevel}, lines 25--29)\par}

{\footnotesize\raggedright\noindent \ref{B6.06}.\quad \texttt{UpperTheory.\allowbreak{}V.\allowbreak{}Z.\allowbreak{}Tf.\allowbreak{}SR.\allowbreak{}descWeak} (\texttt{Development.\allowbreak{}AutomorphicArithmetic.\allowbreak{}AutomorphicVarieties.\allowbreak{}ModularityArgument.\allowbreak{}TraceFormula.\allowbreak{}SplitRigidity.\allowbreak{}HasOfSplit}, lines 29--30)\\
\texttt{UpperTheory.\allowbreak{}V.\allowbreak{}Z.\allowbreak{}Bc.\allowbreak{}cycC\_\allowbreak{}holds} (\texttt{Development.\allowbreak{}AutomorphicArithmetic.\allowbreak{}AutomorphicVarieties.\allowbreak{}ModularityArgument.\allowbreak{}BaseChangeUnits.\allowbreak{}CycFree}, lines 87--88)\par}

{\footnotesize\raggedright\noindent \ref{B6.07}.\quad \texttt{UpperTheory.\allowbreak{}V.\allowbreak{}HV.\allowbreak{}Core.\allowbreak{}of\_\allowbreak{}four\_\allowbreak{}alt} (\texttt{Development.\allowbreak{}AutomorphicArithmetic.\allowbreak{}AutomorphicVarieties.\allowbreak{}ModularityArgument.\allowbreak{}ThresholdArguments.\allowbreak{}DenseThresholdCases.\allowbreak{}HVCore}, lines 18--22)\par}

\section*{Chapter~\ref{ch:vwin-three-length-bounds}: Three length bounds}

{\footnotesize\raggedright\noindent \ref{B6.08}.\quad \texttt{UpperTheory.\allowbreak{}V.\allowbreak{}shaLen\_\allowbreak{}eq\_\allowbreak{}add} (\texttt{Development.\allowbreak{}AutomorphicArithmetic.\allowbreak{}AutomorphicVarieties.\allowbreak{}LengthEqAddEigen}, lines 88--104)\\
\texttt{UpperTheory.\allowbreak{}V.\allowbreak{}vvle\_\allowbreak{}of\_\allowbreak{}parts} (\texttt{Development.\allowbreak{}AutomorphicArithmetic.\allowbreak{}AutomorphicVarieties.\allowbreak{}LengthEqAddEigen}, lines 127--190)\par}

{\footnotesize\raggedright\noindent \ref{B6.09}.\quad \texttt{GlobalTheory.\allowbreak{}Published.\allowbreak{}PoitouTateDuality} (\texttt{Development.\allowbreak{}AbelianVarietyTheory.\allowbreak{}PublishedTheorems.\allowbreak{}GaloisCohomology₂.\allowbreak{}ExistsPrimitiveRootPow}, lines 13--28)\\
\texttt{UpperTheory.\allowbreak{}V.\allowbreak{}HV.\allowbreak{}U.\allowbreak{}pt\_\allowbreak{}s} (\texttt{Development.\allowbreak{}AutomorphicArithmetic.\allowbreak{}AutomorphicVarieties.\allowbreak{}ModularityArgument.\allowbreak{}ThresholdArguments.\allowbreak{}DenseThresholdCases.\allowbreak{}HVFinalU}, lines 11--11)\par}

{\footnotesize\raggedright\noindent \ref{p6c2:bernoulli}.\quad \texttt{UpperTheory.\allowbreak{}V.\allowbreak{}Bern.\allowbreak{}b\_\allowbreak{}primmem} (\texttt{Development.\allowbreak{}AutomorphicArithmetic.\allowbreak{}AutomorphicVarieties.\allowbreak{}GaloisLibrary.\allowbreak{}BernoulliNumbers.\allowbreak{}Primmem}, lines 17--55)\\
\texttt{UpperTheory.\allowbreak{}V.\allowbreak{}Bern.\allowbreak{}b\_\allowbreak{}nzC} (\texttt{Development.\allowbreak{}AutomorphicArithmetic.\allowbreak{}AutomorphicVarieties.\allowbreak{}GaloisLibrary.\allowbreak{}BernoulliNumbers.\allowbreak{}BNzC}, lines 10--70)\\
\texttt{UpperTheory.\allowbreak{}V.\allowbreak{}Bern.\allowbreak{}b\_\allowbreak{}denupper} (\texttt{Development.\allowbreak{}AutomorphicArithmetic.\allowbreak{}AutomorphicVarieties.\allowbreak{}GaloisLibrary.\allowbreak{}BernoulliNumbers.\allowbreak{}Denupper}, lines 28--39)\par}

{\footnotesize\raggedright\noindent \ref{B6.10}.\quad \texttt{UpperTheory.\allowbreak{}V.\allowbreak{}ML.\allowbreak{}iv} (\texttt{Development.\allowbreak{}AutomorphicArithmetic.\allowbreak{}AutomorphicVarieties.\allowbreak{}PrimitiveFrobeniusFixed.\allowbreak{}Iv}, lines 4--7)\par}

{\footnotesize\raggedright\noindent \ref{B6.10a}.\quad \texttt{UpperTheory.\allowbreak{}V.\allowbreak{}Stick.\allowbreak{}coef} (\texttt{Development.\allowbreak{}AutomorphicArithmetic.\allowbreak{}AutomorphicVarieties.\allowbreak{}GaloisLibrary.\allowbreak{}StickelbergerTheorem.\allowbreak{}ImportsAlt}, lines 9--17)\\
\texttt{UpperTheory.\allowbreak{}V.\allowbreak{}Odd.\allowbreak{}clKill} (\texttt{Development.\allowbreak{}AutomorphicArithmetic.\allowbreak{}AutomorphicVarieties.\allowbreak{}OddResidueBounds.\allowbreak{}ClKill}, lines 17--56)\\
\texttt{UpperTheory.\allowbreak{}V.\allowbreak{}Odd.\allowbreak{}homKill} (\texttt{Development.\allowbreak{}AutomorphicArithmetic.\allowbreak{}AutomorphicVarieties.\allowbreak{}OddResidueBounds.\allowbreak{}HomKill}, lines 9--50)\par}

{\footnotesize\raggedright\noindent \ref{B6.11}.\quad \texttt{UpperTheory.\allowbreak{}V.\allowbreak{}Odd.\allowbreak{}primeBound\_\allowbreak{}of} (\texttt{Development.\allowbreak{}AutomorphicArithmetic.\allowbreak{}AutomorphicVarieties.\allowbreak{}OddResidueBounds.\allowbreak{}ThmBcOfPsiBound}, lines 293--294)\par}

{\footnotesize\raggedright\noindent \ref{B6.12}.\quad \texttt{UpperTheory.\allowbreak{}V.\allowbreak{}Odd.\allowbreak{}killPow} (\texttt{Development.\allowbreak{}AutomorphicArithmetic.\allowbreak{}AutomorphicVarieties.\allowbreak{}OddResidueBounds.\allowbreak{}KillPow}, lines 31--105)\\
\texttt{UpperTheory.\allowbreak{}V.\allowbreak{}Odd.\allowbreak{}stickBound} (\texttt{Development.\allowbreak{}AutomorphicArithmetic.\allowbreak{}AutomorphicVarieties.\allowbreak{}OddResidueBounds.\allowbreak{}ThmBcOfPsiBound}, lines 291--292)\\
\texttt{UpperTheory.\allowbreak{}V.\allowbreak{}Odd.\allowbreak{}deg} (\texttt{Development.\allowbreak{}AutomorphicArithmetic.\allowbreak{}AutomorphicVarieties.\allowbreak{}OddResidueBounds.\allowbreak{}DegExistsLift}, lines 33--100)\par}

{\footnotesize\raggedright\noindent \ref{B6.13}.\quad \texttt{UpperTheory.\allowbreak{}V.\allowbreak{}Odd.\allowbreak{}thmBcFree\_\allowbreak{}of} (\texttt{Development.\allowbreak{}AutomorphicArithmetic.\allowbreak{}AutomorphicVarieties.\allowbreak{}OddResidueBounds.\allowbreak{}ThmBcOfPsiBound}, lines 295--296)\\
\texttt{UpperTheory.\allowbreak{}V.\allowbreak{}minusBound\_\allowbreak{}le\_\allowbreak{}exp\_\allowbreak{}of\_\allowbreak{}log\_\allowbreak{}le} (\texttt{Development.\allowbreak{}AutomorphicArithmetic.\allowbreak{}AutomorphicVarieties.\allowbreak{}CMFieldExponentialBounds.\allowbreak{}minusle}, lines 6--76)\par}

{\footnotesize\raggedright\noindent \ref{p6c2:joint-window}.\quad \texttt{UpperTheory.\allowbreak{}V.\allowbreak{}JointWin} (\texttt{Development.\allowbreak{}AutomorphicArithmetic.\allowbreak{}AutomorphicVarieties.\allowbreak{}PrimeSetBoundDataOfPrimitiveChar}, lines 88--94)\par}

{\footnotesize\raggedright\noindent \ref{B6.14}.\quad \texttt{UpperTheory.\allowbreak{}V.\allowbreak{}vwinAt\_\allowbreak{}of\_\allowbreak{}joint} (\texttt{Development.\allowbreak{}AutomorphicArithmetic.\allowbreak{}AutomorphicVarieties.\allowbreak{}CMFieldExponentialBounds.\allowbreak{}AdmissibleOfPageN}, lines 41--130)\\
\texttt{UpperTheory.\allowbreak{}V.\allowbreak{}log\_\allowbreak{}conductor\_\allowbreak{}mul\_\allowbreak{}height\_\allowbreak{}le\_\allowbreak{}of\_\allowbreak{}pageN} (\texttt{Development.\allowbreak{}AutomorphicArithmetic.\allowbreak{}AutomorphicVarieties.\allowbreak{}CMFieldExponentialBounds.\allowbreak{}LogConductorMulHeightLeOfPageN}, lines 6--57)\\
\texttt{UpperTheory.\allowbreak{}V.\allowbreak{}error\_\allowbreak{}terms\_\allowbreak{}le\_\allowbreak{}sqrt\_\allowbreak{}exp\_\allowbreak{}rpow} (\texttt{Development.\allowbreak{}AutomorphicArithmetic.\allowbreak{}AutomorphicVarieties.\allowbreak{}CMFieldExponentialBounds.\allowbreak{}ErrorTermsLeSqrtExpRpow}, lines 6--89)\par}

\section*{Chapter~\ref{ch:B6-cyclotomic-window}: The cyclotomic tower and a common window}

{\footnotesize\raggedright\noindent \ref{B6.15}.\quad \texttt{UpperTheory.\allowbreak{}V.\allowbreak{}ZK.\allowbreak{}lemmaK\_\allowbreak{}litN\_\allowbreak{}alt} (\texttt{Development.\allowbreak{}AutomorphicArithmetic.\allowbreak{}AutomorphicVarieties.\allowbreak{}ModularityArgument.\allowbreak{}ThresholdArguments.\allowbreak{}DenseThresholdCases.\allowbreak{}KFinal}, lines 11--11)\\
\texttt{UpperTheory.\allowbreak{}V.\allowbreak{}LemmaK} (\texttt{Development.\allowbreak{}AutomorphicArithmetic.\allowbreak{}AutomorphicVarieties.\allowbreak{}GeneratorCountTuples.\allowbreak{}LemmaK}, lines 13--18)\par}

{\footnotesize\raggedright\noindent \ref{B6.15}.\quad \texttt{UpperTheory.\allowbreak{}V.\allowbreak{}ZK.\allowbreak{}sFinal} (\texttt{Development.\allowbreak{}AutomorphicArithmetic.\allowbreak{}AutomorphicVarieties.\allowbreak{}ModularityArgument.\allowbreak{}ClassCounting.\allowbreak{}SFinal}, lines 100--163)\\
\texttt{UpperTheory.\allowbreak{}V.\allowbreak{}LDataCardLeFinrank.\allowbreak{}vFormula\_\allowbreak{}of\_\allowbreak{}classicalTheorems} (\texttt{Development.\allowbreak{}AutomorphicArithmetic.\allowbreak{}AutomorphicVarieties.\allowbreak{}GaloisLibrary.\allowbreak{}GaloisCohomologyLength.\allowbreak{}Formula}, lines 15--24)\par}

{\footnotesize\raggedright\noindent \ref{B6.16}.\quad \texttt{UpperTheory.\allowbreak{}V.\allowbreak{}ZK.\allowbreak{}sCount} (\texttt{Development.\allowbreak{}AutomorphicArithmetic.\allowbreak{}AutomorphicVarieties.\allowbreak{}ModularityArgument.\allowbreak{}ClassCounting.\allowbreak{}SelmerOrderCounts}, lines 16--111)\\
\texttt{UpperTheory.\allowbreak{}V.\allowbreak{}ZK.\allowbreak{}sNum} (\texttt{Development.\allowbreak{}AutomorphicArithmetic.\allowbreak{}AutomorphicVarieties.\allowbreak{}ModularityArgument.\allowbreak{}ClassCounting.\allowbreak{}NumThree}, lines 211--378)\par}

{\footnotesize\raggedright\noindent \ref{B6.16}.\quad \texttt{UpperTheory.\allowbreak{}V.\allowbreak{}ZK.\allowbreak{}lemmaK\_\allowbreak{}of\_\allowbreak{}steps} (\texttt{Development.\allowbreak{}AutomorphicArithmetic.\allowbreak{}AutomorphicVarieties.\allowbreak{}ModularityArgument.\allowbreak{}ClassCounting.\allowbreak{}LemmaKOfSteps}, lines 170--175)\\
\texttt{UpperTheory.\allowbreak{}V.\allowbreak{}ZK.\allowbreak{}lemmaK} (\texttt{Development.\allowbreak{}AutomorphicArithmetic.\allowbreak{}AutomorphicVarieties.\allowbreak{}ModularityArgument.\allowbreak{}ClassCounting.\allowbreak{}LemmaKOfLitN}, lines 28--31)\par}

{\footnotesize\raggedright\noindent \ref{B6.16a}.\quad \texttt{GlobalTheory.\allowbreak{}HomStab.\allowbreak{}kStability\_\allowbreak{}main} (\texttt{Development.\allowbreak{}AutomorphicArithmetic.\allowbreak{}AutomorphicVarieties.\allowbreak{}LiteratureResults.\allowbreak{}KTheoryStability.\allowbreak{}HomologicalStability.\allowbreak{}LinkExactnessStability.\allowbreak{}HomologicalStabilityTheorem}, lines 10--27)\\
\texttt{UpperTheory.\allowbreak{}V.\allowbreak{}KL.\allowbreak{}dedekindSR} (\texttt{Development.\allowbreak{}AutomorphicArithmetic.\allowbreak{}AutomorphicVarieties.\allowbreak{}TensorHomCubes.\allowbreak{}DedekindSR}, lines 50--111)\\
\texttt{UpperTheory.\allowbreak{}V.\allowbreak{}ZK.\allowbreak{}Up.\allowbreak{}torsStab} (\texttt{Development.\allowbreak{}AutomorphicArithmetic.\allowbreak{}AutomorphicVarieties.\allowbreak{}ModularityArgument.\allowbreak{}ClassCounting.\allowbreak{}TorsionRankCounts.\allowbreak{}TorsStab}, lines 16--39)\par}

{\footnotesize\raggedright\noindent \ref{B6.16b}.\quad \texttt{UpperTheory.\allowbreak{}V.\allowbreak{}KT.\allowbreak{}sx\_\allowbreak{}exact} (\texttt{Development.\allowbreak{}AutomorphicArithmetic.\allowbreak{}AutomorphicVarieties.\allowbreak{}HermiteNormalForms.\allowbreak{}Assembly}, lines 14--17)\\
\texttt{UpperTheory.\allowbreak{}V.\allowbreak{}KT.\allowbreak{}gramdet} (\texttt{Development.\allowbreak{}AutomorphicArithmetic.\allowbreak{}AutomorphicVarieties.\allowbreak{}HermiteNormalForms.\allowbreak{}Gramdet}, lines 49--100)\\
\texttt{UpperTheory.\allowbreak{}V.\allowbreak{}KT.\allowbreak{}orb} (\texttt{Development.\allowbreak{}AutomorphicArithmetic.\allowbreak{}AutomorphicVarieties.\allowbreak{}HermiteNormalForms.\allowbreak{}ColsMinSmul}, lines 34--75)\\
\texttt{UpperTheory.\allowbreak{}V.\allowbreak{}KT.\allowbreak{}orbnum} (\texttt{Development.\allowbreak{}AutomorphicArithmetic.\allowbreak{}AutomorphicVarieties.\allowbreak{}HermiteNormalForms.\allowbreak{}Orbnum}, lines 80--92)\\
\texttt{UpperTheory.\allowbreak{}V.\allowbreak{}KT.\allowbreak{}stab} (\texttt{Development.\allowbreak{}AutomorphicArithmetic.\allowbreak{}AutomorphicVarieties.\allowbreak{}HermiteNormalForms.\allowbreak{}ImportsAlt}, lines 8--26)\par}

{\footnotesize\raggedright\noindent \ref{B6.16c}.\quad \texttt{UpperTheory.\allowbreak{}V.\allowbreak{}KT.\allowbreak{}tors\_\allowbreak{}le} (\texttt{Development.\allowbreak{}AutomorphicArithmetic.\allowbreak{}AutomorphicVarieties.\allowbreak{}HermiteNormalForms.\allowbreak{}Assembly}, lines 25--30)\\
\texttt{UpperTheory.\allowbreak{}V.\allowbreak{}KH.\allowbreak{}shift} (\texttt{Development.\allowbreak{}AutomorphicArithmetic.\allowbreak{}AutomorphicVarieties.\allowbreak{}TorsionBoundsChains.\allowbreak{}LogLeExp}, lines 14--47)\\
\texttt{UpperTheory.\allowbreak{}V.\allowbreak{}ZK.\allowbreak{}torsAll} (\texttt{Development.\allowbreak{}AutomorphicArithmetic.\allowbreak{}AutomorphicVarieties.\allowbreak{}ModularityArgument.\allowbreak{}ClassCounting.\allowbreak{}ExponentLe}, lines 70--118)\par}

{\footnotesize\raggedright\noindent \ref{B6.17}.\quad \texttt{UpperTheory.\allowbreak{}V.\allowbreak{}Published.\allowbreak{}FerreroWashington} (\texttt{Development.\allowbreak{}AutomorphicArithmetic.\allowbreak{}AutomorphicVarieties.\allowbreak{}PublishedTheorems.\allowbreak{}FerreroWashington}, lines 7--12)\\
\texttt{UpperTheory.\allowbreak{}V.\allowbreak{}FW.\allowbreak{}ferreroWashington} (\texttt{Development.\allowbreak{}AutomorphicArithmetic.\allowbreak{}AutomorphicVarieties.\allowbreak{}GaloisLibrary.\allowbreak{}FerreroWashingtonIndex.\allowbreak{}FerreroWashington}, lines 104--119)\par}

{\footnotesize\raggedright\noindent \ref{B6.18}.\quad \texttt{UpperTheory.\allowbreak{}V.\allowbreak{}Two.\allowbreak{}towerFixed} (\texttt{Development.\allowbreak{}AutomorphicArithmetic.\allowbreak{}AutomorphicVarieties.\allowbreak{}QuadraticExtensions.\allowbreak{}IsAlgClosureRat}, lines 98--111)\\
\texttt{UpperTheory.\allowbreak{}V.\allowbreak{}Two.\allowbreak{}towerTwo} (\texttt{Development.\allowbreak{}AutomorphicArithmetic.\allowbreak{}AutomorphicVarieties.\allowbreak{}QuadraticExtensions.\allowbreak{}IsAlgClosureRat}, lines 87--97)\\
\texttt{UpperTheory.\allowbreak{}V.\allowbreak{}Two.\allowbreak{}baseRk} (\texttt{Development.\allowbreak{}AutomorphicArithmetic.\allowbreak{}AutomorphicVarieties.\allowbreak{}QuadraticExtensions.\allowbreak{}IsAlgClosureRat}, lines 69--86)\par}

{\footnotesize\raggedright\noindent \ref{B6.19}.\quad \texttt{UpperTheory.\allowbreak{}V.\allowbreak{}Two.\allowbreak{}shaTwo} (\texttt{Development.\allowbreak{}AutomorphicArithmetic.\allowbreak{}AutomorphicVarieties.\allowbreak{}QuadraticExtensions.\allowbreak{}ShaTwo}, lines 11--28)\\
\texttt{UpperTheory.\allowbreak{}V.\allowbreak{}Two.\allowbreak{}twoBase} (\texttt{Development.\allowbreak{}AutomorphicArithmetic.\allowbreak{}AutomorphicVarieties.\allowbreak{}QuadraticExtensions.\allowbreak{}IsAlgClosureRat}, lines 112--129)\par}

{\footnotesize\raggedright\noindent \ref{B6.20}.\quad \texttt{UpperTheory.\allowbreak{}V.\allowbreak{}PropA.\allowbreak{}descends} (\texttt{Development.\allowbreak{}AutomorphicArithmetic.\allowbreak{}AutomorphicVarieties.\allowbreak{}AdditiveValuationNorms.\allowbreak{}Descends}, lines 8--15)\\
\texttt{UpperTheory.\allowbreak{}V.\allowbreak{}PropA.\allowbreak{}Descends} (\texttt{Development.\allowbreak{}AutomorphicArithmetic.\allowbreak{}AutomorphicVarieties.\allowbreak{}AdditiveValuationNorms.\allowbreak{}ABoundOfDescends}, lines 16--19)\\
\texttt{UpperTheory.\allowbreak{}V.\allowbreak{}PropA.\allowbreak{}twLen} (\texttt{Development.\allowbreak{}AutomorphicArithmetic.\allowbreak{}AutomorphicVarieties.\allowbreak{}AdditiveValuationNorms.\allowbreak{}CokLenLeOfRangeLe}, lines 26--27)\par}

{\footnotesize\raggedright\noindent \ref{B6.20}.\quad \texttt{UpperTheory.\allowbreak{}V.\allowbreak{}Kum.\allowbreak{}resInfOnto} (\texttt{Development.\allowbreak{}AutomorphicArithmetic.\allowbreak{}AutomorphicVarieties.\allowbreak{}GaloisLibrary.\allowbreak{}KummerTheory.\allowbreak{}ResInfOnto}, lines 102--166)\par}

{\footnotesize\raggedright\noindent \ref{B6.21}.\quad \texttt{UpperTheory.\allowbreak{}V.\allowbreak{}PropA.\allowbreak{}exists\_\allowbreak{}centres\_\allowbreak{}level} (\texttt{Development.\allowbreak{}AutomorphicArithmetic.\allowbreak{}AutomorphicVarieties.\allowbreak{}AdditiveValuationNorms.\allowbreak{}ExistsCentresLevel}, lines 23--33)\\
\texttt{UpperTheory.\allowbreak{}V.\allowbreak{}PropA.\allowbreak{}exists\_\allowbreak{}centres} (\texttt{Development.\allowbreak{}AutomorphicArithmetic.\allowbreak{}AutomorphicVarieties.\allowbreak{}AdditiveValuationNorms.\allowbreak{}ExistsCentresLevel}, lines 34--50)\\
\texttt{UpperTheory.\allowbreak{}V.\allowbreak{}PropA.\allowbreak{}exists\_\allowbreak{}centres\_\allowbreak{}norm} (\texttt{Development.\allowbreak{}AutomorphicArithmetic.\allowbreak{}AutomorphicVarieties.\allowbreak{}AdditiveValuationNorms.\allowbreak{}ExistsCentresLevel}, lines 57--70)\par}

{\footnotesize\raggedright\noindent \ref{B6.22}.\quad \texttt{UpperTheory.\allowbreak{}V.\allowbreak{}PropA.\allowbreak{}aBound\_\allowbreak{}of\_\allowbreak{}isTower} (\texttt{Development.\allowbreak{}AutomorphicArithmetic.\allowbreak{}AutomorphicVarieties.\allowbreak{}AdditiveValuationNorms.\allowbreak{}ABoundOfIsTower}, lines 7--11)\\
\texttt{UpperTheory.\allowbreak{}V.\allowbreak{}PropA.\allowbreak{}aBound\_\allowbreak{}of\_\allowbreak{}descends} (\texttt{Development.\allowbreak{}AutomorphicArithmetic.\allowbreak{}AutomorphicVarieties.\allowbreak{}AdditiveValuationNorms.\allowbreak{}ABoundOfDescends}, lines 36--56)\\
\texttt{UpperTheory.\allowbreak{}V.\allowbreak{}Win.\allowbreak{}SecondBound} (\texttt{Development.\allowbreak{}AutomorphicArithmetic.\allowbreak{}AutomorphicVarieties.\allowbreak{}LayeredIsomorphisms.\allowbreak{}SecondBound}, lines 13--18)\par}

{\footnotesize\raggedright\noindent \ref{B6.22}.\quad \texttt{UpperTheory.\allowbreak{}V.\allowbreak{}PropA.\allowbreak{}shaLeV} (\texttt{Development.\allowbreak{}AutomorphicArithmetic.\allowbreak{}AutomorphicVarieties.\allowbreak{}AdditiveValuationNorms.\allowbreak{}ShaLeV}, lines 7--16)\par}

{\footnotesize\raggedright\noindent \ref{B6.23}.\quad \texttt{UpperTheory.\allowbreak{}V.\allowbreak{}Win.\allowbreak{}exc} (\texttt{Development.\allowbreak{}AutomorphicArithmetic.\allowbreak{}AutomorphicVarieties.\allowbreak{}LayeredIsomorphisms.\allowbreak{}Exc}, lines 7--42)\\
\texttt{UpperTheory.\allowbreak{}V.\allowbreak{}Win.\allowbreak{}FirstBound} (\texttt{Development.\allowbreak{}AutomorphicArithmetic.\allowbreak{}AutomorphicVarieties.\allowbreak{}LayeredIsomorphisms.\allowbreak{}SecondBound}, lines 27--31)\par}

{\footnotesize\raggedright\noindent \ref{B6.24}.\quad \texttt{UpperTheory.\allowbreak{}V.\allowbreak{}Win.\allowbreak{}window1} (\texttt{Development.\allowbreak{}AutomorphicArithmetic.\allowbreak{}AutomorphicVarieties.\allowbreak{}LayeredIsomorphisms.\allowbreak{}Window1}, lines 6--58)\\
\texttt{UpperTheory.\allowbreak{}V.\allowbreak{}Win.\allowbreak{}window} (\texttt{Development.\allowbreak{}AutomorphicArithmetic.\allowbreak{}AutomorphicVarieties.\allowbreak{}LayeredIsomorphisms.\allowbreak{}Window}, lines 7--43)\par}

{\footnotesize\raggedright\noindent \ref{B6.25}.\quad \texttt{UpperTheory.\allowbreak{}V.\allowbreak{}Win.\allowbreak{}point2} (\texttt{Development.\allowbreak{}AutomorphicArithmetic.\allowbreak{}AutomorphicVarieties.\allowbreak{}LayeredIsomorphisms.\allowbreak{}Point2}, lines 8--49)\\
\texttt{UpperTheory.\allowbreak{}V.\allowbreak{}Win.\allowbreak{}final} (\texttt{Development.\allowbreak{}AutomorphicArithmetic.\allowbreak{}AutomorphicVarieties.\allowbreak{}LayeredIsomorphisms.\allowbreak{}Final}, lines 6--68)\par}

{\footnotesize\raggedright\noindent \ref{B6.26}.\quad \texttt{UpperTheory.\allowbreak{}V.\allowbreak{}Win.\allowbreak{}partA} (\texttt{Development.\allowbreak{}AutomorphicArithmetic.\allowbreak{}AutomorphicVarieties.\allowbreak{}LayeredIsomorphisms.\allowbreak{}PartA}, lines 24--78)\\
\texttt{UpperTheory.\allowbreak{}V.\allowbreak{}Win.\allowbreak{}partB} (\texttt{Development.\allowbreak{}AutomorphicArithmetic.\allowbreak{}AutomorphicVarieties.\allowbreak{}LayeredIsomorphisms.\allowbreak{}PartB}, lines 50--95)\par}

{\footnotesize\raggedright\noindent \ref{B6.26}.\quad \texttt{UpperTheory.\allowbreak{}V.\allowbreak{}Win.\allowbreak{}layeremb} (\texttt{Development.\allowbreak{}AutomorphicArithmetic.\allowbreak{}AutomorphicVarieties.\allowbreak{}LayeredIsomorphisms.\allowbreak{}Layeremb}, lines 7--34)\\
\texttt{UpperTheory.\allowbreak{}V.\allowbreak{}Win.\allowbreak{}disc} (\texttt{Development.\allowbreak{}AutomorphicArithmetic.\allowbreak{}AutomorphicVarieties.\allowbreak{}LayeredIsomorphisms.\allowbreak{}Imports}, lines 6--44)\par}

{\footnotesize\raggedright\noindent \ref{B6.27}.\quad \texttt{UpperTheory.\allowbreak{}V.\allowbreak{}Win.\allowbreak{}thmBab\_\allowbreak{}of\_\allowbreak{}lemmaK} (\texttt{Development.\allowbreak{}AutomorphicArithmetic.\allowbreak{}AutomorphicVarieties.\allowbreak{}LayeredIsomorphisms.\allowbreak{}ThmBabOfLemmaK}, lines 9--13)\\
\texttt{UpperTheory.\allowbreak{}V.\allowbreak{}Win.\allowbreak{}thmBab\_\allowbreak{}of} (\texttt{Development.\allowbreak{}AutomorphicArithmetic.\allowbreak{}AutomorphicVarieties.\allowbreak{}LayeredIsomorphisms.\allowbreak{}ThmBabOf}, lines 11--103)\par}

{\footnotesize\raggedright\noindent \ref{B6.27}.\quad \texttt{UpperTheory.\allowbreak{}V.\allowbreak{}HV.\allowbreak{}DescentVanishingReductions.\allowbreak{}vwin\_\allowbreak{}of\_\allowbreak{}hypotheses\_\allowbreak{}H\_\allowbreak{}E\_\allowbreak{}alt} (\texttt{Development.\allowbreak{}AutomorphicArithmetic.\allowbreak{}AutomorphicVarieties.\allowbreak{}ModularityArgument.\allowbreak{}ThresholdArguments.\allowbreak{}DenseThresholdCases.\allowbreak{}WindowBoundOfLiteratureLevel}, lines 25--29)\\
\texttt{UpperTheory.\allowbreak{}V.\allowbreak{}HV.\allowbreak{}Core.\allowbreak{}of\_\allowbreak{}four\_\allowbreak{}alt} (\texttt{Development.\allowbreak{}AutomorphicArithmetic.\allowbreak{}AutomorphicVarieties.\allowbreak{}ModularityArgument.\allowbreak{}ThresholdArguments.\allowbreak{}DenseThresholdCases.\allowbreak{}HVCore}, lines 18--22)\par}

\section*{Chapter~\ref{ch:b6-classwise}: Summing the plus component by classes}

{\footnotesize\raggedright\noindent \ref{B6.28}.\quad \texttt{UpperTheory.\allowbreak{}V.\allowbreak{}HV.\allowbreak{}Host.\allowbreak{}X} (\texttt{Development.\allowbreak{}AutomorphicArithmetic.\allowbreak{}AutomorphicVarieties.\allowbreak{}HeightAndVanishing.\allowbreak{}DirichletPrimeToLPartMulLPowerPart}, lines 188--189)\\
\texttt{UpperTheory.\allowbreak{}V.\allowbreak{}HV.\allowbreak{}chiP} (\texttt{Development.\allowbreak{}AutomorphicArithmetic.\allowbreak{}AutomorphicVarieties.\allowbreak{}HeightAndVanishing.\allowbreak{}DirichletPrimeToLPartMulLPowerPart}, lines 86)\\
\texttt{UpperTheory.\allowbreak{}V.\allowbreak{}Bchi} (\texttt{Development.\allowbreak{}AutomorphicArithmetic.\allowbreak{}AutomorphicVarieties.\allowbreak{}StmtFields}, lines 66--67)\\
\texttt{UpperTheory.\allowbreak{}V.\allowbreak{}HV.\allowbreak{}cases} (\texttt{Development.\allowbreak{}AutomorphicArithmetic.\allowbreak{}AutomorphicVarieties.\allowbreak{}HeightAndVanishing.\allowbreak{}Cases}, lines 8--48)\\
\texttt{UpperTheory.\allowbreak{}V.\allowbreak{}HV.\allowbreak{}CountV} (\texttt{Development.\allowbreak{}AutomorphicArithmetic.\allowbreak{}AutomorphicVarieties.\allowbreak{}HeightAndVanishing.\allowbreak{}Counts}, lines 69)\\
\texttt{UpperTheory.\allowbreak{}V.\allowbreak{}HV.\allowbreak{}CountVs} (\texttt{Development.\allowbreak{}AutomorphicArithmetic.\allowbreak{}AutomorphicVarieties.\allowbreak{}HeightAndVanishing.\allowbreak{}Counts}, lines 71)\par}

{\footnotesize\raggedright\noindent \ref{B6.28a}.\quad \texttt{UpperTheory.\allowbreak{}V.\allowbreak{}HV.\allowbreak{}hostIndex} (\texttt{Development.\allowbreak{}AutomorphicArithmetic.\allowbreak{}AutomorphicVarieties.\allowbreak{}HeightAndVanishing.\allowbreak{}HostIndex}, lines 67--165)\\
\texttt{UpperTheory.\allowbreak{}V.\allowbreak{}HV.\allowbreak{}hostTh1} (\texttt{Development.\allowbreak{}AutomorphicArithmetic.\allowbreak{}AutomorphicVarieties.\allowbreak{}HeightAndVanishing.\allowbreak{}OrdProjNDvd}, lines 37--118)\\
\texttt{UpperTheory.\allowbreak{}V.\allowbreak{}HV.\allowbreak{}hostComp} (\texttt{Development.\allowbreak{}AutomorphicArithmetic.\allowbreak{}AutomorphicVarieties.\allowbreak{}HeightAndVanishing.\allowbreak{}HostComp}, lines 21--35)\\
\texttt{UpperTheory.\allowbreak{}V.\allowbreak{}HV.\allowbreak{}hostOdd} (\texttt{Development.\allowbreak{}AutomorphicArithmetic.\allowbreak{}AutomorphicVarieties.\allowbreak{}HeightAndVanishing.\allowbreak{}HostOdd}, lines 70--114)\par}

{\footnotesize\raggedright\noindent \ref{B6.29}.\quad \texttt{UpperTheory.\allowbreak{}V.\allowbreak{}Chr.\allowbreak{}rho} (\texttt{Development.\allowbreak{}AutomorphicArithmetic.\allowbreak{}AutomorphicVarieties.\allowbreak{}CharacterLevelBounds.\allowbreak{}ShaFamMono}, lines 17--18)\\
\texttt{UpperTheory.\allowbreak{}V.\allowbreak{}Chr.\allowbreak{}VS} (\texttt{Development.\allowbreak{}AutomorphicArithmetic.\allowbreak{}AutomorphicVarieties.\allowbreak{}CharacterLevelBounds.\allowbreak{}ShaFamMono}, lines 30)\\
\texttt{UpperTheory.\allowbreak{}V.\allowbreak{}Chr.\allowbreak{}V'} (\texttt{Development.\allowbreak{}AutomorphicArithmetic.\allowbreak{}AutomorphicVarieties.\allowbreak{}CharacterLevelBounds.\allowbreak{}ShaFamMono}, lines 34)\\
\texttt{UpperTheory.\allowbreak{}V.\allowbreak{}Chr.\allowbreak{}nrm} (\texttt{Development.\allowbreak{}AutomorphicArithmetic.\allowbreak{}AutomorphicVarieties.\allowbreak{}CharacterLevelBounds.\allowbreak{}ShaFamMono}, lines 71)\\
\texttt{UpperTheory.\allowbreak{}V.\allowbreak{}HV.\allowbreak{}VLe} (\texttt{Development.\allowbreak{}AutomorphicArithmetic.\allowbreak{}AutomorphicVarieties.\allowbreak{}HeightAndVanishing.\allowbreak{}DirichletPrimeToLPartMulLPowerPart}, lines 120)\par}

{\footnotesize\raggedright\noindent \ref{B6.29a}.\quad \texttt{UpperTheory.\allowbreak{}V.\allowbreak{}HV.\allowbreak{}coefIndep} (\texttt{Development.\allowbreak{}AutomorphicArithmetic.\allowbreak{}AutomorphicVarieties.\allowbreak{}HeightAndVanishing.\allowbreak{}CoefIndep}, lines 68--77)\par}

{\footnotesize\raggedright\noindent \ref{def:b6-cuspidal-witness}.\quad \texttt{UpperTheory.\allowbreak{}V.\allowbreak{}Af.\allowbreak{}AttC} (\texttt{Development.\allowbreak{}AutomorphicArithmetic.\allowbreak{}AutomorphicVarieties.\allowbreak{}AutomorphicForms.\allowbreak{}BlochKatoSelmerVanishes}, lines 69--82)\\
\texttt{UpperTheory.\allowbreak{}V.\allowbreak{}Z.\allowbreak{}Win.\allowbreak{}WitA} (\texttt{Development.\allowbreak{}AutomorphicArithmetic.\allowbreak{}AutomorphicVarieties.\allowbreak{}ModularityArgument.\allowbreak{}AutomorphyLifting.\allowbreak{}WeakOfAtt}, lines 19--30)\\
\texttt{UpperTheory.\allowbreak{}V.\allowbreak{}Z.\allowbreak{}ImaginaryQuadraticEtaCharacter.\allowbreak{}WinAt} (\texttt{Development.\allowbreak{}AutomorphicArithmetic.\allowbreak{}AutomorphicVarieties.\allowbreak{}ModularityArgument.\allowbreak{}LeviSmoothLocalization.\allowbreak{}WinUpDown}, lines 38--41)\\
\texttt{UpperTheory.\allowbreak{}V.\allowbreak{}Z.\allowbreak{}Win.\allowbreak{}IsAutQA} (\texttt{Development.\allowbreak{}AutomorphicArithmetic.\allowbreak{}AutomorphicVarieties.\allowbreak{}ModularityArgument.\allowbreak{}AutomorphyLifting.\allowbreak{}IsAutCsdA}, lines 16--20)\par}

{\footnotesize\raggedright\noindent \ref{B6.30aa}.\quad \texttt{UpperTheory.\allowbreak{}V.\allowbreak{}Pr.\allowbreak{}irrNormal} (\texttt{Development.\allowbreak{}AutomorphicArithmetic.\allowbreak{}AutomorphicVarieties.\allowbreak{}CurveDVRSpecialization.\allowbreak{}IrrNormal}, lines 62--180)\\
\texttt{UpperTheory.\allowbreak{}V.\allowbreak{}Th.\allowbreak{}Det.\allowbreak{}irrSmall} (\texttt{Development.\allowbreak{}AutomorphicArithmetic.\allowbreak{}AutomorphicVarieties.\allowbreak{}TheoremsMuLinked.\allowbreak{}CharacteristicPolynomials.\allowbreak{}IrrSmall}, lines 33--65)\par}

{\footnotesize\raggedright\noindent \ref{B6.30ac}.\quad \texttt{UpperTheory.\allowbreak{}V.\allowbreak{}Pr.\allowbreak{}matIrr\_\allowbreak{}open\_\allowbreak{}gd} (\texttt{Development.\allowbreak{}AutomorphicArithmetic.\allowbreak{}AutomorphicVarieties.\allowbreak{}CurveDVRSpecialization.\allowbreak{}MatIrrOpenGd}, lines 342--383)\\
\texttt{UpperTheory.\allowbreak{}V.\allowbreak{}U3.\allowbreak{}typeII\_\allowbreak{}irr\_\allowbreak{}cyc} (\texttt{Development.\allowbreak{}AutomorphicArithmetic.\allowbreak{}AutomorphicVarieties.\allowbreak{}UnitaryThreeGroups.\allowbreak{}ExistsUnitPowNe}, lines 72--143)\par}

{\footnotesize\raggedright\noindent \ref{B6.30ab}.\quad \texttt{UpperTheory.\allowbreak{}V.\allowbreak{}Z.\allowbreak{}J.\allowbreak{}rigIUE\_\allowbreak{}alt} (\texttt{Development.\allowbreak{}AutomorphicArithmetic.\allowbreak{}AutomorphicVarieties.\allowbreak{}ModularityArgument.\allowbreak{}ThresholdArguments.\allowbreak{}DenseThresholdCases.\allowbreak{}ERigIU}, lines 80--146)\\
\texttt{UpperTheory.\allowbreak{}V.\allowbreak{}Z.\allowbreak{}Win.\allowbreak{}RR.\allowbreak{}rigII\_\allowbreak{}a\_\allowbreak{}alt} (\texttt{Development.\allowbreak{}AutomorphicArithmetic.\allowbreak{}AutomorphicVarieties.\allowbreak{}ModularityArgument.\allowbreak{}ThresholdArguments.\allowbreak{}DenseThresholdCases.\allowbreak{}RRRigII}, lines 25--102)\par}

{\footnotesize\raggedright\noindent \ref{B6.30ad}.\quad \texttt{UpperTheory.\allowbreak{}V.\allowbreak{}Z.\allowbreak{}Desc.\allowbreak{}V0.\allowbreak{}numJE\_\allowbreak{}alt} (\texttt{Development.\allowbreak{}AutomorphicArithmetic.\allowbreak{}AutomorphicVarieties.\allowbreak{}ModularityArgument.\allowbreak{}ThresholdArguments.\allowbreak{}DenseThresholdCases.\allowbreak{}TNumJE}, lines 171--198)\\
\texttt{UpperTheory.\allowbreak{}V.\allowbreak{}Z.\allowbreak{}Desc.\allowbreak{}V0.\allowbreak{}z\_\allowbreak{}eq\_\allowbreak{}dim\_\allowbreak{}e\_\allowbreak{}of\_\allowbreak{}count\_\allowbreak{}e\_\allowbreak{}res\_\allowbreak{}eq\_\allowbreak{}inv\_\allowbreak{}loc\_\allowbreak{}kir} (\texttt{Development.\allowbreak{}AutomorphicArithmetic.\allowbreak{}AutomorphicVarieties.\allowbreak{}ModularityArgument.\allowbreak{}ThresholdArguments.\allowbreak{}DenseThresholdCases.\allowbreak{}TZEqDimE}, lines 57--105)\par}

{\footnotesize\raggedright\noindent \ref{B6.30ae}.\quad \texttt{UpperTheory.\allowbreak{}V.\allowbreak{}Z.\allowbreak{}Desc.\allowbreak{}V0.\allowbreak{}unramified\_\allowbreak{}at\_\allowbreak{}of\_\allowbreak{}dense\_\allowbreak{}purity\_\allowbreak{}num\_\allowbreak{}e\_\allowbreak{}power\_\allowbreak{}series} (\texttt{Development.\allowbreak{}AutomorphicArithmetic.\allowbreak{}AutomorphicVarieties.\allowbreak{}ModularityArgument.\allowbreak{}ThresholdArguments.\allowbreak{}DenseThresholdCases.\allowbreak{}UnrE}, lines 21--54)\\
\texttt{UpperTheory.\allowbreak{}V.\allowbreak{}Z.\allowbreak{}Desc.\allowbreak{}V0.\allowbreak{}unrAlgW} (\texttt{Development.\allowbreak{}AutomorphicArithmetic.\allowbreak{}AutomorphicVarieties.\allowbreak{}ModularityArgument.\allowbreak{}CuspidalDescent.\allowbreak{}GLDenseRanges.\allowbreak{}InertiaPartMaps.\allowbreak{}UnrAlgW}, lines 141--204)\par}

{\footnotesize\raggedright\noindent \ref{B6.30af}.\quad \texttt{UpperTheory.\allowbreak{}Gal.\allowbreak{}projectedFormsSystemPart.\allowbreak{}red2\_\allowbreak{}alt} (\texttt{Development.\allowbreak{}AutomorphicArithmetic.\allowbreak{}AutomorphicVarieties.\allowbreak{}ModularityArgument.\allowbreak{}ThresholdArguments.\allowbreak{}DenseThresholdCases.\allowbreak{}EndPtOf}, lines 160--172)\\
\texttt{UpperTheory.\allowbreak{}V.\allowbreak{}J.\allowbreak{}redAbs} (\texttt{Development.\allowbreak{}AutomorphicArithmetic.\allowbreak{}AutomorphicVarieties.\allowbreak{}ModularityArgument.\allowbreak{}LinkedUnramifiedChains.\allowbreak{}ScalarIrreducibility.\allowbreak{}RedAbs}, lines 15--48)\par}

{\footnotesize\raggedright\noindent \ref{B6.30ag}.\quad \texttt{UpperTheory.\allowbreak{}V.\allowbreak{}Z.\allowbreak{}J.\allowbreak{}unitOf} (\texttt{Development.\allowbreak{}AutomorphicArithmetic.\allowbreak{}AutomorphicVarieties.\allowbreak{}ModularityArgument.\allowbreak{}LinkedUnramifiedChains.\allowbreak{}CrystallineCardinalities.\allowbreak{}UnitOf}, lines 245--305)\\
\texttt{UpperTheory.\allowbreak{}V.\allowbreak{}Z.\allowbreak{}J.\allowbreak{}locHigh} (\texttt{Development.\allowbreak{}AutomorphicArithmetic.\allowbreak{}AutomorphicVarieties.\allowbreak{}ModularityArgument.\allowbreak{}LinkedUnramifiedChains.\allowbreak{}CrystallineCardinalities.\allowbreak{}LocHigh}, lines 18--87)\\
\texttt{UpperTheory.\allowbreak{}V.\allowbreak{}Z.\allowbreak{}J.\allowbreak{}selLow} (\texttt{Development.\allowbreak{}AutomorphicArithmetic.\allowbreak{}AutomorphicVarieties.\allowbreak{}ModularityArgument.\allowbreak{}LinkedUnramifiedChains.\allowbreak{}CrystallineCardinalities.\allowbreak{}SelLow}, lines 60--98)\\
\texttt{UpperTheory.\allowbreak{}V.\allowbreak{}Z.\allowbreak{}J.\allowbreak{}tan} (\texttt{Development.\allowbreak{}AutomorphicArithmetic.\allowbreak{}AutomorphicVarieties.\allowbreak{}ModularityArgument.\allowbreak{}LinkedUnramifiedChains.\allowbreak{}CrystallineCardinalities.\allowbreak{}BoxUnit}, lines 128--160)\par}

{\footnotesize\raggedright\noindent \ref{B6.30ah}.\quad \texttt{UpperTheory.\allowbreak{}Gal.\allowbreak{}projectedFormsSystemPart.\allowbreak{}UnitaryRankTwoFlagsAlt} (\texttt{Development.\allowbreak{}AutomorphicArithmetic.\allowbreak{}AutomorphicVarieties.\allowbreak{}ModularityArgument.\allowbreak{}ThresholdArguments.\allowbreak{}DenseThresholdCases.\allowbreak{}UnitaryRankTwoFlags}, lines 36--60)\\
\texttt{UpperTheory.\allowbreak{}V.\allowbreak{}Z.\allowbreak{}J.\allowbreak{}unlEnd3E} (\texttt{Development.\allowbreak{}AutomorphicArithmetic.\allowbreak{}AutomorphicVarieties.\allowbreak{}ModularityArgument.\allowbreak{}LinkedUnramifiedChains.\allowbreak{}LogarithmicNormBounds.\allowbreak{}BlockValuationRep}, lines 234--346)\\
\texttt{UpperTheory.\allowbreak{}V.\allowbreak{}J.\allowbreak{}endJ} (\texttt{Development.\allowbreak{}AutomorphicArithmetic.\allowbreak{}AutomorphicVarieties.\allowbreak{}ModularityArgument.\allowbreak{}LinkedUnramifiedChains.\allowbreak{}DerivativeDivisibility.\allowbreak{}CoefInj}, lines 12--73)\\
\texttt{UpperTheory.\allowbreak{}V.\allowbreak{}J.\allowbreak{}tor} (\texttt{Development.\allowbreak{}AutomorphicArithmetic.\allowbreak{}AutomorphicVarieties.\allowbreak{}ModularityArgument.\allowbreak{}LinkedUnramifiedChains.\allowbreak{}DerivativeDivisibility.\allowbreak{}PowNeZero}, lines 124--143)\par}

{\footnotesize\raggedright\noindent \ref{B6.30a}.\quad \texttt{UpperTheory.\allowbreak{}V.\allowbreak{}Z.\allowbreak{}J.\allowbreak{}admissible\_\allowbreak{}typeI\_\allowbreak{}even\_\allowbreak{}level\_\allowbreak{}alt} (\texttt{Development.\allowbreak{}AutomorphicArithmetic.\allowbreak{}AutomorphicVarieties.\allowbreak{}ModularityArgument.\allowbreak{}ThresholdArguments.\allowbreak{}DenseThresholdCases.\allowbreak{}AdmissibleTypeIEvenLevel}, lines 76--151)\\
\texttt{UpperTheory.\allowbreak{}V.\allowbreak{}Z.\allowbreak{}J.\allowbreak{}rigIUE\_\allowbreak{}alt} (\texttt{Development.\allowbreak{}AutomorphicArithmetic.\allowbreak{}AutomorphicVarieties.\allowbreak{}ModularityArgument.\allowbreak{}ThresholdArguments.\allowbreak{}DenseThresholdCases.\allowbreak{}ERigIU}, lines 80--155)\\
\texttt{UpperTheory.\allowbreak{}V.\allowbreak{}Z.\allowbreak{}Win.\allowbreak{}RR.\allowbreak{}rigII\_\allowbreak{}a\_\allowbreak{}alt} (\texttt{Development.\allowbreak{}AutomorphicArithmetic.\allowbreak{}AutomorphicVarieties.\allowbreak{}ModularityArgument.\allowbreak{}ThresholdArguments.\allowbreak{}DenseThresholdCases.\allowbreak{}RRRigII}, lines 25--111)\par}

{\footnotesize\raggedright\noindent \ref{B6.30}.\quad \texttt{UpperTheory.\allowbreak{}V.\allowbreak{}U3.\allowbreak{}countI} (\texttt{Development.\allowbreak{}AutomorphicArithmetic.\allowbreak{}AutomorphicVarieties.\allowbreak{}UnitaryThreeGroups.\allowbreak{}CountI}, lines 22--136)\\
\texttt{UpperTheory.\allowbreak{}V.\allowbreak{}Z.\allowbreak{}J.\allowbreak{}count\_\allowbreak{}i\_\allowbreak{}of\_\allowbreak{}printed\_\allowbreak{}bound\_\allowbreak{}k\_\allowbreak{}kato\_\allowbreak{}automorphy} (\texttt{Development.\allowbreak{}AutomorphicArithmetic.\allowbreak{}AutomorphicVarieties.\allowbreak{}ModularityArgument.\allowbreak{}ThresholdArguments.\allowbreak{}DenseThresholdCases.\allowbreak{}LinkedUnlE}, lines 49--54)\par}

{\footnotesize\raggedright\noindent \ref{B6.30b}.\quad \texttt{UpperTheory.\allowbreak{}V.\allowbreak{}U3.\allowbreak{}countI} (\texttt{Development.\allowbreak{}AutomorphicArithmetic.\allowbreak{}AutomorphicVarieties.\allowbreak{}UnitaryThreeGroups.\allowbreak{}CountI}, lines 22--136)\\
\texttt{UpperTheory.\allowbreak{}V.\allowbreak{}HV.\allowbreak{}countZ} (\texttt{Development.\allowbreak{}AutomorphicArithmetic.\allowbreak{}AutomorphicVarieties.\allowbreak{}HeightAndVanishing.\allowbreak{}CountZ}, lines 5--32)\par}

{\footnotesize\raggedright\noindent \ref{B6.31b}.\quad \texttt{UpperTheory.\allowbreak{}V.\allowbreak{}GradeZeroInterpolatingElement.\allowbreak{}Raw.\allowbreak{}Good} (\texttt{Development.\allowbreak{}AutomorphicArithmetic.\allowbreak{}AutomorphicVarieties.\allowbreak{}SymplecticGSp4.\allowbreak{}DecompSigmaOrd}, lines 94--120)\\
\texttt{UpperTheory.\allowbreak{}V.\allowbreak{}Aut.\allowbreak{}OfChar} (\texttt{Development.\allowbreak{}AutomorphicArithmetic.\allowbreak{}AutomorphicVarieties.\allowbreak{}AutomorphicCases.\allowbreak{}SymplecticGSp4}, lines 12--17)\\
\texttt{UpperTheory.\allowbreak{}V.\allowbreak{}Aut.\allowbreak{}OfCase} (\texttt{Development.\allowbreak{}AutomorphicArithmetic.\allowbreak{}AutomorphicVarieties.\allowbreak{}AutomorphicCases.\allowbreak{}SymplecticGSp4}, lines 19--20)\\
\texttt{UpperTheory.\allowbreak{}V.\allowbreak{}Ws.\allowbreak{}Ctx2} (\texttt{Development.\allowbreak{}AutomorphicArithmetic.\allowbreak{}AutomorphicVarieties.\allowbreak{}WeakShimuraTheory.\allowbreak{}NearSplit}, lines 55--57)\\
\texttt{UpperTheory.\allowbreak{}V.\allowbreak{}Ws.\allowbreak{}Ctx3} (\texttt{Development.\allowbreak{}AutomorphicArithmetic.\allowbreak{}AutomorphicVarieties.\allowbreak{}WeakShimuraTheory.\allowbreak{}NearSplit}, lines 59--61)\par}

{\footnotesize\raggedright\noindent \ref{B6.31c}.\quad \texttt{UpperTheory.\allowbreak{}V.\allowbreak{}GradeZeroInterpolatingElement.\allowbreak{}OnComp} (\texttt{Development.\allowbreak{}AutomorphicArithmetic.\allowbreak{}AutomorphicVarieties.\allowbreak{}SymplecticGSp4.\allowbreak{}DecompSigmaOrd}, lines 27--31)\\
\texttt{UpperTheory.\allowbreak{}V.\allowbreak{}GradeZeroInterpolatingElement.\allowbreak{}Ord} (\texttt{Development.\allowbreak{}AutomorphicArithmetic.\allowbreak{}AutomorphicVarieties.\allowbreak{}SymplecticGSp4.\allowbreak{}MapRepApply}, lines 99--102)\\
\texttt{UpperTheory.\allowbreak{}V.\allowbreak{}GradeZeroInterpolatingElement.\allowbreak{}Raw.\allowbreak{}IsPt} (\texttt{Development.\allowbreak{}AutomorphicArithmetic.\allowbreak{}AutomorphicVarieties.\allowbreak{}SymplecticGSp4.\allowbreak{}DecompSigmaOrd}, lines 155--160)\\
\texttt{UpperTheory.\allowbreak{}V.\allowbreak{}GradeZeroInterpolatingElement.\allowbreak{}Raw.\allowbreak{}PtIn} (\texttt{Development.\allowbreak{}AutomorphicArithmetic.\allowbreak{}AutomorphicVarieties.\allowbreak{}SymplecticGSp4.\allowbreak{}DecompSigmaOrd}, lines 164--167)\\
\texttt{UpperTheory.\allowbreak{}V.\allowbreak{}GradeZeroInterpolatingElement.\allowbreak{}Raw.\allowbreak{}Fin0} (\texttt{Development.\allowbreak{}AutomorphicArithmetic.\allowbreak{}AutomorphicVarieties.\allowbreak{}SymplecticGSp4.\allowbreak{}DecompSigmaOrd}, lines 172--174)\\
\texttt{UpperTheory.\allowbreak{}V.\allowbreak{}GradeZeroInterpolatingElement.\allowbreak{}Raw.\allowbreak{}RigidAt} (\texttt{Development.\allowbreak{}AutomorphicArithmetic.\allowbreak{}AutomorphicVarieties.\allowbreak{}SymplecticGSp4.\allowbreak{}DecompSigmaOrd}, lines 182--184)\par}

{\footnotesize\raggedright\noindent \ref{B6.31a}.\quad \texttt{UpperTheory.\allowbreak{}V.\allowbreak{}X.\allowbreak{}soBar} (\texttt{Development.\allowbreak{}AutomorphicArithmetic.\allowbreak{}AutomorphicVarieties.\allowbreak{}LevelCohomology.\allowbreak{}FdAdjoinRange}, lines 27--96)\\
\texttt{UpperTheory.\allowbreak{}V.\allowbreak{}X.\allowbreak{}soVec} (\texttt{Development.\allowbreak{}AutomorphicArithmetic.\allowbreak{}AutomorphicVarieties.\allowbreak{}LevelCohomology.\allowbreak{}IntIffNorm}, lines 143--148)\par}

{\footnotesize\raggedright\noindent \ref{B6.31d}.\quad \texttt{UpperTheory.\allowbreak{}V.\allowbreak{}GradeZeroInterpolatingElement.\allowbreak{}dataEx} (\texttt{Development.\allowbreak{}AutomorphicArithmetic.\allowbreak{}AutomorphicVarieties.\allowbreak{}SymplecticGSp4.\allowbreak{}DataEx}, lines 366--450)\par}

{\footnotesize\raggedright\noindent \ref{B6.31e}.\quad \texttt{UpperTheory.\allowbreak{}V.\allowbreak{}GradeZeroInterpolatingElement.\allowbreak{}aSmooth} (\texttt{Development.\allowbreak{}AutomorphicArithmetic.\allowbreak{}AutomorphicVarieties.\allowbreak{}SymplecticGSp4.\allowbreak{}ASmooth}, lines 1049--1135)\\
\texttt{UpperTheory.\allowbreak{}V.\allowbreak{}GradeZeroInterpolatingElement.\allowbreak{}exModel} (\texttt{Development.\allowbreak{}AutomorphicArithmetic.\allowbreak{}AutomorphicVarieties.\allowbreak{}SymplecticGSp4.\allowbreak{}ExModel}, lines 263--268,304--332,425--453)\par}

{\footnotesize\raggedright\noindent \ref{B6.31}.\quad \texttt{UpperTheory.\allowbreak{}V.\allowbreak{}Th.\allowbreak{}sp4Main\_\allowbreak{}alt} (\texttt{Development.\allowbreak{}AutomorphicArithmetic.\allowbreak{}AutomorphicVarieties.\allowbreak{}ModularityArgument.\allowbreak{}ThresholdArguments.\allowbreak{}DenseThresholdCases.\allowbreak{}ConductorAtMapInjectiveCoeff}, lines 131--171)\\
\texttt{UpperTheory.\allowbreak{}V.\allowbreak{}HV.\allowbreak{}DescentVanishingReductions.\allowbreak{}countII\_\allowbreak{}and\_\allowbreak{}countIII\_\allowbreak{}of\_\allowbreak{}printsH\_\allowbreak{}of\_\allowbreak{}sp4In} (\texttt{Development.\allowbreak{}AutomorphicArithmetic.\allowbreak{}AutomorphicVarieties.\allowbreak{}ModularityArgument.\allowbreak{}ThresholdArguments.\allowbreak{}DenseThresholdCases.\allowbreak{}AdjointLiftingOfCohomology}, lines 33--39)\par}

{\footnotesize\raggedright\noindent \ref{B6.32a}.\quad \texttt{UpperTheory.\allowbreak{}V.\allowbreak{}Base.\allowbreak{}kerResAb} (\texttt{Development.\allowbreak{}AutomorphicArithmetic.\allowbreak{}AutomorphicVarieties.\allowbreak{}GaloisLibrary.\allowbreak{}SelmerTransport.\allowbreak{}KerResAb}, lines 46--151)\\
\texttt{UpperTheory.\allowbreak{}V.\allowbreak{}Chr.\allowbreak{}upV} (\texttt{Development.\allowbreak{}AutomorphicArithmetic.\allowbreak{}AutomorphicVarieties.\allowbreak{}CharacterLevelBounds.\allowbreak{}LiftGens}, lines 24--57)\par}

{\footnotesize\raggedright\noindent \ref{B6.32}.\quad \texttt{UpperTheory.\allowbreak{}V.\allowbreak{}HV.\allowbreak{}LayerLe} (\texttt{Development.\allowbreak{}AutomorphicArithmetic.\allowbreak{}AutomorphicVarieties.\allowbreak{}HeightAndVanishing.\allowbreak{}LayerLe}, lines 12--14)\\
\texttt{UpperTheory.\allowbreak{}V.\allowbreak{}HV.\allowbreak{}CountIV} (\texttt{Development.\allowbreak{}AutomorphicArithmetic.\allowbreak{}AutomorphicVarieties.\allowbreak{}HeightAndVanishing.\allowbreak{}LayerLe}, lines 16--20)\\
\texttt{UpperTheory.\allowbreak{}V.\allowbreak{}HV.\allowbreak{}countIV} (\texttt{Development.\allowbreak{}AutomorphicArithmetic.\allowbreak{}AutomorphicVarieties.\allowbreak{}HeightAndVanishing.\allowbreak{}CountIV}, lines 39--89)\par}

{\footnotesize\raggedright\noindent \ref{B6.33}.\quad \texttt{UpperTheory.\allowbreak{}V.\allowbreak{}HV.\allowbreak{}ivField} (\texttt{Development.\allowbreak{}AutomorphicArithmetic.\allowbreak{}AutomorphicVarieties.\allowbreak{}HeightAndVanishing.\allowbreak{}IvField}, lines 32--82)\par}

{\footnotesize\raggedright\noindent \ref{B6.34}.\quad \texttt{UpperTheory.\allowbreak{}V.\allowbreak{}Z.\allowbreak{}J.\allowbreak{}count\_\allowbreak{}v\_\allowbreak{}of\_\allowbreak{}printed\_\allowbreak{}bound\_\allowbreak{}k\_\allowbreak{}reduced\_\allowbreak{}unitary\_\allowbreak{}lift} (\texttt{Development.\allowbreak{}AutomorphicArithmetic.\allowbreak{}AutomorphicVarieties.\allowbreak{}ModularityArgument.\allowbreak{}ThresholdArguments.\allowbreak{}DenseThresholdCases.\allowbreak{}LinkedUnlE}, lines 55--61)\\
\texttt{UpperTheory.\allowbreak{}V.\allowbreak{}UnitaryThreeInertia.\allowbreak{}countV\_\allowbreak{}of\_\allowbreak{}finiteness\_\allowbreak{}duality\_\allowbreak{}euler\_\allowbreak{}host} (\texttt{Development.\allowbreak{}AutomorphicArithmetic.\allowbreak{}AutomorphicVarieties.\allowbreak{}ModularityArgument.\allowbreak{}ThresholdArguments.\allowbreak{}DenseThresholdCases.\allowbreak{}CountVOf}, lines 33--60)\par}

{\footnotesize\raggedright\noindent \ref{B6.35}.\quad \texttt{UpperTheory.\allowbreak{}V.\allowbreak{}HV.\allowbreak{}countVs} (\texttt{Development.\allowbreak{}AutomorphicArithmetic.\allowbreak{}AutomorphicVarieties.\allowbreak{}HeightAndVanishing.\allowbreak{}CountVs}, lines 86--171)\par}

{\footnotesize\raggedright\noindent \ref{B6.36}.\quad \texttt{UpperTheory.\allowbreak{}V.\allowbreak{}X.\allowbreak{}propZ} (\texttt{Development.\allowbreak{}AutomorphicArithmetic.\allowbreak{}AutomorphicVarieties.\allowbreak{}LevelCohomology.\allowbreak{}PropZ}, lines 11--16)\\
\texttt{UpperTheory.\allowbreak{}V.\allowbreak{}X.\allowbreak{}zOne} (\texttt{Development.\allowbreak{}AutomorphicArithmetic.\allowbreak{}AutomorphicVarieties.\allowbreak{}LevelCohomology.\allowbreak{}TorsClass}, lines 97--205)\\
\texttt{UpperTheory.\allowbreak{}V.\allowbreak{}X.\allowbreak{}zQuad} (\texttt{Development.\allowbreak{}AutomorphicArithmetic.\allowbreak{}AutomorphicVarieties.\allowbreak{}LevelCohomology.\allowbreak{}ZqField}, lines 228--233)\\
\texttt{UpperTheory.\allowbreak{}V.\allowbreak{}HV.\allowbreak{}countZ} (\texttt{Development.\allowbreak{}AutomorphicArithmetic.\allowbreak{}AutomorphicVarieties.\allowbreak{}HeightAndVanishing.\allowbreak{}CountZ}, lines 5--32)\par}

{\footnotesize\raggedright\noindent \ref{B6.36a}.\quad \texttt{UpperTheory.\allowbreak{}V.\allowbreak{}HV.\allowbreak{}redAll} (\texttt{Development.\allowbreak{}AutomorphicArithmetic.\allowbreak{}AutomorphicVarieties.\allowbreak{}HeightAndVanishing.\allowbreak{}RedAll}, lines 6--21)\\
\texttt{UpperTheory.\allowbreak{}V.\allowbreak{}HV.\allowbreak{}redPrime} (\texttt{Development.\allowbreak{}AutomorphicArithmetic.\allowbreak{}AutomorphicVarieties.\allowbreak{}HeightAndVanishing.\allowbreak{}RedPrime}, lines 11--57)\par}

{\footnotesize\raggedright\noindent \ref{B6.37}.\quad \texttt{UpperTheory.\allowbreak{}V.\allowbreak{}HV.\allowbreak{}sumIn} (\texttt{Development.\allowbreak{}AutomorphicArithmetic.\allowbreak{}AutomorphicVarieties.\allowbreak{}HeightAndVanishing.\allowbreak{}FinalLe}, lines 86--263)\par}

{\footnotesize\raggedright\noindent \ref{B6.38}.\quad \texttt{UpperTheory.\allowbreak{}V.\allowbreak{}HV.\allowbreak{}sumOut} (\texttt{Development.\allowbreak{}AutomorphicArithmetic.\allowbreak{}AutomorphicVarieties.\allowbreak{}HeightAndVanishing.\allowbreak{}SumOut}, lines 50--231)\par}

{\footnotesize\raggedright\noindent \ref{B6.39}.\quad \texttt{UpperTheory.\allowbreak{}V.\allowbreak{}HV.\allowbreak{}pointHV\_\allowbreak{}of\_\allowbreak{}counts} (\texttt{Development.\allowbreak{}AutomorphicArithmetic.\allowbreak{}AutomorphicVarieties.\allowbreak{}HeightAndVanishing.\allowbreak{}JointWinOfCounts}, lines 16--18)\\
\texttt{UpperTheory.\allowbreak{}V.\allowbreak{}HV.\allowbreak{}point} (\texttt{Development.\allowbreak{}AutomorphicArithmetic.\allowbreak{}AutomorphicVarieties.\allowbreak{}HeightAndVanishing.\allowbreak{}PointArith}, lines 36--92)\\
\texttt{UpperTheory.\allowbreak{}V.\allowbreak{}BConst.\allowbreak{}J} (\texttt{Development.\allowbreak{}AutomorphicArithmetic.\allowbreak{}AutomorphicVarieties.\allowbreak{}PrimeSetBoundDataOfPrimitiveChar}, lines 18--23)\\
\texttt{UpperTheory.\allowbreak{}V.\allowbreak{}IsBWindow} (\texttt{Development.\allowbreak{}AutomorphicArithmetic.\allowbreak{}AutomorphicVarieties.\allowbreak{}PrimeSetBoundDataOfPrimitiveChar}, lines 30--41)\\
\texttt{UpperTheory.\allowbreak{}V.\allowbreak{}HV.\allowbreak{}thmHV} (\texttt{Development.\allowbreak{}AutomorphicArithmetic.\allowbreak{}AutomorphicVarieties.\allowbreak{}HeightAndVanishing.\allowbreak{}ThmHV}, lines 6--40)\\
\texttt{UpperTheory.\allowbreak{}V.\allowbreak{}HV.\allowbreak{}jointWin\_\allowbreak{}of\_\allowbreak{}counts} (\texttt{Development.\allowbreak{}AutomorphicArithmetic.\allowbreak{}AutomorphicVarieties.\allowbreak{}HeightAndVanishing.\allowbreak{}JointWinOfCounts}, lines 25--28)\par}

{\footnotesize\raggedright\noindent \ref{B6.39}.\quad \texttt{UpperTheory.\allowbreak{}V.\allowbreak{}HV.\allowbreak{}hostCond} (\texttt{Development.\allowbreak{}AutomorphicArithmetic.\allowbreak{}AutomorphicVarieties.\allowbreak{}HeightAndVanishing.\allowbreak{}HostCond}, lines 256--323)\par}

\section*{Chapter~\ref{ch:vwin-automorphic}: Automorphic counts and deformation rigidity}

{\footnotesize\raggedright\noindent Chapter introduction.\quad \texttt{UpperTheory.\allowbreak{}V.\allowbreak{}Af.\allowbreak{}Level} (\texttt{Development.\allowbreak{}AutomorphicArithmetic.\allowbreak{}AutomorphicVarieties.\allowbreak{}AutomorphicForms.\allowbreak{}BlochKatoSelmerVanishes}, lines 111--115)\\
\texttt{UpperTheory.\allowbreak{}V.\allowbreak{}Th.\allowbreak{}Unl} (\texttt{Development.\allowbreak{}AutomorphicArithmetic.\allowbreak{}AutomorphicVarieties.\allowbreak{}TheoremsMuLinked.\allowbreak{}AutP}, lines 10--14)\par}

{\footnotesize\raggedright\noindent \ref{def:p6c5-witness}.\quad \texttt{UpperTheory.\allowbreak{}V.\allowbreak{}Z.\allowbreak{}Win.\allowbreak{}WitA} (\texttt{Development.\allowbreak{}AutomorphicArithmetic.\allowbreak{}AutomorphicVarieties.\allowbreak{}ModularityArgument.\allowbreak{}AutomorphyLifting.\allowbreak{}WeakOfAtt}, lines 27--31)\\
\texttt{UpperTheory.\allowbreak{}V.\allowbreak{}Z.\allowbreak{}Win.\allowbreak{}WitAQ} (\texttt{Development.\allowbreak{}AutomorphicArithmetic.\allowbreak{}AutomorphicVarieties.\allowbreak{}ModularityArgument.\allowbreak{}AutomorphyLifting.\allowbreak{}AAcrossOf}, lines 17--21)\\
\texttt{UpperTheory.\allowbreak{}V.\allowbreak{}Z.\allowbreak{}Win.\allowbreak{}IsAutQA} (\texttt{Development.\allowbreak{}AutomorphicArithmetic.\allowbreak{}AutomorphicVarieties.\allowbreak{}ModularityArgument.\allowbreak{}AutomorphyLifting.\allowbreak{}IsAutCsdA}, lines 16--21)\par}

{\footnotesize\raggedright\noindent \ref{def:p6c5-witness}.\quad \texttt{UpperTheory.\allowbreak{}V.\allowbreak{}Z.\allowbreak{}J.\allowbreak{}Adm3U\_\allowbreak{}s} (\texttt{Development.\allowbreak{}AutomorphicArithmetic.\allowbreak{}AutomorphicVarieties.\allowbreak{}ModularityArgument.\allowbreak{}ThresholdArguments.\allowbreak{}DenseThresholdCases.\allowbreak{}NumberU3US}, lines 27--29)\\
\texttt{UpperTheory.\allowbreak{}V.\allowbreak{}Z.\allowbreak{}Win.\allowbreak{}Dn.\allowbreak{}Adm2A} (\texttt{Development.\allowbreak{}AutomorphicArithmetic.\allowbreak{}AutomorphicVarieties.\allowbreak{}ModularityArgument.\allowbreak{}AutomorphyLifting.\allowbreak{}AutomorphyDescent.\allowbreak{}AutNewformW}, lines 14--20)\par}

{\footnotesize\raggedright\noindent \ref{B6.40m}.\quad \texttt{UpperTheory.\allowbreak{}V.\allowbreak{}U3.\allowbreak{}CohInp} (\texttt{Development.\allowbreak{}AutomorphicArithmetic.\allowbreak{}AutomorphicVarieties.\allowbreak{}UnitaryThreeGroups.\allowbreak{}KatoUse}, lines 19--25)\\
\texttt{UpperTheory.\allowbreak{}V.\allowbreak{}Loc.\allowbreak{}LocalPublishedTheorems} (\texttt{Development.\allowbreak{}AutomorphicArithmetic.\allowbreak{}AutomorphicVarieties.\allowbreak{}PadicHodgeTheory.\allowbreak{}LocalPublishedTheorems}, lines 12--20)\\
\texttt{UpperTheory.\allowbreak{}V.\allowbreak{}U3.\allowbreak{}Dat} (\texttt{Development.\allowbreak{}AutomorphicArithmetic.\allowbreak{}AutomorphicVarieties.\allowbreak{}UnitaryThreeGroups.\allowbreak{}RhoKApply}, lines 16--34)\\
\texttt{UpperTheory.\allowbreak{}V.\allowbreak{}U3.\allowbreak{}Marks} (\texttt{Development.\allowbreak{}AutomorphicArithmetic.\allowbreak{}AutomorphicVarieties.\allowbreak{}UnitaryThreeGroups.\allowbreak{}BaseRepGL3NormalisedOfPsiEqOne}, lines 68--74)\\
\texttt{UpperTheory.\allowbreak{}V.\allowbreak{}U3.\allowbreak{}r0} (\texttt{Development.\allowbreak{}AutomorphicArithmetic.\allowbreak{}AutomorphicVarieties.\allowbreak{}UnitaryThreeGroups.\allowbreak{}UnitaryThreeReflectionRight}, lines 291--298)\\
\texttt{UpperTheory.\allowbreak{}V.\allowbreak{}U3.\allowbreak{}BoxCond} (\texttt{Development.\allowbreak{}AutomorphicArithmetic.\allowbreak{}AutomorphicVarieties.\allowbreak{}UnitaryThreeGroups.\allowbreak{}GRepresents}, lines 100--103)\\
\texttt{UpperTheory.\allowbreak{}V.\allowbreak{}U3.\allowbreak{}NormCond} (\texttt{Development.\allowbreak{}AutomorphicArithmetic.\allowbreak{}AutomorphicVarieties.\allowbreak{}UnitaryThreeGroups.\allowbreak{}GRepresents}, lines 104--106)\\
\texttt{UpperTheory.\allowbreak{}V.\allowbreak{}U3.\allowbreak{}SRing} (\texttt{Development.\allowbreak{}AutomorphicArithmetic.\allowbreak{}AutomorphicVarieties.\allowbreak{}UnitaryThreeGroups.\allowbreak{}ConjStable}, lines 49--50)\par}

{\footnotesize\raggedright\noindent \ref{B6.40n}.\quad \texttt{UpperTheory.\allowbreak{}V.\allowbreak{}U3.\allowbreak{}LocGood} (\texttt{Development.\allowbreak{}AutomorphicArithmetic.\allowbreak{}AutomorphicVarieties.\allowbreak{}UnitaryThreeGroups.\allowbreak{}ConjStable}, lines 34--38)\\
\texttt{UpperTheory.\allowbreak{}V.\allowbreak{}U3.\allowbreak{}ConjStable} (\texttt{Development.\allowbreak{}AutomorphicArithmetic.\allowbreak{}AutomorphicVarieties.\allowbreak{}UnitaryThreeGroups.\allowbreak{}ConjStable}, lines 18--21)\\
\texttt{UpperTheory.\allowbreak{}V.\allowbreak{}U3.\allowbreak{}LocChart} (\texttt{Development.\allowbreak{}AutomorphicArithmetic.\allowbreak{}AutomorphicVarieties.\allowbreak{}UnitaryThreeGroups.\allowbreak{}ConjStable}, lines 25--33)\\
\texttt{UpperTheory.\allowbreak{}V.\allowbreak{}U3.\allowbreak{}locDims} (\texttt{Development.\allowbreak{}AutomorphicArithmetic.\allowbreak{}AutomorphicVarieties.\allowbreak{}UnitaryThreeGroups.\allowbreak{}ConjStable}, lines 22--24)\\
\texttt{UpperTheory.\allowbreak{}V.\allowbreak{}Loc.\allowbreak{}rtChart} (\texttt{Development.\allowbreak{}AutomorphicArithmetic.\allowbreak{}AutomorphicVarieties.\allowbreak{}PadicHodgeTheory.\allowbreak{}RtChart}, lines 41--74)\\
\texttt{UpperTheory.\allowbreak{}V.\allowbreak{}Loc.\allowbreak{}rtPoints} (\texttt{Development.\allowbreak{}AutomorphicArithmetic.\allowbreak{}AutomorphicVarieties.\allowbreak{}PadicHodgeTheory.\allowbreak{}RtPoints}, lines 11--15)\\
\texttt{UpperTheory.\allowbreak{}V.\allowbreak{}Loc.\allowbreak{}lemma2c} (\texttt{Development.\allowbreak{}AutomorphicArithmetic.\allowbreak{}AutomorphicVarieties.\allowbreak{}PadicHodgeTheory.\allowbreak{}ValuationRingClearLastRowEntry}, lines 219--247)\par}

{\footnotesize\raggedright\noindent \ref{B6.40p}.\quad \texttt{UpperTheory.\allowbreak{}V.\allowbreak{}U3.\allowbreak{}TypeI} (\texttt{Development.\allowbreak{}AutomorphicArithmetic.\allowbreak{}AutomorphicVarieties.\allowbreak{}UnitaryThreeGroups.\allowbreak{}BaseRepGL3NormalisedOfPsiEqOne}, lines 20--21)\\
\texttt{UpperTheory.\allowbreak{}V.\allowbreak{}U3.\allowbreak{}TypeII} (\texttt{Development.\allowbreak{}AutomorphicArithmetic.\allowbreak{}AutomorphicVarieties.\allowbreak{}UnitaryThreeGroups.\allowbreak{}BaseRepGL3NormalisedOfPsiEqOne}, lines 30--32)\\
\texttt{UpperTheory.\allowbreak{}V.\allowbreak{}U3.\allowbreak{}typeII\_\allowbreak{}facts} (\texttt{Development.\allowbreak{}AutomorphicArithmetic.\allowbreak{}AutomorphicVarieties.\allowbreak{}UnitaryThreeGroups.\allowbreak{}CoboundaryLemmaNoColumn}, lines 135--216)\\
\texttt{UpperTheory.\allowbreak{}V.\allowbreak{}U3.\allowbreak{}char\_\allowbreak{}eq\_\allowbreak{}one} (\texttt{Development.\allowbreak{}AutomorphicArithmetic.\allowbreak{}AutomorphicVarieties.\allowbreak{}UnitaryThreeGroups.\allowbreak{}PowEqOneOfUnr}, lines 103--130)\\
\texttt{UpperTheory.\allowbreak{}V.\allowbreak{}U3.\allowbreak{}pt\_\allowbreak{}inertia\_\allowbreak{}above} (\texttt{Development.\allowbreak{}AutomorphicArithmetic.\allowbreak{}AutomorphicVarieties.\allowbreak{}UnitaryThreeGroups.\allowbreak{}PtIneL}, lines 80--129)\\
\texttt{UpperTheory.\allowbreak{}V.\allowbreak{}U3.\allowbreak{}pt\_\allowbreak{}inertia\_\allowbreak{}away} (\texttt{Development.\allowbreak{}AutomorphicArithmetic.\allowbreak{}AutomorphicVarieties.\allowbreak{}UnitaryThreeGroups.\allowbreak{}PtInertiaAway}, lines 54--110)\par}

{\footnotesize\raggedright\noindent \ref{B6.40}.\quad \texttt{UpperTheory.\allowbreak{}V.\allowbreak{}Z.\allowbreak{}J.\allowbreak{}oneCharU\_\allowbreak{}of\_\allowbreak{}katoTwoW\_\allowbreak{}rig\_\allowbreak{}admI\_\allowbreak{}automorphy} (\texttt{Development.\allowbreak{}AutomorphicArithmetic.\allowbreak{}AutomorphicVarieties.\allowbreak{}ModularityArgument.\allowbreak{}ThresholdArguments.\allowbreak{}DenseThresholdCases.\allowbreak{}FOneCharU}, lines 75--95)\\
\texttt{UpperTheory.\allowbreak{}V.\allowbreak{}Z.\allowbreak{}J.\allowbreak{}OneCharU\_\allowbreak{}s} (\texttt{Development.\allowbreak{}AutomorphicArithmetic.\allowbreak{}AutomorphicVarieties.\allowbreak{}ModularityArgument.\allowbreak{}ThresholdArguments.\allowbreak{}DenseThresholdCases.\allowbreak{}NumberU3US}, lines 49--59)\par}

{\footnotesize\raggedright\noindent \ref{B6.40a}.\quad \texttt{UpperTheory.\allowbreak{}V.\allowbreak{}U3.\allowbreak{}sha\_\allowbreak{}inj\_\allowbreak{}gen} (\texttt{Development.\allowbreak{}AutomorphicArithmetic.\allowbreak{}AutomorphicVarieties.\allowbreak{}UnitaryThreeGroups.\allowbreak{}ShaInjGen}, lines 338--362)\\
\texttt{UpperTheory.\allowbreak{}exists\_\allowbreak{}base\_\allowbreak{}pointDerivation\_\allowbreak{}of\_\allowbreak{}weight\_\allowbreak{}homogeneous} (\texttt{Development.\allowbreak{}AutomorphicArithmetic.\allowbreak{}AutomorphicVarieties.\allowbreak{}UnitaryThreeGroups.\allowbreak{}ExistsBasePointDerivationOfWeightHomogeneous}, lines 48--135)\\
\texttt{UpperTheory.\allowbreak{}V.\allowbreak{}U3.\allowbreak{}model\_\allowbreak{}lower} (\texttt{Development.\allowbreak{}AutomorphicArithmetic.\allowbreak{}AutomorphicVarieties.\allowbreak{}UnitaryThreeGroups.\allowbreak{}ModelLower}, lines 373--402)\par}

{\footnotesize\raggedright\noindent \ref{B6.41}.\quad \texttt{UpperTheory.\allowbreak{}GalN.\allowbreak{}V.\allowbreak{}Win.\allowbreak{}N.\allowbreak{}numberU3U\_\allowbreak{}of\_\allowbreak{}prints} (\texttt{Development.\allowbreak{}AutomorphicArithmetic.\allowbreak{}AutomorphicVarieties.\allowbreak{}ModularityArgument.\allowbreak{}ThresholdArguments.\allowbreak{}DenseThresholdCases.\allowbreak{}KeepNumU}, lines 79--80)\\
\texttt{UpperTheory.\allowbreak{}V.\allowbreak{}Z.\allowbreak{}J.\allowbreak{}NumberU3U\_\allowbreak{}s} (\texttt{Development.\allowbreak{}AutomorphicArithmetic.\allowbreak{}AutomorphicVarieties.\allowbreak{}ModularityArgument.\allowbreak{}ThresholdArguments.\allowbreak{}DenseThresholdCases.\allowbreak{}NumberU3US}, lines 30--45)\par}

{\footnotesize\raggedright\noindent \ref{B6.41b}.\quad \texttt{UpperTheory.\allowbreak{}V.\allowbreak{}Keep.\allowbreak{}levelCore\_\allowbreak{}alt} (\texttt{Development.\allowbreak{}AutomorphicArithmetic.\allowbreak{}AutomorphicVarieties.\allowbreak{}ModularityArgument.\allowbreak{}ThresholdArguments.\allowbreak{}DenseThresholdCases.\allowbreak{}KeepLevel}, lines 30--82)\\
\texttt{UpperTheory.\allowbreak{}V.\allowbreak{}Keep.\allowbreak{}levelKA\_\allowbreak{}alt} (\texttt{Development.\allowbreak{}AutomorphicArithmetic.\allowbreak{}AutomorphicVarieties.\allowbreak{}ModularityArgument.\allowbreak{}ThresholdArguments.\allowbreak{}DenseThresholdCases.\allowbreak{}KeepLevel}, lines 95--101)\\
\texttt{UpperTheory.\allowbreak{}V.\allowbreak{}Keep.\allowbreak{}levelPA\_\allowbreak{}alt} (\texttt{Development.\allowbreak{}AutomorphicArithmetic.\allowbreak{}AutomorphicVarieties.\allowbreak{}ModularityArgument.\allowbreak{}ThresholdArguments.\allowbreak{}DenseThresholdCases.\allowbreak{}KeepLevel}, lines 83--90)\\
\texttt{UpperTheory.\allowbreak{}V.\allowbreak{}Z.\allowbreak{}Win.\allowbreak{}A4.\allowbreak{}admissible\_\allowbreak{}restriction\_\allowbreak{}level\_\allowbreak{}n\_\allowbreak{}of\_\allowbreak{}printed\_\allowbreak{}results} (\texttt{Development.\allowbreak{}AutomorphicArithmetic.\allowbreak{}AutomorphicVarieties.\allowbreak{}ModularityArgument.\allowbreak{}ThresholdArguments.\allowbreak{}DenseThresholdCases.\allowbreak{}LevelAdmissibleRestrictionHolds}, lines 16--23)\\
\texttt{UpperTheory.\allowbreak{}V.\allowbreak{}Lev.\allowbreak{}levelFRes} (\texttt{Development.\allowbreak{}AutomorphicArithmetic.\allowbreak{}AutomorphicVarieties.\allowbreak{}ModularityArgument.\allowbreak{}SemisimpleLevels.\allowbreak{}LevelFRes}, lines 46--74)\\
\texttt{UpperTheory.\allowbreak{}V.\allowbreak{}Z.\allowbreak{}Win.\allowbreak{}A4.\allowbreak{}safeUP'\_\allowbreak{}alt} (\texttt{Development.\allowbreak{}AutomorphicArithmetic.\allowbreak{}AutomorphicVarieties.\allowbreak{}ModularityArgument.\allowbreak{}ThresholdArguments.\allowbreak{}DenseThresholdCases.\allowbreak{}SafeUP}, lines 16--19)\\
\texttt{UpperTheory.\allowbreak{}V.\allowbreak{}Z.\allowbreak{}Win.\allowbreak{}A4.\allowbreak{}potQ} (\texttt{Development.\allowbreak{}AutomorphicArithmetic.\allowbreak{}AutomorphicVarieties.\allowbreak{}ModularityArgument.\allowbreak{}AutomorphyLifting.\allowbreak{}AdmissibleLevelRestriction.\allowbreak{}PotQOf}, lines 63--73)\\
\texttt{UpperTheory.\allowbreak{}V.\allowbreak{}Z.\allowbreak{}Win.\allowbreak{}SplitIrr} (\texttt{Development.\allowbreak{}AutomorphicArithmetic.\allowbreak{}AutomorphicVarieties.\allowbreak{}ModularityArgument.\allowbreak{}AutomorphyLifting.\allowbreak{}SplitIrr}, lines 46--51)\\
\texttt{UpperTheory.\allowbreak{}V.\allowbreak{}Z.\allowbreak{}Win.\allowbreak{}A4.\allowbreak{}safeRes'} (\texttt{Development.\allowbreak{}AutomorphicArithmetic.\allowbreak{}AutomorphicVarieties.\allowbreak{}ModularityArgument.\allowbreak{}AutomorphyLifting.\allowbreak{}AdmissibleLevelRestriction.\allowbreak{}UnlAllOf}, lines 30--35)\\
\texttt{UpperTheory.\allowbreak{}V.\allowbreak{}Lev.\allowbreak{}unlGap} (\texttt{Development.\allowbreak{}AutomorphicArithmetic.\allowbreak{}AutomorphicVarieties.\allowbreak{}ModularityArgument.\allowbreak{}SemisimpleLevels.\allowbreak{}UnlGap}, lines 8--30)\\
\texttt{UpperTheory.\allowbreak{}V.\allowbreak{}Z.\allowbreak{}Bc.\allowbreak{}polarization\_\allowbreak{}k\_\allowbreak{}of\_\allowbreak{}printed\_\allowbreak{}results} (\texttt{Development.\allowbreak{}AutomorphicArithmetic.\allowbreak{}AutomorphicVarieties.\allowbreak{}ModularityArgument.\allowbreak{}ThresholdArguments.\allowbreak{}DenseThresholdCases.\allowbreak{}UnderCPl}, lines 21--55)\par}

{\footnotesize\raggedright\noindent \ref{B6.41a}.\quad \texttt{UpperTheory.\allowbreak{}V.\allowbreak{}Af.\allowbreak{}card\_\allowbreak{}k\_\allowbreak{}of\_\allowbreak{}printed\_\allowbreak{}results} (\texttt{Development.\allowbreak{}AutomorphicArithmetic.\allowbreak{}AutomorphicVarieties.\allowbreak{}ModularityArgument.\allowbreak{}ThresholdArguments.\allowbreak{}DenseThresholdCases.\allowbreak{}TCardK}, lines 18--36)\\
\texttt{UpperTheory.\allowbreak{}GalN.\allowbreak{}V.\allowbreak{}Win.\allowbreak{}N.\allowbreak{}countQ\_\allowbreak{}alt} (\texttt{Development.\allowbreak{}AutomorphicArithmetic.\allowbreak{}AutomorphicVarieties.\allowbreak{}ModularityArgument.\allowbreak{}ThresholdArguments.\allowbreak{}DenseThresholdCases.\allowbreak{}FrobeniusTraceCountBoundOfCmEigensystemCount}, lines 255--335)\\
\texttt{UpperTheory.\allowbreak{}GalN.\allowbreak{}V.\allowbreak{}Win.\allowbreak{}traces\_\allowbreak{}from\_\allowbreak{}eigenvalues\_\allowbreak{}rank3} (\texttt{Development.\allowbreak{}AutomorphicArithmetic.\allowbreak{}Foundations.\allowbreak{}GaloisRepresentations.\allowbreak{}AdmissibleResidues.\allowbreak{}FrobeniusTraceCountBoundOfRatEigensystemCount}, lines 78--146)\\
\texttt{UpperTheory.\allowbreak{}GalN.\allowbreak{}V.\allowbreak{}Win.\allowbreak{}traces\_\allowbreak{}from\_\allowbreak{}eigenvalues\_\allowbreak{}rank2} (\texttt{Development.\allowbreak{}AutomorphicArithmetic.\allowbreak{}Foundations.\allowbreak{}GaloisRepresentations.\allowbreak{}AdmissibleResidues.\allowbreak{}FrobeniusTraceCountBoundOfRatEigensystemCount}, lines 169--231)\par}

{\footnotesize\raggedright\noindent \ref{B6.42}.\quad \texttt{UpperTheory.\allowbreak{}V.\allowbreak{}Z.\allowbreak{}J.\allowbreak{}propGK\_\allowbreak{}of\_\allowbreak{}numberU3U\_\allowbreak{}of\_\allowbreak{}oneCharU} (\texttt{Development.\allowbreak{}AutomorphicArithmetic.\allowbreak{}AutomorphicVarieties.\allowbreak{}ModularityArgument.\allowbreak{}ThresholdArguments.\allowbreak{}DenseThresholdCases.\allowbreak{}FPropGKU}, lines 110--172)\\
\texttt{UpperTheory.\allowbreak{}V.\allowbreak{}U3.\allowbreak{}PropGK} (\texttt{Development.\allowbreak{}AutomorphicArithmetic.\allowbreak{}AutomorphicVarieties.\allowbreak{}UnitaryThreeGroups.\allowbreak{}KatoUse}, lines 36--43)\par}

{\footnotesize\raggedright\noindent \ref{B6.42}.\quad \texttt{GlobalTheory.\allowbreak{}Ord2.\allowbreak{}Nf} (\texttt{Development.\allowbreak{}AbelianVarietyTheory.\allowbreak{}Definitions.\allowbreak{}Foundations.\allowbreak{}FinalTheorems}, lines 53--67)\\
\texttt{GlobalTheory.\allowbreak{}Ord2.\allowbreak{}Gr} (\texttt{Development.\allowbreak{}AbelianVarietyTheory.\allowbreak{}Definitions.\allowbreak{}Foundations.\allowbreak{}FinalTheorems}, lines 72--73)\\
\texttt{GlobalTheory.\allowbreak{}Ord2.\allowbreak{}Unr} (\texttt{Development.\allowbreak{}AbelianVarietyTheory.\allowbreak{}Definitions.\allowbreak{}Foundations.\allowbreak{}FinalTheorems}, lines 68--70)\par}

{\footnotesize\raggedright\noindent \ref{B6.42}.\quad \texttt{UpperTheory.\allowbreak{}V.\allowbreak{}Z.\allowbreak{}Kato.\allowbreak{}SphAt} (\texttt{Development.\allowbreak{}AutomorphicArithmetic.\allowbreak{}AutomorphicVarieties.\allowbreak{}ModularityArgument.\allowbreak{}KatoQuadraticCharacters.\allowbreak{}KatoTwoSOfTwo}, lines 9--17)\par}

{\footnotesize\raggedright\noindent \ref{B6.42}.\quad \texttt{UpperTheory.\allowbreak{}V.\allowbreak{}Af.\allowbreak{}IsBK} (\texttt{Development.\allowbreak{}AutomorphicArithmetic.\allowbreak{}AutomorphicVarieties.\allowbreak{}AutomorphicForms.\allowbreak{}BlochKatoSelmerVanishes}, lines 135--143)\\
\texttt{UpperTheory.\allowbreak{}V.\allowbreak{}Loc.\allowbreak{}finiteCocycles} (\texttt{Development.\allowbreak{}AutomorphicArithmetic.\allowbreak{}AutomorphicVarieties.\allowbreak{}PadicHodgeTheory.\allowbreak{}FiniteCohomology}, lines 112--126)\\
\texttt{UpperTheory.\allowbreak{}V.\allowbreak{}Loc.\allowbreak{}tiedP} (\texttt{Development.\allowbreak{}AutomorphicArithmetic.\allowbreak{}AutomorphicVarieties.\allowbreak{}PadicHodgeTheory.\allowbreak{}FiniteCohomology}, lines 84--85)\par}

{\footnotesize\raggedright\noindent \ref{B6.43b}.\quad \texttt{UpperTheory.\allowbreak{}V.\allowbreak{}Af.\allowbreak{}Newform} (\texttt{Development.\allowbreak{}AutomorphicArithmetic.\allowbreak{}AutomorphicVarieties.\allowbreak{}AutomorphicForms.\allowbreak{}Newform}, lines 21--34)\\
\texttt{UpperTheory.\allowbreak{}V.\allowbreak{}Af.\allowbreak{}Newform.\allowbreak{}Att} (\texttt{Development.\allowbreak{}AutomorphicArithmetic.\allowbreak{}AutomorphicVarieties.\allowbreak{}AutomorphicForms.\allowbreak{}Newform}, lines 43--48)\\
\texttt{UpperTheory.\allowbreak{}V.\allowbreak{}Lit.\allowbreak{}KatoLit} (\texttt{Development.\allowbreak{}AutomorphicArithmetic.\allowbreak{}AutomorphicVarieties.\allowbreak{}LiteratureResults.\allowbreak{}CoefFieldEq}, lines 58--67)\\
\texttt{UpperTheory.\allowbreak{}V.\allowbreak{}U3.\allowbreak{}KatoUse} (\texttt{Development.\allowbreak{}AutomorphicArithmetic.\allowbreak{}AutomorphicVarieties.\allowbreak{}UnitaryThreeGroups.\allowbreak{}KatoUse}, lines 13--18)\par}

{\footnotesize\raggedright\noindent \ref{B6.43}.\quad \texttt{UpperTheory.\allowbreak{}V.\allowbreak{}Z.\allowbreak{}Kato.\allowbreak{}WeilSphericalRecords.\allowbreak{}KatoTwoW} (\texttt{Development.\allowbreak{}AutomorphicArithmetic.\allowbreak{}AutomorphicVarieties.\allowbreak{}ModularityArgument.\allowbreak{}KatoQuadraticCharacters.\allowbreak{}SphericalDecomposition.\allowbreak{}KatoTwoW}, lines 10--15)\\
\texttt{UpperTheory.\allowbreak{}V.\allowbreak{}Z.\allowbreak{}Kato.\allowbreak{}Used} (\texttt{Development.\allowbreak{}AutomorphicArithmetic.\allowbreak{}AutomorphicVarieties.\allowbreak{}ModularityArgument.\allowbreak{}KatoQuadraticCharacters.\allowbreak{}KatoTwoOfUseN}, lines 11--42)\\
\texttt{UpperTheory.\allowbreak{}V.\allowbreak{}Z.\allowbreak{}Kato.\allowbreak{}WeilSphericalRecords.\allowbreak{}katoTwoW\_\allowbreak{}of\_\allowbreak{}printsH} (\texttt{Development.\allowbreak{}AutomorphicArithmetic.\allowbreak{}AutomorphicVarieties.\allowbreak{}ModularityArgument.\allowbreak{}ThresholdArguments.\allowbreak{}DenseThresholdCases.\allowbreak{}KatoTwoWOfPublishedH}, lines 48--49)\\
\texttt{UpperTheory.\allowbreak{}V.\allowbreak{}Z.\allowbreak{}Kato.\allowbreak{}WeilSphericalRecords.\allowbreak{}katoTwoW\_\allowbreak{}of\_\allowbreak{}printsH\_\allowbreak{}of\_\allowbreak{}rT\_\allowbreak{}of\_\allowbreak{}quadBC2} (\texttt{Development.\allowbreak{}AutomorphicArithmetic.\allowbreak{}AutomorphicVarieties.\allowbreak{}ModularityArgument.\allowbreak{}ThresholdArguments.\allowbreak{}DenseThresholdCases.\allowbreak{}KatoTwoWOfPublishedH}, lines 38--47)\par}

{\footnotesize\raggedright\noindent \ref{B6.43a}.\quad \texttt{UpperTheory.\allowbreak{}V.\allowbreak{}Z.\allowbreak{}Kato.\allowbreak{}WeilSphericalRecords.\allowbreak{}srcS\_\allowbreak{}alt} (\texttt{Development.\allowbreak{}AutomorphicArithmetic.\allowbreak{}AutomorphicVarieties.\allowbreak{}ModularityArgument.\allowbreak{}ThresholdArguments.\allowbreak{}DenseThresholdCases.\allowbreak{}TwistNewform}, lines 116--118)\\
\texttt{GlobalTheory.\allowbreak{}Ord2.\allowbreak{}exists\_\allowbreak{}nf\_\allowbreak{}srcUr} (\texttt{Development.\allowbreak{}AbelianVarietyTheory.\allowbreak{}Proofs.\allowbreak{}FoundationProofs.\allowbreak{}ExistsOrdinaryNormalFormUnramifiedSource}, lines 36--140)\par}

{\footnotesize\raggedright\noindent \ref{B6.44}.\quad \texttt{UpperTheory.\allowbreak{}V.\allowbreak{}Z.\allowbreak{}Kato.\allowbreak{}WeilSphericalRecords.\allowbreak{}twoS} (\texttt{Development.\allowbreak{}AutomorphicArithmetic.\allowbreak{}AutomorphicVarieties.\allowbreak{}ModularityArgument.\allowbreak{}KatoQuadraticCharacters.\allowbreak{}SphericalDecomposition.\allowbreak{}SphAtTw}, lines 16--35)\\
\texttt{GlobalTheory.\allowbreak{}Ord2.\allowbreak{}vanishOrdL\_\allowbreak{}of\_\allowbreak{}zetaUr} (\texttt{Development.\allowbreak{}AbelianVarietyTheory.\allowbreak{}Proofs.\allowbreak{}FoundationProofs.\allowbreak{}VanishOrdLOfZetaUr}, lines 20--26)\par}

{\footnotesize\raggedright\noindent \ref{B6.46d}.\quad \texttt{GlobalTheory.\allowbreak{}EisM.\allowbreak{}lseries\_\allowbreak{}ne\_\allowbreak{}zero\_\allowbreak{}edge} (\texttt{Development.\allowbreak{}AbelianVarietyTheory.\allowbreak{}Proofs.\allowbreak{}FoundationProofs.\allowbreak{}LseriesNeZeroReGeOneOfPositivity}, lines 21--117)\\
\texttt{GlobalTheory.\allowbreak{}EisM.\allowbreak{}twisted\_\allowbreak{}lvalue\_\allowbreak{}ne\_\allowbreak{}zero\_\allowbreak{}edge} (\texttt{Development.\allowbreak{}AbelianVarietyTheory.\allowbreak{}Proofs.\allowbreak{}FoundationProofs.\allowbreak{}TwistedLvalueNeZeroEdge}, lines 167--220)\par}

{\footnotesize\raggedright\noindent \ref{B6.46d}.\quad \texttt{UpperTheory.\allowbreak{}V.\allowbreak{}GradeZeroInterpolatingElement.\allowbreak{}Raw.\allowbreak{}PtIn} (\texttt{Development.\allowbreak{}AutomorphicArithmetic.\allowbreak{}AutomorphicVarieties.\allowbreak{}SymplecticGSp4.\allowbreak{}DecompSigmaOrd}, lines 164--167)\\
\texttt{UpperTheory.\allowbreak{}V.\allowbreak{}GradeZeroInterpolatingElement.\allowbreak{}Raw.\allowbreak{}EPt} (\texttt{Development.\allowbreak{}AutomorphicArithmetic.\allowbreak{}AutomorphicVarieties.\allowbreak{}SymplecticGSp4.\allowbreak{}DecompSigmaOrd}, lines 168--169)\par}

{\footnotesize\raggedright\noindent \ref{B6.48j}.\quad \texttt{UpperTheory.\allowbreak{}V.\allowbreak{}Loc.\allowbreak{}spPtL} (\texttt{Development.\allowbreak{}AutomorphicArithmetic.\allowbreak{}AutomorphicVarieties.\allowbreak{}PadicHodgeTheory.\allowbreak{}SplAdjUnit}, lines 264--328)\par}

{\footnotesize\raggedright\noindent \ref{B6.48h}.\quad \texttt{UpperTheory.\allowbreak{}V.\allowbreak{}Loc.\allowbreak{}prC} (\texttt{Development.\allowbreak{}AutomorphicArithmetic.\allowbreak{}AutomorphicVarieties.\allowbreak{}PadicHodgeTheory.\allowbreak{}SnocZeroLastCoordinateLast}, lines 345--352)\par}

{\footnotesize\raggedright\noindent \ref{B6.48k}.\quad \texttt{UpperTheory.\allowbreak{}V.\allowbreak{}Th.\allowbreak{}spPt} (\texttt{Development.\allowbreak{}AutomorphicArithmetic.\allowbreak{}AutomorphicVarieties.\allowbreak{}TheoremsMuLinked.\allowbreak{}GaloisRepresentationTypes.\allowbreak{}MapRepComp}, lines 220--286)\\
\texttt{UpperTheory.\allowbreak{}V.\allowbreak{}Ws.\allowbreak{}P1.\allowbreak{}condC} (\texttt{Development.\allowbreak{}AutomorphicArithmetic.\allowbreak{}AutomorphicVarieties.\allowbreak{}WeakShimuraTheory.\allowbreak{}ProjectiveLineFlags.\allowbreak{}CondC}, lines 10--93)\\
\texttt{UpperTheory.\allowbreak{}V.\allowbreak{}Th.\allowbreak{}nice0four\_\allowbreak{}alt} (\texttt{Development.\allowbreak{}AutomorphicArithmetic.\allowbreak{}AutomorphicVarieties.\allowbreak{}ModularityArgument.\allowbreak{}ThresholdArguments.\allowbreak{}DenseThresholdCases.\allowbreak{}TNiceFour}, lines 114--130)\\
\texttt{adjoint\_\allowbreak{}vanishing\_\allowbreak{}gsp4\_\allowbreak{}rigid\_\allowbreak{}point} (\texttt{Development.\allowbreak{}AutomorphicArithmetic.\allowbreak{}AutomorphicVarieties.\allowbreak{}ModularityArgument.\allowbreak{}ThresholdArguments.\allowbreak{}DenseThresholdCases.\allowbreak{}AdjointVanishingGsp4RigidPoint}, lines 31--101)\par}

{\footnotesize\raggedright\noindent \ref{B6.48}.\quad \texttt{UpperTheory.\allowbreak{}V.\allowbreak{}Th.\allowbreak{}rigSp4\_\allowbreak{}alt} (\texttt{Development.\allowbreak{}AutomorphicArithmetic.\allowbreak{}AutomorphicVarieties.\allowbreak{}ModularityArgument.\allowbreak{}ThresholdArguments.\allowbreak{}DenseThresholdCases.\allowbreak{}RigSp4}, lines 181--194)\\
\texttt{UpperTheory.\allowbreak{}V.\allowbreak{}GradeZeroInterpolatingElement.\allowbreak{}Raw.\allowbreak{}RigidAt} (\texttt{Development.\allowbreak{}AutomorphicArithmetic.\allowbreak{}AutomorphicVarieties.\allowbreak{}SymplecticGSp4.\allowbreak{}DecompSigmaOrd}, lines 182--184)\par}

{\footnotesize\raggedright\noindent \ref{B6.48i}.\quad \texttt{UpperTheory.\allowbreak{}V.\allowbreak{}AlgebraTools.\allowbreak{}hkDualEx} (\texttt{Development.\allowbreak{}AutomorphicArithmetic.\allowbreak{}AutomorphicVarieties.\allowbreak{}AlgebraTools.\allowbreak{}HkDualEx}, lines 14--42)\\
\texttt{UpperTheory.\allowbreak{}V.\allowbreak{}AlgebraTools.\allowbreak{}hkSandwich} (\texttt{Development.\allowbreak{}AutomorphicArithmetic.\allowbreak{}AutomorphicVarieties.\allowbreak{}AlgebraTools.\allowbreak{}HkSandwich}, lines 92--129)\\
\texttt{UpperTheory.\allowbreak{}V.\allowbreak{}Pa.\allowbreak{}Lev.\allowbreak{}multiTwo} (\texttt{Development.\allowbreak{}AutomorphicArithmetic.\allowbreak{}AutomorphicVarieties.\allowbreak{}PatchingModules.\allowbreak{}NilpotentLevelLowering.\allowbreak{}MultiTwo}, lines 174--315)\\
\texttt{UpperTheory.\allowbreak{}V.\allowbreak{}Pa.\allowbreak{}Lev.\allowbreak{}twoOf\_\allowbreak{}alt} (\texttt{Development.\allowbreak{}AutomorphicArithmetic.\allowbreak{}AutomorphicVarieties.\allowbreak{}ModularityArgument.\allowbreak{}ThresholdArguments.\allowbreak{}DenseThresholdCases.\allowbreak{}LevTwoOf}, lines 103--117)\par}

{\footnotesize\raggedright\noindent \ref{B6.48ca}.\quad \texttt{UpperTheory.\allowbreak{}V.\allowbreak{}Loc.\allowbreak{}ordFlag\_\allowbreak{}alt} (\texttt{Development.\allowbreak{}AutomorphicArithmetic.\allowbreak{}AutomorphicVarieties.\allowbreak{}ModularityArgument.\allowbreak{}ThresholdArguments.\allowbreak{}DenseThresholdCases.\allowbreak{}LocOrdFlag}, lines 158--168)\par}

{\footnotesize\raggedright\noindent \ref{B6.48cb}.\quad \texttt{UpperTheory.\allowbreak{}V.\allowbreak{}Hi.\allowbreak{}ordStruct} (\texttt{Development.\allowbreak{}AutomorphicArithmetic.\allowbreak{}AutomorphicVarieties.\allowbreak{}HeckeAlgebras.\allowbreak{}OrdinaryHeckeCharacters.\allowbreak{}OrdFree}, lines 295--438)\\
\texttt{UpperTheory.\allowbreak{}V.\allowbreak{}Hi.\allowbreak{}mem\_\allowbreak{}ord\_\allowbreak{}iff} (\texttt{Development.\allowbreak{}AutomorphicArithmetic.\allowbreak{}AutomorphicVarieties.\allowbreak{}HeckeAlgebras.\allowbreak{}OrdinaryHeckeCharacters.\allowbreak{}OrdFree}, lines 263--294)\\
\texttt{UpperTheory.\allowbreak{}exists\_\allowbreak{}section\_\allowbreak{}automorphicSpace\_\allowbreak{}with\_\allowbreak{}prescribed\_\allowbreak{}values} (\texttt{Development.\allowbreak{}AutomorphicArithmetic.\allowbreak{}AutomorphicVarieties.\allowbreak{}BuildingsResidueChars.\allowbreak{}ExistsSectionAutomorphicSpaceWithPrescribedValue}, lines 8--104)\\
\texttt{UpperTheory.\allowbreak{}V.\allowbreak{}Hi.\allowbreak{}HSet} (\texttt{Development.\allowbreak{}AutomorphicArithmetic.\allowbreak{}AutomorphicVarieties.\allowbreak{}HeckeAlgebras.\allowbreak{}CoefficientRepTrivialOnUnramified}, lines 75--108)\par}

{\footnotesize\raggedright\noindent \ref{B6.48c}.\quad \texttt{UpperTheory.\allowbreak{}V.\allowbreak{}CA.\allowbreak{}universalDeformationRing\_\allowbreak{}nonempty\_\allowbreak{}of\_\allowbreak{}nonempty\_\allowbreak{}pair} (\texttt{Development.\allowbreak{}AutomorphicArithmetic.\allowbreak{}AutomorphicVarieties.\allowbreak{}CommutativeAlgebraHeight.\allowbreak{}UniversalDeformationRingNonemptyOfNonemptyPair}, lines 243--320)\\
\texttt{UpperTheory.\allowbreak{}V.\allowbreak{}Th.\allowbreak{}AnnihilatorQuotientFinite.\allowbreak{}levelX\_\allowbreak{}of\_\allowbreak{}isJ\_\allowbreak{}unitaryHeckeData} (\texttt{Development.\allowbreak{}AutomorphicArithmetic.\allowbreak{}AutomorphicVarieties.\allowbreak{}ModularityArgument.\allowbreak{}ThresholdArguments.\allowbreak{}DenseThresholdCases.\allowbreak{}HidaLevelX}, lines 192--198)\\
\texttt{UpperTheory.\allowbreak{}V.\allowbreak{}Th.\allowbreak{}has\_\allowbreak{}operator\_\allowbreak{}statement} (\texttt{Development.\allowbreak{}AutomorphicArithmetic.\allowbreak{}AutomorphicVarieties.\allowbreak{}ModularityArgument.\allowbreak{}ThresholdArguments.\allowbreak{}DenseThresholdCases.\allowbreak{}THasOp}, lines 18--45)\\
\texttt{UpperTheory.\allowbreak{}V.\allowbreak{}W.\allowbreak{}HidaSetting} (\texttt{Development.\allowbreak{}AutomorphicArithmetic.\allowbreak{}AutomorphicVarieties.\allowbreak{}WeilGroupCharacters.\allowbreak{}ContinuousResP}, lines 82--92)\par}

{\footnotesize\raggedright\noindent \ref{B6.48d}.\quad \texttt{UpperTheory.\allowbreak{}V.\allowbreak{}Hi.\allowbreak{}spec} (\texttt{Development.\allowbreak{}AutomorphicArithmetic.\allowbreak{}AutomorphicVarieties.\allowbreak{}HeckeAlgebras.\allowbreak{}OrdinaryHeckeCharacters.\allowbreak{}SpecIntIff}, lines 100--202)\\
\texttt{UpperTheory.\allowbreak{}V.\allowbreak{}Hi.\allowbreak{}ctrlOrd} (\texttt{Development.\allowbreak{}AutomorphicArithmetic.\allowbreak{}AutomorphicVarieties.\allowbreak{}HeckeAlgebras.\allowbreak{}OrdinaryHeckeCharacters.\allowbreak{}CtrlOrd}, lines 220--372)\\
\texttt{UpperTheory.\allowbreak{}V.\allowbreak{}Hi.\allowbreak{}ordBC} (\texttt{Development.\allowbreak{}AutomorphicArithmetic.\allowbreak{}AutomorphicVarieties.\allowbreak{}HeckeAlgebras.\allowbreak{}OrdinaryHeckeCharacters.\allowbreak{}OrdBCMapOrd}, lines 198--233)\par}

{\footnotesize\raggedright\noindent \ref{B6.48a}.\quad \texttt{UpperTheory.\allowbreak{}Gal.\allowbreak{}projectedFormsSystemPart.\allowbreak{}irrArithPointsAutomorphicAway\_\allowbreak{}of\_\allowbreak{}isJ} (\texttt{Development.\allowbreak{}AutomorphicArithmetic.\allowbreak{}AutomorphicVarieties.\allowbreak{}ModularityArgument.\allowbreak{}ThresholdArguments.\allowbreak{}DenseThresholdCases.\allowbreak{}EndoscopyRank2Ordinary}, lines 48--72)\par}

{\footnotesize\raggedright\noindent \ref{B6.48ab}.\quad \texttt{UpperTheory.\allowbreak{}V.\allowbreak{}Pa.\allowbreak{}linAlg} (\texttt{Development.\allowbreak{}AutomorphicArithmetic.\allowbreak{}AutomorphicVarieties.\allowbreak{}PatchingModules.\allowbreak{}SelmerScaling.\allowbreak{}LinAlg}, lines 45--127)\par}

{\footnotesize\raggedright\noindent \ref{B6.48ac}.\quad \texttt{UpperTheory.\allowbreak{}V.\allowbreak{}Pa.\allowbreak{}Core.\allowbreak{}pmDim} (\texttt{Development.\allowbreak{}AutomorphicArithmetic.\allowbreak{}AutomorphicVarieties.\allowbreak{}PatchingModules.\allowbreak{}AnnihilatorBaseChange.\allowbreak{}PmDimQuot}, lines 99--150)\par}

{\footnotesize\raggedright\noindent \ref{B6.48ad}.\quad \texttt{UpperTheory.\allowbreak{}V.\allowbreak{}AlgebraTools.\allowbreak{}exists\_\allowbreak{}least\_\allowbreak{}prime\_\allowbreak{}detIdeal} (\texttt{Development.\allowbreak{}AutomorphicArithmetic.\allowbreak{}AutomorphicVarieties.\allowbreak{}AlgebraTools.\allowbreak{}CaDetMin}, lines 11--51)\\
\texttt{UpperTheory.\allowbreak{}V.\allowbreak{}AlgebraTools.\allowbreak{}le\_\allowbreak{}of\_\allowbreak{}detIdeal\_\allowbreak{}le} (\texttt{Development.\allowbreak{}AutomorphicArithmetic.\allowbreak{}AutomorphicVarieties.\allowbreak{}AlgebraTools.\allowbreak{}CaDetLocus}, lines 10--59)\par}

{\footnotesize\raggedright\noindent \ref{B6.48aad}.\quad \texttt{UpperTheory.\allowbreak{}V.\allowbreak{}Pa.\allowbreak{}locLemma} (\texttt{Development.\allowbreak{}AutomorphicArithmetic.\allowbreak{}AutomorphicVarieties.\allowbreak{}PatchingModules.\allowbreak{}SelmerScaling.\allowbreak{}LocLemma}, lines 225--356)\par}

{\footnotesize\raggedright\noindent \ref{B6.48aae}.\quad \texttt{UpperTheory.\allowbreak{}V.\allowbreak{}Pa.\allowbreak{}Core.\allowbreak{}pmodEx} (\texttt{Development.\allowbreak{}AutomorphicArithmetic.\allowbreak{}AutomorphicVarieties.\allowbreak{}PatchingModules.\allowbreak{}AnnihilatorBaseChange.\allowbreak{}PmodEx}, lines 32--96)\\
\texttt{UpperTheory.\allowbreak{}V.\allowbreak{}Pa.\allowbreak{}Core.\allowbreak{}exists\_\allowbreak{}ultraLimit\_\allowbreak{}pair\_\allowbreak{}of\_\allowbreak{}eventual\_\allowbreak{}quotient\_\allowbreak{}maps} (\texttt{Development.\allowbreak{}AutomorphicArithmetic.\allowbreak{}AutomorphicVarieties.\allowbreak{}PatchingModules.\allowbreak{}AnnihilatorBaseChange.\allowbreak{}ExistsUltraLimitPairOfEventualQuotientMaps}, lines 48--106)\\
\texttt{UpperTheory.\allowbreak{}V.\allowbreak{}Pa.\allowbreak{}Core.\allowbreak{}quotientRegularAt\_\allowbreak{}of\_\allowbreak{}quotient\_\allowbreak{}free\_\allowbreak{}over\_\allowbreak{}powerSeries} (\texttt{Development.\allowbreak{}AutomorphicArithmetic.\allowbreak{}AutomorphicVarieties.\allowbreak{}PatchingModules.\allowbreak{}AnnihilatorBaseChange.\allowbreak{}QuotientRegularAtOfQuotientFreeOverPowerSeries}, lines 9--56)\par}

{\footnotesize\raggedright\noindent \ref{B6.48aa}.\quad \texttt{UpperTheory.\allowbreak{}V.\allowbreak{}W.\allowbreak{}modU\_\allowbreak{}of\_\allowbreak{}printsH} (\texttt{Development.\allowbreak{}AutomorphicArithmetic.\allowbreak{}AutomorphicVarieties.\allowbreak{}ModularityArgument.\allowbreak{}ThresholdArguments.\allowbreak{}DenseThresholdCases.\allowbreak{}CoefIsCoef}, lines 16--22)\\
\texttt{UpperTheory.\allowbreak{}V.\allowbreak{}Pa.\allowbreak{}pFromAut\_\allowbreak{}of} (\texttt{Development.\allowbreak{}AutomorphicArithmetic.\allowbreak{}AutomorphicVarieties.\allowbreak{}PatchingModules.\allowbreak{}ChebotarevFamilies.\allowbreak{}PFromAutOf}, lines 116--132)\par}

{\footnotesize\raggedright\noindent \ref{B6.48aaa}.\quad \texttt{UpperTheory.\allowbreak{}V.\allowbreak{}Pa.\allowbreak{}W1\_\allowbreak{}bound\_\allowbreak{}arithmetic\_\allowbreak{}rank\_\allowbreak{}two} (\texttt{Development.\allowbreak{}AutomorphicArithmetic.\allowbreak{}AutomorphicVarieties.\allowbreak{}PatchingModules.\allowbreak{}EpsilonEtaMatrices.\allowbreak{}W1BoundArithmeticRankTwo}, lines 40--166)\\
\texttt{UpperTheory.\allowbreak{}V.\allowbreak{}Pa.\allowbreak{}B2.\allowbreak{}obstructionBound\_\allowbreak{}of\_\allowbreak{}duality} (\texttt{Development.\allowbreak{}AutomorphicArithmetic.\allowbreak{}AutomorphicVarieties.\allowbreak{}PatchingModules.\allowbreak{}ArchimedeanH0Cards.\allowbreak{}ObstructionBoundOfDuality}, lines 145--217)\par}

{\footnotesize\raggedright\noindent \ref{B6.48aab}.\quad \texttt{UpperTheory.\allowbreak{}V.\allowbreak{}Pa.\allowbreak{}cohLim} (\texttt{Development.\allowbreak{}AutomorphicArithmetic.\allowbreak{}AutomorphicVarieties.\allowbreak{}PatchingModules.\allowbreak{}AutomorphicDetection.\allowbreak{}CohLim}, lines 12--20)\par}

{\footnotesize\raggedright\noindent \ref{B6.48aac}.\quad \texttt{UpperTheory.\allowbreak{}V.\allowbreak{}Pa.\allowbreak{}Core.\allowbreak{}pmodEx} (\texttt{Development.\allowbreak{}AutomorphicArithmetic.\allowbreak{}AutomorphicVarieties.\allowbreak{}PatchingModules.\allowbreak{}AnnihilatorBaseChange.\allowbreak{}PmodEx}, lines 32--96)\\
\texttt{UpperTheory.\allowbreak{}V.\allowbreak{}Pa.\allowbreak{}Core.\allowbreak{}pmDim} (\texttt{Development.\allowbreak{}AutomorphicArithmetic.\allowbreak{}AutomorphicVarieties.\allowbreak{}PatchingModules.\allowbreak{}AnnihilatorBaseChange.\allowbreak{}PmDimQuot}, lines 99--150)\par}

{\footnotesize\raggedright\noindent \ref{B6.48bb}.\quad \texttt{UpperTheory.\allowbreak{}V.\allowbreak{}Ws.\allowbreak{}Ch.\allowbreak{}nice} (\texttt{Development.\allowbreak{}AutomorphicArithmetic.\allowbreak{}AutomorphicVarieties.\allowbreak{}WeakShimuraTheory.\allowbreak{}ValuationCharacters.\allowbreak{}ConstOfFinite}, lines 101--125)\\
\texttt{UpperTheory.\allowbreak{}V.\allowbreak{}Pr.\allowbreak{}nicePrimeEx} (\texttt{Development.\allowbreak{}AutomorphicArithmetic.\allowbreak{}AutomorphicVarieties.\allowbreak{}CurveDVRSpecialization.\allowbreak{}NicePrimeEx}, lines 87--126)\par}

{\footnotesize\raggedright\noindent \ref{B6.48be}.\quad \texttt{UpperTheory.\allowbreak{}V.\allowbreak{}GradeZeroInterpolatingElement.\allowbreak{}CohomologyHyp} (\texttt{Development.\allowbreak{}AutomorphicArithmetic.\allowbreak{}AutomorphicVarieties.\allowbreak{}SymplecticGSp4.\allowbreak{}EntryMem}, lines 161--168)\par}

{\footnotesize\raggedright\noindent \ref{B6.48bf}.\quad \texttt{UpperTheory.\allowbreak{}V.\allowbreak{}Ws.\allowbreak{}x1} (\texttt{Development.\allowbreak{}AutomorphicArithmetic.\allowbreak{}AutomorphicVarieties.\allowbreak{}WeakShimuraTheory.\allowbreak{}CompactIntegralFields.\allowbreak{}IncLoc}, lines 97--118)\par}

{\footnotesize\raggedright\noindent \ref{B6.48b}.\quad \texttt{UpperTheory.\allowbreak{}V.\allowbreak{}Cu.\allowbreak{}chain\_\allowbreak{}rank\_\allowbreak{}four\_\allowbreak{}of\_\allowbreak{}red\_\allowbreak{}a\_\allowbreak{}y\_\allowbreak{}coin\_\allowbreak{}fib\_\allowbreak{}y} (\texttt{Development.\allowbreak{}AutomorphicArithmetic.\allowbreak{}AutomorphicVarieties.\allowbreak{}ModularityArgument.\allowbreak{}ThresholdArguments.\allowbreak{}DenseThresholdCases.\allowbreak{}ReductionTypeChainOfRedAYCoinFibY}, lines 213--268)\\
\texttt{UpperTheory.\allowbreak{}V.\allowbreak{}Cu.\allowbreak{}RankFourReductionTypeChain.\allowbreak{}ch\_\allowbreak{}core4} (\texttt{Development.\allowbreak{}AutomorphicArithmetic.\allowbreak{}AutomorphicVarieties.\allowbreak{}ModularityArgument.\allowbreak{}ThresholdArguments.\allowbreak{}DenseThresholdCases.\allowbreak{}ReductionTypeChainOfRedAYCoinFibY}, lines 78--211)\\
\texttt{UpperTheory.\allowbreak{}V.\allowbreak{}Ws.\allowbreak{}residual\_\allowbreak{}charpoly\_\allowbreak{}rank\_\allowbreak{}four} (\texttt{Development.\allowbreak{}AutomorphicArithmetic.\allowbreak{}AutomorphicVarieties.\allowbreak{}ModularityArgument.\allowbreak{}ThresholdArguments.\allowbreak{}DenseThresholdCases.\allowbreak{}ClassicalFixedByIwahori}, lines 90--126)\par}

{\footnotesize\raggedright\noindent \ref{B6.48bc}.\quad \texttt{UpperTheory.\allowbreak{}V.\allowbreak{}Ws.\allowbreak{}redTan} (\texttt{Development.\allowbreak{}AutomorphicArithmetic.\allowbreak{}AutomorphicVarieties.\allowbreak{}WeakShimuraTheory.\allowbreak{}CompactIntegralFields.\allowbreak{}RedTan}, lines 319--353)\\
\texttt{UpperTheory.\allowbreak{}V.\allowbreak{}Cu.\allowbreak{}redTanA} (\texttt{Development.\allowbreak{}AutomorphicArithmetic.\allowbreak{}AutomorphicVarieties.\allowbreak{}ModularityArgument.\allowbreak{}CuspidalTwists.\allowbreak{}PRedTanA}, lines 258--268)\\
\texttt{UpperTheory.\allowbreak{}V.\allowbreak{}Cu.\allowbreak{}yCoin} (\texttt{Development.\allowbreak{}AutomorphicArithmetic.\allowbreak{}AutomorphicVarieties.\allowbreak{}ModularityArgument.\allowbreak{}CuspidalTwists.\allowbreak{}PYCoin}, lines 117--153)\par}

{\footnotesize\raggedright\noindent \ref{B6.48bd}.\quad \texttt{UpperTheory.\allowbreak{}V.\allowbreak{}Ws.\allowbreak{}off\_\allowbreak{}diagonal\_\allowbreak{}nilpotent\_\allowbreak{}rank\_\allowbreak{}four} (\texttt{Development.\allowbreak{}AutomorphicArithmetic.\allowbreak{}AutomorphicVarieties.\allowbreak{}ModularityArgument.\allowbreak{}ThresholdArguments.\allowbreak{}DenseThresholdCases.\allowbreak{}OffDiagonalNilpotentRankFour}, lines 10--100)\\
\texttt{UpperTheory.\allowbreak{}V.\allowbreak{}Ws.\allowbreak{}fib} (\texttt{Development.\allowbreak{}AutomorphicArithmetic.\allowbreak{}AutomorphicVarieties.\allowbreak{}WeakShimuraTheory.\allowbreak{}CompactIntegralFields.\allowbreak{}FibLow}, lines 68--111)\\
\texttt{UpperTheory.\allowbreak{}V.\allowbreak{}Cu.\allowbreak{}fibY\_\allowbreak{}of} (\texttt{Development.\allowbreak{}AutomorphicArithmetic.\allowbreak{}AutomorphicVarieties.\allowbreak{}ModularityArgument.\allowbreak{}CuspidalTwists.\allowbreak{}PFibY}, lines 67--109)\par}

{\footnotesize\raggedright\noindent \ref{B6.48ba}.\quad \texttt{UpperTheory.\allowbreak{}V.\allowbreak{}Sd.\allowbreak{}propE} (\texttt{Development.\allowbreak{}AutomorphicArithmetic.\allowbreak{}AutomorphicVarieties.\allowbreak{}IdeleGenerators.\allowbreak{}HolomorphicIrreducibility.\allowbreak{}SumDivHecke}, lines 153--246)\\
\texttt{UpperTheory.\allowbreak{}V.\allowbreak{}Z.\allowbreak{}Sym.\allowbreak{}Unl.\allowbreak{}propEC} (\texttt{Development.\allowbreak{}AutomorphicArithmetic.\allowbreak{}AutomorphicVarieties.\allowbreak{}ModularityArgument.\allowbreak{}SymmetricPowerTransfer.\allowbreak{}UnlinkedHalfIntegral.\allowbreak{}UnlinkedHeckeRoots.\allowbreak{}PropEC}, lines 191--285)\par}

{\footnotesize\raggedright\noindent \ref{B6.49a}.\quad \texttt{UpperTheory.\allowbreak{}V.\allowbreak{}Pr.\allowbreak{}FinCrit} (\texttt{Development.\allowbreak{}AutomorphicArithmetic.\allowbreak{}AutomorphicVarieties.\allowbreak{}CurveDVRSpecialization.\allowbreak{}IsUnitOrIsUnitOfAddOne}, lines 114--122)\\
\texttt{UpperTheory.\allowbreak{}V.\allowbreak{}Pr.\allowbreak{}finCrit} (\texttt{Development.\allowbreak{}AutomorphicArithmetic.\allowbreak{}AutomorphicVarieties.\allowbreak{}CurveDVRSpecialization.\allowbreak{}FinCrit}, lines 122--186)\par}

{\footnotesize\raggedright\noindent \ref{B6.49b}.\quad \texttt{UpperTheory.\allowbreak{}Esk.\allowbreak{}pf\_\allowbreak{}TrCong} (\texttt{PositivityTheory.\allowbreak{}AutomorphicForms.\allowbreak{}PublishedInputs.\allowbreak{}JacquetLanglands.\allowbreak{}TorusKernels.\allowbreak{}PfTrCong}, lines 101--128)\\
\texttt{UpperTheory.\allowbreak{}Esk.\allowbreak{}pf\_\allowbreak{}HolWC} (\texttt{PositivityTheory.\allowbreak{}AutomorphicForms.\allowbreak{}PublishedInputs.\allowbreak{}JacquetLanglands.\allowbreak{}TorusKernels.\allowbreak{}PfHolWC}, lines 174--305)\par}

{\footnotesize\raggedright\noindent \ref{B6.49c}.\quad \texttt{UpperTheory.\allowbreak{}V.\allowbreak{}Ws.\allowbreak{}eptV} (\texttt{Development.\allowbreak{}AutomorphicArithmetic.\allowbreak{}AutomorphicVarieties.\allowbreak{}WeakShimuraTheory.\allowbreak{}CompactIntegralFields.\allowbreak{}IsDvrCoef}, lines 25--55)\\
\texttt{UpperTheory.\allowbreak{}V.\allowbreak{}Ws.\allowbreak{}ptPair} (\texttt{Development.\allowbreak{}AutomorphicArithmetic.\allowbreak{}AutomorphicVarieties.\allowbreak{}WeakShimuraTheory.\allowbreak{}CompactIntegralFields.\allowbreak{}PtPair}, lines 22--75)\\
\texttt{UpperTheory.\allowbreak{}V.\allowbreak{}Ws.\allowbreak{}labels} (\texttt{Development.\allowbreak{}AutomorphicArithmetic.\allowbreak{}AutomorphicVarieties.\allowbreak{}WeakShimuraTheory.\allowbreak{}CompactIntegralFields.\allowbreak{}Labels}, lines 22--98)\par}

{\footnotesize\raggedright\noindent \ref{B6.49d}.\quad \texttt{UpperTheory.\allowbreak{}V.\allowbreak{}Ws.\allowbreak{}lat} (\texttt{Development.\allowbreak{}AutomorphicArithmetic.\allowbreak{}AutomorphicVarieties.\allowbreak{}WeakShimuraTheory.\allowbreak{}CompactIntegralFields.\allowbreak{}LatSix}, lines 24--61)\par}

{\footnotesize\raggedright\noindent \ref{B6.49e}.\quad \texttt{UpperTheory.\allowbreak{}V.\allowbreak{}Ws.\allowbreak{}restrQ} (\texttt{Development.\allowbreak{}AutomorphicArithmetic.\allowbreak{}AutomorphicVarieties.\allowbreak{}WeakShimuraTheory.\allowbreak{}CompactIntegralFields.\allowbreak{}RestrQ}, lines 67--118)\\
\texttt{UpperTheory.\allowbreak{}V.\allowbreak{}Ws.\allowbreak{}goodQ} (\texttt{Development.\allowbreak{}AutomorphicArithmetic.\allowbreak{}AutomorphicVarieties.\allowbreak{}WeakShimuraTheory.\allowbreak{}CompactIntegralFields.\allowbreak{}FlagFacts}, lines 64--141)\par}

{\footnotesize\raggedright\noindent \ref{B6.49f}.\quad \texttt{UpperTheory.\allowbreak{}V.\allowbreak{}Ws.\allowbreak{}bl} (\texttt{Development.\allowbreak{}AutomorphicArithmetic.\allowbreak{}AutomorphicVarieties.\allowbreak{}WeakShimuraTheory.\allowbreak{}CompactIntegralFields.\allowbreak{}NotOver}, lines 51--128)\par}

{\footnotesize\raggedright\noindent \ref{B6.49g}.\quad \texttt{UpperTheory.\allowbreak{}V.\allowbreak{}AlgebraTools.\allowbreak{}idCompMat} (\texttt{Development.\allowbreak{}AutomorphicArithmetic.\allowbreak{}AutomorphicVarieties.\allowbreak{}AlgebraTools.\allowbreak{}MonomialRepresentations.\allowbreak{}IdCompMat}, lines 46--136)\\
\texttt{UpperTheory.\allowbreak{}V.\allowbreak{}AlgebraTools.\allowbreak{}regCosetFour} (\texttt{Development.\allowbreak{}AutomorphicArithmetic.\allowbreak{}AutomorphicVarieties.\allowbreak{}AlgebraTools.\allowbreak{}MonomialRepresentations.\allowbreak{}RegCosetFour}, lines 337--394)\\
\texttt{UpperTheory.\allowbreak{}V.\allowbreak{}Pr.\allowbreak{}regCoset} (\texttt{Development.\allowbreak{}AutomorphicArithmetic.\allowbreak{}AutomorphicVarieties.\allowbreak{}CurveDVRSpecialization.\allowbreak{}RegCoset}, lines 143--161)\par}

{\footnotesize\raggedright\noindent \ref{B6.49h}.\quad \texttt{UpperTheory.\allowbreak{}V.\allowbreak{}Af.\allowbreak{}newformRepOf} (\texttt{Development.\allowbreak{}AutomorphicArithmetic.\allowbreak{}AutomorphicVarieties.\allowbreak{}AutomorphicForms.\allowbreak{}TemperedLocalComponents.\allowbreak{}NewformRepOf}, lines 114--139)\\
\texttt{UpperTheory.\allowbreak{}Gal.\allowbreak{}SV.\allowbreak{}galois\_\allowbreak{}rep\_\allowbreak{}exists\_\allowbreak{}totally\_\allowbreak{}real\_\allowbreak{}rank\_\allowbreak{}two\_\allowbreak{}of\_\allowbreak{}base\_\allowbreak{}change} (\texttt{Development.\allowbreak{}AutomorphicArithmetic.\allowbreak{}AutomorphicVarieties.\allowbreak{}ModularityArgument.\allowbreak{}ThresholdArguments.\allowbreak{}DenseThresholdCases.\allowbreak{}SVPatch}, lines 103--165)\\
\texttt{UpperTheory.\allowbreak{}V.\allowbreak{}Pr.\allowbreak{}kaOf} (\texttt{Development.\allowbreak{}AutomorphicArithmetic.\allowbreak{}AutomorphicVarieties.\allowbreak{}CurveDVRSpecialization.\allowbreak{}ValIntLe}, lines 102--140)\par}

{\footnotesize\raggedright\noindent \ref{B6.49i}.\quad \texttt{UpperTheory.\allowbreak{}Gal.\allowbreak{}projectedFormsSystemPart.\allowbreak{}cRealAv} (\texttt{Development.\allowbreak{}AutomorphicArithmetic.\allowbreak{}Foundations.\allowbreak{}GaloisRepresentations.\allowbreak{}PlainTorsionPairs.\allowbreak{}CRealAv}, lines 150--333)\\
\texttt{UpperTheory.\allowbreak{}V.\allowbreak{}Pr.\allowbreak{}p1Of} (\texttt{Development.\allowbreak{}AutomorphicArithmetic.\allowbreak{}AutomorphicVarieties.\allowbreak{}CurveDVRSpecialization.\allowbreak{}FirstPrimeConditionTransfer}, lines 5--20)\par}

{\footnotesize\raggedright\noindent \ref{B6.49}.\quad \texttt{UpperTheory.\allowbreak{}V.\allowbreak{}HV.\allowbreak{}DescentVanishingReductions.\allowbreak{}sp4In0\_\allowbreak{}alt} (\texttt{Development.\allowbreak{}AutomorphicArithmetic.\allowbreak{}AutomorphicVarieties.\allowbreak{}ModularityArgument.\allowbreak{}ThresholdArguments.\allowbreak{}DenseThresholdCases.\allowbreak{}Sp4In0}, lines 18--19)\\
\texttt{UpperTheory.\allowbreak{}V.\allowbreak{}Aut.\allowbreak{}SymplecticFourImageDense} (\texttt{Development.\allowbreak{}AutomorphicArithmetic.\allowbreak{}AutomorphicVarieties.\allowbreak{}ModularityArgument.\allowbreak{}ThresholdArguments.\allowbreak{}DenseThresholdCases.\allowbreak{}IsAut4}, lines 14--17)\par}

{\footnotesize\raggedright\noindent \ref{B6.50}.\quad \texttt{UpperTheory.\allowbreak{}GalN.\allowbreak{}V.\allowbreak{}numberSp4\_\allowbreak{}n\_\allowbreak{}alt} (\texttt{Development.\allowbreak{}AutomorphicArithmetic.\allowbreak{}AutomorphicVarieties.\allowbreak{}ModularityArgument.\allowbreak{}ThresholdArguments.\allowbreak{}DenseThresholdCases.\allowbreak{}VWLevel}, lines 18--20)\\
\texttt{UpperTheory.\allowbreak{}V.\allowbreak{}Af.\allowbreak{}NumberSp4Alt} (\texttt{Development.\allowbreak{}AutomorphicArithmetic.\allowbreak{}AutomorphicVarieties.\allowbreak{}ModularityArgument.\allowbreak{}ThresholdArguments.\allowbreak{}DenseThresholdCases.\allowbreak{}AfStmt}, lines 34--43)\par}

{\footnotesize\raggedright\noindent \ref{B6.51fin}.\quad \texttt{UpperTheory.\allowbreak{}V.\allowbreak{}Pr.\allowbreak{}FinCrit} (\texttt{Development.\allowbreak{}AutomorphicArithmetic.\allowbreak{}AutomorphicVarieties.\allowbreak{}CurveDVRSpecialization.\allowbreak{}IsUnitOrIsUnitOfAddOne}, lines 114--122)\\
\texttt{UpperTheory.\allowbreak{}V.\allowbreak{}Pr.\allowbreak{}finCrit} (\texttt{Development.\allowbreak{}AutomorphicArithmetic.\allowbreak{}AutomorphicVarieties.\allowbreak{}CurveDVRSpecialization.\allowbreak{}FinCrit}, lines 122--186)\par}

{\footnotesize\raggedright\noindent \ref{B6.51}.\quad \texttt{UpperTheory.\allowbreak{}V.\allowbreak{}Aut.\allowbreak{}ThmS6Core\_\allowbreak{}s} (\texttt{Development.\allowbreak{}AutomorphicArithmetic.\allowbreak{}AutomorphicVarieties.\allowbreak{}ModularityArgument.\allowbreak{}ThresholdArguments.\allowbreak{}DenseThresholdCases.\allowbreak{}ThmS6CoreS}, lines 27--36)\\
\texttt{UpperTheory.\allowbreak{}V.\allowbreak{}Pr.\allowbreak{}symmetric\_\allowbreak{}cube\_\allowbreak{}theorem\_\allowbreak{}of\_\allowbreak{}endomorphisms\_\allowbreak{}alt} (\texttt{Development.\allowbreak{}AutomorphicArithmetic.\allowbreak{}AutomorphicVarieties.\allowbreak{}ModularityArgument.\allowbreak{}ThresholdArguments.\allowbreak{}DenseThresholdCases.\allowbreak{}SymmetricCubeTheoremOfEndomorphisms}, lines 76--131)\par}

{\footnotesize\raggedright\noindent \ref{B6.51frame}.\quad \texttt{UpperTheory.\allowbreak{}V.\allowbreak{}Fr.\allowbreak{}Frame.\allowbreak{}Ok} (\texttt{Development.\allowbreak{}AutomorphicArithmetic.\allowbreak{}AutomorphicVarieties.\allowbreak{}FrobeniusSmoothVectors.\allowbreak{}ContinuousCoboundaries2}, lines 176--210)\\
\texttt{UpperTheory.\allowbreak{}V.\allowbreak{}Fr.\allowbreak{}H1C} (\texttt{Development.\allowbreak{}AutomorphicArithmetic.\allowbreak{}AutomorphicVarieties.\allowbreak{}FrobeniusSmoothVectors.\allowbreak{}ContinuousCoboundaries2}, lines 57--58)\\
\texttt{UpperTheory.\allowbreak{}V.\allowbreak{}Fr.\allowbreak{}Frame.\allowbreak{}EC} (\texttt{Development.\allowbreak{}AutomorphicArithmetic.\allowbreak{}AutomorphicVarieties.\allowbreak{}FrobeniusSmoothVectors.\allowbreak{}ContinuousCoboundaries2}, lines 167--175)\\
\texttt{UpperTheory.\allowbreak{}V.\allowbreak{}Pr.\allowbreak{}SeedClause} (\texttt{Development.\allowbreak{}AutomorphicArithmetic.\allowbreak{}AutomorphicVarieties.\allowbreak{}CurveDVRSpecialization.\allowbreak{}HostFrameEx}, lines 58--62)\\
\texttt{UpperTheory.\allowbreak{}V.\allowbreak{}Sd.\allowbreak{}FieldPairCharacterTripleConditions} (\texttt{Development.\allowbreak{}AutomorphicArithmetic.\allowbreak{}AutomorphicVarieties.\allowbreak{}IdeleGenerators.\allowbreak{}SlopeTestRemainingChoices}, lines 136--150)\par}

{\footnotesize\raggedright\noindent \ref{B6.51framea}.\quad \texttt{UpperTheory.\allowbreak{}V.\allowbreak{}GC.\allowbreak{}ec} (\texttt{Development.\allowbreak{}AutomorphicArithmetic.\allowbreak{}AutomorphicVarieties.\allowbreak{}GaloisCochains.\allowbreak{}ResH1SurjectiveOfIsomorphism}, lines 84--179)\par}

{\footnotesize\raggedright\noindent \ref{B6.51aa}.\quad \texttt{UpperTheory.\allowbreak{}V.\allowbreak{}Fr.\allowbreak{}detsec} (\texttt{Development.\allowbreak{}AutomorphicArithmetic.\allowbreak{}AutomorphicVarieties.\allowbreak{}FrobeniusSmoothVectors.\allowbreak{}Detsec}, lines 175--192)\par}

{\footnotesize\raggedright\noindent \ref{B6.51ab}.\quad \texttt{UpperTheory.\allowbreak{}V.\allowbreak{}Fr.\allowbreak{}epres} (\texttt{Development.\allowbreak{}AutomorphicArithmetic.\allowbreak{}AutomorphicVarieties.\allowbreak{}FrobeniusSmoothVectors.\allowbreak{}SumLe}, lines 42--116)\\
\texttt{UpperTheory.\allowbreak{}V.\allowbreak{}Fr.\allowbreak{}ep\_\allowbreak{}pad} (\texttt{Development.\allowbreak{}AutomorphicArithmetic.\allowbreak{}AutomorphicVarieties.\allowbreak{}FrobeniusSmoothVectors.\allowbreak{}ExistsPaddedFamily}, lines 26--64)\par}

{\footnotesize\raggedright\noindent \ref{B6.51a}.\quad \texttt{UpperTheory.\allowbreak{}Gal.\allowbreak{}projectedFormsSystemPart.\allowbreak{}linked\_\allowbreak{}endomorphism\_\allowbreak{}dense} (\texttt{Development.\allowbreak{}AutomorphicArithmetic.\allowbreak{}AutomorphicVarieties.\allowbreak{}ModularityArgument.\allowbreak{}ThresholdArguments.\allowbreak{}DenseThresholdCases.\allowbreak{}LinkedEndE}, lines 18--67)\\
\texttt{UpperTheory.\allowbreak{}V.\allowbreak{}Pr.\allowbreak{}absMain} (\texttt{Development.\allowbreak{}AutomorphicArithmetic.\allowbreak{}AutomorphicVarieties.\allowbreak{}CurveDVRSpecialization.\allowbreak{}AbsMain}, lines 12--30)\\
\texttt{UpperTheory.\allowbreak{}Gal.\allowbreak{}projectedFormsSystemPart.\allowbreak{}twoRuns\_\allowbreak{}alt} (\texttt{Development.\allowbreak{}AutomorphicArithmetic.\allowbreak{}AutomorphicVarieties.\allowbreak{}ModularityArgument.\allowbreak{}ThresholdArguments.\allowbreak{}DenseThresholdCases.\allowbreak{}TwoRuns}, lines 80--109)\par}

{\footnotesize\raggedright\noindent \ref{B6.52a}.\quad \texttt{UpperTheory.\allowbreak{}V.\allowbreak{}Z.\allowbreak{}Win.\allowbreak{}Dn.\allowbreak{}autEnd2\_\allowbreak{}alt} (\texttt{Development.\allowbreak{}AutomorphicArithmetic.\allowbreak{}AutomorphicVarieties.\allowbreak{}ModularityArgument.\allowbreak{}ThresholdArguments.\allowbreak{}DenseThresholdCases.\allowbreak{}AutEnd2}, lines 39--168)\\
\texttt{UpperTheory.\allowbreak{}V.\allowbreak{}Z.\allowbreak{}Win.\allowbreak{}RR.\allowbreak{}autNewformA\_\allowbreak{}of\_\allowbreak{}exRS\_\allowbreak{}of\_\allowbreak{}compatT\_\allowbreak{}of\_\allowbreak{}repNewAE} (\texttt{Development.\allowbreak{}AutomorphicArithmetic.\allowbreak{}AutomorphicVarieties.\allowbreak{}ModularityArgument.\allowbreak{}ThresholdArguments.\allowbreak{}DenseThresholdCases.\allowbreak{}AutNewformAOfEx}, lines 173--205)\\
\texttt{UpperTheory.\allowbreak{}V.\allowbreak{}Z.\allowbreak{}Win.\allowbreak{}Dn.\allowbreak{}across2'\_\allowbreak{}alt} (\texttt{Development.\allowbreak{}AutomorphicArithmetic.\allowbreak{}AutomorphicVarieties.\allowbreak{}ModularityArgument.\allowbreak{}ThresholdArguments.\allowbreak{}DenseThresholdCases.\allowbreak{}Across2}, lines 15--21)\par}

{\footnotesize\raggedright\noindent \ref{B6.52}.\quad \texttt{UpperTheory.\allowbreak{}V.\allowbreak{}Z.\allowbreak{}Win.\allowbreak{}Dn.\allowbreak{}ends0\_\allowbreak{}alt} (\texttt{Development.\allowbreak{}AutomorphicArithmetic.\allowbreak{}AutomorphicVarieties.\allowbreak{}ModularityArgument.\allowbreak{}ThresholdArguments.\allowbreak{}DenseThresholdCases.\allowbreak{}DnEndsTop}, lines 15--16)\\
\texttt{UpperTheory.\allowbreak{}V.\allowbreak{}Z.\allowbreak{}Win.\allowbreak{}Dn.\allowbreak{}ends\_\allowbreak{}of\_\allowbreak{}alt} (\texttt{Development.\allowbreak{}AutomorphicArithmetic.\allowbreak{}AutomorphicVarieties.\allowbreak{}ModularityArgument.\allowbreak{}ThresholdArguments.\allowbreak{}DenseThresholdCases.\allowbreak{}DnEnds}, lines 35--39)\\
\texttt{UpperTheory.\allowbreak{}V.\allowbreak{}Pr.\allowbreak{}GdEndAlt} (\texttt{Development.\allowbreak{}AutomorphicArithmetic.\allowbreak{}AutomorphicVarieties.\allowbreak{}ModularityArgument.\allowbreak{}ThresholdArguments.\allowbreak{}DenseThresholdCases.\allowbreak{}AutEnd3AltAlt}, lines 35--40)\par}

{\footnotesize\raggedright\noindent \ref{B6.53}.\quad \texttt{UpperTheory.\allowbreak{}V.\allowbreak{}Th.\allowbreak{}AdZero} (\texttt{Development.\allowbreak{}AutomorphicArithmetic.\allowbreak{}AutomorphicVarieties.\allowbreak{}TheoremsMuLinked.\allowbreak{}OrdZero}, lines 68--73)\\
\texttt{UpperTheory.\allowbreak{}V.\allowbreak{}Th.\allowbreak{}AdjointLiftingTheorem} (\texttt{Development.\allowbreak{}AutomorphicArithmetic.\allowbreak{}AutomorphicVarieties.\allowbreak{}ModularityArgument.\allowbreak{}ThresholdArguments.\allowbreak{}DenseThresholdCases.\allowbreak{}AutomorphicPointGalois}, lines 35--39)\\
\texttt{UpperTheory.\allowbreak{}V.\allowbreak{}HV.\allowbreak{}DescentVanishingReductions.\allowbreak{}adjoint\_\allowbreak{}lifting\_\allowbreak{}of\_\allowbreak{}cohomology} (\texttt{Development.\allowbreak{}AutomorphicArithmetic.\allowbreak{}AutomorphicVarieties.\allowbreak{}ModularityArgument.\allowbreak{}ThresholdArguments.\allowbreak{}DenseThresholdCases.\allowbreak{}AdjointLiftingOfCohomology}, lines 31--31)\par}

{\footnotesize\raggedright\noindent \ref{B6.53a}.\quad \texttt{UpperTheory.\allowbreak{}V.\allowbreak{}AlgebraTools.\allowbreak{}cocycle\_\allowbreak{}trace\_\allowbreak{}zero} (\texttt{Development.\allowbreak{}AutomorphicArithmetic.\allowbreak{}AutomorphicVarieties.\allowbreak{}AlgebraTools.\allowbreak{}CocycleTraceZero}, lines 6--94)\par}

{\footnotesize\raggedright\noindent \ref{B6.54a}.\quad \texttt{UpperTheory.\allowbreak{}V.\allowbreak{}Z.\allowbreak{}Desc.\allowbreak{}W.\allowbreak{}jPT\_\allowbreak{}of\_\allowbreak{}alt} (\texttt{Development.\allowbreak{}AutomorphicArithmetic.\allowbreak{}AutomorphicVarieties.\allowbreak{}ModularityArgument.\allowbreak{}ThresholdArguments.\allowbreak{}DenseThresholdCases.\allowbreak{}DualOfWeights}, lines 243--339)\par}

{\footnotesize\raggedright\noindent \ref{B6.54}.\quad \texttt{UpperTheory.\allowbreak{}V.\allowbreak{}Th.\allowbreak{}OrdVan2Alt} (\texttt{Development.\allowbreak{}AutomorphicArithmetic.\allowbreak{}AutomorphicVarieties.\allowbreak{}ModularityArgument.\allowbreak{}ThresholdArguments.\allowbreak{}DenseThresholdCases.\allowbreak{}OrdVan2}, lines 13--20)\\
\texttt{UpperTheory.\allowbreak{}V.\allowbreak{}Th.\allowbreak{}ordinary\_\allowbreak{}vanishing\_\allowbreak{}two\_\allowbreak{}of\_\allowbreak{}printed\_\allowbreak{}hida} (\texttt{Development.\allowbreak{}AutomorphicArithmetic.\allowbreak{}AutomorphicVarieties.\allowbreak{}ModularityArgument.\allowbreak{}ThresholdArguments.\allowbreak{}DenseThresholdCases.\allowbreak{}TOrdVan}, lines 159--192)\par}

{\footnotesize\raggedright\noindent \ref{B6.55}.\quad \texttt{UpperTheory.\allowbreak{}V.\allowbreak{}Th.\allowbreak{}core\_\allowbreak{}o\_\allowbreak{}of\_\allowbreak{}literal\_\allowbreak{}automorphy\_\allowbreak{}n} (\texttt{Development.\allowbreak{}AutomorphicArithmetic.\allowbreak{}AutomorphicVarieties.\allowbreak{}ModularityArgument.\allowbreak{}ThresholdArguments.\allowbreak{}DenseThresholdCases.\allowbreak{}TCoreO}, lines 129--159)\\
\texttt{UpperTheory.\allowbreak{}V.\allowbreak{}Th.\allowbreak{}OrdinaryMainCase} (\texttt{Development.\allowbreak{}AutomorphicArithmetic.\allowbreak{}AutomorphicVarieties.\allowbreak{}ModularityArgument.\allowbreak{}ThresholdArguments.\allowbreak{}DenseThresholdCases.\allowbreak{}ThDefs}, lines 53--64)\par}

{\footnotesize\raggedright\noindent \ref{B6.56}.\quad \texttt{UpperTheory.\allowbreak{}V.\allowbreak{}Pa.\allowbreak{}arPl} (\texttt{Development.\allowbreak{}AutomorphicArithmetic.\allowbreak{}AutomorphicVarieties.\allowbreak{}PatchingModules.\allowbreak{}AutomorphicDetection.\allowbreak{}ArPl}, lines 13--180)\\
\texttt{UpperTheory.\allowbreak{}V.\allowbreak{}Pa.\allowbreak{}ArPl} (\texttt{Development.\allowbreak{}AutomorphicArithmetic.\allowbreak{}AutomorphicVarieties.\allowbreak{}PatchingModules.\allowbreak{}TaylorWilesObstructionBound}, lines 119--125)\par}

{\footnotesize\raggedright\noindent \ref{B6.55a}.\quad \texttt{UpperTheory.\allowbreak{}V.\allowbreak{}Th.\allowbreak{}b1O} (\texttt{Development.\allowbreak{}AutomorphicArithmetic.\allowbreak{}AutomorphicVarieties.\allowbreak{}TheoremsMuLinked.\allowbreak{}GaloisRepresentationTypes.\allowbreak{}OrdinaryHeckeAlgebraIntegral}, lines 31--54)\\
\texttt{UpperTheory.\allowbreak{}V.\allowbreak{}Th.\allowbreak{}b1absO} (\texttt{Development.\allowbreak{}AutomorphicArithmetic.\allowbreak{}AutomorphicVarieties.\allowbreak{}TheoremsMuLinked.\allowbreak{}CyclicReductionCohomology.\allowbreak{}ResidualRepOne}, lines 44--130)\\
\texttt{UpperTheory.\allowbreak{}V.\allowbreak{}Th.\allowbreak{}oneLevelO} (\texttt{Development.\allowbreak{}AutomorphicArithmetic.\allowbreak{}AutomorphicVarieties.\allowbreak{}TheoremsMuLinked.\allowbreak{}CyclicReductionCohomology.\allowbreak{}OneLevelO}, lines 52--135)\\
\texttt{UpperTheory.\allowbreak{}V.\allowbreak{}Th.\allowbreak{}dualNO} (\texttt{Development.\allowbreak{}AutomorphicArithmetic.\allowbreak{}AutomorphicVarieties.\allowbreak{}TheoremsMuLinked.\allowbreak{}CyclicReductionCohomology.\allowbreak{}DualNO}, lines 50--145)\par}

{\footnotesize\raggedright\noindent \ref{B6.57a}.\quad \texttt{UpperTheory.\allowbreak{}V.\allowbreak{}Th.\allowbreak{}coc} (\texttt{Development.\allowbreak{}AutomorphicArithmetic.\allowbreak{}AutomorphicVarieties.\allowbreak{}TheoremsMuLinked.\allowbreak{}GaloisRepresentationTypes.\allowbreak{}LiftSmul}, lines 39--119)\\
\texttt{UpperTheory.\allowbreak{}V.\allowbreak{}Th.\allowbreak{}trPsi} (\texttt{Development.\allowbreak{}AutomorphicArithmetic.\allowbreak{}AutomorphicVarieties.\allowbreak{}TheoremsMuLinked.\allowbreak{}GaloisRepresentationTypes.\allowbreak{}TpSingle}, lines 135--178)\par}

{\footnotesize\raggedright\noindent \ref{B6.57}.\quad \texttt{UpperTheory.\allowbreak{}V.\allowbreak{}Th.\allowbreak{}count} (\texttt{Development.\allowbreak{}AutomorphicArithmetic.\allowbreak{}AutomorphicVarieties.\allowbreak{}TheoremsMuLinked.\allowbreak{}GaloisRepresentationTypes.\allowbreak{}MapRingHomSingleMulPartialCharpoly}, lines 131--249)\\
\texttt{UpperTheory.\allowbreak{}V.\allowbreak{}Th.\allowbreak{}Count} (\texttt{Development.\allowbreak{}AutomorphicArithmetic.\allowbreak{}AutomorphicVarieties.\allowbreak{}TheoremsMuLinked.\allowbreak{}OrdinaryHeckeAlgebraIntegral}, lines 134--143)\\
\texttt{UpperTheory.\allowbreak{}V.\allowbreak{}Pa.\allowbreak{}Lev.\allowbreak{}OpMod} (\texttt{Development.\allowbreak{}AutomorphicArithmetic.\allowbreak{}AutomorphicVarieties.\allowbreak{}PatchingModules.\allowbreak{}NilpotentLevelLowering.\allowbreak{}LevelModule}, lines 40--74)\par}

{\footnotesize\raggedright\noindent \ref{B6.58c}.\quad \texttt{UpperTheory.\allowbreak{}V.\allowbreak{}Alg.\allowbreak{}M.\allowbreak{}fin\_\allowbreak{}holds} (\texttt{Development.\allowbreak{}AutomorphicArithmetic.\allowbreak{}AutomorphicVarieties.\allowbreak{}CommutativeAlgebra.\allowbreak{}GradedModulesLength.\allowbreak{}FinHolds}, lines 49--99)\\
\texttt{UpperTheory.\allowbreak{}V.\allowbreak{}Alg.\allowbreak{}M.\allowbreak{}rk\_\allowbreak{}holds} (\texttt{Development.\allowbreak{}AutomorphicArithmetic.\allowbreak{}AutomorphicVarieties.\allowbreak{}CommutativeAlgebra.\allowbreak{}GradedModulesLength.\allowbreak{}LinPartApply}, lines 127--201)\\
\texttt{UpperTheory.\allowbreak{}V.\allowbreak{}Alg.\allowbreak{}M.\allowbreak{}sm\_\allowbreak{}holds} (\texttt{Development.\allowbreak{}AutomorphicArithmetic.\allowbreak{}AutomorphicVarieties.\allowbreak{}CommutativeAlgebra.\allowbreak{}GradedModulesLength.\allowbreak{}SmHolds}, lines 55--113)\\
\texttt{UpperTheory.\allowbreak{}V.\allowbreak{}Alg.\allowbreak{}M.\allowbreak{}cut\_\allowbreak{}holds} (\texttt{Development.\allowbreak{}AutomorphicArithmetic.\allowbreak{}AutomorphicVarieties.\allowbreak{}CommutativeAlgebra.\allowbreak{}GradedModulesLength.\allowbreak{}CutHolds}, lines 101--132)\\
\texttt{UpperTheory.\allowbreak{}V.\allowbreak{}Alg.\allowbreak{}M.\allowbreak{}nz\_\allowbreak{}holds} (\texttt{Development.\allowbreak{}AutomorphicArithmetic.\allowbreak{}AutomorphicVarieties.\allowbreak{}CommutativeAlgebra.\allowbreak{}GradedModulesLength.\allowbreak{}NzCoeffMul}, lines 33--90)\par}

{\footnotesize\raggedright\noindent \ref{B6.58}.\quad \texttt{UpperTheory.\allowbreak{}V.\allowbreak{}Alg.\allowbreak{}lemmaP} (\texttt{Development.\allowbreak{}AutomorphicArithmetic.\allowbreak{}AutomorphicVarieties.\allowbreak{}CommutativeAlgebra.\allowbreak{}LemmaP}, lines 256--425)\\
\texttt{UpperTheory.\allowbreak{}V.\allowbreak{}Alg.\allowbreak{}M.\allowbreak{}mprime\_\allowbreak{}holds} (\texttt{Development.\allowbreak{}AutomorphicArithmetic.\allowbreak{}AutomorphicVarieties.\allowbreak{}CommutativeAlgebra.\allowbreak{}GradedModulesLength.\allowbreak{}Mprime}, lines 38--81)\\
\texttt{UpperTheory.\allowbreak{}V.\allowbreak{}Alg.\allowbreak{}M.\allowbreak{}main\_\allowbreak{}holds} (\texttt{Development.\allowbreak{}AutomorphicArithmetic.\allowbreak{}AutomorphicVarieties.\allowbreak{}CommutativeAlgebra.\allowbreak{}GradedModulesLength.\allowbreak{}MainCotSpan}, lines 32--71)\par}

{\footnotesize\raggedright\noindent \ref{B6.58a}.\quad \texttt{UpperTheory.\allowbreak{}V.\allowbreak{}Alg.\allowbreak{}M.\allowbreak{}low\_\allowbreak{}holds} (\texttt{Development.\allowbreak{}AutomorphicArithmetic.\allowbreak{}AutomorphicVarieties.\allowbreak{}CommutativeAlgebra.\allowbreak{}GradedModulesLength.\allowbreak{}AdmissibleMonoIdeal}, lines 191--323)\\
\texttt{UpperTheory.\allowbreak{}V.\allowbreak{}Alg.\allowbreak{}M.\allowbreak{}hf\_\allowbreak{}holds} (\texttt{Development.\allowbreak{}AutomorphicArithmetic.\allowbreak{}AutomorphicVarieties.\allowbreak{}CommutativeAlgebra.\allowbreak{}GradedModulesLength.\allowbreak{}ReesGradedPieceSurjection}, lines 199--271)\\
\texttt{UpperTheory.\allowbreak{}V.\allowbreak{}Alg.\allowbreak{}M.\allowbreak{}num\_\allowbreak{}holds} (\texttt{Development.\allowbreak{}AutomorphicArithmetic.\allowbreak{}AutomorphicVarieties.\allowbreak{}CommutativeAlgebra.\allowbreak{}GradedModulesLength.\allowbreak{}FrobeniusInclusionXEqMulQ}, lines 206--284)\par}

{\footnotesize\raggedright\noindent \ref{B6.58b}.\quad \texttt{UpperTheory.\allowbreak{}V.\allowbreak{}Alg.\allowbreak{}M.\allowbreak{}zii\_\allowbreak{}holds} (\texttt{Development.\allowbreak{}AutomorphicArithmetic.\allowbreak{}AutomorphicVarieties.\allowbreak{}CommutativeAlgebra.\allowbreak{}GradedModulesLength.\allowbreak{}SameDegreeZeroOfRestriction}, lines 15--43)\\
\texttt{UpperTheory.\allowbreak{}V.\allowbreak{}Alg.\allowbreak{}M.\allowbreak{}up\_\allowbreak{}holds} (\texttt{Development.\allowbreak{}AutomorphicArithmetic.\allowbreak{}AutomorphicVarieties.\allowbreak{}CommutativeAlgebra.\allowbreak{}GradedModulesLength.\allowbreak{}TwoSigned}, lines 17--96)\par}

{\footnotesize\raggedright\noindent \ref{B6.59}.\quad \texttt{UpperTheory.\allowbreak{}V.\allowbreak{}X.\allowbreak{}residual\_\allowbreak{}bernoulli\_\allowbreak{}nonvanishing} (\texttt{Development.\allowbreak{}AutomorphicArithmetic.\allowbreak{}AutomorphicVarieties.\allowbreak{}LevelCohomology.\allowbreak{}ResidualBernoulliNonvanishing}, lines 174--264)\\
\texttt{UpperTheory.\allowbreak{}V.\allowbreak{}X.\allowbreak{}corX} (\texttt{Development.\allowbreak{}AutomorphicArithmetic.\allowbreak{}AutomorphicVarieties.\allowbreak{}LevelCohomology.\allowbreak{}ChrInvNeg}, lines 197--211)\par}

{\footnotesize\raggedright\noindent \ref{B6.59a}.\quad \texttt{UpperTheory.\allowbreak{}V.\allowbreak{}X.\allowbreak{}lx\_\allowbreak{}partial} (\texttt{Development.\allowbreak{}AutomorphicArithmetic.\allowbreak{}AutomorphicVarieties.\allowbreak{}LevelCohomology.\allowbreak{}LxPartial}, lines 15--115)\\
\texttt{UpperTheory.\allowbreak{}V.\allowbreak{}X.\allowbreak{}lx\_\allowbreak{}core} (\texttt{Development.\allowbreak{}AutomorphicArithmetic.\allowbreak{}AutomorphicVarieties.\allowbreak{}LevelCohomology.\allowbreak{}LxCore}, lines 65--221)\\
\texttt{UpperTheory.\allowbreak{}V.\allowbreak{}X.\allowbreak{}lemmaX} (\texttt{Development.\allowbreak{}AutomorphicArithmetic.\allowbreak{}AutomorphicVarieties.\allowbreak{}LevelCohomology.\allowbreak{}LemmaX}, lines 78--303)\par}

{\footnotesize\raggedright\noindent \ref{B6.59a}.\quad \texttt{UpperTheory.\allowbreak{}V.\allowbreak{}Z.\allowbreak{}J.\allowbreak{}count\_\allowbreak{}i\_\allowbreak{}of\_\allowbreak{}printed\_\allowbreak{}bound\_\allowbreak{}k\_\allowbreak{}kato\_\allowbreak{}automorphy} (\texttt{Development.\allowbreak{}AutomorphicArithmetic.\allowbreak{}AutomorphicVarieties.\allowbreak{}ModularityArgument.\allowbreak{}ThresholdArguments.\allowbreak{}DenseThresholdCases.\allowbreak{}LinkedUnlE}, lines 49--54)\par}

{\footnotesize\raggedright\noindent \ref{B6.59a}.\quad \texttt{UpperTheory.\allowbreak{}V.\allowbreak{}HV.\allowbreak{}DescentVanishingReductions.\allowbreak{}countII\_\allowbreak{}and\_\allowbreak{}countIII\_\allowbreak{}of\_\allowbreak{}printsH\_\allowbreak{}of\_\allowbreak{}sp4In} (\texttt{Development.\allowbreak{}AutomorphicArithmetic.\allowbreak{}AutomorphicVarieties.\allowbreak{}ModularityArgument.\allowbreak{}ThresholdArguments.\allowbreak{}DenseThresholdCases.\allowbreak{}AdjointLiftingOfCohomology}, lines 33--39)\par}

{\footnotesize\raggedright\noindent \ref{B6.59a}.\quad \texttt{UpperTheory.\allowbreak{}V.\allowbreak{}Z.\allowbreak{}J.\allowbreak{}count\_\allowbreak{}v\_\allowbreak{}of\_\allowbreak{}printed\_\allowbreak{}bound\_\allowbreak{}k\_\allowbreak{}reduced\_\allowbreak{}unitary\_\allowbreak{}lift} (\texttt{Development.\allowbreak{}AutomorphicArithmetic.\allowbreak{}AutomorphicVarieties.\allowbreak{}ModularityArgument.\allowbreak{}ThresholdArguments.\allowbreak{}DenseThresholdCases.\allowbreak{}LinkedUnlE}, lines 55--61)\par}

\section*{Chapter~\ref{ch:B6-6}: The exponent 401 and the downstream window}

{\footnotesize\raggedright\noindent \ref{B6.60}.\quad \texttt{Formalism.\allowbreak{}MainNotions.\allowbreak{}PageFix} (\texttt{Formalism.\allowbreak{}SharpRestriction.\allowbreak{}PageFix}, lines 9--18)\\
\texttt{GlobalTheory.\allowbreak{}EX.\allowbreak{}TOv} (\texttt{Development.\allowbreak{}AbelianVarietyTheory.\allowbreak{}Definitions.\allowbreak{}GlobalTwisting.\allowbreak{}PairRowIffQdRows}, lines 17--17)\\
\texttt{GlobalTheory.\allowbreak{}EX.\allowbreak{}nWv} (\texttt{Development.\allowbreak{}AbelianVarietyTheory.\allowbreak{}Definitions.\allowbreak{}GlobalTwisting.\allowbreak{}PairRowIffQdRows}, lines 19--20)\\
\texttt{GlobalTheory.\allowbreak{}EX.\allowbreak{}RankVal} (\texttt{Development.\allowbreak{}AbelianVarietyTheory.\allowbreak{}Definitions.\allowbreak{}GlobalTwisting.\allowbreak{}PairRowIffQdRows}, lines 28--32)\\
\texttt{UpperTheory.\allowbreak{}Top.\allowbreak{}eventually\_\allowbreak{}rank\_\allowbreak{}le\_\allowbreak{}rpow\_\allowbreak{}of\_\allowbreak{}lt\_\allowbreak{}one\_\allowbreak{}div\_\allowbreak{}fiftytwo} (\texttt{Development.\allowbreak{}AutomorphicArithmetic.\allowbreak{}CharacterBoxes.\allowbreak{}EventuallyRankLeRpowOfLtOneDivFiftytwo}, lines 5--126)\par}

{\footnotesize\raggedright\noindent \ref{B6.62}.\quad \texttt{UpperTheory.\allowbreak{}Top.\allowbreak{}WinDatum} (\texttt{Development.\allowbreak{}AutomorphicArithmetic.\allowbreak{}CharacterBoxes.\allowbreak{}TwistForm}, lines 19--26)\\
\texttt{UpperTheory.\allowbreak{}Top.\allowbreak{}winPairs} (\texttt{Development.\allowbreak{}AutomorphicArithmetic.\allowbreak{}CharacterBoxes.\allowbreak{}WinPairs}, lines 5--81)\par}

\section*{Chapter~\ref{ch:green}: Green's classification and its use}

{\footnotesize\raggedright\noindent \ref{B7.01}.\quad \texttt{GlobalTheory.\allowbreak{}GLq.\allowbreak{}IsCusp} (\texttt{Development.\allowbreak{}AbelianVarietyTheory.\allowbreak{}Definitions.\allowbreak{}NonGaloisDescent.\allowbreak{}LabelDatum}, lines 64--66)\par}

{\footnotesize\raggedright\noindent \ref{B7.02}.\quad \texttt{GlobalTheory.\allowbreak{}GLq.\allowbreak{}IsPrimeLab} (\texttt{Development.\allowbreak{}AbelianVarietyTheory.\allowbreak{}Definitions.\allowbreak{}NonGaloisDescent.\allowbreak{}LabelDatum}, lines 91--92)\\
\texttt{GlobalTheory.\allowbreak{}Published.\allowbreak{}IsRegularChar} (\texttt{Development.\allowbreak{}AbelianVarietyTheory.\allowbreak{}PublishedTheorems.\allowbreak{}GaloisCohomology₂.\allowbreak{}GreenRegularCharacters}, lines 10--12)\par}

{\footnotesize\raggedright\noindent \ref{B7.03}.\quad \texttt{GlobalTheory.\allowbreak{}GLq.\allowbreak{}LabelDatum} (\texttt{Development.\allowbreak{}AbelianVarietyTheory.\allowbreak{}Definitions.\allowbreak{}NonGaloisDescent.\allowbreak{}LabelDatum}, lines 81--87)\\
\texttt{GlobalTheory.\allowbreak{}GLq.\allowbreak{}thetaOf} (\texttt{Development.\allowbreak{}AbelianVarietyTheory.\allowbreak{}Definitions.\allowbreak{}NonGaloisDescent.\allowbreak{}LabelDatum}, lines 88--90)\par}

{\footnotesize\raggedright\noindent \ref{B7.04}.\quad \texttt{GlobalTheory.\allowbreak{}GLq.\allowbreak{}IsCuspOf} (\texttt{Development.\allowbreak{}AbelianVarietyTheory.\allowbreak{}Definitions.\allowbreak{}NonGaloisDescent.\allowbreak{}LabelDatum}, lines 93--97)\par}

{\footnotesize\raggedright\noindent \ref{B7.05}.\quad \texttt{GlobalTheory.\allowbreak{}Published.\allowbreak{}GreenRegularCharacters} (\texttt{Development.\allowbreak{}AbelianVarietyTheory.\allowbreak{}PublishedTheorems.\allowbreak{}GaloisCohomology₂.\allowbreak{}GreenRegularCharacters}, lines 13--21)\par}

{\footnotesize\raggedright\noindent \ref{B7.06}.\quad \texttt{GlobalTheory.\allowbreak{}GLq.\allowbreak{}Green} (\texttt{Development.\allowbreak{}AbelianVarietyTheory.\allowbreak{}Definitions.\allowbreak{}NonGaloisDescent.\allowbreak{}LabelDatum}, lines 108--113)\par}

{\footnotesize\raggedright\noindent \ref{B7.07}.\quad \texttt{GlobalTheory.\allowbreak{}GreenExists.\allowbreak{}pow\_\allowbreak{}eq\_\allowbreak{}self\_\allowbreak{}iff} (\texttt{Development.\allowbreak{}AbelianVarietyTheory.\allowbreak{}Proofs.\allowbreak{}InertiaEigenvalues.\allowbreak{}CuspidalCharExistsOfGreen}, lines 60--67)\\
\texttt{GlobalTheory.\allowbreak{}GreenExists.\allowbreak{}rootsBar\_\allowbreak{}eq} (\texttt{Development.\allowbreak{}AbelianVarietyTheory.\allowbreak{}Proofs.\allowbreak{}InertiaEigenvalues.\allowbreak{}CuspidalCharExistsOfGreen}, lines 68--95)\par}

{\footnotesize\raggedright\noindent \ref{B7.08}.\quad \texttt{GlobalTheory.\allowbreak{}CuspGL.\allowbreak{}TeichEmb.\allowbreak{}exists\_\allowbreak{}teichmuller} (\texttt{Development.\allowbreak{}AbelianVarietyTheory.\allowbreak{}Proofs.\allowbreak{}CharacterTheory.\allowbreak{}ExistsTeichmullerMapAlgebraicClosureFiniteField}, lines 24--77)\par}

{\footnotesize\raggedright\noindent \ref{B7.09}.\quad \texttt{GlobalTheory.\allowbreak{}CuspGL.\allowbreak{}rootSum} (\texttt{Development.\allowbreak{}AbelianVarietyTheory.\allowbreak{}Definitions.\allowbreak{}CommutativeAlgebra.\allowbreak{}InvariantSubspaceMobiusInverseRecEqOfFinrankLt}, lines 65--67)\\
\texttt{GlobalTheory.\allowbreak{}CuspGL.\allowbreak{}rootSum\_\allowbreak{}isAdditive\_\allowbreak{}and\_\allowbreak{}pow\_\allowbreak{}X\_\allowbreak{}sub\_\allowbreak{}one} (\texttt{Development.\allowbreak{}AbelianVarietyTheory.\allowbreak{}Proofs.\allowbreak{}CharacterTheory.\allowbreak{}RootPowerSumCharacterAdditiveAndXSubOnePow}, lines 9--21)\par}

{\footnotesize\raggedright\noindent \ref{B7.10}.\quad \texttt{GlobalTheory.\allowbreak{}CuspGL.\allowbreak{}psi} (\texttt{Development.\allowbreak{}AbelianVarietyTheory.\allowbreak{}Definitions.\allowbreak{}CommutativeAlgebra.\allowbreak{}InvariantSubspaceMobiusInverseRecEqOfFinrankLt}, lines 35--38)\\
\texttt{GlobalTheory.\allowbreak{}CuspGL.\allowbreak{}psi\_\allowbreak{}eq} (\texttt{Development.\allowbreak{}AbelianVarietyTheory.\allowbreak{}Definitions.\allowbreak{}CommutativeAlgebra.\allowbreak{}InvariantSubspaceMobiusInverseRecEqOfFinrankLt}, lines 55--64)\par}

{\footnotesize\raggedright\noindent \ref{B7.11}.\quad \texttt{GlobalTheory.\allowbreak{}CuspGL.\allowbreak{}psi\_\allowbreak{}conj} (\texttt{Development.\allowbreak{}AbelianVarietyTheory.\allowbreak{}Proofs.\allowbreak{}CharacterTheory.\allowbreak{}InvariantSubspaceMobiusInverseCharpolyLinearEqui}, lines 43--52)\\
\texttt{GlobalTheory.\allowbreak{}CuspGL.\allowbreak{}isPrimitive\_\allowbreak{}psi} (\texttt{Development.\allowbreak{}AbelianVarietyTheory.\allowbreak{}Proofs.\allowbreak{}CharacterTheory.\allowbreak{}IsPrimitiveEndClassFunctionInvariantSubspaceMobi}, lines 188--200)\par}

{\footnotesize\raggedright\noindent \ref{B7.12}.\quad \texttt{GlobalTheory.\allowbreak{}CuspGL.\allowbreak{}psi\_\allowbreak{}one} (\texttt{Development.\allowbreak{}AbelianVarietyTheory.\allowbreak{}Proofs.\allowbreak{}CharacterTheory.\allowbreak{}InvariantSubspaceMobiusInverseOneProdIco}, lines 74--78)\\
\texttt{GlobalTheory.\allowbreak{}CuspGL.\allowbreak{}psi\_\allowbreak{}of\_\allowbreak{}irreducible\_\allowbreak{}charpoly} (\texttt{Development.\allowbreak{}AbelianVarietyTheory.\allowbreak{}Proofs.\allowbreak{}CharacterTheory.\allowbreak{}InvariantSubspaceMobiusInverseCharpolyOfIrreduci}, lines 52--57)\par}

{\footnotesize\raggedright\noindent \ref{B7.13}.\quad \texttt{GlobalTheory.\allowbreak{}CuspGL.\allowbreak{}rootSum\_\allowbreak{}charpoly\_\allowbreak{}mem\_\allowbreak{}virt} (\texttt{Development.\allowbreak{}AbelianVarietyTheory.\allowbreak{}Proofs.\allowbreak{}CharacterTheory.\allowbreak{}RootPowerSumCharpolyMemVirtualCharactersGL}, lines 213--234)\\
\texttt{GlobalTheory.\allowbreak{}CuspGL.\allowbreak{}psi\_\allowbreak{}rootSum\_\allowbreak{}charpoly\_\allowbreak{}mem\_\allowbreak{}virt} (\texttt{Development.\allowbreak{}AbelianVarietyTheory.\allowbreak{}Proofs.\allowbreak{}CharacterTheory.\allowbreak{}IsCharacterCuspidalRootSumRegular}, lines 55--89)\par}

{\footnotesize\raggedright\noindent \ref{B7.14}.\quad \texttt{GlobalTheory.\allowbreak{}CuspGL.\allowbreak{}mass\_\allowbreak{}psi\_\allowbreak{}eq\_\allowbreak{}glCard\_\allowbreak{}mul\_\allowbreak{}mangoldt\_\allowbreak{}div} (\texttt{Development.\allowbreak{}AbelianVarietyTheory.\allowbreak{}Proofs.\allowbreak{}CharacterTheory.\allowbreak{}EndomorphismMassMobiusInverseMonicEq}, lines 27--53)\par}

{\footnotesize\raggedright\noindent \ref{B7.15}.\quad \texttt{GlobalTheory.\allowbreak{}CuspGL.\allowbreak{}ipGL\_\allowbreak{}rootSum\_\allowbreak{}eq\_\allowbreak{}sum\_\allowbreak{}torus} (\texttt{Development.\allowbreak{}AbelianVarietyTheory.\allowbreak{}Proofs.\allowbreak{}CharacterTheory.\allowbreak{}InvertibleEndInnerProductRootPowerSumEqTorusSum}, lines 12--25)\\
\texttt{GlobalTheory.\allowbreak{}CuspGL.\allowbreak{}torus\_\allowbreak{}orthogonality} (\texttt{Development.\allowbreak{}AbelianVarietyTheory.\allowbreak{}Proofs.\allowbreak{}CharacterTheory.\allowbreak{}FrobeniusTorusCharacterOrthogonality}, lines 66--93)\par}

{\footnotesize\raggedright\noindent \ref{B7.16}.\quad \texttt{GlobalTheory.\allowbreak{}CuspGL.\allowbreak{}ipGL\_\allowbreak{}psi\_\allowbreak{}rootSum} (\texttt{Development.\allowbreak{}AbelianVarietyTheory.\allowbreak{}Proofs.\allowbreak{}CharacterTheory.\allowbreak{}InvertibleEndInnerProductRootPowerSumMobius}, lines 30--97)\par}

{\footnotesize\raggedright\noindent \ref{B7.17}.\quad \texttt{GlobalTheory.\allowbreak{}CuspGL.\allowbreak{}expand\_\allowbreak{}primitive} (\texttt{Development.\allowbreak{}AbelianVarietyTheory.\allowbreak{}Proofs.\allowbreak{}CharacterTheory.\allowbreak{}PrimitiveClassFunctionExpansionRootSums}, lines 156--221)\par}

{\footnotesize\raggedright\noindent \ref{B7.18}.\quad \texttt{GlobalTheory.\allowbreak{}CuspGL.\allowbreak{}PrimOrth.\allowbreak{}ipGL\_\allowbreak{}psi\_\allowbreak{}eq} (\texttt{Development.\allowbreak{}AbelianVarietyTheory.\allowbreak{}Proofs.\allowbreak{}CharacterTheory.\allowbreak{}InvertibleEndomorphismInnerProductPrimitiveMobiu}, lines 92--162)\par}

{\footnotesize\raggedright\noindent \ref{B7.19}.\quad \texttt{GlobalTheory.\allowbreak{}CuspGL.\allowbreak{}isChar\_\allowbreak{}psi\_\allowbreak{}rootSum\_\allowbreak{}of\_\allowbreak{}regular} (\texttt{Development.\allowbreak{}AbelianVarietyTheory.\allowbreak{}Proofs.\allowbreak{}CharacterTheory.\allowbreak{}IsCharacterCuspidalRootSumRegular}, lines 92--185)\par}

{\footnotesize\raggedright\noindent \ref{B7.20}.\quad \texttt{GlobalTheory.\allowbreak{}ParamLabel.\allowbreak{}isCusp\_\allowbreak{}psi} (\texttt{Development.\allowbreak{}AbelianVarietyTheory.\allowbreak{}Proofs.\allowbreak{}DivisorProofs.\allowbreak{}GreenCuspidalCharIsCuspidalOfMinpolyOfParam}, lines 18--49)\par}

{\footnotesize\raggedright\noindent \ref{B7.21}.\quad \texttt{GlobalTheory.\allowbreak{}GreenRegular.\allowbreak{}exists\_\allowbreak{}cusp\_\allowbreak{}of\_\allowbreak{}regular} (\texttt{Development.\allowbreak{}AbelianVarietyTheory.\allowbreak{}Proofs.\allowbreak{}InertiaEigenvalues.\allowbreak{}ExistsCuspidalCharacterOfRegularCharacterGLn}, lines 65--226)\par}

{\footnotesize\raggedright\noindent \ref{B7.22}.\quad \texttt{GlobalTheory.\allowbreak{}PrimeLabOrbit.\allowbreak{}primeLab\_\allowbreak{}to\_\allowbreak{}regularChar} (\texttt{Development.\allowbreak{}AbelianVarietyTheory.\allowbreak{}Proofs.\allowbreak{}InertiaEigenvalues.\allowbreak{}MonicIrreducibleInjectRegularCharacterOrbits}, lines 125--152)\par}

{\footnotesize\raggedright\noindent \ref{B7.23}.\quad \texttt{GlobalTheory.\allowbreak{}GreenRegularCharactersInertiaEigenvalues.\allowbreak{}family} (\texttt{Development.\allowbreak{}AbelianVarietyTheory.\allowbreak{}Proofs.\allowbreak{}InertiaEigenvalues.\allowbreak{}GreenCuspidalCharsByRegularCharsProof}, lines 23--34)\par}

{\footnotesize\raggedright\noindent \ref{B7.24}.\quad \texttt{GlobalTheory.\allowbreak{}GLq.\allowbreak{}GenericConstituent.\allowbreak{}exists\_\allowbreak{}generic\_\allowbreak{}constituent} (\texttt{Development.\allowbreak{}AbelianVarietyTheory.\allowbreak{}Proofs.\allowbreak{}InertiaEigenvalues.\allowbreak{}ExistsGenericConstituentHarishChandraCuspidal}, lines 98--126)\par}

{\footnotesize\raggedright\noindent \ref{B7.25}.\quad \texttt{GlobalTheory.\allowbreak{}GLq.\allowbreak{}IpCuspProd.\allowbreak{}ip\_\allowbreak{}hcL\_\allowbreak{}cusp} (\texttt{Development.\allowbreak{}AbelianVarietyTheory.\allowbreak{}Proofs.\allowbreak{}InertiaEigenvalues.\allowbreak{}CharacterPairingHarishChandraCuspidalMackey}, lines 259--264)\par}

{\footnotesize\raggedright\noindent \ref{B7.26}.\quad \texttt{GlobalTheory.\allowbreak{}GLq.\allowbreak{}GGNorm.\allowbreak{}gg\_\allowbreak{}norm} (\texttt{Development.\allowbreak{}AbelianVarietyTheory.\allowbreak{}Proofs.\allowbreak{}InertiaEigenvalues.\allowbreak{}GelfandGraevCharacterGLnNormAndReciprocity}, lines 88--92)\par}

{\footnotesize\raggedright\noindent \ref{B7.27}.\quad \texttt{GlobalTheory.\allowbreak{}GLq.\allowbreak{}card\_\allowbreak{}generic\_\allowbreak{}irr\_\allowbreak{}le} (\texttt{Development.\allowbreak{}AbelianVarietyTheory.\allowbreak{}Proofs.\allowbreak{}InertiaEigenvalues.\allowbreak{}NcardGenericIrreducibleCharactersGLnLe}, lines 49--89)\par}

{\footnotesize\raggedright\noindent \ref{B7.28}.\quad \texttt{GlobalTheory.\allowbreak{}GLq.\allowbreak{}CuspMultisetCount.\allowbreak{}card\_\allowbreak{}le} (\texttt{Development.\allowbreak{}AbelianVarietyTheory.\allowbreak{}Proofs.\allowbreak{}InertiaEigenvalues.\allowbreak{}CardCuspidalMultisetsTotalDegreeLe}, lines 30--62)\par}

{\footnotesize\raggedright\noindent \ref{B7.29}.\quad \texttt{GlobalTheory.\allowbreak{}GrPolys.\allowbreak{}card\_\allowbreak{}monic\_\allowbreak{}coeff\_\allowbreak{}zero\_\allowbreak{}ne\_\allowbreak{}zero} (\texttt{Development.\allowbreak{}AbelianVarietyTheory.\allowbreak{}Proofs.\allowbreak{}InertiaEigenvalues.\allowbreak{}NcardMonicNatDegreeCoeffZeroNeZero}, lines 38--63)\par}

{\footnotesize\raggedright\noindent \ref{B7.30}.\quad \texttt{GlobalTheory.\allowbreak{}GreenRegularCharactersInertiaEigenvalues.\allowbreak{}count} (\texttt{Development.\allowbreak{}AbelianVarietyTheory.\allowbreak{}Proofs.\allowbreak{}InertiaEigenvalues.\allowbreak{}GreenCuspidalCharsByRegularCharsProof}, lines 56--132)\par}

{\footnotesize\raggedright\noindent \ref{B7.31}.\quad \texttt{GlobalTheory.\allowbreak{}greenRegularCharacters\_\allowbreak{}holds} (\texttt{Development.\allowbreak{}AbelianVarietyTheory.\allowbreak{}Proofs.\allowbreak{}InertiaEigenvalues.\allowbreak{}GreenCuspidalCharsByRegularCharsProof}, lines 137--169)\par}

{\footnotesize\raggedright\noindent \ref{B7.32}.\quad \texttt{GlobalTheory.\allowbreak{}GreenBridge.\allowbreak{}exists\_\allowbreak{}emb\_\allowbreak{}eq} (\texttt{Development.\allowbreak{}AbelianVarietyTheory.\allowbreak{}Proofs.\allowbreak{}InertiaEigenvalues.\allowbreak{}SumRangeEqRoots}, lines 32--40)\\
\texttt{GlobalTheory.\allowbreak{}GreenBridge.\allowbreak{}exists\_\allowbreak{}unit\_\allowbreak{}gen} (\texttt{Development.\allowbreak{}AbelianVarietyTheory.\allowbreak{}Proofs.\allowbreak{}InertiaEigenvalues.\allowbreak{}SumRangeEqRoots}, lines 53--69)\par}

{\footnotesize\raggedright\noindent \ref{B7.33}.\quad \texttt{GlobalTheory.\allowbreak{}GreenBridge.\allowbreak{}th\_\allowbreak{}apply} (\texttt{Development.\allowbreak{}AbelianVarietyTheory.\allowbreak{}Proofs.\allowbreak{}InertiaEigenvalues.\allowbreak{}SumRangeEqRoots}, lines 122--125)\\
\texttt{GlobalTheory.\allowbreak{}GreenBridge.\allowbreak{}th\_\allowbreak{}inj} (\texttt{Development.\allowbreak{}AbelianVarietyTheory.\allowbreak{}Proofs.\allowbreak{}InertiaEigenvalues.\allowbreak{}SumRangeEqRoots}, lines 142--147)\\
\texttt{GlobalTheory.\allowbreak{}GreenBridge.\allowbreak{}th\_\allowbreak{}surj} (\texttt{Development.\allowbreak{}AbelianVarietyTheory.\allowbreak{}Proofs.\allowbreak{}InertiaEigenvalues.\allowbreak{}SumRangeEqRoots}, lines 148--161)\par}

{\footnotesize\raggedright\noindent \ref{B7.34}.\quad \texttt{GlobalTheory.\allowbreak{}GreenBridge.\allowbreak{}th\_\allowbreak{}pow} (\texttt{Development.\allowbreak{}AbelianVarietyTheory.\allowbreak{}Proofs.\allowbreak{}InertiaEigenvalues.\allowbreak{}SumRangeEqRoots}, lines 137--141)\\
\texttt{GlobalTheory.\allowbreak{}GreenBridge.\allowbreak{}regular\_\allowbreak{}iff} (\texttt{Development.\allowbreak{}AbelianVarietyTheory.\allowbreak{}Proofs.\allowbreak{}InertiaEigenvalues.\allowbreak{}SumRangeEqRoots}, lines 207--212)\par}

{\footnotesize\raggedright\noindent \ref{B7.35}.\quad \texttt{GlobalTheory.\allowbreak{}GreenBridge.\allowbreak{}rootOf\_\allowbreak{}spec} (\texttt{Development.\allowbreak{}AbelianVarietyTheory.\allowbreak{}Proofs.\allowbreak{}InertiaEigenvalues.\allowbreak{}SumRangeEqRoots}, lines 189--196)\\
\texttt{GlobalTheory.\allowbreak{}GreenBridge.\allowbreak{}root\_\allowbreak{}unit} (\texttt{Development.\allowbreak{}AbelianVarietyTheory.\allowbreak{}Proofs.\allowbreak{}InertiaEigenvalues.\allowbreak{}SumRangeEqRoots}, lines 197--206)\\
\texttt{GlobalTheory.\allowbreak{}GreenBridge.\allowbreak{}pow\_\allowbreak{}inj} (\texttt{Development.\allowbreak{}AbelianVarietyTheory.\allowbreak{}Proofs.\allowbreak{}InertiaEigenvalues.\allowbreak{}SumRangeEqRoots}, lines 180--186)\par}

{\footnotesize\raggedright\noindent \ref{B7.36}.\quad \texttt{GlobalTheory.\allowbreak{}GreenBridge.\allowbreak{}sum\_\allowbreak{}range\_\allowbreak{}eq\_\allowbreak{}roots} (\texttt{Development.\allowbreak{}AbelianVarietyTheory.\allowbreak{}Proofs.\allowbreak{}InertiaEigenvalues.\allowbreak{}SumRangeEqRoots}, lines 213--217)\\
\texttt{GlobalTheory.\allowbreak{}GreenBridge.\allowbreak{}thetaOf\_\allowbreak{}pow} (\texttt{Development.\allowbreak{}AbelianVarietyTheory.\allowbreak{}Proofs.\allowbreak{}InertiaEigenvalues.\allowbreak{}SumRangeEqRoots}, lines 218--224)\par}

{\footnotesize\raggedright\noindent \ref{B7.37}.\quad \texttt{GlobalTheory.\allowbreak{}GreenBridge.\allowbreak{}green\_\allowbreak{}of\_\allowbreak{}regular} (\texttt{Development.\allowbreak{}AbelianVarietyTheory.\allowbreak{}Proofs.\allowbreak{}InertiaEigenvalues.\allowbreak{}SumRangeEqRoots}, lines 225--305)\par}

{\footnotesize\raggedright\noindent \ref{B7.37}.\quad \texttt{UpperTheory.\allowbreak{}LiteratureTheorems.\allowbreak{}green\_\allowbreak{}regular\_\allowbreak{}characters\_\allowbreak{}of\_\allowbreak{}literature} (\texttt{Development.\allowbreak{}AutomorphicArithmetic.\allowbreak{}ThetaCorrespondence.\allowbreak{}ThetaPairLifts.\allowbreak{}ThirdPlaceSeries.\allowbreak{}EveryCharacterOfLiteratureOfVWin}, lines 399--400)\\
\texttt{UpperTheory.\allowbreak{}LiteratureTheorems.\allowbreak{}published\_\allowbreak{}results\_\allowbreak{}eigen} (\texttt{Development.\allowbreak{}AutomorphicArithmetic.\allowbreak{}ThetaCorrespondence.\allowbreak{}ThetaPairLifts.\allowbreak{}ThirdPlaceSeries.\allowbreak{}EveryCharacterOfLiteratureOfVWin}, lines 666--726)\\
\texttt{UpperTheory.\allowbreak{}Lit.\allowbreak{}every\_\allowbreak{}character\_\allowbreak{}of\_\allowbreak{}theta\_\allowbreak{}integrals\_\allowbreak{}local\_\allowbreak{}matching} (\texttt{Development.\allowbreak{}AutomorphicArithmetic.\allowbreak{}ThetaCorrespondence.\allowbreak{}ThetaPairLifts.\allowbreak{}ThirdPlaceSeries.\allowbreak{}SiegelEisensteinSecondTermHypotheses}, lines 69--73)\par}

{\footnotesize\raggedright\noindent \ref{B7.38}.\quad \texttt{GlobalTheory.\allowbreak{}GLm.\allowbreak{}card\_\allowbreak{}proper\_\allowbreak{}subfield\_\allowbreak{}elements\_\allowbreak{}lt} (\texttt{Development.\allowbreak{}AbelianVarietyTheory.\allowbreak{}Proofs.\allowbreak{}FlagMeasures.\allowbreak{}CardProperSubfieldElementsLt}, lines 67--103)\par}

{\footnotesize\raggedright\noindent \ref{B7.39}.\quad \texttt{GlobalTheory.\allowbreak{}GLq.\allowbreak{}LabelUnique.\allowbreak{}exists\_\allowbreak{}eq\_\allowbreak{}of\_\allowbreak{}sum\_\allowbreak{}eq} (\texttt{Development.\allowbreak{}AbelianVarietyTheory.\allowbreak{}Proofs.\allowbreak{}InertiaEigenvalues.\allowbreak{}CuspidalCharacterLabelUniqueGLn}, lines 28--110)\par}

{\footnotesize\raggedright\noindent \ref{B7.40}.\quad \texttt{GlobalTheory.\allowbreak{}GLq.\allowbreak{}LabelUnique.\allowbreak{}label\_\allowbreak{}eq} (\texttt{Development.\allowbreak{}AbelianVarietyTheory.\allowbreak{}Proofs.\allowbreak{}InertiaEigenvalues.\allowbreak{}CuspidalCharacterLabelUniqueGLn}, lines 174--258)\\
\texttt{GlobalTheory.\allowbreak{}GLq.\allowbreak{}isCuspOf\_\allowbreak{}label\_\allowbreak{}unique} (\texttt{Development.\allowbreak{}AbelianVarietyTheory.\allowbreak{}Proofs.\allowbreak{}InertiaEigenvalues.\allowbreak{}CuspidalCharacterLabelUniqueGLn}, lines 263--274)\par}

{\footnotesize\raggedright\noindent \ref{B7.41}.\quad \texttt{GlobalTheory.\allowbreak{}GLq.\allowbreak{}HasLabel} (\texttt{Development.\allowbreak{}AbelianVarietyTheory.\allowbreak{}Definitions.\allowbreak{}NonGaloisDescent.\allowbreak{}LabelDatum}, lines 98--101)\\
\texttt{GlobalTheory.\allowbreak{}GLq.\allowbreak{}eSerC} (\texttt{Development.\allowbreak{}AbelianVarietyTheory.\allowbreak{}Definitions.\allowbreak{}NonGaloisDescent.\allowbreak{}LabelDatum}, lines 102--104)\\
\texttt{GlobalTheory.\allowbreak{}GLq.\allowbreak{}wtC} (\texttt{Development.\allowbreak{}AbelianVarietyTheory.\allowbreak{}Definitions.\allowbreak{}NonGaloisDescent.\allowbreak{}LabelDatum}, lines 105--107)\par}

{\footnotesize\raggedright\noindent \ref{B7.42}.\quad \texttt{GlobalTheory.\allowbreak{}GLq.\allowbreak{}Label.\allowbreak{}label\_\allowbreak{}main} (\texttt{Development.\allowbreak{}AbelianVarietyTheory.\allowbreak{}Proofs.\allowbreak{}InertiaEigenvalues.\allowbreak{}HarishChandraSeriesLabelExistsUniqueGLn}, lines 271--279)\\
\texttt{GlobalTheory.\allowbreak{}gl\_\allowbreak{}list\_\allowbreak{}product\_\allowbreak{}perm\_\allowbreak{}char} (\texttt{Development.\allowbreak{}AbelianVarietyTheory.\allowbreak{}Proofs.\allowbreak{}InertiaEigenvalues.\allowbreak{}HarishChandraProductListIsCharacterAndPermInvari}, lines 74--77)\par}

{\footnotesize\raggedright\noindent \ref{B7.43}.\quad \texttt{GlobalTheory.\allowbreak{}Top.\allowbreak{}Lq.\allowbreak{}exists\_\allowbreak{}perm} (\texttt{Development.\allowbreak{}AbelianVarietyTheory.\allowbreak{}Proofs.\allowbreak{}GlobalInterpolation.\allowbreak{}AdmissibleLabelWeightsOfGreenCuspidalParametrisa}, lines 89--101)\\
\texttt{GlobalTheory.\allowbreak{}Top.\allowbreak{}Lq.\allowbreak{}ext\_\allowbreak{}back} (\texttt{Development.\allowbreak{}AbelianVarietyTheory.\allowbreak{}Proofs.\allowbreak{}GlobalInterpolation.\allowbreak{}AdmissibleLabelWeightsOfGreenCuspidalParametrisa}, lines 129--134)\\
\texttt{GlobalTheory.\allowbreak{}Top.\allowbreak{}Lq.\allowbreak{}eSer\_\allowbreak{}eq} (\texttt{Development.\allowbreak{}AbelianVarietyTheory.\allowbreak{}Proofs.\allowbreak{}GlobalInterpolation.\allowbreak{}AdmissibleLabelWeightsOfGreenCuspidalParametrisa}, lines 233--250)\par}

{\footnotesize\raggedright\noindent \ref{B7.44}.\quad \texttt{GlobalTheory.\allowbreak{}T.\allowbreak{}whittakerJacquetIdempotents\_\allowbreak{}properties\_\allowbreak{}primeField} (\texttt{Development.\allowbreak{}AbelianVarietyTheory.\allowbreak{}Proofs.\allowbreak{}InertiaEigenvalues.\allowbreak{}WhittakerJacquetIdempotentsPropertiesPrimeFieldStatement}, lines 8--9)\\
\texttt{GlobalTheory.\allowbreak{}T.\allowbreak{}harishChandraSeriesWeight\_\allowbreak{}nonneg\_\allowbreak{}of\_\allowbreak{}idempotent\_\allowbreak{}glnPrime} (\texttt{Development.\allowbreak{}AbelianVarietyTheory.\allowbreak{}Proofs.\allowbreak{}InertiaEigenvalues.\allowbreak{}HarishChandraSeriesWeightNonnegOfIdempotentGlnPrStatement}, lines 8--9)\par}

{\footnotesize\raggedright\noindent \ref{B7.45}.\quad \texttt{GlobalTheory.\allowbreak{}GLq.\allowbreak{}wh\_\allowbreak{}cusp\_\allowbreak{}eq\_\allowbreak{}one} (\texttt{Development.\allowbreak{}AbelianVarietyTheory.\allowbreak{}Proofs.\allowbreak{}InertiaEigenvalues.\allowbreak{}WhittakerDimensionCuspidalGreenDegreeEqOne}, lines 66--80)\par}

{\footnotesize\raggedright\noindent \ref{B7.46}.\quad \texttt{ProductIntervalSplit.\allowbreak{}CompanionTheory.\allowbreak{}Member.\allowbreak{}harishChandraSeriesWeight\_\allowbreak{}squarefree\_\allowbreak{}label\_\allowbreak{}unipotent} (\texttt{Development.\allowbreak{}AbelianVarietyTheory.\allowbreak{}Proofs.\allowbreak{}InertiaEigenvalues.\allowbreak{}HarishChandraSeriesWeightSquarefreeLabelGelfandG}, lines 238--294)\par}

{\footnotesize\raggedright\noindent \ref{B7.47a}.\quad \texttt{GlobalTheory.\allowbreak{}GLq.\allowbreak{}tr\_\allowbreak{}eJ\_\allowbreak{}eq\_\allowbreak{}tr\_\allowbreak{}eJ'} (\texttt{Development.\allowbreak{}AbelianVarietyTheory.\allowbreak{}Proofs.\allowbreak{}InertiaEigenvalues.\allowbreak{}GroupAlgebraTracePatternIdempotentExchange}, lines 13--28)\\
\texttt{GlobalTheory.\allowbreak{}RootExch.\allowbreak{}step} (\texttt{Development.\allowbreak{}AbelianVarietyTheory.\allowbreak{}Proofs.\allowbreak{}FlagMeasures.\allowbreak{}PhIUnionTrans}, lines 327--343)\par}

{\footnotesize\raggedright\noindent \ref{B7.47b}.\quad \texttt{GlobalTheory.\allowbreak{}GLq.\allowbreak{}ExactMult.\allowbreak{}tr\_\allowbreak{}eJ'\_\allowbreak{}eq\_\allowbreak{}sum\_\allowbreak{}wh\_\allowbreak{}der} (\texttt{Development.\allowbreak{}AbelianVarietyTheory.\allowbreak{}Proofs.\allowbreak{}InertiaEigenvalues.\allowbreak{}TraceUnitriangularColumnVanishingIdempotentEqSum}, lines 69--93)\\
\texttt{GlobalTheory.\allowbreak{}GLm.\allowbreak{}tr\_\allowbreak{}eJ'\_\allowbreak{}eq\_\allowbreak{}wh\_\allowbreak{}tail} (\texttt{Development.\allowbreak{}AbelianVarietyTheory.\allowbreak{}Proofs.\allowbreak{}FlagMeasures.\allowbreak{}IsUnitOfIsUniUp}, lines 225--239)\\
\texttt{GlobalTheory.\allowbreak{}GLm.\allowbreak{}MirSum.\allowbreak{}eq\_\allowbreak{}sum\_\allowbreak{}indR\_\allowbreak{}der} (\texttt{Development.\allowbreak{}AbelianVarietyTheory.\allowbreak{}Proofs.\allowbreak{}FlagMeasures.\allowbreak{}MirabolicConjInvariantEqSumInducedDerivatives}, lines 110--133)\par}

{\footnotesize\raggedright\noindent \ref{B7.47}.\quad \texttt{GlobalTheory.\allowbreak{}GLq.\allowbreak{}lowerWtComparison} (\texttt{Development.\allowbreak{}AbelianVarietyTheory.\allowbreak{}Proofs.\allowbreak{}InertiaEigenvalues.\allowbreak{}HarishChandraSeriesWeightLowerComparison}, lines 49--98)\par}

{\footnotesize\raggedright\noindent \ref{B7.48}.\quad \texttt{GlobalTheory.\allowbreak{}GLq.\allowbreak{}Leibniz.\allowbreak{}main} (\texttt{Development.\allowbreak{}AbelianVarietyTheory.\allowbreak{}Proofs.\allowbreak{}InertiaEigenvalues.\allowbreak{}BernsteinZelevinskyDerivativeHarishChandraProduct}, lines 102--143)\par}

{\footnotesize\raggedright\noindent \ref{B7.49}.\quad \texttt{GlobalTheory.\allowbreak{}GLq.\allowbreak{}UpperWeight.\allowbreak{}comb\_\allowbreak{}bound} (\texttt{Development.\allowbreak{}AbelianVarietyTheory.\allowbreak{}Proofs.\allowbreak{}InertiaEigenvalues.\allowbreak{}HarishChandraSeriesWeightLeSquarefree}, lines 30--81)\\
\texttt{GlobalTheory.\allowbreak{}GLq.\allowbreak{}UpperWeight.\allowbreak{}wtC\_\allowbreak{}eJ\_\allowbreak{}le} (\texttt{Development.\allowbreak{}AbelianVarietyTheory.\allowbreak{}Proofs.\allowbreak{}InertiaEigenvalues.\allowbreak{}HarishChandraSeriesWeightLeSquarefree}, lines 88--140)\par}

{\footnotesize\raggedright\noindent \ref{B7.50}.\quad \texttt{GlobalTheory.\allowbreak{}Top.\allowbreak{}Lq.\allowbreak{}admissibleLabelWeights\_\allowbreak{}of\_\allowbreak{}greenCuspidalParametrisation} (\texttt{Development.\allowbreak{}AbelianVarietyTheory.\allowbreak{}Proofs.\allowbreak{}GlobalInterpolation.\allowbreak{}AdmissibleLabelWeightsOfGreenCuspidalParametrisa}, lines 255--300)\\
\texttt{Formalism.\allowbreak{}MainNotions.\allowbreak{}ArithmeticCylinders.\allowbreak{}admLq} (\texttt{Formalism.\allowbreak{}SharpRestriction.\allowbreak{}LevelLetterCocycles.\allowbreak{}ArithmeticCylinders}, lines 340--352)\\
\texttt{Formalism.\allowbreak{}L2.\allowbreak{}GlobalTheoryTheory.\allowbreak{}ConcNWeight.\allowbreak{}sum\_\allowbreak{}wtOf} (\texttt{Formalism.\allowbreak{}LocalToGlobal.\allowbreak{}GlobalTheory.\allowbreak{}ConcNWeight}, lines 74--87)\par}

{\footnotesize\raggedright\noindent \ref{B7.50a}.\quad \texttt{Formalism.\allowbreak{}IrreducibleCountingLemmas.\allowbreak{}GaloisIrreducibility.\allowbreak{}type\_\allowbreak{}set\_\allowbreak{}weight\_\allowbreak{}bound} (\texttt{Formalism.\allowbreak{}PadicSubsetSums.\allowbreak{}StableGradedCounts.\allowbreak{}LocalEstimateIVOfTypeSetWeight}, lines 1851--1893)\\
\texttt{Formalism.\allowbreak{}IrreducibleCountingLemmas.\allowbreak{}GaloisIrreducibility.\allowbreak{}type\_\allowbreak{}weight\_\allowbreak{}exact} (\texttt{Formalism.\allowbreak{}PadicSubsetSums.\allowbreak{}StableGradedCounts.\allowbreak{}LocalEstimateIVOfTypeSetWeight}, lines 1282--1315)\\
\texttt{Formalism.\allowbreak{}IrreducibleCountingLemmas.\allowbreak{}GaloisIrreducibility.\allowbreak{}divisor\_\allowbreak{}square\_\allowbreak{}weight\_\allowbreak{}sum\_\allowbreak{}le} (\texttt{Formalism.\allowbreak{}PadicSubsetSums.\allowbreak{}StableGradedCounts.\allowbreak{}LocalEstimateIVOfTypeSetWeight}, lines 1908--1953)\\
\texttt{Formalism.\allowbreak{}IrreducibleCountingLemmas.\allowbreak{}GaloisIrreducibility.\allowbreak{}square\_\allowbreak{}mass\_\allowbreak{}lower\_\allowbreak{}bound} (\texttt{Formalism.\allowbreak{}PadicSubsetSums.\allowbreak{}StableGradedCounts.\allowbreak{}LocalEstimateIVOfTypeSetWeight}, lines 1243--1268)\par}

{\footnotesize\raggedright\noindent \ref{B7.50b}.\quad \texttt{Formalism.\allowbreak{}IrreducibleCountingLemmas.\allowbreak{}GaloisIrreducibility.\allowbreak{}local\_\allowbreak{}estimate\_\allowbreak{}fourth} (\texttt{Formalism.\allowbreak{}PadicSubsetSums.\allowbreak{}StableGradedCounts.\allowbreak{}LocalEstimateIVOfTypeSetWeight}, lines 2451--2453)\\
\texttt{Formalism.\allowbreak{}IrreducibleCountingLemmas.\allowbreak{}GaloisIrreducibility.\allowbreak{}divisor\_\allowbreak{}degree\_\allowbreak{}subsum} (\texttt{Formalism.\allowbreak{}PadicSubsetSums.\allowbreak{}StableGradedCounts.\allowbreak{}IRRDivisorDegreeSubSum}, lines 14--35)\\
\texttt{Formalism.\allowbreak{}IrreducibleCountingLemmas.\allowbreak{}GaloisIrreducibility.\allowbreak{}prodMaxInvLeSq} (\texttt{Formalism.\allowbreak{}PadicSubsetSums.\allowbreak{}StableGradedCounts.\allowbreak{}IRRLemmaP}, lines 133--214)\par}

{\footnotesize\raggedright\noindent \ref{B7.50c}.\quad \texttt{Formalism.\allowbreak{}IrreducibleCountingLemmas.\allowbreak{}GaloisIrreducibility.\allowbreak{}prime\_\allowbreak{}reciprocal\_\allowbreak{}sum\_\allowbreak{}bound} (\texttt{Formalism.\allowbreak{}PadicSubsetSums.\allowbreak{}StableGradedCounts.\allowbreak{}IRRPrimeRecipSum}, lines 333--340)\\
\texttt{Formalism.\allowbreak{}IrreducibleCountingLemmas.\allowbreak{}GaloisIrreducibility.\allowbreak{}PrimeCountingLemmas.\allowbreak{}S\_\allowbreak{}small} (\texttt{Formalism.\allowbreak{}PadicSubsetSums.\allowbreak{}StableGradedCounts.\allowbreak{}IRRPrimeRecipSum}, lines 122--128)\\
\texttt{Formalism.\allowbreak{}IrreducibleCountingLemmas.\allowbreak{}GaloisIrreducibility.\allowbreak{}PrimeCountingLemmas.\allowbreak{}upper\_\allowbreak{}card} (\texttt{Formalism.\allowbreak{}PadicSubsetSums.\allowbreak{}StableGradedCounts.\allowbreak{}IRRPrimeRecipSum}, lines 211--270)\par}

{\footnotesize\raggedright\noindent \ref{B7.50d}.\quad \texttt{Formalism.\allowbreak{}IrreducibleCountingLemmas.\allowbreak{}GaloisIrreducibility.\allowbreak{}local\_\allowbreak{}estimate\_\allowbreak{}fifth} (\texttt{Formalism.\allowbreak{}PadicSubsetSums.\allowbreak{}StableGradedCounts.\allowbreak{}LocalEstimateIVOfTypeSetWeight}, lines 2454--2456)\\
\texttt{Formalism.\allowbreak{}IrreducibleCountingLemmas.\allowbreak{}GaloisIrreducibility.\allowbreak{}prime\_\allowbreak{}cycle\_\allowbreak{}perm\_\allowbreak{}count} (\texttt{Formalism.\allowbreak{}PadicSubsetSums.\allowbreak{}StableGradedCounts.\allowbreak{}IRRCycleProbPrime}, lines 32--112)\\
\texttt{Formalism.\allowbreak{}IrreducibleCountingLemmas.\allowbreak{}GaloisIrreducibility.\allowbreak{}prime\_\allowbreak{}cycle\_\allowbreak{}probability\_\allowbreak{}sum} (\texttt{Formalism.\allowbreak{}PadicSubsetSums.\allowbreak{}StableGradedCounts.\allowbreak{}IRRCycleProbSum}, lines 6--7)\par}

{\footnotesize\raggedright\noindent \ref{B7.50e}.\quad \texttt{Formalism.\allowbreak{}IrreducibleCountingLemmas.\allowbreak{}GaloisIrreducibility.\allowbreak{}local\_\allowbreak{}estimate\_\allowbreak{}second} (\texttt{Formalism.\allowbreak{}PadicSubsetSums.\allowbreak{}StableGradedCounts.\allowbreak{}SquareSetMemExceptionalUnion}, lines 429--584)\\
\texttt{Formalism.\allowbreak{}IrreducibleCountingLemmas.\allowbreak{}GaloisIrreducibility.\allowbreak{}same\_\allowbreak{}factor\_\allowbreak{}count} (\texttt{Formalism.\allowbreak{}PadicSubsetSums.\allowbreak{}StableGradedCounts.\allowbreak{}CoeffZeroNeZeroOfIrreducible}, lines 152--288)\\
\texttt{Formalism.\allowbreak{}IrreducibleCountingLemmas.\allowbreak{}GaloisIrreducibility.\allowbreak{}cross\_\allowbreak{}factor\_\allowbreak{}count} (\texttt{Formalism.\allowbreak{}PadicSubsetSums.\allowbreak{}StableGradedCounts.\allowbreak{}CrossFactorCount}, lines 51--103)\\
\texttt{Formalism.\allowbreak{}IrreducibleCountingLemmas.\allowbreak{}GaloisIrreducibility.\allowbreak{}local\_\allowbreak{}estimate\_\allowbreak{}third} (\texttt{Formalism.\allowbreak{}PadicSubsetSums.\allowbreak{}StableGradedCounts.\allowbreak{}LocalEstimateThird}, lines 354--384)\par}

{\footnotesize\raggedright\noindent \ref{B7.50y}.\quad \texttt{Formalism.\allowbreak{}IrreducibleCountingLemmas.\allowbreak{}GaloisIrreducibility.\allowbreak{}square\_\allowbreak{}mass\_\allowbreak{}lower\_\allowbreak{}bound} (\texttt{Formalism.\allowbreak{}PadicSubsetSums.\allowbreak{}StableGradedCounts.\allowbreak{}LocalEstimateIVOfTypeSetWeight}, lines 1243--1267)\\
\texttt{Formalism.\allowbreak{}IrreducibleCountingLemmas.\allowbreak{}GaloisIrreducibility.\allowbreak{}local\_\allowbreak{}estimate\_\allowbreak{}second} (\texttt{Formalism.\allowbreak{}PadicSubsetSums.\allowbreak{}StableGradedCounts.\allowbreak{}SquareSetMemExceptionalUnion}, lines 429--583)\\
\texttt{Formalism.\allowbreak{}IrreducibleCountingLemmas.\allowbreak{}GaloisIrreducibility.\allowbreak{}local\_\allowbreak{}estimate\_\allowbreak{}third} (\texttt{Formalism.\allowbreak{}PadicSubsetSums.\allowbreak{}StableGradedCounts.\allowbreak{}LocalEstimateThird}, lines 354--383)\\
\texttt{Formalism.\allowbreak{}IrreducibleCountingLemmas.\allowbreak{}GaloisIrreducibility.\allowbreak{}localEstimateIV\_\allowbreak{}of} (\texttt{Formalism.\allowbreak{}PadicSubsetSums.\allowbreak{}StableGradedCounts.\allowbreak{}LocalEstimateIVOfTypeSetWeight}, lines 1590--1607)\\
\texttt{Formalism.\allowbreak{}IrreducibleCountingLemmas.\allowbreak{}GaloisIrreducibility.\allowbreak{}localEstimateV\_\allowbreak{}of} (\texttt{Formalism.\allowbreak{}PadicSubsetSums.\allowbreak{}StableGradedCounts.\allowbreak{}LocalEstimateIVOfTypeSetWeight}, lines 1609--1679)\\
\texttt{Formalism.\allowbreak{}IrreducibleCountingLemmas.\allowbreak{}GaloisIrreducibility.\allowbreak{}prodMaxInvLeSq} (\texttt{Formalism.\allowbreak{}PadicSubsetSums.\allowbreak{}StableGradedCounts.\allowbreak{}IRRLemmaP}, lines 133--214)\\
\texttt{Formalism.\allowbreak{}IrreducibleCountingLemmas.\allowbreak{}GaloisIrreducibility.\allowbreak{}prime\_\allowbreak{}reciprocal\_\allowbreak{}sum\_\allowbreak{}bound} (\texttt{Formalism.\allowbreak{}PadicSubsetSums.\allowbreak{}StableGradedCounts.\allowbreak{}IRRPrimeRecipSum}, lines 333--340)\par}

{\footnotesize\raggedright\noindent \ref{B7.50y}.\quad \texttt{Formalism.\allowbreak{}CyclotomicNorms.\allowbreak{}IsGood} (\texttt{Formalism.\allowbreak{}CyclotomicNorms.\allowbreak{}StatementT}, lines 19--30)\\
\texttt{Formalism.\allowbreak{}MainNotions.\allowbreak{}ArithmeticCylinders.\allowbreak{}goodAt} (\texttt{Formalism.\allowbreak{}SharpRestriction.\allowbreak{}LevelLetterCocycles.\allowbreak{}ArithmeticCylinders}, lines 355--359)\par}

{\footnotesize\raggedright\noindent \ref{B7.51}.\quad \texttt{Formalism.\allowbreak{}L2.\allowbreak{}GlobalTheoryTheory.\allowbreak{}ConcNWeight.\allowbreak{}weightAllRows\_\allowbreak{}of} (\texttt{Formalism.\allowbreak{}LocalToGlobal.\allowbreak{}GlobalTheory.\allowbreak{}ConcNWeight}, lines 345--356)\\
\texttt{Formalism.\allowbreak{}L2.\allowbreak{}GlobalTheoryTheory.\allowbreak{}AllJoin.\allowbreak{}weight\_\allowbreak{}of\_\allowbreak{}join} (\texttt{Formalism.\allowbreak{}LocalToGlobal.\allowbreak{}HilbertRowConditions.\allowbreak{}AllJoinP}, lines 22--25)\\
\texttt{Formalism.\allowbreak{}IrreducibleCountingLemmas.\allowbreak{}GaloisIrreducibility.\allowbreak{}property\_\allowbreak{}W\_\allowbreak{}of\_\allowbreak{}estimates} (\texttt{Formalism.\allowbreak{}PadicSubsetSums.\allowbreak{}StableGradedCounts.\allowbreak{}PropertyWOfEstimates}, lines 7--8)\\
\texttt{Formalism.\allowbreak{}MainNotions.\allowbreak{}N.\allowbreak{}PlaceLabelTeichmuller.\allowbreak{}good\_\allowbreak{}of} (\texttt{Formalism.\allowbreak{}SharpRestriction.\allowbreak{}LevelLetterCocycles.\allowbreak{}PlaceLabelTeichmuller}, lines 660--681)\par}

{\footnotesize\raggedright\noindent \ref{def:p7c1-compatible-rows}.\quad \texttt{Formalism.\allowbreak{}L2.\allowbreak{}GlobalTheoryTheory.\allowbreak{}AllJoin.\allowbreak{}AllJoinA} (\texttt{Formalism.\allowbreak{}LocalToGlobal.\allowbreak{}HilbertRowConditions.\allowbreak{}AllJoin}, lines 18--23)\\
\texttt{Formalism.\allowbreak{}MainNotions.\allowbreak{}PlaceDegreeData.\allowbreak{}RowsTie} (\texttt{Formalism.\allowbreak{}SharpRestriction.\allowbreak{}CuspSpanNormBounds.\allowbreak{}PlaceDegreeData}, lines 380--383)\\
\texttt{Formalism.\allowbreak{}MainNotions.\allowbreak{}S.\allowbreak{}AutTie} (\texttt{Formalism.\allowbreak{}SharpRestriction.\allowbreak{}CuspSpanNormBounds.\allowbreak{}DstarEq}, lines 10--23)\\
\texttt{Formalism.\allowbreak{}L2.\allowbreak{}GlobalTheoryTheory.\allowbreak{}AllJoin.\allowbreak{}LevelSharp} (\texttt{Formalism.\allowbreak{}LocalToGlobal.\allowbreak{}HilbertRowConditions.\allowbreak{}AllJoin}, lines 14--17)\par}

{\footnotesize\raggedright\noindent \ref{B7.52}.\quad \texttt{GlobalTheory.\allowbreak{}Top.\allowbreak{}Irr.\allowbreak{}linkIRR\_\allowbreak{}of\_\allowbreak{}join\_\allowbreak{}ex} (\texttt{Development.\allowbreak{}AbelianVarietyTheory.\allowbreak{}Proofs.\allowbreak{}GlobalInterpolation.\allowbreak{}LinkIRROfJoinEx}, lines 17--33)\\
\texttt{Formalism.\allowbreak{}L2.\allowbreak{}GlobalTheoryTheory.\allowbreak{}ConcN.\allowbreak{}CountAllRows} (\texttt{Formalism.\allowbreak{}LocalToGlobal.\allowbreak{}GlobalTheory.\allowbreak{}IrreducibleOfCountSizes}, lines 183--187)\\
\texttt{Formalism.\allowbreak{}MainNotions.\allowbreak{}ArithmeticCylinders.\allowbreak{}AllRows} (\texttt{Formalism.\allowbreak{}SharpRestriction.\allowbreak{}LevelLetterCocycles.\allowbreak{}ArithmeticCylinders}, lines 244--260)\\
\texttt{Formalism.\allowbreak{}MainNotions.\allowbreak{}ArithmeticCylinders.\allowbreak{}LabelCountGoodInjective} (\texttt{Formalism.\allowbreak{}SharpRestriction.\allowbreak{}LevelLetterCocycles.\allowbreak{}ArithmeticCylinders}, lines 271--274)\par}

{\footnotesize\raggedright\noindent \ref{B7.53}.\quad \texttt{GlobalTheory.\allowbreak{}Top.\allowbreak{}Irr.\allowbreak{}closed\_\allowbreak{}linkIRR\_\allowbreak{}of\_\allowbreak{}join\_\allowbreak{}count} (\texttt{Development.\allowbreak{}AutomorphicArithmetic.\allowbreak{}ThetaCorrespondence.\allowbreak{}ThetaPairLifts.\allowbreak{}ThirdPlaceSeries.\allowbreak{}LinkIrr}, lines 22--43)\\
\texttt{GlobalTheory.\allowbreak{}Top.\allowbreak{}Irr.\allowbreak{}sizesRec\_\allowbreak{}at} (\texttt{Development.\allowbreak{}AbelianVarietyTheory.\allowbreak{}Proofs.\allowbreak{}GlobalInterpolation.\allowbreak{}LinkIRROfJoinEx}, lines 39--137)\par}

{\footnotesize\raggedright\noindent \ref{B7.54}.\quad \texttt{UpperTheory.\allowbreak{}Top.\allowbreak{}shape\_\allowbreak{}closed\_\allowbreak{}link\_\allowbreak{}n\_\allowbreak{}of\_\allowbreak{}green} (\texttt{Development.\allowbreak{}AutomorphicArithmetic.\allowbreak{}ThetaCorrespondence.\allowbreak{}ThetaPairLifts.\allowbreak{}ThirdPlaceSeries.\allowbreak{}TopShape}, lines 33--55)\par}

\section*{Chapter~\ref{ch:p8c1}: A common witness for the 46 conditions}

{\footnotesize\raggedright\noindent Chapter introduction.\quad \texttt{GlobalTheory.\allowbreak{}Cp.\allowbreak{}D3} (\texttt{Development.\allowbreak{}AbelianVarietyTheory.\allowbreak{}Definitions.\allowbreak{}GlobalTwisting.\allowbreak{}CountsIrreducibilityConstants}, lines 35--35)\\
\texttt{Formalism.\allowbreak{}MainNotions.\allowbreak{}Des} (\texttt{Formalism.\allowbreak{}SharpRestriction.\allowbreak{}SmoothDesignExistence.\allowbreak{}DesignExistsDes}, lines 10--11)\\
\texttt{Formalism.\allowbreak{}MainNotions.\allowbreak{}Pk.\allowbreak{}DGamSmooth} (\texttt{Formalism.\allowbreak{}SharpRestriction.\allowbreak{}SmoothDesignExistence.\allowbreak{}DGamSmooth}, lines 23--26)\par}

{\footnotesize\raggedright\noindent \ref{B8.01}.\quad \texttt{UpperTheory.\allowbreak{}Lit.\allowbreak{}nonempty\_\allowbreak{}hyps\_\allowbreak{}definite\_\allowbreak{}side\_\allowbreak{}of\_\allowbreak{}closed\_\allowbreak{}bound} (\texttt{Development.\allowbreak{}AutomorphicArithmetic.\allowbreak{}ThetaCorrespondence.\allowbreak{}ThetaPairLifts.\allowbreak{}ThirdPlaceSeries.\allowbreak{}SiegelEisensteinSecondTermHypotheses}, lines 44--68)\par}

{\footnotesize\raggedright\noindent \ref{B8.01ka}.\quad \texttt{Formalism.\allowbreak{}BasicNotions.\allowbreak{}NTw} (\texttt{Formalism.\allowbreak{}SharpRestriction.\allowbreak{}LevelLetterCocycles.\allowbreak{}DataN}, lines 256--427)\\
\texttt{Formalism.\allowbreak{}MainNotions.\allowbreak{}rPrAv} (\texttt{Formalism.\allowbreak{}SharpRestriction.\allowbreak{}LevelLetterCocycles.\allowbreak{}TwGood}, lines 114--116)\\
\texttt{Formalism.\allowbreak{}MainNotions.\allowbreak{}q0Av} (\texttt{Formalism.\allowbreak{}SharpRestriction.\allowbreak{}LevelLetterCocycles.\allowbreak{}TwGood}, lines 118--118)\\
\texttt{Formalism.\allowbreak{}MainNotions.\allowbreak{}avPrimes} (\texttt{Formalism.\allowbreak{}SharpRestriction.\allowbreak{}LevelLetterCocycles.\allowbreak{}TwGood}, lines 120--124)\\
\texttt{Formalism.\allowbreak{}MainNotions.\allowbreak{}mem\_\allowbreak{}avRec} (\texttt{Formalism.\allowbreak{}SharpRestriction.\allowbreak{}LevelLetterCocycles.\allowbreak{}TwGood}, lines 146--147)\par}

{\footnotesize\raggedright\noindent \ref{B8.01ka}.\quad \texttt{Formalism.\allowbreak{}BasicNotions.\allowbreak{}G0g} (\texttt{Formalism.\allowbreak{}SharpRestriction.\allowbreak{}LevelLetterCocycles.\allowbreak{}DataN}, lines 761--763)\\
\texttt{Formalism.\allowbreak{}BasicNotions.\allowbreak{}NExp} (\texttt{Formalism.\allowbreak{}SharpRestriction.\allowbreak{}LevelLetterCocycles.\allowbreak{}DataN}, lines 769--854)\\
\texttt{Formalism.\allowbreak{}BasicNotions.\allowbreak{}ThData} (\texttt{Formalism.\allowbreak{}LocalRingCardinality.\allowbreak{}DataNChTie}, lines 37--40)\\
\texttt{Formalism.\allowbreak{}MainNotions.\allowbreak{}IntegralHeckeForms.\allowbreak{}AO} (\texttt{Formalism.\allowbreak{}SharpRestriction.\allowbreak{}LevelLetterCocycles.\allowbreak{}IntegralHeckeForms}, lines 23--25)\\
\texttt{Formalism.\allowbreak{}MainNotions.\allowbreak{}IntegralHeckeForms.\allowbreak{}V0} (\texttt{Formalism.\allowbreak{}SharpRestriction.\allowbreak{}LevelLetterCocycles.\allowbreak{}IntegralHeckeForms}, lines 36--37)\\
\texttt{Formalism.\allowbreak{}MainNotions.\allowbreak{}IntegralHeckeForms.\allowbreak{}eTh} (\texttt{Formalism.\allowbreak{}SharpRestriction.\allowbreak{}LevelLetterCocycles.\allowbreak{}IntegralHeckeForms}, lines 52--54)\\
\texttt{Formalism.\allowbreak{}MainNotions.\allowbreak{}IntegralHeckeForms.\allowbreak{}eThPsi} (\texttt{Formalism.\allowbreak{}SharpRestriction.\allowbreak{}LevelLetterCocycles.\allowbreak{}IntegralHeckeForms}, lines 55--59)\\
\texttt{Formalism.\allowbreak{}MainNotions.\allowbreak{}IntegralHeckeForms.\allowbreak{}eThPMt} (\texttt{Formalism.\allowbreak{}SharpRestriction.\allowbreak{}LevelLetterCocycles.\allowbreak{}IntegralHeckeForms}, lines 60--63)\par}

{\footnotesize\raggedright\noindent \ref{B8.01ka}.\quad \texttt{Formalism.\allowbreak{}MainNotions.\allowbreak{}IntegralHeckeForms.\allowbreak{}sIdx} (\texttt{Formalism.\allowbreak{}SharpRestriction.\allowbreak{}LevelLetterCocycles.\allowbreak{}IntegralHeckeForms}, lines 68--71)\\
\texttt{Formalism.\allowbreak{}MainNotions.\allowbreak{}IntegralHeckeForms.\allowbreak{}IdxAt} (\texttt{Formalism.\allowbreak{}SharpRestriction.\allowbreak{}LevelLetterCocycles.\allowbreak{}IntegralHeckeForms}, lines 201--212)\par}

{\footnotesize\raggedright\noindent \ref{B8.01ka}.\quad \texttt{Formalism.\allowbreak{}MainNotions.\allowbreak{}BlochKatoReduction.\allowbreak{}Gen.\allowbreak{}thEval} (\texttt{Formalism.\allowbreak{}SharpRestriction.\allowbreak{}LevelLetterCocycles.\allowbreak{}BlochKatoPhases}, lines 211--215)\\
\texttt{Formalism.\allowbreak{}MainNotions.\allowbreak{}BlochKatoReduction.\allowbreak{}Gen.\allowbreak{}c0} (\texttt{Formalism.\allowbreak{}SharpRestriction.\allowbreak{}LevelLetterCocycles.\allowbreak{}BlochKatoPhases}, lines 224--225)\\
\texttt{Formalism.\allowbreak{}MainNotions.\allowbreak{}BlochKatoReduction.\allowbreak{}Gen.\allowbreak{}tameVal} (\texttt{Formalism.\allowbreak{}SharpRestriction.\allowbreak{}LevelLetterCocycles.\allowbreak{}BlochKatoPhases}, lines 226--229)\\
\texttt{GlobalTheory.\allowbreak{}LamTie.\allowbreak{}sum\_\allowbreak{}th\_\allowbreak{}empty\_\allowbreak{}ne\_\allowbreak{}zero} (\texttt{Development.\allowbreak{}AbelianVarietyTheory.\allowbreak{}Proofs.\allowbreak{}InertiaEigenvalues.\allowbreak{}SumThEmptyNeZero}, lines 10--24)\par}

{\footnotesize\raggedright\noindent \ref{B8.01k}.\quad \texttt{GlobalTheory.\allowbreak{}Cp.\allowbreak{}Duty} (\texttt{Development.\allowbreak{}AbelianVarietyTheory.\allowbreak{}Definitions.\allowbreak{}GlobalTwisting.\allowbreak{}CountsIrreducibilityConstants}, lines 48--49)\\
\texttt{Formalism.\allowbreak{}MainNotions.\allowbreak{}PageFix} (\texttt{Formalism.\allowbreak{}SharpRestriction.\allowbreak{}PageFix}, lines 9--18)\\
\texttt{Formalism.\allowbreak{}MainNotions.\allowbreak{}PageWin} (\texttt{Formalism.\allowbreak{}SharpRestriction.\allowbreak{}PageFix}, lines 20--25)\\
\texttt{Formalism.\allowbreak{}MainNotions.\allowbreak{}TameDistinct} (\texttt{Formalism.\allowbreak{}SharpRestriction.\allowbreak{}LevelLetterCocycles.\allowbreak{}AlignN}, lines 14--19)\\
\texttt{Formalism.\allowbreak{}MainNotions.\allowbreak{}TwGoodAtAll} (\texttt{Formalism.\allowbreak{}SharpRestriction.\allowbreak{}LevelLetterCocycles.\allowbreak{}TwAllF}, lines 22--23)\\
\texttt{Formalism.\allowbreak{}MainNotions.\allowbreak{}IntegralHeckeForms.\allowbreak{}GoodTw} (\texttt{Formalism.\allowbreak{}SharpRestriction.\allowbreak{}LevelLetterCocycles.\allowbreak{}IntegralHeckeForms}, lines 219--231)\\
\texttt{Formalism.\allowbreak{}MainNotions.\allowbreak{}BlochKatoReduction.\allowbreak{}Gen.\allowbreak{}goodE} (\texttt{Formalism.\allowbreak{}SharpRestriction.\allowbreak{}LevelLetterCocycles.\allowbreak{}BlochKatoPhases}, lines 100--109)\\
\texttt{Formalism.\allowbreak{}MainNotions.\allowbreak{}BlochKatoReduction.\allowbreak{}Gen.\allowbreak{}bk} (\texttt{Formalism.\allowbreak{}SharpRestriction.\allowbreak{}LevelLetterCocycles.\allowbreak{}BlochKatoPhases}, lines 167--179)\\
\texttt{Formalism.\allowbreak{}MainNotions.\allowbreak{}BlochKatoReduction.\allowbreak{}Gen.\allowbreak{}selT} (\texttt{Formalism.\allowbreak{}SharpRestriction.\allowbreak{}LevelLetterCocycles.\allowbreak{}BlochKatoPhases}, lines 184--188)\\
\texttt{Formalism.\allowbreak{}MainNotions.\allowbreak{}BlochKatoReduction.\allowbreak{}Gen.\allowbreak{}sLen} (\texttt{Formalism.\allowbreak{}SharpRestriction.\allowbreak{}LevelLetterCocycles.\allowbreak{}BlochKatoPhases}, lines 193--196)\\
\texttt{Formalism.\allowbreak{}MainNotions.\allowbreak{}BlochKatoReduction.\allowbreak{}Gen.\allowbreak{}partAT} (\texttt{Formalism.\allowbreak{}SharpRestriction.\allowbreak{}LevelLetterCocycles.\allowbreak{}BlochKatoPhases}, lines 304--308)\\
\texttt{Formalism.\allowbreak{}MainNotions.\allowbreak{}BlochKatoReduction.\allowbreak{}Gen.\allowbreak{}partC} (\texttt{Formalism.\allowbreak{}SharpRestriction.\allowbreak{}LevelLetterCocycles.\allowbreak{}BlochKatoPhases}, lines 317--321)\\
\texttt{Formalism.\allowbreak{}MainNotions.\allowbreak{}BlochKatoReduction.\allowbreak{}Gen.\allowbreak{}partD} (\texttt{Formalism.\allowbreak{}SharpRestriction.\allowbreak{}LevelLetterCocycles.\allowbreak{}BlochKatoPhases}, lines 323--325)\par}

{\footnotesize\raggedright\noindent \ref{def:p8c1-design-bounds}.\quad \texttt{Formalism.\allowbreak{}MainNotions.\allowbreak{}ThickSetEventually} (\texttt{Formalism.\allowbreak{}SharpRestriction.\allowbreak{}GlobalSchurBounds.\allowbreak{}Chosen}, lines 65--66)\\
\texttt{Formalism.\allowbreak{}BasicNotions.\allowbreak{}Host.\allowbreak{}secondSet\_\allowbreak{}card\_\allowbreak{}lower\_\allowbreak{}bound} (\texttt{Formalism.\allowbreak{}LocalRingCardinality.\allowbreak{}DataHost}, lines 165--167)\\
\texttt{Formalism.\allowbreak{}MainNotions.\allowbreak{}ChTT\_\allowbreak{}at} (\texttt{Formalism.\allowbreak{}SharpRestriction.\allowbreak{}GlobalSchurBounds.\allowbreak{}Chosen}, lines 18--20)\\
\texttt{Formalism.\allowbreak{}MainNotions.\allowbreak{}ChBad\_\allowbreak{}at} (\texttt{Formalism.\allowbreak{}SharpRestriction.\allowbreak{}GlobalSchurBounds.\allowbreak{}Chosen}, lines 34--37)\\
\texttt{Formalism.\allowbreak{}MainNotions.\allowbreak{}ChQD\_\allowbreak{}at} (\texttt{Formalism.\allowbreak{}SharpRestriction.\allowbreak{}GlobalSchurBounds.\allowbreak{}Chosen}, lines 42--46)\\
\texttt{Formalism.\allowbreak{}MainNotions.\allowbreak{}ChCode\_\allowbreak{}at} (\texttt{Formalism.\allowbreak{}SharpRestriction.\allowbreak{}GlobalSchurBounds.\allowbreak{}Chosen}, lines 56--58)\\
\texttt{Formalism.\allowbreak{}MainNotions.\allowbreak{}CodeUpF} (\texttt{Formalism.\allowbreak{}SharpRestriction.\allowbreak{}DesignExistsOfChoices}, lines 18--21)\par}

{\footnotesize\raggedright\noindent \ref{def:p8c1-good-labels}.\quad \texttt{Formalism.\allowbreak{}CyclotomicNorms.\allowbreak{}IsGood} (\texttt{Formalism.\allowbreak{}CyclotomicNorms.\allowbreak{}StatementT}, lines 19--30)\\
\texttt{Formalism.\allowbreak{}BasicNotions.\allowbreak{}NData.\allowbreak{}goodRow} (\texttt{Formalism.\allowbreak{}SharpRestriction.\allowbreak{}LevelLetterCocycles.\allowbreak{}DataN}, lines 234--238)\\
\texttt{Formalism.\allowbreak{}BasicNotions.\allowbreak{}NTypes.\allowbreak{}goodTuple} (\texttt{Formalism.\allowbreak{}SharpRestriction.\allowbreak{}LevelLetterCocycles.\allowbreak{}DataN}, lines 996--1000)\par}

{\footnotesize\raggedright\noindent \ref{B8.02}.\quad \texttt{GlobalTheory.\allowbreak{}Fc.\allowbreak{}BundleCondition} (\texttt{Development.\allowbreak{}AutomorphicArithmetic.\allowbreak{}ThetaCorrespondence.\allowbreak{}ThetaPairLifts.\allowbreak{}ThirdPlaceSeries.\allowbreak{}ColumnChoiceWithTies}, lines 47--54)\\
\texttt{GlobalTheory.\allowbreak{}Fc.\allowbreak{}familySelection} (\texttt{Development.\allowbreak{}AutomorphicArithmetic.\allowbreak{}ThetaCorrespondence.\allowbreak{}ThetaPairLifts.\allowbreak{}ThirdPlaceSeries.\allowbreak{}ColumnChoiceWithTies}, lines 56--56)\\
\texttt{Formalism.\allowbreak{}MainNotions.\allowbreak{}Closed} (\texttt{Formalism.\allowbreak{}SharpRestriction.\allowbreak{}ParamsHeightRankConstant}, lines 159--161)\par}

{\footnotesize\raggedright\noindent \ref{B8.02}.\quad \texttt{GlobalTheory.\allowbreak{}Fc.\allowbreak{}R1.\allowbreak{}SelectionDefiniteRows} (\texttt{Development.\allowbreak{}AutomorphicArithmetic.\allowbreak{}ThetaCorrespondence.\allowbreak{}ThetaPairLifts.\allowbreak{}ThirdPlaceSeries.\allowbreak{}CountsIrreducibilityConstants}, lines 91--93)\par}

{\footnotesize\raggedright\noindent \ref{B8.02}.\quad \texttt{GlobalTheory.\allowbreak{}Fc.\allowbreak{}HypsAlt} (\texttt{Development.\allowbreak{}AutomorphicArithmetic.\allowbreak{}ThetaCorrespondence.\allowbreak{}ThetaPairLifts.\allowbreak{}ThirdPlaceSeries.\allowbreak{}Hyps}, lines 22--90)\\
\texttt{UpperTheory.\allowbreak{}Top.\allowbreak{}FReq.\allowbreak{}cls} (\texttt{Development.\allowbreak{}AutomorphicArithmetic.\allowbreak{}CharacterBoxes.\allowbreak{}Supply}, lines 22--24)\\
\texttt{UpperTheory.\allowbreak{}Top.\allowbreak{}ExS} (\texttt{Development.\allowbreak{}AutomorphicArithmetic.\allowbreak{}CharacterBoxes.\allowbreak{}ExS}, lines 10--13)\\
\texttt{GlobalTheory.\allowbreak{}EX.\allowbreak{}HostExW} (\texttt{Development.\allowbreak{}AbelianVarietyTheory.\allowbreak{}Definitions.\allowbreak{}GlobalTwisting.\allowbreak{}ThinScalesIffHostBound}, lines 120--124)\\
\texttt{Formalism.\allowbreak{}MainNotions.\allowbreak{}MinkMult} (\texttt{Formalism.\allowbreak{}SharpRestriction.\allowbreak{}GlobalSchurBounds.\allowbreak{}MainIn}, lines 51--53)\\
\texttt{Formalism.\allowbreak{}MainNotions.\allowbreak{}DesignChoices} (\texttt{Formalism.\allowbreak{}SharpRestriction.\allowbreak{}DesignExistsOfChoices}, lines 22--27)\\
\texttt{Formalism.\allowbreak{}MainNotions.\allowbreak{}TwGoodE} (\texttt{Formalism.\allowbreak{}SharpRestriction.\allowbreak{}LevelLetterCocycles.\allowbreak{}TwGood}, lines 219--224)\\
\texttt{Formalism.\allowbreak{}MainNotions.\allowbreak{}TwGoodAll} (\texttt{Formalism.\allowbreak{}SharpRestriction.\allowbreak{}LevelLetterCocycles.\allowbreak{}ExistsGoodPlacesEveryConjugate}, lines 38--41)\\
\texttt{Formalism.\allowbreak{}IndefR.\allowbreak{}Links.\allowbreak{}PrimRows3} (\texttt{Formalism.\allowbreak{}LocalToGlobal.\allowbreak{}IndefiniteReal.\allowbreak{}BaseLiftsPackageOfPrimitiveLinks}, lines 199--223)\\
\texttt{Formalism.\allowbreak{}L2.\allowbreak{}Rows.\allowbreak{}RowsHoldLam3} (\texttt{Formalism.\allowbreak{}LocalToGlobal.\allowbreak{}HilbertRowConditions.\allowbreak{}ColumnsFunctionAtSubspaces}, lines 129--140)\\
\texttt{Formalism.\allowbreak{}MainNotions.\allowbreak{}TuSize} (\texttt{Formalism.\allowbreak{}SharpRestriction.\allowbreak{}ParamsHeightRankConstant}, lines 180--183)\\
\texttt{Formalism.\allowbreak{}BasicNotions.\allowbreak{}SizeData₂.\allowbreak{}TuGe} (\texttt{Formalism.\allowbreak{}LocalRingCardinality.\allowbreak{}DataParams₂}, lines 252--253)\\
\texttt{Formalism.\allowbreak{}MainNotions.\allowbreak{}StandardLattice.\allowbreak{}ClassNumberFormula₂} (\texttt{Formalism.\allowbreak{}SharpRestriction.\allowbreak{}GlobalSchurBounds.\allowbreak{}StandardLattice}, lines 453--457)\\
\texttt{Formalism.\allowbreak{}L2.\allowbreak{}SG.\allowbreak{}JoinSI.\allowbreak{}IndefiniteRealisationJoin} (\texttt{Development.\allowbreak{}AutomorphicArithmetic.\allowbreak{}ThetaCorrespondence.\allowbreak{}ThetaPairLifts.\allowbreak{}ThirdPlaceSeries.\allowbreak{}ColJoinSOfSTiesR}, lines 67--81)\\
\texttt{Formalism.\allowbreak{}L2.\allowbreak{}SG.\allowbreak{}JoinSI.\allowbreak{}ColNormT} (\texttt{Formalism.\allowbreak{}LocalToGlobal.\allowbreak{}SiegelLatticeTheory.\allowbreak{}FlagStabilizerTensors.\allowbreak{}JoinSI}, lines 62--64)\\
\texttt{GlobalTheory.\allowbreak{}genusMassAt\_\allowbreak{}two} (\texttt{Development.\allowbreak{}AbelianVarietyTheory.\allowbreak{}Proofs.\allowbreak{}GlobalInterpolation.\allowbreak{}GenusMassAtTwo}, lines 7--9)\\
\texttt{Formalism.\allowbreak{}MainNotions.\allowbreak{}FamP.\allowbreak{}Large} (\texttt{Formalism.\allowbreak{}SharpRestriction.\allowbreak{}ParamsHeightRankConstant}, lines 138--141)\\
\texttt{GlobalTheory.\allowbreak{}Cp.\allowbreak{}SelectionServed} (\texttt{Development.\allowbreak{}AutomorphicArithmetic.\allowbreak{}ThetaCorrespondence.\allowbreak{}ThetaPairLifts.\allowbreak{}ThirdPlaceSeries.\allowbreak{}CountsIrreducibilityConstants}, lines 48--62)\\
\texttt{Formalism.\allowbreak{}L2.\allowbreak{}DIV.\allowbreak{}lPrj} (\texttt{Formalism.\allowbreak{}LocalToGlobal.\allowbreak{}Divisibility.\allowbreak{}RowSideTie}, lines 96--99)\\
\texttt{Formalism.\allowbreak{}L2.\allowbreak{}DIV.\allowbreak{}BlockDat.\allowbreak{}lPr} (\texttt{Formalism.\allowbreak{}LocalToGlobal.\allowbreak{}Divisibility.\allowbreak{}RowLat}, lines 148--150)\\
\texttt{Formalism.\allowbreak{}MainNotions.\allowbreak{}StandardLattice.\allowbreak{}VolTie} (\texttt{Formalism.\allowbreak{}SharpRestriction.\allowbreak{}GlobalSchurBounds.\allowbreak{}StandardLattice}, lines 258--264)\par}

{\footnotesize\raggedright\noindent \ref{B8.02c}.\quad \texttt{Formalism.\allowbreak{}MainNotions.\allowbreak{}S.\allowbreak{}StandingNumericalData} (\texttt{Formalism.\allowbreak{}SharpRestriction.\allowbreak{}CuspSpanNormBounds.\allowbreak{}StandingNumericalData}, lines 84--105)\\
\texttt{Formalism.\allowbreak{}MainNotions.\allowbreak{}S.\allowbreak{}StandingHypotheses} (\texttt{Formalism.\allowbreak{}SharpRestriction.\allowbreak{}CuspSpanNormBounds.\allowbreak{}StandingNumericalData}, lines 106--137)\\
\texttt{Formalism.\allowbreak{}MainNotions.\allowbreak{}S.\allowbreak{}Cnt.\allowbreak{}Hyp} (\texttt{Formalism.\allowbreak{}SharpRestriction.\allowbreak{}CuspSpanNormBounds.\allowbreak{}CntPartners}, lines 36--49)\par}

{\footnotesize\raggedright\noindent \ref{B8.02d}.\quad \texttt{Formalism.\allowbreak{}MainNotions.\allowbreak{}count\_\allowbreak{}hypotheses\_\allowbreak{}of\_\allowbreak{}sizes} (\texttt{Formalism.\allowbreak{}SharpRestriction.\allowbreak{}CuspSpanNormBounds.\allowbreak{}ThresholdCubicInequalitySignPattern}, lines 230--237)\\
\texttt{Formalism.\allowbreak{}MainNotions.\allowbreak{}count\_\allowbreak{}hypotheses\_\allowbreak{}diagonal\_\allowbreak{}words} (\texttt{Formalism.\allowbreak{}SharpRestriction.\allowbreak{}CuspSpanNormBounds.\allowbreak{}ThresholdCubicInequalitySignPattern}, lines 1254--1264)\par}

{\footnotesize\raggedright\noindent \ref{B8.02ba}.\quad \texttt{GlobalTheory.\allowbreak{}Mass.\allowbreak{}sublattice\_\allowbreak{}count} (\texttt{Development.\allowbreak{}AbelianVarietyTheory.\allowbreak{}Proofs.\allowbreak{}FlagMeasures.\allowbreak{}EncardSublatticesOfIndexLe}, lines 10--41)\\
\texttt{GlobalTheory.\allowbreak{}Mass.\allowbreak{}SpaceMass.\allowbreak{}main} (\texttt{Development.\allowbreak{}AbelianVarietyTheory.\allowbreak{}Proofs.\allowbreak{}FlagMeasures.\allowbreak{}MassIntegralLatticesDiscriminantOrderLe}, lines 79--157)\\
\texttt{GlobalTheory.\allowbreak{}Mass.\allowbreak{}MaxFormula.\allowbreak{}main} (\texttt{Development.\allowbreak{}AbelianVarietyTheory.\allowbreak{}Proofs.\allowbreak{}FlagMeasures.\allowbreak{}MaximalHermitianLatticeMassRatio}, lines 111--118)\\
\texttt{GlobalTheory.\allowbreak{}MaxLatStruct.\allowbreak{}main} (\texttt{Development.\allowbreak{}AbelianVarietyTheory.\allowbreak{}Proofs.\allowbreak{}FlagMeasures.\allowbreak{}MaximalHermitianLatticeLocalStructureDiscriminan}, lines 511--551)\\
\texttt{GlobalTheory.\allowbreak{}MaxLatLam.\allowbreak{}main} (\texttt{Development.\allowbreak{}AbelianVarietyTheory.\allowbreak{}Proofs.\allowbreak{}FlagMeasures.\allowbreak{}MaximalLatticeLocalFactorProductPosAndLe}, lines 25--36)\par}

{\footnotesize\raggedright\noindent \ref{B8.02bb}.\quad \texttt{GlobalTheory.\allowbreak{}Mass.\allowbreak{}PrimMean.\allowbreak{}main} (\texttt{Development.\allowbreak{}AbelianVarietyTheory.\allowbreak{}Proofs.\allowbreak{}FlagMeasures.\allowbreak{}PrimitiveSublatticeMeanCountBound}, lines 139--215)\\
\texttt{GlobalTheory.\allowbreak{}Mass.\allowbreak{}FewOrbits.\allowbreak{}main} (\texttt{Development.\allowbreak{}AbelianVarietyTheory.\allowbreak{}Proofs.\allowbreak{}FlagMeasures.\allowbreak{}FewOrbitsSubspacesHermitianLatticeDiscriminant}, lines 40--111)\\
\texttt{GlobalTheory.\allowbreak{}Mass.\allowbreak{}PrimDisc.\allowbreak{}main} (\texttt{Development.\allowbreak{}AbelianVarietyTheory.\allowbreak{}Proofs.\allowbreak{}FlagMeasures.\allowbreak{}PrimitiveHermitianSublatticeDiscriminantGenerato}, lines 150--191)\\
\texttt{GlobalTheory.\allowbreak{}Mass.\allowbreak{}Glue.\allowbreak{}main} (\texttt{Development.\allowbreak{}AbelianVarietyTheory.\allowbreak{}Proofs.\allowbreak{}FlagMeasures.\allowbreak{}CardHermitianLatticesGluedSubspaceComplement}, lines 210--277)\\
\texttt{GlobalTheory.\allowbreak{}Mass.\allowbreak{}FixedSub.\allowbreak{}orbMass\_\allowbreak{}stabilizer\_\allowbreak{}le} (\texttt{Development.\allowbreak{}AbelianVarietyTheory.\allowbreak{}Proofs.\allowbreak{}FlagMeasures.\allowbreak{}OrbitStabilizerMassSubspaceCyclicDiscriminantLe}, lines 100--211)\par}

{\footnotesize\raggedright\noindent \ref{B8.02bc}.\quad \texttt{GlobalTheory.\allowbreak{}SmallRankMass.\allowbreak{}mass\_\allowbreak{}small} (\texttt{Development.\allowbreak{}AbelianVarietyTheory.\allowbreak{}Proofs.\allowbreak{}FlagMeasures.\allowbreak{}ExistsBoundedBasis}, lines 205--255)\\
\texttt{GlobalTheory.\allowbreak{}ClassCount.\allowbreak{}exists\_\allowbreak{}finset\_\allowbreak{}classes} (\texttt{Development.\allowbreak{}AbelianVarietyTheory.\allowbreak{}Proofs.\allowbreak{}CuspFormEvaluations.\allowbreak{}ExistsFinsetSelfDualClassesHerm}, lines 399--442)\\
\texttt{GlobalTheory.\allowbreak{}EtaValueBound.\allowbreak{}lValEta\_\allowbreak{}bounds} (\texttt{Development.\allowbreak{}AbelianVarietyTheory.\allowbreak{}Proofs.\allowbreak{}FlagMeasures.\allowbreak{}LValEtaBounds}, lines 61--71)\\
\texttt{GlobalTheory.\allowbreak{}MassRatioBound.\allowbreak{}mass\_\allowbreak{}mul\_\allowbreak{}mass\_\allowbreak{}le} (\texttt{Development.\allowbreak{}AbelianVarietyTheory.\allowbreak{}Proofs.\allowbreak{}FlagMeasures.\allowbreak{}UnimodularHermitianMassMulLeMass}, lines 131--185)\par}

{\footnotesize\raggedright\noindent \ref{B8.02bd}.\quad \texttt{GlobalTheory.\allowbreak{}Mass.\allowbreak{}CyclicDisc.\allowbreak{}main} (\texttt{Development.\allowbreak{}AbelianVarietyTheory.\allowbreak{}Proofs.\allowbreak{}FlagMeasures.\allowbreak{}MassGenusCyclicDiscriminantSplitPrime}, lines 119--182)\\
\texttt{GlobalTheory.\allowbreak{}DegCnt.\allowbreak{}PrimGenus.\allowbreak{}main} (\texttt{Development.\allowbreak{}AbelianVarietyTheory.\allowbreak{}Proofs.\allowbreak{}FlagMeasures.\allowbreak{}MeanPrimitiveSublatticesGenusDefiniteLattice}, lines 158--198)\par}

{\footnotesize\raggedright\noindent \ref{B8.02be}.\quad \texttt{GlobalTheory.\allowbreak{}HermLat.\allowbreak{}IsModIso} (\texttt{Development.\allowbreak{}AbelianVarietyTheory.\allowbreak{}Definitions.\allowbreak{}SemilocalFrames.\allowbreak{}FlagPairDefinitePartLattice}, lines 115--120)\\
\texttt{GlobalTheory.\allowbreak{}DegCnt.\allowbreak{}IsDat} (\texttt{Development.\allowbreak{}AbelianVarietyTheory.\allowbreak{}Definitions.\allowbreak{}SemilocalFrames.\allowbreak{}NormalFormDualOrthogonalLast}, lines 127--131)\\
\texttt{GlobalTheory.\allowbreak{}DegCnt.\allowbreak{}Trip} (\texttt{Development.\allowbreak{}AbelianVarietyTheory.\allowbreak{}Definitions.\allowbreak{}SemilocalFrames.\allowbreak{}NormalFormDualOrthogonalLast}, lines 132--136)\par}

{\footnotesize\raggedright\noindent \ref{B8.02bf}.\quad \texttt{GlobalTheory.\allowbreak{}DegCnt.\allowbreak{}FixedSpaceCount.\allowbreak{}main} (\texttt{Development.\allowbreak{}AbelianVarietyTheory.\allowbreak{}Proofs.\allowbreak{}FlagMeasures.\allowbreak{}CountIsometryPairsWithGivenFixedSpace}, lines 232--313)\\
\texttt{GlobalTheory.\allowbreak{}DegCnt.\allowbreak{}FixDisc.\allowbreak{}main} (\texttt{Development.\allowbreak{}AbelianVarietyTheory.\allowbreak{}Proofs.\allowbreak{}FlagMeasures.\allowbreak{}DiscriminantOrderCommonFixedLatticeEmbeddedData}, lines 372--409)\\
\texttt{GlobalTheory.\allowbreak{}TraceCount.\allowbreak{}main} (\texttt{Development.\allowbreak{}AbelianVarietyTheory.\allowbreak{}Proofs.\allowbreak{}FlagMeasures.\allowbreak{}EncardSublatticesIndexFixedTraceLe}, lines 153--216)\\
\texttt{GlobalTheory.\allowbreak{}DegCnt.\allowbreak{}DatCount.\allowbreak{}main} (\texttt{Development.\allowbreak{}AbelianVarietyTheory.\allowbreak{}Proofs.\allowbreak{}FlagMeasures.\allowbreak{}CardEmbeddedDataWithGivenFixedLattice}, lines 207--300)\\
\texttt{GlobalTheory.\allowbreak{}IsoCount.\allowbreak{}isoCnt} (\texttt{Development.\allowbreak{}AbelianVarietyTheory.\allowbreak{}Proofs.\allowbreak{}FlagMeasures.\allowbreak{}CardIsometriesFixingSubspaceDefiniteLatticesLe}, lines 115--152)\\
\texttt{GlobalTheory.\allowbreak{}PairNoFix.\allowbreak{}main} (\texttt{Development.\allowbreak{}AbelianVarietyTheory.\allowbreak{}Proofs.\allowbreak{}FlagMeasures.\allowbreak{}PairsIsometryNoFixedVectorCountBound}, lines 58--111)\par}

{\footnotesize\raggedright\noindent \ref{B8.10k}.\quad \texttt{GlobalTheory.\allowbreak{}Amp.\allowbreak{}dyadic\_\allowbreak{}summation} (\texttt{Development.\allowbreak{}AbelianVarietyTheory.\allowbreak{}Proofs.\allowbreak{}FlagMeasures.\allowbreak{}DyadicSummationWithDecay}, lines 104--240)\par}

{\footnotesize\raggedright\noindent \ref{B8.02bg}.\quad \texttt{GlobalTheory.\allowbreak{}TwoSidedMomentsGap.\allowbreak{}moment\_\allowbreak{}of\_\allowbreak{}bessel\_\allowbreak{}gap} (\texttt{Development.\allowbreak{}AbelianVarietyTheory.\allowbreak{}Proofs.\allowbreak{}FlagMeasures.\allowbreak{}MajorantLengthTwoSidedMomentsLargeGap}, lines 57--128)\\
\texttt{GlobalTheory.\allowbreak{}TwoSidedMomentsGap.\allowbreak{}inv\_\allowbreak{}moment\_\allowbreak{}of\_\allowbreak{}bessel\_\allowbreak{}gap} (\texttt{Development.\allowbreak{}AbelianVarietyTheory.\allowbreak{}Proofs.\allowbreak{}FlagMeasures.\allowbreak{}MajorantLengthTwoSidedMomentsLargeGap}, lines 129--199)\\
\texttt{GlobalTheory.\allowbreak{}InvLenMomentGap.\allowbreak{}inv\_\allowbreak{}len\_\allowbreak{}moment\_\allowbreak{}gap} (\texttt{Development.\allowbreak{}AbelianVarietyTheory.\allowbreak{}Proofs.\allowbreak{}FlagMeasures.\allowbreak{}InverseMajorantLengthMomentLargeGap}, lines 24--77)\par}

{\footnotesize\raggedright\noindent \ref{B8.02bh}.\quad \texttt{GlobalTheory.\allowbreak{}DegCnt.\allowbreak{}Absorb.\allowbreak{}exists\_\allowbreak{}absorbing\_\allowbreak{}ideal} (\texttt{Development.\allowbreak{}AbelianVarietyTheory.\allowbreak{}Proofs.\allowbreak{}FlagMeasures.\allowbreak{}ExistsStarStableIdealAbsorbingMinors}, lines 471--599)\\
\texttt{GlobalTheory.\allowbreak{}DegCnt.\allowbreak{}Absorb.\allowbreak{}exists\_\allowbreak{}ideal} (\texttt{Development.\allowbreak{}AbelianVarietyTheory.\allowbreak{}Proofs.\allowbreak{}FlagMeasures.\allowbreak{}ExistsStarStableIdealAbsorbingMinors}, lines 355--403)\par}

{\footnotesize\raggedright\noindent \ref{B8.02b}.\quad \texttt{GlobalTheory.\allowbreak{}rowsPos\_\allowbreak{}typeP} (\texttt{Development.\allowbreak{}AbelianVarietyTheory.\allowbreak{}Proofs.\allowbreak{}FlagMeasures.\allowbreak{}RowsPosTypeP}, lines 9--11)\\
\texttt{Formalism.\allowbreak{}L2.\allowbreak{}SG.\allowbreak{}RowsPos.\allowbreak{}TypeP} (\texttt{Formalism.\allowbreak{}LocalToGlobal.\allowbreak{}SiegelLatticeTheory.\allowbreak{}FlagStabilizerTensors.\allowbreak{}RowsPos}, lines 18--23)\\
\texttt{Formalism.\allowbreak{}L2.\allowbreak{}SG.\allowbreak{}KappaOffDiagonal.\allowbreak{}TypeAt} (\texttt{Formalism.\allowbreak{}LocalToGlobal.\allowbreak{}SiegelLatticeTheory.\allowbreak{}CokernelGLCounts.\allowbreak{}RankPointCounts.\allowbreak{}DimFormC}, lines 40--46)\\
\texttt{GlobalTheory.\allowbreak{}TypeAsm.\allowbreak{}typeBound\_\allowbreak{}of\_\allowbreak{}amplitude\_\allowbreak{}pairs} (\texttt{Development.\allowbreak{}AbelianVarietyTheory.\allowbreak{}Proofs.\allowbreak{}FlagMeasures.\allowbreak{}TermAmplitudeFactorNonneg}, lines 22--67)\par}

{\footnotesize\raggedright\noindent \ref{B8.02yaa}.\quad \texttt{WlocSk.\allowbreak{}ANz.\allowbreak{}aNonzero\_\allowbreak{}of} (\texttt{Formalism.\allowbreak{}LocalToGlobal.\allowbreak{}SiegelLatticeTheory.\allowbreak{}HaarTransversals.\allowbreak{}UnitaryHaarProducts.\allowbreak{}LocalWeightCharts.\allowbreak{}ANonzeroOf}, lines 34--148)\\
\texttt{WlocSk.\allowbreak{}AQ.\allowbreak{}wave\_\allowbreak{}packet\_\allowbreak{}l2\_\allowbreak{}finite} (\texttt{Formalism.\allowbreak{}LocalToGlobal.\allowbreak{}SiegelLatticeTheory.\allowbreak{}HaarTransversals.\allowbreak{}UnitaryHaarProducts.\allowbreak{}LocalWeightCharts.\allowbreak{}WavePacketL1L2OfHypotheses}, lines 397--479)\\
\texttt{WlocSk.\allowbreak{}scWlocShape\_\allowbreak{}of} (\texttt{Formalism.\allowbreak{}LocalToGlobal.\allowbreak{}SiegelLatticeTheory.\allowbreak{}HaarTransversals.\allowbreak{}UnitaryHaarProducts.\allowbreak{}LocalWeightCharts.\allowbreak{}L2SGWfEqOne}, lines 87--125)\par}

{\footnotesize\raggedright\noindent \ref{B8.02yab}.\quad \texttt{Formalism.\allowbreak{}MainNotions.\allowbreak{}BridgeCS3dBOW.\allowbreak{}bow\_\allowbreak{}core} (\texttt{Formalism.\allowbreak{}LocalToGlobal.\allowbreak{}SiegelLatticeTheory.\allowbreak{}FlagStabilizerTensors.\allowbreak{}RootLatticeThresholds.\allowbreak{}BesselOfWPf}, lines 132--222)\\
\texttt{Formalism.\allowbreak{}FourierDecayAlt.\allowbreak{}ThreeDimData.\allowbreak{}DDForm} (\texttt{Formalism.\allowbreak{}HermitianCurvature.\allowbreak{}BesselNonvanishingReduction}, lines 60--75)\par}

{\footnotesize\raggedright\noindent \ref{B8.02ya}.\quad \texttt{ReciprocityTheory.\allowbreak{}ChartRestriction.\allowbreak{}DS.\allowbreak{}Lap.\allowbreak{}TopClosed.\allowbreak{}besselNVG} (\texttt{ReciprocityTheory.\allowbreak{}ReductiveGroups.\allowbreak{}DiscreteSeries.\allowbreak{}LaplaceOperators.\allowbreak{}BesselVectors.\allowbreak{}Closed}, lines 10--11)\\
\texttt{Formalism.\allowbreak{}L2.\allowbreak{}SG.\allowbreak{}OwnFacts.\allowbreak{}factBesselNVGap} (\texttt{Formalism.\allowbreak{}LocalToGlobal.\allowbreak{}SiegelLatticeTheory.\allowbreak{}WeightedGaugeFactors.\allowbreak{}FactsDS}, lines 18--19)\\
\texttt{Formalism.\allowbreak{}L2.\allowbreak{}SG.\allowbreak{}Pf.\allowbreak{}wuDecayOfGap\_\allowbreak{}of\_\allowbreak{}uniqGap} (\texttt{CompanionTheory.\allowbreak{}AutomorphicTheory.\allowbreak{}LocalBesselUniqueness.\allowbreak{}BesselModels.\allowbreak{}WuFinal}, lines 9--12)\\
\texttt{CompanionTheory.\allowbreak{}Lime.\allowbreak{}BM.\allowbreak{}uniqGap\_\allowbreak{}holds} (\texttt{CompanionTheory.\allowbreak{}AutomorphicTheory.\allowbreak{}LocalBesselUniqueness.\allowbreak{}BesselModels.\allowbreak{}MultiplicityOneUnitaryBranchingOfGapHyp}, lines 8--9)\\
\texttt{CompanionTheory.\allowbreak{}Lime.\allowbreak{}BM.\allowbreak{}wuDecayOfGap\_\allowbreak{}holds'} (\texttt{CompanionTheory.\allowbreak{}AutomorphicTheory.\allowbreak{}LocalBesselUniqueness.\allowbreak{}BesselModels.\allowbreak{}MultiplicityOneUnitaryBranchingOfGapHyp}, lines 12--13)\\
\texttt{Formalism.\allowbreak{}L2.\allowbreak{}SG.\allowbreak{}Pf.\allowbreak{}WuDec.\allowbreak{}LocDecayStep} (\texttt{Formalism.\allowbreak{}LocalToGlobal.\allowbreak{}SiegelLatticeTheory.\allowbreak{}FlagStabilizerTensors.\allowbreak{}RootLatticeThresholds.\allowbreak{}WuDecayAll}, lines 433--436)\par}

{\footnotesize\raggedright\noindent \ref{B8.02y}.\quad \texttt{GlobalTheory.\allowbreak{}wuDecayOfGap} (\texttt{Development.\allowbreak{}AbelianVarietyTheory.\allowbreak{}Proofs.\allowbreak{}FlagMeasures.\allowbreak{}WuDecayOfGap}, lines 7--8)\\
\texttt{Formalism.\allowbreak{}L2.\allowbreak{}SG.\allowbreak{}Pf.\allowbreak{}WuDec.\allowbreak{}wuDecayOfGap\_\allowbreak{}of\_\allowbreak{}steps} (\texttt{Formalism.\allowbreak{}LocalToGlobal.\allowbreak{}SiegelLatticeTheory.\allowbreak{}FlagStabilizerTensors.\allowbreak{}RootLatticeThresholds.\allowbreak{}WuDecayAll}, lines 727--734)\par}

{\footnotesize\raggedright\noindent \ref{B8.02a}.\quad \texttt{GlobalTheory.\allowbreak{}Fc.\allowbreak{}sevenDuties\_\allowbreak{}of\_\allowbreak{}class} (\texttt{Development.\allowbreak{}AbelianVarietyTheory.\allowbreak{}Definitions.\allowbreak{}GlobalTwisting.\allowbreak{}ValidForConstructedPointClassesEventuallyWorking}, lines 245--246)\\
\texttt{GlobalTheory.\allowbreak{}EX.\allowbreak{}DQ.\allowbreak{}duty\_\allowbreak{}seven\_\allowbreak{}library} (\texttt{Development.\allowbreak{}AbelianVarietyTheory.\allowbreak{}Proofs.\allowbreak{}GlobalInterpolation.\allowbreak{}DutySevenLibrary}, lines 15--23)\par}

{\footnotesize\raggedright\noindent \ref{B8.02a}.\quad \texttt{GlobalTheory.\allowbreak{}EX.\allowbreak{}DC.\allowbreak{}cls} (\texttt{Development.\allowbreak{}AbelianVarietyTheory.\allowbreak{}Definitions.\allowbreak{}GlobalTwisting.\allowbreak{}AuxiliaryPrimeMemAvoidedPrimes}, lines 106--107)\\
\texttt{GlobalTheory.\allowbreak{}EX.\allowbreak{}DQ.\allowbreak{}weightAdmissible\_\allowbreak{}smallShares\_\allowbreak{}at\_\allowbreak{}point} (\texttt{Development.\allowbreak{}AbelianVarietyTheory.\allowbreak{}Proofs.\allowbreak{}GlobalInterpolation.\allowbreak{}WeightAdmissibleSmallSharesAtPoint}, lines 23--29)\par}

{\footnotesize\raggedright\noindent \ref{B8.02ac}.\quad \texttt{GlobalTheory.\allowbreak{}EX.\allowbreak{}sig\_\allowbreak{}eq\_\allowbreak{}loc} (\texttt{Development.\allowbreak{}AbelianVarietyTheory.\allowbreak{}Proofs.\allowbreak{}GlobalInterpolation.\allowbreak{}WeightRequirementOfValid}, lines 20--30)\\
\texttt{GlobalTheory.\allowbreak{}EX.\allowbreak{}weightWindowRequirement\_\allowbreak{}host\_\allowbreak{}of\_\allowbreak{}point} (\texttt{Development.\allowbreak{}AbelianVarietyTheory.\allowbreak{}Proofs.\allowbreak{}GlobalInterpolation.\allowbreak{}WeightWindowRequirementHostOfPoint}, lines 17--54)\\
\texttt{GlobalTheory.\allowbreak{}EX.\allowbreak{}wreq\_\allowbreak{}of\_\allowbreak{}wok} (\texttt{Development.\allowbreak{}AbelianVarietyTheory.\allowbreak{}Proofs.\allowbreak{}GlobalInterpolation.\allowbreak{}WeightRequirementOfValid}, lines 32--69)\par}

{\footnotesize\raggedright\noindent \ref{B8.02ac}.\quad \texttt{UpperTheory.\allowbreak{}Top.\allowbreak{}hostWW} (\texttt{Development.\allowbreak{}AutomorphicArithmetic.\allowbreak{}CharacterBoxes.\allowbreak{}HostWW}, lines 27--87)\\
\texttt{UpperTheory.\allowbreak{}Top.\allowbreak{}kitS} (\texttt{Development.\allowbreak{}AutomorphicArithmetic.\allowbreak{}CharacterBoxes.\allowbreak{}KitS}, lines 14--62)\\
\texttt{UpperTheory.\allowbreak{}Top.\allowbreak{}designS} (\texttt{Development.\allowbreak{}AutomorphicArithmetic.\allowbreak{}CharacterBoxes.\allowbreak{}DesignS}, lines 21--98)\par}

{\footnotesize\raggedright\noindent \ref{B8.03a}.\quad \texttt{GlobalTheory.\allowbreak{}rowsFinDim\_\allowbreak{}link} (\texttt{Development.\allowbreak{}AbelianVarietyTheory.\allowbreak{}Proofs.\allowbreak{}FlagMeasures.\allowbreak{}RowsFinDimLink}, lines 7--8)\\
\texttt{CompanionTheory.\allowbreak{}D.\allowbreak{}gk.\allowbreak{}RF.\allowbreak{}automorphic\_\allowbreak{}forms\_\allowbreak{}finite\_\allowbreak{}dimensional} (\texttt{CompanionTheory.\allowbreak{}HarishChandraModules.\allowbreak{}KFiniteAnalyticVectors.\allowbreak{}FundamentalDomainCusps.\allowbreak{}AutomorphicFormsFiniteDimensional}, lines 25--27)\par}

{\footnotesize\raggedright\noindent \ref{B8.03}.\quad \texttt{GlobalTheory.\allowbreak{}Fc.\allowbreak{}S0.\allowbreak{}StageU.\allowbreak{}hypotheses\_\allowbreak{}nonempty\_\allowbreak{}of\_\allowbreak{}unramified\_\allowbreak{}exponents\_\allowbreak{}alt} (\texttt{Development.\allowbreak{}AutomorphicArithmetic.\allowbreak{}ThetaCorrespondence.\allowbreak{}ThetaPairLifts.\allowbreak{}ThirdPlaceSeries.\allowbreak{}HypothesesNonemptyOfUnramifiedExponents}, lines 41--61)\par}

{\footnotesize\raggedright\noindent \ref{B8.04}.\quad \texttt{GlobalTheory.\allowbreak{}Fc.\allowbreak{}nonempty\_\allowbreak{}hyps\_\allowbreak{}of\_\allowbreak{}allRows\_\allowbreak{}alt} (\texttt{Development.\allowbreak{}AutomorphicArithmetic.\allowbreak{}ThetaCorrespondence.\allowbreak{}ThetaPairLifts.\allowbreak{}ThirdPlaceSeries.\allowbreak{}NonemptyHypsOfAllRows}, lines 19--101)\par}

{\footnotesize\raggedright\noindent \ref{B8.04}.\quad \texttt{GlobalTheory.\allowbreak{}EX.\allowbreak{}DQ.\allowbreak{}cls7} (\texttt{Development.\allowbreak{}AbelianVarietyTheory.\allowbreak{}Definitions.\allowbreak{}GlobalTwisting.\allowbreak{}FamilyClassSevenDutiesAndStableLattices}, lines 53--54)\\
\texttt{GlobalTheory.\allowbreak{}EX.\allowbreak{}DQ.\allowbreak{}QD} (\texttt{Development.\allowbreak{}AbelianVarietyTheory.\allowbreak{}Definitions.\allowbreak{}GlobalTwisting.\allowbreak{}FamilyClassSevenDutiesAndStableLattices}, lines 32--33)\\
\texttt{GlobalTheory.\allowbreak{}EX.\allowbreak{}DC.\allowbreak{}cls} (\texttt{Development.\allowbreak{}AbelianVarietyTheory.\allowbreak{}Definitions.\allowbreak{}GlobalTwisting.\allowbreak{}AuxiliaryPrimeMemAvoidedPrimes}, lines 106--108)\\
\texttt{GlobalTheory.\allowbreak{}EX.\allowbreak{}DC.\allowbreak{}withObj} (\texttt{Development.\allowbreak{}AbelianVarietyTheory.\allowbreak{}Definitions.\allowbreak{}GlobalTwisting.\allowbreak{}AuxiliaryPrimeMemAvoidedPrimes}, lines 104--105)\\
\texttt{GlobalTheory.\allowbreak{}EX.\allowbreak{}withDuty} (\texttt{Development.\allowbreak{}AbelianVarietyTheory.\allowbreak{}Definitions.\allowbreak{}GlobalTwisting.\allowbreak{}ThinScalesIffHostBound}, lines 244--249)\\
\texttt{Formalism.\allowbreak{}MainNotions.\allowbreak{}DGamChoice\_\allowbreak{}at} (\texttt{Formalism.\allowbreak{}SharpRestriction.\allowbreak{}GlobalSchurBounds.\allowbreak{}MainIn}, lines 57--59)\\
\texttt{GlobalTheory.\allowbreak{}IndefiniteRowSideFullHecke.\allowbreak{}TameGeneratorWindow.\allowbreak{}lJ} (\texttt{Development.\allowbreak{}AbelianVarietyTheory.\allowbreak{}BlockTheory.\allowbreak{}GlobalSignature.\allowbreak{}TameGeneratorWindow}, lines 50--52)\par}

{\footnotesize\raggedright\noindent \ref{B8.05a}.\quad \texttt{GlobalTheory.\allowbreak{}Mem.\allowbreak{}Top.\allowbreak{}DsP} (\texttt{Development.\allowbreak{}AbelianVarietyTheory.\allowbreak{}Definitions.\allowbreak{}AutomorphicMembership.\allowbreak{}CmQuadraticFormEqAlternatingForm}, lines 23--27)\\
\texttt{GlobalTheory.\allowbreak{}Mem.\allowbreak{}Top.\allowbreak{}DsLow.\allowbreak{}trueDsig\_\allowbreak{}eq\_\allowbreak{}dimV} (\texttt{Development.\allowbreak{}AbelianVarietyTheory.\allowbreak{}Proofs.\allowbreak{}IwasawaDecomposition.\allowbreak{}DsPPosAndLower}, lines 15--44)\\
\texttt{GlobalTheory.\allowbreak{}EX.\allowbreak{}WinLow.\allowbreak{}window\_\allowbreak{}lower} (\texttt{Development.\allowbreak{}AbelianVarietyTheory.\allowbreak{}Proofs.\allowbreak{}GlobalInterpolation.\allowbreak{}WindowLower}, lines 177--295)\par}

{\footnotesize\raggedright\noindent \ref{B8.05}.\quad \texttt{UpperTheory.\allowbreak{}Top.\allowbreak{}classSel} (\texttt{Development.\allowbreak{}AutomorphicArithmetic.\allowbreak{}CharacterBoxes.\allowbreak{}ClassSel}, lines 14--29)\par}

{\footnotesize\raggedright\noindent \ref{B8.05}.\quad \texttt{GlobalTheory.\allowbreak{}Fc.\allowbreak{}RowL.\allowbreak{}and} (\texttt{Development.\allowbreak{}AbelianVarietyTheory.\allowbreak{}Definitions.\allowbreak{}GlobalTwisting.\allowbreak{}EventuallyAllScaleContextsAndOfBoth}, lines 76--81)\\
\texttt{GlobalTheory.\allowbreak{}Fc.\allowbreak{}Row.\allowbreak{}elim} (\texttt{Development.\allowbreak{}AbelianVarietyTheory.\allowbreak{}Definitions.\allowbreak{}GlobalTwisting.\allowbreak{}ValidForConstructedPointClassesEventuallyWorking}, lines 364--367)\par}

{\footnotesize\raggedright\noindent \ref{lem:p8c1-row-forms}.\quad \texttt{GlobalTheory.\allowbreak{}T.\allowbreak{}rowFormsMatchPeriodMap\_\allowbreak{}of\_\allowbreak{}furtherProperties} (\texttt{Development.\allowbreak{}AbelianVarietyTheory.\allowbreak{}Proofs.\allowbreak{}SatakeParameters.\allowbreak{}RowFormsMatchPeriodMapOfFurtherPropertiesStatement}, lines 7--9)\par}

{\footnotesize\raggedright\noindent \ref{lem:p8c1-row-forms}.\quad \texttt{GlobalTheory.\allowbreak{}AtG.\allowbreak{}RX.\allowbreak{}lat\_\allowbreak{}of\_\allowbreak{}rx} (\texttt{Development.\allowbreak{}AbelianVarietyTheory.\allowbreak{}Proofs.\allowbreak{}SatakeParameters.\allowbreak{}RowFormsMatchPeriodMapOfFurtherProperties}, lines 71--102)\\
\texttt{GlobalTheory.\allowbreak{}AtG.\allowbreak{}RX.\allowbreak{}rowPin\_\allowbreak{}core} (\texttt{Development.\allowbreak{}AbelianVarietyTheory.\allowbreak{}Proofs.\allowbreak{}SatakeParameters.\allowbreak{}RowFormsMatchPeriodMapOfFurtherProperties}, lines 111--128)\\
\texttt{GlobalTheory.\allowbreak{}AtG.\allowbreak{}RX.\allowbreak{}rowPin\_\allowbreak{}of\_\allowbreak{}rx} (\texttt{Development.\allowbreak{}AbelianVarietyTheory.\allowbreak{}Proofs.\allowbreak{}SatakeParameters.\allowbreak{}RowFormsMatchPeriodMapOfFurtherProperties}, lines 131--141)\par}

{\footnotesize\raggedright\noindent \ref{B8.06}.\quad \texttt{GlobalTheory.\allowbreak{}Fc.\allowbreak{}AllAtMemberStage0.\allowbreak{}allRows\_\allowbreak{}of\_\allowbreak{}open0\_\allowbreak{}alt} (\texttt{Development.\allowbreak{}AutomorphicArithmetic.\allowbreak{}ThetaCorrespondence.\allowbreak{}ThetaPairLifts.\allowbreak{}ThirdPlaceSeries.\allowbreak{}LinkROfIsLinked}, lines 50--86)\\
\texttt{GlobalTheory.\allowbreak{}Fc.\allowbreak{}AllRowsAlt} (\texttt{Development.\allowbreak{}AutomorphicArithmetic.\allowbreak{}ThetaCorrespondence.\allowbreak{}ThetaPairLifts.\allowbreak{}ThirdPlaceSeries.\allowbreak{}AllRows}, lines 12--39)\par}

{\footnotesize\raggedright\noindent \ref{B8.07a}.\quad \texttt{GlobalTheory.\allowbreak{}Mem.\allowbreak{}R3.\allowbreak{}Rw.\allowbreak{}Bare.\allowbreak{}bare\_\allowbreak{}form\_\allowbreak{}saturated} (\texttt{Development.\allowbreak{}AutomorphicArithmetic.\allowbreak{}ThetaCorrespondence.\allowbreak{}ThetaPairLifts.\allowbreak{}LocalTypeDictionary.\allowbreak{}ExistsBareFormOfSaturation}, lines 19--57)\par}

{\footnotesize\raggedright\noindent \ref{B8.07b}.\quad \texttt{GlobalTheory.\allowbreak{}OrderEval.\allowbreak{}order\_\allowbreak{}eval} (\texttt{Development.\allowbreak{}AbelianVarietyTheory.\allowbreak{}Proofs.\allowbreak{}InertiaEigenvalues.\allowbreak{}ExistsLevelOrderExponentsSmallValuationDetermina}, lines 278--377)\\
\texttt{GlobalTheory.\allowbreak{}TwistChoice.\allowbreak{}exists\_\allowbreak{}level} (\texttt{Development.\allowbreak{}AbelianVarietyTheory.\allowbreak{}Proofs.\allowbreak{}InertiaEigenvalues.\allowbreak{}IsAuxiliaryPlaceFamilyAtLevelMono}, lines 321--354)\par}

{\footnotesize\raggedright\noindent \ref{B8.07c}.\quad \texttt{GlobalTheory.\allowbreak{}SelDetLength.\allowbreak{}twisted\_\allowbreak{}selmer\_\allowbreak{}lengths\_\allowbreak{}eq\_\allowbreak{}determinant\_\allowbreak{}valuation} (\texttt{Development.\allowbreak{}AbelianVarietyTheory.\allowbreak{}Proofs.\allowbreak{}InertiaEigenvalues.\allowbreak{}TwistedSelmerLengthsEqDeterminantValuation}, lines 375--386)\\
\texttt{GlobalTheory.\allowbreak{}TwistChoiceAll.\allowbreak{}admissibleIndexBoundAllPowers\_\allowbreak{}of\_\allowbreak{}conditions} (\texttt{Development.\allowbreak{}AbelianVarietyTheory.\allowbreak{}Proofs.\allowbreak{}InertiaEigenvalues.\allowbreak{}AdmissibleIndexBoundAllPowersOfConditions}, lines 92--223)\par}

{\footnotesize\raggedright\noindent \ref{B8.07d}.\quad \texttt{GlobalTheory.\allowbreak{}Mem.\allowbreak{}Top.\allowbreak{}ExportAt} (\texttt{Development.\allowbreak{}AbelianVarietyTheory.\allowbreak{}Definitions.\allowbreak{}AutomorphicMembership.\allowbreak{}AmbientBaseHeq}, lines 36--41)\\
\texttt{GlobalTheory.\allowbreak{}TwAll.\allowbreak{}goodTwistExistsAllCoprimePowers} (\texttt{Development.\allowbreak{}AbelianVarietyTheory.\allowbreak{}Definitions.\allowbreak{}NonGaloisDescent.\allowbreak{}AdmissibleIndexBoundAllPowers}, lines 23--30)\\
\texttt{Formalism.\allowbreak{}MainNotions.\allowbreak{}BlochKatoReduction.\allowbreak{}Gen.\allowbreak{}SetHyp} (\texttt{Formalism.\allowbreak{}SharpRestriction.\allowbreak{}LevelLetterCocycles.\allowbreak{}BlochKatoPhases}, lines 327--330)\\
\texttt{Formalism.\allowbreak{}MainNotions.\allowbreak{}IntegralHeckeForms.\allowbreak{}GoodTw} (\texttt{Formalism.\allowbreak{}SharpRestriction.\allowbreak{}LevelLetterCocycles.\allowbreak{}IntegralHeckeForms}, lines 219--230)\par}

{\footnotesize\raggedright\noindent \ref{B8.07e}.\quad \texttt{GlobalTheory.\allowbreak{}TripleFrobDual.\allowbreak{}frobDual\_\allowbreak{}of\_\allowbreak{}tripleAtt\_\allowbreak{}n} (\texttt{Development.\allowbreak{}AutomorphicArithmetic.\allowbreak{}Foundations.\allowbreak{}GaloisRepresentations.\allowbreak{}StringRepresentations.\allowbreak{}FrobDualOfTripleAttN}, lines 139--183)\\
\texttt{GlobalTheory.\allowbreak{}LatticeSelfDual.\allowbreak{}selfDual\_\allowbreak{}of\_\allowbreak{}charDual} (\texttt{Development.\allowbreak{}AbelianVarietyTheory.\allowbreak{}Proofs.\allowbreak{}InertiaEigenvalues.\allowbreak{}SelfDualLatticeOfConjDualCharpoly}, lines 76--88)\\
\texttt{GlobalTheory.\allowbreak{}ChooseForms.\allowbreak{}nu\_\allowbreak{}pos\_\allowbreak{}of\_\allowbreak{}resIso} (\texttt{Development.\allowbreak{}AbelianVarietyTheory.\allowbreak{}Proofs.\allowbreak{}InertiaEigenvalues.\allowbreak{}ExistsElementsSelmerFormsSmallCorank}, lines 410--425)\\
\texttt{UpperTheory.\allowbreak{}KStar.\allowbreak{}exists\_\allowbreak{}eps\_\allowbreak{}eq\_\allowbreak{}pow\_\allowbreak{}nu} (\texttt{Development.\allowbreak{}AutomorphicArithmetic.\allowbreak{}CyclotomicBounds.\allowbreak{}CyclotomicOrders.\allowbreak{}ExistsCycZPowNeOne}, lines 21--29)\par}

{\footnotesize\raggedright\noindent \ref{B8.07f}.\quad \texttt{GlobalTheory.\allowbreak{}AdmChars.\allowbreak{}exists\_\allowbreak{}ray\_\allowbreak{}chars} (\texttt{Development.\allowbreak{}AbelianVarietyTheory.\allowbreak{}Proofs.\allowbreak{}InertiaEigenvalues.\allowbreak{}NotValuationNatCastLtOne}, lines 329--386)\\
\texttt{GlobalTheory.\allowbreak{}AdmChars.\allowbreak{}exists\_\allowbreak{}admissible\_\allowbreak{}family} (\texttt{Development.\allowbreak{}AbelianVarietyTheory.\allowbreak{}Proofs.\allowbreak{}InertiaEigenvalues.\allowbreak{}NotValuationNatCastLtOne}, lines 453--467)\\
\texttt{GlobalTheory.\allowbreak{}ChosenPrimes.\allowbreak{}stepAB} (\texttt{Development.\allowbreak{}AbelianVarietyTheory.\allowbreak{}Proofs.\allowbreak{}InertiaEigenvalues.\allowbreak{}ExistsAuxiliaryPrimesPrescribedFrobenius}, lines 274--419)\par}

{\footnotesize\raggedright\noindent \ref{B8.07g}.\quad \texttt{GlobalTheory.\allowbreak{}SelFun.\allowbreak{}oneMatrix\_\allowbreak{}of\_\allowbreak{}galoisCohomology} (\texttt{Development.\allowbreak{}AbelianVarietyTheory.\allowbreak{}Proofs.\allowbreak{}InertiaEigenvalues.\allowbreak{}SelmerPresentationMatrixOfGaloisCohomology}, lines 20--35)\\
\texttt{GlobalTheory.\allowbreak{}LocalConditionCount.\allowbreak{}locCount} (\texttt{Development.\allowbreak{}AbelianVarietyTheory.\allowbreak{}Proofs.\allowbreak{}InertiaEigenvalues.\allowbreak{}LocalCountOfSetup}, lines 491--516)\\
\texttt{GlobalTheory.\allowbreak{}SelObstruction.\allowbreak{}selObs} (\texttt{Development.\allowbreak{}AbelianVarietyTheory.\allowbreak{}Proofs.\allowbreak{}InertiaEigenvalues.\allowbreak{}SelmerObstructionLeResidualSelmerSimpleExtension}, lines 521--561)\\
\texttt{GlobalTheory.\allowbreak{}KRep.\allowbreak{}kernelRep} (\texttt{Development.\allowbreak{}AbelianVarietyTheory.\allowbreak{}Proofs.\allowbreak{}InertiaEigenvalues.\allowbreak{}TorsionLeftExactFunctorKernelSquareMatrix}, lines 340--465)\par}

{\footnotesize\raggedright\noindent \ref{B8.07h}.\quad \texttt{GlobalTheory.\allowbreak{}ResOrder.\allowbreak{}core} (\texttt{Development.\allowbreak{}AbelianVarietyTheory.\allowbreak{}Proofs.\allowbreak{}InertiaEigenvalues.\allowbreak{}ResidualSelmerDeterminantOrderBound}, lines 118--322)\\
\texttt{GlobalTheory.\allowbreak{}ResOrder.\allowbreak{}resBound} (\texttt{Development.\allowbreak{}AbelianVarietyTheory.\allowbreak{}Proofs.\allowbreak{}InertiaEigenvalues.\allowbreak{}ResidualSelmerDeterminantOrderBound}, lines 324--350)\\
\texttt{GlobalTheory.\allowbreak{}ResRecip.\allowbreak{}reciprocity} (\texttt{Development.\allowbreak{}AbelianVarietyTheory.\allowbreak{}Proofs.\allowbreak{}InertiaEigenvalues.\allowbreak{}ReciprocityAuxiliaryPlacesSecondHalf}, lines 326--429)\\
\texttt{GlobalTheory.\allowbreak{}ResSel.\allowbreak{}exists\_\allowbreak{}unit\_\allowbreak{}coeff\_\allowbreak{}of\_\allowbreak{}not\_\allowbreak{}dvd\_\allowbreak{}line} (\texttt{Development.\allowbreak{}AbelianVarietyTheory.\allowbreak{}Proofs.\allowbreak{}InertiaEigenvalues.\allowbreak{}ExistsUnitCoeffOfNotXPowDvdLine}, lines 41--79)\par}

{\footnotesize\raggedright\noindent \ref{B8.07i}.\quad \texttt{KStar.\allowbreak{}Tree.\allowbreak{}twist\_\allowbreak{}mul\_\allowbreak{}coord\_\allowbreak{}eq\_\allowbreak{}zero} (\texttt{Development.\allowbreak{}AutomorphicArithmetic.\allowbreak{}CyclotomicBounds.\allowbreak{}ClassNumberBounds.\allowbreak{}CyclotomicTotient.\allowbreak{}PolarTraceMulVecMulVecEq}, lines 28--89)\\
\texttt{KStar.\allowbreak{}Tree.\allowbreak{}twist\_\allowbreak{}smul\_\allowbreak{}eq\_\allowbreak{}zero} (\texttt{Development.\allowbreak{}AutomorphicArithmetic.\allowbreak{}CyclotomicBounds.\allowbreak{}ClassNumberBounds.\allowbreak{}CyclotomicTotient.\allowbreak{}TwistTrace}, lines 78--111)\\
\texttt{GlobalTheory.\allowbreak{}NewPrime.\allowbreak{}FormVanish.\allowbreak{}two\_\allowbreak{}mul\_\allowbreak{}values\_\allowbreak{}eq\_\allowbreak{}zero} (\texttt{Development.\allowbreak{}AbelianVarietyTheory.\allowbreak{}Proofs.\allowbreak{}InertiaEigenvalues.\allowbreak{}TwoMulCocycleValuesZeroOnSmallSubgroup}, lines 49--123)\par}

{\footnotesize\raggedright\noindent \ref{B8.07j}.\quad \texttt{GlobalTheory.\allowbreak{}NewPrime.\allowbreak{}LocalRing.\allowbreak{}exists\_\allowbreak{}finite\_\allowbreak{}local\_\allowbreak{}ring} (\texttt{Development.\allowbreak{}AbelianVarietyTheory.\allowbreak{}Proofs.\allowbreak{}InertiaEigenvalues.\allowbreak{}ExistsFiniteLocalRingPolyDividesVarPow}, lines 282--430)\par}

{\footnotesize\raggedright\noindent \ref{B8.07k}.\quad \texttt{GlobalTheory.\allowbreak{}TwGalInputs.\allowbreak{}galois\_\allowbreak{}inputs\_\allowbreak{}alt} (\texttt{Development.\allowbreak{}AutomorphicArithmetic.\allowbreak{}ThetaCorrespondence.\allowbreak{}ThetaPairLifts.\allowbreak{}ThirdPlaceSeries.\allowbreak{}GaloisInputs}, lines 19--29)\\
\texttt{GlobalTheory.\allowbreak{}TripleLocal.\allowbreak{}triple\_\allowbreak{}att\_\allowbreak{}n\_\allowbreak{}of\_\allowbreak{}is\_\allowbreak{}triple\_\allowbreak{}with} (\texttt{Development.\allowbreak{}AutomorphicArithmetic.\allowbreak{}ThetaCorrespondence.\allowbreak{}ThetaPairLifts.\allowbreak{}ThirdPlaceSeries.\allowbreak{}TripleAttOfIsTripleWith}, lines 54--116)\\
\texttt{GlobalTheory.\allowbreak{}TripleAdm.\allowbreak{}irredOn\_\allowbreak{}rbar\_\allowbreak{}and\_\allowbreak{}inertia\_\allowbreak{}of\_\allowbreak{}residual} (\texttt{Development.\allowbreak{}AutomorphicArithmetic.\allowbreak{}ThetaCorrespondence.\allowbreak{}ThetaPairLifts.\allowbreak{}ThirdPlaceSeries.\allowbreak{}SGalRead}, lines 199--246)\\
\texttt{GlobalTheory.\allowbreak{}TripleGal.\allowbreak{}adm1\_\allowbreak{}adm8\_\allowbreak{}selfDual\_\allowbreak{}noJump\_\allowbreak{}of\_\allowbreak{}tripleAttN} (\texttt{Development.\allowbreak{}AutomorphicArithmetic.\allowbreak{}ThetaCorrespondence.\allowbreak{}ThetaPairLifts.\allowbreak{}ThirdPlaceSeries.\allowbreak{}GaloisHypothesesOfAttached}, lines 23--34)\\
\texttt{GlobalTheory.\allowbreak{}SettingHypotheses.\allowbreak{}setHyp\_\allowbreak{}of\_\allowbreak{}isTriple\_\allowbreak{}of\_\allowbreak{}nonzeroIndexPlace} (\texttt{Development.\allowbreak{}AutomorphicArithmetic.\allowbreak{}ThetaCorrespondence.\allowbreak{}ThetaPairLifts.\allowbreak{}ThirdPlaceSeries.\allowbreak{}SetHypOfIsTripleOfGalHodge}, lines 14--24)\par}

{\footnotesize\raggedright\noindent \ref{B8.07l}.\quad \texttt{GlobalTheory.\allowbreak{}ResDual.\allowbreak{}eqDim} (\texttt{Development.\allowbreak{}AbelianVarietyTheory.\allowbreak{}Proofs.\allowbreak{}InertiaEigenvalues.\allowbreak{}ResidualSelmerFinrankEqDual}, lines 257--271)\\
\texttt{GlobalTheory.\allowbreak{}ResSel.\allowbreak{}SpanOddPf.\allowbreak{}spanOdd} (\texttt{Development.\allowbreak{}AbelianVarietyTheory.\allowbreak{}Proofs.\allowbreak{}InertiaEigenvalues.\allowbreak{}CocycleFormsSpanOddResidueChar}, lines 231--361)\\
\texttt{GlobalTheory.\allowbreak{}ResSel.\allowbreak{}AltSpan.\allowbreak{}spanTwo\_\allowbreak{}of\_\allowbreak{}galHyp} (\texttt{Development.\allowbreak{}AbelianVarietyTheory.\allowbreak{}Proofs.\allowbreak{}InertiaEigenvalues.\allowbreak{}SpanTwoOfGaloisHypothesis}, lines 196--274)\\
\texttt{GlobalTheory.\allowbreak{}ChooseForms.\allowbreak{}single\_\allowbreak{}field} (\texttt{Development.\allowbreak{}AbelianVarietyTheory.\allowbreak{}Proofs.\allowbreak{}InertiaEigenvalues.\allowbreak{}ExistsElementsSelmerFormsSmallCorank}, lines 427--465)\\
\texttt{GlobalTheory.\allowbreak{}ChooseForms.\allowbreak{}chooseH} (\texttt{Development.\allowbreak{}AbelianVarietyTheory.\allowbreak{}Proofs.\allowbreak{}InertiaEigenvalues.\allowbreak{}ExistsElementsSelmerFormsSmallCorank}, lines 557--577)\par}

{\footnotesize\raggedright\noindent \ref{B8.07m}.\quad \texttt{GlobalTheory.\allowbreak{}Growth.\allowbreak{}compare} (\texttt{Development.\allowbreak{}AbelianVarietyTheory.\allowbreak{}Proofs.\allowbreak{}InertiaEigenvalues.\allowbreak{}ChebotarevInductionGrowBoundOfPresentation}, lines 315--453)\\
\texttt{GlobalTheory.\allowbreak{}Growth.\allowbreak{}grow} (\texttt{Development.\allowbreak{}AbelianVarietyTheory.\allowbreak{}Proofs.\allowbreak{}InertiaEigenvalues.\allowbreak{}ChebotarevInductionGrowBoundOfPresentation}, lines 462--503)\\
\texttt{GlobalTheory.\allowbreak{}OrdMul.\allowbreak{}coeff\_\allowbreak{}mem\_\allowbreak{}of\_\allowbreak{}pow\_\allowbreak{}mul} (\texttt{Development.\allowbreak{}AbelianVarietyTheory.\allowbreak{}Proofs.\allowbreak{}InertiaEigenvalues.\allowbreak{}CoeffMemOfPowMul}, lines 81--93)\par}

{\footnotesize\raggedright\noindent \ref{B8.07n}.\quad \texttt{GlobalTheory.\allowbreak{}NewPrime.\allowbreak{}step} (\texttt{Development.\allowbreak{}AbelianVarietyTheory.\allowbreak{}Proofs.\allowbreak{}InertiaEigenvalues.\allowbreak{}AntiOfKer}, lines 153--353)\\
\texttt{GlobalTheory.\allowbreak{}NewPrime.\allowbreak{}DualClass.\allowbreak{}exists\_\allowbreak{}not\_\allowbreak{}cob} (\texttt{Development.\allowbreak{}AbelianVarietyTheory.\allowbreak{}Proofs.\allowbreak{}InertiaEigenvalues.\allowbreak{}ExistsDualSelmerClassNotKilled}, lines 169--198)\\
\texttt{GlobalTheory.\allowbreak{}DualSub.\allowbreak{}dualSel\_\allowbreak{}of\_\allowbreak{}smul\_\allowbreak{}cob} (\texttt{Development.\allowbreak{}AbelianVarietyTheory.\allowbreak{}Proofs.\allowbreak{}InertiaEigenvalues.\allowbreak{}DualSelmerClassFromCoeffsKilledByMul}, lines 48--125)\\
\texttt{GlobalTheory.\allowbreak{}NewPrime.\allowbreak{}FormNonzero.\allowbreak{}exists\_\allowbreak{}bform\_\allowbreak{}ne\_\allowbreak{}zero} (\texttt{Development.\allowbreak{}AbelianVarietyTheory.\allowbreak{}Proofs.\allowbreak{}InertiaEigenvalues.\allowbreak{}ExistsSmallSubgroupCocycleFormNeZero}, lines 39--96)\par}

{\footnotesize\raggedright\noindent \ref{B8.07o}.\quad \texttt{GlobalTheory.\allowbreak{}NewPrime.\allowbreak{}step} (\texttt{Development.\allowbreak{}AbelianVarietyTheory.\allowbreak{}Proofs.\allowbreak{}InertiaEigenvalues.\allowbreak{}AntiOfKer}, lines 153--353)\\
\texttt{GlobalTheory.\allowbreak{}NewPrime.\allowbreak{}det\_\allowbreak{}ne\_\allowbreak{}zero} (\texttt{Development.\allowbreak{}AbelianVarietyTheory.\allowbreak{}Proofs.\allowbreak{}InertiaEigenvalues.\allowbreak{}AntiOfKer}, lines 85--152)\\
\texttt{GlobalTheory.\allowbreak{}LowDet.\allowbreak{}det\_\allowbreak{}expand} (\texttt{Development.\allowbreak{}AbelianVarietyTheory.\allowbreak{}Proofs.\allowbreak{}InertiaEigenvalues.\allowbreak{}DetNeZeroOfLowestFormNeZeroPresentation}, lines 188--231)\\
\texttt{GlobalTheory.\allowbreak{}LowDet.\allowbreak{}det\_\allowbreak{}ne\_\allowbreak{}zero} (\texttt{Development.\allowbreak{}AbelianVarietyTheory.\allowbreak{}Proofs.\allowbreak{}InertiaEigenvalues.\allowbreak{}DetNeZeroOfLowestFormNeZeroPresentation}, lines 234--373)\par}

{\footnotesize\raggedright\noindent \ref{B8.07p}.\quad \texttt{GlobalTheory.\allowbreak{}ResSel.\allowbreak{}OddBound.\allowbreak{}residual\_\allowbreak{}bound\_\allowbreak{}of\_\allowbreak{}parts} (\texttt{Development.\allowbreak{}AbelianVarietyTheory.\allowbreak{}Proofs.\allowbreak{}InertiaEigenvalues.\allowbreak{}ResidualBoundOfComponents}, lines 28--57)\\
\texttt{GlobalTheory.\allowbreak{}StepTwo.\allowbreak{}StartInd.\allowbreak{}start\_\allowbreak{}of\_\allowbreak{}galHyp} (\texttt{Development.\allowbreak{}AbelianVarietyTheory.\allowbreak{}Proofs.\allowbreak{}InertiaEigenvalues.\allowbreak{}ContinuousCocycleChangeScalarsBack}, lines 452--496)\\
\texttt{GlobalTheory.\allowbreak{}StepTwo.\allowbreak{}TwoBound.\allowbreak{}induct} (\texttt{Development.\allowbreak{}AbelianVarietyTheory.\allowbreak{}Proofs.\allowbreak{}InertiaEigenvalues.\allowbreak{}AddValLeOfNotMem}, lines 60--86)\\
\texttt{GlobalTheory.\allowbreak{}StepTwo.\allowbreak{}TwoBound.\allowbreak{}thmC\_\allowbreak{}of\_\allowbreak{}parts} (\texttt{Development.\allowbreak{}AbelianVarietyTheory.\allowbreak{}Proofs.\allowbreak{}InertiaEigenvalues.\allowbreak{}AddValLeOfNotMem}, lines 87--144)\par}

{\footnotesize\raggedright\noindent \ref{B8.07r}.\quad \texttt{GlobalTheory.\allowbreak{}TwInputs.\allowbreak{}admT4\_\allowbreak{}of\_\allowbreak{}full\_\allowbreak{}compatibility} (\texttt{Development.\allowbreak{}AbelianVarietyTheory.\allowbreak{}Proofs.\allowbreak{}InertiaEigenvalues.\allowbreak{}AdmT4OfRootsInField}, lines 110--125)\\
\texttt{GlobalTheory.\allowbreak{}TwTraceCx.\allowbreak{}trace\_\allowbreak{}cx} (\texttt{Development.\allowbreak{}AbelianVarietyTheory.\allowbreak{}Proofs.\allowbreak{}InertiaEigenvalues.\allowbreak{}TwistedThetaGeneratorEqClassSumIntegralForm}, lines 217--242)\par}

{\footnotesize\raggedright\noindent \ref{B8.07s}.\quad \texttt{GlobalTheory.\allowbreak{}XiBound.\allowbreak{}main} (\texttt{Development.\allowbreak{}AbelianVarietyTheory.\allowbreak{}Proofs.\allowbreak{}InertiaEigenvalues.\allowbreak{}UniversalSeriesCoeffBoundedInSupOfTwistSeriesSid}, lines 40--95)\\
\texttt{GlobalTheory.\allowbreak{}TwInputs.\allowbreak{}univSup\_\allowbreak{}of\_\allowbreak{}xi} (\texttt{Development.\allowbreak{}AbelianVarietyTheory.\allowbreak{}Proofs.\allowbreak{}InertiaEigenvalues.\allowbreak{}UnivSupOfXi}, lines 13--45)\par}

{\footnotesize\raggedright\noindent \ref{B8.07q}.\quad \texttt{GlobalTheory.\allowbreak{}ExportFirst.\allowbreak{}exportF\_\allowbreak{}isBuiltE\_\allowbreak{}alt\_\allowbreak{}alt} (\texttt{Development.\allowbreak{}AutomorphicArithmetic.\allowbreak{}ThetaCorrespondence.\allowbreak{}ThetaPairLifts.\allowbreak{}ThirdPlaceSeries.\allowbreak{}SExport}, lines 20--32)\\
\texttt{GlobalTheory.\allowbreak{}ExportClose.\allowbreak{}exportAt\_\allowbreak{}of\_\allowbreak{}bundles\_\allowbreak{}alt} (\texttt{Development.\allowbreak{}AutomorphicArithmetic.\allowbreak{}ThetaCorrespondence.\allowbreak{}ThetaPairLifts.\allowbreak{}ThirdPlaceSeries.\allowbreak{}ExportAtOfBundles}, lines 17--35)\\
\texttt{GlobalTheory.\allowbreak{}TwistChoiceAll.\allowbreak{}admissibleIndexBoundAllPowers\_\allowbreak{}of\_\allowbreak{}conditions} (\texttt{Development.\allowbreak{}AbelianVarietyTheory.\allowbreak{}Proofs.\allowbreak{}InertiaEigenvalues.\allowbreak{}AdmissibleIndexBoundAllPowersOfConditions}, lines 92--223)\\
\texttt{GlobalTheory.\allowbreak{}ExportFromInputs.\allowbreak{}twistModule\_\allowbreak{}all\_\allowbreak{}character\_\allowbreak{}powers\_\allowbreak{}of\_\allowbreak{}admissible} (\texttt{Development.\allowbreak{}AbelianVarietyTheory.\allowbreak{}Proofs.\allowbreak{}InertiaEigenvalues.\allowbreak{}TwistModuleAllCharacterPowersOfAdmissible}, lines 17--53)\par}

{\footnotesize\raggedright\noindent \ref{B8.07t}.\quad \texttt{Formalism.\allowbreak{}L2.\allowbreak{}Fam.\allowbreak{}NPC.\allowbreak{}LocT} (\texttt{Formalism.\allowbreak{}LocalToGlobal.\allowbreak{}NumberFieldFamilies.\allowbreak{}NPartCNd}, lines 135--138)\\
\texttt{GlobalTheory.\allowbreak{}HeckeLocN.\allowbreak{}ndSpec\_\allowbreak{}exists\_\allowbreak{}alt} (\texttt{Development.\allowbreak{}AutomorphicArithmetic.\allowbreak{}ThetaCorrespondence.\allowbreak{}ThetaPairLifts.\allowbreak{}ThirdPlaceSeries.\allowbreak{}NdSpecExists}, lines 21--103)\\
\texttt{GlobalTheory.\allowbreak{}HeckeLocN.\allowbreak{}mkX} (\texttt{Development.\allowbreak{}AbelianVarietyTheory.\allowbreak{}Proofs.\allowbreak{}IwasawaDecomposition.\allowbreak{}UnitsLiftSpec}, lines 143--214)\\
\texttt{GlobalTheory.\allowbreak{}HeightCounting.\allowbreak{}Fin3.\allowbreak{}lv3\_\allowbreak{}of\_\allowbreak{}third\_\allowbreak{}alt\_\allowbreak{}alt} (\texttt{Development.\allowbreak{}AutomorphicArithmetic.\allowbreak{}ThetaCorrespondence.\allowbreak{}ThetaPairLifts.\allowbreak{}ThirdPlaceSeries.\allowbreak{}FinitenessOfHyp}, lines 61--83)\par}

{\footnotesize\raggedright\noindent \ref{B8.07}.\quad \texttt{GlobalTheory.\allowbreak{}T.\allowbreak{}large\_\allowbreak{}unramified\_\allowbreak{}exponent\_\allowbreak{}bundles\_\allowbreak{}exist\_\allowbreak{}alt} (\texttt{Development.\allowbreak{}AutomorphicArithmetic.\allowbreak{}ThetaCorrespondence.\allowbreak{}ThetaPairLifts.\allowbreak{}ThirdPlaceSeries.\allowbreak{}MemberBundlesLargeUnramifiedExponents}, lines 69--70)\\
\texttt{GlobalTheory.\allowbreak{}Fc.\allowbreak{}StageOneU.\allowbreak{}member\_\allowbreak{}bundles\_\allowbreak{}large\_\allowbreak{}unramified\_\allowbreak{}exponents\_\allowbreak{}alt} (\texttt{Development.\allowbreak{}AutomorphicArithmetic.\allowbreak{}ThetaCorrespondence.\allowbreak{}ThetaPairLifts.\allowbreak{}ThirdPlaceSeries.\allowbreak{}MemberBundlesLargeUnramifiedExponents}, lines 34--63)\par}

{\footnotesize\raggedright\noindent \ref{B8.07}.\quad \texttt{GlobalTheory.\allowbreak{}Mem.\allowbreak{}LiteralAutomorphy.\allowbreak{}memberE\_\allowbreak{}extensionWithAtProperty\_\allowbreak{}of\_\allowbreak{}memberE} (\texttt{Development.\allowbreak{}AutomorphicArithmetic.\allowbreak{}ThetaCorrespondence.\allowbreak{}ThetaPairLifts.\allowbreak{}ThirdPlaceSeries.\allowbreak{}IsBuiltEPrmAIff}, lines 35--53)\\
\texttt{GlobalTheory.\allowbreak{}Mem.\allowbreak{}LiteralAutomorphy.\allowbreak{}exists\_\allowbreak{}weilBlock\_\allowbreak{}ne\_\allowbreak{}at\_\allowbreak{}member\_\allowbreak{}alt} (\texttt{Development.\allowbreak{}AutomorphicArithmetic.\allowbreak{}ThetaCorrespondence.\allowbreak{}ThetaPairLifts.\allowbreak{}ThirdPlaceSeries.\allowbreak{}ExistsWeilBlockNeAtMember}, lines 24--33)\par}

{\footnotesize\raggedright\noindent \ref{B8.08a}.\quad \texttt{GlobalTheory.\allowbreak{}MiTop.\allowbreak{}rel\_\allowbreak{}of\_\allowbreak{}links\_\allowbreak{}alt} (\texttt{Development.\allowbreak{}AutomorphicArithmetic.\allowbreak{}ThetaCorrespondence.\allowbreak{}ThetaPairLifts.\allowbreak{}ThirdPlaceSeries.\allowbreak{}RelOfLinks}, lines 22--90)\\
\texttt{GlobalTheory.\allowbreak{}MiTop.\allowbreak{}served'} (\texttt{Development.\allowbreak{}AutomorphicArithmetic.\allowbreak{}ThetaCorrespondence.\allowbreak{}ThetaPairLifts.\allowbreak{}ThirdPlaceSeries.\allowbreak{}NotMemQBase}, lines 47--80)\par}

{\footnotesize\raggedright\noindent \ref{B8.08d}.\quad \texttt{GlobalTheory.\allowbreak{}Mem.\allowbreak{}SRe.\allowbreak{}exists\_\allowbreak{}sRest\_\allowbreak{}top} (\texttt{Development.\allowbreak{}AbelianVarietyTheory.\allowbreak{}Proofs.\allowbreak{}IwasawaDecomposition.\allowbreak{}ExistsSpectralLevelObjectsReplacedStages}, lines 24--40)\\
\texttt{Formalism.\allowbreak{}L2.\allowbreak{}SG.\allowbreak{}RowDict.\allowbreak{}range\_\allowbreak{}sEE} (\texttt{Formalism.\allowbreak{}LocalToGlobal.\allowbreak{}SiegelLatticeTheory.\allowbreak{}FlagStabilizerTensors.\allowbreak{}RowDict}, lines 148--163)\par}

{\footnotesize\raggedright\noindent \ref{B8.08g}.\quad \texttt{GlobalTheory.\allowbreak{}Chain.\allowbreak{}ChainPar} (\texttt{Development.\allowbreak{}AbelianVarietyTheory.\allowbreak{}Definitions.\allowbreak{}TwistedInnerProducts.\allowbreak{}IsStepWeil}, lines 32--37)\\
\texttt{GlobalTheory.\allowbreak{}StepRule.\allowbreak{}rulePar\_\allowbreak{}step} (\texttt{Development.\allowbreak{}AbelianVarietyTheory.\allowbreak{}Proofs.\allowbreak{}DecoratedLatticeProofs.\allowbreak{}FirstOccurrenceRuleThetaChainStepBothDirections}, lines 58--80)\\
\texttt{GlobalTheory.\allowbreak{}LowerMembersRule.\allowbreak{}lowFock} (\texttt{Development.\allowbreak{}AbelianVarietyTheory.\allowbreak{}Proofs.\allowbreak{}TwistedOperators.\allowbreak{}LowerTowerLiftsVanishOnFock}, lines 91--214)\\
\texttt{GlobalTheory.\allowbreak{}PairLift.\allowbreak{}cuspAll\_\allowbreak{}of\_\allowbreak{}lowAll} (\texttt{Development.\allowbreak{}AbelianVarietyTheory.\allowbreak{}Proofs.\allowbreak{}TwistedOperators.\allowbreak{}ThetaLiftsCuspidalOfLowerTowerLiftsVanish}, lines 142--153)\par}

{\footnotesize\raggedright\noindent \ref{B8.08e}.\quad \texttt{GlobalTheory.\allowbreak{}ChainOnto.\allowbreak{}chainLW\_\allowbreak{}holds} (\texttt{Development.\allowbreak{}AbelianVarietyTheory.\allowbreak{}Proofs.\allowbreak{}DecoratedLatticeProofs.\allowbreak{}ThetaChainLValuesNonvanishing}, lines 51--118)\\
\texttt{GlobalTheory.\allowbreak{}ChainOnto.\allowbreak{}QuotL.\allowbreak{}quotL\_\allowbreak{}entire\_\allowbreak{}ne\_\allowbreak{}zero} (\texttt{Development.\allowbreak{}AbelianVarietyTheory.\allowbreak{}Proofs.\allowbreak{}DecoratedLatticeProofs.\allowbreak{}ThetaChainCharQuotientLFunctionEntireNeZero}, lines 45--82)\par}

{\footnotesize\raggedright\noindent \ref{B8.08j}.\quad \texttt{GlobalTheory.\allowbreak{}Fam.\allowbreak{}MemberCount.\allowbreak{}rankMu\_\allowbreak{}of\_\allowbreak{}bundle\_\allowbreak{}alt} (\texttt{Development.\allowbreak{}AutomorphicArithmetic.\allowbreak{}ThetaCorrespondence.\allowbreak{}ThetaPairLifts.\allowbreak{}ThirdPlaceSeries.\allowbreak{}RankMuOfBundle}, lines 17--22)\\
\texttt{GlobalTheory.\allowbreak{}Diag.\allowbreak{}PG.\allowbreak{}duWin\_\allowbreak{}three\_\allowbreak{}of\_\allowbreak{}count} (\texttt{Development.\allowbreak{}AbelianVarietyTheory.\allowbreak{}Proofs.\allowbreak{}FlagMeasures.\allowbreak{}DuWinThreeOfCount}, lines 60--93)\\
\texttt{GlobalTheory.\allowbreak{}Core.\allowbreak{}Rk.\allowbreak{}ramShare\_\allowbreak{}of\_\allowbreak{}trueD} (\texttt{Development.\allowbreak{}AbelianVarietyTheory.\allowbreak{}Proofs.\allowbreak{}DecoratedLatticeProofs.\allowbreak{}RamShareOfFreeT}, lines 36--38)\par}

{\footnotesize\raggedright\noindent \ref{B8.08k}.\quad \texttt{GlobalTheory.\allowbreak{}Core.\allowbreak{}Deg.\allowbreak{}rowsND\_\allowbreak{}one\_\allowbreak{}record} (\texttt{Development.\allowbreak{}AbelianVarietyTheory.\allowbreak{}Proofs.\allowbreak{}DecoratedLatticeProofs.\allowbreak{}NondegenerateCollisionRowBoundAtParameters}, lines 18--84)\\
\texttt{GlobalTheory.\allowbreak{}Core.\allowbreak{}Deg.\allowbreak{}rowsDeg\_\allowbreak{}at\_\allowbreak{}record} (\texttt{Development.\allowbreak{}AbelianVarietyTheory.\allowbreak{}Proofs.\allowbreak{}DecoratedLatticeProofs.\allowbreak{}DegenerateCollisionRowBoundAtParams}, lines 15--91)\\
\texttt{GlobalTheory.\allowbreak{}Core.\allowbreak{}Deg.\allowbreak{}corankSums\_\allowbreak{}expLoss} (\texttt{Development.\allowbreak{}AbelianVarietyTheory.\allowbreak{}Proofs.\allowbreak{}DecoratedLatticeProofs.\allowbreak{}CorankSumsBoundWithExpLoss}, lines 184--220)\\
\texttt{GlobalTheory.\allowbreak{}Core.\allowbreak{}Deg.\allowbreak{}collS\_\allowbreak{}of\_\allowbreak{}rowBounds} (\texttt{Development.\allowbreak{}AbelianVarietyTheory.\allowbreak{}Proofs.\allowbreak{}DecoratedLatticeProofs.\allowbreak{}CollisionAbsSumBoundOfTwoRowBounds}, lines 16--91)\par}

{\footnotesize\raggedright\noindent \ref{B8.08l}.\quad \texttt{GlobalTheory.\allowbreak{}diag\_\allowbreak{}abstract\_\allowbreak{}collisions\_\allowbreak{}once} (\texttt{Development.\allowbreak{}AbelianVarietyTheory.\allowbreak{}Proofs.\allowbreak{}DecoratedLatticeProofs.\allowbreak{}AbstractDiagonalBoundCollisionsCountedOnce}, lines 21--88)\\
\texttt{Formalism.\allowbreak{}FourierDecayAlt.\allowbreak{}P.\allowbreak{}GrowthPolynomialOperator.\allowbreak{}t\_\allowbreak{}PriceOfIota} (\texttt{Formalism.\allowbreak{}HermitianCurvature.\allowbreak{}HermitianCurvatureProofs.\allowbreak{}GrowthOfBasic}, lines 384--393)\\
\texttt{GlobalTheory.\allowbreak{}Core.\allowbreak{}Deg.\allowbreak{}pgC\_\allowbreak{}le\_\allowbreak{}six\_\allowbreak{}fifths} (\texttt{Development.\allowbreak{}AbelianVarietyTheory.\allowbreak{}Proofs.\allowbreak{}DecoratedLatticeProofs.\allowbreak{}EqualIdentificationPairPriceLeSixFifths}, lines 55--75)\\
\texttt{GlobalTheory.\allowbreak{}Core.\allowbreak{}Diag.\allowbreak{}diag\_\allowbreak{}arith} (\texttt{Development.\allowbreak{}AbelianVarietyTheory.\allowbreak{}Proofs.\allowbreak{}DecoratedLatticeProofs.\allowbreak{}DiagonalTheoremOfFourPointBounds}, lines 36--48)\\
\texttt{GlobalTheory.\allowbreak{}Core.\allowbreak{}Rk.\allowbreak{}diagonalRow\_\allowbreak{}bound\_\allowbreak{}from\_\allowbreak{}windowWidthThree} (\texttt{Development.\allowbreak{}AbelianVarietyTheory.\allowbreak{}Proofs.\allowbreak{}DecoratedLatticeProofs.\allowbreak{}DiagonalRowBoundFromWindowWidthThree}, lines 38--98)\par}

{\footnotesize\raggedright\noindent \ref{B8.08m}.\quad \texttt{Formalism.\allowbreak{}L2.\allowbreak{}SG.\allowbreak{}TermOp.\allowbreak{}vec} (\texttt{Formalism.\allowbreak{}LocalToGlobal.\allowbreak{}SiegelLatticeTheory.\allowbreak{}ClosedLocalOrbits.\allowbreak{}TermOp}, lines 343--346)\par}

{\footnotesize\raggedright\noindent \ref{B8.08n}.\quad \texttt{GlobalTheory.\allowbreak{}Core.\allowbreak{}One.\allowbreak{}stepClass1\_\allowbreak{}of} (\texttt{Development.\allowbreak{}AbelianVarietyTheory.\allowbreak{}Proofs.\allowbreak{}DecoratedLatticeProofs.\allowbreak{}StepClass1Of}, lines 9--40)\\
\texttt{GlobalTheory.\allowbreak{}Core.\allowbreak{}One.\allowbreak{}scT\_\allowbreak{}sc\_\allowbreak{}genSet} (\texttt{Development.\allowbreak{}AbelianVarietyTheory.\allowbreak{}Proofs.\allowbreak{}DecoratedLatticeProofs.\allowbreak{}RowGen1POfSteps}, lines 16--22)\\
\texttt{GlobalTheory.\allowbreak{}Core.\allowbreak{}One.\allowbreak{}genericTerms\_\allowbreak{}points\_\allowbreak{}agree\_\allowbreak{}groupSums\_\allowbreak{}corankOne} (\texttt{Development.\allowbreak{}AbelianVarietyTheory.\allowbreak{}Proofs.\allowbreak{}DecoratedLatticeProofs.\allowbreak{}GenericTermsPointsAgreeGroupSumsCorankOne}, lines 8--59)\\
\texttt{Formalism.\allowbreak{}L2.\allowbreak{}SG.\allowbreak{}Pf.\allowbreak{}MainEqCut.\allowbreak{}meanAmplitudeGeneric\_\allowbreak{}holds} (\texttt{Formalism.\allowbreak{}LocalToGlobal.\allowbreak{}SiegelLatticeTheory.\allowbreak{}CokernelGLCounts.\allowbreak{}RankPointCounts.\allowbreak{}MainEqCut}, lines 389--461)\par}

{\footnotesize\raggedright\noindent \ref{B8.08oa}.\quad \texttt{kappa\_\allowbreak{}off\_\allowbreak{}diagonal\_\allowbreak{}dim\_\allowbreak{}form\_\allowbreak{}C} (\texttt{Formalism.\allowbreak{}PeriodIntegrals.\allowbreak{}OrbitalBoundsConclusion.\allowbreak{}UnitaryDimensionFormula}, lines 6--8)\\
\texttt{Formalism.\allowbreak{}L2.\allowbreak{}SG.\allowbreak{}KappaOffDiagonal.\allowbreak{}DimFormC} (\texttt{Formalism.\allowbreak{}LocalToGlobal.\allowbreak{}SiegelLatticeTheory.\allowbreak{}CokernelGLCounts.\allowbreak{}RankPointCounts.\allowbreak{}DimFormC}, lines 80--83)\\
\texttt{PI.\allowbreak{}UnitaryDimensionFormula.\allowbreak{}dimFormC\_\allowbreak{}of\_\allowbreak{}three} (\texttt{Formalism.\allowbreak{}PeriodIntegrals.\allowbreak{}UnitaryGroupMassFormula.\allowbreak{}RankTwoDimensionFormula.\allowbreak{}DimFormCOfThree}, lines 78--80)\\
\texttt{PI.\allowbreak{}UnitaryDimensionFormula.\allowbreak{}dimFormC\_\allowbreak{}of\_\allowbreak{}steps} (\texttt{Formalism.\allowbreak{}PeriodIntegrals.\allowbreak{}UnitaryGroupMassFormula.\allowbreak{}RankTwoDimensionFormula.\allowbreak{}DimFormCOfSteps}, lines 117--120)\\
\texttt{PI.\allowbreak{}UnitaryDimensionFormula.\allowbreak{}dimFormC\_\allowbreak{}of} (\texttt{Formalism.\allowbreak{}PeriodIntegrals.\allowbreak{}UnitaryGroupMassFormula.\allowbreak{}RankTwoDimensionFormula.\allowbreak{}GammaOneEqProd}, lines 106--121)\\
\texttt{PI.\allowbreak{}UnitaryDimensionFormula.\allowbreak{}fibreBound\_\allowbreak{}of} (\texttt{Formalism.\allowbreak{}PeriodIntegrals.\allowbreak{}UnitaryGroupMassFormula.\allowbreak{}RankTwoDimensionFormula.\allowbreak{}DimFormCOfSteps}, lines 79--116)\\
\texttt{PI.\allowbreak{}UnitaryDimensionFormula.\allowbreak{}orbitCount} (\texttt{Formalism.\allowbreak{}PeriodIntegrals.\allowbreak{}UnitaryGroupMassFormula.\allowbreak{}RankTwoDimensionFormula.\allowbreak{}DimFormCOfThree}, lines 13--59)\\
\texttt{PI.\allowbreak{}UnitaryDimensionFormula.\allowbreak{}kuStab} (\texttt{Formalism.\allowbreak{}PeriodIntegrals.\allowbreak{}UnitaryGroupMassFormula.\allowbreak{}RankTwoDimensionFormula.\allowbreak{}DimFormCOfThree}, lines 60--77)\\
\texttt{PI.\allowbreak{}UnitaryDimensionFormula.\allowbreak{}transport\_\allowbreak{}holds} (\texttt{Formalism.\allowbreak{}PeriodIntegrals.\allowbreak{}UnitaryGroupMassFormula.\allowbreak{}RankTwoDimensionFormula.\allowbreak{}Transp}, lines 101--152)\\
\texttt{PI.\allowbreak{}UnitaryDimensionFormula.\allowbreak{}stabBound\_\allowbreak{}holds} (\texttt{Formalism.\allowbreak{}PeriodIntegrals.\allowbreak{}UnitaryGroupMassFormula.\allowbreak{}RankTwoDimensionFormula.\allowbreak{}StabBoundHolds}, lines 261--319)\\
\texttt{PI.\allowbreak{}UnitaryDimensionFormula.\allowbreak{}indexGe} (\texttt{Formalism.\allowbreak{}PeriodIntegrals.\allowbreak{}UnitaryGroupMassFormula.\allowbreak{}RankTwoDimensionFormula.\allowbreak{}IndexGeFull}, lines 10--11)\\
\texttt{PI.\allowbreak{}UnitaryDimensionFormula.\allowbreak{}indexGe\_\allowbreak{}of} (\texttt{Formalism.\allowbreak{}PeriodIntegrals.\allowbreak{}UnitaryGroupMassFormula.\allowbreak{}RankTwoDimensionFormula.\allowbreak{}IndexGe}, lines 117--166)\\
\texttt{PI.\allowbreak{}UnitaryDimensionFormula.\allowbreak{}gamma\_\allowbreak{}one\_\allowbreak{}eq\_\allowbreak{}card} (\texttt{Formalism.\allowbreak{}PeriodIntegrals.\allowbreak{}UnitaryGroupMassFormula.\allowbreak{}RankTwoDimensionFormula.\allowbreak{}IndexGe}, lines 78--116)\\
\texttt{PI.\allowbreak{}UnitaryDimensionFormula.\allowbreak{}idxHom} (\texttt{Formalism.\allowbreak{}PeriodIntegrals.\allowbreak{}UnitaryGroupMassFormula.\allowbreak{}RankTwoDimensionFormula.\allowbreak{}IdxKHom}, lines 25--32)\\
\texttt{PI.\allowbreak{}UnitaryDimensionFormula.\allowbreak{}idxSurj\_\allowbreak{}of\_\allowbreak{}frame} (\texttt{Formalism.\allowbreak{}PeriodIntegrals.\allowbreak{}UnitaryGroupMassFormula.\allowbreak{}RankTwoDimensionFormula.\allowbreak{}IdxDChk}, lines 12--14)\\
\texttt{PI.\allowbreak{}UnitaryDimensionFormula.\allowbreak{}idxDFrame} (\texttt{Formalism.\allowbreak{}PeriodIntegrals.\allowbreak{}UnitaryGroupMassFormula.\allowbreak{}RankTwoDimensionFormula.\allowbreak{}IdxDcvFrameAsm}, lines 83--84)\par}

{\footnotesize\raggedright\noindent \ref{B8.08ob}.\quad \texttt{Formalism.\allowbreak{}L2.\allowbreak{}SG.\allowbreak{}KappaOffDiagonal.\allowbreak{}irrVec\_\allowbreak{}nohyp} (\texttt{Formalism.\allowbreak{}LocalToGlobal.\allowbreak{}SiegelLatticeTheory.\allowbreak{}CokernelGLCounts.\allowbreak{}RankPointCounts.\allowbreak{}CovolumeSpans}, lines 28--34)\\
\texttt{Formalism.\allowbreak{}L2.\allowbreak{}SG.\allowbreak{}KappaOffDiagonal.\allowbreak{}IrrVecC} (\texttt{Formalism.\allowbreak{}LocalToGlobal.\allowbreak{}SiegelLatticeTheory.\allowbreak{}CokernelGLCounts.\allowbreak{}RankPointCounts.\allowbreak{}DimFormC}, lines 47--52)\\
\texttt{Formalism.\allowbreak{}L2.\allowbreak{}SG.\allowbreak{}KappaOffDiagonal.\allowbreak{}vector\_\allowbreak{}eq\_\allowbreak{}zero\_\allowbreak{}of\_\allowbreak{}orbital\_\allowbreak{}zero} (\texttt{Formalism.\allowbreak{}LocalToGlobal.\allowbreak{}SiegelLatticeTheory.\allowbreak{}CokernelGLCounts.\allowbreak{}RankPointCounts.\allowbreak{}CovolumeSpans}, lines 17--27)\\
\texttt{Formalism.\allowbreak{}L2.\allowbreak{}SG.\allowbreak{}TermOp.\allowbreak{}Vanish.\allowbreak{}ooVanish\_\allowbreak{}of} (\texttt{Formalism.\allowbreak{}LocalToGlobal.\allowbreak{}SiegelLatticeTheory.\allowbreak{}ClosedLocalOrbits.\allowbreak{}GaugeLatticeShifts.\allowbreak{}Vanish}, lines 512--570)\\
\texttt{Formalism.\allowbreak{}L2.\allowbreak{}SG.\allowbreak{}NonRelPf.\allowbreak{}nonRelCharC\_\allowbreak{}holds} (\texttt{Formalism.\allowbreak{}LocalToGlobal.\allowbreak{}SiegelLatticeTheory.\allowbreak{}ClosedLocalOrbits.\allowbreak{}GaugeLatticeShifts.\allowbreak{}LieAveraging.\allowbreak{}NonRelCharCHolds}, lines 118--262)\\
\texttt{Formalism.\allowbreak{}L2.\allowbreak{}SG.\allowbreak{}TermOp.\allowbreak{}Vanish.\allowbreak{}rational\_\allowbreak{}cocompact\_\allowbreak{}holds} (\texttt{Formalism.\allowbreak{}LocalToGlobal.\allowbreak{}SiegelLatticeTheory.\allowbreak{}ClosedLocalOrbits.\allowbreak{}GaugeLatticeShifts.\allowbreak{}RCocompactPfFinal}, lines 7--8)\\
\texttt{Formalism.\allowbreak{}L2.\allowbreak{}SG.\allowbreak{}TermOp.\allowbreak{}Vanish.\allowbreak{}rCocompactC\_\allowbreak{}of} (\texttt{Formalism.\allowbreak{}LocalToGlobal.\allowbreak{}SiegelLatticeTheory.\allowbreak{}ClosedLocalOrbits.\allowbreak{}GaugeLatticeShifts.\allowbreak{}RCocompactPfMain}, lines 235--332)\\
\texttt{Formalism.\allowbreak{}L2.\allowbreak{}SG.\allowbreak{}Pf.\allowbreak{}ginfUnimod\_\allowbreak{}holds} (\texttt{Formalism.\allowbreak{}LocalToGlobal.\allowbreak{}SiegelLatticeTheory.\allowbreak{}ClosedLocalOrbits.\allowbreak{}GaugeLatticeShifts.\allowbreak{}NilradicalLatticeZMemSpan}, lines 335--336)\\
\texttt{Formalism.\allowbreak{}L2.\allowbreak{}SG.\allowbreak{}Pf.\allowbreak{}nLamNinf\_\allowbreak{}holds} (\texttt{Formalism.\allowbreak{}LocalToGlobal.\allowbreak{}SiegelLatticeTheory.\allowbreak{}ClosedLocalOrbits.\allowbreak{}GaugeLatticeShifts.\allowbreak{}NilradicalLatticeZMemSpan}, lines 294--315)\\
\texttt{Formalism.\allowbreak{}L2.\allowbreak{}SG.\allowbreak{}Pf.\allowbreak{}psiInfChar\_\allowbreak{}holds} (\texttt{Formalism.\allowbreak{}LocalToGlobal.\allowbreak{}SiegelLatticeTheory.\allowbreak{}ClosedLocalOrbits.\allowbreak{}GaugeLatticeShifts.\allowbreak{}NilradicalLatticeZMemSpan}, lines 471--488)\\
\texttt{Formalism.\allowbreak{}L2.\allowbreak{}SG.\allowbreak{}Pf.\allowbreak{}rClosedC\_\allowbreak{}holds} (\texttt{Formalism.\allowbreak{}LocalToGlobal.\allowbreak{}SiegelLatticeTheory.\allowbreak{}ClosedLocalOrbits.\allowbreak{}GaugeLatticeShifts.\allowbreak{}NilradicalLatticeZMemSpan}, lines 489--496)\\
\texttt{Formalism.\allowbreak{}L2.\allowbreak{}SG.\allowbreak{}Pf.\allowbreak{}roDiscreteC\_\allowbreak{}holds} (\texttt{Formalism.\allowbreak{}LocalToGlobal.\allowbreak{}SiegelLatticeTheory.\allowbreak{}ClosedLocalOrbits.\allowbreak{}GaugeLatticeShifts.\allowbreak{}NilradicalLatticeZMemSpan}, lines 337--345)\\
\texttt{Formalism.\allowbreak{}L2.\allowbreak{}SG.\allowbreak{}Closed.\allowbreak{}ChiTrivialC\_\allowbreak{}holds} (\texttt{Formalism.\allowbreak{}LocalToGlobal.\allowbreak{}SiegelLatticeTheory.\allowbreak{}ClosedVolumeStatements.\allowbreak{}UnipotentInLevelSubgroup}, lines 55--55)\\
\texttt{Formalism.\allowbreak{}L2.\allowbreak{}SG.\allowbreak{}Closed.\allowbreak{}AvgCore\_\allowbreak{}holds} (\texttt{Formalism.\allowbreak{}LocalToGlobal.\allowbreak{}SiegelLatticeTheory.\allowbreak{}ClosedVolumeStatements.\allowbreak{}UnipotentInLevelSubgroup}, lines 56--56)\par}

{\footnotesize\raggedright\noindent \ref{B8.08oc}.\quad \texttt{Formalism.\allowbreak{}L2.\allowbreak{}SG.\allowbreak{}NormCut.\allowbreak{}normExpansionOn\_\allowbreak{}tOf} (\texttt{Formalism.\allowbreak{}LocalToGlobal.\allowbreak{}SiegelLatticeTheory.\allowbreak{}ClosedLocalOrbits.\allowbreak{}GaugeLatticeShifts.\allowbreak{}NormCut}, lines 202--267)\\
\texttt{Formalism.\allowbreak{}L2.\allowbreak{}SG.\allowbreak{}NormCut.\allowbreak{}normExp\_\allowbreak{}abs} (\texttt{Formalism.\allowbreak{}LocalToGlobal.\allowbreak{}SiegelLatticeTheory.\allowbreak{}ClosedLocalOrbits.\allowbreak{}GaugeLatticeShifts.\allowbreak{}NormAsm}, lines 184--343)\\
\texttt{Formalism.\allowbreak{}FourierDecayAlt.\allowbreak{}PB.\allowbreak{}Terms.\allowbreak{}NormExpansionOn} (\texttt{Formalism.\allowbreak{}HermitianCurvature.\allowbreak{}StatementPB}, lines 86--88)\\
\texttt{hilbert\_\allowbreak{}schmidt\_\allowbreak{}blocks} (\texttt{Formalism.\allowbreak{}PeriodIntegrals.\allowbreak{}OrbitalBoundsConclusion.\allowbreak{}NormHSBlocks}, lines 5--7)\\
\texttt{NormHilbertSchmidtBlocks.\allowbreak{}hSBlocks\_\allowbreak{}main} (\texttt{Formalism.\allowbreak{}PeriodIntegrals.\allowbreak{}UnitaryGroupMassFormula.\allowbreak{}HilbertSchmidtBlocks.\allowbreak{}HSBlocksMain}, lines 13--36)\\
\texttt{chi\_\allowbreak{}n\_\allowbreak{}bilinear} (\texttt{Formalism.\allowbreak{}PeriodIntegrals.\allowbreak{}OrbitalBoundsConclusion.\allowbreak{}NormChiNBi}, lines 3--5)\\
\texttt{PIW.\allowbreak{}NormChiNBi.\allowbreak{}chiNBi\_\allowbreak{}holds} (\texttt{Formalism.\allowbreak{}PeriodIntegrals.\allowbreak{}UnitaryGroupMassFormula.\allowbreak{}CharacterNormMatrixInverse.\allowbreak{}FmatMulInv}, lines 265--276)\\
\texttt{Formalism.\allowbreak{}L2.\allowbreak{}SG.\allowbreak{}NormCut.\allowbreak{}colFree\_\allowbreak{}steps} (\texttt{Formalism.\allowbreak{}LocalToGlobal.\allowbreak{}SiegelLatticeTheory.\allowbreak{}ClosedLocalOrbits.\allowbreak{}GaugeLatticeShifts.\allowbreak{}NormLinks}, lines 34--38)\\
\texttt{Formalism.\allowbreak{}L2.\allowbreak{}SG.\allowbreak{}NormCut.\allowbreak{}colFreeBij\_\allowbreak{}of} (\texttt{Formalism.\allowbreak{}LocalToGlobal.\allowbreak{}SiegelLatticeTheory.\allowbreak{}ClosedLocalOrbits.\allowbreak{}GaugeLatticeShifts.\allowbreak{}ColFreePf}, lines 58--149)\\
\texttt{Formalism.\allowbreak{}L2.\allowbreak{}SG.\allowbreak{}Pf.\allowbreak{}colLatK} (\texttt{Formalism.\allowbreak{}LocalToGlobal.\allowbreak{}SiegelLatticeTheory.\allowbreak{}FlagStabilizerTensors.\allowbreak{}RootLatticeThresholds.\allowbreak{}LatDictK}, lines 59--66)\\
\texttt{Formalism.\allowbreak{}L2.\allowbreak{}SG.\allowbreak{}Pf.\allowbreak{}NormSteps.\allowbreak{}levi\_\allowbreak{}onto\_\allowbreak{}reduction} (\texttt{Formalism.\allowbreak{}LocalToGlobal.\allowbreak{}SiegelLatticeTheory.\allowbreak{}ClosedLocalOrbits.\allowbreak{}GaugeLatticeShifts.\allowbreak{}NormStepsLeviOntoSEnd}, lines 23--25)\\
\texttt{Formalism.\allowbreak{}L2.\allowbreak{}SG.\allowbreak{}Pf.\allowbreak{}NormSteps.\allowbreak{}LeviOntoSPf.\allowbreak{}exact\_\allowbreak{}dual\_\allowbreak{}frame\_\allowbreak{}exists} (\texttt{Formalism.\allowbreak{}LocalToGlobal.\allowbreak{}SiegelLatticeTheory.\allowbreak{}ClosedLocalOrbits.\allowbreak{}GaugeLatticeShifts.\allowbreak{}NormStepsLeviFrame}, lines 563--580)\\
\texttt{Formalism.\allowbreak{}L2.\allowbreak{}SG.\allowbreak{}Pf.\allowbreak{}NormSteps.\allowbreak{}column\_\allowbreak{}lattice\_\allowbreak{}representatives} (\texttt{Formalism.\allowbreak{}LocalToGlobal.\allowbreak{}SiegelLatticeTheory.\allowbreak{}ClosedLocalOrbits.\allowbreak{}GaugeLatticeShifts.\allowbreak{}NormStepsColLatRepS}, lines 203--251)\\
\texttt{Formalism.\allowbreak{}L2.\allowbreak{}SG.\allowbreak{}Pf.\allowbreak{}massAC\_\allowbreak{}holds} (\texttt{Formalism.\allowbreak{}LocalToGlobal.\allowbreak{}SiegelLatticeTheory.\allowbreak{}ClosedLocalOrbits.\allowbreak{}GaugeLatticeShifts.\allowbreak{}ActMassC}, lines 7--9)\\
\texttt{Formalism.\allowbreak{}L2.\allowbreak{}SG.\allowbreak{}Pf.\allowbreak{}SatGram.\allowbreak{}satGram\_\allowbreak{}of} (\texttt{Formalism.\allowbreak{}LocalToGlobal.\allowbreak{}SiegelLatticeTheory.\allowbreak{}ClosedLocalOrbits.\allowbreak{}GaugeLatticeShifts.\allowbreak{}SatGram}, lines 529--543)\\
\texttt{Formalism.\allowbreak{}L2.\allowbreak{}SG.\allowbreak{}Pf.\allowbreak{}SatGram.\allowbreak{}zs\_\allowbreak{}gram} (\texttt{Formalism.\allowbreak{}LocalToGlobal.\allowbreak{}SiegelLatticeTheory.\allowbreak{}ClosedLocalOrbits.\allowbreak{}GaugeLatticeShifts.\allowbreak{}SatGram}, lines 486--495)\\
\texttt{Formalism.\allowbreak{}L2.\allowbreak{}SG.\allowbreak{}NormCut.\allowbreak{}leftActC'\_\allowbreak{}of\_\allowbreak{}PD} (\texttt{Formalism.\allowbreak{}LocalToGlobal.\allowbreak{}SiegelLatticeTheory.\allowbreak{}ClosedLocalOrbits.\allowbreak{}GaugeLatticeShifts.\allowbreak{}NormExpSatHolds}, lines 23--36)\\
\texttt{Formalism.\allowbreak{}L2.\allowbreak{}SG.\allowbreak{}NormCut.\allowbreak{}rightActC'\_\allowbreak{}holds} (\texttt{Formalism.\allowbreak{}LocalToGlobal.\allowbreak{}SiegelLatticeTheory.\allowbreak{}ClosedLocalOrbits.\allowbreak{}GaugeLatticeShifts.\allowbreak{}NormExpSatHolds}, lines 37--43)\par}

{\footnotesize\raggedright\noindent \ref{B8.08o}.\quad \texttt{GlobalTheory.\allowbreak{}Core.\allowbreak{}Deg.\allowbreak{}PairsDegOfType} (\texttt{Development.\allowbreak{}AbelianVarietyTheory.\allowbreak{}Definitions.\allowbreak{}LatticeFrames.\allowbreak{}TermSumBoundOfPairsAmplitude}, lines 32--37)\\
\texttt{GlobalTheory.\allowbreak{}Core.\allowbreak{}Deg.\allowbreak{}pairsDegOfType\_\allowbreak{}holds} (\texttt{Development.\allowbreak{}AbelianVarietyTheory.\allowbreak{}Proofs.\allowbreak{}DecoratedLatticeProofs.\allowbreak{}PairsAmplitudeBoundOfTypeBound}, lines 20--92)\\
\texttt{GlobalTheory.\allowbreak{}Core.\allowbreak{}Deg.\allowbreak{}sum\_\allowbreak{}labels\_\allowbreak{}one\_\allowbreak{}record} (\texttt{Development.\allowbreak{}AbelianVarietyTheory.\allowbreak{}Proofs.\allowbreak{}DecoratedLatticeProofs.\allowbreak{}SumLabelsOneRecord}, lines 12--38)\\
\texttt{Formalism.\allowbreak{}L2.\allowbreak{}SG.\allowbreak{}KappaOffDiagonal.\allowbreak{}kloostermanCosetSum\_\allowbreak{}le} (\texttt{Formalism.\allowbreak{}LocalToGlobal.\allowbreak{}SiegelLatticeTheory.\allowbreak{}CokernelGLCounts.\allowbreak{}RankPointCounts.\allowbreak{}KloostermanCosetSumLe}, lines 102--143)\\
\texttt{Formalism.\allowbreak{}L2.\allowbreak{}SG.\allowbreak{}KappaOffDiagonal.\allowbreak{}delta\_\allowbreak{}pos\_\allowbreak{}le\_\allowbreak{}height\_\allowbreak{}sq} (\texttt{Formalism.\allowbreak{}LocalToGlobal.\allowbreak{}SiegelLatticeTheory.\allowbreak{}CokernelGLCounts.\allowbreak{}RankPointCounts.\allowbreak{}RelevantSetZeroGNotIdentity}, lines 60--111)\\
\texttt{Formalism.\allowbreak{}MainNotions.\allowbreak{}KappaEndpoints.\allowbreak{}rank\_\allowbreak{}contribution\_\allowbreak{}le} (\texttt{Formalism.\allowbreak{}SharpRestriction.\allowbreak{}CuspSpanNormBounds.\allowbreak{}KappaEndpoints}, lines 3161--3224)\\
\texttt{Formalism.\allowbreak{}MainNotions.\allowbreak{}KappaEndpoints.\allowbreak{}final\_\allowbreak{}bound\_\allowbreak{}closes} (\texttt{Formalism.\allowbreak{}SharpRestriction.\allowbreak{}CuspSpanNormBounds.\allowbreak{}KappaEndpoints}, lines 2593--2602)\\
\texttt{GlobalTheory.\allowbreak{}Core.\allowbreak{}Deg.\allowbreak{}pgC\_\allowbreak{}le\_\allowbreak{}two\_\allowbreak{}at\_\allowbreak{}record} (\texttt{Development.\allowbreak{}AbelianVarietyTheory.\allowbreak{}Proofs.\allowbreak{}DecoratedLatticeProofs.\allowbreak{}PgCLeTwoAtRecord}, lines 10--30)\par}

{\footnotesize\raggedright\noindent \ref{B8.08p}.\quad \texttt{GlobalTheory.\allowbreak{}Core.\allowbreak{}Deg.\allowbreak{}TermSumBoundOfPairsAmplitude} (\texttt{Development.\allowbreak{}AbelianVarietyTheory.\allowbreak{}Definitions.\allowbreak{}LatticeFrames.\allowbreak{}TermSumBoundOfPairsAmplitude}, lines 38--42)\\
\texttt{GlobalTheory.\allowbreak{}Core.\allowbreak{}Deg.\allowbreak{}termSumBound\_\allowbreak{}of\_\allowbreak{}pairsAmplitude} (\texttt{Development.\allowbreak{}AbelianVarietyTheory.\allowbreak{}Proofs.\allowbreak{}DecoratedLatticeProofs.\allowbreak{}TermSumBoundOfPairsAmplitude}, lines 11--27)\\
\texttt{Formalism.\allowbreak{}L2.\allowbreak{}SG.\allowbreak{}RowsOther.\allowbreak{}kppDiag\_\allowbreak{}of\_\allowbreak{}pairs} (\texttt{Formalism.\allowbreak{}LocalToGlobal.\allowbreak{}SiegelLatticeTheory.\allowbreak{}CokernelGLCounts.\allowbreak{}RowsOther}, lines 35--47)\\
\texttt{Formalism.\allowbreak{}L2.\allowbreak{}SG.\allowbreak{}RowsOther.\allowbreak{}pairsDeg\_\allowbreak{}of\_\allowbreak{}amp} (\texttt{Formalism.\allowbreak{}LocalToGlobal.\allowbreak{}SiegelLatticeTheory.\allowbreak{}CokernelGLCounts.\allowbreak{}RowsOther}, lines 122--125)\\
\texttt{Formalism.\allowbreak{}FourierDecayAlt.\allowbreak{}PB.\allowbreak{}Terms.\allowbreak{}NormExpansionOn} (\texttt{Formalism.\allowbreak{}HermitianCurvature.\allowbreak{}StatementPB}, lines 86--88)\\
\texttt{Formalism.\allowbreak{}L2.\allowbreak{}SG.\allowbreak{}TermOp.\allowbreak{}ggBound} (\texttt{Formalism.\allowbreak{}LocalToGlobal.\allowbreak{}SiegelLatticeTheory.\allowbreak{}ClosedLocalOrbits.\allowbreak{}TermOpStmt}, lines 135--148)\\
\texttt{Formalism.\allowbreak{}L2.\allowbreak{}SG.\allowbreak{}TermOp.\allowbreak{}rowOrth} (\texttt{Formalism.\allowbreak{}LocalToGlobal.\allowbreak{}SiegelLatticeTheory.\allowbreak{}ClosedLocalOrbits.\allowbreak{}TermOpStmt}, lines 83--100)\par}

{\footnotesize\raggedright\noindent \ref{B8.08q}.\quad \texttt{GlobalTheory.\allowbreak{}Core.\allowbreak{}Deg.\allowbreak{}vanishingGramBoundWindow\_\allowbreak{}of\_\allowbreak{}termPairTypeBound} (\texttt{Development.\allowbreak{}AbelianVarietyTheory.\allowbreak{}Proofs.\allowbreak{}DecoratedLatticeProofs.\allowbreak{}VanishingGramBoundWindowOfTermPairTypeBound}, lines 9--21)\\
\texttt{GlobalTheory.\allowbreak{}Core.\allowbreak{}Deg.\allowbreak{}VanishingGramTermsBound} (\texttt{Development.\allowbreak{}AbelianVarietyTheory.\allowbreak{}Definitions.\allowbreak{}LatticeFrames.\allowbreak{}VanishingGramTermsBoundParams}, lines 65--68)\\
\texttt{GlobalTheory.\allowbreak{}Core.\allowbreak{}Deg.\allowbreak{}VanishingGramTermsBoundWindow} (\texttt{Development.\allowbreak{}AbelianVarietyTheory.\allowbreak{}Definitions.\allowbreak{}LatticeFrames.\allowbreak{}VanishingGramTermsBoundParams}, lines 104--105)\\
\texttt{GlobalTheory.\allowbreak{}TypeAsm.\allowbreak{}TypeP} (\texttt{Development.\allowbreak{}AbelianVarietyTheory.\allowbreak{}Definitions.\allowbreak{}SemilocalFrames.\allowbreak{}PerRowC}, lines 30--35)\\
\texttt{Formalism.\allowbreak{}MainNotions.\allowbreak{}NonGenericAndKTwoTheorems.\allowbreak{}Hyp\_\allowbreak{}of} (\texttt{Formalism.\allowbreak{}SharpRestriction.\allowbreak{}CuspSpanNormBounds.\allowbreak{}ThresholdCubicInequalitySignPattern}, lines 238--260)\par}

{\footnotesize\raggedright\noindent \ref{B8.08h}.\quad \texttt{GlobalTheory.\allowbreak{}Core.\allowbreak{}Diag.\allowbreak{}sDiagAt\_\allowbreak{}of\_\allowbreak{}points} (\texttt{Development.\allowbreak{}AbelianVarietyTheory.\allowbreak{}Proofs.\allowbreak{}DecoratedLatticeProofs.\allowbreak{}DiagonalTheoremOfFourPointBounds}, lines 49--105)\\
\texttt{GlobalTheory.\allowbreak{}Core.\allowbreak{}Rk.\allowbreak{}levelPrime\_\allowbreak{}window} (\texttt{Development.\allowbreak{}AbelianVarietyTheory.\allowbreak{}Proofs.\allowbreak{}DecoratedLatticeProofs.\allowbreak{}LevelPrimeWindow}, lines 11--92)\\
\texttt{GlobalTheory.\allowbreak{}MiTop.\allowbreak{}puNZ\_\allowbreak{}of\_\allowbreak{}sDiag} (\texttt{Development.\allowbreak{}AbelianVarietyTheory.\allowbreak{}Proofs.\allowbreak{}DecoratedLatticeProofs.\allowbreak{}NonvanishingFromTruePack}, lines 49--65)\par}

{\footnotesize\raggedright\noindent \ref{B8.08f}.\quad \texttt{GlobalTheory.\allowbreak{}Core.\allowbreak{}Rk.\allowbreak{}finrank\_\allowbreak{}ker\_\allowbreak{}Pu\_\allowbreak{}le\_\allowbreak{}hsDev} (\texttt{Development.\allowbreak{}AbelianVarietyTheory.\allowbreak{}Proofs.\allowbreak{}DecoratedLatticeProofs.\allowbreak{}FinrankKerPuLeHsDev}, lines 9--17)\\
\texttt{GlobalTheory.\allowbreak{}MiTop.\allowbreak{}puNZ\_\allowbreak{}at} (\texttt{Development.\allowbreak{}AbelianVarietyTheory.\allowbreak{}Proofs.\allowbreak{}DecoratedLatticeProofs.\allowbreak{}NonvanishingFromTruePack}, lines 26--47)\\
\texttt{GlobalTheory.\allowbreak{}MiTop.\allowbreak{}anchor\_\allowbreak{}pairs\_\allowbreak{}alt} (\texttt{Development.\allowbreak{}AutomorphicArithmetic.\allowbreak{}ThetaCorrespondence.\allowbreak{}ThetaPairLifts.\allowbreak{}ThirdPlaceSeries.\allowbreak{}AnchorPairs}, lines 24--71)\par}

{\footnotesize\raggedright\noindent \ref{B8.08i}.\quad \texttt{GlobalTheory.\allowbreak{}MiTop.\allowbreak{}rel\_\allowbreak{}of\_\allowbreak{}links\_\allowbreak{}alt} (\texttt{Development.\allowbreak{}AutomorphicArithmetic.\allowbreak{}ThetaCorrespondence.\allowbreak{}ThetaPairLifts.\allowbreak{}ThirdPlaceSeries.\allowbreak{}RelOfLinks}, lines 22--90)\\
\texttt{GlobalTheory.\allowbreak{}MiTop.\allowbreak{}AtReal.\allowbreak{}anchor\_\allowbreak{}word\_\allowbreak{}alt} (\texttt{Development.\allowbreak{}AutomorphicArithmetic.\allowbreak{}ThetaCorrespondence.\allowbreak{}ThetaPairLifts.\allowbreak{}ThirdPlaceSeries.\allowbreak{}MiStar2AtRealisation}, lines 48--79)\par}

{\footnotesize\raggedright\noindent \ref{B8.08}.\quad \texttt{GlobalTheory.\allowbreak{}T.\allowbreak{}open0\_\allowbreak{}of\_\allowbreak{}unitary\_\allowbreak{}group\_\allowbreak{}hypotheses} (\texttt{Development.\allowbreak{}AutomorphicArithmetic.\allowbreak{}ThetaCorrespondence.\allowbreak{}ThetaPairLifts.\allowbreak{}ThirdPlaceSeries.\allowbreak{}OpenOrbitOfUnramifiedExponents}, lines 75--76)\\
\texttt{GlobalTheory.\allowbreak{}Fc.\allowbreak{}S0.\allowbreak{}Open0ProofU.\allowbreak{}open\_\allowbreak{}orbit\_\allowbreak{}of\_\allowbreak{}unramified\_\allowbreak{}exponents\_\allowbreak{}alt} (\texttt{Development.\allowbreak{}AutomorphicArithmetic.\allowbreak{}ThetaCorrespondence.\allowbreak{}ThetaPairLifts.\allowbreak{}ThirdPlaceSeries.\allowbreak{}OpenOrbitOfUnramifiedExponents}, lines 39--69)\\
\texttt{Formalism.\allowbreak{}IndefR.\allowbreak{}Cx.\allowbreak{}CxData.\allowbreak{}MIid} (\texttt{Formalism.\allowbreak{}LocalToGlobal.\allowbreak{}IndefiniteReal.\allowbreak{}IsometricPeriodsOfHypotheses}, lines 99--103)\\
\texttt{Formalism.\allowbreak{}IndefR.\allowbreak{}Cx.\allowbreak{}CxData.\allowbreak{}MIstar2} (\texttt{Formalism.\allowbreak{}LocalToGlobal.\allowbreak{}IndefiniteReal.\allowbreak{}IsometricPeriodsOfUnitHypotheses}, lines 27--28)\par}

{\footnotesize\raggedright\noindent \ref{B8.08b}.\quad \texttt{GlobalTheory.\allowbreak{}BigTop.\allowbreak{}lamLow\_\allowbreak{}of\_\allowbreak{}hypotheses\_\allowbreak{}alt\_\allowbreak{}alt} (\texttt{Development.\allowbreak{}AutomorphicArithmetic.\allowbreak{}ThetaCorrespondence.\allowbreak{}ThetaPairLifts.\allowbreak{}ThirdPlaceSeries.\allowbreak{}LamLowOfPublished}, lines 17--35)\par}

{\footnotesize\raggedright\noindent \ref{B8.08c}.\quad \texttt{GlobalTheory.\allowbreak{}BigOn.\allowbreak{}canonical\_\allowbreak{}error\_\allowbreak{}real} (\texttt{Development.\allowbreak{}AbelianVarietyTheory.\allowbreak{}Proofs.\allowbreak{}DivisorProofs.\allowbreak{}OnOfOnWithCanonical}, lines 8--62)\\
\texttt{GlobalTheory.\allowbreak{}BigOn.\allowbreak{}OntoPath.\allowbreak{}onto\_\allowbreak{}path\_\allowbreak{}alt} (\texttt{Development.\allowbreak{}AutomorphicArithmetic.\allowbreak{}ThetaCorrespondence.\allowbreak{}ThetaPairLifts.\allowbreak{}ThirdPlaceSeries.\allowbreak{}OntoPath}, lines 18--76)\par}

{\footnotesize\raggedright\noindent \ref{B8.08c}.\quad \texttt{GlobalTheory.\allowbreak{}Glob.\allowbreak{}open\_\allowbreak{}and\_\allowbreak{}linked\_\allowbreak{}of\_\allowbreak{}stronger\_\allowbreak{}extension\_\allowbreak{}alt} (\texttt{Development.\allowbreak{}AutomorphicArithmetic.\allowbreak{}ThetaCorrespondence.\allowbreak{}ThetaPairLifts.\allowbreak{}ThirdPlaceSeries.\allowbreak{}OpenAndLinkedOfStrongerExtension}, lines 28--33)\par}

{\footnotesize\raggedright\noindent \ref{B8.09}.\quad \texttt{GlobalTheory.\allowbreak{}T.\allowbreak{}period\_\allowbreak{}entries\_\allowbreak{}tie\_\allowbreak{}and\_\allowbreak{}level\_\allowbreak{}valuations\_\allowbreak{}at\_\allowbreak{}built} (\texttt{Development.\allowbreak{}AutomorphicArithmetic.\allowbreak{}ThetaCorrespondence.\allowbreak{}ThetaPairLifts.\allowbreak{}ThirdPlaceSeries.\allowbreak{}SixRowsAtMemberStatement}, lines 12--13)\\
\texttt{GlobalTheory.\allowbreak{}Fc.\allowbreak{}six\_\allowbreak{}rows\_\allowbreak{}at\_\allowbreak{}member\_\allowbreak{}alt} (\texttt{Development.\allowbreak{}AutomorphicArithmetic.\allowbreak{}ThetaCorrespondence.\allowbreak{}ThetaPairLifts.\allowbreak{}ThirdPlaceSeries.\allowbreak{}SixRowsAtMember}, lines 11--24)\par}

{\footnotesize\raggedright\noindent \ref{B8.10}.\quad \texttt{GlobalTheory.\allowbreak{}LocalRows.\allowbreak{}lg1\_\allowbreak{}sLp\_\allowbreak{}sO\_\allowbreak{}hTF} (\texttt{Development.\allowbreak{}AutomorphicArithmetic.\allowbreak{}ThetaCorrespondence.\allowbreak{}ThetaPairLifts.\allowbreak{}ThirdPlaceSeries.\allowbreak{}AtChoice}, lines 27--28)\\
\texttt{Formalism.\allowbreak{}MainNotions.\allowbreak{}S.\allowbreak{}sC'} (\texttt{Formalism.\allowbreak{}SharpRestriction.\allowbreak{}CuspSpanNormBounds.\allowbreak{}LeSCSCl}, lines 11--12)\\
\texttt{Formalism.\allowbreak{}MainNotions.\allowbreak{}S.\allowbreak{}sCl} (\texttt{Formalism.\allowbreak{}SharpRestriction.\allowbreak{}CuspSpanNormBounds.\allowbreak{}LeSCSCl}, lines 13--15)\\
\texttt{Formalism.\allowbreak{}MainNotions.\allowbreak{}FamP.\allowbreak{}cLast} (\texttt{Formalism.\allowbreak{}SharpRestriction.\allowbreak{}ParamsHeightRankConstant}, lines 114--118)\\
\texttt{Formalism.\allowbreak{}MainNotions.\allowbreak{}S.\allowbreak{}le\_\allowbreak{}sC'\_\allowbreak{}sCl} (\texttt{Formalism.\allowbreak{}SharpRestriction.\allowbreak{}CuspSpanNormBounds.\allowbreak{}LeSCSCl}, lines 56--62)\\
\texttt{Formalism.\allowbreak{}L2.\allowbreak{}SG.\allowbreak{}ThmSOwn.\allowbreak{}LoopRowS} (\texttt{Formalism.\allowbreak{}LocalToGlobal.\allowbreak{}SiegelLatticeTheory.\allowbreak{}FlagStabilizerTensors.\allowbreak{}ThmSOwn}, lines 1015--1025)\par}

{\footnotesize\raggedright\noindent \ref{B8.10}.\quad \texttt{GlobalTheory.\allowbreak{}LocalRows.\allowbreak{}at\_\allowbreak{}choice\_\allowbreak{}alt} (\texttt{Development.\allowbreak{}AutomorphicArithmetic.\allowbreak{}ThetaCorrespondence.\allowbreak{}ThetaPairLifts.\allowbreak{}ThirdPlaceSeries.\allowbreak{}AtChoice}, lines 16--23)\\
\texttt{Formalism.\allowbreak{}L2.\allowbreak{}SG.\allowbreak{}ThmSOwn.\allowbreak{}offDiagonalK0\_\allowbreak{}of\_\allowbreak{}rows} (\texttt{Formalism.\allowbreak{}LocalToGlobal.\allowbreak{}SiegelLatticeTheory.\allowbreak{}FlagStabilizerTensors.\allowbreak{}LinkSFinalP}, lines 58--115)\par}

{\footnotesize\raggedright\noindent \ref{B8.10c}.\quad \texttt{Formalism.\allowbreak{}L2.\allowbreak{}SG.\allowbreak{}DataCut.\allowbreak{}dom\_\allowbreak{}points} (\texttt{Formalism.\allowbreak{}LocalToGlobal.\allowbreak{}SiegelLatticeTheory.\allowbreak{}CokernelGLCounts.\allowbreak{}RankPointCounts.\allowbreak{}DataCut}, lines 761--895)\\
\texttt{Formalism.\allowbreak{}L2.\allowbreak{}SG.\allowbreak{}DataCut.\allowbreak{}rowsDomAtP\_\allowbreak{}of} (\texttt{Formalism.\allowbreak{}LocalToGlobal.\allowbreak{}SiegelLatticeTheory.\allowbreak{}CokernelGLCounts.\allowbreak{}RankPointCounts.\allowbreak{}RowsInputPOfTermsFiniteSig}, lines 67--163)\\
\texttt{Formalism.\allowbreak{}L2.\allowbreak{}SG.\allowbreak{}DataCut.\allowbreak{}rowsInputP\_\allowbreak{}of} (\texttt{Formalism.\allowbreak{}LocalToGlobal.\allowbreak{}SiegelLatticeTheory.\allowbreak{}CokernelGLCounts.\allowbreak{}RankPointCounts.\allowbreak{}RowsInputPOfTermsFiniteSig}, lines 164--170)\par}

{\footnotesize\raggedright\noindent \ref{B8.10ca}.\quad \texttt{CompanionTheory.\allowbreak{}WaveGrid.\allowbreak{}grids\_\allowbreak{}PStar} (\texttt{CompanionTheory.\allowbreak{}AutomorphicTheory.\allowbreak{}BesselWaveGrids.\allowbreak{}GridsText}, lines 52--64)\\
\texttt{Formalism.\allowbreak{}FourierDecayAlt.\allowbreak{}P.\allowbreak{}TBBx5.\allowbreak{}GridOK} (\texttt{Formalism.\allowbreak{}HermitianCurvature.\allowbreak{}HermitianCurvatureProofs.\allowbreak{}OnePointBasisFundamentalDomainVolume}, lines 162--170)\par}

{\footnotesize\raggedright\noindent \ref{B8.10o}.\quad \texttt{Formalism.\allowbreak{}FourierDecayAlt.\allowbreak{}P.\allowbreak{}CutBP.\allowbreak{}family\_\allowbreak{}expansion\_\allowbreak{}of\_\allowbreak{}parts} (\texttt{Formalism.\allowbreak{}HermitianCurvature.\allowbreak{}HermitianCurvatureProofs.\allowbreak{}NormSqUnitNormalizationLe}, lines 200--339)\\
\texttt{Formalism.\allowbreak{}FourierDecayAlt.\allowbreak{}P.\allowbreak{}CutBP.\allowbreak{}coefficient\_\allowbreak{}growth\_\allowbreak{}le\_\allowbreak{}power} (\texttt{Formalism.\allowbreak{}HermitianCurvature.\allowbreak{}HermitianCurvatureProofs.\allowbreak{}NormSqUnitNormalizationLe}, lines 340--367)\par}

{\footnotesize\raggedright\noindent \ref{B8.10ha}.\quad \texttt{Formalism.\allowbreak{}MainNotions.\allowbreak{}S.\allowbreak{}LocBessel} (\texttt{Formalism.\allowbreak{}SharpRestriction.\allowbreak{}CuspSpanNormBounds.\allowbreak{}RelevantZeroGramTerms}, lines 460--494)\\
\texttt{Formalism.\allowbreak{}MainNotions.\allowbreak{}S.\allowbreak{}Point} (\texttt{Formalism.\allowbreak{}SharpRestriction.\allowbreak{}CuspSpanNormBounds.\allowbreak{}Point}, lines 11--13)\\
\texttt{Formalism.\allowbreak{}MainNotions.\allowbreak{}S.\allowbreak{}ptWeight} (\texttt{Formalism.\allowbreak{}SharpRestriction.\allowbreak{}CuspSpanNormBounds.\allowbreak{}RelevantZeroGramTerms}, lines 246--247)\\
\texttt{Formalism.\allowbreak{}MainNotions.\allowbreak{}BesselShellMomentBounds.\allowbreak{}GPoint\_\allowbreak{}nonneg\_\allowbreak{}of\_\allowbreak{}rep} (\texttt{Formalism.\allowbreak{}SharpRestriction.\allowbreak{}CuspSpanNormBounds.\allowbreak{}BesselShellMomentBounds}, lines 690--697)\\
\texttt{Formalism.\allowbreak{}MainNotions.\allowbreak{}S.\allowbreak{}OrbInt} (\texttt{Formalism.\allowbreak{}SharpRestriction.\allowbreak{}CuspSpanNormBounds.\allowbreak{}RootMeasBig}, lines 61--65)\\
\texttt{Formalism.\allowbreak{}MainNotions.\allowbreak{}S.\allowbreak{}ThB.\allowbreak{}orbSq} (\texttt{Formalism.\allowbreak{}SharpRestriction.\allowbreak{}CuspSpanNormBounds.\allowbreak{}LeadMinorK}, lines 103--107)\par}

{\footnotesize\raggedright\noindent \ref{B8.10hba}.\quad \texttt{GlobalTheory.\allowbreak{}WigSh.\allowbreak{}WigMap} (\texttt{Development.\allowbreak{}AbelianVarietyTheory.\allowbreak{}Definitions.\allowbreak{}Jets.\allowbreak{}WigMap}, lines 12--21)\\
\texttt{GlobalTheory.\allowbreak{}WigSh.\allowbreak{}exists\_\allowbreak{}shifted\_\allowbreak{}types\_\allowbreak{}lower\_\allowbreak{}bound} (\texttt{Development.\allowbreak{}AbelianVarietyTheory.\allowbreak{}Proofs.\allowbreak{}JetHarmonics.\allowbreak{}ExistsShiftedTypesLowerBound}, lines 19--66)\par}

{\footnotesize\raggedright\noindent \ref{B8.10hb}.\quad \texttt{PI.\allowbreak{}BesselPeriodGGP.\allowbreak{}eps48} (\texttt{Formalism.\allowbreak{}PeriodIntegrals.\allowbreak{}BesselPeriodGGP}, lines 24--26)\\
\texttt{CompanionTheory.\allowbreak{}WaveGrid.\allowbreak{}WEnd.\allowbreak{}WBp} (\texttt{CompanionTheory.\allowbreak{}AutomorphicTheory.\allowbreak{}BesselWaveGrids.\allowbreak{}PowFourLeRpowNineHalves}, lines 85--104)\\
\texttt{Formalism.\allowbreak{}L2.\allowbreak{}SG.\allowbreak{}Pf.\allowbreak{}BesselExpansion.\allowbreak{}CoefCC1} (\texttt{Formalism.\allowbreak{}LocalToGlobal.\allowbreak{}SiegelLatticeTheory.\allowbreak{}FlagStabilizerTensors.\allowbreak{}RootLatticeThresholds.\allowbreak{}TBBTie}, lines 30--35)\\
\texttt{CompanionTheory.\allowbreak{}WaveGrid.\allowbreak{}WGlue.\allowbreak{}wLoc1N\_\allowbreak{}of} (\texttt{CompanionTheory.\allowbreak{}AutomorphicTheory.\allowbreak{}BesselWaveGrids.\allowbreak{}WeightedUniformRotatedOfExpansion}, lines 30--160)\\
\texttt{CompanionTheory.\allowbreak{}WaveGrid.\allowbreak{}weighted\_\allowbreak{}bound\_\allowbreak{}rotated\_\allowbreak{}uniform\_\allowbreak{}of\_\allowbreak{}hypsG} (\texttt{CompanionTheory.\allowbreak{}AutomorphicTheory.\allowbreak{}BesselWaveGrids.\allowbreak{}WeightedBoundRotatedUniformOfHypsG}, lines 7--9)\\
\texttt{CompanionTheory.\allowbreak{}PieriExpansion.\allowbreak{}hypsG\_\allowbreak{}of\_\allowbreak{}zuck} (\texttt{CompanionTheory.\allowbreak{}AutomorphicTheory.\allowbreak{}PieriRules.\allowbreak{}InfinitesimalCharacterDiscreteSeries}, lines 51--54)\\
\texttt{CompanionTheory.\allowbreak{}Zuck.\allowbreak{}zuckerman\_\allowbreak{}gap} (\texttt{CompanionTheory.\allowbreak{}AutomorphicTheory.\allowbreak{}ZuckermanCommutation.\allowbreak{}ZuckermanFunctors.\allowbreak{}ZuckermanGapSmall}, lines 9--10)\par}

{\footnotesize\raggedright\noindent \ref{B8.10hc}.\quad \texttt{Formalism.\allowbreak{}MainNotions.\allowbreak{}FrameBoundFromBesselTorus.\allowbreak{}OneBesselTor} (\texttt{Formalism.\allowbreak{}SharpRestriction.\allowbreak{}CuspSpanNormBounds.\allowbreak{}IsPtRepTrDat}, lines 14--30)\\
\texttt{Formalism.\allowbreak{}MainNotions.\allowbreak{}OnePlaceExpansion.\allowbreak{}frame\_\allowbreak{}bound\_\allowbreak{}of\_\allowbreak{}bessel\_\allowbreak{}torus} (\texttt{Formalism.\allowbreak{}SharpRestriction.\allowbreak{}CuspSpanNormBounds.\allowbreak{}OnePlaceExpansion}, lines 71--116)\\
\texttt{CompanionTheory.\allowbreak{}WaveGrid.\allowbreak{}WGlue.\allowbreak{}wLoc1\_\allowbreak{}of\_\allowbreak{}N} (\texttt{CompanionTheory.\allowbreak{}AutomorphicTheory.\allowbreak{}BesselWaveGrids.\allowbreak{}WeightedUniformRotatedOfExpansion}, lines 161--187)\\
\texttt{CompanionTheory.\allowbreak{}WaveGrid.\allowbreak{}WGlue.\allowbreak{}w\_\allowbreak{}in\_\allowbreak{}frame} (\texttt{CompanionTheory.\allowbreak{}AutomorphicTheory.\allowbreak{}BesselWaveGrids.\allowbreak{}WeightedUniformRotatedOfExpansion}, lines 218--274)\\
\texttt{Formalism.\allowbreak{}MainNotions.\allowbreak{}TheoremBConstants.\allowbreak{}num\_\allowbreak{}Fr} (\texttt{Formalism.\allowbreak{}SharpRestriction.\allowbreak{}CuspSpanNormBounds.\allowbreak{}TheoremBConstants}, lines 796--827)\\
\texttt{Formalism.\allowbreak{}MainNotions.\allowbreak{}TheoremBConstants.\allowbreak{}GsC} (\texttt{Formalism.\allowbreak{}SharpRestriction.\allowbreak{}CuspSpanNormBounds.\allowbreak{}TheoremBConstants}, lines 133--134)\\
\texttt{Formalism.\allowbreak{}MainNotions.\allowbreak{}TheoremBConstants.\allowbreak{}FrC} (\texttt{Formalism.\allowbreak{}SharpRestriction.\allowbreak{}CuspSpanNormBounds.\allowbreak{}TheoremBConstants}, lines 129--132)\par}

{\footnotesize\raggedright\noindent \ref{B8.10h}.\quad \texttt{CompanionTheory.\allowbreak{}WaveGrid.\allowbreak{}weighted\_\allowbreak{}bound\_\allowbreak{}rotated\_\allowbreak{}uniform} (\texttt{CompanionTheory.\allowbreak{}AutomorphicTheory.\allowbreak{}BesselWaveGrids.\allowbreak{}ThmBFinal}, lines 7--9)\\
\texttt{CompanionTheory.\allowbreak{}WaveGrid.\allowbreak{}WGlue.\allowbreak{}weighted\_\allowbreak{}uniform\_\allowbreak{}rotated\_\allowbreak{}of\_\allowbreak{}expansion\_\allowbreak{}alt\_\allowbreak{}alt} (\texttt{CompanionTheory.\allowbreak{}AutomorphicTheory.\allowbreak{}BesselWaveGrids.\allowbreak{}WeightedUniformRotatedOfExpansion}, lines 295--431)\\
\texttt{Formalism.\allowbreak{}MainNotions.\allowbreak{}S.\allowbreak{}ThB.\allowbreak{}orbSq} (\texttt{Formalism.\allowbreak{}SharpRestriction.\allowbreak{}CuspSpanNormBounds.\allowbreak{}LeadMinorK}, lines 103--107)\par}

{\footnotesize\raggedright\noindent \ref{B8.10j}.\quad \texttt{GlobalTheory.\allowbreak{}HL.\allowbreak{}QuadricShare.\allowbreak{}shares} (\texttt{Development.\allowbreak{}AbelianVarietyTheory.\allowbreak{}Proofs.\allowbreak{}FlagMeasures.\allowbreak{}QuadricCharacterLeakSharesFinite}, lines 285--373)\\
\texttt{GlobalTheory.\allowbreak{}HL.\allowbreak{}GaussSummand.\allowbreak{}norm\_\allowbreak{}gauss\_\allowbreak{}sum\_\allowbreak{}le} (\texttt{Development.\allowbreak{}AbelianVarietyTheory.\allowbreak{}Proofs.\allowbreak{}FlagMeasures.\allowbreak{}NormGaussSumLe}, lines 48--191)\\
\texttt{GlobalTheory.\allowbreak{}two\_\allowbreak{}geometric\_\allowbreak{}sums\_\allowbreak{}index\_\allowbreak{}bound} (\texttt{Development.\allowbreak{}AbelianVarietyTheory.\allowbreak{}Proofs.\allowbreak{}FlagMeasures.\allowbreak{}TwoGeometricSumsIndexBound}, lines 10--107)\par}

{\footnotesize\raggedright\noindent \ref{B8.10n}.\quad \texttt{Formalism.\allowbreak{}L2.\allowbreak{}SG.\allowbreak{}NGCut.\allowbreak{}TorsionTotientBound.\allowbreak{}card\_\allowbreak{}mul\_\allowbreak{}conj\_\allowbreak{}eq\_\allowbreak{}le\_\allowbreak{}of\_\allowbreak{}fractionalIdeal} (\texttt{Formalism.\allowbreak{}LocalToGlobal.\allowbreak{}SiegelLatticeTheory.\allowbreak{}CokernelGLCounts.\allowbreak{}RankPointCounts.\allowbreak{}CardMulConjEqLeOfFractionalIdeal}, lines 10--103)\\
\texttt{Formalism.\allowbreak{}L2.\allowbreak{}SG.\allowbreak{}NGCut.\allowbreak{}NormalInvariants.\allowbreak{}count\_\allowbreak{}bound} (\texttt{Formalism.\allowbreak{}LocalToGlobal.\allowbreak{}SiegelLatticeTheory.\allowbreak{}CokernelGLCounts.\allowbreak{}RankPointCounts.\allowbreak{}NormalInvariants}, lines 97--126)\par}

{\footnotesize\raggedright\noindent \ref{B8.10da}.\quad \texttt{Formalism.\allowbreak{}FourierDecayAlt.\allowbreak{}PB.\allowbreak{}Terms.\allowbreak{}allPairs} (\texttt{Formalism.\allowbreak{}HermitianCurvature.\allowbreak{}StatementPB}, lines 65--66)\\
\texttt{Formalism.\allowbreak{}MainNotions.\allowbreak{}S.\allowbreak{}Point} (\texttt{Formalism.\allowbreak{}SharpRestriction.\allowbreak{}CuspSpanNormBounds.\allowbreak{}Point}, lines 11--13)\\
\texttt{Formalism.\allowbreak{}MainNotions.\allowbreak{}S.\allowbreak{}ThB.\allowbreak{}leadMinorK} (\texttt{Formalism.\allowbreak{}SharpRestriction.\allowbreak{}CuspSpanNormBounds.\allowbreak{}LeadMinorK}, lines 135--137)\\
\texttt{Formalism.\allowbreak{}MainNotions.\allowbreak{}S.\allowbreak{}ThB.\allowbreak{}Gglob} (\texttt{Formalism.\allowbreak{}SharpRestriction.\allowbreak{}CuspSpanNormBounds.\allowbreak{}LeadMinorK}, lines 145--149)\\
\texttt{Formalism.\allowbreak{}L2.\allowbreak{}SG.\allowbreak{}RowsPos.\allowbreak{}RowNGP} (\texttt{Formalism.\allowbreak{}LocalToGlobal.\allowbreak{}SiegelLatticeTheory.\allowbreak{}FlagStabilizerTensors.\allowbreak{}RowsPos}, lines 37--42)\\
\texttt{Formalism.\allowbreak{}L2.\allowbreak{}SG.\allowbreak{}RowsPos.\allowbreak{}RowMainEqP} (\texttt{Formalism.\allowbreak{}LocalToGlobal.\allowbreak{}SiegelLatticeTheory.\allowbreak{}FlagStabilizerTensors.\allowbreak{}RowsPos}, lines 43--48)\par}

{\footnotesize\raggedright\noindent \ref{B8.10dba}.\quad \texttt{HeckeCharacterZeroLocus.\allowbreak{}SiegelWeilU.\allowbreak{}IsContDensityOn} (\texttt{Formalism.\allowbreak{}UnipotentAdelicQuotients.\allowbreak{}SWXSDefs}, lines 167--172)\\
\texttt{GlobalTheory.\allowbreak{}CompanionTheory.\allowbreak{}SW.\allowbreak{}TopDens.\allowbreak{}siegelDensities} (\texttt{CompanionTheory.\allowbreak{}SpectralGeometry.\allowbreak{}LocalGlobalProofs.\allowbreak{}AdelicAnalysis.\allowbreak{}SiegelDensitiesOfHermitianForm}, lines 47--107)\par}

{\footnotesize\raggedright\noindent \ref{B8.10db}.\quad \texttt{GlobalTheory.\allowbreak{}CompanionTheory.\allowbreak{}SW.\allowbreak{}siegelWeilUnitaryConvergent} (\texttt{CompanionTheory.\allowbreak{}SpectralGeometry.\allowbreak{}LocalGlobalProofs.\allowbreak{}AdelicAnalysis.\allowbreak{}Imports}, lines 10--13)\\
\texttt{GlobalTheory.\allowbreak{}CompanionTheory.\allowbreak{}SW.\allowbreak{}finVol\_\allowbreak{}unitaryGroup} (\texttt{CompanionTheory.\allowbreak{}SpectralGeometry.\allowbreak{}LocalGlobalProofs.\allowbreak{}AdelicAnalysis.\allowbreak{}UnitaryGroupFiniteVolume}, lines 8--10)\\
\texttt{GlobalTheory.\allowbreak{}T.\allowbreak{}siegel\_\allowbreak{}densities\_\allowbreak{}of\_\allowbreak{}hermitian\_\allowbreak{}form} (\texttt{CompanionTheory.\allowbreak{}SpectralGeometry.\allowbreak{}LocalGlobalProofs.\allowbreak{}AdelicAnalysis.\allowbreak{}SiegelDensitiesOfHermitianFormStatement}, lines 7--9)\\
\texttt{Formalism.\allowbreak{}L2.\allowbreak{}SG.\allowbreak{}RealObj.\allowbreak{}EnergyConvJ.\allowbreak{}BM.\allowbreak{}realObj\_\allowbreak{}inInfinityX\_\allowbreak{}jsz} (\texttt{CompanionTheory.\allowbreak{}AutomorphicTheory.\allowbreak{}LocalBesselUniqueness.\allowbreak{}BesselModels.\allowbreak{}RealObjInInfinityXJsz}, lines 17--21)\\
\texttt{GlobalTheory.\allowbreak{}St.\allowbreak{}stepChoiceDataC\_\allowbreak{}of\_\allowbreak{}siegelWeil\_\allowbreak{}inInfinity} (\texttt{Development.\allowbreak{}AbelianVarietyTheory.\allowbreak{}Proofs.\allowbreak{}FlagMeasures.\allowbreak{}StepChoiceConstantsCOfSiegelWeilAtInfinity}, lines 9--16)\par}

{\footnotesize\raggedright\noindent \ref{B8.10dca}.\quad \texttt{Formalism.\allowbreak{}L2.\allowbreak{}SG.\allowbreak{}Pf.\allowbreak{}BDatR.\allowbreak{}jointAmplitudeFactorizationRotated\_\allowbreak{}of} (\texttt{Formalism.\allowbreak{}LocalToGlobal.\allowbreak{}SiegelLatticeTheory.\allowbreak{}FlagStabilizerTensors.\allowbreak{}RootLatticeThresholds.\allowbreak{}BDataOfR}, lines 117--173)\\
\texttt{Formalism.\allowbreak{}L2.\allowbreak{}SG.\allowbreak{}Pf.\allowbreak{}ThmBChainG.\allowbreak{}jointAmplitudeRotated\_\allowbreak{}generic} (\texttt{Formalism.\allowbreak{}LocalToGlobal.\allowbreak{}SiegelLatticeTheory.\allowbreak{}FlagStabilizerTensors.\allowbreak{}RootLatticeThresholds.\allowbreak{}ThmBChainG}, lines 20--72)\\
\texttt{Formalism.\allowbreak{}MainNotions.\allowbreak{}S.\allowbreak{}RMeas} (\texttt{Formalism.\allowbreak{}SharpRestriction.\allowbreak{}CuspSpanNormBounds.\allowbreak{}NormalizedLocalScale}, lines 163--167)\\
\texttt{Formalism.\allowbreak{}MainNotions.\allowbreak{}S.\allowbreak{}GoodW} (\texttt{Formalism.\allowbreak{}SharpRestriction.\allowbreak{}CuspSpanNormBounds.\allowbreak{}NormalizedLocalScale}, lines 127--136)\\
\texttt{Formalism.\allowbreak{}MainNotions.\allowbreak{}S.\allowbreak{}IsQuotMeas} (\texttt{Formalism.\allowbreak{}SharpRestriction.\allowbreak{}CuspSpanNormBounds.\allowbreak{}RelevantZeroGramTerms}, lines 208--211)\\
\texttt{Formalism.\allowbreak{}L2.\allowbreak{}SG.\allowbreak{}NGData.\allowbreak{}a2C} (\texttt{Formalism.\allowbreak{}LocalToGlobal.\allowbreak{}SiegelLatticeTheory.\allowbreak{}CokernelGLCounts.\allowbreak{}NGData}, lines 43--45)\\
\texttt{Formalism.\allowbreak{}L2.\allowbreak{}SG.\allowbreak{}Pf.\allowbreak{}MainEqCut.\allowbreak{}a2GC} (\texttt{Formalism.\allowbreak{}LocalToGlobal.\allowbreak{}SiegelLatticeTheory.\allowbreak{}CokernelGLCounts.\allowbreak{}RankPointCounts.\allowbreak{}MainEqCut}, lines 122--125)\par}

{\footnotesize\raggedright\noindent \ref{B8.10dc}.\quad \texttt{Formalism.\allowbreak{}MainStatements.\allowbreak{}RowRest.\allowbreak{}analytic\_\allowbreak{}pinnedTheoremB\_\allowbreak{}uniform} (\texttt{CompanionTheory.\allowbreak{}AutomorphicTheory.\allowbreak{}UniformLocalFormulas.\allowbreak{}Closed}, lines 20--23)\\
\texttt{Formalism.\allowbreak{}MainStatements.\allowbreak{}RowRest.\allowbreak{}bDat\_\allowbreak{}theoremB\_\allowbreak{}global} (\texttt{CompanionTheory.\allowbreak{}AutomorphicTheory.\allowbreak{}UniformLocalFormulas.\allowbreak{}Closed}, lines 24--27)\par}

{\footnotesize\raggedright\noindent \ref{B8.10d}.\quad \texttt{GlobalTheory.\allowbreak{}rowNGP\_\allowbreak{}of\_\allowbreak{}siegelWeil} (\texttt{Development.\allowbreak{}AbelianVarietyTheory.\allowbreak{}Proofs.\allowbreak{}FlagMeasures.\allowbreak{}RowPositivityOfSiegelWeil}, lines 10--16)\\
\texttt{GlobalTheory.\allowbreak{}rowMainEqP\_\allowbreak{}of\_\allowbreak{}siegelWeil} (\texttt{Development.\allowbreak{}AbelianVarietyTheory.\allowbreak{}Proofs.\allowbreak{}FlagMeasures.\allowbreak{}RowMainEqPOfSiegelWeil}, lines 10--17)\par}

{\footnotesize\raggedright\noindent \ref{B8.10l}.\quad \texttt{Formalism.\allowbreak{}HermitianLatticeCounting.\allowbreak{}P.\allowbreak{}nondegenerateRowBound} (\texttt{Formalism.\allowbreak{}FlagUnipotentLattices.\allowbreak{}MinkowskiFactorialBounds.\allowbreak{}LatticeSumBounds}, lines 914--968)\\
\texttt{Formalism.\allowbreak{}HermitianLatticeCounting.\allowbreak{}P.\allowbreak{}degenerateRowBound} (\texttt{Formalism.\allowbreak{}FlagUnipotentLattices.\allowbreak{}MinkowskiFactorialBounds.\allowbreak{}LatticeSumBounds}, lines 969--1029)\par}

{\footnotesize\raggedright\noindent \ref{B8.10m}.\quad \texttt{Formalism.\allowbreak{}HermitianLatticeCounting.\allowbreak{}P.\allowbreak{}prodMinLeOfProdRatio} (\texttt{Formalism.\allowbreak{}FlagUnipotentLattices.\allowbreak{}MinkowskiFactorialBounds.\allowbreak{}LatticeSumBounds}, lines 1150--1218)\\
\texttt{Formalism.\allowbreak{}MainNotions.\allowbreak{}BridgeLemma₂.\allowbreak{}rowsT\_\allowbreak{}small\_\allowbreak{}one} (\texttt{Formalism.\allowbreak{}SharpRestriction.\allowbreak{}CuspSpanNormBounds.\allowbreak{}CayleyDegreeNondegeneracy.\allowbreak{}RowsTrm}, lines 839--894)\\
\texttt{Formalism.\allowbreak{}MainNotions.\allowbreak{}BridgeLemma₂.\allowbreak{}rowsT\_\allowbreak{}Mt\_\allowbreak{}small} (\texttt{Formalism.\allowbreak{}SharpRestriction.\allowbreak{}CuspSpanNormBounds.\allowbreak{}CayleyDegreeNondegeneracy.\allowbreak{}RowsTrm}, lines 1057--1128)\\
\texttt{Formalism.\allowbreak{}MainNotions.\allowbreak{}BridgeLemma₂.\allowbreak{}small\_\allowbreak{}term\_\allowbreak{}pow\_\allowbreak{}le} (\texttt{Formalism.\allowbreak{}SharpRestriction.\allowbreak{}CuspSpanNormBounds.\allowbreak{}CayleyDegreeNondegeneracy.\allowbreak{}RowsOneC}, lines 152--196)\par}

{\footnotesize\raggedright\noindent \ref{B8.10ab}.\quad \texttt{Formalism.\allowbreak{}FourierDecayAlt.\allowbreak{}PB.\allowbreak{}Terms.\allowbreak{}allPairs} (\texttt{Formalism.\allowbreak{}HermitianCurvature.\allowbreak{}StatementPB}, lines 65--66)\\
\texttt{Formalism.\allowbreak{}L2.\allowbreak{}SG.\allowbreak{}TermOp.\allowbreak{}wt} (\texttt{Formalism.\allowbreak{}LocalToGlobal.\allowbreak{}SiegelLatticeTheory.\allowbreak{}ClosedLocalOrbits.\allowbreak{}TermOp}, lines 341--342)\\
\texttt{Formalism.\allowbreak{}L2.\allowbreak{}SG.\allowbreak{}TermOp.\allowbreak{}vec} (\texttt{Formalism.\allowbreak{}LocalToGlobal.\allowbreak{}SiegelLatticeTheory.\allowbreak{}ClosedLocalOrbits.\allowbreak{}TermOp}, lines 343--346)\\
\texttt{Formalism.\allowbreak{}MainNotions.\allowbreak{}S.\allowbreak{}ThB.\allowbreak{}leadMinorK} (\texttt{Formalism.\allowbreak{}SharpRestriction.\allowbreak{}CuspSpanNormBounds.\allowbreak{}LeadMinorK}, lines 135--137)\\
\texttt{Formalism.\allowbreak{}MainNotions.\allowbreak{}S.\allowbreak{}ThB.\allowbreak{}Gglob} (\texttt{Formalism.\allowbreak{}SharpRestriction.\allowbreak{}CuspSpanNormBounds.\allowbreak{}LeadMinorK}, lines 145--149)\par}

{\footnotesize\raggedright\noindent \ref{B8.10ac}.\quad \texttt{GlobalTheory.\allowbreak{}RowsFive.\allowbreak{}rowsND\_\allowbreak{}of} (\texttt{Development.\allowbreak{}AbelianVarietyTheory.\allowbreak{}Proofs.\allowbreak{}FlagMeasures.\allowbreak{}AdmissibleFamilySatisfiesRowEstimates}, lines 29--132)\\
\texttt{GlobalTheory.\allowbreak{}RowsFive.\allowbreak{}rowsDeg\_\allowbreak{}of} (\texttt{Development.\allowbreak{}AbelianVarietyTheory.\allowbreak{}Proofs.\allowbreak{}FlagMeasures.\allowbreak{}AdmissibleFamilySatisfiesRowEstimates}, lines 133--252)\\
\texttt{Formalism.\allowbreak{}MainNotions.\allowbreak{}BridgeLemma₂.\allowbreak{}rows\_\allowbreak{}ND\_\allowbreak{}record} (\texttt{Formalism.\allowbreak{}SharpRestriction.\allowbreak{}CuspSpanNormBounds.\allowbreak{}CayleyDegreeNondegeneracy.\allowbreak{}RowsND}, lines 324--366)\\
\texttt{Formalism.\allowbreak{}MainNotions.\allowbreak{}BridgeLemma₂.\allowbreak{}rows\_\allowbreak{}nd\_\allowbreak{}prod} (\texttt{Formalism.\allowbreak{}SharpRestriction.\allowbreak{}CuspSpanNormBounds.\allowbreak{}CayleyDegreeNondegeneracy.\allowbreak{}RowsND}, lines 191--277)\\
\texttt{Formalism.\allowbreak{}MainNotions.\allowbreak{}BridgeLemma₂.\allowbreak{}rows\_\allowbreak{}deg\_\allowbreak{}count} (\texttt{Formalism.\allowbreak{}SharpRestriction.\allowbreak{}CuspSpanNormBounds.\allowbreak{}CayleyDegreeNondegeneracy.\allowbreak{}RowsDeg}, lines 185--232)\\
\texttt{Formalism.\allowbreak{}MainNotions.\allowbreak{}BridgeLemma₂.\allowbreak{}rows\_\allowbreak{}deg\_\allowbreak{}prod} (\texttt{Formalism.\allowbreak{}SharpRestriction.\allowbreak{}CuspSpanNormBounds.\allowbreak{}CayleyDegreeNondegeneracy.\allowbreak{}RowsDeg}, lines 233--330)\\
\texttt{Formalism.\allowbreak{}MainNotions.\allowbreak{}BridgeLemma₂.\allowbreak{}rows\_\allowbreak{}cF\_\allowbreak{}le} (\texttt{Formalism.\allowbreak{}SharpRestriction.\allowbreak{}CuspSpanNormBounds.\allowbreak{}CayleyDegreeNondegeneracy.\allowbreak{}RowsNDn}, lines 61--80)\\
\texttt{Formalism.\allowbreak{}MainNotions.\allowbreak{}BridgeLemma₂.\allowbreak{}rows\_\allowbreak{}q\_\allowbreak{}bound} (\texttt{Formalism.\allowbreak{}SharpRestriction.\allowbreak{}CuspSpanNormBounds.\allowbreak{}CayleyDegreeNondegeneracy.\allowbreak{}RowsNDF}, lines 11--40)\par}

{\footnotesize\raggedright\noindent \ref{B8.10ada}.\quad \texttt{Formalism.\allowbreak{}MainNotions.\allowbreak{}BridgeLemma₂.\allowbreak{}rowsT\_\allowbreak{}small\_\allowbreak{}one} (\texttt{Formalism.\allowbreak{}SharpRestriction.\allowbreak{}CuspSpanNormBounds.\allowbreak{}CayleyDegreeNondegeneracy.\allowbreak{}RowsTrm}, lines 839--894)\par}

{\footnotesize\raggedright\noindent \ref{B8.10ad}.\quad \texttt{Formalism.\allowbreak{}MainNotions.\allowbreak{}S.\allowbreak{}Cnt.\allowbreak{}Hyp} (\texttt{Formalism.\allowbreak{}SharpRestriction.\allowbreak{}CuspSpanNormBounds.\allowbreak{}CntPartners}, lines 36--49)\\
\texttt{Formalism.\allowbreak{}MainNotions.\allowbreak{}BridgeLemma₂.\allowbreak{}small\_\allowbreak{}term\_\allowbreak{}pow\_\allowbreak{}le} (\texttt{Formalism.\allowbreak{}SharpRestriction.\allowbreak{}CuspSpanNormBounds.\allowbreak{}CayleyDegreeNondegeneracy.\allowbreak{}RowsOneC}, lines 152--196)\\
\texttt{Formalism.\allowbreak{}MainNotions.\allowbreak{}BridgeLemma.\allowbreak{}RowsOne} (\texttt{Formalism.\allowbreak{}SharpRestriction.\allowbreak{}CuspSpanNormBounds.\allowbreak{}BruhatCellsKronecker.\allowbreak{}FamilyPacketOfEquivalentStatement}, lines 389--389)\par}

{\footnotesize\raggedright\noindent \ref{B8.10a}.\quad \texttt{GlobalTheory.\allowbreak{}LoopRowBoundsGap.\allowbreak{}rows\_\allowbreak{}along\_\allowbreak{}family} (\texttt{Development.\allowbreak{}AbelianVarietyTheory.\allowbreak{}Proofs.\allowbreak{}FlagMeasures.\allowbreak{}LoopRowBoundsAlongFamilyLargeGaps}, lines 25--44)\\
\texttt{Formalism.\allowbreak{}L2.\allowbreak{}SG.\allowbreak{}LoopsTies.\allowbreak{}loop\_\allowbreak{}evaluation\_\allowbreak{}threshold\_\allowbreak{}zero} (\texttt{Formalism.\allowbreak{}LocalToGlobal.\allowbreak{}SiegelLatticeTheory.\allowbreak{}FlagStabilizerTensors.\allowbreak{}RootLatticeThresholds.\allowbreak{}LoopEvaluationThresholdZero}, lines 44--120)\\
\texttt{Formalism.\allowbreak{}FourierDecayAlt.\allowbreak{}P.\allowbreak{}GrowthPolynomialOperator.\allowbreak{}t\_\allowbreak{}LoopsJoint} (\texttt{Formalism.\allowbreak{}HermitianCurvature.\allowbreak{}HermitianCurvatureProofs.\allowbreak{}GrowthOfBasic}, lines 1191--1229)\par}

{\footnotesize\raggedright\noindent \ref{B8.10aa}.\quad \texttt{Formalism.\allowbreak{}FourierDecayAlt.\allowbreak{}P.\allowbreak{}Jx2.\allowbreak{}t\_\allowbreak{}LatticeSumsJoint} (\texttt{Formalism.\allowbreak{}HermitianCurvature.\allowbreak{}HermitianCurvatureProofs.\allowbreak{}LevelSizesJointMean}, lines 124--167)\\
\texttt{Formalism.\allowbreak{}FourierDecayAlt.\allowbreak{}S6.\allowbreak{}perLat} (\texttt{Formalism.\allowbreak{}HermitianCurvature.\allowbreak{}PrimeToDifferent}, lines 51--53)\\
\texttt{Formalism.\allowbreak{}FourierDecayAlt.\allowbreak{}S6.\allowbreak{}liftLat} (\texttt{Formalism.\allowbreak{}HermitianCurvature.\allowbreak{}PrimeToDifferent}, lines 54--56)\par}

{\footnotesize\raggedright\noindent \ref{B8.10g}.\quad \texttt{Formalism.\allowbreak{}FourierDecayAlt.\allowbreak{}P.\allowbreak{}change\_\allowbreak{}of\_\allowbreak{}variables\_\allowbreak{}xy} (\texttt{Formalism.\allowbreak{}HermitianCurvature.\allowbreak{}HermitianCurvatureProofs.\allowbreak{}ChangeOfVariablesXy}, lines 22--24)\\
\texttt{Formalism.\allowbreak{}FourierDecayAlt.\allowbreak{}P.\allowbreak{}germ\_\allowbreak{}rate} (\texttt{Formalism.\allowbreak{}HermitianCurvature.\allowbreak{}HermitianCurvatureProofs.\allowbreak{}ChangeOfVariablesXy}, lines 34--38)\par}

{\footnotesize\raggedright\noindent \ref{B8.10i}.\quad \texttt{Formalism.\allowbreak{}FourierDecayAlt.\allowbreak{}P.\allowbreak{}PBK.\allowbreak{}memLp\_\allowbreak{}bz\_\allowbreak{}of\_\allowbreak{}sup} (\texttt{Formalism.\allowbreak{}HermitianCurvature.\allowbreak{}HermitianCurvatureProofs.\allowbreak{}MemLpBzOfSup}, lines 16--79)\\
\texttt{Formalism.\allowbreak{}FourierDecayAlt.\allowbreak{}P.\allowbreak{}PBK.\allowbreak{}GermD.\allowbreak{}tendsto\_\allowbreak{}cross} (\texttt{Formalism.\allowbreak{}HermitianCurvature.\allowbreak{}HermitianCurvatureProofs.\allowbreak{}TendstoRatioPow}, lines 77--122)\\
\texttt{Formalism.\allowbreak{}FourierDecayAlt.\allowbreak{}P.\allowbreak{}PBK.\allowbreak{}GermD.\allowbreak{}germLoc\_\allowbreak{}tail} (\texttt{Formalism.\allowbreak{}HermitianCurvature.\allowbreak{}HermitianCurvatureProofs.\allowbreak{}GermLocTail}, lines 15--48)\par}

{\footnotesize\raggedright\noindent \ref{B8.10p}.\quad \texttt{Formalism.\allowbreak{}FourierDecayAlt.\allowbreak{}P.\allowbreak{}FormNormSumIdentities.\allowbreak{}cutN\_\allowbreak{}smooth} (\texttt{Formalism.\allowbreak{}HermitianCurvature.\allowbreak{}HermitianCurvatureProofs.\allowbreak{}PbGermRateLowBOfParts}, lines 132--141)\\
\texttt{Formalism.\allowbreak{}FourierDecayAlt.\allowbreak{}P.\allowbreak{}FormNormSumIdentities.\allowbreak{}cutN\_\allowbreak{}supp} (\texttt{Formalism.\allowbreak{}HermitianCurvature.\allowbreak{}HermitianCurvatureProofs.\allowbreak{}PbGermRateLowBOfParts}, lines 142--176)\\
\texttt{Formalism.\allowbreak{}FourierDecayAlt.\allowbreak{}P.\allowbreak{}JFxG.\allowbreak{}cptInCC\_\allowbreak{}G} (\texttt{Formalism.\allowbreak{}HermitianCurvature.\allowbreak{}HermitianCurvatureProofs.\allowbreak{}JfBetaAdjCOfLemmaBNorm}, lines 24--44)\par}

{\footnotesize\raggedright\noindent \ref{B8.10e}.\quad \texttt{Formalism.\allowbreak{}FourierDecayAlt.\allowbreak{}P.\allowbreak{}GrowthPolynomialOperator.\allowbreak{}t\_\allowbreak{}MainRow} (\texttt{Formalism.\allowbreak{}HermitianCurvature.\allowbreak{}HermitianCurvatureProofs.\allowbreak{}GrowthOfBasic}, lines 1070--1184)\par}

{\footnotesize\raggedright\noindent \ref{B8.10f}.\quad \texttt{Formalism.\allowbreak{}L2.\allowbreak{}SG.\allowbreak{}GGMirabolic.\allowbreak{}termsGGBound\_\allowbreak{}of\_\allowbreak{}matrix} (\texttt{Formalism.\allowbreak{}LocalToGlobal.\allowbreak{}SiegelLatticeTheory.\allowbreak{}ClosedLocalOrbits.\allowbreak{}GGMirabolic}, lines 136--145)\par}

{\footnotesize\raggedright\noindent \ref{B8.10f}.\quad \texttt{Formalism.\allowbreak{}FourierDecayAlt.\allowbreak{}PB.\allowbreak{}Terms.\allowbreak{}NormExpansionOn} (\texttt{Formalism.\allowbreak{}HermitianCurvature.\allowbreak{}StatementPB}, lines 86--87)\\
\texttt{Formalism.\allowbreak{}FourierDecayAlt.\allowbreak{}PB.\allowbreak{}Terms.\allowbreak{}RowOrth} (\texttt{Formalism.\allowbreak{}HermitianCurvature.\allowbreak{}StatementPB}, lines 92--94)\par}

{\footnotesize\raggedright\noindent \ref{B8.10b}.\quad \texttt{Formalism.\allowbreak{}FourierDecayAlt.\allowbreak{}P.\allowbreak{}GrowthPolynomialOperator.\allowbreak{}t\_\allowbreak{}MainRow} (\texttt{Formalism.\allowbreak{}HermitianCurvature.\allowbreak{}HermitianCurvatureProofs.\allowbreak{}GrowthOfBasic}, lines 1070--1184)\\
\texttt{Formalism.\allowbreak{}FourierDecayAlt.\allowbreak{}P.\allowbreak{}GrowthPolynomialOperator.\allowbreak{}t\_\allowbreak{}PairBound} (\texttt{Formalism.\allowbreak{}HermitianCurvature.\allowbreak{}HermitianCurvatureProofs.\allowbreak{}GrowthOfBasic}, lines 525--551)\\
\texttt{Formalism.\allowbreak{}L2.\allowbreak{}SG.\allowbreak{}TheoremSConclusion.\allowbreak{}theoremS\_\allowbreak{}of\_\allowbreak{}conj} (\texttt{Formalism.\allowbreak{}LocalToGlobal.\allowbreak{}SiegelLatticeTheory.\allowbreak{}FlagStabilizerTensors.\allowbreak{}TheoremSConclusion}, lines 68--83)\\
\texttt{Formalism.\allowbreak{}L2.\allowbreak{}SG.\allowbreak{}FixedDecomposition.\allowbreak{}kppOffC\_\allowbreak{}of\_\allowbreak{}steps0} (\texttt{Formalism.\allowbreak{}LocalToGlobal.\allowbreak{}SiegelLatticeTheory.\allowbreak{}CokernelGLCounts.\allowbreak{}RankPointCounts.\allowbreak{}FixedDecomposition}, lines 755--775)\par}

\section*{Chapter~\ref{ch:p8c2}: Theta correspondence, test vectors and evaluation}

{\footnotesize\raggedright\noindent \ref{eq:p8c2-split-matching}.\quad \texttt{GlobalTheory.\allowbreak{}SplitFlip.\allowbreak{}bigP} (\texttt{Development.\allowbreak{}AbelianVarietyTheory.\allowbreak{}Definitions.\allowbreak{}Divisors.\allowbreak{}RankEqCardIffExistsFrameMatrixFirstColumnBlock}, lines 716--719)\\
\texttt{GlobalTheory.\allowbreak{}SplitFlip.\allowbreak{}smallP} (\texttt{Development.\allowbreak{}AbelianVarietyTheory.\allowbreak{}Definitions.\allowbreak{}Divisors.\allowbreak{}RankEqCardIffExistsFrameMatrixFirstColumnBlock}, lines 720--722)\par}

{\footnotesize\raggedright\noindent \ref{lem:p8c2-hecke-nonvanishing}.\quad \texttt{GlobalTheory.\allowbreak{}ThK.\allowbreak{}Dbl.\allowbreak{}heckeL\_\allowbreak{}ne\_\allowbreak{}zero} (\texttt{Development.\allowbreak{}WeilBound.\allowbreak{}ThetaLifts.\allowbreak{}LatticeBasics.\allowbreak{}ExistsIsHeckeLOfCompleted}, lines 134--161)\\
\texttt{GlobalTheory.\allowbreak{}HeckePartial.\allowbreak{}removed\_\allowbreak{}ne\_\allowbreak{}zero} (\texttt{Development.\allowbreak{}AbelianVarietyTheory.\allowbreak{}PublishedTheorems.\allowbreak{}TwistedTraceFormula.\allowbreak{}HeckeLPartialContinuationHolds}, lines 40--43)\par}

{\footnotesize\raggedright\noindent \ref{lem:p8c2-hecke-nonvanishing}.\quad \texttt{Formalism.\allowbreak{}BasicNotions.\allowbreak{}TH.\allowbreak{}fockAct1} (\texttt{Formalism.\allowbreak{}LocalRingCardinality.\allowbreak{}DataTheta}, lines 210--211)\\
\texttt{Formalism.\allowbreak{}BasicNotions.\allowbreak{}TH.\allowbreak{}fockAct2} (\texttt{Formalism.\allowbreak{}LocalRingCardinality.\allowbreak{}DataTheta}, lines 213--214)\par}

{\footnotesize\raggedright\noindent \ref{B8.11i}.\quad \texttt{PositivityTheory.\allowbreak{}KReal.\allowbreak{}nontrivial\_\allowbreak{}of\_\allowbreak{}pair} (\texttt{Development.\allowbreak{}WeilBound.\allowbreak{}CoreStatements.\allowbreak{}GaussianSeminormBounds.\allowbreak{}RealSeminormFiltration.\allowbreak{}CoreNonv}, lines 160--238)\\
\texttt{PositivityTheory.\allowbreak{}KReal.\allowbreak{}pair\_\allowbreak{}pos} (\texttt{Development.\allowbreak{}WeilBound.\allowbreak{}CoreStatements.\allowbreak{}GaussianSeminormBounds.\allowbreak{}RealSeminormFiltration.\allowbreak{}CorePair}, lines 101--114)\\
\texttt{PositivityTheory.\allowbreak{}KReal.\allowbreak{}pair\_\allowbreak{}neg} (\texttt{Development.\allowbreak{}WeilBound.\allowbreak{}CoreStatements.\allowbreak{}GaussianSeminormBounds.\allowbreak{}RealSeminormFiltration.\allowbreak{}CorePair}, lines 115--127)\\
\texttt{PositivityTheory.\allowbreak{}KReal.\allowbreak{}side\_\allowbreak{}of} (\texttt{Development.\allowbreak{}WeilBound.\allowbreak{}CoreStatements.\allowbreak{}GaussianSeminormBounds.\allowbreak{}RealSeminormFiltration.\allowbreak{}CoreSide}, lines 159--211)\\
\texttt{PositivityTheory.\allowbreak{}KReal.\allowbreak{}la\_\allowbreak{}eq} (\texttt{Development.\allowbreak{}WeilBound.\allowbreak{}CoreStatements.\allowbreak{}GaussianSeminormBounds.\allowbreak{}RealSeminormFiltration.\allowbreak{}CoreComb}, lines 355--370)\par}

{\footnotesize\raggedright\noindent \ref{B8.11j}.\quad \texttt{PositivityTheory.\allowbreak{}KReal.\allowbreak{}star\_\allowbreak{}th} (\texttt{Development.\allowbreak{}WeilBound.\allowbreak{}CoreStatements.\allowbreak{}GaussianSeminormBounds.\allowbreak{}RealSeminormFiltration.\allowbreak{}CoreStar}, lines 77--166)\\
\texttt{PositivityTheory.\allowbreak{}KReal.\allowbreak{}root\_\allowbreak{}of\_\allowbreak{}hw} (\texttt{Development.\allowbreak{}WeilBound.\allowbreak{}CoreStatements.\allowbreak{}GaussianSeminormBounds.\allowbreak{}RealSeminormFiltration.\allowbreak{}CoreRoot}, lines 18--135)\\
\texttt{PositivityTheory.\allowbreak{}KReal.\allowbreak{}comb\_\allowbreak{}pos} (\texttt{Development.\allowbreak{}WeilBound.\allowbreak{}CoreStatements.\allowbreak{}GaussianSeminormBounds.\allowbreak{}RealSeminormFiltration.\allowbreak{}CoreComb}, lines 119--208)\\
\texttt{PositivityTheory.\allowbreak{}KReal.\allowbreak{}comb\_\allowbreak{}neg} (\texttt{Development.\allowbreak{}WeilBound.\allowbreak{}CoreStatements.\allowbreak{}GaussianSeminormBounds.\allowbreak{}RealSeminormFiltration.\allowbreak{}CoreComb}, lines 209--304)\par}

{\footnotesize\raggedright\noindent \ref{B8.11k}.\quad \texttt{PositivityTheory.\allowbreak{}KReal.\allowbreak{}cas\_\allowbreak{}bound\_\allowbreak{}gl} (\texttt{Development.\allowbreak{}WeilBound.\allowbreak{}CoreStatements.\allowbreak{}GaussianSeminormBounds.\allowbreak{}RealSeminormFiltration.\allowbreak{}CoreCasGl}, lines 158--228)\\
\texttt{PositivityTheory.\allowbreak{}KReal.\allowbreak{}common\_\allowbreak{}quot} (\texttt{Development.\allowbreak{}WeilBound.\allowbreak{}CoreStatements.\allowbreak{}GaussianSeminormBounds.\allowbreak{}RealSeminormFiltration.\allowbreak{}CoreMain}, lines 583--713)\\
\texttt{PositivityTheory.\allowbreak{}KReal.\allowbreak{}coreLM} (\texttt{Development.\allowbreak{}WeilBound.\allowbreak{}CoreStatements.\allowbreak{}GaussianSeminormBounds.\allowbreak{}RealSeminormFiltration.\allowbreak{}CoreLM}, lines 14--30)\\
\texttt{PositivityTheory.\allowbreak{}KRealA.\allowbreak{}onePlace} (\texttt{Development.\allowbreak{}WeilBound.\allowbreak{}CoreStatements.\allowbreak{}GaussianSeminormBounds.\allowbreak{}RealSeminormFiltration.\allowbreak{}KRealAsmOne}, lines 143--155)\par}

{\footnotesize\raggedright\noindent \ref{B8.11l}.\quad \texttt{GlobalTheory.\allowbreak{}ThK.\allowbreak{}Dbl.\allowbreak{}bump} (\texttt{Development.\allowbreak{}WeilBound.\allowbreak{}ThetaLifts.\allowbreak{}LatticeBasics.\allowbreak{}NearUnits}, lines 113--185)\\
\texttt{GlobalTheory.\allowbreak{}ThK.\allowbreak{}Dbl.\allowbreak{}nonZeroBottom\_\allowbreak{}of\_\allowbreak{}doubling\_\allowbreak{}integral} (\texttt{Development.\allowbreak{}WeilBound.\allowbreak{}ThetaLifts.\allowbreak{}LatticeBasics.\allowbreak{}BotAsm}, lines 117--147)\\
\texttt{GlobalTheory.\allowbreak{}ThK.\allowbreak{}Dbl.\allowbreak{}pair\_\allowbreak{}ne\_\allowbreak{}zero} (\texttt{Development.\allowbreak{}WeilBound.\allowbreak{}ThetaLifts.\allowbreak{}LatticeBasics.\allowbreak{}BotAsm}, lines 58--116)\\
\texttt{GlobalTheory.\allowbreak{}ThK.\allowbreak{}Dbl.\allowbreak{}badFactor\_\allowbreak{}of} (\texttt{Development.\allowbreak{}WeilBound.\allowbreak{}ThetaLifts.\allowbreak{}LatticeBasics.\allowbreak{}BadFactor}, lines 96--192)\par}

{\footnotesize\raggedright\noindent \ref{B8.11m}.\quad \texttt{GlobalTheory.\allowbreak{}ThK.\allowbreak{}Dbl.\allowbreak{}LD.\allowbreak{}eq\_\allowbreak{}case} (\texttt{Development.\allowbreak{}WeilBound.\allowbreak{}ThetaLifts.\allowbreak{}LatticeBasics.\allowbreak{}LocDistE}, lines 106--144)\\
\texttt{GlobalTheory.\allowbreak{}ThK.\allowbreak{}Dbl.\allowbreak{}TC.\allowbreak{}inv\_\allowbreak{}loc} (\texttt{Development.\allowbreak{}WeilBound.\allowbreak{}ThetaLifts.\allowbreak{}LatticeBasics.\allowbreak{}ThConstInv}, lines 86--153)\\
\texttt{GlobalTheory.\allowbreak{}ThK.\allowbreak{}Dbl.\allowbreak{}thConst\_\allowbreak{}of'} (\texttt{Development.\allowbreak{}WeilBound.\allowbreak{}ThetaLifts.\allowbreak{}LatticeBasics.\allowbreak{}ThConstFin}, lines 13--16)\\
\texttt{GlobalTheory.\allowbreak{}ThK.\allowbreak{}Dbl.\allowbreak{}normOne\_\allowbreak{}mem} (\texttt{Development.\allowbreak{}WeilBound.\allowbreak{}ThetaLifts.\allowbreak{}LatticeBasics.\allowbreak{}ConstIdele}, lines 134--182)\par}

{\footnotesize\raggedright\noindent \ref{B8.11}.\quad \texttt{UpperTheory.\allowbreak{}LiteratureTheorems.\allowbreak{}theta\_\allowbreak{}integrals\_\allowbreak{}agree\_\allowbreak{}and\_\allowbreak{}local\_\allowbreak{}matching\_\allowbreak{}of\_\allowbreak{}literature} (\texttt{Development.\allowbreak{}AutomorphicArithmetic.\allowbreak{}ThetaCorrespondence.\allowbreak{}ThetaPairLifts.\allowbreak{}ThirdPlaceSeries.\allowbreak{}EveryCharacterOfLiteratureOfVWin}, lines 565--566)\\
\texttt{PositivityTheory.\allowbreak{}ThKR.\allowbreak{}theoremK\_\allowbreak{}head2} (\texttt{Development.\allowbreak{}WeilBound.\allowbreak{}ThetaLifts.\allowbreak{}Reductions.\allowbreak{}FiniteClosure.\allowbreak{}TheoremKOfMod}, lines 26--27)\\
\texttt{GlobalTheory.\allowbreak{}IndefiniteRowSideFullHecke.\allowbreak{}WeilTensorModels.\allowbreak{}ThetaIntegralsAgree} (\texttt{Development.\allowbreak{}AbelianVarietyTheory.\allowbreak{}BlockTheory.\allowbreak{}DivisorBase.\allowbreak{}WeilTensorModels}, lines 585--592)\\
\texttt{GlobalTheory.\allowbreak{}IndefiniteRowSideFullHecke.\allowbreak{}WeilTensorModels.\allowbreak{}LocalMatching} (\texttt{Development.\allowbreak{}AbelianVarietyTheory.\allowbreak{}BlockTheory.\allowbreak{}DivisorBase.\allowbreak{}WeilTensorModels}, lines 581--583)\par}

{\footnotesize\raggedright\noindent \ref{B8.11}.\quad \texttt{GlobalTheory.\allowbreak{}IndefiniteRowSideFullHecke.\allowbreak{}WeilTensorModels.\allowbreak{}StepBeta} (\texttt{Development.\allowbreak{}AbelianVarietyTheory.\allowbreak{}BlockTheory.\allowbreak{}DivisorBase.\allowbreak{}WeilTensorModels}, lines 559--567)\\
\texttt{GlobalTheory.\allowbreak{}IndefiniteRowSideFullHecke.\allowbreak{}WeilTensorModels.\allowbreak{}StepAlpha} (\texttt{Development.\allowbreak{}AbelianVarietyTheory.\allowbreak{}BlockTheory.\allowbreak{}DivisorBase.\allowbreak{}WeilTensorModels}, lines 571--579)\par}

{\footnotesize\raggedright\noindent \ref{B8.11}.\quad \texttt{PositivityTheory.\allowbreak{}ThKG.\allowbreak{}Leaves.\allowbreak{}const\_\allowbreak{}of} (\texttt{Development.\allowbreak{}WeilBound.\allowbreak{}ThetaLifts.\allowbreak{}ConeHypotheses.\allowbreak{}CtSmallCone}, lines 88--136)\\
\texttt{PositivityTheory.\allowbreak{}ThKG.\allowbreak{}Leaves.\allowbreak{}isConst\_\allowbreak{}of} (\texttt{Development.\allowbreak{}WeilBound.\allowbreak{}ThetaLifts.\allowbreak{}ConeHypotheses.\allowbreak{}CtSmallCone}, lines 179--189)\\
\texttt{PositivityTheory.\allowbreak{}ThKG.\allowbreak{}Leaves.\allowbreak{}kConst} (\texttt{Development.\allowbreak{}WeilBound.\allowbreak{}ThetaLifts.\allowbreak{}ConeHypotheses.\allowbreak{}CtSmallCone}, lines 191--198)\\
\texttt{PositivityTheory.\allowbreak{}ThKG.\allowbreak{}Leaves.\allowbreak{}kBeta} (\texttt{Development.\allowbreak{}WeilBound.\allowbreak{}ThetaLifts.\allowbreak{}ConeHypotheses.\allowbreak{}CtSmallCone}, lines 200--218)\\
\texttt{PositivityTheory.\allowbreak{}ThKG.\allowbreak{}Leaves.\allowbreak{}kAlpha} (\texttt{Development.\allowbreak{}WeilBound.\allowbreak{}ThetaLifts.\allowbreak{}ConeHypotheses.\allowbreak{}CtSmallCone}, lines 220--231)\\
\texttt{PositivityTheory.\allowbreak{}ThKEnd.\allowbreak{}Two.\allowbreak{}twoStep\_\allowbreak{}of'} (\texttt{Development.\allowbreak{}WeilBound.\allowbreak{}ThetaLifts.\allowbreak{}EigenvalueEndgame.\allowbreak{}SignatureSplitting.\allowbreak{}TwoStepOf}, lines 22--102)\par}

{\footnotesize\raggedright\noindent \ref{B8.11e}.\quad \texttt{PositivityTheory.\allowbreak{}ThKEst.\allowbreak{}sqInt} (\texttt{Development.\allowbreak{}WeilBound.\allowbreak{}ThetaLifts.\allowbreak{}FrameEstimates.\allowbreak{}WeilConv}, lines 68--70)\\
\texttt{PositivityTheory.\allowbreak{}ThKEst.\allowbreak{}iCont} (\texttt{Development.\allowbreak{}WeilBound.\allowbreak{}ThetaLifts.\allowbreak{}FrameEstimates.\allowbreak{}WeilConv}, lines 71--73)\par}

{\footnotesize\raggedright\noindent \ref{B8.11f}.\quad \texttt{PositivityTheory.\allowbreak{}Sph.\allowbreak{}SplitPl.\allowbreak{}irrAt\_\allowbreak{}ne} (\texttt{Development.\allowbreak{}WeilBound.\allowbreak{}ThetaLifts.\allowbreak{}SphericalSplitting.\allowbreak{}SphSplitGeo}, lines 21--46)\\
\texttt{GlobalTheory.\allowbreak{}CompanionTheory.\allowbreak{}SoCoinvNarrow.\allowbreak{}soCoinv\_\allowbreak{}of} (\texttt{Development.\allowbreak{}AutomorphicSpectrum.\allowbreak{}Proofs.\allowbreak{}BlockDecomposition.\allowbreak{}SoCoinvOfNarrow}, lines 197--246)\par}

{\footnotesize\raggedright\noindent \ref{B8.11g}.\quad \texttt{PositivityTheory.\allowbreak{}ThKF.\allowbreak{}factor\_\allowbreak{}of} (\texttt{Development.\allowbreak{}WeilBound.\allowbreak{}ThetaLifts.\allowbreak{}CartanFactorisation.\allowbreak{}Factor}, lines 117--179)\\
\texttt{PositivityTheory.\allowbreak{}Sph.\allowbreak{}theta\_\allowbreak{}kills\_\allowbreak{}cusp\_\allowbreak{}forms\_\allowbreak{}spherical} (\texttt{Development.\allowbreak{}WeilBound.\allowbreak{}ThetaLifts.\allowbreak{}SphericalLattices.\allowbreak{}ThetaKillsCuspFormsSpherical}, lines 154--204)\par}

{\footnotesize\raggedright\noindent \ref{lem:p8c2-real-vanishing}.\quad \texttt{PositivityTheory.\allowbreak{}ThKEnd.\allowbreak{}Comp.\allowbreak{}IZero} (\texttt{Development.\allowbreak{}WeilBound.\allowbreak{}ThetaLifts.\allowbreak{}EigenvalueEndgame.\allowbreak{}EigenKilling.\allowbreak{}ConeIntegralFactorisation}, lines 76--80)\\
\texttt{PositivityTheory.\allowbreak{}ThKEnd.\allowbreak{}Comp.\allowbreak{}realVanish\_\allowbreak{}of} (\texttt{Development.\allowbreak{}WeilBound.\allowbreak{}ThetaLifts.\allowbreak{}EigenvalueEndgame.\allowbreak{}EigenKilling.\allowbreak{}RealVanishOf}, lines 60--112)\par}

{\footnotesize\raggedright\noindent \ref{B8.11d}.\quad \texttt{GlobalTheory.\allowbreak{}ThK.\allowbreak{}Dbl.\allowbreak{}coefLine\_\allowbreak{}of} (\texttt{Development.\allowbreak{}WeilBound.\allowbreak{}ThetaLifts.\allowbreak{}LatticeBasics.\allowbreak{}CoefLine}, lines 13--62)\par}

{\footnotesize\raggedright\noindent \ref{B8.11a}.\quad \texttt{PositivityTheory.\allowbreak{}ThKG.\allowbreak{}Tow.\allowbreak{}Pln.\allowbreak{}exists\_\allowbreak{}plane} (\texttt{Development.\allowbreak{}WeilBound.\allowbreak{}ThetaLifts.\allowbreak{}TowerModels.\allowbreak{}PlaneMain}, lines 26--136)\\
\texttt{PositivityTheory.\allowbreak{}ThKG.\allowbreak{}Tow.\allowbreak{}towerB\_\allowbreak{}of} (\texttt{Development.\allowbreak{}WeilBound.\allowbreak{}ThetaLifts.\allowbreak{}TowerModels.\allowbreak{}TowerCodimensionTwoOf}, lines 32--73)\\
\texttt{PositivityTheory.\allowbreak{}ThKG.\allowbreak{}Tow.\allowbreak{}towerA\_\allowbreak{}of} (\texttt{Development.\allowbreak{}WeilBound.\allowbreak{}ThetaLifts.\allowbreak{}TowerModels.\allowbreak{}TowerCodimensionTwoOf}, lines 75--126)\\
\texttt{PositivityTheory.\allowbreak{}ThKG.\allowbreak{}Tow.\allowbreak{}towerCodimensionTwo\_\allowbreak{}of} (\texttt{Development.\allowbreak{}WeilBound.\allowbreak{}ThetaLifts.\allowbreak{}TowerModels.\allowbreak{}TowerCodimensionTwoOf}, lines 128--172)\par}

{\footnotesize\raggedright\noindent \ref{B8.11b}.\quad \texttt{PositivityTheory.\allowbreak{}ThKCT.\allowbreak{}genSmall\_\allowbreak{}of} (\texttt{Development.\allowbreak{}WeilBound.\allowbreak{}ThetaLifts.\allowbreak{}SmallCones.\allowbreak{}GenSmallOf}, lines 18--128)\\
\texttt{PositivityTheory.\allowbreak{}ThKG.\allowbreak{}Loc.\allowbreak{}Cone.\allowbreak{}Id.\allowbreak{}idealId'} (\texttt{Development.\allowbreak{}WeilBound.\allowbreak{}ThetaLifts.\allowbreak{}LocalTheory.\allowbreak{}ConeShifts.\allowbreak{}IdealId0}, lines 49--73)\\
\texttt{PositivityTheory.\allowbreak{}ThKG.\allowbreak{}Std.\allowbreak{}eps\_\allowbreak{}eq} (\texttt{Development.\allowbreak{}WeilBound.\allowbreak{}ThetaLifts.\allowbreak{}LocalTheory.\allowbreak{}StandardValuations.\allowbreak{}ExistsVol}, lines 85--99)\\
\texttt{PositivityTheory.\allowbreak{}ThKG.\allowbreak{}Loc.\allowbreak{}locB\_\allowbreak{}of} (\texttt{Development.\allowbreak{}WeilBound.\allowbreak{}ThetaLifts.\allowbreak{}LocalTheory.\allowbreak{}LocBOf}, lines 17--249)\\
\texttt{PositivityTheory.\allowbreak{}ThKG.\allowbreak{}Loc.\allowbreak{}locA\_\allowbreak{}of} (\texttt{Development.\allowbreak{}WeilBound.\allowbreak{}ThetaLifts.\allowbreak{}LocalTheory.\allowbreak{}LocAOf}, lines 17--173)\\
\texttt{PositivityTheory.\allowbreak{}ThKG.\allowbreak{}keyB\_\allowbreak{}of} (\texttt{Development.\allowbreak{}WeilBound.\allowbreak{}ThetaLifts.\allowbreak{}ConeHypotheses.\allowbreak{}KeyAOf}, lines 37--70)\\
\texttt{PositivityTheory.\allowbreak{}ThKG.\allowbreak{}keyA\_\allowbreak{}of} (\texttt{Development.\allowbreak{}WeilBound.\allowbreak{}ThetaLifts.\allowbreak{}ConeHypotheses.\allowbreak{}KeyAOf}, lines 88--126)\par}

{\footnotesize\raggedright\noindent \ref{B8.11c}.\quad \texttt{GlobalTheory.\allowbreak{}ThK.\allowbreak{}Dbl.\allowbreak{}ct\_\allowbreak{}zero} (\texttt{Development.\allowbreak{}WeilBound.\allowbreak{}ThetaLifts.\allowbreak{}LatticeBasics.\allowbreak{}CuspLift}, lines 52--162)\\
\texttt{GlobalTheory.\allowbreak{}ThK.\allowbreak{}Dbl.\allowbreak{}cusp\_\allowbreak{}lift} (\texttt{Development.\allowbreak{}WeilBound.\allowbreak{}ThetaLifts.\allowbreak{}LatticeBasics.\allowbreak{}OrthAn}, lines 73--201)\\
\texttt{GlobalTheory.\allowbreak{}ThK.\allowbreak{}Dbl.\allowbreak{}orthAn\_\allowbreak{}of} (\texttt{Development.\allowbreak{}WeilBound.\allowbreak{}ThetaLifts.\allowbreak{}LatticeBasics.\allowbreak{}OrthAn}, lines 333--409)\\
\texttt{PositivityTheory.\allowbreak{}ThKEnd.\allowbreak{}Cusp.\allowbreak{}Key.\allowbreak{}keyThree\_\allowbreak{}of\_\allowbreak{}locThree} (\texttt{Development.\allowbreak{}WeilBound.\allowbreak{}ThetaLifts.\allowbreak{}EigenvalueEndgame.\allowbreak{}CuspidalAverages.\allowbreak{}ActsOfReads}, lines 495--505)\\
\texttt{PositivityTheory.\allowbreak{}ThKEnd.\allowbreak{}Cusp.\allowbreak{}Key.\allowbreak{}keyTwo\_\allowbreak{}of\_\allowbreak{}locTwo\_\allowbreak{}vacBar} (\texttt{Development.\allowbreak{}WeilBound.\allowbreak{}ThetaLifts.\allowbreak{}EigenvalueEndgame.\allowbreak{}CuspidalAverages.\allowbreak{}ActsOfReads}, lines 629--704)\\
\texttt{PositivityTheory.\allowbreak{}ThKEnd.\allowbreak{}Cusp.\allowbreak{}kill\_\allowbreak{}coef} (\texttt{Development.\allowbreak{}WeilBound.\allowbreak{}ThetaLifts.\allowbreak{}EigenvalueEndgame.\allowbreak{}CuspidalAverages.\allowbreak{}KillCoef}, lines 14--60)\\
\texttt{PositivityTheory.\allowbreak{}ThKEnd.\allowbreak{}Cusp.\allowbreak{}cuspOfLine\_\allowbreak{}of\_\allowbreak{}keys} (\texttt{Development.\allowbreak{}WeilBound.\allowbreak{}ThetaLifts.\allowbreak{}EigenvalueEndgame.\allowbreak{}CuspidalAverages.\allowbreak{}CuspOfLineOfKeys}, lines 169--198)\\
\texttt{PositivityTheory.\allowbreak{}ThKEnd.\allowbreak{}zero\_\allowbreak{}of\_\allowbreak{}line} (\texttt{Development.\allowbreak{}WeilBound.\allowbreak{}ThetaLifts.\allowbreak{}EigenvalueEndgame.\allowbreak{}TowerCodimensionTwo}, lines 73--85)\par}

{\footnotesize\raggedright\noindent \ref{B8.12}.\quad \texttt{GlobalTheory.\allowbreak{}ThK.\allowbreak{}Dbl.\allowbreak{}EisBL} (\texttt{Development.\allowbreak{}WeilBound.\allowbreak{}ThetaLifts.\allowbreak{}LatticeBasics.\allowbreak{}HypothesesE}, lines 46--53)\\
\texttt{GlobalTheory.\allowbreak{}ThK.\allowbreak{}Dbl.\allowbreak{}SiegelWeilStd} (\texttt{Development.\allowbreak{}WeilBound.\allowbreak{}ThetaLifts.\allowbreak{}LatticeBasics.\allowbreak{}HypothesesS}, lines 11--21)\\
\texttt{UpperTheory.\allowbreak{}LiteratureTheorems.\allowbreak{}eis\_\allowbreak{}bl\_\allowbreak{}of\_\allowbreak{}literature} (\texttt{Development.\allowbreak{}AutomorphicArithmetic.\allowbreak{}ThetaCorrespondence.\allowbreak{}ThetaPairLifts.\allowbreak{}ThirdPlaceSeries.\allowbreak{}EveryCharacterOfLiteratureOfVWin}, lines 561--562)\\
\texttt{UpperTheory.\allowbreak{}LiteratureTheorems.\allowbreak{}siegel\_\allowbreak{}weil\_\allowbreak{}std\_\allowbreak{}of\_\allowbreak{}literature} (\texttt{Development.\allowbreak{}AutomorphicArithmetic.\allowbreak{}ThetaCorrespondence.\allowbreak{}ThetaPairLifts.\allowbreak{}ThirdPlaceSeries.\allowbreak{}EveryCharacterOfLiteratureOfVWin}, lines 563--564)\par}

{\footnotesize\raggedright\noindent \ref{B8.12}.\quad \texttt{PositivityTheory.\allowbreak{}SiegelWeilFormula.\allowbreak{}Fin.\allowbreak{}siegelWeilStd} (\texttt{Development.\allowbreak{}WeilBound.\allowbreak{}CoreStatements.\allowbreak{}AutomorphicKernels.\allowbreak{}SiegelWeilFormula.\allowbreak{}SiegelWeilStandard.\allowbreak{}SiegelWeilStd}, lines 8--8)\par}

{\footnotesize\raggedright\noindent \ref{B8.13}.\quad \texttt{GlobalTheory.\allowbreak{}Published.\allowbreak{}SiegelWeilUnitaryConvergent} (\texttt{Development.\allowbreak{}AbelianVarietyTheory.\allowbreak{}PublishedTheorems.\allowbreak{}SiegelWeil₂.\allowbreak{}SiegelWeilUnitaryConvergent}, lines 18--32)\\
\texttt{HeckeCharacterZeroLocus.\allowbreak{}SiegelWeilU.\allowbreak{}SiegelWeilUnitaryConvergent} (\texttt{Formalism.\allowbreak{}UnipotentAdelicQuotients.\allowbreak{}SiegelWeilUnitaryConvergent}, lines 46--49)\\
\texttt{HeckeCharacterZeroLocus.\allowbreak{}SiegelWeilU.\allowbreak{}SiegelDensities} (\texttt{Formalism.\allowbreak{}UnipotentAdelicQuotients.\allowbreak{}SiegelWeilUnitaryConvergent}, lines 18--40)\\
\texttt{GlobalTheory.\allowbreak{}Published.\allowbreak{}RallisInnerProductConvergent} (\texttt{Development.\allowbreak{}AbelianVarietyTheory.\allowbreak{}PublishedTheorems.\allowbreak{}TwistedTraceFormula.\allowbreak{}LowerLiftsZeroOfWittIdxEqZero}, lines 41--59)\\
\texttt{UpperTheory.\allowbreak{}LiteratureTheorems.\allowbreak{}siegel\_\allowbreak{}weil\_\allowbreak{}unitary\_\allowbreak{}convergent\_\allowbreak{}of\_\allowbreak{}literature} (\texttt{Development.\allowbreak{}AutomorphicArithmetic.\allowbreak{}ThetaCorrespondence.\allowbreak{}ThetaPairLifts.\allowbreak{}ThirdPlaceSeries.\allowbreak{}EveryCharacterOfLiteratureOfVWin}, lines 481--482)\\
\texttt{UpperTheory.\allowbreak{}LiteratureTheorems.\allowbreak{}siegelWeilU\_\allowbreak{}unitary\_\allowbreak{}convergent\_\allowbreak{}of\_\allowbreak{}literature} (\texttt{Development.\allowbreak{}AutomorphicArithmetic.\allowbreak{}ThetaCorrespondence.\allowbreak{}ThetaPairLifts.\allowbreak{}ThirdPlaceSeries.\allowbreak{}EveryCharacterOfLiteratureOfVWin}, lines 581--582)\\
\texttt{UpperTheory.\allowbreak{}LiteratureTheorems.\allowbreak{}rallis\_\allowbreak{}inner\_\allowbreak{}product\_\allowbreak{}convergent\_\allowbreak{}of\_\allowbreak{}literature} (\texttt{Development.\allowbreak{}AutomorphicArithmetic.\allowbreak{}ThetaCorrespondence.\allowbreak{}ThetaPairLifts.\allowbreak{}ThirdPlaceSeries.\allowbreak{}EveryCharacterOfLiteratureOfVWin}, lines 485--486)\par}

{\footnotesize\raggedright\noindent \ref{B8.13}.\quad \texttt{UpperTheory.\allowbreak{}Lit.\allowbreak{}Rallis.\allowbreak{}RallisIPProof.\allowbreak{}rallis\_\allowbreak{}of\_\allowbreak{}dm} (\texttt{Development.\allowbreak{}AutomorphicArithmetic.\allowbreak{}ThetaCorrespondence.\allowbreak{}FockModels.\allowbreak{}RallisDM}, lines 24--110)\par}

{\footnotesize\raggedright\noindent \ref{B8.14}.\quad \texttt{GlobalTheory.\allowbreak{}Hk3.\allowbreak{}howeUnram\_\allowbreak{}holds} (\texttt{Development.\allowbreak{}AbelianVarietyTheory.\allowbreak{}Proofs.\allowbreak{}SatakeParameters.\allowbreak{}RhoMemOfMkEq}, lines 250--256)\\
\texttt{GlobalTheory.\allowbreak{}Hk3.\allowbreak{}HoweUnram} (\texttt{Development.\allowbreak{}AbelianVarietyTheory.\allowbreak{}PublishedTheorems.\allowbreak{}GaloisGeometry₂.\allowbreak{}HoweUnram}, lines 8--19)\par}

{\footnotesize\raggedright\noindent \ref{B8.14}.\quad \texttt{GlobalTheory.\allowbreak{}Hk3.\allowbreak{}locVac\_\allowbreak{}span\_\allowbreak{}classes\_\allowbreak{}indep} (\texttt{Development.\allowbreak{}AbelianVarietyTheory.\allowbreak{}Proofs.\allowbreak{}SatakeParameters.\allowbreak{}LatticeModelVacuumTranslateClassesIndependent}, lines 8--98)\\
\texttt{GlobalTheory.\allowbreak{}Hk3.\allowbreak{}HoweG.\allowbreak{}gramClass\_\allowbreak{}orbit\_\allowbreak{}mod\_\allowbreak{}lower} (\texttt{Development.\allowbreak{}AbelianVarietyTheory.\allowbreak{}Proofs.\allowbreak{}SatakeParameters.\allowbreak{}ExistsConjugateModelExtremalDualParameter}, lines 16--129)\\
\texttt{GlobalTheory.\allowbreak{}Hk3.\allowbreak{}HoweDuality.\allowbreak{}fixed\_\allowbreak{}mem\_\allowbreak{}span} (\texttt{Development.\allowbreak{}AbelianVarietyTheory.\allowbreak{}Proofs.\allowbreak{}SatakeParameters.\allowbreak{}RhoMemOfMkEq}, lines 145--240)\par}

{\footnotesize\raggedright\noindent \ref{def:p8c2-point-sizes}.\quad \texttt{Formalism.\allowbreak{}MainNotions.\allowbreak{}SizesAt} (\texttt{Formalism.\allowbreak{}SharpRestriction.\allowbreak{}ElevenLe}, lines 10--139)\par}

{\footnotesize\raggedright\noindent \ref{B8.14a}.\quad \texttt{PositivityTheory.\allowbreak{}Sph.\allowbreak{}isVacSph\_\allowbreak{}of} (\texttt{Development.\allowbreak{}WeilBound.\allowbreak{}ThetaLifts.\allowbreak{}SphericalLattices.\allowbreak{}ExistsFixedSum}, lines 90--134)\\
\texttt{PositivityTheory.\allowbreak{}Sph.\allowbreak{}howeInert} (\texttt{Development.\allowbreak{}WeilBound.\allowbreak{}ThetaLifts.\allowbreak{}SphericalLattices.\allowbreak{}HoweDuality}, lines 13--29)\par}

{\footnotesize\raggedright\noindent \ref{B8.15d}.\quad \texttt{GlobalTheory.\allowbreak{}MoyalTheta.\allowbreak{}orthogonality} (\texttt{Development.\allowbreak{}AbelianVarietyTheory.\allowbreak{}Proofs.\allowbreak{}TwistedOperators.\allowbreak{}MoyalOrthogonalityHeisenbergThetaFunctions}, lines 25--76)\\
\texttt{GlobalTheory.\allowbreak{}Asc.\allowbreak{}annVac} (\texttt{Development.\allowbreak{}AbelianVarietyTheory.\allowbreak{}Proofs.\allowbreak{}TwistedOperators.\allowbreak{}HeisenbergVacuumCommutantPreservesLatticeVacua}, lines 36--78)\\
\texttt{GlobalTheory.\allowbreak{}Irr.\allowbreak{}irr} (\texttt{Development.\allowbreak{}AbelianVarietyTheory.\allowbreak{}Proofs.\allowbreak{}TwistedOperators.\allowbreak{}TameFunctionalVacuumIrreducibility}, lines 57--77)\par}

{\footnotesize\raggedright\noindent \ref{B8.15e}.\quad \texttt{GlobalTheory.\allowbreak{}ConjDataTotal.\allowbreak{}fockTot} (\texttt{Development.\allowbreak{}AbelianVarietyTheory.\allowbreak{}Proofs.\allowbreak{}TwistedOperators.\allowbreak{}FockConjugatesSpanDepthData}, lines 190--271)\\
\texttt{GlobalTheory.\allowbreak{}FirstRowDepth.\allowbreak{}datumDep} (\texttt{Development.\allowbreak{}AbelianVarietyTheory.\allowbreak{}Proofs.\allowbreak{}TwistedOperators.\allowbreak{}FirstRowDataHaveDepth}, lines 19--30)\par}

{\footnotesize\raggedright\noindent \ref{B8.15f}.\quad \texttt{GlobalTheory.\allowbreak{}Asc.\allowbreak{}ascent\_\allowbreak{}of\_\allowbreak{}depFixed} (\texttt{Development.\allowbreak{}AbelianVarietyTheory.\allowbreak{}Proofs.\allowbreak{}TwistedOperators.\allowbreak{}BesselPeriodNonvanishingAscentOfDepthData}, lines 21--26)\\
\texttt{GlobalTheory.\allowbreak{}Asc.\allowbreak{}ascent\_\allowbreak{}of\_\allowbreak{}rungs} (\texttt{Development.\allowbreak{}AbelianVarietyTheory.\allowbreak{}Proofs.\allowbreak{}TwistedOperators.\allowbreak{}BesselPeriodAscentOfTowerRungs}, lines 15--42)\par}

{\footnotesize\raggedright\noindent \ref{def:p8c2-series}.\quad \texttt{GlobalTheory.\allowbreak{}GLq.\allowbreak{}ESerTie} (\texttt{Development.\allowbreak{}AbelianVarietyTheory.\allowbreak{}Definitions.\allowbreak{}NonGaloisDescent.\allowbreak{}LabelDatum}, lines 114--117)\\
\texttt{GlobalTheory.\allowbreak{}GLq.\allowbreak{}eSerC} (\texttt{Development.\allowbreak{}AbelianVarietyTheory.\allowbreak{}Definitions.\allowbreak{}NonGaloisDescent.\allowbreak{}LabelDatum}, lines 102--104)\par}

{\footnotesize\raggedright\noindent \ref{B8.15}.\quad \texttt{UpperTheory.\allowbreak{}Lit.\allowbreak{}functionalsSeparate\_\allowbreak{}of\_\allowbreak{}howeUnram\_\allowbreak{}of\_\allowbreak{}siegelWeil} (\texttt{Development.\allowbreak{}AutomorphicArithmetic.\allowbreak{}ThetaCorrespondence.\allowbreak{}ThetaPairLifts.\allowbreak{}LocalTypeDictionary.\allowbreak{}FunctionalsSeparateOfHoweUnramOfSiegelWeil}, lines 17--19)\\
\texttt{UpperTheory.\allowbreak{}Lit.\allowbreak{}real\_\allowbreak{}functionals\_\allowbreak{}separate\_\allowbreak{}of\_\allowbreak{}hypothesis} (\texttt{Development.\allowbreak{}AutomorphicArithmetic.\allowbreak{}ThetaCorrespondence.\allowbreak{}ThetaPairLifts.\allowbreak{}LocalTypeDictionary.\allowbreak{}RealFunctionalsSeparateOfHypothesis}, lines 15--30)\\
\texttt{GlobalTheory.\allowbreak{}IndefiniteRowSideFullHecke.\allowbreak{}LevelPlusSubgroup.\allowbreak{}FunctionalsSeparate} (\texttt{Development.\allowbreak{}AbelianVarietyTheory.\allowbreak{}BlockTheory.\allowbreak{}DivisorBase.\allowbreak{}LevelPlusSubgroup}, lines 83--91)\par}

{\footnotesize\raggedright\noindent \ref{B8.15}.\quad \texttt{GlobalTheory.\allowbreak{}SepRows.\allowbreak{}weakSep\_\allowbreak{}of\_\allowbreak{}steps} (\texttt{Development.\allowbreak{}AbelianVarietyTheory.\allowbreak{}Proofs.\allowbreak{}TwistedOperators.\allowbreak{}RowSeparationByRealFunctionalsOfChainSteps}, lines 71--85)\\
\texttt{UpperTheory.\allowbreak{}Lit.\allowbreak{}Bur.\allowbreak{}separate\_\allowbreak{}of\_\allowbreak{}burGapPos} (\texttt{Development.\allowbreak{}AutomorphicArithmetic.\allowbreak{}ThetaCorrespondence.\allowbreak{}BesselModels.\allowbreak{}SepGap}, lines 222--257)\par}

{\footnotesize\raggedright\noindent \ref{B8.15g}.\quad \texttt{PositivityTheory.\allowbreak{}UnitaryDecomposition.\allowbreak{}DecB.\allowbreak{}dec\_\allowbreak{}U} (\texttt{Development.\allowbreak{}WeilBound.\allowbreak{}CoreStatements.\allowbreak{}BranchingTraces.\allowbreak{}CasimirDecomposition.\allowbreak{}UnitaryDecompositionExists}, lines 45--177)\\
\texttt{PositivityTheory.\allowbreak{}UnitaryDecomposition.\allowbreak{}Com.\allowbreak{}comat\_\allowbreak{}of} (\texttt{Development.\allowbreak{}WeilBound.\allowbreak{}CoreStatements.\allowbreak{}BranchingTraces.\allowbreak{}BranchingComatrices.\allowbreak{}OneLeDiff}, lines 136--276)\\
\texttt{CompanionTheory.\allowbreak{}Lime.\allowbreak{}T3.\allowbreak{}kFiniteWordBound\_\allowbreak{}indefUnitary} (\texttt{Development.\allowbreak{}WeilBound.\allowbreak{}CoreStatements.\allowbreak{}OrbitDecayEstimates.\allowbreak{}CurveDerivativeDecay.\allowbreak{}KFiniteWordBoundIndefUnitary}, lines 245--358)\\
\texttt{PositivityTheory.\allowbreak{}UnitaryDecomposition.\allowbreak{}fdeg\_\allowbreak{}shift} (\texttt{Development.\allowbreak{}WeilBound.\allowbreak{}CoreStatements.\allowbreak{}BranchingTraces.\allowbreak{}ContinuousEmb}, lines 190--239)\par}

{\footnotesize\raggedright\noindent \ref{B8.15h}.\quad \texttt{PositivityTheory.\allowbreak{}RgHe.\allowbreak{}A.\allowbreak{}Pf.\allowbreak{}bigProjP\_\allowbreak{}of} (\texttt{Development.\allowbreak{}WeilBound.\allowbreak{}RepresentationTheoryGL.\allowbreak{}HolomorphicDiscreteSeries.\allowbreak{}SchurEigenProjections.\allowbreak{}SchurFiniteDimensionality.\allowbreak{}IsFdOfSchurK}, lines 126--239)\par}

{\footnotesize\raggedright\noindent \ref{B8.15i}.\quad \texttt{PositivityTheory.\allowbreak{}UnitaryDecomposition.\allowbreak{}momentHeA\_\allowbreak{}of} (\texttt{Development.\allowbreak{}WeilBound.\allowbreak{}RepresentationTheoryGL.\allowbreak{}HolomorphicDiscreteSeries.\allowbreak{}SchurEigenProjections.\allowbreak{}HolomorphicMoments.\allowbreak{}MomentHigherEstimateHolds}, lines 28--179)\\
\texttt{PositivityTheory.\allowbreak{}UnitaryDecomposition.\allowbreak{}moment\_\allowbreak{}higher\_\allowbreak{}estimate\_\allowbreak{}holds} (\texttt{Development.\allowbreak{}WeilBound.\allowbreak{}RepresentationTheoryGL.\allowbreak{}HolomorphicDiscreteSeries.\allowbreak{}SchurEigenProjections.\allowbreak{}HolomorphicMoments.\allowbreak{}MomentHigherEstimateHolds}, lines 181--182)\par}

{\footnotesize\raggedright\noindent \ref{B8.15a}.\quad \texttt{UpperTheory.\allowbreak{}Lit.\allowbreak{}Branch.\allowbreak{}packEx\_\allowbreak{}of\_\allowbreak{}eigen} (\texttt{Development.\allowbreak{}AutomorphicArithmetic.\allowbreak{}ThetaCorrespondence.\allowbreak{}RealBands.\allowbreak{}BranchPack}, lines 205--221)\\
\texttt{UpperTheory.\allowbreak{}Lit.\allowbreak{}Cusp.\allowbreak{}exists\_\allowbreak{}cusp\_\allowbreak{}form\_\allowbreak{}theta\_\allowbreak{}period\_\allowbreak{}ne\_\allowbreak{}zero} (\texttt{Development.\allowbreak{}AutomorphicArithmetic.\allowbreak{}ThetaCorrespondence.\allowbreak{}RealBands.\allowbreak{}Period}, lines 18--108)\par}

{\footnotesize\raggedright\noindent \ref{B8.15b}.\quad \texttt{UpperTheory.\allowbreak{}Lit.\allowbreak{}Branch.\allowbreak{}periodNe\_\allowbreak{}of\_\allowbreak{}eigen} (\texttt{Development.\allowbreak{}AutomorphicArithmetic.\allowbreak{}ThetaCorrespondence.\allowbreak{}RealBands.\allowbreak{}BranchPeriod}, lines 14--24)\par}

{\footnotesize\raggedright\noindent \ref{B8.15j}.\quad \texttt{GlobalTheory.\allowbreak{}ChainOnto.\allowbreak{}heckeVec\_\allowbreak{}howe} (\texttt{Development.\allowbreak{}AbelianVarietyTheory.\allowbreak{}Proofs.\allowbreak{}DecoratedLatticeProofs.\allowbreak{}UnramifiedHeckeVacuumLargerEqSmaller}, lines 23--53)\\
\texttt{GlobalTheory.\allowbreak{}SphTheta.\allowbreak{}satU\_\allowbreak{}eq\_\allowbreak{}satU\_\allowbreak{}liftPar\_\allowbreak{}of\_\allowbreak{}heckeVec\_\allowbreak{}eq} (\texttt{Development.\allowbreak{}AbelianVarietyTheory.\allowbreak{}Proofs.\allowbreak{}TwistedOperators.\allowbreak{}SatakeTransformEqAtThetaLiftOfVacuumEq}, lines 81--177)\par}

{\footnotesize\raggedright\noindent \ref{lem:p8c2-spherical-transfer}.\quad \texttt{GlobalTheory.\allowbreak{}EvenOnto.\allowbreak{}parBound\_\allowbreak{}of\_\allowbreak{}rapidDecay} (\texttt{Development.\allowbreak{}AbelianVarietyTheory.\allowbreak{}Proofs.\allowbreak{}TwistedOperators.\allowbreak{}EvenRankChainParamsTrivialBoundRapidDecay}, lines 28--63)\\
\texttt{GlobalTheory.\allowbreak{}FinZeta.\allowbreak{}ParamBound.\allowbreak{}exists\_\allowbreak{}param\_\allowbreak{}at} (\texttt{Development.\allowbreak{}AbelianVarietyTheory.\allowbreak{}Proofs.\allowbreak{}TwistedOperators.\allowbreak{}ExistsSatakeParametersBoundedEigenfunction}, lines 79--132)\\
\texttt{GlobalTheory.\allowbreak{}ChainOnto.\allowbreak{}transferOn\_\allowbreak{}all} (\texttt{Development.\allowbreak{}AbelianVarietyTheory.\allowbreak{}Proofs.\allowbreak{}DecoratedLatticeProofs.\allowbreak{}SphericalParameterTransferAllFinitePlaces}, lines 22--26)\\
\texttt{GlobalTheory.\allowbreak{}EvenOnto.\allowbreak{}upPar\_\allowbreak{}of\_\allowbreak{}transfer} (\texttt{Development.\allowbreak{}AbelianVarietyTheory.\allowbreak{}Proofs.\allowbreak{}TwistedOperators.\allowbreak{}EvenThetaChainUpwardLiftHeckeParamTransferred}, lines 48--73)\par}

{\footnotesize\raggedright\noindent \ref{lem:p8c2-spherical-transfer}.\quad \texttt{GlobalTheory.\allowbreak{}LGrp.\allowbreak{}bcParam} (\texttt{Development.\allowbreak{}AbelianVarietyTheory.\allowbreak{}Definitions.\allowbreak{}Divisors.\allowbreak{}LineOppositeParabolicUnitaryLeviSectionEq}, lines 870--871)\par}

{\footnotesize\raggedright\noindent \ref{lem:p8c2-spherical-transfer}.\quad \texttt{GlobalTheory.\allowbreak{}ChainOnto.\allowbreak{}stepPar} (\texttt{Development.\allowbreak{}AbelianVarietyTheory.\allowbreak{}Definitions.\allowbreak{}LatticeFrames.\allowbreak{}ChainParSucc}, lines 122--124)\\
\texttt{GlobalTheory.\allowbreak{}ChainOnto.\allowbreak{}SplitTransfer.\allowbreak{}transferOn\_\allowbreak{}isSplit} (\texttt{Development.\allowbreak{}AbelianVarietyTheory.\allowbreak{}Proofs.\allowbreak{}DecoratedLatticeProofs.\allowbreak{}SphericalParameterTransferSplitPlace}, lines 115--150)\\
\texttt{GlobalTheory.\allowbreak{}ChainOnto.\allowbreak{}transferOn\_\allowbreak{}inert} (\texttt{Development.\allowbreak{}AbelianVarietyTheory.\allowbreak{}Proofs.\allowbreak{}DecoratedLatticeProofs.\allowbreak{}SphericalTransferWeilPairingInertPlaces}, lines 89--115)\par}

{\footnotesize\raggedright\noindent \ref{B8.15c}.\quad \texttt{GlobalTheory.\allowbreak{}CoreOnto.\allowbreak{}coreOnto\_\allowbreak{}of\_\allowbreak{}chain} (\texttt{Development.\allowbreak{}AbelianVarietyTheory.\allowbreak{}Proofs.\allowbreak{}TwistedOperators.\allowbreak{}CoreChainLiftsSurjectiveOfEvenRank}, lines 28--55)\\
\texttt{GlobalTheory.\allowbreak{}FunPr.\allowbreak{}ThetaLocalInputs} (\texttt{Development.\allowbreak{}AutomorphicArithmetic.\allowbreak{}ThetaCorrespondence.\allowbreak{}ThetaPairLifts.\allowbreak{}LocalTypeDictionary.\allowbreak{}BernsteinUnitaryProductIrreducible}, lines 6--59)\\
\texttt{GlobalTheory.\allowbreak{}Published.\allowbreak{}WeilSiegelFormulaUnitary} (\texttt{Development.\allowbreak{}AbelianVarietyTheory.\allowbreak{}PublishedTheorems.\allowbreak{}TwistedTraceFormula.\allowbreak{}WeilSiegelFormulaUnitary}, lines 56--69)\\
\texttt{GlobalTheory.\allowbreak{}Published.\allowbreak{}DoublingMethodLocalZetaIntegrals} (\texttt{Development.\allowbreak{}AutomorphicArithmetic.\allowbreak{}ThetaCorrespondence.\allowbreak{}ThetaPairLifts.\allowbreak{}LocalTypeDictionary.\allowbreak{}DiscreteSpectrumResidualClassification}, lines 6--6)\\
\texttt{GlobalTheory.\allowbreak{}Published.\allowbreak{}UnramifiedDoublingGL} (\texttt{Development.\allowbreak{}AbelianVarietyTheory.\allowbreak{}PublishedTheorems.\allowbreak{}TwistedTraceFormula.\allowbreak{}PairwiseDisjointCell}, lines 314--324)\\
\texttt{GlobalTheory.\allowbreak{}Published.\allowbreak{}HarishChandraFinite} (\texttt{Development.\allowbreak{}AbelianVarietyTheory.\allowbreak{}PublishedTheorems.\allowbreak{}AutomorphicForms₂.\allowbreak{}HarishChandraFinite}, lines 50--58)\par}

{\footnotesize\raggedright\noindent \ref{B8.16c}.\quad \texttt{GlobalTheory.\allowbreak{}SquareComm.\allowbreak{}IsTriple} (\texttt{Development.\allowbreak{}AbelianVarietyTheory.\allowbreak{}Proofs.\allowbreak{}TwistedOperators.\allowbreak{}BlkGLMem}, lines 106--129)\\
\texttt{GlobalTheory.\allowbreak{}SquareComm.\allowbreak{}IsMatch} (\texttt{Development.\allowbreak{}AbelianVarietyTheory.\allowbreak{}Proofs.\allowbreak{}TwistedOperators.\allowbreak{}BlkGLMem}, lines 130--144)\\
\texttt{Formalism.\allowbreak{}IndefR.\allowbreak{}Datum.\allowbreak{}IsWeilDatumO} (\texttt{Formalism.\allowbreak{}LocalToGlobal.\allowbreak{}IndefiniteReal.\allowbreak{}LocalNewformFunction}, lines 53--76)\\
\texttt{GlobalTheory.\allowbreak{}SquareAgree.\allowbreak{}kapX} (\texttt{Development.\allowbreak{}AbelianVarietyTheory.\allowbreak{}Proofs.\allowbreak{}TwistedOperators.\allowbreak{}ThetaEqB}, lines 40--45)\par}

{\footnotesize\raggedright\noindent \ref{B8.16d}.\quad \texttt{GlobalTheory.\allowbreak{}T.\allowbreak{}exists\_\allowbreak{}placewise\_\allowbreak{}seesaw\_\allowbreak{}triple\_\allowbreak{}over\_\allowbreak{}hermitian} (\texttt{Development.\allowbreak{}AbelianVarietyTheory.\allowbreak{}Proofs.\allowbreak{}TwistedOperators.\allowbreak{}ExistsPlacewiseSeesawTripleOverHermitianStatement}, lines 16--28)\\
\texttt{GlobalTheory.\allowbreak{}SquareAgree.\allowbreak{}kapX} (\texttt{Development.\allowbreak{}AbelianVarietyTheory.\allowbreak{}Proofs.\allowbreak{}TwistedOperators.\allowbreak{}ThetaEqB}, lines 40--45)\\
\texttt{GlobalTheory.\allowbreak{}SquareAgree.\allowbreak{}thetaInt\_\allowbreak{}eq\_\allowbreak{}modelI} (\texttt{Development.\allowbreak{}AbelianVarietyTheory.\allowbreak{}Proofs.\allowbreak{}TwistedOperators.\allowbreak{}ConjUMeasurableEmbedding}, lines 59--72)\par}

{\footnotesize\raggedright\noindent \ref{B8.16e}.\quad \texttt{GlobalTheory.\allowbreak{}T.\allowbreak{}twisted\_\allowbreak{}theta\_\allowbreak{}kernels\_\allowbreak{}agree\_\allowbreak{}up\_\allowbreak{}to\_\allowbreak{}scalar} (\texttt{Development.\allowbreak{}AutomorphicArithmetic.\allowbreak{}ThetaCorrespondence.\allowbreak{}ThetaPairLifts.\allowbreak{}LocalTypeDictionary.\allowbreak{}ThetaKernelMatchOnOppositeParabolic}, lines 189--203)\par}

{\footnotesize\raggedright\noindent \ref{B8.16}.\quad \texttt{GlobalTheory.\allowbreak{}WeilConv.\allowbreak{}twisted\_\allowbreak{}weil\_\allowbreak{}representation\_\allowbreak{}convergence} (\texttt{Development.\allowbreak{}AutomorphicArithmetic.\allowbreak{}ThetaCorrespondence.\allowbreak{}ThetaPairLifts.\allowbreak{}LocalTypeDictionary.\allowbreak{}TwistedWeilRepresentationConvergence}, lines 16--17)\\
\texttt{GlobalTheory.\allowbreak{}T.\allowbreak{}twisted\_\allowbreak{}local\_\allowbreak{}factors\_\allowbreak{}match} (\texttt{Development.\allowbreak{}AutomorphicArithmetic.\allowbreak{}ThetaCorrespondence.\allowbreak{}ThetaPairLifts.\allowbreak{}LocalTypeDictionary.\allowbreak{}TwistedSquareLocalMatching}, lines 88--89)\\
\texttt{GlobalTheory.\allowbreak{}T.\allowbreak{}twisted\_\allowbreak{}square\_\allowbreak{}theta\_\allowbreak{}integrals\_\allowbreak{}eq} (\texttt{Development.\allowbreak{}AutomorphicArithmetic.\allowbreak{}ThetaCorrespondence.\allowbreak{}ThetaPairLifts.\allowbreak{}LocalTypeDictionary.\allowbreak{}TwistedSquareThetaIntegralsEq}, lines 48--49)\par}

{\footnotesize\raggedright\noindent \ref{B8.16}.\quad \texttt{GlobalTheory.\allowbreak{}WeilConv.\allowbreak{}weilConvergence\_\allowbreak{}of\_\allowbreak{}siegelSet} (\texttt{Development.\allowbreak{}AbelianVarietyTheory.\allowbreak{}Proofs.\allowbreak{}TwistedOperators.\allowbreak{}WeilConvergenceOfSiegelSet}, lines 24--82)\\
\texttt{GlobalTheory.\allowbreak{}SiegelSet.\allowbreak{}lintegral\_\allowbreak{}cone\_\allowbreak{}modulus\_\allowbreak{}lt\_\allowbreak{}top} (\texttt{Development.\allowbreak{}AbelianVarietyTheory.\allowbreak{}Proofs.\allowbreak{}TwistedOperators.\allowbreak{}LintegralSiegelConeInvModulusLtTopWeil}, lines 160--207)\\
\texttt{GlobalTheory.\allowbreak{}SquareLM.\allowbreak{}twisted\_\allowbreak{}square\_\allowbreak{}local\_\allowbreak{}matching} (\texttt{Development.\allowbreak{}AutomorphicArithmetic.\allowbreak{}ThetaCorrespondence.\allowbreak{}ThetaPairLifts.\allowbreak{}LocalTypeDictionary.\allowbreak{}TwistedSquareLocalMatching}, lines 18--81)\\
\texttt{GlobalTheory.\allowbreak{}SquareAgree.\allowbreak{}twisted\_\allowbreak{}theta\_\allowbreak{}integrals\_\allowbreak{}agree} (\texttt{Development.\allowbreak{}AutomorphicArithmetic.\allowbreak{}ThetaCorrespondence.\allowbreak{}ThetaPairLifts.\allowbreak{}LocalTypeDictionary.\allowbreak{}TwistedSquareThetaIntegralsEq}, lines 17--41)\\
\texttt{GlobalTheory.\allowbreak{}SquareAgree.\allowbreak{}CoinvPass.\allowbreak{}coinv\_\allowbreak{}passage} (\texttt{Development.\allowbreak{}AbelianVarietyTheory.\allowbreak{}Proofs.\allowbreak{}TwistedOperators.\allowbreak{}ConjUInvMemGfC}, lines 233--279)\par}

{\footnotesize\raggedright\noindent \ref{B8.16}.\quad \texttt{GlobalTheory.\allowbreak{}RowTest.\allowbreak{}NumbersIn} (\texttt{Development.\allowbreak{}AbelianVarietyTheory.\allowbreak{}Definitions.\allowbreak{}TwistedInnerProducts.\allowbreak{}MemHalfCosetSelf}, lines 245--261)\par}

{\footnotesize\raggedright\noindent \ref{B8.17a}.\quad \texttt{GlobalTheory.\allowbreak{}RowEx.\allowbreak{}BotField.\allowbreak{}sums\_\allowbreak{}mem\_\allowbreak{}bottomField} (\texttt{Development.\allowbreak{}AutomorphicArithmetic.\allowbreak{}ThetaCorrespondence.\allowbreak{}ThetaPairLifts.\allowbreak{}LocalTypeDictionary.\allowbreak{}BottomFieldIsNumberField}, lines 18--28)\par}

{\footnotesize\raggedright\noindent \ref{B8.17}.\quad \texttt{GlobalTheory.\allowbreak{}T.\allowbreak{}theta\_\allowbreak{}lift\_\allowbreak{}extension\_\allowbreak{}field\_\allowbreak{}exists} (\texttt{Development.\allowbreak{}AutomorphicArithmetic.\allowbreak{}ThetaCorrespondence.\allowbreak{}ThetaPairLifts.\allowbreak{}LocalTypeDictionary.\allowbreak{}ThetaLiftExtensionFieldExists}, lines 50--51)\\
\texttt{GlobalTheory.\allowbreak{}RowEx.\allowbreak{}FieldFor} (\texttt{Development.\allowbreak{}AbelianVarietyTheory.\allowbreak{}Definitions.\allowbreak{}TwistedInnerProducts.\allowbreak{}TwistPlacesOffPrimeQZero}, lines 231--236)\\
\texttt{GlobalTheory.\allowbreak{}RowEx.\allowbreak{}ExField} (\texttt{Development.\allowbreak{}AbelianVarietyTheory.\allowbreak{}Definitions.\allowbreak{}TwistedInnerProducts.\allowbreak{}TwistPlacesOffPrimeQZero}, lines 241--242)\par}

{\footnotesize\raggedright\noindent \ref{B8.17}.\quad \texttt{GlobalTheory.\allowbreak{}RowEx.\allowbreak{}exKey\_\allowbreak{}of} (\texttt{Development.\allowbreak{}AbelianVarietyTheory.\allowbreak{}Proofs.\allowbreak{}TwistedOperators.\allowbreak{}ExKeyOf}, lines 23--89)\\
\texttt{GlobalTheory.\allowbreak{}RowEx.\allowbreak{}exField\_\allowbreak{}of\_\allowbreak{}keyFin\_\allowbreak{}exKey} (\texttt{Development.\allowbreak{}AbelianVarietyTheory.\allowbreak{}Proofs.\allowbreak{}TwistedOperators.\allowbreak{}ExFieldOfKeyFinExKey}, lines 22--36)\par}

{\footnotesize\raggedright\noindent \ref{def:p8c2-test-stage}.\quad \texttt{GlobalTheory.\allowbreak{}Mem.\allowbreak{}Top.\allowbreak{}Bl.\allowbreak{}KPx} (\texttt{Development.\allowbreak{}AbelianVarietyTheory.\allowbreak{}Definitions.\allowbreak{}AutomorphicMembership.\allowbreak{}PositionOfIntro}, lines 17--22)\\
\texttt{GlobalTheory.\allowbreak{}Mem.\allowbreak{}Fld.\allowbreak{}listField} (\texttt{Development.\allowbreak{}AbelianVarietyTheory.\allowbreak{}Definitions.\allowbreak{}AutomorphicMembership.\allowbreak{}FixedCoefficientFieldInComplex}, lines 501--503)\\
\texttt{GlobalTheory.\allowbreak{}Mem.\allowbreak{}Top.\allowbreak{}TypesTrue} (\texttt{Development.\allowbreak{}AbelianVarietyTheory.\allowbreak{}Proofs.\allowbreak{}IwasawaDecomposition.\allowbreak{}TypeIdempotentCoeffAcrossLayers}, lines 32--37)\\
\texttt{GlobalTheory.\allowbreak{}Mem.\allowbreak{}R3.\allowbreak{}GrpAt} (\texttt{Development.\allowbreak{}AbelianVarietyTheory.\allowbreak{}Definitions.\allowbreak{}AutomorphicMembership.\allowbreak{}IsAbsIrrCharCompIff}, lines 248--256)\\
\texttt{Formalism.\allowbreak{}IndefR.\allowbreak{}Links.\allowbreak{}RealDat.\allowbreak{}ExAt} (\texttt{Formalism.\allowbreak{}LocalToGlobal.\allowbreak{}IndefiniteReal.\allowbreak{}LinksW}, lines 105--139)\\
\texttt{GlobalTheory.\allowbreak{}RealGeo.\allowbreak{}InGeo} (\texttt{Development.\allowbreak{}AbelianVarietyTheory.\allowbreak{}Definitions.\allowbreak{}AutomorphicMembership.\allowbreak{}JEvExists}, lines 119--128)\par}

{\footnotesize\raggedright\noindent \ref{B8.18}.\quad \texttt{GlobalTheory.\allowbreak{}T.\allowbreak{}test\_\allowbreak{}vector\_\allowbreak{}assembly\_\allowbreak{}holds} (\texttt{Development.\allowbreak{}AutomorphicArithmetic.\allowbreak{}ThetaCorrespondence.\allowbreak{}ThetaPairLifts.\allowbreak{}LocalTypeDictionary.\allowbreak{}TestVectorAtDatumOfJoin}, lines 39--40)\\
\texttt{GlobalTheory.\allowbreak{}RowTestAssembly.\allowbreak{}test\_\allowbreak{}vector\_\allowbreak{}at\_\allowbreak{}datum\_\allowbreak{}of\_\allowbreak{}join} (\texttt{Development.\allowbreak{}AutomorphicArithmetic.\allowbreak{}ThetaCorrespondence.\allowbreak{}ThetaPairLifts.\allowbreak{}LocalTypeDictionary.\allowbreak{}TestVectorAtDatumOfJoin}, lines 20--31)\\
\texttt{Formalism.\allowbreak{}IndefR.\allowbreak{}Rat.\allowbreak{}RowTest} (\texttt{Formalism.\allowbreak{}LocalToGlobal.\allowbreak{}IndefiniteReal.\allowbreak{}RationalAdjoints.\allowbreak{}RatZOfTests}, lines 18--20)\par}

{\footnotesize\raggedright\noindent \ref{B8.18}.\quad \texttt{GlobalTheory.\allowbreak{}RowTestFromClaims.\allowbreak{}rowTest\_\allowbreak{}of\_\allowbreak{}choice} (\texttt{Development.\allowbreak{}AbelianVarietyTheory.\allowbreak{}Proofs.\allowbreak{}TwistedOperators.\allowbreak{}RowTestHilbertModelsOfChoiceClaims}, lines 104--133)\\
\texttt{GlobalTheory.\allowbreak{}Mem.\allowbreak{}R3.\allowbreak{}Rw.\allowbreak{}RowAt.\allowbreak{}atDat\_\allowbreak{}rowTest} (\texttt{Development.\allowbreak{}AbelianVarietyTheory.\allowbreak{}Proofs.\allowbreak{}IwasawaDecomposition.\allowbreak{}AtDatRowTest}, lines 44--64)\\
\texttt{GlobalTheory.\allowbreak{}Mem.\allowbreak{}LiteralAutomorphy.\allowbreak{}rowRatW\_\allowbreak{}of\_\allowbreak{}rowTest} (\texttt{Development.\allowbreak{}AbelianVarietyTheory.\allowbreak{}Proofs.\allowbreak{}IwasawaDecomposition.\allowbreak{}RowRatWOfRowTest}, lines 9--12)\par}

{\footnotesize\raggedright\noindent \ref{B8.20}.\quad \texttt{GlobalTheory.\allowbreak{}Mem.\allowbreak{}LiteralAutomorphy.\allowbreak{}Full} (\texttt{Development.\allowbreak{}AbelianVarietyTheory.\allowbreak{}Definitions.\allowbreak{}AutomorphicMembership.\allowbreak{}GKeyFin}, lines 43--47)\\
\texttt{GlobalTheory.\allowbreak{}Mem.\allowbreak{}LiteralAutomorphy.\allowbreak{}GridAt} (\texttt{Development.\allowbreak{}AbelianVarietyTheory.\allowbreak{}Definitions.\allowbreak{}AutomorphicMembership.\allowbreak{}GKeyFin}, lines 49--54)\\
\texttt{GlobalTheory.\allowbreak{}Mem.\allowbreak{}LiteralAutomorphy.\allowbreak{}ExGrid} (\texttt{Development.\allowbreak{}AbelianVarietyTheory.\allowbreak{}Definitions.\allowbreak{}AutomorphicMembership.\allowbreak{}GKeyFin}, lines 56--59)\\
\texttt{GlobalTheory.\allowbreak{}Mem.\allowbreak{}LiteralAutomorphy.\allowbreak{}GridA} (\texttt{Development.\allowbreak{}AbelianVarietyTheory.\allowbreak{}Definitions.\allowbreak{}AutomorphicMembership.\allowbreak{}GKeyFin}, lines 61--62)\\
\texttt{GlobalTheory.\allowbreak{}Mem.\allowbreak{}LiteralAutomorphy.\allowbreak{}exGrid\_\allowbreak{}of} (\texttt{Development.\allowbreak{}AbelianVarietyTheory.\allowbreak{}Proofs.\allowbreak{}IwasawaDecomposition.\allowbreak{}GridExistenceOfPieces}, lines 34--143)\par}

{\footnotesize\raggedright\noindent \ref{lem:p8c2-square-field}.\quad \texttt{GlobalTheory.\allowbreak{}SqTop.\allowbreak{}SqField.\allowbreak{}sqIn\_\allowbreak{}of} (\texttt{Development.\allowbreak{}AbelianVarietyTheory.\allowbreak{}Proofs.\allowbreak{}DivisorProofs.\allowbreak{}IsAlgebraicDivOfBeta}, lines 73--127)\\
\texttt{GlobalTheory.\allowbreak{}ToricSumValue.\allowbreak{}toricSum\_\allowbreak{}eq} (\texttt{Development.\allowbreak{}AbelianVarietyTheory.\allowbreak{}Proofs.\allowbreak{}DivisorProofs.\allowbreak{}ToricSumEqHeckeLValueMulAuxiliaryFactors}, lines 95--130)\par}

{\footnotesize\raggedright\noindent \ref{lem:p8c2-square-field}.\quad \texttt{GlobalTheory.\allowbreak{}FunObj.\allowbreak{}SqNum.\allowbreak{}Model.\allowbreak{}kap\_\allowbreak{}add} (\texttt{Development.\allowbreak{}AbelianVarietyTheory.\allowbreak{}Definitions.\allowbreak{}Divisors.\allowbreak{}IsValCharTwist}, lines 68--82)\par}

{\footnotesize\raggedright\noindent \ref{lem:p8c2-square-field}.\quad \texttt{GlobalTheory.\allowbreak{}SqTop.\allowbreak{}SqField.\allowbreak{}isAlgebraic\_\allowbreak{}norm} (\texttt{Development.\allowbreak{}AbelianVarietyTheory.\allowbreak{}Proofs.\allowbreak{}DivisorProofs.\allowbreak{}IsAlgebraicDivOfBeta}, lines 32--58)\par}

{\footnotesize\raggedright\noindent \ref{B8.21}.\quad \texttt{GlobalTheory.\allowbreak{}Mem.\allowbreak{}LiteralAutomorphy.\allowbreak{}evaluation\_\allowbreak{}grid\_\allowbreak{}of\_\allowbreak{}join} (\texttt{Development.\allowbreak{}AutomorphicArithmetic.\allowbreak{}ThetaCorrespondence.\allowbreak{}ThetaPairLifts.\allowbreak{}LocalTypeDictionary.\allowbreak{}EvaluationGridOfJoin}, lines 23--44)\\
\texttt{GlobalTheory.\allowbreak{}Mem.\allowbreak{}LiteralAutomorphy.\allowbreak{}evLgG\_\allowbreak{}gridA\_\allowbreak{}of\_\allowbreak{}exGrid} (\texttt{Development.\allowbreak{}AbelianVarietyTheory.\allowbreak{}Proofs.\allowbreak{}IwasawaDecomposition.\allowbreak{}EventuallyLargeGrid}, lines 14--41)\par}

{\footnotesize\raggedright\noindent \ref{B8.21a}.\quad \texttt{GlobalTheory.\allowbreak{}FunObj.\allowbreak{}SqNum.\allowbreak{}IsTorVal} (\texttt{Development.\allowbreak{}AbelianVarietyTheory.\allowbreak{}Definitions.\allowbreak{}Divisors.\allowbreak{}FinitePrimesUnder}, lines 274--280)\\
\texttt{GlobalTheory.\allowbreak{}SqTop.\allowbreak{}thbSt\_\allowbreak{}of\_\allowbreak{}inputs} (\texttt{Development.\allowbreak{}AbelianVarietyTheory.\allowbreak{}Proofs.\allowbreak{}DivisorProofs.\allowbreak{}ToricValueHeightBoundOfDenominatorsMeans}, lines 158--242)\\
\texttt{GlobalTheory.\allowbreak{}BigTop.\allowbreak{}Fun.\allowbreak{}thb\_\allowbreak{}tie\_\allowbreak{}alt} (\texttt{Development.\allowbreak{}AutomorphicArithmetic.\allowbreak{}ThetaCorrespondence.\allowbreak{}ThetaPairLifts.\allowbreak{}ThirdPlaceSeries.\allowbreak{}ThbTie}, lines 16--60)\par}

{\footnotesize\raggedright\noindent \ref{B8.21a}.\quad \texttt{GlobalTheory.\allowbreak{}DenSumPf.\allowbreak{}denSum} (\texttt{Development.\allowbreak{}AbelianVarietyTheory.\allowbreak{}Proofs.\allowbreak{}DivisorProofs.\allowbreak{}LinearCombinationLogBoundRealInequality}, lines 266--392)\\
\texttt{GlobalTheory.\allowbreak{}SqTop.\allowbreak{}lBound\_\allowbreak{}of} (\texttt{Development.\allowbreak{}AbelianVarietyTheory.\allowbreak{}Proofs.\allowbreak{}DivisorProofs.\allowbreak{}HeckeLValueUpperBoundPermutedWeightsOfFE}, lines 16--92)\\
\texttt{GlobalTheory.\allowbreak{}SqTop.\allowbreak{}betaArch\_\allowbreak{}of\_\allowbreak{}transport} (\texttt{Development.\allowbreak{}AbelianVarietyTheory.\allowbreak{}Proofs.\allowbreak{}DivisorProofs.\allowbreak{}BetaConjugateMeanLogAbsBoundOfTransport}, lines 43--145)\\
\texttt{GlobalTheory.\allowbreak{}HeightMeanValueBounds.\allowbreak{}height\_\allowbreak{}meanVal} (\texttt{Development.\allowbreak{}AbelianVarietyTheory.\allowbreak{}Proofs.\allowbreak{}DivisorProofs.\allowbreak{}ToricValueHeightBoundOfDenominatorsMeans}, lines 36--120)\par}

{\footnotesize\raggedright\noindent \ref{B8.21b}.\quad \texttt{Formalism.\allowbreak{}BasicNotions.\allowbreak{}Ranks} (\texttt{Formalism.\allowbreak{}LocalRingCardinality.\allowbreak{}Params}, lines 6--22)\\
\texttt{Formalism.\allowbreak{}BasicNotions.\allowbreak{}HCParam} (\texttt{Formalism.\allowbreak{}LocalRingCardinality.\allowbreak{}Params}, lines 50--69)\\
\texttt{Formalism.\allowbreak{}BasicNotions.\allowbreak{}sExp} (\texttt{Formalism.\allowbreak{}LocalRingCardinality.\allowbreak{}Params}, lines 71--71)\\
\texttt{Formalism.\allowbreak{}BasicNotions.\allowbreak{}sExp'} (\texttt{Formalism.\allowbreak{}LocalRingCardinality.\allowbreak{}Params}, lines 73--73)\\
\texttt{Formalism.\allowbreak{}BasicNotions.\allowbreak{}wG} (\texttt{Formalism.\allowbreak{}LocalRingCardinality.\allowbreak{}Params}, lines 75--75)\\
\texttt{Formalism.\allowbreak{}BasicNotions.\allowbreak{}Host} (\texttt{Formalism.\allowbreak{}LocalRingCardinality.\allowbreak{}DataHost}, lines 39--109)\\
\texttt{Formalism.\allowbreak{}BasicNotions.\allowbreak{}HSpace} (\texttt{Formalism.\allowbreak{}LocalRingCardinality.\allowbreak{}DataSpaces}, lines 26--46)\\
\texttt{Formalism.\allowbreak{}BasicNotions.\allowbreak{}Spaces} (\texttt{Formalism.\allowbreak{}LocalRingCardinality.\allowbreak{}DataSpaces}, lines 77--139)\\
\texttt{Formalism.\allowbreak{}BasicNotions.\allowbreak{}sigYD} (\texttt{Formalism.\allowbreak{}LocalRingCardinality.\allowbreak{}DataSpaces}, lines 75--75)\\
\texttt{Formalism.\allowbreak{}BasicNotions.\allowbreak{}Rows.\allowbreak{}Data} (\texttt{Formalism.\allowbreak{}LocalRingCardinality.\allowbreak{}DataRows}, lines 36--168)\\
\texttt{Formalism.\allowbreak{}BasicNotions.\allowbreak{}HLat} (\texttt{Formalism.\allowbreak{}LocalRingCardinality.\allowbreak{}DataArith}, lines 61--72)\\
\texttt{Formalism.\allowbreak{}BasicNotions.\allowbreak{}DivData} (\texttt{Formalism.\allowbreak{}LocalRingCardinality.\allowbreak{}DataArith}, lines 139--159)\\
\texttt{Formalism.\allowbreak{}BasicNotions.\allowbreak{}D7} (\texttt{Formalism.\allowbreak{}LocalRingCardinality.\allowbreak{}RhoDetPrimeEqMuTildeMulChiPow}, lines 76--335)\\
\texttt{Formalism.\allowbreak{}BasicNotions.\allowbreak{}Member} (\texttt{Formalism.\allowbreak{}LocalRingCardinality.\allowbreak{}RhoDetPrimeEqMuTildeMulChiPow}, lines 71--74)\\
\texttt{Formalism.\allowbreak{}BasicNotions.\allowbreak{}ColumnChain} (\texttt{Formalism.\allowbreak{}LocalRingCardinality.\allowbreak{}DataKernels}, lines 14--47)\\
\texttt{Formalism.\allowbreak{}BasicNotions.\allowbreak{}ColumnsAtT} (\texttt{Formalism.\allowbreak{}LocalRingCardinality.\allowbreak{}DataKernels}, lines 104--138)\\
\texttt{Formalism.\allowbreak{}BasicNotions.\allowbreak{}PeriodMap} (\texttt{Formalism.\allowbreak{}LocalRingCardinality.\allowbreak{}DataKernels}, lines 163--183)\\
\texttt{Formalism.\allowbreak{}BasicNotions.\allowbreak{}MemberForms} (\texttt{Formalism.\allowbreak{}LocalRingCardinality.\allowbreak{}DataForms}, lines 107--215)\\
\texttt{Formalism.\allowbreak{}BasicNotions.\allowbreak{}MemberForms.\allowbreak{}ipC} (\texttt{Formalism.\allowbreak{}LocalRingCardinality.\allowbreak{}DataForms}, lines 227--227)\\
\texttt{Formalism.\allowbreak{}BasicNotions.\allowbreak{}MemberForms.\allowbreak{}ramSub} (\texttt{Formalism.\allowbreak{}LocalRingCardinality.\allowbreak{}DataForms}, lines 239--239)\\
\texttt{Formalism.\allowbreak{}MainNotions.\allowbreak{}S.\allowbreak{}Draft.\allowbreak{}SData} (\texttt{Formalism.\allowbreak{}SharpRestriction.\allowbreak{}CuspSpanNormBounds.\allowbreak{}DataS}, lines 13--84)\\
\texttt{Formalism.\allowbreak{}BasicNotions.\allowbreak{}SizeData₂} (\texttt{Formalism.\allowbreak{}LocalRingCardinality.\allowbreak{}DataParams₂}, lines 76--182)\\
\texttt{GlobalTheory.\allowbreak{}Mem.\allowbreak{}Stages} (\texttt{Development.\allowbreak{}AbelianVarietyTheory.\allowbreak{}Definitions.\allowbreak{}AutomorphicMembership.\allowbreak{}AssembleLattice}, lines 188--196)\\
\texttt{GlobalTheory.\allowbreak{}Mem.\allowbreak{}Top.\allowbreak{}Ctx} (\texttt{Development.\allowbreak{}AbelianVarietyTheory.\allowbreak{}Definitions.\allowbreak{}AutomorphicMembership.\allowbreak{}CmQuadraticFormEqAlternatingForm}, lines 28--43)\\
\texttt{GlobalTheory.\allowbreak{}Mem.\allowbreak{}Top.\allowbreak{}TwDat} (\texttt{Development.\allowbreak{}AbelianVarietyTheory.\allowbreak{}Definitions.\allowbreak{}AutomorphicMembership.\allowbreak{}CmQuadraticFormEqAlternatingForm}, lines 79--90)\\
\texttt{GlobalTheory.\allowbreak{}Mem.\allowbreak{}Top.\allowbreak{}V4.\allowbreak{}Core} (\texttt{Development.\allowbreak{}AbelianVarietyTheory.\allowbreak{}Definitions.\allowbreak{}AutomorphicMembership.\allowbreak{}RationalRootsDef}, lines 38--62)\\
\texttt{GlobalTheory.\allowbreak{}EX.\allowbreak{}OB.\allowbreak{}Lay} (\texttt{Development.\allowbreak{}AbelianVarietyTheory.\allowbreak{}Definitions.\allowbreak{}GlobalTwisting.\allowbreak{}SetUnipotentTranslate}, lines 83--120)\\
\texttt{GlobalTheory.\allowbreak{}EX.\allowbreak{}OB.\allowbreak{}LedPart} (\texttt{Development.\allowbreak{}AbelianVarietyTheory.\allowbreak{}Definitions.\allowbreak{}GlobalTwisting.\allowbreak{}SetUnipotentTranslate}, lines 148--160)\\
\texttt{GlobalTheory.\allowbreak{}Fc.\allowbreak{}BPt} (\texttt{Development.\allowbreak{}AbelianVarietyTheory.\allowbreak{}Definitions.\allowbreak{}GlobalTwisting.\allowbreak{}ValidForConstructedPointClassesEventuallyWorking}, lines 67--83)\\
\texttt{GlobalTheory.\allowbreak{}Fc.\allowbreak{}Built} (\texttt{Development.\allowbreak{}AbelianVarietyTheory.\allowbreak{}Definitions.\allowbreak{}GlobalTwisting.\allowbreak{}ValidForConstructedPointClassesEventuallyWorking}, lines 250--252)\\
\texttt{GlobalTheory.\allowbreak{}BigTop.\allowbreak{}IndefiniteRowSideFullHecke} (\texttt{Development.\allowbreak{}AbelianVarietyTheory.\allowbreak{}Definitions.\allowbreak{}Divisors.\allowbreak{}TupleReadBaseTupleEqFamily}, lines 61--109)\\
\texttt{Formalism.\allowbreak{}IndefR.\allowbreak{}Links.\allowbreak{}RealDat3} (\texttt{Formalism.\allowbreak{}LocalToGlobal.\allowbreak{}IndefiniteReal.\allowbreak{}ZeroPartEigenvalueConvention}, lines 89--97)\\
\texttt{Formalism.\allowbreak{}IndefR.\allowbreak{}FWT.\allowbreak{}FWTie} (\texttt{Formalism.\allowbreak{}LocalToGlobal.\allowbreak{}IndefiniteReal.\allowbreak{}FWTie}, lines 17--59)\\
\texttt{Formalism.\allowbreak{}IndefR.\allowbreak{}FWT.\allowbreak{}OMDat} (\texttt{Formalism.\allowbreak{}LocalToGlobal.\allowbreak{}IndefiniteReal.\allowbreak{}OneMapT}, lines 32--59)\\
\texttt{GlobalTheory.\allowbreak{}BigTop.\allowbreak{}PtAt} (\texttt{Development.\allowbreak{}AbelianVarietyTheory.\allowbreak{}Definitions.\allowbreak{}Divisors.\allowbreak{}TupleReadBaseTupleEqFamily}, lines 145--258)\\
\texttt{Formalism.\allowbreak{}L2.\allowbreak{}DIV.\allowbreak{}frameOf} (\texttt{Formalism.\allowbreak{}LocalToGlobal.\allowbreak{}Divisibility.\allowbreak{}Frames}, lines 22--23)\\
\texttt{Formalism.\allowbreak{}L2.\allowbreak{}DIV.\allowbreak{}frameOfSpec} (\texttt{Formalism.\allowbreak{}LocalToGlobal.\allowbreak{}Divisibility.\allowbreak{}FramesLemmas}, lines 77--90)\\
\texttt{Formalism.\allowbreak{}L2.\allowbreak{}DIV.\allowbreak{}quotFrame} (\texttt{Formalism.\allowbreak{}LocalToGlobal.\allowbreak{}Divisibility.\allowbreak{}Frames}, lines 41--43)\\
\texttt{Formalism.\allowbreak{}L2.\allowbreak{}DIV.\allowbreak{}BlockDat.\allowbreak{}XQ} (\texttt{Formalism.\allowbreak{}LocalToGlobal.\allowbreak{}Divisibility.\allowbreak{}RowLat}, lines 144--145)\par}

{\footnotesize\raggedright\noindent \ref{B8.21c}.\quad \texttt{Formalism.\allowbreak{}L2.\allowbreak{}Host.\allowbreak{}PiBar} (\texttt{Formalism.\allowbreak{}LocalToGlobal.\allowbreak{}UnitaryNorms.\allowbreak{}CharHost}, lines 44--46)\\
\texttt{Formalism.\allowbreak{}L2.\allowbreak{}Host.\allowbreak{}Omg} (\texttt{Formalism.\allowbreak{}LocalToGlobal.\allowbreak{}UnitaryNorms.\allowbreak{}CharHost}, lines 51--53)\\
\texttt{Formalism.\allowbreak{}L2.\allowbreak{}Host.\allowbreak{}IsMu0} (\texttt{Formalism.\allowbreak{}LocalToGlobal.\allowbreak{}UnitaryNorms.\allowbreak{}CharHost}, lines 57--64)\\
\texttt{Formalism.\allowbreak{}L2.\allowbreak{}Host.\allowbreak{}IsEtaT} (\texttt{Formalism.\allowbreak{}LocalToGlobal.\allowbreak{}UnitaryNorms.\allowbreak{}CharHost}, lines 26--30)\\
\texttt{Formalism.\allowbreak{}L2.\allowbreak{}Host.\allowbreak{}IsChiPsi} (\texttt{Formalism.\allowbreak{}LocalToGlobal.\allowbreak{}UnitaryNorms.\allowbreak{}CharHost}, lines 38--42)\\
\texttt{Formalism.\allowbreak{}L2.\allowbreak{}Host.\allowbreak{}IsDelta} (\texttt{Formalism.\allowbreak{}LocalToGlobal.\allowbreak{}UnitaryNorms.\allowbreak{}CharHost}, lines 68--73)\\
\texttt{Formalism.\allowbreak{}L2.\allowbreak{}Host.\allowbreak{}CharFrame} (\texttt{Formalism.\allowbreak{}LocalToGlobal.\allowbreak{}UnitaryNorms.\allowbreak{}CharTied}, lines 16--31)\\
\texttt{Formalism.\allowbreak{}IndefR.\allowbreak{}Links.\allowbreak{}Yc} (\texttt{Formalism.\allowbreak{}LocalToGlobal.\allowbreak{}IndefiniteReal.\allowbreak{}LinksDefs}, lines 358--358)\\
\texttt{Formalism.\allowbreak{}IndefR.\allowbreak{}Links.\allowbreak{}Vc} (\texttt{Formalism.\allowbreak{}LocalToGlobal.\allowbreak{}IndefiniteReal.\allowbreak{}LinksDefs}, lines 366--366)\\
\texttt{Formalism.\allowbreak{}IndefR.\allowbreak{}Links.\allowbreak{}Wc} (\texttt{Formalism.\allowbreak{}LocalToGlobal.\allowbreak{}IndefiniteReal.\allowbreak{}LinksDefs}, lines 362--362)\\
\texttt{Formalism.\allowbreak{}IndefR.\allowbreak{}Links.\allowbreak{}LevCh} (\texttt{Formalism.\allowbreak{}LocalToGlobal.\allowbreak{}IndefiniteReal.\allowbreak{}LinksDefs}, lines 388--419)\\
\texttt{GlobalTheory.\allowbreak{}ChiGrid.\allowbreak{}Consts} (\texttt{Development.\allowbreak{}AbelianVarietyTheory.\allowbreak{}Definitions.\allowbreak{}Divisors.\allowbreak{}BasePlaceNotMemChiGridBadPlaces}, lines 22--28)\\
\texttt{GlobalTheory.\allowbreak{}ChiGrid.\allowbreak{}S1} (\texttt{Development.\allowbreak{}AbelianVarietyTheory.\allowbreak{}Definitions.\allowbreak{}Divisors.\allowbreak{}LeastSplitAuxiliaryPrimeThetaChainSpec}, lines 262--264)\\
\texttt{Formalism.\allowbreak{}L2.\allowbreak{}Theta.\allowbreak{}IsPlaceIdem} (\texttt{Formalism.\allowbreak{}LocalToGlobal.\allowbreak{}BesselHermitianTheta.\allowbreak{}TiedWeil}, lines 68--69)\\
\texttt{Formalism.\allowbreak{}L2.\allowbreak{}Theta.\allowbreak{}IsSplitIdem} (\texttt{Formalism.\allowbreak{}LocalToGlobal.\allowbreak{}BesselHermitianTheta.\allowbreak{}TiedWeil}, lines 71--72)\\
\texttt{Formalism.\allowbreak{}L2.\allowbreak{}NTheta.\allowbreak{}IsIdemAt} (\texttt{Formalism.\allowbreak{}LocalToGlobal.\allowbreak{}HeckeThetaElements.\allowbreak{}IsHeckeElt}, lines 68--70)\\
\texttt{GlobalTheory.\allowbreak{}ChiGrid.\allowbreak{}Places} (\texttt{Development.\allowbreak{}AbelianVarietyTheory.\allowbreak{}Definitions.\allowbreak{}Divisors.\allowbreak{}BasePlaceNotMemChiGridBadPlaces}, lines 43--62)\\
\texttt{GlobalTheory.\allowbreak{}FunObj.\allowbreak{}SqNum.\allowbreak{}condExp} (\texttt{Development.\allowbreak{}AbelianVarietyTheory.\allowbreak{}Definitions.\allowbreak{}Divisors.\allowbreak{}FinitePrimesUnder}, lines 212--212)\\
\texttt{GlobalTheory.\allowbreak{}FunObj.\allowbreak{}SqNum.\allowbreak{}condF} (\texttt{Development.\allowbreak{}AbelianVarietyTheory.\allowbreak{}Definitions.\allowbreak{}Divisors.\allowbreak{}FinitePrimesUnder}, lines 214--216)\\
\texttt{GlobalTheory.\allowbreak{}FunObj.\allowbreak{}SqNum.\allowbreak{}condE} (\texttt{Development.\allowbreak{}AbelianVarietyTheory.\allowbreak{}Definitions.\allowbreak{}Divisors.\allowbreak{}FinitePrimesUnder}, lines 217--218)\\
\texttt{GlobalTheory.\allowbreak{}FunObj.\allowbreak{}SqNum.\allowbreak{}levQ} (\texttt{Development.\allowbreak{}AbelianVarietyTheory.\allowbreak{}Definitions.\allowbreak{}Divisors.\allowbreak{}FinitePrimesUnder}, lines 219--220)\\
\texttt{GlobalTheory.\allowbreak{}FunObj.\allowbreak{}SqNum.\allowbreak{}log\_\allowbreak{}levQ} (\texttt{Development.\allowbreak{}AbelianVarietyTheory.\allowbreak{}Definitions.\allowbreak{}Divisors.\allowbreak{}FinitePrimesUnder}, lines 223--228)\\
\texttt{Formalism.\allowbreak{}L2.\allowbreak{}Host.\allowbreak{}HChar.\allowbreak{}TrivOn} (\texttt{Formalism.\allowbreak{}LocalToGlobal.\allowbreak{}UnitaryNorms.\allowbreak{}IsUnitaryNormOf}, lines 194--195)\\
\texttt{Formalism.\allowbreak{}L2.\allowbreak{}Host.\allowbreak{}HChar.\allowbreak{}HasInf} (\texttt{Formalism.\allowbreak{}LocalToGlobal.\allowbreak{}UnitaryNorms.\allowbreak{}IsUnitaryNormOf}, lines 90--91)\\
\texttt{Formalism.\allowbreak{}L2.\allowbreak{}Host.\allowbreak{}HChar.\allowbreak{}UnrAt} (\texttt{Formalism.\allowbreak{}LocalToGlobal.\allowbreak{}UnitaryNorms.\allowbreak{}IsUnitaryNormOf}, lines 71--71)\\
\texttt{Formalism.\allowbreak{}L2.\allowbreak{}Host.\allowbreak{}HChar.\allowbreak{}LevelLe} (\texttt{Formalism.\allowbreak{}LocalToGlobal.\allowbreak{}UnitaryNorms.\allowbreak{}IsUnitaryNormOf}, lines 67--69)\\
\texttt{GlobalTheory.\allowbreak{}ChiGrid.\allowbreak{}cen} (\texttt{Development.\allowbreak{}AbelianVarietyTheory.\allowbreak{}Definitions.\allowbreak{}Divisors.\allowbreak{}LeastSplitAuxiliaryPrimeThetaChainSpec}, lines 150--150)\\
\texttt{GlobalTheory.\allowbreak{}ChiGrid.\allowbreak{}Chars} (\texttt{Development.\allowbreak{}AbelianVarietyTheory.\allowbreak{}Definitions.\allowbreak{}Divisors.\allowbreak{}BasePlaceNotMemChiGridBadPlaces}, lines 117--139)\\
\texttt{Formalism.\allowbreak{}L2.\allowbreak{}Theta.\allowbreak{}frameMat} (\texttt{Formalism.\allowbreak{}LocalToGlobal.\allowbreak{}BesselHermitianTheta.\allowbreak{}IsAutFunDef}, lines 79--80)\\
\texttt{GlobalTheory.\allowbreak{}ChiGrid.\allowbreak{}antiDiagonalSignGramMatrix} (\texttt{Development.\allowbreak{}AbelianVarietyTheory.\allowbreak{}Definitions.\allowbreak{}Divisors.\allowbreak{}CmHostImaginaryGeneratorNeZero}, lines 102--103)\\
\texttt{GlobalTheory.\allowbreak{}ChiGrid.\allowbreak{}HermitianOrthogonalExtensionData} (\texttt{Development.\allowbreak{}AbelianVarietyTheory.\allowbreak{}Definitions.\allowbreak{}Divisors.\allowbreak{}CmHostImaginaryGeneratorNeZero}, lines 105--118)\\
\texttt{GlobalTheory.\allowbreak{}ChiGrid.\allowbreak{}stairId} (\texttt{Development.\allowbreak{}AbelianVarietyTheory.\allowbreak{}Definitions.\allowbreak{}Divisors.\allowbreak{}CmHostImaginaryGeneratorNeZero}, lines 133--135)\\
\texttt{GlobalTheory.\allowbreak{}ChiGrid.\allowbreak{}Vertices} (\texttt{Development.\allowbreak{}AbelianVarietyTheory.\allowbreak{}Definitions.\allowbreak{}Divisors.\allowbreak{}BasePlaceNotMemChiGridBadPlaces}, lines 168--197)\\
\texttt{GlobalTheory.\allowbreak{}ChiGrid.\allowbreak{}Node} (\texttt{Development.\allowbreak{}AbelianVarietyTheory.\allowbreak{}Definitions.\allowbreak{}Divisors.\allowbreak{}BasePlaceNotMemChiGridBadPlaces}, lines 199--206)\\
\texttt{GlobalTheory.\allowbreak{}ChiGrid.\allowbreak{}Node.\allowbreak{}rank} (\texttt{Development.\allowbreak{}AbelianVarietyTheory.\allowbreak{}Definitions.\allowbreak{}Divisors.\allowbreak{}BasePlaceNotMemChiGridBadPlaces}, lines 215--218)\\
\texttt{GlobalTheory.\allowbreak{}ChiGrid.\allowbreak{}GridData} (\texttt{Development.\allowbreak{}AbelianVarietyTheory.\allowbreak{}Definitions.\allowbreak{}Divisors.\allowbreak{}BasePlaceNotMemChiGridBadPlaces}, lines 335--341)\\
\texttt{GlobalTheory.\allowbreak{}ChiGrid.\allowbreak{}sLen} (\texttt{Development.\allowbreak{}AbelianVarietyTheory.\allowbreak{}Definitions.\allowbreak{}Divisors.\allowbreak{}LeastSplitAuxiliaryPrimeThetaChainSpec}, lines 121--121)\\
\texttt{GlobalTheory.\allowbreak{}ChiGrid.\allowbreak{}rankV} (\texttt{Development.\allowbreak{}AbelianVarietyTheory.\allowbreak{}Definitions.\allowbreak{}Divisors.\allowbreak{}LeastSplitAuxiliaryPrimeThetaChainSpec}, lines 127--127)\\
\texttt{GlobalTheory.\allowbreak{}ChiGrid.\allowbreak{}rankP} (\texttt{Development.\allowbreak{}AbelianVarietyTheory.\allowbreak{}Definitions.\allowbreak{}Divisors.\allowbreak{}LeastSplitAuxiliaryPrimeThetaChainSpec}, lines 129--130)\\
\texttt{GlobalTheory.\allowbreak{}ChiGrid.\allowbreak{}rankX} (\texttt{Development.\allowbreak{}AbelianVarietyTheory.\allowbreak{}Definitions.\allowbreak{}Divisors.\allowbreak{}LeastSplitAuxiliaryPrimeThetaChainSpec}, lines 131--132)\\
\texttt{GlobalTheory.\allowbreak{}BigTop.\allowbreak{}Fun.\allowbreak{}Grid} (\texttt{Development.\allowbreak{}AbelianVarietyTheory.\allowbreak{}Definitions.\allowbreak{}Divisors.\allowbreak{}TwistedPlaceChoiceLiesOverPair}, lines 32--32)\par}

{\footnotesize\raggedright\noindent \ref{B8.21ca}.\quad \texttt{Formalism.\allowbreak{}MainNotions.\allowbreak{}TaylorLogMetric.\allowbreak{}IsCurveIdeal} (\texttt{Formalism.\allowbreak{}SharpRestriction.\allowbreak{}GlobalSchurBounds.\allowbreak{}GeneralLinearActions.\allowbreak{}IsDiscriminantTwelfthRootOfIdealLattice}, lines 56--57)\\
\texttt{Formalism.\allowbreak{}MainNotions.\allowbreak{}TaylorLogMetric.\allowbreak{}IsDiscriminantTwelfthRootOfIdealLattice} (\texttt{Formalism.\allowbreak{}SharpRestriction.\allowbreak{}GlobalSchurBounds.\allowbreak{}GeneralLinearActions.\allowbreak{}IsDiscriminantTwelfthRootOfIdealLattice}, lines 59--60)\\
\texttt{Formalism.\allowbreak{}MainNotions.\allowbreak{}StandardLatticesGL3.\allowbreak{}Crv} (\texttt{Formalism.\allowbreak{}SharpRestriction.\allowbreak{}GlobalSchurBounds.\allowbreak{}StandardLatticesGL3}, lines 28--29)\\
\texttt{Formalism.\allowbreak{}MainNotions.\allowbreak{}TaylorLogMetric.\allowbreak{}lamOf} (\texttt{Formalism.\allowbreak{}SharpRestriction.\allowbreak{}GlobalSchurBounds.\allowbreak{}GeneralLinearActions.\allowbreak{}IsDiscriminantTwelfthRootOfIdealLattice}, lines 53--55)\\
\texttt{Formalism.\allowbreak{}MainNotions.\allowbreak{}TaylorLogMetric.\allowbreak{}area} (\texttt{Formalism.\allowbreak{}SharpRestriction.\allowbreak{}GlobalSchurBounds.\allowbreak{}GeneralLinearActions.\allowbreak{}IsDiscriminantTwelfthRootOfIdealLattice}, lines 50--52)\\
\texttt{Formalism.\allowbreak{}MainNotions.\allowbreak{}GL3Theta.\allowbreak{}IsTwoTorsionPointUnkilled} (\texttt{Formalism.\allowbreak{}SharpRestriction.\allowbreak{}GlobalSchurBounds.\allowbreak{}GL3Theta}, lines 198--201)\\
\texttt{GlobalTheory.\allowbreak{}NormSum.\allowbreak{}SplitIn} (\texttt{Development.\allowbreak{}AbelianVarietyTheory.\allowbreak{}Definitions.\allowbreak{}Foundations.\allowbreak{}NormSumPairRankZeroOfCurveCoordNorm}, lines 112--112)\\
\texttt{GlobalTheory.\allowbreak{}NormSum.\allowbreak{}GoodLevel} (\texttt{Development.\allowbreak{}AbelianVarietyTheory.\allowbreak{}Definitions.\allowbreak{}Foundations.\allowbreak{}NormSumPairRankZeroOfCurveCoordNorm}, lines 166--167)\\
\texttt{Formalism.\allowbreak{}MainNotions.\allowbreak{}StandardLatticesGL3.\allowbreak{}GA} (\texttt{Formalism.\allowbreak{}SharpRestriction.\allowbreak{}GlobalSchurBounds.\allowbreak{}StandardLatticesGL3}, lines 374--376)\\
\texttt{GlobalTheory.\allowbreak{}NormSum.\allowbreak{}IntegralOver} (\texttt{Development.\allowbreak{}AbelianVarietyTheory.\allowbreak{}Statements.\allowbreak{}CuspidalUnitaryTwists.\allowbreak{}ExistsSmallSpherical}, lines 18--19)\\
\texttt{Formalism.\allowbreak{}MainNotions.\allowbreak{}TaylorLogMetric.\allowbreak{}VC} (\texttt{Formalism.\allowbreak{}SharpRestriction.\allowbreak{}GlobalSchurBounds.\allowbreak{}GeneralLinearActions.\allowbreak{}IsDiscriminantTwelfthRootOfIdealLattice}, lines 66--66)\\
\texttt{Formalism.\allowbreak{}MainNotions.\allowbreak{}StandardLatticesGL3.\allowbreak{}oddE0} (\texttt{Formalism.\allowbreak{}SharpRestriction.\allowbreak{}GlobalSchurBounds.\allowbreak{}StandardLatticesGL3}, lines 496--500)\\
\texttt{Formalism.\allowbreak{}MainNotions.\allowbreak{}StandardLatticesGL3.\allowbreak{}WeilO} (\texttt{Formalism.\allowbreak{}SharpRestriction.\allowbreak{}GlobalSchurBounds.\allowbreak{}StandardLatticesGL3}, lines 443--470)\\
\texttt{GlobalTheory.\allowbreak{}NormSum.\allowbreak{}Pres} (\texttt{Development.\allowbreak{}AbelianVarietyTheory.\allowbreak{}Definitions.\allowbreak{}Foundations.\allowbreak{}NormSumPairRankZeroOfCurveCoordNorm}, lines 114--119)\\
\texttt{GlobalTheory.\allowbreak{}NormSum.\allowbreak{}Cls} (\texttt{Development.\allowbreak{}AbelianVarietyTheory.\allowbreak{}Definitions.\allowbreak{}Foundations.\allowbreak{}NormSumPairRankZeroOfCurveCoordNorm}, lines 16--18)\\
\texttt{GlobalTheory.\allowbreak{}NormSum.\allowbreak{}Pair} (\texttt{Development.\allowbreak{}AbelianVarietyTheory.\allowbreak{}Definitions.\allowbreak{}Foundations.\allowbreak{}NormSumPairRankZeroOfCurveCoordNorm}, lines 121--150)\\
\texttt{GlobalTheory.\allowbreak{}NormSum.\allowbreak{}Pair.\allowbreak{}sq} (\texttt{Development.\allowbreak{}AbelianVarietyTheory.\allowbreak{}Definitions.\allowbreak{}Foundations.\allowbreak{}NormSumPairRankZeroOfCurveCoordNorm}, lines 159--160)\\
\texttt{GlobalTheory.\allowbreak{}NormSum.\allowbreak{}Pair.\allowbreak{}norm} (\texttt{Development.\allowbreak{}AbelianVarietyTheory.\allowbreak{}Definitions.\allowbreak{}Foundations.\allowbreak{}NormSumPairRankZeroOfCurveCoordNorm}, lines 161--163)\\
\texttt{GlobalTheory.\allowbreak{}NormSum.\allowbreak{}Pair.\allowbreak{}Small} (\texttt{Development.\allowbreak{}AbelianVarietyTheory.\allowbreak{}Definitions.\allowbreak{}Foundations.\allowbreak{}NormSumPairRankZeroOfCurveCoordNorm}, lines 164--165)\\
\texttt{GlobalTheory.\allowbreak{}FunObj.\allowbreak{}SqNum.\allowbreak{}botBig} (\texttt{Development.\allowbreak{}AbelianVarietyTheory.\allowbreak{}Definitions.\allowbreak{}Divisors.\allowbreak{}FinitePrimesUnder}, lines 171--171)\\
\texttt{GlobalTheory.\allowbreak{}FunObj.\allowbreak{}SqNum.\allowbreak{}botSmall} (\texttt{Development.\allowbreak{}AbelianVarietyTheory.\allowbreak{}Definitions.\allowbreak{}Divisors.\allowbreak{}FinitePrimesUnder}, lines 173--173)\\
\texttt{GlobalTheory.\allowbreak{}FunObj.\allowbreak{}SqNum.\allowbreak{}alphaH} (\texttt{Development.\allowbreak{}AbelianVarietyTheory.\allowbreak{}Definitions.\allowbreak{}Divisors.\allowbreak{}FinitePrimesUnder}, lines 235--235)\\
\texttt{GlobalTheory.\allowbreak{}FunObj.\allowbreak{}SqNum.\allowbreak{}betaH} (\texttt{Development.\allowbreak{}AbelianVarietyTheory.\allowbreak{}Definitions.\allowbreak{}Divisors.\allowbreak{}FinitePrimesUnder}, lines 237--237)\\
\texttt{GlobalTheory.\allowbreak{}FunObj.\allowbreak{}SqNum.\allowbreak{}HasChars} (\texttt{Development.\allowbreak{}AbelianVarietyTheory.\allowbreak{}Definitions.\allowbreak{}Divisors.\allowbreak{}FinitePrimesUnder}, lines 57--59)\\
\texttt{GlobalTheory.\allowbreak{}FunObj.\allowbreak{}SqNum.\allowbreak{}IsStdDatum} (\texttt{Development.\allowbreak{}AbelianVarietyTheory.\allowbreak{}Definitions.\allowbreak{}Divisors.\allowbreak{}FinitePrimesUnder}, lines 49--51)\\
\texttt{GlobalTheory.\allowbreak{}FunObj.\allowbreak{}SqNum.\allowbreak{}IsBlockLevel} (\texttt{Development.\allowbreak{}AbelianVarietyTheory.\allowbreak{}Definitions.\allowbreak{}Divisors.\allowbreak{}FinitePrimesUnder}, lines 53--55)\\
\texttt{GlobalTheory.\allowbreak{}FunObj.\allowbreak{}SqNum.\allowbreak{}IsBottom} (\texttt{Development.\allowbreak{}AbelianVarietyTheory.\allowbreak{}Definitions.\allowbreak{}Divisors.\allowbreak{}FinitePrimesUnder}, lines 256--269)\\
\texttt{GlobalTheory.\allowbreak{}FunObj.\allowbreak{}SqNum.\allowbreak{}wtA} (\texttt{Development.\allowbreak{}AbelianVarietyTheory.\allowbreak{}Definitions.\allowbreak{}Divisors.\allowbreak{}FinitePrimesUnder}, lines 192--193)\\
\texttt{GlobalTheory.\allowbreak{}FunObj.\allowbreak{}SqNum.\allowbreak{}wtA'} (\texttt{Development.\allowbreak{}AbelianVarietyTheory.\allowbreak{}Definitions.\allowbreak{}Divisors.\allowbreak{}FinitePrimesUnder}, lines 195--196)\\
\texttt{GlobalTheory.\allowbreak{}FunObj.\allowbreak{}SqNum.\allowbreak{}maxWt} (\texttt{Development.\allowbreak{}AbelianVarietyTheory.\allowbreak{}Definitions.\allowbreak{}Divisors.\allowbreak{}FinitePrimesUnder}, lines 204--207)\\
\texttt{GlobalTheory.\allowbreak{}NormSum.\allowbreak{}intBound} (\texttt{Development.\allowbreak{}AbelianVarietyTheory.\allowbreak{}Statements.\allowbreak{}CuspidalUnitaryTwists.\allowbreak{}ExistsSmallSpherical}, lines 21--24)\\
\texttt{GlobalTheory.\allowbreak{}NormSum.\allowbreak{}Exponents} (\texttt{Development.\allowbreak{}AbelianVarietyTheory.\allowbreak{}Statements.\allowbreak{}CuspidalUnitaryTwists.\allowbreak{}ExistsSmallSpherical}, lines 25--27)\\
\texttt{GlobalTheory.\allowbreak{}NormSum.\allowbreak{}CoordInt} (\texttt{Development.\allowbreak{}AbelianVarietyTheory.\allowbreak{}Statements.\allowbreak{}CuspidalUnitaryTwists.\allowbreak{}ExistsSmallSpherical}, lines 28--35)\\
\texttt{GlobalTheory.\allowbreak{}NormSum.\allowbreak{}NormInt} (\texttt{Development.\allowbreak{}AbelianVarietyTheory.\allowbreak{}Statements.\allowbreak{}CuspidalUnitaryTwists.\allowbreak{}ExistsSmallSpherical}, lines 36--43)\\
\texttt{Formalism.\allowbreak{}PairingDivisionLocalBalance.\allowbreak{}wt} (\texttt{Formalism.\allowbreak{}PairingDivisors.\allowbreak{}StatementH}, lines 16--18)\\
\texttt{GlobalTheory.\allowbreak{}FunObj.\allowbreak{}SqNum.\allowbreak{}IsIntExp} (\texttt{Development.\allowbreak{}AbelianVarietyTheory.\allowbreak{}Definitions.\allowbreak{}Divisors.\allowbreak{}FinitePrimesUnder}, lines 285--292)\\
\texttt{GlobalTheory.\allowbreak{}FunObj.\allowbreak{}SqNum.\allowbreak{}exists\_\allowbreak{}isIntExp} (\texttt{Development.\allowbreak{}AbelianVarietyTheory.\allowbreak{}Definitions.\allowbreak{}Divisors.\allowbreak{}FinitePrimesUnder}, lines 303--311)\par}

{\footnotesize\raggedright\noindent \ref{B8.21d}.\quad \texttt{Formalism.\allowbreak{}BasicNotions.\allowbreak{}SizeData₂} (\texttt{Formalism.\allowbreak{}LocalRingCardinality.\allowbreak{}DataParams₂}, lines 76--182)\\
\texttt{GlobalTheory.\allowbreak{}IndefiniteRowSideFullHecke.\allowbreak{}LevelPlusSubgroup.\allowbreak{}SizeRows} (\texttt{Development.\allowbreak{}AbelianVarietyTheory.\allowbreak{}BlockTheory.\allowbreak{}DivisorBase.\allowbreak{}LevelPlusSubgroup}, lines 20--31)\\
\texttt{Formalism.\allowbreak{}IndefR.\allowbreak{}Links.\allowbreak{}HostHyp} (\texttt{Formalism.\allowbreak{}LocalToGlobal.\allowbreak{}IndefiniteReal.\allowbreak{}LinksLemmas}, lines 496--504)\\
\texttt{GlobalTheory.\allowbreak{}SqTop.\allowbreak{}Rec} (\texttt{Development.\allowbreak{}AbelianVarietyTheory.\allowbreak{}Statements.\allowbreak{}PlaceShiftStatements.\allowbreak{}NormSmall}, lines 24--32)\\
\texttt{Formalism.\allowbreak{}IndefR.\allowbreak{}Links.\allowbreak{}Face3.\allowbreak{}lv2} (\texttt{Formalism.\allowbreak{}LocalToGlobal.\allowbreak{}IndefiniteReal.\allowbreak{}ZeroPartEigenvalueConvention}, lines 55--62)\\
\texttt{Formalism.\allowbreak{}L2.\allowbreak{}Theta.\allowbreak{}IsStdChar} (\texttt{Formalism.\allowbreak{}LocalToGlobal.\allowbreak{}BesselHermitianTheta.\allowbreak{}TiedWeil}, lines 59--62)\\
\texttt{Formalism.\allowbreak{}IndefR.\allowbreak{}Mi.\allowbreak{}SignsAtG} (\texttt{Formalism.\allowbreak{}LocalToGlobal.\allowbreak{}IndefiniteReal.\allowbreak{}RealSeparation.\allowbreak{}SignsG}, lines 23--31)\\
\texttt{GlobalTheory.\allowbreak{}BigTop.\allowbreak{}Fun.\allowbreak{}SettingAtPoint} (\texttt{Development.\allowbreak{}AutomorphicArithmetic.\allowbreak{}ThetaCorrespondence.\allowbreak{}ThetaPairLifts.\allowbreak{}ThirdPlaceSeries.\allowbreak{}SettingAtPoint}, lines 15--16)\\
\texttt{GlobalTheory.\allowbreak{}SqTop.\allowbreak{}GBig} (\texttt{Development.\allowbreak{}AbelianVarietyTheory.\allowbreak{}Statements.\allowbreak{}PlaceShiftStatements.\allowbreak{}NormSmall}, lines 34--36)\\
\texttt{Formalism.\allowbreak{}IndefR.\allowbreak{}Datum.\allowbreak{}IsWeilDatumO} (\texttt{Formalism.\allowbreak{}LocalToGlobal.\allowbreak{}IndefiniteReal.\allowbreak{}LocalNewformFunction}, lines 53--75)\\
\texttt{GlobalTheory.\allowbreak{}FunObj.\allowbreak{}locPl} (\texttt{Development.\allowbreak{}AbelianVarietyTheory.\allowbreak{}Definitions.\allowbreak{}Divisors.\allowbreak{}ExistsGivenStageEqOne}, lines 114--114)\\
\texttt{GlobalTheory.\allowbreak{}FunObj.\allowbreak{}cosExp} (\texttt{Development.\allowbreak{}AbelianVarietyTheory.\allowbreak{}Definitions.\allowbreak{}Divisors.\allowbreak{}ExistsGivenStageEqOne}, lines 118--119)\\
\texttt{GlobalTheory.\allowbreak{}FunObj.\allowbreak{}depthAt} (\texttt{Development.\allowbreak{}AbelianVarietyTheory.\allowbreak{}Definitions.\allowbreak{}Divisors.\allowbreak{}ExistsGivenStageEqOne}, lines 120--121)\\
\texttt{Formalism.\allowbreak{}L2.\allowbreak{}DIV.\allowbreak{}IsFrame} (\texttt{Formalism.\allowbreak{}LocalToGlobal.\allowbreak{}Divisibility.\allowbreak{}DefGroup}, lines 467--469)\\
\texttt{GlobalTheory.\allowbreak{}FunObj.\allowbreak{}IsSumFrame} (\texttt{Development.\allowbreak{}AbelianVarietyTheory.\allowbreak{}Definitions.\allowbreak{}Divisors.\allowbreak{}InvVolConjLevelOne}, lines 72--76)\\
\texttt{GlobalTheory.\allowbreak{}FunObj.\allowbreak{}Std} (\texttt{Development.\allowbreak{}AbelianVarietyTheory.\allowbreak{}Definitions.\allowbreak{}Divisors.\allowbreak{}KernelNormConstantSpec}, lines 203--235)\\
\texttt{GlobalTheory.\allowbreak{}FunObj.\allowbreak{}Normal} (\texttt{Development.\allowbreak{}AbelianVarietyTheory.\allowbreak{}Definitions.\allowbreak{}Divisors.\allowbreak{}KernelNormConstantSpec}, lines 257--282)\\
\texttt{GlobalTheory.\allowbreak{}FunObj.\allowbreak{}GridChoice} (\texttt{Development.\allowbreak{}AbelianVarietyTheory.\allowbreak{}Definitions.\allowbreak{}Divisors.\allowbreak{}InvVolConjLevelOne}, lines 185--207)\\
\texttt{Formalism.\allowbreak{}IndefR.\allowbreak{}Bridge.\allowbreak{}FamTie} (\texttt{Formalism.\allowbreak{}LocalToGlobal.\allowbreak{}IndefiniteReal.\allowbreak{}BridgeTie}, lines 79--103)\\
\texttt{GlobalTheory.\allowbreak{}FunObj.\allowbreak{}SumChoice} (\texttt{Development.\allowbreak{}AbelianVarietyTheory.\allowbreak{}Definitions.\allowbreak{}Divisors.\allowbreak{}InvVolConjLevelOne}, lines 388--396)\\
\texttt{GlobalTheory.\allowbreak{}RowTest.\allowbreak{}NumbersIn} (\texttt{Development.\allowbreak{}AbelianVarietyTheory.\allowbreak{}Definitions.\allowbreak{}TwistedInnerProducts.\allowbreak{}MemHalfCosetSelf}, lines 245--261)\\
\texttt{GlobalTheory.\allowbreak{}FunObj.\allowbreak{}LaP} (\texttt{Development.\allowbreak{}AbelianVarietyTheory.\allowbreak{}Definitions.\allowbreak{}Divisors.\allowbreak{}GridTestLatticeLargestSlopeBound}, lines 68--71)\\
\texttt{GlobalTheory.\allowbreak{}FunObj.\allowbreak{}TestIx} (\texttt{Development.\allowbreak{}AbelianVarietyTheory.\allowbreak{}Definitions.\allowbreak{}Divisors.\allowbreak{}GridTestLatticeLargestSlopeBound}, lines 80--81)\\
\texttt{GlobalTheory.\allowbreak{}FunObj.\allowbreak{}ClassReps} (\texttt{Development.\allowbreak{}AbelianVarietyTheory.\allowbreak{}Definitions.\allowbreak{}Divisors.\allowbreak{}GridTestLatticeLargestSlopeBound}, lines 82--87)\\
\texttt{Formalism.\allowbreak{}GradedWeylDivisibility.\allowbreak{}ssIdx} (\texttt{Formalism.\allowbreak{}GradedWeylDivisibility.\allowbreak{}StatementD}, lines 9--10)\\
\texttt{Formalism.\allowbreak{}L2.\allowbreak{}DIV.\allowbreak{}VVIdx} (\texttt{Formalism.\allowbreak{}LocalToGlobal.\allowbreak{}Divisibility.\allowbreak{}CoefMod}, lines 256--257)\\
\texttt{GlobalTheory.\allowbreak{}FunObj.\allowbreak{}Vol} (\texttt{Development.\allowbreak{}AbelianVarietyTheory.\allowbreak{}Definitions.\allowbreak{}Divisors.\allowbreak{}BplusOfEquivariant}, lines 45--53)\\
\texttt{GlobalTheory.\allowbreak{}RowTest.\allowbreak{}ChainVec} (\texttt{Development.\allowbreak{}AbelianVarietyTheory.\allowbreak{}Definitions.\allowbreak{}TwistedInnerProducts.\allowbreak{}ExistsMatchOver}, lines 67--73)\\
\texttt{GlobalTheory.\allowbreak{}FunObj.\allowbreak{}RealMatch} (\texttt{Development.\allowbreak{}AbelianVarietyTheory.\allowbreak{}Definitions.\allowbreak{}Divisors.\allowbreak{}ExistsNotRelZero}, lines 65--84)\\
\texttt{GlobalTheory.\allowbreak{}BigTop.\allowbreak{}Fun.\allowbreak{}SlopeTestRemainingChoices} (\texttt{Development.\allowbreak{}AbelianVarietyTheory.\allowbreak{}Definitions.\allowbreak{}Divisors.\allowbreak{}SlopeTestRemainingChoicesAlt}, lines 45--63)\\
\texttt{GlobalTheory.\allowbreak{}BigTop.\allowbreak{}Fun.\allowbreak{}Fld} (\texttt{Development.\allowbreak{}AbelianVarietyTheory.\allowbreak{}Definitions.\allowbreak{}Divisors.\allowbreak{}SlopeTestRemainingChoicesAlt}, lines 81--81)\\
\texttt{Formalism.\allowbreak{}L2.\allowbreak{}Host.\allowbreak{}ValLL0} (\texttt{Formalism.\allowbreak{}LocalToGlobal.\allowbreak{}UnitaryNorms.\allowbreak{}CharTied}, lines 35--40)\\
\texttt{GlobalTheory.\allowbreak{}BigTop.\allowbreak{}Fun.\allowbreak{}FieldHas} (\texttt{Development.\allowbreak{}AbelianVarietyTheory.\allowbreak{}Definitions.\allowbreak{}Divisors.\allowbreak{}TwistedPlaceChoiceLiesOverPair}, lines 148--153)\\
\texttt{GlobalTheory.\allowbreak{}BigTop.\allowbreak{}Fun.\allowbreak{}TwistedPlaceChoiceLiesOverPair} (\texttt{Development.\allowbreak{}AbelianVarietyTheory.\allowbreak{}Definitions.\allowbreak{}Divisors.\allowbreak{}TwistedPlaceChoiceLiesOverPair}, lines 195--198)\\
\texttt{Formalism.\allowbreak{}MainNotions.\allowbreak{}BlochKatoReduction.\allowbreak{}Gen.\allowbreak{}QOk} (\texttt{Formalism.\allowbreak{}SharpRestriction.\allowbreak{}LevelLetterCocycles.\allowbreak{}BlochKatoPhases}, lines 250--252)\\
\texttt{Formalism.\allowbreak{}MainNotions.\allowbreak{}BlochKatoReduction.\allowbreak{}Gen.\allowbreak{}partC} (\texttt{Formalism.\allowbreak{}SharpRestriction.\allowbreak{}LevelLetterCocycles.\allowbreak{}BlochKatoPhases}, lines 317--321)\\
\texttt{Formalism.\allowbreak{}L2.\allowbreak{}Host.\allowbreak{}PsiChoiceCompatible} (\texttt{Formalism.\allowbreak{}LocalToGlobal.\allowbreak{}UnitaryNorms.\allowbreak{}CharGal}, lines 129--134)\\
\texttt{GlobalTheory.\allowbreak{}LamTie.\allowbreak{}IsLamTieFull} (\texttt{Development.\allowbreak{}AbelianVarietyTheory.\allowbreak{}Definitions.\allowbreak{}NonGaloisDescent.\allowbreak{}NotTeichTrivOfIsDelta}, lines 106--113)\\
\texttt{GlobalTheory.\allowbreak{}BigTop.\allowbreak{}Fun.\allowbreak{}GridFit} (\texttt{Development.\allowbreak{}AbelianVarietyTheory.\allowbreak{}Definitions.\allowbreak{}Divisors.\allowbreak{}TwistedPlaceChoiceLiesOverPair}, lines 204--219)\\
\texttt{GlobalTheory.\allowbreak{}BigTop.\allowbreak{}Fun.\allowbreak{}HasGrid} (\texttt{Development.\allowbreak{}AbelianVarietyTheory.\allowbreak{}Definitions.\allowbreak{}Divisors.\allowbreak{}TwistedPlaceChoiceLiesOverPair}, lines 247--248)\\
\texttt{GlobalTheory.\allowbreak{}BigTop.\allowbreak{}Fun.\allowbreak{}Fitted} (\texttt{Development.\allowbreak{}AbelianVarietyTheory.\allowbreak{}Definitions.\allowbreak{}Divisors.\allowbreak{}TwistedPlaceChoiceLiesOverPair}, lines 250--250)\\
\texttt{GlobalTheory.\allowbreak{}BigTop.\allowbreak{}Fun.\allowbreak{}FunIn} (\texttt{Development.\allowbreak{}AbelianVarietyTheory.\allowbreak{}Definitions.\allowbreak{}Divisors.\allowbreak{}TwistedPlaceChoiceLiesOverPair}, lines 37--75)\\
\texttt{GlobalTheory.\allowbreak{}FunObj.\allowbreak{}IsTestLat} (\texttt{Development.\allowbreak{}AbelianVarietyTheory.\allowbreak{}Definitions.\allowbreak{}Divisors.\allowbreak{}GridTestLatticeLargestSlopeBound}, lines 116--118)\\
\texttt{GlobalTheory.\allowbreak{}BigTop.\allowbreak{}Fun.\allowbreak{}eV0} (\texttt{Development.\allowbreak{}AbelianVarietyTheory.\allowbreak{}Definitions.\allowbreak{}Divisors.\allowbreak{}TieTest}, lines 16--19)\\
\texttt{GlobalTheory.\allowbreak{}FunObj.\allowbreak{}SquareNonzero} (\texttt{Development.\allowbreak{}AbelianVarietyTheory.\allowbreak{}Definitions.\allowbreak{}Divisors.\allowbreak{}ExistsNotRelZero}, lines 217--219)\\
\texttt{GlobalTheory.\allowbreak{}FunObj.\allowbreak{}IsSquareConst} (\texttt{Development.\allowbreak{}AbelianVarietyTheory.\allowbreak{}Definitions.\allowbreak{}Divisors.\allowbreak{}ExistsNotRelZero}, lines 213--215)\\
\texttt{GlobalTheory.\allowbreak{}IndefiniteRowSideFullHecke.\allowbreak{}ClassZetaValuationBounds.\allowbreak{}IsClassZetaVals} (\texttt{Development.\allowbreak{}AbelianVarietyTheory.\allowbreak{}BlockTheory.\allowbreak{}DivisorBase.\allowbreak{}VpBounds}, lines 23--24)\\
\texttt{GlobalTheory.\allowbreak{}IndefiniteRowSideFullHecke.\allowbreak{}ClassZetaValuationBounds.\allowbreak{}betaBound\_\allowbreak{}vp} (\texttt{Development.\allowbreak{}AbelianVarietyTheory.\allowbreak{}BlockTheory.\allowbreak{}DivisorBase.\allowbreak{}VpBounds}, lines 66--67)\\
\texttt{GlobalTheory.\allowbreak{}FunObj.\allowbreak{}SqNum.\allowbreak{}IsDenom} (\texttt{Development.\allowbreak{}AbelianVarietyTheory.\allowbreak{}Definitions.\allowbreak{}Divisors.\allowbreak{}FinitePrimesUnder}, lines 281--284)\\
\texttt{GlobalTheory.\allowbreak{}FunObj.\allowbreak{}IsSqOf} (\texttt{Development.\allowbreak{}AbelianVarietyTheory.\allowbreak{}Definitions.\allowbreak{}Divisors.\allowbreak{}FinitePrimesUnder}, lines 318--347)\\
\texttt{GlobalTheory.\allowbreak{}BigTop.\allowbreak{}Fun.\allowbreak{}TieTest} (\texttt{Development.\allowbreak{}AbelianVarietyTheory.\allowbreak{}Definitions.\allowbreak{}Divisors.\allowbreak{}TieTest}, lines 25--26)\\
\texttt{GlobalTheory.\allowbreak{}BigTop.\allowbreak{}Fun.\allowbreak{}TieV} (\texttt{Development.\allowbreak{}AbelianVarietyTheory.\allowbreak{}Definitions.\allowbreak{}Divisors.\allowbreak{}TieTest}, lines 27--29)\\
\texttt{GlobalTheory.\allowbreak{}BigTop.\allowbreak{}Fun.\allowbreak{}TieSq} (\texttt{Development.\allowbreak{}AbelianVarietyTheory.\allowbreak{}Definitions.\allowbreak{}Divisors.\allowbreak{}TieTest}, lines 30--32)\\
\texttt{GlobalTheory.\allowbreak{}BigTop.\allowbreak{}Fun.\allowbreak{}TieAt} (\texttt{Development.\allowbreak{}AbelianVarietyTheory.\allowbreak{}Definitions.\allowbreak{}Divisors.\allowbreak{}TieTest}, lines 33--41)\\
\texttt{GlobalTheory.\allowbreak{}BigTop.\allowbreak{}Fun.\allowbreak{}Comparison₂} (\texttt{Development.\allowbreak{}AbelianVarietyTheory.\allowbreak{}Definitions.\allowbreak{}Divisors.\allowbreak{}TieTest}, lines 43--43)\par}

{\footnotesize\raggedright\noindent \ref{B8.21e}.\quad \texttt{Formalism.\allowbreak{}L2.\allowbreak{}Indef.\allowbreak{}Tree.\allowbreak{}FunData} (\texttt{Formalism.\allowbreak{}LocalToGlobal.\allowbreak{}IndefiniteForms.\allowbreak{}IndefiniteTree.\allowbreak{}FunctionalLowerBoundOfConditionsNumeric}, lines 20--92)\\
\texttt{Formalism.\allowbreak{}L2.\allowbreak{}Indef.\allowbreak{}Tree.\allowbreak{}FunData.\allowbreak{}s} (\texttt{Formalism.\allowbreak{}LocalToGlobal.\allowbreak{}IndefiniteForms.\allowbreak{}IndefiniteTree.\allowbreak{}FunctionalLowerBoundOfConditionsNumeric}, lines 102--103)\\
\texttt{Formalism.\allowbreak{}L2.\allowbreak{}Indef.\allowbreak{}Tree.\allowbreak{}FunData.\allowbreak{}bDel} (\texttt{Formalism.\allowbreak{}LocalToGlobal.\allowbreak{}IndefiniteForms.\allowbreak{}IndefiniteTree.\allowbreak{}FunctionalLowerBoundOfConditionsNumeric}, lines 139--142)\\
\texttt{Formalism.\allowbreak{}L2.\allowbreak{}Indef.\allowbreak{}Tree.\allowbreak{}FunData.\allowbreak{}bArch} (\texttt{Formalism.\allowbreak{}LocalToGlobal.\allowbreak{}IndefiniteForms.\allowbreak{}IndefiniteTree.\allowbreak{}FunctionalLowerBoundOfConditionsNumeric}, lines 143--146)\\
\texttt{Formalism.\allowbreak{}L2.\allowbreak{}Indef.\allowbreak{}Tree.\allowbreak{}FunData.\allowbreak{}bTH} (\texttt{Formalism.\allowbreak{}LocalToGlobal.\allowbreak{}IndefiniteForms.\allowbreak{}IndefiniteTree.\allowbreak{}FunctionalLowerBoundOfConditionsNumeric}, lines 147--148)\\
\texttt{GlobalTheory.\allowbreak{}BigTop.\allowbreak{}Fun.\allowbreak{}Close.\allowbreak{}cloudIcc} (\texttt{Development.\allowbreak{}AbelianVarietyTheory.\allowbreak{}Proofs.\allowbreak{}DivisorProofs.\allowbreak{}ZetaDenominatorsFitFunction}, lines 65--81)\\
\texttt{GlobalTheory.\allowbreak{}BigTop.\allowbreak{}Fun.\allowbreak{}Lc\_\allowbreak{}le\_\allowbreak{}of\_\allowbreak{}cloud\_\allowbreak{}le} (\texttt{Development.\allowbreak{}AbelianVarietyTheory.\allowbreak{}Proofs.\allowbreak{}DivisorProofs.\allowbreak{}RamifiedSumsFifthEstimate}, lines 30--45)\\
\texttt{GlobalTheory.\allowbreak{}BigTop.\allowbreak{}Fun.\allowbreak{}RamSum.\allowbreak{}leaves\_\allowbreak{}of\_\allowbreak{}sum} (\texttt{Development.\allowbreak{}AbelianVarietyTheory.\allowbreak{}Proofs.\allowbreak{}DivisorProofs.\allowbreak{}RamifiedSumsSeventhEstimate}, lines 39--69)\\
\texttt{GlobalTheory.\allowbreak{}BigTop.\allowbreak{}Fun.\allowbreak{}LeafMMTie.\allowbreak{}leafMM\_\allowbreak{}of\_\allowbreak{}test} (\texttt{Development.\allowbreak{}AbelianVarietyTheory.\allowbreak{}Proofs.\allowbreak{}DivisorProofs.\allowbreak{}MultiplicityAtFittedPoints}, lines 94--139)\\
\texttt{GlobalTheory.\allowbreak{}SizeLeafCheck.\allowbreak{}tree\_\allowbreak{}leaf\_\allowbreak{}parameter\_\allowbreak{}arithmetic\_\allowbreak{}of\_\allowbreak{}bounds} (\texttt{Development.\allowbreak{}AbelianVarietyTheory.\allowbreak{}Counterexamples.\allowbreak{}DegenerateExamples.\allowbreak{}TreeLeafParameterArithmeticOfBounds}, lines 66--202)\\
\texttt{Formalism.\allowbreak{}L2.\allowbreak{}Indef.\allowbreak{}Tree.\allowbreak{}PrimePowerModuliSumBound} (\texttt{Formalism.\allowbreak{}LocalToGlobal.\allowbreak{}IndefiniteForms.\allowbreak{}IndefiniteTree.\allowbreak{}FunctionalLowerBoundOfConditionsNumeric}, lines 219--220)\\
\texttt{Formalism.\allowbreak{}L2.\allowbreak{}Indef.\allowbreak{}Tree.\allowbreak{}DegreeAndSizeBounds} (\texttt{Formalism.\allowbreak{}LocalToGlobal.\allowbreak{}IndefiniteForms.\allowbreak{}IndefiniteTree.\allowbreak{}FunctionalLowerBoundOfConditionsNumeric}, lines 215--217)\par}

{\footnotesize\raggedright\noindent \ref{B8.21fa}.\quad \texttt{GlobalTheory.\allowbreak{}SumSize.\allowbreak{}Loss} (\texttt{Development.\allowbreak{}AbelianVarietyTheory.\allowbreak{}Statements.\allowbreak{}PlaceShiftStatements.\allowbreak{}InfClassAt}, lines 20--22)\\
\texttt{GlobalTheory.\allowbreak{}SumSize.\allowbreak{}LossLe} (\texttt{Development.\allowbreak{}AbelianVarietyTheory.\allowbreak{}Statements.\allowbreak{}PlaceShiftStatements.\allowbreak{}InfClassAt}, lines 23--24)\\
\texttt{GlobalTheory.\allowbreak{}SumSize.\allowbreak{}InLat} (\texttt{Development.\allowbreak{}AbelianVarietyTheory.\allowbreak{}Statements.\allowbreak{}PlaceShiftStatements.\allowbreak{}InfClassAt}, lines 29--32)\\
\texttt{DenominatorLossAdditivity.\allowbreak{}GlobalTheory.\allowbreak{}CompanionTheory.\allowbreak{}FinDen.\allowbreak{}core\_\allowbreak{}loss} (\texttt{Development.\allowbreak{}AbelianVarietyTheory.\allowbreak{}Proofs.\allowbreak{}DivisorProofs.\allowbreak{}FiniteSumDenominatorBoundOfFactorClaims}, lines 280--339)\\
\texttt{DenominatorLossAdditivity.\allowbreak{}GlobalTheory.\allowbreak{}CompanionTheory.\allowbreak{}FinDen.\allowbreak{}loss\_\allowbreak{}comb} (\texttt{Development.\allowbreak{}AbelianVarietyTheory.\allowbreak{}Proofs.\allowbreak{}DivisorProofs.\allowbreak{}FiniteSumDenominatorBoundOfFactorClaims}, lines 459--497)\\
\texttt{DenominatorLossAdditivity.\allowbreak{}GlobalTheory.\allowbreak{}CompanionTheory.\allowbreak{}FinDen.\allowbreak{}fin\_\allowbreak{}at} (\texttt{Development.\allowbreak{}AbelianVarietyTheory.\allowbreak{}Proofs.\allowbreak{}DivisorProofs.\allowbreak{}FiniteSumDenominatorBoundOfFactorClaims}, lines 507--532)\par}

{\footnotesize\raggedright\noindent \ref{B8.21f}.\quad \texttt{GlobalTheory.\allowbreak{}T.\allowbreak{}sum\_\allowbreak{}size\_\allowbreak{}bound\_\allowbreak{}all\_\allowbreak{}places} (\texttt{Development.\allowbreak{}AutomorphicArithmetic.\allowbreak{}ThetaCorrespondence.\allowbreak{}ThetaPairLifts.\allowbreak{}LocalTypeDictionary.\allowbreak{}SumSizeFiniteAwayArchimedeanCombined}, lines 50--51)\\
\texttt{GlobalTheory.\allowbreak{}SumSize.\allowbreak{}Join.\allowbreak{}sum\_\allowbreak{}size\_\allowbreak{}finite\_\allowbreak{}away\_\allowbreak{}archimedean\_\allowbreak{}combined} (\texttt{Development.\allowbreak{}AutomorphicArithmetic.\allowbreak{}ThetaCorrespondence.\allowbreak{}ThetaPairLifts.\allowbreak{}LocalTypeDictionary.\allowbreak{}SumSizeFiniteAwayArchimedeanCombined}, lines 26--44)\\
\texttt{GlobalTheory.\allowbreak{}SumSize.\allowbreak{}BigC} (\texttt{Development.\allowbreak{}AbelianVarietyTheory.\allowbreak{}Statements.\allowbreak{}PlaceShiftStatements.\allowbreak{}InfClassAt}, lines 83--88)\\
\texttt{GlobalTheory.\allowbreak{}SumSize.\allowbreak{}FinAt} (\texttt{Development.\allowbreak{}AbelianVarietyTheory.\allowbreak{}Statements.\allowbreak{}PlaceShiftStatements.\allowbreak{}InfClassAt}, lines 135--138)\\
\texttt{GlobalTheory.\allowbreak{}SumSize.\allowbreak{}AwayAt} (\texttt{Development.\allowbreak{}AbelianVarietyTheory.\allowbreak{}Statements.\allowbreak{}PlaceShiftStatements.\allowbreak{}InfClassAt}, lines 139--140)\\
\texttt{GlobalTheory.\allowbreak{}SumSize.\allowbreak{}FinK} (\texttt{Development.\allowbreak{}AbelianVarietyTheory.\allowbreak{}Statements.\allowbreak{}PlaceShiftStatements.\allowbreak{}InfClassAt}, lines 203--203)\\
\texttt{GlobalTheory.\allowbreak{}SumSize.\allowbreak{}AwayK} (\texttt{Development.\allowbreak{}AbelianVarietyTheory.\allowbreak{}Statements.\allowbreak{}PlaceShiftStatements.\allowbreak{}InfClassAt}, lines 205--205)\par}

{\footnotesize\raggedright\noindent \ref{B8.21g}.\quad \texttt{GlobalTheory.\allowbreak{}T.\allowbreak{}leafFixedPoint\_\allowbreak{}of\_\allowbreak{}finK\_\allowbreak{}awayK} (\texttt{Development.\allowbreak{}AutomorphicArithmetic.\allowbreak{}ThetaCorrespondence.\allowbreak{}ThetaPairLifts.\allowbreak{}ThirdPlaceSeries.\allowbreak{}LeafFixedPointOfFinKAwayKStatement}, lines 12--13)\\
\texttt{GlobalTheory.\allowbreak{}Adopt.\allowbreak{}leafFixedPoint\_\allowbreak{}of\_\allowbreak{}finK\_\allowbreak{}awayK} (\texttt{Development.\allowbreak{}AutomorphicArithmetic.\allowbreak{}ThetaCorrespondence.\allowbreak{}ThetaPairLifts.\allowbreak{}ThirdPlaceSeries.\allowbreak{}LeafFixedPointOfFinKAwayK}, lines 53--54)\\
\texttt{GlobalTheory.\allowbreak{}BigTop.\allowbreak{}Fun.\allowbreak{}LeafFixedPoint} (\texttt{Development.\allowbreak{}AutomorphicArithmetic.\allowbreak{}ThetaCorrespondence.\allowbreak{}ThetaPairLifts.\allowbreak{}ThirdPlaceSeries.\allowbreak{}LeafInfinitePlaceBound}, lines 28--28)\\
\texttt{Formalism.\allowbreak{}L2.\allowbreak{}Indef.\allowbreak{}Tree.\allowbreak{}LeafFP} (\texttt{Formalism.\allowbreak{}LocalToGlobal.\allowbreak{}IndefiniteForms.\allowbreak{}IndefiniteTree.\allowbreak{}FunctionalLowerBoundOfConditionsNumeric}, lines 168--173)\\
\texttt{GlobalTheory.\allowbreak{}finitePlaceBound\_\allowbreak{}inv\_\allowbreak{}rootOfUnity\_\allowbreak{}smul} (\texttt{Development.\allowbreak{}AbelianVarietyTheory.\allowbreak{}Proofs.\allowbreak{}DivisorProofs.\allowbreak{}MatLatSmulLeOfIsIntegral}, lines 24--32)\par}

{\footnotesize\raggedright\noindent \ref{B8.21ha}.\quad \texttt{GlobalTheory.\allowbreak{}QuotientNormCoisometry.\allowbreak{}htMap\_\allowbreak{}quot} (\texttt{Development.\allowbreak{}AbelianVarietyTheory.\allowbreak{}Proofs.\allowbreak{}DivisorProofs.\allowbreak{}InfinitePartIdentity}, lines 64--84)\\
\texttt{Formalism.\allowbreak{}HN.\allowbreak{}htMap} (\texttt{Formalism.\allowbreak{}HarderNarasimhan.\allowbreak{}StatementPl}, lines 60--63)\par}

{\footnotesize\raggedright\noindent \ref{B8.21h}.\quad \texttt{GlobalTheory.\allowbreak{}T.\allowbreak{}leafHeightMapBound\_\allowbreak{}of\_\allowbreak{}infK} (\texttt{Development.\allowbreak{}AutomorphicArithmetic.\allowbreak{}ThetaCorrespondence.\allowbreak{}ThetaPairLifts.\allowbreak{}ThirdPlaceSeries.\allowbreak{}IpaTieStatement}, lines 12--13)\\
\texttt{GlobalTheory.\allowbreak{}QuotientNormCoisometry.\allowbreak{}leafHeightMapBound\_\allowbreak{}tie\_\allowbreak{}of\_\allowbreak{}infK} (\texttt{Development.\allowbreak{}AutomorphicArithmetic.\allowbreak{}ThetaCorrespondence.\allowbreak{}ThetaPairLifts.\allowbreak{}ThirdPlaceSeries.\allowbreak{}IpaTie}, lines 22--50)\\
\texttt{GlobalTheory.\allowbreak{}BigTop.\allowbreak{}Fun.\allowbreak{}LeafHeightMapBound} (\texttt{Development.\allowbreak{}AutomorphicArithmetic.\allowbreak{}ThetaCorrespondence.\allowbreak{}ThetaPairLifts.\allowbreak{}ThirdPlaceSeries.\allowbreak{}LeafInfinitePlaceBound}, lines 29--33)\\
\texttt{GlobalTheory.\allowbreak{}SumSize.\allowbreak{}InfAt} (\texttt{Development.\allowbreak{}AbelianVarietyTheory.\allowbreak{}Statements.\allowbreak{}PlaceShiftStatements.\allowbreak{}InfClassAt}, lines 176--180)\\
\texttt{GlobalTheory.\allowbreak{}SumSize.\allowbreak{}InfK} (\texttt{Development.\allowbreak{}AbelianVarietyTheory.\allowbreak{}Statements.\allowbreak{}PlaceShiftStatements.\allowbreak{}InfClassAt}, lines 225--225)\\
\texttt{Formalism.\allowbreak{}HN.\allowbreak{}htMap} (\texttt{Formalism.\allowbreak{}HarderNarasimhan.\allowbreak{}StatementPl}, lines 60--62)\par}

{\footnotesize\raggedright\noindent \ref{B8.21i}.\quad \texttt{GlobalTheory.\allowbreak{}T.\allowbreak{}leafMinimalMass\_\allowbreak{}of\_\allowbreak{}setAt} (\texttt{Development.\allowbreak{}AutomorphicArithmetic.\allowbreak{}ThetaCorrespondence.\allowbreak{}ThetaPairLifts.\allowbreak{}ThirdPlaceSeries.\allowbreak{}MmTieStatement}, lines 11--12)\\
\texttt{GlobalTheory.\allowbreak{}BigTop.\allowbreak{}Fun.\allowbreak{}LeafMMTie.\allowbreak{}mm\_\allowbreak{}tie\_\allowbreak{}alt} (\texttt{Development.\allowbreak{}AutomorphicArithmetic.\allowbreak{}ThetaCorrespondence.\allowbreak{}ThetaPairLifts.\allowbreak{}ThirdPlaceSeries.\allowbreak{}MmTie}, lines 23--26)\\
\texttt{GlobalTheory.\allowbreak{}BigTop.\allowbreak{}Fun.\allowbreak{}LeafMinimalMass} (\texttt{Development.\allowbreak{}AutomorphicArithmetic.\allowbreak{}ThetaCorrespondence.\allowbreak{}ThetaPairLifts.\allowbreak{}ThirdPlaceSeries.\allowbreak{}LeafInfinitePlaceBound}, lines 23--23)\\
\texttt{Formalism.\allowbreak{}L2.\allowbreak{}Indef.\allowbreak{}Tree.\allowbreak{}LeafMM} (\texttt{Formalism.\allowbreak{}LocalToGlobal.\allowbreak{}IndefiniteForms.\allowbreak{}IndefiniteTree.\allowbreak{}FunctionalLowerBoundOfConditionsNumeric}, lines 158--159)\par}

{\footnotesize\raggedright\noindent \ref{B8.21ja}.\quad \texttt{GlobalTheory.\allowbreak{}SqTop.\allowbreak{}OneSqProof.\allowbreak{}vw\_\allowbreak{}ratio\_\allowbreak{}le} (\texttt{Development.\allowbreak{}AbelianVarietyTheory.\allowbreak{}Proofs.\allowbreak{}DivisorProofs.\allowbreak{}OneSqStOfClaims}, lines 16--72)\par}

{\footnotesize\raggedright\noindent \ref{B8.21j}.\quad \texttt{GlobalTheory.\allowbreak{}SqTop.\allowbreak{}OneSqSt} (\texttt{Development.\allowbreak{}AbelianVarietyTheory.\allowbreak{}Statements.\allowbreak{}PlaceShiftStatements.\allowbreak{}NormSmall}, lines 159--167)\\
\texttt{GlobalTheory.\allowbreak{}SqTop.\allowbreak{}ThA} (\texttt{Development.\allowbreak{}AbelianVarietyTheory.\allowbreak{}Statements.\allowbreak{}PlaceShiftStatements.\allowbreak{}NormSmall}, lines 66--70)\\
\texttt{GlobalTheory.\allowbreak{}SqTop.\allowbreak{}DenSum} (\texttt{Development.\allowbreak{}AbelianVarietyTheory.\allowbreak{}Statements.\allowbreak{}PlaceShiftStatements.\allowbreak{}NormSmall}, lines 89--90)\\
\texttt{GlobalTheory.\allowbreak{}SqTop.\allowbreak{}square\_\allowbreak{}top\_\allowbreak{}term\_\allowbreak{}joined\_\allowbreak{}estimates} (\texttt{Development.\allowbreak{}AutomorphicArithmetic.\allowbreak{}ThetaCorrespondence.\allowbreak{}ThetaPairLifts.\allowbreak{}LocalTypeDictionary.\allowbreak{}SquareTopTermJoinedEstimates}, lines 13--35)\\
\texttt{GlobalTheory.\allowbreak{}ToricTransport.\allowbreak{}toricRatioConjugates\_\allowbreak{}of\_\allowbreak{}transport} (\texttt{Development.\allowbreak{}AbelianVarietyTheory.\allowbreak{}Proofs.\allowbreak{}DivisorProofs.\allowbreak{}ToricRatioConjugatesOfTransport}, lines 57--86)\\
\texttt{GlobalTheory.\allowbreak{}ThmAUpper.\allowbreak{}thmA\_\allowbreak{}upper\_\allowbreak{}alt} (\texttt{Development.\allowbreak{}AutomorphicArithmetic.\allowbreak{}ThetaCorrespondence.\allowbreak{}ThetaPairLifts.\allowbreak{}ThirdPlaceSeries.\allowbreak{}ThmAUpper}, lines 17--29)\\
\texttt{GlobalTheory.\allowbreak{}BigTop.\allowbreak{}Fun.\allowbreak{}oneSq\_\allowbreak{}of\_\allowbreak{}log\_\allowbreak{}derivative\_\allowbreak{}bound\_\allowbreak{}thA} (\texttt{Development.\allowbreak{}AutomorphicArithmetic.\allowbreak{}ThetaCorrespondence.\allowbreak{}ThetaPairLifts.\allowbreak{}ThirdPlaceSeries.\allowbreak{}OneSqTie}, lines 23--79)\\
\texttt{GlobalTheory.\allowbreak{}BigTop.\allowbreak{}Fun.\allowbreak{}den\_\allowbreak{}tie\_\allowbreak{}alt} (\texttt{Development.\allowbreak{}AutomorphicArithmetic.\allowbreak{}ThetaCorrespondence.\allowbreak{}ThetaPairLifts.\allowbreak{}ThirdPlaceSeries.\allowbreak{}DenTie}, lines 14--26)\\
\texttt{GlobalTheory.\allowbreak{}T.\allowbreak{}oneSq\_\allowbreak{}of\_\allowbreak{}log\_\allowbreak{}derivative\_\allowbreak{}bound\_\allowbreak{}and\_\allowbreak{}oneSqSt} (\texttt{Development.\allowbreak{}AutomorphicArithmetic.\allowbreak{}ThetaCorrespondence.\allowbreak{}ThetaPairLifts.\allowbreak{}ThirdPlaceSeries.\allowbreak{}OneSqTieStatement}, lines 13--14)\\
\texttt{GlobalTheory.\allowbreak{}T.\allowbreak{}leafDensity\_\allowbreak{}of\_\allowbreak{}logDeriv\_\allowbreak{}bound\_\allowbreak{}and\_\allowbreak{}denSum} (\texttt{Development.\allowbreak{}AutomorphicArithmetic.\allowbreak{}ThetaCorrespondence.\allowbreak{}ThetaPairLifts.\allowbreak{}ThirdPlaceSeries.\allowbreak{}DenTieStatement}, lines 13--14)\par}

{\footnotesize\raggedright\noindent \ref{B8.21ka}.\quad \texttt{GlobalTheory.\allowbreak{}NormSum.\allowbreak{}DenomBound.\allowbreak{}exists\_\allowbreak{}exponents} (\texttt{Development.\allowbreak{}AbelianVarietyTheory.\allowbreak{}Proofs.\allowbreak{}FoundationProofs.\allowbreak{}ClassSumDenominatorBoundOfCmThetaHypotheses}, lines 48--111)\\
\texttt{GlobalTheory.\allowbreak{}NormSum.\allowbreak{}intBound} (\texttt{Development.\allowbreak{}AbelianVarietyTheory.\allowbreak{}Statements.\allowbreak{}CuspidalUnitaryTwists.\allowbreak{}ExistsSmallSpherical}, lines 21--24)\par}

{\footnotesize\raggedright\noindent \ref{B8.21k}.\quad \texttt{GlobalTheory.\allowbreak{}T.\allowbreak{}cmThetaNormBoundHypotheses\_\allowbreak{}satisfied} (\texttt{Development.\allowbreak{}AbelianVarietyTheory.\allowbreak{}Proofs.\allowbreak{}FoundationProofs.\allowbreak{}CmThetaNormBoundHypothesesSatisfiedStatement}, lines 8--9)\\
\texttt{GlobalTheory.\allowbreak{}NormSum.\allowbreak{}inp\_\allowbreak{}all} (\texttt{Development.\allowbreak{}AbelianVarietyTheory.\allowbreak{}Proofs.\allowbreak{}FoundationProofs.\allowbreak{}CmThetaNormBoundHypothesesSatisfied}, lines 29--51)\\
\texttt{GlobalTheory.\allowbreak{}T.\allowbreak{}classSumDenominatorBound\_\allowbreak{}of\_\allowbreak{}cmThetaHypotheses} (\texttt{Development.\allowbreak{}AbelianVarietyTheory.\allowbreak{}Proofs.\allowbreak{}FoundationProofs.\allowbreak{}ClassSumDenominatorBoundOfCmThetaHypothesesStatement}, lines 9--10)\\
\texttt{GlobalTheory.\allowbreak{}NormSum.\allowbreak{}DenomBound.\allowbreak{}intSt\_\allowbreak{}of\_\allowbreak{}forall\_\allowbreak{}inp} (\texttt{Development.\allowbreak{}AbelianVarietyTheory.\allowbreak{}Proofs.\allowbreak{}FoundationProofs.\allowbreak{}ClassSumDenominatorBoundOfCmThetaHypotheses}, lines 118--148)\\
\texttt{GlobalTheory.\allowbreak{}NormSum.\allowbreak{}IntSt} (\texttt{Development.\allowbreak{}AbelianVarietyTheory.\allowbreak{}Statements.\allowbreak{}CuspidalUnitaryTwists.\allowbreak{}ExistsSmallSpherical}, lines 44--47)\\
\texttt{GlobalTheory.\allowbreak{}FunObj.\allowbreak{}SqNum.\allowbreak{}exists\_\allowbreak{}isIntExp} (\texttt{Development.\allowbreak{}AbelianVarietyTheory.\allowbreak{}Definitions.\allowbreak{}Divisors.\allowbreak{}FinitePrimesUnder}, lines 303--311)\\
\texttt{GlobalTheory.\allowbreak{}BigTop.\allowbreak{}Fun.\allowbreak{}TieModel.\allowbreak{}isIntExp\_\allowbreak{}mono} (\texttt{Development.\allowbreak{}AbelianVarietyTheory.\allowbreak{}Definitions.\allowbreak{}Divisors.\allowbreak{}MulVecEqOfBasis}, lines 38--51)\par}

{\footnotesize\raggedright\noindent \ref{B8.21l}.\quad \texttt{GlobalTheory.\allowbreak{}SqTop.\allowbreak{}BotSmall} (\texttt{Development.\allowbreak{}AbelianVarietyTheory.\allowbreak{}Statements.\allowbreak{}PlaceShiftStatements.\allowbreak{}NormSmall}, lines 110--112)\\
\texttt{GlobalTheory.\allowbreak{}SqTop.\allowbreak{}BotNe} (\texttt{Development.\allowbreak{}AbelianVarietyTheory.\allowbreak{}Statements.\allowbreak{}PlaceShiftStatements.\allowbreak{}NormSmall}, lines 122--123)\\
\texttt{GlobalTheory.\allowbreak{}BigTop.\allowbreak{}Fun.\allowbreak{}int\_\allowbreak{}tie\_\allowbreak{}alt} (\texttt{Development.\allowbreak{}AutomorphicArithmetic.\allowbreak{}ThetaCorrespondence.\allowbreak{}ThetaPairLifts.\allowbreak{}ThirdPlaceSeries.\allowbreak{}IntTie}, lines 17--76)\\
\texttt{GlobalTheory.\allowbreak{}T.\allowbreak{}leafIntegrality\_\allowbreak{}of\_\allowbreak{}logDeriv\_\allowbreak{}bound} (\texttt{Development.\allowbreak{}AutomorphicArithmetic.\allowbreak{}ThetaCorrespondence.\allowbreak{}ThetaPairLifts.\allowbreak{}ThirdPlaceSeries.\allowbreak{}IntTieStatement}, lines 13--14)\\
\texttt{Formalism.\allowbreak{}L2.\allowbreak{}Indef.\allowbreak{}Tree.\allowbreak{}LeafINT} (\texttt{Formalism.\allowbreak{}LocalToGlobal.\allowbreak{}IndefiniteForms.\allowbreak{}IndefiniteTree.\allowbreak{}FunctionalLowerBoundOfConditionsNumeric}, lines 197--200)\par}

{\footnotesize\raggedright\noindent \ref{B8.21ma}.\quad \texttt{GlobalTheory.\allowbreak{}PositiveMeans.\allowbreak{}means\_\allowbreak{}le} (\texttt{Development.\allowbreak{}AbelianVarietyTheory.\allowbreak{}Proofs.\allowbreak{}FoundationProofs.\allowbreak{}ExistsBoundMeansPosLogAreaNormalisedTheta}, lines 233--297)\\
\texttt{GlobalTheory.\allowbreak{}ArchFree.\allowbreak{}means.\allowbreak{}means\_\allowbreak{}le} (\texttt{Development.\allowbreak{}AbelianVarietyTheory.\allowbreak{}Proofs.\allowbreak{}FoundationProofs.\allowbreak{}ConjugateLatticeMeansPosLogAreaThetaBounded}, lines 55--119)\\
\texttt{GlobalTheory.\allowbreak{}LatticeAreaBound.\allowbreak{}area\_\allowbreak{}bounded\_\allowbreak{}of\_\allowbreak{}discr\_\allowbreak{}eq\_\allowbreak{}one} (\texttt{Development.\allowbreak{}AbelianVarietyTheory.\allowbreak{}Proofs.\allowbreak{}FoundationProofs.\allowbreak{}AreaBoundedOfDiscrEqOne}, lines 59--83)\par}

{\footnotesize\raggedright\noindent \ref{B8.21m}.\quad \texttt{GlobalTheory.\allowbreak{}ArchNorm.\allowbreak{}norm\_\allowbreak{}meanInf\_\allowbreak{}posLog\_\allowbreak{}le} (\texttt{Development.\allowbreak{}AbelianVarietyTheory.\allowbreak{}Proofs.\allowbreak{}FoundationProofs.\allowbreak{}LatticeUnitNormSqNonneg}, lines 225--342)\\
\texttt{GlobalTheory.\allowbreak{}ArchTop.\allowbreak{}archSt\_\allowbreak{}of\_\allowbreak{}inp} (\texttt{Development.\allowbreak{}AbelianVarietyTheory.\allowbreak{}Proofs.\allowbreak{}FoundationProofs.\allowbreak{}ThetaJetClassSumArchimedeanBoundOfCmTheta}, lines 12--18)\\
\texttt{GlobalTheory.\allowbreak{}NormSum.\allowbreak{}ArchSt} (\texttt{Development.\allowbreak{}AbelianVarietyTheory.\allowbreak{}Statements.\allowbreak{}CuspidalUnitaryTwists.\allowbreak{}ExistsSmallSpherical}, lines 49--59)\\
\texttt{GlobalTheory.\allowbreak{}BigTop.\allowbreak{}Fun.\allowbreak{}arch\_\allowbreak{}tie\_\allowbreak{}of\_\allowbreak{}logDeriv\_\allowbreak{}bound} (\texttt{Development.\allowbreak{}AutomorphicArithmetic.\allowbreak{}ThetaCorrespondence.\allowbreak{}ThetaPairLifts.\allowbreak{}ThirdPlaceSeries.\allowbreak{}ArchTie}, lines 14--47)\\
\texttt{GlobalTheory.\allowbreak{}T.\allowbreak{}arch\_\allowbreak{}tie\_\allowbreak{}of\_\allowbreak{}setAt\_\allowbreak{}logDeriv} (\texttt{Development.\allowbreak{}AutomorphicArithmetic.\allowbreak{}ThetaCorrespondence.\allowbreak{}ThetaPairLifts.\allowbreak{}ThirdPlaceSeries.\allowbreak{}ArchTieStatement}, lines 13--14)\\
\texttt{Formalism.\allowbreak{}L2.\allowbreak{}Indef.\allowbreak{}Tree.\allowbreak{}LeafARCH} (\texttt{Formalism.\allowbreak{}LocalToGlobal.\allowbreak{}IndefiniteForms.\allowbreak{}IndefiniteTree.\allowbreak{}FunctionalLowerBoundOfConditionsNumeric}, lines 202--204)\par}

{\footnotesize\raggedright\noindent \ref{B8.21n}.\quad \texttt{GlobalTheory.\allowbreak{}T.\allowbreak{}existsFunctionIn\_\allowbreak{}of\_\allowbreak{}sqRatio\_\allowbreak{}range} (\texttt{Development.\allowbreak{}AutomorphicArithmetic.\allowbreak{}ThetaCorrespondence.\allowbreak{}ThetaPairLifts.\allowbreak{}ThirdPlaceSeries.\allowbreak{}ExOfClaimsStatement}, lines 13--14)\\
\texttt{GlobalTheory.\allowbreak{}BigTop.\allowbreak{}Fun.\allowbreak{}ex\_\allowbreak{}of\_\allowbreak{}claims\_\allowbreak{}alt} (\texttt{Development.\allowbreak{}AutomorphicArithmetic.\allowbreak{}ThetaCorrespondence.\allowbreak{}ThetaPairLifts.\allowbreak{}ThirdPlaceSeries.\allowbreak{}ExOfClaims}, lines 16--74)\\
\texttt{GlobalTheory.\allowbreak{}BigTop.\allowbreak{}Fun.\allowbreak{}square\_\allowbreak{}ratio\_\allowbreak{}at\_\allowbreak{}point\_\allowbreak{}of\_\allowbreak{}leading\_\allowbreak{}term\_\allowbreak{}alt} (\texttt{Development.\allowbreak{}AutomorphicArithmetic.\allowbreak{}ThetaCorrespondence.\allowbreak{}ThetaPairLifts.\allowbreak{}ThirdPlaceSeries.\allowbreak{}SquareRatioAtPointOfLeadingTerm}, lines 19--45)\par}

{\footnotesize\raggedright\noindent \ref{B8.21o}.\quad \texttt{GlobalTheory.\allowbreak{}BigTop.\allowbreak{}Fun.\allowbreak{}oneMap\_\allowbreak{}of\_\allowbreak{}bundle\_\allowbreak{}setAt\_\allowbreak{}kernel\_\allowbreak{}conditions} (\texttt{Development.\allowbreak{}AutomorphicArithmetic.\allowbreak{}ThetaCorrespondence.\allowbreak{}ThetaPairLifts.\allowbreak{}ThirdPlaceSeries.\allowbreak{}OneMapOfRows}, lines 128--133)\\
\texttt{GlobalTheory.\allowbreak{}BigTop.\allowbreak{}Fun.\allowbreak{}leafKernel\_\allowbreak{}of\_\allowbreak{}linK\_\allowbreak{}injK\_\allowbreak{}constK\_\allowbreak{}formula2K} (\texttt{Development.\allowbreak{}AutomorphicArithmetic.\allowbreak{}ThetaCorrespondence.\allowbreak{}ThetaPairLifts.\allowbreak{}ThirdPlaceSeries.\allowbreak{}KerEqOfFormula}, lines 87--101)\\
\texttt{GlobalTheory.\allowbreak{}RowClaim.\allowbreak{}JoinF.\allowbreak{}joined\_\allowbreak{}local\_\allowbreak{}sum\_\allowbreak{}estimates} (\texttt{Development.\allowbreak{}AutomorphicArithmetic.\allowbreak{}ThetaCorrespondence.\allowbreak{}ThetaPairLifts.\allowbreak{}LocalTypeDictionary.\allowbreak{}JoinedLocalSumEstimates}, lines 12--16)\\
\texttt{GlobalTheory.\allowbreak{}RowTest.\allowbreak{}injK\_\allowbreak{}of\_\allowbreak{}separates} (\texttt{Development.\allowbreak{}AbelianVarietyTheory.\allowbreak{}Proofs.\allowbreak{}TwistedOperators.\allowbreak{}InjKOfSeparates}, lines 91--104)\\
\texttt{GlobalTheory.\allowbreak{}BigTop.\allowbreak{}Fun.\allowbreak{}bundle\_\allowbreak{}of\_\allowbreak{}rows\_\allowbreak{}alt} (\texttt{Development.\allowbreak{}AutomorphicArithmetic.\allowbreak{}ThetaCorrespondence.\allowbreak{}ThetaPairLifts.\allowbreak{}ThirdPlaceSeries.\allowbreak{}BundleOfRows}, lines 13--17)\par}

{\footnotesize\raggedright\noindent \ref{B8.21p}.\quad \texttt{GlobalTheory.\allowbreak{}BigTop.\allowbreak{}Fun.\allowbreak{}existsIntegralMultiple\_\allowbreak{}holds\_\allowbreak{}alt} (\texttt{Development.\allowbreak{}AutomorphicArithmetic.\allowbreak{}ThetaCorrespondence.\allowbreak{}ThetaPairLifts.\allowbreak{}ThirdPlaceSeries.\allowbreak{}ExistsIntegralMultipleHolds}, lines 12--26)\\
\texttt{GlobalTheory.\allowbreak{}BigTop.\allowbreak{}Fun.\allowbreak{}Close.\allowbreak{}desHLeaf\_\allowbreak{}alt} (\texttt{Development.\allowbreak{}AutomorphicArithmetic.\allowbreak{}ThetaCorrespondence.\allowbreak{}ThetaPairLifts.\allowbreak{}ThirdPlaceSeries.\allowbreak{}DesHLeaf}, lines 13--55)\\
\texttt{GlobalTheory.\allowbreak{}BigTop.\allowbreak{}Fun.\allowbreak{}rankGapBound\_\allowbreak{}alt} (\texttt{Development.\allowbreak{}AutomorphicArithmetic.\allowbreak{}ThetaCorrespondence.\allowbreak{}ThetaPairLifts.\allowbreak{}ThirdPlaceSeries.\allowbreak{}RamifiedSumsFifthBound}, lines 16--127)\\
\texttt{GlobalTheory.\allowbreak{}BigTop.\allowbreak{}Fun.\allowbreak{}ramified\_\allowbreak{}sums\_\allowbreak{}seventh\_\allowbreak{}bound\_\allowbreak{}alt} (\texttt{Development.\allowbreak{}AutomorphicArithmetic.\allowbreak{}ThetaCorrespondence.\allowbreak{}ThetaPairLifts.\allowbreak{}ThirdPlaceSeries.\allowbreak{}RamifiedSumsSeventhBound}, lines 18--38)\\
\texttt{GlobalTheory.\allowbreak{}BigTop.\allowbreak{}Fun.\allowbreak{}wt\_\allowbreak{}tie\_\allowbreak{}alt} (\texttt{Development.\allowbreak{}AutomorphicArithmetic.\allowbreak{}ThetaCorrespondence.\allowbreak{}ThetaPairLifts.\allowbreak{}ThirdPlaceSeries.\allowbreak{}WtTie}, lines 20--41)\\
\texttt{GlobalTheory.\allowbreak{}BigTop.\allowbreak{}Fun.\allowbreak{}desPol\_\allowbreak{}alt} (\texttt{Development.\allowbreak{}AutomorphicArithmetic.\allowbreak{}ThetaCorrespondence.\allowbreak{}ThetaPairLifts.\allowbreak{}ThirdPlaceSeries.\allowbreak{}DesPol}, lines 17--58)\par}

{\footnotesize\raggedright\noindent \ref{prop:p8c2-fitted-evaluation}.\quad \texttt{GlobalTheory.\allowbreak{}BigTop.\allowbreak{}Fun.\allowbreak{}Close.\allowbreak{}fittedFunction\_\allowbreak{}at\_\allowbreak{}zero\_\allowbreak{}close\_\allowbreak{}alt} (\texttt{Development.\allowbreak{}AutomorphicArithmetic.\allowbreak{}ThetaCorrespondence.\allowbreak{}ThetaPairLifts.\allowbreak{}ThirdPlaceSeries.\allowbreak{}FittedFunctionAtZeroClose}, lines 42--67)\\
\texttt{GlobalTheory.\allowbreak{}BigTop.\allowbreak{}Fun.\allowbreak{}fun0Fit\_\allowbreak{}of\_\allowbreak{}leaves\_\allowbreak{}alt} (\texttt{Development.\allowbreak{}AutomorphicArithmetic.\allowbreak{}ThetaCorrespondence.\allowbreak{}ThetaPairLifts.\allowbreak{}ThirdPlaceSeries.\allowbreak{}LeafInfinitePlaceBound}, lines 83--120)\\
\texttt{Formalism.\allowbreak{}L2.\allowbreak{}Indef.\allowbreak{}Tree.\allowbreak{}LeafOneSquare} (\texttt{Formalism.\allowbreak{}LocalToGlobal.\allowbreak{}IndefiniteForms.\allowbreak{}IndefiniteTree.\allowbreak{}FunctionalLowerBoundOfConditionsNumeric}, lines 182--191)\\
\texttt{Formalism.\allowbreak{}L2.\allowbreak{}Indef.\allowbreak{}Tree.\allowbreak{}PrimeCountingSumBound} (\texttt{Formalism.\allowbreak{}LocalToGlobal.\allowbreak{}IndefiniteForms.\allowbreak{}IndefiniteTree.\allowbreak{}FunctionalLowerBoundOfConditionsNumeric}, lines 193--195)\par}

{\footnotesize\raggedright\noindent \ref{prop:p8c2-fitted-evaluation}.\quad \texttt{GlobalTheory.\allowbreak{}T.\allowbreak{}oneMap\_\allowbreak{}of\_\allowbreak{}bundle\_\allowbreak{}setAt\_\allowbreak{}lin\_\allowbreak{}inj\_\allowbreak{}const\_\allowbreak{}formula} (\texttt{Development.\allowbreak{}AutomorphicArithmetic.\allowbreak{}ThetaCorrespondence.\allowbreak{}ThetaPairLifts.\allowbreak{}ThirdPlaceSeries.\allowbreak{}OneMapOfRowsStatement}, lines 11--12)\\
\texttt{GlobalTheory.\allowbreak{}T.\allowbreak{}leafKernel\_\allowbreak{}tie\_\allowbreak{}of\_\allowbreak{}claim\_\allowbreak{}inputs} (\texttt{Development.\allowbreak{}AutomorphicArithmetic.\allowbreak{}ThetaCorrespondence.\allowbreak{}ThetaPairLifts.\allowbreak{}ThirdPlaceSeries.\allowbreak{}KerEqOfFormulaStatement}, lines 11--12)\par}

{\footnotesize\raggedright\noindent \ref{prop:p8c2-fitted-evaluation}.\quad \texttt{GlobalTheory.\allowbreak{}T.\allowbreak{}sum\_\allowbreak{}size\_\allowbreak{}bound\_\allowbreak{}all\_\allowbreak{}places} (\texttt{Development.\allowbreak{}AutomorphicArithmetic.\allowbreak{}ThetaCorrespondence.\allowbreak{}ThetaPairLifts.\allowbreak{}LocalTypeDictionary.\allowbreak{}SumSizeFiniteAwayArchimedeanCombined}, lines 50--51)\\
\texttt{GlobalTheory.\allowbreak{}T.\allowbreak{}leafFixedPoint\_\allowbreak{}of\_\allowbreak{}finK\_\allowbreak{}awayK} (\texttt{Development.\allowbreak{}AutomorphicArithmetic.\allowbreak{}ThetaCorrespondence.\allowbreak{}ThetaPairLifts.\allowbreak{}ThirdPlaceSeries.\allowbreak{}LeafFixedPointOfFinKAwayKStatement}, lines 12--13)\\
\texttt{GlobalTheory.\allowbreak{}T.\allowbreak{}leafHeightMapBound\_\allowbreak{}of\_\allowbreak{}infK} (\texttt{Development.\allowbreak{}AutomorphicArithmetic.\allowbreak{}ThetaCorrespondence.\allowbreak{}ThetaPairLifts.\allowbreak{}ThirdPlaceSeries.\allowbreak{}IpaTieStatement}, lines 12--13)\\
\texttt{GlobalTheory.\allowbreak{}T.\allowbreak{}leafMinimalMass\_\allowbreak{}of\_\allowbreak{}setAt} (\texttt{Development.\allowbreak{}AutomorphicArithmetic.\allowbreak{}ThetaCorrespondence.\allowbreak{}ThetaPairLifts.\allowbreak{}ThirdPlaceSeries.\allowbreak{}MmTieStatement}, lines 11--12)\par}

{\footnotesize\raggedright\noindent \ref{prop:p8c2-fitted-evaluation}.\quad \texttt{GlobalTheory.\allowbreak{}BigTop.\allowbreak{}Fun.\allowbreak{}oneSq\_\allowbreak{}of\_\allowbreak{}log\_\allowbreak{}derivative\_\allowbreak{}bound\_\allowbreak{}thA} (\texttt{Development.\allowbreak{}AutomorphicArithmetic.\allowbreak{}ThetaCorrespondence.\allowbreak{}ThetaPairLifts.\allowbreak{}ThirdPlaceSeries.\allowbreak{}OneSqTie}, lines 23--79)\\
\texttt{GlobalTheory.\allowbreak{}BigTop.\allowbreak{}Fun.\allowbreak{}den\_\allowbreak{}tie\_\allowbreak{}alt} (\texttt{Development.\allowbreak{}AutomorphicArithmetic.\allowbreak{}ThetaCorrespondence.\allowbreak{}ThetaPairLifts.\allowbreak{}ThirdPlaceSeries.\allowbreak{}DenTie}, lines 14--26)\par}

{\footnotesize\raggedright\noindent \ref{prop:p8c2-fitted-evaluation}.\quad \texttt{GlobalTheory.\allowbreak{}BigTop.\allowbreak{}Fun.\allowbreak{}int\_\allowbreak{}tie\_\allowbreak{}alt} (\texttt{Development.\allowbreak{}AutomorphicArithmetic.\allowbreak{}ThetaCorrespondence.\allowbreak{}ThetaPairLifts.\allowbreak{}ThirdPlaceSeries.\allowbreak{}IntTie}, lines 17--76)\\
\texttt{GlobalTheory.\allowbreak{}BigTop.\allowbreak{}Fun.\allowbreak{}arch\_\allowbreak{}tie\_\allowbreak{}of\_\allowbreak{}logDeriv\_\allowbreak{}bound} (\texttt{Development.\allowbreak{}AutomorphicArithmetic.\allowbreak{}ThetaCorrespondence.\allowbreak{}ThetaPairLifts.\allowbreak{}ThirdPlaceSeries.\allowbreak{}ArchTie}, lines 14--47)\par}

{\footnotesize\raggedright\noindent \ref{prop:p8c2-fitted-evaluation}.\quad \texttt{GlobalTheory.\allowbreak{}BigTop.\allowbreak{}Fun.\allowbreak{}desPol\_\allowbreak{}alt} (\texttt{Development.\allowbreak{}AutomorphicArithmetic.\allowbreak{}ThetaCorrespondence.\allowbreak{}ThetaPairLifts.\allowbreak{}ThirdPlaceSeries.\allowbreak{}DesPol}, lines 17--58)\par}

{\footnotesize\raggedright\noindent \ref{prop:p8c2-fitted-evaluation}.\quad \texttt{GlobalTheory.\allowbreak{}T.\allowbreak{}leafIntegrality\_\allowbreak{}of\_\allowbreak{}logDeriv\_\allowbreak{}bound} (\texttt{Development.\allowbreak{}AutomorphicArithmetic.\allowbreak{}ThetaCorrespondence.\allowbreak{}ThetaPairLifts.\allowbreak{}ThirdPlaceSeries.\allowbreak{}IntTieStatement}, lines 13--14)\\
\texttt{GlobalTheory.\allowbreak{}T.\allowbreak{}arch\_\allowbreak{}tie\_\allowbreak{}of\_\allowbreak{}setAt\_\allowbreak{}logDeriv} (\texttt{Development.\allowbreak{}AutomorphicArithmetic.\allowbreak{}ThetaCorrespondence.\allowbreak{}ThetaPairLifts.\allowbreak{}ThirdPlaceSeries.\allowbreak{}ArchTieStatement}, lines 13--14)\par}

{\footnotesize\raggedright\noindent \ref{B8.22}.\quad \texttt{GlobalTheory.\allowbreak{}BigTop.\allowbreak{}Fun.\allowbreak{}Close.\allowbreak{}fourier\_\allowbreak{}coefficient\_\allowbreak{}fit\_\allowbreak{}of\_\allowbreak{}join\_\allowbreak{}alt} (\texttt{Development.\allowbreak{}AutomorphicArithmetic.\allowbreak{}ThetaCorrespondence.\allowbreak{}ThetaPairLifts.\allowbreak{}ThirdPlaceSeries.\allowbreak{}FourierCoefficientFitOfJoin}, lines 12--17)\\
\texttt{GlobalTheory.\allowbreak{}BigTop.\allowbreak{}Fun.\allowbreak{}ConstantFunctionFits} (\texttt{Development.\allowbreak{}AutomorphicArithmetic.\allowbreak{}ThetaCorrespondence.\allowbreak{}ThetaPairLifts.\allowbreak{}ThirdPlaceSeries.\allowbreak{}LeafInfinitePlaceBound}, lines 52--53)\\
\texttt{Formalism.\allowbreak{}L2.\allowbreak{}Indef.\allowbreak{}IndefData.\allowbreak{}FUN} (\texttt{Formalism.\allowbreak{}LocalToGlobal.\allowbreak{}IndefiniteForms.\allowbreak{}IndefSide}, lines 331--332)\par}

{\footnotesize\raggedright\noindent \ref{B8.22}.\quad \texttt{Formalism.\allowbreak{}L2.\allowbreak{}Indef.\allowbreak{}Tree.\allowbreak{}functional\_\allowbreak{}lower\_\allowbreak{}bound\_\allowbreak{}of\_\allowbreak{}conditions} (\texttt{Formalism.\allowbreak{}LocalToGlobal.\allowbreak{}IndefiniteForms.\allowbreak{}IndefiniteTree.\allowbreak{}FunctionalLowerBoundOfConditionsNumeric}, lines 940--1033)\\
\texttt{Formalism.\allowbreak{}L2.\allowbreak{}Indef.\allowbreak{}Tree.\allowbreak{}size\_\allowbreak{}arith} (\texttt{Formalism.\allowbreak{}LocalToGlobal.\allowbreak{}IndefiniteForms.\allowbreak{}IndefiniteTree.\allowbreak{}FunctionalLowerBoundOfConditionsNumeric}, lines 243--397)\\
\texttt{GlobalTheory.\allowbreak{}T.\allowbreak{}polynomial\_\allowbreak{}descent\_\allowbreak{}estimate'} (\texttt{Development.\allowbreak{}AutomorphicArithmetic.\allowbreak{}ThetaCorrespondence.\allowbreak{}ThetaPairLifts.\allowbreak{}ThirdPlaceSeries.\allowbreak{}DesPolStatement}, lines 11--12)\\
\texttt{GlobalTheory.\allowbreak{}BigTop.\allowbreak{}Fun.\allowbreak{}desPol\_\allowbreak{}alt} (\texttt{Development.\allowbreak{}AutomorphicArithmetic.\allowbreak{}ThetaCorrespondence.\allowbreak{}ThetaPairLifts.\allowbreak{}ThirdPlaceSeries.\allowbreak{}DesPol}, lines 17--58)\par}

{\footnotesize\raggedright\noindent \ref{B8.23}.\quad \texttt{GlobalTheory.\allowbreak{}Fc.\allowbreak{}StageOne.\allowbreak{}rowFun\_\allowbreak{}at\_\allowbreak{}gridA\_\allowbreak{}alt} (\texttt{Development.\allowbreak{}AutomorphicArithmetic.\allowbreak{}ThetaCorrespondence.\allowbreak{}ThetaPairLifts.\allowbreak{}ThirdPlaceSeries.\allowbreak{}RowFunAtGridA}, lines 32--51)\\
\texttt{GlobalTheory.\allowbreak{}Fc.\allowbreak{}S0.\allowbreak{}FunctionalCondition} (\texttt{Development.\allowbreak{}AutomorphicArithmetic.\allowbreak{}ThetaCorrespondence.\allowbreak{}ThetaPairLifts.\allowbreak{}ThirdPlaceSeries.\allowbreak{}GeometricWeakTies}, lines 34--35)\par}

{\footnotesize\raggedright\noindent \ref{B8.23}.\quad \texttt{GlobalTheory.\allowbreak{}Fc.\allowbreak{}Ind.\allowbreak{}evWC\_\allowbreak{}fun\_\allowbreak{}of\_\allowbreak{}fun0Fit\_\allowbreak{}alt} (\texttt{Development.\allowbreak{}AutomorphicArithmetic.\allowbreak{}ThetaCorrespondence.\allowbreak{}ThetaPairLifts.\allowbreak{}ThirdPlaceSeries.\allowbreak{}EvWCFunOfFun0Fit}, lines 14--21)\par}

{\footnotesize\raggedright\noindent \ref{B8.23}.\quad \texttt{UpperTheory.\allowbreak{}Lit.\allowbreak{}nonempty\_\allowbreak{}hyps\_\allowbreak{}definite\_\allowbreak{}side\_\allowbreak{}of\_\allowbreak{}closed\_\allowbreak{}bound} (\texttt{Development.\allowbreak{}AutomorphicArithmetic.\allowbreak{}ThetaCorrespondence.\allowbreak{}ThetaPairLifts.\allowbreak{}ThirdPlaceSeries.\allowbreak{}SiegelEisensteinSecondTermHypotheses}, lines 44--68)\par}

\section*{Chapter~\ref{ch:B8-covolume}: Covolume bounds}

{\footnotesize\raggedright\noindent \ref{eq:p8c3-map-height}.\quad \texttt{Formalism.\allowbreak{}HN.\allowbreak{}ht} (\texttt{Formalism.\allowbreak{}HarderNarasimhan.\allowbreak{}StatementPl}, lines 29--31)\\
\texttt{Formalism.\allowbreak{}HN.\allowbreak{}htMap} (\texttt{Formalism.\allowbreak{}HarderNarasimhan.\allowbreak{}StatementPl}, lines 60--62)\\
\texttt{GlobalTheory.\allowbreak{}DescentFamily.\allowbreak{}htMap\_\allowbreak{}le\_\allowbreak{}of\_\allowbreak{}points} (\texttt{Development.\allowbreak{}AbelianVarietyTheory.\allowbreak{}Proofs.\allowbreak{}DivisorProofs.\allowbreak{}HtMapLeOfPoints}, lines 34--56)\\
\texttt{Formalism.\allowbreak{}DivisionValues.\allowbreak{}lcv} (\texttt{Formalism.\allowbreak{}QuotientMetricDivisors.\allowbreak{}Statement}, lines 34--37)\par}

{\footnotesize\raggedright\noindent \ref{eq:p8c3-period}.\quad \texttt{Formalism.\allowbreak{}L2.\allowbreak{}DIV.\allowbreak{}ZscOfNum} (\texttt{Formalism.\allowbreak{}LocalToGlobal.\allowbreak{}Divisibility.\allowbreak{}ZscOfNum}, lines 59--72)\\
\texttt{Formalism.\allowbreak{}L2.\allowbreak{}DIV.\allowbreak{}ofRecord\_\allowbreak{}logOm} (\texttt{Formalism.\allowbreak{}LocalToGlobal.\allowbreak{}Divisibility.\allowbreak{}NumbersLemmas}, lines 786--789)\\
\texttt{GlobalTheory.\allowbreak{}CovTop.\allowbreak{}unitE} (\texttt{Development.\allowbreak{}AbelianVarietyTheory.\allowbreak{}Definitions.\allowbreak{}Divisors.\allowbreak{}CovolXOfCovol}, lines 22--23)\\
\texttt{GlobalTheory.\allowbreak{}CovTop.\allowbreak{}mainT} (\texttt{Development.\allowbreak{}AbelianVarietyTheory.\allowbreak{}Definitions.\allowbreak{}Divisors.\allowbreak{}CovolXOfCovol}, lines 24--25)\\
\texttt{Formalism.\allowbreak{}L2.\allowbreak{}Indef.\allowbreak{}IndefScale.\allowbreak{}num} (\texttt{Formalism.\allowbreak{}LocalToGlobal.\allowbreak{}IndefiniteForms.\allowbreak{}IndefSide}, lines 224--232)\par}

{\footnotesize\raggedright\noindent \ref{eq:p8c3-period}.\quad \texttt{GlobalTheory.\allowbreak{}BigTop.\allowbreak{}FamHyp} (\texttt{Development.\allowbreak{}AbelianVarietyTheory.\allowbreak{}Definitions.\allowbreak{}Divisors.\allowbreak{}TupleReadBaseTupleEqFamily}, lines 503--507)\\
\texttt{GlobalTheory.\allowbreak{}CovTop.\allowbreak{}RowsTAlt} (\texttt{Development.\allowbreak{}AutomorphicArithmetic.\allowbreak{}ThetaCorrespondence.\allowbreak{}ThetaPairLifts.\allowbreak{}ThirdPlaceSeries.\allowbreak{}LiftLambdaBound}, lines 17--17)\\
\texttt{GlobalTheory.\allowbreak{}CovTop.\allowbreak{}RecTied} (\texttt{Development.\allowbreak{}AbelianVarietyTheory.\allowbreak{}Definitions.\allowbreak{}Divisors.\allowbreak{}CovolXOfCovol}, lines 161--161)\par}

{\footnotesize\raggedright\noindent \ref{def:p8c3-admissible}.\quad \texttt{GlobalTheory.\allowbreak{}BigTop.\allowbreak{}PtAt.\allowbreak{}PointConditions'} (\texttt{Development.\allowbreak{}AutomorphicArithmetic.\allowbreak{}ThetaCorrespondence.\allowbreak{}ThetaPairLifts.\allowbreak{}ThirdPlaceSeries.\allowbreak{}Realised}, lines 63--80)\\
\texttt{GlobalTheory.\allowbreak{}BigTop.\allowbreak{}PtAt.\allowbreak{}CompatibilityConditions} (\texttt{Development.\allowbreak{}AutomorphicArithmetic.\allowbreak{}ThetaCorrespondence.\allowbreak{}ThetaPairLifts.\allowbreak{}ThirdPlaceSeries.\allowbreak{}Realised}, lines 50--62)\\
\texttt{GlobalTheory.\allowbreak{}FinalJoin2.\allowbreak{}FieldSmall} (\texttt{Development.\allowbreak{}AbelianVarietyTheory.\allowbreak{}Definitions.\allowbreak{}HeightFunctions.\allowbreak{}NormComparisonConstants}, lines 32--36)\\
\texttt{Formalism.\allowbreak{}MainNotions.\allowbreak{}WeightAlgebra.\allowbreak{}PO} (\texttt{Formalism.\allowbreak{}SharpRestriction.\allowbreak{}Divisibility.\allowbreak{}WeightAlgebra}, lines 271--272)\\
\texttt{GlobalTheory.\allowbreak{}GLq.\allowbreak{}ESerTie} (\texttt{Development.\allowbreak{}AbelianVarietyTheory.\allowbreak{}Definitions.\allowbreak{}NonGaloisDescent.\allowbreak{}LabelDatum}, lines 114--117)\\
\texttt{GlobalTheory.\allowbreak{}GLq.\allowbreak{}eSerC} (\texttt{Development.\allowbreak{}AbelianVarietyTheory.\allowbreak{}Definitions.\allowbreak{}NonGaloisDescent.\allowbreak{}LabelDatum}, lines 102--104)\par}

{\footnotesize\raggedright\noindent \ref{B8.24}.\quad \texttt{GlobalTheory.\allowbreak{}CovTop.\allowbreak{}CovolumeBoundOfX} (\texttt{Development.\allowbreak{}AutomorphicArithmetic.\allowbreak{}ThetaCorrespondence.\allowbreak{}ThetaPairLifts.\allowbreak{}ThirdPlaceSeries.\allowbreak{}LiftLambdaBound}, lines 60--61)\\
\texttt{GlobalTheory.\allowbreak{}CovTop.\allowbreak{}EventuallyAtAllPoints} (\texttt{Development.\allowbreak{}AutomorphicArithmetic.\allowbreak{}ThetaCorrespondence.\allowbreak{}ThetaPairLifts.\allowbreak{}ThirdPlaceSeries.\allowbreak{}LiftLambdaBound}, lines 18--19)\\
\texttt{Formalism.\allowbreak{}L2.\allowbreak{}Indef.\allowbreak{}IndefData.\allowbreak{}L3} (\texttt{Formalism.\allowbreak{}LocalToGlobal.\allowbreak{}IndefiniteForms.\allowbreak{}IndefSide}, lines 315--317)\\
\texttt{Formalism.\allowbreak{}MainNotions.\allowbreak{}FamP.\allowbreak{}Large} (\texttt{Formalism.\allowbreak{}SharpRestriction.\allowbreak{}ParamsHeightRankConstant}, lines 138--140)\par}

{\footnotesize\raggedright\noindent \ref{B8.24}.\quad \texttt{GlobalTheory.\allowbreak{}CovTop.\allowbreak{}LiftsCondition} (\texttt{Development.\allowbreak{}AutomorphicArithmetic.\allowbreak{}ThetaCorrespondence.\allowbreak{}ThetaPairLifts.\allowbreak{}ThirdPlaceSeries.\allowbreak{}LiftLambdaBound}, lines 26--28)\\
\texttt{GlobalTheory.\allowbreak{}CovTop.\allowbreak{}CoversCondition} (\texttt{Development.\allowbreak{}AutomorphicArithmetic.\allowbreak{}ThetaCorrespondence.\allowbreak{}ThetaPairLifts.\allowbreak{}ThirdPlaceSeries.\allowbreak{}LiftLambdaBound}, lines 29--32)\\
\texttt{Formalism.\allowbreak{}MainNotions.\allowbreak{}Closed} (\texttt{Formalism.\allowbreak{}SharpRestriction.\allowbreak{}ParamsHeightRankConstant}, lines 159--162)\\
\texttt{GlobalTheory.\allowbreak{}CovTop.\allowbreak{}EventuallyAtAllPoints.\allowbreak{}and} (\texttt{Development.\allowbreak{}AutomorphicArithmetic.\allowbreak{}ThetaCorrespondence.\allowbreak{}ThetaPairLifts.\allowbreak{}ThirdPlaceSeries.\allowbreak{}LiftLambdaBound}, lines 20--25)\par}

{\footnotesize\raggedright\noindent \ref{B8.24}.\quad \texttt{GlobalTheory.\allowbreak{}CovTop.\allowbreak{}uM} (\texttt{Development.\allowbreak{}AbelianVarietyTheory.\allowbreak{}Definitions.\allowbreak{}Divisors.\allowbreak{}CovolXOfCovol}, lines 28--36)\\
\texttt{GlobalTheory.\allowbreak{}CovTop.\allowbreak{}cZeroN} (\texttt{Development.\allowbreak{}AbelianVarietyTheory.\allowbreak{}Definitions.\allowbreak{}Divisors.\allowbreak{}NormNegCmConstantEqAtAllEmbeddings}, lines 83--84)\\
\texttt{GlobalTheory.\allowbreak{}CovTop.\allowbreak{}cZero} (\texttt{Development.\allowbreak{}AbelianVarietyTheory.\allowbreak{}Definitions.\allowbreak{}Divisors.\allowbreak{}NormNegCmConstantEqAtAllEmbeddings}, lines 86--86)\par}

{\footnotesize\raggedright\noindent \ref{def:p8c3-section-laws}.\quad \texttt{GlobalTheory.\allowbreak{}AllWt.\allowbreak{}GoodW} (\texttt{Development.\allowbreak{}AbelianVarietyTheory.\allowbreak{}Definitions.\allowbreak{}Foundations.\allowbreak{}IsProbabilityMeasureMuPart}, lines 806--808)\\
\texttt{GlobalTheory.\allowbreak{}AllWt.\allowbreak{}GoodOn} (\texttt{Development.\allowbreak{}AbelianVarietyTheory.\allowbreak{}Definitions.\allowbreak{}Foundations.\allowbreak{}IsProbabilityMeasureMuPart}, lines 809--811)\\
\texttt{GlobalTheory.\allowbreak{}CovTop.\allowbreak{}DeepGoodAt} (\texttt{Development.\allowbreak{}AbelianVarietyTheory.\allowbreak{}Definitions.\allowbreak{}Divisors.\allowbreak{}CovolXOfCovol}, lines 111--115)\\
\texttt{GlobalTheory.\allowbreak{}FlatLine.\allowbreak{}FlatExistsAt} (\texttt{Development.\allowbreak{}AbelianVarietyTheory.\allowbreak{}Definitions.\allowbreak{}Foundations.\allowbreak{}ExistsFlatSecRankZero}, lines 746--748)\\
\texttt{GlobalTheory.\allowbreak{}FlatLine.\allowbreak{}FlatDegree} (\texttt{Development.\allowbreak{}AbelianVarietyTheory.\allowbreak{}Definitions.\allowbreak{}Foundations.\allowbreak{}ExistsFlatSecRankZero}, lines 751--755)\par}

{\footnotesize\raggedright\noindent \ref{B8.24a}.\quad \texttt{GlobalTheory.\allowbreak{}FlatLine.\allowbreak{}FlatDegree.\allowbreak{}c1\_\allowbreak{}eq} (\texttt{Development.\allowbreak{}AbelianVarietyTheory.\allowbreak{}Definitions.\allowbreak{}Foundations.\allowbreak{}ExistsFlatSecRankZero}, lines 788--790)\\
\texttt{GlobalTheory.\allowbreak{}FlatCore.\allowbreak{}flatDegreeSt\_\allowbreak{}of\_\allowbreak{}prismaticDieudonne} (\texttt{Development.\allowbreak{}AbelianVarietyTheory.\allowbreak{}Proofs.\allowbreak{}DivisorProofs.\allowbreak{}FlatDegreeFromPrismatic}, lines 23--35)\\
\texttt{GlobalTheory.\allowbreak{}FlatSplit.\allowbreak{}twice\_\allowbreak{}twist\_\allowbreak{}ge} (\texttt{Development.\allowbreak{}AbelianVarietyTheory.\allowbreak{}Proofs.\allowbreak{}DivisorProofs.\allowbreak{}FlatBoundOfAnchor}, lines 76--102)\\
\texttt{GlobalTheory.\allowbreak{}CovTop.\allowbreak{}Period.\allowbreak{}period\_\allowbreak{}eq\_\allowbreak{}kronecker} (\texttt{Development.\allowbreak{}AbelianVarietyTheory.\allowbreak{}Proofs.\allowbreak{}DivisorProofs.\allowbreak{}LogPeriodEqKroneckerCharacterLogPeriodOfHost}, lines 13--17)\\
\texttt{GlobalTheory.\allowbreak{}CovTop.\allowbreak{}AnchorRd.\allowbreak{}anchorAt\_\allowbreak{}of} (\texttt{Development.\allowbreak{}AbelianVarietyTheory.\allowbreak{}Proofs.\allowbreak{}DivisorProofs.\allowbreak{}AnchorBoundStatement}, lines 77--209)\\
\texttt{GlobalTheory.\allowbreak{}FlatSplit.\allowbreak{}flatAt\_\allowbreak{}of} (\texttt{Development.\allowbreak{}AbelianVarietyTheory.\allowbreak{}Proofs.\allowbreak{}DivisorProofs.\allowbreak{}FlatBoundOfAnchor}, lines 125--187)\par}

{\footnotesize\raggedright\noindent \ref{B8.25l}.\quad \texttt{UpperTheory.\allowbreak{}LiteratureTheorems.\allowbreak{}toroidal\_\allowbreak{}compactification\_\allowbreak{}of\_\allowbreak{}literature} (\texttt{Development.\allowbreak{}AutomorphicArithmetic.\allowbreak{}ThetaCorrespondence.\allowbreak{}ThetaPairLifts.\allowbreak{}ThirdPlaceSeries.\allowbreak{}EveryCharacterOfLiteratureOfVWin}, lines 521--522)\par}

{\footnotesize\raggedright\noindent \ref{B8.25l}.\quad \texttt{UpperTheory.\allowbreak{}LiteratureTheorems.\allowbreak{}vol\_\allowbreak{}ge\_\allowbreak{}of\_\allowbreak{}literature} (\texttt{Development.\allowbreak{}AutomorphicArithmetic.\allowbreak{}ThetaCorrespondence.\allowbreak{}ThetaPairLifts.\allowbreak{}ThirdPlaceSeries.\allowbreak{}EveryCharacterOfLiteratureOfVWin}, lines 531--532)\par}

{\footnotesize\raggedright\noindent \ref{B8.25l}.\quad \texttt{UpperTheory.\allowbreak{}LiteratureTheorems.\allowbreak{}vol\_\allowbreak{}nef\_\allowbreak{}of\_\allowbreak{}literature} (\texttt{Development.\allowbreak{}AutomorphicArithmetic.\allowbreak{}ThetaCorrespondence.\allowbreak{}ThetaPairLifts.\allowbreak{}ThirdPlaceSeries.\allowbreak{}EveryCharacterOfLiteratureOfVWin}, lines 533--534)\par}

{\footnotesize\raggedright\noindent \ref{B8.25l}.\quad \texttt{UpperTheory.\allowbreak{}LiteratureTheorems.\allowbreak{}ar\_\allowbreak{}calc\_\allowbreak{}of\_\allowbreak{}literature} (\texttt{Development.\allowbreak{}AutomorphicArithmetic.\allowbreak{}ThetaCorrespondence.\allowbreak{}ThetaPairLifts.\allowbreak{}ThirdPlaceSeries.\allowbreak{}EveryCharacterOfLiteratureOfVWin}, lines 249--250)\par}

{\footnotesize\raggedright\noindent \ref{B8.25l}.\quad \texttt{UpperTheory.\allowbreak{}LiteratureTheorems.\allowbreak{}abelian\_\allowbreak{}de\_\allowbreak{}rham\_\allowbreak{}of\_\allowbreak{}literature} (\texttt{Development.\allowbreak{}AutomorphicArithmetic.\allowbreak{}ThetaCorrespondence.\allowbreak{}ThetaPairLifts.\allowbreak{}ThirdPlaceSeries.\allowbreak{}EveryCharacterOfLiteratureOfVWin}, lines 239--240)\par}

{\footnotesize\raggedright\noindent \ref{B8.25l}.\quad \texttt{UpperTheory.\allowbreak{}LiteratureTheorems.\allowbreak{}ray\_\allowbreak{}class\_\allowbreak{}field\_\allowbreak{}st\_\allowbreak{}of\_\allowbreak{}literature} (\texttt{Development.\allowbreak{}AutomorphicArithmetic.\allowbreak{}ThetaCorrespondence.\allowbreak{}ThetaPairLifts.\allowbreak{}ThirdPlaceSeries.\allowbreak{}EveryCharacterOfLiteratureOfVWin}, lines 489--490)\par}

{\footnotesize\raggedright\noindent \ref{B8.25l}.\quad \texttt{UpperTheory.\allowbreak{}LiteratureTheorems.\allowbreak{}semistable\_\allowbreak{}reduction\_\allowbreak{}st\_\allowbreak{}of\_\allowbreak{}literature} (\texttt{Development.\allowbreak{}AutomorphicArithmetic.\allowbreak{}ThetaCorrespondence.\allowbreak{}ThetaPairLifts.\allowbreak{}ThirdPlaceSeries.\allowbreak{}EveryCharacterOfLiteratureOfVWin}, lines 501--502)\par}

{\footnotesize\raggedright\noindent \ref{B8.25l}.\quad \texttt{UpperTheory.\allowbreak{}LiteratureTheorems.\allowbreak{}strong\_\allowbreak{}approx\_\allowbreak{}su\_\allowbreak{}of\_\allowbreak{}literature} (\texttt{Development.\allowbreak{}AutomorphicArithmetic.\allowbreak{}ThetaCorrespondence.\allowbreak{}ThetaPairLifts.\allowbreak{}ThirdPlaceSeries.\allowbreak{}EveryCharacterOfLiteratureOfVWin}, lines 517--518)\par}

{\footnotesize\raggedright\noindent \ref{B8.25l}.\quad \texttt{UpperTheory.\allowbreak{}LiteratureTheorems.\allowbreak{}green\_\allowbreak{}series\_\allowbreak{}st\_\allowbreak{}of\_\allowbreak{}literature} (\texttt{Development.\allowbreak{}AutomorphicArithmetic.\allowbreak{}ThetaCorrespondence.\allowbreak{}ThetaPairLifts.\allowbreak{}ThirdPlaceSeries.\allowbreak{}EveryCharacterOfLiteratureOfVWin}, lines 401--402)\par}

{\footnotesize\raggedright\noindent \ref{B8.25l}.\quad \texttt{UpperTheory.\allowbreak{}LiteratureTheorems.\allowbreak{}good\_\allowbreak{}chern\_\allowbreak{}weil\_\allowbreak{}of\_\allowbreak{}literature} (\texttt{Development.\allowbreak{}AutomorphicArithmetic.\allowbreak{}ThetaCorrespondence.\allowbreak{}ThetaPairLifts.\allowbreak{}ThirdPlaceSeries.\allowbreak{}EveryCharacterOfLiteratureOfVWin}, lines 395--396)\par}

{\footnotesize\raggedright\noindent \ref{B8.25l}.\quad \texttt{UpperTheory.\allowbreak{}LiteratureTheorems.\allowbreak{}landherr\_\allowbreak{}st\_\allowbreak{}of\_\allowbreak{}literature} (\texttt{Development.\allowbreak{}AutomorphicArithmetic.\allowbreak{}ThetaCorrespondence.\allowbreak{}ThetaPairLifts.\allowbreak{}ThirdPlaceSeries.\allowbreak{}EveryCharacterOfLiteratureOfVWin}, lines 453--454)\par}

{\footnotesize\raggedright\noindent \ref{B8.25l}.\quad \texttt{UpperTheory.\allowbreak{}LiteratureTheorems.\allowbreak{}hasse\_\allowbreak{}norm\_\allowbreak{}st\_\allowbreak{}of\_\allowbreak{}literature} (\texttt{Development.\allowbreak{}AutomorphicArithmetic.\allowbreak{}ThetaCorrespondence.\allowbreak{}ThetaPairLifts.\allowbreak{}ThirdPlaceSeries.\allowbreak{}EveryCharacterOfLiteratureOfVWin}, lines 417--418)\par}

{\footnotesize\raggedright\noindent \ref{B8.25l}.\quad \texttt{UpperTheory.\allowbreak{}LiteratureTheorems.\allowbreak{}semi\_\allowbreak{}abelian\_\allowbreak{}hom\_\allowbreak{}extension\_\allowbreak{}of\_\allowbreak{}literature} (\texttt{Development.\allowbreak{}AutomorphicArithmetic.\allowbreak{}ThetaCorrespondence.\allowbreak{}ThetaPairLifts.\allowbreak{}ThirdPlaceSeries.\allowbreak{}EveryCharacterOfLiteratureOfVWin}, lines 497--498)\par}

{\footnotesize\raggedright\noindent \ref{B8.25l}.\quad \texttt{UpperTheory.\allowbreak{}LiteratureTheorems.\allowbreak{}induced\_\allowbreak{}simple\_\allowbreak{}lowest\_\allowbreak{}alcove\_\allowbreak{}of\_\allowbreak{}literature} (\texttt{Development.\allowbreak{}AutomorphicArithmetic.\allowbreak{}ThetaCorrespondence.\allowbreak{}ThetaPairLifts.\allowbreak{}ThirdPlaceSeries.\allowbreak{}EveryCharacterOfLiteratureOfVWin}, lines 441--442)\par}

{\footnotesize\raggedright\noindent \ref{B8.25l}.\quad \texttt{UpperTheory.\allowbreak{}LiteratureTheorems.\allowbreak{}principal\_\allowbreak{}polarisation\_\allowbreak{}one\_\allowbreak{}section\_\allowbreak{}of\_\allowbreak{}literature} (\texttt{Development.\allowbreak{}AutomorphicArithmetic.\allowbreak{}ThetaCorrespondence.\allowbreak{}ThetaPairLifts.\allowbreak{}ThirdPlaceSeries.\allowbreak{}EveryCharacterOfLiteratureOfVWin}, lines 409--410)\par}

{\footnotesize\raggedright\noindent \ref{B8.25l}.\quad \texttt{UpperTheory.\allowbreak{}LiteratureTheorems.\allowbreak{}ab\_\allowbreak{}mul\_\allowbreak{}n\_\allowbreak{}st\_\allowbreak{}of\_\allowbreak{}literature} (\texttt{Development.\allowbreak{}AutomorphicArithmetic.\allowbreak{}ThetaCorrespondence.\allowbreak{}ThetaPairLifts.\allowbreak{}ThirdPlaceSeries.\allowbreak{}EveryCharacterOfLiteratureOfVWin}, lines 237--238)\par}

{\footnotesize\raggedright\noindent \ref{B8.25l}.\quad \texttt{UpperTheory.\allowbreak{}LiteratureTheorems.\allowbreak{}grothendieck\_\allowbreak{}messing\_\allowbreak{}of\_\allowbreak{}literature} (\texttt{Development.\allowbreak{}AutomorphicArithmetic.\allowbreak{}ThetaCorrespondence.\allowbreak{}ThetaPairLifts.\allowbreak{}ThirdPlaceSeries.\allowbreak{}EveryCharacterOfLiteratureOfVWin}, lines 407--408)\par}

{\footnotesize\raggedright\noindent \ref{B8.25l}.\quad \texttt{UpperTheory.\allowbreak{}LiteratureTheorems.\allowbreak{}prismatic\_\allowbreak{}dieudonne\_\allowbreak{}hodge\_\allowbreak{}of\_\allowbreak{}literature} (\texttt{Development.\allowbreak{}AutomorphicArithmetic.\allowbreak{}ThetaCorrespondence.\allowbreak{}ThetaPairLifts.\allowbreak{}ThirdPlaceSeries.\allowbreak{}EveryCharacterOfLiteratureOfVWin}, lines 475--476)\par}

{\footnotesize\raggedright\noindent \ref{B8.25l}.\quad \texttt{UpperTheory.\allowbreak{}LiteratureTheorems.\allowbreak{}kawamata\_\allowbreak{}viehweg\_\allowbreak{}of\_\allowbreak{}literature} (\texttt{Development.\allowbreak{}AutomorphicArithmetic.\allowbreak{}ThetaCorrespondence.\allowbreak{}ThetaPairLifts.\allowbreak{}ThirdPlaceSeries.\allowbreak{}EveryCharacterOfLiteratureOfVWin}, lines 447--448)\par}

{\footnotesize\raggedright\noindent \ref{B8.25l}.\quad \texttt{UpperTheory.\allowbreak{}LiteratureTheorems.\allowbreak{}nef\_\allowbreak{}numerics\_\allowbreak{}of\_\allowbreak{}literature} (\texttt{Development.\allowbreak{}AutomorphicArithmetic.\allowbreak{}ThetaCorrespondence.\allowbreak{}ThetaPairLifts.\allowbreak{}ThirdPlaceSeries.\allowbreak{}EveryCharacterOfLiteratureOfVWin}, lines 469--470)\par}

{\footnotesize\raggedright\noindent \ref{B8.25l}.\quad \texttt{UpperTheory.\allowbreak{}LiteratureTheorems.\allowbreak{}l2\_\allowbreak{}extension\_\allowbreak{}log\_\allowbreak{}canonical\_\allowbreak{}of\_\allowbreak{}literature} (\texttt{Development.\allowbreak{}AutomorphicArithmetic.\allowbreak{}ThetaCorrespondence.\allowbreak{}ThetaPairLifts.\allowbreak{}ThirdPlaceSeries.\allowbreak{}EveryCharacterOfLiteratureOfVWin}, lines 451--452)\par}

{\footnotesize\raggedright\noindent \ref{B8.25l}.\quad \texttt{UpperTheory.\allowbreak{}LiteratureTheorems.\allowbreak{}arith\_\allowbreak{}hilbert\_\allowbreak{}samuel\_\allowbreak{}of\_\allowbreak{}literature} (\texttt{Development.\allowbreak{}AutomorphicArithmetic.\allowbreak{}ThetaCorrespondence.\allowbreak{}ThetaPairLifts.\allowbreak{}ThirdPlaceSeries.\allowbreak{}EveryCharacterOfLiteratureOfVWin}, lines 255--256)\par}

{\footnotesize\raggedright\noindent \ref{B8.25l}.\quad \texttt{UpperTheory.\allowbreak{}LiteratureTheorems.\allowbreak{}chow\_\allowbreak{}form\_\allowbreak{}exists\_\allowbreak{}of\_\allowbreak{}literature} (\texttt{Development.\allowbreak{}AutomorphicArithmetic.\allowbreak{}ThetaCorrespondence.\allowbreak{}ThetaPairLifts.\allowbreak{}ThirdPlaceSeries.\allowbreak{}EveryCharacterOfLiteratureOfVWin}, lines 325--326)\par}

{\footnotesize\raggedright\noindent \ref{B8.25l}.\quad \texttt{UpperTheory.\allowbreak{}LiteratureTheorems.\allowbreak{}cycle\_\allowbreak{}heights\_\allowbreak{}of\_\allowbreak{}literature} (\texttt{Development.\allowbreak{}AutomorphicArithmetic.\allowbreak{}ThetaCorrespondence.\allowbreak{}ThetaPairLifts.\allowbreak{}ThirdPlaceSeries.\allowbreak{}EveryCharacterOfLiteratureOfVWin}, lines 347--348)\par}

{\footnotesize\raggedright\noindent \ref{B8.25l}.\quad \texttt{UpperTheory.\allowbreak{}LiteratureTheorems.\allowbreak{}serre\_\allowbreak{}holomorphic\_\allowbreak{}sections\_\allowbreak{}of\_\allowbreak{}literature} (\texttt{Development.\allowbreak{}AutomorphicArithmetic.\allowbreak{}ThetaCorrespondence.\allowbreak{}ThetaPairLifts.\allowbreak{}ThirdPlaceSeries.\allowbreak{}EveryCharacterOfLiteratureOfVWin}, lines 503--504)\par}

{\footnotesize\raggedright\noindent \ref{B8.25l}.\quad \texttt{UpperTheory.\allowbreak{}LiteratureTheorems.\allowbreak{}nadel\_\allowbreak{}vanishing\_\allowbreak{}and\_\allowbreak{}non\_\allowbreak{}klt\_\allowbreak{}divisor\_\allowbreak{}of\_\allowbreak{}literature} (\texttt{Development.\allowbreak{}AutomorphicArithmetic.\allowbreak{}ThetaCorrespondence.\allowbreak{}ThetaPairLifts.\allowbreak{}ThirdPlaceSeries.\allowbreak{}EveryCharacterOfLiteratureOfVWin}, lines 467--468)\par}

{\footnotesize\raggedright\noindent \ref{B8.25l}.\quad \texttt{UpperTheory.\allowbreak{}LiteratureTheorems.\allowbreak{}principalisation\_\allowbreak{}of\_\allowbreak{}literature} (\texttt{Development.\allowbreak{}AutomorphicArithmetic.\allowbreak{}ThetaCorrespondence.\allowbreak{}ThetaPairLifts.\allowbreak{}ThirdPlaceSeries.\allowbreak{}EveryCharacterOfLiteratureOfVWin}, lines 253--254)\par}

{\footnotesize\raggedright\noindent \ref{B8.25l}.\quad \texttt{UpperTheory.\allowbreak{}LiteratureTheorems.\allowbreak{}formal\_\allowbreak{}algebraisation\_\allowbreak{}of\_\allowbreak{}literature} (\texttt{Development.\allowbreak{}AutomorphicArithmetic.\allowbreak{}ThetaCorrespondence.\allowbreak{}ThetaPairLifts.\allowbreak{}ThirdPlaceSeries.\allowbreak{}EveryCharacterOfLiteratureOfVWin}, lines 381--382)\par}

{\footnotesize\raggedright\noindent \ref{B8.25l}.\quad \texttt{GlobalTheory.\allowbreak{}CovTop.\allowbreak{}PublishedHodgeBundleTheorems} (\texttt{Development.\allowbreak{}AutomorphicArithmetic.\allowbreak{}ThetaCorrespondence.\allowbreak{}ThetaPairLifts.\allowbreak{}HodgeBundleWeights.\allowbreak{}ToKnownTheoremsOf}, lines 7--35)\\
\texttt{UpperTheory.\allowbreak{}LiteratureTheorems.\allowbreak{}published\_\allowbreak{}results\_\allowbreak{}arithmetic} (\texttt{Development.\allowbreak{}AutomorphicArithmetic.\allowbreak{}ThetaCorrespondence.\allowbreak{}ThetaPairLifts.\allowbreak{}ThirdPlaceSeries.\allowbreak{}EveryCharacterOfLiteratureOfVWin}, lines 781--809)\par}

{\footnotesize\raggedright\noindent \ref{B8.24b}.\quad \texttt{GlobalTheory.\allowbreak{}FlatCore.\allowbreak{}archEx\_\allowbreak{}always} (\texttt{Development.\allowbreak{}AbelianVarietyTheory.\allowbreak{}Proofs.\allowbreak{}DivisorProofs.\allowbreak{}ArchimedeanExistenceEverywhere}, lines 38--105)\\
\texttt{GlobalTheory.\allowbreak{}CompanionTheory.\allowbreak{}Ar.\allowbreak{}Flat.\allowbreak{}A0.\allowbreak{}flat\_\allowbreak{}zero\_\allowbreak{}of\_\allowbreak{}resolution\_\allowbreak{}gaga} (\texttt{Development.\allowbreak{}AutomorphicSpectrum.\allowbreak{}Proofs.\allowbreak{}GeometryProofs.\allowbreak{}FlatZeroFromResolutionAndGaga}, lines 39--56)\\
\texttt{GlobalTheory.\allowbreak{}CompanionTheory.\allowbreak{}Ar.\allowbreak{}Flat.\allowbreak{}FlatFirstTheorem.\allowbreak{}flat\_\allowbreak{}first\_\allowbreak{}of\_\allowbreak{}resolution\_\allowbreak{}gaga} (\texttt{Development.\allowbreak{}AutomorphicSpectrum.\allowbreak{}Proofs.\allowbreak{}GeometryProofs.\allowbreak{}FlatFirstFromResolutionAndGaga}, lines 103--122)\par}

{\footnotesize\raggedright\noindent \ref{B8.24c}.\quad \texttt{GlobalTheory.\allowbreak{}CompanionTheory.\allowbreak{}Ar.\allowbreak{}Flat.\allowbreak{}det\_\allowbreak{}adapted\_\allowbreak{}eq\_\allowbreak{}unit\_\allowbreak{}mul\_\allowbreak{}const} (\texttt{Development.\allowbreak{}AutomorphicSpectrum.\allowbreak{}Proofs.\allowbreak{}AlgebraProofs.\allowbreak{}CmTypeValuationEmbedding}, lines 301--303)\\
\texttt{GlobalTheory.\allowbreak{}CompanionTheory.\allowbreak{}Ar.\allowbreak{}Flat.\allowbreak{}CPoint.\allowbreak{}value\_\allowbreak{}at\_\allowbreak{}point} (\texttt{Development.\allowbreak{}AutomorphicSpectrum.\allowbreak{}Proofs.\allowbreak{}GeometryProofs.\allowbreak{}EigenLinePAdicFrame}, lines 72--339)\par}

{\footnotesize\raggedright\noindent \ref{B8.24dc}.\quad \texttt{GlobalTheory.\allowbreak{}CompanionTheory.\allowbreak{}Ar.\allowbreak{}Flat.\allowbreak{}hodgeLat\_\allowbreak{}model} (\texttt{Development.\allowbreak{}AutomorphicSpectrum.\allowbreak{}Proofs.\allowbreak{}GeometryProofs.\allowbreak{}EigenLinePAdicBound}, lines 170--256)\\
\texttt{GlobalTheory.\allowbreak{}CompanionTheory.\allowbreak{}Ar.\allowbreak{}Flat.\allowbreak{}GermSys} (\texttt{Development.\allowbreak{}AutomorphicSpectrum.\allowbreak{}Definitions.\allowbreak{}GlobalGeometry.\allowbreak{}GermSys}, lines 19--46)\\
\texttt{ReciprocityTheory.\allowbreak{}ChartRestriction.\allowbreak{}G.\allowbreak{}Geo.\allowbreak{}Tor.\allowbreak{}SAB} (\texttt{ReciprocityTheory.\allowbreak{}ReductiveGroups.\allowbreak{}ArithmeticGeometry.\allowbreak{}BirationalGeometry.\allowbreak{}ToroidalCompactification.\allowbreak{}TwoWeightEqOneAddRoot}, lines 72--83)\\
\texttt{ReciprocityTheory.\allowbreak{}ChartRestriction.\allowbreak{}G.\allowbreak{}Geo.\allowbreak{}Tor.\allowbreak{}Germ.\allowbreak{}FB} (\texttt{ReciprocityTheory.\allowbreak{}ReductiveGroups.\allowbreak{}ArithmeticGeometry.\allowbreak{}BirationalGeometry.\allowbreak{}ToroidalCompactification.\allowbreak{}CocycleRightZeroThd}, lines 86--159)\\
\texttt{GlobalTheory.\allowbreak{}CompanionTheory.\allowbreak{}Ar.\allowbreak{}Conj.\allowbreak{}GeomPF} (\texttt{Development.\allowbreak{}AutomorphicSpectrum.\allowbreak{}Definitions.\allowbreak{}GlobalGeometry.\allowbreak{}ToFirst}, lines 36--40)\\
\texttt{GlobalTheory.\allowbreak{}CompanionTheory.\allowbreak{}Ar.\allowbreak{}Flat.\allowbreak{}HodgeLat} (\texttt{Development.\allowbreak{}AutomorphicSpectrum.\allowbreak{}Definitions.\allowbreak{}GlobalGeometry.\allowbreak{}LinearIndependentEigProj}, lines 21--33)\\
\texttt{GlobalTheory.\allowbreak{}CompanionTheory.\allowbreak{}Ar.\allowbreak{}Conj.\allowbreak{}transH} (\texttt{Development.\allowbreak{}AutomorphicSpectrum.\allowbreak{}Definitions.\allowbreak{}CommutativeAlgebra.\allowbreak{}GermComp1ZeroPointLeft}, lines 99--103)\\
\texttt{GlobalTheory.\allowbreak{}CompanionTheory.\allowbreak{}Ar.\allowbreak{}Flat.\allowbreak{}DegZ.\allowbreak{}vmap\_\allowbreak{}rhoH} (\texttt{Development.\allowbreak{}AutomorphicSpectrum.\allowbreak{}Proofs.\allowbreak{}GeometryProofs.\allowbreak{}ExistsSubMemEigW}, lines 87--102)\par}

{\footnotesize\raggedright\noindent \ref{B8.24dd}.\quad \texttt{GlobalTheory.\allowbreak{}CompanionTheory.\allowbreak{}Ar.\allowbreak{}Flat.\allowbreak{}DegZ.\allowbreak{}GermFam.\allowbreak{}germ\_\allowbreak{}family} (\texttt{Development.\allowbreak{}AutomorphicSpectrum.\allowbreak{}Proofs.\allowbreak{}GeometryProofs.\allowbreak{}FibreCurveGermPowers}, lines 215--271)\\
\texttt{GlobalTheory.\allowbreak{}CompanionTheory.\allowbreak{}Ar.\allowbreak{}Flat.\allowbreak{}DegProd.\allowbreak{}curve\_\allowbreak{}deg} (\texttt{Development.\allowbreak{}AutomorphicSpectrum.\allowbreak{}Proofs.\allowbreak{}GeometryProofs.\allowbreak{}ValuationBoundedFrame}, lines 326--364)\\
\texttt{GlobalTheory.\allowbreak{}CompanionTheory.\allowbreak{}Ar.\allowbreak{}Flat.\allowbreak{}germ\_\allowbreak{}degZero} (\texttt{Development.\allowbreak{}AutomorphicSpectrum.\allowbreak{}Proofs.\allowbreak{}GeometryProofs.\allowbreak{}HodgeLatticeDvrDegree}, lines 117--238)\\
\texttt{GlobalTheory.\allowbreak{}CompanionTheory.\allowbreak{}Ar.\allowbreak{}Flat.\allowbreak{}exists\_\allowbreak{}trans} (\texttt{Development.\allowbreak{}AutomorphicSpectrum.\allowbreak{}Proofs.\allowbreak{}GeometryProofs.\allowbreak{}HodgeLatticeDvrDegree}, lines 72--110)\\
\texttt{GlobalTheory.\allowbreak{}CompanionTheory.\allowbreak{}Ar.\allowbreak{}Flat.\allowbreak{}DegZ.\allowbreak{}Frame.\allowbreak{}exists\_\allowbreak{}basis\_\allowbreak{}frame} (\texttt{Development.\allowbreak{}AutomorphicSpectrum.\allowbreak{}Proofs.\allowbreak{}GeometryProofs.\allowbreak{}HodgeLatticeGermEigen}, lines 130--149)\\
\texttt{GlobalTheory.\allowbreak{}CompanionTheory.\allowbreak{}Ar.\allowbreak{}Flat.\allowbreak{}DegZ.\allowbreak{}charpoly\_\allowbreak{}lmul\_\allowbreak{}splits\_\allowbreak{}of\_\allowbreak{}conjugates} (\texttt{Development.\allowbreak{}AutomorphicSpectrum.\allowbreak{}Proofs.\allowbreak{}GeometryProofs.\allowbreak{}CharpolySplitsOfEmbeddings}, lines 34--64)\\
\texttt{GlobalTheory.\allowbreak{}CompanionTheory.\allowbreak{}Ar.\allowbreak{}Flat.\allowbreak{}DegCurve.\allowbreak{}step} (\texttt{Development.\allowbreak{}AutomorphicSpectrum.\allowbreak{}Proofs.\allowbreak{}AlgebraProofs.\allowbreak{}SplitHermitianGeneratorsExist}, lines 242--251)\\
\texttt{GlobalTheory.\allowbreak{}CompanionTheory.\allowbreak{}Ar.\allowbreak{}Flat.\allowbreak{}DegCurve.\allowbreak{}deg\_\allowbreak{}count} (\texttt{Development.\allowbreak{}AutomorphicSpectrum.\allowbreak{}Proofs.\allowbreak{}AlgebraProofs.\allowbreak{}SplitHermitianGeneratorsExist}, lines 319--352)\par}

{\footnotesize\raggedright\noindent \ref{B8.24d}.\quad \texttt{GlobalTheory.\allowbreak{}CompanionTheory.\allowbreak{}Ar.\allowbreak{}Flat.\allowbreak{}germ\_\allowbreak{}degZero} (\texttt{Development.\allowbreak{}AutomorphicSpectrum.\allowbreak{}Proofs.\allowbreak{}GeometryProofs.\allowbreak{}HodgeLatticeDvrDegree}, lines 117--238)\\
\texttt{GlobalTheory.\allowbreak{}CompanionTheory.\allowbreak{}Ar.\allowbreak{}Flat.\allowbreak{}COrder.\allowbreak{}order\_\allowbreak{}on\_\allowbreak{}component} (\texttt{Development.\allowbreak{}AutomorphicSpectrum.\allowbreak{}Proofs.\allowbreak{}GeometryProofs.\allowbreak{}HodgeLatticeValuationAgree}, lines 12--44)\\
\texttt{GlobalTheory.\allowbreak{}CompanionTheory.\allowbreak{}Ar.\allowbreak{}Flat.\allowbreak{}FinAll.\allowbreak{}finConstAll} (\texttt{Development.\allowbreak{}AutomorphicSpectrum.\allowbreak{}Proofs.\allowbreak{}GeometryProofs.\allowbreak{}FiniteConstantAllHolds}, lines 137--196)\\
\texttt{GlobalTheory.\allowbreak{}CompanionTheory.\allowbreak{}Ar.\allowbreak{}Br.\allowbreak{}FinGood.\allowbreak{}finGood} (\texttt{Development.\allowbreak{}AutomorphicSpectrum.\allowbreak{}Proofs.\allowbreak{}GeometryProofs.\allowbreak{}FrameFunctionFiniteNorm}, lines 23--41)\par}

{\footnotesize\raggedright\noindent \ref{B8.25}.\quad \texttt{SelfIntersectionVolumeEstimates.\allowbreak{}topCovolSlim2\_\allowbreak{}use\_\allowbreak{}alt\_\allowbreak{}alt} (\texttt{Development.\allowbreak{}AutomorphicArithmetic.\allowbreak{}ThetaCorrespondence.\allowbreak{}ThetaPairLifts.\allowbreak{}ThirdPlaceSeries.\allowbreak{}TopCovolSlim2Use}, lines 16--22)\\
\texttt{PositivityTheory.\allowbreak{}MbHsUse.\allowbreak{}HSUse} (\texttt{Development.\allowbreak{}WeilBound.\allowbreak{}CoreStatements.\allowbreak{}AutomorphicKernels.\allowbreak{}MoriwakiBigness.\allowbreak{}HilbertSamuelApplications.\allowbreak{}SmallGeneratedCountUpperOfUse}, lines 12--13)\par}

{\footnotesize\raggedright\noindent \ref{B8.25}.\quad \texttt{SelfIntersectionVolumeEstimates.\allowbreak{}topCovolG\_\allowbreak{}alt\_\allowbreak{}alt} (\texttt{Development.\allowbreak{}AutomorphicArithmetic.\allowbreak{}ThetaCorrespondence.\allowbreak{}ThetaPairLifts.\allowbreak{}ThirdPlaceSeries.\allowbreak{}TopCovolG}, lines 39--58)\par}

{\footnotesize\raggedright\noindent \ref{B8.25h}.\quad \texttt{ReciprocityTheory.\allowbreak{}ChartRestriction.\allowbreak{}G.\allowbreak{}Geo.\allowbreak{}chartLat} (\texttt{ReciprocityTheory.\allowbreak{}ReductiveGroups.\allowbreak{}ArithmeticGeometry.\allowbreak{}BirationalGeometry.\allowbreak{}ArithmeticHilbertSamuelLowerBound}, lines 59--64)\\
\texttt{ReciprocityTheory.\allowbreak{}ChartRestriction.\allowbreak{}G.\allowbreak{}Geo.\allowbreak{}chartLatNm} (\texttt{ReciprocityTheory.\allowbreak{}ReductiveGroups.\allowbreak{}ArithmeticGeometry.\allowbreak{}BirationalGeometry.\allowbreak{}ArithmeticHilbertSamuelLowerBound}, lines 66--71)\\
\texttt{GroundTheory.\allowbreak{}Heights₂.\allowbreak{}Heights.\allowbreak{}Model.\allowbreak{}FiniteIntegralClosure.\allowbreak{}proof} (\texttt{GroundTheory.\allowbreak{}HeightsAndTheta.\allowbreak{}Heights.\allowbreak{}IntegralModels.\allowbreak{}PFiniteIntegralClosure}, lines 9--17)\\
\texttt{GroundTheory.\allowbreak{}Heights₂.\allowbreak{}Heights.\allowbreak{}HilbertSamuel.\allowbreak{}Conductor.\allowbreak{}exists\_\allowbreak{}chart\_\allowbreak{}conductor} (\texttt{GroundTheory.\allowbreak{}HeightsAndTheta.\allowbreak{}Heights.\allowbreak{}HilbertSamuel.\allowbreak{}PConductor}, lines 54--88)\\
\texttt{GroundTheory.\allowbreak{}Heights₂.\allowbreak{}Heights.\allowbreak{}HilbertSamuel.\allowbreak{}Conductor.\allowbreak{}proof} (\texttt{GroundTheory.\allowbreak{}HeightsAndTheta.\allowbreak{}Heights.\allowbreak{}HilbertSamuel.\allowbreak{}PConductor}, lines 102--160)\\
\texttt{GroundTheory.\allowbreak{}Heights₂.\allowbreak{}Heights.\allowbreak{}HilbertSamuel.\allowbreak{}LeviCompHol.\allowbreak{}proof} (\texttt{GroundTheory.\allowbreak{}HeightsAndTheta.\allowbreak{}Heights.\allowbreak{}HilbertSamuel.\allowbreak{}PLeviCompHol}, lines 11--61)\par}

{\footnotesize\raggedright\noindent \ref{B8.25i}.\quad \texttt{GlobalTheory.\allowbreak{}Izumi.\allowbreak{}exists\_\allowbreak{}dominating\_\allowbreak{}eq\_\allowbreak{}mul\_\allowbreak{}of\_\allowbreak{}not\_\allowbreak{}mem\_\allowbreak{}span} (\texttt{Development.\allowbreak{}AbelianVarietyTheory.\allowbreak{}Proofs.\allowbreak{}DivisorProofs.\allowbreak{}ExistsDominatingLocalDomainEqMulMaximal}, lines 70--99)\\
\texttt{GlobalTheory.\allowbreak{}CurvePoint.\allowbreak{}exists\_\allowbreak{}ringHom\_\allowbreak{}padicAlgCl\_\allowbreak{}ne\_\allowbreak{}zero} (\texttt{Development.\allowbreak{}AbelianVarietyTheory.\allowbreak{}Proofs.\allowbreak{}DivisorProofs.\allowbreak{}ExistsRingHomPadicAlgClNeZero}, lines 10--101)\\
\texttt{GlobalTheory.\allowbreak{}CplDiv.\allowbreak{}dvd\_\allowbreak{}of\_\allowbreak{}norm\_\allowbreak{}le\_\allowbreak{}at\_\allowbreak{}classical} (\texttt{Development.\allowbreak{}AbelianVarietyTheory.\allowbreak{}Proofs.\allowbreak{}DivisorProofs.\allowbreak{}ExistsNormLeOfMemMapPow}, lines 142--208)\par}

{\footnotesize\raggedright\noindent \ref{B8.25j}.\quad \texttt{GlobalTheory.\allowbreak{}ZH.\allowbreak{}exists\_\allowbreak{}normal\_\allowbreak{}proj\_\allowbreak{}hasDiv\_\allowbreak{}family} (\texttt{Development.\allowbreak{}AbelianVarietyTheory.\allowbreak{}Proofs.\allowbreak{}DegenerationProofs.\allowbreak{}ExistsAdmissibleAffineModelDividingTuple}, lines 21--48)\\
\texttt{GlobalTheory.\allowbreak{}Poinc.\allowbreak{}LabConst.\allowbreak{}main} (\texttt{Development.\allowbreak{}AbelianVarietyTheory.\allowbreak{}Proofs.\allowbreak{}DivisorProofs.\allowbreak{}TorsionLabelsConstantCompletedLocalRingClosedPoi}, lines 64--144)\\
\texttt{GlobalTheory.\allowbreak{}ThetaSpan.\allowbreak{}eq\_\allowbreak{}sum\_\allowbreak{}translates} (\texttt{Development.\allowbreak{}AbelianVarietyTheory.\allowbreak{}Proofs.\allowbreak{}UniformizationProofs.\allowbreak{}ValuationNonnegIffSpanMultiplicationTranslates}, lines 32--91)\\
\texttt{GlobalTheory.\allowbreak{}UT.\allowbreak{}nuGe\_\allowbreak{}zero\_\allowbreak{}iff\_\allowbreak{}span\_\allowbreak{}translates} (\texttt{Development.\allowbreak{}AbelianVarietyTheory.\allowbreak{}Proofs.\allowbreak{}UniformizationProofs.\allowbreak{}ValuationNonnegIffSpanMultiplicationTranslates}, lines 219--331)\\
\texttt{GlobalTheory.\allowbreak{}UT.\allowbreak{}CoefUnit.\allowbreak{}exists\_\allowbreak{}fam\_\allowbreak{}coef\_\allowbreak{}val} (\texttt{Development.\allowbreak{}AbelianVarietyTheory.\allowbreak{}Proofs.\allowbreak{}UniformizationProofs.\allowbreak{}ExistsFamCoefVal}, lines 130--161)\\
\texttt{GlobalTheory.\allowbreak{}ZH.\allowbreak{}WeightBound.\allowbreak{}main} (\texttt{Development.\allowbreak{}AbelianVarietyTheory.\allowbreak{}Proofs.\allowbreak{}DegenerationProofs.\allowbreak{}WeightSpaceValuesIntegralTorsionBoundShortestVec}, lines 319--341)\\
\texttt{GlobalTheory.\allowbreak{}TestFibre.\allowbreak{}exists\_\allowbreak{}norm\_\allowbreak{}FF\_\allowbreak{}eq} (\texttt{Development.\allowbreak{}AbelianVarietyTheory.\allowbreak{}Proofs.\allowbreak{}DegenerationProofs.\allowbreak{}ExistsNormTorsionSumShortestVectorEq}, lines 88--191)\\
\texttt{GlobalTheory.\allowbreak{}Poinc.\allowbreak{}CNorm.\allowbreak{}main} (\texttt{Development.\allowbreak{}AbelianVarietyTheory.\allowbreak{}Proofs.\allowbreak{}DivisorProofs.\allowbreak{}ThetaConstantFourierPartsBoundedClassMinima}, lines 100--163)\\
\texttt{GlobalTheory.\allowbreak{}QUnit.\allowbreak{}main} (\texttt{Development.\allowbreak{}AbelianVarietyTheory.\allowbreak{}Proofs.\allowbreak{}DivisorProofs.\allowbreak{}ExistsMinimalPsiCoordinateAssociatedFourierPower}, lines 25--145)\par}

{\footnotesize\raggedright\noindent \ref{B8.25k}.\quad \texttt{GlobalTheory.\allowbreak{}TraitLift.\allowbreak{}exists\_\allowbreak{}lift\_\allowbreak{}adicCompletion} (\texttt{Development.\allowbreak{}AbelianVarietyTheory.\allowbreak{}Proofs.\allowbreak{}DivisorProofs.\allowbreak{}ExistsLiftAdicCompletion}, lines 11--36)\\
\texttt{GlobalTheory.\allowbreak{}TorusOntoSpread.\allowbreak{}restrict\_\allowbreak{}surjective} (\texttt{Development.\allowbreak{}AbelianVarietyTheory.\allowbreak{}Proofs.\allowbreak{}DivisorProofs.\allowbreak{}RestrictSurjectiveAlt}, lines 18--90)\\
\texttt{GlobalTheory.\allowbreak{}Poinc.\allowbreak{}TorCoord.\allowbreak{}coordMap\_\allowbreak{}of\_\allowbreak{}torMatch} (\texttt{Development.\allowbreak{}AbelianVarietyTheory.\allowbreak{}Proofs.\allowbreak{}DivisorProofs.\allowbreak{}CoordMapOfTorMatch}, lines 22--76)\\
\texttt{GlobalTheory.\allowbreak{}UPoinc.\allowbreak{}TraitAtClosed.\allowbreak{}exists\_\allowbreak{}trait\_\allowbreak{}lattice\_\allowbreak{}map} (\texttt{Development.\allowbreak{}AbelianVarietyTheory.\allowbreak{}Proofs.\allowbreak{}DivisorProofs.\allowbreak{}ExistsTraitLatticeMapAtClosedPoint}, lines 24--87)\\
\texttt{GlobalTheory.\allowbreak{}UPoinc.\allowbreak{}TraitPull.\allowbreak{}exists\_\allowbreak{}pull\_\allowbreak{}unif} (\texttt{Development.\allowbreak{}AbelianVarietyTheory.\allowbreak{}Proofs.\allowbreak{}DivisorProofs.\allowbreak{}ExistsLinearUniformisationPulledBackTrait}, lines 18--35)\\
\texttt{GlobalTheory.\allowbreak{}UPoinc.\allowbreak{}uOrdEq\_\allowbreak{}stalkHom\_\allowbreak{}of\_\allowbreak{}associated} (\texttt{Development.\allowbreak{}AbelianVarietyTheory.\allowbreak{}Proofs.\allowbreak{}DivisorProofs.\allowbreak{}UniformiserOrderEqOfAssociatedStalkImage}, lines 17--45)\par}

{\footnotesize\raggedright\noindent \ref{B8.25m}.\quad \texttt{GlobalTheory.\allowbreak{}UPoinc.\allowbreak{}CellCones.\allowbreak{}exists\_\allowbreak{}finite\_\allowbreak{}cones} (\texttt{Development.\allowbreak{}AbelianVarietyTheory.\allowbreak{}Proofs.\allowbreak{}DivisorProofs.\allowbreak{}ExistsFiniteConesPolarizedFanLevelStructure}, lines 99--215)\\
\texttt{GlobalTheory.\allowbreak{}Poinc.\allowbreak{}MinSysLat.\allowbreak{}finite\_\allowbreak{}minimiser\_\allowbreak{}systems\_\allowbreak{}lat} (\texttt{Development.\allowbreak{}AbelianVarietyTheory.\allowbreak{}Proofs.\allowbreak{}DivisorProofs.\allowbreak{}FiniteMinimiserSystemsPosCombinationDefiniteSum}, lines 24--70)\par}

{\footnotesize\raggedright\noindent \ref{B8.25n}.\quad \texttt{GlobalTheory.\allowbreak{}Poinc.\allowbreak{}exists\_\allowbreak{}idealSheaf\_\allowbreak{}comap\_\allowbreak{}eq\_\allowbreak{}of\_\allowbreak{}approx} (\texttt{Development.\allowbreak{}AbelianVarietyTheory.\allowbreak{}Proofs.\allowbreak{}DivisorProofs.\allowbreak{}ExistsIdealSheafPullbackEqDivisorOfApprox}, lines 16--40)\par}

{\footnotesize\raggedright\noindent \ref{B8.25o}.\quad \texttt{GlobalTheory.\allowbreak{}NormSheaf.\allowbreak{}normDescentSheaf} (\texttt{Development.\allowbreak{}AbelianVarietyTheory.\allowbreak{}Proofs.\allowbreak{}DivisorProofs.\allowbreak{}NormDescentIdealSheafFiniteOpenClosedImage}, lines 16--42)\\
\texttt{GlobalTheory.\allowbreak{}NormSheaf.\allowbreak{}normIdeal\_\allowbreak{}map\_\allowbreak{}eq\_\allowbreak{}prod\_\allowbreak{}lifts} (\texttt{Development.\allowbreak{}AbelianVarietyTheory.\allowbreak{}Proofs.\allowbreak{}DivisorProofs.\allowbreak{}NormCoefficientIdealMapCompleteDvrEqProdLifts}, lines 13--90)\par}

{\footnotesize\raggedright\noindent \ref{B8.25a}.\quad \texttt{UpperTheory.\allowbreak{}LiteratureTheorems.\allowbreak{}published\_\allowbreak{}results\_\allowbreak{}arithmetic} (\texttt{Development.\allowbreak{}AutomorphicArithmetic.\allowbreak{}ThetaCorrespondence.\allowbreak{}ThetaPairLifts.\allowbreak{}ThirdPlaceSeries.\allowbreak{}EveryCharacterOfLiteratureOfVWin}, lines 781--809)\\
\texttt{UpperTheory.\allowbreak{}LiteratureTheorems.\allowbreak{}every\_\allowbreak{}character\_\allowbreak{}of\_\allowbreak{}literature\_\allowbreak{}of\_\allowbreak{}vWin} (\texttt{Development.\allowbreak{}AutomorphicArithmetic.\allowbreak{}ThetaCorrespondence.\allowbreak{}ThetaPairLifts.\allowbreak{}ThirdPlaceSeries.\allowbreak{}EveryCharacterOfLiteratureOfVWin}, lines 810--813)\par}

{\footnotesize\raggedright\noindent \ref{B8.25b}.\quad \texttt{GlobalTheory.\allowbreak{}T.\allowbreak{}cov\_\allowbreak{}size\_\allowbreak{}bound} (\texttt{Development.\allowbreak{}AutomorphicArithmetic.\allowbreak{}ThetaCorrespondence.\allowbreak{}ThetaPairLifts.\allowbreak{}ThirdPlaceSeries.\allowbreak{}SizeAltBisStatement}, lines 11--12)\par}

{\footnotesize\raggedright\noindent \ref{B8.25b}.\quad \texttt{GlobalTheory.\allowbreak{}CovTop.\allowbreak{}SizeRows.\allowbreak{}exists\_\allowbreak{}isDiv} (\texttt{Development.\allowbreak{}AbelianVarietyTheory.\allowbreak{}Proofs.\allowbreak{}DivisorProofs.\allowbreak{}SizeBoundStatement}, lines 38--51)\\
\texttt{GlobalTheory.\allowbreak{}CovTop.\allowbreak{}SizeRows.\allowbreak{}rho\_\allowbreak{}bounds} (\texttt{Development.\allowbreak{}AbelianVarietyTheory.\allowbreak{}Proofs.\allowbreak{}DivisorProofs.\allowbreak{}SizeBoundStatement}, lines 52--107)\\
\texttt{GlobalTheory.\allowbreak{}CovTop.\allowbreak{}SizeRows.\allowbreak{}sizeAt\_\allowbreak{}of} (\texttt{Development.\allowbreak{}AbelianVarietyTheory.\allowbreak{}Proofs.\allowbreak{}DivisorProofs.\allowbreak{}SizeBoundStatement}, lines 260--359)\\
\texttt{GlobalTheory.\allowbreak{}CovTop.\allowbreak{}SizeRows.\allowbreak{}size\_\allowbreak{}of\_\allowbreak{}alt} (\texttt{Development.\allowbreak{}AutomorphicArithmetic.\allowbreak{}ThetaCorrespondence.\allowbreak{}ThetaPairLifts.\allowbreak{}ThirdPlaceSeries.\allowbreak{}SizeAltBis}, lines 18--23)\par}

{\footnotesize\raggedright\noindent \ref{B8.25c}.\quad \texttt{GlobalTheory.\allowbreak{}CovTop.\allowbreak{}DeepExist.\allowbreak{}exists\_\allowbreak{}deep\_\allowbreak{}covering\_\allowbreak{}level\_\allowbreak{}alt} (\texttt{Development.\allowbreak{}AutomorphicArithmetic.\allowbreak{}ThetaCorrespondence.\allowbreak{}ThetaPairLifts.\allowbreak{}ThirdPlaceSeries.\allowbreak{}ExistsDeepCoveringLevel}, lines 16--24)\par}

{\footnotesize\raggedright\noindent \ref{B8.25d}.\quad \texttt{GlobalTheory.\allowbreak{}CovTop.\allowbreak{}covering\_\allowbreak{}gram\_\allowbreak{}matrix\_\allowbreak{}of\_\allowbreak{}local\_\allowbreak{}traces\_\allowbreak{}alt} (\texttt{Development.\allowbreak{}AutomorphicArithmetic.\allowbreak{}ThetaCorrespondence.\allowbreak{}ThetaPairLifts.\allowbreak{}ThirdPlaceSeries.\allowbreak{}CoveringGramMatrixOfLocalTraces}, lines 17--20)\\
\texttt{GlobalTheory.\allowbreak{}CovTop.\allowbreak{}MeasRd.\allowbreak{}covering\_\allowbreak{}measure\_\allowbreak{}reduction\_\allowbreak{}alt} (\texttt{Development.\allowbreak{}AutomorphicArithmetic.\allowbreak{}ThetaCorrespondence.\allowbreak{}ThetaPairLifts.\allowbreak{}ThirdPlaceSeries.\allowbreak{}CoveringMeasureReductionBound}, lines 18--27)\\
\texttt{GlobalTheory.\allowbreak{}CovTop.\allowbreak{}LamAll.\allowbreak{}covering\_\allowbreak{}lambda\_\allowbreak{}all\_\allowbreak{}places\_\allowbreak{}alt} (\texttt{Development.\allowbreak{}AutomorphicArithmetic.\allowbreak{}ThetaCorrespondence.\allowbreak{}ThetaPairLifts.\allowbreak{}ThirdPlaceSeries.\allowbreak{}CoveringLambdaInvariantBound}, lines 21--37)\par}

{\footnotesize\raggedright\noindent \ref{B8.25d}.\quad \texttt{GlobalTheory.\allowbreak{}CovHt.\allowbreak{}WtGap} (\texttt{Development.\allowbreak{}AbelianVarietyTheory.\allowbreak{}Definitions.\allowbreak{}Divisors.\allowbreak{}SemiposAtOfPositive}, lines 88--90)\\
\texttt{GlobalTheory.\allowbreak{}CovHt.\allowbreak{}bndB} (\texttt{Development.\allowbreak{}AbelianVarietyTheory.\allowbreak{}Definitions.\allowbreak{}Divisors.\allowbreak{}SemiposAtOfPositive}, lines 91--93)\\
\texttt{GlobalTheory.\allowbreak{}CovHt.\allowbreak{}errX} (\texttt{Development.\allowbreak{}AbelianVarietyTheory.\allowbreak{}Definitions.\allowbreak{}Divisors.\allowbreak{}SemiposAtOfPositive}, lines 94--96)\\
\texttt{GlobalTheory.\allowbreak{}CompHt.\allowbreak{}ValFam} (\texttt{Development.\allowbreak{}AbelianVarietyTheory.\allowbreak{}Definitions.\allowbreak{}Divisors.\allowbreak{}IsHomogeneousCompIdeal}, lines 245--247)\\
\texttt{GlobalTheory.\allowbreak{}CompHt.\allowbreak{}recVals} (\texttt{Development.\allowbreak{}AbelianVarietyTheory.\allowbreak{}Definitions.\allowbreak{}Divisors.\allowbreak{}IsHomogeneousCompIdeal}, lines 250--251)\\
\texttt{GlobalTheory.\allowbreak{}CompHt.\allowbreak{}ValFam.\allowbreak{}patch} (\texttt{Development.\allowbreak{}AbelianVarietyTheory.\allowbreak{}Definitions.\allowbreak{}Divisors.\allowbreak{}IsHomogeneousCompIdeal}, lines 254--259)\\
\texttt{GlobalTheory.\allowbreak{}CompHt.\allowbreak{}htVals} (\texttt{Development.\allowbreak{}AbelianVarietyTheory.\allowbreak{}Definitions.\allowbreak{}Divisors.\allowbreak{}IsLatBasisHtLatIffSpanEq}, lines 485--490)\\
\texttt{GlobalTheory.\allowbreak{}CompHt.\allowbreak{}htPrimes} (\texttt{Development.\allowbreak{}AbelianVarietyTheory.\allowbreak{}Definitions.\allowbreak{}Divisors.\allowbreak{}IsLatBasisHtLatIffSpanEq}, lines 515--516)\\
\texttt{GlobalTheory.\allowbreak{}CompHt.\allowbreak{}mixVals} (\texttt{Development.\allowbreak{}AbelianVarietyTheory.\allowbreak{}Definitions.\allowbreak{}Divisors.\allowbreak{}IsLatBasisHtLatIffSpanEq}, lines 534--535)\\
\texttt{GlobalTheory.\allowbreak{}CompHt.\allowbreak{}compVol} (\texttt{Development.\allowbreak{}AbelianVarietyTheory.\allowbreak{}Definitions.\allowbreak{}Divisors.\allowbreak{}IsHomogeneousCompIdeal}, lines 273--275)\\
\texttt{GlobalTheory.\allowbreak{}CompHt.\allowbreak{}compVolSlope} (\texttt{Development.\allowbreak{}AbelianVarietyTheory.\allowbreak{}Definitions.\allowbreak{}Divisors.\allowbreak{}IsHomogeneousCompIdeal}, lines 280--282)\par}

{\footnotesize\raggedright\noindent \ref{B8.26p}.\quad \texttt{GroundTheory.\allowbreak{}Heights₂.\allowbreak{}Heights.\allowbreak{}LowerBoundFree.\allowbreak{}ZhangLowerX.\allowbreak{}Statement} (\texttt{GroundTheory.\allowbreak{}HeightsAndTheta.\allowbreak{}Heights.\allowbreak{}HodgeLowerBound.\allowbreak{}StatementsZhangLowerX}, lines 11--18)\\
\texttt{GroundTheory.\allowbreak{}Heights₂.\allowbreak{}Heights.\allowbreak{}IndefiniteRowSideFullHecke.\allowbreak{}HeightsSet.\allowbreak{}height} (\texttt{GroundTheory.\allowbreak{}HeightsAndTheta.\allowbreak{}Heights.\allowbreak{}HeightsOfSets.\allowbreak{}IsSaturatedFirstCoordinateLine}, lines 19--22)\\
\texttt{GlobalTheory.\allowbreak{}CompanionTheory.\allowbreak{}Ar.\allowbreak{}Giv.\allowbreak{}lineOfLargeDegree\_\allowbreak{}of\_\allowbreak{}use} (\texttt{Development.\allowbreak{}WeilBound.\allowbreak{}CoreStatements.\allowbreak{}AutomorphicKernels.\allowbreak{}MoriwakiBigness.\allowbreak{}HilbertSamuelApplications.\allowbreak{}LineOfLargeDegreeOfUse}, lines 10--89)\par}

{\footnotesize\raggedright\noindent \ref{B8.26q}.\quad \texttt{PositivityTheory.\allowbreak{}HtPt.\allowbreak{}Ah.\allowbreak{}ais} (\texttt{Development.\allowbreak{}WeilBound.\allowbreak{}CoreStatements.\allowbreak{}HodgeTatePoints.\allowbreak{}AhSlim}, lines 22--25)\\
\texttt{GlobalTheory.\allowbreak{}CompanionTheory.\allowbreak{}Ar.\allowbreak{}Tw.\allowbreak{}Main.\allowbreak{}poincareLelong\_\allowbreak{}of\_\allowbreak{}inv} (\texttt{Development.\allowbreak{}AutomorphicSpectrum.\allowbreak{}Proofs.\allowbreak{}MetricProofs.\allowbreak{}PoincareLelongOfInv}, lines 262--332)\\
\texttt{GlobalTheory.\allowbreak{}CompanionTheory.\allowbreak{}Ar.\allowbreak{}Pub.\allowbreak{}arithIntersectionSystem\_\allowbreak{}of\_\allowbreak{}poincareLelong\_\allowbreak{}resK} (\texttt{Development.\allowbreak{}AutomorphicSpectrum.\allowbreak{}Proofs.\allowbreak{}MetricProofs.\allowbreak{}ArithIntersectionSystemOfPoincareLelong}, lines 42--49)\\
\texttt{GlobalTheory.\allowbreak{}CompanionTheory.\allowbreak{}SAH.\allowbreak{}semiAbelianHomExtension} (\texttt{Development.\allowbreak{}AutomorphicSpectrum.\allowbreak{}Proofs.\allowbreak{}GlobalProofs.\allowbreak{}SemiAbelianHomExtension}, lines 13--20)\\
\texttt{GlobalTheory.\allowbreak{}HomExtNormal.\allowbreak{}existsUnique\_\allowbreak{}hom\_\allowbreak{}extension\_\allowbreak{}of\_\allowbreak{}dvr} (\texttt{Development.\allowbreak{}AbelianVarietyTheory.\allowbreak{}Proofs.\allowbreak{}IwasawaDecomposition.\allowbreak{}IsOpenImmersionQuasiFiniteLocusOfIsIntegral}, lines 263--557)\\
\texttt{GlobalTheory.\allowbreak{}existsUnique\_\allowbreak{}hom\_\allowbreak{}extension\_\allowbreak{}semiAbelian} (\texttt{Development.\allowbreak{}AbelianVarietyTheory.\allowbreak{}Proofs.\allowbreak{}CoherentCohomologyProofs.\allowbreak{}ExistsUniqueHomExtensionSemiAbelian}, lines 11--15)\\
\texttt{PositivityTheory.\allowbreak{}Xgn.\allowbreak{}farguesHodgeTate\_\allowbreak{}of\_\allowbreak{}farguesExact} (\texttt{Development.\allowbreak{}WeilBound.\allowbreak{}CoreStatements.\allowbreak{}TorusUniformisation.\allowbreak{}ChowMahlerFamilies.\allowbreak{}XXgzFarOfExact}, lines 50--61)\\
\texttt{GlobalTheory.\allowbreak{}CompanionTheory.\allowbreak{}Gaga.\allowbreak{}gagaSections} (\texttt{Development.\allowbreak{}AutomorphicSpectrum.\allowbreak{}Proofs.\allowbreak{}GlobalProofs.\allowbreak{}GagaSections}, lines 104--139)\\
\texttt{GlobalTheory.\allowbreak{}CompanionTheory.\allowbreak{}Ar.\allowbreak{}GR.\allowbreak{}QuotientByFiniteFlatStatement} (\texttt{Development.\allowbreak{}AutomorphicSpectrum.\allowbreak{}PublishedTheorems.\allowbreak{}Algebra₂.\allowbreak{}QuotientByFiniteFlatStatement}, lines 21--30)\\
\texttt{GlobalTheory.\allowbreak{}CompanionTheory.\allowbreak{}Ar.\allowbreak{}GR.\allowbreak{}ZarhinStatement} (\texttt{Development.\allowbreak{}AutomorphicSpectrum.\allowbreak{}PublishedTheorems.\allowbreak{}Algebra₂.\allowbreak{}ZarhinStatement}, lines 21--28)\par}

{\footnotesize\raggedright\noindent \ref{B8.26u}.\quad \texttt{GlobalTheory.\allowbreak{}CompanionTheory.\allowbreak{}Ar.\allowbreak{}Link.\allowbreak{}good\_\allowbreak{}tower\_\allowbreak{}and\_\allowbreak{}volume\_\allowbreak{}final} (\texttt{Development.\allowbreak{}AutomorphicSpectrum.\allowbreak{}Proofs.\allowbreak{}GeometryProofs.\allowbreak{}GoodTowerVolumeTheoremAlt}, lines 32--43)\\
\texttt{GlobalTheory.\allowbreak{}CompanionTheory.\allowbreak{}Ar.\allowbreak{}Top.\allowbreak{}goodTower\_\allowbreak{}and\_\allowbreak{}volume} (\texttt{Development.\allowbreak{}AutomorphicSpectrum.\allowbreak{}Proofs.\allowbreak{}DivisibilityProofs.\allowbreak{}GoodTowerVolumeTheorem}, lines 36--51)\\
\texttt{GlobalTheory.\allowbreak{}CompanionTheory.\allowbreak{}Ar.\allowbreak{}Top.\allowbreak{}GoodTowerExistsX} (\texttt{Development.\allowbreak{}AutomorphicSpectrum.\allowbreak{}Statements.\allowbreak{}TowerExistence.\allowbreak{}GoodTowerExistsX}, lines 22--27)\\
\texttt{GlobalTheory.\allowbreak{}T.\allowbreak{}exists\_\allowbreak{}good\_\allowbreak{}root\_\allowbreak{}tower} (\texttt{Development.\allowbreak{}AutomorphicSpectrum.\allowbreak{}Proofs.\allowbreak{}GeometryProofs.\allowbreak{}ExistsGoodRootTowerStatement}, lines 12--14)\par}

{\footnotesize\raggedright\noindent \ref{B8.26}.\quad \texttt{PositivityTheory.\allowbreak{}HtPt.\allowbreak{}Ah.\allowbreak{}deep\_\allowbreak{}slim\_\allowbreak{}use} (\texttt{Development.\allowbreak{}WeilBound.\allowbreak{}CoreStatements.\allowbreak{}AutomorphicKernels.\allowbreak{}MoriwakiBigness.\allowbreak{}HilbertSamuelApplications.\allowbreak{}AhSlimU}, lines 9--15)\\
\texttt{GlobalTheory.\allowbreak{}CompanionTheory.\allowbreak{}Ar.\allowbreak{}Plug.\allowbreak{}DeepStR} (\texttt{Development.\allowbreak{}AutomorphicSpectrum.\allowbreak{}Definitions.\allowbreak{}GlobalGeometry.\allowbreak{}ModelSmall}, lines 63--67)\\
\texttt{GlobalTheory.\allowbreak{}CovHt.\allowbreak{}MeanGe} (\texttt{Development.\allowbreak{}AbelianVarietyTheory.\allowbreak{}Definitions.\allowbreak{}Divisors.\allowbreak{}SemiposAtOfPositive}, lines 97--99)\\
\texttt{GlobalTheory.\allowbreak{}CovHt.\allowbreak{}MeanAll} (\texttt{Development.\allowbreak{}AbelianVarietyTheory.\allowbreak{}Definitions.\allowbreak{}Divisors.\allowbreak{}SemiposAtOfPositive}, lines 101--101)\par}

{\footnotesize\raggedright\noindent \ref{B8.26}.\quad \texttt{GlobalTheory.\allowbreak{}CompanionTheory.\allowbreak{}Ar.\allowbreak{}Plug.\allowbreak{}deepStR\_\allowbreak{}of\_\allowbreak{}use} (\texttt{Development.\allowbreak{}WeilBound.\allowbreak{}CoreStatements.\allowbreak{}AutomorphicKernels.\allowbreak{}MoriwakiBigness.\allowbreak{}HilbertSamuelApplications.\allowbreak{}DeepStROfUse}, lines 12--24)\\
\texttt{GlobalTheory.\allowbreak{}CompanionTheory.\allowbreak{}Ar.\allowbreak{}Link.\allowbreak{}good\_\allowbreak{}tower\_\allowbreak{}volume\_\allowbreak{}theorem\_\allowbreak{}bis} (\texttt{Development.\allowbreak{}WeilBound.\allowbreak{}CoreStatements.\allowbreak{}AutomorphicKernels.\allowbreak{}MoriwakiBigness.\allowbreak{}HilbertSamuelApplications.\allowbreak{}GoodTowerVolumeTheoremBis}, lines 9--20)\\
\texttt{GlobalTheory.\allowbreak{}CompanionTheory.\allowbreak{}Ar.\allowbreak{}Plug.\allowbreak{}top\_\allowbreak{}meanAll} (\texttt{Development.\allowbreak{}AutomorphicSpectrum.\allowbreak{}Proofs.\allowbreak{}GeometryProofs.\allowbreak{}AbelianCmCompLawsBound}, lines 54--138)\\
\texttt{GlobalTheory.\allowbreak{}CompanionTheory.\allowbreak{}Ar.\allowbreak{}Plug.\allowbreak{}deepStR\_\allowbreak{}of\_\allowbreak{}volume} (\texttt{Development.\allowbreak{}AutomorphicSpectrum.\allowbreak{}Proofs.\allowbreak{}GeometryProofs.\allowbreak{}DeepStROfVolume}, lines 21--59)\\
\texttt{GlobalTheory.\allowbreak{}CompanionTheory.\allowbreak{}Ar.\allowbreak{}hostLval} (\texttt{Development.\allowbreak{}AutomorphicSpectrum.\allowbreak{}Definitions.\allowbreak{}GlobalGeometry.\allowbreak{}TowerOrigin}, lines 123--126)\\
\texttt{GlobalTheory.\allowbreak{}CompanionTheory.\allowbreak{}Ar.\allowbreak{}Plug.\allowbreak{}costTwo} (\texttt{Development.\allowbreak{}AutomorphicSpectrum.\allowbreak{}Definitions.\allowbreak{}GlobalGeometry.\allowbreak{}ModelSmall}, lines 54--54)\par}

{\footnotesize\raggedright\noindent \ref{B8.26}.\quad \texttt{GlobalTheory.\allowbreak{}CompanionTheory.\allowbreak{}Ar.\allowbreak{}Top.\allowbreak{}GoodOrigin.\allowbreak{}Extra} (\texttt{Development.\allowbreak{}AutomorphicSpectrum.\allowbreak{}Statements.\allowbreak{}TowerExistence.\allowbreak{}GoodTowerExistsX}, lines 11--20)\par}

{\footnotesize\raggedright\noindent \ref{B8.26r}.\quad \texttt{GlobalTheory.\allowbreak{}CompanionTheory.\allowbreak{}Ar.\allowbreak{}Top.\allowbreak{}RootModel} (\texttt{Development.\allowbreak{}AutomorphicSpectrum.\allowbreak{}Definitions.\allowbreak{}DivisibilityTowers.\allowbreak{}Birational}, lines 11--42)\\
\texttt{GlobalTheory.\allowbreak{}CompanionTheory.\allowbreak{}Ar.\allowbreak{}Top.\allowbreak{}RootTower} (\texttt{Development.\allowbreak{}AutomorphicSpectrum.\allowbreak{}Definitions.\allowbreak{}DivisibilityTowers.\allowbreak{}Proportional}, lines 10--27)\\
\texttt{GlobalTheory.\allowbreak{}CompanionTheory.\allowbreak{}Ar.\allowbreak{}Top.\allowbreak{}RootTower.\allowbreak{}PPc} (\texttt{Development.\allowbreak{}AutomorphicSpectrum.\allowbreak{}Definitions.\allowbreak{}DivisibilityTowers.\allowbreak{}Proportional}, lines 45--46)\\
\texttt{GlobalTheory.\allowbreak{}CompanionTheory.\allowbreak{}Ar.\allowbreak{}Top.\allowbreak{}RootTower.\allowbreak{}PPt} (\texttt{Development.\allowbreak{}AutomorphicSpectrum.\allowbreak{}Definitions.\allowbreak{}DivisibilityTowers.\allowbreak{}Proportional}, lines 47--48)\\
\texttt{GlobalTheory.\allowbreak{}CompanionTheory.\allowbreak{}Ar.\allowbreak{}Top.\allowbreak{}RootTower.\allowbreak{}volT} (\texttt{Development.\allowbreak{}AutomorphicSpectrum.\allowbreak{}Definitions.\allowbreak{}DivisibilityTowers.\allowbreak{}Proportional}, lines 43--44)\\
\texttt{GlobalTheory.\allowbreak{}CompanionTheory.\allowbreak{}Ar.\allowbreak{}Top.\allowbreak{}RootTower.\allowbreak{}slope} (\texttt{Development.\allowbreak{}AutomorphicSpectrum.\allowbreak{}Definitions.\allowbreak{}DivisibilityTowers.\allowbreak{}Proportional}, lines 62--64)\\
\texttt{GlobalTheory.\allowbreak{}CompanionTheory.\allowbreak{}Ar.\allowbreak{}Top.\allowbreak{}degFl} (\texttt{Development.\allowbreak{}AutomorphicSpectrum.\allowbreak{}Definitions.\allowbreak{}DivisibilityTowers.\allowbreak{}CompactGapsGE}, lines 22--25)\\
\texttt{GlobalTheory.\allowbreak{}CompanionTheory.\allowbreak{}Ar.\allowbreak{}Top.\allowbreak{}linRoot} (\texttt{Development.\allowbreak{}AutomorphicSpectrum.\allowbreak{}Definitions.\allowbreak{}DivisibilityTowers.\allowbreak{}CompactGapsGE}, lines 26--27)\\
\texttt{GlobalTheory.\allowbreak{}CompanionTheory.\allowbreak{}Ar.\allowbreak{}Top.\allowbreak{}flagDeg} (\texttt{Development.\allowbreak{}AutomorphicSpectrum.\allowbreak{}Definitions.\allowbreak{}DivisibilityTowers.\allowbreak{}CompactGapsGE}, lines 28--29)\\
\texttt{GlobalTheory.\allowbreak{}CompanionTheory.\allowbreak{}Ar.\allowbreak{}Top.\allowbreak{}Dominant} (\texttt{Development.\allowbreak{}AutomorphicSpectrum.\allowbreak{}Definitions.\allowbreak{}DivisibilityTowers.\allowbreak{}CompactGapsGE}, lines 52--53)\\
\texttt{GlobalTheory.\allowbreak{}CompanionTheory.\allowbreak{}Ar.\allowbreak{}Top.\allowbreak{}RegDominant} (\texttt{Development.\allowbreak{}AutomorphicSpectrum.\allowbreak{}Definitions.\allowbreak{}DivisibilityTowers.\allowbreak{}CompactGapsGE}, lines 54--55)\\
\texttt{GlobalTheory.\allowbreak{}CompanionTheory.\allowbreak{}Ar.\allowbreak{}Top.\allowbreak{}I9} (\texttt{Development.\allowbreak{}AutomorphicSpectrum.\allowbreak{}Definitions.\allowbreak{}DivisibilityTowers.\allowbreak{}Proportional}, lines 133--138)\\
\texttt{GlobalTheory.\allowbreak{}CompanionTheory.\allowbreak{}Ar.\allowbreak{}Top.\allowbreak{}FrobShape} (\texttt{Development.\allowbreak{}AutomorphicSpectrum.\allowbreak{}Definitions.\allowbreak{}DivisibilityTowers.\allowbreak{}HeckeSectionsIsotropicCommutatorPairingPower}, lines 15--16)\\
\texttt{GlobalTheory.\allowbreak{}CompanionTheory.\allowbreak{}Ar.\allowbreak{}Top.\allowbreak{}deltaTerm} (\texttt{Development.\allowbreak{}AutomorphicSpectrum.\allowbreak{}Definitions.\allowbreak{}DivisibilityTowers.\allowbreak{}CompactGapsGE}, lines 83--86)\par}

{\footnotesize\raggedright\noindent \ref{B8.26s}.\quad \texttt{GlobalTheory.\allowbreak{}CompanionTheory.\allowbreak{}Ar.\allowbreak{}Top.\allowbreak{}flagDeg\_\allowbreak{}eq\_\allowbreak{}weyl} (\texttt{Development.\allowbreak{}AutomorphicSpectrum.\allowbreak{}Proofs.\allowbreak{}DivisibilityProofs.\allowbreak{}FlagDegreeWeightFormula}, lines 9--33)\\
\texttt{GlobalTheory.\allowbreak{}CompanionTheory.\allowbreak{}Ar.\allowbreak{}Top.\allowbreak{}form\_\allowbreak{}top\_\allowbreak{}succ\_\allowbreak{}eq\_\allowbreak{}sum\_\allowbreak{}pk} (\texttt{Development.\allowbreak{}AutomorphicSpectrum.\allowbreak{}Proofs.\allowbreak{}DivisibilityProofs.\allowbreak{}HomogeneousPowerSumDecomposition}, lines 9--11)\\
\texttt{GlobalTheory.\allowbreak{}CompanionTheory.\allowbreak{}Ar.\allowbreak{}Top.\allowbreak{}sh\_\allowbreak{}tot\_\allowbreak{}pk\_\allowbreak{}mul} (\texttt{Development.\allowbreak{}AutomorphicSpectrum.\allowbreak{}Proofs.\allowbreak{}DivisibilityProofs.\allowbreak{}FlagShiftTotalPowerSum}, lines 8--10)\\
\texttt{GlobalTheory.\allowbreak{}CompanionTheory.\allowbreak{}Ar.\allowbreak{}Top.\allowbreak{}testForms\_\allowbreak{}sh\_\allowbreak{}tot} (\texttt{Development.\allowbreak{}AutomorphicSpectrum.\allowbreak{}Proofs.\allowbreak{}DivisibilityProofs.\allowbreak{}ExistsFlagDegreeOnePolys}, lines 9--11)\\
\texttt{GlobalTheory.\allowbreak{}T.\allowbreak{}power\_\allowbreak{}sum\_\allowbreak{}delta\_\allowbreak{}decomposition} (\texttt{Development.\allowbreak{}AutomorphicSpectrum.\allowbreak{}Proofs.\allowbreak{}DivisibilityProofs.\allowbreak{}PowerSumDeltaDecompositionStatement}, lines 11--13)\\
\texttt{GlobalTheory.\allowbreak{}CompanionTheory.\allowbreak{}Ar.\allowbreak{}Top.\allowbreak{}restTerm\_\allowbreak{}bound\_\allowbreak{}and\_\allowbreak{}tendsto} (\texttt{Development.\allowbreak{}AutomorphicSpectrum.\allowbreak{}Proofs.\allowbreak{}DivisibilityProofs.\allowbreak{}RemainderTermSummable}, lines 14--16)\par}

{\footnotesize\raggedright\noindent \ref{B8.26a}.\quad \texttt{GlobalTheory.\allowbreak{}CompanionTheory.\allowbreak{}Ar.\allowbreak{}Top.\allowbreak{}slope\_\allowbreak{}lower\_\allowbreak{}of\_\allowbreak{}base\_\allowbreak{}settings} (\texttt{Development.\allowbreak{}AutomorphicSpectrum.\allowbreak{}Proofs.\allowbreak{}DivisibilityProofs.\allowbreak{}SlopeLowerBaseSettings}, lines 110--118)\par}

{\footnotesize\raggedright\noindent \ref{B8.26d}.\quad \texttt{GroundTheory.\allowbreak{}Heights₂.\allowbreak{}Heights.\allowbreak{}RestTerms.\allowbreak{}SameFactorBound.\allowbreak{}proof} (\texttt{GroundTheory.\allowbreak{}HeightsAndTheta.\allowbreak{}Heights.\allowbreak{}RemainderTerms.\allowbreak{}PSameFactorBound}, lines 43--157)\\
\texttt{GroundTheory.\allowbreak{}Heights₂.\allowbreak{}Heights.\allowbreak{}RestTerms.\allowbreak{}OtherFactorBound.\allowbreak{}proof} (\texttt{GroundTheory.\allowbreak{}HeightsAndTheta.\allowbreak{}Heights.\allowbreak{}RemainderTerms.\allowbreak{}POtherFactorBound}, lines 83--188)\par}

{\footnotesize\raggedright\noindent \ref{B8.26b}.\quad \texttt{GlobalTheory.\allowbreak{}CompanionTheory.\allowbreak{}Ar.\allowbreak{}Top.\allowbreak{}exists\_\allowbreak{}identAt} (\texttt{Development.\allowbreak{}AutomorphicSpectrum.\allowbreak{}Proofs.\allowbreak{}DivisibilityProofs.\allowbreak{}TowerIdentityAtCutoff}, lines 155--190)\\
\texttt{GlobalTheory.\allowbreak{}CompanionTheory.\allowbreak{}Ar.\allowbreak{}Top.\allowbreak{}test\_\allowbreak{}height\_\allowbreak{}bound\_\allowbreak{}slope\_\allowbreak{}flat} (\texttt{Development.\allowbreak{}AutomorphicSpectrum.\allowbreak{}Proofs.\allowbreak{}DivisibilityProofs.\allowbreak{}TestHeightBound}, lines 32--142)\par}

{\footnotesize\raggedright\noindent \ref{B8.26c}.\quad \texttt{RootFieldNumerics.\allowbreak{}mean\_\allowbreak{}slope\_\allowbreak{}of\_\allowbreak{}proportional\_\allowbreak{}volume} (\texttt{Development.\allowbreak{}WeilBound.\allowbreak{}CoreStatements.\allowbreak{}FlagImmersions.\allowbreak{}MeanSlopeOfProportionalVolume}, lines 13--39)\par}

{\footnotesize\raggedright\noindent \ref{B8.26c}.\quad \texttt{GlobalTheory.\allowbreak{}CompanionTheory.\allowbreak{}Ar.\allowbreak{}Cusp.\allowbreak{}volumeTheoremX\_\allowbreak{}of\_\allowbreak{}inputs} (\texttt{Development.\allowbreak{}AutomorphicSpectrum.\allowbreak{}Proofs.\allowbreak{}DivisibilityProofs.\allowbreak{}VolumeTheoremOfInputs}, lines 72--121)\\
\texttt{GlobalTheory.\allowbreak{}CompanionTheory.\allowbreak{}Ar.\allowbreak{}Top.\allowbreak{}tower\_\allowbreak{}T2\_\allowbreak{}of\_\allowbreak{}inputs} (\texttt{Development.\allowbreak{}AutomorphicSpectrum.\allowbreak{}Proofs.\allowbreak{}DivisibilityProofs.\allowbreak{}TowerT2}, lines 69--164)\\
\texttt{RootFieldNumerics.\allowbreak{}rest\_\allowbreak{}sum\_\allowbreak{}total\_\allowbreak{}bound} (\texttt{Development.\allowbreak{}WeilBound.\allowbreak{}CoreStatements.\allowbreak{}FlagImmersions.\allowbreak{}TwoMulThreePowLeFactorial}, lines 112--163)\\
\texttt{SheetDecomposition.\allowbreak{}volEq\_\allowbreak{}of} (\texttt{Development.\allowbreak{}WeilBound.\allowbreak{}CoreStatements.\allowbreak{}VolumeEquality.\allowbreak{}VolEqOf}, lines 57--97)\par}

{\footnotesize\raggedright\noindent \ref{B8.27}.\quad \texttt{ReciprocityTheory.\allowbreak{}ChartRestriction.\allowbreak{}G.\allowbreak{}ArakelovGeometry.\allowbreak{}MoriwakiBigStatement} (\texttt{ReciprocityTheory.\allowbreak{}ReductiveGroups.\allowbreak{}ArithmeticGeometry.\allowbreak{}ArakelovGeometry.\allowbreak{}MoriwakiSt}, lines 29--37)\\
\texttt{PositivityTheory.\allowbreak{}Pf1.\allowbreak{}HilbertSamuelMoriwaki.\allowbreak{}moriwakiBig\_\allowbreak{}outright} (\texttt{Development.\allowbreak{}WeilBound.\allowbreak{}CoreStatements.\allowbreak{}AutomorphicKernels.\allowbreak{}MoriwakiBigness.\allowbreak{}MoriwakiBigTheorem.\allowbreak{}Outright}, lines 9--11)\\
\texttt{GroundTheory.\allowbreak{}Heights₂.\allowbreak{}Heights.\allowbreak{}HilbertSamuel.\allowbreak{}ArithmeticVolumeDegreeZero.\allowbreak{}proof} (\texttt{GroundTheory.\allowbreak{}HeightsAndTheta.\allowbreak{}Heights.\allowbreak{}HilbertSamuel.\allowbreak{}CoeffLeOfEvalLeMul}, lines 79--94)\\
\texttt{GroundTheory.\allowbreak{}Heights₂.\allowbreak{}Heights.\allowbreak{}Targets.\allowbreak{}MoriwakiBig.\allowbreak{}of\_\allowbreak{}found} (\texttt{GroundTheory.\allowbreak{}HeightsAndTheta.\allowbreak{}Heights.\allowbreak{}OrbitLimits.\allowbreak{}MoriwakiBig}, lines 11--16)\\
\texttt{UpperTheory.\allowbreak{}LiteratureTheorems.\allowbreak{}every\_\allowbreak{}character\_\allowbreak{}of\_\allowbreak{}literature\_\allowbreak{}of\_\allowbreak{}vWin} (\texttt{Development.\allowbreak{}AutomorphicArithmetic.\allowbreak{}ThetaCorrespondence.\allowbreak{}ThetaPairLifts.\allowbreak{}ThirdPlaceSeries.\allowbreak{}EveryCharacterOfLiteratureOfVWin}, lines 810--813)\par}

{\footnotesize\raggedright\noindent \ref{B8.27}.\quad \texttt{GlobalTheory.\allowbreak{}CovTop.\allowbreak{}lenT} (\texttt{Development.\allowbreak{}AbelianVarietyTheory.\allowbreak{}Definitions.\allowbreak{}Divisors.\allowbreak{}CovolXOfCovol}, lines 42--42)\\
\texttt{GlobalTheory.\allowbreak{}CovTop.\allowbreak{}zed} (\texttt{Development.\allowbreak{}AbelianVarietyTheory.\allowbreak{}Definitions.\allowbreak{}Divisors.\allowbreak{}CovolXOfCovol}, lines 44--45)\\
\texttt{GlobalTheory.\allowbreak{}CovTop.\allowbreak{}Quo} (\texttt{Development.\allowbreak{}AbelianVarietyTheory.\allowbreak{}Definitions.\allowbreak{}Divisors.\allowbreak{}CovolXOfCovol}, lines 49--49)\\
\texttt{GlobalTheory.\allowbreak{}CovTop.\allowbreak{}SizeAt} (\texttt{Development.\allowbreak{}AbelianVarietyTheory.\allowbreak{}Definitions.\allowbreak{}Divisors.\allowbreak{}CovolXOfCovol}, lines 75--79)\par}

{\footnotesize\raggedright\noindent \ref{B8.25p}.\quad \texttt{GlobalTheory.\allowbreak{}HtLowProof.\allowbreak{}htLow} (\texttt{Development.\allowbreak{}AbelianVarietyTheory.\allowbreak{}Proofs.\allowbreak{}DivisorProofs.\allowbreak{}MultiformHeightLowerBound}, lines 107--157)\\
\texttt{GlobalTheory.\allowbreak{}HSShort.\allowbreak{}deg\_\allowbreak{}ge\_\allowbreak{}ht\_\allowbreak{}det} (\texttt{Development.\allowbreak{}AbelianVarietyTheory.\allowbreak{}Proofs.\allowbreak{}DivisorProofs.\allowbreak{}LatticeDegreeGeHeightPolynomialVectorsDet}, lines 91--122)\par}

{\footnotesize\raggedright\noindent \ref{B8.25e}.\quad \texttt{GlobalTheory.\allowbreak{}CovTop.\allowbreak{}DegProof.\allowbreak{}covering\_\allowbreak{}degree\_\allowbreak{}eq\_\allowbreak{}alt} (\texttt{Development.\allowbreak{}AutomorphicArithmetic.\allowbreak{}ThetaCorrespondence.\allowbreak{}ThetaPairLifts.\allowbreak{}ThirdPlaceSeries.\allowbreak{}CoveringDegreeMultiplicative}, lines 25--55)\\
\texttt{GlobalTheory.\allowbreak{}CovTop.\allowbreak{}ShareRowsNef.\allowbreak{}share\_\allowbreak{}alt\_\allowbreak{}alt} (\texttt{Development.\allowbreak{}AutomorphicArithmetic.\allowbreak{}ThetaCorrespondence.\allowbreak{}ThetaPairLifts.\allowbreak{}ThirdPlaceSeries.\allowbreak{}ShareStatementOfKnownTheorems}, lines 30--66)\\
\texttt{GlobalTheory.\allowbreak{}CovTop.\allowbreak{}DegAt} (\texttt{Development.\allowbreak{}AbelianVarietyTheory.\allowbreak{}Definitions.\allowbreak{}Divisors.\allowbreak{}CovolXOfCovol}, lines 140--141)\\
\texttt{GlobalTheory.\allowbreak{}CovTop.\allowbreak{}ShareAt} (\texttt{Development.\allowbreak{}AbelianVarietyTheory.\allowbreak{}Definitions.\allowbreak{}Divisors.\allowbreak{}CovolXOfCovol}, lines 143--144)\\
\texttt{GlobalTheory.\allowbreak{}CovTop.\allowbreak{}HSAt} (\texttt{Development.\allowbreak{}AbelianVarietyTheory.\allowbreak{}Definitions.\allowbreak{}Divisors.\allowbreak{}CovolXOfCovol}, lines 153--156)\\
\texttt{GlobalTheory.\allowbreak{}CovTop.\allowbreak{}SumAt} (\texttt{Development.\allowbreak{}AbelianVarietyTheory.\allowbreak{}Definitions.\allowbreak{}Divisors.\allowbreak{}CovolXOfCovol}, lines 158--159)\par}

{\footnotesize\raggedright\noindent \ref{B8.25e}.\quad \texttt{GlobalTheory.\allowbreak{}CovTop.\allowbreak{}ShortRoad.\allowbreak{}hilbertSamuel\_\allowbreak{}of\_\allowbreak{}component\_\allowbreak{}estimates\_\allowbreak{}alt} (\texttt{Development.\allowbreak{}AutomorphicArithmetic.\allowbreak{}ThetaCorrespondence.\allowbreak{}ThetaPairLifts.\allowbreak{}ThirdPlaceSeries.\allowbreak{}HilbertSamuelLowerBoundOfComponents}, lines 25--113)\\
\texttt{GlobalTheory.\allowbreak{}CovHSNum.\allowbreak{}deg\_\allowbreak{}bound\_\allowbreak{}of\_\allowbreak{}claim\_\allowbreak{}numbers} (\texttt{Development.\allowbreak{}AbelianVarietyTheory.\allowbreak{}Proofs.\allowbreak{}DivisorProofs.\allowbreak{}DegBoundOfClaimNumbers}, lines 9--130)\\
\texttt{GlobalTheory.\allowbreak{}CovSum.\allowbreak{}sum\_\allowbreak{}deg\_\allowbreak{}div\_\allowbreak{}le\_\allowbreak{}mu} (\texttt{Development.\allowbreak{}AbelianVarietyTheory.\allowbreak{}Proofs.\allowbreak{}DivisorProofs.\allowbreak{}SumComponentDegreesDivLeLatticeSlope}, lines 61--77)\par}

{\footnotesize\raggedright\noindent \ref{B8.25f}.\quad \texttt{GlobalTheory.\allowbreak{}exists\_\allowbreak{}isShortGen\_\allowbreak{}of\_\allowbreak{}short\_\allowbreak{}span} (\texttt{Development.\allowbreak{}AbelianVarietyTheory.\allowbreak{}Proofs.\allowbreak{}DivisorProofs.\allowbreak{}ExistsShortGeneratorsOfShortSpanningFamily}, lines 13--153)\par}

{\footnotesize\raggedright\noindent \ref{B8.25q}.\quad \texttt{ReciprocityTheory.\allowbreak{}ChartRestriction.\allowbreak{}G.\allowbreak{}I2.\allowbreak{}exists\_\allowbreak{}maximal} (\texttt{ReciprocityTheory.\allowbreak{}ReductiveGroups.\allowbreak{}ArithmeticGeometry.\allowbreak{}SquareIntegrability.\allowbreak{}ExistsMaximal}, lines 196--223)\\
\texttt{ReciprocityTheory.\allowbreak{}ChartRestriction.\allowbreak{}G.\allowbreak{}I2.\allowbreak{}split\_\allowbreak{}pos} (\texttt{ReciprocityTheory.\allowbreak{}ReductiveGroups.\allowbreak{}ArithmeticGeometry.\allowbreak{}SquareIntegrability.\allowbreak{}SplitG}, lines 109--212)\\
\texttt{ReciprocityTheory.\allowbreak{}ChartRestriction.\allowbreak{}G.\allowbreak{}I2.\allowbreak{}cs\_\allowbreak{}bound} (\texttt{ReciprocityTheory.\allowbreak{}ReductiveGroups.\allowbreak{}ArithmeticGeometry.\allowbreak{}SquareIntegrability.\allowbreak{}SplitVB}, lines 125--245)\\
\texttt{ReciprocityTheory.\allowbreak{}ChartRestriction.\allowbreak{}G.\allowbreak{}I2.\allowbreak{}sum\_\allowbreak{}sq\_\allowbreak{}norm\_\allowbreak{}ge\_\allowbreak{}of\_\allowbreak{}span} (\texttt{ReciprocityTheory.\allowbreak{}ReductiveGroups.\allowbreak{}ArithmeticGeometry.\allowbreak{}SquareIntegrability.\allowbreak{}SplitV}, lines 32--160)\par}

{\footnotesize\raggedright\noindent \ref{B8.25r}.\quad \texttt{GlobalTheory.\allowbreak{}LeviZero.\allowbreak{}levi\_\allowbreak{}phla\_\allowbreak{}zero} (\texttt{Development.\allowbreak{}AbelianVarietyTheory.\allowbreak{}Proofs.\allowbreak{}TorusHaarMeasure.\allowbreak{}LeviFormWeightMetricOrthonormalDiagonal}, lines 417--471)\\
\texttt{GlobalTheory.\allowbreak{}LeviOriginMargin.\allowbreak{}levi\_\allowbreak{}margin\_\allowbreak{}origin} (\texttt{Development.\allowbreak{}AbelianVarietyTheory.\allowbreak{}Proofs.\allowbreak{}DivisorProofs.\allowbreak{}LeviMarginOrthonormalChartOrigin}, lines 99--116)\\
\texttt{GlobalTheory.\allowbreak{}frobenius\_\allowbreak{}lower\_\allowbreak{}comparison} (\texttt{Development.\allowbreak{}AbelianVarietyTheory.\allowbreak{}Proofs.\allowbreak{}DivisorProofs.\allowbreak{}FrobeniusStrictLowerComparisonNearIdentity}, lines 33--136)\\
\texttt{GlobalTheory.\allowbreak{}PeakWeight.\allowbreak{}levi\_\allowbreak{}lower} (\texttt{Development.\allowbreak{}AbelianVarietyTheory.\allowbreak{}Proofs.\allowbreak{}DivisorProofs.\allowbreak{}ExistsHormanderPeakCutoffSingularWeights}, lines 160--233)\\
\texttt{GlobalTheory.\allowbreak{}HormanderPeak.\allowbreak{}hormander\_\allowbreak{}peak} (\texttt{Development.\allowbreak{}AbelianVarietyTheory.\allowbreak{}Proofs.\allowbreak{}DivisorProofs.\allowbreak{}ExistsHormanderPeakOrthonormalFrameL2Bound}, lines 34--317)\par}

{\footnotesize\raggedright\noindent \ref{B8.25s}.\quad \texttt{GlobalTheory.\allowbreak{}BallSum.\allowbreak{}ballSum\_\allowbreak{}le} (\texttt{Development.\allowbreak{}AbelianVarietyTheory.\allowbreak{}Proofs.\allowbreak{}DivisorProofs.\allowbreak{}OrbitSumKernelBallMeasureLeIntegral}, lines 192--254)\\
\texttt{GlobalTheory.\allowbreak{}KernelIntegral.\allowbreak{}lintegral\_\allowbreak{}kaF\_\allowbreak{}le} (\texttt{Development.\allowbreak{}AbelianVarietyTheory.\allowbreak{}Proofs.\allowbreak{}DivisorProofs.\allowbreak{}LintegralUnitaryKernelLeExpPowEight}, lines 279--302)\\
\texttt{GlobalTheory.\allowbreak{}BallVol.\allowbreak{}ballvol} (\texttt{Development.\allowbreak{}AbelianVarietyTheory.\allowbreak{}Proofs.\allowbreak{}DivisorProofs.\allowbreak{}InvariantMeasureChartBallOrthonormalLowerBound}, lines 91--144)\\
\texttt{GlobalTheory.\allowbreak{}PoincareL2.\allowbreak{}poincare\_\allowbreak{}l2\_\allowbreak{}of\_\allowbreak{}orthogonal} (\texttt{Development.\allowbreak{}AbelianVarietyTheory.\allowbreak{}Proofs.\allowbreak{}DivisorProofs.\allowbreak{}PoincareL2InequalityFundamentalSet}, lines 87--216)\par}

{\footnotesize\raggedright\noindent \ref{B8.25t}.\quad \texttt{GlobalTheory.\allowbreak{}Dict.\allowbreak{}thInCusp\_\allowbreak{}cusp} (\texttt{Development.\allowbreak{}AbelianVarietyTheory.\allowbreak{}Proofs.\allowbreak{}SiegelThetaTransformation.\allowbreak{}ThInCuspCusp}, lines 24--61)\\
\texttt{GlobalTheory.\allowbreak{}ThFr.\allowbreak{}IsAdapted.\allowbreak{}flowCoeff\_\allowbreak{}eq\_\allowbreak{}zero} (\texttt{Development.\allowbreak{}AbelianVarietyTheory.\allowbreak{}Proofs.\allowbreak{}CuspFormEvaluations.\allowbreak{}IsHarmonicAndMemMinorSpan}, lines 591--630)\\
\texttt{GlobalTheory.\allowbreak{}Dict.\allowbreak{}CuspAny.\allowbreak{}mem\_\allowbreak{}SS\_\allowbreak{}of\_\allowbreak{}eq\_\allowbreak{}sum\_\allowbreak{}thFr} (\texttt{Development.\allowbreak{}AbelianVarietyTheory.\allowbreak{}Proofs.\allowbreak{}SiegelThetaTransformation.\allowbreak{}ExistsExpansionOfRowHighest}, lines 68--96)\par}

{\footnotesize\raggedright\noindent \ref{B8.25ga}.\quad \texttt{ReciprocityTheory.\allowbreak{}ChartRestriction.\allowbreak{}G.\allowbreak{}I2.\allowbreak{}PT} (\texttt{ReciprocityTheory.\allowbreak{}ReductiveGroups.\allowbreak{}ArithmeticGeometry.\allowbreak{}SquareIntegrability.\allowbreak{}GCauchyStatement}, lines 29--47)\\
\texttt{ReciprocityTheory.\allowbreak{}ChartRestriction.\allowbreak{}G.\allowbreak{}I2.\allowbreak{}PT.\allowbreak{}la1} (\texttt{ReciprocityTheory.\allowbreak{}ReductiveGroups.\allowbreak{}ArithmeticGeometry.\allowbreak{}SquareIntegrability.\allowbreak{}GCauchyStatement}, lines 64--65)\\
\texttt{ReciprocityTheory.\allowbreak{}ChartRestriction.\allowbreak{}G.\allowbreak{}I2.\allowbreak{}th} (\texttt{ReciprocityTheory.\allowbreak{}ReductiveGroups.\allowbreak{}ArithmeticGeometry.\allowbreak{}SquareIntegrability.\allowbreak{}GCauchyStatement}, lines 144--145)\\
\texttt{ReciprocityTheory.\allowbreak{}ChartRestriction.\allowbreak{}G.\allowbreak{}I2.\allowbreak{}Par} (\texttt{ReciprocityTheory.\allowbreak{}ReductiveGroups.\allowbreak{}ArithmeticGeometry.\allowbreak{}SquareIntegrability.\allowbreak{}GCauchyStatement}, lines 156--177)\\
\texttt{ReciprocityTheory.\allowbreak{}ChartRestriction.\allowbreak{}G.\allowbreak{}I2.\allowbreak{}phiFD} (\texttt{ReciprocityTheory.\allowbreak{}ReductiveGroups.\allowbreak{}ArithmeticGeometry.\allowbreak{}SquareIntegrability.\allowbreak{}TorusAt}, lines 20--22)\\
\texttt{GlobalTheory.\allowbreak{}EtaAn.\allowbreak{}RunOK} (\texttt{Development.\allowbreak{}AbelianVarietyTheory.\allowbreak{}Definitions.\allowbreak{}TorusSections.\allowbreak{}IsTorusOf}, lines 56--58)\\
\texttt{GlobalTheory.\allowbreak{}EtaAn.\allowbreak{}p3} (\texttt{Development.\allowbreak{}AbelianVarietyTheory.\allowbreak{}Definitions.\allowbreak{}TorusSections.\allowbreak{}IsTorusOf}, lines 61--64)\\
\texttt{GlobalTheory.\allowbreak{}EtaAn.\allowbreak{}p16} (\texttt{Development.\allowbreak{}AbelianVarietyTheory.\allowbreak{}Definitions.\allowbreak{}TorusSections.\allowbreak{}IsTorusOf}, lines 67--71)\par}

{\footnotesize\raggedright\noindent \ref{B8.25gb}.\quad \texttt{GlobalTheory.\allowbreak{}RunParametersOfLower.\allowbreak{}real\_\allowbreak{}bounds} (\texttt{Development.\allowbreak{}AbelianVarietyTheory.\allowbreak{}Proofs.\allowbreak{}CuspFormEvaluations.\allowbreak{}RealBounds}, lines 18--111)\\
\texttt{GlobalTheory.\allowbreak{}RunParametersOfLower.\allowbreak{}runOK\_\allowbreak{}of\_\allowbreak{}le} (\texttt{Development.\allowbreak{}AbelianVarietyTheory.\allowbreak{}Proofs.\allowbreak{}CuspFormEvaluations.\allowbreak{}RealBounds}, lines 112--162)\\
\texttt{GlobalTheory.\allowbreak{}ThetaSizeBounds.\allowbreak{}archimedeanBoundTotalConstant\_\allowbreak{}nonneg\_\allowbreak{}le\_\allowbreak{}bound} (\texttt{Development.\allowbreak{}AbelianVarietyTheory.\allowbreak{}Proofs.\allowbreak{}DivisorProofs.\allowbreak{}EtaBoundLengthUnit}, lines 26--71)\\
\texttt{GlobalTheory.\allowbreak{}EtaAn.\allowbreak{}archimedeanBoundTotalConstant\_\allowbreak{}size} (\texttt{Development.\allowbreak{}AbelianVarietyTheory.\allowbreak{}Proofs.\allowbreak{}TorusSplitting.\allowbreak{}ArchimedeanBoundTotalConstantSize}, lines 15--167)\\
\texttt{GlobalTheory.\allowbreak{}T.\allowbreak{}exists\_\allowbreak{}borel\_\allowbreak{}orthonormal\_\allowbreak{}flag\_\allowbreak{}frame} (\texttt{Development.\allowbreak{}AbelianVarietyTheory.\allowbreak{}Proofs.\allowbreak{}TorusSplitting.\allowbreak{}ExistsBorelOrthonormalFlagFrameStatement}, lines 11--13)\\
\texttt{GlobalTheory.\allowbreak{}EtaAn.\allowbreak{}exists\_\allowbreak{}unit\_\allowbreak{}frame} (\texttt{Development.\allowbreak{}AbelianVarietyTheory.\allowbreak{}Proofs.\allowbreak{}TorusSplitting.\allowbreak{}ExistsBorelOrthonormalFlagFrame}, lines 15--36)\\
\texttt{GlobalTheory.\allowbreak{}EtaAn.\allowbreak{}AHFrame} (\texttt{Development.\allowbreak{}AbelianVarietyTheory.\allowbreak{}Definitions.\allowbreak{}TorusSections.\allowbreak{}IsTorusOf}, lines 23--38)\\
\texttt{GlobalTheory.\allowbreak{}EtaAn.\allowbreak{}IsFnOf} (\texttt{Development.\allowbreak{}AbelianVarietyTheory.\allowbreak{}Definitions.\allowbreak{}TorusSections.\allowbreak{}IsTorusOf}, lines 39--44)\\
\texttt{GlobalTheory.\allowbreak{}T.\allowbreak{}exists\_\allowbreak{}torus\_\allowbreak{}of\_\allowbreak{}period\_\allowbreak{}hermitian} (\texttt{Development.\allowbreak{}AbelianVarietyTheory.\allowbreak{}Proofs.\allowbreak{}TorusSplitting.\allowbreak{}ExistsTorusOfPeriodHermitianStatement}, lines 11--13)\\
\texttt{GlobalTheory.\allowbreak{}T.\allowbreak{}periodTorus\_\allowbreak{}isTorusAt\_\allowbreak{}frame} (\texttt{Development.\allowbreak{}AbelianVarietyTheory.\allowbreak{}Proofs.\allowbreak{}TorusSplitting.\allowbreak{}PeriodTorusIsTorusAtFrameStatement}, lines 11--13)\\
\texttt{GlobalTheory.\allowbreak{}T.\allowbreak{}theta\_\allowbreak{}function\_\allowbreak{}torus\_\allowbreak{}l2\_\allowbreak{}norm} (\texttt{Development.\allowbreak{}AbelianVarietyTheory.\allowbreak{}Proofs.\allowbreak{}TorusSplitting.\allowbreak{}ThetaFunctionTorusL2NormStatement}, lines 11--13)\\
\texttt{GlobalTheory.\allowbreak{}T.\allowbreak{}fullLinearLevel\_\allowbreak{}eq\_\allowbreak{}exp\_\allowbreak{}periods\_\allowbreak{}complexPoint} (\texttt{Development.\allowbreak{}AbelianVarietyTheory.\allowbreak{}Proofs.\allowbreak{}TorusSplitting.\allowbreak{}FullLinearLevelEqExpPeriodsComplexPointStatement}, lines 11--13)\\
\texttt{GlobalTheory.\allowbreak{}T.\allowbreak{}norm\_\allowbreak{}nullVal\_\allowbreak{}eq\_\allowbreak{}norm\_\allowbreak{}theta\_\allowbreak{}alt} (\texttt{Development.\allowbreak{}AbelianVarietyTheory.\allowbreak{}Proofs.\allowbreak{}TorusSplitting.\allowbreak{}NormNullValEqNormThetaStatement}, lines 11--13)\\
\texttt{GlobalTheory.\allowbreak{}EtaAn.\allowbreak{}ArchPt.\allowbreak{}gen\_\allowbreak{}mod} (\texttt{Development.\allowbreak{}AbelianVarietyTheory.\allowbreak{}Proofs.\allowbreak{}TorusSplitting.\allowbreak{}NormEtaProductFormLeAtUnitFrame}, lines 36--54)\\
\texttt{GlobalTheory.\allowbreak{}nullValue\_\allowbreak{}existsUnique\_\allowbreak{}natural\_\allowbreak{}linear} (\texttt{Development.\allowbreak{}AbelianVarietyTheory.\allowbreak{}Proofs.\allowbreak{}CuspFormEvaluations.\allowbreak{}ThetaNullValueExistsUniqueNaturalLinear}, lines 159--169)\\
\texttt{GlobalTheory.\allowbreak{}ThetaConst.\allowbreak{}tauValAt\_\allowbreak{}self} (\texttt{Development.\allowbreak{}AbelianVarietyTheory.\allowbreak{}Definitions.\allowbreak{}Divisors.\allowbreak{}ExistsIsTauZeroNonzeroAll}, lines 114--120)\\
\texttt{GlobalTheory.\allowbreak{}ThetaConst.\allowbreak{}ArchBound.\allowbreak{}norm\_\allowbreak{}pow\_\allowbreak{}le\_\allowbreak{}at\_\allowbreak{}unit\_\allowbreak{}frame} (\texttt{Development.\allowbreak{}AbelianVarietyTheory.\allowbreak{}Proofs.\allowbreak{}DivisorProofs.\allowbreak{}NormSqThetaConstantSpanningClassesLeExp}, lines 46--148)\\
\texttt{GlobalTheory.\allowbreak{}ThetaConst.\allowbreak{}normSq\_\allowbreak{}thetaConst\_\allowbreak{}le\_\allowbreak{}exp} (\texttt{Development.\allowbreak{}AbelianVarietyTheory.\allowbreak{}Proofs.\allowbreak{}DivisorProofs.\allowbreak{}NormSqThetaConstantSpanningClassesLeExp}, lines 155--191)\par}

{\footnotesize\raggedright\noindent \ref{B8.25gc}.\quad \texttt{GlobalTheory.\allowbreak{}AllWt.\allowbreak{}bergE} (\texttt{Development.\allowbreak{}AbelianVarietyTheory.\allowbreak{}Definitions.\allowbreak{}Foundations.\allowbreak{}IsProbabilityMeasureMuPart}, lines 648--650)\\
\texttt{GlobalTheory.\allowbreak{}AllWt.\allowbreak{}bergOn} (\texttt{Development.\allowbreak{}AbelianVarietyTheory.\allowbreak{}Definitions.\allowbreak{}Foundations.\allowbreak{}IsProbabilityMeasureMuPart}, lines 753--754)\\
\texttt{GlobalTheory.\allowbreak{}ThinBerg.\allowbreak{}pnorm\_\allowbreak{}pos'} (\texttt{Development.\allowbreak{}AbelianVarietyTheory.\allowbreak{}Proofs.\allowbreak{}DivisorProofs.\allowbreak{}FrameExponentialVolumeLarge}, lines 22--31)\\
\texttt{GlobalTheory.\allowbreak{}ThinBerg.\allowbreak{}le\_\allowbreak{}bergOn\_\allowbreak{}of\_\allowbreak{}ratio} (\texttt{Development.\allowbreak{}AbelianVarietyTheory.\allowbreak{}Proofs.\allowbreak{}DivisorProofs.\allowbreak{}FrameExponentialVolumeLarge}, lines 32--60)\\
\texttt{GlobalTheory.\allowbreak{}BergLower.\allowbreak{}le\_\allowbreak{}bergOn} (\texttt{Development.\allowbreak{}AbelianVarietyTheory.\allowbreak{}Proofs.\allowbreak{}DivisorProofs.\allowbreak{}CuspFormLowerBoundBergmanFunctionOnComponent}, lines 59--153)\\
\texttt{GlobalTheory.\allowbreak{}T.\allowbreak{}cuspForm\_\allowbreak{}lower\_\allowbreak{}bound\_\allowbreak{}bergmanFunction\_\allowbreak{}on\_\allowbreak{}component} (\texttt{Development.\allowbreak{}AbelianVarietyTheory.\allowbreak{}Proofs.\allowbreak{}DivisorProofs.\allowbreak{}CuspFormLowerBoundBergmanFunctionOnComponentStatement}, lines 15--17)\par}

{\footnotesize\raggedright\noindent \ref{B8.25gca}.\quad \texttt{GlobalTheory.\allowbreak{}LevGeo.\allowbreak{}Laws} (\texttt{Development.\allowbreak{}AbelianVarietyTheory.\allowbreak{}Definitions.\allowbreak{}Divisors.\allowbreak{}UnifOfSpaceEq}, lines 20--42)\\
\texttt{UpperTheory.\allowbreak{}Dv.\allowbreak{}Thin.\allowbreak{}theta\_\allowbreak{}sup\_\allowbreak{}le} (\texttt{Development.\allowbreak{}AutomorphicArithmetic.\allowbreak{}LocalEmbeddings.\allowbreak{}ThinCore}, lines 26--106)\\
\texttt{GlobalTheory.\allowbreak{}Dict.\allowbreak{}ThetaPt.\allowbreak{}theta\_\allowbreak{}bound\_\allowbreak{}ortho} (\texttt{Development.\allowbreak{}AbelianVarietyTheory.\allowbreak{}Proofs.\allowbreak{}CuspFormEvaluations.\allowbreak{}ThetaGeneratorPointwiseBoundOrthonormalFrame}, lines 46--102)\par}

{\footnotesize\raggedright\noindent \ref{B8.25gcb}.\quad \texttt{GlobalTheory.\allowbreak{}LevGeo.\allowbreak{}Laws} (\texttt{Development.\allowbreak{}AbelianVarietyTheory.\allowbreak{}Definitions.\allowbreak{}Divisors.\allowbreak{}UnifOfSpaceEq}, lines 20--42)\\
\texttt{Formalism.\allowbreak{}L2.\allowbreak{}Pj.\allowbreak{}FromTest} (\texttt{Formalism.\allowbreak{}LocalToGlobal.\allowbreak{}ProjectiveLattices.\allowbreak{}AlgebraMapCAt}, lines 51--53)\\
\texttt{GlobalTheory.\allowbreak{}ThetaLower.\allowbreak{}thetaRule\_\allowbreak{}sum\_\allowbreak{}lower} (\texttt{Development.\allowbreak{}AbelianVarietyTheory.\allowbreak{}Proofs.\allowbreak{}CuspFormEvaluations.\allowbreak{}SumLeFinsumOfSurjectsOntoFinset}, lines 240--338)\\
\texttt{ReciprocityTheory.\allowbreak{}ChartRestriction.\allowbreak{}G.\allowbreak{}I2.\allowbreak{}Par.\allowbreak{}Jbound} (\texttt{ReciprocityTheory.\allowbreak{}ReductiveGroups.\allowbreak{}ArithmeticGeometry.\allowbreak{}SquareIntegrability.\allowbreak{}GCauchyStatement}, lines 182--186)\\
\texttt{GlobalTheory.\allowbreak{}BlockTh.\allowbreak{}p14} (\texttt{Development.\allowbreak{}AbelianVarietyTheory.\allowbreak{}Definitions.\allowbreak{}ThetaValues.\allowbreak{}IsBlockProd}, lines 41--44)\\
\texttt{GlobalTheory.\allowbreak{}AnVal.\allowbreak{}neg\_\allowbreak{}log\_\allowbreak{}sqrt\_\allowbreak{}le\_\allowbreak{}of\_\allowbreak{}jbound} (\texttt{Development.\allowbreak{}AbelianVarietyTheory.\allowbreak{}Proofs.\allowbreak{}CuspFormEvaluations.\allowbreak{}NegLogSqrtLeOfJbound}, lines 21--76)\\
\texttt{Formalism.\allowbreak{}L2.\allowbreak{}Pj.\allowbreak{}I2FramePointLemmas.\allowbreak{}jbound\_\allowbreak{}pos} (\texttt{Formalism.\allowbreak{}LocalToGlobal.\allowbreak{}ProjectiveLattices.\allowbreak{}ContDiffAtOfDiff}, lines 140--149)\\
\texttt{UpperTheory.\allowbreak{}Dv.\allowbreak{}Thin.\allowbreak{}theta\_\allowbreak{}lower} (\texttt{Development.\allowbreak{}AutomorphicArithmetic.\allowbreak{}LocalEmbeddings.\allowbreak{}ThinCore}, lines 134--200)\par}

{\footnotesize\raggedright\noindent \ref{B8.25g}.\quad \texttt{GlobalTheory.\allowbreak{}ThetaSizeBounds.\allowbreak{}run\_\allowbreak{}params\_\allowbreak{}alt} (\texttt{Development.\allowbreak{}AutomorphicArithmetic.\allowbreak{}ThetaCorrespondence.\allowbreak{}ThetaPairLifts.\allowbreak{}ThirdPlaceSeries.\allowbreak{}RunParams}, lines 15--166)\\
\texttt{GlobalTheory.\allowbreak{}ThetaConst.\allowbreak{}normSq\_\allowbreak{}thetaConst\_\allowbreak{}le\_\allowbreak{}exp} (\texttt{Development.\allowbreak{}AbelianVarietyTheory.\allowbreak{}Proofs.\allowbreak{}DivisorProofs.\allowbreak{}NormSqThetaConstantSpanningClassesLeExp}, lines 155--191)\\
\texttt{GlobalTheory.\allowbreak{}T.\allowbreak{}bergman\_\allowbreak{}kernel\_\allowbreak{}thin\_\allowbreak{}part\_\allowbreak{}estimate\_\allowbreak{}alt} (\texttt{Development.\allowbreak{}AutomorphicArithmetic.\allowbreak{}ThetaCorrespondence.\allowbreak{}ThetaPairLifts.\allowbreak{}ThirdPlaceSeries.\allowbreak{}BergmanKernelThinPartEstimate}, lines 58--59)\\
\texttt{GlobalTheory.\allowbreak{}ThickBerg.\allowbreak{}exp\_\allowbreak{}neg\_\allowbreak{}B\_\allowbreak{}le\_\allowbreak{}bergOn\_\allowbreak{}of\_\allowbreak{}mem\_\allowbreak{}framesOn} (\texttt{Development.\allowbreak{}AutomorphicArithmetic.\allowbreak{}ThetaCorrespondence.\allowbreak{}ThetaPairLifts.\allowbreak{}ThirdPlaceSeries.\allowbreak{}BergmanKernelThickPartEstimate}, lines 16--81)\par}

{\footnotesize\raggedright\noindent \ref{B8.28a}.\quad \texttt{GlobalTheory.\allowbreak{}ThetaSteep.\allowbreak{}tau\_\allowbreak{}bound} (\texttt{Development.\allowbreak{}AbelianVarietyTheory.\allowbreak{}Proofs.\allowbreak{}DivisorProofs.\allowbreak{}ExistsNonzeroSectionWeightBound}, lines 109--127)\\
\texttt{GlobalTheory.\allowbreak{}ThetaSteep.\allowbreak{}phiC\_\allowbreak{}le\_\allowbreak{}valuation} (\texttt{Development.\allowbreak{}AbelianVarietyTheory.\allowbreak{}Proofs.\allowbreak{}DivisorProofs.\allowbreak{}ExistsNonzeroSectionWeightBound}, lines 36--107)\par}

{\footnotesize\raggedright\noindent \ref{B8.28b}.\quad \texttt{GlobalTheory.\allowbreak{}SumBound.\allowbreak{}primeSum\_\allowbreak{}eOf} (\texttt{Development.\allowbreak{}AbelianVarietyTheory.\allowbreak{}Proofs.\allowbreak{}DivisorProofs.\allowbreak{}WeightedPrimeLogSumOscillationBoundLe}, lines 145--267)\\
\texttt{GlobalTheory.\allowbreak{}RefSum.\allowbreak{}sumAt\_\allowbreak{}of\_\allowbreak{}sizes} (\texttt{Development.\allowbreak{}AbelianVarietyTheory.\allowbreak{}Proofs.\allowbreak{}DivisorProofs.\allowbreak{}SumPointOfWeightNumerical}, lines 55--98)\\
\texttt{GlobalTheory.\allowbreak{}RefSum.\allowbreak{}ToWeight.\allowbreak{}primeSum\_\allowbreak{}add} (\texttt{Development.\allowbreak{}AbelianVarietyTheory.\allowbreak{}Proofs.\allowbreak{}DivisorProofs.\allowbreak{}ReferenceWeightOfTheta}, lines 51--70)\par}

{\footnotesize\raggedright\noindent \ref{B8.28c}.\quad \texttt{GlobalTheory.\allowbreak{}ShortPhi.\allowbreak{}exists\_\allowbreak{}shortBasis} (\texttt{Development.\allowbreak{}AbelianVarietyTheory.\allowbreak{}Proofs.\allowbreak{}DivisorProofs.\allowbreak{}ExistsShortPeriodLatticeBasisBalancedSlope}, lines 18--51)\\
\texttt{Formalism.\allowbreak{}MainNotions.\allowbreak{}ModuliIntegralCharts.\allowbreak{}H1.\allowbreak{}gram\_\allowbreak{}form\_\allowbreak{}phi\_\allowbreak{}holds} (\texttt{Formalism.\allowbreak{}SharpRestriction.\allowbreak{}Divisibility.\allowbreak{}GeneratorPowers.\allowbreak{}GramSpace}, lines 55--194)\par}

{\footnotesize\raggedright\noindent \ref{B8.28d}.\quad \texttt{GlobalTheory.\allowbreak{}SmallJets.\allowbreak{}small\_\allowbreak{}characters\_\allowbreak{}span\_\allowbreak{}jets} (\texttt{Development.\allowbreak{}AbelianVarietyTheory.\allowbreak{}Proofs.\allowbreak{}JetHarmonics.\allowbreak{}SmallCharactersSpanJets}, lines 425--507)\\
\texttt{GlobalTheory.\allowbreak{}UnifJets.\allowbreak{}jets} (\texttt{Development.\allowbreak{}AbelianVarietyTheory.\allowbreak{}Proofs.\allowbreak{}JetHarmonics.\allowbreak{}ExistsPowFrame}, lines 92--160)\par}

{\footnotesize\raggedright\noindent \ref{B8.28e}.\quad \texttt{GlobalTheory.\allowbreak{}RefAway.\allowbreak{}refExp\_\allowbreak{}le\_\allowbreak{}of\_\allowbreak{}oscLe} (\texttt{Development.\allowbreak{}AbelianVarietyTheory.\allowbreak{}Proofs.\allowbreak{}DivisorProofs.\allowbreak{}BaseLocusExponentLeOscillationNonTypePrime}, lines 125--230)\\
\texttt{GlobalTheory.\allowbreak{}RefSum.\allowbreak{}iInf\_\allowbreak{}thRule\_\allowbreak{}le\_\allowbreak{}nuLocOf} (\texttt{Development.\allowbreak{}AbelianVarietyTheory.\allowbreak{}Proofs.\allowbreak{}DivisorProofs.\allowbreak{}IInfThRuleLeNuLocOf}, lines 25--93)\\
\texttt{GlobalTheory.\allowbreak{}RefSum.\allowbreak{}refExp\_\allowbreak{}vp\_\allowbreak{}cover\_\allowbreak{}le} (\texttt{Development.\allowbreak{}AbelianVarietyTheory.\allowbreak{}Proofs.\allowbreak{}DivisorProofs.\allowbreak{}MemSecOfMemCuspWWtZ}, lines 30--65)\par}

{\footnotesize\raggedright\noindent \ref{B8.28f}.\quad \texttt{GlobalTheory.\allowbreak{}exists\_\allowbreak{}coset\_\allowbreak{}minimiser} (\texttt{Development.\allowbreak{}AbelianVarietyTheory.\allowbreak{}Proofs.\allowbreak{}JetHarmonics.\allowbreak{}ExistsPosZsmulMemRange}, lines 117--185)\\
\texttt{GlobalTheory.\allowbreak{}DepthBound.\allowbreak{}exists\_\allowbreak{}eigen\_\allowbreak{}column} (\texttt{Development.\allowbreak{}AbelianVarietyTheory.\allowbreak{}Proofs.\allowbreak{}JetHarmonics.\allowbreak{}ExistsEigenColumn}, lines 123--253)\par}

{\footnotesize\raggedright\noindent \ref{B8.28g}.\quad \texttt{GlobalTheory.\allowbreak{}FactorDepth.\allowbreak{}nuFac} (\texttt{Development.\allowbreak{}AbelianVarietyTheory.\allowbreak{}Definitions.\allowbreak{}Divisors.\allowbreak{}NonemptyValExtOfOverLa}, lines 276--278)\\
\texttt{GlobalTheory.\allowbreak{}FGDepth.\allowbreak{}detDepth} (\texttt{Development.\allowbreak{}AbelianVarietyTheory.\allowbreak{}Definitions.\allowbreak{}Foundations.\allowbreak{}IsSatFrame}, lines 49--51)\\
\texttt{GlobalTheory.\allowbreak{}RefSum.\allowbreak{}OscLe} (\texttt{Development.\allowbreak{}AbelianVarietyTheory.\allowbreak{}Definitions.\allowbreak{}Divisors.\allowbreak{}EightDvdNCol}, lines 33--35)\\
\texttt{Formalism.\allowbreak{}L2.\allowbreak{}Pj.\allowbreak{}harmA} (\texttt{Formalism.\allowbreak{}LocalToGlobal.\allowbreak{}ProjectiveLattices.\allowbreak{}IsThRule}, lines 86--88)\\
\texttt{GlobalTheory.\allowbreak{}FactorDepth.\allowbreak{}nuLocOf} (\texttt{Development.\allowbreak{}AbelianVarietyTheory.\allowbreak{}Definitions.\allowbreak{}Divisors.\allowbreak{}NonemptyValExtOfOverLa}, lines 260--265)\\
\texttt{Formalism.\allowbreak{}L2.\allowbreak{}Pj.\allowbreak{}nabValN} (\texttt{Formalism.\allowbreak{}LocalToGlobal.\allowbreak{}ProjectiveLattices.\allowbreak{}PlaceBelow}, lines 30--35)\\
\texttt{GlobalTheory.\allowbreak{}SlopeRec.\allowbreak{}refExp\_\allowbreak{}vp} (\texttt{Development.\allowbreak{}AbelianVarietyTheory.\allowbreak{}Definitions.\allowbreak{}Foundations.\allowbreak{}NonnegOfMemCStarSet}, lines 272--275)\par}

{\footnotesize\raggedright\noindent \ref{B8.28h}.\quad \texttt{GlobalTheory.\allowbreak{}FGDepth.\allowbreak{}detDepth\_\allowbreak{}bounds} (\texttt{Development.\allowbreak{}AbelianVarietyTheory.\allowbreak{}Proofs.\allowbreak{}DivisorProofs.\allowbreak{}FormalGroupDeterminantPowDepthBounds}, lines 29--45)\\
\texttt{GlobalTheory.\allowbreak{}FGDepth.\allowbreak{}detDepth\_\allowbreak{}eq\_\allowbreak{}zero\_\allowbreak{}of\_\allowbreak{}lt} (\texttt{Development.\allowbreak{}AbelianVarietyTheory.\allowbreak{}Proofs.\allowbreak{}DivisorProofs.\allowbreak{}FormalGroupDeterminantPowDepthEqZeroOfLtPrime}, lines 55--85)\\
\texttt{GlobalTheory.\allowbreak{}RefSum.\allowbreak{}facBetween\_\allowbreak{}zero\_\allowbreak{}boundAll} (\texttt{Development.\allowbreak{}AbelianVarietyTheory.\allowbreak{}Proofs.\allowbreak{}DivisorProofs.\allowbreak{}FactorDepthBetweenZeroAndBound}, lines 14--42)\par}

{\footnotesize\raggedright\noindent \ref{B8.28i}.\quad \texttt{GlobalTheory.\allowbreak{}RefSum.\allowbreak{}nuBetween\_\allowbreak{}sum\_\allowbreak{}and\_\allowbreak{}eLoc\_\allowbreak{}le} (\texttt{Development.\allowbreak{}AbelianVarietyTheory.\allowbreak{}Proofs.\allowbreak{}DivisorProofs.\allowbreak{}ProductDepthSumRoundingOscillation}, lines 16--47)\\
\texttt{GlobalTheory.\allowbreak{}RefSum.\allowbreak{}oscLe\_\allowbreak{}zero\_\allowbreak{}above\_\allowbreak{}and\_\allowbreak{}degree} (\texttt{Development.\allowbreak{}AbelianVarietyTheory.\allowbreak{}Proofs.\allowbreak{}DivisorProofs.\allowbreak{}DepthOscillationZeroAboveAndDegreeBound}, lines 78--100)\par}

{\footnotesize\raggedright\noindent \ref{B8.28j}.\quad \texttt{GlobalTheory.\allowbreak{}RefSum.\allowbreak{}Clear.\allowbreak{}clear} (\texttt{Development.\allowbreak{}AbelianVarietyTheory.\allowbreak{}Proofs.\allowbreak{}DivisorProofs.\allowbreak{}CuspFormClearDenominatorsAwayFromPlace}, lines 16--48)\\
\texttt{GlobalTheory.\allowbreak{}RefSum.\allowbreak{}ScalApprox.\allowbreak{}exists\_\allowbreak{}scal} (\texttt{Development.\allowbreak{}AbelianVarietyTheory.\allowbreak{}Proofs.\allowbreak{}DivisorProofs.\allowbreak{}ExistsScal}, lines 35--116)\par}

{\footnotesize\raggedright\noindent \ref{B8.28k}.\quad \texttt{GlobalTheory.\allowbreak{}IsoCols.\allowbreak{}exists\_\allowbreak{}isIsoCols} (\texttt{Development.\allowbreak{}AbelianVarietyTheory.\allowbreak{}Proofs.\allowbreak{}DivisorProofs.\allowbreak{}ExistsNormEqNegOne}, lines 194--243)\\
\texttt{GlobalTheory.\allowbreak{}IsoCols.\allowbreak{}exists\_\allowbreak{}pair} (\texttt{Development.\allowbreak{}AbelianVarietyTheory.\allowbreak{}Proofs.\allowbreak{}DivisorProofs.\allowbreak{}ExistsNormEqNegOne}, lines 97--154)\\
\texttt{GlobalTheory.\allowbreak{}IsoCols.\allowbreak{}exists\_\allowbreak{}norm\_\allowbreak{}eq\_\allowbreak{}neg\_\allowbreak{}one} (\texttt{Development.\allowbreak{}AbelianVarietyTheory.\allowbreak{}Proofs.\allowbreak{}DivisorProofs.\allowbreak{}ExistsNormEqNegOne}, lines 11--41)\\
\texttt{GlobalTheory.\allowbreak{}RefGlob.\allowbreak{}exists\_\allowbreak{}harmA\_\allowbreak{}congr\_\allowbreak{}polyT} (\texttt{Development.\allowbreak{}AbelianVarietyTheory.\allowbreak{}Proofs.\allowbreak{}DivisorProofs.\allowbreak{}ExistsGlobalHarmonicCongrColumnPolynomial}, lines 97--119)\\
\texttt{GlobalTheory.\allowbreak{}totalDegree\_\allowbreak{}le\_\allowbreak{}and\_\allowbreak{}exists\_\allowbreak{}combination\_\allowbreak{}of\_\allowbreak{}mem\_\allowbreak{}FFV} (\texttt{Development.\allowbreak{}AbelianVarietyTheory.\allowbreak{}Proofs.\allowbreak{}JetHarmonics.\allowbreak{}TotalDegreeLeAndExistsCombinationOfMemFFV}, lines 21--51)\par}

{\footnotesize\raggedright\noindent \ref{B8.28l}.\quad \texttt{GlobalTheory.\allowbreak{}RefSum.\allowbreak{}refExp\_\allowbreak{}vp\_\allowbreak{}cover\_\allowbreak{}le} (\texttt{Development.\allowbreak{}AbelianVarietyTheory.\allowbreak{}Proofs.\allowbreak{}DivisorProofs.\allowbreak{}MemSecOfMemCuspWWtZ}, lines 30--65)\\
\texttt{GlobalTheory.\allowbreak{}CoverDown.\allowbreak{}exists\_\allowbreak{}intFr\_\allowbreak{}below} (\texttt{Development.\allowbreak{}AbelianVarietyTheory.\allowbreak{}Proofs.\allowbreak{}DivisorProofs.\allowbreak{}ExistsIntFrBelow}, lines 94--106)\par}

{\footnotesize\raggedright\noindent \ref{B8.28m}.\quad \texttt{GlobalTheory.\allowbreak{}RefSum.\allowbreak{}LowP.\allowbreak{}thRule\_\allowbreak{}lower} (\texttt{Development.\allowbreak{}AbelianVarietyTheory.\allowbreak{}Proofs.\allowbreak{}DivisorProofs.\allowbreak{}ThRuleLower}, lines 26--81)\\
\texttt{GlobalTheory.\allowbreak{}RefSum.\allowbreak{}TermLoc.\allowbreak{}harmonic\_\allowbreak{}value\_\allowbreak{}ge\_\allowbreak{}local\_\allowbreak{}depth} (\texttt{Development.\allowbreak{}AbelianVarietyTheory.\allowbreak{}Proofs.\allowbreak{}DivisorProofs.\allowbreak{}HarmonicValueGeLocalDepthPrimeToP}, lines 24--63)\\
\texttt{GlobalTheory.\allowbreak{}RefSum.\allowbreak{}iInf\_\allowbreak{}thRule\_\allowbreak{}le\_\allowbreak{}nuLocOf} (\texttt{Development.\allowbreak{}AbelianVarietyTheory.\allowbreak{}Proofs.\allowbreak{}DivisorProofs.\allowbreak{}IInfThRuleLeNuLocOf}, lines 25--94)\par}

{\footnotesize\raggedright\noindent \ref{B8.28n}.\quad \texttt{GlobalTheory.\allowbreak{}RefAway.\allowbreak{}refExp\_\allowbreak{}le\_\allowbreak{}of\_\allowbreak{}oscLe} (\texttt{Development.\allowbreak{}AbelianVarietyTheory.\allowbreak{}Proofs.\allowbreak{}DivisorProofs.\allowbreak{}BaseLocusExponentLeOscillationNonTypePrime}, lines 125--230)\\
\texttt{GlobalTheory.\allowbreak{}thetaRules\_\allowbreak{}exist\_\allowbreak{}and\_\allowbreak{}cusp} (\texttt{Development.\allowbreak{}AbelianVarietyTheory.\allowbreak{}Proofs.\allowbreak{}DivisorProofs.\allowbreak{}ThetaRulesExistAndCusp}, lines 18--28)\par}

{\footnotesize\raggedright\noindent \ref{B8.28ua}.\quad \texttt{GlobalTheory.\allowbreak{}LimInv.\allowbreak{}limSub} (\texttt{Development.\allowbreak{}AbelianVarietyTheory.\allowbreak{}Proofs.\allowbreak{}DivisorProofs.\allowbreak{}LimitEntriesMatrixInverseClosedProof}, lines 9--28)\\
\texttt{GlobalTheory.\allowbreak{}LimInv.\allowbreak{}inv\_\allowbreak{}mem\_\allowbreak{}limSub} (\texttt{Development.\allowbreak{}AbelianVarietyTheory.\allowbreak{}Proofs.\allowbreak{}DivisorProofs.\allowbreak{}LimitEntriesMatrixInverseClosedProof}, lines 29--71)\\
\texttt{GlobalTheory.\allowbreak{}IsoTriv.\allowbreak{}limitEntriesMatrixInverseClosed\_\allowbreak{}proof} (\texttt{Development.\allowbreak{}AbelianVarietyTheory.\allowbreak{}Proofs.\allowbreak{}DivisorProofs.\allowbreak{}LimitEntriesMatrixInverseClosedProof}, lines 75--99)\par}

{\footnotesize\raggedright\noindent \ref{B8.28ub}.\quad \texttt{GlobalTheory.\allowbreak{}ColumnGerm.\allowbreak{}exists\_\allowbreak{}nabX\_\allowbreak{}eq} (\texttt{Development.\allowbreak{}AbelianVarietyTheory.\allowbreak{}Proofs.\allowbreak{}DivisorProofs.\allowbreak{}ExistsIntegralGermOnCopiesSameValue}, lines 47--87)\par}

{\footnotesize\raggedright\noindent \ref{B8.28uc}.\quad \texttt{GlobalTheory.\allowbreak{}IsoTriv.\allowbreak{}PrimeIdempotent.\allowbreak{}epsSt} (\texttt{Development.\allowbreak{}AbelianVarietyTheory.\allowbreak{}Proofs.\allowbreak{}DivisorProofs.\allowbreak{}PrimeIdempotentApproximationsDVRProof}, lines 40--82)\\
\texttt{GlobalTheory.\allowbreak{}IsoTriv.\allowbreak{}LimAct.\allowbreak{}cauchy\_\allowbreak{}eps\_\allowbreak{}mul} (\texttt{Development.\allowbreak{}AbelianVarietyTheory.\allowbreak{}Proofs.\allowbreak{}DivisorProofs.\allowbreak{}LimitActionLieAlgebraOfPrimeIdempotentApprox}, lines 107--135)\par}

{\footnotesize\raggedright\noindent \ref{B8.28ud}.\quad \texttt{GlobalTheory.\allowbreak{}IsoTriv.\allowbreak{}limActSt\_\allowbreak{}of\_\allowbreak{}epsSt} (\texttt{Development.\allowbreak{}AbelianVarietyTheory.\allowbreak{}Proofs.\allowbreak{}DivisorProofs.\allowbreak{}LimitActionLieAlgebraOfPrimeIdempotentApprox}, lines 144--194)\\
\texttt{GlobalTheory.\allowbreak{}IsoTriv.\allowbreak{}limPt\_\allowbreak{}of\_\allowbreak{}steps} (\texttt{Development.\allowbreak{}AbelianVarietyTheory.\allowbreak{}Proofs.\allowbreak{}DivisorProofs.\allowbreak{}LimitActionIntegralPointOfSteps}, lines 13--53)\par}

{\footnotesize\raggedright\noindent \ref{B8.28ue}.\quad \texttt{GlobalTheory.\allowbreak{}IsoTriv.\allowbreak{}limFrame} (\texttt{Development.\allowbreak{}AbelianVarietyTheory.\allowbreak{}Proofs.\allowbreak{}DivisorProofs.\allowbreak{}LimitIsotropicFrameIntegralOfSteps}, lines 44--114)\par}

{\footnotesize\raggedright\noindent \ref{B8.28u}.\quad \texttt{GlobalTheory.\allowbreak{}FactorDepth.\allowbreak{}nonempty\_\allowbreak{}valExt\_\allowbreak{}of\_\allowbreak{}overLa} (\texttt{Development.\allowbreak{}AbelianVarietyTheory.\allowbreak{}Definitions.\allowbreak{}Divisors.\allowbreak{}NonemptyValExtOfOverLa}, lines 320--322)\\
\texttt{GlobalTheory.\allowbreak{}ColApprox.\allowbreak{}exists\_\allowbreak{}germ\_\allowbreak{}cols} (\texttt{Development.\allowbreak{}AbelianVarietyTheory.\allowbreak{}Proofs.\allowbreak{}DivisorProofs.\allowbreak{}ExistsIntegralGermOnColumnsSameValue}, lines 90--167)\\
\texttt{GlobalTheory.\allowbreak{}ColApprox.\allowbreak{}exists\_\allowbreak{}germ\_\allowbreak{}close} (\texttt{Development.\allowbreak{}AbelianVarietyTheory.\allowbreak{}Proofs.\allowbreak{}DivisorProofs.\allowbreak{}SubstTwoPointsPowerSeriesZeroOfVanishFirstAxis}, lines 331--429)\\
\texttt{GlobalTheory.\allowbreak{}ColumnGerm.\allowbreak{}exists\_\allowbreak{}germ\_\allowbreak{}on\_\allowbreak{}copies} (\texttt{Development.\allowbreak{}AbelianVarietyTheory.\allowbreak{}Proofs.\allowbreak{}DivisorProofs.\allowbreak{}ExistsIntegralGermOnCopiesSameValue}, lines 97--116)\\
\texttt{GlobalTheory.\allowbreak{}ColumnGerm.\allowbreak{}exists\_\allowbreak{}nabX\_\allowbreak{}eq} (\texttt{Development.\allowbreak{}AbelianVarietyTheory.\allowbreak{}Proofs.\allowbreak{}DivisorProofs.\allowbreak{}ExistsIntegralGermOnCopiesSameValue}, lines 47--87)\par}

{\footnotesize\raggedright\noindent \ref{B8.28o}.\quad \texttt{GlobalTheory.\allowbreak{}ColDepth.\allowbreak{}nuLocOf\_\allowbreak{}eq\_\allowbreak{}nuP} (\texttt{Development.\allowbreak{}AbelianVarietyTheory.\allowbreak{}Proofs.\allowbreak{}DivisorProofs.\allowbreak{}ColumnPolynomialLocalDepthEqPlaceDepth}, lines 60--77)\\
\texttt{GlobalTheory.\allowbreak{}ColDepth.\allowbreak{}iInf\_\allowbreak{}eq\_\allowbreak{}of\_\allowbreak{}approx} (\texttt{Development.\allowbreak{}AbelianVarietyTheory.\allowbreak{}Proofs.\allowbreak{}DivisorProofs.\allowbreak{}ColumnPolynomialLocalDepthEqPlaceDepth}, lines 49--57)\\
\texttt{GlobalTheory.\allowbreak{}DepthDescent.\allowbreak{}nuLocOf\_\allowbreak{}eq\_\allowbreak{}nuLocX\_\allowbreak{}vp} (\texttt{Development.\allowbreak{}AbelianVarietyTheory.\allowbreak{}Proofs.\allowbreak{}DivisorProofs.\allowbreak{}LocalDepthOverDvrEqOverClosedField}, lines 19--96)\\
\texttt{GlobalTheory.\allowbreak{}ColApprox.\allowbreak{}exists\_\allowbreak{}germ\_\allowbreak{}cols} (\texttt{Development.\allowbreak{}AbelianVarietyTheory.\allowbreak{}Proofs.\allowbreak{}DivisorProofs.\allowbreak{}ExistsIntegralGermOnColumnsSameValue}, lines 90--167)\\
\texttt{GlobalTheory.\allowbreak{}ColumnGerm.\allowbreak{}exists\_\allowbreak{}germ\_\allowbreak{}on\_\allowbreak{}copies} (\texttt{Development.\allowbreak{}AbelianVarietyTheory.\allowbreak{}Proofs.\allowbreak{}DivisorProofs.\allowbreak{}ExistsIntegralGermOnCopiesSameValue}, lines 97--116)\\
\texttt{GlobalTheory.\allowbreak{}IsoFactorDepth.\allowbreak{}nuH\_\allowbreak{}vp\_\allowbreak{}eq\_\allowbreak{}nuP} (\texttt{Development.\allowbreak{}AbelianVarietyTheory.\allowbreak{}Proofs.\allowbreak{}DivisorProofs.\allowbreak{}IsotropicFactorDepthEqSumFactorDepths}, lines 47--84)\par}

{\footnotesize\raggedright\noindent \ref{B8.28va}.\quad \texttt{ReciprocityTheory.\allowbreak{}ChartRestriction.\allowbreak{}G.\allowbreak{}TB.\allowbreak{}harmSecond\_\allowbreak{}holds} (\texttt{ReciprocityTheory.\allowbreak{}ReductiveGroups.\allowbreak{}ArithmeticGeometry.\allowbreak{}TaylorDiscriminants.\allowbreak{}WickVanishing.\allowbreak{}HarmSecond}, lines 23--60)\\
\texttt{ReciprocityTheory.\allowbreak{}ChartRestriction.\allowbreak{}G.\allowbreak{}TB.\allowbreak{}pairMin} (\texttt{ReciprocityTheory.\allowbreak{}ReductiveGroups.\allowbreak{}ArithmeticGeometry.\allowbreak{}TaylorDiscriminants.\allowbreak{}GermSC}, lines 22--25)\\
\texttt{ReciprocityTheory.\allowbreak{}ChartRestriction.\allowbreak{}G.\allowbreak{}PullbackSections.\allowbreak{}Pf.\allowbreak{}sum\_\allowbreak{}pairMin\_\allowbreak{}le} (\texttt{ReciprocityTheory.\allowbreak{}ReductiveGroups.\allowbreak{}ArithmeticGeometry.\allowbreak{}PullbackSections.\allowbreak{}HolomorphicFamilies.\allowbreak{}ChartValueRamifiedOfLifts}, lines 110--127)\\
\texttt{ReciprocityTheory.\allowbreak{}ChartRestriction.\allowbreak{}G.\allowbreak{}PullbackSections.\allowbreak{}Pf.\allowbreak{}Jv} (\texttt{ReciprocityTheory.\allowbreak{}ReductiveGroups.\allowbreak{}ArithmeticGeometry.\allowbreak{}PullbackSections.\allowbreak{}HolomorphicFamilies.\allowbreak{}ChartValueRamifiedOfLifts}, lines 35--46)\\
\texttt{ReciprocityTheory.\allowbreak{}ChartRestriction.\allowbreak{}G.\allowbreak{}PullbackSections.\allowbreak{}Pf.\allowbreak{}mem\_\allowbreak{}Jv\_\allowbreak{}pow} (\texttt{ReciprocityTheory.\allowbreak{}ReductiveGroups.\allowbreak{}ArithmeticGeometry.\allowbreak{}PullbackSections.\allowbreak{}HolomorphicFamilies.\allowbreak{}ChartValueRamifiedOfLifts}, lines 56--81)\par}

{\footnotesize\raggedright\noindent \ref{B8.28vb}.\quad \texttt{ReciprocityTheory.\allowbreak{}ChartRestriction.\allowbreak{}G.\allowbreak{}TB.\allowbreak{}ordNabla} (\texttt{ReciprocityTheory.\allowbreak{}ReductiveGroups.\allowbreak{}ArithmeticGeometry.\allowbreak{}TaylorDiscriminants.\allowbreak{}GermSC}, lines 14--21)\\
\texttt{ReciprocityTheory.\allowbreak{}ChartRestriction.\allowbreak{}G.\allowbreak{}TB.\allowbreak{}LemmaHColsSt} (\texttt{ReciprocityTheory.\allowbreak{}ReductiveGroups.\allowbreak{}ArithmeticGeometry.\allowbreak{}TaylorDiscriminants.\allowbreak{}GermSCCols}, lines 24--38)\\
\texttt{ReciprocityTheory.\allowbreak{}ChartRestriction.\allowbreak{}G.\allowbreak{}TB.\allowbreak{}SCColsSt} (\texttt{ReciprocityTheory.\allowbreak{}ReductiveGroups.\allowbreak{}ArithmeticGeometry.\allowbreak{}TaylorDiscriminants.\allowbreak{}GermSCCols}, lines 10--23)\\
\texttt{ReciprocityTheory.\allowbreak{}ChartRestriction.\allowbreak{}G.\allowbreak{}TB.\allowbreak{}SCPf.\allowbreak{}ordNabla\_\allowbreak{}value\_\allowbreak{}expand} (\texttt{ReciprocityTheory.\allowbreak{}ReductiveGroups.\allowbreak{}ArithmeticGeometry.\allowbreak{}TaylorDiscriminants.\allowbreak{}WickVanishing.\allowbreak{}LemmaH}, lines 58--68)\\
\texttt{ReciprocityTheory.\allowbreak{}ChartRestriction.\allowbreak{}G.\allowbreak{}TB.\allowbreak{}WickVanishing.\allowbreak{}HCols.\allowbreak{}colAlg} (\texttt{ReciprocityTheory.\allowbreak{}ReductiveGroups.\allowbreak{}ArithmeticGeometry.\allowbreak{}TaylorDiscriminants.\allowbreak{}WickVanishing.\allowbreak{}LemmaHCols}, lines 153--208)\\
\texttt{ReciprocityTheory.\allowbreak{}ChartRestriction.\allowbreak{}G.\allowbreak{}TB.\allowbreak{}WickVanishing.\allowbreak{}HCols.\allowbreak{}lemmaHCols\_\allowbreak{}of} (\texttt{ReciprocityTheory.\allowbreak{}ReductiveGroups.\allowbreak{}ArithmeticGeometry.\allowbreak{}TaylorDiscriminants.\allowbreak{}WickVanishing.\allowbreak{}LemmaHCols}, lines 223--239)\\
\texttt{ReciprocityTheory.\allowbreak{}ChartRestriction.\allowbreak{}G.\allowbreak{}TB.\allowbreak{}SCPf.\allowbreak{}sc\_\allowbreak{}of} (\texttt{ReciprocityTheory.\allowbreak{}ReductiveGroups.\allowbreak{}ArithmeticGeometry.\allowbreak{}TaylorDiscriminants.\allowbreak{}WickVanishing.\allowbreak{}SCMain}, lines 28--84)\\
\texttt{ReciprocityTheory.\allowbreak{}ChartRestriction.\allowbreak{}G.\allowbreak{}TB.\allowbreak{}WickVanishing.\allowbreak{}Closed.\allowbreak{}scaled\_\allowbreak{}connection\_\allowbreak{}columns} (\texttt{ReciprocityTheory.\allowbreak{}ReductiveGroups.\allowbreak{}ArithmeticGeometry.\allowbreak{}TaylorDiscriminants.\allowbreak{}WickVanishing.\allowbreak{}Closed}, lines 47--63)\par}

{\footnotesize\raggedright\noindent \ref{B8.28vc}.\quad \texttt{ReciprocityTheory.\allowbreak{}ChartRestriction.\allowbreak{}G.\allowbreak{}TB.\allowbreak{}sgOf} (\texttt{ReciprocityTheory.\allowbreak{}ReductiveGroups.\allowbreak{}ArithmeticGeometry.\allowbreak{}TaylorDiscriminants.\allowbreak{}ConnectionKillsSeries}, lines 18--20)\\
\texttt{ReciprocityTheory.\allowbreak{}ChartRestriction.\allowbreak{}G.\allowbreak{}TB.\allowbreak{}sg1} (\texttt{ReciprocityTheory.\allowbreak{}ReductiveGroups.\allowbreak{}ArithmeticGeometry.\allowbreak{}TaylorDiscriminants.\allowbreak{}ConnectionKillsSeries}, lines 21--23)\\
\texttt{ReciprocityTheory.\allowbreak{}ChartRestriction.\allowbreak{}G.\allowbreak{}TB.\allowbreak{}Hset} (\texttt{ReciprocityTheory.\allowbreak{}ReductiveGroups.\allowbreak{}ArithmeticGeometry.\allowbreak{}TaylorDiscriminants.\allowbreak{}ConnectionKillsSeries}, lines 41--44)\\
\texttt{ReciprocityTheory.\allowbreak{}ChartRestriction.\allowbreak{}G.\allowbreak{}TB.\allowbreak{}WickVanishing.\allowbreak{}Nab.\allowbreak{}nabla\_\allowbreak{}eq} (\texttt{ReciprocityTheory.\allowbreak{}ReductiveGroups.\allowbreak{}ArithmeticGeometry.\allowbreak{}TaylorDiscriminants.\allowbreak{}WickVanishing.\allowbreak{}ConnectionKillsSeries}, lines 101--107)\\
\texttt{ReciprocityTheory.\allowbreak{}ChartRestriction.\allowbreak{}G.\allowbreak{}TB.\allowbreak{}WickVanishing.\allowbreak{}Nab.\allowbreak{}connection\_\allowbreak{}second\_\allowbreak{}of\_\allowbreak{}commute} (\texttt{ReciprocityTheory.\allowbreak{}ReductiveGroups.\allowbreak{}ArithmeticGeometry.\allowbreak{}TaylorDiscriminants.\allowbreak{}WickVanishing.\allowbreak{}ConnectionSecondOfCommute}, lines 38--49)\\
\texttt{ReciprocityTheory.\allowbreak{}ChartRestriction.\allowbreak{}G.\allowbreak{}TB.\allowbreak{}WickVanishing.\allowbreak{}Nab.\allowbreak{}connection\_\allowbreak{}kills\_\allowbreak{}series} (\texttt{ReciprocityTheory.\allowbreak{}ReductiveGroups.\allowbreak{}ArithmeticGeometry.\allowbreak{}TaylorDiscriminants.\allowbreak{}WickVanishing.\allowbreak{}ConnectionKillsSeries}, lines 117--121)\\
\texttt{ReciprocityTheory.\allowbreak{}ChartRestriction.\allowbreak{}G.\allowbreak{}TB.\allowbreak{}WickVanishing.\allowbreak{}Val.\allowbreak{}harmonic\_\allowbreak{}value\_\allowbreak{}integral} (\texttt{ReciprocityTheory.\allowbreak{}ReductiveGroups.\allowbreak{}ArithmeticGeometry.\allowbreak{}TaylorDiscriminants.\allowbreak{}WickVanishing.\allowbreak{}HarmonicValueIntegral}, lines 144--148)\\
\texttt{ReciprocityTheory.\allowbreak{}ChartRestriction.\allowbreak{}G.\allowbreak{}TB.\allowbreak{}WickVanishing.\allowbreak{}Nab.\allowbreak{}connection\_\allowbreak{}first\_\allowbreak{}identity} (\texttt{ReciprocityTheory.\allowbreak{}ReductiveGroups.\allowbreak{}ArithmeticGeometry.\allowbreak{}TaylorDiscriminants.\allowbreak{}WickVanishing.\allowbreak{}ExpConjConnectionIntegral}, lines 84--91)\par}

{\footnotesize\raggedright\noindent \ref{B8.28v}.\quad \texttt{GlobalTheory.\allowbreak{}PlaceInt.\allowbreak{}chBound\_\allowbreak{}zero\_\allowbreak{}of\_\allowbreak{}isUnit\_\allowbreak{}discr} (\texttt{Development.\allowbreak{}AbelianVarietyTheory.\allowbreak{}Proofs.\allowbreak{}RingsOfIntegers.\allowbreak{}HarmonicValueIntegralOfUnitDiscriminant}, lines 168--235)\\
\texttt{GlobalTheory.\allowbreak{}PlaceInt.\allowbreak{}ChartRam.\allowbreak{}chBound\_\allowbreak{}ramified} (\texttt{Development.\allowbreak{}AbelianVarietyTheory.\allowbreak{}Proofs.\allowbreak{}RingsOfIntegers.\allowbreak{}HarmonicValueBoundIntegralChartRamified}, lines 55--90)\\
\texttt{GlobalTheory.\allowbreak{}PlaceInt.\allowbreak{}host\_\allowbreak{}galois\_\allowbreak{}discr\_\allowbreak{}moves} (\texttt{Development.\allowbreak{}AbelianVarietyTheory.\allowbreak{}Proofs.\allowbreak{}RingsOfIntegers.\allowbreak{}HostFieldGaloisDegreeDiscrConjMovesPrimes}, lines 26--55)\\
\texttt{GlobalTheory.\allowbreak{}PlaceInt.\allowbreak{}HostChart.\allowbreak{}chBound\_\allowbreak{}host} (\texttt{Development.\allowbreak{}AbelianVarietyTheory.\allowbreak{}Proofs.\allowbreak{}RingsOfIntegers.\allowbreak{}HarmonicValuationLowerBoundHostFieldAllPlaces}, lines 36--70)\par}

{\footnotesize\raggedright\noindent \ref{B8.28p}.\quad \texttt{GlobalTheory.\allowbreak{}PlaceInt.\allowbreak{}HostChart.\allowbreak{}chBound\_\allowbreak{}host} (\texttt{Development.\allowbreak{}AbelianVarietyTheory.\allowbreak{}Proofs.\allowbreak{}RingsOfIntegers.\allowbreak{}HarmonicValuationLowerBoundHostFieldAllPlaces}, lines 36--70)\\
\texttt{GlobalTheory.\allowbreak{}PlaceInt.\allowbreak{}host\_\allowbreak{}galois\_\allowbreak{}discr\_\allowbreak{}moves} (\texttt{Development.\allowbreak{}AbelianVarietyTheory.\allowbreak{}Proofs.\allowbreak{}RingsOfIntegers.\allowbreak{}HostFieldGaloisDegreeDiscrConjMovesPrimes}, lines 26--55)\\
\texttt{GlobalTheory.\allowbreak{}PlaceInt.\allowbreak{}chBound\_\allowbreak{}zero\_\allowbreak{}of\_\allowbreak{}isUnit\_\allowbreak{}discr} (\texttt{Development.\allowbreak{}AbelianVarietyTheory.\allowbreak{}Proofs.\allowbreak{}RingsOfIntegers.\allowbreak{}HarmonicValueIntegralOfUnitDiscriminant}, lines 168--235)\\
\texttt{GlobalTheory.\allowbreak{}PlaceInt.\allowbreak{}ChartRam.\allowbreak{}chBound\_\allowbreak{}ramified} (\texttt{Development.\allowbreak{}AbelianVarietyTheory.\allowbreak{}Proofs.\allowbreak{}RingsOfIntegers.\allowbreak{}HarmonicValueBoundIntegralChartRamified}, lines 55--90)\\
\texttt{GlobalTheory.\allowbreak{}RefSum.\allowbreak{}GlobalPoly.\allowbreak{}exists\_\allowbreak{}polyBetween} (\texttt{Development.\allowbreak{}AbelianVarietyTheory.\allowbreak{}Proofs.\allowbreak{}DivisorProofs.\allowbreak{}ExistsHarmonicPolynomialLocalDepthBetween}, lines 62--93)\\
\texttt{GlobalTheory.\allowbreak{}RefSum.\allowbreak{}GlobalPoly.\allowbreak{}vV\_\allowbreak{}prime\_\allowbreak{}le} (\texttt{Development.\allowbreak{}AbelianVarietyTheory.\allowbreak{}Proofs.\allowbreak{}DivisorProofs.\allowbreak{}ExistsHarmonicPolynomialLocalDepthBetween}, lines 35--59)\\
\texttt{GlobalTheory.\allowbreak{}nuLocOf\_\allowbreak{}pert} (\texttt{Development.\allowbreak{}AbelianVarietyTheory.\allowbreak{}Proofs.\allowbreak{}DivisorProofs.\allowbreak{}LocalDepthStableUnderPrimePowPerturbation}, lines 63--84)\par}

{\footnotesize\raggedright\noindent \ref{B8.28wa}.\quad \texttt{GlobalTheory.\allowbreak{}PlaceInt.\allowbreak{}uu\_\allowbreak{}eq\_\allowbreak{}pull\_\allowbreak{}reach} (\texttt{Development.\allowbreak{}AbelianVarietyTheory.\allowbreak{}Proofs.\allowbreak{}RingsOfIntegers.\allowbreak{}MixedTranslatePullbackEqReachedSection}, lines 17--65)\\
\texttt{GlobalTheory.\allowbreak{}ReachRegular.\allowbreak{}regular\_\allowbreak{}of\_\allowbreak{}reach} (\texttt{Development.\allowbreak{}AbelianVarietyTheory.\allowbreak{}Proofs.\allowbreak{}RingsOfIntegers.\allowbreak{}ReachedSectionRegularOnIntegralModel}, lines 71--101)\\
\texttt{GlobalTheory.\allowbreak{}RegularGerm.\allowbreak{}regular\_\allowbreak{}germ} (\texttt{Development.\allowbreak{}AbelianVarietyTheory.\allowbreak{}Proofs.\allowbreak{}RingsOfIntegers.\allowbreak{}RegularSectionIntegralGermUniqueScale}, lines 250--253)\\
\texttt{GlobalTheory.\allowbreak{}PlaceInt.\allowbreak{}hv\_\allowbreak{}eq\_\allowbreak{}nabValN} (\texttt{Development.\allowbreak{}AbelianVarietyTheory.\allowbreak{}Proofs.\allowbreak{}RingsOfIntegers.\allowbreak{}HarmonicValueIntegralPointEqValueOnBiextensionCo}, lines 76--94)\\
\texttt{GlobalTheory.\allowbreak{}PlaceInt.\allowbreak{}hv\_\allowbreak{}pulls} (\texttt{Development.\allowbreak{}AbelianVarietyTheory.\allowbreak{}Proofs.\allowbreak{}RingsOfIntegers.\allowbreak{}HarmonicValueChainRuleNormalizedPullback}, lines 15--54)\\
\texttt{GlobalTheory.\allowbreak{}PlaceInt.\allowbreak{}TermBound.\allowbreak{}term\_\allowbreak{}bound} (\texttt{Development.\allowbreak{}AbelianVarietyTheory.\allowbreak{}Proofs.\allowbreak{}RingsOfIntegers.\allowbreak{}HarmonicValueTranslatedProductValuationBound}, lines 76--124)\par}

{\footnotesize\raggedright\noindent \ref{B8.28w}.\quad \texttt{GlobalTheory.\allowbreak{}ThetaBound.\allowbreak{}bound\_\allowbreak{}of\_\allowbreak{}chBound} (\texttt{Development.\allowbreak{}AbelianVarietyTheory.\allowbreak{}Proofs.\allowbreak{}RingsOfIntegers.\allowbreak{}ThetaGeneratorPlaceValuationBound}, lines 88--149)\\
\texttt{GlobalTheory.\allowbreak{}TorsFac.\allowbreak{}factorisation} (\texttt{Development.\allowbreak{}AbelianVarietyTheory.\allowbreak{}Proofs.\allowbreak{}RingsOfIntegers.\allowbreak{}DvdOfCastDvd}, lines 91--172)\\
\texttt{GlobalTheory.\allowbreak{}T.\allowbreak{}harmonicValue\_\allowbreak{}translated\_\allowbreak{}product\_\allowbreak{}valuation\_\allowbreak{}bound} (\texttt{Development.\allowbreak{}AbelianVarietyTheory.\allowbreak{}Proofs.\allowbreak{}RingsOfIntegers.\allowbreak{}HarmonicValueTranslatedProductValuationBoundStatement}, lines 11--13)\par}

{\footnotesize\raggedright\noindent \ref{B8.28q}.\quad \texttt{GlobalTheory.\allowbreak{}ThetaLow.\allowbreak{}thetaRule\_\allowbreak{}bounded\_\allowbreak{}below} (\texttt{Development.\allowbreak{}AbelianVarietyTheory.\allowbreak{}Proofs.\allowbreak{}DivisorProofs.\allowbreak{}ThetaRuleBoundedBelow}, lines 43--121)\\
\texttt{GlobalTheory.\allowbreak{}ThetaBound.\allowbreak{}bound\_\allowbreak{}of\_\allowbreak{}chBound} (\texttt{Development.\allowbreak{}AbelianVarietyTheory.\allowbreak{}Proofs.\allowbreak{}RingsOfIntegers.\allowbreak{}ThetaGeneratorPlaceValuationBound}, lines 88--149)\\
\texttt{GlobalTheory.\allowbreak{}ThetaBound.\allowbreak{}Prep.\allowbreak{}exists\_\allowbreak{}good\_\allowbreak{}chart\_\allowbreak{}allRat} (\texttt{Development.\allowbreak{}AbelianVarietyTheory.\allowbreak{}Proofs.\allowbreak{}RingsOfIntegers.\allowbreak{}ExistsExtAllModelsUniformised}, lines 318--336)\\
\texttt{GlobalTheory.\allowbreak{}TorsFac.\allowbreak{}factorisation} (\texttt{Development.\allowbreak{}AbelianVarietyTheory.\allowbreak{}Proofs.\allowbreak{}RingsOfIntegers.\allowbreak{}DvdOfCastDvd}, lines 91--172)\\
\texttt{GlobalTheory.\allowbreak{}SplitLevel.\allowbreak{}level\_\allowbreak{}halfSplit} (\texttt{Development.\allowbreak{}AbelianVarietyTheory.\allowbreak{}Proofs.\allowbreak{}DivisorProofs.\allowbreak{}LevelHalfSplit}, lines 15--24)\\
\texttt{GlobalTheory.\allowbreak{}RefLevel.\allowbreak{}coprime\_\allowbreak{}level} (\texttt{Development.\allowbreak{}AbelianVarietyTheory.\allowbreak{}Proofs.\allowbreak{}DivisorProofs.\allowbreak{}SupLatticeReferenceExponentThetaWeightLevelLeOne}, lines 66--70)\par}

{\footnotesize\raggedright\noindent \ref{B8.28q}.\quad \texttt{ReciprocityTheory.\allowbreak{}ChartRestriction.\allowbreak{}G.\allowbreak{}Geo.\allowbreak{}Tor.\allowbreak{}Datum.\allowbreak{}mult} (\texttt{ReciprocityTheory.\allowbreak{}ReductiveGroups.\allowbreak{}ArithmeticGeometry.\allowbreak{}BirationalGeometry.\allowbreak{}ToroidalCompactification.\allowbreak{}FlagBasis}, lines 76--77)\par}

{\footnotesize\raggedright\noindent \ref{B8.28xe}.\quad \texttt{UpperTheory.\allowbreak{}Depth.\allowbreak{}SplitPlaces} (\texttt{Development.\allowbreak{}AutomorphicArithmetic.\allowbreak{}CMTraceForms.\allowbreak{}Places}, lines 19--30)\\
\texttt{GlobalTheory.\allowbreak{}RefSum.\allowbreak{}oscLe\_\allowbreak{}eOf} (\texttt{Development.\allowbreak{}AbelianVarietyTheory.\allowbreak{}Proofs.\allowbreak{}DivisorProofs.\allowbreak{}DepthOscillationLeBoundEveryPrime}, lines 17--62)\\
\texttt{GlobalTheory.\allowbreak{}T.\allowbreak{}depth\_\allowbreak{}oscillation\_\allowbreak{}le\_\allowbreak{}at\_\allowbreak{}two\_\allowbreak{}cmField} (\texttt{Development.\allowbreak{}AbelianVarietyTheory.\allowbreak{}Proofs.\allowbreak{}DivisorProofs.\allowbreak{}DepthOscillationLeAtTwoCmFieldStatement}, lines 11--13)\\
\texttt{GlobalTheory.\allowbreak{}RefSum.\allowbreak{}oscLe\_\allowbreak{}two} (\texttt{Development.\allowbreak{}AbelianVarietyTheory.\allowbreak{}Proofs.\allowbreak{}DivisorProofs.\allowbreak{}DepthOscillationLeAtTwoCmField}, lines 27--97)\\
\texttt{GlobalTheory.\allowbreak{}OscSplit.\allowbreak{}oscLe\_\allowbreak{}split} (\texttt{Development.\allowbreak{}AbelianVarietyTheory.\allowbreak{}Proofs.\allowbreak{}DivisorProofs.\allowbreak{}DepthOscillationLeOddSplitUnramifiedPrime}, lines 57--97)\\
\texttt{GlobalTheory.\allowbreak{}OscSplit.\allowbreak{}sumBig\_\allowbreak{}le} (\texttt{Development.\allowbreak{}AbelianVarietyTheory.\allowbreak{}Proofs.\allowbreak{}DivisorProofs.\allowbreak{}DepthOscillationLeOddSplitUnramifiedPrime}, lines 39--56)\par}

{\footnotesize\raggedright\noindent \ref{B8.28xfa}.\quad \texttt{GlobalTheory.\allowbreak{}T.\allowbreak{}frobenius\_\allowbreak{}orbit\_\allowbreak{}of\_\allowbreak{}unramified\_\allowbreak{}prime} (\texttt{Development.\allowbreak{}AbelianVarietyTheory.\allowbreak{}Proofs.\allowbreak{}DivisorProofs.\allowbreak{}FrobeniusOrbitOfUnramifiedPrimeStatement}, lines 10--12)\\
\texttt{GlobalTheory.\allowbreak{}ChF.\allowbreak{}unramified\_\allowbreak{}companion\_\allowbreak{}chart\_\allowbreak{}characters} (\texttt{Development.\allowbreak{}AbelianVarietyTheory.\allowbreak{}Proofs.\allowbreak{}UnramifiedDepth.\allowbreak{}UnramifiedCompanionChartCharacters}, lines 32--76)\\
\texttt{GlobalTheory.\allowbreak{}CharMinor.\allowbreak{}colength\_\allowbreak{}bound} (\texttt{Development.\allowbreak{}AbelianVarietyTheory.\allowbreak{}Proofs.\allowbreak{}JetHarmonics.\allowbreak{}ExistsCharactersBoundedMinorEigenspace}, lines 82--152)\\
\texttt{GlobalTheory.\allowbreak{}RefSum.\allowbreak{}hasChars\_\allowbreak{}frame\_\allowbreak{}of\_\allowbreak{}hasAF} (\texttt{Development.\allowbreak{}AbelianVarietyTheory.\allowbreak{}Proofs.\allowbreak{}DivisorProofs.\allowbreak{}FlagStepCharactersOfColength}, lines 20--67)\par}

{\footnotesize\raggedright\noindent \ref{B8.28xf}.\quad \texttt{GlobalTheory.\allowbreak{}T.\allowbreak{}factorDepth\_\allowbreak{}upperBound\_\allowbreak{}oddUnramifiedPrime} (\texttt{Development.\allowbreak{}AbelianVarietyTheory.\allowbreak{}Proofs.\allowbreak{}DivisorProofs.\allowbreak{}FactorDepthUpperBoundOddUnramifiedPrimeStatement}, lines 10--12)\\
\texttt{GlobalTheory.\allowbreak{}RefSum.\allowbreak{}facBetween\_\allowbreak{}zero\_\allowbreak{}boundAll} (\texttt{Development.\allowbreak{}AbelianVarietyTheory.\allowbreak{}Proofs.\allowbreak{}DivisorProofs.\allowbreak{}FactorDepthBetweenZeroAndBound}, lines 14--42)\par}

{\footnotesize\raggedright\noindent \ref{B8.28xga}.\quad \texttt{GlobalTheory.\allowbreak{}T.\allowbreak{}formalGroup\_\allowbreak{}rootsOfUnity\_\allowbreak{}typed\_\allowbreak{}lieRanks} (\texttt{Development.\allowbreak{}AbelianVarietyTheory.\allowbreak{}Proofs.\allowbreak{}DivisorProofs.\allowbreak{}FormalGroupRootsOfUnityTypedLieRanksStatement}, lines 5--7)\\
\texttt{GlobalTheory.\allowbreak{}PrimePart.\allowbreak{}formalGroupTypes} (\texttt{Development.\allowbreak{}AbelianVarietyTheory.\allowbreak{}Proofs.\allowbreak{}DivisorProofs.\allowbreak{}FormalGroupRootsOfUnityTypedLieRanks}, lines 221--223)\par}

{\footnotesize\raggedright\noindent \ref{B8.28xgb}.\quad \texttt{GlobalTheory.\allowbreak{}RefSum.\allowbreak{}facBetween\_\allowbreak{}fine} (\texttt{Development.\allowbreak{}AbelianVarietyTheory.\allowbreak{}Proofs.\allowbreak{}DivisorProofs.\allowbreak{}FactorDepthBoundedOddInertPrimeCmType}, lines 28--84)\par}

{\footnotesize\raggedright\noindent \ref{B8.28xg}.\quad \texttt{GlobalTheory.\allowbreak{}IndefiniteRowSideFullHecke.\allowbreak{}FactorDepthInertBounds.\allowbreak{}boundFine} (\texttt{Development.\allowbreak{}AbelianVarietyTheory.\allowbreak{}BlockTheory.\allowbreak{}DivisorBase.\allowbreak{}FactorDepthInertBounds}, lines 38--41)\\
\texttt{GlobalTheory.\allowbreak{}RefSum.\allowbreak{}oscLe\_\allowbreak{}eSmall\_\allowbreak{}of\_\allowbreak{}inert} (\texttt{Development.\allowbreak{}AbelianVarietyTheory.\allowbreak{}Proofs.\allowbreak{}DivisorProofs.\allowbreak{}DepthOscillationLeAtSmallOddInertPrime}, lines 180--195)\\
\texttt{GlobalTheory.\allowbreak{}RefSum.\allowbreak{}FacBetween} (\texttt{Development.\allowbreak{}AbelianVarietyTheory.\allowbreak{}Definitions.\allowbreak{}Divisors.\allowbreak{}EightDvdNCol}, lines 21--25)\par}

{\footnotesize\raggedright\noindent \ref{B8.28xha}.\quad \texttt{GlobalTheory.\allowbreak{}detPow\_\allowbreak{}content\_\allowbreak{}const} (\texttt{Development.\allowbreak{}AbelianVarietyTheory.\allowbreak{}Proofs.\allowbreak{}DivisorProofs.\allowbreak{}DividedPowerContentDetPowDependsOnPIK}, lines 51--77)\\
\texttt{GlobalTheory.\allowbreak{}T.\allowbreak{}formalGroupDeterminantPowDepth\_\allowbreak{}bounds} (\texttt{Development.\allowbreak{}AbelianVarietyTheory.\allowbreak{}Proofs.\allowbreak{}DivisorProofs.\allowbreak{}FormalGroupDeterminantPowDepthBoundsStatement}, lines 8--10)\par}

{\footnotesize\raggedright\noindent \ref{B8.28xhb}.\quad \texttt{GlobalTheory.\allowbreak{}FGDepth.\allowbreak{}InertBound.\allowbreak{}content\_\allowbreak{}le\_\allowbreak{}detDepth\_\allowbreak{}add} (\texttt{Development.\allowbreak{}AbelianVarietyTheory.\allowbreak{}Proofs.\allowbreak{}DivisorProofs.\allowbreak{}ContentLeDetDepthAdd}, lines 151--300)\\
\texttt{GlobalTheory.\allowbreak{}T.\allowbreak{}exists\_\allowbreak{}typicalCurve\_\allowbreak{}basis\_\allowbreak{}adapted\_\allowbreak{}to\_\allowbreak{}saturated\_\allowbreak{}frame} (\texttt{Development.\allowbreak{}AbelianVarietyTheory.\allowbreak{}Proofs.\allowbreak{}DivisorProofs.\allowbreak{}ExistsTypicalCurveBasisAdaptedToSaturatedFrameStatement}, lines 8--10)\\
\texttt{ReciprocityTheory.\allowbreak{}ChartRestriction.\allowbreak{}G.\allowbreak{}Sh.\allowbreak{}pt5v'} (\texttt{ReciprocityTheory.\allowbreak{}ReductiveGroups.\allowbreak{}ArithmeticGeometry.\allowbreak{}WeightedSumBounds.\allowbreak{}WtSubOneLe}, lines 99--190)\par}

{\footnotesize\raggedright\noindent \ref{B8.28xh}.\quad \texttt{GlobalTheory.\allowbreak{}T.\allowbreak{}factorDepth\_\allowbreak{}bounds\_\allowbreak{}inert\_\allowbreak{}prime} (\texttt{Development.\allowbreak{}AbelianVarietyTheory.\allowbreak{}Proofs.\allowbreak{}DivisorProofs.\allowbreak{}FactorDepthBoundsInertPrimeStatement}, lines 10--12)\\
\texttt{GlobalTheory.\allowbreak{}IndefiniteRowSideFullHecke.\allowbreak{}FactorDepthInertBounds.\allowbreak{}boundInert} (\texttt{Development.\allowbreak{}AbelianVarietyTheory.\allowbreak{}BlockTheory.\allowbreak{}DivisorBase.\allowbreak{}FactorDepthInertBounds}, lines 28--29)\\
\texttt{GlobalTheory.\allowbreak{}RefSum.\allowbreak{}InertOsc.\allowbreak{}big\_\allowbreak{}arith} (\texttt{Development.\allowbreak{}AbelianVarietyTheory.\allowbreak{}Proofs.\allowbreak{}DivisorProofs.\allowbreak{}DepthOscillationLargeInertPrime}, lines 32--37)\\
\texttt{GlobalTheory.\allowbreak{}RefSum.\allowbreak{}InertOsc.\allowbreak{}big\_\allowbreak{}bound} (\texttt{Development.\allowbreak{}AbelianVarietyTheory.\allowbreak{}Proofs.\allowbreak{}DivisorProofs.\allowbreak{}DepthOscillationLargeInertPrime}, lines 38--55)\par}

{\footnotesize\raggedright\noindent \ref{B8.28x}.\quad \texttt{GlobalTheory.\allowbreak{}RefSum.\allowbreak{}oscLe\_\allowbreak{}eOf} (\texttt{Development.\allowbreak{}AbelianVarietyTheory.\allowbreak{}Proofs.\allowbreak{}DivisorProofs.\allowbreak{}DepthOscillationLeBoundEveryPrime}, lines 17--62)\\
\texttt{GlobalTheory.\allowbreak{}OscSplit.\allowbreak{}oscLe\_\allowbreak{}split} (\texttt{Development.\allowbreak{}AbelianVarietyTheory.\allowbreak{}Proofs.\allowbreak{}DivisorProofs.\allowbreak{}DepthOscillationLeOddSplitUnramifiedPrime}, lines 57--97)\\
\texttt{GlobalTheory.\allowbreak{}T.\allowbreak{}factorDepth\_\allowbreak{}upperBound\_\allowbreak{}oddUnramifiedPrime} (\texttt{Development.\allowbreak{}AbelianVarietyTheory.\allowbreak{}Proofs.\allowbreak{}DivisorProofs.\allowbreak{}FactorDepthUpperBoundOddUnramifiedPrimeStatement}, lines 10--12)\\
\texttt{GlobalTheory.\allowbreak{}RefSum.\allowbreak{}oscLe\_\allowbreak{}eSmall\_\allowbreak{}of\_\allowbreak{}inert} (\texttt{Development.\allowbreak{}AbelianVarietyTheory.\allowbreak{}Proofs.\allowbreak{}DivisorProofs.\allowbreak{}DepthOscillationLeAtSmallOddInertPrime}, lines 180--195)\\
\texttt{GlobalTheory.\allowbreak{}RefSum.\allowbreak{}OscInert.\allowbreak{}sum\_\allowbreak{}bound\_\allowbreak{}le} (\texttt{Development.\allowbreak{}AbelianVarietyTheory.\allowbreak{}Proofs.\allowbreak{}DivisorProofs.\allowbreak{}DepthOscillationLeAtSmallOddInertPrime}, lines 112--175)\\
\texttt{GlobalTheory.\allowbreak{}RefSum.\allowbreak{}InertOsc.\allowbreak{}oscLe\_\allowbreak{}inert\_\allowbreak{}large} (\texttt{Development.\allowbreak{}AbelianVarietyTheory.\allowbreak{}Proofs.\allowbreak{}DivisorProofs.\allowbreak{}DepthOscillationLargeInertPrime}, lines 62--118)\\
\texttt{GlobalTheory.\allowbreak{}T.\allowbreak{}factorDepth\_\allowbreak{}bounds\_\allowbreak{}inert\_\allowbreak{}prime} (\texttt{Development.\allowbreak{}AbelianVarietyTheory.\allowbreak{}Proofs.\allowbreak{}DivisorProofs.\allowbreak{}FactorDepthBoundsInertPrimeStatement}, lines 10--12)\\
\texttt{GlobalTheory.\allowbreak{}RefSum.\allowbreak{}InertOsc.\allowbreak{}big\_\allowbreak{}bound} (\texttt{Development.\allowbreak{}AbelianVarietyTheory.\allowbreak{}Proofs.\allowbreak{}DivisorProofs.\allowbreak{}DepthOscillationLargeInertPrime}, lines 38--55)\par}

{\footnotesize\raggedright\noindent \ref{B8.28r}.\quad \texttt{GlobalTheory.\allowbreak{}RefSum.\allowbreak{}ThetaExp.\allowbreak{}expBd\_\allowbreak{}thA} (\texttt{Development.\allowbreak{}AbelianVarietyTheory.\allowbreak{}Proofs.\allowbreak{}DivisorProofs.\allowbreak{}ThetaExponentsBigSmallFactorBounds}, lines 234--252)\\
\texttt{GlobalTheory.\allowbreak{}RefSum.\allowbreak{}oscLe\_\allowbreak{}eOf} (\texttt{Development.\allowbreak{}AbelianVarietyTheory.\allowbreak{}Proofs.\allowbreak{}DivisorProofs.\allowbreak{}DepthOscillationLeBoundEveryPrime}, lines 17--62)\\
\texttt{GlobalTheory.\allowbreak{}RefSum.\allowbreak{}sumAt\_\allowbreak{}of\_\allowbreak{}sizes} (\texttt{Development.\allowbreak{}AbelianVarietyTheory.\allowbreak{}Proofs.\allowbreak{}DivisorProofs.\allowbreak{}SumPointOfWeightNumerical}, lines 55--98)\\
\texttt{GlobalTheory.\allowbreak{}CharDepth.\allowbreak{}twoPlaces\_\allowbreak{}host} (\texttt{Development.\allowbreak{}AbelianVarietyTheory.\allowbreak{}Proofs.\allowbreak{}JetHarmonics.\allowbreak{}HostFieldTwoConjugatePlacesOverTwo}, lines 16--27)\\
\texttt{UpperTheory.\allowbreak{}KR.\allowbreak{}small\_\allowbreak{}vd\_\allowbreak{}sum} (\texttt{Development.\allowbreak{}AutomorphicArithmetic.\allowbreak{}CyclotomicBounds.\allowbreak{}RamifiedDiscriminants.\allowbreak{}SmallVdSum}, lines 7--56)\\
\texttt{UpperTheory.\allowbreak{}splitPlaces\_\allowbreak{}discriminantDepth\_\allowbreak{}of\_\allowbreak{}splitAt} (\texttt{Development.\allowbreak{}AutomorphicArithmetic.\allowbreak{}CMTraceForms.\allowbreak{}SplitPlacesDiscriminantDepthOfSplitAt}, lines 59--150)\\
\texttt{GlobalTheory.\allowbreak{}OscSplit.\allowbreak{}oscLe\_\allowbreak{}split} (\texttt{Development.\allowbreak{}AbelianVarietyTheory.\allowbreak{}Proofs.\allowbreak{}DivisorProofs.\allowbreak{}DepthOscillationLeOddSplitUnramifiedPrime}, lines 57--97)\\
\texttt{GlobalTheory.\allowbreak{}RefSum.\allowbreak{}oscLe\_\allowbreak{}eSmall\_\allowbreak{}of\_\allowbreak{}inert} (\texttt{Development.\allowbreak{}AbelianVarietyTheory.\allowbreak{}Proofs.\allowbreak{}DivisorProofs.\allowbreak{}DepthOscillationLeAtSmallOddInertPrime}, lines 180--195)\\
\texttt{GlobalTheory.\allowbreak{}RefSum.\allowbreak{}InertOsc.\allowbreak{}oscLe\_\allowbreak{}inert\_\allowbreak{}large} (\texttt{Development.\allowbreak{}AbelianVarietyTheory.\allowbreak{}Proofs.\allowbreak{}DivisorProofs.\allowbreak{}DepthOscillationLargeInertPrime}, lines 62--118)\par}

{\footnotesize\raggedright\noindent \ref{B8.28ya}.\quad \texttt{GlobalTheory.\allowbreak{}LevelPrimePoint.\allowbreak{}exists\_\allowbreak{}minor\_\allowbreak{}harm} (\texttt{Development.\allowbreak{}AbelianVarietyTheory.\allowbreak{}Proofs.\allowbreak{}JetHarmonics.\allowbreak{}ThetaGeneratorsLeastValueOverLevelPrime}, lines 57--131)\\
\texttt{GlobalTheory.\allowbreak{}LevelPrimePoint.\allowbreak{}crP\_\allowbreak{}mem\_\allowbreak{}harmK\_\allowbreak{}of\_\allowbreak{}disjoint} (\texttt{Development.\allowbreak{}AbelianVarietyTheory.\allowbreak{}Proofs.\allowbreak{}JetHarmonics.\allowbreak{}ThetaGeneratorsLeastValueOverLevelPrime}, lines 31--54)\\
\texttt{Formalism.\allowbreak{}L2.\allowbreak{}Pj.\allowbreak{}Crude.\allowbreak{}crP\_\allowbreak{}rowInv} (\texttt{Formalism.\allowbreak{}LocalToGlobal.\allowbreak{}ProjectiveLattices.\allowbreak{}CrudeMin}, lines 143--194)\\
\texttt{Formalism.\allowbreak{}L2.\allowbreak{}Pj.\allowbreak{}Crude.\allowbreak{}crP\_\allowbreak{}rowDeg} (\texttt{Formalism.\allowbreak{}LocalToGlobal.\allowbreak{}ProjectiveLattices.\allowbreak{}CrudeMin}, lines 344--350)\par}

{\footnotesize\raggedright\noindent \ref{B8.28yb}.\quad \texttt{GlobalTheory.\allowbreak{}PowGerm.\allowbreak{}lemmaBT\_\allowbreak{}of\_\allowbreak{}characters} (\texttt{Development.\allowbreak{}AbelianVarietyTheory.\allowbreak{}Proofs.\allowbreak{}JetHarmonics.\allowbreak{}IsUnitIffNotIsUnitOfCoprime}, lines 50--140)\\
\texttt{GlobalTheory.\allowbreak{}T.\allowbreak{}exists\_\allowbreak{}characters\_\allowbreak{}boundedMinors\_\allowbreak{}of\_\allowbreak{}hodgeTateChartBound} (\texttt{Development.\allowbreak{}AbelianVarietyTheory.\allowbreak{}Proofs.\allowbreak{}JetHarmonics.\allowbreak{}ExistsCharactersBoundedMinorsOfHodgeTateChartBouStatement}, lines 19--21)\\
\texttt{GlobalTheory.\allowbreak{}T.\allowbreak{}depth\_\allowbreak{}bound\_\allowbreak{}from\_\allowbreak{}characters\_\allowbreak{}pointDatum} (\texttt{Development.\allowbreak{}AbelianVarietyTheory.\allowbreak{}Proofs.\allowbreak{}JetHarmonics.\allowbreak{}DepthBoundFromCharactersPointDatumStatement}, lines 13--15)\par}

{\footnotesize\raggedright\noindent \ref{B8.28yc}.\quad \texttt{GlobalTheory.\allowbreak{}T.\allowbreak{}pointStatements\_\allowbreak{}on\_\allowbreak{}copies\_\allowbreak{}hold} (\texttt{Development.\allowbreak{}AbelianVarietyTheory.\allowbreak{}Proofs.\allowbreak{}JetHarmonics.\allowbreak{}PointStatementsOnCopiesHoldStatement}, lines 15--17)\\
\texttt{GlobalTheory.\allowbreak{}PowGerm.\allowbreak{}GenPointT} (\texttt{Development.\allowbreak{}AbelianVarietyTheory.\allowbreak{}Definitions.\allowbreak{}Jets.\allowbreak{}JetsOnCopiesCompatibleFrames}, lines 341--368)\\
\texttt{GlobalTheory.\allowbreak{}PtTie.\allowbreak{}harmonicValue\_\allowbreak{}translate} (\texttt{Development.\allowbreak{}AbelianVarietyTheory.\allowbreak{}Proofs.\allowbreak{}JetHarmonics.\allowbreak{}HarmonicValueTranslateEqUnitMulNormaliserThetaVa}, lines 62--131)\\
\texttt{GlobalTheory.\allowbreak{}LevelPrimePoint.\allowbreak{}exists\_\allowbreak{}val\_\allowbreak{}le} (\texttt{Development.\allowbreak{}AbelianVarietyTheory.\allowbreak{}Proofs.\allowbreak{}JetHarmonics.\allowbreak{}ThetaGeneratorsLeastValueOverLevelPrime}, lines 156--165)\par}

{\footnotesize\raggedright\noindent \ref{B8.28y}.\quad \texttt{GlobalTheory.\allowbreak{}LevelPrimePoint.\allowbreak{}least\_\allowbreak{}value\_\allowbreak{}le} (\texttt{Development.\allowbreak{}AbelianVarietyTheory.\allowbreak{}Proofs.\allowbreak{}JetHarmonics.\allowbreak{}ThetaGeneratorsLeastValueOverLevelPrime}, lines 196--303)\\
\texttt{GlobalTheory.\allowbreak{}ChF.\allowbreak{}hodgeTate\_\allowbreak{}inclusion\_\allowbreak{}chart} (\texttt{Development.\allowbreak{}AbelianVarietyTheory.\allowbreak{}Proofs.\allowbreak{}UnramifiedDepth.\allowbreak{}HodgeTateInclusionChartCharacters}, lines 49--93)\par}

{\footnotesize\raggedright\noindent \ref{B8.28s}.\quad \texttt{GlobalTheory.\allowbreak{}RefLevel.\allowbreak{}refExp\_\allowbreak{}level\_\allowbreak{}le\_\allowbreak{}one} (\texttt{Development.\allowbreak{}AbelianVarietyTheory.\allowbreak{}Proofs.\allowbreak{}DivisorProofs.\allowbreak{}SupLatticeReferenceExponentThetaWeightLevelLeOne}, lines 163--172)\\
\texttt{GlobalTheory.\allowbreak{}RefLevel.\allowbreak{}refExp\_\allowbreak{}le\_\allowbreak{}of\_\allowbreak{}halfSplit} (\texttt{Development.\allowbreak{}AbelianVarietyTheory.\allowbreak{}Proofs.\allowbreak{}DivisorProofs.\allowbreak{}SupLatticeReferenceExponentThetaWeightLevelLeOne}, lines 71--162)\\
\texttt{GlobalTheory.\allowbreak{}RefLevel.\allowbreak{}lattice\_\allowbreak{}terms\_\allowbreak{}le} (\texttt{Development.\allowbreak{}AbelianVarietyTheory.\allowbreak{}Proofs.\allowbreak{}DivisorProofs.\allowbreak{}SupLatticeReferenceExponentThetaWeightLevelLeOne}, lines 27--35)\\
\texttt{GlobalTheory.\allowbreak{}LevelPrimePoint.\allowbreak{}least\_\allowbreak{}value\_\allowbreak{}le} (\texttt{Development.\allowbreak{}AbelianVarietyTheory.\allowbreak{}Proofs.\allowbreak{}JetHarmonics.\allowbreak{}ThetaGeneratorsLeastValueOverLevelPrime}, lines 196--303)\\
\texttt{GlobalTheory.\allowbreak{}ChF.\allowbreak{}hodgeTate\_\allowbreak{}inclusion\_\allowbreak{}chart} (\texttt{Development.\allowbreak{}AbelianVarietyTheory.\allowbreak{}Proofs.\allowbreak{}UnramifiedDepth.\allowbreak{}HodgeTateInclusionChartCharacters}, lines 49--93)\par}

{\footnotesize\raggedright\noindent \ref{B8.28z}.\quad \texttt{GlobalTheory.\allowbreak{}T.\allowbreak{}refW\_\allowbreak{}of\_\allowbreak{}prints\_\allowbreak{}size\_\allowbreak{}flatDeg\_\allowbreak{}refTheta} (\texttt{Development.\allowbreak{}AutomorphicArithmetic.\allowbreak{}ThetaCorrespondence.\allowbreak{}ThetaPairLifts.\allowbreak{}ThirdPlaceSeries.\allowbreak{}RefinedSumWeightOfRefinedTheta}, lines 48--49)\\
\texttt{GlobalTheory.\allowbreak{}RefSum.\allowbreak{}ToWeight.\allowbreak{}primeSum\_\allowbreak{}add} (\texttt{Development.\allowbreak{}AbelianVarietyTheory.\allowbreak{}Proofs.\allowbreak{}DivisorProofs.\allowbreak{}ReferenceWeightOfTheta}, lines 51--71)\\
\texttt{GlobalTheory.\allowbreak{}RefSum.\allowbreak{}ToWeight.\allowbreak{}sum\_\allowbreak{}different\_\allowbreak{}le} (\texttt{Development.\allowbreak{}AbelianVarietyTheory.\allowbreak{}Proofs.\allowbreak{}DivisorProofs.\allowbreak{}ReferenceWeightOfTheta}, lines 32--45)\\
\texttt{GlobalTheory.\allowbreak{}RefSum.\allowbreak{}refinedSum\_\allowbreak{}weight\_\allowbreak{}of\_\allowbreak{}size\_\allowbreak{}alt} (\texttt{Development.\allowbreak{}AutomorphicArithmetic.\allowbreak{}ThetaCorrespondence.\allowbreak{}ThetaPairLifts.\allowbreak{}ThirdPlaceSeries.\allowbreak{}RefinedSumWeightOfSizeLe}, lines 16--17)\par}

{\footnotesize\raggedright\noindent \ref{B8.28t}.\quad \texttt{GlobalTheory.\allowbreak{}RefSum.\allowbreak{}refCore\_\allowbreak{}of\_\allowbreak{}size\_\allowbreak{}sumPt} (\texttt{Development.\allowbreak{}AutomorphicArithmetic.\allowbreak{}ThetaCorrespondence.\allowbreak{}ThetaPairLifts.\allowbreak{}ThirdPlaceSeries.\allowbreak{}RefCoreOfSumPt}, lines 15--37)\\
\texttt{GlobalTheory.\allowbreak{}RefSum.\allowbreak{}primeSum\_\allowbreak{}add\_\allowbreak{}pair} (\texttt{Development.\allowbreak{}AbelianVarietyTheory.\allowbreak{}Proofs.\allowbreak{}DivisorProofs.\allowbreak{}ReferenceCoreOfSumPoint}, lines 15--38)\\
\texttt{GlobalTheory.\allowbreak{}RefSum.\allowbreak{}ToWeight.\allowbreak{}primeSum\_\allowbreak{}add} (\texttt{Development.\allowbreak{}AbelianVarietyTheory.\allowbreak{}Proofs.\allowbreak{}DivisorProofs.\allowbreak{}ReferenceWeightOfTheta}, lines 51--71)\\
\texttt{GlobalTheory.\allowbreak{}RefSum.\allowbreak{}ToWeight.\allowbreak{}sum\_\allowbreak{}different\_\allowbreak{}le} (\texttt{Development.\allowbreak{}AbelianVarietyTheory.\allowbreak{}Proofs.\allowbreak{}DivisorProofs.\allowbreak{}ReferenceWeightOfTheta}, lines 32--45)\\
\texttt{GlobalTheory.\allowbreak{}RefSum.\allowbreak{}refinedSum\_\allowbreak{}weight\_\allowbreak{}of\_\allowbreak{}size\_\allowbreak{}alt} (\texttt{Development.\allowbreak{}AutomorphicArithmetic.\allowbreak{}ThetaCorrespondence.\allowbreak{}ThetaPairLifts.\allowbreak{}ThirdPlaceSeries.\allowbreak{}RefinedSumWeightOfSizeLe}, lines 16--17)\par}

{\footnotesize\raggedright\noindent \ref{B8.28}.\quad \texttt{SelfIntersectionVolumeEstimates.\allowbreak{}height\_\allowbreak{}lower\_\allowbreak{}bound\_\allowbreak{}general\_\allowbreak{}alt} (\texttt{Development.\allowbreak{}AutomorphicArithmetic.\allowbreak{}ThetaCorrespondence.\allowbreak{}ThetaPairLifts.\allowbreak{}ThirdPlaceSeries.\allowbreak{}HeightLowerBoundGeneral}, lines 21--22)\\
\texttt{GlobalTheory.\allowbreak{}CovTop.\allowbreak{}HeightAt} (\texttt{Development.\allowbreak{}AbelianVarietyTheory.\allowbreak{}Definitions.\allowbreak{}Divisors.\allowbreak{}CovolXOfCovol}, lines 146--152)\par}

{\footnotesize\raggedright\noindent \ref{B8.28}.\quad \texttt{GlobalTheory.\allowbreak{}CovHt.\allowbreak{}HtTopG.\allowbreak{}slope\_\allowbreak{}le\_\allowbreak{}height} (\texttt{Development.\allowbreak{}AutomorphicArithmetic.\allowbreak{}ThetaCorrespondence.\allowbreak{}ThetaPairLifts.\allowbreak{}HodgeBundleWeights.\allowbreak{}HeightAddLogNonneg}, lines 29--65)\\
\texttt{GlobalTheory.\allowbreak{}CovHt.\allowbreak{}HtTopG.\allowbreak{}heightAt\_\allowbreak{}eq\_\allowbreak{}of\_\allowbreak{}cover} (\texttt{Development.\allowbreak{}AutomorphicArithmetic.\allowbreak{}ThetaCorrespondence.\allowbreak{}ThetaPairLifts.\allowbreak{}HodgeBundleWeights.\allowbreak{}HeightAddLogNonneg}, lines 66--105)\\
\texttt{GlobalTheory.\allowbreak{}CovHt.\allowbreak{}HtTopG.\allowbreak{}height\_\allowbreak{}of\_\allowbreak{}meanAll\_\allowbreak{}bound} (\texttt{Development.\allowbreak{}AutomorphicArithmetic.\allowbreak{}ThetaCorrespondence.\allowbreak{}ThetaPairLifts.\allowbreak{}ThirdPlaceSeries.\allowbreak{}HeightOfSteps}, lines 19--38)\\
\texttt{GlobalTheory.\allowbreak{}CompHt.\allowbreak{}le\_\allowbreak{}compVolSlope\_\allowbreak{}iff} (\texttt{Development.\allowbreak{}AbelianVarietyTheory.\allowbreak{}Definitions.\allowbreak{}Divisors.\allowbreak{}IsHomogeneousCompIdeal}, lines 283--286)\par}

{\footnotesize\raggedright\noindent \ref{B8.29}.\quad \texttt{GlobalTheory.\allowbreak{}T.\allowbreak{}uniform\_\allowbreak{}lower\_\allowbreak{}bound\_\allowbreak{}of\_\allowbreak{}sums'} (\texttt{Development.\allowbreak{}AutomorphicArithmetic.\allowbreak{}ThetaCorrespondence.\allowbreak{}ThetaPairLifts.\allowbreak{}ThirdPlaceSeries.\allowbreak{}ULowOfStatement}, lines 11--12)\\
\texttt{GlobalTheory.\allowbreak{}T.\allowbreak{}liftLowBound\_\allowbreak{}of\_\allowbreak{}cover\_\allowbreak{}estimates} (\texttt{Development.\allowbreak{}AutomorphicArithmetic.\allowbreak{}ThetaCorrespondence.\allowbreak{}ThetaPairLifts.\allowbreak{}ThirdPlaceSeries.\allowbreak{}LowOfCoverHypothesesStatement}, lines 11--12)\par}

{\footnotesize\raggedright\noindent \ref{B8.29}.\quad \texttt{GlobalTheory.\allowbreak{}CovTop.\allowbreak{}DefLow.\allowbreak{}covering\_\allowbreak{}low\_\allowbreak{}defect\_\allowbreak{}alt} (\texttt{Development.\allowbreak{}AutomorphicArithmetic.\allowbreak{}ThetaCorrespondence.\allowbreak{}ThetaPairLifts.\allowbreak{}ThirdPlaceSeries.\allowbreak{}CoveringLowDefectBound}, lines 17--46)\\
\texttt{GlobalTheory.\allowbreak{}T.\allowbreak{}two\_\allowbreak{}definiteBound\_\allowbreak{}ge\_\allowbreak{}anchoredMainTerm\_\allowbreak{}sub\_\allowbreak{}crude\_\allowbreak{}errors} (\texttt{Development.\allowbreak{}AbelianVarietyTheory.\allowbreak{}Proofs.\allowbreak{}DivisorProofs.\allowbreak{}TwoDefiniteBoundGeAnchoredMainTermSubCrudeErrorsStatement}, lines 12--13)\\
\texttt{GlobalTheory.\allowbreak{}T.\allowbreak{}weight\_\allowbreak{}dimension\_\allowbreak{}log\_\allowbreak{}bound'} (\texttt{Development.\allowbreak{}AutomorphicArithmetic.\allowbreak{}ThetaCorrespondence.\allowbreak{}ThetaPairLifts.\allowbreak{}ThirdPlaceSeries.\allowbreak{}ParametersScaleRankGapBoundsStatement}, lines 13--14)\\
\texttt{GlobalTheory.\allowbreak{}CovTop.\allowbreak{}DefSizes.\allowbreak{}parameters\_\allowbreak{}scale\_\allowbreak{}rank\_\allowbreak{}gap\_\allowbreak{}bounds} (\texttt{Development.\allowbreak{}AutomorphicArithmetic.\allowbreak{}ThetaCorrespondence.\allowbreak{}ThetaPairLifts.\allowbreak{}ThirdPlaceSeries.\allowbreak{}ParametersScaleRankGapBounds}, lines 24--41)\\
\texttt{GlobalTheory.\allowbreak{}CovTop.\allowbreak{}DefSizes.\allowbreak{}sizes\_\allowbreak{}at} (\texttt{Development.\allowbreak{}AbelianVarietyTheory.\allowbreak{}Proofs.\allowbreak{}DivisorProofs.\allowbreak{}WeightDimensionLogBound}, lines 188--240)\\
\texttt{GlobalTheory.\allowbreak{}CovTop.\allowbreak{}FromComp.\allowbreak{}core} (\texttt{Development.\allowbreak{}AbelianVarietyTheory.\allowbreak{}Proofs.\allowbreak{}DivisorProofs.\allowbreak{}UniformLowerBoundOfSums}, lines 15--33)\\
\texttt{Formalism.\allowbreak{}ExtensionFiltration.\allowbreak{}sumComponents'} (\texttt{Formalism.\allowbreak{}MeanHeightDescent.\allowbreak{}AverageDegreeBounds.\allowbreak{}SumComponents}, lines 8--119)\\
\texttt{GlobalTheory.\allowbreak{}DescentFamily.\allowbreak{}descent} (\texttt{Development.\allowbreak{}AbelianVarietyTheory.\allowbreak{}Proofs.\allowbreak{}DivisorProofs.\allowbreak{}HtMapLeOfPoints}, lines 69--144)\\
\texttt{GlobalTheory.\allowbreak{}CovTop.\allowbreak{}LowCover.\allowbreak{}lowAt\_\allowbreak{}of\_\allowbreak{}cover} (\texttt{Development.\allowbreak{}AbelianVarietyTheory.\allowbreak{}Proofs.\allowbreak{}DivisorProofs.\allowbreak{}LowerBoundOfAllEstimates}, lines 65--106)\\
\texttt{GlobalTheory.\allowbreak{}CovTop.\allowbreak{}LowCover.\allowbreak{}low\_\allowbreak{}of\_\allowbreak{}cover\_\allowbreak{}hypotheses} (\texttt{Development.\allowbreak{}AutomorphicArithmetic.\allowbreak{}ThetaCorrespondence.\allowbreak{}ThetaPairLifts.\allowbreak{}ThirdPlaceSeries.\allowbreak{}LowOfCoverHypotheses}, lines 16--32)\par}

{\footnotesize\raggedright\noindent \ref{B8.29}.\quad \texttt{GlobalTheory.\allowbreak{}CovTop.\allowbreak{}latGood} (\texttt{Development.\allowbreak{}AbelianVarietyTheory.\allowbreak{}Definitions.\allowbreak{}Divisors.\allowbreak{}NormNegCmConstantEqAtAllEmbeddings}, lines 44--46)\\
\texttt{GlobalTheory.\allowbreak{}CovTop.\allowbreak{}latBad} (\texttt{Development.\allowbreak{}AbelianVarietyTheory.\allowbreak{}Definitions.\allowbreak{}Divisors.\allowbreak{}NormNegCmConstantEqAtAllEmbeddings}, lines 47--49)\par}

{\footnotesize\raggedright\noindent \ref{B8.30}.\quad \texttt{GlobalTheory.\allowbreak{}T.\allowbreak{}liftBadSet\_\allowbreak{}of\_\allowbreak{}deep\_\allowbreak{}size\_\allowbreak{}flatDeg\_\allowbreak{}share} (\texttt{Development.\allowbreak{}AutomorphicArithmetic.\allowbreak{}ThetaCorrespondence.\allowbreak{}ThetaPairLifts.\allowbreak{}ThirdPlaceSeries.\allowbreak{}CoveringBadLocusOfDescentSteps}, lines 54--55)\\
\texttt{GlobalTheory.\allowbreak{}CovTop.\allowbreak{}BadAt} (\texttt{Development.\allowbreak{}AbelianVarietyTheory.\allowbreak{}Definitions.\allowbreak{}Divisors.\allowbreak{}CovolXOfCovol}, lines 69--71)\par}

{\footnotesize\raggedright\noindent \ref{B8.30}.\quad \texttt{GlobalTheory.\allowbreak{}CovTop.\allowbreak{}LocShareNum.\allowbreak{}local\_\allowbreak{}share\_\allowbreak{}count\_\allowbreak{}le\_\allowbreak{}exp\_\allowbreak{}alt} (\texttt{Development.\allowbreak{}AutomorphicArithmetic.\allowbreak{}ThetaCorrespondence.\allowbreak{}ThetaPairLifts.\allowbreak{}ThirdPlaceSeries.\allowbreak{}LocalShareCountLeExp}, lines 31--56)\\
\texttt{GlobalTheory.\allowbreak{}CovTop.\allowbreak{}LocShareNum.\allowbreak{}wg\_\allowbreak{}sum\_\allowbreak{}le} (\texttt{Development.\allowbreak{}AbelianVarietyTheory.\allowbreak{}Proofs.\allowbreak{}DivisorProofs.\allowbreak{}LocalShareSumBound}, lines 271--283)\\
\texttt{GlobalTheory.\allowbreak{}BadEmb.\allowbreak{}embReg\_\allowbreak{}deep} (\texttt{Development.\allowbreak{}AbelianVarietyTheory.\allowbreak{}Proofs.\allowbreak{}DivisorProofs.\allowbreak{}CoveringRepresentationTypeFrame}, lines 185--208)\\
\texttt{GlobalTheory.\allowbreak{}BadEmb.\allowbreak{}regular\_\allowbreak{}embeddings\_\allowbreak{}cover\_\allowbreak{}bad\_\allowbreak{}embedding\_\allowbreak{}locus\_\allowbreak{}alt} (\texttt{Development.\allowbreak{}AutomorphicArithmetic.\allowbreak{}ThetaCorrespondence.\allowbreak{}ThetaPairLifts.\allowbreak{}ThirdPlaceSeries.\allowbreak{}RegularEmbeddingsCoverBadEmbeddingLocus}, lines 38--86)\\
\texttt{GlobalTheory.\allowbreak{}BadIndef.\allowbreak{}bad\_\allowbreak{}components\_\allowbreak{}indefinite\_\allowbreak{}alt} (\texttt{Development.\allowbreak{}AutomorphicArithmetic.\allowbreak{}ThetaCorrespondence.\allowbreak{}ThetaPairLifts.\allowbreak{}ThirdPlaceSeries.\allowbreak{}BadComponentsIndefinite}, lines 24--54)\\
\texttt{GlobalTheory.\allowbreak{}CovTop.\allowbreak{}definite\_\allowbreak{}bad\_\allowbreak{}place\_\allowbreak{}lift\_\allowbreak{}alt} (\texttt{Development.\allowbreak{}AutomorphicArithmetic.\allowbreak{}ThetaCorrespondence.\allowbreak{}ThetaPairLifts.\allowbreak{}ThirdPlaceSeries.\allowbreak{}ExistsLiftAtDefiniteBadPlaces}, lines 21--53)\\
\texttt{GlobalTheory.\allowbreak{}CovTop.\allowbreak{}covering\_\allowbreak{}bad\_\allowbreak{}locus\_\allowbreak{}of\_\allowbreak{}descent\_\allowbreak{}steps\_\allowbreak{}alt} (\texttt{Development.\allowbreak{}AutomorphicArithmetic.\allowbreak{}ThetaCorrespondence.\allowbreak{}ThetaPairLifts.\allowbreak{}ThirdPlaceSeries.\allowbreak{}CoveringBadLocusOfDescentSteps}, lines 21--40)\par}

{\footnotesize\raggedright\noindent \ref{B8.31}.\quad \texttt{GlobalTheory.\allowbreak{}CutNum.\allowbreak{}slope\_\allowbreak{}good\_\allowbreak{}ge} (\texttt{Development.\allowbreak{}AbelianVarietyTheory.\allowbreak{}Proofs.\allowbreak{}DivisorProofs.\allowbreak{}MajAtAddRight}, lines 111--143)\par}

{\footnotesize\raggedright\noindent \ref{B8.31}.\quad \texttt{GlobalTheory.\allowbreak{}CutNum.\allowbreak{}numeric} (\texttt{Development.\allowbreak{}AbelianVarietyTheory.\allowbreak{}Proofs.\allowbreak{}DivisorProofs.\allowbreak{}MajAtAddRight}, lines 93--109)\\
\texttt{Formalism.\allowbreak{}ExtensionFiltration.\allowbreak{}goodBadCut} (\texttt{Formalism.\allowbreak{}MeanHeightDescent.\allowbreak{}DivisionInequalityWithCovolume}, lines 23--23)\par}

{\footnotesize\raggedright\noindent \ref{B8.31}.\quad \texttt{GlobalTheory.\allowbreak{}CovTop.\allowbreak{}GramAt} (\texttt{Development.\allowbreak{}AbelianVarietyTheory.\allowbreak{}Definitions.\allowbreak{}Divisors.\allowbreak{}CovolXOfCovol}, lines 53--55)\\
\texttt{GlobalTheory.\allowbreak{}CovTop.\allowbreak{}RecAt} (\texttt{Development.\allowbreak{}AbelianVarietyTheory.\allowbreak{}Definitions.\allowbreak{}Divisors.\allowbreak{}CovolXOfCovol}, lines 57--59)\\
\texttt{GlobalTheory.\allowbreak{}CovTop.\allowbreak{}CutAt} (\texttt{Development.\allowbreak{}AbelianVarietyTheory.\allowbreak{}Definitions.\allowbreak{}Divisors.\allowbreak{}CovolXOfCovol}, lines 61--67)\\
\texttt{GlobalTheory.\allowbreak{}CovTop.\allowbreak{}CoreAt} (\texttt{Development.\allowbreak{}AbelianVarietyTheory.\allowbreak{}Definitions.\allowbreak{}Divisors.\allowbreak{}CovolXOfCovol}, lines 92--94)\\
\texttt{GlobalTheory.\allowbreak{}CovTop.\allowbreak{}LowAt} (\texttt{Development.\allowbreak{}AbelianVarietyTheory.\allowbreak{}Definitions.\allowbreak{}Divisors.\allowbreak{}CovolXOfCovol}, lines 96--96)\\
\texttt{GlobalTheory.\allowbreak{}CovTop.\allowbreak{}recAt\_\allowbreak{}of\_\allowbreak{}gramAt\_\allowbreak{}alt} (\texttt{Development.\allowbreak{}AutomorphicArithmetic.\allowbreak{}ThetaCorrespondence.\allowbreak{}ThetaPairLifts.\allowbreak{}ThirdPlaceSeries.\allowbreak{}RecAtOfGramAt}, lines 17--46)\\
\texttt{GlobalTheory.\allowbreak{}T.\allowbreak{}group\_\allowbreak{}idempotent\_\allowbreak{}coefficients\_\allowbreak{}alt} (\texttt{Development.\allowbreak{}AutomorphicArithmetic.\allowbreak{}ThetaCorrespondence.\allowbreak{}ThetaPairLifts.\allowbreak{}ThirdPlaceSeries.\allowbreak{}TiedGoodIdemStatement}, lines 12--18)\\
\texttt{GlobalTheory.\allowbreak{}CovTop.\allowbreak{}cutAt\_\allowbreak{}of\_\allowbreak{}rows\_\allowbreak{}alt} (\texttt{Development.\allowbreak{}AutomorphicArithmetic.\allowbreak{}ThetaCorrespondence.\allowbreak{}ThetaPairLifts.\allowbreak{}ThirdPlaceSeries.\allowbreak{}CutAtOfRows}, lines 18--26)\\
\texttt{GlobalTheory.\allowbreak{}CoreSlope.\allowbreak{}coreLe\_\allowbreak{}of\_\allowbreak{}gram\_\allowbreak{}meas\_\allowbreak{}rat\_\allowbreak{}alt} (\texttt{Development.\allowbreak{}AutomorphicArithmetic.\allowbreak{}ThetaCorrespondence.\allowbreak{}ThetaPairLifts.\allowbreak{}ThirdPlaceSeries.\allowbreak{}CoreLeOfGramMeasRat}, lines 16--24)\par}

{\footnotesize\raggedright\noindent \ref{B8.32}.\quad \texttt{GlobalTheory.\allowbreak{}CovTop.\allowbreak{}Assembly.\allowbreak{}lcv\_\allowbreak{}le} (\texttt{Development.\allowbreak{}AbelianVarietyTheory.\allowbreak{}Proofs.\allowbreak{}DivisorProofs.\allowbreak{}CovolumeFromGramAndBounds}, lines 39--85)\par}

{\footnotesize\raggedright\noindent \ref{B8.33}.\quad \texttt{GlobalTheory.\allowbreak{}CovTop.\allowbreak{}Assembly.\allowbreak{}covolX\_\allowbreak{}of\_\allowbreak{}lift\_\allowbreak{}bounds\_\allowbreak{}and\_\allowbreak{}coreLe} (\texttt{Development.\allowbreak{}AutomorphicArithmetic.\allowbreak{}ThetaCorrespondence.\allowbreak{}ThetaPairLifts.\allowbreak{}ThirdPlaceSeries.\allowbreak{}CovolXOfSteps}, lines 16--40)\par}

{\footnotesize\raggedright\noindent \ref{B8.33}.\quad \texttt{GlobalTheory.\allowbreak{}CovTop.\allowbreak{}Assembly.\allowbreak{}exists\_\allowbreak{}loss\_\allowbreak{}le} (\texttt{Development.\allowbreak{}AbelianVarietyTheory.\allowbreak{}Proofs.\allowbreak{}DivisorProofs.\allowbreak{}CovolumeFromGramAndBounds}, lines 25--35)\par}

{\footnotesize\raggedright\noindent \ref{B8.34a}.\quad \texttt{GlobalTheory.\allowbreak{}AtMember.\allowbreak{}fieldTies\_\allowbreak{}of\_\allowbreak{}ksExact} (\texttt{Development.\allowbreak{}AbelianVarietyTheory.\allowbreak{}Proofs.\allowbreak{}IwasawaDecomposition.\allowbreak{}IndexBoundCoefficientFieldOfExactSmallField}, lines 28--77)\\
\texttt{GlobalTheory.\allowbreak{}FieldSize.\allowbreak{}sz\_\allowbreak{}Om\_\allowbreak{}poR} (\texttt{Development.\allowbreak{}AbelianVarietyTheory.\allowbreak{}Proofs.\allowbreak{}IwasawaDecomposition.\allowbreak{}OmegaFieldDegreeRootDiscrBoundedByRingPolynomial}, lines 52--141)\par}

{\footnotesize\raggedright\noindent \ref{B8.34}.\quad \texttt{GlobalTheory.\allowbreak{}Glob.\allowbreak{}CovolRow.\allowbreak{}linked\_\allowbreak{}bundles\_\allowbreak{}l3\_\allowbreak{}of\_\allowbreak{}open\_\allowbreak{}alt} (\texttt{Development.\allowbreak{}AutomorphicArithmetic.\allowbreak{}ThetaCorrespondence.\allowbreak{}ThetaPairLifts.\allowbreak{}ThirdPlaceSeries.\allowbreak{}LinkedBundlesL3OfOpen}, lines 31--46)\\
\texttt{GlobalTheory.\allowbreak{}Fc.\allowbreak{}S0.\allowbreak{}CovolumeBoundThirdPlaceSeries} (\texttt{Development.\allowbreak{}AutomorphicArithmetic.\allowbreak{}ThetaCorrespondence.\allowbreak{}ThetaPairLifts.\allowbreak{}ThirdPlaceSeries.\allowbreak{}GeometricWeakTies}, lines 32--33)\\
\texttt{GlobalTheory.\allowbreak{}Fc.\allowbreak{}Row} (\texttt{Development.\allowbreak{}AbelianVarietyTheory.\allowbreak{}Definitions.\allowbreak{}GlobalTwisting.\allowbreak{}ValidForConstructedPointClassesEventuallyWorking}, lines 268--269)\\
\texttt{GlobalTheory.\allowbreak{}Fc.\allowbreak{}W3} (\texttt{Development.\allowbreak{}AbelianVarietyTheory.\allowbreak{}Definitions.\allowbreak{}GlobalTwisting.\allowbreak{}ValidForConstructedPointClassesEventuallyWorking}, lines 371--372)\\
\texttt{GlobalTheory.\allowbreak{}Fc.\allowbreak{}Ind.\allowbreak{}EvWC} (\texttt{Development.\allowbreak{}AbelianVarietyTheory.\allowbreak{}Definitions.\allowbreak{}Divisors.\allowbreak{}BaseRankDimension}, lines 90--91)\par}

{\footnotesize\raggedright\noindent \ref{B8.34}.\quad \texttt{GlobalTheory.\allowbreak{}Fc.\allowbreak{}Ind.\allowbreak{}eventual\_\allowbreak{}L3\_\allowbreak{}of\_\allowbreak{}covolume\_\allowbreak{}alt} (\texttt{Development.\allowbreak{}AutomorphicArithmetic.\allowbreak{}ThetaCorrespondence.\allowbreak{}ThetaPairLifts.\allowbreak{}ThirdPlaceSeries.\allowbreak{}EventualL3OfCovolume}, lines 14--20)\\
\texttt{GlobalTheory.\allowbreak{}Fc.\allowbreak{}Ind.\allowbreak{}evWC\_\allowbreak{}onPt\_\allowbreak{}of\_\allowbreak{}guarded\_\allowbreak{}alt} (\texttt{Development.\allowbreak{}AutomorphicArithmetic.\allowbreak{}ThetaCorrespondence.\allowbreak{}ThetaPairLifts.\allowbreak{}ThirdPlaceSeries.\allowbreak{}EvWCOnPtOfGuarded}, lines 12--32)\\
\texttt{GlobalTheory.\allowbreak{}Fc.\allowbreak{}Ind.\allowbreak{}summandL3Bound} (\texttt{Development.\allowbreak{}AutomorphicArithmetic.\allowbreak{}ThetaCorrespondence.\allowbreak{}ThetaPairLifts.\allowbreak{}ThirdPlaceSeries.\allowbreak{}SummandFunctionalBound}, lines 54--54)\par}

{\footnotesize\raggedright\noindent \ref{B8.34}.\quad \texttt{UpperTheory.\allowbreak{}Lit.\allowbreak{}nonempty\_\allowbreak{}hyps\_\allowbreak{}definite\_\allowbreak{}side\_\allowbreak{}of\_\allowbreak{}closed\_\allowbreak{}bound} (\texttt{Development.\allowbreak{}AutomorphicArithmetic.\allowbreak{}ThetaCorrespondence.\allowbreak{}ThetaPairLifts.\allowbreak{}ThirdPlaceSeries.\allowbreak{}SiegelEisensteinSecondTermHypotheses}, lines 44--68)\\
\texttt{Formalism.\allowbreak{}L2.\allowbreak{}Indef.\allowbreak{}IndefSide.\allowbreak{}boundsMu} (\texttt{Formalism.\allowbreak{}LocalToGlobal.\allowbreak{}IndefiniteForms.\allowbreak{}IndefUse}, lines 271--335)\\
\texttt{Formalism.\allowbreak{}L2.\allowbreak{}DIV.\allowbreak{}wordFacts\_\allowbreak{}of\_\allowbreak{}face} (\texttt{Formalism.\allowbreak{}LocalToGlobal.\allowbreak{}Divisibility.\allowbreak{}RecBridge}, lines 204--226)\par}

\section*{Chapter~\ref{ch:B8-definite-local}: The definite side and local errors}

{\footnotesize\raggedright\noindent \ref{def:B8-word-numbers}.\quad \texttt{Formalism.\allowbreak{}L2.\allowbreak{}DIV.\allowbreak{}NumDatum.\allowbreak{}NumFacts} (\texttt{Formalism.\allowbreak{}LocalToGlobal.\allowbreak{}Divisibility.\allowbreak{}NumbersLink}, lines 43--83)\\
\texttt{GlobalTheory.\allowbreak{}Fc.\allowbreak{}DefiniteSideCondition} (\texttt{Development.\allowbreak{}AutomorphicArithmetic.\allowbreak{}ThetaCorrespondence.\allowbreak{}ThetaPairLifts.\allowbreak{}ThirdPlaceSeries.\allowbreak{}DefiniteSideCondition}, lines 25--35)\par}

{\footnotesize\raggedright\noindent \ref{B8.35}.\quad \texttt{Formalism.\allowbreak{}L2.\allowbreak{}DIV.\allowbreak{}HtNums} (\texttt{Formalism.\allowbreak{}LocalToGlobal.\allowbreak{}Divisibility.\allowbreak{}Inputs}, lines 210--311)\\
\texttt{Formalism.\allowbreak{}L2.\allowbreak{}DIV.\allowbreak{}GrpTie} (\texttt{Formalism.\allowbreak{}LocalToGlobal.\allowbreak{}Divisibility.\allowbreak{}RowSideTie}, lines 156--176)\\
\texttt{Formalism.\allowbreak{}L2.\allowbreak{}DIV.\allowbreak{}LvlTie} (\texttt{Formalism.\allowbreak{}LocalToGlobal.\allowbreak{}Divisibility.\allowbreak{}RowSideTie}, lines 128--154)\\
\texttt{Formalism.\allowbreak{}DimensionVariation.\allowbreak{}liftedMapLogSize} (\texttt{Formalism.\allowbreak{}ThetaSupremumBounds.\allowbreak{}DivIneqExplicit6}, lines 15--23)\\
\texttt{GlobalTheory.\allowbreak{}Fc.\allowbreak{}DefiniteSideCondition} (\texttt{Development.\allowbreak{}AutomorphicArithmetic.\allowbreak{}ThetaCorrespondence.\allowbreak{}ThetaPairLifts.\allowbreak{}ThirdPlaceSeries.\allowbreak{}DefiniteSideCondition}, lines 25--35)\\
\texttt{Formalism.\allowbreak{}L2.\allowbreak{}DIV.\allowbreak{}DivTied3} (\texttt{Formalism.\allowbreak{}LocalToGlobal.\allowbreak{}Divisibility.\allowbreak{}DivisionCompatibleToSub2Indef2}, lines 172--196)\\
\texttt{Formalism.\allowbreak{}L2.\allowbreak{}DIV.\allowbreak{}ColTie} (\texttt{Formalism.\allowbreak{}LocalToGlobal.\allowbreak{}Divisibility.\allowbreak{}RowSideTie}, lines 247--257)\\
\texttt{Formalism.\allowbreak{}L2.\allowbreak{}DIV.\allowbreak{}Cons.\allowbreak{}nuOf} (\texttt{Formalism.\allowbreak{}LocalToGlobal.\allowbreak{}Divisibility.\allowbreak{}DivisibilityConstruction.\allowbreak{}ShiftOfColTie}, lines 14--16)\\
\texttt{Formalism.\allowbreak{}L2.\allowbreak{}RowDataTransport.\allowbreak{}RowsU} (\texttt{Formalism.\allowbreak{}LocalToGlobal.\allowbreak{}PtKOfReal}, lines 141--159)\\
\texttt{Formalism.\allowbreak{}L2.\allowbreak{}DIV.\allowbreak{}ColChain} (\texttt{Formalism.\allowbreak{}LocalToGlobal.\allowbreak{}Divisibility.\allowbreak{}Inputs}, lines 319--341)\\
\texttt{Formalism.\allowbreak{}L2.\allowbreak{}DIV.\allowbreak{}FirstCols} (\texttt{Formalism.\allowbreak{}LocalToGlobal.\allowbreak{}Divisibility.\allowbreak{}DivisionHeightBoundStatement}, lines 504--507)\par}

{\footnotesize\raggedright\noindent \ref{B8.36a}.\quad \texttt{GlobalTheory.\allowbreak{}CoreHt.\allowbreak{}HtRowsS} (\texttt{Development.\allowbreak{}AbelianVarietyTheory.\allowbreak{}Definitions.\allowbreak{}LatticeFrames.\allowbreak{}RecHeightSWith}, lines 11--18)\\
\texttt{GlobalTheory.\allowbreak{}CoreHt.\allowbreak{}DelNat} (\texttt{Development.\allowbreak{}AbelianVarietyTheory.\allowbreak{}Definitions.\allowbreak{}LatticeFrames.\allowbreak{}RecHeightWith}, lines 16--17)\\
\texttt{GlobalTheory.\allowbreak{}CoreHt.\allowbreak{}ThickUp} (\texttt{Development.\allowbreak{}AbelianVarietyTheory.\allowbreak{}Definitions.\allowbreak{}LatticeFrames.\allowbreak{}RecHeightWith}, lines 18--20)\par}

{\footnotesize\raggedright\noindent \ref{B8.35a}.\quad \texttt{GlobalTheory.\allowbreak{}CoreHt.\allowbreak{}HtRowsS} (\texttt{Development.\allowbreak{}AbelianVarietyTheory.\allowbreak{}Definitions.\allowbreak{}LatticeFrames.\allowbreak{}RecHeightSWith}, lines 11--19)\\
\texttt{GlobalTheory.\allowbreak{}CoreHt.\allowbreak{}RowsOfDesign.\allowbreak{}rows\_\allowbreak{}of\_\allowbreak{}design} (\texttt{Development.\allowbreak{}AbelianVarietyTheory.\allowbreak{}Proofs.\allowbreak{}DecoratedLatticeProofs.\allowbreak{}TwoLeOfLogPos}, lines 24--104)\\
\texttt{GlobalTheory.\allowbreak{}EX.\allowbreak{}szRow\_\allowbreak{}sizeNodes} (\texttt{Development.\allowbreak{}AbelianVarietyTheory.\allowbreak{}Proofs.\allowbreak{}GlobalInterpolation.\allowbreak{}ExplicitChoiceSizeNodes}, lines 11--73)\\
\texttt{GlobalTheory.\allowbreak{}CoreHt.\allowbreak{}famRowsSizes} (\texttt{Development.\allowbreak{}AbelianVarietyTheory.\allowbreak{}Proofs.\allowbreak{}DivisorProofs.\allowbreak{}FamRowsSizes}, lines 13--69)\par}

{\footnotesize\raggedright\noindent \ref{def:B8-column-comparisons}.\quad \texttt{Formalism.\allowbreak{}L2.\allowbreak{}Theta.\allowbreak{}RealTied} (\texttt{Formalism.\allowbreak{}LocalToGlobal.\allowbreak{}BesselHermitianTheta.\allowbreak{}MinfDensity}, lines 82--91)\\
\texttt{Formalism.\allowbreak{}MainNotions.\allowbreak{}RT.\allowbreak{}ChiFin0} (\texttt{Formalism.\allowbreak{}SharpRestriction.\allowbreak{}GlobalSchurBounds.\allowbreak{}IntegralityOfRationalityLocal}, lines 15--19)\\
\texttt{Formalism.\allowbreak{}L2.\allowbreak{}DIV.\allowbreak{}ColumnCharactersCompatible} (\texttt{Development.\allowbreak{}AutomorphicArithmetic.\allowbreak{}ThetaCorrespondence.\allowbreak{}ThetaPairLifts.\allowbreak{}ThirdPlaceSeries.\allowbreak{}DIVColPins}, lines 26--72)\\
\texttt{Formalism.\allowbreak{}L2.\allowbreak{}Ker.\allowbreak{}KerEntTied.\allowbreak{}WtTie} (\texttt{Formalism.\allowbreak{}LocalToGlobal.\allowbreak{}KernelLattices.\allowbreak{}JetFieldLeKR}, lines 108--112)\\
\texttt{GlobalTheory.\allowbreak{}ColRows.\allowbreak{}KerTies} (\texttt{Development.\allowbreak{}AbelianVarietyTheory.\allowbreak{}Definitions.\allowbreak{}Divisors.\allowbreak{}KerSizes}, lines 17--26)\\
\texttt{GlobalTheory.\allowbreak{}ColRows.\allowbreak{}KerSizes} (\texttt{Development.\allowbreak{}AbelianVarietyTheory.\allowbreak{}Definitions.\allowbreak{}Divisors.\allowbreak{}KerSizes}, lines 27--31)\\
\texttt{Formalism.\allowbreak{}L2.\allowbreak{}Col.\allowbreak{}Sizes} (\texttt{Formalism.\allowbreak{}LocalToGlobal.\allowbreak{}Columns.\allowbreak{}AbelianHostFieldConditions}, lines 486--495)\\
\texttt{GlobalTheory.\allowbreak{}Core.\allowbreak{}Ht2.\allowbreak{}SizeRowsS} (\texttt{Development.\allowbreak{}AbelianVarietyTheory.\allowbreak{}Definitions.\allowbreak{}LatticeFrames.\allowbreak{}ColHeightSharpWith}, lines 12--27)\par}

{\footnotesize\raggedright\noindent \ref{B8.36}.\quad \texttt{GlobalTheory.\allowbreak{}Fc.\allowbreak{}RL.\allowbreak{}DefSide.\allowbreak{}definiteSideCondition\_\allowbreak{}of\_\allowbreak{}nprimeDef\_\allowbreak{}of\_\allowbreak{}nuFull} (\texttt{Development.\allowbreak{}AutomorphicArithmetic.\allowbreak{}ThetaCorrespondence.\allowbreak{}ThetaPairLifts.\allowbreak{}ThirdPlaceSeries.\allowbreak{}PointEvaluationSide}, lines 49--64)\\
\texttt{GlobalTheory.\allowbreak{}ColRows.\allowbreak{}rowsU\_\allowbreak{}of\_\allowbreak{}colPins\_\allowbreak{}alt} (\texttt{Development.\allowbreak{}AutomorphicArithmetic.\allowbreak{}ThetaCorrespondence.\allowbreak{}ThetaPairLifts.\allowbreak{}ThirdPlaceSeries.\allowbreak{}RowsUOfColPins}, lines 17--84)\par}

{\footnotesize\raggedright\noindent \ref{B8.37}.\quad \texttt{GlobalTheory.\allowbreak{}Fc.\allowbreak{}RL.\allowbreak{}DefSide.\allowbreak{}linkedAt\_\allowbreak{}df\_\allowbreak{}of\_\allowbreak{}parameters} (\texttt{Development.\allowbreak{}AutomorphicArithmetic.\allowbreak{}ThetaCorrespondence.\allowbreak{}ThetaPairLifts.\allowbreak{}ThirdPlaceSeries.\allowbreak{}PointEvaluationSide}, lines 65--83)\\
\texttt{GlobalTheory.\allowbreak{}Fc.\allowbreak{}RL.\allowbreak{}DefSide.\allowbreak{}full\_\allowbreak{}rank} (\texttt{Development.\allowbreak{}AbelianVarietyTheory.\allowbreak{}Proofs.\allowbreak{}GlobalInterpolation.\allowbreak{}WeilUnitaryEntries}, lines 25--32)\par}

{\footnotesize\raggedright\noindent \ref{def:B8-reference-points}.\quad \texttt{GlobalTheory.\allowbreak{}BigTop.\allowbreak{}PtAt.\allowbreak{}PointConditions'} (\texttt{Development.\allowbreak{}AutomorphicArithmetic.\allowbreak{}ThetaCorrespondence.\allowbreak{}ThetaPairLifts.\allowbreak{}ThirdPlaceSeries.\allowbreak{}Realised}, lines 63--80)\\
\texttt{GlobalTheory.\allowbreak{}BigTop.\allowbreak{}PtAt.\allowbreak{}CompatibilityConditions} (\texttt{Development.\allowbreak{}AutomorphicArithmetic.\allowbreak{}ThetaCorrespondence.\allowbreak{}ThetaPairLifts.\allowbreak{}ThirdPlaceSeries.\allowbreak{}Realised}, lines 51--62)\\
\texttt{GlobalTheory.\allowbreak{}BigOn.\allowbreak{}More} (\texttt{Development.\allowbreak{}AbelianVarietyTheory.\allowbreak{}Statements.\allowbreak{}PlaceShiftStatements.\allowbreak{}SingleIndexPlacesAreTwisted}, lines 82--83)\\
\texttt{GlobalTheory.\allowbreak{}BigOn.\allowbreak{}CoreTies} (\texttt{Development.\allowbreak{}AbelianVarietyTheory.\allowbreak{}Definitions.\allowbreak{}Divisors.\allowbreak{}IsDeltaDelAtOwnAt}, lines 19--33)\\
\texttt{GlobalTheory.\allowbreak{}BigOn.\allowbreak{}RecTies} (\texttt{Development.\allowbreak{}AbelianVarietyTheory.\allowbreak{}Statements.\allowbreak{}PlaceShiftStatements.\allowbreak{}SingleIndexPlacesAreTwisted}, lines 63--81)\par}

{\footnotesize\raggedright\noindent \ref{B8.38}.\quad \texttt{GlobalTheory.\allowbreak{}Fc.\allowbreak{}StageZeroRow.\allowbreak{}row\_\allowbreak{}rRows\_\allowbreak{}stage0\_\allowbreak{}alt} (\texttt{Development.\allowbreak{}AutomorphicArithmetic.\allowbreak{}ThetaCorrespondence.\allowbreak{}ThetaPairLifts.\allowbreak{}ThirdPlaceSeries.\allowbreak{}TupAllUsedAtPart}, lines 46--67)\\
\texttt{GlobalTheory.\allowbreak{}Fc.\allowbreak{}row\_\allowbreak{}rRows\_\allowbreak{}of\_\allowbreak{}parts\_\allowbreak{}alt} (\texttt{Development.\allowbreak{}AutomorphicArithmetic.\allowbreak{}ThetaCorrespondence.\allowbreak{}ThetaPairLifts.\allowbreak{}ThirdPlaceSeries.\allowbreak{}RowRRowsOfParts}, lines 21--63)\\
\texttt{GlobalTheory.\allowbreak{}Fc.\allowbreak{}R1.\allowbreak{}SelectionDefiniteRows} (\texttt{Development.\allowbreak{}AutomorphicArithmetic.\allowbreak{}ThetaCorrespondence.\allowbreak{}ThetaPairLifts.\allowbreak{}ThirdPlaceSeries.\allowbreak{}CountsIrreducibilityConstants}, lines 91--93)\par}

{\footnotesize\raggedright\noindent \ref{B8.38b}.\quad \texttt{Formalism.\allowbreak{}L2.\allowbreak{}Indef.\allowbreak{}IndefData.\allowbreak{}ON} (\texttt{Formalism.\allowbreak{}LocalToGlobal.\allowbreak{}IndefiniteForms.\allowbreak{}IndefSide}, lines 318--323)\\
\texttt{Formalism.\allowbreak{}L2.\allowbreak{}Indef.\allowbreak{}Tree.\allowbreak{}on\_\allowbreak{}core} (\texttt{Formalism.\allowbreak{}LocalToGlobal.\allowbreak{}IndefiniteForms.\allowbreak{}IndefiniteTree.\allowbreak{}ExampleSatisfiesConditionsAtFive}, lines 603--872)\\
\texttt{GlobalTheory.\allowbreak{}BigOn.\allowbreak{}canonical\_\allowbreak{}error\_\allowbreak{}real} (\texttt{Development.\allowbreak{}AbelianVarietyTheory.\allowbreak{}Proofs.\allowbreak{}DivisorProofs.\allowbreak{}OnOfOnWithCanonical}, lines 8--62)\\
\texttt{GlobalTheory.\allowbreak{}BigOn.\allowbreak{}OntoPath.\allowbreak{}onto\_\allowbreak{}path\_\allowbreak{}alt} (\texttt{Development.\allowbreak{}AutomorphicArithmetic.\allowbreak{}ThetaCorrespondence.\allowbreak{}ThetaPairLifts.\allowbreak{}ThirdPlaceSeries.\allowbreak{}OntoPath}, lines 18--76)\par}

{\footnotesize\raggedright\noindent \ref{B8.38a}.\quad \texttt{GlobalTheory.\allowbreak{}Published.\allowbreak{}CasselmanWallach} (\texttt{Development.\allowbreak{}AbelianVarietyTheory.\allowbreak{}PublishedTheorems.\allowbreak{}AutomorphicForms₂.\allowbreak{}HasDRepIffHasDerivAt}, lines 98--111)\\
\texttt{PositivityTheory.\allowbreak{}CWB.\allowbreak{}casselmanWallach} (\texttt{Development.\allowbreak{}WeilBound.\allowbreak{}CoreStatements.\allowbreak{}PublishedTheorems.\allowbreak{}CasselmanWallach.\allowbreak{}FinalB}, lines 8--9)\par}

{\footnotesize\raggedright\noindent \ref{B8.38a}.\quad \texttt{PositivityTheory.\allowbreak{}CW.\allowbreak{}e3d} (\texttt{Development.\allowbreak{}WeilBound.\allowbreak{}CoreStatements.\allowbreak{}PublishedTheorems.\allowbreak{}CasselmanWallach.\allowbreak{}E3D}, lines 10--56)\\
\texttt{GlobalTheory.\allowbreak{}CompanionTheory.\allowbreak{}SmoothClosed.\allowbreak{}isClosed\_\allowbreak{}image\_\allowbreak{}of\_\allowbreak{}sobolev\_\allowbreak{}lower} (\texttt{Development.\allowbreak{}AutomorphicSpectrum.\allowbreak{}Proofs.\allowbreak{}BlockDecomposition.\allowbreak{}IsClosedImageOfSobolevLower}, lines 26--60)\\
\texttt{PositivityTheory.\allowbreak{}CW.\allowbreak{}cw\_\allowbreak{}of\_\allowbreak{}steps} (\texttt{Development.\allowbreak{}WeilBound.\allowbreak{}CoreStatements.\allowbreak{}PublishedTheorems.\allowbreak{}CasselmanWallach.\allowbreak{}CasselmanWallachAdmissible}, lines 61--108)\par}

{\footnotesize\raggedright\noindent \ref{B8.39}.\quad \texttt{Formalism.\allowbreak{}MainNotions.\allowbreak{}Closed} (\texttt{Formalism.\allowbreak{}SharpRestriction.\allowbreak{}ParamsHeightRankConstant}, lines 159--161)\\
\texttt{GlobalTheory.\allowbreak{}Fc.\allowbreak{}S0.\allowbreak{}NL} (\texttt{Development.\allowbreak{}AbelianVarietyTheory.\allowbreak{}Definitions.\allowbreak{}GlobalTwisting.\allowbreak{}GeometricWeakTiesOfStrong}, lines 36--36)\par}

{\footnotesize\raggedright\noindent \ref{B8.40}.\quad \texttt{Formalism.\allowbreak{}MainNotions.\allowbreak{}LinkN} (\texttt{Formalism.\allowbreak{}SharpRestriction.\allowbreak{}GlobalTheoremInputs}, lines 45--47)\\
\texttt{Formalism.\allowbreak{}MainNotions.\allowbreak{}ShaFinOdd} (\texttt{Formalism.\allowbreak{}SharpRestriction.\allowbreak{}GlobalTheoremInputs}, lines 25--26)\\
\texttt{Formalism.\allowbreak{}MainNotions.\allowbreak{}vvTerm} (\texttt{Formalism.\allowbreak{}SharpRestriction.\allowbreak{}GlobalTheoremInputs}, lines 31--32)\\
\texttt{Formalism.\allowbreak{}MainNotions.\allowbreak{}etaSum} (\texttt{Formalism.\allowbreak{}SharpRestriction.\allowbreak{}GlobalTheoremInputs}, lines 36--38)\\
\texttt{Formalism.\allowbreak{}MainNotions.\allowbreak{}LinkN\_\allowbreak{}at} (\texttt{Formalism.\allowbreak{}SharpRestriction.\allowbreak{}GlobalTheoremInputs}, lines 40--43)\par}

{\footnotesize\raggedright\noindent \ref{B8.40b}.\quad \texttt{Formalism.\allowbreak{}L2.\allowbreak{}DIV.\allowbreak{}sum\_\allowbreak{}le\_\allowbreak{}etaSum\_\allowbreak{}num} (\texttt{Formalism.\allowbreak{}LocalToGlobal.\allowbreak{}Divisibility.\allowbreak{}SumLeEtaSumNum}, lines 28--58)\par}

{\footnotesize\raggedright\noindent \ref{B8.40a}.\quad \texttt{UpperTheory.\allowbreak{}Top.\allowbreak{}vvRecord} (\texttt{Development.\allowbreak{}AutomorphicArithmetic.\allowbreak{}CharacterBoxes.\allowbreak{}DivinNOfHypV}, lines 17--57)\\
\texttt{UpperTheory.\allowbreak{}Top.\allowbreak{}divinN\_\allowbreak{}at} (\texttt{Development.\allowbreak{}AutomorphicArithmetic.\allowbreak{}CharacterBoxes.\allowbreak{}DivinNOfHypV}, lines 76--115)\\
\texttt{UpperTheory.\allowbreak{}Top.\allowbreak{}divinN\_\allowbreak{}of\_\allowbreak{}hypV} (\texttt{Development.\allowbreak{}AutomorphicArithmetic.\allowbreak{}CharacterBoxes.\allowbreak{}DivinNOfHypV}, lines 117--130)\\
\texttt{Formalism.\allowbreak{}L2.\allowbreak{}DIV.\allowbreak{}sum\_\allowbreak{}le\_\allowbreak{}etaSum\_\allowbreak{}num} (\texttt{Formalism.\allowbreak{}LocalToGlobal.\allowbreak{}Divisibility.\allowbreak{}SumLeEtaSumNum}, lines 28--58)\\
\texttt{Formalism.\allowbreak{}L2.\allowbreak{}DIV.\allowbreak{}inEta3\_\allowbreak{}num} (\texttt{Formalism.\allowbreak{}LocalToGlobal.\allowbreak{}Divisibility.\allowbreak{}SumLeEtaSumNum}, lines 78--85)\par}

{\footnotesize\raggedright\noindent \ref{B8.41p}.\quad \texttt{ReciprocityTheory.\allowbreak{}Gal.\allowbreak{}LocalFiniteness} (\texttt{ReciprocityTheory.\allowbreak{}Hypotheses.\allowbreak{}GalCoh}, lines 195--199)\\
\texttt{ReciprocityTheory.\allowbreak{}Gal.\allowbreak{}LocalTateDuality} (\texttt{ReciprocityTheory.\allowbreak{}Hypotheses.\allowbreak{}GalCoh}, lines 201--205)\\
\texttt{ReciprocityTheory.\allowbreak{}Gal.\allowbreak{}LocalEulerChar} (\texttt{ReciprocityTheory.\allowbreak{}Hypotheses.\allowbreak{}GalCoh}, lines 213--217)\par}

{\footnotesize\raggedright\noindent \ref{B8.41h}.\quad \texttt{GlobalTheory.\allowbreak{}NLink.\allowbreak{}Hypotheses₂} (\texttt{Development.\allowbreak{}AbelianVarietyTheory.\allowbreak{}Definitions.\allowbreak{}NonGaloisDescent.\allowbreak{}BoundAtTwoAndThreeForFamily}, lines 18--25)\par}

{\footnotesize\raggedright\noindent \ref{B8.41h}.\quad \texttt{GlobalTheory.\allowbreak{}NLink.\allowbreak{}SetupAlt} (\texttt{Development.\allowbreak{}AutomorphicArithmetic.\allowbreak{}ThetaCorrespondence.\allowbreak{}ThetaPairLifts.\allowbreak{}ThirdPlaceSeries.\allowbreak{}BoundAtTwoThreeFamilyFon}, lines 13--21)\\
\texttt{GlobalTheory.\allowbreak{}HubX.\allowbreak{}Canon} (\texttt{Development.\allowbreak{}AbelianVarietyTheory.\allowbreak{}Definitions.\allowbreak{}NonGaloisDescent.\allowbreak{}DecompositionInvariantsCondition}, lines 40--47)\\
\texttt{GlobalTheory.\allowbreak{}HubLoc.\allowbreak{}cBnum} (\texttt{Development.\allowbreak{}AbelianVarietyTheory.\allowbreak{}Definitions.\allowbreak{}NonGaloisDescent.\allowbreak{}FinrankMaxUnramifiedSubfieldPadicFinExtPos}, lines 376--377)\\
\texttt{Formalism.\allowbreak{}MainNotions.\allowbreak{}N.\allowbreak{}KillingConditions.\allowbreak{}DataAD} (\texttt{Formalism.\allowbreak{}SharpRestriction.\allowbreak{}LevelLetterCocycles.\allowbreak{}JointConditionsZpowers.\allowbreak{}CommutingEndomorphismCondition}, lines 168--179)\\
\texttt{GlobalTheory.\allowbreak{}NLink.\allowbreak{}rowC} (\texttt{Development.\allowbreak{}AbelianVarietyTheory.\allowbreak{}Definitions.\allowbreak{}NonGaloisDescent.\allowbreak{}BoundAtTwoAndThreeForFamily}, lines 29--30)\par}

{\footnotesize\raggedright\noindent \ref{B8.41}.\quad \texttt{UpperTheory.\allowbreak{}N.\allowbreak{}at23\_\allowbreak{}family\_\allowbreak{}alt} (\texttt{Development.\allowbreak{}AutomorphicArithmetic.\allowbreak{}ThetaCorrespondence.\allowbreak{}ThetaPairLifts.\allowbreak{}ThirdPlaceSeries.\allowbreak{}At23OfAll}, lines 19--20)\\
\texttt{UpperTheory.\allowbreak{}N.\allowbreak{}boundAtTwoThree\_\allowbreak{}linked\_\allowbreak{}of\_\allowbreak{}recT} (\texttt{Development.\allowbreak{}AutomorphicArithmetic.\allowbreak{}ThetaCorrespondence.\allowbreak{}ThetaPairLifts.\allowbreak{}ThirdPlaceSeries.\allowbreak{}At23OfAll}, lines 24--26)\\
\texttt{GlobalTheory.\allowbreak{}NLink.\allowbreak{}BoundAtTwoThree} (\texttt{Development.\allowbreak{}AbelianVarietyTheory.\allowbreak{}Definitions.\allowbreak{}NonGaloisDescent.\allowbreak{}BoundAtTwoThreeFamily}, lines 52--53)\par}

{\footnotesize\raggedright\noindent \ref{B8.41ca}.\quad \texttt{GlobalTheory.\allowbreak{}PrismS.\allowbreak{}exists\_\allowbreak{}prism} (\texttt{Development.\allowbreak{}AbelianVarietyTheory.\allowbreak{}Proofs.\allowbreak{}InertiaEigenvalues.\allowbreak{}ExistsWachPrismOverIntegers}, lines 331--370)\\
\texttt{GlobalTheory.\allowbreak{}Prism.\allowbreak{}LocalValue.\allowbreak{}main} (\texttt{Development.\allowbreak{}AbelianVarietyTheory.\allowbreak{}Proofs.\allowbreak{}InertiaEigenvalues.\allowbreak{}PrismaticFCrystalValueLocalPrismAinfWittTensor}, lines 141--151)\\
\texttt{GlobalTheory.\allowbreak{}WachLattice.\allowbreak{}build\_\allowbreak{}prism} (\texttt{Development.\allowbreak{}AbelianVarietyTheory.\allowbreak{}Proofs.\allowbreak{}InertiaEigenvalues.\allowbreak{}WachLatticesExistAllLayersOfBhattScholze}, lines 26--169)\\
\texttt{GlobalTheory.\allowbreak{}HWL.\allowbreak{}lattice\_\allowbreak{}free} (\texttt{Development.\allowbreak{}AbelianVarietyTheory.\allowbreak{}Proofs.\allowbreak{}InertiaEigenvalues.\allowbreak{}WachLatticeFreeOfCrystalline}, lines 43--74)\\
\texttt{GlobalTheory.\allowbreak{}NL.\allowbreak{}LatLow.\allowbreak{}lat\_\allowbreak{}low} (\texttt{Development.\allowbreak{}AbelianVarietyTheory.\allowbreak{}Proofs.\allowbreak{}InertiaEigenvalues.\allowbreak{}WachLatticeExpansionsStdBelowLevel}, lines 186--219)\\
\texttt{GlobalTheory.\allowbreak{}NL.\allowbreak{}GamModX.\allowbreak{}sub\_\allowbreak{}mem\_\allowbreak{}X\_\allowbreak{}of\_\allowbreak{}val} (\texttt{Development.\allowbreak{}AbelianVarietyTheory.\allowbreak{}Proofs.\allowbreak{}InertiaEigenvalues.\allowbreak{}ExistsSmulOfDvdRepr}, lines 70--96)\par}

{\footnotesize\raggedright\noindent \ref{B8.41cb}.\quad \texttt{GlobalTheory.\allowbreak{}NL.\allowbreak{}InvX.\allowbreak{}sand\_\allowbreak{}of} (\texttt{Development.\allowbreak{}AbelianVarietyTheory.\allowbreak{}Proofs.\allowbreak{}InertiaEigenvalues.\allowbreak{}WachFreeLatticeEqWachLatticeInvertX}, lines 74--205)\\
\texttt{GlobalTheory.\allowbreak{}NL.\allowbreak{}InvX.\allowbreak{}compare} (\texttt{Development.\allowbreak{}AbelianVarietyTheory.\allowbreak{}Proofs.\allowbreak{}InertiaEigenvalues.\allowbreak{}WachFreeLatticeEqWachLatticeInvertX}, lines 207--229)\\
\texttt{GlobalTheory.\allowbreak{}NL.\allowbreak{}FrobeniusMatrixUniqueness.\allowbreak{}main} (\texttt{Development.\allowbreak{}AbelianVarietyTheory.\allowbreak{}Proofs.\allowbreak{}InertiaEigenvalues.\allowbreak{}WachModuleUniquenessCharacterisation}, lines 247--269)\par}

{\footnotesize\raggedright\noindent \ref{B8.41c}.\quad \texttt{GlobalTheory.\allowbreak{}NL.\allowbreak{}CokerKill.\allowbreak{}X\_\allowbreak{}pow\_\allowbreak{}kills\_\allowbreak{}coker} (\texttt{Development.\allowbreak{}AbelianVarietyTheory.\allowbreak{}Proofs.\allowbreak{}InertiaEigenvalues.\allowbreak{}CokernelCrystallineSurjectionKilledByXPow}, lines 411--437)\\
\texttt{Formalism.\allowbreak{}MainNotions.\allowbreak{}LiftingTheory.\allowbreak{}coreWB\_\allowbreak{}of} (\texttt{Formalism.\allowbreak{}SharpRestriction.\allowbreak{}LevelLetterCocycles.\allowbreak{}FrobeniusPowUniformizerFactorisation}, lines 3449--3457)\par}

{\footnotesize\raggedright\noindent \ref{B8.41aa}.\quad \texttt{GlobalTheory.\allowbreak{}Cyc.\allowbreak{}TwistLift.\allowbreak{}lift\_\allowbreak{}cocycle\_\allowbreak{}of\_\allowbreak{}duality} (\texttt{Development.\allowbreak{}AbelianVarietyTheory.\allowbreak{}Proofs.\allowbreak{}InertiaEigenvalues.\allowbreak{}ContinuousFixingSubgroupEquiv}, lines 307--383)\par}

{\footnotesize\raggedright\noindent \ref{B8.41ab}.\quad \texttt{GlobalTheory.\allowbreak{}NL.\allowbreak{}WCont.\allowbreak{}letterwise} (\texttt{Development.\allowbreak{}AbelianVarietyTheory.\allowbreak{}Proofs.\allowbreak{}InertiaEigenvalues.\allowbreak{}ExistsCocycleSolutionConv}, lines 238--323)\\
\texttt{GlobalTheory.\allowbreak{}NL.\allowbreak{}exists\_\allowbreak{}cocycle\_\allowbreak{}solution\_\allowbreak{}conv} (\texttt{Development.\allowbreak{}AbelianVarietyTheory.\allowbreak{}Proofs.\allowbreak{}InertiaEigenvalues.\allowbreak{}ExistsCocycleSolutionConv}, lines 405--423)\par}

{\footnotesize\raggedright\noindent \ref{B8.41a}.\quad \texttt{UpperTheory.\allowbreak{}N.\allowbreak{}locAll} (\texttt{Development.\allowbreak{}AutomorphicArithmetic.\allowbreak{}ResidueCharacters.\allowbreak{}LocAll}, lines 9--13)\\
\texttt{UpperTheory.\allowbreak{}N.\allowbreak{}LiftK} (\texttt{Development.\allowbreak{}AutomorphicArithmetic.\allowbreak{}ResidueCharacters.\allowbreak{}InvariantsKilledOverCyclotomicAlt}, lines 31--34)\\
\texttt{UpperTheory.\allowbreak{}N.\allowbreak{}InvariantsKilledOverCyclotomic} (\texttt{Development.\allowbreak{}AutomorphicArithmetic.\allowbreak{}ResidueCharacters.\allowbreak{}InvariantsKilledOverCyclotomicAlt}, lines 38--41)\\
\texttt{UpperTheory.\allowbreak{}N.\allowbreak{}TBK} (\texttt{Development.\allowbreak{}AutomorphicArithmetic.\allowbreak{}ResidueCharacters.\allowbreak{}InvariantsKilledOverCyclotomicAlt}, lines 43--45)\\
\texttt{Formalism.\allowbreak{}MainNotions.\allowbreak{}N.\allowbreak{}TBAt} (\texttt{Formalism.\allowbreak{}SharpRestriction.\allowbreak{}LevelLetterCocycles.\allowbreak{}QuotientHomNonvanishingBound}, lines 171--187)\par}

{\footnotesize\raggedright\noindent \ref{B8.41b}.\quad \texttt{UpperTheory.\allowbreak{}N.\allowbreak{}conclusion\_\allowbreak{}split\_\allowbreak{}alt} (\texttt{Development.\allowbreak{}AutomorphicArithmetic.\allowbreak{}ThetaCorrespondence.\allowbreak{}ThetaPairLifts.\allowbreak{}ThirdPlaceSeries.\allowbreak{}ConclusionSplitOfHub}, lines 21--82)\\
\texttt{Formalism.\allowbreak{}MainNotions.\allowbreak{}EmbeddingTypes.\allowbreak{}boundSplit} (\texttt{Formalism.\allowbreak{}SharpRestriction.\allowbreak{}LevelLetterCocycles.\allowbreak{}EmbeddingTypes}, lines 293--304)\par}

{\footnotesize\raggedright\noindent \ref{B8.42p}.\quad \texttt{GlobalTheory.\allowbreak{}Published.\allowbreak{}cheb\_\allowbreak{}holds} (\texttt{Development.\allowbreak{}AbelianVarietyTheory.\allowbreak{}Proofs.\allowbreak{}InertiaEigenvalues.\allowbreak{}ChebotarevDegreeOne}, lines 66--68)\par}

{\footnotesize\raggedright\noindent \ref{B8.42pa}.\quad \texttt{Formalism.\allowbreak{}MainNotions.\allowbreak{}AdmissibleThreePrime.\allowbreak{}residual\_\allowbreak{}irreducible\_\allowbreak{}of\_\allowbreak{}theorem\_\allowbreak{}three\_\allowbreak{}prime\_\allowbreak{}alt} (\texttt{Development.\allowbreak{}AutomorphicArithmetic.\allowbreak{}ThetaCorrespondence.\allowbreak{}ThetaPairLifts.\allowbreak{}ThirdPlaceSeries.\allowbreak{}ResidualIrreducibleOfTheoremThreePrime}, lines 12--14)\\
\texttt{Formalism.\allowbreak{}MainNotions.\allowbreak{}N.\allowbreak{}GlobalDatumResidualRepCondition} (\texttt{Formalism.\allowbreak{}SharpRestriction.\allowbreak{}LevelLetterCocycles.\allowbreak{}JointConditionsZpowers.\allowbreak{}ResidualIrreducibleOfTheorem}, lines 13--19)\\
\texttt{Formalism.\allowbreak{}MainNotions.\allowbreak{}N.\allowbreak{}IsRhoBar} (\texttt{Formalism.\allowbreak{}SharpRestriction.\allowbreak{}LevelLetterCocycles.\allowbreak{}DataN}, lines 1325--1331)\par}

{\footnotesize\raggedright\noindent \ref{B8.42pa}.\quad \texttt{UpperTheory.\allowbreak{}GalC.\allowbreak{}GaloisPreconditions} (\texttt{Development.\allowbreak{}AutomorphicArithmetic.\allowbreak{}ThetaCorrespondence.\allowbreak{}ThetaPairLifts.\allowbreak{}ThirdPlaceSeries.\allowbreak{}GaloisPreconditions}, lines 13--31)\\
\texttt{Formalism.\allowbreak{}BasicNotions.\allowbreak{}NTw.\allowbreak{}NonzeroIndexPlace} (\texttt{Formalism.\allowbreak{}SharpRestriction.\allowbreak{}LevelLetterCocycles.\allowbreak{}DataN}, lines 455--460)\\
\texttt{Formalism.\allowbreak{}MainNotions.\allowbreak{}BlochKatoReduction.\allowbreak{}Gen.\allowbreak{}goodE} (\texttt{Formalism.\allowbreak{}SharpRestriction.\allowbreak{}LevelLetterCocycles.\allowbreak{}BlochKatoPhases}, lines 100--109)\par}

{\footnotesize\raggedright\noindent \ref{B8.42}.\quad \texttt{UpperTheory.\allowbreak{}GalC.\allowbreak{}exists\_\allowbreak{}good\_\allowbreak{}cut\_\allowbreak{}facts\_\allowbreak{}tie4} (\texttt{Development.\allowbreak{}AutomorphicArithmetic.\allowbreak{}ThetaCorrespondence.\allowbreak{}ThetaPairLifts.\allowbreak{}ThirdPlaceSeries.\allowbreak{}CFacts}, lines 58--62)\\
\texttt{UpperTheory.\allowbreak{}GalC.\allowbreak{}linkN\_\allowbreak{}of\_\allowbreak{}twGood\_\allowbreak{}galoisPreconditions} (\texttt{Development.\allowbreak{}AutomorphicArithmetic.\allowbreak{}ThetaCorrespondence.\allowbreak{}ThetaPairLifts.\allowbreak{}ThirdPlaceSeries.\allowbreak{}LinkNFam}, lines 37--53)\par}

{\footnotesize\raggedright\noindent \ref{B8.42a}.\quad \texttt{GlobalTheory.\allowbreak{}NLink.\allowbreak{}twist\_\allowbreak{}kills\_\allowbreak{}selmer} (\texttt{Development.\allowbreak{}AbelianVarietyTheory.\allowbreak{}Proofs.\allowbreak{}InertiaEigenvalues.\allowbreak{}ChosenTwistAnnihilatesSelmerGroups}, lines 8--32)\\
\texttt{UpperTheory.\allowbreak{}KStar.\allowbreak{}sum\_\allowbreak{}nuL\_\allowbreak{}le} (\texttt{Development.\allowbreak{}AutomorphicArithmetic.\allowbreak{}CyclotomicBounds.\allowbreak{}CyclotomicOrders.\allowbreak{}NuHost}, lines 37--39)\\
\texttt{GlobalTheory.\allowbreak{}NLink.\allowbreak{}multLe\_\allowbreak{}of\_\allowbreak{}rank} (\texttt{Development.\allowbreak{}AbelianVarietyTheory.\allowbreak{}Proofs.\allowbreak{}InertiaEigenvalues.\allowbreak{}GoodPointMultiplicityLeRankProof}, lines 70--113)\par}

{\footnotesize\raggedright\noindent \ref{B8.42ba}.\quad \texttt{Formalism.\allowbreak{}RepresentationBoundsCohomology.\allowbreak{}N2.\allowbreak{}second\_\allowbreak{}descent\_\allowbreak{}statement\_\allowbreak{}holds} (\texttt{Formalism.\allowbreak{}WachLatticeTwists.\allowbreak{}LagrangeExtensionCocycles.\allowbreak{}ExtensionCocycleProofs.\allowbreak{}SecondDescentStatementHolds}, lines 7--77)\\
\texttt{Formalism.\allowbreak{}RepresentationBoundsCohomology.\allowbreak{}N2.\allowbreak{}unit\_\allowbreak{}lemma\_\allowbreak{}prime\_\allowbreak{}case\_\allowbreak{}holds} (\texttt{Formalism.\allowbreak{}WachLatticeTwists.\allowbreak{}LagrangeExtensionCocycles.\allowbreak{}ExtensionCocycleProofs.\allowbreak{}Letters}, lines 10--96)\\
\texttt{Formalism.\allowbreak{}MainNotions.\allowbreak{}WachCoordinates.\allowbreak{}rF\_\allowbreak{}le} (\texttt{Formalism.\allowbreak{}SharpRestriction.\allowbreak{}LevelLetterCocycles.\allowbreak{}BoundedNumeratorPowerSummability}, lines 100--148)\par}

{\footnotesize\raggedright\noindent \ref{B8.42bb}.\quad \texttt{Formalism.\allowbreak{}ExponentWindows.\allowbreak{}E.\allowbreak{}Full.\allowbreak{}consts\_\allowbreak{}DD} (\texttt{Formalism.\allowbreak{}DistinguishedUnitIdeals.\allowbreak{}DistinguishedOrderProofs.\allowbreak{}ThmEFull}, lines 52--67)\\
\texttt{Formalism.\allowbreak{}ExponentWindows.\allowbreak{}levels\_\allowbreak{}mod\_\allowbreak{}const\_\allowbreak{}iff\_\allowbreak{}of\_\allowbreak{}down\_\allowbreak{}down} (\texttt{Formalism.\allowbreak{}DistinguishedUnitIdeals.\allowbreak{}DistinguishedOrderProofs.\allowbreak{}LevelsModConstIffOfDownDown}, lines 477--634)\\
\texttt{Formalism.\allowbreak{}ExponentWindows.\allowbreak{}E.\allowbreak{}CL.\allowbreak{}lifts\_\allowbreak{}mod\_\allowbreak{}const\_\allowbreak{}of\_\allowbreak{}down\_\allowbreak{}down} (\texttt{Formalism.\allowbreak{}DistinguishedUnitIdeals.\allowbreak{}DistinguishedOrderProofs.\allowbreak{}ArcLiftOfLocalizationData}, lines 395--408)\par}

{\footnotesize\raggedright\noindent \ref{B8.42bc}.\quad \texttt{Formalism.\allowbreak{}MainNotions.\allowbreak{}ElementDenominator.\allowbreak{}CongruenceModules.\allowbreak{}Ynum} (\texttt{Formalism.\allowbreak{}SharpRestriction.\allowbreak{}LevelLetterCocycles.\allowbreak{}LocalTheoremOfGoodExtensions}, lines 42--138)\\
\texttt{Formalism.\allowbreak{}IrreducibleCountingLemmas.\allowbreak{}WindowSums.\allowbreak{}valuationSumBound} (\texttt{Formalism.\allowbreak{}PadicSubsetSums.\allowbreak{}StableGradedCounts.\allowbreak{}LiftingTheExponent}, lines 34--68)\\
\texttt{Formalism.\allowbreak{}MainNotions.\allowbreak{}ElementDenominator.\allowbreak{}CongruenceModules.\allowbreak{}thmA\_\allowbreak{}of\_\allowbreak{}claim} (\texttt{Formalism.\allowbreak{}SharpRestriction.\allowbreak{}LevelLetterCocycles.\allowbreak{}LocalTheoremOfGoodExtensions}, lines 171--202)\par}

{\footnotesize\raggedright\noindent \ref{B8.42bd}.\quad \texttt{GlobalTheory.\allowbreak{}LocalFieldsQD.\allowbreak{}main} (\texttt{Development.\allowbreak{}AbelianVarietyTheory.\allowbreak{}Proofs.\allowbreak{}InertiaEigenvalues.\allowbreak{}LocalFieldsAtAuxiliaryPrimes}, lines 19--29)\\
\texttt{GlobalTheory.\allowbreak{}CrysFromGTwo.\allowbreak{}crys\_\allowbreak{}loc\_\allowbreak{}alt\_\allowbreak{}alt} (\texttt{Development.\allowbreak{}AutomorphicArithmetic.\allowbreak{}ThetaCorrespondence.\allowbreak{}ThetaPairLifts.\allowbreak{}ThirdPlaceSeries.\allowbreak{}CrysLoc}, lines 15--21)\\
\texttt{GlobalTheory.\allowbreak{}HubFon.\allowbreak{}main\_\allowbreak{}alt\_\allowbreak{}alt} (\texttt{Development.\allowbreak{}AutomorphicArithmetic.\allowbreak{}ThetaCorrespondence.\allowbreak{}ThetaPairLifts.\allowbreak{}ThirdPlaceSeries.\allowbreak{}MainAltAltMore}, lines 11--41)\par}

{\footnotesize\raggedright\noindent \ref{B8.42bm}.\quad \texttt{Formalism.\allowbreak{}WachCohomologyFacts.\allowbreak{}V2.\allowbreak{}sgN} (\texttt{Formalism.\allowbreak{}FrobeniusContraction.\allowbreak{}HighestWeightResidueAvoids}, lines 40--42)\\
\texttt{Formalism.\allowbreak{}MainNotions.\allowbreak{}tableLineDegreeTwo.\allowbreak{}PositionHyps} (\texttt{Formalism.\allowbreak{}SharpRestriction.\allowbreak{}LevelLetterCocycles.\allowbreak{}TableLineDegreeTwo}, lines 25--32)\\
\texttt{Formalism.\allowbreak{}WachCohomologyFacts.\allowbreak{}V2.\allowbreak{}ResidueAvoidance} (\texttt{Formalism.\allowbreak{}FrobeniusContraction.\allowbreak{}HighestWeightResidueAvoids}, lines 72--75)\par}

{\footnotesize\raggedright\noindent \ref{B8.42bn}.\quad \texttt{GlobalTheory.\allowbreak{}PositionData.\allowbreak{}position\_\allowbreak{}conditions} (\texttt{Development.\allowbreak{}AbelianVarietyTheory.\allowbreak{}Proofs.\allowbreak{}InertiaEigenvalues.\allowbreak{}JumpPositionConditionsAtLetter}, lines 57--216)\\
\texttt{GlobalTheory.\allowbreak{}PositionData.\allowbreak{}far\_\allowbreak{}from\_\allowbreak{}block} (\texttt{Development.\allowbreak{}AbelianVarietyTheory.\allowbreak{}Proofs.\allowbreak{}InertiaEigenvalues.\allowbreak{}JumpPositionConditionsAtLetter}, lines 42--52)\\
\texttt{GlobalTheory.\allowbreak{}T.\allowbreak{}shiftedJumpUpperBound\_\allowbreak{}le\_\allowbreak{}twist\_\allowbreak{}bound} (\texttt{Development.\allowbreak{}AbelianVarietyTheory.\allowbreak{}Proofs.\allowbreak{}InertiaEigenvalues.\allowbreak{}ShiftedJumpUpperBoundLeTwistBoundStatement}, lines 7--14)\\
\texttt{GlobalTheory.\allowbreak{}PositionData.\allowbreak{}count\_\allowbreak{}US} (\texttt{Development.\allowbreak{}AbelianVarietyTheory.\allowbreak{}Proofs.\allowbreak{}InertiaEigenvalues.\allowbreak{}MemIntPts}, lines 40--47)\par}

{\footnotesize\raggedright\noindent \ref{B8.42bo}.\quad \texttt{Formalism.\allowbreak{}MainNotions.\allowbreak{}tableLineDegreeTwo.\allowbreak{}TopCore} (\texttt{Formalism.\allowbreak{}SharpRestriction.\allowbreak{}LevelLetterCocycles.\allowbreak{}TableLineDegreeTwo}, lines 193--213)\\
\texttt{Formalism.\allowbreak{}MainNotions.\allowbreak{}tableLineDegreeTwo.\allowbreak{}TopHyp} (\texttt{Formalism.\allowbreak{}SharpRestriction.\allowbreak{}LevelLetterCocycles.\allowbreak{}TableLineDegreeTwo}, lines 180--192)\\
\texttt{Formalism.\allowbreak{}MainNotions.\allowbreak{}tableLineDegreeTwo.\allowbreak{}IsAdmCoc} (\texttt{Formalism.\allowbreak{}SharpRestriction.\allowbreak{}LevelLetterCocycles.\allowbreak{}TableLineDegreeTwo}, lines 172--179)\\
\texttt{Formalism.\allowbreak{}MainNotions.\allowbreak{}tableLineDegreeTwo.\allowbreak{}AdmQ} (\texttt{Formalism.\allowbreak{}SharpRestriction.\allowbreak{}LevelLetterCocycles.\allowbreak{}TableLineDegreeTwo}, lines 93--99)\\
\texttt{Formalism.\allowbreak{}MainNotions.\allowbreak{}tableLineDegreeTwo.\allowbreak{}IsFull} (\texttt{Formalism.\allowbreak{}SharpRestriction.\allowbreak{}LevelLetterCocycles.\allowbreak{}TableLineDegreeTwo}, lines 59--64)\\
\texttt{Formalism.\allowbreak{}MainNotions.\allowbreak{}tableLineDegreeTwo.\allowbreak{}TopTheoremLocal} (\texttt{Formalism.\allowbreak{}SharpRestriction.\allowbreak{}LevelLetterCocycles.\allowbreak{}TableLineDegreeTwo}, lines 398--401)\\
\texttt{GlobalTheory.\allowbreak{}TopClosing.\allowbreak{}blochKatoTop\_\allowbreak{}of\_\allowbreak{}localTateDuality\_\allowbreak{}eulerChar\_\allowbreak{}alt} (\texttt{Development.\allowbreak{}AbelianVarietyTheory.\allowbreak{}Proofs.\allowbreak{}InertiaEigenvalues.\allowbreak{}BlochKatoTopOfLocalTateDualityEulerChar}, lines 25--43)\par}

{\footnotesize\raggedright\noindent \ref{B8.42boa}.\quad \texttt{GlobalTheory.\allowbreak{}TopSteps.\allowbreak{}TopSit} (\texttt{Development.\allowbreak{}AbelianVarietyTheory.\allowbreak{}Definitions.\allowbreak{}NonGaloisDescent.\allowbreak{}ResBound}, lines 36--46)\\
\texttt{GlobalTheory.\allowbreak{}AdmLoc.\allowbreak{}Ea} (\texttt{Development.\allowbreak{}AbelianVarietyTheory.\allowbreak{}Definitions.\allowbreak{}NonGaloisDescent.\allowbreak{}MkMemAdmissibleAtIffCocycle}, lines 308--312)\\
\texttt{GlobalTheory.\allowbreak{}FinExtIntegers.\allowbreak{}dvr\_\allowbreak{}complete} (\texttt{Development.\allowbreak{}AbelianVarietyTheory.\allowbreak{}Proofs.\allowbreak{}InertiaEigenvalues.\allowbreak{}DvrComplete}, lines 10--30)\\
\texttt{GlobalTheory.\allowbreak{}T.\allowbreak{}exists\_\allowbreak{}pow\_\allowbreak{}smul\_\allowbreak{}eq\_\allowbreak{}zero\_\allowbreak{}of\_\allowbreak{}finite\_\allowbreak{}of\_\allowbreak{}not\_\allowbreak{}isUnit} (\texttt{Development.\allowbreak{}AbelianVarietyTheory.\allowbreak{}Proofs.\allowbreak{}InertiaEigenvalues.\allowbreak{}FiniteModuleKilledByPowOfNotIsUnitStatement}, lines 21--24)\par}

{\footnotesize\raggedright\noindent \ref{B8.42bod}.\quad \texttt{GlobalTheory.\allowbreak{}PhiGam.\allowbreak{}RingPG} (\texttt{Development.\allowbreak{}AbelianVarietyTheory.\allowbreak{}Definitions.\allowbreak{}NonGaloisDescent.\allowbreak{}PowerSeriesOneAddConstMulXMulInv}, lines 9--24)\\
\texttt{GlobalTheory.\allowbreak{}PhiGam.\allowbreak{}IsLat} (\texttt{Development.\allowbreak{}AbelianVarietyTheory.\allowbreak{}Definitions.\allowbreak{}NonGaloisDescent.\allowbreak{}PowerSeriesOneAddConstMulXMulInv}, lines 29--30)\\
\texttt{GlobalTheory.\allowbreak{}IndefiniteRowSideFullHecke.\allowbreak{}OnePlusPowerCharacters.\allowbreak{}IsCyclotomic} (\texttt{Development.\allowbreak{}AbelianVarietyTheory.\allowbreak{}BlockTheory.\allowbreak{}NonGaloisCusps.\allowbreak{}OnePlusPowerAssembly}, lines 28--35)\\
\texttt{GlobalTheory.\allowbreak{}IndefiniteRowSideFullHecke.\allowbreak{}FrobeniusBinomialReduction.\allowbreak{}IsCyclotomic} (\texttt{Development.\allowbreak{}AbelianVarietyTheory.\allowbreak{}BlockTheory.\allowbreak{}NonGaloisCusps.\allowbreak{}KernelCharacterDelta}, lines 31--37)\\
\texttt{GlobalTheory.\allowbreak{}IndefiniteRowSideFullHecke.\allowbreak{}FrobeniusBinomialReduction.\allowbreak{}frobπ} (\texttt{Development.\allowbreak{}AbelianVarietyTheory.\allowbreak{}BlockTheory.\allowbreak{}NonGaloisCusps.\allowbreak{}KernelCharacterDelta}, lines 21--22)\\
\texttt{GlobalTheory.\allowbreak{}IndefiniteRowSideFullHecke.\allowbreak{}FrobeniusBinomialReduction.\allowbreak{}gamπ} (\texttt{Development.\allowbreak{}AbelianVarietyTheory.\allowbreak{}BlockTheory.\allowbreak{}NonGaloisCusps.\allowbreak{}KernelCharacterDelta}, lines 23--24)\\
\texttt{GlobalTheory.\allowbreak{}IndefiniteRowSideFullHecke.\allowbreak{}FrobeniusBinomialReduction.\allowbreak{}kerPsi} (\texttt{Development.\allowbreak{}AbelianVarietyTheory.\allowbreak{}BlockTheory.\allowbreak{}NonGaloisCusps.\allowbreak{}KernelCharacterDelta}, lines 53--54)\par}

{\footnotesize\raggedright\noindent \ref{B8.42boe}.\quad \texttt{GlobalTheory.\allowbreak{}PsiShiftProof.\allowbreak{}lemmaPsi} (\texttt{Development.\allowbreak{}AbelianVarietyTheory.\allowbreak{}Proofs.\allowbreak{}InertiaEigenvalues.\allowbreak{}ExistsPullBasis}, lines 239--295)\\
\texttt{GlobalTheory.\allowbreak{}IndefiniteRowSideFullHecke.\allowbreak{}OnePlusPowerCharacters.\allowbreak{}PsiMod} (\texttt{Development.\allowbreak{}AbelianVarietyTheory.\allowbreak{}BlockTheory.\allowbreak{}NonGaloisCusps.\allowbreak{}OnePlusPowerAssembly}, lines 52--59)\par}

{\footnotesize\raggedright\noindent \ref{B8.42bof}.\quad \texttt{GlobalTheory.\allowbreak{}PhiGam.\allowbreak{}ModPG} (\texttt{Development.\allowbreak{}AbelianVarietyTheory.\allowbreak{}Definitions.\allowbreak{}NonGaloisDescent.\allowbreak{}PowerSeriesOneAddConstMulXMulInv}, lines 54--73)\\
\texttt{GlobalTheory.\allowbreak{}PhiGam.\allowbreak{}CondC} (\texttt{Development.\allowbreak{}AbelianVarietyTheory.\allowbreak{}Definitions.\allowbreak{}NonGaloisDescent.\allowbreak{}PowerSeriesOneAddConstMulXMulInv}, lines 43--44)\\
\texttt{GlobalTheory.\allowbreak{}Models.\allowbreak{}DFun} (\texttt{Development.\allowbreak{}AbelianVarietyTheory.\allowbreak{}Definitions.\allowbreak{}NonGaloisDescent.\allowbreak{}TwistedModelAdmissibleCocycles}, lines 46--90)\\
\texttt{GlobalTheory.\allowbreak{}Models.\allowbreak{}Lats} (\texttt{Development.\allowbreak{}AbelianVarietyTheory.\allowbreak{}Definitions.\allowbreak{}NonGaloisDescent.\allowbreak{}TwistedModelAdmissibleCocycles}, lines 19--20)\\
\texttt{GlobalTheory.\allowbreak{}Models.\allowbreak{}IsModel} (\texttt{Development.\allowbreak{}AbelianVarietyTheory.\allowbreak{}Definitions.\allowbreak{}NonGaloisDescent.\allowbreak{}TwistedModelAdmissibleCocycles}, lines 99--103)\\
\texttt{GlobalTheory.\allowbreak{}BasePairs.\allowbreak{}RedN} (\texttt{Development.\allowbreak{}AbelianVarietyTheory.\allowbreak{}Proofs.\allowbreak{}InertiaEigenvalues.\allowbreak{}ExistsModThOfPres}, lines 21--27)\\
\texttt{GlobalTheory.\allowbreak{}Present.\allowbreak{}Pres} (\texttt{Development.\allowbreak{}AbelianVarietyTheory.\allowbreak{}Definitions.\allowbreak{}NonGaloisDescent.\allowbreak{}JSmul}, lines 28--52)\\
\texttt{GlobalTheory.\allowbreak{}Models.\allowbreak{}resCount} (\texttt{Development.\allowbreak{}AbelianVarietyTheory.\allowbreak{}Definitions.\allowbreak{}NonGaloisDescent.\allowbreak{}TwistedModelAdmissibleCocycles}, lines 110--112)\\
\texttt{GlobalTheory.\allowbreak{}Present.\allowbreak{}RedLat.\allowbreak{}red\_\allowbreak{}isLat} (\texttt{Development.\allowbreak{}AbelianVarietyTheory.\allowbreak{}Proofs.\allowbreak{}InertiaEigenvalues.\allowbreak{}WachDeltaInvariantsReductionIsLattice}, lines 277--287)\par}

{\footnotesize\raggedright\noindent \ref{B8.42boba}.\quad \texttt{Formalism.\allowbreak{}IrreducibleCountingLemmas.\allowbreak{}TOP.\allowbreak{}carryFloorLeGeom} (\texttt{Formalism.\allowbreak{}PadicSubsetSums.\allowbreak{}StableGradedCounts.\allowbreak{}TOPSgLe}, lines 19--41)\par}

{\footnotesize\raggedright\noindent \ref{B8.42bobb}.\quad \texttt{GlobalTheory.\allowbreak{}BasePairs.\allowbreak{}eq\_\allowbreak{}of\_\allowbreak{}dvd} (\texttt{Development.\allowbreak{}AbelianVarietyTheory.\allowbreak{}Proofs.\allowbreak{}InertiaEigenvalues.\allowbreak{}NonemptyEquivOfSemilinear}, lines 322--341)\par}

{\footnotesize\raggedright\noindent \ref{B8.42bobc}.\quad \texttt{GlobalTheory.\allowbreak{}TopSteps.\allowbreak{}mem\_\allowbreak{}cob\_\allowbreak{}of\_\allowbreak{}repr\_\allowbreak{}mem\_\allowbreak{}herrB} (\texttt{Development.\allowbreak{}AbelianVarietyTheory.\allowbreak{}Proofs.\allowbreak{}InertiaEigenvalues.\allowbreak{}MemCoboundariesOfHerrPairCoboundary}, lines 9--104)\\
\texttt{GlobalTheory.\allowbreak{}ResidueBound.\allowbreak{}exists\_\allowbreak{}independent\_\allowbreak{}admissible\_\allowbreak{}cocycles\_\allowbreak{}residue} (\texttt{Development.\allowbreak{}AbelianVarietyTheory.\allowbreak{}Proofs.\allowbreak{}InertiaEigenvalues.\allowbreak{}ExistsIndependentAdmissibleCocyclesResidue}, lines 175--208)\\
\texttt{GlobalTheory.\allowbreak{}Models.\allowbreak{}DFun} (\texttt{Development.\allowbreak{}AbelianVarietyTheory.\allowbreak{}Definitions.\allowbreak{}NonGaloisDescent.\allowbreak{}TwistedModelAdmissibleCocycles}, lines 46--90)\par}

{\footnotesize\raggedright\noindent \ref{B8.42bobd}.\quad \texttt{GlobalTheory.\allowbreak{}TopVanish.\allowbreak{}vanish\_\allowbreak{}of\_\allowbreak{}lattice} (\texttt{Development.\allowbreak{}AbelianVarietyTheory.\allowbreak{}Proofs.\allowbreak{}InertiaEigenvalues.\allowbreak{}VanishOfLattice}, lines 191--217)\\
\texttt{GlobalTheory.\allowbreak{}TopVanish.\allowbreak{}vanish\_\allowbreak{}of\_\allowbreak{}real} (\texttt{Development.\allowbreak{}AbelianVarietyTheory.\allowbreak{}Proofs.\allowbreak{}InertiaEigenvalues.\allowbreak{}VanishOfLattice}, lines 234--256)\par}

{\footnotesize\raggedright\noindent \ref{B8.42bobe}.\quad \texttt{GlobalTheory.\allowbreak{}UnramAdm.\allowbreak{}unr\_\allowbreak{}subset} (\texttt{Development.\allowbreak{}AbelianVarietyTheory.\allowbreak{}Proofs.\allowbreak{}InertiaEigenvalues.\allowbreak{}AdmissibleQuotientOfFullRepresentation}, lines 480--528)\\
\texttt{GlobalTheory.\allowbreak{}UnramAdm.\allowbreak{}unramified\_\allowbreak{}admissible} (\texttt{Development.\allowbreak{}AbelianVarietyTheory.\allowbreak{}Proofs.\allowbreak{}InertiaEigenvalues.\allowbreak{}AdmissibleQuotientOfFullRepresentation}, lines 537--545)\\
\texttt{GlobalTheory.\allowbreak{}TopSteps.\allowbreak{}ChiOne.\allowbreak{}eq\_\allowbreak{}hom\_\allowbreak{}of\_\allowbreak{}res} (\texttt{Development.\allowbreak{}AbelianVarietyTheory.\allowbreak{}Proofs.\allowbreak{}InertiaEigenvalues.\allowbreak{}ClsEqHomModInertia}, lines 146--253)\\
\texttt{GlobalTheory.\allowbreak{}TopSteps.\allowbreak{}ChiOne.\allowbreak{}cls\_\allowbreak{}eq\_\allowbreak{}homModInertia} (\texttt{Development.\allowbreak{}AbelianVarietyTheory.\allowbreak{}Proofs.\allowbreak{}InertiaEigenvalues.\allowbreak{}ClsEqHomModInertia}, lines 278--290)\par}

{\footnotesize\raggedright\noindent \ref{B8.42bobf}.\quad \texttt{GlobalTheory.\allowbreak{}T.\allowbreak{}admissibleLetterData\_\allowbreak{}window\_\allowbreak{}arithmetic} (\texttt{Development.\allowbreak{}AbelianVarietyTheory.\allowbreak{}Proofs.\allowbreak{}InertiaEigenvalues.\allowbreak{}AdmissibleLetterDataWindowArithmeticStatement}, lines 8--10)\par}

{\footnotesize\raggedright\noindent \ref{B8.42bobg}.\quad \texttt{GlobalTheory.\allowbreak{}Models.\allowbreak{}AdmBound.\allowbreak{}main} (\texttt{Development.\allowbreak{}AbelianVarietyTheory.\allowbreak{}Proofs.\allowbreak{}InertiaEigenvalues.\allowbreak{}AdmissibleCocyclesCardLeSumDims}, lines 73--82)\\
\texttt{GlobalTheory.\allowbreak{}T.\allowbreak{}exists\_\allowbreak{}largest\_\allowbreak{}model'} (\texttt{Development.\allowbreak{}AbelianVarietyTheory.\allowbreak{}Proofs.\allowbreak{}InertiaEigenvalues.\allowbreak{}ExistsLargestModelStatement}, lines 8--10)\\
\texttt{GlobalTheory.\allowbreak{}Models.\allowbreak{}EPolyModel.\allowbreak{}model\_\allowbreak{}ePoly} (\texttt{Development.\allowbreak{}AbelianVarietyTheory.\allowbreak{}Proofs.\allowbreak{}InertiaEigenvalues.\allowbreak{}MapAevalNab}, lines 50--61)\par}

{\footnotesize\raggedright\noindent \ref{B8.42bobh}.\quad \texttt{GlobalTheory.\allowbreak{}Models.\allowbreak{}ModTh.\allowbreak{}hom\_\allowbreak{}chr\_\allowbreak{}eq\_\allowbreak{}zero} (\texttt{Development.\allowbreak{}AbelianVarietyTheory.\allowbreak{}Proofs.\allowbreak{}InertiaEigenvalues.\allowbreak{}HomToCharacterLineEqZero}, lines 25--64)\\
\texttt{GlobalTheory.\allowbreak{}T.\allowbreak{}senWeights\_\allowbreak{}close\_\allowbreak{}of\_\allowbreak{}height\_\allowbreak{}gamma\_\allowbreak{}conditions} (\texttt{Development.\allowbreak{}AbelianVarietyTheory.\allowbreak{}Proofs.\allowbreak{}InertiaEigenvalues.\allowbreak{}SenWeightsCloseOfHeightGammaConditionsStatement}, lines 8--10)\\
\texttt{GlobalTheory.\allowbreak{}T.\allowbreak{}senMultiplicity\_\allowbreak{}subquotient\_\allowbreak{}lattice\_\allowbreak{}le} (\texttt{Development.\allowbreak{}AbelianVarietyTheory.\allowbreak{}Proofs.\allowbreak{}InertiaEigenvalues.\allowbreak{}SenMultiplicitySubquotientLatticeLeStatement}, lines 8--10)\par}

{\footnotesize\raggedright\noindent \ref{B8.42bobi}.\quad \texttt{GlobalTheory.\allowbreak{}CharSplit.\allowbreak{}admissibleCocycle\_\allowbreak{}negCharPower\_\allowbreak{}mem\_\allowbreak{}coboundaries} (\texttt{Development.\allowbreak{}AbelianVarietyTheory.\allowbreak{}Proofs.\allowbreak{}InertiaEigenvalues.\allowbreak{}AdmissibleCocycleNegCharPowerMemCoboundaries}, lines 14--82)\\
\texttt{GlobalTheory.\allowbreak{}TopSteps.\allowbreak{}TrivLine.\allowbreak{}main} (\texttt{Development.\allowbreak{}AbelianVarietyTheory.\allowbreak{}Proofs.\allowbreak{}InertiaEigenvalues.\allowbreak{}TrivialLineAdmissibleCocyclesNotIndependent}, lines 141--175)\par}

{\footnotesize\raggedright\noindent \ref{B8.42bob}.\quad \texttt{GlobalTheory.\allowbreak{}AdmClosure.\allowbreak{}admissible\_\allowbreak{}closed} (\texttt{Development.\allowbreak{}AbelianVarietyTheory.\allowbreak{}Proofs.\allowbreak{}InertiaEigenvalues.\allowbreak{}AdmissibleOfSurjective}, lines 176--180)\\
\texttt{GlobalTheory.\allowbreak{}TopSteps.\allowbreak{}CrysRules.\allowbreak{}crysCalc} (\texttt{Development.\allowbreak{}AbelianVarietyTheory.\allowbreak{}Proofs.\allowbreak{}InertiaEigenvalues.\allowbreak{}CrystallineSumsTwistsCharactersExpectedJumps}, lines 205--235)\\
\texttt{GlobalTheory.\allowbreak{}TopSteps.\allowbreak{}StdAdm.\allowbreak{}stdAdm} (\texttt{Development.\allowbreak{}AbelianVarietyTheory.\allowbreak{}Proofs.\allowbreak{}InertiaEigenvalues.\allowbreak{}AdmissibleStd}, lines 380--385)\\
\texttt{GlobalTheory.\allowbreak{}AdmClasses.\allowbreak{}admissibleLocalSituation\_\allowbreak{}admissible\_\allowbreak{}classes} (\texttt{Development.\allowbreak{}AbelianVarietyTheory.\allowbreak{}Proofs.\allowbreak{}InertiaEigenvalues.\allowbreak{}AdmissibleClassesPropsOfFlatCoefficient}, lines 360--366)\\
\texttt{GlobalTheory.\allowbreak{}T.\allowbreak{}admissible\_\allowbreak{}quotientExtension\_\allowbreak{}iff\_\allowbreak{}admissible\_\allowbreak{}ext} (\texttt{Development.\allowbreak{}AbelianVarietyTheory.\allowbreak{}Proofs.\allowbreak{}InertiaEigenvalues.\allowbreak{}AdmissibleQuotientExtensionIffAdmissibleExtStatement}, lines 21--29)\\
\texttt{GlobalTheory.\allowbreak{}T.\allowbreak{}cocycle\_\allowbreak{}class\_\allowbreak{}admissible} (\texttt{Development.\allowbreak{}AbelianVarietyTheory.\allowbreak{}Proofs.\allowbreak{}InertiaEigenvalues.\allowbreak{}CocycleClassAdmissibleStatement}, lines 21--30)\\
\texttt{GlobalTheory.\allowbreak{}RedFun.\allowbreak{}realD} (\texttt{Development.\allowbreak{}AbelianVarietyTheory.\allowbreak{}Proofs.\allowbreak{}InertiaEigenvalues.\allowbreak{}ExistsPresTensor}, lines 556--571)\\
\texttt{GlobalTheory.\allowbreak{}BasePairs.\allowbreak{}real\_\allowbreak{}of\_\allowbreak{}presented} (\texttt{Development.\allowbreak{}AbelianVarietyTheory.\allowbreak{}Proofs.\allowbreak{}InertiaEigenvalues.\allowbreak{}ExistsModThOfPres}, lines 72--82)\\
\texttt{GlobalTheory.\allowbreak{}TopBridge.\allowbreak{}stmtAt\_\allowbreak{}of\_\allowbreak{}real} (\texttt{Development.\allowbreak{}AbelianVarietyTheory.\allowbreak{}Proofs.\allowbreak{}InertiaEigenvalues.\allowbreak{}TwistedLatticeCohomologyTheoremTOP}, lines 44--66)\\
\texttt{GlobalTheory.\allowbreak{}TopSteps.\allowbreak{}CharVanish.\allowbreak{}stmt\_\allowbreak{}of\_\allowbreak{}real} (\texttt{Development.\allowbreak{}AbelianVarietyTheory.\allowbreak{}Proofs.\allowbreak{}InertiaEigenvalues.\allowbreak{}StmtOfReal}, lines 24--30)\\
\texttt{GlobalTheory.\allowbreak{}TopSteps.\allowbreak{}admissibleUntwistedCharUnramified\_\allowbreak{}of\_\allowbreak{}residueRealisation} (\texttt{Development.\allowbreak{}AbelianVarietyTheory.\allowbreak{}Proofs.\allowbreak{}InertiaEigenvalues.\allowbreak{}AdmissibleUntwistedCharUnramifiedOfResidueRealis}, lines 16--30)\par}

{\footnotesize\raggedright\noindent \ref{B8.42bo1}.\quad \texttt{Formalism.\allowbreak{}MainNotions.\allowbreak{}tableLineDegreeTwo.\allowbreak{}TopTheoremLocal} (\texttt{Formalism.\allowbreak{}SharpRestriction.\allowbreak{}LevelLetterCocycles.\allowbreak{}TableLineDegreeTwo}, lines 398--401)\\
\texttt{GlobalTheory.\allowbreak{}TopBridge.\allowbreak{}twistedLatticeCohomology\_\allowbreak{}theoremTOP} (\texttt{Development.\allowbreak{}AbelianVarietyTheory.\allowbreak{}Proofs.\allowbreak{}InertiaEigenvalues.\allowbreak{}TwistedLatticeCohomologyTheoremTOP}, lines 67--83)\\
\texttt{GlobalTheory.\allowbreak{}TopDict.\allowbreak{}topCore\_\allowbreak{}of\_\allowbreak{}typings} (\texttt{Development.\allowbreak{}AbelianVarietyTheory.\allowbreak{}Proofs.\allowbreak{}InertiaEigenvalues.\allowbreak{}ExactLocalConditionsTopCore}, lines 75--172)\par}

{\footnotesize\raggedright\noindent \ref{B8.42bp}.\quad \texttt{GlobalTheory.\allowbreak{}LocalPositionData.\allowbreak{}position\_\allowbreak{}data\_\allowbreak{}exist} (\texttt{Development.\allowbreak{}AbelianVarietyTheory.\allowbreak{}Proofs.\allowbreak{}InertiaEigenvalues.\allowbreak{}ExistsLocalSlotDataOffExceptionalPrimes}, lines 21--117)\\
\texttt{GlobalTheory.\allowbreak{}IndefiniteRowSideFullHecke.\allowbreak{}BaseSizes.\allowbreak{}Exc} (\texttt{Development.\allowbreak{}AbelianVarietyTheory.\allowbreak{}BlockTheory.\allowbreak{}NonGaloisCusps.\allowbreak{}BaseSizes}, lines 31--33)\\
\texttt{Formalism.\allowbreak{}BasicNotions.\allowbreak{}LevelStarExponents.\allowbreak{}lStar} (\texttt{Formalism.\allowbreak{}LocalRingCardinality.\allowbreak{}DataParams}, lines 38--41)\\
\texttt{GlobalTheory.\allowbreak{}Pass.\allowbreak{}cutCoc} (\texttt{Development.\allowbreak{}AbelianVarietyTheory.\allowbreak{}Definitions.\allowbreak{}NonGaloisDescent.\allowbreak{}CentralIdemCoeff}, lines 122--134)\\
\texttt{GlobalTheory.\allowbreak{}UseLoc.\allowbreak{}Cases.\allowbreak{}main} (\texttt{Development.\allowbreak{}AbelianVarietyTheory.\allowbreak{}Proofs.\allowbreak{}InertiaEigenvalues.\allowbreak{}UnramifiedEnlargementCutCocyclesBlochKato}, lines 126--147)\par}

{\footnotesize\raggedright\noindent \ref{B8.42bq}.\quad \texttt{Formalism.\allowbreak{}MainNotions.\allowbreak{}N.\allowbreak{}Ln34.\allowbreak{}killA} (\texttt{Formalism.\allowbreak{}SharpRestriction.\allowbreak{}LevelLetterCocycles.\allowbreak{}LocalFormOfThreeInputs}, lines 28--39)\\
\texttt{Formalism.\allowbreak{}MainNotions.\allowbreak{}ElementDenominator.\allowbreak{}CongruenceModules.\allowbreak{}YA} (\texttt{Formalism.\allowbreak{}SharpRestriction.\allowbreak{}LevelLetterCocycles.\allowbreak{}LocalTheoremOfGoodExtensions}, lines 150--152)\\
\texttt{Formalism.\allowbreak{}MainNotions.\allowbreak{}ElementDenominator.\allowbreak{}thmA} (\texttt{Formalism.\allowbreak{}SharpRestriction.\allowbreak{}LevelLetterCocycles.\allowbreak{}ElementDenominator}, lines 36--48)\\
\texttt{Formalism.\allowbreak{}WeaklyAdmissibleFiltrations.\allowbreak{}LatHom.\allowbreak{}ImF} (\texttt{Formalism.\allowbreak{}WeaklyAdmissibleFiltrations.\allowbreak{}Statement}, lines 260--261)\\
\texttt{GlobalTheory.\allowbreak{}theoremA\_\allowbreak{}closing} (\texttt{Development.\allowbreak{}AbelianVarietyTheory.\allowbreak{}Proofs.\allowbreak{}InertiaEigenvalues.\allowbreak{}LocalLiftingTheoremAOfBlochKatoTateDualityEulerC}, lines 14--17)\par}

{\footnotesize\raggedright\noindent \ref{B8.42bqa}.\quad \texttt{Formalism.\allowbreak{}MainNotions.\allowbreak{}N.\allowbreak{}Ln34.\allowbreak{}DvdPf.\allowbreak{}dvdCondition} (\texttt{Formalism.\allowbreak{}SharpRestriction.\allowbreak{}LevelLetterCocycles.\allowbreak{}DvrLatticeRank.\allowbreak{}DvdCondition}, lines 85--102)\\
\texttt{Formalism.\allowbreak{}MainNotions.\allowbreak{}N.\allowbreak{}Ln34.\allowbreak{}quotient\_\allowbreak{}jumps\_\allowbreak{}subset} (\texttt{Formalism.\allowbreak{}SharpRestriction.\allowbreak{}LevelLetterCocycles.\allowbreak{}DvrLatticeRank.\allowbreak{}QuotPf}, lines 62--94)\\
\texttt{Formalism.\allowbreak{}MainNotions.\allowbreak{}N.\allowbreak{}Ln34.\allowbreak{}normR\_\allowbreak{}holds} (\texttt{Formalism.\allowbreak{}SharpRestriction.\allowbreak{}LevelLetterCocycles.\allowbreak{}NormAHolds}, lines 55--58)\\
\texttt{Formalism.\allowbreak{}MainNotions.\allowbreak{}N.\allowbreak{}Ln34.\allowbreak{}uniformiser\_\allowbreak{}ring\_\allowbreak{}of\_\allowbreak{}hypotheses} (\texttt{Formalism.\allowbreak{}SharpRestriction.\allowbreak{}LevelLetterCocycles.\allowbreak{}DvrLatticeRank.\allowbreak{}UniformiserRingOfHypotheses}, lines 9--26)\\
\texttt{Formalism.\allowbreak{}MainNotions.\allowbreak{}N.\allowbreak{}Ln34.\allowbreak{}ClassMapPf.\allowbreak{}classmap\_\allowbreak{}pf} (\texttt{Formalism.\allowbreak{}SharpRestriction.\allowbreak{}LevelLetterCocycles.\allowbreak{}DvrLatticeRank.\allowbreak{}ClassMapPf}, lines 248--251)\\
\texttt{Formalism.\allowbreak{}MainNotions.\allowbreak{}N.\allowbreak{}Ln34.\allowbreak{}LamPS} (\texttt{Formalism.\allowbreak{}SharpRestriction.\allowbreak{}LevelLetterCocycles.\allowbreak{}LamPS}, lines 33--36)\\
\texttt{Formalism.\allowbreak{}MainNotions.\allowbreak{}N.\allowbreak{}Ln34.\allowbreak{}JZ} (\texttt{Formalism.\allowbreak{}SharpRestriction.\allowbreak{}LevelLetterCocycles.\allowbreak{}LamPS}, lines 57--59)\\
\texttt{Formalism.\allowbreak{}MainNotions.\allowbreak{}DeltaInvariantLattices.\allowbreak{}p2aK} (\texttt{Formalism.\allowbreak{}SharpRestriction.\allowbreak{}LevelLetterCocycles.\allowbreak{}DeltaInvariantLattices}, lines 308--313)\\
\texttt{Formalism.\allowbreak{}MainNotions.\allowbreak{}N.\allowbreak{}Ln34.\allowbreak{}UnivLinkPf.\allowbreak{}univLinkPf\_\allowbreak{}local} (\texttt{Formalism.\allowbreak{}SharpRestriction.\allowbreak{}LevelLetterCocycles.\allowbreak{}DvrLatticeRank.\allowbreak{}UnivLinkPfB}, lines 186--277)\\
\texttt{Formalism.\allowbreak{}MainNotions.\allowbreak{}N.\allowbreak{}Ln34.\allowbreak{}splitR} (\texttt{Formalism.\allowbreak{}SharpRestriction.\allowbreak{}LevelLetterCocycles.\allowbreak{}MixedPositionHypothesis}, lines 44--54)\par}

{\footnotesize\raggedright\noindent \ref{B8.42bqd}.\quad \texttt{Formalism.\allowbreak{}BasicNotions.\allowbreak{}LocWach} (\texttt{Formalism.\allowbreak{}SharpRestriction.\allowbreak{}LevelLetterCocycles.\allowbreak{}DataNLoc}, lines 164--312)\\
\texttt{Formalism.\allowbreak{}MainNotions.\allowbreak{}DeltaInvariantLattices.\allowbreak{}ptK} (\texttt{Formalism.\allowbreak{}SharpRestriction.\allowbreak{}LevelLetterCocycles.\allowbreak{}DeltaInvariantLattices}, lines 504--509)\\
\texttt{Formalism.\allowbreak{}MainNotions.\allowbreak{}N.\allowbreak{}Ln34.\allowbreak{}valAt} (\texttt{Formalism.\allowbreak{}SharpRestriction.\allowbreak{}LevelLetterCocycles.\allowbreak{}DownSlotProposition}, lines 22--27)\\
\texttt{Formalism.\allowbreak{}MainNotions.\allowbreak{}DeltaInvariantLattices.\allowbreak{}FilC} (\texttt{Formalism.\allowbreak{}SharpRestriction.\allowbreak{}LevelLetterCocycles.\allowbreak{}DeltaInvariantLattices}, lines 156--158)\\
\texttt{Formalism.\allowbreak{}MainNotions.\allowbreak{}N.\allowbreak{}Ln34.\allowbreak{}LcQ} (\texttt{Formalism.\allowbreak{}SharpRestriction.\allowbreak{}LevelLetterCocycles.\allowbreak{}LamPS}, lines 49--51)\\
\texttt{Formalism.\allowbreak{}MainNotions.\allowbreak{}N.\allowbreak{}Ln34.\allowbreak{}LA} (\texttt{Formalism.\allowbreak{}SharpRestriction.\allowbreak{}LevelLetterCocycles.\allowbreak{}DownSlotProposition}, lines 36--37)\\
\texttt{Formalism.\allowbreak{}MainNotions.\allowbreak{}N.\allowbreak{}Ln34.\allowbreak{}LamLA} (\texttt{Formalism.\allowbreak{}SharpRestriction.\allowbreak{}LevelLetterCocycles.\allowbreak{}DownSlotProposition}, lines 50--53)\\
\texttt{Formalism.\allowbreak{}MainNotions.\allowbreak{}N.\allowbreak{}Ln34.\allowbreak{}QLamLA} (\texttt{Formalism.\allowbreak{}SharpRestriction.\allowbreak{}LevelLetterCocycles.\allowbreak{}DownSlotProposition}, lines 96--98)\\
\texttt{Formalism.\allowbreak{}MainNotions.\allowbreak{}N.\allowbreak{}Ln34.\allowbreak{}QLamP} (\texttt{Formalism.\allowbreak{}SharpRestriction.\allowbreak{}LevelLetterCocycles.\allowbreak{}DownSlotProposition}, lines 99--101)\par}

{\footnotesize\raggedright\noindent \ref{B8.42bqb}.\quad \texttt{Formalism.\allowbreak{}MainNotions.\allowbreak{}N.\allowbreak{}Ln34.\allowbreak{}IsRepL} (\texttt{Formalism.\allowbreak{}SharpRestriction.\allowbreak{}LevelLetterCocycles.\allowbreak{}DownSlotProposition}, lines 54--58)\\
\texttt{Formalism.\allowbreak{}MainNotions.\allowbreak{}N.\allowbreak{}Ln34.\allowbreak{}expS} (\texttt{Formalism.\allowbreak{}SharpRestriction.\allowbreak{}LevelLetterCocycles.\allowbreak{}DownSlotProposition}, lines 65--68)\\
\texttt{Formalism.\allowbreak{}MainNotions.\allowbreak{}N.\allowbreak{}Ln34.\allowbreak{}NLset} (\texttt{Formalism.\allowbreak{}SharpRestriction.\allowbreak{}LevelLetterCocycles.\allowbreak{}DownSlotProposition}, lines 69--73)\\
\texttt{Formalism.\allowbreak{}MainNotions.\allowbreak{}N.\allowbreak{}Ln34.\allowbreak{}IsFrame} (\texttt{Formalism.\allowbreak{}SharpRestriction.\allowbreak{}LevelLetterCocycles.\allowbreak{}DownSlotProposition}, lines 122--152)\\
\texttt{Formalism.\allowbreak{}MainNotions.\allowbreak{}N.\allowbreak{}Ln34.\allowbreak{}InvL} (\texttt{Formalism.\allowbreak{}SharpRestriction.\allowbreak{}LevelLetterCocycles.\allowbreak{}DownSlotProposition}, lines 82--85)\\
\texttt{Formalism.\allowbreak{}MainNotions.\allowbreak{}N.\allowbreak{}Ln34.\allowbreak{}IIrest} (\texttt{Formalism.\allowbreak{}SharpRestriction.\allowbreak{}LevelLetterCocycles.\allowbreak{}DownSlotProposition}, lines 86--89)\\
\texttt{Formalism.\allowbreak{}MainNotions.\allowbreak{}N.\allowbreak{}Ln34.\allowbreak{}IsCoord} (\texttt{Formalism.\allowbreak{}SharpRestriction.\allowbreak{}LevelLetterCocycles.\allowbreak{}DownSlotProposition}, lines 163--171)\\
\texttt{Formalism.\allowbreak{}MainNotions.\allowbreak{}N.\allowbreak{}Ln34.\allowbreak{}FilL} (\texttt{Formalism.\allowbreak{}SharpRestriction.\allowbreak{}LevelLetterCocycles.\allowbreak{}DownSlotProposition}, lines 28--31)\par}

{\footnotesize\raggedright\noindent \ref{B8.42bqe}.\quad \texttt{Formalism.\allowbreak{}MainNotions.\allowbreak{}N.\allowbreak{}Ln34.\allowbreak{}ModelL} (\texttt{Formalism.\allowbreak{}SharpRestriction.\allowbreak{}LevelLetterCocycles.\allowbreak{}ModelKStep}, lines 29--41)\\
\texttt{Formalism.\allowbreak{}MainNotions.\allowbreak{}N.\allowbreak{}Ln34.\allowbreak{}modell\_\allowbreak{}pf} (\texttt{Formalism.\allowbreak{}SharpRestriction.\allowbreak{}LevelLetterCocycles.\allowbreak{}DvrLatticeRank.\allowbreak{}ModelLPf}, lines 11--35)\\
\texttt{Formalism.\allowbreak{}MainNotions.\allowbreak{}N.\allowbreak{}Ln34.\allowbreak{}ModelF} (\texttt{Formalism.\allowbreak{}SharpRestriction.\allowbreak{}LevelLetterCocycles.\allowbreak{}DownSlotProposition}, lines 198--212)\\
\texttt{Formalism.\allowbreak{}MainNotions.\allowbreak{}N.\allowbreak{}Ln34.\allowbreak{}modelflink\_\allowbreak{}pf} (\texttt{Formalism.\allowbreak{}SharpRestriction.\allowbreak{}LevelLetterCocycles.\allowbreak{}DvrLatticeRank.\allowbreak{}ModelFPf}, lines 9--68)\\
\texttt{Formalism.\allowbreak{}MainNotions.\allowbreak{}N.\allowbreak{}Ln34.\allowbreak{}modelqlink\_\allowbreak{}pf} (\texttt{Formalism.\allowbreak{}SharpRestriction.\allowbreak{}LevelLetterCocycles.\allowbreak{}DvrLatticeRank.\allowbreak{}ModelQPf}, lines 14--93)\par}

{\footnotesize\raggedright\noindent \ref{B8.42bqca}.\quad \texttt{Formalism.\allowbreak{}MainNotions.\allowbreak{}N.\allowbreak{}Ln34.\allowbreak{}LevelExpansions.\allowbreak{}dv} (\texttt{Formalism.\allowbreak{}SharpRestriction.\allowbreak{}LevelLetterCocycles.\allowbreak{}DvrLatticeRank.\allowbreak{}ExpansionLowCoeffsIntegral}, lines 199--220)\\
\texttt{Formalism.\allowbreak{}MainNotions.\allowbreak{}N.\allowbreak{}Ln34.\allowbreak{}LevelExpansions.\allowbreak{}simple} (\texttt{Formalism.\allowbreak{}SharpRestriction.\allowbreak{}LevelLetterCocycles.\allowbreak{}DvrLatticeRank.\allowbreak{}ExpansionLowCoeffsIntegral}, lines 234--249)\par}

{\footnotesize\raggedright\noindent \ref{B8.42bqcb}.\quad \texttt{Formalism.\allowbreak{}MainNotions.\allowbreak{}N.\allowbreak{}Ln34.\allowbreak{}LevelExpansions.\allowbreak{}fact2} (\texttt{Formalism.\allowbreak{}SharpRestriction.\allowbreak{}LevelLetterCocycles.\allowbreak{}DvrLatticeRank.\allowbreak{}ExpansionLowCoeffsIntegral}, lines 252--328)\\
\texttt{Formalism.\allowbreak{}MainNotions.\allowbreak{}N.\allowbreak{}Ln34.\allowbreak{}LevelExpansions.\allowbreak{}caseC} (\texttt{Formalism.\allowbreak{}SharpRestriction.\allowbreak{}LevelLetterCocycles.\allowbreak{}DvrLatticeRank.\allowbreak{}PolyExpand}, lines 417--452)\par}

{\footnotesize\raggedright\noindent \ref{B8.42bqc}.\quad \texttt{Formalism.\allowbreak{}MainNotions.\allowbreak{}N.\allowbreak{}Ln34.\allowbreak{}u1link\_\allowbreak{}pf} (\texttt{Formalism.\allowbreak{}SharpRestriction.\allowbreak{}LevelLetterCocycles.\allowbreak{}DvrLatticeRank.\allowbreak{}U1linkPf}, lines 16--68)\\
\texttt{Formalism.\allowbreak{}MainNotions.\allowbreak{}N.\allowbreak{}Ln34.\allowbreak{}FirstPropositionFrame} (\texttt{Formalism.\allowbreak{}SharpRestriction.\allowbreak{}LevelLetterCocycles.\allowbreak{}DownSlotProposition}, lines 216--232)\\
\texttt{Formalism.\allowbreak{}MainNotions.\allowbreak{}N.\allowbreak{}Ln34.\allowbreak{}SecondPropositionFree} (\texttt{Formalism.\allowbreak{}SharpRestriction.\allowbreak{}LevelLetterCocycles.\allowbreak{}DownSlotProposition}, lines 233--238)\\
\texttt{Formalism.\allowbreak{}MainNotions.\allowbreak{}N.\allowbreak{}Ln34.\allowbreak{}ThirdPropositionSlot} (\texttt{Formalism.\allowbreak{}SharpRestriction.\allowbreak{}LevelLetterCocycles.\allowbreak{}DownSlotProposition}, lines 244--251)\\
\texttt{Formalism.\allowbreak{}MainNotions.\allowbreak{}N.\allowbreak{}Ln34.\allowbreak{}kapPiPf\_\allowbreak{}main} (\texttt{Formalism.\allowbreak{}SharpRestriction.\allowbreak{}LevelLetterCocycles.\allowbreak{}DvrLatticeRank.\allowbreak{}KapPiPf}, lines 217--284)\\
\texttt{Formalism.\allowbreak{}MainNotions.\allowbreak{}N.\allowbreak{}Ln34.\allowbreak{}UnivLinkPf.\allowbreak{}univLinkPf\_\allowbreak{}cong} (\texttt{Formalism.\allowbreak{}SharpRestriction.\allowbreak{}LevelLetterCocycles.\allowbreak{}DvrLatticeRank.\allowbreak{}UnivLinkPfB}, lines 39--110)\par}

{\footnotesize\raggedright\noindent \ref{B8.42bq1}.\quad \texttt{Formalism.\allowbreak{}MainNotions.\allowbreak{}N.\allowbreak{}Ln34.\allowbreak{}killA} (\texttt{Formalism.\allowbreak{}SharpRestriction.\allowbreak{}LevelLetterCocycles.\allowbreak{}LocalFormOfThreeInputs}, lines 28--39)\\
\texttt{Formalism.\allowbreak{}MainNotions.\allowbreak{}N.\allowbreak{}Ln34.\allowbreak{}killR\_\allowbreak{}of\_\allowbreak{}steps} (\texttt{Formalism.\allowbreak{}SharpRestriction.\allowbreak{}LevelLetterCocycles.\allowbreak{}MixedPositionHypothesis}, lines 105--171)\\
\texttt{Formalism.\allowbreak{}MainNotions.\allowbreak{}N.\allowbreak{}Ln34.\allowbreak{}lifting\_\allowbreak{}theorem\_\allowbreak{}of\_\allowbreak{}uniform\_\allowbreak{}link} (\texttt{Formalism.\allowbreak{}SharpRestriction.\allowbreak{}LevelLetterCocycles.\allowbreak{}DvrLatticeRank.\allowbreak{}MainPf}, lines 38--47)\\
\texttt{GlobalTheory.\allowbreak{}TheoremAFinal.\allowbreak{}integralityFromMetric\_\allowbreak{}of\_\allowbreak{}le} (\texttt{Development.\allowbreak{}AbelianVarietyTheory.\allowbreak{}Proofs.\allowbreak{}InertiaEigenvalues.\allowbreak{}LocalLiftingOfEtaleExactDeRhamCrystallineConst}, lines 16--21)\\
\texttt{GlobalTheory.\allowbreak{}theoremA\_\allowbreak{}closing} (\texttt{Development.\allowbreak{}AbelianVarietyTheory.\allowbreak{}Proofs.\allowbreak{}InertiaEigenvalues.\allowbreak{}LocalLiftingTheoremAOfBlochKatoTateDualityEulerC}, lines 14--17)\par}

{\footnotesize\raggedright\noindent \ref{B8.42br}.\quad \texttt{Formalism.\allowbreak{}MainNotions.\allowbreak{}ElementDenominator.\allowbreak{}CongruenceModules.\allowbreak{}thmA\_\allowbreak{}of\_\allowbreak{}claim} (\texttt{Formalism.\allowbreak{}SharpRestriction.\allowbreak{}LevelLetterCocycles.\allowbreak{}LocalTheoremOfGoodExtensions}, lines 171--202)\\
\texttt{Formalism.\allowbreak{}MainNotions.\allowbreak{}N.\allowbreak{}Ln34.\allowbreak{}claimA\_\allowbreak{}of\_\allowbreak{}steps} (\texttt{Formalism.\allowbreak{}SharpRestriction.\allowbreak{}LevelLetterCocycles.\allowbreak{}LocalFormOfThreeInputs}, lines 40--55)\\
\texttt{Formalism.\allowbreak{}MainNotions.\allowbreak{}ElementDenominator.\allowbreak{}CongruenceModules.\allowbreak{}claimA} (\texttt{Formalism.\allowbreak{}SharpRestriction.\allowbreak{}LevelLetterCocycles.\allowbreak{}LocalTheoremOfGoodExtensions}, lines 157--167)\par}

{\footnotesize\raggedright\noindent \ref{B8.42bs}.\quad \texttt{Formalism.\allowbreak{}MainNotions.\allowbreak{}PairedCharacterCondition.\allowbreak{}PairedCharacterConditionAlt} (\texttt{Development.\allowbreak{}AutomorphicArithmetic.\allowbreak{}ThetaCorrespondence.\allowbreak{}ThetaPairLifts.\allowbreak{}ThirdPlaceSeries.\allowbreak{}LocalGlobalCompatibility}, lines 15--27)\\
\texttt{Formalism.\allowbreak{}MainNotions.\allowbreak{}PairedCharacterCondition.\allowbreak{}rhs} (\texttt{Formalism.\allowbreak{}SharpRestriction.\allowbreak{}LevelLetterCocycles.\allowbreak{}PairedCharacterCondition}, lines 151--158)\\
\texttt{Formalism.\allowbreak{}MainNotions.\allowbreak{}PairedCharacterCondition.\allowbreak{}HostJ} (\texttt{Formalism.\allowbreak{}SharpRestriction.\allowbreak{}LevelLetterCocycles.\allowbreak{}PairedCharacterCondition}, lines 132--143)\\
\texttt{GlobalTheory.\allowbreak{}ClosingLocalInputs.\allowbreak{}main\_\allowbreak{}alt\_\allowbreak{}alt} (\texttt{Development.\allowbreak{}AutomorphicArithmetic.\allowbreak{}ThetaCorrespondence.\allowbreak{}ThetaPairLifts.\allowbreak{}ThirdPlaceSeries.\allowbreak{}AdmG14OfAux}, lines 15--32)\par}

{\footnotesize\raggedright\noindent \ref{B8.42bs1a}.\quad \texttt{GlobalTheory.\allowbreak{}EtaFromCut.\allowbreak{}lenUnit\_\allowbreak{}le\_\allowbreak{}of\_\allowbreak{}family} (\texttt{Development.\allowbreak{}AbelianVarietyTheory.\allowbreak{}Proofs.\allowbreak{}InertiaEigenvalues.\allowbreak{}LenUnitLeOfFamily}, lines 64--94)\par}

{\footnotesize\raggedright\noindent \ref{B8.42bs1b}.\quad \texttt{GlobalTheory.\allowbreak{}CutFam.\allowbreak{}perPoint} (\texttt{Development.\allowbreak{}AbelianVarietyTheory.\allowbreak{}Proofs.\allowbreak{}InertiaEigenvalues.\allowbreak{}ExistsOfMemCcAt}, lines 114--154)\par}

{\footnotesize\raggedright\noindent \ref{B8.42bs1c}.\quad \texttt{GlobalTheory.\allowbreak{}ClosingStatement.\allowbreak{}main} (\texttt{Development.\allowbreak{}AbelianVarietyTheory.\allowbreak{}Proofs.\allowbreak{}InertiaEigenvalues.\allowbreak{}LocalAbsorptionLengthBoundOfBlochKato}, lines 15--26)\\
\texttt{Formalism.\allowbreak{}MainNotions.\allowbreak{}BlochKatoDuality.\allowbreak{}second\_\allowbreak{}proposition\_\allowbreak{}every\_\allowbreak{}localisation} (\texttt{Formalism.\allowbreak{}SharpRestriction.\allowbreak{}LevelLetterCocycles.\allowbreak{}RamificationTwoAndConjugationAtQAndD}, lines 1405--1408)\par}

{\footnotesize\raggedright\noindent \ref{B8.42bs1d}.\quad \texttt{GlobalTheory.\allowbreak{}InertiaEigenvalues.\allowbreak{}theoremT\_\allowbreak{}all\_\allowbreak{}records} (\texttt{Development.\allowbreak{}AbelianVarietyTheory.\allowbreak{}Proofs.\allowbreak{}InertiaEigenvalues.\allowbreak{}CrystallinePeriodBoundAllConstantGeForty}, lines 13--20)\\
\texttt{Formalism.\allowbreak{}MainNotions.\allowbreak{}ElementDenominator.\allowbreak{}thmT} (\texttt{Formalism.\allowbreak{}SharpRestriction.\allowbreak{}LevelLetterCocycles.\allowbreak{}ElementDenominator}, lines 178--198)\\
\texttt{Formalism.\allowbreak{}BasicNotions.\allowbreak{}pc} (\texttt{Formalism.\allowbreak{}SharpRestriction.\allowbreak{}LevelLetterCocycles.\allowbreak{}DataNLoc}, lines 657--694)\\
\texttt{Formalism.\allowbreak{}WeaklyAdmissibleFiltrations.\allowbreak{}PotCrysData.\allowbreak{}InertiaThrough} (\texttt{Formalism.\allowbreak{}WeaklyAdmissibleFiltrations.\allowbreak{}StatementPC}, lines 73--76)\par}

{\footnotesize\raggedright\noindent \ref{B8.42bs1e}.\quad \texttt{GlobalTheory.\allowbreak{}StepLiftPf.\allowbreak{}liftFam\_\allowbreak{}of\_\allowbreak{}eq} (\texttt{Development.\allowbreak{}AbelianVarietyTheory.\allowbreak{}Proofs.\allowbreak{}InertiaEigenvalues.\allowbreak{}LiftFMonoLoc}, lines 156--199)\\
\texttt{GlobalTheory.\allowbreak{}StepLiftPf.\allowbreak{}liftFam\_\allowbreak{}core} (\texttt{Development.\allowbreak{}AbelianVarietyTheory.\allowbreak{}Proofs.\allowbreak{}InertiaEigenvalues.\allowbreak{}LiftFMonoLoc}, lines 111--145)\par}

{\footnotesize\raggedright\noindent \ref{B8.42bs1f}.\quad \texttt{GlobalTheory.\allowbreak{}Pass.\allowbreak{}LocDirect.\allowbreak{}main} (\texttt{Development.\allowbreak{}AbelianVarietyTheory.\allowbreak{}Proofs.\allowbreak{}InertiaEigenvalues.\allowbreak{}LocalCountDirectLiftingProof}, lines 20--77)\\
\texttt{GlobalTheory.\allowbreak{}Pass.\allowbreak{}LocDual.\allowbreak{}main} (\texttt{Development.\allowbreak{}AbelianVarietyTheory.\allowbreak{}Proofs.\allowbreak{}InertiaEigenvalues.\allowbreak{}LocalCountDualLiftingProof}, lines 30--56)\par}

{\footnotesize\raggedright\noindent \ref{B8.42bs1}.\quad \texttt{GlobalTheory.\allowbreak{}ClosingLocalInputs.\allowbreak{}main\_\allowbreak{}alt\_\allowbreak{}alt} (\texttt{Development.\allowbreak{}AutomorphicArithmetic.\allowbreak{}ThetaCorrespondence.\allowbreak{}ThetaPairLifts.\allowbreak{}ThirdPlaceSeries.\allowbreak{}AdmG14OfAux}, lines 15--32)\\
\texttt{GlobalTheory.\allowbreak{}T.\allowbreak{}admT\_\allowbreak{}of\_\allowbreak{}artin\_\allowbreak{}reciprocity\_\allowbreak{}and\_\allowbreak{}local\_\allowbreak{}conditions} (\texttt{Development.\allowbreak{}AutomorphicArithmetic.\allowbreak{}ThetaCorrespondence.\allowbreak{}ThetaPairLifts.\allowbreak{}ThirdPlaceSeries.\allowbreak{}ConclusionOfAdmissiblesStatement}, lines 11--12)\\
\texttt{GlobalTheory.\allowbreak{}Pass.\allowbreak{}cutBound\_\allowbreak{}of\_\allowbreak{}steps\_\allowbreak{}alt} (\texttt{Development.\allowbreak{}AutomorphicArithmetic.\allowbreak{}ThetaCorrespondence.\allowbreak{}ThetaPairLifts.\allowbreak{}ThirdPlaceSeries.\allowbreak{}CutBoundOfSteps}, lines 15--78)\\
\texttt{GlobalTheory.\allowbreak{}EtaFromCut.\allowbreak{}etaBound\_\allowbreak{}of\_\allowbreak{}cut} (\texttt{Development.\allowbreak{}AbelianVarietyTheory.\allowbreak{}Proofs.\allowbreak{}InertiaEigenvalues.\allowbreak{}LenUnitLeOfFamily}, lines 107--109)\par}

{\footnotesize\raggedright\noindent \ref{B8.42bt}.\quad \texttt{GlobalTheory.\allowbreak{}NLinkTyp.\allowbreak{}etaLoc\_\allowbreak{}le\_\allowbreak{}rhs\_\allowbreak{}alt} (\texttt{Development.\allowbreak{}AutomorphicArithmetic.\allowbreak{}ThetaCorrespondence.\allowbreak{}ThetaPairLifts.\allowbreak{}ThirdPlaceSeries.\allowbreak{}EtaLocLeRhs}, lines 21--57)\par}

{\footnotesize\raggedright\noindent \ref{B8.42bu}.\quad \texttt{Formalism.\allowbreak{}MainNotions.\allowbreak{}RamifiedCoefficientChains.\allowbreak{}MassData} (\texttt{Formalism.\allowbreak{}SharpRestriction.\allowbreak{}LevelLetterCocycles.\allowbreak{}ExistsPosConstLagrangeCondition}, lines 269--283)\\
\texttt{Formalism.\allowbreak{}MainNotions.\allowbreak{}RamifiedCoefficientChains.\allowbreak{}SetCS} (\texttt{Formalism.\allowbreak{}SharpRestriction.\allowbreak{}LevelLetterCocycles.\allowbreak{}RamifiedCoefficientChains}, lines 60--72)\\
\texttt{Formalism.\allowbreak{}MainNotions.\allowbreak{}RamifiedCoefficientChains.\allowbreak{}eBd\_\allowbreak{}parts} (\texttt{Formalism.\allowbreak{}SharpRestriction.\allowbreak{}LevelLetterCocycles.\allowbreak{}ExistsPosConstLagrangeCondition}, lines 153--167)\\
\texttt{GlobalTheory.\allowbreak{}ramifiedCoefficientChains\_\allowbreak{}prv\_\allowbreak{}twenty\_\allowbreak{}five} (\texttt{Development.\allowbreak{}AbelianVarietyTheory.\allowbreak{}Proofs.\allowbreak{}InertiaEigenvalues.\allowbreak{}RamifiedCoefficientChainsPrvTwentyFive}, lines 17--29)\par}

{\footnotesize\raggedright\noindent \ref{B8.42bua}.\quad \texttt{Formalism.\allowbreak{}MainNotions.\allowbreak{}RamifiedCoefficientChains.\allowbreak{}S9.\allowbreak{}AdB.\allowbreak{}refine} (\texttt{Formalism.\allowbreak{}SharpRestriction.\allowbreak{}LevelLetterCocycles.\allowbreak{}RingClausesAtUniformizer}, lines 29--117)\par}

{\footnotesize\raggedright\noindent \ref{B8.42buc}.\quad \texttt{Formalism.\allowbreak{}MainNotions.\allowbreak{}CrystallineLiftingInputs.\allowbreak{}DL} (\texttt{Formalism.\allowbreak{}SharpRestriction.\allowbreak{}LevelLetterCocycles.\allowbreak{}CrystallineLiftingInputs}, lines 29--31)\\
\texttt{Formalism.\allowbreak{}MainNotions.\allowbreak{}CrystallineLiftingInputs.\allowbreak{}InPN} (\texttt{Formalism.\allowbreak{}SharpRestriction.\allowbreak{}LevelLetterCocycles.\allowbreak{}CrystallineLiftingInputs}, lines 32--34)\\
\texttt{Formalism.\allowbreak{}MainNotions.\allowbreak{}CrystallineLiftingInputs.\allowbreak{}PolesFin} (\texttt{Formalism.\allowbreak{}SharpRestriction.\allowbreak{}LevelLetterCocycles.\allowbreak{}CrystallineLiftingInputs}, lines 35--37)\\
\texttt{Formalism.\allowbreak{}MainNotions.\allowbreak{}CrystallineLiftingInputs.\allowbreak{}InI} (\texttt{Formalism.\allowbreak{}SharpRestriction.\allowbreak{}LevelLetterCocycles.\allowbreak{}CrystallineLiftingInputs}, lines 119--121)\\
\texttt{Formalism.\allowbreak{}MainNotions.\allowbreak{}CrystallineLiftingInputs.\allowbreak{}Dlat} (\texttt{Formalism.\allowbreak{}SharpRestriction.\allowbreak{}LevelLetterCocycles.\allowbreak{}CrystallineLiftingInputs}, lines 38--40)\\
\texttt{Formalism.\allowbreak{}MainNotions.\allowbreak{}CrystallineLiftingInputs.\allowbreak{}DlatImg} (\texttt{Formalism.\allowbreak{}SharpRestriction.\allowbreak{}LevelLetterCocycles.\allowbreak{}CrystallineLiftingInputs}, lines 41--44)\\
\texttt{Formalism.\allowbreak{}MainNotions.\allowbreak{}CrystallineLiftingInputs.\allowbreak{}pullJ} (\texttt{Formalism.\allowbreak{}SharpRestriction.\allowbreak{}LevelLetterCocycles.\allowbreak{}CrystallineLiftingInputs}, lines 143--146)\\
\texttt{Formalism.\allowbreak{}MainNotions.\allowbreak{}LiftingTheory.\allowbreak{}qp} (\texttt{Formalism.\allowbreak{}SharpRestriction.\allowbreak{}LevelLetterCocycles.\allowbreak{}LiftingTheory}, lines 328--329)\\
\texttt{GlobalTheory.\allowbreak{}AdmHerrLocal.\allowbreak{}main} (\texttt{Development.\allowbreak{}AbelianVarietyTheory.\allowbreak{}Proofs.\allowbreak{}InertiaEigenvalues.\allowbreak{}ImportsBisRest}, lines 9--14)\\
\texttt{Formalism.\allowbreak{}MainNotions.\allowbreak{}RamifiedCoefficientChains.\allowbreak{}S9.\allowbreak{}PStep} (\texttt{Formalism.\allowbreak{}SharpRestriction.\allowbreak{}LevelLetterCocycles.\allowbreak{}LatticeBoundsOfThreeParts}, lines 15--20)\\
\texttt{Formalism.\allowbreak{}MainNotions.\allowbreak{}RamifiedCoefficientChains.\allowbreak{}S9.\allowbreak{}PChain} (\texttt{Formalism.\allowbreak{}SharpRestriction.\allowbreak{}LevelLetterCocycles.\allowbreak{}LatticeBoundsOfThreeParts}, lines 26--30)\\
\texttt{Formalism.\allowbreak{}MainNotions.\allowbreak{}RamifiedCoefficientChains.\allowbreak{}S9.\allowbreak{}cocyclesToWachVectorsExponent} (\texttt{Formalism.\allowbreak{}SharpRestriction.\allowbreak{}LevelLetterCocycles.\allowbreak{}LatticeBoundsOfThreeParts}, lines 172--175)\\
\texttt{Formalism.\allowbreak{}MainNotions.\allowbreak{}RamifiedCoefficientChains.\allowbreak{}S9.\allowbreak{}EG} (\texttt{Formalism.\allowbreak{}SharpRestriction.\allowbreak{}LevelLetterCocycles.\allowbreak{}LatticeBoundsOfThreeParts}, lines 199--201)\\
\texttt{Formalism.\allowbreak{}MainNotions.\allowbreak{}RamifiedCoefficientChains.\allowbreak{}S9.\allowbreak{}wachVectorLiftingExponent} (\texttt{Formalism.\allowbreak{}SharpRestriction.\allowbreak{}LevelLetterCocycles.\allowbreak{}LatticeBoundsOfThreeParts}, lines 202--205)\\
\texttt{Formalism.\allowbreak{}MainNotions.\allowbreak{}RamifiedCoefficientChains.\allowbreak{}S9.\allowbreak{}B6} (\texttt{Formalism.\allowbreak{}SharpRestriction.\allowbreak{}LevelLetterCocycles.\allowbreak{}LatticeBoundsOfThreeParts}, lines 206--207)\par}

{\footnotesize\raggedright\noindent \ref{B8.42bud}.\quad \texttt{Formalism.\allowbreak{}MainNotions.\allowbreak{}RamifiedCoefficientChains.\allowbreak{}S9.\allowbreak{}CocyclesToWachVectorsBound} (\texttt{Formalism.\allowbreak{}SharpRestriction.\allowbreak{}LevelLetterCocycles.\allowbreak{}LatticeBoundsOfThreeParts}, lines 250--265)\\
\texttt{Formalism.\allowbreak{}MainNotions.\allowbreak{}RamifiedCoefficientChains.\allowbreak{}S9.\allowbreak{}ExtensionLemmaAlt.\allowbreak{}cocyclesToWachVectorsBound\_\allowbreak{}of\_\allowbreak{}extTool} (\texttt{Formalism.\allowbreak{}SharpRestriction.\allowbreak{}LevelLetterCocycles.\allowbreak{}CocyclesToWachVectorsBoundOfExtTool}, lines 75--94)\\
\texttt{Formalism.\allowbreak{}MainNotions.\allowbreak{}RamifiedCoefficientChains.\allowbreak{}S9.\allowbreak{}LocalLemmasFromLiterature.\allowbreak{}local\_\allowbreak{}lemmas\_\allowbreak{}of\_\allowbreak{}literature} (\texttt{Formalism.\allowbreak{}SharpRestriction.\allowbreak{}LevelLetterCocycles.\allowbreak{}UniformizerClauseOfAdmittedAlt}, lines 53--114)\par}

{\footnotesize\raggedright\noindent \ref{B8.42buea}.\quad \texttt{Formalism.\allowbreak{}MainNotions.\allowbreak{}RamifiedCoefficientChains.\allowbreak{}S9.\allowbreak{}lemGA\_\allowbreak{}la\_\allowbreak{}of\_\allowbreak{}adm} (\texttt{Formalism.\allowbreak{}SharpRestriction.\allowbreak{}LevelLetterCocycles.\allowbreak{}LiftRingOfUniformizerClause}, lines 31--38)\\
\texttt{Formalism.\allowbreak{}MainNotions.\allowbreak{}RamifiedCoefficientChains.\allowbreak{}S9.\allowbreak{}lemGA\_\allowbreak{}x\_\allowbreak{}of} (\texttt{Formalism.\allowbreak{}SharpRestriction.\allowbreak{}LevelLetterCocycles.\allowbreak{}QPowSmulFrobeniusMem}, lines 70--242)\\
\texttt{Formalism.\allowbreak{}MainNotions.\allowbreak{}RamifiedCoefficientChains.\allowbreak{}S9.\allowbreak{}lemGA0\_\allowbreak{}of\_\allowbreak{}tools} (\texttt{Formalism.\allowbreak{}SharpRestriction.\allowbreak{}LevelLetterCocycles.\allowbreak{}LemGA0ConvVanish}, lines 106--177)\\
\texttt{Formalism.\allowbreak{}MainNotions.\allowbreak{}RamifiedCoefficientChains.\allowbreak{}S9.\allowbreak{}lemGB\_\allowbreak{}of} (\texttt{Formalism.\allowbreak{}SharpRestriction.\allowbreak{}LevelLetterCocycles.\allowbreak{}DerivativeProdSubEval}, lines 58--240)\par}

{\footnotesize\raggedright\noindent \ref{B8.42bue}.\quad \texttt{Formalism.\allowbreak{}MainNotions.\allowbreak{}RamifiedCoefficientChains.\allowbreak{}S9.\allowbreak{}WachVectorLiftingBound} (\texttt{Formalism.\allowbreak{}SharpRestriction.\allowbreak{}LevelLetterCocycles.\allowbreak{}LatticeBoundsOfThreeParts}, lines 386--399)\\
\texttt{Formalism.\allowbreak{}MainNotions.\allowbreak{}RamifiedCoefficientChains.\allowbreak{}Asm.\allowbreak{}frobenius\_\allowbreak{}lift\_\allowbreak{}lemmas\_\allowbreak{}of\_\allowbreak{}admitted} (\texttt{Formalism.\allowbreak{}SharpRestriction.\allowbreak{}LevelLetterCocycles.\allowbreak{}FrobeniusLiftLemmasOfAdmitted}, lines 251--269)\\
\texttt{Formalism.\allowbreak{}MainNotions.\allowbreak{}N.\allowbreak{}Pv.\allowbreak{}LocalLengthBound.\allowbreak{}length\_\allowbreak{}bound\_\allowbreak{}of\_\allowbreak{}hypotheses} (\texttt{Formalism.\allowbreak{}SharpRestriction.\allowbreak{}LevelLetterCocycles.\allowbreak{}QuotientImageDivisibility.\allowbreak{}LengthBoundOfHypotheses}, lines 101--314)\par}

{\footnotesize\raggedright\noindent \ref{B8.42bufa}.\quad \texttt{Formalism.\allowbreak{}MainNotions.\allowbreak{}N.\allowbreak{}Pv.\allowbreak{}LocalLengthBound.\allowbreak{}pchain\_\allowbreak{}DN} (\texttt{Formalism.\allowbreak{}SharpRestriction.\allowbreak{}LevelLetterCocycles.\allowbreak{}QuotientImageDivisibility.\allowbreak{}PchainDNAlt}, lines 69--142)\par}

{\footnotesize\raggedright\noindent \ref{B8.42bufb}.\quad \texttt{Formalism.\allowbreak{}MainNotions.\allowbreak{}N.\allowbreak{}Pv.\allowbreak{}LocalLengthBound.\allowbreak{}exists\_\allowbreak{}inI\_\allowbreak{}inPN\_\allowbreak{}pow\_\allowbreak{}smul\_\allowbreak{}sub\_\allowbreak{}map} (\texttt{Formalism.\allowbreak{}SharpRestriction.\allowbreak{}LevelLetterCocycles.\allowbreak{}QuotientImageDivisibility.\allowbreak{}ExistsInIInPNPowSmulSubMap}, lines 44--109)\par}

{\footnotesize\raggedright\noindent \ref{B8.42buf}.\quad \texttt{Formalism.\allowbreak{}MainNotions.\allowbreak{}RamifiedCoefficientChains.\allowbreak{}S9.\allowbreak{}LatticeBoundSix} (\texttt{Formalism.\allowbreak{}SharpRestriction.\allowbreak{}LevelLetterCocycles.\allowbreak{}LatticeBoundsOfThreeParts}, lines 400--413)\\
\texttt{Formalism.\allowbreak{}MainNotions.\allowbreak{}N.\allowbreak{}Pv.\allowbreak{}LocalLengthBound.\allowbreak{}frobenius\_\allowbreak{}sq\_\allowbreak{}bound\_\allowbreak{}of\_\allowbreak{}local} (\texttt{Formalism.\allowbreak{}SharpRestriction.\allowbreak{}LevelLetterCocycles.\allowbreak{}QuotientImageDivisibility.\allowbreak{}FrobeniusSqBoundOfLocal}, lines 55--166)\\
\texttt{Formalism.\allowbreak{}MainNotions.\allowbreak{}N.\allowbreak{}Pv.\allowbreak{}LocalLengthBound.\allowbreak{}rhoLift\_\allowbreak{}of} (\texttt{Formalism.\allowbreak{}SharpRestriction.\allowbreak{}LevelLetterCocycles.\allowbreak{}QuotientImageDivisibility.\allowbreak{}RhoLiftOf}, lines 20--150)\\
\texttt{Formalism.\allowbreak{}MainNotions.\allowbreak{}N.\allowbreak{}Pv.\allowbreak{}LocalLengthBound.\allowbreak{}rhoLift\_\allowbreak{}D} (\texttt{Formalism.\allowbreak{}SharpRestriction.\allowbreak{}LevelLetterCocycles.\allowbreak{}QuotientImageDivisibility.\allowbreak{}RhoLiftOf}, lines 151--166)\par}

{\footnotesize\raggedright\noindent \ref{B8.42bu1a}.\quad \texttt{Formalism.\allowbreak{}MainNotions.\allowbreak{}UnusedSlot80.\allowbreak{}UnivF.\allowbreak{}Sat} (\texttt{Formalism.\allowbreak{}SharpRestriction.\allowbreak{}LevelLetterCocycles.\allowbreak{}UniversalGeneralOfRestriction}, lines 17--19)\\
\texttt{Formalism.\allowbreak{}MainNotions.\allowbreak{}UnusedSlot80.\allowbreak{}UnivF.\allowbreak{}adapted} (\texttt{Formalism.\allowbreak{}SharpRestriction.\allowbreak{}LevelLetterCocycles.\allowbreak{}UniversalGeneralOfRestriction}, lines 20--47)\par}

{\footnotesize\raggedright\noindent \ref{B8.42bu1b}.\quad \texttt{Formalism.\allowbreak{}MainNotions.\allowbreak{}RamifiedCoefficientChains.\allowbreak{}S9.\allowbreak{}AdB.\allowbreak{}refine} (\texttt{Formalism.\allowbreak{}SharpRestriction.\allowbreak{}LevelLetterCocycles.\allowbreak{}RingClausesAtUniformizer}, lines 29--117)\\
\texttt{Formalism.\allowbreak{}MainNotions.\allowbreak{}RamifiedCoefficientChains.\allowbreak{}S9.\allowbreak{}AdB.\allowbreak{}canon} (\texttt{Formalism.\allowbreak{}SharpRestriction.\allowbreak{}LevelLetterCocycles.\allowbreak{}RingClausesAtUniformizer}, lines 19--28)\par}

{\footnotesize\raggedright\noindent \ref{B8.42bu1c}.\quad \texttt{Formalism.\allowbreak{}MainNotions.\allowbreak{}RamifiedCoefficientChains.\allowbreak{}S9.\allowbreak{}AdM.\allowbreak{}adaptBas\_\allowbreak{}of} (\texttt{Formalism.\allowbreak{}SharpRestriction.\allowbreak{}LevelLetterCocycles.\allowbreak{}RingClausesAtUniformizer}, lines 393--602)\\
\texttt{Formalism.\allowbreak{}MainNotions.\allowbreak{}RamifiedCoefficientChains.\allowbreak{}S9.\allowbreak{}ShB.\allowbreak{}AdaptData} (\texttt{Formalism.\allowbreak{}SharpRestriction.\allowbreak{}LevelLetterCocycles.\allowbreak{}DualLemmaOfDeterminant}, lines 48--57)\par}

{\footnotesize\raggedright\noindent \ref{B8.42bu1}.\quad \texttt{GlobalTheory.\allowbreak{}ramifiedCoefficientChains\_\allowbreak{}prv\_\allowbreak{}twenty\_\allowbreak{}five} (\texttt{Development.\allowbreak{}AbelianVarietyTheory.\allowbreak{}Proofs.\allowbreak{}InertiaEigenvalues.\allowbreak{}RamifiedCoefficientChainsPrvTwentyFive}, lines 17--29)\\
\texttt{GlobalTheory.\allowbreak{}T.\allowbreak{}ramified\_\allowbreak{}coefficient\_\allowbreak{}chains\_\allowbreak{}of\_\allowbreak{}hypotheses} (\texttt{Development.\allowbreak{}AbelianVarietyTheory.\allowbreak{}Proofs.\allowbreak{}InertiaEigenvalues.\allowbreak{}TheoremCConst25OfHypothesesStatement}, lines 12--14)\\
\texttt{Formalism.\allowbreak{}MainNotions.\allowbreak{}RamifiedCoefficientChains.\allowbreak{}Asm.\allowbreak{}theorem\_\allowbreak{}of\_\allowbreak{}literature\_\allowbreak{}hypotheses} (\texttt{Formalism.\allowbreak{}SharpRestriction.\allowbreak{}LevelLetterCocycles.\allowbreak{}FrobeniusLiftLemmasOfAdmitted}, lines 292--298)\par}

{\footnotesize\raggedright\noindent \ref{B8.42bv}.\quad \texttt{GlobalTheory.\allowbreak{}T.\allowbreak{}theoremE\_\allowbreak{}corollaryA\_\allowbreak{}closing\_\allowbreak{}all\_\allowbreak{}data} (\texttt{Development.\allowbreak{}AbelianVarietyTheory.\allowbreak{}Proofs.\allowbreak{}InertiaEigenvalues.\allowbreak{}TheoremECorollaryAClosingAllDataStatement}, lines 8--10)\\
\texttt{GlobalTheory.\allowbreak{}NLink.\allowbreak{}localBoundTableFromData\_\allowbreak{}rowB\_\allowbreak{}rowEntries} (\texttt{Development.\allowbreak{}AbelianVarietyTheory.\allowbreak{}Proofs.\allowbreak{}InertiaEigenvalues.\allowbreak{}LocalBoundTableFromDataRowBRowEntries}, lines 90--114)\par}

{\footnotesize\raggedright\noindent \ref{B8.42bw}.\quad \texttt{GlobalTheory.\allowbreak{}NLink.\allowbreak{}rowC} (\texttt{Development.\allowbreak{}AbelianVarietyTheory.\allowbreak{}Definitions.\allowbreak{}NonGaloisDescent.\allowbreak{}BoundAtTwoAndThreeForFamily}, lines 29--30)\\
\texttt{GlobalTheory.\allowbreak{}T.\allowbreak{}star\_\allowbreak{}ad\_\allowbreak{}and\_\allowbreak{}bb\_\allowbreak{}le\_\allowbreak{}of\_\allowbreak{}joint\_\allowbreak{}cond} (\texttt{Development.\allowbreak{}AutomorphicArithmetic.\allowbreak{}ThetaCorrespondence.\allowbreak{}ThetaPairLifts.\allowbreak{}ThirdPlaceSeries.\allowbreak{}RowsADInSettingStatement}, lines 11--12)\\
\texttt{GlobalTheory.\allowbreak{}T.\allowbreak{}mediumInert\_\allowbreak{}smallInert\_\allowbreak{}lengthBound\_\allowbreak{}tabX} (\texttt{Development.\allowbreak{}AbelianVarietyTheory.\allowbreak{}Proofs.\allowbreak{}InertiaEigenvalues.\allowbreak{}LocalBoundTableRowsCDAdmissibleStatement}, lines 7--9)\par}

{\footnotesize\raggedright\noindent \ref{B8.42bxa}.\quad \texttt{GlobalTheory.\allowbreak{}NLink.\allowbreak{}BBSize.\allowbreak{}blockBoundDeclaration\_\allowbreak{}le\_\allowbreak{}middleRange} (\texttt{Development.\allowbreak{}AbelianVarietyTheory.\allowbreak{}Proofs.\allowbreak{}InertiaEigenvalues.\allowbreak{}NonPlacePartLocalBoundLeMiddleRange}, lines 100--118)\\
\texttt{GlobalTheory.\allowbreak{}NLink.\allowbreak{}BBSize.\allowbreak{}blockBoundAtZero\_\allowbreak{}le\_\allowbreak{}middleRange} (\texttt{Development.\allowbreak{}AbelianVarietyTheory.\allowbreak{}Proofs.\allowbreak{}InertiaEigenvalues.\allowbreak{}NonPlacePartLocalBoundLeMiddleRange}, lines 119--150)\\
\texttt{GlobalTheory.\allowbreak{}NLink.\allowbreak{}BBSize.\allowbreak{}main} (\texttt{Development.\allowbreak{}AbelianVarietyTheory.\allowbreak{}Proofs.\allowbreak{}InertiaEigenvalues.\allowbreak{}NonPlacePartLocalBoundLeMiddleRange}, lines 151--174)\par}

{\footnotesize\raggedright\noindent \ref{B8.42bx}.\quad \texttt{GlobalTheory.\allowbreak{}HTW.\allowbreak{}RecS.\allowbreak{}fam\_\allowbreak{}jointBCD\_\allowbreak{}closed\_\allowbreak{}alt} (\texttt{Development.\allowbreak{}AutomorphicArithmetic.\allowbreak{}ThetaCorrespondence.\allowbreak{}ThetaPairLifts.\allowbreak{}ThirdPlaceSeries.\allowbreak{}FamJointBCDClosed}, lines 13--19)\\
\texttt{GlobalTheory.\allowbreak{}T.\allowbreak{}star\_\allowbreak{}ad\_\allowbreak{}and\_\allowbreak{}bb\_\allowbreak{}le\_\allowbreak{}of\_\allowbreak{}joint\_\allowbreak{}cond} (\texttt{Development.\allowbreak{}AutomorphicArithmetic.\allowbreak{}ThetaCorrespondence.\allowbreak{}ThetaPairLifts.\allowbreak{}ThirdPlaceSeries.\allowbreak{}RowsADInSettingStatement}, lines 11--12)\\
\texttt{GlobalTheory.\allowbreak{}NLink.\allowbreak{}RowA.\allowbreak{}main} (\texttt{Development.\allowbreak{}AbelianVarietyTheory.\allowbreak{}Proofs.\allowbreak{}InertiaEigenvalues.\allowbreak{}TwoEALe}, lines 49--79)\par}

{\footnotesize\raggedright\noindent \ref{B8.42b}.\quad \texttt{GlobalTheory.\allowbreak{}NLinkClose.\allowbreak{}main\_\allowbreak{}alt\_\allowbreak{}alt} (\texttt{Development.\allowbreak{}AutomorphicArithmetic.\allowbreak{}ThetaCorrespondence.\allowbreak{}ThetaPairLifts.\allowbreak{}ThirdPlaceSeries.\allowbreak{}ImportsFurtherTail}, lines 22--109)\\
\texttt{Formalism.\allowbreak{}MainNotions.\allowbreak{}TwistedPlaces.\allowbreak{}etaSum\_\allowbreak{}le} (\texttt{Formalism.\allowbreak{}SharpRestriction.\allowbreak{}LevelLetterCocycles.\allowbreak{}MemPrimesOverOfMem}, lines 558--647)\\
\texttt{Formalism.\allowbreak{}MainNotions.\allowbreak{}TwistedPlaces.\allowbreak{}point\_\allowbreak{}bound} (\texttt{Formalism.\allowbreak{}SharpRestriction.\allowbreak{}LevelLetterCocycles.\allowbreak{}ExportBoundOfAdmissible}, lines 893--996)\par}

{\footnotesize\raggedright\noindent \ref{B8.43i}.\quad \texttt{GlobalTheory.\allowbreak{}ToricTransport.\allowbreak{}toricRatioConjugates\_\allowbreak{}of\_\allowbreak{}transport} (\texttt{Development.\allowbreak{}AbelianVarietyTheory.\allowbreak{}Proofs.\allowbreak{}DivisorProofs.\allowbreak{}ToricRatioConjugatesOfTransport}, lines 57--87)\\
\texttt{GlobalTheory.\allowbreak{}ToricRatioBridge.\allowbreak{}ratioConj\_\allowbreak{}of\_\allowbreak{}toricRatioConjugates} (\texttt{Development.\allowbreak{}AbelianVarietyTheory.\allowbreak{}Proofs.\allowbreak{}DivisorProofs.\allowbreak{}RatioConjOfToricRatioConjugates}, lines 84--137)\par}

{\footnotesize\raggedright\noindent \ref{B8.43j}.\quad \texttt{GlobalTheory.\allowbreak{}LevelFourier.\allowbreak{}gens} (\texttt{Development.\allowbreak{}AbelianVarietyTheory.\allowbreak{}Proofs.\allowbreak{}InertiaEigenvalues.\allowbreak{}VacuumPairUniformLevelBound}, lines 44--110)\\
\texttt{GlobalTheory.\allowbreak{}LevelFourier.\allowbreak{}smul\_\allowbreak{}mem\_\allowbreak{}trSp} (\texttt{Development.\allowbreak{}AbelianVarietyTheory.\allowbreak{}Proofs.\allowbreak{}InertiaEigenvalues.\allowbreak{}VacuumPairUniformLevelBound}, lines 131--301)\par}

{\footnotesize\raggedright\noindent \ref{B8.43k}.\quad \texttt{GlobalTheory.\allowbreak{}LevelSeries.\allowbreak{}exists\_\allowbreak{}denom} (\texttt{Development.\allowbreak{}AbelianVarietyTheory.\allowbreak{}Proofs.\allowbreak{}InertiaEigenvalues.\allowbreak{}SerIntAtOfSeriesRat}, lines 90--134)\\
\texttt{GlobalTheory.\allowbreak{}ConstBound.\allowbreak{}const\_\allowbreak{}dvd} (\texttt{Development.\allowbreak{}AbelianVarietyTheory.\allowbreak{}Proofs.\allowbreak{}InertiaEigenvalues.\allowbreak{}IsTorsionFreeIntOf}, lines 99--150)\par}

{\footnotesize\raggedright\noindent \ref{B8.43l}.\quad \texttt{GlobalTheory.\allowbreak{}TowerPlace.\allowbreak{}placeIn\_\allowbreak{}unique} (\texttt{Development.\allowbreak{}AbelianVarietyTheory.\allowbreak{}Proofs.\allowbreak{}InertiaEigenvalues.\allowbreak{}PlaceInUnique}, lines 9--118)\\
\texttt{GlobalTheory.\allowbreak{}InvDiffLocal.\allowbreak{}valuation\_\allowbreak{}mul\_\allowbreak{}different\_\allowbreak{}le\_\allowbreak{}one} (\texttt{Development.\allowbreak{}AbelianVarietyTheory.\allowbreak{}Proofs.\allowbreak{}InertiaEigenvalues.\allowbreak{}ValuationMulLeOneOfTraceApproxIdempotentIntegral}, lines 117--184)\\
\texttt{GlobalTheory.\allowbreak{}TowerDifferent.\allowbreak{}relDifferent\_\allowbreak{}order\_\allowbreak{}bounded} (\texttt{Development.\allowbreak{}AbelianVarietyTheory.\allowbreak{}Proofs.\allowbreak{}InertiaEigenvalues.\allowbreak{}RelativeDifferentOrderBoundedAbsoluteRamificatio}, lines 56--61)\par}

{\footnotesize\raggedright\noindent \ref{B8.43m}.\quad \texttt{GlobalTheory.\allowbreak{}GoodThetaBound.\allowbreak{}theta\_\allowbreak{}value\_\allowbreak{}bounded} (\texttt{Development.\allowbreak{}AbelianVarietyTheory.\allowbreak{}Proofs.\allowbreak{}InertiaEigenvalues.\allowbreak{}ShortThetaValueBoundedGoodReduction}, lines 191--231)\\
\texttt{GlobalTheory.\allowbreak{}GaugeBoundedDenom.\allowbreak{}exists\_\allowbreak{}canonGauge\_\allowbreak{}nabVal\_\allowbreak{}bounded\_\allowbreak{}denom} (\texttt{Development.\allowbreak{}AbelianVarietyTheory.\allowbreak{}Proofs.\allowbreak{}InertiaEigenvalues.\allowbreak{}ExistsCanonGaugeNabValBoundedDenom}, lines 239--255)\par}

{\footnotesize\raggedright\noindent \ref{B8.43n}.\quad \texttt{GlobalTheory.\allowbreak{}GoodTransfer.\allowbreak{}bddAll\_\allowbreak{}of\_\allowbreak{}bddGood} (\texttt{Development.\allowbreak{}AbelianVarietyTheory.\allowbreak{}Proofs.\allowbreak{}InertiaEigenvalues.\allowbreak{}CuspFormBoundGoodPointsExtendsAllPoints}, lines 497--502)\par}

{\footnotesize\raggedright\noindent \ref{B8.43o}.\quad \texttt{GlobalTheory.\allowbreak{}LevelSeries.\allowbreak{}series\_\allowbreak{}bound} (\texttt{Development.\allowbreak{}AbelianVarietyTheory.\allowbreak{}Proofs.\allowbreak{}InertiaEigenvalues.\allowbreak{}SerIntAtOfSeriesRat}, lines 137--239)\\
\texttt{GlobalTheory.\allowbreak{}LevelSeries.\allowbreak{}serIntAt\_\allowbreak{}record} (\texttt{Development.\allowbreak{}AbelianVarietyTheory.\allowbreak{}Proofs.\allowbreak{}InertiaEigenvalues.\allowbreak{}SerIntAtOfSeriesRat}, lines 274--375)\par}

{\footnotesize\raggedright\noindent \ref{B8.43p}.\quad \texttt{GlobalTheory.\allowbreak{}DescIntRoad.\allowbreak{}descIntAt\_\allowbreak{}of\_\allowbreak{}laws} (\texttt{Development.\allowbreak{}AbelianVarietyTheory.\allowbreak{}Proofs.\allowbreak{}InertiaEigenvalues.\allowbreak{}UnramifiedDescentIntegralFormsTowerCover}, lines 21--31)\\
\texttt{GlobalTheory.\allowbreak{}XiIntSteps.\allowbreak{}intForm\_\allowbreak{}descent\_\allowbreak{}below} (\texttt{Development.\allowbreak{}AbelianVarietyTheory.\allowbreak{}Proofs.\allowbreak{}InertiaEigenvalues.\allowbreak{}IntFormDescentBelow}, lines 95--175)\par}

{\footnotesize\raggedright\noindent \ref{B8.43q}.\quad \texttt{GlobalTheory.\allowbreak{}ThetaTransport.\allowbreak{}existsUnique\_\allowbreak{}trace} (\texttt{Development.\allowbreak{}AbelianVarietyTheory.\allowbreak{}Proofs.\allowbreak{}TwistedOperators.\allowbreak{}HeckeIsogenyThetaGeneratorTransport}, lines 668--704)\\
\texttt{GlobalTheory.\allowbreak{}HeckePointLaw.\allowbreak{}exists\_\allowbreak{}hkRelFr} (\texttt{Development.\allowbreak{}AbelianVarietyTheory.\allowbreak{}Proofs.\allowbreak{}TwistedOperators.\allowbreak{}HeckePointLawsOfTransport}, lines 228--245)\par}

{\footnotesize\raggedright\noindent \ref{B8.43r}.\quad \texttt{GlobalTheory.\allowbreak{}HkIntegral.\allowbreak{}heckeFrameInt} (\texttt{Development.\allowbreak{}AbelianVarietyTheory.\allowbreak{}Proofs.\allowbreak{}TwistedOperators.\allowbreak{}IsUnitNatCastUp}, lines 327--374)\\
\texttt{GlobalTheory.\allowbreak{}HkIntegral.\allowbreak{}heckeValueInt} (\texttt{Development.\allowbreak{}AbelianVarietyTheory.\allowbreak{}Proofs.\allowbreak{}TwistedOperators.\allowbreak{}IsUnitNatCastUp}, lines 376--401)\par}

{\footnotesize\raggedright\noindent \ref{B8.43s}.\quad \texttt{GlobalTheory.\allowbreak{}HeckeEx.\allowbreak{}secTransform} (\texttt{Development.\allowbreak{}AbelianVarietyTheory.\allowbreak{}Proofs.\allowbreak{}TwistedOperators.\allowbreak{}HeckeTransformOfPointLaws}, lines 225--254)\\
\texttt{GlobalTheory.\allowbreak{}HeckeEx.\allowbreak{}cuspTransform} (\texttt{Development.\allowbreak{}AbelianVarietyTheory.\allowbreak{}Proofs.\allowbreak{}TwistedOperators.\allowbreak{}HeckeTransformOfPointLaws}, lines 404--417)\par}

{\footnotesize\raggedright\noindent \ref{B8.43t}.\quad \texttt{GlobalTheory.\allowbreak{}HeckeReadProof.\allowbreak{}heckeReadTie} (\texttt{Development.\allowbreak{}AbelianVarietyTheory.\allowbreak{}Proofs.\allowbreak{}TwistedOperators.\allowbreak{}HeckeReadTie}, lines 34--78)\par}

{\footnotesize\raggedright\noindent \ref{B8.43h}.\quad \texttt{GlobalTheory.\allowbreak{}HeightCounting.\allowbreak{}Up.\allowbreak{}PublishedInputsThirdLayerAlt} (\texttt{Development.\allowbreak{}AutomorphicArithmetic.\allowbreak{}ThetaCorrespondence.\allowbreak{}ThetaPairLifts.\allowbreak{}ThirdPlaceSeries.\allowbreak{}PublishedInputsThirdLayer}, lines 14--28)\\
\texttt{GlobalTheory.\allowbreak{}NRows.\allowbreak{}ThetaLocalInputsAlt} (\texttt{Development.\allowbreak{}AutomorphicArithmetic.\allowbreak{}ThetaCorrespondence.\allowbreak{}ThetaPairLifts.\allowbreak{}ThirdPlaceSeries.\allowbreak{}ThetaLocalInputs}, lines 13--28)\\
\texttt{GlobalTheory.\allowbreak{}NRows.\allowbreak{}GeometricGaloisRepCompatibility} (\texttt{Development.\allowbreak{}AutomorphicArithmetic.\allowbreak{}ThetaCorrespondence.\allowbreak{}ThetaPairLifts.\allowbreak{}LocalTypeDictionary.\allowbreak{}GeometricGaloisRepCompatibility}, lines 6--18)\\
\texttt{GlobalTheory.\allowbreak{}HeightCounting.\allowbreak{}PhtOn} (\texttt{Development.\allowbreak{}AbelianVarietyTheory.\allowbreak{}Definitions.\allowbreak{}AutomorphicMembership.\allowbreak{}NonemptyAllIff}, lines 97--101)\\
\texttt{GlobalTheory.\allowbreak{}HeightCounting.\allowbreak{}Up.\allowbreak{}SerN} (\texttt{Development.\allowbreak{}AbelianVarietyTheory.\allowbreak{}Definitions.\allowbreak{}AutomorphicMembership.\allowbreak{}CenUsedOfCenReal}, lines 56--62)\\
\texttt{GlobalTheory.\allowbreak{}MemberClose.\allowbreak{}exists\_\allowbreak{}hyp\_\allowbreak{}n\_\allowbreak{}of\_\allowbreak{}geometric\_\allowbreak{}galois\_\allowbreak{}compat} (\texttt{Development.\allowbreak{}AutomorphicArithmetic.\allowbreak{}ThetaCorrespondence.\allowbreak{}ThetaPairLifts.\allowbreak{}ThirdPlaceSeries.\allowbreak{}ThetaLiftHypothesesOfNonvanishing}, lines 24--46)\par}

{\footnotesize\raggedright\noindent \ref{B8.43}.\quad \texttt{UpperTheory.\allowbreak{}GalC.\allowbreak{}final\_\allowbreak{}alt\_\allowbreak{}alt} (\texttt{Development.\allowbreak{}AutomorphicArithmetic.\allowbreak{}ThetaCorrespondence.\allowbreak{}ThetaPairLifts.\allowbreak{}ThirdPlaceSeries.\allowbreak{}CFinal}, lines 19--22)\\
\texttt{UpperTheory.\allowbreak{}GalC.\allowbreak{}final\_\allowbreak{}of\_\allowbreak{}alt} (\texttt{Development.\allowbreak{}AutomorphicArithmetic.\allowbreak{}ThetaCorrespondence.\allowbreak{}ThetaPairLifts.\allowbreak{}ThirdPlaceSeries.\allowbreak{}LinkNFam}, lines 75--86)\\
\texttt{UpperTheory.\allowbreak{}GalC.\allowbreak{}linkN\_\allowbreak{}of\_\allowbreak{}boundAtTwoThree\_\allowbreak{}slot} (\texttt{Development.\allowbreak{}AutomorphicArithmetic.\allowbreak{}ThetaCorrespondence.\allowbreak{}ThetaPairLifts.\allowbreak{}ThirdPlaceSeries.\allowbreak{}LinkNFam}, lines 54--63)\par}

\section*{Chapter~\ref{ch:p8c5}: Primes in short intervals, masses and counts}

{\footnotesize\raggedright\noindent \ref{B8.44}.\quad \texttt{Formalism.\allowbreak{}MainNotions.\allowbreak{}LitThin} (\texttt{Formalism.\allowbreak{}SharpRestriction.\allowbreak{}GlobalSchurBounds.\allowbreak{}MainSup}, lines 88--90)\\
\texttt{GlobalTheory.\allowbreak{}ThinScaleLiterature.\allowbreak{}litThin} (\texttt{Development.\allowbreak{}AbelianVarietyTheory.\allowbreak{}Proofs.\allowbreak{}GlobalInterpolation.\allowbreak{}LitThin}, lines 8--9)\par}

{\footnotesize\raggedright\noindent \ref{B8.45}.\quad \texttt{GlobalTheory.\allowbreak{}ZerosNearOne.\allowbreak{}density\_\allowbreak{}and\_\allowbreak{}regions} (\texttt{Development.\allowbreak{}AbelianVarietyTheory.\allowbreak{}Proofs.\allowbreak{}GlobalInterpolation.\allowbreak{}DensityAndRegions}, lines 10--16)\\
\texttt{Formalism.\allowbreak{}MainNotions.\allowbreak{}ShortIntervalPrimes.\allowbreak{}adm6} (\texttt{Formalism.\allowbreak{}SharpRestriction.\allowbreak{}GlobalSchurBounds.\allowbreak{}ShortIntervalPrimes}, lines 128--139)\par}

{\footnotesize\raggedright\noindent \ref{B8.46}.\quad \texttt{GlobalTheory.\allowbreak{}ZeroBounds.\allowbreak{}explicit\_\allowbreak{}formula\_\allowbreak{}and\_\allowbreak{}window\_\allowbreak{}count} (\texttt{Development.\allowbreak{}AbelianVarietyTheory.\allowbreak{}Proofs.\allowbreak{}GlobalInterpolation.\allowbreak{}ExplicitFormulaAndWindowCount}, lines 9--13)\\
\texttt{Formalism.\allowbreak{}MainNotions.\allowbreak{}ShortIntervalPrimes.\allowbreak{}adm7} (\texttt{Formalism.\allowbreak{}SharpRestriction.\allowbreak{}GlobalSchurBounds.\allowbreak{}ShortIntervalPrimes}, lines 141--147)\\
\texttt{Formalism.\allowbreak{}MainNotions.\allowbreak{}ShortIntervalPrimes.\allowbreak{}adm11} (\texttt{Formalism.\allowbreak{}SharpRestriction.\allowbreak{}GlobalSchurBounds.\allowbreak{}ShortIntervalPrimes}, lines 158--161)\par}

{\footnotesize\raggedright\noindent \ref{B8.47}.\quad \texttt{Formalism.\allowbreak{}MainNotions.\allowbreak{}chosenThin\_\allowbreak{}of\_\allowbreak{}boxesN} (\texttt{Formalism.\allowbreak{}SharpRestriction.\allowbreak{}GlobalSchurBounds.\allowbreak{}MainSup}, lines 105--114)\\
\texttt{Formalism.\allowbreak{}MainNotions.\allowbreak{}chosenThin\_\allowbreak{}closed} (\texttt{Formalism.\allowbreak{}SharpRestriction.\allowbreak{}GlobalSchurBounds.\allowbreak{}MainSup}, lines 124--127)\\
\texttt{Formalism.\allowbreak{}MainNotions.\allowbreak{}FamCore} (\texttt{Formalism.\allowbreak{}SharpRestriction.\allowbreak{}FamCore}, lines 71--136)\par}

{\footnotesize\raggedright\noindent \ref{B8.48}.\quad \texttt{GlobalTheory.\allowbreak{}Mass.\allowbreak{}massRatio\_\allowbreak{}all\_\allowbreak{}unramified\_\allowbreak{}cmFields} (\texttt{Development.\allowbreak{}AbelianVarietyTheory.\allowbreak{}Proofs.\allowbreak{}GlobalInterpolation.\allowbreak{}MassRatioAllUnramifiedCmFields}, lines 8--9)\\
\texttt{Formalism.\allowbreak{}HermitianLattices.\allowbreak{}SiegelMassBounds.\allowbreak{}siegel\_\allowbreak{}mass\_\allowbreak{}ratio} (\texttt{Formalism.\allowbreak{}GL3MeanSquares.\allowbreak{}HookFormulaTorsion.\allowbreak{}MassUpperBoundClosed}, lines 12--12)\\
\texttt{Formalism.\allowbreak{}MainNotions.\allowbreak{}StandardLattice.\allowbreak{}SiegelMassRatioCM} (\texttt{Formalism.\allowbreak{}SharpRestriction.\allowbreak{}GlobalSchurBounds.\allowbreak{}StandardLattice}, lines 281--284)\par}

{\footnotesize\raggedright\noindent \ref{B8.48}.\quad \texttt{Formalism.\allowbreak{}HermitianLattices.\allowbreak{}Ext.\allowbreak{}L3.\allowbreak{}massConst\_\allowbreak{}succ} (\texttt{Formalism.\allowbreak{}GL3MeanSquares.\allowbreak{}GenusMeanAsymptotics.\allowbreak{}UnimodularClassRepresentatives.\allowbreak{}MassConstSucc}, lines 99--131)\\
\texttt{Formalism.\allowbreak{}HermitianLattices.\allowbreak{}Ext.\allowbreak{}L3.\allowbreak{}siegelMassRatio\_\allowbreak{}of\_\allowbreak{}LI} (\texttt{Formalism.\allowbreak{}GL3MeanSquares.\allowbreak{}GenusMeanAsymptotics.\allowbreak{}UnimodularClassRepresentatives.\allowbreak{}SiegelMassRatioOfLI}, lines 35--55)\\
\texttt{GlobalTheory.\allowbreak{}Mass.\allowbreak{}binary\_\allowbreak{}hermitian\_\allowbreak{}represents} (\texttt{Development.\allowbreak{}AbelianVarietyTheory.\allowbreak{}Proofs.\allowbreak{}FlagMeasures.\allowbreak{}BinaryHermitianRepresents}, lines 12--66)\par}

{\footnotesize\raggedright\noindent \ref{B8.49}.\quad \texttt{GlobalTheory.\allowbreak{}VolRatioFromSiegelWeil.\allowbreak{}volumeRatio} (\texttt{Development.\allowbreak{}AbelianVarietyTheory.\allowbreak{}Proofs.\allowbreak{}FlagMeasures.\allowbreak{}UnitaryHyperplaneVolumeRatioSiegelWeil}, lines 10--14)\\
\texttt{Formalism.\allowbreak{}MainNotions.\allowbreak{}StandardLattice.\allowbreak{}VolumeRatio₂} (\texttt{Formalism.\allowbreak{}SharpRestriction.\allowbreak{}GlobalSchurBounds.\allowbreak{}ThickPlacesCorollaryOfAnyTranslate}, lines 2330--2340)\par}

{\footnotesize\raggedright\noindent \ref{B8.50}.\quad \texttt{GlobalTheory.\allowbreak{}DefiniteMass.\allowbreak{}mass\_\allowbreak{}eq\_\allowbreak{}std} (\texttt{Development.\allowbreak{}AbelianVarietyTheory.\allowbreak{}Proofs.\allowbreak{}GlobalInterpolation.\allowbreak{}UnramifiedOfIsCMUnram}, lines 147--166)\\
\texttt{GlobalTheory.\allowbreak{}genusMassAt\_\allowbreak{}two} (\texttt{Development.\allowbreak{}AbelianVarietyTheory.\allowbreak{}Proofs.\allowbreak{}GlobalInterpolation.\allowbreak{}GenusMassAtTwo}, lines 7--8)\par}

{\footnotesize\raggedright\noindent \ref{B8.50a}.\quad \texttt{PI.\allowbreak{}UnitaryGroupMassFormula.\allowbreak{}HCFdegRGap.\allowbreak{}hCFdegRGap\_\allowbreak{}holds} (\texttt{Formalism.\allowbreak{}PeriodIntegrals.\allowbreak{}UnitaryGroupMassFormula.\allowbreak{}FormalDegreeGap.\allowbreak{}HCFdegRGapHolds}, lines 108--111)\\
\texttt{Formalism.\allowbreak{}MainNotions.\allowbreak{}GappedMultiplicityBounds.\allowbreak{}theorem\_\allowbreak{}three\_\allowbreak{}rows\_\allowbreak{}of\_\allowbreak{}real\_\allowbreak{}place} (\texttt{Formalism.\allowbreak{}SharpRestriction.\allowbreak{}GlobalSchurBounds.\allowbreak{}TheoremThreeFamilyOfGappedMultiplicity}, lines 131--193)\\
\texttt{Formalism.\allowbreak{}MainNotions.\allowbreak{}GappedMultiplicityBounds.\allowbreak{}explicit\_\allowbreak{}bound\_\allowbreak{}family} (\texttt{Formalism.\allowbreak{}SharpRestriction.\allowbreak{}GlobalSchurBounds.\allowbreak{}TheoremThreeFamilyOfGappedMultiplicity}, lines 200--201)\par}

{\footnotesize\raggedright\noindent \ref{B8.51b}.\quad \texttt{Formalism.\allowbreak{}L2.\allowbreak{}GlobalTheoryTheory.\allowbreak{}GenusMassVolumes.\allowbreak{}vol\_\allowbreak{}indef} (\texttt{Formalism.\allowbreak{}LocalToGlobal.\allowbreak{}GlobalTheory.\allowbreak{}GL5.\allowbreak{}DiscreteSeriesOfFormalDegree}, lines 57--84)\\
\texttt{Formalism.\allowbreak{}L2.\allowbreak{}GlobalTheoryTheory.\allowbreak{}GenusMassVolumes.\allowbreak{}volume\_\allowbreak{}eq\_\allowbreak{}mass\_\allowbreak{}rank\_\allowbreak{}two} (\texttt{Formalism.\allowbreak{}LocalToGlobal.\allowbreak{}GlobalTheory.\allowbreak{}GL5.\allowbreak{}DiscreteSeriesOfFormalDegree}, lines 104--110)\par}

{\footnotesize\raggedright\noindent \ref{B8.51}.\quad \texttt{GlobalTheory.\allowbreak{}Fc.\allowbreak{}row\_\allowbreak{}hTu\_\allowbreak{}cTmin\_\allowbreak{}alt} (\texttt{Development.\allowbreak{}AutomorphicArithmetic.\allowbreak{}ThetaCorrespondence.\allowbreak{}ThetaPairLifts.\allowbreak{}ThirdPlaceSeries.\allowbreak{}RowHTuCTmin}, lines 11--21)\\
\texttt{Formalism.\allowbreak{}L2.\allowbreak{}GlobalTheoryTheory.\allowbreak{}GenusMassVolumes.\allowbreak{}volume\_\allowbreak{}eq\_\allowbreak{}mass\_\allowbreak{}rank\_\allowbreak{}two} (\texttt{Formalism.\allowbreak{}LocalToGlobal.\allowbreak{}GlobalTheory.\allowbreak{}GL5.\allowbreak{}DiscreteSeriesOfFormalDegree}, lines 104--110)\par}

{\footnotesize\raggedright\noindent \ref{B8.51a}.\quad \texttt{UpperTheory.\allowbreak{}T3.\allowbreak{}levelPrimeN} (\texttt{Development.\allowbreak{}AutomorphicArithmetic.\allowbreak{}ShortNorms.\allowbreak{}levelPrimeN}, lines 36--161)\\
\texttt{GlobalTheory.\allowbreak{}EX.\allowbreak{}OB.\allowbreak{}exists\_\allowbreak{}levCh\_\allowbreak{}of\_\allowbreak{}levelPrime} (\texttt{Development.\allowbreak{}AbelianVarietyTheory.\allowbreak{}Proofs.\allowbreak{}GlobalInterpolation.\allowbreak{}ExistsLevChOfLevelPrime}, lines 20--112)\par}

{\footnotesize\raggedright\noindent \ref{B8.52a}.\quad \texttt{Formalism.\allowbreak{}L2.\allowbreak{}SG.\allowbreak{}Pf.\allowbreak{}AllRowsLab.\allowbreak{}rowsLab\_\allowbreak{}holds} (\texttt{Formalism.\allowbreak{}LocalToGlobal.\allowbreak{}SiegelLatticeTheory.\allowbreak{}FlagStabilizerTensors.\allowbreak{}RootLatticeThresholds.\allowbreak{}AllRowsLab}, lines 19--24)\\
\texttt{Formalism.\allowbreak{}L2.\allowbreak{}SG.\allowbreak{}Pf.\allowbreak{}AllRowsAddFin.\allowbreak{}eHom\_\allowbreak{}fin} (\texttt{Formalism.\allowbreak{}LocalToGlobal.\allowbreak{}SiegelLatticeTheory.\allowbreak{}FlagStabilizerTensors.\allowbreak{}RootLatticeThresholds.\allowbreak{}AllRowsAddFin}, lines 15--24)\\
\texttt{Formalism.\allowbreak{}L2.\allowbreak{}SG.\allowbreak{}Pf.\allowbreak{}RowsGoodRed.\allowbreak{}rowsGood\_\allowbreak{}of} (\texttt{Formalism.\allowbreak{}LocalToGlobal.\allowbreak{}SiegelLatticeTheory.\allowbreak{}FlagStabilizerTensors.\allowbreak{}RootLatticeThresholds.\allowbreak{}RowsGoodRed}, lines 39--105)\par}

{\footnotesize\raggedright\noindent \ref{B8.52}.\quad \texttt{GlobalTheory.\allowbreak{}CountRowByName.\allowbreak{}row\_\allowbreak{}counts\_\allowbreak{}of\_\allowbreak{}labels\_\allowbreak{}alt} (\texttt{Development.\allowbreak{}AutomorphicArithmetic.\allowbreak{}ThetaCorrespondence.\allowbreak{}ThetaPairLifts.\allowbreak{}ThirdPlaceSeries.\allowbreak{}RowCountsOfLabels}, lines 13--14)\\
\texttt{GlobalTheory.\allowbreak{}CountRow.\allowbreak{}row\_\allowbreak{}counts\_\allowbreak{}of\_\allowbreak{}ties\_\allowbreak{}alt} (\texttt{Development.\allowbreak{}AutomorphicArithmetic.\allowbreak{}ThetaCorrespondence.\allowbreak{}ThetaPairLifts.\allowbreak{}ThirdPlaceSeries.\allowbreak{}RowCountsOfTies}, lines 12--21)\par}

{\footnotesize\raggedright\noindent \ref{B8.53}.\quad \texttt{GlobalTheory.\allowbreak{}CountBridge.\allowbreak{}count\_\allowbreak{}unbundled} (\texttt{Development.\allowbreak{}AbelianVarietyTheory.\allowbreak{}Proofs.\allowbreak{}FlagMeasures.\allowbreak{}UnitaryAutomorphicCountingUnbundled}, lines 40--74)\\
\texttt{GlobalTheory.\allowbreak{}AllRowsAtRecord.\allowbreak{}countAllRows\_\allowbreak{}of} (\texttt{Development.\allowbreak{}AbelianVarietyTheory.\allowbreak{}Proofs.\allowbreak{}FlagMeasures.\allowbreak{}AllRowsAtRecord}, lines 37--65)\\
\texttt{Formalism.\allowbreak{}MainNotions.\allowbreak{}S.\allowbreak{}CountE.\allowbreak{}CountIdempotentsClass} (\texttt{Formalism.\allowbreak{}SharpRestriction.\allowbreak{}CuspSpanNormBounds.\allowbreak{}BridgeCountE}, lines 45--59)\\
\texttt{Formalism.\allowbreak{}L2.\allowbreak{}GlobalTheoryTheory.\allowbreak{}ConcN.\allowbreak{}CountAllRows} (\texttt{Formalism.\allowbreak{}LocalToGlobal.\allowbreak{}GlobalTheory.\allowbreak{}IrreducibleOfCountSizes}, lines 183--186)\par}

{\footnotesize\raggedright\noindent \ref{B8.53a-jet}.\quad \texttt{GlobalTheory.\allowbreak{}Jet.\allowbreak{}Germ.\allowbreak{}jets\_\allowbreak{}vanish} (\texttt{Development.\allowbreak{}AbelianVarietyTheory.\allowbreak{}Proofs.\allowbreak{}FlagMeasures.\allowbreak{}HighestWeightGapBound}, lines 244--284)\\
\texttt{GlobalTheory.\allowbreak{}Jet.\allowbreak{}UniqAlg.\allowbreak{}uniqAlg} (\texttt{Development.\allowbreak{}AbelianVarietyTheory.\allowbreak{}Proofs.\allowbreak{}FlagMeasures.\allowbreak{}JetUniquenessAlgebraic}, lines 249--319)\par}

{\footnotesize\raggedright\noindent \ref{B8.53a}.\quad \texttt{GlobalTheory.\allowbreak{}DecGapBound.\allowbreak{}decGap\_\allowbreak{}of\_\allowbreak{}formulaGap} (\texttt{Development.\allowbreak{}AbelianVarietyTheory.\allowbreak{}Proofs.\allowbreak{}FlagMeasures.\allowbreak{}LowestKTypeDecayGapOfCompressionFormulaGap}, lines 11--32)\\
\texttt{GlobalTheory.\allowbreak{}Jet.\allowbreak{}FormGapFree.\allowbreak{}formulaGap} (\texttt{Development.\allowbreak{}AbelianVarietyTheory.\allowbreak{}Proofs.\allowbreak{}FlagMeasures.\allowbreak{}LowestKTypeCompressionFormulaOfStrictLowestKType}, lines 22--99)\\
\texttt{GlobalTheory.\allowbreak{}AutBall.\allowbreak{}autBallCount} (\texttt{Development.\allowbreak{}AbelianVarietyTheory.\allowbreak{}Proofs.\allowbreak{}FlagMeasures.\allowbreak{}AutBallCount}, lines 28--48)\par}

{\footnotesize\raggedright\noindent \ref{B8.53b}.\quad \texttt{GlobalTheory.\allowbreak{}FlagMeas.\allowbreak{}AbsVol.\allowbreak{}absoluteVolume\_\allowbreak{}of\_\allowbreak{}volumeStep} (\texttt{Development.\allowbreak{}AbelianVarietyTheory.\allowbreak{}Proofs.\allowbreak{}FlagMeasures.\allowbreak{}HermAbsoluteVolumeBoundOfVolumeRatioInduction}, lines 37--69)\\
\texttt{GlobalTheory.\allowbreak{}FlagMeas.\allowbreak{}constantBound} (\texttt{Development.\allowbreak{}AbelianVarietyTheory.\allowbreak{}Proofs.\allowbreak{}FlagMeasures.\allowbreak{}FlagExplicitConstantLeExpNegHalf}, lines 82--228)\\
\texttt{GlobalTheory.\allowbreak{}NonzeroFlag.\allowbreak{}nonzeroFlagBound} (\texttt{Development.\allowbreak{}AbelianVarietyTheory.\allowbreak{}Proofs.\allowbreak{}FlagMeasures.\allowbreak{}NonzeroFlagProbabilityBoundOfVolumeRatioInductio}, lines 84--148)\par}

{\footnotesize\raggedright\noindent \ref{B8.53c}.\quad \texttt{GlobalTheory.\allowbreak{}deep\_\allowbreak{}region\_\allowbreak{}exponents} (\texttt{Development.\allowbreak{}AbelianVarietyTheory.\allowbreak{}Proofs.\allowbreak{}FlagMeasures.\allowbreak{}DeepRegionExponents}, lines 54--204)\par}

{\footnotesize\raggedright\noindent \ref{B8.53d}.\quad \texttt{GlobalTheory.\allowbreak{}DeepFamilyGap.\allowbreak{}deep\_\allowbreak{}bound\_\allowbreak{}family} (\texttt{Development.\allowbreak{}AbelianVarietyTheory.\allowbreak{}Proofs.\allowbreak{}FlagMeasures.\allowbreak{}DeepPointBoundCuspidalIsometricIntertwinerFamily}, lines 118--183)\par}

{\footnotesize\raggedright\noindent \ref{B8.53e}.\quad \texttt{GlobalTheory.\allowbreak{}RowsFinMult.\allowbreak{}rows\_\allowbreak{}finiteDimensional} (\texttt{Development.\allowbreak{}AbelianVarietyTheory.\allowbreak{}Proofs.\allowbreak{}FlagMeasures.\allowbreak{}FiniteDimensionalRowsSpaceDiscreteSeriesCount}, lines 238--344)\\
\texttt{GlobalTheory.\allowbreak{}DeepKernel.\allowbreak{}deep\_\allowbreak{}kernel\_\allowbreak{}bound\_\allowbreak{}of\_\allowbreak{}ptBound} (\texttt{Development.\allowbreak{}AbelianVarietyTheory.\allowbreak{}Proofs.\allowbreak{}FlagMeasures.\allowbreak{}DeepPointKernelBoundOfPointwiseBound}, lines 22--48)\\
\texttt{GlobalTheory.\allowbreak{}DeepRegion.\allowbreak{}deepBound} (\texttt{Development.\allowbreak{}AbelianVarietyTheory.\allowbreak{}Proofs.\allowbreak{}FlagMeasures.\allowbreak{}DeepRegionKernelDiagonalBound}, lines 28--150)\par}

{\footnotesize\raggedright\noindent \ref{B8.54}.\quad \texttt{GlobalTheory.\allowbreak{}Fc.\allowbreak{}R0.\allowbreak{}classCountBound\_\allowbreak{}paramsLambda\_\allowbreak{}constructedPoints} (\texttt{Development.\allowbreak{}AbelianVarietyTheory.\allowbreak{}Proofs.\allowbreak{}GlobalInterpolation.\allowbreak{}ClassCountBoundParamsLambdaConstructedPoints}, lines 17--25)\\
\texttt{Formalism.\allowbreak{}MainNotions.\allowbreak{}StandardLattice.\allowbreak{}ClassNumberFormula₂} (\texttt{Formalism.\allowbreak{}SharpRestriction.\allowbreak{}GlobalSchurBounds.\allowbreak{}StandardLattice}, lines 453--456)\\
\texttt{Formalism.\allowbreak{}MainNotions.\allowbreak{}J.\allowbreak{}level\_\allowbreak{}prime\_\allowbreak{}splits\_\allowbreak{}totally\_\allowbreak{}real} (\texttt{Formalism.\allowbreak{}SharpRestriction.\allowbreak{}AdmissibleParameterJoins.\allowbreak{}DIVfromGLz}, lines 13--18)\par}

{\footnotesize\raggedright\noindent \ref{B8.55a}.\quad \texttt{GlobalTheory.\allowbreak{}ClassSumSpan.\allowbreak{}mem\_\allowbreak{}span\_\allowbreak{}of\_\allowbreak{}map\_\allowbreak{}mem} (\texttt{Development.\allowbreak{}AbelianVarietyTheory.\allowbreak{}Proofs.\allowbreak{}DivisorProofs.\allowbreak{}ClassSumsSpan}, lines 32--54)\par}

{\footnotesize\raggedright\noindent \ref{B8.55b}.\quad \texttt{GlobalTheory.\allowbreak{}ClassSumSpan.\allowbreak{}atDict\_\allowbreak{}indep} (\texttt{Development.\allowbreak{}AbelianVarietyTheory.\allowbreak{}Proofs.\allowbreak{}DivisorProofs.\allowbreak{}ClassSumsSpan}, lines 165--196)\\
\texttt{GlobalTheory.\allowbreak{}reading\_\allowbreak{}linearIndependent\_\allowbreak{}and\_\allowbreak{}injective} (\texttt{Development.\allowbreak{}AbelianVarietyTheory.\allowbreak{}Proofs.\allowbreak{}UnitaryExistence.\allowbreak{}ReadingLinearIndependentAndInjective}, lines 22--55)\par}

{\footnotesize\raggedright\noindent \ref{B8.55}.\quad \texttt{GlobalTheory.\allowbreak{}Published.\allowbreak{}DiscreteSeriesCover} (\texttt{Development.\allowbreak{}AbelianVarietyTheory.\allowbreak{}PublishedTheorems.\allowbreak{}TwistedTraceFormula.\allowbreak{}AdmissibleDiscreteSeriesOfHasParam}, lines 35--39)\\
\texttt{Formalism.\allowbreak{}BasicNotions.\allowbreak{}TH.\allowbreak{}GrpData} (\texttt{Formalism.\allowbreak{}LocalRingCardinality.\allowbreak{}DataTheta}, lines 93--120)\\
\texttt{Formalism.\allowbreak{}BasicNotions.\allowbreak{}TH.\allowbreak{}DSRep} (\texttt{Formalism.\allowbreak{}LocalRingCardinality.\allowbreak{}DataTheta}, lines 165--179)\\
\texttt{Formalism.\allowbreak{}BasicNotions.\allowbreak{}TH.\allowbreak{}IsParam} (\texttt{Formalism.\allowbreak{}LocalRingCardinality.\allowbreak{}DataTheta}, lines 382--385)\\
\texttt{GlobalTheory.\allowbreak{}CoverDS.\allowbreak{}HasTypes} (\texttt{Development.\allowbreak{}AbelianVarietyTheory.\allowbreak{}PublishedTheorems.\allowbreak{}TwistedTraceFormula.\allowbreak{}AdmissibleDiscreteSeriesOfHasParam}, lines 24--28)\\
\texttt{GlobalTheory.\allowbreak{}CoverDS.\allowbreak{}UnitEquiv} (\texttt{Development.\allowbreak{}AbelianVarietyTheory.\allowbreak{}PublishedTheorems.\allowbreak{}TwistedTraceFormula.\allowbreak{}AdmissibleDiscreteSeriesOfHasParam}, lines 16--18)\par}

{\footnotesize\raggedright\noindent \ref{B8.55}.\quad \texttt{GlobalTheory.\allowbreak{}Published.\allowbreak{}HarishChandraModules} (\texttt{Development.\allowbreak{}AbelianVarietyTheory.\allowbreak{}PublishedTheorems.\allowbreak{}TwistedTraceFormula.\allowbreak{}HarishChandraModulesOn}, lines 43--45)\\
\texttt{GlobalTheory.\allowbreak{}Published.\allowbreak{}HarishChandraModulesOn} (\texttt{Development.\allowbreak{}AbelianVarietyTheory.\allowbreak{}PublishedTheorems.\allowbreak{}TwistedTraceFormula.\allowbreak{}HarishChandraModulesOn}, lines 24--36)\\
\texttt{GlobalTheory.\allowbreak{}Published.\allowbreak{}KFiniteAnalyticUJ} (\texttt{Development.\allowbreak{}AbelianVarietyTheory.\allowbreak{}PublishedTheorems.\allowbreak{}TwistedTraceFormula.\allowbreak{}HarishChandraModulesOn}, lines 37--42)\par}

{\footnotesize\raggedright\noindent \ref{B8.55}.\quad \texttt{GlobalTheory.\allowbreak{}Published.\allowbreak{}CuspRapidDecay} (\texttt{Development.\allowbreak{}AbelianVarietyTheory.\allowbreak{}PublishedTheorems.\allowbreak{}AutomorphicForms₂.\allowbreak{}ExistsNormLeHeightPow}, lines 28--33)\\
\texttt{Formalism.\allowbreak{}L2.\allowbreak{}Theta.\allowbreak{}IsAutFun} (\texttt{Formalism.\allowbreak{}LocalToGlobal.\allowbreak{}BesselHermitianTheta.\allowbreak{}IsAutFunDef}, lines 43--47)\\
\texttt{Formalism.\allowbreak{}L2.\allowbreak{}Theta.\allowbreak{}IsCuspFun} (\texttt{Formalism.\allowbreak{}LocalToGlobal.\allowbreak{}BesselHermitianTheta.\allowbreak{}IsAutFunDef}, lines 67--71)\\
\texttt{GlobalTheory.\allowbreak{}IndefiniteRowSideFullHecke.\allowbreak{}LieAlgebraExponentialFamilies.\allowbreak{}IsUnifModerate} (\texttt{Development.\allowbreak{}AbelianVarietyTheory.\allowbreak{}BlockTheory.\allowbreak{}HeightDefect.\allowbreak{}WordDerivatives}, lines 122--125)\\
\texttt{GlobalTheory.\allowbreak{}CuspDecay.\allowbreak{}IsRapidDecay} (\texttt{Development.\allowbreak{}AbelianVarietyTheory.\allowbreak{}PublishedTheorems.\allowbreak{}AutomorphicForms₂.\allowbreak{}ExistsNormLeHeightPow}, lines 21--23)\par}

{\footnotesize\raggedright\noindent \ref{B8.55}.\quad \texttt{GlobalTheory.\allowbreak{}Published.\allowbreak{}CasselmanWallach} (\texttt{Development.\allowbreak{}AbelianVarietyTheory.\allowbreak{}PublishedTheorems.\allowbreak{}AutomorphicForms₂.\allowbreak{}HasDRepIffHasDerivAt}, lines 98--112)\\
\texttt{GlobalTheory.\allowbreak{}SmoothGlob.\allowbreak{}IsSmoothModRep} (\texttt{Development.\allowbreak{}AbelianVarietyTheory.\allowbreak{}PublishedTheorems.\allowbreak{}AutomorphicForms₂.\allowbreak{}HasDRepIffHasDerivAt}, lines 22--31)\\
\texttt{GlobalTheory.\allowbreak{}SmoothGlob.\allowbreak{}sobNorm} (\texttt{Development.\allowbreak{}AbelianVarietyTheory.\allowbreak{}PublishedTheorems.\allowbreak{}AutomorphicForms₂.\allowbreak{}HasDRepIffHasDerivAt}, lines 34--36)\\
\texttt{GlobalTheory.\allowbreak{}SmoothGlob.\allowbreak{}IsGKMap} (\texttt{Development.\allowbreak{}AbelianVarietyTheory.\allowbreak{}PublishedTheorems.\allowbreak{}AutomorphicForms₂.\allowbreak{}HasDRepIffHasDerivAt}, lines 37--41)\\
\texttt{GlobalTheory.\allowbreak{}SmoothGlob.\allowbreak{}IsSmoothGMap} (\texttt{Development.\allowbreak{}AbelianVarietyTheory.\allowbreak{}PublishedTheorems.\allowbreak{}AutomorphicForms₂.\allowbreak{}HasDRepIffHasDerivAt}, lines 42--47)\\
\texttt{Formalism.\allowbreak{}REAL.\allowbreak{}hs} (\texttt{Formalism.\allowbreak{}RealPlaceAnalysis.\allowbreak{}Complexify}, lines 27--28)\par}

{\footnotesize\raggedright\noindent \ref{B8.55}.\quad \texttt{GlobalTheory.\allowbreak{}Fc.\allowbreak{}R1.\allowbreak{}row\_\allowbreak{}hRB\_\allowbreak{}univ\_\allowbreak{}alt} (\texttt{Development.\allowbreak{}AutomorphicArithmetic.\allowbreak{}ThetaCorrespondence.\allowbreak{}ThetaPairLifts.\allowbreak{}ThirdPlaceSeries.\allowbreak{}RowHRBUniv}, lines 19--24)\\
\texttt{Formalism.\allowbreak{}IndefR.\allowbreak{}IndefiniteRealisationBoundedRows} (\texttt{Development.\allowbreak{}AutomorphicArithmetic.\allowbreak{}ThetaCorrespondence.\allowbreak{}ThetaPairLifts.\allowbreak{}ThirdPlaceSeries.\allowbreak{}IndefRRealTB}, lines 39--41)\par}

{\footnotesize\raggedright\noindent \ref{B8.55}.\quad \texttt{GlobalTheory.\allowbreak{}AdTh.\allowbreak{}IoExtCore.\allowbreak{}ioExt\_\allowbreak{}core} (\texttt{Development.\allowbreak{}AbelianVarietyTheory.\allowbreak{}Proofs.\allowbreak{}DivisorProofs.\allowbreak{}IoExtCore}, lines 22--52)\\
\texttt{GlobalTheory.\allowbreak{}RowsAtRealTBClose.\allowbreak{}rowsB\_\allowbreak{}of\_\allowbreak{}decay\_\allowbreak{}alt} (\texttt{Development.\allowbreak{}AutomorphicArithmetic.\allowbreak{}ThetaCorrespondence.\allowbreak{}ThetaPairLifts.\allowbreak{}ThirdPlaceSeries.\allowbreak{}RowsBOfDecay}, lines 16--20)\par}

{\footnotesize\raggedright\noindent \ref{B8.55}.\quad \texttt{Formalism.\allowbreak{}IndefR.\allowbreak{}FWT.\allowbreak{}OMPins} (\texttt{Formalism.\allowbreak{}LocalToGlobal.\allowbreak{}IndefiniteReal.\allowbreak{}OneMapT}, lines 143--145)\\
\texttt{GlobalTheory.\allowbreak{}OneMapAll.\allowbreak{}ClassSumsSpan} (\texttt{Development.\allowbreak{}AbelianVarietyTheory.\allowbreak{}Proofs.\allowbreak{}DivisorProofs.\allowbreak{}OneMapOfOfExtends}, lines 38--40)\\
\texttt{GlobalTheory.\allowbreak{}ClassSumRule.\allowbreak{}exists\_\allowbreak{}cusp\_\allowbreak{}isThClsD\_\allowbreak{}core} (\texttt{Development.\allowbreak{}AbelianVarietyTheory.\allowbreak{}Proofs.\allowbreak{}DivisorProofs.\allowbreak{}ExistsCuspFormClassSumNormalisedData}, lines 17--49)\par}

{\footnotesize\raggedright\noindent \ref{B8.55}.\quad \texttt{GlobalTheory.\allowbreak{}ClassSumSpan.\allowbreak{}rdU\_\allowbreak{}mem\_\allowbreak{}span\_\allowbreak{}raw} (\texttt{Development.\allowbreak{}AbelianVarietyTheory.\allowbreak{}Proofs.\allowbreak{}DivisorProofs.\allowbreak{}ClassSumsSpan}, lines 199--236)\\
\texttt{GlobalTheory.\allowbreak{}RawSpan.\allowbreak{}holSub0\_\allowbreak{}le\_\allowbreak{}span\_\allowbreak{}raw} (\texttt{Development.\allowbreak{}AbelianVarietyTheory.\allowbreak{}Proofs.\allowbreak{}DivisorProofs.\allowbreak{}HolSub0LeSpanRaw}, lines 146--177)\\
\texttt{GlobalTheory.\allowbreak{}DatumMove.\allowbreak{}rawSeries\_\allowbreak{}move\_\allowbreak{}datum} (\texttt{Development.\allowbreak{}AbelianVarietyTheory.\allowbreak{}Proofs.\allowbreak{}DivisorProofs.\allowbreak{}RawSeriesMoveDatum}, lines 209--240)\par}

{\footnotesize\raggedright\noindent \ref{B8.55}.\quad \texttt{GlobalTheory.\allowbreak{}ClassSumSpan.\allowbreak{}exists\_\allowbreak{}split\_\allowbreak{}move} (\texttt{Development.\allowbreak{}AbelianVarietyTheory.\allowbreak{}Proofs.\allowbreak{}DivisorProofs.\allowbreak{}ClassSumsSpan}, lines 76--116)\\
\texttt{GlobalTheory.\allowbreak{}DatumMove.\allowbreak{}exists\_\allowbreak{}move\_\allowbreak{}const} (\texttt{Development.\allowbreak{}AbelianVarietyTheory.\allowbreak{}Proofs.\allowbreak{}DivisorProofs.\allowbreak{}RawSeriesMoveDatum}, lines 148--197)\par}

{\footnotesize\raggedright\noindent \ref{B8.55}.\quad \texttt{GlobalTheory.\allowbreak{}ClassSumSpanCore.\allowbreak{}kernel\_\allowbreak{}eq\_\allowbreak{}shortVal} (\texttt{Development.\allowbreak{}AbelianVarietyTheory.\allowbreak{}Proofs.\allowbreak{}DivisorProofs.\allowbreak{}ClassSumsSpanCuspFormsBaseTupleSecondField}, lines 169--265)\\
\texttt{GlobalTheory.\allowbreak{}ClassSumSpanCore.\allowbreak{}lattice\_\allowbreak{}scalar\_\allowbreak{}value} (\texttt{Development.\allowbreak{}AbelianVarietyTheory.\allowbreak{}Proofs.\allowbreak{}DivisorProofs.\allowbreak{}ClassSumsSpanCuspFormsBaseTupleSecondField}, lines 113--166)\\
\texttt{GlobalTheory.\allowbreak{}PtId.\allowbreak{}cmpConst} (\texttt{Development.\allowbreak{}AbelianVarietyTheory.\allowbreak{}Proofs.\allowbreak{}IwasawaDecomposition.\allowbreak{}LevelMapExtension}, lines 29--31)\par}

{\footnotesize\raggedright\noindent \ref{B8.55}.\quad \texttt{GlobalTheory.\allowbreak{}KernelSpanCore.\allowbreak{}kernel\_\allowbreak{}mem\_\allowbreak{}span\_\allowbreak{}core} (\texttt{Development.\allowbreak{}AbelianVarietyTheory.\allowbreak{}Proofs.\allowbreak{}DivisorProofs.\allowbreak{}ThetaKernelFinitePointMemSpanRationalFormLifts}, lines 22--58)\par}

{\footnotesize\raggedright\noindent \ref{B8.55}.\quad \texttt{GlobalTheory.\allowbreak{}AdTh.\allowbreak{}SerOf} (\texttt{Development.\allowbreak{}AbelianVarietyTheory.\allowbreak{}Definitions.\allowbreak{}Divisors.\allowbreak{}MarkAd}, lines 47--58)\\
\texttt{GlobalTheory.\allowbreak{}T.\allowbreak{}exists\_\allowbreak{}adelicLift\_\allowbreak{}eq\_\allowbreak{}geometric\_\allowbreak{}classSum\_\allowbreak{}all\_\allowbreak{}tuples} (\texttt{Development.\allowbreak{}AbelianVarietyTheory.\allowbreak{}Proofs.\allowbreak{}DivisorProofs.\allowbreak{}ExistsAdelicLiftEqGeometricClassSumAllTuplesStatement}, lines 14--16)\par}

{\footnotesize\raggedright\noindent \ref{B8.55}.\quad \texttt{GlobalTheory.\allowbreak{}ClassSumSpanCore.\allowbreak{}classSumsSpan\_\allowbreak{}of\_\allowbreak{}pins\_\allowbreak{}core} (\texttt{Development.\allowbreak{}AbelianVarietyTheory.\allowbreak{}Proofs.\allowbreak{}DivisorProofs.\allowbreak{}ClassSumsSpanCuspFormsBaseTupleSecondField}, lines 284--352)\par}

{\footnotesize\raggedright\noindent \ref{B8.55}.\quad \texttt{GlobalTheory.\allowbreak{}OneMapCore.\allowbreak{}oneMapRowR\_\allowbreak{}of\_\allowbreak{}core} (\texttt{Development.\allowbreak{}AbelianVarietyTheory.\allowbreak{}Proofs.\allowbreak{}DivisorProofs.\allowbreak{}OneMapRowAllTuplesOfCoreClauses}, lines 22--31)\\
\texttt{GlobalTheory.\allowbreak{}OneMapCore.\allowbreak{}oneMapRowR\_\allowbreak{}of\_\allowbreak{}core\_\allowbreak{}of\_\allowbreak{}size} (\texttt{Development.\allowbreak{}AbelianVarietyTheory.\allowbreak{}Proofs.\allowbreak{}DivisorProofs.\allowbreak{}OneMapAllTuplesOverCoreOfRankBound}, lines 21--29)\\
\texttt{GlobalTheory.\allowbreak{}RowsAtRealTB.\allowbreak{}rowsB\_\allowbreak{}of\_\allowbreak{}clauses\_\allowbreak{}alt} (\texttt{Development.\allowbreak{}AutomorphicArithmetic.\allowbreak{}ThetaCorrespondence.\allowbreak{}ThetaPairLifts.\allowbreak{}ThirdPlaceSeries.\allowbreak{}RowsBOfClauses}, lines 16--19)\\
\texttt{Formalism.\allowbreak{}IndefR.\allowbreak{}FWT.\allowbreak{}OneMapRowR} (\texttt{Formalism.\allowbreak{}LocalToGlobal.\allowbreak{}IndefiniteReal.\allowbreak{}OneMapT}, lines 151--156)\par}

{\footnotesize\raggedright\noindent \ref{B8.56f}.\quad \texttt{GlobalTheory.\allowbreak{}Corev.\allowbreak{}CmTwoClosing.\allowbreak{}cmGoodModel\_\allowbreak{}of\_\allowbreak{}globalArtinReciprocity} (\texttt{Development.\allowbreak{}AbelianVarietyTheory.\allowbreak{}Proofs.\allowbreak{}DecoratedLatticeProofs.\allowbreak{}CmGoodModelOfGlobalArtinReciprocity}, lines 10--12)\\
\texttt{GlobalTheory.\allowbreak{}HondaCM3.\allowbreak{}cm3} (\texttt{Development.\allowbreak{}AbelianVarietyTheory.\allowbreak{}Proofs.\allowbreak{}DecoratedLatticeProofs.\allowbreak{}OneSubTwoOmegaSq}, lines 77--159)\par}

{\footnotesize\raggedright\noindent \ref{B8.56}.\quad \texttt{GlobalTheory.\allowbreak{}T.\allowbreak{}hServ\_\allowbreak{}of\_\allowbreak{}linkR} (\texttt{Development.\allowbreak{}AutomorphicArithmetic.\allowbreak{}ThetaCorrespondence.\allowbreak{}ThetaPairLifts.\allowbreak{}ThirdPlaceSeries.\allowbreak{}RowHServOfLinkStatement}, lines 14--15)\\
\texttt{GlobalTheory.\allowbreak{}Cp.\allowbreak{}SelectionServed} (\texttt{Development.\allowbreak{}AutomorphicArithmetic.\allowbreak{}ThetaCorrespondence.\allowbreak{}ThetaPairLifts.\allowbreak{}ThirdPlaceSeries.\allowbreak{}CountsIrreducibilityConstants}, lines 48--62)\par}

{\footnotesize\raggedright\noindent \ref{B8.56e}.\quad \texttt{Formalism.\allowbreak{}L2.\allowbreak{}KerInert.\allowbreak{}beProp\_\allowbreak{}ofA} (\texttt{Formalism.\allowbreak{}LocalToGlobal.\allowbreak{}InertKernelTheory.\allowbreak{}FinalA}, lines 18--22)\\
\texttt{Formalism.\allowbreak{}L2.\allowbreak{}KerInert.\allowbreak{}BeCob.\allowbreak{}coboundary\_\allowbreak{}exact} (\texttt{Formalism.\allowbreak{}LocalToGlobal.\allowbreak{}InertKernelTheory.\allowbreak{}InertCoprimeLemmas.\allowbreak{}ExistsWittSeriesInverseFormalGroupIso}, lines 111--206)\\
\texttt{Formalism.\allowbreak{}L2.\allowbreak{}Th1.\allowbreak{}C.\allowbreak{}norm\_\allowbreak{}adjoint\_\allowbreak{}all\_\allowbreak{}places} (\texttt{Formalism.\allowbreak{}LocalToGlobal.\allowbreak{}TorsionIntegrality.\allowbreak{}SecondCoefficientNorms.\allowbreak{}BinomialSeriesCoefficients.\allowbreak{}UnitaryInverseIdentity}, lines 424--461)\par}

{\footnotesize\raggedright\noindent \ref{B8.56d}.\quad \texttt{GlobalTheory.\allowbreak{}OrbitBoundStdRec.\allowbreak{}orbit\_\allowbreak{}bound\_\allowbreak{}rec} (\texttt{Development.\allowbreak{}AbelianVarietyTheory.\allowbreak{}Proofs.\allowbreak{}SatakeParameters.\allowbreak{}OrbitBoundStandardSeriesInertPlaces}, lines 22--42)\par}

{\footnotesize\raggedright\noindent \ref{B8.56a}.\quad \texttt{Formalism.\allowbreak{}L2.\allowbreak{}KerInert.\allowbreak{}kernels\_\allowbreak{}integral\_\allowbreak{}inert} (\texttt{Formalism.\allowbreak{}LocalToGlobal.\allowbreak{}InertKernelTheory.\allowbreak{}KernelsIntegralInert}, lines 7--9)\par}

{\footnotesize\raggedright\noindent \ref{B8.56b}.\quad \texttt{Formalism.\allowbreak{}L2.\allowbreak{}KerInert.\allowbreak{}Rep.\allowbreak{}k1gen\_\allowbreak{}of\_\allowbreak{}orth} (\texttt{Formalism.\allowbreak{}LocalToGlobal.\allowbreak{}InertKernelTheory.\allowbreak{}InertCoprimeLemmas.\allowbreak{}RepAsm}, lines 17--24)\par}

{\footnotesize\raggedright\noindent \ref{B8.56c}.\quad \texttt{GlobalTheory.\allowbreak{}TwV.\allowbreak{}site\_\allowbreak{}cmpLite\_\allowbreak{}siteAdm} (\texttt{Development.\allowbreak{}AbelianVarietyTheory.\allowbreak{}Proofs.\allowbreak{}SatakeParameters.\allowbreak{}JointSiteChainComparedCarriesAdmittedSystem}, lines 77--89)\\
\texttt{GlobalTheory.\allowbreak{}TwV.\allowbreak{}Lt.\allowbreak{}AdmAtEnt2} (\texttt{Development.\allowbreak{}AbelianVarietyTheory.\allowbreak{}Definitions.\allowbreak{}GaloisGeometry.\allowbreak{}ThetaSeriesDataHypothesesBelowMaximalLevel}, lines 106--129)\par}

\section*{Chapter~\ref{ch:p8c6}: The chain length}

{\footnotesize\raggedright\noindent \ref{B8.57}.\quad \texttt{GlobalTheory.\allowbreak{}Mem.\allowbreak{}Num.\allowbreak{}evenCeil} (\texttt{Development.\allowbreak{}AbelianVarietyTheory.\allowbreak{}Definitions.\allowbreak{}AutomorphicMembership.\allowbreak{}EvenCeil}, lines 11--11)\\
\texttt{GlobalTheory.\allowbreak{}Mem.\allowbreak{}Num.\allowbreak{}mAt} (\texttt{Development.\allowbreak{}AbelianVarietyTheory.\allowbreak{}Definitions.\allowbreak{}AutomorphicMembership.\allowbreak{}EvenCeil}, lines 21--21)\\
\texttt{Formalism.\allowbreak{}BasicNotions.\allowbreak{}Scale.\allowbreak{}delta1} (\texttt{Formalism.\allowbreak{}LocalRingCardinality.\allowbreak{}DataParams₂}, lines 45--45)\par}

{\footnotesize\raggedright\noindent \ref{B8.58}.\quad \texttt{GlobalTheory.\allowbreak{}Mem.\allowbreak{}Num.\allowbreak{}RChosen} (\texttt{Development.\allowbreak{}AbelianVarietyTheory.\allowbreak{}Definitions.\allowbreak{}AutomorphicMembership.\allowbreak{}EventuallyThickCardLeCyclicOrder}, lines 60--62)\\
\texttt{Formalism.\allowbreak{}HC.\allowbreak{}Frame.\allowbreak{}o} (\texttt{Formalism.\allowbreak{}SharpRestriction.\allowbreak{}RamificationFrobenius.\allowbreak{}TotientGaloisClasses.\allowbreak{}ZsmulGeneratorEqZeroOfProjection}, lines 70--70)\\
\texttt{GlobalTheory.\allowbreak{}Mem.\allowbreak{}Num.\allowbreak{}kAt} (\texttt{Development.\allowbreak{}AbelianVarietyTheory.\allowbreak{}Definitions.\allowbreak{}AutomorphicMembership.\allowbreak{}EvenCeil}, lines 23--23)\\
\texttt{GlobalTheory.\allowbreak{}Mem.\allowbreak{}Num.\allowbreak{}nWAt} (\texttt{Development.\allowbreak{}AbelianVarietyTheory.\allowbreak{}Definitions.\allowbreak{}AutomorphicMembership.\allowbreak{}EvenCeil}, lines 27--27)\\
\texttt{GlobalTheory.\allowbreak{}Mem.\allowbreak{}Num.\allowbreak{}tAt} (\texttt{Development.\allowbreak{}AbelianVarietyTheory.\allowbreak{}Definitions.\allowbreak{}AutomorphicMembership.\allowbreak{}EvenCeil}, lines 33--33)\\
\texttt{GlobalTheory.\allowbreak{}Mem.\allowbreak{}Num.\allowbreak{}RanksOK} (\texttt{Development.\allowbreak{}AbelianVarietyTheory.\allowbreak{}Definitions.\allowbreak{}AutomorphicMembership.\allowbreak{}EvenCeil}, lines 41--43)\\
\texttt{GlobalTheory.\allowbreak{}Mem.\allowbreak{}Num.\allowbreak{}ranksOf} (\texttt{Development.\allowbreak{}AbelianVarietyTheory.\allowbreak{}Definitions.\allowbreak{}AutomorphicMembership.\allowbreak{}EvenCeil}, lines 45--59)\\
\texttt{Formalism.\allowbreak{}BasicNotions.\allowbreak{}Ranks} (\texttt{Formalism.\allowbreak{}LocalRingCardinality.\allowbreak{}Params}, lines 6--23)\par}

{\footnotesize\raggedright\noindent \ref{B8.59}.\quad \texttt{GlobalTheory.\allowbreak{}EX.\allowbreak{}exists\_\allowbreak{}ranks\_\allowbreak{}rankVal\_\allowbreak{}frame} (\texttt{Development.\allowbreak{}AbelianVarietyTheory.\allowbreak{}Proofs.\allowbreak{}GlobalInterpolation.\allowbreak{}ExistsRanksRankValFrame}, lines 45--99)\\
\texttt{GlobalTheory.\allowbreak{}EX.\allowbreak{}TOv} (\texttt{Development.\allowbreak{}AbelianVarietyTheory.\allowbreak{}Definitions.\allowbreak{}GlobalTwisting.\allowbreak{}PairRowIffQdRows}, lines 17--17)\\
\texttt{GlobalTheory.\allowbreak{}EX.\allowbreak{}nWv} (\texttt{Development.\allowbreak{}AbelianVarietyTheory.\allowbreak{}Definitions.\allowbreak{}GlobalTwisting.\allowbreak{}PairRowIffQdRows}, lines 19--20)\\
\texttt{GlobalTheory.\allowbreak{}EX.\allowbreak{}RankVal} (\texttt{Development.\allowbreak{}AbelianVarietyTheory.\allowbreak{}Definitions.\allowbreak{}GlobalTwisting.\allowbreak{}PairRowIffQdRows}, lines 28--32)\par}

{\footnotesize\raggedright\noindent \ref{B8.59}.\quad \texttt{GlobalTheory.\allowbreak{}EX.\allowbreak{}exists\_\allowbreak{}even\_\allowbreak{}between} (\texttt{Development.\allowbreak{}AbelianVarietyTheory.\allowbreak{}Proofs.\allowbreak{}GlobalInterpolation.\allowbreak{}ExistsRanksRankValFrame}, lines 10--16)\\
\texttt{GlobalTheory.\allowbreak{}EX.\allowbreak{}exists\_\allowbreak{}ranks\_\allowbreak{}of\_\allowbreak{}frame} (\texttt{Development.\allowbreak{}AbelianVarietyTheory.\allowbreak{}Proofs.\allowbreak{}GlobalInterpolation.\allowbreak{}ExistsRanksRankValFrame}, lines 24--43)\par}

\section*{Chapter~\ref{ch:galois-polarized}: Galois representations attached to polarized cuspidal representations}

{\footnotesize\raggedright\noindent \ref{B9.01}.\quad \texttt{UpperTheory.\allowbreak{}LiteratureTheorems.\allowbreak{}OpenLiteratureStatements} (\texttt{Development.\allowbreak{}AutomorphicArithmetic.\allowbreak{}ThetaCorrespondence.\allowbreak{}ThetaPairLifts.\allowbreak{}ThirdPlaceSeries.\allowbreak{}EveryCharacterOfLiteratureOfVWin}, lines 215--218)\\
\texttt{UpperTheory.\allowbreak{}FinalTheorems.\allowbreak{}everyCharacter\_\allowbreak{}of\_\allowbreak{}row} (\texttt{Development.\allowbreak{}AutomorphicArithmetic.\allowbreak{}CharacterBoxes.\allowbreak{}MainConclusion}, lines 7--9)\\
\texttt{UpperTheory.\allowbreak{}FinalTheorems.\allowbreak{}everyCharacter} (\texttt{Development.\allowbreak{}AutomorphicArithmetic.\allowbreak{}CharacterBoxes.\allowbreak{}MainConclusion}, lines 10--12)\par}

{\footnotesize\raggedright\noindent \ref{B9.02}.\quad \texttt{UpperTheory.\allowbreak{}Gal.\allowbreak{}Int.\allowbreak{}galoisLocalGlobalStrong\_\allowbreak{}recParameter} (\texttt{Development.\allowbreak{}AutomorphicArithmetic.\allowbreak{}Foundations.\allowbreak{}GaloisRepresentations.\allowbreak{}IntegralStructures.\allowbreak{}RowSrc}, lines 6--7)\par}

{\footnotesize\raggedright\noindent \ref{B9.03}.\quad \texttt{UpperTheory.\allowbreak{}Gal.\allowbreak{}Int.\allowbreak{}galoisLocalGlobalStrong\_\allowbreak{}of\_\allowbreak{}geometry\_\allowbreak{}lift} (\texttt{Development.\allowbreak{}AutomorphicArithmetic.\allowbreak{}Foundations.\allowbreak{}GaloisRepresentations.\allowbreak{}IntegralStructures.\allowbreak{}Development}, lines 30--32)\\
\texttt{UpperTheory.\allowbreak{}Gal.\allowbreak{}Int.\allowbreak{}galoisRepresentations\_\allowbreak{}of\_\allowbreak{}thetaLift} (\texttt{Development.\allowbreak{}AutomorphicArithmetic.\allowbreak{}Foundations.\allowbreak{}GaloisRepresentations.\allowbreak{}IntegralStructures.\allowbreak{}GaloisRepresentationsOfThetaLift}, lines 6--7)\par}

{\footnotesize\raggedright\noindent \ref{B9.04}.\quad \texttt{UpperTheory.\allowbreak{}GalStar.\allowbreak{}PropsB} (\texttt{Development.\allowbreak{}AutomorphicArithmetic.\allowbreak{}Foundations.\allowbreak{}GaloisRepresentations.\allowbreak{}TypeConditions.\allowbreak{}WeightsConjBis}, lines 54--68)\\
\texttt{UpperTheory.\allowbreak{}GalStar.\allowbreak{}Row0b} (\texttt{Development.\allowbreak{}AutomorphicArithmetic.\allowbreak{}Foundations.\allowbreak{}GaloisRepresentations.\allowbreak{}TypeConditions.\allowbreak{}WeightsConjBis}, lines 69--74)\par}

{\footnotesize\raggedright\noindent \ref{B9.05}.\quad \texttt{UpperTheory.\allowbreak{}GalStar.\allowbreak{}AwayB} (\texttt{Development.\allowbreak{}AutomorphicArithmetic.\allowbreak{}Foundations.\allowbreak{}GaloisRepresentations.\allowbreak{}TypeConditions.\allowbreak{}WeightsConjBis}, lines 26--32)\\
\texttt{UpperTheory.\allowbreak{}GalStar.\allowbreak{}Sqf} (\texttt{Development.\allowbreak{}AutomorphicArithmetic.\allowbreak{}Foundations.\allowbreak{}GaloisRepresentations.\allowbreak{}TypeConditions.\allowbreak{}WeightsConjAlt}, lines 29--34)\par}

{\footnotesize\raggedright\noindent \ref{B9.06}.\quad \texttt{Formalism.\allowbreak{}BasicNotions.\allowbreak{}LocTied} (\texttt{Formalism.\allowbreak{}LocalRingCardinality.\allowbreak{}AdicCompletionLocal.\allowbreak{}LocTied}, lines 11--37)\par}

{\footnotesize\raggedright\noindent \ref{B9.06}.\quad \texttt{UpperTheory.\allowbreak{}GalStar.\allowbreak{}AboveB} (\texttt{Development.\allowbreak{}AutomorphicArithmetic.\allowbreak{}Foundations.\allowbreak{}GaloisRepresentations.\allowbreak{}TypeConditions.\allowbreak{}WeightsConjBis}, lines 39--53)\par}

{\footnotesize\raggedright\noindent \ref{B9.07a}.\quad \texttt{GlobalTheory.\allowbreak{}RegAlgCharAlg.\allowbreak{}isAlgPos\_\allowbreak{}of\_\allowbreak{}polarizedPair} (\texttt{Development.\allowbreak{}AbelianVarietyTheory.\allowbreak{}Proofs.\allowbreak{}SatakeParameters.\allowbreak{}IsAlgebraicHeckeCharOfRegularAlgebraicPolarizedP}, lines 37--63)\\
\texttt{UpperTheory.\allowbreak{}heckeChar\_\allowbreak{}parallel\_\allowbreak{}weight} (\texttt{Development.\allowbreak{}AutomorphicArithmetic.\allowbreak{}AutomorphicVarieties.\allowbreak{}BuildingsResidueChars.\allowbreak{}HeckeCharParallelWeight}, lines 234--290)\\
\texttt{UpperTheory.\allowbreak{}Gal.\allowbreak{}Int.\allowbreak{}Row.\allowbreak{}exists\_\allowbreak{}weights} (\texttt{Development.\allowbreak{}AutomorphicArithmetic.\allowbreak{}Foundations.\allowbreak{}GaloisRepresentations.\allowbreak{}IntegralStructures.\allowbreak{}InertCharpoly.\allowbreak{}GaloisLocalGlobalOfLocalGlobalCompatibility}, lines 127--149)\par}

{\footnotesize\raggedright\noindent \ref{B9.07}.\quad \texttt{UpperTheory.\allowbreak{}Gal.\allowbreak{}Int.\allowbreak{}Row.\allowbreak{}galoisLocalGlobal\_\allowbreak{}of\_\allowbreak{}placeClassCompatibility} (\texttt{Development.\allowbreak{}AutomorphicArithmetic.\allowbreak{}Foundations.\allowbreak{}GaloisRepresentations.\allowbreak{}IntegralStructures.\allowbreak{}InertCharpoly.\allowbreak{}GaloisLocalGlobalOfLocalGlobalCompatibility}, lines 157--172)\\
\texttt{UpperTheory.\allowbreak{}Gal.\allowbreak{}Int.\allowbreak{}Row.\allowbreak{}away} (\texttt{Development.\allowbreak{}AutomorphicArithmetic.\allowbreak{}Foundations.\allowbreak{}GaloisRepresentations.\allowbreak{}IntegralStructures.\allowbreak{}InertCharpoly.\allowbreak{}GaloisLocalGlobalOfLocalGlobalCompatibility}, lines 96--106)\\
\texttt{UpperTheory.\allowbreak{}Gal.\allowbreak{}Int.\allowbreak{}Row.\allowbreak{}above} (\texttt{Development.\allowbreak{}AutomorphicArithmetic.\allowbreak{}Foundations.\allowbreak{}GaloisRepresentations.\allowbreak{}IntegralStructures.\allowbreak{}InertCharpoly.\allowbreak{}GaloisLocalGlobalOfLocalGlobalCompatibility}, lines 108--126)\\
\texttt{UpperTheory.\allowbreak{}Gal.\allowbreak{}Int.\allowbreak{}rest} (\texttt{Development.\allowbreak{}AutomorphicArithmetic.\allowbreak{}Foundations.\allowbreak{}GaloisRepresentations.\allowbreak{}IntegralStructures.\allowbreak{}Development}, lines 24--26)\par}

{\footnotesize\raggedright\noindent \ref{B9.08}.\quad \texttt{UpperTheory.\allowbreak{}Gal.\allowbreak{}Int.\allowbreak{}Geo.\allowbreak{}GeoOne} (\texttt{Development.\allowbreak{}AutomorphicArithmetic.\allowbreak{}Foundations.\allowbreak{}GaloisRepresentations.\allowbreak{}IntegralStructures.\allowbreak{}GeometricFlagBases.\allowbreak{}HodgeTrAt}, lines 79--91)\\
\texttt{UpperTheory.\allowbreak{}Gal.\allowbreak{}Int.\allowbreak{}Geo.\allowbreak{}geoOne\_\allowbreak{}holds} (\texttt{Development.\allowbreak{}AutomorphicArithmetic.\allowbreak{}Foundations.\allowbreak{}GaloisRepresentations.\allowbreak{}IntegralStructures.\allowbreak{}GeometricFlagBases.\allowbreak{}WeightOneGeometricRealization}, lines 11--12)\par}

{\footnotesize\raggedright\noindent \ref{B9.09a}.\quad \texttt{UpperTheory.\allowbreak{}Gal.\allowbreak{}CertificateRingHomMap.\allowbreak{}seedGU\_\allowbreak{}holds} (\texttt{Development.\allowbreak{}AutomorphicArithmetic.\allowbreak{}Foundations.\allowbreak{}GaloisRepresentations.\allowbreak{}ProjectiveAutomorphicHecke.\allowbreak{}NhdsBasisLevelSubgroup}, lines 82--83)\\
\texttt{UpperTheory.\allowbreak{}Gal.\allowbreak{}Int.\allowbreak{}Geo.\allowbreak{}Pl.\allowbreak{}vecEx\_\allowbreak{}of} (\texttt{Development.\allowbreak{}AutomorphicArithmetic.\allowbreak{}Foundations.\allowbreak{}GaloisRepresentations.\allowbreak{}IntegralStructures.\allowbreak{}GeometricFlagBases.\allowbreak{}PrimeLevelLocalization.\allowbreak{}PVecWt}, lines 13--24)\par}

{\footnotesize\raggedright\noindent \ref{B9.09ab}.\quad \texttt{UpperTheory.\allowbreak{}Gal.\allowbreak{}Int.\allowbreak{}Geo.\allowbreak{}Pl.\allowbreak{}master} (\texttt{Development.\allowbreak{}AutomorphicArithmetic.\allowbreak{}Foundations.\allowbreak{}GaloisRepresentations.\allowbreak{}IntegralStructures.\allowbreak{}GeometricFlagBases.\allowbreak{}PrimeLevelLocalization.\allowbreak{}SameKTrans}, lines 86--206)\\
\texttt{UpperTheory.\allowbreak{}Gal.\allowbreak{}Spc.\allowbreak{}eigEx} (\texttt{Development.\allowbreak{}AutomorphicArithmetic.\allowbreak{}Foundations.\allowbreak{}GaloisRepresentations.\allowbreak{}ScalarCommutation.\allowbreak{}FrameTh}, lines 52--92)\\
\texttt{UpperTheory.\allowbreak{}Gal.\allowbreak{}W2.\allowbreak{}Seed.\allowbreak{}pHasEigOf} (\texttt{Development.\allowbreak{}AutomorphicArithmetic.\allowbreak{}Foundations.\allowbreak{}GaloisRepresentations.\allowbreak{}SphericalCoordinates.\allowbreak{}AlgebraEquivSetup.\allowbreak{}PHasEigOf}, lines 45--87)\par}

{\footnotesize\raggedright\noindent \ref{B9.09ac}.\quad \texttt{UpperTheory.\allowbreak{}Gal.\allowbreak{}Spc.\allowbreak{}specTopG} (\texttt{Development.\allowbreak{}AutomorphicArithmetic.\allowbreak{}Foundations.\allowbreak{}GaloisRepresentations.\allowbreak{}ScalarCommutation.\allowbreak{}SpecTopG}, lines 9--22)\\
\texttt{UpperTheory.\allowbreak{}Gal.\allowbreak{}Spc.\allowbreak{}specOfF} (\texttt{Development.\allowbreak{}AutomorphicArithmetic.\allowbreak{}Foundations.\allowbreak{}GaloisRepresentations.\allowbreak{}ScalarCommutation.\allowbreak{}SpecOfF}, lines 179--228)\par}

{\footnotesize\raggedright\noindent \ref{B9.09}.\quad \texttt{UpperTheory.\allowbreak{}Gal.\allowbreak{}Int.\allowbreak{}Geo.\allowbreak{}PlanEx} (\texttt{Development.\allowbreak{}AutomorphicArithmetic.\allowbreak{}Foundations.\allowbreak{}GaloisRepresentations.\allowbreak{}IntegralStructures.\allowbreak{}GeometricFlagBases.\allowbreak{}PackExOf}, lines 82--92)\\
\texttt{UpperTheory.\allowbreak{}Gal.\allowbreak{}Int.\allowbreak{}Geo.\allowbreak{}Pl.\allowbreak{}planEx\_\allowbreak{}holds} (\texttt{Development.\allowbreak{}AutomorphicArithmetic.\allowbreak{}Foundations.\allowbreak{}GaloisRepresentations.\allowbreak{}IntegralStructures.\allowbreak{}GeometricFlagBases.\allowbreak{}PrimeLevelLocalization.\allowbreak{}SpecTopHolds}, lines 9--10)\\
\texttt{UpperTheory.\allowbreak{}Gal.\allowbreak{}Int.\allowbreak{}Geo.\allowbreak{}Plan.\allowbreak{}pack} (\texttt{Development.\allowbreak{}AutomorphicArithmetic.\allowbreak{}Foundations.\allowbreak{}GaloisRepresentations.\allowbreak{}IntegralStructures.\allowbreak{}GeometricFlagBases.\allowbreak{}PackExOf}, lines 101--126)\par}

{\footnotesize\raggedright\noindent \ref{B9.10}.\quad \texttt{GlobalTheory.\allowbreak{}UPel.\allowbreak{}HT.\allowbreak{}Gen.\allowbreak{}GObj} (\texttt{Development.\allowbreak{}AbelianVarietyTheory.\allowbreak{}Definitions.\allowbreak{}Divisors.\allowbreak{}ExistsGDatumOneOneOne}, lines 139--168)\\
\texttt{GlobalTheory.\allowbreak{}UPel.\allowbreak{}HT.\allowbreak{}Gen.\allowbreak{}GenModel} (\texttt{Development.\allowbreak{}AbelianVarietyTheory.\allowbreak{}Definitions.\allowbreak{}Divisors.\allowbreak{}ExistsGDatumOneOneOne}, lines 181--186)\par}

{\footnotesize\raggedright\noindent \ref{B9.10a}.\quad \texttt{GlobalTheory.\allowbreak{}UPel.\allowbreak{}HT.\allowbreak{}model\_\allowbreak{}nonempty} (\texttt{Development.\allowbreak{}AbelianVarietyTheory.\allowbreak{}Proofs.\allowbreak{}DivisorProofs.\allowbreak{}HarrisTaylorIntegralModelNonempty}, lines 11--15)\\
\texttt{GlobalTheory.\allowbreak{}UPel.\allowbreak{}HT.\allowbreak{}exists\_\allowbreak{}repFam\_\allowbreak{}qproj} (\texttt{Development.\allowbreak{}AbelianVarietyTheory.\allowbreak{}Proofs.\allowbreak{}DivisorProofs.\allowbreak{}ExistsQuasiProjectiveHarrisTaylorFamily}, lines 13--36)\\
\texttt{GlobalTheory.\allowbreak{}UPel.\allowbreak{}HT.\allowbreak{}RepFam.\allowbreak{}smoothOfRelativeDimension} (\texttt{Development.\allowbreak{}AbelianVarietyTheory.\allowbreak{}Proofs.\allowbreak{}DivisorProofs.\allowbreak{}SmoothOfRelativeDimension}, lines 15--37)\\
\texttt{GlobalTheory.\allowbreak{}UPel.\allowbreak{}HT.\allowbreak{}Model.\allowbreak{}isProjOver\_\allowbreak{}of\_\allowbreak{}compactB} (\texttt{Development.\allowbreak{}AbelianVarietyTheory.\allowbreak{}Proofs.\allowbreak{}DivisorProofs.\allowbreak{}IsProjOverOfCompactB}, lines 10--14)\par}

{\footnotesize\raggedright\noindent \ref{B9.11}.\quad \texttt{GlobalTheory.\allowbreak{}T.\allowbreak{}unitaryCharZeroGenericModel\_\allowbreak{}nonempty} (\texttt{Development.\allowbreak{}AbelianVarietyTheory.\allowbreak{}Proofs.\allowbreak{}DivisorProofs.\allowbreak{}UnitaryCharZeroGenericModelNonemptyStatement}, lines 8--10)\\
\texttt{GlobalTheory.\allowbreak{}UPel.\allowbreak{}HT.\allowbreak{}Gen.\allowbreak{}genModel\_\allowbreak{}nonempty} (\texttt{Development.\allowbreak{}AbelianVarietyTheory.\allowbreak{}Proofs.\allowbreak{}DivisorProofs.\allowbreak{}UnitaryCharZeroGenericModelNonempty}, lines 118--123)\par}

{\footnotesize\raggedright\noindent \ref{B9.12}.\quad \texttt{UpperTheory.\allowbreak{}Gal.\allowbreak{}Int.\allowbreak{}Geo.\allowbreak{}KS} (\texttt{Development.\allowbreak{}AutomorphicArithmetic.\allowbreak{}Foundations.\allowbreak{}GaloisRepresentations.\allowbreak{}IntegralStructures.\allowbreak{}GeometricFlagBases.\allowbreak{}IntegOfMem}, lines 48--60)\\
\texttt{UpperTheory.\allowbreak{}Gal.\allowbreak{}G3.\allowbreak{}pKSEx} (\texttt{Development.\allowbreak{}AutomorphicArithmetic.\allowbreak{}Foundations.\allowbreak{}GaloisRepresentations.\allowbreak{}PAdicWittLattices.\allowbreak{}ExistsDeeperNormalLevel}, lines 100--173)\\
\texttt{UpperTheory.\allowbreak{}Gal.\allowbreak{}G3.\allowbreak{}kSNice\_\allowbreak{}holds} (\texttt{Development.\allowbreak{}AutomorphicArithmetic.\allowbreak{}Foundations.\allowbreak{}GaloisRepresentations.\allowbreak{}PAdicWittLattices.\allowbreak{}FibrePowerTowerAdmissibleLevel}, lines 152--153)\\
\texttt{UpperTheory.\allowbreak{}Gal.\allowbreak{}G3.\allowbreak{}kSEx\_\allowbreak{}holds} (\texttt{Development.\allowbreak{}AutomorphicArithmetic.\allowbreak{}Foundations.\allowbreak{}GaloisRepresentations.\allowbreak{}PAdicWittLattices.\allowbreak{}FibrePowerTowerAdmissibleLevel}, lines 172--173)\par}

{\footnotesize\raggedright\noindent \ref{B9.12a}.\quad \texttt{UpperTheory.\allowbreak{}baseExtension\_\allowbreak{}holds} (\texttt{Development.\allowbreak{}AutomorphicArithmetic.\allowbreak{}CommutativeAlgebra.\allowbreak{}BaseExtensionHolds}, lines 76--138)\\
\texttt{UpperTheory.\allowbreak{}site\_\allowbreak{}exists\_\allowbreak{}of\_\allowbreak{}splitting\_\allowbreak{}and\_\allowbreak{}power\_\allowbreak{}projector} (\texttt{Development.\allowbreak{}AutomorphicArithmetic.\allowbreak{}CommutativeAlgebra.\allowbreak{}SiteExistsOfSplittingAndPowerProjector}, lines 11--38)\\
\texttt{UpperTheory.\allowbreak{}Gal.\allowbreak{}crystallineComparison.\allowbreak{}Mod.\allowbreak{}esNorm\_\allowbreak{}of} (\texttt{Development.\allowbreak{}AutomorphicArithmetic.\allowbreak{}Foundations.\allowbreak{}GaloisRepresentations.\allowbreak{}AnnihilatorModels.\allowbreak{}MEsNorm}, lines 16--43)\\
\texttt{UpperTheory.\allowbreak{}Gal.\allowbreak{}crystallineComparison.\allowbreak{}Mod.\allowbreak{}ordKSE\_\allowbreak{}of} (\texttt{Development.\allowbreak{}AutomorphicArithmetic.\allowbreak{}Foundations.\allowbreak{}GaloisRepresentations.\allowbreak{}AnnihilatorModels.\allowbreak{}MOrdKSE}, lines 67--149)\\
\texttt{UpperTheory.\allowbreak{}Gal.\allowbreak{}FiniteEtaleModels.\allowbreak{}ksGood\_\allowbreak{}holds} (\texttt{Development.\allowbreak{}AutomorphicArithmetic.\allowbreak{}Foundations.\allowbreak{}GaloisRepresentations.\allowbreak{}FiniteEtaleModels.\allowbreak{}KsGoodHolds}, lines 45--47)\par}

{\footnotesize\raggedright\noindent \ref{B9.13}.\quad \texttt{GlobalTheory.\allowbreak{}EtC.\allowbreak{}HlSub} (\texttt{Development.\allowbreak{}AbelianVarietyTheory.\allowbreak{}Definitions.\allowbreak{}Divisors.\allowbreak{}PullHlComp}, lines 216--224)\\
\texttt{GlobalTheory.\allowbreak{}EtC.\allowbreak{}HQ} (\texttt{Development.\allowbreak{}AbelianVarietyTheory.\allowbreak{}Definitions.\allowbreak{}Divisors.\allowbreak{}PullHlComp}, lines 239--240)\\
\texttt{GlobalTheory.\allowbreak{}GalV.\allowbreak{}Htot} (\texttt{Development.\allowbreak{}AbelianVarietyTheory.\allowbreak{}Definitions.\allowbreak{}Divisors.\allowbreak{}IsCorrOpOfAvg}, lines 11--12)\\
\texttt{GlobalTheory.\allowbreak{}GalV.\allowbreak{}galTotRep} (\texttt{Development.\allowbreak{}AbelianVarietyTheory.\allowbreak{}Definitions.\allowbreak{}Divisors.\allowbreak{}IsCorrOpOfAvg}, lines 58--63)\par}

{\footnotesize\raggedright\noindent \ref{B9.14}.\quad \texttt{UpperTheory.\allowbreak{}Gal.\allowbreak{}G3.\allowbreak{}pHopSpec} (\texttt{Development.\allowbreak{}AutomorphicArithmetic.\allowbreak{}Foundations.\allowbreak{}GaloisRepresentations.\allowbreak{}PAdicWittLattices.\allowbreak{}PHopSpec}, lines 17--51)\\
\texttt{UpperTheory.\allowbreak{}Gal.\allowbreak{}Int.\allowbreak{}Geo.\allowbreak{}hopSpec} (\texttt{Development.\allowbreak{}AutomorphicArithmetic.\allowbreak{}Foundations.\allowbreak{}GaloisRepresentations.\allowbreak{}IntegralStructures.\allowbreak{}GeometricFlagBases.\allowbreak{}RepMain}, lines 56--58)\par}

{\footnotesize\raggedright\noindent \ref{B9.15}.\quad \texttt{UpperTheory.\allowbreak{}Gal.\allowbreak{}Int.\allowbreak{}Geo.\allowbreak{}repEx\_\allowbreak{}of} (\texttt{Development.\allowbreak{}AutomorphicArithmetic.\allowbreak{}Foundations.\allowbreak{}GaloisRepresentations.\allowbreak{}IntegralStructures.\allowbreak{}GeometricFlagBases.\allowbreak{}RepAsm}, lines 13--60)\\
\texttt{UpperTheory.\allowbreak{}Gal.\allowbreak{}Int.\allowbreak{}Geo.\allowbreak{}repEx\_\allowbreak{}core} (\texttt{Development.\allowbreak{}AutomorphicArithmetic.\allowbreak{}Foundations.\allowbreak{}GaloisRepresentations.\allowbreak{}IntegralStructures.\allowbreak{}GeometricFlagBases.\allowbreak{}RepMain}, lines 74--76)\par}

{\footnotesize\raggedright\noindent \ref{B9.16}.\quad \texttt{UpperTheory.\allowbreak{}Gal.\allowbreak{}crystallineComparison.\allowbreak{}ordDecomp\_\allowbreak{}of} (\texttt{Development.\allowbreak{}AutomorphicArithmetic.\allowbreak{}Foundations.\allowbreak{}GaloisRepresentations.\allowbreak{}AnnihilatorModels.\allowbreak{}POrdDecomp}, lines 98--170)\\
\texttt{UpperTheory.\allowbreak{}Gal.\allowbreak{}crystallineComparison.\allowbreak{}telescope\_\allowbreak{}of} (\texttt{Development.\allowbreak{}AutomorphicArithmetic.\allowbreak{}Foundations.\allowbreak{}GaloisRepresentations.\allowbreak{}AnnihilatorModels.\allowbreak{}PTelescope}, lines 47--77)\\
\texttt{UpperTheory.\allowbreak{}Gal.\allowbreak{}crystallineComparison.\allowbreak{}Mod.\allowbreak{}congKS\_\allowbreak{}link} (\texttt{Development.\allowbreak{}AutomorphicArithmetic.\allowbreak{}Foundations.\allowbreak{}GaloisRepresentations.\allowbreak{}AnnihilatorModels.\allowbreak{}QuotObjLink}, lines 13--13)\\
\texttt{UpperTheory.\allowbreak{}Gal.\allowbreak{}crystallineComparison.\allowbreak{}ordCount\_\allowbreak{}holds} (\texttt{Development.\allowbreak{}AutomorphicArithmetic.\allowbreak{}Foundations.\allowbreak{}GaloisRepresentations.\allowbreak{}AnnihilatorModels.\allowbreak{}GenericMulOfProdHolds}, lines 51--51)\par}

{\footnotesize\raggedright\noindent \ref{B9.17}.\quad \texttt{UpperTheory.\allowbreak{}Gal.\allowbreak{}Mod.\allowbreak{}esStep} (\texttt{Development.\allowbreak{}AutomorphicArithmetic.\allowbreak{}Foundations.\allowbreak{}GaloisRepresentations.\allowbreak{}BasicOpenSections.\allowbreak{}PEsStep}, lines 49--95)\\
\texttt{UpperTheory.\allowbreak{}Gal.\allowbreak{}Mod.\allowbreak{}towES} (\texttt{Development.\allowbreak{}AutomorphicArithmetic.\allowbreak{}Foundations.\allowbreak{}GaloisRepresentations.\allowbreak{}BasicOpenSections.\allowbreak{}PTowES}, lines 147--221)\\
\texttt{UpperTheory.\allowbreak{}Gal.\allowbreak{}Int.\allowbreak{}Geo.\allowbreak{}CongFor} (\texttt{Development.\allowbreak{}AutomorphicArithmetic.\allowbreak{}Foundations.\allowbreak{}GaloisRepresentations.\allowbreak{}IntegralStructures.\allowbreak{}GeometricFlagBases.\allowbreak{}HodgeTrAt}, lines 54--57)\\
\texttt{UpperTheory.\allowbreak{}Gal.\allowbreak{}Int.\allowbreak{}Geo.\allowbreak{}Pack.\allowbreak{}rep} (\texttt{Development.\allowbreak{}AutomorphicArithmetic.\allowbreak{}Foundations.\allowbreak{}GaloisRepresentations.\allowbreak{}IntegralStructures.\allowbreak{}GeometricFlagBases.\allowbreak{}MapInvMul}, lines 43--87)\par}

{\footnotesize\raggedright\noindent \ref{B9.18}.\quad \texttt{UpperTheory.\allowbreak{}Gal.\allowbreak{}LocalSmoothExtension.\allowbreak{}HolomorphicLefschetzFormula} (\texttt{Development.\allowbreak{}AutomorphicArithmetic.\allowbreak{}Foundations.\allowbreak{}GaloisRepresentations.\allowbreak{}LefschetzSigmaSums.\allowbreak{}HodgeTraceFromLefschetzAndGroupSide}, lines 102--112)\\
\texttt{UpperTheory.\allowbreak{}Gal.\allowbreak{}Int.\allowbreak{}Geo.\allowbreak{}weight\_\allowbreak{}one\_\allowbreak{}geometric\_\allowbreak{}realization} (\texttt{Development.\allowbreak{}AutomorphicArithmetic.\allowbreak{}Foundations.\allowbreak{}GaloisRepresentations.\allowbreak{}IntegralStructures.\allowbreak{}GeometricFlagBases.\allowbreak{}WeightOneGeometricRealization}, lines 7--9)\par}

{\footnotesize\raggedright\noindent \ref{B9.18b}.\quad \texttt{UpperTheory.\allowbreak{}Gal.\allowbreak{}Cnt.\allowbreak{}weighted\_\allowbreak{}count\_\allowbreak{}definite\_\allowbreak{}unitary} (\texttt{Development.\allowbreak{}AutomorphicArithmetic.\allowbreak{}Foundations.\allowbreak{}GaloisRepresentations.\allowbreak{}DefiniteTraceForms.\allowbreak{}WeightedCountDefiniteUnitary}, lines 10--11)\\
\texttt{UpperTheory.\allowbreak{}Cb.\allowbreak{}complex\_\allowbreak{}count\_\allowbreak{}regular\_\allowbreak{}marked\_\allowbreak{}points} (\texttt{Development.\allowbreak{}AutomorphicArithmetic.\allowbreak{}Foundations.\allowbreak{}EtaleClopenDescent.\allowbreak{}ComplexCountRegularMarkedPoints}, lines 552--571)\\
\texttt{GlobalTheory.\allowbreak{}LKDef.\allowbreak{}cls\_\allowbreak{}equiv\_\allowbreak{}dcls\_\allowbreak{}irr} (\texttt{Development.\allowbreak{}AbelianVarietyTheory.\allowbreak{}Proofs.\allowbreak{}DivisorProofs.\allowbreak{}MarkedFieldPairsEquivDefiniteFieldPairsOdd}, lines 81--187)\par}

{\footnotesize\raggedright\noindent \ref{B9.18a}.\quad \texttt{UpperTheory.\allowbreak{}Gal.\allowbreak{}LocalSmoothExtension.\allowbreak{}groupTop\_\allowbreak{}of} (\texttt{Development.\allowbreak{}AutomorphicArithmetic.\allowbreak{}Foundations.\allowbreak{}GaloisRepresentations.\allowbreak{}LefschetzSigmaSums.\allowbreak{}PGroupTop}, lines 12--66)\\
\texttt{UpperTheory.\allowbreak{}Gal.\allowbreak{}LocalSmoothExtension.\allowbreak{}hodgeEig\_\allowbreak{}of} (\texttt{Development.\allowbreak{}AutomorphicArithmetic.\allowbreak{}Foundations.\allowbreak{}GaloisRepresentations.\allowbreak{}LefschetzSigmaSums.\allowbreak{}PHodgeEig}, lines 6--68)\\
\texttt{UpperTheory.\allowbreak{}Gal.\allowbreak{}LocalSmoothExtension.\allowbreak{}endPol} (\texttt{Development.\allowbreak{}AutomorphicArithmetic.\allowbreak{}Foundations.\allowbreak{}GaloisRepresentations.\allowbreak{}LefschetzSigmaSums.\allowbreak{}InvPol}, lines 43--44)\\
\texttt{UpperTheory.\allowbreak{}Gal.\allowbreak{}LocalSmoothExtension.\allowbreak{}markOf} (\texttt{Development.\allowbreak{}AutomorphicArithmetic.\allowbreak{}Foundations.\allowbreak{}GaloisRepresentations.\allowbreak{}LefschetzSigmaSums.\allowbreak{}InvPol}, lines 45--46)\par}

{\footnotesize\raggedright\noindent \ref{B9.19}.\quad \texttt{UpperTheory.\allowbreak{}Gal.\allowbreak{}Int.\allowbreak{}Geo.\allowbreak{}PhiK} (\texttt{Development.\allowbreak{}AutomorphicArithmetic.\allowbreak{}Foundations.\allowbreak{}GaloisRepresentations.\allowbreak{}IntegralStructures.\allowbreak{}GeometricFlagBases.\allowbreak{}HodgeTrAt}, lines 32--37)\\
\texttt{UpperTheory.\allowbreak{}Gal.\allowbreak{}Int.\allowbreak{}Geo.\allowbreak{}CrysStep} (\texttt{Development.\allowbreak{}AutomorphicArithmetic.\allowbreak{}Foundations.\allowbreak{}GaloisRepresentations.\allowbreak{}IntegralStructures.\allowbreak{}GeometricFlagBases.\allowbreak{}HodgeTrAt}, lines 112--115)\\
\texttt{UpperTheory.\allowbreak{}Gal.\allowbreak{}Int.\allowbreak{}Geo.\allowbreak{}JumpStep} (\texttt{Development.\allowbreak{}AutomorphicArithmetic.\allowbreak{}Foundations.\allowbreak{}GaloisRepresentations.\allowbreak{}IntegralStructures.\allowbreak{}GeometricFlagBases.\allowbreak{}HodgeTrAt}, lines 117--120)\\
\texttt{UpperTheory.\allowbreak{}Gal.\allowbreak{}Int.\allowbreak{}Geo.\allowbreak{}PhiStep} (\texttt{Development.\allowbreak{}AutomorphicArithmetic.\allowbreak{}Foundations.\allowbreak{}GaloisRepresentations.\allowbreak{}IntegralStructures.\allowbreak{}GeometricFlagBases.\allowbreak{}HodgeTrAt}, lines 122--126)\\
\texttt{UpperTheory.\allowbreak{}Gal.\allowbreak{}Int.\allowbreak{}Geo.\allowbreak{}jumpStep\_\allowbreak{}holds} (\texttt{Development.\allowbreak{}AutomorphicArithmetic.\allowbreak{}Foundations.\allowbreak{}GaloisRepresentations.\allowbreak{}IntegralStructures.\allowbreak{}GeometricFlagBases.\allowbreak{}GeometricRealizationOfWeightOne}, lines 18--19)\\
\texttt{UpperTheory.\allowbreak{}Ru.\allowbreak{}crystallineComparison.\allowbreak{}phiStep} (\texttt{Development.\allowbreak{}AutomorphicArithmetic.\allowbreak{}CommutativeAlgebra.\allowbreak{}EtaleCohomology.\allowbreak{}SmoothModelCrystalline}, lines 13--14)\par}

{\footnotesize\raggedright\noindent \ref{B9.20}.\quad \texttt{UpperTheory.\allowbreak{}Gal.\allowbreak{}KunnethLefschetz.\allowbreak{}core\_\allowbreak{}holds} (\texttt{Development.\allowbreak{}AutomorphicArithmetic.\allowbreak{}Foundations.\allowbreak{}GaloisRepresentations.\allowbreak{}KunnethLefschetz.\allowbreak{}CoreHolds}, lines 5--7)\par}

{\footnotesize\raggedright\noindent \ref{B9.21}.\quad \texttt{GlobalTheory.\allowbreak{}EtC.\allowbreak{}etC\_\allowbreak{}c2} (\texttt{Development.\allowbreak{}AbelianVarietyTheory.\allowbreak{}Proofs.\allowbreak{}DivisorProofs.\allowbreak{}SmoothAffineCurveEtaleCohomologicalDimensionLeOn}, lines 9--15)\par}

{\footnotesize\raggedright\noindent \ref{B9.22}.\quad \texttt{UpperTheory.\allowbreak{}Cb.\allowbreak{}etale\_\allowbreak{}cohomology\_\allowbreak{}ind\_\allowbreak{}etale\_\allowbreak{}base\_\allowbreak{}change} (\texttt{Development.\allowbreak{}AutomorphicArithmetic.\allowbreak{}Foundations.\allowbreak{}EtaleClopenDescent.\allowbreak{}EtaleCohomologyIndEtaleBaseChange}, lines 272--309)\\
\texttt{UpperTheory.\allowbreak{}Cb.\allowbreak{}C3.\allowbreak{}local\_\allowbreak{}eqv} (\texttt{Development.\allowbreak{}AutomorphicArithmetic.\allowbreak{}Foundations.\allowbreak{}EtaleClopenDescent.\allowbreak{}EtaleCohomologyIndEtaleBaseChange}, lines 210--247)\par}

{\footnotesize\raggedright\noindent \ref{B9.23}.\quad \texttt{GlobalTheory.\allowbreak{}EtC.\allowbreak{}HomotopyInv.\allowbreak{}homotopy\_\allowbreak{}invariance\_\allowbreak{}etale} (\texttt{Development.\allowbreak{}AbelianVarietyTheory.\allowbreak{}Proofs.\allowbreak{}DivisorProofs.\allowbreak{}HomotopyInvarianceEtale}, lines 36--46)\par}

{\footnotesize\raggedright\noindent \ref{B9.24}.\quad \texttt{GlobalTheory.\allowbreak{}EtC.\allowbreak{}laurentPolynomial\_\allowbreak{}cohomology\_\allowbreak{}strictlyHenselian\_\allowbreak{}vanishing} (\texttt{Development.\allowbreak{}AbelianVarietyTheory.\allowbreak{}Proofs.\allowbreak{}DivisorProofs.\allowbreak{}LaurentPolynomialCohomologyVanishing}, lines 8--9)\par}

{\footnotesize\raggedright\noindent \ref{B9.25}.\quad \texttt{GlobalTheory.\allowbreak{}EtC.\allowbreak{}projectiveBundleFormula\_\allowbreak{}vectorBundle\_\allowbreak{}any\_\allowbreak{}space} (\texttt{Development.\allowbreak{}AbelianVarietyTheory.\allowbreak{}Proofs.\allowbreak{}DivisorProofs.\allowbreak{}ProjectiveBundleFormulaVectorBundleAnySpace}, lines 11--12)\\
\texttt{GlobalTheory.\allowbreak{}T.\allowbreak{}projectiveBundleFormula\_\allowbreak{}discrete\_\allowbreak{}space} (\texttt{Development.\allowbreak{}AbelianVarietyTheory.\allowbreak{}Proofs.\allowbreak{}DivisorProofs.\allowbreak{}ProjectiveBundleFormulaDiscreteSpaceStatement}, lines 5--6)\\
\texttt{GlobalTheory.\allowbreak{}PBFg.\allowbreak{}pbf\_\allowbreak{}of\_\allowbreak{}disc} (\texttt{Development.\allowbreak{}AbelianVarietyTheory.\allowbreak{}Proofs.\allowbreak{}DivisorProofs.\allowbreak{}ProjectiveBundleFormulaOfDiscreteSpaces}, lines 65--77)\par}

{\footnotesize\raggedright\noindent \ref{B9.26}.\quad \texttt{UpperTheory.\allowbreak{}Gal.\allowbreak{}KunnethLefschetz.\allowbreak{}QDual} (\texttt{Development.\allowbreak{}AutomorphicArithmetic.\allowbreak{}Foundations.\allowbreak{}GaloisRepresentations.\allowbreak{}KunnethLefschetz.\allowbreak{}OneMul}, lines 191--193)\\
\texttt{UpperTheory.\allowbreak{}Gal.\allowbreak{}KunnethLefschetz.\allowbreak{}Core.\allowbreak{}qDual} (\texttt{Development.\allowbreak{}AutomorphicArithmetic.\allowbreak{}Foundations.\allowbreak{}GaloisRepresentations.\allowbreak{}KunnethLefschetz.\allowbreak{}EtSideGood}, lines 45--46)\par}

{\footnotesize\raggedright\noindent \ref{B9.27}.\quad \texttt{UpperTheory.\allowbreak{}Gal.\allowbreak{}Mod.\allowbreak{}spSpecIf} (\texttt{Development.\allowbreak{}AutomorphicArithmetic.\allowbreak{}Foundations.\allowbreak{}GaloisRepresentations.\allowbreak{}BasicOpenSections.\allowbreak{}PSpSpecIf}, lines 18--108)\\
\texttt{UpperTheory.\allowbreak{}Gal.\allowbreak{}Mod.\allowbreak{}specGood} (\texttt{Development.\allowbreak{}AutomorphicArithmetic.\allowbreak{}Foundations.\allowbreak{}GaloisRepresentations.\allowbreak{}BasicOpenSections.\allowbreak{}EsGenAll}, lines 8--9)\par}

{\footnotesize\raggedright\noindent \ref{B9.28}.\quad \texttt{GlobalTheory.\allowbreak{}CongL.\allowbreak{}SeedWU4} (\texttt{Development.\allowbreak{}AbelianVarietyTheory.\allowbreak{}Definitions.\allowbreak{}GaloisGeometry.\allowbreak{}LimWU4OfLimWU3}, lines 15--36)\\
\texttt{UpperTheory.\allowbreak{}Gal.\allowbreak{}Int.\allowbreak{}Div.\allowbreak{}seedFinOdd\_\allowbreak{}of\_\allowbreak{}geoOne} (\texttt{Development.\allowbreak{}AutomorphicArithmetic.\allowbreak{}Foundations.\allowbreak{}GaloisRepresentations.\allowbreak{}IntegralStructures.\allowbreak{}ZetaCharacterDivisibility.\allowbreak{}MainFin}, lines 13--16)\\
\texttt{UpperTheory.\allowbreak{}Gal.\allowbreak{}Int.\allowbreak{}Div.\allowbreak{}seedFinOdd\_\allowbreak{}of} (\texttt{Development.\allowbreak{}AutomorphicArithmetic.\allowbreak{}Foundations.\allowbreak{}GaloisRepresentations.\allowbreak{}IntegralStructures.\allowbreak{}ZetaCharacterDivisibility.\allowbreak{}SeedFinOddOf}, lines 44--79)\\
\texttt{UpperTheory.\allowbreak{}Gal.\allowbreak{}Int.\allowbreak{}Div.\allowbreak{}divOne\_\allowbreak{}of} (\texttt{Development.\allowbreak{}AutomorphicArithmetic.\allowbreak{}Foundations.\allowbreak{}GaloisRepresentations.\allowbreak{}IntegralStructures.\allowbreak{}ZetaCharacterDivisibility.\allowbreak{}PDivOne}, lines 35--73)\par}

{\footnotesize\raggedright\noindent \ref{B9.29}.\quad \texttt{UpperTheory.\allowbreak{}Gal.\allowbreak{}Int.\allowbreak{}Div.\allowbreak{}DivSet} (\texttt{Development.\allowbreak{}AutomorphicArithmetic.\allowbreak{}Foundations.\allowbreak{}GaloisRepresentations.\allowbreak{}IntegralStructures.\allowbreak{}ZetaCharacterDivisibility.\allowbreak{}PhiOfDivOf}, lines 121--130)\\
\texttt{LateralTheory.\allowbreak{}DivDb.\allowbreak{}Key} (\texttt{Development.\allowbreak{}EtaleCohomology.\allowbreak{}DivisorCycles.\allowbreak{}CycleCount}, lines 51--55)\\
\texttt{LateralTheory.\allowbreak{}DivDb.\allowbreak{}plus\_\allowbreak{}of} (\texttt{Development.\allowbreak{}EtaleCohomology.\allowbreak{}DivisorCycles.\allowbreak{}DivPlus}, lines 37--70)\\
\texttt{UpperTheory.\allowbreak{}Gal.\allowbreak{}Int.\allowbreak{}Div.\allowbreak{}divSet\_\allowbreak{}holds} (\texttt{Development.\allowbreak{}AutomorphicArithmetic.\allowbreak{}Foundations.\allowbreak{}GaloisRepresentations.\allowbreak{}IntegralStructures.\allowbreak{}ZetaCharacterDivisibility.\allowbreak{}DivisionOneGeometricRepresentation}, lines 52--52)\\
\texttt{LateralTheory.\allowbreak{}DivDb.\allowbreak{}KeyPf.\allowbreak{}Key\_\allowbreak{}of} (\texttt{Development.\allowbreak{}EtaleCohomology.\allowbreak{}DivisorCycles.\allowbreak{}SumMulSum}, lines 90--176)\par}

{\footnotesize\raggedright\noindent \ref{B9.30}.\quad \texttt{UpperTheory.\allowbreak{}Gal.\allowbreak{}Int.\allowbreak{}Div.\allowbreak{}genPin\_\allowbreak{}of} (\texttt{Development.\allowbreak{}AutomorphicArithmetic.\allowbreak{}Foundations.\allowbreak{}GaloisRepresentations.\allowbreak{}IntegralStructures.\allowbreak{}ZetaCharacterDivisibility.\allowbreak{}PGenPin}, lines 28--48)\\
\texttt{UpperTheory.\allowbreak{}Gal.\allowbreak{}Int.\allowbreak{}Div.\allowbreak{}divAll\_\allowbreak{}holds} (\texttt{Development.\allowbreak{}AutomorphicArithmetic.\allowbreak{}Foundations.\allowbreak{}GaloisRepresentations.\allowbreak{}IntegralStructures.\allowbreak{}ZetaCharacterDivisibility.\allowbreak{}DivisionOneGeometricRepresentation}, lines 54--58)\\
\texttt{UpperTheory.\allowbreak{}Gal.\allowbreak{}Int.\allowbreak{}Div.\allowbreak{}detEx\_\allowbreak{}holds} (\texttt{Development.\allowbreak{}AutomorphicArithmetic.\allowbreak{}Foundations.\allowbreak{}GaloisRepresentations.\allowbreak{}IntegralStructures.\allowbreak{}ZetaCharacterDivisibility.\allowbreak{}DivisionOneGeometricRepresentation}, lines 49--49)\\
\texttt{UpperTheory.\allowbreak{}Gal.\allowbreak{}Int.\allowbreak{}Div.\allowbreak{}phiSeed\_\allowbreak{}of} (\texttt{Development.\allowbreak{}AutomorphicArithmetic.\allowbreak{}Foundations.\allowbreak{}GaloisRepresentations.\allowbreak{}IntegralStructures.\allowbreak{}ZetaCharacterDivisibility.\allowbreak{}PPhiSeed}, lines 20--62)\par}

{\footnotesize\raggedright\noindent \ref{B9.30a}.\quad \texttt{UpperTheory.\allowbreak{}Gal.\allowbreak{}Gen.\allowbreak{}DenseGen} (\texttt{Development.\allowbreak{}AutomorphicArithmetic.\allowbreak{}Foundations.\allowbreak{}GaloisRepresentations.\allowbreak{}SeparatingAlgebras.\allowbreak{}GenDataEmpty}, lines 65--70)\\
\texttt{UpperTheory.\allowbreak{}Gal.\allowbreak{}Gen.\allowbreak{}denseGen\_\allowbreak{}pf} (\texttt{Development.\allowbreak{}AutomorphicArithmetic.\allowbreak{}Foundations.\allowbreak{}GaloisRepresentations.\allowbreak{}SeparatingAlgebras.\allowbreak{}DenseGen}, lines 14--51)\\
\texttt{UpperTheory.\allowbreak{}Gal.\allowbreak{}Gen.\allowbreak{}step\_\allowbreak{}pf} (\texttt{Development.\allowbreak{}AutomorphicArithmetic.\allowbreak{}Foundations.\allowbreak{}GaloisRepresentations.\allowbreak{}SeparatingAlgebras.\allowbreak{}StepPf}, lines 8--64)\par}

{\footnotesize\raggedright\noindent \ref{B9.31}.\quad \texttt{GlobalTheory.\allowbreak{}CongL.\allowbreak{}CongUW4} (\texttt{Development.\allowbreak{}AbelianVarietyTheory.\allowbreak{}Definitions.\allowbreak{}GaloisGeometry.\allowbreak{}LimWU4OfLimWU3}, lines 37--56)\\
\texttt{GlobalTheory.\allowbreak{}CongL.\allowbreak{}PowSatU} (\texttt{Development.\allowbreak{}AbelianVarietyTheory.\allowbreak{}Definitions.\allowbreak{}GaloisGeometry.\allowbreak{}HasLevelOffOf}, lines 51--53)\\
\texttt{UpperTheory.\allowbreak{}Gal.\allowbreak{}Int.\allowbreak{}Cong.\allowbreak{}congFinOdd\_\allowbreak{}of\_\allowbreak{}ex} (\texttt{Development.\allowbreak{}AutomorphicArithmetic.\allowbreak{}Foundations.\allowbreak{}GaloisRepresentations.\allowbreak{}IntegralStructures.\allowbreak{}CongruenceOddLevel.\allowbreak{}CongFinOddOfLvLac}, lines 15--16)\\
\texttt{UpperTheory.\allowbreak{}Gal.\allowbreak{}Int.\allowbreak{}Cong.\allowbreak{}congUWG\_\allowbreak{}of} (\texttt{Development.\allowbreak{}AutomorphicArithmetic.\allowbreak{}Foundations.\allowbreak{}GaloisRepresentations.\allowbreak{}IntegralStructures.\allowbreak{}CongruenceOddLevel.\allowbreak{}NotApartOfOverPair}, lines 111--250)\\
\texttt{UpperTheory.\allowbreak{}Gal.\allowbreak{}Int.\allowbreak{}Cong.\allowbreak{}CongUWG.\allowbreak{}congUW4} (\texttt{Development.\allowbreak{}AutomorphicArithmetic.\allowbreak{}Foundations.\allowbreak{}GaloisRepresentations.\allowbreak{}IntegralStructures.\allowbreak{}CongruenceOddLevel.\allowbreak{}NotApartOfOverPair}, lines 252--264)\par}

{\footnotesize\raggedright\noindent \ref{B9.32}.\quad \texttt{GlobalTheory.\allowbreak{}Cong.\allowbreak{}PolyApprox} (\texttt{Development.\allowbreak{}AbelianVarietyTheory.\allowbreak{}Definitions.\allowbreak{}GaloisGeometry.\allowbreak{}PolyApproxZeroPadic}, lines 10--13)\par}

{\footnotesize\raggedright\noindent \ref{B9.33}.\quad \texttt{GlobalTheory.\allowbreak{}Cong.\allowbreak{}exists\_\allowbreak{}limit\_\allowbreak{}rep} (\texttt{Development.\allowbreak{}AbelianVarietyTheory.\allowbreak{}Proofs.\allowbreak{}InertiaEigenvalues.\allowbreak{}ExistsLimitRep}, lines 14--29)\par}

{\footnotesize\raggedright\noindent \ref{B9.34}.\quad \texttt{GlobalTheory.\allowbreak{}PsChar.\allowbreak{}exists\_\allowbreak{}ssRep} (\texttt{Development.\allowbreak{}AbelianVarietyTheory.\allowbreak{}Proofs.\allowbreak{}InertiaEigenvalues.\allowbreak{}FiniteSemisimple}, lines 64--146)\\
\texttt{GlobalTheory.\allowbreak{}SSCont.\allowbreak{}continuous\_\allowbreak{}of\_\allowbreak{}continuous\_\allowbreak{}trace} (\texttt{Development.\allowbreak{}AbelianVarietyTheory.\allowbreak{}Proofs.\allowbreak{}InertiaEigenvalues.\allowbreak{}ContinuousOfContinuousTrace}, lines 60--70)\par}

{\footnotesize\raggedright\noindent \ref{B9.35}.\quad \texttt{UpperTheory.\allowbreak{}Gal.\allowbreak{}Lim.\allowbreak{}limFinOdd\_\allowbreak{}of} (\texttt{Development.\allowbreak{}AutomorphicArithmetic.\allowbreak{}Foundations.\allowbreak{}GaloisRepresentations.\allowbreak{}TraceLimits.\allowbreak{}LimFinOddOf}, lines 197--198)\\
\texttt{UpperTheory.\allowbreak{}Gal.\allowbreak{}Lim.\allowbreak{}exists\_\allowbreak{}trace\_\allowbreak{}limit} (\texttt{Development.\allowbreak{}AutomorphicArithmetic.\allowbreak{}Foundations.\allowbreak{}GaloisRepresentations.\allowbreak{}TraceLimits.\allowbreak{}ExistsTraceLimit}, lines 15--170)\\
\texttt{UpperTheory.\allowbreak{}Gal.\allowbreak{}Lim.\allowbreak{}limWU4\_\allowbreak{}of} (\texttt{Development.\allowbreak{}AutomorphicArithmetic.\allowbreak{}Foundations.\allowbreak{}GaloisRepresentations.\allowbreak{}TraceLimits.\allowbreak{}LimFinOddOf}, lines 16--196)\\
\texttt{GlobalTheory.\allowbreak{}CongL.\allowbreak{}crysDRU\_\allowbreak{}of\_\allowbreak{}linApprox} (\texttt{Development.\allowbreak{}AbelianVarietyTheory.\allowbreak{}Proofs.\allowbreak{}SatakeParameters.\allowbreak{}BlowUpFibreSemilocal}, lines 67--97)\par}

{\footnotesize\raggedright\noindent \ref{B9.35aa}.\quad \texttt{GlobalTheory.\allowbreak{}BKLim.\allowbreak{}exists\_\allowbreak{}frame\_\allowbreak{}sst} (\texttt{Development.\allowbreak{}AbelianVarietyTheory.\allowbreak{}Proofs.\allowbreak{}InertiaEigenvalues.\allowbreak{}ExistsBreuilKisinFrameSemistableBoundedJumps}, lines 109--135)\\
\texttt{GlobalTheory.\allowbreak{}BKLim.\allowbreak{}Frame.\allowbreak{}isCr\_\allowbreak{}of\_\allowbreak{}wSeed} (\texttt{Development.\allowbreak{}AbelianVarietyTheory.\allowbreak{}Proofs.\allowbreak{}InertiaEigenvalues.\allowbreak{}DvdGalSub}, lines 55--111)\\
\texttt{GlobalTheory.\allowbreak{}BKLim.\allowbreak{}WDR.\allowbreak{}frame\_\allowbreak{}isDR} (\texttt{Development.\allowbreak{}AbelianVarietyTheory.\allowbreak{}Proofs.\allowbreak{}InertiaEigenvalues.\allowbreak{}SemistablePeriodFrameIsDeRham}, lines 18--41)\\
\texttt{GlobalTheory.\allowbreak{}BKLim.\allowbreak{}exists\_\allowbreak{}beta\_\allowbreak{}of\_\allowbreak{}disc\_\allowbreak{}cocycle} (\texttt{Development.\allowbreak{}AbelianVarietyTheory.\allowbreak{}Proofs.\allowbreak{}InertiaEigenvalues.\allowbreak{}ExistsBetaOfDiscCocycle}, lines 103--243)\par}

{\footnotesize\raggedright\noindent \ref{B9.35ab}.\quad \texttt{GlobalTheory.\allowbreak{}BKLim.\allowbreak{}beta\_\allowbreak{}bdd\_\allowbreak{}denominators} (\texttt{Development.\allowbreak{}AbelianVarietyTheory.\allowbreak{}Proofs.\allowbreak{}InertiaEigenvalues.\allowbreak{}BreuilKisinBetaBoundedDenominators}, lines 41--105)\\
\texttt{GlobalTheory.\allowbreak{}BKLim.\allowbreak{}exists\_\allowbreak{}isZ\_\allowbreak{}not\_\allowbreak{}root} (\texttt{Development.\allowbreak{}AbelianVarietyTheory.\allowbreak{}Proofs.\allowbreak{}InertiaEigenvalues.\allowbreak{}ExistsGaloisElemCyclotomicPeriodNotRootOfUnity}, lines 108--162)\\
\texttt{GlobalTheory.\allowbreak{}BKLim.\allowbreak{}beta\_\allowbreak{}step} (\texttt{Development.\allowbreak{}AbelianVarietyTheory.\allowbreak{}Proofs.\allowbreak{}InertiaEigenvalues.\allowbreak{}BreuilKisinBetaNextCoeffThetaMem}, lines 63--138)\\
\texttt{GlobalTheory.\allowbreak{}LiftDiv.\allowbreak{}dvd\_\allowbreak{}mod\_\allowbreak{}E} (\texttt{Development.\allowbreak{}AbelianVarietyTheory.\allowbreak{}Proofs.\allowbreak{}InertiaEigenvalues.\allowbreak{}ExistsPrimePowDvdModEisensteinOfTheta}, lines 14--47)\\
\texttt{GlobalTheory.\allowbreak{}BKLim.\allowbreak{}BetaDiv.\allowbreak{}isBeta\_\allowbreak{}div} (\texttt{Development.\allowbreak{}AbelianVarietyTheory.\allowbreak{}Proofs.\allowbreak{}InertiaEigenvalues.\allowbreak{}BreuilKisinIsBetaSubOfLpowSmul}, lines 14--73)\par}

{\footnotesize\raggedright\noindent \ref{B9.35a}.\quad \texttt{UpperTheory.\allowbreak{}Gal.\allowbreak{}ClosedPhiFromClosedTrace.\allowbreak{}coverJ\_\allowbreak{}holds} (\texttt{Development.\allowbreak{}AutomorphicArithmetic.\allowbreak{}Foundations.\allowbreak{}GaloisRepresentations.\allowbreak{}ClosedPhiFunctoriality.\allowbreak{}CoverJ}, lines 23--125)\\
\texttt{UpperTheory.\allowbreak{}Gal.\allowbreak{}ClosedPhiFromClosedTrace.\allowbreak{}frEx\_\allowbreak{}holds} (\texttt{Development.\allowbreak{}AutomorphicArithmetic.\allowbreak{}Foundations.\allowbreak{}GaloisRepresentations.\allowbreak{}ClosedPhiFunctoriality.\allowbreak{}FrExHolds}, lines 11--38)\\
\texttt{UpperTheory.\allowbreak{}Gal.\allowbreak{}ClosedPhiFromClosedTrace.\allowbreak{}opEx\_\allowbreak{}holds} (\texttt{Development.\allowbreak{}AutomorphicArithmetic.\allowbreak{}Foundations.\allowbreak{}GaloisRepresentations.\allowbreak{}ClosedPhiFunctoriality.\allowbreak{}OpExHolds}, lines 72--103)\\
\texttt{GlobalTheory.\allowbreak{}AcrisPow.\allowbreak{}main} (\texttt{Development.\allowbreak{}AbelianVarietyTheory.\allowbreak{}Proofs.\allowbreak{}InertiaEigenvalues.\allowbreak{}PrimePowDvdTiltUniformizerPowAcrisEventually}, lines 70--83)\\
\texttt{UpperTheory.\allowbreak{}Gal.\allowbreak{}ClosedPhiFromClosedTrace.\allowbreak{}secU\_\allowbreak{}holds} (\texttt{Development.\allowbreak{}AutomorphicArithmetic.\allowbreak{}Foundations.\allowbreak{}GaloisRepresentations.\allowbreak{}ClosedPhiFunctoriality.\allowbreak{}SecUHolds}, lines 12--86)\par}

{\footnotesize\raggedright\noindent \ref{B9.36}.\quad \texttt{GlobalTheory.\allowbreak{}CongPhi.\allowbreak{}ClosedPhi} (\texttt{Development.\allowbreak{}AbelianVarietyTheory.\allowbreak{}Definitions.\allowbreak{}NonGaloisDescent.\allowbreak{}JointApproxOfMem}, lines 79--94)\\
\texttt{GlobalTheory.\allowbreak{}LimitTop.\allowbreak{}closedAtL} (\texttt{Development.\allowbreak{}AbelianVarietyTheory.\allowbreak{}Proofs.\allowbreak{}SatakeParameters.\allowbreak{}ClosedAtL}, lines 25--117)\\
\texttt{UpperTheory.\allowbreak{}Gal.\allowbreak{}ClosedPhiFromClosedTrace.\allowbreak{}trStep\_\allowbreak{}holds} (\texttt{Development.\allowbreak{}AutomorphicArithmetic.\allowbreak{}Foundations.\allowbreak{}GaloisRepresentations.\allowbreak{}ClosedPhiFunctoriality.\allowbreak{}TrStep}, lines 83--162)\\
\texttt{UpperTheory.\allowbreak{}Gal.\allowbreak{}ClosedPhiFromClosedTrace.\allowbreak{}trLim\_\allowbreak{}holds} (\texttt{Development.\allowbreak{}AutomorphicArithmetic.\allowbreak{}Foundations.\allowbreak{}GaloisRepresentations.\allowbreak{}ClosedPhiFunctoriality.\allowbreak{}TrLimHolds}, lines 8--58)\\
\texttt{UpperTheory.\allowbreak{}Gal.\allowbreak{}ClosedPhiFromClosedTrace.\allowbreak{}congPhi\_\allowbreak{}holds} (\texttt{Development.\allowbreak{}AutomorphicArithmetic.\allowbreak{}Foundations.\allowbreak{}GaloisRepresentations.\allowbreak{}ClosedPhiFunctoriality.\allowbreak{}ClosedPhiHolds}, lines 44--45)\par}

{\footnotesize\raggedright\noindent \ref{B9.36c}.\quad \texttt{UpperTheory.\allowbreak{}ConvolutionTheory.\allowbreak{}D.\allowbreak{}muEx} (\texttt{Development.\allowbreak{}AutomorphicArithmetic.\allowbreak{}Foundations.\allowbreak{}Convolution.\allowbreak{}HermitianDescent.\allowbreak{}OffUnitsLeUnitIdl}, lines 357--402)\par}

{\footnotesize\raggedright\noindent \ref{B9.36b}.\quad \texttt{UpperTheory.\allowbreak{}V.\allowbreak{}Z.\allowbreak{}Tf.\allowbreak{}GlobalTraceComparison.\allowbreak{}Asm.\allowbreak{}rUEval} (\texttt{Development.\allowbreak{}AutomorphicArithmetic.\allowbreak{}AutomorphicVarieties.\allowbreak{}ModularityArgument.\allowbreak{}TraceFormula.\allowbreak{}GlobalTraceComparison.\allowbreak{}TraceAsymptotics.\allowbreak{}RUEval}, lines 197--331)\\
\texttt{UpperTheory.\allowbreak{}Zero.\allowbreak{}ConvolutionTheory.\allowbreak{}EB.\allowbreak{}rEndo\_\allowbreak{}of} (\texttt{Development.\allowbreak{}AutomorphicArithmetic.\allowbreak{}Foundations.\allowbreak{}Convolution.\allowbreak{}EndomorphismSlices.\allowbreak{}REndoOf}, lines 53--101)\\
\texttt{UpperTheory.\allowbreak{}Zero.\allowbreak{}ConvolutionTheory.\allowbreak{}EB.\allowbreak{}rStable} (\texttt{Development.\allowbreak{}AutomorphicArithmetic.\allowbreak{}Foundations.\allowbreak{}Convolution.\allowbreak{}EndomorphismSlices.\allowbreak{}FiniteAdelicEstimates.\allowbreak{}RStable}, lines 75--133)\\
\texttt{UpperTheory.\allowbreak{}V.\allowbreak{}Z.\allowbreak{}Tf.\allowbreak{}GlobalTraceComparison.\allowbreak{}Asm.\allowbreak{}ends} (\texttt{Development.\allowbreak{}AutomorphicArithmetic.\allowbreak{}AutomorphicVarieties.\allowbreak{}ModularityArgument.\allowbreak{}TraceFormula.\allowbreak{}GlobalTraceComparison.\allowbreak{}TraceAsymptotics.\allowbreak{}EvaluationCombinedTestFunction}, lines 322--347)\par}

{\footnotesize\raggedright\noindent \ref{B9.36a}.\quad \texttt{UpperTheory.\allowbreak{}V.\allowbreak{}Z.\allowbreak{}Tf.\allowbreak{}GaloisFromTraceFormula.\allowbreak{}asmTh\_\allowbreak{}of} (\texttt{Development.\allowbreak{}AutomorphicArithmetic.\allowbreak{}AutomorphicVarieties.\allowbreak{}ModularityArgument.\allowbreak{}TraceFormula.\allowbreak{}GaloisFromTraceFormula.\allowbreak{}AsmThOf}, lines 10--113)\\
\texttt{UpperTheory.\allowbreak{}Gal.\allowbreak{}Desc.\allowbreak{}weakAll} (\texttt{Development.\allowbreak{}AutomorphicArithmetic.\allowbreak{}Foundations.\allowbreak{}GaloisRepresentations.\allowbreak{}BaseChangeDescent.\allowbreak{}BcBackOffT}, lines 26--26)\par}

{\footnotesize\raggedright\noindent \ref{B9.37}.\quad \texttt{GlobalTheory.\allowbreak{}Cong.\allowbreak{}DescOffT} (\texttt{Development.\allowbreak{}AbelianVarietyTheory.\allowbreak{}Definitions.\allowbreak{}GaloisGeometry.\allowbreak{}MatchOfMatchOffTEmpty}, lines 32--37)\\
\texttt{UpperTheory.\allowbreak{}Gal.\allowbreak{}Desc.\allowbreak{}descOffT} (\texttt{Development.\allowbreak{}AutomorphicArithmetic.\allowbreak{}Foundations.\allowbreak{}GaloisRepresentations.\allowbreak{}BaseChangeDescent.\allowbreak{}BcBackOffT}, lines 31--32)\par}

{\footnotesize\raggedright\noindent \ref{B9.37a}.\quad \texttt{GlobalTheory.\allowbreak{}QuadLiftBridge.\allowbreak{}quadLift} (\texttt{Development.\allowbreak{}AbelianVarietyTheory.\allowbreak{}Proofs.\allowbreak{}SatakeParameters.\allowbreak{}QuadLift}, lines 14--24)\par}

{\footnotesize\raggedright\noindent \ref{B9.38}.\quad \texttt{UpperTheory.\allowbreak{}Gal.\allowbreak{}Int.\allowbreak{}Core.\allowbreak{}toCoreK} (\texttt{Development.\allowbreak{}AutomorphicArithmetic.\allowbreak{}Foundations.\allowbreak{}GaloisRepresentations.\allowbreak{}IntegralStructures.\allowbreak{}CrystallineDescent.\allowbreak{}ToCoreK}, lines 7--10)\\
\texttt{UpperTheory.\allowbreak{}Gal.\allowbreak{}Int.\allowbreak{}Inp.\allowbreak{}coreAll\_\allowbreak{}of\_\allowbreak{}core} (\texttt{Development.\allowbreak{}AutomorphicArithmetic.\allowbreak{}Foundations.\allowbreak{}GaloisRepresentations.\allowbreak{}IntegralStructures.\allowbreak{}CrystallineInputs.\allowbreak{}CoreAllOfRowStarCore}, lines 4--6)\\
\texttt{UpperTheory.\allowbreak{}Gal.\allowbreak{}Int.\allowbreak{}Core.\allowbreak{}patchOne\_\allowbreak{}holds} (\texttt{Development.\allowbreak{}AutomorphicArithmetic.\allowbreak{}Foundations.\allowbreak{}GaloisRepresentations.\allowbreak{}IntegralStructures.\allowbreak{}CrystallineDescent.\allowbreak{}PatchOne}, lines 83--128)\\
\texttt{UpperTheory.\allowbreak{}Gal.\allowbreak{}Int.\allowbreak{}Core.\allowbreak{}stepDown\_\allowbreak{}of} (\texttt{Development.\allowbreak{}AutomorphicArithmetic.\allowbreak{}Foundations.\allowbreak{}GaloisRepresentations.\allowbreak{}IntegralStructures.\allowbreak{}CrystallineDescent.\allowbreak{}JoinEndHolds}, lines 92--135)\\
\texttt{UpperTheory.\allowbreak{}Gal.\allowbreak{}Int.\allowbreak{}Inp.\allowbreak{}coreAll\_\allowbreak{}of\_\allowbreak{}core\_\allowbreak{}inp} (\texttt{Development.\allowbreak{}AutomorphicArithmetic.\allowbreak{}Foundations.\allowbreak{}GaloisRepresentations.\allowbreak{}IntegralStructures.\allowbreak{}CrystallineInputs.\allowbreak{}CoreAllOfRowStarCoreInp}, lines 538--551)\\
\texttt{UpperTheory.\allowbreak{}Gal.\allowbreak{}Int.\allowbreak{}Core.\allowbreak{}joinU\_\allowbreak{}of} (\texttt{Development.\allowbreak{}AutomorphicArithmetic.\allowbreak{}Foundations.\allowbreak{}GaloisRepresentations.\allowbreak{}IntegralStructures.\allowbreak{}CrystallineDescent.\allowbreak{}JoinEndHolds}, lines 19--62)\\
\texttt{GlobalTheory.\allowbreak{}GoodCsdTwist.\allowbreak{}exists\_\allowbreak{}goodCsd\_\allowbreak{}twist} (\texttt{Development.\allowbreak{}AbelianVarietyTheory.\allowbreak{}Proofs.\allowbreak{}SatakeParameters.\allowbreak{}ExistsGoodConjSelfDualTwistSplitUnramified}, lines 65--111)\\
\texttt{UpperTheory.\allowbreak{}Gal.\allowbreak{}Int.\allowbreak{}Core.\allowbreak{}joinUp\_\allowbreak{}of'} (\texttt{Development.\allowbreak{}AutomorphicArithmetic.\allowbreak{}Foundations.\allowbreak{}GaloisRepresentations.\allowbreak{}IntegralStructures.\allowbreak{}CrystallineDescent.\allowbreak{}BadFinite}, lines 46--93)\par}

{\footnotesize\raggedright\noindent \ref{B9.38a}.\quad \texttt{GalTw.\allowbreak{}GlobalTheory.\allowbreak{}RoadU.\allowbreak{}window} (\texttt{Development.\allowbreak{}AutomorphicArithmetic.\allowbreak{}Foundations.\allowbreak{}GaloisRepresentations.\allowbreak{}SphericalGaloisCompatibility.\allowbreak{}ToralDeRham.\allowbreak{}NormEqOneOfIsUnitarisable}, lines 100--119)\par}

{\footnotesize\raggedright\noindent \ref{B9.38e}.\quad \texttt{UpperTheory.\allowbreak{}V.\allowbreak{}Z.\allowbreak{}ImaginaryQuadraticEtaCharacter.\allowbreak{}gamLocX} (\texttt{Development.\allowbreak{}AutomorphicArithmetic.\allowbreak{}AutomorphicVarieties.\allowbreak{}ModularityArgument.\allowbreak{}LeviSmoothLocalization.\allowbreak{}PrincipalSeriesMatching.\allowbreak{}GamLocX}, lines 56--86)\par}

{\footnotesize\raggedright\noindent \ref{B9.38b}.\quad \texttt{GalTw.\allowbreak{}GlobalTheory.\allowbreak{}RoadU.\allowbreak{}exists\_\allowbreak{}sph\_\allowbreak{}lift\_\allowbreak{}of\_\allowbreak{}toral} (\texttt{Development.\allowbreak{}AutomorphicArithmetic.\allowbreak{}Foundations.\allowbreak{}GaloisRepresentations.\allowbreak{}SphericalGaloisCompatibility.\allowbreak{}ToralDeRham.\allowbreak{}FrobeniusCharpolyOfUnlinkedUnramifiedSupport}, lines 32--46)\\
\texttt{GalTw.\allowbreak{}GlobalTheory.\allowbreak{}RoadU.\allowbreak{}TorLiftC.\allowbreak{}exists\_\allowbreak{}cm\_\allowbreak{}lift\_\allowbreak{}fix\_\allowbreak{}charKer} (\texttt{Development.\allowbreak{}AutomorphicArithmetic.\allowbreak{}Foundations.\allowbreak{}GaloisRepresentations.\allowbreak{}SphericalGaloisCompatibility.\allowbreak{}ToralDeRham.\allowbreak{}WeaklyAttachedToralOfCmLift}, lines 29--92)\\
\texttt{UpperTheory.\allowbreak{}Gal.\allowbreak{}GamX.\allowbreak{}charSupp\_\allowbreak{}of} (\texttt{Development.\allowbreak{}AutomorphicArithmetic.\allowbreak{}Foundations.\allowbreak{}GaloisRepresentations.\allowbreak{}MixedUnlinkedExponents.\allowbreak{}CharLiftPrimeOf}, lines 93--115)\\
\texttt{GalTw.\allowbreak{}GlobalTheory.\allowbreak{}RoadU.\allowbreak{}TorLift.\allowbreak{}hasSat\_\allowbreak{}charPol\_\allowbreak{}above} (\texttt{Development.\allowbreak{}AutomorphicArithmetic.\allowbreak{}Foundations.\allowbreak{}GaloisRepresentations.\allowbreak{}SphericalGaloisCompatibility.\allowbreak{}ToralDeRham.\allowbreak{}FrobeniusCharpolyEqSatakeOfPlaceAbove}, lines 28--93)\par}

{\footnotesize\raggedright\noindent \ref{B9.38c}.\quad \texttt{UpperTheory.\allowbreak{}Gal.\allowbreak{}ToralGalois.\allowbreak{}galoisCompat\_\allowbreak{}unlinkedSupport\_\allowbreak{}of\_\allowbreak{}spherical} (\texttt{Development.\allowbreak{}AutomorphicArithmetic.\allowbreak{}Foundations.\allowbreak{}GaloisRepresentations.\allowbreak{}SphericalGaloisCompatibility.\allowbreak{}GaloisCompatUnlinkedSupportOfSpherical}, lines 10--11)\\
\texttt{GalTw.\allowbreak{}GlobalTheory.\allowbreak{}RoadU.\allowbreak{}toralGal\_\allowbreak{}of\_\allowbreak{}holds} (\texttt{Development.\allowbreak{}AutomorphicArithmetic.\allowbreak{}Foundations.\allowbreak{}GaloisRepresentations.\allowbreak{}SphericalGaloisCompatibility.\allowbreak{}ToralDeRham.\allowbreak{}ToralGalois2}, lines 22--57)\\
\texttt{GalTw.\allowbreak{}GlobalTheory.\allowbreak{}RoadU.\allowbreak{}toralAbove\_\allowbreak{}of\_\allowbreak{}holds} (\texttt{Development.\allowbreak{}AutomorphicArithmetic.\allowbreak{}Foundations.\allowbreak{}GaloisRepresentations.\allowbreak{}SphericalGaloisCompatibility.\allowbreak{}ToralDeRham.\allowbreak{}ToralAboveOfHolds}, lines 31--72)\\
\texttt{GalTw.\allowbreak{}GlobalTheory.\allowbreak{}DeRhamOfPlainTorus.\allowbreak{}deRham\_\allowbreak{}of\_\allowbreak{}plainTor} (\texttt{Development.\allowbreak{}AutomorphicArithmetic.\allowbreak{}Foundations.\allowbreak{}GaloisRepresentations.\allowbreak{}SphericalGaloisCompatibility.\allowbreak{}ToralDeRham.\allowbreak{}GaloisCompatUnlinkedSupportOfSpherical}, lines 27--68)\par}

{\footnotesize\raggedright\noindent \ref{B9.38d}.\quad \texttt{UpperTheory.\allowbreak{}Gal.\allowbreak{}ToralGalois.\allowbreak{}R2.\allowbreak{}toTor} (\texttt{Development.\allowbreak{}AutomorphicArithmetic.\allowbreak{}Foundations.\allowbreak{}GaloisRepresentations.\allowbreak{}SphericalGaloisCompatibility.\allowbreak{}PAdicSplitting.\allowbreak{}PToTor}, lines 61--99)\\
\texttt{UpperTheory.\allowbreak{}Gal.\allowbreak{}InertialTypesAlt.\allowbreak{}suppAny} (\texttt{Development.\allowbreak{}AutomorphicArithmetic.\allowbreak{}Foundations.\allowbreak{}GaloisRepresentations.\allowbreak{}RegularSemisimpleFibres.\allowbreak{}PSuppAny}, lines 30--102)\\
\texttt{UpperTheory.\allowbreak{}Gal.\allowbreak{}InertialTypesAlt.\allowbreak{}mixUnlinked} (\texttt{Development.\allowbreak{}AutomorphicArithmetic.\allowbreak{}Foundations.\allowbreak{}GaloisRepresentations.\allowbreak{}RegularSemisimpleFibres.\allowbreak{}PMixUnlinked}, lines 10--16)\par}

{\footnotesize\raggedright\noindent \ref{B9.39}.\quad \texttt{UpperTheory.\allowbreak{}TP.\allowbreak{}TLiftK} (\texttt{Development.\allowbreak{}AutomorphicArithmetic.\allowbreak{}Foundations.\allowbreak{}GaloisRepresentations.\allowbreak{}ThetaParseval.\allowbreak{}HalfIndexLattices.\allowbreak{}TLiftK}, lines 12--24)\\
\texttt{UpperTheory.\allowbreak{}TP.\allowbreak{}W.\allowbreak{}GoodK} (\texttt{Development.\allowbreak{}AutomorphicArithmetic.\allowbreak{}Foundations.\allowbreak{}GaloisRepresentations.\allowbreak{}ThetaParseval.\allowbreak{}HalfIndexLattices.\allowbreak{}SplitStab}, lines 18--25)\\
\texttt{UpperTheory.\allowbreak{}TP.\allowbreak{}LocPair} (\texttt{Development.\allowbreak{}AutomorphicArithmetic.\allowbreak{}Foundations.\allowbreak{}GaloisRepresentations.\allowbreak{}ThetaParseval.\allowbreak{}SchurStep}, lines 93--101)\\
\texttt{UpperTheory.\allowbreak{}TP.\allowbreak{}IsLiftWt} (\texttt{Development.\allowbreak{}AutomorphicArithmetic.\allowbreak{}Foundations.\allowbreak{}GaloisRepresentations.\allowbreak{}ThetaParseval.\allowbreak{}SchurStep}, lines 103--104)\\
\texttt{UpperTheory.\allowbreak{}TP.\allowbreak{}IsLift} (\texttt{Development.\allowbreak{}AutomorphicArithmetic.\allowbreak{}Foundations.\allowbreak{}GaloisRepresentations.\allowbreak{}ThetaParseval.\allowbreak{}SchurStep}, lines 106--112)\\
\texttt{UpperTheory.\allowbreak{}TP.\allowbreak{}K.\allowbreak{}tLiftK} (\texttt{Development.\allowbreak{}AutomorphicArithmetic.\allowbreak{}Foundations.\allowbreak{}GaloisRepresentations.\allowbreak{}ThetaParseval.\allowbreak{}ThetaKernels.\allowbreak{}TLiftK}, lines 4--6)\par}

{\footnotesize\raggedright\noindent \ref{B9.40}.\quad \texttt{UpperTheory.\allowbreak{}TP.\allowbreak{}Rallis\_\allowbreak{}k} (\texttt{Development.\allowbreak{}AutomorphicArithmetic.\allowbreak{}Foundations.\allowbreak{}GaloisRepresentations.\allowbreak{}ThetaParseval.\allowbreak{}ThetaKernels.\allowbreak{}GalTPDefs}, lines 61--68)\\
\texttt{UpperTheory.\allowbreak{}TP.\allowbreak{}NonVan\_\allowbreak{}k} (\texttt{Development.\allowbreak{}AutomorphicArithmetic.\allowbreak{}Foundations.\allowbreak{}GaloisRepresentations.\allowbreak{}ThetaParseval.\allowbreak{}ThetaKernels.\allowbreak{}GalTPDefs}, lines 114--121)\\
\texttt{UpperTheory.\allowbreak{}TP.\allowbreak{}rallis\_\allowbreak{}k} (\texttt{Development.\allowbreak{}AutomorphicArithmetic.\allowbreak{}Foundations.\allowbreak{}GaloisRepresentations.\allowbreak{}ThetaParseval.\allowbreak{}ThetaKernels.\allowbreak{}TPGRallis}, lines 99--225)\\
\texttt{UpperTheory.\allowbreak{}TP.\allowbreak{}nonVan\_\allowbreak{}k} (\texttt{Development.\allowbreak{}AutomorphicArithmetic.\allowbreak{}Foundations.\allowbreak{}GaloisRepresentations.\allowbreak{}ThetaParseval.\allowbreak{}ThetaKernels.\allowbreak{}TPGNonVan}, lines 26--154)\\
\texttt{UpperTheory.\allowbreak{}TP.\allowbreak{}IsBump} (\texttt{Development.\allowbreak{}AutomorphicArithmetic.\allowbreak{}Foundations.\allowbreak{}GaloisRepresentations.\allowbreak{}ThetaParseval.\allowbreak{}SchurStep}, lines 44--50)\par}

{\footnotesize\raggedright\noindent \ref{B9.40a}.\quad \texttt{UpperTheory.\allowbreak{}Gal.\allowbreak{}TP.\allowbreak{}Loc.\allowbreak{}liftIso\_\allowbreak{}of} (\texttt{Development.\allowbreak{}AutomorphicArithmetic.\allowbreak{}Foundations.\allowbreak{}GaloisRepresentations.\allowbreak{}ThetaParseval.\allowbreak{}SatakeLifts.\allowbreak{}IsUnitaryCharInv}, lines 22--28)\\
\texttt{UpperTheory.\allowbreak{}Gal.\allowbreak{}TP.\allowbreak{}Loc.\allowbreak{}pairSat\_\allowbreak{}holds} (\texttt{Development.\allowbreak{}AutomorphicArithmetic.\allowbreak{}Foundations.\allowbreak{}GaloisRepresentations.\allowbreak{}ThetaParseval.\allowbreak{}SatakeLifts.\allowbreak{}LiftSatUHolds}, lines 31--32)\par}

{\footnotesize\raggedright\noindent \ref{B9.41}.\quad \texttt{UpperTheory.\allowbreak{}Gal.\allowbreak{}TP.\allowbreak{}twoUpOK\_\allowbreak{}of\_\allowbreak{}steps} (\texttt{Development.\allowbreak{}AutomorphicArithmetic.\allowbreak{}Foundations.\allowbreak{}GaloisRepresentations.\allowbreak{}ThetaParseval.\allowbreak{}TwoUpOKOfSteps}, lines 13--32)\\
\texttt{UpperTheory.\allowbreak{}Gal.\allowbreak{}TP.\allowbreak{}TwoUpOK} (\texttt{Development.\allowbreak{}AutomorphicArithmetic.\allowbreak{}Foundations.\allowbreak{}GaloisRepresentations.\allowbreak{}ThetaParseval.\allowbreak{}AllCoreOf}, lines 18--22)\\
\texttt{UpperTheory.\allowbreak{}Gal.\allowbreak{}TP.\allowbreak{}InertOK} (\texttt{Development.\allowbreak{}AutomorphicArithmetic.\allowbreak{}Foundations.\allowbreak{}GaloisRepresentations.\allowbreak{}ThetaParseval.\allowbreak{}AllCoreOf}, lines 12--16)\\
\texttt{UpperTheory.\allowbreak{}Gal.\allowbreak{}TP.\allowbreak{}Sub.\allowbreak{}IsUp} (\texttt{Development.\allowbreak{}AutomorphicArithmetic.\allowbreak{}Foundations.\allowbreak{}GaloisRepresentations.\allowbreak{}ThetaParseval.\allowbreak{}SplitLineSubtraction.\allowbreak{}RowDefOfOdd}, lines 85--95)\\
\texttt{UpperTheory.\allowbreak{}Gal.\allowbreak{}TP.\allowbreak{}twoUpOK\_\allowbreak{}of\_\allowbreak{}k} (\texttt{Development.\allowbreak{}AutomorphicArithmetic.\allowbreak{}Foundations.\allowbreak{}GaloisRepresentations.\allowbreak{}ThetaParseval.\allowbreak{}TwoUpOKOfSteps}, lines 33--33)\\
\texttt{UpperTheory.\allowbreak{}Gal.\allowbreak{}TP.\allowbreak{}Chr.\allowbreak{}chrPairs} (\texttt{Development.\allowbreak{}AutomorphicArithmetic.\allowbreak{}Foundations.\allowbreak{}GaloisRepresentations.\allowbreak{}ThetaParseval.\allowbreak{}IdeleNormCharacters.\allowbreak{}ChrPairsPf}, lines 26--92)\par}

{\footnotesize\raggedright\noindent \ref{B9.42}.\quad \texttt{UpperTheory.\allowbreak{}Gal.\allowbreak{}TP.\allowbreak{}frobIdQ} (\texttt{Development.\allowbreak{}AutomorphicArithmetic.\allowbreak{}Foundations.\allowbreak{}GaloisRepresentations.\allowbreak{}ThetaParseval.\allowbreak{}CoreAtQIffRowQ}, lines 52--87)\par}

{\footnotesize\raggedright\noindent \ref{B9.43}.\quad \texttt{UpperTheory.\allowbreak{}Gal.\allowbreak{}TP.\allowbreak{}Sub.\allowbreak{}peelAll\_\allowbreak{}holds} (\texttt{Development.\allowbreak{}AutomorphicArithmetic.\allowbreak{}Foundations.\allowbreak{}GaloisRepresentations.\allowbreak{}ThetaParseval.\allowbreak{}SplitLineSubtraction.\allowbreak{}PPeelAll}, lines 17--52)\\
\texttt{UpperTheory.\allowbreak{}Gal.\allowbreak{}TP.\allowbreak{}Sub.\allowbreak{}eigTransfer\_\allowbreak{}holds} (\texttt{Development.\allowbreak{}AutomorphicArithmetic.\allowbreak{}Foundations.\allowbreak{}GaloisRepresentations.\allowbreak{}ThetaParseval.\allowbreak{}SplitLineSubtraction.\allowbreak{}PEigTransfer}, lines 13--35)\\
\texttt{UpperTheory.\allowbreak{}Gal.\allowbreak{}TP.\allowbreak{}Sub.\allowbreak{}splitLine\_\allowbreak{}holds} (\texttt{Development.\allowbreak{}AutomorphicArithmetic.\allowbreak{}Foundations.\allowbreak{}GaloisRepresentations.\allowbreak{}ThetaParseval.\allowbreak{}SplitLineSubtraction.\allowbreak{}PSplitLine}, lines 100--133)\par}

{\footnotesize\raggedright\noindent \ref{B9.44}.\quad \texttt{UpperTheory.\allowbreak{}Gal.\allowbreak{}TP.\allowbreak{}atPlaceQ} (\texttt{Development.\allowbreak{}AutomorphicArithmetic.\allowbreak{}Foundations.\allowbreak{}GaloisRepresentations.\allowbreak{}ThetaParseval.\allowbreak{}CoreAtQIffRowQ}, lines 103--175)\par}

{\footnotesize\raggedright\noindent \ref{B9.45}.\quad \texttt{UpperTheory.\allowbreak{}Gal.\allowbreak{}TP.\allowbreak{}subtractQ\_\allowbreak{}of} (\texttt{Development.\allowbreak{}AutomorphicArithmetic.\allowbreak{}Foundations.\allowbreak{}GaloisRepresentations.\allowbreak{}ThetaParseval.\allowbreak{}CoreAtQIffRowQ}, lines 200--210)\\
\texttt{UpperTheory.\allowbreak{}Gal.\allowbreak{}TP.\allowbreak{}weakIsoQ} (\texttt{Development.\allowbreak{}AutomorphicArithmetic.\allowbreak{}Foundations.\allowbreak{}GaloisRepresentations.\allowbreak{}ThetaParseval.\allowbreak{}CoreAtQIffRowQ}, lines 88--102)\\
\texttt{UpperTheory.\allowbreak{}Gal.\allowbreak{}Int.\allowbreak{}Core.\allowbreak{}subtractQ\_\allowbreak{}holds} (\texttt{Development.\allowbreak{}AutomorphicArithmetic.\allowbreak{}Foundations.\allowbreak{}GaloisRepresentations.\allowbreak{}IntegralStructures.\allowbreak{}CrystallineDescent.\allowbreak{}LimOddOfLimFinOdd}, lines 31--34)\par}

\section*{Chapter~\ref{ch:weak-local-langlands}: A weak local Langlands correspondence}

{\footnotesize\raggedright\noindent \ref{B9.46}.\quad \texttt{UpperTheory.\allowbreak{}Zero.\allowbreak{}LLC.\allowbreak{}recTExists} (\texttt{Development.\allowbreak{}AutomorphicArithmetic.\allowbreak{}Foundations.\allowbreak{}LocalLanglands.\allowbreak{}LocalCentralUnits.\allowbreak{}RecTExists}, lines 20--66)\par}

{\footnotesize\raggedright\noindent \ref{B9.47}.\quad \texttt{UpperTheory.\allowbreak{}Zero.\allowbreak{}LLC.\allowbreak{}LawsT} (\texttt{Development.\allowbreak{}AutomorphicArithmetic.\allowbreak{}Foundations.\allowbreak{}LocalLanglands.\allowbreak{}IsRecTOfIsRec}, lines 55--72)\\
\texttt{UpperTheory.\allowbreak{}Zero.\allowbreak{}LLC.\allowbreak{}IsRecT} (\texttt{Development.\allowbreak{}AutomorphicArithmetic.\allowbreak{}Foundations.\allowbreak{}LocalLanglands.\allowbreak{}IsRecTOfIsRec}, lines 76--76)\par}

{\footnotesize\raggedright\noindent \ref{B9.48}.\quad \texttt{UpperTheory.\allowbreak{}Zero.\allowbreak{}LLC.\allowbreak{}GoodCusp} (\texttt{Development.\allowbreak{}AutomorphicArithmetic.\allowbreak{}Foundations.\allowbreak{}LocalLanglands.\allowbreak{}IsRecTOfIsRec}, lines 20--22)\\
\texttt{UpperTheory.\allowbreak{}Zero.\allowbreak{}LLC.\allowbreak{}Good} (\texttt{Development.\allowbreak{}AutomorphicArithmetic.\allowbreak{}Foundations.\allowbreak{}LocalLanglands.\allowbreak{}IsRecTOfIsRec}, lines 24--24)\par}

{\footnotesize\raggedright\noindent \ref{B9.49}.\quad \texttt{UpperTheory.\allowbreak{}Zero.\allowbreak{}LLC.\allowbreak{}TypeLaw} (\texttt{Development.\allowbreak{}AutomorphicArithmetic.\allowbreak{}Foundations.\allowbreak{}LocalLanglands.\allowbreak{}IsRecTOfIsRec}, lines 35--53)\par}

{\footnotesize\raggedright\noindent \ref{B9.50}.\quad \texttt{UpperTheory.\allowbreak{}Zero.\allowbreak{}LLC.\allowbreak{}artEx} (\texttt{Development.\allowbreak{}AutomorphicArithmetic.\allowbreak{}Foundations.\allowbreak{}LocalLanglands.\allowbreak{}LocalCentralUnits.\allowbreak{}ArtEx}, lines 7--15)\\
\texttt{GlobalTheory.\allowbreak{}LocLanglands.\allowbreak{}IsLocalArtin} (\texttt{Development.\allowbreak{}AbelianVarietyTheory.\allowbreak{}PublishedTheorems.\allowbreak{}AutomorphicForms₂.\allowbreak{}HasExponentZeroIff}, lines 145--152)\par}

{\footnotesize\raggedright\noindent \ref{B9.51}.\quad \texttt{UpperTheory.\allowbreak{}Zero.\allowbreak{}LLC.\allowbreak{}parEx} (\texttt{Development.\allowbreak{}AutomorphicArithmetic.\allowbreak{}Foundations.\allowbreak{}LocalLanglands.\allowbreak{}LocalCentralUnits.\allowbreak{}CentralUnit}, lines 39--200)\\
\texttt{UpperTheory.\allowbreak{}Zero.\allowbreak{}LLC.\allowbreak{}parUniq} (\texttt{Development.\allowbreak{}AutomorphicArithmetic.\allowbreak{}Foundations.\allowbreak{}LocalLanglands.\allowbreak{}LocalCentralUnits.\allowbreak{}ParUniq}, lines 15--54)\\
\texttt{UpperTheory.\allowbreak{}Zero.\allowbreak{}LLC.\allowbreak{}goodMap} (\texttt{Development.\allowbreak{}AutomorphicArithmetic.\allowbreak{}Foundations.\allowbreak{}LocalLanglands.\allowbreak{}LocalCentralUnits.\allowbreak{}GoodMap}, lines 144--164)\\
\texttt{UpperTheory.\allowbreak{}Zero.\allowbreak{}LLC.\allowbreak{}twTame} (\texttt{Development.\allowbreak{}AutomorphicArithmetic.\allowbreak{}Foundations.\allowbreak{}LocalLanglands.\allowbreak{}LocalCentralUnits.\allowbreak{}TwTame}, lines 190--194)\par}

{\footnotesize\raggedright\noindent \ref{B9.52}.\quad \texttt{UpperTheory.\allowbreak{}Zero.\allowbreak{}LLC.\allowbreak{}junkEx} (\texttt{Development.\allowbreak{}AutomorphicArithmetic.\allowbreak{}Foundations.\allowbreak{}LocalLanglands.\allowbreak{}LocalCentralUnits.\allowbreak{}JunkEx}, lines 257--262)\\
\texttt{UpperTheory.\allowbreak{}tameInertiaCharacter\_\allowbreak{}frobenius\_\allowbreak{}conj\_\allowbreak{}eq\_\allowbreak{}pow} (\texttt{Development.\allowbreak{}AutomorphicArithmetic.\allowbreak{}WhittakerHecke.\allowbreak{}TameInertiaCharacterFrobeniusConjEqPow}, lines 111--177)\par}

{\footnotesize\raggedright\noindent \ref{B9.53}.\quad \texttt{UpperTheory.\allowbreak{}Zero.\allowbreak{}LLC.\allowbreak{}cuspEx} (\texttt{Development.\allowbreak{}AutomorphicArithmetic.\allowbreak{}Foundations.\allowbreak{}LocalLanglands.\allowbreak{}LocalCentralUnits.\allowbreak{}CuspEx}, lines 168--266)\\
\texttt{UpperTheory.\allowbreak{}Zero.\allowbreak{}LLC.\allowbreak{}expo} (\texttt{Development.\allowbreak{}AutomorphicArithmetic.\allowbreak{}Foundations.\allowbreak{}LocalLanglands.\allowbreak{}LocalCentralUnits.\allowbreak{}ExponentMulAbsPow}, lines 63--96)\par}

{\footnotesize\raggedright\noindent \ref{B9.54}.\quad \texttt{UpperTheory.\allowbreak{}Zero.\allowbreak{}LLC.\allowbreak{}Pr.\allowbreak{}exh} (\texttt{Development.\allowbreak{}AutomorphicArithmetic.\allowbreak{}Foundations.\allowbreak{}LocalLanglands.\allowbreak{}Multi}, lines 11--11)\\
\texttt{UpperTheory.\allowbreak{}Zero.\allowbreak{}LLC.\allowbreak{}Pr.\allowbreak{}multi} (\texttt{Development.\allowbreak{}AutomorphicArithmetic.\allowbreak{}Foundations.\allowbreak{}LocalLanglands.\allowbreak{}Multi}, lines 12--12)\\
\texttt{GlobalTheory.\allowbreak{}LocLanglands.\allowbreak{}segSum} (\texttt{Development.\allowbreak{}AbelianVarietyTheory.\allowbreak{}Definitions.\allowbreak{}NonGaloisDescent.\allowbreak{}Langlands}, lines 17--18)\\
\texttt{GlobalTheory.\allowbreak{}LocLanglands.\allowbreak{}exists\_\allowbreak{}isLExt} (\texttt{Development.\allowbreak{}AbelianVarietyTheory.\allowbreak{}Proofs.\allowbreak{}InertiaEigenvalues.\allowbreak{}LanglandsExtensionUniqueOnCuspidal}, lines 34--56)\par}

{\footnotesize\raggedright\noindent \ref{B9.55}.\quad \texttt{UpperTheory.\allowbreak{}Zero.\allowbreak{}LLC.\allowbreak{}recTUnique} (\texttt{Development.\allowbreak{}AutomorphicArithmetic.\allowbreak{}Foundations.\allowbreak{}LocalLanglands.\allowbreak{}LocalCentralUnits.\allowbreak{}RecTUnique}, lines 28--53)\par}

{\footnotesize\raggedright\noindent \ref{B9.55a}.\quad \texttt{UpperTheory.\allowbreak{}Gal.\allowbreak{}Int.\allowbreak{}Row.\allowbreak{}plainTor\_\allowbreak{}of\_\allowbreak{}tor} (\texttt{Development.\allowbreak{}AutomorphicArithmetic.\allowbreak{}Foundations.\allowbreak{}GaloisRepresentations.\allowbreak{}IntegralStructures.\allowbreak{}InertCharpoly.\allowbreak{}CharPolOfInert}, lines 97--106)\\
\texttt{UpperTheory.\allowbreak{}Gal.\allowbreak{}Int.\allowbreak{}Row.\allowbreak{}dz\_\allowbreak{}of\_\allowbreak{}cusp} (\texttt{Development.\allowbreak{}AutomorphicArithmetic.\allowbreak{}Foundations.\allowbreak{}GaloisRepresentations.\allowbreak{}IntegralStructures.\allowbreak{}InertCharpoly.\allowbreak{}CharPolOfInert}, lines 114--139)\\
\texttt{UpperTheory.\allowbreak{}Gal.\allowbreak{}Int.\allowbreak{}Row.\allowbreak{}mixed\_\allowbreak{}of\_\allowbreak{}tame} (\texttt{Development.\allowbreak{}AutomorphicArithmetic.\allowbreak{}Foundations.\allowbreak{}GaloisRepresentations.\allowbreak{}IntegralStructures.\allowbreak{}InertCharpoly.\allowbreak{}CharPolOfInert}, lines 151--181)\par}

\section*{Chapter~\ref{ch:p1c5}: The window criterion in the archive}

{\footnotesize\raggedright\noindent \ref{B1.47}.\quad \texttt{Deduction.\allowbreak{}landau\_\allowbreak{}window\_\allowbreak{}key} (\texttt{Development.\allowbreak{}AutomorphicArithmetic.\allowbreak{}DedekindLFunctions.\allowbreak{}LandauTheorems.\allowbreak{}landauWindow}, lines 40--74)\par}

{\footnotesize\raggedright\noindent \ref{B1.48}.\quad \texttt{Deduction.\allowbreak{}iteratedDeriv\_\allowbreak{}laplace\_\allowbreak{}hasDerivAt} (\texttt{Development.\allowbreak{}AutomorphicArithmetic.\allowbreak{}DedekindLFunctions.\allowbreak{}LandauTheorems.\allowbreak{}IteratedDerivLaplace}, lines 48--83)\\
\texttt{Deduction.\allowbreak{}iteratedDeriv\_\allowbreak{}laplace} (\texttt{Development.\allowbreak{}AutomorphicArithmetic.\allowbreak{}DedekindLFunctions.\allowbreak{}LandauTheorems.\allowbreak{}IteratedDerivLaplace}, lines 84--106)\par}

{\footnotesize\raggedright\noindent \ref{B1.49}.\quad \texttt{Deduction.\allowbreak{}landau\_\allowbreak{}coeff} (\texttt{Development.\allowbreak{}AutomorphicArithmetic.\allowbreak{}DedekindLFunctions.\allowbreak{}LandauTheorems.\allowbreak{}LandauCoeff}, lines 10--29)\par}

{\footnotesize\raggedright\noindent \ref{B1.50}.\quad \texttt{Deduction.\allowbreak{}kernel\_\allowbreak{}integral} (\texttt{Development.\allowbreak{}AutomorphicArithmetic.\allowbreak{}DedekindLFunctions.\allowbreak{}LandauTheorems.\allowbreak{}kernelIntegral}, lines 9--21)\\
\texttt{Deduction.\allowbreak{}landau\_\allowbreak{}kernel\_\allowbreak{}max} (\texttt{Development.\allowbreak{}AutomorphicArithmetic.\allowbreak{}DedekindLFunctions.\allowbreak{}LandauTheorems.\allowbreak{}LandauKernelMax}, lines 9--20)\\
\texttt{Deduction.\allowbreak{}landau\_\allowbreak{}tails\_\allowbreak{}integrable} (\texttt{Development.\allowbreak{}AutomorphicArithmetic.\allowbreak{}DedekindLFunctions.\allowbreak{}LandauTheorems.\allowbreak{}landauTails}, lines 10--17)\par}

{\footnotesize\raggedright\noindent \ref{B1.51}.\quad \texttt{Deduction.\allowbreak{}landau\_\allowbreak{}tails\_\allowbreak{}chernoff} (\texttt{Development.\allowbreak{}AutomorphicArithmetic.\allowbreak{}DedekindLFunctions.\allowbreak{}LandauTheorems.\allowbreak{}landauTails}, lines 18--61)\\
\texttt{Deduction.\allowbreak{}landau\_\allowbreak{}tails\_\allowbreak{}pow} (\texttt{Development.\allowbreak{}AutomorphicArithmetic.\allowbreak{}DedekindLFunctions.\allowbreak{}LandauTheorems.\allowbreak{}landauTails}, lines 62--79)\par}

{\footnotesize\raggedright\noindent \ref{B1.52}.\quad \texttt{Deduction.\allowbreak{}landau\_\allowbreak{}tails} (\texttt{Development.\allowbreak{}AutomorphicArithmetic.\allowbreak{}DedekindLFunctions.\allowbreak{}LandauTheorems.\allowbreak{}landauTails}, lines 80--132)\par}

{\footnotesize\raggedright\noindent \ref{B1.53}.\quad \texttt{Deduction.\allowbreak{}landau\_\allowbreak{}window\_\allowbreak{}ineq} (\texttt{Development.\allowbreak{}AutomorphicArithmetic.\allowbreak{}DedekindLFunctions.\allowbreak{}LandauTheorems.\allowbreak{}landauWindowIneq}, lines 27--112)\\
\texttt{Deduction.\allowbreak{}landau\_\allowbreak{}window\_\allowbreak{}ineq\_\allowbreak{}int} (\texttt{Development.\allowbreak{}AutomorphicArithmetic.\allowbreak{}DedekindLFunctions.\allowbreak{}LandauTheorems.\allowbreak{}landauWindowIneq}, lines 10--26)\par}

{\footnotesize\raggedright\noindent \ref{B1.56}.\quad \texttt{Deduction.\allowbreak{}landau\_\allowbreak{}dichotomy} (\texttt{Development.\allowbreak{}AutomorphicArithmetic.\allowbreak{}DedekindLFunctions.\allowbreak{}LandauTheorems.\allowbreak{}LandauDichotomy}, lines 24--88)\par}

{\footnotesize\raggedright\noindent \ref{B1.56}.\quad \texttt{Deduction.\allowbreak{}eq\_\allowbreak{}ofReal\_\allowbreak{}sub\_\allowbreak{}of\_\allowbreak{}dist\_\allowbreak{}le\_\allowbreak{}of\_\allowbreak{}le\_\allowbreak{}sub\_\allowbreak{}re} (\texttt{Development.\allowbreak{}AutomorphicArithmetic.\allowbreak{}DedekindLFunctions.\allowbreak{}LandauTheorems.\allowbreak{}LandauDichotomy}, lines 10--23)\par}

{\footnotesize\raggedright\noindent \ref{B1.57}.\quad \texttt{Deduction.\allowbreak{}landau\_\allowbreak{}nearest} (\texttt{Development.\allowbreak{}AutomorphicArithmetic.\allowbreak{}DedekindLFunctions.\allowbreak{}LandauTheorems.\allowbreak{}LandauNearest}, lines 36--110)\par}

{\footnotesize\raggedright\noindent \ref{B1.57}.\quad \texttt{Deduction.\allowbreak{}landau\_\allowbreak{}nearest\_\allowbreak{}equi} (\texttt{Development.\allowbreak{}AutomorphicArithmetic.\allowbreak{}DedekindLFunctions.\allowbreak{}LandauTheorems.\allowbreak{}LandauNearest}, lines 10--35)\par}

{\footnotesize\raggedright\noindent \ref{B1.58}.\quad \texttt{Deduction.\allowbreak{}landau\_\allowbreak{}cauchy} (\texttt{Development.\allowbreak{}AutomorphicArithmetic.\allowbreak{}DedekindLFunctions.\allowbreak{}LandauTheorems.\allowbreak{}LandauCauchy}, lines 9--34)\par}

{\footnotesize\raggedright\noindent \ref{B1.59}.\quad \texttt{Deduction.\allowbreak{}landau\_\allowbreak{}pole\_\allowbreak{}asym} (\texttt{Development.\allowbreak{}AutomorphicArithmetic.\allowbreak{}DedekindLFunctions.\allowbreak{}LandauTheorems.\allowbreak{}LandauPoleAsym}, lines 19--80)\par}

{\footnotesize\raggedright\noindent \ref{B1.59}.\quad \texttt{Deduction.\allowbreak{}landau\_\allowbreak{}pole\_\allowbreak{}asym\_\allowbreak{}deriv} (\texttt{Development.\allowbreak{}AutomorphicArithmetic.\allowbreak{}DedekindLFunctions.\allowbreak{}LandauTheorems.\allowbreak{}LandauPoleAsym}, lines 10--18)\par}

{\footnotesize\raggedright\noindent \ref{B1.60}.\quad \texttt{Deduction.\allowbreak{}landau\_\allowbreak{}growth\_\allowbreak{}contra} (\texttt{Development.\allowbreak{}AutomorphicArithmetic.\allowbreak{}DedekindLFunctions.\allowbreak{}LandauTheorems.\allowbreak{}LandauGrowthContra}, lines 9--41)\\
\texttt{Deduction.\allowbreak{}landau\_\allowbreak{}window\_\allowbreak{}contra} (\texttt{Development.\allowbreak{}AutomorphicArithmetic.\allowbreak{}DedekindLFunctions.\allowbreak{}LandauTheorems.\allowbreak{}landauWindow}, lines 20--39)\par}

{\footnotesize\raggedright\noindent \ref{B1.61}.\quad \texttt{Deduction.\allowbreak{}landau\_\allowbreak{}window} (\texttt{Development.\allowbreak{}AutomorphicArithmetic.\allowbreak{}DedekindLFunctions.\allowbreak{}LandauTheorems.\allowbreak{}landauWindow}, lines 92--133)\par}

{\footnotesize\raggedright\noindent \ref{B1.62}.\quad \texttt{Deduction.\allowbreak{}landau\_\allowbreak{}window} (\texttt{Development.\allowbreak{}AutomorphicArithmetic.\allowbreak{}DedekindLFunctions.\allowbreak{}LandauTheorems.\allowbreak{}landauWindow}, lines 134--166)\par}

\begin{thebibliography}{999}\small

\bibitem{sutherland-2026} A.~V.~Sutherland, \emph{Verification files for the Lean proof of the Riemann Hypothesis}, \url{https://math.mit.edu/~drew/grh/}, 2026.

\bibitem{leanfro-2026} Lean and Mathlib communities, \emph{Lean verification of the Riemann Hypothesis}, joint statement endorsed by the Lean FRO and the Mathlib Initiative, \url{https://leanprover-community.github.io/blog/posts/riemann-hypothesis/}, 21 November 2026.

\bibitem{alpoge-2026} L.~Alpöge, remark on the proof and Lean archive, \url{https://alpo.ge/remark.pdf}, 20 November 2026.

\bibitem{apostol-1976} T.~M.~Apostol, \emph{Introduction to Analytic Number Theory}, Undergraduate Texts in Mathematics, Springer, New York, 1976, Chapter~12.

\bibitem{olver-2010} F.~W.~J.~Olver, D.~W.~Lozier, R.~F.~Boisvert and C.~W.~Clark (eds.), \emph{NIST Handbook of Mathematical Functions}, Cambridge University Press, 2010; online companion, DLMF, \S\S25.4, 25.15 and 5.2. % [Zeta reflection](https://dlmf.nist.gov/25.4), [Dirichlet L-functions](https://dlmf.nist.gov/25.15), [Gamma function](https://dlmf.nist.gov/5.2).

\bibitem{howe-1989} R.~Howe, Transcending classical invariant theory, \emph{Journal of the American Mathematical Society} 2 (1989), 535--552. \url{https://doi.org/10.1090/S0894-0347-1989-0985172-6}.

\bibitem{shimura-1971} G.~Shimura, \emph{Introduction to the Arithmetic Theory of Automorphic Functions}, Publications of the Mathematical Society of Japan 11, Iwanami Shoten and Princeton University Press, 1971.

\bibitem{chapdelaine-2009} H.~Chapdelaine, Functional equation for partial zeta functions twisted by additive characters, \emph{Acta Arithmetica} 136 (2009), 213--228. \url{https://www.impan.pl/en/publishing-house/journals-and-series/acta-arithmetica/all/136/3/82809/functional-equation-for-partial-zeta-functions-twisted-by-additive-characters}.

\bibitem{serre-1979} J.-P.~Serre, \emph{Local Fields}, Graduate Texts in Mathematics 67, Springer, New York, 1979.

\bibitem{klingen-1962} H.~Klingen, \"Uber die Werte der Dedekindschen Zetafunktion, \emph{Mathematische Annalen} 145 (1962), 265--272.

\bibitem{deligne-ribet-1980} P.~Deligne and K.~A.~Ribet, Values of abelian $L$-functions at negative integers over totally real fields, \emph{Inventiones Mathematicae} 59 (1980), 227--286.

\bibitem{tate-1967} J.~T.~Tate, Fourier analysis in number fields and Hecke's zeta-functions, in \emph{Algebraic Number Theory}, J.~W.~S.~Cassels and A.~Fr\"ohlich (eds.), Academic Press, London, 1967, 305--347.

\bibitem{jacquet-shalika-1981} H.~Jacquet and J.~A.~Shalika, On Euler products and the classification of automorphic representations I and II, \emph{American Journal of Mathematics} 103 (1981), 499--558 and 777--815.

\bibitem{bernstein-zelevinsky-1977} I.~N.~Bernstein and A.~V.~Zelevinsky, Induced representations of reductive $p$-adic groups I, \emph{Annales Scientifiques de l'\'Ecole Normale Sup\'erieure} (4) 10 (1977), 441--472.

\bibitem{bernstein-1984} J.~Bernstein, $P$-invariant distributions on $\mathrm{GL}(N)$ and the classification of unitary representations of $\mathrm{GL}(N)$ (non-Archimedean case), in \emph{Lie Group Representations II}, Lecture Notes in Mathematics 1041, Springer, Berlin, 1984, 50--102.

\bibitem{zelevinsky-1980} A.~V.~Zelevinsky, Induced representations of reductive $p$-adic groups II. On irreducible representations of $\mathrm{GL}(n)$, \emph{Annales Scientifiques de l'\'Ecole Normale Sup\'erieure} (4) 13 (1980), 165--210. \url{https://numdam.org/articles/10.24033/asens.1379/}.

\bibitem{dixmier-malliavin-1978} J.~Dixmier and P.~Malliavin, Factorisations de fonctions et de vecteurs ind\'efiniment diff\'erentiables, \emph{Bulletin des Sciences Math\'ematiques} (2) 102 (1978), 307--330.

\bibitem{flath-1979} D.~Flath, Decomposition of representations into tensor products, in \emph{Automorphic Forms, Representations and $L$-Functions}, Proceedings of Symposia in Pure Mathematics 33, Part 1, American Mathematical Society, Providence, RI, 1979, 179--183.

\bibitem{platonov-rapinchuk-1994} V.~Platonov and A.~Rapinchuk, \emph{Algebraic Groups and Number Theory}, Pure and Applied Mathematics 139, Academic Press, Boston, 1994.

\bibitem{green-1955} J.~A.~Green, The characters of the finite general linear groups, \emph{Transactions of the American Mathematical Society} 80 (1955), 402--447. \url{https://doi.org/10.1090/S0002-9947-1955-0072878-2}.

\bibitem{clozel-1990} L.~Clozel, Motifs et formes automorphes: applications du principe de fonctorialit\'e, in \emph{Automorphic Forms, Shimura Varieties, and $L$-Functions}, Vol.~I, Perspectives in Mathematics 10, Academic Press, Boston, 1990, 77--159.

\bibitem{borel-1976} A.~Borel, Admissible representations of a semi-simple group over a local field with vectors fixed under an Iwahori subgroup, \emph{Inventiones Mathematicae} 35 (1976), 233--259. \url{https://doi.org/10.1007/BF01390139}.

\bibitem{casselman-1980} W.~Casselman, The unramified principal series of $p$-adic groups I. The spherical function, \emph{Compositio Mathematica} 40 (1980), 387--406.

\bibitem{radford-1979} D.~E.~Radford, A natural ring basis for the shuffle algebra and an application to group schemes, \emph{Journal of Algebra} 58 (1979), 432--454.

\bibitem{wallach-1992} N.~R.~Wallach, \emph{Real Reductive Groups II}, Pure and Applied Mathematics 132-II, Academic Press, Boston, 1992.

\bibitem{gross-1980} B.~H.~Gross, \emph{Arithmetic on Elliptic Curves with Complex Multiplication}, Lecture Notes in Mathematics 776, Springer, Berlin, 1980. \url{https://link.springer.com/book/10.1007/BFb0096754}.

\bibitem{miller-schmid-2012} S.~D.~Miller and W.~Schmid, On the rapid decay of cuspidal automorphic forms, \emph{Advances in Mathematics} 231 (2012), 940--964. \url{https://arxiv.org/abs/1106.2149}.

\bibitem{thorner-2024} J.~Thorner and A.~Zaman, Refinements to the prime number theorem for arithmetic progressions, \emph{Mathematische Zeitschrift} 306 (2024), no.~3, Paper No.~54. \url{https://arxiv.org/abs/2108.10878}.

\bibitem{kudla-1994} S.~S.~Kudla, Splitting metaplectic covers of dual reductive pairs, \emph{Israel Journal of Mathematics} 87 (1994), 361--401.

\bibitem{folland-1999} G.~B.~Folland, \emph{Real Analysis: Modern Techniques and Their Applications}, 2nd ed., John Wiley \& Sons (1999), Theorem~2.47, pp.~74--76. % [Text](https://upload.xuanyan.ws/Real%20Analysis%20Modern%20Techniques%20and%20Their%20Applications.pdf#page=87)

\bibitem{ahlfors-1979} L.~V.~Ahlfors, \emph{Complex Analysis}, 3rd ed., McGraw--Hill, 1979.

\bibitem{davenport-2000} H.~Davenport, \emph{Multiplicative Number Theory}, 3rd ed., revised by H.~L.~Montgomery, Springer, 2000.

\bibitem{rosser-schoenfeld-1962} J.~Barkley Rosser and L.~Schoenfeld, Approximate formulas for some functions of prime numbers, \emph{Illinois Journal of Mathematics} 6 (1962), 64--94.

\bibitem{montgomery-vaughan-2007} H.~L.~Montgomery and R.~C.~Vaughan, \emph{Multiplicative Number Theory I: Classical Theory}, Cambridge Studies in Advanced Mathematics 97, Cambridge University Press, 2007, \S2.2, Theorem~2.7(a). % [Publisher record](https://assets.cambridge.org/97805218/49036/frontmatter/9780521849036_frontmatter.pdf)

\bibitem{serre-1977} J.-P.~Serre, \emph{Linear Representations of Finite Groups}, Graduate Texts in Mathematics 42, Springer, New York, 1977, \S2.5, Theorem~7, and \S3.1, Theorem~9, pp.~10--26. % [Publisher records](https://link.springer.com/chapter/10.1007/978-1-4684-9458-7_2), (https://link.springer.com/chapter/10.1007/978-1-4684-9458-7_3)

\bibitem{benli-2026} K.~Benli, S.~Goel, H.~Twiss, and A.~Zaman, Explicit Deuring--Heilbronn phenomenon for Dirichlet \(L\)-functions, \emph{Proc. Amer. Math. Soc.} 154 (2026), no.~2, 509--525; revised text, arXiv:2410.06082v3.

\bibitem{thorner-2019} J.~Thorner and A.~Zaman, A unified and improved Chebotarev density theorem, \emph{Algebra Number Theory} 13 (2019), 1039--1068.

\bibitem{page-1935} A.~Page, On the number of primes in an arithmetic progression, \emph{Proc. London Math. Soc.} (2) 39 (1935), 116--141.

\bibitem{montgomery-1971} H.~L.~Montgomery, \emph{Topics in Multiplicative Number Theory}, Lecture Notes in Mathematics 227, Springer, 1971.

\bibitem{neukirch-1999} J.~Neukirch, \emph{Algebraic Number Theory}, Grundlehren der mathematischen Wissenschaften 322, Springer, 1999.

\bibitem{mueller-2024} W.~M\"uller, The Weyl law for congruence subgroups and arbitrary \(K_\infty\)-types, \emph{Journal of Functional Analysis} 287 (2024), article 110444. \url{https://doi.org/10.1016/j.jfa.2024.110444}

\bibitem{haines-kottwitz-prasad-2010} T.~J.~Haines, R.~E.~Kottwitz and A.~Prasad, Iwahori--Hecke algebras, \emph{Journal of the Ramanujan Mathematical Society} 25 (2010), 113--145.

\bibitem{elkies-2019} N.~D.~Elkies, Statements of general modularity results for weighted theta functions, \emph{Rational Lattices and their Theta Functions}, Harvard course notes, 11 November 2019. \url{https://people.math.harvard.edu/~elkies/M272.19/nov11.pdf}.

\bibitem{weil-1964} A.~Weil, Sur certains groupes d'op\'erateurs unitaires, \emph{Acta Mathematica} 111 (1964), 143--211. \url{https://webhomes.maths.ed.ac.uk/~v1ranick/papers/weil2.pdf}.

\bibitem{weil-1965} A.~Weil, Sur la formule de Siegel dans la th\'eorie des groupes classiques, \emph{Acta Mathematica} 113 (1965), 1--87. \url{https://doi.org/10.1007/BF02391774}.

\bibitem{waldspurger-1990} J.-L.~Waldspurger, D\'emonstration d'une conjecture de dualit\'e de Howe dans le cas \(p\)-adique, \(p\ne2\), in \emph{Festschrift in Honor of I.~I.~Piatetski-Shapiro}, Part I, Israel Mathematical Conference Proceedings 2 (1990), 267--324.

\bibitem{gan-takeda-2016} W.~T.~Gan and S.~Takeda, A proof of the Howe duality conjecture, \emph{Journal of the American Mathematical Society} 29 (2016), 473--493. \url{https://arxiv.org/abs/1407.1995}.

\bibitem{kudla-1996} S.~S.~Kudla, \emph{Notes on the Local Theta Correspondence}, European School in Group Theory lecture notes, Schlo{\ss} Hirschberg, 1996.

\bibitem{prasad-1996} D.~Prasad, \emph{A Brief Survey on the Theta Correspondence}, MPIM Preprint 1996 (138). \url{https://archive.mpim-bonn.mpg.de/id/eprint/3012/}.

\bibitem{rallis-1984} S.~Rallis, On the Howe duality conjecture, \emph{Compositio Mathematica} 51 (1984), 333--399. \url{https://numdam.org/item/CM_1984__51_3_333_0/}.

\bibitem{rallis-1987} S.~Rallis, \emph{\(L\)-Functions and the Oscillator Representation}, Lecture Notes in Mathematics 1245, Springer, 1987. \url{https://doi.org/10.1007/BFb0077894}.

\bibitem{kudla-rallis-1994} S.~S.~Kudla and S.~Rallis, A regularized Siegel--Weil formula: the first term identity, \emph{Annals of Mathematics} 140 (1994), 1--80. \url{https://annals.math.princeton.edu/1994/140-1/p01}.

\bibitem{yamana-2013} S.~Yamana, The Siegel--Weil formula for unitary groups, \emph{Pacific Journal of Mathematics} 264 (2013), 235--256. \url{https://doi.org/10.2140/pjm.2013.264.235}.

\bibitem{gan-qiu-takeda-2014} W.~T.~Gan, Y.~Qiu, and S.~Takeda, The regularized Siegel--Weil formula (the second term identity) and the Rallis inner product formula, \emph{Inventiones Mathematicae} 198 (2014), 739--831. \url{https://arxiv.org/abs/1207.4709}.

\bibitem{hazewinkel-1979} M.~Hazewinkel, On formal groups. The functional equation lemma and some of its applications, \emph{Ast\'erisque} 63 (1979), 73--82.

\bibitem{krattenthaler-2015} C.~Krattenthaler and T.~W.~M\"uller, Truncated versions of Dwork's lemma for exponentials of power series and \(p\)-divisibility of arithmetic functions, \emph{Advances in Mathematics} 283 (2015), 489--529.

\bibitem{akin-1982} K.~Akin, D.~A.~Buchsbaum and J.~Weyman, Schur functors and Schur complexes, \emph{Advances in Mathematics} 44 (1982), 207--278.

\bibitem{koblitz-1984} N.~Koblitz, \emph{\(p\)-adic Numbers, \(p\)-adic Analysis, and Zeta-Functions}, second edition, Graduate Texts in Mathematics 58, Springer-Verlag, New York (1984).

\bibitem{bannai-kobayashi-2010} K.~Bannai and S.~Kobayashi, Algebraic theta functions and the \(p\)-adic interpolation of Eisenstein--Kronecker numbers, \emph{Duke Math. J.} 153 (2010), 229--295. \url{https://doi.org/10.1215/00127094-2010-024}.

\bibitem{cohen-1946} I.~S.~Cohen and A.~Seidenberg, Prime ideals and integral dependence, \emph{Bull. Amer. Math. Soc.} 52 (1946), 252--261.

\bibitem{deconcini-1980} C.~De Concini, D.~Eisenbud and C.~Procesi, Young diagrams and determinantal varieties, \emph{Inventiones Mathematicae} 56 (1980), 129--165. \url{https://eisenbud.github.io/papers/pdfs/1980-003.pdf}

\bibitem{bruhat-tits-1987} F.~Bruhat and J.~Tits, Sch\'emas en groupes et immeubles des groupes classiques sur un corps local. II. Groupes unitaires, \emph{Bull. Soc. Math. France} 115 (1987), 141--195. \url{https://doi.org/10.24033/bsmf.2073}.

\bibitem{hardy-2008} G.~H.~Hardy and E.~M.~Wright, \emph{An Introduction to the Theory of Numbers}, sixth edition, Oxford University Press, 2008.

\bibitem{jantzen-2003} J.~C.~Jantzen, \emph{Representations of Algebraic Groups}, second edition, Mathematical Surveys and Monographs 107, American Mathematical Society, 2003. \url{https://doi.org/10.1090/surv/107}.

\bibitem{mahler-1958} K.~Mahler, An interpolation series for continuous functions of a $p$-adic variable, \emph{J. Reine Angew. Math.} 199 (1958), 23--34. \url{https://doi.org/10.1515/crll.1958.199.23}.

\bibitem{silverman-2009} J.~H.~Silverman, \emph{The Arithmetic of Elliptic Curves}, second edition, Graduate Texts in Mathematics 106, Springer, 2009. \url{https://doi.org/10.1007/978-0-387-09494-6}.

\bibitem{stromberg-2013} F.~Str\"omberg, Weil representations associated with finite quadratic modules, \emph{Math. Z.} 275 (2013), 509--527. \url{https://doi.org/10.1007/s00209-013-1145-x}.

\bibitem{schrijver-1986} A.~Schrijver, \emph{Theory of Linear and Integer Programming}, Wiley, Chichester, 1986, Corollary~7.1g.

\bibitem{borel-2007} A.~Borel, Automorphic forms on reductive groups, in \emph{Automorphic Forms and Applications}, IAS/Park City Mathematics Series 12, American Mathematical Society, 2007, 5--40. \url{https://www.math.ias.edu/~goresky/pdf/BorelFinal.pdf}.

\bibitem{wang-2015} S.~Wang, Grunwald--Wang theorem, an effective version, \emph{Science China Mathematics} 58 (2015), 1589--1606. \url{https://doi.org/10.1007/s11425-015-4977-5}.

\bibitem{voight-2021} J.~Voight, \emph{Quaternion Algebras}, Graduate Texts in Mathematics 288, Springer, Cham (2021). % [Publisher text](https://link.springer.com/book/10.1007/978-3-030-56694-4)

\bibitem{stacks-2026} The Stacks Project Authors, \emph{The Stacks Project} (2026).

\bibitem{fonseca-2019} T.~J.~Fonseca, Algebraic independence for values of integral curves, \emph{Algebra \& Number Theory} 13 (2019), 643--694. \url{https://msp.org/ant/2019/13-3/ant-v13-n3-p03-s.pdf}.

\bibitem{stacks-2026b} The Stacks Project Authors, \emph{The Stacks Project}, More on Algebra, Remark 15.113.6, Tag 09E8 (2026). \url{https://stacks.math.columbia.edu/tag/09E8}.

\bibitem{bost-chen-2013} J.-B.~Bost and H.~Chen, Concerning the semistability of tensor products in Arakelov geometry, \emph{Journal de Mathématiques Pures et Appliquées} 99 (2013), 436--488. \url{https://webusers.imj-prg.fr/~huayi.chen/Recherche/bost_chen.pdf}.

\bibitem{selberg-1967} A.~Selberg and S.~Chowla, On Epstein's zeta-function, \emph{J. Reine Angew. Math.} 227 (1967), 86--110. \url{https://doi.org/10.1515/crll.1967.227.86}.

\bibitem{banaszczyk-1993} W.~Banaszczyk, New bounds in some transference theorems in the geometry of numbers, \emph{Math. Ann.} 296 (1993), 625--635. \url{https://doi.org/10.1007/BF01445125}.

\bibitem{goodman-wallach-2009} R.~Goodman and N.~R.~Wallach, \emph{Symmetry, Representations, and Invariants}, Graduate Texts in Mathematics 255, Springer, 2009. \url{https://sites.math.rutgers.edu/~goodman/repbook.html}.

\bibitem{howe-umeda-1991} R.~Howe and T.~Umeda, The Capelli identity, the double commutant theorem, and multiplicity-free actions, \emph{Math. Ann.} 290 (1991), 565--619. \url{https://doi.org/10.1007/BF01459261}.

\bibitem{macdonald-1971} I.~G.~Macdonald, \emph{Spherical Functions on a Group of \(p\)-adic Type}, Publications of the Ramanujan Institute 2, University of Madras, 1971. \url{https://math.soimeme.org/~arunram/Resources/SFOAGOPTContents.html}.

\bibitem{harish-chandra-1966} Harish-Chandra, Discrete series for semisimple Lie groups. II. Explicit determination of the characters, \emph{Acta Math.} 116 (1966), 1--111. \url{https://repository.ias.ac.in/30833/1/336.pdf}.

\bibitem{lapid-rallis-2005} E.~M.~Lapid and S.~Rallis, On the local factors of representations of classical groups, in \emph{Automorphic Representations, \(L\)-Functions and Applications: Progress and Prospects}, Ohio State University Mathematical Research Institute Publications 11, de Gruyter, 2005, 309--359. \url{https://doi.org/10.1515/9783110892703.309}.

\bibitem{ash-2007} R.~B.~Ash, \emph{Basic Abstract Algebra}, Dover, 2007. \url{https://store.doverpublications.com/products/9780486453569}.

\bibitem{atiyah-1969} M.~F.~Atiyah and I.~G.~Macdonald, \emph{Introduction to Commutative Algebra}, Addison-Wesley, 1969. \url{https://books.google.com/books?id=senlelzZPBsC}.

\bibitem{fulton-1991} W.~Fulton and J.~Harris, \emph{Representation Theory: A First Course}, Graduate Texts in Mathematics 129, Springer, 1991. \url{https://doi.org/10.1007/978-1-4612-0979-9}.

\bibitem{gaudron-2008} \'E.~Gaudron, Pentes des fibr\'es vectoriels ad\'eliques sur un corps global, \emph{Rendiconti del Seminario Matematico della Universit\`a di Padova} 119 (2008), 21--95. \url{https://arxiv.org/abs/math/0605408}.

\bibitem{lieb-loss-2001} E.~H.~Lieb and M.~Loss, \emph{Analysis}, second edition, Graduate Studies in Mathematics 14, American Mathematical Society, Providence, RI, 2001, Theorem~1.13. \url{https://www.math.utoronto.ca/almut/MAT1000/LL-1.pdf}.

\bibitem{davenport-1980} H.~Davenport, \emph{Multiplicative Number Theory}, second edition, revised by H.~L.~Montgomery, Graduate Texts in Mathematics 74, Springer, New York, 1980. \url{https://doi.org/10.1007/978-1-4757-5927-3}.

\bibitem{ireland-rosen-1990} K.~Ireland and M.~Rosen, \emph{A Classical Introduction to Modern Number Theory}, second edition, Graduate Texts in Mathematics 84, Springer, New York, 1990. \url{https://doi.org/10.1007/978-1-4757-2103-4}.

\bibitem{beuzart-plessis-2021} R.~Beuzart-Plessis, A new proof of Jacquet--Rallis's fundamental lemma, \emph{Duke Mathematical Journal} 170 (2021), 2805--2814. \url{https://arxiv.org/abs/1901.02653}

\bibitem{borel-labesse-schwermer-1996} A.~Borel, J.-P.~Labesse and J.~Schwermer, On the cuspidal cohomology of $S$-arithmetic subgroups of reductive groups over number fields, \emph{Compositio Mathematica} 102 (1996), 1--40. \url{https://www.numdam.org/item/CM_1996__102_1_1_0/}

\bibitem{ferrero-washington-1979} B.~Ferrero and L.~C.~Washington, The Iwasawa invariant $\mu_p$ vanishes for abelian number fields, \emph{Annals of Mathematics} (2) 109 (1979), 377--395. \url{https://doi.org/10.2307/1971116}

\bibitem{arthur-clozel-1989} J.~Arthur and L.~Clozel, \emph{Simple Algebras, Base Change, and the Advanced Theory of the Trace Formula}, Annals of Mathematics Studies 120, Princeton University Press, 1989. \url{https://www.claymath.org/library/cw/arthur/pdf/30.pdf}

\bibitem{labesse-2011} J.-P.~Labesse, Changement de base CM et s\'eries discr\`etes, in \emph{On the Stabilization of the Trace Formula}, International Press, 2011, 429--470. \url{https://webusers.imj-prg.fr/~michael.harris/CBCM.pdf}

\bibitem{moeglin-waldspurger-1989} C.~M{\oe}glin and J.-L.~Waldspurger, Le spectre r\'esiduel de $\mathrm{GL}(n)$, \emph{Annales scientifiques de l'\'Ecole Normale Sup\'erieure} (4) 22 (1989), 605--674. \url{https://doi.org/10.24033/asens.1595}

\bibitem{scharlau-1985} W.~Scharlau, \emph{Quadratic and Hermitian Forms}, Grundlehren der mathematischen Wissenschaften 270, Springer, 1985.

\bibitem{borel-1963} A.~Borel, Some finiteness properties of adele groups over number fields, \emph{Publications Math\'ematiques de l'IH\'ES} 16 (1963), 5--30. \url{https://numdam.org/articles/10.1007/BF02684289/}

\bibitem{harish-chandra-1968} Harish-Chandra, \emph{Automorphic Forms on Semisimple Lie Groups}, notes by G.~van Dijk, Lecture Notes in Mathematics 62, Springer, 1968. \url{https://link.springer.com/book/10.1007/BFb0098434}

\bibitem{murphy-1990} G.~J.~Murphy, \emph{\(C^*\)-Algebras and Operator Theory}, Academic Press, 1990, Theorem~2.1.10, p.~41.

\bibitem{saito-2009} T.~Saito, Hilbert modular forms and \(p\)-adic Hodge theory, \emph{Compositio Mathematica} 145 (2009), 1081--1113.

\bibitem{henniart-1993} G.~Henniart, Caract\'erisation de la correspondance de Langlands locale par les facteurs \(\epsilon\) de paires, \emph{Inventiones Mathematicae} 113 (1993), 339--350.

\bibitem{chevalley-1951} C.~Chevalley, Deux th\'eor\`emes d'arithm\'etique, \emph{Journal of the Mathematical Society of Japan} 3 (1951), 36--44.

\bibitem{berger-2002} L.~Berger, Repr\'esentations \(p\)-adiques et \'equations diff\'erentielles, \emph{Inventiones Mathematicae} 148 (2002), 219--284.

\bibitem{grothendieck-1974} A.~Grothendieck, \emph{Groupes de Barsotti--Tate et cristaux de Dieudonn\'e}, S\'eminaire de math\'ematiques sup\'erieures 45, Presses de l'Universit\'e de Montr\'eal, 1974.

\bibitem{berthelot-breen-messing-1982} P.~Berthelot, L.~Breen and W.~Messing, \emph{Th\'eorie de Dieudonn\'e cristalline II}, Lecture Notes in Mathematics 930, Springer, 1982.

\bibitem{birkenhake-lange-2004} C.~Birkenhake and H.~Lange, \emph{Complex Abelian Varieties}, second edition, Grundlehren der mathematischen Wissenschaften 302, Springer, 2004.

\bibitem{katz-oda-1968} N.~M.~Katz and T.~Oda, On the differentiation of de Rham cohomology classes with respect to parameters, \emph{Journal of Mathematics of Kyoto University} 8 (1968), 199--213.

\bibitem{serre-tate-1968} J.-P.~Serre and J.~Tate, Good reduction of abelian varieties, \emph{Annals of Mathematics} 88 (1968), 492--517.

\bibitem{sorensen-2020} C.~M.~Sorensen, A patching lemma, in \emph{Shimura Varieties}, London Mathematical Society Lecture Note Series 457, Cambridge University Press, 2020, 297--305.

\bibitem{borel-wallach-2000} A.~Borel and N.~Wallach, \emph{Continuous Cohomology, Discrete Subgroups, and Representations of Reductive Groups}, second edition, Mathematical Surveys and Monographs 67, American Mathematical Society, 2000. \url{https://doi.org/10.1090/surv/067}.

\bibitem{ji-1997} L.~Ji, The trace class conjecture and the Weyl upper bound on the discrete spectrum of arithmetic groups, \emph{Mathematical Research Letters} 4 (1997), 501--511. \url{https://intlpress.com/site/pub/files/_fulltext/journals/mrl/1997/0004/0004/MRL-1997-0004-0004-a007.pdf}.

\bibitem{arthur-2005} J.~Arthur, An introduction to the trace formula, in \emph{Harmonic Analysis, the Trace Formula, and Shimura Varieties}, Clay Mathematics Proceedings 4 (2005), 1--263.

\bibitem{milne-2006} J.~S. Milne, \emph{Arithmetic Duality Theorems}, second edition, BookSurge, 2006, Chapter~I, \S\S2,4. \url{https://www.jmilne.org/math/Books/ADTnot.pdf}.

\bibitem{washington-1997} L.~C. Washington, \emph{Introduction to Cyclotomic Fields}, second edition, Graduate Texts in Mathematics 83, Springer, New York, 1997, Chapters~4--6 and Appendix, \S3. \url{https://doi.org/10.1007/978-1-4612-1934-7}.

\bibitem{borel-1974} A.~Borel, Stable real cohomology of arithmetic groups, \emph{Annales scientifiques de l'\'Ecole Normale Sup\'erieure} (4) 7 (1974), 235--272. \url{https://numdam.org/articles/10.24033/asens.1269/}

\bibitem{kallen-1980} W.~van der Kallen, Homology stability for linear groups, \emph{Inventiones Mathematicae} 60 (1980), 269--295. \url{https://doi.org/10.1007/BF01390018}

\bibitem{kisin-2010} M.~Kisin, The structure of potentially semi-stable deformation rings, in \emph{Proceedings of the International Congress of Mathematicians, Hyderabad 2010}, vol.~II, 294--311. \url{https://people.math.harvard.edu/~kisin/dvifiles/licm.pdf}.

\bibitem{chenevier-harris-2013} G.~Chenevier and M.~Harris, Construction of automorphic Galois representations, II, \emph{Cambridge Journal of Mathematics} 1 (2013), 53--73. \url{https://doi.org/10.4310/CJM.2013.v1.n1.a2}.

\bibitem{deligne-serre-1974} P.~Deligne and J.-P.~Serre, Formes modulaires de poids \(1\), \emph{Annales scientifiques de l'\'Ecole Normale Sup\'erieure} (4) 7 (1974), 507--530. \url{https://doi.org/10.24033/asens.1277}.

\bibitem{jacquet-langlands-1970} H.~Jacquet and R.~P. Langlands, \emph{Automorphic Forms on \(\mathrm{GL}(2)\)}, Lecture Notes in Mathematics 114, Springer, 1970. \url{https://doi.org/10.1007/BFb0058988}.

\bibitem{wiles-1995} A.~Wiles, Modular elliptic curves and Fermat's Last Theorem, \emph{Annals of Mathematics} (2) 141 (1995), 443--551. \url{https://annals.math.princeton.edu/1995/141-3/p01}.

\bibitem{ax-1970} J.~Ax, Zeros of polynomials over local fields---the Galois action, \emph{J. Algebra} 15 (1970), 417--428. \url{https://doi.org/10.1016/0021-8693(70)90069-4}.

\bibitem{bloch-kato-1990} S.~Bloch and K.~Kato, \(L\)-functions and Tamagawa numbers of motives, in \emph{The Grothendieck Festschrift}, vol.~I, Progress in Mathematics 86, Birkh\"auser, 1990, 333--400. \url{https://math.stanford.edu/~conrad/BSDseminar/refs/BKTamagawa.pdf}.

\bibitem{dirichlet-1837} P.~G.~L.~Dirichlet, Beweis des Satzes, dass jede unbegrenzte arithmetische Progression, deren erstes Glied und Differenz ganze Zahlen ohne gemeinschaftlichen Factor sind, unendlich viele Primzahlen enth\"alt, \emph{Abhandlungen der K\"oniglich Preussischen Akademie der Wissenschaften zu Berlin} (1837). \url{https://sites.math.rutgers.edu/~alexk/2022S572/Dirichlet1837.pdf}.

\bibitem{fontaine-1982} J.-M.~Fontaine, Formes diff\'erentielles et modules de Tate des vari\'et\'es ab\'eliennes sur les corps locaux, \emph{Invent. Math.} 65 (1982), 379--409. \url{https://www.imo.universite-paris-saclay.fr/~fontaine/varab.pdf}.

\bibitem{serre-2002} J.-P. Serre, \emph{Galois Cohomology}, Springer Monographs in Mathematics, Springer, Berlin, 2002.

\bibitem{grothendieck-1968} A.~Grothendieck, \emph{Cohomologie locale des faisceaux coh\'erents et th\'eor\`emes de Lefschetz locaux et globaux (SGA~2)}, Advanced Studies in Pure Mathematics 2, North-Holland, Amsterdam, 1968. \url{https://arxiv.org/abs/math/0511279}.

\bibitem{gelbart-jacquet-1978} S.~Gelbart and H.~Jacquet, A relation between automorphic representations of \(\mathrm{GL}(2)\) and \(\mathrm{GL}(3)\), \emph{Annales scientifiques de l'\'Ecole Normale Sup\'erieure} (4) 11 (1978), 471--542. \url{https://doi.org/10.24033/asens.1355}.

\bibitem{kato-2004} K.~Kato, \(p\)-adic Hodge theory and values of zeta functions of modular forms, \emph{Ast\'erisque} 295 (2004), 117--290.

\bibitem{chenevier-2014} G.~Chenevier, The \(p\)-adic analytic space of pseudocharacters of a profinite group and pseudorepresentations over arbitrary rings, in \emph{Automorphic Forms and Galois Representations}, vol.~1, London Mathematical Society Lecture Note Series 414, Cambridge University Press, 2014, 221--285. \url{https://doi.org/10.1017/CBO9781107446335.008}.

\bibitem{neukirch-2008} J.~Neukirch, A.~Schmidt and K.~Wingberg, \emph{Cohomology of Number Fields}, second edition, Grundlehren der mathematischen Wissenschaften 323, Springer, Berlin, 2008. \url{https://doi.org/10.1007/978-3-540-37889-1}.

\bibitem{gille-szamuely-2006} P.~Gille and T.~Szamuely, \emph{Central Simple Algebras and Galois Cohomology}, Cambridge Studies in Advanced Mathematics 101, Cambridge University Press, 2006, Theorem~6.2.8.

\bibitem{weibel-2013} C.~A.~Weibel, \emph{The \(K\)-book: An introduction to algebraic \(K\)-theory}, Graduate Studies in Mathematics 145, American Mathematical Society, 2013. \url{https://sites.math.rutgers.edu/~weibel/Kbook/Kbook.VI.pdf}.

\bibitem{thorne-2022} J.~A.~Thorne, On the vanishing of adjoint Bloch--Kato Selmer groups of irreducible automorphic Galois representations, \emph{Pure Appl. Math. Q.} 18 (2022), no.~5, 2159--2202. \url{https://doi.org/10.4310/pamq.2022.v18.n5.a5}.

\bibitem{stein-2007} W.~A. Stein, \emph{Modular Forms, a Computational Approach}, Graduate Studies in Mathematics 79, American Mathematical Society, Providence, RI, 2007, \S9.2, Theorem~9.4. \url{https://wstein.org/books/modform/modform/newforms.html}.

\bibitem{geraghty-2019} D.~Geraghty, Modularity lifting theorems for ordinary Galois representations, \emph{Math. Ann.} 373 (2019), 1341--1427. \url{https://doi.org/10.1007/s00208-018-1742-4}.

\bibitem{jacquet-liu-2018} H.~Jacquet and B.~Liu, On the local converse theorem for \(p\)-adic \(\mathrm{GL}_n\), \emph{American Journal of Mathematics} 140 (2018), 1399--1422.

\bibitem{jiang-nien-stevens-2015} D.~Jiang, C.~Nien and S.~Stevens, Towards the Jacquet conjecture on the local converse problem for \(p\)-adic \(\mathrm{GL}_n\), \emph{Journal of the European Mathematical Society} 17 (2015), 991--1007.

\bibitem{delorme-2023} P.~Delorme, On the spectral theorem of Langlands, \emph{Advances in Mathematics} 426 (2023), 109107.

\bibitem{goren-2001} E.~Z.~Goren, Hasse invariants for Hilbert modular varieties, \emph{Israel Journal of Mathematics} 122 (2001), 157--174.

\bibitem{deo-dimitrov-wiese-2024} S.~V.~Deo, M.~Dimitrov and G.~Wiese, Unramifiedness of weight \(1\) Hilbert Hecke algebras, \emph{Algebra \& Number Theory} 18 (2024), 1465--1496.

\bibitem{colmez-fontaine-2000} P.~Colmez and J.-M.~Fontaine, Construction des repr\'esentations \(p\)-adiques semi-stables, \emph{Invent. Math.} 140 (2000), 1--43.

\bibitem{lidl-niederreiter-1997} R.~Lidl and H.~Niederreiter, \emph{Finite Fields}, second edition, Encyclopedia of Mathematics and its Applications 20, Cambridge University Press, 1997.

\bibitem{zelevinsky-1981} A.~V.~Zelevinsky, \emph{Representations of Finite Classical Groups: A Hopf Algebra Approach}, Lecture Notes in Mathematics 869, Springer, 1981.

\bibitem{deriziotis-gotsis-1998} D.~I.~Deriziotis and C.~P.~Gotsis, The cuspidal modules of the finite general linear groups, \emph{LMS Journal of Computation and Mathematics} 1 (1998), 75--108.

\bibitem{brauer-1947} R.~Brauer, On Artin's $L$-series with general group characters, \emph{Annals of Mathematics} 48 (1947), 502--514. \href{https://doi.org/10.2307/1969183}{doi:10.2307/1969183}.

\bibitem{serre-2007} J.-P.~Serre, Bounds for the orders of the finite subgroups of $G(k)$, in \emph{Group Representation Theory}, EPFL Press, 2007, 405--450. \href{https://arxiv.org/abs/1011.0346}{arXiv:1011.0346}.

\bibitem{jiang-sun-zhu-2010} D.~Jiang, B.~Sun and C.-B.~Zhu, Uniqueness of Bessel models: the Archimedean case, \emph{Geometric and Functional Analysis} 20 (2010), 690--709. \url{https://doi.org/10.1007/s00039-010-0077-4}.

\bibitem{duflo-moore-1976} M.~Duflo and C.~C.~Moore, On the regular representation of a nonunimodular locally compact group, \emph{Journal of Functional Analysis} 21 (1976), 209--243. \url{https://www.sciencedirect.com/science/article/pii/0022123676900793}.

\bibitem{rosser-1962} J.~Barkley Rosser and L.~Schoenfeld, Approximate formulas for some functions of prime numbers, \emph{Illinois Journal of Mathematics} 6 (1962), 64--94.

\bibitem{stein-shakarchi-2003} E.~M.~Stein and R.~Shakarchi, \emph{Fourier Analysis: An Introduction}, Princeton Lectures in Analysis I, Princeton University Press, 2003. \url{https://press.princeton.edu/books/hardcover/9780691113845/fourier-analysis}.

\bibitem{borel-1969} A.~Borel, \emph{Introduction aux groupes arithm\'etiques}, Hermann, Paris, 1969, Th\'eor\`emes 8.7 et 13.1.

\bibitem{gross-2021} B.~H.~Gross, Incoherent definite spaces and Shimura varieties, in \emph{Relative Trace Formulas}, Simons Symposia, Springer, Cham (2021), 187--215. \url{https://doi.org/10.1007/978-3-030-68506-5_5}.

\bibitem{hormander-2003} L.~Hörmander, \emph{The Analysis of Linear Partial Differential Operators I: Distribution Theory and Fourier Analysis}, second edition, Classics in Mathematics, Springer, Berlin, 2003, Theorem~7.7.5. \url{https://doi.org/10.1007/978-3-642-61497-2}.

\bibitem{atobe-2020} H.~Atobe, On the non-vanishing of theta liftings of tempered representations of \(\mathrm U(p,q)\), \emph{Advances in Mathematics} 363 (2020), 106984, Definition~1.6 and Theorem~1.7.

\bibitem{langlands-1976} R.~P.~Langlands, \emph{On the Functional Equations Satisfied by Eisenstein Series}, Lecture Notes in Mathematics 544, Springer-Verlag, Berlin--Heidelberg, 1976. \url{https://publications.ias.edu/node/39}.

\bibitem{eischen-harris-li-skinner-2020} E.~Eischen, M.~Harris, J.~Li and C.~Skinner, $p$-adic $L$-functions for unitary groups, \emph{Forum of Mathematics, Pi} 8 (2020), e9. \url{https://doi.org/10.1017/fmp.2020.4}.

\bibitem{platonov-1969} V.~P.~Platonov, The problem of strong approximation and the Kneser--Tits conjecture for algebraic groups, \emph{Mathematics of the USSR-Izvestiya} 3 (1969), 1139--1147. \url{https://doi.org/10.1070/IM1969v003n06ABEH000838}.

\bibitem{he-2017} H.~He, On the Gan--Gross--Prasad conjecture for $U(p,q)$, \emph{Inventiones Mathematicae} 209 (2017), 837--884. \url{https://www.math.lsu.edu/~hongyu/pub/ggp_hhe_final.pdf}.

\bibitem{selberg-chowla-1967} A.~Selberg and S.~Chowla, On Epstein's zeta-function, \emph{Journal f\"ur die reine und angewandte Mathematik} 227 (1967), 86--110. \url{https://doi.org/10.1515/crll.1967.227.86}.

\bibitem{schmid-vilonen-1998} W.~Schmid and K.~Vilonen, Two geometric character formulas for reductive Lie groups, \emph{Journal of the American Mathematical Society} 11 (1998), 799--867. \url{https://arxiv.org/abs/math/9801081}.

\bibitem{molev-2003} A.~I.~Molev, Yangians and their applications, in \emph{Handbook of Algebra}, vol.~3, North-Holland, Amsterdam, 2003, 907--959. \url{https://arxiv.org/abs/math/0211288}.

\bibitem{liu-2022} Y.~Liu, Theta correspondence for almost unramified representations of unitary groups, \emph{Journal of Number Theory} 230 (2022), 196--224. \url{https://arxiv.org/abs/2105.14223}.

\bibitem{paul-1999} A.~Paul, First occurrence for the dual pairs \((U(p,q),U(r,s))\), \emph{Canad. J. Math.} 51 (1999), 636--657.

\bibitem{mazur-messing-1974} B.~Mazur and W.~Messing, \emph{Universal Extensions and One Dimensional Crystalline Cohomology}, Lecture Notes in Mathematics 370, Springer, 1974.

\bibitem{fargues-2011} L.~Fargues, La filtration canonique des points de torsion des groupes \(p\)-divisibles, \emph{Ann. Sci. \'Ec. Norm. Sup\'er. (4)} 44 (2011), 905--961.

\bibitem{bost-1996} J.-B.~Bost, Intrinsic heights of stable varieties and abelian varieties, \emph{Duke Math. J.} 82 (1996), 21--70.

\bibitem{gabriel-1970} P.~Gabriel, Construction de pr\'esch\'emas quotients, in \emph{Sch\'emas en Groupes I (SGA~3)}, Exp.~V, Lecture Notes in Mathematics 151, Springer, 1970, 250--286.

\bibitem{grothendieck-dieudonne-1961} A.~Grothendieck, avec la collaboration de J.~Dieudonn\'e, \'El\'ements de g\'eom\'etrie alg\'ebrique II, \'Etude globale \'el\'ementaire de quelques classes de morphismes, \emph{Publ. Math. IH\'ES} 8 (1961), 5--222.

\bibitem{hironaka-1964} H.~Hironaka, Resolution of singularities of an algebraic variety over a field of characteristic zero, I and II, \emph{Ann. of Math.} 79 (1964), 109--203 and 205--326. \href{https://doi.org/10.2307/1970486}{Part I}, \href{https://doi.org/10.2307/1970547}{Part II}.

\bibitem{scholze-weinstein-2013} P.~Scholze and J.~Weinstein, Moduli of \(p\)-divisible groups, \emph{Cambridge J. Math.} 1 (2013), 145--237. \href{https://doi.org/10.4310/CJM.2013.v1.n2.a1}{Article}.

\bibitem{lan-2013} K.-W.~Lan, \emph{Arithmetic Compactifications of PEL-Type Shimura Varieties}, London Mathematical Society Monographs 36, Princeton University Press, 2013. \href{https://www.jstor.org/stable/j.ctt24hpwv}{Book}.

\bibitem{fujino-2010} O.~Fujino, Effective base point free theorem for log canonical pairs, II. Angehrn--Siu type theorems, \emph{Michigan Math. J.} 59 (2010), 303--312. \href{https://www.math.kyoto-u.ac.jp/~fujino/effective2-final.pdf}{Full text}.

\bibitem{bost-gillet-soule-1994} J.-B.~Bost, H.~Gillet and C.~Soul\'e, Heights of projective varieties and positive Green forms, \emph{J. Amer. Math. Soc.} 7 (1994), 903--1027. \href{https://doi.org/10.1090/S0894-0347-1994-1260106-X}{Article}.

\bibitem{moriwaki-2009} A.~Moriwaki, Continuity of volumes on arithmetic varieties, \emph{Journal of Algebraic Geometry} 18 (2009), 407--457. \href{https://arxiv.org/html/math/0612269v2}{Full text}.

\bibitem{zhang-1995} S.~Zhang, Positive line bundles on arithmetic varieties, \emph{J. Amer. Math. Soc.} 8 (1995), 187--221.

\bibitem{kawamata-1982} Y.~Kawamata, A generalization of Kodaira--Ramanujam's vanishing theorem, \emph{Math. Ann.} 261 (1982), 43--46. \href{https://doi.org/10.1007/BF01456407}{Article}.

\bibitem{serre-1956} J.-P.~Serre, G\'eom\'etrie alg\'ebrique et analytique, \emph{Ann. Inst. Fourier} 6 (1956), 1--42. \href{https://www.numdam.org/item/AIF_1956__6__1_0/}{Article}.

\bibitem{demailly-2016} J.-P.~Demailly, Extension of holomorphic functions defined on non reduced analytic subvarieties, in \emph{The Legacy of Bernhard Riemann After One Hundred and Fifty Years}, Vol.~I, Adv. Lect. Math. 35, International Press, 2016, 191--222. \href{https://arxiv.org/abs/1510.05230}{Full text}.

\bibitem{faltings-chai-1990} G.~Faltings and C.-L.~Chai, \emph{Degeneration of Abelian Varieties}, Ergebnisse der Mathematik und ihrer Grenzgebiete 22, Springer, 1990. \href{https://doi.org/10.1007/978-3-662-02632-8}{Book}.

\bibitem{moriwaki-2009b} A.~Moriwaki, Continuous extension of arithmetic volumes, \emph{Int. Math. Res. Not. IMRN} 2009 (2009), 3598--3638. \url{https://doi.org/10.1093/imrn/rnp068}.

\bibitem{anschutz-lebras-2023} J.~Ansch\"utz and A.-C.~Le Bras, Prismatic Dieudonn\'e theory, \emph{Forum Math. Pi} 11 (2023), e2. \url{https://doi.org/10.1017/fmp.2022.22}.

\bibitem{scholze-weinstein-2020} P.~Scholze and J.~Weinstein, \emph{Berkeley Lectures on \(p\)-adic Geometry}, Annals of Mathematics Studies 207, Princeton University Press, 2020. \url{https://www.math.uni-bonn.de/people/scholze/Berkeley.pdf}.

\bibitem{grothendieck-1972} A.~Grothendieck, Mod\`eles de N\'eron et monodromie, in \emph{Groupes de monodromie en g\'eom\'etrie alg\'ebrique}, I, Exp.~IX, Lecture Notes in Mathematics 288, Springer, 1972, 313--523. \url{https://grothendiecksga.com/read/sga7/fr/9-3.html}.

\bibitem{grothendieck-1961} A.~Grothendieck, avec la collaboration de J.~Dieudonn\'e, \'El\'ements de g\'eom\'etrie alg\'ebrique III, \'Etude cohomologique des faisceaux coh\'erents, premi\`ere partie, \emph{Publ. Math. IH\'ES} 11 (1961), 5--167. \url{https://www.numdam.org/item/PMIHES_1961__11__5_0/}.

\bibitem{bost-2006} J.-B.~Bost, Evaluation maps, slopes, and algebraicity criteria, in \emph{Proceedings of the International Congress of Mathematicians, Madrid, August 22--30, 2006}, vol.~II, European Mathematical Society, Z\"urich, 2007, 537--562.

\bibitem{macdonald-1995} I.~G.~Macdonald, \emph{Symmetric Functions and Hall Polynomials}, second edition, Oxford Mathematical Monographs, Clarendon Press, Oxford, 1995.

\bibitem{blocki-2014} Z.~B{\l}ocki, Cauchy--Riemann meet Monge--Amp\`ere, \emph{Bull. Math. Sci.} 4 (2014), 433--480.

\bibitem{lan-2017} K.-W.~Lan, Integral models of toroidal compactifications with projective cone decompositions, \emph{Int. Math. Res. Not. IMRN} 2017 (2017), 3237--3280.

\bibitem{deligne-1970} P.~Deligne, \emph{\'Equations diff\'erentielles \`a points singuliers r\'eguliers}, Lecture Notes in Mathematics 163, Springer, 1970. \url{https://doi.org/10.1007/BFb0061194}.

\bibitem{fukuda-yoshikawa-2017} K.~Fukuda and S.~Yoshikawa, The 12th roots of the discriminant of an elliptic curve and the torsion points, \emph{International Journal of Number Theory} 13 (2017), 1037--1060. \url{https://doi.org/10.1142/S1793042117500555}.

\bibitem{lan-2016} K.-W.~Lan, Compactifications of PEL-type Shimura varieties in ramified characteristics, \emph{Forum Math. Sigma} 4 (2016), e1. \url{https://doi.org/10.1017/fms.2015.31}.

\bibitem{vandergeer-looijenga-2021} G.~van der Geer and E.~Looijenga, Lifting Chern classes by means of Ekedahl--Oort strata, \emph{Tunisian J. Math.} 3 (2021), 469--480. \url{https://doi.org/10.2140/tunis.2021.3.469}.

\bibitem{mumford-1970} D.~Mumford, \emph{Abelian Varieties}, Tata Institute of Fundamental Research Studies in Mathematics 5, Oxford University Press, 1970.

\bibitem{bernstein-2014} J.~Bernstein and B.~Kr\"otz, Smooth Fr\'echet globalizations of Harish-Chandra modules, \emph{Israel J. Math.} 199 (2014), 45--111. \url{https://www.math.tau.ac.il/~bernstei/Publication_list/publication_texts/Bern-Kroetz-2014.pdf}.

\bibitem{miller-2012} S.~D.~Miller and W.~Schmid, On the rapid decay of cuspidal automorphic forms, \emph{Adv. Math.} 231 (2012), 940--964. \url{https://arxiv.org/abs/1106.2149}.

\bibitem{fontaine-1994} J.-M.~Fontaine, Le corps des p\'eriodes $p$-adiques, \emph{Ast\'erisque} 223 (1994), 59--111. \url{https://numdam.org/item/AST_1994__223__59_0/}.

\bibitem{bhatt-scholze-2023} B.~Bhatt and P.~Scholze, Prismatic $F$-crystals and crystalline Galois representations, \emph{Cambridge J. Math.} 11 (2023), 507--562. \url{https://people.mpim-bonn.mpg.de/scholze/PrismaticCrystals.pdf}.

\bibitem{brauer-1941} R.~Brauer and C.~Nesbitt, On the modular characters of groups, \emph{Ann. of Math.} (2) 42 (1941), 556--590. \url{https://doi.org/10.2307/1968918}.

\bibitem{jacobson-1956} N.~Jacobson, \emph{Structure of Rings}, American Mathematical Society Colloquium Publications 37, American Mathematical Society, 1956. \url{https://doi.org/10.1090/coll/037}.

\bibitem{herr-1998} L.~Herr, Sur la cohomologie galoisienne des corps \(p\)-adiques, \emph{Bull. Soc. Math. France} 126 (1998), 563--600. \url{https://doi.org/10.24033/bsmf.2337}.

\bibitem{gee-liu-2022} T.~Gee and T.~Liu, Breuil--Kisin--Fargues modules and potentially semistable representations, Appendix~F in M.~Emerton and T.~Gee, \emph{Moduli Stacks of \'Etale $(\varphi,\Gamma)$-Modules and the Existence of Crystalline Lifts}, Annals of Mathematics Studies 215, Princeton University Press, 279--288. \url{https://doi.org/10.23943/princeton/9780691241357.005.0006}.

\bibitem{wu-2021} Z.~Wu, Galois representations, $(\varphi,\Gamma)$-modules and prismatic $F$-crystals, \emph{Doc. Math.} 26 (2021), 1771--1798. \url{https://doi.org/10.4171/DM/855}.

\bibitem{gan-hanke-yu-2001} W.~T.~Gan, J.~P.~Hanke and J.-K.~Yu, On an exact mass formula of Shimura, \emph{Duke Mathematical Journal} 107 (2001), 101--131. \url{https://doi.org/10.1215/S0012-7094-01-10716-3}.

\bibitem{jacobowitz-1962} R.~Jacobowitz, Hermitian forms over local fields, \emph{American Journal of Mathematics} 84 (1962), 441--465. \url{https://www.jstor.org/stable/2372982}.

\bibitem{cox-2013} D.~A.~Cox, \emph{Primes of the Form \(x^2+ny^2\)}, second edition, Wiley, Hoboken, 2013.

\bibitem{reinhardt-walker-2010} W.~P.~Reinhardt and P.~L.~Walker, Weierstrass elliptic and modular functions, in \emph{NIST Handbook of Mathematical Functions}, Cambridge University Press, 2010, Chapter~23. \url{https://dlmf.nist.gov/23}.

\bibitem{fontaine-2004} J.-M.~Fontaine, Arithm\'etique des repr\'esentations galoisiennes $p$-adiques, \emph{Ast\'erisque} 295 (2004), 1--115. \url{https://numdam.org/item/AST_2004__295__1_0/}

\bibitem{serre-1998} J.-P.~Serre, \emph{Abelian $\ell$-Adic Representations and Elliptic Curves}, Research Notes in Mathematics 7, A K Peters, 1998.

\bibitem{shalika-1974} J.~A.~Shalika, The multiplicity one theorem for $\GL_n$, \emph{Annals of Mathematics} 100 (1974), 171--193. \url{https://annals.math.princeton.edu/1974/100-2/p01}

\bibitem{sen-1980} S.~Sen, Continuous cohomology and \(p\)-adic Galois representations, \emph{Inventiones Mathematicae} 62 (1980), 89--116.

\bibitem{kisin-2006} M.~Kisin, Crystalline representations and \(F\)-crystals, in \emph{Algebraic Geometry and Number Theory}, Progress in Mathematics 253, Birkh\"auser, 2006, 459--496.

\bibitem{shin-2011} S.~W.~Shin, Galois representations arising from some compact Shimura varieties, \emph{Annals of Mathematics} 173 (2011), 1645--1741.

\bibitem{caraiani-2012} A.~Caraiani, Local-global compatibility and the action of monodromy on nearby cycles, \emph{Duke Mathematical Journal} 161 (2012), 2311--2413.

\bibitem{caraiani-2014} A.~Caraiani, Monodromy and local-global compatibility for \(l=p\), \emph{Algebra \& Number Theory} 8 (2014), 1597--1646.

\bibitem{harris-taylor-2001} M.~Harris and R.~Taylor, \emph{The Geometry and Cohomology of Some Simple Shimura Varieties}, Annals of Mathematics Studies 151, Princeton University Press, 2001.

\bibitem{atiyah-bott-1968} M.~F.~Atiyah and R.~Bott, A Lefschetz fixed point formula for elliptic complexes, II: Applications, \emph{Annals of Mathematics} 88 (1968), 451--491.

\bibitem{milne-1980} J.~S.~Milne, \emph{Étale Cohomology}, Princeton Mathematical Series 33, Princeton University Press, 1980.

\bibitem{artin-grothendieck-verdier-1973} M.~Artin, A.~Grothendieck, and J.-L.~Verdier, \emph{Théorie des topos et cohomologie étale des schémas, Tome 3}, Lecture Notes in Mathematics 305, Springer, 1973.

\bibitem{bhatt-morrow-scholze-2018} B.~Bhatt, M.~Morrow, and P.~Scholze, Integral \(p\)-adic Hodge theory, \emph{Publications Mathématiques de l'IHÉS} 128 (2018), 219--397.

\bibitem{rouquier-1996} R.~Rouquier, Caractérisation des caractères et pseudo-caractères, \emph{Journal of Algebra} 180 (1996), 571--586.

\bibitem{taylor-1991} R.~Taylor, Galois representations associated to Siegel modular forms of low weight, \emph{Duke Mathematical Journal} 63 (1991), 281--332.

\bibitem{tadic-1986} M.~Tadi\'c, Classification of unitary representations in irreducible representations of general linear group (non-Archimedean case), \emph{Annales Scientifiques de l'École Normale Supérieure} (4) 19 (1986), 335--382.

\bibitem{silberger-zink-2008} A.~J.~Silberger and E.-W.~Zink, Explicit Shintani base change and the Macdonald correspondence for characters of \(\GL_n(k)\), \emph{Journal of Algebra} 319 (2008), 4147--4176.

\bibitem{moy-prasad-1996} A.~Moy and G.~Prasad, Jacquet functors and unrefined minimal \(K\)-types, \emph{Commentarii Mathematici Helvetici} 71 (1996), 98--121.

\bibitem{sorensen-2010} C.~M.~Sorensen, Galois representations attached to Hilbert--Siegel modular forms, \emph{Documenta Mathematica} 15 (2010), 623--670.

\bibitem{artin-tate-2009} E.~Artin and J.~Tate, \emph{Class Field Theory}, AMS Chelsea Publishing 366, American Mathematical Society, 2009. \url{https://www.ams.org/books/chel/366/}.

\bibitem{drinfeld-1974} V.~G.~Drinfel'd, Elliptic modules, \emph{Math. USSR-Sb.} 23 (1974), 561--592. \url{https://doi.org/10.1070/SM1974v023n04ABEH001731}.

\bibitem{godement-jacquet-1972} R.~Godement and H.~Jacquet, \emph{Zeta Functions of Simple Algebras}, Lecture Notes in Mathematics 260, Springer, 1972. \url{https://doi.org/10.1007/BFb0070263}.

\bibitem{jacquet-piatetskii-shapiro-shalika-1983} H.~Jacquet, I.~I.~Piatetskii-Shapiro and J.~A.~Shalika, Rankin--Selberg convolutions, \emph{Amer. J. Math.} 105 (1983), 367--464. \url{https://www.math.columbia.edu/~hj/Rankin%20Selberg%20convolutions.pdf}.

\bibitem{katz-1981} N.~M.~Katz, Serre--Tate local moduli, in \emph{Surfaces algébriques}, Lecture Notes in Mathematics 868, Springer, 1981, 138--202. \url{https://doi.org/10.1007/BFb0090648}.

\bibitem{messing-1972} W.~Messing, \emph{The Crystals Associated to Barsotti--Tate Groups: With Applications to Abelian Schemes}, Lecture Notes in Mathematics 264, Springer, 1972. \url{https://doi.org/10.1007/BFb0058301}.

\bibitem{scholze-2013} P.~Scholze, The local Langlands correspondence for \(\GL_n\) over \(p\)-adic fields, \emph{Invent. Math.} 192 (2013), 663--715. \url{https://doi.org/10.1007/s00222-012-0420-5}.

\bibitem{rodier-1982} F.~Rodier, Repr\'esentations de \(\GL(n,k)\) o\`u \(k\) est un corps \(p\)-adique, \emph{Ast\'erisque} 92--93 (1982), 201--218. \url{https://www.numdam.org/item/SB_1981-1982__24__201_0/}.

\bibitem{gelbart-piatetski-shapiro-rallis-1987} S.~Gelbart, I.~Piatetski-Shapiro and S.~Rallis, \emph{Explicit Constructions of Automorphic $L$-Functions}, Lecture Notes in Mathematics 1254, Springer, Berlin, 1987.

\bibitem{widder-1941} D.~V.~Widder, \emph{The Laplace Transform}, Princeton Mathematical Series 6, Princeton University Press (1941), Chapter II, pp.~35--99.

\end{thebibliography}

\end{document}
